
%%%%%%%%%%%%%%%%%% PREAMBULE %%%%%%%%%%%%%%%%%%

\documentclass[11pt,a4paper]{article}

\usepackage{amsfonts,amsmath,amssymb,amsthm}
\usepackage[T1]{fontenc}
\usepackage[french]{babel}
\usepackage{mathptmx}
\usepackage{fancybox}
\usepackage{graphicx}
\usepackage{ifthen}

\usepackage{tikz}   

\usepackage{hyperref}
\hypersetup{colorlinks=true, linkcolor=blue, urlcolor=blue,
pdftitle={Exo7 - Exercices de mathématiques}, pdfauthor={Exo7}}

\usepackage{geometry}
\geometry{top=2cm, bottom=2cm, left=2cm, right=2cm}

%----- Ensembles : entiers, reels, complexes -----
\newcommand{\Nn}{\mathbb{N}} \newcommand{\N}{\mathbb{N}}
\newcommand{\Zz}{\mathbb{Z}} \newcommand{\Z}{\mathbb{Z}}
\newcommand{\Qq}{\mathbb{Q}} \newcommand{\Q}{\mathbb{Q}}
\newcommand{\Rr}{\mathbb{R}} \newcommand{\R}{\mathbb{R}}
\newcommand{\Cc}{\mathbb{C}} \newcommand{\C}{\mathbb{C}}
\newcommand{\Kk}{\mathbb{K}} \newcommand{\K}{\mathbb{K}}

%----- Modifications de symboles -----
\renewcommand{\epsilon}{\varepsilon}
\renewcommand{\Re}{\mathop{\mathrm{Re}}\nolimits}
\renewcommand{\Im}{\mathop{\mathrm{Im}}\nolimits}
\renewcommand{\ge}{\geqslant} \renewcommand{\geq}{\geqslant}
\renewcommand{\le}{\leqslant} \renewcommand{\leq}{\leqslant}
\newcommand{\dd}{\mathrm{d}} \newcommand{\Ker}{\mathop{\mathrm{Ker}}\nolimits}

%----- Fonctions usuelles -----
\newcommand{\ch}{\mathop{\mathrm{ch}}\nolimits}
\newcommand{\sh}{\mathop{\mathrm{sh}}\nolimits}
\renewcommand{\tanh}{\mathop{\mathrm{th}}\nolimits}
\newcommand{\cotan}{\mathop{\mathrm{cotan}}\nolimits}
\newcommand{\Arcsin}{\mathop{\mathrm{arcsin}}\nolimits}
\newcommand{\Arccos}{\mathop{\mathrm{arccos}}\nolimits}
\newcommand{\Arctan}{\mathop{\mathrm{arctan}}\nolimits}
\newcommand{\Argsh}{\mathop{\mathrm{argsh}}\nolimits}
\newcommand{\Argch}{\mathop{\mathrm{argch}}\nolimits}
\newcommand{\Argth}{\mathop{\mathrm{argth}}\nolimits}
\newcommand{\pgcd}{\mathop{\mathrm{pgcd}}\nolimits} 

%----- Structure des exercices ------

\newcommand{\exercice}[1]{\video{0}}
\newcommand{\finexercice}{}
\newcommand{\noindication}{}
\newcommand{\nocorrection}{}

\newcounter{exo}
\newcommand{\enonce}[2]{\refstepcounter{exo}\hypertarget{exo7:#1}{}\label{exo7:#1}{\bf Exercice \arabic{exo}}\ \  #2\vspace{1mm}\hrule\vspace{1mm}}

\newcommand{\finenonce}[1]{
\ifthenelse{\equal{\ref{ind7:#1}}{\ref{bidon}}\and\equal{\ref{cor7:#1}}{\ref{bidon}}}{}{\par{\footnotesize
\ifthenelse{\equal{\ref{ind7:#1}}{\ref{bidon}}}{}{\hyperlink{ind7:#1}{\texttt{Indication} $\blacktriangledown$}\qquad}
\ifthenelse{\equal{\ref{cor7:#1}}{\ref{bidon}}}{}{\hyperlink{cor7:#1}{\texttt{Correction} $\blacktriangledown$}}}}
\ifthenelse{\equal{\myvideo}{0}}{}{{\footnotesize\qquad\texttt{\href{https://www.youtube.com/watch?v=\myvideo}{Vidéo $\blacksquare$}}}}
\hfill{\scriptsize\texttt{[#1]}}\vspace{1mm}\hrule\vspace*{7mm}}

\newcommand{\indication}[1]{\hypertarget{ind7:#1}{}\label{ind7:#1}{\bf Indication pour \hyperlink{exo7:#1}{l'exercice \ref{exo7:#1} $\blacktriangle$}}\vspace{1mm}\hrule\vspace{1mm}}
\newcommand{\finindication}{\vspace{1mm}\hrule\vspace*{7mm}}
\newcommand{\correction}[1]{\hypertarget{cor7:#1}{}\label{cor7:#1}{\bf Correction de \hyperlink{exo7:#1}{l'exercice \ref{exo7:#1} $\blacktriangle$}}\vspace{1mm}\hrule\vspace{1mm}}
\newcommand{\fincorrection}{\vspace{1mm}\hrule\vspace*{7mm}}

\newcommand{\finenonces}{\newpage}
\newcommand{\finindications}{\newpage}


\newcommand{\fiche}[1]{} \newcommand{\finfiche}{}
%\newcommand{\titre}[1]{\centerline{\large \bf #1}}
\newcommand{\addcommand}[1]{}

% variable myvideo : 0 no video, otherwise youtube reference
\newcommand{\video}[1]{\def\myvideo{#1}}

%----- Presentation ------

\setlength{\parindent}{0cm}

\definecolor{myred}{rgb}{0.93,0.26,0}
\definecolor{myorange}{rgb}{0.97,0.58,0}
\definecolor{myyellow}{rgb}{1,0.86,0}

\newcommand{\LogoExoSept}[1]{  % input : echelle       %% NEW
{\usefont{U}{cmss}{bx}{n}
\begin{tikzpicture}[scale=0.1*#1,transform shape]
  \fill[color=myorange] (0,0)--(4,0)--(4,-4)--(0,-4)--cycle;
  \fill[color=myred] (0,0)--(0,3)--(-3,3)--(-3,0)--cycle;
  \fill[color=myyellow] (4,0)--(7,4)--(3,7)--(0,3)--cycle;
  \node[scale=5] at (3.5,3.5) {Exo7};
\end{tikzpicture}}
}


% titre
\newcommand{\titre}[1]{%
\vspace*{-4ex} \hfill \hspace*{1.5cm} \hypersetup{linkcolor=black, urlcolor=black} 
\href{https://exo7.emath.fr}{\LogoExoSept{3}} 
 \vspace*{-5.7ex}\newline 
\hypersetup{linkcolor=blue, urlcolor=blue}  {\Large \bf #1} \newline 
 \rule{12cm}{1mm} \vspace*{3ex}}

%----- Commandes supplementaires ------

\usepackage{graphics}
\usepackage[all]{xy}


\begin{document}

%%%%%%%%%%%%%%%%%% EXERCICES %%%%%%%%%%%%%%%%%%
\fiche{f00000, exo7, 2027/01/01}

\titre{Exo7 - Tous les exercices}


\vspace*{2cm}

\tableofcontents

\newpage

\pdfcompresslevel=9

\section{100.01 Logique}


\exercice{104, bodin, 1998/09/01}

\enonce{000104}{}
Soient $R$ et $S$ des relations. Donner la n\'egation de $R \Rightarrow S$.
\finenonce{000104}



\finexercice



\exercice{105, bodin, 1998/09/01}

\enonce{000105}{}

 D\'emontrer que $(1=2) \Rightarrow (2=3)$.

\finenonce{000105}



\finexercice



\exercice{106, bodin, 1998/09/01}
\video{nfJizOU7DbA}
\enonce{000106}{}

Soient les quatre assertions suivantes :
$$ (a) \ \exists x\in \Rr \quad \forall y\in \Rr \quad x+y > 0 \quad ; \quad
 (b) \ \forall x\in \Rr \quad \exists y\in \Rr \quad x+y > 0  \ ;$$
$$ (c) \ \forall x\in \Rr \quad \forall y\in \Rr \quad x+y > 0  \quad ; \quad
(d) \ \exists x\in \Rr \quad \forall y\in \Rr \quad y^2 > x .$$
\begin{enumerate}
    \item Les assertions $a$, $b$, $c$, $d$ sont-elles vraies ou fausses ?
    \item Donner leur n\'egation.
\end{enumerate}
\finenonce{000106} 


\finexercice


\exercice{107, bodin, 1998/09/01}
\video{ZCLCoGR3hcc}
\enonce{000107}{}
 Soit  $f$  une application de  ${\Rr}$  dans
${\Rr}$. Nier, de la mani\`ere la plus pr\' ecise possible, les \'
enonc\' es qui suivent :
\begin{enumerate}
    \item Pour tout  $x\in {\Rr}\ f(x)\leq 1$.
    \item L'application  $f$  est croissante.
    \item L'application  $f$  est croissante et positive.
    \item Il existe  $x\in {\Rr}^+$  tel que  $f(x)\leq 0$.
    \item  Il existe $x\in \Rr$ tel que quel que soit $y\in \Rr$, si $x<y$ alors $f(x)>f(y)$.
\end{enumerate}
\noindent On ne demande pas de d\' emontrer quoi que ce soit, juste d'\' ecrire le contraire
d'un \' enonc\' e.

\finenonce{000107} 


\finexercice


\exercice{108, bodin, 1998/09/01}
\video{2tmyTz5_2Rw}
\enonce{000108}{}

Compl\'eter les pointill\'es par le connecteur logique qui s'impose : $\Leftrightarrow ,\ \Leftarrow,\ \Rightarrow.$
\begin{enumerate}
    \item $  x\in \Rr\ \ x^2=4 \ \ldots\ldots \ x=2$ ;
    \item $  z\in \Cc\ \ z=\overline{z} \ \ldots\ldots \ z\in\Rr$ ;
    \item $ x\in \Rr\ \ x=\pi \ \ldots\ldots \ e^{2ix}=1$.
\end{enumerate}
\finenonce{000108} 


\finexercice


\exercice{109, bodin, 1998/09/01}
\video{q07OjCbDg3o}
\enonce{000109}{}

Dans $\R^2$, on d\'efinit les ensembles $F_{1}=\{(x,y)\in\R^2,\ y\leq0\}$ et
$F_{2}=\{(x,y)\in\R^2,\ xy\geq1,\ x\geq0\}$. On note $M_1M_2$ la distance usuelle entre deux points $M_1$ et $M_2$ de $\R^2$.
\'Evaluer les propositions suivantes :
\begin{enumerate}
\item $
 \forall\epsilon\in]0,+\infty[\quad
 \exists M_{1}\in F_{1}\quad
 \exists M_{2}\in F_{2}
 \qquad M_{1}M_{2}<\epsilon$
\item $
 \exists M_{1}\in F_{1}\quad
 \exists M_{2}\in F_{2}\quad
 \forall\epsilon\in]0,+\infty[
 \qquad M_{1}M_{2}<\epsilon$
\item $
 \exists\epsilon\in]0,+\infty[\quad 
 \forall M_{1}\in F_{1}\quad
 \forall M_{2}\in F_{2}
 \qquad M_{1}M_{2}<\epsilon$
\item $
 \forall M_{1}\in F_{1}\quad
 \forall M_{2}\in F_{2}\quad
 \exists\epsilon\in]0,+\infty[
 \qquad M_{1}M_{2}<\epsilon$
\end{enumerate}
Quand elles sont fausses, donner leur n\'egation.
\finenonce{000109} 


\finexercice


\exercice{110, gourio, 2001/09/01}
\video{QSbHxGnf-9M}
\enonce{000110}{}
Nier la proposition: ``tous les habitants de la rue du Havre qui ont les
yeux bleus gagneront au loto et prendront leur retraite avant 50 ans''.

\finenonce{000110} 


\finexercice


\exercice{111, bodin, 1998/09/01}

\enonce{000111}{}
\'Ecrire la n\'egation des assertions suivantes o\`u $P, Q, R, S$ sont des
propositions.
\begin{enumerate}
    \item $P \Rightarrow Q$,
    \item $P$ et non $Q$,
    \item $P$ et ($Q$ et $R$),
    \item $P$ ou ($Q$ et $R$),
    \item ($P$ et $Q$) $\Rightarrow \ (R \Rightarrow S)$.
\end{enumerate}

\finenonce{000111}



\finexercice



\exercice{112, bodin, 1998/09/01}
\video{FH3vTlYT23M}
\enonce{000112}{}
 Nier les assertions suivantes :
\begin{enumerate}
    \item tout triangle rectangle poss\`ede un angle droit ;
    \item dans toutes les \'ecuries, tous les chevaux sont noirs ;
    \item pour tout entier $x$, il existe un entier $y$ tel que, pour tout entier $z$, 
la relation $z<x$ implique la relation $z<x+1$ ;
    \item $\forall \epsilon >0 \quad \exists \alpha >0 \qquad (|x-7/5|<\alpha \Rightarrow |5x-7|<\epsilon)$.
\end{enumerate}


\finenonce{000112} 


\finexercice


\exercice{113, bodin, 1998/09/01}

\enonce{000113}{Le missionnaire et les cannibales}
Les cannibales d'une tribu se pr\'eparent \`a manger un missionnaire. D\'esirant lui prouver une derni\`ere fois leur respect de la dignit\'e et de la libert\'e humaine, les cannibales proposent au missionnaire de d\'ecider lui-m\^eme de son sort en faisant une courte d\'eclaration : si celle-ci est vraie, le missionnaire sera r\^oti, et il sera bouilli dans le cas contraire. Que doit dire le missionnaire pour sauver sa vie ? (d'apr\`es Cervant\`es)
\finenonce{000113}



\finexercice



\exercice{114, ridde, 1999/11/01}

\enonce{000114}{}
La proposition $ \big( P \wedge Q \Rightarrow (\neg P) \vee Q \big)$ est-elle vraie ?
\finenonce{000114}



\finexercice



\exercice{115, ridde, 1999/11/01}

\enonce{000115}{}
On
suppose que la proposition $P$ est vraie ainsi que les propositions suivantes :
\begin{enumerate}
\item $ (\neg Q) \wedge P \Rightarrow \neg S$.
\item $S \Rightarrow (\neg P) \vee Q$.
\item $P \Rightarrow R \vee S$.
\item $ S \wedge Q \Rightarrow \neg P$.
\item $R \wedge \neg (S \vee Q) \Rightarrow T$.
\item $R \Rightarrow (\neg P) \vee (\neg Q)$.
\end{enumerate}
La proposition $T$ est-elle vraie ?
\finenonce{000115}



\finexercice



\exercice{116, ridde, 1999/11/01}

\enonce{000116}{}

Ecrire la n\'egation des phrases suivantes :
\begin{enumerate}
\item $ (\forall x) (\exists n)/ (x \leq n)$.
\item $ (\exists M) / (\forall n) (\left|u_n\right| \leq M)$.
\item $ (\forall x) (\forall y) (xy = yx)$.
\item $ (\forall x) (\exists y)/ (yxy^{-1} = x)$.
\item $ (\forall \epsilon>0) (\exists N \in \Nn)/ (\forall n \geq N) (\left|u_n\right|
< \epsilon)$.
\item $(\forall x \in \Rr) (\forall \epsilon >0) (\exists \alpha >0)/ (\forall f \in \mathcal{F})
 (\forall y \in \Rr) (\left|x-y\right|<\alpha \Rightarrow \left|f (x)-f (y)\right|< \epsilon)$.
\end{enumerate}
\finenonce{000116}



\finexercice



\exercice{117, ridde, 1999/11/01}

\enonce{000117}{}
Comparer les diff\'erentes phrases (sont-elles \'equivalentes, contraires, quelles sont
celles qui impliquent les autres...)
\begin{enumerate}
\item $ (\forall x) (\exists y)/ (x \leq y)$.
\item $ (\forall x) (\forall y) (x \leq y)$.
\item $ (\exists x) (\exists y)/(x \leq y)$.
\item $ (\exists x)/(\forall y) (x \leq y)$.
\item $ (\exists x)/(\forall y) (y<x)$.
\item $ (\exists x) (\exists y)/ (y<x)$.
\item $ (\forall x) (\exists y)/(x  =  y)$.
\end{enumerate}
\finenonce{000117}



\finexercice



\exercice{118, ridde, 1999/11/01}

\enonce{000118}{}
Si $P (x)$ est une proposition d\'ependant de $x \in X$, on note $\overline{P} =
\left\{ x \in X / P (x) \text{ est vraie}\right\}$.
Exprimer en fonction de $\overline{P}$ et $\overline{Q}$ les ensembles $\overline{\neg P},
 \overline{P \wedge Q}, \overline{P \vee Q}, \overline{P \Rightarrow Q},
\overline{P \Leftrightarrow Q}$.
\finenonce{000118}



\finexercice



\exercice{119, ridde, 1999/11/01}
\video{UK22qba7LV0}
\enonce{000119}{}

Montrer que 
$$\forall \epsilon >0 \quad \exists N \in \Nn \text{ tel que }
 (n \geq N \Rightarrow 2-\epsilon < \frac{2n + 1}{n + 2} < 2 + \epsilon).$$
\finenonce{000119} 


\finexercice


\exercice{120, bodin, 1998/09/01}
\video{-Kid1jaNv5A}
\enonce{000120}{}

Soient $f,g$ deux fonctions de $\Rr$ dans $\Rr$. Traduire
en termes de quantificateurs les expressions suivantes :
\begin{enumerate}
  \item $f$ est major\'ee;
  \item $f$ est born\'ee;
  \item $f$ est paire;
  \item $f$ est impaire;
  \item $f$ ne s'annule jamais;
  \item $f$ est p\'eriodique;
  \item $f$ est croissante;
  \item $f$ est strictement d\'ecroissante;
  \item $f$ n'est pas la fonction nulle;
  \item $f$ n'a jamais les m\^{e}mes valeurs en deux points distincts;
  \item $f$ atteint toutes les valeurs de $\Nn$;
  \item $f$ est inf\'erieure \`a $g$;
  \item $f$ n'est pas inf\'erieure \`a $g$.
\end{enumerate}
\finenonce{000120} 


\finexercice


\exercice{5103, rouget, 2010/06/30}
\enonce{005103}{**IT}
Exprimer à l'aide de quantificateurs les phrases suivantes puis donner leur négation.

\begin{enumerate}
\item  ($f$ étant une application du plan dans lui-même)
\begin{enumerate}
\item $f$ est l'identité du plan.
\item $f$ a au moins un point invariant (on dit aussi point fixe).
\end{enumerate}
\item  ($f$ étant une application de $\Rr$ dans $\Rr$)
\begin{enumerate}
\item $f$ est l'application nulle.
\item L'équation $f(x)=0$ a une solution.
\item L'équation $f(x)=0$ a exactement une solution.
\end{enumerate}
\item  ($(u_n)_{n\in\Nn}$ étant une suite réelle)
\begin{enumerate}
\item La suite $(u_n)_{n\in\Nn}$ est bornée.
\item La suite $(u_n)_{n\in\Nn}$ est croissante.
\item La suite $(u_n)_{n\in\Nn}$ est monotone.
\end{enumerate}
\end{enumerate}
\finenonce{005103}


\finexercice


\exercice{5104, rouget, 2010/06/30}
\enonce{005104}{*IT}
Donner la négation des phrases suivantes

\begin{enumerate}
\item  $x\geq3$
\item  $0<x\leq2$.
\end{enumerate}

\finenonce{005104}


\finexercice


\exercice{5105, rouget, 2010/06/30}
\enonce{005105}{**IT}
Les phrases suivantes sont-elles équivalentes~?

\begin{enumerate}
\item  \og~$\forall x\in\Rr,\;(f(x)=0\;\mbox{et}\;g(x)=0)$~\fg~et \og~$(\forall
x\in\Rr,\;f(x)=0)\;\mbox{et}\;(\forall x\in\Rr,\;g(x)=0)$~\fg.
\item  \og~$\forall x\in\Rr,\;(f(x)=0\;\mbox{ou}\;g(x)=0)$~\fg~et \og~$(\forall
x\in\Rr,\;f(x)=0)\;\mbox{ou}\;(\forall x\in\Rr,\;g(x)=0)$~\fg.
\end{enumerate}
Donner un exemple de fonctions $f$ et $g$ de $\Rr$ dans $\Rr$, toutes deux non nulles et dont le produit est nul.
\finenonce{005105}


\finexercice

\section{100.02 Ensemble}


\exercice{121, bodin, 1998/09/01}

\enonce{000121}{}
Montrer que $\emptyset \subset X$, pour tout ensemble $X$.
\finenonce{000121}



\finexercice



\exercice{122, bodin, 1998/09/01}
\video{MxucdpW0WDI}
\enonce{000122}{}

 Montrer par contraposition les assertions suivantes, $E$ \'etant
un ensemble :
\begin{enumerate}
\item $\forall A,B \in \mathcal{P}(E) \quad (A\cap B=A\cup B)\Rightarrow A=B$,
\item $\forall A,B,C \in \mathcal{P}(E) \quad
(A\cap B=A\cap C \text{ et } A\cup B=A\cup C)\Rightarrow B=C$.
\end{enumerate}
\finenonce{000122} 


\finexercice


\exercice{123, bodin, 1998/09/01}
\video{I6E-O_PNk1Y}
\enonce{000123}{}
 Soit $A,B$ deux ensembles, montrer
$\complement(A\cup B) =\complement A\cap \complement B$ et
$\complement (A\cap B) =\complement A\cup \complement B$.

\finenonce{000123} 


\finexercice


\exercice{124, bodin, 1998/09/01}
\video{64sBWnSD4nM}
\enonce{000124}{}

Soient $E$ et $F$ deux ensembles, $f:E\rightarrow F$. D\'emontrer que :\\
$\forall A,B \in \mathcal{P}(E) \quad (A\subset B)\Rightarrow (f(A)\subset f(B))$,\\
$\forall A,B \in \mathcal{P}(E) \quad f(A\cap B)\subset f(A)\cap f(B)$,\\
$\forall A,B \in \mathcal{P}(E) \quad f(A\cup B) = f(A)\cup f(B)$,\\
$\forall A,B \in \mathcal{P}(F) \quad f^{-1}(A\cup B) = f^{-1}(A)\cup f^{-1}(B)$,\\
$\forall A \in \mathcal{P}(F) \quad f^{-1}(F\setminus A)=E\setminus f^{-1}(A)$.

\finenonce{000124} 


\finexercice


\exercice{125, cousquer, 2003/10/01}

\enonce{000125}{}
 $A$ et $B$ étant des parties d'un ensemble~$E$,
démontrer les lois de Morgan~:
$$\complement A \cup\complement B = \complement(A \cap B) \quad\mbox{et}
\quad\complement A \cap\complement B = \complement(A \cup B).$$
\finenonce{000125}



\finexercice



\exercice{126, cousquer, 2003/10/01}

\enonce{000126}{}
 Démontrer les relations suivantes~:
$$ A \cup(B \cap C) = (A \cup B) \cap(A \cup C) \quad\mbox{et} \quad
 A \cap(B \cup C) = (A \cap B) \cup(A \cap C).$$
\finenonce{000126}



\finexercice



\exercice{127, cousquer, 2003/10/01}

\enonce{000127}{}
Montrer que si $F$ et $G$ sont des sous-ensembles de $E$~:
$$ (F \subset G \iff F \cup G = G )\quad\mbox{et} \quad(F \subset G \iff
\complement F \cup G = E).$$
En déduire que~: 
$$ (F \subset G \iff F \cap G = F)
\quad\mbox{et} \quad(F \subset G \iff F \cap\complement G = \emptyset).$$
\finenonce{000127}



\finexercice



\exercice{128, cousquer, 2003/10/01}

\enonce{000128}{}
 Soit $E$ et $F$ des ensembles. Si $A \subset E$ et $B \subset F$
montrer que $A \times B \subset E\times F$.
\finenonce{000128}



\finexercice



\exercice{129, cousquer, 2003/10/01}

\enonce{000129}{}
 Soit $A= \{a_1, a_2, a_3, a_4\}$ et $B= \{b_1, b_2, b_3, b_4 ,b_5\}$.
Écrire le produit cartésien $A \times B$.
Quel est le nombre de parties de $A \times B$~?
\finenonce{000129}



\finexercice



\exercice{130, cousquer, 2003/10/01}

\enonce{000130}{}
 Soit $E$ un ensemble à $n$ éléments. Quel est le nombre d'éléments
de $E^p$~? Quel est le nombre de parties de $E^p$~?
\finenonce{000130}



\finexercice



\exercice{131, cousquer, 2003/10/01}

\enonce{000131}{}
$x$, $y$, $z$ étant des nombres réels, résoudre le système~:
$$\left\{\begin{array}{rcl}
 (x-1)(y-2)z & = & 0  \\
 (x-2)(y-3)  & = & 0
 \end{array}\right.$$
Représenter graphiquement l'ensemble des solutions.
\finenonce{000131}



\finexercice



\exercice{132, cousquer, 2003/10/01}

\enonce{000132}{}
Soit $A$ une partie de $E$, on appelle fonction caractéristique de~$A$
l'application~$f$ de~$E$ dans l'ensemble à deux éléments $\{0, 1\}$, telle
que~:
$$f(x)=\begin{cases}
0&  \text{ si } x\notin A \cr 1& \text{ si } x \in A \cr
\end{cases}$$
Soit $A$ et $B$ deux parties de $E$, $f$ et $g$ leurs fonctions
caractéristiques. Montrer que les fonctions suivantes sont les fonctions
caractéristiques d'ensembles que l'on déterminera~:
\begin{enumerate}
\item $1-f$.
\item $fg$.
\item $f+g-fg$.
\end{enumerate}
\finenonce{000132}



\finexercice



\exercice{133, cousquer, 2003/10/01}

\enonce{000133}{}
Soit un ensemble $E$ et deux parties $A$ et $B$ de~$E$. On désigne par
$A \triangle  B $ l'ensemble $ (A \cup B) \setminus(A \cap B)$.
\textit{Dans les questions ci-après il pourra être commode 
d'utiliser la notion de fonction caractéristique.}
\begin{enumerate}
\item Démontrer que $A \triangle B = (A \setminus B) \cup(B \setminus A)$.
\item Démontrer que pour toutes les parties $A$, $B$, $C$ de $E$ on a
$(A \bigtriangleup B) \bigtriangleup C = A \triangle(B \triangle C)$.
\item Démontrer qu'il existe une unique partie~$X$ de~$E$ telle que 
pour toute partie~$A$ de~$E$, $A \triangle X = X \triangle A = A$.
\item Démontrer que pour toute partie $A$ de $E$, il existe une partie~$A'$
de~$E$ et une seule telle que $A \triangle A' = A' \triangle A = X$.
\end{enumerate}
\finenonce{000133}



\finexercice



\exercice{134, cousquer, 2003/10/01}

\enonce{000134}{}
\begin{enumerate}
\item Écrire l'ensemble de définition de chacune des fonctions numériques
suivantes : $x\mapsto\sqrt x$, $x\mapsto\frac{1}{x-1}$, $x\mapsto\sqrt 
x+\frac{1}{x-1}$.
\item Simplifier $[1,3]\cap[2,4]$ et $[1,3]\cup[2,4]$.
\item Pour tout $n\in\Nn$, on note $n\Zz$ l'ensemble des entiers 
relatifs multiples de $n$ : $n\Zz=\{np \mid\nolinebreak p\in\Zz\}$. 
Simplifier $2\Zz\cap3\Zz$.
\end{enumerate}
\finenonce{000134}



\finexercice



\exercice{135, cousquer, 2003/10/01}

\enonce{000135}{}
On définit les cinq ensembles suivants :
\begin{eqnarray*}
A_1 & = & \left\{(x,y)\in\mathbb{R}^2\,,\; x+y<1\right\}\\
A_2 & = & \left\{(x,y)\in\mathbb{R}^2\,,\; |x+y|<1\right\}\\
A_3 & = & \left\{(x,y)\in\mathbb{R}^2\,,\; |x|+|y|<1\right\}\\
A_4 & = & \left\{(x,y)\in\mathbb{R}^2\,,\; x+y>-1\right\}\\
A_5 & = & \left\{(x,y)\in\mathbb{R}^2\,,\; |x-y|<1\right\}\\
\end{eqnarray*}

\begin{enumerate}
\item Représenter ces cinq ensembles.
\item En déduire une démonstration géométrique de
$$(|x+y|<1\;\mbox{ et }\;|x-y|<1) \Leftrightarrow |x|+|y|<1.$$
\end{enumerate}
\finenonce{000135}



\finexercice



\exercice{136, cousquer, 2003/10/01}

\enonce{000136}{}
Montrer que chacun des ensembles suivants est un intervalle,
éventuellement vide ou réduit à un point
$$I_1=\bigcap_{n=1}^{+\infty}\left[3,3+\frac{1}{n^2}\right[\;\mbox{ et
}\; I_2=\bigcap_{n=1}^{+\infty}\left]-2-\frac{1}{n},4+n^2\right].$$

\finenonce{000136}



\finexercice



\exercice{137, cousquer, 2003/10/01}
\video{a7iULB800tg}
\enonce{000137}{}

 Montrez que chacun des ensembles suivants est un intervalle que vous calculerez.
$$I=\bigcap_{n=1}^{+\infty}\left[-\frac{1}{n},2+\frac{1}{n}\right[\ \quad \mbox{
et }\quad \ J=\bigcup_{n=2}^{+\infty}\left[1+\frac{1}{n},n\right]$$

\finenonce{000137} 


\finexercice


\exercice{138, ridde, 1999/11/01}

\enonce{000138}{}
Soient $E$ un ensemble et $A, B, C$ trois parties de $E$ telles que
$A \cup B = A \cup C$ et $A \cap B = A \cap C$. Montrer que $B = C$.
\finenonce{000138}



\finexercice



\exercice{139, ridde, 1999/11/01}

\enonce{000139}{}
Soient $E$ un ensemble et $A, B, C$ trois parties de $E$.\\
Montrer que $ (A \cup B) \cap (B \cup C) \cap (C \cup A) =
(A \cap B) \cup (B \cap C) \cup (C \cap A)$.
\finenonce{000139}



\finexercice



\exercice{140, ridde, 1999/11/01}

\enonce{000140}{}
Donner les positions relatives de $A, B, C \subset E$ si $A \cup B = B \cap C$.
\finenonce{000140}



\finexercice



\exercice{141, ridde, 1999/11/01}

\enonce{000141}{}
Est-il vrai que $\mathcal{P} (A \cap B) = \mathcal{P} (A) \cap \mathcal{P} (B)$ ?
Et $\mathcal{P} (A \cup B) = \mathcal{P} (A) \cup \mathcal{P} (B)$ ?
\finenonce{000141}



\finexercice



\exercice{142, ridde, 1999/11/01}

\enonce{000142}{}
Montrer que $A \cap B = A \cap C \Leftrightarrow A \cap \complement B =
A \cap \complement C  $.
\finenonce{000142}



\finexercice



\exercice{143, ridde, 1999/11/01}

\enonce{000143}{}
Donner la liste des \'el\'ements de $\mathcal{P} (\mathcal{P} (\left\{
1, 2\right\}))$.
\finenonce{000143}



\finexercice



\exercice{144, ridde, 1999/11/01}
\enonce{000144}{}
Soient $A, B \subset E$. R\'esoudre les \'equations \`a l'inconnue $X \subset E$
\begin{enumerate}
\item $A \cup X = B$.
\item $A \cap X = B$.
\end{enumerate}

\finenonce{000144} 


\finexercice


\exercice{145, ridde, 1999/11/01}

\enonce{000145}{}
Soient $E, F, G$ trois ensembles. Montrer que $ (E \times G) \cup (F \times G) =
 (E \cup F)\times G$.
\finenonce{000145}



\finexercice



\exercice{146, ridde, 1999/11/01}

\enonce{000146}{}
Soient $E, F, G, H$ quatre ensembles. Comparer les ensembles
$ (E \times F) \cap (G \times H)$ et $ (E \cap G) \times (F \cap H)$.
\finenonce{000146}



\finexercice



\exercice{147, ridde, 1999/11/01}

\enonce{000147}{}
Soit $E$ l'ensemble des fonctions de $\Nn$ dans $\left\{ 1, 2, 3\right\}$.
Pour $i = 1, 2, 3$ on pose $A_i = \left\{ f \in E / f (0) = i \right\}$.
Montrer que les $A_i$ forment une partition de $E$.
\finenonce{000147}



\finexercice



\exercice{5112, rouget, 2010/06/30}
\enonce{005112}{**T}
$A$ et $B$ sont des parties d'un ensemble $E$. Montrer que :

\begin{enumerate}
\item  $(A\Delta B=A\cap B)\Leftrightarrow(A=B=\varnothing)$.
\item  $(A\cup B)\cap(B\cup C)\cap(C\cup A)=(A\cap B)\cup(B\cap C)\cup(C\cap A)$.
\item  $A\Delta B=B\Delta A$.
\item  $(A\Delta B)\Delta C=A\Delta(B\Delta C)$.
\item  $A\Delta B=\varnothing\Leftrightarrow A=B$.
\item  $A\Delta C=B\Delta C\Leftrightarrow A=B$.
\end{enumerate}
\finenonce{005112}


\finexercice


\exercice{5113, rouget, 2010/06/30}
\enonce{005113}{***IT}
Soient $(A_i)_{i\in I}$ une famille de parties d'un ensemble $E$ indéxée par un ensemble $I$ et $(B_i)_{i\in I}$ une
famille de parties d'un ensemble $F$ indéxée par un ensemble $I$. Soit $f$ une
application de $E$ vers $F$. Comparer du point de vue de l'inclusion les parties suivantes~:

\begin{enumerate}
\item  $f(\bigcup_{i\in I}A_i)$ et $\bigcup_{i\in I}f(A_i)$ (recommencer par $f(A\cup B)$ si on n'a pas les 
idées claires).
\item  $f(\bigcap_{i\in I}A_i)$ et $\bigcap_{i\in I}f(A_i)$.
\item  $f(E\setminus A_i)$ et $F\setminus f(A_i)$.
\item  $f^{-1}(\bigcap_{i\in I}B_i)$ et $\bigcap_{i\in I}f^{-1}(B_i)$.
\item  $f^{-1}(\bigcup_{i\in I}B_i)$ et $\bigcup_{i\in I}f^{-1}(B_i)$.
\item  $f^{-1}(F\setminus B_i)$ et $E\setminus f^{-1}(Bi)$.
\end{enumerate}
\finenonce{005113}


\finexercice


\exercice{5117, rouget, 2010/06/30}
\enonce{005117}{***I Théorème de \textsc{Cantor}}
\begin{enumerate}
\item   Montrer qu'il existe une injection de $E$ dans $\mathcal{P}(E)$.
\item  En considérant la partie $A=\{x\in E/\;x\notin f(x)\}$, montrer qu'il n'existe pas de bijection $f$ de $E$
sur $\mathcal{P}(E)$.
\end{enumerate}
\finenonce{005117}


\finexercice


\exercice{7186, megy, 2019/07/17}

\enonce{007186}{}
Soit $E$ un ensemble et  $\mathcal O$ une partie de $\mathcal P(E)$. On dit que $\mathcal O$ est une \emph{topologie sur $E$} si les conditions suivantes sont vérifiées
\begin{itemize}
\item $\mathcal O$ est stable par intersection finie, autrement dit : pour tout $n\in \N^*$ et toute famille $U_1, \cdots U_n$ d'éléments de $\mathcal O$, on a $\bigcap_{i=1}^n U_i\in \mathcal O$.
\item $\mathcal O$ est stable par union quelconque, autrement dit : pour tout ensemble $I$ et toute famille $(U_i)_{i\in I}$ d'éléments de $\mathcal O$, $\bigcup_{i\in I}U_i \in \mathcal O$.
\item Les parties $\emptyset$ et $E$ sont des éléments de $\mathcal O$.
\end{itemize}
\begin{enumerate}
\item Montrer que $\mathcal O_1=\{\emptyset, E\}$ et $\mathcal O_2=\mathcal P(E)$ sont des topologies sur $E$.
\item Montrer que 
\[ \mathcal O_3 = \left\{U\in \mathcal P(E)\:\middle|\: U=\emptyset \text{ ou }{}^cU\text{ est fini}\right\}
\]
est une topologie sur $E$.
\item Combien de topologies différentes y a-t-il si $E$ est l'ensemble vide ? S'il n'a qu'un seul élément ? Deux éléments ? Trois éléments ?
\end{enumerate}
\finenonce{007186}



\finexercice


\exercice{7187, megy, 2019/07/17}
\enonce{007187}{}
Dans l'ensemble $\R$, il existe une notion de \emph{partie bornée} : c'est une partie qui est incluse dans un segment du type $[-M,M]$, pour un certain $M$. Cet exercice montre comment généraliser cette notion de \emph{partie bornée} à un ensemble quelconque.

Soit $E$ un ensemble et  $\mathcal B$ une partie de $\mathcal P(E)$. On dit que $\mathcal B$ est une \emph{bornologie sur $E$} si les conditions suivantes sont vérifiées
\begin{itemize}
\item Si $A\in \mathcal B$ et $B\subseteq A$, alors $B\in \mathcal B$.
\item Si $A\in \mathcal B$ et $B \in \mathcal B$, alors $A\cup B\in \mathcal B$.
\item Pour tout $x\in E$, on a  $\{x\} \in \mathcal B$.
\end{itemize}
Les éléments de $\mathcal B$ sont dits \emph{$\mathcal B$-bornés}, ou simplement \emph{bornés} s'il n'y a pas d'ambiguïté sur la bornologie utilisée.

Dans la suite, on fixe un ensemble $E$.
\begin{enumerate}
\item Montrer que $\mathcal B_1=\{\emptyset, E\}$ est une bornologie de $E$. On l'appelle la \emph{bornologie triviale (ou : grossière)}.
\item Montrer que l'ensemble $\mathcal B_2$ des parties finies de $E$ est une bornologie de $E$. On l'appelle la \emph{bornologie discrète}.
\item Combien de bornologies différentes y a-t-il si $E$ est vide ? S'il contient (exactement) un élément ? Deux ? Trois ?
\item On suppose maintenant que $E=\R$. Soit $\mathcal B_3$ l'ensemble des parties $A\subseteq \R$ bornées au sens classique, autrement dit 
\[ A\in \mathcal B_3 \iff \exists M\in \R, \forall a\in A, |a|\leq M\]
Montrer que $\mathcal B_3$ est une bornologie. On l'appelle la \emph{bornologie usuelle sur $\R$}, et lorsqu'on parle de bornés de $\R$, il est implicite qu'on se réfère à cette bornologie (et non aux deux premières par exemple).
\end{enumerate}
\finenonce{007187}
\finexercice


\exercice{7188, megy, 2019/07/17}

\enonce{007188}{}
Soit $E$ un ensemble et  $\mathcal A$ une partie de $\mathcal P(E)$. On dit que $\mathcal A$ est une \emph{algèbre de parties $E$} si les conditions suivantes sont vérifiées:
\begin{itemize}
\item $\mathcal A$ n'est pas vide.
\item Si $X\in \mathcal A$, alors $E\setminus X$ aussi.
\item $\mathcal A$ est stable par union finie, autrement dit : pour tout $n\in \N^*$ et toute famille $U_1, \cdots U_n$ d'éléments de $\mathcal A$, on a $\bigcup_{i=1}^n U_i\in \mathcal A$.
\end{itemize}
\begin{enumerate}
\item Montrer que $\mathcal P(E)$ est une algèbre de parties de $E$.
\item Montrer  qu'une algèbre de parties de $E$ est stable par intersection finie.
\item Combien d'algèbres de parties y a-t-il si $E$ a (exactement) un, deux, ou trois éléments ?
\end{enumerate}
\finenonce{007188}



\finexercice

\section{100.03 Absurde et contraposée}


\exercice{148, bodin, 1998/09/01}

\enonce{000148}{}
Montrer que $\sqrt{2} \notin \Qq$.
\finenonce{000148}



\finexercice



\exercice{149, bodin, 1998/09/01}

\enonce{000149}{}
Soit $X$ un ensemble et $f$ une application de $X$ dans l'ensemble $\mathcal{P}(X)$ des
parties de $X$. On note $A$ l'ensemble des $x \in X$ v\'erifiant $x \notin f(x)$. D\'emontrer qu'il n'existe aucun $x \in X$ tel que $A=f(x)$.
\finenonce{000149}



\finexercice



\exercice{150, bodin, 1998/09/01}
\video{RDbMQTxMIuA}
\enonce{000150}{}
 Soit $(f_n)_{n\in\Nn}$ une suite d'applications de
l'ensemble $\Nn$ dans lui-m\^eme. On d\'efinit une application $f$
de $\Nn$ dans $\Nn$ en posant $f(n)=f_n(n)+1$. D\'emontrer qu'il
n'existe aucun $p\in \Nn$ tel que $f=f_p$.

\finenonce{000150} 


\finexercice


\exercice{151, bodin, 1998/09/01}
\video{RHiN60xhmoo}
\enonce{000151}{}
\begin{enumerate}
\item
Soit $p_{1},p_{2},\ldots ,p_{r}$, $r$ nombres premiers. Montrer que l'entier
$N=p_{1}p_{2}\ldots p_{r}+1$ n'est divisible par aucun des entiers $p_{i}$.
\item
Utiliser la question pr\'ec\'edente pour montrer par l'absurde qu'il existe une infinit\'e de
nombres premiers.
\end{enumerate}

\finenonce{000151} 


\finexercice

\section{100.04 Récurrence}


\exercice{152, bodin, 1998/09/01}

\enonce{000152}{}
D\'emontrer, en raisonnant par r\'ecurrence, que $10^{6n+2}+10^{3n+1}+1$ est divisible
par $111$ quel que soit $n\in \Nn$. (Indication : $1000=9\times 111 +1$ ).
\finenonce{000152}



\finexercice



\exercice{153, bodin, 1998/09/01}
\video{SUMAdDoMPto}
\enonce{000153}{}
\label{ex104}
Montrer :
\begin{enumerate}
    \item $\displaystyle{\sum_{k=1}^{n}k=\frac{n(n+1)}{2} \quad \forall n \in \Nn^{*} \,.}$
    \item $\displaystyle{\sum_{k=1}^{n}k^2=\frac{n(n+1)(2n+1)}{6} \quad \forall n \in \Nn^{*} \,.}$
\end{enumerate}

\finenonce{000153} 


\finexercice


\exercice{154, bodin, 1998/09/01}

\enonce{000154}{}
\textbf{En quoi le raisonnement suivant est-il faux ?} \\
Soit $\mathcal{P}(n)$ : $n$ crayons de couleurs sont tous de la m\^eme couleur.
\begin{itemize}
    \item $\mathcal{P}(1)$ est vraie car un crayon de couleur est de la m\^eme couleur
que lui-m\^eme.
    \item Supposons $\mathcal{P}(n)$. Soit $n+1$ crayons. On en retire $1$. Les $n$
crayons restants sont de la m\^eme couleur par hypoth\`ese de r\'ecurrence.

Reposons ce crayon et retirons-en un autre ; les $n$ nouveaux crayons sont \`a nouveau
de la m\^eme couleur. Le premier crayon retir\'e \'etait donc bien de la m\^eme couleur que les $n$ autres. La proposition est donc vraie au rang $n+1$.
    \item On a donc d\'emontr\'e que tous les crayons en nombre infini d\'enombrable
sont de la m\^eme couleur.
\end{itemize}
\finenonce{000154}



\finexercice



\exercice{155, bodin, 1998/09/01}
\video{6iVRvfkMLTw}
\enonce{000155}{}
Soit la suite $(x_n)_{n\in \Nn}$ d\'efinie par
$x_0=4$ et $\displaystyle{x_{n+1}=\frac{2x_n^2-3}{x_n+2}}$.
\begin{enumerate}
    \item Montrer que : $\forall n\in\Nn\quad x_n>3$.
    \item Montrer que : $\forall n\in\Nn \quad x_{n+1}-3>\frac{3}{2}(x_n-3)$.
    \item Montrer que : $\forall n\in\Nn \quad x_n \geqslant \left(\frac{3}{2}\right)^n+3$.
    \item La suite $(x_n)_{n\in\Nn}$ est-elle convergente ?

\end{enumerate}

\finenonce{000155} 


\finexercice


\exercice{156, bodin, 1998/09/01}
\enonce{000156}{}

\begin{enumerate}
\item
Dans le plan, on consid\`ere trois droites $\Delta_{1},\Delta_{2},\Delta_{3}$ formant un
``vrai'' triangle : elles ne sont pas concourantes, et il n'y en a pas deux parall\`eles.
Donner le nombre $R_{3}$ de r\'egions (zones blanches) d\'ecoup\'ees par ces trois droites.
\item
On consid\`ere quatre droites $\Delta_{1},\ldots,\Delta_{4}$, telles qu'il n'en existe pas
trois concourantes, ni deux parall\`eles. Donner le nombre $R_{4}$ de r\'egions d\'ecoup\'ees par
ces quatre droites.
\item
On consid\`ere $n$ droites $\Delta_{1},\ldots,\Delta_{n}$, telles qu'il n'en existe pas
trois concourantes, ni deux parall\`eles. Soit $R_{n}$ le nombre de r\'egions d\'elimit\'ees par
$\Delta_{1}\ldots\Delta_{n}$, et $R_{n-1}$ le nombre de r\'egions d\'elimit\'ees par
$\Delta_{1}\ldots\Delta_{n-1}$. Montrer que $R_{n}=R_{n-1}+n$.
\item
Calculer par r\'ecurrence le nombre de r\'egions d\'elimit\'ees par $n$ droites en position
g\'en\'erale, c'est-\`a-dire telles qu'il n'en existe pas trois concourantes ni deux parall\`eles.
\end{enumerate}
\finenonce{000156} 


\finexercice


\exercice{157, ridde, 1999/11/01}
\video{okAsZpWWB0M}
\enonce{000157}{}
Soit $X$ un ensemble. Pour $f \in \mathcal{F} (X, X)$, on d\'efinit $f^0 = id$ et
par r\'ecurrence pour $n \in \Nn$ $f^{n + 1} = f^n \circ f$.
\begin{enumerate}
\item Montrer que $\forall n \in \Nn$ $f^{n + 1} = f \circ f^n$.
\item Montrer que si $f$ est bijective alors $\forall n \in \Nn$ $ (f^{-1})^n
 = (f^n)^{-1}$.
 \end{enumerate}

\finenonce{000157} 


\finexercice


\exercice{158, gourio, 2001/09/01}

\enonce{000158}{}
Montrer que
$$\forall n\geq 2,n!\leq \left( \frac{n+1}{2}\right) ^{n}.$$
\finenonce{000158}



\finexercice



\exercice{159, cousquer, 2003/10/01}

\enonce{000159}{}
Pour tout entier naturel $n$, on pose 
$$ S_{n} = 1\cdot 2 + 2 \cdot 3+\cdots + (n-1)\cdot n$$
Démontrer que l'on a 
$$S_{n} = \frac{1}{3} n(n-1)(n+1)$$
\finenonce{000159}



\finexercice



\exercice{160, cousquer, 2003/10/01}

\enonce{000160}{}
Pour $n\in \mathbb{N}$ on considère la propriété suivante~:
$$P_{n}: \quad 2^n>n^2$$
\begin{enumerate}
\item Pour quelles valeurs de $n$ l'implication $P_{n} \Longrightarrow P_{n+1}$
est-elle vraie~?
\item Pour quelles valeurs de $n$ la propriété $P_{n}$ est-elle vraie~?
\end{enumerate}
\finenonce{000160}



\finexercice



\exercice{161, cousquer, 2003/10/01}

\enonce{000161}{}
Que pensez-vous de la démonstration suivante~?

\begin{enumerate}
\item Pour tout $n\geq2$, on considère la propriété~:
$$P(n) :\quad \mbox{$n$ points distincts du plan sont toujours alignés}$$
\item Initialisation~: $P(2)$ est vraie car deux points distincts 
sont toujours alignés.
\item Hérédité~: On suppose que $P(n)$ est vraie et on 
va démontrer $P(n+1)$.

 \noindent Soit donc $A_{1}, A_{2},\ldots,A_{n}, A_{n+1}$ des points distincts. 
D'après l'hypothèse de récurrence, $A_{1}, A_{2},\ldots, A_{n}$ sont 
alignés sur une droite~$d$, et $A_{2},\ldots, A_{n}, A_{n+1}$ sont 
alignés sur une droite~$d'$. Les deux droites $d$ et $d'$ ayant $n-1$
points communs $ A_{2},\ldots, A_{n}$ sont confondues. Donc
$A_{1}, A_{2},\ldots, A_{n}, A_{n+1}$ sont alignés,
ce qui montre l'hérédité de la propriété.
\item Conclusion~: la propriété $P(n)$ est vraie pour tout $n \geq 2$.
\end{enumerate}
\finenonce{000161}



\finexercice



\exercice{162, cousquer, 2003/10/01}

\enonce{000162}{}
\begin{enumerate}
\item Démontrer que pour tout entier naturel $n$, $9$ divise $10^n -1$.
\item Soit $k$ un entier strictement positif. Étudier la propriété suivante~:
pour tout entier naturel $n$, $k$ divise $(k+1)^n + 2$.
\end{enumerate}
\finenonce{000162}



\finexercice



\exercice{163, cousquer, 2003/10/01}

\enonce{000163}{}
Démontrer que pour $n\geq 1$, le produit de $n$ entiers impairs est un 
entier impair.
\finenonce{000163}



\finexercice



\exercice{164, cousquer, 2003/10/01}

\enonce{000164}{}
On considère une suite $(u_n)_{n\in \mathbb{N}}$ telle que ~:
$$u_0=0 \quad \mbox{et}\quad u_1=1 \quad\mbox{et}\quad 
\forall n\geq1,\; u_{n+1}=u_n+2u_{n-1}$$
Démontrer que~:
\begin{enumerate}
\item  $\forall n\in \mathbb{N},\; u_n\in \mathbb{N}$,
\item  $\forall n\in \mathbb{N},\; u_n=\frac{1}{3}(2^n-(-1)^n)$.
\end{enumerate}
\finenonce{000164}



\finexercice



\exercice{165, cousquer, 2003/10/01}

\enonce{000165}{}
 Soit $b\geq 2$ un entier fixé. Démontrer que pour tout 
$N\in \mathbb{N}^{\ast}$, il existe un entier $n\in \mathbb{N}$ 
et des entiers $a_0,a_1,\ldots, a_n$ appartenant à 
$\{\,0,1,\ldots,b-1\,\}$ tels que~;
$$N=a_0+a_1b+\cdots+a_nb^n\quad\mbox{et}\quad a_n\neq 0$$

 \noindent Démontrer que pour chaque $N$, le système $(n,a_0,a_1,\ldots, a_n)$ est déterminé par la propriété ci-dessus.

 \noindent On dit que $a_0,a_1,\ldots,a_n$ sont les chiffres de l'écriture du nombre $N$ suivant la base~$b$.
\finenonce{000165}



\finexercice



\exercice{166, cousquer, 2003/10/01}

\enonce{000166}{}
Démontrer par récurrence que pour tout $k\in\mathbb{N}$, $k!$ 
divise le produit de $k$ entiers consécutifs~:
$$\forall n\in \mathbb{N},\; k! \mid n(n+1)\cdots(n+k-1)$$
\finenonce{000166}



\finexercice



\exercice{167, cousquer, 2003/10/01}

\enonce{000167}{}
Les propriétés 
$$P_n \;:\; 3 \mid 4^n-1 \;,\; \forall n \in \mathbb{N},$$
et 
$$Q_n \;:\; 3 \mid 4^n+1 \;,\; \forall n \in \mathbb{N},$$
sont-elles vraies ou fausses ?
\finenonce{000167}



\finexercice



\exercice{168, cousquer, 2003/10/01}

\enonce{000168}{}
\begin{enumerate}
\item  Calculer les restes de la division euclidienne de $1,4,4^2,4^3$ par
  $3$. 
\item  Formuler, pour tout $n\in\mathbb{N}$, une hypothèse $\mathcal{P}(n)$ concernant le
  reste de la division euclidienne de $4^n$ par $3$. Démontrer que $\mathcal{P}(n)$
  est vérifiée pour tout $n\in\mathbb{N}$.
\item  Pour tout $n\in\mathbb{N}$, le nombre $16^n+4^n +3$ est-il divisible par $3$.
\end{enumerate}
\finenonce{000168}



\finexercice



\exercice{169, cousquer, 2003/10/01}

\enonce{000169}{}
D\'emontrer, en raisonnant par r\'ecurrence, que $3^{2n+2}-2^{n+1}$ est divisible
par $7$ quel que soit $n\in \mathbb{N}$. 
\finenonce{000169}



\finexercice



\exercice{170, cousquer, 2003/10/01}

\enonce{000170}{}
\begin{enumerate}
\item  Démontrer par récurrence:
$$ \sum_{k=0}^n k = \frac{n(n+1)}{2} $$
\item  Calculer de deux manières différentes:
$$ \sum_{k=1}^{n+1} k^3 - \sum_{k=0}^n (k+1)^3. $$
\item  En déduire:
 $$ \sum_{k=0}^n k^2 = {1 \over 6} (2n^3+3n^2+3n). $$
\end{enumerate}
\finenonce{000170}



\finexercice



\exercice{171, cousquer, 2003/10/01}

\enonce{000171}{}
 Montrer que pour tout entier $n\geq 1$:
$$\frac 1{1.2} + \frac 1{2.3} + \ldots + \frac 1{n.(n+1)} = \frac n{n+1}.$$
\finenonce{000171}



\finexercice



\exercice{172, cousquer, 2003/10/01}

\enonce{000172}{}
Démontrer, en le déterminant qu'il existe un entier $n_0$ tel que
$$
\forall n\geq n_0,\  2^n\geq (n+2)^2.
$$
\finenonce{000172}



\finexercice



\exercice{173, cousquer, 2003/10/01}

\enonce{000173}{}
D\'emontrer par r\'ecurrence sur $n$ que pour tout $n\geq 2$
l'implication $$ [x>-1,x\not =0] \Rightarrow [ (1+x)^n >1+nx]$$
est vraie.
\finenonce{000173}



\finexercice



\exercice{174, cousquer, 2003/10/01}

\enonce{000174}{}
\begin{enumerate}
\item  Soit $n\in\mathbb{N}$; montrer que pour tout entier $k\geq 1$ on a
$$n^k+kn^{k-1}\leq (n+1)^k.$$
\item  Soit $b$ un réel positif ou nul. Montrer par récurrence, que
pour tout $n\geq 1$ on a
$$ (1+b)^n\leq 1+\frac{nb}{1!}+\frac{(nb)^2}{2!}+...+\frac{(nb)^n}{n!}.$$
\end{enumerate}
\finenonce{000174}



\finexercice



\exercice{175, cousquer, 2003/10/01}

\enonce{000175}{}
Montrer par récurrence que pour tout entier $n\in\mathbb{N}$, 
$$\left(a+b\right)^n=\sum_{k=0}^{n}C_n^k a^kb^{n-k},$$
pour tout réel $a$ et $b$.
\finenonce{000175}



\finexercice



\exercice{176, cousquer, 2003/10/01}

\enonce{000176}{}
 On définit une suite $(F_n)$ de la fa\c con suivante :
$$F_{n+1} = F_n+F_{n-1} ; \quad F_0=1 , F_1=1 \ .$$

\begin{enumerate}
\item  Calculer $F_n$ pour $1<n<10$.
\item  Montrer que l'équation $x^2=x+1$ admet une unique solution positive 
$a$ que l'on calculera.
\item  Montrer que, pour tout $n\geq 2$, on a
$$a^{n-2}< F_n < a^{n-1} \ .$$
\end{enumerate}
\finenonce{000176}



\finexercice



\exercice{177, cousquer, 2003/10/01}

\enonce{000177}{}
Montrer que :
$$2\cos \frac{\pi}{2^n} = \sqrt{2+\sqrt{2+\ldots\sqrt{2}}}.$$
\finenonce{000177}



\finexercice



\exercice{178, cousquer, 2003/10/01}

\enonce{000178}{}
 Pour $n \in \mathbb{N},$ $n \geq 2,$ trouver une loi simplifiant le produit :
$$(1-\frac 14)...(1-\frac 1n).$$
\finenonce{000178}



\finexercice



\exercice{179, cousquer, 2003/10/01}

\enonce{000179}{}
Pour $n \in \mathbb{N},$   soient $a_0,\ldots,a_n$ des nombres réels de même signe tel que $a_i > -1,$ montrer que:
$$(1+a_0)...(1+a_n) > 1+a_0+\ldots+a_n.$$ 
\finenonce{000179}



\finexercice



\exercice{7011, megy, 2016/08/25}

\enonce{007011}{} 
Montrer $\forall n \in \N,\: \sum_{k=0}^n k^3=\frac{n^2(n+1)^2}{4}$.
\finenonce{007011}



\finexercice


\exercice{7012, megy, 2016/08/25}

\enonce{007012}{}
Montrer que pour tout entier $n$ positif, l'entier $10^n - (-1)^n$ est divisible par $11$.
\finenonce{007012}



\finexercice


\exercice{7013, megy, 2016/08/25}

\enonce{007013}{}
Soit $(u_n)_{n\in\N}$ la suite de nombres réels définie par  $u_0=0$ et pour tout $n$ positif, $u_{n+1} = \sqrt{3u_n+4}$. Montrer que la suite est majorée par $4$.
\finenonce{007013}



\finexercice


\exercice{7014, megy, 2016/08/25}

\enonce{007014}{}
Soit $(u_n)_{n\in\N}$ la suite de nombres réels définie par  $u_0=0$ et pour tout $n$ positif, $u_{n+1} = 2u_n+1$. Calculer  $u_n$ en fonction de $n$.
\finenonce{007014}



\finexercice


\exercice{7015, megy, 2016/08/25}

\enonce{007015}{}
Soit $(u_n)_{n\in\N}$ la suite de nombres réels définie par $u_0=1$, $u_1=2$ et pour tout $n$ positif, $u_{n+2} = 5u_{n+1}-6u_n$. Calculer  $u_n$ en fonction de $n$.
\finenonce{007015}



\finexercice


\exercice{7016, megy, 2016/08/25}

\enonce{007016}{}
Soit $(u_n)_{n\in\N}$ la suite de nombres réels définie par $u_0=1$, $u_1=1$ et pour tout $n$ positif, $u_{n+2} = u_{n+1}+\frac{2}{n+2}u_n$. Montrer : $\forall n\in \N^*, \: 1\leq u_n \leq n^2$.
\finenonce{007016}



\finexercice


\exercice{7017, megy, 2016/08/25}

\enonce{007017}{} 
Montrer que pour tout $n\in \N$, la somme des $n$ premiers entiers positifs impairs est toujours le carré d'un entier.
\finenonce{007017}



\finexercice


\exercice{7018, megy, 2016/08/25}

\enonce{007018}{} 
Montrer : $\forall u\in \R, \forall n\in \N, |\sin(nu)|\leq n|\sin(u)|$.
\finenonce{007018}



\finexercice


\exercice{7019, megy, 2016/08/25}

\enonce{007019}{}
% Bernoulli amélioré
\begin{enumerate}
\item Soit $a \in \R_+$. Montrer $ \forall n\in\N^*,\:
(1+a)^n \geq 1+na + \frac{n(n-1)}{2}a^2$.

\item Soit $(u_n)_{n\in\N}$ la suite définie par $u_n = \frac{3n}{3^n}$. Montrer que pour tout $n\in\N^*$, on a $0\leq u_n \leq\frac{3n}{2n^2+1}$. 
\end{enumerate}
\finenonce{007019}



\finexercice


\exercice{7020, megy, 2016/08/25}

\enonce{007020}{}
Soit $a\in]0,\pi/2[$, et définissons une suite réelle par $u_0=2\cos(a)$ et pour tout $n\in\N$, $u_{n+1}=\sqrt{2+u_n}$. Montrer que pour tout $n\in\N$, on a $u_n=2\cos\left(\frac{a}{2^n}\right)$.
\finenonce{007020}



\finexercice


\exercice{7021, megy, 2016/08/25}
\enonce{007021}{} 
%(Bac Antilles-Guyane 2005)
Définissons une suite par $u_0=1$ et pour tout $n \in\N$, $u_{n+1} = \frac12 u_n+n-1$. 
\begin{enumerate}
\item Démontrer que pour tout $n\geq 3$, $u_n$ est positif. En déduire que pour tout $n\geq 4$, on a $u_n\geq n-2$. En déduire la limite de la suite.
\item Définissons maintenant la suite $v_n=4u_n-8n+24$. Montrer que la suite $(v_n)$ est une suite géométrique, donner son premier terme et sa raison. Montrer que pour tout $n\in \N, u_n = 7\left(\frac{1}{2}\right)^n+2n-6$. Remarquer que $u_n$ est la somme d'une suite géométrique et d'une suite arithmétique dont on précisera les raisons et les premiers termes. En déduire une formule pour la quantité $u_0+u_1+...+u_n$ en fonction de l'entier $n$.
\end{enumerate}
\finenonce{007021}
\finexercice


\exercice{7022, megy, 2016/08/25}

\enonce{007022}{}  
On considère la suite réelle $(u_n)_{n\in\N}$ définie par $u_0=2$ et pour tout $n\in\N,\: u_{n+1}=\sqrt{u_n}$.
\begin{enumerate}
\item Montrer que pour tout $n\in\N$, $u_n>1$.
\item Montrer que pour tout réel $a\in]1;+\infty[$, on a $\frac{1}{\sqrt a+1} \leq \frac12$.
\item En déduire que pour tout $n\in\N$, on a $u_{n+1}-1\leq \frac12\left(u_n-1\right)$.
\item Montrer que pour tout $n\in\N$, $u_n-1 \leq \left(\frac12\right)^n$. En déduire la limite de la suite $(u_n)$.
\end{enumerate}
\finenonce{007022}



\finexercice



\exercice{7023, megy, 2016/08/25}

\enonce{007023}{Récurrence de Cauchy et application}
% https://www.artofproblemsolving.com/wiki/index.php?title=Cauchy_Induction
Soit $A$ une partie de $\N^*$ contenant $1$ et telle que
\begin{enumerate}
\item $\forall n\in \N^*,\: n\in A \Rightarrow 2n \in A$;
\item $\forall n\in \N^*,\: n+1\in A \Rightarrow n \in A$.
\end{enumerate}
Montrer que $A = \N^*$.\\
En déduire l'inégalité arithmético-géométrique : si $a_1, ..., a_n$ sont des réels positifs, alors on a 
\[  \sqrt[n]{a_1 a_2  ...   a_n} \leq \frac{a_1+...+a_n}{n} \]
\finenonce{007023}



\finexercice



\exercice{7024, megy, 2016/08/25}

\enonce{007024}{}
Démontrer que tout entier $n\geq 1$ peut s'écrire comme somme de puissances de deux distinctes.
\finenonce{007024}



\finexercice



\exercice{7025, megy, 2016/08/25}

\enonce{007025}{}
Démontrer que tout entier $n\geq 1$ peut s'écrire de façon unique sous la forme $2^p(2q+1)$, avec $p$ et $q$ entiers.
\finenonce{007025}



\finexercice


\exercice{7035, megy, 2016/11/20}

\enonce{007035}{} 
Montrer que pour tout $n>1$  on a l'inégalité 
\[\frac 1{\sqrt 1} + \frac 1{\sqrt 2} + \cdots + \frac 1{\sqrt n} > \sqrt n.\]
\finenonce{007035}



\finexercice


\exercice{7036, megy, 2016/11/20}

\enonce{007036}{Une récurrence descendante}
Montrer que pour tout entier $N\geq 2$,
\[ \sqrt{2\sqrt{3\sqrt{4\sqrt{...\sqrt{(N-1)\sqrt{N}}}}}} <  3.\]
\finenonce{007036}



\finexercice


\exercice{7037, megy, 2016/11/20}

\enonce{007037}{Inégalité du binôme}
Montrer que pour tous $a,b > 0$ distincts et tout $n >1$, on a l'inégalité
\[ 2^{n-1} (a^n + b^n) > (a+b)^n.\]
\finenonce{007037}



\finexercice


\exercice{7038, megy, 2016/11/20}

\enonce{007038}{Variantes du raisonnement par récurrence}
Parmi les énoncés suivants, lesquels permettent d'en déduire que $P_n$ est vraie pour tout $n \in \N$ ?
\begin{enumerate}
\item $P_0$ et 
$\forall n\in\N, P_n \Rightarrow 
(P_{2n}\wedge P_{2n+1})$;
\item $P_0$, $P_1$ et 
$\forall n\geq 1, P_n \Rightarrow 
(P_{2n}\wedge P_{2n+1})$;
\item $P_0$, $P_1$, $P_2$ et 
$\forall n\geq 2, P_n \Rightarrow 
(P_{2n}\wedge P_{2n+1})$;
\item $P_0$, $P_1$ et 
$\forall n\geq 1, P_n \Rightarrow 
(P_{n-1}\wedge P_{n+1})$.
\end{enumerate}
\finenonce{007038}



\finexercice


\exercice{7039, megy, 2016/11/20}

\enonce{007039}{Conducteur d'un sous-monoïde}
Soit $P_n$ une assertion dépendant de $n \in \N$ telle que:
\begin{enumerate}
\item $P_0$ est vraie;
\item $\forall n \in \N, P_n \Rightarrow 
(P_{n+3} \text{ et } P_{n+4})$.
\end{enumerate}
La propriété est-elle vraie pour tout $n \in \N$ ? Pour $n$ assez grand  ? (Et si oui à partir de quel rang ?) Pour quels entiers est-elle vraie ?

Répondre aux deux premières questions en remplaçant dans l'énoncé les nombres $3$ et $4$ par des paramètres entiers positifs $a$ et $b$ quelconques.
% si a et b sont premiers entre eux, le conducteur est (a-1)(b-1), voir un livre d'algèbre commutative, chapitre cloture intégrale ?
\finenonce{007039}



\finexercice


\exercice{7040, megy, 2016/11/20}

\enonce{007040}{Nombres de Catalan}
On définit une suite $(C_n)_{n\in\N}$ par $C_0=1$ et pour tout naturel $n$,  
$C_{n+1} = \sum_{k=0}^{n}C_kC_{n-k}$.
\begin{enumerate}
\item Calculer les cinq premiers termes de la suite;
\item Montrer par récurrence que pour tout $n\geq 0$,  $C_n \geq 2^{n-1}$;
\item Montrer par récurrence forte que pour tout $n\geq 0$,  $C_n \geq 3^{n-2}$;
\item Tenter de montrer par une récurrence similaire à la précédente que pour tout $n\geq 0$,  $C_n \geq 4^{n-2}$. À quel endroit ceci échoue-t-il ? Pourquoi est-il heureux que cela échoue ?
\end{enumerate}
\finenonce{007040}



\finexercice


\exercice{7041, megy, 2016/11/20}

\enonce{007041}{} 
% récurrence forte
Soit $x$ un réel tel que $x+\frac{1}{x}$ soit entier. Montrer que pour tout $n\in \N$, $x^n+\frac{1}{x^n}$ est entier.
\finenonce{007041}



\finexercice


\exercice{7042, megy, 2016/11/20}

\enonce{007042}{} 
% valeurs particulières, conjecturer une formule,
% récurrence double ou forte
Soit $(u_n)_{n\in\N}$ la suite réelle définie par $u_0=2$, $u_1=3$, et pour tout $n\in\N$, $u_{n+2}=3u_{n+1}-2u_n$. Déterminer $u_n$ en fonction de $n$.
\finenonce{007042}



\finexercice


\exercice{7043, megy, 2016/11/20}

\enonce{007043}{} 
% valeurs particulières, conjecturer une formule,
% récurrence forte
Soit $(u_n)_{n\in\N}$ la suite réelle définie par $u_0=1$ et pour tout $n\in\N$, $u_{n+1}=\sum_{k=0}^n u_k$. Déterminer $u_n$ en fonction de $n$.
\finenonce{007043}



\finexercice


\exercice{7044, megy, 2016/11/20}

\enonce{007044}{}
Soit $n\in \N^*$. On trace $n$ cercles dans le plan. Montrer que l'on peut colorier chaque région du plan ainsi délimitée avec exactement deux couleurs, de manière à ce que deux régions séparées par un arc de cercle soient toujours de couleur différente.
\finenonce{007044}



\finexercice


\exercice{7045, megy, 2016/11/20}

\enonce{007045}{} 
Soit $n$ un entier supérieur ou égal à $2$. On place $2n$ points dans l'espace, et on trace $n^2+1$ segments entre ces points. Montrer que l'on a tracé au moins un triangle.
\finenonce{007045}



\finexercice


\exercice{7046, megy, 2016/11/20}

\enonce{007046}{} 
% Un peu difficile. Récurrence forte !!!!!
Déterminer les valeurs de $n$ pour lesquelles le nombre
\[ u_n := 1+\frac12 + \frac 13 + ... + \frac1n\]
est entier.
\finenonce{007046}



\finexercice

\section{100.05 Relation d'équivalence, relation d'ordre}


\exercice{207, bodin, 1998/09/01}

\enonce{000207}{}
\begin{enumerate}
    \item Soit $E=\Nn\times\Nn$, on d\'efinit $\mathcal{R}$ par :
$(a,b)\mathcal{R}(a',b')\Leftrightarrow a+b'=b+a'$. Montrer que $\mathcal{R}$ est une
relation d'\'equivalence. Identifier $E/\mathcal{R}$.
    \item M\^emes questions avec $E=\Zz\times\Nn^*$ et $(p,q)\mathcal{R}(p',q') \Leftrightarrow pq'=p'q$.
\end{enumerate}
\finenonce{000207}



\finexercice



\exercice{208, bodin, 1998/09/01}

\enonce{000208}{}
Dans $\Rr^2$ on d\'efinit la relation $\mathcal{R}$ par :
$$(x,y)\mathcal{R}(x',y')\Leftrightarrow y=y'.$$
\begin{enumerate}
    \item Montrer que $\mathcal{R}$ est une relation d'\'equivalence.
    \item D\'eterminer la classe d'\'equivalence de $(x,y)\in \Rr^2$.
\end{enumerate}
\finenonce{000208}



\finexercice



\exercice{209, bodin, 1998/09/01}
\video{vf0_R5F6cfc}
\enonce{000209}{}

 Dans $\Cc$ on d\'efinit la relation $\mathcal{R}$
par :
$$z\mathcal{R}z'\Leftrightarrow |z|=|z'|.$$
\begin{enumerate}
    \item Montrer que $\mathcal{R}$ est une relation d'\'equivalence.
    \item D\'eterminer la classe d'\'equivalence de chaque $z \in \Cc$.
\end{enumerate}
\finenonce{000209} 


\finexercice


\exercice{210, bodin, 1998/09/01}
\enonce{000210}{}

 Soit $\mathcal{R}$ une relation binaire sur un
ensemble $E$, sym\'etrique et transitive. Que penser du
raisonnement suivant ?
\begin{center}
``$x\mathcal{R} y \Rightarrow y\mathcal{R} x$ car $\mathcal{R}$ est sym\'etrique, \\
or $(x \mathcal{R} y \text{ et } y\mathcal{R} x) \Rightarrow x \mathcal{R} x$ car $\mathcal{R}$
est transitive,\\
donc $\mathcal{R}$ est r\'eflexive.''
\end{center}
\finenonce{000210} 


\finexercice


\exercice{211, gourio, 2001/09/01}

\enonce{000211}{}
\'Etudier la relation $\mathcal{R}$ d\'{e}finie sur $\R^{\R}$ (l'ensemble des
applications de $\R$ dans $\R$) par:
$$f\mathcal{R} g\Longleftrightarrow \exists A>0,\forall x\in \R,\left| x\right|>A\Rightarrow f(x)=g(x).$$
\finenonce{000211}



\finexercice



\exercice{212, gourio, 2001/09/01}
\video{t7QwqxS5RzQ}
\enonce{000212}{}

Montrer que la relation $\mathcal{R}$ d\'{e}finie sur $\Rr$ par :
$$x\mathcal{R} y\Longleftrightarrow xe^{y}=ye^{x}$$
est une relation d'\'{e}quivalence.
Pr\'{e}ciser, pour $x$ fix\'{e} dans $\Rr$, le nombre d'\'{e}l\'{e}ments de
la classe de $x$ modulo $\mathcal{R}$.
\finenonce{000212} 


\finexercice


\exercice{213, bodin, 1998/09/01}

\enonce{000213}{}
La relation ``divise'' est-elle une relation d'ordre sur $\Nn$ ?
sur $\Zz$ ? Si oui, est-ce une relation d'ordre total ?
\finenonce{000213}



\finexercice



\exercice{214, bodin, 1998/09/01}

\enonce{000214}{}
\'Etudier les propri\'et\'es des relations suivantes. Dans le cas
d'une relation d'\'equivalence, pr\'eciser les classes ; dans le cas
d'une relation d'ordre, pr\'eciser si elle est totale, si l'ensemble admet
un plus petit ou plus grand \'el\'ement.
\begin{enumerate}
    \item Dans $\mathcal{P}(E)$ : $A\mathcal{R}_1 B \Leftrightarrow A\subset B\quad ;
\quad A\mathcal{R}_2 B \Leftrightarrow A\cap B=\emptyset$.
    \item Dans $\Zz$ : $a\mathcal{R}_3 b \Leftrightarrow a$ et $b$ ont la m\^eme parit\'e $\quad;
\quad a\mathcal{R}_4 b \Leftrightarrow \exists n\in \Nn \ \,a-b=3n \quad ;
\quad a\mathcal{R}_5 b \Leftrightarrow a-b$ est divisible par $3$.
\end{enumerate}
\finenonce{000214}



\finexercice



\exercice{215, ridde, 1999/11/01}

\enonce{000215}{}
Soient $ (X, \leq)$ et $ (Y, \leq)$ deux ensembles ordonn\'es (on note abusivement
les deux ordres de la même façon). On d\'efinit sur $X \times Y$ la relation
$ (x, y) \leq (x', y')$ ssi $ (x < x')$ ou $ (x = x'$ et $y \leq y')$.
Montrer que c'est un ordre et qu'il est total ssi $X$ et $Y$ sont totalement
ordonn\'es.
\finenonce{000215}



\finexercice



\exercice{216, ridde, 1999/11/01}

\enonce{000216}{}
Un ensemble est dit bien ordonn\'e si toute partie non vide admet un plus petit
\'el\'ement.
\begin{enumerate}
\item Donner un exemple d'ensemble bien ordonn\'e et un exemple d'ensemble qui ne
l'est pas.
\item Montrer que bien ordonn\'e implique totalement ordonn\'e.
\item La r\'eciproque est-elle vraie ?
\end{enumerate}
\finenonce{000216}



\finexercice



\exercice{217, ridde, 1999/11/01}
\video{6f56RN5_li0}
\enonce{000217}{}

Soit $ (E, \leq)$ un ensemble ordonn\'e. On d\'efinit sur
$\mathcal{P} (E)\setminus\left\{ \emptyset \right\}$ la relation
$\prec$ par 
$$X \prec Y \quad \text{ ssi } \quad (X = Y \ \text{ ou } \ 
\forall x \in X \  \forall y \in Y \  x \leq y).$$ 
V\'erifier que c'est une relation d'ordre.
\finenonce{000217} 


\finexercice


\exercice{218, ridde, 1999/11/01}

\enonce{000218}{}
Montrer que $a*b = \dfrac{a + b}{1 + ab}$ est une l.c.i  sur $]-1, 1[$ et d\'eterminer
ses propri\'et\'es.
\finenonce{000218}



\finexercice



\exercice{3030, quercia, 2010/03/08}

\enonce{003030}{Congruence des carr{\'e}s modulo 5}
On d{\'e}finit la relation $\sim$ sur $\Z$ par
$x \sim y \iff x^2 \equiv y^2 \mathrm{mod}\, 5$.
\begin{enumerate}
  \item  D{\'e}terminer l'ensemble quotient.
  \item  Peut-on d{\'e}finir une addition quotient ? une multiplication quotient ?
\end{enumerate}
\finenonce{003030}



\finexercice 


\exercice{3031, quercia, 2010/03/08}

\enonce{003031}{Produit cart{\'e}sien}
Soient deux relations d'{\'e}quivalence : ${\cal R}$ sur $E$, et ${\cal S}$ sur $F$.
On d{\'e}finit sur $E \times F$ :
$$(x,y) \sim (x',y') \iff x {\cal R} x' \text{ et } y {\cal S} y'.$$
\begin{enumerate}
  \item   V{\'e}rifier que $\sim$ est une relation d'{\'e}quivalence.
  \item   Soit $\phi : {E\times F} \to 
            {\left(E/{\cal R}\right) \times \left(F/{\cal S}\right)},
            {(x,y)}  \mapsto {(\dot x, \dot y)}$

      D{\'e}montrer que $\phi$ est compatible avec $\sim$, et que l'application
      quotient associ{\'e}e est une bijection.
\end{enumerate}
\finenonce{003031}



\finexercice 


\exercice{3032, quercia, 2010/03/08}

\enonce{003032}{$X \cup A = Y \cup A$}
Soit $E$ un ensemble et $A \subset E$. On d{\'e}finit la relation sur ${\cal P}(E)$ :
$$X \sim Y \iff X \cup A = Y \cup A.$$
\begin{enumerate}
\item Montrer que c'est une relation d'{\'e}quivalence.
\item Soit $\phi : {{\cal P}(E)} \to {{\cal P}(E\setminus A)}, X
  \mapsto {X \setminus A}$.
  \\
  Montrer que $\phi$ est compatible avec $\sim$, et que l'application
  quotient associ{\'e}e est une bijection.
\end{enumerate}
\finenonce{003032}



\finexercice 


\exercice{3033, quercia, 2010/03/08}

\enonce{003033}{{\'E}quivalences sur $E^E$}
Soit $E$ un ensemble non vide. On consid{\`e}re les relations sur $F = E^E$:
\begin{align*}
  f \sim g &\iff \exists\ n \in \N^* \text{ tq } f^n = g^n,\cr f
  \approx g &\iff \exists\ m,n \in \N^* \text{ tq } f^n = g^m,\cr f
  \equiv g &\iff f(E) = g(E). 
\end{align*}
\begin{enumerate}
  \item Montrer que $\sim$, $\approx$, $\equiv$ sont des relations d'{\'e}quivalence.
  \item Pour $f \in F$, on note $f^\sim$, $f^\approx$, $f^\equiv$ les classes
    d'{\'e}quivalence de $f$ modulo $\sim$, $\approx$, $\equiv$.
  \begin{enumerate}
    \item Comparer $f^\sim$, $f^\approx$.
    \item Montrer que toute classe d'{\'e}quivalence pour $\approx$ est r{\'e}union de
        classes d'{\'e}quivalence pour $\sim$.
    \item Que pouvez-vous dire de $f$ s'il existe $g \in f^\approx$ injective ?
        surjective ?
    \item M{\^e}me question avec $f^\equiv$.
   \end{enumerate}
\end{enumerate}
\finenonce{003033}



\finexercice 


\exercice{3034, quercia, 2010/03/08}

\enonce{003034}{Relation d'{\'e}quivalence quotient}
Soient ${\cal R}$ et ${\cal S}$ deux relations d'{\'e}quivalence sur un ensemble $E$,
telles que : $$\forall\ x,y \in E,\ x {\cal R} y  \Rightarrow  x {\cal S} y.$$

On d{\'e}finit $\dot{{\cal S}}$ sur $E/{\cal R}$ par :
$\dot x \dot{{\cal S}} \dot y \iff x {\cal S} y$.

V{\'e}rifier que $\dot{{\cal S}}$ est une relation d'{\'e}quivalence, puis d{\'e}finir une
bijection entre $(E/{\cal R})/\dot{{\cal S}}$ et $E/{\cal S}$.
\finenonce{003034}



\finexercice 


\exercice{3035, quercia, 2010/03/08}

\enonce{003035}{Compl{\'e}tion d'une relation r{\'e}flexive et transitive}
Soit ${\cal R}$ une relation binaire sur un ensemble $E$ r{\'e}flexive et transitive.
On d{\'e}finit les deux relations :
$$\begin{matrix}  x {\cal S} y &\iff &( x {\cal R} y \text{ et } y {\cal R} x),\cr
            x {\cal T} y &\iff &( x {\cal R} y \text{ ou } y {\cal R} x).\end{matrix}$$

Est-ce que ${\cal S}$ et ${\cal T}$ sont des relations d'{\'e}quivalence ?
\finenonce{003035}



\finexercice 


\exercice{3036, quercia, 2010/03/08}

\enonce{003036}{Parties satur{\'e}es pour une relation d'{\'e}quivalence}
Soit $\sim$ une relation d'{\'e}quivalence sur un ensemble $E$. Pour $A \subset E$,
on d{\'e}finit $s(A) = \bigcup_{x \in A} \dot x$.
\begin{enumerate}
  \item   Comparer $A$ et $s(A)$.
  \item   Simplifier $s(s(A))$.
  \item   Montrer que : $\forall\ x \in E$, on a
      $(x \in s(A)) \iff (\dot x \cap s(A) \ne \varnothing)$.
      En d{\'e}duire $s(E\setminus s(A))$.
  \item   D{\'e}montrer que $s\left(\bigcup_{i \in I} A_i \right) = \bigcup_{i \in I} s(A_i)$
      et $s\left(\bigcap_{i \in I} A_i \right) \subset \bigcap_{i \in I} s(A_i)$.
  \item   Donner un exemple d'inclusion stricte.
\end{enumerate}
\finenonce{003036}



\finexercice 


\exercice{3037, quercia, 2010/03/08}

\enonce{003037}{Ordre sur les fonctions}
Soit $X$ un ensemble et $E = \R^X$. On ordonne $E$ par :
$f \le g \iff \forall\ x \in X,\ f(x) \le g(x)$.
\begin{enumerate}
  \item  V{\'e}rifier que c'est une relation d'ordre.
  \item  L'ordre est-il total ?
  \item  Comparer les {\'e}nonc{\'e}s : {\it``$f$ est major{\'e}e''}, et {\it``$\{f\}$ est major{\'e}''}.
  \item  Soit $(f_i)_{i\in I}$ une famille major{\'e}e de fonctions de $E$. Montrer
     qu'elle admet une borne sup{\'e}rieure.
\end{enumerate}
\finenonce{003037}



\finexercice 


\exercice{3038, quercia, 2010/03/08}

\enonce{003038}{$\sup \circ \inf$ et  $\inf \circ \sup$}
Soit $f : {\R^2} \to {\R}$ une fonction born{\'e}e. On d{\'e}finit les
fonctions :
$$g : {\R} \to {\R},  t \mapsto {\sup\{ f(t,y)\text{ tq } y\in \R\}}$$
 $$ h : {\R} \to {\R}, t \mapsto {\inf\{ f(x,t)\text{ tq } x\in \R\}}$$

Montrer que $g$ et $h$ sont born{\'e}es, puis comparer $\sup h$ et $\inf g$.
\finenonce{003038}



\finexercice


\exercice{3039, quercia, 2010/03/08}

\enonce{003039}{Ordre lexicographique}
On note $E = [-1, 1]^2$, et on d{\'e}finit sur $E$ la relation :\\
$(x,y) \preceq (x',y') \iff
  \Bigl( (x < x') \text{ ou }( x = x' \text{ et } y \le y')\Bigr)$
\hfill({\it ordre lexicographique}).
\begin{enumerate}
  \item  Pour $(a,b) \in E$, repr{\'e}senter graphiquement l'ensemble des majorants
     de $(a,b)$.
  \item  Soit $A$ une partie non vide de $E$. Montrer que $A$ admet une borne sup{\'e}\-rieu\-re.
\end{enumerate}
\finenonce{003039}



\finexercice 


\exercice{3040, quercia, 2010/03/08}

\enonce{003040}{Distance entre un point et une partie}
Pour $A \subset \R$ non vide et born{\'e}e, et $x \in \R$, on note :
$$d(x,A) = \inf\{ |x-a| \text{ tq } a \in A\}  \qquad (\text{\it distance de $x$ {\`a} $A$}).$$

Montrer que $\bigl|d(x,A) - d(y,A)\bigr| \le |x-y|$.
\finenonce{003040}



\finexercice 


\exercice{3041, quercia, 2010/03/08}

\enonce{003041}{Parties adjacentes}
Soient $A,B \subset \R$ v{\'e}rifiant :
$$\begin{cases}\forall\ a \in A,\ \forall\ b \in B,\ a \le b \cr
  \forall\ \epsilon > 0,\ \exists\ a \in A,\ \exists\ b \in B \text{
    tq } b-a \le \epsilon\end{cases}$$
(on dit que $A$ et $B$ sont {\it adjacentes\/}). Montrer que $\sup(A) = \inf(B)$.
\finenonce{003041}



\finexercice 


\exercice{3042, quercia, 2010/03/08}

\enonce{003042}{borne sup $ \Rightarrow $ borne inf}
Soit $E$ ordonn{\'e} tel que toute partie non vide et major{\'e}e admet une borne sup{\'e}rieure.
Montrer que toute partie non vide et minor{\'e}e admet une borne inf{\'e}rieure.
\finenonce{003042}



\finexercice


\exercice{3043, quercia, 2010/03/08}

\enonce{003043}{Ordre sur $\R^2$}
On d{\'e}finit sur $\R^2$ : $(x,y) \ll (x',y') \iff |x'-x| \le y'-y$.
\begin{enumerate}
  \item  V{\'e}rifier que c'est une relation d'ordre.
  \item  Dessiner les ensembles des majorants et des minorants d'un couple $(a,b)$.
  \item  L'ordre est-il total ?
  \item  Soit $A = \{ (x,y) \in \R^2$ tq $x^2 + y^2 \le 1 \}$.
     D{\'e}terminer $\sup(A)$.
\end{enumerate}
\finenonce{003043}


\finexercice 


\exercice{3044, quercia, 2010/03/08}

\enonce{003044}{Propri{\'e}t{\'e}s de sup et inf}
Un treillis est un ensemble ordonn{\'e} $E$ dans lequel
pour tous $x,y \in E,\ \sup(x,y)$ et $\inf(x,y)$ existent.
Soit $E$ un treillis.

\begin{enumerate}
  \item  Montrer que $\sup$ et $\inf$ sont des op{\'e}rations associatives.
  \item  A quelle condition ont-elles des {\'e}l{\'e}ments neutres ?
  \item  Montrer que :
        \begin{align*}
          &\forall\ x,y \in E,& &\sup\bigl(x,\inf(x,y)\bigr) =
          \inf\bigl(x,\sup(x,y)\bigr) = x,\\
          &\forall\ x,y,z \in E,& &x \le z \Rightarrow
          \sup\bigl(x,\inf(y,z)\bigr) \le
          \inf\bigl(\sup(x,y),z\bigr),\\
          &\forall\ x,y,z \in E,& &\inf\bigl(x,\sup(y,z)\bigr) \ge
          \sup\bigl(\inf(x,y),\inf(x,z)\bigr).
    \end{align*}
\end{enumerate}
\finenonce{003044}



\finexercice 


\exercice{3045, quercia, 2010/03/08}

\enonce{003045}{Ordre d{\'e}duit d'une loi idempotente}
Soit $\cdot$ une op{\'e}ration commutative et associative sur $E$, telle que :
$\forall\ x \in E,\ x\cdot x = x$.

On d{\'e}finit la relation $\le$ sur $E$ par : $x \le y \iff x\cdot y = x$

\begin{enumerate}
  \item  Reconna{\^\i}tre $\le$ quand $\cdot$ est $\cap$ sur ${\cal P}(X)$\quad
     (resp $\cup$).
  \item  Montrer que $\le$ est une relation d'ordre.
  \item  D{\'e}montrer que : $\forall\ x,y \in E,\ x\cdot y = \inf(x,y)$.
\end{enumerate}
\finenonce{003045}



\finexercice 


\exercice{3046, quercia, 2010/03/08}

\enonce{003046}{Borne sup{\'e}rieure parmi les intervalles}
Soit $E$ l'ensemble des intervalles de $\R$ (y compris $\varnothing$)
ordonn{\'e} par l'inclusion.

Soient $I,J$ deux intervalles. Qu'est-ce que $\inf(I,J)$ ? $\sup(I,J)$ ?
\finenonce{003046}



\finexercice 


\exercice{3047, quercia, 2010/03/08}

\enonce{003047}{Prolongement d'applications}
Soit $E$ un ensemble et
${\cal E} = \{ (A,f)$ tq $A \subset E,\ A\ne\varnothing,$ et $f\in E^A \}$.
On ordonne $\cal E$ par :
$$(A,f) \preceq (B,g) \iff \begin{cases}A \subset B \cr
                                       \forall\ x \in A,\ f(x) = g(x)\end{cases}$$

(c'est-{\`a}-dire que la fonction $g$, d{\'e}finie sur $B$, prolonge la fonction $f$,
d{\'e}finie seulement sur $A$).
\begin{enumerate}
  \item  Montrer que $\preceq$ est une relation d'ordre. L'ordre est-il total ?
  \item  Soient $(A,f)$ et $(B,g)$ deux {\'e}l{\'e}ments de $\cal E$. Trouver une CNS
     pour que la partie $\{ (A,f), (B,g) \}$ soit major{\'e}e.
     Quelle est alors sa borne sup{\'e}rieure ?
  \item  M{\^e}me question avec minor{\'e}e.
\end{enumerate}
\finenonce{003047}



\finexercice 


\exercice{3048, quercia, 2010/03/08}

\enonce{003048}{Point fixe d'une fonction croissante}
Soit $f : {[0,1]} \to {[0,1]}$ croissante.
On note $A = \{ x \in {[0,1]}$ tq $f(x) \le x \}$.
\begin{enumerate}
  \item  D{\'e}montrer que $A$ n'est pas vide.
  \item  D{\'e}montrer que $f(A) \subset A$.
  \item  Soit $a = \inf(A)$. Montrer que $f(a)$ minore $A$.
  \item  En d{\'e}duire que $f(a) = a$.
\end{enumerate}
Cela prouve que toute application croissante de $[0,1]$ dans lui-m{\^e}me admet
un point fixe. Montrer que c'est faux pour l'intervalle $[0,1[$.
\finenonce{003048}



\finexercice 


\exercice{3049, quercia, 2010/03/08}

\enonce{003049}{Relation d'ordre sur un ensemble quotient}
Soit $\cal R$ une relation sur $E$ r{\'e}flexive et transitive. On d{\'e}finit la
relation : $x \sim y \iff x {\cal R} y$ et $y {\cal R} x$.
\begin{enumerate}
  \item Montrer que $\sim$ est une relation d'{\'e}quivalence sur $E$.\par

Sur $E/\sim$ on pose : $\dot x \le \dot y \iff x {\cal R} y$.

  \item Montrer que cette d{\'e}finition est ind{\'e}pendante des repr{\'e}sentants $x$ et $y$
    choisis.
  \item Montrer que $\le$ est une relation d'ordre sur $E/\sim$.
\end{enumerate}
\finenonce{003049}



\finexercice 


\exercice{3050, quercia, 2010/03/08}

\enonce{003050}{Pas de borne sup{\'e}rieure dans $\Q$}
Dans cet exercice, on admet que : $\forall\ x \in \Q,\ x^2 \ne 2$.
\begin{enumerate}
  \item Soient $A = \{\ x \in \Z^{+*}$ tq $x^2 < 2\ \}$
        et $B = \{\ x \in \Z^{+*}$ tq $x^2 > 2\ \}$.
    D{\'e}terminer $\sup(A)$ et $\inf(B)$.

  \item Soient $A = \{\ x \in \Q^{+*}$ tq $x^2 < 2\ \}$
        et $B = \{\ x \in \Q^{+*}$ tq $x^2 > 2\ \}$.
    On veut d{\'e}montrer que $A$ n'admet pas de borne sup{\'e}rieure
    {\it dans $\Q$}. Pour cela, on suppose au contraire que
    $\alpha = \sup(A)$ existe ($\alpha \in \Q$), et on pose~$\beta = \frac 2\alpha$.
  \begin{enumerate}
    \item Montrer que $\beta = \inf(B)$.
    \item Montrer que : $\forall\ a \in A,\ \forall\ b \in B,$ ona $a \le b$.
        Que pouvez-vous en d{\'e}duire pour $\alpha$ et $\beta$ ?
    \item Obtenir une contradiction en consid{\'e}rant
        $\gamma = \frac {\alpha + \beta}2$.
   \end{enumerate}
\end{enumerate}
\finenonce{003050}



\finexercice 


\exercice{7189, megy, 2019/07/23}

\enonce{007189}{}

Soit $E$ l'ensemble des droites du plan. Le parallélisme et l'orthogonalité sont-elles des relations réflexives, symétriques, antisymétriques, transitives ?

\finenonce{007189}



\finexercice


\exercice{7190, megy, 2019/07/23}

\enonce{007190}{}

Soit $E$ un ensemble fini, de cardinal $n$. Combien de relations binaires y a-t-il sur $E$ ? De relations symétriques ? Réflexives ?

\finenonce{007190}



\finexercice


\exercice{7191, megy, 2019/07/23}

\enonce{007191}{}

Soit $\leq$ une relation d'ordre sur un ensemble $E$, et $<$ la relation d'ordre strict associée, c'est-à-dire par définition : $x<y \iff x\leq y \text{ et } x\neq y$. Est-ce que le contraire de $x\leq y$ est $y<x$ ?
% non. C'est vrai si l'ordre est total

\finenonce{007191}



\finexercice


\exercice{7192, megy, 2019/07/23}

\enonce{007192}{}

Soit $E$ un ensemble fini et $f : E\to E$ une involution, c'est-à-dire une application vérifiant $f\circ f=\operatorname{Id}$. Montrer que si $f$ n'a pas de points fixes, alors $|E|$ est pair. Plus généralement, montrer que la parité de $|E|$ est celle du nombre de points fixes de $f$.

\finenonce{007192}



\finexercice


\exercice{7193, megy, 2019/07/23}

\enonce{007193}{}

Soit $\mathcal R$ la relation d'équivalence la plus fine sur $\{0,1,2\}$ vérifiant $0\mathcal R 1$. Décrire le graphe de $\mathcal R$ (donner tous ses éléments).

\finenonce{007193}



\finexercice


\exercice{7194, megy, 2019/07/23}

\enonce{007194}{}
(Coordonnées polaires) 
Soit $\sim$ la relation d'équivalence la plus fine sur $\R_+\times [-\pi,\pi]$ vérifiant les conditions: 
\[
\begin{cases}
\forall \theta,\theta'\in [-\pi,\pi], (0,\theta)\sim (0,\theta')\\
\forall r\in \R_+^*, (r,-\pi)\sim (r,\pi)
\end{cases}
\]
Décrire le graphe de $\sim$, ainsi que ses classes d'équivalence.

\finenonce{007194}



\finexercice


\exercice{7195, megy, 2019/07/23}

\enonce{007195}{}
Soit $l>0$ un réel et $X=[0,l]\times [-1,1]$, et $\sim$ la relation d'équivalence la plus fine sur $X$ telle que $(0,y)\sim(l,-y)$ pour tout $y\in [-1,1]$.  Décrire le graphe et les classes d'équivalence de la relation. 


Note : l'ensemble quotient $\mathcal M = X/\sim$ est donc l'ensemble obtenu en recollant le rectangle $X=[0,l]\times [-1,1]$ le long de deux bords opposés, en suivant une orientation opposée. On l'appelle le \emph{ruban de Möbius} (de longueur $l$).

\finenonce{007195}



\finexercice


\exercice{7196, megy, 2019/07/23}

\enonce{007196}{}
Soient $\mathcal R$ et $\mathcal S$ des relations binaires sur $E$. On dit que $\mathcal R$ est plus fine que $\mathcal S$, ou encore que c'est est un raffinement, si $\forall x, y\in E, x\mathcal R y \implies x\mathcal S y$.  De façon équivalente, $\mathcal R$ est plus fine que $\mathcal S$ si on a l'inclusion de graphes $\Gamma_{\mathcal R} \subseteq \Gamma_{\mathcal S}$.
\begin{enumerate}
\item Montrer que \og être plus fine que \fg{} est une relation d'ordre sur l'ensemble des relations binaires sur $E$.
\item Soient $\mathcal R$ et $\mathcal S$ des relations binaires sur $E$. Montrer qu'il existe une relation binaire sur $E$ qui raffine à la fois $\mathcal R$ et $\mathcal S$, et  qu'il existe aussi une relation binaire sur $E$ simultanément moins fine que $\mathcal R$ et $\mathcal S$.
\end{enumerate}
\finenonce{007196}



\finexercice


\exercice{7197, megy, 2019/07/23}

\enonce{007197}{}
Soit $f : \R\to \mathbb U, t\mapsto e^{it}$, et soit $\mathcal R$ la relation d'équivalence sur $\R$ définie par $x\mathcal R y \iff x\equiv y \pmod{2\pi}$. On note $\R/2\pi\Z$ l'ensemble quotient $\R/\mathcal R$. Montrer que l'application $f$ descend au quotient en une application $[f] :\R/2\pi\Z \to \mathbb U$ qui est une bijection.

\finenonce{007197}



\finexercice


\exercice{7198, megy, 2019/07/23}

\enonce{007198}{}
(Produit de deux relations) 
Soient $\mathcal R$ et $\mathcal S$ deux relations sur $E$. Leur \emph{produit}, noté $\mathcal R \mathcal S$, est la relation binaire définie par:
\[ \forall x,y\in E, x \mathcal R \mathcal S y 
\iff \exists a\in E, \: (x \mathcal R a \text{ et } a \mathcal S y)
\]
\begin{enumerate}
\item Prouver par un exemple qu'en général, les relations $\mathcal R \mathcal S$ et $\mathcal S \mathcal R$ sont distinctes.
\item Montrer que le produit de relations est néanmoins associatif, autrement dit si $\mathcal R$, $\mathcal S$ et $\mathcal T$ sont trois relations, on a 
\[ (\mathcal R \mathcal S) \mathcal T = \mathcal R (\mathcal S \mathcal T)\]
\end{enumerate}
\finenonce{007198}



\finexercice


\exercice{7199, megy, 2019/07/23}

\enonce{007199}{}
(Clôture transitive. Cet exercice utilise la notion de produit de relations)
Soit $\mathcal R$ une relation sur $E$. Pour $n\in \N$ et $\mathcal R$ est une relation sur $E$, on définit alors par récurrence la relation $\mathcal R^n$ (en définissant $\mathcal R^0$ comme l'égalité, puis $\mathcal R^{n+1} = \mathcal R \mathcal R^n$).

Montrer que toutes les relations suivantes sont égales:
\begin{enumerate}
\item $\bigvee_{n\geq 0} \mathcal R^n$;
\item la relation dont le graphe est $\bigcup_{n\geq 0} \Gamma_{\mathcal R^{n}}$;
\item la relation dont le graphe est l'intersection de tous les graphes de relations transitives qui contiennent $\Gamma_{\mathcal R}$.
\item la plus fine relation  parmi toutes les relations transitives moins fines que $\mathcal R$.
\end{enumerate}

Cette relation binaire (qui est donc transitive) est appelée \emph{clôture transitive} de $\mathcal R$.
Montrer que si $\mathcal R$ est symétrique (resp. réflexive), sa clôture transitive l'est également.
\finenonce{007199}



\finexercice


\exercice{7200, megy, 2019/07/23}

\enonce{007200}{}
Soit $E$ l'ensemble des couples de la forme $(I,f)$, où $I$ est un intervalle de $\R$ et $f$ est une fonction de $I$ dans $\R$.

La relation $\preceq$ sur $E$ est définie par 
\[ (I,f)~\preceq~(J,g) ~\iff~ (I\subseteq J \text{ et } f=g|_I).\]
Montrer qu'il s'agit d'une relation d'ordre.
\finenonce{007200}



\finexercice


\exercice{7201, megy, 2019/07/23}

\enonce{007201}{}
Soit $E=\R^\R$ l'ensemble des fonctions de $\R$ dans $\R$ et $f, g\in E$. On dit que $f$ et $g$ ont \og même germe en zéro\fg{} et on note $f \underset{0}{=} g$ si:
\[  \exists \epsilon>0, f|_{]-\epsilon,\epsilon[} = g|_{]-\epsilon,\epsilon[} \]
\begin{enumerate}
\item Montrer que $\underset{0}{=}$ est une relation d'équivalence sur $E$.
\item Montrer  que si $f \underset{0}{=} g$ alors $f(0)=g(0)$, mais que la réciproque est fausse.
\item Montrer également que pour tout $a\in \R^*$, il existe deux fonctions $f$ et $g$ avec $f \underset{0}{=} g$ et $f(a)\neq g(a)$.
\end{enumerate}

La classe d'équivalence d'une fonction $f$ pour cette relation d'équivalence s'appelle le \emph{germe de $f$ en zéro}.

Attention, cette relation d'équivalence n'est \textbf{pas} \og l'équivalence en zéro\fg{} qui sera par la suite introduite dans le cours d'analyse.
\finenonce{007201}



\finexercice


\exercice{7202, megy, 2019/07/23}

\enonce{007202}{}
Soit $f : E\to F$, soit $\equiv_f$ la relation d'équivalence sur $E$ dont les classes d'équivalence sont les fibres de $f$, et soit $Q = E/\equiv_f$ l'ensemble quotient. 

\begin{enumerate}
\item Montrer que $f$ passe au quotient en une application $\bar f : Q\to F$ qui est injective.
\item Montrer qu'une relation d'équivalence $\mathcal R$ sur $E$ est plus fine que $\equiv_f$ si et seulement si $f$ passe au quotient par $\mathcal R$.
\item En déduire quelles sont les relations d'équivalence les plus et moins fines telles que $f$ passe au quotient par $\mathcal R$.
\end{enumerate}
\finenonce{007202}



\finexercice


\exercice{7203, megy, 2019/07/23}

\enonce{007203}{}
(Coégalisateur) 
Soient $A$ et $B$ deux ensembles et $f$ et $g$ deux applications entre $A$ et $B$. On définit sur $B$ la relation binaire suivante : $\mathcal R$ est la relation d'équivalence la plus fine telle que $\forall a\in A, f(a)\mathcal R g(a)$. Le \emph{coégalisateur de $f$ et $g$} est par définition l'ensemble quotient $C = B/\mathcal R$. On note $\pi : B \to C$ la surjection canonique sur le quotient. On a alors $\pi\circ f = \pi \circ g$.

Montrer que $C$ et $\pi$ vérifient la propriété suivante (dite \emph{propriété universelle du coégalisateur}):


Pour tout ensemble $X$ et application $\phi : B \to X$ vérifiant $\phi\circ f = \phi\circ g$,  il existe une unique application $h : C\to X$ telle que $\phi = h \circ \pi $.

\finenonce{007203}



\finexercice


\exercice{7204, megy, 2019/07/23}

\enonce{007204}{}
(Somme amalgammée d'ensembles. Cet exercice utilise la notion de coégalisateur.) 
Soient $A$,  $B$ et $C$  des ensembles et $f :  C\to A$, $g : C\to B$ des applications.

Soit $A\coprod B$ l'union disjointe de $A$ et $B$ et $i_A$ et $i_B$ les injections canoniques de $A$ et $B$ dans $A\coprod B$.

Les deux applications $i_A \circ f$ et $i_B \circ g$ vont toutes deux de $C$ dans $A\coprod B$.  Leur coégalisateur est appelé \emph{la somme amalgamée de $A$ et $B$ sous $C$}, est noté $A\coprod_C B$. La surjection canonique  $A\coprod B \to A\coprod_C B$ est notée $\pi$ et on note $j_A = \pi \circ i_A$ et $j_B = \pi \circ i_B$.

Montrer que $A\coprod_C B$ vérifie la propriété universelle suivante:

Pour tout ensemble  $D$ muni d'applications $\phi : A\to D$ et $\psi : B\to D$, il existe une unique application $h : A\coprod_C B \to D$ telle que $\phi = h\circ j_A$ et $\psi = h\circ j_B$.

\finenonce{007204}



\finexercice


\exercice{7205, megy, 2019/07/23}

\enonce{007205}{}
(Écrasement d'une partie d'un ensemble) 
Soit $X$ un ensemble. Pour tout sous-ensemble $A\subseteq X$, on définit la relation binaire $\sim_A$ sur $X$ comme suit:
\[ \forall (x,y)\in X^2, \: x\sim_A y \iff \left(x=y \text{ ou } \left(x\in A\text{ et } y\in A\right)\right).\]
\begin{enumerate}
\item Montrer que c'est une relation d'équivalence sur $X$. Quelles sont ses classes d'équivalence ?
\item Soit $f$ une fonction de $X$ dans un ensemble $E$, constante sur $A$. Montrer qu'elle descend au quotient en une application $[f] : X/\sim_A \to E$.
\item Montrer que pour tout ensemble $E$, l'application 
\[
\phi : \left\{f\in \mathcal F(X,E),\:\middle| \: f\text{ est constante sur } A\right\} \to \mathcal F(X/\sim_A, E),
\]
qui à $f$ associe $[f]$ est surjective.
\item Identifier, parmi les relations d'équivalence étudiées dans le cours et les exercices du chapitre, celles qui sont des cas particuliers d'écrasements de parties.
\end{enumerate}
\finenonce{007205}



\finexercice


\exercice{7206, megy, 2019/07/23}

\enonce{007206}{}
(Cône sur un ensemble) 
Soit $X$ un ensemble et $Y = X\times [0,1]$. Soit $\mathcal R$ la relation d'équivalence la plus fine  sur $Y$ telle que $\forall x,x'\in X, (x,0)\mathcal R (x',0)$.
\begin{enumerate}
\item Montrer que $(x,t)\mathcal R (x',t') \iff \left(t=t'=0)\text{ ou } (x,t)=(x',t')\right)$.
\item Le cône sur $X$, noté $\operatorname{Cone}(X)$, est par définition $Y/\mathcal R$. Le nom de \og cône\fg{} peut s'expliquer à l'aide de l'exemple suivant. Définir une bijection entre $\operatorname{Cone}(\mathbb S^1)$ et l'ensemble
\[ \left\{(x,y,z)\in \R^3\:\middle|\: z^2=x^2+y^2, \text{ et } 0\leq z \leq 1\right\}\]
(qui est un vrai cône au sens usuel: faire un dessin). 
\end{enumerate}
\finenonce{007206}



\finexercice


\exercice{7207, megy, 2019/07/23}

\enonce{007207}{}
(Suspension d'un ensemble)
Soit $X$ un ensemble. Sur l'ensemble $X\times [-1,1]$, on considère la relation d'équivalence la plus fine vérifiant:
\[ \begin{cases}
\forall x,x'\in X, (x,-1)\mathcal R (x',-1)\\
\forall x,x'\in X, (x,1)\mathcal R (x',1)
\end{cases}
\] 
\begin{enumerate}
\item Montrer que 
\[ (x,t)\mathcal R (x',t') \iff \big( t=t'=-1 \text{ ou } t=t'=1\text{ ou } (x,t)= (x',t')\big)\] L'ensemble quotient est appelé \emph{suspension de $X$}, et est noté $S(X)$.
\item Soit $X=\{-1,1\}$. Montrer que l'application $f : X\times[-1,1] \to \R^2, \: (x,t) \mapsto (t,x\sqrt{1-t^2})$ est à valeurs dans le cercle unité du plan, noté $\mathbb S^1$, et passe au quotient en application injective de $S(X)$ vers $\R^2$ dont l'image est $\mathbb S^1$. Ceci formalise la phrase \og la suspension de deux points est un cercle.\fg
\end{enumerate}

(Note : plus généralement, on peut montrer que pour tout $n\in\N$, la suspension de la sphère $\mathbb S^n$ est en bijection naturelle avec la sphère $\mathbb S^{n+1}$. Cet exercice traite le cas $n=0$.
)
\finenonce{007207}



\finexercice


\exercice{7208, megy, 2019/07/23}

\enonce{007208}{}
(Union/disjonction et intersection/conjonction de deux relations)
Soient $\mathcal R$ et $\mathcal S$ deux relations sur $E$. On définit la disjonction (ou union), notée $\mathcal R \vee \mathcal S$, par : 
\[ x (\mathcal R \vee \mathcal S) y \iff (x \mathcal R y\text{ ou } x \mathcal S y)\]
De façon équivalente, le graphe de $\mathcal R \vee \mathcal S$ est l'union des graphes de $\mathcal R$ et de $\mathcal S$. De même, on définit la conjonction (ou intersection) $\mathcal R \wedge \mathcal S$ comme la relation dont le graphe est l'intersection des deux graphes de $\mathcal R$ et $\mathcal S$, c'est-à-dire 
\[ x (\mathcal R \wedge \mathcal S) y \iff (x \mathcal R y\text{ et } x \mathcal S y).\]

Si $\mathcal R$ et $\mathcal S$ sont des relations d'équivalence, montrer que $\mathcal R \wedge \mathcal S$ est une relation d'équivalence, mais pas forcément $\mathcal R \vee \mathcal S$.

(Note : on peut définir la conjonction ou la disjonction d'un nombre quelconque de relations, à l'aide de l'union ou de l'intersection des graphes associés.)
\finenonce{007208}



\finexercice
\section{100.99 Autre}


\exercice{180, gourio, 2001/09/01}

\enonce{000180}{}
Quels sont les entiers $n$ $ $tels que $4^{n}\leq n!$ ?
\finenonce{000180}



\finexercice



\exercice{181, gourio, 2001/09/01}

\enonce{000181}{}
Montrer que :
$$\forall n\geq 2,u_{n}=\sum\limits_{k=1}^{n}\frac{1}{k}\notin \Nn.$$
\emph{Indication} : montrer que
$$\forall n\geq 2,\exists (p_{n},q_{n})\in (\Nn^{*})^{2},u_{n}=\frac{2p_{n}+1}{2q_{n}}.$$
\finenonce{000181}



\finexercice



\exercice{182, gourio, 2001/09/01}

\enonce{000182}{}
Soit $f:\Nn^{*}\rightarrow \Nn^{*}$ une application v\'{e}rifiant :
$$\forall n\in \Nn^{*},f(n+1)>f(f(n)).$$
Montrer que $f=Id_{\Nn^{*}}.$
\emph{Indications} : que dire de $k\in \Nn$ tel que $f(k)=\inf \{f(n)|n\in \Nn\}$? En
d\'{e}duire que $\forall n>0,f(n)>f(0).$
Montrer ensuite que $\forall n\in \Nn, $ on a: $\forall m>n,f(m)>f(n)$ et $%
\forall m\leq n,f(m)\geq m $ (on pourra introduire $k$ tel que $f(k)$ soit
le plus petit entier de la forme $f(m)  $ avec $m>n$).
En d\'{e}duire que $f$ est strictement croissante et qu'il n'existe qu'une
seule solution au probl\`{e}me. Laquelle ?
\finenonce{000182}



\finexercice



\exercice{183, ridde, 1999/11/01}

\enonce{000183}{}
Pour $p \in \left\{ 1, 2, 3\right\}$ on note $S_{p} = \sum\limits_{k = 0}^n k^p$.
\begin{enumerate}
\item A l'aide du changement d'indice $i = n-k$ dans $S_{1}$, calculer $S_{1}$.
\item Faire de même avec $S_{2}$. Que se passe-t-il ?
\item Faire de même avec $S_{3}$ pour l'exprimer en fonction de $n$ et $S_{2}$.
\item En utilisant l'exercice \ref{ex104}, calculer $S_{3}$.
\end{enumerate}
\finenonce{000183}



\finexercice



\exercice{184, ridde, 1999/11/01}

\enonce{000184}{}
Pour calculer des sommes portant sur deux indices, on a int\'erêt
\`a repr\'esenter la z{o}ne du plan couverte par ces indices et
\`a sommer en lignes, colonnes ou diagonales... Calculer :
\begin{enumerate}
\item $\sum\limits_{1\leq i \leq j \leq n}ij$.
\item $\sum\limits_{1\leq i < j \leq n}i (j-1)$.
\item $\sum\limits_{1\leq i < j \leq n} (i-1)j$.
\item $\sum\limits_{1\leq i \leq j \leq n} (n-i) (n-j)$.
\item $\sum\limits_{1\leq p, q \leq n} (p + q)^2$ (on posera $k = p + q$).
\end{enumerate}
\finenonce{000184}



\finexercice


\section{101.01 Application}


\exercice{185, bodin, 1998/09/01}
\video{MJjD9ZZkUuE}
\enonce{000185}{}
Soient $f : \Rr \rightarrow \Rr$ et $g : \Rr \rightarrow \Rr$ telles que $f(x) = 3x+1$ et $g(x)=x^2-1$.
A-t-on $f\circ g=g\circ f$ ?

\finenonce{000185} 


\finexercice


\exercice{186, cousquer, 2003/10/01}

\enonce{000186}{}
Soit l'application de $\mathbb{R}$ dans $\mathbb{R}$, 
$f\colon x\mapsto x^2$.
\begin{enumerate}
\item Déterminer les ensembles suivants : $f([-3,-1])$, $f([-2,1])$,
$f([-3,-1]\cup[-2,1])$ et $f([-3,-1]\cap[-2,1])$. Les comparer.
\item Mêmes questions avec les ensembles $f^{-1}(\mathopen]-\infty,2])$,
$f^{-1}([1,+\infty\mathclose[)$,
$f^{-1}(\mathopen]-\infty,2]\cup\nolinebreak{}[1,+\infty\mathclose[)$
et $f^{-1}(\mathopen]-\infty,2]\cap\nolinebreak{}[1,+\infty\mathclose[)$.
\end{enumerate}
\finenonce{000186}



\finexercice



\exercice{2889, quercia, 2010/03/08}

\enonce{002889}{Images directes et r{\'e}ciproques}
Soit $f : E \to F$ une application, $A,A' \subset E$ et $B,B' \subset F$.
\begin{enumerate}
  \item Simplifier $f(f^{-1}(f(A)))$ et $f^{-1}(f(f^{-1}(B)))$.
  \item Montrer que $f(A\cap f^{-1}(B)) = f(A)\cap B$.
  \item Comparer $f(A \mathop{\Delta} A')$ et $f(A) \mathop{\Delta} f(A')$.
  \item Comparer $f^{-1}(B \mathop{\Delta} B')$ et
    $f^{-1}(B) \mathop{\Delta} f^{-1}(B')$.
  \item A quelle condition sur $f$ a-t-on :
    $\forall\ A \subset E,\ f(E\setminus A) = F\setminus f(A)$ ?
\end{enumerate}
\finenonce{002889}



\finexercice 


\exercice{2890, quercia, 2010/03/08}

\enonce{002890}{$(X \cap A, X \cap B)$}
Soit $E$ un ensemble, et $A,B$ deux parties fix{\'e}es de $E$.
Soit $ \phi : {{\cal P}(E)} \to {{\cal P}(A) \times {\cal P}(B)},
      X \mapsto {(X \cap A, X \cap B).}$
\begin{enumerate}
  \item Qu'est-ce que $\phi(\varnothing)$ ? $\phi( E\setminus(A \cup B))$ ?
  \item A quelle condition sur $A$ et $B$, $\phi$ est-elle injective ?
  \item Est-ce que le couple $(\varnothing, B)$ poss{\`e}de un ant{\'e}c{\'e}dent par $\phi$ ?
  \item A quelle condition sur $A$ et $B$, $\phi$ est-elle surjective ?
\end{enumerate}
\finenonce{002890}



\finexercice 


\exercice{2891, quercia, 2010/03/08}

\enonce{002891}{Partie stable par une application}
Soit $f : E \to E$. Pour $n \in \N^*$, on note
$f^n = \underbrace{f \circ f \circ \dots \circ f}_{n \text{ fois}}$\ ,
et $f^0 = \text{id}_E$.

Soit $A \subset E$, $A_n = f^n(A)$, et $B = \bigcup_{n\in \N} A_n$.
\begin{enumerate}
  \item Montrer que $f(B) \subset B$.
  \item Montrer que $B$ est la plus petite partie de $E$ stable par $f$ et
    contenant $A$.
\end{enumerate}
\finenonce{002891}



\finexercice 


\exercice{2892, quercia, 2010/03/08}

\enonce{002892}{Factorisation d'une application}
\begin{enumerate}
  \item Soit $f : F \to  E$ et $g : G  \to E$ deux applications.
    Montrer qu'il existe une
    application $h : G \to  F$ telle que $g = f\circ h$ si et seulement si :
    $g(G) \subset f(F)$.\\
    A quelle condition $h$ est-elle unique ?
  \item Soit $f: E \to F$ et $g : E \to G$ deux applications.
    Montrer qu'il existe une
    application $h: F\to G$ telle que $g = h\circ f$ si et seulement si :
    $\forall\ x,y \in E,\ \bigl( f(x) = f(y)  \Rightarrow  g(x) = g(y) \bigr)$.
    \\
    A quelle condition $h$ est-elle unique ?
\end{enumerate}
\finenonce{002892}



\finexercice 


\exercice{2893, quercia, 2010/03/08}

\enonce{002893}{Propri{\'e}t{\'e}s des applications $A  \mapsto f(A)$ et $B  \mapsto f^{-1}(B)$}
Soit $f: E\to F$. On consid{\`e}re les applications
$$\Phi: {{\cal P}(E)}\to {{\cal P}(F)}, A \mapsto{f(A)}
\quad \text{ et } \quad
\Psi :{{\cal P}(F)}\to {{\cal P}(E)}, B \mapsto {f^{-1}(B).}$$
Montrer que :
\par\leavevmode\vtop{\halign{\bf#) \hfil&&#\hfil\cr
1&$f$ est injective  &$\iff \Phi$ est injective  &$\iff \Psi$ est surjective. \cr
2&$f$ est surjective &$\iff \Phi$ est surjective &$\iff \Psi$ est injective.  \cr}}
\finenonce{002893}



\finexercice 


\exercice{2894, quercia, 2010/03/08}

\enonce{002894}{${\varphi}  \mapsto {f\circ\varphi}$ et ${\varphi}  \mapsto {\varphi\circ f}$}
Soit $f: E \to F$ une application, et $G$ un troisi{\`e}me ensemble
ayant au moins deux {\'e}l{\'e}ments.
On construit deux nouvelles applications :
$$
{f_*} : {E^G} \to  {F^G},  {\varphi} \mapsto {f\circ\varphi}
     \qquad\text{ et }  \qquad
 {f^*} {G^F} \to {G^E}, {\varphi}\mapsto {\varphi\circ f}
$$
Montrer que :
\begin{enumerate}
  \item $f$ est injective  $\iff f_*$ est injective  $\iff f^*$ est surjective.
  \item $f$ est surjective $\iff f_*$ est surjective $\iff f^*$ est injective.
\end{enumerate}
\finenonce{002894}



\finexercice 


\exercice{2895, quercia, 2010/03/08}

\enonce{002895}{}[$h\circ g\circ f$, $g\circ f\circ h$ injectives et $f\circ
h\circ g$ surjective] Soient $E \xrightarrow{f} F \xrightarrow{g} G
\xrightarrow{h} E$ trois applications telles que $h\circ g\circ f$ et
$g\circ f\circ h$ sont injectives et $f\circ h\circ g$ est surjective.
Montrer que $f,g,h$ sont bijectives.  \finenonce{002895}



\finexercice 


\exercice{2896, quercia, 2010/03/08}

\enonce{002896}{Parties satur{\'e}es pour la relation d'{\'e}quivalence associ{\'e}e {\`a} $f$}

Soit $f:E\to F$ une application, et
${\cal S} = \{ X \subset E$ tq $f^{-1}(f(X)) = X \}$.

\begin{enumerate}
  \item Pour $A \subset E$, montrer que $f^{-1}(f(A)) \in {\cal S}$.
  \item Montrer que ${\cal S}$ est stable par intersection et r{\'e}union.
  \item Soient $X \in {\cal S}$ et $A \subset E$ tels que $X \cap A = \varnothing$.
    Montrer que $X \cap f^{-1}(f(A)) = \varnothing$.
  \item Soient $X$ et $Y \in {\cal S}$.
    Montrer que $\overline X$ et $Y \setminus X$ appartienent {\`a} ${\cal S}$.
  \item Montrer que l'application $ {\cal S} \to {{{\cal P}}(f(E))}, A \mapsto {f(A)}$
    est une bijection.
\end{enumerate}
\finenonce{002896}



\finexercice 


\exercice{2897, quercia, 2010/03/08}

\enonce{002897}{Conjugaison}
Soit $E$ un ensemble et $f : E \to E$ bijective.

La conjugaison par $f$ est l'application
${\Phi_f} : {E^E} \to  {E^E},  \phi \mapsto {f\circ \phi\circ f^{-1}}$
\begin{enumerate}
  \item Montrer que $\Phi_f$ est une bijection de $E^E$.
  \item Simplifier $\Phi_f \circ \Phi_g$.
  \item Simplifier $\Phi_f(\phi) \circ \Phi_f(\psi)$.
  \item Soient $\cal I$, $\cal S$, les sous-ensembles de $E^E$ constitu{\'e}s
    des injections et des surjections.
    Montrer que $\cal I$ et $\cal S$ sont invariants par $\Phi_f$.
  \item Lorsque $\phi$ est bijective, qu'est-ce que $\Bigl(\Phi_f(\phi)\Bigr)^{-1}$ ?
\end{enumerate}
\finenonce{002897}



\finexercice 


\exercice{2898, quercia, 2010/03/08}

\enonce{002898}{Ensembles {\'e}quipotents}
Soient $E,F$ deux ensembles. On dit que :
\begin{tabular}{lll}
$E$ est moins puissant que $F$
         &s'il existe une injection
         &$f:\ E \to F$\\
$E$ est plus puissant que  $F$
         &s'il existe une surjection
         &$f:\ E \to F$\\
$E$ et $F$ sont {\'e}quipotents
         &s'il existe une bijection
         &$f:\ E \to F$.      
    \end{tabular}
\begin{enumerate}
  \item  D{\'e}montrer que : ($E$ est moins puissant que $F$) $\iff$
                     ($F$ est plus puissant que $E$).
  \item  Montrer que $\N$, $\N^*$,
     $\{n \in \N$ tq $n$ est divisible par $3 \}$, et $\Z$ sont
     deux {\`a} deux {\'e}quipotents.
  \item  D{\'e}montrer que $E$ est moins puissant que ${\cal P}(E)$.
  \item  Soit $f : E \to {{\cal P}(E)}$ quelconque et
     $A = \{ x \in E$ tq $x \notin f(x) \}$. Prouver que $A \notin f(E)$.
  \item  Est-ce que $E$ et ${\cal P}(E)$ peuvent {\^e}tre {\'e}quipotents ?
  \item  Soit $G$ un troisi{\`e}me ensemble. Si $E$ est moins puissant que $F$,
     d{\'e}montrer que $E^G$ est moins puissant que~$F^G$.
\end{enumerate}
\finenonce{002898}



\finexercice 


\exercice{2899, quercia, 2010/03/08}

\enonce{002899}{Affirmations}
Soit $f: E\to F$.
Que pensez-vous des affirmations suivantes ?
\begin{enumerate}
  \item $\forall\ x \in E     \quad \forall\ y \in F     \quad f(x)=y.$
  \item $\forall\ x \in E     \quad \exists\ y \in F \text{ tel que } \quad f(x)=y.$
  \item $\exists\ x \in E \text{ tel que } \quad \forall\ y \in F     \quad f(x)=y.$
  \item $\exists\ x \in E \text{ tel que } \quad \exists\ y \in F \text{ tel que } \quad f(x)=y.$
  \item $\forall\ y \in F     \quad \forall\ x \in E     \quad f(x)=y.$
  \item $\forall\ y \in F     \quad \exists\ x \in E \text{ tel que } \quad f(x)=y.$
  \item $\exists\ y \in F \text{ tel que } \quad \forall\ x \in E     \quad f(x)=y.$
  \item $\exists\ y \in F \text{ tel que } \quad \exists\ x \in E \text{ tel que } \quad f(x)=y.$
\end{enumerate}
\finenonce{002899}



\finexercice 

\section{101.02 Injection, surjection}


\exercice{187, gourio, 2001/09/01}

\enonce{000187}{}
Donner des exemples d'applications de $\Rr$ dans $\Rr$ (puis de $\Rr^{2}$ dans $\Rr$)
 injective et non surjective, puis surjective et non injective.
\finenonce{000187}



\finexercice



\exercice{188, bodin, 1998/09/01}

\enonce{000188}{}
Soit $f  : \Rr \rightarrow \Rr$ d\'efinie par $f(x) = x^3-x$.\\
$f$ est-elle injective ? surjective ? D\'eterminer $f^{-1}([-1,1])$ et $f(\Rr_+)$.
\finenonce{000188}



\finexercice



\exercice{189, bodin, 1998/09/01}

\enonce{000189}{}
Les fonctions suivantes sont-elles injectives ? surjectives ? bijectives ?
$$ f : \Zz\rightarrow\Zz, \ n\mapsto 2n \quad ; \quad f : \Zz\rightarrow\Zz ,\ n\mapsto -n $$
$$ f:\Rr\rightarrow\Rr ,\ x\mapsto x^2 \quad ; \quad f : \Rr\rightarrow\Rr_+ ,\ x\mapsto x^2 $$
$$   f : \Cc\rightarrow\Cc ,\ z\mapsto z^2.$$
\finenonce{000189}



\finexercice



\exercice{190, ridde, 1999/11/01}
\video{GVJXQpK7lpY}
\enonce{000190}{}
Les applications suivantes sont-elles injectives, surjectives, bijectives ?
\begin{enumerate}
\item $f : {\Nn} \to {\Nn}, {n} \mapsto {n + 1}$
\item $g : {\Zz} \to {\Zz}, {n}\mapsto{n + 1}$
\item $h : {\Rr^2} \to {\Rr^2}, {(x, y)}\mapsto{ (x + y, x-y)}$
\item $k : {\Rr \setminus \left\{ 1\right\}} \to {\Rr}, {x}\mapsto{\frac{x + 1}{x - 1}}$
\end{enumerate}
\finenonce{000190}



\finexercice


\exercice{191, bodin, 1998/09/01}
\video{n2-T6hM33AM}
\enonce{000191}{}

Soit $f  : \Rr \rightarrow \Rr$ d\'efinie par $f(x) = 2x/(1+x^2)$.
\begin{enumerate}
    \item $f$ est-elle injective ? surjective ?
    \item Montrer que $f(\Rr)=[-1,1]$.
    \item Montrer que la restriction $g  : [-1,1] \rightarrow [-1,1]$  $g(x) = f(x)$
est une bijection.
    \item Retrouver ce r\'esultat en \'etudiant les variations de $f$.
\end{enumerate}

\finenonce{000191} 


\finexercice


\exercice{192, bodin, 1998/09/01}

\enonce{000192}{}
L'application $ f : \Cc\setminus{\{0\}} \rightarrow \Cc, \ z\mapsto z+1/z$ est-elle injective ?
surjective ? bijective ?  \\
Donner l'image par $f$ du cercle de centre $0$ et de rayon $1$. \\
Donner l'image r\'eciproque par $f$ de la droite $i\Rr$.
\finenonce{000192}



\finexercice



\exercice{193, bodin, 1998/09/01}
\video{hBxRKn9zxFs}
\enonce{000193}{}

On consid\`ere quatre ensembles $A,B,C$ et $D$ et des applications $f:A\rightarrow B$, $g:B\rightarrow
C$, $h:C\rightarrow D$. Montrer que :
$$g\circ f\text{ injective } \Rightarrow f\text{ injective,}$$
$$g\circ f\text{ surjective } \Rightarrow g\text{ surjective.}$$
Montrer que :
$$\big(\text{$g\circ f$ et $h\circ g$ sont bijectives }\big) \Leftrightarrow
\big(\text{$f,g$ et $h$ sont bijectives}\big).$$
\finenonce{000193} 


\finexercice


\exercice{194, ridde, 1999/11/01}

\enonce{000194}{}
Soit $f : X \rightarrow Y$. Montrer que
\begin{enumerate}
\item $\forall B \subset Y \, \, f (f^{-1} (B)) = B \cap f (X)$.
\item $f$ est surjective ssi $\forall B \subset Y \, \, f (f^{-1} (B)) = B $.
\item $f$ est injective ssi $\forall A \subset X \, \, f^{-1} (f(A)) = A $.
\item $f$ est bijective ssi $\forall A \subset X \, \,f (\complement A) = \complement f (A).$
\end{enumerate}
\finenonce{000194}



\finexercice



\exercice{195, ridde, 1999/11/01}

\enonce{000195}{}
Soit $f : X \rightarrow Y$. Montrer que les trois propositions suivantes sont
\'equivalentes :
\begin{enumerate}
\item[i. ] $f$ est injective.
\item[ii. ] $\forall A, B \subset X \, \, f (A \cap B) = f (A) \cap f (B)$.
\item[iii. ]$\forall A, B \subset X \, \, A \cap B = \emptyset \Rrightarrow
f (A) \cap f (B) = \emptyset$.
\end{enumerate}
\finenonce{000195}



\finexercice



\exercice{196, ridde, 1999/11/01}

\enonce{000196}{}
Soit $f : X \rightarrow Y$.On note ${\hat{f}} : {\mathcal{P} (X)} \to {\mathcal{P} (Y)},
{A} \mapsto {f (A)}$ et ${\tilde{f}} : {\mathcal{P} (Y)} \to {\mathcal{P} (X)},
{B} \mapsto {f^{-1} (B)}$. Montrer que :
\begin{enumerate}
\item $f$ est injective ssi $\hat{f}$ est injective.
\item $f$ est surjective ssi $\tilde{f}$ est injective.
\end{enumerate}
\finenonce{000196}



\finexercice



\exercice{197, bodin, 1998/09/01}
\video{h9jrsWR1bYw}
\enonce{000197}{Exponentielle complexe}
 Si $z=x+iy$, $(x,y)\in \Rr^2$, on pose $e^z=e^x \times e^{iy}$.
\begin{enumerate}
    \item D\'eterminer le module et l'argument de $e^z$.
    \item Calculer $e^{z+z'}, e^{\overline{z}}, e^{-z}, \left( e^z \right)^n \text{ pour } n \in \Zz$.
    \item L'application $\exp : \Cc \rightarrow \Cc, z \mapsto e^z$, est-elle
injective ?, surjective ?
\end{enumerate}

\finenonce{000197}



\finexercice



\exercice{5110, rouget, 2010/06/30}
\enonce{005110}{*IT}
\label{exo:suprou8}
Montrer que~:~($g\circ f\;\mbox{injective}\Rightarrow f\;\mbox{injective}$) et ($g\circ
f\;\mbox{surjective}\Rightarrow g\;\mbox{surjective}$).
\finenonce{005110}


\finexercice


\exercice{5114, rouget, 2010/06/30}
\enonce{005114}{***IT}
Montrer que les assertions suivantes sont équivalentes ($f$ est une application d'un ensemble $E$ dans lui-même)~:

\begin{enumerate}
\item  $f$ est injective.
\item  $\forall X\in\mathcal{P}(E),\;f^{-1}(f(X))=X$.
\item  $\forall(X,Y)\in\mathcal{P}(E)^2,\;f(X\cap Y)=f(X)\cap f(Y)$.
\item  $\forall(X,Y)\in\mathcal{P}(E)^2,\;X\cap Y=\varnothing\Rightarrow f(X)\cap f(Y)=\varnothing$.
\item  $\forall(X,Y)\in\mathcal{P}(E)^2,\;Y\subset X\Rightarrow f(X\setminus Y)=f(X)\setminus f(Y)$.
\end{enumerate}
\finenonce{005114}


\finexercice

\section{101.03 Bijection}


\exercice{198, bodin, 1998/09/01}

\enonce{000198}{}
Soient $a,\ b \in \Rr$ avec $a \not= 0$, et $f_{a,b} : \Rr \rightarrow \Rr$ telle que
$f_{a,b}(x) = ax+b$. D\'emontrer que $f_{a,b}$ est une permutation et d\'eterminer sa
r\'eciproque.


\finenonce{000198}



\finexercice



\exercice{199, bodin, 1998/09/01}
\video{7avL2IaR9fg}
\enonce{000199}{}
Soit $f: [0,1] \rightarrow [0,1]$ telle que
$$f(x) =    \begin{cases}
            x     & \text{si}\ x\in[0,1]\cap\Qq,\\
            1-x  & \text{sinon.}
        \end{cases}  $$
D\'emontrer que $f \circ f= id$.

\finenonce{000199}



\finexercice



\exercice{200, bodin, 1998/09/01}
\video{j8JoMNuw-vc}
\enonce{000200}{}
 Soit $ f:\Rr \rightarrow \Cc, \ t\mapsto e^{it}$. 
Changer les ensembles de départ et d'arrivée afin que
(la restriction de) $f$ devienne bijective.
\finenonce{000200} 


\finexercice


\exercice{201, bodin, 1998/09/01}

\enonce{000201}{}
On appelle \emph{demi-plan de Poincar\'e} l'ensemble $\mathcal{P}$ des nombres
complexes $z$ tels que $\Im{z}>0$, et \emph{disque unit\'e} l'ensemble $\mathcal{D}$
des nombres complexes $z$ tels que $|z|<1$.
D\'emontrer que $z\mapsto\frac{z-i}{z+i}$ est une bijection de $\mathcal{P}$ sur
$\mathcal{D}$.
\finenonce{000201}



\finexercice



\exercice{202, bodin, 1998/09/01}
\video{lZLGkwlniv4}
\enonce{000202}{}

Soit $f  : [1,+\infty[\rightarrow[0,+\infty[$ telle que
$f(x)=x^2-1$. $f$ est-elle bijective ?
\finenonce{000202} 


\finexercice


\exercice{203, ridde, 1999/11/01}

\enonce{000203}{}
Soient $A \xrightarrow{f} B \xrightarrow{g} C \xrightarrow{h} D$.
Montrer que si $g \circ f$ et $h \circ g$ sont bijectives alors $f, g$ et $h$
le sont \'egalement.
\finenonce{000203}



\finexercice



\exercice{204, ridde, 1999/11/01}

\enonce{000204}{}
Soient $A \xrightarrow{f} B \xrightarrow{g} C \xrightarrow{h} A$.
Montrer que si $h \circ g \circ f$ et $g \circ f \circ h$ sont injectives
et $f \circ h \circ g$ surjective alors $f, g$ et $h$ sont bijectives.
\finenonce{000204}



\finexercice



\exercice{205, ridde, 1999/11/01}

\enonce{000205}{}
Soit $X$ un ensemble. Si $A \subset X$ on note $\chi _A$ la fonction caract\'eristique
associ\'ee. Montrer que ${\Phi} : {\mathcal{P} (X)} \to {\mathcal{F} (X, \left\{ 0, 1 \right\})}$, ${A} \mapsto {\chi _A}$ est bijective.
\finenonce{000205}



\finexercice



\exercice{206, gourio, 2001/09/01}

\enonce{000206}{}
Soit $E$ un ensemble non vide.\ On se donne deux parties $A$ et $B$ de $E$
et on d\'{e}finit l'application $f:\wp (E)\rightarrow \wp (E),$
$X\mapsto (A\cap X)\cup (B\cap X^{c}).$
Discuter et r\'{e}soudre l'\'{e}quation $f(X)=\emptyset$. En d\'{e}duire
une condition n\'ecessaire pour que $f$ soit bijective.

On suppose maintenant $B=A^{c}$. Exprimer $f$ \`{a} l'aide de la
diff\'{e}rence sym\'{e}trique $\Delta$. Montrer que $f$ est bijective,
pr\'{e}ciser $f^{-1}$.  $f$ est-elle involutive (i.e. $f^{2}=id$) ? Quelle
propri\'{e}t\'{e} en d\'{e}duit-on ?
\finenonce{000206}



\finexercice



\exercice{5106, rouget, 2010/06/30}
\enonce{005106}{**IT}
Dans chacun des cas suivants, déterminer $f(I)$ puis vérifier que $f$ réalise une bijection
de $I$ sur $J=f(I)$ puis préciser $f^{-1}$~:

\begin{enumerate}
\item  $f(x)=x^2-4x+3$, $I=]-\infty,2]$.
\item  $f(x)=\frac{2x-1}{x+2}$, $I=]-2,+\infty[$.
\item  $f(x)=\sqrt{2x+3}-1$, $I=\left[-\frac{3}{2},+\infty\right[$.
\item  $f(x)=\frac{x}{1+|x|}$, $I=\Rr$.
\end{enumerate}
\finenonce{005106}


\finexercice


\exercice{5107, rouget, 2010/06/30}
\enonce{005107}{**IT}
Pour $z\neq i$, on pose $f(z)=\frac{z+i}{z-i}$. Montrer que $f$ réalise une bijection
de $D=\{z\in\Cc/\;|z|<1\}$ sur $P=\{z\in\Cc/\;\Re(z)< 0\}$. Préciser $f^{-1}$.
\finenonce{005107}


\finexercice


\exercice{5111, rouget, 2010/06/30}
\enonce{005111}{**T}
Parmi $f\circ g\circ h$, $g\circ h\circ f$ et $h\circ f\circ g$ deux sont injectives et une est surjective. Montrer que
$f$, $g$ et $h$ sont bijectives.
\finenonce{005111}


\finexercice

\exercice{5118, rouget, 2010/06/30}
\enonce{005118}{**** Une bijection entre $\Nn^2$ et $\Nn$}
Soit $\begin{array}[t]{cccc}
f~:&\Nn^2&\rightarrow&\Nn\\
 &(x,y)&\mapsto&y+\frac{(x+y)(x+y+1)}{2}
\end{array}$. Montrer que $f$ est une bijection. Préciser, pour $n\in\Nn$ donné, le couple $(x,y)$ dont il est l'image.
\finenonce{005118}


\finexercice

\section{101.99 Autre}

\section{102.01 Binôme de Newton et combinaison}


\exercice{219, bodin, 1998/09/01}

\enonce{000219}{}
D\'emontrer que si $p$ est un nombre premier, $p$ divise
$C_p^k$ pour $1\le k \le p-1$.
\finenonce{000219}



\finexercice



\exercice{220, bodin, 1998/09/01}
\video{kV-ZtFtGAWI}
\enonce{000220}{}

En utilisant la fonction $x\mapsto(1+x)^n$, calculer :
$$ \sum_{k=0}^{n}C_n^k \quad ; \quad \sum_{k=0}^{n}(-1)^k C_n^k \quad ; \quad
 \sum_{k=1}^{n}kC_n^k \quad ; \quad
 \sum_{k=0}^{n}\frac{1}{k+1}C_n^k.$$

\finenonce{000220} 


\finexercice


\exercice{221, bodin, 1998/09/01}

\enonce{000221}{}
D\'emontrer que $C_n^kC_{n-k}^{p-k}=C_p^kC_n^p$ (pour $0\le k \le p \le n$).
En d\'eduire que
$$ \sum_{k=0}^{n}C_n^kC_{n-k}^{p-k}=2^pC_n^p.$$
\finenonce{000221}



\finexercice



\exercice{222, bodin, 1998/09/01}
\video{GNMFCMiipyM}
\enonce{000222}{}

 En utilisant la formule du bin\^ome, d\'emontrer
que :
\begin{enumerate}
    \item $2^n+1$ est divisible par $3$ si et seulement si $n$ est impair ;
    \item $3^{2n+1}+2^{4n+2}$ est divisible par $7$.
\end{enumerate}
\finenonce{000222} 


\finexercice


\exercice{223, bodin, 1998/09/01}

\enonce{000223}{}
D\'emontrer que $C_n^p=C_{n-1}^p+C_{n-1}^{p-1}$ pour $1\le p \le n-1$.
\finenonce{000223}



\finexercice



\exercice{224, cousquer, 2003/10/01}

\enonce{000224}{}
 Montrer que, pour $p$ et $n$  entiers naturels non nuls tels que 
$1\le p\le n$, on a : $$p C_n^p=nC_{n-1}^{p-1}.$$
\finenonce{000224}



\finexercice



\exercice{225, cousquer, 2003/10/01}

\enonce{000225}{}
\begin{enumerate}
\item Montrer que:
    $$\sum_{k=0}^{p}   C_n^kC_{n-k}^{p-k}= 2^pC_n^p , $$ où 
    $p$ et $n$  sont des entiers naturels avec $0\le p\le n$.
  \item Avec les m\^emes notations, montrer que
    $$\sum_{k=0}^{p} (-1)^{k} C_n^k C_{n-k}^{p-k}= 0.$$
\end{enumerate}
\finenonce{000225}



\finexercice



\exercice{226, cousquer, 2003/10/01}

\enonce{000226}{}
\begin{enumerate}
\item Soient $n$,  $p$ et $q$  des entiers naturels tels que $0\le p,q\le n$.
\item Montrer que l'on a $C_n^p=C_n^q$ si et seulement si $p=q$ ou $p+q=n$.
 \item Résoudre l'équation $$C_{2n+4}^{3n-1}=C_{2n+4}^{n^2-2n+3}.$$
\end{enumerate}
\finenonce{000226}



\finexercice



\exercice{227, bodin, 1998/09/01}

\enonce{000227}{}
Soient $m,n\in\Nn^*$ et $p\in\Nn$. En utilisant la formule du bin\^ome,
d\'emontrer que $m^{2p+1}+n^{2p+1}$ est divisible par $m+n$.
\finenonce{000227}



\finexercice



\exercice{228, bodin, 1998/09/01}

\enonce{000228}{}
 En utilisant la formule du bin\^ome montrer :
$$(a)\ \sum_{k=0}^{n}(-1)^kC_n^k=0
\qquad (b)\ \sum_{k=0}^{n}k^2C_n^k=n(n-1)2^{n-2}+n2^{n-1}. $$

\finenonce{000228}



\finexercice



\exercice{229, cousquer, 2003/10/01}
\enonce{000229}{}

Calculer le module et l'argument de $(1+i)^n$.
En d\'eduire les valeurs de
\begin{eqnarray*}
S_1 & = & 1-C_n^2+C_n^4-C_n^6+\cdots\\
S_2 & = & C_n^1-C_n^3+C_n^5-\cdots
\end{eqnarray*}

\finenonce{000229} 


\finexercice


\exercice{230, bodin, 1998/09/01}

\enonce{000230}{}
D\'emontrer les formules suivantes :
\begin{enumerate}
    \item $C_n^m=C_m^{n-m}$ (on pourra utiliser le fait que $\mathcal{P}(E)
    \longrightarrow \mathcal{P}(E) A\mapsto A^c$ est une bijection.)
    \item $C_n^m=C_{n-1}^m+C_{n-1}^{m-1},$
    \item $C_{n}^{m}=C_{n-2}^{m} +2C_{n-2}^{m-1}+C_{n-2}^{m-2}.$
\end{enumerate}

\finenonce{000230}



\finexercice



\exercice{231, bodin, 1998/09/01}

\enonce{000231}{}
Soient $E$ un ensemble non vide et $X ,Y$ une partition de $E$.
\begin{enumerate}
    \item Montrer que l'application suivante est une bijection :
    $$\mathcal{P}(E)\longrightarrow \mathcal{P}(X)\times \mathcal{P}(Y)$$
    $$A\mapsto(A \cap X, A \cap Y)$$
    \item Montrer que pour $p,q,r\in \Nn$ tel que $r \le p+q$ on a :
    $$\sum_{i+j=r}C_p^iC_q^j=C_{p+q}^r.$$
    \item En d\'eduire que :
    $$\sum_{k=0}^{n}(C_n^k)^2=C_{2n}^n.$$
\end{enumerate}
\finenonce{000231}



\finexercice



\exercice{232, ridde, 1999/11/01}

\enonce{000232}{}
Soit $E$ un ensemble, $a \in E$ et
$f : \begin{cases}
\mathcal{P} (E) \rightarrow \mathcal{P} (E) \\
X \mapsto X \cup \left\{ a\right\} \text{ si } a \notin X\\
X \mapsto X - \left\{ a\right\} \text{ si } a \in X
\end{cases}$
\begin{enumerate}
\item Montrer que $f$ est une bijection.
\item On suppose d\'esormais que $E$ est fini et $\mathrm{Card} (E) = n$. On pose
$\mathcal{P}_0 (E)$ l'ensemble des parties de $E$ de cardinal pair et
$\mathcal{P}_1 (E)$ l'ensemble des parties de $E$ de cardinal impair.
Montrer que $\mathrm{Card} (\mathcal{P}_0 (E)) = \mathrm{Card} (\mathcal{P}_1 (E))$.
\item Calculer ces cardinaux et en d\'eduire la valeur de
$\sum\limits_{k = 0}^n (-1)^k C_n^k$.
\end{enumerate}
\finenonce{000232}



\finexercice



\exercice{233, ridde, 1999/11/01}

\enonce{000233}{}
En utilisant la formule du binôme de Newton, montrer que
$\sum\limits_{k = 0}^n (-1)^k C_n^k = 0$. En d\'eduire la valeur de
$\sum\limits_{0 \leq 2k \leq n} C_n^{2k}$.

\finenonce{000233}



\finexercice



\exercice{234, ridde, 1999/11/01}

\enonce{000234}{}
Soient $0 \leq p \leq n$.
\begin{enumerate}
\item Montrer par r\'ecurrence sur $n$ que $\sum\limits_{k = p}^n C_k^p  =
C_{n + 1}^{p + 1}$.
\item \'Ecrire ces \'egalit\'es pour $p = 2$ et $p = 3$.
\item En d\'eduire les sommes
$$ S_2' = 1. 2 + 2. 3 + \ldots  + (n-1).n
\qquad S_2 = 1^2 + 2^2 + \ldots  + n^2 $$
$$ S_3' = 1^2. 2 + 2^2. 3 + \ldots + (n-1)^2.n
\qquad S_3 = 1^3 + 2^3 + \ldots + n^3 $$
\end{enumerate}
\finenonce{000234}



\finexercice



\exercice{2900, quercia, 2010/03/08}

\enonce{002900}{Calcul de sommes}
Calculer $\sum_{k=0}^n\, kC_n^k$ et $\sum_{k=0}^n \frac{C_n^k}{k+1}$.
\finenonce{002900}



\finexercice


\exercice{2901, quercia, 2010/03/08}

\enonce{002901}{Calcul de sommes}
Soient $n,p \in \N^*$ avec $n \ge p$.
\begin{enumerate}
  \item V{\'e}rifier que $C_n^kC_k^p = C_n^pC_{n-p}^{k-p}$ pour $p \le k \le n$.
  \item Calculer $\sum_{k=0}^n \,(-1)^kC_n^kC_k^p$.    
  \item En d{\'e}duire $\sum_{k=0}^n (-1)^kC_n^kk^p = 0$ si $p < n$.
\end{enumerate}
\finenonce{002901}


\finexercice 


\exercice{2902, quercia, 2010/03/08}

\enonce{002902}{Calcul de sommes}
Soient $n,p \in \N^*$. Simplifier $\sum_{k=0}^p\,(-1)^kC_n^k$.
\finenonce{002902}


\finexercice 


\exercice{2903, quercia, 2010/03/08}

\enonce{002903}{Sommes de cardinaux}
Soit $E$ un ensemble fini de cardinal $n$.
Calculer $\sum_{A \subset E} \mathrm{Card}\,(A)$,
         $\sum_{A,B \subset E} \mathrm{Card}\,(A\cap B)$,
         $\sum_{A,B \subset E} \mathrm{Card}\,(A\cup B)$.
\finenonce{002903}


\finexercice 


\exercice{2904, quercia, 2010/03/08}

\enonce{002904}{Sommes d'entiers}
Soit $n \in \N$. Calculer $\displaystyle\sum_{i+j=n}ij$ et $\displaystyle\sum_{i+j+k=n}ijk$.
\finenonce{002904}


\finexercice 


\exercice{2905, quercia, 2010/03/08}

\enonce{002905}{Combinaisons avec r{\'e}p{\'e}titions}
Soient $n,p \in \N$. On note $\Gamma_n^p$ le nombre de $n$-uplets
$(x_1,\dots, x_n) \in \N^n$ tels que $x_1 + \dots + x_n = p$.
\begin{enumerate}
  \item D{\'e}terminer $\Gamma_n^0$, $\Gamma_n^1$, $\Gamma_n^2$, $\Gamma_2^n$.
  \item D{\'e}montrer que $\Gamma_{n+1}^{p+1} = \Gamma_{n+1}^p + \Gamma_n^{p+1}$
    (on classera les $(n+1)$-uplets tels que $x_1 + \dots + x_{n+1} = p + 1$
    suivant que $x_1 = 0$ ou non).
  \item En d{\'e}duire que $\Gamma_n^p = C_{n+p-1}^p$.
\end{enumerate}
\finenonce{002905}



\finexercice 


\exercice{2906, quercia, 2010/03/08}

\enonce{002906}{Sommes de coefficients du bin{\^o}me}
Soient $n,p \in \N$.
Montrer que $\sum_{k=0}^n C_{p+k}^p = C_{p+n+1}^{p+1}$.
\finenonce{002906}



\finexercice 


\exercice{2907, quercia, 2010/03/08}

\enonce{002907}{$C_n^p$ maximal}
Soit $n \in \N$ fix{\'e}. D{\'e}terminer pour quelle valeur de $p$ le nombre
$C_n^p$ est maximal (on {\'e}tudiera le rapport $C_n^p/C_n^{p+1}$).
\finenonce{002907}


\finexercice 


\exercice{2908, quercia, 2010/03/08}

\enonce{002908}{Parit{\'e} de $C_n^p$}
Soit $p \in \N^*$, et $n = 2^p$.
\begin{enumerate}
  \item Soit $ k \in \{1,\dots, n-1\}$. V{\'e}rifier que $kC_n^k = nC_{n-1}^{k-1}$.
  \item En d{\'e}duire que : $\forall\ k \in \{1,\dots, n-1\}$, $C_n^k$ est pair.
  \item En d{\'e}duire que : $\forall\ k \in \{0,\dots, n-1\}$, $C_{n-1}^k$ est impair.
\end{enumerate}
\finenonce{002908}



\finexercice 


\exercice{2909, quercia, 2010/03/08}

\enonce{002909}{Formule de Vandermonde}
Soient $a,b,c \in \N$.
D{\'e}montrer que $\sum_{k=0}^c\,C_a^kC_b^{c-k} = C_{a+b}^c$ $\ldots$
\begin{enumerate}
  \item En calculant de deux mani{\`e}res $(1+x)^a(1+x)^b$.
  \item En cherchant le nombre de parties de cardinal $c$ dans $E\cup F$, o{\`u}
    $E$ et $F$ sont des ensembles disjoints de cardinaux $a$ et $b$.
  \item Application : Soient $n,p,q \in \N$.
    Montrer que $\sum_{k=0}^q\,C_q^kC_n^{p+k} = C_{n+q}^{p+q}$.
\end{enumerate}
\finenonce{002909}



\finexercice 


\exercice{2910, quercia, 2010/03/08}

\enonce{002910}{Formule d'inversion}
Soit $(x_n)$ une suite de r{\'e}els. On pose $y_n = \sum_{k=0}^n C_n^kx_k$.
Montrer que $(-1)^nx_n = \sum_{k=0}^n \,(-1)^kC_n^ky_k$.
\finenonce{002910}



\finexercice 


\exercice{2911, quercia, 2010/03/08}

\enonce{002911}{Suite de Fibonacci}
Soit $u_n = \sum_{p=0}^n C_{n-p}^p$. Montrer que $u_0 = u_1 = 1$ et :
$\forall\ n \in \N,\ u_{n+2} = u_{n+1} + u_n$
({\it suite de Fibonacci\/}).
\finenonce{002911}



\finexercice 


\exercice{5137, rouget, 2010/06/30}
\enonce{005137}{IT Identités combinatoires}
 \emph{La difficulté va en augmentant graduellement de facile à assez difficile sans être
insurmontable.}

\begin{enumerate}
\item  Calculer $\binom{n}{0}+\binom{n}{1}+...+\binom{n}{n}$.
\item  Montrer que $\binom{n}{0}+\binom{n}{2}+\binom{n}{4}+...=\binom{n}{1}+\binom{n}{3}+\binom{n}{5}+...$ et trouver la valeur commune des
deux sommes.
\item  Calculer les sommes $\binom{n}{0}+\binom{n}{3}+\binom{n}{6}+...$ et $\binom{n}{0}+\binom{n}{4}+\binom{n}{8}+...$.

\item  Montrer que $\forall n\in\Nn^*,\;\forall k\in\llbracket1,n\rrbracket,\;k\binom{n}{k}=n\binom{n-1}{k-1}$.
\item  Montrer que $\binom{n}{0}^2 +\binom{n}{1}^2 + ... +\binom{n}{n}^2 =\binom{2n}{n}$ (utiliser le polynôme $(1+x)^{2n}$).
\item  Calculer les sommes $0.\binom{n}{0}+1.\binom{n}{1}+ ...+n.\binom{n}{n}$ et $\frac{\binom{n}{0}}{1}+\frac{\binom{n}{1}}{2}+...
+\frac{\binom{n}{n}}{n+1}$ (considérer dans chaque cas un certain polynôme astucieusement choisi).
\item  Montrer que $\binom{p}{p}+\binom{p+1}{p}... +\binom{n}{p}=\binom{n+1}{p+1}$ où $0\leq p\leq n$. Interprétation dans le triangle de \textsc{Pascal}~?
\item 
\begin{enumerate}
\item Soit $I_n=\int_{0}^{1}(1-x^2)^n\;dx$. Trouver une relation de récurrence liant $I_n$ et $I_{n+1}$ et en déduire
$I_n$ en fonction de $n$ (faire une intégration par parties dans $I_n-I_{n+1}$).
\item Démontrer l'identité valable pour
$n\geq1$~:~$1-\frac{\binom{n}{1}}{3}+\frac{\binom{n}{2}}{5}+...+(-1)^n\frac{\binom{n}{n}}{2n+1}
=\frac{2.4.....(2n)}{1.3...(2n+1)}$.
\end{enumerate}
\end{enumerate}
\finenonce{005137}


\finexercice


\exercice{5138, rouget, 2010/06/30}
\enonce{005138}{**}
Quel est le coefficient de $a^4b^2c^3$ dans le développement de $(a-b+2c)^9$.
\finenonce{005138}


\finexercice


\exercice{5139, rouget, 2010/06/30}
\enonce{005139}{**I}
Développer $(a+b+c+d)^2$ et $(a+b+c)^3$.
\finenonce{005139}


\finexercice


\exercice{5140, rouget, 2010/06/30}
\enonce{005140}{***}
Soit $(n,a,b)\in\Nn^*\times]0,+\infty[\times]0,+\infty[$. Quel est le plus grand terme du développement de $(a+b)^n$~?
\finenonce{005140}


\finexercice


\exercice{5141, rouget, 2010/06/30}
\enonce{005141}{*}
Résoudre dans $\Nn*$ l'équation $\binom{n}{1}+\binom{n}{2}+\binom{n}{3}= 5n$.
\finenonce{005141}


\finexercice


\exercice{5147, rouget, 2010/06/30}
\enonce{005147}{*I Inégalité de \textsc{Bernoulli}}
Montrer que, pour $a$ réel positif et $n$ entier naturel donnés, $(1+a)^n\geq1+na$.
\finenonce{005147}


\finexercice

\exercice{5158, rouget, 2010/06/30}
\enonce{005158}{****I}
Soit $n\in\Nn^*$.
\begin{enumerate}
\item  Montrer qu'il existe $(a_n,b_n)\in(\Nn^*)^2$ tel que $(2+\sqrt{3})^n=a_n+b_n\sqrt{3}$, puis que
$3b_n^2=a_n^2-1$.
\item  Montrer que $E((2+\sqrt{3})^n)$ est un entier impair (penser à $(2-\sqrt{3})^n)$).
\end{enumerate}
\finenonce{005158}


\finexercice

\exercice{5278, rouget, 2010/07/04}
\enonce{005278}{IT}
\begin{enumerate}
\item (***) Trouver une démonstration combinatoire de l'identité $\sum_{}^{}C_n^{2k}=\sum_{}^{}C_n^{2k+1}$ ou encore démontrer directement qu'un ensemble à $n$ éléments contient autant de parties de cardinal pair que de parties de cardinal impair.
\item (****) Trouver une démonstration combinatoire de l'identité $kC_n^k=nC_{n-1}^{k-1}$.
\item (****) Trouver une démonstration combinatoire de l'identité $C_{2n}^n=\sum_{k=0}^{n}(C_{n}^k)^2$.
\end{enumerate}
\finenonce{005278}


\finexercice


\exercice{5280, rouget, 2010/07/04}
\enonce{005280}{***}
Combinaisons avec répétitions. Montrer que le nombre de solutions en nombres entiers $x_i\geq0$ de l'équation $x_1+x_2+...+x_n=k$ ($k$ entier naturel donné) est $C_{n+k-1}^k$. (Noter $a_{n,k}$ le nombre de solutions et procéder par récurrence.) 
\finenonce{005280}


\finexercice

\section{102.02 Cardinal}


\exercice{235, bodin, 1998/09/01}

\enonce{000235}{}
Montrer que $\Zz$ est d\'enombrable en utilisant l'application :
$$ \phi:\Zz\rightarrow\Nn \begin{cases}
            n \mapsto 2n-1 & \text{ si } n>0 \, ;\\
            n \mapsto -2n  & \text{ sinon}.
        \end{cases}
 $$
\finenonce{000235}



\finexercice



\exercice{236, bodin, 1998/09/01}
\video{zkb9m4IjyfM}
\enonce{000236}{}

 Pour $A,B$ deux ensembles de $E$ on note $A\Delta
B= (A\cup B )\setminus(A\cap B)$. Pour $E$ un ensemble fini,
montrer :
$$\text{Card\,} A\Delta B = \text{Card\,} A + \text{Card\,} B -2\text{Card\,} A\cap B.$$
\finenonce{000236} 


\finexercice


\exercice{237, bodin, 1998/09/01}
\video{YVY_dCYJms8}
\enonce{000237}{}

 Soit $E$ un ensemble \`a $n$ \'el\'ements, et
$A\subset E$ un sous-ensemble \`a $p$ \'el\'ements. Quel est le
nombre de parties de $E$ qui contiennent un et un seul \'el\'ement
de $A$ ?
\finenonce{000237} 


\finexercice


\exercice{238, bodin, 1998/09/01}

\enonce{000238}{}
D\'eterminer le nombre de mots distincts que l'on peut former avec $6$
voyelles et $20$ consonnes, chaque mot \'etant compos\'e de $3$ consonnes
et $2$ voyelles, en excluant les mots qui renferment $3$ consonnes cons\'ecutives.
\finenonce{000238}



\finexercice



\exercice{239, cousquer, 2003/10/01}
\video{eXJ80BfW16o}
\enonce{000239}{}
On considère les mains de $5$~cartes que l'on peut extraire d'un jeu de
$52$~cartes.
\begin{enumerate}
\item Combien y a-t-il de mains différentes~?
\item Combien y a-t-il de mains comprenant exactement un as~?
\item Combien y a-t-il de mains comprenant au moins un valet~?
\item Combien y a-t-il de mains comprenant (à la fois) au moins un roi et au
moins une dame~?
\end{enumerate}
\finenonce{000239}




\finexercice



\exercice{240, bodin, 1998/09/01}

\enonce{000240}{}
Soient $A,A' ,B ,B'$ quatre ensembles tels que :
$$\mathrm{Card}(A)=\mathrm{Card}(A')=a \text{ et } \mathrm{Card}(B)=\mathrm{Card}(B')=b.$$
\begin{enumerate}
    \item D\'eterminer le nombre de bijections de $A \times B $ sur $A'\times B'$.
    \item Supposons maintenant que $\{A,B\},\  \{A' ,B'\}$ forment deux partitions
     de $E$, un ensemble.
    D\'eterminer le nombre de bijections $f :E \longrightarrow E$ telles que $f(A)=A'$
    et $f(B)=B'$.
\end{enumerate}
\finenonce{000240}



\finexercice



\exercice{241, bodin, 1998/09/01}

\enonce{000241}{}
Soient $A$ et $B$ deux sous ensembles finis d'un ensemble $E$.
\begin{enumerate}
    \item Montrer que : $\mathrm{Card}(A \cup B)=\mathrm{Card}(A)+\mathrm{Card}(B)-\mathrm{Card}(A \cap B) $.
    \item Montrer par r\'ecurrence que si $(F_i)_{1\le i\le n}$ est une famille de sous-ensembles finis de E alors :
    $$\mathrm{Card}({\displaystyle{\bigcup_{i=1}^{n}}}F_i)\le \sum_{i=1}^{n}\mathrm{Card}(F_i)$$
    avec \'egalit\'e si les $F_i$ sont deux \`a deux disjoints.
\end{enumerate}
\finenonce{000241}



\finexercice



\exercice{242, ridde, 1999/11/01}

\enonce{000242}{}
Soient $1 \leq k \leq n$. D\'eterminer le nombre de $k$-uplets $ (i_1, \ldots, i_k)$
tels que $1 \leq i_1 < \ldots < i_k \leq n$.
\finenonce{000242}



\finexercice



\exercice{2912, quercia, 2010/03/08}
\video{kXaDLtNT8hI}
\enonce{002912}{Permutations}
Combien y a-t-il de bijections $f$ de $\{1,\dots,12\}$ dans lui-m{\^e}me poss{\'e}dant :
\begin{enumerate}
  \item la propri{\'e}t{\'e} : $n$ est pair $ \Rightarrow $ $f(n)$ est pair ?
  \item la propri{\'e}t{\'e} : $n$ est divisible par 3 $ \Rightarrow $ $f(n)$ est divisible par 3 ?
  \item ces deux propri{\'e}t{\'e}s {\`a} la fois ?
  \item Reprendre les questions pr{\'e}c{\'e}dentes en rempla{\c c}ant {\it bijection\/} par
    {\it application.}
\end{enumerate}
\finenonce{002912}


\finexercice 


\exercice{2913, quercia, 2010/03/08}

\enonce{002913}{Permutations de couples}
On doit placer autour d'une table ronde un groupe de $2n$ personnes, $n$ hommes
et $n$ femmes, qui constituent $n$ couples.
Combien existe-t-il de dispositions $\ldots$
\begin{enumerate}
  \item au total ?                                          
  \item en respectant l'alternance des sexes ?              
  \item sans s{\'e}parer les couples ?                         
  \item en remplissant les deux conditions pr{\'e}c{\'e}dentes ?   
\end{enumerate}
\finenonce{002913}



\finexercice 


\exercice{2914, quercia, 2010/03/08}

\enonce{002914}{Nombre d'op{\'e}rations}
\begin{enumerate}
  \item Combien existe-t-il d'op{\'e}rations internes sur un ensemble {\`a} $n$ {\'e}l{\'e}ments ?
  \item Combien sont commutatives ?
  \item Combien ont un {\'e}l{\'e}ment neutre ?
  \item Combien sont commutatives et ont un {\'e}l{\'e}ment neutre ?
\end{enumerate}
\finenonce{002914}


\finexercice 


\exercice{2915, quercia, 2010/03/08}

\enonce{002915}{Formule du crible}
Soient $A_1, \dots, A_n$\ $n$ ensembles finis.
\begin{enumerate}
\item
  \begin{enumerate}
  \item Calculer $\mathrm{Card}\,(A_1\cup A_2 \cup A_3)$ et $\mathrm{Card}\,(A_1\cup A_2\cup A_3\cup A_4)$.
  \item Sugg{\'e}rer une formule pour $\mathrm{Card}\,(A_1\cup \dots \cup A_n)$.
    \label{devine}
  \end{enumerate}
\item D{\'e}monstration de la formule : On note $E = \bigcup_{i=1}^n A_i$,
  et pour $x \in E$ on pose $f_i(x) = \begin{cases}1 &\text{si }x \in
    A_i \cr 0 &\text{ sinon.}\end{cases}$
  \begin{enumerate}
  \item Soient $x_1,\dots,x_n \in \R$. D{\'e}velopper compl{\`e}tement
    $p = (1-x_1)\times \dots \times (1-x_n)$.
  \item En consid{\'e}rant la somme $\sum_{x \in E}\, (1-f_1(x))  \dots (1-f_n(x))$,
    d{\'e}montrer la formule \ref{devine}.
  \end{enumerate}
\item Applications :
  \begin{enumerate}
     \item D{\'e}terminer le nombre d'applications $f : {\{1,\dots,p\}} \to {\{1,\dots,n\}}$
    non surjectives.
     \item D{\'e}terminer le nombre de permutations d'un ensemble {\`a} $n$ {\'e}l{\'e}ments ayant au moins
    un point fixe.
  \end{enumerate}  
\end{enumerate}
\finenonce{002915}



\finexercice 


\exercice{2916, quercia, 2010/03/08}

\enonce{002916}{In{\'e}galit{\'e}s pour la formule du crible}
Soient $A_1, \dots, A_n$\ $n$ ensembles finis, et $E = \bigcup_{i=1}^n A_i$.
\begin{enumerate}
  \item Montrer que $\mathrm{Card}\,(E) \le \sum_{i=1}^n \mathrm{Card}\,(A_i)$. Cas d'{\'e}galit{\'e} ?
  \item Montrer que $\mathrm{Card}\,(E) \ge \sum_{i=1}^n \mathrm{Card}\,(A_i) -
    \sum_{1\le i< j \le n} \mathrm{Card}\,(A_i\cap A_j)$. Cas d'{\'e}galit{\'e} ?
\end{enumerate}
\finenonce{002916}



\finexercice 


\exercice{2917, quercia, 2010/03/08}

\enonce{002917}{Couples $(A,B)$ tels que  $A\cup B = E$}
Soit $E$ un ensemble fini {\`a} $n$ {\'e}l{\'e}ments, et
${\cal E} = \{ (A,B) \in ({\cal P}(E))^2$ tq $A \cup B = E \}$.
Chercher card(${\cal E}$).
\finenonce{002917}


\finexercice 


\exercice{2918, quercia, 2010/03/08}

\enonce{002918}{Parties ne contenant pas d'{\'e}l{\'e}ments cons{\'e}cutifs}
\begin{enumerate}
\item Quel est le nombre de parties {\`a} $p$ {\'e}l{\'e}ments de $\{1, \dots, n\}$
  ne contenant pas d'{\'e}l{\'e}ments cons{\'e}cutifs ?
\item Soit $t_n$ le nombre de parties de $\{1, \dots, n\}$ de cardinal
  quelconque sans {\'e}l{\'e}ments cons{\'e}cutifs.
  \begin{enumerate}
  \item Montrer que $t_{n+2} = t_{n+1} + t_n$, $t_{2n+1} = t_n^2 +
    t_{n-1}^2$, et $t_{2n} = t_n^2 - t_{n-2}^2$.
  \item Calculer $t_{50}$.
  \end{enumerate}
\end{enumerate}
\finenonce{002918}



\finexercice 


\exercice{2919, quercia, 2010/03/08}

\enonce{002919}{Nombre de relations d'{\'e}quivalence}
Soit $R_n$ le nombre de relations d'{\'e}quivalence sur un ensemble {\`a} $n$ {\'e}l{\'e}ments.
\begin{enumerate}
  \item Trouver une relation de r{\'e}currence entre $R_n$ et les $R_k$, $k < n$
    \par
    (fixer un {\'e}l{\'e}ment, et raisonner sur la classe d'{\'e}quivalence de cet
    {\'e}l{\'e}ment).
  \item Calculer $R_n$ pour $n \le 6$.
\end{enumerate}
\finenonce{002919}


\finexercice 


\exercice{2920, quercia, 2010/03/08}

\enonce{002920}{Equivalence entre fonctions}
Soient $E,F$, deux ensembles non vides. On d{\'e}finit deux relations sur $X = F^E$
par :
$$\begin{matrix}
  f \sim g &\iff &\exists\ \phi : F \to  F \text{ bijective tq } g=\phi\circ f,\hfill\cr
  f \equiv g &\iff &\bigl( \forall\ x,y \in E,\ f(x)=f(y)\iff g(x)=g(y)\bigr).\cr
\end{matrix}$$
\begin{enumerate}
  \item  Montrer que ce sont des relations d'{\'e}quivalence.
  \item  Montrer que $f \sim g  \Rightarrow  f \equiv g$.
  \item  On suppose $f\equiv g$. Montrer que $f \sim g$ dans les cas suivants :
  \begin{enumerate}
\item  $F$ est fini et $f$ est surjective.
\item  $F$ est fini et $f$ est quelconque.
\item  $E$ est fini.
  \end{enumerate}
  \item  Chercher un contrexemple pour $E = F = \N$.
\end{enumerate}
\finenonce{002920}



\finexercice 


\exercice{2921, quercia, 2010/03/08}

\enonce{002921}{Tr{\`e}s bon ordre}
Soit $E$ un ensemble ordonn{\'e} dans lequel toute partie non vide poss{\`e}de un
plus grand et un plus petit {\'e}l{\'e}ment.
\\
Montrer que $E$ est totalement ordonn{\'e} et fini.
\finenonce{002921}



\finexercice 


\exercice{2922, quercia, 2010/03/08}

\enonce{002922}{{\'E}l{\'e}ment maximal}
Soit $E$ un ensemble ordonn{\'e}. Un {\'e}l{\'e}ment $a \in E$ est dit {\it maximal\/}
s'il n'existe pas de $b \in E$ tq $b > a$.
\begin{enumerate}
  \item   Si $E$ est totalement ordonn{\'e}, montrer que :
      {\it maximal\/} $\iff$ {\it maximum.}
  \item   $E = \{1,2,3,4,5,6\}$ ordonn{\'e} par la divisibilit{\'e}. Chercher les
      {\'e}l{\'e}ments maximaux.
  \item   Si $E$ est fini, montrer qu'il existe un {\'e}l{\'e}ment maximal.
  \item   Si $E$ est fini et n'a qu'un seul {\'e}l{\'e}ment maximal, montrer que cet {\'e}l{\'e}ment est
      maximum.
\end{enumerate}
\finenonce{002922}



\finexercice 


\exercice{2923, quercia, 2010/03/08}

\enonce{002923}{Nombres de Catalan}
Soient $x_1,\dots,x_n$ $n$ r{\'e}els. Pour calculer la somme $x_1+\dots+x_n$,
on place des parenth{\`e}ses de fa{\c c}on {\`a} n'avoir que des additions de deux nombres
{\`a} effectuer. Soit $t_n$ le nombre de mani{\`e}res de placer les parenth{\`e}ses
(on pose $t_1 = 1$).
\begin{enumerate}
  \item D{\'e}terminer $t_2,t_3,t_4$.
  \item Trouver une relation de r{\'e}currence entre $t_n$ et $t_1,\dots,t_{n-1}$.
\end{enumerate}
\finenonce{002923}


\finexercice 


\exercice{5279, rouget, 2010/07/04}
\enonce{005279}{***}
Combien y a-t-il de partitions d'un ensemble à $pq$ éléments en $p$ classes ayant chacune $q$ éléments~? (Si $E$ est un ensemble à $pq$ éléments et si $A_1$,..., $A_p$ sont $p$ parties de $E$, $A_1$,..., $A_p$ forment une partition de $E$ si et seulement si tout élément de $E$ est dans une et une seule des parties $A_i$. Il revient au même de dire que la réunion des $A_i$ est $E$ et que les $A_i$ sont deux à deux disjoints.)
\finenonce{005279}


\finexercice

\exercice{5281, rouget, 2010/07/04}
\enonce{005281}{*}
Combien y a-t-il de nombres de $5$ chiffres où $0$ figure une fois et une seule~?
\finenonce{005281}


\finexercice


\exercice{5284, rouget, 2010/07/04}
\video{4EJcG9wA3cI}
\enonce{005284}{**I}

On part du point de coordonnées $(0,0)$ pour rejoindre le point de coordonnées $(p,q)$ ($p$ et $q$ entiers naturels donnés) en se déplaçant à chaque étape d'une unité vers la droite ou vers le haut. Combien y a-t-il de chemins possibles~?
\finenonce{005284}





\finexercice



\exercice{5285, rouget, 2010/07/04}
\enonce{005285}{***I}
De combien de façons peut-on payer $100$ euros avec des pièces de $10$, $20$ et $50$ centimes~?
\finenonce{005285}


\finexercice


\exercice{5286, rouget, 2010/07/04}
\enonce{005286}{****}
\begin{enumerate}
\item  Soit $E$ un ensemble fini et non vide. Soient $n$ un entier naturel non nul et $A_1$,..., $A_n$, $n$ parties de $E$. Montrer la \og~formule du crible~\fg~:

\begin{align*}\ensuremath
\mbox{card}(A_1\cup...\cup A_n)&=\sum_{i=1}^{n}\mbox{card}(A_i)-\sum_{1\leq i_1< i_2\leq n}^{}\mbox{card}(A_{i_1}\cap A_{i_2})\\
 &+...+(-1)^{k-1}\sum_{1\leq i_1<i_2<...<i_k\leq n}^{}\mbox{card}(A_{i_1}\cap A_{i_2}\cap...\cap A_{i_k})\\
 &+...+(-1)^{n-1}\mbox{card}(A_1\cap...\cap A_n).
\end{align*}

\item  Combien y a-t-il de permutations $\sigma$ de $\{1,...,n\}$ vérifiant $\forall i\in\{1,...,n\},\;\sigma(i)\neq i$~?~(Ces permutations sont appelées dérangements (permutations sans point fixe)). Indication~:~noter $A_i$ l'ensemble des permutations qui fixent $i$ et utiliser 1).

On peut alors résoudre un célèbre problème de probabilité, le problème des chapeaux. $n$ personnes laissent leur chapeau à un vestiaire. En repartant, chaque personne reprend un chapeau au hasard. Montrer que la probabilité qu'aucune de ces personnes n'ait repris son propre chapeau est environ $\frac{1}{e}$ quand $n$ est grand.
\end{enumerate}
\finenonce{005286}


\finexercice


\exercice{5287, rouget, 2010/07/04}
\enonce{005287}{**}
Combien y a-t-il de surjections de $\{1,...,n+1\}$ sur $\{1,...,n\}$~?
\finenonce{005287}


\finexercice


\exercice{5288, rouget, 2010/07/04}
\enonce{005288}{***}
Soit $(P)$ un polygone convexe à $n$ sommets. Combien ce polygone a-t-il de diagonales~?~En combien de points distincts des sommets se coupent-elles au maximum~?
\finenonce{005288}


\finexercice


\exercice{5289, rouget, 2010/07/04}
\enonce{005289}{***}
\begin{enumerate}
\item  On donne $n$ droites du plan. On suppose qu'il n'en existe pas deux qui soient parallèles, ni trois qui soient concourantes. Déterminer le nombre $P(n)$ de régions délimitées par ces droites.
\item  On donne $n$ plans de l'espace. On suppose qu'il n'en existe pas deux qui soient parallèles, ni trois qui soient concourants en une droite, ni quatre qui soient concourants en un point. Déterminer le nombre $Q(n)$ de régions délimitées par ces plans.
\end{enumerate}
\finenonce{005289}


\finexercice


\exercice{5290, rouget, 2010/07/04}
\enonce{005290}{***}
Soit $P_n^k$ le nombre de partitions d'un ensemble à $n$ éléments en $k$ classes.

Montrer que $P_n^k=P_{n-1}^{k-1}+kP_{n-1}^k$ pour $2\leq k\leq n-1$.

Dresser un tableau pour $1\leq k,n\leq 5$.

Calculer en fonction de $P_n^k$ le nombre de surjections d'un ensemble à $n$ éléments sur un ensemble à $p$ éléments.
\finenonce{005290}


\finexercice

\section{102.99 Autre}


\exercice{243, bodin, 1998/09/01}

\enonce{000243}{}
\begin{enumerate}
    \item ({\it principe des bergers}) Soient $E ,F$ deux ensembles avec $F$  ensemble fini,
    et $f$ une surjection de $E$ sur $F$ v\'erifiant :
    $$\forall y\in F ,\ \mathrm{Card}(f^{-1}({y}))=p$$
    Montrer que E est alors un ensemble fini et $\mathrm{Card}(E) =p\mathrm{Card}(F)$.
    \item ({\it principe des tiroirs}) Soient
     $\alpha _1,\alpha _2,\ldots,\alpha _p ,$ $p$ \'elements distincts d'un ensemble
    $E$, r\'epartis entre une famille de $n$ sous-ensembles de $E$.
    Si $n<p$ montrer qu'il existe au moins un ensemble de la famille contenant
     au moins deux \'el\'ements parmi les $\alpha _i$.(on pourra raisonner par l'absurde)
\end{enumerate}
\finenonce{000243}



\finexercice



\exercice{244, ridde, 1999/11/01}

\enonce{000244}{}
\label{ex105}
Montrer par r\'ecurrence sur $n$ que si $A_1, \ldots, A_n \subset E$ alors
$\mathrm{Card} (A_1 \cup \ldots \cup A_n) = \sum\limits_{k = 1}^n (-1)^{k + 1}
\sum\limits_{1 \leq i_1 < \ldots < i_k \leq n}\mathrm{Card} (A_{i_1} \cap \ldots \cap A_{i_k})$.
\finenonce{000244}



\finexercice



\exercice{245, gourio, 2001/09/01}

\enonce{000245}{}
Soit $p_{n}(k)$ le nombre de permutations de $\{1,...,n\}$ ayant $k$ points
fixes, montrer alors que :
$$\sum\limits_{k=0}^{n}kp_{n}(k)=n!.$$
Interpr\'{e}ter.
\finenonce{000245}



\finexercice



\exercice{246, gourio, 2001/09/01}

\enonce{000246}{}
Soit $E$ un ensemble de cardinal $nm\in \Nn^{*}$, o\`{u} $(n,m)\in (\Nn^{*})^{2}$,
 et $P_{n,m} $l'ensemble des partitions de $E$ en $n$ parties \`{a} $m $
\'{e}l\'{e}ments chacune. Montrer que :
$$N_{n,m}=card(P_{n,m})=\frac{(nm)!}{n!(m!)^{n}}.$$
(\emph{Indication} : on peut proc\'{e}der par r\'{e}currence.)
\finenonce{000246}



\finexercice



\exercice{247, ridde, 1999/11/01}

\enonce{000247}{}
L'histoire: $n$ personnes apportent chacune un cadeau \`{a} une f\^{e}te, et
chacun tire au sort un cadeau dans le tas form\'{e} par tous les
pr\'{e}sents apport\'{e}s.\ Quelle est la probabilit\'{e} qu'au moins une
personne reparte avec son cadeau? Que devient cette probabilit\'{e} quand le
nombre de personnes devient tr\`{e}s grand, i.e.: $n\rightarrow \infty $?
(On remarquera que l'intuition met en \'{e}vidence deux effets
contradictoires: plus de personnes c'est plus de proba qu'une personne ait
son cadeau car... il y a plus de personnes, mais c'est aussi plus de
cadeaux, donc une proportion plus \'{e}lev\'{e}e de cadeaux ``acceptables'').


Soit $S_n  = \sigma (\left\{ 1, \ldots, n\right\})$. On dit que $\sigma
\in S_n$ est un d\'erangement si $\forall i \in \left\{ 1, \ldots, n\right\}
\, \, \sigma (i) \neq i$. On note $A_i = \left\{ \sigma \in S_n /
\sigma (i) = i\right\}$ et $D_n$ l'ensemble des d\'erangements.
\begin{enumerate}
\item Calculer $\mathrm{Card} (A_i)$.
\item Exprimer $S_n - D_n$ en fonction des $A_i$.
\item En d\'eduire $\mathrm{Card} (D_n)$ (on pourra utiliser l'exercice \ref{ex105}).
\item D\'eterminer la limite de $\dfrac{\mathrm{Card}{D_n}}{\mathrm{Card}{S_n}}$. (on rappelle
que $\lim\limits_{n \rightarrow  + \infty} (1 + x + \ldots + \frac{x^n}{n ! }) = e^x$).
\end{enumerate}
\finenonce{000247}



\finexercice



\exercice{248, gourio, 2001/09/01}

\enonce{000248}{}
Soit $ E$ un ensemble de cardinal $ n, \Re  $ une relation d'\'{e}quivalence
sur $ E$, avec $ k$ classes d'\'{e}quivalences et $ r  $ couples $ \left(
x,y\right) \in E^{2} $ tels que $ x\Re y.$
Montrer que $ n^{2}\leq kr.$
\finenonce{000248}



\finexercice



\exercice{3051, quercia, 2010/03/08}

\enonce{003051}{D{\'e}nombrement de $\N^2$}
Soit 
\begin{align*}
  f : {\N^2} &\to {\N}, \\ {(p,q)} &\mapsto {{\frac12}(p+q)(p+q+1) + p}.
\end{align*}
 \begin{enumerate}
  \item  Montrer pour $q > 0$ : $f(p+1,q-1) = f(p,q)+1$ et $f(0,p+1) = f(p,0)+1$.
  \item  Montrer que : $f(0,p+q) \le f(p,q) < f(0,p+q+1)$.
  \item  Montrer que $g$ : $n \mapsto f(0,n)$ est strictement croissante.
  \item  Montrer que $f$ est injective (on supposera $f(p,q) = f(p',q')$ et on montrera
     dans un premier temps que $p+q = p'+q'$).
  \item  Montrer que $f$ est surjective.
\end{enumerate}
\finenonce{003051}



\finexercice 


\exercice{3052, quercia, 2010/03/08}

\enonce{003052}{Parties d{\'e}nombrables}
Soit $(n_k)$ une suite d'entiers naturels. On dit que la suite est :

\indent\vbox{\halign{&#\hfil\cr
- presque nulle &s'il existe $p \in \N$ tq $\forall\ k \ge p,\ n_k = 0$\cr
- stationnaire  &s'il existe $p \in \N$ tq $\forall\ k \ge p,\ n_k = n_p$.\cr
}}

Montrer que les ensembles des suites presque nulles et des suites stationnaires sont
d{\'e}nombrables.
\finenonce{003052}



\finexercice


\exercice{3053, quercia, 2010/03/08}

\enonce{003053}{Propri{\'e}t{\'e}s du pgcd et du ppcm}
Soient $a,b \in \N$. On pose $m = \text{ppcm}(a,b)$ et
$d = \text{pgcd}(a,b)$.
\begin{enumerate}
  \item  Soit $x$ un multiple commun {\`a} $a$ et $b$. En {\'e}crivant la division
     euclidienne de $x$ par $m$, montrer que $m \mid x$.
  \item  Soit $x$ un diviseur commun {\`a} $a$ et $b$. Montrer que $\text{ppcm}(x,d)$ est
     aussi un diviseur commun {\`a} $a$ et $b$. En d{\'e}duire $x \mid d$.
  \item  Comment qualifier $m$ et $d$ pour la relation d'ordre de divisibilit{\'e} ?
\end{enumerate}
\finenonce{003053}



\finexercice


\exercice{3054, quercia, 2010/03/08}

\enonce{003054}{Bases de num{\'e}ration}
Soit $b \in \N\setminus \{0,1\}$ et $p \in \N$.
Montrer que pour tout entier $n \in \{0,\dots,b^p-1\}$, il existe
un unique $p$-uplet $(n_0,\dots,n_{p-1})$ d'entiers naturels tel que :
$$\forall\ k < p,\ n_k \in \{0,\dots,b-1\},\quad\text{et}\quad
  n = \sum_{k=0}^{p-1} n_kb^k.$$
\finenonce{003054}



\finexercice 


\exercice{3055, quercia, 2010/03/08}

\enonce{003055}{Bases de num{\'e}ration}
Soit $n \in \N^*$. Montrer qu'il existe $p \in \N$ et $n_0,n_1,\dots,n_p
\in \{1,2\}$ uniques tels que $n = \sum_{k=0}^p n_k2^k$.
\finenonce{003055}



\finexercice 


\exercice{3056, quercia, 2010/03/08}

\enonce{003056}{Bases de num{\'e}ration}
Soient $n,p \in \N^*$ avec $n < p!$. Montrer qu'il existe
un unique $p$-uplet $(n_1,\dots,n_p)$ d'entiers naturels tel que
$$\forall\ k \le p,\ n_k \le k,\quad\text{et}\quad
  n = \sum_{k=1}^p n_kk!.$$
\finenonce{003056}



\finexercice 


\exercice{3057, quercia, 2010/03/08}

\enonce{003057}{R{\'e}currence d'ordre 2}
On note $a_n = 25^n + 2^{3n+4}$.
\begin{enumerate}
  \item  Trouver $a,b\in\Z$ tels que :
     $\forall\ n\in\N,\ a_{n+2} = a.a_{n+1} + b.a_n$.

  \item  En d{\'e}duire que : $\forall\ n \in \N,\ a_n$ est divisible par $17$.
\end{enumerate}
\finenonce{003057}



\finexercice 


\exercice{3058, quercia, 2010/03/08}

\enonce{003058}{Ordre sur $\N^{\N}$}
Soit $E = {\N}^{\N}$. Pour $f,g \in E$ avec $f\ne g$,
on note $n_{f,g} = \min\{k$ tq $f(k) \ne g(k)\}$.

On ordonne $E$ par :
$$\forall\ f,g\in E,\
  f \ll g \iff (f = g) \text{ ou }\bigl(f(n_{f,g}) < g(n_{f,g})\bigr).$$
\begin{enumerate}
  \item  Montrer que c'est une relation d'ordre total.
  \item  Montrer que toute partie de $E$ non vide admet une borne inf{\'e}rieure
      et toute partie de $E$ non vide et major{\'e}e admet une borne sup{\'e}rieure.
\end{enumerate}
\finenonce{003058}



\finexercice 


\exercice{3059, quercia, 2010/03/08}

\enonce{003059}{$f\circ f(n) = n+k$}
On veut montrer qu'il n'existe pas d'application $f: {\N} \to {\N}$
v{\'e}rifiant : $\forall\ n \in \N,\ f(f(n)) = n + 1987$.\\
({\it Olympiades 1987})

Soit $f$ une telle application. On pose :
$$E = \{0,\dots,1986\},  \qquad F = \N\setminus E,  \qquad G = f(\N) \cap E,
  \qquad H = E\setminus G.$$

D{\'e}montrer successivement :
\begin{enumerate}
  \item  $f$ est injective,
  \item  $f(F) \subset F$,
  \item  $f^{-1}(F) = F \cup G$,
  \item  $f^{-1}(G) = H$,
\end{enumerate}
puis obtenir une contradiction.
\finenonce{003059}



\finexercice 


\exercice{3060, quercia, 2010/03/08}

\enonce{003060}{$f(f(n)) < f(n+1)$}
Soit $f : {\N} \to {\N}$ telle que :
$\forall\ n \in \N,\ f(f(n)) < f(n+1)$.
On veut montrer que $f = \text{id}_{\N}$.
\par {\it (Olympiades 1977)}
\begin{enumerate}
  \item Montrer que $\forall\ n\in \N,\ \forall\ x \ge n,\ f(x) \ge n$.
  \item Soit $n \in \N$ et $a \ge n$ tel que $f(a) = \min\{f(x)$ tq $x\ge n\}$.
    Montrer que $a=n$.
  \item En d{\'e}duire que $f$ est strictement croissante, puis conclure.
\end{enumerate}
\finenonce{003060}



\finexercice 


\exercice{5282, rouget, 2010/07/04}
\enonce{005282}{***I}
Quelle est la probabilité $p_n$ pour que dans un groupe de $n$ personnes choisies au hasard, deux personnes au moins aient le même anniversaire (on considèrera que l'année a toujours $365$ jours, tous équiprobables). Montrer que pour $n\geq23$, on a $p_n\geq\frac{1}{2}$.
\finenonce{005282}


\finexercice


\exercice{5283, rouget, 2010/07/04}
\enonce{005283}{***}
Montrer que le premier de l'an tombe plus souvent un dimanche qu'un samedi.
\finenonce{005283}


\finexercice

\section{103.01 Divisibilité, division euclidienne}


\exercice{249, bodin, 1998/09/01}
\video{akuNAnpeHZM}
\enonce{000249}{}
Combien $15!$ admet-il de diviseurs ?

\finenonce{000249} 


\finexercice


\exercice{250, bodin, 1998/09/01}
\video{6MywGfpcPic}
\enonce{000250}{}

Trouver le reste de la division par $13$ du nombre $100^{1000}$.
\finenonce{000250} 


\finexercice


\exercice{251, bodin, 1998/09/01}
\video{NL9J5okSGCw}
\enonce{000251}{}

Sachant que l'on a $96842=256\times 375 + 842$, d\'eterminer, sans faire
la division, le reste de la division du nombre $96842$ par chacun des nombres
$256$ et $375$.
\finenonce{000251} 


\finexercice


\exercice{252, bodin, 1998/09/01}

\enonce{000252}{}
Soient $m\ge 1$ et $n\ge 2$ des entiers ; montrer que :
\begin{enumerate}
    \item $n-1 |n^m-1$ ;
    \item $(n-1)^2|n^m-1 \text{ si et seulement si } n-1|m$.
\end{enumerate}
\finenonce{000252}



\finexercice



\exercice{253, bodin, 1998/09/01}

\enonce{000253}{}
Soit $a$  un entier relatif quelconque, d\'emontrer que le nombre
$a(a^2-1)$ et, plus g\'en\'eralement, $a(a^{2n}-1)$ est divisible par $6$.
\finenonce{000253}



\finexercice



\exercice{254, bodin, 1998/09/01}
\video{I_npnj8Lb0M}
\enonce{000254}{}
 D\'emontrer que le nombre $7^n+1$ est divisible par
$8$ si $n$ est impair ; dans le cas $n$ pair, donner le reste de
sa division par $8$.

\finenonce{000254} 


\finexercice


\exercice{255, bodin, 1998/09/01}

\enonce{000255}{}
Quel est le plus petit entier naturel qui, divis\'e par $8,15,18$ et $24$, donne
respectivement pour reste $7,14,17$ et $23$ ?
\finenonce{000255}



\finexercice



\exercice{256, bodin, 1998/09/01}

\enonce{000256}{}
Montrer que si $x$ et $y$ sont des entiers naturels tels que $x^2$
divise $y^2$, alors $x$ divise $y$.  Application : d\'emontrer, par
l'absurde, que $\sqrt{2}$ n'est pas rationnel.
\finenonce{000256}



\finexercice



\exercice{257, bodin, 1998/09/01}
\video{wEt81lOn-LA}
\enonce{000257}{}

 Montrer que $\forall n\in \Nn$ :
$$n(n+1)(n+2)(n+3)  \text {  est divisible par }24,$$
$$n(n+1)(n+2)(n+3)(n+4) \text{  est divisible par }120.$$
\finenonce{000257} 


\finexercice


\exercice{258, bodin, 1998/09/01}

\enonce{000258}{}
Trouver tous les entiers relatifs $n$ tels que $n^2+n+7$ soit divisible par $13$.
\finenonce{000258}



\finexercice



\exercice{259, bodin, 1998/09/01}

\enonce{000259}{}
On consid\`ere le nombre $m = 2^np$, dans lequel $n$ d\'esigne
un entier naturel quelconque et $p$ un nombre premier. Dresser
la liste des diviseurs de $m$, y compris $1$ et $m$ lui-m\^eme,
et calculer, en fonction de $m$ et $p$, la somme $S$ de tous
ces diviseurs.
\finenonce{000259}



\finexercice



\exercice{260, bodin, 1998/09/01}

\enonce{000260}{}
Le diviseur d'une division est \'egal \`a $45$ ; le reste est le carr\'e du quotient.
Calculer le dividende entier naturel.
\finenonce{000260}



\finexercice



\exercice{261, bodin, 1998/09/01}

\enonce{000261}{}
Trouver le plus petit entier naturel $n$ telle que
le d\'eveloppement d\'ecimal de $1/n$ admette une plus petite p\'eriode de
longueur $5$,
c'est-\`a-dire $1/n = 0,abcde\, abcde \,ab\ldots$ avec $a,b,\ldots,e \in \{0,1,2,\ldots,9\}$.
\finenonce{000261}



\finexercice



\exercice{262, cousquer, 2003/10/01}

\enonce{000262}{}
Les nombres $a$, $b$, $c$, $d$ étant des éléments non nuls de $\mathbb{Z}$, 
dire si les propriétés suivantes sont vraies ou 
fausses, en justifiant la réponse.
\begin{enumerate}
    \item Si $a$ divise $b$ et $c$, alors $c^2-2b$ est multiple de~$a$.
    \item S'il existe $u$ et $v$ entiers tels que $au+bv=d$ alors 
$\mbox{pgcd}(a,b)=\vert d\vert$.
    \item Si $a$ est premier avec $b$, alors $a$ est premier avec $b^3$.
    \item Si $a$ divise $b+c$ et $b-c$, alors $a$ divise $b$ et $a$ divise 
$c$.
    \item Si $19$ divise $ab$, alors $19$ divise $a$ ou $19$ divise $b$.
    \item Si $a$ est multiple de~$b$ et si $c$ est multiple de~$d$, alors 
$a+c$ est multiple de $b+d$.
    \item Si $4$ ne divise pas $bc$, alors $b$ ou $c$ est impair.
    \item Si $a$ divise $b$ et $b$ ne divise pas $c$, alors $a$ ne divise 
pas $c$.
    \item Si $5$ divise $b^2$, alors $25$ divise $b^2$.
    \item Si $12$ divise $b^2$, alors $4$ divise $b$.
    \item Si $12$ divise $b^2$, alors $36$ divise $b^2$.
    \item Si $91$ divise $ab$, alors $91$ divise $a$ ou $91$ divise $b$.
\end{enumerate}
\finenonce{000262}



\finexercice



\exercice{263, cousquer, 2003/10/01}

\enonce{000263}{}
On définit les trois ensembles suivants :
\begin{eqnarray*}
E_1 & = & \left\{7n\,,\; n\in\mathbb{N}\right\}\\
E_2 & = & \left\{n\in\mathbb{N}\,\mbox{ tel que }\; n \mbox{ est multiple de }
4\right\}\\
E_3 & = & \left\{28n\,,\; n\in\mathbb{N}\right\}
\end{eqnarray*}
\begin{enumerate}
\item  Pour $1\leq i,j\leq 3$, déterminer si on a l'inclusion
$E_i\subset E_j$.
\item  Ecrire $E_1\cap E_2$ sous la forme $E=\left\{n\in\mathbb{N}\,,\;
\mathcal{P}(n)\right\}$. Montrer que $E_1\cap E_2=E_3$. 
\end{enumerate}
\finenonce{000263}



\finexercice



\exercice{264, cousquer, 2003/10/01}

\enonce{000264}{}
 Montrer que si $r$ et $s$ sont deux nombres entiers naturels somme de deux carrés d'entiers alors il en est de même pour le produit
$rs$.
\finenonce{000264}



\finexercice



\exercice{265, cousquer, 2003/10/01}

\enonce{000265}{}
Soit $n$ un entier relatif. Montrer que soit $8$ divise $n^2$, soit $8$ divise
$n^2-1$, soit $8$ divise $n^2-4$.
\finenonce{000265}



\finexercice



\exercice{266, cousquer, 2003/10/01}

\enonce{000266}{}
 Étant donnés deux nombres relatifs $n$ et $p$ montrer
que soit $np$ est pair, soit $n^2-p^2$ est divisible par $8$.
\finenonce{000266}



\finexercice



\exercice{267, cousquer, 2003/10/01}
\video{0ClhoFc-5jQ}
\enonce{000267}{}

 Montrer que si $n$ est un entier naturel somme de deux carr\'es d'entiers 
alors le reste de la division euclidienne de $n$ par $4$ n'est jamais \'egal \`a $3$.

\finenonce{000267} 


\finexercice


\exercice{268, cousquer, 2003/10/01}

\enonce{000268}{}
\begin{enumerate}
\item  Soit $n$ un entier naturel dont le reste de la division euclidienne par 5 vaut
$2$ ou $3$, montrer que $n^2 +1$ est divisible par $5$.
\item  Montrer que pour tout entier naturel $n$, l'entier $n^5-n$ est
divisible par $5$.
\end{enumerate}
\finenonce{000268}



\finexercice



\exercice{269, cousquer, 2003/10/01}

\enonce{000269}{}
Soit $n\in\mathbb{N}^*$. Montrer que parmi les trois entiers $n.(n+1)$,
$n.(n+2)$ et $(n+1).(n+2)$, il y en a exactement deux qui sont divisibles par $3$. 
\finenonce{000269}



\finexercice



\exercice{270, cousquer, 2003/10/01}
\enonce{000270}{}

\begin{enumerate}
\item  Pour tout couple de nombres r\'eels $(x,y)$ montrer, par r\'ecurrence, que
pour tout $n \in \mathbb{N}^{*}$ on a la relation
$$(*)\; x^{n}-y^{n} = (x-y).\sum_{k=0}^{n-1}x^{k}y^{n-1-k}.$$
Indication: on pourra \'ecrire de deux mani\`eres diff\'erentes la quantit\'e
$y(x^{n}-y^{n}) + (x-y)x^{n}.$
\item  \label{div} Soit $(a,b,p)$ des entiers \'el\'ements de $\mathbb{N}$. En utilisant
la formule $(*),$ montrer que s'il existe un entier $l \in \mathbb{N}$ tel que $b = a+pl,$
alors pour tout $n \in \mathbb{N}^{*},$ il existe un entier $m \in \mathbb{N}$ tel que $b^{n} =
a^{n}+pm.$
\item  \label{divv} Soient $a,b,p$ des entiers \'el\'ements de $\mathbb{N}$, en
utilisant la question \ref{div}, montrer que si $a-b$ est divisible par $p,$
$$\sum_{k=0}^{p-1}a^{k}b^{p-k-1}$$
est aussi divisible par $p.$
En d\'eduire, \`a l'aide de la question \ref{div} et de la formule $(*),$ que si $a-b$
est divisible par $p^{n}$ i.e. il existe un entier $l \in \mathbb{N}$ tel que $a-b =
l.p^{n},$ alors $a^{p}-b^{p}$ est divisible par $p^{n+1}.$
\end{enumerate}
\finenonce{000270} 


\finexercice


\exercice{271, cousquer, 2003/10/01}

\enonce{000271}{}
 Calculer $2000^{2000}$ modulo $7$ et $2^{500}$ modulo $3$. 
\finenonce{000271}



\finexercice



\exercice{272, cousquer, 2003/10/01}

\enonce{000272}{}
 Soit $a,b \in {\mathbb{Z}}^2$  dont les restes modulo $11$
sont $7$ et $2$ respectivement. Donner le reste modulo $11$
de $a^2-b^2$.
\finenonce{000272}



\finexercice



\exercice{273, cousquer, 2003/10/01}

\enonce{000273}{}
\begin{enumerate}
  \item Montrer que $7$ divise $2222^{5555}+5555^{2222}$;
  \item montrer que que $11$ divise 
$$5^{10^{5^{10^{5^{10}}}}}+10^{5^{10^{5^{10^5}}}};$$ 
  \item trouver un critère de divisibilité par $8$ puis par $6$.
\end{enumerate}
\finenonce{000273}



\finexercice



\exercice{274, cousquer, 2003/10/01}

\enonce{000274}{}
 Montrer que  pour  tout $n> 0$ :
\begin{enumerate}
  \item  $7 \hbox{ divise } 3^{2n+1}+2^{n+2}$
  \item   $11\hbox{ divise }2^{6n+3} +3^{2n+1}$
  \item   $6\hbox{ divise }5n^3+n$
  \item   $8 \hbox{ divise } 5^n +2.3^{n-1} +1$ .
\end{enumerate}
\finenonce{000274}



\finexercice



\exercice{275, cousquer, 2003/10/01}

\enonce{000275}{}
\begin{enumerate}
\item  Déterminer la somme des chiffres 
de la somme des chiffres de la somme des chiffres de $3^{500}$.
\item  On se donne $51$ nombres compris entre $1$ et $100$. Montrer que parmi 
ces nombres il y en a nécessairement au moins deux tels que l'un divise l'autre.
Montrer que l'on peut toujours trouver un ensemble de $50$ nombres compris entre 
entre $1$ et $100$ ne vérifiant pas la propriété de divisibilité 
ci-dessus. 
\end{enumerate}
\finenonce{000275}



\finexercice



\exercice{276, cousquer, 2003/10/01}

\enonce{000276}{}
Trouver les entiers positifs $n$ tels que $n-1$ divise $n^2+1$.
\finenonce{000276}



\finexercice



\exercice{277, cousquer, 2003/10/01}

\enonce{000277}{}
 Montrer que pour chaque $n \in \mathbb{N}$, $4$ ne divise pas $n^2+1$.
\finenonce{000277}



\finexercice



\exercice{278, cousquer, 2003/10/01}

\enonce{000278}{}
 Montrer que pour chaque entier positif $n$, $49$ divise $2^{3n+3}-7n-8$.
\finenonce{000278}



\finexercice



\exercice{279, cousquer, 2003/10/01}

\enonce{000279}{}
 Trouver tous les entiers positifs $a$ tels que $a^{10}+1$ est 
divisible par $10$.
\finenonce{000279}



\finexercice



\exercice{280, cousquer, 2003/10/01}

\enonce{000280}{}
 Quel est le chiffre des unités de $19971997^{10}$~?
\finenonce{000280}



\finexercice



\exercice{281, cousquer, 2003/10/01}

\enonce{000281}{}
 Montrer que~:
\begin{enumerate}

\item  Si un entier est de la forme $6k+5$, alors il est 
    nécessairement de la forme $3k-1$, alors que la réciproque est 
    fausse.
\item Le carré d'un entier de la forme $5k+1$ est aussi de cette 
    forme.
\item Le carré d'un entier est de la forme $3k$ ou $3k+1$, mais
    jamais de la forme $3k+2$.
\item Le carré d'un entier est de la forme $4k$ ou $4k+1$, mais jamais 
de la forme $4k+2$ ni de la forme $4k+3$.
\item Le cube de tout entier est de la forme $9k$, $9k+1$ ou $9k+8$.    
\item Si un entier est à la fois un carré et un cube, alors c'est une 
puissance sixième, et il est de la forme $7k$ ou $7k+1$.
\end{enumerate}
\finenonce{000281}



\finexercice



\exercice{282, ridde, 1999/11/01}

\enonce{000282}{}
D\'eterminer les entiers $n \in \Nn$ tels que :
\begin{enumerate}
\item $n|n + 8$.
\item $n-1|n + 11$.
\item $n-3|n^3-3$.
\end{enumerate}
\finenonce{000282}



\finexercice



\exercice{283, ridde, 1999/11/01}

\enonce{000283}{}
Soit $k \in \Zz$. D\'eterminer les entiers $n \in \Nn^*$ tels que
$ (n|2k + 1 \text{ et } n|9k + 4)$.
\finenonce{000283}



\finexercice



\exercice{284, ridde, 1999/11/01}

\enonce{000284}{}
Montrer que $\forall (a, b) \in \Nn \times \Nn^*$ il existe un unique $r(a)
\in \left\{ 0, \ldots, b-1\right\}$ tel qu'il existe $q \in \Nn$ avec
$a = bq + r (a)$.
\begin{enumerate}
\item En utilisant ceci pour $b = 13$, d\'eterminer les entiers $n\in \Nn$ tels
que $13|n^2 + n + 7$.
\item Si $a \in \Nn$ et $b = 7$, d\'eterminer les valeurs possibles de $r (a^2)$
 (on rappelle que $r (a^2)$ doit appartenir \`a $\left\{ 0, \ldots, b-1\right\}$).\\
Montrer alors que $\forall (x, y) \in \Nn^2$ $ (7|x^2 + y^2) \text{ ssi }(7|x \text{ et }
7|y)$.
\item Montrer qu'un entier positif de la forme $8k + 7$ ne peut pas être la
somme de trois carr\'es d'entiers.
\end{enumerate}
\finenonce{000284}



\finexercice



\exercice{285, gourio, 2001/09/01}
\video{goNvkUJGM88}
\enonce{000285}{}
\begin{enumerate}
\item Montrer que le reste de la division euclidienne par $8$ du carr\'{e} de tout
nombre impair est $1$.
\item Montrer de m\^{e}me que tout nombre pair v\'{e}rifie $x^{2}=0 \pmod{8} $ ou
$x^{2}=4 \pmod{8}.$
\item Soient $a,b,c$ trois entiers impairs. D\'{e}terminer le reste modulo $8$ de
$a^{2}+b^{2}+c^{2}$ et celui de $2(ab+bc+ca).$
\item En d\'{e}duire que ces deux nombres ne sont pas des carr\'{e}s puis que
$ab+bc+ca$ non plus.
\end{enumerate}

\finenonce{000285} 


\finexercice


\exercice{3090, quercia, 2010/03/08}

\enonce{003090}{Sommes de nombres impairs}
Soit $n\in\N$, $n\ge 2$. Montrer que si $N$ est la somme de $n$ nombres impairs
cons{\'e}cutifs, alors $N$ n'est pas premier.
\finenonce{003090}



\finexercice


\exercice{3091, quercia, 2010/03/08}

\enonce{003091}{Petit th{\'e}or{\`e}me de Fermat}
Soit $p \in \N$ premier. Montrer que pour $1 \le k \le p-1$, $p$ divise $C_p^k$.

En d{\'e}duire que $\forall\ n \in \Z$, $n^p \equiv n (\mathrm{mod}\, p)$.
\label{Fermat}
\finenonce{003091}



\finexercice 


\exercice{3092, quercia, 2010/03/08}

\enonce{003092}{$(p-1)(p-2)...(p-n)/n!$}
Soit $p \in \N^*$ premier et $n \in \N^*$, $n < p$.
Montrer que $\frac {(p-1)(p-2) \dots (p-n)}{n!} - (-1)^n$ est  un entier
divisible par~$p$.
\finenonce{003092}



\finexercice 


\exercice{3093, quercia, 2010/03/08}

\enonce{003093}{$n^7 \equiv n (\mathrm{mod}\, 42)$}
Montrer que : $\forall\ n \in \Z$, $n^7 \equiv n (\mathrm{mod}\, {42})$.
\finenonce{003093}



\finexercice 


\exercice{3094, quercia, 2010/03/08}

\enonce{003094}{Puissances de 10 modulo 7}

\begin{enumerate}
  \item V{\'e}rifier $10^6 \equiv 1 (\mathrm{mod}\, 7)$.
  \item Montrer que $\sum_{k=1}^{10} 10^{10^k} \equiv 5 (\mathrm{mod}\, 7)$.
\end{enumerate}
\finenonce{003094}



\finexercice 


\exercice{3095, quercia, 2010/03/08}

\enonce{003095}{Puissances de 7}
Quel est le dernier chiffre de $7^{7^{7^{7^{7^{7^7}}}}}$ ?
\finenonce{003095}


\finexercice 


\exercice{3096, quercia, 2010/03/08}

\enonce{003096}{$3^x = 2^y + 1$}
\begin{enumerate}
  \item Soient $x,y\in \N$, $y \ge 3$.
    Montrer par r{\'e}\-cur\-rence sur $y$ que :
    $3^x \equiv 1 (\mathrm{mod}\, {2^y}) \iff 2^{y-2}\mid x$.

  \item Trouver tous les couples d'entiers $x,y \in \N$ tels que
    $3^x = 2^y + 1$.
\end{enumerate}
\finenonce{003096}


\finexercice 


\exercice{3097, quercia, 2010/03/08}

\enonce{003097}{Suites r{\'e}currentes lin{\'e}aires}
Montrer que pour tout $n \in \N$, $3^{2n+1} + 2^{n+2}$ est divisible par $7$.
\finenonce{003097}



\finexercice 


\exercice{3098, quercia, 2010/03/08}

\enonce{003098}{Suites r{\'e}currentes lin{\'e}aires}
D{\'e}terminer le reste de la division euclidienne de $2^{10n-7} + 3^{5n-2}$ par
$11$.
\finenonce{003098}


\finexercice 


\exercice{3099, quercia, 2010/03/08}

\enonce{003099}{$a \equiv b (\mathrm{mod}\, n)  \Rightarrow  a^n \equiv b^n (\mathrm{mod}\, {n^2})$}
Soient $a,b \in \Z$ et $n \in \N^*$.
Montrer que : $a \equiv b (\mathrm{mod}\, n)  \Rightarrow  a^n \equiv b^n (\mathrm{mod}\, {n^2})$.
\finenonce{003099}



\finexercice 


\exercice{3100, quercia, 2010/03/08}

\enonce{003100}{$a^2 + b^2 + c^2 + 1 \not \equiv 0 (\mathrm{mod}\, 8)$}
Montrer que :
$\forall\ a,b,c \in \Z$, $a^2 + b^2 + c^2 + 1 \not \equiv 0 \mathrm{mod}\, 8$.
\finenonce{003100}



\finexercice 


\exercice{3101, quercia, 2010/03/08}

\enonce{003101}{Cubes cons{\'e}cutifs}
Montrer que la somme de trois cubes con\-s{\'e}\-cutifs est toujours divisible par 9.
\finenonce{003101}



\finexercice 


\exercice{3102, quercia, 2010/03/08}

\enonce{003102}{$n^2 + 3n + 5$ mod 121}
Montrer que : $\forall\ n \in \Z$, $n^2 + 3n + 5$ n'est pas divisible par 121.
\finenonce{003102}


\finexercice 


\exercice{3103, quercia, 2010/03/08}

\enonce{003103}{$n\in\Z$, $n(2n+1)(7n+1)$ est divisible par 6}
Montrer que pour tout entier $n\in\Z$, $n(2n+1)(7n+1)$ est divisible par 6.
\finenonce{003103}



\finexercice 


\exercice{3104, quercia, 2010/03/08}

\enonce{003104}{$2^{32} + 1$ est divisible par 641}
Montrer sans calculatrice que $2^{32} + 1$ est divisible par 641.
\finenonce{003104}



\finexercice 


\exercice{3105, quercia, 2010/03/08}

\enonce{003105}{$3^x.7^y$ mod 10}
Trouver tous les couples $(x,y) \in \N^2$ tels que $3^x7^y$ se termine par
1 en base 10.
\finenonce{003105}


\finexercice 


\exercice{3106, quercia, 2010/03/08}

\enonce{003106}{$a^3 = \dots123456789$}
Soit $a \in \N$ premier {\`a} 10.
\begin{enumerate}
  \item Montrer que $a^4 \equiv 1 (\mathrm{mod}\,{10})$.
  \item Montrer que pour tout entier $k \in \N$, $a^{4\times10^k} \equiv 1 (\mathrm{mod}\, {10^{k+1}})$.
  \item En d{\'e}duire qu'il existe un nombre $x \in \N$ tel que $x^3$ se termine par
    $123456789$ en base 10.
\end{enumerate}
\finenonce{003106}


\finexercice 


\exercice{3107, quercia, 2010/03/08}

\enonce{003107}{$mn(m^{60} - n^{60})$ est divisible par 56786730}
Montrer que pour tous entiers $m$ et $n$, le nombre
$mn(m^{60} - n^{60})$ est divisible par 56786730.
\finenonce{003107}


\finexercice 


\exercice{3108, quercia, 2010/03/08}

\enonce{003108}{$q\mid2^p-1$}
Soient $p,q$ premiers impairs tels que $q\mid2^p-1$.
Montrer que $q\equiv 1 (\mathrm{mod}\,{2p})$.
\finenonce{003108}


\finexercice


\exercice{3109, quercia, 2010/03/08}

\enonce{003109}{Divisibilit{\'e} par 7}
Infirmer ou justifier le crit{\`e}re de divisibilit{\'e} par $7$ suivant retrouv{\'e} dans
un vieux grimoire~:
\emph{S{\'e}pare en unit{\'e}s et dizaines puis
cherche la diff{\'e}rence entre le double des unit{\'e}s et les dizaines.
Agis ainsi tant que tu as des dizaines et obtiens z{\'e}ro ou sept.
Ainsi 364 devient 28 puis 14 puis enfin 7.}
\finenonce{003109}



\finexercice 


\exercice{5115, rouget, 2010/06/30}
\enonce{005115}{*****}
$k$ est un entier impair. Montrer par récurrence que, pour $n\geq1$, la somme $1^k+2^k+...+n^k$ est un entier divisible
par $\frac{n(n+1)}{2}.$
\finenonce{005115}



\finexercice


\exercice{5116, rouget, 2010/06/30}
\enonce{005116}{****}
Pour $n\geq1$, on pose $H_n=\sum_{k=1}^{n}\frac{1}{k}$. Montrer que, pour $n\geq2$, $H_n$ n'est jamais un entier
(indication~:~montrer par récurrence que $H_n$ est le quotient d'un entier impair par un entier pair en distingant
les cas où $n$ est pair et $n$ est impair).
\finenonce{005116}


\finexercice

\exercice{5157, rouget, 2010/06/30}
\enonce{005157}{****}
Montrer que pour $n\in\Nn$, $E(\frac{1}{3}(n+2-E(\frac{n}{25})))=E(\frac{8n+24}{25})$.
\finenonce{005157}


\finexercice
\section{103.02 Sous-groupes de Z}


\exercice{286, cousquer, 2003/10/01}

\enonce{000286}{}
 Montrer qu'il est équivalent dans $\mathbb{Z}$ de dire $m$ divise~$n$, ou
$n\mathbb{Z}\subset m\mathbb{Z}$.
\finenonce{000286}



\finexercice



\exercice{287, cousquer, 2003/10/01}

\enonce{000287}{}
\begin{enumerate}
 \item Montrer que l'intersection de deux sous-groupes de $\mathbb{Z}$ est un
sous-groupe de $\mathbb{Z}$. Caractériser le sous-groupe
$a \mathbb{Z}\cap b \mathbb{Z}$. Caractériser les sous-groupes
suivants~: 
$$2\mathbb{Z}\cap 3\mathbb{Z}\;; \quad 5 \mathbb{Z}\cap 13\mathbb{Z}\;; 
    \quad 5 \mathbb{Z}\cap 25\mathbb{Z}.$$
\item Montrer que toute intersection de sous-groupes de $\mathbb{Z}$ est un 
sous-groupe de $\mathbb{Z}$. Caractériser l'intersection d'une famille finie de
sous-groupes. Caractériser les sous-groupes suivants~:
$$\bigcap_{n=1}^{17} 2^n \mathbb{Z}\;;\quad
4 \mathbb{Z}\cap 6\mathbb{Z}\cap 8\mathbb{Z}\cap 19\mathbb{Z}\cap 35\mathbb{Z}.$$
\end{enumerate}
\finenonce{000287}



\finexercice



\exercice{288, cousquer, 2003/10/01}

\enonce{000288}{}
\begin{enumerate}
\item Déterminer $2\mathbb{Z}\cup 3\mathbb{Z}$. Est-ce un sous-groupe de
$\mathbb{Z}$~?
\item Déterminer~:
$7\mathbb{Z}\cup 49\mathbb{Z}$~; $5\mathbb{Z}\cup 45\mathbb{Z}$~;
$\bigcup_{n=1}^{28} 2^n \mathbb{Z}$.
Ces ensembles sont-ils des sous-groupes de $\mathbb{Z}$~?
\item Trouver une condition nécessaire et suffisante pour qu'une réunion de
deux sous-groupes de $\mathbb{Z}$ soit un sous-groupe de~$\mathbb{Z}$.
\end{enumerate}
\finenonce{000288}



\finexercice



\exercice{289, cousquer, 2003/10/01}

\enonce{000289}{}
\begin{enumerate}
\item Soit $A$ une partie non vide de $\mathbb{Z}$~; montrer que la famille 
des sous-groupes contenant~$A$ n'est pas vide. Soit $H$ une partie 
contenant~$A$. Montrer l'équivalence des conditions suivantes~:
\begin{enumerate}
    \item[i)]  $H$ est l'intersection des sous-groupes de $\mathbb{Z}$ 
    qui contiennent~$A$,
    \item[ii)]  $H$ est le plus petit sous-groupe de $\mathbb{Z}$ qui 
    contient~$A$,
    \item[iii)]  $H$ est l'ensemble des sommes finies d'éléments de $A$ 
    ou d'éléments dont l'opposé est dans~$A$.
\end{enumerate}
Si ces conditions sont vérifiées on dit que $H$ est le sous-groupe
engendré par~$A$.
\item Soient $m\mathbb{Z}$ et $n\mathbb{Z}$ deux sous-groupes de $\mathbb{Z}$.
Montrer que
$$m\mathbb{Z}+n\mathbb{Z}= \{mu+nv \mid u,v \in\mathbb{Z}\}$$
\begin{enumerate}
\item[a)] est un sous-groupe de $\mathbb{Z}$,
\item[b)] contient $m\mathbb{Z}$ et $n\mathbb{Z}$,
\item[c)] est contenu dans tout sous-groupe de $\mathbb{Z}$ qui 
contient $m\mathbb{Z}$ et $n\mathbb{Z}$.
\item[d)] Si $m\mathbb{Z}+n\mathbb{Z}= d\mathbb{Z}$, 
que peut-on dire de $d$~?
\end{enumerate}
\item Déterminer les sous-groupes engendrés par~:
$14\mathbb{Z}\cup 35\mathbb{Z}$~;
$4\mathbb{Z}\cup 8\mathbb{Z}\cup 6\mathbb{Z}\cup 64\mathbb{Z}$~;
$2\mathbb{Z}\cup 3\mathbb{Z}$~;
$4\mathbb{Z}\cup 21\mathbb{Z}$~;
$5\mathbb{Z}\cup 25\mathbb{Z}\cup 7\mathbb{Z}$~;
$\{70 ,4\}$. 
\end{enumerate}
\finenonce{000289}



\finexercice


\section{103.03 Pgcd, ppcm, algorithme d'Euclide}


\exercice{290, bodin, 1998/09/01}
\video{uMmuJ-TCa6U}
\enonce{000290}{}

Calculer le pgcd des nombres suivants :
\begin{enumerate}
    \item 126, 230.
    \item 390, 720, 450.
    \item 180, 606, 750.
\end{enumerate}
\finenonce{000290} 


\finexercice


\exercice{291, bodin, 1998/09/01}

\enonce{000291}{}
\begin{enumerate}
    \item Calculer le ppcm des nombres  : 108 et 144 ; 128 et 230 ; 6, 16 et 50.
    \item Montrer que si $a\ge 1$ et $b\ge 1$ sont des entiers
    de pgcd $d$ et, si on pose $a=da' ; b=db' $, le ppcm de $a$ et $b$ est $da'b'$.
    \item Montrer que si $a,b,c$ sont des entiers sup\'erieurs \`a  $1$,
    on a :
    $$\text{ppcm}(a,b,c)=\text{ppcm}(\text{ppcm}(a,b),c).$$

\end{enumerate}
\finenonce{000291}



\finexercice



\exercice{292, bodin, 1998/09/01}
\video{haNOPZqNwMY}
\enonce{000292}{}
 D\'eterminer les couples d'entiers naturels de pgcd
18 et de somme 360. De m\^eme avec pgcd 18 et produit 6480.
\finenonce{000292} 


\finexercice


\exercice{293, bodin, 1998/09/01}

\enonce{000293}{}
Si $a,b,c,d$ sont des entiers sup\'erieurs \`a  $1$,
montrer que l'on a :
$$(a,b,c,d)=((a,b),(c,d))$$ o\`u ( , ) d\'esigne le pgcd .
\finenonce{000293}



\finexercice



\exercice{294, bodin, 1998/09/01}

\enonce{000294}{}
\begin{enumerate}
\item Soient $a,b,c$ des entiers relatifs tels que $(a,b)\not =(0,0)$,
montrer que pour que l'\'equation
$$ax+by=c$$
ait une solution $(x,y)$ en entiers relatifs $x$ et $y$,
il faut et il suffit que le pgcd de $a$ et $b$ divise $c$.
\item R\'esoudre en entiers relatifs les \'equations suivantes :
$$7x-9y=1,$$
$$7x-9y=6,$$
$$11x+17y=5.$$
\end{enumerate}
\finenonce{000294}



\finexercice



\exercice{295, bodin, 1998/09/01}

\enonce{000295}{}
Soient $a$ et $b$ deux entiers tels que $a\ge b\ge 1$ et $\mathrm{pgcd}(a,b)=1$.
\begin{enumerate}
    \item Montrer que $\text{ pgcd}(a+b,a-b)=1$ ou $2$,
    \item Si pgcd$(a,b)=1 $, montrer que pgcd$(a+b,ab)=1$,
    \item Si pgcd$(a,b)=1 $, montrer que pgcd$(a+b,a^2+b^2)=1$ ou $2$.
\end{enumerate}
\finenonce{000295}



\finexercice



\exercice{296, bodin, 1998/09/01}
\video{uekZIfTccPg}
\enonce{000296}{}

Calculer par l'algorithme d'Euclide : $\pgcd(18480,9828)$.
En d\'eduire une \'ecriture de $84$ comme combinaison lin\'eaire de $18480$ et
$9828$.
\finenonce{000296} 


\finexercice


\exercice{297, cousquer, 2003/10/01}

\enonce{000297}{}
 Déterminer le pgcd de $99\,099$ et $43\,928$.
Déterminer le pgcd de $153\,527$ et $245\,479$.
\finenonce{000297}



\finexercice



\exercice{298, bodin, 1998/09/01}

\enonce{000298}{}
D\'eterminer l'ensemble de tous les couples $(m,n)$ tels que
$$ 955m+183n=1.$$

\finenonce{000298}



\finexercice



\exercice{299, bodin, 1998/09/01}

\enonce{000299}{}
Calculer, en pr\'ecisant la m\'ethode suivie,
$$ a=\text{pgcd} (720,252) \quad\quad b = \text{ppcm} (720,252)$$
ainsi que deux entiers $u$ et $v$ tels que $720u+252v=a$.
\finenonce{000299}



\finexercice



\exercice{300, bodin, 1998/09/01}

\enonce{000300}{}
D\'emontrer :
$$ a\wedge(b_1b_2)=1 \Leftrightarrow (a\wedge b_1=1 \text{ et } a\wedge b_2 =1), $$
puis par r\'ecurence :
$$ a\wedge(b_1 \ldots b_n)=1 \Leftrightarrow \forall i=1,\ldots,n\ \ a\wedge b_i =1. $$
\finenonce{000300}



\finexercice



\exercice{301, bodin, 1998/09/01}

\enonce{000301}{}
D\'emontrer pour $m,n \in \Nn^*$ :
$$a^m \wedge b^n=1 \Rightarrow a \wedge b=1.$$
\finenonce{000301}



\finexercice



\exercice{302, bodin, 1998/09/01}

\enonce{000302}{}
D\'eteminer deux entiers naturels connaissant leur somme, $1008$,
et leur pgcd, $24$.
\finenonce{000302}



\finexercice



\exercice{303, bodin, 1998/09/01}
\video{atC2uUlIQ64}
\enonce{000303}{}
Notons $a=1\;111\;111\;111$ et $b=123\;456\;789$.
\begin{enumerate}
    \item Calculer le quotient et le reste de la division euclidienne de $a$ par $b$.
    \item Calculer $p=\, \text{pgcd}(a,b)$.
    \item D\'eterminer deux entiers relatifs $u$ et $v$ tels que $au+bv=p$.
\end{enumerate}

\finenonce{000303} 


\finexercice


\exercice{304, bodin, 1998/09/01}

\enonce{000304}{}
Soient $m$ et $n$ deux entiers $(m>n>0)$ et $a \ge 2$ un entier.
Montrer que le reste de la division euclidienne de $a^m-1$ par $a^n-1$ est
$a^r-1$ o\`u $r$ est le reste  de la division euclidienne de $m$ par $n$, et
que le pgcd de $a^m-1$ et $a^n-1$ est $a^d-1$, o\`u $d$ est le pgcd de $m$ et $n$.
\finenonce{000304}



\finexercice



\exercice{305, gourio, 2001/09/01}
\video{IXE-IGeo-ts}
\enonce{000305}{}

R\'{e}soudre dans ${\Zz}:1665x+1035y=45.$
\finenonce{000305} 




\exercice{306, cousquer, 2003/10/01}

\enonce{000306}{}
 Montrer qu'il n'existe pas d'entiers 
$m$ et $n$ tels que $$m+n=101\quad \mbox{et} \quad \mbox{pgcd}(m,n)=3$$
\finenonce{000306}



\finexercice



\exercice{307, cousquer, 2003/10/01}

\enonce{000307}{}
 Soit $m$ et $n$ deux entiers positifs.
\begin{enumerate}
\item  Si $\mbox{pgcd}(m,4)=2$ et $\mbox{pgcd}(n,4)=2$, montrer que $\mbox{pgcd}(m+n,4)=4$.
\item  Montrer que pour chaque entier $n$, $6$ divise $n^3-n$. 
\item  Montrer que pour chaque entier $n$, $30$ divise $n^5-n$. 
\item  Montrer que si $m$ et $n$ sont des entiers impairs, 
$m^2+n^2$ est pair mais non divisible par~$4$. 
\item  Montrer que le produit de quatre entiers consécutifs est 
divisible par~$24$.
\item  Montrer que si  $\mbox{pgcd}(a,b)=1$, alors
\begin{itemize}
\item   $\mbox{pgcd}(a+b,a-b)\in \{1,2\}$,
\item   $\mbox{pgcd}(2a+b,a+2b)\in \{1,3\}$,
\item  $\mbox{pgcd}(a^2+b^2,a+b)\in \{1,2\}$, 
\item  $\mbox{pgcd}(a+b,a^2-3ab+b^2)\in \{1,5\}$.
\end{itemize}
\end{enumerate}
\finenonce{000307}



\finexercice



\exercice{308, cousquer, 2003/10/01}

\enonce{000308}{}
Trouver une CNS pour que $ax+b\equiv 0 \ \rm{mod} \ n$ ait une solution.
\finenonce{000308}



\finexercice



\exercice{309, cousquer, 2003/10/01}

\enonce{000309}{}
\begin{enumerate}
  \item Calculer $\mbox{pgcd}(18,385)$ par l'algorithme d'Euclide,  en déduire un couple
$(u_0,v_0)\in\mathbb{Z}^2$ solution de l'équation $18u+385v=1$, avec $(u,v)\in\mathbb{Z}^2$.
  \item Fournir enfin l'ensemble des solutions entières de $$18u+385v=1;\quad
18u+385v=3;\quad 54u+1155v=3;\quad 54u+1155v=5.$$
\end{enumerate} 
\finenonce{000309}



\finexercice



\exercice{310, cousquer, 2003/10/01}

\enonce{000310}{}
\noindent Trouver $a$ et $b$ entiers naturels tels que
\begin{enumerate}
  \item $a+b=2070$ et $\mbox{ppcm}(a,b)=9180$ ;
  \item $a^2+b^2=5409$ et $\mbox{ppcm}(a,b)=360$ (on pourra commencer par montrer que 
$\mbox{pgcd}(a,b)$
divise $\mbox{pgcd}(5409,360)$ et considérer ensuite différents cas).
\end{enumerate}
\finenonce{000310}



\finexercice



\exercice{311, cousquer, 2003/10/01}

\enonce{000311}{}
 Résoudre  dans $\mathbb{Z}$ les équations\,:
$35x \equiv 7 \ \rm{mod} \ 4; \  \ 22x \equiv 33\ \rm{mod}\ 5 $
\finenonce{000311}



\finexercice



\exercice{312, cousquer, 2003/10/01}

\enonce{000312}{}
 Résoudre  dans $\mathbb{Z}$ le système suivant\,:
$$S : \left\{
\begin{array}{cccc}
x&\equiv &4&\rm{mod}\ 6\\
x&\equiv &7&\rm{mod}\ 9
\end{array}
\right.
$$
On recherchera d'abord une solution particulière.
\finenonce{000312}



\finexercice



\exercice{313, cousquer, 2003/10/01}

\enonce{000313}{}
\begin{enumerate}
  \item  Résoudre dans $\mathbb{Z}$ les équations\,:
$\ x^2\equiv 2\ \rm{mod}\ 6;\quad \ x^3\equiv 3\ \rm{mod}\ 9$.
  \item  Résoudre dans $\mathbb{Z}^2$ les équations  suivantes\,: 
$5x^2 +2xy-3=0\ ;\quad y^2+4xy-2=0$.
\end{enumerate}
\finenonce{000313}



\finexercice



\exercice{314, cousquer, 2003/10/01}

\enonce{000314}{}
 Résoudre dans $\mathbb{Z}^2$ les équations  suivantes\,:
$$\begin{array}{lcr}
\hbox{a) } 17x+6y=1 &\qquad &  \hbox{b) } 27x+25y=1\\
\hbox{c) } 118x+35y=1 &\qquad &  \hbox{d) } 39 x+26 y=1
\end{array}$$
\finenonce{000314}



\finexercice



\exercice{315, cousquer, 2003/10/01}

\enonce{000315}{}
 Montrer que si $a$ divise $42n+37$ et $7n+4$, pour une valeur de $n$ donnée, 
alors $a$ divise $13$. Quelles sont les valeurs possibles 
pour~$n$~?
\finenonce{000315}



\finexercice



\exercice{316, cousquer, 2003/10/01}

\enonce{000316}{}
  Trouver $\mbox{pgcd}(-357,629)$ et trouver des entiers $x$ et $y$ 
tels que $$\mbox{pgcd}(-357,629)=-357x+629y$$ 
\finenonce{000316}



\finexercice



\exercice{317, cousquer, 2003/10/01}

\enonce{000317}{}
 Trouver $\mbox{pgcd}(2183,6313)=d$ et trouver des entiers $x$ et $y$ tels que $$d=2183x+6313y$$
\finenonce{000317}



\finexercice



\exercice{318, cousquer, 2003/10/01}

\enonce{000318}{}
Supposons $\mbox{pgcd}(a,b)=d$ et soit $x_{0}$ et $y_{0}$
 des entiers tels que $d=ax_{0}+by_{0}$. Montrer que~:
 \begin{enumerate}
\item  $\mbox{pgcd}(x_{0},y_{0})=1$,
\item  $x_{0}$ et $y_{0}$ ne sont pas uniques.
\end{enumerate}
\finenonce{000318}



\finexercice



\exercice{319, cousquer, 2003/10/01}

\enonce{000319}{}
Soit $a$, $b$, $c$ des entiers.
\begin{enumerate}
\item Montrer que $\mbox{pgcd}(ca,cb)=\vert c\vert\;\mbox{pgcd}(a,b)$.
\item Montrer que $\mbox{pgcd}(a^2,b^2)=(\mbox{pgcd}(a,b))^2$.
\item Montrer que si $\mbox{pgcd}(a,b)=1$ et si $c$ divise $a$, alors $\mbox{pgcd}(c,b)=1$.
\item Montrer que  $\mbox{pgcd}(a,bc)=1 \iff \mbox{pgcd}(a,b)=\mbox{pgcd}(a,c)=1$.
\item Montrer que si $\mbox{pgcd}(b,c)=1$ alors $\mbox{pgcd}(a,bc)=\mbox{pgcd}(a,b) \mbox{pgcd}(a,c)$.
\item  Montrer que $\mbox{pgcd}(a,b)=\mbox{pgcd}(a+b,\mbox{ppcm}(a,b))$. 
\end{enumerate}
\finenonce{000319}



\finexercice



\exercice{320, cousquer, 2003/10/01}

\enonce{000320}{}

 En divisant un nombre par $8$, un élève  a obtenu $4$ pour reste~; en divisant
ce même  nombre par $12$, il a obtenu $3$ pour reste. Qu'en pensez-vous~?

 Le fort en calcul de la classe, qui ne fait jamais d'erreur, a divisé le
millésime de l'année par $29$, il a trouvé $25$ pour reste~; il a divisé le
même millésime par $69$, il a trouvé $7$ pour reste. En quelle année cela
se passait-il~?
\finenonce{000320}



\finexercice



\exercice{321, cousquer, 2003/10/01}

\enonce{000321}{}
 Trouver deux nombres sachant que leur somme est $581$ et que le quotient de leur
PPCM par leur pgcd est $240$.
\finenonce{000321}



\finexercice



\exercice{322, cousquer, 2003/10/01}

\enonce{000322}{}
 Trouver les solutions entières de l'équation~:
$$102x-18018y=18.$$
Combien y a-t-il de solutions telles que $x$ et $y$ soient
compris entre entre $0$ et $4000$~?
\finenonce{000322}



\finexercice



\exercice{323, cousquer, 2003/10/01}

\enonce{000323}{}
 Le pgcd de deux nombres est $12$~; les quotients successifs obtenus dans le
calcul de ce pgcd par l'algorithme d'Euclide sont $8$, $2$ et~$7$. 
Trouver ces deux nombres.
\finenonce{000323}



\finexercice



\exercice{324, cousquer, 2003/10/01}

\enonce{000324}{}
 Trouver les couples de nombres $a$ et $b$, divisibles par~$3$, vérifiant les
propriétés suivantes~: leur ppcm est $7560$, et si on augmente chacun de ces
nombres d'un tiers de sa valeur, le pgcd des deux nombres obtenus 
est $84$.
\finenonce{000324}



\finexercice



\exercice{325, cousquer, 2003/10/01}

\enonce{000325}{}
 Un terrain rectangulaire dont les dimensions en mètres $a$ et $b$ sont des
nombres entiers, a pour aire $3024\; \mbox{m}^2$. Calculer son périmètre
sachant que le pgcd de $a$ et $b$ est~$6$. Combien y a-t-il de solutions
possibles~?
\finenonce{000325}



\finexercice



\exercice{326, cousquer, 2003/10/01}

\enonce{000326}{}
\begin{enumerate}
  \item Dans $\mathbb{Z}/n\mathbb{Z}$, écrire l'ensemble des multiples de $\bar x$, classe
de $x$, pour $x$ variant de $0$ à $n-1$ dans chacun des cas suivants~:
$\mathbb{Z}/5\mathbb{Z}$, $\mathbb{Z}/6\mathbb{Z}$, $\mathbb{Z}/8\mathbb{Z}$.
    \item Dans $\mathbb{Z}/n\mathbb{Z}$, montrer l'équivalence des trois propositions~:
\begin{itemize}
\item[i)] $\bar x$ est inversible~;
\item[ii)] $x$ et $n$ sont premiers entre eux~;
\item[iii)] $\bar x$ engendre  $\mathbb{Z}/n\mathbb{Z}$, c'est à dire que
l'ensemble des multiples de $\bar x$ est $\mathbb{Z}/n\mathbb{Z}$.
\end{itemize}
    \item La classe de $18$ est-elle inversible dans $\mathbb{Z}/49\mathbb{Z}$~? Si oui, quel
est son inverse~? (On pourra utiliser le théorème de Bézout).
\end{enumerate}
\finenonce{000326}



\finexercice



\exercice{327, cousquer, 2003/10/01}

\enonce{000327}{}
 Résoudre dans $\mathbb{Z}$ les équations suivantes~:
\begin{enumerate}
\item $91x-65y=156$.
\item $135 x-54y=63$.
\item $72x+35y=13$.
\end{enumerate}
\finenonce{000327}



\finexercice



\exercice{328, cousquer, 2003/10/01}

\enonce{000328}{}
 Résoudre dans $\mathbb{N}$ les équations suivantes~:
\begin{enumerate}
\item $31x-13y=1$.
\item $31x-13y=-1$.
\end{enumerate}
\textit{Application~:}
Au bord d'une piscine pleine d'eau, on dispose d'une cuve fixe de 31 
litres munie à sa base d'un robinet de vidange, et d'un seau de 13 
litres. Expliquer comment opérer pour obtenir exactement 1~litre dans 
le seau.
\finenonce{000328}



\finexercice



\exercice{329, cousquer, 2003/10/01}

\enonce{000329}{}
 Résoudre dans $\mathbb{N}$ l'équation $77x+105y=2401$.
\finenonce{000329}



\finexercice



\exercice{330, cousquer, 2003/10/01}

\enonce{000330}{}
 Dans un pays nommé ASU, dont l'unité monétaire est le rallod, 
la banque nationale émet seulement des 
billets de 95 rallods et des pièces de 14 rallods.
\begin{enumerate}
\item Montrer qu'il est possible de payer n'importe quelle somme 
entière (à condition bien sûr que les deux parties disposent chacune 
d'assez de pièces et de billets).
\item On suppose que vous devez payer une somme~$S$, que vous avez une 
quantité illimitée de pièces et de billets, mais que votre créancier 
ne puisse pas rendre la monnaie. Ainsi, il est possible de payer si 
$S=14$, mais pas si $S=13$ ou si $S=15$\dots{} 
Montrer qu'il est toujours possible de payer si $S$ est assez grande.
Quelle est la plus grande valeur de~$S$ telle qu'il soit impossible 
de payer~$S$~?
\end{enumerate}
\finenonce{000330}



\finexercice



\exercice{331, cousquer, 2003/10/01}

\enonce{000331}{}
 Trouver tous les points à coordonnées entières du plan 
d'équation $6x+10y+15z=1997$. Combien y a-t-il de solutions dans 
$\mathbb{N}^3$~?
\finenonce{000331}



\finexercice



\exercice{332, cousquer, 2003/10/01}

\enonce{000332}{}
\begin{enumerate}
\item Trouver tous les points à coordonnées entières de la droite 
de l'espace d'équations 
$\left\{
\begin{array}{rcl}
    4x-2y-z-5 & = & 0  \\
    x+3y-4z-7 & = & 0
\end{array}\right.$.
\item Même question avec la droite 
$\left\{
\begin{array}{rcl}
    x+3y-5z -5 & = & 0  \\
    4x-2y+z+13 & = & 0
\end{array}\right.$.
\end{enumerate}
\finenonce{000332}



\finexercice



\exercice{333, cousquer, 2003/10/01}

\enonce{000333}{}
 Résoudre dans $\mathbb{N}$ et dans $\mathbb{Z}$ l'équation 
$$\frac{1}{x}+\frac{1}{y}=\frac{1}{15}$$
\finenonce{000333}



\finexercice



\exercice{334, cousquer, 2003/10/01}

\enonce{000334}{}
 Un coq coûte $5$ pièces d'argent, une poule $3$ pièces, et un lot de 
quatre poussins $1$ pièce. Quelqu'un a acheté $100$ volailles pour $100$ 
pièces~; combien en a-t-il acheté de chaque sorte~?
\finenonce{000334}



\finexercice



\exercice{335, cousquer, 2003/10/01}

\enonce{000335}{}
Soient $a$ et $b$ deux nombres entiers relatifs.
On note $d$ leur pgcd.
Construisons les suites $a_n$ et $b_n $ $n \in \mathbb{N},$ à valeurs dans $\mathbb{Z}
$de la manière suivante :
\begin{eqnarray*}  
a_0 &=& a\\
b_0 &=& b\\
\end{eqnarray*}
et pour tout $n \in \mathbb{N},$ on pose $a_{n+1}= b_n$ et $b_{n+1}= r$ 
où $r$ est le reste de la division euclidienne de $a_n$ par $b_n.$

\begin{enumerate}
\item 
Montrer que si $d_n$ est le pgcd de $a_n$ et $b_n$ alors $d_n$ est également le pgcd de $a_{n+1}$ et $b_{n+1}.$
\item 
Déduire de la questionh précédente que $d$ est le pgcd des nombres $a_n$ et $b_n$ pour tout $n \in \mathbb{N}.$
\item 
Montrer que la suite $b_n$ est strictement décroissante.
Que peut-on en déduire?
\item 
Déduire de ce qui précède que pour tout couple d'entiers relatifs $(a,b)$ 
il existe un couple d'entier relatifs $(u,v)$ tel que:
$$
d = au+bv.
$$
\end{enumerate}
\finenonce{000335}



\finexercice



\exercice{3110, quercia, 2010/03/08}

\enonce{003110}{}
Soient $a,b,c \in \Z$ tels que $a\wedge b = 1$.
Montrer que $a\wedge(bc) = a\wedge c$.
\finenonce{003110}



\finexercice 


\exercice{3111, quercia, 2010/03/08}

\enonce{003111}{pgcd$( a+b, ppcm(a,b) )$}
Soient $a,b$ entiers, $d = a \wedge b$, $m = a \vee b$.
Chercher $(a+b) \wedge m$.
\finenonce{003111}



\finexercice 


\exercice{3112, quercia, 2010/03/08}

\enonce{003112}{pgcd$((a-b)^3, a^3-b^3)$}
Soient $a,b \in \Z$. Chercher $(a-b)^3 \wedge (a^3 - b^3)$.
\finenonce{003112}


\finexercice 


\exercice{3113, quercia, 2010/03/08}

\enonce{003113}{pgcd$( n^3 + n, 2n + 1 )$}
Soit $n \in \N$. Chercher $(n^3+n) \wedge (2n+1)$.
\finenonce{003113}


\finexercice 


\exercice{3114, quercia, 2010/03/08}

\enonce{003114}{pgcd$( 15n^2+8n+6, 30n^2+21n+13 )$}
Soit $n \in \N$. Chercher $(15n^2+8n+6) \wedge (30n^2+21n+13)$.
\finenonce{003114}


\finexercice 


\exercice{3115, quercia, 2010/03/08}

\enonce{003115}{pgcd et ppcm impos{\'e}s}
Soient $d,m \in \N^*$. Donner une condition n{\'e}c{\'e}ssaire et suffisante sur
$d$ et $m$ pour qu'il existe $a,b \in \Z$ tels que $a\wedge b = d$
et $a\vee b = m$.

R{\'e}soudre ce probl{\`e}me pour $d = 50$ et $m = 600$.

\finenonce{003115}


\finexercice 


\exercice{3116, quercia, 2010/03/08}

\enonce{003116}{ppcm$(x,y) + 11$pgcd$(x,y) = 203$}
Trouver les couples d'entiers $(x,y) \in \Z^2$ tels que :
$x\vee y + 11(x\wedge y) = 203$.
\finenonce{003116}


\finexercice 


\exercice{3117, quercia, 2010/03/08}

\enonce{003117}{$x^2+y^2 = 85113$, ppcm$(x,y) = 1764$}
R{\'e}soudre : $$\begin{cases}x^2+y^2 = 85113\cr x\vee y = 1764.\end{cases}$$
\finenonce{003117}


\finexercice 


\exercice{3118, quercia, 2010/03/08}

\enonce{003118}{ppcm$(x,y) = 210$ pgcd$(x,y)$, $y - x =$ pgcd$(x,y)$}
R{\'e}soudre : $\begin{cases} x \vee y = 210(x \wedge y)\cr y-x = x \wedge y.\end{cases}$
\finenonce{003118}


\finexercice 


\exercice{3119, quercia, 2010/03/08}

\enonce{003119}{pgcd$(x,y) = x+y-1$}
R{\'e}soudre dans $\Z$ : $x \wedge y = x + y - 1$.
\finenonce{003119}


\finexercice 


\exercice{3120, quercia, 2010/03/08}

\enonce{003120}{ppcm$(x,y) = x+y-1$}
R{\'e}soudre dans $\Z^*$ : $x \vee y = x + y - 1$.
\finenonce{003120}


\finexercice 


\exercice{3121, quercia, 2010/03/08}

\enonce{003121}{pgcd$(x,y) = x-y$, ppcm$(x,y) = 300$}
R{\'e}soudre dans $\N^*$ :
$\begin{cases}x \wedge y = x-y \cr x \vee y = 300.\end{cases}$
\finenonce{003121}


\finexercice 


\exercice{3122, quercia, 2010/03/08}

\enonce{003122}{pgcd$( a^n - 1, a^m - 1 )$}
Soient $a,m,n\ \in \N^{*}$, $a\ge 2$, et $d = (a^n - 1) \wedge (a^m - 1)$.
\begin{enumerate}
  \item   Soit $n = qm + r$ la division euclidienne de $n$ par $m$.
      D{\'e}montrer que $a^n \equiv a^r (\mathrm{mod}\, {a^m - 1})$.

  \item   En d{\'e}duire que $d = (a^r - 1) \wedge (a^m - 1)$,
      puis $d = a^{(n \wedge m)} - 1$.

  \item   A quelle condition $a^m - 1$ divise-t-il $a^n-1$ ?
\end{enumerate}
\finenonce{003122}



\finexercice 


\exercice{3123, quercia, 2010/03/08}

\enonce{003123}{pgcd multiple}
Soient $a_1, \dots, a_n \in \N^{*}$ et $b_i = \prod_{j \ne i} a_j$.
Montrer que $a_1, \dots, a_n$ sont {\it deux {\`a} deux\/} premiers entre eux
si et seulement si
$b_1, \dots, b_n$ sont premiers entre eux {\it dans leur ensemble}.
\finenonce{003123}



\finexercice


\exercice{3124, quercia, 2010/03/08}

\enonce{003124}{{\'E}quations {\`a} coefficients entiers}
Soient $a,b,c$ trois entiers relatifs. On consid{\`e}re l'{\'e}quation : $ax+by=c$,
dont on recherche les solutions dans $\Z^2$.
\begin{enumerate}
  \item Donner une condition n{\'e}c{\'e}ssaire et suffisante pour que cette {\'e}quation admette
    une solution.
  \item Soit $(x_0,y_0)$ une solution du probl{\`e}me de B{\'e}zout : $ax_0 + by_0 = d$.
    D{\'e}ter\-miner toutes les solutions de $ax + by = c$ en fonction de $a$, $b$,
    $c$, $d$, $x_0$ et $y_0$.
  \item R{\'e}soudre dans $\Z^2$ : $2520x - 3960y = 6480$.
\end{enumerate}
\finenonce{003124}


\finexercice 


\exercice{3125, quercia, 2010/03/08}

\enonce{003125}{{\'E}quations {\`a} coefficients entiers}

R{\'e}soudre dans $\Z$ :
\begin{enumerate}
  \item   $95x + 71y = 46$.
  \item   $20x - 53y = 3$.
  \item   $12x + 15y + 20z = 7$.
\end{enumerate}
\finenonce{003125}


\finexercice 


\exercice{3126, quercia, 2010/03/08}

\enonce{003126}{Congruences simultan{\'e}es}
\begin{enumerate}
  \item Soient $a,b,a',b' \in \Z$ avec $b \wedge b' = 1$.
    Montrer que le syst{\`e}me : $\begin{cases}x \equiv a  (\mathrm{mod}\, b) \cr
                             x \equiv a' (\mathrm{mod}\, {b'}) \end{cases}$
    poss{\`e}de des solutions et qu'elles sont congrues entre elles modulo $bb'$.
  \item G{\'e}n{\'e}raliser.
\end{enumerate}
\finenonce{003126}



\finexercice 


\exercice{3127, quercia, 2010/03/08}

\enonce{003127}{Congruences simultan{\'e}es}
R{\'e}soudre :
\begin{enumerate}
  \item $\begin{cases}x\equiv 2  \hfill(\mathrm{mod}\, {140}) \cr
            x\equiv -3 \hfill(\mathrm{mod}\, {99}).\end{cases}$

  \item $\begin{cases}x\equiv 3 \hfill(\mathrm{mod}\, 4)\cr
            x\equiv-2 \hfill(\mathrm{mod}\, 3)\cr
            x\equiv 7 \hfill(\mathrm{mod}\, 5).\end{cases}$
\end{enumerate}
\finenonce{003127}


\finexercice 


\exercice{3128, quercia, 2010/03/08}

\enonce{003128}{Congruences simultan{\'e}es}
Une bande de 17 pirates dispose d'un butin compos{\'e} de $N$ pi{\`e}ces d'or
d'{\'e}gale valeur. Ils d{\'e}cident de se le partager {\'e}galement et de donner le
reste au cuisinier (non pirate). Celui ci re{\c c}oit 3 pi{\`e}ces.

Mais une rixe {\'e}clate et 6 pirates sont tu{\'e}s. Tout le butin est reconstitu{\'e} et
partag{\'e} entre les survivants comme pr{\'e}c{\'e}demment; le cuisinier re{\c c}oit alors 4
pi{\`e}ces.

Dans un naufrage ult{\'e}rieur, seuls le butin, 6 pirates et le cuisinier sont
sauv{\'e}s. Le butin est {\`a} nouveau partag{\'e} de la m{\^e}me mani{\`e}re et le cuisinier
re{\c c}oit 5 pi{\`e}ces.

Quelle est alors la fortune minimale que peut esp{\'e}rer le cuisinier lorsqu'il
d{\'e}cide d'empoisonner le reste des pirates~?
\finenonce{003128}


\finexercice 


\exercice{3129, quercia, 2010/03/08}

\enonce{003129}{D{\'e}composition {\`a} coefficients positifs}
Soient $a,b \in \N^*$ premiers entre eux.
Montrer que~: $\forall\ x \ge ab$, $\exists\ u,v \in \N$ tels que $au + bv = x$.
\finenonce{003129}



\finexercice 

\section{103.04 Nombres premiers, nombres premiers entre eux}


\exercice{336, bodin, 1998/09/01}
\video{xFq39ocyF_c}
\enonce{000336}{}
 Soient $a,b$ des entiers sup\'erieurs ou \'egaux
\`a $1$. Montrer :
\begin{enumerate}
    \item $(2^a-1) | (2^{ab}-1)$ ;
    \item $2^p-1 \text{ premier}\ \   \Rightarrow \ \  p \text { premier }$ ;
    \item $\pgcd(2^a-1,2^b-1) =  2^{\pgcd(a,b)}-1$.
\end{enumerate}
\finenonce{000336} 


\finexercice

\exercice{337, bodin, 1998/09/01}
\video{POYqKNGJr44}
\enonce{000337}{}
D\'emontrer que, si $a$ et $b$ sont des entiers
premiers entre eux, il en est de m\^eme des entiers $a+b$ et $ab$.

\finenonce{000337} 


\finexercice


\exercice{338, bodin, 1998/09/01}

\enonce{000338}{}
R\'esoudre l'\'equation $29x-11y=1$ dans $\Zz$.

On consid\`ere maintenant l'\'equation $29x-11y=5$.
D\'eduire de ce qui pr\'ec\`ede une solution particuli\`ere de
cette \'equation, puis en donner la solution g\'en\'erale.
\finenonce{000338}



\finexercice



\exercice{339, bodin, 1998/09/01}
\video{RR5gG5ZLCGs}
\enonce{000339}{}
Soit $p$ un nombre premier.
\begin{enumerate}
    \item Montrer que $\forall i\in \Nn, 0< i < p$ on a : $$C_p^i\text{  est divisible par }p .$$
    \item Montrer par r\'ecurence que : $$\forall p \text{ premier}, \forall a\in \Nn^* ,\text{ on a } a^p-a \text{ est divisible par } p.$$
\end{enumerate}

\finenonce{000339} 


\finexercice


\exercice{340, bodin, 1998/09/01}

\enonce{000340}{}
\begin{enumerate}
\item Soit $(x,y,z)\in \Nn^3$ .
Montrer que :
\begin{eqnarray*}
    x^2+y^2=z^2\Leftrightarrow \exists (x',y',z')\in \Nn^3,\exists n\in \Nn \text{ tq } \\
     \text{ pgcd}(x',y',z')=1 \\
    {x'}^2+{y'}^2={z'}^2 \\
    x=nx' \text{ et }y=ny' \text{ et } z=nz'.
\end{eqnarray*}
\item Soit$(x,y,z)\in \Nn^3$ tels que $x^2+y^2=z^2$.
On suppose que pgcd$(x,y,z)=1$
    \begin{enumerate}
    \item Montrer que $x$ et $y$ ne sont pas de m\^emes parit\'e.
    \item On suppose $x$ pair et $y$ impair.
    On pose :
    $$x=2u,\ z-y=2v,\ z+y=2w$$
    avec $(u,v)\in \N^*$. Montrer que $v$ et $w$ sont premiers entre eux.
    \item
    Montrer que
    $$x=2mn,\ y=m^2-n^2,\ z=m^2+n^2$$
    avec $m$ et $n$ entiers naturels de parit\'e diff\'erentes.
    \item Montrer que si
    $$x=2mn,\ y=m^2-n^2,\ z=m^2+n^2$$
    alors
    $$x^2+y^2=z^2.$$
    \end{enumerate}
\end{enumerate}
\finenonce{000340}



\finexercice



\exercice{341, bodin, 1998/09/01}
\video{I-7jctXUUZk}
\enonce{000341}{}
\begin{enumerate}
\item
Montrer par r\'ecurrence que $\forall n\in \Nn ,\forall k\geqslant 1$ on a :
$$2^{2^{n+k}}-1=\left( 2^{2^n}-1 \right) \times \prod_{i=0}^{k-1}(2^{2^{n+i}}+1).$$
\item
On pose $F_n=2^{2^n}+1$. Montrer que pour $m\not= n$,
 $F_n$ et $F_m$ sont premiers entre eux.
\item En d\'eduire qu'il y a une infinit\'e de nombres premiers.
\end{enumerate}


\finenonce{000341} 


\finexercice


\exercice{342, cousquer, 2003/10/01}

\enonce{000342}{}
 Les nombres $a$, $b$, $c$, $d$ étant des éléments non nuls de $\mathbb{Z}$, 
dire si les propriétés suivantes sont vraies ou 
fausses, en justifiant la réponse.
\begin{enumerate}
    \item Si $a$ divise $b$ et $b$ divise $c$, alors $a$ divise $c$.
    \item Si $a$ divise $b$ et $c$, alors $a$ divise $2b+3c$.
    \item S'il existe $u$ et $v$ entiers tels que $au+bv=4$ alors $\mbox{pgcd}(a,b)=4$.
    \item Si $7a-9b=1$ alors $a$ et $b$ sont premiers entre eux.
    \item Si $a$ divise $b$ et $b$ divise $c$ et $c$ divise $a$, alors 
$\vert a\vert = \vert b\vert$.
    \item \og $a$ et $b$ premiers entre eux\fg{} équivaut à \og 
$\mbox{ppcm}(a,b)=\vert ab\vert$\fg.
    \item Si $a$ divise $c$ et $b$ divise $d$, alors $ab$ divise $cd$.
    \item Si $9$ divise $ab$ et si $9$ ne divise pas $a$, alors $9$ divise $b$.
    \item Si $a$ divise $b$ ou $a$ divise $c$, alors $a$ divise $bc$.
    \item \og $a$ divise $b$\fg{} équivaut à \og $\mbox{ppcm}(a,b)=\vert b\vert$\fg.
    \item Si $a$ divise $b$, alors $a$ n'est pas premier avec $b$.
    \item Si $a$ n'est pas premier avec $b$, alors $a$ divise $b$ ou $b$ 
divise $a$.
\end{enumerate}
\finenonce{000342}



\finexercice



\exercice{343, cousquer, 2003/10/01}

\enonce{000343}{}
\begin{enumerate}
  \item Soit $p \in \mathbb{Z}$ un nombre premier. 
Montrer que si $a\in \mathbb{Z}$ n'est pas congru à $0$ modulo $p$
 alors $p $ ne divise pas $a$ et donc $\mbox{pgcd}(a,p)=1$.
  \item Soit $a\in \mathbb{Z}$ non congru à $0$ modulo $p$ avec $p$ premier. 
 Montrer en utilisant
 le a) qu'il existe $u \in \mathbb{Z}$ non congru à $0$ modulo $p$ vérifiant
 $au \equiv 1[p]$. (Remarquer que cela donne un inverse de $a$ modulo $p$).
  \item Montrer que si $p$ n'est pas premier, 
il existe des éléments $a,u\in \mathbb{Z}$ non nuls modulo $p$
tels que $au\equiv 0[p]$.
\end{enumerate}
\finenonce{000343}



\finexercice



\exercice{344, cousquer, 2003/10/01}

\enonce{000344}{}
\begin{enumerate}
 \item Montrer que deux entiers non nuls consécutifs sont toujours premiers entre eux.
 \item Montrer que pour tout entier naturel $n$, $\mbox{pgcd}((n+1)^2,n+2)=1$.
\end{enumerate}
\finenonce{000344}



\finexercice



\exercice{345, cousquer, 2003/10/01}

\enonce{000345}{}
Prouver que pour vérifier qu'un entier $p$ est premier, il suffit de vérifier
qu'il n'a pas de diviseurs inférieurs ou égaux à $\sqrt{p}$.
\finenonce{000345}



\finexercice



\exercice{346, cousquer, 2003/10/01}

\enonce{000346}{Théorème de Wilson}
\noindent Démontrer que tout nombre premier $p$ divise $(p-1)!+1$.
\finenonce{000346}



\finexercice



\exercice{347, ridde, 1999/11/01}

\enonce{000347}{}
Montrer que les nombres suivants ne sont pas premiers :
\begin{enumerate}
\item $n^4-20n^2 + 4$ pour $n \in \Nn$.
\item $\frac14 (n^3 + (n + 2)^3)$ pour $n \geq 2$.
\item $a^4 + 4b^4$ pour $a, b \geq 2$.
\end{enumerate}
\finenonce{000347}



\finexercice



\exercice{348, ridde, 1999/11/01}
\video{x73t94ZLjSU}
\enonce{000348}{}
Soit $X$ l'ensemble des nombres premiers de la forme $4k + 3$ avec $k \in \Nn$.
\begin{enumerate}
\item Montrer que $X$ est non vide.
\item Montrer que le produit de nombres de la forme $4k + 1$ est encore de cette forme.
\item On suppose que $X$ est fini et on l'\'ecrit alors $X = \left\{
p_1, \ldots, p_n\right\}$.\\  Soit $a = 4p_1 p_2 \ldots p_n  - 1$. Montrer par l'absurde
que $a$ admet un diviseur premier de la forme $4k + 3$.
\item Montrer que ceci est impossible et donc que $X$ est infini.
\end{enumerate}

\finenonce{000348} 


\finexercice


\exercice{349, gourio, 2001/09/01}
\video{HuY48EkNke0}
\enonce{000349}{}
Soit $a\in \Nn  $ tel que $a^{n}+1 $ soit premier, montrer que $\exists
 k\in \Nn,n=2^{k}.$
Que penser de la conjecture : $\forall  n\in \Nn,2^{2^{n}}+1$ est
premier ?

\finenonce{000349} 


\finexercice


\exercice{350, gourio, 2001/09/01}

\enonce{000350}{}
Soit $n $ un nombre premier et $p\in \{1,...,n-1\}, $ montrer que $n $%
divise $C_{n}^{p}.$
\finenonce{000350}



\finexercice



\exercice{351, gourio, 2001/09/01}

\enonce{000351}{}
Soient $a$ et $b$ deux entiers sup\'{e}rieurs \`{a} 2 premiers entre eux,
montrer que :
$$\exists N_{0}\in \Nn,\forall n\geq N_{0},n\in \left\{ ax+by|(x,y)\in
\Nn^{2}\right\} . $$
\finenonce{000351}



\finexercice



\exercice{3130, quercia, 2010/03/08}

\enonce{003130}{pgcd $\times$ ppcm}
Soient $a,b,c \in \N^*$. Quand a-t-on
$\text{pgcd}(a,b,c) \times \text{ppcm}(a,b,c) = abc$ ?
\finenonce{003130}


\finexercice 


\exercice{3131, quercia, 2010/03/08}

\enonce{003131}{pgcd $\times$ ppcm}
Soient $a_1, \dots, a_n \in \N^{*}$ et $b_i = \prod_{j \ne i} a_j$.

Montrer que : \\$\text{pgcd}( a_1, \dots, a_n ) \times \text{ppcm}( b_1, \dots, b_n) =
       \text{ppcm}( a_1, \dots, a_n ) \times \text{pgcd}( b_1, \dots, b_n) =
       \prod a_i$.
\finenonce{003131}


\finexercice 


\exercice{3132, quercia, 2010/03/08}

\enonce{003132}{$ab$ est un carr{\'e} parfait}
Soient $a,b \in \N^*$ premiers entre eux tels que $ab$ est un carr{\'e} parfait.
Montrer que $a$ et $b$ sont des carr{\'e}s parfaits.
\finenonce{003132}



\finexercice 


\exercice{3133, quercia, 2010/03/08}

\enonce{003133}{$a^n = b^m$}
Soient $a,b \in \N^*$ et $m,n$ premiers entre eux tels que $a^n = b^m$.
Montrer qu'il existe $c \in \N^*$ tel que $a = c^m$ et $b = c^n$.
\finenonce{003133}



\finexercice 


\exercice{3134, quercia, 2010/03/08}

\enonce{003134}{Valuation 2-adique de $5^{2^n}-1$}
Montrer que la plus grande puissance de $2$ divisant $5^{(2^n)} - 1$ est $2^{n+2}$.
\finenonce{003134}


\finexercice 


\exercice{3135, quercia, 2010/03/08}

\enonce{003135}{$a^r-1$ premier~?}
On suppose que $a^r-1$ est un nombre premier. Montrez
que $r$ est premier, puis que $a$ vaut 2. R{\'e}ciproque ?
\finenonce{003135}


\finexercice 


\exercice{3136, quercia, 2010/03/08}

\enonce{003136}{Nombres de Mersenne}
On note $M_n = 2^n - 1$ ($n$-i{\`e}me nombre de Mersenne).
\begin{enumerate}
  \item Montrer que : $M_n$ est premier $ \Rightarrow $ $n$ est premier.
  \item V{\'e}rifier que $M_{11}$ n'est pas premier.
\end{enumerate}
\finenonce{003136}


\finexercice 


\exercice{3137, quercia, 2010/03/08}

\enonce{003137}{$a^n + 1$ est premier}
Soient $a,n \in \N$ tels que $a \ge 2$, $n \ge 1$, et $a^n + 1$ est premier.
Montrer que $n$ est une puissance de $2$.
\finenonce{003137}



\finexercice 


\exercice{3138, quercia, 2010/03/08}

\enonce{003138}{Nombre de diviseurs d'un nombre entier}
Pour $n \in \N^*$, on note $d_n$ le nombre de diviseurs positifs de $n$.
\begin{enumerate}
  \item   Montrer que si $n = ab$ avec $a \wedge b = 1$, alors $d_n$ = $d_a d_b$.
  \item   Montrer que $n$ est un carr{\'e} parfait si et seulement si $d_n$ est impair.
  \item   Montrer que : $\prod_{d|n}d = {\sqrt n}^{d_n}$.
\end{enumerate}
\finenonce{003138}



\finexercice 


\exercice{3139, quercia, 2010/03/08}

\enonce{003139}{Nombres premiers congrus {\`a} 3 modulo 4}
Montrer qu'il y a une infinit{\'e} de nombres premiers $p$ tels que $p \equiv -1 (\mathrm{mod}\, 4)$.
\finenonce{003139}



\finexercice 


\exercice{3140, quercia, 2010/03/08}

\enonce{003140}{Nombres premiers congrus {\`a} 1 modulo 4}
On rappelle que si $p$ est premier et $n\wedge p = 1$, alors
$n^{p-1} \equiv 1(\mathrm{mod}\, p)$.
\begin{enumerate}
  \item Soit $n\in \N$ et $p\ge 3$ un diviseur premier de $n^2+1$.
    Montrer que $p\equiv 1(\mathrm{mod}\, 4)$.
  \item En d{\'e}duire qu'il y a une infinit{\'e} de nombres premiers de la forme $4k+1$.
\end{enumerate}
\finenonce{003140}


\finexercice 


\exercice{3141, quercia, 2010/03/08}

\enonce{003141}{Intervalle sans nombres premiers}
Trouver 1000 entiers cons{\'e}cutifs non premiers.
\finenonce{003141}



\finexercice 


\exercice{3142, quercia, 2010/03/08}

\enonce{003142}{Factorisation de $1000!$}
Quelle est la plus grande puissance de 6 divisant $1000\,!$ ?
\finenonce{003142}


\finexercice 


\exercice{3143, quercia, 2010/03/08}

\enonce{003143}{$1/2 + 1/3 + \dots + 1/n$ n'est pas entier}
Soit $n \in \N$, $n \ge 2$.
Montrer que $x_n = 1 + \frac 12 + \frac 13 + \dots + \frac 1n$ est de la forme :
$\frac {p_n}{2q_n}$ avec $p_n,q_n \in \N^*$ et $p_n$ impair.
\finenonce{003143}


\finexercice 

\section{103.99 Autre}


\exercice{352, bodin, 1998/09/01}

\enonce{000352}{}
R\'esoudre en nombres entiers naturels l'\'equation :
$$(x+1)(y+2)=2xy.$$
\finenonce{000352}



\finexercice



\exercice{353, bodin, 1998/09/01}

\enonce{000353}{}
Montrer que $(0,0,0)$ est le seul triplet $(x,y,z)$ d'entiers naturels tels que l'on ait :
$$x^2+y^2=3z^2.$$
\finenonce{000353}



\finexercice



\exercice{354, cousquer, 2003/10/01}

\enonce{000354}{}
 Déterminer les solutions  des équations\,:
$$x^2-5x-11 \equiv 0\ \rm{mod}\ 17;\ \  \cos((n^2-8n+2)\pi/7)=1$$
\finenonce{000354}



\finexercice



\exercice{355, cousquer, 2003/10/01}

\enonce{000355}{}
Un groupe de $N\geq 2$ personnes se réunit. 
Montrer qu'au moins deux personnes ont serré le meme nombre de mains. On 
pourra séparer les deux cas suivants : soit tout le monde a serré au 
moins une main, soit il existe quelqu'un qui n'a serré aucune main.
\finenonce{000355}



\finexercice


\section{104.01 Forme cartésienne, forme polaire}


\exercice{1, bodin, 1998/09/01}
\video{g5vcj-DIa2M}
\enonce{000001}{}

 Mettre sous la forme $a+ib$ ($a,b \in \Rr$) les nombres :
$$ \frac{3+6i}{3-4i}  \quad ; \quad \left(\frac{1+i}{2-i}\right)^2 + \frac{3+6i}{3-4i} \quad;  \quad
\frac{2+5i}{1-i} + \frac{2-5i}{1+i}. $$

\finenonce{000001}



\finexercice



\exercice{2, bodin, 1998/09/01}

\enonce{000002}{}
\'Ecrire les nombres complexes suivants sous la forme $a+ib$ ($a,b \in \Rr$) :
$$ \frac{5+2i}{1-2i} \quad  ; \quad \left(-\frac{1}{2}+i\frac{\sqrt{3}}{2}\right)^3 \quad ;  \quad
\frac{(1+i)^9}{(1-i)^7}. $$
\finenonce{000002}



\finexercice



\exercice{3, cousquer, 2003/10/01}
\video{2kur52c78pA}
\enonce{000003}{}
 Écrire sous la forme $a+ib$ les nombres complexes suivants :
\begin{enumerate}
\item Nombre de module $2$ et d'argument $\pi /3$.
\item Nombre de module $3$ et d'argument $-\pi /8$.
\end{enumerate}

\finenonce{000003}




\finexercice



\exercice{4, liousse, 2003/10/01}

\enonce{000004}{}
  Placer dans le plan cart\'esien,  les points d'affixes suivantes :
$\begin{array}{lccr} 
z_1 =i, & z_2 =1+i, & z_3 = -2+2i, & z_4 = e^{-i\frac{\pi}{3}}.
\end{array}$
\finenonce{000004}



\finexercice



\exercice{5, liousse, 2003/10/01}

\enonce{000005}{}
Mettre chacun des nombres complexes suivants sous la forme $a+ib,$ $a\in\Rr$ 
et $b\in\Rr.$
$$\frac{-2}{1-i\sqrt{3}}\hbox{, }\frac{1}{(1+2i)(3-i)}\hbox{, }
\frac{1+2i}{1-2i},\hbox{ }\frac{2+5i}{1-i}+\frac{2-5i}{1+i}.$$
\finenonce{000005}



\finexercice



\exercice{6, liousse, 2003/10/01}

\enonce{000006}{}
\begin{enumerate}
\item
   Mettre sous forme 
trigonom\'etrique les nombres complexes suivants :
$z_1 =3+3i\hbox{, }z_2 =-1-\sqrt{3} i\hbox{, }z_3 = -\displaystyle\frac{4}{3}i\hbox{, }
z_4=-2\hbox{, }z_5 =e^{i\theta}+e^{2i\theta}.$
\item Calculer $({1+i\sqrt {3}\over 2})^{2000}$.
\end{enumerate}
\finenonce{000006}



\finexercice



\exercice{7, cousquer, 2003/10/01}

\enonce{000007}{}
 Effectuer les calculs suivants :
\begin{enumerate}
\item $(3+2i)(1-3i)$.
\item Produit du nombre complexe de module $2$ et d'argument $\pi /3$ par le
nombre complexe de module $3$ et d'argument $-5\pi /6$.
\item $\frac{3+2i}{1-3i}$.
\item Quotient du nombre complexe de module $2$ et d'argument $\pi /3$ par le
nombre complexe de module $3$ et d'argument $-5\pi /6$.
\end{enumerate}

\finenonce{000007}



\finexercice



\exercice{8, cousquer, 2003/10/01}

\enonce{000008}{}
 Calculer le module et l'argument des nombres complexes suivants, ainsi
que de leurs conjugués :
\begin{enumerate}
\item $1+i(1+\sqrt2)$.
\item $\sqrt{10+2\sqrt5}+i(1-\sqrt5)$.
\item $\frac{\tan\varphi -i}{\tan\varphi +i}$ où $\varphi$ est un angle donné.
\end{enumerate}

\finenonce{000008}



\finexercice



\exercice{9, bodin, 1998/09/01}

\enonce{000009}{}
 Repr\'esenter sous forme trigonom\'etrique les
nombres :
$$ 1+i \quad; \quad 1+i\sqrt{3} \quad; \quad \sqrt{3}+i \quad;\quad \frac{1+i\sqrt{3}}{\sqrt{3}-i}.$$
\finenonce{000009}



\finexercice



\exercice{10, bodin, 1998/09/01}

\enonce{000010}{}
 \'Etablir les \'egalit\'es suivantes :
\begin{enumerate}
    \item $(\cos (\pi/7) + i \sin (\pi/7))(\frac{1-i\sqrt{3}}{2})(1+i) = \sqrt{2}(\cos(5\pi/84)+i\sin(5\pi/84)),$
    \item $(1-i)(\cos (\pi/5) + i \sin (\pi/5))(\sqrt{3}-i) = 2\sqrt{2}(\cos(13\pi/60)-i\sin(13\pi/60)),$
    \item $\frac {\sqrt{2}(\cos (\pi/12) + i \sin (\pi/12))}{1+i} = \frac {\sqrt{3}-i}{2}.$
\end{enumerate}

\finenonce{000010}



\finexercice


\exercice{11, bodin, 1998/09/01}
\video{XzALEyZLQYc}
\enonce{000011}{}
Calculer le module et l'argument de $u =
\frac{\sqrt{6}-i\sqrt{2}}{2}$ et $v = 1 - i$. En d\'eduire le
module et l'argument de $w = \frac{u}{v}$.

\finenonce{000011}



\finexercice


\exercice{12, bodin, 1998/09/01}

\enonce{000012}{}
\'Ecrire sous la forme partie r\'eelle-partie imaginaire, puis sous la forme
module-argument le nombre complexe :
$$ \left(  \frac {1+i-\sqrt3(1-i)}{1+i}   \right)^2.$$
\finenonce{000012}



\finexercice

\exercice{13, bodin, 1998/09/01}
\video{-Kjd4SKjwEo}
\enonce{000013}{}
 D\'eterminer le module et l'argument des nombres
complexes :
$$e^{e^{i\alpha}} \quad \text{ et } \quad e^{i\theta}+e^{2i\theta}.$$

\finenonce{000013}




\finexercice



\exercice{14, bodin, 1998/09/01}

\enonce{000014}{}
 D\'eterminer le module et l'argument de
$\frac{1+i}{1-i}$. Calculer $(\frac{1+i}{1-i})^{32}$.

\finenonce{000014}



\finexercice

\exercice{15, cousquer, 2003/10/01}

\enonce{000015}{}
Calculer $Z=(1+i\sqrt 3)^{2000}$.
\finenonce{000015}



\finexercice

\exercice{16, bodin, 1998/09/01}

\enonce{000016}{}
Calculer $(1+i\sqrt{3})^5 + (1-i\sqrt{3})^5$ et $(1+i\sqrt{3})^5 - (1-i\sqrt{3})^5$.
\finenonce{000016}



\finexercice

\exercice{17, bodin, 1998/09/01}

\enonce{000017}{}
Calculer le module et l'argument de $z = \frac{1}{1+i\tan \alpha}$.
\finenonce{000017}



\finexercice

\exercice{18, bodin, 1998/09/01}

\enonce{000018}{}
Calculer les puissances $n$-i\`emes des nombres complexes :
$$ z_1 = \frac{1+i\sqrt{3}}{1+i} \quad ; \quad z_2=1+j \quad ; \quad z_3=\frac{1+i\tan \theta}{1-i\tan \theta}.$$
\finenonce{000018}



\finexercice


\exercice{19, bodin, 1998/09/01}

\enonce{000019}{}
Comment choisir l'entier naturel $n$ pour que $(\sqrt{3}+i)^n$ soit un r\'eel ?
un imaginaire ?
\finenonce{000019}



\finexercice


\exercice{20, bodin, 1998/09/01}
\video{8nXKgqMsucU}
\enonce{000020}{}
 Soit $z$ un nombre complexe de module $\rho$,
d'argument $\theta$, et soit $\overline{z}$ son conjugu\'e.
Calculer
$(z+\overline{z})(z^2+\overline{z}^2)\ldots(z^n+\overline{z}^n)$
en fonction de $\rho$ et $\theta$.

\finenonce{000020}


\finexercice

\exercice{21, bodin, 1998/09/01}

\enonce{000021}{partiel novembre 88}
Soient $\alpha$ et $\beta$ deux nombres r\'eels. Mettre le nombre complexe $z  =
e^{i\alpha}+e^{i\beta}$ sous forme trigonom\'etrique $z = \rho e^{i\gamma}$ (indication :
poser $u = \frac {\alpha + \beta}{2}$, $v = \frac {\alpha - \beta}{2}$).

En d\'eduire la valeur de
$$ \sum_{p=0}^{n}C_n^p \cos  [p\alpha + (n-p)\beta ].$$


\finenonce{000021}



\finexercice


\exercice{22, liousse, 2003/10/01}

\enonce{000022}{}
\'Ecrire l'expression 
$(1+\cos \phi+i\sin \phi)$ sous forme trigonom\'etrique. En d\'eduire 
l'expression de $(1+\cos \phi+i\sin \phi)^{n}.$
\finenonce{000022}



\finexercice


\exercice{23, ridde, 1999/11/01}

\enonce{000023}{}
Mettre sous forme trigonom\'etrique $1 + e^{i\theta}$ où $\theta \in ]-\pi,\pi [$.
Donner une interpr\'etation g\'eom\'etrique.

\finenonce{000023}



\finexercice

\exercice{24, ridde, 1999/11/01}

\enonce{000024}{}
Montrer que si $\left|z\right| \leq k <1$ alors $1-k \leq \left|1 + z\right| \leq 1 + k$.
Faire un dessin et montrer qu'il peut y avoir \'egalit\'e.
\finenonce{000024}



\finexercice

\exercice{25, ridde, 1999/11/01}

\enonce{000025}{}
Montrer alg\'ebriquement et g\'eom\'etriquement que si $\left|z\right| = 1$ alors
$\left|1 + z\right| \geq 1$ ou $\left|1 + z^2\right| \geq 1$.

\finenonce{000025}



\finexercice

\exercice{26, ridde, 1999/11/01}

\enonce{000026}{}
R\'esoudre l'\'equation $\exp (z) = \sqrt{3}  + 3i$.
\finenonce{000026}



\finexercice

\exercice{2924, quercia, 2010/03/08}

\enonce{002924}{$\sum z_i+z_j$}
\begin{enumerate}
  \item Soient $u,v \in \C$.
    Montrer que $|u+v| + |u-v| \ge |u| + |v|$, et d{\'e}terminer les cas d'{\'e}galit{\'e}.
  \item Soient $z_1,z_2,z_3,z_4 \in \C$. Montrer que
    $\sum_{k=1}^4|z_k| \le \sum_{k=1}^3\sum_{\ell=k+1}^4 |z_k+z_\ell|$.
\end{enumerate}
\finenonce{002924}



\finexercice 


\exercice{2927, quercia, 2010/03/08}

\enonce{002927}{}
Soient $a,b \in \mathbb{U}$ distincts et $z \in \C$.
On note $u = \frac{z+ab\overline{z}-a-b}{a-b}$.
Montrer que $u^2 \in \R$.

\finenonce{002927}


\finexercice

\exercice{5119, rouget, 2010/06/30}
\enonce{005119}{**IT}
Calculer de deux façons les racines carrées de $1+i$ et en déduire les valeurs exactes de $\cos\left(\frac{\pi}{8}\right)$
et $\sin\left(\frac{\pi}{8}\right)$.
\finenonce{005119}


\finexercice


\exercice{5127, rouget, 2010/06/30}
\enonce{005127}{**I}
Déterminer les complexes $z$ tels que $z$, $\frac{1}{z}$ et $z-1$ aient même module.
\finenonce{005127}


\finexercice

\exercice{5128, rouget, 2010/06/30}
\enonce{005128}{**I}
On note $U$ l'ensemble des nombres complexes de module $1$. Montrer que~:

$$\forall z\in\Cc,\;(z\in U\setminus\{-1\}\Leftrightarrow\exists x\in\Rr/\;z=\frac{1+ix}{1-ix}).$$

\finenonce{005128}


\finexercice


\exercice{5129, rouget, 2010/06/30}
\enonce{005129}{**IT}
Forme trigonométrique de $\frac{1+\cos\theta-i\sin\theta}{1-\cos\theta+i\sin\theta}$ et de
$\frac{1+e^{i\theta}}{1-e^{i\theta}}$.
\finenonce{005129}


\finexercice


\exercice{5130, rouget, 2010/06/30}
\enonce{005130}{*T}
Calculer $(1+i\sqrt{3})^9$.
\finenonce{005130}


\finexercice


\exercice{7209, megy, 2019/12/27}

\enonce{007209}{}
Résoudre sur $\C$ l'équation $z^2+\overline{z}-1=0$.
\finenonce{007209}



\finexercice


\exercice{7210, megy, 2019/12/27}

\enonce{007210}{}
Résoudre sur $\C$ l'équation $4z^2+8\lvert z\rvert -3=0$.
\finenonce{007210}



\finexercice


\exercice{7211, megy, 2019/12/27}

\enonce{007211}{}
Résoudre sur $\C$ l'équation $iz^2-2\bar{z}+2-i=0$.
\finenonce{007211}



\finexercice

\exercice{7518, mourougane, 2021/08/10}
\enonce{007518}{Parties réelles et imaginaires}
\begin{enumerate}
    \item Soit $z=a+ib$ un nombre complexe non nul.
    Déterminer la partie réelle et la partie imaginaire de l'inverse de $z$.
    \item Soit $t$ un nombre réel et $z=a+ib$ un nombre complexe.
    Déterminer la partie réelle et la partie imaginaire de $\frac{z-t}{z+t}$.
    \item Soit $z$ un nombre complexe. Démontrer que $z$ est réel si et seulement si $\overline{z}=z$.
\end{enumerate}
\finenonce{007518}
\finexercice


\exercice{7519, mourougane, 2021/08/10}
\enonce{007519}{Modules}
\begin{enumerate}
\item Soit $z$ et $w$ deux nombres complexes. Démontrer que \begin{eqnarray*}
                               |z+w|^2&=&|z|^2+2\Re (z\overline{w})+|w|^2\\
                               |z+w|^2+|z-w|^2&=&2(|z|^2+|w|^2)\\
                               |z+w|&\leq& |z|+|w|\\
                               ||z|-|w||&\leq&|z-w|.
                               \end{eqnarray*}

\item Montrer que l'ensemble $\Cc$ des nombres complexes muni de la distance définie par $$d(z,w):=|w-z|$$
est un espace métrique, c'est à dire que pour tous $(x,y,z)\in\Cc^3$,
\begin{enumerate}
 \item $d(x,y)\geq 0$
 \item $d(x,y)=0\iff x=y$
 \item $d(y,x)=d(x,y)$
 \item $d(x,z)\leq d(x,y)+d(y,z)$.
\end{enumerate}
\end{enumerate}
\finenonce{007519}
\finexercice


\exercice{7520, mourougane, 2021/08/10}
\enonce{007520}{Arguments}
\begin{enumerate}
\item Soit $\theta$ et $\theta'$ deux nombres réels.
En utilisant les formules d'addition pour le calcul de $\cos(\theta+\theta')$ et $\sin(\theta+\theta')$,
montrer que $$e^{i(\theta+\theta')}=e^{i\theta} e^{i\theta'}.$$
\item Soit $z$ et $z'$ deux nombres complexes. Déterminer $\arg (zz')$ en fonction de $\arg (z)$ et $\arg(z')$.
\item En déduire l'écriture polaire de $z^n$, puis ses parties réelles et imaginaires.
\item Déterminer $\cos (5\theta)$ en fonction de $\cos(\theta)$ et de $\sin(\theta)$.
\end{enumerate}
\finenonce{007520}
\finexercice

\section{104.02 Racine carrée, équation du second degré}


\exercice{27, bodin, 1998/09/01}
\video{BPzsFnvZypQ}
\enonce{000027}{}
 Calculer les racines carr\'ees de $1,\ i,\ 3+4i,\
8-6i,$ et $7+24i$.

\finenonce{000027}




\finexercice



\exercice{28, cousquer, 2003/10/01}

\enonce{000028}{}
 Trouver les racines carrées de 
$3-4i$ et de $24-10i$.

\finenonce{000028}



\finexercice

\exercice{29, bodin, 1998/09/01}
\video{Q5kK2BZBffI}
\enonce{000029}{}
\begin{enumerate}
    \item Calculer les racines carr\'ees de $\frac{1+i}{\sqrt{2}}$. En d\'eduire
les valeurs de $\cos(\pi/8)$ et $\sin(\pi/8)$.
    \item Calculer les valeurs de $\cos(\pi/12)$ et $\sin(\pi/12)$.
\end{enumerate}

\finenonce{000029}



\finexercice

\exercice{30, bodin, 1998/09/01}

\enonce{000030}{}
 Montrer que les solutions de $az^2+bz+c=0$ avec $a$,
$b$, $c$ r\'eels, sont r\'eelles ou conjugu\'ees.

\finenonce{000030}



\finexercice

\exercice{31, bodin, 1998/09/01}
\video{TC-3du7GQAQ}
\enonce{000031}{}
 R\'esoudre dans $\Cc$ les \'equations suivantes :
$$ z^2+z+1 = 0 \quad ; \quad z^2-(1+2i)z+i-1 = 0 \quad ; \quad z^2-\sqrt{3}z-i = 0 \quad ;$$
$$z^2-(5-14i)z-2(5i+12)=0 \  ; \  z^2-(3+4i)z-1+5i =0 \  ; \  4z^2-2z+1=0 \  ;$$
$$z^4+10z^2 +169=0 \quad ; \quad z^4+2z^2 +4=0.$$

\finenonce{000031}



\finexercice

\exercice{32, bodin, 1998/09/01}

\enonce{000032}{}
Trouver les racines complexes de l'\'equation suivante :
$$x^4-30x^2+289=0.$$
\finenonce{000032}



\finexercice

\exercice{33, bodin, 1998/09/01}

\enonce{000033}{}
Pour $z \in \Cc \setminus \{2i\}$, on pose
$$f(z)=\frac{2z-i}{z-2i}.$$
\begin{enumerate}
    \item R\'esoudre l'\'equation $z^2=i, \  z \in \Cc.$
    \item R\'esoudre l'\'equation $f(z)=z, \ z \in \Cc \setminus \{2i\}.$
\end{enumerate}
\finenonce{000033}



\finexercice

\exercice{34, liousse, 2003/10/01}

\enonce{000034}{}
 On note $j = e^{2 \pi \over 3}.$
\begin{enumerate}
\item Mettre $j$ et $j^2$ sous forme alg\'ebrique.
\item V\'erifier que $1+j+j^2 = 0$.
\item Factoriser le polyn\^ome $z^3-8i$.
\end{enumerate}
\finenonce{000034}



\finexercice


\exercice{35, liousse, 2003/10/01}

\enonce{000035}{}
\begin{enumerate}
\item Calculer les racines carr\'ees de $1+i$, $7+24i$, $i$, $5+12i$, 
$\frac{1+i\sqrt{3}}{\sqrt{3}+i}$.
\item   R\'esoudre les \'equations 
suivantes :
\begin{enumerate}
\item $z^2 + z +1 =0$ 
\item $z^2 + z -2 =0$
\item   $z^2 -(5-14i)z -2(5i+12) =0$
\item $z^2 + 4z +5 =0$
\item  $z^2 -(3+4i) z-1 +5i=0$
\item  $z^4 -(1-i)z^2 -i =0$
\item   $z^4 + 4z^3 + 6z^2 + 4z -15 =0$
\end{enumerate}
\end{enumerate}
\finenonce{000035}



\finexercice


\exercice{36, cousquer, 2003/10/01}
\enonce{000036}{}

R\'esoudre dans $\mathbb{C}$ les \'equations suivantes :
\begin{enumerate}
\item $z^2-(11-5i)z+24-27i=0$.
\item $z^3+3z-2i=0$.
\end{enumerate}

\finenonce{000036} 


\finexercice


\exercice{37, cousquer, 2003/10/01}

\enonce{000037}{}
On considère dans $\mathbb{C}$ l'équation $(E)$ suivante:
$$ z^2-\left(1+a\right)\left(1+i\right)z+\left(1+a^2\right)i=0,$$
où $a$ est un paramètre réel.
\begin{enumerate}
\item  Calculer en fonction de $a\in\mathbb{R}$ les solutions $z_1$ et $z_2$ de $(E)$
(indication: on pourra déterminer les racines carées complexes de
$-2i(1-a)^2$). 
\item  On désigne par $Z_1$ (resp. $Z_2$) les points du plan complexe
d'affixe $z_1$ (resp. $z_2$) et par $M$ le milieu de
$\left[Z_1,Z_2\right]$. Tracer la courbe du plan complexe décrite par
$M$ lorsque $a$ varie dans $\mathbb{R}$. 
\end{enumerate}
\finenonce{000037}



\finexercice



\exercice{38, liousse, 2003/10/01}

\enonce{000038}{}
\begin{enumerate}
\item
  
Pour $\alpha\in\Rr$, r\'esoudre 
dans $\Cc$ l'\'equation $z^2-2\cos(\alpha)z+1=0.$ 
En d\'eduire la forme trigonom\'etrique des solutions de l'\'equation :
$$z^{2n}-2\cos(\alpha)z^n+1=0,\hbox{ o\`u $n$ est un entier naturel non nul.}$$
 $$P_\alpha(z)=z^{2n}-2\cos(\alpha)z^n+1.$$
\begin{enumerate}
\item
 Justifier la factorisation suivante de $P_\alpha$ :
$$P_\alpha(z)=\left(z^2-2\cos\left(\frac{\alpha}{n}\right)+1\right)
\left(z^2-2\cos\left(\frac{\alpha}{n}+\frac{2\pi}{n}\right)+1\right)
\dots\left(z^2-2\cos
\left(\frac{\alpha}{n}+\frac{2(n-1)\pi}{n}\right)+1\right).$$
\item Prouver, \`a l'aide des nombres complexes par exemple, la formule suivante :
$$1-\cos\theta=2\sin^2\left(\frac{\theta}{2}\right),\quad\theta\in\Rr.$$
\item Calculer $P_\alpha(1)$. En d\'eduire 
$$\sin^2\left(\frac{\alpha}{2n}\right)\sin^2
\left(\frac{\alpha}{2n}+\frac{\pi}{n}\right)\dots\sin^2
\left(\frac{\alpha}{2n}+\frac{(n-1)\pi}{n}\right)=
\frac{\sin^2\left(\frac{\alpha}{2}\right)}{4^{n-1}}.$$
\end{enumerate}
\item
 Pour tout $\alpha$ appartenant \`a $]0,\pi[$, et 
pour tout entier naturel $n\geq 2$, on pose :
$$H_n(\alpha)=\sin\left(\frac{\alpha}{2n}+\frac{\pi}{2n}\right)
\sin\left(\frac{\alpha}{2n}+\frac{2\pi}{n}\right)\dots
\sin\left(\frac{\alpha}{2n}+\frac{(n-1)\pi}{n}\right).$$
\begin{enumerate}
\item
 Montrer que, pour tout $\alpha$ non nul, on a :
$$2^{n-1}H_n(\alpha)=\frac{\sin(\alpha/2)}{\sin(\alpha/2n)}.$$
\item Quelle est la limite de $H_n(\alpha)$ lorsque $\alpha$ tend vers $0$?\\
\item En d\'eduire que, pour tout entier naturel $n$ sup\'erieur ou \'egal 
\`a $2$, on a 
$$\sin\left(\frac{\pi}{n}\right)
\sin\left(\frac{2\pi}{n}\right)
\dots\sin\left(\frac{(n-1)\pi}{n}\right)=\frac{n}{2^{n-1}}.$$
\end{enumerate}
\end{enumerate}
\finenonce{000038}



\finexercice



\exercice{2945, quercia, 2010/03/08}

\enonce{002945}{Position des racines carr{\'e}es}
Soit $z \in \C$ et $p,q$ ses racines carr{\'e}es.
A quelle condition $z,p,q$ forment-ils un triangle rectangle en $z$~?
\finenonce{002945}



\finexercice 


\exercice{2946, quercia, 2010/03/08}

\enonce{002946}{{\'E}quations du second degr{\'e}}
R{\'e}soudre dans $\C$ : $z^4 - (5-14i)z^2 -2(5i+12) = 0$.
\finenonce{002946}


\finexercice 


\exercice{2947, quercia, 2010/03/08}

\enonce{002947}{Ensi P 91}   
R{\'e}soudre dans $\C$ : $z^4 + 6z^3 + 9z^2 + 100 = 0$.

\finenonce{002947}


\finexercice 


\exercice{2948, quercia, 2010/03/08}

\enonce{002948}{}
Comment faut-il choisir $m \in \C$ pour que l'{\'e}quation :
$z^2 - (2+im)z - (1+im) = 0$ admette deux racines imaginaires conjugu{\'e}es ?
\finenonce{002948}


\finexercice


\exercice{2949, quercia, 2010/03/08}

\enonce{002949}{}
\begin{enumerate}
   \item Soient $u,v \in \C$. V{\'e}rifier que
     $$\left(|u|^2-|v|^2\right)^2 =
     \left(\frac {|u+v|^2 + |u-v|^2}2\right)^2 - 4|uv|^2.$$
  \item Soient $\alpha,\beta \in \C$. CNS pour que les racines de
    $z^2 + \alpha z + \beta = 0$ aient m{\^e}me module ?
\end{enumerate}
\finenonce{002949}



\finexercice


\exercice{2950, quercia, 2010/03/08}

\enonce{002950}{Moyennes g{\'e}om{\'e}trique et arithm{\'e}tique}
\begin{enumerate}
  \item Soient $u,v \in \C$. Montrer que $|u+v|^2 + |u-v|^2 = 2|u|^2 + 2|v|^2$.
  \item Soient $\alpha,\beta \in \C$, $m = \frac {\alpha+\beta}2$ et
    $\mu$ une racine carr{\'e}e de $\alpha\beta$.
    Montrer que $|\alpha|+|\beta|  = |m+\mu| + |m-\mu|$.
    
\end{enumerate}
\finenonce{002950}


\finexercice 


\exercice{5120, rouget, 2010/06/30}
\enonce{005120}{**T}
Résoudre dans $\Cc$ les équations suivantes~:
\begin{enumerate}
\item  $z^2+z+1=0$
\item  $2z^2+2z+1=0$
\item  $z^2-2z\cos\theta+1=0$, $\theta$ réel donné.
\item  $z^2-(6+i)z+(11+13i)=0$
\item  $2z^2-(7+3i)z+(2+4i)=0$.
\end{enumerate}
\finenonce{005120}


\finexercice

\exercice{5125, rouget, 2010/06/30}
\enonce{005125}{**T}
Résoudre dans $\Cc$ l'équation $z^4-(5-14i)z^2-2(5i+12)=0$.

\finenonce{005125}


\finexercice

\exercice{7212, megy, 2019/12/27}

\enonce{007212}{}
Résoudre sur $\C$ l'équation
\[
iz^3-(1+i)z^2+(1-2i)z+6+8i=0,
\]
sachant qu'elle possède une solution réelle.
\finenonce{007212}



\finexercice


\exercice{7213, megy, 2019/12/27}

\enonce{007213}{}
Soit $a\in\C$. Résoudre sur $\C$ l'équation
\[z^2 -2(1-i) z +a^2 -2i=0.\]
Pour quelles valeurs de $a$ cette équation possède-t-elle au moins une racine réelle ?
\finenonce{007213}



\finexercice
\section{104.03 Racine n-ieme}


\exercice{39, liousse, 2003/10/01}

\enonce{000039}{}
\begin{enumerate}
\item
  Pour quelles valeurs 
de $z\in\Cc$ a-t-on $\vert 1+iz\vert=\vert 1-iz\vert.$\\
  On consid\`ere dans $\Cc$ l'\'equation 
$\left( \frac{1+iz}{1-iz}\right)^{n}=\frac{1+ia}{1-ia},$ o\`u $a\in\Rr.$ 
Montrer, sans les calculer, que les solutions de 
cette \'equation sont r\'eelles. Trouver alors les solutions.\\
  Calculer les racines cubiques 
de $\frac{\sqrt{3}+i}{\sqrt{3}-i}$.
\end{enumerate}
\finenonce{000039}



\finexercice



\exercice{40, cousquer, 2003/10/01}

\enonce{000040}{}
Pour tout nombre complexe $Z$, on pose $P(Z)=Z^4-1$.
\begin{enumerate}
\item  Factoriser $P(Z)$ et en déduire les solutions dans $\mathbb{C}$ de 
l'équation $P(Z)=0$.
\item  Déduire de 1. les solutions de l'équation d'inconnue $z$ :
$$((2z+1)/(z-1))^4=1$$
\end{enumerate}
\finenonce{000040}



\finexercice



\exercice{41, cousquer, 2003/10/01}

\enonce{000041}{}
Résoudre dans $\mathbb{C}$ l'équation suivante : $\quad z^4 =
\left(1-i\right)/\left(1+i\sqrt{3}\right).$
\finenonce{000041}



\finexercice



\exercice{42, cousquer, 2003/10/01}
\enonce{000042}{}

R\'esoudre dans $\Cc$ l'\'equation  $z^3 = \frac 14(-1+i)$
et montrer qu'une seule de ses solutions a une puissance quatri\`eme
r\'eelle.

\finenonce{000042} 


\finexercice


\exercice{43, cousquer, 2003/10/01}
\video{6B7bztuiRh8}
\enonce{000043}{}
 Trouver les racines cubiques de 
$2-2i$ et de $11+2i$.

\finenonce{000043} 


\finexercice


\exercice{44, cousquer, 2003/10/01}
\enonce{000044}{}

 Calculer $\frac{\frac{1+i\sqrt3}{2}}{\frac{\sqrt2(1+i)}{2}}$ alg\'ebriquement,
puis trigonom\'etriquement. En d\'eduire $\cos\frac{\pi}{12}$, 
$\sin\frac{\pi}{12}$, $\tan\frac{\pi}{12}$, $\tan\frac{5\pi}{12}$. 
R\'esoudre dans $\Cc$ l'\'equation $z^{24}=1$.
\finenonce{000044} 


\finexercice


\exercice{45, cousquer, 2003/10/01}

\enonce{000045}{}
 Trouver les racines quatrièmes de $81$ et de $-81$.

\finenonce{000045}



\finexercice



\exercice{46, cousquer, 2003/10/01}

\enonce{000046}{}
\begin{enumerate}
\item  Montrer que, pour tout $n\in\mathbb{N}^*$ et tout nombre $z\in\mathbb{C}$, on a:
$$\left(z-1\right)\left(1+z+z^2+...+z^{n-1}\right)=z^n-1,$$
et en déduire que, si $z\neq 1$, on a:
$$1+z+z^2+...+z^{n-1}=\frac{z^n-1}{z-1}.$$
\item  \label{Q2} Vérifier que pour tout $x\in\mathbb{R}$ , on a
$  
\exp(ix)-1=2i\exp\left(\frac{ix}{2}\right)\sin\left(\frac{x}{2}\right).$
\item  \label{Q3} Soit $n\in\mathbb{N}^*$. Calculer pour tout $x\in\mathbb{R}$ la somme:
$$Z_n=1+\exp(ix)+\exp(2ix)+...+\exp((n-1)ix),$$ 
et en déduire les valeurs de
\begin{eqnarray*}
X_n&=&1+\cos(x)+\cos(2x)+...+\cos((n-1)x)\\
Y_n&=&\sin(x)+\sin(2x)+...+\sin((n-1)x).
\end{eqnarray*} 
\end{enumerate}

\finenonce{000046}



\finexercice



\exercice{47, bodin, 1998/09/01}
\video{MJT-lI-lt-M}
\enonce{000047}{}

 Calculer la somme $S_n = 1+z+z^2+\cdots+z^n$.
\finenonce{000047} 


\finexercice


\exercice{48, bodin, 1998/09/01}
\video{b63m6T3TlsY}
\enonce{000048}{}

\begin{enumerate}
    \item R\'esoudre $z^3 = 1$ et montrer que les racines s'\'ecrivent $1$, $j$, $j^2$.
Calculer $1+j+j^2$ et en d\'eduire les racines de $1+z+z^2 =0$.
    \item R\'esoudre $z^n = 1$ et montrer que les racines s'\'ecrivent
 $1,\epsilon,\ldots,\epsilon^{n-1}$. En d\'eduire les racines de $1+z+z^2+\cdots+z^{n-1} =0$.
Calculer, pour $p \in \Nn$, $1+\epsilon^p+\epsilon^{2p}+\cdots+\epsilon^{(n-1)p}$.
\end{enumerate}
\finenonce{000048} 


\finexercice


\exercice{49, liousse, 2003/10/01}

\enonce{000049}{}
R\'esoudre dans  $\Cc$ :
\begin{enumerate}
\item $z^5 =1$.
\item $z^5 =1-i$.
\item $z^3 =-2+2i$.
\item $z^{5}=\bar{z}.$
\end{enumerate}
\finenonce{000049}



\finexercice



\exercice{50, bodin, 1998/09/01}

\enonce{000050}{}
\begin{enumerate}
    \item Calculer les racines $n$-i\`emes de $-i$ et de $1+i$.
    \item R\'esoudre $z^2-z+1-i=0$.
    \item En d\'eduire les racines de $z^{2n}-z^{n}+1-i=0$.
\end{enumerate}
\finenonce{000050}



\finexercice



\exercice{51, bodin, 1998/09/01}

\enonce{000051}{}
Soit $\epsilon$ une racine $n$-i\`eme de l'unit\'e ; calculer
$$ S = 1 +2\epsilon+3\epsilon^2+\cdots+n\epsilon^{n-1}.$$
\finenonce{000051}



\finexercice



\exercice{52, bodin, 1998/09/01}

\enonce{000052}{}
R\'esoudre, dans $\Cc$, l'\'equation $(z+1)^n=(z-1)^n$.
\finenonce{000052}



\finexercice



\exercice{53, bodin, 1998/09/01}

\enonce{000053}{}
R\'esoudre, dans $\Cc$, l'\'equation $z^n=\overline{z}$ o\`u $n \ge 1$.
\finenonce{000053}



\finexercice



\exercice{54, bodin, 1998/09/01}

\enonce{000054}{}
R\'esoudre les \'equations suivantes :
$$ z^6 = \frac{1+i\sqrt{3}}{1-i\sqrt{3}} \quad ; \quad z^4 = \frac{1-i}{1+i\sqrt{3}}.$$
\finenonce{000054}



\finexercice



\exercice{55, bodin, 1998/09/01}

\enonce{000055}{}
R\'esoudre $ z^6 +27 =0$. ($z \in \Cc$)
\finenonce{000055}



\finexercice



\exercice{56, bodin, 1998/09/01}
\video{sVzO_9JjHTM}
\enonce{000056}{}
\begin{enumerate}
        \item Soient $z_1$, $z_2$, $z_3$ trois nombres complexes distincts
        ayant le m\^eme cube.

        Exprimer $z_2$ et $z_3$ en fonction de $z_1$.

        \item Donner, sous forme polaire, les solutions dans $\Cc$ de :
        $$ z^6+(7-i)z^3 -8 -8i = 0.$$
        (Indication : poser $Z=z^3$ ; calculer $(9+i)^2$)
\end{enumerate}

\finenonce{000056} 


\finexercice


\exercice{57, ridde, 1999/11/01}

\enonce{000057}{}
R\'esoudre dans $\Cc$ l'\'equation $27 (z-1)^6 + (z + 1)^6 = 0$.
\finenonce{000057}



\finexercice



\exercice{58, ridde, 1999/11/01}

\enonce{000058}{}
D\'eterminer les racines quatri\`emes de $-7-24i$.
\finenonce{000058}



\finexercice



\exercice{59, liousse, 2003/10/01}

\enonce{000059}{}
  Soit 
$\beta \in \Cc$ tel que $\beta^7=1$ et $\beta \not=1$. Montrer 
$${\beta \over 1+ \beta^2} + {\beta^2 \over 1+ \beta^4} + {\beta^3 \over 1+ \beta^6}
=-2$$
\finenonce{000059}



\finexercice



\exercice{2939, quercia, 2010/03/08}

\enonce{002939}{Racines de l'unit{\'e}}
R{\'e}soudre :
\begin{enumerate}
  \item $(z+1)^n = (z-1)^n$.
  \item $(z+1)^n = z^n = 1$.
  \item $z^4 - z^3 + z^2 - z + 1 = 0$.
  \item $1 + 2z + 2z^2 + \dots + 2z^{n-1} + z^n = 0$.
  \item $\left(\frac{1+ix}{1-ix}\right)^n = \frac{1+i\tan a}{1-i\tan a}$.
  \item $\overline x = x^{n-1}$.
  \item $\left(\frac{z+1}{z-1}\right)^3 + \left(\frac{z-1}{z+1}\right)^3 = 0$.
\end{enumerate}
\finenonce{002939}



\finexercice 


\exercice{2940, quercia, 2010/03/08}

\enonce{002940}{Sommes sur les racines de l'unit{\'e}}
Soit $\omega = \exp\frac {2i\pi}n$. Calculer :
\begin{enumerate}
  \item $\sum_{k=0}^{n-1} (1+\omega^k)^n$.

  \item $\sum_{k=0}^{n-1}\sum_{\ell=k}^{n-1} C_\ell^k \omega^{k+\ell}$.


\end{enumerate}
\finenonce{002940}



\finexercice 


\exercice{2941, quercia, 2010/03/08}

\enonce{002941}{Somme des puissances $p$-{\`e}mes des racines de l'unit{\'e}}
Soient $n,p \in \N^*$ et $\mathbb{U}_n$ le groupe des racines $n$-{\`e}mes de 1.
\begin{enumerate}
  \item Calculer $\sum_{x \in \mathbb{U}_n} x^p$.

  \item Soit $P$ un polyn{\^o}me {\`a} coefficients complexes de degr{\'e} inf{\'e}rieur ou {\'e}gal
    {\`a} $n-1$ et $M = \max\{|P(x)|,\ x \in\mathbb{U}_n\}$.
    Montrer que tous les coefficients de $P$ sont born{\'e}s par $M$.

\end{enumerate}
\finenonce{002941}



\finexercice 


\exercice{2942, quercia, 2010/03/08}

\enonce{002942}{$\sum \omega^{k^2}$}
Soient $n \in \N^*$, $\omega = e^{2i\pi/n}$ et
$Z = \sum_{k=0}^{n-1}\,\omega^{k^2}$. On demande de calculer $|Z|^2$.
Pour cela $\ldots$
\begin{enumerate}
   \item {\'E}crire $|Z|^2$ comme une somme double.
  \item Regrouper les termes diagonalement en tenant compte de la p{\'e}riodicit{\'e} de
    la fonction $k \mapsto \omega^k$.
  \item Terminer le calcul.
\end{enumerate}
\finenonce{002942}


\finexercice 


\exercice{2943, quercia, 2010/03/08}

\enonce{002943}{ $e^{2i\pi/7}$}
Soit $z = \exp\frac {2i\pi}7$ et $u = z + z^2 + z^4$, $v = z^3 + z^5 + z^6$.
\begin{enumerate}
  \item Calculer $u+v$ et $u^2$.
  \item En d{\'e}duire $\sin\frac {2\pi}7 + \sin\frac {4\pi}7 + \sin\frac {8\pi}7$.
\end{enumerate}
\finenonce{002943}


\finexercice


\exercice{2944, quercia, 2010/03/08}

\enonce{002944}{Calcul de produit}
Simplifier $x = \prod_{p=2}^n\frac {p^3-1}{p^3+1}$ en utilisant $1,j,j^2$.
\finenonce{002944}


\finexercice 


\exercice{5122, rouget, 2010/06/30}
\enonce{005122}{***}
 Soit $\alpha\in\left]-\frac{\pi}{2},\frac{\pi}{2}\right[$ donné. Résoudre dans $\Cc$ l'équation
$\left(\frac{1+iz}{1-iz}\right)^3=\frac{1+i\tan\alpha}{1-i\tan\alpha}$.
\finenonce{005122}


\finexercice

\exercice{5126, rouget, 2010/06/30}
\enonce{005126}{**}
Résoudre dans $\Cc$ l'équation $(z^2+1)^n-(z-1)^{2n}=0$.
\finenonce{005126}


\finexercice


\exercice{5131, rouget, 2010/06/30}
\enonce{005131}{**T}
Déterminer les racines quatrièmes de $i$ et les racines sixièmes de $\frac{-4}{1+i\sqrt{3}}$.
\finenonce{005131}


\finexercice

\exercice{5135, rouget, 2010/06/30}
\enonce{005135}{**I}
On considère l'équation $(E)~:~(z-1)^n-(z+1)^n=0$ où $n$ est un entier naturel supérieur ou
égal à $2$ donné.

\begin{enumerate}
\item  Montrer que les solutions de $(E)$ sont imaginaires pures.
\item  Montrer que les solutions de $(E)$ sont deux à deux opposées.
\item  Résoudre $(E)$.
\end{enumerate}
\finenonce{005135}


\finexercice

\exercice{5313, rouget, 2010/07/04}
\enonce{005313}{***I}
Calculer $a_n=\prod_{k=1}^{n}\sin\frac{k\pi}{n}$, $b_n=\prod_{k=1}^{n}\cos(a+\frac{k\pi}{n})$ et $c_n=\prod_{k=1}^{n}\tan(a+\frac{k\pi}{n})$ en éliminant tous les cas particuliers concernant $a$.
\finenonce{005313}


\finexercice


\exercice{7521, mourougane, 2021/08/10}
\enonce{007521}{Racines}
\begin{enumerate}
\item Déterminer les racines carrées et les racines cubiques de $i$.
\item Déterminer les racines cubiques de $1$.
\item Déterminer les racines carrées de $\sqrt{3}+3i$.
\item Résoudre les équations $ z^2 + 2z -2+4i = 0 $ et $ z^6 - z^3 + 1 = 0$. 
\end{enumerate}
\finenonce{007521}
\finexercice
\section{104.04 Géométrie}


\exercice{60, bodin, 1998/09/01}
\video{iiReM2y_jKk}
\enonce{000060}{}
D\'eterminer l'ensemble des nombres complexes $z$
tels que :
\begin{enumerate}
    \item $\displaystyle{\left|\frac{z-3}{z-5}\right|=1},$
    \item $\displaystyle{\left|\frac{z-3}{z-5}\right|= \frac{\sqrt{2}}{2}}.$
\end{enumerate}

\finenonce{000060} 


\finexercice


\exercice{61, cousquer, 2003/10/01}

\enonce{000061}{}
\begin{enumerate}
\item  Résoudre dans $\mathbb{C}$ l'équation (1) $ (z-2)/(z-1)=i.$ On donnera la solution sous forme algébrique.
\item  Soit $M,A,$ et $B$ les points d'affixes respectives $z,1,2$. On suppose que
$M\neq A$ et que $ M\neq B$. Interpréter géométriquement le module et un argument de $(z-2)/(z-1)$ et
retrouver la solution de l'équation (1).
\end{enumerate}
\finenonce{000061}



\finexercice



\exercice{62, cousquer, 2003/10/01}

\enonce{000062}{}
Le plan $P$ est rapporté à un repère orthonormé et identifié à l'ensemble $\C$ des nombres complexes par
$$
M(x,y) \mapsto x+iy = z,
$$
où $z$ est appelé l'affixe de $M.$
Soit $f : P \mathrm{rg} P$ qui à tout point $M$ d'affixe $z$ associe $M'$ d'affixe $z' = \frac{z-i}{z+i}.$
\begin{enumerate}
\item Sur quel sous ensemble de $P$, $f$ est-elle définie?
\item Calculer $|z'|$ pour $z$ affixe d'un point $M$ situé dans le demi plan ouvert 
$$
H := \{ M(x,y) \in P \mid y > 0.\}?
$$
\item En déduire l'image par $f$ de $H.$
\end{enumerate}
\finenonce{000062}



\finexercice



\exercice{63, cousquer, 2003/10/01}

\enonce{000063}{}
 Le plan $P$ est rapporté à un repère orthonormé et on identifie $P$ à l'ensemble des nombres complexes $\C$ par 
$$
M(x,y) \mapsto x+iy = z,
$$
où $z$ est appelé l'affixe de $M.$
Soit $g : P \mathrm{rg} P$ qui à tout point $M$ d'fixe $z \not= -1$ associe $g(M)$ d'affixe $z' = \frac{1-z}{1+z}$.
\begin{enumerate}
\item Calculer $z' + \bar{z'}$ pour $|z| = 1$. 
\item En déduire l'image du cercle de rayon $1$ de centre $0$  privé du point de coordonnées $(-1,0)$ par l'application $g.$
\end{enumerate}
\finenonce{000063}



\finexercice



\exercice{64, cousquer, 2003/10/01}

\enonce{000064}{}
Soit $C$ la courbe d'équation $x^2-xy+y^2 = 0$ dans le plan $P$ rapporté à un repère orthonormé.
\begin{enumerate}
\item
La courbe $C$ a-t-elle des points d'intersection avec le rectangle ouvert $R$ dont les sommets sont :
\begin{eqnarray*}
  A &=& (-3,2)\\
B &=& (4,2)\\
C &=& (4,-1)\\
D &=& (-3,-1).\\
\end{eqnarray*}
\item
Même question pour le rectangle fermé $R'$ de sommets :
\begin{eqnarray*}
  A' &=& (-1,4)\\
B' &=& (2,4)\\
C' &=& (2,1)\\
D' &=& (-1,1).\\
\end{eqnarray*}
\end{enumerate}
\finenonce{000064}



\finexercice



\exercice{65, cousquer, 2003/10/01}
\enonce{000065}{}
\label{exo:compl}
 D\'eterminer par le calcul et g\'eom\'etriquement les nombres complexes $z$ tels
que $\Bigl\vert\frac{z-3}{z-5}\Bigr\vert=1$. G\'en\'eraliser pour
$\Bigl\vert\frac{z-a}{z-b}\Bigr\vert=1$.

\finenonce{000065} 


\finexercice


\exercice{66, cousquer, 2003/10/01}
\enonce{000066}{}
 D\'eterminer par le calcul et g\'eom\'etriquement les nombres complexes $z$ tels
que $\Bigl\vert\frac{z-3}{z-5}\Bigr\vert= k$ ($k>0$, $k\neq1$).
G\'en\'eraliser pour $\Bigl\vert\frac{z-a}{z-b}\Bigr\vert=k$.

\finenonce{000066} 


\finexercice


\exercice{67, cousquer, 2003/10/01}
\enonce{000067}{}
\begin{enumerate}
\item Soit $A$, $B$, $C$ trois points du plan complexe dont les affixes sont
respectivement $a$, $b$, $c$. On suppose que $a+jb+j^2c=0$ ; montrer que $ABC$
est un triangle \'equilat\'eral ($j$ et $j^2$ sont les racines cubiques complexes
de $1$ --- plus pr\'ecis\'ement $j=\frac{-1+i\sqrt3}{2}$). R\'eciproque ?
\item $ABC$ \'etant un triangle \'equilat\'eral direct du plan complexe, on construit
les triangles \'equilat\'eraux directs $BOD$ et $OCE$, ce qui d\'etermine les points
$D$ et $E$ ($O$ est l'origine du plan complexe). Quelle est la nature du
quadrilat\`ere $ADOE$ ? Comparer les triangles $OBC$, $DBA$ et $EAC$.
\end{enumerate}

\finenonce{000067} 


\finexercice


\exercice{68, cousquer, 2003/10/01}

\enonce{000068}{}
 Soit $H$ une hyperbole équilatère de centre $O$, et $M$ un point de $H$.
Montrer que le cercle de centre $M$ qui passe par le symétrique de $M$ par
rapport à $O$ recoupe $H$ en trois points qui sont les sommets d'un triangle
équilatéral.

\emph{Indications :} en choisissant un repère adéquat, $H$ a une équation du type
$xy=1$, autrement dit en identifiant le plan de $H$ au plan complexe, $z^2-\bar
z^2=4i$. En notant $a$ l'affixe de $M$, le cercle a pour équation $\vert
z-a\vert^2=4a\bar a$. On pose $Z=z-a$ et on élimine $\bar Z$ entre les
équations du cercle et de l'hyperbole. En divisant par $Z+2a$ pour éliminer la
solution déjà connue du symétrique de $M$, on obtient une équation du type
$Z^3-A=0$.
\finenonce{000068}



\finexercice



\exercice{69, bodin, 1998/09/01}
\video{MNGaSR9bknI}
\enonce{000069}{}

Montrer que pour $u,v \in \Cc$, on a $|u+v|^2+|u-v|^2=2(|u|^2+|v|^2).$
Donner une interprétation géométrique.
\finenonce{000069} 


\finexercice


\exercice{70, bodin, 1998/09/01}

\enonce{000070}{}
Soient $z,z' \in \Cc$ tels que $\text{Arg}(z)-\text{Arg}(z')=\frac{\pi}{2}$.
\begin{enumerate}
    \item Montrer que $z\overline{z'}+\overline{z}z' = 0$.
    \item Montrer que $|z+z'|^2=|z-z'|^2=|z|^2+|z'|^2$.
\end{enumerate}
\finenonce{000070}



\finexercice



\exercice{71, bodin, 1998/09/01}

\enonce{000071}{}
\begin{enumerate}
    \item D\'eterminer l'ensemble des points $M$ du plan complexe,
d'affixe $z$ tels que : $\overline{z}(z-1)=z^2(\overline{z}-1)$.
    \item D\'eterminer l'ensemble des points $M$ du plan complexe,
d'affixe $z$ tels que les images de $1$, $z$, $1+z^2$ soient align\'ees.
\end{enumerate}
\finenonce{000071}



\finexercice



\exercice{72, bodin, 1998/09/01}

\enonce{000072}{}
Soit $s=(1-z)(1-iz)$.
\begin{enumerate}
    \item D\'eterminer l'ensemble des images des nombres complexes $z$ tel que $s$ soit r\'eel.
    \item D\'eterminer l'ensemble des images des nombres complexes $z$ tel que $s$ soit imaginaire pur.


\end{enumerate}
\finenonce{000072}



\finexercice



\exercice{73, liousse, 2003/10/01}

\enonce{000073}{}
\begin{enumerate}
\item
  Soit $A$ un point du plan 
d'affixe $\alpha = a+ib$. D\'eterminer l'ensemble 
des points $M$ du plan dont l'affixe $z$ v\'erifie 
$|z|^2 = \alpha \bar{z} + \bar{\alpha} z.$
\item  Quelles conditions doivent v\'erifier les points $M_1$ 
et $M_2$ d'affixes $z_1$ et $z_2$ pour que $\frac{z_1}{z_2}$ soit 
r\'eel ?
\item   D\'eterminer les nombres 
complexes $z$ tels que les points du plan complexe d'affixes 
$z,$ $iz,$ $i$ forment un triangle \'equilat\'eral.  
\item Soit $z=a+ib$, mettre l'expression $\frac{z-1}{z+1}$ sous forme $A+iB$, 
. D\'eterminer l'ensemble des points du plan complexe d'affixe 
$z$  telle que l'argument de $\frac{z-1}{z+1}$ soit $\pi\over 2$.
\end{enumerate}
\finenonce{000073}



\finexercice



\exercice{74, gourio, 2001/09/01}

\enonce{000074}{}
D\'{e}terminer les nombres complexes $z $ tels que le triangle ayant pour
sommets les points d'affixes $z,z^{2},z^{3} $ soit rectangle au point
d'affixe $z$.
\finenonce{000074}



\finexercice



\exercice{75, gourio, 2001/09/01}

\enonce{000075}{}
D\'{e}terminer les nombres complexes $z\in \Cc^{*}$ tels que les points
d'affixes $z,\frac{1}{z}$ et $(1-z)$ soient sur un m\^{e}me cercle de centre
O.
\finenonce{000075}



\finexercice



\exercice{76, gourio, 2001/09/01}

\enonce{000076}{}
R\'{e}soudre dans $\Cc$ le syst\`{e}me:
$$\left| z-1\right| \leq 1,\left| z+1\right| \leq 1.$$
\finenonce{000076}



\finexercice



\exercice{77, bodin, 1998/09/01}
\video{wHlb0IsMB7Q}
\enonce{000077}{}

Soit $(A_{0},A_{1},A_{2},A_{3},A_{4})$ un pentagone r\'egulier. On note $O$ son centre et on
choisit un rep\`ere orthonorm\'e $(O,\overrightarrow{u},\overrightarrow{v})$ avec
$\overrightarrow{u}=\overrightarrow{OA_{0}}$, qui nous permet d'identifier le plan avec
l'ensemble des nombres complexes $\Cc$.
%dessin
$$
\includegraphics{../images/img000077-1}
$$

\begin{enumerate}
\item
Donner les affixes $\omega_{0},\ldots,\omega_{4}$ des points $A_{0},\ldots,A_{4}$. Montrer
que  $\omega_{k}={\omega_{1}}^k$ pour $ k\in\{0,1,2,3,4\}$. Montrer que
$1+\omega_{1}+\omega_{1}^2+\omega_{1}^3+\omega_{1}^4=0$.

\item
En d\'eduire que $\cos(\frac{2\pi}{5})$ est l'une des solutions de l'\'equation $4z^2+2z-1=0$.
En d\'eduire la valeur de $\cos(\frac{2\pi}{5})$.

\item
On consid\`ere le point $B$ d'affixe $-1$. Calculer la longueur $BA_{2}$ en fonction de
$\sin\frac{\pi}{10}$ puis de $\sqrt{5}$ (on remarquera que
$\sin\frac{\pi}{10}=\cos\frac{2\pi}{5}$).

\item
On consid\`ere le point $I$ d'affixe $\frac{i}{2}$, le cercle $\mathcal{C}$ de centre $I$ de
rayon $\frac{1}{2}$ et enfin le point $J$ d'intersection de $\mathcal{C}$ avec le segment
$[BI]$. Calculer la longueur $BI$ puis la longueur $BJ$.

\item
\textbf{Application:} Dessiner un pentagone r\'egulier \`a la r\`egle et au compas. Expliquer.
\end{enumerate}



\finenonce{000077} 


\finexercice


\exercice{2925, quercia, 2010/03/08}

\enonce{002925}{{\'E}quations affines}
\begin{enumerate}
  \item Montrer que toute droite du plan admet pour {\'e}quation complexe :
    $az + \overline{az} = b$ avec $a \in \C^*, b \in \R$.
  \item Soient $a,b,c \in \C$, $a,b$ non tous deux nuls. Discuter la nature
    de $E = \{ z \in \C$ tq $az + b\overline{z} = c \}$.
\end{enumerate}
\finenonce{002925}


\finexercice 


\exercice{2926, quercia, 2010/03/08}

\enonce{002926}{Transformation homographique}
Soit $f : {\C\setminus\{i\}} \to {\C\setminus\{1\}}, z \mapsto {\frac {z+i}{z-i}}$
\begin{enumerate}
  \item Montrer que $f$ est bijective.
  \item D{\'e}terminer $f(\R)$, $f(\mathbb{U}\setminus\{i\})$, $f(i\R\setminus\{i\})$.
\end{enumerate}
\finenonce{002926}



\finexercice 


\exercice{2928, quercia, 2010/03/08}

\enonce{002928}{Triangle {\'e}quilat{\'e}ral}
Soient $a,b,c \in \C$ distincts. Montrer que les propositions suivantes sont
{\'e}quivalentes :
\begin{enumerate}
  \item $\{a,b,c\}$ est un triangle {\'e}quilat{\'e}ral.
  \item $j$ ou $j^2$ est racine de $az^2 + bz + c = 0$.
  \item $a^2 + b^2 + c^2 = ab + ac + bc$.
  \item $\frac 1{a-b} + \frac 1{b-c} + \frac 1{c-a} = 0$.
\end{enumerate}
\finenonce{002928}



\finexercice 


\exercice{2929, quercia, 2010/03/08}

\enonce{002929}{Sommets d'un carr{\'e}}
Soient $a,b,c,d \in \C$ tels que $$\begin{cases}a+ib &= c+id\cr a+c &= b+d.\end{cases}$$

Que pouvez-vous dire des points d'affixes $a,b,c,d$ ?

En d{\'e}duire qu'il existe $z \in \C$ tel que
$(z-a)^4 = (z-b)^4 = (z-c)^4 = (z-d)^4$.


\finenonce{002929}


\finexercice 


\exercice{2930, quercia, 2010/03/08}

\enonce{002930}{Configuration de points}
D{\'e}terminer les nombres $z \in \C$ tels que $\dots$
\begin{enumerate}
  \item $z,z^2,z^4$ sont align{\'e}s.
  \item $1,z,z^2$ forment un triangle rectangle.
  \item $z,\frac 1z,-i$ sont align{\'e}s.
\end{enumerate}
\finenonce{002930}



\finexercice 


\exercice{2931, quercia, 2010/03/08}

\enonce{002931}{$a+b+c=1$}
Trouver $a,b,c \in \mathbb{U}$ tels que $\begin{cases} a+b+c = 1\cr abc = 1.\end{cases}$
\finenonce{002931}


\finexercice 


\exercice{2932, quercia, 2010/03/08}

\enonce{002932}{$u+v+w = 0$}
Soient $u,v,w$ trois complexes unitaires tels que $u+v+w = 0$.
Montrer que $u=jv=j^2w$ ou $u=jw=j^2v$.
\finenonce{002932}



\finexercice 


\exercice{2933, quercia, 2010/03/08}

\enonce{002933}{$z+1/z = 2$}
Trouver les complexes $z \in \C^*$ tels que $\left|z + \frac 1z\right| = 2$.
\finenonce{002933}


\finexercice 


\exercice{2934, quercia, 2010/03/08}

\enonce{002934}{Symétrique par rapport à une droite}
Les points $A,B,M$ ayant pour affixes $a,b,z$, calculer l'affixe du symétrique $M'$ 
de $M$ par rapport à la droite $(AB)$.
\finenonce{002934}



\finexercice 


\exercice{2935, quercia, 2010/03/08}

\enonce{002935}{Orthocentre}
Soient $a,b,c,d \in \C$ deux {\`a} deux distincts.
Montrer que si deux des rapports $\frac {d-a}{b-c}, \frac {d-b}{c-a},
\frac {d-c}{a-b}$ sont imaginaires purs, alors le troisi{\`e}me l'est aussi.
\finenonce{002935}


\finexercice 


\exercice{2936, quercia, 2010/03/08}

\enonce{002936}{Similitudes dans un triangle}
On donne un triangle $ABC$, un r{\'e}el positif $k$ et un angle $\theta$.
On note $S_M$ la similitude directe de centre $M$, de rapport $k$ et
d'angle $\theta$.
Soit $C_1$ d{\'e}duit de $C$ par $S_A$,
     $B_1$ d{\'e}duit de $B$ par $S_C$,
     $A_1$ d{\'e}duit de $A$ par $S_B$.
Montrer que les deux triangles $ABC$ et $A_1B_1C_1$ ont m{\^e}me centre de gravit{\'e}.
\finenonce{002936}



\finexercice 


\exercice{2937, quercia, 2010/03/08}

\enonce{002937}{Centre du cercle circonscrit}
Soient $a,b,c \in \C$, affixes de points $A,B,C$ non align{\'e}s.
Calculer l'affixe du centre du cercle circonscrit {\`a} $ABC$ en fonction
de $a,b,c$.
\finenonce{002937}


\finexercice 


\exercice{2938, quercia, 2010/03/08}

\enonce{002938}{Sph{\`e}re de $\R^3$}
Soient $u,v \in \C$ tels que $u+v \ne 0$.
On pose $x = \frac {1+uv}{u+v}$,
        $y = i\frac {1-uv}{u+v}$,
        $z = \frac {u-v}{u+v}$.
\begin{enumerate}
  \item CNS sur $u$ et $v$ pour que $x,y,z$ soient r{\'e}els ?
  \item On suppose cette condition r{\'e}alis{\'e}e. Montrer que le point $M(x,y,z)$ dans
    l'espace appartient {\`a} la sph{\`e}re de centre $O$ et de rayon 1.
  \item A-t-on ainsi tous les points de cette sph{\`e}re ?
\end{enumerate}
\finenonce{002938}


\finexercice 


\exercice{5121, rouget, 2010/06/30}
\enonce{005121}{**IT Une construction du pentagone régulier à la règle et au compas}
\begin{enumerate}
\item  On pose $z=e^{2i\pi/5}$ puis $a=z+z^4$ et $b=z^2+z^3$. Déterminer une équation du second degré dont les
solutions sont $a$ et $b$ et en déduire les valeurs exactes de $\cos\left(\frac{2\pi}{5}\right)$, $\sin\left(\frac{2\pi}{5}\right)$,
$\cos\left(\frac{4\pi}{5}\right)$, $\sin\left(\frac{4\pi}{5}\right)$, $\cos\left(\frac{\pi}{5}\right)$ et $\sin\left(\frac{\pi}{5}\right)$.

\item   Le cercle de centre $\Omega$ d'affixe $-\frac{1}{2}$ passant par le point $M$ d'affixe $i$ recoupe 
$(Ox)$ en deux points $I$ et $J$. Montrer que
$\overline{OI}+\overline{OJ}=\overline{OI}.\overline{OJ}=-1$ et en déduire une construction à la règle et au
compas, du pentagone régulier
inscrit dans le cercle de centre $O$ et de rayon $1$ dont un des sommets est le point d'affixe $1$.\rule[-1mm]{0mm}{0mm}

\item La diagonale $[AC]$ d'un pentagone régulier $(ABCDE)$ est recoupée par deux autres diagonales en deux points
$F$ et $G$. Calculer les rapports $\frac{AF}{AC}$ et $\frac{FG}{AF}$.

\end{enumerate}
\finenonce{005121}


\finexercice


\exercice{5123, rouget, 2010/06/30}
\enonce{005123}{****}
\begin{enumerate}
\item  Soit $(ABC)$ un triangle dont les longueurs des côtés $BC$, $CA$ et $AB$ sont notées respectivement $a$, $b$
et $c$. Soit $I$ le centre du cercle inscrit au triangle $(ABC)$. Montrer que $I=\mbox{bar}\{A(a),B(b),C(c)\}$.
\item  Déterminer $z$ complexe tel que $O$ soit le centre du cercle inscrit au triangle $(PQR)$ dont les sommets
ont pour affixes respectives $z$, $z^2$ et $z^3$.
\end{enumerate}
\finenonce{005123}


\finexercice


\exercice{5124, rouget, 2010/06/30}
\enonce{005124}{***I}
Soient $A$, $B$ et $C$ trois points du plan, deux à deux distincts, d'affixes respectives
$a$, $b$ et $c$. Montrer que~:~
\begin{align*}
ABC\;\mbox{équilatéral}&\Leftrightarrow j\;\mbox{ou}\;j^2\;\mbox{est racine de l'équation}\;
az^2+bz+c=0\\
 &\Leftrightarrow a^2+b^2+c^2=ab+ac+bc\Leftrightarrow\frac{1}{b-c}+\frac{1}{c-a}+\frac{1}{a-b}=0.
\end{align*}
\finenonce{005124}


\finexercice


\exercice{5133, rouget, 2010/06/30}
\enonce{005133}{**T}
 Pour $z\in\Cc\setminus\{1\}$, on pose $Z=\frac{1+z}{1-z}$. Déterminer et construire
l'ensemble des points $M$ d'affixes $z$ tels que
\begin{enumerate}
\item  $|Z|=1$.
\item  $|Z|=2$.
\item  $Z\in\Rr$.
\item  $Z\in i\Rr$.
\end{enumerate}
\finenonce{005133}


\finexercice


\exercice{5134, rouget, 2010/06/30}
\enonce{005134}{*T}
Nature et éléments caractéristiques de la transformation d'expression complexe~:

\begin{enumerate}
\item  $z'=z+3-i$
\item  $z'=2z+3$
\item  $z'=iz+1$
\item  $z'=(1-i)z+2+i$
\end{enumerate}

\finenonce{005134}


\finexercice


\exercice{7004, megy, 2016/04/26}

\enonce{007004}{Théorèmes de Thébault et de Van Aubel}
Soit $ABCD$ un quadrilatère convexe direct. On construit quatre carrés qui s'appuient extérieurement sur les côtés $[AB]$, $[BC]$, $[CD]$ et $[DA]$. Les centres respectifs de ces carrés sont notés $P$, $Q$, $R$ et $S$.
\begin{enumerate}
\item Montrer que dans le carré construit sur $[AB]$, on a $p=\frac{a-ib}{1-i}$. Démontrer des relations analogues pour les autres carrés.
\item Montrer le théorème de Van Aubel : $PQRS$ est un \emph{pseudo-carré}, c'est-à-dire que ses diagonales sont de même longueur et se croisent à angle droit. Pour cela, calculer $\frac{s-q}{r-p}$.
\item (Théorème de Thébault) Dans le cas particulier où $ABCD$ est un parallélogramme, montrer que $PQRS$ est un carré.
\end{enumerate}
\finenonce{007004}



\finexercice



\exercice{7005, megy, 2016/04/26}


\enonce{007005}{Point de Vecten}

Soit $ABC$ un triangle direct. On construit trois carrés qui s'appuient extérieurement sur les côtés $[AB]$, $[BC]$ et $[CA]$. Les centres respectifs de ces carrés sont notés $P$, $Q$ et $R$. Le but est de montrer que $(AQ)$, $(BR)$ et $(CP)$ sont concourantes. Le point de concours est appelé \emph{point de Vecten} du triangle.

\begin{enumerate}

\item Montrer que dans le carré construit sur $[AB]$, on a $p=\frac{a-ib}{1-i}$.

 Démontrer des relations analogues pour les autres carrés.

\item Montrer que $ABC$ et $PQR$ ont même centre de gravité.

\item Montrer que $(AQ)$ et $(PR)$ sont perpendiculaires. Conclure.

\end{enumerate}

\finenonce{007005}







\finexercice



\exercice{7006, megy, 2016/04/26}


\enonce{007006}{Théorème de Napoléon}

Soit $ABC$ un triangle direct. Soient $P, Q, R$ tels que $CBP$, $ACQ$ et $BAR$ soient des triangles équilatéraux directs. On note $U, V, W$ les centres de gravité respectifs de ces trois triangles équilatéraux. Montrer que $UVW$ est équilatéral, de même centre de gravité que $ABC$, en utilisant la caractérisation des triangles équilatéraux.

\finenonce{007006}






\finexercice



\exercice{7007, megy, 2016/04/26}
% source possible : http://culturemath.ens.fr/maths/pdf/geometrie/cercles.pdf

\enonce{007007}{Théorème de Ptolémée}
On admet le résultat suivant:\\
\emph{Quatre points distincts d'affixes $a, b, c, d$ sont cocycliques ou alignés (resp. cocycliques ou alignés dans cet ordre) si et seulement si leur birapport  \[
[a,b,c,d]:= \frac{(a-c)(b-d)}{(b-c)(a-d)}
\]
est un réel (resp. réel positif).}\\
Le but de l'exercice est de démontrer le théorème de Ptolémée dans sa version suivante: \\
{\bf Théorème} (Ptolémée) \emph{Soient $A$, $B$, $C$, $D$ quatre points du plan non alignés. Alors on a 
\[AC\cdot BD \leq AB\cdot CD + AD\cdot BC,\]
avec égalité si et seulement si $A, B, C, D$ sont cocycliques dans cet ordre.}

\begin{enumerate}
\item (\'Echauffement) Montrer que pour tous $x, y, z \in \C$, 
\[ |x|\cdot |y-z| \leq |y| \cdot |z-x| + |z|\cdot|x-y|.\]
\item Prouver le théorème si deux des points sont égaux.
\item Dans la suite on suppose les points distincts deux à deux. En utilisant les affixes $a, b, c, d$ des points, prouver l'inégalité.
\item \'Etudier le cas d'égalité  et conclure.
\end{enumerate}
\finenonce{007007}
\finexercice


\exercice{7008, megy, 2016/04/26}
% source possible: http://culturemath.ens.fr/maths/pdf/geometrie/cercles.pdf

\enonce{007008}{Théorème des quatre cercles de Miquel}
On admet le résultat suivant:\\
\emph{Quatre points distincts d'affixes $a, b, c, d$ sont cocycliques ou alignés  si et seulement si leur birapport  \[
[a,b,c,d]:= \frac{(a-c)(b-d)}{(b-c)(a-d)}
\]
est réel.}\\
Soient $\mathcal{C}_1 , \mathcal{C}_2 , \mathcal{C}_3$ et $\mathcal{C}_4$ quatre cercles du plan vérifiant la condition suivante :\\
$\mathcal{C}_1$ coupe $\mathcal{C}_2$ en deux points distincts $z_1$ et $w_1$, qui coupe $\mathcal{C}_3$ en deux points distincts $z_2$ et $w_2$, qui coupe $\mathcal{C}_4$ en deux points distincts $z_3$ et $w_3$, qui coupe $\mathcal{C}_1$ en deux points distincts $z_4$ et $w_4$.

On suppose les huit points ci-dessus tous \emph{distincts}.

\begin{enumerate}
\item Démontrer que 
\[
\frac{[z_1 , w_2 , z_2 , w_1] \cdot [z_3, w_4, z_4, w_3]}{[z_2,w_3,z_3,w_2] \cdot [z_4 , w_1 ,z_1 ,  w_4]}=[z_1,z_3,z_2,z_4] \cdot [w_1,w_3,w_2,w_4].
\]


\item En déduire que si $Z_1 , Z_2 , Z_3 , Z_4$ sont alignés ou cocycliques, alors il en est de même de $W_1 , W_2 , W_3 , W_4$. 
\end{enumerate} 
\finenonce{007008}
\finexercice


\exercice{7009, megy, 2016/04/29}
\enonce{007009}{}
  Déterminer une équation complexe de la droite
  \begin{enumerate}
    \item contenant les points d'affixes $i$ et $1+2i$;
    \item contenant le point d'affixe $1+i$ et de vecteur normal d'affixe $2+i$;
    \item contenant le point d'affixe $1+i$ et de vecteur directeur d'affixe $2+i$.
  \end{enumerate}
\finenonce{007009}
 
\finexercice


\exercice{7010, megy, 2016/04/29}
\enonce{007010}{}
Soient $A$ et $B$ deux points distincts d'affixes $a$ et $b$, et $\theta \in \R$. Déterminer l'ensemble des points $M$ d'affixe $z$ tels que 
\[ \mathrm{Arg} \frac{z-b}{z-a} \equiv \theta [\pi].\]
\finenonce{007010}



\finexercice

\exercice{7145, megy, 2017/05/01}
\enonce{007145}{}
Déterminer les éléments caractéristiques des transformations représentées par:
\begin{enumerate} 
\item $z\mapsto (1-i )z + i$;
\item $z\mapsto i \bar z + 1-i$;
\item $z\mapsto 2i\bar z +3$;
\item $ z\mapsto \bar z + 1$.
\end{enumerate}
\finenonce{007145}

    
\finexercice  


\exercice{7146, megy, 2017/05/01}
\enonce{007146}{}
\'Ecrire en coordonnée complexe :
\begin{enumerate}
\item la rotation d'angle $\pi/4$ et de centre d'affixe $2+3i$;
\item la réflexion d'axe d'équation  $y=2x+1$.
\end{enumerate}
\finenonce{007146}

\finexercice  


\exercice{7147, megy, 2017/05/01}
\enonce{007147}{}
\'Ecrire en coordonnée complexe  les deux similitudes (directe et indirecte) envoyant les points d'affixes $2$ et $3$ sur ceux d'affixes $i$ et $3i$ et trouver leurs éléments caractéristiques.

\finenonce{007147}


\finexercice  


\exercice{7148, megy, 2017/05/01}
\enonce{007148}{}
Soit $a\in \C^*$, et soit $f$ la similitude directe du plan représentée par $z\mapsto a^2z+a-1$.

Déterminer l'ensemble des paramètres $a$ pour lesquels $f$ est :
\begin{enumerate}
\item une translation; 
\item  une homothétie de rapport $-4$; 
\item une rotation d'angle $\pi/2$.
\end{enumerate}
\finenonce{007148}

     
\finexercice  


\exercice{7149, megy, 2017/05/01}
\enonce{007149}{}
Soit $ABC$ un triangle tel que $C$ soit l'image de $B$ par la rotation de centre $A$ et d'angle $\pi/2$.
Soit $s$ une similitude envoyant $A$ sur $B$ et $B$ sur $C$.
\begin{enumerate}
\item Que peut valoir $s(C)$ ?
\item On suppose que $s$ est directe. Déterminer son centre $\Omega$. On l'exprimera comme barycentre de $A$, $B$ et $C$.
\item Si la similitude est indirecte, déterminer son centre et son axe.
\end{enumerate}
\finenonce{007149}

     
\finexercice  


\exercice{7150, megy, 2017/05/01}
\enonce{007150}{}

Fixons un repère orthonormé direct du plan.
À quelle condition sur les réels $a$, $b$, $c$, $d$, $e$ et $f$ la transformation
$\begin{pmatrix}x\\y\end{pmatrix} \mapsto  \begin{pmatrix}ax+cy+e\\ bx+dy+f\end{pmatrix}$ est-elle une similitude directe ? Indirecte ?

Application : écrire en coordonnée complexe les applications

\[
\phi : \begin{pmatrix}x\\y\end{pmatrix} \mapsto  \begin{pmatrix}-2x-y-1\\ x-2y+1\end{pmatrix}
\text{ et }
\psi : \begin{pmatrix}x\\y\end{pmatrix} \mapsto  \begin{pmatrix}-x+y\sqrt 3\\ x\sqrt 3+y\end{pmatrix}
\]

\finenonce{007150}

     
\finexercice  


\exercice{7151, megy, 2017/05/01}
\enonce{007151}{}

%[d'après bac France, juin 2008]
Le plan est muni d'un repère orthonormé direct $(O,\vec u, \vec v)$. On note $A$ et $B$ les points d'affixes $z_A=1-i$ et $z_B = 7+\frac72 i$.

\begin{enumerate}

\item Soit $s$ l'application du plan dans lui-même qui envoie un point d'affixe $z$ sur celui d'affixe 
\[  \frac23 iz +\frac13-\frac53 i.\]
Déterminer la nature et les éléments caractéristiques de $s$. % centre $A$...
\item On note $B_0=B$ et pour tout $n\in \N$, 
on note $B_{n+1} = s(B_n)$.
\begin{enumerate}
\item Calculer la distance $AB_{n+1}$ en fonction de $AB_n$.
\item Déterminer le plus petit entier $N$ vérifiant la propriété suivante :  pour tout $n\geq N$, le point $B_n$ appartient au disque de centre $A$ et de rayon $10^{-2}$. On demande une formule exacte pour cet entier mais pas son écriture explicite en base $10$. 
\item Déterminer l'ensemble des entiers $n$ tels que les points $A$, $B$ et $B_n$ soient alignés.
\end{enumerate}
\end{enumerate} 

\finenonce{007151}

     
\finexercice  


\exercice{7163, megy, 2017/06/11}
\enonce{007163}{}
\begin{enumerate}
\item Soit $ABC$ un triangle. Montrer qu'il est équilatéral direct ssi les affixes (un repère orthonormé direct ayant été fixé) des sommets vérifient $a+bj+cj^2=0$.
\item (Notations réinitialisées) Soit $O$ un point du plan,  $\mathcal C$ un cercle de centre $O$, et $A$, $B$, $C$, $D$, $E$ et $F$ des points distincts de $\mathcal C$ vérifiant (dans $\R/2\pi\Z$) l'égalité:
\[ (\overrightarrow{OA},\overrightarrow{OB})
=(\overrightarrow{OC},\overrightarrow{OD})
=(\overrightarrow{OE},\overrightarrow{OF})
=\pi/3.\]

On note $M$ (resp. $N$, $P$) le milieu de $[BC]$ (resp. $[DE]$, $[FA]$). Montrer que $MNP$ est équilatéral direct.
\end{enumerate}
\begin{center}
\includegraphics{../images/img007163-1}
\end{center}
\finenonce{007163}


\finexercice  


\exercice{7164, megy, 2017/06/11}
\enonce{007164}{}
Soit $ABC$ un triangle équilatéral direct, et $M$ un point. On note $A'$ (resp. $B'$ et $C'$) le symétrique orthogonal de $M$ par rapport à la droite $(BC)$ (resp. $(CA)$ et $(AB)$). Le but de l'exercice est de démontrer que $ABC$ et $A'B'C'$ ont le même centre de gravité.
\begin{enumerate}
\item \'Ecrire en coordonnée complexe (relativement à un repère que l'on choisira judicieusement) la réflexion $\sigma_{AB}$ par rapport à l'axe $(AB)$. 
\item Conclure.
\end{enumerate}
\finenonce{007164}

\finexercice  


\exercice{7166, megy, 2017/07/08}
\enonce{007166}{}
% parallélogramme de Varignon, minimisation de périmètre
Soit $M_1M_2M_3M_4$ un parallélogramme direct du plan, de centre $O$, et $A$ un point quelconque du plan. On considère $B$ le symétrique de $A$ par rapport à $M_1$, $C$ le symétrique de $B$ par rapport à $M_2$, $D$ le symétrique de $C$ par rapport à $M_3$ et $E$ le symétrique de $D$ par rapport à $M_4$.
\begin{center}
\includegraphics{../images/img007166-1}
\end{center}
\begin{enumerate}
\item Montrer que $E=A$.
\item Montrer que si $z$ et $z'$ sont deux complexes, alors $|z+z'|+|z-z'|\geq 2|z|$.
\item On fixe un repère orthonormé direct de centre $O$. Exprimer $a$, $b$, $c$ et $d$ puis le périmètre de $ABCD$  en fonction de $m_1$, $m_2$ et de $t=a-m_1+m_2$.
\item On fait maintenant varier le point $A$. Montrer que  le périmètre du quadrilatère $ABCD$ est minimal lorsque  $AM_1OM_4$ est un parallélogramme.
\end{enumerate}
\finenonce{007166}
\finexercice

\section{104.05 Trigonométrie}


\exercice{78, cousquer, 2003/10/01}

\enonce{000078}{}
On rappelle la formule ($\theta\in\mathbb{R}$) :
$$ e^{i\theta}=\cos \theta + i\sin \theta.$$
\begin{enumerate}
\item  Etablir les formules d'Euler ($\theta \in \mathbb{R}$) :
\begin{eqnarray*}
  \cos \theta = \frac{e^{i\theta}+e^{-i\theta}}{2}\;\mbox{ et }\;
  \sin \theta = \frac{e^{i\theta}-e^{-i\theta}}{2i}.
\end{eqnarray*}

\item  En utilisant les formules d'Euler, linéariser (ou transformer de produit en
somme) ($a,\,b\in \mathbb{R}$) :
\begin{eqnarray*}
  2\cos a \cos b \;\;\; ; \;\;\; 2\sin a \sin b \;\;\; ; \;\;\; \cos^2 a
  \;\;\; ; \;\;\; \sin^2 a.
\end{eqnarray*}

\item  A l'aide de la formule : $e^{ix}e^{iy}=e^{i\left(x+y\right)}$
($x,\,y\in\mathbb{R}$), retrouver celles pour  $\sin(x+y)$,
$\cos(x+y)$ et $\tan(x+y)$ en fonction de sinus, cosinus et tangente
de $x$ ou de $y$; en déduire les formules de calcul pour $\sin(2x)$,
$\cos(2x)$ et $\tan (2x)$ ($x,\,y \in \mathbb{R}$).

\item  Calculer $\cos x$ et $\sin x$ en fonction de $\tan {\displaystyle \frac{x}{2}}$
($x \ne \pi + 2k\pi\,,\; k \in \mathbb{Z}$).

\item  Etablir la formule de Moivre ($\theta\in\mathbb{R}$) :
$$
(\cos \theta + i \sin \theta)^n=\cos(n\theta)+i\sin(n\theta).
$$

\item  En utilisant la formule de Moivre, calculer $\cos(3x)$ et $\sin(3x)$ en
fonction de $\sin x$ et $\cos x$.
\end{enumerate}
\finenonce{000078}



\finexercice



\exercice{79, bodin, 1998/09/01}

\enonce{000079}{}
\begin{enumerate}
    \item Calculer $\cos 5\theta$, $\cos 8\theta$, $\sin 6\theta$, $\sin 9\theta$, en fonction des lignes trigonom\'etriques de l'angle $\theta$.
    \item Calculer $\sin^3 \theta$, $\sin^4 \theta$, $\cos^5 \theta$, $\cos^6 \theta$,
\`a l'aide des lignes trigonom\'etriques des multiples entiers de $\theta$.
\end{enumerate}
\finenonce{000079}



\finexercice



\exercice{80, cousquer, 2003/10/01}
\video{yE-CGGcOrYA}
\enonce{000080}{}

 En utilisant les nombres complexes, calculer $\cos 5\theta$ et
$\sin5\theta$ en fonction de $\cos\theta$ et $\sin\theta$.
\finenonce{000080} 


\finexercice


\exercice{81, liousse, 2003/10/01}

\enonce{000081}{}
\begin{enumerate}
\item
  Soit $\theta \in \Rr$. 
A l'aide de la formule de Moivre exprimer en fonction 
de $\cos \theta$ et de $\sin \theta$~:\\
\begin{enumerate}
\item
 $\cos (2\theta)$ et $\sin (2\theta)$.\\
\item $\cos (3\theta)$ et $\sin (3\theta)$. En d\'eduire une \'equation 
du troisi\`eme degr\'e admettant pour solution $\cos(\frac{\pi}{3})$ et 
la r\'esoudre.
\end{enumerate}
\item Lin\'eariser les polynomes trigonom\'etriques suivants :
$1+\cos ^2 x$, $ \cos ^3 x +2\sin^2 x$.
\end{enumerate}
\finenonce{000081}



\finexercice



\exercice{82, ridde, 1999/11/01}

\enonce{000082}{}
Exprimer $ (\cos 5x) (\sin 3x)$ en fonction de $\sin x$ et $\cos x$.
\finenonce{000082}



\finexercice



\exercice{83, liousse, 2003/10/01}

\enonce{000083}{}
  Soit $x$ un nombre r\'eel. 
On note $C = 1 + \cos x + \cos 2x + \ldots + \cos nx 
= \sum_{k=0}^n \cos kx$, et $S = \sin x + \sin 2x + \ldots + \sin nx 
= \sum_{k=0}^n \sin kx$. Calculer $C$ et $S$. 
\finenonce{000083}



\finexercice



\exercice{84, cousquer, 2003/10/01}

\enonce{000084}{}
 Résoudre dans $\mathbb{R}$ les équations : 
$$\sin x=\frac{1}{2}\,,\;\cos x=-\frac{1}{2}\,,\;
\tan x=-1,$$
et placer sur le cercle trigonométrique les images des solutions; résoudre
dans $\mathbb{R}$ l'équation 
$$\cos (5x)=\cos \left(\frac{2\pi}{3}-x\right).$$
\finenonce{000084}



\finexercice



\exercice{85, cousquer, 2003/10/01}

\enonce{000085}{}
Calculer
$\sin (25 \pi/3), \cos (19 \pi/4), \tan (37 \pi/6).$
\finenonce{000085}



\finexercice



\exercice{86, cousquer, 2003/10/01}

\enonce{000086}{}
 Résoudre l'équation : $2\sin^2x-3\sin x-2=0$, puis l'inéquation :
$2\sin^2x-3\sin x-2>0$.
\finenonce{000086}



\finexercice



\exercice{87, cousquer, 2003/10/01}

\enonce{000087}{}
 Etudier le signe de la fonction donnée par $f(x)=\cos
3x+\cos 5x.$
\finenonce{000087}



\finexercice



\exercice{88, cousquer, 2003/10/01}

\enonce{000088}{}
 Simplifier, suivant la valeur de $x\in [-\pi,\pi]$, l'expression
$\sqrt {1+\cos x}+|\sin x/2|$.
\finenonce{000088}



\finexercice



\exercice{89, cousquer, 2003/10/01}

\enonce{000089}{}
 Résoudre dans $\mathbb{R}$ les équations suivantes:
(donner les valeurs des solutions appartenant à 
$\left]-\pi, \pi\right]$ et les placer sur le cercle trigonométrique). 
\begin{enumerate}
 \item  $ \sin\left(5x\right)=\sin\left(\frac{2\pi}{3}+x\right)$,
\item  $  \sin\left(2x-\frac{\pi}{3}\right)=\cos\left(\frac{x}{3}\right)$,
 \item  $\cos\left(3x\right)=\sin\left(x\right)$.
\end{enumerate}


\finenonce{000089}



\finexercice



\exercice{90, cousquer, 2003/10/01}

\enonce{000090}{}
A quelle condition sur le réel $m$ l'équation
$\sqrt{3}\cos(x)+\sin(x)=m$ a-t-elle une solution réelle ? Résoudre
cette équation pour $m=\sqrt{2}$.

\finenonce{000090}



\finexercice



\exercice{91, cousquer, 2003/10/01}

\enonce{000091}{}
Résoudre dans $\mathbb{R}$ les inéquations suivantes:
\begin{eqnarray*}
\cos(5x) + \cos(3x) \leq \cos(x) \\
2\cos^2(x) -9\cos(x)
+4 >0.
\end{eqnarray*}  

\finenonce{000091}



\finexercice



\exercice{92, cousquer, 2003/10/01}
\enonce{000092}{}

R\'esoudre dans $\mathbb{R}$ les \'equations suivantes: 
\begin{enumerate}
\item  $\cos^2(x)-\sin^2(x)=\sin(3x)$. 
\item  $\cos^4(x)-\sin^4(x)=1$. 
\end{enumerate}
\finenonce{000092} 


\finexercice


\exercice{2951, quercia, 2010/03/08}

\enonce{002951}{Somme de coefficients binomiaux}
A l'aide de formules du bin{\^o}me, simplifier :
\begin{enumerate}
  \item $\sum_{k=0}^{[n/3]}\,C_n^{3k}$.
  \item $\sum_{k=0}^{[n/2]}\,C_n^{2k}(-3)^k$.
  \item $\sum_{k=0}^n\,C_n^k\cos(k\theta)$.
  \item $\sum_{k=0}^n\,C_n^k\sin\bigl((k+1)\theta\bigr)$.
  \item $\cos a + C_n^1\cos(a+b) + C_n^2\cos(a+2b) + \dots + C_n^n\cos(a+nb)$.
\end{enumerate}
\finenonce{002951}


\finexercice 


\exercice{2952, quercia, 2010/03/08}

\enonce{002952}{Sommes trigonom{\'e}triques}
Simplifier :
\begin{enumerate}
  \item $\sum_{k=0}^n\,k\cos(k\theta)$.
  \item $\sum_{k=1}^n\, \sin^3(k\theta)$.
\end{enumerate}
\finenonce{002952}


\finexercice 


\exercice{2953, quercia, 2010/03/08}

\enonce{002953}{{\'E}quation trigonom{\'e}trique}
Soit $a \in \R$. R{\'e}soudre :
$$\begin{cases}\cos(a) + \cos(a+x) + \cos(a+y) = 0\cr
        \sin(a) + \sin(a+x) + \sin(a+y) = 0.\end{cases}$$
\finenonce{002953}


\finexercice 


\exercice{2954, quercia, 2010/03/08}

\enonce{002954}{$\sum \cos^{2p}(x+k\pi/2p)$}
Soit $\theta \in \R$.
\begin{enumerate}
  \item Simplifier $\cos^4\theta + \cos^4\left(\theta + \frac \pi4\right) +
                \cos^4\left(\theta + \frac {2\pi}4\right) +
                \cos^4\left(\theta + \frac {3\pi}4\right)$.
  \item Simplifier $\cos^6\theta + \cos^6\left(\theta + \frac \pi6\right) + \dots +
                \cos^6\left(\theta + \frac {5\pi}6\right)$.
  \item Simplifier $\cos^{2p}\theta + \cos^{2p}\left(\theta + \frac \pi{2p}\right)
                + \dots + \cos^{2p}\left(\theta + \frac {(2p-1)\pi}{2p}\right)$.
\end{enumerate}
\finenonce{002954}



\finexercice 


\exercice{2955, quercia, 2010/03/08}

\enonce{002955}{$\sum\cos(kx)/\cos x^k = 0$}
R{\'e}soudre : $\sum_{k=0}^{n-1}\,\frac {\cos(kx)}{\cos^kx} = 0$.
\finenonce{002955}



\finexercice 


\exercice{2956, quercia, 2010/03/08}

\enonce{002956}{$\sum C_n^kx^{n-k}\cos(k\alpha) = 0$}
R{\'e}soudre en $x$ :
$x^n + C_n^1x^{n-1}\cos\alpha + \dots + C_n^n\cos(n\alpha) = 0$.
\finenonce{002956}


\finexercice 


\exercice{2957, quercia, 2010/03/08}

\enonce{002957}{$\sum 2^{-k}/\cos\theta\dots\cos(2^k\theta)$}
Simplifier $$\sum_{k=1}^n\, \frac 1
{2^k \cos\theta \cos2\theta \cos4\theta \dots \cos2^{k-1}\theta}.$$
\finenonce{002957}


\finexercice 


\exercice{2958, quercia, 2010/03/08}

\enonce{002958}{Calcul de $\tan(nx)$}
Soit $n \in \N$, et $x \in \R$.
Exprimer $\tan(nx)$ en fonction de $\tan x$.
\finenonce{002958}



\finexercice 


\exercice{2959, quercia, 2010/03/08}

\enonce{002959}{$z = (1+ia)/(1-ia)$}
Soit $z \in \mathbb{U}$. Peut-on trouver $a \in \R$ tel que $z = \frac {1+ia}{1-ia}$ ?
\finenonce{002959}


\finexercice 


\exercice{5063, rouget, 2010/06/30}
\enonce{005063}{*IT}
Résoudre dans $\Rr$ puis dans $[0,2\pi]$ les équations suivantes~:
\begin{enumerate}
 \item $\sin x=0$,
 \item $\sin x=1$,
 \item $\sin x=-1$,
 \item $\cos x=1$,
 \item $\cos x=-1$,
 \item $\cos x=0$,
 \item $\tan x=0$,
 \item $\tan x=1$.
\end{enumerate}

\finenonce{005063}


\finexercice


\exercice{5064, rouget, 2010/06/30}
\enonce{005064}{*IT}
Résoudre dans $\Rr$ puis dans $[0,2\pi]$ les équations suivantes~:
\begin{enumerate}
 \item $\sin x=\frac{1}{2}$,
 \item $\sin x=-\frac{1}{\sqrt{2}}$,
 \item $\tan x=-1$,
 \item $\tan x=\frac{1}{\sqrt{3}}$,
 \item $\cos x=\frac{\sqrt{3}}{2}$,
 \item $\cos x=-\frac{1}{\sqrt{2}}$.
\end{enumerate}
\finenonce{005064}


\finexercice


\exercice{5065, rouget, 2010/06/30}
\enonce{005065}{**IT}
Résoudre dans $\Rr$ puis dans $I$ les équations suivantes~:

\begin{enumerate}
 \item $\sin(2x)=\frac{1}{2},\;I=[0,2\pi]$,
 \item $\sin\left(\frac{x}{2}\right)=-\frac{1}{\sqrt{2}},\;I=[0,4\pi]$,
 \item $\tan(5x)=1,\;I=[0,\pi]$,
 \item $\cos(2x)=\cos^2x,\;I=[0,2\pi]$,
 \item $2\cos^2 x-3\cos x+1=0,\;I=[0,2\pi]$,
 \item $\cos(nx)=0\;(n\in\Nn^*)$,
 \item $|\cos(nx)|=1$,
 \item $\sin(nx)=0$,
 \item $|\sin(nx)|=1$,
 \item $\sin x=\tan x,\;I=[0,2\pi]$,
 \item $\sin(2x)+\sin x=0,\;I=[0,2\pi]$,
 \item $12\cos^2x-8\sin^2x=2,\;I=[-\pi,\pi]$.
\end{enumerate}
\finenonce{005065}


\finexercice


\exercice{5066, rouget, 2010/06/30}
\enonce{005066}{**IT}
Résoudre dans $I$ les inéquations suivantes~:
\begin{enumerate}
 \item $\cos x\leq\frac{1}{2},\;I=[-\pi,\pi]$,
 \item $\sin x\geq-\frac{1}{\sqrt{2}},\;I=\Rr$,
 \item $\cos x>\cos\frac{x}{2},\;I=[0,2\pi]$,
 \item $\cos^2x\geq\cos(2x),\;I=[-\pi,\pi]$,
 \item $\cos^2x\leq\frac{1}{2},\;I=[0,2\pi]$,
 \item $\cos\frac{x}{3}\leq\sin\frac{x}{3},\;I=[0,2\pi]$.
\end{enumerate}
\finenonce{005066}


\finexercice


\exercice{5067, rouget, 2010/06/30}
\enonce{005067}{*I}
Calculer $\cos\frac{\pi}{8}$ et $\sin\frac{\pi}{8}$.
\finenonce{005067}


\finexercice


\exercice{5068, rouget, 2010/06/30}
\enonce{005068}{*I}
Calculer $\cos\frac{\pi}{12}$ et $\sin\frac{\pi}{12}$.
\finenonce{005068}


\finexercice


\exercice{5069, rouget, 2010/06/30}
\enonce{005069}{***}
Montrer que $\sum_{}^{}\cos\left(\pm a_1\pm a_2\pm...\pm
a_n\right)=2^n\cos a_1\cos a_2...\cos a_n$ (la somme comporte $2^n$ termes).
\finenonce{005069}


\finexercice


\exercice{5070, rouget, 2010/06/30}
\enonce{005070}{***I}
\begin{enumerate}
 \item  Calculer $\prod_{k=1}^{n}\cos\left(\frac{a}{2^k}\right)$ pour $a$ élément donné de $]0,\pi[$ (penser à $\sin(2x)=2\sin x\cos x$).

 \item  Déterminer $\lim_{n\rightarrow +\infty}\sum_{k=1}^{n}\ln\left(\cos(\frac{a}{2^k})\right)$.
\end{enumerate}
\finenonce{005070}


\finexercice


\exercice{5071, rouget, 2010/06/30}
\enonce{005071}{**}
Résoudre dans $\Rr$ l'équation $2^{4\cos^2x+1}+16.2^{4\sin^2x-3}=20$.
\finenonce{005071}


\finexercice


\exercice{5072, rouget, 2010/06/30}
\enonce{005072}{***}
 Soit $a$ un réel distinct de $\frac{1}{\sqrt{3}}$ et $-\frac{1}{\sqrt{3}}$.
\begin{enumerate}
 \item  Calculer $\tan(3\theta)$ en fonction de $\tan\theta$.
 \item  Résoudre dans $\Rr$ l'équation~:

$$\frac{3x-x^3}{1-3x^2}=\frac{3a-a^3}{1-3a^2}.$$
On trouvera deux méthodes, l'une algébrique et l'autre utilisant la formule de trigonométrie établie en 1).
\end{enumerate}
\finenonce{005072}


\finexercice


\exercice{5073, rouget, 2010/06/30}
\enonce{005073}{****}
On veut calculer $S=\tan9^\circ-\tan27^\circ-\tan63^\circ+\tan81^\circ$.\\
\begin{enumerate}
 \item  Calculer $\tan(5x)$ en fonction de $\tan x$.
 \item  En déduire un polynôme de degré $4$ dont les racines sont $\tan9^\circ$, $-\tan27^\circ$, $-\tan63^\circ$ et
$\tan81^\circ$ puis la valeur de $S$.
\end{enumerate}
\finenonce{005073}


\finexercice


\exercice{5074, rouget, 2010/06/30}
\enonce{005074}{***}
Combien l'équation
$$\tan x+\tan(2x)+\tan(3x)+\tan(4x)=0,$$
possède-t-elle de solutions dans $[0,\pi]$~?
\finenonce{005074}


\finexercice


\exercice{5075, rouget, 2010/06/30}
\enonce{005075}{**I}
On veut calculer $\cos\frac{2\pi}{5}$ et $\sin\frac{2\pi}{5}$. Pour cela, on pose
$a=2\cos\frac{2\pi}{5}$, $b=2\cos\frac{4\pi}{5}$ et $z=e^{2i\pi/5}$.
\begin{enumerate}
 \item  Vérifier que $a=z+z^4$ et $b=z^2+z^3$.
 \item  Vérifier que $1+z+z^2+z^3+z^4=0$.
 \item  En déduire un polynôme de degré $2$ dont les racines sont $a$ et $b$ puis les valeurs exactes de $\cos\frac{2\pi}{5}$ et $\sin\frac{2\pi}{5}$.
\end{enumerate}
\finenonce{005075}


\finexercice


\exercice{5076, rouget, 2010/06/30}
\enonce{005076}{**I}
Calculer une primitive de chacune des fonctions suivantes~:
\begin{enumerate}
 \item $x\mapsto\cos^2x$,
 \item $x\mapsto\cos^4x$,
 \item $x\mapsto\sin^4x$,
 \item $x\mapsto\cos^2x\sin^2x$,
 \item $x\mapsto\sin^6x$,
 \item $x\mapsto\cos x\sin^6x$,
 \item $x\mapsto\cos^5x\sin^2x$,
 \item $x\mapsto\cos^3x$.
\end{enumerate}
\finenonce{005076}


\finexercice


\exercice{5077, rouget, 2010/06/30}
\enonce{005077}{**}
Calculer $I=\int_{\pi/6}^{\pi/3}\cos^4x\sin^6x\;dx$ et
$J=\int_{\pi/6}^{\pi/3}\cos^4x\sin^7x\;dx$.
\finenonce{005077}


\finexercice


\exercice{5078, rouget, 2010/06/30}
\enonce{005078}{**}
Démontrer les identités suivantes, en précisant à chaque fois leur domaine de validité~:
\begin{enumerate}
 \item $\frac{1-\cos x}{\sin x}=\tan\frac{x}{2}$,
 \item $\sin\left(x-\frac{2\pi}{3}\right)+\sin x+\sin\left(x+\frac{2\pi}{3}\right)=0$,
 \item $\tan\left(\frac{\pi}{4}+x\right)+\tan\left(\frac{\pi}{4}-x\right)=\frac{2}{\cos(2x)}$,
 \item $\frac{1}{\tan x}-\tan x=\frac{2}{\tan(2x)}$.
\end{enumerate}

\finenonce{005078}


\finexercice


\exercice{5079, rouget, 2010/06/30}
\enonce{005079}{***}
Soit $k$ un réel distinct de $-1$ et de $1$.
\begin{enumerate}
 \item  Etudier les variations de $f_k~:~x\mapsto\frac{\sin x}{\sqrt{1-2k\cos x+k^2}}$.
 \item  Calculer $\int_{0}^{\pi}f_k(x)\;dx$.
\end{enumerate}
\finenonce{005079}


\finexercice


\exercice{5080, rouget, 2010/06/30}
\enonce{005080}{***I}
Calculer les sommes suivantes~:

\begin{enumerate}
 \item  $\sum_{k=0}^{n}\cos(kx)$ et $\sum_{k=0}^{n}\sin(kx)$, ($x\in\Rr$ et $n\in\Nn$ donnés).
 \item  $\sum_{k=0}^{n}\cos^2(kx)$ et $\sum_{k=0}^{n}\sin^2(kx)$, ($x\in\Rr$ et $n\in\Nn$ donnés).
 \item  $\sum_{k=0}^{n}\dbinom{n}{k}\cos(kx)$ et $\sum_{k=0}^{n}\dbinom{n}{k}\sin(kx)$, ($x\in\Rr$ et $n\in\Nn$
donnés).
\end{enumerate}
\finenonce{005080}


\finexercice


\exercice{5081, rouget, 2010/06/30}
\enonce{005081}{***}
Résoudre le système $\left\{
\begin{array}{l}
\cos a+\cos b+\cos c=0\\
\sin a+\sin b+\sin c=0
\end{array}\right.$ où $a$, $b$ et $c$ sont trois réels.
\finenonce{005081}


\finexercice


\exercice{5082, rouget, 2010/06/30}
\enonce{005082}{**}
Montrer que
$\cos^4\frac{\pi}{8}+\cos^4\frac{3\pi}{8}+\cos^4\frac{5\pi}{8}+\cos^4\frac{7\pi}{8}=\frac{3}{2}$.
\finenonce{005082}


\finexercice


\exercice{5083, rouget, 2010/06/30}
\enonce{005083}{***}
\begin{enumerate}
 \item  Résoudre dans $\Rr$ l'équation $\cos(3x)=\sin(2x)$.
 \item  En déduire les valeurs de $\sin x$ et $\cos x$ pour $x$ élément de
$\left\{\frac{\pi}{10},\frac{\pi}{5},\frac{3\pi}{10}\right\}$.
\end{enumerate}
\finenonce{005083}


\finexercice

\exercice{5162, rouget, 2010/06/30}
\enonce{005162}{***}
Montrer que $\forall n\in\Nn^*,\;\sum_{k=1}^{n}|\cos k|\geq\frac{n}{4}$ (remarquer que si $x\in[0;1],\;x^2\leq x$).
\finenonce{005162}


\finexercice


\exercice{7255, mourougane, 2021/08/10}
\enonce{007255}{}
Soit $\theta$ un nombre réel.
\begin{enumerate}
\item \`A l'aide des formules d'addition, calculer $\cos 2\theta$ et $\sin 2\theta$ en fonction de $\cos\theta$ et $\sin\theta$.
\item Vérifier la validité des formules obtenues pour $\theta=\pi/2$ et $\theta=\pi/3$.
\item Calculer $\cos 3\theta$ et $\sin 3\theta$ en fonction de $\cos\theta$ et $\sin\theta$.
\item Vérifier la validité des formules obtenues pour $\theta=\pi/2$ et $\theta=\pi/3$.
\end{enumerate}
\finenonce{007255}
\finexercice


\exercice{7256, mourougane, 2021/08/10}
\enonce{007256}{}
\begin{enumerate}

\item Exprimer $\cos(a)\cos(b)$ en fonction de $\cos(a+b)$ et $\cos(a-b)$.

\item En effectuant un changement de variables à préciser, montrer que pour tous réels $p$ et $q$ on a :

\[ \cos(p)+\cos(q)=2\cos\left(\frac{p+q}{2}\right)\cos\left(\frac{p-q}{2}\right). \]

\item En déduire les solutions de l'équation suivante :

\[ \cos(x)+\cos(2x)+\cos(3x)=0. \]

\end{enumerate}
\finenonce{007256}
\finexercice


\exercice{7257, mourougane, 2021/08/10}
\enonce{007257}{}
\begin{enumerate}

\item Résoudre dans $\Rr$ l'équation $\sqrt{3} \cos(x) + \sin(x) = \sqrt{2}$.

\item \`A l'aide une méthode similaire, résoudre l'équation $\cos(x) + \sin(x)=1$.

\end{enumerate}
\finenonce{007257}
\finexercice

\section{104.99 Autre}


\exercice{93, bodin, 1998/09/01}

\enonce{000093}{}
Montrer que tout nombre complexe $z$ non r\'eel de module $1$ peut se mettre sous
la forme $\frac{1+ir}{1-ir}$, o\`u $r \in \Rr$.
\finenonce{000093}



\finexercice



\exercice{94, bodin, 1998/09/01}

\enonce{000094}{}
Soit $u$, $v$ des nombres complexes non r\'eels tels que $|u|=|v|=1$ et $uv\not= -1$.
Montrer que $\frac{u+v}{1+uv}$ est r\'eel.
\finenonce{000094}



\finexercice



\exercice{95, bodin, 1998/09/01}

\enonce{000095}{}
Calculer les sommes suivantes :
$$\sum_{k=0}^{n}\cos(kx) \quad ; \quad \sum_{k=0}^{n}C_n^k\cos(kx).$$
\finenonce{000095}



\finexercice



\exercice{96, bodin, 1998/09/01}
\video{JdTvqdw6lRE}
\enonce{000096}{}
Soit $\Zz[i] = \{ a+ib \  ; \ a,b \in \Zz \}$.
\begin{enumerate}
        \item Montrer que si $\alpha$ et $\beta$ sont dans $\Zz[i]$ alors
        $\alpha + \beta$ et $\alpha\beta$ le sont aussi.

        \item Trouver les \'elements inversibles de $\Zz[i]$, c'est-\`a-dire les
        \'el\'ements $\alpha \in \Zz[i]$ tels qu'il existe $\beta \in \Zz[i]$ avec
        $\alpha\beta = 1$.

        \item V\'erifier que quel que soit $\omega \in \Cc$ il existe $\alpha \in \Zz[i]$
        tel que $|\omega - \alpha| < 1$.

        \item Montrer qu'il existe sur $\Zz[i]$ une division euclidienne,
 c'est-\`a-dire que,
        quels que soient $\alpha$ et $\beta$ dans $\Zz[i]$ il existe $q$ et $r$ dans
        $\Zz[i]$ v\'erifiant :
        $$ \alpha = \beta q + r \qquad \text{avec} \qquad |r| < |\beta|.$$
        (Indication : on pourra consid\'erer le complexe $\frac{\alpha}{\beta}$)

\end{enumerate}

\finenonce{000096} 


\finexercice


\exercice{97, ridde, 1999/11/01}

\enonce{000097}{}
Montrer que $\forall z \in \Cc$ $\dfrac{\left|\Re (z)\right| + \left|\Im (z)\right|}
{\sqrt 2}
 \leq \left|z\right| \leq \left|\Re (z)\right| + \left|\Im (z)\right| $. \'Etudier les cas
 d'\'egalit\'e.

\finenonce{000097}



\finexercice



\exercice{98, ridde, 1999/11/01}

\enonce{000098}{}
Soit $ (a, b, c, d) \in \Rr^4$ tel que $ad-bc = 1$ et $c \neq 0$. Montrer que
si $z \neq -\dfrac dc$ alors $\Im (\dfrac{az + b}{cz + d}) = \dfrac{\Im (z)}
{ \left|(cz + d)\right|^2}$.
\finenonce{000098}



\finexercice



\exercice{99, ridde, 1999/11/01}

\enonce{000099}{}
Que dire de trois complexes $a$, $b$, $c$ non nuls tels que $\left|a + b + c\right|
 = \left|a\right| + \left|b\right| + \left|c\right|$.
\finenonce{000099}



\finexercice



\exercice{100, gourio, 2001/09/01}

\enonce{000100}{}
\begin{enumerate}
  \item  \'Etudier la suite $(z_{n})_{n\in \Nn} $ d\'{e}finie par: $z_{0}=4,$ $%
z_{n+1}=f(z_{n})$  o\`{u} $f $ est l'application de $\Cc $ sur
lui-m\^{e}me d\'{e}finie par :
$$\forall z\in \Cc,f(z)=i+\frac{1}{4}(1-i\sqrt{3})z.$$
\emph{Indication }: on commencera par rechercher les coordonn\'{e}es
cart\'{e}siennes de l'unique point $\alpha $ tel que $f(\alpha )=\alpha $,
puis on s'int\'{e}ressera \`{a} la suite $(x_{n})_{n\in \Nn}$ d\'{e}finie par :
$$\forall n\in \Nn,x_{n}=z_{n}-\alpha .$$
  \item On pose $\forall n\in \Nn,l_{n}=\left| z_{n+1}-z_{n}\right|$. Calculer
$$\lim_{n\rightarrow \infty }\sum_{k=0}^{n}l_{k}$$
et interpr\'{e}ter g\'{e}om\'{e}triquement.
\end{enumerate}
\finenonce{000100}



\finexercice



\exercice{101, liousse, 2003/10/01}

\enonce{000101}{Examen octobre 1999}
On d\'efinit une fonction $f$ de $\mathbb C-\{i\}$ dans $\mathbb C-\{1\}$
en posant $$f(z)={z+i \over z-i}.$$
\begin{enumerate}
\item On suppose $z$ r\'eel. Quel est le module de $f(z)$~?
\item  Trouver les nombres complexes $z$ tels que $f(z)=z$.
\end{enumerate}
\finenonce{000101}



\finexercice



\exercice{102, liousse, 2003/10/01}

\enonce{000102}{Examen novembre 2001}

Soit  $f$ la fonction de $\mathbb C$ dans $\mathbb C$ d\'efinie
 par $f(z) ={1+z\over 1-z}$.
\begin{enumerate}
\item Calculer les points fixes de la fonction $f$, c'est \`a dire les nombres 
complexes $z$ tels que $f(z) =z$.
\item D\'eterminer les nombres complexes $z$ pour lesquels 
$f(z)$ est r\'eel.
\end{enumerate}
\finenonce{000102}



\finexercice



\exercice{103, cousquer, 2003/10/01}
\enonce{000103}{}

\begin{enumerate}
\item Montrer que si $x+y+z=a$, $yz+zx+xy=b$, $xyz=c$, alors $x$, $y$ et $z$
sont solutions de l'\'equation $Z^3-aZ^2+bZ-c=0$. Trouver $x$, $y$ et $z$ si on
suppose $a=b=0$ et $c=-8$.
\item R\'esoudre le syst\`eme
$$\left\{\begin{array}{rcl}
x+y+z & = & 4 \\
x^2+y^2+z^2 & = & 4 \\
x^3+y^3+z^3 & = & 1
\end{array}\right.$$
\end{enumerate}
\finenonce{000103} 


\finexercice


\exercice{5132, rouget, 2010/06/30}
\enonce{005132}{***}
Montrer que les solutions de l'équation $1+z+z^2+...+z^{n-1}-nz^n=0$ sont de module inférieur
ou égal à $1$.

\finenonce{005132}


\finexercice

\exercice{5136, rouget, 2010/06/30}
\enonce{005136}{***T ESIM 1993}
Pour $z\in\Cc$, on pose $\ch z=\frac{1}{2}(e^z+e^{-z})$, $\sh z=\frac{1}{2}(e^z-e^{-z})$ et $\tanh z=\frac{\sh z}{\ch
z}$.
\begin{enumerate}
\item  Quels sont les nombres complexes $z$ pour lesquels $\tanh z$ existe~?
\item  Résoudre dans $\Cc$ l'équation $\tanh z=0$.
\item  Résoudre dans $\Cc$ le système $\left\{
\begin{array}{l}
|\Im z|<\frac{\pi}{2}\\
|\tanh z|<1
\end{array}\right.$.
\item  Montrer que la fonction $\tanh$ réalise une bijection de $\Delta=\{z\in\Cc/\;|\Im z|<\frac{\pi}{4}\}$ sur
$U=\{z\in\Cc/\;|z|<1\}$.
\end{enumerate}
\finenonce{005136}


\finexercice

\section{105.01 Division euclidienne}


\exercice{356, bodin, 1998/09/01}

\enonce{000356}{}
Effectuer la division euclidienne du polyn\^ome $P=X^5-X^4+2X^3+X^2 +4$ par
$Q=X^2 -1$.
M\^eme exercice lorsque $P=X^4-2X\cos(2\varphi )+1$ et $Q= X^2 -2X\cos(\varphi )
+1$.
\finenonce{000356}



\finexercice



\exercice{357, bodin, 1998/09/01}

\enonce{000357}{}
 Soit $P$ un polyn\^ome. Sachant que le reste de la division
 euclidienne de $P$ par $X-a $ est 1 et celui de la division de $P$ par
$X-b $ est $-1$, $(a \not =b)$, quel est le reste de la division euclidienne de $P
$ par $(X-a)(X-b)$ ?
\finenonce{000357}



\finexercice



\exercice{358, bodin, 1998/09/01}

\enonce{000358}{}
Calculer le reste de la division euclidienne du polyn\^ome  %$X^n$ par
%le polyn\^ome $Q= X^2+1$
$X^n+X+1$ par le polyn\^ome $(X-1)^2$.
\finenonce{000358}



\finexercice



\exercice{359, bodin, 1998/09/01}

\enonce{000359}{}
 Pour quelles valeurs de $m $ le polyn\^ome $P=(X+1)^m -X^m-1$ est-il
divisible par le polyn\^ome $Q=X^2+X+1$ ?
\finenonce{000359}



\finexercice



\exercice{360, bodin, 1998/09/01}

\enonce{000360}{}
 Montrer que le polyn\^ome $P(X)-X$ divise le polyn\^ome $P(P(X))-X$.
\finenonce{000360}



\finexercice



\exercice{361, bodin, 1998/09/01}

\enonce{000361}{}
D\'eterminer $ a,b  \in \Zz$ de fa\c con \`a ce que le polyn\^ome $aX^{n+1}-bX^n+1$
soit divisible par le polyn\^ome $(X-1)^2$. Calculer alors le quotient des deux polyn\^omes.
\finenonce{000361}



\finexercice



\exercice{362, bodin, 1998/09/01}

\enonce{000362}{}
Existe-t-il un polyn\^ome $P$ de degr\'e 7 tel que  $(X-1)^4 $ divise
$P(X)+1$ et $(X+1)^4$ divise $P(X)-1$ ?
\finenonce{000362}



\finexercice



\exercice{363, bodin, 1998/09/01}

\enonce{000363}{}
Effectuer les divisions par puissances croissantes de :
\begin{enumerate}
    \item $P=1$ par $Q=1-X$, \`a l'ordre $n$,
    \item $P=1+X$ par $Q=1+X^2$ \`a l'ordre $5$,
    \item $P=X-\frac{X^3}{6}+\frac{X^5}{12}$ par $Q=1-2X^2 +X^4$ \`a l'ordre 5.
\end{enumerate}
\finenonce{000363}



\finexercice



\exercice{364, bodin, 1998/09/01}
\enonce{000364}{}

Effectuer les divisions euclidiennes de \\
$3X^5+4X^2+1 \ \text{ par }\  X^2+2X+3$,\\
$3X^5+2X^4-X^2+1 \  \text{ par }\  X^3+X+2$,\\
$X^4-X^3+X-2 \ \text{ par }\  X^2-2X+4$.
\finenonce{000364} 


\finexercice


\exercice{365, bodin, 1998/09/01}

\enonce{000365}{}
Dans $\Cc[X]$, effectuer les divisions euclidiennes de\\
$X^2-3iX-5(1+i) \ \text{ par }\  X-1+i$,\\
$4X^3+X^2$ \ \text{ par }\  $X+1+i$.
\finenonce{000365}



\finexercice



\exercice{366, bodin, 1998/09/01}
\enonce{000366}{}

Effectuer la division selon les puissances
croissantes de :
$$X^{4}+X^{3}-2X+1\text{ par }X^{2}+X+1 \text{ \`a l'ordre }2.$$
\finenonce{000366} 


\finexercice


\exercice{367, cousquer, 2003/10/01}

\enonce{000367}{}
 Soit $a$ et $b$ deux nombres complexes distincts, $m$ et~$n$ deux entiers
naturels. Montrer que si les polynômes $(X-a)^m$ et $(X-b)^n$ divisent un
polynôme $P$, alors le polynôme $(X-\nolinebreak a)^m(X-b)^n$ divise~$P$.
\finenonce{000367}



\finexercice



\exercice{368, cousquer, 2003/10/01}

\enonce{000368}{}
 Pour $n\in\mathbb{N}$, quel est le reste de la division de 
$X^n+X+b$ par $(X-a)^2$~?
\finenonce{000368}



\finexercice



\exercice{369, cousquer, 2003/10/01}

\enonce{000369}{}
 Pour $n\in \mathbb{N}$, montrer que le polynôme $(X-1)^{n+2}+X^{2n+1}$ est 
divisible par $X^2-X+1$. Trouver le quotient si $n=2$.
\finenonce{000369}



\finexercice



\exercice{370, cousquer, 2003/10/01}
\video{ndl3qcjN2vw} 
\enonce{000370}{}
Chercher tous les polynômes $P$ tels que $P+1$ soit divisible par $(X-1)^4$ et $P-1$ par $(X+1)^4$. 

\medskip
 

\emph{Indications.} 
Commencer par trouver une solution particulière $P_0$ avec l'une des méthode suivantes :
\begin{enumerate}
\item à partir de la relation de Bézout entre  $(X-1)^4$ et $(X+1)^4$;
\item en considérant le polynôme dérivé $P_0'$ et en cherchant un polynôme de degré minimal.
\end{enumerate}
Montrer que $P$ convient si et seulement si le polynôme $P-P_0$ est divisible par 
$(X-1)^4(X+1)^4$, et en déduire toutes les solutions du problème.
\finenonce{000370} 


\finexercice

\exercice{371, cousquer, 2003/10/01}
\enonce{000371}{}
  Effectuer la division de $A=X^6-2X^4+X^3+1$ par $B=X^3+X^2+1$~:
\begin{enumerate}
    \item Suivant les puissances d\'ecroissantes.
    \item \`A l'ordre~$4$ (c'est-\`a-dire tel que le reste soit divisible par 
$X^5$) suivant les puissances croissantes.
\end{enumerate}
\finenonce{000371} 
\finexercice


\exercice{372, ridde, 1999/11/01}

\enonce{000372}{}
D\'eterminer $a$ et $b$ dans $\Rr$ tels que $X^2 + 2$ divise $X^4 + X^3 + aX^2 +
bX + 2$.
\finenonce{000372}



\finexercice





\exercice{374, ridde, 1999/11/01}

\enonce{000374}{}
Soit $P$ un polynôme dont le reste de la division euclidienne par $X-1$ est $7$
et par $X + 5$ est $3$. Quel est le reste de la division euclidienne de $P$
par $X^2 + 4X -5$ ?
\finenonce{000374}



\finexercice



\exercice{375, ridde, 1999/11/01}
\enonce{000375}{}
 Effectuer la division euclidienne de
$X^5-7X^4-X^2-9X + 9$ par $X^2-5X + 4$.
\finenonce{000375} 
\finexercice


\exercice{376, ridde, 1999/11/01}

\enonce{000376}{}
Soit $n \geq 1$. D\'eterminer le reste de la division euclidienne de
$nX^{n + 1}- (n + 1)X^n + 1$ par $ (X-1)^2$.
\finenonce{000376}



\finexercice



\exercice{377, ortiz, 1999/04/01}

\enonce{000377}{}
Soient $P,Q\in K[X]$ tels que $X^2+X+1$ divise
$P(X^3)+XQ(X^3)$. Montrer que $P(1)=Q(1)=0$.
R\'eciproque ?
\finenonce{000377}



\finexercice



\exercice{378, gourio, 2001/09/01}
\video{ypZfDoacMUs}
\enonce{000378}{}
Quels sont les polynômes $P\in\Cc[X]$ tels que $P'$ divise $P$?
\finenonce{000378} 


\finexercice


\exercice{3196, quercia, 2010/03/08}

\enonce{003196}{D{\'e}composition en puissances croissantes}
Soit $A \in { K[X]}$ de degr{\'e} $> 0$.
Montrer que pour tout polyn{\^o}me $P \in { K_n[X]}$, il existe des polyn{\^o}mes $P_0,P_1,\dots,
P_n$ uniques v{\'e}rifiant :
$$\begin{cases}\deg P_i < \deg A\cr
P = P_0 + P_1A + \dots + P_nA^n.\end{cases}$$
\finenonce{003196}



\finexercice 


\exercice{3197, quercia, 2010/03/08}

\enonce{003197}{Lin{\'e}arit{\'e} du reste et du quotient}
Soit $B \in { K[X]}$ de degr{\'e} $n > 0$. On consid{\`e}re les applications :
$$\Phi : { K[X]} \to { K_{n-1}[X]}, P \mapsto R$$
et
 $$\Psi : { K[X]} \to { K[X]}, P \mapsto Q\quad
  \text{ avec }P = QB + R.$$
\begin{enumerate}
  \item Montrer que $\Phi$ et $\Psi$ sont lin{\'e}aires. Chercher leurs noyaux et leurs images.
  \item Simplifier $\Phi(P_1P_2)$.
\end{enumerate}
\finenonce{003197}



\finexercice 


\exercice{3198, quercia, 2010/03/08}

\enonce{003198}{Endomorphisme $P  \mapsto AP \bmod B$}
Soit $E =  K_3[X],\ A = X^4-1,\ B = X^4-X,$
et $\varphi : E \to E, P \mapsto
    {\text{reste de la div. euclid. de } AP \text{ par } B.}$

Chercher $\mathrm{Ker}\varphi$, $\Im\varphi$.

\finenonce{003198}


\finexercice 


\exercice{3199, quercia, 2010/03/08}

\enonce{003199}{Congruences}
Soient $P \in { K[X]}$, $a,b\in K$ distincts, et $\alpha = P(a)$, $\beta = P(b)$.
\begin{enumerate}
  \item Quel est le reste de la division euclidienne de $P$ par $(X-a)(X-b)$ ?
    
  \item Trouver le reste de la division euclidienne de
    $(\cos\theta + X\sin\theta)^n$ par $X^2+1$.
    
\end{enumerate}
\finenonce{003199}


\finexercice 


\exercice{3200, quercia, 2010/03/08}

\enonce{003200}{Congruences}
D{\'e}terminer les polyn{\^o}mes $P \in \Q_3[X]$ divisibles par $X+1$ et
dont les restes des divisions par $X+2,X+3,X+4$ sont {\'e}gaux.
\finenonce{003200}


\finexercice 


\exercice{3201, quercia, 2010/03/08}

\enonce{003201}{Calcul de pgcd}
Calculer le pgcd de $P$ et $Q$ pour :
\begin{enumerate}
  \item $P=X^4 + X^3 - 3X^2 - 4X - 1$\par
    $Q=X^3 + X^2 - X - 1$
    
  \item $P=X^4 - 10X^2 + 1$\par
    $Q=X^4 -4X^3 + 6X^2 - 4X + 1$
    
  \item $P=X^5 - iX^4 + X^3 - X^2 + iX - 1$\par
    $Q=X^4 - iX^3 + 3X^2 -2iX + 2$
    
\end{enumerate}
\finenonce{003201}


\finexercice 


\exercice{3202, quercia, 2010/03/08}

\enonce{003202}{Coefficients de B{\'e}zout}
Montrer que les polyn{\^o}mes $P$ et $Q$ suivants sont premiers entre eux.
Trouver $U,V \in { K[X]}$ tels que ${UP + VQ = 1}$.
\begin{enumerate}
  \item $P=X^4 + X^3 -2X +1$\par
    $Q=X^2 + X + 1$
    
  \item $P=X^3 + X^2 +1$\par
    $Q=X^3 + X + 1$
\end{enumerate}
\finenonce{003202}


\finexercice 


\exercice{3203, quercia, 2010/03/08}

\enonce{003203}{Division de $(X+1)^n - X^n - 1$ par $X^2 + X + 1$}
Chercher le reste de la division euclidienne de
$(X+1)^n - X^n - 1$ par $X^2 + X + 1$.
\finenonce{003203}


\finexercice 


\exercice{3204, quercia, 2010/03/08}

\enonce{003204}{Ensi P 90}
    Pour quels $n\in \N$ le polyn{\^o}me $(1+X^4)^n - X^n$ est-il divisible par
    $1+X+X^2$ dans $\R[X]$ ?
\finenonce{003204}


\finexercice 


\exercice{3205, quercia, 2010/03/08}

\enonce{003205}{Division de $(X-2)^2n + (X-1)^n - 1$ par $(X-1)(X-2)$}
Soit $P_n = (X-2)^{2n} + (X-1)^n - 1$.
\begin{enumerate}
  \item Montrer que $P_n$ est divisible par $X-1$ et par $X-2$.
    On note $Q_1$ et $Q_2$ les quotients correspondant.
  \item Montrer que $P_n$ est divisible par $(X-1)(X-2)$ et que le quotient est $Q_2-Q_1$.
  \item Montrer que ce quotient est {\'e}gal {\`a} :
    $$\Bigl((X-2)^{2n-2} - (X-2)^{2n-3} + \dots - (X-2) + 1\Bigr)
    + \Bigl((X-1)^{n-2} + (X-1)^{n-3} + \dots + (X-1) + 1\Bigr).$$
    
\end{enumerate}
\finenonce{003205}


\finexercice 


\exercice{3206, quercia, 2010/03/08}

\enonce{003206}{Calcul de restes}
Trouver les restes des divisions euclidiennes :
\begin{enumerate}
  \item de $X^{50}$ par $X^2-3X+2$.
    
  \item de $\bigl(X+\sqrt3\bigr)^{17}$ par $X^2+1$.
    
  \item de $X^8 - 32X^2 + 48$ par $\bigl(X-\sqrt2\bigr)^3$.
    
\end{enumerate}
\finenonce{003206}


\finexercice 


\exercice{3207, quercia, 2010/03/08}

\enonce{003207}{Divisibilit{\'e}}
Trouver $\lambda, \mu \in \C$ tels que $X^2+X+1$ divise
$X^5 + \lambda X^3 + \mu X^2 + 1$.
\finenonce{003207}


\finexercice 


\exercice{3208, quercia, 2010/03/08}

\enonce{003208}{Congruences}
Soit $P \in { K[X]}$ tel que les restes des divisions de $P$ par $X^2+1$
et $X^2-1$ valent respectivement $2X-2$ et $-4X$.
Quel est le reste de la division de $P$ par $X^4-1$ ?
\finenonce{003208}


\finexercice 


\exercice{3209, quercia, 2010/03/08}

\enonce{003209}{pgcd$(X^n-1,X^m-1)$}
Soient $m,n \in \N^*$. Chercher pgcd$( X^n-1,X^m-1)$.
\finenonce{003209}


\finexercice 


\exercice{3210, quercia, 2010/03/08}

\enonce{003210}{Degr{\'e} minimal dans la formule de B{\'e}zout}
Soient $P,Q \in { K[X]}$ non nuls et $D = \text{pgcd}(P,Q)$.
\begin{enumerate}
  \item   D{\'e}montrer qu'il existe $U,V \in { K[X]}$ uniques tels que :
      $\begin{cases} UP + VQ = D \cr
               \deg  U < \deg  Q - \deg  D \hfill\cr
               \deg  V < \deg  P - \deg  D .\end{cases}$

  \item   Montrer que la m{\'e}thode des divisions euclidiennes fournit $U$ et $V$.
      
\end{enumerate}
\finenonce{003210}


\finexercice 


\exercice{3211, quercia, 2010/03/08}

\enonce{003211}{Application $(U,V)  \mapsto UA + VB$}
Soient $A,B \in { K[X]}$, $p = \deg A$, $q = \deg B$.
On consid{\`e}re l'application :
    $${\Phi} : {{ K_{q-1}[X]} \times { K_{p-1}[X]}} \to 
                          { K_{p+q-1}[X]},  {(U,V)}  \mapsto {UA + VB}$$


D{\'e}montrer que : $A \wedge B = 1 \iff \Phi$ est bijective.
\finenonce{003211}



\finexercice 


\exercice{3212, quercia, 2010/03/08}

\enonce{003212}{pgcd$(P(X),P(-X))$ et ppcm$(P(X),P(-X))$}
Soit $P \in { K[X]}$.
D{\'e}montrer que $\pgcd(P(X),P(-X))$ et $\mathrm{ppcm}(P(X),P(-X))$
sont pairs ou impairs.
\finenonce{003212}



\finexercice 


\exercice{3213, quercia, 2010/03/08}

\enonce{003213}{$A\circ P | B\circ P  \Rightarrow  A | B$}
Soient $A,B,P \in { K[X]}$ avec $P$ non constant.
Montrer que si $A\circ P$ divise $B\circ P$, alors $A$ divise $B$.
\finenonce{003213}



\finexercice 


\exercice{5323, rouget, 2010/07/04}
\enonce{005323}{***}
Division euclidienne de $P=\sin aX^n-\sin(na)X+\sin((n-1)a)$ par $Q=X^2-2X\cos a+1$, $a$ réel donné.
\finenonce{005323}


\finexercice


\exercice{6955, exo7, 2014/04/01}
\video{c08DgsqArHw}
% Mélange de 364 (bodin), 375 (ridde), 366 (bodin), 371 (cousquer)
\enonce{006955}{}
\begin{enumerate}
\item Effectuer la division euclidienne de $A$ par $B$ :
\begin{enumerate}
\item $A=3X^5+4X^2+1,\ B=X^2+2X+3$
\item $A=3X^5+2X^4-X^2+1,\ B=X^3+X+2$
\item $A=X^4-X^3+X-2,\ B=X^2-2X+4$
\item $A=X^5-7X^4-X^2-9X+9,\ B=X^2-5X+4$
\end{enumerate}
\item Effectuer la division selon les puissances croissantes de $A$ 
par $B$ à l'ordre $k$ (c'est-à-dire tel que le reste soit divisible 
par $X^{k+1}$) :
\begin{enumerate}
\item $A=1-2X+X^3+X^4,\ B=1+2X+X^2,\ k=2$
\item $A=1+X^3-2X^4+X^6,\ B=1+X^2+X^3,\ k=4$
\end{enumerate}
\end{enumerate}
\finenonce{006955} 


\finexercice

\exercice{6956, blanc-centi, 2014/04/01}
\video{J1GwUtSQ9D4}
\enonce{006956}{}
\`A quelle condition sur $a,b,c\in\Rr$ le polynôme 
$X^4+aX^2+bX+c$ est-il divisible par $X^2+X+1$ ?
\finenonce{006956} 


\finexercice
\section{105.02 Pgcd}


\exercice{379, bodin, 1998/09/01}
\enonce{000379}{}
 Calculer pgcd$(P,Q)$ lorsque :
\begin{enumerate}
    \item  $P=X^3-X^2-X-2$ et $Q=X^5-2X^4+X^2-X-2$,
    \item  $P=X^4+X^3-2X+1$ et $Q=X^3+X+1$.
\end{enumerate}
\finenonce{000379} 
\finexercice


\exercice{380, bodin, 1998/09/01}
\enonce{000380}{}
D\'eterminer le pgcd des polyn\^omes suivants :\\
$X^5+3X^4+X^3+X^2+3X+1 \text{ et } X^4+2X^3+X+2$,\\
$X^4+X^3-3X^2-4X-1 \text{ et } X^3+X^2-X-1$,\\
$X^5+5X^4+9X^3+7X^2+5X+3 \text{ et } X^4+2X^3+2X^2+X+1$.
\finenonce{000380} 
\finexercice


\exercice{381, bodin, 1998/09/01}

\enonce{000381}{}
D\'eterminer $A,B \in \Rr[X]$ tels que $(X^3+1)A+(X^2+X+1)B=1$.
\finenonce{000381}



\finexercice



\exercice{382, bodin, 1998/09/01}

\enonce{000382}{}
Montrer qu'il existe  deux polyn\^omes : $U, V$, v\'erifiant  :
 $ (\star) \ (X-1)^nU+X^nV=1 $.
D\'eterminer $U_1$  et $V_1$ de degr\'e strictement inf\'erieur \`a $n$,
 satisfaisant cette \'egalit\'e. En d\'eduire tous les polyn\^omes $U,V$
 v\'erifiant $(\star)$.
\finenonce{000382}



\finexercice



\exercice{383, bodin, 1998/09/01}

\enonce{000383}{}
Soient $P,Q$ deux polyn\^omes premiers entre eux.
\begin{enumerate}
    \item Montrer qu'alors $P^n$ et $Q^m$ sont premiers entre eux
     o\`u $n,m$ sont deux entiers positifs.
    \item Montrer de m\^eme que $P+Q$ et $PQ$ sont premiers entre eux.
\end{enumerate}
\finenonce{000383}



\finexercice



\exercice{384, bodin, 1998/09/01}

\enonce{000384}{}
Soit $n$ un entier positif.
\begin{enumerate}
    \item D\'eterminer le pgcd des polyn\^omes $(X^n-1)$ et $(X-1)^n$.
    \item Pour $n=3$ d\'emontrer qu'il existe un couple de polyn\^omes $(U,V)$
 tel que $(X^3-1) U+ (X-1)^3V=X-1$. En donner un.
\end{enumerate}
\finenonce{000384}



\finexercice



\exercice{385, cousquer, 2003/10/01}

\enonce{000385}{}
 Montrer que les éléments $X^2+X$, $X^2-X$, $X^2-1$ de  $\mathbb{R}[X]$ 
sont premiers entre eux, mais ne sont pas premiers entre eux deux à deux.
\finenonce{000385}



\finexercice



\exercice{386, cousquer, 2003/10/01}

\enonce{000386}{}
 Trouver tous les polynômes $U$ et $V$ de ${\bf R}[X]$ tels que $AU+BV$ soit un
pgcd de~$A$ et~$B$ avec
$A=X^4-2X^3-2X^2+10X-7$ et $B=X^4-2X^3-3X^2+13X-10$.
\finenonce{000386}



\finexercice



\exercice{387, cousquer, 2003/10/01}
\enonce{000387}{}
 Calculer le pgcd $D$ des polyn\^omes $A$
et $B$ d\'efinis ci-dessous. Trouver des polyn\^omes
$U$ et~$V$ tels que $D=AU+BV$.
\begin{enumerate}
    \item $A=X^5+3X^4+2X^3-X^2-3X-2$ \quad et\quad $B=X^4+2X^3+2X^2+7X+6$.
    \item $A=X^6-2X^5+2X^4-3X^3+3X^2-2X$ \quad et\quad $B=X^4-2X^3+X^2-X+1$.
\end{enumerate}
\finenonce{000387} 
\finexercice


\exercice{388, cousquer, 2003/10/01}

\enonce{000388}{}
 Trouver le pgcd des trois polynômes~:
\begin{eqnarray*}
A & = & X^5+4X^4+6X^3+6X^2+5X+2 \\
B & = & X^2+3X+2 \\
C & = & X^3+2X^2+X+2.
\end{eqnarray*}
\finenonce{000388}



\finexercice



\exercice{389, cousquer, 2003/10/01}

\enonce{000389}{}
 Soit les polynômes de $\mathbb{R}[X]$~:
\begin{eqnarray*}
A & = & (X+3)^2(X+1)(X^2+1)^3 \\
B & = & (X+3)^2(X+2)^2(X^2+1) \\
C & = & (X+3)(X+2)(X^2+1)^2.
\end{eqnarray*}
\begin{enumerate}
\item Combien $A$ possède-t-il de diviseurs normalisés~? et $B$~? et $C$~?
\item Écrire le pgcd et le ppcm de $A$ et~$B$.
\item Écrire le pgcd et le ppcm des trois polynômes~$A$, $B$ et~$C$.
\end{enumerate}
\finenonce{000389}



\finexercice



\exercice{390, cousquer, 2003/10/01}

\enonce{000390}{}
\begin{enumerate}
\item Trouver le pgcd de $X^{24}-1$ et $X^{15}-1$~; 
le pgcd de $X^{280}-1$ et $X^{60}-1$.
\item Montrer que quels que soient les entiers positifs $b$ et~$q$, $X^{b}-1$
divise $X^{bq}-1$. En déduire que le reste de la division de $X^{a}-1$ par
$X^{b}-1$ est $X^{r}-1$ où $r$ est le reste de la division dans $\mathbb{N}$
de $a$ par~$b$. Quel est alors le pgcd de $X^{a}-1$ et $X^{b}-1$~?
Application~: trouver le pgcd de $ X^{5400}-1 $ et $ X^{1920}-1 $.
\item $P$ étant un polynôme quelconque de $\mathbb{C}[X]$, et $a$ et~$b$ deux
entiers naturels, quel est le pgcd de $P^a-1$ et $P^b-1$~? Indication~:
utiliser le théorème de Bézout dans~$\mathbb{Z}$ et dans $\mathbb{C}[X]$.
\end{enumerate}
\finenonce{000390}



\finexercice



\exercice{391, cousquer, 2003/10/01}

\enonce{000391}{}
 Soit $A\in \mathbb{C}[X]$ et $B\in\mathbb{C}[X]$.
\begin{enumerate}
\item A-t-on $\mbox{pgcd}(A,B)=1\iff\mbox{pgcd}(A+B,AB)=1$~?
\item A-t-on $\mbox{pgcd}(A,B)=\mbox{pgcd}(A+B,AB)$~?
\end{enumerate}
\finenonce{000391}



\finexercice



\exercice{392, cousquer, 2003/10/01}

\enonce{000392}{}
 Soit $n$ un entier strictement positif.
\begin{enumerate}
\item Démontrer qu'il existe un unique couple de polynômes $P$ et~$Q$ de degrés
strictement inférieurs à~$n$ tels que $(1-X)^nP(X)+X^nQ(X)=1$.
\item Démontrer que $P(1-X)=Q(X)$ et $Q(1-X)=P(X)$.
\item Démontrer qu'il existe une constante $a$ telle que
$$(1-X)P'(X)-nP(X)=aX^{n-1}.$$ En déduire les coefficients de~$P$ et la valeur
de~$a$.
\end{enumerate}
\noindent{\sl Réponse~:\/} $a=-(2n-1)C_{2n-2}^{n-1}$.
\finenonce{000392}



\finexercice



\exercice{393, cousquer, 2003/10/01}

\enonce{000393}{}
 Déterminer les polynômes $P\in \mathbb{R}[X]$ et $Q\in \mathbb{R}[X]$, 
premiers entre eux,
tels que $P^2+Q^2=(X^2+1)^2$. En déduire que l'équation $x^2+y^2=z^2$ a une
infinité de solutions (non proportionnelles) dans $\mathbb{Z}$.
\finenonce{000393}



\finexercice



\exercice{394, ortiz, 1999/04/01}

\enonce{000394}{}
\begin{enumerate}
    \item  Montrer que les polyn\^omes $X-1$ et $X-2$ sont premiers entre eux et en d\'eduire
$d=\text{pgcd} ((X-1)^2,(X-2)^3)$ et des $U$
et $V$ polyn\^omes tels que
$${U(X-1)^2+V(X-2)^3=d.}$$
    \item  D\'eterminer le polyn\^ome $P$, de degr\'e minimal,
tel que le reste de la division euclidienne de $P$
par $(X-1)^2$ est $2X$ et le reste de la division
euclidienne de $P$ par $(X-2)^3$ est $3X$.
\end{enumerate}
\finenonce{000394}



\finexercice



\exercice{395, ortiz, 1999/04/01}

\enonce{000395}{}
  Montrer que les polyn\^omes complexes $P=X^{1998}+X+1$ et $Q=X^5+X+1$
  sont premiers entre eux.
\finenonce{000395}



\finexercice



\exercice{5317, rouget, 2010/07/04}
\enonce{005317}{**IT}
Déterminer le PGCD de $X^6-7X^4+8X^3-7X+7$ et $3X^5-7X^3+3X^2-7$.
\finenonce{005317}


\finexercice


\exercice{6957, exo7, 2014/04/01}
\video{DYun3S4_zgw}
% Mélange de 379 (bodin), 380 (bodin), 387 (cousquer)
\enonce{006957}{}
\begin{enumerate}
\item Déterminer les pgcd des polynômes suivants:
\begin{enumerate}
\item $X^3-X^2-X-2$ et $X^5-2X^4+X^2-X-2$
\item $X^4+X^3-2X+1$ et $X^3+X+1$
\item $X^5+3X^4+X^3+X^2+3X+1$ et $X^4+2X^3+X+2$
\item $nX^{n+1}-(n+1)X^n+1$ et $X^n-nX+n-1$ ($n\in\N^*)$
\end{enumerate}
\item Calculer le pgcd $D$ des polynômes $A$ et $B$ ci-dessous. 
Trouver des polynômes $U$ et $V$ tels que $AU+BV=D$.
\begin{enumerate}
\item $A=X^5+3X^4+2X^3-X^2-3X-2$\\ et $B=X^4+2X^3+2X^2+7X+6$
\item $A=X^6-2X^5+2X^4-3X^3+3X^2-2X$\\ et $B=X^4-2X^3+X^2-X+1$
\end{enumerate}
\end{enumerate}
\finenonce{006957} 


\finexercice

\exercice{6958, blanc-centi, 2014/04/01}
\video{FljXmG2ie3Q}
\enonce{006958}{}

\
\begin{enumerate}
\item Montrer que si $A$ et $B$ sont deux polynômes à coefficients 
dans $\Q$, alors le quotient et le reste de la division 
euclidienne de $A$ par $B$, ainsi que $\pgcd(A,B)$, sont 
aussi à coefficients dans $\Qq$.

\item Soit $a,b,c\in\Cc^*$ distincts, et $0<p<q<r$ des entiers. 
Montrer que si $P(X)=(X-a)^p(X-b)^q(X-c)^r$ est à coefficients 
dans $\Qq$, alors $a,b,c \in \Qq$.
\end{enumerate}
\finenonce{006958} 



\section{105.03 Racine, décomposition en facteurs irréductibles}


\exercice{396, legall, 2003/10/01}

\enonce{000396}{}
\begin{enumerate}
\item  Montrer que le polyn\^ome $P(X)= X^5-X^2+1$ admet une unique 
racine r\'eelle et que celle-ci est irationnelle.
\item Montrer que le polyn\^ome $Q(X)=2X^3-X^2-X-3$ a une racine 
rationnelle (qu'on calculera). En d\'eduire sa d\'ecomposition en 
produit de facteurs
irr\'eductibles dans $\Cc [X].$
\end{enumerate}
\finenonce{000396}



\finexercice



\exercice{397, legall, 2003/10/01}

\enonce{000397}{}
Soit $P(X)= a_nX^n+\cdots + a_0$ un polyn\^ome \`a 
coefficients entiers premiers entre eux (c'est \`a dire tels que les 
seuls diviseurs
communs \`a tous les $a_i$ soient $-1$ et $1$). Montrer que si 
$r=\dfrac{p}{q}$ avec $p$ et $q$ premiers entre eux est une racine 
rationnelle de $P$
alors $p$ divise $a_0$ et $q$ divise $a_n.$
\finenonce{000397}



\finexercice



\exercice{398, legall, 2003/10/01}

\enonce{000398}{}
Soit $P\in \Qq [X]$ un polyn\^ome 
de degr\'e $n$.
\begin{enumerate} 
\item Montrer que si $P$ est irr\'eductible dans $\Qq$ alors il n'a 
que des racines simples dans $\Cc $.
\item Soit $\lambda \in \Cc $ une racine de $P$, de multiplicit\'e 
strictement plus grande que $\dfrac n 2$.Montrer que $\lambda $ est 
rationnel.
\end{enumerate}
\finenonce{000398}



\finexercice



\exercice{399, bodin, 1998/09/01}

\enonce{000399}{}
 Montrer que le polyn\^ome $nX^{n+2}-(n+2)X^{n+1}+(n+2)X-n$ admet
une racine multiple. Application : d\'eterminer les racines du polyn\^ome
$3X^5-5X^4+5X-3$.

\finenonce{000399}



\finexercice



\exercice{400, bodin, 1998/09/01}

\enonce{000400}{}
Soit $P=(X^2-X+1)^2+1$.
\begin{enumerate}
    \item V\'erifier que $i $ est racine de $P$.
    \item  En d\'eduire alors la d\'ecomposition en produit de facteurs
     irr\'eductibles de $P$ sur ${\R}[X]$
    \item  Factoriser sur ${\C}[X]$ et sur ${\R}[X]$ les polyn\^omes suivants en
produit de polyn\^omes irr\'eductibles :
    $P=X^4+X^2+1,\  Q=X^{2n}+1,\ R=X^6-X^5+X^4-X^3+X^2-X+1$,
 $S=X^5-13X^4+67X^3-171 X^2+216X-108$ (on cherchera les racines doubles de $S$).
\end{enumerate}
\finenonce{000400}



\finexercice



\exercice{401, bodin, 1998/09/01}
\enonce{000401}{}
D\'ecomposer dans $\Rr[X]$, sans d\'eterminer ses racines,
le polyn\^ome $P = X^4+1$, en produit de facteurs irr\'eductibles.
\finenonce{000401} 
\finexercice


\exercice{402, bodin, 1998/09/01}

\enonce{000402}{}
Pour tout $a\in\Rr$ et tout $n\in\Nn^*$, d\'emontrer que
$X-a$ divise $X^n-a^n$.
\finenonce{000402}



\finexercice



\exercice{403, bodin, 1998/09/01}

\enonce{000403}{}
D\'ecomposer $X^{12}-1$ en produit de facteurs
irr\'eductibles dans $\Rr[X]$.
\finenonce{000403}



\finexercice



\exercice{404, bodin, 1998/09/01}

\enonce{000404}{}
Prouver que $B$ divise $A$, où : \\
$A=X^{3n+2}+X^{3m+1}+X^{3p} \text{ et } B=X^2+X+1$,\\
$A=(X+1)^{2n}-X^{2n}-2X-1 \text{ et } B = X(X+1)(2X+1)$,\\
$A=nX^{n+1}-(n+1)X^n+1 \text{ et } B=(X-1)^2$.
\finenonce{000404}



\finexercice



\exercice{405, bodin, 1998/09/01}

\enonce{000405}{}
Soit $P \in {\Z}[X]$ et $n \in {\Z}$ ; notons $m=P(n)$ ;
 (deg($P$) $\ge 1$).
\begin{enumerate}
    \item Montrer que : $\forall k \in {\Z}, m$ divise $P(n+km)$.
    \item Montrer qu'il n'existe pas de polyn\^ome $P$ dans ${\Z}[X]$, non constant,
    tel que pour tout $n \in {\Z}$, $P(n)$ soit premier.
\end{enumerate}
\finenonce{000405}



\finexercice



\exercice{406, bodin, 1998/09/01}

\enonce{000406}{}
Soit $P$ un polyn\^ome de ${\R}[X]$ tel que $P(x) \ge 0$
 pour tout $x \in {\R}$.

Montrer qu'il existe $S,T \in {\R}[X]$ tels que
$P=S^2+T^2$ (on utilisera la factorisation dans ${\C}[X]$).
\textit{Indications :}
\begin{enumerate}
    \item Soient $a,b \in \mathbb{R}$, d\'eterminer $c,d  \in
    \mathbb{R}$ tels que : $ab = c^{2}-d^{2}$,
     v\'erifier que $(a^{2}+b^{2})(c^{2}+d^{2}) = (ac+bd)^{2}+(bc-ad)^{2}$.
    \item R\'esoudre le probl\`eme pour $P$ de degr\'e 2.
    \item Conclure.
\end{enumerate}
\finenonce{000406}



\finexercice



\exercice{407, cousquer, 2003/10/01}

\enonce{000407}{}
Soit $\theta \in\mathbb{R}$~; on suppose $\sin n\theta \neq 0$. 
Déterminer les racines du polynôme $P=\sum_{k=1}^n C_n^k\sin k\theta \,X^k$. 
Vérifier que ces racines sont toutes réelles.
\finenonce{000407}



\finexercice



\exercice{408, cousquer, 2003/10/01}

\enonce{000408}{}
 Soit $a\in\mathbb{C}$, $P \in \mathbb{C}[X]$ et $Q\in \mathbb{C}[X]$, 
premiers entre eux. On suppose
que $a$ est racine double de $P^2+Q^2$. Montrer que $a$ est racine de
${P'}^2+{Q'}^2$.
\finenonce{000408}



\finexercice



\exercice{409, cousquer, 2003/10/01}
\enonce{000409}{}
 Pour $n\in \mathbb{N}^{\ast}$, 
quel est l'ordre de multiplicit\'e de~2 comme racine du polyn\^ome
$$nX^{n+2}-(4n+1)X^{n+1}+4(n+1)X^n-4X^{n-1}$$

\finenonce{000409} 
\finexercice


\exercice{410, cousquer, 2003/10/01}
\video{ywVn-MdHV0A}
\enonce{000410}{}
Pour quelles valeurs de $a$ le polynôme $(X+1)^7-X^7-a$ admet-il une racine multiple réelle?
\finenonce{000410} 


\finexercice


\exercice{411, cousquer, 2003/10/01}

\enonce{000411}{}
Montrer que le polynôme $X^3+2$ est irréductible dans $\mathbb{Q}[X]$. 
Factoriser ce polynôme dans $\mathbb{R}[X]$ et dans $\mathbb{C}[X]$.
\finenonce{000411}



\finexercice



\exercice{412, cousquer, 2003/10/01}
\enonce{000412}{}
  Dans $\mathbb{R}[X]$ et dans $\mathbb{C}[X]$, 
d\'ecomposer les polyn\^omes suivants en facteurs
irr\'eductibles.
\begin{enumerate}
    \item $X^3-3$.
    \item $X^{12}-1$.
\end{enumerate}
\finenonce{000412} 
\finexercice


\exercice{413, cousquer, 2003/10/01}

\enonce{000413}{}
 Quelle est la décomposition de $X^6+1$ en facteurs irréductibles dans 
$\mathbb{C}[X]$~? Dans $\mathbb{R}[X]$~?
\finenonce{000413}



\finexercice



\exercice{414, cousquer, 2003/10/01}

\enonce{000414}{}
 Soit $P$ le polynôme $X^4+2X^2+1$. Déterminer les multiplicités des
racines $i$ et $-i$, de deux façons différentes~: soit en décomposant~$P$
dans $\mathbb{C}[X]$, soit en utilisant le polynôme dérivé de~$P$.
\finenonce{000414}



\finexercice



\exercice{415, cousquer, 2003/10/01}

\enonce{000415}{}
 Soit le polynôme $P=X^8+2X^6+3X^4+2X^2+1$.
\begin{enumerate}
\item Montrer que $j$ est racine de ce polynôme. 
Déterminer son ordre de multiplicité. 
\item Quelle conséquence peut-on tirer de la parité de $P$~? 
\item Décomposer $P$ en facteurs irréductibles dans $\mathbb{C}[X]$ et dans 
$\mathbb{R}[X]$.
\end{enumerate}
\finenonce{000415}



\finexercice



\exercice{416, cousquer, 2003/10/01}

\enonce{000416}{}
 Soit $E$ le polynôme du troisième degré~: $aX^3+bX^2+cX+d$ avec
$a, b, c, d \in \mathbb{R}$ et $a\neq0$, et soit $x_1$, $x_2$, $x_3$ ses trois
racines dans~$\mathbb{C}$. Trouver un polynôme ayant pour racines
$x_1x_2$, $x_2x_3$ et~$x_3x_1$.
\finenonce{000416}



\finexercice



\exercice{417, ridde, 1999/11/01}

\enonce{000417}{}
Soient $x_1, x_2, x_3$ les racines de $X^3-2X^2 + X + 3$. Calculer
$x_1^3 + x_2^3 + x_3^3$.
\finenonce{000417}



\finexercice



\exercice{418, ridde, 1999/11/01}

\enonce{000418}{}
Soit $n \in \Nn$ fix\'e. Montrer qu'il y a un nombre fini de polynômes unitaires
de degr\'e $n$ \`a coefficients entiers ayant toutes leurs racines de module
inf\'erieur ou \'egal \`a $1$.
\finenonce{000418}



\finexercice



\exercice{419, ridde, 1999/11/01}

\enonce{000419}{}
Soit $n \geq 2$ et $P_n (X) = \sum\limits_{k = 0}^n \frac 1{k ! }X^k$.
$P_n$ a-t-il une racine double ?
\finenonce{000419}



\finexercice



\exercice{420, ridde, 1999/11/01}

\enonce{000420}{}
R\'esoudre les \'equations :
\begin{enumerate}
  \item $P'P'' = 18P$ où $P \in \Rr[X]$.
  \item $P (X^2) = (X^2 + 1)P (X)$ où $P \in \Cc[X]$.
\end{enumerate}
\finenonce{000420}



\finexercice



\exercice{421, ridde, 1999/11/01}

\enonce{000421}{}
Soit $P \in \Rr[X]$ scind\'e sur $\Rr$ \`a racines simples.
\begin{enumerate}
\item Montrer qu'il en est de même de $P'$.
\item Montrer que le polynôme $P^2 + 1$ n'a que des racines simples dans $\Cc$.
\end{enumerate}
\finenonce{000421}



\finexercice



\exercice{422, ridde, 1999/11/01}

\enonce{000422}{}
Soit $n \in \Nn^*$ et $P (X) = (X + 1)^n- (X-1)^n$.
\begin{enumerate}
\item Quel est le degr\'e de $P$ ?
\item Factoriser $P$ dans $\Cc[X]$.
\item Montrer que $\forall p \in \Nn^* \, \, \prod\limits_{k = 1}^p
\text{cotan} (\dfrac{k \pi}{2p + 1}) = \dfrac 1{\sqrt{2p + 1}}$.
\end{enumerate}
\finenonce{000422}



\finexercice



\exercice{423, ridde, 1999/11/01}
\enonce{000423}{}
Factoriser dans $\Rr[X]$ :
\begin{enumerate}
\item $X^6 + 1$.
\item $X^9 + X^6 + X^3 + 1$.
\end{enumerate}

\finenonce{000423} 
\finexercice


\exercice{3214, quercia, 2010/03/08}

\enonce{003214}{Factorisation de $X^n-1$}
Factoriser $X^n-1$ sur $\C$.
\begin{enumerate}
  \item En d{\'e}duire $\prod_{k=1}^{n-1} \sin\left(\frac {k\pi}n\right)$.

  \item Calculer {\'e}galement $\prod_{k=0}^{n-1} \sin\left(\frac {k\pi}n + \theta\right)$.

  \item On note $\omega = e^{2i\pi/n}$. Calculer
    $\prod_{0 \le k,\ell < n, k \ne \ell} (\omega^k - \omega^\ell)$.
\end{enumerate}
\finenonce{003214}


\finexercice 


\exercice{3215, quercia, 2010/03/08}

\enonce{003215}{Mines MP 1999}
Montrer que $\prod_{k=0}^{n-1} (\omega ^{2k}-2\omega ^k \cos \theta +1)=2(1-\cos (n\theta))$ avec
$\omega = e^{\frac{2i\pi}{n}}$.
\finenonce{003215}


\finexercice 


\exercice{3216, quercia, 2010/03/08}

\enonce{003216}{Racines de $j$ et $j^2$}
Montrer que si $p \le n$, alors $X^{2^p} + X^{2^{p-1}} + 1$ divise
$X^{2^n} + X^{2^{n-1}} + 1$.
\finenonce{003216}



\finexercice 


\exercice{3217, quercia, 2010/03/08}

\enonce{003217}{$X^2-2X\cos\theta+1$ divise $X^{2n}-2X^n\cos(n\theta)+1$}
Montrer que $X^2-2X\cos\theta+1$ divise $X^{2n}-2X^n\cos n\theta+1$.
Pour $\sin\theta \ne 0$, chercher le quotient.
\finenonce{003217}


\finexercice 


\exercice{3218, quercia, 2010/03/08}

\enonce{003218}{$X^2-2X\cos\theta+1$ divise $X^{n+1}\cos(n-1)\theta-X^n\cos n\theta-X\cos\theta+1$} 
Montrer que $X^2-2X\cos\theta+1$ divise
$X^{n+1}\cos(n-1)\theta-X^n\cos n\theta-X\cos\theta+1$, puis
d{\'e}terminer le quotient.  \finenonce{003218}



\finexercice 


\exercice{3219, quercia, 2010/03/08}

\enonce{003219}{$X^8 + X^4 + 1$ divise $X^{8n} + pX^{4n} + q$} Donner une CNS
sur $p,q \in \C$ pour que $X^8 + X^4 + 1$ divise $X^{8n} + pX^{4n} +
q$ ($n \in \N^*$ fix{\'e}).
 \finenonce{003219}



\finexercice 


\exercice{3220, quercia, 2010/03/08}

\enonce{003220}{Racines rationnelles}
Factoriser $P(X) = 3X^4 + 11X^3 + 20X^2 + 7X - 5$, sachant
qu'il existe des racines rationnelles.
\finenonce{003220}


\finexercice 


\exercice{3221, quercia, 2010/03/08}

\enonce{003221}{{\'E}quation de degr{\'e} 4 tq $x_1x_2 = 5$}
Trouver les racines de $P(X) =X^4-3X^3+6X^2-15X+5$
sachant que deux racines, $x_1$ et $x_2$, v{\'e}rifient~: $x_1x_2 = 5$
(on introduira le polyn{\^o}me $Q = X^4P(5/X)$).
\finenonce{003221}


\finexercice 


\exercice{3222, quercia, 2010/03/08}

\enonce{003222}{Racines multiples}
Factoriser $P = X^5 - 13X^4 + 67X^3 - 171X^2 + 216X - 108$
sachant qu'il admet une racine triple.
\finenonce{003222}


\finexercice 


\exercice{3223, quercia, 2010/03/08}

\enonce{003223}{Recherche d'une racine triple}
Soit $P = X^5 + aX^2 + 15X - 6i$. Trouver $a \in \C$ tel que
$P$ a une racine triple dans $\C$. Factoriser alors $P$.
\finenonce{003223}


\finexercice 


\exercice{3224, quercia, 2010/03/08}

\enonce{003224}{Ensi P 90}
Donner une condition sur $\lambda$ pour que l'{\'e}quation :
$x^4-2x^3+\lambda x^2+2x-1 = 0$ ait une racine au moins triple.
\finenonce{003224}


\finexercice 


\exercice{3225, quercia, 2010/03/08}

\enonce{003225}{$x_1+x_2 = 1$}
Soient $p,q\in\C$ et $P(X) = X^5 + pX + q$.
Donner une CNS sur $p$ et $q$ pour que deux des racines de $P$ aient
pour somme 1.
\finenonce{003225}


\finexercice 


\exercice{3226, quercia, 2010/03/08}

\enonce{003226}{Factorisation}
Factoriser $$1 -\frac X{1!} + \frac {X(X-1)}{2!} - \dots +
            (-1)^{n+1}\frac {X(X-1)\cdots(X-n)}{(n+1)!}.$$
\finenonce{003226}


\finexercice 


\exercice{3227, quercia, 2010/03/08}

\enonce{003227}{$X-1 \mid P(X^n)  \Rightarrow  X-1 \mid P$}
Soient $P,Q \in { K[X]}$.
\begin{enumerate}
  \item Montrer que si $P(X^n)$ est divisible par $X - 1$, alors $P$ est divisible par $X-1$\quad $(n\in\N)$.
  \item Montrer que si $P(X^3) + XQ(X^3)$ est divisible par $X^2+X+1$, alors $P$ et $Q$
    sont divisibles par $X-1$.
\end{enumerate}
\finenonce{003227}



\finexercice  


\exercice{3228, quercia, 2010/03/08}

\enonce{003228}{Racines de $\sum_{k = 0}^n C_n^k(\sin k\theta)X^k$}
Soit $\theta \in \R$ tel que $\sin n\theta \ne 0$.
D{\'e}montrer que le polyn{\^o}me $P = \sum_{k = 0}^n C_n^k(\sin k\theta)X^k$ a toutes ses
racines r{\'e}elles.
\finenonce{003228}


\finexercice 


\exercice{3229, quercia, 2010/03/08}

\enonce{003229}{}
D{\'e}montrer que $1 + X + X^n$ n'a que des racines simples.
\finenonce{003229}



\finexercice


\exercice{3230, quercia, 2010/03/08}

\enonce{003230}{$P'$ divise $P$}
Quels sont les polyn{\^o}mes $P \in { K[X]}$ tels que $P'$ divise $P$ ?
\finenonce{003230}


\finexercice 


\exercice{3231, quercia, 2010/03/08}

\enonce{003231}{{\'E}quations fonctionnelles}
Trouver tous les polyn{\^o}mes $P \in {\C[X]}$ tels que \dots

\begin{enumerate}
  \item $P(X^2) = P(X-1)P(X+1)$.
  \item $P(X^2) = P(X)P(X-1)$.
  \item $P(X)P(X+2) + P(X^2) = 0$.
\end{enumerate}
\finenonce{003231}


\finexercice 


\exercice{3232, quercia, 2010/03/08}

\enonce{003232}{$P$ {\`a} racines r{\'e}elles simples $=> P^2+a^2$ {\`a} racines simples}
Soit $P \in {\R[X]}$ dont toutes les racines sont r{\'e}elles.
\begin{enumerate}
  \item D{\'e}montrer que les racines de $P'$ sont aussi r{\'e}elles.
  \item En d{\'e}duire que : $\forall\ a \in \R^*$, les racines de $P^2+a^2$ sont
    simples.
\end{enumerate}
\finenonce{003232}



\finexercice 


\exercice{3233, quercia, 2010/03/08}

\enonce{003233}{$P$ et $Q$ ont m{\^e}me module}
Soient $P,Q \in {\C[X]}$ tels que : $\forall\ z \in \C,\ |P(z)| = |Q(z)|$.
D{\'e}montrer qu'il existe $u \in \C, |u| = 1$ tel que $P = uQ$.
\finenonce{003233}


\finexercice 


\exercice{3234, quercia, 2010/03/08}

\enonce{003234}{Valeur moyenne}

Soient $z_0,z_1,\dots,z_n \in \C$ tels que :
$\forall\ P \in \C_{n-1}[X]$, on a $P(z_0) = \frac {P(z_1) +\dots+ P(z_n)}n$.
\\
On note $\Phi(X) = \prod_{i=1}^n (X-z_i)$.
\begin{enumerate}
  \item  Calculer $\frac {\Phi(z_0)}{z_0-z_k}$.
  \item  En d{\'e}duire que $\Phi(X) = \frac {(X-z_0)\Phi'(X)}n + \Phi(z_0)$.
  \item  D{\'e}montrer que $z_1,\dots,z_n$ sont les sommets d'un polygone r{\'e}gulier de centre
     $z_0$.
  \item  R{\'e}ciproque ?
\end{enumerate}
\finenonce{003234}


\finexercice 


\exercice{3235, quercia, 2010/03/08}

\enonce{003235}{$P(x)\ne14$}
Soit $P \in {\Z[X]}$ tel que $P(x) = 7$ pour au moins 4 valeurs distinctes
$x \in \Z$.
\\
D{\'e}montrer que : $\forall\ x \in \Z$, on a $P(x) \ne 14$.
\finenonce{003235}



\finexercice 


\exercice{3236, quercia, 2010/03/08}

\enonce{003236}{Nombre alg{\'e}brique rationnel}
Soit $\alpha \in \C$. On dit que $\alpha$ est {\it alg{\'e}brique\/}
s'il existe un polyn{\^o}me $P \in {\Q[X]}$ tel que $P(\alpha) = 0$. \\
Le polyn{\^o}me unitaire de plus bas degr{\'e} v{\'e}rifiant $P(\alpha) = 0$ est appel{\'e} :
{\it polyn{\^o}me minimal de $\alpha$}.
\begin{enumerate}
\item Soit $\alpha$ alg{\'e}brique de polyn{\^o}me minimal $P$.  D{\'e}montrer que
  $P$ est irr{\'e}\-duc\-tible dans ${\Q[X]}$ et que $\alpha$ est racine
  simple de $P$.
  \item  Soit $\alpha$ alg{\'e}brique, et $P \in {\Q[X]}$ tel que $P(\alpha) = 0$.
     On suppose que la multiplicit{\'e} de $\alpha$ dans $P$
     est strictement sup{\'e}rieure {\`a} $\frac 12\deg P$. D{\'e}montrer que $\alpha \in \Q$.
\end{enumerate}
\finenonce{003236}



\finexercice 


\exercice{3237, quercia, 2010/03/08}

\enonce{003237}{$P(\sqrt 2) = 0$}
Soit $P \in {\Q[X]}$ tel que $P\bigl(\sqrt2\bigr) = 0$. D{\'e}montrer que $-\sqrt2$ est aussi racine de $P$
avec la m{\^e}me multiplicit{\'e} que~$\sqrt2$.
\finenonce{003237}



\finexercice 


\exercice{3238, quercia, 2010/03/08}

\enonce{003238}{Polyn{\^o}me minimal de $2\cos(2\pi/7)$}
Montrer que $x = 2\cos\frac {2\pi}7$ est racine de $X^3 + X^2 - 2X - 1$.
Quelles sont les autres racines ?
\finenonce{003238}


\finexercice 


\exercice{3239, quercia, 2010/03/08}

\enonce{003239}{Racines r{\'e}elles simples}
Soit $P = \sum_{k=0}^n a_kX^k \in {\R[X]}$ dont les racines sont r{\'e}elles simples.

\begin{enumerate}
  \item   D{\'e}montrer que : $\forall\ x \in \R$, on a $P(x)P''(x) \le P'^2(x)$.
  \item   D{\'e}montrer que : $\forall\ k \in \{1,\dots, n-1\},\ a_{k-1}a_{k+1} \le a_k^2$.
\end{enumerate}
\finenonce{003239}


\finexercice 


\exercice{3240, quercia, 2010/03/08}

\enonce{003240}{M{\'e}thode de Ferrari}
Soit $P = X^4 - 6X^3 +7X^2 -18X - 8$.
\par
Trouver $Q \in {\R[X]}$ tel que deg$(Q) = {}$deg$(P-Q^2) = 2$, et $P-Q^2$ a une
racine double. Factoriser alors $P$ sur $\R$.
\finenonce{003240}


\finexercice 


\exercice{3241, quercia, 2010/03/08}

\enonce{003241}{Pgcd $\ne 1 \Leftrightarrow$ racine commune}
Soient $P,Q \in {\Q[X]}$. Montrer que $P$ et $Q$ sont premiers entre eux si et seulement
si $P$ et $Q$ n'ont pas de racine en commun dans $\C$.
\finenonce{003241}



\finexercice 


\exercice{3242, quercia, 2010/03/08}

\enonce{003242}{Mines MP 2001}
Soit $ K$ un corps de caract{\'e}ristique $p$.
\begin{enumerate}
  \item Montrer que $\sigma\ :\ x \mapsto x^p$ est un morphisme de corps.

  \item Montrer que $\sigma$ est surjectif si et seulement si tout polyn{\^o}me
    $P\in K[X]$ irr{\'e}ductible v{\'e}rifie $P'\ne 0$.
\end{enumerate}
\finenonce{003242}


\finexercice 


\exercice{3243, quercia, 2010/03/08}

\enonce{003243}{Centrale MP 2001}
Soit $P \in \R_n[X]\setminus\{0\}$.

Pour $x\in\R$ on note $V(x)$ le
nombre de changements de signe dans la suite $(P(x),P'(x),\dots,P^{(n)}(x))$
en convenant de retirer les termes nuls.
Soient $\alpha<\beta$ deux r{\'e}els non racines de~$P$. Montrer que le nombre
de racines de~$P$ dans $[\alpha,\beta]$, compt{\'e}es avec leur ordre de multiplicit{\'e},
a m{\^e}me parit{\'e} que~$V(\alpha)-V(\beta)$ et que $V(\alpha)-V(\beta)\ge 0$.
\finenonce{003243}


\finexercice 


\exercice{3244, quercia, 2010/03/08}

\enonce{003244}{X MP$^*$ 2004}
Soit $P\in\C[X]$ de degr{\'e} $d$ dont toutes les racines sont de module
strictement inf{\'e}rieur {\`a}~$1$. Pour $\omega\in\mathbb{U}$ on note $\overline P$ le polyn{\^o}me dont les
coefficients sont les conjugu{\'e}s de ceux de~$P$ et $Q(X) = P(X) + \omega X^d\overline P(1/X)$.
Montrer que les racines de~$Q$ sont de module~$1$.
\finenonce{003244}


\finexercice 


\exercice{3245, quercia, 2010/03/08}

\enonce{003245}{X MP$^*$ 2005}
Soient $a_0,\dots,a_n\in\R$ tels que $|a_0| + \dots + |a_{n-1}| < a_n$.
Soit $f(x) = a_0 + a_1\cos x + \dots + a_n\cos(nx)$. Montrer que les
z{\'e}ros de~$f$ sont tous r{\'e}els (cad. si $x\in\C\setminus\R$, alors $f(x)\ne 0$).
\finenonce{003245}


\finexercice 


\exercice{3246, quercia, 2010/03/08}

\enonce{003246}{Factorisation sur $\R$ de $X^8 + X^4 + 1$}
Factoriser $X^8 + X^4 + 1$ sur $\R$.
\finenonce{003246}


\finexercice 


\exercice{3247, quercia, 2010/03/08}

\enonce{003247}{Polyn{\^o}me irr{\'e}ductible sur $\Q$}
D{\'e}montrer que $1 + (X-1)^2(X-3)^2$ est irr{\'e}ductible dans ${\Q[X]}$.
\finenonce{003247}


\finexercice 


\exercice{3248, quercia, 2010/03/08}

\enonce{003248}{Polyn{\^o}mes positifs sur $\R$}
Soit ${\cal E} = \{P \in {\R[X]}$ tq $\exists\ Q,R \in {\R[X]}$ tq $P = Q^2 + R^2 \}$.
\begin{enumerate}
  \item   Montrer que ${\cal E}$ est stable par multiplication.

  \item   Montrer que ${\cal E} = \{P \in {\R[X]}$ tq $\forall\ x \in \R,\ P(x) \ge 0 \}$.

  \item (Centrale MP 2000, avec Maple)
    $P=65X^4-134X^3+190X^2-70X+29$. Trouver $A$ et $B$ dans $\Z [X]$ tels que $P=A^2+B^2$.
\end{enumerate}
\finenonce{003248}


\finexercice 


\exercice{3249, quercia, 2010/03/08}

\enonce{003249}{Lemme de Gauss}
Soit $P \in {\Z[X]}$. On appelle {\it contenu de $P$\/} le pgcd des coefficients
de $P$ (notation : cont($P$)).
\begin{enumerate}
  \item Soient $P,Q \in {\Z[X]}$ avec cont$(P)=1$, et $R = PQ$. Soit $p$ un facteur
    premier de cont($R$).
  \begin{enumerate}
    \item   Si $p$ est premier avec le coefficient constant de $P$, D{\'e}montrer que
          $p$ divise tous les coefficients de $Q$.
    \item   Si $p$ divise le coefficient constant de $P$, se ramener au cas
          pr{\'e}c{\'e}dent.
    \item   En d{\'e}duire que cont($Q$) = cont($R$).
  \end{enumerate}
  \item  Lorsque cont$(P) \ne 1$, trouver cont($PQ$).
  \item  Application : Soit $R \in {\Z[X]}$, et $P,Q \in {\Q[X]}$ tels que $R = PQ$.
     Montrer qu'il existe $P_1,Q_1 \in {\Z[X]}$ proportionnels {\`a} $P$ et $Q$ et tels que
     $R = P_1Q_1$.

     ({\it cad : un polyn{\^o}me {\`a} coefficients entiers r{\'e}ductible sur $\Q$
     est aussi r{\'e}ductible sur $\Z$})
\end{enumerate}
\finenonce{003249}



\finexercice 


\exercice{3250, quercia, 2010/03/08}

\enonce{003250}{Polyn{\^o}mes irr{\'e}ductibles sur $\Z$}
D{\'e}montrer que $X^4 + X + 1$ et $X^6 + X^2 + 1$ sont irr{\'e}ductibles dans ${\Z[X]}$.
\finenonce{003250}



\finexercice 


\exercice{3251, quercia, 2010/03/08}

\enonce{003251}{Polyn{\^o}mes irr{\'e}ductibles sur $\Z$}
Soient $a_1, \dots, a_n \in \Z$ distincts.
\begin{enumerate}
  \item Montrer que $(X-a_1) \dots (X-a_n) - 1$ est irr{\'e}ductible dans ${\Z[X]}$.
  \item M{\^e}me question avec $(X-a_1) \dots (X-a_n) + 1$, $n$ impair.
\end{enumerate}
\finenonce{003251}


\finexercice 


\exercice{3252, quercia, 2010/03/08}

\enonce{003252}{Crit{\`e}re d'irr{\'e}ductibilit{\'e} d'Eisenstein}
Soit $P\in\Z[X]$, $P=X^n+a_{n-1}X^{n-1}+\dots+a_0X^0$ et $p$ un
nombre premier tel que~:
$$a_0\equiv 0 (\mathrm{mod}\, p),\quad\dots,\quad a_{n-1}\equiv 0 (\mathrm{mod}\, p),\quad
a_0\not\equiv 0 (\mathrm{mod}\, {p^2}).$$
Montrer que $P$ est irr{\'e}ductible dans $\Z[X]$.

\finenonce{003252}


\finexercice 


\exercice{3253, quercia, 2010/03/08}

\enonce{003253}{Irr{\'e}ductibilit{\'e} de $X^p-a$}
Soit $ K$ un sous-corps de~$\C$, $a\in K$ et $p\in\N$ premier.
Montrer que le polyn{\^o}me $X^p-a$ est irr{\'e}ductible sur~$ K$ si et seulement
s'il n'a pas de racine dans~$ K$.
\finenonce{003253}


\finexercice 


\exercice{5318, rouget, 2010/07/04}
\enonce{005318}{**T}
Pour quelles valeurs de l'entier naturel $n$ le polynôme $(X+1)^n-X^n-1$ est-il divisible par $X^2+X+1$~?
\finenonce{005318}


\finexercice


\exercice{5319, rouget, 2010/07/04}
\enonce{005319}{***}
Soit $P$ un polynôme à coefficients réels tel que $\forall x\in\Rr,\;P(x)\geq0$. Montrer qu'il existe deux polynômes $R$ et $S$ à coefficients réels tels que $P=R^2+S^2$.
\finenonce{005319}


\finexercice


\exercice{5324, rouget, 2010/07/04}
\enonce{005324}{****I Théorème de \textsc{Lucas}}
Soit $P\in\Cc[X]$ de degré supérieur ou égal à $1$.
Montrer que les racines de $P'$ sont barycentres à coefficients positifs des racines de $P$ (on dit que les racines de $P'$ sont dans l'enveloppe convexe des racines de $P$). Indication~:~calculer $\frac{P'}{P}$.
\finenonce{005324}


\finexercice


\exercice{5325, rouget, 2010/07/04}
\enonce{005325}{***}
Trouver tous les polynômes divisibles par leur dérivée.
\finenonce{005325}


\finexercice


\exercice{5328, rouget, 2010/07/04}
\enonce{005328}{**T}
Déterminer $a\in\Cc$ tel que $P=X^5-209X+a$ admette deux zéros dont le produit vaut $1$.
\finenonce{005328}


\finexercice


\exercice{5329, rouget, 2010/07/04}
\enonce{005329}{***T}
Soit $(a_k)_{1\leq k\leq 5}$ la famille des racines de $P=X^5+2X^4-X-1$. Calculer $\sum_{k=1}^{5}\frac{a_k+2}{a_k-1}$.
\finenonce{005329}


\finexercice


\exercice{5342, rouget, 2010/07/04}
\enonce{005342}{}
Décomposer en produit de facteurs irréductibles dans $\Rr[X]$ le polynôme $X^6-2X^3\cos a+1$ où $a$ est un réel donné dans $[0,\pi]$.
\finenonce{005342}


\finexercice


\exercice{5345, rouget, 2010/07/04}
\enonce{005345}{}
Former une équation du sixième degré dont les racines sont les $\sin\frac{k\pi}{7}$ où $k\in\{-3,-2,-1,1,2,3\}$ puis montrer que ces six nombres sont irrationnels.
\finenonce{005345}


\finexercice


\exercice{5349, rouget, 2010/07/04}
\enonce{005349}{}
Déterminer $\lambda$ et $\mu$ complexes tels que les zéros de $z^4-4z^3-36z^2+\lambda z+\mu$ soient en progression arithmétique. Résoudre alors l'équation.
\finenonce{005349}


\finexercice


\exercice{5350, rouget, 2010/07/04}
\enonce{005350}{}
Soient $x_1$, $x_2$, $x_3$ les zéros de $X^3+2X-1$. Calculer $x_1^4+x_2^4+x_3^4$.
\finenonce{005350}


\finexercice


\exercice{5351, rouget, 2010/07/04}
\enonce{005351}{}
Soient $x_1$,..., $x_8$ les zéros de $X^8+X^7-X+3$. Calculer $\sum_{}^{}\frac{x_1}{x_2x_3}$ (168 termes).
\finenonce{005351}


\finexercice


\exercice{6959, exo7, 2014/04/01}
\video{aYIfJ9ze69o}
% Mélange 412 (cousquer), 423 (ridde)
\enonce{006959}{}
\begin{enumerate}
\item Factoriser dans $\Rr[X]$ et $\Cc[X]$ les polynômes suivants :
$$a)\ X^3-3\quad\quad b)\ X^{12}-1\quad\quad c)\ X^6+1\quad\quad d)\ X^9+X^6+X^3+1$$

\item Factoriser les polynômes suivants :
$$a)\ X^2+(3i-1)X-2-i\quad\quad b)\ X^3+(4+i)X^2+(5-2i)X+2-3i$$
\end{enumerate}
\finenonce{006959} 


\finexercice

\exercice{6960, blanc-centi, 2014/04/01}
\video{nC2gu_mKDR4}
\enonce{006960}{}
Trouver tous les polynômes $P$ qui vérifient la relation 
$$P(X^2)=P(X)P(X+1)$$
\finenonce{006960} 


\finexercice


\exercice{6961, blanc-centi, 2014/04/01}
\video{8EYEgxhsGFA}
\enonce{006961}{}
Soit $n\in\Nn$. Montrer qu'il existe un unique $P\in\Cc[X]$ tel que 
$$\forall z\in\Cc^* \qquad  P\left(z+\frac{1}{z}\right) = z^n+\frac{1}{z^n}$$

Montrer alors que toutes les racines de $P$ sont réelles, simples, 
et appartiennent à l'intervalle $[-2,2]$.
\finenonce{006961} 


\finexercice


\exercice{6962, blanc-centi, 2014/04/01}
\video{uMsrg-zPUko}
\enonce{006962}{}
\begin{enumerate}
\item Soit $P=X^n+a_{n-1}X^{n-1}+\cdots+a_1X+a_0$ un polynôme de degré $n\ge 1$ 
à coefficients dans $\Zz$. Démontrer que si $P$ admet une racine dans $\Zz$, 
alors celle-ci divise $a_0$.

\item Les polynômes $X^3-X^2-109X-11$ et $X^{10}+X^5+1$ ont-ils des racines dans $\Z$?
\end{enumerate}
\finenonce{006962} 


\finexercice


\exercice{6963, blanc-centi, 2014/04/01}
\video{Ez8FgWWwzVE}
\enonce{006963}{}

Soient $a_0,\ldots,a_n$ des réels deux à deux distincts.
 Pour tout $i=0,\ldots,n$, on pose
$$L_i(X)=\prod_{\substack{1\le j\le n \\ j\not= i}}\frac{X-a_j}{a_i-a_j}$$
(les $L_i$ sont appelés \emph{polynômes interpolateurs de Lagrange}).
Calculer $L_i(a_j)$.

Soient $b_0,\ldots,b_n$ des réels fixés. 
Montrer que $P(X)=\sum_{i=0}^nb_iL_i(X)$ est l'unique polynôme de degré inférieur ou égal à $n$ qui vérifie:
$$P(a_j)=b_j  \quad \text{ pour tout }j=0,\ldots,n.$$


\emph{Application.} Trouver le polynôme $P$ de degré inférieur ou égal à $3$ tel que 
$$P(0)=1\quad\text{et}\quad P(1)=0\quad\text{et}\quad P(-1)=-2\quad\text{et}\quad P(2)=4.$$
\finenonce{006963} 


\finexercice

\section{105.04 Fraction rationnelle}


\exercice{443, ridde, 1999/11/01}

\enonce{000443}{}
D\'ecomposer les fractions rationnelles suivantes :
$$\frac{3}{X^3 + 1} \text{ sur } \Cc \text{ puis sur } \Rr$$
 $$\frac{X^3}{X^3-1} \text{ sur } \Rr $$
 $$\frac{X^2 + X + 1}{ (X-1)^2 (X + 1)^2} \text{ sur } \Rr $$
 $$F (X) = \frac{1}{ (X^3-1)^2} \text{ sur } \Cc \text{ en remarquant que }
F (jX) = F (X)$$
$$\frac{X^7 + 1}{ (X^2 + 1) (X^2 + X + 1)} \text{ sur } \Rr $$
$$\frac{3X^5 + 2X^4 + X^2 + 3X + 2}{X^4 + 1} \text{ sur } \Rr $$
$$\frac{1}{X^{2n} + 1} \text{ sur } \Cc \text{ puis sur } \Rr$$
$$\frac{X^3 + X}{ (X^2 + X + 1)^2} \text{ sur } \Rr $$
\finenonce{000443}



\finexercice



\exercice{444, cousquer, 2003/10/01}
\enonce{000444}{}
\begin{enumerate}
  \item
  D\'ecomposer 
$\frac{X^3-3X^2+X-4}{X-1}$
en \'el\'ements simples sur $\mathbb{R}$.
   \item D\'ecomposer 
$\frac{2X^3+X^2-X+1}{X^2-3X+2}$
en \'el\'ements simples sur $\mathbb{R}$.
  \item
 D\'ecomposer 
$\frac{2X^3+X^2-X+1}{X^2-2X+1}$
en \'el\'ements simples sur $\mathbb{R}$.
 \item
D\'ecomposer 
$\frac{X^4+2X^2+1}{X^2-1}$
en \'el\'ements simples sur $\mathbb{R}$.
 \item D\'ecomposer 
$\frac{X}{X^2-4}$
en \'el\'ements simples sur $\mathbb{R}$.
  \item D\'ecomposer 
$\frac{X^5+X^4+1}{X^3-X}$
en \'el\'ements simples sur $\mathbb{R}$.
 \item D\'ecomposer 
$\frac{X^5+X^4+1}{X(X-1)^4}$
en \'el\'ements simples sur $\mathbb{R}$.
  \item D\'ecomposer 
$\frac{X^5+X^4+1}{(X-1)^3(X+1)^2}$
en \'el\'ements simples sur $\mathbb{R}$.
 \item  D\'ecomposer 
$\frac{X^7+3}{(X^2+X+2)^3}$
en \'el\'ements simples sur $\mathbb{R}$.
 \item D\'ecomposer 
$\frac{(3-2 i )X-5+3 i }{X^2+ i  X+2}$
en \'el\'ements simples sur $\mathbb{C}$.
 \item D\'ecomposer 
$\frac{X+ i }{X^2+ i }$
en \'el\'ements simples sur $\mathbb{C}$.
 \item D\'ecomposer 
$\frac{X}{(X+ i )^2}$
en \'el\'ements simples sur $\mathbb{C}$.
 \item D\'ecomposer $\frac{X^2+1}{X^4+1}$ en \'el\'ements simples sur $\mathbb{R}$ et sur $\mathbb{C}$.
\item  D\'ecomposer 
$\frac{X}{X^4+1}$
en \'el\'ements simples sur $\mathbb{R}$ et sur $\mathbb{C}$.
\item  D\'ecomposer 
$\frac{X^2+X+1}{X^4+1}$
en \'el\'ements simples sur $\mathbb{R}$ et sur $\mathbb{C}$.
\item  D\'ecomposer 
$\frac{X^5+X+1}{X^4-1}$
en \'el\'ements simples sur $\mathbb{R}$ et sur $\mathbb{C}$.
\item  D\'ecomposer 
$\frac{X^5+X+1}{X^6-1}$
en \'el\'ements simples sur $\mathbb{R}$ et sur $\mathbb{C}$.
\item  D\'ecomposer 
$\frac{X^3-2}{X^4(X^2+X+1)^2}$
en \'el\'ements simples sur $\mathbb{R}$ et sur $\mathbb{C}$.
 \item D\'ecomposer 
$\frac{X}{(X^2+1)(X^2+4)}$
en \'el\'ements simples sur $\mathbb{R}$ et sur $\mathbb{C}$.
 \item D\'ecomposer 
$\frac{X^2-3}{(X^2+1)(X^2+4)}$
en \'el\'ements simples sur $\mathbb{R}$ et sur $\mathbb{C}$.
\end{enumerate}

\finenonce{000444} 
\finexercice


\exercice{445, cousquer, 2003/10/01}
\enonce{000445}{}
D\'ecomposition en \'el\'ements simples
 $\displaystyle\Phi={2x^4+x^3+3x^2-6x+1\over2x^3-x^2}.$ 

\finenonce{000445} 
\finexercice


\exercice{446, cousquer, 2003/10/01}
\enonce{000446}{}
D\'ecomposition en \'el\'ements simples
 $\displaystyle\Phi={2x^5-8x^3+8x^2-4 x+1\over x^3(x-1)^2}.$
\finenonce{000446} 
\finexercice


\exercice{447, cousquer, 2003/10/01}
\enonce{000447}{}
D\'ecomposition en \'el\'ements simples
 $\displaystyle\Phi={4x^6-2x^5+11x^4-x^3+11x^2+2x+3\over x(x^2+1)^3}.$

\finenonce{000447} 
\finexercice


\exercice{448, ridde, 1999/11/01}

\enonce{000448}{}
Soient $a$ et $b$ deux r\'eels distincts et $F (X) = \dfrac1{ (X-a)^n (X-b)^n}$.
En utilisant la formule de Taylor en $a$ pour $f (X) = (X-a)^n F (X)$, d\'ecomposer
$F$ sur $\Rr$.
\finenonce{000448}



\finexercice



\exercice{449, gourio, 2001/09/01}

\enonce{000449}{}
Donner une CNS sur $f\in \Cc(X)$ pour qu'il existe $g\in \Cc(X)$ tel que 
$f=g^{\prime }.$
\finenonce{000449}



\finexercice



\exercice{450, gourio, 2001/09/01}

\enonce{000450}{}
On appelle valuation une application $v :\Cc(X)\rightarrow {\Zz}\cup
\{\infty \}  $ telle que :
$\lambda \in \Cc^{*}\Rrightarrow v(\lambda )=0,v(0)=\infty ,\exists
 a\in \Cc(X):v(a)=1$
$$\forall (f,g)\in \Cc(X)^{2},v(fg)=v(f)+v(g) $$
$$\forall (f,g)\in \Cc(X)^{2},v(f+g)\geq \min (v(f),v(g)) $$
(avec les convention \'{e}videntes $k+\infty =\infty ,\forall k\geq
1:k\infty =\infty ,0\infty =0$, etc.)
D\'{e}terminer toutes les valuations de $\Cc(X)$ et montrer la formule
(la somme portant sur toutes les valuations) :
$$\forall f\in \Cc(X)-\{0\},\sum\limits_{v}v(f)=0. $$
\finenonce{000450}



\finexercice



\exercice{3270, quercia, 2010/03/08}

\enonce{003270}{Substitution de fractions}
Soit $F \in { K(X)}$ non constante et $P \in { K[X]}$, $P \ne 0$.
\begin{enumerate}
  \item Montrer que $P\circ F \ne 0$.
  \item Montrer que l'application${ K(X)} \to { K(X)},  G \mapsto  {G\circ F}$ est un morphisme
    injectif d'alg{\`e}bre.
  \item A quelle condition est-il surjectif ?
  \item Montrer que tous les isomorphismes de corps de ${ K(X)}$ sont de cette forme.
\end{enumerate}
\finenonce{003270}



\finexercice 


\exercice{3271, quercia, 2010/03/08}

\enonce{003271}{Multiplicit{\'e} des p{\^o}les}
Soient $F, G_0, \dots, G_{n-1} \in { K(X)}$ telles que
$F^n + G_{n-1}F^{n-1} + \dots + G_0 = 0$.

Montrer que l'ensemble des p{\^o}les de $F$ est inclus dans la r{\'e}union des ensembles
des p{\^o}les des $G_i$.
\finenonce{003271}



\finexercice 


\exercice{3272, quercia, 2010/03/08}

\enonce{003272}{Ensemble image d'une fonction rationelle}
\label{Imf}
Soit $F \in {\C(X)}$. {\'E}tudier $F( \C\setminus \{\text{p{\^o}les}\} )$.
\finenonce{003272}


\finexercice 


\exercice{3273, quercia, 2010/03/08}

\enonce{003273}{$F \circ G$ est un polyn{\^o}me}
Trouver tous les couples $(F,G) \in \bigl({\C(X)}\bigr)^2$ tels que
$F\circ G \in {\C[X]}$ (utiliser l'exercice \ref{Imf}).

\finenonce{003273}


\finexercice 


\exercice{3274, quercia, 2010/03/08}

\enonce{003274}{Fractions invariantes}
\begin{enumerate}
  \item Soit $F \in {\C(X)}$ telle que $F(e^{2i\pi/n}X) = F(X)$.
    Montrer qu'il existe une unique fraction $G \in {\C(X)}$ telle que $F(X) = G(X^n)$.
  \item Application : Simplifier
    $\sum_{k=0}^{n-1} \frac {X+e^{2ik\pi/n}}{X-e^{2ik\pi/n}}$.
\end{enumerate}
\finenonce{003274}


\finexercice 


\exercice{3275, quercia, 2010/03/08}

\enonce{003275}{Fractions invariantes}
Soit $H = \{ F \in { K(X)}$ tel que $F(\scriptstyle X) = F(\frac 1X) \}$.
\begin{enumerate}
  \item Montrer que : $F \in H \Leftrightarrow \exists\ G \in { K(X)}$ tel que
    $F(\scriptstyle X) = G\Bigl(\scriptstyle X+\frac 1X\Bigr)$.
  \item Montrer que $H$ est un sous-corps de ${ K(X)}$.
  \item Que vaut $\dim_H({ K(X)})$ ? Donner une base de ${ K(X)}$ sur $H$.
\end{enumerate}
\finenonce{003275}



\finexercice 


\exercice{3276, quercia, 2010/03/08}

\enonce{003276}{Formule de Taylor}
Soit $F \in { K(X)}$ d{\'e}finie en $a \in  K$.
D{\'e}montrer qu'il existe une fraction $G_n$ d{\'e}finie en $a$ telle que :
$$
   F(X) = F(a) + (X-a)F'(a) + \dots + (X-a)^{n-1} \frac {F^{(n-1)}(a)}{(n-1)!}
   + (X-a)^nG_n(X).
$$
\finenonce{003276}



\finexercice 


\exercice{3277, quercia, 2010/03/08}

\enonce{003277}{D{\'e}riv{\'e}e de $1/(x^2+1)$}
Soit $F = \frac 1{X^2+1}$. Montrer qu'il existe un polyn{\^o}me $P_n \in {\Z_n[X]}$
tel que $F^{(n)} = \frac {P_n}{(X^2+1)^n}$.

Montrer que les racines de $P_n$ sont r{\'e}elles et simples.
\finenonce{003277}



\finexercice 


\exercice{3278, quercia, 2010/03/08}

\enonce{003278}{Fractions de degr{\'e} n{\'e}gatif}
Soit $A = \{ F \in { K(X)}$ tels que $\deg F \le 0 \}$. D{\'e}montrer que $A$ est une sous-alg{\`e}bre
de ${ K(X)}$.

Chercher ses id{\'e}aux.
\finenonce{003278}


\finexercice 


\exercice{3279, quercia, 2010/03/08}
\enonce{003279}{D{\'e}compositions pratiques des fractions rationnelles}

          %+--------------------------------------------------------+
          %|  D{\'e}compositions pratiques des fractions rationnelles   |
          %+--------------------------------------------------------+


\def \FFF#1/#2 {\frac {#1}{#2}}
\def \AAA#1{(x-1)^{#1}}      
\def \BBB#1{(x+1)^{#1}}      
\def \CCC#1{(x^2+1)^{#1}}    
\def \DDD#1{(x^2+x+1)^{#1}}  

                        %+----------------------------+
                        %|  {\'e}l{\'e}ments de 1{\`e}re esp{\`e}ce   |
                        %+----------------------------+

\bigskip 
\textbf{{\'E}l{\'e}ments de 1{\`e}re esp{\`e}ce}

\begin{align*}
\FFF 1/(x^2-1)^5         =& \FFF1/32\AAA5 -\FFF5/64\AAA4 +\FFF15/128\AAA3
-\FFF35/256\AAA2 +\FFF35/256({x-1}) \\ -&\FFF35/256({x+1}) -\FFF35/256\BBB2
-\FFF15/128\BBB3 -\FFF15/64\BBB4 -\FFF1/32\BBB5 \cr  
\FFF \CCC2/\AAA6             =& \FFF4/\AAA6 + \FFF8/\AAA5 + \FFF8/\AAA4 + \FFF4/\AAA3 + \FFF1/\AAA2 \cr
\FFF x^3+x+1/x^4\AAA3      =& -\FFF1/x^4 -\FFF4/x^3 -\FFF9/x^2 -\FFF17/x +\FFF3/\AAA3 -\FFF8/\AAA2 +\FFF17/{{x-1}} \cr
\FFF (x^2-x+1)^2/x^2\AAA2  =& 1 + \FFF1/x^2 + \FFF1/\AAA2 \cr
\FFF x^2/(x^2-1)^2       =& \FFF1/4\AAA2 + \FFF1/4({x-1}) + \FFF1/4\BBB2 + \FFF1/4({x+1}) \cr
\end{align*}

\bigskip
                                  %+----------+
                                  %|  x^2+1   |
                                  %+----------+

\textbf{Du type $x^2+1$}

\begin{align*}
\FFF x^2/(x^2+1)^2       =& \FFF -1/\CCC2 + \FFF 1/{{x^2+1}} \cr
\FFF x/(x^4-1)^2         =& \FFF1/16\AAA2 -\FFF1/8({x-1}) - \FFF1/16\BBB2 -\FFF1/8({x+1}) + \FFF x/4\CCC2 + \FFF x/4({x^2+1}) \cr
\FFF x/({x-1})\CCC2        =& \FFF 1/4({x-1}) + \FFF 1-x/2\CCC2 - \FFF x+1/4({x^2+1}) \cr
\FFF x^6/\CCC2\BBB2          =& 1 +\FFF1/4\BBB2 - \FFF1/{{x+1}} + \FFF x/2\CCC2 - \FFF {x+1/4}/{{x^2+1}} \cr
\FFF x^6/({x^2+1})\AAA3    =& x + 3 + \FFF x-1/4({x^2+1}) + \FFF 1/2\AAA3 + \FFF 5/2\AAA2 + \FFF 19/4({x-1})  \cr
\end{align*}

\bigskip
                                 %+------------+
                                 %|  x^2+x+1   |
                                 %+------------+

\textbf{Du type $x^2+x+1$}

\begin{align*}
\FFF x/x^4+x^2+1         =& \FFF 1/2(x^2-x+1) - \FFF 1/2(x^2+x+1) \cr
\FFF x^4+1/x^4+x^2+1     =& 1 + \FFF x/2(x^2+x+1) - \FFF x/2(x^2-x+1) \cr
\FFF x^4+1/x^2\DDD2       = & \FFF1/x^2 -\FFF2/x -\FFF1/\DDD2 + \FFF2x+2/{{x^2+x+1}} \cr
\FFF 3x^5-5x^4+4x^2-11x+1/\DDD6  =& -\FFF 23x+6/\DDD6 + \FFF 13x+18/\DDD5 + \FFF 3x-11/\DDD4 \cr
\end{align*}

\bigskip
                     %+-----------------------------------+
                     %|  autres {\'e}l{\'e}ments de 2{\`e}me esp{\`e}ce   |
                     %+-----------------------------------+

\textbf{Autres {\'e}l{\'e}ments de 2{\`e}me esp{\`e}ce}

\begin{align*}
\FFF x^8/x^6-1  =& x^2 + \FFF1/6 \biggl( \FFF1/x-1 -\FFF1/x+1 +\FFF2x+1/x^2+x+1 -\FFF2x-1/x^2-x+1 \biggr) \cr
\FFF 1/x^4+1   =& \FFF1/2\sqrt2  \biggl( \FFF x+\sqrt2/x^2+x\sqrt2+1 - \FFF x-\sqrt2/x^2-x\sqrt2+1 \biggr) \cr
\FFF x/x^4+1   =& \FFF 1/2\sqrt2 \biggl( \FFF 1/x^2-x\sqrt2+1 - \FFF 1/x^2+x\sqrt2+1 \biggr) \cr
\FFF 1/x^5+1   =& \FFF 1/5({x+1}) - \FFF1/5 \biggl( \FFF {\omega x-2}/{x^2-\omega x+1}
                   + \FFF {\omega'x-2}/{x^2-\omega'x+1} \biggr),\quad
                   \omega = \FFF1+\sqrt5/2 ,\omega' = \FFF1-\sqrt5/2 \cr
\end{align*}

\bigskip
                           %+-----------------------+
                           %|  racines de l'unit{\'e}   |
                           %+-----------------------+

\textbf{Racines de l'unit{\'e}}

\begin{align*}
\FFF x^n+1/x^n-1         =& 1 + 2\sum_{k=0}^{n-1} \FFF \omega^k/n(x-\omega^k) ,\quad \omega = e^{2i\pi/n}\cr
\FFF 1/x^n-1             =& \sum_{k=1 ; 2k \ne n}^{n-1} \FFF 2x\cos\alpha_k-2/n(x^2-2x\cos\alpha_k+1)
                            + \FFF1/n(x-1) \ \left[ - \FFF1/n(x+1) \text{ si n est pair} \right],
                            \quad \alpha_k = \FFF2k\pi/n \cr
\sum_{k=0}^{n-1} \FFF 1/x-\omega^k        =& \FFF nx^{n-1}/x^n-1 ,\quad \omega = e^{2i\pi/n}\cr
\sum_{k=0}^{n-1} \FFF 1/(x-\omega^k)^2    =& \FFF nx^{2n-2}+n(n-1)x^{n-2}/(x^n-1)^2
                                             ,\quad \omega = e^{2i\pi/n}\quad(\text{d{\'e}riv{\'e}e})\cr
\end{align*}

\bigskip
                        %+-----------------------------+
                        %|  polyn{\^o}mes de Tchebychev    |
                        %+-----------------------------+

\textbf{Polyn{\^o}mes de Tchebychev}

\begin{align*}
\FFF 1/{\cos(n\arccos x)}  =& \FFF1/n \sum_{k=0}^{n-1} \FFF (-1)^k\sin\beta_k/x-\cos\beta_k ,\quad \beta_k = \FFF(2k+1)\pi/2n \cr
\tan(n\arctan x)         =& \FFF1/n \sum_{k=0 ; 2k\ne n-1}^{n-1} \FFF1/\cos^2\beta_k(\tan\beta_k-x)
                              \left[ + \FFF x/n \ \text{si n est impair} \right],\quad \beta_k = \FFF(2k+1)\pi/2n \cr
\end{align*}

\bigskip
                                 %+-----------+
                                 %|  divers   |
                                 %+-----------+

\textbf{Divers}

\begin{align*}
\FFF x^{2n}/{\CCC n}  =& \sum_{k=0}^n \FFF (-1)^kC_n^k/{\CCC k} \cr
\FFF 1/(x^2-1)^n    =& \sum_{k=0}^{n-1} \FFF \Gamma_n^k/2^{n+k} \biggl( \FFF (-1)^k/(x-1)^{n-k} + \FFF (-1)^n/(x+1)^{n-k} \biggr) \cr
\FFF 1/{\CCC n}      = & \sum_{k=0}^{n-1} \FFF (-1)^n\Gamma_n^k/2^{n+k} \biggl( \FFF {i^{k+n}}/(x-i)^{n-k} + \FFF {(-i)^{k+n}}/(x+i)^{n-k} \biggr) \cr
\FFF n!/(x+1)(x+2)\dots(x+n)  =& \sum_{k=1}^n \FFF (-1)^{k-1}kC_n^k/x+k \cr
\FFF x^2/x^4-2x^2\cos\alpha+1 =& \FFF1/4\cos(\alpha/2) \biggl(\FFF x/x^2-2x\cos(\alpha/2)+1 - \FFF x/x^2+2x\cos(\alpha/2)+1 \biggr), \quad \alpha\not\equiv0(\mathrm{mod}\,\pi)\cr
\end{align*}


\finenonce{003279}



\finexercice 

\exercice{3280, quercia, 2010/03/08}

\enonce{003280}{Ensi PC 1999}
D{\'e}composer en {\'e}l{\'e}ments simples sur $\R$ puis sur $\C$~:
$\frac1{(X^2+2X+1)(X^3-1)}$.
\finenonce{003280}


\finexercice 


\exercice{3281, quercia, 2010/03/08}

\enonce{003281}{Calcul de d{\'e}riv{\'e}es}
Calculer les d{\'e}riv{\'e}es $p$-i{\`e}mes des fractions suivantes :
\begin{enumerate}
  \item $\frac 1{X(X+1)\dots(X+n)}$.
  \item $\frac 1{X^2-2X\cos\alpha+1}$ ($\alpha\not\equiv0(\mathrm{mod}\,{\pi})$).
  \item $\frac 1{X^2-2X\sh\alpha-1}$ ($\alpha\in\R$).
\end{enumerate}
\finenonce{003281}



\finexercice 


\exercice{3282, quercia, 2010/03/08}

\enonce{003282}{Sommation de s{\'e}ries}
A l'aide de d{\'e}composition en {\'e}l{\'e}ments simples, calculer :
\begin{enumerate}
  \item $\sum_{n=1}^\infty \frac 1{n(n + 1)}$.      
  \item $\sum_{n=1}^\infty \frac 1{n(n + 1)(n+2)}$.
  \item $\sum_{n=1}^\infty \frac n{n^4 + n^2 + 1}$. 
\end{enumerate}
\finenonce{003282}


\finexercice 


\exercice{3283, quercia, 2010/03/08}

\enonce{003283}{Partie polaire pour un p{\^o}le d'ordre 2}
Soit $F(X) = \frac 1{R(X)} = \frac 1{(X-a)^2Q(X)}$ avec $Q(a) \ne 0$.
Chercher la partie polaire de $F$ en $a$ en fonction de $Q$ puis en fonction
de~$R$.
\finenonce{003283}


\finexercice 


\exercice{3284, quercia, 2010/03/08}

\enonce{003284}{}
Soient $a_1,\dots,a_n \in  K$ distincts et $P = (X-a_1)\dots(X-a_n)$.
\begin{enumerate}
  \item D{\'e}composer en {\'e}l{\'e}ments simples la fraction $\frac {(1+X^2)^n}{P^2}$.

  \item Montrer que les coefficients des $\frac1{X-a_i}$ sont tous nuls si et seulement si :
    $(1+X^2)P'' - 2nXP' + n(n+1)P = 0$.
\end{enumerate}
\finenonce{003284}


\finexercice


\exercice{3285, quercia, 2010/03/08}

\enonce{003285}{$P$ {\`a} racines $x_i$ simples $ \Rightarrow \sum x_i^k/P'(x_i) = 0$}
Soit $P \in {\C_n[X]}$ $(n \ge 2)$ ayant $n$ racines distinctes :
$x_1,\dots,x_n$.
\begin{enumerate}
  \item D{\'e}montrer que $\sum_{i=1}^n \frac 1{P'(x_i)} = 0$.
  \item Calculer $\sum_{i=1}^n \frac {x_i^k}{P'(x_i)}$ pour $0 \le k \le n-1$.
\end{enumerate}
\finenonce{003285}


\finexercice 


\exercice{3286, quercia, 2010/03/08}

\enonce{003286}{Les racines de $P'$ sont des barycentres des racines de $P$}
Soit $P \in {\C[X]}$ de racines $x_1,x_2,\dots,x_n$ avec les multiplicit{\'e}s
$m_1,m_2,\dots,m_n$.
\begin{enumerate}
  \item D{\'e}composer en {\'e}l{\'e}ments simples $\frac {P'}P$.
  \item En d{\'e}duire que les racines de $P'$ sont dans l'enveloppe convexe
    de $x_1,\dots,x_n$.
\end{enumerate}
\finenonce{003286}


\finexercice 


\exercice{3287, quercia, 2010/03/08}

\enonce{003287}{$F'(X)/F(X) = \dots$}
Soient $a_1, \dots, a_n \in  K$ distincts et
$\alpha_1, \dots, \alpha_n \in  K$.
Existe-t-il $F \in { K(X)}$ telle que :
$\frac {F'(X)}{F(X)} = \sum_{k=1}^n \frac {\alpha_k}{X-a_k}$ ?
\finenonce{003287}



\finexercice 


\exercice{3288, quercia, 2010/03/08}

\enonce{003288}{$F(X+1) - F(X) = \dots$}
Trouver les fractions $F \in {\R(X)}$ telles que :
$F(X+1) - F(X) = \frac {X+3}{X(X-1)(X+1)}$.
\finenonce{003288}


\finexercice 


\exercice{3289, quercia, 2010/03/08}

\enonce{003289}{Inversion de la matrice $(1/(a_i-b_j))$}
Soient $a_1, \dots, a_n, b_1, \dots, b_n$, et $c$ des scalaires distincts.
On note $A$ la matrice carr{\'e}e $\left( \frac 1{a_i-b_j} \right)$ et
$B$ la matrice colonne $\left( \frac 1{a_i-c} \right)$.
Montrer que l'{\'e}quation $AX = B$ poss{\`e}de une solution unique en consid{\'e}rant une
fraction rationnelle bien choisie.
\finenonce{003289}


\finexercice 


\exercice{3290, quercia, 2010/03/08}

\enonce{003290}{Racines de $(X^2+1)PP' + X(P^2+P'^2)$}
Soit $P \in {\R[X]}$ ayant $n$ racines positives distinctes (entre autres).

Factoriser le polyn{\^o}me $Q = (X^2+1)PP' + X(P^2+P'^2)$ en deux termes,
faire appara{\^\i}tre $\frac {P'}P$, et D{\'e}montrer que $Q$ admet au moins $2n-2$ racines positives.
\finenonce{003290}


\finexercice 


\exercice{3291, quercia, 2010/03/08}

\enonce{003291}{In{\'e}galit{\'e}}
Soit $P\in\R[X]$ unitaire de degr{\'e}~$n$ et $Q(X) = X(X-1)\dots(X-n)$.

Calculer $\sum_{k=0}^n\frac{P(k)}{\prod_{i\ne k}(k-i)}$ et en d{\'e}duire
l'existence de $k\in{[[0,n]]}$ tel que $|P(k)|\ge \frac{n!}{2^n}$.

\finenonce{003291}


\finexercice 


\exercice{3292, quercia, 2010/03/08}

\enonce{003292}{ENS MP 2002}
  \item Soit $P\in\C[X]$ admettant deux racines distinctes et tel que $P''$
    divise $P$. Montrer que $P$ est {\`a} racines simples.

  \item Soit $P\in\R[X]$ admettant deux racines r{\'e}elles distinctes, et tel que
    $P''$ divise $P$. Montrer que $P$ est scind{\'e} sur~$\R$ et {\`a} racines simples.
\finenonce{003292}


\finexercice 


\exercice{3293, quercia, 2010/03/08}

\enonce{003293}{Division de $X^3-1$ par $X^2+1$}
\begin{enumerate}
  \item Effectuer la division suivant les puissances croissantes de
    $X^3-1$ par $X^2+1$ {\`a} l'ordre 3.
  \item En d{\'e}duire une primitive de $f$ : $x \mapsto \frac{x^3-1}{x^4(x^2+1)}$.
\end{enumerate}
\finenonce{003293}


\finexercice 


\exercice{3294, quercia, 2010/03/08}

\enonce{003294}{Division de $1$ par $(1-X)^2$}
\begin{enumerate}
  \item Effectuer la division suivant les puissances croissantes {\`a} un ordre $n$
    quelconque de $1$ par $(1-X)^2$.
  \item En d{\'e}duire $1 + 2\cos\theta + 3\cos2\theta + \dots + n\cos (n-1)\theta$,
    $n \in \N^*$, $\theta \in \R$.
\end{enumerate}
\finenonce{003294}


\finexercice 


\exercice{3295, quercia, 2010/03/08}

\enonce{003295}{Division de $1-X^2$ par $1 - 2X \cos t + X^2$}
\begin{enumerate}
  \item Effectuer la division suivant les puissances croissantes {\`a} un ordre
    queclonque de $1-X^2$ par $1 - 2X\cos\theta + X^2$.
  \item En d{\'e}duire la valeur de $1 + 2\sum_{k=1}^n \cos k\theta$,
    ($\theta \not\equiv 0 (\mathrm{mod}\,{2\pi})$).
\end{enumerate}
\finenonce{003295}


\finexercice 


\exercice{3296, quercia, 2010/03/08}

\enonce{003296}{Coefficients de B{\'e}zout}
Soient $P = 1 + 2X + 3X^2 + 3X^3 + 2X^4 + X^5$ et $Q = X^5$.
\begin{enumerate}
  \item V{\'e}rifier que $P$ et $Q$ sont premiers entre eux.
  \item Trouver $U,V\in{ K[X]}$ tels que $UP + VQ = 1$ (utiliser une division
    suivant les puissances croissantes).
\end{enumerate}
\finenonce{003296}


\finexercice 


\exercice{5335, rouget, 2010/07/04}
\enonce{005335}{}
Décomposer en éléments simples dans $C(X)$ les fractions rationnelles suivantes

$$\begin{array}{lll}
1)\;\frac{X^2+3X+5}{X^2-3X+2}&2)\;\frac{X^2+1}{(X-1)(X-2)(X-3)}&3)\;\frac{1}{X(X-1)^2}\\
4)\;\frac{X^2+1}{(X-1)^2(X+1)^2}&5)\;\frac{1}{(X-2)^3(X+2)^3}&6)\;\frac{X^6}{(X^3-1)^2}\\
7/\;\frac{1}{X^6+1}&8)\;\frac{X^2+3}{X^5-3X^4+5X^3-7X^2+6X-2}&9)\;\frac{X}{(X^2+1)^3(X^2-1)}\\
10)\;\frac{X^6+1}{X^5-X^4+X^3-X^2+X-1}&11)\;\frac{X^7+1}{(X^2+X+1)^3}&12)\;\frac{X^2+1}{X(X-1)^4(X^2-2)^2}\\ 
13)\;\frac{1}{(X+1)^7-X^7-1}.
\end{array}$$
\finenonce{005335}


\finexercice


\exercice{5336, rouget, 2010/07/04}
\enonce{005336}{}
Décomposer en éléments simples dans $C(X)$ les fractions rationnelles suivantes

$$\begin{array}{lll}
1)\;\frac{1}{X^n-1}&2)\;\frac{1}{(X-1)(X^n-1)}&3)\;\frac{n!}{(X-1)(X-2)...(X-n)}\\
4)\;\frac{X^2}{X^4-2X^2\cos(2a)+1}&5)\;\frac{1}{X^{2n}+1}.
\end{array}$$
\finenonce{005336}


\finexercice


\exercice{5337, rouget, 2010/07/04}
\enonce{005337}{}
Soit $U_n$ l'ensemble des racines $n$-ièmes de l'unité dans $\Cc$. Ecrire sous forme d'une fraction rationnelle (ou encore réduire au même dénominateur) $F=\sum_{\omega\in U_n}^{}\frac{\omega X+1}{\omega^2X^2+\omega X+1}$.
\finenonce{005337}


\finexercice


\exercice{5338, rouget, 2010/07/04}
\enonce{005338}{}
Soit $F=\frac{P}{Q}$ où $P$ et $Q$ sont des polynômes tous deux non nuls et premiers entre eux. Montrer que $F$ est paire si et seulement si $P$ et $Q$ sont pairs. Etablir un résultat analogue pour $F$ impaire.
\finenonce{005338}


\finexercice


\exercice{5339, rouget, 2010/07/04}
\enonce{005339}{}
Montrer que $(\frac{1}{X-a})_{a\in\Cc}$ est libre dans $K(X)$.
\finenonce{005339}


\finexercice


\exercice{5340, rouget, 2010/07/04}
\enonce{005340}{}
Calculer la dérivée $n$-ième de $\frac{1}{X^2+1}$.
\finenonce{005340}


\finexercice


\exercice{5341, rouget, 2010/07/04}
\enonce{005341}{}
On pose $P=a(X-x_1) ... (X-x_n)$ où les $x_i$ sont des complexes non nécessairement deux à deux distincts et $a$ est un complexe non nul.

Calculer $\frac{P'}{P}$. De manière générale, déterminer la décomposition en éléments simples de $\frac{P'}{P}$ quand $P$ est un polynôme scindé. Une application~:~déterminer tous les polynômes divisibles par leur dérivées.
\finenonce{005341}


\finexercice


\exercice{6964, blanc-centi, 2014/04/08}
\video{f93zdMxazWw}
\enonce{006964}{}

Existe-t-il une fraction rationnelle $F$ telle que 
$$\big(F(X)\big)^2=(X^2+1)^3 \ \ ?$$

\finenonce{006964} 



\finexercice

\exercice{6965, blanc-centi, 2014/04/08}
\video{7jn7oh2DdnQ}
\enonce{006965}{}
Soit $F=\frac{P}{Q}$ une fraction rationnelle écrite sous forme irréductible. 
On suppose qu'il existe une fraction rationnelle $G$ telle que 
$$G\left(\frac{P(X)}{Q(X)}\right)=X$$
\begin{enumerate}
\item Si $G=\frac{a_nX^n+\cdots+a_1X+a_0}{b_nX^n+\cdots+b_1X+b_0}$, montrer que $P$ divise $(a_0-b_0X)$ et que $Q$ divise $(a_n-b_nX)$.
\item En déduire que $F=\frac{P}{Q}$ est de la forme $F(X)=\frac{aX+b}{cX+d}$.
\item Pour $Y=\frac{aX+b}{cX+d}$, exprimer $X$ en fonction de $Y$. En déduire l'expression de $G$.
\end{enumerate}
\finenonce{006965} 


\finexercice



\exercice{6966, blanc-centi, 2014/04/08}
\video{BDJWi5Od7uA}
\enonce{006966}{}
Soit $n\in\Nn^*$ et $P(X)=c(X-a_1)\cdots(X-a_n)$ 
(où les $a_i$ sont des nombres complexes et où $c\not=0$).
\begin{enumerate}
\item Exprimer à l'aide de $P$ et de ses dérivées les sommes suivantes:
$$\sum_{k=1}^n\frac{1}{X-a_k}\quad\quad
\sum_{k=1}^n\frac{1}{(X-a_k)^2}\quad\quad
\sum_{\substack{1\le k,\ell\le n \\ k\not= \ell}}\frac{1}{(X-a_k)(X-a_\ell)}$$
\item Montrer que si $z$ est racine de $P'$ mais pas de $P$, alors il existe 
$\lambda_1,\hdots,\lambda_n$ des réels positifs ou nuls tels que 
$\sum_{k=1}^n\lambda_k=1$ et $z=\sum_{k=1}^n\lambda_ka_k$. 
Si toutes les racines de $P$ sont réelles, que peut-on en déduire sur les racines de $P'$ ?
\end{enumerate}
\finenonce{006966} 


\finexercice

\exercice{6967, exo7, 2014/04/08}
\video{KHEyahxXAKk}
\enonce{006967}{}
% De 444, cousquer
Décomposer les fractions suivantes en éléments simples sur $\Rr$, 
par identification des coefficients.
\begin{enumerate}
\item $F=\frac{X}{X^2-4}$
\item $G=\frac{X^3-3X^2+X-4}{X-1}$
\item $H=\frac{2X^3+X^2-X+1}{X^2-2X+1}$
\item $K=\frac{X+1}{X^4+1}$
\end{enumerate}
\finenonce{006967} 


\finexercice



\exercice{6968, exo7, 2014/04/08}
\video{cMc2VRPKaKI}
\enonce{006968}{}
Décomposer les fractions suivantes en éléments simples sur $\Rr$, 
en raisonnant par substitution pour obtenir les coefficients.
\begin{enumerate}
\item % De 444, cousquer
$F=\frac{X^5+X^4+1}{X^3-X}$

\item % De 444, cousquer
$G=\frac{X^3+X+1}{(X-1)^3(X+1)}$

\item % De 444, cousquer
$H=\frac{X}{(X^2+1)(X^2+4)}$

\item % De 445, cousquer
$K=\frac{2X^4+X^3+3X^2-6X+1}{2X^3-X^2}$
\end{enumerate}

\finenonce{006968} 


\finexercice



\exercice{6969, exo7, 2014/04/08}
\video{S53f12bRBhE}
\enonce{006969}{}
Décomposer les fractions suivantes en éléments simples sur $\Rr$.
\begin{enumerate}
\item \`A l'aide de divisions euclidiennes successives :
% De 447, cousquer
$$F=\frac{4X^6-2X^5+11X^4-X^3+11X^2+2X+3}{X(X^2+1)^3}$$

\item \`A l'aide d'une division selon les puissances croissantes :
% De 446, cousquer
$$G=\frac{4X^4-10X^3+8X^2-4X+1}{X^3(X-1)^2}$$

\item Idem pour :
% De 444, cousquer
$$H=\frac{X^4+2X^2+1}{X^5-X^3}$$

\item A l'aide du changement d'indéterminée $X=Y+1$ :
% De 444, cousquer
$$K=\frac{X^5+X^4+1}{X(X-1)^4}$$

\end{enumerate}
\finenonce{006969} 


\finexercice


\exercice{6970, exo7, 2014/04/08}
\video{ocbw3l8Wwl0}
\enonce{006970}{}
% De 444, cousquer
\begin{enumerate}
\item Décomposer les fractions suivantes en éléments simples sur $\Cc$.
$$\frac{(3-2 i )X-5+3 i }{X^2+ i  X+2} \qquad\qquad \frac{X+ i }{X^2+ i } \qquad\qquad \frac{2X}{(X+ i )^2}$$

\item Décomposer les fractions suivantes en éléments simples sur $\Rr$, puis sur $\Cc$.
$$\frac{X^5+X+1}{X^4-1} \qquad\qquad \frac{X^2-3}{(X^2+1)(X^2+4)} \qquad\qquad \frac{X^2+1}{X^4+1}$$
\end{enumerate}
\finenonce{006970} 


\finexercice

\exercice{6971, blanc-centi, 2014/04/08}
\video{FPyfOQL32Hg}
\enonce{006971}{}
On pose $Q_0=(X-1)(X-2)^2$, $Q_1=X(X-2)^2$ et $Q_2=X(X-1)$. 
\`A l'aide de la décomposition en éléments simples de $\frac{1}{X(X-1)(X-2)^2}$, 
trouver des polynômes $A_0,\ A_1,\ A_2$ tels que $A_0Q_0+A_1Q_1+A_2Q_2=1$. 
Que peut-on en déduire sur $Q_1$, $Q_2$ et $Q_3$?
\finenonce{006971} 


\finexercice


\exercice{6972, blanc-centi, 2014/04/08}
\enonce{006972}{}
Soit $T_n(x)=\cos\big(n \arccos(x)\big)$ pour $x\in [-1,1]$.
\begin{enumerate}
  \item 
  \begin{enumerate}
    \item Montrer que pour tout $\theta\in[0,\pi]$, $T_n(\cos\theta)=\cos(n\theta)$.
    \item Calculer $T_0$ et $T_1$.
    \item Montrer la relation de récurrence $T_{n+2}(x) = 2xT_{n+1}(x)-T_n(x)$, pour tout $n \ge0$.
    \item En déduire que $T_n$ une fonction polynomiale de degré $n$.
  \end{enumerate}

  \item Soit $P(X)=\lambda(X-a_1)\cdots(X-a_n)$ un polynôme, où les $a_k$ sont deux à deux distincts 
et $\lambda\not=0$. Montrer que 
$$\frac{1}{P(X)}=\sum_{k=1}^n\frac{\frac{1}{P'(a_k)}}{X-a_k}$$

\item Décomposer $\frac{1}{T_n}$ en éléments simples.
\end{enumerate}
\finenonce{006972} 


\finexercice

\section{105.05 Définition, degré, produit}

\section{105.99 Autre}


\exercice{424, cousquer, 2003/10/01}

\enonce{000424}{}
 Montrer que pour tout $n\in\mathbb{N}^{\ast}$ 
il existe un polynôme~$P_n$ et
un seul tel que
$$\forall \theta \in \mathbb{R},\; P_n(2\cos \theta )=2\cos n\theta.$$
Montrer que $P_n$ est unitaire et que ses coefficients sont entiers.
En déduire les $r$~rationnels tels que $\cos r\pi $ soit rationnel.
\finenonce{000424}



\finexercice



\exercice{425, cousquer, 2003/10/01}

\enonce{000425}{}
Déterminer, s'il en existe, tous les idéaux $J$ de $\mathbb{R}[X]$ 
tels que~:
$I(P) \subset J \subset\mathbb{R}[X]$, avec $I(P)$ idéal engendré par $P$ dans
les cas suivants~:
$$ P=X^2+X+1, \quad P=X^2+2X+1, \quad P=X^3+3X-4.$$
\finenonce{000425}



\finexercice



\exercice{426, cousquer, 2003/10/01}
\enonce{000426}{}
 Trouver un polyn\^ome~$P$ de degr\'e~$\leq2$ tel que 
$$P(1)=-2 \quad \mbox{et} \quad P(-2)=3 \quad \mbox{et} \quad P(0)=-1$$

\finenonce{000426} 
\finexercice


\exercice{427, cousquer, 2003/10/01}
\video{59DItk_pxuc}
\enonce{000427}{}

Trouver le polynôme $P$ de degré inférieur ou égal à $3$ tel que :
$$P(0)=1\quad\text{et}\quad P(1)=0\quad\text{et}\quad P(-1)=-2\quad\text{et}\quad P(2)=4.$$
\finenonce{000427} 


\finexercice

\exercice{428, ridde, 1999/11/01}

\enonce{000428}{}
Trouver les polynômes $P$ de $\Rr[X]$ tels que $\forall k \in \Zz \, \,
\int_{k}^{k + 1}P (t)dt  = k + 1$ (on pourra utiliser le polynôme
$Q (x) = \int_0^x P (t)dt$).
\finenonce{000428}



\finexercice



\exercice{429, ridde, 1999/11/01}

\enonce{000429}{}
Soit $ (P_0, P_1, \ldots, P_n)$ une famille de polynômes de $\mathbb{K}[X]$ telle
que $\forall k \in \left\{ 0, \ldots, n\right\} \, \, \text{deg}P_k  = k$.
Montrer \`a l'aide d'une r\'ecurrence soigneuse que cette famille est libre.
\finenonce{000429}



\finexercice



\exercice{430, ridde, 1999/11/01}

\enonce{000430}{}
Soit $n \in \Nn^*$ fix\'e et ${\Delta} : {\Rr_n [X]} \mapsto {\Rr_n [X]}$, ${P (X)} \mapsto {P (X + 1)
-P (X)}$.
\begin{enumerate}
\item Montrer que $\Delta$ est lin\'eaire, i.e. que $\forall (a, b) \in \Rr^2$ et
$ (P, Q) \in \Rr_n [X] \, \, \Delta (aP  + bQ) = a\Delta (P) + b\Delta (Q)$.
\item D\'eterminer $\ker (\Delta) = \left\{ P \in \Rr_n[X] / \Delta (P) = 0\right\}$.
\item Soient $H_0 = 1$ et pour $k \in \left\{ 1, \ldots, n\right\} \, \,
H_k = \dfrac 1{k ! } X (X-1)\ldots (X-k + 1)$. Calculer $\Delta (H_k)$.
\item Soit $Q \in \Rr_{n-1}[X]$. Comment trouver $P \in \Rr_{n}[X]$ tel que
$\Delta (P) = Q$. %(on pourra utiliser l'exercice $14$).
\item D\'eterminer $P$ pour $Q = X^2$ tel que $P (1) = 0$.
\item En d\'eduire la somme $1^2 + 2^2 + \ldots  + n^2$.
\end{enumerate}
\finenonce{000430}



\finexercice



\exercice{431, gourio, 2001/09/01}

\enonce{000431}{}
R\'{e}soudre l'\'{e}quation d'inconnue $P\in \Cc[X]:P(X+1)P(X)=-P(X^{2}).$
\finenonce{000431}



\finexercice



\exercice{432, gourio, 2001/09/01}

\enonce{000432}{}
Soit $(P,Q)\in \Rr_{n}[X]^{2} $ tels que $\exists (a,A)\in
(\Rr^{+*})^{2},\forall x\in ]-a,a[,\left| P(x)-Q(x)\right| \leq A\left|x^{n+1}\right| .$
Que dire de $P$ et $Q$ ?
\finenonce{000432}



\finexercice



\exercice{433, gourio, 2001/09/01}

\enonce{000433}{}
Soient $W_{n}=(X^{2}-1)^{n}, \ L_{n}=\frac{1}{2^{n}n!}W_{n}^{(n)}.$
\begin{enumerate}
\item
Donner le degr\'{e} de $L_{n}$, son coefficient dominant, sa parit\'{e},
calculer $L_{n}(1).$
Donner $L_{0},L_{1},L_{2}.$

\item D\'{e}montrer : $\forall n\geq 1,(X^{2}-1)W_{n}^{^{\prime }}=2nXW_{n},$ en
d\'{e}duire :
$$\forall n\in \Nn,(X^{2}-1)L_{n}^{^{\prime \prime }}+
2XL_{n}^{^{\prime}}-n(n+1)L_{n}=0.$$

\item Montrer ensuite : $\forall n\geq 1,L_{n}^{\prime }=XL_{n-1}^{^{\prime
}}+nL_{n-1},$ puis $nL_{n}=XL_{n}^{^{\prime }}-L_{n-1}^{^{\prime }}.$

\item Montrer enfin que les polyn\^{o}mes $L_{n} $ peuvent \^{e}tre d\'{e}finis
par la r\'{e}currence :
$$(n+1)L_{n+1}=(2n+1)XL_{n}-nL_{n-1}.$$
\end{enumerate}
\finenonce{000433}



\finexercice



\exercice{434, gourio, 2001/09/01}

\enonce{000434}{}
Montrer que si $n\geq 3,$ l'\'{e}quation $x^{n}+y^{n}=z^{n}$ n'a pas de
solution non triviale (i.e. $xyz\neq 0$) dans $\Cc[X].$

\emph{Indication} : on peut supposer $x,y,z,$ sans facteurs communs.\ D\'{e}river la
relation, la multiplier par $z$, \'{e}tudier le degr\'{e}.
\finenonce{000434}



\finexercice



\exercice{435, gourio, 2001/09/01}

\enonce{000435}{}
Soit $n\in \Nn^{*},P\in \Cc[X]$ de degr\'{e} $n$, avec $P(0)=1,P(1)=0,$ montrer :
$$\sup\limits_{\left| z\right| =1}\left| P(z)\right| \geq 1+\frac{1}{n}.$$

\emph{Indication} : $w_{k}=e^{\frac{2ik\pi }{n+1}},$ montrer $\sum
\limits_{k=0}^{n}P(w_{k})=(n+1)a_{0.}$
\finenonce{000435}



\finexercice



\exercice{436, gourio, 2001/09/01}

\enonce{000436}{}
\begin{enumerate}
\item Lemme : Soit $P\in \Cc[X]$ non constant, $z_{0}\in \Cc, $ montrer que
$$\forall \epsilon >0,\exists z\in D(z_{0},\epsilon )=\{z\in \Cc|\left|
z-z_{0}\right| \leq \epsilon \},\left| P(z)\right| >\left|
P(z_{0})\right| .$$

\emph{Indications} : Ecrire $P(z_{0}+h)=P(z_{0})+\sum_{m=k}^{\deg P}\frac{h^{m}}{m!}%
P^{(m)}(z_{0}) $o\`{u} $k$ est le plus petit entier strictement positif tel
que $P^{(i)}(z_{0})\neq 0.$

On se propose de d\'{e}montrer le th\'{e}or\`{e}me de d'Alembert-Gauss : tout
polyn\^{o}me non constant \`{a} coefficients complexes admet une racine
complexe.

\item Expliquer pourquoi le minimum de la fonction $z\rightarrow \left|
P(z)\right| $ est atteint sur un disque centr\'{e} en $0$, mettons $D(0,\Rr),$
et expliquer pourquoi :
$$\exists z_{0}\in \Cc,\left| P(z_{0})\right| =
\inf\limits_{z\in \Cc}\left|P(z)\right| .$$

\item Montrer avec le lemme que $P(z_{0})=0.$
\end{enumerate}
\finenonce{000436}



\finexercice



\exercice{437, gourio, 2001/09/01}

\enonce{000437}{}
Soit $n\in \Nn^{*},$ et $P(X)=(X+1)^{n}-(X-1)^{n}.$
Quel est le degr\'{e} de $P$ ? Le factoriser dans $\Cc[X].$
\finenonce{000437}



\finexercice



\exercice{438, gourio, 2001/09/01}

\enonce{000438}{}
Soit $P\in {\Rr}[X]$ un polyn\^{o}me dont tous les z\'{e}ros sont
r\'{e}els et distincts, montrer que $\phi =(P^{\prime })^{2}-PP''
$ n'a pas de z\'{e}ro r\'{e}el.
\finenonce{000438}



\finexercice



\exercice{439, ortiz, 1999/04/01}

\enonce{000439}{}
Soit $K\subseteq\Cc$ un corps pour les lois
usuelles sur $\Cc$ et $P\in K[X]$ non constant.
\begin{enumerate}
    \item  Montrer que si $\alpha$ est racine de $P$ de
multiplicit\'e $m\in [1,+\infty [ $ alors $\alpha$
est racine du polyn\^ome $P'$ avec la
multiplicit\'e $m-1$.
    \item  On suppose $K=\Rr$ et $P$ scind\'e sur $\Rr$.
Montrer que $P'$ est scind\'e sur $\Rr$ (on
utilisera le th\'eor\`eme de Rolle).
\end{enumerate}
\finenonce{000439}



\finexercice



\exercice{440, ortiz, 1999/04/01}

\enonce{000440}{}
 Soient $m,n\in [1,+\infty[$, $d=\mathrm{pgcd} (m,n)$ et $P=X^m-1, Q=X^n-1, D=X^d-1 \in\Cc[X]$.
\begin{enumerate}
    \item
    \begin{enumerate}
        \item Montrer que si $x\in\Cc$ est racine commune de $P$ et $Q$ alors $x$ est racine de $D$
 (on pourra utiliser l'\'egalit\'e de B\'ezout dans $\Zz$).
        \item Montrer que si $y\in\Cc$ est racine de $D$ alors $y$ est racine commune de $P$ et $Q$
(utiliser la d\'efinition de $d$).
    \end{enumerate}
    \item
    \begin{enumerate}
        \item Soient $A,B\in \Cc[X]$ tels que toute racine de $A$ est racine de $B.$
Peut-on en d\'eduire que $A$ divise $B$? M\^eme
question si les racines de $A$ sont simples.
         \item Montrer que les racines de $D$ et $P$ sont simples et en d\'eduire que $\mathrm{pgcd} (P,Q)=D$.
    \end{enumerate}
\end{enumerate}
\finenonce{000440}



\finexercice



\exercice{441, ortiz, 1999/04/01}

\enonce{000441}{}
Soient les polyn\^omes complexes $P_1=X^3-2$,
$P_2=X^4+4$ et  $P_3=X^4+4X^3+8$.
\begin{enumerate}
    \item \'Etudier leur irr\'eductibilit\'e sur $\Cc$ et sur $\Rr$.
    \item  Montrer que $P_1$ est irr\'eductible sur $\Q$ (on utilisera que $\sqrt[3]{2}\notin \Q$).
    \item  Montrer que $P_2$ est r\'eductible sur $\Zz$.
    \item  Montrer que $P_3$ est irr\'eductible sur $\Zz$.
\end{enumerate}
\finenonce{000441}



\finexercice



\exercice{442, ortiz, 1999/04/01}

\enonce{000442}{}
Soit $P=X^4-5X^3+9X^2-15X+18\in\Cc[X]$.
D\'eterminer toutes les racines complexes de $P$
sachant que deux d'entre elles ont 6 pour produit.
\finenonce{000442}



\finexercice



\exercice{3164, quercia, 2010/03/08}

\enonce{003164}{Familles libres de polyn{\^o}mes}
Soit $a,b \in  K,\ a \ne b$. On pose $P_k = (X-a)^k(X-b)^{n-k}$.
D{\'e}montrer que la famille $(P_0,\dots,P_n)$ est libre.
\finenonce{003164}



\finexercice 


\exercice{3165, quercia, 2010/03/08}

\enonce{003165}{Formule de Van der Monde}
Soit $n\in\N^*$. Pour $k\in{[[0,n]]}$ on pose $P_k = X^k(1-X)^{n-k}$.
D{\'e}montrer que ${\cal B} = (P_0,\dots,P_n)$ est une base de~$\R_n[X]$.
Calculer les composantes dans $\cal B$ de
$\frac{d^n}{d x^n} \bigl(X^n(1-X)^n\bigr)$. En d{\'e}duire la valeur de
$\sum_{k=0}^n (C_n^k)^2$.
\finenonce{003165}



\finexercice 


\exercice{3166, quercia, 2010/03/08}

\enonce{003166}{Famille libre de polyn{\^o}mes}
Soient $U,V \in  K[X]$ non constants. On pose $P_k = U^kV^{n-k}$.
Montrer que $(P_0,\dots,P_n)$ est libre \dots
\begin{enumerate}
  \item lorsque $U \wedge V = 1$.
  \item lorsque $(U,V)$ est libre.
\end{enumerate}
\finenonce{003166}



\finexercice 


\exercice{3167, quercia, 2010/03/08}

\enonce{003167}{Ensi PC 1999}
D{\'e}terminer les poly{\^o}mes $P \in \R_{2n-1}(X)$ tels que
$P(X)+1$ est multiple de $(X-1)^n$ et
$P(X)-1$ est multiple de $(X+1)^n$.
\finenonce{003167}


\finexercice 


\exercice{3168, quercia, 2010/03/08}

\enonce{003168}{Op{\'e}rateur diff{\'e}rence}
\label{opdiff}
On note $U_p = \frac { X (X-1) \cdots (X-p+1) }{p!},\quad p \in \N$,
et $\Delta : { K[X]} \to { K[X]}, P \mapsto {P(X+1) - P(X)}$
\begin{enumerate}
  \item   D{\'e}montrer que la famille $(U_p)_{p \in \N}$ est une base de ${ K[X]}$.
  \item   Calculer $\Delta^n(U_p)$.
  \item   En d{\'e}duire que : $\forall\ P \in { K_n[X]}$, ona
             $P = P(0) + (\Delta P)(0)U_1 + (\Delta^2 P)(0)U_2 + \dots
                       + (\Delta^n P)(0)U_n$.
  \item   Soit $P \in { K[X]}$. D{\'e}montrer que :\par\indent
      $\bigl( \forall\ n \in \Z$, on a $P(n) \in \Z \bigr) \Leftrightarrow \bigl($
      les coordonn{\'e}es de $P$ dans la base $(U_p)$ sont enti{\`e}res$\bigr)$.
  \item   Soit $f : \Z  \to \Z$ une fonction quelconque.
      D{\'e}montrer que $f$ est polynomiale si et seulement si~:
      $\exists\ n \in \N$ tq $\Delta^n(f) = 0$.
\end{enumerate}
\finenonce{003168}



\finexercice 


\exercice{3169, quercia, 2010/03/08}

\enonce{003169}{Libert{\'e} de $P(X), \dots, P(X+n)$}
Soit $P \in { K[X]}$ de degr{\'e} $n$.
D{\'e}montrer que la famille $\bigl(P(X), P(X+1), \dots, P(X+n)\bigr)$ est une base de~${ K_n[X]}$.

(Utiliser l'op{\'e}rateur $\Delta$ de l'exercice \ref{opdiff})
\finenonce{003169}


\finexercice 


\exercice{3170, quercia, 2010/03/08}

\enonce{003170}{$(X+z_0)^n,\dots,(X+z_k)^n$ (Centrale MP 2003)}
Soit $k\in\N^*$ et $z_0,\dots,z_k$ des complexes.
Soient les polyn{\^o}mes $P_0 = (X+z_0)^n,\dots,P_k=(X+z_k)^n$.
Donner une condition n{\'e}cessaire et suffisante pour que~$(P_0,\dots,P_k)$ soit une
base de~$C_n[X]$.

\finenonce{003170}


\finexercice 


\exercice{3171, quercia, 2010/03/08}

\enonce{003171}{$P-X \mid P\circ P - X$}
\begin{enumerate}
  \item Soit $P \in { K[X]}$. D{\'e}montrer que $P-X$ divise $P\circ P-X$.
  \item R{\'e}soudre dans $\C$ : $(z^2+3z+1)^2 + 3z^2+8z+4 = 0$.
\end{enumerate}
\finenonce{003171}


\finexercice 


\exercice{3172, quercia, 2010/03/08}

\enonce{003172}{$P  \mapsto P(X+1) + P(X-1) - 2P(X)$}
Soit $\Phi : { K[X]} \to { K[X]},  P \mapsto {P(X+1)+P(X-1)-2P(X)}$
\begin{enumerate}
  \item  Chercher $\deg(\Phi(P))$ en fonction de $\deg P$.
  \item  En d{\'e}duire $\mathrm{Ker}\Phi$ et $\Im\Phi$.
  \item  Montrer que : $\forall\ Q \in { K[X]},\ \exists!\ P \in { K[X]}$ tq
            $\begin{cases}\Phi(P) = Q \cr P(0) = P'(0) = 0.\end{cases}$
\end{enumerate}
\finenonce{003172}



\finexercice 


\exercice{3173, quercia, 2010/03/08}

\enonce{003173}{$P  \mapsto (X-a)(P'(X)+P'(a)) + P(X)-P(a)$}
Soit $a\in K$ et $ \Phi : { K_n[X]} \to { K_n[X]}, P \mapsto {(X-a)(P'(X)+P'(a)) + P(X)-P(a).}$

Chercher $\mathrm{Ker} \Phi$ et $\Im \Phi$.

\finenonce{003173}


\finexercice 


\exercice{3174, quercia, 2010/03/08}

\enonce{003174}{$A^3 + B = C^3 + D$}
Soient $A,B,C,D \in {\R[X]}$ tels que :
$\begin{cases} \deg A = \deg C = m \cr
         \deg B < 2m,\ \deg D < 2m \cr
         A^3 + B = C^3 + D.\cr \end{cases}$

Montrer que $A = C$ et $B = D$.

Trouver un contre-exemple avec des polyn{\^o}mes {\`a} coefficients complexes.
\finenonce{003174}



\finexercice 


\exercice{3175, quercia, 2010/03/08}

\enonce{003175}{$P(n) \mid P(n+P(n))$}
Soit $P \in {\Z[X]},\ n \in \Z$, et $p = P(n)$. Montrer que $p$ divise $P(n+p)$.
\finenonce{003175}


\finexercice 


\exercice{3176, quercia, 2010/03/08}

\enonce{003176}{$P(a/b) = 0  \Rightarrow  a-kb$ divise $P(k)$}
Soit $P \in {\Z[X]}$ et $a,b \in \Z^*$ premiers entre eux tels que
$P\Bigl(\frac ab\Bigr) = 0$.
\begin{enumerate}
  \item Montrer que $a$ divise le coefficient constant de $P$.
  \item Montrer que pour tout $k\in\Z$, $a-kb$ divise $P(k)$.
\end{enumerate}
\finenonce{003176}


\finexercice 


\exercice{3177, quercia, 2010/03/08}

\enonce{003177}{Automorphismes des polyn{\^o}mes}
Pour $A \in { K[X]}$ on note ${\Phi_A} :{ K[X]} \to  { K[X]}, P \mapsto {P\circ A}$
\begin{enumerate}
  \item  D{\'e}montrer que les applications $\Phi_A$ sont les seuls endomorphismes d'alg{\`e}\-bre
     de ${ K[X]}$.
  \item  A quelle condition $\Phi_A$ est-il un isomorphisme ?
\end{enumerate}
\finenonce{003177}



\finexercice 


\exercice{3178, quercia, 2010/03/08}

\enonce{003178}{Sous anneau non principal des polyn{\^o}mes}
Soit $A = \{P \in { K[X]} \text { dont le coefficient de $X$ est nul} \}$.
D{\'e}montrer que $A$ est un sous anneau non principal de~${ K[X]}$.
\finenonce{003178}



\finexercice 


\exercice{3179, quercia, 2010/03/08}

\enonce{003179}{{\'E}quation $P^2+Q^2 = (X^2+1)^2$}
Trouver $P,Q \in {\R[X]}$ premiers entre eux tels que $P^2 + Q^2 = (X^2+1)^2$.
\finenonce{003179}


\finexercice 


\exercice{3180, quercia, 2010/03/08}

\enonce{003180}{{\'E}quation $X(X-1)P' + P^2 -(2X+1)P + 2X = 0$}
Trouver tous les polyn{\^o}mes $P \in { K[X]}$ tels que :
$X(X-1)P' + P^2 -(2X+1)P + 2X = 0$.
\finenonce{003180}


\finexercice 


\exercice{3181, quercia, 2010/03/08}

\enonce{003181}{$P(X) + P(X+1) = 2X^n$}
\begin{enumerate}
  \item  Montrer qu'il existe un uni\-que polyn{\^o}me
     $P_n \in { K[X]}$ tel que $P_n(X) + P_n(X+1) = 2X^n$.
  \item  Chercher une relation de r{\'e}currence entre $P_n'$ et $P_{n-1}$.
  \item  D{\'e}composer $P_n(X+1)$ sur la base $(P_k)_{k \in \N}$.
  \item  D{\'e}montrer que $P_n(1-X) = (-1)^nP_n(X)$.
\end{enumerate}
\finenonce{003181}


\finexercice 


\exercice{3182, quercia, 2010/03/08}

\enonce{003182}{$(1-X)^nP + X^nQ = 1$}
\begin{enumerate}
  \item  D{\'e}montrer qu'il existe $P,Q \in  K_{n-1}[X]$ uniques tels que
     $(1-X)^nP + X^nQ = 1$.
  \item  Montrer que $Q = P(1-X)$.
  \item  Montrer que : $\exists\ \lambda \in  K$ tel que $(1-X)P' - nP = \lambda X^{n-1}$.
  \item  En d{\'e}duire $P$.
\end{enumerate}
\finenonce{003182}


\finexercice 


\exercice{3183, quercia, 2010/03/08}

\enonce{003183}{Endomorphismes qui commutent avec la d{\'e}rivation}
Soit $\Phi \in \mathcal{L} { K[X]}$ commutant avec la d{\'e}rivation, c'est {\`a} dire :
$\forall\ P \in { K[X]}$, on a $\Phi(P') = \Phi(P)'$.
\begin{enumerate}
  \item  D{\'e}montrer qu'il existe un unique suite $(a_k)_{k \in \N}$ de
     scalaires tels que :
       $$\forall\ P \in { K_n[X]}, \text{ on a } \Phi(P) = \sum_{k=0}^n a_kP^{(k)}.$$
     (On {\'e}crit {\it formellement\/} :
     $\Phi = \sum_{k=0}^\infty a_k \text{D}^k$
     avec $\text{D}(P) = P'$)
  \item  D{\'e}composer ainsi l'endomorphisme $\Phi : P  \mapsto P(X+1)$.
\end{enumerate}
\finenonce{003183}



\finexercice 


\exercice{3184, quercia, 2010/03/08}

\enonce{003184}{$P$ est positif $\Rightarrow P + P' + P" + \dots$ aussi}
Soit $P \in {\R[X]}$ tel que : $\forall\ x \in \R$, on a $P(x) \ge 0$.
D{\'e}montrer que : $\forall\ x \in \R$, on a $(P + P' + P'' + \dots )(x) \ge 0$.
\finenonce{003184}




\exercice{3185, quercia, 2010/03/08}

\enonce{003185}{$P(\tan\alpha) = Q\Bigl(\frac 1{\cos\alpha}\Bigr)$}
Soit $P \in {\R[X]}$. Existe-t-il $Q \in {\R[X]}$ tel que
$\forall\ \alpha \in \left]-\frac\pi2, \frac\pi2\right[,\
 P(\tan\alpha) = Q\Bigl(\frac 1{\cos\alpha}\Bigr)$ ?
\finenonce{003185}


\finexercice 


\exercice{3186, quercia, 2010/03/08}

\enonce{003186}{$X^n + 1/X^n = P_n(X+1/X)$}
\begin{enumerate}
  \item  Montrer que pour tout entier $n \in \N$ il existe un unique polyn{\^o}me
     $P_n \in {\Z[X]}$ v{\'e}rifiant :
     $$\forall\ z \in \C^*,\ z^n + z^{-n} = P_n(z + z^{-1}).$$

  \item  D{\'e}terminer le degr{\'e}, le coefficient dominant, et les racines de $P_n$.

  \item  Pour $P \in {\C[X]}$, on note $\tilde P$ le polyn{\^o}me tel que :
     $$\forall\ z \in \C^*,\ P(z) + P(z^{-1}) = \tilde P(z + z^{-1}).$$

     {\'E}tudier l'application $P  \mapsto \tilde P$.
\end{enumerate}
\finenonce{003186}


\finexercice 


\exercice{3187, quercia, 2010/03/08}

\enonce{003187}{Polytechnique MP$^*$ 2000}
\begin{enumerate}
  \item Donner un isomorphisme $f$ en\-tre $\C^{n+1}$ et $\C_n[X]$.
  \item Montrer que $\sigma:{\C^{n+1}}\to{\C^{n+1}},
    {(a_0,\dots,a_n)}\mapsto{(a_n,a_0,\dots,a_{n-1})}$ est lin{\'e}aire.
  \item Si $(P,Q)\in(\C[X])^2$, on d{\'e}finit le produit $\overline{PQ}$
    comme le reste de la division euclidienne de $PQ$ par $X^{n+1}-1$.
    Montrer que l'application induite par~$\sigma$ sur~$\C_n[X]$
    (c'est-{\`a}-dire $f\circ \sigma\circ f^{-1}$) est l'application
    qui {\`a}~$P$ associe $\overline{XP}$.

  \item Soit~$F$ un sous-espace de~$\C^{n+1}$ stable par~$\sigma$.

    Montrer qu'il existe un polyn{\^o}me $Q$ tel que
    $f(F) = \{\overline{RQ},\ R\in\C_n[X]\}$.
\end{enumerate}
\finenonce{003187}


\finexercice 


\exercice{3188, quercia, 2010/03/08}

\enonce{003188}{Centrale MP 2002}
D{\'e}terminer tous les polyn{\^o}mes $P$ tels que $P(\C)\subset\R$
puis tels que $P(\Q )\subset\Q$ et enfin tels que $P(\Q )=\Q$.
\finenonce{003188}


\finexercice 


\exercice{3189, quercia, 2010/03/08}

\enonce{003189}{Polytechnique MP 2002}
Soient $x_1,\dots,x_n\in\C$ distincts et $y_1,\dots,y_n\in\C$.
Trouver $E = \{P\in\C[X]\text{ tq }\forall\ i,\ P^{-1}(\{y_i\}) = \{x_i\}\}$.
\finenonce{003189}


\finexercice


\exercice{3190, quercia, 2010/03/08}

\enonce{003190}{ENS Ulm MP 2002}
Soit $S\subset \N$ fini et $P=\sum_{s\in S} a_sX^s\in\C[X]$.
\begin{enumerate}
  \item  On suppose que les $a_s$ sont r{\'e}els. Montrer que $P$ a moins de racines strictement positives distinctes que la suite 
$(a_s)$ n'a de changement de signe.

  \item  On suppose que $P$ v{\'e}rifie~: $\forall\ s\in S,\ P(s)=0$. Montrer que $P$ est nul.
\end{enumerate}
\finenonce{003190}


\finexercice 


\exercice{3191, quercia, 2010/03/08}

\enonce{003191}{$\sum_{k=1}^{100}\frac k{x-k}\ge 1$ (Ens Ulm-Lyon-Cachan MP$^*$ 2003)}
Montrer que l'ensemble des solutions de l'in{\'e}quation $\sum_{k=1}^{100}\frac k{x-k}\ge 1$
est une r{\'e}union finie d'intervalles disjoints. Calculer la somme des
longueurs de ces intervalles.
\finenonce{003191}


\finexercice 


\exercice{3192, quercia, 2010/03/08}

\enonce{003192}{Polyn{\^o}me positif (Ens Ulm MP$^*$ 2003)}
Soit~$P\in\R[X]$. Montrer~:

($\forall\ x\ge 0,\ P(x) > 0$) $\Leftrightarrow$ ($\exists\ \ell \in\N$  tq $(X+1)^\ell P(X)$ est {\`a} coefficients strictement positifs).
\finenonce{003192}


\finexercice 


\exercice{3193, quercia, 2010/03/08}

\enonce{003193}{Diviseurs premiers de la suite $(P(n))$ (Ens ULM-Lyon-Cachan MP$^*$ 2003)}
Soit $P\in\Z[X]$ non constant et $E$ l'ensemble des diviseurs premiers d'au moins un $P(n)$, $n\in\Z$.
Montrer que~$E$ est infini.
\finenonce{003193}


\finexercice 


\exercice{3194, quercia, 2010/03/08}

\enonce{003194}{Centrale MP 2004}
Soit~$n\in\N^*$. Montrer l'existence de~$P_n\in\R[X]$ tel que $1+X-P_n^2$ est divisible
par~$X^n$.
\finenonce{003194}


\finexercice 


\exercice{3195, quercia, 2010/03/08}

\enonce{003195}{Polyn{\^o}mes {\`a} coefficients entiers, ULM-Lyon-Cachan MP$^*$ 2004}
On donne un entier $n\ge 0$.

Montrer qu'il existe des polyn{\^o}mes
$P_0,\dots,P_n$ dans $\Z_n[X]$ tels que
$\forall\ i,j\in{[[0,n]]},\  \int_{t=0}^1 t^iP_j(t)\,dt = \delta_{ij}$.
\finenonce{003195}


\finexercice 


\exercice{3254, quercia, 2010/03/08}

\enonce{003254}{$a/b + b/c + c/a$}
Soient $a,b,c$ les racines de $X^3 + pX + q$, $q \ne 0$. Calculer :
$\sum_{\sigma \in S_3} \left( \frac {\sigma(a)}{\sigma(b)}
                      + \frac {\sigma(b)}{\sigma(c)}
                      + \frac {\sigma(c)}{\sigma(a)} \right)$.
\finenonce{003254}


\finexercice 


\exercice{3255, quercia, 2010/03/08}

\enonce{003255}{$1/(x_i-1)$}
Soient $x_1,x_2,x_3,x_4$ les racines de $X^4 + X + 1$.
Calculer $\sum_{i=1}^4 \frac 1{x_i-1}$.
\finenonce{003255}


\finexercice 


\exercice{3256, quercia, 2010/03/08}

\enonce{003256}{$x_i/(x_jx_k)$}
Soient $x_1,\dots,x_8$ les racines de $X^8 + X^7 - X^2 + 3$.
Calculer $\sum_{1\le i\le8 , 1\le j<k \le8} \frac {x_i}{x_jx_k}$.
\finenonce{003256}



\finexercice 


\exercice{3257, quercia, 2010/03/08}

\enonce{003257}{$x_i^7$}
Soient $a,b,c$ les racines de $X^3 - X + 1$. Calculer $a^7 + b^7 + c^7$.
\finenonce{003257}



\finexercice 


\exercice{3258, quercia, 2010/03/08}

\enonce{003258}{$a+b+c$, $a^2+b^2+c^2$, $1/a+1/b+1/c$ donn{\'e}s}
R{\'e}soudre $$\left\{\begin{aligned} a         &+ &b         &+ &c         &= &0  \cr
                  a^2       &+ &b^2       &+ &c^2       &= &1  \cr
                  \frac 1a &+ &\frac 1b &+ &\frac 1c &= &-1.\cr\end{aligned}\right.$$
\finenonce{003258}


\finexercice 


\exercice{3259, quercia, 2010/03/08}

\enonce{003259}{Ensi P 90}
    R{\'e}soudre dans $\C$ le syst{\`e}me :
    $\begin{cases} x+y+z=2\cr x^2+y^2+z^2=6\cr
	     \frac1x+\frac1y+\frac1z=\frac12.\cr \end{cases}$
\finenonce{003259}


\finexercice 


\exercice{3260, quercia, 2010/03/08}

\enonce{003260}{$ \int_{t=-1}^1 P(t)\,dt = d\bigl(P(a)+P(b)+P(c)\bigr)$}
Trouver $a,b,c,d \in \R$ tels que : $\forall\ P \in \R_3[X],\
 \int_{t=-1}^1 P(t)\,dt = d\bigl(P(a)+P(b)+P(c)\bigr)$.
\finenonce{003260}


\finexercice 


\exercice{3261, quercia, 2010/03/08}

\enonce{003261}{$a,b,c$ en progression g{\'e}om{\'e}trique}
Soient $a,b,c \in \C$.

Montrer que ces nombres sont en progression g{\'e}om{\'e}trique
si et seulement si $(ab+ac+bc)^3 = abc(a+b+c)^3$.
\finenonce{003261}



\finexercice 


\exercice{3262, quercia, 2010/03/08}

\enonce{003262}{Condition liant les racines}
Soit $P = X^3 + pX + q$ de racines $a,b,c$.
\begin{enumerate}
  \item   CNS pour ces racines soient aux sommets d'un carr{\'e} ?
  \item   CNS pour que $a^2 + b^2 = 1 + c^2$ ?
\end{enumerate}
\finenonce{003262}


\finexercice 


\exercice{3263, quercia, 2010/03/08}

\enonce{003263}{Condition liant les racines}
Soient $A,B,C$ les points dont les affixes sont les racines de
$X^3 + pX + q$, $p,q \in \C$.
A quelle condition sur $p$ et $q$ a t-on $AB = AC = 2BC$ ?
\finenonce{003263}


\finexercice 


\exercice{3264, quercia, 2010/03/08}

\enonce{003264}{Condition liant les racines}
Soit $P = X^4 + aX^2 + bX + c$ de racines $\alpha,\beta,\gamma,\delta$.
CNS pour ces racines soient en progression arithm{\'e}tique~?
\finenonce{003264}


\finexercice 


\exercice{3265, quercia, 2010/03/08}

\enonce{003265}{Transformation d'{\'e}quation}
Soient $x_1,x_2,x_3$ les racines de $X^3 + 2X^2 + 3X + 4$.

Calculer le polyn{\^o}me unitaire de $\R_3[X]$ dont
$x_1+x_2,x_2+x_3,x_3+x_1$ sont {\it les\/} racines.
\finenonce{003265}


\finexercice 


\exercice{3266, quercia, 2010/03/08}

\enonce{003266}{Transformation d'{\'e}quation}
Soient $x_1,x_2,x_3$ les racines de $X^3 + aX^2 + bX + c$.

Calculer le polyn{\^o}me unitaire de $\R_3[X]$ dont
$x_1^2,x_2^2,x_3^2$ sont {\it les\/} racines.
\finenonce{003266}


\finexercice 


\exercice{3267, quercia, 2010/03/08}

\enonce{003267}{$2X^3 + 5X^2 - X + \lambda$ a une racine de module 1}
Trouver $\lambda \in \R$ tel que $2X^3 + 5X^2 - X + \lambda$
ait une racine de module 1.
\finenonce{003267}


\finexercice 


\exercice{3268, quercia, 2010/03/08}

\enonce{003268}{Polyn{\^o}mes dont les racines sont de module 1}
Soit $n \in \N^*$ et ${\cal E}$ l'ensemble des polyn{\^o}mes {\`a} coefficients
entiers, unitaires de degr{\'e} $n$ et dont toutes les racines sont de module 1.
\begin{enumerate}
  \item   D{\'e}montrer que $\cal E$ est fini.
  \item   Pour $P \in \cal E$ de racines $x_1,\dots,x_n$, on note $\widetilde P$ le
      polyn{\^o}me unitaire de racines $x_1^2,\dots,x_n^2$.

      D{\'e}montrer que $\widetilde P \in \cal E$.
  \item   En d{\'e}duire que : $\forall\ P \in \cal E$, les racines de $P$ sont des racines de
      l'unit{\'e}.
\end{enumerate}
\finenonce{003268}


\finexercice 


\exercice{3269, quercia, 2010/03/08}

\enonce{003269}{Centrale MP 2001}
Soit $f(x) = x^4 + ax^3 + bx^2 + cx + d$ avec $a,b,c,d$ r{\'e}els.
Donner une condition n{\'e}cessaire et suffisante portant sur~$a,b,c,d$
pour qu'il existe une droite coupant la courbe repr{\'e}sentative
de~$f$ en quatre points distincts $M_1,M_2,M_3,M_4$ tels que
$M_1M_2 = M_2M_3 = M_3M_4$.
\finenonce{003269}


\finexercice 


\exercice{5314, rouget, 2010/07/04}
\enonce{005314}{***}
On pose $\omega_k=e^{2ik\pi/n}$ et $Q=1+2X+...+nX^{n-1}$. Calculer $\prod_{k=0}^{n-1}Q(\omega_k)$.
\finenonce{005314}


\finexercice

\exercice{5320, rouget, 2010/07/04}
\enonce{005320}{**}
Soit $P$ un polynôme différent de $X$. Montrer que $P(X)-X$ divise $P(P(X))-X$.
\finenonce{005320}


\finexercice


\exercice{5321, rouget, 2010/07/04}
\enonce{005321}{***}
Soit $P$ un polynôme à coefficients entiers relatifs de degré supérieur ou égal à $1$. Soit $n$ un entier relatif 
et $m=P(n)$.
\begin{enumerate}
\item  Montrer que $\forall k\in\Zz,\;P(n+km)$ est un entier divisible par $m$.
\item  Montrer qu'il n'existe pas de polynômes non constants à coefficients entiers tels que $P(n)$ soit premier pour tout entier $n$.
\end{enumerate}
\finenonce{005321}


\finexercice


\exercice{5322, rouget, 2010/07/04}
\enonce{005322}{*** Polynômes $P$ vérifiant $P(\Zz)\subset\Zz$}
Soit $E$ la partie de $\Cc[X]$ formée des polynômes $P$ vérifiant $\forall a\in\Zz,\;P(a)\in\Zz$.
\begin{enumerate}
\item  On pose $P_0=1$ et pour $n$ entier naturel non nul, $P_n=\frac{1}{n!}\prod_{k=1}^{n}(X+k)$ (on peut définir la notation $P_n=C_{X+n}{n}$). Montrer que $\forall n\in\Nn,\;P_n\in E$.
\item  Montrer que toute combinaison linéaire à coefficients entiers relatifs des $P_n$ est encore un élément de $E$.
\item  Montrer que $E$ est l'ensemble des combinaisons linéaires à coefficients entiers relatifs des $P_n$.
\end{enumerate}
\finenonce{005322}


\finexercice


\exercice{5326, rouget, 2010/07/04}
\enonce{005326}{***T}
Trouver un polynôme de degré $5$ tel que $P(X)+10$ soit divisible par $(X+2)^3$ et $P(X)-10$ soit divisible par 
$(X-2)^3$.
\finenonce{005326}


\finexercice


\exercice{5327, rouget, 2010/07/04}
\enonce{005327}{***I}
Trouver les polynômes $P$ de $\Rr[X]$ vérifiant $P(X^2)=P(X)P(X+1)$ (penser aux racines de $P$).
\finenonce{005327}


\finexercice


\exercice{5330, rouget, 2010/07/04}
\enonce{005330}{**}
Résoudre dans $\Cc^3$ (resp. $\Cc^4$) le système~:

$$1)\;\left\{
\begin{array}{l}
x+y+z=1\\
\frac{1}{x}+\frac{1}{y}+\frac{1}{z}=1\\
xyz=-4
\end{array}
\right.
\quad2)\;\left\{
\begin{array}{l}
x+y+z+t=0\\
x^2+y^2+z^2+t^2=10\\
x^3+y^3+z^3+t^3=0\\
x^4+y^4+z^4+t^4=26
\end{array}
\right.
.$$ 
\finenonce{005330}


\finexercice


\exercice{5331, rouget, 2010/07/04}
\enonce{005331}{**T}
Trouver tous les polynômes $P$ vérifiant $P(2X)=P'(X)P''(X)$.
\finenonce{005331}


\finexercice


\exercice{5332, rouget, 2010/07/04}
\enonce{005332}{**}
Factoriser dans $\Cc[X]$ le polynôme $12X^4+X^3+15X^2-20X+4$.
\finenonce{005332}


\finexercice


\exercice{5333, rouget, 2010/07/04}
\enonce{005333}{***}
Soit $n\in\Nn^*$. Montrer que $(X-1)^{2n}-X^{2n}+2X-1$ est divisible par $2X^3-3X^2+X$ puis déterminer le quotient.
\finenonce{005333}


\finexercice


\exercice{5334, rouget, 2010/07/04}
\enonce{005334}{**I}
Déterminer deux polynômes $U$ et $V$ vérifiant $UX^n+V(1-X)^m=1$ et $\mbox{deg}(U)<m$ et $\mbox{deg}(V)<n$.
\finenonce{005334}


\finexercice


\exercice{5343, rouget, 2010/07/04}
\enonce{005343}{}
\label{exo:suprou9}
Soit $P=$ où $n$ est un entier naturel non nul, les $a_i$ sont des entiers relatifs et $a_0$ et $a_n$ sont non nuls. Soient $p$ un entier relatif non nul et $q$ un entier naturel non nul tels que $p\wedge q=1$.

Montrer que, si $r=\frac{p}{q}$ est une racine (rationnelle) de $P$ alors $p$ divise $a_0$ et $q$ divise $a_n$.

Application. Résoudre dans $\Cc$ l'équation $9z^4-3z^3+16z^2-6z-4=0$.
\finenonce{005343}


\finexercice


\exercice{5344, rouget, 2010/07/04}
\enonce{005344}{Equations réciproques}
Résoudre dans $\Cc$ les équations suivantes~:
\begin{enumerate}
\item  $z^4+2z^3+3z^2+2z+1=0$ en posant $Z=z+\frac{1}{z}$ (ou autrement).
\item  $z^6-5z^5+5z^4-5z^2+5z-1=0$.
\item  $z^7-z^6-7z^5+7z^4+7z^3-7z^2-z+1=0$.
\end{enumerate}
\finenonce{005344}


\finexercice

\exercice{5346, rouget, 2010/07/04}
\enonce{005346}{}
Soit $P$ un polynôme à coefficients complexes de degré $4$.

Montrer que les images dans le plan complexe des racines de $P$ forment un parallélogramme si et seulement si $P'$ et $P^{(3)}$ ont une racine commune
\finenonce{005346}


\finexercice


\exercice{5347, rouget, 2010/07/04}
\enonce{005347}{}
Résoudre dans $\Cc^3$ le système $\left\{
\begin{array}{l}
y^2+yz+z^2=7\\
z^2+zx+x^2=13\\
x^2+xy+y^2=3
\end{array}
\right.$.
\finenonce{005347}


\finexercice


\exercice{5348, rouget, 2010/07/04}
\enonce{005348}{}
Soit $n$ un entier naturel supérieur ou égal à $2$. Pour $k\in\Zz$, on pose $\omega_k=e^{2ik\pi/n}$.
\begin{enumerate}
\item  Calculer $\prod_{k=0}^{n-1}\left(1+\frac{2}{2-\omega_k}\right)$.
\item  Montrer que, pour tout réel $a$, $\prod_{k=0}^{n-1}(\omega_k^2-2\omega_k\cos a+1)=2(1-\cos(na))$ (questions indépendantes.)
\end{enumerate}
\finenonce{005348}


\finexercice


\exercice{5352, rouget, 2010/07/04}
\enonce{005352}{}
Résoudre dans $\Cc$ l'équation $z^4-21z+8=0$ sachant qu'il existe deux des solutions sont inverses l'une de l'autre.
\finenonce{005352}


\finexercice

\section{106.01 Définition, sous-espace}


\exercice{886, legall, 1998/09/01}
\video{1rxZQKF8yc0}
\enonce{000886}{}
D\' eterminer lesquels des
ensembles $E_1$, $E_2$, $E_3$ et $E_4$ sont des sous-espaces
vectoriels de ${\Rr}^3$. 

 $E_1 =\{ (x,y,z)\in {\Rr}^3\ \mid \ 3x-7y = z \} $ 

 $E_2 =\{(x,y,z)\in {\Rr}^3\ \mid \ x^2-z^2=0 \} $  

 $E_3=\{ (x,y,z)\in {\Rr}^3\ \mid \ x+y-z=x+y+z=0 \} $ 

 $E_4 =\{ (x,y,z)\in {\Rr}^3\ \mid \ z(x^2+y^2)=0 \} $
\finenonce{000886} 




\finexercice


\exercice{887, legall, 1998/09/01}

\enonce{000887}{}
Soit $\R_+^*$  muni de la loi interne $\oplus$ d\'efinie par $a \oplus b = ab, \forall a,
b \in \R_+^*$ et de la loi externe $\otimes$ telle que $\lambda\otimes a = a^\lambda, \forall
a \in \R_+^*, \forall \lambda \in \R$.
    Montrer que $E=(\R_+^*, \oplus , \otimes)$ est un $\R$-espace vectoriel.
\finenonce{000887}



\finexercice



\exercice{888, legall, 1998/09/01}
\video{-rEusJplGVE}
\enonce{000888}{}
Parmi les ensembles suivants reconna\^\i tre ceux qui sont des
sous-espaces vectoriels.

$ E_1 =\left\{ (x,y,z)\in \R^3 \mid x+y+a=0 \hbox{ et }  x +3az =0\right\}$

$ E_2 =\left\{f \in {\mathcal F}(\R,\R) \mid f(1)=0\right\}$

$ E_3 =\left\{f \in {\mathcal F}(\R,\R) \mid  f(0)=1\right\}$

$E_4 =\left\{(x,y)\in \R^2 \mid x + \alpha y +1 \geqslant 0\right\}$
\finenonce{000888} 


\finexercice


\exercice{889, legall, 1998/09/01}

\enonce{000889}{}
Parmi les ensembles suivants, reconna\^{\i}tre ceux qui sont des sous-espaces
vectoriels :\\
$
\begin{array}{ll}
E_{1}=\{(x,y,z)\in { \Rr}^{3}/ x+y=0\}; & E_{1}'=\{(x,y,z)\in {\Rr}^{3}/xy=0\}.\\
E_{2}=\{(x,y,z,t)\in { \Rr}^{4}/ x=0,   y=z\}; & E_{2}'=\{(x,y,z)\in { \Rr}^{3}/x=1\}.\\
E_{3}=\{(x,y)\in { \Rr}^{2}/ x^{2}+xy\ge 0\}; & E_{3}'=\{(x,y)\in {\Rr}^{2}/x^{2}+xy+y^{2}\ge 0\}.\\
E_{4}=\{f\in { \Rr}^{ \Rr}/ f(1)=0\}; & E_{4}'=\{f\in { \Rr}^{ \Rr}/ f(0)=1\};\\
& E_{4}"=\{f\in { \Rr}^{ \Rr}/ f \mbox{ est croissante}\}.
\end{array}
$
\finenonce{000889}



\finexercice



\exercice{890, legall, 1998/09/01}

\enonce{000890}{}
D\'eterminer si ${ \Rr}^{2}$, muni des lois internes et externes suivantes, est
ou n'est pas un ${ \Rr}$-espace vectoriel :
\begin{enumerate}
    \item  $(a,b)+(c,d)=(a+c,b+d);  \lambda(a,b) =(a,\lambda b),   \lambda \in { \Rr}.  $
    \item $ (a,b)+(c,d)=(a+c,b+d);  \lambda(a,b) = (\lambda^{2}a,\lambda^{2} b),    \lambda \in { \Rr}.  $
    \item  $(a,b)+(c,d)=(c,d);  \lambda(a,b) = (\lambda a,\lambda b),    \lambda \in { \Rr}. $
\end{enumerate}

\finenonce{000890}



\finexercice



\exercice{891, cousquer, 2003/10/01}

\enonce{000891}{}
 Dire si les objets suivants sont des espaces vectoriels~:
\begin{enumerate}


\item L'ensemble des fonctions réelles sur
$\lbrack0,1 \rbrack$, continues, positives ou nulles, pour
l'addition et le produit par un réel.

\item L'ensemble des fonctions réelles sur ${ \bf R}$~ vérifiant
$\lim_{x \to+\infty} f(x)=0$ pour les mêmes opérations.

\item L'ensemble des solutions $(x_1,x_2,x_3)$ du système~:
$\left\{\begin{array}{rcl}
    2x_1-x_2+x_3 &=&0 \\
    x_1-4x_2+7x_3 &=&0 \\ 
    x_1+3x_2-6x_3 &=&0.
\end{array}\right.$

\item L'ensemble des fonctions continues sur $\lbrack0,1 \rbrack$
vérifiant $f(1/2)=0$.

\item L'ensemble $\mathbb{R}_+^*$ pour les opérations $x \oplus y=xy$ et 
$\lambda\cdot x=x^{\lambda}$, $(\lambda\in\mathbb{R})$.

\item L'ensemble des fonctions impaires sur $\mathbb{R}$.

\item L'ensemble des fonctions sur $\lbrack a,b \rbrack$ continues, vérifiant
$f(a)=7f(b)+\int_a^b t^3f(t)\,dt$.

\item L'ensemble des fonctions sur $\mathbb{R}$ qui sont nulle en~$1$
ou nulle en~$4$.

\item L'ensemble des fonctions sur $\mathbb{R}$ qui peuvent s'écrire comme somme
d'une fonction nulle en~$1$ et d'une fonction nulle en~$4$. Identifier cet
ensemble.

\item L'ensemble des polynômes de degré exactement~$n$.


\item L'ensemble des fonctions de classe $C^2$ vérifiant $f''+\omega^2f=0$.

\item L'ensemble des fonctions sur $\mathbb{R}$ telles que $f(3)=7$.

\item L'ensemble des primitives de la fonction $xe^x$ sur $\mathbb{R}$.

\item L'ensemble des nombres complexes d'argument $\pi/4+k\pi$,
$(k\in\mathbb{Z})$.

\item L'ensemble des points $(x,y)$ de $\mathbb{R}^2$,  vérifiant
$\sin(x+y)=0$.

\item L'ensemble des vecteurs $(x,y,z)$ de $\mathbb{R}^3$ orthogonaux
au vecteur $(-1,3,-2)$.

\item L'ensemble des fonctions continues sur $\lbrack0,1 \rbrack$
vérifiant $\int_0^1\sin x f(x)\,dx=0$.

\item L'ensemble des polynômes ne comportant pas de terme de degré~$7$.

\item L'ensemble des fonctions paires sur $\mathbb{R}$.

\end{enumerate}
\finenonce{000891}



\finexercice



\exercice{892, legall, 1998/09/01}

\enonce{000892}{}
Montrer que l'ensemble ${\mathcal E}=\{ f\in { \Rr}^{ \Rr}/(\exists (a,\varphi) \in
{ \Rr}^{2}) (\forall x \in { \Rr})    f(x)=a\cos (x-\varphi)\}$ est un ${ \Rr}$-espace
vectoriel.
\finenonce{000892}



\finexercice



\exercice{893, legall, 1998/09/01}
\video{GNQ77W5XwGs}
\enonce{000893}{}
Soit $E$ un espace vectoriel.
\begin{enumerate}
\item Soient $F$ et $G$ deux sous-espaces de $E$. Montrer que
  $$F\cup G \hbox{ est un sous-espace vectoriel de } E
 \quad \Longleftrightarrow \quad F\subset G \hbox{ ou } G \subset F.$$
\item Soit $H$ un troisi\`eme sous-espace vectoriel de $E$. Prouver
  que
  $$G \subset F \Longrightarrow F\cap(G+H) = G + (F\cap H) .$$
 \end{enumerate}
\finenonce{000893} 


\finexercice


\exercice{894, ridde, 1999/11/01}

\enonce{000894}{}
On munit $\Rr^2$ de l'addition usuelle et de la loi externe $\lambda (x, y)
 = (\lambda x, y)$. Est-ce un $\Rr$-espace vectoriel ?
\finenonce{000894}



\finexercice



\exercice{895, ridde, 1999/11/01}

\enonce{000895}{}
Montrer que $\left\{ (x, y, z) \in \Rr^3 /x + y + z = 0 \text{ et }
2x-y + 3z = 0\right\}$ est un sous-espace vectoriel de $\Rr^3$.
\finenonce{000895}



\finexercice



\exercice{896, gourio, 2001/09/01}

\enonce{000896}{}
Montrer que $$
F=\{f\in C(\Rr,\Rr)|\exists (A,\phi )\in \Rr^{2},\forall x\in \Rr,f(x)=A\cos (x+\phi
)\}$$
est un espace vectoriel.
\finenonce{000896}



\finexercice



\exercice{2425, matexo1, 2002/02/01}

\enonce{002425}{}
\begin{center}
{\sc VRAI ou FAUX}
\end{center}

\begin{enumerate}
\item L'ensemble $\{0\}$ est un espace vectoriel r\'eel.
\item L'ensemble $\{0,1\}$ est un espace vectoriel r\'eel.
\item Tout sous-espace vectoriel autre que $\{0\}$ poss\`ede un
sous-espace strict.
\item L'intersection de deux sous-espaces vectoriels (d'un m\^eme
espace plus grand) est un espace vectoriel.
\item La r\'eunion de deux sous-espaces vectoriels est un
espace vectoriel.
\item La somme de deux sous-espaces vectoriels est un espace
vectoriel.
\item Le produit cart\'esien $E\times F$ de deux espaces vectoriels
est un espace vectoriel.
\end{enumerate}
\finenonce{002425}
\finexercice


\exercice{2427, matexo1, 2002/02/01}

\enonce{002427}{}
On note $\R^n$ l'ensemble des $n$-uplets $(x_1, \ldots, x_n)$
de nombres r\'eels\,; $\R [X]$ l'ensemble des polyn\^omes \`a 
coefficients r\'eels en la variable X\,; $\R[X]_p$ le sous-ensemble
des polyn\^omes de degr\'e $\le p$\,; $\R(X)$ l'ensemble des
fractions rationnelles \`a coefficients r\'eels en la variable X\,;
$\R(X)_p$ le sous-ensemble des fractions rationnelles de 
degr\'e $\le p$\,; $C^k(\R)$ l'ensemble des fonctions r\'eelles
d\'efinies sur $\R$ et $k$ fois contin\^ument d\'erivables ($k
\ge 0$ entier)\,; $C^\infty(\R)$ l'ensemble des fonctions
ind\'efiniment d\'erivables sur $\R$.  
\begin{enumerate}
\item Dot\'es des op\'erations
d'addition et de multiplication usuelles, lesquels de ces
ensembles sont des espaces vectoriels\,?
\item Montrer que $\R[X]_p \subset
\R[X] \subset \R(X)$ et que $C^\infty(\R) \subset C^k(\R) \subset
C^0(\R)$,
et que ce sont des sous-espaces vectoriels. 
\item Si l'on identifie les
polyn\^omes et les fractions rationnelles aux fonctions
correspondantes, a-t-on
$\R[X] \subset C^\infty(\R)$ et $\R(X) \subset C^\infty(\R)$\,?
\end{enumerate}
\finenonce{002427}


\finexercice


\exercice{2428, matexo1, 2002/02/01}

\enonce{002428}{}
Soit $E$ un espace vectoriel de dimension finie, $F, G$ deux
sous-espaces
de $E$. Montrer que $\dim(F + G) = \dim F + \dim G - \dim(F \cap
G)$.
\finenonce{002428}


\finexercice


\exercice{2429, matexo1, 2002/02/01}

\enonce{002429}{}
Soit $E = \R[X]_n$ (polyn\^omes de degr\'e $\le n$), et $P
\in E$.

\begin{enumerate}
\item Montrer que l'ensemble $F_P$ des polyn\^omes de
$E$ multiples de $P$ est un sous-espace vectoriel de $E$. Quelle en
est la dimension en fonction du degr\'e de $P$\,?
\item Soit $Q \in E$ un polyn\^ome sans racine commune avec
$P$, et tel que $\deg P + \deg Q = n+1$. Montrer que $E = F_P
\oplus F_Q$.
\item En d\'eduire qu'il existe deux polyn\^omes $U$ et $V$
tels que $U P + V Q = 1$.
\end{enumerate}
\finenonce{002429}


\finexercice


\exercice{2447, matexo1, 2002/02/01}

\enonce{002447}{}
Soit $E$ l'espace vectoriel des fonctions r\'eelles
ind\'efiniment d\'erivables \`a valeurs dans $\R$.

\begin{enumerate}
\item Montrer que les quatre fonctions d\'efinies par
\begin{eqnarray*}
&&x_1(t) = \cos t \cosh t, \\
&&x_2(t) = \sin t \cosh t, \\
&&x_3(t) = \cos t \sinh t, \\
&&x_4(t) =\sin t \sinh t
\end{eqnarray*}
appartiennent \`a $E$ et sont lin\'eairement ind\'ependantes.
\item Soit $F$ le sous-espace vectoriel de $E$ engendr\'e par
ces quatres vecteurs, et $u$ l'endomor\-phisme de $E$ d\'efini par
$u(f) = f'$. Montrer que $F$ est stable par $u$ et d\'eterminer
la matrice $M$ de $u$ dans la bae $(x_1, x_2, x_3, x_4)$ de $F$.
\item Calculer $M^n$.
\end{enumerate}
\finenonce{002447}


\finexercice


\exercice{2778, tumpach, 2009/10/25}
\enonce{002778}{}
\begin{enumerate}
\item
En utilisant les op\'erations d'addition $+$ et de multiplication $\cdot$ de deux nombres, d\'efinir, pour chaque ensemble $E$ de la liste ci-dessous~:
\begin{itemize}
\item une addition $\oplus~: E \times E \rightarrow E$ ;
\item une multiplication par un nombre r\'eel $\odot~: \mathbb{R}\times E \rightarrow E$. 
\end{itemize}
\begin{enumerate}
\item $E = \mathbb{R}^n$ ;
\item $E = $ l'ensemble des trajectoires d'une particule  ponctuelle dans l'espace $\mathbb{R}^3$ ;
\item $E = \textrm{l'ensemble des solutions}\,\,\,(x, y, z)\in \mathbb{R}^3\,\,\, \textrm{de l'\'equation} \quad  
\mathcal{S}_1~:
x -  2y  +  3z =  0 ;
$
\item $E = \textrm{l'ensemble des solutions}\,\,\,(x, y, z)\in \mathbb{R}^3\,\,\,\textrm{du syst\`eme d'\'equations}$.
$$\mathcal{S} _2~: \left\{
\begin{array}{l} 
2x  +  4y  -  6z   =  0\\
y + z =  0\end{array}   
 \right.
 ;$$
  \item $E = $ l'ensemble des solutions de l'\'equation diff\'erentielle $y'' + 2 y' + 3 y = 0$ ;
 \item $E = $ l'ensemble des fonctions $y(x)$ telles que $$ y''(x)\sin x + x^3 y'(x) +  y(x)\log x = 0,\,\,\, \forall x >0 ;$$ 
 \item $E = $ l'ensemble des fonctions $\Psi(t, x)$, \`a valeurs complexes, solutions de l'\'equation de Schr\"odinger~: $$i \hbar\frac{\partial}{\partial t}\Psi(t, x) = -\frac{\hbar}{2m}\frac{\partial^2}{\partial x^2}\Psi(x, t) + x^2\Psi(t, x)$$ o\`u $\hbar$ et $m$ sont des constantes ;
  \item $E = \textrm{ l'ensemble des suites}\,\,\,  (x_n)_{n\in\mathbb{N}}$ de nombres r\'eels ;
\item $E = $ l'ensemble des polyn\^omes $P(x)$ \`a coefficients r\'eels ;
\item $E = $ l'ensemble des polyn\^omes $P(x)$ \`a coefficients r\'eels de degr\'e inf\'erieur ou \'egal \`a $3$ ;
\item $E = $ l'ensemble des polyn\^omes $P(x)$ \`a coefficients r\'eels divisibles par $(x - 1)$ ;
\item $E = $ l'ensemble des fonctions continues sur l'intervalle $[0, 1]$ \`a valeurs r\'eelles ;
\item $E = $ l'ensemble des fonctions continues sur l'intervalle $[0, 1]$ \`a valeurs r\'eelles et d'int\'egrale nulle ;
\item $E = $ l'ensemble des fonctions d\'erivables sur l'intervalle $]0, 1[$ \`a valeurs r\'eelles ;
\item $E = $ l'ensemble des fonctions r\'eelles qui s'annulent en $0 \in \mathbb{R}$.
\item $E = $ l'ensemble des fonctions r\'eelles qui tendent vers $0$ lorsque $x$ tend vers $+\infty$ ; 

\end{enumerate}


\item Pour les op\'erations d'addition $\oplus$ construites, montrer que  $E$ poss\`ede un \'el\'ement neutre (terme \`a d\'efinir), et que chaque \'el\'ement de $E$ poss\`ede un inverse.
\end{enumerate}
\finenonce{002778}



\finexercice


\exercice{2779, tumpach, 2009/10/25}
\enonce{002779}{} 
Qu'est -ce qui emp\^eche de d\'efinir les m\^emes op\'erations que dans l'exercice pr\'ec\'edent sur les ensembles suivants~?
\begin{enumerate}
\item[(a)] $E = \textrm{l'ensemble des solutions}\,\,\, (x, y, z)\in \mathbb{R}^3\,\,\, \textrm{de l'\'equation} \quad  
\mathcal{S}_3~: 
x -  2y  +  3z  =  3
$ ;
 \item[(b)] $E = $ l'ensemble des fonctions $y(x)$ telles que $y''(x)\sin x  + x^3 y^2(x) + y(x) \log x = 0, \forall x >0$ ; 
\item[(c)]  $E = \mathbb{N}$ ;
\item[(d)] $E = \mathbb{Z}$ ;
\item[(e)] $E = \mathbb{R}^{+}$ ;
\item[(f)] $ E = \mathbb{Q}^n$ ;
\item[(g)] $E = \textrm{ l'ensemble des suites}\,\,\,  (x_n)_{n\in\mathbb{N}}$ de nombres positifs ;
\item[(h)] $E = $ l'ensemble des fonctions r\'eelles qui prennent la valeur $1$ en $0$ ;
\item[(i)] $E = $ l'ensemble des fonctions r\'eelles qui tendent vers $+\infty$ lorsque $x$ tend vers $+\infty$ ; 
\end{enumerate}

\finenonce{002779}



\finexercice


\exercice{3298, quercia, 2010/03/09}
\enonce{003298}{Somme de sous-espaces}

Soient $F,G,H$ trois sous-espaces d'un espace vectoriel $E$.
Comparer $F \cap (G + (F \cap H))$ et $(F \cap G) + (F \cap H)$.

\finenonce{003298}



\finexercice


\exercice{3299, quercia, 2010/03/09}
\enonce{003299}{$F\cap G = F'\cap G'$}

Soient $F,G,F',G'$ des sev d'un ev $E$.

Montrer que si $F\cap G = F'\cap G'$ alors
$\bigl( F+(G\cap F')\bigr) \cap \bigl( F+(G\cap G')\bigr) = F$.

\finenonce{003299}



\finexercice


\exercice{3300, quercia, 2010/03/09}
\enonce{003300}{$E$ n'est pas union de sous-espaces stricts}

Soit $E$ un $ K$-ev non nul et $F_1,\dots,F_n$ des sev stricts de $E$.
On veut montrer que $E \ne F_1 \cup \dots \cup F_n$ :

\begin{enumerate}
  \item Traiter le cas $n = 2$.
  \item Cas général : on suppose $F_n \not\subset F_1 \cup \dots \cup F_{n-1}$
     et on choisit
     $\vec x \in F_n \setminus (F_1 \cup \dots \cup F_{n-1})$ et
     $\vec y \notin F_n$.
  \begin{enumerate}
     \item Montrer que : $\forall\ \lambda \in  K$, $\lambda\vec x + \vec y \notin F_n$.
     \item Montrer que : $\forall\ i \le n-1$, il existe au plus un $\lambda \in  K$ tel que
         $\lambda\vec x + \vec y \in F_i$.
     \item Conclure.
  \end{enumerate}
\end{enumerate}
\finenonce{003300}



\finexercice


\exercice{3323, quercia, 2010/03/09}
\enonce{003323}{Intersection et somme de sev}
Soit $E$ un ev de dimension finie et
$(F_i)_{i\in I}$ une famille de sous-espaces de $E$.

On note $H = \bigcap_{i\in I} F_i$ et
$S = \sum_{i\in I} F_i = {}$vect$\Big(\bigcup_{i\in I} F_i\Big)$.

Montrer qu'il existe une partie finie, $J$, de $I$ telle que :
    $H = \bigcap_{i\in J} F_i$ et $S = \sum_{i\in J} F_i$.
\finenonce{003323}



\finexercice


\exercice{5164, rouget, 2010/06/30}
\enonce{005164}{*T}
Soit $E$ le $\Rr$-espace vectoriel des applications de $[0,1]$ dans $\Rr$ (muni de $f+g$ et $\lambda.f$ usuels)
(ne pas hésiter à redémontrer que $E$ est un $\Rr$ espace vectoriel). Soit $F$ l'ensemble des applications de $[0,1]$
dans $\Rr$ vérifiant l'une des conditions suivantes~:

$$\begin{array}{llll}
1)\;f(0)+f(1)=0&2)\;f(0)=0&3)\;f(\frac{1}{2})=\frac{1}{4}&4)\;\forall x\in[0,1],\;f(x)+f(1-x)=0\\
5)\;\forall x\in[0,1],\;f(x)\geq0&6)\;2f(0)=f(1)+3
\end{array}$$

Dans quel cas $F$ est-il un sous-espace vectoriel de $E$~?
\finenonce{005164}


\finexercice


\exercice{5165, rouget, 2010/06/30}
\enonce{005165}{**T}
On munit $\Rr^n$ des lois produit usuelles. Parmi les sous-ensembles suivants $F$ de $\Rr^n$, lesquels
sont des sous-espaces vectoriels~?

$$\begin{array}{lll}
1)\;F=\{(x_1,...,x_n)\in\Rr^n/\;x_1=0\}&2)\;F=\{(x_1,...,x_n)\in\Rr^n/\;x_1=1\}\\
3)\;F=\{(x_1,...,x_n)\in\Rr^n/\;x_1=x_2\}&4)\;F=\{(x_1,...,x_n)\in\Rr^n/\;x_1+...+x_n=0\}\\
5)\;F=\{(x_1,...,x_n)\in\Rr^n/\;x_1.x_2=0\}
\end{array}$$
\finenonce{005165}


\finexercice


\exercice{5166, rouget, 2010/06/30}
\enonce{005166}{**}
Soit $E$ un $\Kk$-espace vectoriel. Soient $A$, $B$ et $C$ trois sous-espaces vectoriels de $E$ vérifiant $A\cap
B=A\cap C$, $A+B=A+C$ et $B\subset C$. Montrer que $B=C$.
\finenonce{005166}


\finexercice

\exercice{5168, rouget, 2010/06/30}
\enonce{005168}{**T}
Soit $F$ le sous-espace vectoriel de $\Rr^4$ engendré par $u=(1,2,-5,3)$ et $v=(2,-1,4,7)$. Déterminer $\lambda$ et
$\mu$ réels tels que $(\lambda,\mu,-37,-3)$ appartienne à $F$.
\finenonce{005168}


\finexercice


\exercice{5169, rouget, 2010/06/30}
\enonce{005169}{**T}
Montrer que $a=(1,2,3)$ et $b=(2,-1,1)$ engendrent le même sous espace de $\Rr^3$ que $c=(1,0,1)$ et $d=(0,1,1)$.
\finenonce{005169}


\finexercice

\exercice{5172, rouget, 2010/06/30}
\enonce{005172}{**}
 Soient $E$ un $\Kk$-espace vectoriel et $A$, $B$ et $C$ trois sous-espaces de $E$.
\begin{enumerate}
\item  Montrer que~:~$(A\cap B)+(A\cap C)\subset A\cap(B+C)$.
\item  A-t-on toujours l'égalité~?
\item  Montrer que~:~$(A\cap B)+(A\cap C)=A\cap(B+(A\cap C))$.
\end{enumerate}
\finenonce{005172}


\finexercice


\exercice{5173, rouget, 2010/06/30}
\enonce{005173}{**T}
Dans $E=\Rr^4$, on considère $V=\{(x,y,z,t)\in E/\;x-2y=0\;\mbox{et}\;y-2z=0\}$ et $W=\{(x,y,z,t)\in E/\;x+z=y+t\}$.
\begin{enumerate}
\item  Montrer que $V$ et $W$ sont des sous espaces vectoriels de $E$.
\item  Donner une base de $V$, $W$ et $V\cap W$.
\item  Montrer que $E=V+W$.
\end{enumerate}
\finenonce{005173}


\finexercice


\exercice{5174, rouget, 2010/06/30}
\enonce{005174}{***}
Soit $\begin{array}[t]{cccc}
f~:&[0,+\infty[\times[0,2\pi[&\rightarrow&\Rr^2\\
 &(x,y)&\mapsto&(x\cos y,x\sin y)
\end{array}$.

\begin{enumerate}
\item  $f$ est-elle injective~?~surjective~?
\item  Soient $a$, $b$, $\alpha$ et $\beta$ quatre réels. Montrer qu'il existe $(c,\gamma)\in\Rr^2$ tel que~:~
$\forall x\in\Rr,\;a\cos(x-\alpha)+b\cos(x-\beta)=c\cos(x-\gamma)$.
\item  Soit $E$ le $\Rr$-espace vectoriel des applications de $\Rr$ dans $\Rr$. Soit $F=\{u\in
E/\;\exists(a,b,\alpha,\beta)\in\Rr^4\;\mbox{tel que}\;\forall x\in\Rr,\;u(x)=a\cos(x-\alpha)+b\cos(2x-\beta)\}$.
Montrer que $F$ est un sous-espace vectoriel de $E$.
\item  Déterminer $\{\cos x,\sin x,\cos(2x),\sin(2x),1,\cos^2x,\sin^2x\}\cap F$.
\item  Montrer que $(\cos x,\sin x,\cos(2x),\sin(2x))$ est une famille libre de $F$.
\end{enumerate}
\finenonce{005174}


\finexercice


\exercice{5175, rouget, 2010/06/30}
\enonce{005175}{**}
Soit $C$ l'ensemble des applications de $\Rr$ dans $\Rr$, croissantes sur $\Rr$.
\begin{enumerate}
\item  $C$ est-il un espace vectoriel (pour les opérations usuelles)~?
\item  Montrer que $V=\{f\in\Rr^\Rr/\;\exists(g,h)\in C^2\;\mbox{tel que}\;f= g-h\}$ est un $\Rr$-espace
vectoriel.
\end{enumerate}
\finenonce{005175}


\finexercice


\exercice{5176, rouget, 2010/06/30}
\enonce{005176}{**}
Montrer que la commutativité de la loi $+$ est une conséquence des autres axiomes de la structure d'espace vectoriel.
\finenonce{005176}


\finexercice


\exercice{5177, rouget, 2010/06/30}
\enonce{005177}{***}
Soient $E$ un $\Kk$-espace vectoriel et $A$, $B$ et $C$ trois sous-espaces vectoriels de $E$.
Montrer que $$(A\cap B)+(B\cap C)+(C\cap A)\subset(A+B)\cap(B+C)\cap(C+A).$$
\finenonce{005177}


\finexercice

\exercice{5563, rouget, 2010/10/16}
\enonce{005563}{** I}
\label{ex:rou1}
Soient $F$ et $G$ deux sous-espaces vectoriels d'un espace vectoriel $E$.

Montrer que : $\left[(F\cup G\;\text{sous-espace de}\;E)\Leftrightarrow(F\subset G\;\text{ou}\;G\subset F)\right]$.
\finenonce{005563}


\finexercice


\exercice{5564, rouget, 2010/10/16}
\enonce{005564}{****}
Généralisation de l'exercice \ref{ex:rou1}. Soient $n$ un entier supérieur ou égal à $2$ puis $F_1$, ... , $F_n$ $n$ sous-espaces de $E$ où $E$ est un espace vectoriel sur un sous-corps $\Kk$ de $\Cc$.
Montrer que $\left[(F_1\cup ... \cup F_n\;\text{sous-espace de}\; E)\Leftrightarrow(\text{il existe}\;i\in\llbracket1,n\rrbracket/\;\displaystyle\bigcup_{j\neq i}F_j\subset F_i)\right]$.
\finenonce{005564}


\finexercice


\exercice{6868, chataur, 2012/05/13}
\video{QSafAt0zfGk}
\enonce{006868}{}
Montrer que les ensembles ci-dessous sont des espaces vectoriels (sur $\Rr$) :
\begin{itemize}
  \item  $E_1 = \big\{ f : [0,1] \to \Rr \big\}$ : l'ensemble
des fonctions à valeurs r\'eelles d\'efinies sur l'intervalle $[0,1]$, 
muni de l'addition $f+g$ des fonctions et de la multiplication par un nombre r\'eel $\lambda \cdot f$.

  \item $E_2 = \big\{ (u_n) : \Nn \to \Rr \big\}$ : l'ensemble
des suites r\'eelles muni de l'addition des suites définie par $(u_n)+(v_n)=(u_n+v_n)$
et de la multiplication par un nombre r\'eel $\lambda \cdot (u_n) = (\lambda \times u_n)$.

  \item $E_3 = \big\{ P \in \Rr[x] \mid \deg P \le n \big\}$ : l'ensemble des polynômes
à coefficients réels de degré inférieur ou égal à $n$
muni de l'addition $P+Q$ des polynômes et de la multiplication par un nombre r\'eel $\lambda \cdot P$. 
\end{itemize}
\finenonce{006868}



\finexercice


\exercice{6869, chataur, 2012/05/13}
\video{grzoTdCre2E}
\enonce{006869}{}
\begin{enumerate}
  \item D\'ecrire les sous-espaces vectoriels de $\Rr$ ;  puis de $\Rr^2$ et $\Rr^3$.
  \item Dans $\mathbb{R}^3$ donner un exemple de deux sous-espaces dont l'union 
n'est pas un sous-espace vectoriel.
\end{enumerate}


\finenonce{006869}



\finexercice

\section{106.02 Système de vecteurs}


\exercice{897, liousse, 2003/10/01}

\enonce{000897}{}
Soient dans $\Rr^3$ les vecteurs $\vec{v_1}(1,1,0)$, $\vec{v_2}(4,1,4)$ et 
$\vec{v_3}(2,-1,4)$.
\begin{enumerate}
\item Montrer que $\vec{v_1}$ et $\vec{v_2}$ ne sont pas colin\'eaires. 
Faire de m\^eme avec 
$\vec{v_1}$ et $\vec{v_3}$, puis avec $\vec{v_2}$ et $\vec{v_3}$.
\item La famille $(\vec{v_1}, \vec{v_2}, \vec{v_3})$ est-elle libre ?
 \end{enumerate}
\finenonce{000897}



\finexercice



\exercice{898, liousse, 2003/10/01}

\enonce{000898}{}
Les familles suivantes sont-elles libres ?
\begin{enumerate}
  \item  $\vec{v_1}(1,0,1)$, $\vec{v_2}(0,2,2)$ et $\vec{v_3}(3,7,1)$ dans $\Rr^3$.
 \item  $\vec{v_1}(1,0,0)$, $\vec{v_2}(0,1,1)$ et $\vec{v_3}(1,1,1)$ dans $\Rr^3$.
 \item $\vec{v_1}(1,2,1,2,1)$, $\vec{v_2}(2,1,2,1,2)$, $\vec{v_3}(1,0,1,1,0)$ 
et $\vec{v_4}(0,1,0,0,1)$ dans $\Rr^5$.
 \item  $\vec{v_1}(2,4,3,-1,-2,1)$, $\vec{v_2}(1,1,2,1,3,1)$ et 
$\vec{v_3}(0,-1,0,3,6,2)$ dans $\Rr^6$.
 \item  $\vec{v_1}(2,1,3,-1,4,-1)$, $\vec{v_2}(-1,1,-2,2,-3,3)$ et 
$\vec{v_3}(1,5,0,4,-1,7)$ dans $\Rr^6$.
 \end{enumerate}
\finenonce{000898}



\finexercice



\exercice{899, liousse, 2003/10/01}

\enonce{000899}{}
On consid\`ere dans $\Rr^n$ 
une famille de $4$ vecteurs lin\'eairement ind\'ependants : 
$(\vec{e_1}, \vec{e_2}, \vec{e_3}, \vec{e_4})$.
Les familles suivantes sont-elles libres ?
 \begin{enumerate}
  \item $(\vec{e_1}, 2\vec{e_2}, \vec{e_3})$. 
 \item $(\vec{e_1}, \vec{e_3})$.
 \item $(\vec{e_1}, 2\vec{e_1}+\vec{e_4}, \vec{e_4})$.
 \item $(3\vec{e_1}+\vec{e_3}, \vec{e_3}, \vec{e_2}+\vec{e_3})$.
 \item $(2\vec{e_1}+\vec{e_2}, \vec{e_1}-3\vec{e_2}, 
\vec{e_4}, \vec{e_2}-\vec{e_1})$.
 \end{enumerate}
\finenonce{000899}



\finexercice



\exercice{900, liousse, 2003/10/01}
\video{aMookydUKT4}
\enonce{000900}{}
Soient dans $\Rr^4$ les
vecteurs $v_1=(1,2,3,4)$ et $v_2=(1,-2,3,-4)$. Peut-on
d\'eterminer $x$ et $y$ pour que $(x,1,y,1) \in \text{Vect}\{ v_1, v_2 \}$ ? 
Et pour que $(x,1,1,y) \in \text{Vect}\{v_1,v_2 \}$ ?
\finenonce{000900} 



\finexercice


\exercice{901, liousse, 2003/10/01}
\video{A-ieYnIs2KA}
\enonce{000901}{}
 Dans $\Rr^4$ on consid\`ere
l'ensemble $E$ des vecteurs $(x_1,x_2,x_3,x_4)$ v\'erifiant
$x_1+x_2+x_3+x_4=0$. L'ensemble $E$ est-il un sous-espace vectoriel de
$\Rr^4$ ? Si oui, en donner une base.
\finenonce{000901} 


\finexercice


\exercice{902, cousquer, 2003/10/01}

\enonce{000902}{}
 Dans l'espace $\mathbb{R}^4$, on se donne cinq vecteurs~:
$V_1=(1,1,1,1)$, $V_2=(1,2,3,4)$, $V_3=(3,1,4,2)$, $V_4=(10,4,13,7)$,
$V_5=(1,7,8,14)$.
Chercher les relations de dépendance linéaires entre ces
vecteurs. Si ces vecteurs sont dépendants, en extraire au moins une famille
libre engendrant le même sous-espace.
\finenonce{000902}



\finexercice



\exercice{903, cousquer, 2003/10/01}

\enonce{000903}{}
Dans l'espace $\mathbb{R}^4$, on se donne cinq vecteurs~:
$V_1=(1,1,1,1)$, $V_2=(1,2,3,4)$, $V_3=(3,1,4,2)$, $V_4=(10,4,13,7)$,
$V_5=(1,7,8,14)$.
À quelle(s) condition(s) un vecteur $B=(b_1,b_2,b_3,b_4)$
appartient-il au sous-espace engendré par les vecteurs $V_1$, $V_2$, 
$V_3$, $V_4$, $V_5$~? Définir ce sous-espace par une ou des équations. 
\finenonce{000903}



\finexercice



\exercice{904, cousquer, 2003/10/01}

\enonce{000904}{}
 Soient  les vecteurs $e_1=(1,2,3,4)$, $e_2=(1,-2,3,-4)$ de $\mathbb{R}^4$.
Peut-on déterminer $x$ et $y$ pour que $(x,1,y,1) \in\mbox{Vect}\{e_1,e_2\}$~?
pour que $(x,1,1,y)\in \mbox{Vect}\{e_1,e_2\}$~?
\finenonce{000904}



\finexercice



\exercice{905, cousquer, 2003/10/01}

\enonce{000905}{}
 Soit $E$ un espace vectoriel sur $\mathbb{R}$ et $x,y,z,t$ une famille libre
d'éléments de $E$, les familles suivantes sont-elles libres~?
\begin{enumerate}
\item  $x,2y,z$.
\item  $x,z$.
\item  $x,2x+t,t$.
\item  $3x+z, z, y+z$.
\item  $2x+y,x-3y,t, y-x$. 
\end{enumerate}
\finenonce{000905}



\finexercice



\exercice{906, cousquer, 2003/10/01}

\enonce{000906}{}
 Dans $\mathbb{R}^4$, comparer les sous-espaces $F$ et $G$ suivants~:
\begin{eqnarray*}
    F & = & \mbox{Vect}\{(1,0,1,1), (-1,-2,3,-1), (-5,-3,1,-5)\}\\
    G & = & \mbox{Vect}\{(-1,-1,1,-1), (4,1,2,4)\}
\end{eqnarray*}
\finenonce{000906}



\finexercice



\exercice{907, cousquer, 2003/10/01}

\enonce{000907}{}
 On suppose que $v_1, v_2, v_3, \ldots,v_n$ sont des vecteurs indépendants
de $\mathbb{R}^n$.
\begin{enumerate}


\item Les vecteurs $v_1-v_2,  v_2-v_3, v_3-v_4, \ldots, v_n-v_1$ sont-ils
linéairement indépendants~?
\item Les vecteurs $v_1+v_2,  v_2+v_3, v_3+v_4, \ldots, v_n+v_1$ sont-ils
linéairement indépendants~?
\item Les vecteurs $v_1, v_1+v_2, v_1+v_2+v_3, v_1+v_2+v_3+v_4, \ldots, v_1+v_2+
\cdots+v_n$ sont-ils linéairement indépendants~?
\end{enumerate}
\finenonce{000907}



\finexercice



\exercice{908, legall, 1998/09/01}
\video{SmI00EAW31k}
\enonce{000908}{}
Soit $E$ le 
sous-espace vectoriel de ${\Rr}^3$ engendr\'e par
les vecteurs $v_1=(2, 3, -1)$ et $v_2=(1, -1, -2)$  et $F$ celui engendré par
 $w_1=(3, 7, 0)$ et $w_2=(5, 0, -7)$. Montrer que  $E$  et  $F$  sont
 \'egaux.
\finenonce{000908} 


\finexercice


\exercice{909, legall, 1998/09/01}

\enonce{000909}{}
Prouver que dans ${ \Rr}^{3}$, les vecteurs $u_{1}=(2,3,-1)$ et $u_{2}=(1,-1,-2)$
engendrent le m\^eme s.e.v. que les vecteurs $v_{1}=(3,7,0)$ et $v_{2}=(5,0,-7)$.\\
\finenonce{000909}



\finexercice



\exercice{910, legall, 1998/09/01}

\enonce{000910}{}
\begin{enumerate}
    \item Montrer que les syst\`emes  : $S_{1}=(1;\sqrt{2})$ et
$S_{2}=(1;\sqrt{2};\sqrt{3})$ sont libres dans ${ \Rr}$ consid\'er\'e comme ${ \Qq}$-espace
vectoriel.
    \item Soient, dans ${ \Rr}^{2}$, les vecteurs $u_{1}=(3+\sqrt{5},2+3\sqrt{5})$ et
$u_{2}=(4,7\sqrt{5}-9)$. Montrer que le syst\`eme $(u_{1},u_{2})$ est ${ \Qq}$-libre et ${ \Rr}$-li\'e.
    \item Soient les vecteurs
$v_{1}=(1-i,i)$ et $v_{2}=(2,-1+i)$ dans
${ \Cc}^{2}$.
    \begin{enumerate}
        \item Montrer que le syst\`eme $(v_{1},v_{2})$ est ${ \Rr}$-libre et ${\Cc}$-li\'e.
        \item V\'erifier que le syst\`eme $S=\{(1,0),(i,0),(0,1),(0,i)\}$ est une base de l'e.v. ${ \Cc}^{2}$ sur ${ \Rr}$, et donner les composantes des vecteurs $v_{1}, v_{2}$ par rapport \`a cette base.
    \end{enumerate}
\end{enumerate}
\finenonce{000910}



\finexercice



\exercice{911, legall, 1998/09/01}

\enonce{000911}{}
\begin{enumerate}
    \item On d\'efinit les fonctions suivantes : $f_{1} : t \mapsto \cos t .\mbox{ch}t,
f_{2} : t \mapsto \cos t .\mbox{sh}t,   f_{3} : t \mapsto \sin t .\mbox{ch}t,   f_{4} : t \mapsto
\sin t .\mbox{sh} t $. Montrer que le syst\`eme $(f_{1},f_{2},f_{3},f_{4})$ est libre dans ${ \Rr}^{ \Rr}$.
    \item M\^eme question pour la famille ${\mathcal F}=\{f_{\lambda} : t \mapsto e^{\lambda t},   \lambda \in { \Rr}\}$.
\end{enumerate}
\finenonce{000911}



\finexercice



\exercice{912, cousquer, 2003/10/01}

\enonce{000912}{}
 Dans $\mathcal{F}(\mathbb{R},\mathbb{R})$, les trois fonctions 
$x\mapsto\sin x$,
$x\mapsto\sin 2x$, $x\mapsto\sin 3x$, sont-elles linéairement
indépendantes~? Généraliser.
\finenonce{000912}



\finexercice



\exercice{913, legall, 1998/09/01}

\enonce{000913}{}
Soit $E$ un ${ \Cc}$-espace vectoriel et $S_{1}=(e_{1},e_{2},...,e_{n})$ un
syst\`eme libre dans $E$, $n\ge 2$.
\begin{enumerate}
    \item On consid\`ere le syst\`eme $S_{2}=(e_{1}',e_{2}',...,e_{n}')$ d\'efini par :
$e_{j}'=\sum_{k=1}^{j}e_{k},   1\le j \le n$. $S_{2}$ est-il libre ?
    \item  On consid\`ere le syst\`eme $S_{3} = (\epsilon_{1},\epsilon_{2},...,\epsilon_{n})$ d\'efini par
: $\epsilon_{j}=e_{j}+e_{j+1},   1\le j \le n-1$ et $\epsilon_{n}=e_{n}+e_{1}$. Montrer les
r\'esultats suivants :
    \begin{enumerate}
        \item  $S_{3}$ libre $\Rightarrow$ $S_{1}$ libre.
        \item $n$ impair : $S_{3}$
libre $\Leftrightarrow$ $S_{1}$ libre.
\item $n$ pair : $S_{3}$ li\'e.
    \end{enumerate}
\end{enumerate}
\finenonce{000913}



\finexercice



\exercice{914, legall, 1998/09/01}
\enonce{000914}{}
Peut-on d\'eterminer des
r\'eels $x, y$ pour que le vecteur $v=(-2,x,y,3)$ appartienne au
s.e.v. engendr\'e dans ${ \Rr}^{4}$ par le syst\`eme $(e_{1},e_{2})$
o\`u $e_{1}=(1,-1,1,2)$ et $e_{2}=(-1,2,3,1)$ ?
\finenonce{000914} 


\finexercice


\exercice{915, ridde, 1999/11/01}

\enonce{000915}{}
Soient $f (x) = \cos (x)$, $g (x) = \cos (x) \cos (2x)$ et $h (x) = \sin (x) \sin (2x)$.
D\'eterminer $\text{vect} (f, g, h)$.
\finenonce{000915}



\finexercice



\exercice{916, ridde, 1999/11/01}
\enonce{000916}{}
Soit $\alpha \in \Rr$ et
soit $f_{\alpha} : \Rr\rightarrow \Rr$ la fonction définie par 
$$
\begin{cases}
f_{\alpha}(x)= 1 & \text{ si } x = \alpha \\
f_{\alpha}(x)= 0 & \text{ si } x \neq \alpha
 \end{cases}
.$$ 
Montrer que la famille $ (f_{\alpha})_{\alpha \in\Rr}$ est libre.
\finenonce{000916} 



\finexercice

\exercice{917, ridde, 1999/11/01}
\video{x8fXA1UP7LU}
\enonce{000917}{}
Soit $\alpha \in \Rr$ et $f_\alpha : \Rr \to \Rr$, $x\mapsto e^{\alpha x}$.
Montrer que la famille $(f_\alpha)_{\alpha \in \Rr}$  est libre.
\finenonce{000917}





\finexercice


\exercice{918, gourio, 2001/09/01}

\enonce{000918}{}
Montrer que les familles suivantes sont libres dans ${\Rr}^{{\Rr}},$ et
ce quelque soit $N\in \Nn^{*}$ :
$$\left( x\rightarrow \left| x-a\right| \right) _{a=1,3,5,...,2N+1};\left(
x\rightarrow \cos nx\right) _{n=1,2,...,N};\left( x\rightarrow e^{ax}\right)
_{a=1,...,N} $$
\finenonce{000918}



\finexercice



\exercice{2430, matexo1, 2002/02/01}

\enonce{002430}{}
Soit $E$ l'ensemble des suites r\'eelles $(u_n)_{n \ge 0}$.
Pour tout $k \in \N$, on note $\delta_k$ l'\'el\'ement de $E$
dont les coordonn\'ees sont toutes nulles, sauf $\delta_{k,k} =
1$. Montrer que la famille infinie ${\cal B} = \{\delta_k\}_{k \in
\N}$ est libre (en ce sens que toute sous-famille finie est
libre).
Soit $E_0$ le sous-ensemble des suites qui convergent vers z\'ero.
Montrer que c'est un sous-espace vectoriel de $E$. Montrer que
$\cal B$ est aussi une famille libre de $E_0$.

\finenonce{002430}


\finexercice


\exercice{2435, matexo1, 2002/02/01}

\enonce{002435}{}
Soit $E$ un espace vectoriel r\'eel et $u$ un endomorphisme
de $E$ tel que $u^2 = -I$.

\begin{enumerate}
\item Montrer que $u$ est bijectif.
\item On suppose que les $2p-1$ vecteurs $x_1,\ldots,x_p,
u(x_1),\ldots,u(x_{p-1})$ sont lin\'eairement ind\'ependants.
Montrer que les $2p$ vecteurs $x_1,\ldots,x_p,$ 
$u(x_1),\ldots,u(x_p)$ sont lin\'eairement ind\'ependants.
\item On suppose que $E$ est de dimension finie. Montrer que
$E$ poss\`ede une base de la forme $x_1,\ldots,x_p,
u(x_1),\ldots,u(x_p)$ et est de dimension paire. Donner la
matrice de $u$ dans cette base.
\end{enumerate}

\finenonce{002435}


\finexercice


\exercice{2450, matexo1, 2002/02/01}

\enonce{002450}{}
Dans $\R^4$, on note $v_1 = {\null}^t(1, 2, 0, -1)$, $v_2 = {\null}^t(3, 2, -1, -1)$,
$v_3 = {\null}^t(-1, 2, 1, -3)$ et $v_4 = {\null}^t(1, -1, 1, -1)$. Sont-ils
lin\'eairement ind\'ependants\,? Trouver une relation de d\'ependance lin\'eaire entre
eux.
\finenonce{002450}


\finexercice


\exercice{2459, matexo1, 2002/02/01}

\enonce{002459}{$\mathbb C$ isomorphe \`a un sous-espace de $\mathcal M_2 (\mathbb{R})$}

Soit $E$ l'espace vectoriel des matrices carr\'ees r\'eelles
d'ordre 2.

\begin{enumerate}
\item Montrer que les ``vecteurs"
$$ I = \left( \begin{array}{cc} 1&0\\ 0&1 \end{array} \right), \,
J = \left( \begin{array}{cc} 0&1 \\ -1&0 \end{array} \right), \,
K = \left( \begin{array}{cc} 0&1 \\ 1&0 \end{array} \right), \,
L = \left( \begin{array}{cc} 1&0 \\ 0&-1 \end{array} \right)$$
de $E$ sont lin\'eairement ind\'ependants.
\item Montrer que tout \'el\'ement $X = \left( \begin{array}{cc} a&b\\ c&d \end{array} \right)$
de $E$ s'\'ecrit de fa\c con unique sous la forme $X = x_1 I +
x_2 J + x_3 K + x_4 L$ et calculer $x_1, x_2, x_3, x_4$ en
fonction de $a,b,c,d$.
\item V\'erifier la relation $J^2 = -I$. Calculer $JX$ et $XJ$.
Montrer que l'\'equation $XJ = JX$ est \'equivalente \`a $x_3 =
x_4 = 0$. En d\'eduire que le sous-espace de $E$ engendr\'e par $I,
J$ est isomorphe au corps des complexes $\mathbb C$.
\end{enumerate}

\finenonce{002459}


\finexercice


\exercice{2780, tumpach, 2009/10/25}
\enonce{002780}{}
Soient dans $\mathbb{R}^3$ les vecteurs $\vec{v_1}= (1,1,0)$, $\vec{v_2}= (4,1,4)$ et 
$\vec{v_3}= (2,-1,4)$.
\begin{enumerate}
\item Montrer que $\vec{v_1}$ et $\vec{v_2}$ ne sont pas colin\'eaires. 
Faire de m\^eme avec 
$\vec{v_1}$ et $\vec{v_3}$, puis avec $\vec{v_2}$ et $\vec{v_3}$.

\item La famille $(\vec{v_1}, \vec{v_2}, \vec{v_3})$ est-elle libre ?
\end{enumerate}
\finenonce{002780}



\finexercice


\exercice{2781, tumpach, 2009/10/25}
\enonce{002781}{}
Les familles suivantes sont-elles libres ?
\begin{enumerate}
\item $\vec{v_1} = (1,0,1)$, $\vec{v_2} = (0,2,2)$ et $\vec{v_3} = (3,7,1)$ dans $\mathbb{R}^3$.

\item $\vec{v_1}= (1,0,0)$, $\vec{v_2} = (0,1,1)$ et $\vec{v_3}= (1,1,1)$ dans $\mathbb{R}^3$.

\item $\vec{v_1} = (1,2,1,2,1)$, $\vec{v_2} = (2,1,2,1,2)$, $\vec{v_3} = (1,0,1,1,0)$ 
et $\vec{v_4} = (0,1,0,0,1)$ dans $\mathbb{R}^5$.

\item $\vec{v_1} = (2,4,3,-1,-2,1)$, $\vec{v_2} = (1,1,2,1,3,1)$ et 
$\vec{v_3} = (0,-1,0,3,6,2)$ dans $\mathbb{R}^6$.

\item $\vec{v_1} = (2,1,3,-1,4,-1)$, $\vec{v_2} = (-1,1,-2,2,-3,3)$ et 
$\vec{v_3} = (1,5,0,4,-1,7)$ dans $\mathbb{R}^6$.\\
\end{enumerate}
\finenonce{002781}



\finexercice


\exercice{2782, tumpach, 2009/10/25}
\enonce{002782}{}
 On suppose que $v_1, v_2, v_3, \ldots,v_n$ sont des vecteurs ind\'ependants
de $\mathbb{R}^n$.

\begin{enumerate}
\item Les vecteurs $v_1-v_2,  v_2-v_3, v_3-v_4, \ldots, v_n-v_1$ sont-ils
lin\'eairement ind\'ependants~?

\item  Les vecteurs $v_1+v_2,  v_2+v_3, v_3+v_4, \ldots, v_n+v_1$ sont-ils
lin\'eairement ind\'ependants~?

\item Les vecteurs $v_1, v_1+v_2, v_1+v_2+v_3, v_1+v_2+v_3+v_4, \ldots, v_1+v_2+
\cdots+v_n$ sont-ils lin\'eairement ind\'ependants?\\
\end{enumerate}
\finenonce{002782}



\finexercice


\exercice{3297, quercia, 2010/03/09}
\enonce{003297}{Sev de $ K^3$ engendrés par deux vecteurs}

On considère les vecteurs de $ K^3$ :
$\vec a = (1, 2, 1)$,
$\vec b = (1, 3, 2)$,
$\vec c = (1, 1, 0)$,
$\vec d = (3, 8, 5)$.

Soient $F = \text{vect}(\vec a,\vec b\,)$ et
       $G = \text{vect}(\vec c,\vec d\,)$.
Comparer $F$ et $G$.

\finenonce{003297}



\finexercice

\exercice{3301, quercia, 2010/03/09}
\enonce{003301}{\'Etude de liberté}

\'Etudier la liberté des familles suivantes :
\begin{enumerate}
  \item $E=\{  {\text{fcts}} : \R    \to    \R \}$, ${\cal F}=(\sin, \cos)$.
  \item $E=\{ {\text{fcts}} : {\R^{+*}} \to \R \}$, ${\cal F}=(f_a : x \longmapsto x^a),\ a\in\R$.
  \item $E=\{ {\text{fcts}} : \R      \to  \R \}$, ${\cal F}=(f_a : x \longmapsto |x-a|),\ a\in\R$.
\end{enumerate}
\finenonce{003301}



\finexercice


\exercice{3302, quercia, 2010/03/09}
\enonce{003302}{Nombres algébriques}

On considère que $\R$ est un $\Q$-espace vectoriel.

\begin{enumerate}
  \item Montrer que la famille $\big(1,\sqrt2,\sqrt3\big)$ est libre.
  \item Montrer que la famille $(\ln p)$ où $p$ décrit l'ensemble des nombres
    premiers positifs est libre.
\end{enumerate}
\finenonce{003302}



\finexercice


\exercice{3303, quercia, 2010/03/09}
\enonce{003303}{Modification des vecteurs d'une famille libre}

Soit $E$ un espace vectoriel, $(\vec x_1,\dots,\vec x_n)$ une famille libre de
vecteurs de $E$, et $\alpha_1,\dots,\alpha_n$ des scalaires.

{\def\v#1{{\vec x_{\rlap{$\scriptstyle#1$}}\hskip 1pt}'}
On pose $\vec y = \sum_{i=1}^n \alpha_i\vec x_i$, et $\v i = \vec x_i + \vec y$.
\'Etudier à quelle condition la famille $(\v1,\dots,\v n\,)$ est libre.
}
\finenonce{003303}



\finexercice


\exercice{3304, quercia, 2010/03/09}
\enonce{003304}{Polynômes trigonométriques}

Soit $E$ l'ev $\R^\R$,
$F$ le sev engendré par les fonctions $f_n\,:\,x \mapsto \cos(nx)$, $n \in \N$, et
$G$ le sev engendré par les fonctions $g_n\,:\,x \mapsto \cos^nx$, $n \in \N$.
Montrer que $F$ = $G$.
\finenonce{003304}



\finexercice

\exercice{5167, rouget, 2010/06/30}
\enonce{005167}{**T}
Soit $\Rr^\Nn$ le $\Rr$-espace vectoriel des suites réelles (muni des opérations usuelles). On considère les trois
éléments de $E$ suivants~:~$u=(\cos(n\theta))_{n\in\Nn}$, $v=(\cos(n\theta+a))_{n\in\Nn}$ et
$w=(\cos(n\theta+b))_{n\in\Nn}$ où $\theta$, $a$ et $b$ sont des réels donnés. Montrer que $(u,v,w)$ est une famille
liée.
\finenonce{005167}


\finexercice

\exercice{5180, rouget, 2010/06/30}
\enonce{005180}{***T}
Dans $E=\Rr^\Rr$, étudier la liberté des familles suivantes $A$ de vecteurs de $E$~:
\begin{enumerate}
\item  $a$, $b$ et $c$ étant trois réels donnés, $A=(f_a,f_b,f_c)$ où, pour tout réel $x$, $f_u(x)=\sin(x+u)$.
\item  $A=(f_n)_{n\in\Zz}$ où, pour tout réel $x$, $f_n(x)=nx+n^2+1$.
\item  $A=(x\mapsto x^\alpha)_{\alpha\in\Rr}$ (ici $E=(]0;+\infty[)^2$).
\item  $A=(x\mapsto|x-a|)_{a\in\Rr}$.
\end{enumerate}
\finenonce{005180}


\finexercice

\exercice{5566, rouget, 2010/10/16}
\enonce{005566}{**}
Les familles suivantes de $\Rr^4$ sont-elles libres ou liées ? Fournir des relations de dépendance linéaire quand ces relations existent.

\begin{enumerate}
 \item  $(e_1,e_2,e_3)$ où $e_1=(3,0,1,-2)$, $e_2=(1,5,0,-1)$ et $e_3=(7,5,2,1)$.

\item  $(e_1,e_2,e_3,e_4)$ où $e_1=(1,1,1,1)$, $e_2=(1,1,1,-1)$, $e_3=(1,1,-1,1)$ et $e_4=(1,-1,1,1)$.

\item  $(e_1,e_2,e_3,e_4)$ où $e_1=(0,0,1,0)$, $e_2=(0,0,0,1)$, $e_3=(1,0,0,0)$ et $e_4=(0,1,0,0)$.

\item  $(e_1,e_2,e_3,e_4)$ où $e_1=(2,-1,3,1)$, $e_2=(1,1,1,1)$, $e_3=(4,1,5,3)$ et $e_4=(1,-2,2,0)$.
\end{enumerate}
\finenonce{005566}


\finexercice


\exercice{5567, rouget, 2010/10/16}
\enonce{005567}{***}
Montrer que $(1,\sqrt{2},\sqrt{3})$ est une famille libre du $\Qq$-espace vectoriel $\Rr$.
\finenonce{005567}


\finexercice


\exercice{5568, rouget, 2010/10/16}
\enonce{005568}{**}
Soit $f(x) =\ln(1+x)$ pour $x$ réel positif. Soient $f_1 = f$, $f_2=f\circ f$ et $f_3= f\circ f\circ f$. Etudier la liberté de $(f_1,f_2,f_3)$ dans $[0,+\infty[^{[0,+\infty[}$.
\finenonce{005568}


\finexercice


\exercice{5569, rouget, 2010/10/16}
\enonce{005569}{**}
Soit $f_a(x)=|x-a|$ pour $a$ et $x$ réels. Etudier la liberté de la famille $(f_a)_{a\in\Rr}$.
\finenonce{005569}


\finexercice


\exercice{5570, rouget, 2010/10/16}
\enonce{005570}{**I}
On pose $f_a(x)=e^{ax}$ pour $a$ et $x$ réels. Etudier la liberté de la famille de fonctions $(f_a)_{a\in\Rr}$.
\finenonce{005570}


\finexercice


\exercice{5571, rouget, 2010/10/16}
\enonce{005571}{**}
Montrer que toute suite de polynômes non nuls de degrés deux à deux distincts est libre.

Montrer que toute suite de polynômes non nuls de valuations deux à deux distinctes est libre.
\finenonce{005571}


\finexercice


\exercice{5574, rouget, 2010/10/16}
\enonce{005574}{**}
 \begin{enumerate}
 \item  Calculer pour $p$ et $q$ entiers naturels donnés les intégrales suivantes :

\begin{center}
$J(p,q)=\int_{0}^{2\pi}\cos(px)\cos(qx)\;dx$, $K(p,q)=\int_{0}^{2\pi}\cos(px)\sin(qx)\;dx$ et $L(p,q)=\int_{0}^{2\pi}\sin(px)\sin(qx)\;dx$.
\end{center}

\item  Montrer que la famille de fonctions $(\cos(px))_{p\in\Nn}\cup(\sin(qx))_{q\in\Nn^*}$ est libre.
\end{enumerate}
\finenonce{005574}


\finexercice


\exercice{6870, chataur, 2012/05/13}
\video{5yA5SIqdVAQ}
\enonce{006870}{}
\begin{enumerate}
  \item Soient $v_1=(2,1,4)$, $v_2=(1,-1,2)$ et $v_3=(3,3,6)$ des vecteurs de $\Rr^3$, 
trouver trois r\'eels non tous nuls $\alpha,\beta,\gamma$ tels que $\alpha v_1+ \beta v_2 + \gamma v_3=0$.

  \item On considère deux plans vectoriels
$$P_1=\{(x,y,z) \in \Rr^3 \mid x-y+z=0\}$$
$$P_2=\{(x,y,z) \in \Rr^3 \mid x-y=0\}$$
trouver un vecteur directeur de la droite $D=P_1\cap P_2$ ainsi qu'une \'equation param\'etr\'ee.
\end{enumerate}
\finenonce{006870}





\finexercice


\exercice{7410, mourougane, 2021/08/10}
\enonce{007410}{}
Dans $\Rr^4$ on considère les quatre vecteurs
$$v_1=(1,0,-1,1),\ \ v_2=(2,1,0,1), \ \ v_3 =(1,1,1,0), \ \ v_4=(3,1,-1,2). $$ 
Soit $V=Vect(v_1,v_2,v_3,v_4)$. De plus, soit 
$$H=\{(x,y,z,t)\in\Rr^4 \ / \ \ -3x+y+2z-t=0 \}.$$

\begin{enumerate}
 \item Montrer que $\dim V =2$. Le systême $(v_1,v_2,v_3,v_4)$ est-il libre? Est-il générateur de $\Rr^4$?

\item Donner une base de $V$, la compléter en une base de $\Rr^4$.

\item Calculer des équations cartésiennes pour $V$.

\item Montrer que $H$ est un sous-espace vectoriel de $\Rr^4$.

\item Trouver une représentation paramétrique de $H$, et en déduire une base de $H$. Que vaut $\dim H$? 

\item Montrer que $v_3\in H$ et que $v_1\notin H$. En déduire $\dim(V\cap H)$ et $\dim(V+H)$.

\item Donner une base de $V\cap H$.
 \end{enumerate}
\finenonce{007410}
\finexercice

\section{106.03 Somme directe}


\exercice{919, liousse, 2003/10/01}
\video{lE3VGElRQrs}
\enonce{000919}{}
Soient $v_1=(0,1,-2,1),
v_2=(1,0,2,-1), v_3=(3,2,2,-1), v_4 = (0,0,1,0)$ et
$v_5=(0,0,0,1)$ des vecteurs de $\Rr^4$.  Les propositions
suivantes sont-elles vraies ou fausses ?  Justifier votre r\'eponse.
 \begin{enumerate}
 \item $\text{Vect}\{ v_1, v_2, v_3 \} =  \text{Vect}\{(1,1,0,0),(-1,1,-4,2)\}$.
 \item $(1,1,0,0) \in \text{Vect}\{ v_1, v_2 \} \cap \text{Vect}\{ v_2, v_3, v_4 \}$.
 \item $\dim(\text{Vect}\{ v_1, v_2 \} \cap \text{Vect}\{ v_2, v_3, v_4 \})=1$ (c'est-à-dire c'est une droite vectorielle).
 \item $\text{Vect}\{ v_1, v_2 \} + \text{Vect}\{ v_2, v_3, v_4 \}= \Rr^4$.
 \item $\text{Vect}\{ v_4, v_5 \}$ est un sous-espace vectoriel 
   suppl\'ementaire de $\text{Vect}\{ v_1, v_2, v_3 \}$ dans $\Rr^4$.
 \end{enumerate}

\finenonce{000919} 


\finexercice


\exercice{920, cousquer, 2003/10/01}
\video{ByyFihncOvA}
\enonce{000920}{}
 On consid\`ere les vecteurs
$v_1=(1,0,0,1)$, $v_2=(0,0,1,0)$, $v_3=(0,1,0,0)$, $v_4=(0,0,0,1)$,
$v_5=(0,1,0,1)$ dans $\mathbb{R}^4$.
\begin{enumerate}
\item $\mbox{Vect}\{v_1,v_2\}$ et $\mbox{Vect}\{v_3\}$ sont-ils
  suppl\'ementaires dans $\mathbb{R}^4$ ?
\item $\mbox{Vect}\{v_1,v_2\}$ et $\mbox{Vect}\{v_4, v_5\}$ sont-ils
  suppl\'ementaires dans $\mathbb{R}^4$ ?
\item $\mbox{Vect}\{v_1,v_3,v_4\}$ et  $\mbox{Vect}\{v_2,v_5\}$ sont-ils
  suppl\'ementaires dans $\mathbb{R}^4$ ?
\item $\mbox{Vect}\{v_1,v_4\}$ et $\mbox{Vect}\{v_3, v_5\}$ sont-ils
  suppl\'ementaires dans $\mathbb{R}^4$ ?
\end{enumerate}
\finenonce{000920} 


\finexercice


\exercice{921, cousquer, 2003/10/01}

\enonce{000921}{}
  Si $L,M,N$ sont trois sous-espaces vectoriels de $E$, a-t-on~:
$$L\cap(M+N) = L\cap M + L \cap N\;?$$
\finenonce{000921}



\finexercice



\exercice{922, legall, 1998/09/01}

\enonce{000922}{}
Soit $E = \R[X]$ l'espace vectoriel des polyn\^omes. On d\'efinit
$$E_a =\left\{  P\in E ;  (X-a)/P\right\} $$
pour $a\in \R$. Montrer que si $a\not = b$ alors $E = E_a
+ E_b$. La somme est-elle directe?
\finenonce{000922}



\finexercice



\exercice{923, ridde, 1999/11/01}
\video{40VUaSvJ4DY}
\enonce{000923}{}
Soit $E = \Delta^1 (\Rr, \Rr)$ l'espace des fonctions dérivables
et $F = \left\{ f \in E \mid f (0) = f' (0) = 0\right\}$. Montrer que $F$
est un sous-espace vectoriel de $E$ et d\'eterminer un
suppl\'ementaire de $F$ dans $E$.
\finenonce{000923} 


\finexercice


\exercice{924, legall, 1998/09/01}

\enonce{000924}{}
Soient  $E$  un espace vectoriel, $F$  et  $G$  deux sous-espaces
vectoriels de  $E$. On dit que  $F$  et  $G$  sont
{\em suppl\' ementaires} dans  $E$  lorsque  $F\cap G =\{ 0 \} $  et  $E=F+G$. On
note  $E=F\oplus G $.
\begin{enumerate}
    \item Soient  $e_1=\begin{pmatrix} 1 \\ 1 \\ 0 \\ 0 \\ \end{pmatrix}  ,
e_2=\begin{pmatrix} 0 \\ 1 \\ 1 \\ 0 \\  \end{pmatrix}  ,
e_3=\begin{pmatrix} 1 \\ 1 \\ 0 \\ 1 \\  \end{pmatrix}  ,
e_4=\begin{pmatrix} 1 \\ 0 \\ 0 \\ 0 \\  \end{pmatrix}
\hbox{  et  }
e_5=\begin{pmatrix} 1 \\ 1 \\ 1 \\ 1 \\  \end{pmatrix} $  des vecteurs de  ${ \Rr}^4$. Posons  $F=
\hbox{Vect }\{ e_1   ,   e_2  \}  ,
G=
\hbox{Vect }\{ e_3   ,   e_4  \}   ,
G'=
\hbox{Vect }\{ e_3   ,   e_4   ,   e_5 \} $. Montrer
que  $E=F\oplus G $  et  $E\not =F\oplus G' $.
    \item Supposons que  $E$  est de dimension
finie  $n$, que  $\hbox{dim }(F)=p$  et  $E=F\oplus G $.
\begin{enumerate}
    \item Calculer  $\hbox{dim }(G)$.
    \item Montrer que tout \' el\' ement  $x$  de  $E$  se d\' ecompose d'une mani\`ere
{\em unique} en une somme  $x=y+z $  avec  $y\in F$  et  $z\in G$.
    \item Soient  $\mathcal{F} =\{ f_1  , \cdots ,  f_k \}$  une famille libre de  $F$  et  $\mathcal{G}
=\{ g_1  ,\cdots ,  g_l \}$  une famille libre de  $G$. Montrer que la famille  $\mathcal{F} \cup \mathcal{G} $  est libre.
    \item Soit  $\varphi $  une application lin\' eaire
de  $E$  dans  ${ \Rr}^q$, $q\in { \Nn}$. Construire deux applications
lin\' eaires  $\psi$  et  $\psi '$  de  $ E$  dans  ${ \Rr}^q$  telles que :  $\forall
y \in F   :   \psi '(y)=0  ,    \forall
z \in G   :   \psi (z)=0$  et  $\forall x \in E   :   \varphi (x)=\psi (x)+\psi ' (x)$.
\end{enumerate}
\end{enumerate}
\finenonce{000924}



\finexercice



\exercice{925, legall, 1998/09/01}

\enonce{000925}{Caract\'erisation de la somme directe de trois s.e.v.}
Soient $U, V, W$ des s.e.v. d'un e.v. $E$, v\'erifiant $(I) :   U\cap V=\{0\}=(U+V)\cap W$.
\begin{enumerate}
    \item D\'emontrer que $V\cap W=\{0\}=U\cap(V+W)$.\\
     \item Montrer que $(I)$ \'equivaut \`a $$(II) :  (\forall x \in U+V+W)    (\exists ! (u,v,w) \in U\times V \times W)   (x=u+v+w).$$
\end{enumerate}
\finenonce{000925}



\finexercice



\exercice{926, gourio, 2001/09/01}
\video{uLIAnT8Tc2w}
\enonce{000926}{}
Soit $$E=\big\{(u_{n})_{n\in
  \Nn}\in \Rr^{\Nn}\ |\ (u_{n})_{n} \text{ converge }\big\}.$$
Montrer que
l'ensemble des suites constantes et l'ensemble des suites convergeant
vers $0$ sont des sous-espaces suppl\'{e}mentaires dans $E.$
\finenonce{000926} 


\finexercice


\exercice{2433, matexo1, 2002/02/01}
\video{F1PmEC1haKY}
\enonce{002433}{}
Soit $f$ l'endomorphisme de $\R^3$ dont la matrice par
rapport \`a la base canonique $(e_1, e_2, e_3)$ est
$$A= \left( 
       \begin{array}{ccc}
        15 & -11 & 5 \\
        20 & -15 & 8 \\
        8 &  -7 & 6
       \end{array}
       \right).$$
Montrer que les vecteurs
$$ e'_1 = 2e_1+3e_2+e_3,\quad e'_2 = 3e_1+4e_2+e_3,\quad e'_3 =
e_1+2e_2+2e_3$$
forment une base de $\R^3$ et calculer la matrice de $f$ par
rapport \`a cette base.
\finenonce{002433}


\finexercice



\exercice{3305, quercia, 2010/03/09}
\enonce{003305}{Supplémentaire commun, X MP$^*$ 2005}

\begin{enumerate}
  \item Soit $A = \{P\in\R[X]$ tq $P=(1-X)Q(X^2)$ avec $Q\in\R[X]\}$.
  \begin{enumerate}
    \item Montrer que $A$ est un $\R$-ev et que l'on a
    $R[X] = A \oplus \{$ polynômes pairs $\}$.

    A-t-on $R[X] = A \oplus \{$ polynômes impairs $\}$~?
    \item Que peut-on dire si l'on remplace $Q(X^2)$ par une fonction $f$ paire~?
  \end{enumerate}
  \item Soient $E_1$, $E_2$ deux sev d'un ev~$E$ tels que $E_1$ et $E_2$ sont isomorphes
et $E = E_1 \oplus E_2$. Montrer que $E_1$ et $E_2$ ont un supplémentaire commun.
\end{enumerate}
\finenonce{003305}



\finexercice


\exercice{3652, quercia, 2010/03/10}
\enonce{003652}{}
Soit $E =  K_3[X]$,
$F = \{ P \in E$ tq $P(0) = P(1) = P(2) = 0 \}$,
$G = \{ P \in E$ tq $P(1) = P(2) = P(3) = 0 \}$,
et $H = \{ P \in E$ tq $P(X) = P(-X) \}$.

\begin{enumerate}
  \item Montrer que $F \oplus G = \{ P \in E$ tq $P(1) = P(2) = 0 \}$.
  \item Montrer que $F \oplus G \oplus H = E$.
\end{enumerate}
\finenonce{003652}



\finexercice


\exercice{3653, quercia, 2010/03/10}
\enonce{003653}{Caractérisation des sommes directes}

Soient $F_1$, $F_2$, $F_3$ trois sev de $E$.
Montrer que $F_1+F_2+F_3$ est directe si et seulement si :
$F_1 \cap F_2 = \{\vec 0\}$ et $(F_1 + F_2) \cap F_3 = \{\vec 0\}$.

Généraliser.

\finenonce{003653}



\finexercice


\exercice{3654, quercia, 2010/03/10}
\enonce{003654}{Somme directe dans $E$ $ \Rightarrow $ somme directe dans $\mathcal{L}(E)$}

Soit $E$ un $ K$-ev de dimension finie $n$ et
${\cal B} = ({\vec e}_1,\dots,{\vec e}_n)$ une base de $E$.

On note $F_i = \{ u \in \mathcal{L}(E)$ tq $\Im u \subset \text{vect}({\vec e}_i)\}$.

\begin{enumerate}
  \item Caractériser matriciellement les éléments de $F_i$.
  \item Montrer que $F_1 \oplus F_2 \oplus \dots \oplus F_n = \mathcal{L}(E)$.

\end{enumerate}
\finenonce{003654}



\finexercice


\exercice{3655, quercia, 2010/03/10}
\enonce{003655}{Toute somme peut être rendue directe en réduisant les sev}

Soit $E$ un $ K$-ev de dimension finie, $F_1$, $F_2$, \dots, $F_n$
des sev de $E$ tels que $F_1 + \dots + F_n = E$.
Montrer qu'il existe des sev $G_1 \subset F_1$, \dots, $G_n \subset F_n$
tels que $G_1 \oplus G_2 \oplus \dots \oplus G_n = E$.
\finenonce{003655}



\finexercice


\exercice{3656, quercia, 2010/03/10}
\enonce{003656}{Somme et intersection}

Soit $E$ un $ K$-ev, $E_1,\dots,E_n$ des sev tels que
$E_1 \oplus \dots \oplus E_n = E$, $F$ un autre sev de $E$, et
$F_i = E_i \cap F$.

\begin{enumerate}
  \item Montrer que la somme $G = F_1 + \dots + F_n$ est directe.
  \item Comparer $F$ et $G$.

\end{enumerate}
\finenonce{003656}



\finexercice


\exercice{3657, quercia, 2010/03/10}
\enonce{003657}{Somme directe d'endomorphismes}

Soit $E$ un $ K$-ev, $E_1,\dots,E_n$ des sev tels que
$E_1 \oplus \dots \oplus E_n = E$.
Soient $u_1 \in \mathcal{L}({E_1}),\dots,u_n \in \mathcal{L}({E_n})$.

\begin{enumerate}
  \item Montrer qu'il existe un unique endomorphisme $u \in \mathcal{L}(E)$ tel que pour tout
    $i$ : $u_{|F_i} = u_i$.
  \item Montrer que $\mathrm{Ker}(u) = \mathrm{Ker}(u_1) \oplus \dots \oplus \mathrm{Ker}(u_n)$ et
        $\Im (u) = \Im (u_1) \oplus \dots \oplus \Im (u_n)$.

\end{enumerate}
\finenonce{003657}



\finexercice


\exercice{3658, quercia, 2010/03/10}
\enonce{003658}{Somme de projecteurs}

Soit $E$ un $ K$-ev de dimension finie et $p_1,\dots,p_n$ des projecteurs
tels que $p_1 + \dots + p_n = \mathrm{id}_E$.

\begin{enumerate}
  \item Montrer que $\mathrm{tr}(p_i) = \mathrm{rg}(p_i)$.
  \item Montrer que $E = \Im(p_1) \oplus \dots \oplus \Im(p_n)$.
\end{enumerate}
\finenonce{003658}



\finexercice


\exercice{3659, quercia, 2010/03/10}
\enonce{003659}{Projecteurs}

Soit $E$ un espace vectoriel de dimension $n$, et $f_1,\ldots,f_n$ $n$ applications linéaires
toutes non nulles.
On suppose que~: $\forall\ (i,j)\in{[[1,n]]^2},\ f_i\circ f_j=\delta_{i,j}f_i$.
Montrer les $f_i$ sont toutes de rang un.

\finenonce{003659}



\finexercice

\exercice{5178, rouget, 2010/06/30}
\enonce{005178}{**IT}
Soient $u=(1,1,...,1)$ et $F=\mbox{Vect}(u)$ puis $G=\{(x_1,...,x_n)\in\Rr^n/\;x_1+...+x_n=0\}$. Montrer que $G$ est
un sous-espace vectoriel de $\Rr^n$ et que $\Rr^n=F\oplus G$.
\finenonce{005178}


\finexercice

\exercice{5185, rouget, 2010/06/30}
\enonce{005185}{**}
Soit $\Kk$ un sous-corps de $\Cc$ et $E$ un $\Kk$-espace vectoriel de dimension finie. Soient $f$ et $g$ deux endomorphismes de $E$ vérifiant
$E=\mbox{Ker}f+\mbox{Ker}g=\mbox{Im}f+\mbox{Im}g$. Montrer que ces sommes sont directes.
\finenonce{005185}


\finexercice

\exercice{5565, rouget, 2010/10/16}
\enonce{005565}{** I}
$E=\Kk^n$ où $\Kk$ est un sous-corps de $\Cc$.

Soient $F=\{(x_1,...,x_n)\in  E/\;x_1+...+x_n=0\}$ et $G=\text{Vect}\left((1,...,1)\right)$. Montrer que $F$ est un sous-espace vectoriel de $E$. Montrer que $F$ et $G$ sont supplémentaires dans $E$. Préciser le projeté d'un vecteur $x$ de $E$ sur $F$ parallèlement à $G$ et sur $G$ parallèlement à $F$.
\finenonce{005565}


\finexercice


\exercice{6871, chataur, 2012/05/13}
\video{u7JkO7A-9ZA}
\enonce{006871}{}
Par des considérations géométriques répondez aux questions suivantes :
\begin{enumerate}
  \item Deux droites vectorielles de $\mathbb{R}^3$ sont-elles suppl\'ementaires ?
  \item Deux plans vectoriels de $\mathbb{R}^3$ sont-ils suppl\'ementaires ?
  \item A quelle condition un plan vectoriel et une droite vectorielle de $\mathbb{R}^3$ sont-ils supplémentaires ?
\end{enumerate}
\finenonce{006871}




\finexercice


\exercice{7414, mourougane, 2021/08/10}
\enonce{007414}{}
Soit $E$ l'espace vectoriel des fonctions de classe $\mathcal C^1$ de $\Rr$ dans $\Rr$. Soit $F$ l'ensemble de toutes les fonctions de $\Rr$ dans $\Rr$ qui s'\'{e}crivent $x\mapsto ax+b$, pour certains r\'{e}els $a$ et $b$.
Soit enfin $G=\{f\in E | f(0)=0, f'(0)=0\}$.
\begin{enumerate}
	\item Montrer que $F$ et $G$ sont des sous-espaces vectoriels de $E$.
	\item Montrer que $$E=F\oplus G.$$
\end{enumerate}
\finenonce{007414}
\finexercice

\section{106.04 Base}


\exercice{979, legall, 1998/09/01}
\enonce{000979}{}
Montrer que les vecteurs  $\{ \begin{pmatrix}1 \\ 1 \\ 1 \\ \end{pmatrix} ,
 \begin{pmatrix} -1 \cr 1 \cr 0\cr \end{pmatrix}  ,  \begin{pmatrix}1 \cr 0 \cr -1 \cr
\end{pmatrix} \}$
forment une base
de   ${\Rr}^3$. Calculer les coordonn\' ees respectives des vecteurs  $
\begin{pmatrix}1 \cr 0 \cr 0 \cr\end{pmatrix} ,
 \begin{pmatrix}1 \cr 0 \cr 1\cr \end{pmatrix}  ,  \begin{pmatrix}0 \cr 0 \cr 1 \cr \end{pmatrix}$  dans cette base.
\finenonce{000979} 


\finexercice


\exercice{980, liousse, 2003/10/01}

\enonce{000980}{}
Soient $\vec{v_1}(1,2,3,4), \vec{v_2}(2,2,2,6), 
\vec{v_3}(0,2,4,4), \vec{v_4}(1,0,-1,2), 
\vec{v_5}(2,3,0,1)$ dans $\Rr^4$. 
Soient $F=Vect\{\vec{v_1}, \vec{v_2}, \vec{v_3}\}$ 
et $G=Vect\{\vec{v_4}, \vec{v_5} \}$. D\'eterminer une base des 
sous-espaces 
$F\cap G,$ $F,$ $G$ et $F+G.$
\finenonce{000980}



\finexercice



\exercice{981, liousse, 2003/10/01}
\video{vlr-D0jhb-Q}
\enonce{000981}{}
\begin{enumerate}
\item Montrer que les vecteurs 
$v_1=(0,1,1)$, $v_2=(1,0,1)$ et $v_3=(1,1,0)$ forment une base de 
$\Rr^3$. Trouver les composantes du vecteur $w=(1,1,1)$ dans cette base $(v_1,v_2,v_3)$.

\item Montrer que les vecteurs 
$v_1=(1,1,1)$, $v_2=(-1,1,0)$ et $v_3=(1,0,-1)$ forment une base de 
$\Rr^3$. Trouver les composantes du vecteur $e_1=(1,0,0)$,
$e_2=(0,1,0)$, $e_3=(0,0,1)$ et $w=(1,2,-3)$ dans cette base $(v_1,v_2,v_3)$.

\item Dans $\Rr^3$, donner un exemple de famille libre qui n'est pas
g\'en\'eratrice.

\item Dans $\Rr^3$, donner un exemple de famille g\'en\'eratrice qui n'est pas libre.
\end{enumerate}
\finenonce{000981} 


\finexercice


\exercice{982, cousquer, 2003/10/01}

\enonce{000982}{}
 On considère dans $\mathbb{R}^4$, $F=\mbox{lin}\{a,b,c\}$ et $G=\mbox{lin}\{d,e\}$,
avec $a=(1,2,3,4)$, $b=(2,2,2,6)$, $c=(0,2,4,4)$, $d=(1,0,-1,2)$ et
$e=(2,3,0,1)$. Déterminer des bases des sous-espaces $F\cap G$, $F$,
$G$, $F+G$.
\finenonce{000982}



\finexercice



\exercice{983, cousquer, 2003/10/01}

\enonce{000983}{}
 Dans l'espace $\mathcal{P}_5$ des polynômes de degré $\leq 5$, on définit 
les sous-ensembles~:\newline
$E_1=\{P\in\mathcal{P}_5 \mid P(0)=0\}$\newline
$E_2=\{P\in\mathcal{P}_5 \mid P'(1)=0\}$\newline
$E_3=\{P\in\mathcal{P}_5 \mid x^2+1 \mbox{ divise } P\}$\newline
$E_4=\{P\in\mathcal{P}_5 \mid x\mapsto P(x) 
    \mbox{ est une fonction paire}\}$\newline
$E_5=\{P\in\mathcal{P}_5 \mid \forall x,\; P(x)=xP'(x)\}$.
\begin{enumerate}
\item  Déterminer des bases des sous-espaces vectoriels $E_1$, $E_2$, $E_3$,
$E_4$, $E_5$, $E_1\cap E_2$, $E_1\cap E_3$, $E_1\cap E_2\cap E_3$,
$E_1\cap E_2\cap E_3\cap E_4$.
\item  Déterminer dans $\mathcal{P}_5$ des sous-espaces supplémentaires de 
$E_4$ et de $E_1\cap E_3$. 
\end{enumerate}
\finenonce{000983}



\finexercice



\exercice{984, cousquer, 2003/10/01}

\enonce{000984}{}
 Dans $\mathbb{R}^4$ on considère l'ensemble $E$ des vecteurs 
$(x_1,x_2,x_3,x_4)$ vérifiant l'équation $x_1+x_2+x_3+x_4 =0$. L'ensemble $E$ est-il
un sous-espace vectoriel de $\mathbb{R}^4$~? Si oui, en donner une base.
\finenonce{000984}



\finexercice



\exercice{985, cousquer, 2003/10/01}
\enonce{000985}{}
 Vrai ou faux ?
On d\'esigne par $E$ un $\mathbb{R}$-espace vectoriel de dimension finie.
\begin{enumerate}
\item  Si les vecteurs $x,y,z$ sont deux \`a deux non colin\'eaires, alors la
famille $x,y,z$ est libre.
\item  Soit $x_1,x_2, \dots,x_p$ une famille de vecteurs. Si aucun n'est une
combinaison lin\'eaire des autres, la famille est libre.
\end{enumerate}
\finenonce{000985} 


\finexercice


\exercice{986, cousquer, 2003/10/01}

\enonce{000986}{}
 Étudier l'indépendance linéaire des listes de vecteurs suivantes,
et trouver à chaque fois une base du sous-espace engendré.
\begin{enumerate}
\item  $(1,0,1)$, $(0,2,2)$, $(3,7,1)$ dans $\mathbb{R}^3$.
\item  $(1,0,0)$, $(0,1,1)$, $(1,1,1)$ dans $\mathbb{R}^3$.
\item  $(1,2,1,2,1)$, $(2,1,2,1,2)$, $(1,0,1,1,0)$, $(0,1,0,0,1)$ dans $\mathbb{R}^5$.
\item  $(2,4,3,-1,-2,1)$, $(1,1,2,1,3,1)$, $(0,-1,0,3,6,2)$ dans $\mathbb{R}^6$.
\item  $(2,1,3,-1,4,-1)$, $(-1,1,-2,2,-3,3)$, $(1,5,0,4,-1,7)$ dans $\mathbb{R}^6$.
\end{enumerate}
\finenonce{000986}



\finexercice



\exercice{987, cousquer, 2003/10/01}
\enonce{000987}{}
 Dans $\mathbb{R}^3$, les vecteurs suivants forment-ils une base~? Sinon
d\'ecrire le sous-espace qu'ils engendrent.
\begin{enumerate}

\item $v_1 =(1,1,1), v_2=(3,0,-1),v_3=(-1,1,-1).$
\item $v_1 =(1,2,3), v_2=(3,0,-1),v_3=(1,8,13).$
\item $v_1 =(1,2,-3), v_2=(1,0,-1),v_3=(1,10,-11).$

\end{enumerate}

\finenonce{000987} 


\finexercice


\exercice{988, cousquer, 2003/10/01}

\enonce{000988}{}
 Dans $\mathbb{R}^3$, comparer les sous-espaces $F$ et $G$ suivants~:\\
$F=\mbox{lin}\{(2,3,-1), (1,-1,-2)\}$ et $G=\mbox{lin}\{(3,7,0), (5,0,-7)\}$.
\finenonce{000988}



\finexercice



\exercice{989, cousquer, 2003/10/01}

\enonce{000989}{}
 Dans $\mathbb{R}^4$, on considère les familles de vecteurs suivantes
\\ $v_1=(1,1,1,1)$, $v_2=(0,1,2,-1)$, $v_3=(1,0,-2,3)$, 
$v_4=(2,1,0,-1)$, $v_5=(4,3,2,1)$.
\\  $v_1=(1,2,3,4)$, $v_2=(0,1,2,-1)$, $v_3=(3,4,5,16)$.
\\ $v_1=(1,2,3,4)$, $v_2=(0,1,2,-1)$, $v_3=(2,1,0,11)$, $v_4=(3,4,5,14)$.\\
Ces vecteurs  forment-ils~:
\begin{enumerate}
    \item Une famille libre~? Si oui, la compléter pour obtenir une base de 
    $\mathbb{R}^4$.
    Si non donner des relations de dépendance entre eux et extraire de cette
    famille au moins une famille libre. 
    \item Une famille génératrice~? Si oui, en
    extraire au moins une base de l'espace. Si non, donner la dimension du
    sous-espace qu'ils engendrent.
\end{enumerate}
\finenonce{000989}



\finexercice



\exercice{990, cousquer, 2003/10/01}

\enonce{000990}{}
 Si $E$ est un espace vectoriel de dimension finie, $F$ et $G$ deux
sous-espaces de $E$, montrer que $F\cup G$ est
un sous-espace vectoriel si et seulement si $F\subset G$ ou $G\subset F$.
\finenonce{000990}



\finexercice



\exercice{991, cousquer, 2003/10/01}

\enonce{000991}{}
 On désigne par $E$ un $\mathbb{R}$-espace vectoriel de dimension finie. Les
propriétés suivantes sont-elles vraies ou fausses~?
\begin{enumerate}
\item  Soient $D_1,D_2,D_3$ des droites vectorielles de $\mathbb{R}^3$ distinctes 
deux à deux. Alors $\mathbb{R}^3$ est somme de $D_1,D_2,D_3$.
\item  Soient $F$ et $G$ des hyperplans vectoriels de $E$. Alors $E \neq F \cup G$.
\item  Soient $P_1$ et $P_2$ des plans vectoriels de $E$ tels que $P_1 \cap P_2 =
\{0\}$. Alors $\dim E \ge 4$.
\item  Soient $F$ et $G$ des sous-espaces de dimension~$3$ de $\mathbb{R}^5$.
Alors $F \cap G \ne\{0\}$.
\item  Soit $(e_1,e_2,e_3,e_4)$ la base canonique de $\mathbb{R}^4$ et
$F=\mbox{lin}\{e_1,e_3\}$. Tout sous-espace vectoriel supplémentaire de $F$
contient $e_2$.
\end{enumerate}
\finenonce{000991}



\finexercice



\exercice{992, cousquer, 2003/10/01}
\video{FZf39lWtpZw}
\enonce{000992}{}
\begin{enumerate}
  \item Soit $E=\Rr_n[X]$ l'espace vectoriel des polynômes de degré inférieur ou égal à $n$.
Montrer que toute famille de polyn\^omes $\{P_0,P_1,\ldots,P_n\}$ avec $\deg P_i = i$ 
(pour $i=0,1,\ldots,n$) forme une base de $E$.

\item  \'Ecrire le polyn\^ome $F=3X-X^2+8X^3$ sous la forme $F=a+b(1-X)+c(X-X^2)+d(X^2-X^3)$
($a,b,c,d \in \Rr$) puis sous la forme
$F=\alpha+\beta(1+X)+\gamma(1+X+X^2)+\delta(1+X+X^2+X^3)$
($\alpha,\beta,\gamma,\delta \in \Rr$).

\end{enumerate}
\finenonce{000992} 


\finexercice


\exercice{993, cousquer, 2003/10/01}

\enonce{000993}{}
 Dans l'espace vectoriel $\mathcal{P}_2$ des polynômes de degré 
$\leq 2$, on considère les polynômes
$P_1=X^2+X(1-X)+(1-X)^2$, $P_2=X^2+(1-X)^2$, $P_3=X^2+1+(1-X)^2$,
$P_4=X(1-X)$.
Peut-on extraire de $\{P_1,P_2,P_3,P_4\}$ des bases de $\mathcal{P}_2$~?
Si oui, les trouver toutes.
\finenonce{000993}



\finexercice



\exercice{994, cousquer, 2003/10/01}

\enonce{000994}{}
 Soit $E$ l'ensemble des fractions rationnelles $F$ qui peuvent s'écrire
$$ F=\frac{P}{(X-1)^3(X^2+1)^2},
\qquad P \mbox{ polynôme de degré }\leq 6.$$
Les fractions $\frac{1}{(X-1)}$,
$\frac{1}{(X-1)^2}$, $\frac{1}{(X-1)^3}$,
$\frac{1}{X^2+1}$, $\frac{X}{X^2+1}$,
$\frac{1}{(X^2+1)^2}$, $\frac{X}{(X^2+1)^2}$
forment-elles une base de $E$~?\\
Que se passe-t-il si on suppose que $P$ décrit l'ensemble des
polynômes de degré $\leq 9$~?
\finenonce{000994}



\finexercice



\exercice{995, cousquer, 2003/10/01}

\enonce{000995}{}
 Problème de l'interpolation~: soit les cinq points
$(x_1,y_1)=(-2,3)$, $(x_2,y_2)=(0,-2)$,
$(x_3,y_3)=(1,5)$, $(x_4,y_4)=(5,1)$, $(x_5,y_5)=(6,7)$
de $\mathbb{R}^2$, et $\mathcal{P}_4$ l'espace vectoriel des polynômes
de degré $\leq 4$. On veut trouver un polynôme~$F$
dans $\mathcal{P}_4$ tel que pour $i=1,\ldots,5$ on ait
$F(x_i)=y_i$.
\begin{enumerate}
\item  Sans effectuer les calculs, indiquer
comment on pourrait calculer $a,b,c,d,e$ exprimant
$F=a+bX+cX^2+dX^3+eX^4$ selon la base 
$\{1,X,X^2,X^3,X^4\}$ de $\mathcal{P}_4$.

\item  Montrer que $\{1,X+2,(X+2)X,(X+2)X(X-1),
(X+2)X(X-1)(X-5)\}$ est une base de $\mathcal{P}_4$. Calculer directement
(indépendamment de la question précédente) les coordonnées de $F$
dans cette base.

\item  Montrer que l'ensemble des polynômes 
$X(X-1)(X-5)(X-6),(X+2)(X-1)(X-5)(X-6),
(X+2)X(X-5)(X-6),
(X+2)X(X-1)(X-6),(X+2)X(X-1)(X-5)$
forment une base de $\mathcal{P}_4$. Calculer directement (indépendamment 
des questions précédentes) les coordonnées de~$F$ dans cette base.

\item  Dans laquelle des diverses bases ci-dessus le calcul de~$F$  vous
paraît-il le plus simple~?
\end{enumerate}
\finenonce{000995}



\finexercice



\exercice{996, legall, 1998/09/01}
\video{Pmhg9aR10aA}
\enonce{000996}{}
D\'eterminer pour quelles valeurs de $ t\in {\Rr} $ les
vecteurs  
$$\big\{(1, 0, t),  (1, 1, t), (t,0,1)\big\}$$
forment une base de  $\Rr^3$.
\finenonce{000996} 


\finexercice


\exercice{997, legall, 1998/09/01}

\enonce{000997}{}
Soit  $(\Sigma )$  le syst\`eme d'\' equations lin\' eaires :
$$\left\lbrace \begin{array}{l}
x+3y+2z= 0 \\
x+y+z+t= 0 \\
x-t= 0
\end{array}\right.$$


\noindent Montrer que l'ensemble des solutions de  $(\Sigma )$  forme
un sous-espace vectoriel  $F$  de  ${\Rr}^4$. D\' eterminer la dimension
et une base de  $F$.
\finenonce{000997}



\finexercice



\exercice{998, legall, 1998/09/01}

\enonce{000998}{}
Soit $ a\in \R .$ On pose, pour tout $ p\in \Nn
: A_p(X)=(X-a)^p $
et $ B_p(X)=X^p .$
\begin{enumerate}
    \item  Montrer que $\epsilon = \{ A_0, \ldots, A_n\}  $ est une base de
$ \R _n[X] .$
    \item Soit $ P\in \R _n[X] .$ Montrer que $ \displaystyle{ P(X)= \sum _{k=0}^n
\frac{1}{ k!}P^{(k)}(a)A_k(X)} .$ (On pourra montrer que l'ensemble $ E $
des \'el\'ement de $ \R _n[X] $ qui satisfont \`a cette \'egalit\'e
est un sous-espace vectoriel de $ \R _n[X] $ et contient une base.)
\end{enumerate}
\finenonce{000998}



\finexercice



\exercice{999, legall, 1998/09/01}

\enonce{000999}{}
On munit $E={ \Rr}^{*}_{+}\times { \Rr}$ de la loi interne
``addition'' + : $(a,b)+(a',b')=(aa',b+b')$, et de la loi externe
. \`a coefficients r\'eels : $(\forall \lambda \in { \Rr}) \forall
(a,b) \in E   \lambda.(a,b)=(a^{\lambda},\lambda b)$.
\begin{enumerate}
    \item V\'erifier que $(E,+,.)$ est un ${ \Rr}$-e.v.
    \item Les syst\`emes suivants sont-ils libres ou li\'es : ((1,0),(1,1)) ? ((2,1),(8,3)) ?
((2,1),(6,3)) ?
    \item V\'erifier que le syst\`eme $b=((2,0),(2,1))$ est une base de $E$ et d\'eterminer les
composantes du vecteur $v=(x,y) \in E$ par rapport \`a la base $b$.
\end{enumerate}
\finenonce{000999}



\finexercice



\exercice{1000, legall, 1998/09/01}

\enonce{001000}{}
Pour $k=2,3,4$ montrer que $V_{k}$ est un s.e.v. de ${ \Cc}^{k}$, et en donner
une base :
$$V_{2}=\{(a,b)\in { \Cc}^{2}/a+ib=0\},   V_{3}=\{(a,b,c)\in { \Cc}^{3} /a+2b+3c=0\},$$
$$V_{4}=\{(a,b,c,d)\in { \Cc}^{4} /a+ib=b+ic=c+id\}.$$
\finenonce{001000}



\finexercice



\exercice{1001, legall, 1998/09/01}

\enonce{001001}{}
Soit $n \in { \Nn}$ et $E={ \Rr}_{n}[X]$, l'espace vectoriel des polyn\^omes \`a
coefficients r\'eels, de degr\'e $\le n$.
\begin{enumerate}
    \item  Soit $\beta=(P_{0},P_{1},...,P_{n})$ un syst\`eme de $(n+1)$ polyn\^omes tels que, $\forall k$,
$0\le k \le n$, $\text{deg}\,P_{k}=k$. Montrer que $\beta$ est une base de $E$.
    \item Soit $P$ un polyn\^ome de degr\'e $n$. Montrer que : $\gamma=(P,P',\ldots,P^{(n)})$ est une base
de $E$ et d\'eterminer les composantes du polyn\^ome $Q$ d\'efini par : $Q(X)=P(X+a)$, ($a$ r\'eel
fix\'e), dans la base $\gamma$.
    \item D\'emontrer que le syst\`eme $S=(X^{k}(1-X)^{n-k})_{0\le k \le n}$ est une base de $E$, et
d\'eterminer, pour tout $p \in \{0,1,\ldots,n\}$, les composantes du polyn\^ome $X^{p}$ dans la base $S$.
\end{enumerate}
\finenonce{001001}



\finexercice



\exercice{1002, legall, 1998/09/01}

\enonce{001002}{}
Soient $\mathbf{v}_1 = (1,0,0,-1), \mathbf{v}_2 = (2,1, 0,1), \mathbf{v}_3 = (1,-1,1,-1), \mathbf{v}_4=(7,2,0,-1)$. Donner une base du sous-espace vectoriel $F =<\mathbf{v}_1,\mathbf{v}_2,\mathbf{v}_3,\mathbf{v}_4>$.
D\'eterminer un suppl\'ementaire $G$ de $F$ dans $\R^4$.
\finenonce{001002}



\finexercice



\exercice{1003, legall, 1998/09/01}

\enonce{001003}{}
Soient le triplet $ {\mathbf v}_1 =(1,2,3,0), {\mathbf v}_2
=(-1,1,2,1), {\mathbf v}_3 =(1,5,8,1)$ et le triplet
$ {\mathbf w}_1 =(0,3,5,1), {\mathbf w}_2
=(1,-1,1,0), {\mathbf w}_3 =(0,0,3,1)$. On consid\`ere les sous-espaces vectoriels $F= < {\mathbf v}_1
,{\mathbf v}_2,{\mathbf v}_3>$ et
$G= < {\mathbf w}_1 ,{\mathbf w}_2,{\mathbf w}_3>$. Donner une base des sous-espaces suivants
$F,G, F\cap G$ et $F+G$.
\finenonce{001003}



\finexercice



\exercice{1004, legall, 1998/09/01}

\enonce{001004}{}
Soit
$$ E =\left\{ f_{\alpha,A} \in {\mathcal F}(\R,\R);  (\alpha,A)\in \R^2,\quad f_{\alpha,A}(x) =
A\cos(x+\alpha)\right\}.
$$
Montrer que $E$ est un sous-espace vectoriel de    ${\mathcal F}(\R,\R)$ et en donner une base.
\finenonce{001004}



\finexercice



\exercice{1005, legall, 1998/09/01}

\enonce{001005}{}
Soit $E =\R^3$. On d\'efinit le syst\`eme $$S=\left\{ {\mathbf e}_1 =(1,1,1), {\mathbf e}_2
=(1,1,2), {\mathbf e}_3 =(1,2,3)\right\}$$
\begin{enumerate}
    \item Montrer que $S$ est une base de $E$.
    \item Calculer les coordonn\'ees de ${\mathbf v} = (5,7,12)$ dans cette base.
\end{enumerate}
\finenonce{001005}



\finexercice



\exercice{1006, legall, 1998/09/01}
\video{G-Xa2ZmQUwQ}
\enonce{001006}{}
\begin{enumerate}
    \item Montrer que les vecteurs  $v_1 =(1,-1,i)$, $v_2=(-1,i,1)$, $v_3 =(i,1,-1)$ forment une base de $\Cc^3$.
    \item Calculer les coordonnées de $v = (1+i,1-i,i)$ dans cette base.
\end{enumerate}
\finenonce{001006} 


\finexercice


\exercice{1007, legall, 1998/09/01}

\enonce{001007}{}
\begin{enumerate}
\item Montrer que les syst\`emes ${\mathbf s}_1= (1,\sqrt2)$ et ${\mathbf s}_2= (1,\sqrt2,\sqrt3)$
sont libres dans $\R$ consid\'er\'e comme un espace vectoriel sur $\Qq$.
    \item Soient dans $\R^2$, les vecteurs ${\mathbf u}_1 = (3+\sqrt5, 2+3\sqrt5)$ et
${\mathbf u}_2 = (4, 7\sqrt5 -9)$. Montrer que le syst\`eme $({\mathbf u}_1 ,{\mathbf u}_2)$ est $\Qq$--libre et $\R$--li\'e.
    \item Soient dans $\C^2$, les vecteurs ${\mathbf r}_1 = (1+i, 1-2i)$
et
${\mathbf r}_2 = (3i-1, 5)$. Montrer que le syst\`eme $({\mathbf r}_1 ,{\mathbf r}_2)$ est $\Rr$--libre et
$\C$--li\'e.
\end{enumerate}
\finenonce{001007}



\finexercice



\exercice{1008, legall, 1998/09/01}

\enonce{001008}{}
D\' eterminer pour quelles valeurs
de  $t\in {\Rr}$  les polyn\^omes $ X^2+t/2\ , \ X-t \ ,\ (X+t+1)^2 $
 forment une base de  ${\Rr}_2[X]$.
\finenonce{001008}



\finexercice



\exercice{1009, ridde, 1999/11/01}

\enonce{001009}{}
Etudier la libert\'e des familles
\begin{enumerate}
\item $ (1, 1), (1, 2)$.
\item $ (2, 3), (-6, 9)$.
\item $ (1, 3, 1), (1, 3, 0), (0, 3, 1)$.
\item $ (1, 3), (-1, -2), (0, 1)$.
\end{enumerate}
\finenonce{001009}



\finexercice



\exercice{1010, ridde, 1999/11/01}

\enonce{001010}{}
Les familles suivantes sont-elles g\'en\'eratrices ?
\begin{enumerate}
\item $ (1, 1), (3, 1)$ dans $\Rr^2$.
\item $ (1, 0, 2), (1, 2, 1)$ dans $\Rr^3$.
\end{enumerate}
\finenonce{001010}



\finexercice



\exercice{1011, ridde, 1999/11/01}

\enonce{001011}{}
On consid\`ere dans $\Rr^3$,  $\Pi = \text{vect}\left\{ (1, 1, 1), (1, 1, -1)\right\}$
et $D = \text{vect} \left\{ (0, 1, -1)\right\}$. Montrer que $\Rr^3 = \Pi \oplus D$.
\finenonce{001011}



\finexercice



\exercice{1012, ridde, 1999/11/01}

\enonce{001012}{}
D\'eterminer une base de $\left\{ (x, y, z) \in \Rr^3 / x + y + z = 0\right\}$.
\finenonce{001012}



\finexercice



\exercice{1013, ridde, 1999/11/01}

\enonce{001013}{}
D\'eterminer une base de $D = \left\{ (x, y, z) \in \Rr^3 / x + y = 0, x-y + z = 0\right\}$.
\finenonce{001013}



\finexercice



\exercice{3317, quercia, 2010/03/09}
\enonce{003317}{Essai de bases}
Montrer que dans $\R^3$, les trois vecteurs $\vec a = (1,0,1)$,
$\vec b = (-1,-1,2)$ et $\vec c = (-2,1,-2)$ forment une base, et calculer
les coordonnées dans cette base d'un vecteur $\vec x = (x,y,z)$.
\finenonce{003317}



\finexercice


\exercice{3318, quercia, 2010/03/09}
\enonce{003318}{Rang de vecteurs}
Dans $\R^4$, trouver le rang de la famille de vecteurs :

$$\vec a = (3,2,1,0),\quad
  \vec b = (2,3,4,5),\quad
  \vec c = (0,1,2,3),\quad
  \vec d = (1,2,1,2),\quad
  \vec e = (0,-1,2,1).$$



\finenonce{003318}



\finexercice


\exercice{3319, quercia, 2010/03/09}
\enonce{003319}{Fonctions affines par morceaux}
Soit $0=x_0 < x_1 < \dots < x_n = 1$ une subdivision de $[0,1]$
et $F$ l'ensemble des fonctions $f : {[0,1]} \to \R$ continues dont
la restriction à chaque intervalle $[x_i,x_{i+1}]$ est affine.

Montrer que $F$ est de dimension finie et trouver une base de $F$.

\finenonce{003319}



\finexercice


\exercice{3320, quercia, 2010/03/09}
\enonce{003320}{Projection et symétrie dans $ K^3$}
Dans $ K^3$, on donne les sous espaces :
$\begin{cases} H = \{\smash{\overrightarrow X} = (x,y,z) \text{ tq } x + y + z = 0 \} \cr
         K = \text{vect}(\overrightarrow U = (1,1,2)).\hfill    \cr \end{cases}$

\begin{enumerate}
  \item Déterminer $\dim H$ et en donner une base.
  \item Démontrer que $H \oplus K =  K^3$.
  \item Donner les expressions analytiques des projection et symétrie associées :
    $\pi_H$ et $s_H$.
\end{enumerate}
\finenonce{003320}



\finexercice


\exercice{3321, quercia, 2010/03/09}
\enonce{003321}{Supplémentaires}
Soit $E = H \oplus K$ et $(\vec e_1,\dots,\vec e_k)$ une base de $K$.

\begin{enumerate}
  \item Montrer que pour tout $\vec a \in H,\
    K_{\vec a} = \text{vect}(\vec e_1 + \vec a, \dots, \vec e_k + \vec a)$
    est un supplémentaire de $H$.
  \item Montrer que si $\vec a \ne \vec b$, alors $K_{\vec a} \ne K_{\vec b}$.

\end{enumerate}
\finenonce{003321}



\finexercice

\exercice{5179, rouget, 2010/06/30}
\enonce{005179}{****}
\begin{enumerate}
\item  Soit $n$ un entier naturel. Montrer que si $n$ n'est pas un carré parfait alors $\sqrt{n}\notin\Qq$.
\item  Soit $E=\{a+b\sqrt{2}+c\sqrt{3}+d\sqrt{6},\;(a,b,c,d)\in\Qq^4\}$. Vérifier que $E$ est un $\Qq$-espace
vectoriel puis déterminer une base de $E$.
\end{enumerate}
\finenonce{005179}


\finexercice


\exercice{5572, rouget, 2010/10/16}
\enonce{005572}{**I}
$E=\Rr_n[X]$. Pour $0\leqslant k\leqslant n$, on pose $P_k=X^k(1-X)^{n-k}$. Montrer que la famille $(P_k)_{0\leqslant k\leqslant n}$ est une base de $E$.
\finenonce{005572}


\finexercice


\exercice{5573, rouget, 2010/10/16}
\enonce{005573}{**I Polynômes d'interpolation de \textsc{Lagrange}}
Soient $a_0$,..., $a_n$ $n+1$ nombres complexes deux à deux distincts et $b_0$,..., $b_n$ $n+1$  nombres complexes.

Montrer qu'il existe une unique famille de $n+1$ polynômes à coefficients complexes de degré $n$ exactement vérifiant $\forall(i,j)\in\llbracket0,n\rrbracket$,  $L_i(a_j)=1$ si $i=j$ et $0$ sinon. 

Montrer que la famille $(L_i)_{0\leqslant i\leqslant n}$ est une base de $\Cc_n[X]$.

Montrer qu'il existe un unique polynôme $P$ de degré inférieur ou égal à $n$ vérifiant $\forall i\in\llbracket0,n\rrbracket$, $P(a_i) =b_i$. Expliciter $P$ puis déterminer tous les polynômes vérifiant les égalités précédentes.
\finenonce{005573}


\finexercice

\section{106.05 Dimension}


\exercice{1014, liousse, 2003/10/01}

\enonce{001014}{}
Calculer la dimension du sous-espace vectoriel de $\mathbb R^4$ engendr\'e par les vecteurs
$ V_1=(0,1,2,3)$, $V_2=(1,2,3,4)$ et $V_3=(2,3,4,5)$.
\finenonce{001014}



\finexercice



\exercice{1015, cousquer, 2003/10/01}
\video{dSfb2GXKgeU}
\enonce{001015}{}
Soit $E$ est un espace vectoriel de dimension finie et $F$ et $G$ deux
sous-espaces vectoriels de $E$. Montrer que : 
$$\dim(F+G) = \dim F+\dim G - \dim(F\cap G).$$
\finenonce{001015} 


\finexercice


\exercice{1016, legall, 1998/09/01}
\video{QdZN5i_YucI}
\enonce{001016}{}
Montrer que tout sous-espace vectoriel d'un espace vectoriel de dimension finie
est de dimension finie.
\finenonce{001016} 


\finexercice


\exercice{1017, legall, 1998/09/01}

\enonce{001017}{}
Soient $P_0$, $P_1$, $P_2$ et $P_3\in\mathbb{R}_2[X]$ d\'efinis par
$$P_0(X)=\frac{(X-1)(X-2)}{2},\;\;P_1(X)=\frac{X(X-1)}{2},$$
$$P_2(X)=2X(X-2),\;\;P_3(X)=\frac{(X-1)(X-3)}{3}.$$
Exprimer $1,\;X,\;X^2$ en fonction de $P_0,\;P_1$ et $P_2$. On note
$F=Vect\{P_0,P_1\}$ et $G=Vect\{P_2,P_3\}$. Calculer $\dim F$, $\dim G$,
$\dim(F+G)$ et
$\dim(F\cap G)$. V\'erifier que
$$\dim(F+G)=\dim(F)+\dim(G)-\dim(F\cap G).$$
\finenonce{001017}



\finexercice



\exercice{1018, legall, 1998/09/01}

\enonce{001018}{}
Donner la dimension du sous-espace $F$ de ${\mathcal F}(\R,\R)$ engendr\'e par $f_1(x) =
\sin^2x, f_2(x) =\cos^2x, f_3(x) =\sin2x$ et $ f_4(x) =\cos2x$.
\finenonce{001018}



\finexercice



\exercice{1019, legall, 1998/09/01}
\video{mHpyqFvkVLI}
\enonce{001019}{}
On consid\`ere, dans  ${\Rr}^4$, les vecteurs : 
$$v_1=(1, 2, 3, 4),\quad v_2= (1, 1, 1, 3),\quad v_3= (2, 1, 1, 1),\quad v_4=(-1, 0, -1, 2),\quad v_5=(2, 3, 0, 1).$$

Soit  $F$  l'espace vectoriel engendr\'e par  $\{v_1, v_2  ,  v_3\} $  et soit $G$  celui engendr\'e par  $\{v_4, v_5\}$.
Calculer les dimensions respectives de  $F$, $G$, $F\cap G$, $F+G$.
\finenonce{001019} 


\finexercice


\exercice{1020, ridde, 1999/11/01}

\enonce{001020}{}
Soient $E = \left\{ (x, y, z, t) \in \Rr^4 / x + y + z + t = 0\right\}$
et $F = \left\{ (x, y, z, t) \in \Rr^4 / x + y  =  z + t \right\}$.
D\'eterminer $\dim E, \dim F, \dim (E + F), \dim (E \cap F)$.
\finenonce{001020}



\finexercice



\exercice{1021, ridde, 1999/11/01}

\enonce{001021}{}
Montrer que ${f} : {\Rr^3} \to {\Rr^3}$, ${(x, y, z)} \mapsto { (z, x-y, y + z)}$ est un automorphisme.
\finenonce{001021}



\finexercice



\exercice{1022, ridde, 1999/11/01}

\enonce{001022}{}
Soit $E$ un $\Qq$-espace vectoriel de dimension $n$. Montrer que
$$ n \text{ est pair} \Leftrightarrow \exists f \in \mathcal{L} (E) / \text{Im}f
 = \ker f$$
\finenonce{001022}



\finexercice



\exercice{1023, ridde, 1999/11/01}

\enonce{001023}{}
Montrer qu'il existe une unique forme lin\'eaire $f$ sur $\Rr^2$ telle que
$f (1, 2) = 2$ et $f (-2, 1) = 5$. D\'eterminer le noyau et l'image de $f$.
\finenonce{001023}



\finexercice



\exercice{1024, ridde, 1999/11/01}

\enonce{001024}{}
D\'eterminer suivant la valeur de $x \in \Rr$ le rang de la famille de vecteurs
$e_1 = (1, x, -1), e_2 = (x, 1, x), e_3 = (-1, x, 1)$.
\finenonce{001024}



\finexercice



\exercice{1025, ridde, 1999/11/01}

\enonce{001025}{}
Soit $E$ un espace vectoriel de dimension $3$ et $f \in \mathcal{L} (E)$ telle
que $f^2 \neq 0$ et $f^3 = 0$. Soit $x_0 \in E / f^2 (x_0) \neq 0$.
\begin{enumerate}
\item Montrer que $ (x_0, f (x_0), f^2 (x_0))$ est une base.
\item Montrer que l'ensemble des endomorphismes qui commutent avec $f$ est un
sous-espace vectoriel de $\mathcal{L} (E)$ de base $ (id, f, f^2)$.
\end{enumerate}
\finenonce{001025}



\finexercice



\exercice{1026, ridde, 1999/11/01}

\enonce{001026}{}
Soit $E$ de dimension finie et $f \in \mathcal{L} (E)$. Montrer l'\'equivalence
des trois propri\'et\'es :
\begin{enumerate}
\item[(i)] $\ker f = \ker f^2$.
\item[(ii)] $\text{Im}f = \text{Im}f^2$.
\item[(iii)] $E = \ker f \oplus \text{Im}f$.
\end{enumerate}
\finenonce{001026}



\finexercice



\exercice{1027, ridde, 1999/11/01}
\video{GU9xDUgfz48}
\enonce{001027}{}
Soit $E$ et $F$ de dimensions finies et $u, v \in \mathcal{L} (E, F)$.
\begin{enumerate}
\item Montrer que $\text{rg} (u + v) \leq \text{rg} (u) + \text{rg} (v)$.
\item En d\'eduire que $\left|\text{rg} (u) - \text{rg} (v)\right| \leq
\text{rg} (u + v)$.
\end{enumerate}

\finenonce{001027} 


\finexercice


\exercice{1028, gourio, 2001/09/01}

\enonce{001028}{}
Soit $(f,g)\in \left( L(E)\right) ^{2}$ o\`{u} $E$ est un ${K}$-espace
vectoriel de dimension finie $n$, montrer les in\'{e}galit\'{e}s :
$$\mathrm{rg}(f)+\mathrm{rg}(g)-n\leq \mathrm{rg}(f\circ g)\leq \inf (\mathrm{rg}(f),\mathrm{rg}(g)) $$
 (on pourra utiliser $g_{|\ker (f \circ g)} = h$ dont on d\'eterminera le noyau)
\finenonce{001028}



\finexercice



\exercice{1029, gourio, 2001/09/01}

\enonce{001029}{}
Soit $(f,g)\in \left( L(E)\right) ^{2}$ o\`{u} $E$ est un ${K}$-espace
vectoriel de dimension finie $n$, tel que : $(f+g)$ est inversible et $fg=0$.
Montrer que : $$rg(f)+rg(g)=n. $$
\finenonce{001029}



\finexercice



\exercice{1030, gourio, 2001/09/01}

\enonce{001030}{}
Soit $U$ un sous-espace vectoriel de $E$ espace vectoriel, et
$$A=\{f\in L(E)|U\subset Ker(f)\}. $$
Montrer que $A$ est un sous-espace vectoriel de $L(E).$ Si $E$ est de
dimension finie, quelle est la dimension de $A$ ?
\finenonce{001030}



\finexercice



\exercice{1031, gourio, 2001/09/01}

\enonce{001031}{}
Soient $E_{0},E_{1},...,E_{n}$ $n+1$ espaces vectoriels sur un m\^{e}me
corps commutatif ${K}$, de dimensions respectives $\alpha _{0},\alpha
_{1},...,\alpha _{n}.$ On suppose qu'il existe $n$ applications
lin\'{e}aires $f_{0},f_{1},...,f_{n-1}$ telles que :
$$\forall k\in \{0,...,n-1\},f_{k}\in L(E_{k},E_{k+1}). $$
et de plus :
\begin{itemize}
\item  $f_{0}$ est injective;
\item $\forall j\in \{1,...,n-1\},\mathrm{Im} f_{j-1}=Ker(f_{j});$
\item $f_{n-1}$ est surjective.
\end{itemize}
Montrer que
$$\sum\limits_{j=0}^{n}(-1)^{j}\alpha _{j}=0. $$
\finenonce{001031}



\finexercice



\exercice{1032, gourio, 2001/09/01}

\enonce{001032}{}
Soient $H_{1}$ et $H_{2}$ deux hyperplans de $E,$ espace vectoriel de
dimension $n.$
Montrer que :
$$\dim (H_{1}\cap H_{2})\geq n-2. $$
G\'{e}n\'{e}raliser.
\finenonce{001032}



\finexercice



\exercice{1033, gourio, 2001/09/01}

\enonce{001033}{}
Donner un exemple d'endomorphisme d'un espace vectoriel injectif et non
surjectif, puis d'un endomorphisme surjectif et non injectif.
\finenonce{001033}



\finexercice



\exercice{1034, gourio, 2001/09/01}

\enonce{001034}{}
Soit $E $ un espace vectoriel de dimension finie et $f\in L(E)$, montrer
l'\'{e}quivalence :
$$E=\mathrm{Ker}(f)\oplus \mathrm{Im} (f)\Leftrightarrow \mathrm{Im} f=\mathrm{Im} f^{2}.$$

Donner un contre-exemple quand $\dim E=+\infty .$
\finenonce{001034}



\finexercice



\exercice{1035, gourio, 2001/09/01}

\enonce{001035}{}
Soit $(f,g)\in L(E,F)^{2}$ avec $E,F$ de dimension finie. On suppose
$$\mathrm{rg}(f+g)=\mathrm{rg}(f)+\mathrm{rg}(g). $$
 Montrer que :
$$E=\mathrm{Ker}(f)+\mathrm{Im} f ; $$
$$\mathrm{Im} f\cap \mathrm{Im} g=\{0\}. $$
\finenonce{001035}



\finexercice



\exercice{1036, gourio, 2001/09/01}

\enonce{001036}{}
Soit $E$ un espace vectoriel de dimension finie, et $(f,g)\in L(E)^{2}$ avec
$E=\mathrm{Im} f+\mathrm{Im} g=\mathrm{Ker}(f)+\mathrm{Ker}(g). $ Montrer que ces sommes sont directes.
\finenonce{001036}



\finexercice



\exercice{1037, gourio, 2001/09/01}

\enonce{001037}{}
Soit $E$ un espace vectoriel de dimension finie, et $(f_{1},...,f_{k})$ des
projecteurs de $E. $ Montrer l'\'{e}quivalence :
$$\left[ \forall (i,j)\in \{1,...,k\}^{2},i\neq j\Rightarrow
f_{i}f_{j}=0\right] \Leftrightarrow \sum\limits_{i=1}^{k}f_{i} \text{ est un
projecteur}. $$
\finenonce{001037}



\finexercice



\exercice{1038, gourio, 2001/09/01}

\enonce{001038}{}
Soit $f\in L(E)$ o\`{u} $E$ est un ${K}$-espace vectoriel de dimension $n $,
 tel que : $$f^{2}=-Id. $$
\begin{enumerate}
\item Montrer que $f$ est inversible et que la dimension de $E$ est paire, donc $n=2p.$
\item
Soit $x\neq 0,$ monter que $x$ et $f(x)$ sont lin\'{e}airement
ind\'{e}pendants, et qu'ils engendrent un sous-espace stable de $E.$
\item Montrer qu'il existe $p$ sous-espaces de dimension deux stables par $f$, $%
E_{1}...E_{p}$ tels que :
$E=\bigoplus\limits_{i=1}^{p}E_{i}.$ En d\'{e}duire une ``bonne'' formule de
calcul de $f.$
\end{enumerate}
\finenonce{001038}



\finexercice



\exercice{1039, gourio, 2001/09/01}

\enonce{001039}{}
Soit $E$ un ${K}$ espace vectoriel de dimension finie $n\geq 1.$ Soit $f\in L(E)$
nilpotente. On note $q\in \Nn^{*}$ l'indice de nilpotence de $f,$ i.e.:
$$q=\inf \{j\in \Nn^{*}|f^{j}=0\}. $$
\begin{enumerate}
\item  Montrer que :
$\exists x_{0}\in E$ tel que $\{x_{0},f(x_{0}),...,f^{q-1}(x_{o})\}$ soit
libre.
En d\'{e}duire $q\leq n.$
\item Soit $r=\dim Ker(f). $ Montrer que $r>0$ et que
$$\frac{n}{r}\leq q\leq n+1-r. $$
\end{enumerate}
\finenonce{001039}



\finexercice



\exercice{3322, quercia, 2010/03/09}
\enonce{003322}{$\dim H = \dim K \Leftrightarrow H$ et $K$ ont un supplémentaire commun}
Soient $H,K$ deux sev d'un ev $E$ de dimension finie.
Montrer que $\dim H = \dim K$ si et seulement si $H \text{ et } K$ ont un supplémentaire commun
(par récurrence sur codim$\,H$).
\finenonce{003322}



\finexercice

\exercice{5183, rouget, 2010/06/30}
\enonce{005183}{**IT}
$E$ désigne l'espace vectoriel $\Rr^4$ (muni des opérations usuelles). On considère les vecteurs $e_1=(1,2,3,4)$,
$e_2=(1,1,1,3)$, $e_3=(2,1,1,1)$, $e_4=(-1,0,-1,2)$ et $e_5=(2,3,0,1)$. Soient alors $F=\mbox{Vect}(e_1,e_2,e_3)$ et
$G=\mbox{Vect}(e_4,e_5)$. Quelles sont les dimensions de $F$, $G$, $F\cap G$ et $F+G$~?
\finenonce{005183}


\finexercice


\exercice{5184, rouget, 2010/06/30}
\enonce{005184}{**IT}
Soit $\Kk$ un sous-corps de $\Cc$, $E$ un $\Kk$-espace vectoriel de dimension finie $n\geq2$. Soient $H_1$ et $H_2$ deux hyperplans de $E$.
Déterminer $\mbox{dim}_\Kk(H_1\cap H_2)$. Interprétez le résultat quand $n=2$ ou $n=3$.
\finenonce{005184}


\finexercice

\exercice{5575, rouget, 2010/10/16}
\enonce{005575}{***I}
Soient $F$ et $G$ deux sous-espaces vectoriels d'un espace vectoriel de dimension finie sur $\Kk$.

Démontrer que $\text{dim}(F+G)=\text{dim}F+\text{dim}G-\text{dim}(F\cap G)$.
\finenonce{005575}


\finexercice


\exercice{5576, rouget, 2010/10/16}
\enonce{005576}{**}
\label{ex:rou14}
Soient $F$, $G$ et $H$ trois sous-espaces d'un espace vectoriel $E$ de dimension finie sur $\Kk$.

Montrer que : $\text{dim}(F+G+H)\leqslant\text{dim}F+\text{dim}G+\text{dim}H-\text{dim}(F\cap G)-\text{dim}(G\cap H)-\text{dim}(H\cap F)+\text{dim}(F\cap G\cap H)$.

Trouver un exemple où l'inégalité est stricte.
\finenonce{005576}


\finexercice


\exercice{5577, rouget, 2010/10/16}
\enonce{005577}{***}
\label{ex:rou15}
Soient $F_1$, $F_2$,..., $F_n$ $n$ sous-espaces vectoriels d'un espace $E$ de dimension finie sur $\Kk$ $(n\geqslant2)$.

Montrer que $\text{dim}(F_1+...+F_n)\leqslant\text{dim}F_1+ ... +\text{dim}F_n$ avec égalité si et seulement si la somme est directe.
\finenonce{005577}


\finexercice


\exercice{5578, rouget, 2010/10/16}
\enonce{005578}{**I}
Soit $E$ un $\Kk$-espace vectoriel de dimension $n\geqslant3$. Montrer que l'intersection de $n-1$ hyperplans de $E$ est non nulle.
\finenonce{005578}


\finexercice


\exercice{5579, rouget, 2010/10/16}
\enonce{005579}{**}
Soient $(x_1,..,x_n)$ une famille de $n$ vecteurs de rang $r$ et $(x_1,...,x_m)$ une sous famille de rang $s$ ($m\leqslant n$ et $s\leqslant r$). Montrer que $s\geqslant r+m-n$. Cas d'égalité ?
\finenonce{005579}


\finexercice


\exercice{5580, rouget, 2010/10/16}
\enonce{005580}{**}
Soient $E$ et $F$ deux espaces vectoriels de dimension finie et soient $f$ et $g$ deux applications linéaires de $E$ dans $F$. Montrer que $|\text{rg}f-\text{rg}g|\leqslant\text{rg}(f+g)\leqslant\text{rg}f+\text{rg}g$.
\finenonce{005580}


\finexercice


\exercice{5581, rouget, 2010/10/16}
\enonce{005581}{**}
Soient $E$, $F$ et $G$, trois $\Kk$-espaces vectoriels puis $f\in\mathcal{L}(E,F)$ et $g\in\mathcal{L}(F,G)$.

Montrer que $\text{rg}f+\text{rg}g-\text{dim}F\leqslant\text{rg}(g\circ f)\leqslant\text{Min}\{\text{rg}f,\text{rg}g\}$.

\finenonce{005581}


\finexercice

\section{106.99 Autre}

\section{107.01 Définition}


\exercice{927, cousquer, 2003/10/01}

\enonce{000927}{}

\emph{Notations :\\
$\mathcal{C}$~: ensemble des fonctions numériques continues sur $[0,1]$.\\ 
$\mathcal{C}_d$~: ensemble des fonctions numériques ayant
une dérivée continue sur $[0,1]$.\\ 
$\mathcal{C}(\mathbb{R})$ et $\mathcal{C}^1(\mathbb{R})$~: définis de façon analogue
pour les fonctions définies sur $\mathbb{R}$.\\ 
$\mathcal{P}$~: ensemble des polynômes sur $\mathbb{R}$.\\ 
$\mathcal{P}_n$~: ensemble des polynômes sur $\mathbb{R}$, de degré $\leq n$.}


\medskip

Dire si les applications suivantes sont des applications linéaires~:

\begin{enumerate}

\item $\mathbb{R} \rightarrow \mathbb{R} : x\mapsto 2x^2$.

\item $\mathbb{R} \rightarrow \mathbb{R} : x\mapsto 4x-3$.

\item $\mathbb{R} \rightarrow \mathbb{R} : x\mapsto \sqrt{x^2}$.

\item $\mathbb{R}^2 \rightarrow \mathbb{R} ^2  :(x,y)\mapsto (y,x)$.

\item $\mathcal{C}\rightarrow \mathcal{C} : f \mapsto \{t\mapsto {f(t)\over1+t^2}\}$.

\item $\mathcal{C} \rightarrow \mathbb{R} : f\mapsto  f(3/4)$.

\item $\mathcal{C} \rightarrow \mathbb{R} : f\mapsto  f(1/4)-\int_{1/2}^1f(t)\,dt$.

\item $\mathbb{R}^2 \rightarrow \mathbb{R} :(x,y)\mapsto 3x+5y$.

\item $\mathbb{R}^2 \rightarrow \mathbb{R} : (x,y)\mapsto \sqrt{3x^2+5y^2}$.

\item $\mathbb{R}^2 \rightarrow \mathbb{R} : (x,y)\mapsto \sin(3x+5y)$.

\item $\mathbb{R}^2 \rightarrow \mathbb{R}^2 : (x,y)\mapsto (-x,y)$.

\item $\mathbb{R}^2 \rightarrow \mathbb{R} :(x,y)\mapsto  xy$.

\item $\mathbb{R}^2 \rightarrow \mathbb{R} : (x,y)\mapsto 
{x^2 y\over x^2+y^2} \mbox{ si } x^2+y^2 \ne0 \mbox{ et} 0 \mbox{ sinon}.$

\item $\mathcal{C} \rightarrow \mathcal{C}_d : f\mapsto \{x\mapsto  e^{-x}\int_0^1f(t)\,dt\}$.

\item $\mathcal{P} \rightarrow \mathcal{P}_n : A \mapsto $ quotient de $A$ par $B$
à l'ordre $n$ selon les puissances croissantes ($B$ et $n$ fixés,
avec $B(0)\ne0$).

\item $\mathbb{R}^2 \rightarrow \mathbb{R}^2 : M \mapsto  M'$ défini par: ${\overrightarrow{OM'}=
{\overrightarrow{OM}\over\bigl\Vert\overrightarrow{OM}\bigr\Vert} \mbox{ si } 
\overrightarrow{OM}\ne\overrightarrow 0 \mbox{ et  } 0 \mbox{ sinon}.}$

\item $\mathbb{R}^3 \rightarrow \mathbb{R} : M\mapsto \overrightarrow{OM}\cdot \overrightarrow V$
où $\overrightarrow V=(4,-1,1/2)$.

\item $\mathbb{R} \rightarrow \mathbb{R}^3 : x\mapsto (2x,x/\pi,x\sqrt2)$.

\item $\mathcal{C} \rightarrow \mathbb{R} : f\mapsto \max_{t\in[0,1]}f(t)$.

\item $\mathcal{C} \rightarrow \mathbb{R} : f\mapsto \max_{t\in[0,1]}f(t)-\min_{t\in[0,1]}f(t)$.

\item $\mathbb{R}^2 \rightarrow \mathbb{R}^2 : (x,y)\mapsto $ la solution du système d'équations
en $(u,v)$~:
$$\left\{
 \begin{array}{rcl}
   3u-v  & = & x \\
   6u+2v & = & y.
 \end{array}
\right.$$

\item $\mathbb{R}^2 \rightarrow \mathbb{R}^2 : (x,y)\mapsto $ le symétrique de $(x,y)$ par rapport à
la droite d'équation $x+y-a=0$ (discuter selon les valeurs de $a$).

\item $\mathbb{R}^3 \rightarrow \mathbb{R}^3 : (x,y,z)\mapsto $ la projection de $(x,y,z)$ sur le plan
$x+y+z-a=0$ parallèlement à $Oz$ (discuter selon les valeurs de $a$).

\item $\mathcal{C}_d \rightarrow \mathcal{C} : f\mapsto f'$.

\item $\mathbb{R}^3 \rightarrow \mathbb{R}^2 : (x,y,z)\mapsto (2x-3y+z,x-y+z/3)$.

\item $\mathbb{R} \rightarrow \mathcal{C}d : \lambda\mapsto $ la solution de l'équation
différentielle $y'-{y\over x^2+1}=0$ valant $\lambda$ en $x_0=1$.

\item $\mathcal{C}\rightarrow \mathbb{R} : f\mapsto \int_0^1\ln(1+\vert f(t)\vert)\,dt$.

\item $\mathbb{R} \rightarrow \mathbb{R} : x\mapsto $ la $17$-ième décimale de $x$
(en écriture décimale).

\item $\mathcal{C}_d \rightarrow \mathbb{R} : f \mapsto  f'(1/2)+\int_0^1{f(t)\,dt}$.

\item $\mathbb{R} \rightarrow \mathbb{R} : x\mapsto \ln(3^{x\sqrt2})$.

\item $\mathbb{R} \times \mathcal{C}(\mathbb{R}) \rightarrow \mathcal{C}(\mathbb{R}) : (\lambda,f)\mapsto $
la primitive de $f$ qui vaut $\lambda$ en $x_0=\pi$.

\item $\mathcal{C}^1(\mathbb{R}) \rightarrow \mathcal{C}(\mathbb{R}): f\mapsto \{x\mapsto 
f'(x)+f(x) \cdot
\sin x\}$.

\end{enumerate}
\finenonce{000927}



\finexercice



\exercice{928, legall, 1998/09/01}

\enonce{000928}{}
Soient $f$ et $g$, applications de ${ \Cc}$ dans ${ \Cc}$, d\'efinies par
$f(z)=\bar{z}$ et $g(z)=\Re (z)$. Montrer que $f$ et $g$ sont lin\'eaires sur ${ \Cc}$ en tant
que ${ \Rr}$-e.v., et non lin\'eaires sur ${ \Cc}$ en tant que ${ \Cc}$-e.v.
\finenonce{000928}



\finexercice



\exercice{929, legall, 1998/09/01}
\video{4CS7MiS5AQA}
\enonce{000929}{}
D\'eterminer si les applications $f_{i}$ suivantes sont lin\'eaires :
$$\begin{array}{rl}
f_1 : \Rr^2 \to \Rr^2 & f_1(x,y)=(2x+y,x-y)  \\
f_2 : \Rr^3 \to \Rr^3 & f_2(x,y,z)=(xy,x,y) \\
f_3 : \Rr^3 \to \Rr^3 & f_3(x,y,z)=(2x+y+z,y-z,x+y) \\
f_4 : \Rr^2 \to \Rr^4 & f_4(x,y)=(y,0,x-7y,x+y) \\
f_5 : \Rr_3[X] \to \Rr^3 & f_5(P) = \big( P(-1), P(0), P(1) \big) \\
\end{array}
$$
\finenonce{000929} 


\finexercice


\exercice{930, legall, 1998/09/01}
\video{EQJ4bUFDuiQ}
\enonce{000930}{}
Soit $E$ un espace vectoriel de dimension $n$ et
$\phi $  une application lin\'eaire de  $E$  dans lui-m\^eme telle que  $\phi ^n=0$  et
$\phi ^{n-1}\not = 0$.
Soit  $x\in E$  tel que  $\phi ^{n-1}(x )\not = 0$. Montrer que la
famille  $\{ x,\phi(x),\phi^2(x), \ldots ,\phi ^{n-1}(x)\} $  est une base de $E$.

\finenonce{000930} 


\finexercice


\exercice{2431, matexo1, 2002/02/01}

\enonce{002431}{}
Parmi les applications suivantes, lesquelles sont des formes
lin\'eaires sur $C^\infty(\R)$\,:
$$ f \mapsto f(0), \quad f \mapsto f(1) -1, \quad
f \mapsto f''(3), \quad f \mapsto (f'(2))^2,\quad
f \mapsto \int_0^1{f(t) dt}. $$
\finenonce{002431}


\finexercice


\exercice{2740, tumpach, 2009/10/25}
\enonce{002740}{}
\begin{enumerate}
\item On munit $\mathbb{R}^2$ d'un rep\`ere orthonorm\'e $(O, \vec{i}, \vec{j})$. Montrer qu'une application lin\'eaire de $\mathbb{R}^{2}$ dans $\mathbb{R}^2$ est uniquement d\'etermin\'ee par ses valeurs sur les vecteurs $\vec{i}$ et $\vec{j}$.
\item Quelle est la matrice de la sym\'etrie axiale par rapport \`a l'axe des abscisses dans la base $\{\vec{i}, \vec{j}\}$~?
\item Quelle est la matrice de la projection orthogonale sur l'axe des abscisses dans la base $\{\vec{i}, \vec{j}\}$~?
\item Quelle est la matrice de la rotation d'angle $\theta$ et de centre $O$ dans la base $\{\vec{i}, \vec{j}\}$~?
\item Quelle est la matrice de l'homoth\'etie de centre $O$ et de rapport $k$ dans la base $\{\vec{i}, \vec{j}\}$~?
\item Quelle est la matrice de la sym\'etrie centrale de centre $O$  dans la base $\{\vec{i}, \vec{j}\}$~?
\item Est-ce qu'une translation est une application lin\'eaire~?
\end{enumerate}
\finenonce{002740}



\finexercice


\exercice{2741, tumpach, 2009/10/25}
\enonce{002741}{}
Soit $f$ la fonction de $\mathbb{R}^4$ dans $\mathbb{R}^4$ d\'efinie par~:
$$
f(x, y, z, t) = \left(x + y + z + t, x + y + z + t, x + y + z + t, 2x + 2y + 2z + 2t\right).
$$
\begin{enumerate}
\item Montrer que $f$ est lin\'eaire et d\'eterminer sa matrice dans la base canonique de $\mathbb{R}^{4}$.
\item V\'erifier que les vecteurs $\vec{a} = (1, -1, 0, 0)$, $\vec{b} = (0, 1, -1, 0)$ et $\vec{c} = (0, 0, 1, -1)$ appartiennent \`a $\ker f$.
\item V\'erifier que le vecteur $\vec{d} = (5, 5, 5, 10)$ appartient \`a $\textrm{Im} f$. 
\end{enumerate}
\finenonce{002741}



\finexercice


\exercice{2742, tumpach, 2009/10/25}
\enonce{002742}{}
Soit l'application $f~:\mathbb{R}^3\rightarrow \mathbb{R}^3$ donn\'ee par~:
$$
f(x, y, z) = (x + 2y + z, 2x  + y + 3z, -x -y -z).
$$
\begin{enumerate}
\item Justifier que $f$ est lin\'eaire.
\item Donner la matrice de $A$ dans la base canonique de $\mathbb{R}^3$.
\item \begin{enumerate}
\item[(a)] D\'eterminer une base et la dimension du noyau de $f$, not\'e $\ker f$.
\item [(b)] L'application $f$ est-elle injective~?
\end{enumerate}
\item \begin{enumerate}
\item Donner le rang de $f$ et une base de $\textrm{Im} f$.
\item L'application $f$ est-elle surjective~?
\end{enumerate}
\end{enumerate}
\finenonce{002742}



\finexercice


\exercice{3309, quercia, 2010/03/09}
\enonce{003309}{Endomorphisme tel que tout vecteur non nul est propre}

Soit $E$ un espace vectoriel et $f \in \mathcal{L}(E)$ tel que pour tout
$\vec x \in E$, la famille $\bigl(\vec x, f(\vec x)\bigr)$ est liée.

\begin{enumerate}
  \item Montrer que si $\vec x \ne \vec 0$, il existe un unique scalaire
    $\lambda_{\vec x}$ tel que $f(\vec x) = \lambda_{\vec x}\vec x$.
  \item Comparer $\lambda_{\vec x}$ et $\lambda_{\vec y}$ lorsque
    $(\vec x, \vec y)$ est libre.
  \item Montrer que $f$ est une homothétie.
\end{enumerate}
\finenonce{003309}



\finexercice

\exercice{3312, quercia, 2010/03/09}
\enonce{003312}{Applications $\R$-linéaires sur $\C$}

On considère que $\C$ est un $\R$-espace vectoriel.

\begin{enumerate}
  \item Donner une base de $\C$.
  \item Montrer que tout endomorphisme de $\C$ peut se mettre sous la forme :
    $f(z) = az + b\overline{z}$, avec $a,b \in \C$.
  \item CNS sur $a$ et $b$ pour que $f$ soit bijectif ?
    
\end{enumerate}
\finenonce{003312}



\finexercice


\exercice{5170, rouget, 2010/06/30}
\enonce{005170}{**T}
\begin{enumerate}
\item  Vérifier qu'il existe une unique application linéaire de $\Rr^3$ dans $\Rr^2$ vérifiant  $f((1,0,0))=(1,1)$
puis $f((0,1,0))=(0,1)$ et $f((0,0,1))=(-1,1)$. Calculer $f((3,-1,4))$ et $f((x,y,z))$ en général.
\item  Déterminer $\mbox{Ker}f$. En fournir une base. Donner un supplémentaire de $\mbox{Ker}f$ dans $\Rr^3$ et
vérifier qu'il est isomorphe à $\mbox{Im}f$.
\end{enumerate}
\finenonce{005170}


\finexercice


\exercice{5186, rouget, 2010/06/30}
\enonce{005186}{***I}
Soit $E=\Rr_n[X]$, le $\Rr$-espace vectoriel des polynômes à coefficients réels de degré inférieur ou égal à $n$ ($n$
entier naturel donné). Soit $\varphi$ l'application définie par~:~$\forall P\in E,\;\varphi(P)=P(X+1)-P(X)$.
\begin{enumerate}
\item  Vérifier que $\varphi$ est un endomorphisme de $E$.
\item  Déterminer $\mbox{Ker}\varphi$ et $\mbox{Im}\varphi$.
\end{enumerate}
\finenonce{005186}


\finexercice

\exercice{5188, rouget, 2010/06/30}
\enonce{005188}{**}
Soit $\begin{array}[t]{cccc}
f~:&\Cc&\rightarrow&\Cc\\
 &z&\mapsto&z+a{\bar z}
\end{array}
$ où $a$ est un nombre complexe donné non nul.
Montrer que $f$ est un endomorphisme du $\Rr$-espace vectoriel $\Cc$. $f$ est-il un endomorphisme du $\Cc$-espace
vectoriel $\Cc$~?~Déterminer le noyau et l'image de $f$.
\finenonce{005188}


\finexercice


\exercice{5189, rouget, 2010/06/30}
\enonce{005189}{**}
Soit $f\in\mathcal{L}(\Rr^2)$. Pour $(x,y)\in\Rr^2$, on pose $f((x,y))=(x',y')$.
\begin{enumerate}
\item  Rappeler l'écriture générale de $(x',y')$ en fonction de $(x,y)$.
\item  Si on pose $z=x+iy$ et $z'=x'+iy'$ (où $i^2=-1$), montrer que~:~$\exists(a,b)\in\Cc^2/\;\forall
z\in\Cc,\;z'=az+b{\bar z}$.

\item  Réciproquement, montrer que l'expression ci-dessus définit un unique endomorphisme de
$\Rr^2$ (en clair, l'expression complexe d'un endomorphisme de $\Rr^2$ est $z'=az+b{\bar z}$).
\end{enumerate}
\finenonce{005189}


\finexercice

\section{107.02 Image et noyau, théorème du rang}


\exercice{931, liousse, 2003/10/01}

\enonce{000931}{}
Soient $F$ et $G$ deux sous-espaces vectoriels de $\Rr^n$, 
on d\'efinit l'application 
$f:F \times G \rightarrow \Rr^n$ par $f(x_1,x_2)= x_1+x_2$.
 \begin{enumerate}
  \item Montrer que f est lin\'eaire.
 \item  D\'eterminer le noyau et l'image de $f$.
 \end{enumerate}
\finenonce{000931}



\finexercice



\exercice{932, liousse, 2003/10/01}

\enonce{000932}{}
Soit $f$ une application lin\'eaire de $\Rr^n$ dans $\Rr^n$. 
Montrer que les 
propri\'et\'es $(1)$ \`a $(3)$ sont \'equivalentes.
$$\Rr^n=Im (f) \bigoplus Ker (f)\leqno(1)$$
$$Im (f)= Im (f^2)\leqno(2)$$
$$Ker (f)= Ker (f^2)\leqno(3)$$
\finenonce{000932}



\finexercice



\exercice{933, liousse, 2003/10/01}

\enonce{000933}{}
Soient :
$E$, $F$ et $G$ trois sous espaces vectoriels de $\mathbb R^N$,
 $f$ une application lin\'eaire de $E$ dans $F$ et 
 $g$ une application lin\'eaire de $F$ dans $G$.
On rappelle que  $g\circ f$ est l'application de $E$ dans $G$ d\'efinie par $
g\circ f(v) = g(f(v))$, pour tout vecteur $v$ de $E$.
\begin{enumerate}
\item Montrer que $g\circ f$ est une application lin\'eaire.
\item Montrer que $f\big({\text Ker}(g\circ f) \big) = {\text Ker} g \cap {\text Im} f$.
\end{enumerate}
\finenonce{000933}



\finexercice



\exercice{934, cousquer, 2003/10/01}
\video{pNH1v_tXHvg}
\enonce{000934}{}
Soit $E$ un espace vectoriel et soient $E_1$ et $E_2$ deux sous-espaces vectoriels de dimension finie
de $E$, on d\'efinit l'application $f\colon E_1\times E_2 \to E$ par $f(x_1,x_2)=x_1+x_2$.
\begin{enumerate}
\item Montrer que $f$ est lin\'eaire.
\item D\'eterminer le noyau et l'image de $f$.
\item Que donne le th\'eor\`eme du rang ?
\end{enumerate}

\finenonce{000934} 


\finexercice


\exercice{935, cousquer, 2003/10/01}

\enonce{000935}{}
Soit $E$ l'espace vectoriel des polynômes de degré inférieur ou
égal à $n$. Pour $p\leq n$ on note $e_p$ le polynôme
$x\mapsto x^p$. Soit $f$ l'application définie sur $E$ par $f(P)=Q$
avec $Q(x)=P(x+1)+P(x-1)-2P(x)$.
\begin{enumerate}
\item  Montrer que $f$ est une application linéaire de $E$ dans $E$.
\item  Calculer $f(e_p)$~; quel est son degré~? En déduire $\ker f$,
$\mbox{Im }f$ et le rang de $f$.
\item  Soit $Q$ un polynôme de $\mbox{Im }f$~; montrer qu'il existe un polynôme
unique $P$ tel que~: $f(P)=Q$ et $P(0)=P'(0)=0$.
\end{enumerate}
\finenonce{000935}



\finexercice



\exercice{936, cousquer, 2003/10/01}

\enonce{000936}{}
 Soit $E$, $F$, $G$ trois espaces vectoriels, $f$ et $g$ deux applications
linéaires $E\buildrel f\over\rightarrow\nolinebreak F\buildrel g\over\rightarrow
\nolinebreak G$~; montrer que~:
$$\ker(g\circ f) = f^{-1}(\ker g\cap\mbox{Im }f)=f^{-1}(\ker g).$$
\finenonce{000936}



\finexercice



\exercice{937, cousquer, 2003/10/01}

\enonce{000937}{}
 Soit $E$ un espace vectoriel de dimension finie, $F$ et $N$ deux
sous-espaces vectoriels de $E$~; donner une condition nécessaire
et suffisante pour qu'il existe une application linéaire $f$ de $E$
dans~$E$ vérifiant~: $f(E)=F$ et $\ker f=N$.
\finenonce{000937}



\finexercice



\exercice{938, cousquer, 2003/10/01}

\enonce{000938}{}
 Soit $E$, $F$, $G$ trois espaces vectoriels de dimensions respectives $n$,
$p$, $q$, $f$ et $g$ deux applications linéaires
$E\buildrel f\over\rightarrow F\buildrel g\over\rightarrow G$ telles que
$g\circ f=0$. Quelle relation existe-t-il entre le rang de $f$ et celui de $g$~?
\finenonce{000938}



\finexercice



\exercice{939, cousquer, 2003/10/01}

\enonce{000939}{}
 Soit $E$ un espace vectoriel de dimension finie, $f$ une application
linéaire de~$E$ dans~$E$~; montrer que les propriétés (1) à (3) sont
équivalentes~:
\begin{enumerate}
    \item[(1)] $ E=\mbox{Im }f \oplus\ker f$,
    \item[(2)] $ \mbox{Im }f = \mbox{Im }f^2$,
    \item[(3)] $\ker f=\ker f^2$.
\end{enumerate}
\finenonce{000939}



\finexercice



\exercice{940, cousquer, 2003/10/01}

\enonce{000940}{}
 Soit $E$ un espace vectoriel, et $u$ une application linéaire de~$E$ 
dans~$E$. Dire si les
propriétés suivantes sont vraies ou fausses~:
\begin{enumerate}
\item  Si $e_1,e_2,\ldots,e_p$ est libre, il en est de même de
$u(e_1),u(e_2),\ldots,u(e_p)$.
\item  Si $u(e_1),u(e_2),\ldots,u(e_p)$ est libre, il en est de même de
$e_1,e_2,\ldots,e_p$.
\item  Si $e_1,e_2,\ldots,e_p$ est génératrice, il en est de même de
$u(e_1),u(e_2),\ldots,u(e_p)$.
\item  Si $u(e_1),u(e_2),\ldots,u(e_p)$ est génératrice, il en est de même de
$e_1,e_2,\ldots,e_p$.
\item  Si $u(e_1),u(e_2),\ldots,u(e_p)$ est une base de $\Im u$, alors
$e_1,e_2,\ldots,e_p$ est une base d'un sous-espace vectoriel supplémentaire
de $\mathrm{Ker} u$.
\end{enumerate}
\finenonce{000940}



\finexercice



\exercice{941, legall, 1998/09/01}
\enonce{000941}{}
 Soient  $E$  un espace vectoriel et  $\varphi $
une application lin\' eaire de  $E$  dans  $E$. On suppose que
$\hbox {Ker } (\varphi ) \cap \hbox{Im }(\varphi )=\{ 0\}$.
Montrer que, si  $x \not \in \hbox {Ker } (\varphi )$  alors, pour
tout  $n\in { \Nn }  : \varphi ^n(x)\not = 0$.

\finenonce{000941} 


\finexercice


\exercice{942, legall, 1998/09/01}

\enonce{000942}{}
Pour des applications lin\'eaires $f: E\rightarrow F$, $g:
F\rightarrow G$, \'etablir l'\'equivalence
$$g\circ f=0\Longleftrightarrow \text{Im} f\subset \text{Ker} g.$$

Soit $f$ un endomorphisme d'un e.v. $E$, v\'erifiant l'identit\'e
$f^2+f-2i_E=0$. Etablir
$\text{Im} (f-i_E)\subset \text{Ker} (f+2i_E)$;  $\;\;\text{Im} (f+2i_E)\subset \text{Ker}
(f-i_E)$;  $\;\;E=\text{Ker}(f-i_E)\oplus \text{Ker}(f+2i_E)$.
\finenonce{000942}



\finexercice



\exercice{943, legall, 1998/09/01}
\video{yGZYazY1EhM}
\enonce{000943}{}
Soit $E$  un espace vectoriel de dimension  $n$
et $f$ une application lin\' eaire de  $E$  dans lui-m\^eme.
Montrer que les deux assertions qui suivent sont \' equivalentes :
\begin{enumerate}
\item[(i)] $\ker f = \Im f$
\item[(ii)] $f^2=0 \ \text{ et } \  n=2\cdot \text{rg}(f)$
\end{enumerate}
\finenonce{000943}
 

\finexercice


\exercice{944, legall, 1998/09/01}

\enonce{000944}{}
Soient  $E$  un espace vectoriel et  $F$  un sous-espace vectoriel
de  $E$  de dimension finie. Soit  $f$  une application lin\' eaire de
$E$  dans lui-m\^eme.
\begin{enumerate}
    \item Montrer que, si $F\subset f(F)$  alors  $f(F)=F$.
    \item Montrer que, si  $f$  est injective et
$f(F)\subset F$  alors  $f(F)=F$.
\end{enumerate}
\finenonce{000944}



\finexercice



\exercice{945, ridde, 1999/11/01}

\enonce{000945}{}
Soient $f : E \rightarrow F$ et $g : F \rightarrow G$ deux applications lin\'eaires.
Montrer que $\ker (f) \subset \ker (g \circ f)$ et $\text{Im} (g \circ f)
\subset \text{Im} (f)$.
\finenonce{000945}



\finexercice



\exercice{946, legall, 1998/09/01}

\enonce{000946}{}
 Soit  $E$  un espace vectoriel de dimension finie et  $\varphi $
une application lin\' eaire de  $E$  dans lui-m\^eme. Posons
$K_n=\hbox{Ker }(\varphi ^n)$  et $I_n=\hbox{Im }(\varphi ^n)$.
Montrer qu'il existe  $n_0\in {\Nn}$ tel que  pour tout  $n\geq
n_0$  on ait  $K_n=K_{n_0}$. D\' eduiser en que pour tout  $n\geq
n_0$  on a \' egalement  $I_n=I_{n_0}$.
\finenonce{000946}



\finexercice



\exercice{947, ridde, 1999/11/01}
\video{qUT7yUQxk4w}
\enonce{000947}{}
Soient $f$ et $g$ deux endomorphismes de $E$ tels que $f \circ g = g \circ f$.
Montrer que $\ker f$ et $\Im f$ sont stables par $g$.
\finenonce{000947} 


\finexercice


\exercice{948, ridde, 1999/11/01}

\enonce{000948}{}
Soit $f \in \mathcal{L} (E)$ telle que $f^3 = f^2 + f$. Montrer que
$E = \ker (f) \oplus \text{Im} (f)$ (on remarquera que $f \circ (f^2-f-id) = 0$).
\finenonce{000948}



\finexercice



\exercice{949, ridde, 1999/11/01}
\enonce{000949}{}
 Soit $f \in \mathcal{L} (E)$. Montrer que $\ker (f)
\cap \text{Im} (f) = f (\ker (f\circ f))$.
\finenonce{000949} 


\finexercice


\exercice{950, gourio, 2001/09/01}

\enonce{000950}{}
Soit $U$ un sous-espace vectoriel de $E$ espace vectoriel, et
$$A=\{f\in L(E)|U\subset \mathrm{Ker}(f)\}.$$
Montrer que $A$ est un sous-espace vectoriel de $L(E).$
\finenonce{000950}



\finexercice



\exercice{951, gourio, 2001/09/01}
\enonce{000951}{}
Donner des exemples d'applications lin\'{e}aires de $\Rr^{2}$ dans $\Rr^{2}$
v\'{e}rifiant :
\begin{enumerate}
\item  $\mathop{\mathrm{Ker}}\nolimits(f)=\mathop{\mathrm{Im}}\nolimits (f).$
\item $\mathop{\mathrm{Ker}}\nolimits(f)$ inclus strictement dans $\mathop{\mathrm{Im}}\nolimits (f).$
\item $\mathop{\mathrm{Im}}\nolimits (f)$ inclus strictement dans $\mathop{\mathrm{Ker}}\nolimits(f).$
\end{enumerate}

\finenonce{000951} 


\finexercice


\exercice{952, gourio, 2001/09/01}

\enonce{000952}{}
 Soit $(u,v)\in (L(E))^{2}$, tels que $u^{2}=u$ et
$vu=0. $ Montrer que
$$\mathrm{Im} (u+v)=\mathrm{Im} (u)+\mathrm{Im} (v).$$
\finenonce{000952}



\finexercice



\exercice{953, liousse, 2003/10/01}

\enonce{000953}{}
Soit $(\vec{e_1}, \vec{e_2}, \vec{e_3})$ une base 
de $\Rr^3$, et $\lambda$ un nombre 
r\'eel. D\'emontrer que la donn\'ee de 
$$\left\{\begin{array}{rcl}\phi(\vec{e_1})&=&\vec{e_1}+\vec{e_2} 
\\ \phi(\vec{e_2})&=&\vec{e_1}-\vec{e_2} \\ 
\phi(\vec{e_3})&=&\vec{e_1}+\lambda \vec{e_3} \end{array}\right.$$
d\'efinit une application lin\'eaire de $\Rr^3$ dans $\Rr^3$. 
Ecrire l'image du 
vecteur $\vec{v}=a_1 \vec{e_1}+a_2 \vec{e_2}+ a_3 \vec{e_3}$. 
Comment choisir $\lambda$ 
pour que $\phi$ soit injective ? surjective ?
\finenonce{000953}



\finexercice



\exercice{954, cousquer, 2003/10/01}
\video{nAFU6wBhyrE}
\enonce{000954}{}
 Soit $E$ un espace vectoriel de dimension $3$, $\{e_1,e_2,e_3\}$ une base
de $E$, et $t$ un param\`etre r\'eel. \\
D\'emontrer que la donn\'ee de
$\left\{
\begin{array}{rcl}
    \phi(e_1) & = & e_1+e_2  \\
    \phi(e_2) & = & e_1-e_2  \\
    \phi(e_3) & = & e_1+t e_3
\end{array}\right.$
d\'efinit une application lin\'eaire
$\phi$ de $E$ dans $E$. \'Ecrire le transform\'e du vecteur 
$x=\alpha_1e_1+\alpha_2e_2+\alpha_3e_3$. Comment choisir $t$ pour que 
$\phi$ soit injective ? surjective ?
\finenonce{000954} 


\finexercice


\exercice{955, cousquer, 2003/10/01}

\enonce{000955}{}
 $E$ étant un espace vectoriel de dimension $n$ sur $\mathbb{R}$, $f$ une
application linéaire de~$E$ dans~$E$, construire dans
les trois cas suivants deux applications linéaires bijectives $u$ et $v$ de~$E$
dans~$E$ telles que $f=u-v$.
\begin{itemize}
\item $f$ est bijective.
\item $\mathrm{Ker} f+\mathrm{Im} f=E$.
\item $f$ est quelconque.
\end{itemize}
\finenonce{000955}



\finexercice



\exercice{956, legall, 1998/09/01}
\video{DY3GrL-j6C4}
\enonce{000956}{}
Pour les applications lin\'eaires suivantes,  d\'eterminer $\ker f_i$ et 
$\Im f_i$. En d\'eduire si $f_i$ est injective, surjective, bijective.
$$\begin{array}{rl}
f_1 : \Rr^2 \to \Rr^2 & f_1(x,y)=(2x+y,x-y)  \\
f_2 : \Rr^3 \to \Rr^3 & f_2(x,y,z)=(2x+y+z,y-z,x+y) \\
f_3 : \Rr^2 \to \Rr^4 & f_3(x,y)=(y,0,x-7y,x+y) \\
f_4 : \Rr_3[X] \to \Rr^3 & f_4(P) = \big( P(-1), P(0), P(1) \big) \\
\end{array}
$$
\finenonce{000956} 


\finexercice


\exercice{957, gourio, 2001/09/01}

\enonce{000957}{}
Soit $f\in L(E)$ non nul; montrer que $f$ est injective si et seulement si
pour tout couple $(E_{1},E_{2})$ de sous-espaces suppl\'{e}mentaires de $E$,
la somme $f(E_{1})+f(E_{2})$ est directe (i.e. $f(E_{1})$ et $f(E_{2})$ sont
suppl\'{e}mentaires).
\finenonce{000957}



\finexercice



\exercice{958, gourio, 2001/09/01}

\enonce{000958}{}
Soit $f\in L(E)$ o\`{u} $E$ est un $K-$espace vectoriel.\ On suppose :
$$\forall x\in E,\exists \lambda \in K,f(x)=\lambda x.$$
Montrer :
$$\exists \mu \in K,f=\mu id.$$
\finenonce{000958}



\finexercice



\exercice{959, legall, 1998/09/01}
\video{ce0mL82x6Y0}
\enonce{000959}{}
Soit $E = \Rr_n[X]$ et soient $A$ et $B$ deux polyn\^omes \`a coefficients réels de
degr\'e $n+1$. On consid\`ere l'application $f$ qui \`a tout polyn\^ome $P$ de $E$, associe
le reste de la division euclidienne de $AP$ par $B$.
\begin{enumerate}
    \item Montrer que $f$ est un endomorphisme de $E$.
    \item Montrer l'\'equivalence
$$
f \hbox{ est bijective} \Longleftrightarrow \hbox{$A$ et $B$ sont premiers entre eux}.
$$
\end{enumerate}
\finenonce{000959} 


\finexercice


\exercice{960, ridde, 1999/11/01}

\enonce{000960}{}
Soit $f \in \mathcal{L} (E)$ telle que $f^3 = f^2 + f + id$. Montrer que $f$
est un automorphisme.
\finenonce{000960}



\finexercice



\exercice{961, legall, 1998/09/01}

\enonce{000961}{}
Soit $E$ un $\C$--espace vectoriel et $f \in{\mathcal L}  (E)$ tel que $f^2 - 3f + 2Id
=0_{{\mathcal L}(E)}$.
\begin{enumerate}
    \item Montrer que $f$ est un automorphisme.
    \item Montrer que $E = \ker(f-Id) \oplus \ker(f-2Id)$.
    \item D\'eduire de 2. que si $E$ est de dimension finie $n$, il existe une base $\beta =
(\epsilon_i)_{1\le i\le n}$, telle que $\forall i, f(\epsilon_i) = \lambda_i
\epsilon_i$ avec $\lambda_i = 1$ ou $\lambda_i=2$.
\end{enumerate}
\finenonce{000961}



\finexercice



\exercice{962, legall, 1998/09/01}

\enonce{000962}{}
Montrer que si  $p<q$  il n'existe pas d'application lin\' eaire
surjective de  ${\Rr}^p$  dans  ${\Rr}^q$. Montrer que si  $q<p $  il
n'existe pas non plus
d'application lin\' eaire injective de  ${\Rr}^p$  dans  ${\Rr}^q$.
\finenonce{000962}



\finexercice



\exercice{963, legall, 1998/09/01}
\video{ZIq3tKZZgfg}
\enonce{000963}{}
Soit  $E$  et  $F$  deux espaces vectoriels de dimension finie
et  $\phi $  une application lin\' eaire  de  $E$  dans  $F$.
Montrer que  $\phi $  est un isomorphisme si et seulement si l'image par  $\phi $  de
toute base de  $E$  est une base de  $F$.
\finenonce{000963} 


\finexercice


\exercice{964, legall, 1998/09/01}

\enonce{000964}{}
\begin{enumerate}
    \item Soient $ E $ et $ F $ deux espaces
vectoriels et $ \varphi  $ une application lin\'eaire  bijective de
$ E $ dans $ F .$
Montrer que la bijection r\'eciproque $ \varphi ^{-1} $ est
lin\'eaire. Une telle
application est dite un isomorphisme d'espaces vectoriels.
    \item Soient $ E $ et $ F $ deux espaces
vectoriels de dimension finie. Montrer qu'il existe un
isomorphisme d'espaces
vectoriels de $ E $ \`a valeurs dans $ F $ si et seulement si $
\hbox{dim} (E)=\hbox{dim} (F) .$
\end{enumerate}
\finenonce{000964}



\finexercice



\exercice{965, legall, 1998/09/01}

\enonce{000965}{}
Soit $ E $ un espace vectoriel de dimension finie $ \varphi  $ et
$ \psi  $ deux applications lin\'eaires  de $ E $ dans lui-m\^eme
telles que $ \varphi \circ \psi =\hbox{id}_{E} .$ Montrer que $ \psi \circ
\varphi =\hbox{id}_{E} .$
\finenonce{000965}



\finexercice



\exercice{2441, matexo1, 2002/02/01}

\enonce{002441}{}
Soit $E$ un espace vectoriel de dimension
finie $n$, et $u, v$ deux endomorphismes de $E$.

\begin{enumerate}
\item Montrer que $u \circ v =0$ si et seulement si l'image de
$v$ est contenue dans le noyau de $u$.
\item Soit $(e_1, \ldots, e_n)$ une base de $E$. On suppose
dans cette question que $u$ et $v$ s'expriment dans cette base par
$$u(e_1) = e_1, \qquad u(e_i) = 0\quad\hbox{si}\quad i \neq 1,$$
$$v(e_2) = e_2, \qquad v(e_i) = 0\quad\hbox{si}\quad i \neq 2.$$
Trouver les matrices de $u$, $v$ et $u\circ v$ dans cette base.
\item Si $u$ est un endomorphisme quelconque non nul de $E$, quelle
condition doit v\'erifier le noyau de $u$ pour qu'il existe un
endomorphisme non nul $v$ tel que $u\circ v =0$\,? Dans ce cas, $u$
est-il bijectif\,?
\end{enumerate}

\finenonce{002441}


\finexercice


\exercice{2743, tumpach, 2009/10/25}
\enonce{002743}{}
\begin{enumerate}
\item
Soit $f$ une application lin\'eaire surjective de $\mathbb{R}^{4}$ dans $\mathbb{R}^2$.  Quelle est la dimension du noyau de $f$~?
\item Soit $g$ une application injective de $\mathbb{R}^{26}$ dans $\mathbb{R}^{100}$. Quelle est la dimension de l'image de $g$~?
\item Existe-t-il une application lin\'eaire bijective entre $\mathbb{R}^{50}$ et $\mathbb{R}^{72}$~?
\end{enumerate}
\finenonce{002743}



\finexercice


\exercice{2744, tumpach, 2009/10/25}
\enonce{002744}{}
Soit la matrice $$A = \left( \begin{array}{ccc}2 & 7 & 1\\-1 & 2 & 0\\ 3 & 5 &  1 \end{array}\right).$$
\begin{enumerate}
\item D\'eterminer une base du noyau de $A$.
\item D\'eterminer une base de l'image de $A$.
\end{enumerate}
\finenonce{002744}



\finexercice


\exercice{2745, tumpach, 2009/10/25}
\enonce{002745}{}
Soit la matrice $$B = \left( \begin{array}{cccc}1 & 2 & 3 & 1\\-1 & 2 & -1 &-3\\ -3 & 5 &  2 &-3\end{array}\right).$$
\begin{enumerate}
\item D\'eterminer une base du noyau de $B$.
\item D\'eterminer une base de l'image de $B$.
\end{enumerate}
\finenonce{002745}



\finexercice


\exercice{2746, tumpach, 2009/10/25}
\enonce{002746}{}
Soit la matrice $$C = \left( \begin{array}{ccc}-1 & 3 & 1\\1 & 2 & 0\\ 2 & -1 &  -1\\ 2 & 4 & 0\\1 & 7& 1\end{array}\right).$$
\begin{enumerate}
\item D\'eterminer une base du noyau de $C$.
\item D\'eterminer une base de l'image de $C$.
\end{enumerate}
\finenonce{002746}



\finexercice


\exercice{2770, tumpach, 2009/10/25}
\enonce{002770}{}
Pour chaque couple de matrices $(A_i, b_i)$, $1\leq i\leq 5$, ci-dessous
\begin{enumerate}
\item donner la nature de l'ensemble des solutions du syst\`eme $A_i X = b_i$ ;
\item donner une repr\'esentation param\'etrique de l'ensemble des solutions de $A_i X = b_i$ ;
\item donner une base de l'image et une base du noyau de $A_i$.
\end{enumerate}
{\footnotesize
\begin{tabular}{ll}
a) $A_1 = \left(\begin{array}{cccc}1 & 2 & 3 & 4\\ 0 & 1 & 2 & 3\\ 0 & 0 & 1 & 2 \\ 0 & 0 & 0 &1\end{array}\right)  \quad\quad b_1 = \left(\begin{array}{c}1\\1\\1\\1\end{array}\right) ;$ &
b) $A_2 = \left(\begin{array}{ccccc}1 & 2 & 0 & 1 & 3\\ 0 & 1 & 1 & 1&2\\ 0 & 0 & 1 & 2& 3 \\ 0 & 0 & 0 &1 & 1\end{array}\right)  \quad\quad b_2 = \left(\begin{array}{c}1\\1\\1\\1\end{array}\right) ;$\\ & \\
c) $A_3 = \left(\begin{array}{cccc}1 & 2 & 3 & 4\\ 0 & 1 & 2 & 3\\ 0 & 0 & 1 & 2 \\ 0 & 0 & 0 &1\\0 & 0 & 0& 0\end{array}\right)  \quad\quad b_3 = \left(\begin{array}{c}1\\1\\1\\1\\1\end{array}\right) ;$ &
d) $A_4 = \left(\begin{array}{ccccc}1 & 2 & 0 & 1 & 1\\ 0 & 1 & 1 & 2&2\\ 0 & 0 & 1 & 2& 1 \\ 0 & 0 & 0 &1 & 1 \\ 0 & 0 & 0 & 0& 0\end{array}\right)  \quad\quad b_4 = \left(\begin{array}{c}1\\1\\1\\1 \\ 1\end{array}\right) ;$\\ & \\
e) $A_5 = \left(\begin{array}{ccccc}1 & 2 & 0 & 1 & 1\\ 0 & 1 & 1 & 2&2\\ 0 & 0 & 1 & 2& 1 \\ 0 & 0 & 0 &1 & 1 \\ 0 & 0 & 0 & 0& 0\end{array}\right)  \quad\quad b_5 = \left(\begin{array}{c}1\\1\\1\\1 \\ 0\end{array}\right) ;$
\end{tabular}}
\finenonce{002770}



\finexercice


\exercice{2771, tumpach, 2009/10/25}
\enonce{002771}{}
Calculer une base de l'image et une base du noyau de l'application lin\'eaire 
$$
\begin{array}{cccl}f~:&\mathbb{R}^3&\longrightarrow &\mathbb{R}^5\\
& (x, y, z)&\longmapsto & (x + y, x + y + z, 2x + y + z, 2x + 2y + z, y + z)
\end{array}
$$
Quel est le rang de $f$~?
\finenonce{002771}



\finexercice


\exercice{3310, quercia, 2010/03/09}
\enonce{003310}{$f\circ g \circ f = f$ et $g \circ f\circ g = g$}

Soient $f,g \in \mathcal{L}(E)$ tels que $f\circ g \circ f = f$ et $g \circ f\circ g = g$.
\begin{enumerate}
  \item Montrer que $E = \mathrm{Ker} f \oplus \Im g$.
    
  \item Montrer que $f(\Im g) = \Im f$.
    
\end{enumerate}
\finenonce{003310}



\finexercice


\exercice{3311, quercia, 2010/03/09}
\enonce{003311}{$f^3 = \mathrm{id}$}

Soit $f \in \mathcal{L}(E)$ tel que $f^3 = \mathrm{id}_E$.

\begin{enumerate}
  \item  Montrer que $\mathrm{Ker}(f-\mathrm{id}) \oplus \Im(f-\mathrm{id}) = E$.
  \item  Montrer que $\mathrm{Ker}(f-\mathrm{id}) = \Im(f^2 + f + \mathrm{id})$
 et $\mathrm{Im}(f-\mathrm{id}) = \mathrm{Ker}(f^2 + f + \mathrm{id})$.
\end{enumerate}
\finenonce{003311}



\finexercice

\exercice{3313, quercia, 2010/03/09}
\enonce{003313}{Supplémentaire d'un hyperplan}

Soit $E$ un $ K$-ev et $f : E \to { K}$ une forme linéaire non identiquement
nulle. On note $H = \mathrm{Ker} f$.

\begin{enumerate}
  \item Montrer que $\Im f =  K$.
  \item Soit $\vec u \in E \setminus H$ et $F = \text{vect}(\vec u)$.
    Montrer que $F \oplus H = E$.
\end{enumerate}
\finenonce{003313}



\finexercice


\exercice{3316, quercia, 2010/03/09}
\enonce{003316}{Commutants itérés}
Soit $u\in\mathcal{L}(E)$. On pose pour $v\in\mathcal{L}(E)$~: $\varphi(v) = v\circ u - u\circ v$,
et on note $c_i = \mathrm{Ker}\varphi^i$ ($c_0 = \{0\}$, $c_1$ est le commutant
de~$u$, $c_2$ est l'ensemble des $v$ tels que $v\circ u - u\circ v$ commute
avec $u$,\dots).

\begin{enumerate}
  \item Calculer $\varphi(v\circ w)$ en fonction de $v,w,\varphi(v)$ et $\varphi(w)$.

  \item Montrer que $c = \bigcup_{i\in\N}c_i$ est une sous-algèbre de~$\mathcal{L}(E)$.


\end{enumerate}
\finenonce{003316}



\finexercice

\exercice{3327, quercia, 2010/03/09}
\enonce{003327}{Applications du thm du rang}

Soient $E,F$ deux $ K$-ev et $f \in \mathcal{L}(E,F)$.

\begin{enumerate}
  \item Montrer que si $H$ est un sev de $E$, alors \\
    $\dim f(H) = \dim H - \dim (H\cap\mathrm{Ker} f)$.
  \item Montrer que si $K$ est un sev de $F$, alors\\
    $\dim f^{-1}(K) = \dim (K\cap\Im f) + \dim(\mathrm{Ker} f)$.
\end{enumerate}
\finenonce{003327}



\finexercice


\exercice{3328, quercia, 2010/03/09}
\enonce{003328}{Application du thm du rang}

Soient $E$, $F$ deux ev de dimensions finies et $u,v \in \mathcal{L}(E,F)$.

Montrer que
$\dim(\mathrm{Ker}(u+v)) \le \dim(\mathrm{Ker} u\cap\mathrm{Ker} v) + \dim(\Im u\cap \Im v)$.

(considérer $w = u_{|\mathrm{Ker}(u+v)}$)

\finenonce{003328}



\finexercice


\exercice{3329, quercia, 2010/03/09}
\enonce{003329}{Rang de $f\circ g$}

Soit $E$ un ev de dimension finie et $f,g \in \mathcal{L}(E)$. \'Etablir :

\begin{enumerate}
  \item $\dim\mathrm{Ker}(f\circ g) \le \dim\mathrm{Ker} f + \dim\mathrm{Ker} g$.
  \item $\dim(\Im f \cap \mathrm{Ker} g) = \mathrm{rg}(f) - \mathrm{rg}(g\circ f)$.
  \item $\mathrm{rg}(f) + \mathrm{rg}(g) - \dim E \le \mathrm{rg}(f\circ g) \le \min(\mathrm{rg}(f),\mathrm{rg}(g))$.

\end{enumerate}
\finenonce{003329}



\finexercice


\exercice{3330, quercia, 2010/03/09}
\enonce{003330}{CNS pour que $\mathrm{Ker} f$ et $\Im f$ soient supplémentaires}

Soit $E$ un ev de dimension finie et $f \in \mathcal{L}(E)$.
Montrer que les propriétés suivantes sont équivalentes :


\begin{enumerate}
  \item $\mathrm{Ker} f^2 = \mathrm{Ker} f$.
  \item $\Im  f^2 = \Im  f$.
  \item $\mathrm{Ker} f \oplus \Im f = E$.
  \item $\mathrm{Ker} f \cap \Im f = \{\vec 0\}$.
  \item $\mathrm{Ker} f + \Im f = E$.
\end{enumerate}
\finenonce{003330}



\finexercice


\exercice{3331, quercia, 2010/03/09}
\enonce{003331}{$fog = 0$}

Soit $E$ un ev de dimension finie et $f,g \in \mathcal{L}(E)$ tels que $f\circ g = 0$.
Trouver une inégalité liant les rangs de $f$ et de $g$.
Peut-on avoir égalité ?
\finenonce{003331}



\finexercice


\exercice{3332, quercia, 2010/03/09}
\enonce{003332}{Rang de $f+g$}

Soient $E,F$ deux ev, $E$ de dimension finie, et $f,g \in \mathcal{L}(E,F)$.

\begin{enumerate}
  \item Démontrer que $\mathrm{rg}(f+g) \le \mathrm{rg}(f) + \mathrm{rg}(g)$.
  \item Montrer qu'il y a égalité si et seulement si
    $\Im f \cap \Im g = \{\vec 0_F \}$ et $\mathrm{Ker} f + \mathrm{Ker} g = E$.
\end{enumerate}
\finenonce{003332}



\finexercice


\exercice{3333, quercia, 2010/03/09}
\enonce{003333}{$\mathrm{Ker} f + \mathrm{Ker} g = \Im f + \Im g = E$}

Soient $E$ un ev de dimension finie et $f,g \in \mathcal{L}(E)$ tels que
$\mathrm{Ker} f + \mathrm{Ker} g = \Im f + \Im g = E$.

Montrer que les sommes sont directes.

\finenonce{003333}



\finexercice


\exercice{3334, quercia, 2010/03/09}
\enonce{003334}{$f^3 = 0$}

Soit $f \in \mathcal{L}(E)$ tel que $f^3 = 0$.

\begin{enumerate}
  \item  Montrer que $\mathrm{rg} f + \mathrm{rg} f^2 \le \dim E$.
  \item  Montrer que $2\mathrm{rg} f^2 \le \mathrm{rg} f$
     (appliquer le théorème du rang à $f_{|\Im f}$).
\end{enumerate}
\finenonce{003334}



\finexercice


\exercice{3335, quercia, 2010/03/09}
\enonce{003335}{$f\circ g = 0$ et $f+g \in GL(E)$}

Soit $E$ de dimension finie et $f,g \in \mathcal{L}(E)$ tels que :
$\begin{cases} f\circ g = 0 \cr f+g \in GL(E).\end{cases}$

Montrer que $\mathrm{rg} f + \mathrm{rg} g = \dim E$.

\finenonce{003335}



\finexercice


\exercice{3336, quercia, 2010/03/09}
\enonce{003336}{$f$ tq $\Im f$ et $\mathrm{Ker} f$ sont imposés}

Soit $E$ un $ K$-ev de dimension finie et $H,K$ deux sev fixés de $E$.

\begin{enumerate}
  \item A quelle condition existe-t-il un endomorphisme $f \in \mathcal{L}(E)$ tel que
     $\Im f = H \text{ et } \mathrm{Ker} f = K$ ?
  \item On note ${\cal E} = \{ f \in \mathcal{L}(E) \text{ tq } \Im f = H \text{ et } \mathrm{Ker} f = K \}$.
     Montrer que $\cal E$ est un groupe pour $\circ$ si et seulement si
     $H \oplus K = E$.
\end{enumerate}
\finenonce{003336}



\finexercice


\exercice{3337, quercia, 2010/03/09}
\enonce{003337}{Thms de factorisation}

Soient $E,F,G$ trois $ K$-ev avec $\dim(G)$ finie.

\begin{enumerate}
  \item Soient $u \in \mathcal{L}(F,E)$ et $v \in \mathcal{L}(G,E)$.
    Montrer qu'il existe $h \in \mathcal{L}({G,F})$
    tel que $v = u\circ h$ si et seulement si $\Im v \subset \Im u$.
  \item Soient $u \in \mathcal{L}(E,F)$ et $v \in \mathcal{L}(E,G)$.
    Montrer qu'il existe $h \in \mathcal{L}({G,F})$
    tel que $u = h\circ v$ si et seulement si $\mathrm{Ker} v \subset \mathrm{Ker} u$.
\end{enumerate}
\finenonce{003337}



\finexercice

\exercice{3349, quercia, 2010/03/09}
\enonce{003349}{Noyaux itérés}

Soit $E$ un ev de dimension finie et $f \in \mathcal{L}(E)$.
On pose $N_k = \mathrm{Ker}(f^k)$ et $I_k = \Im(f^k)$.

\begin{enumerate}
  \item Montrer que la suite $(N_k)$ est croissante (pour l'inclusion)
    et que la suite $(I_k)$ est décroissante.
  \item Soit $p$ tel que $N_p = N_{p+1}$. Justifier l'existence de $p$ et
    montrer que $N_{p+1} = N_{p+2} = \dots = N_{p+k} = \dots$
  \item Montrer que les suites $(N_k)$ et $(I_k)$ sont stationnaires à partir du
    même rang $p$.
  \item Montrer que $N_p \oplus I_p = E$.
  \item Montrer que la suite $(\dim(N_{k+1})-\dim(N_k))$ est décroissante.
\end{enumerate}
\finenonce{003349}



\finexercice


\exercice{3350, quercia, 2010/03/09}
\enonce{003350}{Dimension des $g$ tq $f\circ g = 0$ et/ou $g\circ f = 0$}

Soit $f \in \mathcal{L}(E)$.
On pose $K = \mathrm{Ker} f$, $I = \Im f$,
${\cal K} = \{ g \in \mathcal{L}(E) \text{ tq } f\circ g = 0 \}$ et
${\cal I} = \{ g \in \mathcal{L}(E) \text{ tq } g\circ f = 0 \}$.

\begin{enumerate}
  \item Montrer que ${\cal K}$ et ${\cal I}$ sont des sev de $\mathcal{L}(E)$.

  \item Soit $g \in \mathcal{L}(E)$.
    Montrer que : $g \in {\cal K} \iff \Im g \subset K$,
    et  : $g \in {\cal I} \iff \mathrm{Ker} g \supset I$.

  \item
  \begin{enumerate}
   \item Montrer que l' application
        $\Phi : {{\cal K}} \to {\mathcal{L}({E,K})}, g \mapsto {g^{|K}}$ est un isomorphisme
        d'ev. En déduire $\dim {\cal K}$.
   \item Chercher de même $\dim {\cal I}$ en introduisant un supplémentaire $I'$
        de $I$.
   \item Chercher aussi $\dim({\cal K} \cap {\cal I})$.

  \end{enumerate}
\end{enumerate}
\finenonce{003350}



\finexercice


\exercice{3351, quercia, 2010/03/09}
\enonce{003351}{Rang de $f  \mapsto u\circ f\circ v$}

Soient $u,v \in \mathcal{L}(E)$. Déterminer le rang de l'endomorphisme de $\mathcal{L}(E)$ :
$f  \mapsto u \circ f \circ v$.
\finenonce{003351}



\finexercice

\exercice{3353, quercia, 2010/03/09}
\enonce{003353}{Idéaux de $\mathcal{L}(E)$}

Un idéal à gauche de $\mathcal{L}(E)$ est un sev ${\cal I}$ de $\mathcal{L}(E)$ tel que :
$\forall\ f \in {\cal I},\ \forall\ g \in \mathcal{L}(E),\ f\circ g \in {\cal I}$.

Soit ${\cal I}$ un idéal à gauche.

\begin{enumerate}
  \item  Montrer que si $f \in {\cal I}$ et $\Im g \subset \Im f$,
     alors $g \in {\cal I}$.
  \item  Soient $f_1, f_2 \in {\cal I}$. Montrer qu'il existe $g_1, g_2 \in \mathcal{L}(E)$
     tels que $\Im(f_1\circ g_1 + f_2\circ g_2) = \Im f_1 + \Im f_2$.
  \item  Soit $f \in {\cal I}$ tel que rg($f$) soit maximal.
     Montrer que ${\cal I} = \{g \in \mathcal{L}(E) \text{ tq } \Im g \subset \Im f\}
                   = \{ f\circ g \text{ tq } g \in \mathcal{L}(E) \}$.
\end{enumerate}
\finenonce{003353}



\finexercice

\exercice{5171, rouget, 2010/06/30}
\enonce{005171}{**I}
Soit $E$ un $\Kk$-espace vectoriel et $f$ un élément de $\mathcal{L}(E)$.
\begin{enumerate}
\item  Montrer que $[\mbox{Ker}f=\mbox{Ker}f^2\Leftrightarrow\mbox{Ker}f\cap\mbox{Im}f=\{0\}]$ et $[\mbox{Im}f=\mbox{Im}f^2\Leftrightarrow
E=\mbox{Ker}f+\mbox{Im}f]$ (où $f^2=f\circ f)$.
\item  Par définition, un endomorphisme $p$ de $E$ est un projecteur si et seulement si $p^2=p$.

Montrer que
$$[p\;\mbox{projecteur}\Leftrightarrow Id-p\;\mbox{projecteur}]$$ puis que
$$[p\;\mbox{projecteur}\Rightarrow\mbox{Im}p=\mbox{Ker}(Id-p)\;\mbox{et}\;\mbox{Ker}p=\mbox{Im}(Id-p)\;\mbox{et}\;
E=\mbox{Ker}p\oplus\mbox{Im}p].$$

\item  Soient $p$ et $q$ deux projecteurs, montrer que~:~$[\mbox{Ker}p=\mbox{Ker}q\Leftrightarrow p=p\circ
q\;\mbox{et}\;q=q\circ p]$.
\item  $p$ et $q$ étant deux projecteurs vérifiant $p\circ q+q\circ p=0$, montrer que $p\circ q=q\circ p=0$. Donner
une condition nécessaire et suffisante pour que $p+q$ soit un projecteur lorsque $p$ et $q$ le sont. Dans ce
cas, déterminer $\mbox{Im}(p+q)$ et $\mbox{Ker}(p+q)$ en fonction de $\mbox{Ker}p$, $\mbox{Ker}q$, $\mbox{Im}p$ et
$\mbox{Im}q$.
\end{enumerate}
\finenonce{005171}


\finexercice


\exercice{5181, rouget, 2010/06/30}
\enonce{005181}{****}
Soit $E$ un $\Kk$-espace vectoriel et soit $(u,v)\in(\mathcal{L}(E))^2$.
\begin{enumerate}
\item  Montrer que $[\mbox{Ker}v\subset\mbox{Ker}u\Leftrightarrow\exists w\in\mathcal{L}(E)/\;u=w\circ v]$.
\item  En déduire que $[v\;\mbox{injectif}\Leftrightarrow\exists w\in\mathcal{L}(E)/\;w\circ v=Id_E]$.
\end{enumerate}
\finenonce{005181}


\finexercice


\exercice{5182, rouget, 2010/06/30}
\enonce{005182}{***}
Soit $E=\Rr[X]$ le $\Rr$-espace vectoriel des polynômes à coefficients réels.

\begin{itemize}
\item  Soit $\begin{array}[t]{cccc}f~:&E&\rightarrow&E\\ &P&\mapsto&P'\end{array}$. $f$ est-elle linéaire, injective,
surjective~?~Fournir un supplémentaire de $\mbox{Ker}f$.
\item  Mêmes questions avec $\begin{array}[t]{cccc}g~:&E&\rightarrow&E\\ &P&\mapsto&\int_{0}^{x}P(t)\;dt\end{array}$.
\end{itemize}
\finenonce{005182}


\finexercice

\exercice{5187, rouget, 2010/06/30}
\enonce{005187}{**T}
Soient $(e_i)_{1\leq i\leq 4}$ la base canonique de $\Rr^4$ et $f$ l'endomorphisme de $\Rr^4$ défini
par~:~$f(e_1)=2e_1+e_3$, $f(e_2)=-e_2+e_4$, $f(e_3)=e_1+2e_3$ et $f(e_4)=e_2-e_4$. Déterminer $\mbox{Ker}f$ et
$\mbox{Im}f$.

\finenonce{005187}


\finexercice


\exercice{5190, rouget, 2010/06/30}
\enonce{005190}{**I}
Soient $\Kk$ un sous-corps de $\Cc$ et $E$ et $F$ deux $\Kk$-espaces vectoriels de dimensions finies sur $\Kk$ et $u$ et $v$
deux applications linéaires de $E$ dans $F$. Montrer que ~:~
$|\mbox{rg}u-\mbox{rg}v|\leq\mbox{rg}(u+v)\leq\mbox{rg}u+\mbox{rg}v$.
\finenonce{005190}


\finexercice


\exercice{5191, rouget, 2010/06/30}
\enonce{005191}{****}
Soient $\Kk$ un sous-corps de $\Cc$ et $E$ un $\Kk$-espace vectoriel de dimension finie $n$.
\begin{enumerate}
\item  Montrer que, pour tout endomorphisme $f$ de $\Rr^2$, on a~:

\begin{center}
$(\mbox{Ker}f=\mbox{Im}f)\Leftrightarrow(f^2=0\;\mbox{et}\;n=2\mbox{rg}f)\Leftrightarrow(f^2=0\;\mbox{et}\;
\exists g\in\mathcal{L}(E)/\;f\circ g+g\circ f=Id_E)$.
\end{center}

\item  On suppose $\mbox{Ker}f=\mbox{Im}f$. Montrer qu'il existe une base $(u_1,...,u_p,v_1,...,v_p)$ de $E$ telle
que~:

\begin{center}
$\forall i\in\{1,...,p\},\;f(u_i)=0\;\mbox{et}\;f(v_i)=u_i$.
\end{center}
\end{enumerate}
\finenonce{005191}


\finexercice


\exercice{5192, rouget, 2010/06/30}
\enonce{005192}{***I Le théorème des noyaux itérés}
\label{exo:suprou10}
Soient $\Kk$ un sous-corps de $\Cc$, $E$ un $\Kk$-espace vectoriel de dimension finie $n$ et $f$ un endomorphisme de $E$ non injectif. Pour $k$
entier naturel donné, on pose $N_k=\mbox{Ker}f^k$ et $I_k=\mbox{Im}f^k$ (avec la convention $f^0=Id_E$).
\begin{enumerate}
\item  Montrer que~:~$\forall k\in\Nn,\;(N_k\subset N_{k+1}\;\mbox{et}\;I_{k+1}\subset I_k)$.
\item 
\begin{enumerate}
\item Montrer que~:~$(\forall k\in\Nn,\;(N_k=N_{k+1}\Rightarrow N_{k+1}=N_{k+2})$.
\item Montrer que~:~$\exists p\in\Nn/\;\forall k\in\Nn,\;(k<p\Rightarrow N_k\neq N_{k+1}\;\mbox{et}\;k\geq
p\Rightarrow N_k=N_{k+1})$.
\item Montrer que $p\leq n$.
\end{enumerate}
\item  Montrer que si $k<p$, $I_k=I_{k+1}$ et si $k\geq p$, $I_k=I_{k+1}$.
\item  Montrer que $E=I_p\oplus N_p$ et que $f$ induit un automorphisme de $I_p$.
\item  Soit $d_k=\mbox{dim}I_k$. Montrer que la suite $(d_k-d_{k+1})_{k\in\Nn}$ est décroissante (en d'autres
termes la suite des images itérées $I_k$ décroît de moins en moins vite).
\end{enumerate}
\finenonce{005192}


\finexercice


\exercice{5194, rouget, 2010/06/30}
\enonce{005194}{***I}
\label{exo:suprou12}
Soient $\Kk$ un sous-corps de $\Cc$, $E$ un $\Kk$-espace vectoriel de dimension quelconque sur $\Kk$ et $f$ un endomorphisme de $E$ vérifiant
$f^2-5f+6Id_E=0$. Montrer que $E=\mbox{Ker}(f-2Id)\oplus\mbox{Ker}(f-3Id)$.
\finenonce{005194}


\finexercice


\exercice{5582, rouget, 2010/10/16}
\enonce{005582}{***}
Soient $E$ un espace de dimension finie et $F$ et $G$ deux sous-espaces de $E$. Condition nécessaire et suffisante sur $F$ et $G$ pour qu'il existe un endomorphisme $f$ de $E$ tel que $F=\text{Ker}f$ et $G=\text{Im}f$.
\finenonce{005582}


\finexercice


\exercice{5583, rouget, 2010/10/16}
\enonce{005583}{***}
Soient $E$ un espace vectoriel non nul de dimension finie et $f$ un endomorphisme de $E$.

Montrer que : 
\begin{enumerate}
 \item  $(f\;\text{non injective})\Leftrightarrow(f=0\;\text{ou}\;f\;\text{diviseur de zéro à gauche})$.

 \item  $(f\;\text{non surjective})\Leftrightarrow(f=0\;\text{ou}\;f\;\text{diviseur de zéro à droite})$.
\end{enumerate}
\finenonce{005583}


\finexercice


\exercice{5586, rouget, 2010/10/16}
\enonce{005586}{**I Noyaux itérés}
Soient $E$ un espace vectoriel et $f$ un endomorphisme de $E$. Pour $k\in\Nn$, on pose $N_k= \text{Ker}(f^k)$ et $I_k=\text{Im}(f^k)$ puis $N=\displaystyle\bigcup_{k\in\Nn}N_k$ et $I=\displaystyle\bigcap_{k\in\Nn}I_k$. ($N$ est le nilespace de $f$ et $I$ le c\oe ur de $f$)

\begin{enumerate}
 \item  

 \begin{enumerate}
   \item Montrer que les suites $(N_k)_{k\in\Nn}$ et $(I_k)_{k\in\Nn}$ sont respectivement croissante et décroissante pour l'inclusion.
   \item Montrer que $N$ et $I$ sont stables par $f$.
   \item Montrer que $\forall k\in\Nn$, $(N_k =N_{k+1})\Rightarrow(N_{k+1}=N_{k+2})$.
  \end{enumerate}

\item  On suppose de plus que $\text{dim}E=n$ entier naturel non nul.
 \begin{enumerate}
   \item Soit $A=\{k\in\Nn/\;N_k=N_{k+1}\}$ et $B=\{k\in\Nn/\;I_k =I_{k+1}\}$. Montrer qu'il existe un entier $p\leqslant n$ tel que $A=B=\{k\in\Nn/\;k\geqslant p\}$.
   \item Montrer que $E=N_p\oplus I_p$.
   \item Montrer que $f_{/N}$ est nilpotent et que $f_{/I}\in GL(I)$.
  \end{enumerate}

\item  Trouver des exemples où 
 \begin{enumerate}
   \item $A$ est vide et $B$ est non vide,
   \item  $A$ est non vide et $B$ est vide,
   \item (****) $A$ et $B$ sont vides.
  \end{enumerate}

\item  Pour $k\in\Nn$, on pose $d_k=\text{dim}(I_k)$. Montrer que la suite $(d_k-d_{k+1})_{k\in\Nn}$ est décroissante.
\end{enumerate}
\finenonce{005586}


\finexercice


\exercice{5601, rouget, 2010/10/16}
\enonce{005601}{***}
Soient $E$ et $F$ deux $\Kk$-espaces vectoriels et $f$ une application linéaire de $E$ vers $F$.

\begin{enumerate}
 \item  Montrer que $\left[\left(\forall g\in\mathcal{L}(F,E),\;f\circ g\circ f=0\Rightarrow g = 0\right)\Rightarrow f\;\text{bijective}\right]$.

\item  On pose $\text{dim}E=p$, $\text{dim}F=n$ et $\text{rg}f=r$. Calculer la dimension de $\{g\in\mathcal{L}(F,E)/\;f\circ g\circ f =0\}$.
\end{enumerate}
\finenonce{005601}


\finexercice


\exercice{5602, rouget, 2010/10/16}
\enonce{005602}{**I}
Soit $E=\Kk_n[X]$. $u$ est l'endomorphisme de $E$ défini par : $\forall P\in E,\; u(P)=P(X+1)-P$.

\begin{enumerate}
 \item  Déterminer $\text{Ker}u$  et $\text{Im}u$.

\item  Déterminer explicitement une base dans laquelle la matrice de $u$ est  $\left(
\begin{array}{ccccc}
0&1&0&\ldots&0\\
\vdots&\ddots&\ddots&\ddots&\vdots\\
 & & &\ddots&0\\
\vdots& & &\ddots&1\\
0&\ldots& &\ldots&0
\end{array}
\right)$.
\end{enumerate}
\finenonce{005602}


\finexercice

\section{107.03 Morphismes particuliers}


\exercice{966, legall, 1998/09/01}

\enonce{000966}{}

Soient $  U   $ et $  V   $ deux ensembles non vides et $  f   $ une application de
$  U   $ \`a valeurs dans $  V  . $ Le {\em graphe} de $  f  $ est le sous-ensemble
de $  U\times V   $ d\' efini par $  \mathcal{G} _f=\{ (x, y) \in U\times V
 \hbox{ tels que } y=f(x) \}  .$
\begin{enumerate}
    \item On suppose maintenant que $  U   $ et $  V   $ sont des espaces vectoriels. Rappeler la d\' efinition de la structure d'espace vectoriel
de $  U\times V  .$
    \item Montrer qu'une partie $  H   $ de
$  U\times V  $ est le graphe d'une application lin\' eaire de $  U  $ dans $
V  $ si et seulement si les trois conditions qui suivent
sont satisfaites :

 {\em i)} La projection canonique $  H \rightarrow U   $ d\' efinie par $  (x,y) \mapsto x   $ est surjective.

{\em ii)} $  H   $ est un sous-espace vectoriel de $  U\times V  .$

{\em iii)} $  H \cap \left( \{ 0_U \}  )\times V \right) =\{ 0_{U\times V}\} .  $ ($  0_U   $
et $  0_{U\times V}  $ sont les \' el\' ements neutres respectifs de $  U   $ et $  U\times V  .)$

    \item On identifie $  { \Rr}^4  $ \`a $  { \Rr}^2\times { \Rr}^2  $ par l'isomorphisme
$  (x,y,z,t)\mapsto \left( (x,y),(z,t)\right)   .$ Enoncer des conditions n\' ec\' essaires et suffisantes
pour que $  E  $ soit le graphe d'une application lin\' eaire de $  { \Rr}^2  $ dans lui-m\^eme.
    \item Montrer que $  E  $ est le graphe d'une application lin\' eaire $  \varphi   $
de $  { \Rr}^2  $ dans lui-m\^eme. D\' eterminer sa matrice dans une base que l'on d\' efinira au  pr\' ealabe.
\end{enumerate}
\finenonce{000966}



\finexercice



\exercice{967, legall, 1998/09/01}

\enonce{000967}{Projecteur et involution}
Soit $E$ un espace vectoriel; on note $i_E$ l'identit\'e sur $E$. Un
endomorphisme $u$ de $E$ est un \textbf{projecteur} si $u\circ u=u$.
\begin{enumerate}
    \item Montrer que si $u$ est un projecteur alors $i_E-u$ est un
projecteur. V\'erifier aussi que  $\text{Im} u=\{x\in E;\;\;u(x)=x\}$ et que
$E=\text{Ker} u\oplus \text{Im} u$.\\
Un endomorphisme $u$ de $E$ est appel\'e \textit{involutif} si $u\circ
u=i_E$.
    \item Montrer que si $u$ est involutif alors $u$ est bijectif et
$E=\text{Im}(i_E+u)\oplus \text{Im}(i_E-u)$.\\
     Soit $E=F\oplus G$ et soit $x\in E$ qui s'\'ecrit donc de fa\c{c}on
unique $x=f+g$, $f\in F$, $g\in $G. Soit $u:E\ni x\mapsto f-g\in E$.
    \item Montrer que $u$ est involutif, $F=\{x\in E;\;u(x)=x\}$ et
$G=\{x\in E;\;u(x)=-x\}$.
    \item Montrer que si $u$ est un projecteur, $2u-i_E$ est involutif
et que tout endomorphisme involutif peut se mettre sous cette forme.
\end{enumerate}
\finenonce{000967}



\finexercice



\exercice{968, legall, 1998/09/01}

\enonce{000968}{}
Soient $ P=\{ (x, y, z)\in \R ^3  ;  2x+y-z=0\}  $
et $ D=\{ (x, y, z)\in \R ^3  ;  2x-2y+z=0 ,  x-y-z=0\}.$
 On  d\'esigne par $ \epsilon  $ la base canonique de $ \R ^3 .$
\vskip1mm
\begin{enumerate}
    \item  Donner une base $ \{ e_1, e_2\} $ de $ P $ et $ \{ e_3\}
$ une base de $ D .$ Montrer
 que $ \R ^3=P\oplus D $ puis que $ \epsilon '=\{ e_1,e_2,e_3\}
 $ est une base de $ \R ^3 .$
    \item Soit $ p $ la projection de $ \R ^3 $ sur $ P $
parall\'element \`a $ D .$
D\'eterminer $ \hbox{Mat}(p, \epsilon ',\epsilon ') $ puis
$ A=\hbox{Mat}(p, \epsilon ,\epsilon ) .$ V\'erifier $ A^2=A .$
    \item Soit $ s $ la sym\'etrie de $ \R ^3 $ par rapport \`a $ P $
parall\'element \`a $ D .$
D\'eterminer $ \hbox{Mat}(s, \epsilon ',\epsilon ') $ puis
$ B=\hbox{Mat}(s, \epsilon ,\epsilon ) .$ V\'erifier
$ B^2=I ,$ $ AB=A $ et $ BA=A .$
\end{enumerate}
\finenonce{000968}



\finexercice



\exercice{969, legall, 1998/09/01}

\enonce{000969}{}
\begin{enumerate}
    \item Soit  $E$  un espace vectoriel de
dimension  $n$. Un {\em hyperplan} de  $E$ est un sous-espace vectoriel de
dimension  $n-1$.
Montrer que l'intersection de deux hyperplans de  $E$  a une dimension
sup\' erieure ou
\' egale \`a  $n-2$. Montrer que, pour tout  $p\leq n$, l'intersection de
$p$  hyperplans
a une dimension sup\' erieure ou \' egale \`a  $n-p$.
    \item Montrer que, pour tout  $n\in {\Nn}$  et pour
tout  $y\in {\Rr}$,  l'application  $e_{y}$  de  ${\Rr}_n[X]$  \`a
valeurs dans  ${\Rr}$  d\' efinie en posant  $e_y(P(X))=P(y)$  ( i.e.
l'application  $e_y$  est l'\' evaluation en  $y$)  est lin\' eaire.
Calculer la dimension de son noyau.

    \item M\^eme question avec l'application  $e'_{y}$  de
 ${\Rr}_n[X]$  \`a valeurs
dans  ${\Rr}$  d\' efinie en posant  $e'_y(P(X))=P'(y)$  (en d\' esignant par  $P'$  le
polyn\^ome d\' eriv\' e de  $P$).
    \item D\' emontrer, \`a l'aide de ces deux r\' esultats,
qu'il existe dans
${\Rr}_6[X]$  un  polyn\^ome  $P$  non nul et ayant les propri\' et\' es
suivantes :  $P(0)=P(1)=P(2)=0$
et $ P'(4) =P'(5)=P'(6)=0 .$
\end{enumerate}
\finenonce{000969}



\finexercice



\exercice{970, ridde, 1999/11/01}

\enonce{000970}{}
Soit ${f} : {\Rr^2} \to {\Rr^2}$, ${ (x, y)}\mapsto {\frac 13 (-x + 2y, -2x + 4y)}$.
Montrer que $f$ est la b\^\i\^\i \^\i \^\i p par rapport \`a
b\^\i\^\i \^\i \^\i p parall\`element \`a b\^\i\^\i \^\i \^\i p.
\finenonce{000970}



\finexercice



\exercice{971, gourio, 2001/09/01}

\enonce{000971}{}
$E$ est un $\Rr-$espace vectoriel, $F$ et $G$ deux sous-espaces
suppl\'{e}mentaires de $E$: $E=F\bigoplus G.$ On pose $s(u)=u_{F}-u_{G} $
o\`{u} $u=u_{F}+u_{G}$ est la d\'{e}composition (unique) obtenue gr\^{a}ce
\`{a} $E=F\bigoplus G.$ $s$ est la sym\'{e}trie par-rapport \`{a} $F$ de
direction $G.$
\begin{enumerate}
\item Montrer que $s\in L(E),$ que $u\in F\Leftrightarrow s(u)=u,u\in
G\Leftrightarrow s(u)=-u,$ donner $\mathrm{Ker}(s)$ et calculer $s^{2}.$
\item
R\'{e}ciproquement si $f\in L(E)$ v\'{e}rifie $f^{2}=id_{E}.$ On pose
$p=\frac{f+id_{E}}{2}.$ Calculer $f(u)$ en fonction de $p(u)$ et $u.$
V\'{e}rifier que $p$ est un projecteur, calculer son noyau et son image.
Montrer que $f $ est la sym\'{e}trie par rapport \`{a} $F=\{u\in E|f(u)=u\}$
de direction $ G=\{u\in E|f(u)=-u\}.$
\end{enumerate}
\finenonce{000971}



\finexercice



\exercice{972, gourio, 2001/09/01}

\enonce{000972}{}
Soient $p$ et $q$ deux projecteurs de $E$, espace vectoriel, tels que $pq=qp$
($p$ et $q$ commutent). Montrer que $pq$ et $(p+q-pq)$ sont deux projecteurs de
$E$, et que :
$$\mathrm{Im} (pq)=\mathrm{Im} p\cap \mathrm{Im} q,$$
$$\mathrm{Im} (p+q-pq)=\mathrm{Im} p+\mathrm{Im} q.$$
\finenonce{000972}



\finexercice



\exercice{973, gourio, 2001/09/01}

\enonce{000973}{}
Soient $p$ et $q$ deux projecteurs de $E$, espace vectoriel; donner une condition
n\'ecessaire et suffisante pour que $p+q$ soit un projecteur de $E$ ;
donner alors $\mathrm{Im} (p+q)$ et $\mathrm{Ker}(p+q).$

\emph{Indication} : on montrera que $\mathrm{Im} (p+q)=\mathrm{Im} p\bigoplus \mathrm{Im} q$
et que   $\mathrm{Ker}(p+q)=\mathrm{Ker}(p)\cap \mathrm{Ker}(q).$
\finenonce{000973}



\finexercice



\exercice{974, gourio, 2001/09/01}
\video{6N7D6lPLPHc}
\enonce{000974}{}
Soit $E$ l'espace vectoriel des fonctions de $\Rr$ dans $\Rr$. Soient $P$ le
sous-espace des fonctions paires et $I$ le sous-espace des fonctions
impaires. Montrer que $E=P\bigoplus I$. Donner l'expression du projecteur sur
$P$ de direction $I$.

\finenonce{000974} 


\finexercice


\exercice{975, gourio, 2001/09/01}

\enonce{000975}{}
Soit $E=\Rr[X]$ l'espace vectoriel des polyn\^{o}mes, et $f:E\rightarrow E$
d\'{e}finie par :
$$\forall P\in E,\text{ }f(P)(X)=\frac{P(-X)-P(X)}{2}.$$

Montrer que $f\in L(E)$, que $E=\mathrm{Im} f\bigoplus \mathrm{Ker}(f)$ mais que
$f^{2}=-f.$ Quel th\'{e}or\`{e}me cet exemple illustre t-il ?
\finenonce{000975}



\finexercice



\exercice{976, gourio, 2001/09/01}
\video{yivtpWQkPFg}
\enonce{000976}{}
Soit $E={\Rr}_{n}[X]$ l'espace vectoriel des polyn\^{o}mes de degr\'{e} $%
\leq n$, et $f:E\rightarrow E$ d\'{e}finie par:
$$f(P)=P+(1-X)P'. $$
Montrer que $f$ est une application linéaire et donner une base de $\Im f$ et de $\ker f.$
\finenonce{000976} 


\finexercice


\exercice{977, gourio, 2001/09/01}

\enonce{000977}{}
Soit $E=C({\Rr}^{+},{\Rr})$ et  $U : E \rightarrow E$ d\'efinie par $f\mapsto U(f)$ telle que :
$$\forall x\in {\Rr}^{+*},U(f)(x)=\frac{1}{x}\int_{0}^{x}f(t)dt. $$
et $U(f)(0)=f(0).$
Montrer que $U\in L(E),$ d\'{e}terminer $\mathrm{Ker}(U)$ et $\mathrm{Im} (U).$
\finenonce{000977}



\finexercice



\exercice{978, cousquer, 2003/10/01}

\enonce{000978}{}
On désigne par $\mathcal{P}_q$ l'espace vectoriel des polynômes à
coefficients réels de degré inférieur ou égal à $q$, et $\mathcal{O}_q$
l'espace vectoriel des polynômes d'ordre supérieur ou égal à $q$,
c'est-à-dire divisibles par $x^q$. $P$ étant un polynôme, on note $T(P)$ le
polynôme défini par~: 
$$T(P)(x)=xP(0)-{1\over20}x^5P^{(4)}(0)+
\int_0^x t^2[P(t+1)-P(t)-P'(t)]\,dt.$$
\begin{enumerate}
\item  Montrer que $T$ est linéaire. Déterminer $T(e_i)$ où
$e_0=1$, $e_1=x$, $e_2=x^2$, $e_3=x^3$, $e_4=x^4$, et vérifier que
$T(\mathcal{P}_4) \subset\mathcal{P}_4$. Désormais, on considère $T$ comme application
linéaire de $\mathcal{P}_4$ dans $\mathcal{P}_4$. Écrire sa matrice par 
rapport à la base $(e_0, e_1, e_2, e_3, e_4)$.

\item  Déterminer soigneusement les espaces $T(\mathcal{P}_4\cap\mathcal{O}_3)$ et
$T(\mathcal{P}_4\cap\mathcal{O}_2)$.

\item  La restriction $T'$ de $T$ à $\mathcal{P}_4\cap\mathcal{O}_2$ est-elle injective~?
Sinon déterminer une base du noyau de $T'$.

\item  Montrer que $\Im T = (\mathcal{O}_1\cap\mathcal{P}_1)\oplus(\mathcal{O}_3\cap\mathcal{P}_4)$. Quel est le 
rang de $T$~?

\item  Montrer que $\mathrm{Ker} T$ peut s'écrire sous la forme
$(\mathcal{O}_1\cap\mathcal{P}_1)\oplus V$~; expliciter un sous-espace $V$ possible.
Déterminer $\mathrm{Ker} T\cap\Im T$.

\item  On cherche un vecteur non nul $u=ae_3+be_4$ de $\mathcal{O}_3\cap\mathcal{P}_4$, et un
nombre réel~$\lambda$, tels que $T(u)=\lambda u$.
Écrire les équations que doivent vérifier $a,b,\lambda$.
Montrer qu'il existe deux valeurs possibles de $\lambda$,
$\lambda_1$ et $\lambda_2$, telles $0<\lambda_1<\lambda_2$~; les
calculer. Trouver deux vecteurs non nuls $u_3$ et $u_4$ de
$\mathcal{O}_3\cap\mathcal{P}_4$ tels que $T(u_3)=\lambda_1u_3$ et
$T(u_4)=\lambda_2u_4$.

\item  On pose $u_0=e_1$, $u_1=e_2-4e_3+3e_4$, $u_2=e_0$. Montrer que
$\{u_0,u_1,u_2,u_3,u_4\}$ est une base de $\mathcal{P}_4$. Écrire la matrice
de $T$ dans cette base.
\end{enumerate}
\finenonce{000978}



\finexercice



\exercice{2432, matexo1, 2002/02/01}

\enonce{002432}{}
Parmi les applications suivantes, lesquelles sont des
endomorphismes de $C^\infty(\R)$ ($\phi \in C^\infty(\R)$ est
fix\'e)\,: 
$$ f \mapsto f + \phi, \quad f \mapsto \phi f, \quad
f \mapsto f\circ\phi, \quad f\mapsto \phi\circ f, \quad
f \mapsto \int{f}, \quad f \mapsto f'. $$
Lesquelles sont des endomorphismes de $C^0(\R)$\,?
\par Pour quelles valeurs de $\phi$ les endomorphismes 
$\Phi : f \mapsto f\circ\phi$ et $D : f \mapsto f'$ commutent-ils
(c'est-\`a-dire v\'erifient $D(\Phi f) = \Phi(D f), \forall f$)\,?
\finenonce{002432}


\finexercice


\exercice{3306, quercia, 2010/03/09}
\enonce{003306}{Image d'une somme, d'une intersection}

Soit $f : E \to F$ une application linéaire et $E_1$, $E_2$ deux
sous-espaces vectoriels de $E$, $F_1$, $F_2$ deux sous-espaces vectoriels
de $F$.
Que pouvez-vous-dire de $f(E_1+E_2)$, $f(E_1\cap E_2)$, $f^{-1}(F_1+F_2)$, $f^{-1}(F_1\cap F_2)$ ?
\finenonce{003306}



\finexercice


\exercice{3307, quercia, 2010/03/09}
\enonce{003307}{Effet sur les familles libres et génératrices}

Soient $E,F$ deux espaces vectoriels et $f : E \to F$ linéaire.

\begin{enumerate}
  \item Montrer que $f$ est injective si et seulement si $f$ transforme
    toute famille libre de $E$ en une famille libre de~$F$.
  \item Montrer que $f$ est surjective si et seulement s'il existe une famille
    génératrice de $E$ transformée par $f$ en une famille génératrice de $F$.
\end{enumerate}
\finenonce{003307}



\finexercice


\exercice{3308, quercia, 2010/03/09}
\enonce{003308}{$f(\mathrm{Ker} (g\circ f))$}

Soit $E$ un espace vectoriel et $f,g \in \mathcal{L}(E)$.
Montrer que $f(\mathrm{Ker} (g\circ f)) = \mathrm{Ker} g \cap \Im f$.
\finenonce{003308}



\finexercice

\exercice{3314, quercia, 2010/03/09}
\enonce{003314}{Permutation de coordonnées dans $ K^n$}

Soit $\sigma \in S_n$ (groupe symétrique) et
${f_\sigma} : { K^n} \to { K^n}, {(x_1,\dots x_n)} \mapsto
              {(x_{\sigma(1)},\dots,x_{\sigma(n)})}$

On munit $ K^n$ de la structure d'algèbre pour les opérations
composante par composante.

\begin{enumerate}
  \item Montrer que $f_{\sigma}$ est un automorphisme d'algèbre.

  \item Soit $\varphi$ un automorphisme d'algèbre de $ K^n$.
  \begin{enumerate}
    \item  Montrer que la base canonique de $ K^n$ est invariante par $\varphi$
         (étudier $\varphi(e_i^2)$ et $\varphi(e_i\times e_j)$).
    \item  En déduire qu'il existe $\sigma \in S_n$ tel que $\varphi = f_\sigma$.
  \end{enumerate}
  \item Montrer que $\{0\}$, $ K(1,\dots, 1)$,
        $\{(x_1,\dots,x_n) \text{ tq } x_1 + \dots + x_n = 0 \}$
        et $ K^n$ sont les seuls sev stables par tous les endomorphismes
        $f_\sigma$.
\end{enumerate}
\finenonce{003314}



\finexercice


\exercice{3338, quercia, 2010/03/09}
\enonce{003338}{Isomorphisme $\circ$ projecteur}

Soient $E$ en ev de dimension finie et $f \in \mathcal{L}(E)$.

\begin{enumerate}
  \item Montrer qu'il existe un projecteur $p\in \mathcal{L}(E)$
    et un isomorphisme $g\in GL(E)$ tels que $f = g\circ p$.
  \item Montrer qu'il existe un projecteur $p\in \mathcal{L}(E)$
    et un isomorphisme $g\in GL(E)$ tels que $f = p\circ g$.
\end{enumerate}
\finenonce{003338}



\finexercice


\exercice{3339, quercia, 2010/03/09}
\enonce{003339}{Centre de $\mathcal{L}(E)$}

Soit $E$ un $ K$-ev de dimension finie.
Le centre de $\mathcal{L}(E)$ est :
$Z = \{f \in \mathcal{L}(E) \text{ tq } \forall\ g \in \mathcal{L}(E),\ f\circ g = g\circ f\}$.

\begin{enumerate}
  \item Soit $f \in \mathcal{L}(E)$ et $\vec x \in E$. Si $(\vec x,f(\vec x))$ est libre,
    montrer qu'il existe $g \in \mathcal{L}(E)$ telle que $g(\vec x) = \vec x$ et
    $g\circ f(\vec x) = -f(\vec x)$.
  \item En déduire que $Z$ est l'ensemble des homothéties.
  \item Déterminer $Z' =
    \{f \in \mathcal{L}(E) \text{ tq } \forall\ g \in GL(E),\ f\circ g = g\circ f\}.$
\end{enumerate}
\finenonce{003339}



\finexercice


\exercice{3340, quercia, 2010/03/09}
\enonce{003340}{\'Eléments réguliers dans $\mathcal{L}(E)$}
Soit $f \in \mathcal{L}(E,F)$.

\begin{enumerate}
  \item Montrer que :  ($f$ est \rlap{injectif)}\phantom{surjectif)}
    $\iff$ ($\forall\ g \in \mathcal{L}(E),\ f\circ g = 0  \Rightarrow  g = 0$).
  \item Montrer que :  ($f$ est surjectif)
    $\iff$ ($\forall\ g \in \mathcal{L}(F),\ g\circ f = 0  \Rightarrow  g = 0$).

\end{enumerate}
\finenonce{003340}



\finexercice


\exercice{3341, quercia, 2010/03/09}
\enonce{003341}{$f^2 = -\mathrm{id}$}

Soit $E$ un $\R$-ev et $f \in \mathcal{L}(E)$ tel que $f\circ f = -\mathrm{id}_E$.
Pour $z = x + iy \in \C$ et $\vec u \in E$, on pose :
$z\vec u = x\vec u + yf(\vec u)$.

\begin{enumerate}
  \item Montrer qu'on définit ainsi une structure de $\C$-ev sur $E$.
  \item En déduire que $\dim_{\R}(E)$ est paire.
\end{enumerate}
\finenonce{003341}



\finexercice


\exercice{3342, quercia, 2010/03/09}
\enonce{003342}{$f\circ f = 0$ et $f\circ g + g\circ f = \mathrm{id}$}
\begin{enumerate}
  \item Soit $E$ un $ K$-ev et $f,g \in \mathcal{L}(E)$ tels que :
    $\begin{cases} f^2 = 0 \cr f\circ g + g\circ f = \mathrm{id}_E.\end{cases}$ \par
    Montrer que $\mathrm{Ker} f = \Im f$.

  \item Réciproquement, soit $f \in \mathcal{L}(E)$ tel que $\mathrm{Ker} f = \Im f$, et $F$ un supplémentaire de $\mathrm{Ker} f$.
   Montrer que
  \begin{enumerate}
   \item $f^2 = 0$.
   \item $\forall\ \vec x \in E$, il existe $\vec y,\vec z \in F$ uniques tels
       que $\vec x = \vec y + f(\vec z)$.
   \item Il existe $g \in \mathcal{L}(E)$ tel que $f\circ g + g\circ f = \mathrm{id}_E$.
  \end{enumerate}
\end{enumerate}
\finenonce{003342}



\finexercice


\exercice{3343, quercia, 2010/03/09}
\enonce{003343}{Endomorphisme nilpotent}

Un endomorphisme $f \in \mathcal{L}(E)$ est dit {\it nilpotent\/} s'il existe
$p \in \N$ tel que $f^p = 0$. Dans ce cas, {\it l'indice\/} de $f$ est le
plus petit entier $p$ tel que $f^p = 0$.
On considère $f\in\mathcal{L}(E)$ nilpotent d'indice $p$.

\begin{enumerate}
  \item Soit $\vec u \in E \setminus \mathrm{Ker} f^{p-1}$. Montrer que la famille
    $\bigl(\vec u, f(\vec u), \dots, f^{p-1}(\vec u) \bigr)$ est libre.
  \item En déduire que si $E$ est de dimension finie $n$, alors $f^n = 0$.
  \item Soit $g \in GL(E)$ tel que $f\circ g = g\circ f$.
    Montrer que $f + g \in GL(E) \dots$
  \begin{enumerate}
    \item en dimension finie.
    \item pour $E$ quelconque.
  \end{enumerate}
  \item Dans $\mathcal{L}({ K^2})$, soient $f,g$ de matrices :
     $\left(\begin{smallmatrix}0&0\cr1&0\end{smallmatrix}\right)$ et $\left(\begin{smallmatrix}0&1\cr-1&0\end{smallmatrix}\right)$.
     Vérifier que $f$ est nilpotent, $g \in GL( K^2)$, mais ${f+g \notin GL( K^2)}$.
\end{enumerate}
\finenonce{003343}



\finexercice


\exercice{3344, quercia, 2010/03/09}
\enonce{003344}{Matexo}

Soit $E$ un $K$ espace vectoriel de dimension finie et
$f\in\mathcal{L}(E)$ tel que
$\forall\ x\in E,\ \exists p_x\in\N^*,\ f^{p_x}(x) = \vec 0$.
Montrer que $f$ est nilpotent.
Donner un contre-exemple en dimension infinie.
\finenonce{003344}



\finexercice


\exercice{3345, quercia, 2010/03/09}
\enonce{003345}{Mines P' 1995}

Soit $E$ un $ K$-espace vectoriel de dimension finie et $f\in \mathcal{L}(E)$ nilpotente
d'indice~$n$.

Soit $\phi : {\mathcal{L}(E)} \to {\mathcal{L}(E)}, g \mapsto{f\circ g - g\circ f.}$
\begin{enumerate}
  \item Montrer que $\phi^p(g) = \sum_{k=0}^p (-1)^kC_p^k f^{p-k}\circ g \circ f^k$.
    En déduire que $\phi$ est nilpotente.

  \item Soit $a\in \mathcal{L}(E)$. Montrer qu'il existe $b\in \mathcal{L}(E)$ tel que
    $a\circ b \circ a = a$. En déduire l'indice de nilpotence de~$\phi$.
\end{enumerate}
\finenonce{003345}



\finexercice


\exercice{3346, quercia, 2010/03/09}
\enonce{003346}{Endomorphisme cyclique}

Soit $E$ un ev de dimension $n$ et
$f \in \mathcal{L}(E)$. On suppose qu'il existe un vecteur $\vec u \in E$ tel
que la famille $\bigl(f^k(\vec u)\bigr)_{k \in \N}$ engendre~$E$.

\begin{enumerate}
  \item  Montrer que $\bigl(\vec u, f(\vec u), \dots, f^{n-1}(\vec u)\bigr)$
     est une base de $E$.
     (Considérer $p$ maximal tel que
     ${\cal F} = \bigl(\vec u,\dots, f^{p-1}(\vec u)\bigr)$
     est libre, et prouver que $f^k(\vec u)$ est combinaison linéaire de
     $\cal F$ pour tout entier $k$)
  \item  Montrer qu'un endomorphisme $g \in \mathcal{L}(E)$ commute avec $f$ si et seulement
     si c'est un polynôme en $f$.
\end{enumerate}
\finenonce{003346}



\finexercice


\exercice{3347, quercia, 2010/03/09}
\enonce{003347}{$u^2 = 0$ en dimension 3}

Soit $E$ un ev de dimension 3 et $u\in \mathcal{L}(E)$ tel que $u^2 = 0$.
Montrer qu'il existe $f \in E^*$ et $\vec a \in E$ tels que :
$\forall\ \vec x \in E,\ u(\vec x\,) = f(\vec x\,)\vec a$.

\finenonce{003347}



\finexercice


\exercice{3348, quercia, 2010/03/09}
\enonce{003348}{$(u,x,f(x))$ liée}

Soit $E$ un ev de dimension supérieure ou égale à 3 et
$\vec u \in E \setminus \{\vec 0\}$.
Trouver tous les endomorphismes $f \in \mathcal{L}(E)$ tels que :
$\forall\ \vec x \in E,\text{ la famille }(\vec u, \vec x, f(\vec x))
 \text{ est liée.}$
\finenonce{003348}



\finexercice

\exercice{3354, quercia, 2010/03/09}
\enonce{003354}{Automorphismes de $\mathcal{L}(E)$}

Soit $E$ un ev de dimension $n$ et
$\Phi : {\mathcal{L}(E)} \to {\mathcal{L}(E)}$ un automorphisme d'algèbre. On note
$(\vec e_1, \dots, \vec e_n)$ une base fixée de $E$, $(\varphi_{ij})$ la
base de $\mathcal{L}(E)$ associée
$\bigl(\varphi_{ij}(\vec e_k) = \delta_{jk}\vec e_i\bigr)$ et
$\psi_{ij} = \Phi(\varphi_{ij})$.

\begin{enumerate}
  \item  Simplifier $\psi_{ij} \circ \psi_{k\ell}$.
     
  \item  En déduire qu'il existe $\vec u_1 \in E\setminus\{\vec 0\}$ tel que
     $\psi_{11}(\vec u_1) = \vec u_1$.
     
  \item  On note $\vec u_i = \psi_{i1}(\vec u_1)$.
     Montrer que $\psi_{ij}(\vec u_k) = \delta_{jk}\vec u_i$ et en déduire que
     $(\vec u_i)$ est une base de $E$.
     
  \item  Soit $f \in GL(E)$ définie par : $f(\vec e_i) = \vec u_i$.
     Montrer que : $\forall\ g \in \mathcal{L}(E),\ \Phi(g) = f\circ g \circ f^{-1}$.
     
\end{enumerate}
\finenonce{003354}



\finexercice


\exercice{3355, quercia, 2010/03/09}
\enonce{003355}{$f^2 = 0  \Rightarrow  f=g\circ h$ avec $h\circ g=0$}

Soit $f\in \mathcal{L}(E)$ telle que $f^2 = 0$.
Montrer qu'il existe $g,h \in \mathcal{L}(E)$ tels que $f=g\circ h$ et $h\circ g = 0$.

\finenonce{003355}



\finexercice

\exercice{3485, quercia, 2010/03/10}
\enonce{003485}{Barycentre de projections}

Soient $p,q$ deux projections de même base $H$ et de directions $F$,$G$.
Soit $\lambda \in  K$. Montrer que $\lambda p + (1 - \lambda)q$ est encore une projection de base $H$.
\finenonce{003485}



\finexercice


\exercice{3486, quercia, 2010/03/10}
\enonce{003486}{Valeurs propres d'une projection}

Soit $E$ un espace vectoriel et $p \in \mathcal{L}(E)$ une projection.
Montrer que pour tout $\lambda \in  K \setminus \{-1\}$, $\mathrm{id}_E + \lambda p$ est
un isomorphisme de $E$.
\finenonce{003486}



\finexercice


\exercice{3487, quercia, 2010/03/10}
\enonce{003487}{Projections ayant même base ou même direction}

Soit $E$ un espace vectoriel et $p,q \in \mathcal{L}(E)$ deux projections.

\begin{enumerate}
  \item Montrer que $p$ et $q$ ont même base si et seulement si :
    $p\circ q = q$ et $q \circ p = p$.
  \item Donner une condition analogue pour que $p$ et $q$ aient même direction.

\end{enumerate}
\finenonce{003487}



\finexercice


\exercice{3488, quercia, 2010/03/10}
\enonce{003488}{Somme de deux projecteurs}

Soient $p,q$ deux projections. Montrer les équivalences :

$$p+q \text{ est une projection } \Leftrightarrow
p \circ q + q \circ p = 0         \Leftrightarrow
\begin{cases} \text{Base}(p) \subset \text{Dir}(q) \cr
        \text{Base}(q) \subset \text{Dir}(p).\cr \end{cases}$$

Chercher alors la base et la direction de $p+q$.
\finenonce{003488}



\finexercice


\exercice{3489, quercia, 2010/03/10}
\enonce{003489}{$f\circ g = f$ et $g\circ f = g$}

Soit $E$ un $ K$-ev. Trouver tous les couples $(f,g)$ d'endomorphismes
de $E$ tels que : $\begin{cases}f \circ g = f \cr g\circ f = g. \cr \end{cases}$

\finenonce{003489}



\finexercice


\exercice{3490, quercia, 2010/03/10}
\enonce{003490}{$f\circ g = \mathrm{id}$}

Soit $E$ un espace vectoriel et $f,g \in \mathcal{L}(E)$ tels que $f\circ g = \mathrm{id}_E$.
Montrer que $g\circ f$ est une projection et déterminer ses éléments.
\finenonce{003490}



\finexercice


\exercice{3491, quercia, 2010/03/10}
\enonce{003491}{Projection $p + q - q\circ p$}

Soient $p,q$ deux projections telles que $p \circ q = 0$.
Montrer que $p + q - q \circ p$ est une projection, et déterminer ses éléments.


\finenonce{003491}



\finexercice


\exercice{3492, quercia, 2010/03/10}
\enonce{003492}{Endomorphisme de rang 1}

Soit $f \in \mathcal{L}(E)$ de rang 1. Montrer qu'il existe un unique
$\lambda \in  K$ tel que $f^2 = \lambda f$.

Montrer que : $\lambda = 1 \iff \mathrm{id} - f$ est non injective
                  $\iff \mathrm{id} - f$ est non surjective
(même en dimension infinie).

\finenonce{003492}



\finexercice


\exercice{3493, quercia, 2010/03/10}
\enonce{003493}{Relation d'ordre sur les projecteurs}

On munit l'ensemble des projections d'un ev $E$ de la relation :
$p \ll q \iff p \circ q = q \circ p = p$.

\begin{enumerate}
  \item   Montrer que c'est une relation d'ordre.
  \item   Soient $p,q$ deux projections permutables.
      Montrer que $\sup(p,q) = p+q-p \circ q$ et $\inf(p,q) = p \circ q$.
\end{enumerate}
\finenonce{003493}



\finexercice


\exercice{3494, quercia, 2010/03/10}
\enonce{003494}{Expressions analytiques}

Soit $E =  K^3$, $F = \{ \vec X = (x,y,z) \text{ tq }  x + 2y + z = 0 \}$
et $G = \text{vect}\bigl(\vec U = (1,1,1)\bigr)$.

\begin{enumerate}
  \item Vérifier que $F \oplus G = E$.
  \item Soit $s$ la symétrie de base $F$ de direction $G$ et $\vec X = (x,y,z)$.
    Déterminer $s(\vec X)$.


\end{enumerate}
\finenonce{003494}



\finexercice


\exercice{3495, quercia, 2010/03/10}
\enonce{003495}{Trace nulle}
Soit $E$ un $\R$-ev de dimension finie et $A$ une partie finie de $GL(E)$
stable par composition. On pose $u = \sum_{f\in A} f$. Montrer que
$\mathrm{tr}(u)=0  \Rightarrow  u=0$.

\finenonce{003495}



\finexercice

\exercice{5193, rouget, 2010/06/30}
\enonce{005193}{***I}
Soient $\Kk$ un sous-corps de $\Cc$ et $E$ un $\Kk$-espace vectoriel de dimension finie notée $n$. Soit $u$ un endomorphisme de $E$.
On dit que $u$ est nilpotent si et seulement si $\exists k\in\Nn^*/\;u^k=0$ et on appelle alors indice de
nilpotence de $u$ le plus petit de ces entiers $k$ (par exemple, le seul endomorphisme $u$, nilpotent d'indice $1$ est
$0$).
\begin{enumerate}
\item  Soit $u$ un endomorphisme nilpotent d'indice $p$. Montrer qu'il existe un vecteur $x$ de $E$ tel que la
famille

$(x,\;u(x),...,\;u^{p-1}(x))$ soit libre.
\item  Soit $u$ un endomorphisme nilpotent. Montrer que $u^n=0$.
\item  On suppose dans cette question que $u$ est nilpotent d'indice $n$. Déterminer $\mbox{rg}u$.
\end{enumerate}
\finenonce{005193}


\finexercice


\exercice{5584, rouget, 2010/10/16}
\enonce{005584}{*** I}
Soient $E$ un espace de dimension finie $n$ non nulle et $f$ un endomorphisme nilpotent de $E$. Montrer que $f^n=0$.
\finenonce{005584}


\finexercice

\exercice{5587, rouget, 2010/10/16}
\enonce{005587}{***I}
\label{ex:rou25}
Soit $E$ un espace vectoriel non nul. Soit $f$ un endomorphisme de $E$ tel que pour tout vecteur $x$ de $E$  la famille $(x,f(x))$ soit liée. Montrer que $f$ est une homothétie.
\finenonce{005587}


\finexercice


\exercice{5588, rouget, 2010/10/16}
\enonce{005588}{***I}
Soit $E$ un espace de dimension finie. Trouver les endomorphismes (resp. automorphismes) de $E$ qui commutent avec tous les endomorphismes (resp. automorphismes) de $E$.
\finenonce{005588}


\finexercice


\exercice{5589, rouget, 2010/10/16}
\enonce{005589}{**I}
Soient $p$ et $q$ deux projecteurs d'un $\Cc$-espace vectoriel $E$.

Montrer que $(p+q\;\text{projecteur})\Leftrightarrow(p\circ q=q\circ p=0)\Leftrightarrow(\text{Im}(p)\subset\text{Ker}(q)\;\text{et}\;\text{Im}(q)\subset\text{Ker}(p))$.

Dans le cas où $p+q$ est un projecteur, déterminer $\text{Ker}(p+q)$ et $\text{Im}(p+q)$.
\finenonce{005589}


\finexercice


\exercice{5590, rouget, 2010/10/16}
\enonce{005590}{**I}
\label{ex:rou28}
Soit $E$ un espace de dimension finie. Montrer que la trace d'un projecteur est son rang.
\finenonce{005590}


\finexercice


\exercice{5591, rouget, 2010/10/16}
\enonce{005591}{****}
\label{ex:rou29}
Soient $p_1$,..., $p_n$ $n$ projecteurs d'un $\Cc$-espace de dimension finie. Montrer que $(p_1+...+p_n\;\text{projecteur})\Leftrightarrow\forall i\neq j,\;p_i\circ p_j=0$.
\finenonce{005591}


\finexercice


\exercice{5592, rouget, 2010/10/16}
\enonce{005592}{***}
Soit $E$ un $\Cc$-espace de dimension finie $n$. Soient $p_1$,..., $p_n$ $n$ projecteurs non nuls de $E$ tels que $\forall i\neq j$, $p_i\circ p_j=0$.

\begin{enumerate}
 \item  Montrer que tous les $p_i$ sont de rang $1$.

\item  Soient $q_1$,..., $q_n$ $n$ projecteurs vérifiant les mêmes égalités. Montrer qu'il existe un automorphisme $f$ de $E$ tel que $\forall i\in\llbracket1,n\rrbracket$, $q_i=f\circ p_i\circ f^{-1}$.
\end{enumerate}
\finenonce{005592}


\finexercice


\exercice{5593, rouget, 2010/10/16}
\enonce{005593}{***}
Soit $E$ un espace vectoriel. Soit $G$ un sous-groupe fini de $\mathcal{GL}(E)$ de cardinal $n$. Soit $F$ un sous-espace de $E$ stable par tous les éléments de $G$ et $p$ un projecteur d'image $F$. Montrer que $\frac{1}{n}\sum_{g\in G}^{}g\circ p\circ g^{-1}$ est un projecteur d'image $F$.
\finenonce{005593}


\finexercice


\exercice{5598, rouget, 2010/10/16}
\enonce{005598}{***}
Soient $E$ un $\Cc$-espace vectoriel de dimension finie et $f$ un endomorphisme de $E$. Montrer qu'il existe un projecteur $p$ et un automorphisme $g$ de $E$ tel que $f=g\circ p$.
\finenonce{005598}


\finexercice


\exercice{5599, rouget, 2010/10/16}
\enonce{005599}{**I}
Soient $E$ un $\Cc$-espace vectoriel non nul de dimension finie $n$ et $f$ un endomorphisme de $E$ tel que $\forall x\in E$, $\exists p\in\Nn^*$ tel que $f^p(x)=0$. Montrer que $f$ est nilpotent.
\finenonce{005599}


\finexercice


\exercice{5600, rouget, 2010/10/16}
\enonce{005600}{***}
 Soit $E$ un $\Cc$-espace vectoriel de dimension finie non nulle. Soient $f$ et $g$ deux projecteurs distincts et non nuls de $E$ tels qu'il existe deux complexes $a$ et $b$ tels que :

\begin{center}
$fg - gf = af + bg$.
\end{center}

\begin{enumerate}
 \item  Montrer que si $a\neq0$ et $a\neq1$ on a : $\text{Im}(f)\subset\text{Im}(g)$. En déduire que $gf=f$ puis que $a+b =0$ puis que $a=-1$.

\item  Montrer que si $a\neq0$ et $a\neq-1$, on a $\text{Ker}(g)\subset\text{Ker}(f)$. Que peut-on en déduire ?

\item  Montrer que si $f$ et $g$ sont deux projecteurs qui ne commutent pas et vérifient de plus 
$fg-gf=af+bg$ alors 

$(a,b)$ est élément de $\{(-1,1),(1,-1)\}$. Caractériser alors chacun de ces cas.
\end{enumerate}
\finenonce{005600}


\finexercice

\section{107.99 Autre}


\exercice{3315, quercia, 2010/03/09}
\enonce{003315}{$\mathcal{L}({E\times F})$, Chimie P 1996}

Est-il vrai que $\mathcal{L}({E\times F})$ et $\mathcal{L}(E)\times \mathcal{L}(F)$ sont isomorphes~?
($E$ et $F$ espaces vectoriels de dimensions finies).

\finenonce{003315}



\finexercice


\exercice{3356, quercia, 2010/03/09}
\enonce{003356}{Centrale MP 2001}

Soit $f$ un endomorphisme  donné de $E$ de dimension $n$ et
$F=\{g\in \mathcal{L}(E)\mid g\circ f=f\circ g=0\}$. Trouver la dimension
de~$F$. 

\finenonce{003356}



\finexercice


\exercice{3357, quercia, 2010/03/09}
\enonce{003357}{X MP$^*$ 2001}

Soit $G$ un sous-groupe fini de $GL(\R^n)$ et $F=\bigcap_{g\in G} \mathrm{Ker}(g-\mathrm{id})$. 
Montrer que $\mathrm{Card}\,(G)\times\dim F= \sum_{g\in G} \mathrm{tr}(g)$.

\finenonce{003357}



\finexercice

\section{108.01 Propriétés élémentaires, généralités}


\exercice{1040, liousse, 2003/10/01}
\video{XwtvirsK2HU}
\enonce{001040}{}
Effectuer le produit des  matrices  :
$$\left( \begin{array}{cc} 2 & 1  \\ 3& 2 \end{array} \right)\times
 \left( \begin{array}{cc}1 & -1  \\ 1& 2 \end{array} \right) \ \ \ \  \left( \begin{array}
 {ccc} 1 & 2 & 0  \\ 3 & 1 & 4 \end{array} \right) 
\times
\left( \begin{array}{ccc} -1 &-1& 0 \\ 1 & 4 & -1 \\ 2 & 1 & 2\end{array} \right)
 \ \ \ \ \left( \begin{array}{ccc} a &b& c \\ c & b & a \\ 1 & 1 & 1\end{array} \right)
\times
 \left( \begin{array}{ccc} 1 &a& c \\ 1 & b & b \\ 1 & c & a\end{array} \right)$$
\finenonce{001040}


\finexercice



\exercice{1041, liousse, 2003/10/01}

\enonce{001041}{}
On consid\`ere 
la matrice suivante :
$$M = \left( \begin{array}{cccc} 0 & a & b & c \\ 0 & 0 & d & e \\ 0 & 0 & 0 & f \\ 
0 & 0 & 0 & 0 \end{array} \right)$$
Calculer $M^2, M^3, M^4, M^5.$
\finenonce{001041}



\finexercice



\exercice{1042, liousse, 2003/10/01}

\enonce{001042}{}
On consid\`ere 
les trois matrices suivantes :
$$A = \left( \begin{array}{cccc} 2 & -3 & 1 & 0 \\ 5 & 4 & 1 & 3 \\ 
6 & -2 & -1 & 7 \end{array} \right)\ \ \ \ \ 
B = \left( \begin{array}{cc} 7 & 2  \\ -5 & 2  \\ 3 & 1  \\ 
6 & 0 \end{array} \right) \ \ {\hbox { et } } \ \ 
C = \left( \begin{array}{ccc} -1 & 2 & 6  \\ 3 & 5 & 7 \end{array} \right)$$ 

\begin{enumerate}
\item Calculer $AB$ puis $(AB)C$.\\
\item Calculer $BC$ puid $A(BC)$.\\
\item Que remarque-t-on ?
\end{enumerate}
\finenonce{001042}



\finexercice



\exercice{1043, liousse, 2003/10/01}

\enonce{001043}{}
On consid\`ere 
les deux matrices suivantes : 
$$A = \left( \begin{array}{cccc} 2 & 3 & -4 & 1 \\ 5 & 2 & 1 & 0 \\ 3 & 1 & -6 & 7 \\ 
2 & 4 & 0 & 1\end{array} \right), \ \ \ \ \ \ B = \left( \begin{array}{cccc} 3 & -1 & -3 & 7 \\ 
4 & 0 & 2 & 1 \\ 2 & 3 & 0 & -5 \\ 1 & 6 & 6 & 1 \end{array} \right)$$
\begin{enumerate}
\item Calculer $AB$.\\
\item Calculer $BA$.\\
\item Que remarque-t-on ?
\end{enumerate}
\finenonce{001043}



\finexercice



\exercice{1044, liousse, 2003/10/01}

\enonce{001044}{}
Trouver les matrices qui 
commutent avec $A=\left( \begin{array}{ccc} 1 & 0 & 0\\ 
 0 &1 & 1 \\ 3 & 1 & 2 \end{array} \right).$ De m\^eme avec $A=
\left( \begin{array}{cc} a & b\\ 
 0 &a \end{array} \right).$
\finenonce{001044}



\finexercice



\exercice{1045, liousse, 2003/10/01}

\enonce{001045}{}
Soit $A=
\left( \begin{array}{ccc} 0 & 1 & 1\\ 
 1 &0 & 1 \\ 1 & 1 & 0 \end{array} \right).$ Calculer $A^{2}$ et v\'erifier que 
$A^{2}=A+2I_3,$ o\`u $I_3$ est la matrice identit\'e $3\times 3.$ En d\'eduire que 
$A$ est inversible et calculer son inverse.
\finenonce{001045}



\finexercice



\exercice{1046, liousse, 2003/10/01}

\enonce{001046}{}
\begin{enumerate}
\item Soit $A= \left( \begin{array}{ccc} 1 & 1 & 0\\ 
0& 1 &1 \\ 0&0&1  \end{array} \right)$ et soit $B=A-I_3$.

\begin{enumerate}
\item Calculer $B^2$, $B^3$ en d\'eduire une formule de r\'ecurrence que l'on
 d\'emontrera pour $B ^n$, pour tout entier $n$.

\item D\'evelopper $(B+I_3)^n$ par la formule du binome et simplifier.

\item En d\'eduire $A^{n}$ Pour tout entier $n$.
\end{enumerate}
\item Soit $A=
\left( \begin{array}{cccc} 1 & 1 & 1&1\\ 
0& 1 &1& 1 \\ 0&0&1 & 1 \\0&0&0&1 \end{array} \right).$ Pour tout entier $n,$ 
calculer $A^{n}$ en utilisant $A-I_4.$
\end{enumerate}
\finenonce{001046}



\finexercice



\exercice{1047, liousse, 2003/10/01}

\enonce{001047}{}
\begin{enumerate}
\item On consid\`ere la matrice  $A=\left( \begin{array}{ccc} 1 & 0 & 0\\ 
 0 &1 & 1 \\ 3 & 1 & 1 \end{array} \right).$ 
\begin{enumerate}
\item  Soient $B=\left( \begin{array}{ccc} 1 & 1 & 1\\ 
 0 &1 & 0 \\ 1 & 0 & 0 \end{array} \right)$ et $C=\left( \begin{array}{ccc} 1 &1 & 1\\ 
 1 &2 & 1 \\ 0 & -1 & -1 \end{array} \right)$

Montrer que $AB=AC$, a-t-on $A=C$ ? $A$ peut-elle \^etre inversible ?
\item D\'eterminer toutes les matrices $F$ telles que $A\times F =O$ ($O$ \'etant la matrice dont 
tous les coefficients sont nuls).
\end{enumerate}
\item Soit $A = \left( \begin{array}{cc} 1 & 2  \\ 3 & 4  \\ -1 & 4  \end{array} \right)$.
D\'eterminer toutes les matrices $B$ telles que $BA=I_2$.
\item  Soient $A$ et $B$ deux matrices carr\'ees $n\times n$ telles que $AB=A+I_n$.

Montrer que $A$ est inversible et d\'eterminer son inverse (en fonction de $B$).
\end{enumerate}
\finenonce{001047}



\finexercice



\exercice{1048, liousse, 2003/10/01}

\enonce{001048}{}
suivantes.
$$\left[\begin{array}{ccccc}1&1&2&1&1\\2&1&1&1&1\\1&1&1&2&1
\\2&1&1&1&1\\1&1&1&1&2\end{array}\right]\hbox{, }
\left[\begin{array}{ccccc}1&0&0&0&1\\1&1&0&0&0\\1&0&1&0&0
\\1&0&0&1&0\\1&0&0&0&1\end{array}\right]\hbox{, }
\left[\begin{array}{ccccc}1&1&1&1&3\\0&2&1&1&2\\1&1&1&2&2
\\2&1&1&1&3\\1&-1&1&1&0\end{array}\right]
\hbox{, }
\left[\begin{array}{ccccc}1&1&1&1&2\\0&2&1&1&2\\1&1&1&2&2
\\2&1&1&1&3\\1&-1&1&1&0\end{array}\right]$$
\finenonce{001048}



\finexercice



\exercice{1049, liousse, 2003/10/01}

\enonce{001049}{}

Soit $A$ une matrice carr\'{e}e d'ordre $n$ ; on suppose que $A^{2}$ est
une combinaison lin\'{e}aire de $A$ et $I_{n}$ : $A^{2}=\alpha A+\beta I_{n}.$
\begin{enumerate}
\item Montrer que $A^{n}$ est \'{e}galement une combinaison lin\'{e}aire de 
$A$ et $I_{n}$ pour tout $n\in \Nn^{\ast }.$\\
\item Montrer que si $\beta $ est non nul, alors $A$ est inversible et que $
A^{-1}$ est encore combinaison lin\'{e}aire de $A$ et $I_{n}.$\\
\item Application 1 : soit $A=J_{n}-I_{n},$ o\`{u} $J_{n}$ est la matrice
Attila (envahie par les uns...), avec $n\geq 1$. Montrer que $
A^{2}=\left( n-2\right) A+\left( n-1\right) I_{n}$ ; en d\'{e}duire que $A$
est inversible, et d\'{e}terminer son inverse.\\
\item Application 2 : montrer que si $n=2$, $A^{2}$ est toujours une
combinaison lin\'{e}aire de $A$ et $I_{2},$ et retrouver la formule donnant $
A^{-1}$ en utilisant 2.
\end{enumerate}
\finenonce{001049}



\finexercice



\exercice{1050, liousse, 2003/10/01}

\enonce{001050}{}
Soit $A=  \begin{pmatrix} -1 & 1 &1   \\ 1 & -1&1&\\ 1 & 1&-1 \end{pmatrix}$

Calculer $A^2$ et montrer que $A^2= 2I -A$, en d\'eduire que $A$ est inversible et calculer $A^{-1}$.
\finenonce{001050}



\finexercice



\exercice{1051, legall, 1998/09/01}

\enonce{001051}{}
Rappeler la structure d'espace vectoriel de  $M_n({\Rr})$. D\' eterminer une base de  $M_n({\Rr})$.
Donner sa dimension.
\finenonce{001051}



\finexercice



\exercice{1052, legall, 1998/09/01}
\video{7AwRBOj2CYY}
\enonce{001052}{}
Soit 
$A=\begin{pmatrix} 
1 & 0 & 2 \cr
0 & -1 & 1 \cr
1 & -2 & 0 \cr
\end{pmatrix}$. 
Calculer $A^3-A$. En déduire que $A$ est inversible puis déterminer $A^{-1}$.
\finenonce{001052}




\finexercice


\exercice{1053, legall, 1998/09/01}

\enonce{001053}{}
 D\' eterminer deux \' el\' ements  $A$  et  $B$  de
$M_2({\Rr})$  tels que : $AB=0$  et $BA\not = 0$.

\finenonce{001053}



\finexercice



\exercice{1054, legall, 1998/09/01}

\enonce{001054}{}
Soit $ E  $ le sous ensemble de $ M_3({\Rr}) $
d\' efini par
$ E = \Bigl \{  M(a,b,c)=\begin{pmatrix} a & 0 & c \cr
                    0 & b & 0 \cr
                    c & 0 & a \cr \end{pmatrix}  a , b  , c \in {\Rr} \Bigr \} .$

\begin{enumerate}
    \item Montrer que $ E  $ est un sous-espace vectoriel de $ M_3(\Rr) $ stable
pour la multiplication des matrices. Calculer $ \hbox{dim } (E) .$
    \item Soit $ M(a,b,c)  $ un \' el\' ement de $ E  .$
D\' eterminer, suivant les valeurs des param\`etres $  a , b \hbox{ et }
 c \in {\Rr} $ son rang. Calculer (lorsque cela est possible) l'inverse $ M(a,b,c)
^{-1} $ de $ M(a,b,c)  .$
    \item Donner une base de $ E  $ form\' ee de matrices inversibles et une
autre form\' ee de matrices de rang $ 1 .$
\end{enumerate}
\finenonce{001054}



\finexercice



\exercice{1055, legall, 1998/09/01}

\enonce{001055}{}
Soit $ A\in M_2(\R ) .$ On nomme commutant de $
A $ et on note
 $ C(A) $ l'ensemble des $ B\in M_2(\R ) $ telles que $ AB=BA .$
\begin{enumerate}
    \item Montrer que  $ C(A) $ et un sous espace vectoriel de $ M_2(\R ) .$
    \item Montrer que pour tout $ k \in \N  , A^k\in C(A) .$
\end{enumerate}
\finenonce{001055}



\finexercice





\exercice{1057, legall, 1998/09/01}

\enonce{001057}{}
Montrer que  $F=\{ M\in M_2({\Rr }) ;  tr(M)=0\} $  est un
sous-espace vectoriel de $  M_2({\Rr }) .$ D\' eterminer une base de $ F $
et la compl\'eter en une base de $  M_2({\Rr }) .$

\finenonce{001057}



\finexercice



\exercice{1058, legall, 1998/09/01}

\enonce{001058}{}
Soient $  A  $ et $  B \in M_n(\mathbb{K} )  $ deux matrices triangulaires
sup\'erieures.
\begin{enumerate}
    \item Montrer (en calculant les coefficients) que $  AB  $ est triangulaire sup\'erieure.
    \item Soit $  \varphi   $ un endomorphisme bijectif de $  \mathbb{K} ^n  $ et $  F  $
un sous-espace vectoriel de $  \mathbb{K} ^n  $ tel que $  \varphi (F) \subset F  .$ Montrer que
 que $  \varphi ^{-1}(F) \subset F  .$
    \item En d\'eduire une nouvelle d\'emonstration de 1. Montrer que si $  A  $ est inversible, $  A^{-1}  $ est triangulaire sup\'erieure.
\end{enumerate}
\finenonce{001058}



\finexercice



\exercice{1059, legall, 1998/09/01}

\enonce{001059}{}
Soit $  N\in M_n(\C )  $ une matrice nilpotente. Calculer $
\hbox{det}(I+N)  .$ Si $  A\in M_n(\C )  $ commute avec $  N  ,$ montrer que $
\hbox{det}(A+N)=\hbox{det}(A)  .$ (on pourra commencer par \'etudier le cas o\`u $  A  $ est inversible.)
\finenonce{001059}



\finexercice



\exercice{1060, ridde, 1999/11/01}

\enonce{001060}{}
Soit $G = \left\{ \begin{pmatrix}
2^x&0&0 \\ 0&1&x \\ 0&0&1
\end{pmatrix}, x \in \Rr\right\}$. Montrer que $G$ est un groupe multiplicatif.
\finenonce{001060}



\finexercice



\exercice{1061, ridde, 1999/11/01}
\video{O1CDw2RqXUw}
\enonce{001061}{}
Soit $A(\theta) = \begin{pmatrix} \cos \theta & -\sin \theta \\ \sin \theta
& \cos \theta \end{pmatrix}$ pour $\theta \in \Rr$. Calculer $A(\theta) \times A(\theta')$ et $\big(A(\theta)\big)^n$ pour
$n \ge 1$.
\finenonce{001061}


\finexercice


\exercice{1062, ridde, 1999/11/01}

\enonce{001062}{}
Soit $A = \begin{pmatrix} 0&0&0 \\ -2&1&-1 \\ 2&0&2 \end{pmatrix}$.
\begin{enumerate}
\item Calculer $A^3-3A^2 + 2A$.
\item Quel est le reste de la division euclidienne de $X^n$ par $X^3-3X^2 + 2X$ ?
\item Calculer $A^n$ pour $n \in \Nn$.
\item $A$ est-elle inversible ?
\end{enumerate}
\finenonce{001062}



\finexercice



\exercice{1063, ridde, 1999/11/01}
\video{-ew4tp-RHJY}
\enonce{001063}{}
Soient $A$ et $B \in \mathcal{M}_n(\Rr)$ telles que $\forall X \in \mathcal{M}_n (\Rr)$,
$\text{tr} (AX) = \text{tr} (BX)$. Montrer que $A = B$.
\finenonce{001063}



\finexercice


\exercice{1064, ridde, 1999/11/01}
\video{IfhgaXB6Oh0}
\enonce{001064}{}
Que peut-on dire d'une matrice $A \in \mathcal{M}_n (\Rr)$ qui vérifie $\text{tr} (A \ {}^{t}\!{A}) = 0$ ?
\finenonce{001064}



\finexercice

\exercice{1065, ridde, 1999/11/01}

\enonce{001065}{}
Discuter suivant les valeurs de $\lambda \in \Rr$ le rang de la matrice
$\begin{pmatrix}
1&\frac 12&\frac 13\\\frac 12&\frac 13&\frac 14\\\frac 13 & \frac 14 &\lambda
\end{pmatrix}$.
\finenonce{001065}



\finexercice



\exercice{1066, ridde, 1999/11/01}

\enonce{001066}{}
Calculer l'inverse de $\begin{pmatrix} 1&2&1\\1&2&-1\\-2&-2&-1 \end{pmatrix}$.
\finenonce{001066}



\finexercice



\exercice{1067, gourio, 2001/09/01}

\enonce{001067}{}
D\'{e}terminer l'ensemble des matrices $M\in M_{n}({\Rr})$ telles que :
$$\forall H\in M_{n}({\Rr}),MH=HM. $$
\finenonce{001067}



\finexercice



\exercice{1068, gourio, 2001/09/01}

\enonce{001068}{}
Soit $M\in M_{n}({\Rr})$ telle que $M-I_{n}$ soit nilpotente (ie $\exists
k\in \Nn,(M-I_{n})^{k}=0$). Montrer que $M$ est inversible.
\finenonce{001068}



\finexercice



\exercice{1069, gourio, 2001/09/01}
\video{XdgqhVnKjTg}
\enonce{001069}{}
$A=\left( a_{i,j}\right) \in \mathcal{M}_{n}(\Rr)$
telle que :
$$\forall i =1,\ldots,n \qquad \left| a_{i,i}\right| >\sum_{j\neq i}\left|a_{i,j}\right| . $$
Montrer que $A$ est inversible.
\finenonce{001069}


\finexercice


\exercice{1070, gourio, 2001/09/01}

\enonce{001070}{}
Montrer que si $(A,B)\in M_{n}({\Rr})$ et $AB=A+B$ alors $AB=BA$.
\finenonce{001070}



\finexercice



\exercice{1071, gourio, 2001/09/01}

\enonce{001071}{}
Soit $M=\left( a_{i,j}\right) _{(i,j)\in \{1,...,n\}^{2}}\in M_{n}({\Rr}), $
montrer :
$$\min\limits_{j}\max\limits_{i}a_{i,j}\geq
\max\limits_{i}\min\limits_{j}a_{i,j}. $$
\finenonce{001071}



\finexercice



\exercice{1072, gourio, 2001/09/01}

\enonce{001072}{}
Soit $J\in M_{n}({\Rr})$ une matrice telle que : $J^{2}=I$ et
$$E=\{A\in M_{n}({\Rr})|\exists (a,b)\in {\Rr}^{2};A=aI+bJ\}. $$
\begin{enumerate}
\item
Montrer que $E $ est un espace vectoriel stable par multiplication (Est-ce
une alg\`{e}bre ?). En d\'{e}duire que :
$$
\forall A\in E,\forall n\in \Nn,\exists (a_{n},b_{n})\in {\Rr}%
^{2};A^{n}=a_{n}I+b_{n}J $$
et calculer les coefficients $a_{n}$ et $b_{n}.$

\item Soit $S_{n}=\sum\limits_{k=0}^{n}\frac{A^{k}}{k!}.$ Calculer $(u_{n},v_{n})$
tel que $S_{n}=u_{n}I+v_{n}J$ en fonction de $a$ et de $b.$ Calculer les
limites de $(u_{n})_{n\in \Nn}$ et de $(v_{n})_{n\in \Nn}.$ On pose $%
e^{A}=uI+vJ$ o\`{u} $u=\lim\limits_{n\rightarrow \infty }u_{n},$ $%
v=\lim\limits_{n\rightarrow \infty }v_{n}. $ Calculer $e^{-A\text{ }}$ et le
produit $e^{-A\text{ }}e^{A\text{ }}.$

\item Application :
$$
J=\left(
\begin{array}{cc}
0 & 1 \\
1 & 0
\end{array}
\right) ,A=\left(
\begin{array}{cc}
a & b \\
b & a
\end{array}
\right) .
$$
Calculer $e^{A}.$
\end{enumerate}
\finenonce{001072}



\finexercice



\exercice{1073, gourio, 2001/09/01}

\enonce{001073}{}
Soit $(A,B)\in (M_{n}(\Cc))^{2}$ tel que $\forall X\in M_{n}(\Cc),$ $AXB=0.$
Montrer que $A=0$ ou $B=0.$
\finenonce{001073}



\finexercice



\exercice{1074, gourio, 2001/09/01}

\enonce{001074}{}
Soit $(A,B)\in (M_{n}(\Cc))^{2}$ tel que $AB=I+A+A^{2}. $ Montrer que
$AB=BA$ (\emph{Indication} : voir d'abord que $A$ est inversible).
\finenonce{001074}



\finexercice



\exercice{1075, gourio, 2001/09/01}

\enonce{001075}{}
Soit $A\in M_{n}({\Rr}) $une matrice triangulaire \`{a} \'{e}l\'{e}ments
diagonaux nuls, montrer que :
$$A^{n}=0. $$
\finenonce{001075}



\finexercice



\exercice{1076, gourio, 2001/09/01}

\enonce{001076}{}
Calculer les puissances de :
$$\left(
\begin{array}{ll}
a & b \\
0 & a
\end{array}
\right) ,\ \left(
\begin{array}{ll}
a & b \\
b & a
\end{array}
\right) ,\ \left(
\begin{array}{lll}
1 & 1 & 1 \\
0 & 1 & 1 \\
0 & 0 & 1
\end{array}
\right) .
$$
\finenonce{001076}



\finexercice



\exercice{1077, gourio, 2001/09/01}

\enonce{001077}{}
Soit $A\in M_{n}({\Rr}) $ nilpotente, on d\'{e}finit :
$$\exp A=\sum_{i\geq 0}\frac{A^{i}}{i!}, $$
la somme \'{e}tant finie et s'arr\^{e}tant par exemple au premier indice $i $
tel que $A^{i}=0.$
Montrer que si $A$ et $B$ sont nilpotentes et commutent, alors $\exp(A+B)=\exp(A)\exp(B).$
En d\'{e}duire que $\exp(A)$ est toujours inversible et calculer son inverse.
\finenonce{001077}



\finexercice



\exercice{1078, gourio, 2001/09/01}

\enonce{001078}{}
Calculer l'inverse de :
$$
\left(
\begin{array}{llll}
1 & ... & ... & 1 \\
0 & 1 & ... & ... \\
... & 0 & 1 & ... \\
0 & ... & 0 & 1
\end{array}
\right) ,\left(
\begin{array}{llll}
1 & 2 & ... & n \\
0 & 1 & 2 & ... \\
... & 0 & 1 & 2 \\
0 & ... & 0 & 1
\end{array}
\right) .
$$
\finenonce{001078}



\finexercice



\exercice{1079, gourio, 2001/09/01}

\enonce{001079}{}
Calculer l'inverse de :
$$
\left(
\begin{array}{llll}
1 & a & ... & a \\
0 & 1 & a & ... \\
... & 0 & 1 & a \\
0 & ... & 0 & 1
\end{array}
\right) ,a\in {\Rr}.
$$
\finenonce{001079}



\finexercice



\exercice{1080, monthub, 2001/11/01}

\enonce{001080}{Examen}
 Soient $(x_n)_{n \in \N}$ et  $(y_n)_{n \in \N}$
deux suites réelles, vérifiant la relation de récurrence linéaire
suivante :
$$\Big\{
  \begin{array}{rccc}
x_{n+1}&=&-9 x_n &-18 y_n \\
y_{n+1}&=&6 x_n &+ 12y_n\\
  \end{array}$$
avec $x_0=-137$ et $y_0=18$. On se propose dans ce problème de
trouver les termes généraux de ces deux suites.

\begin{enumerate} 
\item Montrer qu'il existe une matrice $A \in M_2(\R)$ telle que la relation
  de récurrence linéaire ci-dessus soit équivalente à la relation
  $U_{n+1}=AU_n$, où $U_n=  \begin{pmatrix}    x_n\\y_n  \end{pmatrix}$.
\item Trouver une expression de $U_n$ en fonction de $A$ et de $U_0$.
\item Trouver le noyau de $A$, et en donner une base $B_1$. Calculer le rang de $A$.
\item Montrer que l'ensemble des vecteurs $X \in \R^2$ tels que $AX=3X$ est un
  sous-espace vectoriel de $\R^2$. Quelle
  est sa dimension ? En donner une base, qu'on notera $B_2$.

\item Montrer que la réunion $B_1 \cup B_2$ forme une base $B$ de $\R^2$. Soit
  $P$ la matrice formée des composantes des vecteurs de $B$ relativement à la
  base canonique de $\R^2$. Montrer que $P$ est inversible, et que  le produit
  $P^{-1}AP$ est une matrice diagonale $D$ qu'on calculera.

\item Montrer que $A^n=PD^nP^{-1}$. Calculer $D^n$, et en déduire $A^n$, pour tout $n \in \N$.
\item Donner les termes généraux  $x_n$ et $y_n$.

\end{enumerate}

\finenonce{001080}



\finexercice




\exercice{2442, matexo1, 2002/02/01}
\video{ATfy0zSe_04}
\enonce{002442}{}
Pour toute matrice carrée $A$ de dimension $n$, 
on appelle trace de $A$, et l'on note $\textrm{tr}\, A$, la
somme des éléments diagonaux de $A$\,:
$$\textrm{tr}\, A = \sum_{i = 1}^n a_{i,i}$$
\begin{enumerate}
\item Montrer que si $A, B$ sont deux matrices carrées
d'ordre $n$, alors $\textrm{tr}(AB)=\textrm{tr} (BA)$.

\item Montrer que si $f$ est un endomorphisme d'un espace
vectoriel $E$ de dimension $n$, $M$ sa matrice par rapport à
une base $e$, $M'$ sa matrice par rapport à une base $e'$,
alors $\textrm{tr}\, M = \textrm{tr}\, M'$. 
On note $\textrm{tr}\, f$ la valeur commune de ces quantités.

\item Montrer que si $g$ est un autre endomorphisme de $E$,
$\textrm{tr}(f\circ g - g\circ f) = 0$.
\end{enumerate}
\finenonce{002442}




\finexercice


\exercice{2443, matexo1, 2002/02/01}

\enonce{002443}{}
On rappelle qu'une matrice carr\'ee $A$ d'ordre $n$ est dite
{\em sym\'etrique} si $a_{i,j}=a_{j,i}, \forall\, i,j$, et {\em
antisym\'etrique} si $a_{i,j}=-a_{j,i}$.

\begin{enumerate}
\item Combien y a-t-il de matrices antisym\'etriques diagonales\,?
\item Montrer que $A {\,}^t A$ est sym\'etrique pour toute matrice
carr\'ee $A$.
\item Montrer que si $A, B$ sont sym\'etriques, leur produit
$C = AB$ est sym\'etrique si et seulement si $AB = BA$. Que dire
si elles sont antisym\'etriques\,? Si l'une est sym\'etrique et
l'autre antisym\'etrique\,?
\item Soit $P$ un polyn\^ome. Montrer que si $A$ est
sym\'etrique, $P(A)$ l'est aussi. Que dire si $A$ est
antisym\'etrique\,?
\end{enumerate}
\finenonce{002443}


\finexercice



\exercice{2444, matexo1, 2002/02/01}
\video{GsaYq45QDFE}
\enonce{002444}{}
Soient $A, B$ deux matrices semblables (i.e. il existe $P$
inversible telle que $B = P^{-1} A P$). Montrer que si l'une est
inversible, l'autre aussi\,; que si l'une est idempotente, l'autre
aussi\,; que si l'une est nilpotente, l'autre aussi\,; que si $A =
\lambda I$, alors $A = B$.
\finenonce{002444}



\finexercice



\exercice{2445, matexo1, 2002/02/01}

\enonce{002445}{}
Soit $A$ une matrice carr\'ee d'ordre $n$ v\'erifiant pour
tout $i\in\{1,\ldots,n\}$
$$ | a_{i,i} | > |a_{i,1}| + |a_{i,2}| +\ldots+ |a_{i,i-1}|+
|a_{i,i+1}| +\ldots+|a_{i,n}|.$$
Montrer que $A$ est inversible.
\finenonce{002445}


\finexercice


\exercice{2446, matexo1, 2002/02/01}

\enonce{002446}{}
Une matrice carr\'ee r\'eelle $A$ est dite {\em
stochastique} si $0 \le a_{i,j} \le 1, \forall\,i,j$ et
$\displaystyle{\sum_{i=1}^n a_{i,j} = 1}$, $\forall j$.

\begin{enumerate}
\item Montrer que le produit de deux matrices stochastique est
aussi une matrice stochastique.
\item Soit 
$B = A^2$, $A_i = \sup_j a_{i,j}$, $a_i = \inf_j a_{i,j}$.
Montrer que $a_i \le b_{i,j} \le A_i, \forall j$.
\end{enumerate}
\finenonce{002446}


\finexercice


\exercice{2747, tumpach, 2009/10/25}
\enonce{002747}{}
On consid\`ere les matrices suivantes~:
$$
A = \left(\begin{array}{ccc}1& 2 & -1\\ 2 & 3 & -2\\ 0 &0 & 0\end{array}\right); \quad\quad B = \left(\begin{array}{ccc}1 & 0 & -1\\ 2 & 0 & 4\\ 1 & 0 & -2\end{array}\right); \quad\quad  C = \left(\begin{array}{cc} 2 &-2 \\ 1 & 1\\ 3 & 1\end{array}\right);$$
$$D = \left(\begin{array}{c} 2 \\ 1 \\ 3\end{array}\right);\quad\quad E = \left(\begin{array}{ccc}0 & 1 & 2 \end{array}\right).
$$
Calculer \textit{lorsque cela est bien d\'efini} les produits de matrices suivants~: $AB$, $BA$, $AC$, $CA$, $AD$, $AE$, $BC$, $BD$, $BE$, $CD$, $DE$.
\finenonce{002747}



\finexercice


\exercice{2748, tumpach, 2009/10/25}
\enonce{002748}{}
Soient les matrices suivantes~:
$$
A = \left(\begin{array}{ccc}2& 5 & -1\\ 0 & 1 & 3\\ 0 &-2 & 4\end{array}\right); \quad\quad B = \left(\begin{array}{ccc}1 & 7 & -1\\ 2 & 3 & 4\\ 0 & 0 & 0\end{array}\right); \quad\quad  C = \left(\begin{array}{cc} 1 & 2 \\ 0 & 4\\ -1 & 0\end{array}\right).
$$
Calculer~:
 $(A - 2 B)C$,
$C^{T}A$,
$C^T B$,
$C^{T}(A^T - 2 B^T)$,
o\`u $C^T$ d\'esigne la matrice transpos\'ee de $C$.
\finenonce{002748}



\finexercice


\exercice{2749, tumpach, 2009/10/25}
\enonce{002749}{}
Calculer $A^n$ pour tout $n\in\mathbb{Z}$, avec successivement
$$
A = \left(\begin{array}{cc}\cos(a) & -\sin(a)\\ \sin(a) & \cos(a)\end{array}\right), \quad\quad\left(\begin{array}{cc}\cosh(a) & \sinh(a)\\ \sinh(a) & \cosh(a)\end{array}\right).
$$
\finenonce{002749}



\finexercice


\exercice{2750, tumpach, 2009/10/25}
\enonce{002750}{}
Les matrices suivantes sont-elles inversibles? Si oui, calculer leurs inverses.
$$
\left(\begin{array}{ccc}1 & 2 & 3\\ 2 & 3 & 1\\ 3 & 1 & 2\end{array}\right),
\left(\begin{array}{ccc} 1 & 0 & -1\\ 2 & 0  & 1\\1 & 1 & 3\end{array}\right),
\left(\begin{array}{ccc} 2 & 1 & -1\\0 & 3 & 0\\0 & 2 & 1\end{array}\right).
$$
\finenonce{002750}



\finexercice


\exercice{2751, tumpach, 2009/10/25}
\enonce{002751}{}
Inverser les matrices suivantes~:
$$
\left(\begin{array}{cccc} 1 & 1 & 1 & 1\\ 1 & 1 & -1 & -1\\ 1 & -1 & 1 & -1\\1 & -1 & -1 & -1\end{array}\right),
\left(\begin{array}{cccc} 1 & a & a^2 & a^3\\ 0 & 1 & a & a^2\\ 0 & 0 & 1 & a\\  0 & 0 & 0 &1\end{array}\right),
\left(\begin{array}{cccc}1 & 2 & 3 & 4\\ 0 & 1 & 2 & 3\\ 0 & 0 & 1& 2\\ 0 & 0& 0&1\end{array}\right).
$$
\finenonce{002751}



\finexercice


\exercice{2752, tumpach, 2009/10/25}
\enonce{002752}{}
L'\textit{exponentielle} d'une matrice carr\'ee $M$ est, par d\'efinition, la limite de la s\'erie 
$$
e^M = 1 + M + \frac{M^2}{2!} + \dots = \lim_{n\rightarrow+\infty} \sum_{k=0}^{n}\frac{M^k}{k!}.
$$
On admet que cette limite existe en vertu d'un th\'eor\`eme d'analyse.
\begin{enumerate}
\item Montrer que si $AB = BA$ alors $e^{A+B} = e^A e^B$. On est autoris\'e, pour traiter cette question, \`a passer \`a la limite sans pr\'ecautions.
\item Calculer $e^M$ pour les quatre matrices suivantes~:
$$
\left(\begin{array}{ccc}
a & 0 & 0\\0 & b & 0\\0 & 0& c
\end{array}\right),
\left(\begin{array}{ccc}
0 & a & b\\0 & 0 & c\\ 0 & 0 & 0
\end{array}\right),
\left(\begin{array}{cc}
0 & 1\\-1 & 0
\end{array}\right),
\left(\begin{array}{cc}
1 & 0 \\ 0  & 0
\end{array}\right).
$$
\item Chercher un exemple simple o\`u $e^{A + B}\neq e^A e^B$.
\end{enumerate}
\finenonce{002752}



\finexercice


\exercice{2772, tumpach, 2009/10/25}
\enonce{002772}{}
On consid\`ere la matrice $A = \left(\begin{array}{ccc}1 & 0 & 0\\ 0 & 1 & 1\\ 3 & 1 & 1\end{array}\right)$.
\begin{enumerate}
\item Soient $B = \left(\begin{array}{ccc}1 & 1 & 1\\ 0 & 1 & 0\\ 1 & 0 & 0\end{array}\right)$ et $C = \left(\begin{array}{ccc}1 & 1 & 1\\ 1 & 2 & 1\\ 0 & -1 & -1\end{array}\right)$.
Montrer que $A B = A C$. La matrice $A$ peut-elle \^etre inversible~?
\item D\'eterminer toutes les matrices $F$ de taille $(3, 3)$ telles que $A F = 0$, (o\`u $0$ est la matrice dont tous les coefficients sont nuls).
\end{enumerate}
\finenonce{002772}



\finexercice


\exercice{2773, tumpach, 2009/10/25}
\enonce{002773}{}
Pour quelles valeurs de $a$ la matrice 
$$
A = \left(\begin{array}{ccc} 1 & 1 & 1\\ 1 & 2 & 4\\ 1 & 3 & a\end{array}\right)
$$
est-elle inversible~? Calculer dans ce cas son inverse.
\finenonce{002773}



\finexercice



\exercice{2774, tumpach, 2009/10/25}
\video{JhulB3S-Wf4}
\enonce{002774}{} 
Soit $a $ et $b$ deux réels et $A$ la matrice
$$
A = \left(
\begin{array}{cccc}
a & 2 & -1 & b\\
3 & 0 & 1 & -4\\ 
5 & 4 & -1 & 2
\end{array}
\right)
$$
Montrer que $\textrm{rg}(A) \geq 2$. 
Pour quelles valeurs de $a$ et $b$ a-t-on $\textrm{rg}(A) = 2$~?
\finenonce{002774}



\finexercice



\exercice{2775, tumpach, 2009/10/25}
\enonce{002775}{}
Calculer l'inverse de la matrice suivante~:
$$
A = \left(\begin{array}{cccc} 4 & 8 & 7 & 4\\1 & 3 & 2 & 1\\
1 & 2 & 3 & 2 \\
0 & 0 & 1 & 1   \end{array}\right)
$$
\finenonce{002775}



\finexercice


\exercice{3358, quercia, 2010/03/09}
\enonce{003358}{Matrices en damier}

Soit $M = (a_{ij}) \in \mathcal{M}_n(K)$. On dit que $M$ est {\it en damier\/} si
$a_{ij} = 0$ pour $j-i$ impair.
On note ${\cal D}$ l'ensemble des matrices $n\times n$ en damier.
Montrer que ${\cal D}$ est une sous-algèbre de $\mathcal{M}_n(K)$.
Quelle est sa dimension ?

\finenonce{003358}



\finexercice


\exercice{3359, quercia, 2010/03/09}
\enonce{003359}{Matrices stochastiques}

Soit\\ 
${\cal D} = \left\{ A=(a_{ij}) \in \mathcal{M}_n(\R) \text{ tq } \forall\ i,j,\ a_{ij} \ge 0
      \text{ et } \forall\ i,\ \sum_{j=1}^n a_{ij} = 1 \right\}$.

\begin{enumerate}
  \item Montrer que ${\cal D}$ est stable par multiplication.
  \item Déterminer les matrices $A \in {\cal D}$ inversibles telles que
    $A^{-1} \in {\cal D}$.
\end{enumerate}
\finenonce{003359}



\finexercice


\exercice{3360, quercia, 2010/03/09}
\enonce{003360}{Matrices centrosymétriques}

Soit $A = (a_{ij}) \in \mathcal{M}_n(K)$. On dit que $A$ est {\it centro-symétrique\/}
si pour tous $i,j$ : $a_{n+1-i,n+1-j} = a_{ij}$.
Montrer que si $A$ et $B$ sont centro-symétriques, il en est de même de $AB$.
Montrer aussi que si $A$ est centro-symétrique et inversible alors $A^{-1}$ est aussi
centro-symétrique.
\finenonce{003360}



\finexercice

\exercice{3393, quercia, 2010/03/10}
\enonce{003393}{\'Equation $AX = B$}

Soit $A = \begin{pmatrix} 1&2&3\cr 2&3&4\cr 3&4&5 \cr\end{pmatrix}$.

\begin{enumerate}
  \item Montrer que l'équation en $X$ : $AX = B$, $X,B \in \mathcal{M}_{3,n}(K)$,
    a des solutions si et seulement si
   les colonnes de $B$ sont des progressions arithmétiques
   (traiter d'abord le cas $n=1$).

  \item Résoudre $AX = \begin{pmatrix}3 &3 \cr 4 &5 \cr 5 &7 \cr\end{pmatrix}$.
    
\end{enumerate}
\finenonce{003393}



\finexercice


\exercice{3394, quercia, 2010/03/10}
\enonce{003394}{\'Equation $AX = B$}

Soient $A = \begin{pmatrix}-2&1&1 \cr 8&1&-5 \cr 4&3&-3 \cr\end{pmatrix}$ et
       $C = \begin{pmatrix} 1&2&-1 \cr 2&-1&-1 \cr -5&0&3 \cr\end{pmatrix}$.
Existe-t-il une matrice $B$ telle que $BC = A$ ?
\finenonce{003394}



\finexercice


\exercice{3395, quercia, 2010/03/10}
\enonce{003395}{Calcul de $A^n$ par la formule du binôme}

Soit $A = \begin{pmatrix} 1&0&0\cr 0&1&1\cr 1&0&1 \cr\end{pmatrix}$.
En écrivant $A = I + J$, calculer $A^n$, $n \in \Z$.

\finenonce{003395}



\finexercice


\exercice{3396, quercia, 2010/03/10}
\enonce{003396}{Calcul de $A^n$ par polynôme annulateur}

Soit $A = \begin{pmatrix} 1 &2 &3\cr 2 &3 &1\cr 3 &1 &2 \cr\end{pmatrix}$.

\begin{enumerate}
  \item Vérifier que $(A-6I)(A^2-3I) = 0$.
  \item Soit $n\in\N$ et $P_n$ le polynôme de degré inférieur ou égal à 2
    tel que
    $${P(6) = 6^n},\quad {P\bigl(\sqrt3\bigr) = \bigl(\sqrt3\bigr)^n},\quad
    \text{et } {P\bigl(-\sqrt3\bigr) = \bigl(-\sqrt3\bigr)^n}.$$
    Montrer que $A^n = P_n(A)$.
  \item Même question pour $n \in \Zz$.
\end{enumerate}
\finenonce{003396}



\finexercice


\exercice{3397, quercia, 2010/03/10}
\enonce{003397}{Calcul de $A^k$}

Calculer $A^k$ pour $k \in \N$ :
\begin{enumerate}
  \item $A = \begin{pmatrix}1 &       &(2)\cr
                     &\ddots     \cr
                  (2)&       &1  \cr\end{pmatrix}$.
   

  \item $A = \begin{pmatrix} 1 &2 &3 &4 \cr
                    0 &1 &2 &3 \cr
                    0 &0 &1 &2 \cr
                    0 &0 &0 &1 \cr \end{pmatrix}$.
   

  \item $A = \begin{pmatrix} x^2 &xy  &xz  \cr
                    xy  &y^2 &yz  \cr
                    xz  &yz  &z^2 \cr \end{pmatrix}$.
   
\end{enumerate}
\finenonce{003397}



\finexercice

\exercice{5258, rouget, 2010/07/04}
\enonce{005258}{**}
Pour $x$ réel, on pose~:

$$A(x)=
\left(
\begin{array}{cc}
\ch x&\sh x\\
\sh x&\ch x
\end{array}
\right)
.$$

Déterminer $(A(x))^n$ pour $x$ réel et $n$ entier relatif.
\finenonce{005258}


\finexercice


\exercice{5262, rouget, 2010/07/04}
\enonce{005262}{**}
Soit $A=\left(
\begin{array}{ccccc}
0&0&\ldots&0&1\\
0& & &1&0\\
\vdots& & & &\vdots\\
0&1&0& &0\\
1&0&\ldots&\ldots&0 
\end{array}
\right)
\in\mathcal{M}_p(\Rr)$. Calculer $A^n$ pour $n$ entier relatif.
\finenonce{005262}


\finexercice


\exercice{5263, rouget, 2010/07/04}
\enonce{005263}{**}
Montrer que $\{\frac{1}{\sqrt{1-x^2}}\left(
\begin{array}{cc}
1&x\\
x&1
\end{array}
\right),\;x\in]-1,1[\}$ est un groupe pour la multiplication des matrices.
\finenonce{005263}


\finexercice


\exercice{5611, rouget, 2010/10/16}
\enonce{005611}{**}
Soient $A=(a_{i,j})_{1\leqslant i,j\leqslant n}$ et $B=(b_{i,j})_{1\leqslant i,j\leqslant n}$ deux matrices carrées de format $n$ telles que 
$a_{i,j}=0$ si $j\leqslant i+r -1$ et $b_{i,j}= 0$ si $j\leqslant i+s-1$ où $r$ et $s$ sont deux entiers donnés entre $1$ et $n$.
Montrer que si $AB=(c_{i,j})_{1\leqslant i,j\leqslant n}$ alors $c_{i,j}=0$ si $j\leqslant i+r+s-1$.
\finenonce{005611}


\finexercice
\section{108.02 Noyau, image}


\exercice{2436, matexo1, 2002/02/01}

\enonce{002436}{}
 Soit  $e = (e_1, e_2, e_3)$ la base canonique de $\R^3$ et $u$
l'endomorphisme dont la matrice dans cette base est 
$$ M = \left( 
       \begin{array}{ccc}
        1 & 1 & 1 \\
        1 & 1 & 1 \\
        1 & 1 & 1
       \end{array}
       \right).$$
Chercher le noyau et l'image de $u$. Calculer son rang de deux
mani\`eres. Calculer la matrice de $u^2$ dans la base $e$.
Montrer que $u^2-3u=0$.

\finenonce{002436}


\finexercice


\exercice{2449, matexo1, 2002/02/01}

\enonce{002449}{}
Calculer le rang de la matrice
$$ A = 
\left( \begin{array}{ccc}
-1& 1& 2 \\ -2& 2 & 4 \\ -1& 1& 2
\end{array} \right). $$
\finenonce{002449}


\finexercice


\exercice{5257, rouget, 2010/07/04}
\enonce{005257}{**T}
Soit $u$ l'endomorphisme de $\Rr^3$ dont la matrice dans la base canonique $(i,j,k)$ de $\Rr^3$ est~:

$$M=\left(
\begin{array}{ccc}
2&1&0\\
-3&-1&1\\
1&0&-1
\end{array}
\right)
.$$

\begin{enumerate}
\item  Déterminer $u(2i-3j+5k)$.
\item  Déterminer $\mbox{Ker}u$ et $\mbox{Im}u$.
\item  Calculer $M^2$ et $M^3$.
\item  Déterminer $\mbox{Ker}u^2$ et $\mbox{Im}u^2$.
\item  Calculer $(I-M)(I+M+M^2)$ et en déduire que $I-M$ est inversible. Préciser $(I-M)^{-1}$.
\end{enumerate}
\finenonce{005257}


\finexercice


\exercice{5260, rouget, 2010/07/04}
\enonce{005260}{**}
Soit $\begin{array}[t]{cccc}
f~:&\Rr_n[X]&\rightarrow&\Rr_{n+1}[X]\\
 &P&\mapsto&Q=e^{X^2}(Pe^{-X^2})'
\end{array}$.

\begin{enumerate}
\item  Vérifier que $f\in(\mathcal{L}(\Rr_n[X],\Rr_{n+1}[X])$.
\item  Déterminer la matrice de $f$ relativement aux bases canoniques de $\Rr_n[X]$ et $\Rr_{n+1}[X]$.
\item  Déterminer $\mbox{Ker}f$ et $\mbox{rg}f$.
\end{enumerate}
\finenonce{005260}


\finexercice


\exercice{5269, rouget, 2010/07/04}
\enonce{005269}{***T}
Déterminer le rang des matrices suivantes~:

$$\begin{array}{llll}
1)
\left(
\begin{array}{ccc}
1&1/2&1/3\\
1/2&1/3&1/4\\
1/3&1/4&m
\end{array}
\right)
&
2)\left(
\begin{array}{ccc}
1&1&1\\
b+c&c+a&a+b\\
bc&ca&ab
\end{array}
\right)
&
3)\left(
\begin{array}{cccc}
1&a&1&b\\
a&1&b&1\\
1&b&1&a\\
b&1&a&1
\end{array}
\right)
&
4)\;(i+j+ij)_{1\leq i,j\leq n}\\
5)\;(\sin(i+j))_{1\leq i,j\leq n}
&
6)\;
\left(
\begin{array}{ccccc}
a&b&0&\ldots&0\\
0&a&\ddots&\ddots&\vdots\\
\vdots&\ddots&\ddots&\ddots&0\\
0& &\ddots&\ddots&b\\
b&0&\ldots&0&a
\end{array}
\right).
\end{array}
$$
\finenonce{005269}


\finexercice


\exercice{5603, rouget, 2010/10/16}
\enonce{005603}{***}
Rang de la matrice $\left(
\begin{array}{cccc}
1&\cos(a)&\cos(2a)&\cos(3a)\\
\cos(a)&\cos(2a)&\cos(3a)&\cos(4a)\\
\cos(2a)&\cos(3a)&\cos(4a)&\cos(5a)\\
\cos(3a)&\cos(4a)&\cos(5a)&\cos(6a)
\end{array}
\right)$.
\finenonce{005603}


\finexercice


\exercice{5607, rouget, 2010/10/16}
\enonce{005607}{**}
Rang de la matrice $(i+j+ij)_{1\leqslant i,j\leqslant n}$.
\finenonce{005607}


\finexercice


\exercice{5622, rouget, 2010/10/16}
\enonce{005622}{***}
Soient $A\in\mathcal{M}_n(\Cc)$ et $B$ l'élément de $\mathcal{M}_{np}(\Cc)$ défini par blocs par $B=\left(
\begin{array}{cccc}
A&0&\ldots&0\\
0&\ddots&\ddots&\vdots\\
\vdots&\ddots&\ddots&0\\
0&\ldots&0&A
\end{array}
\right)$. Déterminer le rang de $B$ en fonction du rang de $A$.
\finenonce{005622}


\finexercice


\exercice{5623, rouget, 2010/10/16}
\enonce{005623}{***}
Soit $H$ un élément de $\mathcal{M}_n(\Cc)$ tel que $\forall A\in\mathcal{M}_n(\Cc),\;\exists\lambda_A\in\Cc/\;HAH=\lambda_AH$. Montrer que $\text{rg}H\leqslant 1$.
\finenonce{005623}


\finexercice


\exercice{5624, rouget, 2010/10/16}
\enonce{005624}{***}
Soit $M\in\mathcal{M}_3(\Rr)$. Montrer que les deux propriétés suivantes sont équivalentes :

\begin{center}
(1) $M^2 = 0$ et (2) $\text{rg}M\leqslant 1$ et $\text{tr}M = 0$.
\end{center}
\finenonce{005624}


\finexercice

\section{108.03 Matrice et application linéaire}


\exercice{1081, cousquer, 2003/10/01}

\enonce{001081}{}
Soit $h$ l'homomorphisme de $\mathbb{R}^3$ dans $\mathbb{R}^2$ défini par
rapport à deux bases $(e_1,e_2,e_3)$ et $(f_1,f_2)$ par la
matrice
$A=\begin{pmatrix}
    2 & -1 & 1\cr
    3 & 2  &-3
 \end{pmatrix}$.
\begin{enumerate}
\item On prend dans $\mathbb{R}^3$ la nouvelle base~:
$$e'_1=e_2+e_3,\quad e'_2=e_3+e_1,\quad e'_3=e_1+e_2.$$
Quelle est la nouvelle matrice $A_1$ de $h$~?

\item On choisit pour base de $\mathbb{R}^2$ les vecteurs~:
$$f'_1=\frac{1}{2}(f_1+f_2), \quad f'_2=\frac{1}{2}(f_1-f_2)$$
en conservant la base  $(e'_1,e'_2,e'_3)$ de $\mathbb{R}^3$.
Quelle est la nouvelle matrice $A_2$ de $h$~?
\end{enumerate}
\finenonce{001081}



\finexercice



\exercice{1082, cousquer, 2003/10/01}

\enonce{001082}{}
 Soit $h$ une application linéaire de rang $r$, de $E$, espace vectoriel de
dimension~$n$, dans~$F$, espace vectoriel de dimension~$m$.
\begin{enumerate}
\item Préciser comment obtenir une base $(e_i)_{i=1}^n$ de  $E$, et une base
$(f_j)_{j=1}^m$ de $F$, telles que $h(e_k)=f_k$ pour $k=1,\ldots,r$ et
$h(e_k)=0$ pour $k>r$. Quelle est la matrice de~$h$ dans un tel couple de bases~?

\item Déterminer un tel couple de bases pour l'homomorphisme
de $\mathbb{R}^4$ dans $\mathbb{R}^3$ défini dans les bases canoniques
par~:
$$h(x_1,x_2,x_3,x_4) = (y_1,y_2,y_3)\quad\mbox{avec}\quad
\left\{\begin{array}{rcl}
    y_1 & = & 2x_1-x_2+x_3-x_4 \\
    y_2 & = & x_2+x_3-2x_4 \\
    y_3 & = & x_1+2x_2+x_3+x_4
    \end{array}\right.$$

\item Même question pour l'application $f$ de $\mathbb{R}^3$ dans lui-même
définie par~:
$$ f(x,y,z)=(2x+y+z,-y+z,x+y).$$
\end{enumerate}
\finenonce{001082}



\finexercice



\exercice{1083, cousquer, 2003/10/01}

\enonce{001083}{}
 On désigne par $\mathcal{P}_2$ l'espace des polynômes sur~$\mathbb{R}$ 
de degré inférieur ou égal à~$2$. On désigne par $(e_0,e_1,e_2)$
la base canonique de $\mathcal{P}_2$ et on pose 
$$p_0=e_0, \quad p_1=e_1-\frac{1}{2}e_0,\quad p_2=e_2-e_1+\frac{1}{2}e_0.$$
\begin{enumerate}
\item Montrer que tout polynôme de $\mathcal{P}_2$ peut s'écrire de
façon unique sous la forme
$p=b_0p_0 +b_1p_1+b_2p_2$.

\item Écrire sous cette forme les polynômes~: $p'_0$, $p'_1$, $p'_2$,
$p'$, $Xp'$, $p''$.

\item Montrer que l'application $\varphi:\mathcal{P}_2 \to\mathcal{P}_2$ 
définie par $\varphi(p)=Xp'-\frac{1}{2}p'+\frac{1}{4}p''$
est une application linéaire.
Préciser le noyau et l'image de cette application.
Écrire les matrices de cette application par rapport à la base
canonique $(e_i)$ et par rapport à la base $(p_i)$. Écrire la matrice de
passage de la base $(e_i)$ à la base $(p_i)$~; quelle relation
lie cette matrice aux deux précédentes~?
\end{enumerate}
\finenonce{001083}



\finexercice



\exercice{1084, legall, 1998/09/01}

\enonce{001084}{}
Soit $ f : \C \rightarrow \C  $
l'application $ z \mapsto e^{i\theta }\bar{z} .$ On consid\`ere $ \C $ comme
un $\R $-espace vectoriel et on fixe la base $ \epsilon =\{ 1,i\} .$
\begin{enumerate}
    \item  Montrer que $ f $ est $\R $-lin\'eaire.
    \item Calculer $ A=\hbox{Mat}(f, \epsilon , \epsilon ) .$
    \item Existent-ils $ x $ et $ y \in \C -\{ 0\}  $ tels que $
f(x)=x $ et $ f(y) =-y ?$ Si c'est le cas d\'eterminer un tel $ x $ et un tel $ y  .$
\item D\'ecrire g\'eom\'etriquement $ f .$
\item Soit $ g : \C \rightarrow \C  $
l'application $ z \mapsto e^{i\rho }\bar{z} .$ Calculer
$ A=\hbox{Mat}(g\circ f, \epsilon , \epsilon ) $ et d\'ecrire
 g\'eom\'etriquement $ g\circ f .$
\end{enumerate}
\finenonce{001084}



\finexercice



\exercice{1085, legall, 1998/09/01}

\enonce{001085}{}
Soit $ f\in\mathcal{L}(\R ^3) $ telle que $ f^3=-f $ et
$ f\not= 0 .$
\begin{enumerate}
    \item Montrer que $ \hbox{Ker}(f)\cap  \hbox{Ker}(f^2+I)
=\{ 0\} ,$ $ \hbox{Ker}(f)\not= \{ 0\} $ et $
\hbox{Ker}(f^2+I)\not= \{ 0\}. $
    \item Soit $ x $ un \'el\'ement distinct de $ 0 $ de $
\hbox{Ker}(f^2+I) .$ Montrer
qu'il n'existe pas $ \alpha \in \R $ tel que $ f(x)=\alpha x  . $
En d\'eduire que $ \{ x, f(x)\}  $ est libre.
    \item Calculer $ \hbox{dim}(\hbox{Ker}(f)) $ et $
\hbox{dim}(\hbox{Ker}(f^2+I)) .$
    \item D\'eterminer une base $ \epsilon  $ de $ \R ^3 $ telle que~:
$ \hbox{Mat}(f, \epsilon )=\begin{pmatrix} 0 & 0 & 0 \cr
                                      0 & 0 & -1 \cr
                                      0 & 1 & 0 \cr \end{pmatrix} .$
\end{enumerate}
\finenonce{001085}



\finexercice



\exercice{1086, legall, 1998/09/01}

\enonce{001086}{}
Soient  $E$  un espace vectoriel de dimension  $n$, $f$  une application
lin\' eaire de  $E$  dans lui-m\^eme et  $x$  un \' el\' ement de  $E$  tel que la
famille  $f(x), ...,f^n(x)$  soit libre.
\begin{enumerate}
    \item Montrer que la famille  $x,f(x),\ldots,f^{n-1}(x)$  est
une base de  $E$. D\' eduiser-en que  $f$  est bijective.
    \item On suppose maintenant que  $f^n(x)=x$. D\' eterminer la
matrice de  $f$  dans la
base   $x,f(x),\ldots,f^{n-1}(x)$.
\end{enumerate}
\finenonce{001086}



\finexercice



\exercice{1087, ridde, 1999/11/01}
\video{SRrzcCzYHo0}
\enonce{001087}{}
Soit $\Rr^2$ muni de la base canonique $\mathcal{B}=(\vec{i}, \vec{j})$.
Soit $f : \Rr^2 \to \Rr^2$ la projection sur l'axe des abscisses $\Rr \vec{i}$ 
parall\`element à $\Rr (\vec{i} + \vec{j})$.
Déterminer $\textrm{Mat}_{\mathcal{B},\mathcal{B}}(f)$, la matrice de $f$ dans la base $(\vec{i}, \vec{j})$.

Même question avec $\textrm{Mat}_{\mathcal{B}',\mathcal{B}}(f)$ où $\mathcal{B'}$ est la base 
$(\vec{i} - \vec{j}, -2\vec{i}+3\vec{j})$ de $\Rr^2$.
Même question avec $\textrm{Mat}_{\mathcal{B}',\mathcal{B}'}(f)$.
\finenonce{001087}



\finexercice

\exercice{1088, legall, 1998/09/01}

\enonce{001088}{}
Soit  ${\Rr}[X]$  l'espace vectoriel des polyn\^omes
\`a coefficients r\' eels.
\begin{enumerate}
    \item Soit  $n\in {\Nn}$. Montrer que  ${\Rr}_n[X]$,
ensemble des polyn\^omes \`a
coefficients r\' eels et de degr\' e inf\' erieur ou \' egal \`a  $n$, est
un sous-espace vectoriel
de  ${\Rr}[X]$. Montrer que la famille  $1,X,\ldots,X^n$  est une base de
${\Rr}_n[X]$.
    \item Soient  $f$, $g$  et  $h$  les applications de  ${\Rr}[X]$
 dans lui-m\^eme d\' efinies par :
\setbox1=\vbox{\hbox{$f(P(X))=XP(X),$}
\vskip1mm
\hbox{$g(P(X))=P'(X),$}
\vskip1mm
\hbox{$h(P(X))=(P(X))^2.$}
}
\vskip1mm
\centerline{\box1}
\vskip1mm
\noindent
Montrer que les applications  $f$  et  $g$  sont lin\' eaires,
mais que  $h$  ne l'est pas.
 $f$  et  $g$  sont-elles injectives ? Surjectives ? D\' eterminer la
dimension de leurs noyaux respectifs. D\' eterminer l'image de  $f$.
    \item On d\' esigne par  $f_n$  et  $g_n$  les restrictions
de  $f$  et de  $g$  \`a  ${\Rr}_n[X]$. Montrer que l'image de  $g_n$
est incluse
dans  ${\Rr}_n[X]$  et celle de  $f_n$  est incluse dans  ${\Rr}_{n+1}[X]$.
D\' eterminer
la matrice de  $g_n$  dans la base  $1,X, ...,X^n$  de  ${\Rr}_n[X]$.
D\' eterminer la matrice de  $f_n$  de la base  $1,X, ...,X^n$  dans la
base  $1,X, ...,X^{n+1}$.
Calculer les dimensions respectives des images de  $f_n$  et de  $g_n$.
\end{enumerate}
\finenonce{001088}



\finexercice



\exercice{1089, legall, 1998/09/01}

\enonce{001089}{}
Soient $ A=\begin{pmatrix}-1 & 2 \cr
                                      1 & 0 \cr \end{pmatrix} $ et $ f $
l'application de $ M_2(\R ) $
dans lui-m\^eme $ M\mapsto AM .$ Montrer que $ f $ est lin\'eaire.
D\'eterminer sa matrice
dans la base canonique de $ M_2(\R ) .$
\finenonce{001089}



\finexercice



\exercice{1090, legall, 1998/09/01}

\enonce{001090}{}
Soit  $\varphi $  une application lin\' eaire de  ${\Rr}^2$  dans lui-m\^eme
telle que  $\varphi \not = 0$  et  $\varphi ^2  = 0$.
\begin{enumerate}
    \item Construire des exemples de telles applications.
    \item Soit  $x\in {\Rr}^2$  tel que  $\varphi (x)\not = 0$. Montrer que
$\{ x  ,  \varphi (x)\} $  est une base de  ${\Rr}^2$. D\' eterminer
la matrice de  $\varphi $  dans cette base.
\end{enumerate}
\finenonce{001090}



\finexercice



\exercice{1091, legall, 1998/09/01}

\enonce{001091}{}
Soit $ E $ un espace vectoriel et $ \varphi \in \mathcal{L} (E) . $
\begin{enumerate}
    \item On suppose que $ \hbox{Ker} (\varphi )=\hbox{Ker} (\varphi ^2 ) .$
Soit $ p\geq 1 $ et $ x \in \hbox{Ker}(\varphi ^p) .$
Montrer que $ x \in \hbox{Ker}(\varphi ^{p-1}) .$
En d\'eduire que $ \hbox{Ker} (\varphi ^p) =\hbox{Ker} (\varphi ) $
pour tout $ p\geq 1 .$
    \item Montrer de m\^eme que si $ \hbox{Ker} (\varphi ^2)=
\hbox{Ker} (\varphi ^3 ) $ alors $ \hbox{Ker} (\varphi ^p) =
\hbox{Ker} (\varphi ^2) $ pour tout $ p\geq 2 .$
    \item  On suppose d\'esormais que  $\varphi $ est une application lin\' eaire de
${\Rr}^3$  dans lui-m\^eme telle que $ \varphi ^2  \not = 0 $.
Soit $x\in {\Rr}^3$ tel que  $\varphi ^2(x)\not = 0$. Montrer que
$ \{ x  ,  \varphi (x) ,  \varphi ^2(x)\}  $ est une base de
${\Rr}^3$. D\' eterminer
 la matrice de  $\varphi $  dans cette base.
\end{enumerate}
\finenonce{001091}



\finexercice



\exercice{1092, legall, 1998/09/01}

\enonce{001092}{}
Soient  $E$  un espace vectoriel de dimension  $3$  et
$\varphi $  une application lin\' eaire  de  $E$  dans  $E$  telle
que  $\varphi ^2=0$  et  $\varphi \not =0$. Posons  $r= \hbox{rg} (\varphi
)$.
\begin{enumerate}
    \item Montrer que  $
\hbox{Im } (\varphi ) \subset \hbox {Ker } (\varphi )$. D\' eduiser-en
que  $r\leq 3-r$.  Calculer $r$.
    \item Soit  $e_1\in E$  tel que  $\varphi (e_1)\not =0$. Posons
$e_2=\varphi (e_1)$.
Montrer qu'il existe  $e_3\in \hbox {Ker } (\varphi )$  tel que la
famille  $\{ e_2 ,  e_3\} $  soit libre. Montrer que  $\{ e_1 ,
e_2 ,  e_3\} $  est une base  de  $E$.
    \item D\' eterminer la matrice de  $\varphi $  dans la base  $\{ e_1
,  e_2 ,  e_3\} $.
\end{enumerate}
\finenonce{001092}



\finexercice




\exercice{1093, legall, 1998/09/01}
\video{A2_9r3WOfTM}
\enonce{001093}{}
\label{exo1093}
Soit  $E$  un espace vectoriel et  $f$  une application linéaire de
$E$  dans  lui-m\^eme telle que  $f^2=f$.
\begin{enumerate}
    \item Montrer que  $E= \Ker f \oplus \Im f$.

    \item Supposons que  $E$ soit de dimension finie  $n$. 
Posons  $r= \dim \Im f$. 
Montrer qu'il existe une base 
$\mathcal{B}= ( e_1, \ldots ,e_n)$ de  $E$  telle que : 
 $f(e_i)=e_i$ si $i\le r$ et $f(e_i)=0$ si $i>r$. 
Déterminer la matrice de  $f$ dans cette base $\mathcal{B}$.
\end{enumerate}
\finenonce{001093}



\finexercice



\exercice{1094, legall, 1998/09/01}
\video{10aMCsQRMG4}
\enonce{001094}{}
Soit $f$ l'application de $\Rr_n[X]$  dans  $\Rr[X]$
définie en posant pour tout  $P(X)\in \Rr_n[X]$ : $f(P(X))=P(X+1)+P(X-1)-2P(X).$
\begin{enumerate}
    \item Montrer que  $f$  est linéaire et que son image est
incluse dans  $\Rr_n[X]$.
    \item Dans le cas o\`u  $n=3$, donner la matrice de  $f$  dans
la base  $1,X, X^2, X^3$. Déterminer ensuite, pour une valeur de  $n$
quelconque, la matrice de  $f$  dans la base  $1,X,\ldots,X^n$.
    \item Déterminer le noyau et l'image de  $f$. Calculer
leur dimension respective.
    \item Soit  $Q$  un élément de l'image de  $f$.
Montrer qu'il existe un unique  $P\in \Rr_n[X]$
tel que : $f(P)=Q$  et  $P(0)=P'(0)=0$.
\end{enumerate}
\finenonce{001094}



\finexercice

\exercice{1095, cousquer, 2003/10/01}

\enonce{001095}{}
 Soit $(e_1,e_2,e_3)$ une base de l'espace $E$ à trois dimensions sur un
corps $K$. $I_E$ désigne l'application identique de $E$. On considère
l'application  linéaire $f$ de $E$ dans $E$ telle que~:
$$f(e_1)= 2e_2+3e_3, \quad f(e_2)=2e_1-5e_2-8e_3, \quad f(e_3)=-e_1+4e_2+6e_3.$$
\begin{enumerate}
\item Étudier le sous-espace $\ker(f-I_E)$~: dimension, base.

\item Étudier le sous-espace $\ker(f^2+I_E)$~: dimension, base.

\item Montrer que la réunion des bases précédentes constitue
une base de~$E$. Quelle est la matrice de $f$ dans cette nouvelle
base~? et celle de $f^2$~?
\end{enumerate}
\finenonce{001095}



\finexercice



\exercice{1096, cousquer, 2003/10/01}

\enonce{001096}{}
 Soit $E$ un espace à $n$ dimensions et $f$ un endomorphisme
de $E$.
\begin{enumerate}


\item Montrer que la condition $f^2=0$ est équivalente à
$\mbox{Im} f \subset\ker f$. Quelle condition vérifie alors le
rang de $f$~? On suppose dans le reste de
l'exercice que $f^2=0$.

\item \label{qu2}Soit $E_1$ un supplémentaire de $\ker f$ dans $E$ et soit
$(e_1,e_2,\ldots,e_r)$ une base de $E_1$. Montrer que la famille des
vecteurs $(e_1,e_2,\ldots,e_r,f(e_1),f(e_2),\ldots,f(e_r))$ est libre.
Montrer comment on peut la compléter, si nécessaire, par des vecteurs de
$\ker f$ de façon à obtenir une base de~$E$. Quelle est la matrice
de~$f$ dans cette base~?

\item Sous quelle condition nécessaire et suffisante a-t-on
$\mbox{Im} f=\ker f$~?

\item Exemple~:
Soit $f$ l'endomorphisme de $\mathbb{R}^3$ dont la matrice dans
la base canonique est
$M(f)=\begin{pmatrix}
    1  & 0  & 1\cr
    2  & 0  & 2\cr
    -1  & 0  &  -1\end{pmatrix}$.
Montrer que $f^2=0$.
Déterminer une nouvelle base dans laquelle la matrice de $f$ a la forme
indiquée dans la question~\ref{qu2}).
\end{enumerate}
\finenonce{001096}



\finexercice



\exercice{1097, cousquer, 2003/10/01}
\video{jQcF6rnWgyI}
\enonce{001097}{}
Soient trois vecteurs $e_1,e_2,e_3$ formant une base de $\Rr^3$.
On note $\phi$ l'application linéaire définie par
$\phi(e_1)=e_3$, $\phi(e_2)=-e_1+e_2+e_3$ et $\phi(e_3)=e_3$.

\begin{enumerate}
\item Écrire la matrice $A$ de $\phi$ dans la base $(e_1,e_2,e_3)$.
Déterminer le noyau de cette application. 

\item On pose $f_1=e_1-e_3$, $f_2=e_1-e_2$,  $f_3=-e_1+e_2+e_3$.
Calculer $e_1,e_2,e_3$ en fonction de $f_1,f_2,f_3$.
Les vecteurs $f_1,f_2,f_3$ forment-ils une base de $\Rr^3$ ?

\item Calculer $\phi(f_1), \phi(f_2), \phi(f_3)$ en fonction de $f_1,f_2,f_3$.
Écrire la matrice $B$ de $\phi$ dans la base $(f_1,f_2,f_3)$ et trouver la nature
de l'application $\phi$.

\item On pose $P=\begin{pmatrix}1&1&-1\cr 0&-1&1\cr-1&0&1\cr\end{pmatrix}$. Vérifier que $P$ est
inversible et calculer $P^{-1}$. Quelle relation lie $A$, $B$, $P$ et $P^{-1}$ ?
\end{enumerate}
\finenonce{001097}



\finexercice

\exercice{1098, legall, 1998/09/01}

\enonce{001098}{}
Soit  $M_{\alpha , \beta }$  la matrice : $ M_{\alpha ,
\beta }= \begin{pmatrix}
1 &3
& \alpha & \beta \cr
2 & -1  & 2 & 1 \cr
-1 & 1 & 2 & 0 \cr \end{pmatrix}\in M_{3,4}({\Rr})$. D\' eterminer pour quelles
valeurs de  $\alpha $  et
de  $\beta $  l'application lin\' eaire qui lui est associ\' ee est  surjective.

\finenonce{001098}



\finexercice



\exercice{1099, legall, 1998/09/01}
\video{96dAjfnOIhI}
\enonce{001099}{}
Soient  
$A=\begin{pmatrix} 
1 & 2 & 1 \cr
3 & 4 & 1 \cr
5 & 6 & 1 \cr
7 & 8 & 1 \cr
\end{pmatrix},\ 
B=\begin{pmatrix} 
2 & 2 & -1 & 7  \cr
4 & 3 & -1 & 11 \cr
0 & -1 & 2 & -4 \cr
3 & 3 & -2 & 11 \cr 
\end{pmatrix} $.
Calculer $\textrm{rg}(A)$ et $\textrm{rg}(B)$. Déterminer une base du
noyau et une base de l'image pour chacune des applications linéaires associées $f_A$ et $f_B$.
\finenonce{001099}



\finexercice

\exercice{1100, legall, 1998/09/01}

\enonce{001100}{}
Soit  $E$  un espace vectoriel de dimension  $n$  et
$\varphi $  une
application lin\' eaire de  $E$  dans  $E$. Montrer qu'il existe un
polyn\^ome  $P\in {\Rr} [X]$  tel que  $P(f)=0$. (On pourra utiliser le
fait que  $\mathcal{L} (E)$  est isomorphe \`a  $M_n({\Rr})$.)

\finenonce{001100}



\finexercice





\exercice{1102, ridde, 1999/11/01}

\enonce{001102}{}
Même chose avec $A = \begin{pmatrix} 0&1& \dots &0 \\ \vdots & \ddots & \ddots & \vdots\\
\vdots & & \ddots &1\\ 0 & \dots & \dots &0 \end{pmatrix}$.
\finenonce{001102}



\finexercice



\exercice{1103, ridde, 1999/11/01}

\enonce{001103}{}
Soit $f \in \mathcal{L} (\Rr^3)$ de matrice $\begin{pmatrix} 3&-1&1\\0&2&0\\1&-1&3
\end{pmatrix}$ dans la base canonique. D\'eterminer la matrice de $f$ dans la base
$ (1, 0, -1), (0, 1, 1), (1, 0, 1)$.
\finenonce{001103}



\finexercice



\exercice{1104, ridde, 1999/11/01}
\video{VTTwQGImicw}
\enonce{001104}{}
Soit $f$ l'endomorphisme de $\Rr^2$ de matrice $A=\begin{pmatrix} 2&\frac 23\\
-\frac 52&-\frac 23 \end{pmatrix}$ dans la base canonique. Soient 
$e_1 = \begin{pmatrix} -2 \\ 3\end{pmatrix}$
et $e_2 = \begin{pmatrix} -2 \\ 5 \end{pmatrix}$.
\begin{enumerate}
\item Montrer que $\mathcal{B}'= (e_1, e_2)$ est une base de $\Rr^2$ et déterminer 
$\text{Mat}_{\mathcal{B}'}(f)$.
\item Calculer $A^n$ pour $n \in \Nn$.
\item Déterminer l'ensemble des suites réelles qui vérifient $\forall n \in \Nn$
$\begin{cases} x_{n + 1} = 2x_n + \dfrac 23 y_n \\ y_{n + 1} = -\dfrac 52 x_n -
\dfrac 23 y_n \end{cases}$.
\end{enumerate}
\finenonce{001104}


\finexercice

\exercice{1105, ridde, 1999/11/01}

\enonce{001105}{}
Soit $E = \text{vect} (AB-BA, (A, B)\in M_n (\Qq)^2)$.
\begin{enumerate}
\item Montrer que $E = \ker \text{tr}$ (pour l'inclusion non triviale, on trouvera
une base de $\ker \text{tr}$ form\'ee de matrices de la forme $AB-BA$).
\item Soit $f \in M_n (\Qq)^*$ telle que $\forall (A, B)\in M_n (\Qq)^2 $
$f (AB) = f (BA)$. Montrer qu'il existe $\alpha \in \Rr$ tel que $f = \alpha \text{tr}$.
\end{enumerate}
\finenonce{001105}



\finexercice



\exercice{1106, ridde, 1999/11/01}

\enonce{001106}{}
Soient $A = \begin{pmatrix} 1&1\\0&1 \end{pmatrix}$ et ${\Phi} : {M_2 (\Rr)} \to 
{M_2 (\Rr)}$, ${M} \mapsto {AM-MA}$. Montrer que $\Phi$ est lin\'eaire, d\'eterminer sa matrice
dans la base canonique et calculer $\ker \Phi$  et $\text{Im} \Phi$.
\finenonce{001106}



\finexercice



\exercice{2565, delaunay, 2009/05/19}
\enonce{002565}{}
Soit $A=(a_{ij})_{1\leq i,j\leq n}$ une matrice carr\'ee $n\times n$. On veut d\'emontrer le
 r\'esultat suivant d\^u \`a Hadamard : Supposons que pour tout $i\in\{1,\cdots,n\}$, on ait
$$|a_{ii}|>\sum_{j=1,j\neq i}^{n}|a_{ij}| $$
alors $A$ est inversible.
\begin{enumerate}  
\item Montrer le r\'esultat pour $n=2$.
  \item Soit $B$, la matrice obtenue en rempla\c cant, pour $j\geq 2$, chaque colonne $c_j$ de $A$ par la colonne
$$ c_j-{\frac{a_{1j}}{a_{11}}}c_1 ,$$ Calculer les $b_{ij}$ en fonction des $a_{ij}$. Montrer que si les
coefficients de $A$ satisfont les in\'egalit\'es ci-dessus, alors pour $i\geq 2$, on a
$$|b_{ii}|>\sum_{j=2,j\neq i}^{n}|b_{ij}| .$$
  \item D\'emontrer le r\'esultat de Hadamard pour $n$ quelconque.
\end{enumerate}

\finenonce{002565} 



\finexercice

\exercice{2585, delaunay, 2009/05/19}
\enonce{002585}{}
Soient $A$ et $B$ des matrices non nulles de $M_n(\R)$. On suppose que $A.B=0$.
\begin{enumerate}
 \item D\'emontrer que $\Im B\subset \ker A$.
 \item On suppose que le rang de $A$ est \'egal \`a $n-1$, d\'eterminer le rang de $B$.
\end{enumerate}
\finenonce{002585} 


\finexercice

\exercice{2776, tumpach, 2009/10/25}
\enonce{002776}{}
On d\'esigne par $\{e_1, e_2, \dots, e_n\}$ la base canonique de $\mathbb{R}^n$. \`A une permutation $\sigma\in \mathcal{S}_n$, on associe l'endomorphisme $u_\sigma$ de $\mathbb{R}^n$ suivant~:
$$
\begin{array}{cccl}
u_{\sigma}~:&  \mathbb{R}^n & \longrightarrow & \mathbb{R}^n\\ 
 & \left(\begin{array}{c}x_1\\ \vdots \\x_n\end{array}\right) & \longmapsto &  \left(\begin{array}{c}x_{\sigma(1)}\\ \vdots \\x_{\sigma(n)}\end{array}\right)
\end{array}
$$ 
\begin{enumerate}
\item Soit $\tau = (i j)$ une transposition. \'Ecrire la matrice de $u_{\tau}$ dans la base canonique. Montrer que $\textrm{det}(u_\tau) = -1$.
\item Montrer que $\forall \sigma, \sigma'\in \mathcal{S}_n$, $u_{\sigma} \circ u_{\sigma'} = u _{\sigma'\circ\sigma}$.
\item En d\'eduire que $\forall \sigma \in \mathcal{S}_n$, $\textrm{det}\, u_{\sigma} = \varepsilon(\sigma)$ o\`u $\varepsilon$ d\'esigne la signature.
\end{enumerate}
\finenonce{002776}



\finexercice


\exercice{3406, quercia, 2010/03/10}
\enonce{003406}{Coefficients du binôme}

Soit $A \in \mathcal{M}{n+1}(\Q)$ telle que $a_{ij} = C_{j-1}^{i-1}$.
Interpréter $A$ comme la matrice d'un endomorphisme simple de $\Q_n[X]$.
En déduire la matrice $A^{-1}$.
\finenonce{003406}



\finexercice


\exercice{3407, quercia, 2010/03/10}
\enonce{003407}{Coefficients du binôme}

Soit $A \in \mathcal{M}_n(K)$ telle que $a_{ij} = (-1)^{n-j}C_{n-j}^{i-1}$.

\begin{enumerate}
  \item Interpréter $A$ comme la matrice d'un endomorphisme de $ K_{n-1}[X]$.
    
  \item En déduire $A^3$.
\end{enumerate}
\finenonce{003407}



\finexercice


\exercice{5261, rouget, 2010/07/04}
\enonce{005261}{***I}
Soit $f$ un endomorphisme de $\Rr^3$, nilpotent d'indice $2$. Montrer qu'il existe une base de $\Rr^3$ dans 
laquelle la matrice de $f$ s'écrit $\left(
\begin{array}{ccc}
0&0&0\\
1&0&0\\
0&0&0
\end{array}
\right)$.
\finenonce{005261}


\finexercice


\exercice{7411, mourougane, 2021/08/10}
\enonce{007411}{}
Soit $f : \Rr^3 \rightarrow \Rr^3$ l'application linéaire définie par
$$f(x,y,z) = (4x + y + z, 4x + 7y +2z, -6x -6y-z).$$ 
\begin{enumerate}
 \item \'Ecrire la matrice $A$ de $f$ dans la base canonique $\mathcal{B}$ de $\Rr^3$.

 \item Montrer que les vecteurs $v_1 = (1,0,-2)$, $v_2 = (1,-1,0)$ et $v_3 = (0,-1,1)$ forment une base de $\Rr^3$ que l'on note $\mathcal{B}'$.

 \item Déterminer la matrice de passage $[Id_{R^3}]_{\mathcal{B}'}^{\mathcal{B}}$ de $\mathcal{B}$ à $\mathcal{B}'$, et la matrice de passage $[Id_{R^3}]_{\mathcal{B}}^{\mathcal{B}'}$ de $\mathcal{B}'$ à $\mathcal{B}$.

 \item Déterminer la matrice $M$ de $f$ dans la base $\mathcal{B}'$.

 \item En déduire $A^n$ pour tout $n \in \mathbb{N}^{*}$.
\end{enumerate}
\finenonce{007411}
\finexercice


\exercice{7412, mourougane, 2021/08/10}
\enonce{007412}{}
Soit $\R[X]$ l'ensemble de polynômes dans $\R$ et soit $\R_{n}[X]$ l'ensemble des polynômes de degré au plus $n$.

On considère l'application $\Phi:\R_{2}[X]\longrightarrow\R[X]$ définie par :
\begin{equation*}
\Phi(P)(X):=(2X+1)P(X)-(X^2-1)P'(X).
\end{equation*}
\begin{enumerate}
	\item Montrer que l'application $\Phi$ est un endomorphisme.

	\item Montrer que la matrice $A$ de l'application $\Phi$ dans la base $(1,X,X^2)$ s'écrit :
\begin{equation*}
A=\begin{pmatrix}
1&1&0\\
2&1&2\\
0&1&1
\end{pmatrix}.
\end{equation*}

	\item Déterminer le rang de $\Phi$.

	\item En déduire que $\Phi$ est une application inversible.

	\item Déterminer une base du noyau de l'application $\Phi-Id$, où $Id$ désigne l'application identité.

	\item Montrer que la dimension de l'image de $\Phi-Id$ est de dimension $2$. Déterminer une base de l'image.

	\item Déterminer l'ensemble des polynômes de $\R[X]$ vérifiant l'identité : $2X\,P=(X^2-1)\,P'$.
\end{enumerate}
\finenonce{007412}
\finexercice

\section{108.04 Exemples géométriques}


\exercice{3364, quercia, 2010/03/09}
\enonce{003364}{Homographies}

Pour $M = \begin{pmatrix} a&b\cr c&d \cr\end{pmatrix} \in GL_2(\R)$, on note
${f_M} : {\R \cup \{\infty\}} \to {\R \cup \{\infty\}},
 x \mapsto  {\frac {ax+b}{cx+d}}$

Montrer que $M  \mapsto f_M$ est un morphisme de groupes. Quel est son noyau ?

\finenonce{003364}



\finexercice
\section{108.05 Inverse, méthode de Gauss}


\exercice{3361, quercia, 2010/03/09}
\enonce{003361}{Conservation de l'inverse sur un sous-corps}

Soit $M \in \mathcal{M}_n(\Q)$. Comparer les énoncés :
\par {\bf1:} $M$ est inversible dans $\mathcal{M}_n(\Q)$.
\par {\bf2:} $M$ est inversible dans $\mathcal{M}_n(\C)$.



\finenonce{003361}



\finexercice


\exercice{3362, quercia, 2010/03/09}
\enonce{003362}{Algèbre de matrices}

On note
$U = \begin{pmatrix} 1   &\dots &1   \cr
               \vdots &&\vdots \cr
               1   &\dots &1   \cr \end{pmatrix} \in \mathcal{M}_n(\R)$
     et ${\cal A} = \{ aU + bI,\ a,b \in \R \}\quad (n \ge 2)$.

\begin{enumerate}
  \item Montrer que ${\cal A}$ est une sous algèbre commutative de $\mathcal{M}_n(\R)$.
  \item Soit $M = aU+bI \in {\cal A}$. Montrer que $M$ possède un inverse dans
    ${\cal A}$ si et seulement si
    $b(b+na) \ne 0$, et le cas échéant, donner $M^{-1}$.
    
  \item Montrer que si $b(b+na) = 0$, alors $M$ n'est pas inversible dans $\mathcal{M}_n(\R)$.
  \item Trouver les matrices $M \in {\cal A}$ vérifiant : $M^n = I$.
    

\end{enumerate}
\finenonce{003362}



\finexercice

\exercice{3365, quercia, 2010/03/09}
\enonce{003365}{Opérations par blocs}

\begin{enumerate}
  \item Soient $A_1 \in \mathcal{M}_{n,p_1}(K)$, $A_2 \in \mathcal{M}_{n,p_2}(K)$,
    $B_1 \in \mathcal{M}_{p_1,q}(K)$, $B_2 \in \mathcal{M}_{p_2,q}(K)$.

    On pose $A = \begin{pmatrix} A_1 &A_2 \end{pmatrix} \in \mathcal{M}_{n,p_1+p_2}(K)$
    et $B = \begin{pmatrix} B_1 \cr B_2 \cr \end{pmatrix} \in \mathcal{M}_{p_1+p_2,q}(K)$.
    Montrer que $AB = A_1B_1 + A_2B_2$.

  \item Soit $M = \begin{pmatrix}A & B \cr 0 & C \cr\end{pmatrix}$
    où $A,B,0,C$ sont des matrices de tailles $p\times p$, $p \times q$,
    $q\times p$, $q \times q$ (matrice triangulaire par blocs).
    Montrer que $M$ est inversible si et seulement si $A$ et $C$ le sont.
    Le cas échéant, donner $M^{-1}$ sous la même forme.
  \item En déduire une nouvelle démonstration de la propriété :
    {\it L'inverse d'une matrice triangulaire est triangulaire.}

\end{enumerate}
\finenonce{003365}



\finexercice


\exercice{3366, quercia, 2010/03/09}
\enonce{003366}{Décomposition d'une matrice en matrices inversibles}

Soit $A \in \mathcal{M}_n(K)$. Montrer qu'il existe $U,V \in GL_n( K)$ telles que $A = U+V$.

\finenonce{003366}



\finexercice

\exercice{3376, quercia, 2010/03/09}
\enonce{003376}{Tout hyperplan de $\mathcal{M}_n(K)$ contient une matrice inversible}

Soit $H$ un hyperplan de $\mathcal{M}_n(K)$ $(n\ge 2)$.

\begin{enumerate}
  \item Montrer qu'il existe $A \in \mathcal{M}_n(K)$ telle que
    $H = \{ M \text{ tq } \mathrm{tr}(AM) = 0 \}$.
  \item En déduire que $H$ contient une matrice inversible.
    

\end{enumerate}
\finenonce{003376}



\finexercice

\exercice{3380, quercia, 2010/03/09}
\video{dad4PGxRe4U}
\enonce{003380}{$M$ antisymétrique $ \Rightarrow  I+M$ est inversible}

Soit $M \in \mathcal{M}_n(\R)$ antisymétrique.

\begin{enumerate}
  \item Montrer que $I+M$ est inversible (si $(I+M)X = 0$, calculer ${}^t\!(MX)(MX)$).
  \item Soit $A = (I-M)(I+M)^{-1}$. Montrer que ${}^t\!A = A^{-1}$.
\end{enumerate}
\finenonce{003380}



\finexercice



\exercice{3381, quercia, 2010/03/09}
\enonce{003381}{\'Equation $X^2 + X = A$}

Soit $A = \begin{pmatrix} 1&1\cr 1&1 \cr\end{pmatrix}$. On veut résoudre l'équation dans
$\mathcal{M}_{2}(K)$ : $X^2 + X = A$.

Soit $X$ une solution et $\phi_A$, $\phi_X$ les endomorphismes de $ K^2$
de matrices $A$ et $X$ dans la base canonique.

\begin{enumerate}
  \item Montrer que $X$ ou $X+I$ n'est pas inversible.
  \item Si $X$ n'est pas inversible, montrer que $X$ est proportionnelle à $A$
     (on montrera que $\mathrm{Ker}\phi_X = \mathrm{Ker}\phi_A$ et $\Im\phi_X = \Im\phi_A$).
  \item Résoudre l'équation.
    

\end{enumerate}
\finenonce{003381}



\finexercice


\exercice{3382, quercia, 2010/03/09}
\enonce{003382}{Groupes de matrices}

Soit ${\cal G} \subset \mathcal{M}_n(K)$ tel que pour la multiplication, $\cal G$ soit un
groupe. On note $J$ l'élément neutre et pour $M \in \cal G$,
$\phi_M$ l'endomorphisme de $ K^n$ canoniquement associé à $M$.

\begin{enumerate}
  \item Montrer que $\phi_J$ est une projection.
    
  \item Montrer que : $\forall\ M \in \cal G$, ${\phi_M}_{|\mathrm{Ker}\phi_J} = 0$ et
    ${\phi_M}_{|\Im\phi_J}$ est un isomorphisme de $\Im\phi_J$.
    
  \item En déduire que $\cal G$ est isomorphe à un groupe $GL_k( K)$.
    
\end{enumerate}
\finenonce{003382}



\finexercice

\exercice{3398, quercia, 2010/03/10}
\enonce{003398}{Inversion de matrices}

Inverser les matrices suivantes :


\begin{enumerate}
  \item $\begin{pmatrix}0 &       &(1) \cr
                 &\ddots      \cr
              (1)&       &0   \cr\end{pmatrix}$.
    

  \item $\begin{pmatrix}a &       &(b) \cr
                 &\ddots      \cr
              (b)&       &a   \cr\end{pmatrix}$.
    

  \item $\begin{pmatrix} 1 &1      &       &(0) \cr
                  &\ddots &\ddots      \cr
                  &       &\ddots &1   \cr
             \span(0)     &       &1   \cr\end{pmatrix}$.
   

  \item $\begin{pmatrix} 1        &\bar\alpha &\bar\alpha^2 \cr
               \alpha   &1          &\bar\alpha   \cr
               \alpha^2 &\alpha     &1            \cr \end{pmatrix}$, $\alpha \in \C$.
   

  \item $\begin{pmatrix} (0) &&a_n \cr &\cdots\cr a_1 &&(0) \cr\end{pmatrix}$.
    

  \item $\begin{pmatrix}  1+\frac 1{\lambda_1} &&(1)                 \cr
                     &\ddots                               \cr
               (1)   &               &1+\frac 1{\lambda_n} \cr \end{pmatrix}$.
   
\end{enumerate}
\finenonce{003398}


\finexercice


\exercice{3399, quercia, 2010/03/10}
\enonce{003399}{Effet des arrondis}

Soient {\normalbaselineskip=15pt
       $A = \begin{pmatrix} 1       &\frac12 &\frac13 \cr
                      \frac12 &\frac13 &\frac14 \cr
                      \frac13 &\frac14 &\frac15 \cr \end{pmatrix}$}
    et $B = \begin{pmatrix} 1\phantom{.00}       &0.5\phantom{0}     &0.33    \cr
                      0.5\phantom{0}     &0.33    &0.25    \cr
                      0.33    &0.25    &0.20    \cr \end{pmatrix}$.
Calculer $A^{-1}$ et $B^{-1}$.

\finenonce{003399}



\finexercice

\exercice{5267, rouget, 2010/07/04}
\enonce{005267}{***}
Soit $A=(a_{i,j})_{1\leq i,j\leq n}$ $(n\geq2)$ définie par 

$$\forall i\in\{1,...,n\},\;a_{i,j}=
\left\{
\begin{array}{l}
i\;\mbox{si}\;i=j\\
1\;\mbox{si}\;i>j\\
0\;\mbox{si}\;i<j
\end{array}
\right..$$

Montrer que $A$ est inversible et calculer son inverse.
\finenonce{005267}


\finexercice


\exercice{5272, rouget, 2010/07/04}
\enonce{005272}{***I Théorème de \textsc{Hadamard}}
Soit $A\in\mathcal{M}_n(\Cc)$ telle que~:~$\forall i\in\{1,...,n\},\;|a_{i,i}|>\sum_{j\neq i}^{}|a_{i,j}|$. Montrer que $A$ est inversible.
\finenonce{005272}


\finexercice


\exercice{5274, rouget, 2010/07/04}
\enonce{005274}{***I Matrice de \textsc{Vandermonde} des racines $n$-ièmes de l'unité}
Soit $\omega=e^{2i\pi/n}$, $(n\geq 2)$. Soit $A=(\omega^{(j-1)(k-1)})_{1\leq j,k\leq n}$. Montrer que $A$ est inversible et calculer $A^{-1}$ (calculer d'abord $A\overline{A}$).
\finenonce{005274}


\finexercice


\exercice{5276, rouget, 2010/07/04}
\enonce{005276}{***I}
Soit $A=(a_{i,j})_{1\leq i,j\leq n+1}$ définie par $a_{i,j}=0$ si $i>j$ et $a_{i,j}=C_{j-1}^{i-1}$ si $i\leq j$.

Montrer que $A$ est inversible et déterminer son inverse. (Indication~:~considérer l'endomorphisme de $\Rr_n[X]$ qui à un polynôme $P$ associe le polynôme $P(X+1)$).
\finenonce{005276}


\finexercice

\exercice{5597, rouget, 2010/10/16}
\enonce{005597}{****}
Montrer que tout hyperplan de $M_n(\Rr)$ contient des matrices inversibles.
\finenonce{005597}


\finexercice

\exercice{5604, rouget, 2010/10/16}
\enonce{005604}{***}
Soit $A=(a_{i,j})_{1\leqslant i,j\leqslant n}$ définie par $a_{i,j}=1$ si $i=j$, $j$ si $i=j-1$ et $0$ sinon.
Montrer que $A$ est inversible et calculer $A^{-1}$.
\finenonce{005604}


\finexercice


\exercice{5610, rouget, 2010/10/16}
\enonce{005610}{***}
Soient $a_1$,..., $a_n$ $n$ réels tous non nuls et $A=\left(
\begin{array}{ccccc}
1+a_1&1&\ldots&\ldots&1\\
1&\ddots&\ddots& &\vdots\\
\vdots&\ddots& &\ddots&\vdots\\
\vdots& &\ddots&\ddots&1\\
1&\ldots&\ldots&1&1+a_n
\end{array}
\right)$.

Inverse de $A$ en cas d'existence ?
\finenonce{005610}


\finexercice


\exercice{5612, rouget, 2010/10/16}
\enonce{005612}{**I}
Calculer l'inverse de $\left(
\begin{array}{cccccc}
\dbinom{0}{0}&\dbinom{1}{0}&\dbinom{2}{0}&\ldots&\dbinom{n-1}{0}&\dbinom{n}{0}\\
\rule{0mm}{7mm}0&\dbinom{1}{1}&\dbinom{2}{1}&\ldots&\ldots&\dbinom{n}{1}\\
\rule{0mm}{7mm}\vdots&\ddots&\dbinom{2}{2}& & &\vdots\\
 & & &\ddots& & \\
\vdots& & &\ddots&\dbinom{n-1}{n-1}&\vdots\\
0&\ldots& &\ldots&0&\dbinom{n}{n}
\end{array}
\right)$.
\finenonce{005612}


\finexercice


\exercice{5613, rouget, 2010/10/16}
\enonce{005613}{***I}
Soit $n$ un entier naturel supérieur ou égal à $2$ et $\omega=e^{2i\pi/n}$.

Soit $A =(\omega^{(j-1)(k-1)})_{1\leqslant j,k\leqslant n}$. Montrer que $A$ est inversible et calculer $A^{-1}$.
\finenonce{005613}


\finexercice


\exercice{5617, rouget, 2010/10/16}
\enonce{005617}{***I Théorème de \textsc{Hadamard}}
Soit $A=(a_{i,j})_{1\leqslant i,j\leqslant n}\in\mathcal{M}_n(\Cc)$ telle que $\forall i\in\llbracket1,n\rrbracket$, $|a_{i,i}| >\sum_{j\neq i}^{}|a_{i,j}|$. Montrer que $A\in\mathcal{GL}_n(\Cc)$. (Une matrice à diagonale strictement dominante est inversible.)
\finenonce{005617}


\finexercice



\exercice{6872, bodin, 2012/05/29}
\video{SbWnWREUACY}
\enonce{006872}{}
Calculer (s'il existe) l'inverse des matrices :
$$
\begin{pmatrix} a & b \\ c & d \\ \end{pmatrix}
\qquad 
\begin{pmatrix} 1&2&1 \\ 1&2&-1\\ -2&-2&-1 \end{pmatrix}
\qquad 
\begin{pmatrix} 
1 &\bar\alpha &\bar\alpha^2 \cr
\alpha   &1          &\bar\alpha   \cr
\alpha^2 &\alpha     &1            \cr \end{pmatrix} (\alpha \in \Cc)
\qquad
\begin{pmatrix} 
0 & 1 & 1 & 1 \\ 1 & 0 & 1 & 1 \\ 1 & 1 & 0 & 1 \\ 1 & 1 & 1 & 0 \\ 
\end{pmatrix}
$$

$$
\begin{pmatrix}
1 & 1 & \cdots &\cdots & 1 \\
0 & 1 & \ddots & & \vdots  \\
 & \ddots & \ddots & \ddots &  \vdots \\
& \cdots & 0 & 1 & 1 \\
0 & & \cdots & 0 & 1
\end{pmatrix}
\qquad 
\begin{pmatrix}
1 & 2 & 3 & \cdots & n \\
0 & 1 & 2 & \cdots & \vdots \\
  & \ddots & \ddots & \ddots &  \vdots \\
\vdots &  & 0 & 1 & 2 \\
 0 & \cdots & & 0 & 1
\end{pmatrix}
$$
\finenonce{006872}


\finexercice



\exercice{7413, mourougane, 2021/08/10}
\enonce{007413}{}
Pour un entier $n\geq 2$ et $x$ dans $\Rr$, considérons la matrice d'ordre $n$:
\[
D_n =
\left ( {\begin{array}{cccc}
 x & 1 & \ldots & 1 \\
 1 & x & \ddots & \vdots \\
 \vdots & \ddots & \ddots & 1 \\
 1 & \ldots & 1 & x
 \end{array} } \right ).
\]
\begin{enumerate}
	\item Calculer $det(D_2)$ et $det(D_3)$.

	\item Montrer d'abord que
 \[
 det(D_n)= \left | {\begin{array}{ccccc}
 x+n-1 & x+n-1 & \ldots & \ldots & x+n-1 \\
 1 & x & 1& \ldots & 1 \\
 \vdots & \ddots & \ddots &\ddots& \vdots \\
 \vdots &\ldots & \ddots & \ddots & 1 \\
 1 & \ldots & \ldots & 1 & x
 \end{array} } \right |,
\]
 
 et ensuite
 
 \[
 det(D_n)= (x+n-1)\left | {\begin{array}{cccc}
 1 & 1 & \ldots & 1 \\
 1 & x & \ddots & \vdots \\
 \vdots & \ddots & \ddots & 1 \\
 1 & \ldots & 1 & x
 \end{array} } \right |.
\]

	\item En utilisant la méthode du pivot de Gauss, calculer $det(D_n)$ pour tout $n$.

	\item Pour chaque $n$, pour quelles valeurs de $x$ la matrice $D_n$ est-elle inversible?
\end{enumerate}
\finenonce{007413}
\finexercice

\section{108.06 Changement de base, matrice de passage}


\exercice{3367, quercia, 2010/03/09}
\enonce{003367}{Conjugaison}

\begin{enumerate}
  \item Soit $P \in GL_n( K)$. Montrer que l'application
    ${\phi_P} : {\mathcal{M}_n(K)} \to {\mathcal{M}_n(K)}, M \mapsto {P^{-1}MP}$ est un isomorphisme
    d'algèbre.

  \item Soit $\phi$ : $A = (a_{ij}) \longmapsto A' = (a_{n+1-i,n+1-j})$.
  \begin{enumerate}
    \item Montrer que $\phi$ est un isomorphisme d'algèbre de $\mathcal{M}_n(K)$.
    \item Trouver une matrice $P \in GL_n( K)$ telle que $\phi = \phi_P$.
  \end{enumerate}
\end{enumerate}
\finenonce{003367}



\finexercice


\exercice{3368, quercia, 2010/03/09}
\enonce{003368}{Chimie P' 1996}

Soit $u\in\mathcal{L}({\R^3})$ ayant pour matrice dans la base canonique
$M = \frac 12\begin{pmatrix}5 &1 &?\cr 2 &-1 &?\cr 1 &0 &0\cr\end{pmatrix}$ et
$M' = \begin{pmatrix}1 &1 &0\cr 0&1 &1\cr 0 &0 &1\cr\end{pmatrix}$ dans une autre base.
Donner la matrice de passage.

\finenonce{003368}



\finexercice


\exercice{3400, quercia, 2010/03/10}
\enonce{003400}{Changement de base}

Soit $f$ l'application linéaire de $\R^4$ dans $\R^3$ dont la matrice
relativement aux bases canoniques, $(\vec I, \vec J, \vec K, \vec L\,)$
et $(\vec i, \vec j, \vec k\,)$ est
$\begin{pmatrix} 4 &5  &-7 &\phantom-7 \cr
            2 &1  &-1 &3          \cr
            1 &-1 &2  &1          \cr \end{pmatrix}$.

On définit deux nouvelles bases :
${\cal B} = (\vec I, \vec J, 4\vec I+\vec J-3\vec L, -7\vec I+\vec K+5\vec L\,)$
et ${\cal B}' = (4\vec i+2\vec j+\vec k, 5\vec i+\vec j-\vec k, \vec k\,)$.

Quelle est la matrice de $f$ relativement à ${\cal B}$ et ${\cal B}'$ ?

	
\finenonce{003400}



\finexercice


\exercice{3401, quercia, 2010/03/10}
\enonce{003401}{Matrices semblables}

Soient $A = \begin{pmatrix} 1 &1 &0 &0 \cr
                      0 &1 &1 &0 \cr
                      0 &0 &1 &1 \cr
                      0 &0 &0 &1 \cr \end{pmatrix}$
et     $B = \begin{pmatrix} 1 &2 &3 &4 \cr
                      0 &1 &2 &3 \cr
                      0 &0 &1 &2 \cr
                      0 &0 &0 &1 \cr \end{pmatrix}$.
Montrer que $A$ et $B$ sont semblables.

(On cherchera $P$ inversible telle que $PB = AP$)

\finenonce{003401}



\finexercice


\exercice{3402, quercia, 2010/03/10}
\enonce{003402}{Matrices semblables}
\normalbaselineskip=15pt
\def\D{\phantom{\scriptstyle-}\frac 12}
\def\d{{\scriptstyle-}\frac12}%
Montrer que $M = \begin{pmatrix} 2   &\d  &\d \cr
                    0   &\D  &\D \cr
                    1   &\d  &\D \cr \end{pmatrix}$
et  $N = \begin{pmatrix} 1   &1   &0   \cr
                   0   &1   &1   \cr
                   0   &0   &1   \cr\end{pmatrix}$
sont semblables.


\finenonce{003402}



\finexercice


\exercice{3403, quercia, 2010/03/10}
\enonce{003403}{Matrices non semblables}

Montrer que $A = \begin{pmatrix} 1 &2 &3 \cr
                    0 &1 &2 \cr
                    0 &0 &1 \cr\end{pmatrix}$
et  $B = \begin{pmatrix} 3 &1 &2 \cr
                    2 &0 &1 \cr
                    1 &0 &0 \cr\end{pmatrix}$
ne sont pas semblables.


\finenonce{003403}



\finexercice


\exercice{3404, quercia, 2010/03/10}
\enonce{003404}{Matrices non semblables}

Soient $A = \begin{pmatrix} 29 & 38 &-18 \cr
                      -11 &-14 &  7 \cr
                       20 & 27 &-12 \cr \end{pmatrix}$
et     $B = \begin{pmatrix}  7 & -8 &  4 \cr
                        3 & -3 &  2 \cr
                       -3 &  4 & -1 \cr \end{pmatrix}$.

Montrer que $A$ et $B$ ont même rang, même déterminant, même trace mais ne sont
pas semblables (calculer $(A-I)^2$ et $(B-I)^2$).
\finenonce{003404}



\finexercice


\exercice{3405, quercia, 2010/03/10}
\enonce{003405}{Ensi Physique P 1995}

Les matrices
$\left(\begin{smallmatrix}0&1&1&0\cr 0&0&1&0\cr 0&0&0&0\cr 0&0&0&0\cr\end{smallmatrix}\right)$ et
$\left(\begin{smallmatrix}0&0&0&0\cr 0&0&1&1\cr 0&0&0&1\cr 0&0&0&0\cr\end{smallmatrix}\right)$
sont-elles semblables~?
\finenonce{003405}



\finexercice

\exercice{3433, quercia, 2010/03/10}
\enonce{003433}{Comatrice}

Soit $n\ge 2$ et $A \in \mathcal{M}_n(K)$.
\begin{enumerate}
  \item Si $A$ et $B$ sont inversibles, démontrer que
    $\text{com}\,(AB) = (\text{com}\, A)(\text{com}\, B)$.
  \item Démontrer le même résultat dans le cas général, en considérant
    les scalaires $\lambda$ tels que $A-\lambda I$ et $B-\lambda I$
    soient inversibles.
  \item En déduire que si $A$ et $B$ sont semblables,
    alors $\text{com}\, A$ et $\text{com}\, B$ le sont.
\end{enumerate}
\finenonce{003433}



\finexercice

\exercice{3577, quercia, 2010/03/10}
\enonce{003577}{Matrices réelles semblables sur $\C$}

Soient $A,B \in \mathcal{M}_n(\R)$ semblables sur $\C$ :
Il existe $P,Q \in \mathcal{M}_n(\R)$ telles que :
$\begin{cases} P+iQ \in GL_n(\C) \cr (P+iQ)A = B(P+iQ). \cr \end{cases}$

\begin{enumerate}
  \item Montrer que : $\forall\ \lambda \in \R,\ (P+\lambda Q)A = B(P+\lambda Q)$.
  \item En déduire que $A$ et $B$ sont semblables sur $\R$.
\end{enumerate}
\finenonce{003577}



\finexercice


\exercice{5259, rouget, 2010/07/04}
\enonce{005259}{***T}
Soit $u$ l'endomorphisme de $\Rr^3$ dont la matrice dans la base canonique $(i,j,k)$ de $\Rr^3$ est~:

$$M=\left(
\begin{array}{ccc}
0&1&0\\
0&0&1\\
1&-3&3
\end{array}
\right)
.$$

\begin{enumerate}
\item  Montrer que $u$ est un automorphisme de $\Rr^3$ et déterminer $u^{-1}$.
\item  Déterminer une base $(e_1,e_2,e_3)$ de $\Rr^3$ telle que $u(e_1)=e_1$, $u(e_2)=e_1+e_2$ et $u(e_3)=e_2+e_3$.
\item  Déterminer $P$ la matrice de passage de $(i,j,k)$ à $(e_1,e_2,e_3)$ ainsi que $P^{-1}$.
\item  En déduire $u^n(i)$, $u^n(j)$ et $u^n(k)$ pour $n$ entier relatif.
\end{enumerate}
\finenonce{005259}


\finexercice


\exercice{5626, rouget, 2010/10/16}
\enonce{005626}{**}
Soient $M(a)=\left(
\begin{array}{ccc}
4-a&1&-1\\
-6&-1-a&2\\
2&1&1-a
\end{array}
\right)$  et $N(a)=\left(
\begin{array}{ccc}
1-a&1&0\\
0&1-a&0\\
0&0&2-a
\end{array}
\right)$.
$M(a)$ et $N(a)$ sont-elles semblables ?
\finenonce{005626}


\finexercice


\exercice{5627, rouget, 2010/10/16}
\enonce{005627}{***I}
Soient $A$ et $B$ deux éléments de $\mathcal{M}_n(\Rr)$. Montrer que si $A$ et $B$ sont semblables dans $\mathcal{M}_n(\Cc)$, elles le sont dans $\mathcal{M}_n(\Rr)$.
\finenonce{005627}


\finexercice

\section{108.99 Autre}


\exercice{2434, matexo1, 2002/02/01}

\enonce{002434}{}
Soit $A$ une matrice carr\'ee qui commute avec toutes les
matrices carr\'ees. Montrer que c'est une matrice scalaire.
\finenonce{002434}


\finexercice


\exercice{2437, matexo1, 2002/02/01}

\enonce{002437}{}
Soit $A$ une matrice carr\'ee.

\begin{enumerate}
\item Montrer que $A^2 = I$ si et seulement si $(I-A)(I+A)=0$.
Montrer que dans ce cas $A$ est inversible.
\item Montrer que si $A$ est idempotente ($A^2 = A$), alors
$B=I-A$ l'est aussi et que $AB=BA=0$.
\item Montrer que $I$ est la seule matrice idempotente
inversible. 
\end{enumerate}
\finenonce{002437}


\finexercice


\exercice{2475, matexo1, 2002/02/01}
\video{zf0ETNslBnc}
\enonce{002475}{}
\label{exo2475}
Trouver toutes les matrices de $\mathcal{M}_3(\Rr)$ qui vérifient
\begin{enumerate}
\item $M^2 = 0$ ;
\item $M^2 = M$ ; 
\item $M^2 = I$. 
\end{enumerate}
\finenonce{002475}



\finexercice

\exercice{2693, matexo1, 2002/02/01}
\enonce{002693}{}
Un train qui ralentit avec une d{\'e}c{\'e}l{\'e}ration constante met 20s pour
parcourir le premier km et 30s pour parcourir le deuxi{\`e}me km. On veut calculer
la distance qu'il devra parcourir pour parvenir {\`a} l'arr{\^e}t.
\begin{itemize}
\item

En prenant pour origine la position initiale du train, {\'e}crire l'{\'e}quation
g{\'e}n{\'e}rale d'un mouvement uniform{\'e}ment   d{\'e}c{\'e}l{\'e}r{\'e}.

\item

En d{\'e}duire un syst{\`e}me de deux {\'e}quations dont les inconnues sont la
d{\'e}c{\'e}l{\'e}ration et la vitesse initiale du train, et r{\'e}soudre ce syst{\`e}me.
\item

Conclure.
\end{itemize}
\finenonce{002693}
\finexercice


\exercice{3363, quercia, 2010/03/09}
\enonce{003363}{Quaternions}

Montrer que
${\cal C} = \left\{ M = \begin{pmatrix} a&b\cr-b&a\cr\end{pmatrix} \in \mathcal{M}_2(\R) \right\}$
est un corps isomorphe à $\C$.


Montrer que
${\cal H} = \left\{ M = \begin{pmatrix} a&b\cr-\overline b&\overline a \cr\end{pmatrix}
\in \mathcal{M}_2(\C) \right\}$ est un corps non commutatif. 

\finenonce{003363}



\finexercice

\exercice{3369, quercia, 2010/03/09}
\enonce{003369}{Centre de $GL_n( K)$}

On note $(E_{ij})$ la base canonique de $\mathcal{M}_n(K)$.

\begin{enumerate}
  \item Montrer que $F_{ij} = I + E_{ij}$ est inversible.
  \item En déduire que $\text{vect}\big(GL_n( K)\big) = \mathcal{M}_n(K)$.
  \item Quel est le centre de $GL_n( K)$ ?
\end{enumerate}
\finenonce{003369}



\finexercice


\exercice{3370, quercia, 2010/03/09}
\enonce{003370}{Centre de $GL_n( K)$}

Soit $f \in \mathcal{L}(E)$ ayant même matrice dans toutes les bases de $E$.
Montrer que $f$ est une homothétie.

\finenonce{003370}



\finexercice


\exercice{3371, quercia, 2010/03/09}
\enonce{003371}{Centre des matrices triangulaires unipotentes}

On note ${\cal G} = \{ A = (a_{ij}) \in \mathcal{M}_n(K)$ tq $a_{ij} = 0$ si $i>j$ et
        $a_{ii} = 1 \}$.

\begin{enumerate}
  \item Montrer que $\cal G$ est un sous-groupe de $GL_n( K)$.
  \item En utilisant la base canonique de $\mathcal{M}_n(K)$, déterminer le centre de
    $\cal G$, et montrer que c'est un groupe commutatif isomorphe à
    $( K,+)$.
\end{enumerate}
\finenonce{003371}



\finexercice


\exercice{3372, quercia, 2010/03/09}
\enonce{003372}{\'Equation $aX + (trX)A = B$}

Soit $\alpha \in  K$, et $A,B \in \mathcal{M}_n(K)$. \'Etudier l'équation d'inconnue
$X \in \mathcal{M}_n(K)$ : $\alpha X + (\mathrm{tr} X)A = B$.
\finenonce{003372}



\finexercice


\exercice{3373, quercia, 2010/03/09}
\enonce{003373}{Commutant d'une matrice diagonale}

Soit $A \in \mathcal{M}_n(K)$ et ${\cal C}_A = \{ M \in \mathcal{M}_n(K) \text{ tq } AM = MA \}$
({\it commutant de $A$}).

\begin{enumerate}
  \item  Montrer que ${\cal C}_A$ est une sous-algèbre de $\mathcal{M}_n(K)$.
  \item Soit $A = \text{diag}(\lambda_1, \lambda_2, \dots, \lambda_n)$ une matrice
    diagonale dont tous les $\lambda_i$ sont distincts.
  \begin{enumerate}
     \item Chercher ${\cal C}_A$.
     \item Soit $\phi : {\mathcal{M}_n(K)} \to {\mathcal{M}_n(K)}, M \mapsto {MA-AM}$
         \par
         Montrer que $\Im \phi$ est l'ensemble des matrices à diagonale nulle.
  \end{enumerate}

\end{enumerate}
\finenonce{003373}



\finexercice


\exercice{3374, quercia, 2010/03/09}
\enonce{003374}{Matrices de trace nulle}

Soit $M \in \mathcal{M}_n(K)$ {\it non scalaire\/} telle que $\mathrm{tr} M = 0$.

\begin{enumerate}
  \item Montrer qu'il existe une matrice colonne $X_1$ telle que $MX_1$ ne soit
    pas colinéaire à $X_1$.
  \item En déduire que $M$ est semblable à une matrice
    $N = \begin{pmatrix} 0      &\dots  \cr
                   \vdots &M_1    \cr
    \end{pmatrix}$
    où $M_1 \in \mathcal{M}_{n-1}(K)$ et $\mathrm{tr} M_1 = 0$.
  \item Montrer que $M$ est semblable à une matrice à diagonale nulle.
  \item Montrer qu'il existe
    $A,B \in \mathcal{M}_n(K)$ telles que $M = AB-BA$.
\end{enumerate}
\finenonce{003374}



\finexercice


\exercice{3375, quercia, 2010/03/09}
\enonce{003375}{Forme bilinéaire trace}

\begin{enumerate}
  \item Soit $A \in \mathcal{M}_{n,p}(K)$ non nulle. Montrer que l'application
    ${f_A} : {\mathcal{M}_{p,n}(K)} \to { K}, X \mapsto {\mathrm{tr}(AX)}$ est une forme linéaire
    non nulle sur $\mathcal{M}_{p,n}(K)$.

  \item Réciproquement : Soit $\phi  : {\mathcal{M}_{p,n}(K)} \to { K}$ une forme linéaire
    quelconque. Montrer qu'il existe une unique matrice $A \in \mathcal{M}_{n,p}(K)$
    telle que $\phi = f_A$ (on pourra considérer l'application $A  \mapsto f_A$).

  \item Soit $\phi  : {\mathcal{M}_n(K)} \to  K$ une forme linéaire vérifiant :
    $\forall\ X,Y \in \mathcal{M}_n(K),\ \phi(XY) = \phi(YX)$.
    \par
    Montrer qu'il existe $\lambda \in  K$ tel que $\phi = \lambda\mathrm{tr}$.

\end{enumerate}
\finenonce{003375}



\finexercice


\exercice{3377, quercia, 2010/03/09}
\enonce{003377}{Matrices magiques}

Une matrice carrée $M$ est dite {\it magique} si les sommes des coefficients
de $M$ par ligne et par colonne sont constantes.
On note $s(M)$ leur valeur commune.


Soit $U = \begin{pmatrix} 1     &\dots&1     \cr
                     \vdots&     &\vdots\cr
                     1     &\dots&1     \cr\end{pmatrix}$
et ${\cal M} = \{ \text{matrices } n\times n \text{ magiques}\}$.

\begin{enumerate}
  \item Montrer que $\cal M$ est une sous-algèbre de $\mathcal{M}_n(K)$ et
    $s : {\cal M} \to  K$ est un morphisme d'algèbre (calculer $MU$ et $UM$).

  \item Si $M$ est magique inversible, montrer que $M^{-1}$ est aussi magique.

  \item Montrer que $\cal M$ est la somme directe du sev des matrices magiques
    symétriques et du sev des matrices magiques antisymétriques.

  \item Pour $M \in \mathcal{M}_n(K)$, on note $\phi_M$ l'endomorphisme de $ K^n$
    canoniquement associé à $M$.

    Soit ${\cal H} = \{ (x_1, \dots ,x_n)  \in  K^n \text{ tq }
                        x_1 + \dots + x_n = 0 \}$
    et   ${\cal K} = \{ ( x, \dots, x) \in  K^n \}$.
  \begin{enumerate}
     \item Montrer que : $M \in {\cal M} \iff {\cal H} \text{ et } {\cal K}$
        sont stables par $\phi_M$.
     \item En déduire dim($\cal M$).
  \end{enumerate}
\end{enumerate}
\finenonce{003377}



\finexercice


\exercice{3378, quercia, 2010/03/09}
\enonce{003378}{Matrices triangulaires nilpotentes}

\begin{enumerate}
  \item Soit $A$ une matrice triangulaire à diagonale nulle.
    Montrer que $A$ est nilpotente.
  \item Soit $A \in \mathcal{M}_n(K)$ une matrice nilpotente d'indice $n$ et $\phi$
    l'endomorphisme de $ K^n$ associé.

    On note $E_i = \mathrm{Ker} \phi^i$, et $\vec e_i$ un vecteur quelconque choisi
    dans $E_i \setminus E_{i-1}$ $(\vec e_1 \in E_1 \setminus \{\vec 0 \})$.
  \begin{enumerate}
     \item Justifier l'existence de $\vec e_i$.
     \item Montrer que la famille $(\vec e_i)$ est une base de $ K^n$.
     \item En déduire que $A$ est semblable à une matrice triangulaire à diagonale nulle.
  \end{enumerate}
\end{enumerate}
\finenonce{003378}



\finexercice

\exercice{3383, quercia, 2010/03/09}
\enonce{003383}{Matrice vérifiant $A^k = I$}

Soit $A \in \mathcal{M}_n(K)$ telle que $A^k = I\ (k \ne 0)$.
On pose $B = I + A + A^2 + \dots + A^{k-1}$.
Soient $u,v$ les endomorphismes de $ K^n$ matrices $A$ et $B$ dans la base
canonique.

\begin{enumerate}
  \item Montrer que : $\mathrm{Ker} (u-\mathrm{id}) = \Im v$,\quad $\Im(u-\mathrm{id}) = \mathrm{Ker} v$,\quad
    $\mathrm{Ker} v \oplus \Im v =  K^n$.
  \item En déduire : $\mathrm{tr} B = k\mathrm{rg} B$.
     

\end{enumerate}
\finenonce{003383}



\finexercice


\exercice{3384, quercia, 2010/03/09}
\enonce{003384}{$A > 0, X > 0$ et $A^kX = X$}

Soit $A \in \mathcal{M}_{n,p}(\R)$. On dit que $A$ est positive si tous ses
coefficients sont strictement positifs.

Soit $M \in \mathcal{M}_n(\R)$ positive. On suppose qu'il existe $X \in \mathcal{M}_{n,1}(\R)$
positif et $k\in\N^*$ tels que $M^kX = X$.
Montrer qu'il existe $Y \in \mathcal{M}_{n,1}(\R)$ positif tel que $MY = Y$.
\finenonce{003384}



\finexercice


\exercice{3385, quercia, 2010/03/09}
\enonce{003385}{Suite récurrente linéaire matricielle}

Soient $A,B \in \mathcal{M}_n(K)$. Exprimer en fonction de $k$ le terme général de la
suite $(M_k)$ de matrices de $\mathcal{M}_n(K)$ définie par :
$\begin{cases} M_0 \text{ est donnée,}\cr M_{k+1} = AM_k + B. \cr\end{cases}$
\finenonce{003385}



\finexercice


\exercice{3386, quercia, 2010/03/09}
\enonce{003386}{$A, A^2, A^3$ données $ \Rightarrow  A^p$}

Soit $A \in \mathcal{M}_n(K)$. On suppose qu'il existe $\lambda,\mu \in  K$ et
$U,V \in \mathcal{M}_n(K)$ tels que :
$\begin{cases} A   = \lambda  U + \mu  V \cr
         A^2 = \lambda^2U + \mu^2V \cr
         A^3 = \lambda^3U + \mu^3V.\cr\end{cases}$

\begin{enumerate}
  \item Montrer que : $\forall\ p \in \N^*,\  A^p = \lambda^pU + \mu^pV$
    (chercher une relation linéaire entre $A$, $A^2$, $A^3$).
    

  \item On suppose ici $\lambda \ne \mu$, $\lambda\ne 0$ et $\mu\ne 0$.
    Soit $X$ un vecteur propre de $A$. Montrer que $X$ est vecteur propre de
    $U$ et de $V$ avec les valeurs propres $0,0$ ou $1,0$, ou $0,1$.
    
\end{enumerate}
\finenonce{003386}



\finexercice


\exercice{3387, quercia, 2010/03/09}
\enonce{003387}{Idéaux de $\mathcal{M}_n(K)$}

Une partie ${\cal I} \subset \mathcal{M}_n(K)$ est appelée {\it idéal à droite
de $\mathcal{M}_n(K)$} si c'est un sous-groupe additif vérifiant :
$$\forall\ A \in {\cal I},\ \forall\ B \in \mathcal{M}_n(K),\ AB \in {\cal I}.$$

Pour $A \in \mathcal{M}_n(K)$, on note ${\cal H}_A$ le sev de $\mathcal{M}_{n,1}(K)$ engendré par
les colonnes de $A$, et ${\cal I}_A$ l'idéal à droite engendré par $A$ :
${\cal I}_A = \{ AM \text{ tq } M \in \mathcal{M}_n(K) \}$.


\begin{enumerate}
  \item Soient $A,M \in \mathcal{M}_n(K)$. Montrer que :
    $M \in {\cal I}_A \iff {\cal H}_M \subset {\cal H}_A$.
  \item Soient $A,B \in \mathcal{M}_n(K)$. Montrer qu'il existe $C \in \mathcal{M}_n(K)$ telle que
    ${\cal H}_A + {\cal H}_B = {\cal H}_C$.
    Simplifier ${\cal I}_A + {\cal I}_B$.
  \item Soit $\cal I$ un idéal à droite de $\mathcal{M}_n(K)$. Montrer que $\cal I$ est un sev
    de $\mathcal{M}_n(K)$, puis qu'il existe $A \in \mathcal{M}_n(K)$ telle que
    ${\cal I} = {\cal I}_A$.
  \item Que peut-on dire des idéaux {\it à gauche\/} de $\mathcal{M}_n(K)$ ?


\end{enumerate}
\finenonce{003387}



\finexercice


\exercice{3388, quercia, 2010/03/09}
\enonce{003388}{Classes d'équivalence dans $\mathcal{M}_{n,1}(\Z)$}

\begin{enumerate}
  \item Soit $M \in \mathcal{M}_n(\Z)$. Montrer que $M \in GL_n(\Z)$ si et seulement si
    $|\det M| = 1$.
  \item Soit $X = \begin{pmatrix} x_1 \cr \vdots \cr x_n \cr \end{pmatrix} \in \mathcal{M}_{n,1}(\Z)$ et
    $d$ le pgcd de $x_1, \dots, x_n$.
    Montrer qu'il existe $A \in GL_n(\Z)$ telle que
    $AX = \begin{pmatrix}d \cr 0\cr \vdots \cr 0 \cr \end{pmatrix}$ (par récurrence sur $n$).
  \item Soient $X,Y \in \mathcal{M}_{n,1}(\Z)$. CNS pour qu'il existe $A \in GL_n(\Z)$
    telle que $AX = Y$ ?
\end{enumerate}
\finenonce{003388}



\finexercice


\exercice{3389, quercia, 2010/03/09}
\enonce{003389}{Rayon spectral d'une matrice à coefficients positifs}

Soit $A = (a_{ij}) \in \mathcal{M}_n(\R)$ avec : $\forall\ i,j,\ a_{ij} > 0$.
On munit $\mathcal{M}_{n,1}(\R)$ de la relation d'ordre :
$$(X \ge Y) \iff (\forall\ i,\ x_i \ge y_i),$$
et on pose pour $X \in \mathcal{M}_{n,1}(\R)$, $X \ge 0$, $X\ne 0$ :
$$\left\{
\begin{array}{lll} R(X)    &=& \text{sup}\{ r \ge 0 \text{ tq } AX \ge rX \},    \cr
         \hfil R &=& \text{sup}\{R(X) \text{ tq } X \ge 0, X \ne 0 \}. \cr
         \end{array}\right.$$

\begin{enumerate}
  \item Montrer que $R$ est fini et qu'il existe $X_0 \in \R^n$ tel que
    $R(X_0) = R$.
    

  \item Montrer que toutes les coordonnées de $X_0$ sont strictement positives.
  \item On pose $AX_0 = RX_0 + Y$. Montrer que $Y = 0$.
  \item Soit $\lambda$ une valeur propre complexe de $A$.
    Montrer que $|\lambda| \le R$, et $(|\lambda| = R) \Leftrightarrow (\lambda = R)$.
    
\end{enumerate}
\finenonce{003389}



\finexercice


\exercice{3390, quercia, 2010/03/09}
\enonce{003390}{INT ingénieurs 93}

Soit $E = \{ \text{matrices de } \mathcal{M}_n(\R) \text{ antisymétriques}\}$ et
$f : E \to  E,  M \mapsto {^t\!AM+MA}$ où $A \in \mathcal{M}_n(\R)$.

\begin{enumerate}
  \item Montrer que $f$ est un endomorphisme.
  \item Quelle est la trace de $f$ ?

\end{enumerate}
\finenonce{003390}



\finexercice


\exercice{3391, quercia, 2010/03/09}
\enonce{003391}{Ensam PSI 1998}

Soit $A = \begin{pmatrix}a_1    &1      &0       &\dots  &0      \cr
                   a_2    &0      &1       &\ddots &\vdots \cr
                   \vdots &\vdots &\ddots  &\ddots &0      \cr
                   \vdots &\vdots &        &\ddots &1      \cr
                   a_n    &0      &\dots   &\dots  &0      \cr\end{pmatrix}
\in \mathcal{M}_n(K)$
et $\mathcal{C}(A)$ son commutant.

Montrer que pour $M,N \in \mathcal{C}(A)$ on a~: $M=N \Leftrightarrow M$ et $N$ ont la même
dernière colonne.

En déduire que $\mathcal{C}(A) =  K_{n-1}[A]$.
\finenonce{003391}



\finexercice


\exercice{3392, quercia, 2010/03/09}
\enonce{003392}{ENS MP 2002}
Que dire des morphismes de groupe $\varphi : {GL_n(\R)} \to {\Z/p\Z}$~?
\finenonce{003392}



\finexercice


\exercice{5264, rouget, 2010/07/04}
\enonce{005264}{***}
\begin{enumerate}
\item  Montrer qu'une matrice triangulaire supérieure est inversible si et seulement si ses coefficients diagonaux sont tous non nuls.
\item  Montrer que toute matrice triangulaire supérieure est semblable à une matirce triangulaire inférieure.
\end{enumerate}
\finenonce{005264}


\finexercice


\exercice{5265, rouget, 2010/07/04}
\enonce{005265}{***}
Soient $I=\left(
\begin{array}{cc}
1&0\\
0&1
\end{array}
\right)$ et $J=\left(
\begin{array}{cc}
1&1\\
0&1
\end{array}
\right)$ puis $E=\{M(x,y)=xI+yJ,\;(x,y)\in\Rr^2\}$.

\begin{enumerate}
\item  Montrer que $(E,+,.)$ est un sous-espace vectoriel de $\mathcal{M}_2(\Rr)$. Déterminer une base de $E$ et sa dimension.
\item  Montrer que $(E,+,\times)$ est un anneau commutatif.
\item  Quels sont les inversibles de $E$~?
\item  Résoudre dans $E$ les équations suivantes~:

$$a)\;X^2=I\quad b)\;X^2=0\quad c)\;X^2 = X.$$

\item  Calculer $(M(x,y))^n$ pour $n$ entier naturel non nul.
\end{enumerate}
\finenonce{005265}


\finexercice


\exercice{5266, rouget, 2010/07/04}
\enonce{005266}{****}
Soit $A\in\mathcal{M}_{3,2}(\Rr)$ et $B\in\mathcal{M}_{2,3}(\Rr)$ telles que~:

$$AB=
\left(
\begin{array}{ccc}
0&-1&-1\\
-1&0&-1\\
1&1&2
\end{array}
\right)
.$$

Montrer l'existence d'au moins un couple $(A,B)$ vérifiant les conditions de l'énoncé puis calculer $BA$. (Indication. Calculer $(AB)^2$ et utiliser le rang.)
\finenonce{005266}


\finexercice


\exercice{5268, rouget, 2010/07/04}
\enonce{005268}{***I}
Déterminer le centre de $\mathcal{M}_n(\Kk)$, c'est à dire l'ensemble des éléments de $\mathcal{M}_n(\Kk)$ qui commutent avec tous les éléments de $\mathcal{M}_n(\Kk)$ (utiliser les matrices élémentaires).
\finenonce{005268}


\finexercice


\exercice{5270, rouget, 2010/07/04}
\enonce{005270}{****}
Montrer que tout hyperplan de $\mathcal{M}_n(\Kk)$ $(n\geq 2)$ contient au moins une matrice inversible.
\finenonce{005270}


\finexercice


\exercice{5271, rouget, 2010/07/04}
\enonce{005271}{***}
Soit $f$ qui, à $P\in\Rr_{2n}[X]$ associe $f(P)=X(X+1)P'-2kXP$.
Trouver $k$ tel que $f\in\mathcal{L}(\Rr_{2n}[X])$ puis, pour cette valeur de $k$, trouver tous les polynômes $P$ non nuls tels que la famille $(P,f(P))$ soit liée.
\finenonce{005271}



\finexercice


\exercice{5273, rouget, 2010/07/04}
\enonce{005273}{***I}
Calculs par blocs.
\begin{enumerate}
\item  Soit $M=
\left(
\begin{array}{cc}
A&B\\
C&D
\end{array}
\right)
$ et $N=\left(
\begin{array}{cc}
A'&B'\\
C'&D'
\end{array}
\right)$ avec $(A,A')\in(\mathcal{M}_{p,r}(\Kk))^2$, $(B,B')\in(\mathcal{M}_{p,s}(\Kk))^2$, $(C,C')\in(\mathcal{M}_{q,r}(\Kk))^2$ et $(D,D')\in(\mathcal{M}_{q,s}(\Kk))^2$. Calculer $M+N$ en fonction de $A$, $B$, $C$, $D$, $A'$, $B'$, $C'$ et $D'$.
\item  Question analogue pour $MN$ en analysant précisément les formats de chaque matrice.
\end{enumerate}
\finenonce{005273}


\finexercice

\exercice{5275, rouget, 2010/07/04}
\enonce{005275}{***}
Soit $A=
\left(
\begin{array}{cccc}
7&4&0&0\\
-12&-7&0&0\\
20&11&-6&-12\\
-12&-6&6&11
\end{array}
\right)
$ et $u$ l'endomorphisme de $\Cc^4$ de matrice $A$ dans la base canonique de $\Cc^4$.
\begin{enumerate}
\item  Déterminer une base de $\Cc^4$ formée de vecteurs colinéaires à leurs images.
\item  Ecrire les formules de changement de base correspondantes.
\item  En déduire le calcul de $A^n$ pour $n$ entier naturel.
\end{enumerate}
\finenonce{005275}


\finexercice


\exercice{5585, rouget, 2010/10/16}
\enonce{005585}{***I}
Soient $A\in\mathcal{M}_{3,2}(\Rr)$ et $B\in\mathcal{M}_{2,3}(\Rr)$ telles que $AB=\left(
\begin{array}{ccc}
8&2&-2\\
2&5&4\\
-2&4&5
\end{array}
\right)$. Justifier l'existence de $A$ et $B$ puis calculer $BA$.
\finenonce{005585}


\finexercice

\exercice{5594, rouget, 2010/10/16}
\enonce{005594}{***}
\label{ex:rou32}
Soit $G$ un sous-groupe fini de $GL_n(\Rr)$ tel que $\sum_{M\in G}^{}\text{Tr}(M)=0$. Montrer que $\sum_{M\in G}^{}M=0$. 
\finenonce{005594}


\finexercice


\exercice{5595, rouget, 2010/10/16}
\enonce{005595}{***}
Soit $G$ un sous-groupe de $GL(E)$ avec $\text{dim}E=n$ et $\text{card}G=p$.
Soit $F=\{x\in E/\;\forall g\in G,\;g(x)=x\}$.

Montrer que $\text{dim}F= \frac{1}{p}\sum_{g\in G}^{}\text{Tr}g$.
\finenonce{005595}


\finexercice


\exercice{5596, rouget, 2010/10/16}
\enonce{005596}{***I}
Soient $A_1$,..., $A_p$ $p$ matrices distinctes et inversibles de $M_n(\Rr)$ telles que $G=\{A_1,...,A_p\}$ soit stable pour la multiplication. Soit $A=A_1 + ... + A_p$. Montrer que $\text{Tr}A$ est un entier divisible par $p$.
\finenonce{005596}


\finexercice


\exercice{5605, rouget, 2010/10/16}
\enonce{005605}{***}
Soient $n$ un entier naturel non nul puis $A\in\mathcal{M}_n(\Kk)$.
Soit $f$ l'endomorphisme de $M_n(\Kk)$ qui à une matrice $X$ associe $AX+XA$. Calculer $\text{Tr}(f)$.
\finenonce{005605}


\finexercice


\exercice{5606, rouget, 2010/10/16}
\enonce{005606}{**}
Soient $a$ un réel non nul et $A$ et $B$ deux éléments de $\mathcal{M}_n(\Rr)$.

Résoudre dans $\mathcal{M}_n(\Rr)$ l'équation d'inconnue $M$ : $aM+ \text{Tr}(M)A=B$.
\finenonce{005606}


\finexercice


\exercice{5608, rouget, 2010/10/16}
\enonce{005608}{**}
Soient $I=\left(
\begin{array}{cc}
1&0\\
0&1
\end{array}
\right)$ et $J=\left(
\begin{array}{cc}
1&1\\
0&1
\end{array}
\right)$. Soit $E=\{M(x,y)=xI+yJ,\;(x,y)\in\Rr^2\}$.

\begin{enumerate}
 \item  Montrer que $(E,+,.)$ est un $\Rr$-espace vectoriel et préciser sa dimension.

\item  Montrer que $(E,+,\times)$ est un anneau commutatif.

\item  Quels sont les éléments inversibles de l'anneau $(E,+,\times)$ ?

\item  Résoudre dans $E$ les équations : 
  \begin{enumerate}
  \item $X^2=I$ 
  \item $X^2=0$ 
  \item $X^2=X$.
  \end{enumerate}
\item  Calculer $(M(x,y))^n$ pour $n$ entier naturel et $x$ et $y$ réels.
\end{enumerate}
\finenonce{005608}


\finexercice


\exercice{5609, rouget, 2010/10/16}
\enonce{005609}{***}
On appelle idéal bilatère de l'anneau $(\mathcal{M}_n(\Kk),+,\times)$ tout sous-ensemble I de $\mathcal{M}_n(\Kk)$ tel que

\begin{center}
a) $(I,+)$ est un groupe et b) $\forall A\in I$,  $\forall M\in\mathcal{M}_n(\Kk)$, $AM\in I$ et $MA\in I$.
\end{center}

Déterminer tous les idéaux bilatères de l'anneau $(\mathcal{M}_n(\Kk),+,\times)$.
\finenonce{005609}


\finexercice


\exercice{5618, rouget, 2010/10/16}
\enonce{005618}{*I}
Existe-t-il deux matrices carrées $A$ et $B$ telles que $AB-BA=I_n$.
\finenonce{005618}


\finexercice


\exercice{5619, rouget, 2010/10/16}
\enonce{005619}{**I}
Soit $f$ une forme linéaire sur $\mathcal{M}_n(\Cc)$ telle que $\forall (A,B)\in(\mathcal{M}_n(\Cc))^2$, $f(AB) = f(BA)$. Montrer qu'il existe un complexe $a$ tel que $f=a\text{Tr}$.
\finenonce{005619}


\finexercice


\exercice{5620, rouget, 2010/10/16}
\enonce{005620}{***}
Soit $A_n=
\left(
\begin{array}{cc}
1&-\frac{a}{n}\\
\frac{a}{n}&1
\end{array}
\right)$ ($a$ réel donné). Calculer $\lim_{n \rightarrow +\infty}A_n^n$.
\finenonce{005620}


\finexercice


\exercice{5621, rouget, 2010/10/16}
\enonce{005621}{**}
Soient $A$ une matrice carrée de format $n$ et $f$ l'application de $\mathcal{M}_n(\Cc)$ dans lui-même qui à une matrice $M$ associe $MA$. Trouver la matrice de $f$ dans la base canonique de $\mathcal{M}_n(\Cc)$ (ordonnée par l'ordre lexicographique).
\finenonce{005621}


\finexercice


\exercice{5625, rouget, 2010/10/16}
\enonce{005625}{***I}
Soient $A$ et $B$ deux matrices carrées de format $n$ telles que $AB-BA=A$. Calculer la trace de $A^{2010}$.
\finenonce{005625}


\finexercice


\exercice{5628, rouget, 2010/10/16}
\enonce{005628}{**I Exponentielle d'une matrice nilpotente}

Pour $A$ matrice nilpotente donnée, on pose $\text{exp}A=\sum_{k=0}^{+\infty}\frac{A^k}{k!}$.

\begin{enumerate}
 \item  Montrer que si $A$ et $B$ commutent et sont nilpotentes alors $A+B$ est nilpotente et $\text{exp}(A+B) =\text{exp}A\times\text{exp}B$.

\item  Montrer que $\text{exp}A$ est inversible.

\item  Calculer $\text{exp}A$ où $A=\left(
\begin{array}{ccccc}
0&1&0&\ldots&0\\
\vdots&\ddots&\ddots&\ddots&\vdots\\
 & & &\ddots&0\\
\vdots& & &\ddots&1\\
0&\ldots& &\ldots&0
\end{array}
\right)$.
\end{enumerate}
\finenonce{005628}


\finexercice

\section{120.01 Les rationnels}


\exercice{451, bodin, 1998/09/01}
\video{1d2LZ6zGdjg}
\enonce{000451}{}
\begin{enumerate}
    \item D\'emontrer que si $r \in \Q$ et $ x \notin \Q $ alors $ r+x
\notin \Q $ et si $r\not= 0$ alors $ r.x \notin \Q $.
    \item Montrer que $\sqrt 2 \not\in\Q$,
    \item En d\'eduire : entre deux nombres rationnels il y a toujours un nombre irrationnel.
\end{enumerate}
\finenonce{000451} 


\finexercice

\exercice{452, cousquer, 2003/10/01}

\enonce{000452}{}
 Les nombres suivants sont-ils des rationnels ? des décimaux ?
\\$a=1/3$,\quad 
$b=1/15$,\quad
$c=1/25$,\quad
$d=1/125$,\quad
$e$,\quad
$f=0,333\cdots3\cdots$,\quad 
$g=\sqrt2$,\quad\\
$h=0{,}123\,456\,789\,123\,456\,789\,123\cdots$,\quad 
$i=0{,}123\,456\,789\,101\,112\,131\,4\cdots$,\quad
$j=\pi$,\quad 
$k=13/7$,\quad 
$l=27/17$.
\finenonce{000452}



\finexercice



\exercice{453, cousquer, 2003/10/01}

\enonce{000453}{Un procédé géométrique d'approximation de $\sqrt2$}
Dans le plan $xOy$, on porte sur $Ox$ une suite de points
$a_1, a_2,\ldots,a_n,\ldots$
et sur $Oy$ une suite de points $b_1, b_2,\ldots,b_n,\ldots\,$, construites
de la manière suivante :
\begin{description}
\item[(i)]  $a_1=2$  et $b_1=1$, 

\item[(ii)]  $a_n= {a_{n-1}+b_{n-1} \over2}$,

\item[(iii)]  $a_nb_n=2$ (le rectangle de côtés $a_n$ et $b_n$
           a pour aire 2).
\end{description}
\begin{enumerate}
\item Représentez cette suite de rectangles de côtés $a_n$ et $b_n$.
\item Démontrez successivement que :
$\forall n,\; b_n < a_n$ ; $(a_n)_{n \in\bf N}$ décroissante ;
$(b_n)_{n\in\bf N}$ croissante.
\item Calculez $a_n-b_n$ en fonction de $a_{n-1}-b_{n-1}$ et $a_n$.
Montrez que l'on a l'inégalité :
$$ a_n-b_n < {{(a_{n-1}-b_{n-1})}^2 \over4}.$$
\item Calculez les premiers termes de la suite $a_1, a_2,\ldots,a_6$.
Combien de décimales exactes de $\sqrt2$ obtenez-vous à chaque pas ?
Utilisez l'inégalité précédente pour montrer que le nombre
de décimales  exactes obtenues double \emph{grosso modo} à chaque pas.
\end{enumerate}
\finenonce{000453}



\finexercice



\exercice{454, cousquer, 2003/10/01}

\enonce{000454}{}
 Calculer avec une calculette :
$\frac{1}{3}+\frac{1}{3}+\frac{1}{3}\quad\mbox{et}\quad
1-\frac{1}{3}-\frac{1}{3}-\frac{1}{3}.$
Expliquer le résultat.
\finenonce{000454}



\finexercice



\exercice{455, cousquer, 2003/10/01}

\enonce{000455}{}
 On considère les nombres rationnels inférieurs à $\sqrt{2}$.
Y a-t-il un nombre rationnel juste avant $\sqrt{2}$, plus grand que tous les
nombres rationnels inférieurs à $\sqrt{2}$ ?

\noindent Une suite de nombres rationnels a-t-elle pour limite un nombre rationnel ?

\noindent Une suite de nombres décimaux a-t-elle pour limite un nombre décimal ?
\finenonce{000455}



\finexercice



\exercice{456, ridde, 1999/11/01}

\enonce{000456}{}
Soient $a$ et $b$ deux rationnels positifs tels que $\sqrt{a}$ et $\sqrt{b}$
soient irrationnels. Montrer que $\sqrt{a} + \sqrt{b}$ est irrationnel.
\finenonce{000456}



\finexercice



\exercice{457, bodin, 1998/09/01}
\video{KX375CPpZjU}
\enonce{000457}{}
 Soit $p(x) = \sum_{i=0}^{n} a_i \cdot x^i$. On suppose que tous les
$a_i$ sont des entiers.
\begin{enumerate}
    \item  Montrer que si $p$ a une racine rationnelle $\frac{\alpha}{\beta}$ (avec $\alpha$ et $\beta$ premiers entre eux)
alors $\alpha$ divise $a_0$ et $\beta$ divise $a_n$.

    \item On consid\`ere le nombre $\sqrt 2+\sqrt 3$. En calculant son carr\'e, montrer que ce
carr\'e est racine d'un polyn\^ome de degr\'e 2. En d\'eduire, \`a
l'aide du r\'esultat pr\'ec\'edent qu'il n'est pas rationnel.
\end{enumerate}
\finenonce{000457} 


\finexercice

\exercice{458, bodin, 1998/09/01}

\enonce{000458}{}
Trouver sous la forme $\frac{p}{q}$ des rationnels $x$
dont les d\'evelopements d\'ecimaux p\'eriodiques sont donn\'es par :
\par $3,14\!\buildrel\frown\over{14}...\ ;\ \ \ \
0,99\!\!\buildrel\frown\over{9}...\ ;\ \ \ \ 3,149\!\!\buildrel\frown\over{9}...  $
\finenonce{000458}



\finexercice



\exercice{459, bodin, 1998/09/01}
\video{0MkdQiQ4ceI}
\enonce{000459}{}
\begin{enumerate}
    \item Soit $N_n = 0,1997\,1997\ldots 1997$ ($n$ fois).
Mettre $N_n$ sous la forme $\frac{p}{q}$ avec $p,q \in \Nn^*$.
    \item Soit $M = 0,1997\,1997\,1997\ldots\ldots$ Donner
le rationnel dont l'\'ecriture d\'ecimale est $M$.
    \item M\^eme question avec :
$ P = 0,11111\ldots + 0,22222\ldots +0,33333\ldots
+0,44444\ldots+0,55555\ldots+0,66666\ldots
+0,77777\ldots + 0,88888\ldots+0,99999\ldots $
\end{enumerate}
\finenonce{000459} 


\finexercice

\exercice{460, bodin, 1998/09/01}

\enonce{000460}{}
 Montrer que l'ensemble
$\{r^3\ ;\ r\in \Q\}$ est dense dans $\R$.
\finenonce{000460}



\finexercice



\exercice{461, gourio, 2001/09/01}
\video{cHENuXePV9g}
\enonce{000461}{}
Montrer que $\frac{\ln 3}{\ln 2}$ est irrationnel.
\finenonce{000461} 


\finexercice

\exercice{462, gourio, 2001/09/01}

\enonce{000462}{}
Soit $a\in \Rr$, montrer :
$$\exists (p,q)\in \Zz\times \Nn^{*};\left| a-\frac{p}{q}\right| \leq \frac{1}{q^{2}}.$$
\emph{Indication} : consid\'{e}rer les parties fractionnaires de $0,a,2a,...,qa $ et
la partition $[0,\frac{1}{q}[,[\frac{1}{q},\frac{2}{q}[,...[\frac{q-1}{q},1[$
de $[0,1[.$
\finenonce{000462}



\finexercice



\exercice{463, gourio, 2001/09/01}

\enonce{000463}{}
Montrer que l'ensemble des nombres dyadiques :
$$\left\{ \frac{a}{2^{k}},(a,k)\in \Zz\times \Nn\right\}$$
est dense dans $\Rr$.
\finenonce{000463}



\finexercice



\exercice{506, bodin, 1998/09/01}
\video{lyrZpWFC8AM}
\enonce{000506}{}
 Montrer que toute suite convergente est born\'ee.
\finenonce{000506} 


\finexercice

\exercice{507, bodin, 1998/09/01}
\video{HI2i2rdz3_A}
\enonce{000507}{}
  Montrer que la suite $(u_n)_{n\in\Nn}$ d\'efinie par
$$u_n = (-1)^n+\frac{1}{n}$$
n'est pas convergente.
\finenonce{000507} 


\finexercice


\exercice{508, cousquer, 2003/10/01}

\enonce{000508}{}
 Étudier la suite $ u_n=\frac{a^n-b^n}{a^n+b^n}$,
$a$ et $b$ étant donnés dans $\mathbb{R}_+^*$.

\finenonce{000508}



\finexercice



\exercice{509, cousquer, 2003/10/01}

\enonce{000509}{}
 Les énoncés suivants sont-ils vrais ou faux~?

\begin{enumerate}
\item Si une suite positive est non majorée, elle tend vers $+\infty$.

\item Si une suite d'entiers converge, elle est stationnaire.

\item Si une suite a un nombre fini de valeurs, elle converge si et seulement
si elle est stationnaire.

\item Une suite est convergente si et seulement si elle est bornée.

\item Si une suite n'est pas majorée, elle est minorée.
\end{enumerate}
\finenonce{000509}



\finexercice



\exercice{510, cousquer, 2003/10/01}

\enonce{000510}{}
 Soit $l$ un nombre réel. Peut-on dire qu'une suite qui vérifie 
$$\forall\epsilon\in\mathopen]0,1\mathclose[,\; 
\exists N\in\mathbb{N},\;
\forall n>N,\; \vert u_n-l\vert<\epsilon$$
converge vers $l$~?
\finenonce{000510}



\finexercice



\exercice{511, cousquer, 2003/10/01}

\enonce{000511}{}
 Construire une suite $u_n=v_nw_n$ (resp. $v_n+w_n$) convergente
et telle que l'une au moins des suites $(v_n)$ et $(w_n)$ diverge.
\finenonce{000511}



\finexercice



\exercice{3066, quercia, 2010/03/08}

\enonce{003066}{Nombres irrationnels}
Soit $a \in \Q^+$ tel que $\sqrt a \notin \Q$.

Montrer qu'il existe $C>0$ tel que pour tout rationnel $r=\frac pq$, on a :
$\bigl|r-\sqrt a\bigr| \ge \frac C{q^2}$.

\finenonce{003066}


\finexercice 


\exercice{3067, quercia, 2010/03/08}

\enonce{003067}{Nombres irrationnels}
Soient $a,b \in \Q^+$ tels que $\sqrt b \notin \Q^+$. Montrer qu'il existe
$x,y \in \Q^+$ tels que $\sqrt x + \sqrt y = \sqrt{a+\sqrt b}$ si et seulement
si $a^2-b$ est un carr{\'e} dans $\Q$.
\finenonce{003067}


\finexercice 


\exercice{3144, quercia, 2010/03/08}

\enonce{003144}{Parties fractionnaires}
Soit $x=\frac pq \in \Q^*$ avec $p,q$ entiers, $q\ge 1$, $p\wedge q = 1$.
Calculer $\sum_{k=0}^{q-1} \text{frac}(kx)$.
\finenonce{003144}


\finexercice 


\exercice{3145, quercia, 2010/03/08}

\enonce{003145}{D{\'e}nominateurs dans un sous-anneau}
Soit $A$ un sous-anneau de $\Q$. On {\'e}crit les {\'e}l{\'e}ments de $A$ sous forme
irr{\'e}ductible; soit $P$ l'ensemble des d{\'e}nominateurs.
Montrer que $A = \left\{ \frac mp \text{ tels que } m \in \Z,\ p \in P\right\}$.
\finenonce{003145}


\finexercice 


\exercice{3146, quercia, 2010/03/08}

\enonce{003146}{Les sous-anneaux de $\Q$ sont principaux}
Soit $A$ un sous-anneau de $\Q$. Montrer que $A$ est principal
(si $I$ est un id{\'e}al de $A$, consid{\'e}rer $I\cap \Z$).
\finenonce{003146}



\finexercice 


\exercice{3147, quercia, 2010/03/08}

\enonce{003147}{D{\'e}composition en inverses}
Soit $x\in\Q$, $0<x<1$. On d{\'e}finit une suite $(x_n)$ de rationnels par
r{\'e}currence :
\begin{itemize}
 \item $x_0 = x$,
 \item  Si $x_n$ existe et est non nul, soit $k_n \in \N^*$ le plus
              petit entier tel que $\frac1{k_n}\le x_n$. On pose
              $x_{n+1} = x_n - \frac1{k_n}$,
 \item Si $x_n = 0$, on s'arr{\^e}te. Dans ce cas,
              $x = \frac 1{k_0} + \frac 1{k_1} + \dots + \frac1{k_{n-1}}$.
\end{itemize}

\begin{enumerate}
  \item Montrer que la suite est toujours finie.

  \item Montrer que si $k_{i+1}$ existe, alors $k_{i+1} > k_i(k_i-1)$.

  \item R{\'e}ciproquement, soit une d{\'e}composition :
    $x = \frac1{n_0} + \dots + \frac1{n_p}$ avec $n_i \in \N^*$ et
    $n_{i+1} > n_i(n_i-1)$. Montrer que pour tout $i$, on a $n_i = k_i$.
\end{enumerate}
\finenonce{003147}


\finexercice 


\exercice{3148, quercia, 2010/03/08}

\enonce{003148}{Combinaison de fractions}
Soient $\frac ab < \frac cd$ deux rationnels avec $a,c \in \Z$, et $b,d \in \N^*$.
\begin{enumerate}
  \item Prouver que tout rationnel s'{\'e}crit : $x = \frac{ma+nc}{mb+nd}$ avec $m,n\in \Z$,
    et $mb+nd \ne 0$.

  \item {\'E}tudier l'unicit{\'e} d'une telle {\'e}criture.

  \item Montrer que $\frac{ma+nc}{mb+nd}$ est compris entre $\frac ab$ et $\frac cd$
    si et seulement si $m$ et $n$ ont m{\^e}me signe.
\end{enumerate}
\finenonce{003148}


\finexercice 


\exercice{3149, quercia, 2010/03/08}

\enonce{003149}{{\'E}quations alg{\'e}briques}
D{\'e}terminer $x \in \Q$ sachant que :
\begin{enumerate}
  \item $2x^3-x^2+x+1 = 0$.        
  \item $6x^5+11x^4-x^3+5x-6 = 0$. 
  \item $2x^3-x-4 = 0$.   
\end{enumerate}
\finenonce{003149}


\finexercice 


\exercice{3150, quercia, 2010/03/08}

\enonce{003150}{$x^y = y^x$}
On cherche les couples $(x,y) \in \bigl(\Q^{+*}\bigr)^2$ tels que $x < y$
et $x^y = y^x$ ($x^y,y^x\in\R$).

On pose $x = \frac pq$, $y = \frac {p'}{q'}$ (formes irr{\'e}ductibles),
$d = pq'\wedge p'q$,\quad $pq' = ad$ et $p'q = bd$.
\begin{enumerate}
  \item Montrer qu'il existe $m,n\in\N^*$ tels que :
    $p=m^a$, $p'=m^b$, $q=n^a$ et $q'=n^b$.

  \item En d{\'e}duire : $b-a = m^{b-a} - n^{b-a}$.

  \item Montrer que $b-a \le 1$ et conclure.
\end{enumerate}
\finenonce{003150}


\finexercice 


\exercice{5209, rouget, 2010/06/30}
\enonce{005209}{I}
Montrer que les nombres suivants sont irrationnels.

\begin{enumerate}
\item (**) $\sqrt{2}$ et plus généralement $\sqrt[n]{m}$ où $n$ est un entier supérieur ou égal à $2$ et $m$ est un entier naturel supérieur ou égal à $2$, qui n'est pas une puissance $n$-ième parfaite. 
\item (**) $\log 2$.
\item (****) $\pi$ (\textsc{Lambert} a montré en 1761 que $\pi$ est irrationnel, \textsc{Legendre} a démontré en 1794 que $\pi^2$ est irrationnel, \textsc{Lindemann} a démontré en 1882 que $\pi$ est transcendant).

Pour cela, supposer par l'absurde que $\pi=\frac{p}{q}$ avec $p$ et $q$ entiers naturels non nuls et premiers entre eux. Considérer alors $I_n=\int_{0}^{p/q}\frac{x^n(p-qx)^n}{n!}\sin x\;dx$, $n\in\Nn^*$ et montrer que $I_n$ vérifie 
\begin{enumerate}
\item $I_n$ est un entier relatif~;
\item $I_n>0$~;
\item $\lim_{n\rightarrow +\infty}I_n=0$ (voir devoir).
\end{enumerate}
\item (***) $e$ (\textsc{Hermite} a démontré en 1873 que $e$ est transcendant. C'est historiquement le premier \og~vrai~\fg~nombre dont on a réussi à démontrer la transcendance).

Pour cela, établir que pour tout entier naturel $n$, $e=\sum_{k=0}^{n}\frac{1}{k!}+\int_{0}^{1}\frac{(1-t)^n}{n!}e^t\;dt$, puis que \textbf{pour tout} entier naturel non nul $n$, 
$0<e-\sum_{k=0}^{n}\frac{1}{k!}<\frac{3}{(n+1)!}$. Raisonner alors par l'absurde.
\item (***) $\cos(\frac{2\pi}{7})$. Pour cela trouver une équation du troisième degré à coefficients entiers dont les solutions sont $\cos(\frac{2\pi}{7})$, $\cos(\frac{4\pi}{7})$ et $\cos(\frac{6\pi}{7})$, puis vérifier que cette équation n'a pas de racine rationnelle (supposer par l'absurde qu'il y a une racine rationnelle $\frac{p}{q}$ avec $p\in\Zz^*$, $q\in\Nn^*$ et $\mbox{PGCD}(p,q)=1$ et montrer que $p$ divise $1$ et $q$ divise $8$). (On rappelle le théorème de \textsc{Gauss}~:~soient $a$, $b$ et $c$ trois entiers relatifs tous non nuls. Si $a$ divise $bc$ et $a$ et $b$ sont premiers entre eux, alors $a$ divise $c$).

\item (***) $\sqrt{2}+\sqrt{3}+\sqrt{5}$.
\end{enumerate}
\finenonce{005209}


\finexercice

\exercice{5214, rouget, 2010/06/30}
\enonce{005214}{****}
Soit $u_n$ le chiffre des unités de $C_n^k$, $k$ entier naturel fixé non nul et $n$ entier naturel supèrieur ou égal à $k$. Montrer que le nombre $0,u_ku_{k+1}u_{k+2}...$ est rationnel.
\finenonce{005214}


\finexercice

\exercice{5243, rouget, 2010/06/30}
\enonce{005243}{****}
Soit $(u_n)=\left(\frac{p_n}{q_n}\right)$ avec $p_n\in\Zz$ et $q_n\in\Nn^*$, une suite de rationnels convergeant vers un irrationnel $x$. Montrer que les suites $(|p_n|)$ et $(q_n)$ tendent vers $+\infty$ quand $n$ tend vers $+\infty$.
\finenonce{005243}


\finexercice

\section{120.02 Maximum, minimum, borne supérieure}




\exercice{465, monthub, 2001/11/01}
\video{icn0oS9sQ-0}
\enonce{000465}{}
D{\'e}terminer la borne sup{\'e}rieure et inf{\'e}rieure
(si elles existent) de : $A=\{u_n \mid n\in\N\}$ en posant
$u_n=2^n$ si $n$ est pair et  $u_n=2^{-n}$ sinon.
\finenonce{000465} 


\finexercice

\exercice{466, bodin, 1998/09/01}
\video{P4ovnPBvMNo}
\enonce{000466}{}
 D\'eterminer (s'ils existent) : les majorants, les
minorants, la borne sup\'erieure, la borne inf\'erieure, le plus
grand \'el\'ement, le plus petit \'el\'ement des ensembles
suivants :
$$
[0,1]\cap \Qq \ , \quad ]0,1[\cap\Qq \ ,\quad \Nn \ ,\quad \left\lbrace (-1)^n+\frac{1}{n^2} \mid n\in \Nn^* \right\rbrace.
$$
\finenonce{000466} 


\finexercice

\exercice{467, bodin, 1998/09/01}

\enonce{000467}{}
Soit
$$
I = \left\lbrace x \in \Rr \  | -2 < x +\frac{1}{2x}  \le 2 \right\rbrace.
$$
\begin{enumerate}
    \item Montrer que $I$ est la r\'eunion de deux
intervalles.
    \item D\'eterminer (s'ils existent) : les majorants,
les minorants, la borne sup\'erieure, la borne
inf\'erieure, le plus grand \'el\'ement, le plus petit
\'el\'ement de $I$.
\end{enumerate}
\finenonce{000467}



\finexercice



\exercice{468, cousquer, 2003/10/01}

\enonce{000468}{}
Les ensembles suivants ont-ils une borne supérieure, un plus 
grand élément, une borne inférieure, un plus petit élément, 
dans $\mathbb{D}$, dans $\mathbb{Q}$, dans $\mathbb{R}$, 
(si la question se pose) ?
\begin{enumerate}
\item $[0,3\mathclose[$,
\item $\{0\}\cup\mathopen]1,2]$,
\item $\mathbb{D} \cap[0,1/3]$,
\item $\{x \mid \exists n \in\mathbb{N},\; x=1/n\}$,
\item $\{x\in\mathbb{Q}\mid x^2<2 \}$.
\end{enumerate}
\finenonce{000468}



\finexercice



\exercice{469, cousquer, 2003/10/01}

\enonce{000469}{}
 On considère l'ensemble des nombres de la forme $1+\frac{1}{n}$,
où $n$ décrit l'ensemble des entiers strictement positifs. Cet ensemble
est-il majoré ? Minoré ? A-t-il un plus petit élément ? 
Un plus grand élément ? Justifier vos réponses.
\finenonce{000469}



\finexercice



\exercice{470, cousquer, 2003/10/01}

\enonce{000470}{}
 Étant donné un ensemble $A \subset\mathbb{R}$,
écrire avec des quantificateurs
les propriétés suivantes :
\begin{enumerate}
\item $10$ est un majorant de $A$,
\item $m$ est un minorant de $A$,
\item $P$ n'est pas un majorant de $A$,
\item $A$ est majoré,
\item $A$ n'est pas minoré,
\item $A$ est borné,
\item $A$ n'est pas borné.
\end{enumerate}
\finenonce{000470}



\finexercice



\exercice{471, cousquer, 2003/10/01}

\enonce{000471}{}
 Soit $E$ l'ensemble des réels de la forme $\frac{n-1/n}{n + 1/n}$ avec
$n \in\mathbb{N}^*$. L'ensemble $E$ admet-il une borne inférieure, une borne
supérieure, un plus grand élément, un plus petit élément ?
\finenonce{000471}



\finexercice



\exercice{472, cousquer, 2003/10/01}

\enonce{000472}{}
Soit $E=\{\frac{1}{n}\cos n \mid  n \in\mathbb{N}^* \}$ ; calculer
$\inf E$ et $\sup E$.
\finenonce{000472}



\finexercice



\exercice{473, cousquer, 2003/10/01}

\enonce{000473}{}
 Soient $A$ et $B$ deux parties non vides de $\mathbb{R}$ telles que
pour tout $x$ de $A$ et tout $y$ de $B$ on ait $x \leq y$. Démontrer
que $\sup A$ et $\inf B$  existent et que $\sup A \leq\inf B$.
\finenonce{000473}



\finexercice



\exercice{474, cousquer, 2003/10/01}

\enonce{000474}{}
Soit $\bigl( a_{ij}\bigr)_{(i,j) \in I \times J}$ une famille non vide
et bornée de réels ; comparer :
$$\inf_i (\sup_j a_{ij}) \quad\hbox{avec} \quad\sup_j(\inf_i a_{ij}).$$
\finenonce{000474}



\finexercice



\exercice{475, cousquer, 2003/10/01}

\enonce{000475}{}
 Soit $A$ une partie majorée de $\mathbb{R}$ d'au moins deux éléments et $x$
un élément de $A$.
\begin{enumerate}
\item Montrer que si $x<\sup A$, alors $\sup(A\setminus\{x\})=\sup A$.
\item Montrer que si $\sup(A\setminus\{x\})<\sup A$, alors $x=\sup A$.
\end{enumerate}
\finenonce{000475}



\finexercice



\exercice{476, bodin, 1998/09/01}
\video{72sAcDZMmL8}
\enonce{000476}{}
Soient $A$ et $B$ deux parties born\'ees de $\Rr$.
On note $A+B = \{ a+b \mid (a,b)\in A\times B \}$.
\begin{enumerate}
    \item Montrer que $\sup A + \sup B$ est un majorant de $A+B$.
    \item Montrer que $\sup(A+B)=\sup A + \sup B$.
\end{enumerate}
\finenonce{000476} 


\finexercice

\exercice{477, bodin, 1998/09/01}
\video{HlWsIpnoVLI}
\enonce{000477}{}
 Soit $A$ et $B$ deux parties born\'ees de $\Rr$.
\textbf{Vrai} ou \textbf{faux} ?
\begin{enumerate}
    \item  $A \subset B \Rightarrow \sup A \leqslant \sup B$,
    \item $A \subset B \Rightarrow \inf A \leqslant \inf B$,
    \item $\sup (A\cup B) = \max(\sup A,\sup B)$,
    \item $\sup(A+B) < \sup A + \sup B$,
    \item $\sup(-A) = -\inf A$,
    \item $\sup A +\inf B \leqslant \sup(A+B)$.
\end{enumerate}
\finenonce{000477} 


\finexercice

\exercice{478, gourio, 2001/09/01}

\enonce{000478}{}
Donner la borne sup\'{e}rieure et la borne inf\'{e}rieure (si elles
existent) de l'ensemble:
$$D =\left\lbrace \frac{n-\frac{1}{n}}{n+\frac{1}{n}}|n\in \Nn^{*}\right\rbrace.$$
Cet ensemble admet-il un maximum, un minimum ?
\finenonce{000478}



\finexercice



\exercice{479, gourio, 2001/09/01}

\enonce{000479}{}
Soient $n\in \Nn^{*}$ et $a_{1}\leq a_{2}\leq ...\leq a_{n}$, $n$ nombres r\'{e}els.
Calculer :
$$\inf\limits_{x\in \Rr}\sum\limits_{k=1}^{n}\left| x-a_{i}\right| .$$
\finenonce{000479}



\finexercice



\exercice{480, bodin, 1998/09/01}

\enonce{000480}{}
Soit $f : \Rr \rightarrow \Rr$, $f(x) = x^3-3x$.
Tracer les graphes des fonctions $f,|f|,f_+,f_-$ où :
$f_+ = \text{max}(f,0), f_- = \text{min}(f,0)$.
\finenonce{000480}



\finexercice



\exercice{481, cousquer, 2003/10/01}

\enonce{000481}{}
Si $a= \sup A$, montrer qu'il existe une suite d'éléments de $A$ qui converge vers $a$. Réciproque.
\finenonce{000481}



\finexercice



\exercice{482, cousquer, 2003/10/01}

\enonce{000482}{}
 Soit $A= \mathbb{Q}\cap\mathopen]0,1 \mathclose[$ et $a,b\in\mathbb{R}_+$.
On considère les applications suivantes de $A$ dans $\mathbb{R}_+$ :
$$ f \colon\frac{p}{q} \mapsto\frac{q-p}{q+p}\;; \qquad
g \colon\frac{p}{q} \mapsto\frac{aq+bp}{p+q}$$
Déterminer la borne supérieure et la borne inférieure de $f(A)$
et de $g(A)$.
\finenonce{000482}



\finexercice



\exercice{483, cousquer, 2003/10/01}

\enonce{000483}{}
 Soit $A$ l'ensemble des nombres réels qui peuvent s'écrire
$ x= \frac{2p^2-3q}{p^2 +q}$ pour $p$ et $q$ entiers vérifiant
$0<p<q$.
\begin{enumerate}
\item Montrer que $A$ est minorée par $-3$ et majorée par $2$.
\item Déterminer $\inf A$ et $\sup A$ (pour la borne supérieure on pourra
prendre $q=p+1$).
\end{enumerate}
\finenonce{000483}



\finexercice



\exercice{484, cousquer, 2003/10/01}

\enonce{000484}{}
 Soit $(u_n)_{n\in \mathbb{N}}$ une suite bornée.
On pose $A_p = \sup_{n>p} u_n$ et $B_p= \inf_{n>p} u_n$.
Montrer que $(A_p)_{p\in\mathbb{N}}$ est une suite décroissante bornée
et que $(B_p)_{p\in \mathbb{N}}$ est une suite croissante bornée.
Soit $L = \lim_{p \to\infty} A_p$ et $l= \lim_{p \to\infty} B_p$.
\begin{enumerate}
\item Dans le cas particulier où $u_n= {n+2\over n+1} \cos{n\pi\over3}$,
calculer $L$ et $l$.
\item Montrer que :
 $$\displaylines{
\forall\epsilon>0,\;\exists p\in\mathbb{N},\;\forall n \geq p,\;
u_n > l- \epsilon\cr
\forall\epsilon>0,\;\forall p\in\mathbb{N},\;\exists n \geq p,\;
u_n < l+ \epsilon \cr}$$
\item Interpréter ces propriétés. Énoncer des propriétés analogues pour $L$.
Démontrez-les.
\item Que peut-on dire de $(u_n)$ si $L=l$ ?
\end{enumerate}
\finenonce{000484}



\finexercice



\exercice{485, ridde, 1999/11/01}

\enonce{000485}{}
Soient $x$ et $y$ deux r\'eels strictement positifs. On pose
$$ a = \frac{x + y}2 \qquad g = \sqrt{xy} \qquad h = \frac{2xy}{x + y}
\qquad q = \sqrt{\frac12 (x^2 + y^2)}$$
Montrer que $a, g, h, q$ sont rang\'es dans un ordre ind\'ependant de $x$ et $y$.
\finenonce{000485}



\finexercice



\exercice{486, ridde, 1999/11/01}

\enonce{000486}{}
Soient $A$ et $B$ deux parties non vides born\'ees de $\Rr$.
\begin{enumerate}
\item Montrer que $A \cup B$ est born\'ee et que $\sup (A \cup B) = \max (
\sup (A), \sup (B))$.
\item Enoncer un r\'esultat analogue pour $\inf (A \cup B)$.
\item Qu'en est-il pour $A \cap B$ ?
\end{enumerate}
\finenonce{000486}



\finexercice



\exercice{5210, rouget, 2010/06/30}
\enonce{005210}{**IT}
\label{exo:suprou2bis}
Soient $A$ et $B$ deux parties de $\Rr$, non vides et bornées. Montrer que $\mbox{sup }A$, $\mbox{sup }B$, $\mbox{sup}(A+B)$, $\mbox{inf }A$, $\mbox{inf B}$, $\mbox{inf }(A+B)$ existent et que l'on a $\mbox{sup }(A+B)=\mbox{sup }A+\mbox{sup }B$ et $\mbox{inf }(A+B)=\mbox{inf }A+\mbox{inf }B$. ($A+B$ désigne l'ensemble des sommes d'un élément de $A$ et d'un élément de $B$).
\finenonce{005210}


\finexercice


\exercice{5211, rouget, 2010/06/30}
\enonce{005211}{**}
Soit $A=\left\{\frac{1}{n}+(-1)^n,\;n\in\Nn^*\right\}$. Déterminer $\mbox{sup }A$ et $\mbox{inf }A$.
\finenonce{005211}


\finexercice


\exercice{5212, rouget, 2010/06/30}
\enonce{005212}{**IT}
Soit $A$ une partie non vide et bornée de $\Rr$. Montrer que $\mbox{sup}\{|x-y|,\;(x,y)\in A^2\}=\mbox{sup }A-\mbox{inf }A$.
\finenonce{005212}


\finexercice


\exercice{5213, rouget, 2010/06/30}
\enonce{005213}{***IT}
Soient $A$ et $B$ deux parties non vides et majorées de $\Rr$. Que dire de $\mbox{sup}(A\cap B)$, $\mbox{sup}(A\cup B)$, $\mbox{sup}(A+B)$ et $\mbox{sup}(AB)$~?~($A+B$ (resp. $AB$) désigne l'ensemble des sommes (resp. des produits) d'un élément de $A$ et d'un élément de $B$).
\finenonce{005213}


\finexercice
\section{120.03 Propriétés des nombres réels}

\section{120.04 Intervalle, densité}

\section{120.99 Autre}


\exercice{487, bodin, 1998/09/01}

\enonce{000487}{}
D\'emontrer par r\'ecurrence sur $n$ que pour tout $n\geq 2$
l'implication $$ [x>-1,x\not =0] \Rightarrow [ (1+x)^n >1+nx]$$
est vraie.
\finenonce{000487}



\finexercice



\exercice{488, bodin, 1998/09/01}

\enonce{000488}{}
Soient $a_1,\ldots,a_n,\ b_1,\ldots,b_n \in \Rr$, les $a_i$
n'\'etant pas tous nuls. Soit $p(x) = \sum_{i=1}^n(a_i+xb_i)^2$.
Montrer que le discriminant de cette \'equation du second degr\'e
est $\le 0$. En d\'eduire que :
$$ \left| \sum_{i=1}^n a_ib_i \right| \le \left( \sum_{i=1}^n a_i^2 \right)^{1/2}\left( \sum_{i=1}^n b_i^2 \right)^{1/2},$$
et que
$$ \left( \sum_{i=1}^n (a_i+b_i )^2 \right)^{1/2} \le \left( \sum_{i=1}^n a_i^2 \right)^{1/2} + \left( \sum_{i=1}^n b_i^2 \right)^{1/2}.$$
\finenonce{000488}



\finexercice



\exercice{489, bodin, 1998/09/01}

\enonce{000489}{}
Deux entiers naturels distincts peuvent-ils v\'erifier la
relation $a^b=b^a$ ?
\finenonce{000489}



\finexercice



\exercice{490, bodin, 1998/09/01}

\enonce{000490}{}
 R\'esoudre l'\'equation $\sqrt[4]{41+x}+\sqrt[4]{41-x}=4$,
$x$ \'etant un r\'eel positif.
\finenonce{000490}



\finexercice



\exercice{491, bodin, 1998/09/01}
\enonce{000491}{}
 Si $a$ et $b$ sont  des r\'eels positifs ou nuls, montrer
que :
$$
\sqrt{a}+\sqrt{b} \leqslant 2\sqrt{a+b}.
$$
\finenonce{000491} 



\finexercice

\exercice{492, bodin, 1998/09/01}

\enonce{000492}{}
Soient $x=(x_1,\ldots,x_n)$ et $y=(y_1,\ldots,y_n)\in\Rr^n$.
On note $\|x\|_1=\sum_{i=1}^n|x_i|$ et $\|x\|_\infty=\text{max}_{1\le i \le n}|x_i|$.

Montrer que dans les deux cas on a :
$$ \|x+y\| \le \|x\| + \|y\|.$$
\finenonce{000492}



\finexercice



\exercice{493, bodin, 1998/09/01}

\enonce{000493}{}
Pout tout $x\in \Rr$ on note $E(x)$ sa partie enti\`ere et
$\{x\}$ sa partie d\'ecimale.
\begin{enumerate}
    \item Tracer les graphes des fonctions
$x \mapsto E(x)$ et $x\mapsto \{x\}$.
    \item Montrer les relations suivantes :
$E(x)+E(y) \le E(x+y)$, $E(x+n)=E(x)+n$ pour tout $n\in \Zz$,
$E \left( \frac{E(nx)}{n} \right) = E(x)$ pour tout $n\in\Nn^*$.
    \item D\'eterminer $\lim E(x)$ et $\lim \{x\}$ lorsque
$x \rightarrow -1_+$ et $x \rightarrow -1_-$. Ces fonctions ont-elles
une limites lorsque $x \rightarrow -1$ ?
\end{enumerate}
\finenonce{000493}



\finexercice



\exercice{494, bodin, 1998/09/01}

\enonce{000494}{}
Pour tout $x,y\in \Rr$ et $\lambda > 0$ montrer
que :
$$
2xy \le \frac{x^2}{\lambda}+\lambda y^2.
$$
\finenonce{000494}



\finexercice



\exercice{495, cousquer, 2003/10/01}

\enonce{000495}{}
 Soit deux nombres réels $a$ et $b$ vérifiant :
$-1<a<4 \quad\mbox{et}\quad-3<b<-1.$
Donner un encadrement de $a-b$ et de $a/b$.
\finenonce{000495}



\finexercice



\exercice{496, ridde, 1999/11/01}

\enonce{000496}{}
On note $E (x)$ la partie enti\`ere d'un r\'eel $x$.
\begin{enumerate}
\item Montrer que $\forall (x, y) \in \Rr^2 \; \; E (x) + E (y) \leq E (x + y)
\leq E (x) + E (y) + 1$.
\item Calculer $E (x) + E (-x)$ pour $x \in \Rr$.
\item Montrer que $\forall n \in \Nn^*$ et $\forall x \in \Rr$ $E (x) =E(\dfrac{E (nx)}
{n})$.
\end{enumerate}
\finenonce{000496}



\finexercice



\exercice{497, ridde, 1999/11/01}
\video{TLQnc9s8vkc}
\enonce{000497}{}
 Soit $f : \Rr \rightarrow \Rr$  telle que
$$\forall (x, y)\in \Rr^2 \quad  f(x + y) = f(x) + f(y).$$ 
Montrer que
\begin{enumerate}
\item $\forall n\in \Nn \qquad f(n) = n \cdot f(1)$.
\item $\forall n\in \Zz \qquad f(n) = n \cdot f(1)$.
\item $\forall q\in \Qq \qquad f(q) = q \cdot f(1)$.
\item $\forall x\in \Rr \qquad f(x) = x \cdot f(1)$ si $f$ est croissante.
\end{enumerate}
\finenonce{000497} 


\finexercice


\exercice{498, gourio, 2001/09/01}

\enonce{000498}{}
Soient $n\in \Nn^{*},$ et $(x_{1},x_{2},...,x_{n})\in \Rr^{n} $ tels que $%
\sum\limits_{i=1}^{n}x_{i}=\sum\limits_{i=1}^{n}x_{i}^{2}=n.$
Montrer que
$$\forall i\in \{1,...,n\},x_{i}=1.$$
\finenonce{000498}



\finexercice



\exercice{499, gourio, 2001/09/01}

\enonce{000499}{}
Soient $n\in \Nn^{*},$ et $(x_{1},x_{2},...,x_{n})\in [0,1]^{n}$, montrer que :
$$\prod\limits_{i=1}^{n}(1-x_{i})\geq 1-\sum\limits_{i=1}^{n}x_{i}.$$
\finenonce{000499}



\finexercice



\exercice{500, gourio, 2001/09/01}

\enonce{000500}{}
Soit $A$ une partie de $\Rr$ v\'{e}rifiant :
$$A\neq \emptyset ,$$
$$\forall x\in A,\exists \epsilon _{x}>0,]x-\epsilon _{x},x+\epsilon
_{x}[\subset A,$$
$$\forall x\in \Rr:(\forall \epsilon >0,]x-\epsilon ,x+\epsilon [\cap
A\neq \emptyset )\Rightarrow x\in A.$$
Montrer que $A=\Rr.$
\finenonce{000500}



\finexercice



\exercice{501, gourio, 2001/09/01}

\enonce{000501}{}
Montrer :
$$\forall n\geq 1,\forall x\in \Rr,\sum\limits_{k=0}^{n-1}E(x+\frac{k}{n})=E(nx).$$
\finenonce{000501}



\finexercice



\exercice{502, gourio, 2001/09/01}

\enonce{000502}{}
Soient $A$ et $B$ deux parties denses de $\Rr$, $AB$ et $A+B$ sont-elles
denses ?
\'Etude de la r\'{e}ciproque.
\finenonce{000502}



\finexercice



\exercice{503, gourio, 2001/09/01}

\enonce{000503}{}
D\'{e}montrer que :
$$\forall n\in \Nn^{*},E(\sqrt{n}+\sqrt{n+1})=E(\sqrt{4n+2}).$$
\finenonce{000503}



\finexercice



\exercice{3061, quercia, 2010/03/08}

\enonce{003061}{Morphismes de $\R$}
Soit $f : \R \to \R$ un morphisme de corps.
\begin{enumerate}
  \item Montrer que : $\forall\ x \in \Q,\ f(x) = x$.
  \item Montrer que $f$ est une application croissante.
  \item En d{\'e}duire que $f = \mathrm{id}_\R$.
\end{enumerate}
\finenonce{003061}



\finexercice 


\exercice{3062, quercia, 2010/03/08}

\enonce{003062}{Parties denses}
Soit $A \subset \R$ v{\'e}rifiant :
$$\begin{cases}\forall\ x \in \R,\ \exists\ a,b\in A \text{ tq }a < x < b\cr
        \forall\ a,b \in A,\ \frac{a+b}2 \in A.\cr\end{cases}$$
Montrer que $A$ est dense dans $\R$.
\finenonce{003062}



\finexercice 


\exercice{3063, quercia, 2010/03/08}

\enonce{003063}{Parties denses}
Soir $A$ un sous-anneau de $\R$. Montrer que $A$ est dense dans $\R$ si et
seulement si $A \cap ]0,1[ \ne \varnothing$.
\finenonce{003063}



\finexercice 


\exercice{3064, quercia, 2010/03/08}

\enonce{003064}{Sous-groupes de $\R$}
Soit $H$ un sous-groupe additif de $\R$, $H \ne \{0\}$.
On pose $H^{+*} = H\cap\R^{+*}$, et $\alpha = \inf(H^{+*})$.
\begin{enumerate}
  \item Si $\alpha \in H^{+*}$, montrer que $H = \alpha \Z$.
  \item Si $\alpha \notin H^{+*}$, montrer que $\alpha = 0$ et en d{\'e}duire que $H$
    est dense dans $\R$.
\end{enumerate}
\finenonce{003064}



\finexercice 


\exercice{3065, quercia, 2010/03/08}

\enonce{003065}{Partie enti{\`e}re}

\begin{enumerate}
\def\f#1{\left[\frac{#1}b\right]}
  \item Soient $a\in\Z$ et $b\in\N^*$.
    Montrer que : $\f{a} + \f{a+1} + \dots + \f{a+b-1} = a$.

  \item Soient $a\in\R$ et $b\in\N^*$.
    Montrer que : $\f{a} + \f{a+1} + \dots + \f{a+b-1} = [a]$.
\end{enumerate}
\finenonce{003065}




\finexercice 


\exercice{5146, rouget, 2010/06/30}
\enonce{005146}{**I Moyennes arithmétique, géométrique et harmonique}
Soient $x$ et $y$ deux réels tels que $0<x\leq y$. On pose $m=\frac{x+y}{2}$ (moyenne arithmétique), $g=\sqrt{xy}$
(moyenne géométrique) et $\frac{1}{h}=\frac{1}{2}(\frac{1}{x}+\frac{1}{y})$ (moyenne harmonique). Montrer que $x\leq
h\leq g\leq m\leq y$.
\finenonce{005146}


\finexercice


\exercice{5151, rouget, 2010/06/30}
\enonce{005151}{***}
Soient $a$, $b$ et $c$ trois réels positifs. Montrer que l'un au moins des trois réels $a(1-b)$, $b(1-c)$, $c(1-a)$ est
inférieur ou égal à $\frac{1}{4}$.
\finenonce{005151}


\finexercice


\exercice{5152, rouget, 2010/06/30}
\enonce{005152}{**I}
\begin{enumerate}
\item  Montrer que~:~$\forall x\in\Rr,\;E(x+1)=E(x)+1$.
\item  Montrer que~:~$\forall(x,y)\in\Rr^2,\;E(x)+E(y)\leq E(x+y)$.
\item  Montrer que~:~$\forall(x,y)\in\Rr^2,\;E(x)+E(y)+E(x+y)\leq E(2x)+E(2y)$.
\end{enumerate}
\finenonce{005152}


\finexercice


\exercice{5153, rouget, 2010/06/30}
\enonce{005153}{**I}
Tout entier naturel non nul $n$ s'écrit de manière unique sous la forme

$$n=a_0+10a_1+...+10^pa_p,$$

où $p$ est un entier naturel et les $a_i$ sont des entiers éléments de $\{0,...,9\}$, $a_p$ étant non nul. Déterminer
$p$ en fonction de $n$.
\finenonce{005153}


\finexercice

\exercice{5155, rouget, 2010/06/30}
\enonce{005155}{**I}
Soient $n$ un entier naturel et $x$ un réel positif.
\begin{enumerate}
\item  Combien y a-t-il d'entiers naturels entre $1$ et $n$~?~entre $1$ et $x$~?
\item  Combien y a-t-il d'entiers naturels entre $0$ et $n$~?~entre $0$ et $x$~?
\item  Combien y a-t-il d'entiers naturels pairs entre $0$ et $x$~?~Combien y a-t-il d'entiers naturels impairs entre
$0$ et $x$~?
\item  Combien y a-t-il de multiples de $3$ entre $0$ et $x$~?
\item  Combien l'équation $x+2y=n$, $n$ entier naturel donné et $x$ et $y$ entiers naturels inconnus, a-t-elle de
couples solutions~?
\item  De combien de façons peut-on payer $10$ euros avec des pièces de $10$ et $20$ centimes d'euros~?
\item (***) Combien l'équation $2x+3y=n$, $n$ entier naturel donné et $x$ et $y$ entiers naturels inconnus, a-t-elle
de couples solutions~?
\end{enumerate}
\finenonce{005155}


\finexercice


\exercice{5156, rouget, 2010/06/30}
\enonce{005156}{****}
Montrer que~:~$\forall n\in\Nn^*,\;\forall x\in\Rr,\;\sum_{k=0}^{n-1}E(x+\frac{k}{n})=E(nx)$ (poser
la division euclidienne de $E(nx)$ par $n$).
\finenonce{005156}


\finexercice

\exercice{5159, rouget, 2010/06/30}
\enonce{005159}{**}
Montrer que $\forall n\in\Nn^*,\;\forall x\in\Rr,\;E(\frac{E(nx)}{n})=E(x)$.
\finenonce{005159}


\finexercice


\exercice{5160, rouget, 2010/06/30}
\enonce{005160}{***}
Soit $n\in\Nn^*$ et $(x_1,x_2,...,x_n )\in[-1,1]^n$ tels que $x_1+x_2+...+x_n=0$.

Montrer que $|x_1+2x_2+...+nx_n|\leq E(\frac{n^2}{4})$.
\finenonce{005160}


\finexercice


\exercice{5215, rouget, 2010/06/30}
\enonce{005215}{** Identité de \textsc{Catalan}}
Montrer que pour tout entier naturel non nul $n$, $\sum_{k=0}^{2n-1}\frac{(-1)^k}{k+1}=\sum_{k=n+1}^{2n}\frac{1}{k}$.
\finenonce{005215}


\finexercice


\exercice{5216, rouget, 2010/06/30}
\enonce{005216}{**I Inégalités de \textsc{Cauchy}-\textsc{Schwarz} et de \textsc{Minkowski}}
Soient $a_1$,..., $a_n$, $b_1$,..., $b_n$ des nombres réels.
\begin{enumerate}
\item  En considérant la fonction $f~:~x\mapsto\sum_{k=1}^{n}(a_k+xb_k)^2$, montrer que
$|\sum_{k=1}^{n}a_kb_k|\leq\sqrt{\sum_{k=1}^{n}a_k^2}\sqrt{\sum_{k=1}^{n}b_k^2}$ (inégalité de \textsc{Cauchy}-\textsc{Schwarz}).
\item  En déduire l'inégalité de \textsc{Minkowski}~:~$\sqrt{\sum_{k=1}^{n}(a_k+b_k)^2}\leq\sqrt{\sum_{k=1}^{n}a_k^2}+\sqrt{\sum_{k=1}^{n}b_k^2}$.

(l'inégalité de \textsc{Cauchy}-\textsc{Schwarz} affirme que le produit scalaire de deux vecteurs est inférieur ou égal au produit de leurs normes et l'inégalité de \textsc{Minkowski} est l'inégalité triangulaire).
\end{enumerate}
\finenonce{005216}


\finexercice


\exercice{5217, rouget, 2010/06/30}
\enonce{005217}{**}
Résoudre dans $\Rr$ l'équation $\sqrt{x+2\sqrt{x-1}}+\sqrt{x-2\sqrt{x-1}}=1$.
\finenonce{005217}


\finexercice


\exercice{5218, rouget, 2010/06/30}
\enonce{005218}{**** Sous groupes de $(\Rr,+)$}
\begin{enumerate}
\item  Montrer que les sous groupes du groupe $(\Rr,+)$ sont soit de la forme $a\Zz$, $a$ réel donné, soit denses dans $\Rr$.

Indication~:~pour $G$ sous-groupe donné de $(\Rr,+)$, non réduit à $\{0\}$, considérer $a=\mbox{Inf }(G\cap]0;+\infty[)$ puis envisager les deux cas $a=0$ et $a>0$.

(Definition~:~$G$ est dense dans $\Rr$ si et seulement si~:~$(\forall x\in\Rr,\;\forall\varepsilon>0,\;\exists y\in G/\;|y-x|<\varepsilon)$.
\item  Application 1. Montrer que $\{a+b\sqrt{2},\;(a,b)\in\Zz^2\}$ est dense dans $\Rr$.
\item  Application 2 (groupe des périodes d'une fonction).
\begin{enumerate}
\item Soit $f$ une fonction définie sur $\Rr$ à valeurs dans $\Rr$. Montrer que l'ensemble des périodes de $f$ est un sous groupe de $(\Rr,+)$ (ce sous-groupe est réduit à $\{0\}$ si $f$ n'est pas périodique).
\item Montrer qu'une fonction continue sur $\Rr$ qui admet $1$ et $\sqrt{2}$ pour périodes, est constante sur $\Rr$.
\end{enumerate}
\end{enumerate}
\finenonce{005218}


\finexercice


\exercice{5219, rouget, 2010/06/30}
\enonce{005219}{**}
Montrer que $\{r^3,\;r\in\Qq\}$ est dense dans $\Rr$.
\finenonce{005219}


\finexercice


\exercice{5982, bodin, 2010/12/06}
\video{DJNMuwA-0Ts}
\enonce{005982}{}

Soit $x$ un r\'eel.
\begin{enumerate}
  \item Donner l'encadrement qui définit la partie entière $E(x)$.
  \item Soit $(u_n)_{n\in \Nn^*}$ la suite définie par $u_n = \dfrac{E (x) + E (2x) + \ldots + E (nx)}{n^2}$. \\
  Donner un encadrement simple de $n^2 \times u_n$, qui utilise $\sum_{k=1}^n k$.
  \item En déduire que $(u_n)$ converge et calculer sa limite.
  \item En d\'eduire que $\Qq$ est dense dans $\Rr$.
\end{enumerate}
\finenonce{005982}



\finexercice

\exercice{7169, megy, 2017/07/26}
\enonce{007169}{}
%[application directe d'AM>GM]
Soient $a$ et $b$ deux réels.
Montrer
\[ ab \leq \frac{a^2}{4}+b^2.\]
Généraliser.
\finenonce{007169}

\finexercice  


\exercice{7170, megy, 2017/07/26}
\enonce{007170}{}
%[application directe de'AM>GM]
On considère deux réels positifs dont le produit vaut $100$. Leur somme a-t-elle une valeur minimale et si oui laquelle et dans quel(s) cas?
\finenonce{007170}
\finexercice  


\exercice{7171, megy, 2017/07/26}
\enonce{007171}{}
%[application directe de'AM>GM]
Soit $n>0$ un entier. On considère $n$ réels positifs dont le produit vaut  $1$. Leur somme a-t-elle une valeur minimale et si oui laquelle et dans quel(s) cas?
\finenonce{007171}
\finexercice  


\exercice{7172, megy, 2017/07/26}
\enonce{007172}{}
%[application directe d'AM>GM]
Un magasin vend au même prix deux lots de trois cristaux. Le premier lot comporte trois cristaux cubiques de côté $a$, $b$ et $c$ respectivement. Le second lot comporte trois cristaux identiques en forme de parallélépipède de dimensions $a\times b \times c$. Quel lot est-il préférable d'acheter ?
\finenonce{007172}

\finexercice  


\exercice{7173, megy, 2017/07/26}
\enonce{007173}{}
%[Application directe]
Soient $a$, $b$ et $c$ des réels positifs. Montrer que 
\[ 
\frac{a^2}{bc}+\frac{b^2}{ca}+\frac{c^2}{ab} \geq 3.
\]
\finenonce{007173}

\finexercice  


\exercice{7174, megy, 2017/07/26}
\enonce{007174}{}
%[application directe]
Soient $a_1$, ... $a_n$ des réels strictement positifs et $b_1$, ... $b_n$ les mêmes $n$ réels mais numérotés dans un ordre différent. Montrer que
\[
\frac{a_1}{b_1} + ... + \frac{a_n}{b_n} \geq n.
\]
\finenonce{007174}

\finexercice  


\exercice{7175, megy, 2017/07/26}
\enonce{007175}{}
%[séparer les termes]
Soient $a, b \in \R^*$. Montrer que
\[
(1+a^2)(1+b^2) \geq 4ab.
\]
\finenonce{007175}

\finexercice  


\exercice{7176, megy, 2017/07/26}
\enonce{007176}{}
%[séparer les termes]
Soient $a, b, c \in \R_+^*$. Montrer que
\[
\left(a^2b+b^2c+c^2a\right)\left(ab^2+bc^2+ca^2\right)
\geq
9 a^2b^2c^2
\]
\finenonce{007176}

\finexercice  


\exercice{7177, megy, 2017/07/26}
\enonce{007177}{}
%[application directe d'AM>GM après avoir développé et séparé les termes]
% question ouverte
Soient $a$ et $b$ deux réels positifs tels que $a+b=8$. Déterminer la valeur minimale de
\[ \left(1+\frac1a\right)\left(1+\frac1b\right)\]
et préciser pour quelles valeurs elle est atteinte.
\finenonce{007177}

\finexercice  


\exercice{7178, megy, 2017/07/26}
\enonce{007178}{Cauchy-Schwarz}
% application directe d'AM>GM après avoir développé
Soient $a$, $b$, $c$ et $d$ quatre réels. Montrer que
\[ (a^2+b^2)(c^2+d^2) \geq (ac+bd)^2.\]
\finenonce{007178}

\finexercice  


\exercice{7179, megy, 2017/07/26}
\enonce{007179}{}
Résoudre sur $\R$ l'équation
\[ 2^x + x^2 = 2-\frac{1}{2^x}.\]
\finenonce{007179}

\finexercice  


\exercice{7180, megy, 2017/07/26}
\enonce{007180}{}
Résoudre le système
\[
\left\{
\begin{matrix}
4x &+& \frac{18}{y} &=&14\\
2y &+& \frac{9}{z} &=&15\\
9z &+& \frac{16}{x} &=&17\\
\end{matrix}
\right.
\]
\finenonce{007180}

\finexercice  


\exercice{7181, megy, 2017/07/26}
\enonce{007181}{}
Soient $a$ et $b$ deux réels positifs. Montrer que
\[ 2a^3+b^3 \geq 3a^2b.\]
\finenonce{007181}

\finexercice  


\exercice{7182, megy, 2017/07/26}
\enonce{007182}{}
%[séparation de termes]
Soient $a$, $b$ et $c$ trois réels. Montrer que
\[ a^3+b^3+c^3 \geq a^2b+b^2c+c^2a.\]
\finenonce{007182}

\finexercice  


\exercice{7183, megy, 2017/07/26}
\enonce{007183}{}
%[séparation de termes]
Soient $a$, $b$ et $c$ trois réels. Montrer que
\[ a^4+b^4+c^4 \geq abc(a+b+c).\]
\finenonce{007183}

\finexercice  


\exercice{7184, megy, 2017/07/26}
\enonce{007184}{}
%[séparation et groupement de termes]
Soient $a$, $b$ et $c$ trois réels non nuls. Montrer que
\[ \frac{a^2}{b^2}+\frac{b^2}{c^2}+\frac{c^2}{a^2}
\geq
\frac{b}{a}+\frac{c}{b}+\frac{a}{c}.
\]
\finenonce{007184}

\finexercice  


\exercice{7185, megy, 2017/07/26}
\enonce{007185}{}
%[modifier pour se ramener à un cas simple]
Soient $a$, $b$ et $c$ des réels strictement positifs. Montrer que
\[ \frac{c}{a} + \frac{a}{b+c} +\frac{b}{c} \geq 2.
\]
\finenonce{007185}

\finexercice  

\section{121.01 Convergence}


\exercice{504, cousquer, 2003/10/01}

\enonce{000504}{}
\begin{enumerate}
\item
Dessiner les suites suivantes~:
\begin{enumerate}
  \item $\displaystyle u_n = \frac{n^2-25}{2n^2+1}$\qquad
(prendre 2~cm comme unité sur $Oy$)
  \item $u_n= (-1)^n$
  \item $\displaystyle u_n =\frac{1}{n}\cos n \qquad
v_n=\frac{1}{n} \vert\cos n\vert$ \qquad($n$ en radians)
  \item $u_n = \cos n$
  \item $u_1=1$~; $u_2=2$~; $u_3=3$~; $u_4=-1$~; $u_n=2$ pour $n\ge 5$.
  \item $\displaystyle u_n=\frac{(-1)^n}{n^2+1}$ \qquad
(prendre 10 cm comme unité sur $Oy$)
  \item $\displaystyle u_n= \cos\frac{n\pi}{6}$
  \item $\displaystyle u_n = \sin\frac{1}{\sqrt{n}}$ \qquad
(prendre 1 cm comme unité sur $Oy$)
  \item $u_n = n^2+1$
  \item $\displaystyle u_n=\frac{1}{n+(-1)^n\sqrt{n}}$ \quad(pour $n\ge 2$)
\end{enumerate}


\item
Classer les dessins par paquets en précisant vos critères.

\item
Pour chaque suite, pouvez-vous trouver
$l$ et $n$ tels que $\vert u_n-l\vert<\frac{1}{10}$ ou $\frac{1}{100}$~?
Mettre en relation avec le classement précédent.

\item

Les énoncés suivants sont-ils vrais ou faux ?

\begin{enumerate}
\item Une suite à termes positifs qui tend vers 0 est décroissante
à partir d'un certain rang.

\item Si une suite a une limite strictement positive, tous ses termes sont
strictement positifs à partir d'un certain rang. Réciproque~?

\end{enumerate}
\end{enumerate}
\finenonce{000504}



\finexercice



\exercice{505, bodin, 1998/09/01}
\video{FRpMiQ8DOwI}
\enonce{000505}{}
 Soit $(u_n)_{n\in\N}$ une suite de $\R$. Que pensez-vous des
propositions suivantes :
\par\noindent $\bullet$ Si $(u_{n})_n$ converge vers un r\'eel $\ell$ alors $(u_{2n})_n$ et $(u_{2n+1})_n$
convergent vers $\ell$.
\par\noindent  $\bullet$ Si $(u_{2n})_n$ et $(u_{2n+1})_n$ sont convergentes, il en est
de m\^{e}me de $(u_{n})_n$.
\par\noindent $\bullet$ Si $(u_{2n})_n$ et $(u_{2n+1})_n$ sont convergentes, de m\^{e}me
limite $\ell$, il en est de m\^{e}me de $(u_{n})_n$.
\finenonce{000505} 


\finexercice


\exercice{512, cousquer, 2003/10/01}

\enonce{000512}{}
 Vrai ou faux : il existe une suite $(u_n)$ telle que $(u_{n+1}-u_n)$
tend vers~$0$ et qui diverge.
\finenonce{000512}



\finexercice



\exercice{513, cousquer, 2003/10/01}

\enonce{000513}{}
 Encadrer la suite $(u_n)$ définie par 
$u_n=\sum_{k=1}^n \frac{1}{n^2+k^2}$. Que peut-on en 
déduire ?
\finenonce{000513}



\finexercice



\exercice{514, cousquer, 2003/10/01}

\enonce{000514}{}
\begin{enumerate}
\item Que peut-on dire d'une suite qui vérifie $\lim_{n\to\infty}nu_n=0$~?
\item Que peut-on dire d'une suite qui vérifie $\lim_{n\to\infty}nu_n=1$~?
\item Que peut-on dire d'une suite qui vérifie $\lim_{n\to\infty}nu_n=+\infty$~?
\end{enumerate}
\finenonce{000514}



\finexercice



\exercice{515, cousquer, 2003/10/01}

\enonce{000515}{}
 Étant donné $k\in\mathbb{R}_+$, que peut-on dire d'une suite 
$(u_n)$ qui vérifie
$\lim_{n\to\infty} \frac{u_{n+1}}{u_n} = k$ ?
\textit{Application} : Étudier $u_n=\frac{1\cdot 2 \cdots n}{1\cdot 4\cdots(3n-2)}$.
\finenonce{000515}



\finexercice



\exercice{516, ridde, 1999/11/01}

\enonce{000516}{}
Montrer qu'une partie $D$ est dense dans $\Rr$ ssi tout r\'eel est limite d'une suite
de points de $D$.
\finenonce{000516}



\finexercice



\exercice{517, ridde, 1999/11/01}

\enonce{000517}{}
Soit $A$ une partie born\'ee de $\Rr$ et $x$ un r\'eel.
\begin{enumerate}
\item Montrer que $x = \sup (A)$ ssi ($x$ majore $A$ et il existe une suite
$ (x_n)_{n \in \Nn}$ de A qui converge vers $x$).
\item \'Enoncer un r\'esultat analogue pour $\inf (A)$.
\end{enumerate}
\finenonce{000517}



\finexercice



\exercice{518, ridde, 1999/11/01}

\enonce{000518}{}
\'Etudier la convergence des suites :\\

$\sqrt{n^2 + n + 1}-\sqrt n \qquad \dfrac{n\sin (n)}{n^2 + 1} \qquad \dfrac1n  + (-1)^n
\qquad n \sum\limits_{k = 1}^{2n + 1}\dfrac 1{n^2 + k} \qquad
\dfrac 1n \sum\limits_{k = 0}^{n-1}\cos (\dfrac1{\sqrt{n + k}})$

\finenonce{000518}




\finexercice



\exercice{519, ridde, 1999/11/01}
\video{d6n_rWtLv1Y}
\enonce{000519}{}
 Montrer qu'une suite d'entiers qui converge est
constante \`a partir d'un certain rang.
\finenonce{000519} 


\finexercice

\exercice{520, ridde, 1999/11/01}
\video{HZBJUfv7fSA}
\enonce{000520}{}
 Soit $H_n = 1 + \dfrac12 + \cdots + \dfrac1n$.
\begin{enumerate}
\item En utilisant une int\'egrale, montrer que pour tout $n>0$ : $\dfrac1{n + 1} \leq
\ln (n + 1)-\ln (n) \leq \dfrac1n$.
\item En d\'eduire que $\ln (n + 1) \leq H_n \leq \ln (n) + 1$.
\item D\'eterminer la limite de $H_n$.
\item Montrer que $u_n = H_n-\ln (n)$ est d\'ecroissante et positive.
\item Conclusion ?
\end{enumerate}
\finenonce{000520} 


\finexercice

\exercice{521, ridde, 1999/11/01}

\enonce{000521}{}
Montrer qu'une suite monotone dont une suite extraite converge est convergente.
\finenonce{000521}



\finexercice



\exercice{522, ridde, 1999/11/01}

\enonce{000522}{}
Montrer que $ (u_{n})$ converge ssi $ (u_{2n}), (u_{2n + 1}), (u_{3n})$ convergent
 (leurs limites n'\'etant pas n\'ecessairement \'egales).
\finenonce{000522}



\finexercice



\exercice{523, monthub, 2001/11/01}

\enonce{000523}{}
Etudier la convergence de la suite $u_n=(-1)^n\dfrac{n+1}{n}$.
\finenonce{000523}



\finexercice



\exercice{524, monthub, 2001/11/01}
\video{WW2PHzsJGkw}
\enonce{000524}{}
Soit $q$ un entier au moins {\'e}gal {\`a} $2$. Pour tout $n\in
\N$, on pose $u_n=\cos{\dfrac{2n\pi}{q}}$.
\begin{enumerate}
\item Montrer que $u_{n+q}=u_n$ pour tout $n \in \N$.
\item Calculer $u_{nq}$ et $u_{nq+1}$. En d{\'e}duire que la suite $(u_n)$
  n'a pas de limite.
\end{enumerate}
\finenonce{000524} 


\finexercice


\exercice{525, gourio, 2001/09/01}

\enonce{000525}{}
Soit $(u_{n})_{n\in \Nn}$ une suite r\'{e}elle prenant toute les valeurs
rationnelles. Montrer que $(u_{n})_{n\in \Nn}$ n'admet pas de limite.
\finenonce{000525}



\finexercice



\exercice{526, gourio, 2001/09/01}

\enonce{000526}{}
Soit $(u_{n})_{n\in \Nn}$ une suite r\'{e}elle telle que $\lim\limits_{n
\rightarrow \infty }u_{n}^{2}=\lambda$. Que dire de $(u_{n})_{n\in \Nn}$ ?
\finenonce{000526}



\finexercice



\exercice{527, gourio, 2001/09/01}

\enonce{000527}{}
\begin{enumerate}
\item Donner un exemple de suite born\'{e}e divergente, puis de suite divergente
telle que
$$\forall k\in \Nn,\lim_{n\rightarrow \infty }x_{n+k}-x_{n}=0.$$

\item Donner un exemple de suite divergente qui a une seule valeur d'adh\'{e}rence
(i.e. telle qu'il existe une seule extraction $\phi  $ telle que
$x_{\phi(n)\text{ }}$ converge).

\item Donner un exemple de suite $(x_{n})_{n\in \Nn} $ divergente telle que $\forall
k\geq 2,(x_{nk})_{n\in \Nn}$ converge.
\end{enumerate}
\finenonce{000527}



\finexercice



\exercice{528, cousquer, 2003/10/01}

\enonce{000528}{}
 Que peut-on dire des nombres réels $a$ et $b$ si
$$\forall n \in\mathbb{N}^\ast,\; a-\frac{1}{n}\leq b \leq 
a+\frac{1}{n}~?$$
\finenonce{000528}



\finexercice



\exercice{529, cousquer, 2003/10/01}

\enonce{000529}{}
 Étudier la suite $(u_n)$ définie par~:
$$u_n=
\begin{cases}
0 & \text{ si } $n$ \text{ est premier }\cr
67+1/n &  \text{ sinon }.\cr
\end{cases}$$
Si cette suite converge, montrer que sa limite est inférieure
à~$72$. Étudier la convergence de cette suite.
\finenonce{000529}



\finexercice



\exercice{530, cousquer, 2003/10/01}

\enonce{000530}{}
On donne la suite $(u_n)$ définie par~:
$$ u_1= \sqrt2 \quad\mbox{et}\quad u_n = \sqrt{2-u_{n-1}}.$$
En étudiant les suites $(u_{2n})$ et $(u_{2n+1})$,
montrer que la suite $(u_n)$ est convergente.
\finenonce{000530}



\finexercice



\exercice{531, cousquer, 2003/10/01}

\enonce{000531}{}
\begin{enumerate}
\item Soit $(u_n)$, $(v_n)$, $(w_n)$ trois suites telles que pour $n$ assez
grand on ait $v_n\leq u_n\leq w_n$. On suppose que $(v_n)$ et $(w_n)$
sont convergentes, et on note $v=\lim v_n$ et $w=\lim w_n$.
Montrer que pour tout $\epsilon$ positif, on a
$v-\epsilon\leq u_n \leq w+\epsilon$ pour $n$~assez grand
(\emph{théorème d'encadrement}).
Que peut-on en déduire si $v=w$~?

\item Soit $(u_n)$ une suite convergente de limite~$l$.
Montrer que la suite
$$v_n = \frac{u_1+u_2+ \cdots+u_n}{n}$$
est convergente et a pour limite~$l$.
Pour cela, encadrer $u_n$ à $\epsilon$~près pour $n$~assez grand,
et en déduire un encadrement de~$v_n$.
\end{enumerate}
\finenonce{000531}



\finexercice



\exercice{532, cousquer, 2003/10/01}

\enonce{000532}{}
 Soit $\alpha$ un nombre irrationnel positif et $(p_n)$ et $(q_n)$
deux suites d'éléments de $\mathbb{N}^\ast$ telles que 
$\alpha= \lim_{n\to\infty}\frac{p_n}{q_n}$. Montrer que 
$$\lim_{n\to\infty} q_n =\lim_{n\to\infty} p_n=+\infty.$$
\finenonce{000532}



\finexercice



\exercice{533, cousquer, 2003/10/01}

\enonce{000533}{}
 Étudier la suite $u_n=\ln(1+\ln(2+\ln(3+\cdots+\ln(n-1+\ln n)\cdots)))$.
\finenonce{000533}



\finexercice



\exercice{534, cousquer, 2003/10/01}

\enonce{000534}{}
 Montrer que pour $n\geq1$, l'équation $x^n+x^{n-1}+x^2+x-{n+1 \over n}=0$
admet une unique racine positive~; on la note $u_n$.
Étudier la suite $(u_n)$.
\finenonce{000534}



\finexercice



\exercice{535, cousquer, 2003/10/01}

\enonce{000535}{}
Un ivrogne part à un instant donné d'un point donné. À chaque seconde,
il fait un pas dans une direction inconnue (et qui peut changer de façon
arbitraire à chaque pas). Comme il se fatigue, ses pas sont de plus en plus
courts. Peut-on prévoir qu'au bout d'un certain temps il restera à moins
d'un mètre d'une certaine position si on admet que la longueur de son
$n$-ième pas est~:

\begin{enumerate}
\item $1/n$ mètre~?

\item $1/n^2$ mètre~?
\end{enumerate}
\finenonce{000535}



\finexercice



\exercice{1193, legall, 1998/09/01}

\enonce{001193}{}
Soient $  (u_n)_{n\geq 2}  $
d\'efinie par $  \displaystyle{ u_n =\prod_{k=2}^n
\hbox{cos} (\frac{\pi}{2^k})}  $ et $  \displaystyle{v_n =u_n \hbox{sin}(\frac{\pi} { 2^n})}  .$
\begin{enumerate}
    \item Montrer que $  (u_n)_{n\geq 2}  $ est convergente.
    \item Montrer que $  (v_n)_{n\geq 2}  $ est une suite g\'eom\'etrique. En d\'eduire
la limite de $  (u_n)_{n\geq 2}  .$
\end{enumerate}
\finenonce{001193}



\finexercice



\exercice{1194, legall, 1998/09/01}

\enonce{001194}{}
Soit $  (u_n)_{n\in { \Nn}}  $ une suite born\' ee de nombres r\' eels telle que
$  \displaystyle{\lim _{n\rightarrow \infty }(u_{n+1}-u_{n})=0}  .$ Montrer que les valeurs
d'adh\`erence
de la suite $  (u_n)_{n\in { \Nn}}  $ forment un intervalle de $  { \Rr}  .$
\finenonce{001194}



\finexercice



\exercice{1195, legall, 1998/09/01}

\enonce{001195}{}
On d\' efinit par r\' ecurrence les suites $  (u_n)_{n\in { \Nn}}  $
et $  (v_n)_{n\in { \Nn}}  $ par~:
$$u_0=1  ,   v_0=2  ,   u_{n+1}=\frac{ (u_n)^2}{ u_n+v_n}  ,   v_{n+1}=\frac{ (v_n)^2}{ u_n+v_n}
.$$
\begin{enumerate}
    \item Montrer par r\' ecurrence que l'on a $  u_n>0  $ et $  v_n>0   .$
    \item Montrer que les suites $  (u_n)_{n\in { \Nn}}  $
et $  (v_n)_{n\in { \Nn}}  $ d\' ecroissent.
En d\' eduire qu'elles convergent vers $  \ell   $ et $  \ell '  $ respectivement.
Montrer que l'on a $  \ell \ell '=0  .$
\item Montrer que la suite $  (v_n-u_n)_{n\in { \Nn}}  $ est constante.
En d\' eduire $  \ell   $ et $  \ell '  .$
\end{enumerate}
\finenonce{001195}



\finexercice



\exercice{1196, legall, 1998/09/01}

\enonce{001196}{}
Soient $  (u_n)_{n\in { \Nn}}  $ et $  (v_n)_{n\in { \Nn}}  $
deux suites de nombres r\' eels telles que $  0<u_1 <v_1  $ et $  u_{n+1}= \sqrt{ u_nv_n}  $
et $  \displaystyle{ v_{n+1}=\frac{u_n+v_n} { 2}}  .$ Montrer qu'elles convergent vers la m\^eme limite.
\finenonce{001196}



\finexercice



\exercice{1197, legall, 1998/09/01}

\enonce{001197}{}
\begin{enumerate}
    \item Soit $  (u_n)_{n\in { \Nn}}  $ une suite de nombres r\'eels non nuls
convergeant vers une limite $  \ell   $ diff\'erente de z\'ero. Montrer que la suite $
\displaystyle{(\frac{1}{ u_n})_{n\in { \Nn}}}  $ converge vers $  \displaystyle{\frac{ 1}{\ell} }  .$
    \item Soit $  (u_n)_{n\in { \Nn}}  $ une suite de nombres r\'eels positifs
convergeant vers une limite $  \ell   $ diff\'erente de z\'ero. Montrer que la suite $
(\sqrt{u_n})_{n\in { \Nn}}  $ converge vers $  \sqrt{\ell }  .$
\end{enumerate}
\finenonce{001197}



\finexercice



\exercice{1198, legall, 1998/09/01}

\enonce{001198}{}
\begin{enumerate}
    \item Soit $  (u_n)_{n\in { \Nn}}  $ une suite de nombres r\'eels
telle que les suites extraites $  (u_{2n})_{n\in { \Nn}}  $ et $  (u_{2n+1})_{n\in { \Nn}}  $
convergent vers une m\^eme limite $  \ell   .$ Montrer que $  (u_n)_{n\in { \Nn}}  $ converge
\'egale\-ment vers $  \ell   .$
    \item En d\'eduire que la suite $  (u_n)_{n\in { \Nn}}  $ de terme g\'en\'eral $  \displaystyle{
u_n =\sum _{k=0}^n \frac{ (-1)^k }{ (2k)!}}  $ converge.
\end{enumerate}
\finenonce{001198}



\finexercice



\exercice{1199, legall, 1998/09/01}

\enonce{001199}{}
Soit $  (u_n)_{n\in { \Nn}}  $ une suite de nombres r\'eels et $
\displaystyle{v_n=\frac{u_1+u_2+\cdots+u_n}{ n}}  $ o\`u $  n\in { \Nn}^*  .$
\begin{enumerate}
    \item Montrer que si $  (u_n)_{n\geq 1}  $ converge vers $  \ell   ,$ alors
$  (v_n)_{n\geq 1}  $ converge vers $  \ell   .$ La r\'eciproque est elle vraie~?
    \item Calculer $  \displaystyle{ \lim_{n \rightarrow +\infty } \sum _{k=1}^n \frac{k+1 }{ 2nk+k}}  .$
    \item Soit $  (a_n)_{n\geq 0}  $ une suite telle que
$  \displaystyle{ \lim_{n\rightarrow +\infty}(a_{n+1}-a_n)=\ell }  .$ Prouver que $
\displaystyle{\lim_{n\rightarrow +\infty}\frac{a_n}{ n}=\ell}  .$
    \item Soit $  (u_n)_{n\geq 1}  $ une suite strictement positive telle
que $  \displaystyle{\lim_{n\rightarrow
+\infty}\frac{u_{n+1}}{ u_n}=\ell}  .$ D\'emontrer que
$  \displaystyle{\lim_{n\rightarrow +\infty}(u_n)^{1/n} =\ell }  .$
\end{enumerate}
\finenonce{001199}



\finexercice



\exercice{1200, legall, 1998/09/01}

\enonce{001200}{}
Pour tout $  n \in { \Nn}^*  $ on note $  \displaystyle{ u_n =\sum _{k=1}^n
\frac{1}{ k!}}  $ et $  \displaystyle{ v_n=u_n +\frac{1}{ n!n}}  .$ On rapelle que $  \displaystyle{
e=\lim _{n \rightarrow \infty }u_n }  .$
\begin{enumerate}
    \item Montrer que les suites $  (u_n)_{n\geq 1}  $ et $  (v_n)_{n\geq 1}  $
sont adjacentes. En d\'eduire une valeur approch\'ee de $  e   $ \`a $  \displaystyle{\frac{1}{
1000}}  .$

    \item D\'emontrer que $  e  $ est irrationnel.
\end{enumerate}
\finenonce{001200}



\finexercice



\exercice{1201, legall, 2003/10/01}

\enonce{001201}{}
Une suite $(u_n)_{n\in \Nn }$ est dite de Cauchy 
lorsque, pour tout $\epsilon >0$ il existe $N\in \Nn $ tel que, si
$m,n \geq N$ alors $\vert u_n-u_m \vert < \epsilon .$
\begin{enumerate}
\item Montrer que toute suite convergente est de Cauchy. Montrer que 
toute suite de Cauchy est born\'ee.
\item Soit $u_n = 1 + \dfrac12 + \ldots + \dfrac1n$.
  Montrer que, pour tout $p\in \Nn ,$ $u_{2^p }\geq \dfrac {p+2}2 .$ 
En d\'eduire que $(u_n)_{n\in \Nn}$ tend vers l'infini.
\item Une suite $(u_n)_{n\in \Nn }$ satisfait au crit\`ere $\cal{C 
'}$ lorsque, pour tout $\epsilon >0$ il existe $N\in \Nn $ tel 
que, si
$n \geq N$ alors $\vert u_n-u_{n+1} \vert < \epsilon .$ Une suite 
satisfaisant au crit\`ere $\cal{C '}$ est-elle de Cauchy~?

\item  Montrer que les trois assertions qui suivent sont \'equivalentes~:
\begin{enumerate}
\item Toute partie major\'ee de $\Rr $ admet une borne sup\'erieure 
et toute partie minor\'ee de $\Rr $ admet une borne inf\'erieure.
\item Toute suite de Cauchy est convergente.
\item Deux suites adjacentes sont convergentes.
\end{enumerate}

\end{enumerate}
\finenonce{001201}



\finexercice



\exercice{4668, quercia, 2010/03/16}
\enonce{004668}{Limite de la partie enti{\`e}re d'une suite}

Soit $(u_n)$ une suite r{\'e}elle convergeant vers $\ell \in \R$.
La suite $([u_n])$ est-elle convergente ?


\finenonce{004668}



\finexercice


\exercice{4669, quercia, 2010/03/16}
\enonce{004669}{Limites doubles diff{\'e}rentes}

Comparer $\lim_{n\to\infty} \left( \lim_{k\to\infty} \frac{n^k}{(n+1)^k} \right)$
et       $\lim_{k\to\infty} \left( \lim_{n\to\infty} \frac{n^k}{(n+1)^k} \right)$.
\finenonce{004669}



\finexercice


\exercice{4670, quercia, 2010/03/16}
\enonce{004670}{Suites convergeant vers 0}

\begin{enumerate}
  \item Soit $(u_n)$ une suite r{\'e}elle telle que $\frac {u_n}{1+u_n} \xrightarrow[n\to\infty]{} 0$.
    Montrer que $u_n \xrightarrow[n\to\infty]{} 0$.
  \item M{\^e}me question avec
    $\begin{cases} \frac {u_n}{1+u_n^2} \xrightarrow[n\to\infty]{} 0\cr
             (u_n) \text{ est born{\'e}e.}\cr\end{cases}$

\end{enumerate}
\finenonce{004670}



\finexercice


\exercice{4671, quercia, 2010/03/16}
\enonce{004671}{$u_nv_n \to 1$}
                  

Soient $(u_n)$ et $(v_n)$ deux suites v{\'e}rifiant :
$\begin{cases} 0 \le u_n \le 1 \cr 0 \le v_n \le 1 \cr u_nv_n \xrightarrow[n\to\infty]{} 1.\cr\end{cases}$
Que pouvez-vous dire de ces suites ?

\finenonce{004671}



\finexercice


\exercice{4672, quercia, 2010/03/16}
\enonce{004672}{S{\'e}rie altern{\'e}e}

On pose $u_n = \frac {96\times(-1)^n}{(2n-3)(2n-1)(2n+1)(2n+3)(2n+5)}$
et $v_n = \sum_{k=0}^n\,u_k$.

\begin{enumerate}
  \item {\'E}tudier les suites $(v_{2n})$ et $(v_{2n+1})$ et
    montrer que la suite $(v_n)$ est convergente.
  \item Calculer $\ell = \lim_{n\to\infty} v_n$ {\`a} $10^{-5}$ pr{\`e}s.


\end{enumerate}
\finenonce{004672}



\finexercice


\exercice{4673, quercia, 2010/03/16}
\enonce{004673}{Croissance compar{\'e}e}

Montrer que l'ensemble des entiers $n$ tels que $2^{n^2} < (4n)!$ est fini.

\finenonce{004673}



\finexercice


\exercice{4674, quercia, 2010/03/16}
\enonce{004674}{Limite de $n^{1/n}$}

D{\'e}montrer, sans utiliser la fonction $\ln$, que $\sqrt[n]n \xrightarrow[n\to\infty]{} 1$.

Chercher $\lim_{n\to\infty} \sqrt[n]{n!}$.
\finenonce{004674}



\finexercice


\exercice{4675, quercia, 2010/03/16}
\enonce{004675}{Croissance logarithmique compar{\'e}e}

Soient $(a_n)$, $(b_n)$ deux suites strictement positives telles que :
$\forall\ n \in \N,\ \frac {a_{n+1}}{a_n} \le \frac {b_{n+1}}{b_n}$.


Montrer que si $b_n \xrightarrow[n\to\infty]{} 0$, alors $a_n \xrightarrow[n\to\infty]{} 0$.

\finenonce{004675}



\finexercice


\exercice{4676, quercia, 2010/03/16}
\enonce{004676}{Somme de parties enti{\`e}res}

Soit $x\in\R$.
Chercher $\lim_{n\to\infty} \frac {[x] + [2x] + \dots + [nx]}{n^2}$.

\finenonce{004676}



\finexercice


\exercice{4677, quercia, 2010/03/16}
\enonce{004677}{Divergence de $\cos(nt)$ et $\sin(nt)$}

Soit $\theta \in \R$. Montrer que si $\theta \not \equiv 0 (\mathrm{mod}\, {\pi})$,
les suites $(\cos(n\theta))$ et $(\sin(n\theta))$ sont toutes les deux
divergentes (montrer que si l'une converge, alors l'autre aussi, puis
obtenir une contradiction).

\finenonce{004677}



\finexercice


\exercice{4678, quercia, 2010/03/16}
\enonce{004678}{Somme des $1/k^{1/2}$}

Soit $u_n = 1 + \frac1{\sqrt2} + \dots + \frac1{\sqrt n}$.

\begin{enumerate}
  \item Chercher $\lim_{n\to\infty} (u_{2n}-u_n)$, puis $\lim_{n\to\infty} u_n$.
  \item Comparer $\frac 1{2\sqrt k}$, $\sqrt{k+1} - \sqrt k$,
    et $\sqrt k - \sqrt{k-1}$.
    En d{\'e}duire que la suite $\bigl(u_n - 2\sqrt n\bigr)$ est convergente.


\end{enumerate}
\finenonce{004678}



\finexercice


\exercice{4679, quercia, 2010/03/16}
\enonce{004679}{Limite de $(1+1/n)^n$}

\begin{enumerate}
  \item On pose $u_n = \frac 1{0!} + \frac 1{1!} + \dots + \frac 1{n!}$.
  \begin{enumerate}
    \item Montrer que la suite $(u_n)$ est convergente.
    \item Calculer le nombre $e = \lim_{n\to\infty}u_n$ {\`a} $10^{-7}$ pr{\`e}s.
  \end{enumerate}
  \item  On note $v_n = \left(1+\frac1n\right)^{1/n}$.
  \begin{enumerate}
     \item D{\'e}velopper $v_n$ et montrer que $v_n \le e$.
     \item On fixe $p \in \N$ et $\varepsilon > 0$.
     Montrer que pour $n$ suffisament grand, $v_n \ge \sum_{k=0}^p \frac 1{k!}
     - \varepsilon$.
     \item Que pouvez-vous en d{\'e}duire ?
  \end{enumerate}
\end{enumerate}
\finenonce{004679}



\finexercice


\exercice{4680, quercia, 2010/03/16}
\enonce{004680}{{\'E}tude de $C_{2n}^n/4^n$}
                        
On pose $u_n = \frac {1\times 3\times 5 \times \dots \times (2n-1)}
                      {2\times 4\times 6 \times \dots \times (2n)}$.

\begin{enumerate}
  \item Exprimer $u_n$ {\`a} l'aide de factorielles.
    
  \item Montrer que la suite $(u_n)$ est convergente.
  \item Soit $v_n = (n+1)u_n^2$. Montrer que la suite $(v_n)$ converge.
    Que pouvez-vous en d{\'e}duire pour $\lim_{n\to\infty}u_n$ ?
  \item On note $\alpha = \lim_{n\to\infty}v_n$. En {\'e}tudiant la suite $(nu_n^2)$,
    montrer que $\alpha > 0$.

\end{enumerate}
\finenonce{004680}



\finexercice


\exercice{4681, quercia, 2010/03/16}
\enonce{004681}{Suite $a^n/\prod(1+a^k)$}
                        
Soit $a \in \C\setminus \mathbb{U}$. {\'E}tudier la suite de terme g{\'e}n{\'e}ral :
$u_n = \frac {a^n}{(1+a)(1+a^2)\dots(1+a^n)}$.
\finenonce{004681}



\finexercice


\exercice{4682, quercia, 2010/03/16}
\enonce{004682}{Lemme de C{\'e}saro}

Soit $(u_n)$ une suite r{\'e}elle. On pose $v_n = \frac {u_1 + \dots + u_n}n$.

\begin{enumerate}
  \item  Montrer que si $u_n \xrightarrow[n\to\infty]{} 0$, alors $v_n \xrightarrow[n\to\infty]{} 0$.
  \item  Montrer que si $u_n \xrightarrow[n\to\infty]{} \ell$, alors $v_n \xrightarrow[n\to\infty]{} \ell$.
     ($\ell \in \overline{\R}$)
  \item  Donner un exemple o{\`u} $(v_n)$ converge mais $(u_n)$ diverge.

\end{enumerate}
\finenonce{004682}



\finexercice


\exercice{4683, quercia, 2010/03/16}
\enonce{004683}{Lemme de C{\'e}saro}
\begin{enumerate}
  \item Soit $(b_n)$ une suite r{\'e}elle strictement croissante tendant vers $+\infty$,
    et $(a_n)$ une suite r{\'e}elle telle que :\par
    $\frac{a_n-a_{n-1}}{b_n-b_{n-1}} \xrightarrow[n\to\infty]{} \ell \in \R$.
    Montrer que $\frac{a_n}{b_n} \xrightarrow[n\to\infty]{} \ell$.
  \item Application : Quelle est la limite de
    $\frac{1^k + 2^k + \dots + n^k}{n^{k+1}}$ $(k \in \N)$ ?
\end{enumerate}
\finenonce{004683}



\finexercice


\exercice{4684, quercia, 2010/03/16}
\enonce{004684}{C{\'e}saro g{\'e}n{\'e}ralis{\'e}}

Soit $(u_n)$ une suite r{\'e}elle convergente, et
$S_n = \frac 1{2^n} \sum_{p=0}^n\,C_n^pu_p$.
{\'E}tudier la suite $(S_n)$.


\finenonce{004684}



\finexercice


\exercice{4685, quercia, 2010/03/16}
\enonce{004685}{Produit de Cauchy}

Soient $(a_n)$, $(b_n)$ deux suites convergeant vers $a,b$.
Montrer que $\frac {a_0b_n + a_1b_{n-1} + \dots + a_nb_0}{n+1}
\xrightarrow[n\to\infty]{} ab$.

\finenonce{004685}



\finexercice


\exercice{4686, quercia, 2010/03/16}
\enonce{004686}{$x_n-ax_{n-1} \to 0$}

Soit $(x_n)$ une suite r{\'e}elle et $\alpha \in\, ]0,1[$.
On pose $$\left\{\begin{aligned}y_0 &=& x_0\hfill \cr y_n &=& x_n - \alpha x_{n-1}\text{ pour }n\ge 1.\end{aligned}\right.$$
\par
Montrer que : $(x_n \xrightarrow[n\to\infty]{} 0) \Leftrightarrow (y_n \xrightarrow[n\to\infty]{} 0)$.

\finenonce{004686}



\finexercice


\exercice{4687, quercia, 2010/03/16}
\enonce{004687}{$x_n + x_{2n}/2 \to 1$}
                         

Soit $(x_n)$ une suite born{\'e}e telle que $x_n + \frac{x_{2n}}2 \xrightarrow[n\to\infty]{} 1$.
Montrer que $x_n\xrightarrow[n\to\infty]{}\frac23$.
\finenonce{004687}



\finexercice


\exercice{4688, quercia, 2010/03/16}
\enonce{004688}{Approximation d'un irrationnel}

Soit $x \in \R^*$ et $(r_n)$ une suite de rationnels convergeant vers $x$.
On {\'e}crit $r_n = \frac {p_n}{q_n}$ avec $p_n\in\Z$, $q_n\in\N^*$.

\begin{enumerate}
  \item Montrer que si l'une des suites $(p_n)$, $(q_n)$ est born{\'e}e, alors l'autre l'est
    aussi, et $x \in \Q$.

  \item En d{\'e}duire que si $x \in \R\setminus \Q$, alors $|p_n| \xrightarrow[n\to\infty]{} +\infty$
     et $q_n \xrightarrow[n\to\infty]{} +\infty$.
\end{enumerate}
\finenonce{004688}



\finexercice


\exercice{4689, quercia, 2010/03/16}
\enonce{004689}{Somme des chiffres de $n$}

Pour $n \in \N^*$, on note $S(n)$ la somme des chiffres de l'{\'e}criture d{\'e}cimale de $n$.

\begin{enumerate}
  \item Encadrer $S(n+1)$ en fonction de $S(n)$.
    En d{\'e}duire que la suite $\left(\frac {S(n+1)}{S(n)}\right)$ est born{\'e}e.
    

  \item Chercher $\inf\left\{\frac {S(n+1)}{S(n)} \text{ tq } n \in \N^*\right\}$, et
             $\sup\left\{\frac {S(n+1)}{S(n)} \text{ tq } n \in \N^*\right\}$.
             

  \item La suite $\left(\frac {S(n+1)}{S(n)}\right)$ est-elle convergente ?
\end{enumerate}
\finenonce{004689}



\finexercice


\exercice{4690, quercia, 2010/03/16}
\enonce{004690}{{\'E}quation $x^n + x^{n-1} + \dots+ x - 1 = 0$}

On consid{\`e}re l'{\'e}quation : $x^n + x^{n-1} + \dots + x - 1 = 0$.

\begin{enumerate}
  \item Prouver qu'il existe une unique racine positive, $a_n$.
  \item Montrer que la suite $(a_n)$ est d{\'e}croissante.
  \item Montrer que $a_n \xrightarrow[n\to\infty ]{} \frac12$ (calculer $a_n^{n+1} - 1$).

\end{enumerate}
\finenonce{004690}



\finexercice


\exercice{4691, quercia, 2010/03/16}
\enonce{004691}{Suite n'ayant qu'une valeur d'adh{\'e}rence}

Soit $(u_n)$ une suite r{\'e}elle. On appelle {\it valeur d'adh{\'e}rence\/}
toute limite d'une sous-suite convergente extraite de~$(u_n)$.

\begin{enumerate}
  \item  Quelles sont les valeurs d'adh{\'e}rence d'une suite convergente ?
  \item  Quelles sont les valeurs d'adh{\'e}rence de la suite $(\cos(n\pi/3))$ ?
  \item  Montrer que si la suite $(u_n)$ est born{\'e}e et diverge, elle a au moins deux
     valeurs d'adh{\'e}rence.

\end{enumerate}
\finenonce{004691}



\finexercice


\exercice{4692, quercia, 2010/03/16}
\enonce{004692}{Limites sup et inf}

Soit $(x_n)$ une suite born{\'e}e de r{\'e}els.
On pose : $\begin{cases} y_n = \sup\{x_p \text{ tq } p\ge n\}  \cr
                   z_n = \inf\{x_p \text{ tq } p\ge n\}. \cr \end{cases}$

\begin{enumerate}
  \item Montrer que les suites $(y_n)$ et $(z_n)$ convergent.
  \item Montrer que $(x_n)$ converge si et seulement si $(y_n)$ et $(z_n)$
    ont m{\^e}me limite.

\end{enumerate}
\finenonce{004692}



\finexercice


\exercice{4693, quercia, 2010/03/16}
\enonce{004693}{Convergence vers $0$ et monotonie}

Soit $(x_n)$ une suite de r{\'e}els strictement positifs convergeant vers 0.

\begin{enumerate}
  \item Montrer qu'il existe une infinit{\'e} d'indices $n$ tels que
    $x_n = \max(x_n,x_{n+1},x_{n+2},\dots)$.
    
  \item Montrer qu'il existe une infinit{\'e} d'indices $n$ tels que
    $x_n = \min(x_0,x_1,\dots,x_n)$.
    
\end{enumerate}
\finenonce{004693}



\finexercice


\exercice{4694, quercia, 2010/03/16}
\enonce{004694}{Convergence vers $0$ et monotonie}
Soit $(u_n)$ une suite de r{\'e}els strictement positifs convergeant vers 0.
Montrer qu'il existe une bijection $\sigma : {\N} \to {\N}$ telle que
la suite $(u_{\sigma(n)})$ converge vers $0$ en d{\'e}croissant.

\finenonce{004694}



\finexercice


\exercice{4695, quercia, 2010/03/16}
\enonce{004695}{Fonction $\N \to \N$ injective}

Soit $f : {\N} \to {\N}$ injective. Montrer que $f(n) \xrightarrow[n\to\infty]{} +\infty$.
\finenonce{004695}



\finexercice


\exercice{4696, quercia, 2010/03/16}
\enonce{004696}{Fonction $\N \to \N$ injective}
Soit $f : {\N} \to {\N}$ injective.
Montrer que $\frac{f(1)}{1^2} + \frac{f(2)}{2^2} + \dots + \frac{f(n)}{n^2} \xrightarrow[n\to\infty]{} +\infty$.

\finenonce{004696}



\finexercice


\exercice{4697, quercia, 2010/03/16}
\enonce{004697}{Radicaux it{\'e}r{\'e}s}

Soit $u_n = \sqrt{ n + \sqrt{ n-1 + \dots + \sqrt 1 }}$.

\begin{enumerate}
  \item Montrer que la suite $\left(\frac{u_n}{\sqrt n}\right)$ est born{\'e}e.
    
  \item D{\'e}terminer $\lim_{n\to\infty} \left(\frac{u_n}{\sqrt n}\right)$.
    
  \item D{\'e}terminer $\lim_{n\to\infty}(u_n-\sqrt n\,)$.
    
\end{enumerate}
\finenonce{004697}



\finexercice


\exercice{4698, quercia, 2010/03/16}
\enonce{004698}{Ensae MP$^*$ 2000}
Soit $(a_n)$ une suite de r{\'e}els sup{\'e}rieurs ou {\'e}gaux {\`a} $1$ telle que pour tous $n,m$, 
$a_{n+m}\le a_n\, a_m$. On pose $b_n=\frac{\ln a_n}{n}\cdotp$
Montrer que $(b_n)$ converge vers $\inf \{b_n\,|\, n\in\N^*\}$.

\finenonce{004698}



\finexercice


\exercice{4699, quercia, 2010/03/16}
\enonce{004699}{Polytechnique MP$^*$ 2000}
Soit $h$ croissante de~$\R^+$ dans $\R^+$, tendant vers~$+\infty$ en~$+\infty$,
et telle que $h(x+1)-h(x)$ tend vers~$0$ en~$+\infty$.
Soit~$V$ l'ensemble des valeurs d'adh{\'e}rence de la suite de terme g{\'e}n{\'e}ral~$e^{ih(n)}$
Montrer que $V$ est exactement le cercle trigonom{\'e}trique
(i.e. $\{z\in\C,\ |z|=1\}$).


\finenonce{004699}



\finexercice


\exercice{4700, quercia, 2010/03/16}
\enonce{004700}{$u_n^2+u_n-u_{n+1} \to 0$ (X MP$^*$ 2000)}

Soit $u_n$ une suite r{\'e}elle born{\'e}e. On suppose que
$u_n^2+u_n-u_{n+1}\xrightarrow[n\to\infty]{} 0$. Montrer que $u_n\rightarrow 0$. 
\finenonce{004700}



\finexercice


\exercice{4701, quercia, 2010/03/16}
\enonce{004701}{Point fixe (Ensae MP$^*$ 2003)}
Soit une fonction continue~$f$ de~$\R$ dans~$\R$ et~$x_0\in\R$.
On d{\'e}finit $(x_n)_{n\in\N}$ par la relation de r{\'e}currence~: ${x_{n+1} = f(x_n)}$.
Montrer que si la suite $(x_n)$ admet une unique valeur d'adh{\'e}rence alors elle est convergente.

\finenonce{004701}



\finexercice


\exercice{4702, quercia, 2010/03/16}
\enonce{004702}{Suite r{\'e}curente}
Soit $u_0\in\N^*$ et $(u_n)$ la suite d{\'e}finie par la relation
de r{\'e}currence~: $u_{n+1} = u_n^2+1$. Montrer qu'il exitste $a\in\R$ tel
que $u_n = [a^{2^n}]$ pour tout~$n$ o{\`u} $[\ ]$ d{\'e}signe la partie enti{\`e}re.

\finenonce{004702}



\finexercice


\exercice{5221, rouget, 2010/06/30}
\enonce{005221}{***}
Soit $(u_n)_{n\in\Nn}$ une suite réelle. Montrer que si la suite $(u_n)_{n\in\Nn}$ converge au sens de \textsc{Césaro} et est monotone, alors la suite $(u_n)_{n\in\Nn}$ converge.
\finenonce{005221}


\finexercice


\exercice{5222, rouget, 2010/06/30}
\enonce{005222}{**IT}
\label{exo:suprou3bis}
Pour $n$ entier naturel non nul, on pose $H_n=\sum_{k=1}^{n}\frac{1}{k}$ (série harmonique).

\begin{enumerate}
\item  Montrer que~:~$\forall n\in\Nn^*,\;\ln(n+1)<H_n<1+\ln(n)$ et en déduire $\lim_{n\rightarrow +\infty}H_n$.
\item  Pour $n$ entier naturel non nul, on pose $u_n=H_n-\ln(n)$ et $v_n=H_n-\ln(n+1)$. Montrer que les suites $(u_n)$ et $(v_n)$ convergent vers un réel $\gamma\in\left[\frac{1}{2},1\right]$ ($\gamma$ est appelée la constante d'\textsc{Euler}). Donner une valeur approchée de $\gamma$ à $10^{-2}$ près.
\end{enumerate}
\finenonce{005222}


\finexercice


\exercice{5225, rouget, 2010/06/30}
\enonce{005225}{***}
Soient $a$ et $b$ deux réels tels que $0<a<b$. On pose $u_0=a$ et $v_0=b$ puis, pour $n$ entier naturel donné, $u_{n+1}= \frac{u_n+v_n}{2}$ et $v_{n+1}=\sqrt{u_{n+1}v_n}$.
Montrer que les suites $(u_n)$ et $(v_n)$ sont adjacentes et que leur limite commune est égale à $\frac{b\sin(\Arccos(\frac{a}{b}))}{\Arccos(\frac{a}{b})}$.
\finenonce{005225}


\finexercice


\exercice{5226, rouget, 2010/06/30}
\enonce{005226}{**}
Limite quand $n$ tend vers $+\infty$ de 
\begin{enumerate}
 \item $\frac{\sin n}{n}$,
 \item $\left(1+\frac{1}{n}\right)^n$,
 \item $\frac{n!}{n^n}$,
 \item $\frac{E\left((n+\frac{1}{2})^2\right)}{E\left((n-\frac{1}{2})^2\right))}$,
 \item $\sqrt[n]{n^2}$,
 \item $\sqrt{n+1}-\sqrt{n}$,
 \item $\frac{\sum_{k=1}^{n}k^2}{n^3}$,
 \item $\prod_{k=1}^{n}2^{k/2^{2^k}}$.
\end{enumerate}

\finenonce{005226}


\finexercice


\exercice{5227, rouget, 2010/06/30}
\enonce{005227}{**}
Etudier la suite $(u_n)$ définie par $\sqrt{n+1}-\sqrt{n}=\frac{1}{2\sqrt{n+u_n}}$.
\finenonce{005227}


\finexercice


\exercice{5232, rouget, 2010/06/30}
\enonce{005232}{***}
Soit $u$ une suite complexe et $v$ la suite définie par $v_n=|u_n|$. On suppose que la suite $(\sqrt[n]{v_n})$ converge vers un réel positif $l$. Montrer que si $0\leq\ell<1$, la suite $(u_n)$ converge vers $0$ et si $\ell>1$, la suite $(v_n)$ tend vers $+\infty$.
Montrer que si $\ell=1$, tout est possible.
\finenonce{005232}


\finexercice


\exercice{5233, rouget, 2010/06/30}
\enonce{005233}{***}
\begin{enumerate}
\item  Soit $u$ une suite de réels strictement positifs. Montrer que si la suite $(\frac{u_{n+1}}{u_n})$ converge vers un réel $\ell$, alors $(\sqrt[n]{u_n})$ converge et a même limite.
\item  Etudier la réciproque.
\item  Application~:~limites de 
  \begin{enumerate}
  \item $\sqrt[n]{C_{2n}^n}$,
  \item $\frac{n}{\sqrt[n]{n!}}$,
  \item $\frac{1}{n^2}\sqrt[n]{\frac{(3n)!}{n!}}$.
  \end{enumerate}
\end{enumerate}
\finenonce{005233}


\finexercice


\exercice{5234, rouget, 2010/06/30}
\enonce{005234}{*}
Soient $u$ et $v$ deux suites de réels de $[0,1]$ telles que $\lim_{n\rightarrow +\infty}u_nv_n=1$. Montrer que $(u_n)$ et $(v_n)$ convergent vers $1$.

\finenonce{005234}


\finexercice


\exercice{5235, rouget, 2010/06/30}
\enonce{005235}{**}
Montrer que si les suites $(u_n^2)$ et $(u_n^3)$ convergent alors $(u_n)$ converge.
\finenonce{005235}


\finexercice


\exercice{5236, rouget, 2010/06/30}
\enonce{005236}{***T}
Etudier les deux suites $u_n=\left(1+\frac{1}{n}\right)^n$  et $v_n=\left(1+\frac{1}{n}\right)^{n+1}$.
\finenonce{005236}


\finexercice


\exercice{5237, rouget, 2010/06/30}
\enonce{005237}{**T}
Etudier les deux suites $u_n=\sum_{k=0}^{n}\frac{1}{k!}$ et $v_n=u_n+\frac{1}{n.n!}$.

\finenonce{005237}


\finexercice


\exercice{5238, rouget, 2010/06/30}
\enonce{005238}{}
Etudier les deux suites $u_n=\left(\sum_{k=1}^{n}\frac{1}{\sqrt{k}}\right)-2\sqrt{n+1}$ et $v_n=\left(\sum_{k=1}^{n}\frac{1}{\sqrt{k}}\right)-2\sqrt{n}$.
\finenonce{005238}


\finexercice


\exercice{5241, rouget, 2010/06/30}
\enonce{005241}{***}
Montrer que, pour $n\geq2$, 

\begin{center}
$\cos\left(\frac{\pi}{2^{n}}\right)=\frac{1}{2}\sqrt{2+\sqrt{2+...+\sqrt{2}}}$ ($n-1$ radicaux) et $\sin\left(\frac{\pi}{2^{n}}\right)=\frac{1}{2}\sqrt{2-\sqrt{2+...+\sqrt{2}}}$ ($n-1$ radicaux).
\end{center}
En déduire $\lim_{n\rightarrow +\infty}2^n\sqrt{2-\sqrt{2+...+\sqrt{2}}}$ ($n$ radicaux).
\finenonce{005241}


\finexercice


\exercice{5242, rouget, 2010/06/30}
\enonce{005242}{***}
\begin{enumerate}
\item  Montrer que pour $x$ réel strictement positif, on a~:~$\ln(1+x)<x<(1+x)\ln(1+x)$.
\item  Montrer que $\prod_{k=1}^{n}\left(1+\frac{1}{k}\right)^k<e^n<\prod_{k=1}^{n}\left(1+\frac{1}{k}\right)^{k+1}$ et en déduire la limite quand $n$ tend vers $+\infty$ de $\frac{\sqrt[n]{n!}}{n}$.
\end{enumerate}
\finenonce{005242}


\finexercice


\exercice{5244, rouget, 2010/06/30}
\enonce{005244}{**}
Donner un exemple de suite $(u_n)$ divergente, telle que $\forall k\in\Nn^*\setminus\{1\}$, la suite $(u_{kn})$ converge.
\finenonce{005244}


\finexercice


\exercice{5245, rouget, 2010/06/30}
\enonce{005245}{***I}
Soit $f$ une application injective de $\Nn$ dans $\Nn$. Montrer que $\lim_{n\rightarrow +\infty}f(n)=+\infty$.
\finenonce{005245}


\finexercice


\exercice{5247, rouget, 2010/07/04}
\enonce{005247}{****I}
Etude des suites $(u_n)=(\cos na)$ et $(v_n)=(\sin na)$ où $a$ est un réel donné.

\begin{enumerate}
\item  Montrer que si $\frac{a}{2\pi}$ est rationnel, les suites $u$ et $v$ sont périodiques et montrer dans ce cas que $(u_n)$ et $(v_n)$ convergent si et seulement si $a\in2\pi\Zz$.
\item  On suppose dans cette question que $\frac{a}{2\pi}$ est irrationnel .
\begin{enumerate}
\item Montrer que $(u_n)$ converge si et seulement si $(v_n)$ converge .
\item En utilisant différentes formules de trigonométrie fournissant des relations entre $u_n$ et $v_n$, montrer par l'absurde que $(u_n)$ et $(v_n)$ divergent.
\end{enumerate}
\item  On suppose toujours que $\frac{a}{2\pi}$ est irrationnel. On veut montrer que l'ensemble des valeurs de la suite $(u_n)$ (ou $(v_n)$) est dense dans $[-1,1]$, c'est-à-dire que $\forall x\in[-1,1],\;\forall\varepsilon>0,\;\exists n\in\Nn/\;|u_n-x|<\varepsilon$ (et de même pour $v$).
\begin{enumerate}
\item Montrer que le problème se ramène à démontrer que $\{na+2k\pi,\;n\in\Nn\;\mbox{et}\;k\in\Zz\}$ est dense dans $\Rr$.
\item Montrer que $E=\{na+2k\pi,\;n\in\Nn\;\mbox{et}\;k\in\Zz\}$ est dense dans $\Rr$ (par l'absurde en supposant que $\mbox{inf}(E\cap\Rr_+^*)>0$ pour en déduire que $\frac{a}{2\pi}\in\Qq$).
\item Conclure.
\end{enumerate}
\end{enumerate}
\finenonce{005247}


\finexercice


\exercice{5250, rouget, 2010/07/04}
\enonce{005250}{**I}
Soit $(u_n)$ une suite réelle non majorée. Montrer qu'il existe une suite extraite de $(u_n)$ tendant vers $+\infty$.

\finenonce{005250}


\finexercice


\exercice{5251, rouget, 2010/07/04}
\enonce{005251}{***}
Soit $(u_n)$ une suite de réels éléments de $]0,1[$ telle que $\forall n\in\Nn,\;(1-u_n)u_{n+1}>\frac{1}{4}$. Montrer que $(u_n)$ converge vers $\frac{1}{2}$.
\finenonce{005251}


\finexercice


\exercice{5316, rouget, 2010/07/04}
\enonce{005316}{****I}
\begin{enumerate}
\item  Soient $p$ un entier naturel et $a$ un réel. Donner le développement de $(\cos a+i\sin a)^{2p+1}$ puis en choisissant astucieusement $a$, déterminer $\sum_{k=1}^{p}\cotan^2\frac{k\pi}{2p+1}$. En déduire alors $\sum_{k=1}^{p}\frac{1}{\sin^2\frac{k\pi}{2p+1}}$.
\item  Pour $n$ entier naturel non nul, on pose $u_n=\sum_{k=1}^{n}\frac{1}{k^2}$. Montrer que la suite $(u_n)_{n\in\Nn^*}$ converge (pour majorer $u_n$, on remarquera que $\frac{1}{k^2}\leq\frac{1}{k(k-1)}$).
\item  Montrer que pour tout réel $x$ de $]0,\frac{\pi}{2}[$, on a $\cotan x<\frac{1}{x}<\frac{1}{\sin x}$.
\item  En déduire un encadrement de $u_n$ puis la limite de $(u_n)$.
\end{enumerate}
\finenonce{005316}


\finexercice


\exercice{5459, rouget, 2010/07/10}
\enonce{005459}{**}
Déterminer les limites quand $n$ tend vers $+\infty$ de
 
$$1)\;u_n=\frac{1}{n!}\int_{0}^{1}\Arcsin^nx\;dx\;2)\;\int_{0}^{1}\frac{x^n}{1+x}\;dx\;3)\;\int_{0}^{\pi}\frac{n\sin x}{x+n}\;dx.$$ 

\finenonce{005459}


\finexercice
\section{121.02 Suite définie par une relation de récurrence}


\exercice{536, monthub, 2001/11/01}

\enonce{000536}{} Soit $(u_n)$ la suite réelle définie par récurrence en posant $u_0=1$ et $u_{n+1}=\sqrt{1+u_n}$ si $n\in\N^{*}$. \begin{enumerate} \item Montrer que $(u_n)$ est croissante et majorée. \item Montrer que $(u_n)$ converge vers le nombre réel positif $\ell$ qui

vérifie $\ell^2-\ell-1=0$ et calculer $\ell$.

\end{enumerate} \finenonce{000536}



\finexercice 



\exercice{537, monthub, 2001/11/01}

\enonce{000537}{}
Etudier la suite $(u_n)$ d{\'e}finie par
\[
u_{n+1}=\dfrac{1}{2}u_n(u_n^2-3u_n+4) \ \forall n\geq0.
\]
\finenonce{000537}



\finexercice



\exercice{538, ridde, 1999/11/01}

\enonce{000538}{}
\'Etudier les suites :
\begin{enumerate}
\item $u_0 = 0$ et $u_{n + 1} = \sqrt{u_n + 2}$.
\item $u_0 \in \Rr$ et $u_{n + 1} = u_n-u_n^2$.
\end{enumerate}
\finenonce{000538}



\finexercice



\exercice{539, bodin, 2001/11/01}
\video{WJahdgKkrdc}
\enonce{000539}{}
 On consid\`ere la fonction $ f : \R \longrightarrow \R$ d\'efinie
par
$$f (x) = \frac{x^{3}}{9} + \frac{2 x}{3} + \frac{1}{9}$$
et on d\'efinit la suite $(x_{n})_{n \geq 0}$ en posant $x_{0} =
0$ et $x_{n + 1} = f (x_{n})$ pour $n \in \N.$
\begin{enumerate}
  \item Montrer que l'\'equation $x^{3} - 3 x + 1 = 0$ poss\`ede une
solution unique $\alpha \in ]0, 1\slash 2[.$
  \item Montrer que l'\'equation $f (x) = x$ est \'equivalente \`a
l'\'equation $x^{3} - 3 x + 1 = 0$ et en d\'eduire que $\alpha $
est l'unique solution de l'\'equation $f (x) = x$ dans
l'intervalle $[0 , 1\slash 2].$
  \item Montrer que  la fonction $f$ est croissante sur $\R^{+}$ et que $f (\R^{+}) \subset
\R^{+}$. En
d\'eduire que la suite $(x_{n})$ est croissante.
   \item Montrer que
$f (1\slash 2) < 1 \slash 2$ et en d\'eduire que $ 0 \leq  x_{n} <
1 \slash 2$ pour tout $ n \geq 0.$
  \item  Montrer que la suite
$(x_{n})_{n \geq 0}$ converge vers $\alpha.$
 \end{enumerate}
\finenonce{000539} 


\finexercice

\exercice{540, cousquer, 2003/10/01}

\enonce{000540}{}
 Soit $a\in \mathbb{R}$. On considère la suite $(u_n)$ définie par 
$u_0=a$ et $u_{n+1}=e^{u_n}-2$ pour $n\geq 0$.
\begin{enumerate}
\item Étudier cette suite si $a=0$.
\item Étudier cette suite si $a=-10$.
\item Étudier cette suite si $a=3$.
\item Généraliser en discutant selon la valeur de~$a$.
\end{enumerate}
\finenonce{000540}



\finexercice



\exercice{541, cousquer, 2003/10/01}

\enonce{000541}{}
 Étudier la suite définie par $u_{n+1}=1+\frac{u_n^3}{10}$ dans 
les cas suivants~:

\begin{enumerate}
\item $u_0=-4$.
\item $u_0=-2$.
\item $u_0=2$.
\item $u_0=3$.
\end{enumerate}
\finenonce{000541}



\finexercice



\exercice{542, cousquer, 2003/10/01}

\enonce{000542}{}
 Étudier la suite $(u_n)$ définie par $u_0=0$ et 
$u_{n+1}=\frac{(u_n-3)^2}{4}$.
\finenonce{000542}



\finexercice



\exercice{543, cousquer, 2003/10/01}

\enonce{000543}{}
 Étudier la suite définie par $u_{n+1}= e^{-u_n}$ et $u_0=0$.
\finenonce{000543}



\finexercice



\exercice{544, cousquer, 2003/10/01}

\enonce{000544}{}
 Étudier la suite définie par $u_{n+1}= \cos u_n$ et $u_0=-8$.
\finenonce{000544}



\finexercice



\exercice{545, cousquer, 2003/10/01}

\enonce{000545}{}
Étudier la suite définie par $u_{n+1}= \frac{2u_n^3+7}{3(u_n^2+1)}$ 
et $u_0=2$. En déduire une valeur approchée à $10^{-8}$ près de la 
racine réelle du polynôme $X^3+3X-7$.
\finenonce{000545}



\finexercice



\exercice{546, cousquer, 2003/10/01}

\enonce{000546}{}
Étudier la suite définie par $u_0=0$ et 
$u_{n+1}=\frac{-u_n^2-u_n+24}{6}$ pour $n\geq 0$.
\finenonce{000546}



\finexercice



\exercice{547, cousquer, 2003/10/01}

\enonce{000547}{}
 Étudier la suite définie par $u_0=0$ et 
$u_{n+1}=-\frac{3}{13}u_n^2 -\frac{1}{9}u_n+3$ pour $n\geq 0$.
\finenonce{000547}



\finexercice



\exercice{548, cousquer, 2003/10/01}

\enonce{000548}{}
 Étudier la suite définie par $u_0=0$ et 
$u_{n+1}=-\frac{1}{5}u_n^2-\frac{1}{6}u_n+\frac{33}{10}$ pour $n\geq 0$.
\finenonce{000548}



\finexercice



\exercice{549, cousquer, 2003/10/01}

\enonce{000549}{}
 Étudier la suite définie par $u_0=0$ et $u_{n+1}=\ln(e-1+u_n)$.
\finenonce{000549}



\finexercice



\exercice{550, cousquer, 2003/10/01}

\enonce{000550}{}
 Discuter suivant les valeurs de $u_0$ la nature de la suite
$u_{n+1}=e^{u_n}-2$.
\finenonce{000550}



\finexercice



\exercice{551, cousquer, 2003/10/01}

\enonce{000551}{}
 Soient $a$ et $b$ deux réels strictement positifs~; on définit
une suite $(u_n)$ par~:
$$ u_0 \ge0 \quad{\rm et} \quad u_{n+1}= \sqrt{au_n+b}$$
\begin{enumerate}
\item Montrer qu'il existe une valeur de $u_0$ pour laquelle cette suite
est stationnaire.
\item Montrer que si $u_0$ est distinct de cette valeur, $(u_n)$ est monotone et
bornée. Trouver $\lim_{n \to\infty} u_n $.
\end{enumerate}
\finenonce{000551}



\finexercice



\exercice{552, cousquer, 2003/10/01}

\enonce{000552}{}
 Étudier suivant les valeurs données à $u_0$ appartenant à $\mathbb{C}$ les
suites~:
\begin{eqnarray*}
  u_{n+1} & = & \frac{u_n-2}{u_n +4}\\
  u_{n+1} & = & \frac{u_n+2}{u_n+1}\\
  u_{n+1} & = & \frac{-1}{u_n+1}
\end{eqnarray*}
\finenonce{000552}



\finexercice



\exercice{553, cousquer, 2003/10/01}

\enonce{000553}{}
Soit $f\colon [0,1]\rightarrow [0,1]$. On considère $a\in [0,1]$ 
et la suite $(u_n)_{n\in \mathbb{N}}$ vérifiant $u_0=a$ et $\forall n\in 
\mathbb{N},\; u_{n+1}=f(u_n)$. Les propriétés suivantes sont-elles 
vraies ou fausses~:

\begin{enumerate}
\item Si $f$ est croissante, alors $(u_n)$ est croissante.

\item Si $(u_n)$ est croissante, alors $f$ est 
croissante.

\item Si $(u_n)$ est croissante et $f$ monotone, alors $f$ est 
croissante.

\item Si $(u_n)$ converge vers une limite~$l$, alors $l$ est point fixe 
de~$f$.

\item Si $f$ est dérivable, alors $(u_n)$ est bornée.

\item Si le graphe de~$f$ est au dessus de la droite d'équation $y=x$, 
alors $(u_n)$ est croissante.

\item Si $(u_n)$ converge vers un point fixe~$l$ de~$f$, alors $f$ est 
continue en~$l$.
\end{enumerate}
\finenonce{000553}



\finexercice



\exercice{554, cousquer, 2003/10/01}

\enonce{000554}{}
 Étudier la suite définie par $u_{n+1}= (1-u_n)^2$ (discuter 
suivant les valeurs de $u_0$).
\finenonce{000554}



\finexercice



\exercice{555, cousquer, 2003/10/01}

\enonce{000555}{}
 Soit $f(x)=-x^3+x^2-x+1$ et $a\in [0,1]$. Étudier la suite 
définie par $u_0=a$ et $u_{n+1}=f(u_n)$.
\finenonce{000555}



\finexercice



\exercice{556, cousquer, 2003/10/01}

\enonce{000556}{}
 Étudier la suite définie par $u_0=0$ et 
$u_{n+1}=\frac{1}{2}(1+u_n+\mathrm{E}(u_n))$ où $\mathrm{E}$ 
désigne la fonction \og partie entière\fg.
\finenonce{000556}



\finexercice



\exercice{557, cousquer, 2003/10/01}

\enonce{000557}{}
\begin{enumerate}
\item Étudier la suite définie par récurrence par $u_0=a$ et $u_{n+1}=\cos
u_n$, où $a$ est un nombre réel donné.

\item Étudier la suite définie pour
$n\geq1$ par $u_n=\underbrace{\cos(\cos(\cos(\cdots(\cos}
_{n \; \rm fois\; cos} n)\cdots)))$.
\end{enumerate}
\finenonce{000557}



\finexercice



\exercice{558, cousquer, 2003/10/01}

\enonce{000558}{}
\begin{enumerate}
\item Étudier dans $\mathbb{C}$ une suite $(u_n)$ telle que 
$\forall n \in\mathbb{N}^\ast,\;u_{n+1}=u_n^2$. Discuter suivant~$u_0$.

\item On considère dans $\mathbb{C}$ une suite $(v_n)$ telle que 
$\forall n,\; v_{n+1}={1\over2}\bigl(v_n+{A\over v_n}\bigr)$
où $A$ est un nombre complexe non nul donné. Étudier l'existence et la
convergence de cette suite suivant les valeurs de $v_0$. On pourra noter $a$
une des racines carrées de $A$ et poser $w_n={v_n-a\over v_n+a}$.
\end{enumerate}
\finenonce{000558}



\finexercice



\exercice{559, cousquer, 2003/10/01}

\enonce{000559}{}
\begin{enumerate}
\item
On donne $A\geq 0$, $B\geq 0$, $u_0\geq 0$~; étudier la suite définie
par la relation de récurrence
$\displaystyle u_{n+1}=\frac{A}{n+1} + B u_n$.
\item Étudier la suite définie par $u_0=0$ et
$\displaystyle u_{n+1}=\frac{\frac{4n}{n+1} -u_n}{2+u_n}$
(on pourra utiliser la question précédente pour terminer). 
\end{enumerate}
\finenonce{000559}



\finexercice



\exercice{560, cousquer, 2003/10/01}

\enonce{000560}{}
 On considère la suite réelle définie par~:
$$x_0=1 \quad\mbox{et} \quad x_{n+1}= \sqrt{2x_n+1}.$$

\begin{enumerate}
\item Montrer que $x_n$ est supérieur ou égal à~$1$ pour tout~$n$.

\item Montrer que si $(x_n)$ converge, sa limite $l$ vérifie
$$l= \sqrt{2l+1}.$$

\item $l$ étant définie par l'égalité de 2), est-il
possible de trouver $k\in\mathopen]0,1\mathclose[$ tel que
$$\vert x_n-l\vert \leq k \vert x_{n-1}-l\vert.$$
Si oui en déduire que $\vert x_n-l\vert \leq k^n \vert x_0-l\vert$.
Conclure.
\end{enumerate}
\finenonce{000560}



\finexercice



\exercice{561, cousquer, 2003/10/01}

\enonce{000561}{}
 En utilisant les méthodes de l'exercice précédent,
étudier les suites définies par~:
$$\displaylines{y_0=3\;;\quad y_{n+1}=\frac{4+3y_n}{3+2y_n},\cr
z_0=1\;;\quad z_{n+1}=1+ \frac{1}{z_n}.}$$
\finenonce{000561}



\finexercice



\exercice{562, cousquer, 2003/10/01}

\enonce{000562}{}
 Soit une suite qui vérifie une relation de récurrence
\begin{eqnarray}
u_n & = & \frac{au_{n-1}+b}{cu_{n-1}+d}
\label{eq1}
\end{eqnarray}

\begin{enumerate}
\item Montrer que si la transformation homographique~:
$x \mapsto y=\frac{ax+b}{cx+d}$
a deux points fixes distincts, $\alpha$ et~$\beta$, on peut
écrire la relation (\ref{eq1}) sous la forme~:
$\frac{u_n-\alpha}{u_n-\beta} = k\frac{u_{n-1}-\alpha}{u_{n-1}-\beta}$. 
Calculer $\frac{u_n-\alpha}{u_n-\beta}$ en fonction de
$\frac{u_1-\alpha}{u_1-\beta}$.

\item Montrer que si la transformation homographique a un seul point fixe
$\gamma$, on peut mettre la relation (\ref{eq1}) sous la forme~:
$\frac{1}{u_n-\gamma} = \frac{1}{u_{n-1}-\gamma}+k$.
Calculer $\frac{1}{u_n-\gamma}$ en fonction de~$u_1$.


\item Utiliser la méthode précédente pour étudier les suites $(u_n)$
définies par~:
\begin{center}
\renewcommand{\arraystretch}{2}
\renewcommand{\tabcolsep}{2 em}
\begin{tabular}{ll}
a) $\displaystyle u_{n+1}=\frac{4u_n+2}{u_n+3}$,&
b) $\displaystyle u_{n+1}=\frac{-3u_n-1}{u_n-3}$,\\ 
c) $\displaystyle u_{n+1}=\frac{5u_n-3}{u_n+1}$, &
d) $\displaystyle u_{n+1}=\frac{2u_n-1}{u_n+4}$.
\end{tabular}
\end{center}

Discuter suivant les valeurs de $u_1$~; préciser pour quelles valeurs de $u_1$
chaque suite est définie.
\end{enumerate}
\finenonce{000562}



\finexercice



\exercice{4703, quercia, 2010/03/16}
\enonce{004703}{Étude de suites}

Étudier la convergence de la suite $(u_n)$ définie par :
\begin{enumerate}
  \item $u_0 = a > 1$, $u_{n+1} = \frac 12\left(u_n + \frac a{u_n}\right)$.
  \item $0 < u_0 < \frac {\sqrt5-1}2$, $u_{n+1} = 1 - u_n^2$.
  \item $u_{n+1} = u_n-u_n^2$.
  \item $u_0 = 0$, $u_{n+1} = u_n^2+\alpha$.
  \item $u_{n+1} = u_n + \frac{1+u_n}{1+2u_n}$.
  \item $u_0 \in\, [0,1]$, $u_{n+1} = \frac{\sqrt{u_n}}{\sqrt{u_n}+\sqrt{1-u_n}}$.
  \item $u_{n+1} = \sqrt{2-u_n}$.
  \item $u_{n+1} = \sqrt{4-3u_n}$.
  \item $u_{n+1} = \frac{u_n-\ln(1+u_n)}{u_n^2}$.
  \item $u_{n+1} = \frac3{2u_n^2+1}$.
  \item $u_0 > 0$, $u_{n+1} = u_n^\alpha$.
  \item $u_0 > 0$, $u_{n+1} = \alpha^{u_n}$.
\end{enumerate}
\finenonce{004703}



\finexercice


\exercice{4704, quercia, 2010/03/16}
\enonce{004704}{Convergence quadratique}

Soit $k\in\C$ fix{\'e}. {\'E}tudier la convergence de la suite $(a_n)$ d{\'e}finie par :
$a_0\in \C$, $a_{n+1} = k a_n^2$.


\finenonce{004704}



\finexercice


\exercice{4705, quercia, 2010/03/16}
\enonce{004705}{$u_{n+1}(1-u_n) > 1/4$}

Soit $(u_n)$ une suite r{\'e}elle telle que pour tout entier $n$ :
$u_n \in\,[0,1]$ et $u_{n+1}(1-u_n) > \frac 14$.\par
Montrer que cette suite converge vers $\frac 12$.


\finenonce{004705}



\finexercice


\exercice{4706, quercia, 2010/03/16}
\enonce{004706}{Radicaux it{\'e}r{\'e}s}

Trouver $\lim_{n\to\infty} \sqrt{1 + \sqrt{1 + \dots + \sqrt{1}}}$
($n$ radicaux).


\finenonce{004706}



\finexercice


\exercice{4707, quercia, 2010/03/16}
\enonce{004707}{Radicaux it{\'e}r{\'e}s}

On consid{\`e}re la suite $(u_n)$ d{\'e}finie par :
$u_0 > 0$, $u_{n+1} = \sqrt{ u_0 + \dots + u_n }$.
Montrer que $u_n \xrightarrow[n\to\infty]{} +\infty$.
\finenonce{004707}



\finexercice


\exercice{4708, quercia, 2010/03/16}
\enonce{004708}{Radicaux it{\'e}r{\'e}s}

On pose $u_n = \sqrt{1 + \sqrt{2 + \sqrt{ \dots + \sqrt{n-1 +\sqrt n}}}}$
     et $v_n = \sqrt{1 + \sqrt{2 + \sqrt{ \dots + \sqrt{n-1 +\sqrt{2n}}}}}$.

\begin{enumerate}
  \item Montrer que ces suites sont convergentes.
  \item On note $\lambda = \lim_{n\to\infty} u_n$.
    Montrer que $\lambda-u_n \le \frac n{2^n\sqrt{n!}}$.
\end{enumerate}
\finenonce{004708}



\finexercice


\exercice{4709, quercia, 2010/03/16}
\enonce{004709}{Suites homographiques}

Soient $a,b \in \R^*$.
On d{\'e}finit la suite $(u_n)$ par :
$\begin{cases} u_0 \in \R^* \cr u_{n+1} = a + \frac b{u_n}.\end{cases}$

On suppose $u_0$ choisi de sorte que pour tout $n \in \N$, $u_n \ne 0$.

\begin{enumerate}
  \item Quelles sont les limites possibles pour $(u_n)$ ?
  \item On suppose que l'{\'e}quation $x^2 = ax + b$ poss{\`e}de deux racines r{\'e}elles
    $\alpha,\beta$ avec $|\alpha| > |\beta|$.

    {\'E}tudier la suite $(v_n) = \left(\frac{u_n-\alpha}{u_n-\beta}\right)$
    et en d{\'e}duire $\lim u_n$.
\end{enumerate}
\finenonce{004709}



\finexercice


\exercice{4710, quercia, 2010/03/16}
\enonce{004710}{Syst{\`e}me d'ordre 1}

Soient $0 < x_0 < y_0$ et $(x_n)$, $(y_n)$ les suites d{\'e}finies par :
 $\begin{cases}x_{n+1} = \frac {x_n^2}{x_n+y_n}\cr
         y_{n+1} = \frac {y_n^2}{x_n+y_n}.\cr\end{cases}$

Montrer qu'elles sont convergentes et calculer leurs limites.


\finenonce{004710}



\finexercice


\exercice{4711, quercia, 2010/03/16}
\enonce{004711}{Syst{\`e}me d'ordre 1}

{\'E}tudier la convergence des suites $(x_n)$, $(y_n)$ d{\'e}finies par :
$0 < x_0 < y_0$ et $\begin{cases} x_{n+1} = \frac{2x_n+y_n}3 \cr
                            y_{n+1} = \frac{2y_n+x_n}3.\cr\end{cases}$
\finenonce{004711}



\finexercice


\exercice{4712, quercia, 2010/03/16}
\enonce{004712}{Syst{\`e}me d'ordre 1}

{\'E}tudier la convergence des suites $(x_n)$, $(y_n)$ d{\'e}finies par :
$0 < y_0 < x_0$ et $\begin{cases}x_{n+1} = \frac{x_n+y_n}2 \cr
                           y_{n+1} = \frac{2x_ny_n}{x_n+y_n} .\cr\end{cases}$


\finenonce{004712}



\finexercice


\exercice{4713, quercia, 2010/03/16}
\enonce{004713}{Syst{\`e}me d'ordre 1}

Soient $0 < a < b$ et $(x_n),(y_n)$ les suites d{\'e}finies par :
$\begin{cases}x_0 = a\cr y_0 = b\cr\end{cases}$ \quad et \quad
$\begin{cases}x_{n+1} = \frac {x_n+y_n}2 \cr
        y_{n+1} = \sqrt{ x_{n+1}y_n }.\cr\end{cases}$

\begin{enumerate}
  \item Montrer que ces suites convergent vers la m{\^e}me limite.
  \item On pose $a = b\cos\varphi$. Exprimer cette limite en fonction de $b$ et
    $\varphi$.
    
\end{enumerate}
\finenonce{004713}



\finexercice


\exercice{4714, quercia, 2010/03/16}
\enonce{004714}{Moyennes arithm{\'e}tique, g{\'e}om{\'e}trique, harmonique}

\begin{enumerate}
  \item Soient $x,y,z \ge 0$. Montrer que $x^3+y^3+z^3 - 3xyz \ge 0$
    (mettre $x+y+z$ en facteur).
  \item {\'E}tudier la convergence des suites $(a_n)$, $(b_n)$, $(c_n)$ d{\'e}finies par :
    $$0 < a_0 < b_0 < c_0,\quad \text{et}\quad
     \begin{cases}\frac 3{a_{n+1}} = \frac 1{a_n} + \frac 1{b_n} + \frac 1{c_n} \cr
            b_{n+1} = \sqrt[3]{a_nb_nc_n} \cr
            3c_{n+1} = a_n + b_n + c_n. \cr\end{cases}$$
\end{enumerate}
\finenonce{004714}



\finexercice


\exercice{4715, quercia, 2010/03/16}
\enonce{004715}{Centrale MP 2000}

On consid{\`e}re la fonction $f$ : $x \mapsto\ln \bigl(\frac{e^x-1}{x}\bigr)$ et la suite d{\'e}finie par
$\begin{cases}u_0\in\R^*\cr u_{n+1}=f(u_n).\cr\end{cases}$
{\'E}tudier la suite $(u_n)$, puis la s{\'e}rie $\sum u_n$.

\finenonce{004715}



\finexercice


\exercice{4716, quercia, 2010/03/16}
\enonce{004716}{$u_{n+1}-u_n \to0$}
Soit $f : {[a,b]} \to {[a,b]}$ continue et la suite $(u_n)$ d{\'e}finie par
$u_0\in{[a,b]}$ et $u_{n+1} = f(u_n)$. Montrer que si $\lim(u_{n+1}-u_n) = 0$
alors la suite $(u_n)$ converge.

\finenonce{004716}



\finexercice


\exercice{5228, rouget, 2010/06/30}
\enonce{005228}{**T Récurrences homographiques}
Déterminer $u_n$ en fonction de $n$ quand la suite $u$ vérifie~:
\begin{enumerate}
 \item $\forall n\in\Nn,\;u_{n+1}=\frac{u_n}{3-2u_n}$,
 \item $\forall n\in\Nn,\;u_{n+1}=\frac{4(u_n-1)}{u_n}$ (ne pas se poser de questions d'existence).
\end{enumerate}
\finenonce{005228}


\finexercice


\exercice{5229, rouget, 2010/06/30}
\enonce{005229}{**}
Soient $(u_n)$ et $(v_n)$ les suites définies par la donnée de $u_0$ et $v_0$ et les relations de récurrence 

$$u_{n+1}=\frac{2u_n+v_n}{3}\;\mbox{et}\;v_{n+1}=\frac{u_n+2v_n}{3}.$$
Etudier les suites $u$ et $v$ puis déterminer $u_n$ et $v_n$ en fonction de $n$ en recherchant des combinaisons linéaires intéressantes de $u$ et $v$. En déduire $\lim_{n\rightarrow +\infty}u_n$ et $\lim_{n\rightarrow +\infty}v_n$.
\finenonce{005229}


\finexercice


\exercice{5230, rouget, 2010/06/30}
\enonce{005230}{**}
Soient $(u_n)$, $(v_n)$ et $(w_n)$ les suites définies par la donnée de $u_0$,
$v_0$ et $w_0$ et les relations de récurrence 
$$u_{n+1}=\frac{v_n+w_n}{2},\,v_{n+1}=\frac{u_n+w_n}{2}\;\mbox{et}\;
w_{n+1}=\frac{u_n+v_n}{2}.$$

Etudier les suites $u$, $v$ et $w$ puis déterminer $u_n$, $v_n$ et $w_n$ en fonction de $n$
en recherchant des combinaisons linéaires intéressantes de $u$, $v$ et $w$. En
déduire $\lim_{n\rightarrow +\infty}u_n$, $\lim_{n\rightarrow +\infty}v_n$et $\lim_{n\rightarrow +\infty}w_n$.

\finenonce{005230}


\finexercice


\exercice{5231, rouget, 2010/06/30}
\enonce{005231}{***}
Montrer que les suites définies par la donnée de $u_0$, $v_0$ et $w_0$ réels tels que $0<u_0<v_0<w_0$ et les relations de récurrence~:
 
$$\frac{3}{u_{n+1}}=\frac{1}{u_n}+\frac{1}{v_n}+\frac{1}{w_n}\;\mbox{et}\;v_{n+1}=\sqrt[3]{u_nv_nw_n}\;\mbox{et}\;w_{n+1}=\frac{u_n+v_n+w_n}{3},$$  
 
 
ont une limite commune que l'on ne cherchera pas à déterminer.

\finenonce{005231}


\finexercice

\exercice{5240, rouget, 2010/06/30}
\enonce{005240}{****}
On pose $u_1=1$ et, $\forall n\in\Nn^*,\;u_{n+1}=1+\frac{n}{u_n}$. Montrer que $\lim_{n\rightarrow +\infty}(u_n-\sqrt{n})=\frac{1}{2}$.
\finenonce{005240}


\finexercice


\exercice{5277, rouget, 2010/07/04}
\enonce{005277}{**I}
On pose $u_0=1$, $v_0=0$, puis, pour $n\in\Nn$, $u_{n+1}=2u_n+v_n$ et $v_{n+1}=u_n+2v_n$.
\begin{enumerate}
\item  Soit $A=\left(
\begin{array}{cc}
2&1\\
1&2
\end{array}
\right)$. Pour $n\in\Nn$, calculer $A^n$. En déduire $u_n$ et $v_n$ en fonction de $n$.
\item  En utilisant deux combinaisons linéaires intéressantes des suites $u$ et $v$, calculer directement $u_n$ et $v_n$ en fonction de $n$.
\end{enumerate}

\finenonce{005277}


\finexercice


\exercice{5425, rouget, 2010/07/06}
\enonce{005425}{***I}
Etudier la suite $(u_n)$ dans chacun des cas suivants~:

$$\begin{array}{ll}
1)\;u_0\geq-1\;\mbox{et}\;\forall n\in\Nn,\;u_{n+1}=\sqrt{1+u_n},&2)\;u_0>-1\;\mbox{et}\;\forall n\in\Nn,\;u_{n+1}=\ln(1+u_n)\\
3)\;u_0\in\Rr\;\mbox{et}\;\forall n\in\Nn,\;u_{n+1}=\sin u_n,&4)\;u_0\in\Rr\;\mbox{et}\;\forall n\in\Nn,\;u_{n+1}=\cos(u_n),\\
5)\;u_0\in\Rr\;\mbox{et}\;\forall n\in\Nn,\;u_{n+1}=\sin(2u_n),&6)\;u_0\in\Rr\;\mbox{et}\;\forall n\in\Nn,\;u_{n+1}=u_n^2-2u_n+2.
\end{array}
$$
\finenonce{005425}


\finexercice

\exercice{5435, rouget, 2010/07/06}
\enonce{005435}{***I}
Soit $u_0\in]0,\frac{\pi}{2}]$. Pour $n\in\Nn$, on pose $u_{n+1}=\sin(u_n)$.
\begin{enumerate}
\item  Montrer brièvement que la suite $u$ est strictement positive et converge vers $0$.
\item 
\begin{enumerate}
\item Déterminer un réel $\alpha$ tel que la suite $u_{n+1}^\alpha-u_n^\alpha$ ait une limite finie non nulle.
\item En utilisant le lemme de \textsc{Cesaro}, déterminer un équivalent simple de $u_n$.
\end{enumerate}
\end{enumerate}

\finenonce{005435}


\finexercice


\exercice{5436, rouget, 2010/07/06}
\enonce{005436}{**I}
Soit $u$ la suite définie par la donnée de son premier terme $u_0>0$ et la relation $\forall n\in\Nn,\;u_{n+1}=u_ne^{-u_n}$. Equivalent simple de $u_n$ quand $n$ tend vers $+\infty$.
\finenonce{005436}


\finexercice

\section{121.03 Suites équivalentes, suites négligeables}


\exercice{563, monthub, 2001/11/01}
\video{FBzl-Zyr1e0}
\enonce{000563}{}
Posons $  u_2=1-\frac{1}{2^2}$ et pour tout entier $n\geq 3$,
\[
  u_n=\left(1-\frac{1}{2^2}\right)\left(1-\frac{1}{3^2}\right)\cdots\left(1-\frac{1}{n^2}\right).\]
Calculer $u_n$. En d{\'e}duire que l'on a
$\lim{u_n}=\dfrac{1}{2}$.
\finenonce{000563} 


\finexercice

\exercice{564, monthub, 2001/11/01}

\enonce{000564}{}
Calculer, lorsqu'elles convergent, les limites des suites
d{\'e}finies par
:\\
$u_n=n-\sqrt{n^2-n}$ \qquad $u_n=\sqrt{n(n+a)}-n$ \qquad
 $u_n=\dfrac{n}{2} \sin{\dfrac{n\pi}{2}}$
\qquad $u_n=\dfrac{\sin{n}^2-\cos{n}^3}{n}$.
\finenonce{000564}



\finexercice



\exercice{565, monthub, 2001/11/01}

\enonce{000565}{}
Montrer que les suites d{\'e}finies pour $n\geq1$ par :\\
$u_n=\dfrac{n+1}{n}$ \qquad $u_n=\dfrac{n}{n+1}$ \qquad
$u_n=\dfrac{1}{n^2+1}$ \qquad $u_n=\frac{n}{n^2+1} $\\[0.2cm]
admettent toutes des limites que l'on calculera.
\finenonce{000565}



\finexercice



\exercice{566, monthub, 2001/11/01}

\enonce{000566}{}
Soit $(u_n)_{n\in\N}$ la suite de nombres r{\'e}els d{\'e}finie en
posant $u_0=0$ et $\forall n\geq1, \ u_{n+1}=\sqrt{6+u_n}$.
Montrer que la suite $(u_n)_{n\in\N}$ est convergente et
d{\'e}terminer sa limite.
\finenonce{000566}



\finexercice



\exercice{567, monthub, 2001/11/01}

\enonce{000567}{}
Etudier la limite des suites suivantes :
$a_n=\cos{(\dfrac{2^n}{n!})}$; $b_n=\sqrt[n]{3-\sin{n}^2}$;
$c_n=\dfrac{n^3+2^n}{3^n}$;
$d_n=\dfrac{n^2+(-1)^n}{n^2+\sqrt{n}}$;
$e_n=(\cos{n})\sin{\dfrac{(-1)^n}{\sqrt{n}}}$.
\finenonce{000567}



\finexercice



\exercice{568, cousquer, 2003/10/01}
\video{M2V-1XVn_ig}
\enonce{000568}{}
D\'eterminer les limites lorsque $n$ tend vers l'infini des suites
ci-dessous~; pour chacune, essayer de pr\'eciser en quelques mots la
m\'ethode employ\'ee.
\begin{enumerate}
\item $\displaystyle 1\;;~-\frac{1}{2}\;;~\frac{1}{3}\;;~ 
\ldots\;;~\frac{(-1)^{n-1}}{n}\;;~\ldots$
\item $2/1$~; $4/3$~; $6/5$~; $\ldots$~; $2n/(2n-1)$~; $\ldots$
\item $0{,}23\;;~0{,}233\;;~\ldots\;;~0{,}233\cdots3\;;~\ldots$
\item $\displaystyle\frac{1}{n^2}+\frac{2}{n^2}+\cdots+\frac{n-1}{n^2}$
\item $\displaystyle\frac{(n+1)(n+2)(n+3)}{n^3}$
\item $\displaystyle\biggl\lbrack \frac{1+3+5+\cdots+(2n-1)}{n+1} -
\frac{2n+1}{2}\biggr\rbrack$
\item $\displaystyle\frac{n+(-1)^n}{n-(-1)^n}$
\item  $\displaystyle\frac{2^{n+1}+3^{n+1}}{2^n + 3^n}$
\item $\displaystyle\bigl(1/2+1/4+1/8+\cdots+1/2^n\bigr)$\quad 
puis\quad $\displaystyle \sqrt{2}\;;~\sqrt{2\sqrt{2}}\;;~
\sqrt{2\sqrt{2\sqrt{2}}}\;;~\ldots$
\item  $\displaystyle\biggl(1-\frac{1}{3}+\frac{1}{9}-\frac{1}{27}+\cdots
+\frac{(-1)^n}{3^n} \biggr)$
\item $\bigl( \sqrt{n+1}-\sqrt{n}\bigr)$
\item $\displaystyle\frac{n\sin(n!)}{n^2+1}$
\item  D\'emontrer la formule $1+2^2+3^2+\cdots+n^2 = 
\frac{1}{6}
n(n+1)(2n+1)$ ; en d\'eduire $\lim_{n \to\infty}
\frac{1+2^2+3^2+\cdots+n^2}{n^3}$.
\end{enumerate}
\finenonce{000568} 


\finexercice

\exercice{569, bodin, 2001/11/01}
\video{sHmOnEw9rxQ}
\enonce{000569}{}
 Soit $a>0$. On d\'efinit la suite $(u_n)_{n\geq 0}$ par
$u_0$ un r\'eel v\'erifiant $u_0>0$ et par la relation
$$u_{n+1}= \frac12 \left( u_n+\frac{a}{u_n}\right).$$
On se propose de montrer que $(u_n)$ tend vers $\sqrt a$.
\begin{enumerate}
  \item Montrer que
$${u_{n+1}}^2-a= \frac{({u_n}^2-a)^2}{4{u_n}^2}.$$
\item Montrer que si $n\geq 1$ alors $u_n \geq \sqrt a$ puis que
la suite $(u_n)_{n\geq 1}$ est d\'ecroissante.
\item En d\'eduire que la suite $(u_n)$ converge vers $\sqrt a$.
\item En utilisant la relation
${u_{n+1}}^2-a= ({u_{n+1}}-\sqrt{a})({u_{n+1}}+\sqrt{a})$ donner
une majoration de ${u_{n+1}}-\sqrt{a}$ en fonction de
${u_{n}}-\sqrt{a}$.
\item Si $u_1-\sqrt a \leq k$ et pour $n\geq 1$ montrer que
$$u_n - \sqrt a \leq 2\sqrt a \left( \frac k {2\sqrt a}\right)^{2^{n-1}}.$$
\item Application : Calculer $\sqrt{10}$ avec une pr\'ecision de 8 chiffres apr\`es la virgule,
en prenant $u_0 = 3$.
\end{enumerate}
\finenonce{000569} 


\finexercice

\exercice{570, bodin, 1998/09/01}
\video{k8gZEvm5MQs}
\enonce{000570}{}
 On consid\`ere les deux suites :
$$u_n =1+\frac{1}{2!}+\frac{1}{3!}+\cdots+\frac{1}{n!}\ ;\ n\in\N,$$
$$v_n = u_n+\frac{1}{n!}\ ;\ n\in\N.$$
\noindent Montrer que $(u_n)_n$ et $(v_n)_n$ convergent vers une m\^{e}me
limite. Et montrer que cette limite est un \'el\'ement de $\R\backslash\Q$.
\finenonce{000570} 


\finexercice

\exercice{571, bodin, 1998/09/01}
\video{4ehFn1cnohc}
\enonce{000571}{}
 Soient $a$ et $b$ deux r\'eels, $a<b$. On consid\`ere la
fonction $f:\lbrack a,b\rbrack\longrightarrow \lbrack a,b\rbrack$ suppos\'ee continue
et une suite r\'ecurrente $(u_n)_n$ d\'efinie par :
$$u_0\in\lbrack a,b\rbrack\ \ \text{et pour tout }\ n\in\N,\ \ u_{n+1}=f(u_n).$$
\begin{enumerate}
    \item On suppose ici que $f$ est croissante. Montrer que $(u_n)_n$
est monotone et en d\'eduire sa convergence vers une solution de l'\'equation
$f(x)=x$.
    \item  \emph{Application.} Calculer la limite de la suite définie par :
$$u_0=4\ \ \text{et pour tout }\ n\in\N,\ \ u_{n+1}=\frac{4u_n+5}{u_n+3}.$$
    \item  On suppose maintenant que $f$ est d\'ecroissante. Montrer que les suites
$(u_{2n})_n$ et $(u_{2n+1})_n$ sont monotones et convergentes.
    \item  \emph{Application.} Soit
$$u_0=\frac{1}{2}\ \ \text{et pour tout }\  n\in\N,\ \ u_{n+1}=(1-u_n)^2.$$
\noindent Calculer les limites des suites $(u_{2n})_n$ et $(u_{2n+1})_n$.
\end{enumerate}
\finenonce{000571} 


\finexercice

\exercice{572, bodin, 1998/09/01}
\video{ZTrWGqYQ0R8}
\enonce{000572}{}
\begin{enumerate}
    \item Soient $a,b > 0$. Montrer que $\sqrt{ab} \leqslant \frac{a+b}{2}$.
    \item Montrer les in\'egalit\'es suivantes ($b \geqslant a > 0$) :
$$ a \leqslant \frac{a+b}{2} \leqslant b \qquad \text{et} \qquad a \leqslant \sqrt{ab} \leqslant b.$$
    \item Soient $u_0$ et $v_0$ des r\'eels strictement positifs avec
$u_0 < v_0$. On d\'efinit deux
suites $(u_n)$ et $(v_n)$ de la fa\c{c}on suivante :
$$ u_{n+1} = \sqrt{u_nv_n} \quad \text{et}\quad v_{n+1}=\frac{u_n+v_n}{2}.$$
    \begin{enumerate}
        \item Montrer que $u_n \leqslant v_n$ quel que soit $n\in\Nn$.
        \item Montrer que $(v_n)$ est une suite d\'ecroissante.
        \item Montrer que $(u_n)$ est croissante En d\'eduire que
les suites $(u_n)$ et $(v_n)$ sont convergentes et quelles ont m\^eme limite.
    \end{enumerate}
\end{enumerate}
\finenonce{000572} 


\finexercice


\exercice{573, ridde, 1999/11/01}

\enonce{000573}{}
Soit $x$ un r\'eel.
\begin{enumerate}
\item D\'eterminer la limite de $u_n = \dfrac{E (x) + E (2x) + \ldots + E (nx)}{n^2}$.
\item En d\'eduire que $\Qq$ est dense dans $\Rr$.
\end{enumerate}
\finenonce{000573}



\finexercice



\exercice{574, ridde, 1999/11/01}
\video{1x3LSvkKLsM}
\enonce{000574}{}
 Soit $n \geq 1$.
\begin{enumerate}
\item Montrer que l'\'equation $\sum\limits_{k = 1}^n{x^k} = 1$ admet une unique solution,
notée $a_n$, dans $[0, 1]$.
\item Montrer que $ (a_n)_{n \in \Nn}$ est d\'ecroissante minor\'ee par $\frac12$.
\item Montrer que $ (a_n)$ converge vers $\frac 12$.
\end{enumerate}
\finenonce{000574} 


\finexercice

\exercice{575, cousquer, 2003/10/01}

\enonce{000575}{}
 Calculer suivant les valeurs de $x$~:
$$\lim_{n\to\infty}\Bigl\lbrack
\lim_{m\to\infty}\lbrack\cos(n!\pi x)\rbrack^{2m}\Bigr\rbrack.$$
\finenonce{000575}



\finexercice



\exercice{576, ridde, 1999/11/01}

\enonce{000576}{}
Soient $a_0$ et $b_0$ deux r\'eels fix\'es. On d\'efinit par r\'ecurrence les suites
$ (a_n)$ et $ (b_n)$ par $a_{n + 1} = \dfrac{2a_n + b_n}{3}$ et
 $b_{n + 1} = \dfrac{a_n + 2b_n}{3}$.
 \begin{enumerate}
 \item Montrer que ces deux suites sont adjacentes.
 \item En calculant $a_n + b_n$, montrer qu'elles convergent vers $\dfrac{a_0 + b_0}2$.
 \end{enumerate}
\finenonce{000576}



\finexercice



\exercice{577, ridde, 1999/11/01}

\enonce{000577}{}
Soit $ (u_n)$ une suite qui tend vers $0$. On pose $x_n = \dfrac1n \sum\limits_{k = 0}
^{n-1}u_k$. Montrer que $ (x_n)$ converge vers $0$ ( on pourra fixer $\epsilon$ puis
s\'eparer la somme en deux et enfin choisir $N$... ).
\finenonce{000577}



\finexercice



\exercice{578, ridde, 1999/11/01}

\enonce{000578}{}
D\'eterminer les limites de $\dfrac{n^{\ln (n)}}{\ln^n (n)} $ et $\sqrt[n]{n^2}$.
\finenonce{000578}



\finexercice



\exercice{579, gourio, 2001/09/01}

\enonce{000579}{}
Soit $(u_{n})_{n\in \Nn\text{ }}$ une suite r\'{e}elle dont tous les termes
sont non nuls et telle que :
$$\lim\limits_{n\rightarrow \infty }\left| \frac{u_{n+1}}{u_{n}}\right| =0.$$
Montrer que $\lim\limits_{n\rightarrow \infty }u_{n}=0.$
\finenonce{000579}



\finexercice



\exercice{580, gourio, 2001/09/01}

\enonce{000580}{}
\'Etudier la suite d\'{e}finie par r\'{e}currence :
$$u_{0}=a>0,u_{n+1}=\sqrt{1+u_{n}.}$$
\finenonce{000580}



\finexercice



\exercice{581, gourio, 2001/09/01}

\enonce{000581}{}
\'Etudier la convergence et calculer la limite \'{e}ventuelle de la suite
$(u_{n})_{n\in \Nn^{*}}$ d\'{e}finie par :
$$\forall n\in \Nn^{*},u_{n}=\prod\limits_{k=1}^{n}(1+\frac{k}{n^{2}}).$$
\finenonce{000581}



\finexercice



\exercice{582, gourio, 2001/09/01}

\enonce{000582}{}
\'Etudier la convergence et calculer la limite \'{e}ventuelle de la suite
$(u_{n})_{n\in \Nn}$ d\'{e}finie par :
$$\forall n\in \Nn,u_{n}=\sum\limits_{k=0}^{n}\frac{1}{C_{n}^{k}}.$$
\finenonce{000582}



\finexercice



\exercice{583, gourio, 2001/09/01}

\enonce{000583}{}
Soit $\phi :\Nn\rightarrow \Nn$ bijective, telle que $\lim\limits_{n\rightarrow
\infty }\frac{\phi (n)}{n}=\ell.$ Calculer $\ell$.
\finenonce{000583}



\finexercice



\exercice{584, gourio, 2001/09/01}

\enonce{000584}{}
Soit $\phi :\Nn\rightarrow \Nn$ injective; montrer que $\lim\limits_{n%
\rightarrow \infty }\phi (n)=+\infty .$
\finenonce{000584}



\finexercice



\exercice{585, gourio, 2001/09/01}

\enonce{000585}{}
Soit $(u_{n})_{n\in \Nn}$ une suite born\'{e}e.\ On pose $v_{n}=u_{n+1}-u_{n}$ et
$w_{n}=v_{n+1}-v_{n}$, et on suppose que $(w_{n})_{n\in \Nn} $ converge. Montrer que
$\lim\limits_{n\rightarrow \infty }w_{n}=0$, puis que
$\lim\limits_{n\rightarrow \infty }v_{n}=0.$
\finenonce{000585}



\finexercice



\exercice{586, gourio, 2001/09/01}

\enonce{000586}{}
Soit $(u_{n})_{n\in \Nn}$ une suite r\'{e}elle convergeant vers $\ell$ et $\phi$
une bijection de $\Nn$ sur $\Nn.$ (pas n\'{e}cessairement strictement
croissante !). Montrer que $\lim\limits_{n\rightarrow \infty }u_{\phi(n)}=\ell.$
\finenonce{000586}



\finexercice



\exercice{587, gourio, 2001/09/01}

\enonce{000587}{}
Soient $(u_{n})_{n\in \Nn} $ et $(v_{n})_{n\in \Nn}$ deux suites r\'{e}elles
telles que :
$$\lim_{n\rightarrow \infty }u_{n}+v_{n}=\lim_{n\rightarrow \infty}u_{n}v_{n}=0.$$
Montrer que
$$\lim_{n\rightarrow \infty }v_{n}=\lim_{n\rightarrow \infty }u_{n}=0.$$
\finenonce{000587}



\finexercice



\exercice{588, gourio, 2001/09/01}

\enonce{000588}{}
Soient $(u_{n})_{n\in \Nn} $ et $(v_{n})_{n\in \Nn}$ deux suites r\'{e}elles
telles que :
$$\lim_{n\rightarrow \infty }u_{n}=\lim_{n\rightarrow \infty }v_{n}=+\infty
,\lim_{n\rightarrow \infty }u_{n+1}-u_{n}=0.$$
Montrer que
$$E=\{u_{n}-v_{m}|(n,m)\in \Nn^{2}\}$$ est dense dans $\Rr$.
\finenonce{000588}



\finexercice



\exercice{589, gourio, 2001/09/01}

\enonce{000589}{}
Soient $(u_{n})_{n\in \Nn} $ et $(v_{n})_{n\in \Nn}$ deux suites \`{a} valeurs
dans $[0,1]$ telles que :
$$\lim_{n\rightarrow \infty }u_{n}v_{n}=1.$$
Montrer que :
$$\lim_{n\rightarrow \infty }u_{n}=\lim_{n\rightarrow \infty }v_{n}=1.$$
\finenonce{000589}



\finexercice



\exercice{590, gourio, 2001/09/01}

\enonce{000590}{}
Soient $(u_{n})_{n\in \Nn} $ et $(v_{n})_{n\in \Nn}$ deux suites convergeant
respectivement vers $\ell$ et $L$. \'Etudier la suite $(w_{n})_{n\in \Nn}$
d\'{e}finie par :
$$\forall n\in \Nn,w_{n}=\frac{1}{n}\sum\limits_{k=0}^{n}u_{k}v_{n-k}.$$
\finenonce{000590}



\finexercice



\exercice{591, gourio, 2001/09/01}

\enonce{000591}{}
Soit $(u_{n})_{n\in \Nn}$ une suite born\'{e}e telle que :
$$\lim\limits_{n\rightarrow \infty }(u_{n}+\frac{u_{2n}}{2})=1.$$
Que dire de $(u_{n})_{n\in \Nn}$?
\finenonce{000591}



\finexercice



\exercice{592, gourio, 2001/09/01}

\enonce{000592}{}
Soit $f:\Cc\rightarrow \Cc$ d\'{e}finie par :
$$\forall z\in \Cc,\text{ }f(z)=\frac{z+\left| z\right| }{2}.$$
\'Etudier la suite d\'{e}finie par :
$$z_{0}\in \Cc,\forall n\in \Nn,z_{n+1}=f(z_{n}).$$
\emph{Indication} : on \'{e}crira $z_{n}=\rho _{n}e^{i\phi _{n}}, $o\`{u}
$(\rho_{n},\phi _{n})\in \Rr^{+*}\times ]-\pi ,\pi [$ et on utilisera :
$$\sin \phi =2^{n}\sin \frac{\phi }{2^{n}}\prod\limits_{i=1}^{n}\cos \frac{\phi }{2^{i}}.$$
\finenonce{000592}



\finexercice



\exercice{5252, rouget, 2010/07/04}
\enonce{005252}{***I}
 Déterminer un équivalent le plus simple possible de chacune des suites suivantes quand $n$ tend vers $+\infty$.

$$
\begin{array}{lllll}
1)\;\Arccos\frac{n-1}{n}&2)\;\Arccos\frac{1}{n}&3)\;\ch(\sqrt{n})&4)\;\left(1+\frac{1}{n}\right)^n&5)\;\frac{\Argch n}{\sqrt{n^4+n^2-1}}\\
6)\;(1+\sqrt{n})^{-\sqrt{n}}&7)\;\ln(\cos\frac{1}{n})(\ln\sin\frac{1}{n})&8)\;(\frac{\pi}{2})^{3/5}-(\Arctan n)^{3/5}
&9)\;\sqrt{1+\frac{(-1)^n}{\sqrt{n}}}-1
\end{array}
$$
\finenonce{005252}


\finexercice


\exercice{5253, rouget, 2010/07/04}
\enonce{005253}{***I}
Montrer que  $\sum_{k=0}^{n}k!\sim n!$.
\finenonce{005253}


\finexercice


\exercice{5254, rouget, 2010/07/04}
\enonce{005254}{***I}
\begin{enumerate}
\item  Soient $u$ et $v$ deux suites réelles strictement positives. Pour $n\in\Nn$, on pose $U_n=\sum_{k=0}^{n}u_k$ et $V_n=\sum_{k=0}^{n}v_k$. Montrer que si $u_n\sim v_n$ et si $\lim_{n\rightarrow +\infty}V_n=+\infty$, alors $U_n\sim V_n$.
\item  Application. Trouver un équivalent de $\sum_{k=1}^{n}\frac{1}{\sqrt{k}}$ et $\sum_{k=1}^{n}\ln(k)$.
\end{enumerate}
\finenonce{005254}


\finexercice


\exercice{5255, rouget, 2010/07/04}
\enonce{005255}{****}
Soit $(u_n)$ une suite réelle de limite nulle. Montrer que si $u_n+u_{2n}\sim\frac{3}{2n}$, alors $u_n\sim\frac{1}{n}$. A-t-on~:~si $u_n+u_{n+1}\sim\frac{2}{n}$, alors $u_n\sim\frac{1}{n}$~?
\finenonce{005255}


\finexercice


\exercice{5256, rouget, 2010/07/04}
\enonce{005256}{***I}
Soit $u$ la suite définie par $u_0=\frac{\pi}{2}$ et, $\forall n\in\Nn,\;u_{n+1}=\sin(u_n)$.
\begin{enumerate}
\item  Montrer que la suite $u$ est strictement positive, décroissante de limite nulle.
\item  On admet que si $u$ est une suite de limite nulle, alors, quand $n$ tend vers $+\infty$, $\sin(u_n)=u_n-\frac{u_n^3}{6}+o(u_n^3)$.

Déterminer un réel $\alpha$ tel que la suite $v_n=u_{n+1}^\alpha-u_n^\alpha$ ait une limite réelle non nulle. En appliquant le lemme de \textsc{Césaro} à la suite $(v_n)$, en déduire un équivalent simple de $u_n$ quand $n$ tend vers $+\infty$.
\end{enumerate}
\finenonce{005256}


\finexercice

\section{121.04 Suite récurrente linéaire}


\exercice{593, ridde, 1999/11/01}

\enonce{000593}{}
Que penser-vous de l'\'enonc\'e suivant : si $ (u_n) \sim (v_n)$ alors $ (e^{u_n})
\sim (e^{v_n})$. Donner un \'enonc\'e correct.
\finenonce{000593}



\finexercice



\exercice{594, ridde, 1999/11/01}

\enonce{000594}{}
\begin{enumerate}
\item Montrer que si $\forall n \in \Nn$ $u_n \neq 0$ et si $ (u_n) \rightarrow 0$
alors $ \ln (1 + u_n) \sim u_n$.
\item Soit $a$ un r\'eel. D\'eterminer la limite de $ (1 + \dfrac an)^n$.
\end{enumerate}
\finenonce{000594}



\finexercice



\exercice{595, ridde, 1999/11/01}

\enonce{000595}{}
Comparer les suites suivantes :
$$a_n = n^n, \qquad b_n = n^{\ln (n)}, \qquad c_n = e^{n^2}, \qquad d_n =
 (\ln n)^{n\ln n}$$
\finenonce{000595}



\finexercice



\exercice{596, gourio, 2001/09/01}

\enonce{000596}{}
Soient $(u_{n})_{n\in \Nn}$ et $(v_{n})_{n\in \Nn} $ deux suites r\'{e}elles de
limite $+\infty  $ telles que $u_{n}=o(v_{n}).$
Montrer qu'il existe une suite $(w_{n})_{n\in \Nn} $ de limite $+\infty  $
telle que $u_{n}=o(w_{n})$ et $w_{n}=o(v_{n}).$
\finenonce{000596}



\finexercice



\exercice{597, gourio, 2001/09/01}

\enonce{000597}{}
Donner un exemple de suites $(u_{n})_{n\in \Nn}$ et $(v_{n})_{n\in \Nn} $ telles
que $u_{n}=O(v_{n}) $ mais qu'on n'ait ni $u_{n}=o(v_{n}), $ ni   $v_{n}=O(u_{n}).$
\finenonce{000597}



\finexercice



\exercice{598, gourio, 2001/09/01}

\enonce{000598}{}
\'Etude de $(u_{n})_{n\in \Nn }$ d\'{e}finie par :
$$u_{0}\in [0,1],u_{n+1}=u_{n}^{2}.$$
Donner un \'{e}quivalent de $u_{n}$ quand $n\rightarrow \infty .$
\finenonce{000598}



\finexercice



\exercice{599, gourio, 2001/09/01}

\enonce{000599}{}
Montrer la r\'{e}ciproque du th\'{e}or\`{e}me de C\'{e}saro (i.e.
$\lim\limits_{n\rightarrow \infty }u_{n}=l$) :
\begin{enumerate}
\item  dans le cas o\`{u} $\lim\limits_{n\rightarrow \infty }v_{n}=l $ et
$$u_{n+1}-u_{n}=O(\frac{1}{n}),$$
\item dans le cas o\`{u} $(u_{n})_{n\in \Nn}$ est croissante.
\end{enumerate}
\finenonce{000599}



\finexercice



\exercice{600, gourio, 2001/09/01}

\enonce{000600}{}
\'Etudier la suite $(u_{n})_{n\in \Nn}$ d\'{e}finie par $u_{0}=1 $ et
$\forall n\in \Nn$ $u_{n+1}=u_{n}+\frac{2}{u_{n}}.$
En utilisant $v_{n}=\frac{u_{n}^{2}}{4},$ donner un \'{e}quivalent de $%
u_{n}.$
\emph{Indication} : on montrera que $\lim\limits_{n\rightarrow \infty
}v_{n+1}-v_{n}=1, $ on en d\'{e}duira un \'{e}quivalent de $v_{n} $ puis de
$u_{n}.$
\finenonce{000600}



\finexercice



\exercice{601, gourio, 2001/09/01}

\enonce{000601}{}
Soit $(u_{n})_{n\in \Nn}$ la suite d\'{e}finie par $u_{n+1}=u_{n}+u_{n}^{2}. $
L'\'{e}tudier et, en utilisant $v_{n}=\frac{1}{u_{n}},$ en donner un
\'{e}quivalent dans le cas $u_{0}\in ]-1;0]$. Que dire dans le cas $u_{0}\in
]0;\infty [$ ? (On \'{e}tudiera $v_{n}=\frac{\ln (u_{n})}{2^{n}}$.)
\finenonce{000601}



\finexercice



\exercice{602, gourio, 2001/09/01}

\enonce{000602}{}
Soient $f$ et $g$ deux formes lin\'{e}aires sur un espace vectoriel $E$
telles que $fg=0.$ Montrer que $f=0$ ou $g=0.$
\'Etudier la suite $(x_{n})_{n\in \Nn}$ d\'{e}finie par
$x_{0}=1,x_{n+1}=\frac{x_{n}}{1+nx_{n}^{2}}. $
En \'{e}tudiant $y_{n}=\frac{1}{x_{n+1}}-\frac{1}{x_{n}}$, en donner un \'{e}quivalent.
\finenonce{000602}



\finexercice



\exercice{603, gourio, 2001/09/01}

\enonce{000603}{}
\'Etudier la suite $(u_{n})_{n\in \Nn}$ d\'{e}finie par :
$$u_{0}\in ]0,\frac{\pi }{2}[,u_{n+1}=\sin u_{n}.$$
Donner $\lim_{x\rightarrow 0}\frac{1}{\sin ^{2}x}-\frac{1}{x^{2}}$
,(r\'{e}ponse: $\frac{1}{3}$) en d\'{e}duire un \'{e}quivalent de $u_{n}^{-2} $
donc de $u_{n}.$
\finenonce{000603}



\finexercice



\exercice{604, gourio, 2001/09/01}

\enonce{000604}{}
Montrer que $\forall n\in \Nn^{*},\exists !x_{n}\in [n,n+1[$ solution de $
x-E(x)=\frac{1}{x^{2}}.$
Donner un \'{e}quivalent de $x_{n}$ puis faire un d\'{e}veloppement
asymptotique de $x_{n}-n$ \`{a} l'ordre 5 en fonction de $\frac{1}{n}.$
\finenonce{000604}



\finexercice



\exercice{605, gourio, 2001/09/01}

\enonce{000605}{}
\'Etudier la convergence et calculer la limite \'{e}ventuelle de la suite
$(u_{n})_{n\in \Nn^{*}}$ d\'{e}finie par :
$$\forall n\in \Nn^{*},u_{n}=1+\frac{1}{2}+...+\frac{1}{n}-\frac{1}{n+1}
-...-\frac{1}{n^{2}}. $$
On montrera pr\'{e}alablement que :
$$1+\frac{1}{2}+...+\frac{1}{n}=\ln n+\gamma +o(1) $$
quand $n\rightarrow \infty .$
\finenonce{000605}



\finexercice



\exercice{1202, ridde, 1999/11/01}

\enonce{001202}{}
 Soit $ (u_n)$ d\'efinie par $u_0$ et $u_1$
strictement positifs et $u_{n + 1}
 = u_n + u_{n-1}$ pour $n \geq 1$.
\begin{enumerate}
\item Montrer que $\lim (\dfrac{u_{n + 1}}{u_n})$ existe et la d\'eterminer. Que
remarquez-vous ?
\item Soit $a_n = \dfrac{u_{n + 1}}{u_n}$. Exprimer $a_{n + 1}$ en fonction de $a_n$.
\item Montrer que $a_{2n}$ et $a_{2n + 1}$ sont adjacentes.
\item D\'eterminer un rationnel $r$ tel que $\left|r-\frac{1 + \sqrt 5}2\right|< 10^{-3}$.
\end{enumerate}

\finenonce{001202}



\finexercice



\exercice{1203, ridde, 1999/11/01}

\enonce{001203}{}
D\'eterminer $ (u_n)$ telle que
\begin{enumerate}
\item $u_0 = 1$, $u_1 = 3$, $u_{n + 2} = 4u_{n + 1}-4u_n$.
\item $u_0 = 1$, $u_1 = i$, $u_{n + 2} = 4u_{n + 1}-5u_n$.
\end{enumerate}
\finenonce{001203}



\finexercice



\exercice{1204, ridde, 1999/11/01}

\enonce{001204}{}
 D\'eterminer les suites born\'ees qui v\'erifient
$u_{n + 2} = 3u_{n + 1}-2u_n$.

\finenonce{001204}



\finexercice



\exercice{1205, ridde, 1999/11/01}

\enonce{001205}{}
D\'eterminer les suites convergentes qui v\'erifient $2u_{n + 2} = 7u_{n + 1}-3u_n$.
\finenonce{001205}



\finexercice



\exercice{1206, ridde, 1999/11/01}

\enonce{001206}{}
Montrer que la suite $u_0 = 1$, $u_1 = 2$ et $u_{n  +  2} = \sqrt{u_{n + 1}u_n}$
est bien d\'efinie et la d\'eterminer.
\finenonce{001206}



\finexercice



\exercice{1207, ridde, 1999/11/01}

\enonce{001207}{}
D\'eterminer les suites $ (u_n)$ et $ (v_n)$ qui v\'erifient $\begin{cases}u_0 = 2 \\
v_0 = -2 \end{cases}$ et
$\begin{cases} u_{n + 1} = u_n + v_n \\ v_{n + 1} = 3u_n-v_n \end{cases}$
\finenonce{001207}



\finexercice



\exercice{3068, quercia, 2010/03/08}

\enonce{003068}{Ensi Chimie P' 93}
\begin{enumerate}
  \item R{\'e}soudre $\begin{cases} u_{n+2} = \frac 12(u_{n+1} + u_n)\cr
                      u_0 = a,\ u_1 = b.\cr\end{cases}$
  \item Si $a=0$, trouver $\lim u_n$.
  \item R{\'e}soudre : $v_{n+2} = \sqrt{v_{n+1} v_n}$.
\end{enumerate}
\finenonce{003068}


\finexercice 


\exercice{3069, quercia, 2010/03/08}

\enonce{003069}{{\'E}quations de r{\'e}currence lin{\'e}aire}
\begin{enumerate}
  \item R{\'e}soudre : $$\begin{cases} u_{n+2} - u_n = n-1 \cr u_0 = u_1 = 0.\end{cases}$$
  \item R{\'e}soudre : $u_{n+2} + u_{n+1} + u_n = n$.
\end{enumerate}
\finenonce{003069}


\finexercice 


\exercice{3070, quercia, 2010/03/08}

\enonce{003070}{Syst{\`e}me r{\'e}current}
On donne $u_0,v_0$. R{\'e}soudre le syst{\`e}me :
$\begin{cases} 5u_n = 2u_{n-1} + 3v_{n-1} \cr 5v_n = 2v_{n-1} + 3u_{n-1}.\end{cases}$

\finenonce{003070}



\finexercice 


\exercice{3071, quercia, 2010/03/08}

\enonce{003071}{Caract{\'e}risation des suites polynomiales}
Soit $(u_n)$ une suite de r{\'e}els. On d{\'e}finit les suites d{\'e}riv{\'e}es de $(u_n)$ :
$$\begin{cases}(u'_n)  = (u_{n+1} - u_n)   \cr
         (u''_n) = (u'_{n+1} - u'_n) \cr
         \dots \cr 
         \left(u^{(k+1)}_n\right) = \left(u^{(k)}_{n+1} - u^{(k)}_n\right).\end{cases}$$
\begin{enumerate}
  \item Exprimer $u^{(k)}_n$ en fonction de $u_n, u_{n+1}, \dots, u_{n+k}$.
  \item Montrer que la suite $(u_n)$ est polynomiale si et seulement s'il existe $k \in \N$
    tel que $\left(u^{(k)}_n\right) = (0)$.
\end{enumerate}
\finenonce{003071}


\finexercice 


\exercice{3072, quercia, 2010/03/08}

\enonce{003072}{Nombre de nombres ne comportant pas 13}
Soit $T_n$ le nombre d'entiers naturels de $n$ chiffres exactement ne comportant
pas la s{\'e}quence 13 en nu\-m{\'e}\-ra\-tion d{\'e}cimale.
\begin{enumerate}
  \item   Montrer que $T_{n+2} = 10T_{n+1} - T_n$.
  \item   Calculer $T_n$ en fonction de $n$. 
\end{enumerate}
\finenonce{003072}


\finexercice 


\exercice{3073, quercia, 2010/03/08}

\enonce{003073}{$(\sqrt3+1)^{2n+1} - (\sqrt3-1)^{2n+1}$}
On note $x_n = (\sqrt3 + 1)^{2n+1}$, $y_n = (\sqrt3 - 1)^{2n+1}$, et
$z_n = [x_n]$.
\begin{enumerate}
  \item   Montrer que $z_n = x_n - y_n$.
  \item   En d{\'e}duire que $2^{n+1}$ divise $z_n$.
\end{enumerate}
\finenonce{003073}



\finexercice 


\exercice{5239, rouget, 2010/06/30}
\enonce{005239}{**T}
Déterminer $u_n$ en fonction de $n$ et de ses premiers termes dans chacun des cas suivants~:

\begin{enumerate}
\item  $\forall n\in\Nn,\;4u_{n+2}=4u_{n+1}+3u_n$.
\item  $\forall n\in\Nn,\;4u_{n+2}=u_n$.
\item  $\forall n\in\Nn,\;4u_{n+2}=4u_{n+1}+3u_n+12$.
\item  $\forall n\in\Nn,\;\frac{2}{u_{n+2}}=\frac{1}{u_{n+1}}-\frac{1}{u_n}$.
\item  $\forall n\geq2,\;u_n= 3u_{n-1}-2u_{n-2}+n^3$.
\item  $\forall n\in\Nn,\;u_{n+3}-6u_{n+2}+11u_{n+1}-6u_n=0$.
\item  $\forall n\in\Nn,\;u_{n+4}-2u_{n+3}+2u_{n+2}-2u_{n+1}+u_n=n^5$.
\end{enumerate}

\finenonce{005239}


\finexercice

\section{121.05 Suite de Cauchy}


\exercice{1208, vignal, 2001/09/01}

\enonce{001208}{}
Montrer que la suite $  \left(\frac{\sin n}{2^n}\right)_{n\in \Nn}$ est de Cauchy et que la suite $ \left((-1)^n+\frac{1}{n}\right)_{n\in \Nn}$
ne l'est pas.
\finenonce{001208}



\finexercice



\exercice{1209, cousquer, 2003/10/01}

\enonce{001209}{}
 Montrer que la suite définie par
$$u_n=1+\frac{\cos1}{1!} + \frac{\cos2}{2!}+\cdots+\frac{\cos n}{n!}$$
est une suite de Cauchy. En déduire sa convergence.
\finenonce{001209}



\finexercice



\exercice{1210, cousquer, 2003/10/01}

\enonce{001210}{}
 Montrer que toute sous-suite extraite d'une suite de Cauchy est
aussi une suite de Cauchy.

Montrer que si $(u_n)$ est une suite de Cauchy, on peut trouver une sous-suite
$(u_{n_k})$ de $(u_n)$ telle que~:
$$\forall p\in\mathbf{N},\;\forall q\geq p,\; \vert 
u_{n_p}-u_{n_q}\vert\leq \frac{1}{2^p}.$$
\finenonce{001210}



\finexercice



\exercice{1211, cousquer, 2003/10/01}

\enonce{001211}{}
 Une suite $(x_n)$ est définie par une relation de récurrence
$x_{n+1}=a\sin x_n+b$ où $a$ est un nombre réel de $\mathopen]0,1\mathclose[$ et
$b$ un nombre réel quelconque. Montrer que
pour tout $p\in\mathbf{N}$, $\vert x_{p+1}-x_p\vert\leq a^p\vert x_1-x_0\vert$.
En déduire que la suite $(x_n)$ est une suite de Cauchy.

Combien de termes faut-il calculer pour obtenir une valeur approchée
de $\lim x_n$ à $10^{-10}$ près si on suppose $a=1/2$, $b=5$, $x_0=1$~?
\finenonce{001211}



\finexercice


\section{121.06 Suite dans Rn}


\exercice{1901, maillot, 2001/09/01}

\enonce{001901}{}
Soit $x_n$ une suite de $\Rr^d$.
Montrer que l'ensemble $A$ des valeurs d'adh\'erence de $x_n$
est ferm\'e. Indication: prouver que le compl\'ement de $A$
est ouvert.
\finenonce{001901}



\finexercice



\exercice{1902, maillot, 2001/09/01}

\enonce{001902}{}
Soit $x_n$ une suite born\'ee de $\Rr^d$.
Montrer que $x_n$ converge si et seulement si
$A$ est un singleton. Indication: pour prouver
la convergence, utiliser qu'une suite born\'ee
de $\Rr^d$ a au moins une valeur d'adh\'erence.
\finenonce{001902}



\finexercice



\exercice{1903, maillot, 2001/09/01}

\enonce{001903}{}
Soit $f:\Rr^d\rightarrow\Rr^d$ continue.
Soit $x_0\in\Rr^d$. Soit $x_n$ la suite d\'efinie par
$$x_{n+1}=f(x_n).$$
Supposons que $||x_n-x_{n+1}||\rightarrow 0$.
Montrer que si $a\in\ A$ alors $f(a)=a$.

Indication: appliquer la d\'efinition de la continuit\'e
de $f$ en $a$ en termes de limites.
\finenonce{001903}



\finexercice



\exercice{1904, maillot, 2001/09/01}

\enonce{001904}{}
Soit $x_n$ une suite born\'ee de $\Rr^d$.
Supposons que $||x_n-x_{n+1}||\rightarrow 0$.
Montrer que l'ensemble $A$ est non-vide, compact, connexe.

Indication: pour la connexit\'e, supposer que $A=A_1\cup A_2$ avec
$A_1$ et $A_2$ non-vides, disjoints, ferm\'es.

Si $d=1$ conclure que $A=[a,b]$ avec
$a\leq b$.

\finenonce{001904}



\finexercice



\exercice{1905, maillot, 2001/09/01}

\enonce{001905}{}
Soit $f:\Rr\rightarrow\Rr$ continue.
Soit $x_0\in\Rr$. Soit $x_n$ la suite d\'efinie par
$$x_{n+1}=f(x_n).$$
Supposons que $x_n$ est born\'ee.
Montrer que $x_n$ converge
si et seulement si
$$||x_n-x_{n+1}||\rightarrow 0.$$

Indication. Montrer qu'il suffit de prouver que $a=b$ dans
$[a,b]=A$. Si $a<b$ montrer  que la suite est stationnaire.
\finenonce{001905}



\finexercice



\exercice{1906, maillot, 2001/09/01}

\enonce{001906}{}
Soit $s_n=\Sigma_{k=1}^n1/k$ et $x_n=\cos (s_n)$.
Montrer qu'il n'existe pas d'application
 $f:\Rr\rightarrow\Rr$ continue telle que
$$x_{n+1}=f(x_n).$$

Indication: montrer que  $||x_n-x_{n+1}||\rightarrow 0$ mais que
$x_n$ ne converge pas.
\finenonce{001906}



\finexercice


\section{121.99 Autre}


\exercice{5145, rouget, 2010/06/30}
\enonce{005145}{I}
\begin{enumerate}
\item (*) Calculer $\prod_{k=1}^{n}(1+\frac{1}{k})$, $n\in\Nn^*$.
\item (***) Calculer $\prod_{k=1}^{n}\cos\frac{a}{2^k}$, $a\in]0,\pi[$, $n\in\Nn^*$.
\end{enumerate}
\finenonce{005145}


\finexercice

\exercice{5148, rouget, 2010/06/30}
\enonce{005148}{***}
On veut montrer de manière élémentaire (c'est-à-dire en se passant du logarithme népérien et en ne travaillant qu'avec
les deux opérations $+$ et $\times$) que pour $n\in\Nn^*$, $(1+\frac{1}{n})^n<3$.

Pour cela développer, puis majorer $u_k=\frac{C_n^k}{n^k}$ en commençant par majorer $v_k=\frac{u_{k+1}}{u_k}$ par
$\frac{1}{2}$.
\finenonce{005148}


\finexercice

\exercice{5154, rouget, 2010/06/30}
\enonce{005154}{**I}
Soient $x$ un réel. Déterminer $\lim_{n\rightarrow +\infty}\frac{E(x)+E(2x)+...+E(nx)}{n^2}$.
\finenonce{005154}


\finexercice


\exercice{5220, rouget, 2010/06/30}
\enonce{005220}{***IT}
Soient $(u_n)_{n\in\Nn}$ une suite réelle et $(v_n)_{n\in\Nn}$ la suite définie par~:~$\forall n\in\Nn,\;v_n=\frac{u_0+u_1+...+u_n}{n+1}$.

\begin{enumerate}
\item  Montrer que si la suite $(u_n)_{n\in\Nn}$ vers un réel $\ell$, la suite $(v_n)_{n\in\Nn}$ converge et a pour limite $\ell$. Réciproque~?
\item  Montrer que si la suite $(u_n)_{n\in\Nn}$ est bornée, la suite $(v_n)_{n\in\Nn}$ est bornée. Réciproque~?
\item  Montrer que si la suite $(u_n)_{n\in\Nn}$ est croissante alors la suite $(v_n)_{n\in\Nn}$ l'est aussi.
\end{enumerate}
\finenonce{005220}


\finexercice


\exercice{5246, rouget, 2010/06/30}
\enonce{005246}{***I}
Soit $u_n$ l'unique racine positive de l'équation $x^n+x-1=0$. Etudier la suite $(u_n)$.
\finenonce{005246}


\finexercice


\exercice{5248, rouget, 2010/07/04}
\enonce{005248}{****}
Montrer que l'ensemble $E$ des réels de la forme $u_n=\sin(\ln(n))$, $n$ entier naturel non nul, est dense dans 
$[-1,1]$.
\finenonce{005248}


\finexercice


\exercice{5249, rouget, 2010/07/04}
\enonce{005249}{***}
Calculer $\mbox{inf}_{\alpha\in]0,\pi[}(\mbox{sup}_{n\in\Nn}(|\sin(n\alpha)|))$.
\finenonce{005249}


\finexercice

\section{122.01 Série à  termes positifs}


\exercice{1930, roussel, 2001/09/01}

\enonce{001930}{}
Soient, pour $n>0$,
$u_n=\displaystyle{\frac{n!e^n}{n^{n+\frac{1}{2}}}}$ et $v_n=\ln u_n$.
\begin{enumerate}
\item Etudier la serie de terme g\'en\'eral ~$w_n$ o\`u, pour $n \geq 2,~w_n = v_n-v_{n-1}$
et $w_1=v_1$.

\item En d\'eduire, en utilisant la convergence de la suite des sommes
partielles de $w_n$, que la suite $u_n$ converge vers $\lambda >0$.
\item D\'eterminer $\lambda$ en utilisant la formule de Wallis :
$\lim _{n \rightarrow + \infty} \displaystyle{\frac{2^{2n}(n!)^2}{\sqrt{n}(2n)!}}=
\sqrt{\pi}$. En d\'eduire
un \'equivalent de $n!$.

\emph{Indication} : Exprimer $n!$ (respectivement $(2n)!$) en fonction de
$u_n$ (resp. de $u_{2n}$) et remplacer-les dans la formule de Wallis.
\end{enumerate}
\finenonce{001930}



\finexercice



\exercice{1932, gineste, 2001/11/01}

\enonce{001932}{}
Etudier la série de terme général
$$ u_n=\frac{a^n2^{\sqrt{n}}}{2^{\sqrt{n}} + b^n } \mbox{ o\`u } a>0, \; b>0. $$
Indication : Chercher un équivalent suivant les valeurs de $b$.
\finenonce{001932}



\finexercice



\exercice{1934, gineste, 2001/11/01}

\enonce{001934}{Comparaison à des séries de Riemann et équivalent}
Etudier les séries de termes généraux
\begin{enumerate}
\item $$ u_n=\cos(\frac{ \pi n^2 }{2n^2+an+1}) \mbox{ avec } a>0 $$
\item $$ v_n=e^{-\sqrt{n}} $$
\item $$ w_n=(1 - \frac{1}{n^2})^n $$
\end{enumerate}
\finenonce{001934}



\finexercice



\exercice{1935, gineste, 2001/11/01}

\enonce{001935}{}
Soit $(u_n)$ une suite de réels strictement positifs, on suppose que $ \displaystyle \lim(\frac{u_{n+1}}{u_n})=1 $ et que $$ \frac{u_{n+1}}{u_n}=1 - \frac{\alpha}{n} + O(\frac{1}{n^{\beta}}) \mbox{ , o\`u } \alpha > 0 \ \  \beta > 1.$$
On pose $v_n=n^{\alpha}u_n$. Etudier  $ \displaystyle  \frac{v_{n+1}}{v_n} $ et montrer que $(v_n)$ a une limite finie. \
Application : Etudier la série de terme général $$u_n = \sqrt{n!} \sin 1 \sin \frac{1}{\sqrt{2}} \cdots \sin \frac{1}{\sqrt{n}} . $$
\finenonce{001935}



\finexercice



\exercice{1936, gineste, 2001/11/01}

\enonce{001936}{}
Déterminer la nature des séries de terme général:
\[
\begin{array}{clclcl}
1. & \displaystyle  \frac{n!}{n^n} &
2. & (\ch \sqrt{\ln n})^{-2} &
3. & n^{-(1+(1/n))}
\\
\\
4. & \displaystyle \frac{1}{\sqrt{n}} \ln \left(1+\frac{1}{\sqrt{n}} \right) &
5. &  \displaystyle \frac{\ln n}{\ln (e^n-1)} &
6. & n^{\ln n}e^{-\sqrt{n}}
\end{array}
\]
\finenonce{001936}



\finexercice



\exercice{1937, gineste, 2001/11/01}

\enonce{001937}{}
Étudier, suivant les valeurs de $p \in \Nn$, la nature de la série de terme général 
\[
u_n = \frac{1! + 2! + \cdots + n!}{(n+p)!} \cdotp
\]
\finenonce{001937}



\finexercice


\exercice{1938, gineste, 2001/11/01}
 
\enonce{001938}{}
Calculer les sommes des séries suivantes, en montrant leur convergence :
\begin{enumerate}
  \item  $\sum_{n \ge 0} (n+1)3^{-n}$
  \item  $\displaystyle \sum_{n \ge 0} \frac{n}{n^4+n^2+1}$
  \item  $\displaystyle \sum_{n \ge 3} \frac{2n-1}{n^3-4n}$
\end{enumerate}
\finenonce{001938}
 
 
 
\finexercice


\exercice{1939, gineste, 2001/11/01}

\enonce{001939}{}
Soit $(u_n)$ une suite réelle positive et $ \displaystyle S_n=\sum_{p=0}^{n}u_p.$ Comparer la nature des séries $(\sum u_n)$ et $\displaystyle (\sum \frac{u_n}{S_n}).$
\finenonce{001939}



\finexercice



\exercice{1941, gineste, 2001/11/01}

\enonce{001941}{Utilisation d'une série}
Le but de cet exercice est de montrer la convergence de l'intégrale généralisée suivante $ \displaystyle \int_0^{\infty} \frac{dx}{1 + x^4 \sin^2 x}.$

Pour cela, on considère la série de terme général $$u_n=\int_{n\pi}^{(n+1)\pi} \frac{dx}{1 + x^4 \sin^2 x} .$$
Par un changement de variable, transformer $u_n$ en
$$ u_n =\int_{0}^{\pi} \frac{dx}{1 + (n\pi + x)^4 \sin^2 x} $$
Encadrer ensuite $u_n$ par les termes de la suite $v_n$ o\`u
$$ v_n =\int_{0}^{\pi} \frac{dx}{1 + (n\pi)^4 \sin^2 x} $$
Calculer explicitement l'intégrale $v_n$ et en déduire un équivalent de $u_n$. Conclure.
\finenonce{001941}



\finexercice



\exercice{1942, gineste, 2001/11/01}

\enonce{001942}{}
Soit $u_n$ une suite décroissante à termes positifs. On suppose $(\sum u_n)$ converge. Montrer que $$ \lim_{n \rightarrow \infty} (nu_n)=0 .$$
Indication : Encadrer $\sum_{p+1}^n u_k$ pour $n>p$. Puis revenir aux définitions des limites avec les epsilons.
\finenonce{001942}



\finexercice



\exercice{1943, gineste, 2001/11/01}

\enonce{001943}{}
Soient $\sum_{n \ge 0} u_n$, $\sum_{n \ge 0} v_n$ deux séries à termes réels strictement positifs. On suppose que $\sum_{n \ge 0} v_n$ converge, et que
$$ \forall n \in \N, \; \frac{u_{n+2}}{u_n} \le \frac{v_{n+2}}{v_n} .$$
Montrer que $\sum_{n \ge 2} u_n$ converge.
\finenonce{001943}



\finexercice



\exercice{1945, gineste, 2001/11/01}

\enonce{001945}{Examen 2000}
\begin{enumerate}
\item On rappelle que la s\'erie harmonique altern\'ee converge et a pour somme $$\sum_{n=1}^{\infty}\frac{(-1)^n}{n}=-\log 2.$$
Montrer la convergence des deux s\'eries $\sum_{k=1}^{\infty}\left(\frac{1}{2k-1}-\frac{1}{2k}\right)$ et $\sum_{k=1}^{\infty}\left(\frac{1}{2k+1}-\frac{1}{2k}\right)$ et calculer leur somme \`a l'aide du rappel ci dessus.
\item
D\'ecomposer en \'el\'ements simples la fraction rationnelle
$\frac{1}{4x^3-x}$.
\item Montrer la convergence de la s\'erie $\sum_{k=1}^{\infty}\frac{1}{4k^3-k}$et calculer sa somme \`a l'aide de ce qui pr\'ec\`ede.
\item L'int\'egrale impropre $\int_1^{\infty}\frac{dx}{4x^3-x}$ converge t-elle? Si oui, la calculer.
\end{enumerate}
\finenonce{001945}



\finexercice



\exercice{1949, gineste, 2001/11/01}
 
\enonce{001949}{}
Soit $0 < a < b$ et $(u_n)_{n\geq 0}$ d\'efini par $u_0=1$ et
$\frac{u_{n+1}}{u_n}=\frac{n+a}{n+b}$ pour $n\geq 0.$ Montrer que
la limite de la suite
$W_n=\log (n^{b-a}u_n)$ existe et est finie. En d\'eduire les valeurs de
$a$ et $b$ telles que la s\'erie $\sum_{j=0}^{\infty} u_j$ converge.
Calculer alors sa somme: pour cela expliciter sa somme partielle $s_n,$
en montrant d'abord que pour tout $n$ on a
$$\sum_{j=0}^n[(j+1)+b-1]u_{j+1}=\sum_{j=0}^n[j +a]u_j.$$
\finenonce{001949}
 
 
\finexercice


\exercice{2718, matexo1, 2002/02/01}
\enonce{002718}{}
Par un calcul direct, montrer que les sommes partielles de la s\'erie
harmonique
$$S_n = \sum_{k=1}^n k^{-1}, \qquad n \ge 1$$
ne forment pas une suite de Cauchy. En d\'eduire que cette 
s\'erie diverge.
\finenonce{002718}
\finexercice


\exercice{2719, matexo1, 2002/02/01}
\enonce{002719}{}
En discutant \'eventuellement selon la valeur des param\`etres
r\'eels $\alpha $ et $\beta $, \'etudier
les s\'eries de termes g\'en\'eraux positifs ($n\geq 2$) :

$$\begin{array}{ll}
\displaystyle \frac{n+\alpha}{n+\beta}, & \displaystyle \frac 1{n (n^2-1)}, \\
\displaystyle \sqrt{n^4+2n+1} -\sqrt{n^4+\alpha n}, \quad \alpha \leq 2, &
\displaystyle \tan \left({\frac 1 n} \right) + \ln {n^2+\frac{\sqrt n}{n^2-n}}, \\
\displaystyle \left(\frac{2n+1}{3n+1}\right)^{\frac n 2}, &
\displaystyle \frac 1{\left(1+1/\sqrt n \right)^{n \sqrt n}}, \\
\displaystyle \frac{n^n \alpha^n}{n!}, \phantom{\int} & \displaystyle \sqrt[n]{n} -1, \\
\displaystyle n^{\alpha} (\ln n)^{\beta}, &
\displaystyle \int_n^{n+1/2} \frac 1{\sqrt{t^4+1}}dt, \\
\displaystyle \frac{1!+2!+\cdots+n!}{(n+k)!}, \quad k \in \Z, &
\displaystyle n^{\alpha} \left[ (n+1)^{(n+1) / n} - (n-1)^{(n-1)/ n}\right], \\
\displaystyle \int_1^\infty \!\! \exp(-x^n)\,dx \text{ (indication : changer de variable $t=x^n$)}.&
\end{array}$$
\finenonce{002719}
\finexercice


\exercice{2720, matexo1, 2002/02/01}
\enonce{002720}{}
Soit $(a_n)$ une suite de r\'eels strictement positifs tels que, au voisinage de $+\infty$,
on ait
$$ \frac{a_{n+1}}{a_n} = 1 + \frac{\alpha}{n} +o\left(\frac 1 n \right).$$
\begin{enumerate}
\item Montrer que la s\'erie de terme g\'en\'eral $n^\alpha$ est de ce type ; rappeler pour
quelles valeurs de $\alpha $ elle converge.

\item Montrer que si $\alpha >-1$, la s\'erie de terme g\'en\'eral $a_n$ diverge, et que si
$\alpha <-1$ elle converge.

\item Application : \'etudier la s\'erie 
$$\sum_{n\geq 1}{1.3.5\ldots(2n-1)\over2.4.6\ldots(2n+2)}.$$

\item Montrer que si l'on a au voisinage de $+\infty$,
$$ \frac{a_{n+1}}{a_n} = 1 - \frac 1 n +o \left(\frac 1{n \ln n} \right),$$
alors la s\'erie de terme g\'en\'eral $a_n$ diverge.

\item Application : \'etudier la s\'erie 
$$\sum_{n\geq 1} \left(1-\exp (-1/n) \right).$$
\end{enumerate}
\finenonce{002720}
\finexercice


\exercice{2721, matexo1, 2002/02/01}
\enonce{002721}{}
Soit $(u_n)$ la suite d\'efinie par $u_0 \in ]0,1[$ donn\'e et $u_{n+1}=u_n-{u_n}^2$.
Montrer que cette suite converge et en donner la limite. 
Montrer que la s\'erie de terme g\'en\'eral ${u_n}^2$ converge et en donner la limite.
Montrer que les s\'eries de terme g\'en\'eraux $u_n$ et $\ln(u_{n+1}/u_n)$ divergent.
\finenonce{002721}
\finexercice


\exercice{2722, matexo1, 2002/02/01}
\enonce{002722}{}
Montrer qu'il existe deux r\'eels $\alpha$, $\beta$, tels que pour tout $n\in\N^*$, 
$$ \int_0^\pi  (\alpha t+\beta t^2) \cos(nt) \,dt = \frac 1{n^2}.$$
En d\'eduire la valeur de
$$S = \sum_{n=1}^{+ \infty} \frac 1{n^2}.$$
\finenonce{002722}


\finexercice


\exercice{2723, matexo1, 2002/02/01}
\enonce{002723}{}
Soit $(a_n)$ un suite de r\'eels strictement positifs telle que $\sum a_n$ converge.
\'Etudier les s\'eries
$$\sum a_n^2, \qquad \sum \frac{a_n}{1+a_n}, \qquad \sum a_n a_{2n}, \qquad \sum_{n\geq 1}
\frac{\sqrt{a_n}}n.$$
\finenonce{002723}
\finexercice


\exercice{2725, matexo1, 2002/02/01}
\enonce{002725}{}
Justifier la convergence et calculer les sommes des s\'eries suivantes 
$$\begin{array}{ll}
\displaystyle \sum_{n\geq 1} \frac 1{n(n+k)} \quad(k\in\N^*), &
\displaystyle\sum \frac 9{(3n+1)(3n+4)}, \\
\displaystyle\sum \frac{n^2+n-3}{n!}, &
\displaystyle\sum_{n\geq 2} \ln \frac{n^2}{n^2-1}.
\end{array}$$
\finenonce{002725}
\finexercice


\exercice{5108, rouget, 2010/06/30}
\enonce{005108}{**T}
Montrer par récurrence que, pour tout $n\in\Nn^*$,
$\sum_{k=1}^{n}\frac{1}{k(k+1)(k+2)}=\frac{n(n+3)}{4(n+1)(n+2)}$. Trouver une démonstration directe.
\finenonce{005108}


\finexercice


\exercice{5109, rouget, 2010/06/30}
\enonce{005109}{***I}
\begin{enumerate}
\item  Montrer par récurrence que, pour tout naturel non nul $n$, $\sum_{k=1}^{n}k=\frac{n(n+1)}{2}$. En calculant 
la différence $(k+1)^2-k^2$, trouver une démonstration directe de ce résultat.\rule[-1mm]{0mm}{0mm}
\item  Calculer de même les sommes $\sum_{k=1}^{n}k^2$, $\sum_{k=1}^{n}k^3$ et $\sum_{k=1}^{n}k^4$ (et 
mémoriser les résultats).
\item  On pose $S_p=\sum_{k=1}^{n}k^p$. Déterminer une relation de récurrence permettant de calculer les $S_p$ 
de proche en proche.
\end{enumerate}
\finenonce{005109}


\finexercice


\exercice{5143, rouget, 2010/06/30}
\enonce{005143}{Sommes télescopiques}
\label{exo:suprou7}
Calculer les sommes suivantes~:
\begin{enumerate}
\item (**) $\sum_{k=1}^{n}\frac{1}{k(k+1)}$ et $\sum_{k=1}^{n}\frac{1}{k(k+1)(k+2)}$
\item (***) Calculer $S_p=\sum_{k=1}^{n}k^p$ pour $n\in\Nn^*$ et $p\in\{1,2,3,4\}$ (dans chaque cas, chercher un
polynôme $P_p$ de degré $p+1$ tel que $P_p(x+1)-P_p(x)=x^p$).
\item (**) Calculer $\sum_{k=1}^{n}\Arctan\frac{1}{k^2+k+1}$ (aller relire certaines formules établies dans une planche précédente).
\item (**) Calculer $\sum_{k=1}^{n}\Arctan\frac{2}{k^2}$.
\end{enumerate}
\finenonce{005143}


\finexercice


\exercice{5144, rouget, 2010/06/30}
\enonce{005144}{I}
Calculer les sommes suivantes~:
\begin{enumerate}
\item (**) $\sum_{1\leq i<j\leq n}^{}1$.
\item (**) $\sum_{1\leq i,j\leq n}^{}j$ et $\sum_{1\leq i<j\leq n}^{}j$.
\item (*) $\sum_{1\leq i,j\leq n}^{}ij$.
\item (***) Pour $n\in\Nn^*$, on pose $u_n=\frac{1}{n^5}\sum_{k=1}^{n}\sum_{h=1}^{n}(5h^4-18h^2k^2+5k^4)$.
Déterminer$\lim_{n\rightarrow +\infty}u_n$ (utiliser les résultats de l'exercice \ref{exo:suprou7}, 2)).
\end{enumerate}
\finenonce{005144}


\finexercice


\exercice{5149, rouget, 2010/06/30}
\enonce{005149}{***I}
\label{exo:suprou4bis}
Soient $n\in\Nn^*$ et $a_1$, $a_2$,..., $a_n$, $n$ réels strictement positifs.

Montrer que $(a_1+a_2+ ... +a_n)(\frac{1}{a_1}+...+\frac{1}{a_n})\geq n^2$ (développer et penser à
$f(x)=x+\frac{1}{x}$).
\finenonce{005149}


\finexercice

\exercice{5223, rouget, 2010/06/30}
\enonce{005223}{**}
Soit $(u_n)_{n\in\Nn}$ une suite arithmétique ne s'annulant pas. Montrer que pour tout entier naturel $n$, on a $\sum_{k=0}^{n}\frac{1}{u_ku_{k+1}}=\frac{n+1}{u_0u_{n+1}}$.
\finenonce{005223}


\finexercice


\exercice{5224, rouget, 2010/06/30}
\enonce{005224}{**}
Calculer $\lim_{n\rightarrow +\infty}\sum_{k=1}^{n}\frac{1}{1^2+2^2+...+k^2}$.
\finenonce{005224}


\finexercice

\exercice{5697, rouget, 2010/10/16}
\enonce{005697}{**}
Soit $(u_n)_{n\in\Nn}$ une suite positive telle que la série de terme général $u_n$ converge. Etudier la nature de la série de terme général $\frac{\sqrt{u_n}}{n}$.
\finenonce{005697}


\finexercice


\exercice{5698, rouget, 2010/10/16}
\enonce{005698}{***}
Soit $(u_n)_{n\in\Nn}$ une suite de réels positifs. Trouver la nature de la série de terme général $v_n =\frac{u_n}{(1+u_1)\ldots(1+u_n)}$, $n\geqslant1$,  connaissant la nature de la série de terme général $u_n$ puis en calculer la somme en cas de convergence.
\finenonce{005698}


\finexercice


\exercice{5699, rouget, 2010/10/16}
\enonce{005699}{****}
Soit  $(u_n)_{n\in\Nn}$ une suite de réels strictement positifs telle que la série de terme général $u_n$ diverge.

Pour $n\in\Nn$, on pose $S_n = u_0+...+u_n$. Etudier en fonction de $\alpha> 0$ la nature de la série de terme général $\frac{u_n}{(S_n)^\alpha}$.

\finenonce{005699}


\finexercice


\exercice{5704, rouget, 2010/10/16}
\enonce{005704}{*}
Nature de la série de terme général $u_n=\sum_{k=1}^{n}\frac{1}{(n+k)^p}$, $p\in]0,+\infty[$.
\finenonce{005704}


\finexercice


\exercice{5705, rouget, 2010/10/16}
\enonce{005705}{**}
Déterminer un équivalent simple de $\frac{n!}{(a+1)(a+2)\ldots(a+n)}$ quand $n$ tend vers l'infini ($a$ réel positif donné).
\finenonce{005705}


\finexercice

\section{122.02 Convergence absolue}


\exercice{1933, gineste, 2001/11/01}

\enonce{001933}{Utilisation des règles de Cauchy et d'Alembert}
Etudier les séries de termes généraux
\begin{enumerate}
\item $$ u_n = \sqrt{n!} \sin x \sin \frac{x}{\sqrt{2}} \cdots \sin \frac{x}{\sqrt{n}} \mbox{ avec } x>0 .$$
\item $$ v_n=e^{an^2}(1 - \frac{a}{n})^{n^3} $$
\end{enumerate}
\finenonce{001933}



\finexercice



\exercice{1940, gineste, 2001/11/01}

\enonce{001940}{Séries à termes quelconques}
Etudier les séries de termes généraux
\begin{enumerate}
\item $$ u_n=\frac{(-1)^n}{(\ln n) (n^{1/n})} $$
\item $$ v_n=\frac{(-1)^n}{\sqrt{n^\alpha + (-1)^n}} \mbox{ o\`u } \alpha > 0 $$
\item $$ w_n=\ln(1 + \frac{(-1)^n}{n^\alpha}) \mbox{ o\`u } \alpha > 0 $$
\end{enumerate}
Indication : Des calculs de D.L. peuvent etre fructueux ...
\finenonce{001940}



\finexercice



\exercice{1944, gineste, 2001/11/01}

\enonce{001944}{Examen 2000}
 En justifiant votre r\'eponse, classer les dix s\'eries $\sum u_n$
suivantes en 4 cat\'egories
\begin{itemize}
\item GD: celles telles que $u_n$ ne tend pas vers 0;
\item ZD: celles qui divergent et telles que $\lim u_n=0;$
\item AC: celles qui convergent absolument;
\item SC: celles qui convergent, mais non absolument.
\end{itemize}
(Attention: pour pouvoir r\'epondre, certaines s\'eries demandent deux
d\'emonstrations: par exemple pour montrer que $\sum u_n$ est SC, il
faut montrer que $\sum u_n$ converge \emph{et} que $\sum|u_n|$ diverge.
$$\sum_{n=1}^{\infty}\left(\frac{(-1)^n}{n}+\frac{1}{n^2}\right);\
\sum_{n=1}^{\infty}\left(\sqrt{n+1}-\sqrt{n}\right);\
\sum_{n=1}^{\infty}\frac{1}{\sqrt{n}}\left(\sqrt{n+1}-\sqrt{n}\right)^2;$$

$$
\sum_{n=1}^{\infty}\left[\frac{1}{n}-\log(1+\frac{1}{n})\right];\
\sum_{n=1}^{\infty}\frac{n!}{n^n};
\sum_{n=1}^{\infty}\left(1-(1-\frac{1}{n})^n\right);
$$

$$
\sum_{n=1}^{\infty}\frac{2^n+1000}{3^n+1};\
\sum_{n=1}^{\infty}(1-\cos\frac{\pi}{n});\
\sum_{n=1}^{\infty}\sin(\pi n)\sin(\frac{\pi}{n});\
\sum_{n=0}^{\infty}\left(\sum_{k=0}^{n}\frac{1}{2^k}\frac{1}{3^{n-k}}\right).
$$
\finenonce{001944}



\finexercice



\exercice{2724, matexo1, 2002/02/01}
\enonce{002724}{}
D\'eterminer, en fonction des param\`etres r\'eels $\alpha$, $\beta$, la nature
des s\'eries de termes g\'en\'eraux ($n\geq 2$)
$$\begin{array}{ll}
\displaystyle (-1)^n n^\alpha, \phantom{\int}  & 
\displaystyle n^\beta \left(1-(-1)^n n^\alpha\right), \\  
\displaystyle \frac{(-1)^n}{n-\ln n}, &  
\displaystyle \exp \left( \frac{-1}{\sqrt n} -1 \right), \\ 
\displaystyle \ln \left(1- \frac{(-1)^n}n \right), &
\displaystyle \frac{(-1)^n}{\sqrt n+(-1)^n}, \\
\displaystyle \sin \left( 2\pi  \frac{n!}e \right) &
\!\!\!\!\!\!\!\!\! (\mbox{on pourra utiliser que\,: }
\displaystyle 1/e = \sum_{k=0}^{+\infty} \frac{(-1)^k}{k!}).
\end{array}$$
\finenonce{002724}
\finexercice

\section{122.03 Séries semi-convergentes}


\exercice{1944, gineste, 2001/11/01}

\enonce{001944}{Examen 2000}
 En justifiant votre r\'eponse, classer les dix s\'eries $\sum u_n$
suivantes en 4 cat\'egories
\begin{itemize}
\item GD: celles telles que $u_n$ ne tend pas vers 0;
\item ZD: celles qui divergent et telles que $\lim u_n=0;$
\item AC: celles qui convergent absolument;
\item SC: celles qui convergent, mais non absolument.
\end{itemize}
(Attention: pour pouvoir r\'epondre, certaines s\'eries demandent deux
d\'emonstrations: par exemple pour montrer que $\sum u_n$ est SC, il
faut montrer que $\sum u_n$ converge \emph{et} que $\sum|u_n|$ diverge.
$$\sum_{n=1}^{\infty}\left(\frac{(-1)^n}{n}+\frac{1}{n^2}\right);\
\sum_{n=1}^{\infty}\left(\sqrt{n+1}-\sqrt{n}\right);\
\sum_{n=1}^{\infty}\frac{1}{\sqrt{n}}\left(\sqrt{n+1}-\sqrt{n}\right)^2;$$

$$
\sum_{n=1}^{\infty}\left[\frac{1}{n}-\log(1+\frac{1}{n})\right];\
\sum_{n=1}^{\infty}\frac{n!}{n^n};
\sum_{n=1}^{\infty}\left(1-(1-\frac{1}{n})^n\right);
$$

$$
\sum_{n=1}^{\infty}\frac{2^n+1000}{3^n+1};\
\sum_{n=1}^{\infty}(1-\cos\frac{\pi}{n});\
\sum_{n=1}^{\infty}\sin(\pi n)\sin(\frac{\pi}{n});\
\sum_{n=0}^{\infty}\left(\sum_{k=0}^{n}\frac{1}{2^k}\frac{1}{3^{n-k}}\right).
$$
\finenonce{001944}



\finexercice


\section{122.04 Séries alternées}


\exercice{1931, roussel, 2001/09/01}

\enonce{001931}{}
Soit $S= \displaystyle{\sum_{n=1}^{\infty}} \displaystyle{\frac{(-1)^{n-1}}{n^3}}$. Donner une
valeur approch\'ee de $S$ en garantissant une erreur inf\'erieure ou \'egale \`a $10^{-3}$.
\finenonce{001931}



\finexercice



\exercice{1945, gineste, 2001/11/01}

\enonce{001945}{Examen 2000}
\begin{enumerate}
\item On rappelle que la s\'erie harmonique altern\'ee converge et a pour somme $$\sum_{n=1}^{\infty}\frac{(-1)^n}{n}=-\log 2.$$
Montrer la convergence des deux s\'eries $\sum_{k=1}^{\infty}\left(\frac{1}{2k-1}-\frac{1}{2k}\right)$ et $\sum_{k=1}^{\infty}\left(\frac{1}{2k+1}-\frac{1}{2k}\right)$ et calculer leur somme \`a l'aide du rappel ci dessus.
\item
D\'ecomposer en \'el\'ements simples la fraction rationnelle
$\frac{1}{4x^3-x}$.
\item Montrer la convergence de la s\'erie $\sum_{k=1}^{\infty}\frac{1}{4k^3-k}$et calculer sa somme \`a l'aide de ce qui pr\'ec\`ede.
\item L'int\'egrale impropre $\int_1^{\infty}\frac{dx}{4x^3-x}$ converge t-elle? Si oui, la calculer.
\end{enumerate}
\finenonce{001945}



\finexercice



\exercice{1948, gineste, 2001/11/01}

\enonce{001948}{Permutation dans la série harmonique alternée : Pringsheim (1883)}
Pour tout
entier $n>0$, soit $u(n)=(-1)^n/n$ .
Soit $\sigma$ une permutation des entiers
$>0$ et soit $\tau$ la permutation réciproque. On suppose de plus que

(1) pour tout entier $p>0$ on a $\tau(2p-1)<\tau(2p+1)$ et
$\tau(2p)<\tau(2p+2).$

(2) Notant par $p(n)$ le nombre d'entiers $k$ tels que $1\leq k\leq n$ et
$\sigma(k)$ est pair, alors $\alpha=\lim _{n\infty}p(n)/n$ existe et est dans
$]0,1[.$

\begin{enumerate}
\item
Dans le cas particulier o\`u $\sigma$ est définie par
$$\sigma(3p)=2p,\ \sigma(3p+1)=4p+1,\ \sigma(3p+2)=4p+3$$
pour tout entier $p>0,$ calculer explicitement $\tau$, et vérifier que $\sigma$ satisfait
(1) et (2), en calculant $p(n)$ pour tout $n$ ainsi que $\alpha.$

\item
On note $f(n)=\sum_{k=1}^n1/k-\log n,$ et on rappelle le fait,
vu en cours, que  $\lim _{n\infty}f(n)=\gamma$ existe (Constante d'Euler).
On revient au cas général pour $\sigma$, on considère la série de terme
général $v_n=u(\sigma(n))$ et on note $s_n=v_1+\cdots+v_n$.

\item
Montrer par récurrence  que
$s_n=\sum_{k=1}^{p(n)}\frac{1}{2k}-\sum_{k=1}^{n-p(n)}\frac{1}{2k-1}$ et que
$$s_n=\frac{1}{2}f(p(n))+\frac{1}{2}f(n-p(n))-f(2n-2p(n))
+\frac{1}{2}\log\frac{p(n)}{n-p(n)}-\log 2.$$
En déduire que $\sum _{n=1}^{\infty}v_n$ converge
et calculer sa somme en fonction de $\alpha.$
\end{enumerate}
\finenonce{001948}



\finexercice



\exercice{5710, rouget, 2010/10/16}
\enonce{005710}{**}
 Calculer $\sum_{n=0}^{+\infty}\frac{(-1)^n}{3n+1}$.
\finenonce{005710}


\finexercice

\section{122.05 Familles sommables}


\exercice{4487, quercia, 2010/03/14}
\enonce{004487}{Dénombrabilité}
$A$ étant un ensemble infini dénombrable, les ensembles suivants sont-ils
dénombrables~:

\begin{enumerate}
  \item ${\cal P}(A)$~?
  \item $\{\text{parties finies de }A\}$~?
  \item $\{\text{suites périodiques à valeurs dans }A\}$~?
  \item $\{\text{suites ultimement périodiques à valeurs dans }A\}$~?
  \item $\{\text{relations d'ordre total sur } A\}$~?
\end{enumerate}
\finenonce{004487}



\finexercice


\exercice{4488, quercia, 2010/03/14}
\enonce{004488}{Discontinuités d'une fonction monotone}

\begin{enumerate}
  \item Soit $f:\R \to \R$ croissante. Montrer que l'ensemble des points
de discontinuité de~$f$ est dénombrable (pour $[a,b]\subset \R$,
considérer la famille $(f(x^+)-f(x^-))_{x\in{[a,b]}}$).

  \item Donner un exemple de fonction $f:\R \to \R$ croissante ayant
une infinité dénombrable de discontinuités.

  \item $(**)$ Trouver une fonction $f:\R \to \R$ {\it strictement\/} croissante dont
l'ensemble des points de discontinuité est égal à~$\Q$.


\end{enumerate}
\finenonce{004488}



\finexercice


\exercice{4489, quercia, 2010/03/14}
\enonce{004489}{Ensemble non vide~?}
Soit $(r_n)_{n\ge 1}$ une énumération des rationnels.
On note $I_n = {\Bigl]r_n-\frac1{n^2}, r_n+\frac1{n^2}\Bigr[}$, $E = \bigcup_{n=1}^\infty I_n$ et
$F = \R\setminus E$. Montrer que $F\ne \varnothing$ (ceci est choquant vu que les
éléments de~$F$ sont, par définition, "loin" de chaque rationnel, pourtant c'est vrai) .

\finenonce{004489}



\finexercice


\exercice{4490, quercia, 2010/03/14}
\enonce{004490}{\'Etude de convergence}

\'Etudier la finitude des sommes suivantes~:

\begin{enumerate}
  \item $\sum_{(i,j)\in(\N^*)^2}\frac1{(i+j)^\alpha}$.
    

  \item $\sum_{(i,j)\in(\N^*)^2}\frac1{\strut i^\alpha +j^\alpha}$.
    
  \item $\sum_{x\in\Q\cap[1,+\infty[}\frac1{x^2}$.
    
  \item $\sum_{(p,q)\in\N^2}\frac1{a^p+b^q}$, ${a>1,b>1}$.
    
\end{enumerate}
\finenonce{004490}



\finexercice


\exercice{4491, quercia, 2010/03/14}
\enonce{004491}{Série des restes}

Calculer $\sum_{n=0}^{+\infty}\sum_{k=n}^{+\infty}\frac1{k!}$.
\finenonce{004491}



\finexercice


\exercice{4492, quercia, 2010/03/14}
\enonce{004492}{Série des restes}

Calculer $\sum_{p=1}^{+\infty}\sum_{q=p}^{+\infty}\frac{(-1)^p}{q^3}$ en fonction de 
$\zeta(3)$.

\finenonce{004492}



\finexercice


\exercice{4493, quercia, 2010/03/14}
\enonce{004493}{Non interversion des sommations}

On pose $a_{n,p} = \frac1{\strut n^2-p^2}$ si $n\ne p$ et $a_{n,n} = 0$.

\begin{enumerate}
  \item Expliquer simplement pourquoi la suite double $(a_{n,p})_{(n,p)\in\N^2}$
    n'est pas sommable.
    
  \item Calculer $\sum_{n=0}^\infty\sum_{p=0}^\infty a_{n,p}$ et
    $\sum_{p=0}^\infty\sum_{n=0}^\infty a_{n,p}$.
    

\end{enumerate}
\finenonce{004493}



\finexercice


\exercice{4494, quercia, 2010/03/14}
\enonce{004494}{Identité remarquable}

Montrer que pour $x\in\C$, $|x|<1$, on a l'égalité :
$\sum_{n=0}^{+\infty}\frac{x^{2n+1}}{1-x^{2n+1}}=
\sum_{n=1}^{+\infty}\frac{x^{n}}{1-x^{2n}}$.

\finenonce{004494}



\finexercice


\exercice{4495, quercia, 2010/03/14}
\enonce{004495}{Calcul de somme}

Soit $z\in\C$ tel que $|z| < 1$.
Montrer que $\sum_{n=1}^\infty \frac{z^n}{\strut 1-z^n} = \sum_{n=1}^\infty d(n)z^n$
où $d(n)$ est le nombre de diviseurs positifs de $n$.
\finenonce{004495}



\finexercice


\exercice{4496, quercia, 2010/03/14}
\enonce{004496}{Centrale MP 2000}
Soit $S(t) = \sum_{n=1}^\infty \frac{t^n}{1+t^n}$.

\begin{enumerate}
  \item Pour quelles valeurs de~$t$ $S(t)$ a-t-elle un sens~?
    
  \item Montrer que $S(t) = \sum_{k=1}^\infty (-1)^{k-1}\frac{t^k}{1-t^k}$.
    
  \item Soit $F_m(t) = \sum_{k=1}^m (-1)^{k-1}\frac{t^k(1-t)}{1-t^k}$.
    Montrer que $(F_m(t))$ converge uniformément vers $(1-t)S(t)$ sur $[0,1]$.
    En déduire la limite en~$1$ de $(1-t)S(t)$.
    On rappelle que $\ln2 = \sum_{m=1}^\infty \frac{(-1)^{m-1}}m$.
    
  \item Calculer le développement en série entière de~$S(t)$.
    Donner une interprétation arithmétique des coefficients de ce développement
    et préciser leur signe en fonction de~$n$.
    
\end{enumerate}
\finenonce{004496}



\finexercice


\exercice{4497, quercia, 2010/03/14}
\enonce{004497}{Centrale MP 2002}

Soient $a,b,c\in\N^*$. On pose $f(z) = \sum_{n=0}^\infty \frac{z^{bn}}{1-z^{an+c}}$.

\begin{enumerate}
  \item Étudier la convergence de la série et montrer qu'on peut intervertir
    $b$ et $c$ dans la formule.


  \item Développer en série entière~: $\sum_{m=1}^\infty\frac{z^{2m}}{1-z^m}$.

\end{enumerate}
\finenonce{004497}



\finexercice


\exercice{4498, quercia, 2010/03/14}
\enonce{004498}{Calcul de sommes}

Calculer les sommes suivantes : $A = \sum_{(p,q)\in(\N^*)^2}\frac1{p^2q^2}$,
$B = \sum_{(p,q)\in(\N^*)^2 ; p|q}\frac1{p^2q^2}$ et
$C = \sum_{(p,q)\in(\N^*)^2 ; p\wedge q=1}\frac1{p^2q^2}$.

\finenonce{004498}



\finexercice


\exercice{4499, quercia, 2010/03/14}
\enonce{004499}{Série harmonique alternée}

On réordonne les termes de la série harmonique alternée en prenant tour
à tour $p$ termes positifs puis $q$ termes négatifs, $p,q\ge 1$.
Calculer la somme de la série correspondante.


\finenonce{004499}



\finexercice


\exercice{4500, quercia, 2010/03/14}
\enonce{004500}{Familles de carrés sommable}

\begin{enumerate}
  \item Soit $P\in\R[X]$. Vérifier que~:
    $ \int_{t=-1}^1 P(t)\,d t + i \int_{\theta=0}^\pi P(e^{i\theta})e^{i\theta}\, d\theta=0$.

    En déduire~: $ \int_{t=0}^1P^2(t)\,d t \le \frac12 \int_{\theta=-\pi}^\pi |P(e^{i\theta})|^2\, d\theta$.

  \item Soient $2n$ réels positifs $a_1,\dots,a_n,b_1,\dots,b_n$.
    Montrer que $\sum_{k=1}^n\sum_{\ell=1}^n \frac{a_kb_\ell}{\strut k+\ell}
    \le \pi\sqrt{\sum_{k=1}^n a_k^2}\,\sqrt{\sum_{\ell=1}^n b_\ell^2}$.

  \item Soient $(a_k)_{k\in\N}$ et $(b_\ell)_{\ell\in\N}$ deux suites complexes
    de carrés sommables.

    Montrer que la suite double
    $\left(\frac{a_kb_\ell}{\strut k+\ell}\right)_{(k,\ell)\in\N^2}$ est sommable.
\end{enumerate}
\finenonce{004500}



\finexercice


\exercice{4501, quercia, 2010/03/14}
\enonce{004501}{Associativité générale}

Soit $(a_i)_{i\in I}$ une famille sommable et $(I_n)_{n\in\N}$ une
suite croissante de parties de~$I$, non nécéssairement finies, telle que $\bigcup_{n\in\N}I_n = I$.
Montrer que $\sum_{i\in I_n} a_i \to \sum_{i\in I} a_i$ lorsque $n\to\infty$.
En déduire que si $(J_n)_{n\in\N}$ est une partition dénombrable de~$I$ alors
$\sum_{i\in I}a_i = \sum_{n=0}^\infty\sum_{i\in J_n}a_i$.
\finenonce{004501}



\finexercice


\exercice{4502, quercia, 2010/03/14}
\enonce{004502}{Mines MP 2001}

Déterminer l'ensemble de définition de $f(x) = \sum_{k=2}^\infty\frac{(-1)^k}{x+k}$.
Montrer que $f$ est de classe $\mathcal{C}^\infty$ sur son domaine et la développer
en série entière.


\finenonce{004502}



\finexercice

\section{122.06 Fonction exponentielle complexe}


\exercice{4403, quercia, 2010/03/12}
\enonce{004403}{$\cos z$}
Quels sont les complexes $z$ tels que $\cos z \in {[-1,1]}$ ?
\finenonce{004403}



\finexercice


\exercice{4404, quercia, 2010/03/12}
\enonce{004404}{$\lim((1+z/n)^n)$}
Soit $z\in\C$. Montrer que $\left(1+\frac zn\right)^n \to e^z$ lorsque $n\to\infty$.


\finenonce{004404}



\finexercice


\exercice{4405, quercia, 2010/03/12}
\enonce{004405}{Inégalité}
Soit $z \in \C$. Montrer que $|e^z-1| \le e^{|z|}-1 \le |z|e^{|z|}$.


\finenonce{004405}



\finexercice


\exercice{4406, quercia, 2010/03/12}
\enonce{004406}{Inégalité, Polytechnique MP$^*$ 2006}
Soit $z=x+iy\in\C$ avec $x,y\in\R$ et $x\ne 0$.
Montrer que $\Bigl|\frac{e^z-1}z\Bigr|\le\Bigl|\frac{e^x-1}x\Bigr|$.
Que dire en cas d'égalité~?


\finenonce{004406}



\finexercice


\exercice{4407, quercia, 2010/03/12}
\enonce{004407}{Morphismes $(\R,+) \to (\C,*)$}
Soit $f : \R \to {\C^*}$ telle que :
$\forall\ x,y\in\R,\ f(x+y) = f(x)f(y)$.

\begin{enumerate}
  \item Si $f$ est dérivable, montrer qu'il existe $\lambda \in \C$ tel que :
    $\forall\ x \in \R,\ f(x) = e^{\lambda x}$.

  \item Obtenir le même résultat si $f$ est seulement supposée continue
    (prendre une primitive, $F$, de $f$ et montrer qu'elle est de classe $\C^2$).
\end{enumerate}
\finenonce{004407}



\finexercice


\exercice{4408, quercia, 2010/03/12}
\enonce{004408}{$e^z = z$}
Montrer qu'il existe une infinité de complexes $z$ tels que $e^z = z$
(on calculera $x$ en fonction de $y$, et on étudiera l'équation obtenue).



\finenonce{004408}



\finexercice


\exercice{4409, quercia, 2010/03/12}
\enonce{004409}{\'Equations trigonométriques}
Résoudre dans $\C$ :

\begin{enumerate}
  \item $\cos z = 2$.
    

  \item $\ch z = -1$.
    

  \item $\sin z + \sin jz + \sin j^2z = 0$.
    

  \item $8\cos z + 4i\sin z = 7 + 5i$.
     

\end{enumerate}
\finenonce{004409}



\finexercice


\exercice{4410, quercia, 2010/03/12}
\enonce{004410}{$|\cos|$ et $|\sin|$ sur le cercle unité}
Calculer $\sup\{|\cos z| \text{ tel que } |z|\le 1\}$ et
$\sup\{|\sin z| \text{ tel que } |z|\le 1\}$.



\finenonce{004410}



\finexercice


\exercice{4411, quercia, 2010/03/12}
\enonce{004411}{Courbes}
Soient $M,M'$ deux points du plan d'affixes $z = x+iy$ et $z' = x'+iy'$.

\begin{enumerate}
  \item On suppose que $z$ et $z'$ sont liés par la relation : $z' = e^z$.
    \'Etudier la courbe décrite par $M'$ lorsque $M$ décrit~:
  \begin{enumerate}
    \item une droite $x = $ cste. \label{droite verticale}
    \item une droite $y = $ cste. \label{droite horizontale}
    \item une droite quelconque.
  \end{enumerate}
  \item Reprendre les questions \ref{droite verticale} et \ref{droite horizontale}
avec $z' = \cos z$.
\end{enumerate}
\finenonce{004411}



\finexercice


\exercice{4412, quercia, 2010/03/12}
\enonce{004412}{Centrale MP 2002}

Résoudre dans $\mathcal{M}_2(\C)$~: $\exp(M) = \begin{pmatrix}2i&1+i\cr0&2i\cr\end{pmatrix}$.

                           
\finenonce{004412}



\finexercice

\section{122.99 Autre}


\exercice{1946, gineste, 2001/11/01}

\enonce{001946}{Examen 2000}
 Soit $a>0$ fix\'e. Pour $n$
entier positif ou nul on définit $P_n(a)$ par $P_0(a)=1$,
$P_1(a)=a$, $P_2(a)=a(a+1)$ et, plus g\'en\'eralement
$P_{n+1}(a)=(n+a)P_n(a)$. Montrer que
$$L(a)=\lim _{n\infty}\frac{P_n(a)}{n!n^{a-1}}$$
existe et est un nombre strictement positif. M\'ethode:
consid\'erer la s\'erie de terme g\'en\'eral pour $n>0$: $u_n=
\log(n+a)-a\log(n+1)+(a-1)\log n,$
 comparer sa somme partielle d'ordre $n-1$ avec
$\log \frac{P_n(a)}{n!n^{a-1}},$ et, ... l'aide d'un
d\'eveloppement limit\'e en $1/n$ d'ordre convenable, montrer que,
$\sum _{n=1}^{\infty}u_n$ converge.
\finenonce{001946}



\finexercice



\exercice{1947, gineste, 2001/11/01}

\enonce{001947}{}
Soit $\alpha$ et $\beta$ deux nombres r\'eels ou complexes
tels que $\alpha\beta=-1$ et $|\alpha|>1>|\beta|.$ Pour $n$ dans
l'ensemble ${\bf Z}$ des entiers positifs ou n\'egatifs on pose
$F_n=\frac{1}{\alpha-\beta}(\alpha^n-\beta^n)$ et $L_n=\alpha^n+\beta^n$ (si $\alpha+\beta=1$
ces nombres sont appel\'es entiers de Fibonacci (1225) et de Lucas (1891)).
\begin{enumerate}
\item Montrer par le crit\`ere de D'Alembert que
la s\'erie $\sum_{n=1}^{\infty}\frac{1}{F_{2n+1}+1}$ converge et
calculer la limite de $Q_n=L_n/F_n$ si
$n\rightarrow +\infty.$ .
\item On admet (identit\'e de Backstrom (1981)) que pour
tous $n$ et $k$ de ${\bf Z}$ on a
$$
\frac{1}{F_{4n-2k-1}+F_{2k+1}}+\frac{1}{F_{4n+2k+1}+F_{2k+1}}=\frac{1}{2L_{2k+1}}
\left(Q_{2n+2k+1}-Q_{2n-2k-1}\right).$$
En faisant $k=0$ dans cette identit\'e,
calculer la somme partielle d'ordre $2n$ de la s\'erie initiale,
c'est \`a dire $s_{2n}=\sum_{j=1}^{2n}\frac{1}{F_{2j+1}+1}$
 en montrant par r\'ecurrence sur $n$ que $s_{2n}=\frac{1}{2L_1}
(Q_{2n+1}-Q_1).$ En d\'eduire la
somme de la s\'erie en termes de $\alpha$ et $\beta.$ Donner une expression
simple du terme g\'en\'eral de la s\'erie et de sa somme si $\alpha=\exp t$
et $\beta=-\exp (-t)$ si $t $ est r\'eel.
\item
 Montrer que la s\'erie  $\sum_{n=1}^{\infty}\frac{1}{F_{2n+1}+F_3}$
converge et calculer sa somme.
\end{enumerate}
\finenonce{001947}



\finexercice



\exercice{2656, matexo1, 2002/02/01}
\enonce{002656}{}
Indiquer pour quelles valeurs de $\alpha \in\R$ la s\'erie
$$\sum_{n\geq 0}\alpha ^n$$
converge, et calculer sa somme.

En d\'eduire l'\'ecriture en base 10 des nombres $1/9$ et $1/11$ et
plus g\'en\'eralement en base $k$, du nombre $1/(k-1)$ et du nombre
$1/(k+1)$.

\finenonce{002656}

\finexercice


\exercice{2726, matexo1, 2002/02/01}
\enonce{002726}{}
En comparant
avec les int\'egrales de Wallis
$$I_n = \int_0^{\pi /2} \cos^n t\,dt,$$
d\'eterminer la nature et calculer
la somme de la s\'erie
$$\displaystyle \sum_{n\geq 1} (-1)^n 2^{-2n} C_{2n}^n.$$ 
\finenonce{002726}
\finexercice


\exercice{4413, quercia, 2010/03/14}
\enonce{004413}{\'Etude de convergence}

\'Etudier la convergence des séries de terme général : 


\begin{enumerate}
  \item $\left(1+\frac 1n\right)^n - e$.
  \item $\ch^\alpha n - \sh^\alpha n$.
  \item $2\ln(n^3+1) - 3\ln(n^2+1)$.
  \item $\sqrt[n] {n+1} - \sqrt[n] n$.
  \item $\arccos\left(\frac {n^3+1}{n^3+2} \right)$.
  \item $\frac {a^n}{1+a^{2n}}$.
  \item $\frac{(-1)^n}{\sqrt{n^2+n}}$.
  \item $\frac{(-1)^n}{\ln n}$.
  \item $\frac{1+(-1)^n\sqrt n}n$.
  \item $\frac{2.4.6\dots(2n)}{n^n}$.
  \item $\frac{1!+2!+\dots+n!}{(n+2)!}$.
  \item $\frac{1!-2!+\dots\pm n!}{(n+1)!}$.
  \item $\frac{(-1)^n}{\ln n + \sin(2n\pi/3)}$.
  \item $\sqrt{1+\frac{(-1)^n}{\sqrt n}}-1$.
  \item $\frac{(-1)^n}{\sqrt n + (-1)^n}$.
  \item $\frac{(-1)^{[\sqrt n\,]}}n$.
  \item $\frac{(\ln n)^n}{n^{\ln n}}$.
  \item $\frac1{(\ln n)^{\ln n}}$.
\end{enumerate}
\finenonce{004413}



\finexercice


\exercice{4414, quercia, 2010/03/14}
\enonce{004414}{Centrale PC 1999}

Soit la suite de terme général~:
$u_n = (n^4+n^2)^{1/4}-P(n)^{1/3}$ où $P$ est un polynôme.
A quelle condition sur~$P$ la série $\sum u_n$ converge-t-elle~?

\finenonce{004414}



\finexercice


\exercice{4415, quercia, 2010/03/14}
\enonce{004415}{Ensi PC 1999}

Quelle est la nature de la série de terme général
$\ln\Bigl(1+\frac{(-1)^n}{\sqrt{n(n+1)}}\Bigr)$~?

\finenonce{004415}



\finexercice


\exercice{4416, quercia, 2010/03/14}
\enonce{004416}{Mines MP 2000}

Soit $\alpha >0$. Étudier la série $\sum u_n$, avec 
$u_n=\frac{(-1)^n}{\sqrt{n^{\alpha}+(-1)^n}}\cdotp$
\finenonce{004416}



\finexercice


\exercice{4417, quercia, 2010/03/14}
\enonce{004417}{Mines MP 2003}
Si $\alpha>0$, donner la nature des séries $\sum_{n\ge 2}\frac{(-1)^n}{(-1)^n+n^\alpha}$,
$\sum_{n\ge 2}\ln\Bigl(1+\frac{(-1)^n}{n^\alpha}\Bigr)$ et $\sum_{n\ge 2}\frac1{n\ln n}$.

\finenonce{004417}



\finexercice


\exercice{4418, quercia, 2010/03/14}
\enonce{004418}{Ensi PC 1999}

Soit $(u_n)$ une suite réelle telle que
$\frac{u_{2n+1}}{u_{2n}}\to a$ et
$\frac{u_{2n}}{u_{2n-1}}\to b$ lorsque $n\to\infty$.
\'Etudier la convergence de $\sum u_n$.
\finenonce{004418}



\finexercice


\exercice{4419, quercia, 2010/03/14}
\enonce{004419}{Encadrement}

Soient $\sum u_n$, $\sum v_n$, $\sum w_n$ trois séries réelles telles que
$\sum u_n$ et $\sum w_n$ convergent, et
${u_n \le v_n \le w_n}$ pour tout~$n$. Montrer que $\sum v_n$ converge.
\finenonce{004419}



\finexercice


\exercice{4420, quercia, 2010/03/14}
\enonce{004420}{Calcul approché}

Montrer que la série $\sum_{n=1}^\infty \Bigl(n\sin(0.4/n)\Bigr)^n$ converge.
Calculer à la machine une valeur approchée à $10^{-8}$ près de sa somme.

\finenonce{004420}



\finexercice


\exercice{4421, quercia, 2010/03/14}
\enonce{004421}{Ensi MP 2002}
On suppose que la série à termes positifs de terme général $u_n$ est divergente et on pose 
$S_n=\sum_{k=0}^n u_k$.

Soit $f : {\R^+} \to {\R^+}$ une application continue décroissante.
Comparer les énoncés~:

1. $f$ est intégrable

2. La série de terme général $u_nf(S_n)$ converge.

\finenonce{004421}



\finexercice


\exercice{4422, quercia, 2010/03/14}
\enonce{004422}{Centrale P' 1996}

Montrer que la série $\sum_{n=1}^\infty \frac{n^2}{(1+n^2)^2}$ converge.
Calculer une valeur approchée à $10^{-4}$ près de sa somme.

\finenonce{004422}



\finexercice


\exercice{4423, quercia, 2010/03/14}
\enonce{004423}{${C_{2n}^n}/{n4^n}$}

L'une au moins des deux séries : $\sum \frac{C_{2n}^n}{n4^n}$ et
$\sum \frac{n4^n}{C_{2n}^n}$ diverge. Dire pourquoi et dire laquelle.

\finenonce{004423}



\finexercice


\exercice{4424, quercia, 2010/03/14}
\enonce{004424}{$1/(1+n^2u_n)$, Mines-Ponts MP 2005}

Soit $(u_n)$ une suite réelle positive et $v_n = \frac1{1+n^2u_n}$.
Montrer que $\sum u_n$ converge $ \Rightarrow  \sum v_n$ diverge.
\'Etudier le cas où $\sum u_n$ diverge.



\finenonce{004424}



\finexercice


\exercice{4425, quercia, 2010/03/14}
\enonce{004425}{${a_n}/{(1+a_1)(1+a_2)\dots(1+a_n)}$}

Soit $(a_n)$ une suite réelle positive.
On pose $u_n = \frac {a_n}{(1+a_1)(1+a_2)\dots(1+a_n)}$.

\begin{enumerate}
  \item Montrer que la série $\sum u_n$ converge.
    
  \item Calculer $\sum_{n=1}^\infty u_n$ lorsque $a_n = \frac 1{\sqrt n}$.
    


\end{enumerate}
\finenonce{004425}



\finexercice


\exercice{4426, quercia, 2010/03/14}
\enonce{004426}{$1/a^{\text{nb de chiffres de }n}$}

Pour $n \in \N^*$ on note $p_n$ le nombre de chiffres de l'écriture décimale
de $n$ (sans zéros inutiles).
Soit $a > 0$. \'Etudier la convergence et déterminer la somme éventuelle
de la série $\sum_{k=1}^\infty \frac1{a^{p_k}}$.

\finenonce{004426}



\finexercice


\exercice{4427, quercia, 2010/03/14}
\enonce{004427}{Cauchy-Schwarz}

Soient $(u_n)$, $(v_n)$ deux suites réelles telles que $\sum u_n^2$
et $\sum v_n^2$ convergent.

\begin{enumerate}
  \item Montrer que $\sum u_nv_n$ converge.
  \item Montrer que $\sum (u_n+v_n)^2$ converge et :
    $\sqrt{\sum(u_n+v_n)^2} \le \sqrt{\sum u_n^2} + \sqrt{\sum v_n^2}$.

\end{enumerate}
\finenonce{004427}



\finexercice


\exercice{4428, quercia, 2010/03/14}
\enonce{004428}{${(-1)^n}/(n^{3/4}+\cos n)$}

Soit $u_n = \frac{(-1)^n}{n^{3/4}+\cos n}$.

\begin{enumerate}
  \item La série $\sum u_n$ est-elle absolument convergente ?
  \item En écrivant $u_n = \frac{(-1)^n}{n^{3/4}} + v_n$, étudier la convergence
    de $\sum u_n$.
    

\end{enumerate}
\finenonce{004428}



\finexercice


\exercice{4429, quercia, 2010/03/14}
\enonce{004429}{Reste d'une série alternée}

On pose $u_n = \sum_{k=n}^\infty \frac {(-1)^k}{\sqrt{k+1}}$.
\'Etudier la convergence de la série $\sum u_n$.

\finenonce{004429}



\finexercice


\exercice{4430, quercia, 2010/03/14}
\enonce{004430}{Calcul de sommes}

Calculer les sommes des séries suivantes : 
\begin{enumerate}
  \item $\sum_{k=2}^\infty \frac 1{k^2-1}$.
  \item $\sum_{k=1}^\infty \frac 1{k(k+1)(k+2)}$.
  \item $\sum_{k=1}^\infty \frac 1{k(k+1)\dots(k+p)}$.
  \item $\sum_{k=0}^\infty \frac 1{k^3+8k^2+17k+10}$.
  \item $\sum_{k=1}^\infty \ln\left(1+\frac2{k(k+3)}\right)$.
  \item $\sum_{k=2}^\infty \ln\left(1-\frac1{k^2}\right)$.
  \item $\sum_{k=0}^\infty \ln\left(\cos\frac\alpha{2^k}\right)$.
  \item $\sum_{k=0}^\infty 2^{-k}\tan(2^{-k}\alpha)$.
  \item $\sum_{k=0}^\infty \frac{2k^3-3k^2+1}{(k+3)!}$.
  \item $\sum_{n=p}^\infty C_n^p x^n$.
  \item $\sum_{k=1}^\infty \frac{x^k}{(1-x^k)(1-x^{k+1})}$.
  \item $\sum_{k=1}^\infty \frac{k-n[k/n]}{k(k+1)}$.
\end{enumerate}
\finenonce{004430}



\finexercice


\exercice{4431, quercia, 2010/03/14}
\enonce{004431}{} 

Convergence et somme de la série de terme général
$u_n=\frac{\strut\lfloor\sqrt{n+1}\rfloor-\lfloor\sqrt{n}\rfloor}n$.

\finenonce{004431}



\finexercice


\exercice{4432, quercia, 2010/03/14}
\enonce{004432}{Chimie P 90}

\begin{enumerate}
  \item Résoudre les équations différentielles :
    $y'' + 2y' + 2y = 0$, $y'' + 4y' + 4y = 2e^{-x}\cos x$.
    
  \item Soit $f$ la solution commune.
    On définit la série de terme général
    $u_n =  \int_{x=n\pi}^{(n+1)\pi} f(x)\,dx$.
    Montrer que $\sum u_n$ converge et calculer sa somme.
    

\end{enumerate}
\finenonce{004432}



\finexercice


\exercice{4433, quercia, 2010/03/14}
\enonce{004433}{$1/n^2(n+1)^2$}

On admet que $\sum_{k=1}^\infty \frac1{k^2} = \frac{\pi^2}6$.
Calculer $\sum_{k=1}^\infty \frac 1{k^2(k+1)^2}$.

\finenonce{004433}



\finexercice


\exercice{4434, quercia, 2010/03/14}
\enonce{004434}{$1/(1^2+2^2+...+n^2)$}

On admet que $\sum_{k=1}^\infty \frac {(-1)^{k+1}}k = \ln 2$.
Montrer que la série $\sum_{k=1}^\infty \frac1{1^2 + 2^2 + \dots + k^2}$ est
convergente et calculer sa somme.


\finenonce{004434}



\finexercice


\exercice{4435, quercia, 2010/03/14}
\enonce{004435}{$\ln(n) + a\ln(n+1) + b\ln(n+2)$}

Pour quelles valeurs de $a,b\in \R$ la série de terme général
$\ln(n) + a\ln(n+1) + b\ln(n+2)$ est-elle convergente~?
Calculer alors la somme de la série.


\finenonce{004435}



\finexercice


\exercice{4436, quercia, 2010/03/14}
\enonce{004436}{$\Arctan(1/(k^2+k+1))$}

Montrer que
$\sum_{k=0}^\infty \Arctan\Bigl(\frac1{k^2+k+1}\Bigr) = \frac \pi2$.
(On pourra calculer $\tan s_n$)

\finenonce{004436}



\finexercice


\exercice{4437, quercia, 2010/03/14}
\enonce{004437}{$\Arctan(n+a) - \Arctan n$}

Soit $a \in \R$.

\begin{enumerate}
  \item Montrer que la série de terme général $\Arctan(n+a) - \Arctan n$ est
    convergente.
    
  \item On pose $S(a) = \sum_{k=0}^\infty (\Arctan(k+a) - \Arctan k)$.
    Trouver $\lim_{a\to+\infty} S(a)$.
    

\end{enumerate}
\finenonce{004437}



\finexercice


\exercice{4438, quercia, 2010/03/14}
\enonce{004438}{Pile en porte à faux}

Peut-on empiler 100 pièces de 1F de sorte que la dernière soit complètement
en porte à faux ?
(c'est-à-dire que sa projection sur un plan horizontal ne rencontre pas la projection de
la première pièce)

\finenonce{004438}



\finexercice


\exercice{4439, quercia, 2010/03/14}
\enonce{004439}{Recherche d'équivalents}

Par comparaison à une intégrale, donner un équivalent de :
\begin{enumerate}
  \item $\sum_{k=n+1}^{2n} \frac 1{\sqrt k}$. 
  \item $\sum_{k=2}^{n} \frac 1{k\ln k}$.     


\end{enumerate}
\finenonce{004439}



\finexercice


\exercice{4440, quercia, 2010/03/14}
\enonce{004440}{$\ln^2(k)$}

Par comparaison à une intégrale, donner un équivalent de
$u_n = \sum_{k=1}^n \ln^2k $. La série de terme général $\frac1{u_n}$
est-elle convergente ?


\finenonce{004440}



\finexercice


\exercice{4441, quercia, 2010/03/14}
\enonce{004441}{$k^{-2/3}$}

Trouver la partie entière de $\sum_{k=1}^{10^9} k^{-2/3}$.


\finenonce{004441}



\finexercice


\exercice{4442, quercia, 2010/03/14}
\enonce{004442}{$(-1)^k\sqrt k$}

On pose $u_n = \sum_{k=1}^{2n} (-1)^k\sqrt k$. Donner un équivalent de $u_n$
quand $n\to\infty$.
(Regrouper les termes deux par deux puis comparer à une intégrale)


\finenonce{004442}



\finexercice


\exercice{4443, quercia, 2010/03/14}
\enonce{004443}{Constante d'Euler}

Soit $f : {\R^+} \to {\R^+}$ décroissante.
On pose $u_n = f(n)$ et $s_n = u_0 + \dots + u_n$.

Montrer que la suite de terme général $s_n -  \int_{t=0}^{n+1} f(t)\,d t$ est
convergente. Donner une interprétation graphique de ce fait.

Application : On pose $\gamma = \lim_{n\to\infty}
\left(1 + \frac 12 + \dots + \frac 1n - \ln n\right)$.
Justifier l'existence de $\gamma$ et montrer que $\frac 12 \le \gamma \le 1$.
\finenonce{004443}



\finexercice


\exercice{4444, quercia, 2010/03/14}
\enonce{004444}{Constante d'Euler (Centrale MP 2003)}
Soit $S_n = \sum_{k=1}^{n-1}\frac1k - \frac1n - \ln n$ et
$T_n = \sum_{k=1}^{n-1}\frac1k + \frac1n - \ln n$.
Les suites $(S_n)$ et $(T_n)$  sont-elles adjacentes~?
\finenonce{004444}



\finexercice


\exercice{4445, quercia, 2010/03/14}
\enonce{004445}{Constante d'Euler, Mines-Ponts MP 2005}

Soit $u_{n,k}$ le reste de la division du $n$ par $k$.
Quelle est la limite de $\frac1n\sum_{k=1}^n\frac{u_{n,k}}k$~?


\finenonce{004445}



\finexercice


\exercice{4446, quercia, 2010/03/14}
\enonce{004446}{Mines MP 2003}
Soit la suite de terme général $u_n = \frac{\ln 2}2 + \frac{\ln 3}3 + \dots + \frac{\ln n}n$.

\begin{enumerate}
  \item Donner un équivalent de~$u_n$ en~$+\infty$.
    
  \item Montrer que la suite de terme général~: $v_n = u_n - \frac{\ln^2n}2$ est convergente.
    
  \item Soit $\ell = \lim_{n\to\infty} v_n$. Donner un équivalent de~$v_n-\ell$.
    
\end{enumerate}
\finenonce{004446}



\finexercice


\exercice{4447, quercia, 2010/03/14}
\enonce{004447}{Centrale MP 2001}

Donner un équivalent simple de $\sum_{k=0}^{n-1}\frac1{n^2-k^2}$.

\finenonce{004447}



\finexercice


\exercice{4448, quercia, 2010/03/14}
\enonce{004448}{$1/n\ln^2(n)$}

\begin{enumerate}
  \item Prouver la convergence de la série de terme général
    $u_n = \frac1{n\ln^2 n}$.
  \item On note $S_n = \sum_{k=2}^n u_k$ et $S = \sum_{k=2}^\infty u_k$.
    Montrer que $\frac1{\ln(n+1)} \le S-S_n \le \frac1{\ln n}$ pour $n \ge 2$.
  \item Montrer que si $S_n$ est une valeur approchée de $S$ à $10^{-3}$ près
    alors $n > 10^{434}$.
  \item On suppose disposer d'une machine calculant un million de termes de la
    série par seconde avec 12 chiffres significatifs. Peut-on obtenir une valeur
    approchée de $S$ à $10^{-3}$ près ?
    (Remarque : 1 an $\approx$ 32 millions de secondes)
  \item Donner une valeur approchée de $S$ à $10^{-3}$ près.
    

\end{enumerate}
\finenonce{004448}



\finexercice


\exercice{4449, quercia, 2010/03/14}
\enonce{004449}{$(x-1)\zeta(x) \to 1$}
 
Pour $x > 1$ on note $\zeta(x) = \sum_{k=1}^\infty \frac 1{k^x}$.
En comparant $\zeta(x)$ à une intégrale,
trouver $\lim_{x\to1^+} (x-1)\zeta(x)$.
 
\finenonce{004449}
 
 
 
\finexercice


\exercice{4450, quercia, 2010/03/14}
\enonce{004450}{$u_n/(1+u_n)$}

Soit $\sum u_n$ une série à termes positifs et $v_n = \frac{u_n}{1+u_n}$.
Montrer que $\sum u_n$ et $\sum v_n$ ont même nature.


\finenonce{004450}



\finexercice


\exercice{4451, quercia, 2010/03/14}
\enonce{004451}{Série des restes}

\begin{enumerate}
  \item Soit $(u_n)$ une suite réelle
    telle que $\sum |u_n|$ et $\sum n|u_n|$ convergent.
    On note $v_n = \sum_{k=n}^\infty u_k$.
 \begin{enumerate} 
    \item Montrer que $nv_n \to 0$ lorsque $n\to\infty$.
    \item Montrer que $\sum_{n=1}^\infty v_n = \sum_{n=1}^\infty nu_n$.
  \end{enumerate} 
  \item Application : Calculer lorsque c'est possible : $\sum_{k=1}^\infty kr^k$.
    


\end{enumerate}
\finenonce{004451}



\finexercice


\exercice{4452, quercia, 2010/03/14}
\enonce{004452}{X MP$^*$ 2001}

Soit $(u_n)$ une suite réelle positive, $U_n=\sum_{i=0}^n u_i$ et
$\alpha>0$ un réel donné. On suppose $\frac{U_n}{nu_n}\to \alpha$ lorsque $n\to\infty$. 
Étudier la suite de terme général $\frac1{n^2u_n}\sum_{k=0}^n ku_k$.
\finenonce{004452}



\finexercice


\exercice{4453, quercia, 2010/03/14}
\enonce{004453}{$\sum nu_n$ converge}

On considère une suite $(u_n)_{n\ge 1}$ telle que la
série $\sum_{n\ge 1} n\, u_n$ converge. Montrer
que la série $\sum_{n\ge 1} u_n$ converge.
\finenonce{004453}



\finexercice


\exercice{4454, quercia, 2010/03/14}
\enonce{004454}{$(u_n)$ décroit}

Soit $(u_n)_{n\ge 1}$ une suite réelle positive décroissante telle que
$\sum u_n$ converge.

\begin{enumerate}
  \item Montrer que $nu_n \to 0$ lorsque $n\to\infty$.
    $\Bigl($considérer $\sum_{k=n+1}^{2n} u_k \Bigr)$
  \item Montrer que $\sum_{n=1}^\infty n(u_n - u_{n+1})$ converge et a même somme que $\sum_{n=1}^\infty u_n$.
  \item Application : calculer pour $0 \le r < 1$ :
    $\sum_{k=1}^\infty kr^k$ et $\sum_{k=1}^\infty k^2r^k$.
\end{enumerate}
\finenonce{004454}



\finexercice


\exercice{4455, quercia, 2010/03/14}
\enonce{004455}{$u_n/S_n$}

Soit $(u_n)$ une suite à termes strictement positifs convergeant vers~$0$. On pose
$S_n = \sum_{k=0}^n u_k$.

\begin{enumerate}
  \item Si la série $\sum u_n$ converge, que dire de la série
    $\sum \frac{u_n}{S_n}$ ?
  \item Si la série $\sum u_n$ diverge, montrer que la série
    $\sum \frac{u_n}{S_n}$ diverge aussi.
    On pourra considérer $p_n = \prod_{k=1}^n \left(1-\frac{u_k}{S_k}\right)$.

    

\end{enumerate}
\finenonce{004455}



\finexercice


\exercice{4456, quercia, 2010/03/14}
\enonce{004456}{Polytechnique MP$^*$ 2000}

On donne une suite de réels strictement positifs $(a_n)$, décroissante
et de limite nulle. Montrer que la série de terme général
$\frac{a_n-a_{n+1}}{a_n}$ diverge.
\finenonce{004456}



\finexercice


\exercice{4457, quercia, 2010/03/14}
\enonce{004457}{$(u_n + u_{n+1} + \dots + u_{2n-1})/n$}

Soit $\sum u_n$ une série à termes positifs.
On pose $v_n = \frac{ u_n + u_{n+1} + \dots + u_{2n-1} }n$.
Montrer que $\sum v_n$ a même nature que $\sum u_n$.



\finenonce{004457}



\finexercice


\exercice{4458, quercia, 2010/03/14}
\enonce{004458}{$\sum ku_k/n(n+1)$}

Soit $(u_n)_{n\ge 1}$ une suite positive.
On pose $v_n = \frac{1}{n(n+1)}\sum_{k=1}^n ku_k$.
Montrer que les séries $\sum u_n$ et $\sum v_n$ ont même nature et
éventuellement même somme.


\finenonce{004458}



\finexercice


\exercice{4459, quercia, 2010/03/14}
\enonce{004459}{$\sum ku_k/n^2$}

Soit $\sum u_n$ une série à termes positifs convergente.

\'Etudier la convergence de la série de terme général
$v_n = \frac1{n^2}\sum_{k=1}^n ku_k$.


\finenonce{004459}



\finexercice


\exercice{4460, quercia, 2010/03/14}
\enonce{004460}{Principe d'accumulation}

Soit $(u_n)$ une suite réelle positive décroissante. On pose $v_n = 2^nu_{2^n}$.
Montrer que les séries $\sum u_n$ et $\sum v_n$ ont même nature.


Applications : 
Retrouver la convergence des séries de Riemann $\sum\frac1{n^\alpha}$.\par
\'Etudier la convergence des séries de Bertrand~:
$\sum\frac1{n(\ln n)^\alpha}$.
\finenonce{004460}



\finexercice


\exercice{4461, quercia, 2010/03/14}
\enonce{004461}{$u_{n+1} = 1/{ne^{u_n}}$. Ensi P 90}

Soit $(u_n)$ définie par : $u_1 \in \R$, $u_{n+1} = \frac1{ne^{u_n}}$.
Quelle est la nature de la série $\sum u_n$ ?


\finenonce{004461}



\finexercice


\exercice{4462, quercia, 2010/03/14}
\enonce{004462}{$x_{n+1} = x_n + x_n^2$}

Soit $(x_n)$ une suite définie par : $x_0 > 0$,
$\forall\ n \in \N,\ x_{n+1} = x_n + x_n^2$.

\begin{enumerate}
  \item Montrer que $x_n \to +\infty$ lorsque $n\to\infty$.
  \item On pose $u_n = 2^{-n}\ln x_n$. Montrer que la suite $(u_n)$ est convergente.
    (On étudiera la série $\sum u_{n+1}-u_n$)
  \item En déduire qu'il existe $\alpha > 0$ tel que $x_n \sim \alpha^{2^n}$.

\end{enumerate}
\finenonce{004462}



\finexercice


\exercice{4463, quercia, 2010/03/14}
\enonce{004463}{$u_{n+1} = u_n - u_n^2$}

On considère la suite $(u_n)$ définie par : $0 < u_0 < 1$
et $\forall\ n\in\N,\ u_{n+1} = u_n - u_n^2$.

\begin{enumerate}
  \item Montrer que la suite $(u_n)$ converge. Quelle est sa limite ?
  \item Montrer que la série de terme général $u_n^2$ converge.
  \item Montrer que les séries de termes généraux
    $\ln\left(\frac{u_{n+1}}{u_n}\right)$ et $u_n$ divergent.
  \item Montrer que $u_n < \frac 1{n+1}$ et que la suite $(nu_n)$ est croissante.
    On note $\ell$ sa limite.
  \item On pose $u_n = \frac{\ell-v_n}n$. Montrer que la série de terme général
    $v_{n+1}-v_n$ converge.
  \item En déduire que $u_n$ est équivalent à $\frac 1n$.
\end{enumerate}
\finenonce{004463}



\finexercice


\exercice{4464, quercia, 2010/03/14}
\enonce{004464}{$u_{n+1}/u_n = (n+a)/(n+b)$}

Soit $(u_n)$ une suite définie par la donnée de $u_0 \in \R^*$ et la
relation : $\forall\ n\in\N,\ \frac{u_{n+1}}{u_n} = \frac{n+a}{n+b}$
où $a$, $b$ sont deux constantes réelles ($-a,-b\notin\N$).

\begin{enumerate}
  \item Montrer que $u_n$ est de signe constant à partir d'un certain rang.
  \item On pose $v_n = (n+b-1)u_n$. \'Etudier la convergence de la suite $(v_n)$
    (on introduira la série de terme général $\ln(v_{n+1})-\ln(v_n)$).
    
  \item En déduire que la série $\sum u_n$ converge si et seulement si $a-b+1 < 0$
    et calculer sa somme en fonction de $a,b,u_0$.
    

\end{enumerate}
\finenonce{004464}



\finexercice


\exercice{4465, quercia, 2010/03/14}
\enonce{004465}{} 

On se donne $u_1$ et $a$ deux réels strictement positifs et l'on
définit par récurrence la suite $(u_n)$ par $u_{n+1}=u_n+\frac1{n^a u_n}\cdotp$
Étudiez la limite de la suite $(u_n)$, et, quand $a\le1$, en donner un
équivalent.

\finenonce{004465}



\finexercice


\exercice{4466, quercia, 2010/03/14}
\enonce{004466}{$1/{k^\alpha(n-k)^\alpha}$}

Soit $\alpha > 0$.
On pose $u_n = \sum_{k=1}^{n-1} \frac 1{k^\alpha(n-k)^\alpha}$.
\'Etudier la convergence de $\sum u_n$.


\finenonce{004466}



\finexercice


\exercice{4467, quercia, 2010/03/14}
\enonce{004467}{Produit de Cauchy de trois séries}

Soient $\sum a_n$, $\sum b_n$, $\sum c_n$ trois séries absolument convergentes
de sommes $A$, $B$, $C$.

On pose $u_n = \sum_{i+j+k=n} a_ib_jc_k$. Montrer que $\sum u_n = ABC$.

\finenonce{004467}



\finexercice


\exercice{4468, quercia, 2010/03/14}
\enonce{004468}{Produit de séries géométriques}

Soient $a \in {[0,1[}$. \'Ecrire $\frac1{(1-a)^2}$ comme produit de deux
séries. En déduire la somme de la série $\sum_{k=0}^\infty ka^k$.
Calculer par la même méthode $\sum_{k=0}^\infty k^2a^k$.


\finenonce{004468}



\finexercice


\exercice{4469, quercia, 2010/03/14}
\enonce{004469}{Produit de séries géométriques}

Pour $n \in \N$ on note $T_n$ le nombre de manières de décomposer
$n$ francs avec des pièces de 1, 2, 5 et 10 francs ($T_0 = 1$).
Montrer que :
$$\forall\ x \in {[0,1[},\
\sum_{k=0}^\infty T_kx^k = \frac1{(1-x)(1-x^2)(1-x^5)(1-x^{10})}.$$
\finenonce{004469}



\finexercice


\exercice{4470, quercia, 2010/03/14}
\enonce{004470}{$\sum u_k/2^{n-k}$}

Soit $\sum u_n$ une série convergente.
On pose $v_n = \frac{u_n}{1} + \frac{u_{n-1}}{2} + \dots + \frac{u_0}{2^n}$.

\begin{enumerate}
  \item Montrer que $v_n \to 0$ lorsque $n\to\infty$.
    
  \item Montrer que $\sum v_n$ converge et donner sa valeur.
    

\end{enumerate}
\finenonce{004470}



\finexercice


\exercice{4471, quercia, 2010/03/14}
\enonce{004471}{$\sum a_n/n^p = 0$}

Soit $(a_n)$ une suite bornée telle que pour tout entier $p \ge 2$ :
$\sum_{n=1}^\infty \frac{a_n}{n^p} = 0$.
Montrer que : $\forall\ n\in\N^*,\ a_n = 0$.


\finenonce{004471}



\finexercice


\exercice{4472, quercia, 2010/03/14}
\enonce{004472}{$\sum x_{kn} = 0$}

Soit $\sum_{n\ge1} x_n$ une série absolument convergente telle que
pour tout entier $k \ge 1$ on a $\sum_{n=1}^\infty x_{kn} = 0$.

Montrer que : $\forall\ n\in\N^*,\ x_n = 0$.

\finenonce{004472}



\finexercice


\exercice{4473, quercia, 2010/03/14}
\enonce{004473}{Césaro}

\begin{enumerate}
  \item Soient $k,p \in \N$ avec $k \le p$.
    Montrer que $\sum_{n=k}^p \frac{{C_n^k}-C_n^{k+1}}{2^n}
    = \frac{C_{p+1}^{k+1}}{2^p}$.
    
  \item Soit $(u_n)$ une série convergente.
    On pose $v_n = \frac1{2^n} \sum_{p=0}^n C_n^pu_p$.
    Montrer que la série $(v_n)$ est convergente.



\end{enumerate}
\finenonce{004473}



\finexercice


\exercice{4474, quercia, 2010/03/14}
\enonce{004474}{$nu_n\to 0$}

Soit $(u_n)$ une série convergente à termes positifs décroissants.

\begin{enumerate}
  \item Montrer que $nu_n\to 0$ lorsque $n\to\infty$.
    
  \item Montrer que $\sum_{u_k \ge 1/n} \frac 1{u_k} = \text{o}(n^2)$.
    


\end{enumerate}
\finenonce{004474}



\finexercice


\exercice{4475, quercia, 2010/03/14}
\enonce{004475}{$u_n/R_n^p$}

Soit $(a_n)$ une série positive convergente, $A = \sum_{k=0}^\infty a_k$,
$R_n = \sum_{k=n}^\infty a_k$ et $p\in{]0,1[}$.

\begin{enumerate}
  \item Montrer qu'il existe $C_p \in \R$ tel que
    $\sum_{n=0}^\infty \frac {a_n}{R_n^p} \le C_p A^{1-p}$.
    
  \item Trouver la meilleure constante $C_p$.
    

\end{enumerate}
\finenonce{004475}



\finexercice


\exercice{4476, quercia, 2010/03/14}
\enonce{004476}{$u_{n+1} = u_n + {a_n}/{u_n}$}

Soit $(a_n)$ une suite réelle positive et $(u_n)$ la suite
définie par la relation de récurrence~:
$u_{n+1} = u_n + \frac{a_n}{u_n}$ avec $u_0 > 0$.
Montrer que la suite $(u_n)$ converge si et seulement si la série
$\sum a_n$ converge.


\finenonce{004476}



\finexercice


\exercice{4477, quercia, 2010/03/14}
\enonce{004477}{Raabe-Duhamel}

Soit $(u_n)$ une suite réelle positive telle que
$\frac{u_{n+1}}{u_n} = 1 - \frac\alpha n + O\left(\frac 1{n^2}\right)$.
Montrer qu'il existe $A > 0$ tel que $u_n\sim \frac A{n^\alpha}$.
\finenonce{004477}



\finexercice


\exercice{4478, quercia, 2010/03/14}
\enonce{004478}{Stirling$++$}

Montrer que $n! = \left(\frac ne\right)^n\sqrt{2\pi n}\left(1+\frac 1{12n} +  O\left(\frac 1{n^2}\right)\right)$.
\finenonce{004478}



\finexercice


\exercice{4479, quercia, 2010/03/14}
\enonce{004479}{Développement factoriel}

Soit $\cal S$ l'ensemble des suites croissantes d'entiers $(q_i)$ telles que $q_0\ge 2$.

\begin{enumerate}
  \item Si $s = (q_i) \in \cal S$, montrer que la série
$\sum_{k=0}^\infty \frac1{q_0\dots q_k}$ converge. On note $\Phi(s)$ sa somme.

  \item Montrer que l'application $\Phi : {\cal S} \to {]0,1]}$ est bijective.
  \item Soit $s = (q_i) \in \cal S$. Montrer que $\Phi(s) \in\Q$ si et seulement
    si $s$ est stationnaire.
\end{enumerate}
\finenonce{004479}



\finexercice


\exercice{4480, quercia, 2010/03/14}
\enonce{004480}{Développement asymptotique}

\begin{enumerate}
  \item Montrer qu'il existe $C\in\R$ tel que $\sum_{k=1}^n \frac{\ln k}k = \frac12\ln^2(n) + C +  o(1)$.
  \item Prouver~: $\frac{\ln2}2 -  \int_{t=1}^3\frac{\ln t}t\,d t \le C \le
               \frac{\ln2}2 + \frac{\ln3}3 -  \int_{t=1}^3\frac{\ln t}t\,d t$.
  \item Prouver~: $\sum_{k=1}^n \frac{\ln k}k = \frac12\ln^2(n) + C + \frac{\ln n}{2n} +  o\Bigl(\frac{\ln n}{n}\Bigr)$.
\end{enumerate}
\finenonce{004480}



\finexercice


\exercice{4481, quercia, 2010/03/14}
\enonce{004481}{} 

Soit $(u_n)$ une suite de complexes telle que $\frac{u_1+\dots+u_n}n\to\ell\in\C$ lorsque $n\to\infty$.
Montrer que $\frac{1}{\ln(n)}\Bigl(\frac{u_1}1+\dots+\frac{u_n}n\Bigr)\to \ell$ lorsque $n\to\infty$.


\finenonce{004481}



\finexercice


\exercice{4482, quercia, 2010/03/14}
\enonce{004482}{} 

Soit $(u_n)$ une suite de complexes qui converge au sens
de Césaro vers zéro.

Étudiez la suite de terme général
$v_n=\sum_{k=0}^n \frac{u_k}{n+k+1}\cdotp$

\finenonce{004482}



\finexercice


\exercice{4483, quercia, 2010/03/14}
\enonce{004483}{Centrale MP 2000}

Soient deux suites de termes généraux $u_n$ et $v_n$ définies
par la donnée de $u_1$ et~$v_1$, tous deux réels, et les relations~:
$$u_{n+1} = u_n - \frac{v_n}{n(n+1)},
  \qquad
  v_{n+1} = v_n + \frac{u_n}{n(n+1)}.$$
Montrer que ces suites sont définies et bornées.
\finenonce{004483}



\finexercice


\exercice{4484, quercia, 2010/03/14}
\enonce{004484}{Produits infinis, Polytechnique 2000}

On considère une suite $(a_n)$ de réels et on définit $P_N=\prod_{n=1}^N(1+a_n)$
et $S_N=\sum_{n=1}^N a_n$.

\begin{enumerate}
  \item  On suppose que pour tout $n$, $a_n\ge 0$. 
  \begin{enumerate}
    \item  Montrer que, pour tout $N$, $1+S_N\le P_N \le e^{S_N}$.
    \item Comparer les convergences respectives des suites $(S_N)$ et $(P_N)$.
  \end{enumerate}
  \item On suppose maintenant que pour tout $n$, $-1\le a_n\le 0$. 
  \begin{enumerate}
    \item La relation précédente est-elle encore vérifiée ?
    
    \item Discuter de la convergence des suites $(S_N)$ et $(P_N)$.
 \end{enumerate} 
  \item On suppose que $(a_n)$ est de signe quelconque et que pour tout $n$, $1+a_n> 0$.
    On suppose de plus que la série $\sum a_n$ converge. 
    Montrer que $(P_N)$ a une limite et que cette limite est
    nulle si et seulement si $\sum a_n^2$ diverge.
    
  \item Complément. On suppose que la suite $(a_n)$ est complexe, que pour tout $n$ $|a_n|<1$ et que la 
    série $\sum |a_n|$ est convergente.
  \begin{enumerate}
     \item Montrer que $\prod_{n=1}^{\infty} (1+|a_n|)$ existe, puis que $\prod_{n=1}^{\infty} (1+a_n)$
    existe $\Bigl($on pourra  démontrer et utiliser l'inégalité $\Bigl|\prod_{n=1}^{N} (1+a_n)-1\Bigr|\le
    \prod_{n=1}^{N} (1+|a_n|)-1\Bigr)$.
    
     \item Montrer que $\prod_{n=1}^{\infty} (1+a_n)$ n'est pas nul.
  \end{enumerate}
\end{enumerate}
\finenonce{004484}



\finexercice


\exercice{4485, quercia, 2010/03/14}
\enonce{004485}{Polytechnique MP 2002}
Trouver les fonctions $f : {[0,1]} \to \R$ continues vérifiant~:
$\forall\ x\in{[0,1]},\ f(x) = \sum_{n=1}^\infty \frac{f(x^n)}{2^n}$.


\finenonce{004485}



\finexercice


\exercice{4486, quercia, 2010/03/14}
\enonce{004486}{ENS Cachan MP$^*$ 2005}

Soit $P(n)=\max\{p \text{ premier},\, p\mid n\}$. Montrer que $\sum_{n}\frac{1}{nP(n)}$ converge.
\finenonce{004486}



\finexercice


\exercice{5142, rouget, 2010/06/30}
\enonce{005142}{IT}
Cet exercice est consacré aux sommes de termes consécutifs d'une suite arithmétique ou d'une suite géométrique.
\begin{enumerate}
\item (*) Calculer $\sum_{i=3}^{n}i$, $n\in\Nn\setminus\{0,1,2\}$, $\sum_{i=1}^{n}(2i-1)$, $n\in\Nn^*$, et
$\sum_{k=4}^{n+1}(3k+7)$, $n\in\Nn\setminus\{0,1,2\}$.

\item (*) Calculer le nombre $1,1111...=\lim_{n\rightarrow +\infty}1,\underbrace{11...1}_n$ et le nombre
$0,9999...=\lim_{n\rightarrow +\infty}0,\underbrace{99...9}_n$.

\item (*) Calculer $\underbrace{1-1+1-...+(-1)^{n-1}}_n$,
$n\in\Nn^*$.

\item (*) Calculer
$\frac{1}{2}+\frac{1}{4}+\frac{1}{8}+...=\lim_{n\rightarrow +\infty}\sum_{k=1}^{n}\frac{1}{2^k}$.

\item (**) Calculer
$\sum_{k=0}^{n}\cos\frac{k\pi}{2}$, $n\in\Nn$.

\item (**) Soient $n\in\Nn$ et $\theta\in\Rr$. Calculer $\sum_{k=0}^{n}\cos(k\theta)$ et
$\sum_{k=0}^{n}\sin(k\theta)$.

\item (***) Pour $x\in[0,1]$ et $n\in\Nn^*$, on pose
$S_n=\sum_{k=1}^{n}(-1)^{k-1}\frac{x^k}{k}$. Déterminer$\lim_{n\rightarrow +\infty}S_n$.

\item (**) On pose $u_0=1$ et, pour $n\in\Nn$, $u_{n+1}=2u_n-3$.
\begin{enumerate}
\item Calculer la suite $(u_n-3)_{n\in\Nn}$.
\item Calculer $\sum_{k=0}^{n}u_k$.
\end{enumerate}
\end{enumerate}
\finenonce{005142}


\finexercice

\exercice{5150, rouget, 2010/06/30}
\enonce{005150}{***I Inégalité de \textsc{Cauchy}-\textsc{Schwarz}}
Soient $n\in\Nn^*$ et $a_1$, $a_2$,..., $a_n$, $b_1$, $b_2$,..., $b_n$, $2n$ réels. Montrer que

$$|\sum_{k=1}^{n}a_kb_k|\leq\sum_{k=1}^{n}|a_k|.|b_k|\leq\sqrt{\sum_{k=1}^{n}a_k^2}\sqrt{\sum_{k=1}^{n}b_k^2}.$$

(Indication. Considérer le polynôme $f(x)=\sum_{k=1}^{n}(a_k+b_kx)^2$, développer puis ordonner suivant les puissances
décroissantes puis utiliser, dans le cas général, les connaissances sur le second degré). Retrouver alors le résultat de l'exercice \ref{exo:suprou4bis}.
\finenonce{005150}


\finexercice

\exercice{5458, rouget, 2010/07/10}
\enonce{005458}{**}
Montrer que  $\sum_{k=1}^{n}\sin\frac{k}{n}=\frac{1}{2}+\frac{1}{2n}+o(\frac{1}{n})$.
\finenonce{005458}


\finexercice


\exercice{5688, rouget, 2010/10/16}
\enonce{005688}{}
\label{ex:rou1ter}
Nature de la série de terme général 

\begin{center}
\begin{tabular}{llll}
\textbf{1) (*)} $\ln\left(\frac{n^2+n+1}{n^2+n-1}\right)$&\textbf{2) (*)}  $\frac{1}{n+(-1)^n\sqrt{n}}$&\textbf{3) (**)} $\left(\frac{n+3}{2n+1}\right)^{\ln n}$  &\textbf{4) (**)} $\frac{1}{\ln(n)\ln(\ch n)}$\\
\textbf{5) (**)} $\Arccos\sqrt[3]{1-\frac{1}{n^2}}$&\textbf{6) (*)} $\frac{n^2}{(n-1)!}$&\textbf{7)} $\left(\cos\frac{1}{\sqrt{n}}\right)^n-\frac{1}{\sqrt{e}}$&\textbf{8) (**)} $\ln\left(\frac{2}{\pi}\Arctan\frac{n^2+1}{n}\right)$\\
\textbf{9) (*)} $\int_{0}^{\pi/2}\frac{\cos^2x}{n^2+\cos^2x}\;dx$&  
\textbf{10) (**)} $n^{-\sqrt{2}\sin(\frac{\pi}{4}+\frac{1}{n})}$&
\textbf{11) (**)} $e-\left(1+\frac{1}{n}\right)^n$
\end{tabular}
\end{center}
\finenonce{005688}


\finexercice


\exercice{5689, rouget, 2010/10/16}
\enonce{005689}{}
Nature de la série de terme général 

\textbf{1) (***)} $\sqrt[4]{n^4+2n^2}-\sqrt[3]{P(n)}$  où $P$ est un polynôme. \qquad\textbf{2) (**)} $\frac{1}{n^\alpha}S(n)$ où $S(n) =\sum_{p=2}^{+\infty}\frac{1}{p^n}$. 

\textbf{3) (**)} $u_n$ où $\forall n\in\Nn^*$, $u_n=\frac{1}{n}e^{-u_{n-1}}$.

\textbf{4) (****)} $u_n=\frac{1}{p_n}$ où $p_n$ est le $n$-ème nombre premier

(indication : considérer 
$\sum_{n=1}^{N}\ln\left(\frac{1}{1-\frac{1}{p_n}}\right)=\sum_{n=1}^{N}\ln(1+p_n+p_n^2+\ldots)$).

\textbf{5) (***)} $u_n=\frac{1}{n(c(n))^\alpha}$  où $c(n)$ est le nombre de chiffres de $n$ en 
base $10$.

\textbf{6) (*)} $\frac{\left(\prod_{k=2}^{n}\ln k\right)^a}{(n!)^b}$ $a > 0$ et $b> 0$.\qquad \textbf{7) (**)} $\Arctan\left(\left(1+\frac{1}{n}\right)^a\right) -\Arctan\left(\left(1-\frac{1}{n}\right)^a\right)$.

\textbf{8) (**)} $\frac{1}{n^\alpha}\sum_{k=1}^{n}k^{3/2}$.\qquad \textbf{9) (***) } $\left(\prod_{k=1}^{n}\left(1+\frac{k}{n^\alpha}\right)\right)-1$.

\finenonce{005689}


\finexercice


\exercice{5690, rouget, 2010/10/16}
\enonce{005690}{}
Nature de la série de terme général 

\textbf{1) (***)} $\sin\left(\frac{\pi n^2}{n+1}\right)$\qquad\textbf{2) (**)} $\frac{(-1)^n}{n+(-1)^{n-1}}$\qquad\textbf{3) (**)} $\ln\left(1+\frac{(-1)^n}{\sqrt{n}}\right)$\qquad\textbf{4) (***)} $\frac{e^{in\alpha}}{n}$, $\frac{\cos(n\alpha)}{n}$ et $\frac{\sin(n\alpha)}{n}$

\textbf{5) (**)} $(-1)^n\frac{\ln n}{n}$\qquad\item  $(-1)^n\frac{P(n)}{Q(n)}$ où $P$ et $Q$ sont deux polynômes non nuls\qquad

\textbf{7) (****)} $(\sin(n!\pi e))^p$ $p$ entier naturel non nul.
\finenonce{005690}


\finexercice


\exercice{5691, rouget, 2010/10/16}
\enonce{005691}{}
\label{ex:rou4}
Calculer les sommes des séries suivantes après avoir vérifié leur convergence.

\begin{center}
\begin{tabular}{lll}
\textbf{1) (**)} $\sum_{n=0}^{+\infty}\frac{n+1}{3^n}$&\textbf{2) (**)} $\sum_{n=3}^{+\infty}\frac{2n-1}{n^3-4n}$&\textbf{3) (***)} $\sum_{n=0}^{+\infty}\frac{1}{(3n)!}$\\
\textbf{4) (*)} $\sum_{n=2}^{+\infty}\left(\frac{1}{\sqrt{n-1}}+\frac{1}{\sqrt{n+1}}-\frac{2}{\sqrt{n}}\right)$&
\textbf{5) (**)} $\sum_{n=2}^{+\infty}\ln\left(1+\frac{(-1)^n}{n}\right)$&\textbf{6) (***)} $\sum_{n=0}^{+\infty}\ln\left(\cos\frac{a}{2^n}\right)$  $a\in\left]0,\frac{\pi}{2}\right[$\\
textbf{7)}  $\sum_{n=0}^{+\infty}\frac{\tanh\frac{a}{2^n}}{2^n}$
\end{tabular}
\end{center}
\finenonce{005691}


\finexercice


\exercice{5692, rouget, 2010/10/16}
\enonce{005692}{*** I}
Soit $(u_n)_{n\in\Nn}$ une suite décroissante de nombres réels strictement positifs telle que la série de terme général $u_n$ converge. Montrer que $u_n\underset{n\rightarrow+\infty}{=}o\left(\frac{1}{n}\right)$. Trouver un exemple de suite $(u_n)_{n\in\Nn}$ de réels strictement positifs telle que la série de terme général $u_n$ converge mais telle que la suite de terme général $nu_n$ ne tende pas vers $0$.
\finenonce{005692}


\finexercice


\exercice{5693, rouget, 2010/10/16}
\enonce{005693}{***}
Soit $\sigma$ une injection de $\Nn^*$ dans lui-même. Montrer que la série de terme général $\frac{\sigma(n)}{n^2}$ diverge.

\finenonce{005693}


\finexercice


\exercice{5694, rouget, 2010/10/16}
\enonce{005694}{**}
Soit $(u_n)_{n\in\Nn}$ une suite de réels strictement positifs. Montrer que les séries de termes généraux $u_n$, $\frac{u_n}{1+u_n}$, $\ln(1+u_n)$ et $\int_{0}^{u_n}\frac{dx}{1+x^e}$  sont de mêmes natures.
\finenonce{005694}


\finexercice


\exercice{5695, rouget, 2010/10/16}
\enonce{005695}{***}
Trouver un développement limité à l'ordre $4$ quand $n$ tend vers l'infini de $\left(e-\sum_{k=0}^{n}\frac{1}{k!}\right)\times(n+1)!$.
\finenonce{005695}


\finexercice


\exercice{5696, rouget, 2010/10/16}
\enonce{005696}{***}
Nature de la série de terme général $u_n=\sin\left(\pi(2+\sqrt{3})^n\right)$.
\finenonce{005696}


\finexercice


\exercice{5700, rouget, 2010/10/16}
\enonce{005700}{**}
Soit $\alpha\in\Rr$. Nature de la série de terme général $u_n=\frac{1+(-1)^nn^\alpha}{n^{2\alpha}}$, $n\geqslant1$.  
\finenonce{005700}


\finexercice


\exercice{5701, rouget, 2010/10/16}
\enonce{005701}{****}
On sait que $1-\frac{1}{2}+\frac{1}{3}-\frac{1}{4}+\ldots=\ln2$.

A partir de la série précédente, on construit une nouvelle série en prenant $p$ termes positifs, $q$ termes négatifs, $p$ termes positifs ... (Par exemple pour $p = 3$ et $q = 2$, on s'intéresse à $1+\frac{1}{3}+\frac{1}{5}-\frac{1}{2}-\frac{1}{4}+\frac{1}{7}+\frac{1}{9}+\frac{1}{11}-\frac{1}{6}-\frac{1}{8}+\ldots$). Convergence et somme de cette série.
\finenonce{005701}


\finexercice


\exercice{5702, rouget, 2010/10/16}
\enonce{005702}{***}
Nature de la série de terme général $u_n=\sum_{k=1}^{n-1}\frac{1}{(k(n-k))^\alpha}$.  
\finenonce{005702}


\finexercice


\exercice{5703, rouget, 2010/10/16}
\enonce{005703}{}
Convergence et somme éventuelle de la série de terme général

\begin{center}
\textbf{1) (**)} $u_n =\frac{2n^3-3n^2+1}{(n+3)!}$\qquad\textbf{2) (***)} $u_n=\frac{n!}{(a+1)(a+2)\ldots(a+n)}$, $n\geqslant1$, $a\in\Rr^{+*}$ donné.
\end{center} 
\finenonce{005703}


\finexercice


\exercice{5706, rouget, 2010/10/16}
\enonce{005706}{*}
Nature de la série de terme général $u_n=\sum_{k=1}^{n}\frac{1}{(n+k)^p}$, $p\in]0,+\infty[$.
\finenonce{005706}


\finexercice


\exercice{5707, rouget, 2010/10/16}
\enonce{005707}{*** I}
Développement limité à l'ordre $4$ de $\sum_{k=n+1}^{+\infty}\frac{1}{k^2}$ quand $n$ tend vers l'infini.
\finenonce{005707}


\finexercice


\exercice{5708, rouget, 2010/10/16}
\enonce{005708}{}
Partie principale quand $n$ tend vers $+\infty$ de

\begin{center}
\textbf{1) (***)} $\sum_{p=n+1}^{+\infty}(-1)^p\frac{\ln p}{p}$\qquad\textbf{2) (**)} $\sum_{p=1}^{n}p^p$.
\end{center}
\finenonce{005708}


\finexercice


\exercice{5709, rouget, 2010/10/16}
\enonce{005709}{***}
Soit $p\in\Nn^*$, calculer $\sum_{p\in\Nn^*}^{}\left(\sum_{n\in\Nn^*,\;n\neq p}^{}\frac{1}{n^2-p^2}\right)$ et $\sum_{n\in\Nn^*}^{}\left(\sum_{p\in\Nn^*,\;p\neq n}^{}\frac{1}{n^2-p^2}\right)$. Que peut-on en déduire ?
\finenonce{005709}


\finexercice


\exercice{5710, rouget, 2010/10/16}
\enonce{005710}{**}
 Calculer $\sum_{n=0}^{+\infty}\frac{(-1)^n}{3n+1}$.
\finenonce{005710}


\finexercice


\exercice{5711, rouget, 2010/10/16}
\enonce{005711}{****}
Soient $(u_n)_{n\geqslant1}$ une suite réelle. Pour $n\geqslant 1$, on pose $v_n=\frac{u_1+\ldots+u_n}{n}$. Montrer que si la série de terme général $(u_n)^2$ converge alors la série de terme général $(v_n)^2$ converge et que $\sum_{n=1}^{+\infty}(v_n)^2\leqslant4\sum_{n=1}^{+\infty}(u_n)^2$ (indication : majorer $v_n^2 - 2u_nv_n$).
\finenonce{005711}


\finexercice


\exercice{5712, rouget, 2010/10/16}
\enonce{005712}{***}
Convergence et somme de la série de terme général $u_n=\frac{\pi}{4}-\sum_{k=0}^{n}\frac{(-1)^k}{2k+1}$, $n\geqslant 0$.
\finenonce{005712}


\finexercice

\section{123.01 Continuité : théorie}


\exercice{639, bodin, 1998/09/01}
\video{RMlTCZ6T9wc}
\enonce{000639}{}
Soit $I$ un intervalle ouvert de $\Rr$, $f$ et $g$ deux fonctions d\'efinies sur $I$.
\begin{enumerate}
        \item Soit $a\in I$. Donner une raison pour laquelle :
$$\left( \lim_{x\rightarrow a}f(x)=f(a) \right) \Rightarrow
\left( \lim_{x\rightarrow a} |f(x)|=|f(a)| \right). $$
    \item On suppose que $f$ et $g$ sont continues sur $I$. En utilisant
l'implication d\'emontr\'ee ci-dessus, la relation $\sup(f,g)=\frac{1}{2}(f+g+|f-g|)$,
et les propri\'et\'es des fonctions continues, montrer que la fonction $\sup(f,g)$
est continue sur $I$.
\end{enumerate}
\finenonce{000639} 


\finexercice

\exercice{640, bodin, 1998/09/01}

\enonce{000640}{}
 Soit $f$ une fonction de $\lbrack a,b \rbrack$ dans $\lbrack a,b\rbrack$
telle que pour tout $x$ et $x'$ ($x\neq x'$) de $\lbrack a,b \rbrack$ on ait :
$\ \ \vert f(x)-f(x')\vert<\vert x-x'\vert.$
\begin{enumerate}
    \item  Montrer que $f$ est continue sur $\lbrack a,b\rbrack$.
    \item  Montrer que l'\'equation $f(x)=x$ admet une et une seule solution
\mbox{dans $\lbrack a,b\rbrack$.}
 (On pourra introduire la fonction: $x\mapsto g(x)=f(x)-x$).
\end{enumerate}
\finenonce{000640}



\finexercice



\exercice{641, bodin, 1998/09/01}

\enonce{000641}{}
\begin{enumerate}
    \item Soit $f$ une fonction continue sur $]a,b[$ telle que
$f(]a,b[)\subset [a,b]$. Montrer, par consid\'eration de $\phi(x)=f(x)-x$, qu'il
existe $c$ dans $[a,b]$ tel que $f(c)=c$.
    \item Soit $f$ une fonction continue sur $[0,1]$ telle que $f(0)=f(1)$.
Montrer qu'il existe $c$ dans $[0,\frac{1}{2}]$ tel que $f(c) = f(c+\frac{1}{2})$.
    \item Un mobile parcours, \`a vitesse continue, une distance $d$ en
une unit\'e de temps. Montrer qu'il existe un intervalle d'une demi-unit\'e de temps
 pendant lequel il parcourt une distance $\frac{d}{2}$.
\end{enumerate}
\finenonce{000641}



\finexercice



\exercice{642, bodin, 1998/09/01}
\video{Uh9semrdJvk}
\enonce{000642}{}
 Soit $f:[a,b]\longrightarrow\R$ une fonction continue telle que
$f(a)=f(b)$. Montrer que la fonction $g(t)=f(t+\frac{b-a}{2})-f(t)$ s'annule en au moins un point
de $[a,\frac{a+b}{2}]$.

\emph{Application :} une personne parcourt 4 km en 1 heure.
Montrer qu'il existe un intervalle de 30 mn pendant lequel elle parcourt exactement 2 km.

\finenonce{000642}



\finexercice



\exercice{643, ridde, 1999/11/01}

\enonce{000643}{}
Soit $f : \Rr \rightarrow \Rr$ continue telle que $\lim\limits_{ -\infty}f =  -\infty$
et $\lim\limits_{  + \infty}f =   + \infty$. Montrer que $f$ s'annule. Appliquer
ceci aux polynômes de degr\'e impair.

\finenonce{000643}




\finexercice



\exercice{644, ridde, 1999/11/01}

\enonce{000644}{}
Soit $f : \Rr \rightarrow \Rr^ + $ continue telle que $f (0) = 1$,
$\lim\limits_{ -\infty}f = 0$ et $\lim\limits_{  + \infty}f = 0$.
\begin{enumerate}
\item  Montrer qu'il existe $a>0$ tel que si $\left|x\right|>a$ alors $f (x) \leq \frac 12$.
\item Montrer que $f$ est born\'ee et poss\`ede un maximum.
\end{enumerate}
\finenonce{000644}



\finexercice



\exercice{645, ridde, 1999/11/01}
\video{ZeT60IE0owY}
\enonce{000645}{}
Soient $I$ un intervalle de $\Rr$ et $f : I \rightarrow \Rr$ continue, telle que pour chaque $x \in I$, $f (x)^2 = 1$.
Montrer que $f = 1$  ou $f = -1$.
\finenonce{000645} 


\finexercice


\exercice{646, ridde, 1999/11/01}
\video{gC6TqV9IOPg}
\enonce{000646}{}
Soit $f : \Rr^ + \rightarrow \Rr$ continue admettant une limite finie en $ + \infty$.
Montrer que $f$ est born\'ee. Atteint-elle ses bornes ?
\finenonce{000646} 


\finexercice

\exercice{647, ridde, 1999/11/01}

\enonce{000647}{}
Soient $f$ et $g$ continues sur $[0, 1]$ telles que $\forall x \in [0, 1] \, \,
f (x )<g (x)$. Montrer qu'il existe $m>0$ tel que $\forall x \in [0, 1] \, \,
f (x ) + m <g (x)$.
\finenonce{000647}



\finexercice



\exercice{648, ridde, 1999/11/01}

\enonce{000648}{}
Soit $f$  croissante sur $[a, b]$ et prenant toute valeur entre $f (a)$
et $f (b)$. Montrer que $f$ est continue.
\finenonce{000648}



\finexercice



\exercice{649, gourio, 2001/09/01}

\enonce{000649}{}
Soit $f:\Rr\rightarrow \Rr$ continue en $0$ telle que $\forall x\in
\Rr,f(x)=f(2x). $ Montrer que $f$ est constante.
\finenonce{000649}



\finexercice



\exercice{650, gourio, 2001/09/01}

\enonce{000650}{}
Soit $f$ p\'{e}riodique croissante. Que dire de $f$ ?
\finenonce{000650}



\finexercice



\exercice{651, gourio, 2001/09/01}

\enonce{000651}{}
Donner un exemple de fonction continue sur $[0,1]$ non lipschitzienne, puis
de fonction continue en un seul point, puis de fonction discontinue sur les
rationnels et continue sur les irrationnels, enfin de fonction continue
telle que $f(x)\in \Rr\setminus\Qq$ si $x\in \Rr\setminus\Qq$ ou si $x=0,$ et $f(x)\in \Qq$ si $x\in
\Qq\setminus\{0\}. $ Une fonction telle que $\forall x\in \Rr,\lim\limits_{h\rightarrow
0}f(x+h)-f(x-h)=0$ est-elle continue sur $\Rr$ ? Donner un exemple de bijection
de $[0,1]$ sur $[0,1]$ discontinue en tout point.
\finenonce{000651}



\finexercice



\exercice{652, gourio, 2001/09/01}

\enonce{000652}{}
Soit $f$ continue sur $\Rr$ admettant $1$ et $\sqrt{2}$ pour p\'{e}riodes.
Que dire de $f$ ?
\finenonce{000652}



\finexercice



\exercice{653, gourio, 2001/09/01}
\video{ZCT1KDuJcsI}
\enonce{000653}{}
Soit $f:[0,1]\rightarrow [0,1]$ croissante, montrer qu'elle a un point fixe.

\emph{Indication} : \'{e}tudier
$$E=\big\{x\in [0,1] \mid \forall t\in [0,x],f(t)>t\big\}.$$
\finenonce{000653} 


\finexercice


\exercice{654, gourio, 2001/09/01}

\enonce{000654}{}
Soit $f:\Rr^{+*}\rightarrow \Rr$ croissante telle que $x\rightarrow \frac{f(x)}{x}$
 soit d\'{e}croissante ; montrer que $f$ est continue sur $\Rr^{+*}.$
\finenonce{000654}



\finexercice



\exercice{655, gourio, 2001/09/01}

\enonce{000655}{}
Soit $f:\Rr^{+*}\rightarrow \Rr$ une fonction v\'{e}rifiant :
$$\forall x\in \Rr^{+*},f(x)e^{f(x)}=x.$$
Donner les variations de $f$ puis comparer $f$ et $\ln$ au voisinage de $+\infty.$
\finenonce{000655}



\finexercice



\exercice{656, gourio, 2001/09/01}

\enonce{000656}{}
Soit $f:{\Rr}^{+}\rightarrow {\Rr}$ croissante. Construire
$g:{\Rr}^{+}\rightarrow {\Rr}$ continue telle que $f\leq g.$
\finenonce{000656}



\finexercice



\exercice{657, gourio, 2001/09/01}

\enonce{000657}{}
Donner un exemple d'application $f:{\Rr}\rightarrow {\Rr}$ non constante
telle que :
$$\forall x\in {\Rr},  f(x)=f(x^{2}). $$
On suppose $f$ continue en $0$ et en $1$, montrer que $f$ est constante.
\finenonce{000657}



\finexercice



\exercice{658, gourio, 2001/09/01}

\enonce{000658}{}
Soit $f:[0,1]\rightarrow [0,1]$ continue. Montrer que :
$$\forall n\in \Nn^{*},\exists  a_{n}\in [0,1],f(a_{n})=a_{n}^{n}. $$
On suppose $f$ strictement d\'{e}croissante.\ Montrer que $a_{n}$ est unique
et \'{e}tudier la suite $(a_{n})_{n\in \Nn^{*}}.$
\finenonce{000658}



\finexercice



\exercice{659, gourio, 2001/09/01}

\enonce{000659}{}
Existe-t-il une bijection continue de $[0,1[$ sur ${\Rr }$ ?
\finenonce{000659}



\finexercice



\exercice{660, gourio, 2001/09/01}
\enonce{000660}{}
Soit $f:[0,1]\rightarrow [0,1]$ continue telle que $f\circ f=f$.
On note $E_{f}=\{x\in [0,1]|f(x)=x\}$. Montrer que $E_{f}\neq \varnothing$
puis que c'est un intervalle de ${\Rr}$.

Trouver toutes les fonctions $f:[0,1]\rightarrow [0,1]$ continues telles que $f \circ f = f$.
\finenonce{000660} 


\finexercice


\exercice{661, gourio, 2001/09/01}

\enonce{000661}{}
Soit $f:[0,1]\rightarrow {\Rr}$ continue, \'{e}valuer :
$$\lim\limits_{n\rightarrow \infty }\sum\limits_{k=1}^{n}(-1)^{k}f\left(\frac{k}{n}\right). $$
\finenonce{000661}



\finexercice



\exercice{662, gourio, 2001/09/01}
\enonce{000662}{}
Une fonction qui v\'{e}rifie la propri\'{e}t\'{e} des valeurs
interm\'{e}diaires est-elle n\'{e}cessairement continue ?
\finenonce{000662} 


\finexercice

\exercice{663, gourio, 2001/09/01}

\enonce{000663}{}
Soit $f$ uniform\'{e}ment continue sur ${\Rr}^{+}$ telle que $\forall
x\geq 0, $ la suite $\left( f(xn)\right) _{n\in \Nn}  $ tend vers 0
quand $n\rightarrow \infty . $ Montrer $\lim\limits_{x\rightarrow \infty
}f(x)=0.$
\finenonce{000663}



\finexercice



\exercice{664, gourio, 2001/09/01}

\enonce{000664}{}
Soit $f\in C({\Rr}^{+},{\Rr})$ admettant une limite finie en $+\infty $,
montrer qu'alors $f$ est uniform\'{e}ment continue sur ${\Rr}^{+}.$
\finenonce{000664}



\finexercice



\exercice{665, gourio, 2001/09/01}

\enonce{000665}{}
Soit $f$ continue sur $[a,b]$, montrer :
$$\forall \epsilon >0,\exists  k\in {\Rr},\forall (x,y)\in
[a,b]^{2},\left| f(x)-f(y)\right| \leq k\left| x-y\right| +\epsilon . $$
\finenonce{000665}



\finexercice



\exercice{666, gourio, 2001/09/01}

\enonce{000666}{}
Soit $(f,g)\in C([0,1],[0,1])^{2},$ tel que : $fg=gf.$
On veut montrer que $f-g$ s'annulle par deux m\'{e}thodes :
\begin{itemize}
\item  par l'absurde, utiliser le fait que $(f-g)([0,1])$ est un segment ne
contenant pas 0.
\item par l'absurde, en examinant, si  $f-g>0$ par exemple, $\min\{x\in
[0,1]|f(x)=x\}.$
\end{itemize}
Le r\'{e}sultat subsiste-t-il si l'on remplace $[0,1]$ par ${\Rr }$ ?
\finenonce{000666}



\finexercice



\exercice{667, gourio, 2001/09/01}

\enonce{000667}{}
Soit $f:[0,1]\rightarrow {\Rr}$ continue, telle que $f(0)=f(1).$
Montrer que :
$$\forall n\in \Nn^{*},\exists x_{n}\in \left[ 0,1-\frac{1}{n}\right]
,f\left( x_{n}+\frac{1}{n}\right) =f\left( x_{n}\right) . $$
\finenonce{000667}



\finexercice



\exercice{668, gourio, 2001/09/01}

\enonce{000668}{}
Soit $f$ continue de ${\Rr}$ dans ${\Rr}$, montrer que :
$\lim\limits_{\left| x\right| \rightarrow \infty }\left| f(x)\right|
=+\infty \Leftrightarrow  $
l'image r\'{e}ciproque de toute partie
born\'{e}e est born\'{e}e.
\finenonce{000668}



\finexercice



\exercice{669, monthub, 2001/11/01}
\enonce{000669}{}
Soit $f : [a,b] \to \R$ une fonction continue. On veut d\'emontrer
que
$$ \sup_{a<x<b} f(x) =  \sup_{a\leq x \leq b} f(x).$$
\begin{enumerate}
\item Montrer que $$ \sup_{a<x<b} f(x) \leq  \sup_{a\leq x \leq b} f(x).$$

\item Soit $x_0 \in [a,b]$ tel que $f(x_0)=\sup_{a\leq x \leq b} f(x)$.
Montrer que $f(x_0)=\sup_{a< x< b} f(x)$ en distinguant les trois
cas : $x_0=a, x_0=b, x_0\in ]a,b[$.

\item Soit $g:[0,1] \to \R$ la fonction d\'efinie par $g(x)=0$ si $x\in [0,1[$
  et $g(x)=1$ si $x=1$. Montrer que $$ \sup_{0<x<1} g(x) \neq  \sup_{0\leq x \leq 1} g(x).$$
Quelle hypoth\`ese  est essentielle dans la propri\'et\'e d\'emontr\'ee
auparavant ?
\end{enumerate}
\finenonce{000669} 


\finexercice


\exercice{3845, quercia, 2010/03/11}
\enonce{003845}{Fonction périodique}

Soit $f : \R \to \R$ et $T > 0$. On suppose que $f$ est $T$-périodique cad :
$\forall\ x \in \R,\ f(x+T) = f(x)$.

\begin{enumerate}
  \item Si $f$ possède une limite en $+\infty$, montrer que $f$ est constante.
  \item Si $f$ est continue non constante, montrer que $f$ a une plus petite période.
  \item Si $f$ est continue, montrer que $f$ est bornée et atteint ses bornes.
\end{enumerate}
\finenonce{003845}



\finexercice


\exercice{3846, quercia, 2010/03/11}
\enonce{003846}{Fonction ayant des limites à l'infini}

Soit $f : {[0,+\infty[} \to \R$ une fonction continue ayant une limite
finie en $+\infty$.

\begin{enumerate}
  \item Montrer que $f$ est bornée.
  \item Montrer que $f$ admet un maximum ou un minimum absolu, mais pas nécéssairement les deux.
  \item Montrer que $f$ est uniformément continue.
\end{enumerate}
\finenonce{003846}



\finexercice


\exercice{3847, quercia, 2010/03/11}
\enonce{003847}{Permutation de décimales}

Pour $x \in {[0,1[}$, on note $x = \sum_{k=1}^\infty \frac{x_k}{10^k}$ le
développement décimal propre de $x$.

\begin{enumerate}
  \item   Soit $f : {[0,1[} \to {[0,1[}$ définie par :
      $f(x) = \sum_{k=1}^\infty \frac{x_{k+1}}{10^k}$. Montrer que $f$ est continue
      par morceaux.

  \item   Soit $g : {[0,1[} \to {[0,1[}$ définie par :
      $g(x) = \sum_{k=1}^\infty \left( \frac{x_{2k}}{10^{2k-1}} +
                                        \frac{x_{2k-1}}{10^{2k}} \right)$.
      Déterminer les points où $g$ est continue.
\end{enumerate}
\finenonce{003847}



\finexercice

\exercice{3849, quercia, 2010/03/11}
\enonce{003849}{$\max(f,g)$}
               

Soient ${f,g} : {\R} \to {\R}$ deux fonctions continues.
On pose pour $x \in \R$ : $h(x) = \max\bigl(f(x),g(x)\bigr)$.
Montrer que $h$ est continue.

\finenonce{003849}



\finexercice


\exercice{3850, quercia, 2010/03/11}
\enonce{003850}{Prolongement d'inégalités}

\begin{enumerate}
  \item Soient ${f,g} :  \R \to \R$ continues telles que :
    $\forall\ x,y \in \Q,\ f(x) < g(x)$.
  \begin{enumerate}
    \item Montrer que $f\le g$.
    \item Montrer qu'on n'a pas nécéssairement : $\forall\ x,y \in \R,\ f(x) < g(x)$.
  \end{enumerate}
  \item Soit $f : \R \to \R$ continue dont la restriction à $\Q$ est strictement
    croissante. Montrer que $f$ est strictement croissante.

\end{enumerate}
\finenonce{003850}



\finexercice


\exercice{3851, quercia, 2010/03/11}
\enonce{003851}{Étude d'un sup}

Soit $f : \R \to \R$ uniformément continue.
On pose $g(x) = \sup\bigl(f([x,x+1])\bigr)$.
Montrer que $g$ est continue.

Même question en supposant seulement $f$ continue.
\finenonce{003851}



\finexercice


\exercice{3852, quercia, 2010/03/11}
\enonce{003852}{Weierstrass}

Soient ${f,g}: {[a,b]}\to \R$ continues.
On pose $h(t) = \sup\{f(x)+tg(x) \text{ tq } x \in {[a,b]}\}$.

Montrer que $h$ est continue.

\finenonce{003852}



\finexercice


\exercice{3853, quercia, 2010/03/11}
\enonce{003853}{Weierstrass}
Soit $f : {[a,b]} \to \R$ continue.
Montrer que $\mathop{\sup}\limits_{[a,b]} f = \mathop{\sup}\limits_{]a,b[} f$.
\finenonce{003853}



\finexercice


\exercice{3854, quercia, 2010/03/11}
\enonce{003854}{Weierstrass}
Soient ${f,g} : {[a,b]} \to \R$ continues.
On suppose que : $\forall\ x \in {[a,b]},\ f(x) > g(x) > 0$.
Montrer qu'il existe $k > 1$ tel que $f > kg$.

\finenonce{003854}



\finexercice


\exercice{3855, quercia, 2010/03/11}
\enonce{003855}{TVI à l'infini}

Soit $f : {[0,+\infty[} \to \R$ continue ayant une limite
$\ell \in \overline\R$ en $+\infty$.
Montrer que $f$ prend toute valeur comprise entre $f(0)$ et $\ell$ ($\ell$ exclu).

\finenonce{003855}



\finexercice


\exercice{3856, quercia, 2010/03/11}
\enonce{003856}{$f(x)=g(x)$}

\begin{enumerate}
  \item Soit $f : {[0,1]} \to {[0,1]}$ continue. Montrer qu'il existe $x \in {[0,1]}$
tel que $f(x) = x$.

  \item Soient ${f,g} : {[0,1]} \to {[0,1]}$ continues telles que $f\circ g = g\circ f$.
Montrer qu'il existe $x \in {[0,1]}$
tel que $f(x) = g(x)$ (on pourra s'intéresser aux points fixes de~$f$).

\end{enumerate}
\finenonce{003856}



\finexercice


\exercice{3857, quercia, 2010/03/11}
\enonce{003857}{$f$ continue décroissante $ \Rightarrow $ point fixe}

Soit $f : \R \to \R$ continue décroissante. Montrer qu'il existe un unique réel $x$
tel que $f(x) = x$.
\finenonce{003857}



\finexercice


\exercice{3858, quercia, 2010/03/11}
\enonce{003858}{Mines MP 2002}

Soit $f : \R \to \R$ continue telle qu'existe $a$ vérifiant
$f\circ f(a)=a$. $f$ a-t-elle des points fixes~? Généraliser.

\finenonce{003858}



\finexercice


\exercice{3859, quercia, 2010/03/11}
\enonce{003859}{Cordes de longueur $1/n$}

Soit $f : {[0,1]} \to \R$ continue telle que $f(0) = f(1)$.

\begin{enumerate}
  \item   Montrer qu'il existe $x \in \left[0,\frac 12\right]$ tel que
      $f(x) = f\left(x+\frac 12\right)$.
  \item   Pour $n \in \N, n \ge 2$, montrer qu'il existe
      $x \in \left[0,1-\frac 1n\right]$ tel que $f(x) = f\left(x+\frac 1n\right)$.
  \item   Trouver une fonction $f$ telle que : $\forall\ x\in \left[0,\frac35\right]$,
      $f(x) \ne f(x+\frac25)$.
  \item   Montrer qu'il existe $a > 0$ tel que : $\forall\ b \in {]0,a]},\ \exists\ x
      \in {[0, 1-b]} \text{ tq } f(x) = f(x+b)$.

\end{enumerate}
\finenonce{003859}



\finexercice


\exercice{3860, quercia, 2010/03/11}
\enonce{003860}{}

Soient ${f,g} : {[a,b]}  \to \R$ continues. On suppose que :
$\forall\ x \in {[a,b]},\ \exists\ y \in [a,b] \text{ tq } f(x) = g(y)$.

Montrer qu'il existe $x \in {[a,b]}$ tel que $f(x) = g(x)$.

\finenonce{003860}



\finexercice


\exercice{3861, quercia, 2010/03/11}
\enonce{003861}{TVI $+$ injective $ \Rightarrow $ continue}

Soit $f : \R \to \R$. On dit que $f$ vérifie la propriété des valeurs intermédiaires
si :
$$\forall\ a,b\in\R \text{ avec } a < b, \forall\ y \text{ compris entre }
f(a) \text{ et } f(b),\ \exists\ x \in {[a,b]} \text{ tq } f(x) = y.$$

\begin{enumerate}
  \item Montrer que si $f$ vérifie la propriété des valeurs intermédiaires et est
    injective, alors elle est continue.

  \item Trouver une fonction discontinue ayant la propriété des valeurs
    intermédiaires.
\end{enumerate}
\finenonce{003861}



\finexercice


\exercice{3862, quercia, 2010/03/11}
\enonce{003862}{$f$ uc $ \Rightarrow  f$(intervalle borné) $=$ intervalle borné}

Soit $I$ un intervalle borné et $f : I \to \R$ uniformément continue.
Montrer que $f(I)$ est un intervalle borné.
\finenonce{003862}



\finexercice


\exercice{3863, quercia, 2010/03/11}
\enonce{003863}{$f$ uc $ \Rightarrow  |f(x)| \le a + b|x|$}

Soit $f : \R  \to \R$ uniformément continue.
Montrer qu'il existe $a,b \in \R$ tels que : $\forall\ x \in \R,\ |f(x)| \le a + b|x|$.

(prendre $\varepsilon = 1$ et majorer $|f(x) - f(0)|$)

\finenonce{003863}



\finexercice


\exercice{3864, quercia, 2010/03/11}
\enonce{003864}{Composition}

Soient $f : D \to \R$ uniformément continue bornée et $g : \R \to \R$ continue.
Montrer que $g\circ f$ est uniformément continue.
\finenonce{003864}



\finexercice


\exercice{3865, quercia, 2010/03/11}
\enonce{003865}{$\sin(t^2)$}

Montrer que $t \mapsto \sin(t^2)$ n'est pas uniformément continue sur $\R$.

\finenonce{003865}



\finexercice


\exercice{3866, quercia, 2010/03/11}
\enonce{003866}{$f$ uc et $f(n) \to +\infty$}

Soit $f : \R \to \R$ uniformément continue telle que
$f(n) \to +\infty$ lorsque $n\to\infty$.
Montrer que $f(x) \to +\infty$ lorsque $x\to\infty$.


\finenonce{003866}



\finexercice


\exercice{3867, quercia, 2010/03/11}
\enonce{003867}{$f(x+y) = f(x) + f(y)$}

Soit $f : {\R} \to {\R}$ telle que : $\forall\ x,y \in \R,\ f(x+y) = f(x) + f(y)$.
On pose $a = f(1)$.

\begin{enumerate}
  \item Montrer que : $\forall\ x\in\Q$, $f(x) = ax$.
  \item On suppose $f$ continue. Montrer que : $\forall\ x \in \R$, $f(x) = ax$.
  \item On suppose que $f$ est bornée au voisinage de 0.
    Montrer que : $\forall\ x \in \R$, $f(x) = ax$.
\end{enumerate}
\finenonce{003867}



\finexercice


\exercice{3868, quercia, 2010/03/11}
\enonce{003868}{$f(x^2) = f(x)$}

Trouver toutes les fonctions $f : {[0,1]} \to {\R}$ continues telles
que : $\forall\ x \in {[0,1]},\ f(x^2) = f(x)$.

\finenonce{003868}



\finexercice


\exercice{3869, quercia, 2010/03/11}
\enonce{003869}{$f(x+y)f(x-y) = f^2(x)f^2(y)$}

Trouver toutes les fonctions $f : {\R} \to {\R}$ continues telles
que : $\forall\ x,y \in \R,\ f(x+y)f(x-y) = f^2(x)f^2(y)$.
\finenonce{003869}



\finexercice


\exercice{3870, quercia, 2010/03/11}
\enonce{003870}{Polytechnique MP$^*$ 2000}
Soit $f$ continue sur~$[a,b]$, à valeurs dans~$\R$, et $\delta$ un
réel positif. On note $\omega(\delta) = \sup\{|f(x)-f(y)|\text{ tq } |x-y|\le\delta\}$.
Montrer que $\omega(\delta)$ tend vers~$0$ quand $\delta$ tend vers~$0$,
puis que $\omega$ est continue.
\finenonce{003870}



\finexercice


\exercice{3871, quercia, 2010/03/11}
\enonce{003871}{Ensae MP$^*$ 2003}
Soient ${f,g} : {[0,1]} \to {[0,1]}$ continues telles que $f\circ g = g\circ f$.
Montrer qu'il existe $x\in[0,1]$ tel que $f(x) = g(x)$.

\finenonce{003871}



\finexercice


\exercice{3872, quercia, 2010/03/11}
\enonce{003872}{Plus grande fonction lipschitzienne minorant~$f$ (Ens Lyon MP$^*$ 2003)}
\ 
\begin{enumerate}
  \item Existe-t-il toujours~$\varphi$ lipschitzienne telle que $\varphi\le f$
    où $f : {\R^n} \to \R$ est une application continue donnée~?
    
  \item Soit~$k>0$. Trouver une CNS sur~$f$ pour qu'il existe $\varphi : {\R^n}\to \R$
    $k$-lipschitzienne minorant~$f$.
    
  \item On suppose cette CNS vérifiée pour~$k_0>0$. Montrer que si $k\ge k_0$
    alors il existe $\varphi_k$, $k$-lipschitzienne minorant~$f$ et maximale
    pour l'ordre usuel des fonctions.
    
\end{enumerate}
\finenonce{003872}



\finexercice


\exercice{3873, quercia, 2010/03/11}
\enonce{003873}{Suite $(f(nx))$ ENS Cachan MP$^*$ 2004}
Soit $f : \R \to \R$ une fonction uniformément continue telle que pour
tout~$x>0$ la suite $(f(nx))_{n\in\N}$ est convergente. Que peut-on dire
de~$f$~?

\finenonce{003873}



\finexercice


\exercice{5382, rouget, 2010/07/06}
\enonce{005382}{***I}
Soit $f$ une fonction réelle d'une variable réelle définie et continue sur un voisinage de $+\infty$. On suppose que la fonction $f(x+1)-f(x)$ admet dans $\Rr$ une limite $\ell$ quand $x$ tend vers $+\infty$. Etudier l'existence et la valeur eventuelle de $\lim_{x\rightarrow +\infty}\frac{f(x)}{x}$.
\finenonce{005382}


\finexercice


\exercice{5383, rouget, 2010/07/06}
\enonce{005383}{***}
Soit $f$ une fonction définie sur un voisinage de $0$ telle que $\lim_{x\rightarrow 0}f(x)=0$ et $\lim_{x\rightarrow 0}\frac{f(2x)-f(x)}{x}=0$. Montrer que $\lim_{x\rightarrow 0}\frac{f(x)}{x}=0$. (Indication. Considérer $g(x)=\frac{f(2x)-f(x)}{x}$.)
\finenonce{005383}


\finexercice


\exercice{5384, rouget, 2010/07/06}
\enonce{005384}{**I}
Soient $f$ et $g$ deux fonctions continues en $x_0\in\Rr$. Montrer que $\mbox{Min}\{f,g\}$ et $\mbox{Max}\{f,g\}$ sont continues en $x_0$.
\finenonce{005384}


\finexercice


\exercice{5385, rouget, 2010/07/06}
\enonce{005385}{***I Distance d'un point à une partie}
\label{exo:roudist}
Soit $A$ une partie non vide de $\Rr$. Pour $x\in\Rr$, on pose $f(x)=\mbox{Inf}\{|y-x|,\;y\in A\}$. Montrer que $f$ est continue en tout point de $\Rr$.
\finenonce{005385}


\finexercice


\exercice{5386, rouget, 2010/07/06}
\enonce{005386}{**T}
Montrer en revenant à la définition que $f(x)=\frac{3x-1}{x-5}$ est continue en tout point de $\Rr\setminus\{5\}$.
\finenonce{005386}


\finexercice


\exercice{5387, rouget, 2010/07/06}
\enonce{005387}{**IT}
Montrer que la fonction caractéristique de $\Qq$ est discontinue en chacun de ses points.
\finenonce{005387}


\finexercice


\exercice{5388, rouget, 2010/07/06}
\enonce{005388}{****}
Etudier l'existence d'une limite et la continuité éventuelle en chacun de ses points de la fonction définie sur $]0,+\infty[$ par $f(x)=0$ si $x$ est irrationnel et $f(x)=\frac{1}{p+q}$ si $x$ est rationnel égal à $\frac{p}{q}$, la fraction $\frac{p}{q}$ étant irréductible.
\finenonce{005388}


\finexercice


\exercice{5389, rouget, 2010/07/06}
\enonce{005389}{**IT}
Etudier en chaque point de $\Rr$ l'existence d'une limite à droite, à gauche, la continuité de la fonction $f$ définie par $f(x)=xE(\frac{1}{x})$ si $x\neq0$ et $1$ si $x=0$.
\finenonce{005389}


\finexercice


\exercice{5390, rouget, 2010/07/06}
\enonce{005390}{**}
Trouver $f$  bijective de $[0,1]$ sur lui-même et discontinue en chacun de ses points.
\finenonce{005390}


\finexercice


\exercice{5391, rouget, 2010/07/06}
\enonce{005391}{***}
Soit $f$ une fonction continue et périodique sur $\Rr$ à valeurs dans $\Rr$, admettant une limite réelle quand $x$ tend vers $+\infty$. Montrer que $f$ est constante.
\finenonce{005391}


\finexercice


\exercice{5392, rouget, 2010/07/06}
\enonce{005392}{**I}
Soit $A$ une partie non vide de $\Rr$. Pour $x$ réel, on pose $f(x)=d(x,A)=\mbox{Inf}\{|y-x|,\;y\in A\}$. Montrer que $f$ est Lipschitzienne.
\finenonce{005392}


\finexercice


\exercice{5395, rouget, 2010/07/06}
\enonce{005395}{***}
Soit $f$ croissante de $[a,b]$ dans lui-même. Montrer que $f$ a un point fixe.
\finenonce{005395}


\finexercice


\exercice{5396, rouget, 2010/07/06}
\enonce{005396}{****}
Soit $f$ croissante sur $[a,b]$ telle que $f([a,b])=[f(a),f(b)]$. Montrer que $f$ est continue sur $[a,b]$.
\finenonce{005396}


\finexercice


\exercice{5398, rouget, 2010/07/06}
\enonce{005398}{***IT}
Soit $f$ continue sur $\Rr^+$ à valeurs dans $\Rr$ admettant une limite réelle quand $x$ tend vers $+\infty$. Montrer que $f$ est uniformément continue sur $\Rr^+$.
\finenonce{005398}


\finexercice


\exercice{5401, rouget, 2010/07/06}
\enonce{005401}{***}
Soit $f$ périodique et continue sur $\Rr$. Montrer que $f$ est bornée et uniformément continue sur $\Rr$.
\finenonce{005401}


\finexercice

\section{123.02 Continuité : pratique}


\exercice{670, vignal, 2001/09/01}
\video{AqwFsD85iiU}
\enonce{000670}{}
 Soit $f : \R\setminus\{1/3\}\rightarrow \R$ 
telle que $f(x)= \frac{2x+3}{3x-1}$.

Pour tout $\epsilon>0$ d\'eterminer $\delta$ tel que,
($x\not=1/3$ et $|x|\leq\delta)\Rightarrow |f(x)+3|\leq \epsilon$.

Que peut-on en conclure ?
\finenonce{000670}


\finexercice

\exercice{671, vignal, 2001/09/01}
\video{xN3Z9gW5JEs}
\enonce{000671}{}
Soit $f$ la fonction r\'eelle \`a valeurs r\'eelles
d\'efinie par
$$f(x)=\left\{\begin{array}{ll}
x & \hbox{ si } x< 1\\
x^2  & \hbox{ si }1\leq x\leq 4\\
8\,\sqrt{x}  & \hbox{ si } x>4
\end{array}\right.$$
\begin{enumerate}
  \item Tracer le graphe de $f$.
  \item  $f$ est elle continue ?
  \item Donner la formule d\'efinissant $f^{-1}$.
\end{enumerate}
\finenonce{000671}


\finexercice


\exercice{672, vignal, 2001/09/01}
\enonce{000672}{}
 Etudier la continuit\'e de $f$
la fonction r\'eelle \`a valeurs r\'eelles
d\'efinie par
$f(x)= \frac{\sin x}{x}$ si $x\not= 0$ et $f(0)=1$.
\finenonce{000672} 


\finexercice


\exercice{673, vignal, 2001/09/01}

\enonce{000673}{}
\begin{enumerate}
  \item Soit la fonction r\'eelle d\'efinie par $f(x)=1$ si $x\in\Q$
et $f(x)=0$ sinon. Montrer que $f$ n'admet pas de limite en tout point de $\R$.
  \item Soit la fonction r\'eelle d\'efinie par $f(x)=x$ si $x\in\Q$
et $f(x)=1-x$ sinon. En quels points de $\R$ $f$ est elle continue ?
\end{enumerate}
\finenonce{000673}



\finexercice



\exercice{674, cousquer, 2003/10/01}

\enonce{000674}{}
 On admet que pour tout $x\in \mathbb{R} , \vert \sin x \vert \leq \vert x\vert .$
\begin{enumerate}
 \item  Montrer que $x\mapsto \sin x $ est continue en $0$ puis sur
$\mathbb{R}$ tout entier. 
\item  En déduire que $x\mapsto \cos x $ est continue sur $\mathbb{R}$. 
\end{enumerate}
\finenonce{000674}



\finexercice



\exercice{675, cousquer, 2003/10/01}

\enonce{000675}{}
 Etudier la continuité sur $\mathbb{R}$ des fonctions suivantes :
\begin{enumerate}
        \item $f_1(x)=x^2 \cos \frac{1}{x} \mbox{ si } x\not=0, \mbox{ et } 
f_1(0)=0$ ;
        \item $f_2(x)= \sin x\sin \frac{1}{x} \mbox{ si } x\not=0, \mbox{ et } f_2(0)=0$ ;
        \item $f_3(x)=x E(x) $ ;
        \item $f_4(x) = E(x) \sin(\pi x) $.
\end{enumerate}
\finenonce{000675}



\finexercice



\exercice{676, cousquer, 2003/10/01}

\enonce{000676}{}
 Montrer que l'application $f : \mathbb{R} \rightarrow \mathbb{R}$ définie par
$\displaystyle{ f(x)=\frac{x}{1+\vert x\vert}}$
est strictement croissante puis que pour tout $y\in\quad]-1,1[ $ il existe
un unique $x\in \mathbb{R}$
tel que $f(x)=y$.
  
\finenonce{000676}



\finexercice



\exercice{677, bodin, 1998/09/01}
\video{C7uoURImunQ}
\enonce{000677}{}
Les fonctions suivantes sont-elles prolongeables par continuit\'e sur
$\R$ ?
$$ a)\ f(x)=\sin x \cdot \sin \frac{1}{x}\ ;\ \ \
b)\ g(x)=\frac{1}{x}\ln\frac{e^x+e^{-x}}{2}\ ;$$
$$c)\ h(x)=\frac{1}{1-x}-\frac{2}{1-x^2}\ .$$
\finenonce{000677} 


\finexercice


\exercice{678, bodin, 1998/09/01}

\enonce{000678}{}
\'Etudier la continuit\'e sur $\Rr$ des fonctions suivantes :
\begin{enumerate}
    \item $f(x) = E(x)\sin(x)$,

    \item $g(x) = E(x)\sin(\pi x)$.
\end{enumerate}
\finenonce{000678}



\finexercice



\exercice{679, ridde, 1999/11/01}

\enonce{000679}{}
Etudier la continuit\'e de
\begin{enumerate}
\item $f (x) = x + \sqrt{x-\text{E} (x)}$.
\item $g (x) = \text{E} (x) + \sqrt{x-\text{E} (x)}$.
\end{enumerate}
\finenonce{000679}



\finexercice



\exercice{680, ridde, 1999/11/01}
\video{bdGeuvwKDQA}
\enonce{000680}{}
Soit $f : \Rr \rightarrow \Rr$ continue en $0$ telle que pour chaque $x \in \Rr$,
$f(x) = f(2x)$. Montrer que $f$ est constante.
\finenonce{000680} 


\finexercice

\exercice{681, ridde, 1999/11/01}

\enonce{000681}{}
La fonction $\dfrac 1x$ est-elle lipschitzienne sur $]0,  + \infty[$ ? sur
$[1,  + \infty[$ ?
\finenonce{000681}



\finexercice



\exercice{682, bodin, 1998/09/01}

\enonce{000682}{}
 Soit $f : [0,1]\longrightarrow \R$ d\'efinie par $f(0)=0$,
$f(x)=1/2-x$ si $x\in ]0,1/2[$, $f(1/2)=1/2$, $f(x)=3/2-x$ si $x\in ]1/2,1[$ et $f(1)=1$.
\begin{enumerate}
    \item  Tracer le graphe de $f$. \'Etudier sa continuit\'e.\\
    \item  D\'emontrer que $f$ est une bijection de $[0,1]$ sur $[0,1]$.\\
    \item  D\'emontrer que pour tout $x\in [0,1]$, on a
$ f(x)= \frac{1}{2}-x+\frac{1}{2}E(2x)-\frac{1}{2}E(1-2x)$.
\end{enumerate}
\finenonce{000682}



\finexercice



\exercice{683, bodin, 1998/09/01}

\enonce{000683}{}
\'Etudier la continuit\'e des fonctions suivantes :
\begin{enumerate}
    \item $f_1(x)=x^2\cos \frac{1}{x} \text{\ \  si }x\not=0 \qquad f_1(0)=0$ ;
    \item $f_2(x)= \sin x\sin \frac{1}{x} \text{\ \  si }x\not=0 \qquad f_2(0)=0$ ;
    \item $f_3(x)=xE(x) \text{ sur } \Rr$ ;
    \item $f_4(x) = \left[ x-E(x)\right]^2 \text{\ \  et\ \  } f_5(x) = E(x)+f_4(x)$.
\end{enumerate}
\finenonce{000683}



\finexercice



\exercice{684, ridde, 1999/11/01}

\enonce{000684}{}
En \'etudiant la suite
 $u_0 \in \Rr$ et $u_{n  + 1} = \cos (u_n)$, d\'eterminer une valeur approch\'ee \`a $10^{-5}$
 pr\`es de l'unique r\'eel solution de $\cos (x) = x$.
\finenonce{000684}



\finexercice



\exercice{685, gourio, 2001/09/01}

\enonce{000685}{}
Soit $f$ d\'{e}finie par $f(x)=E(x)+\sqrt{x-E(x)}, $ o\`{u} $E$ d\'{e}signe
la partie enti\`{e}re.
Donner le domaine de d\'{e}finition de $f,$ puis une relation entre $f(x+1)$
et $f(x)$. $f$ est-elle monotone?\ $f$ est-elle $k-$lipschitzienne sur $[a,1] $($a>0$) ?
Et sur $[0,1]$ ?
\'Etudier la continuit\'{e} de $f$ sur $[0,1]$ en utilisant la
d\'{e}finition. D\'{e}duisez en la continuit\'{e} sur $\Rr$.
\finenonce{000685}



\finexercice



\exercice{1212, legall, 1998/09/01}

\enonce{001212}{}
Soit $  f   $ une fonction continue de $  [0,1]  $ dans lui-m\^eme telle que $  f(0)=0  $
et pour tout couple $  (x,y)  $ de $  [0,1]\times[0,1]  $ on ait $  \vert f(x)-f(y)\vert \geq
\vert x-y\vert   .$
\begin{enumerate}
    \item Soit $  x   $ un \' el\' ement de $  [0,1]  .$ On pose $  x_0=x  $ et $  x_{n+1}=f(x_n)  .$ Montrer
que la suite $  (x_n)_{n \in { \Nn}}  $ est convergente.
    \item En d\' eduire que $  f(x)=x  $ pour tout $  x\in [0,1]  .$
    \item Le r\' esultat reste-t-il vrai sans l'hypoth\`ese  $  f(0)=0  ?$
\end{enumerate}
\finenonce{001212}



\finexercice


\section{123.03 Limite de fonctions}


\exercice{606, bodin, 1998/09/01}

\enonce{000606}{}

 \'Ecrire les d\'efinitions des limites suivantes :
 $\lim_{x\to -\infty}f(x)=l$, $l\in\R$ ; $\lim_{x\to -\infty}f(x)=+\infty$ ; $\lim_{x\to x_0}f(x)=-\infty$, $x_0\in\R$.\\
(On pr\'ecisera sur quel type d'intervalle la fonction $f$ doit \^etre d\'efinie.)
\finenonce{000606}



\finexercice



\exercice{607, bodin, 1998/09/01}

\enonce{000607}{}

 Soit $f$ une fonction d\'efinie sur un intervalle $I$ contenant
$x_0$ dans son int\'erieur. On suppose que $ \lim_{x\rightarrow x_0}f(x)=u>0$. D\'emontrer qu'il
existe $t>0$ tel que si $0<\vert x-x_0\vert < t$ alors $\vert f(x)\vert \geq \frac{u}{  2}$.
\finenonce{000607}



\finexercice



\exercice{608, bodin, 1998/09/01}

\enonce{000608}{}

 Montrer que si une fonction $f$ d\'efinie sur $E\subset\R$ est continue en
 $x_0$ alors la fonction $\vert f\vert$ est, elle aussi, continue en $x_0$.
Montrer que la r\'eciproque est fausse.
\finenonce{000608}



\finexercice



\exercice{609, bodin, 1998/09/01}
\video{db5yEXsEbYc}
\enonce{000609}{}
\begin{enumerate}
    \item  D\'emontrer que $\displaystyle{ \lim_{x\rightarrow 0}\frac{\sqrt{1+x}-\sqrt{1-x}}{ x}=1}$.
    \item  Soient $m,n$ des entiers positifs. \'Etudier $\displaystyle{\lim_{x\rightarrow 0}\frac{\sqrt{1+x^m}-
\sqrt{1-x^m}}{ x^n}}$.
    \item D\'emontrer que $\displaystyle{ \lim_{x\rightarrow 0}\frac{1}{ x}(\sqrt{1+x+x^2}-1)=
\frac{1}{ 2}}$.
\end{enumerate}
\finenonce{000609} 


\finexercice

\exercice{610, bodin, 1998/09/01}

\enonce{000610}{}

Soit $f$ une fonction de variable r\'eelle telle que $\frac{f(x)}{|x|} \rightarrow \infty$ quand
$x \rightarrow \infty$. Montrer que pour tout r\'eel $\alpha$ il existe $X_\alpha$ tel que
$f(x)-|\alpha x| \ge |x|$ si $|x| \ge X_\alpha$.
En d\'eduire que pour tout $\alpha$ r\'eel $f(x)-\alpha x \rightarrow \infty$ quand  $x\rightarrow \infty$.
\finenonce{000610}



\finexercice



\exercice{611, bodin, 1998/09/01}

\enonce{000611}{}

Soient $f$ et $g$ deux fonctions d\'efinies sur $\Rr_+$ telles que
$$
\forall x\in \Rr_+ \ g(x) > 0 \text{ et } \lim_{x\rightarrow\infty} \frac{f(x)}{g(x)} = L \not= 0.
$$
\begin{enumerate}
    \item Montrer que
$$ \lim_{x \rightarrow \infty} f(x) = 0 \Leftrightarrow  \lim_{x \rightarrow \infty} g(x) = 0. $$
    \item Montrer que si $L>0$,
$$ \lim_{x \rightarrow \infty} f(x) = \infty \Leftrightarrow  \lim_{x \rightarrow \infty} g(x) = \infty. $$
\end{enumerate}
\finenonce{000611}



\finexercice



\exercice{612, bodin, 1998/09/01}
\video{OEzb_HFP0yM}
\enonce{000612}{}
\begin{enumerate}
    \item Montrer que toute fonction p\'eriodique
et non constante n'admet pas de limite en $+\infty$.
    \item Montrer que toute fonction croissante
et major\'ee admet une limite finie en $+\infty$.
\end{enumerate}
\finenonce{000612}
 

\finexercice


\exercice{613, cousquer, 2003/10/01}

\enonce{000613}{}
 Soit $I$ un intervalle de $\mathbb{R}$ et $x_0 \in I.$
Soient $f$ et $g$ deux fonctions de la variable réelle à valeurs réelles
définies sur $\dot{I} := I-\{x_0\}.$ Montrer que si $f$ admet une limite à droite et une limite à gauche en $x_0$ et que de plus ces deux limites coïncident, alors $f$ admet une limite en $x_0$ dont la valeur est la valeur commune des limites à droite et à gauche.
\finenonce{000613}



\finexercice



\exercice{614, cousquer, 2003/10/01}

\enonce{000614}{}
 Soient $P$ et $Q$ deux polynômes à coefficients réels de degré
respectif $d$ et $d'.$ Etudier suivant les valeurs de $d$ et $d',$ et éventuellement de certains des coefficients de $P$ et $Q,$
$$ \lim_{x \rightarrow +\infty}P(x)/Q(x). $$
\finenonce{000614}



\finexercice



\exercice{615, ridde, 1999/11/01}

\enonce{000615}{}
Soit $f : \Rr^+  \rightarrow \Rr^+ $ croissante telle que $ \lim\limits_{
x \rightarrow  + \infty}f (x + 1)-f (x) = 0 $. Montrer que $\lim\limits_{
x \rightarrow  + \infty}\dfrac{f (x)}x = 0$.
 (on pourra utiliser des $\epsilon$, sommer des in\'egalit\'es et utiliser la
 monotonie de $f$ pour montrer qu'elle est born\'ee sur un segment).\\
 Comment g\'en\'eraliser ce r\'esultat ?
\finenonce{000615}



\finexercice



\exercice{616, vignal, 2001/09/01}
\video{OYJj7QRtecs}
\enonce{000616}{}
 Calculer lorsqu'elles existent les limites suivantes
$$\begin{array}{lll}
  a)\ \lim_{x\rightarrow 0}\frac{x^2+2\,|x|}{x} &\quad  b)\
\lim_{x\rightarrow -\infty}\frac{x^2+2\,|x|}{x}&
\quad c) \ \lim_{x\rightarrow 2}\frac{x^2-4}{x^2-3\,x+2}\\
\\
d) \ \lim_{x\rightarrow\pi}\frac{\sin^2x}{1+\cos x}&
\quad  e)\ \lim_{x\rightarrow 0}\frac{\sqrt{1+x}-\sqrt{1+x^2}}{x}&
\quad  f)\ \lim_{x\rightarrow +\infty}\sqrt{x+5}-\sqrt{x-3}\\
\\
  g)\ \lim_{x\rightarrow 0}\frac{\sqrt[3]{1+x^2}-1}{x^2}&
\quad  h)\ \lim_{x\rightarrow 1}\frac{x-1}{x^n-1}
\end{array}$$
\finenonce{000616} 


\finexercice

\exercice{617, cousquer, 2003/10/01}

\enonce{000617}{}
\begin{enumerate}
  \item  Montrer que pour tout $0<\epsilon <1$ et pour $x \in \mathbb{R},$ on a:
$$ |x-1| < \frac{\epsilon}{4} \Rightarrow |x^2+x-2| < \epsilon.$$
  \item En déduire:
$$ \lim_{x \rightarrow 1}x^2+x-1\;\mbox{ et }\; \lim_{x \rightarrow 1}(x^2+x-2)\cos x.$$   
\end{enumerate}
\finenonce{000617}



\finexercice



\exercice{618, cousquer, 2003/10/01}

\enonce{000618}{}
\begin{enumerate}
  \item Montrer que pour tout $a \in \mathbb{R}^{+*},$ et pour tout couple de nombres
    réels $(x,y)$ appartenant à $]-\infty,-a]$ ou à $[a,\infty[,$ on~a:
$$|\frac 1x - \frac 1y| \leq \frac 1{a^2}|x-y|.$$
  \item  En déduire que pour tout $x_0 \in \mathbb{R}^*$ et pour tout $\epsilon > 0$ il existe $\alpha > 0$  tel que :
$$|x-x_0| < \alpha \;\Rightarrow\; |\frac 1x - \frac 1{x_0}| < \epsilon.$$
  \item  En déduire que la fonction $x \mapsto \frac 1x$ est continue en tout point de $\mathbb{R}^*.$ 
\end{enumerate}
\finenonce{000618}



\finexercice



\exercice{619, cousquer, 2003/10/01}

\enonce{000619}{}
\begin{enumerate}
\item  Pour tout $n$ entier naturel et tout couple de réels $(x,y)$, établir la formule :
$$
x^n-y^n = (x-y).\sum_{k=0}^{n-1}x^ky^{n-1-k}.
$$
\item Déduire de la question précédente que pour tout entier $n$ tout réel strictement positif $a$ et tout couple de réels $(x,y)$ tel que $|x| \leq a$ et $|y| \leq a,$
$$
|x^n-y^n| \leq na^{n-1}|x-y|.
$$
\item Déduire de ce qui précède que pour tout $x_0 \in \mathbb{R},$ et pour tout $\epsilon > 0,$ il existe $\alpha > 0$ tel que:
$$
|x-x_0| < \alpha \;\Rightarrow\; |x^n-x_0^n| < \epsilon.
$$
Conclure.
\item Sur quel sous ensemble $D$ de $\mathbb{R},$ la fonction de la variable réelle $f$ donnée par 
$$
f(x) := \frac{1-x^n}{1-x}
$$
est-elle définie? Calculer les limites de $f$ aux bornes de $D.$ 
\end{enumerate}
\finenonce{000619}



\finexercice



\exercice{620, cousquer, 2003/10/01}

\enonce{000620}{}

\begin{enumerate}
  \item Rappeler que pour tout nombre réels $\epsilon > 0$ il existe un entier $n$ tel que: 
\begin{eqnarray*}
      \frac 1{2n\pi} &<& \epsilon\\
\frac 1{(2n+1)\pi} &<& \epsilon.\\
    \end{eqnarray*}
  \item Montrer que pour tout nombre réel $l,$ et pour tout $\epsilon > 0,$ il existe $x \in ]-\epsilon,\epsilon[$ tel que:
$$
|\sin \frac 1x -l | > \frac 12.
$$
\item En déduire que la fonction $x \mapsto \sin \frac 1x$ n'a pas de limite lorsque $x$ tend vers $0.$
\item 
Montrer que la fonction définie par $f(x) =x \sin(\frac 1x)$ pour $x \not= 0$ et $f(0) =0$ est continue sur $\mathbb{R}.$ 
\end{enumerate}
\finenonce{000620}



\finexercice



\exercice{621, cousquer, 2003/10/01}

\enonce{000621}{}
 Déterminer les limites suivantes:
 $$\mbox{a)} \quad \lim_{x \rightarrow +\infty}\sqrt{x^2+1}-x \qquad \mbox{b)} \quad  
\lim_{x \rightarrow 1}\frac 1{x-1}-\frac 2{x^2-1} $$
$$\mbox{c)} \quad \lim_{x \rightarrow 0^+}\sqrt{1+\frac 1x}-\sqrt{\frac 1x} \qquad \mbox{d)} \quad  \lim_{x \rightarrow 4}\frac{\sqrt{2x+1}-3}{\sqrt{x-2}-\sqrt 2}$$ 
 $$\mbox{e)} \quad \lim_{x \rightarrow +\infty}\sqrt{x^2+1}-\sqrt{x^2-1}\qquad \mbox{f)} \quad \lim_{x \rightarrow -\infty}x(\sqrt{1+x^2}-x)$$
\finenonce{000621}



\finexercice



\exercice{622, cousquer, 2003/10/01}

\enonce{000622}{}
 On rappelle les limites :
$\lim_{x \rightarrow 0}\frac{\sin x}x = 1 \;\mbox{ et }\; \lim_{x \rightarrow 0}\frac{1-\cos x}{x^2} = \frac 12. $
\\Calculer les limites suivantes:
$$\mbox{a)} \quad \lim_{x \rightarrow 0^+}\sqrt x.\sin\frac 1{\sqrt x} \qquad \mbox{b)}\quad 
  \lim_{x \rightarrow 0}\frac{\sin 2x}{\sin 3x} $$
$$\mbox{c)} \quad \lim_{x \rightarrow 0}\frac{x\sin x}{1-\cos x}\qquad \mbox{d)}\quad 
\lim_{x \rightarrow 0}\frac{\sin x - \sin 2x}{x^2} $$
$$\mbox{e)} \quad \lim_{x \rightarrow 0}x\frac{\tan x}{\cos^2 x-1} \qquad \mbox{f)}\quad 
 \lim_{x \rightarrow 0}\frac{\tan x-\sin x}{\sin^3(\frac x2)}$$
\finenonce{000622}



\finexercice



\exercice{623, cousquer, 2003/10/01}
\enonce{000623}{}
\noindent D\'eterminer les limites suivantes, en justifiant vos calculs.
\begin{enumerate}
\item $\displaystyle\lim_{x \rightarrow0^+}{x+2 \over x^2 \ln x}$

\item $\displaystyle\lim_{x \rightarrow0^+}2x \ln(x+\sqrt x)$

\item $\displaystyle\lim_{x \rightarrow+ \infty}{x^3-2x^2+3 \over x \ln x}$

\item $\displaystyle\lim_{x \rightarrow+ \infty}{e^{\sqrt x+1} \over x+2}$

\item $\displaystyle\lim_{x \rightarrow0^+}{\ln(3x+1) \over2x}$

\item $\displaystyle\lim_{x \rightarrow0^+}{x^x-1 \over\ln(x+1)}$

\item $\displaystyle\lim_{x \rightarrow- \infty}{2 \over x+1}\ln
\Bigl({x^3+4 \over1-x^2}\Bigr)$

\item $\displaystyle\lim_{x \rightarrow{(-1)}^+}(x^2-1) \ln(7x^3+4x^2+3)$

\item $\displaystyle\lim_{x \rightarrow2^+}{(x-2)}^2 \ln(x^3-8)$

\item $\displaystyle\lim_{x \rightarrow0^+}{x(x^x-1) \over\ln(x+1)}$

\item $\displaystyle\lim_{x \rightarrow+ \infty}(x \ln x -x \ln(x+2))$

\item $\displaystyle\lim_{x \rightarrow+ \infty}{e^x-e^{x^2} \over x^2-x}$ 

\item $\displaystyle\lim_{x \rightarrow0^+}{{(1+x)}^{\ln x}}$

\item $\displaystyle\lim_{x \rightarrow+ \infty}{\Bigl({x+1 \over x-3}\Bigr)^x}$

\item $\displaystyle\lim_{x \rightarrow+ \infty}{\Bigl({x^3+5 \over
x^2+2}\Bigr)^{x+1 \over x^2+1}}$

\item $\displaystyle\lim_{x \rightarrow+ \infty}{\Bigl({e^x+1 \over
x+2}\Bigr)^{1 \over x+1}}$

\item $\displaystyle\lim_{x \rightarrow0^+}\bigl(\ln(1+x)\bigr)^{1\over\ln x}$

\item $\displaystyle\lim_{x \rightarrow+ \infty}{x^{(x^{x-1})} \over
x^{(x^{x})}}$

\item $\displaystyle\lim_{x \rightarrow+ \infty}{(x+1)^x \over x^{x+1}}$

\item $\displaystyle\lim_{x \rightarrow+ \infty}{x \sqrt{\ln(x^2+1)} \over
1+e^{x-3}}$
\end{enumerate}
\finenonce{000623} 


\finexercice


\exercice{624, bodin, 1998/09/01}

\enonce{000624}{}

Soient $a,b$ des r\'eels positifs. $E(x)$ d\'esigne la partie enti\`ere de $x$.
Montrer que :
$$ \lim_{x\to 0^+}\frac{x}{a}E(\frac{b}{x})=\frac{b}{a} \qquad ;\qquad \lim_{x\to 0^+}\frac{b}{x}E(\frac{x}{a})=0.$$
\finenonce{000624}



\finexercice



\exercice{625, bodin, 1998/09/01}

\enonce{000625}{}

 Calculer les limites suivantes :
$$\lim_{x\to 1}\frac{x-1}{x^n-1};\; \lim_{x\to a}\frac{x^m-a^m}{x^p-a^p}\;(a>0,m,p\in \N^*);
\;\lim_{h\to 0}\frac{(x+h)^n-x^n}{h}\;(x\in \R,n\in\N^*)$$
$$\lim_{x \to 0^+}(\sqrt{\frac{1}{x}+1}-\sqrt{\frac{1}{x}-1});\;
\lim_{x\to -\frac{\pi}{4}}\frac{\cos{x}+\sin{x}}{4x+\pi} ;\;\lim_{x \to 0^+}\frac{\sqrt{x}-x}{\sqrt{x}+x}.$$
\finenonce{000625}



\finexercice



\exercice{626, bodin, 1998/09/01}

\enonce{000626}{}

 En utilisant la d\'efinition d'une limite, montrer que :
$$
a)\ \lim_{x\to-\frac{2}{3}}(3x+2)\sin\left(\frac{1}{3x+2}\right)=0\ ;\ \ \
b)\  \lim_{x\to 0^{+}} \frac{2}{1+e^{-\frac{1}{x}}}=2.
$$
\finenonce{000626}



\finexercice



\exercice{627, bodin, 1998/09/01}

\enonce{000627}{}

Calculer les limites suivantes :
$$a)\  \lim_{x\to 0^{+}}xE(\frac{1}{x}) \ ;\ \ \
b)\ \lim_{x\to +\infty}xE(\frac{1}{x})\ ;\ \ \
c)\ \lim_{x\to 0^{+}}\sqrt{x}E(\frac{1}{x})\ ;$$

$$\ d)\ \lim_{x\to +\infty}\frac{\sqrt{x+\sqrt{x}}+\sqrt{x}}{\sqrt{x+1}}\ ;\
\ \ e)\ \lim_{x\to +\infty}(\sqrt{x+5}-\sqrt{x-3}).$$
\finenonce{000627}



\finexercice



\exercice{628, bodin, 1998/09/01}
\video{llaMFH13iLQ}
\enonce{000628}{}
Calculer, lorsqu'elles existent, les limites suivantes :
$$
\lim_{x\rightarrow \alpha} \frac{x^{n+1}-\alpha^{n+1}}{x^n-\alpha^n},
$$

$$
\lim_{x\rightarrow 0} \frac{\tan x - \sin x}{\sin x(\cos 2x - \cos x)},
$$

$$
\lim_{x\rightarrow +\infty} \sqrt{x+\sqrt{x+\sqrt{x}}}-\sqrt{x},
$$

$$
 \lim_{x\rightarrow \alpha^+} \frac{\sqrt{x}-\sqrt{\alpha}-\sqrt{x-\alpha}}{\sqrt{x^2-\alpha^2}}, \quad (\alpha >0)
$$

$$
\lim_{x\rightarrow 0} xE\left(\frac{1}{x}\right),
$$

$$
\lim_{x\rightarrow 2} \frac{e^x-e^2}{x^2+x-6},
$$

$$
\lim_{x\rightarrow +\infty} \frac{x^4}{1+x^\alpha\sin^2x}, \text{ en fonction de $\alpha \in \Rr$.}
$$
\finenonce{000628} 


\finexercice


\exercice{629, ridde, 1999/11/01}

\enonce{000629}{}
D\'eterminer les limites suivantes :
$$\frac{x}{2 + \sin (\frac 1x)} \qquad \text{ en } 0$$

$$\frac{x^3-3x^2 + 5x-3}{4x^4 + x^2 + x-6} \qquad \text{ en } 1$$
$$\frac{\sqrt{1 + \sin x}-\sqrt{1-\sin x}}{x}\qquad \text{ en } 0$$
$$\frac{\tan x}{\sqrt{x^2 + 4} + x-2}\qquad \text{ en } 0$$
$$\frac{1-\cos x}{x^2}\qquad \text{ en } 0$$
$$\frac{1-\sin x + \cos x}{\sin x  + \cos x -1}\qquad \text{ en } \frac {\pi}2$$
$$\frac{\tan (x + \frac{\pi}4)-1}{\sqrt 3 -2\cos (x + \frac{\pi}6)}
\qquad \text{ en } 0$$
\finenonce{000629}



\finexercice



\exercice{630, ridde, 1999/11/01}

\enonce{000630}{}
\'Etudier les asymptotes de $f (x) = e^{\frac 1x}\sqrt{x (x + 2)}$.
\finenonce{000630}



\finexercice



\exercice{631, bodin, 1998/09/01}

\enonce{000631}{}

Montrer que
$$\frac{\ln (x)}{x^\alpha} < \frac{2}{\alpha x^{\alpha/2}} \text{ o\`u } \alpha > 0.$$
En d\'eduire que
$$\lim_{x\rightarrow +\infty} \frac{\ln (x)}{x^\alpha} = 0,\ \ \alpha>0.$$
\finenonce{000631}



\finexercice



\exercice{632, bodin, 1998/09/01}

\enonce{000632}{}

Calculer les limites suivantes :
$$
\text{a) }\ \lim_{x\rightarrow 0}\ \  \frac{x^2+2|x|}{x} \qquad
\text{b) }\ \lim_{x\rightarrow 1}\ \  \frac{x^7-1}{x^6-1}  \qquad
\text{c) }\ \lim_{x\rightarrow 1}\ \  \frac{x^n-1}{x^m-1} \ \  n,m\in \Nn^*
$$

$$
\text{d) }\ \lim_{x\rightarrow 0}\  \frac{1-\cos x}{x^2}  \qquad
\text{e) }\ \lim_{x\rightarrow 0}\  \frac{x\sin x}{1-\cos x}  \qquad
\text{f) }\ \lim_{x\rightarrow 0}\  \frac{\ln (1+x^3)}{x}
$$

$$
\text{g) }\ \lim_{x\rightarrow 0}\  \frac{a^x-b^x}{x} \ \  a,b>0 \qquad
\text{h) }\ \lim_{x\rightarrow 0}\  \frac{e^{-ax}-e^{-bx}}{x} \qquad
\text{i) }\ \lim_{x\rightarrow \alpha^+}\  \frac{\sqrt{x}-\sqrt{\alpha}+\sqrt{x-\alpha}}{\sqrt{x^2-\alpha^2}}.
$$
\finenonce{000632}



\finexercice



\exercice{633, gourio, 2001/09/01}

\enonce{000633}{}
Calculer :
$$\lim\limits_{x\rightarrow \infty }\ln (1+e^{-x})^{\frac{1}{x}
},\lim\limits_{x\rightarrow 0}\frac{x}{2+\sin \frac{1}{x}}
,\lim\limits_{x\rightarrow 0^{+}}x^{\frac{1}{\ln (e^{x}-1)}}.$$
\finenonce{000633}



\finexercice



\exercice{634, gourio, 2001/09/01}

\enonce{000634}{}
 Calculer :
$$\lim_{x\rightarrow 0^{+}}\left(\frac{1}{(\sin x)^{2}}-\frac{1}{(\sinh x)^{2}}\right).$$


\finenonce{000634}



\finexercice



\exercice{635, gourio, 2001/09/01}
\video{rxI8iTwhH6E}
\enonce{000635}{}
 Calculer :
$$\lim\limits_{x\rightarrow 0}\frac{x}{2+\sin \frac{1}{x}}
,\ \ \lim\limits_{x\rightarrow +\infty }(\ln (1+e^{-x}))^{\frac{1}{x}}
,\ \ \lim\limits_{x\rightarrow 0^{+}}x^{\frac{1}{\ln (e^{x}-1)}}.$$
\finenonce{000635} 


\finexercice


\exercice{636, gourio, 2001/09/01}
\enonce{000636}{}
 Trouver :
$$\lim_{x\rightarrow 0^{+}}\frac{x^{x^{x}}\ln x}{x^{x}-1}$$
\finenonce{000636} 


\finexercice

\exercice{637, gourio, 2001/09/01}
\enonce{000637}{}
 Trouver pour $(a,b)\in (\Rr^{+*})^{2}$ :
$$\lim\limits_{x\rightarrow +\infty}\left(\frac{a^{x}+b^{x}}{2}\right)^{\frac{1}{x}}.$$
\finenonce{000637} 


\finexercice


\exercice{638, gourio, 2001/09/01}
\video{UAO0DkHX2EQ}
\enonce{000638}{}
Trouver pour $(a,b)\in (\Rr^{+*})^{2}$ :
$$\lim\limits_{x\rightarrow 0^{+}}\left(\frac{a^{x}+b^{x}}{2}\right)^{\frac{1}{x}}.$$
\finenonce{000638} 


\finexercice

\exercice{3848, quercia, 2010/03/11}
\enonce{003848}{$f(x+1) - f(x) \to a$}

Soit $f : \R \to \R$ continue telle que
$f(x+1) - f(x) \to a \in \overline{\R}$ lorsque $x\to+\infty$.

\begin{enumerate}
  \item   Montrer que $\frac {f(n)}n \to a$ lorsque $n\to\infty$.
  \item   Montrer que $\frac {f(x)}x \to a$ lorsque $x\to\infty$.
\end{enumerate}
\finenonce{003848}



\finexercice

\exercice{5101, rouget, 2010/06/30}
\enonce{005101}{**}
Trouver $\lim_{x\rightarrow +\infty}\frac{(x^x)^x}{x^{(x^x)}}$.
\finenonce{005101}


\finexercice

\section{123.04 Etude de fonctions}


\exercice{686, vignal, 2001/09/01}
\video{XW9frtB7vKc}
\enonce{000686}{}
 D\'eterminer les domaines de d\'efinition des fonctions suivantes

$$f(x)=\sqrt{\frac{2+3\,x}{5-2\,x}}\ ;\quad g(x)=\sqrt{x^2-2\,x-5}\ ;\quad
h(x)=\ln\left(4\,x+3\right).$$
\finenonce{000686} 


\finexercice

\exercice{687, vignal, 2001/09/01}

\enonce{000687}{}
 Montrer que l'\'equation $x^7-3\,x^2+4\,x-1=0$ admet au moins
une solution dans l'intervalle $]-1,1[$. M\^eme question pour l'\'equation
$x^{29}+14\, x^{17}-7\,x^5+2=0$.
\finenonce{000687}



\finexercice



\exercice{688, bodin, 1998/09/01}

\enonce{000688}{}
 Soient $n \in \N^*$ et $d \in \R^+$. D\'emontrer en utilisant
le th\'eor\`eme des valeurs interm\'ediaires que le polyn\^ome
$P(X)=X^n-d$ a au moins une racine dans $\R$.
\finenonce{000688}



\finexercice



\exercice{689, bodin, 1998/09/01}

\enonce{000689}{}
 En \'etudiant les variations de la fonction $f$
d\'efinie sur $\rbrack 0,+\infty\lbrack$ par $\displaystyle{f(x)=x^{\frac{1}{x}}}$, trouver
le plus grand \'el\'ement de l'ensemble $f(\N^*)$.
\par En d\'eduire que quels soient $m$ et $n$ appartenant \`a $\N^*$,
l'un des nombres $\sqrt[n]{m}$, $\sqrt[m]{n}$ est inf\'erieur ou \'egal
\`a $\sqrt[3]{3}$.
\finenonce{000689}



\finexercice



\exercice{690, bodin, 1998/09/01}
\enonce{000690}{}
Soit $f : \Rr \to \Rr$ définie par $f(x) = \frac{\cos x}{1+x^2}.$
Montrer que $f$ est major\'ee sur $\Rr$, minor\'ee sur $\Rr$.
D\'eterminer $\sup_{x\in \Rr} f(x)$.
\finenonce{000690} 


\finexercice

\exercice{691, bodin, 1998/09/01}

\enonce{000691}{}
\begin{enumerate}
    \item Soit la fonction $f:\lbrack -1,+\infty\lbrack\to\R$, d\'efinie par
$f(x)=\frac{1}{\sqrt{x^2+2x+2}}$. Montrer que $f$ admet une r\'eciproque
 que l'on explicitera.
    \item  Trouver un intervalle de $\R$ sur lequel la fonction $g(x)=\tan(x^3)$
admette une fonction r\'eciproque (on pr\'ecisera alors le domaine de
d\'efinition de cette r\'eciproque et son image).
\end{enumerate}
\finenonce{000691}



\finexercice



\exercice{692, gourio, 2001/09/01}

\enonce{000692}{}
Montrer que les fonctions suivantes ne sont pas des polyn\^{o}mes :
$$x\rightarrow e^{x},\ \ x\rightarrow \ln x,\ \ x\rightarrow \sqrt{x^{2}+1},\ \ x\rightarrow \cos x.$$
\finenonce{000692}



\finexercice



\exercice{3874, quercia, 2010/03/11}
\enonce{003874}{$f$ croissante et $f\circ f = \mathrm{id}$}

Soit $f : \R \to \R$ croissante telle que :
$\forall\ x \in \R,\ f\circ f(x) = x$.
Montrer que : $\forall\ x \in \R,\ f(x) = x$.
\finenonce{003874}



\finexercice


\exercice{3875, quercia, 2010/03/11}
\enonce{003875}{$f$ croissante et $x  \mapsto f(x)/x$ est décroissante}

Soit $f : {]0,+\infty[} \to {\R}$ croissante telle que
$g : {]0,+\infty[} \to {\R}, x \mapsto {\frac {f(x)}x}$ est décroissante.
Montrer que $f$ est continue.
\finenonce{003875}



\finexercice


\exercice{3876, quercia, 2010/03/11}
\enonce{003876}{\'Etude de $x(2+\sin(1/x))$}

On pose : 
$$\begin{cases} f(x) = |x|\left({2+\sin\left(\frac1x\right)}\right) &\text { si } x \ne 0,\cr
                   f(0) = 0.\cr\end{cases}$$

Montrer que $f$ est continue, minimale en $0$, mais pour tout $\varepsilon > 0$,
$f_{|[0,\varepsilon]}$ n'est pas monotone.
\finenonce{003876}



\finexercice


\exercice{3877, quercia, 2010/03/11}
\enonce{003877}{Borne supérieure de fonctions croissantes}

Soit $f : {\R} \to {\R}$ une fonction bornée.

On note ${\cal E} = \{g : \R \to \R \text{ croissantes tq } g \le f \}$,
et pour $x\in\R$ : $\tilde f(x) = \sup\{ g(x) \text{ tq } g \in {\cal E} \}$.

\begin{enumerate}
  \item   Montrer que $\tilde f \in {\cal E}$.
  \item   On suppose $f$ continue. Montrer que $\tilde f$ est aussi continue.\par
      (S'il existe un point $x_0 \in \R$ tel que
      $\lim_{x\to x_0^-} \tilde f(x) < \lim_{x\to x_0^+} \tilde f(x)$,
      construire une fonction de $\cal E$ supérieure à~$\tilde f\,$)
\end{enumerate}
\finenonce{003877}



\finexercice


\exercice{3878, quercia, 2010/03/11}
\enonce{003878}{L'ensemble des points de discontinuité est dénombrable}

Soit $f : {[a,b]} \to {\R}$ monotone.
Pour $x \in {]a,b[}$,
on pose $\delta(x) = \bigl|\lim_{y\to x^+} f(y) - \lim_{y\to x^-} f(y)\bigr|$
(saut de $f$ en $x$).

\begin{enumerate}
  \item   Pour $n \in \N^*$, montrer que
      $E_n = \left\{ x \in {]a,b[} \text{ tq } \delta(x) > \frac 1n \right\}$ est fini.
  \item   En déduire que l'ensemble des points de discontinuité de $f$ est au plus dénombrable.
\end{enumerate}
\finenonce{003878}



\finexercice


\exercice{3879, quercia, 2010/03/11}
\enonce{003879}{Fonction localement croissante}

Soit $f : {\R} \to {\R.}$
On dit que $f$ est localement croissante si :
$\forall\ x \in \R,\ \exists\ \varepsilon > 0 \text{ tq }
f_{|]x-\varepsilon, x+\varepsilon[}$ est croissante.

Montrer que : ($f$ est localement croissante) $ \Rightarrow $ ($f$ est croissante).

(\'Etudier $E = \{ x \ge 0 \text{ tq } f_{|[0,x]}\text{ est croissante}\}$ et
         $F = \{ x \le 0 \text{ tq } f_{|[x,0]}\text{ est croissante}\}$)

\finenonce{003879}



\finexercice


\exercice{3880, quercia, 2010/03/11}
\enonce{003880}{Prolongement d'une fonction uniformément continue}

Soit $f : {]0,1]} \to {\R}$ uniformément continue.
Pour $x \in {]0,1]}$ on pose :
$\begin{cases} g(x) = \sup\bigl(f(]0,x])\bigr) \cr
         h(x) = \inf\bigl(f(]0,x])\bigr).\cr\end{cases}$

\begin{enumerate}
  \item Montrer que $g$ et $h$ sont monotones.

On note $\ell = \lim_{x\to0^+} g(x)$ et $m = \lim_{x\to0^+} h(x)$.

  \item   En utilisant la continuité uniforme de $f$, montrer que $\ell = m$.
  \item   En déduire que $f(x) \to \ell$, si $x\to0^+$.
\end{enumerate}
\finenonce{003880}



\finexercice


\exercice{3881, quercia, 2010/03/11}
\enonce{003881}{$f$ continue, croissante sur $\Q$}

Soit $f : {\R} \to {\R}$ continue telle que $f_{|\Q}$ est
strictement croissante. Montrer que $f$ est strictement croissante.
\finenonce{003881}



\finexercice


\exercice{3882, quercia, 2010/03/11}
\enonce{003882}{Morphismes de $\R$}

Soit $f : {\R} \to {\R}$ non identiquement nulle telle que :
$\forall\ x,y \in \R,\ \begin{cases} f(x+y) &= f(x) + f(y) \cr
                               f(xy)  &= f(x)f(y).   \cr\end{cases}$

Montrer que $f$ est croissante, puis $f = \mathrm{id}_{\R}$.
\finenonce{003882}



\finexercice


\exercice{3883, quercia, 2010/03/11}
\enonce{003883}{Point fixe pour une application croissante}

Soit $f : {[0,1]} \to {[0,1]}$ croissante. Montrer qu'il existe $x \in {[0,1]}$
tel que $f(x) = x$.

(\'Etudier $A = \{ x \in [0,1] \text{ tq } f(x) \le x \}$)

\finenonce{003883}



\finexercice


\exercice{3884, quercia, 2010/03/11}
\enonce{003884}{Fonction localement monotone à droite}
Soit $f : \R \to \R$ continue telle que~:
$\forall\ x\in\R,\ \exists\ \delta>0,\text{ tq }\forall\ y\in[x,x+\delta],\ f(y)\ge f(x)$.
Montrer que $f$ est croissante.

\finenonce{003884}



\finexercice


\exercice{3885, quercia, 2010/03/11}
\enonce{003885}{Fonction affine}
Soit $f : \R \to \R$ vérifiant~:
$\forall\ x,y,z\in\R,\ |x-y|<|x-z|  \Rightarrow  |f(x)-f(y)|<|f(x)-f(z)|$.
Montrer successivement que $f$ est injective, monotone, continue, et enfin
affine.

\finenonce{003885}



\finexercice


\exercice{3886, quercia, 2010/03/11}
\enonce{003886}{Calcul de limite}

Montrer que : $\forall\ x > 0,\ x - \frac {x^2}2 < \ln(1+x) < x$.

En déduire
$\lim_{n\to\infty} \left(1+\frac 1{n^2}\right)\left(1+\frac 2{n^2}\right) \dots \left(1+\frac n{n^2}\right)$.

\finenonce{003886}



\finexercice


\exercice{3887, quercia, 2010/03/11}
\enonce{003887}{Dérivées de $\exp(-1/x)$}

On pose $f(x) = \begin{cases}\exp\left(-\frac 1x\right)& \text{ si } x > 0 \cr
                       0                          & \text{ si } x = 0.\cr\end{cases}$

\begin{enumerate}
  \item
Montrer que $f$ est de classe $\mathcal{C}^\infty$ sur $\R^{+*}$, et que $f^{(n)}(x)$ est de la
forme $\frac {P_n(x)}{x^{2n}}\exp\left(-\frac 1x\right)$ où $P_n$ est une fonction
polynomiale de degré inférieur ou égal à $n-1$ ($n \ge 1$).

  \item   Montrer que $f$ est de classe $\mathcal{C}^\infty$ en $0^+$.
  \item   Montrer que le polynôme $P_n$ possède $n-1$ racines dans $\R^{+*}$.

\end{enumerate}
\finenonce{003887}



\finexercice


\exercice{3888, quercia, 2010/03/11}
\enonce{003888}{$(1+1/t)^t$}

\begin{enumerate}
  \item Montrer que : $\forall\ t > 1,\ \left(1+\frac 1t\right)^t < e < \left(1+\frac 1{t-1}\right)^t$.
    
  \item Montrer que : $\forall\ x,y>0,\ \left(1+\frac xy\right)^y < e^x < \left(1+\frac xy\right)^{x+y}$.
    

\end{enumerate}
\finenonce{003888}



\finexercice


\exercice{3889, quercia, 2010/03/11}
\enonce{003889}{$\ln(1+ax)/\ln(1+bx)$}

Soient $0 < a < b$. Montrer que la fonction
$f : {\R^{+*}} \to {\R}, x \mapsto {\frac {\ln(1+ax)}{\ln(1+bx)}}$ est
croissante.
\finenonce{003889}



\finexercice


\exercice{3890, quercia, 2010/03/11}
\enonce{003890}{Inégalité}

Soient $0 < a < b$. Montrer que : $\forall\ x > 0,\ ae^{-bx} - be^{-ax} > a-b$.

\finenonce{003890}



\finexercice


\exercice{3891, quercia, 2010/03/11}
\enonce{003891}{Formules d'addition pour les fonctions hyperboliques}

Calculer $\ch(a+b)$, $\sh(a+b)$, $\tanh(a+b)$ en fonction de
$\ch a$, $\sh a$, $\tanh a$, $\ch b$, $\sh b$, $\tanh b$.
\finenonce{003891}



\finexercice


\exercice{3892, quercia, 2010/03/11}
\enonce{003892}{Simplification de $a\ch x + b\sh x$}

Soient $a,b \in \R$ non tous deux nuls.

\begin{enumerate}
  \item Peut-on trouver $A,\varphi \in \R$ tels que :
    $\forall\ x \in \R,\ a\ch(x)+b\sh(x) = A\ch(x+\varphi)$ ?
    
  \item Peut-on trouver $A,\varphi \in \R$ tels que :
    $\forall\ x \in \R,\ a\ch(x)+b\sh(x) = A\sh(x+\varphi)$ ?
    
\end{enumerate}
\finenonce{003892}



\finexercice


\exercice{3893, quercia, 2010/03/11}
\enonce{003893}{Somme de ch}

Calculer $\sum_{k=0}^n\,\ch(kx)$.


\finenonce{003893}



\finexercice


\exercice{3894, quercia, 2010/03/11}
\enonce{003894}{Somme de sh}

Soit $a \in \R$. Résoudre :
$\sh a + \sh(a+x) + \sh(a+2x) + \sh(a+3x) = 0$.

\finenonce{003894}



\finexercice


\exercice{3895, quercia, 2010/03/11}
\enonce{003895}{Somme de th}

Soit $x \in \R^*$. Vérifier que $\tanh x = 2\coth 2x - \coth x$.
En déduire la convergence et la somme de la série de terme général
$\frac1{2^n}\tanh\left(\frac{x}{2^n}\right)$.


\finenonce{003895}



\finexercice


\exercice{3896, quercia, 2010/03/11}
\enonce{003896}{Somme de $1/\sh$}

Soit $x \in \R^*$. Vérifier que $\frac1{\sh x} = \coth \frac x2 - \coth x$.
En déduire la convergence et la somme de la série de terme général
$\frac1{\sh(2^nx)}$.

\finenonce{003896}



\finexercice


\exercice{3897, quercia, 2010/03/11}
\enonce{003897}{$\ch(nx)$ et $\sh(nx)$}

Montrer que les fonctions : $x  \mapsto \ch(n\Argch(x))$ et
$x  \mapsto \frac{\sh(n\Argch(x))}{\sqrt{x^2-1}}$ $(n \in \N)$ sont
polynomiales.
\finenonce{003897}



\finexercice


\exercice{3898, quercia, 2010/03/11}
\enonce{003898}{$\ch x + \ch y = a$, $\sh x + \sh y = b$}

Soient $a,b\in \R$.
\'Etudier l'existence de solutions pour le système :
$\begin{cases} \ch x+\ch y &= a\cr \sh x+\sh y &= b.\cr\end{cases}$



\finenonce{003898}



\finexercice


\exercice{3899, quercia, 2010/03/11}
\enonce{003899}{Relation entre les fonctions hyperboliques et circulaires}

Soit $y \in \left]-\frac \pi2, \frac \pi2\right[$.
On pose $x = \ln\biggl(\tan\Bigl( \frac y2 + \frac \pi4\Bigr) \biggr)$.

Montrer que $\tanh \frac x2 = \tan \frac y2$,\quad
    $\tanh x = \sin y$,\quad $\ch x = \frac 1{\cos y}$.
\finenonce{003899}



\finexercice


\exercice{3900, quercia, 2010/03/11}
\enonce{003900}{$\Argth((1+3\tanh x)/(3+\tanh x))$}

Simplifier $\Argth\left(\frac {1+3\tanh x}{3+\tanh x}\right)$.

\finenonce{003900}



\finexercice


\exercice{3901, quercia, 2010/03/11}
\enonce{003901}{Équations diverses}

Résoudre $\Argch x = \Argsh(x-\frac12)$.



\finenonce{003901}



\finexercice


\exercice{3902, quercia, 2010/03/11}
\enonce{003902}{Calcul de primitives}

Déterminer des primitives des fonctions suivantes :

\begin{enumerate}
  \item   $f(x) = \frac 1{\sqrt{x^2+x+1}}$.
      
  \item   $f(x) = \frac 1{x^2+x-1}$.
      

\end{enumerate}
\finenonce{003902}



\finexercice


\exercice{3903, quercia, 2010/03/11}
\enonce{003903}{Racine d'une somme d'exponentielles}

Soient $0 < a_1 < a_2 < \dots < a_p$ des réels fixés.

\begin{enumerate}
  \item   Montrer que pour tout réel $a > a_p$ il existe un unique réel $x_a > 0$
      solution de l'équation : $a_1^x + \dots + a_p^x = a^x$.
      
  \item   Pour $a < b$, comparer $x_a$ et $x_b$.
      
  \item   Chercher $\lim_{a\to+\infty} x_a$ puis
      $\lim_{a\to+\infty} x_a\ln a$
      
\end{enumerate}
\finenonce{003903}



\finexercice


\exercice{3904, quercia, 2010/03/11}
\enonce{003904}{Centrale MP 2000}
Soit ${f} : {\R^{+*}} \to {\R^{+*}}$ telle que : 
$\forall\ x,y>0,\ f(xf(y))=yf(x)$ et $f(x)\to +\infty$ lorsque $x\to 0^+$
\begin{enumerate}
  \item Montrer que $f$ est involutive.
    
  \item Montrer que $f$ conserve le produit. Que peut-on dire de la monotonie de $f$, de sa continuité ? 
    
  \item Trouver $f$.
    
\end{enumerate}
\finenonce{003904}



\finexercice


\exercice{3905, quercia, 2010/03/11}
\enonce{003905}{\'Equations trigonométriques}
Résoudre les équations suivantes : 

% ---------------------------------------------- (–6+–2)cos‚ + (–6-–2)sin‚ = 2
\begin{enumerate}
  \item $(\sqrt6 + \sqrt2)\cos\theta + (\sqrt6 - \sqrt2)\sin\theta = 2$.

% ------------------------------------------- sin‚ + sin2‚ + sin3‚ + sin4‚ = 0
  \item $\sin\theta + \sin2\theta + \sin3\theta + \sin4\theta = 0$.

% ---------------------------------------------------- cos‚ + cos2‚ + cos3‚ = 0
  \item $\cos\theta + \cos2\theta + \cos3\theta = 0$.

% -------------------------------------------------------- cos‚ - cos2‚ = sin3‚
  \item $\cos\theta - \cos2\theta = \sin3\theta$.

% -------------------------------------------------------- cos‚ + cos7‚ = cos4‚
  \item $\cos\theta + \cos7\theta = \cos4\theta$.

% ----------------------------------------------------- cos2‚ + cos12‚ = –3cos5‚
  \item $\cos2\theta + \cos12\theta = \sqrt3\cos5\theta$.

% -------------------------------------------------------- sin7‚ - sin‚ = sin3‚
  \item $\sin7\theta - \sin\theta = \sin3\theta$.

% ----------------------------------------------- cos‚ sin3‚ + cos3‚sin‚  = 3/4
  \item $\cos^3\theta\sin3\theta + \cos3\theta\sin^3\theta = \frac 34$.

% ------------------------------------------------------ sin‚sin3‚ = sin5‚sin7‚
  \item $\sin\theta\sin3\theta = \sin5\theta\sin7\theta$.

% --------------------------------------------------------------- 3tan‚ = 2cos‚
  \item $3\tan\theta = 2\cos\theta$.

% --------------------------------------------------------------- tan4‚ = 4tan‚
  \item $\tan4\theta = 4\tan\theta$.

% ------------------------------------------------- cotan‚ - tan‚ = cos‚ + sin‚
  \item $\mathrm{cotan}\theta - \tan\theta = \cos\theta + \sin\theta$.

% ----------------------------------------- tan(x) + tan(y) = 1, tan(x+y) = 4/3
  \item $\begin{cases}\tan x + \tan y = 1\cr \tan(x+y) = 4/3.\end{cases}$
\end{enumerate}
\finenonce{003905}



\finexercice


\exercice{3906, quercia, 2010/03/11}
\enonce{003906}{Inéquations}

% ------------------------------------------------------- cos‚ + cos(‚+ã/3) > 0
\begin{enumerate}
  \item Résoudre : $\cos\theta + \cos(\theta + \pi/3) > 0$.
    

% ------------------------------------------------------------ 2cos‚ + sin‚ < 2
  \item Résoudre : $2\cos\theta + \sin\theta < 2$.
    
\end{enumerate}
\finenonce{003906}



\finexercice


\exercice{3907, quercia, 2010/03/11}
\enonce{003907}{Linéarisation}
\  \\

$2\cos^2\theta = 1 + \cos2\theta$. \\
$2\sin^2\theta = 1 - \cos2\theta$. \\
$4\cos^3\theta = 3\cos\theta + \cos3\theta$. \\
$4\sin^3\theta = 3\sin\theta - \sin3\theta$. \\

$8\cos^4\theta = 3 + 4\cos2\theta + \cos4\theta$. \\
$8\sin^4\theta = 3 - 4\cos2\theta + \cos4\theta$. \\
$32\cos^6\theta = 10 + 15\cos2\theta + 6\cos4\theta + \cos6\theta$. \\
$32\sin^6\theta = 10 - 15\cos2\theta + 6\cos4\theta - \cos6\theta$. \\


$32\cos^4\theta\sin^2\theta = 2 + \cos2\theta - 2\cos4\theta - \cos6\theta$. \\
$32\sin^4\theta\cos^2\theta = 2 - \cos2\theta - 2\cos4\theta + \cos6\theta$. \\
$16\cos\theta\sin^4\theta = \cos5\theta - 3\cos3\theta + 2\cos\theta$. \\
$16\sin\theta\cos^4\theta = \sin5\theta + 3\sin3\theta + 2\sin\theta$. \\

$4\sin\theta \sin\left(\frac\pi3-\theta\right)\sin\left(\frac\pi3+\theta\right) = \sin3\theta$.
\finenonce{003907}



\finexercice


\exercice{3908, quercia, 2010/03/11}
\enonce{003908}{$\alpha+\beta+\gamma = \pi$}

Soient $\alpha,\beta,\gamma \in \R$ tels que $\alpha+\beta+\gamma = \pi$.

\begin{enumerate}
  \item
Démontrer que : $1 - \cos\alpha + \cos\beta + \cos\gamma =
                 4\sin\frac\alpha2 \cos\frac\beta2 \cos\frac\gamma2$.

  \item
 Simplifier $\tan\frac\alpha2 \tan\frac\beta2  +
             \tan\frac\beta2  \tan\frac\gamma2 +
             \tan\frac\gamma2 \tan\frac\alpha2 $.

\end{enumerate}
\finenonce{003908}



\finexercice


\exercice{3909, quercia, 2010/03/11}
\enonce{003909}{$\sin^2(\theta-\alpha), \sin^2\theta, \sin^2(\theta+\alpha)$ en progression arithmétique}

Montrer qu'il existe $\theta \in {]0, \frac \pi2 [}$ tel que pour tout
$\alpha \in \R$, les nombres
$\sin^2(\theta-\alpha), \sin^2\theta, \sin^2(\theta+\alpha)$
soient en progression arithmétique.
\finenonce{003909}



\finexercice


\exercice{3910, quercia, 2010/03/11}
\enonce{003910}{Calcul de somme}

Calculer $\tan p - \tan q$. En déduire la valeur de
$S_n = \sum_{k=1}^n \frac 1{\cos(k\theta)\cos((k+1)\theta)}$, $\theta \in \R$.
\finenonce{003910}



\finexercice


\exercice{3911, quercia, 2010/03/11}
\enonce{003911}{Calcul de somme}

Simplifier $\sum_{k=0}^{n-1} 3^k\sin^3\left(\frac\alpha{3^{k+1}}\right)$.

\finenonce{003911}



\finexercice


\exercice{3912, quercia, 2010/03/11}
\enonce{003912}{Calcul de somme}

Calculer $\mathrm{cotan} x - 2\mathrm{cotan}2x$.
Simplifier $\sum_{k=0}^n \frac 1{2^k}\tan\frac\alpha{2^k}$.

\finenonce{003912}



\finexercice


\exercice{3913, quercia, 2010/03/11}
\enonce{003913}{Heptagone régulier}

Soit $ABCDEFG$ un heptagone (7 côtés) plan régulier. On pose
$\alpha = AB$, $\beta = AC$, $\gamma = AD$ (distances).

Montrer que $\frac 1\alpha = \frac 1\beta + \frac 1\gamma$.

\finenonce{003913}



\finexercice


\exercice{5097, rouget, 2010/06/30}
\enonce{005097}{**I}
\begin{enumerate}
\item  Soit $f$ une fonction dérivable sur $\Rr$ à valeurs dans $\Rr$. Montrer que si $f$ est paire, $f'$ est
impaire et si $f$ est impaire, $f'$ est paire.
\item  Soient $n\in\Nn^*$ et $f$ une fonction $n$ fois dérivable sur $\Rr$ à valeurs dans $\Rr$. $f^{(n)}$
désignant la dérivée $n$-ième de $f$, montrer que si $f$ est paire, $f^{(n)}$ est paire si $n$ est pair et impaire si
$n$ est impair.
\item  Soit $f$ une fonction continue sur $\Rr$ à valeurs dans $\Rr$. A-t-on des résultats analogues concernant les
primitives de $f$~?
\item  Reprendre les questions précédentes en remplaçant la condition \og~$f$ est paire (ou impaire)~\fg~par la
condition \og~$f$ est $T$-périodique~\fg.
\end{enumerate}
\finenonce{005097}


\finexercice


\exercice{5098, rouget, 2010/06/30}
\enonce{005098}{**}
Trouver la plus grande valeur de $\sqrt[n]{n}$, $n\in\Nn^*$.
\finenonce{005098}


\finexercice


\exercice{5099, rouget, 2010/06/30}
\enonce{005099}{**I}
\begin{enumerate}
\item  Etudier brièvement la fontion $x\mapsto\frac{\ln x}{x}$ et tracer son graphe.
\item  Trouver tous les couples $(a,b)$ d'entiers naturels non nuls et distincts vérifiant $a^b=b^a$.
\end{enumerate}
\finenonce{005099}


\finexercice


\exercice{5100, rouget, 2010/06/30}
\enonce{005100}{}
Résoudre dans $\Rr$ les équations ou inéquations suivantes~:
\begin{enumerate}
 \item $(**)\;\ln|x+1|-\ln|2x+1|\leq\ln2$,
 \item $(*)\;x^{\sqrt{x}}=\sqrt{x}^x$,
 \item $(**)\;2\Argsh x=\Argch3-\Argth\frac{7}{9}$,
 \item $(**)\;\mbox{ln}_x(10)+2\mbox{ln}_{10x}(10)+3\mbox{ln}_{100x}(10)=0$,
 \item $(**)\;2^{2x}-3^{x-\frac{1}{2}}=3^{x+\frac{1}{2}}-2^{2x-1}$.
\end{enumerate}
\finenonce{005100}


\finexercice

\exercice{5102, rouget, 2010/06/30}
\enonce{005102}{}
Construire le graphe des fonctions suivantes~:
\begin{enumerate}
\item (*) $f_1(x)=2|2x-1|-|x+2|+3x$.
\item (**) $f_{2}(x)=\ln(\ch x)$.
\item (***) $f_3(x)=x+\sqrt{|x^2-1|}$.
\item (**) $f_4(x)=|\tan x|+\cos x$.
\item (***) $f_5(x)=\left(1+\frac{1}{x}\right)^x$ (à étudier sur $]0,+\infty[$).
\item (**) $f_6(x)=\mbox{log}_2(1-\mbox{log}_{\frac{1}{2}}(x^2-5x+6))$.
\end{enumerate}
\finenonce{005102}


\finexercice


\exercice{5404, rouget, 2010/07/06}
\enonce{005404}{**}
Soit $f$ de $[0,1]$ dans lui-même telle que $\forall(x,y)\in([0,1])^2,\;|f(y)-f(x)|\geq|x-y|$. Montrer que $f=Id$ ou $f=1-Id$.
\finenonce{005404}


\finexercice


\exercice{5443, rouget, 2010/07/10}
\enonce{005443}{}

Etude complète des fonctions suivantes

\begin{enumerate}
\item  $f_1(x)=\frac{1+x^2}{x^3}(\Arctan x-\frac{x}{1+x^2})$.
\item  $f_2(x)=|\tan x|+\cos x$.
\item  $f_3(x)=x-\ln\left|\frac{120+60x+12x^2+x^3}{120-60x+12x^2-x^3}\right|$.
\item  $4_(x)=xe^{\frac{2x}{x^2-1}}$. 
\item  $f_5(x)=\frac{1}{x}\ln\left(\frac{e^x-1}{x}\right)$.
\item  $f_6(x)=x+\sqrt{|x^2-1|}$. 
\item  $f_7(x)=e^{/\ln x}$. 
\item  $f_8(x)=\left(1+\frac{1}{x}\right)^x$. 
\item  $f_9(x)=\mbox{log}_2(1-\mbox{log}_{\frac{1}{2}}(x^2-5x+6))$.
\item  $f_{10}(x)=E(x)+(x-E(x))^2$.
\item  $f_{11}(x)=\Arcsin\sqrt{\frac{1}{2}-x}+\Arcsin\sqrt{\frac{1}{2}+x}$.
\item  $f_{12}(x)=\frac{\Arcsin x}{x}$. 
\item  $f_{13}(x)=e^{1/x}\sqrt{x+4}$. 
\item  $f_{14}(x)=\Arccos(\frac{1}{\ch x})$.
\item  $f_{15}(x)=\ln(y+\sqrt{y^2-1})-\ln(\frac{1+x}{1-x})$ où $y=\frac{1+x^2}{1-x^2}$.
\item  $f_{16}(x)=\ln|\sh x-1|$.
\item  $f_{17}(x)=x^{(x^x)}$. 
\item  $f_{18}(x)=(\cos x+\sin x)^{1/x}$. 
\item  $f_{19}(x)=\sqrt[3]{x^3+1}-\sqrt{x^2-1}$. 
\item  $f_{20}(x)=\Arcsin(2x-1)+2\Arctan\sqrt{\frac{1-x}{x}}$.
\item  $f_{21}(x)=\ln(\ch x)$.
\item  $f_{22}(x)=3^{2x-1}-5.3^{x-1}-x\ln3$.
\item  $f_{23}(x)=\ln\left|\frac{1}{e^x-1}\right|$.
\end{enumerate}
\finenonce{005443}


\finexercice

\section{123.05 Fonction continue par morceaux}


\exercice{693, gourio, 2001/09/01}

\enonce{000693}{}
Soit $g:[a,b]\rightarrow {\Rr}$ une fonction telle que :
$$\forall \epsilon >0,\exists  \phi \in CM\left( [a,b],{\Rr}\right)
,\forall x\in [a,b],\left| g(x)-\phi (x)\right| <\epsilon . $$
Montrer que l'on peut choisir $\phi \in E\left( [a,b],{\Rr}\right)$, ie :
$$\forall \epsilon >0,\exists  \phi \in E\left( [a,b],{\Rr}\right)
,\forall x\in [a,b],\left| g(x)-\phi (x)\right| <\epsilon . $$
NB : CM pour continue par morceaux et E pour escalier.
\finenonce{000693}



\finexercice



\exercice{694, gourio, 2001/09/01}

\enonce{000694}{}
Donner un exemple de fonction qu'on ne puisse approcher \`{a} $\epsilon$
pr\`{e}s par des fonctions en escaliers.
\finenonce{000694}



\finexercice



\exercice{695, gourio, 2001/09/01}

\enonce{000695}{}
On dit qu'un ensemble $A$ de fonctions d\'{e}finies sur un intervalle $I=[a,b]$
de $\Rr$ est dense dans un ensemble $B$ si :
$$\forall f\in B,\forall \epsilon >0,\exists g\in A,\forall x\in I,\left|
f(x)-g(x)\right| <\epsilon . $$
Le cours dit par exemple que l'ensemble des fonctions en escaliers est dense
dans l'ensemble des fonctions continues par morceaux si $I=[a,b].$
Montrer que l'ensemble des fonctions continues affines par morceaux est
dense dans l'ensemble des fonctions continues sur un intervalle $I=[a,b]$.
\finenonce{000695}



\finexercice



\exercice{696, gourio, 2001/09/01}

\enonce{000696}{}
On dit qu'une suite $(f_{n})_{n\in \Nn}$ de fonctions d\'{e}finies sur
$I=[a,b]$ converge uniform\'{e}ment vers $f$ si :
$$\forall \epsilon >0,\exists \Nn\in \Nn,\forall n\geq \Nn,\forall x\in
I,\left| f_{n}(x)-f(x)\right| <\epsilon . $$

On suppose que $(f_{n})_{n\in \Nn}$ converge uniform\'{e}ment vers $f$
sur l'intervalle $[a,b]$, et que toutes les $f_{n}$ sont continues.\ Montrer
que $\forall x\in [a,b],$ la suite $(f_{n}(x))_{n\in \Nn}$ est
convergente, et donner sa limite. Montrer que $f$ est born\'{e}e et
continue.

On ne suppose plus que $(f_{n})_{n}$ converge uniform\'{e}ment mais
seulement point par point (ie, $\forall x\in [a,b],$ la suite $(f_{n}(x))_{n\in \Nn}$
 est convergente vers $f(x)$); de plus toutes les $f_{n}$
sont lipschitziennes de rapport $k$ ; montrer que $f$ est
lipschitzienne de rapport $k$ et qu'il y a converge uniforme.
\finenonce{000696}



\finexercice



\exercice{697, gourio, 2001/09/01}

\enonce{000697}{}
$f:[a,b]\rightarrow {\Rr}$ est \`{a} variation born\'{e}e si et seulement
si :
$$\exists  \mu \in {\Rr}^{+},\forall  d=\{a=x_{0},x_{1},...,x_{n}=b\}\text{
subdivision de }[a,b],\sum\limits_{i=1}^{n}\left| f(x_{i})-f(x_{i-1})\right|
=\sigma (d)\leq \mu .$$
On appelle alors $V(a,b)=\sup\limits_{d \text{ subdivision}}\sigma (d)$ et
on d\'{e}finit une fonction de $[a,b]$ dans ${\Rr}^{+}:x\rightarrow V(a,x). $

Montrer que toute fonction monotone est \`{a} variation born\'{e}e puis que $%
x\rightarrow V(a,x)$ est croissante ainsi que $x\rightarrow V(a,x)-f(x).$ En
d\'{e}duire que toute fonction \`{a} variation born\'ee est la diff\'{e}rence de deux
fonctions croissantes (d'o\`{u} la nature de ses discontinuit\'{e}s). Une
fonction continue, une fonction lipschitzienne sont-elles \`{a} variation born\'ee ?
\finenonce{000697}



\finexercice


\section{123.06 Fonctions équivalentes, fonctions négligeables}


\exercice{1213, ridde, 1999/11/01}

\enonce{001213}{}
\`A quelle condition sur $f$ et $g$ a-t-on $e^f \underset{a}{\sim}  e^g$ ?
\finenonce{001213}



\finexercice



\exercice{1214, ridde, 1999/11/01}

\enonce{001214}{}
Soient $f$ et $g$ \'equivalentes au voisinage de $a$ et strictement positives.
Montrer que si $f$ admet en $a$ une limite dans $\bar{\Rr}$ diff\'erente de $1$ alors
$\ln f \underset{a}{\sim} \ln g$.
\finenonce{001214}



\finexercice



\exercice{1215, ridde, 1999/11/01}

\enonce{001215}{}
Montrer que si $f$ tend vers $0$ en $a$ alors $\ln (1  +  f) \underset{a}{\sim} f$
et $e^f-1 \underset{a}{\sim} f$.
\finenonce{001215}



\finexercice



\exercice{1216, ridde, 1999/11/01}

\enonce{001216}{}
 \'Etudier en $ + \infty$ et $-\infty$ la fonction
$f (x) = \sqrt[3]{x^3 + 1}
 + \sqrt{x^2 + x + 1}$.

\finenonce{001216}



\finexercice



\exercice{1217, ridde, 1999/11/01}

\enonce{001217}{}
 Calculer les limites de
\begin{enumerate}
\item $\dfrac{\sin x \ln (1 + x^2)}{x \tan x} \text{ en } 0$.
\item $\dfrac{\ln (1 + \sin x)}{\tan (6x)} \text{ en } 0$.
\item $ (\ln (e + x))^{\frac 1x} \text{ en } 0$.
\item $(\ln (1 + e^{-x}))^{\frac 1x} \text{ en }  + \infty$.
\end{enumerate}

\finenonce{001217}



\finexercice



\exercice{1218, ridde, 1999/11/01}

\enonce{001218}{}
Trouver un \'equivalent simple en $ + \infty$ de $(\dfrac{\ln (1 + x)}{\ln x})^x-1$.
\finenonce{001218}



\finexercice



\exercice{1219, ridde, 1999/11/01}

\enonce{001219}{}
$$\text{Limite en } + \infty \text{ de }\sqrt [3]{{x}^{3}+{x}^{2}}-\sqrt [3]{{x}^{3}-{x}^{2}}$$

$$\text{\'Equivalent en }  + \infty \text{ de } \sqrt {{x}^{2}+\sqrt {{x}^{4}+1}}-x\sqrt {2}$$

$$\text{Limite en 0 de }{\frac {\tan(ax)-\sin(ax)}{\tan(bx)-\sin(bx)}} $$

$$\text{Limite en } \frac{\pi}4 \text{ de } \left (x-{\frac {\pi }{4}}\right )\tan(x+{\frac {\pi }{4}}) $$

$$ \text{Limite en } \frac{\pi}4 \text{ de }  {\frac {\cos(x)-\sin(x)}{\left (4\,x-\pi \right )\tan(x)}}$$

$$ \text{\'Equivalent en } 0 \text{ de }  {\frac {\tan(x-x\cos(x))}{\sin(x)+\cos(x)-1}} $$

$$ \text{\'Equivalent en } \frac{\pi}4  \text{ de }
\left (\tan(2\,x)+\tan(x+{\frac {\pi }{4}})\right )\left (
\cos(x+{\frac {\pi }{4}})\right )^{2} $$

$$ \text{Limite en 0 de }{x}^{\frac 1{1+2\,\ln (x)}} $$

$$\text{Limite en } \frac 12 \text{ de }\left (2\,{x}^{2}-3\,x+1\right )\tan(\pi \,x)
 $$

$$ \text{Limite en 0 de }{\frac {\left (\sin(x)\right )^{\sin(x)}-1}{\left (\tan(x)
\right )^{\tan(x)}-1}} $$

$$\text{\'Equivalent en }  + \infty \text{ de }\frac{\sqrt{1 + x^2}}{\sin (\frac 1x)} \ln
 (\frac x{x + 1}) $$

\finenonce{001219}



\finexercice



\exercice{1220, gourio, 2001/09/01}

\enonce{001220}{}
Soit $(f_{n})_{n\in \Nn}$ une suite de fonctions r\'{e}elles. Montrer qu'il
existe $f:\Rr\rightarrow \Rr$ telle que $\forall n\in \Nn,f_{n}(t)=o(f(t))$ si
$t\rightarrow \infty .$
\finenonce{001220}



\finexercice


\section{123.99 Autre}


\exercice{5161, rouget, 2010/06/30}
\enonce{005161}{***I}
Montrer que $\forall n\in\Nn,\;(n\geq3\Rightarrow\sqrt{n}<\sqrt[n]{n!})$.

(commencer par vérifier que pour $k=2,3,...,n$,
on a~:~$(n-k+1)k>n$).
\finenonce{005161}


\finexercice

\exercice{5163, rouget, 2010/06/30}
\enonce{005163}{**I}
Montrer que $\forall n\in\Nn,\;\forall x\in\Rr,\;|\sin(nx)|\leq n|\sin x|$.
\finenonce{005163}


\finexercice

\exercice{5393, rouget, 2010/07/06}
\enonce{005393}{**I}
\label{exo:rouconti}
Soit $f$ continue sur $[a,b]$ à valeurs dans $[a,b]$. Montrer que $f$ a un point fixe.
\finenonce{005393}


\finexercice


\exercice{5394, rouget, 2010/07/06}
\enonce{005394}{**I}
Soit $f$ définie sur $[0,+\infty[$ à valeurs dans $[0,+\infty[$, continue sur $[0,+\infty[$ telle que $\frac{f(x)}{x}$ a une limite réelle $\ell\in[0,1[$ quand $x$ tend vers $+\infty$. Montrer que $f$ a un point fixe.
\finenonce{005394}


\finexercice


\exercice{5397, rouget, 2010/07/06}
\enonce{005397}{***}
Soit $f$ continue sur $\Rr^+$ telle que, pour tout réel positif $x$, on ait $f(x^2)=f(x)$. Montrer que $f$ est constante sur $\Rr^+$. Trouver un exemple où $f$ n'est pas constante.
\finenonce{005397}


\finexercice


\exercice{5399, rouget, 2010/07/06}
\enonce{005399}{***I}
Trouver tous les morphismes continus de $(\Rr,+)$.
\finenonce{005399}


\finexercice


\exercice{5400, rouget, 2010/07/06}
\enonce{005400}{***}
Soient $a$ et $b$ deux réels tels que $0<a<b$. Montrer que $\bigcup_{k\geq1}]ka,kb[$ contient un intervalle de la forme $]A,+\infty[$ puis déterminer la plus petite valeur possible de $A$.
\finenonce{005400}


\finexercice


\exercice{5402, rouget, 2010/07/06}
\enonce{005402}{*** Théorème d'homéomorphie}
Soit $f$ une application continue sur un intervalle $I$ de $\Rr$ à valeurs dans $\Rr$. Montrer que $f$ est injective si et seulement si $f$ est strictement monotone et que dans ce cas $f(I)$ est un intervalle de même nature que $I$ (ouvert, semi-ouvert, fermé).
\finenonce{005402}


\finexercice


\exercice{5403, rouget, 2010/07/06}
\enonce{005403}{***}
Trouver un exemple de fonction pèriodique dont le groupe des pèriodes est dense dans $\Rr$ mais pas $\Rr$.
\finenonce{005403}


\finexercice


\exercice{5405, rouget, 2010/07/06}
\enonce{005405}{***}
Trouver les fonctions bijectives de $[0,1]$ sur lui-même vérifiant $\forall x\in[0,1],\;f(2x-f(x))=x$.
\finenonce{005405}


\finexercice


\exercice{5406, rouget, 2010/07/06}
\enonce{005406}{***I}
Soit $f$ une application de $[0,1]$ dans $\Rr$, continue sur $[0,1]$ et vérifiant $f(0)=f(1)$.
\begin{enumerate}
\item  Soit $n$ un entier naturel non nul et soit $a=\frac{1}{n}$. Montrer que l'équation $f(x+a)=f(x)$ admet au moins une solution.
\item  Montrer (en fournissant une fonction précise) que, si $a$ est un réel de $]0,1[$ qui n'est pas de la forme précédente, il est possible que l'équation $f(x+a)=f(x)$ n'ait pas de solution.
\item  Application. Un cycliste parcourt 20 km en une heure.
\begin{enumerate}
\item Montrer qu'il existe au moins un intervalle de temps de durée une demi-heure pendant lequel il a parcouru 10 km.
\item Montrer qu'il existe au moins un intervalle de temps de durée 3 min pendant lequel il a
parcouru 1 km.
\item Montrer qu'il n'existe pas nécessairement un intervalle de temps de durée 45 min pendant lequel il a parcouru 15 km.
\end{enumerate}
\end{enumerate}
\finenonce{005406}


\finexercice
\section{124.01 Calculs}


\exercice{698, bodin, 1998/09/01}
\video{yNbsJ0UALww}
\enonce{000698}{}
\'Etudier  la d\'erivabilit\'e des fonctions suivantes :
$$f_1(x)=x^2\cos \frac{1}{x}, \text{\ \  si }x\not=0 \qquad ; \qquad f_1(0)=0 ;$$

$$f_2(x)= \sin x \cdot \sin \frac{1}{x}, \text{\ \  si }x\not=0 \qquad ; \qquad f_2(0)=0 ;$$

$$f_3(x) = \frac{|x|\sqrt{x^2-2x+1}}{x-1}, \text{\ \  si } x\not= 1 \qquad ; \qquad f_3(1)=1.$$
\finenonce{000698}
 

\finexercice


\exercice{699, bodin, 1998/09/01}
\video{XPnRztcUKK4}
\enonce{000699}{}
D\'eterminer $a,b \in \Rr$ de mani\`ere \`a ce que la fonction $f$ d\'efinie sur $\Rr_+$ par :
$$ f(x)=\sqrt{x} \quad \text{ si } 0\leqslant x \leqslant 1 \quad \text{\ \  et\ \  } \quad f(x) = ax^2+bx+1 \quad \text{ si } x>1$$
soit d\'erivable sur $\Rr_+^*$.
\finenonce{000699} 


\finexercice


\exercice{700, bodin, 1998/09/01}
\video{K9yZt1R0xW8}
\enonce{000700}{}
Soit $f : \Rr^* \longrightarrow \Rr$ d\'efinie par
$\displaystyle{f(x)= x^2\sin \frac{1}{x} }$. Montrer que
$f$ est prolongeable par continuit\'e en $0$ ; on note encore
$f$ la fonction prolong\'ee. Montrer que $f$ est
d\'erivable sur $\Rr$ mais que $f'$ n'est pas continue en $0$.
\finenonce{000700} 


\finexercice


\exercice{701, bodin, 1998/09/01}
\enonce{000701}{}
Calculer la fonction d\'eriv\'ee d'ordre $n$ des fonctions $f,g,h$ d\'efinies par :
$$f(x) = \sin x \quad ; \quad g(x)=\sin^2x \quad ; \quad h(x)=\sin^3x+\cos^3x.$$
\finenonce{000701} 


\finexercice


\exercice{702, bodin, 1998/09/01}

\enonce{000702}{}
Calculer les d\'eriv\'ees d'ordre $n$ des fonctions :
$$f(x) = \frac{2x-5}{(x-2)^2(x+1)(x-3)}\qquad g(x) = \ln(1+x).$$
\finenonce{000702}



\finexercice



\exercice{703, bodin, 1998/09/01}

\enonce{000703}{Formule de Leibnitz}
\'Etant donn\'ees $u$ et $v$ des  fonctions d\'erivables \`a
l'ordre $n$ sur l'intervalle $I$, montrer par r\'ecurrence que la d\'eriv\'ee
d'ordre $n$ du produit $uv$ sur cet intervalle est :
$$(uv)^{(n)}=\sum_{k=0}^nC_n^k u^{(k)}v^{(n-k)}.$$

En d\'eduire les d\'eriv\'ees successives des fonctions :
$$ x\mapsto x^2e^x \quad;\quad x\mapsto x^2(1+x)^n \quad;\quad x\mapsto \frac{x^2+1}{(x+1)^2}\quad;\quad
x\mapsto x^{n-1}\ln x.$$
\finenonce{000703}



\finexercice



\exercice{704, cousquer, 2003/10/01}

\enonce{000704}{}
 Etudier la dérivabilité sur $\mathbb{R}$ des
applications suivantes~:
$$ f : x \mapsto x\vert x\vert, \qquad g : x \mapsto  \frac{x}{1+\vert x\vert}, \qquad  h :  \mapsto  \frac{1}{1+\vert x\vert}.$$
  
\finenonce{000704}



\finexercice



\exercice{705, cousquer, 2003/10/01}

\enonce{000705}{}
 Calculer les dérivées des fonctions~:
\begin{enumerate}
\item  $x \mapsto \sqrt{1+x^2 \sin^2 x},\qquad  x \mapsto \frac{\exp(1/x)+1} {\exp(1/x)-1}.$ 
\item  $x \mapsto \log (\frac{1+\sin(x)} {1-\sin(x)}),\qquad  x \mapsto (x(x-2))^{1/3}.$
\end{enumerate}
\finenonce{000705}



\finexercice



\exercice{706, cousquer, 2003/10/01}

\enonce{000706}{}
Soit $f$ une fonction dérivable sur $\mathbb{R}$.
\begin{enumerate}
\item  Calculer la dérivée de $x\mapsto \hbox{sin} (f(x)^2)$ et de
$ x \mapsto \hbox{sin} (f(x^2)).$
\item  On suppose $f(x)\not =0$ pour tout $x\in \mathbb{R}$. Calculer la
dérivée de $x\mapsto \hbox{log}(\vert f(x)\vert ).$ 
\end{enumerate}
\finenonce{000706}



\finexercice



\exercice{707, ridde, 1999/11/01}

\enonce{000707}{}
Prolonger par continuit\'e en $0$ et \'etudier la d\'erivabilt\'e de
\begin{enumerate}
\item $f (x) = \sqrt{x} \ln x$.
\item $g (x) = \dfrac{e^x-1}{\sqrt x}$.
\end{enumerate}
\finenonce{000707}



\finexercice



\exercice{708, ridde, 1999/11/01}

\enonce{000708}{}
Soit $ f : \begin{cases} \Rr \rightarrow \Rr \\
x \mapsto e^x \text{ si }x<0 \\
x \mapsto ax^2 + bx + c \text{ sinon}
\end{cases}$ \\
D\'eterminer $a, b, c$ pour que $f$ soit $C^2$ (et $C^3$ ?).
\finenonce{000708}



\finexercice



\exercice{709, ridde, 1999/11/01}

\enonce{000709}{}
Soit $f (x) = \exp (-\frac 1{x^2})$ si $x \neq 0$ et $f (0) = 0$.
Montrer que $f$ est $C^{\infty}$ et que $\forall n \in \Nn \, \, f^{ (n)} (0) = 0$.

\finenonce{000709}




\finexercice



\exercice{710, ridde, 1999/11/01}

\enonce{000710}{}
Soient $a$ et $b$ deux r\'eels et $f (x) = (x-a)^n (x-b)^n$.
Calculer $f^{ (n)}$ et en d\'eduire $\sum\limits_{k = 0}^n (C_n^k)^2$.
\finenonce{000710}



\finexercice



\exercice{711, gourio, 2001/09/01}

\enonce{000711}{}
Soit $f:{\Rr}\rightarrow {\Rr}$ d\'{e}finie par :
$$\forall x\neq 0,f(x)=e^{-\frac{1}{x^{2}}},f(0)=0. $$
Montrer que $f\in C^{\infty }({\Rr},{\Rr})$ et calculer ses
d\'{e}riv\'{e}es en 0.
\finenonce{000711}



\finexercice



\exercice{712, gourio, 2001/09/01}

\enonce{000712}{}
Calculer la d\'{e}riv\'{e}e de $x\rightarrow \ln \cos (\pi +\frac{x^{2}-1}{
x^{2}+1}).$
\finenonce{000712}



\finexercice



\exercice{713, gourio, 2001/09/01}

\enonce{000713}{}
La fonction $x\rightarrow \cos \sqrt{x} $ est-elle d\'{e}rivable en $0 $ ?
\finenonce{000713}



\finexercice



\exercice{714, gourio, 2001/09/01}

\enonce{000714}{}
En quels points la fonction $f:{\Rr}\rightarrow {\Rr}$ d\'{e}finie par :
$$\forall x\in \Qq,f(x)=x^{2},\forall x\in {\Rr}-\Qq,f(x)=0, $$
est-elle d\'{e}rivable ?
\finenonce{000714}



\finexercice



\exercice{1221, legall, 1998/09/01}

\enonce{001221}{}
Montrer que pour tout $  x \in { \Rr}_+  ,   \sin(x)  \leq x .$
\finenonce{001221}



\finexercice



\exercice{1222, legall, 1998/09/01}

\enonce{001222}{}
Pour tout $  x \in ] 1, +\infty [ $ on pose
 $  f (x) =x \ln (x) -x  .$
 Montrer que $  f  $ est une bijection de
$  ] 1, +\infty [ $ sur $  ] -1, +\infty [ .$
 On pose $  g=f^{-1}
$ l'application r\'eciproque de $  f  .$ Calculer $  g(0)  $ et $
g'(0)  .$
\finenonce{001222}



\finexercice



\exercice{1223, legall, 1998/09/01}

\enonce{001223}{}
\'Etudier la continuit\'e, la d\'erivabilit\'e, la continuit\'e de la d\'eriv\'ee pour les
applications suivantes~:
\begin{enumerate}
\item $  \displaystyle { f :x \mapsto \hbox{sin }(\frac{1}{ x }) }  $ si $  x\not = 0  $ et $
f(0)=0  .$
\item $  \displaystyle { g :x \mapsto x \hbox{sin }(\frac{1}{ x }) }  $ si $  x\not = 0  $ et $
f(0)=0  .$
\item $  \displaystyle { h :x \mapsto x^2 \hbox{sin }(\frac{1}{ x }) }  $ si $  x\not = 0  $ et $
f(0)=0  .$
\end{enumerate}
\finenonce{001223}



\finexercice



\exercice{1224, legall, 1998/09/01}

\enonce{001224}{}
Soit $  g  $ une fonction $  2   $ fois d\' erivable sur $  [a,b]  $ telle que
$  g(a)=g(b)=0  $ et $  g''(x)\leq 0  $ pour tout $  x\in ]a,b[  .$ Montrer que pour tout
$  x\in ]a,b[  ,   g(x)\geq 0  .$
\finenonce{001224}



\finexercice



\exercice{1225, legall, 1998/09/01}

\enonce{001225}{}
Soit $  f : { \Rr}\rightarrow { \Rr}  $ une fonction deux fois d\'erivable telle
que $  \forall x \in { \Rr}  $ on ait $  f(x)\geq 0  ,$ $  f'(x)\geq 0  $
et $  f''(x)\geq 0  .$ \'Etudier $  \displaystyle{ \lim _{x \rightarrow \infty } f(x)}  $
et $  \displaystyle{ \lim _{x \rightarrow \infty } \frac{f(x)}{ x}}  .$
\finenonce{001225}



\finexercice



\exercice{1226, legall, 1998/09/01}

\enonce{001226}{}
Soit $  f   $ une application continue de
$  [a,b]  $ \`a valeurs dans $  \R   $ d\'erivable sur $  ]a,b]  .$ Montrer que si
$  \displaystyle{ \lim _{x\rightarrow a} f'(x)}  $ existe, $  f  $ est d\'erivable en $  a  . $
\finenonce{001226}



\finexercice



\exercice{1227, legall, 2003/10/01}

\enonce{001227}{}
Soit $f:\Rr_+\rightarrow \Rr_+^*$ une fonction 
born\'ee deux fois d\'erivable et telle qu'il existe $\alpha > 0$
tel que, pour tout $x\in \Rr_+ ,$ on ait $\alpha f(x)\leq f''(x).$
\begin{enumerate}
        \item
\begin{enumerate}
        \item Montrer que $f'$ a une limite en $+\infty .$ Quelle est 
la valeur de cette limite~?
  \item Montrer que $f$ est d\'ecroissante et que $\displaystyle{ \lim 
_{+\infty} f(x)=0}.$
\end{enumerate}

\item
\begin{enumerate}
        \item Soit $g :x\mapsto \alpha f^2(x)-(f'(x))^2.$ Montrer que 
$g$ est croissante et a pour limite $0$ en $\infty$.
  \item En posant $f(x)=h(x)\exp (-\sqrt{\alpha}x)$, montrer que, pour 
tout $x\in \R _+ :f(x)\leq f(0)\exp (-\sqrt{\alpha}x)$.
\end{enumerate}
\end{enumerate}
\finenonce{001227}



\finexercice



\exercice{2689, matexo1, 2002/02/01}
\enonce{002689}{}
Montrer que
$$\forall x\in\R,\qquad |\sin x| \leq  |x|,$$
$$\forall x\in[0,\frac \pi 2], \qquad 1-\cos x \leq  x \sin x,$$
$$\forall x\in[-1,1],\qquad 
|\arcsin x| \leq  \left| \frac x{\sqrt{1-x^2}} \right|.$$
\finenonce{002689}
\finexercice


\exercice{5413, rouget, 2010/07/06}
\enonce{005413}{**}
Déterminer dans chacun des cas suivants la dérivée $n$-ème de la fonction proposée~:

$$1)\;x\mapsto x^{n-1}\ln(1+x)\;2)\;x\mapsto\cos^3x\sin(2x)\;3)\;x\mapsto\frac{x^2+1}{(x-1)^3}\;4)\;x\mapsto(x^3+2x-7)e^x.$$
\finenonce{005413}


\finexercice

\section{124.02 Théorème de Rolle et accroissements finis}


\exercice{715, bodin, 1998/09/01}
\video{J-dlr6vwO8A}
\enonce{000715}{}
Montrer que le polyn\^ome $P_n$ d\'efini par
$$P_n(t)=\left[ \left( 1-t^2 \right)^n \right]^{(n)}$$
est un polyn\^ome de degr\'e $n$ dont les racines sont
 r\'eelles, simples, et appartiennent \`a $[-1,1]$.
\finenonce{000715} 


\finexercice


\exercice{716, cousquer, 2003/10/01}

\enonce{000716}{}
 Etudier la fonction $f: x\mapsto x^5-5x+1$ sur $\mathbb{R}$ et en déduire que
l'équation
 $ x^5-5x+1=0 $ a trois solutions réelles.
\finenonce{000716}



\finexercice



\exercice{717, legall, 1998/09/01}
\video{w_zY_6IuJnE}
\enonce{000717}{}
Montrer que le polyn\^ome $  X^n+aX+b  $,  ($  a  $ et $  b  $ r\' eels) admet au plus trois racines r\' eelles.
\finenonce{000717} 


\finexercice

\exercice{718, legall, 1998/09/01}
\enonce{000718}{}
Soit $  f  $ une fonction $  n   $ fois d\' erivable sur $  ]a,b[  $
s'annulant en $  n+1  $ points de $  ]a,b[  .$ Montrer
que si $  f^{(n)}  $ est
continue, il existe un point $  x_0  $ de $
]a,b[  $ tel que $  f^{(n)}(x_0)=0  .$
\finenonce{000718} 


\finexercice

\exercice{719, bodin, 1998/09/01}

\enonce{000719}{}
\'Etant donn\'e $y$ un r\'eel positif
et $n$ un entier naturel pair, montrer que
$(x+y)^n=x^n+y^n$ si et seulement si $x=0$.
Cas $n$ impair ?
\finenonce{000719}



\finexercice



\exercice{720, bodin, 1998/09/01}

\enonce{000720}{}
Soit $f$ une fonction continue et d\'erivable sur
$[a,+\infty[$ et telle que
$\lim_{x\rightarrow\infty}f(x)=f(a)$.
Montrer qu'il existe un \'el\'ement $c$ dans $]a,+\infty [$
tel que $f'(c)=0$.
\finenonce{000720}



\finexercice



\exercice{721, bodin, 1998/09/01}
\video{wAM4vww5WSo}
\enonce{000721}{}
Dans l'application du th\'eor\`eme des accroissements finis
\`a la fonction
$$ f(x) =\alpha x^2+\beta x+\gamma$$
sur l'intervalle $[a,b]$
pr\'eciser le nombre ``$c$'' de $]a,b[$.
Donner une interpr\'etation g\'eom\'etrique.
\finenonce{000721} 


\finexercice

\exercice{722, bodin, 1998/09/01}

\enonce{000722}{}
Appliquer la formule des accroissements finis \`a la
fonction
$$f(x)=a+bx+ce^{\alpha x}$$
(où $a,b,c,\alpha$ sont r\'eels,
et $c$ et $\alpha$ sont non nuls)
sur l'intervalle $[0,X]$.
\begin{enumerate}
    \item Calculer ``$\theta$'' en fonction de $X$.
    \item En d\'eduire que
$$x\mapsto \frac{1}{\alpha x}\ln \frac{e^{\alpha x}-1}{\alpha x}$$
est born\'ee sur $\Rr$.
\end{enumerate}
\finenonce{000722}



\finexercice



\exercice{723, bodin, 1998/09/01}

\enonce{000723}{}
Soit $f$ une fonction deux fois d\'erivable sur $[a,a+2h]$.
Par introduction de la fonction
$$g(t) = f(a+t+h)-f(a+t)$$
montrer qu'il existe $\alpha$ dans $]0,2[$ tel que
$$f(a)-2f(a+h)+f(a+2h)=h^2f''(a+\alpha h).$$
\finenonce{000723}



\finexercice



\exercice{724, bodin, 1998/09/01}
\video{CfqCgc6zKgg}
\enonce{000724}{}
Soient $x$ et $y$ r\'eels avec $0<x<y$.
\begin{enumerate}
    \item Montrer que
$$ x < \frac{y-x}{\ln y - \ln x} < y.$$
    \item On consid\`ere la fonction $f$ d\'efinie sur
$[0,1]$ par
$$\alpha \mapsto f(\alpha) = \ln (\alpha x +(1-\alpha)y)-\alpha
\ln x -(1-\alpha)\ln y.$$
De l'\'etude de $f$ d\'eduire que pour tout $\alpha$ de $]0,1[$
$$ \alpha
\ln x +(1-\alpha)\ln y < \ln (\alpha x +(1-\alpha)y) .$$
Interpr\'etation g\'eom\'etrique ?
\end{enumerate}
\finenonce{000724} 


\finexercice


\exercice{725, bodin, 1998/09/01}
\enonce{000725}{}
Par application du th\'eor\`eme des accroissements finis \`a
$f(x) = \ln x$ sur $[n,n+1]$
montrer que
$$ S_n = \sum_{k=1}^n \frac{1}{k}$$
tend vers l'infini quand $n$ tend vers l'infini.
\finenonce{000725} 


\finexercice

\exercice{726, bodin, 1998/09/01}

\enonce{000726}{}
\'Etant donn\'e $\alpha$ dans $]0,1[$, montrer que pour
 tout entier naturel $n$
$$
\frac{\alpha}{(n+1)^{1-\alpha}} \ge (n+1)^\alpha - n^\alpha \ge \frac{\alpha}{n^{1-\alpha}}.
$$
En d\'eduire la limite
$$\lim_{n\rightarrow\infty}\sum_{p=1}^{n}\frac{1}{p^\alpha}.$$
\finenonce{000726}



\finexercice



\exercice{727, ridde, 1999/11/01}
\enonce{000727}{}
Montrer que pour tout $x \in \Rr$, $\left|e^x-1-x\right| \leq \frac{x^2}2 e^{\left|x\right|}$.
\finenonce{000727} 


\finexercice

\exercice{2688, matexo1, 2002/02/01}
\enonce{002688}{}
Soit $f:[a,+\infty[ \to \R$ une fonction de classe $C^1$, telle
que

$$\lim_{t\to +\infty}f(t) = f(a).$$

Montrer qu'il existe $c\in
]a,+\infty[$ tel que $f'(c) = 0$.
\finenonce{002688}
\finexercice


\exercice{5408, rouget, 2010/07/06}
\enonce{005408}{*** Formule de \textsc{Taylor}-\textsc{Lagrange}}
Soient $a$ et $b$ deux réels tels que $a<b$ et $n$ un entier naturel. Soit $f$ une fonction élément de $C^n([a,b],\Rr)\cap D^{n+1}(]a,b[,\Rr)$. Montrer qu'il existe $c\in]a,b[$ tel que 

$$f(b)=\sum_{k=0}^{n}\frac{f^{(k)}(a)}{k!}(b-a)^k+\frac{(b-a)^{n+1}f^{(n+1)}(c)}{(n+1)!}.$$

Indication. Appliquer le théorème de \textsc{Rolle} à la fonction $g(x)=f(b)-\sum_{k=0}^{n}\frac{f^{(k)}(x)}{k!}(b-x)^k-A\frac{(b-x)^{n+1}}{(n+1)!}$ où $A$ est intelligemment choisi.
\finenonce{005408}


\finexercice


\exercice{5409, rouget, 2010/07/06}
\enonce{005409}{*** Formule des trapèzes}
Soit $f\in C^2([a,b],\Rr)\cap D^3(]a,b[,\Rr)$. Montrer qu'il existe $c\in]a,b[$ tel que

$$f(b)=f(a)+\frac{b-a}{2}(f'(a)+f'(b))-f^{(3)}(c).$$

Indication. Appliquer le théorème de \textsc{Rolle} à $g'$ puis $g$ où $g(x)=f(x)-f(a)-\frac{x-a}{2}(f'(x)+f'(a))-A(x-a)^3$ où $A$ est intelligemment choisi.

Que devient cette formule si on remplace $f$ par $F$ une primitive d'une fonction $f$ de classe $C^1$ sur $[a,b]$ et deux fois dérivable sur $]a,b[$~?~Interprétez géométriquement.
\finenonce{005409}


\finexercice


\exercice{5412, rouget, 2010/07/06}
\enonce{005412}{***I Polynômes de \textsc{Legendre}}
Pour $n$ entier naturel non nul donné, on pose $L_n=((X^2-1)^n)^{(n)}$.
\begin{enumerate}
\item  Déterminer le degré et le coefficient dominant de $L_n$.
\item  En étudiant le polynôme $A_n=(X^2-1)^n$, montrer que $L_n$ admet $n$ racines réelles simples et toutes dans $]-1;1[$.
\end{enumerate}
\finenonce{005412}


\finexercice


\exercice{5415, rouget, 2010/07/06}
\enonce{005415}{**}
Montrer que pour tout réel strictement positif $x$, on a~:$\left(1+\frac{1}{x}\right)^x<e<\left(1+\frac{1}{x}\right)^{x+1}$.
\finenonce{005415}


\finexercice


\exercice{5416, rouget, 2010/07/06}
\enonce{005416}{**}
Soit $f$ une fonction dérivable sur $\Rr$ à valeurs dans $\Rr$ vérifiant $f(0)=f(a)=f'(0)=0$ pour un certain $a$ non nul. Montrer qu'il existe un point distinct de $O$ de la courbe représentative de $f$ en lequel la tangente passe par l'origine.
\finenonce{005416}


\finexercice


\exercice{5421, rouget, 2010/07/06}
\enonce{005421}{** Généralisation du théorème des accroissements finis}
Soient $f$ et $g$ deux fonctions continues sur $[a,b]$ et dérivables sur $]a,b[$.

Soit $\begin{array}[t]{cccc}
\Delta~:&[a,b]&\rightarrow&\Rr\\
 &x&\mapsto&\left|
 \begin{array}{ccc}
 f(a)&f(b)&f(x)\\
 g(a)&g(b)&g(x)\\
 1&1&1
 \end{array}
 \right|
\end{array}$.

\begin{enumerate}
\item  Montrer que $\Delta$ est continue sur $[a,b]$, dérivable sur $]a,b[$ et calculer sa dérivée.
\item  En déduire qu'il existe $c$ dans $]a,b[$ tel que $(g(b)-g(a))f'(c)=(f(b)-f(a))g'(c)$.
\end{enumerate}
\finenonce{005421}


\finexercice

\section{124.03 Applications}


\exercice{1228, legall, 1998/09/01}

\enonce{001228}{}
Soit $  f   $ une fonction continue de $  [0,1]  $ \`a valeurs dans $
{ \Rr}  .$ Pour chaque $  n \in { \Nn}  $, on note $  g_n  $ la fonction
$  \displaystyle{  x\mapsto f(x+\frac{1 }{ n})-f(x)}  .$
\begin{enumerate}
    \item On suppose $  g_n(x)>0  $ pour tout $  \displaystyle{ x\in [0,1-\frac{1}{ n}[}  .$
Montrer que $  f(1)>f(0)  .$
    \item On suppose d\'esormais que $  f(0)=f(1)  .$ Montrer que, pour chaque $  n \in { \Nn}  ,$
 la fonction $   g_n   $ s'annule en au moins un point de
l'intervalle $  \displaystyle{ [0,1-\frac{1}{ n}]}  .$
\end{enumerate}
\finenonce{001228}



\finexercice



\exercice{1229, legall, 1998/09/01}

\enonce{001229}{}
Pour tout $  n   $ entier sup\'erieur o\`u \'egal \`a $  2  $, on consid\`ere le polyn\^ome de
degr\' e $  n   $ \`a coefficients r\' eels~:
$$P_n(X)=X^n+X^{n-1}+X^2+X-1$$
\begin{enumerate}
    \item Soit $  n \geq 2  .$ Montrer que $  P_n  $ a une unique racine r\' eelle positive que l'on nommera $  \lambda _n  .$ (On pourra \'etudier l'application $  X\mapsto P_n(X)  .$)
    \item Montrer que la suite $  (\lambda _n)_{n\geq 2}  $ est croissante puis qu'elle
converge vers une limite que l'on notera $  \ell   .$
    \item Montrer que $  \ell  $ est racine du polyn\^ome $  X^2+X-1  .$ En d\' eduire sa valeur.
\end{enumerate}

\finenonce{001229}



\finexercice



\exercice{1230, legall, 1998/09/01}

\enonce{001230}{}
Soit $  f  $ une fonction d'un intervalle $  I  $ \`a valeurs dans
$  \R  $ d\'erivable sur $  I  .$ Montrer que les propri\'et\'es suivantes sont \'equivalentes~:
\begin{enumerate}
    \item
$  f  $ est strictement croissante sur $  I  .$
    \item $  f'  $ est positive ou nulle  sur $  I  $ et $  \{ x\in I  ;   f'(x)>0\}   $ est dense
dans $  I  .$
\end{enumerate}
\finenonce{001230}



\finexercice



\exercice{1231, legall, 1998/09/01}

\enonce{001231}{}
\begin{enumerate}
    \item Soit $  f   $ une application de $  { \Rr}  $ dans $  { \Rr}  $
d\'erivable en $  0  .$ Montrer qu'il existe une
application $  \epsilon   $ de $  { \Rr}  $
dans lui-m\^eme telle que
$  \forall x\in { \Rr}  :   f(x)= f(0)+ xf'(0)+x \epsilon (x)  $ et
$  \displaystyle{\lim _{x\rightarrow 0} \epsilon (x)=0 }  .$ Donner une interprtation
g\'eom\'etrique de ce r\'esultat.
    \item En d\'eduire les limites des suites $  (u_n)_{n \geq 1}  $ et $  (v_n)_{n\geq 1}  $
d\'efinies en posant, pour tout $  n \in { \Nn}^*  $~: $  \displaystyle{u_n = (n^3+1)^{\frac{1}{
3}}-n}  $ et $  \displaystyle{v_n=(1+\frac{\alpha}{ n})^{\frac{1}{ n}}}  .$
    \item Construire un exemple de suite $  (w_n)_{n \geq 1}  $ avec, $  u_n<1   $
pour tout $  n \geq 1  $ et telle que $  \displaystyle{ \lim _{n\rightarrow \infty }w_n
=1}  .$ (On pourra s'inpirer de l'exemple de $  (v_n)_{n\geq 1}  $ ci-dessus.)
\end{enumerate}
\finenonce{001231}



\finexercice



\exercice{1232, legall, 1998/09/01}

\enonce{001232}{}
\begin{enumerate}
    \item Montrer que pour tout $  x>0  $ on a~: $  \displaystyle{\frac {1}{ x+1}<\hbox{log}(x+1)
-\hbox{log}(x)<\frac{1}{ x}}  .$
    \item En d\'eduire que pour tout entier $  n\geq 1  :$ $\displaystyle {
\hbox{log}(n+1)< 1+ \frac{1}{ 2}+\cdots +\frac{1}{ n}<1+\hbox{log}(n)}  .$
    \item Posons $  \displaystyle{u_n=1+ \frac{1}{ 2}+\cdots +\frac{1}{ n}-\hbox{log}(n)}  $
Montrer que la suite $  (u_n)_{n \in { \Nn}}  $ est d\'ecroisante et convergente.
\end{enumerate}
\finenonce{001232}



\finexercice



\exercice{1233, legall, 1998/09/01}

\enonce{001233}{}
\begin{enumerate}
    \item Soit $  f   $ une application continue d'un intervalle $  ]a,b[   $ \`a valeurs dans $  { \Rr}  ,$ d\' erivable en $  c \in ]a,b[   .$ Montrer qu'il existe une (unique) application continue
$  \epsilon   $ de $  ]a,b[   $ dans $  { \Rr}  $ telle que $  f(c)=0  $ et,
pour tout $  x \in ]a,b[   $ distinct de $  c  ,$ on ait~:
$$  f (x)=f (c) +(x-c)f'(c)+ (x-c)\epsilon (x)  $$
    \item Montrer que la suite  $  (S_n)_{n\geq 1}  $ de
terme g\'en\'eral~:
$$S_n = \displaystyle\frac{ 1}{ n}+\displaystyle{\frac{ 1}{ n+1}}+\cdots + \displaystyle{\frac{ 1}{2n}}
=\displaystyle{\sum _{\scriptstyle k=0}^{\scriptstyle n}\frac{1}{ n+k}}$$
est d\' ecroissante et qu'elle converge vers une limite que l'on nommera $  S  .$
    \item Pourquoi peut on dire, {\it a priori}, que  $
\displaystyle{\frac{1}{ 2}}\leq S \leq 1  ?$
    \item Soit $f  :  \rbrack -1, 1\lbrack \rightarrow   { \Rr}  $ une application continue, d\'erivable en $
 0  $ et telle que $  f(0)=0  .$ Montrer que la suite $  (\sigma _n(f))_{n\geq 1}  $ de
terme g\'en\'eral~:
$$\sigma _n(f) = \displaystyle f\left( \frac{ 1}{ n}\right) +\displaystyle f\left(\frac{ 1}{ n+1}\right) +\cdots +
\displaystyle f\left(\frac{ 1}{ 2n} \right)
$$
converge vers $  f'(0)S  $ (utiliser 1.).
    \item Montrer que $  \sigma _n(f)=\hbox{log }(2)  $ lorsque $  f  $ est l'application
$  x\mapsto \hbox{log
}(1+x)  $ et en  d\' eduire la valeur de $  S  .$
    \item Calculer la limite de la suite  $  (\sigma _n)_{n\geq 1}  $ de
terme g\'en\'eral~:
$$\sigma _n = \displaystyle \hbox{sin }\frac{ 1}{ n}+\displaystyle \hbox{sin }\frac{ 1}{ n+1}+\cdots +
 \displaystyle \hbox{sin }\frac{ 1}{ 2n}  .
$$
    \item Plus g\' en\' eralement, quelle est la valeur pour $  p\in { \Nn}^*  $ donn\' e, de la
limite $  S_p  $ de la suite $  (\sigma _n(p))_{n\geq 1}  $ de
terme g\'en\'eral~:
$$\sigma _n(p) =\displaystyle{\sum _{\scriptstyle k=0}^{\scriptstyle  pn}\frac{1}{ n+k}}  ?$$
\end{enumerate}

\finenonce{001233}



\finexercice



\exercice{1234, legall, 1998/09/01}

\enonce{001234}{}
Soit $  f   $ une fonction d\' erivable et $  a  $ un r\' eel. Soit $  h>0  $
un nombre r\' eel strictement positif fix\' e.
\begin{enumerate}
    \item Montrer qu'il existe $  \theta \in ]0,1[  $ tel que
$$\displaystyle{\frac {f(a+h)-2f(a)+f(a-h)}{ h}=f'(a+\theta h)-f'(a-\theta h)}  .$$
    \item Pour tout $  h\not = 0  $ on note :
$  \displaystyle{ \varphi (h)=\frac{f(a+h)-2f(a)+f(a-h)}{ h^2}}  .$
Montrer que si $  f''(a)  $ existe, alors $  \displaystyle{\lim _{h\rightarrow 0}\varphi(h)=f''(a) }.$
\end{enumerate}
\finenonce{001234}



\finexercice



\exercice{1235, legall, 2003/10/01}

\enonce{001235}{}
  Soit $I$ un intervalle ouvert contenant $0$ et $1 $ 
et $f: I \rightarrow \Rr $ une fonction d\'erivable.
On pose $p=f(1)-f(0).$

\begin{enumerate}
        \item Soit $g :[0,1]\rightarrow \Rr $ la fonction d\'efinie 
par $g(0)=f'(0)$ et $g(x)=\dfrac{f(x)-f(0)}{x}$ sinon.
Montrer que si $u$ est un r\'eel compris entre $f'(0)$ et $p$ alors 
il existe $a\in [0,1] $ tel que $u=f'(a).$

\item  Soit $h :[0,1]\rightarrow \Rr $ la fonction d\'efinie par 
$h(1)=f'(1)$ et $h(x)=\dfrac{f(x)-f(1)}{x-1}$ sinon.
Montrer que si $v$ est un r\'eel compris entre $f'(1)$ et $p$ alors 
il existe $b\in [0,1] $ tel que $v=f'(b).$
\item Soit $w$ un r\'eel compris entre $f'(0)$ et $f'(1)$. Montrer 
qu'il existe $c\in [0,1] $ tel que $w=f'(c).$
\end{enumerate}
\finenonce{001235}



\finexercice



\exercice{1236, legall, 2003/10/01}

\enonce{001236}{}
Soit $P(X)$ un polyn\^ome \`a coefficients complexes de degr\'e $3$ 
ayant trois racines distinctes. Montrer que les racines de $P'$ sont 
dans le triangle ayant pour sommet les racines
de $P$
\finenonce{001236}



\finexercice


\section{124.04 Fonctions convexes}


\exercice{3981, quercia, 2010/03/11}
\enonce{003981}{Déterminant}

Soit $f : \R \to \R$ convexe et $x < y < z$.
Montrer que $\begin{vmatrix} 1 &x &f(x) \cr 1 &y &f(y) \cr 1 &z &f(z) \cr\end{vmatrix} > 0$.
\finenonce{003981}



\finexercice


\exercice{3982, quercia, 2010/03/11}
\enonce{003982}{Somme de fractions}
Soient $x_1, x_2, \dots, x_n > 0$.
Montrer que $\frac {x_1}{x_2} + \frac{x_2}{x_3} + \dots + \frac{x_n}{x_1} \ge n$.


\finenonce{003982}



\finexercice


\exercice{3983, quercia, 2010/03/11}
\enonce{003983}{Monotonie}

Soit $f : {\R} \to  {\R}$ convexe. Montrer que l'on a :

-- soit $f$ est croissante sur $\R$.\par
-- soit $f$ est décroissante sur $\R$.\par
-- soit il existe $a \in \R$ tel que $f$ est décroissante sur
$]-\infty,a]$, puis croissante sur $[a,+\infty[$.


\finenonce{003983}



\finexercice


\exercice{3984, quercia, 2010/03/11}
\enonce{003984}{Fonction convexe bornée}

\begin{enumerate}
  \item Soit $f : {\R^+} \to {\R}$ convexe et bornée. Montrer que $f$ est décroissante.
  \item Soit $f : {\R} \to {\R}$ convexe et bornée. Montrer que $f$ est constante.

\end{enumerate}
\finenonce{003984}



\finexercice


\exercice{3985, quercia, 2010/03/11}
\enonce{003985}{$f$ convexe majorée par $g$ affine}

Soit  $f  : {\R^{+*}} \to  \R$ convexe et $g : {\R^{+*}} \to \R$ affine.
On suppose : $\begin{cases}\forall\ x > 0,\ f(x) \le g(x),\cr f(1) = g(1).\cr\end{cases}$

Montrer que $f = g$.
\finenonce{003985}



\finexercice


\exercice{3986, quercia, 2010/03/11}
\enonce{003986}{Position par rapport à une asymptote}

Soit $f : {\R} \to {\R}$ convexe telle que ${\cal C}_f$ admet une
asymptote d'équation $y = mx+p$ en $+\infty$.

Montrer que ${\cal C}_f$ est au dessus de cette asymptote.

\finenonce{003986}



\finexercice


\exercice{3987, quercia, 2010/03/11}
\enonce{003987}{Fonction convexe dérivable}

Soit $f : {\R} \to {\R}$ convexe dérivable. Montrer que $f'$ est continue.

\finenonce{003987}



\finexercice


\exercice{3988, quercia, 2010/03/11}
\enonce{003988}{\'Etude à l'infini}

Soit $f : \R \to \R$ deux fois dérivable telle que :
$f \ge 0$, $f' \ge 0$, $f'' \ge 0$.

\begin{enumerate}
  \item \'Etudier l'existence des limites (dans $\overline\R\,$)
    en $+\infty$ de $f(x)$, $f'(x)$, $\frac {f(x)}x$.

  \item Même question pour les limites en $-\infty$ de $f(x)$, $f'(x)$, et $xf'(x)$.
    
\end{enumerate}
\finenonce{003988}



\finexercice


\exercice{3989, quercia, 2010/03/11}
\enonce{003989}{Zéro de $f''$}

Soit $f : {[0,+\infty[} \to {\R}$ deux fois dérivable
telle que $f(x) \to f(0)$
lorsque $x\to{+\infty}$.
Montrer qu'il existe $c \in {]0,+\infty[}$ tel que $f''(c) = 0$.

\finenonce{003989}



\finexercice


\exercice{3990, quercia, 2010/03/11}
\enonce{003990}{$f((x+y)/2) \le (f(x) + f(y))/2$}

Soit $f : {[a,b]} \to {\R}$ continue telle que :
$\forall\ x,y \in {[a,b]},\ f\left(\frac {x+y}2\right) \le \frac {f(x) + f(y)}2$.

Montrer que $f$ est convexe.


\finenonce{003990}



\finexercice


\exercice{3991, quercia, 2010/03/11}
\enonce{003991}{Suites adjacentes}

Soit $f : {[a,b]} \to {[c,d]}$ convexe, bijective, croissante.
On définit les suites $(u_n)$ et $(v_n)$ par :

$$a \le u_0 \le v_0 \le b,\quad u_{n+1} = \frac {u_n + v_n}2,\quad
v_{n+1} = f^{-1}\left(\frac {f(u_n) + f(v_n)}2\right).$$

Montrer que $(u_n)$ et $(v_n)$ convergent vers une même limite.

\finenonce{003991}



\finexercice


\exercice{3992, quercia, 2010/03/11}
\enonce{003992}{Polygone inscrit dans un cercle de périmètre maximum}

Soit $n \ge 3$ et $A_1A_2\dots A_n$ un polygone convexe à $n$ côtés inscrit
dans un cercle fixé.

Montrer que le périmètre de ce polygone est maximal si et seulement si le polygone est régulier.
\finenonce{003992}



\finexercice


\exercice{3993, quercia, 2010/03/11}
\enonce{003993}{Fonctions logarithmiquement convexe}

Soit $f : {\R} \to {\R^{+*}}$.
Montrer que : ($\ln f$ est convexe) $\iff$ ($\forall\ \alpha > 0$, $f^\alpha$ est convexe).
\finenonce{003993}



\finexercice


\exercice{3994, quercia, 2010/03/11}
\enonce{003994}{Limite de $f(x) - xf'(x)$}

Soit $f : \R \to \R$ convexe dérivable.

\begin{enumerate}
  \item Montrer que $p = \lim_{x\to+\infty} (f(x) - xf'(x))$ existe.
    

  \item On suppose $p$ fini. En utilisant le fait que $f(x) - xf'(x)$ est
    bornée au voisinage de $+\infty$, montrer que $\frac{f(x)}x$ et $f'(x)$
    admettent une même limite $m$ finie en $+\infty$.

  \item Montrer alors que $f(x) - mx - p \to 0$ lorsque $x\to {+\infty}$.
    

\end{enumerate}
\finenonce{003994}



\finexercice


\exercice{3995, quercia, 2010/03/11}
\enonce{003995}{Fonction positive concave}

Soit $f : {[0,+\infty[} \to {[0,+\infty[}$ concave.

\begin{enumerate}
  \item Montrer que la fonction $x \longmapsto \frac {f(x)}x$
    est décroissante sur $]0,+\infty[$.
  \item  Montrer que : $\forall\ x,y \ge 0,\ f(x+y) \le f(x) + f(y)$.

\end{enumerate}
\finenonce{003995}



\finexercice


\exercice{3996, quercia, 2010/03/11}
\enonce{003996}{Constante d'Euler}

Soit $f : {[0,+\infty[} \to {\R}$ concave, dérivable, croissante.

\begin{enumerate}
  \item Montrer que : $\forall\ x\ge 1,\ f(x+1)-f(x) \le f'(x) \le f(x)-f(x-1)$.
  \item On pose : $\begin{cases} u_n = f'(1) + f'(2) + \dots + f'(n) - f(n) \cr
                       v_n = f'(1) + f'(2) + \dots + f'(n) - f(n+1).\cr\end{cases}$
    Montrer que ces suites convergent.

  \item On prend $f(x) = \ln x$. Soit $\gamma = \lim_{n\to\infty} u_n$
    (constante d'Euler). Calculer $\gamma$ à $10^{-2}$ près.

\end{enumerate}
\finenonce{003996}



\finexercice


\exercice{3997, quercia, 2010/03/11}
\enonce{003997}{Tangentes passant par un point}

Soit $f : \R \to \R$ convexe dérivable, et $A = (a,b) \in \R^2$.
\'Etudier le nombre maximal de tangentes à ${\cal C}_f$ passant par $A$.

\finenonce{003997}



\finexercice


\exercice{3998, quercia, 2010/03/11}
\enonce{003998}{Caractérisation des fonctions convexes ou concaves par le TAF}

Soit $f : \R \to \R$ dérivable telle que :
$\forall\ a,b \in \R \text{ tq } a < b,\ \exists!\ c \in ]a,b[ \text{ tq }
 f(b) - f(a) = (b-a) f'(c)$.

\begin{enumerate}
  \item
Montrer que pour tout $a \in \R$, la fonction $b \longmapsto \frac {f(b)-f(a)}{b-a}$
est monotone sur $]-\infty,a[$ et sur $]a,+\infty[$.
  \item En déduire que $f$ est strictement convexe ou strictement concave.
\end{enumerate}
\finenonce{003998}



\finexercice


\exercice{3999, quercia, 2010/03/11}
\enonce{003999}{Pseudo-dérivée seconde}

Soit $f : {\R} \to {\R}$ continue. On suppose que :
$\forall\ x \in \R,\
D^2f(x) = \lim_{h \to 0} \frac {f(x+h)+f(x-h)-2f(x)}{h^2}$
existe.

\begin{enumerate}
  \item  Si $f$ est de classe $\mathcal{C}^2$, calculer $D^2f(x)$.

  \item   Soit $f$ quelconque et $a < b < c$ tels que
      $f(a) = f(b) = f(c)$.

      Montrer qu'il existe $x \in {]a,c[}$ tq $D^2f(x) \le 0$.
      


On suppose à présent que : $\forall\ x \in \R, D^2f(x) \ge 0$.

  \item   Soient $a < b < c$ et
      $P$ le polynôme de degré inférieur ou égal à 2 coïncidant avec $f$
      aux points $a,b,c$. Montrer que $P'' \ge 0$.

  \item   Calculer $P''$ en fonction de $a,b,c$ et $f(a),f(b),f(c)$.
      En déduire que $f$ est convexe.
\end{enumerate}
\finenonce{003999}



\finexercice


\exercice{4000, quercia, 2010/03/11}
\enonce{004000}{Fonction convexe non dérivable sur un sous ensemble dénombrable}

Soit $(a_n)$ une suite bornée de réels.
On pose $f(x) = \sum_{n=0}^\infty \frac {|x-a_n|}{3^n}$.


Montrer que $f$ est convexe, et n'est pas dérivable aux points $a_n$.

\finenonce{004000}



\finexercice


\exercice{4001, quercia, 2010/03/11}
\enonce{004001}{Convergence simple $+$ convexité $=>$ convergence uniforme sur un compact}

Soit $(f_n)$ une suite de fonctions convexes sur $[a,b]$ convergeant
simplement vers une fonction $f$ supposée continue.

Soit $\varepsilon > 0$.

\begin{enumerate}
  \item  Montrer qu'il existe $p \in \N^*$ tel que :
      $\forall\ x,y \in {[a,b]},\
      |x-y| \le \frac {b-a}p  \Rightarrow  |f(x) - f(y)| \le \varepsilon$.

On choisit un tel $p$, et on fixe une subdivision $(a_k)$ de $[a,b]$ telle
que $a_k = a + k\frac {b-a}p$.

  \item   Soit $t \in {[0,1]}$. Encadrer $f_n\bigl(ta_k + (1-t)a_{k+1}\bigr)$ par deux
      fonctions affines de t en utilisant la convexité de~$f_n$.

  \item   Montrer que la suite $(f_n)$ converge uniformément vers $f$.

\end{enumerate}
\finenonce{004001}



\finexercice


\exercice{4002, quercia, 2010/03/11}
\enonce{004002}{DL d'une fonction convexe}

Soit $f:\R \to \R$ convexe dérivable telle que
$f(x) = a + bx + \frac {cx^2}2 + o_{x\to0} (x^2)$.

Montrer que $f$ est deux fois dérivable en $0$ et $f''(0) = c$

(encadrer $f'(x)$ par les taux d'accroissements de $f$ entre $x - \varepsilon x$,
$x$ et $x + \varepsilon x$).
\finenonce{004002}



\finexercice


\exercice{4003, quercia, 2010/03/11}
\enonce{004003}{DL d'une fonction convexe}
Soit $f$ continue et croissante sur $\R^+$. On pose
$F(x)= \int_0^x f$, et l'on suppose que $F(x)=x^2+ o(x)$. Montrer que
$f(x)=2x+ o(\sqrt x)$.

\finenonce{004003}



\finexercice

\section{124.99 Autre}


\exercice{728, bodin, 1998/09/01}
\enonce{000728}{}
Soit $f: \Rr \longrightarrow \Rr$ d\'efinie par
$f(x) = (1-k)^3x^2+(1+k)x^3$
o\`u $k$ est un nombre r\'eel. D\'eterminer
les valeurs de $k$ pour lesquelles
l'origine est un extremum local de $f$.
\finenonce{000728} 


\finexercice

\exercice{729, bodin, 1998/09/01}

\enonce{000729}{}
Appliquer la r\`egle de l'Hôpital aux
calculs des limites suivantes :
$$
\lim_{x\rightarrow 0}\ \ \left( \frac{1}{\sin^2x}-\frac{1}{x^2}\right),
$$
$$
\lim_{x\rightarrow 0}\ \  (1-\cos x)\text{cotan}\, x.
$$
\finenonce{000729}



\finexercice



\exercice{730, bodin, 1998/09/01}

\enonce{000730}{}
Calculer
$$
\lim_{x\rightarrow 0}\ \  \frac{\cos (x^4)-1}{x^4e^x} ;
$$

$$
\lim_{x\rightarrow 0}\ \  \frac{\ln{\cos ax}}{\ln {\cos bx}} ;
$$

$$
\lim_{x\rightarrow 0}\ \  x^2\left( \exp \frac{1}{x} - \exp\frac{1}{x+1}  \right).
$$
\finenonce{000730}



\finexercice



\exercice{731, ridde, 1999/11/01}

\enonce{000731}{}
Soit $f \in C^2 (\Rr)$ telle que $\forall (x, y) \in \Rr^2 \, \, f (x + y)f (x-y)
\leq f (x)^2$.
Montrer que $\forall x \in \Rr \, \, f (x)f'' (x) \leq f' (x)^2$.
\finenonce{000731}



\finexercice



\exercice{732, ridde, 1999/11/01}

\enonce{000732}{}
Soit $f : \Rr^ + \rightarrow \Rr$ d\'erivable telle que $\lim\limits_{ + \infty}
f' = l$. Montrer qu'alors $\lim\limits_{ + \infty} \dfrac{f (x)}x  = l$.
\finenonce{000732}



\finexercice



\exercice{733, ridde, 1999/11/01}
\video{_YeLn2_0Tso}
\enonce{000733}{}
D\'eterminer les extremums de $f (x) = x^4-x^3 + 1$ sur $\Rr$.
\finenonce{000733} 


\finexercice

\exercice{734, gourio, 2001/09/01}
\enonce{000734}{}
Quel est le lieu des points d'inflexion (puis des extr\'{e}mums locaux) de $
f_{\lambda } $ quand $\lambda $ d\'{e}crit ${\Rr}, $ o\`{u} :
$$f_{\lambda }:x \mapsto \lambda e^{x}+x^{2}. $$
\finenonce{000734} 


\finexercice

\exercice{735, gourio, 2001/09/01}

\enonce{000735}{}
Trouver les fonctions $f:{\Rr\rightarrow \Rr}$ d\'{e}rivables en $0 $  telles que :
$$\exists \lambda \in {\Rr}^{+}-\{1\},\forall x\in {\Rr},f(\lambda
x)=\lambda f(x).
$$
\finenonce{000735}



\finexercice



\exercice{736, gourio, 2001/09/01}

\enonce{000736}{}
Soit $f$ d\'{e}rivable sur ${\Rr } $ telle que $f(\omega )=\omega . $ On
d\'{e}finit une suite $(x_{n})_{n\in \Nn}$ par la donn\'{e}e de $x_{0}$
et la r\'{e}currence $x_{n+1}=f(x_{n}).$
Montrer que si $\left| f^{\prime }(\omega )\right| <1,\exists \epsilon
>0,\forall x_{0}\in ]\omega -\epsilon ,\omega +\epsilon
[,(x_{n})_{n\in \Nn}$ converge vers $w,$ et que si $\left| f^{\prime
}(\omega )\right| >1 $ la suite $(x_{n})_{n\in \Nn}$ converge vers $w$
si et seulement si elle est stationnaire (i.e.\ $x_{n}=\omega $ \`{a} partir
d'un certain rang). Que dire dans le cas $\left| f^{\prime }(\omega )\right|
=1$ ?
\finenonce{000736}



\finexercice



\exercice{737, gourio, 2001/09/01}

\enonce{000737}{}
Soit $f\in C^{1}([0;1],{\Rr}), $telle que $f(0)=0. $ Calculer :
$$\lim\limits_{n\rightarrow \infty }\sum\limits_{k=1}^{n}f(\frac{k}{n^{2}}). $$
\finenonce{000737}



\finexercice



\exercice{738, bodin, 2001/11/01}
\video{MqjLEkOQD3w}
\enonce{000738}{}
 Soient $f,g :[a , b] \longrightarrow \R$ deux
fonctions continues
 sur $[a, b]$ ($a < b$) et d\'erivables sur $]a , b[.$ On suppose que
 $g' (x) \neq 0$ pour tout $x \in ]a , b[.$

\begin{enumerate}
\item
 Montrer que $g (x) \neq g (a)$ pour tout $x \in ]a , b[.$

\item Posons $p = \frac{f (b) - f (a)}{g (b) - g (a)}$ et
consid\'erons  la fonction $h (x) = f
 (x) - p g (x)$ pour $x \in [a , b].$
 Montrer que $h$ v\'erifie les hypoth\`eses du th\'eor\`eme de Rolle
 et en d\'eduire qu'il existe un nombre r\'eel $c \in ]a , b[$ tel que
 $$ \frac{f (a) - f (b)}{g (a) - g (b)} = \frac{f' (c)}{g' (c)}.$$
\item On suppose que $\lim_{x \to b^{-}} \frac{f' (x)}{g' (x)} =
\ell,$ o\`u $\ell $ est un nombre r\'eel.
 Montrer que  $$ \lim_{x \to b^{-}} \frac{f (x) - f (b)}{g (x) - g (b)} = \ell.$$
\item  \emph{Application.} Calculer la limite suivante:
 $$ \lim_{x \to 1^{-}} \frac{\Arccos x}{\sqrt{1- x^{2}}}.$$
\end{enumerate}
\finenonce{000738} 


\finexercice

\exercice{739, bodin, 2001/11/01}
\video{OFIk3Tj3QIA}
\enonce{000739}{}
 Soit $n \geq 2$ un entier fix\'e et $f:\R^{+} = [0,
+ \infty[ \longrightarrow \R$ la fonction d\'efinie par la formule
suivante:
$$ f (x) = \frac{1 + x^n}{(1 + x)^n}, \ \ x \geq 0.$$
\begin{enumerate}
  \item
  \begin{enumerate}
     \item Montrer que $f$ est d\'erivable sur $\R^{+}$ et calculer $f' (x)$
pour $x \geq 0.$
     \item En \'etudiant le signe de $f' (x)$ sur $\R^{+},$ montrer que $f$
atteint un minimum sur $\R^{+}$ que l'on d\'eterminera.
  \end{enumerate}
  \item
  \begin{enumerate}
     \item En d\'eduire l'in\'egalit\'e suivante:
$$ (1+ x)^n \leq 2^{n - 1} (1+ x^n), \ \ \forall x \in \R^{+}.$$
     \item Montrer que si  $x \in \R^{+}$ et $y \in \R^{+}$ alors on a
$$ (x + y)^n \leq 2^{n - 1} (x^n + y^n).$$
   \end{enumerate}
 \end{enumerate}
\finenonce{000739} 


\finexercice

\exercice{740, monthub, 2001/11/01}
\video{EVKKvovB3ps}
\enonce{000740}{}
On consid\`ere la fonction $f : \R \to \R$ d\'efinie par
\begin{equation*}
f(t) =
\begin{cases}
  e^{1/t} & \mathrm{\ si\ } t<0\\
0  & \mathrm{\ si\ } t \geq 0
\end{cases}
\end{equation*}

\begin{enumerate}
\item D\'emontrer  que $f$ est d\'erivable sur $\R$, en particulier en $t=0$.
\item Etudier l'existence de $f''(0)$.
\item On veut montrer que pour $t<0$, la d\'eriv\'ee $n$-i\`eme de $f$ s'\'ecrit
$$f^{(n)}(t)=\frac{P_n(t)}{t^{2n}}e^{1/t}$$
o\`u $P_n$ est un polyn\^ome.
\begin{enumerate}
\item Trouver $P_1$ et $P_2$.
\item Trouver une relation de r\'ecurrence entre $P_{n+1}, P_n$ et $P'_n$ pour
  $n\in \N^*$.
\end{enumerate}
\item Montrer que $f$ est de classe $C^{\infty}$.
\end{enumerate}
\finenonce{000740}


\finexercice

\exercice{3932, quercia, 2010/03/11}
\enonce{003932}{Limite double}

Soit $f:\R \to \R$ continue en $0$.
Montrer que $f$ est dérivable en $0$, et $f'(0) = \ell$ si et seulement si :
$$\forall\ \varepsilon > 0,\ \exists\ \delta > 0 \text{ tq }
  \forall\ h,k\in {]0,\delta[},\
  \left|\frac {f(h) - f(-k)}{h+k} - \ell\right| \le \varepsilon.$$
\finenonce{003932}



\finexercice


\exercice{3933, quercia, 2010/03/11}
\enonce{003933}{Propriétés de parité et de périodicité}

Soit $f:\R \to \R$ dérivable.

\begin{enumerate}
  \item Que peut-on dire de $f'$ si on sait que $f$ est paire ? impaire ? périodique ?
  \item Que peut-on dire de $f$ si on sait que $f'$ est paire ? impaire ? périodique ?
  \item Montrer que si $f'$ est $T$-périodique et $f(T) \ne f(0)$, alors $f$ n'a pas de
    période (on étudiera $f(nT)$ pour $n\in\N$).
\end{enumerate}
\finenonce{003933}



\finexercice


\exercice{3934, quercia, 2010/03/11}
\enonce{003934}{Propriété de parité}
Soit $f:\R \to \R$ de classe $\mathcal{C}^1$ telle que la fonction
$t  \mapsto 2f(t) - tf'(t)$ est paire. $f$ est-elle paire~?


\finenonce{003934}



\finexercice


\exercice{3935, quercia, 2010/03/11}
\enonce{003935}{Injectivité locale}

Soit $f:\R \to \R$ dérivable et $a \in \R$ tel que $f'(a) \ne 0$.

\begin{enumerate}
  \item Montrer qu'il existe un voisinage $V$ de $a$ tel que
    $\forall\ x \in V\setminus\{a\},\ f(x) \ne f(a)$.
  \item Si $f'$ est continue au point $a$, montrer qu'il existe un voisinage $V$ de $a$
    tel que $f_{|V}$ soit injective.
\end{enumerate}
\finenonce{003935}



\finexercice


\exercice{3936, quercia, 2010/03/11}
\enonce{003936}{Dérivabilité de $|f|$}
Soit $f:\R \to \R$ dérivable. Montrer que $|f|$ admet en tout point une dérivée
à droite et une dérivée à gauche.
\finenonce{003936}



\finexercice


\exercice{3937, quercia, 2010/03/11}
\enonce{003937}{$f'(x) \to \ell$ et $f$ est bornée}

\label{limfprime}
Soit $f:\R \to \R$ dérivable et bornée telle que
$f'(x) \to \ell$ lorsque $x \to +\infty$. Montrer que $\ell = 0$.
\finenonce{003937}



\finexercice


\exercice{3938, quercia, 2010/03/11}
\enonce{003938}{$\lim_\infty f'(x) = \lim_\infty f(x)/x$ }
Soit $f:\R \to \R$ dérivable telle que $f'(x) \to \ell$ lorsque $x\to +\infty$.
Montrer que $\frac {f(x)}x \to \ell$ lorsque $x\to +\infty$.

Chercher un contrexemple pour la réciproque.

\finenonce{003938}



\finexercice


\exercice{3939, quercia, 2010/03/11}
\enonce{003939}{Centrale MP 2006}
Soit $f : {\R^+} \to \R$ continue telle que
$f(x) \int_{t=0}^x f^2(t)\,d t \to \ell\in\R^*$ lorsque $x\to+\infty$.
Montrer qu'il existe $\alpha,\beta\in\R^*$ tels que 
$f(x)\sim \frac\alpha {x^\beta}$ en $+\infty$.


\finenonce{003939}



\finexercice


\exercice{3940, quercia, 2010/03/11}
\enonce{003940}{Propriété des valeurs intermédiaires pour $f'$}
Soit $f : {[a,b]} \to \R$ dérivable.

\begin{enumerate}
  \item   On suppose que : $\forall\ x \in {[a,b]},\ f'(x) \ne 0$. Montrer que $f'$ est
      de signe constant.
  \item   Dans le cas général, montrer que $f'([a,b])$ est un intervalle.

\end{enumerate}
\finenonce{003940}



\finexercice


\exercice{3941, quercia, 2010/03/11}
\enonce{003941}{Propriété des valeurs intermédiaires pour $f'$}
Soit $f : {[a,b]} \to \R$ dérivable.

\begin{enumerate}
  \item On note $E = \{(x,y) \in {[a,b]}^2 \text{ tq } x < y\}$ et pour $(x,y) \in E$ :
    $\varphi(x,y) = \frac{f(x)-f(y)}{x-y}$.
    Montrer que $\varphi(E)$ est un intervalle.
  \item En déduire que $f'([a,b])$ est un intervalle.
\end{enumerate}
\finenonce{003941}



\finexercice


\exercice{3942, quercia, 2010/03/11}
\enonce{003942}{Règle de l'Hospital}

Soient ${f,g} : {[a,b]} \to \R$ dérivables avec :
$\forall\ x \in {]a,b[},\ g'(x) \ne 0$.

\begin{enumerate}
  \item   Montrer qu'il existe $c \in {]a,b[}$ tel que :
      $\frac {f(b) - f(a)}{g(b)-g(a)} = \frac {f'(c)}{g'(c)}$.

      (Appliquer le théorème de Rolle à $f-\lambda g$, où $\lambda$ est un réel
      bien choisi)
  \item   En déduire que si $\frac {f'(x)}{g'(x)} \to \ell$ lorsque $x\to {a^+}$, alors
      $\frac {f(x) - f(a)}{g(x)-g(a)} \to \ell$ lorsque $x\to {a^+}$
      (règle de l'Hospital).

  \item   Application :
      déterminer $\lim_{x\to 0^+} \frac {\cos x - e^x}{(x+1)e^x - 1}$.
\end{enumerate}
\finenonce{003942}



\finexercice


\exercice{3943, quercia, 2010/03/11}
\enonce{003943}{Recherche de limite}

Trouver $\lim_{x\to\pi/4} \frac {\ln(\sin x) - \ln(\cos x)}{\sin x - \cos x}$.

\finenonce{003943}



\finexercice


\exercice{3944, quercia, 2010/03/11}
\enonce{003944}{$f(a)=f(b)=0$, $f'(a)f'(b) > 0  \Rightarrow $ il existe un autre zéro}


Soit $f : {[a,b]} \to {\R}$ dérivable telle que
$f(a) = f(b) = 0$, et $f'(a) > 0$, $f'(b) > 0$.

Montrer qu'il existe $c \in {]a,b[}$ tel que $f(c) = 0$, et $f'(c) \le 0$.
\finenonce{003944}



\finexercice


\exercice{3945, quercia, 2010/03/11}
\enonce{003945}{$f'(a) = f'(b)$}
Soit $f : {[a,b]} \to \R$ dérivable telle que $f'(a) = f'(b)$.

Montrer qu'il existe $c \in {]a,b[}$ tel que $f'(c) = \frac {f(c) - f(a)}{c-a}$.

\finenonce{003945}



\finexercice


\exercice{3946, quercia, 2010/03/11}
\enonce{003946}{Tangentes passant par un point donné}

Soit $f : {[a,b]} \to \R$ dérivable, telle que $f(a) = f(b) = 0$.
Montrer que pour tout $d \in \R\setminus[a,b]$, il existe une tangente à ${\cal C}_f$
passant par le point $(d,0)$.
\finenonce{003946}



\finexercice


\exercice{3947, quercia, 2010/03/11}
\enonce{003947}{Rolle itéré}

Soit $f : {[a,b]} \to \R$ $n$ fois dérivable.

\begin{enumerate}
  \item   Si $f$ s'annule en $n+1$ points distincts dans $[a,b]$,
      montrer qu'il existe $c \in {]a,b[}$ tel que $f^{(n)}(c) = 0$.
  \item   Si $f(a) = f'(a) = \dots = f^{(n-1)}(a) = f(b) = 0$,
      montrer qu'il existe $c \in {]a,b[}$ tel que $f^{(n)}(c) = 0$.
\end{enumerate}
\finenonce{003947}



\finexercice


\exercice{3948, quercia, 2010/03/11}
\enonce{003948}{Rolle à l'infini}

Soit $f : {[a,+\infty[} \to \R$ dérivable telle que $f(x) \to f(a)$ lorsque $x\to {+\infty}$.
Montrer qu'il existe $x \in {]a,+\infty[}$ tel que $f'(x) = 0$.

\finenonce{003948}



\finexercice


\exercice{3949, quercia, 2010/03/11}
\enonce{003949}{Formule des accroissements finis avec $\theta = 1/2$}

Soit $f:\R \to \R$ dérivable telle que :
$\forall\ a,b \in \R,\ f(b) - f(a) = (b-a)f'\left(\frac {a+b}2\right)$.
Montrer que $f$ est polynomiale de degré inférieur ou égal à 2.


\finenonce{003949}



\finexercice


\exercice{3950, quercia, 2010/03/11}
\enonce{003950}{Fonction $\mathcal{C}^\infty$ bornée}
Soit $f:\R \to \R$ de classe $\mathcal{C}^\infty$ bornée.

\begin{enumerate}
  \item Montrer que si une dérivée, $f^{(k)}$, $k\ge 2$, admet un nombre fini de zéros,
    alors les dérivées précédentes, $f^{(p)}$, $1\le p < k$, tendent vers $0$
    en $\pm\infty$.

  \item En déduire que pour tout $k \ge 2$, $f^{(k)}$ s'annule au
    moins $k-1$ fois sur $\R$.
\end{enumerate}
\finenonce{003950}



\finexercice


\exercice{3951, quercia, 2010/03/11}
\enonce{003951}{Distance à la corde}

Soit $f : {[a,b]} \to \R$ de classe $\mathcal{C}^2$.

\begin{enumerate}
  \item   On suppose que $f(a) = f(b) = 0$. Soit $c \in {]a,b[}$.
      Montrer qu'il existe $d \in {]a,b[}$ tel que :
      $$f(c) = -\frac {(c-a)(b-c)}2 f''(d).$$
      (Considérer
      $g(t) = f(t) + \lambda (t-a)(b-t)$ où $\lambda$ est choisi de
      sorte que $g(c) = 0$)

  \item   Cas général : Soit $c \in {]a,b[}$.
      Montrer qu'il existe $d \in {]a,b[}$ tel que :
      $$f(c) = \frac {b-c}{b-a} f(a) + \frac {c-a}{b-a} f(b) -\frac {(c-a)(b-c)}2 f''(d).$$

\end{enumerate}
\finenonce{003951}



\finexercice


\exercice{3952, quercia, 2010/03/11}
\enonce{003952}{\'Ecart à un polynôme interpolateur}
Soit $f:\R \to \R$ de classe $\mathcal{C}^n$, $a_1, \dots, a_n$ $n$ points distincts dans $\R$,
et $P$ le polynôme de Lagrange prenant les mêmes valeurs que $f$ aux points $a_i$.
On pose $Q(x) = \frac 1{n!} \prod_{i=1}^n (x-a_i)$.


Montrer que : $\forall\ b \in \R,\ \exists\ c \in \R \text{ tq } f(b) = P(b) + Q(b)f^{(n)}(c)$

(considérer $g(t) = f(t) - P(t) - \lambda Q(t)$ où $\lambda$ est
choisi de sorte que $g(b) = 0$).


\finenonce{003952}



\finexercice


\exercice{3953, quercia, 2010/03/11}
\enonce{003953}{Polynômes de Legendre}

On pose $f(t) = (t^2-1)^n$.

\begin{enumerate}
  \item Montrer que : $\forall\ k \in \{0,\dots,n-1\},\ f^{(k)}(1) = f^{(k)}(-1) = 0$.
  \item Calculer $f^{(n)}(1)$ et $f^{(n)}(-1)$.
    
  \item Montrer que $f^{(n)}$ s'annule au moins $n$ fois dans l'intervalle $]-1,1[$.

\end{enumerate}
\finenonce{003953}



\finexercice


\exercice{3954, quercia, 2010/03/11}
\enonce{003954}{Racines de $x^n + ax + b$}

Soit $n \in \N,\ n \ge 2$, et $a,b \in \R$.
Montrer que l'équation $x^n + ax + b = 0$ ne peut avoir plus de deux racines réelles
distinctes si $n$ est pair, et plus de trois racines réelles distinctes si
$n$ est impair.

\finenonce{003954}



\finexercice


\exercice{3955, quercia, 2010/03/11}
\enonce{003955}{Racines de $P(x) - e^x$}
Soit $P$ un polynôme. Montrer qu'il existe au plus un nombre fini de réels
$x$ tels que $P(x) = e^x$.


\finenonce{003955}



\finexercice


\exercice{3956, quercia, 2010/03/11}
\enonce{003956}{Limite de $1/(n+1) + \dots + 1/2n$}

On veut calculer $\ell = \lim_{n \to \infty} \frac 1{n+1} + \dots + \frac 1{2n}$.

\begin{enumerate}
  \item   Montrer l'existence de $\ell$.
  \item   Soit $f : {[0,+\infty[} \to \R$ dérivable telle que $f(0) = 0$.
      Montrer que $f\left(\frac 1{n+1}\right) + \dots + f\left(\frac 1{2n}\right) \to \ell f'(0)$ lorsque $n\to \infty$.
  \item   On prend $f(x) = \ln(1+x)$. Déterminer $\ell$.
\end{enumerate}
\finenonce{003956}



\finexercice


\exercice{3957, quercia, 2010/03/11}
\enonce{003957}{Calcul de limite}
Soit $f:\R \to \R$ dérivable telle que $f(0) = 0$.
Chercher $\lim_{n\to\infty}\sum_{k=1}^n f\left(\frac k{n^2}\right)$.


\finenonce{003957}



\finexercice


\exercice{3958, quercia, 2010/03/11}
\enonce{003958}{Somme $1/k\ln k$}

Pour $k\in\N$, $k\ge 2$, appliquer le théorème des accroissements finis à
$x \mapsto\ln(\ln x)$ sur $[k,k+1]$.
En déduire que la série de terme général $\frac1{k\ln k}$ est divergente.
\finenonce{003958}



\finexercice


\exercice{3959, quercia, 2010/03/11}
\enonce{003959}{$f'(x)f'(f(x)) = 1$}


Trouver toutes les applications $f:\R \to \R$ dérivables telles que :
$\begin{cases}\forall\ x \in \R,\ f'(x)f'(f(x)) = 1 \cr
        f(0) = 0 \text{ et } f'(0) > 0.\cr\end{cases}$


\finenonce{003959}



\finexercice


\exercice{3960, quercia, 2010/03/11}
\enonce{003960}{$f\circ f = f$}
Soit $f : {[0,1]} \to  {[0,1]}$ dérivable telle que $f\circ f = f$.
Montrer que $f$ est constante ou bien $f = \text{id}_{[0,1]}$.
\finenonce{003960}



\finexercice


\exercice{3961, quercia, 2010/03/11}
\enonce{003961}{Dérivabilité uniforme}
Soit $f : {[a,b]} \to \R$ de classe $\mathcal{C}^1$. Démontrer que :
    $$\forall\ \varepsilon>0, \exists\ \delta>0 \text{ tq }
      \forall\ x,y \in [a,b], \text{ si } 0 < |x-y| < \delta, \text{ alors }
      \left|\frac {f(x)-f(y)}{x-y} - f'(x)\right| \le \varepsilon.$$
\finenonce{003961}



\finexercice


\exercice{3962, quercia, 2010/03/11}
\enonce{003962}{Formes indéterminées}
Soient ${u,v} : {\R} \to {\R}$ deux fonctions telles que :
$\begin{cases}\forall\ x \in \R,\ u(x) \ne v(x) \cr
        \lim_{x\to0}u(x) = \lim_{x\to0}v(x) = a > 0.\cr\end{cases}$

\begin{enumerate}
  \item Chercher $\lim_{x\to0} \frac {u^v - v^v}{u - v}$.
  \item Chercher $\lim_{x\to0} \frac {u^v - v^u}{u^u - v^v}$.
\end{enumerate}
\finenonce{003962}



\finexercice


\exercice{3963, quercia, 2010/03/11}
\enonce{003963}{$(1+k)(1+k^2)\dots(1+k^n)$}

\begin{enumerate}
  \item Montrer que : $\forall\ x \ge -1,\ \ln(1+x) \le x$.
  \item Soit $k \in {]-1,1[}$. On pose $u_n = (1+k)(1+k^2)\dots(1+k^n)$.
    Montrer que la suite $(u_n)$ est convergente
    (traiter séparément les cas $k \ge 0$, $k < 0$).
\end{enumerate}
\finenonce{003963}



\finexercice


\exercice{3964, quercia, 2010/03/11}
\enonce{003964}{Dérivée $n$-ème de $\cos^3x$}

Calculer la dérivée $n^{\text{ème}}$ de la fonction $x  \mapsto \cos^3 x$.

\finenonce{003964}



\finexercice


\exercice{3965, quercia, 2010/03/11}
\enonce{003965}{Dérivée $n$-ème de $\arctan x$ et $e^{x^3}$}

\'Etablir une formule de récurrence pour les dérivées successives des fonctions :

$f : x \mapsto \arctan x$ et $g : x \mapsto e^{x^3}$.
\finenonce{003965}



\finexercice


\exercice{3966, quercia, 2010/03/11}
\enonce{003966}{Dérivée $n$-ème de $(x^3+2x^2-5)e^{-x}$}

Calculer la dérivée $n^{\text{ème}}$ de $x \mapsto (x^3+2x^2-5)e^{-x}$.

\finenonce{003966}



\finexercice


\exercice{3967, quercia, 2010/03/11}
\enonce{003967}{Ensi Chimie P 94}

    Soit $f(x) = e^{x\sqrt3}\sin x$. Calculer $f^{(n)}(x)$ pour $n\in\N$.
    
\finenonce{003967}



\finexercice


\exercice{3968, quercia, 2010/03/11}
\enonce{003968}{Dérivée $n$-ème de $x^n(1-x)^n$}

Calculer la dérivée $n^{\text{ème}}$ de $x \mapsto x^n(1-x)^n$.
En déduire la valeur de $\sum_{k=0}^n \bigl(C_n^k\bigr)^2$.


\finenonce{003968}



\finexercice


\exercice{3969, quercia, 2010/03/11}
\enonce{003969}{Dérivées $n$-èmes de $t^{n-1}\ln(t)$ et $t^{n-1}e^{1/t}$}

Calculer $\frac {d^n}{d t^n}\left(t^{n-1}\ln t\right)$, et
$\frac {d^n}{d t^n}\left(t^{n-1}\exp(1/t)\right)$
(essayer $n=1,2,3$).
\finenonce{003969}




\finexercice


\exercice{3970, quercia, 2010/03/11}
\enonce{003970}{Dérivée $n$-ème de $f(x^2)$}
Soit $f : \R \to \R$ une fonction de classe $\mathcal{C}^n$.
On pose $g(x) = f(x^2)$.

\begin{enumerate}
  \item Montrer qu'il existe des entiers $a_{n,k}$ tels que :
    $\forall\ x,\ g^{(n)}(x) = \sum_{k= [(n+1)/2]}^n a_{n,k}f^{(k)}(x^2)(2x)^{2k-n}$.
    
  \item Calculer $a_{n,k}$ en fonction de $n$ et $k$.
    
\end{enumerate}
\finenonce{003970}



\finexercice


\exercice{3971, quercia, 2010/03/11}
\enonce{003971}{Dérivée $n$-ème de $f(1/x)$}

Soit $f$ une fonction $n$ fois dérivable sur un intervalle $I$ ne contenant
pas $0$, et $g(x) = f\left(\frac 1x \right)$.

\'Etablir : $g^{(n)}(x) = (-1)^n\sum_{p=0}^{n-1}
           \frac {(n-1)(n-2)\dots(n-p)}{x^{2n-p}} C_n^p f^{(n-p)}\left(\frac 1x \right)$.

\finenonce{003971}



\finexercice


\exercice{3972, quercia, 2010/03/11}
\enonce{003972}{Dérivées de $e^{-1/x^2}$}

Soit $f : \R \to \R$
$$ \begin{cases}x \ne 0 &f(x) = \exp\left(-\frac 1{x^2}\right)\cr
                            &f(0)=0.\cr\end{cases}$$
Montrer que $f$ est de classe $\mathcal{C}^\infty$ en $0$ et : $\forall\ k \in \N,\ f^{(k)}(0) = 0$.


\finenonce{003972}



\finexercice


\exercice{3973, quercia, 2010/03/11}
\enonce{003973}{$(f(2t)-f(t))/t$}

Soit $f:\R \to \R$ continue telle que
$\frac {f(2t)-f(t)}t \to a$ (si $t\to 0$. Montrer que $f$ est dérivable en $0$ et $f'(0) = a$.

\finenonce{003973}



\finexercice


\exercice{3974, quercia, 2010/03/11}
\enonce{003974}{$\sin x -3x/\pi + 4x^3/\pi^3 \ge 0$}

On pose $f(x) = \sin x - \frac {3x}\pi + \frac {4x^3}{\pi^3}$.
Montrer que : $\forall\ x \ge 0,\ f(x) \ge 0$
(chercher le signe de $f^{(4)}$).

\finenonce{003974}



\finexercice


\exercice{3975, quercia, 2010/03/11}
\enonce{003975}{Courbes homothétiques}

Soit $a > 0,\ a \ne 1$. On note $\cal C$ la courbe d'équation : $y = \ln x$,
et ${\cal C}'$ celle d'équation : $y = a\ln x$.

\begin{enumerate}
  \item Montrer que $\cal C$ et ${\cal C}'$ ont une et une seule tangente commune.

  \item Montrer que $\cal C$ et ${\cal C}'$ sont homothétiques.
    

\end{enumerate}
\finenonce{003975}



\finexercice


\exercice{3976, quercia, 2010/03/11}
\enonce{003976}{Matexo}

Soit $f$ une application dérivable de $\R$ dans $\R$ telle que 
$\forall\ x\in\R,\quad f(x)f'(x)\ge 0$.
Montrer que $f^{-1}(\R^*)$ est un intervalle.

\finenonce{003976}



\finexercice


\exercice{3977, quercia, 2010/03/11}
\enonce{003977}{Mines MP 2000}

Montrer que pour tout $x$ réel, il existe $a(x)$ unique tel que
$ \int_{t=x}^{a(x)} e^{t^2}\, d t=1$. Montrez que $a$ est indéfiniment
dérivable, et que son graphe est symétrique par rapport à la deuxième
bissectrice. 
\finenonce{003977}



\finexercice


\exercice{3978, quercia, 2010/03/11}
\enonce{003978}{$\varphi(2x)=2\varphi(x)$ (Centrale MP 2003)}
Trouver toutes les fonctions $\varphi : \R \to \R$ dérivables telles que $\forall\ x\in\R,\ \varphi(2x)=2\varphi(x)$.

\finenonce{003978}



\finexercice


\exercice{3979, quercia, 2010/03/11}
\enonce{003979}{$f' = f^{-1}$ (Ens Cachan MP$^*$ 2003)}

On note~$E$ l'ensemble des fonctions~$f$ de classe~$\mathcal{C}^1$ bijectives
de~$]0,+\infty[$ sur~$]0,+\infty[$ telles que $f' = f^{-1}$.

\begin{enumerate}
  \item Trouver un élément de~$E$ du type $x \mapsto c x^m$, où $c$ et $m$ sont réels.
    
  \item Quelle est la limite en~$0$ de~$f$~?
  \item Montrer que~$f$ est un $\mathcal{C}^\infty$ difféomorphisme de~$]0,+\infty[$ sur~$]0,+\infty[$.
  \item Montrer que~$f$ admet un unique point fixe.
    
  \item Soit~$g$ un deuxième élément de~$E$. Montrer que $g$ admet le même
    point fixe que~$f$.
    
    
\end{enumerate}
\finenonce{003979}



\finexercice


\exercice{3980, quercia, 2010/03/11}
\enonce{003980}{$f^{2} + (1+f')^{2} \le 1$, Polytechnique MP$^*$ 2006}
Soit $f:\R \to \R$ dérivable telle que $f^{2} + (1+f')^{2} \le 1$.
Montrer que $f$ est nulle.


\finenonce{003980}



\finexercice


\exercice{5407, rouget, 2010/07/06}
\enonce{005407}{***}
Soit $f\in C^1([a,b],\Rr)$ telle que $\frac{f(b)-f(a)}{b-a}=\mbox{sup}\{f'(x),\;x\in[a,b]\}$. Montrer que $f$ est affine.
\finenonce{005407}


\finexercice


\exercice{5410, rouget, 2010/07/06}
\enonce{005410}{**}
Soit $f$ une fonction convexe sur un intervalle ouvert $I$ de $\Rr$. Montrer que $f$ est continue sur $I$ et même dérivable à droite et à gauche en tout point de $I$.
\finenonce{005410}


\finexercice


\exercice{5411, rouget, 2010/07/06}
\enonce{005411}{*** Inégalités de convexité}
\begin{enumerate}
\item  Soient $x_1$, $x_2$,..., $x_n$, $n$ réels positifs ou nuls et $\alpha_1$,..., $\alpha_n$, $n$ réels strictement positifs tels que $\alpha_1+...+\alpha_n=1$. Montrer que $x_1^{\alpha_1}..x_n^{\alpha_n}\leq\alpha_1x_1+...+\alpha_nx_n$. En déduire que $\sqrt[n]{x_1...x_n}\leq\frac{x_1+...+x_n}{n}$.
\item  Soient $p$ et $q$ deux réels strictement positifs tels que $\frac{1}{p}+\frac{1}{q}=1$.
\begin{enumerate}
\item Montrer que, pour tous réels $a$ et $b$ positifs ou nuls, $ab\leq\frac{a^p}{p}+\frac{b^q}{q}$ avec égalité si et seulement si $a^p=b^q$. 
\item Soient $a_1$,..., $a_n$ et $b_1$,..., $b_n$, $2n$ nombres complexes. Montrer que~:

$$\left|\sum_{k=1}^{n}a_kb_k\right|\leq\sum_{k=1}^{n}|a_kb_k|\leq\left(\sum_{k=1}^{n}|a_k|^p\right)^{1/p} \left(\sum_{k=1}^{n}|b_k|^q\right)^{1/q}\;(\mbox{Inégalité de \textsc{Hölder}}).$$

\item Montrer que la fonction $x\mapsto x^p$ est convexe et retrouver ainsi l'inégalité de \textsc{Hölder}.
\item Trouver une démonstration directe et simple dans le cas $p=q=2$ (inégalité de \textsc{Cauchy}-\textsc{Schwarz}).
\end{enumerate}
\end{enumerate}
\finenonce{005411}


\finexercice


\exercice{5414, rouget, 2010/07/06}
\enonce{005414}{***I}
Montrer que la fonction définie sur $\Rr$ par $f(x)=e^{-1/x^2}$ si $x\neq0$ et $0$ si $x=0$ est de classe $C^\infty$ sur $\Rr$.
\finenonce{005414}


\finexercice


\exercice{5417, rouget, 2010/07/06}
\enonce{005417}{**** Toute fonction dérivée vérifie le théorème des valeurs intermédiaires}
Soit $f$ une fonction dérivable sur un intervalle ouvert $I$  à valeurs dans $\Rr$. Soient $a$ et $b$ deux points distincts de $I$ vérifiant $f'(a)<f'(b)$ et soit enfin un réel $m$ tel que $f'(a)<m<f'(b)$.
\begin{enumerate}
\item  Montrer qu'il existe $h>0$ tel que $\frac{f(a+h)-f(a)}{h}<m<\frac{f(b+h)-f(b)}{h}$.
\item  Montrer qu'il existe $y$ dans $[a,b]$ tel que $m=\frac{f(y+h)-f(y)}{h}$ puis qu'il exsite $x$ tel que $f'(x)=m$.
\end{enumerate}
\finenonce{005417}


\finexercice


\exercice{5419, rouget, 2010/07/06}
\enonce{005419}{*IT}
Etudier la dérivabilité à droite en $0$ de la fonction $f~:~x\mapsto\cos\sqrt{x}$.
\finenonce{005419}


\finexercice


\exercice{5420, rouget, 2010/07/06}
\enonce{005420}{**}
Soit $P$ un polynôme réel de degré supèrieur ou égal à $2$.
\begin{enumerate}
\item  Montrer que si $P$ n'a que des racines simples et réelles, il en est de même de $P'$.
\item  Montrer que si $P$ est scindé sur $\Rr$, il en est de même de $P'$.
\end{enumerate}
\finenonce{005420}


\finexercice


\exercice{5422, rouget, 2010/07/06}
\enonce{005422}{**}
Soit $f$ de classe $C^1$ sur $\Rr_+^*$ telle que $\lim_{x\rightarrow +\infty}xf'(x)=1$. Montrer que $\lim_{x\rightarrow +\infty}f(x)=+\infty$.
\finenonce{005422}


\finexercice


\exercice{5423, rouget, 2010/07/06}
\enonce{005423}{***}
Soit $f$ de classe $C^1$ sur $\Rr$ vérifiant pour tout $x$ réel, $f\circ f(x)=\frac{x}{2}+3$. En remarquant que $f(\frac{x}{2}+3)=\frac{f(x)}{2}+3$, montrer que $f'$ est constante puis déterminer $f$.
\finenonce{005423}


\finexercice


\exercice{5424, rouget, 2010/07/06}
\enonce{005424}{***I}
Soit $f$ de classe $C^1$ sur $\Rr$ vérifiant $\lim_{x\rightarrow +\infty}(f(x)+f'(x))=0$. Montrer que $\lim_{x\rightarrow +\infty}f(x)=\lim_{x\rightarrow +\infty}f'(x)=0$. (Indication. Considérer $g(x)=e^xf(x)$).
\finenonce{005424}


\finexercice

\section{125.01 Formule de Taylor}


\exercice{1267, legall, 1998/09/01}
\video{MY2kvANRkJo}
\enonce{001267}{}
Soit $f$ l'application de $\Rr$ dans $\Rr$ définie par
$f(x)=\displaystyle{\frac{x^3}{1+x^6}}.$ Calculer $f^{(n)}(0)$ pour tout $n \in \Nn.$
\finenonce{001267}



\finexercice



\exercice{1268, legall, 1998/09/01}
\video{kxEXonvTl5g}
\enonce{001268}{}
Soit $a$ un nombre réel et $f : ] a , +\infty [ \rightarrow \Rr$ une application de classe $C^2$. 
On suppose $f$ et $f''$ bornées ; on pose $\displaystyle  M_0=\sup _{x>a}\vert f(x)\vert$ et 
$\displaystyle  M_2=\sup_{x> a}\vert f''(x)\vert$.
\begin{enumerate}
    \item En appliquant une formule de Taylor reliant $f(x)$ et $f(x+h)$, montrer que, pour tout 
$x>a$ et tout $h>0$, on a~: $\displaystyle \vert f'(x)\vert \leq \frac{h}{2}M_2+\frac{2}{h}M_0$.

    \item En déduire que $f'$ est bornée sur $]a,+\infty[$.

    \item \'Etablir le résultat suivant : soit $g : ]0,+\infty[ \rightarrow \Rr$ une application
de classe $C^2$ à dérivée seconde bornée et telle que $\displaystyle \lim _{x\rightarrow +\infty }g(x)=0$. 
Alors $\displaystyle\lim _{x\rightarrow +\infty }g'(x)=0$.
\end{enumerate}
\finenonce{001268}



\finexercice

\exercice{1269, legall, 1998/09/01}

\enonce{001269}{}
Soient $  a,b,c \in { \Zz}  $ tels que~: $  ae^2+be+c=0  .$
\begin{enumerate}
    \item En appliquant la formule de
 Taylor sur $  [ 0,1]   $ \`a l'application $  \varphi : x \mapsto ae^x+ce^{-x}  $ d\'emontrer
que, pour tout $  n \in { \Nn}  $ il existe $  \theta _n \in ]0,1[   $ tel que~:
$$-b= \frac{ae^{\theta _n} +(-1)^n ce^{-\theta _n}}{ (n+1)!}+\sum _{k=0}^n \frac{a +(-1)^k c}{ k!}
   .$$
    \item En d\'eduire que pour $  n   $ assez grand $
ae^{\theta _n} +(-1)^n ce^{-\theta _n}=0  $ puis
que $  a=b=c=0  .$
(On rappelle que $  \displaystyle{e=\sum _{n=1}^{\infty}\frac{1}{ n!}  }.$)
\end{enumerate}
\finenonce{001269}



\finexercice



\exercice{1270, gourio, 2001/09/01}

\enonce{001270}{}
Soit $f\in C^{\infty }({\Rr},{\Rr})$ telle que $\forall n\in \Nn
,f^{(n)}(0)=0$ et $f^{(n)\text{ }}$ est born\'{e}e sur ${\Rr}$ avec
$\sup\limits_{x\in {\Rr}}\left| f^{(n)}(x)\right| =o(\frac{n!}{a^{n}}),a$
constante fix\'{e}e.
Montrer que $\forall x\in [-a,a],f(x)=0,$ puis que $f=0.$
\finenonce{001270}



\finexercice



\exercice{1271, gourio, 2001/09/01}

\enonce{001271}{}
Soit $P\in {\Rr}_{n}[X]$ tel que $P\geq 0$. On pose $Q=P+P^{\prime}+...
+P^{(n)}.$ Montrer que $Q\geq 0.$
\finenonce{001271}



\finexercice



\exercice{1272, ridde, 1999/11/01}

\enonce{001272}{}
Soient $a$ et $b$ deux r\'eels tels que $a < b$ 
et $f \in C^3 ([a,b], \Rr)$. Montrer qu'il existe 
$c \in ]a, b[$ tel que $f (b) = f(a) + (b-a)f' (\dfrac{a + b}2)
 + \dfrac{ (b-a)^3}{24} f''' (c)$ (on pourra utiliser Taylor-Lagrange entre
 $a, \dfrac{a + b}2, b$).
\finenonce{001272}



\finexercice



\exercice{2690, matexo1, 2002/02/01}
\enonce{002690}{}
Soit $f:[-a,a] \to \R$ une fonction de classe $C^2$. Montrer que
$$\forall x \in [-a, a],\qquad
|f'(x)| \leq  \frac 1{2a} |f(a) -f(-a)| + \frac{a^2+x^2}{2a}
\sup_{t \in [-a,a]} |f''(t)|.$$
Application : montrer que si $0 \leq x \leq \pi /2$, on a 
$\sin x \geq  x \cos x -x^2$.

\finenonce{002690}
\finexercice


\exercice{4004, quercia, 2010/03/11}
\enonce{004004}{Déterminant}
Soit $f:\R \to \R$ trois fois dérivable en $a$.
\'Etudier $\lim_{h\to0}\frac1{h^4}
\begin{vmatrix} 1  &f(a)     &f(a+h)  \cr
          1  &f(a+h)   &f(a+2h) \cr
          1  &f(a+2h)  &f(a+3h) \cr \end{vmatrix}$.


\finenonce{004004}



\finexercice


\exercice{4005, quercia, 2010/03/11}
\enonce{004005}{Dérivées nulles en $0$}

Soit $f:\R \to \R$ de classe $\mathcal{C}^\infty$ telle que :
\begin{align*}&\forall\ n \in \N,\ f^{(n)}(0) = 0, \\
           &\exists\ \lambda>0 \text{ tq } \forall\ n\in\N,\ \sup\limits_\R \bigl|f^{(n)}\bigr| \le \lambda^nn!.\\
\end{align*}
Montrer que $f$ est nulle sur l'intervalle $]-\frac1\lambda,\frac1\lambda[$, puis
sur $\R$.

\finenonce{004005}



\finexercice


\exercice{4006, quercia, 2010/03/11}
\enonce{004006}{Fonctions absolument monotones}

Soit $f:\R \to \R$ de classe $\mathcal{C}^\infty$ telle que pour tous $n\in\N$ et $x\in\R$,
on a $f^{(n)}(x) > 0$.

Montrer que pour tout entier $n$, $\frac{f(x)}{x^n}\to +\infty$ lorsque $x\to+\infty$.

\finenonce{004006}



\finexercice


\exercice{4007, quercia, 2010/03/11}
\enonce{004007}{Fonction $\mathcal{C}^\infty$ à support compact}

Soit $f : {\R^+} \to \R$ de classe $\mathcal{C}^\infty$ telle que $f(0) = 1$,
et : $\forall\ x \ge \frac12,\ f(x) = 0$.

\begin{enumerate}
  \item Montrer que $\sup\limits_{\R^+} \bigl|f^{(n)}\bigr| \ge 2^nn!$.
    
  \item Montrer que pour $n \ge 1$, $\sup\limits_{\R^+} \bigl|f^{(n)}\bigr| > 2^nn!$.
\end{enumerate}
\finenonce{004007}



\finexercice


\exercice{4008, quercia, 2010/03/11}
\enonce{004008}{Formule de Simpson}
\begin{enumerate}
  \item Soit $f:\R \to \R$ de classe $\mathcal{C}^5$, impaire, telle que $f'(0) = 0$
    et : $\forall\ x\in\R,\ |f^{(5)}(x)| \le M$.

    Montrer qu'il existe une constante $\lambda$ telle que :
    $\forall\ x \in \R,\ \left|f(x)-\frac x3f'(x)\right| \le \lambda M|x^5|$.
    

  \item Soit $f : {[a,b]} \to \R$ de classe $\mathcal{C}^5$ telle que :
    $$f'(a) = f'(b) = f'\left(\frac {a+b}2\right) = 0,  \qquad
      \text{et}\quad \forall\ x\in{[a,b]},\ |f^{(5)}(x)| \le M.$$
    Montrer que $\bigl|f(b)-f(a)\bigr| \le \frac {M(b-a)^5}{2880}$.

\end{enumerate}
\finenonce{004008}



\finexercice


\exercice{4009, quercia, 2010/03/11}
\enonce{004009}{$f'(x) - (f(b)-f(a))/(b-a)$}
Soit $f : {[a,b]} \to \R$ de classe $\mathcal{C}^2$. On note $M = \sup|f''|$ et on suppose $M>0$.

\begin{enumerate}
  \item Montrer que : $\forall\ x \in {]a,b[}$, on a $\left|f'(x) - \frac {f(b)-f(a)}{b-a}\right|
    < M\frac {b-a}2$.

  \item Si $\left|f'(a) - \frac {f(b)-f(a)}{b-a}\right| = M\frac {b-a}2$,
    montrer que $f$
    est polynomiale de degré inférieur ou égal à 2.

\end{enumerate}
\finenonce{004009}



\finexercice


\exercice{4010, quercia, 2010/03/11}
\enonce{004010}{Matexo}

Soit $f : {[-a,a]} \to \R$ de classe $\mathcal{C}^2$.
Montrer que~:

$$\forall\ x\in {[-a,a]},\
  |f'(x)| \le \frac1{2a}|f(a)-f(-a)| + \frac{a^2+x^2}{2a}\sup|f''|.$$



Application. Montrer que si $0\le x\le \pi/2$ on a
$\sin x \ge x\cos x - x^2$.
\finenonce{004010}



\finexercice


\exercice{4011, quercia, 2010/03/11}
\enonce{004011}{Limite de $\theta$}

Soit $f:\R \to \R$ de classe $\mathcal{C}^{n+1}$. Pour $a$ fixé, on écrit la
formule de Taylor-Lagrange :
$$f(a+h) = f(a) + \dots + \frac {h^{n-1}}{(n-1)!}f^{(n-1)}(a)
                        + \frac {h^n}{n!}f^{(n)}(a + h\theta_h).$$

Montrer que si $f^{(n+1)}(a) \ne 0$, alors pour $h$ suffisament petit, $\theta_h$ est
unique et $\theta_h \to \frac 1{n+1}$ lorsque $h\to0$.
\finenonce{004011}



\finexercice


\exercice{4012, quercia, 2010/03/11}
\enonce{004012}{}
\label{diffinie}Différences finies
Soit $f:\R \to \R$ de classe $\mathcal{C}^\infty$ et $h > 0$.
On pose : $$\Delta_hf(x) = \frac {f(x+h/2) - f(x-h/2)}h
  \qquad \text{et}\quad
\Delta_h^p = \underbrace{\Delta_h \circ \Delta_h \circ \dots \circ \Delta_h}%
_{p \text{ fois}}.$$

Par exemple, $\Delta_h^2f(x) = \frac {f(x+h) - 2f(x) + f(x-h)}{h^2}$.
\begin{enumerate}
  \item
  \begin{enumerate}
    \item   Montrer que : $\forall\ x \in \R,\ \exists\ \theta \in {]-1,1[} \text{ tq } \Delta_hf(x) = f'\left(x + \frac {\theta h}2\right)$.
    \item   Montrer que : $\forall\ x \in \R,\ \exists\ \theta'\in {]-1,1[} \text{ tq } \Delta_hf(x) = f'(x) + \frac {h^2}{24} f^{(3)}\left( x + \frac{\theta' h}2\right)$.
  \end{enumerate}
  \item   Montrer par récurrence sur $p$ que :
      $$\forall\ x \in \R,\ \exists\ \theta_p \in {]-p,p[} \text{ tq }
      \Delta_h^pf(x) = f^{(p)}(x) + \frac {ph^2}{24} f^{(p+2)}\left( x + \frac{\theta_p h}2\right).$$
\end{enumerate}
\finenonce{004012}



\finexercice


\exercice{4013, quercia, 2010/03/11}
\enonce{004013}{$f$ et $f''$ sont bornées}

Soit $f : \R \to \R$ une fonction de classe $\mathcal{C}^2$.
On suppose : $\forall\ x \in \R,\ |f(x)| \le \alpha$ et $|f''(x)| \le \beta$.

\begin{enumerate}
  \item Montrer que : $\forall\ h > 0,\ \forall\ x \in \R,\ |f'(x)| \le \frac {2\alpha}h + \frac {h\beta}2$.
  \item Pour quelle valeur de $h$ obtient-on la meilleure inégalité ?
    
\end{enumerate}
\finenonce{004013}



\finexercice


\exercice{4014, quercia, 2010/03/11}
\enonce{004014}{Inégalité sur $f'$}
Soit $f : {\R} \to {\R^+}$ une fonction de classe $\mathcal{C}^2$.
On suppose : $\forall\ x \in \R,\ |f''(x)| \le M$.

\begin{enumerate}
  \item Montrer que : $\forall\ x,y \in \R,\ f(x) + yf'(x) + \frac {y^2}2M \ge 0$.

  \item En déduire que : $\forall\ x \in \R,\ |f'(x)| \le \sqrt{2Mf(x)}$.
    
  \item On suppose que : $\forall\ x \in \R,\ |f'(x)| = \sqrt{2Mf(x)}$.
    Que pouvez-vous dire de $f$ ?
    

\end{enumerate}
\finenonce{004014}



\finexercice


\exercice{4015, quercia, 2010/03/11}
\enonce{004015}{Majoration des dérivées de $f$}
Soit $f:\R \to \R$ une fonction de classe $\mathcal{C}^n$ telle
que $f$ et $f^{(n)}$ sont bornées sur $\R$.
On veut montrer que les dérivées intermédiaires sont aussi bornées sur $\R$.

\begin{enumerate}
  \item   Cas $n = 2$ : Utiliser la formule de Taylor-Lagrange à l'ordre 2.
  \item   Cas général : Utiliser l'exercice \ref{diffinie}.


\end{enumerate}
\finenonce{004015}



\finexercice


\exercice{4016, quercia, 2010/03/11}
\enonce{004016}{$f''(x) \ge -\frac k{x^2}$}

Soit $f : {]0,+\infty[} \to \R$ de classe $\mathcal{C}^2$ telle que
$f(x) \to \ell\in\R$ lorsque $x\to0^+$, et : $\forall\ x > 0,\ f''(x) \ge -\frac k{x^2}$.

Montrer que $xf'(x) \to 0$ lorsque $x\to0^+$
(écrire la formule de Taylor-Lagrange à l'ordre 2 entre $x$ et $x+\varepsilon x$).

\finenonce{004016}



\finexercice


\exercice{4017, quercia, 2010/03/11}
\enonce{004017}{Ens PC$^*$ 2001}

Soient $P,Q$ deux polynômes à coefficients réels, non constants,
de coefficients dominants positifs.

On note $x_1 < x_2 < \dots < x_p$
les racines de $P'$ de multiplicités $m_1,\dots,m_p$
et $y_1 < y_2 < \dots < y_q$ celles de $Q'$ de multiplicités $n_1,\dots,n_q$.
Montrer qu'il existe $f$, $\mathcal{C}^1$ difféomorphisme croissant de~$\R$ sur~$\R$,
tel que $P\circ f = Q$ si et seulement si~:
$$p=q,
  \qquad\forall\ i,\ P(x_i) = Q(y_i),  \qquad\forall\ i,\ m_i = n_i.$$
\finenonce{004017}




\finexercice


\exercice{5418, rouget, 2010/07/06}
\enonce{005418}{****}
Soit $f$ une fonction de classe $C^3$ sur $\Rr$ vérifiant~:~$\forall(x,y)\in\Rr^2,\;f(x+y)f(x-y)\leq(f(x))^2$.
Montrer que $\forall x\in\Rr,\;f(x)f''(x)\leq(f'(x))^2$ (Indication. Appliquer la formule de \textsc{Taylor}-\textsc{Laplace} entre $x$ et $x+y$ puis entre $x$ et $x-y$).
\finenonce{005418}


\finexercice


\exercice{5457, rouget, 2010/07/10}
\enonce{005457}{***}
Partie principale quand $n$ tend vers $+\infty$ de $u_n=\sum_{k=1}^{n}\sin\frac{1}{(n+k)^2}$. 
\finenonce{005457}


\finexercice

\section{125.02 Calculs}


\exercice{1237, legall, 1998/09/01}

\enonce{001237}{}
 Donner le d\'eveloppement limit\'e en $  0  $ des
fonctions~:
\begin{enumerate}
    \item $  x\mapsto \ln (\hbox{cos}
(x))  $ (\`a l'ordre $  6  $).
    \item $  x \mapsto \tan(x)  $ (\`a l'ordre $ 7  $).
    \item $  x \mapsto \sin(\tan(x))  $ (\`a l'ordre $ 7  $).
    \item $  x\mapsto (\ln(1+x))^2  $ (\`a l'ordre $ 4  $).
    \item  $  x\mapsto \exp(\sin(x))  $ (\`a l'ordre $ 3  $).
    \item $  x \mapsto \sin^6(x)  $ (\`a l'ordre $ 9  .$)
\end{enumerate}

\finenonce{001237}



\finexercice



\exercice{1238, legall, 1998/09/01}

\enonce{001238}{}
\begin{enumerate}
    \item Soit $  f : { \Rr} \rightarrow { \Rr}   $ la fonction d\'efinie
par $  f(x)=0  $ si $  x \leq 0  $ et $  \displaystyle{ f(x)=\hbox{exp }(\frac{-1}{ x})  }$
sinon. Calculer, pour tout $  n \in { \Nn}  ,$ le d\'eveloppement limit\'e de $  f   $ en $  0  .$
Quelles conclusions en tirer~?
    \item Soit $  g : { \Rr} \rightarrow { \Rr}   $ la fonction d\'efinie
par $  g(0)=0  $ et, si $  x \not = 0  :$  $  \displaystyle{ g(x)=x^3\hbox{sin}
(\frac{1}{ x})  }.$ Montrer que $  g  $ a un d\'eveloppement limit\'e d'ordre $  2  $ en $  0  $
mais n'a pas de d\'eriv\'ee seconde (en $  0  $).
\end{enumerate}
\finenonce{001238}



\finexercice



\exercice{1239, ridde, 1999/11/01}

\enonce{001239}{}
 D\'eterminer la limite en $0$ de $\dfrac{\arctan x
- \sin x}{\tan x - \arcsin x}$.

\finenonce{001239}




\finexercice



\exercice{1240, ridde, 1999/11/01}

\enonce{001240}{}
 Faire un d\'eveloppement limit\'e ou asymptotique
en $a$ \`a l'ordre $n$ de :
\begin{enumerate}
\item $\ln \cos x$ $n = 6$ $a = 0$.

\item $\dfrac{\arctan x -x}{\sin x -x}$ $n = 2$ $a = 0$.

\item $\ln \tan (\frac x2  + \frac{\pi}4)$ $n = 3$ $a = 0$.

\item $\ln \sin x$ $n = 3$ $a = \frac{\pi}4$.

\item $\sqrt[3]{x^3 + x}-\sqrt[3]{x^3-x}$ $n = 4$ $a =  + \infty$.

\item $ (1 + x)^{\frac 1x}$ $n = 3$ $a = 0$.

\item $x (\sqrt{x^2 + \sqrt{x^4 + 1}}-x\sqrt 2)$ $n = 2$ $a =  + \infty$.

\end{enumerate}

\finenonce{001240}



\finexercice



\exercice{1241, roussel, 2001/09/01}

\enonce{001241}{}
D\'eveloppements limit\'es ~en $0$ de :
\begin{enumerate}
\item $\cos x. \ln (1+x)$ \`a l'ordre $4$.
\item $\displaystyle{\frac{1}{\cos x}}$ \`a l'ordre $4$.
\item $\arcsin \left ( \ln(1+x^2) \right )$ \`a l'ordre $6$.
\item $\displaystyle{\frac{\sinh x-x}{x^3}}$ ~\`a l'ordre $4$.
\item $(1+x)^{\frac{1}{1+x}}$ \`a l'ordre $3$.
\end{enumerate}
\finenonce{001241}



\finexercice



\exercice{1242, roussel, 2001/09/01}

\enonce{001242}{}
Pour chacune des fonctions suivantes, donner les conditions sur $\epsilon (x)$
pour que ces fonctions soient des d\'eveloppements limit\'es ~au voisinage d'un point et \`a un ordre
que vous pr\'eciserez.
\begin{enumerate}
\item $f_1(x)=x-\displaystyle{\frac{x^3}{3}}+x^2\epsilon (x)$
\item $f_2(x)=1-\displaystyle{\frac{2}{x^2}}+\displaystyle{\frac{1}{x^3}}+\displaystyle{\frac{1}{x^3}}\epsilon
(x)$
\item $f_3(x)=(x-2)+\displaystyle{\frac{(x-2)^2}{5}}+(x-2)^3\epsilon (x)$
\item $f_4(x)=x^2-x+1+\displaystyle{\frac{1}{x}}+\displaystyle{\frac{1}{x}}\epsilon (x)$
\item $f_5(x)=x^3+3x^2-x+1+(x-1)^2\epsilon (x)$
\item $f_6(x)=(x-2)^2+(x-2)-2+(x-2)\epsilon (x)$
\item $f_7(x)=\{ 2x+x^2+1+x^2\epsilon (x)\}\{-x+3+x^2-x^3\epsilon (x)\}$
\end{enumerate}
\finenonce{001242}



\finexercice



\exercice{1243, roussel, 2001/09/01}
\video{eGl4xUQYh6Q}
\enonce{001243}{}
\begin{enumerate}
\item Développement limité en $1$ à l'ordre $3$ de $f(x)=\sqrt{x}$.

\item Développement limité en $1$ à l'ordre $3$ de $g(x)= e^{\sqrt{x}}$.

\item Développement limité à l'ordre $3$ en $\frac\pi3$ de $h(x)=\ln (\sin x)$.
\end{enumerate}

\finenonce{001243}



\finexercice


\exercice{1244, roussel, 2001/09/01}
\video{KCxPBmuttZE}
\enonce{001244}{}
Donner un développement limité à l'ordre $2$ de $f(x)=
\displaystyle{\frac{\sqrt{1+x^2}}{1+x+\sqrt{1+x^2}}}$ en $0$.
En déduire un développement à l'ordre $2$ en $+\infty$.
Calculer un développement à l'ordre $1$ en $-\infty$.
\finenonce{001244}



\finexercice



\exercice{1245, roussel, 2001/09/01}

\enonce{001245}{}
Donner un d\'eveloppements limit\'e ~en $0$ \`a l'ordre $10$
de :
\begin{enumerate}
\item $x \longmapsto \displaystyle{\int_{0}^{x} \cos t^2 dt}.$
\item $x \longmapsto \displaystyle{\int_{x}^{x^2} \frac{1}{\sqrt{1+t^4}}dt}=
F(x^2)-F(x)$ o\`u $F$ est une primitive de
$t \longmapsto \displaystyle{\frac{1}{\sqrt{1+t^4}}}.$
\end{enumerate}
\finenonce{001245}



\finexercice



\exercice{1246, gourio, 2001/09/01}

\enonce{001246}{}
Donner le DL2 en $+\infty $ de :
$$x\rightarrow \sqrt{\frac{x-2}{x+1}}e^{\frac{x}{x-1}}.$$
\finenonce{001246}



\finexercice



\exercice{2657, matexo1, 2002/02/01}
\video{2r2CIG013os}
\enonce{002657}{}
Calculer
$$\ell = \lim_{x\to+\infty}\left(\frac{\ln(x+1)}{\ln x}\right)^x.$$
Donner un équivalent de 
$$\left(\frac{\ln(x+1)}{\ln x}\right)^x - \ell$$
lorsque $x \to +\infty$.
\finenonce{002657}


\finexercice


\exercice{2683, matexo1, 2002/02/01}

\enonce{002683}{}
A 

\begin{enumerate}
\item Soient $a$ et $z$ deux réels. Soit $f$ une fonction de
classe $C^{n+1}$ sur le segment d'extrémités
$a$ et $z$ et $\phi$ un polynôme de degré $n$. Prouver que pour tout $t$ compris dans
l'intervalle $[0,1]$,
\[
{d\over dt}\sum_{m=1}^n (-1)^m (z-a)^m
                \phi^{(n-m)}(t) f^{(m)}\left(a+t(z-a)\right)
\]\[
  = -(z-a) \phi^{(n)}(t) f'\left(a+t(z-a)\right)
    + (-1)^n (z-a)^{n+1} \phi(t) f^{(n+1)}\left(a+t(z-a)\right) 
\]

\item 

\begin{enumerate}
\item Montrer que la fonction $t\mapsto {t \over e^t-1}$ est prolongeable par continuité
en zéro, que son prolongement est indéfiniment dérivable et admet des développements
limités en zéro de la forme:
$$ 1 - {t/2} + {b_1 t^2 \over 2!}
                                + {b_2 t^4 \over 4!} + \dots + {b_n t^{2n} \over (2n)!}
+o(t^{2n+1}),$$
où les $b_i$ sont des réels qu'on ne cherchera pas à déterminer.

Montrer que la dérivée $n^{i\grave eme}$ en zéro, notée $\phi_n(z)$, de la fonction
$t\mapsto t {e^{zt}-1 \over e^t-1}$
est un polynôme en $z$ de degré $n$ et que
$$\phi_n(z)= z^n - {1\over 2}nz^{n-1} + C_n^2 b_1 z^{n-2}
                + C_n^4 b_2 z^{n-4} + \dots
                + C_n^{2N} b_N z^{n-2N}$$
où $N=E({n-1\over2})$, $E$ désignant la fonction partie entière.

\item
Prouver que $nz^{n-1}=\phi_n(z+1) - \phi_n(z)$
\end{enumerate}

\item 
Prouver que
$$
\begin{matrix}
(i)     &\phi_n^{(n-k)}(1) = \phi_n^{(n-k)}(0)\ \ (2\leq k\leq n) &\quad
& (ii)  &\phi_n^{(n-2k-1)}(0) = 0 \ \ (1\leq k\leq N)             \cr
(iii)   &\phi_n^{(n-2k)}(0) = {n! b_k \over (2k)!}  \ \ (1\leq k\leq N)  &\quad
& (iv)  &\phi_n^{(n-1)}(0) = -{1\over2} n!         \cr
 (v)    &\phi_n^{(n-1)}(1) = +{1\over2} n! & \quad
&(vi)   &\phi_n^{(n)} = n!  \cr
\end{matrix}
$$

\item 

\begin{enumerate}
\item 
On suppose $f$ de classe $C^{2n+1}$. Prouver que
\begin{eqnarray*}
0 &=& f(z) - f(a) - {z-a \over 2} \left[ f'(z) + f'(a) \right] 
\\
&&+ \sum_{m=1}^{n-1} b_m {(z-a)^{2m} \over (2m)!} \left[ f^{(2m)}(z) - f^{(2m)}(a) \right] 
\\
&& - {(z-a)^{2n+1} \over (2n)!} \int_0^1 \phi_{2n}(t) f^{(2n+1)}\Big( a+(z-a)t \Big) dt
 \end{eqnarray*}

\item
En déduire que si $F$ est de classe $C^{2n}$ sur $[a,a+r\omega]$
où $r\in \Nn$ et $\omega>0$, alors
\begin{eqnarray*}
\int_a^{a+r\omega} F(x)dx &=&
\omega \Big[ {1\over2} F(a) + F(a+\omega) + \dots 
+ F(a+{\scriptstyle (r-1)}\omega) + {1\over2} F(a+r\omega) \Big]
\\
&& - \sum_{m=1}^{n-1}      b_m {\omega^{2m} \over (2m)!} \left[ F^{(2m-1)}(a+r\omega) - F^{(2m-1)}(a) \right] 
+ R_n                                               
 \end{eqnarray*} 
où 
\[
R_n={\omega^{2n+1} \over (2n)!} \int_0^1 \phi_{2n}(t)
\sum_{m=0}^{r-1} F^{(2n)}(a+m\omega+\omega t) dt \; .
\]
\end{enumerate}
\end{enumerate}
                        
B

\begin{enumerate}
\item
Soit $u_k: x>0 \mapsto \ln(x+k)-\ln(k)+x\ln\left({k\over k+1}\right)
                                \qquad(k\in \Nn^*)$\par
Montrer que pour tout $x$ strictement positif, la série
$\sum_{k\geq 1}u_k(x)$ est convergente. On pose pour la suite
$G(x)=\ln(x)+\sum_{k=1}^{\infty}u_k(x)$

\item
Prouver que $G$ vérifie l'équation fonctionnelle 
$$\forall x>0 \quad G(x+1)=G(x)-\ln(x).$$

\item
En déduire que $\forall m\in \Nn \quad \exp(-G(m+1))=m!$

\item
Soit $x$ et $y$ deux réels strictement positifs. Montrer que la
série\par
\centerline{$\sum_{k\geq1}\bigl[ \ln(y+k) - \ln(x+k)
                                + (y-x)\ln\left({k\over k+1}\right) \bigr]$}
est convergente et que sa somme est $G(y) - G(x) - \ln y + \ln x$.

\item
Prouver à l'aide de A que pour tous entiers strictement positifs $n$
et $p$ 
\begin{eqnarray*}
\sum_{k=0}^n \ln(y+k) - \ln(x+k) &=&
\int_0^n f(t)dt + {1\over2} \big( f(0)+f(n) \big)
\\
&&+\sum_{h=1}^{p-1} {b_h\over(2h)!} \big( f^{(2h-1)}(n) - f^{(2h-1)}(0) \big) + T_{p,n}(x,y)
 \end{eqnarray*}
où $f: t\mapsto \ln(y+t) -\ln(x+t)$ et
$T_{p,n}(x,y)$ est une expression que l'on précisera.

\item
Prouver que $R_p(x,y)=\lim_{n \to +\infty}T_{p,n}(x,y)$ existe.

\item
On pose $g(z)=z\ln z - z - {1\over2} \ln z + \sum_{h=1}^{p-1} 
                                {b_h\over(2h)(2h-1)} {1\over z^{2h-1}}$.
Montrer que $G(y)+g(y)=G(x)+g(x)+R_p(x,y)$

\item
Montrer que $ R_p(x,y)  = O\left( {1 \over \bigl[ \inf(x,y) \bigr]^{2p-1}} \right)$
quand $\inf(x,y)\to +\infty$.

\item
Prouver à l'aide de la formule de Stirling que
$G(m)+g(m)\to {1\over2}\ln2\pi$ quand
$m \to +\infty$.

\item
Montrer que
$$G(y)= -y\ln y + y + {1\over2}\ln y - {1\over2}\ln2\pi
- \sum_{h=1}^{p-1}  {b_h\over(2h)(2h-1)} {1\over y^{2h-1}}
+ O({1\over y^{2p-1}})$$

\item
Donner un développement asymptotique de $\ln(m!)$ quand $m$ tend vers $+\infty$ à
un $O({1\over m^7})$ près.
\end{enumerate}

\finenonce{002683}
\finexercice


\exercice{4018, quercia, 2010/03/11}
\enonce{004018}{Calculs de DL}

\ \\
\textbf{Fonctions trigonométriques}
%
                    %+------------------------------+
                    %|  Fonctions trigonométriques  |
                    %+------------------------------+
\begin{align*}
x/\sin x              &= 1 + x^2/6 + 7x^4/360 + o(x^4)              \cr
1/\cos x              &= 1 + x^2/2 + 5x^4/24 + o(x^4)               \cr
\ln(\sin x/x)         &= -x^2/6 -x^4/180 -x^6/2835  + o(x^6)        \cr
\exp(\sin x/x)        &= e(1 - x^2/6 + x^4/45) + o(x^4)             \cr
\sqrt{\tan x}         &= 1 + h + h^2/2   + o(h^2), h = x-\pi/4                   \cr
\sin(x+x^2+x^3-x^4)   &= x + x^2 + 5x^3/6 -3x^4/2  + o(x^4)           \cr
\ln(x\tan(1/x))       &= x^{-2}/3 + 7x^{-4}/90  + o(1/x^4)              \cr
(1-\cos x)/(e^x-1)^2  &= 1/2 - x/2 + x^2/6   + o(x^2)                 \cr
\sin((\pi\cos x)/2)   &=1 -\pi^2x^4/32 + \pi^2x^6/192  + o(x^6)     \cr
\cos x\ln(1+x)        &= x - x^2/2 - x^3/6    + o(x^4)                \cr
(\sin x-1)/(\cos x+1) &= -1/2 + x/2 - x^2/8 + o(x^2)                 \cr
\ln(2\cos x+\tan x)   &= \ln2+x/2-5x^2/8+11x^3/24-59x^4/192  + o(x^4) \cr
e^{\cos x}            &= e(1 - x^2/2 + x^4/6)   + o(x^5)              \cr
\end{align*}


\textbf{Fonctions circulaires inverses}
%
                  %+----------------------------------+
                  %|  Fonctions circulaires inverses  |
                  %+----------------------------------+
\begin{align*}
\arcsin^2 x               &= x^2 + x^4/3 + 8x^6/45  + o(x^6)            \cr
1/\arcsin^2 x             &= x^{-2} - 1/3 - x^2/15 + o(x^2)               \cr
\arctan\sqrt{(x+1)/(x+2)} &= \pi/4 - x^{-1}/4 + 3x^{-2}/8      \cr
\arccos(\sin x/x)         &= |x|/\sqrt3(1 - x^2/90)  + o(1/x^3)            \cr
1/\arctan x               &= x^{-1} + x/3 -4x^3/45 +44x^5/945  + o(x^5) \cr
\arcsin\sqrt x            &= \pi/6+1/\sqrt3(2h-4h^2/3+32h^3/9) + o(h^3), h = x-1/4  \cr
\arcsin(\sin^2 x)         &= x^2 -x^4/3 +19x^6/90 -107x^8/630  + o(x^8)   \cr
\arctan(1+x)              &= \pi/4 + x/2 - x^2/4 + x^3/12  + o(x^4)     \cr
\arcsin x/(x-x^2)         &= 1 + x + 7x^2/6   + o(x^2)                  \cr
e^{\arcsin x}             &= e^{\pi/6}(1 + 2h/\sqrt3 + 2(1+\sqrt3)h^2/(3\sqrt3)) + o(h^2), h = x-1/2  \cr
e^{1/x}\arctan x          &= \frac\pi2+ (\frac\pi2-1)x^{-1} + (\frac\pi4-1)x^{-2} + (\frac\pi{12}-\frac16)x^{-3}  + o(1/x^3)  \cr
\end{align*}

\textbf{Exponentielle et logarithme}
%
                   %+-------------------------------+
                   %|  Exponentielle et logarithme  |
                   %+-------------------------------+
\begin{align*}
x/(e^x-1)          &= 1 - x/2 + x^2/12  + o(x^2)                       \cr
\ln x/\sqrt x      &= h - h^2 + 23h^3/24   + o(h^3), h = x-1           \cr
\ln((2-x)/(3-x^2)) &= \ln(2/3) - x/2 + 5x^2/24  + o(x^2)               \cr
\ln(1+x)/(1-x+x^2) &= x + x^2/2 - x^3/6   + o(x^3)                     \cr
\ch x/\ln(1+x)     &= x^{-1} + 1/2 + 5x/12    + o(x)                 \cr
\ln(\ln(1+x)/x)    &= -x/2 + 5x^2/24 - x^3/8 + o(x^3)                  \cr
\ln(a^x+b^x)       &= \ln2 + x\ln\sqrt{ab} + x^2\ln^2(a/b)/8 + o(x^2)  \cr
\exp(1/x)/x^2      &= e(1 - 3h + 13h^2/2 - 73h^3/6)  + o(h^3), h = x-1         \cr
\end{align*}

\textbf{Fonctions hyperboliques inverses}
%
                 %+------------------------------------+
                 %|  Fonctions hyperboliques inverses  |
                 %+------------------------------------+

\begin{align*}
\Argth(\sin x)   &= x + x^3/6 + x^5/24  + o(x^5)                  \cr
\Argsh(e^x)      &= \ln(1+\sqrt2) + 1/\sqrt2(x + x^2/4)  + o(x^2) \cr
\end{align*}


\textbf{Formes exponentielles}
%
                      %+-------------------------+
                      %|  Formes exponentielles  |
                      %+-------------------------+

\begin{align*}
(1-x+x^2)^{1/x}       &= e^{-1}(1 + x/2 + 19x^2/24) + o(x^2)           \cr
((1+x)/(1-x))^\alpha  &= 1 + 2\alpha x + 2\alpha^2x^2 + 2\alpha(2\alpha^2+1)x^3/3 + o(x^3) \cr
(\sin x/x)^{2/x^2}    &= e^{-1/3}(1 - x^2/90) + o(x^3)                \cr
(\sin x/x)^{3/x^2}    &= e^{-1/2}(1 - x^2/60 - 139x^4/151200)+ o(x^4) \cr
(1+\sin x)^{1/x}      &= e(1 - x/2 + 7x^2/24)  + o(x^2)                \cr
(1+\sin x + \cos x)^x &= 1 + x\ln2 + x^2(\ln^22+1)/2  + o(x^2)         \cr
(\sin x)^{\sin x}     &= 1 - h^2/2 + 7h^4/24  + o(h^4), h = x-\pi/2                 \cr
(\tan x)^{\tan2x}     &= e^{-1}(1 + 2h^2/3 + 4h^4/5)   + o(h^4), h = x-\pi/4     \cr
                      &   \text{Développer d'abord $\ln((1+x)/(1-x))$}   \cr
\end{align*}

\textbf{Radicaux}
%
                             %+------------+
                             %|  Radicaux  |
                             %+------------+
\begin{align*}
x\sqrt{(x-1)/(x+1)}   &= 1/\sqrt3(2 + 5h/3 + h^3/54)  + o(h^3), h = x-2              \cr
\sqrt{1+\sqrt{1-x}}   &= \sqrt2(1 - x/8 - 5x^2/128 - 21x^3/1024) + o(x^3)  \cr
\sqrt{1-\sqrt{1-x^2}} &= |x|/\sqrt2(1 + x^2/8 + 7x^4/128) + o(x^5)        \cr
e^x-\sqrt{1+2x}       &= x^2 - x^3/3 + 2x^4/3 - 13x^5/15  + o(x^5)        \cr
(\sqrt[3]{x^3+x^2}+\sqrt[3]{x^3-x^2})/x &= 2 - 2x^{-2}/9  + o(1/x^3)          \cr
\end{align*}

\finenonce{004018}



\finexercice


\exercice{4019, quercia, 2010/03/11}
\enonce{004019}{EIT 1999}
Calculer le développement limité de $\left(\frac{\tan x}x\right)^{1/x^2}$
en~$0$ à l'ordre~$3$.
\finenonce{004019}



\finexercice


\exercice{5426, exo7, 2010/07/06}
\enonce{005426}{IT}
Etudier l'existence et la valeur éventuelle des limites suivantes
\begin{enumerate}
 \item $\lim_{x\rightarrow \pi/2}(\sin x)^{1/(2x-\pi)}$
 \item $\lim_{x\rightarrow \pi/2}|\tan x|^{\cos x}$
 \item $\lim_{n\rightarrow +\infty}\left(\cos(\frac{n\pi}{3n+1})+\sin(\frac{n\pi}{6n+1})\right)^n$
 \item $\lim_{x\rightarrow 0}(\cos x)^{\ln|x|}$
 \item $\lim_{x\rightarrow \pi/2}\cos x.e^{1/(1-\sin x)}$
 \item $\lim_{x\rightarrow \pi/3}\frac{2\cos^2x+\cos x-1}{2\cos^2x-3\cos x+1}$
 \item $\lim_{x\rightarrow 0}\left(\frac{1+\tan x}{1+\tanh x}\right)^{1/\sin x}$
 \item $\lim_{x\rightarrow e,\;x<e}(\ln x)^{\ln(e-x)}$
 \item $\lim_{x\rightarrow 1,\;x>1}\frac{x^x-1}{\ln(1-\sqrt{x^2-1})}$
 \item $\lim_{x\rightarrow +\infty}\frac{x\ln(\ch x-1)}{x^2+1}$
 \item $\lim_{x\rightarrow 0,\;x>0}\frac{(\sin x)^x-x^{\sin x}}{\ln(x-x^2)+x-\ln x}$
 \item $\lim_{x\rightarrow +\infty}\left(\frac{\ln(x+1)}{\ln x}\right)^x$
 \item $\lim_{x \rightarrow 1/\sqrt{2}}\frac{(\Arcsin x)^2-\frac{\pi^2}{16}}{2x^2-1}$
 \item $\lim_{x\rightarrow +\infty}\left(\frac{\cos(a+\frac{1}{x})}{\cos a}\right)^x\;(\mbox{où}\;\cos a\neq0)$
\end{enumerate}
\finenonce{005426}


\finexercice


\exercice{5427, exo7, 2010/07/06}
\enonce{005427}{IT}
Déterminer les développements limités à l'ordre demandé au voisinage des points indiqués~:
\begin{enumerate}
 \item $\frac{1}{1-x^2-x^3}\;(\mbox{ordre}\;7\;\mbox{en}\;0)$
 \item $\frac{1}{\cos x}\;(\mbox{ordre}\;7\;\mbox{en}\;0)$
 \item $\Arccos\sqrt{\frac{x}{\tan x}}\;(\mbox{ordre}\;3\;\mbox{en}\;0)$
 \item $\tan x\;(\mbox{ordre}\;3\;\mbox{en}\;\frac{\pi}{4})$
 \item $(\ch x)^{1/x^2}\;(\mbox{ordre}\;2\;\mbox{en}\;0)$
 \item $\tan^3x(\cos(x^2)-1)\;(\mbox{ordre}\;8\;\mbox{en}\;0)$
 \item $\frac{\ln(1+x)}{x^2}\;(\mbox{ordre}\;3\;\mbox{en}\;1)$ 
 \item $\Arctan(\cos x)\;(\mbox{ordre}\;5\;\mbox{en}\;0)$
 \item $\Arctan\sqrt{\frac{x+1}{x+2}}\;(\mbox{ordre}\;2\;\mbox{en}\;0)$
 \item $\frac{1}{x^2}-\frac{1}{\Arcsin^2x}\;(\mbox{ordre}\;5\;\mbox{en}\;0)$
 \item $\int_{x}^{x^2}\frac{1}{\sqrt{1+t^4}}\;dt\;(\mbox{ordre}\;10\;\mbox{en}\;0)$
 \item $\ln\left(\sum_{k=0}^{99}\frac{x^k}{k!}\right)\;(\mbox{ordre}\;100\;\mbox{en}\;0)$
 \item $\tan\sqrt[3]{4(\pi^3+x^3)}\;(\mbox{ordre}\;3\;\mbox{en}\;\pi)$
\end{enumerate}
\finenonce{005427}


\finexercice


\exercice{5428, exo7, 2010/07/06}
\enonce{005428}{***}
Soit $0<a<b$. Etude complète de la fonction $f(x)=\left(\frac{a^x+b^x}{2}\right)^{1/x}$. 
\finenonce{005428}


\finexercice


\exercice{5429, exo7, 2010/07/06}
\enonce{005429}{**}
Etude au voisinage de $+\infty$ de $\sqrt{x^2-3}-\sqrt[3]{8x^3+7x^2+1}$.
\finenonce{005429}


\finexercice


\exercice{5430, exo7, 2010/07/06}
\enonce{005430}{**}
Soit $f(x)=\frac{x}{1-x^2}$. Calculer $f^{(n)}(0)$ en moins de $10$ secondes puis $f^{(n)}(x)$ pour $|x|\neq1$ en à peine plus de temps).
\finenonce{005430}


\finexercice


\exercice{5431, exo7, 2010/07/06}
\enonce{005431}{IT}
\begin{enumerate}
\item  Equivalent simple en $+\infty$ et $-\infty$ de $\sqrt{x^2+3x+5}-x+1$.
\item  Equivalent simple en $0$, $1$, $2$ et $+\infty$ de $3x^2-6x$
\item Equivalent simple en $0$ de $(\sin x)^{x-x^2}-(x-x^2)^{\sin x}$.
\item  Equivalent simple en $+\infty$ de $x^{\tanh x}$.
\item  Equivalent simple en $0$ de $\tan(\sin x)-\sin(\tan x)$.
\end{enumerate}
\finenonce{005431}


\finexercice


\exercice{5432, exo7, 2010/07/06}
\enonce{005432}{**IT}
Développement asymptotique à la précision $\frac{1}{n^3}$ de $u_n=\frac{1}{n!}\sum_{k=0}^{n}k!$.
\finenonce{005432}


\finexercice


\exercice{5433, exo7, 2010/07/06}
\enonce{005433}{**IT}
\begin{enumerate}
\item  Développement asymptotique à la précision $x^2$ en $0$ de $\frac{1}{x(e^x-1)}-\frac{1}{x^2}$.
\item  Développement asymptotique à la précision $\frac{1}{x^3}$ en $+\infty$ de $x\ln(x+1)-(x+1)\ln x$.
\end{enumerate}
\finenonce{005433}


\finexercice


\exercice{5434, exo7, 2010/07/06}
\enonce{005434}{**}
Soient $a>0$ et $b>0$. Pour $n\in\Nn^*$ et $x\in\Rr$, on pose $f_n(x)=\left(1+\frac{x}{n}\right)^n$.
\begin{enumerate}
\item  Equivalent simple quand $n$ tend vers $+\infty$ de $f_n(a+b)-f_n(a)f_n(b)$.
\item  Même question pour $e^{-a}f_n(a)-1+\frac{a^2}{2n}$.
\end{enumerate}
\finenonce{005434}


\finexercice


\exercice{5437, rouget, 2010/07/06}
\enonce{005437}{***I}
\begin{enumerate}
\item  Montrer que l'équation $\tan x=x$ a une unique solution dans l'intervalle $[n\pi,(n+1)\pi]$ pour $n$ entier naturel donné. On note $x_n$ cette solution.
\item  Trouver un développement asymptotique de $x_n$ à la précision $\frac{1}{n^2}$.
\end{enumerate}
\finenonce{005437}


\finexercice


\exercice{5438, rouget, 2010/07/06}
\enonce{005438}{}
\begin{enumerate}
\item  Montrer que l'équation $x+\ln x=k$ admet, pour $k$ réel donné, une unique solution dans $]0,+\infty[$, notée $x_k$.
\item  Montrer que, quand $k$ tend vers $+\infty$, on a~:~$x_k=ak+b\ln k+c\frac{\ln k}{k}+o\left(\frac{\ln k}{k}\right)$ où $a$, $b$ et $c$ sont des constantes à déterminer.
\end{enumerate}
\finenonce{005438}


\finexercice

\exercice{5439, rouget, 2010/07/10}
\enonce{005439}{**}
Soit $f(x)=1+x+x^2+x^3\sin\frac{1}{x^2}$ si $x\neq0$ et $1$ si $x=0$.
\begin{enumerate}
\item  Montrer que $f$ admet en $0$ un développement limité d'ordre $2$.
\item  Montrer que $f$ est dérivable sur $\Rr$.
\item  Montrer que $f'$ n'admet en $0$ aucun développement limité d'aucun ordre que ce soit.
\end{enumerate}
\finenonce{005439}


\finexercice


\exercice{5440, rouget, 2010/07/10}
\enonce{005440}{**IT}
Etude au voisinage de $0$ de $f(x)=\frac{1}{x}-\frac{1}{\Arcsin x}$ (existence d'une tangente~?)
\finenonce{005440}


\finexercice


\exercice{5441, rouget, 2010/07/10}
\enonce{005441}{**I}
\begin{enumerate}
\item  La fonction $x\mapsto\Arccos x$ admet-elle en $1$ (à gauche) un développement limité d'ordre $0$~?~d'ordre $1$~?
\item  Equivalent simple de $\Arccos x$ en $1$.
\end{enumerate}
\finenonce{005441}


\finexercice


\exercice{5442, rouget, 2010/07/10}
\enonce{005442}{***}
\begin{enumerate}
\item  Développement limité à l'ordre $n$ en $0$ de $f(x)=\frac{1}{(1-x)^2(1+x)}$.
\item  Soit $a_k$ le $k$-ème coefficient. Montrer que $a_k$ est le nombre de solutions dans $\Nn^2$ de l'équation $p+2q=k$.
\end{enumerate}
\finenonce{005442}


\finexercice


\exercice{6888, bodin, 2012/09/05}
\video{oBG5H6nDyhk}
\enonce{006888}{} Donner le développement limité en $0$ des fonctions :

\begin{enumerate}
  \item $\cos x \cdot \exp x$ \quad à l'ordre $3$

  \item $\left( \ln (1+x) \right)^2$ \quad à l'ordre $4$

  \item $\displaystyle{\frac{\sh x-x}{x^3}}$ \quad à l'ordre $6$

  \item $\exp\big(\sin(x)\big)$ \quad à l'ordre $4$

  \item  $\sin^6(x)$ \quad à l'ordre $9$

  \item $\ln \big(\cos(x)\big)$ \quad à l'ordre $6$

  \item $\displaystyle{\frac{1}{\cos x}}$ \quad à l'ordre $4$

  \item $\tan x$ \quad à l'ordre $5$ (ou $7$ pour les plus courageux)

  \item $(1+x)^{\frac{1}{1+x}}$ \quad à l'ordre $3$

  \item $\arcsin \left ( \ln(1+x^2) \right )$ \quad à l'ordre $6$
\end{enumerate}
\finenonce{006888}



\finexercice

\section{125.03 Applications}


\exercice{1247, vignal, 2001/09/01}
\video{HpQ-7NAmPFs}
\enonce{001247}{}
 Calculer les limites suivantes
$$\lim_{x\rightarrow 0}\frac{e^{x^2}-\cos x}{x^2}
\quad\quad\lim_{x\rightarrow 0}\frac{\ln (1+x)-\sin x}{x}
\quad\quad \lim_{x\rightarrow 0}\frac{\cos x-\sqrt{1-x^2}}{x^4}$$
\finenonce{001247}



\finexercice


\exercice{1248, legall, 1998/09/01}

\enonce{001248}{}
 Calculer les limites suivantes~: $\displaystyle{
\lim _{x\rightarrow 0} \frac{e^x-(\cos(x)+x)}{ x^2} ,
 \lim _{x\rightarrow 0} \frac{x^3\arctan(x)-x ^4}{\cos(x^2)-1}.}$

\finenonce{001248}



\finexercice



\exercice{1249, legall, 1998/09/01}
\video{LJLvAjm8KcY}
\enonce{001249}{}
\'Etudier la position du graphe de l'application $x\mapsto \ln(1+x+x^2)$ par rapport 
à sa tangente en $0$ et $1$.
\finenonce{001249}



\finexercice



\exercice{1250, legall, 1998/09/01}

\enonce{001250}{}
Montrer que pour tout $  n \in { \Nn}  ,  \displaystyle {\lim
_{x\rightarrow +\infty }\frac{ e^x}{ x^n}=+\infty}  .$
\finenonce{001250}



\finexercice



\exercice{1251, legall, 1998/09/01}

\enonce{001251}{}
\'Etablir pour tout $  x \in { \Rr}_+^*  $ l'in\'egalit\'e~:
$$\frac{3}{ 2} \sqrt{x}+\frac{3}{ 8 \sqrt{x+1}}<(x+1)^{3/2}-x^{3/2}<\frac{3}{ 2} \sqrt{x}+
\frac{3}{8 \sqrt{x}}.$$
\finenonce{001251}



\finexercice



\exercice{1252, legall, 1998/09/01}

\enonce{001252}{}
Montrer que pour tout $  x \in { \Rr}_+  ,  \displaystyle{\frac{x^2}{ 2}
 \leq e^x-x-1\leq  \frac{x^2}{ 2} e^x}  .$
\finenonce{001252}



\finexercice



\exercice{1253, ridde, 1999/11/01}

\enonce{001253}{}
Soit $f (x) = (\cos x)^{\frac 1x}$ pour $x \in ]-\frac{\pi}2, \frac{\pi}2[-
\left\{ 0\right\}$.
\begin{enumerate}
\item Montrer que $f$ est prolongeable par continuit\'e en $0$.
\item D\'eterminer un DL de $f$ en $0$ \`a l'ordre $2$.
\item Etudier la d\'erivabilit\'e du prolongement de $f$.
\end{enumerate}
\finenonce{001253}



\finexercice



\exercice{1254, ridde, 1999/11/01}

\enonce{001254}{}
\'Etudier les branches infinies des fonctions :
\begin{enumerate}
\item $f (x) = x^2 \arctan (\frac 1{1 + x^2})$.
\item $g (x) = x \sqrt{\frac{x-1}{3x + 1}}$.
\end{enumerate}
\finenonce{001254}



\finexercice



\exercice{1255, ridde, 1999/11/01}

\enonce{001255}{}
Soit $ (1)$ l'\'equation $x-E (x) = \frac 1 {x^2}$.
\begin{enumerate}
\item Montrer que pour tout $n \in \Nn^*$ il existe un unique $x_n \in [n, n + 1[$
solution de $ (1)$.
\item D\'eterminer un \'equivalent de $x_n$.
\item Faire un DAS de $x_n-n$ en $ + \infty$ en fonction de $\frac 1n$ \`a l'ordre $5$.
\end{enumerate}
\finenonce{001255}



\finexercice



\exercice{1256, gourio, 2001/09/01}

\enonce{001256}{}
Calculer pour $a\in \Rr^{+*}$ :
$$\lim\limits_{x\rightarrow a}\frac{x^{a}-a^{x}}{x^{x}-a^{a}}%
,\lim\limits_{n\rightarrow \infty }(3(2)^{\frac{1}{n}}-2(3)^{\frac{1}{n}})^{n}$$
\finenonce{001256}



\finexercice



\exercice{1257, gourio, 2001/09/01}

\enonce{001257}{}
Calculer :
$$\ell=\lim\limits_{x\rightarrow 0 }\left( \frac{\ln x+1}{\ln x}\right)^{x\ln x}$$
et donner un \'{e}quivalent de $\left( \frac{\ln x+1}{\ln x}\right) ^{x\ln x}-\ell$
quand $x\rightarrow 0 .$
\finenonce{001257}



\finexercice



\exercice{1258, gourio, 2001/09/01}

\enonce{001258}{}
Soit $x\in \Rr^{+}, $ on d\'{e}finit $(u_{n}(x))_{n}$ et $(v_{n}(x))_{n} $
par :
$$\forall n\in \Nn,u_{n+1}(x)=\frac{u_{n}(x)+v_{n}(x)}{2},v_{n+1}(x)=
\sqrt{u_{n}(x)v_{n}(x)},u_{0}(x)=1,v_{0}(x)=x.$$
\begin{enumerate}
\item Montrer que ces deux suites convergent vers une m\^{e}me limite $\ell_{x}.$

\item Soit $f:\Rr^{+}\rightarrow \Rr$ d\'{e}finie par: $f(x)=\ell_{x}. $ Calculer
$f(1),f(0),$ donner $f(\frac{1}{x})$ en fonction de $f(x)$ si $x>0.$
Montrer que $f$ est croissante, en d\'{e}duire le sens de variations de
$x\rightarrow \frac{f(x)}{x}.$

\item Montrer que $f$ est d\'{e}rivable en $1$ (on utilisera $\sqrt{x}\leq
f(x)\leq \frac{1+x}{2}$) puis que $\lim_{x\rightarrow \infty }f(x)=+\infty .$

\item Montrer que $f$ est continue sur $\Rr^{+*}, $ puis que $f$ est continue en $0.$

\item Donner l'allure du graphe de $f,$ pr\'{e}ciser la tangente en $0$ ainsi que
le comportement asymptotique en $+\infty .$
\end{enumerate}
\finenonce{001258}



\finexercice



\exercice{1259, gourio, 2001/09/01}

\enonce{001259}{}
Soit $n\in \Nn^{*},x\neq 0,$ on d\'{e}finit :
$$u_{n}(x)=\left( \frac{1^{x}+2^{x}+...+n^{x}}{n}\right) ^{\frac{1}{x}}.$$

D\'{e}terminer $\ell_{n}=\lim\limits_{x\rightarrow 0}u_{n}(x).$
\finenonce{001259}



\finexercice



\exercice{1260, gourio, 2001/09/01}

\enonce{001260}{}
D\'{e}terminer :
$$\lim\limits_{x\rightarrow 0}\frac{2\tan x-\sh 2x}{(1-\cos 3x)\arctan x}.$$
\finenonce{001260}



\finexercice



\exercice{1261, gourio, 2001/09/01}

\enonce{001261}{}
Soient $u,v,f$ d\'{e}finies par :
$$u(x)=(x^{3}-2x^{2}+1)^{\frac{1}{3}},\ v(x)=\sqrt{x^{2}+x+1},\ f(x)=u(x)-v(x).$$
\begin{enumerate}
\item
Donner un \'{e}quivalent de $f$ au voisinage de $-\infty $, en d\'{e}duire
$\lim\limits_{-\infty }f.$

\item D\'{e}teminer $\lim\limits_{x\rightarrow -\infty
}u(x)-x,\lim\limits_{x\rightarrow -\infty }v(x)+x.$ En d\'{e}duire
l'\'{e}quation d'une droite asymptote au graphe de $f $ en $-\infty $ et
positionner $f$ par-rapport \`{a} cette asymptote.

\item M\^{e}me \'{e}tude en $+\infty $.
\end{enumerate}
\finenonce{001261}



\finexercice



\exercice{1262, legall, 2003/10/01}

\enonce{001262}{}
Soit $g$ la
fonction $x\mapsto \dfrac{\arctan x}{(\sin x)^3}-\dfrac {1}{x^2}$.
\begin{enumerate}
\item Donner le domaine de d\'efinition de $g$.
\item Montrer qu'elle se prolonge par continuit\'e en $0$ en
une fonction
d\'erivable.
\item D\'eterminer la tangente en $0$ au graphe de cette fonction et la
position de ce graphe par rapport \`a celle-ci.
\end{enumerate}


\finenonce{001262}



\finexercice



\exercice{1263, legall, 2003/10/01}

\enonce{001263}{}
Soient $ f:x\mapsto \dfrac{x^3+2}{x^2-1}$ et $g
:x\mapsto (x+1)
\exp (\dfrac{1}{x-1})$ deux fonctions. D\'eterminer si leurs graphes
respectifs ont des asymptotes puis la position
de ces graphes par rapport \`a celles-ci.
\finenonce{001263}



\finexercice



\exercice{1264, legall, 2003/10/01}

\enonce{001264}{}
 Montrer que, pour tout $x$ r\'eel v\'erifiant $\vert 
x \vert \leq 1$~:
$$  \displaystyle{ \left \vert \frac
{x+\sin 2x}{x^9+x^2-3
}\right \vert }\leq 2 .$$


\finenonce{001264}



\finexercice




\exercice{1265, legall, 2003/10/01}
\video{c-IqNr313V8}
\enonce{001265}{}
Déterminer:
\begin{enumerate}
\item
\begin{enumerate}
\item $\displaystyle \lim _{x \rightarrow +\infty} \sqrt{x^2+3x+2} +x$
\item $\displaystyle \lim _{x \rightarrow -\infty} \sqrt{x^2+3x+2} +x$
\end{enumerate}

\item $\displaystyle \lim _{x \rightarrow 0^+}(\Arctan x)^{\frac{1}{x^2}}$

\item $\displaystyle \lim _{x \rightarrow 0} \frac{(1+3x)^{\frac{1}{3}}-1-\sin x}{1-\cos x}$
\end{enumerate}

\finenonce{001265}


\finexercice



\exercice{1266, legall, 2003/10/01}

\enonce{001266}{}
\begin{enumerate}
\item Soit $g$ la fonction d\'efinie par~: $$\displaystyle{ g(x) 
=\frac{x+1}{1+x^2} +\Arctan x}.$$
\begin{enumerate}
\item Quel est le domaine de d\'efinition de $g$~?
\item Etudier ses
variations.
\item Montrer que $g$ s'annule une et une seule fois sur $\Rr$ en un 
point $\alpha $ compris entre $-1$ et $0$ (on ne demande pas de 
pr\'eciser la valeur de $\alpha$).
\item Dessiner le graphe de $g$.
\end{enumerate}
\item Soit $f$ la fonction d\'efinie sur $\Rr$ par~: $$f(x)=(x+1)\Arctan x .$$
\begin{enumerate}
\item Calculer la d\'eriv\'ee de $f$ et \'etablir son tableau de variation.
\item Le graphe de $f$ a-t-il des points
d'inflexion~? Si oui, donner les coordonn\'ees de ce (ou ces) point(s).
\end{enumerate}
\item Donner l'\'equation de la tangente au point d'abcisse $x=0$ au 
graphe de $f$ et la position de ce graphe par rapport \`a cette 
tangente (au voisinage de ce point).

\item En utilisant les r\'esultats de l'exercice \textbf{???}, montrer que~:
\begin{enumerate}
\item $\displaystyle{\frac{ f(x) 
}{x}=(1+\frac{1}{x})(\frac{\pi}{2}-\Arctan \frac{1}{x}}) \hbox{ si 
}x>0 .$
\item $\displaystyle{\frac{ f(x) }{x}=(1+\frac{1}{x})(-\frac{\pi}{2}-\Arctan
\frac{1}{x})}\hbox{ si }x<0.$
\end{enumerate}
\item En d\'eduire l'existence d'une fonction $\epsilon $ telle 
que $\displaystyle {\lim _{x\rightarrow +\infty}\epsilon 
(\frac{1}{x})=0}$
et, pour tout $x>0$, on ait~: $$\displaystyle{ f(x)=\frac{\pi}{2}x+ 
(\frac{\pi}{2}-1)-\frac{1}{x}+\frac{1}{x}\epsilon 
(\frac{1}{x})}.$$ Etablir un r\'esultat analogue pour $x<0$.

\item Quelles sont les asymptotes au graphe de $f$~? Pr\'eciser la 
position de ce graphe par rapport \`a ces asymptotes.
\item Dessiner le graphe de $f$.
\end{enumerate}
\finenonce{001266}



\finexercice



\exercice{4020, quercia, 2010/03/11}
\enonce{004020}{Fonctions circulaires et hyperboliques}
\begin{enumerate}
  \item $\frac 1{\sin^2x} - \frac 1{\sh^2x}                      \to \frac 23$ lorsque  $x\to 0$.
  \item $\frac {\sin x\sh x - \tan x\tanh x}{\sh^4x - \tanh^4x}       \to -> -\frac1{12}$ lorsque  $x\to 0$.
  \item $(\ch x)^\alpha - (\sh x)^\alpha                           \to 
  \begin{cases}+\infty &\text{ si }\alpha>2,\cr 1 &\text{ si } \alpha=2,\cr 0 &\text{ si } \alpha<2.\cr\end{cases}$ lorsque  $x\to +\infty$
  \item $\frac {\exp(x^2) - \ch(x\sqrt2\,)}{(\ch x - \cos x)(\ch2x - \cos2x)} \to \frac1{12}$ lorsque  $x\to 0$.
  \item $(2x^2-3x+1)\tan \pi x                                     \to \frac1\pi$ lorsque  $x\to 1/2$.
  \item $\frac{\cos\pi x}{4x^2-9}                                 \to  -> \frac\pi{12}$ lorsque  $x\to 3/2$.
  \item $\frac{\sin 3x}{1-2\cos x}                                \to-\sqrt3$ lorsque  $x\to \pi/3$.
  \item $\frac{e^{\sin x}-e^x}{\sin x-x}                          \to 1$ lorsque  $x\to 0$.
  \item $\frac 1x \ln\ch x                                        \to 1$ lorsque  $x\to +\infty$.
\end{enumerate}
\finenonce{004020}



\finexercice


\exercice{4021, quercia, 2010/03/11}
\enonce{004021}{Logarithme et exponentielle}
\begin{enumerate}
  \item $\frac 1x\ln\Bigl(\frac {e^x-1}x \Bigr)                  \to \frac 12$. lorsque  $x\to 0$
  \item $\frac {x^x-1}{\ln x - x + 1}                             \to \pm\infty$ lorsque  $x\to 1$.
  \item $\frac {x^a - a^x}{\log_a(x) - \log_x(a)}                 \to \frac{a^{a+1}\ln a(1-\ln a)}2$ lorsque  $x\to a$.
  \item $\left(\frac{a^x+b^x}{1+c^x}\right)^{1/x}                 \to \exp\left(\frac{a+b-c}2\right)$ lorsque  $x\to 0$.
  \item $\frac{x^{x^x}}{x^x-1}                                    \to 0$ lorsque  $x\to 0^+$.
\end{enumerate}
\finenonce{004021}



\finexercice


\exercice{4022, quercia, 2010/03/11}
\enonce{004022}{Exposants variables}
\begin{enumerate}
  \item $x^{\arcsin x}                                             \to 1$ lorsque  $x\to 0^+$.
  \item $\frac{(\sin x)^{\sin x}-1}{x^x-1}                         \to 1$ lorsque  $x\to 0^+$.
  \item $(2-x)^{\tan(\pi x/2)}                                     \to e^{2/\pi}$ lorsque  $x\to 1$.
  \item $(2-x)^{\tan(\pi x/2)}                                     \to -> 1$ lorsque  $x\to 2^-$.
  \item $(\sin x + \cos x)^{1/x}                                   \to e$ lorsque  $x\to 0$.
  \item $(\cos 2x - 2\sin x)^{1/x}                                 \to e^{-2}$. lorsque  $x\to 0$
  \item $(\sin x)^{\tan x}                                         \to 1$. lorsque  $x\to \pi/2$
  \item $(\tan x)^{\cos x/\cos 2x}                                 \to e^{-1/\sqrt2}$ lorsque  $x\to \pi/4$.
  \item $(\tan x)^{\cos x/\cos 2x}                                 \to 1$ lorsque  $x\to (\pi/2)^-$.
  \item $(\sin x)^{1/\ln x}                                        \to e$ lorsque  $x\to 0^+$.
  \item $(\ln x)^{x-1}                                             \to 1$ lorsque  $x\to 1^+$.
  \item $(\ln x)^{\ln(e-x)}                                        \to -> 1$ lorsque  $x\to e^-$.
\end{enumerate}
\finenonce{004022}



\finexercice


\exercice{4023, quercia, 2010/03/11}
\enonce{004023}{Radicaux}
\begin{enumerate}
  \item $\frac {\sqrt{x+3} - \sqrt[3]{3x+5}}{1 - \tan(\pi x/4)}     \to 0$ lorsque  $x\to 1$.
  \item $\sh\sqrt{x^2+x} - \sh\sqrt{x^2-x}                          \to +\infty$ lorsque  $x\to +\infty$.
  \item $\frac{\sqrt{3x+1}-\sqrt{x+1}}{\sin x}                      \to 1$ lorsque  $x\to 0$.

\end{enumerate}
\finenonce{004023}



\finexercice


\exercice{4024, quercia, 2010/03/11}
\enonce{004024}{Sommes de cotangentes}

Soient $a_1,\dots,a_n \in \R$.
CNS pour que $\sum_{k=1}^n a_k \mathrm{cotan}(kx)$ ait une limite finie en 0 ?

\finenonce{004024}



\finexercice


\exercice{4025, quercia, 2010/03/11}
\enonce{004025}{$\left({\frac1n \sum_{k=0}^{n-1} \Bigl(1+\frac kn\Bigr)^{1/p}}\right)^p$}

On pose $u_{n,p} = \left({\frac1n \sum_{k=0}^{n-1} \Bigl(1+\frac kn\Bigr)^{1/p}}\right)^p$.
Trouver : $v_p = \lim_{n\to\infty} u_{n,p}$, $v = \lim_{p\to\infty} v_p$,
          $w_n = \lim_{p\to\infty} u_{n,p}$ et $w = \lim_{n\to\infty} w_n$.
\finenonce{004025}



\finexercice


\exercice{4026, quercia, 2010/03/11}
\enonce{004026}{Ensi P 91}

    Calculer $\lim_{n\to\infty} \sum_{k=1}^n \sin \frac k{n^2}$ puis
    $\lim_{n\to\infty} \sum_{k=1}^n f\Bigl(\frac k{n^2}\Bigr)$ où $f$ est une
    fonction de classe $\mathcal{C}^2$ sur $\R$ vérifiant $f(0)=0$.
    

\finenonce{004026}



\finexercice


\exercice{4027, quercia, 2010/03/11}
\enonce{004027}{Recherche de tangentes}

Pour chacune des courbes suivantes,
déterminer la tangente pour $x=0$ et la position de la courbe par rapport à
cette tangente.

\begin{enumerate}
  \item $y = \frac{e^{\sin x}-1}x$.          
  \item $y = \frac1{\sh x} - \frac 1x$.     
  \item $y = \frac1{\arcsin x} - \frac 1x$. 
  \item $y = (2e^x - e^{-x})^{1/x}$.          
  \item $y = \frac 2{e^{2x}-1} - \frac1x$.  
  \item $y = \sqrt{1+\sqrt{1+x}}$.            
\end{enumerate}
\finenonce{004027}



\finexercice


\exercice{4028, quercia, 2010/03/11}
\enonce{004028}{Comparaison de fonctions}

On pose :
$f(x) = 1/(1+x)$,
$g(x) = e^{-x}$,
$h(x)  = \sqrt{1-2\sin x}$,
$k(x)  = \cos(\sqrt{2x})$.

Préciser les positions relatives de $\mathcal{C}_f$, $\mathcal{C}_g$, $\mathcal{C}_h$, $\mathcal{C}_k$ au voisinage
de $0$.

\finenonce{004028}



\finexercice


\exercice{4029, quercia, 2010/03/11}
\enonce{004029}{Recherche d'asymptotes}

Rechercher si les courbes suivantes admettent une asymptote en $+\infty$ et
déterminer la position s'il y a lieu :


\begin{enumerate}
  \item $y = \sqrt{x(x+1)}$.                               
  \item $y = \sqrt{\frac{x^3}{x-1}}$.                     
  \item $y = (x^2-1)\ln\left(\frac{x+1}{x-1}\right)$.     
  \item $y = (x+1)\arctan(1+2/x)$.                         
  \item $y = x.\arctan x.e^{1/x}$.                         
  \item $y = e^{2/x}\sqrt{1+x^2}\arctan x$.                
  \item $y = \sqrt{x^2-x}\exp\left(\frac 1{x+1}\right)$.  
\end{enumerate}
\finenonce{004029}




\section{125.04 Développements limités implicites}


\exercice{4037, quercia, 2010/03/11}
\video{jBtvOPwMUdQ}
\enonce{004037}{tan$(x) = x$}
\begin{enumerate}
  \item  Montrer que l'équation $\tan x = x$ possède une unique solution
   $x_n$ dans
   $\left]n\pi-\frac \pi2, n\pi+\frac \pi2\right[$ $(n\in \N)$.
  \item  Quelle relation lie $x_n$ et $\arctan(x_n)$ ? \label{relation}
  \item  Donner un DL de $x_n$ en fonction de $n$ à l'ordre $0$ pour $n\to\infty$.
  \item  En reportant dans la relation trouvée en \ref{relation},
     obtenir un DL de $x_n$ à l'ordre 2.
\end{enumerate}
\finenonce{004037}



\finexercice


\exercice{4038, quercia, 2010/03/11}
\enonce{004038}{maximum de $x\cos^nx$}
On note $f_n(x) = x\cos^nx$. Soit $x_n \in \left[0,\frac \pi2 \right]$ tel que
$f_n(x_n)$ soit maximal.

\begin{enumerate}
  \item   Existence et unicité de $x_n$ ?
      

  \item   Chercher $\lim_{n \to \infty} x_n$.
      

  \item   Montrer que $x_n^2 \sim \frac 1n$ ($n\to\infty$).


  \item   Trouver un équivalent de $f_n(x_n)$.
      


\end{enumerate}
\finenonce{004038}



\finexercice


\exercice{4039, quercia, 2010/03/11}
\enonce{004039}{Développement asymptotique}
Soit $f : x  \mapsto \frac{x+1}x e^x$.

\begin{enumerate}
  \item Tracer la courbe $\mathcal{C}$ représentative de $f$.
  \item Soit $\lambda \in \R^+$. Si $\lambda$ est assez grand, la droite d'équation
    $y=\lambda$ coupe $\mathcal{C}$ en deux points d'abscisses $a < b$.
 \begin{enumerate} 
    \item Montrer que $a \sim \frac1\lambda$, et
            $e^b \sim \lambda$ pour $\lambda \to +\infty$.
    \item Chercher la limite de $b^a$ quand $\lambda$ tend vers $+\infty$.
        
  \end{enumerate}
\end{enumerate}
\finenonce{004039}



\finexercice


\exercice{4040, quercia, 2010/03/11}
\enonce{004040}{Polytechnique MP$^*$ 2000}
Soit $f(x) = \frac{\ln|x-2|}{\strut\ln|x|}$. Montrer que pour tout~$n\in\N^*$,
il existe un unique $x_n$ vérifiant $f(x_n)=1-\frac1n$. Trouver la limite et un
équivalent de la suite~$(x_n)$ en $+\infty$.

\finenonce{004040}



\finexercice


\exercice{4041, quercia, 2010/03/11}
\enonce{004041}{} 
Soit $u_n$ une suite réelle telle que pour tout $n$ on ait 
$u_n^5+nu_n-1=0$. Trouver un développement asymptotique à deux termes de $u_n$.

\finenonce{004041}



\finexercice


\exercice{4042, quercia, 2010/03/11}
\enonce{004042}{Mines MP 2001}

Montrez que pour $n$ entier ($n>0$) l'équation $e^x=n-x$ admet une
unique solution positive $x_n$. Déterminer les trois premiers termes du
développement asymptotique de $x_n$ en fonction de~$n$.
\finenonce{004042}



\finexercice


\exercice{4043, quercia, 2010/03/11}
\enonce{004043}{Centrale MP 2001}

Pour tout $n$ entier naturel non nul, on donne $f_n(x) = nx^{n+1} - (n+1)x^n - \frac12$.

\begin{enumerate}
  \item Montrer que $f_n$ admet une unique racine positive notée $x_n$.
  \item Montrer que la suite $(x_n)$ converge vers une limite~$\ell$ et trouver
    un équivalent de $x_n-\ell$.
\end{enumerate}
\finenonce{004043}



\finexercice

\section{125.05 Equivalents}



\exercice{4044, quercia, 2010/03/11}

\video{QO_S3C9WRgc}

\enonce{004044}{Recherche d'équivalents}

Donner des équivalents simples pour les fonctions suivantes :

\begin{enumerate}
  \item $2e^x - \sqrt{1+4x} - \sqrt{1+6x^2}$,            en $0$      
  \item $(\cos x)^{\sin x} - (\cos x)^{\tan x}$,         en $0$      
  \item $\arctan x + \arctan \frac 3x -\frac {2\pi}3$, en $\sqrt3$ 
  \item $\sqrt{x^2+1} -2\sqrt[3]{x^3+x} + \sqrt[4]{x^4+x^2}$,  en $+\infty$
  \item $\Argch\left(\frac1{\cos x}\right)$,            en $0$      
\end{enumerate}
\finenonce{004044}



\finexercice



\exercice{4045, quercia, 2010/03/11}
\video{OchDDAXSqog}
\enonce{004045}{Approximation de cos}
Trouver $a,b\in\R$ tels que 
$$\cos x - \frac{1+ax^2}{1+bx^2}$$
soit un $o(x^n)$ en $0$ avec $n$ maximal.
\finenonce{004045}




\exercice{4046, quercia, 2010/03/11}
\enonce{004046}{Approximation de sin}
Trouver $a,b\in\R$ tels que
$\sin x - \frac{x+ax^3}{1+bx^2}$ soit infiniment petit d'ordre maximal.


\finenonce{004046}



\finexercice


\exercice{4047, quercia, 2010/03/11}
\enonce{004047}{Equivalent de $\arccos x$ en $1$}

Simplifier $\arccos(1-2x^2)$, en trouver un équivalent pour $x \to 0$, puis
donner un équivalent de $\arccos(u)$ pour~${u \to 1^-}$.

\finenonce{004047}



\finexercice


\exercice{4048, quercia, 2010/03/11}
\enonce{004048}{$\arcsin \circ \arctan - \arctan \circ \arcsin$}

\begin{enumerate}
  \item Soient $P(X) = X + aX^3 + bX^5 + cX^7$ et $Q(X) = X + \alpha X^3 + \beta X^5 + \gamma X^7$.
    Chercher la partie de degré inférieur ou égale à 7 de $P\circ Q - Q \circ P$.
    

  \item Application : Donner le DL à l'ordre 7 en 0 de
    $\arcsin(\arctan x) - \arctan(\arcsin x)$.
    


\end{enumerate}
\finenonce{004048}



\finexercice


\exercice{4049, quercia, 2010/03/11}
\enonce{004049}{$(u^v - v^u)/(u-v)$}

Soient $u,v$ deux fonctions positives, $u \sim v$, $u \to 0$.
Montrer que $\frac {u^v - v^u}{u-v} \sim -\ln(v)$.

(\'Ecrire $u = v + w$ avec $w/v \to 0$)

\finenonce{004049}



\finexercice


\exercice{4050, quercia, 2010/03/11}
\enonce{004050}{Développement asymptotique d'une réciproque}
Soit $f : {[-1,+\infty[} \to {[-e^{-1},+\infty[}, x \mapsto {xe^x.} $
Montrer que $f$ est bijective et $f^{-1}(x) \mathop{\sim}\limits_{+\infty} \ln(x)$.

\finenonce{004050}



\finexercice


\exercice{4051, quercia, 2010/03/11}
\enonce{004051}{\'Equivalent de $x^{x^{.^{.^{.^{\scriptstyle x}}}}}$}
Chercher un équivalent simple
en $0^+$ de $f_k(x) = x^{x^{.^{.^{.^{\scriptstyle x}}}}}$
($k$ fois $x$).



\finenonce{004051}



\finexercice


\exercice{4052, quercia, 2010/03/11}
\enonce{004052}{$\sum_{k=1}^n \frac1{k^{1-\alpha}}$}

Soit $\alpha \in {]0,1[}$.

\begin{enumerate}
  \item Montrer que : $\forall n \in \N^*,\
           \frac\alpha{(n+1)^{1-\alpha}} \le (n+1)^\alpha - n^\alpha \le
           \frac\alpha{n^{1-\alpha}}$.
  \item En déduire que $\sum_{k=1}^n \frac1{k^{1-\alpha}} \sim \frac{n^\alpha}{\alpha}$ pour $n\to\infty$.

\end{enumerate}
\finenonce{004052}



\finexercice


\exercice{4053, quercia, 2010/03/11}
\enonce{004053}{$(1+a_n/n)^n$}
\begin{enumerate}
  \item Soit $f(x) = \ln(1+x) - x$.
  \begin{enumerate}
  
    \item \'Etudier $f$.
    \item Chercher un équivalent simple de $f$ en $0$.
    \item Soit $(x_n)$ une suite de réels telle que $f(x_n) =  o(1/n)$.
        Montrer que $nx_n^2 \to 0$ lorsque $n\to\infty$.
  \end{enumerate}
  \item Application :
    Soit $(a_n)$ une suite de réels. Montrer que les deux propriétés suivantes sont
    équivalentes :
  \begin{enumerate}
    \item $a_n =  o(\sqrt n)$.
    \item $\left(1+\frac {a_n}n\right)^n \sim e^{a_n}$.
  \end{enumerate}
\end{enumerate}
\finenonce{004053}



\finexercice

\section{125.99 Autre}


\exercice{4030, quercia, 2010/03/11}
\enonce{004030}{DL de $(\ch x)^(1/x)$}

\begin{enumerate}
  \item Montrer que $\frac1x\ln(\ch x)$ admet en $+\infty$ un développement limité
    généralisé à tout ordre.
  \item En déduire le développement limité de $(\ch x)^{1/x}$ en $+\infty$ à
    un ordre $n$ quelconque.
    
\end{enumerate}
\finenonce{004030}



\finexercice


\exercice{4031, quercia, 2010/03/11}
\enonce{004031}{Théorème de division}
Soit $f : \R \to \R$ une fonction de classe $\mathcal{C}^n$.
On pose
$g(x) =\begin{cases}\frac {f(x)-f(0)}x &\text{ si } x \ne 0 \cr f'(0) &\text{ si } x=0. \cr\end{cases}$

\begin{enumerate}
  \item On suppose que $f(x) = o (x^n)$.
  \begin{enumerate}
    \item Démontrer que : $\forall\ p \le n,\ f^{(p)}(x) = o (x^{n-p})$,
        et : $\forall\ p < n,\ g^{(p)}(x) = o (x^{n-p-1})$.
    \item En déduire que $g$ est de classe $\mathcal{C}^{n-1}$ en $0$.
  \end{enumerate}
  \item Démontrer le même résultat dans le cas général.

  \item Soient ${f,g} : {\R} \to {\R}$ deux fonctions $\mathcal{C}^\infty$ telles
    que $f(0)=g(0)=0$ et $g'(0)\ne 0$. Montrer que $f/g$ se prolonge en une
    fonction $\mathcal{C}^\infty$ au voisinage de~$0$.
\end{enumerate}
\finenonce{004031}



\finexercice


\exercice{4032, quercia, 2010/03/11}
\enonce{004032}{DL de $f^{-1}$}
Soit $P \in \R[X]$ de valuation 1. Démontrer que pour tout entier $n \in \N$,
il existe deux polynômes $Q_n$ et $R_n$ uniques tels que :

$$\begin{cases}X = Q_n\circ P + R_n \cr \deg Q_n \le n < \text{v}(R_n).\cr\end{cases}$$

Application : Soit $f : \R \to \R$ bijective telle que
$f(x) = a_1x + a_2x^2 + \dots + a_nx^n + o (x^n)$, avec~$a_1 \ne 0$.
Démontrer que $f^{-1}$ admet un développement limité en $0$ à l'ordre $n$, et donner
les deux premiers termes.


\finenonce{004032}



\finexercice


\exercice{4033, quercia, 2010/03/11}
\enonce{004033}{DL de $(1-e^x)^n$}

Développer de deux manières $(1-e^x)^n$ en $0$ à l'ordre $n+2$.

En déduire $\sum_{k=0}^n(-1)^kC_n^kk^p$ pour $p = 0,1,\dots,n+2$.


\finenonce{004033}



\finexercice


\exercice{4034, quercia, 2010/03/11}
\enonce{004034}{Approximation de $f''$}


Soit $f : \R \to \R$ deux fois dérivable.
Chercher $\lim\limits_{h\to0} \frac {f(x-h) - 2f(x) + f(x+h)}{h^2}$.

\finenonce{004034}



\finexercice


\exercice{4035, quercia, 2010/03/11}
\enonce{004035}{Dérivation d'un DL d'ordre 2}

Soit $f : \R \to \R$ convexe dérivable telle que
$f(a+h) = f(a) + hf'(a) + o (h^2)$.

Démontrer que $f$ est deux fois dérivable en $a$ et $f''(a) = 0$
(comparer $f'(a+h)$ aux taux d'accroissement de $f$ entre $a$ et $a+h$,
et entre $a+h$ et $a+2h$).

\'Etudier le cas où $f(a+h) = f(a) + hf'(a) + \frac {Lh^2}2 + o (h^2)$.

\finenonce{004035}



\finexercice


\exercice{4036, quercia, 2010/03/11}
\enonce{004036}{$f(x+y)f(x-y) \le f^2(x)$}

Soit $f : \R \to \R$ de classe $\mathcal{C}^2$ telle que :
$\forall\ x,y \in \R,\ f(x+y)f(x-y) \le f^2(x)$.

Montrer que : $\forall\ t \in \R,\ f(t)f''(t) \le f'^2(t)$.
\finenonce{004036}



\finexercice

\section{126.01 Fonctions circulaires inverses}


\exercice{741, bodin, 1998/09/01}

\enonce{000741}{}
 \'Ecrire sous la forme $\frac{m}{n}\pi$ avec 
$m\in \Z$, $n\in\N^*$, $\vert m \vert$ et $n$ premiers entre eux,
$\arcsin(\sin\alpha)$, $\arccos(\cos\alpha)$ et $\arctan(\tan\alpha)$ dans les
cas: $\alpha=\frac{59}{5}\pi$; $\alpha=\frac{84}{5}\pi$; 
$\alpha=\frac{76}{5}\pi .$
\finenonce{000741}



\finexercice



\exercice{742, bodin, 1998/09/01}

\enonce{000742}{}
 R\'esoudre les \'equations suivantes :
\begin{enumerate}
    \item  $\arctan(2x)+\arctan x= \frac{\pi}{4}$.
    \item  $\arcsin(2x)-\arcsin(x\sqrt{3})=\arcsin(x)$.
\end{enumerate}
\finenonce{000742}



\finexercice



\exercice{743, gourio, 2001/09/01}

\enonce{000743}{}
R\'{e}soudre dans $\Rr$ l'\'{e}quation:
$$\arctan (x)+ \arctan (\sqrt{3}x)=\frac{7\pi }{12}.$$
\finenonce{000743}



\finexercice



\exercice{744, bodin, 1998/09/01}

\enonce{000744}{}
 Soient les fonctions $f:x\mapsto\arcsin(\sin x)$ et
$g:x\mapsto\arctan\sqrt{\frac{1-\cos x}{1+\cos x}}.$
\begin{enumerate}
    \item Simplifier les expressions de $f(x)$ et $g(x)$.
    \item Construire les graphes de $f$ et $g$.
\end{enumerate}
\finenonce{000744}



\finexercice



\exercice{745, bodin, 1998/09/01}
\video{_HL7BV1u578}
\enonce{000745}{}
 Une statue de hauteur $s$ est placée sur un 
piédestal de hauteur $p$. 
\begin{enumerate}
\item \`A quelle distance $x_0$ doit se placer un observateur
(dont la taille est supposée négligeable) pour voir la statue sous un 
angle maximal $\alpha_0$? 
\item Vérifier que $\alpha_0=\Arctan \frac{s}{2\sqrt{p(p+s)}}$.
\item Application à la statue de la liberté : haute de $46$ mètres avec un piédestal de 
$47$ mètres.
\end{enumerate}
\finenonce{000745}
 

\finexercice

\exercice{746, bodin, 1998/09/01}
\enonce{000746}{}
D\'emontrer les in\'egalit\'es suivantes :
$$
\Arcsin a < \frac{a}{\sqrt{1-a^2}} \quad \text{ si } 0<a<1 ;
$$
$$
\Arctan a > \frac{a}{1+a^2} \quad  \text{ si } a>0.
$$
\finenonce{000746} 


\finexercice

\exercice{747, bodin, 1998/09/01}
\video{elLpf1K7wu4}
\enonce{000747}{}
\'Ecrire sous forme d'expression algébrique
\begin{enumerate}
\item $ \sin(\Arccos x),\quad \cos(\Arcsin x),\quad \cos(2 \Arcsin x)$.
\item $ \sin(\Arctan x),\quad \cos(\Arctan x),\quad \sin(3 \Arctan x)$.
\end{enumerate}
\finenonce{000747} 


\finexercice

\exercice{748, bodin, 1998/09/01}

\enonce{000748}{}
Tracer les courbes repr\'esentatives des fonctions 
$$
x \mapsto f(x) = \sin(\Arcsin x),\qquad x \mapsto f(x) = \Arcsin(\sin x).
$$
\finenonce{000748}



\finexercice



\exercice{749, bodin, 1998/09/01}
\video{s66EQAvtk84}
\enonce{000749}{}
Résoudre les équations suivantes:
\begin{enumerate}
\item $\Arccos x = 2\Arccos \frac{3}{4}$.
\item $\Arcsin x = \Arcsin \frac{2}{5} + \Arcsin \frac{3}{5}$.
\item $\Arctan {2x}+\Arctan x=\frac{\pi}{4}$.
\end{enumerate}
\finenonce{000749} 


\finexercice

\exercice{750, bodin, 1998/09/01}

\enonce{000750}{}
Calculer
$$
\Arctan \frac{1}{2}+\Arctan \frac{1}{5}+\Arctan \frac{1}{8}.
$$
\finenonce{000750}



\finexercice



\exercice{751, bodin, 1998/09/01}

\enonce{000751}{}
Simplifier les expressions suivantes :
$$
\Arctan (\tan x) \quad (-\frac{\pi}{2}<x<\frac{3\pi}{2}),\quad
\Arctan \sqrt{\frac{1-\cos x}{1+\cos x}} \quad (0<x<2\pi),
$$
$$
\Arctan \frac{\sqrt{1-x^2}}{x}.
$$
\finenonce{000751}



\finexercice



\exercice{752, bodin, 1998/09/01}
\video{AjgAXbaAu3E}
\enonce{000752}{}
Vérifier
$$
\Arcsin x + \Arccos x = \frac{\pi}{2}\qquad\text{ et } \quad
\Arctan x + \Arctan\frac{1}{x} = \text{sgn}(x)\frac{\pi}{2}.
$$

\finenonce{000752} 

\finexercice

\exercice{753, ridde, 1999/11/01}

\enonce{000753}{}
Montrer que $\dfrac{\pi}4  = 4\arctan (\dfrac 15) - \arctan (\dfrac 1{239})$
 (on montrera que $0 \leq \arctan (\dfrac 15) \leq \dfrac{\pi}8$
 et $0 \leq \arctan (\dfrac 1{239}) \leq \dfrac{\pi}2$).
\finenonce{000753}



\finexercice



\exercice{754, gourio, 2001/09/01}

\enonce{000754}{}
\'Etudier la suite $(u_{n})_{n\in \Nn}$ d\'{e}finie par :
$$\forall n\in \Nn,u_{n}=\sum\limits_{k=1}^{n}\arctan \frac{1}{k^{2}-k+1}.$$
On montrera qu'elle converge (vers $\ell$) et on \'{e}valuera $\lim_{n\rightarrow \infty }n(u_{n}-\ell)$.

\emph{Indication} : que vaut $\arctan a-\arctan b$ ?
\finenonce{000754}



\finexercice



\exercice{755, gourio, 2001/09/01}

\enonce{000755}{}
\'Etudier la fonction : 
$$\phi :x\rightarrow \arcsin \frac{1-x^{2}}{1+x^{2}}+\arccos \frac{2x}{1+x^{2}}.$$
\finenonce{000755}



\finexercice



\exercice{756, gourio, 2001/09/01}

\enonce{000756}{}
R\'{e}soudre dans $\Rr$ l'\'{e}quation d'inconnue $x$ :
$$\arctan (x-1)+\arctan (x)+\arctan (x+1)=\frac{\pi }{2}.$$
\finenonce{000756}



\finexercice



\exercice{3914, quercia, 2010/03/11}
\enonce{003914}{arcsin et arccos à partir de arctan}

Le langage ``Pascal'' ne dispose pas des fonctions $\arcsin$ et $\arccos$.
Définir $\arcsin x$ et $\arccos x$ à l'aide de la fonction $\arctan$.
\finenonce{003914}



\finexercice


\exercice{3915, quercia, 2010/03/11}
\enonce{003915}{Formules d'addition}

Soient $a,b \in \R$. Simplifier $\arctan a + \arctan b$.

\finenonce{003915}



\finexercice


\exercice{3916, quercia, 2010/03/11}
\enonce{003916}{$\arcsin x = \arccos \frac13 - \arccos \frac14$}

Résoudre l'équation :
$\arcsin x = \arccos \frac13 - \arccos \frac14$.
\finenonce{003916}



\finexercice


\exercice{3917, quercia, 2010/03/11}
\enonce{003917}{$\arcsin \frac{1+\sqrt5}4$}
Soit $x = \arcsin \frac{1+\sqrt5}4$. Calculer $\cos 4x$ et en déduire $x$.



\finenonce{003917}



\finexercice


\exercice{3918, quercia, 2010/03/11}
\enonce{003918}{arctangentes}

\begin{enumerate}
  \item Simplifier $\arctan \frac {1-x}{1+x}$.
  \item Simplifier $\arctan \sqrt{\frac {1-x}{1+x}}$.
  \item Simplifier $\arctan\frac{x-\sqrt{1-x^2}}{x+\sqrt{1-x^2}}$.
  \item Simplifier $\arctan\frac{\sqrt{x^2+1}-1}x + \arctan\bigl(\sqrt{1+x^2} - x\bigr)$.
  \item Simplifier $\arctan \frac1{2x^2} - \arctan\frac x{x-1} + \arctan\frac{x+1}x$.
\end{enumerate}
\finenonce{003918}



\finexercice


\exercice{3919, quercia, 2010/03/11}
\enonce{003919}{$2\arcsin x + \arcsin f(x) = \frac \pi6$}

Existe-t-il une fonction $f : D \to \R$ telle que :
$\forall\ x \in D,\ 2\arcsin x + \arcsin f(x) = \frac \pi6$ ?
\finenonce{003919}



\finexercice


\exercice{3920, quercia, 2010/03/11}
\enonce{003920}{$\cos(3\arctan x)$ et $\cos^2\left(\frac 12\arctan x\right)$.}

Simplifier $\cos(3\arctan x)$ et $\cos^2\left(\frac 12\arctan x\right)$.
\finenonce{003920}



\finexercice


\exercice{3921, quercia, 2010/03/11}
\enonce{003921}{$\arccos(\cos x) - \frac 12\arccos(\cos 2x)$}

Simplifier $\arccos(\cos x) - \frac 12\arccos(\cos 2x)$ pour $x \in {[0,2\pi]}$.

\finenonce{003921}



\finexercice


\exercice{3922, quercia, 2010/03/11}
\enonce{003922}{\'Equation}

Résoudre :
$2\arccos\left(\frac {1-x^2}{1+x^2}\right) + \arcsin\left(\frac {2x}{1+x^2}\right)
 - \arctan\left(\frac {2x}{1-x^2}\right) = \frac {2\pi}3$.

\finenonce{003922}



\finexercice


\exercice{3923, quercia, 2010/03/11}
\enonce{003923}{$\frac x2 - \arcsin\sqrt{\frac{1+\sin x}2}$}

Simplifier
$\frac x2 - \arcsin\sqrt{\frac{1+\sin x}2}$ pour $x \in {[-\pi,\pi]}$.

\finenonce{003923}



\finexercice


\exercice{3924, quercia, 2010/03/11}
\enonce{003924}{$\cos(\arctan(\sin(\arctan \frac1x)))$}

Simplifier $\cos(\arctan(\sin(\arctan \frac1x)))$.

\finenonce{003924}



\finexercice


\exercice{3925, quercia, 2010/03/11}
\enonce{003925}{\'Equations aux arctan}

Résoudre :

\begin{enumerate}
  \item $\arctan 2x + \arctan 3x = \frac \pi4$.
    
  \item $\arctan\left(\frac{x-1}{x-2}\right) + \arctan\left(\frac{x+1}{x+2}\right) = \frac\pi4$.
    
  \item $\arctan\left(\frac1x\right) + \arctan\left(\frac{x-1}{x+1}\right) = \frac\pi4$.
    
  \item $\arctan(x-3) + \arctan(x) + \arctan(x+3) = \frac {5\pi}4$.
    
\end{enumerate}
\finenonce{003925}



\finexercice


\exercice{3926, quercia, 2010/03/11}
\enonce{003926}{Sommes remarquables}

\begin{enumerate}
  \item Montrer que : $4\arctan\frac 15 - \arctan \frac 1{239} = \frac \pi4$.
  \item Montrer que : $\arcsin\frac 45 + \arcsin \frac 5{13} + \arcsin \frac {16}{65} = \frac \pi2$.
\end{enumerate}
\finenonce{003926}



\finexercice


\exercice{3927, quercia, 2010/03/11}
\enonce{003927}{Sommes remarquables}
\begin{enumerate}
  \item Montrer que : $\forall\ x \in \R,\ \arctan x + 2\arctan\bigl(\sqrt{1+x^2} - x\bigr) = \frac \pi2$.
  \item Montrer que : $\forall\ x \in {]0,1]},\ 2\arctan\sqrt{\frac {1-x}x} + \arcsin(2x-1) = \frac \pi2$.
\end{enumerate}
\finenonce{003927}



\finexercice


\exercice{3928, quercia, 2010/03/11}
\enonce{003928}{$\arctan((x-\sin a)/\cos a)$}

Soit $a \in \left[0,\frac \pi2\right[$. On pose
$f(x) = \arcsin \left( \frac {2(x-\sin a)\cos a}{x^2-2x\sin a + 1} \right)$
et $g(x) = \arctan \left( \frac{x-\sin a}{\cos a} \right)$.

Vérifier que $f$ est bien définie,
calculer $\sin\bigl(2g(x)\bigr)$ et comparer $f(x)$ et $g(x)$.
\finenonce{003928}



\finexercice


\exercice{3929, quercia, 2010/03/11}
\enonce{003929}{Polynômes de Chebicheff}

Pour $n \in \N$, on pose $f_n(x) = \cos(n\arccos x )$ et
$g_n(x) = \frac {\sin( n\arccos x )}{\sqrt{1-x^2}}$.
Montrer que $f_n$ et $g_n$ sont des fonctions polynomiales.
\finenonce{003929}



\finexercice


\exercice{3930, quercia, 2010/03/11}
\enonce{003930}{\'Equivalent de arccos$(1-x)$}

A l'aide d'un changement de variable judicieux, trouver
$\lim_{x\to0^+} \frac {\arccos(1-x)}{\sqrt x}$.
\finenonce{003930}



\finexercice


\exercice{3931, quercia, 2010/03/11}
\enonce{003931}{Matexo}

Montrer que~:
$\forall\ x\in{]-1,1[},\ |\arcsin(x)| \le \frac{|x|}{\sqrt{1-x^2}}$.
\finenonce{003931}



\finexercice


\exercice{5084, rouget, 2010/06/30}
\enonce{005084}{***IT}
Domaine de définition et calcul des fonctions suivantes :
\begin{enumerate}
 \item  $x\mapsto\sin(\Arcsin x)$,  \item  $x\mapsto\Arcsin(\sin x)$,
 \item  $x\mapsto\cos(\Arccos x)$,  \item  $x\mapsto\Arccos(\cos x)$,  \item  $x\mapsto\tan(\Arctan x)$,  \item  $x\mapsto\Arctan(\tan
x)$.
\end{enumerate}
\finenonce{005084}


\finexercice


\exercice{5085, rouget, 2010/06/30}
\enonce{005085}{***IT}
\label{exo:suprou2}
\begin{enumerate}
 \item  Calculer $\Arccos x+\Arcsin x$ pour $x$ élément de $[-1,1]$.
 \item  Calculer $\Arctan x+\Arctan\frac{1}{x}$ pour $x$ réel non nul.
 \item  Calculer $\cos(\Arctan a)$ et $\sin(\Arctan a)$ pour $a$ réel donné.
 \item  Calculer, pour $a$ et $b$ réels tels que $ab\neq1$, $\Arctan a+\Arctan b$ en fonction de $\Arctan\frac{a+b}{1-ab}$ (on étudiera d'abord
$\cos(\Arctan a+\Arctan b)$ et on distinguera les cas $ab<1$, $ab>1$ et $a>0$, $ab>1$ et $a<0$).
\end{enumerate}
\finenonce{005085}


\finexercice

\exercice{5087, rouget, 2010/06/30}
\enonce{005087}{***I}
Existence et calcul de $\int_{0}^{\sin^2x}\Arcsin\sqrt{t}\;dt+\int_{0}^{\cos^2x}\Arccos\sqrt{t}\;dt$.
\finenonce{005087}


\finexercice


\exercice{5088, rouget, 2010/06/30}
\enonce{005088}{**}
Simplifier les expressions suivantes :
\begin{enumerate}
\item  $f_1(x)=\Arcsin\left(\frac{x}{\sqrt{1+x^2}}\right)$.
\item  $f_2(x)=\Arccos\left(\frac{1-x^2}{1+x^2}\right)$.
\item  $f_3(x)=\Arcsin\sqrt{1-x^2}-\Arctan\left(\sqrt{\frac{1-x}{1+x}}\right)$.
\item  $f_4(x)=\Arctan\frac{1}{2x^2}-\Arctan\frac{x}{x+1}+\Arctan\frac{x-1}{x}$.
\end{enumerate}
\finenonce{005088}


\finexercice


\exercice{5089, rouget, 2010/06/30}
\enonce{005089}{**I}
Calculer $\Arctan\frac{1}{2}+\Arctan\frac{1}{5}+\Arctan\frac{1}{8}$.
\finenonce{005089}


\finexercice


\exercice{5090, rouget, 2010/06/30}
\enonce{005090}{***I}
Calculer $u_n=\Arctan\frac{2}{1^2}+\Arctan\frac{2}{2^2}+...+\Arctan\frac{2}{n^2}$ pour $n$ entier naturel non nul
donné puis déterminer $\lim_{n\rightarrow +\infty}u_n$. (Utiliser l'exercice \ref{exo:suprou2} 4))
\finenonce{005090}


\finexercice


\exercice{5092, rouget, 2010/06/30}
\enonce{005092}{** Mines de DOUAI 1984}
On considère la fonction numérique $f$ telle que~:

$$f(x)=(x^2-1)\Arctan\frac{1}{2x-1},$$
et on appelle $(\mathcal{C})$ sa courbe représentative dans un repère orthonormé.

\begin{enumerate}
 \item  Quel est l'ensemble de définition $\mathcal{D}$ de $f$~?

 \item  Exprimer, sur $\mathcal{D}\setminus\{0\}$, la dérivée de $f$ sous la forme~:~$f'(x)=2xg(x)$.

 \item  Montrer que~:~$\forall x\in\Rr,\;2x^4-4x^3+9x^2-4x+1>0$ et en déduire le tableau de variation de $g$.

 \item  Dresser le tableau de variation de $f$.
\end{enumerate}
\finenonce{005092}


\finexercice


\exercice{5095, rouget, 2010/06/30}
\enonce{005095}{**}
Simplifier les expressions suivantes
\begin{enumerate}
 \item $\sin(2\Arcsin x)$,
 \item $\cos(2\Arccos x)$,
 \item $\sin^2\left(\frac{\Arccos x}{2}\right)$,
 \item $\ln(\sqrt{x^2+1}+x)+\ln(\sqrt{x^2+1}-x)$,
 \item $\Argsh\left(\frac{x^2-1}{2x}\right)$,
 \item $\Argch(2x^2-1)$,
 \item $\Argth\left(\sqrt{\frac{\ch x-1}{\ch x+1}}\right)$,
 \item $\frac{\ch(\ln x)+\sh(\ln x)}{x}$.
\end{enumerate}
\finenonce{005095}


\finexercice


\exercice{5096, rouget, 2010/06/30}
\enonce{005096}{**}
Résoudre dans $\Rr$ les équations suivantes~:
\begin{enumerate}
 \item $\ch x=2$,
 \item $\Arcsin(2x)=\Arcsin x+\Arcsin(x\sqrt{2})$,
 \item $2\Arcsin x=\Arcsin(2x\sqrt{1-x^2})$.
\end{enumerate}
\finenonce{005096}


\finexercice

\exercice{5315, rouget, 2010/07/04}
\enonce{005315}{***}
Montrer que $\sum_{k=0}^{n-1}\cotan^2(\frac{\pi}{2n}+\frac{k\pi}{n})= n(n-1)$. (Indication. Poser $x_k=\cotan^2(\frac{\pi}{2n}+\frac{k\pi}{n})$ puis trouver un polynôme dont les $x_k$ sont les racines.)
\finenonce{005315}


\finexercice


\exercice{5852, rouget, 2010/10/16}
\enonce{005852}{*** I}
Donner un développement à la précision $ \frac{1}{n^2}$ de la $n$-ième racine positive $x_n$ de l'équation $\tan x = x$.
\finenonce{005852}


\finexercice


\exercice{5853, rouget, 2010/10/16}
\enonce{005853}{*** I}
Soit $z$ un nombre complexe. Déterminer $\lim_{n \rightarrow +\infty}\left(1+ \frac{z}{n}\right)^n$.
\finenonce{005853}


\finexercice


\exercice{6973, blanc-centi, 2014/05/06}
\video{70pEydvJEw4}
\enonce{006973}{}
Montrer que pour tout $x>0$, on a
$$\Arctan\left(\frac{1}{2x^2}\right)=\Arctan\left(\frac{x}{x+1}\right)-\Arctan\left(\frac{x-1}{x}\right).$$
En déduire une expression de $\displaystyle S_n=\sum_{k=1}^n\Arctan\left(\frac{1}{2k^2}\right)$ 
et calculer $\displaystyle\lim_{n\to +\infty}S_n$.
\finenonce{006973} 


\finexercice


\exercice{6974, blanc-centi, 2014/05/06}
\video{pT2xX0g3eDA}
\enonce{006974}{}
Soit $z=x+iy$ un nombre complexe, où $x=\Re z$ et $y=\Im z$. 
On sait que si $z$ est non nul, on peut l'écrire de façon unique sous la forme 
$z=x+iy=re^{i\theta}$, où $\theta\in]-\pi,\pi]$ et $r=\sqrt{x^2+y^2}$. 
\begin{center}
\begin{tikzpicture}[scale=1]
      \draw[->,>=latex, gray] (-0.5,0)--(4.5,0);
       \draw[->,>=latex, gray] (0,-0.5)--(0,3);

      \coordinate (A) at (3,2);
      \coordinate (Ax) at (3,0);
      \coordinate (Ay) at (0,2);

       \draw[thin] (0,0)--(A) node[midway,above] {$r$};
       \draw[dashed,gray] (A)--(Ax);
       \draw[dashed,gray] (A)--(Ay);

       \fill (A) circle (2pt);
       \fill (0,0) circle (2pt);

       \node at (0,0) [below left] {$0$}; 
       \node at (A) [above right] {$z=x+iy$}; 
       \node at (Ax) [below] {$x$}; 
       \node at (Ay) [left] {$y$};
       
       \draw[->,>=latex] (0:1.6) arc (0:33.69:1.6);
       \node[right] at (16.5:1.6) {$\theta$};
\end{tikzpicture}  
\end{center}

\begin{enumerate}
\item Montrer que si $x>0$, alors $\theta=\Arctan\frac{y}{x}$.
\item Montrer que si $\theta\in]-\pi,\pi[$, alors 
$\theta=2\Arctan\left(\frac{\sin\theta}{1+\cos\theta}\right)$. 
\item En déduire que si $z$ n'est pas réel négatif ou nul, on a l'égalité
$$\theta=2\Arctan\left(\frac{y}{x+\sqrt{x^2+y^2}}\right).$$
\end{enumerate}
\finenonce{006974} 


\finexercice
\section{126.02 Fonctions hyperboliques et hyperboliques inverses}


\exercice{757, gourio, 2001/09/01}

\enonce{000757}{}
Soit $f:\Rr^{2}\rightarrow \Rr^{2}$ d\'{e}finie par:
$$\forall (x,y)\in \Rr^{2},f(x,y)=(\cos x+\ch y,\cos x\ch y).$$
Discuter et d\'{e}terminer selon $p\in \Rr$ l'image r\'{e}ciproque de $(4,p)$. 
On exprimera $y$ \`{a} l'aide d'un logarithme.
D\'{e}terminer num\'{e}riquement cette image r\'{e}ciproque si $p=-2$.
\finenonce{000757}



\finexercice



\exercice{758, gourio, 2001/09/01}
\enonce{000758}{}
\begin{enumerate}
  \item Montrer qu'il n'existe pas de fonction $f:[1;+\infty [\rightarrow \Rr$
v\'{e}rifiant:
$$ \forall x\in \Rr, \quad f(\ch x)=e^{x}. $$
  \item D\'{e}terminer toutes les fonctions $f:\Rr^{+*}\rightarrow \Rr$ telles que:
$$ \forall x\in \Rr, \quad f(e^{x})=\ch x.$$
Pr\'{e}ciser le nombre de solutions.
  \item D\'{e}terminer toutes les fonctions $f:\Rr^{+}\rightarrow \Rr$ telles que:
$$\forall x\in \Rr, \quad f(e^{x})=\ch x.$$
Pr\'{e}ciser le nombre de solutions ; y a t-il des solutions continues sur $\Rr^{+}$ ?
 \end{enumerate}
\finenonce{000758} 


\finexercice

\exercice{759, gourio, 2001/09/01}
\enonce{000759}{}
Calculer :
$$\lim_{x\rightarrow +\infty }e^{-x}(\ch ^{3}x-\sh ^{3}x) \quad \text{ et }\quad
\lim\limits_{x\rightarrow +\infty }(x-\ln (\ch x)).$$
\finenonce{000759} 


\finexercice


\exercice{760, bodin, 1998/09/01}

\enonce{000760}{}
 Donner un expression plus simple de :
$$y=\mathrm{argch}\sqrt{\frac{1+\mathrm{ch}\ x}{2}};\ \ \
y=\mathrm{argsh}(2x\sqrt{1+x^2});\ \ y=\mathrm{argth}\frac{x^2-1}{x^2+1}.$$
\finenonce{000760}



\finexercice



\exercice{761, gourio, 2001/09/01}

\enonce{000761}{}
Calculer pour $(n,a,b)\in \Nn^{*}\times \Rr^{2}$ :
$$\sum\limits_{k=0}^{n-1}\ch (a+bk),\quad \sum\limits_{k=0}^{n-1}\sh (a+bk).$$
\finenonce{000761}



\finexercice



\exercice{762, bodin, 1998/09/01}

\enonce{000762}{}
 Soit $(a,b)\in\R^2$, r\'esoudre le syst\`eme 
$\left\lbrace
\begin{array}{c}
 \mathrm{ch}x+\mathrm{sh}y=a  \\
 \mathrm{sh}x+\mathrm{ch}y=b
\end{array}\right.$.
\finenonce{000762}



\finexercice





\exercice{764, bodin, 1998/09/01}
\video{eW51oHrCmVU}
\enonce{000764}{} 
Soit $x\in\R$. On pose $t=\Arctan(\sh x)$.
\begin{enumerate}
\item \'Etablir les relations  
$$\tan t=\sh x \qquad\qquad \frac{1}{\cos t}=\ch x \qquad\qquad \sin t=\tanh x$$
\item Montrer que $x = \ln \big(\tan\big(\frac{t}{2}+\frac{\pi}{4}\big)\big)$.
\end{enumerate}
\finenonce{000764} 


\finexercice

\exercice{765, bodin, 1998/09/01}

\enonce{000765}{}
Montrer que $\ch nx$ et $\sh nx$ peuvent s'exprimer
comme polynômes en $\ch x$ et $\sh x$.
Calculer $\ch 3x$ et $\sh 3x$ en fonctions de $\ch x$ et $\sh x$.
En d\'eduire $\tanh 3x$ en fonction de $\tanh x$.
\finenonce{000765}



\finexercice



\exercice{766, bodin, 1998/09/01}

\enonce{000766}{}
Exprimer $\ch^n x$ et $\sh^n x$ au moyen de
$\{ \sh px, \ch px \ ; \ 1 \le p \le n \}$.
Expliciter $\ch^5x$ et $\sh^5x$.
\finenonce{000766}



\finexercice



\exercice{767, bodin, 1998/09/01}

\enonce{000767}{}
Calculer les sommes
$$1+\ch x + \ch 2x + \cdots + \ch nx \quad \text{et} \quad
1+\sh x + \sh 2x + \cdots + \sh nx.$$
\finenonce{000767}



\finexercice



\exercice{768, bodin, 1998/09/01}

\enonce{000768}{}
Simplifier
$$
\Argth \frac{x^2-1}{x^2+1}.
$$
\finenonce{000768}



\finexercice



\exercice{769, bodin, 1998/09/01}

\enonce{000769}{}
V\'erifier les \'egalit\'es 
$$
2\Argth \tan x = \Argth \sin 2x,\qquad 
\Argsh (3x+4x^3) = 3 \Argsh x.
$$
\finenonce{000769}



\finexercice



\exercice{770, bodin, 1998/09/01}

\enonce{000770}{}
Expliciter au moyen de la fonction logarithme $\Argch \frac{1}{x}$
et $\Argsh \frac{1}{x}$.
\finenonce{000770}



\finexercice



\exercice{771, bodin, 1998/09/01}

\enonce{000771}{}
R\'esoudre
$$
x^{\sqrt{x}} = \sqrt{x}^{x} ;
$$
$$
xy = a^2 \text{ et } \ln^2 x+\ln^2 y = \frac{5}{2} \ln^2 a.
$$
\finenonce{000771}



\finexercice



\exercice{772, bodin, 1998/09/01}

\enonce{000772}{}
Pr\'eciser les comportements \\
de $\displaystyle{x\mapsto \frac{x^2-e^x}{x-e}}$ quand $x \rightarrow e$,\\
de $\displaystyle{x\mapsto \sqrt{\ln(1+x)}-\sqrt{\ln x} }$ quand $x \rightarrow +\infty$,\\
de $\displaystyle{x\mapsto \frac{a^x-b^x}{x}}$ quand $x \rightarrow 0$.
\finenonce{000772}



\finexercice



\exercice{773, bodin, 1998/09/01}

\enonce{000773}{}
D\'emontrer les in\'egalit\'es : 
$$ x- \frac{x^2}{2} < \ln (1+x) \text{ pour } x>0 \quad \text{et} \quad
1+x \le e^x \text{ pour tout $x$ r\'eel}.$$

\finenonce{000773}



\finexercice



\exercice{774, ridde, 1999/11/01}

\enonce{000774}{}
D\'eterminer $\lim\limits_{ + \infty} (x-\ln (\text{ch}x))$.
\finenonce{000774}



\finexercice



\exercice{775, ridde, 1999/11/01}

\enonce{000775}{}
Montrer que $\forall x \in \Rr \, \, \text{ch} (2x)  = 1 + 2\text{sh}^2 x$.
En d\'eduire un \'equivalent de $\text{ch}x -1$ en $0$.
\finenonce{000775}



\finexercice



\exercice{776, ridde, 1999/11/01}
\video{_TunkNZcaeI}
\enonce{000776}{}
R\'esoudre l'\'equation $x^y = y^x$ o\`u $x$ et $y$ sont des entiers positifs non nuls.
\finenonce{000776} 


\finexercice

\exercice{777, ridde, 1999/11/01}

\enonce{000777}{}
R\'esoudre l'\'equation $\tan (3 \arcsin x) = 1$. On exprimera les trois solutions
au moyen de radicaux.
\finenonce{000777}



\finexercice



\exercice{5086, rouget, 2010/06/30}
\enonce{005086}{*IT}
\label{exo:suprou3}
Etablir pour ch, sh et th les formules d'addition, de duplication et de linéarisation.
\finenonce{005086}


\finexercice


\exercice{5091, rouget, 2010/06/30}
\enonce{005091}{*}
Etudier $f~:~x\mapsto\ln(\ch x)-x$.
\finenonce{005091}


\finexercice


\exercice{5093, rouget, 2010/06/30}
\enonce{005093}{**}
Résoudre dans $\Rr$ l'équation $\sh(2+x)+\sh(2+2x)+...+\sh(2+100x)=0$.
\finenonce{005093}


\finexercice


\exercice{5094, rouget, 2010/06/30}
\enonce{005094}{**I}
\begin{enumerate}
 \item  Montrer que pour tout réel $x$ non nul, on a~:~$\tanh x=\frac{2}{\tanh(2x)}-\frac{1}{\tanh x}$.

 \item  En déduire la valeur de $u_n=2^0\tanh(2^0x)+2^1\tanh(2^1x)+ \cdots +2^{n}\tanh(2^{n}x)$ pour $n$ entier naturel non nul et $x$ réel
non nul donnés puis calculer la limite de $(u_n)$.
\end{enumerate}
\finenonce{005094}


\finexercice


\exercice{6975, blanc-centi, 2014/05/06}
\video{g13-StbHolQ}
\enonce{006975}{} 
Simplifier l'expression $\displaystyle\frac{2\ch^2(x)-\sh(2x)}{x-\ln(\ch x)-\ln 2}$ 
et donner ses limites en $-\infty$ et $+\infty$.
\finenonce{006975} 


\finexercice


\exercice{6976, blanc-centi, 2014/05/06}
\video{pMuDpGtaNAM}
\enonce{006976}{} 
Soit $x$ un réel fixé. Pour $n\in\Nn^*$, on pose
$$C_n=\sum_{k=1}^n\ch(kx)\qquad\text{ et }\qquad S_n=\sum_{k=1}^n\sh(kx).$$
Calculer $C_n$ et $S_n$.
\finenonce{006976} 


\finexercice


\exercice{6977, blanc-centi, 2014/05/06}
\video{GEaDEYtHv7U}
\enonce{006977}{} 
Soit $a$ et $b$ deux réels positifs tels que $a^2-b^2=1$. Résoudre le système
$$\left\{\begin{array}{l}
\ch(x)+\ch(y)=2a\\
\sh(x)+\sh(y)=2b
\end{array}\right.$$
\finenonce{006977} 


\finexercice


\exercice{6978, blanc-centi, 2014/05/06}
\video{UkRPGUQfPP8}
\enonce{006978}{} 
Simplifier les expressions suivantes:
\begin{enumerate}
\item $\ch(\Argsh x),\quad \tanh(\Argsh x),\quad \sh(2\Argsh x)$.
\item $ \sh(\Argch x),\quad \tanh(\Argch x),\quad \ch(3\Argch x)$.
\end{enumerate}
\finenonce{006978} 


\finexercice


\exercice{6979, blanc-centi, 2014/05/06}
\video{bRKLAqnlbBM}
\enonce{006979}{} 
\'Etudier le domaine de définition de la fonction $f$ définie par 
$$f(x)=\Argch\left[\frac{1}{2}\left(x+\frac{1}{x}\right)\right]$$
et simplifier son expression lorsqu'elle a un sens.
\finenonce{006979} 


\finexercice


\exercice{6980, blanc-centi, 2014/05/06}
\video{QkvgaXwOwyU}
\enonce{006980}{} 
Montrer que l'équation $\Argsh x+\Argch x=1$ admet une unique solution, puis la déterminer.
\finenonce{006980} 


\finexercice

\section{126.99 Autre}

\section{127.01 Théorie}


\exercice{778, ridde, 1999/11/01}

\enonce{000778}{}
D\'eterminer les fonctions $f$ de $[a, b]$ dans $\Rr$ telles que
$\int_{a}^{b}{f (t)dt} = (b-a) \sup\limits_{[a, b]}\left|f\right|$.
\finenonce{000778}



\finexercice



\exercice{779, ridde, 1999/11/01}

\enonce{000779}{}
Soient $f \in C^{1} ([a, b], \Rr)$ et $I_{n} = \int_{a}^{b}{f (t)\sin (nt)dt}$.
\begin{enumerate}
\item A l'aide d'une int\'egration par parties, montrer que $I_{n} \rightarrow 0$.
\item Montrer que ceci est encore vrai si $f$ est en escalier.
\item En d\'eduire que le r\'esultat subsiste pour $f$ continue par morceaux.
\end{enumerate}
\finenonce{000779}



\finexercice



\exercice{780, ridde, 1999/11/01}

\enonce{000780}{}
Soient $0<a \leq b$. Montrer que $\int_{a}^{b}{\frac{dx}{x}} \leq \frac{b-a}{\sqrt{ab}}$.
\finenonce{000780}



\finexercice



\exercice{781, ridde, 1999/11/01}

\enonce{000781}{}
Soit $f \in C^{0} ([0, 1], \Rr)$ telle que $\int_{0}^{1}{f (t)dt} = \frac 12$.
Montrer qu'il existe $a\in ]0, 1[$ telle que $f (a) = a$.
\finenonce{000781}



\finexercice



\exercice{782, ridde, 1999/11/01}

\enonce{000782}{}
Soit $f \in C^0 (\Rr)$. On d\'efinit ${g} : {\Rr^*} \to {\Rr}$, ${x}\mapsto {\frac 1x \int_0^x f (t)dt}$.
\begin{enumerate}
\item Montrer que $g$ se prolonge par continuit\'e en $0$.
\item Montrer que si $f$ est p\'eriodique, $g$ admet une limite en $ + \infty$.
\end{enumerate}
\finenonce{000782}



\finexercice



\exercice{783, gourio, 2001/09/01}

\enonce{000783}{}
Soit $f$ continue de $[0,1]$ dans ${\Rr}, n\in \Nn$ tels que :
$$\forall k\in \{0,...,n\},\int_{0}^{1}f(u)u^{k}du=0. $$
Montrer que $f$ admet au moins $n+1$ z\'{e}ros distincts dans $]0,1[.$
\finenonce{000783}



\finexercice



\exercice{784, gourio, 2001/09/01}

\enonce{000784}{}
Soit $f:[0,1]\rightarrow {\Rr}$ une application continue strictement
croissante telle que :
$$f(0)=0,\  f(1)=1. $$
Calculer :
$$\lim\limits_{n\rightarrow \infty }\int_{0}^{1}f^{n}(t)dt. $$
\finenonce{000784}



\finexercice



\exercice{785, gourio, 2001/09/01}

\enonce{000785}{}
Soit $f:[0,1]\rightarrow {\Rr}$ une application continue, n'admettant
qu'un nombre fini de z\'{e}ros sur $[0,1]$, et telle que $f(0)=0, f(1)=1.$
Montrer que :
$$\lim\limits_{n\rightarrow \infty }\left| \int_{0}^{1}e^{nt}f(t)dt\right|
=+\infty . $$
\finenonce{000785}



\finexercice



\exercice{786, gourio, 2001/09/01}

\enonce{000786}{Irrationnalit\'{e} de $\pi $}
\begin{enumerate}
\item Soit $(a,b)\in (\Nn^{*})^{2}$, $n\in \Nn^{*},$ montrer que le
polyn\^{o}me $P_{n}=\frac{X^{n}(bX-a)^{n}}{n!}$ et ses d\'{e}riv\'{e}es
successives prennent, en $0$ et $\frac{a}{b}$, des valeurs enti\`{e}res.
\item Montrer que :
$$I_{n}=\int_{0}^{\pi }P_{n}(t)\sin (t)dt\rightarrow 0\text{ quand }
n\rightarrow \infty . $$
\item Montrer par l'absurde que $\pi \in {\Rr}\setminus\Qq.$
\end{enumerate}
\finenonce{000786}



\finexercice



\exercice{787, gourio, 2001/09/01}

\enonce{000787}{}
Soit $f$ continue sur $[0,\pi ]$ telle que $\int_{0}^{\pi }f(u)\cos
(u)du=\int_{0}^{\pi }f(u)\sin (u)du=0,$ montrer que $f$ s'annulle au moins
deux fois sur $]0,\pi [.$
\finenonce{000787}



\finexercice



\exercice{788, gourio, 2001/09/01}

\enonce{000788}{}
Soit $f\in C([0,1],{\Rr})$ telle que :
$$\forall g\in E\left( [0,1],{\Rr}\right) ,\int_{0}^{1}fg=0. $$
Montrer que $f=0.$
\finenonce{000788}



\finexercice



\exercice{789, gourio, 2001/09/01}

\enonce{000789}{}
Soit $f$ une fonction $C^{1}$ sur $[a,b]$ \`{a} valeurs dans ${\Rr}$.
On suppose $f(a)=0.$
Montrer que :
$$\int_{a}^{b}f^{2}(u)du\leq \frac{(b-a)^{2}}{2}\int_{a}^{b}f^{\prime 2}(u)du. $$
\finenonce{000789}



\finexercice



\exercice{790, gourio, 2001/09/01}

\enonce{000790}{}
Soit $f$ continue sur $[0,1]$ \`{a} valeurs dans $[a,b]$. On suppose $a<0<b$
et $\int_{0}^{1}f(t)dt=0.$
Montrer que :
$$\int_{0}^{1}f^{2}(t)dt\leq -ab. $$
\finenonce{000790}



\finexercice



\exercice{791, gourio, 2001/09/01}

\enonce{000791}{}
Soit $(a,b)\in {\Rr}^{2}$ ($a<b$), et $f$ continue positive de $[a,b]$
dans ${\Rr}$. Montrer que
$$\lim\limits_{n\rightarrow \infty }\left( \int_{a}^{b}f^{n}(t)dt\right) ^{%
\frac{1}{n}}=\sup\limits_{t\in [a,b]}\left| f(t)\right| . $$
\finenonce{000791}



\finexercice



\exercice{792, gourio, 2001/09/01}

\enonce{000792}{}
Calculer sans utiliser de primitive, pour $a<b$ :
$$\int_{a}^{b}e^{t}dt. $$
\finenonce{000792}



\finexercice



\exercice{793, gourio, 2001/09/01}

\enonce{000793}{}
Soit $f$ continue de $[0,1]$ dans ${\Rr}$ telle que $\int_{0}^{1}f^{n}(u)du$
ne prenne qu'un nombre fini de valeurs quand $n$ d\'{e}crit $\Nn.$
Montrer que $f=-1$ ou $f=0$ ou $f=1.$
\finenonce{000793}



\finexercice



\exercice{795, gourio, 2001/09/01}

\enonce{000795}{}
Calculer :
$$\lim\limits_{n\rightarrow \infty }\int_{0}^{1}\frac{e^{t}}{1+t^{n}}dt. $$
\'Eventuellemment, en donner un DL en $\frac{1}{n}$.
\finenonce{000795}



\finexercice



\exercice{796, gourio, 2001/09/01}

\enonce{000796}{}
Calculer :
$$\lim\limits_{n\rightarrow \infty }\int_{0}^{1}\frac{e^{-nx}}{1+x}dx. $$
Soit $f:[0,1]\rightarrow {\Rr}$ une application continue ; calculer :
$$\lim\limits_{n\rightarrow \infty }\int_{0}^{1}nx^{n}f(x)dx. $$
\finenonce{000796}



\finexercice



\exercice{797, gourio, 2001/09/01}

\enonce{000797}{}
Soit $f:[0,1]\rightarrow {\Rr}$ une application continue par morceaux,
continue en $0,$ trouver une suite $(g_{n})_{n\in \Nn}$ de fonctions en
escaliers telle que :
$$\lim_{n\rightarrow \infty }\int_{0}^{1}f(t)g_{n}(t)dt=f(0). $$
\finenonce{000797}



\finexercice



\exercice{798, cousquer, 2003/10/01}

\enonce{000798}{}
Dire (avec justification) si les affirmations
suivantes sont vraies ou fausses.

\begin{enumerate}

\item  Toute fonction intégrable sur $[a,b]$ est continue.

\item  Si $f$ est intégrable sur $[a,b]$, $\frac{d}{dx}
\int_a^xf(t)\,dt=f(x)$ pour tout $x$ de $[a,b]$.

\item  Soit $f$ une fonction sur $[a,b]$ vérifiant la propriété~: pour
tout $\epsilon>0$, il existe $g_\epsilon$ intégrable sur $[a,b]$
telle que $\forall x\in[a,b],\; \vert f(x)-g_\epsilon(x)\vert
\leq\epsilon$~; alors $f$ est intégrable.

\item  Si $f$ est intégrable sur $[a,b]$, alors $\vert f\vert$ est
intégrable sur $[a,b]$.

\item  Si $\vert f\vert$ est intégrable sur $[a,b]$, alors $f$ est
intégrable sur $[a,b]$.

\item  Si $f$ et $g$ sont des fonctions intégrables sur $[a,b]$, alors la
fonction $fg$ est intégrable sur $[a,b]$.

\item  Si $f$ et $g$ sont des fonctions continues sur $[a,b]$, alors la fonction
$fg$ est continue sur $[a,b]$, et $\int_a^bf(t)g(t)\,dt=
\int_a^bf(t)\,dt\cdot\int_a^bg(t)\,dt$.

\item  Soit $f$ la fonction définie sur $[0,1]$ par
$$\begin{cases}
f\equiv\lambda_n& \text{ sur } \mathopen]\frac{1}{2^n},\frac{1}{2^{n-1}}]
\text{ pour tout entier } n\geq1\cr
f(0)=\mu\cr \end{cases}$$
où $(\lambda_n)$ est une suite bornée de nombres réels,
et $\mu$ un nombre réel. Alors $f$ est intégrable.

\item  Soit $f$ bornée sur $[0,1]$, continue sauf au point $1/3$~;
alors $f$ est intégrable sur $[0,1]$.

\item  Il existe $f\geq 0$ continue sur $[0,1]$, avec $f(1/2)>0$, et telle que
$\int_0^1f(t)\,dt=0$.

\item  Soit $f$ intégrable sur $[a,b]$. Si $\int_a^bf(t)\,dt>0$ alors
$f\geq 0$ sur $[a,b]$.

\item  Si $f$ est croissante sur $[a,b]$, elle est intégrable sur $[a,b]$
et de plus $F(x)=\int_a^xf(t)\,dt$ est croissante.

\item  Si $f\leq 0$ est continue sur $[a,b]$, alors $G(x)=\int_x^bf(t)\,dt$
est croissante sur $[a,b]$.

\item  Si $f$ est continue sur $[0,1]$, $H(x)=\int_0^{x^2}f(t)\,dt$
est dérivable sur $[0,1]$, et $\forall x\in[0,1],\;\allowbreak H'(x)=f(x^2)$.

\end{enumerate}
\finenonce{000798}



\finexercice



\exercice{799, cousquer, 2003/10/01}

\enonce{000799}{}
 Soit $\varphi$ une fonction bornée sur $[a,b]$~; comparer les
assertions suivantes\footnote%
    {L'une des implications à étudier est très difficile~; on pourra 
    admettre après avoir traité toutes les autres que celle qui reste 
    est fausse.}~:
    
\begin{enumerate}
\item $\varphi$ a une primitive sur $[a,b]$.
\item $\varphi$ est intégrable sur $[a,b]$.
\item $\varphi$ est continue sur $[a,b]$.
\item $\varphi$ est dérivable sur $[a,b]$.
\end{enumerate}
\finenonce{000799}



\finexercice



\exercice{800, cousquer, 2003/10/01}

\enonce{000800}{}
 Soit $f$ une fonction continue et strictement croissante de $[a,b]$ sur
$[\alpha,\beta]$. On note $g$ la fonction réciproque de~$f$. Montrer que
$$\int_a^bf(x)\,dx+\int_\alpha^\beta g(x)\,dx = b\beta-a\alpha$$
\finenonce{000800}



\finexercice



\exercice{801, cousquer, 2003/10/01}

\enonce{000801}{}
Soit $f$ et $g$ deux fonctions intégrables sur $[a,b]$. On suppose
que $f$ est monotone sur $[a,b]$ et que $g$ est positive sur $[a,b]$.
Montrer qu'il existe $c\in[a,b]$ tel que
$$\int_a^bf(t)g(t)\,dt=f(a)\int_a^cg(t)\,dt+f(b)\int_c^bg(t)\,dt$$
(considérer $\varphi(x)=f(a)\int_a^xg(t)\,dt+f(b)\int_x^bg(t)\,dt$).
\finenonce{000801}



\finexercice



\exercice{802, cousquer, 2003/10/01}

\enonce{000802}{}
 Soit $f$ une fonction dérivable sur $[0,1]$, vérifiant~:
\begin{enumerate}
    \item [i)]  $0\leq f'\leq2$~;

    \item [ii)]  $f'$ est décroissante~;

    \item [iii)]  $f(0)=0$ et $f(1)=1$.
\end{enumerate}

Trouver le plus grand nombre $m$ et le plus petit nombre $M$ tels qu'on
soit sûr d'avoir $m\leq\int_0^1f(t)\,dt\leq M$.
Peut-il y avoir égalité~?
\finenonce{000802}



\finexercice



\exercice{803, cousquer, 2003/10/01}

\enonce{000803}{}
 Soit $f$ définie et continue sur $[0,+\infty\mathclose[$, vérifiant
$\lim_{x\to+\infty}f(x)=l$. Montrer que 
$\lim_{x\to+\infty}\frac{1}{x}\int_0^xf(t)\,dt=l$ (étant
donné $\epsilon>0$, choisir $A$ assez grand pour que sur $[A,+\infty\mathclose[$ 
on ait
$l-\epsilon\leq f(t)\leq l+\epsilon$~; puis encadrer 
$\frac{1}{x}\int_A^xf(t)\,dt$,
pour $x>A$~; estimer l'erreur\dots{} et faire un dessin~!).

Pour $x\geq 0$, on pose $F(x)=\int_0^x\sqrt{1+\frac{\sin^2 t}{1+t^2}}\,dt$.
Étudier la branche infinie du graphe de $F$ quand $x\to+\infty$.
\finenonce{000803}



\finexercice



\exercice{804, cousquer, 2003/10/01}

\enonce{000804}{Méthode des trapèzes}
\begin{enumerate}
\item
 Soit $f$ deux fois dérivable sur $[a,b]$, vérifiant
$\vert f''\vert\leq M$ sur $[a,b]$. Soit
$$\varphi(t)=f(t)-f(a)-(t-a)\frac{f(b)-f(a)}{b-a}-A(b-t)(t-a)$$
Soit $x\in\mathopen]a,b\mathclose[$~; on choisit $A=A(x)$ pour que 
$\varphi(x)=\nolinebreak 0$
(dessiner~!).
Montrer qu'il existe $c_1,c_2\in[a,b]$ tels que $c_1<c_2$ et
$\varphi'(c_1)=\varphi'(c_2)=0$, puis qu'il existe $c\in[a,b]$ tel que
$\varphi''(c)=0$. En déduire une majoration de $\vert A\vert$ pour 
$x\in[a,b]$. On convient de poser $A(a)=A(b)=0$.

\item
 On note $E$ l'erreur commise en remplaçant $\int_a^bf(x)\,dx$ par
l'aire du trapèze défini par l'axe des $x$, les droites $x=a$ et
$x=b$ et la corde du graphe joignant les points $(a,f(a))$ et $(b,f(b))$
(dessiner~!).
Montrer que $E=\int_a^bA(x)(b-x)(x-a)\,dx$, et vérifier que
l'intégrale a un sens. En déduire que
$\vert E\vert\leq\frac{M(b-a)^3}{12}$ (utiliser 1)).

\item
 Pour $n\geq1$ on pose $I_n=\frac{b-a}{n}\Bigl[\frac{f(a)}{2}+f(x_1)+
f(x_2)+\cdots+f(x_{n-1})+\frac{f(b)}{2}\Bigr]$ où $x_p=a+p\frac{b-a}{n}$ pour
$p=1,2,\ldots,n-1$. Montrer que $I_n$ est la somme des aires des trapèzes
construits sur les points d'abscisses $a,x_1,x_2,\ldots,x_{n-1},b$ et
les cordes correspondantes du graphe de $f$ (dessiner~!). Montrer que
$$\biggl\vert\int_a^b f(x)\,dx-In\biggr\vert\leq\frac{M(b-a)^3}{12n^2}$$

\item On prend $[a,b]=[0,1]$ et $f(x)=e^{-x^2}$. Calculer
$M=\sup_{[0,1]}\vert f''\vert$.
Déterminer $n$ pour que la méthode des trapèzes avec $n$ intervalles
donne un nombre qui approche $\int_0^1e^{-x^2}\,dx$ à moins de $10^{-2}$
près. En déduire un encadrement de cette intégrale.
\end{enumerate}
\finenonce{000804}



\finexercice



\exercice{2081, bodin, 2008/02/04}
\video{UXp2ntGBZNE}
\enonce{002081}{}
Soit $f$ la fonction définie sur $[0,4]$ par 
\begin{equation*}
  f(x)=
  \begin{cases}
    -1 &\text{ si $x=0$}\\
    1 &\text{ si $0<x<1$}\\
    3 &\text{ si $x=1$}\\
    -2 &\text{ si $1<x\leq 2$}\\
    4 &\text{ si $2<x\leq 4$.}
  \end{cases}
\end{equation*}
\begin{enumerate}
\item Calculer $\int_0^4f(t) \, dt$.
\item Soit $x\in [0,4]$, calculer $F(x)=\int_0^x  f(t) \, dt$.
\item Montrer que $F$ est une fonction continue sur $[0,4]$. La fonction
$F$ est-elle dérivable sur $[0,4]$ ?
\end{enumerate}
\finenonce{002081} 


\finexercice



\exercice{2082, bodin, 2008/02/04}
\video{KDx-xpueG-U}
\enonce{002082}{}
Soient les fonctions définies sur $\R$,
$$f(x)=x \text{ , } g(x)=x^2 \text{ et  } h(x)=e^x,$$
Justifier qu'elles sont intégrables sur tout intervalle fermé borné de $\R$. En utilisant les
sommes de Riemann, calculer les intégrales $\int_0^1f(x)d x$, $\int_1^2 g(x)
d x$ et $\int_0^x h(t) d t$.

\finenonce{002082} 


\finexercice


\exercice{2083, bodin, 2008/02/04}
\enonce{002083}{}
Calculer l'int\'egrale de $f:[a,b]\rightarrow \R$ comme
limite de sommes de Riemann-Darboux dans les cas suivants:
\begin{enumerate}
\item $f(x)=\sin x$ et $f(x)=\cos x$ sur $[0,\frac \pi 2]$ et $x_k=\frac{k\pi }{2n}$, $k=0,1,...,n$,
\item $g(x)=\frac 1x$ sur $\left[ a,b\right] \subset \R_{+}^{*}$ et $%
x_k=aq^k$ , $k=0,1,...,n$ ($q$ \'etant \`a d\'eterminer),
\item $h(x)=\alpha ^x$ sur $[a,b]$ , $\alpha >0$, et $x_k=a+(b-a).\frac kn$, $%
k=0,1,...,n$.
\end{enumerate}
\finenonce{002083} 
\finexercice


\exercice{2084, bodin, 2008/02/04}
\enonce{002084}{}
Les fonctions suivantes sont-elles int\'egrables au sens de Riemann?
\begin{enumerate}
\item $f(x)=[x]$ sur $[0,2]$
\item $g: [0,1] \rightarrow \R,\quad  g(x)=
\begin{cases}
 \left[\frac 1x\right] &\text{ si $0<x\leq 1$, }\\
 1&\text{si $x=0$ } 
\end{cases} $ 
\item $h: [0,1] \rightarrow \R,\quad  h(x)=
\begin{cases}
 \frac 1x\sin \left(\frac 1x\right) &\text{ si $0<x\leq 1$, }\\
 1&\text{si $x=0$ } 
\end{cases}$
\item $k: [0,1] \rightarrow \R,\quad  k_1(x)=
\begin{cases}
1&\text{ si $x\in [0,1]\cap \Q$, }\\
0&\text{ si $x\in [0,1]\backslash \Q$}
\end{cases}$
\end{enumerate}

\finenonce{002084} 
\finexercice



\exercice{2085, bodin, 2008/02/04}
\video{OjBbmNvQoSY}
\enonce{002085}{}
Soit $f:[a,b]\rightarrow \R$ une fonction continue sur $[a,b]$ ($a<b$).
\begin{enumerate}
\item On suppose que $f(x) \ge 0$ pour tout $x\in [a,b]$, et que $f(x_0)>0$ en un point $x_0\in [a,b]$. 
Montrer que $\int_a^b f(x) d x>0$. En déduire que : <<si $f$ est une fonction continue
positive sur $[a,b]$ telle que $\int_a^b f(x) d x=0$ alors $f$ est
identiquement nulle>>.
\item On suppose que $\int_a^b f(x) d x=0$. Montrer qu'il existe $c\in [a,b]$ tel que $f(c)=0$. 
\item Application: on suppose
que $f$ est une fonction continue sur $[0,1]$ telle que $\int_0^1 f(x) dx=\frac 12$. 
Montrer qu'il existe $d\in [0,1]$ tel que $f(d)=d$.
\end{enumerate}
\finenonce{002085} 


\finexercice



\exercice{2086, bodin, 2008/02/04}
\enonce{002086}{}
Soit $f:[a,b]\rightarrow \R$ continue, positive; on pose $%
m=\sup \{f(x),x\in [a,b]\}$. Montrer que 
\[
\lim_{n\rightarrow \infty }\left( \int_a^b(f(x))^n d x\right) _{}^{\frac
1n}=m. 
\]
\finenonce{002086} 
\finexercice


\exercice{2087, bodin, 2008/02/04}
\enonce{002087}{}
Soit $f:[0,1]\rightarrow {\R}$ une application strictement
croissante telle que $f(0)=0,\  f(1)=1$. Calculer :
$$\lim_{n\rightarrow \infty }\int_{0}^{1}f^{n}(t) d t. $$
\finenonce{002087} 
\finexercice



\exercice{2091, bodin, 2008/02/04}
\video{fKCXQ-Fl1Z8}
\enonce{002091}{}
Soit $f\,:\;\R\to\R$ une fonction continue sur $\R$ et $F(x)=\int_0^x
f(t)d t$. Répondre par vrai ou faux aux affirmations suivantes:
\begin{enumerate}
\item $F$ est continue sur $\R$.
\item $F$ est dérivable sur $\R$ de dérivée $f$.
\item Si $f$ est croissante sur $\R$ alors $F$ est croissante sur $\R$.
\item Si $f$ est positive sur $\R$ alors $F$ est positive sur $\R$.
\item Si $f$ est positive sur $\R$ alors $F$ est croissante sur $\R$. 
\item Si $f$ est $T$-périodique sur $\R$ alors $F$ est $T$-périodique sur $\R$.
\item Si $f$ est paire alors $F$ est impaire.
\end{enumerate}
\finenonce{002091} 
\finexercice



\exercice{2092, bodin, 2008/02/04}
\enonce{002092}{}
\label{ex:compint}
Soient $u$ et $v$ deux fonctions d\'erivables sur $\R$ et $f$ une fonction
continue sur $\R$.
\begin{enumerate}
\item On pose $\displaystyle F(x)=\displaystyle \int_{u(x)}^{v(x)}f(t)d t$. Montrer que $F$ est
d\'erivable sur $\R$ et calculer sa d\'eriv\'ee.
\item Calculer la d\'eriv\'ee de $\displaystyle G(x)=\int_x^{2x}\frac{d t}{1+t^2+t^4}$. 
 \end{enumerate}

\finenonce{002092} 
\finexercice


\exercice{2093, bodin, 2008/02/04}
\enonce{002093}{}
Soit $F(x)=\displaystyle \int_x^{x^2}\frac{1}{\ln t}d t$
\begin{enumerate}
\item Quel est l'ensemble de d\'efinition de $F$. $F$ est-elle continue,
d\'erivable sur son ensemble de d\'efinition ?
\item D\'eterminer $\lim_{x\to 1^+} F(x)$ en comparant $F$ \`a $H(x)=\displaystyle \int_x^{x^2}\frac{1}{t\ln t}d t$.
\end{enumerate}

\finenonce{002093} 
\finexercice


\exercice{2315, matexo1, 2002/02/01}
\enonce{002315}{}
\begin{enumerate}
\item Soit $f$ une fonction continue d\'efinie sur un intervalle
born\'e $[a,b] \subset \Rr$, telle que 
$$\int_a^b f(t)\,d t = (b-a) \min_{x\in [a,b]} f(x).$$
Montrer que $f$ est constante.
\item Soient $u$, $v$, deux fonctions continues sur $[a,b]$, \`a
valeurs dans $\Cc$. Montrer l'in\'egalit\'e de Cauchy-Schwarz
$$\int_a^b |u(t) v(t)| \,d t \le 
\left( \int_a^b |u(t)|^2 \right)^{1/2}
\left(\int_a^b |v(t)|^2 \right)^{1/2}.$$
Indication\,: poser, pour $\lambda \in \Cc$ arbitraire, 
$f_\lambda (t) = |\lambda u(t) +v(t)|^2$ 
et appliquer la question pr\'ec\'edente.
\item Dans quels cas cette in\'egalit\'e est-elle une \'egalit\'e\,?
\item Soit $C([a,b])$ l'espace des fonctions continues sur
$[a,b]$, \`a valeurs r\'eelles. Montrer que
$$ u \in C([a,b]) \mapsto 
\left(\int_a^b u(t)^2\,d t \right)^{1/2}$$
est une norme sur $C([a,b])$.

\end{enumerate}
\finenonce{002315}


\finexercice


\exercice{2316, matexo1, 2002/02/01}
\enonce{002316}{}
Soit $f$ une fonction continue sur l'intervalle
$[a,b]\subset\Rr$, \`a valeurs dans $\left]0,+\infty\right[$. Montrer
que
$$ \int_a^b f(t)\,d t \times \int_a^b \frac 1{f(t)}\,d t
\geq  (b-a)^2.$$
Dans quels cas y a-t-il \'egalit\'e ? (On pourra utiliser
l'in\'egalit\'e de Cauchy-Schwarz.)


\finenonce{002316}


\finexercice


\exercice{2317, matexo1, 2002/02/01}
\enonce{002317}{}
Soit $f : [0,1]\to \Rr$ une fonction continue.

\begin{enumerate}
\item En utilisant la formule de la moyenne, montrer que
$$\forall a\in[0,1[, \qquad
\lim_{n\to\infty} \int_0^a f(t^n)\,d t=a f(0).$$
\item Montrer qu'il existe $M>0$, tel que 
$$\forall a\in[0,1[, \qquad
\left| (a-1)f(0) + \int_a^1 f(t^n)\,d t \right| \leq  2 M(1-a).$$
En d\'eduire que 
$$\lim_{n\to\infty} \int_0^1 f(t^n)\,d t= f(0).$$
\end{enumerate}
\finenonce{002317}


\finexercice


\exercice{2324, matexo1, 2002/02/01}
\enonce{002324}{}
Soient $f$ et $g$ deux fonctions r\'eelles p\'eriodiques de 
p\'eriode $T$ continues sur $\Rr$. On appelle {\sl produit de 
convolution} de $f$ et $g$ la fonction $h$ not\'ee $f \star g$
et d\'efinie par
$$ h(x) = {1 \over T} \int_0^T f(t) g(x-t) ~dt. $$

\begin{enumerate}
\item Montrer que $h$ est une fonction p\'eriodique de 
p\'eriode $T$.
\item Montrer
$$h(x)={1 \over T} \int_a^{a+T} f(t) g(x-t) ~dt, \qquad \forall a \in \Rr.$$
\item En d\'eduire que $f \star g = g \star f$.
\end{enumerate}

\finenonce{002324}


\finexercice


\exercice{4220, quercia, 2010/03/12}
\enonce{004220}{Densité des fonctions en escalier}

Soit $f : {[a,b]} \to \R$ continue telle que pour toute fonction
$g : {[a,b]} \to \R$ en escalier, $ \int_{t=a}^b f(t)g(t)\,d t = 0$.
Démontrer que $f = 0$.

\finenonce{004220}



\finexercice


\exercice{4221, quercia, 2010/03/12}
\enonce{004221}{Changements de signe}

Soit $f : {[a,b]} \to \R$ continue non identiquement nulle, telle que :
$\forall\ k \in \{ 0, 1, \dots, n-1 \},\  \int_{t=a}^b t^kf(t)\,d t = 0$.

Démontrer, par récurrence, que $f$ change au moins $n$ fois de signe
sur $]a,b[$ (raisonner par l'absurde).


\finenonce{004221}



\finexercice


\exercice{4222, quercia, 2010/03/12}
\enonce{004222}{Formule de la moyenne généralisée}

Soient ${f,g} : {[a,b]} \to \R$ continues, $f$ positive.

\begin{enumerate}
  \item Démontrer qu'il existe $c \in {[a,b]}$ tel que
    $ \int_{t=a}^b f(t)g(t)\,d t = g(c) \int_{t=a}^b f(t)\,d t$.
  \item Si $f$ ne s'annule pas, montrer que $c \in {]a,b[}$.
  \item Application : Soit $f$ continue au voisinage de $0$. Déterminer
    $\lim_{x\to 0}\frac1{x^2} \int_{t=0}^x tf(t)\,d t$.
\end{enumerate}
\finenonce{004222}



\finexercice


\exercice{4223, quercia, 2010/03/12}
\enonce{004223}{Inégalité de Jensen}

Soit $f : {[a,b]} \to \R$ continue et $g : \R \to \R$ continue convexe.

Démontrer que $g\left(\frac 1{b-a} \int_{t=a}^b f(t)\,d t\right)
      \le \frac 1{b-a} \int_{t=a}^b g(f(t))\,d t$.
\finenonce{004223}



\finexercice


\exercice{4224, quercia, 2010/03/12}
\enonce{004224}{$\sqrt{1+f^2}$}

Soit $f : {[0,1]} \to \R$ continue positive. On pose $A =  \int_{t=0}^1 f(t)\,d t$.

Montrer que $\sqrt{1+A^2} \le  \int_{t=0}^1 \sqrt{1+f^2(t)}\,d t \le 1+A$.


\finenonce{004224}



\finexercice


\exercice{4225, quercia, 2010/03/12}
\enonce{004225}{Calcul de limite}

Chercher
$\lim_{x \to 0^+}  \int_{t=x}^{2x} \frac{\cos t\,\ln(1+t^2)}{\sin^2t\,\sh t}\,d t$.
\finenonce{004225}



\finexercice


\exercice{4226, quercia, 2010/03/12}
\enonce{004226}{Calcul de limite}

Pour $0<a<b$, déterminez $\lim_{x\to 0^+} \int_{t=ax}^{bx} \frac{1-\cos u}{u^3}\,d u$.


\finenonce{004226}



\finexercice


\exercice{4227, quercia, 2010/03/12}
\enonce{004227}{$\int f + \int f^{-1}$}
Soit $f : {[a,b]} \to {[c,d]}$ continue, bijective, strictement croissante.

Calculer $ \int_{t=a}^b f(t)\,d t +  \int _{u=c}^d f^{-1}(u)\,d u$
(faire un dessin, et commencer par le cas où $f$ est de classe $\mathcal{C}^1$).


\finenonce{004227}



\finexercice

\exercice{4232, quercia, 2010/03/12}
\enonce{004232}{Maximum-minimum}

Soient $a,b \in \R$. \'Etudier la convergence des suites $(a_n)$, $(b_n)$
définies par :
$$a_0=a,\ b_0=b,
  \qquad a_{n+1} = \frac12 \int_{x=-1}^1 \min(x,b_n)\,d x,\
                       b_{n+1} = \frac12 \int_{x=-1}^1 \max(x,a_n)\,d x.$$
\finenonce{004232}



\finexercice

\exercice{4234, quercia, 2010/03/12}
\enonce{004234}{Intégrale de $|f|$}

Soit $f : {[a,b]} \to \R$ continue. Pour $n \in \N^*$, on pose
$I_n = \sum_{k=0}^{n-1} \left|  \int_{t=a_k}^{a_{k+1}} f(t)\,d t \right|$
où $a_k = a + k\frac {b-a}n$.

Montrer que $I_n \to  \int_{t=a}^b |f(t)|\,d t$ lorsque $n\to\infty$.

\finenonce{004234}



\finexercice


\exercice{4235, quercia, 2010/03/12}
\enonce{004235}{Usage de symétrie}

Soit $I =  \int_{t=0}^\pi \ \frac {t\sin t}{1 + \cos^2t}\,d t$.
Effectuer dans $I$ le changement de variable $u = \pi - t$, et en déduire
la valeur de~$I$.

\finenonce{004235}



\finexercice


\exercice{4236, quercia, 2010/03/12}
\enonce{004236}{Usage de symétrie}

Calculer $I =  \int_{t=0}^\pi \ \frac t{1+\sin t}\,d t$.
\finenonce{004236}



\finexercice


\exercice{4237, quercia, 2010/03/12}
\enonce{004237}{Usage de symétrie}

Calculer $ \int_{t=0}^{\pi/4} \ln(1+\tan t)\,d t$.
On remarquera que $\cos t + \sin t = \sqrt 2 \cos\left(\frac \pi4 - t\right)$.
\finenonce{004237}



\finexercice

\exercice{5444, rouget, 2010/07/10}
\enonce{005444}{****}
Soient $f$ et $g$ deux fonctions continues et strictement positives sur $[a,b]$. Pour $n$ entier naturel non nul donné, on pose $u_n=\left(\int_{a}^{b}(f(x))^ng(x)\;dx\right)^{1/n}$.

Montrer que la suite $(u_n)$ converge et déterminer sa limite (commencer par le cas $g=1$).
\finenonce{005444}


\finexercice


\exercice{5449, rouget, 2010/07/10}
\enonce{005449}{***T}
Soit $E$ l'ensemble des fonctions continues strictement positives sur $[a,b]$.

Soit $\begin{array}[t]{cccc}
\varphi~:&E&\rightarrow&\Rr\\
 &f&\mapsto&\left(\int_{a}^{b}f(t)\;dt\right)\left(\int_{a}^{b}\frac{1}{f(t)}\;dt\right)
 \end{array}$.
	   
\begin{enumerate}
\item  Montrer que $\varphi(E)$ n'est pas majoré.
\item  Montrer que $\varphi(E)$ est minoré. Trouver $m=\mbox{Inf}\{\varphi(f),\;f\in E\}$. Montrer que cette borne infèrieure est atteinte et trouver toutes les $f$ de $E$ telles que $\varphi(f)=m$.
\end{enumerate}
\finenonce{005449}


\finexercice


\exercice{5917, tumpach, 2010/11/11}
\enonce{005917}{}
En utilisant la d\'efinition d'une fonction int\'egrable au sens
de Riemann, montrer les propri\'et\'es suivantes~:
\begin{enumerate}
\item Si $f$ et $g$ sont Riemann-int\'egrables sur $[a,b]$, alors
 $f + g$ est Riemann-int\'egrable sur $[a,b]$.
\item Si $f$ est Riemann-int\'egrable sur $[a,b]$ et $\lambda \in
\mathbb{R}$, alors $\lambda\,f$ est Riemann-int\'egrable sur
$[a,b]$. 
\item Si $f$ et $g$ sont deux fonctions
Riemann-int\'egrables sur $[a,b]$ telles que, pour tout $t\in
[a,b]$, $f(t)\leq g(t)$, alors
$\int_{a}^{b}f(t)\,dt\leq\int_{a}^{b}g(t)\,dt$. 
\item Une limite
uniforme de fonctions Riemann-int\'egrables sur $[a,b]$ est
Riemann-int\'egrable sur $[a,b]$.
\end{enumerate}
\finenonce{005917}



\finexercice


\exercice{5918, tumpach, 2010/11/11}
\enonce{005918}{}
Montrer qu'une fonction \emph{monotone} sur $[a,b]$ est
Riemann-int\'egrable sur $[a,b]$.
\finenonce{005918}



\finexercice


\exercice{5919, tumpach, 2010/11/11}
\enonce{005919}{}
Montrer qu'une fonction \emph{continue} sur $[a,b]$ est
Riemann-int\'egrable sur $[a,b]$.
\finenonce{005919}



\finexercice


\exercice{5920, tumpach, 2010/11/11}
\enonce{005920}{}
\begin{enumerate}
\item Montrer que la fonction $f~:[0,1]\rightarrow \mathbb{R}$
d\'efinie par~:
$$
 f(x) = \left\{\begin{array}{l} 1 \quad\text{si}\quad  x\in\mathbb{Q}\\
0 \quad\text{si} \quad x\in\mathbb{R}\setminus\mathbb{Q}
\end{array}\right.
$$
n'est pas Riemann-int\'egrable sur $[0,1]$. 

\item Montrer que la
fonction $g~:[0,1]\rightarrow \mathbb{R}$ d\'efinie par~:
$$
 g(x) = \left\{\begin{array}{l} \frac{1}{q} \quad\text{si}\quad
 \, x=\frac{p}{q} \quad\text{avec}\,\, p \,\,\text{et} \,\,q \,\,\text{premiers entre eux}\\
0 \quad\text{si} \quad x\in\mathbb{R}\setminus\mathbb{Q} \text{ ou } x=0
\end{array}\right.
$$
est Riemann-int\'egrable sur $[0,1]$.
\end{enumerate}
\finenonce{005920}



\finexercice


\exercice{5921, tumpach, 2010/11/11}
\enonce{005921}{}
On dit qu'une partie $A$ de $\mathbb{R}$ est \emph{n\'egligeable}
si, pour tout nombre r\'eel $\varepsilon>0$, il existe une suite
$(I_{n})_{n\in\mathbb{N}}$ d'intervalles $I_{n} = ]a_{n},b_{n}[$
telle que~:
\begin{equation*}
A \subset \bigcup_{n\in\mathbb{N}}I_{n}\quad \text{et}\quad
\sum_{n\in\mathbb{N}}(b_{n}-a_{n})\leq\varepsilon.
\end{equation*}
\begin{enumerate}
\item Montrer qu'une r\'eunion d\'enombrable d'ensembles
n\'egligeables est un ensemble n\'egligeable. \item Montrer qu'une
fonction born\'ee $f~:[a,b]\rightarrow\mathbb{R}$ est int\'egrable
au sens de Riemann sur $[a,b]$ si et seulement si l'ensemble des
points o\`u $f$ n'est pas continue est \emph{n\'egligeable}.
\end{enumerate}
\finenonce{005921}



\finexercice


\exercice{5922, tumpach, 2010/11/11}
\enonce{005922}{}
Pour tout $n\in\mathbb{N}$, on d\'efinit
$f_{n}~:]0,1]\rightarrow\mathbb{R}$ par~: $f_{n}(x)= n e^{-nx}$.
Montrer que la suite $(f_{n})_{n\in\mathbb{R}}$ converge
simplement vers une fonction $f$ sur $]0,1]$ mais que
\begin{equation*}
\int_{0}^{1}\lim_{n\rightarrow+\infty}f_{n}(x)\,dx \quad\neq\quad
\lim_{n\rightarrow+\infty}\int_{0}^{1}f_{n}(x)\,dx.
\end{equation*}
V\'erifier que la convergence de $(f_{n})_{n\in\mathbb{N}}$ vers
$f$ n'est pas \emph{uniforme} sur $]0,1]$.
\finenonce{005922}



\finexercice


\exercice{5923, tumpach, 2010/11/11}
\enonce{005923}{}
Montrer que, si $f~:[a,b] \rightarrow \mathbb{R}$ est une fonction
int\'egrable au sens de Riemann, on a~:
\begin{equation*}
\frac{1}{b-a}\int_{a}^{b}f(t)\,dt=\lim_{n\rightarrow+\infty}\frac{1}{n}\sum_{k=1}^{n}f\left(a+k\frac{b-a}{n}\right).
\end{equation*}
En d\'eduire les limites suivantes~:
\begin{equation*}
a)\quad\lim_{n\rightarrow+\infty}
\frac{1}{n}\sum_{k=1}^{n}\tan\frac{k}{n} \quad\quad\quad
b)\quad\lim_{n\rightarrow+\infty}\sum_{k=1}^{n}\frac{n}{n^2+k^2}
\quad\quad\quad
c)\quad\lim_{n\rightarrow+\infty}\sum_{k=1}^{n}\log\left(\frac{n}{n+k}\right)^{\frac{1}{n}}
\end{equation*}
\finenonce{005923}



\finexercice


\exercice{5924, tumpach, 2010/11/11}
\enonce{005924}{}
\begin{enumerate}
\item Montrer que si $f~:[a,b]\rightarrow\mathbb{R}$ est
Riemann-int\'egrable, alors
\begin{equation*}\int_{a}^{b}f(x)\,dx =
\int_{a}^{b}f(a+b-x)\,dx. \end{equation*} \item Calculer (en
utilisant 1.) les int\'egrales suivantes~:
\begin{equation*}
a)\quad\int_{0}^{\pi}\frac{x\sin
x}{1+{\cos}^{2}x}\,dx\quad\quad\quad\quad
b)\quad\int_{0}^{\frac{\pi}{4}}\log\left(1+\tan x\right)\,dx.
\end{equation*}
\end{enumerate}
\begin{equation*}
\textit{Rappel~:}\quad \tan(\alpha-\beta) =
\frac{\tan(\alpha)-\tan(\beta)}{1+\tan(\alpha)\tan(\beta)}
\end{equation*}
\finenonce{005924}



\finexercice

\section{127.02 Somme de Riemann}


\exercice{794, gourio, 2001/09/01}

\enonce{000794}{}
Soient $f$ et $g$ de ${\Rr}^{+}$ dans ${\Rr}$ croissantes.\ Montrer que :
$$\forall x\in {\Rr}^{+},\left( \int_{0}^{x}f\right) \left(
\int_{0}^{x}g\right) \leq x\int_{0}^{x}fg. $$

\emph{Indication} : on \'{e}tablira d'abord que, si $a_{1}\leq a_{2}\leq ...\leq
a_{n}$ et $b_{1}\leq b_{2}\leq ...\leq b_{n},$ alors :
$$\left( \frac{1}{n}\sum\limits_{i=1}^{n}a_{i}\right) \left( \frac{1}{n}
\sum\limits_{i=1}^{n}b_{i}\right) \leq \frac{1}{n}\sum \limits_{i=1}^{n}a_{i}b_{i}. $$
Remarquer que :
$$\sum\limits_{1\leq i\leq j\leq n}(a_{i-}a_{j})(b_{i-}b_{j})\geq 0. $$
\finenonce{000794}



\finexercice



\exercice{840, gourio, 2001/09/01}

\enonce{000840}{}
Calculer :
$$\lim\limits_{n\rightarrow \infty }\sum\limits_{k=1}^{n}\frac{(-1)^{k}}{2k+1}
,\lim\limits_{n\rightarrow \infty }\sum\limits_{k=1}^{n}\frac{(-1)^{k}}{k}. $$
\finenonce{000840}



\finexercice



\exercice{841, gourio, 2001/09/01}

\enonce{000841}{}
Calculer :
$$
\lim\limits_{n\rightarrow \infty }\prod\limits_{k=1}^{n}\left( 1+\frac{k^{2}
}{n^{2}}\right) ^{\frac{1}{n}}\ ;\
\lim\limits_{n\rightarrow \infty
}n\sum\limits_{k=1}^{n}\frac{e^{-\frac{n}{k}}}{k{{}^{2}}}\ ;\
\lim\limits_{n\rightarrow \infty }\sum\limits_{k=1}^{n}\frac{n+k}{n^{2}+k{}}
\ ;\ \lim\limits_{n\rightarrow \infty }\sum\limits_{k=1}^{n-1}\frac{1}{\sqrt{
n^{2}-k{{}^{2}}}}.
$$
Soit $(u_{n})_{n\in \Nn^{*}\text{ }}$ la suite r\'{e}elle d\'{e}finie par :
$$\forall n\in \Nn^{*},\text{ }u_{n}=\sum\limits_{k=1}^{n}\frac{n}{n{{}^{2}}
+k{{}^{2}}}.  $$
Calculer :
$$\ell=\lim\limits_{n\rightarrow \infty }u_{n} $$
et donner un \'{e}quivalent de $u_{n}-\ell.$
\finenonce{000841}



\finexercice



\exercice{842, gourio, 2001/09/01}

\enonce{000842}{}
Soient $f$ et $g$ continues de $[0,1]$ dans ${\Rr}.$
Calculer :
$$\lim\limits_{n\rightarrow \infty }\frac{1}{n}\sum\limits_{k=0}^{n-1}f\left(\frac{k}
{n}\right)
g\left(\frac{k+1}{n}\right). $$
\finenonce{000842}



\finexercice



\exercice{843, gourio, 2001/09/01}

\enonce{000843}{}
Calculer :
$$\lim_{n\rightarrow \infty }\sum\limits_{k=0}^{n^{2}}\frac{n}{n^{2}+k^{2}}. $$
\finenonce{000843}



\finexercice



\exercice{844, cousquer, 2003/10/01}

\enonce{000844}{}
 Calculer les limites suivantes~:
\begin{enumerate}
\item $\displaystyle\lim_{n\to\infty}
    \frac{1+\sqrt2+\sqrt3+\cdots+\sqrt n}{n\sqrt n}$.
\item $\displaystyle\lim_{n\to\infty}\sum_{p=1}^n\frac{n}{n^2+p^2}$.
\item $\displaystyle\lim_{n\to\infty}\frac{1}{n}\sum_{p=1}^n
\ln\frac{(3n+6p-4)(n+2p)^2}{3n^3}$.
\end{enumerate}
\finenonce{000844}



\finexercice



\exercice{2100, bodin, 2008/02/04}
\video{Mw9ODmfCvyM}
\enonce{002100}{}
Calculer la limite des suites suivantes :
  \begin{enumerate}
  \item $\displaystyle u_n=n\sum_{k=0}^{n-1}\frac 1{k^2+n^2}$
  \item $\displaystyle v_n=\prod\limits_{k=1}^n\left(1+\frac{k^2}{n^2}\right) ^{\frac 1n}$
  \end{enumerate}
\finenonce{002100} 


\finexercice


\exercice{4228, quercia, 2010/03/12}
\enonce{004228}{Sommes de Riemann}

\begin{enumerate}
  \item Trouver $\lim_{n\to\infty} \frac 1{n+1} + \frac 1{n+2} + \dots + \frac 1{kn}$
    pour $k$ entier supérieur ou égal à $2$ fixé.

  \item Trouver $\lim_{n\to\infty} \frac 1{n^2} \Bigl( \sqrt{1(n-1)}
    + \sqrt{2(n-2)} + \dots + \sqrt{(n-1)1} \Bigr)$.
    

  \item Trouver $\lim_{n\to\infty} \sqrt[n]{\Bigl(1+\frac 1n\Bigr)\Bigl(1+\frac 2n\Bigr) \dots
             \Bigl(1+\frac nn\Bigr)}$.
    

  \item Trouver $\lim_{n\to\infty} \ln\left(1+\frac\pi n\right)\sum_{k=0}^{n-1} \frac1{2+\cos(3k\pi/n)}$.
    

  \item Donner un équivalent pour $n\to\infty$ de $\sum_{k=1}^n \sqrt k$.
    

  \item Soit $A_1A_2\dots A_n$ un polygone régulier inscrit dans un cercle de
    rayon 1. Chercher $\lim_{n\to\infty} \frac 1n\sum_{k=2}^n A_1A_k$.
    

\end{enumerate}
\finenonce{004228}



\finexercice


\exercice{4229, quercia, 2010/03/12}
\enonce{004229}{Calcul de limite}

Soit $f : {[0,1]} \to \R$ continue.
Chercher $\lim_{n \to \infty} \frac 1{n^2}
          \sum_{1 \le i < j \le n} f\Bigl(\frac in\Bigr)f\Bigl(\frac jn\Bigr)$.


\finenonce{004229}



\finexercice


\exercice{4230, quercia, 2010/03/12}
\enonce{004230}{Moyenne géométrique}

Soit $f : {[0,1]} \to \R$ continue.
Montrer que $\Bigl(1+\frac1nf\bigl(\frac1n\bigr)\Bigr)
     \Bigl(1+\frac1nf\bigl(\frac2n\bigr)\Bigr)\dots
     \Bigl(1+\frac1nf\bigl(\frac nn\bigr)\Bigr)
     \to \exp \int_{t=0}^1 f(t)\,d t$ lorsque $n\to\infty$.

(On pourra utiliser : $\forall\ x\ge -\frac12,\ x-x^2 \le \ln x \le x$)
\finenonce{004230}



\finexercice


\exercice{4231, quercia, 2010/03/12}
\enonce{004231}{}
\begin{enumerate}
  \item Montrer que : $\forall\ x \ge 0,\ x-\frac{x^2}2 \le \ln(1+x) \le x$.
  \item Trouver $\lim_{n\to\infty}\prod_{k=1}^n\left(1+\frac1{k^2+n^2}\right)^n$.
    

\end{enumerate}
\finenonce{004231}



\finexercice

\exercice{4233, quercia, 2010/03/12}
\enonce{004233}{Intégrale de $\ln|x-e^{it}|$}

Pour $x \in \R,\ x \ne \pm1$, on pose $I =  \int_{t=0}^{2\pi} \ln|x-e^{it}|\,d t$.
En utilisant les sommes de Riemann, calculer $I$.

\finenonce{004233}



\finexercice

\exercice{5446, rouget, 2010/07/10}
\enonce{005446}{***IT}
Limites de 
$$\begin{array}{llll}
1)\;\frac{1}{n^3}\sum_{k=1}^{n}k^2\sin\frac{k\pi}{n}&2)\;(\frac{1}{n!}\prod_{k=1}^{n}(a+k))^{1/n}\;(a>0\;\mbox{donné})&3)\;\sum_{k=1}^{n}\frac{n+k}{n^2+k}&4)\;\sum_{k=1}^{n}\frac{1}{\sqrt{n^2-k^2}}\\  
5)\;\frac{1}{n\sqrt{n}}\sum_{k=1}^{n}E(\sqrt{k})&6)\;\sum_{k=1}^{n}\frac{k^2}{8k^3+n^3}&7)\;\sum_{k=n}^{2n-1}\frac{1}{2k+1}&8)\;n\sum_{k=1}^{n}\frac{e^{-n/k}}{k^2}
\end{array}  
$$
\finenonce{005446}


\finexercice


\exercice{5447, rouget, 2010/07/10}
\enonce{005447}{***I}
Soit $f$ une fonction de classe $C^2$ sur $[0,1]$. Déterminer le réel $a$ tel que :
 
$$\int_{0}^{1}f(t)\;dt-\frac{1}{n}\sum_{k=1}^{n-1}f(\frac{k}{n})\underset{n\rightarrow+\infty}{=}\frac{a}{n}+o(\frac{1}{n}).$$

\finenonce{005447}


\finexercice

\section{127.03 Longueur, aire, volume}


\exercice{805, cousquer, 2003/10/01}

\enonce{000805}{}
 Construire la courbe paramétrée
$C\left\lbrace\renewcommand{\arraystretch}{1.2}
\begin{array}{l}
    x = \frac{\cos t}{1+\lambda\cos t}  \\
    y = \frac{\sin t}{1+\lambda\cos t}
\end{array}\right.$
où $\lambda$ est un
paramètre appartenant à $[0,1\mathclose[$. \\Calculer l'aire $S$ limitée par $C$
de deux façons~:
\begin{itemize}
\item En se ramenant au calcul de $\int_0^{2\pi}\frac{dt}{(1+\lambda\cos t)^2}$.

\item En reconnaissant la nature géométrique de $C$.
\end{itemize}

\finenonce{000805}



\finexercice



\exercice{806, cousquer, 2003/10/01}

\enonce{000806}{}
 Représenter la courbe définie par son équation polaire
$\rho=a\sin^3\frac{\theta}{3}$. Calculer sa longueur $L$ et les aires
$A_1$ et $A_2$ limitées par les deux boucles qu'elle forme.

\finenonce{000806}






\finexercice



\exercice{807, cousquer, 2003/10/01}

\enonce{000807}{}
 On appelle \emph{tore} la figure obtenue par révolution d'un cercle
de rayon $r$ autour d'une droite de son plan passant à distance $R$ de son
centre (on suppose $r<R$). Calculer l'aire $A$ du tore, et son volume~$V$.

\finenonce{000807}



\finexercice



\exercice{808, cousquer, 2003/10/01}

\enonce{000808}{}
 On appelle \emph{cycloïde} la courbe décrite par un point d'un
cercle de rayon~$R$, lié à ce cercle, quand celui-ci roule sans glisser
sur une droite en restant dans plan fixe.
Montrer que dans un repère bien choisi,
la cycloïde admet la représentation paramétrique~:
$\left\lbrace
\begin{array}{l}
    x = R(t-\sin t)  \\
    y = R(1-\cos t)
\end{array}\right.$
Représenter la cycloïde et calculer~: la longueur $L$ d'une arche,
l'aire $A$ de la surface $S$ comprise entre cette arche et la droite fixe
$(Ox)$,
les volumes $V_1$ et $V_2$ obtenus par révolution de $S$ autour de $Ox$
et $Oy$ respectivement,
les aires $A_1$ et $A_2$ obtenues par révolution d'une arche
de la cycloïde  autour de $Ox$ et $Oy$ respectivement.

\finenonce{000808}





\finexercice



\exercice{809, cousquer, 2003/10/01}

\enonce{000809}{}
 On appelle \emph{épicycloïde} la courbe décrite par un point d'un
cercle de rayon~$r$, lié à ce cercle, quand celui-ci roule sans glisser
sur un cercle de rayon~$R$ en restant tangent extérieurement à ce dernier, 
et dans son plan. On pose $n=R/r$.
Montrer que dans un repère que l'on précisera, l'épicycloïde
admet la représentation paramétrique~:
$$\left\lbrace
\begin{array}{rcl}
    x & = & r\bigl((n+1)\cos t-\cos(n+1)t\bigr)  \\
    y & = & r\bigl((n+1)\sin t-\sin(n+1)t\bigr)  
\end{array}\right.$$
Représenter la courbe pour $n=1,2,3$.
En supposant $n$ entier, calculer la longueur $L$ de la courbe et l'aire $A$
limitée par celle-ci.
Dans le cas $n=1$ (\emph{cardioïde}), calculer de plus l'aire $S$ de la
surface de révolution obtenue en faisant tourner la courbe autour de son
axe de symétrie, ainsi que le volume $V$ limitée par cette surface.


\finenonce{000809}



\finexercice



\exercice{810, cousquer, 2003/10/01}

\enonce{000810}{}
 Soit $C$ un cercle fixe de rayon $R$. Un cercle $C'$ de même rayon
roule sans glisser sur $C$ en restant dans un plan (variable) perpendiculaire
à celui de~$C$. Un point~$M$ lié au cercle~$C'$ décrit une courbe~$\Gamma$. 
Montrer que suivant un repère convenablement choisi, $\Gamma$
admet la représentation paramétrique~:
$\left\lbrace
\begin{array}{rcl}
    x & = & R(\cos t+\sin^2t)  \\
    y & = & R\sin t(1-\cos t)  \\
    z & = & R(1-\cos t)
\end{array}\right.$.
En déduire la longueur $L$ de $\Gamma$.
Représenter les projections de $\Gamma$ sur chacun des trois plans
de coordonnées.

\finenonce{000810}



\finexercice



\exercice{2098, bodin, 2008/02/04}
\enonce{002098}{}
Calculer $\int_{-R}^R\sqrt{R^2-x^2}d x$ (on posera $\theta =\arcsin \frac xR$
) et en d\'eduire l'aire d'un disque de rayon $R$. 
\finenonce{002098} 
\finexercice


\exercice{2099, bodin, 2008/02/04}
\video{1rhApdE7JPY}
\enonce{002099}{}
Calculer l'aire de la région délimitée par les courbes
d'équation $\displaystyle y=\frac{x^2}2$ et $\displaystyle y=\frac 1{1+x^2}$. 
\finenonce{002099} 


\finexercice


\exercice{4241, quercia, 2010/03/12}
\enonce{004241}{Approximation des rectangles pour une fonction lipchitzienne}

Soit $f : {[a,b]} \to \R$, $K$-lipchitzienne.

Montrer que $\left| \int_{t=a}^b f(t)\,d t - \frac{b-a}n\sum_{k=1}^n
f\Bigl(a+k\frac{b-a}n\Bigr)\right| \le \frac{K(b-a)^2}{2n}$.

\finenonce{004241}



\finexercice


\exercice{4242, quercia, 2010/03/12}
\enonce{004242}{Approximation des tangentes}

Soit $f : {[a,b]} \to \R$ de classe $\mathcal{C}^2$. On fixe $n\in\N^*$ et on note :
$a_k = a + k\frac{b-a}n$, $a_{k+{\frac 1 2}} = \frac{ a_k + a_{k+1} }2$.

Soit $I_n = \frac{b-a}n\sum_{k=0}^{n-1} f(a_{k+{\frac 1 2}})$.

\begin{enumerate}
  \item Donner une interprétation géométrique de $I_n$.
  \item Montrer que $\left| \int_{t=a}^b f(t)\,d t - I_n \right| \le
    \frac{M_2(b-a)^3}{24n^2}$ où $M_2 = \sup\limits_{[a,b]}|f''|$.
\end{enumerate}
\finenonce{004242}



\finexercice


\exercice{4243, quercia, 2010/03/12}
\enonce{004243}{Approximation des trapèzes}

Soit $f : {[a,b]} \to \R$ de classe $\mathcal{C}^2$.


\begin{enumerate}
  \item Montrer que $ \int_{t=a}^b f(t)\,d t = (b-a)\frac {f(a)+f(b)}2
                    +  \int_{t=a}^b \frac {(t-a)(t-b)}2f''(t)\,d t$.

  \item Application : Soit $f : {[a,b]} \to \R$, $I =  \int_{t=a}^b f(t)\,d t$,
    et $I_n$ la valeur approchée de $I$ obtenue par la méthode des trapèzes
    avec $n$ intervalles.
    Démontrer que $|I-I_n| \le \frac {\sup|f''| (b-a)^3}{12n^2}$.

\end{enumerate}
\finenonce{004243}



\finexercice

\exercice{4245, quercia, 2010/03/12}
\enonce{004245}{Aire sous une corde}

Soit $f : {[a,b]} \to {\R}$ de classe $\mathcal{C}^1$ telle que $f(a) = f(b) = 0$.
On pose $M' = \|f'\|_\infty$.

\begin{enumerate}
  \item   En majorant $f$ par une fonction affine par morceaux, démontrer que
      $\left| \int_{t=a}^b f(t)\,d t \right| \le M'\frac {(b-a)^2}4$.
  \item   Quand y a-t-il égalité ?
\end{enumerate}
\finenonce{004245}



\finexercice

\exercice{6863, bodin, 2012/04/13}
\video{U6GrjVSfshM}
\enonce{006863}{}
Calculer l'aire intérieure d'une ellipse d'équation :
$$\frac{x^2}{a^2}+ \frac{y^2}{b^2} = 1.$$

\emph{Indications.}  On pourra calculer seulement la partie de l'ellipse correspondant
à $x\ge 0$, $y\ge 0$. Puis exprimer $y$ en fonction de $x$. Enfin calculer une intégrale.

\finenonce{006863}


\finexercice
\section{127.04 Intégration à  l'aide d'une fonction auxiliaire}


\exercice{811, bodin, 1998/09/01}

\enonce{000811}{}
Calculer les primitives suivantes :
$$
\int \frac{dx}{x^2+5} \quad ; \quad
\int \frac{dx}{\sqrt{x^2-5}} \quad ; \quad
\int e^x\sin(e^x) dx \quad ; \quad
\int \tan^3x dx \quad ;
$$
$$
\int \frac{1}{\tan^3x}dx \quad ; \quad
\int\frac{2x+3}{(x^2+3x+7)^m}dx, \ m\in\Nn \quad ; \quad
\int \frac{\ln x}{x}dx \quad ; \quad \int \frac{\ch x\, dx}{\sh^5x}.
$$
\finenonce{000811}



\finexercice


\section{127.05 Changement de variables}


\exercice{812, bodin, 1998/09/01}

\enonce{000812}{}
Consid\'erons l'int\'egrale  $$I=\int_0^{\ln 2}\sqrt{e^x-1}\, dx$$
Effectuer le changement de variables $u=\sqrt{e^x-1}$
et calculer $I$.\\

\emph{R\'esultat: $I=2-\pi/2$.}
\finenonce{000812}



\finexercice



\exercice{813, bodin, 1998/09/01}

\enonce{000813}{}
 Soit $f:[a,b]\rightarrow \R$ une fonction strictement
croissante et contin\^ument d\'erivable.
On consid\`ere les deux int\'egrales $I_1=\int_a^b f(t)\, dt$ et
$I_2=\int_{f(a)}^{f(b)} f^{-1}(t)\, dt$.
\begin{enumerate}
    \item  Rappeler pourquoi $f$ admet une fonction r\'eciproque $f^{-1}$.
    \item  Faire le changement de variable $t=f (u)$ dans l'int\'egrale $I_2$.
    \item Calculer $I_2$ en fonction de $I_1$.
    \item Faire un dessin faisant appara\^{\i}tre $f$ et $f^{-1}$,
et interpr\'eter ce r\'esultat g\'eom\'etriquement.
\end{enumerate}
\finenonce{000813}



\finexercice



\exercice{814, bodin, 1998/09/01}

\enonce{000814}{}
Calculer les primitives suivantes :
$$
\int \frac{1}{\sqrt{2+x}+\sqrt[3]{2+x}}dx, \quad (t=\sqrt[6]{2+x}) \ ;
$$

$$
\int \frac{1}{((x-1)^2-4)^2}dx, \quad (\frac{x-1}{2} = \tanh u \text{ ou coth}\, u) \ ;
$$

$$
\int (\arcsin x)^2 dx \quad ; \quad
\int x^2 \sqrt{1+x^3}dx.
$$
\finenonce{000814}



\finexercice



\exercice{2318, matexo1, 2002/02/01}
\enonce{002318}{}
 Sans calculer les int\'egrales, montrer que
$$ \int_0^{\pi /2} \sin^n x dx = \int_0^{\pi /2} \cos^n x dx. $$



\finenonce{002318}


\finexercice


\exercice{2321, matexo1, 2002/02/01}
\enonce{002321}{}
Calculer les int\'egrales suivantes\,:
\begin{eqnarray*}
&&\int_0^1 {t  \over {\sqrt{1-t^2}}} ~dt, \\
&&\int_0^a \sqrt{a^2 - t^2} ~dt, \\
&&\int_0^\pi t^2 \sin t ~dt, \\
&&\int_{1-  {\pi^2 \over 4}}^1 \cos \sqrt{1-t}~ dt.
\end{eqnarray*}


\finenonce{002321}


\finexercice


\exercice{2322, matexo1, 2002/02/01}
\enonce{002322}{}
Calculer les int\'egrales suivantes\,:
\begin{eqnarray*}
&& \int_0^\pi {dt \over {(2+\cos^2 t)^2}},  \\
&& \int_0^{\pi / 4} \cos^2 t \cos 3t \sqrt{\cos 2t}~ dt, \\
&& \int_0^1 {dt \over {\sqrt{1+t^2}+\sqrt{1-t^2}}}.
\end{eqnarray*}

\finenonce{002322}


\finexercice


\exercice{2323, matexo1, 2002/02/01}
\enonce{002323}{}
Soit $f$ une fonction continue dans $[0,\pi]$. Montrer,
en utilisant un changement de variables, que l'on a
$$ \int_0^\pi x f(\sin x) ~dx = {{\pi \over 2}} \int_0^\pi f(\sin x) ~dx.$$
En d\'eduire la valeur de
$$ \int_0^\pi {x \sin x \over {1+\cos^2 x}} ~dx. $$

\finenonce{002323}


\finexercice


\exercice{2326, matexo1, 2002/02/01}
\enonce{002326}{}
Calculer les int\'egrales suivantes\,:
\begin{eqnarray*}
&& \int_1^e {t^n \ln^4 t}\,dt, \qquad n \neq -1, \\
&& \int_0^1 {dt \over {\sqrt{x} (1+\sqrt[3]{x})^2}}, \\
&& \int_a^b \sqrt{(t-a)(t-b)}\,dt, \\
&& \int_0^1 2^t \cdot 3^{2t} \cdot 5^{3t}\,dt, \\ 
&& \int_0^1 {dt \over {\sqrt{x^2 +2x +5}}}, \\
&& \int_{-\pi}^\pi \sqrt{1+\cos t}\,dt, \\
&& \int_0^1 t^7 \arctan t\,dt.
\end{eqnarray*}
\finenonce{002326}


\finexercice


\exercice{2329, matexo1, 2002/02/01}
\enonce{002329}{}
Soit $x>0$ un r\'eel. Calculer les valeurs de
$$I(x) = \int_0^x{\arctan t \over 1+t^2} dt \qquad \mbox{\rm et}
\qquad J(x) = \int_0^x{\arctan t \over (1+t)^2} dt. $$
Quelles sont leurs limites quand $x \to  +\infty$\,?
\finenonce{002329}


\finexercice


\exercice{2335, matexo1, 2002/02/01}
\enonce{002335}{}
Trouver les primitives des fonctions suivantes\,:
$${
{1\over x \sqrt{x^2-1}}, \qquad 
{(4x^2+4x+5)^{-1/2}}, \qquad
{\sin^4 x\over \cos^2 x}, \qquad
{\arctan x \over 1+x^2}, \qquad
{1\over x\,(1+\ln^2 x)}. 
}$$
\finenonce{002335}


\finexercice


\exercice{5470, rouget, 2010/07/10}
\enonce{005470}{}
Calculer les intégrales suivantes ($a$, $b$ réels donnés, $p$ et $q$ entiers naturels donnés)

$$
\begin{array}{ll}
1)\;\int_{1/a}^{a}\frac{\ln x}{x^2+1}\;(0<a)&2)\;\int_{0}^{\pi}{2}\cos(px)\cos(qx)\;dx\;\mbox{et}\;\int_{0}^{\pi}{2}\cos(px)\sin(qx)\;dx\;\mbox{et}\;\int_{0}^{\pi}{2}\sin(px)\sin(qx)\;dx\\
3)\;\int_{a}^{b}\sqrt{(x-a)(b-x)}\;dx&4)\;\int_{-2}^{2}(|x-1|+|x|+|x+1|+|x+2|)\;dx\\
5)\;\int_{1/2}^{2}\left(1+\frac{1}{x^2}\right)\Arctan x\;dx&6)\;\int_{-1}^{1}\sqrt{1+|x(1-x)|}\;dx\\  
7)\int_{0}^{\pi}\;\frac{x\sin x}{1+\cos^2x}&8)\;\int_{1}^{x}(\ln t)^n\;dt\;(n\in\Nn^*) 
\end{array}
$$
\finenonce{005470}


\finexercice


\exercice{6865, bodin, 2012/04/13}
\video{qdUaqxk3B2s}
\enonce{006865}{}
Calculer les primitives suivantes par changement de variable.
\begin{enumerate}
  \item $\int (\cos x) ^{1234} \sin x \, d x$

  \item $\int \frac 1{x\ln x} \, dx$

  \item $\int \frac 1{3+\exp \left( -x\right)}dx$

  \item $\int \frac{1}{\sqrt{4x-x^2}}dx$
\end{enumerate}
\finenonce{006865} 


\finexercice

\section{127.06 Intégration par parties}


\exercice{815, bodin, 1998/09/01}

\enonce{000815}{}
Calculer les primitives suivantes :
$$
\int e^x\cos x dx \quad ; \quad
\int \frac{\ln x}{x^n}dx \ n\in\Nn \quad ; \quad
\int x \text{Arctan}\,x dx \quad ; \quad
\int (x^2+x+1)e^x dx.
$$
\finenonce{000815}



\finexercice



\exercice{816, ridde, 1999/11/01}

\enonce{000816}{}
Soit $I_{n} = \int_{0}^{1}{ (1-t^2)^n dt}$.
\begin{enumerate}
\item \'Etablir une relation de r\'ecurrence entre $I_{n}$ et $I_{n + 1}$.
\item Calculer $I_{n}$.
\item En d\'eduire $\sum\limits_{k = 0}^{n}{\frac{ (-1)^k}{2k + 1}C_{n}^{k}}$.
\end{enumerate}
\finenonce{000816}



\finexercice



\exercice{817, ridde, 1999/11/01}

\enonce{000817}{}
Soit $f \in C^{2} ([a, b], \Rr)$.
\begin{enumerate}
\item  Montrer que $\int_{a}^{b}f (t)dt = \frac{b-a}2 (f (a) + f (b)) +
  \frac 12 \int_a^b f'' (x) (a-x) (b-x)dx$.
\item En d\'eduire un encadrement de  $\int_{a}^{b}f (t)dt $ si $\forall x\in [a, b]
 \,\,  m \leq f'' (x) \leq M$.
\end{enumerate}
\finenonce{000817}



\finexercice



\exercice{818, ridde, 1999/11/01}

\enonce{000818}{Int\'egrales de Wallis}
  Soit $I_{n} = \int_0^{\frac{\pi}2} \sin ^n tdt$.
\begin{enumerate}
\item \'Etablir une relation de r\'ecurrence entre $I_{n}$ et $I_{n + 2}$.
\item En d\'eduire $I_{2p}$ et $I_{2p + 1}$.
\item Montrer que $ (I_{n})_{n \in \Nn}$ est d\'ecroissante et strictement positive.
\item En d\'eduire que $I_{n} \sim I_{n + 1}$.
\item Calculer $nI_{n}I_{n + 1}$.
\item Donner alors un \'equivalent simple de $I_{n}$.
\end{enumerate}
\finenonce{000818}



\finexercice



\exercice{819, ridde, 1999/11/01}

\enonce{000819}{}
Soit $I_{n} = \int_0^1 \frac{x^n}{1 + x}dx$.
\begin{enumerate}
\item En majorant la fonction int\'egr\'ee, montrer que $ (I_{n})_{n \in \Nn} \rightarrow 0$.
\item  Calculer $I_n + I_{n + 1}$.
\item D\'eterminer $\lim \limits_{n \rightarrow  + \infty} (\sum\limits_{k = 1}^n \frac{
 (-1)^{k + 1}}k)$.
\end{enumerate}
\finenonce{000819}



\finexercice



\exercice{820, gourio, 2001/09/01}

\enonce{000820}{}
Calculer par r\'{e}currence :
$$I_{n}=\int_{0}^{\frac{\pi }{4}}\frac{du}{\cos ^{n}u}. $$
\finenonce{000820}



\finexercice



\exercice{821, gourio, 2001/09/01}

\enonce{000821}{}
Calculer par r\'{e}currence :
$$J_{n}=\int_{1}^{e}\log (u)^{n}du. $$
\finenonce{000821}



\finexercice



\exercice{2336, matexo1, 2002/02/01}
\enonce{002336}{}
Pour tous $n$, $p$ dans $\Nn$, on d\'efinit
$$J_{n,p} =\int_0^{\pi /2} \sin^n t \cos^p t \,dt.$$
Trouver des relations de r\'ecurrence liant
$J_{n,p}$ et $J_{n,p-2}$, ainsi que $J_{n,p}$ et $J_{n-2,p}$. En
d\'eduire la valeur de $J_{n,p}$.
\finenonce{002336}


\finexercice


\exercice{6864, bodin, 2012/04/13}
\video{15IrPAzKwzc}
\enonce{006864}{}
Calculer les primitives suivantes par intégration par parties.
\begin{enumerate}
  \item $\int x^2 \ln x \, dx$

  \item $\int x \arctan x \, dx$

  \item $\int \ln x \, dx$ \quad  puis \quad  $\int (\ln x)^2 \, dx$

  \item $\int \cos x\exp x \, dx$
\end{enumerate}
\finenonce{006864} 


\finexercice

\section{127.07 Polynôme en sin, cos ou en sh, ch}


\exercice{822, bodin, 1998/09/01}

\enonce{000822}{}
Calculer les primitives suivantes :
$$
\int (\cos x \cos 2x + \sin x \sin 3x)dx \quad ; \quad
\int \cos x \sin^4x dx \quad ; \quad
\int \cos^6x dx \quad ;
$$

$$
\int \sin^3x \cos x dx \quad ; \quad
\int \sin^4x dx \quad ; \quad
\int \sin^3x \cos^2x dx \quad ;
$$

$$
\int \ch^2 x \sh^2 x dx \quad ; \quad
\int \sh x \ch^3 x dx \quad ; \quad
\int \ch x \sh^3x dx.
$$
\finenonce{000822}



\finexercice



\exercice{823, ridde, 1999/11/01}

\enonce{000823}{}
D\'eterminer les intervalles d'\'etude et calculer les primitives des fonctions :
$$ x\cos^2 x$$
$$\cos (2x) \cos^2 x$$
\finenonce{000823}



\finexercice



\exercice{2090, bodin, 2008/02/04}
\enonce{002090}{}
Calculer les primitives suivantes, en pr\'ecisant si n\'ecessaire les
intervalles de validit\'e des calculs:
$$
\begin{array}{llll}
\textbf{a)~} \displaystyle \int \sin ^8x\cos ^3x d x \quad & \textbf{b)~} \displaystyle \int
\cos^4x d x \quad  & \textbf{c)~} \displaystyle \int \cos ^{2003}x\sin x d x & \textbf{d)~}
\displaystyle \int \frac {1}{2+\sin x+\cos x}d x \\
\textbf{e)~}\displaystyle \int \frac 1{\sin x}d x & \textbf{f)~} \displaystyle \int \frac
1{\cos x}d x & \textbf{g)~} \displaystyle \int \frac{3-\sin x}{2\cos x+3\tan x}d x
\quad & \textbf{h)~} \displaystyle \int \frac 1{7+\tan x}d x
\end{array}$$
\finenonce{002090} 
\finexercice


\exercice{2096, bodin, 2008/02/04}
\video{Rp5pIHte82w}
\enonce{002096}{Intégrales de Wallis}
Soit $\displaystyle I_n=\int_0^{\frac \pi 2}(\sin x)^n \, d x$ \ \ pour $n\in \N$.
\begin{enumerate}
\item Montrer que $I_{n+2}=\frac{n+1}{n+2}I_n$. Expliciter $I_n$. En déduire $\int_{-1}^1\left( 1-x^2\right) ^n d x$.
\item Montrer que $\left( I_n\right) _n$ est positive décroissante. Montrer que $I_n\sim I_{n+1}$
\item Simplifier $I_n \cdot I_{n+1}$. Montrer que $I_n\sim \sqrt{%
\frac \pi {2n}}$. En déduire 
$\frac{1 \cdot 3 \cdots \left( 2n+1\right) }{2 \cdot 4 \cdots \left( 2n\right) }\sim 2\sqrt{\frac n \pi }$.
\end{enumerate}
\finenonce{002096} 


\finexercice



\exercice{5474, rouget, 2010/07/10}
\enonce{005474}{Intégrales de \textsc{Wallis}}
Pour $n$ entier naturel, on pose $W_n=\int_{0}^{\pi/2}\sin^nx\;dx$.
\begin{enumerate} 
\item  Calculer $W_0$ et $W_1$. Déterminer une relation entre $W_n$ et $W_{n+2}$ et en déduire $W_{2n}$ et $W_{2n+1}$ en fonction de $n$.
\item  Etudier les variations de la suite $(W_n)$ et en déduire $\lim_{n\rightarrow +\infty}\frac{W_{n+1}}{W_n}$.
\item  Montrer que la suite $(nW_nW_{n-1})_{n\in\Nn^*}$ est constante. En déduire $\lim_{n\rightarrow +\infty}W_n$, puis un équivalent simple de $W_n$. En écrivant $\int_{0}^{\pi/2}=\int_{0}^{\alpha}+\int_{\alpha}^{\pi}{2}$, retrouver directement $\lim_{n\rightarrow +\infty}W_n$.
\item  Montrer que $\lim_{n\rightarrow +\infty}n\left(\frac{1.3....(2n-1)}{2.4....(2n)}\right)^2=\frac{1}{\pi}$. (Formule de \textsc{Wallis})
\end{enumerate}
\finenonce{005474}



\finexercice


\exercice{5475, rouget, 2010/07/10}
\enonce{005475}{}
Pour $n$ entier naturel, on pose $In=\int_{0}^{\pi/4}\tan^nx\;dx$.
\begin{enumerate} 
\item  Calculer $I_0$ et $I_1$. Trouver une relation entre $I_n$ et $I_{n+2}$. En déduire $I_n$ en fonction de $n$.
\item  Montrer que $I_n$ tend vers $0$ quand $n$ tend vers $+\infty$, et en déduire les limites des suites $(u_n)$ et $(v_n)$ définies par~:~$u_n=\sum_{k=1}^{n}\frac{(-1)^{k-1}}{k}$ ($n\in\Nn^*$) et $v_n=\sum_{k=1}^{n}\frac{(-1)^{k-1}}{2k-1}$.
\end{enumerate}  
\finenonce{005475}


\finexercice

\section{127.08 Fraction rationnelle}


\exercice{824, cousquer, 2003/10/01}

\enonce{000824}{}
Décomposer les fractions rationnelles suivantes~; en
calculer les primitives.
\begin{enumerate}

\item  $\displaystyle{1 \over a^2+x^2}$.

\item  ${1 \over{(1+x^2)}^2}$.

\item  $\displaystyle{x^3 \over x^2-4}$.

\item  $\displaystyle{4x \over{(x-2)}^2}$.

\item  $\displaystyle{1 \over x^2+x+1}$.

\item  $\displaystyle{1 \over{(t^2+2t-1)}^2}$.

\item  $\displaystyle{3t+1 \over{(t^2-2t+10)}^2}$.

\item  $\displaystyle{3t+1 \over{t^2-2t+10}}$.

\item  $\displaystyle{ 1 \over{t^3+1}}$.

\item  $\displaystyle{x^3+2 \over{(x+1)}^2}$.

\item  $\displaystyle{x+1 \over{x{(x-2)}^2}}$.

\item  $\displaystyle{(x^2-1)(x^3+3) \over2x+2x^2}$.

\item  $\displaystyle{x^2 \over{{(x^2+3)}^3 (x+1)}}$.

\item  $\displaystyle{x^7+x^3-4x-1 \over x{(x^2+1)}^2}$.

\item  $\displaystyle{3x^4-9x^3+12x^2-11x+7\over(x-1)^3(x^2+1)}$.
\end{enumerate}



\finenonce{000824}



\finexercice



\exercice{825, cousquer, 2003/10/01}

\enonce{000825}{}
Calculer les intégrales de fractions rationnelles suivantes.
\begin{enumerate}


\item  $\displaystyle \int_0^1 \frac{dx}{x^2+2}$.

\item  $\displaystyle \int_{-1/2}^{1/2} \frac{dx}{1-x^2}$.

\item  $\displaystyle \int_2^3 \frac{2x+1}{x^2+x-3}\,dx$.

\item  $\displaystyle \int_0^2 \frac{x\,dx}{x^4+16}$.

\item  $\displaystyle \int_0^3 \frac{x^4+6x^3-5x^2+3x-7}{(x-4)^3}\,dx$.

\item  $\displaystyle \int_{-2}^0 \frac{dx}{x^3-7x+6}$.

\item  $\displaystyle \int_{-1}^1 \frac{2x^4+3x^3+5x^2+17x+30}{x^3+8}\,dx$.

\item  $\displaystyle \int_2^3 \frac{4x^2}{x^4-1}\,dx$.

\item  $\displaystyle \int_{-1}^0 \frac{x^3+2x+1}{x^3-3x+2}\,dx$.

\item  $\displaystyle \int_1^2 \frac{2x^8+5x^6-12x^5+30x^4+36x^2+24}
    {x^4(x^2+2)^3}\,dx$.

\item  $\displaystyle \int_0^a \frac{-2x^2+6x+7}
    {x^4+5x^2+4}\,dx$ pour $a\in \mathbb{R}$. Y a-t-il une 
    limite quand $a\to+\infty$~?

\item  $\displaystyle \int_0^2 \frac{dx}{x^4+1}$.

\end{enumerate}

\finenonce{000825}



\finexercice



\exercice{826, bodin, 1998/09/01}

\enonce{000826}{}
Calculer les primitives suivantes :
$$
\int \frac{x^4+1}{x(x-1)^3}dx  \quad ; \quad
\int \frac{dx}{(x^4+1)^2}  \quad ; \quad
\int \frac{x dx}{x^4+x^2+1}  \quad ; \quad
\int \frac{dx}{(x-1)(x^2-2x-2)^2}.
$$
\finenonce{000826}



\finexercice



\exercice{827, ridde, 1999/11/01}

\enonce{000827}{}
D\'eterminer les intervalles d'\'etude et calculer les primitives des fonctions :
$$ \frac 1{ (x + 2) (x^2 + 2x + 5)}$$
$$ \frac{2x}{ (1-x + x^2)^2}$$
$$ \frac{x^2}{ (x-1)^2 (x^2 + 4)}$$
$$ \frac 1{ (1 + x^3)^3}$$
\finenonce{000827}



\finexercice



\exercice{2097, bodin, 2008/02/04}
\video{JSvF3eC5EuA}
\enonce{002097}{}
Soit $\displaystyle I_{n} = \int_0^1 \frac{x^n}{1 + x}d x$.
\begin{enumerate}
\item En majorant la fonction int\'egr\'ee, montrer que
$\lim_{n\to +\infty} I_{n}=0$.
\item  Calculer $I_n + I_{n + 1}$.
\item D\'eterminer $\displaystyle \lim_{n \rightarrow  + \infty} \left(\sum_{k = 1}^n \frac{
 (-1)^{k + 1}}k\right)$.
\end{enumerate}
\finenonce{002097} 



\finexercice


\exercice{4263, quercia, 2010/03/12}
\enonce{004263}{Fractions rationnelles}
                      %+--------------------------+
                      %|  Fractions rationnelles  |
                      %+--------------------------+
$$\begin{array}{ll}
\frac1{x^3-1}                          & \frac13 \ln|x-1| -\frac16 \ln(x^2+x+1) - \frac1{\sqrt3} \Arctan\Bigl(\frac{2x+1}{\sqrt3} \Bigr) \cr
\frac1{(x^3-1)^2}                      & -\frac29 \ln|x-1| +\frac19 \ln(x^2+x+1) + \frac2{3\sqrt3} \Arctan\Bigl(\frac{2x+1}{\sqrt3} \Bigr) -\frac x{3(x^3-1)} \cr
\frac1{x^3(1+x^3)}                     & -\frac1{2x^2} + \frac16 \ln\Bigl[\frac{x^2-x+1}{(x+1)^2} \Bigr] -\frac1{\sqrt3} \Arctan\bigl[ \frac{2x-1}{\sqrt3} \bigr] \cr
\frac{x^2+x+1}{(x^2-1)^2}              & -\frac3{4(x-1)} - \frac1{4(x+1)} \cr
\frac1{1+x^4}                          & \frac1{4\sqrt2} \ln\Bigl[ \frac{1+x\sqrt2+x^2}{1-x\sqrt2+x^2} \Bigr] + \frac1{2\sqrt2} \bigl[ \Arctan(1+x\sqrt2) - \Arctan(1-x\sqrt2) \bigr] \cr
\frac{x^2}{1+x^4}                      & \frac1{4\sqrt2} \ln\Bigl[ \frac{1-x\sqrt2+x^2}{1+x\sqrt2+x^2} \Bigr] + \frac1{2\sqrt2} \bigl[ \Arctan(1+x\sqrt2) - \Arctan(1-x\sqrt2) \bigr] \cr
\frac{x}{(x^4+1)^2}                    & \frac{\Arctan{x^2}}4 + \frac {x^2}{4(x^4+1)} \cr
\frac{x^2+x+1}{x^3-2x-4}               & \frac7{10} \ln|x-2| + \frac3{20} \ln(x^2+2x+2) -\frac1{10} \Arctan(x+1) \cr
\frac{x^2-4}{x^6-2x^4+x^2}             & \frac4x + \frac{3x}{2(x^2-1)} + \frac{11}4 \ln\Bigl|\frac{x-1}{x+1} \Bigr| \cr
\frac1{x^{20}-1}                       & \frac1{10} \sum_{k=1}^9 \Bigl[ \frac12 \cos k\alpha \ln( x^2-2x\cos k\alpha +1 ) - \sin k\alpha \Arctan\bigl( \frac{x-\cos k\alpha}{\sin k\alpha} \bigr)\Bigr] + \frac1{20} \ln\Bigl|\frac{x-1}{x+1} \Bigr|,\quad \alpha = \frac\pi{10} \cr
\frac1{(x-a)^n(x-b)}                   & \frac1{(b-a)^n} \ln\Bigl|\frac{x-b}{x-a} \Bigr| + \sum_{k=1}^{n-1} \frac1{k(b-a)^{n-k}(x-a)^k} \cr
\end{array}$$
\finenonce{004263}



\finexercice

\section{127.09 Fraction rationnelle en sin, cos ou en sh, ch}


\exercice{828, bodin, 1998/09/01}

\enonce{000828}{}
Calculer les primitives suivantes :
$$
\int \frac{\cos^3 x}{\sin^5x}dx  \quad ; \quad
\int \frac{\sin^3x}{1+\cos x}dx  \quad ; \quad
\int \frac{dx}{\cos^4x+\sin^4x}  \quad ; \quad
\int \frac{\cos x}{1+\sin 2x}dx \quad ;
$$

$$
\int \frac{\tan x -\tan a}{\tan x + \tan a}dx  \quad ; \quad
\int \frac{\sh x \ch x}{\sh^4x+\ch^4x}dx.
$$
\finenonce{000828}



\finexercice



\exercice{829, ridde, 1999/11/01}

\enonce{000829}{}
D\'eterminer les intervalles d'\'etude et calculer les primitives des fonctions :
$$\frac{\cos^3 x}{\sin x}$$
$$ \frac 1{1 + \tan x}$$
$$\frac 1 {\text{th}^2 x}$$
\finenonce{000829}



\finexercice



\exercice{2089, bodin, 2008/02/04}
\enonce{002089}{}
Calculer les primitives suivantes:
$$\int \frac{\sin x}{\sin x+\cos x}d x \quad \mbox{ et } \quad \int \frac{\cos x}{\sin x+\cos x}d x.$$
\finenonce{002089} 
\finexercice


\exercice{2095, bodin, 2008/02/04}
\video{1uLeRF-liOk}
\enonce{002095}{}
Calculer les intégrales suivantes :
$$\int_0^{\frac \pi 2}\frac 1{1+\sin x}d x \quad \mbox{ et } \quad \int_0^{\frac \pi 2}\frac{\sin
x}{1+\sin x}d x.$$
\finenonce{002095} 


\finexercice



\exercice{4264, quercia, 2010/03/12}
\enonce{004264}{Fonctions trigonométriques}
                    %+------------------------------+
                    %|  Fonctions trigonométriques  |
                    %+------------------------------+
$$\begin{array}{ll}
\frac1{\sin x\sin 4x}                  & -\frac1{4\sin x} + \frac18 \ln\Bigl| \frac{1-\sin x}{1+\sin x} \Bigr| - \frac1{2\sqrt2} \ln\biggl| \frac{1-\sqrt2\sin x}{1+\sqrt2\sin x} \biggr| \cr
\frac{\tan x}{1+\tan x}                & \frac x2 - \frac12 \ln|\cos x + \sin x| \cr
\cos x\sqrt{\cos 2x}                    & \frac{\sin x\sqrt{\cos 2x}}2 + \frac1{2\sqrt2} \Arcsin(\sqrt2\sin x) \cr
\frac1{\sin x + \sin2x}                & \frac16 \ln(1-\cos x) + \frac12 \ln(1+\cos x)- \frac23 \ln|1+2\cos x| \cr
\frac1{\cos x\cos 2x}                  & \sqrt2\Argth(\sqrt2\sin x) -\Argth(\sin x) \cr
\frac1{\sin x\sqrt{\sin x(1+\sin x)}}  & -2\sqrt{\frac{1-\sin x}{\sin x} } + \sqrt2\Arctan\sqrt{\frac{1-\sin x}{2\sin x} } \quad(\text{poser } u = 1/\sin x) \cr
\frac{a\sin x}{\cos x\sqrt{\cos^2x-a^2\sin^2x}} & -\Arctan\biggl(\frac{\sqrt{\cos^2x-a^2\sin^2x}}a \biggr) \cr
\end{array}$$
\finenonce{004264}



\finexercice
\section{127.10 Intégrale abélienne}


\exercice{830, bodin, 1998/09/01}

\enonce{000830}{}
Calculer les primitives suivantes :
$$
\int \frac{dx}{x+\sqrt{x-1}} \quad ; \quad
\int \frac{dx}{x\sqrt{x^2+x+1}}  \quad ; \quad
\int \frac{x}{\sqrt{9+4x^4}}dx  \quad ;
$$

$$
\int \frac{\sqrt[3]{x+1}-\sqrt{x+1}}{x+2}dx  \quad ; \quad
\int \frac{x+1}{\sqrt{-4x^2+4x+1}}dx.
$$
\finenonce{000830}



\finexercice



\exercice{831, ridde, 1999/11/01}

\enonce{000831}{}
D\'eterminer les intervalles d'\'etude et calculer les primitives des fonctions :
$$ \frac{8x-3}{\sqrt{12x-4x^2-5}}$$
$$ \sqrt{x^2-1}$$
$$ \frac{x \sqrt x}{x^2-5x + 4}$$
\finenonce{000831}



\finexercice



\exercice{4265, quercia, 2010/03/12}
\enonce{004265}{Radicaux}

                             %+------------+
                             %|  Radicaux  |
                             %+------------+
$$\begin{array}{ll}
\frac{x+1}{\sqrt{x^2-3x+2}}            & \sqrt{x^2-3x+2} + \frac52 \ln\bigl| 2x-3 + 2\sqrt{x^2-3x+2} \bigr| \cr
\frac{4x-3}{\sqrt{-4x^2+12x-5}}        & -\sqrt{-4x^2+12x-5} + \frac32 \Arcsin(x-3/2) \cr
\frac1{2x-x^2+\sqrt{2x-x^2}}           & \frac{1-\sqrt{2x-x^2}}{x-1} \cr
\frac1{2+\sqrt{1+x}+\sqrt{3-x}}        & \sqrt{1+x} - \sqrt{3-x} - \Arcsin\Bigl(\frac{x-1}2 \Bigr) \quad(\text{poser } x = 1+2\cos\varphi) \cr
\frac{2+\sqrt{x+3}}{1+\sqrt{x+4}}      & (\sqrt{x+3}+4)(\sqrt{x+4}-2) -4\ln(1+\sqrt{x+4}\,) +\ln(\sqrt{x+3}+\sqrt{x+4}\,) \cr
 x+\sqrt{a^2+x^2}                       & \frac{(x+\sqrt{a^2+x^2}\,)^2}4 + \frac{a^2}2 \ln( x+ \sqrt{a^2+x^2}\, ) \cr
(x+\sqrt{a^2+x^2}\,)^n                  & \frac{(x+\sqrt{a^2+x^2}\,)^{n+1}}{2(n+1)} + a^2\frac{(x+\sqrt{a^2+x^2}\,)^{n-1}}{2(n-1)} \quad (n\ne1) \cr
\frac1{\sqrt[3]{1+x^3}}                   & \frac16 \ln\Bigl[ \frac{u^2+u+1}{(u-1)^2} \Bigr] - \frac1{\sqrt3} \Arctan \frac{2u+1}{\sqrt3} , u = \sqrt[3]{1+1/x^3}\quad (\text{poser } v=1/x^3) \cr
\end{array}$$
\finenonce{004265}



\finexercice

\section{127.11 Primitives diverses}


\exercice{832, cousquer, 2003/10/01}

\enonce{000832}{}
Calculer les primitives suivantes.
\begin{enumerate}

\item  $\displaystyle\int e^{\sin^2x} \sin2x\,dx$.

\item  $\displaystyle\int\cos^5t\,dt$ ; $\displaystyle\int\cosh^3t\,dt$; 
$\displaystyle\int\cos^4t\,dt$ ; $\displaystyle\int\sinh^4t\,dt$.

\item  $\displaystyle\int x^3 e^x\,dx$.

\item  $\displaystyle\int\ln x \,dx$ ; $\displaystyle\int x\ln x\,dx$; 
$\displaystyle\int\arcsin x \,dx$.

\item  $\displaystyle\int\cosh t \sin t \, dt$.

\item  $\displaystyle\int{dx \over\sin x}$.

\item  $\displaystyle\int\sqrt{a^2-x^2}\,dx$.

\item  $\displaystyle\int{e^{2x}\over\sqrt{e^x+1}}\,dx$.

\item  $\displaystyle\int e^{ax}\cos bx \,dx$~; 
$\displaystyle\int e^{ax}\sin bx \,dx$.

\item  $\displaystyle\int\sqrt{x \over(1-x)^3}\,dx$\quad pour\quad$0<x<1$.

\item  $\displaystyle\int{x^2 \over\sqrt{1-x^2}}\,dx$.

\item  $\displaystyle\int{dx \over\cos x+2\sin x +3}$.

\item  $\displaystyle\int{\sqrt x\,dx \over\sqrt{a^3-x^3}}$ \quad avec 
\quad $0<x<a$.

\item  $\displaystyle\int{\cosh x\over\cosh x+\sinh x}\,dx$.
\end{enumerate}



\finenonce{000832}



\finexercice



\exercice{833, bodin, 1998/09/01}

\enonce{000833}{}
Calculer les primitives suivantes :
$$
\int \frac {dx}{\ch x\sqrt{\ch 2x}} \quad ; \quad
\int \frac{x}{\cos^2x}dx \quad ; \quad
\int \frac{1+\cos 2x}{1-\tan^2x}dx \quad ;
$$

$$
\int \frac {\sin ax + \cos bx}{e^x}dx \quad ; \quad
\int \frac {x(2+\cos x)}{\sin^2x}dx.
$$
\finenonce{000833}



\finexercice



\exercice{834, ridde, 1999/11/01}

\enonce{000834}{}
D\'eterminer les intervalles d'\'etude et calculer les primitives des fonctions :
$$ \text{ch}x \sin (2x)$$
$$ \frac 1{\sqrt{2 + \tan ^2 x}}$$
$$ (x^2 + 2x + 2) \cos (2x)$$
$$ x^2 \cos x \text{ et } x^2 \sin x \text{ en utilisant les complexes}$$
$$\frac 1 { (x^2-1)^3} \text{ et } \frac 1{ (x^2-1)^2}$$
$$ \frac {\sqrt{1 + x}} {x \sqrt{1-x}}$$
\finenonce{000834}



\finexercice



\exercice{835, ridde, 1999/11/01}

\enonce{000835}{}
Calculer $\int_{0}^{1}{\ln (1 + x^2)}$.
\finenonce{000835}



\finexercice



\exercice{836, ridde, 1999/11/01}

\enonce{000836}{}
D\'eterminer $\lim \limits_{n \rightarrow  + \infty} (\sum\limits_{k = 0}^n \frac{
 n}{n^2 + k^2})$.

\finenonce{000836}



\finexercice



\exercice{837, ridde, 1999/11/01}

\enonce{000837}{}
Calculer $\lim \limits_{n \rightarrow  + \infty} ( \frac{ (2n) ! }{n ! n^n})^{\frac 1n}$.

\finenonce{000837}



\finexercice



\exercice{838, ridde, 1999/11/01}

\enonce{000838}{}
Soient $I = \int_0^\pi x\cos^2xdx$ et $J =\int_0^\pi x\sin^2xdx $.
\begin{enumerate}
\item Calculer $I$ et $I + J$.
\item  En d\'eduire $J$.
\end{enumerate}
\finenonce{000838}



\finexercice



\exercice{839, ridde, 1999/11/01}

\enonce{000839}{}
Soit $a_{n} = \int_0^1 t^n e^t dt$.
\begin{enumerate}
\item Calculer $a_0, \ldots, a_4$.
\item Etudier la suite $ (a_n)_{n \in \Nn}$.
\end{enumerate}
\finenonce{000839}



\finexercice



\exercice{2088, bodin, 2008/02/04}
\enonce{002088}{}
Calculer les primitives suivantes, en pr\'ecisant si n\'ecessaire les
intervalles de validit\'e des calculs:
$$
\begin{array}{llll}
\textbf{a)~}\displaystyle\int \arctan x d x & \textbf{b)~}\displaystyle\int \tan ^2x d x &
\textbf{c)~}\displaystyle \int \frac 1{x\ln x} d x & \textbf{d)~}\displaystyle \int \frac
x{\sqrt{x+1}}d x \\
\textbf{e)~} \displaystyle\int \arcsin x d x & \textbf{f)~}\displaystyle \int \frac 1{3+\exp
  \left( -x\right)}d x \quad & \textbf{g)~} \displaystyle \int \frac{-1}{\sqrt{4x-x^2}}d
x & \textbf{h)~} \displaystyle \int \frac 1{x\sqrt{1-\ln ^2x}}d x\\
\textbf{i)~} \displaystyle \int \frac 1{\sqrt{1+\exp x}}d x \quad & \textbf{j)~} \displaystyle \int
\frac{x-1}{x^2+x+1} d x & \textbf{k)~} \displaystyle \int \frac{x+2}{x^2-3x-4}d x
\quad &
\textbf{l)~} \displaystyle\int \cos x\exp xd x 
\end{array}$$
\finenonce{002088} 
\finexercice


\exercice{2094, bodin, 2008/02/04}
\enonce{002094}{}
Calculer les int\'egales suivantes:
$$
\begin{array}{lll}
\textbf{a)~} \displaystyle \int_{0^{}}^1\frac{\arctan x}{1+x^2}d x & \textbf{b)~}
\displaystyle \int_{\frac 12}^2\left( 1+\frac 1{x^2}\right) \arctan x d x \quad &
\textbf{c)~} \displaystyle \int_0^{\frac{\pi}{2}}x\sin x d x \\
\textbf{d)~} \displaystyle \int_{-1}^1\left( \arccos x\right) ^2 d x \quad & \textbf{e)~}
\displaystyle \int_0^1\frac 1{\left( 1+x^2\right) ^2}d x & \textbf{f)~} \displaystyle \int_0^{\sqrt{3}}\frac{x^2}{\sqrt{4-x^2}}d x \\
 \textbf{g)~}\displaystyle \int_1^2x^2\ln xd x  & \textbf{h)~} \displaystyle \int_{-1}^1\frac
1{x^2+4x+7}d x & \textbf{i)~}\displaystyle \int_0^1\frac{%
3x+1}{\left( x+1\right) ^2}d x    
\end{array}$$

\finenonce{002094} 
\finexercice


\exercice{2319, matexo1, 2002/02/01}
\enonce{002319}{}
Calculer les primitives des fonctions suivantes\,:
\begin{eqnarray*}
t &\mapsto &t^2 \exp (t^3), \\ 
t &\mapsto &{\sin^3 t \over {1 + \cos^2 t}}, \\
t &\mapsto &{1 \over {1-t^2+2 \sqrt{1-t^2}}}, \\
t &\mapsto &{\sinh^2 t \over {\cosh t}}, \\
t &\mapsto &{\cos t \over {\cos 2t}}, \\
t &\mapsto &{1 \over {1+\tanh^2 t}}.
\end{eqnarray*}

\finenonce{002319}


\finexercice


\exercice{2320, matexo1, 2002/02/01}
\enonce{002320}{}
Calculer les primitives des fonctions suivantes\,:
\begin{eqnarray*}
t &\mapsto & {\sin t \over {\sin^2 t - \cos t }}, \\
t &\mapsto & {t^2 \over {(\cos t + t \sin t)^2 }}, \\
t &\mapsto & {\sqrt{1+\sqrt{1-t^2}} \over {\sqrt{1-t^2} }}, \\
t &\mapsto & {1 \over {t^8+t^4+1}}.
\end{eqnarray*}
\finenonce{002320}


\finexercice


\exercice{2325, matexo1, 2002/02/01}
\enonce{002325}{}
Calculer les primitives des fonctions suivantes\,:
\begin{eqnarray*}
t &\mapsto & \tan t, \\
t &\mapsto & \arg \sinh t, \\
t &\mapsto & {\tan^3 t \over {\cos^6 t }}, \\
t &\mapsto & {1 \over {\sqrt{t} + \sqrt[3]{t}}}, \\
t &\mapsto & {1 \over {2\sqrt {t} + \sqrt{t+2}}}, \\
t &\mapsto & t \sqrt[n]{1+t}, \\
t &\mapsto & \cosh^3 t, \\
t &\mapsto & {t^3 \over {(a^2-t^2)^{3/2} }}, \\
t &\mapsto & {\sqrt[3]{1+\sqrt[4]t} \over {\sqrt{t}}}, \\
t &\mapsto & \sqrt{t \sqrt{t \sqrt{t}}}, \\
t &\mapsto & { {1 \over {t^2}} \sqrt[4]{t+1 \over {t-1}}}.
\end{eqnarray*}
\finenonce{002325}


\finexercice


\exercice{2330, matexo1, 2002/02/01}
\enonce{002330}{}
{\bf Fonction Gamma - }
Pour tout $x>0$, on pose
$$ \Gamma (x) = \int_0^{+\infty} {t^{x-1} e^{-t}} dt $$
(on admettra que l'int\'egrale converge). Montrer que 
$\Gamma(x+1) = x \Gamma(x)$. Calculer la valeur de 
$\Gamma(1)$. En d\'eduire celle de $\Gamma(n)$, pour tout
entier $n>0$.

\item Soit $a>0$ un r\'eel, et $n>0$ un entier. Montrer que
$$ \int {dt \over (x^2 + a^2)^n} = {1 \over a^{2n-1}}
	\int \cos^{2n-2}\theta d\theta \qquad \mbox{\rm o\`u}\qquad
\theta = \arctan{x \over a}. $$
En d\'eduire la primitive de $\displaystyle{ x+4 \over (x^2+2x+2)^3 }$.

\item Soient $x$ et $y$ deux r\'eels v\'erifiant $1>y>x>0$. Calculer
$$ \lim_{x\to  0 \atop y\to  1} 
 \int_x^y {\ln t \over (1+t)\sqrt{1-t^2}} dt. $$

\item Soit $f$ une fonction continue et positive sur $[0, +\infty[$.
On pose pour tout $x>0$ et tout entier $n>0$
$$ u_n(x) = \left[ \int_0^x f(t)^n\ dt \right]^{1/n} $$
et
$$ M(x) = \sup_{t \in [0, x]} \left| f(t) \right|. $$

\begin{enumerate}
\item Montrer que $u_n(x)  \leq M(x) x^{1/n}$.
\item En utilisant la continuit\'e de $f$, montrer que, quel que 
soit $\varepsilon > 0$, il existe $\delta > 0$ tel que 
$u_n(x) \geq \delta^{1/n} [ M(x)-\varepsilon ]$.
\item En d\'eduire que 
$$\lim_{n \to +\infty} u_n(x) = M(x).$$
\end{enumerate}
\finenonce{002330}


\finexercice


\exercice{4266, quercia, 2010/03/12}
\enonce{004266}{Diverses primitives} 
                              %+----------+
                              %|  Divers  |
                              %+----------+

$$\begin{array}{ll}
x^k\ln x                                & \frac{x^{k+1}}{k+1} \Bigl(\ln x - \frac1{k+1} \Bigr) \cr
\ln(1+x^2)                              & x\ln(1+x^2) -2x + 2\Arctan x \cr
\frac{x^2+a}{x^2+1} \Arctan x          & \frac12 \bigl( (2x+(a-1)\Arctan x)\Arctan x - \ln(1+x^2)\bigr) \cr
\Bigl(1-\frac1x \Bigr)e^{1/x}          & xe^{1/x} \cr
\frac x{\cos^2x}                       & x\tan x + \ln|\cos x| \cr
\frac1{\sqrt{e^x-1}}                   & 2\Arctan\sqrt{e^x-1} \cr
\Arctan\sqrt{\frac{x+1}{x+3}}          & (x+2)\Arctan\sqrt{\frac{x+1}{x+3}} - \ln\bigl( \sqrt{x+1} + \sqrt{x+3}\, \bigr) \cr
\Arcsin\sqrt{\frac x{x+1}}             &  x\Arcsin\sqrt{\frac x{x+1}} - \sqrt x + \Arctan \sqrt x \cr
e^{\Arcsin x}                           & \frac{x+\sqrt{1-x^2}}2 e^{\Arcsin x} \cr
x(\cos^2x)e^{-x}                        & \frac{e^{-x}}{50} \bigl( (3-5x)\cos 2x + (4+10x)\sin 2x - 25(x+1) \bigr) \cr
(x^2+x+1)e^{2x}\cos x                   & (\frac{2x^2}5 + \frac{4x}{25} + \frac{39}{125} )e^{2x}\cos x + (\frac{x^2}5 - \frac{3x}{25} + \frac{27}{125} )e^{2x}\sin x \cr
\end{array}$$
\finenonce{004266}



\finexercice


\exercice{4267, quercia, 2010/03/12}
\enonce{004267}{Intégrales définies}
                          %ÚÄÄÄÄÄÄÄÄÄÄÄÄÄÄÄÄÄÄÄÄÄÄÄÄ¿
                          %³  Intégrales définies   ³
                          %ÀÄÄÄÄÄÄÄÄÄÄÄÄÄÄÄÄÄÄÄÄÄÄÄÄÙ

$ \int_{t=0}      ^{\pi/2} \cos^4t                          \,d t = \frac{3\pi}{16}$\par
$ \int_{t=-\pi/2} ^{\pi/2} \sin^2t\cos^3t                   \,d t = \frac4{15}$\par
$ \int_{t=0}      ^{\pi/2} t^2\cos t                        \,d t = \frac{\pi^2}4 - 2$\par
$ \int_{t=-\pi/2} ^{\pi/2} t^2\sin t\cos^2t                 \,d t = 0$\par
$ \int_{t=0}      ^{\pi/2} \frac{\sin t}{1+\cos^2t}        \,d t = \frac \pi4$\par
$ \int_{t=0}      ^{\pi/2} \frac{d t}{1 + \sin t}                = 1$\par
$ \int_{t=0}      ^{\pi/2} \frac{\sin^2t}{\sin t + \cos t} \,d t = -\frac1{\sqrt2} \ln(\sqrt2-1)$\par
$ \int_{t=0}      ^{\pi/2} \frac{\sin2t}{\sqrt{1-a\sin t}} \,d t = \frac{4(2-(a+2)\sqrt{1-a})}{3a^2}$\par
%

$ \int_{t=0}      ^1       t\ln t                           \,d t = -\frac14$\par
$ \int_{t=0}      ^1       \Arcsin t                        \,d t = \frac\pi2 - 1$\par
$ \int_{t=0}      ^3       \frac{2t}{(1+t^2)(3+t^2)}       \,d t = \frac12 \ln\frac52$\par
$ \int_{t=0}      ^1       \frac{t^2\Arctan t}{1+t^2}      \,d t = \frac\pi4 - \frac{\pi^2}{32} - \ln\sqrt2$\par
$ \int_{t=0}      ^{\ln2}  \sqrt{e^t-1}                     \,d t = 2-\frac\pi2$\par
$ \int_{t=4}      ^9       \frac{d t}{\sqrt t - 1}               = 2 + 2\ln2$\par
$ \int_{t=0}      ^1       \frac{te^t}{\sqrt{e^t+1}}       \,d t = 4\sqrt2 - 2\sqrt{e+1} + 4\ln\Bigl[ \frac{\sqrt{e+1}+1}{\sqrt2+1} \Bigr] - 2$\par
$ \int_{t=0}      ^1       \frac{\ln(1-a^2t^2)}{t^2}       \,d t = a\ln\left|\frac{1-a}{1+a}\right| -\ln(1-a^2)$\par
%

$ \int_{t=0}      ^1       \frac{d t}{2+\sqrt{1-t^2}}            = \frac\pi6 (3 - \frac4{\sqrt3})$\par
$ \int_{t=-1}     ^1       \frac{d t}{t+\sqrt{t^2+1}}            = \ln(1+\sqrt2) + \sqrt2$\par
$ \int_{t=-1}     ^1       \sqrt{1+t^2}                     \,d t = \ln(1+\sqrt2) + \sqrt2$\par
\finenonce{004267}



\finexercice


\exercice{5450, rouget, 2010/07/10}
\enonce{005450}{***}
Etude complète de la fonction $f(x)=\frac{1}{x-1}\int_{1}^{x}\frac{t^2}{\sqrt{1+t^8}}\;dt$. 
\finenonce{005450}


\finexercice


\exercice{5451, rouget, 2010/07/10}
\enonce{005451}{***}
Pour $x$ réel, on pose $f(x)=e^{-x^2}\int_{0}^{x}e^{t^2}\;dt$.
\begin{enumerate}  
\item  Montrer que $f$ est impaire et de classe $C^\infty$ sur $\Rr$. 
\item  Montrer que $f$ est solution de l'équation différentielle $y'+2xy=1$.
\item  Montrer que $\lim_{x\rightarrow +\infty}2xf(x)=1$.
\item  Soit $g(x)=\frac{e^{x^2}}{2x}f'(x)$. Montrer que $g$ est strictement décroissante sur $]0,+\infty[$ et que $g$ admet sur $]0,+\infty[$ un unique zéro noté $x_0$ vérifiant de plus $0<x_0<1$.
\item  Dresser le tableau de variations de $f$.
\end{enumerate}
\finenonce{005451}


\finexercice


\exercice{5465, rouget, 2010/07/10}
\enonce{005465}{****}
Soit $f(t)=\frac{t^2}{e^t-1}$ si $t\neq0$ et $0$ si $t=0$.
\begin{enumerate}
\item  Vérifier que $f$ est continue sur $\Rr$.
\item  Soit $F(x)=\int_{0}^{x}f(t)\;dt$. Montrer que $F$ a une limite réelle $\ell$ quand $x$ tend vers $+\infty$ puis que $\ell=2\lim_{n\rightarrow +\infty}\sum_{k=1}^{n}\frac{1}{k^3}$.
\end{enumerate}
\finenonce{005465}


\finexercice


\exercice{5466, rouget, 2010/07/10}
\enonce{005466}{}
Calculer les primitives des fonctions suivantes en précisant le ou les intervalles considérés~:

$$
\begin{array}{lllll}
1)\;\frac{1}{x^3+1}&2)\;\frac{x^2}{x^3+1}&3)\;\frac{x^5}{x^3-x^2-x+1}&4)\;\frac{1-x}{(x^2+x+1)^5}&5)\;\frac{1}{x(x^2+1)^2}\\
6)\;\frac{x^2+x}{x^6+1}&7)\;\frac{1}{x^4+1}&8)\;\frac{1}{(x^4+1)^2}&9)\;\frac{1}{x^8+x^4+1}&10)\;\frac{x}{(x^4+1)^3}\\
11)\;\frac{1}{(x+1)^7-x^7-1}
\end{array}
$$ 
\finenonce{005466}


\finexercice


\exercice{5467, rouget, 2010/07/10}
\enonce{005467}{}
Calculer les primitives des fonctions suivantes en précisant le ou les intervalles considérés~:
$$
\begin{array}{lllll}
1)\;\frac{1}{\cos x}\;\mbox{et}\;\frac{1}{\ch x}&2)\;\frac{1}{\sin x}\;\mbox{et}\;\frac{1}{\sh x}&3)\;\frac{1}{\tan x}\;\mbox{et}\;\frac{1}{\tanh x}&4)\;\frac{\sin^2(x/2)}{x-\sin x}&5)\;\frac{1}{2+\sin^2x}\\
6)\;\frac{\cos x}{\cos x+\sin x}&7)\;\frac{\cos(3x)}{\sin x+\sin(3x)}&8)\;\frac{1}{\cos^4x+\sin^4x}&9)\;\frac{\sin x\sin(2x)}{\sin^4x+\cos^4x+1}&10)\;\frac{\tan x}{1+\sin(3x)}\\
11)\;\frac{\cos x+2\sin x}{\sin x-\cos x}&12)\;\frac{\sin x}{\cos(3x)}&13)\;\frac{1}{\alpha\cos^2x+\beta\sin^2x}&14)\;\frac{\ch^3 x}{1+\sh x}&15)\;\sqrt{\ch x-1}\\
16)\;\frac{\tanh x}{1+\ch x}&17)\;\frac{1}{\sh^5x}&18)\frac{1}{1-\ch x}
\end{array}
$$  
\finenonce{005467}


\finexercice


\exercice{5468, rouget, 2010/07/10}
\enonce{005468}{}
Calculer les primitives des fonctions suivantes en précisant le ou les intervalles considérés~:

$$
\begin{array}{lllll}
1)\frac{1}{\sqrt{x^2+2x+5}}\;\mbox{et}\;\sqrt{x^2+2x+5}&2)\;\frac{1}{\sqrt{2x-x^2}}&3)\;\frac{\sqrt{1+x^6}}{x}&4)\;
\frac{1}{\sqrt{1+x}+\sqrt{1-x}}&5)\;\sqrt{\frac{x+1}{x-1}}\\
6)\;\frac{x^2+1}{x\sqrt{x^4-x^2+1}}&7)\;\sqrt{\frac{1-\sqrt{x}}{\sqrt{x}}}&8)\;\frac{1}{1+\sqrt{1+x^2}}&9)\;\frac{\sqrt[3]{x^3+1}}{x^2}\;\mbox{et}\;\frac{1}{\sqrt[3]{x^3+1}}&\;\\
10)\;\frac{1}{\sqrt{x+1}+\sqrt[3]{x+1}}
\end{array}
$$    
\finenonce{005468}


\finexercice


\exercice{5469, rouget, 2010/07/10}
\enonce{005469}{}
Calculer les primitives des fonctions suivantes en précisant le ou les intervalles considérés~:
$$
\begin{array}{lllll}
1)\;\frac{1}{x\ln x}&2)\;\Arcsin x&3)\;\Arctan x&4)\;\Arccos x&5)\;\Argsh x\\
6)\;\Argch x&7)\;\Argth x&8)\;\ln(1+x^2)&9)\;e^{\Arccos x}&10)\;\cos x\ln(1+\cos x)\\
11)\;\frac{\Arctan x}{\sqrt{x}}&12)\;\frac{xe^x}{(x+1)^2}&13)\;(\frac{x}{e})^x\ln x&14)\;x^n\ln x\;(n\in\Nn)&15)\;e^{ax}\cos(\alpha x)\;((a,\alpha)\in(\Rr^*)^2)\\
16)\;\sin(\ln x)\;\mbox{et}\;\cos(\ln x)&17)\;\frac{\sqrt{x^n+1}}{x}&18)\;x^2e^x\sin x
\end{array}
$$
\finenonce{005469}


\finexercice


\exercice{5471, rouget, 2010/07/10}
\enonce{005471}{}
Condition nécessaire et suffisante sur $a$, $b$, $c$ et $d$ pour que les primitives de $\frac{(x-a)(x-b)}{x-c)^2(x-d)^2}$ soient rationnelles ($a$, $b$, $c$ et $d$ réels donnés).
\finenonce{005471}


\finexercice


\exercice{6866, bodin, 2012/04/13}
\video{0bvIaJVZNwY}
\enonce{006866}{}
Calculer les primitives suivantes, en précisant si nécessaire les
intervalles de validité des calculs :
\begin{enumerate}
  \item $\int \frac{x+2}{x^2-3x-4}\,dx$

  \item $\int \frac{x-1}{x^2+x+1}\,dx$

  \item $\int \sin ^8x\cos ^3x \, dx$

  \item $\int \frac 1{\sin x} \, dx$

  \item $\int \frac{3-\sin x}{2\cos x+3\tan x}\,dx$ 
\end{enumerate}
\finenonce{006866} 


\finexercice


\exercice{6867, bodin, 2012/04/13}
\video{Bydd17Yz8RA}
\enonce{006867}{}
Calculer les intégrales suivantes :
\begin{enumerate}
  \item $\int_0^{\frac \pi 2}x\sin x \, dx$ \quad (intégration par parties)

  \item $\int_0^1 \frac{e^x}{\sqrt{e^x+1}} \,  dx$ \quad (à l'aide d'un changement de variable simple)

  \item $\int_0^1\frac 1{\left( 1+x^2\right) ^2} \, dx$ \quad (changement de variable $x=\tan t$)

  \item $\int_0^1\frac{3x+1}{\left( x+1\right) ^2} \,  dx$ \quad (décomposition en
éléments simples)

  \item $\int_{\frac 12}^2\left( 1+\frac 1{x^2}\right) \arctan x \, dx$ \quad (changement de variable $u=\frac 1x$)

\end{enumerate}
\finenonce{006867} 


\finexercice

\section{127.12 Intégrale impropre}


\exercice{1280, ridde, 1999/11/01}

\enonce{001280}{}
 Donner la nature des int\'egrales suivantes :
$$\displaystyle{\int _{0}^{\infty }\!{\frac {{e^{-x}}}{\sqrt {x}}}{dx}}.$$
$$ \displaystyle{\int _{1}^{\infty }\!{x}^{x}{dx} }.$$

$$\displaystyle{\int _{0}^{\infty }\!{\frac {\sqrt {x}\sin(\frac 1x)}{\ln (1+x)}}{dx}} .$$

Nature et calcul des int\'egrales suivantes :
$$\displaystyle{\int _{1}^{2}\!{\frac {1}{\sqrt {{x}^{2}-1}}}{dx}}.$$
$$\displaystyle{\int _{0}^{\infty }\!{\frac {{x}^{5}}{{x}^{12}+1}}{dx}}.$$
$$\displaystyle{\int _{0}^{\infty }\!{e^{-\sqrt {x}}}{dx}} .$$
$$\displaystyle{\int _{1}^{\infty }\frac 1{\text{sh} (bile)}d (bile) }.$$

\finenonce{001280}



\finexercice



\exercice{1281, ridde, 1999/11/01}

\enonce{001281}{}
\begin{enumerate}

\item Montrer que $\forall x>-1 \, \,   \ln (1 + x)\leq x.$\\
\item Soit $n \in \Nn^{*}.$ Montrer que $\forall x \in [0, n]\, \,   (1-\frac{x}{n})^{n}\leq
e^{-x} \leq (1 + \frac{x}{n})^{-n}.$\\
\item En d\'eduire que\\
$\displaystyle{\int _{0}^{\sqrt {n}}\!\left (1-{\frac {{t}^{2}}{n}}\right )^{n}{dt}} \leq
\displaystyle{\int _{0}^{\sqrt {n}}\!{e^{-{t}^{2}}}{dt}} \leq
\displaystyle{\int _{0}^{\sqrt {n}}\!\frac{1}{\left (1+{\frac {{t}^{2}}{n}}\right )^n}{dt}}. $ \\

Rappel (int\'egrales de Wallis) : $I_{n} =
\displaystyle{\int _{0}^{{\frac {\pi }{2}}}\!\left (\cos(\theta)\right )^{n}{d\theta}} \sim
\sqrt {\frac{\pi}{2n}}. $\\

\item Montrer que $\displaystyle{\int _{0}^{\infty } \frac 1{(1+{u}^{2})^n}{du}}$
existe et vaut $I_{2n-2}. $\\

\item Montrer que $\displaystyle{\int _{0}^{\infty }\!{e^{-{x}^{2}}}{dx}} $ existe et vaut
 $\frac{\sqrt\pi}{2} $.
\end{enumerate}
\finenonce{001281}



\finexercice



\exercice{1282, gourio, 2001/09/01}

\enonce{001282}{}
\'Etude de :
$$f:{\Rr}\rightarrow {\Rr} $$
$$x\mapsto \int_{1}^{x}\frac{e^{t}}{t}dt. $$
Donner un \'{e}quivalent de $f$ en 0 et en $+\infty $.
\finenonce{001282}



\finexercice



\exercice{1283, gourio, 2001/09/01}

\enonce{001283}{}
Soit $f$ une application $C^{2}$ de ${\Rr}$ dans ${\Rr}$ telle que
$f+f^{^{\prime \prime }}\geq 0.$ Montrer que :
$$\forall x\in {\Rr},f(x)+f(x+\pi )\geq 0. $$
\finenonce{001283}



\finexercice



\exercice{1284, gourio, 2001/09/01}

\enonce{001284}{}
Soit $f$ une application continue de ${\Rr}^{+}$ dans ${\Rr}$ et $F$ de
${\Rr}^{+*}$ dans ${\Rr}$ d\'{e}finie par :
$$\forall x\in {\Rr}^{+*},F(x)=\frac{1}{x}\int_{0}^{x}f(t)dt. $$
\begin{enumerate}
\item Montrer que si $f$ admet une limite $\ell$ en $+\infty $, alors $F$\ a aussi la
limite $\ell$ en $+\infty .$
\item Donner un exemple o\`{u} $f$ n'a pas de limite en $+\infty $ mais o\`{u} $F$
tend vers $0$.
\item Montrer que si $f\rightarrow \infty $ quand $x\rightarrow \infty $, alors
$F\rightarrow \infty $ quand $x\rightarrow \infty .$
\end{enumerate}
\finenonce{001284}



\finexercice



\exercice{1285, gourio, 2001/09/01}

\enonce{001285}{}
\'Etudier la fonction :
$$h:x\rightarrow \int_{x}^{x^{2}}\frac{dt}{\log t}. $$
Domaine de d\'{e}finition, continuit\'{e} et d\'{e}rivabilit\'{e},
variations, limites aux bornes de ce domaine, et $\lim\limits_{x\rightarrow
\infty }\frac{h(x)}{x},\lim\limits_{x\rightarrow 0}\frac{h(x)}{x}, $
\'{e}ventuellement convexit\'{e}.
\finenonce{001285}



\finexercice



\exercice{1286, gourio, 2001/09/01}

\enonce{001286}{}
Donner un exemple d'une fonction continue positive telle que :
$$\int_{0}^{\infty }f(u)du $$
existe mais telle qu'on n'ait pas :
$$\lim\limits_{x\rightarrow \infty }f(x)=0. $$
Donner un exemple de fonction continue positive telle que :
$$\int_{0}^{\infty }f(u)du $$
existe mais telle que :
$$\int_{0}^{\infty }f^{2}(u)du $$
n'existe pas.
\finenonce{001286}



\finexercice



\exercice{1287, gourio, 2001/09/01}

\enonce{001287}{}
Soit $f$ une fonction positive d\'{e}croissante de ${\Rr}^{+}$ dans ${\Rr},$
 telle que $\int_{0}^{\infty }f$ existe. Montrer que :
$$f(x)=o(\frac{1}{x}) $$
quand $x\rightarrow \infty .$
\finenonce{001287}



\finexercice



\exercice{1288, gourio, 2001/09/01}

\enonce{001288}{}
Soit $f$ une application continue par morceaux de ${\Rr}^{+}$ dans ${\Rr}$
 poss\'{e}dant une limite $\ell$ en $+\infty $, telle que $\int_{0}^{\infty }f$
existe ; montrer que $\ell=0.$

Soit $f$ une application uniform\'{e}ment continue de ${\Rr}^{+}$ dans ${\Rr}$
 telle que $\int_{0}^{\infty }f$ existe.\ Montrer que :
$$\lim\limits_{x\rightarrow \infty }f(x)=0. $$
\finenonce{001288}



\finexercice



\exercice{1289, gourio, 2001/09/01}

\enonce{001289}{}
Soit $f$ une application continue de ${\Rr}^{+}$ dans ${\Rr}$ telle
que $\int_{0}^{\infty }f^{2}$ existe.\ Montrer que quand $x\rightarrow\infty :$
$$\int_{0}^{x}f(t)dt=o(\sqrt{x}). $$
\finenonce{001289}



\finexercice



\exercice{1290, gourio, 2001/09/01}

\enonce{001290}{}
\'Etudier la nature de
$$\int_{0}^{\infty }\frac{\sin t}{t^{\alpha }}dt $$
selon $\alpha \in {\Rr}.$
\finenonce{001290}



\finexercice



\exercice{1291, gourio, 2001/09/01}

\enonce{001291}{}
Convergence et calcul de :
\[
\int_{0}^{1}\frac{\ln (1+t^{2})dt}{t^{2}}, \qquad
\int_{0}^{\infty }\ln \left( 1+\frac{1}{t^{2}}\right) dt, \qquad
\int_{1}^{\infty }\frac{\ln t}{t^{n}}dt.
\]
\finenonce{001291}


\finexercice



\exercice{1292, gourio, 2001/09/01}

\enonce{001292}{}
Soit $f:[1,\infty [\rightarrow {\Rr}^{+} $ continue telle que
$$\int_{1}^{\infty }f(t)dt $$ converge.
Montrer que
$$\lim\limits_{x\rightarrow \infty }\frac{1}{x}\int_{1}^{x}tf(t)dt=0. $$
\finenonce{001292}



\finexercice



\exercice{1293, gourio, 2001/09/01}

\enonce{001293}{}
Soit $f\in C([1,\infty [,{\Rr}^{+})$ d\'{e}croissante, on pose :
$$x_{n}=\sum_{k=1}^{n}f(k)-\int_{1}^{n}f(t)dt. $$
\begin{enumerate}
\item Montrer que la suite $\left( x_{n}\right) _{n\in \Nn} $ converge.
\item Montrer que la suite $S_{n}=\sum_{k=1}^{n}f(k) $ a une limite quand
$n\rightarrow \infty $ si et seulement si $\int_{1}^{\infty }f $ converge, et
que dans ce cas :
$$\int_{n+1}^{\infty }f\leq \lim_{m\rightarrow \infty
}\sum_{k=n+1}^{m}f(k)\leq \int_{n}^{\infty }f. $$
\item Montrer que si $\int_{1}^{\infty }f $ diverge on a: $S_{n}\backsim
\int_{1}^{n}f $ quand $n\rightarrow \infty .$
\end{enumerate}
\finenonce{001293}



\finexercice



\exercice{1294, gourio, 2001/09/01}

\enonce{001294}{}
Soit $f:]0;1] \rightarrow {\Rr } $ continue et monotone, telle que $\int_{0}^{1}f$
existe. Calculer
$$\lim\limits_{n\rightarrow \infty }\frac{1}{n}\sum_{k=1}^{n}f\left( \frac{k}{n}\right).$$
\finenonce{001294}



\finexercice



\exercice{1295, gourio, 2001/09/01}

\enonce{001295}{}
Montrer que si $f:{\Rr}^{+}\rightarrow {\Rr}$ est uniform\'{e}ment
continue, alors
$$\int_{0}^{\infty }\exp (if(t)) dt $$
n'existe pas.
\finenonce{001295}



\finexercice



\exercice{1296, gourio, 2001/09/01}

\enonce{001296}{}
Nature de :
$$\int_{0}^{\infty }\sin t\sin \frac{1}{t}dt,\ \int_{0}^{\infty }\frac{e^{\sin t}}{t}dt,\
\int_{2}^{\infty }\frac{\sin t}{\sqrt{t}+\sin t}dt,\ \int_{0}^{1}\cos \ln t dt,\int_{0}^{\infty }
\cos \exp t dt. $$
\finenonce{001296}



\finexercice



\exercice{1297, gourio, 2001/09/01}

\enonce{001297}{}
Nature et calcul de :
$$\int_{0}^{\infty }\ln t\ln \left( 1+\frac{a^{2}}{t^{2}}\right)
dt,a>0\ ;\
\int_{0}^{\infty }\exp \left( -t^{\frac{1}{n}}\right) dt,n\in \Nn^{*}\ ;\
\int_{0}^{1}\left( \frac{1}{t}-E(\frac{1}{t})\right) dt. $$
\finenonce{001297}



\finexercice



\exercice{1298, gourio, 2001/09/01}

\enonce{001298}{}
Convergence et calcul de :
$$
\int_{0}^{\infty }\frac{dx}{1+\cosh ^{2}x}, \
\int_{1}^{\infty }\frac{dx}{\sinh x}, \
\int_{-\infty }^{\infty }\frac{dt}{\cosh t}.$$
\finenonce{001298}



\finexercice



\exercice{1299, gourio, 2001/09/01}

\enonce{001299}{}
Soient $f$ et $g $ deux fonctions de ${\Rr}^{+}$ dans ${\Rr}$ telles que
$f\geq 0, g\geq 0, g=o(f) $ en $+\infty , $ et $ \int_{0}^{\infty }f $
    n'existe pas. Montrer alors :
$$\int_{0}^{x}g(u)du=o\left( \int_{0}^{x}f(u)du\right) $$ quand $x\rightarrow \infty .$
\finenonce{001299}



\finexercice



\exercice{1300, gourio, 2001/09/01}

\enonce{001300}{}
Soit $f:{\Rr}^{+}\rightarrow {\Rr}$ continue, tendant vers $\ell$ en $+\infty , $
montrer alors :
$$\lim\limits_{n\rightarrow \infty }\int_{0}^{\infty }\frac{f(t)n}{n^{2}+t^{2}}dt
=\frac{\pi }{2}\ell. $$
\finenonce{001300}



\finexercice



\exercice{1301, gourio, 2001/09/01}

\enonce{001301}{}
Calculer :
$$
\lim\limits_{a\rightarrow 0^{+}}\int_{a}^{3a}\frac{\tan t}{t^{2}}dt \ ,\
\lim\limits_{n\rightarrow \infty }\int_{0}^{1}\frac{x^{n}-x^{2n}}{1-x}dx.$$
\finenonce{001301}



\finexercice



\exercice{1302, gourio, 2001/09/01}

\enonce{001302}{}
Soit $f\in C({\Rr},{\Rr}) $ telle que $\int_{-\infty }^{\infty }f$
existe, montrer que $ F(x)=\int_{-\infty }^{\infty }f(t)\cos tx dt$ est
uniform\'{e}ment continue sur ${\Rr}.$
\finenonce{001302}



\finexercice



\exercice{2331, matexo1, 2002/02/01}
\enonce{002331}{}
Sans les calculer, dire si les int\'egrales suivantes sont
convergentes ou divergentes\,:
\begin{eqnarray*}
&&\int_0^1 {dt\over \sqrt[3]{t \sqrt{1-t}}}, \\
&&\int_0^{\pi/2} \tan t\ dt, \\
&&\int_0^1{dt\over\sqrt{\arcsin t} \ln(1-t)}.
\end{eqnarray*}
\finenonce{002331}


\finexercice


\exercice{2332, matexo1, 2002/02/01}
\enonce{002332}{}
Sans les calculer, dire si les int\'egrales suivantes sont
convergentes ou divergentes\,:
\begin{eqnarray*}
&&\int_0^\infty{t^3-5t^2+1 \over 2t^5-2t^3+t^2+1}\ dt, \\
&&\int_1^\infty{{1\over t^2}\ln{t-1 \over t+1}\ dt},\\
&&\int_0^\infty{dt\over t \arg \cosh t},  \\
&&\int_0^1 \sin {1\over t}\ dt,
\end{eqnarray*}
\finenonce{002332}


\finexercice


\exercice{2333, matexo1, 2002/02/01}
\enonce{002333}{}
Soit $n \ge 0$ un entier. Montrer que l'int\'egrale
$$I_n = \int_0^\infty t^n \exp(-t^2) dt$$
est convergente. La calculer en fonction de $n$, sachant que $I_0
= \sqrt{\pi}/2$.


\finenonce{002333}


\finexercice


\exercice{2334, matexo1, 2002/02/01}
\enonce{002334}{}
On d\'efinit 
$$ F(x) = \int_0^x {\sin t\over t}\ dt $$
et pour tout $n \in \Nn$, on note 
$$u_n = F((n+1)\pi) - F(n\pi) = \int_{n\pi}^{(n+1)\pi}
{\sin t\over t} dt.$$

\begin{enumerate}
\item Montrer que $F(x)$ est bien d\'efinie pour tout $x \in \Rr$.
\item Montrer que si $k \ge 1$, alors
$${2\over(2k+1)\pi} < u_{2k} < {1\over k\pi}.$$
Trouver une in\'egalit\'e similaire pour $u_{2k+1}$, puis pour
$u_{2k}+u_{2k+1}$.
\item Montrer que la suite de terme g\'en\'eral
$\displaystyle{v_n = \sum_{i=1}^n {1\over i^2}}$
admet une limite finie. En d\'eduire que 
$$ I = \int_0^\infty
{\sin t\over t}\ dt$$
est convergente.
\end{enumerate}
\finenonce{002334}


\finexercice


\exercice{2337, matexo1, 2002/02/01}
\enonce{002337}{}
Soit $\varphi$ la fonction d\'efinie sur $\left[0,1\right[$ par 
$$\varphi (x) =
\int_x^{x^2} {dt\over \ln t}.$$
Montrer que $\varphi (x)$ a
une limite quand $x$ tend vers $1$ et la calculer. 
(Indication\,: comparer \`a 
$\displaystyle \int_x^{x^2} {1/ (t\ln t)}~dt$).
\finenonce{002337}


\finexercice


\exercice{2338, matexo1, 2002/02/01}
\enonce{002338}{}
Dire si les int\'egrales suivantes sont convergentes (en discutant
\'eventuellement suivant la valeur des param\`etres)\,:
$$\int_0^1 {dt\over \sqrt t \sqrt[3]{1-t}},\quad
\int_0^{\pi /2} \tan t\,dt,\quad
\int_0^1 {dt\over t^\alpha  |\ln t|^\beta } ,\quad
\int_0^1 \cos(\ln t)\,dt,\quad
\int_0^1 \sin{1\over t}\, dt,$$
$$ 
\int_0^\infty  {t^2+t-1\over \sqrt t\,(t^3 -2t^2+3t-6)}\,dt,\quad
\int_0^\infty  t^\alpha \,[1-e^{-1/\sqrt t}]\,dt,\quad
\int_0^\infty  {\ln t - \ln(1-e^{-t})\over t} e^{-\alpha  t}\, dt.
$$
\finenonce{002338}


\finexercice


\exercice{2339, matexo1, 2002/02/01}
\enonce{002339}{}
Montrer la convergence des int\'egrales suivantes puis les calculer\,:
$$
\int_0^{\pi /2} {dt\over \cos\alpha  \cos t +1}, \quad 
\int_0^1 {dt\over \sqrt t  +\root3\of t}, \quad 
\int_a^b {dt\over \sqrt{(t-a)(b-t)}},$$
$$ 
\int_{-\infty }^0 {dt\over e^{2t} +e^t-6}     ,\quad
\int_a^\infty  {dt\over t^2\sqrt{t^2+a^2}},\ (a>0), \quad
\int_1^{+\infty } {t^3-t^2-1 \over t^6+2t^4+t^2}\,dt.
$$
\finenonce{002339}


\finexercice


\exercice{4268, quercia, 2010/03/12}
\enonce{004268}{\'Etude de convergence}
%
  $\int_{-\infty}^{+\infty} \frac{ d t}{e^t+t^2e^{-t}} $ (cv)\par
  $\int_1^{+\infty} \frac{ e^{\sin t}}{t} \,d t $ (dv)\par
  $\int_0^1 \frac{ t^\alpha-1}{\ln{t}} \,d t$ (cv ssi  $\alpha > -1 $)\par
  $\int_{e^2}^{+\infty} \frac{ d t}{t(\ln t)(\ln\ln t)} $ (dv)\par
  $\int_0^{+\infty} \ln\Bigl( \frac{1+t^2}{1+t^3} \Bigr) \,d t $ (dv)\par
  $\int_0^{+\infty} \Bigl(2+(t+3)\ln\bigl(\frac{ t+2}{t+4} \bigr)\Bigr) d t $ (cv)\par
  $\int_0^{+\infty} \frac{ {t\ln t}}{(1+t^2)^\alpha{}} \,d t$ (cv  ssi $\alpha > 1 $)\par
  $\int_0^1 \frac{ d t}{1-\sqrt t} $ (dv)\par
  $\int_0^{+\infty} \frac{ (t+1)^\alpha-t^\alpha}{t^\beta} \,d t$ (cv ssi $0 < \beta-\alpha < 1$ ou $\alpha = 0 $)\par
  $\int_0^{+\infty} \sin(t^2)\,d t $ (cv)\par
  $\int_0^1 \frac{ d t}{\arccos{t}} $ (cv)\par
  $\int_0^{+\infty} \frac{ \ln(\Arctan{t})}{t^\alpha} \,d t$ (dv)\par
  $\int_1^{+\infty} \frac{ {\ln(1+1/t)\,d t}}{(t^2-1)^\alpha}$ (cv  ssi  $0<\alpha<1$)\par
  $\int_0^1 \frac{ |\ln{t}|^\beta}{(1-t)^\alpha} \,d t$  (cv ssi $\alpha < \beta+1$)\par
  $\int_0^{+\infty} t^\alpha\bigl(1-e^{-1}{\sqrt t}\bigr)\,d t$ (cv ssi $-1<\alpha<-\frac 12$)\par
  $\int_0^1 \sin\bigl(\frac1t\bigr)e^{-1}{t}t^{-k}\,d t$ (cv)\par
\finenonce{004268}



\finexercice


\exercice{4269, quercia, 2010/03/12}
\enonce{004269}{Fractions rationnelles}
%
  $\int_0^{+\infty} \frac{ d t}{(1+t^2)^2} = \frac{\pi}{4} $\par
  $\int_{-\infty}^{+\infty} \frac{ d t}{t^2+2t+2} = \pi $\par
  $\int_0^{+\infty} \frac{ d t}{(1+t^2)^4} = \frac{ 5\pi}{32} $\par
  $\int_{-\infty}^{+\infty} \frac{ d t}{(t^2+1)(t^2-2t\cos\alpha+1)} = \frac{\pi}{2|\sin\alpha|} $\par
  $\int_0^{+\infty} \frac{ 2t^2+1}{(t^2+1)^2} \,d t = \frac{ 3\pi}{4} $\par
  $\int_{-\infty}^{+\infty} \frac{ t^2\,d t}{(t^2+1)(t^2+a^2)} = \frac{ \pi}{1+|a|} $\par
  $\int_0^{+\infty} \frac{ d t}{1+t^4} = \frac{\pi}{2\sqrt2} $\par
  $\int_0^{+\infty} \frac{ t^2\,d t}{1+t^4} = \frac{\pi}{2\sqrt2} $\par
  $\int_1^{+\infty} \frac{ d t}{t^6(1+t^{10})} = \frac{ 4-\pi}{20} $\par
\finenonce{004269}



\finexercice


\exercice{4270, quercia, 2010/03/12}
\enonce{004270}{Fonctions trigonométriques}

%
  $\int_0^{2\pi} \frac{ d t}{2+\sin{t}} = \frac{ 2\pi}{\sqrt3} $\par
  $\int_{-\pi}^\pi \frac{ 2d t}{2+\sin t+\cos t} = 2\pi\sqrt2 $\par
  $\int_0^{\pi/2} \sqrt{\tan t}\, d t = \int_0{+\infty} \frac{ 2t^2\,d t}{1+t^4} = \frac{ \pi}{\sqrt2} $\par
  $\int_0^{\pi/2} \frac{ d t}{3\tan{t}+2} = \frac{ {\pi+3\ln(3/2)}}{13} $\par
  $\int_0^\pi \frac{ d t}{(a\sin^2t+b\cos^2t)^2} = \frac{ \pi(a+b)}{2\sqrt{ab}^3} $\par
  $\int_0^{\pi/4} \cos t\ln(\tan t)\,d t = -\ln(1+\sqrt2\,)$\par
\finenonce{004270}



\finexercice


\exercice{4271, quercia, 2010/03/12}
\enonce{004271}{Radicaux}
%
  $\int_0^1 \sqrt{ \frac{ t}{1-t} }\,d t = \frac{ \pi}{2} $\par
  $\int_1^{10} \frac{d t}{\sqrt[3]{t-2}} = \frac{9}{2} $\par
  $\int_a^b \frac{d t}{\sqrt{(t-a)(b-t)}} = \pi $\par
  $\int_0^1 \frac{ t^5\,d t}{\sqrt{1-t^2}} = \frac{ 8}{15} $\par
  $\int_{-1}^1 \frac{ d t}{(1+t^2)\sqrt{1-t^2}} = \frac{ \pi}{\sqrt2} $\par\penalty5000
  $\int_0^1 \frac{ d t}{(4-t^2)\sqrt{1-t^2}} = \frac{ \pi}{4\sqrt3} $\par
  $\int_0^1 \frac{ t\,d t}{\sqrt{(1-t)(1+3t)}} = \frac{ 2\pi}{9\sqrt3} + \frac{1}{3} $\par
  $\int_0^1 \frac{ d t}{(1+t){\sqrt[3]{t^2-t^3}}} = \frac{ \pi{\root 3 \of 4}}{\sqrt3} $\par
  $\int_0^1 \Arctan\sqrt {1-t^2} \,d t = \frac{ \pi(\sqrt2-1)}{2} $\par
  $\int_1^{+\infty} \frac{ d t}{t\sqrt{t^{10}+t^5+1}} = \frac{1}{5} \ln\biggl(1+\frac{2}{\sqrt3} \biggr)$\par
  $\int_0^{+\infty} \frac{ d t}{(1+t^2)\sqrt{t}} = \frac{\pi}{\sqrt2} $\par
\finenonce{004271}



\finexercice


\exercice{4272, quercia, 2010/03/12}
\enonce{004272}{Exponentielles}
%
  $\int_2^{+\infty} \frac{ e^t\,d t}{(e^{2t}-5e^t+6)(e^t-1)}
   = \ln\biggl( \frac{ e^2-2}{\sqrt{e^4-4e^2+3}} \biggr) $\par
  $\int_0^{+\infty} \frac{ d t}{\ch^4t+\sh^4t} = \frac{ \ln(\sqrt2+1)}{\sqrt2} $\par
\finenonce{004272}



\finexercice


\exercice{4273, quercia, 2010/03/12}
\enonce{004273}{Divers}
%
  $\int_0^{+\infty} te^{-\sqrt t} \,d t = 12 $\par
  $\int_0^1 \Arcsin t \,d t = \frac{\pi}{2} - 1 $\par
  $\int_0^1 \frac{ \ln(1-t^2)}{t^2} \,d t = -2\ln 2 $\par
  $\int_0^{+\infty} \frac{ t^3\ln{t}}{(1+t^4)^3} \,d t = -\frac{1}{32} $\par
  $\int_0^{\pi/2} \ln\sin t\,d t = -\frac{ \pi\ln2}{2} $\par
  $\int_0^1 \frac{ \ln{t}}{\sqrt{1-t}} \,d t = 4\ln2 - 4 \ (u = \sqrt{1-t}\,) $\par
  $\int_0^{+\infty} \frac{ \ln{t}}{1+t^2} \,d t = 0\ (u = 1/t)$\par
  $\int_0^1 \frac{ \ln{t}}{(1+t)\sqrt{1-t^2}} \,d t = \ln2 - \frac{\pi}{2} \biggl( u = \sqrt{{\frac{1-t}{1+t}}}\,\biggr)$\par
  $\int_0^1 \frac{ d t}{\sqrt{1+t}+\sqrt{1-t}} = \sqrt2 + \ln(\sqrt2 - 1) $\par
  $\int_0^{+\infty} \ln\Bigl(1+\frac{ a^2}{t^2} \Bigr)\,d t = a\pi$\par
  $\int_0^{+\infty}\ln\left|\frac{1+t}{1-t} \right|\frac{ t\,d t}{(a^2+t^2)^2} = \frac{\pi}{2|a|(a^2+1)} $\par
\finenonce{004273}



\finexercice


\exercice{4274, quercia, 2010/03/12}
\enonce{004274}{Centrale PC 1999}

Soit $(a_k)$ une suite de réels telle que $\sum_{k=0}^n a_k = 0$.
\'Etudier la convergence de
$ \int_{t=0}^{+\infty} \sum_{k=0}^n a_k\cos(a_kt)\frac{d t}t$.
\finenonce{004274}



\finexercice


\exercice{4275, quercia, 2010/03/12}
\enonce{004275}{Chimie P 91}

    Existence et calcul de $f(x) =  \int_0^\pi\frac{dt}{1-x\cos t}$.
    

\finenonce{004275}



\finexercice


\exercice{4276, quercia, 2010/03/12}
\enonce{004276}{Chimie P 1996}

    Convergence et calcul de $ \int_{t=0}^{+\infty}\frac{t\,d t}{\sh t}$
    (on pourra décomposer l'intégrande en somme d'une série de fonctions).

\finenonce{004276}



\finexercice


\exercice{4277, quercia, 2010/03/12}
\enonce{004277}{Calcul par récurrence}

On pose $I_n =  \int_{t=0}^{\pi/2} \cos(2nt)\ln(\sin t)\,d t$ ($n\in\N^*$).
Calculer $2nI_n - (2n+2)I_{n+1}$ et en déduire $I_n$ en fonction de~$n$.

\finenonce{004277}



\finexercice


\exercice{4278, quercia, 2010/03/12}
\enonce{004278}{Calcul par récurrence}

Soit $\alpha \in\, ]0,\pi[$ et
$I_n =  \int_{t=0}^\pi \frac{\cos nt\,d t}{1-\sin\alpha\cos t} $.

Calculer $I_n+I_{n+2}$ en fonction de $I_{n+1}$ puis
exprimer $I_n$ en fonction de $\alpha$ et $n$.

\finenonce{004278}



\finexercice


\exercice{4279, quercia, 2010/03/12}
\enonce{004279}{Calcul par récurrence}

Calculer par récurrence :
$I_n =  \int_{t=0}^1 \frac{t^n\,d t}{\sqrt[4]{t^3(1-t)}} $.

\finenonce{004279}



\finexercice


\exercice{4280, quercia, 2010/03/12}
\enonce{004280}{Mines-Ponts 1999}

Calculer $I_n =  \int_{t=0}^{+\infty}\frac{d t}{(t+1)(t+2)\dots(t+n)}$.

\finenonce{004280}



\finexercice


\exercice{4281, quercia, 2010/03/12}
\enonce{004281}{Calcul de $\int_0^\infty \sin t/t\,d t$}

\begin{enumerate}
  \item A l'aide d'une intégration par parties, montrer que
$ \int_{t=0}^{+\infty} \frac{\sin t}t \,d t =  \int_{t=0}^{+\infty} \frac{\sin^2 t}{t^2} \,d t$.

  \item Montrer que $I_n =  \int_{t=0}^{\pi/2} \frac{\sin^2 nt}{t^2}\,d t$ est comprise
    entre $A_n =  \int_{t=0}^{\pi/2} \frac{\sin^2 nt}{\sin^2 t}\,d t$ et
    $B_n =  \int_{t=0}^{\pi/2} \mathrm{cotan}^2 t\sin^2 nt\,d t$.
  \item Calculer $A_n+A_{n+2} - 2A_{n+1}$ et $A_n-B_n$.
    En déduire les valeurs de $A_n$ et $B_n$ en fonction de $n$.
    
  \item Lorsque  $n\to\infty$ montrer que  $\frac{I_n}n \to J =  \int_{t=0}^{+\infty} \frac{\sin^2 t}{t^2}\,d t$
    et donner la valeur de cette dernière intégrale.
    

\end{enumerate}
\finenonce{004281}



\finexercice


\exercice{4282, quercia, 2010/03/12}
\enonce{004282}{$\int_0^\infty \text{périodique}/t\,d t$}

Soit $f:\R \to \R$ continue, périodique de période $T > 0$.
On note $m = \frac1T \int_{t=0}^T f(t)\,d t$. Montrer que
$ \int_{t=T}^{+\infty} \frac{f(t)}t\,d t$ converge si et seulement
si $m=0$.
\finenonce{004282}



\finexercice


\exercice{4283, quercia, 2010/03/12}
\enonce{004283}{$\int_0^\infty f(t)/t\,d t$}

Soit $f$ une application continue de $[1,+\infty[$ dans $\R$.
Montrer que si l'intégrale $ \int_{t=1}^{+\infty} f(t)\,d t$ converge, il en est de même
de l'intégrale $ \int_{t=1}^{+\infty} \frac{f(t)}t\,d t$.
On pourra introduire la fonction $F(x) =  \int_{t=1}^x f(t)\,d t$.
\finenonce{004283}



\finexercice


\exercice{4284, quercia, 2010/03/12}
\enonce{004284}{Polynôme$\times e^{-t}$}

Soit $\varphi : {\R_n[X]} \to {\R^{n+1}}, P \mapsto {(a_0,\dots,a_n)}$
avec $a_k =  \int_{t=0}^{+\infty} e^{-t}t^kP(t)\,d t$.

\begin{enumerate}
  \item Justifier l'existence de $\varphi$.
  \item Montrer que $\varphi$ est un isomorphisme d'espace vectoriel.
\end{enumerate}
\finenonce{004284}



\finexercice


\exercice{4285, quercia, 2010/03/12}
\enonce{004285}{Constante d'Euler}

Calculer $ \int_{t=1}^{+\infty} \frac{t-[t]}{t^2} \,d t$ en fonction de la constante d'Euler.

\finenonce{004285}



\finexercice


\exercice{4286, quercia, 2010/03/12}
\enonce{004286}{Constante d'Euler}

Soit $\gamma$ la constante d'Euler. Montrer que~\dots

\begin{enumerate}
  \item $ \int_{t=0}^{+\infty} e^{-t}\ln t\,d t = -\gamma$.


  \item $ \int_{t=0}^1 \frac{1-e^{-t}-e^{-1/t}}t\,d t = \gamma$.


  \item $ \int_{t=0}^1 \left(\frac 1t + \frac 1{\ln(1-t)}\right)d t = \gamma$.

\end{enumerate}
\finenonce{004286}



\finexercice


\exercice{4287, quercia, 2010/03/12}
\enonce{004287}{Sommes de Riemann}

Soit $f : {[a,b[} \to {\R^+}$ continue croissante.
On pose $S_n = \frac{b-a}n \sum_{k=0}^{n-1} f\Bigl(a+k\frac{b-a}n \Bigr)$.

\begin{enumerate}
  \item Si $ \int_{t=a}^b f(t)\,d t$ converge, montrer que $S_n \to  \int_{t=a}^b f(t)d t$ lorsque $n\to\infty$.
  \item Si $ \int_{t=a}^b f(t)\,d t$ diverge, montrer que  $S_n \to +\infty$ lorsque $n\to\infty$.
\end{enumerate}
\finenonce{004287}



\finexercice


\exercice{4288, quercia, 2010/03/12}
\enonce{004288}{Sommes de Riemann}

Calculer $\lim_{n\to\infty} \frac 1{\sqrt{n^2-1}} + \frac 1{\sqrt{n^2-4}}
          + \dots + \frac 1{\sqrt{n^2-(n-1)^2}}$.


\finenonce{004288}



\finexercice


\exercice{4289, quercia, 2010/03/12}
\enonce{004289}{Comparaison série-intégrale}

Soit $f : {[0,+\infty[} \to \R$ continue décroissante telle que
$ \int_{}^{+\infty} f(t)\,d t$ converge.

\begin{enumerate}
  \item Montrer que la série $\sum_{k=0}^\infty f(k)$ converge et encadrer le reste :
     $\sum_{k=n}^\infty f(k)$ à l'aide d'intégrales de $f$.
  \item Application : Pour $\alpha > 1$, donner un équivalent pour $n \to \infty$ de
    $\sum_{k=n}^\infty \frac1{k^\alpha} $.

\end{enumerate}
\finenonce{004289}



\finexercice


\exercice{4290, quercia, 2010/03/12}
\enonce{004290}{Comparaison série-intégrale}

Soit $f : {\R^+} \to \R$. On pose, sous réserve de convergence,
$g(t) = \sum_{n=0}^\infty f(nt)$ pour $t>0$.

\begin{enumerate}
  \item Si $f$ est monotone et intégrable, montrer que $g(t)$ existe
pour tout $t>0$ et que l'on a $\smash{tg(t)\to  \int_{u=0}^{+\infty} f(u)\,d u}$ lorsque $t\to0^+$.


  \item Même question en supposant $f$ de classe $\mathcal{C}^1$ et $f$,$f'$ intégrables.


  \item On suppose maintenant $f$ de classe $\mathcal{C}^2$ et $f$, $f'$, $f''$ intégrables.

Montrer que $g(t) = \frac1t \int_{u=0}^{+\infty} f(u)\,d u + \frac{f(0)}2 + O_{t\to0^+}(t)$.


\end{enumerate}
\finenonce{004290}



\finexercice


\exercice{4291, quercia, 2010/03/12}
\enonce{004291}{Valeur moyenne d'une variable aléatoire à densité}

Soit $f : {[0,+\infty[} \to {\R^+}$ continue telle que
$ \int_{t=0}^{+\infty} tf(t) \,d t$ converge.
On pose $F(x) =  \int_{t=x}^{+\infty} f(t)\,d t$.

\begin{enumerate}
  \item  Justifier l'existence de $F(x)$, et montrer que $F(x) = o \Bigl({\frac1x }\Bigr)$ pour $x\to+\infty$.
  \item  Montrer que $ \int_{t=0}^{+\infty} F(t) \,d t =  \int_{t=0}^{+\infty} tf(t) d t$.
\end{enumerate}
\finenonce{004291}



\finexercice


\exercice{4292, quercia, 2010/03/12}
\enonce{004292}{$\int_0^\infty f(t)/t^2\,d t$}

Soit $f : {\R^+} \to {\R^+}$ une fonction de classe $\mathcal{C}^1$ vérifiant :
$\exists\ \alpha > 0 \text{ tel que } \forall\ x \ge 0,\ f'(x) \ge \alpha$.
Montrer que $ \int_{t=1}^{+\infty} \frac{f(t)}t\,d t$ diverge.

\finenonce{004292}



\finexercice


\exercice{4293, quercia, 2010/03/12}
\enonce{004293}{$x(f(x)-f(x+1))$}

Soit $f : {[1,+\infty[} \to {\R^+}$ une fonction décroissante telle que
$ \int_{t=1}^{+\infty} f(t)\,d t$ converge.

Montrer que $xf(x) \to 0$ lorsque $x\to+\infty$, puis que
$ \int_{t=1}^{+\infty} t(f(t)-f(t+1))\,d t$ converge, et calculer la valeur
de cette intégrale.

\finenonce{004293}



\finexercice


\exercice{4294, quercia, 2010/03/12}
\enonce{004294}{$f(|t-1/t|)$}

Soit $f : {[0,+\infty[} \to \R^+$ une fonction continue telle que
$ \int_{t=0}^{+\infty} f(t) \,d t$ converge.

Montrer que $ \int_{t=0}^{+\infty} f(t) \,d t =
     \int_{u=0}^{+\infty} f\Bigl(\Bigl| u - \frac1u \Bigr|\Bigr) \,d u$.

\finenonce{004294}



\finexercice


\exercice{4295, quercia, 2010/03/12}
\enonce{004295}{$(f(ax) - f(x))/x$}

\begin{enumerate}
  \item Soit $f : {]0,{+\infty}[} \to \R$ une fonction continue telle que
$\begin{cases} f(x) \to \ell \cr f(x) \text{ si } x\to{0^+}\to L \text{ si } x\to{+\infty} .\cr\end{cases}$

Pour $a > 0$, établir la convergence et calculer la valeur de
$ \int_{t=0}^{+\infty} \frac {f(at) - f(t)}t \,d t$.

  \item Application : Calculer $ \int_{t=0}^1 \frac{t-1}{\ln t}\,d t$.

\end{enumerate}
\finenonce{004295}



\finexercice


\exercice{4296, quercia, 2010/03/12}
\enonce{004296}{$f(t+a)-f(t)$, Ensi PC 1999}

\begin{enumerate}
  \item Soit $f : {\R^+} \to {\R^+}$ continue ayant une limite finie en~$+\infty$.
    Montrer que $ \int_{t=0}^{+\infty} (f(t+a) - f(t))\,d t$ converge.
  \item Calculer $ \int_{t=0}^{+\infty} (\arctan(t+1) - \arctan(t))\,d t$.
    

\end{enumerate}
\finenonce{004296}



\finexercice


\exercice{4297, quercia, 2010/03/12}
\enonce{004297}{Valeur moyenne}

Soit $f:\R \to \R$ continue par morceaux telle que $ \int_{t=-\infty}^{+\infty}|f(t)|\,d t$ converge.
On pose $F(x) = \frac12 \int_{t=x-1}^{x+1}f(t)\,d t$.

Montrer que $ \int_{t=-\infty}^{+\infty}F(t)\,d t =  \int_{t=-\infty}^{+\infty}f(t)\,d t$.

Démontrer le même résultat en supposant seulement la convergence de
$ \int_{t=-\infty}^{+\infty}f(t)\,d t$.
\finenonce{004297}



\finexercice


\exercice{4298, quercia, 2010/03/12}
\enonce{004298}{$(\int tf(t)\,d t)/x$}

Soit $f : {[0,+\infty[} \to \R$ continue telle que $ \int_{t=0}^{+\infty} f(t)\,d t$
converge. Montrer que $\frac1x \int_{t=0}^x tf(t)\,d t \to 0$ lorsque $x\to+\infty$.



\finenonce{004298}



\finexercice


\exercice{4299, quercia, 2010/03/12}
\enonce{004299}{$f$ uniformément continue}

Soit $f : {[0,+\infty[} \to {\R}$ uniformément continue telle que
$ \int_{t=0}^{+\infty} f(t)\,d t$ converge.

\begin{enumerate}
  \item Montrer que $f(t) \to 0$ lorsque $t\to{+\infty}$(raisonner par l'absurde).
  \item Si $f$ est positive, montrer que $ \int_{t=0}^{+\infty} f^2(t)\,d t$ converge.
  \item Donner un contre-exemple si $f$ n'est pas de signe constant.

\end{enumerate}
\finenonce{004299}



\finexercice


\exercice{4300, quercia, 2010/03/12}
\enonce{004300}{$f$ décroissante $ \Rightarrow  xf(x) \to 0$}

Soit $f : {[0,+\infty[} \to \R$ continue telle que $ \int_{t=0}^{+\infty} f(t)\,d t$
converge.

\begin{enumerate}
  \item  Si $f(x) \to L$ lorsque $x\to{+\infty}$, combien vaut $L$ ?
  \item  Donner un exemple où $f$ n'a pas de limite en $+\infty$.
  \item  Si $f$ est décroissante, montrer que $xf(x) \to 0$ lorsque $x\to{+\infty}$.

\end{enumerate}
\finenonce{004300}



\finexercice


\exercice{4301, quercia, 2010/03/12}
\enonce{004301}{$\int e^{-t}/t,d t$}

On pose $f(x) =  \int_{t=x}^{+\infty} \frac{e^{-t}}t \,d t$.


\begin{enumerate}
  \item   Chercher $\lim_{x\to{+\infty}} f(x)$.
  \item   A l'aide d'une intégration par parties, donner un équivalent de $f(x)$
      pour $x\to{+\infty}$.
      
  \item   Donner un équivalent de $f(x)$ pour $x\to0^+$.
      
\end{enumerate}
\finenonce{004301}



\finexercice


\exercice{4302, quercia, 2010/03/12}
\enonce{004302}{Intégrale de Gauss}

\begin{enumerate}
  \item Montrer que pour $0 \le x \le \sqrt n$ on a : $\Bigl(1-\frac{x^2}n\Bigr)^n \le e^{-x^2}$
    et pour $x$ quelconque : $e^{-x^2} \le \Bigl(1+\frac{x^2}n\Bigr)^{-n}$.
  \item Calculer les intégrales
    $I_n =  \int_{t=0}^{\sqrt n} \Bigl(1-\frac{t^2}n\Bigr)^n\, d t$ et
    $J_n =  \int_{t=0}^{+\infty} \Bigl(1+\frac{t^2}n\Bigr)^{-n}\, d t$
    en fonction des intégrales :
    $K_p =  \int_{t=0}^{\pi/2} \cos^pt\,d t$.
    
  \item On admet que $K_p \sim \sqrt{\frac\pi{2p}}$ quand $p\to\infty$.
    Calculer $ \int_{t=0}^{+\infty} e^{-t^2}\,d t$.
    
\end{enumerate}
\finenonce{004302}



\finexercice


\exercice{4303, quercia, 2010/03/12}
\enonce{004303}{Intégrales de Gauss}

On admet que $ \int_{t=0}^{+\infty} e^{-t^2}\,d t =\frac{\sqrt\pi}2$.
Calculer les intégrales :
$I_n =  \int_{t=0}^{+\infty} e^{-t^2}t^{2n}\,d t$ pour $n \in \N$.

\finenonce{004303}



\finexercice


\exercice{4304, quercia, 2010/03/12}
\enonce{004304}{Mines-Ponts MP 2005}

Nature et calcul de $ \int_{x=0}^{+\infty} \exp\Bigl(-\bigl(x-\frac 1x\bigr)^2\Bigr)\,d x$~?

\finenonce{004304}



\finexercice


\exercice{4305, quercia, 2010/03/12}
\enonce{004305}{} 

Existence de $ \int_{x=0}^{+\infty} \sin(x^4+x^2+x)\,d x$.

\finenonce{004305}



\finexercice


\exercice{4306, quercia, 2010/03/12}
\enonce{004306}{$\cos(P(t))$}

Soit $P$ un polynôme à coefficients réels de degré supérieur ou égal à 2.
Montrer que $ \int_{t=0}^{+\infty} \cos(P(t))\,d t$ converge.


\finenonce{004306}



\finexercice


\exercice{4307, quercia, 2010/03/12}
\enonce{004307}{Ensi PC 1999}

Soient $I =  \int_{u=0}^{+\infty} \frac{d u}{(1+u^2)(1+u^n)}$
et $J =  \int_{u=0}^{+\infty} \frac{u^n\,d u}{(1+u^2)(1+u^n)}$
($n\in\N$).


Prouver que ces intégrales convergent, qu'elles sont égales et les calculer.

\finenonce{004307}



\finexercice


\exercice{4308, quercia, 2010/03/12}
\enonce{004308}{$f$ et $f''$ de carrés sommables}

Soit $f:\R \to \R$ de classe $\mathcal{C}^2$ telle que $ \int_{t=0}^{+\infty}f^2(t)\,d t$
et  $ \int_{t=0}^{+\infty}f''^2(t)\,d t$ convergent.
Montrer que $ \int_{t=0}^{+\infty}f'^2(t)\,d t$ converge.
\finenonce{004308}



\finexercice


\exercice{4309, quercia, 2010/03/12}
\enonce{004309}{$f' \le 1$, Ulm 1999}

Soit $f : {\R^+} \to {\R^+}$ de classe $\mathcal{C}^1$, intégrable.
\begin{enumerate}
  \item On suppose $f' \le 1$. Montrer que $f(x)\to0$ lorsque $x\to+\infty$.
  \item Est-ce encore vrai si on suppose seulement $f'\le 1+g$ avec $g$ intégrable~?
\end{enumerate}
\finenonce{004309}



\finexercice


\exercice{4310, quercia, 2010/03/12}
\enonce{004310}{Intégrales emboitées}

\'Etablir la convergence et calculer la valeur de
$ \int_{x=0}^{+\infty} \int_{t=x}^{+\infty}\frac{\sin t}t\,d t\,d x$.
\finenonce{004310}



\finexercice


\exercice{4311, quercia, 2010/03/12}
\enonce{004311}{Centrale MP 2001}

Soit $f$ de classe $\mathcal{C}^2$ sur $\R^+$ à valeurs dans $\R$ telle que
$f^2$ et $f''^2$ sont intégrables sur $\R^+$.
Montrer que $ff''$ et $f'^2$ sont intégrables sur $\R^+$, que $f$ est
uniformément continue et qu'elle tend vers zéro en $+\infty$.
\finenonce{004311}



\finexercice


\exercice{4312, quercia, 2010/03/12}
\enonce{004312}{X MP$^*$ 2000}
Donnez un équivalent pour $x\to+\infty$ de $ \int_{t=0}^x \Bigl|\frac{\sin t}{t}\Bigr|\,d t$.
\finenonce{004312}



\finexercice


\exercice{5713, rouget, 2010/10/16}
\enonce{005713}{}
Etudier l'existence des intégrales suivantes

\begin{center}
\begin{tabular}{lll}
\textbf{1) (**)} $\int_{0}^{+\infty}\left(x+2 -\sqrt{x^2+4x+1}\right)\;dx$&\textbf{2) (**)} $\int_{1}^{+\infty}\left(e -\left(1+\frac{1}{x}\right)^x\right)\;dx$&\textbf{3) (**)} $\int_{0}^{+\infty}\frac{\ln x}{x+e^x}\;dx$\\
\rule[-6mm]{0mm}{14mm}\textbf{4) (***)} 
$\int_{0}^{+\infty}\left(\sqrt[3]{x+1}-\sqrt[3]{x}\right)^{\sqrt{x}}\;dx$&
\textbf{5) (**)} $\int_{1}^{+\infty}e^{-\sqrt{x^2-x}}\;dx$&\textbf{6) (**)} $\int_{0}^{+\infty}x^{-\ln x}\;dx$\\
\rule[-6mm]{0mm}{0mm}\textbf{7) (**)} $\int_{0}^{+\infty}\frac{\sin(5x)-\sin(3x)}{x^{5/3}}\;dx$&\textbf{8) (**)} $\int_{0}^{+\infty}\frac{\ln x}{x^2-1}\;dx$&\textbf{9) (**)} $\int_{-\infty}^{+\infty}\frac{e^{-x^2}}{\sqrt{|x|}}\;dx$\\
\textbf{10) (**)} $\int_{-1}^{1}\frac{1}{(1+x^2)\sqrt{1-x^2}}\;dx$&\textbf{11) (**)} $\int_{0}^{1}\frac{1}{\sqrt[3]{x^2-x^3}}\;dx$&\textbf{12) (***)} $\int_{0}^{1}\frac{1}{\Arccos(1-x)}\;dx$.
\end{tabular}
\end{center}
\finenonce{005713}


\finexercice


\exercice{5714, rouget, 2010/10/16}
\enonce{005714}{}
Etudier l'existence des intégrales suivantes.

\begin{center}
\begin{tabular}{ll}
\textbf{1) (***) I} $\int_{2}^{+\infty}\frac{1}{x^a\ln^bx}\;dx$ (Intégrales de \textsc{Bertrand})&\textbf{2) (**)} $\int_{0}^{\pi/2}(\tan x)^a\;dx$\\
\rule[-6mm]{0mm}{14mm}\textbf{3) (**)} $\int_{1}^{+\infty}\left(\left(1+\frac{1}{x}\right)^{1+\frac{1}{x}}- a-\frac{b}{x}\right) dx$& 
\textbf{4) (***)} $\int_{0}^{+\infty}\frac{1}{x^a(1+x^b)}\;dx$
\end{tabular}
\end{center}
\finenonce{005714}


\finexercice


\exercice{5715, rouget, 2010/10/16}
\enonce{005715}{}

\label{ex:rou3}

(Hors programme) Etudier la convergence des intégrales impropres suivantes :
\begin{enumerate}
\item \textbf{(**)} $\displaystyle \int_{0}^{+\infty}\frac{\sin x}{x}\;dx$

\item \textbf{(**)} $\displaystyle \int_{0}^{+\infty}\frac{\sin x}{x^a}\;dx$

\item \textbf{(**)} $\displaystyle \int_{0}^{+\infty}e^{ix^2}\;dx$

\item \textbf{(**)} $\displaystyle \int_{0}^{+\infty}x^3\sin(x^8)\;dx$

\item \textbf{(**)} $\displaystyle \int_{0}^{+\infty}\cos(e^x)\;dx$

\item \textbf{(****)} $\displaystyle \int_{0}^{+\infty}\frac{1}{1+x^3\sin^2x}\;dx$
\end{enumerate}

\end{center}

\finenonce{005715}







\finexercice



\exercice{5716, rouget, 2010/10/16}
\enonce{005716}{}
Existence et calcul de :

\begin{center}
\begin{tabular}{ll}
\textbf{1) (** I)} $I_n=\int_{0}^{+\infty}\frac{1}{(x^2+1)^n}\;dx$&\textbf{2) (très long)} $\int_{2}^{+\infty}\frac{1}{(x-1)^3(x^4+1)}\;dx$\\
\textbf{3) (** I)} $\int_{0}^{+\infty}\frac{1}{x^3+1}\;dx$&
\textbf{4) (***)} $\int_{0}^{+\infty}\frac{1}{(x+1)(x+2)\ldots(x+n)}\;dx$\\
\rule[-6mm]{0mm}{14mm}\textbf{5)(***)} $\int_{0}^{1}\frac{1}{\sqrt{(1-x)(1+ax)}}\;dx$&\textbf{6) (**)} $\int_{0}^{+\infty}\frac{1}{(e^x+1)(e^{-x}+1)}\;dx$\\
\rule[-6mm]{0mm}{0mm}\textbf{7) (**)} $\int_{0}^{+\infty}\frac{1}{5\ch x+3\sh x+4}\;dx$&\textbf{8) (***)} $\int_{0}^{+\infty}\left(2+(t+3)\ln\left(\frac{t+2}{t+4}\right)\right)dt$\\
\textbf{9) (** I)} $\int_{0}^{+\infty}\frac{x\Arctan x}{(1+x^2)^2}\;dx$&
\rule[-6mm]{0mm}{0mm}\textbf{10) (I très long)} $\int_{0}^{+\infty}\frac{x\ln x}{(x^2+1)^a}\;dx$ (calcul pour $a\in\left\{\frac{3}{2},2,3\right\}$)\\
\rule[-7mm]{0mm}{0mm}\textbf{11) (***)} $\int_{0}^{\pi/2}\sqrt{\tan x}\;dx$&\textbf{12) (*** I)} $\int_{0}^{+\infty}\frac{e^{-at}-e^{-bt}}{t}\;dt$ ($0 < a < b$)
\end{tabular}
\end{center}
\finenonce{005716}


\finexercice


\exercice{5717, rouget, 2010/10/16}
\enonce{005717}{}
Deux calculs de $I =\int_{0}^{\pi/2}\ln(\sin x)\;dx$.

\textbf{1) (** I)} En utilisant $J=\int_{0}^{\pi/2}\ln(\cos x)\;dx$, calculer $I$ (et $J$).

\textbf{2) (*** I)} Calculer $P_n=\prod_{k=1}^{n-1}\sin\frac{k\pi}{2n}$ (commencer par $P_n^2$) et en déduire $I$.
\finenonce{005717}


\finexercice


\exercice{5718, rouget, 2010/10/16}
\enonce{005718}{** I}
En utilisant un développement de $\frac{1}{1-t}$, calculer $\int_{0}^{1}\frac{\ln t}{t-1}\;dt$.
\finenonce{005718}


\finexercice


\exercice{5719, rouget, 2010/10/16}
\enonce{005719}{*** I}
\label{ex:rou7bis}
Calculer $\int_{0}^{1}\frac{t-1}{\ln t}\;dt$ (en écrivant $\int_{0}^{x}\frac{t-1}{\ln t}\;dt=\int_{0}^{x}\frac{t}{\ln t}\;dt-\int_{0}^{x}\frac{1}{\ln t}\;dt$).
\finenonce{005719}


\finexercice


\exercice{5720, rouget, 2010/10/16}
\enonce{005720}{}
\textbf{1) (** I)} Trouver un équivalent simple quand $x$ tend vers $+\infty$ de $e^{x^2}\int_{x}^{+\infty}e^{-t^2}\;dt$.

\textbf{2) (***)} Montrer que $\int_{a}^{+\infty}\frac{\cos x}{x}\;dx\underset{a\rightarrow0}{\sim}-\ln a$.

\textbf{3) (*)} Montrer que $\int_{0}^{1}\frac{1}{x^3+a^2}\;dx\underset{a\rightarrow+\infty}{\sim}\frac{1}{a^2}$.
\finenonce{005720}


\finexercice


\exercice{5721, rouget, 2010/10/16}
\enonce{005721}{***}
Etude complète de $f~:~x\mapsto\int_{x}^{x^2}\frac{1}{\ln t}\;dt$.
\finenonce{005721}


\finexercice


\exercice{5722, rouget, 2010/10/16}
\enonce{005722}{***}
(Hors programme) Convergence et calcul de $\int_{1}^{+\infty}\frac{(-1)^{E(x)}}{x}\;dx$.
\finenonce{005722}


\finexercice


\exercice{5723, rouget, 2010/10/16}
\enonce{005723}{***}
Soit $f$ définie, continue, positive et décroissante sur $[1,+\infty[$, intégrable sur $[1,+\infty[$.

\begin{enumerate}
 \item  Montrer que $xf(x)$ tend vers $0$ quand $x$ tend vers $+\infty$.

\item  Existence et calcul de $\int_{1}^{+\infty}x(f(x+1)-f(x))dx$.
\end{enumerate}
\finenonce{005723}


\finexercice


\exercice{5724, rouget, 2010/10/16}
\enonce{005724}{***}
\begin{enumerate}
 \item  Soit $f$ de classe $C^1$ sur $\Rr^+$ à valeurs dans $\Rr$ telle que l'intégrale $\int_{0}^{+\infty}f(x)\;dx$ converge en $+\infty$. Montrer que $\int_{0}^{+\infty}f'(x)\;dx$ converge en $+\infty$ si et seulement si $f(x)$ tend vers $0$ quand $x$ tend vers $+\infty$.

\item
 \begin{enumerate}
 \item On suppose que $f$ est une fonction de classe $C^2$ sur $\Rr^+$ à valeurs dans $\Rr$ telle que $f$ et $f''$ admettent des limites réelles quand $x$ tend vers $+\infty$. Montrer que $f'$ tend vers $0$ quand $x$ tend vers $+\infty$.

 \item En déduire que si les intégrales $\int_{0}^{+\infty}f(x)\;dx$ et $\int_{0}^{+\infty}f''(x)\;dx$ convergent alors $f$ tend vers $0$ quand $x$ tend vers $+\infty$.
 \end{enumerate}
\end{enumerate}
\finenonce{005724}


\finexercice


\exercice{5725, rouget, 2010/10/16}
\enonce{005725}{***}
Soit $f$ de classe $C^2$ sur $\Rr$ à valeurs dans $\Rr$ telle que $f^2$ et $(f'')^2$ soient intégrables sur $\Rr$. Montrer que $f'^2$ est intégrable sur $\Rr$ et que $\left(\int_{-\infty}^{+\infty}f'^2(x)\;dx\right)^2\leqslant\left(\int_{-\infty}^{+\infty}f^2(x)\;dx\right)\left(\int_{-\infty}^{+\infty}f''^2(x)\;dx\right)$. Cas d'égalité ?
\finenonce{005725}


\finexercice


\exercice{5925, tumpach, 2010/11/11}
\enonce{005925}{}
\begin{enumerate}
\item \emph{Le but de cette question est de montrer que
$\int_{0}^{+\infty} \frac{\sin t}{t}\, dt$ n'est pas absolument
convergente.} Pour $n\in\mathbb{N}$, on pose~:
\begin{eqnarray*}
u_{n} = \int_{n\pi}^{(n+1)\pi} \frac{|\sin t|}{t} \, dt.
\end{eqnarray*}
Montrer que pour $n\geq 0$, $\frac{2}{(n+1)\pi} \leq u_{n}$. En
d\'eduire que $\int_{0}^{+\infty} \frac{|\sin t|}{t}\, dt$ est
divergente. 
\item \emph{Deuxi\`eme formule de la moyenne.}
 Soient $f$ et $g$ deux fonctions Riemann-int\'egrables sur $[a, b]$,
 admettant des primitives not\'ees $F$ et $G$ respectivement. Supposons
 que $F$ est positive et d\'ecroissante. Montrer qu'il existe
 $y\in [a,b]$ tel que~:
 \begin{equation*}
\int_{a}^{b}F(x)g(x)\,dx = F(a)\int_{a}^{y}g(x)\,dx.
 \end{equation*}
\item En d\'eduire que $\int_{0}^{+\infty} \frac{\sin t}{t}\, dt$
est convergente.
\item \emph{Le but de cette question est de
calculer la valeur de cette int\'egrale.} Pour tout nombre r\'eel
$\lambda\geq 0$, on pose~:
\begin{equation*}
\left\{
\begin{array}{lll}
f(t, \lambda) &=& e^{-\lambda t}\,\frac{\sin t}{t}\quad\quad
\text{pour}\quad t> 0\\
f(0, \lambda) &=& 1.
\end{array}\right.
\end{equation*}

  \begin{enumerate}
  \item  Pour $0< x \leq y$, d\'emontrer que l'on a~:
\begin{equation*}
\left|\int_{x}^{y} f(t, \lambda)\, dt\right| \leq \frac{2}{x}
e^{-\lambda x}.
\end{equation*}  

  \item   En d\'eduire que les int\'egrales g\'en\'eralis\'ees
$\,\int_{0}^{+\infty} f(t, \lambda)\, dt\,$ sont convergentes,
uniform\'ement pour $\lambda\geq0$. On pose, pour $\lambda \geq
0$,
\begin{equation*}
F(\lambda) = \int_{0}^{+\infty} e^{-\lambda t}\,\frac{\sin
t}{t}\,dt.
\end{equation*}
D\'emontrer que la fonction $F$ est continue pour $\lambda \geq
0$.

  \item D\'emontrer que la fonction $F$ est d\'erivable pour
$\lambda
> 0$ et que sa d\'eriv\'ee est \'egale \`a l'int\'egrale
g\'en\'eralis\'ee convergente
\begin{equation*}
F'(\lambda) = - \int_{0}^{+\infty} e^{-\lambda t}\sin t \,dt.
\end{equation*}

  \item Calculer cette derni\`ere int\'egrale g\'en\'eralis\'ee, par
exemple en int\'egrant par parties sur $[0, x]$ et en calculant la
limite
quand $x\rightarrow +\infty$. 

  \item En d\'eduire la valeur de $F(\lambda)$ pour $\lambda\geq0$
\`a une constante additive pr\`es. D\'emontrer que
$F(\lambda)\rightarrow 0$ quand $\lambda \rightarrow +\infty$. En
d\'eduire la valeur de la constante additive, puis la valeur de
l'int\'egrale $\int_{0}^{+\infty} \frac{\sin t}{t}\, dt$.

  \end{enumerate}
\end{enumerate}
\finenonce{005925}



\finexercice

\section{127.99 Autre}


\exercice{1273, legall, 1998/09/01}

\enonce{001273}{}
Soit $  f :[0,1]\rightarrow { \Rr}  $ une fonction de classe $  C ^1  .$
Montrer que $  \displaystyle{\lim _{n\rightarrow \infty } \int _0 ^1 \hbox{cos}(nt)f(t)dt=0 }  .$
\finenonce{001273}



\finexercice



\exercice{1274, legall, 1998/09/01}

\enonce{001274}{}
Soit $  f :[0,1]\rightarrow { \Rr}  $ une fonction continue telle que
$  f(0)=0  .$
Montrer que $$  \displaystyle{\lim _{n\rightarrow \infty } \int _0 ^1 f(t^n)dt=0 }  .$$
G\'en\'eraliser au cas o\`u $  f(0)  $ est quelconque.
\finenonce{001274}



\finexercice



\exercice{1275, legall, 1998/09/01}

\enonce{001275}{}
Soit $  f :[a,b]\rightarrow { \Rr}  $ une fonction int\'egrable.
\begin{enumerate}
    \item Montrer que $  f  $ est born\'ee. On pose $  \displaystyle{ M=\sup _{x\in [a,b]}\vert
f(x)\vert }  .$
    \item Soient $  x  $ et $  y\in [a,b]  $
Montrer que $  \displaystyle{ \vert \int _x ^y f(t)dt\vert \leq M\vert x-y\vert }  .$
En d\'eduire que l'application $  \displaystyle{F:x\mapsto \int _a ^x f(t)dt }  $ est continue sur
$  [a, b]  .$
    \item Soit $  x_0\in [a,b]  .$ Montrer que si $  f  $ est continue en $  x_0  $ alors $  F  $
est d\'erivable en $  x_0  .$
\end{enumerate}
\finenonce{001275}



\finexercice



\exercice{1276, legall, 1998/09/01}

\enonce{001276}{}
Soit $  f :[0,1]\rightarrow { \Rr}  $ une fonction continue.
Montrer que
$$  \displaystyle{\lim _{n\rightarrow \infty }
\int _0 ^1 nt^nf(t^n)dt = \int _0 ^1 f(t)dt }  .$$
(On pourra faire le changement de variable $  u=t^n  $).
\finenonce{001276}



\finexercice



\exercice{1277, legall, 1998/09/01}

\enonce{001277}{}
Soit $  f :[a,b]\rightarrow { \Rr}  $ une fonction de classe
$  C^1  $ telle que $  f(a)=f(b)=0  .$ Posons ${M=\sup _{x\in [a,b]}
\vert f'(x)\vert }  .$ Montrer que
$  \displaystyle{
\vert \int _a^bf(t)dt\vert \leq M\frac{ (b-a)^2}{ 4} }  .$ (Indication~: faire
des d\'e\-ve\-lop\-pements limit\'es de
$  \displaystyle{
x\mapsto  \int _a^xf(t)dt }  $ et $  \displaystyle{
x\mapsto  \int _x^bf(t)dt }  ).$
\finenonce{001277}



\finexercice



\exercice{1278, gourio, 2001/09/01}

\enonce{001278}{}
Soit $f$ continue sur $[0,1]$ avec $f(1)\neq 0$, montrer :
$$\int_{0}^{1}x^{n}f(t)dt\sim \frac{f(1)}{n}. $$
En d\'{e}duire :
$$\lim\limits_{n\rightarrow \infty }\int_{0}^{n}e^{-2t}\left( 1-\frac{t}{n}\right) ^{n}dt$$
On posera $u=1-\frac{1}{n} $ puis $v=ue^{2(u-1)}.$
\finenonce{001278}



\finexercice



\exercice{1279, gourio, 2001/09/01}

\enonce{001279}{}
Donner un d\'{e}veloppement :
$$\int_{0}^{1}\frac{e^{t}}{1+t^{n}}dt=a+\frac{b}{n}+o(\frac{1}{n}). $$
\finenonce{001279}



\finexercice



\exercice{2327, matexo1, 2002/02/01}
\enonce{002327}{}
Calculer les int\'egrales
$$ I = \int_0^x {\exp 2t \cos 3t ~dt} \qquad \mbox{\rm et}\qquad
J = \int_0^x{\exp 2t \sin3t ~dt}.$$
\finenonce{002327}


\finexercice


\exercice{2328, matexo1, 2002/02/01}
\enonce{002328}{}
Soient $a \neq 0$ un r\'eel, et $y>x>0$.

\begin{enumerate}
\item Calculer la valeur de 
$$ I(x,y) = \int_x^y{ \ln \left( 1+ {a^2 \over x^2} \right) dt}. $$
\item Montrer que $I(x,y)$ a une limite $I_0(y)$ quand $x$ tend
vers z\'ero et la calculer.
\item Montrer que $I_0(y)$ a une limite quand $y$ tend vers
$+\infty$ et la calculer.
\end{enumerate}
\finenonce{002328}


\finexercice


\exercice{4238, quercia, 2010/03/12}
\enonce{004238}{École de l'air 94}

On note $I_n =  \int_0^\pi \frac{\cos nx}{2-\cos x}\,d x$,
        $J_n =  \int_0^{\pi/2} \frac{\cos nx}{2-\cos x}\,d x$,
        $K_n =  \int_0^{\pi/2} \frac{\cos nx}{2+\cos x}\,d x$.\par
Montrer que pour tout $n\in \N$, on a $I_n = J_n + (-1)^nK_n$ et
$I_{n+1} = 4I_n - I_{n-1}$.
En déduire $I_n$ en fonction de $n$.


\finenonce{004238}



\finexercice


\exercice{4239, quercia, 2010/03/12}
\enonce{004239}{Calcul d'intégrale}

Calculer pour tout $n \in \N^*$~:
$I_n =  \int_{x=0}^\pi \frac{dx}{1+\cos^2(nx)}$.


\finenonce{004239}



\finexercice


\exercice{4240, quercia, 2010/03/12}
\enonce{004240}{$\Arcsin$ et $\Arccos$}

Simplifier $ \int_{t=0}^{\sin^2x} \Arcsin\sqrt t\,d t +  \int_{t=0}^{\cos^2x} \Arccos\sqrt t\,d t$.
\finenonce{004240}



\finexercice

\exercice{4244, quercia, 2010/03/12}
\enonce{004244}{Calcul de limite}

Étudiez la limite de la suite définie par
$u_n=\frac n2-\sum_{k=1}^n\frac{n^2}{(n+k)^2}\cdotp$

\finenonce{004244}



\finexercice

\exercice{4246, quercia, 2010/03/12}
\enonce{004246}{\'Echange de décimales}

Soit $f : {[0,1]} \to {[0,1]}$ définie par
$f(0,a_1a_2a_3\dots) = 0,a_2a_1a_3\dots$ (échange des deux $1^{\text{ères}}$
décimales).

Montrer que $f$ est continue par morceaux et calculer $ \int_{t=0}^1 f(t)\,d t$.
\finenonce{004246}



\finexercice


\exercice{4247, quercia, 2010/03/12}
\enonce{004247}{$\int f(t)\cos(t)\,d t$}

Soit $f : {[0,2\pi]} \to \R$ convexe de classe $\mathcal{C}^2$.
Quel est le signe de $I =  \int_{t=0}^{2\pi} f(t)\cos t\,d t$ ?
\finenonce{004247}



\finexercice


\exercice{4248, quercia, 2010/03/12}
\enonce{004248}{Convexité}

Soit $f:\R \to \R$ convexe et $g(x) =  \int_{t=x-1}^{x+1} f(t)\,d t$.
Montrer que $g$ est convexe.
\finenonce{004248}



\finexercice

\exercice{4249, quercia, 2010/03/12}
\enonce{004249}{Expression d'une primitive $n$-ème de $f$}

Soit $f : {[a,b]} \to \R$ continue et
$g(x) =  \int_{t=a}^x \frac {(x-t)^{n-1}}{(n-1)!} f(t)\,d t$.
Montrer que $g^{(n)} = f$.
\finenonce{004249}



\finexercice


\exercice{4250, quercia, 2010/03/12}
\enonce{004250}{Théorème de division}

Soit $f:\R \to \R$ de classe $\mathcal{C}^{n+p}$ telle que
$f(0) = f'(0) = \dots = f^{(n-1)}(0) = 0$.

On pose $g(x) = \frac {f(x)}{x^n}$ pour $x \ne 0$ et
$g(0) = \frac {f^{(n)}(0)}{n!}$.

\begin{enumerate}
  \item \'Ecrire $g(x)$ sous forme d'une intégrale.
  \item En déduire que $g$ est de classe $\mathcal{C}^p$ et
    $|g^{(p)}(x)| \le \frac {p!}{(p+n)!}{\sup\{|f^{(n+p)}(tx)| \text{ tel que } 0 \le t \le 1 \}}$.
\end{enumerate}
\finenonce{004250}



\finexercice


\exercice{4251, quercia, 2010/03/12}
\enonce{004251}{Fonction absolument monotone}

Soit $f : {[0, a[} \to \R$ de classe $\mathcal{C}^\infty$ telle que $f$ et toutes ses dérivées
sont positives sur $[0, a[$.

\begin{enumerate}
  \item   Montrer que la fonction
      $g_n : x \longmapsto \frac 1{x^{n}}\left(f(x) - f(0) - \dots -
      \frac {x^{n-1}}{(n-1)!}f^{(n-1)}(0)\right)$
      est croissante.
      

  \item   On fixe $r \in {]0,a[}$. Montrer que la série de Taylor de $f$ converge
      vers $f$ sur $[0,r[$.

\end{enumerate}
\finenonce{004251}



\finexercice


\exercice{4252, quercia, 2010/03/12}
\enonce{004252}{Deuxième formule de la moyenne}

Soient ${f,g} : {[a,b]} \to \R$ continues, $f$ positive décroissante.

On note $G(x) =  \int_{t=a}^x g(t)\,d t$, et
$$M = \sup\{ G(x),\ x \in {[a,b]} \}\qquad m = \inf\{ G(x),\ x \in {[a,b]} \}.$$

\begin{enumerate}
  \item  On suppose ici que $f$ est de classe $\mathcal{C}^1$.
     Démontrer que $mf(a) \le  \int_{t=a}^b f(t)g(t)\,d t \le Mf(a)$.

  \item  Démontrer la même inégalité si $f$ est seulement continue, en admettant
     qu'elle est limite uniforme de fonctions de classe $\mathcal{C}^1$ décroissantes.

  \item  Démontrer enfin qu'il existe $c \in {[a,b]}$ tel que
     $ \int_{t=a}^b f(t)g(t)\,d t = f(a)  \int_{t=a}^c g(t)\,d t$.

\end{enumerate}
\finenonce{004252}



\finexercice


\exercice{4253, quercia, 2010/03/12}
\enonce{004253}{Inégalité de la moyenne}

Soient ${f,g} : {[a,b]} \to {\R}$ continues, $f$ décroissante, et
$0 \le g \le 1$.
On note $G(x) = a +  \int_{t=a}^x g(t)\,d t$.

Démontrer que $ \int_{t=a}^b fg(t)\,d t \le  \int_{t=a}^{G(b)} f(t)\,d t$.

\finenonce{004253}



\finexercice


\exercice{4254, quercia, 2010/03/12}
\enonce{004254}{Une inégalité}

Soit $f : {[a,b]} \to \R$ de classe $\mathcal{C}^1$ telle que
$f(a) = 0$ et $\forall\ t \in {[a,b]},\ 0 \le f'(t) \le 1$.
Comparer $ \int_{t=a}^b f^3(t)\,d t$ et $\left(  \int_{t=a}^b f(t)\,d t \right)^2$.

On introduira les fonctions : $F(x) =  \int_{t=a}^x f(t)\,d t$,
$G(x) =  \int_{t=a}^x f^3(t)\,d t$, et $H = F^2 - G$.
\finenonce{004254}



\finexercice


\exercice{4255, quercia, 2010/03/12}
\enonce{004255}{Intégrales de Wallis}

On note $I_n =  \int_{t=0}^{\pi/2} \cos^nt\,d t$.

\begin{enumerate}
  \item   Comparer $I_n$ et $ \int_{t=0}^{\pi/2} \sin^nt\,d t$.
  \item   En coupant $\left[0,\frac\pi2\right]$ en $[0,\alpha]$ et
      $\left[\alpha,\frac\pi2\right]$, démontrer que $I_n \to 0$ pour $n\to\infty$.
  \item   Chercher une relation de récurrence entre $I_n$ et $I_{n+2}$.
      En déduire $I_{2k}$ et $I_{2k+1}$ en fonction de $k$.
  \item   Démontrer que $nI_nI_{n-1} = \frac \pi2$.
  \item   Démontrer que $I_n \sim I_{n-1}$ et en déduire un équivalent simple de $I_n$
      puis de $C_{2n}^n$ pour $n \to \infty$.

\end{enumerate}
\finenonce{004255}



\finexercice


\exercice{4256, quercia, 2010/03/12}
\enonce{004256}{Norme $L^\infty$}

Soit $f : {[a,b]} \to {\R^+}$ continue non identiquement nulle.
On pose $I_n =  \int_{t=a}^b f^n(t)\,d t$ et $u_n = \sqrt[n]{I_n}$.

Soit $M = \max\{ f(x)\text{ tel que } a \le x \le b \}$ et $c \in {[a,b]}$ tel que $f(c) = M$.

\begin{enumerate}
  \item   Comparer $M$ et $u_n$.
  \item   En utilisant la continuité de $f$ en $c$, démontrer que :
      $\forall\ \varepsilon \in{]0,M[}$ il existe $\delta > 0$ tel que
      $I_n \ge \delta(M-\varepsilon)^n$.
  \item   En déduire $\lim_{n\to\infty} u_n$.
\end{enumerate}
\finenonce{004256}



\finexercice


\exercice{4257, quercia, 2010/03/12}
\enonce{004257}{Lemme de Lebesgue}

Soit $f : {[a,b]} \to {\R}$ continue. Montrer que $ \int_{t=a}^b f(t)\cos(nt)\,d t
\to 0$, (lorsque $n\to\infty$) $\dots$

\begin{enumerate}
  \item si $f$ est de classe $\mathcal{C}^1$.
  \item si $f$ est en escalier.
  \item si $f$ est continue.

\end{enumerate}
\finenonce{004257}



\finexercice


\exercice{4258, quercia, 2010/03/12}
\enonce{004258}{Plus grande fonction convexe minorant $f$}

\begin{enumerate}
  \item
Soit $(f_i)$ une famille de fonctions convexes sur un intervalle $I$.

On suppose que : $\forall\ x \in I,\ f(x) = \sup(f_i(x))$ existe.
Montrer que $f$ est convexe.

  \item
Soit $f : I \to {\R}$ minorée. Montrer qu'il existe une plus grande fonction
convexe minorant $f$. On la note $\tilde f$.

  \item
Soit $f : {[0,1]} \to {\R^+}$ croissante.
Montrer que $ \int_{t=0}^1 \tilde f(t)d t \ge \frac 12 \int_{t=0}^1 f(t)d t$
(commencer par le cas où $f$ est en escalier).
\end{enumerate}
\finenonce{004258}



\finexercice


\exercice{4259, quercia, 2010/03/12}
\enonce{004259}{Centrale PC 1998}
Soit $f : {[a,b]} \to {\R^{+*}}$ continue.

\begin{enumerate}
  \item Montrer qu'il existe une subdivision de $[a,b]$~:
    ${a=x_0 < x_1 < \dots < x_n = b}$ telle que~:

    $\forall\ k \in [[0,n-1]],\  \int_{t=x_k}^{x_{k+1}}f(t)\,d t =
        \frac1n  \int_{t=a}^{b}f(t)\,d t$.
  \item Étudier $\lim_{n\to\infty} \frac1n \sum_{k=0}^{n-1}f(x_k)$.

\end{enumerate}
\finenonce{004259}



\finexercice


\exercice{4260, quercia, 2010/03/12}
\enonce{004260}{Mines MP 2000}

Soit ${f} : {\R} \to {\C}$ de classe $\mathcal{C}^1$, $2\pi$ périodique, ne s'annulant pas.
Montrer que $I(f)=\frac{1}{2\pi i} \int_0^{2\pi} \frac{f'}{f}$ est un entier.

\finenonce{004260}



\finexercice


\exercice{4261, quercia, 2010/03/12}
\enonce{004261}{Fonctions affines}
Soit $E = \mathcal{C}([a,b])$, et $F=\{f\in \mathcal{C}^2([a,b]),\text{ tel que } f(a)=f'(a)=f(b)=f'(b)=0\}$.

\begin{enumerate}
  \item Soit $f\in E$. Montrer qu'il existe $g\in F$ vérifiant $g''=f$ si et
    seulement si $ \int_{x=a}^b f(x)\,d x =  \int_{x=a}^b xf(x)\,d x = 0$.
    
  \item Soit $f\in E$ telle que $ \int_{x=a}^b f(x)g''(x)\,d x = 0$ pour toute
    fonction $g\in F$. Montrer que $f$ est affine.
    

\end{enumerate}
\finenonce{004261}



\finexercice


\exercice{4262, quercia, 2010/03/12}
\enonce{004262}{Mines MP 2001}

Soit $a<0<b$ et $f$ continue sur $[0,1]$, à valeurs dans $[a,b]$
telle que $\int_0^1 f=0$. Montrer que $\int_0^1 f^2\le -ab$.

\finenonce{004262}



\finexercice


\exercice{5445, rouget, 2010/07/10}
\enonce{005445}{**I}
\begin{enumerate}
\item  Soit $f$ une application de classe $C^1$ sur $[0,1]$ telle que $f(1)\neq0$.

Pour $n\in\Nn$, on pose $u_n=\int_{0}^{1}t^nf(t)\;dt$. Montrer que $\lim_{n\rightarrow +\infty}u_n=0$ puis déterminer un équivalent simple de $u_n$ quand $n$ tend vers $+\infty$ (étudier $\lim_{n\rightarrow +\infty}nu_n$).
\item  Mêmes questions en supposant que $f$ est de classe $C^2$ sur $[0,1]$ et que $f(1)=0$ et $f'(1)\neq0$.
\end{enumerate}
\finenonce{005445}


\finexercice


\exercice{5448, rouget, 2010/07/10}
\enonce{005448}{**I Le lemme de \textsc{Lebesgue}}
\begin{enumerate}
\item  On suppose que $f$ est une fonction de classe $C^1$ sur $[a,b]$. Montrer que $\lim_{\lambda\rightarrow +\infty}\int_{a}^{b}\sin(\lambda t)f(t)\;dt=0$.
\item  (***) Redémontrer le même résultat en supposant simplement que $f$ est continue par morceaux sur $[a,b]$ (commencer par le cas des fonctions en escaliers).
\end{enumerate}
\finenonce{005448}


\finexercice


\exercice{5452, rouget, 2010/07/10}
\enonce{005452}{***}
Soit $f$ une fonction de classe $C^1$ sur $[0,1]$ telle que $f(0)=0$. Montrer que $2\int_{0}^{1}f^2(t)\;dt\leq\int_{0}^{1}{f'}^2(t)\;dt$. 
\finenonce{005452}


\finexercice


\exercice{5454, rouget, 2010/07/10}
\enonce{005454}{***}
Soit $a$ un réel strictement positif et $f$ une application de classe $C^1$ et strictement croissante sur $[0,a]$ telle que $f(0)=0$. Montrer que $\forall x\in[0,a],\;\forall y\in[0,f(a)],\;xy\leq\int_{0}^{x}f(t)\;dt+\int_{0}^{y}f^{-1}(t)\;dt$.
\finenonce{005454}


\finexercice


\exercice{5455, rouget, 2010/07/10}
\enonce{005455}{**}
Soit $f$ continue sur $[0,1]$ telle que $\int_{0}^{1}f(t)\;dt=\frac{1}{2}$. Montrer que $f$ admet un point fixe.
\finenonce{005455}


\finexercice


\exercice{5456, rouget, 2010/07/10}
\enonce{005456}{**}
Soient $f$ et $g$ deux fonctions continues par morceaux et positives sur $[0,1]$ telles que

$\forall x\in[0,1],\;f(x)g(x)\geq1$. Montrer que $(\int_{0}^{1}f(t)\;dt)(\int_{0}^{1}g(t)\;dt)\geq1$.
\finenonce{005456}


\finexercice


\exercice{5461, rouget, 2010/07/10}
\enonce{005461}{***}
Trouver toutes les applications continues sur $\Rr$ vérifiant~:~$\forall(x,y)\in\Rr^2,\;f(x)f(y)=\int_{x-y}^{x+y}f(t)\;dt$.  
\finenonce{005461}


\finexercice


\exercice{5462, rouget, 2010/07/10}
\enonce{005462}{***}
Soit $f$ une fonction de classe $C^1$ sur $[a,b]$ telle que $f(a)=f(b)=0$ et soit $M=\mbox{sup}\{|f '(x)|,\;x\in[a,b]\}$. Montrer que $\left|\int_{a}^{b}f(x)\;dx\right|\leq M\frac{(b-a)^2}{4}$. 
\finenonce{005462}


\finexercice


\exercice{5463, rouget, 2010/07/10}
\enonce{005463}{**}
Déterminer les fonctions $f$ continues sur $[0,1]$ vérifiant $\left|\int_{0}^{1}f(t)\;dt\right|=\int_{0}^{1}|f(t)|\;dt$. 
\finenonce{005463}


\finexercice
\section{140.01 Distance, norme, produit scalaire}

\section{140.02 Droites}

\section{141.01 Produit scalaire, produit vectoriel, déterminant}

\section{141.02 Aire, volume}

\section{141.03 Plans}

\section{141.04 Droites de l'espace}

\section{141.05 Distance}

\section{200.01 Forme multilinéaire}


\exercice{1107, barraud, 2003/09/01}

\enonce{001107}{}

On consid\`{e}re l'espace $\mathcal{M}_{n}(\mathbb{K})$ des matrices carr\'{e}es $n\times n$ \`{a} coefficients dans le
corps $\mathbb{K}$. On rappelle que la \emph{trace} $\mathrm{tr}(A)$ d'une matrice $A\in\mathcal{M}_{n}(\mathbb{K})$ est la
somme de ses coefficients diagonaux.

Pour une matrice $M$ donn\'{e}e, on note $\alpha_{M}$ l'application d\'{e}finie par
$$
  \forall X\in\mathcal{M}_{n}(\mathbb{K}),\quad \alpha_{M}(X)=\mathrm{tr}(MX).
$$

\begin{enumerate}
\item
V\'{e}rifier que $\forall M\in\mathcal{M}_{n}(\mathbb{K}),\quad \alpha_{M}\in(\mathcal{M}_{n}(\mathbb{K}))^{*}$.

\bigskip

On note $\phi$ l'application suivante :
$$
 \phi :
  \begin{array}{ccc}
    \mathcal{M}_{n}(\mathbb{K}) & \rightarrow  & (\mathcal{M}_{n}(\mathbb{K}))^{*} \\[3mm]
         M     & \mapsto & \alpha_{M}
  \end{array}
$$

\item
Etudier l'injectivit\'{e} et la surjectivit\'{e} de $\phi$.
\item
En d\'{e}duire que pour toute forme lin\'{e}aire $\alpha\in(\mathcal{M}_{n}(\mathbb{K}))^{*}$, il existe une
matrice $A\in\mathcal{M}_{n}(\mathbb{K})$ telle que :
$$
 \forall X\in\mathcal{M}_{n}(\mathbb{K}),\quad \alpha(X)=\mathrm{tr}(AX).
$$
\item
D\'{e}terminer toute les formes lin\'{e}aires $\alpha\in(\mathcal{M}_{n}(\mathbb{K}))^{*}$ telles que
$$
 \forall (X,Y)\in(\mathcal{M}_{n}(\mathbb{K}))^{2},\quad \alpha(XY)=\alpha(YX).
$$
\end{enumerate}
\finenonce{001107}



\finexercice



\exercice{1108, barraud, 2003/09/01}

\enonce{001108}{}
On note $\R_{n}[X]$ l'espace vectoriel des polyn\^{o}mes \`{a} coefficients r\'{e}els de degr\'{e}
inf\'{e}rieur ou \'{e}gal \`{a} $n$. 

Pour chaque $i\in\{0,\ldots,n\}$, on note $\alpha_{i}$ l'application
$$
 \alpha_{i} :
  \begin{array}{ccc}
    \R_{n}[X] & \rightarrow  & \R \\[3mm]
        P     & \mapsto & P(x_{i})
  \end{array}
$$
\begin{enumerate}
\item
V\'{e}rifier que chaque $\alpha_{i}$ est une forme lin\'{e}aire sur $\R_{n}[X]$
\item
On note $G$ l'espace engendr\'{e} par $\alpha_{1},\ldots,\alpha_{n}$. D\'{e}terminer $G^{\circ}$. En d\'{e}duire que la famille $(\alpha_{0},\ldots,\alpha_{n})$ est une base de $(\R_{n}[X])^{*}$.
\item
Montrer que la famille $(\alpha_{0},\ldots,\alpha_{n})$ est une base de $(\R_{n}[X])^{*}$.
\item
Montrer qu'il existe des r\'{e}els $\lambda_{0},\ldots,\lambda_{n}$ tels que
$$
 \forall P\in\R_{n}[X]\;\; \int_{0}^{1}{P(t)dt} = \sum_{i=0}^{n}\lambda_{i}P(x_{i})
$$
\item
Montrer qu'il existe une unique famille de polyn\^omes $(P_0,\ldots,P_n)$ de $R_n[X]$ telle que
  $
 \forall (i,j)\in\{0,\ldots,n\}^{2}\;\;\;P_{i}(x_{j})=
         \begin{cases}
                 1 &\text{ si $j=i$}\\
                 0 &\text{ sinon}
         \end{cases}
  $
\item

En d\'eduire que pour toute fonction continue $f$ de $\R$ dans $\R$, il existe un polyn\^ome $P$ de degr\'e $n$,  qui interpole $f$ en chaque point $x_i$, c'est \`a dire qui satisfait :
$$
  \forall i\in\{!,...,n\}\quad  P(x_i)=f(x_i).
$$
\end{enumerate}
\finenonce{001108}



\finexercice



\exercice{1109, barraud, 2003/09/01}

\enonce{001109}{}
Dans chacun des cas ci-dessous, dire si l'application $\phi$ de
$\R^{3}\times\R^{3}\times\R^{3}$ dans $\R$, est multilinéaire.


\begin{align*}
\phi\left( \left(\begin{smallmatrix}  x_{1}\\x_{2}\\x_{3} \end{smallmatrix}\right),
\left(\begin{smallmatrix}  y_{1}\\y_{2}\\y_{3} \end{smallmatrix}\right),
\left(\begin{smallmatrix}  z_{1}\\z_{2}\\z_{3} \end{smallmatrix}\right) \right)
  &= x_{1}+y_{2}+z_{3}
\\
\phi\left(
 \left(\begin{smallmatrix}  x_{1}\\x_{2}\\x_{3} \end{smallmatrix}\right),
  \left(\begin{smallmatrix}  y_{1}\\y_{2}\\y_{3} \end{smallmatrix}\right),
 \left(\begin{smallmatrix}  z_{1}\\z_{2}\\z_{3} \end{smallmatrix}\right)
\right)
  &=  x_{1}y_{3}+y_{2}z_{1}+z_{3}x_{2}
\\
\phi\left(
 \left(\begin{smallmatrix}  x_{1}\\x_{2}\\x_{3} \end{smallmatrix}\right),
  \left(\begin{smallmatrix}  y_{1}\\y_{2}\\y_{3} \end{smallmatrix}\right),
 \left(\begin{smallmatrix}  z_{1}\\z_{2}\\z_{3} \end{smallmatrix}\right)
\right)
  &=  x_{1}y_{2}z_{3}
     +x_{2}y_{3}z_{1}
     +x_{3}y_{1}z_{2}
\\
\phi\left(
 \left(\begin{smallmatrix}  x_{1}\\x_{2}\\x_{3} \end{smallmatrix}\right),
  \left(\begin{smallmatrix}  y_{1}\\y_{2}\\y_{3} \end{smallmatrix}\right),
 \left(\begin{smallmatrix}  z_{1}\\z_{2}\\z_{3} \end{smallmatrix}\right)
\right)
  &=  x_{1}x_{2}x_{3}
     +y_{1}y_{2}y_{3}
     +z_{1}z_{2}z_{3}
\\
\phi\left(
 \left(\begin{smallmatrix}  x_{1}\\x_{2}\\x_{3} \end{smallmatrix}\right),
  \left(\begin{smallmatrix}  y_{1}\\y_{2}\\y_{3} \end{smallmatrix}\right),
 \left(\begin{smallmatrix}  z_{1}\\z_{2}\\z_{3} \end{smallmatrix}\right)
\right)
  &=  x_{1}y_{1}z_{1}
     +x_{2}y_{2}z_{2}
     +x_{3}y_{3}z_{3}
\\
\phi\left(
 \left(\begin{smallmatrix}  x_{1}\\x_{2}\\x_{3} \end{smallmatrix}\right),
  \left(\begin{smallmatrix}  y_{1}\\y_{2}\\y_{3} \end{smallmatrix}\right),
 \left(\begin{smallmatrix}  z_{1}\\z_{2}\\z_{3} \end{smallmatrix}\right)
\right)
  &=  (x_{1}y_{1}+x_{2}y_{2}+x_{3}y_{3})(z_{1}+z_{3})
\\
\phi\left(
 \left(\begin{smallmatrix}  x_{1}\\x_{2}\\x_{3} \end{smallmatrix}\right),
  \left(\begin{smallmatrix}  y_{1}\\y_{2}\\y_{3} \end{smallmatrix}\right),
 \left(\begin{smallmatrix}  z_{1}\\z_{2}\\z_{3} \end{smallmatrix}\right)
\right)
  &=  (x_{1}+2x_{2})(z_{1}+z_{3})
\end{align*}
\finenonce{001109}



\finexercice



\exercice{1110, barraud, 2003/09/01}

\enonce{001110}{}
  Montrer que l'espace des formes bi-linéaires sur $\R^{2}$ est un espace
  vectoriel. En donner une base.
\finenonce{001110}



\finexercice



\exercice{1111, barraud, 2003/09/01}

\enonce{001111}{}
  Donner toutes les formes tri-linéaires alternées sur $\R^{2}$. Plus
  généralement, que dire des formes $m$-linéaires alternées sur un
  espace de dimension $n$ lorsque $m>n$~?
\finenonce{001111}



\finexercice



\exercice{1112, barraud, 2003/09/01}

\enonce{001112}{}
  Soit $A\in\mathcal{M}_{n,n}(\R)$. On considère l'application
  $\Phi_{A}$ suivante~:
  $$
  \Phi_{A}~:
  \begin{array}{ccc}
    (\R^{n})^{n}       &\rightarrow  & \R \\
    M=(C_{1},...,C_{n})&\mapsto &\det(AM)
  \end{array}
  $$
  Montrer que $\Phi_{A}$ est $n$-linéaire.

  Calculer $A \times \left(\begin{array}{c|c}
    \begin{smallmatrix}
      0&1\\1&0
    \end{smallmatrix}
     &0\\
     \hline
    0&\mathrm{id}_{n-2}
  \end{array}\right)
  $. En déduire que
  $\Phi_{A}(e_{2},e_{1},e_{3}...e_{n})=-\Phi_{A}(e_{1},e_{2},e_{3}...e_{n})$.

  Plus généralement, montrer que $\Phi_{A}$ est alternée.

  Montrer que $\Phi_{A}(M)=\det(A)\det(M)$.
  
  En déduire que~:
  $$
  \forall (A,B)\in\mathcal{M}_{n,n}(\R),\qquad \det(AB)=\det(BA)=\det(A)\det(B)
  $$
\finenonce{001112}



\finexercice



\exercice{1113, barraud, 2003/09/01}

\enonce{001113}{}
Dans $\R^{3}$ muni de sa base canonique, on considère les applications
$\omega$ et $\alpha$ suivantes~:
$$
\omega~:
\begin{array}{ccl}
\R^{3}\times\R^{3} & \rightarrow & \R \\
\begin{pmatrix}x_{1}\\x_{2}\\x_{3}\end{pmatrix},
\begin{pmatrix}y_{1}\\y_{2}\\y_{3}\end{pmatrix}
 & \mapsto & x_{1}y_{2}-x_{2}y_{1}
\end{array}
\quad\text{ et }\quad
\alpha~:
\begin{array}{ccl}
\R^{3} & \rightarrow & \R \\
\begin{pmatrix}x_{1}\\x_{2}\\x_{3}\end{pmatrix}
 & \mapsto & x_{3}
\end{array}
$$

\begin{enumerate}
\item 
  Montrer que $\omega$ est antisymétrique et bilinéaire.

A l'aide de $\omega$ et $\alpha$, on définit une nouvelle application,
notée $\omega\wedge\alpha$, de la façon suivante~:
$$
\omega\wedge\alpha~:
\begin{array}{ccl}
\R^{3}\times\R^{3}\times\R^{3} & \rightarrow  & \R \\
(X,Y,Z) & \mapsto &                                     
 \omega(X,Y)\alpha(Z)
+\omega(Y,Z)\alpha(X)
+\omega(Z,X)\alpha(Y)
\end{array}
$$
\item  Montrer que $\omega\wedge\alpha$ est alternée.
\item  Montrer que $\omega\wedge\alpha$ est trilinéaire.
\item  Calculer $\omega\wedge\alpha(e_{1},e_{2},e_{3})$. En déduire que
  $\forall(X,Y,Z)\in(\R^{3})^{3}\ \omega\wedge\alpha(X,Y,Z)=\det(X,Y,Z)$
\end{enumerate}

\finenonce{001113}



\finexercice



\exercice{3427, quercia, 2010/03/10}
\enonce{003427}{Changements de signe}

Soit $A = (a_{ij}) \in \mathcal{M}_n(K)$ et $A' = ((-1)^{i+j}a_{ij})$.
Comparer $\det A$ et $\det A'$.

\finenonce{003427}



\finexercice


\exercice{3428, quercia, 2010/03/10}
\enonce{003428}{Somme des colonnes, Matexo}


Soit $M$ une matrice carrée d'ordre $n$, et $M'$ la matrice
déduite de $M$ en rempla\c{c}ant, pour tout $j$, la $j$-ième colonne
par la somme des colonnes de~$M$ 
d'indices différents de~$j$. Comparer les déterminants de~$M$ et~$M'$.
\finenonce{003428}



\finexercice


\exercice{3438, quercia, 2010/03/10}
\enonce{003438}{Trace d'un endomorphisme}

Soit $E$ un ev de dimension $n$, $f \in \mathcal{L}(E)$,
et $\vec u_1 , \dots, \vec u_n$, $n$ vecteurs de $E$.
On note $\det$ le déterminant dans une base fixée de $E$.
Démontrer que :
$$ \det( f(\vec u_1), \vec u_2, \dots, \vec u_n )
 + \det( \vec u_1, f(\vec u_2), \vec u_3, \dots, \vec u_n )
 + \dots
 + \det( \vec u_1, \vec u_2, \dots, f(\vec u_n ))
 = 
 \det( \vec u_1, \vec u_2, \dots, \vec u_n ) \mathrm{tr}(f).
$$


\finenonce{003438}



\finexercice

\exercice{5635, rouget, 2010/10/16}
\enonce{005635}{**}
Soient $A=(a_{i,j})_{1\leqslant i,j\leqslant n}$ une matrice carrée et $B= (b_{i,j})_{1\leqslant i,j\leqslant n}$ où $b_{i,j}=(-1)^{i+j}a_{i,j}$. Calculer $\text{det}(B)$ en fonction de $\text{det}(A)$. 
\finenonce{005635}


\finexercice


\exercice{5636, rouget, 2010/10/16}
\enonce{005636}{***I}
On définit par blocs une matrice $A$ par $A=\left(
\begin{array}{cc}
B&D\\
0&C
\end{array}
\right)$ où $A$, $B$ et $C$ sont des matrices carrées de formats respectifs $n$, $p$ et $q$ avec $p+q=n$. Montrer que $\text{det}(A)=\text{det}(B)\times\text{det}(C)$.
\finenonce{005636}


\finexercice

\section{200.02 Calcul de déterminants}


\exercice{1114, legall, 1998/09/01}

\enonce{001114}{}
 Calculer les d\'eterminants des matrices suivantes : $
\begin{pmatrix}
 a & c & c & b \cr c & a & b & c \cr
 c & b & a & c \cr b & c & c & a \cr \end{pmatrix} $
$ \begin{pmatrix}
 c & a & b & c \cr a & c & c & b \cr
 b & c & c & a \cr c & b & a & c \cr \end{pmatrix} $
$ \begin{pmatrix}
 a & x & y & z \cr b & x & y & z \cr
 c & x' & y' & z' \cr d & x' & y' & z' \cr \end{pmatrix} $
$ \begin{pmatrix}
 1+a & b & a & b \cr b & 1+a & b & a \cr
 a & b & 1+a & b \cr b & a & b & 1+a \cr \end{pmatrix} $
$ \begin{pmatrix}
 a & a & a^2 & b+c+d \cr a & b & b^2 & c+d+a \cr
 a & c & c^2 & d+a+b \cr a & d & d^2 & a+b+c \cr \end{pmatrix} $
$ \begin{pmatrix}
 1 & 0 & a & a^2 \cr 0 & 1 & b & b^2 \cr
 1 & 0 & c & c^2 \cr 0 & 1 & d & d^2 \cr \end{pmatrix} $
\finenonce{001114}



\finexercice



\exercice{1115, legall, 1998/09/01}

\enonce{001115}{}
 Calculer, pour tout $ t\in {\Rr} $ le rang des
matrices $ M_t=\begin{pmatrix}
 1 & t & 1 \cr t & 1 & 1 \cr
 1 & t & 1 \cr \end{pmatrix} $ et
$ N_t=\begin{pmatrix}
 1 & 1 & t \cr 1 & t & 1 \cr
 t & 1 & 1 \cr \end{pmatrix}.$
\finenonce{001115}



\finexercice



\exercice{1116, legall, 1998/09/01}

\enonce{001116}{}
\begin{enumerate}
    \item Soient $ A \in M_p({\Rr}) $ et $ B \in
M_q({\Rr}).$ Calculer
(en fonction de $ \hbox{det}(A) $ et $ \hbox{det}(B) $) le
d\'eterminant de la matrice $ M =\begin{pmatrix}A & 0\cr 0 & B \cr \end{pmatrix} \in
M_{p+q}({\Rr}).$ (On pourra pour cela
d\'ecomposer $ M $ comme produit de deux matrices de d\'eterminant
\'evident et utiliser la multiplicativit\'e du d\'eterminant.)
    \item Soient $ A \in M_p({\Rr}) ,$ $ B \in M_q({\Rr}) $ et $ C
\in M_{p,q}({\Rr}) .$
Calculer le d\'eterminant de la matrice $ M =\begin{pmatrix}A & C\cr 0 & B \cr
\end{pmatrix} \in M_{p+q}({\Rr}) .$
(On pourra g\'en\'eraliser la m\'ethode de 1.)
\end{enumerate}
\finenonce{001116}



\finexercice



\exercice{1117, ridde, 1999/11/01}

\enonce{001117}{}
Sans calcul, montrer que $\begin{vmatrix}
2 & 0 &4 \\5&2&7\\2&5&5\\
\end{vmatrix}$ est divisible par $17$.
\finenonce{001117}



\finexercice



\exercice{1118, ridde, 1999/11/01}

\enonce{001118}{}
Soit $\Delta (x) = \det (a_{i, j} (x))$ de taille $n = 2$ ou $3$ avec $a_{i, j}$
 des fonctions d\'erivables.
\begin{enumerate}

\item Montrer que $\Delta' (x)$ est la somme des $n$ d\'eterminants obtenus en
remplaçant successivement dans $\Delta (x)$ chaque colonne par sa d\'eriv\'ee.

\item Calculer $\begin{vmatrix} x + a_1 &x&x\\x&x + a_2&x\\x&x&x + a_3 \\
\end{vmatrix}$ et $\begin{vmatrix} 1 + x&1&1\\1&1 + x&1\\1&1&1 + x\\ \end{vmatrix}$.
\end{enumerate}
\finenonce{001118}



\finexercice



\exercice{1119, ridde, 1999/11/01}

\enonce{001119}{}
Calculer $\begin{vmatrix}
1&1&1\\x&y&z\\x^2&y^2&z^2\\
\end{vmatrix}$ et d\'eterminer la condition d'inversibilit\'e de la matrice.
\finenonce{001119}



\finexercice



\exercice{1120, ridde, 1999/11/01}

\enonce{001120}{}
La famille $ (2, 1, 0)$, $ (1, 3, 1)$, $ (5, 2, 1)$ est-elle libre ?
\finenonce{001120}



\finexercice



\exercice{1121, ridde, 1999/11/01}

\enonce{001121}{}
Calculer $\begin{vmatrix}
a&b&c\\c&a&b\\b&c&a
\end{vmatrix}$.
\finenonce{001121}



\finexercice



\exercice{1122, ridde, 1999/11/01}

\enonce{001122}{}
Calculer $\begin{vmatrix}
1&\sin x&\cos x\\ 1&\sin y&\cos y\\1&\sin z&\cos z\\
\end{vmatrix}$
\finenonce{001122}



\finexercice



\exercice{1123, liousse, 2003/10/01}

\enonce{001123}{}
 Soit $n$ un entier sup\'erieur ou 
\'egal \`a 3. On se place dans $\Rr^{n}.$ On note 
$e_i$ le vecteur de $\Rr^{n}$ dont la i-i\`eme  composante est 
\'egale \`a 1 et toutes les autres sont nulles. \'Ecrire la matrice $n\times n$ 
dont les vecteurs colonnes $C_i$ sont donn\'es par 
$C_i= e_i+e_n$ pour $1\leq i\leq n-1$ et $C_n=e_1+e_2+e_n.$ 
Calculer alors son d\'eterminant.
\finenonce{001123}



\finexercice



\exercice{1124, liousse, 2003/10/01}

\enonce{001124}{}
 On note $a,$ $b,$ $c$ des r\'eels. 
Calculer les d\'eterminants suivants.
$$D_1=\left\vert\begin{array}{cccc}1&0&0&1\\0&1&0&0\\1&0&1&1\\
2&3&1&1\end{array}\right\vert\hbox{, }
D_2=\left\vert\begin{array}{cccc}a+b+c&b&b&b\\c&a+b+c&b&b\\c&c&a+b+c&b\\
c&c&c&a+b+c\end{array}\right\vert\hbox{, }
D_3=\left\vert\begin{array}{ccccc}1&0&3&0&0\\0&1&0&3&0\\a&0&a&0&3\\
b&a&0&a&0\\0&b&0&0&a\end{array}\right\vert$$
G\'en\'eraliser le calcul de $D_2$ \`a un d\'eterminant $n\times n$ 
du m\^eme type. 
\finenonce{001124}



\finexercice



\exercice{1125, liousse, 2003/10/01}

\enonce{001125}{}
 On note $a_1,\cdots,a_n$ 
des r\'eels. Calculer les d\'eterminants $n\times n$ suivants.
$$D_1=\left\vert\begin{array}{cccc}1&1&\cdots&1\\a_1&a_2&\cdots&a_n\\a_1^2&
a_2^2&\cdots&a_n^2\\
\vdots&&&\vdots\\a_1^{n-1}&a_2^{n-1}&\cdots&a_n^{n-1}\end{array}\right\vert
\hbox{, }
D_2=\left\vert\begin{array}{ccccc}a_1&a_2&a_3&\cdots&a_n\\a_2&a_2&a_3&\cdots&a_n\\a_3&a_3&
a_3&\cdots&a_n\\
\vdots&&&&\vdots\\a_n&a_n&a_n&\cdots&a_n\end{array}\right\vert$$
\finenonce{001125}



\finexercice



\exercice{1126, liousse, 2003/10/01}

\enonce{001126}{}
 Montrer que 
\[
\left| 
\begin{array}{ccc}
\cos a & \cos b & \cos c \\ 
\sin a & \sin b & \sin c \\ 
1 & 1 & 1
\end{array}
\right| =\sin \left( c-b\right) +\sin \left( b-a\right) +\sin \left(
a-c\right) =4\sin \frac{c-b}{2}\sin \frac{b-a}{2}\sin \frac{a-c}{2} 
\] 
\finenonce{001126}



\finexercice



\exercice{1127, liousse, 2003/10/01}

\enonce{001127}{}
 Soient $a,$ $b$ deux r\'eels distincts. 
Calculer le d\'eterminant suivant.
$$D_1=\left\vert\begin{array}{ccccc}a&b&\cdots&b&b\\b&a&\cdots&b&b\\
\vdots&&&&\vdots\\b&b&\cdots&a&b\\b&b&\cdots&b&a\end{array}\right\vert$$
\finenonce{001127}



\finexercice



\exercice{1128, liousse, 2003/10/01}

\enonce{001128}{}
 Calculer le d\'eterminant 
de la matrice suivante :
$$\left(\begin{array}{cccc}
m&0&1&2m\\
1&m&0&0\\
0&2m+2&m&1\\
m&0&0&m\\
\end{array}
\right).$$
Calculer alors, suivant la valeur du 
param\`etre $m$, le rang de cette matrice.
\finenonce{001128}



\finexercice



\exercice{1129, barraud, 2003/09/01}

\enonce{001129}{}
  Calculer le déterminant
$$\Delta_{n}=\left\vert
  \begin{matrix}
    3 &1     &0     &      &0\\
    0 &3     &1     &\ddots& \\
    -4&0     &3     &\ddots&0\\
      &\ddots&\ddots&\ddots&1\\
    0 &      &4     &0     &3
    \end{matrix}
\right\vert
$$
en fonction de $n$. (vérifier que $-1$ est racine de $X^{3}-3X^{2}+4$)
\finenonce{001129}



\finexercice



\exercice{1130, barraud, 2003/09/01}

\enonce{001130}{}
Calculer les déterminants suivants :
$$
 \Delta_1=
 \left\vert
   \begin{matrix}
     1&2&3&4 \\
     2&3&4&1 \\
     3&4&1&2 \\
     4&1&2&3
   \end{matrix}
 \right\vert
\qquad
 \Delta_2=
 \left\vert
   \begin{matrix}
     a&a&b&0 \\
     a&a&0&b \\
     c&0&a&a \\
     0&c&a&a
   \end{matrix}
 \right\vert
\qquad
 \Delta_3=
 \left\vert
   \begin{matrix}
     a_1   &a_2   &\cdots&a_n    \\
     a_1   &a_1   &\ddots&\vdots \\
     \vdots&\ddots&\ddots&a_2    \\
     a_1   &\cdots&a_1   &a_1
   \end{matrix}
 \right\vert
$$
\finenonce{001130}



\finexercice



\exercice{1131, barraud, 2003/09/01}

\enonce{001131}{}
Soit $(a,x,y)\in\R^3$.
 Pour $n\in\N$, $n\geq2$, on note $A_n$ le déterminant suivant :
$$
 A_n=
 \left\vert
   \begin{matrix}
     a     &x     &\cdots&x   \\
     y     &a     &0     &    \\
     \vdots&0     &\ddots&    \\
     y     &      &      &a
   \end{matrix}
 \right\vert
$$
Montrer que $\forall n\in\N,n\geq3,\;A_n=aA_{n-1}-xya^{n-2}$.
En déduire une expression de $A_n$ en fonction de 
$n,a,x$ et $y$.
\finenonce{001131}



\finexercice



\exercice{1132, barraud, 2003/09/01}

\enonce{001132}{}
Soit $(a,b)\in\R^2$ avec $a\neq b$. Pour $n\in\N$, $n\geq2$, on note $B_n$ le déterminant suivant :
$$
 B_n=
 \left\vert
   \begin{matrix}
     a+b   &a     &      &0   \\
     b     &\ddots&\ddots&    \\
           &\ddots&\ddots&a   \\
     0     &      &b     &a+b
   \end{matrix}
 \right\vert
$$

Montrer que $\forall n\in\N,n\geq4,\;B_n=(a+b)B_{n-1}-abB_{n-2}$
Montrer que 
$$
 \forall n\in\N,n\geq2,\;B_n=\frac{a^{n+1}-b^{n+1}}{a-b}.
$$
\finenonce{001132}



\finexercice



\exercice{1133, barraud, 2003/09/01}

\enonce{001133}{}

On s'intéresse aux suites réelles $(u_{n})_{n\in\N}$ satisfaisant la
relation de récurrence 
$$
    \forall n\in\N\qquad u_{n+2}=\sqrt{2}u_{n+1}-u_{n}
    \rlap{\hspace{1cm}($\star$)}
$$
\begin{enumerate}
\item 
Déterminer toutes les suites complexes satisfaisant la relation $(\star)$.

\item 
Déterminer toutes les suites réelles satisfaisant la relation $(\star)$.


On considère maintenant le déterminant d'ordre $n$ suivant :
$$
 \Delta_{n}=\left|
   \begin{matrix}
     \sqrt{2} &    1   &        &
     \hspace{-3ex}\raisebox{-2ex}[0pt][0pt]{\text{\LARGE{$0$}}} \\
           1  & \ddots & \ddots &           \\
              & \ddots & \ddots &        1  \\
     \hspace{3ex}\raisebox{1ex}[0pt][0pt]{\text{\LARGE{$0$}}} 
              &        &    1   &  \sqrt{2} \\

   \end{matrix}\,
   \right|
$$
\item
Calculer $\Delta_{n+2}$ en fonction de $\Delta_{n+1}$ et $\Delta_{n}$
pour $n\in\N$ (on pose $\Delta_{0}=1$).

En déduire la valeur de $\Delta_{n}$ en fonction de $n$. 
\end{enumerate}
\finenonce{001133}



\finexercice



\exercice{1134, barraud, 2003/09/01}

\enonce{001134}{}
  Calculer les déterminants suivants~:
$$
\left|
  \begin{matrix}
    2 & 3 \\
    -1& 4
  \end{matrix}
\right|
\quad
\left|
  \begin{matrix}
    1 & 0 & 2 \\
    3 & 4 & 5  \\
    5 & 6  & 7
  \end{matrix}
\right|
\quad
\left|
  \begin{matrix}
    1 & 0 & 6 \\
    3 & 4 & 15\\
    5 & 6 & 21
  \end{matrix}
\right|
\quad
\left|
  \begin{matrix}
    1 & 0 & 0 \\
    2 & 3 & 5  \\
    4 & 1 & 3
  \end{matrix}
\right|
$$
\finenonce{001134}



\finexercice



\exercice{1135, barraud, 2003/09/01}

\enonce{001135}{}
  Calculer les déterminants suivants~:
$$
\left|
  \begin{matrix}
    -4 &1 &1 &1 &1 \\
     1 &-4&1 &1 &1 \\
     1 & 1&-4&1 &1 \\
     1 &1 &1 &-4&1 \\
     1 &1 &1 &1 &-4
   \end{matrix}
\right|
\quad
\left|
  \begin{matrix}
    1 & a & b+c \\
    1 & b & c+a \\
    1 & c & a+b
  \end{matrix}
\right|
\quad
\left|
  \begin{matrix}
    1 & a_{1}^{} & a_{1}^{2} & a_{1}^{3} \\
    1 & a_{2}^{} & a_{2}^{2} & a_{2}^{3} \\
    1 & a_{3}^{} & a_{3}^{2} & a_{3}^{3} \\
    1 & a_{4}^{} & a_{4}^{2} & a_{4}^{3} 
  \end{matrix}
\right|
\quad
\left|
  \begin{matrix}
    1 & a_{1}^{} & a_{1}^{2} &\hdots& a_{1}^{n-1} \\
    1 & a_{2}^{} & a_{2}^{2} &\hdots& a_{2}^{n-1} \\
    \vdots&\vdots&\vdots&  &\vdots\\
    1 & a_{n}^{} & a_{n}^{2} &\hdots& a_{n}^{n-1} 
  \end{matrix}
\right|
$$
\finenonce{001135}



\finexercice



\exercice{1136, barraud, 2003/09/01}

\enonce{001136}{}
  Les nombres 119, 153 et 289 sont tous divisibles par 17. Montrer, sans
  le développer que le déterminant $\left|
   \begin{matrix}
     1&1&9\\1&5&3\\2&8&9
   \end{matrix}
  \right| $ est divisible par 17.
\finenonce{001136}



\finexercice



\exercice{1137, barraud, 2003/09/01}

\enonce{001137}{}
Calculer les d\'eterminants suivants :

$$ 
\Delta_1 = 
\left\vert
\begin{matrix}
  a&c&c&b \\
  c&a&b&c \\
  c&b&a&c \\
  b&c&c&a
\end{matrix}-
\right\vert
\qquad
\Delta_2 = 
\left\vert
\begin{matrix}
  c&a&b&c \\
  a&c&c&b \\
  b&c&c&a \\
  c&b&a&c
\end{matrix}
\right\vert
\qquad
\Delta_3 = 
\left\vert
\begin{matrix}
  a & 0 & b & 0 \\
  0 & a & 0 & b \\
  c & 0 & d & 0 \\
  0 & c & 0 & d 
\end{matrix}
\right\vert
$$
\finenonce{001137}



\finexercice



\exercice{1138, barraud, 2003/09/01}

\enonce{001138}{}
Pour $(a_{0},\dots,a_{n-1)}\in\R^{n}$, on note $A_{(a_{0}\dots a_{n})}$
la matrice
$$ 
A_{(a_{0}\dots a_{n-1})}= 
\left\vert       
\begin{matrix}
  0 &  0   & \cdots   & 0        & a_{0}        \\
  1 &  0   & \ddots   & \vdots   & \vdots       \\  
  0 &  1   & \ddots   & 0        & \vdots       \\  
  \vdots   & \ddots & \ddots   & 0        & a_{n-2}      \\    
  0        & \cdots &  0       & 1        & a_{n-1}-\lambda \\
 \end{matrix}
\right\vert
$$
et à $\lambda\in\R$, on associe $\Delta_{(a_{0}\dots
  a_{n-1})}(\lambda)=\det(A_{(a_{0},...,a_{n-1})}-\lambda\mathrm{id})$. Calculer
$\Delta_{(a_{0}\dots a_{n-1})}(\lambda)$ en fonction de
$\Delta_{(a_{1}\dots a_{n-1})}(\lambda)$ et $a_{0}$. En déduire
$\Delta_{(a_{0}\dots a_{n-1})}(\lambda)$.
\finenonce{001138}



\finexercice



\exercice{1139, barraud, 2003/09/01}

\enonce{001139}{}
  Calculer les déterminants suivant~:
$$
\left\vert
\begin{matrix}
  a_{11} &a_{12} &\cdots  &a_{1n}\\
  0      &a_{22} &\ddots  &\vdots\\
  \vdots &\ddots &\ddots  &a_{n-1,n}      \\
  0      &\cdots &0       &a_{nn}
\end{matrix}
\right\vert
\qquad
\left\vert
\begin{matrix}
  p &q      &       &  \hspace{-1ex}\raisebox{-2ex}[0ex][0ex]{\Large 0} \\
  1 &p      &\ddots &  \\
    &\ddots &\ddots &q \\
\hspace{1ex}\raisebox{1ex}[0ex][0ex]{\Large 0}
    &       &1      &p
\end{matrix}
\right\vert
\qquad
\left\vert
\begin{matrix}
  1 &0      &       &1 \\
  1 &1      &\ddots &  \\
    &\ddots &\ddots &0 \\
  0 &       &1      &1
\end{matrix}
\right\vert
\qquad
\left\vert
\begin{matrix}
    a+b  &    a   & \cdots &  a       \\
     a   &   a+b  & \ddots & \vdots   \\
  \vdots & \ddots & \ddots &  a       \\
     a   & \cdots &    a   & a+b 
\end{matrix}
\right\vert
$$
\finenonce{001139}



\finexercice



\exercice{1140, barraud, 2003/09/01}

\enonce{001140}{}
  Soit $B\in\mathcal{M}_{n,m}(\R)$ et $C\in\mathcal{M}_{m,m}(\R)$. On
  considère l'application $\phi$ suivante~:
$$
\phi~:
\begin{array}{ccc}
  \mathcal{M}_{n,n}(\R) & \rightarrow  & \R \\
  A                     & \mapsto &
  \det
  \begin{pmatrix}
    A & B  \\
    0 & C 
  \end{pmatrix}
\end{array}
$$
Etudier la multi-linéarité de $\phi$ par rapport aux colonnes de $A$.
Calculer $\phi(\mathrm{id})$. En déduire que 
$$
\det
  \begin{pmatrix}
    A & B  \\
    0 & C 
  \end{pmatrix}=\det(A)\det(C)
$$
Soit $M=
\begin{pmatrix}
  A_{1}&\cdots &   \\
       &\ddots &\vdots\\
  0    &       & A_{k} \\
\end{pmatrix}
$ une matrice triangulaire par blocs. Montrer que
$\det(M)=\det(A_{1})\cdots\det(A_{k})$
\finenonce{001140}



\finexercice



\exercice{1141, barraud, 2003/09/01}

\enonce{001141}{}
  Calculer le déterminant suivant~:
  $$\Delta=\left|
    \begin{matrix}
     0     & a_{12}& a_{13}& a_{14}& a_{15} \\
    -a_{21}& 0     & a_{23}& a_{24}& a_{25} \\
    -a_{31}&-a_{32}& 0     & a_{34}& a_{35} \\
    -a_{41}&-a_{42}&-a_{43}& 0     & a_{45} \\
    -a_{51}&-a_{52}&-a_{53}&-a_{54}& 0      \\
  \end{matrix}
\right|$$
Comment généraliser ce résultat en dimension plus grande~?
\finenonce{001141}



\finexercice



\exercice{1142, barraud, 2003/09/01}

\enonce{001142}{}
Calculer les déterminants suivants~:
$$
\left|
  \begin{matrix}
  1      &1       &1       &1       \\
  \cos x & \cos y & \cos z & \cos t \\
  \cos2x & \cos2y & \cos2z & \cos2t \\
  \cos3x & \cos3y & \cos3z & \cos3t 
\end{matrix}
\right|
\qquad
\left\vert
\begin{matrix}
     1   &    2   &   3    & \cdots &  n      \\
    -1   &    0   &   3    &        &  n      \\
    -1   &   -2   &   0    &        &  n      \\
  \vdots & \vdots & \vdots &        & \vdots  \\
    -1   &   -2   &   -3   & \cdots &  0 
\end{matrix}
\right\vert
\qquad
\left\vert
\begin{matrix}
  0      & 1      & 2      & \cdots & n-1       \\
  1      & 0      & 1      & \ddots & \cdots    \\
  2      & 1      & \ddots & \ddots & 2         \\ 
  \vdots & \ddots & \ddots & 0      & 1         \\
  n-1    & \cdots & 2      & 1      & 0
\end{matrix}
\right\vert
$$
\finenonce{001142}



\finexercice



\exercice{1143, barraud, 2003/09/01}
\video{64911kcKNmg}
\enonce{001143}{}
  Soit $(a_{0},...,a_{n-1})\in\C^{n}$, $x\in\C$. Calculer
 $$
 \Delta_{n}=
 \left|
   \begin{matrix}
   x &  0    &        & a_{0}   \\
    -1 &\ddots &\ddots  &\vdots  \\
      &\ddots &x      & a_{n-2} \\
    0 &       & -1      & x+a_{n-1}
   \end{matrix}
\right|
 $$
\finenonce{001143}


\finexercice

\exercice{1144, barraud, 2003/09/01}

\enonce{001144}{}
Soit $(a_1,a_2,a_3)\in(\mathbb{K})^{3}$. On note $j=e^{\frac{2i\pi}{3}}$, et on
considère les deux matrices suivantes :
$$
 A=
 \begin{pmatrix}
   a_1 & a_2 & a_3 \\
   a_3 & a_1 & a_2 \\
   a_2 & a_3 & a_1 
 \end{pmatrix}
\qquad\text{ et }
 V=
 \begin{pmatrix}
   1 & 1   & 1   \\
   1 & j   & j^2 \\
   1 & j^2 & j 
 \end{pmatrix}
$$
Calculer le produit $AV$, puis $\det(AV)$ en fonction de $\det(V)$, et en déduire $\det(A)$.
\finenonce{001144}



\finexercice




\exercice{1145, barraud, 2003/09/01}
\video{YveAcpVpK2g}
\enonce{001145}{}
Soit $a$ un réel.
On note $\Delta_n$ le déterminant suivant : 
$$
\Delta_n = 
\left\vert
\begin{matrix}
       a   &    0   & \cdots & 0      & n-1 \\
       0   &    a   & \ddots & \vdots & \vdots \\
    \vdots & \ddots & \ddots & 0      & 2 \\
       0   & \cdots &   0    & a      & 1 \\
      n-1  & \cdots &   2    & 1      & a
\end{matrix}
\right\vert
$$
\begin{enumerate}
  \item Calculer $\Delta_n$ en fonction de $\Delta_{n-1}$.
  \item Démontrer que : $\displaystyle \forall n\geq2\quad
\Delta_n=a^n-a^{n-2}\sum_{i=1}^{n-1}{i^2}$.
\end{enumerate}
 
\finenonce{001145}


\finexercice


\exercice{1146, barraud, 2003/09/01}

\enonce{001146}{}
Soit $a$ un réel différent de 1. Pour $n\in\N$, $n\geq2$, on note 
$$ 
D_n = 
\left\vert       
\begin{matrix}
    1+a^2 &   a    &    0   & \cdots & 0      \\
    a     &  1+a^2 &    a   & \ddots & \vdots \\
    0     &   a    & \ddots & \ddots & 0      \\
   \vdots & \ddots & \ddots & 1+a^2  & a      \\
    0     & \cdots &    0   &  a     & 1+a^2
\end{matrix}
\right\vert
$$
Calculer $D_n$ en foncion de $D_{n-1}$ et $D_{n-2}$. Monter que 
$
  D_n=\frac{1-a^{2\rlap{$\scriptscriptstyle n+2$}}}{1-a^{2}}
  \hphantom{\scriptscriptstyle n+2}.
$
Combien vaut $D_n$ si $a=1$ ?
\finenonce{001146}



\finexercice



\exercice{1147, barraud, 2003/09/01}

\enonce{001147}{}
  Soient $a,b,c$ trois réels et $\Delta_{n}$ le déterminant de taille
  $n$ suivant~:
$$
\Delta_{n}=\left\vert
  \begin{array}{c@{\ }c@{}c@{\ }c}
   a     & b      &        & 0
\\[-1ex]
   c     & \ddots & \ddots &                  \\[-2ex]
         & \ddots & \ddots & b \\[-1ex]
   0 & & c      & a \\
\end{array}
\right\vert
$$
\begin{enumerate}
\item 
On pose $\Delta_{0}=1$, $\Delta_{1}=a$. Montrer que $\forall n\in\N,
\Delta_{n+2}=a\Delta_{n+1}-bc\Delta_{n}$.

\item 
On suppose que $a^{2}=4bc$. Montrer par récurrence que~: 
$$\forall
n\in\N, \Delta_{n}=(n+1)\,\frac{a^{n}}{2^{n}}$$
\end{enumerate}
\finenonce{001147}



\finexercice



\exercice{1148, barraud, 2003/09/01}

\enonce{001148}{}
  Calculer le déterminant suivant~:
$$
\Delta=\left\vert
  \begin{matrix}
    1 &2 &4 &8 \\
    1 &3 &9 &27\\
    1 &4 &16&64\\
    1 &5 &25&125
  \end{matrix}
\right\vert
$$
\finenonce{001148}



\finexercice



\exercice{1149, barraud, 2003/09/01}

\enonce{001149}{}
  Soit $\Delta_{n}$ le déterminant de taille $n$ suivant~:
$$
\Delta_{n}=\left\vert
  \begin{matrix}
   3     & 1    & 0      &\cdots & 0 \\
   2     & 3    & 1      &\ddots &\vdots \\
   0     & 2    & 3      &\ddots & 0 \\
  \vdots &\ddots&\ddots  &\ddots & 1 \\
   0     &\cdots& 0      & 2     & 3
\end{matrix}
\right\vert
$$
\begin{enumerate}
\item Montrer que $\forall n\in\N^{*}, \Delta_{n+2}=3\Delta_{n+1}-2\Delta_{n}$
(avec la convention $\Delta_{0}=1$, $\Delta_{1}=3$).

\item Montrer par récurrence que $\forall n\in\N^{*}, \Delta_{n}=2^{n+1}-1$
\end{enumerate}
\finenonce{001149}



\finexercice



\exercice{1150, barraud, 2003/09/01}

\enonce{001150}{}
  Soit $u$ l'application de $\R_{n}[X]$ dans $\R_{n}[X]$ définie par
  $u(P)=P+P'$. Calculer $\det u$. Même question lorsque $u(P)=XP'+P(1)$.
\finenonce{001150}



\finexercice



\exercice{2448, matexo1, 2002/02/01}

\enonce{002448}{}
Calculer les d\'eterminants suivants
$$\left|
\begin{array}{cc}
4 & 5 \\
-7 & 8
\end{array}
\right|,
\qquad
\left|
\begin{array}{ccc}
1 & 3 & 2 \\
-1 & -2 & 3 \\
4 & 0 & 1
\end{array}
\right|,
\qquad
\left|
\begin{array}{cccc}
1 & 0 & 3 & 2 \\
-1 & 1 & 8 &-5 \\ 
2 & 0 & 6 & 3 \\
-1 & 0 & -3 & -7
\end{array}
\right|.$$
\finenonce{002448}


\finexercice


\exercice{2452, matexo1, 2002/02/01}

\enonce{002452}{}
Calculer par r\'ecurrence le d\'eterminant
$$ \Delta_n = \left|
\begin{array}{ccccc}
a_1+b_1 & b_1     & b_1 & \ldots & b_1 \\
   b_2  & a_2+b_2 & b_2 & \ldots & b_2 \\
  \ldots&\ldots   &\ldots&\ldots&\ldots\\
   b_n  &  b_n    &\ldots&\ldots& a_n+b_n
\end{array}\right|.$$
 
\finenonce{002452}


\finexercice



\exercice{2453, matexo1, 2002/02/01}
\video{NwFtb_SYHA8}
\enonce{002453}{Déterminant de Vandermonde}
Montrer que
$$\left|
\begin{array}{ccccc}
1 & t_1 & t_1^2 & \ldots & t_1^{n-1} \\
1 & t_2 & t_2^2 & \ldots & t_2^{n-1} \\\
\ldots&\ldots&\ldots& \ldots & \ldots \\
1 & t_n & t_n^2 & \ldots & t_n^{n-1}
\end{array}\right|
 = \prod_{1 \le i < j \le n} (t_j - t_i) $$
\finenonce{002453}


\finexercice



\exercice{2454, matexo1, 2002/02/01}

\enonce{002454}{}
Soit $M$ une matrice carr\'ee d'ordre $n$, et $M'$ la matrice
d\'eduite de $M$ en rempla\c{c}ant, pour tout $j$, la $j$-i\`eme colonne
par la somme des colonnes de $M$ 
d'indices diff\'erents de $j$. Montrer que
$\det M' = (-1)^{n-1} (n-1) \det M $.
\finenonce{002454}


\finexercice


\exercice{2455, matexo1, 2002/02/01}

\enonce{002455}{}
Calculer le d\'eterminant d'ordre $n$
$$\left|
\begin{array}{ccccccc}
-1    &   1    &   1   &  0  &   0  & \ldots & 0 \\
0    &    -1   &   1   &  1  &  0   & \ldots  & 0 \\
\ldots &\ldots & \ldots&\ldots&\ldots&\ldots & \ldots \\
0  &      0    &  0     &  0 &  \ldots& -1   &1
\end{array}\right|.$$
\finenonce{002455}


\finexercice


\exercice{2456, matexo1, 2002/02/01}

\enonce{002456}{}
Calculer le d\'eterminant 
$$ D_n = \left|
\begin{array}{ccccc}
3 & 1 & 1 & \ldots & 1 \\
1 & 4 & 1 & \ldots & 1 \\
\ldots&\ldots&\ldots& \ldots & \ldots \\
1 & 1 & 1 & \ldots & n+2
\end{array}\right|.$$
\finenonce{002456}


\finexercice


\exercice{2458, matexo1, 2002/02/01}

\enonce{002458}{}
Soit $E$ l'espace vectoriel des matrices carr\'ees d'ordre $n$
r\'eelles, et $A \in E$ fix\'ee. On d\'efinit une application
$u_A$ de $E$ sur lui-m\^eme par $u_A(B) = AB$. Montrer que c'est
un endomorphisme de $E$ et que $\det u_A = (\det A)^n $.
\finenonce{002458}


\finexercice


\exercice{2466, matexo1, 2002/02/01}

\enonce{002466}{}
Soit la matrice
$$A=\left( \begin{array}{ccc}
m-2 & 2 & -1 \\
2 & m & 2 \\
2m & 2m+2 & m+1
\end{array}\right).$$
\begin{enumerate}
\item Calculer $\det A$.
\item Soit $u$ l'endomorphisme de $\R^3$ dont la matrice par
rapport \`a la base canonique est $A$. Pour quelles valeurs de $m$
est-ce un isomorphisme de $\R^3$\,?
\item On pose $m = 1$. Trouver une base du noyau de $u$.
\end{enumerate}
\finenonce{002466}


\finexercice


\exercice{2578, delaunay, 2009/05/19}
\enonce{002578}{}
Soient $a,b,c$ des r\'eels v\'erifiant $a^2+b^2+c^2=1$ et
$P$ la matrice r\'eelle $3\times3$ suivante :
$$P=\begin{pmatrix}a^2&ab&ac \\  ab&b^2&bc \\  ac&bc&c^2\end{pmatrix}$$
\begin{enumerate}  
  \item Calculer le d\'eterminant de $P$.
  \item D\'eterminer les sous-espaces vectoriels de $\R^3$, $\ker P$ et $\Im P$.
  \item Soit $Q=I-P$, calculer $P^2$, $PQ$, $QP$ et $Q^2$.
  \item Caract\'eriser g\'eom\'etriquement $P$ et $Q$.
 \end{enumerate}
\finenonce{002578} 


\finexercice

\exercice{2582, delaunay, 2009/05/19}
\enonce{002582}{}
Soit $a\in\R$ et $A$ la matrice suivante
$$A=\begin{pmatrix}1&0&a \\  0&a&1 \\  a&1&0\end{pmatrix}$$

\begin{enumerate}
 \item Calculer le d\'eterminant de $A$ et d\'eterminer 
pour quelles valeurs de $a$ la matrice est inversible.
 \item Calculer $A^{-1}$ lorsque $A$ est inversible.
\end{enumerate}
\finenonce{002582} 


\finexercice


\exercice{2753, tumpach, 2009/10/25}
\video{d4ZZDotSUcw}
\enonce{002753}{}
\begin{enumerate}
\item Calculer l'aire du parall\'elogramme construit sur les vecteurs 
$\vec{u} = \left(\begin{array}{c}2\\3\end{array}\right)$ et 
$\vec{v} = \left(\begin{array}{c}1\\4\end{array}\right)$.
\item Calculer le volume du  parall\'el\'epip\`ede construit sur les vecteurs \\
$\vec{u} = \left(\begin{array}{c}1\\2\\0\end{array}\right)$, 
$\vec{v} = \left(\begin{array}{c}0\\1\\3\end{array}\right)$ et 
$\vec{w} = \left(\begin{array}{c}1\\1\\1\end{array}\right)$.
\item Montrer que le volume d'un parall\'el\'epip\`ede dont les sommets sont des points de 
$\Rr^3$ \`a coefficients entiers est un nombre entier.
\end{enumerate}
\finenonce{002753}


\finexercice


\exercice{2754, tumpach, 2009/10/25}
\enonce{002754}{}
Calculer les d\'eterminants des matrices suivantes~:
$$
\left(\begin{array}{ccc}  1 & 2 & 3\\  2 & 3 & 0\\  3 & 0 & 1\end{array}\right),
\left(\begin{array}{cccc}  1 & 1 & 0\\  0 & 0 & 1\\  1 & 0 & 1  \end{array}\right),
\left( \begin{array}{ccc}2 & 7 & 1\\-1 & 2 & 0\\ 3 & 5 &  1 \end{array}\right),$$
$$
\left(\begin{array}{cccc} 2 & 1 & 2\\  3 & 1 & 3\\  1 & 0& 6\end{array}\right),
\left(\begin{array}{ccc}1 & 2 & 3\\ 2 & 3 & 1\\ 3 & 1 & 2\end{array}\right),
\left(\begin{array}{ccc} 1 & 0 & -1\\ 2 & 0  & 1\\1 & 1 & 3\end{array}\right).
$$
\finenonce{002754}



\finexercice


\exercice{2755, tumpach, 2009/10/25}
\enonce{002755}{}
Calculer les d\'eterminants des matrices suivantes~:
$$
\left(\begin{array}{cccc} 0 & 1 & 2 & 3\\ 1 & 2 & 3 & 0\\ 2 & 3 & 0 & 1\\3 & 0 & 1 & 2\end{array}\right),
\left(\begin{array}{cccc} 0 & 1 & 1 & 0\\ 1 & 0 & 0 & 1\\ 1 & 1 & 0 & 1\\  1 & 1 & 1 &0\end{array}\right),
\left(\begin{array}{cccc}1 & 2 & 1 & 2\\ 1 & 3 & 1 & 3\\ 2 & 1 & 0& 6\\ 1 & 1& 1&7\end{array}\right).
$$
\finenonce{002755}



\finexercice


\exercice{2756, tumpach, 2009/10/25}
\enonce{002756}{}
Calculer les d\'eterminants des matrices suivantes~:
$$
\left(\begin{array}{cccc} 0 & 1 & 2 & 3\\ 1 & 3 & 6 & 10\\ 1 & 4 & 10 & 20\\1 & 5 & 15 & 35\end{array}\right),
\left(\begin{array}{cccc} a & b & c & d\\ -a & b & \alpha & \beta\\ -a & -b& c & \gamma\\  -a & -b & -c &d\end{array}\right),
\left(\begin{array}{cccc}-1 & 1 & 1 & 1\\ 1 & -1 & 1 & 1\\ 1 & 1 & -1& 1\\ 1 & 1& 1&-1\end{array}\right).
$$
\finenonce{002756}



\finexercice


\exercice{2757, tumpach, 2009/10/25}
\enonce{002757}{}
Calculer les d\'eterminants des matrices suivantes~:
$$
\left(\begin{array}{cccc} a & b &c & d\\ e & f& g & 0\\h& k & 0 & 0\\l & 0 & 0 & 0\end{array}\right),
\left(\begin{array}{cccc} a & 0 & 0 & 0\\b & c & 0 & 0\\d & e & f & 0\\g & h & k & l\end{array}\right),
\left(\begin{array}{cccc}a & b & c & d\\ 0 & e & f & g\\ 0 &0 &h & k\\ 0&0&0&l\end{array}\right).
$$
\finenonce{002757}



\finexercice


\exercice{2758, tumpach, 2009/10/25}
\enonce{002758}{}
Soit $M = (m_{ij})$ une matrice carr\'ee de taille $n$. On construit \`a partir de $M$ la matrice $N = (n_{ij})$ de la mani\`ere suivante~: pour tout couple d'indices $i, j$, on appelle $M_{ij}$ la matrice obtenue \`a partir de $M$ en rayant la ligne $i$ et la colonne $j$ ; alors $n_{ij} = (-1)^{i+j}\det(M_{ji})$. D\'emontrer que $MN = NM = \det(M)I$, o\`u $I$ d\'esigne la matrice identit\'e. En d\'eduire une m\'ethode d'inversion de matrices passant par le calcul de d\'eterminants, et l'appliquer \`a la matrice
$$
M = \left(\begin{array}{cccc}3 & -2 & 0 & -1\\0 & 2 & 2 & 1\\1& -2 & -3 & -2\\ 0 & 1 & 2 & 1\end{array}\right).
$$
\finenonce{002758}



\finexercice


\exercice{2759, tumpach, 2009/10/25}
\enonce{002759}{}
Calculer les inverses des matrices suivantes de deux mani\`eres diff\'erentes~:
$$
A = \left(\begin{array}{ccc}
1 & 2 & 3 \\2 & 5 & 4\\-1 & 0 & 2
\end{array}\right); \quad\quad
B =  \left(\begin{array}{ccc}
1 & 1 & -1\\1 & 2 &3 \\1 & 3 & 6
\end{array}\right);$$ 
$$
C = \left(\begin{array}{cccc}
0 & 1 & 0 & 0 \\
3 & 0 & 2 & 0\\
0 & 2 & 0 & 3\\
0 & 0 & 1 & 0
\end{array}\right);\quad\quad
D = \left(\begin{array}{cccc}
1 & 1 & 1 & 1\\
0 &-1 & -2 &-3\\ 
0 & 0 &1 & 3\\ 
0 & 0 & 0&-1 
\end{array}\right).
$$
\finenonce{002759}



\finexercice


\exercice{2760, tumpach, 2009/10/25}
\enonce{002760}{}
En utilisant le d\'eterminant montrer que chacun des syst\`emes suivants admet une solution unique.  R\'esoudre chacun de ces syst\`eme en inversant la matrice de ses  coefficients~:  
$$
\left\{\begin{array}{ccccccc}
x & + & y & + & z & = &1\\
2x & +  & 3y & + & 4z &  =  & 2\\
&& y &+ &  4z & = & 3 
\end{array}\right.
$$
$$
\left\{\begin{array}{cccccccc}
x & + & y & + &  z & + &  t  & = 1\\
x & + &y & - & z & - & t & =  0\\
x & - & y & - & z&  +  &t & = 2\\
x & - & y & + & z & - & t & = 3  
\end{array}\right.
$$
\vspace{.2cm}
$$
\left\{
\begin{array}{cccccc}
3x & + & 2y & - &  z &  = -1\\
x&+&y&-&3z& = 1\\
3x& + &2y &&&= 0\\
\end{array}
\right.
$$
\finenonce{002760}



\finexercice


\exercice{3451, quercia, 2010/03/10}
 \enonce{003451}{Calcul de déterminants}
 
 Calculer les déterminants suivants :
 
  \begin{enumerate}
   \item $\begin{vmatrix}  x &a &b &x \cr                                      %+---------+
                 b &x &x &a \cr                                      %|  x,a,b  |
                 x &b &a &x \end{vmatrix}$.                                      %+---------+
 
 
   \item $\begin{vmatrix} a-b-c &2a    &2a    \cr                      %+-----------------+
                2b    &b-c-a &2b    \cr                      %|  a-b-c, 2a, 2a  |
                2c    &2c    &c-a-b \cr \end{vmatrix}$.                  %+-----------------+
     
 
   \item $\begin{vmatrix} a+b     &b+c     &c+a     \cr                  %+---------------+
                a^2+b^2 &b^2+c^2 &c^2+a^2 \cr                  %|  a+b b+c c+a  |
                a^3+b^3 &b^3+c^3 &c^3+a^3 \cr \end{vmatrix}$.              %+---------------+
   
 
   \item $\begin{vmatrix} a+b &ab &a^2+b^2 \cr                         %+-----------------+
                b+c &bc &b^2+c^2 \cr                         %|  a+b ab a¿+b¿   |
                c+a &ca &c^2+a^2 \cr \end{vmatrix}
                $.                     %+-----------------+
 
 
   \item $\begin{vmatrix} a^2 &b^2 &ab  \cr                                %+-------------+
                b^2 &ab  &a^2 \cr                                %|  a¿ b¿ ab   |
                ab  &a^2 &b^2 \cr \end{vmatrix}$.                            %+-------------+
 
   \item $\begin{vmatrix}   a & b      &        &(0)  \cr     %+-------------------------+
                   c & \ddots & \ddots       \cr     %|  Matrice tridiagonale   |
                     & \ddots & \ddots & b   \cr     %+-------------------------+
                  (0)&        & c      & a   \cr \end{vmatrix}$.
     \label{tridiag}
 
   \item $\begin{vmatrix}a      &b       &\dots   &b       \cr               %+-----------+
               b      &\ddots  &(0)     &\vdots  \cr               %|  anneau   |
               \vdots &(0)     &\ddots  &b       \cr               %+-----------+
               b      &\dots   &b       &a       \cr\end{vmatrix}$.
     
                                                   %+---------------------------+
                                                   %|  Déterminants de Pascal   |
                                                   %+---------------------------+
   \item $\begin{vmatrix} C_n^0\hfill  & C_n^1\hfill  & \dots & C_n^p\hfill  \cr
                C_{n+1}^0    & C_{n+1}^1    & \dots & C_{n+1}^p    \cr
                \vdots       & \vdots       &       & \vdots       \cr
                C_{n+p}^0    & C_{n+p}^1    & \dots & C_{n+p}^p    \cr \end{vmatrix}$.
 
 
                                                                %+--------------+
                                                                %|  ai+bi, bi   |
                                                                %+--------------+
   \item $\begin{vmatrix}
               a_1 + b_1 & b_1       & \dots  & \dots  & b_1       \cr
               b_2       & a_2 + b_2 & b_2    & \dots  & b_2       \cr
               \vdots    &           & \ddots &        & \vdots    \cr
               \vdots    &           &        & \ddots & b_{n-1}   \cr
               b_n       & \dots     & \dots  & b_n    & a_n + b_n \cr \end{vmatrix}$.
     
      
   \item $\begin{vmatrix} a_1-b_1 &\dots &a_1-b_n \cr                        %+-----------+
                \vdots  &\dots &\vdots  \cr                        %|   ai-bj   |
                a_n-b_1 &\dots &a_n-b_n \cr \end{vmatrix}$, $(n \ge 3)$.       %+-----------+
     
 
   \item $\begin{vmatrix} 1     & 2     & \dots & n     \cr   %+--------------------------+
                2     & 3     & \dots & 1     \cr   %|  Déterminant circulant   |
                \vdots& \vdots&       & \vdots\cr   %+--------------------------+
                n     & 1     & \dots & n-1   \cr\end{vmatrix}$.
 
   \item $\begin{vmatrix} 0      &1     &2      &\dots  &n-1     \cr         %+-----------+
                1      &0     &1      &       &\vdots  \cr         %|   |i-j|   |
                2      &1     &0      &\ddots &2       \cr         %+-----------+
                \vdots &      &\ddots &\ddots &1       \cr
                n-1    &\dots &2      &1      &0       \cr\end{vmatrix}$.
   
 
 
 Pour \ref{tridiag} : Chercher une relation de récurrence linéaire d'ordre 2.
 On notera $\alpha$ et $\beta$ les racines dans $\C$ de l'équation
 caractéristique, et on exprimera le déterminant en fonction de $\alpha$ et
 $\beta$.
 \end{enumerate}
 \finenonce{003451}
 
 
 
 \finexercice


\exercice{3452, quercia, 2010/03/10}
 \enonce{003452}{Factorisation de polynômes}
 
 Déterminer les cas d'annulation des déterminants suivants, puis les
 calculer :

 
 \begin{enumerate}
   \item $\begin{vmatrix}1      &        &        & (1)    \cr
                      & 1-x    &        &        \cr
                      &        & \ddots &        \cr
               (1)    &        &        & n-x    \cr \end{vmatrix} $.
 
   \item $\begin{vmatrix}x     & a_1  & a_2   & \dots & a_n   \cr
               a_1   & x    & a_2   & \dots & a_n   \cr
               \vdots&      & \ddots&       & \vdots\cr
               \vdots&      &       & \ddots& a_n   \cr
               a_1   & a_2  & \dots & a_n   & x     \cr \end{vmatrix}$.
 
   \item $\begin{vmatrix}a     &b  &c  &\dots &z     \cr
               b     &b  &c  &      &z     \cr
               c     &c  &c  &      &z     \cr
               \vdots&   &   &\ddots&\vdots\cr
               z     &z  &z  &\dots &z     \cr \end{vmatrix}$.
 
   \item $\begin{vmatrix}\frac 1{a+x} & \frac 1{a+y} & \frac 1{a+z} \cr
                \frac 1{b+x} & \frac 1{b+y} & \frac 1{b+z} \cr
                \frac 1{c+x} & \frac 1{c+y} & \frac 1{c+z} \cr\end{vmatrix}$.
     
 \end{enumerate}
 \finenonce{003452}
 
 
 
 \finexercice


\exercice{3453, quercia, 2010/03/10}
\enonce{003453}{Calcul par dérivation}

\begin{enumerate}
  \item Soient ${a,b,c,d} : \R \to \R$ des fonctions dérivables et
    $f(x) = \begin{vmatrix} a(x) &b(x) \cr c(x) & d(x) \end{vmatrix}$.

    Montrer que $f$ est dérivable et que :
    $f'(x) = \begin{vmatrix} a'(x) &b(x)  \cr c'(x) &d(x)  \end{vmatrix}
           + \begin{vmatrix} a(x)  &b'(x) \cr c(x)  &d'(x) \end{vmatrix}$.

  \item Généraliser à un déterminant $n\times n$.
  \item Application :
    Calculer $\begin{vmatrix} 1 & \cos x          & \sin x          \cr
                        1 & \cos (x+\alpha) & \sin (x+\alpha) \cr
                        1 & \cos (x+\beta)  & \sin (x+\beta)  \cr\end{vmatrix}$.
    
\end{enumerate}
\finenonce{003453}



\finexercice


\exercice{3454, quercia, 2010/03/10}
\enonce{003454}{$\det( A + (\alpha) )$}

Soit $A \in \mathcal{M}_n(K)$ et
$U = \begin{pmatrix} 1       &\dots  &1       \cr
                \vdots  &       &\vdots  \cr
                1       &\dots  &1       \cr \end{pmatrix}$.


\begin{enumerate}
  \item Démontrer que :
    $\det(A + \alpha U) = \det A + \alpha \sum \text{cofacteurs de } A$.
    
  \item  En déduire la valeur de
     $D(a,b,c) = \begin{vmatrix}a       & b      & \dots  & b      \cr
                           c       & \ddots & \ddots & \vdots \cr
                           \vdots  & \ddots & \ddots & b      \cr
                           c       & \dots  & c      & a      \end{vmatrix}$,
  \begin{enumerate}
     \item pour $b \ne c$.
     \item pour $b = c$.
  \end{enumerate}
\end{enumerate}
\finenonce{003454}



\finexercice


\exercice{3455, quercia, 2010/03/10}
\enonce{003455}{Déterminant tridiagonal, Matexo}

Soit $n\in\N^*$, $(a_1,a_2,\ldots,a_n)\in{(\R^*_+)}^n$,
$(b_1,b_2,\ldots,b_{n-1})\in{(\R^*_+)}^{n-1}$ et
$(c_1,c_2,\ldots,c_{n-1})\in{(\R^*_-)}^{n-1}$.
Montrer que le déterminant de la matrice suivante est strictement positif :
$$
\begin{vmatrix}
a_1 &b_1     &       &       &        &	       \cr
c_1 &a_2     &b_2    &       & 0 &	       \cr
    &c_2     &a_3    &\ddots &        &	       \cr
    &        &\ddots &\ddots &\ddots  &	       \cr
    & 0 &       &\ddots &\ddots  &b_{n-1} \cr
    &        &       &       &c_{n-1} &a_n     \cr\end{vmatrix}.
$$
\finenonce{003455}



\finexercice


\exercice{3456, quercia, 2010/03/10}
\enonce{003456}{Déterminants de Vandermonde}

Soient $a_1, \dots ,a_n \in  K$. Le déterminant de Vandermonde associé
aux $a_i$ est : $V(a_1,\dots,a_n) = \det\Bigl(a_i^{j-1}\Bigr)$.

\begin{enumerate}
  \item Calculer et factoriser $V(a,b)$ et $V(a,b,c)$.
  \item Pour $x \in  K$,
    montrer que $V(a_1,\dots,a_n,x) = V(a_1,\dots,a_n)\prod_{i=1}^n(x-a_i)$.
  \item En déduire l'expression générale de $V(a_1,\dots,a_n)$.
\end{enumerate}
\finenonce{003456}



\finexercice


\exercice{3457, quercia, 2010/03/10}
\enonce{003457}{Racines de l'unité}

On note $\omega = e^{2i\pi/n}$, $\alpha = e^{i\pi/n}$ et $D$ le déterminant
$n \times n$ : $D = \det\Bigl( \omega^{(k-1)(l-1)} \Bigr)$.

\begin{enumerate}
  \item  Calculer $D^2$.
     
  \item  Montrer que $D = \prod_{k < \ell}(\omega^\ell - \omega^k)
            = \prod_{k < \ell}\left(\alpha^{k+\ell}\cdot
              2i\sin\frac {\ell-k}n \pi\right)$.
  \item  Exprimer $D$ sous forme trigonométrique.
\end{enumerate}
\finenonce{003457}



\finexercice


\exercice{3458, quercia, 2010/03/10}
\enonce{003458}{Cosinus}

Soient $\alpha_1, \dots, \alpha_n \in \R$.
Mettre le déterminant : $\det\Bigl( \cos( (j-1)\alpha_i ) \Bigr)$ sous la forme d'un
déterminant de Vandermonde.


\finenonce{003458}



\finexercice


\exercice{3459, quercia, 2010/03/10}
\enonce{003459}{$(x_i+y_j)^k$}

Soit $k \le n-1$ et $M = \Bigl( (x_i+y_j)^k \Bigr)$.
\'Ecrire $M$ comme produit de deux matrices et calculer $\det M$.


\finenonce{003459}



\finexercice


\exercice{3460, quercia, 2010/03/10}
\enonce{003460}{$P(i+j)$}

Soit $P \in  K_{n-1}[X]$ et $A = \Bigl( P(i+j) \Bigr) \in \mathcal{M}_n(K)$.
Développer $P(i+j)$ par la formule de Taylor et écrire $A$ comme produit
de deux matrices.
En déduire $\det A$.



\finenonce{003460}



\finexercice


\exercice{5362, rouget, 2010/07/06}
\enonce{005362}{**}
Montrer que $\left|
\begin{array}{ccc}
-2a&a+b&a+c\\
b+a&-2b&b+c\\
c+a&c+b&-2c
\end{array}\right|=4(b+c)(c+a)(a+b)$.
\finenonce{005362}


\finexercice


\exercice{5363, rouget, 2010/07/06}
\enonce{005363}{**}
Pour $a$, $b$ et $c$ deux à deux distincts donnés, factoriser $\left|
\begin{array}{cccc}
X&a&b&c\\
a&X&c&b\\
b&c&X&a\\
c&b&a&X
\end{array}\right|$.
\finenonce{005363}


\finexercice


\exercice{5364, rouget, 2010/07/06}
\enonce{005364}{***}
Calculer :

\begin{enumerate}
 \item  $\mbox{det}(|i-j|)_{1\leq i,j\leq n}$  
 \item  $\mbox{det}(\sin(a_i+a_j))_{1\leq i,j\leq n}$ ($a_1$,...,$a_n$ étant $n$ réels donnés) 
 \item  $\left|
\begin{array}{cccccc}
a&0&\ldots& &\ldots&b\\
0&a&\ddots& &b&0\\
\vdots&0&\ddots& &0&\vdots\\
\vdots&0& &\ddots&0&\vdots\\
0&b& & &a&0\\
b&0&\ldots& &\ldots&a
\end{array}
\right|$

 \item  $\left|\begin{array}{ccccc}
1&1&\ldots& &1\\
1&1&0&\ldots&0\\
\vdots&0&\ddots&\ddots&\vdots\\
\vdots&\vdots&\ddots&1&0\\
1&0&\ldots&0&1
\end{array}
\right|$   
  \item  $\mbox{det}(C_{n+i-1}^{j-1})_{1\leq i,j\leq p+1}$   
  \item  $\left|\begin{array}{cccccc}
-X&1&0&\ldots&\ldots&0\\
0&-X&1&\ddots& &\vdots\\
\vdots&\ddots&\ddots&\ddots&\ddots&\vdots\\
\vdots& &\ddots&\ddots&\ddots&0\\
0&\ldots&\ldots&0&-X&1\\
a_0&\ldots& &\ldots&a_{n-2}&a_{n-1}-X
\end{array}
\right|$
\end{enumerate}
\finenonce{005364}


\finexercice


\exercice{5365, rouget, 2010/07/06}
\enonce{005365}{**** Déterminant de \textsc{Cauchy} et déterminant de \textsc{Hilbert}}
Soit $A=\left(\frac{1}{a_i+b_j}\right)_{1\leq i,j\leq n}$ où $a_1$,..., $a_n$, $b_1$,...,$b_n$ sont $2n$ réels tels que toutes les sommes $a_i+b_j$ soient non nulles. Calculer $\mbox{det}A$ (en généralisant l'idée du calcul d'un déterminant de \textsc{Vandermonde} par l'utilisation d'une fraction rationnelle) et en donner une écriture condensée dans le cas $a_i=b_i=i$.
\finenonce{005365}


\finexercice


\exercice{5366, rouget, 2010/07/06}
\enonce{005366}{****}
Soit $A=(a_{i,j})_{1\leq i,j\leq n}$ où, pour tout $i$ et tout $j$, $a_{i,j}\in\{-1,1\}$. Montrer que $\mbox{det}\;A$ est un entier divisible par $2^{n-1}$.
\finenonce{005366}


\finexercice


\exercice{5369, rouget, 2010/07/06}
\enonce{005369}{**I}
Soit $A=(a_{i,j})_{1\leq i,j\leq n}$ et $B=(b_{i,j})_{1\leq i,j\leq n}$ avec $b_{i,j}=(-1)^{i+j}a_{i,j}$. Montrer que $\mbox{det}B=\mbox{det}A$.

\finenonce{005369}


\finexercice


\exercice{5371, rouget, 2010/07/06}
\enonce{005371}{**** Déterminant circulant}
Soit $A=
\left(
\begin{array}{cccccc}
a_1&a_2&\ldots& &\ldots&a_n\\
a_n&a_1&a_2& & &a_{n-1}\\
a_{n-1}&a_n&a_1& & &a_{n-2}\\
\vdots& &\ddots&\ddots& &\vdots\\
 & & &\ddots&\ddots&\vdots\\
a_2&a_3&\ldots& &a_n&a_1
\end{array}
\right)
$ et $P=(\omega^{(k-1)(l-1)})_{1\leq k,l\leq n}$ où $\omega=e^{2i\pi/n}$. Calculer $P^2$ et $PA$. En déduire $\mbox{det}A$.
\finenonce{005371}


\finexercice


\exercice{5372, rouget, 2010/07/06}
\enonce{005372}{***I}
Calculer $\mbox{det}(\mbox{com}A)$ en fonction de $\mbox{det}A$ puis étudier le rang de $\mbox{com}A$ en fonction du rang de $A$.

\finenonce{005372}


\finexercice


\exercice{5373, rouget, 2010/07/06}
\enonce{005373}{***I Dérivée d'un déterminant}
Soient $a_{i,j}$ ($(i,j)$ élément de $\{1,...,n\}^2$) $n^2$ fonctions de $\Rr$ dans $\Rr$, dérivables sur $\Rr$ et $A=(ai,j)_{1\leq i,j\leq n}$. Calculer la dérivée de la fonction $x\mapsto\mbox{det}(A(x))$.

Applications. Calculer \begin{enumerate}
 \item  $\left|\begin{array}{ccccc}
x+1&1&\ldots& &1\\
1&x+1&\ddots& &\vdots\\
\vdots&\ddots&\ddots&\ddots&\vdots\\
\vdots& &\ddots&\ddots&1\\
1&\ldots&\ldots&1&x+1
\end{array}
\right|$   \item  $\left|\begin{array}{ccccc}
x+a_1&x&\ldots& &x\\
x&x+a_2&\ddots& &\vdots\\
\vdots&\ddots&\ddots&\ddots&\vdots\\
\vdots& &\ddots&\ddots&x\\
x&\ldots&\ldots&x&x+a_n
\end{array}
\right|$
\end{enumerate}
\finenonce{005373}


\finexercice


\exercice{5374, rouget, 2010/07/06}
\enonce{005374}{***I}
Calculer

\begin{enumerate}
 \item  $\left|\begin{array}{ccccc}
0&1&\ldots& &1\\
1&0&\ddots& &\vdots\\
\vdots&\ddots&\ddots&\ddots&\vdots\\
\vdots& &\ddots&0&1\\
1&\ldots&\ldots&1&0
\end{array}
\right|$ et $\left|\begin{array}{ccccc}
1&1&\ldots& &1\\
1&0&\ddots& &\vdots\\
\vdots&\ddots&\ddots&\ddots&\vdots\\
\vdots& &\ddots&0&1\\
1&\ldots&\ldots&1&0
\end{array}
\right|$   \item  $\mbox{det}((i+j-1)^2)$  \item  $\left|\begin{array}{ccccc}
a&b&\ldots& &b\\
b&a&\ddots& &\vdots\\
\vdots&\ddots&\ddots&\ddots&\vdots\\
\vdots& &\ddots&a&b\\
b&\ldots&\ldots&b&a
\end{array}
\right|$

 \item  $\left|\begin{array}{ccccc}
a_1+x&c+x&\ldots&\ldots&c+x\\
b+x&a_2+x&\ddots& &\vdots\\
\vdots&\ddots&\ddots&\ddots&\vdots\\
\vdots& &\ddots&a_{n-1}+x&c+x\\
b+x&\ldots&\ldots&b+x&a_n+x
\end{array}
\right|$$b$, $c$ complexes distincts   \item  $\left|\begin{array}{ccccc}
2&1&0&\ldots&0\\
1&2&\ddots&\ddots&\vdots\\
0&\ddots&\ddots&\ddots&0\\
\vdots&\ddots&\ddots&2&1\\
1&\ldots&0&1&2
\end{array}
\right|$.
\end{enumerate}
\finenonce{005374}


\finexercice


\exercice{5376, rouget, 2010/07/06}
\enonce{005376}{}
Donner une base du sous-espace vectoriel de $\Rr^5$ défini par~:

$$\left\{
\begin{array}{l}
x_1+2x_2-x_3+3x_4+x_5=0\\
x_2+x_3-2x_4+2x_5=0\\
2x_1+x_2-5x_3-4x_5=0
\end{array}
\right..$$
\finenonce{005376}


\finexercice


\exercice{5637, rouget, 2010/10/16}
\enonce{005637}{***I Déterminants de \textsc{Vandermonde}}
Soient $x_0$,..., $x_{n-1}$ $n$ nombres complexes. Calculer $\text{Van}(x_0,...,x_{n-1})= \text{det}(x_{j-1}^{i-1})_{1\leqslant i,j\leqslant n}$.
\finenonce{005637}


\finexercice


\exercice{5638, rouget, 2010/10/16}
\enonce{005638}{****I Déterminant de \textsc{Cauchy}}
Soient $a_1$,..., $a_n$, $b_1$,..., $b_n$ $2n$ nombres complexes tels que toutes les sommes $a_i+b_j$, $1\leqslant i,j\leqslant n$, soient non nulles. Calculer $C_n=\text{det}\left(\frac{1}{a_i+b_j}\right)_{1\leqslant i,j\leqslant n}$. Cas particulier : $\forall i\in\llbracket1,n\rrbracket$, $a_i= b_i =i$ (déterminant de \textsc{Hilbert}).
\finenonce{005638}


\finexercice


\exercice{5640, rouget, 2010/10/16}
\enonce{005640}{**}
Calculer $\text{det}(\sin(a_i+a_j))_{1\leqslant i,j\leqslant n}$ où $a_1$,..., $a_n$ sont $n$ réels donnés ($n\geqslant 2$).
\finenonce{005640}


\finexercice


\exercice{5641, rouget, 2010/10/16}
\enonce{005641}{**}
Calculer $\text{det}(a_i+b_j)_{1\leqslant i,j\leqslant n}$ où $a_1$,..., $a_n$, $b_1$,\ldots, $b_n$ sont $2n$ complexes donnés.
\finenonce{005641}


\finexercice


\exercice{5642, rouget, 2010/10/16}
\enonce{005642}{**}
Calculer $\text{det}((a+i+j)^2)_{1\leqslant i,j\leqslant n}$ où $a$ est un complexe donné.
\finenonce{005642}


\finexercice


\exercice{5643, rouget, 2010/10/16}
\enonce{005643}{****}
Soient $x_1$,..., $x_n$ $n$ entiers naturels tels que $x_1<...<x_n$. A l'aide du calcul de $\text{det}(C_{x_j}^{i-1})_{1\leqslant i,j\leqslant n}$, montrer
que $\prod_{1\leqslant i,j\leqslant n}^{}\frac{x_j-x_i}{j-i}$ est un entier naturel.
\finenonce{005643}


\finexercice


\exercice{5644, rouget, 2010/10/16}
\enonce{005644}{**** Déterminants circulants}
Soient $a_0$,...,$a_{n-1}$ $n$ nombres complexes. Calculer $\left|
\begin{array}{ccccc}
a_0&a_1&\ldots&a_{n-2}&a_{n-1}\\
a_{n-1}&a_0&\ddots& &a_{n-2}\\
\vdots&\ddots&\ddots&\ddots&\vdots\\
a_2& &\ddots&\ddots&a_1\\
a_1&a_2&\ldots&a_{n-1}&a_0
\end{array}
\right|=\text{det}A$.
Pour cela, on calculera d'abord $A\Omega$ où $\Omega=(\omega^{(j-1)(k-1)})_{1\leqslant j,k\leqslant n}$ avec $\omega=e^{2i\pi/n}$.
\finenonce{005644}


\finexercice


\exercice{5645, rouget, 2010/10/16}
\enonce{005645}{**I}
\begin{enumerate}
 \item  Soient $a_{i,j}$, $1\leqslant i,j\leqslant n$, $n^2$ fonctions dérivables sur $\Rr$ à valeurs dans $\Cc$. Soit $d=\text{det}(a_{i,j})_{1\leqslant i,j\leqslant n}$.

Montrer que $d$ est dérivable sur $\Rr$ et calculer $d'$.

\item  Application : calculer $d_n(x)=\left|
\begin{array}{cccc}
x+1&1&\ldots&1\\
1&\ddots&\ddots&\vdots\\
\vdots&\ddots&\ddots&1\\
1&\ldots&1&x+1
\end{array}
\right|$.
\end{enumerate}
\finenonce{005645}


\finexercice


\exercice{5646, rouget, 2010/10/16}
\enonce{005646}{***}
Soient $A$ et $B$ deux matrices carrées réelles de format $n$. Montrer que le déterminant de la matrice $\left(
\begin{array}{cc}
A&-B\\
B&A
\end{array}
\right)$ de format $2n$  est un réel positif.
\finenonce{005646}


\finexercice


\exercice{5647, rouget, 2010/10/16}
\enonce{005647}{***}
Soient $A$, $B$, $C$ et $D$ quatre matrices carrées de format $n$. Montrer que si $C$ et $D$ commutent et si $D$ est inversible alors $\text{det}\left(
\begin{array}{cc}
A&B\\
C&D
\end{array}
\right)=\text{det}(AD-BC)$. Montrer que le résultat persiste si $D$ n'est pas inversible.
\finenonce{005647}


\finexercice


\exercice{5649, rouget, 2010/10/16}
\enonce{005649}{**I}
\label{exo:rou15}
Soient $a_0$, ... , $a_{n-1}$ $n$ nombres complexes et $A=\left(
\begin{array}{ccccc}
0&\ldots&\ldots&0&a_0\\
1&\ddots& &\vdots&a_1\\
0&\ddots&\ddots&\vdots&\vdots\\
\vdots&\ddots&\ddots&0&\vdots\\
0&\ldots&0&1&a_{n-1}
\end{array}
\right)$. Calculer $\text{det}(A-xI_n)$.
\finenonce{005649}


\finexercice


\exercice{5650, rouget, 2010/10/16}
\enonce{005650}{**}
Calculer les déterminants suivants :

\begin{enumerate}
 \item  $\text{det}A$ où $A\in M_{2n}(\Kk)$ est telle que $a_{i,i}= a$ et $a_{i,2n+1-i}=b$ et $a_{i,j}= 0$ sinon.

\item  $\left|\begin{array}{cccccc}
1&0&\ldots&\ldots&0&1\\
0&0& & &0&0\\
\vdots& & & & &\vdots\\
\vdots& & & & &\vdots\\
0&0& & &0&0\\
1&0&\ldots&\ldots&0&1
\end{array}
\right|$ \item  $\left|\begin{array}{ccccc}
1&\ldots& &\ldots&1\\
\vdots&0&1&\ldots&1\\
 &1&\ddots&\ddots&\vdots\\
\vdots&\vdots&\ddots&\ddots&1\\
1&1&\ldots&1&0
\end{array}
\right|$ et $\left|\begin{array}{ccccc}
0&1&\ldots&\ldots&1\\
1&\ddots&\ddots& &\vdots\\
\vdots&\ddots& &\ddots&\vdots\\
\vdots& &\ddots&\ddots&1\\
1&\ldots&\ldots&1&0
\end{array}
\right|$ $(n\geqslant2$)

\item (I) $\left|\begin{array}{cccc}
a&b&\ldots&b\\
b&\ddots&\ddots&\vdots\\
\vdots&\ddots&\ddots&b\\
b&\ldots&b&a
\end{array}
\right|$ $(n\geqslant2$).
\end{enumerate}
\finenonce{005650}


\finexercice


\exercice{6885, exo7, 2012/09/05}
\video{ifWv3IZcAKU}
% D'après #1134 barraud et #2755 tumpach
\enonce{006885}{}
  Calculer les déterminants des matrices suivantes :
$$
  \begin{pmatrix}
    7 & 11 \\
    -8 & 4
  \end{pmatrix}
\quad
  \begin{pmatrix}
    1 & 0 & 6 \\
    3 & 4 & 15\\
    5 & 6 & 21
  \end{pmatrix}
\quad
  \begin{pmatrix}
    1 & 0 & 2 \\
    3 & 4 & 5  \\
    5 & 6  & 7
  \end{pmatrix}
\quad
  \begin{pmatrix}
    1 & 0 & -1 \\
    2 & 3 & 5  \\
    4 & 1 & 3
  \end{pmatrix}
$$

$$
  \begin{pmatrix} 0 & 1 & 2 & 3\\ 1 & 2 & 3 & 0\\ 2 & 3 & 0 & 1\\3 & 0 & 1 & 2\end{pmatrix}
  \begin{pmatrix} 0 & 1 & 1 & 0\\ 1 & 0 & 0 & 1\\ 1 & 1 & 0 & 1\\  1 & 1 & 1 &0\end{pmatrix}
  \begin{pmatrix} 1 & 2 & 1 & 2\\ 1 & 3 & 1 & 3\\ 2 & 1 & 0& 6\\ 1 & 1& 1&7\end{pmatrix}
$$


\finenonce{006885}



\finexercice


\exercice{6886, exo7, 2012/09/05}
\video{ibKDAQnNURg}

% D'après #1121 ridde, #2756 tumpach, #1130 barraud; #1124 liousse
\enonce{006886}{}  

Calculer les déterminants des matrices suivantes :

$$
\begin{pmatrix}
a&b&c\\c&a&b\\b&c&a
\end{pmatrix}
\begin{pmatrix}
1&0&0&1 \\ 0&1&0&0 \\ 1&0&1&1 \\ 2&3&1&1
\end{pmatrix}
\begin{pmatrix}
-1 & 1 & 1 & 1\\ 1 & -1 & 1 & 1\\ 1 & 1 & -1& 1\\ 1 & 1& 1&-1
\end{pmatrix}
\begin{pmatrix}
10 & 0 & -5 & 15 \\ -2 & 7 & 3 & 0 \\ 8 & 14 & 0 & 2 \\ 0 & -21 & 1 & -1
\end{pmatrix}
$$
 
$$
\begin{pmatrix}
a&a&b&0 \\  a&a&0&b \\  c&0&a&a \\ 0&c&a&a
\end{pmatrix}
\begin{pmatrix}
1&0&3&0&0 \\ 0&1&0&3&0 \\ a&0&a&0&3 \\ b&a&0&a&0 \\ 0&b&0&0&a  
\end{pmatrix}
\begin{pmatrix}
1&0&0&1&0 \\ 0&-4&3&0&0 \\ -3&0&0&-3&-2 \\ 0&1&7&0&0 \\ 4&0&0&7&1  
\end{pmatrix}
$$
\finenonce{006886}




\finexercice


\exercice{6887, barraud, 2012/09/05}
\video{D-wMy5vRJo4}
% D'apres #1139 barraud
\enonce{006887}{}
  Calculer les déterminants suivant~:
$$
 \left\vert
   \begin{matrix}
     a_1   &a_2   &\cdots&a_n    \\
     a_1   &a_1   &\ddots&\vdots \\
     \vdots&\ddots&\ddots&a_2    \\
     a_1   &\cdots&a_1   &a_1
   \end{matrix}
 \right\vert
\qquad
\left\vert
\begin{matrix}
  1    &      &        &1 \\
  1    &1      & (0)   & \\
       &\ddots &\ddots & \\  
  (0)  &       &1      &1
\end{matrix}
\right\vert
\qquad
\left\vert
\begin{matrix}
    a+b  &    a   & \cdots &  a       \\
     a   &   a+b  & \ddots & \vdots   \\
  \vdots & \ddots & \ddots &  a       \\
     a   & \cdots &    a   & a+b 
\end{matrix}
\right\vert
$$
\finenonce{006887}


\finexercice
\section{200.03 Système linéaire, rang}


\exercice{1163, barraud, 2003/09/01}
\video{I_lQjh6jGgo}
% d'après 1163 + 1164 : barraud
\enonce{001163}{}
  Résoudre les systèmes suivants
  $$
  \left\{
    \begin{array}{rcrcrcl}
         x &+&    y &-&    z &=&  0  \\
         x &-&    y & &      &=&  0  \\
         x &+&    4y &+&   z &=&  0  
    \end{array}
  \right.  
 \qquad
\left\{
    \begin{array}{rcrcrcl}
       x &+&  y &+& 2z &=& 5  \\
       x &-&  y &-&  z &=& 1  \\
       x & &    &+&  z &=& 3    
    \end{array}
  \right. 
\qquad 
  \left\{
    \begin{array}{rcrcrcl}
      3x &-& y  &+& 2z &=& a \\
      -x &+& 2y &-& 3z &=& b \\
       x &+& 2y &+&  z &=& c 
    \end{array}
  \right.
  $$
\finenonce{001163}



\finexercice


\exercice{1164, barraud, 2003/09/01}

\enonce{001164}{}
  Sans chercher à résoudre les systêmes suivants, discuter la nature
  de leurs ensembles de solution :
  $$
  \left\{
    \begin{array}{rc@{}rc@{}rcl}
         x &+&    y &-&    z &=&  0  \\
         x &-&    y & &      &=&  0  \\
         x &+&    y &+&    z &=&  0  
    \end{array}
  \right.  \qquad \left\{
    \begin{array}{rc@{}rc@{}rcl}
         x &+&  3 y &+&  2 z &=&  1  \\
       2 x &-&  2 y & &      &=&  2  \\
         x &+&    y &+&    z &=&  2  
    \end{array}
  \right.  \qquad \left\{
    \begin{array}{rc@{}rc@{}rcl}
         x &+&  3 y &+&  2 z &=&  1  \\
       2 x &-&  2 y & &      &=&  2  \\
         x &+&    y &+&    z &=&  3   
    \end{array}
  \right.
  $$
\finenonce{001164}



\finexercice



\exercice{1165, barraud, 2003/09/01}

\enonce{001165}{}
  Soient $x_0$,$x_1$,...,$x_n$, $n+1$ réels distincts, et
  $y_0$,$y_1$,...,$y_n$, $n+1$ réels (distincts ou non).
  
  Montrer qu'il existe un unique polynôme $P$ tel que :
  $$
  \forall i\in\{0,...,n\}\quad P(x_i)=y_i
  $$
\finenonce{001165}



\finexercice



\exercice{1166, cousquer, 2003/10/01}

\enonce{001166}{}
 Résoudre, suivant les valeurs de $m$~:
$$(S_1)\; \left\{\begin{array}{rcl}
    x+(m+1)y &=& m+2 \\
    mx+(m+4)y &=&3 
\end{array}\right.
\qquad
(S_2)\; \left\{\begin{array}{rcl}
    mx+(m-1)y & =& m+2 \\
    (m+1)x-my &=&5m+3
\end{array}\right.$$

\finenonce{001166}



\finexercice



\exercice{1167, cousquer, 2003/10/01}

\enonce{001167}{}
 Écrire les conditions, portant sur les réels $a$, $b$, $c$, pour que les
systèmes suivants admettent des solutions non nulles~; expliciter
ces solutions.
$$(S_1)\; \left\{\begin{array}{rcl}
    x+y+z &=&0 \\
    (b+c)x+(c+a)y+(a+b)z &=&0 \\
    bcx+acy+abz &=&0
\end{array}\right.
\qquad(S_2)\; \left\{\begin{array}{rcl}
    x-a(y+z) &=&0 \\ 
    y-b(x+z) &=& 0\\
    z-c(x+y)&=&0
\end{array}\right.$$

\finenonce{001167}



\finexercice



\exercice{1168, cousquer, 2003/10/01}

\enonce{001168}{}
 Résoudre et discuter suivant les valeurs de $b_1$, $b_2$, $b_3$~et $b_4$~:
$$\displaylines{
    (S_1)\;\left\{\begin{array}{rcl}
        x+3y+4z+7t &=&b_1 \\ 
        x+3y+4z+5t &=&b_2 \\
        x+3y+3z+2t &=&b_3 \\ 
        x+y+z+t &=&b_4
    \end{array}\right.
    \qquad(S_2)\;\left\{\begin{array}{rcl}
        x+3y+5z+3t &=&b_1 \\ 
        x+4y+7z+3t &=&b_2\\
        y+2z &=& b_3\\ 
        x+2y+3z+2t &=&b_4
    \end{array}\right.\cr
    (S_3)\;\left\{\begin{array}{rcl}
        x+y+2z-t &=&b_1 \\
        -x+3y+t &=&b_2 \\
        2x-2y+2z-2t &=&b_3 \\
        2y+z&=&b_4\\
    \end{array}\right.
    \qquad(S_4)\;\left\{\begin{array}{rcl}
        x+2y+z+2t &=&b_1 \\
        -2x-4y-2z-4t &=&b_2\\
        -x-2y-z-2t &=& b_3\\
        3x+6y+3z+6t &=&b_4 
    \end{array}\right.}$$

\finenonce{001168}



\finexercice



\exercice{1169, cousquer, 2003/10/01}
\video{JfeFfZxwyVI}
\enonce{001169}{}
 Discuter et résoudre suivant les valeurs des réels 
$\lambda$, $a$, $b$, $c$, $d$ le système :
$$ (S)\;\left\{\begin{array}{rcl}
    (1+\lambda)x+y+z+t &=&a\\ 
    x+(1+\lambda)y+z+t &=&b \\
    x+y+(1+\lambda)z+t &=&c \\ 
    x+y+z+(1+\lambda)t &=&d
\end{array}\right.$$

\finenonce{001169}



\finexercice


\exercice{1170, cousquer, 2003/10/01}

\enonce{001170}{}
 Discuter et résoudre suivant les valeurs des réels $\lambda$~et $a$~:
$$(S)\;\left\{\begin{array}{rcl}
    3x+2y-z+t &=&\lambda\\
    2x+y-z &=&\lambda-1 \\
    5x+4y-2z &=&2\lambda\\
    (\lambda+2)x+(\lambda+2)y-z&=&3\lambda+a \\
    3x-z+3t &=& -\lambda^2
\end{array}\right.$$

\finenonce{001170}



\finexercice



\exercice{1171, liousse, 2003/10/01}

\enonce{001171}{}
Mettre sous forme matricielle et r\'esoudre les syst\`emes suivants.\\
\begin{enumerate}
\item $ \left\{\begin{array}{rcl}2x+y+z&=&3\\3x-y-2z&=&0\\
x+y-z&=&-2\\x+2y+z&=&1\end{array}\right.$ \hfill
\item $ \left\{\begin{array}{rcl}x+y+z+t&=&1\\x-y+2z-3t&=&2\\
2x+4z+4t&=&3\\2x+2y+3z+8t&=&2\\5x+3y+9z+19t&=&6\end{array}\right.$ 

\item $ \left\{\begin{array}{rcl}2x+y+z+t&=&1\\x+2y+3z+4t&=&2\\
3x-y-3z+2t&=&5\\5y+9z-t&=&-6\end{array}\right.$ \hfill

\item $ \left\{\begin{array}{rcl}x-y+z+t&=&5\\2x+3y+4z+5t&=&8\\
3x+y-z+t&=&7\end{array}\right.$ 

\item $\left\{\begin{array}{rcl}x+2y+3z&=&0\\2x+3y-z&=&0\\
3x+y+2z&=&0\end{array}\right.$ 
\end{enumerate}
\finenonce{001171}



\finexercice



\exercice{1172, liousse, 2003/10/01}

\enonce{001172}{}
 Calculer les d\'eterminants suivants.
$$D_1=\left\vert\begin{array}{ccc}1&3&2\\1&3&3\\1&2&1\end{array}\right\vert\hbox{, }
D_2=\left\vert\begin{array}{ccc}1&1&1\\3&3&2\\2&3&1\end{array}\right\vert\hbox{, }
D_3=\left\vert\begin{array}{ccc}5&-3&13\\0&-1&-16\\0&0&2\end{array}\right\vert
D_4=\left\vert\begin{array}{ccc}1&0&0\\0&{\sqrt{3}\over2}&-{1\over 2}\\
0&{1\over 2}&{\sqrt{3}\over2}\end{array}\right\vert
D_5=\left\vert\begin{array}{ccc}0&0&1\\1&0&0\\0&1&0\end{array}\right\vert$$

\finenonce{001172}



\finexercice



\exercice{1173, cousquer, 2003/10/01}

\enonce{001173}{}
 Résoudre et discuter le système linéaire suivant~:
$$(S)\;\left\{\begin{array}{rcl}
    x_1+ x_2+3x_3+10x_4+  x_5 & = & b_1 \\
    x_1+2x_2+ x_3+ 4x_4+ 7x_5 & = & b_2 \\
    x_1+3x_2+4x_3+13x_4+ 8x_5 & = & b_3 \\
    x_1+4x_2+2x_3+ 7x_4+14x_5 & = & b_4
\end{array}\right.$$
\finenonce{001173}



\finexercice



\exercice{1174, cousquer, 2003/10/01}

\enonce{001174}{}
On considère l'application $f$ de $\mathbb{R}^5$
dans $\mathbb{R}^4$ qui à un élément $X=(x_1, x_2, x_3, x_4, x_5)$ associe
l'élément $Y=(y_1, y_2, y_3, y_4)$, défini par~:
$$(S)\;\left\{\begin{array}{rcl}
    x_1+ x_2+3x_3+10x_4+  x_5 & = & y_1 \\
    x_1+2x_2+ x_3+ 4x_4+ 7x_5 & = & y_2 \\
    x_1+3x_2+4x_3+13x_4+ 8x_5 & = & y_3 \\
    x_1+4x_2+2x_3+ 7x_4+14x_5 & = & y_4
\end{array}\right.$$

\begin{enumerate}
\item
Montrer que $f$ est linéaire.

\item
 On considère $A$ l'ensemble des solutions de $(S_H)$. 
$$(S_H)\;\left\{\begin{array}{rcl}
    x_1+ x_2+3x_3+10x_4+  x_5 & = & 0 \\
    x_1+2x_2+ x_3+ 4x_4+ 7x_5 & = & 0 \\
    x_1+3x_2+4x_3+13x_4+ 8x_5 & = & 0 \\
    x_1+4x_2+2x_3+ 7x_4+14x_5 & = &0 
\end{array}\right.$$
Quelle est la nature de~$A$~? Que représente~$A$ pour l'application $f$~? 
Donner une base de~$A$~; quelle est la dimension de~$A$~? Donner un système
minimal d'équations qui définissent~$A$.

\item
 Dans l'espace $\mathbb{R}^4$, on considère les cinq vecteurs~:
$V_1=(1,1,1,1)$, $V_2=(1,2,3,4)$, $V_3=(3,1,4,2)$, $V_4=(10,4,13,7)$,
$V_5=(1,7,8,14)$.
Que représentent ces vecteurs pour l'application~$f$~? Trouver une base de
$\mbox{Im} f$.

\item
 On considère le système $(S)$ où les inconnues sont les $x_i$, et 
où les $y_j$ sont des paramètres.
Comment interpréter les conditions de possibilité de ce système du
point de vue de~$f$~?

\item
 Donner une interprétation du théorème du rang relativement à ce
système. Quel est le lien entre le rang de~$f$ et le rang du système~?
\end{enumerate}
\finenonce{001174}



\finexercice



\exercice{1175, legall, 1998/09/01}

\enonce{001175}{}
Pour tout $ a  $ r\' eel, on consid\`ere la matrice
$ A  $ et le syst\`eme lin\' eaire
 $ (S) $ d\' efinis par~:
$$ A=\begin{pmatrix}  a & 1 & 1  & 1  \cr
                                   1 & a & 1 &  1  \cr
                                    1 & 1 & a & 1 \cr
                                     1 & 1 & 1 &  a   \cr \end{pmatrix}  \hskip10mm
(S)\hskip2mm
\left\{ \begin{matrix} ax &+ & y & + & z & + & t & = & 1 \cr
x &+ & ay & + & z & + & t & = & 1 \cr
x &+ & y & + & az & + & t & = & 1 \cr
x &+ & y & + & z & + & at & = & 1 \cr\end{matrix}\right.$$
aux inconnues r\' eelles $ x ,  y  , z ,  t .$
\begin{enumerate}
    \item Discuter le rang de $ A  $ suivant les valeurs de $ a  .$
    \item Pour quelles valeurs de
 $ a $ le syst\`eme $ (S) $ est-il de Cramer~? Compatible~?
Incompatible~?
    \item Lorsqu'il est de Cramer, r\' esoudre $ (S) $ avec un minimum d'op\'
erations (on pourra montrer
d'abord que l'on a n\' ecessairement $ x=y=z=t .).$
    \item  Retrouver 3. par application des formules de Cramer.
\end{enumerate}
\finenonce{001175}



\finexercice



\exercice{1176, ridde, 1999/11/01}

\enonce{001176}{}
D\'eterminer le noyau de la matrice $\begin{pmatrix}
1&-1&1\\0&1&1\\2&3&7
\end{pmatrix}$
\finenonce{001176}



\finexercice



\exercice{1177, ridde, 1999/11/01}

\enonce{001177}{}
Soit $A = \begin{pmatrix}
2&2&0\\1&2&1\\0&2&2\\
\end{pmatrix}$. D\'eterminer les $\lambda \in \Rr$ tels que $ \exists X \in \Rr^3
-\left\{ (0, 0, 0)\right\}$ tel que $AX = \lambda X$. Pour chaque $\lambda$
d\'eterminer $E_{\lambda} = \left\{ X \in \Rr^3 / AX = \lambda X\right\}$.
\finenonce{001177}



\finexercice



\exercice{1178, ridde, 1999/11/01}
\video{LKkeqctcOMw}
\enonce{001178}{}
Trouver les solutions de
$$\left\{ \begin{array}{l} 3x + 2z = 0 \\ 3y + z + 3t = 0 \\ x + y + z + t = 0 \\
2x-y + z-t = 0 \end{array}\right.$$
\finenonce{001178}


\finexercice


\exercice{1179, ridde, 1999/11/01}

\enonce{001179}{}
R\'esoudre suivant les valeurs de $a \in \Rr$
$ \begin{cases} x + ay + a^2z = 0 \\ a^2x + y + az = 0 \\ ax + a^2y + z = 0
\end{cases}$.
\finenonce{001179}



\finexercice



\exercice{1180, ridde, 1999/11/01}

\enonce{001180}{}
R\'esoudre suivant les valeurs de $a \text{ et } \mu \in \Rr$
$ \begin{cases} ax + y + z + t = 1 \\ x + ay + z + t = \mu \\ x + y + az + t =
\mu^2 \\ x + y + z + at = \mu^3
\end{cases}$.
\finenonce{001180}



\finexercice



\exercice{1181, ridde, 1999/11/01}

\enonce{001181}{}
Inverser en utilisant un syst\`eme lin\'eaire la matrice $\begin{pmatrix}
1&1&1 \\ 2&1&1 \\ 1&2&1
\end{pmatrix}$.
\finenonce{001181}



\finexercice



\exercice{1182, ridde, 1999/11/01}

\enonce{001182}{}
R\'esoudre $\begin{cases} x + y + z = 1 \\ ax + by + cz = d \\ a^2x + b^2y + c^2z
 = d^2 \end{cases}$.
\finenonce{001182}



\finexercice



\exercice{1183, ridde, 1999/11/01}

\enonce{001183}{}
R\'esoudre $\begin{cases} -cy + bz = \alpha \\ cx-az = \beta \\ -bx + ay = \gamma
\end{cases}$.
\finenonce{001183}



\finexercice



\exercice{1184, liousse, 2003/10/01}

\enonce{001184}{}
 Soit $F$ le sous-espace 
vectoriel de $\Rr^4$ des \'el\'ements $(x,y,z,t)$ qui satisfont :
$$\left\{
\begin{array}{rcl}
x+y+z+3t&=&0\\
2x+3y+4t&=&0\\
2x+5y-4z&=&0\\
\end{array}\right.$$
Donner une base de $F$ et sa dimension.
\finenonce{001184}



\finexercice



\exercice{1185, liousse, 2003/10/01}

\enonce{001185}{}
On consid\`ere le syst\`eme $$(S) : \left \lbrace
\begin{array}{ll}
   x +  y+ z+  t &=0 \\
  x -  y-2z+2t &=0 \\
 2x + y+  z  &= 0 \\
\end{array} \right.$$
\begin{enumerate}
\item R\'esoudre le syst\`eme $(S)$ puis indiquer son rang.

\item Montrer que l'ensemble des solutions de $(S)$ est un
sous-espace vectoriel de $\mathbb R^4$, indiquer sa dimension  
et en donner une base. 
\end{enumerate}
\finenonce{001185}



\finexercice



\exercice{1186, liousse, 2003/10/01}

\enonce{001186}{}
L'objectif de ce probl\`eme est de r\'esoudre l'\'enigme du berger :\\
Un berger poss\`ede un troupeau de 101 moutons et remarque par hasard la 
propri\'et\'e suivante : pour chaque mouton, il peut trouver une fa\c{c}on 
de scinder 
le troupeau des 100
autres moutons en deux troupeaux de 50 moutons et de m\^eme poids total. Il en
d\'eduit que tous les moutons ont le m\^eme poids. Comment a-t-il fait ? 
On montre, 
dans un premier temps, un r\'esultat utile pour la d\'emonstration finale.\\
\begin{enumerate}
\item
\begin{enumerate}
\item Montrer par r\'ecurrence que le d\'eterminant de 
toute matrice carr\'ee, dont les \'el\'ements diagonaux sont des nombres impairs, 
et dont tous les autres sont des nombres pairs, est un nombre impair.\\
\item En d\'eduire qu'une matrice de cette forme est inversible.\\
\end{enumerate}
\item L'objectif de cette question est de r\'esoudre l'\'enigme du berger.
On note $B$ la matrice carr\'ee de taille 101 construite de la mani\`ere suivante:\\
On num\'erote les moutons de 1 \`a 101. Quand le berger retire le i\`eme mouton du
troupeau, il s\'epare alors le reste du troupeau en deux
troupeaux  \'egaux ( troupeau A, troupeau B) et de
m\^eme poids. On note alors $B_{i,j}$ les coefficients de la i\`eme ligne de la
matrice $B$ obtenu de la fa\c{c}on suivante
$$ B_{i,j}= \left \{ \begin{array}{ll} 
1 & \hbox{ si } j=i \\
0 & \hbox{ si le j-i\`eme mouton se trouve dans le troupeau A }\\
2 & \hbox{ si le j-i\`eme mouton se trouve dans le troupeau B }.
\end{array} \right.$$

On note $X$ la matrice de taille $101\times 1$ constitu\'ee des 
poids des moutons
$$X=\left ( \begin{array}{c}
\hbox{ poids du monton 1} \\
\hbox{ poids du mouton 2} \\
\vdots \\
\hbox{ poids du mouton 100 }\\
\hbox{ poids du mouton 101 }
\end{array} \right).$$

On note $M$ le poids total du troupeau.\\
\begin{enumerate}
\item Calculer 
$$ B\times \left ( \begin{array}{c}
1\\
1\\
\vdots \\
1\\
1\\
\end{array} \right).$$\\
\item Calculer
$$ BX. $$
\item Montrer que $B$ est inversible.\\
\item En d\'eduire $X$ et r\'esoudre l'\'enigme du berger.
\end{enumerate}
\end{enumerate}
\finenonce{001186}



\finexercice



\exercice{1187, barraud, 2003/09/01}

\enonce{001187}{}
  Pour quelles valeurs de $a$ la matrice
  $$
  A=
  \begin{pmatrix}
    1 & 1 & 1  \\
    1 & 2 & 4  \\
    1 & 3 & a
  \end{pmatrix}
  $$
  est-elle inversible ? Calculer dans ce cas son inverse.
\finenonce{001187}



\finexercice



\exercice{1188, barraud, 2003/09/01}

\enonce{001188}{}
  Soient $a$ et $b$ deux réels, et $A$ la matrice
  $$
  A=
  \begin{pmatrix}
    a &  2 & -1 &  b \\
    3 &  0 &  1 & -4 \\
    5 &  4 & -1 &  2
  \end{pmatrix}
  $$
  Montrer que $\mathrm{rg}(A)\geq2$. Pour quelles valeurs de $a$ et $b$
  a-t-on $\mathrm{rg}(A)=2$ ?
\finenonce{001188}



\finexercice



\exercice{1189, barraud, 2003/09/01}

\enonce{001189}{}
  Soient $v_1=\begin{pmatrix} a\\b\\c \end{pmatrix}$ et $v_2=\begin{pmatrix} a'\\b'\\c' \end{pmatrix}$ deux
  vecteurs indépendants de $\R^3$. Donner, sous forme d'équation, une
  condition nécessaire et suffisante pour qu'un vecteur
  $w=\begin{pmatrix} x\\y\\z \end{pmatrix}$ appartienne à l'espace vectoriel engendré par
  $v_1$ et $v_2$.
  
  Même question pour un plan engendré par deux vecteurs de $\R^4$.
\finenonce{001189}



\finexercice



\exercice{1190, barraud, 2003/09/01}

\enonce{001190}{}
  Soit $u$ un endomorphisme de $E$, et $\mathcal{B}$ une base de $E$.
  Discuter dans chacun des cas ci-dessous la dimension du noyau de
  $u$.
  $$
  M_\mathcal{B}(u)=
  \begin{pmatrix}
     2 & 1 & a & 1 \\
    -1 & 1 & 1 & b \\
     0 & 0 & a & 1 \\
     0 & 0 & 1 & b
  \end{pmatrix}
  \qquad M_\mathcal{B}(u)=
  \begin{pmatrix}
     -1-\lambda & 2         &  1         \\
      4         & 1-\lambda & -2         \\
      0         & 0         &  3-\lambda
  \end{pmatrix}
  \qquad M_\mathcal{B}(u)=
  \begin{pmatrix}
     12-\lambda & -6         &  3         \\
     -9         & -5-\lambda &  3         \\
    -12         & -8         &  9-\lambda 
  \end{pmatrix}
  $$
\finenonce{001190}



\finexercice



\exercice{1191, barraud, 2003/09/01}

\enonce{001191}{}
  Discuter le rang de la matrice suivante en fonction des paramètres
  réels $x$ et $y$~:
$$
 A=
 \begin{pmatrix}
   1 & 2 & y \\
   0 & x & 1 \\
   1 & 0 & 2 \\
   1 & 2 & 1
 \end{pmatrix}
$$
\finenonce{001191}



\finexercice



\exercice{1192, barraud, 2003/09/01}

\enonce{001192}{}
  Sans chercher à le résoudre, discuter la nature des solutions du
  système suivant, en fonction de $\alpha,a,b$ et $c$~:
  $$
  \left\{
    \begin{array}{*8{c@{\;}}}
      x&-& y&-&\alpha z&=&a\\
      x&+&2y&+& z&=&b\\
      x&+& y&-& z&=&c
    \end{array}
    \right.
  $$
\finenonce{001192}



\finexercice



\exercice{2462, matexo1, 2002/02/01}

\enonce{002462}{}
{\bf Syst\`emes lin\'eaires.}

\begin{enumerate}
\item R\'esoudre le syst\`eme d'\'equations lin\'eaires
sur $\R$
$$\left\{
\begin{array}{rcrcrcc}
2x   &+&  y    &+&   z  &=&  3 \\
x    &-&  y    &+&   3z &=&  8 \\
x    &+&  2y   &-&   z  &=&  -3
\end{array}\right.$$
\item Le syst\`eme suivant admet-il des solutions sur $\R$\,:
$$\left\{
\begin{array}{rcrcrcc}
2x   &+& y    &+&  z  &=&  3 \\
x    &-& y    &+&  3z &=&  8 \\
x    &+& 2y   &-&  z  &=&  -3 \\
x    &+& y    &+&  2z &=&  -1
\end{array}\right.$$
\item On consid\`ere le syst\`eme pr\'ec\'edent, mais dont
coefficients et inconnues sont dans le corps $\Z/5\Z$. R\'esoudre
ce syst\`eme.
\end{enumerate}
\finenonce{002462}


\finexercice


\exercice{2463, matexo1, 2002/02/01}

\enonce{002463}{}
D\'eterminer les valeurs de $m$ pour lesquelles le syst\`eme
suivant admet des solutions diff\'erentes de $x=y=z=t=0$\,:
$$\left\{
\begin{array}{rcrcrcrcc}
 mx   &+&  y    &+&   z  &+&   t  &=&  0 \\
 x   &+&  (1+m)y &+&   z  &+&   t  &=&  0 \\
 x   &+&  y    &+&   (2+m)z  &+&   t  &=&  0 \\
 x   &+&  y    &+&   z  &+&   t  &=&  0
\end{array}\right.$$

\finenonce{002463}


\finexercice


\exercice{2464, matexo1, 2002/02/01}

\enonce{002464}{}
 Soit $a, b$ deux r\'eels diff\'erents. Montrer que le
syst\`eme lin\'eaire 
$$\left\{\
\begin{array}{rcrcrcccrcc}
 x_1   &+& x_2  &+&  x_3  &+&\ldots&+&  x_n &=&1\\
 bx_1  &+& ax_2 &+& ax_3  &+&\ldots&+& ax_n &=&c_1\\
 bx_1  &+& bx_2 &+& ax_3  &+&\ldots&+& ax_n &=&c_2\\
\ldots  & &\ldots & &\ldots  & &\ldots& &\ldots & &\ldots\\
 bx_1  &+& bx_2 &+& bx_3  &+&\ldots&+& ax_n &=&c_n
\end{array}\right.$$
admet une solution unique que l'on calculera.

\finenonce{002464}


\finexercice


\exercice{2465, matexo1, 2002/02/01}

\enonce{002465}{}
Soit $A$ une matrice carr\'ee d'ordre $n$ tridiagonale,
c'est-\`a-dire telle que $a_{i,j} = 0$ si $|i-j| > 1$. Montrer
qu'il existe une matrice triangulaire inf\'erieure $L$
v\'erifiant  $l_{i,j} = 0$ si $j > i+1$ et une triangulaire
sup\'erieure $U$ v\'erifiant $u_{i,i} = 1$ et $u_{i,j} = 0$ si $i
> j+1$ telles que $A = LU$, et que ces matrices sont uniques. En
d\'eduire la solution du syst\`eme lin\'eaire $Ax = b$, o\`u $b$
est un vecteur donn\'e dans $\R^n$.
\finenonce{002465}


\finexercice


\exercice{2478, matexo1, 2002/02/01}

\enonce{002478}{}
\catcode`\@=11
\def\system#1{\left\{\null\,\vcenter{\openup1\jot\m@th
\ialign{\strut\hfil$##$&$##$\hfil&&\enspace$##$\enspace&
\hfil$##$&$##$\hfil\crcr#1\crcr}}\right.}
\catcode`\@=12

R{\'e}soudre
$$S_1\system{
        &x&+&\cosh a\ &y&+&\cosh 2a\ &z&=&&\cosh 3a\cr
        \cosh a\ &x&+&\cosh 2a\ &y&+&\cosh 3a\ &z&=&&\cosh 4a\cr
        \cosh 2a\ &x&+&\cosh 3a\ &y&+&\cosh 4a\ &z&=&&\cosh 5a\cr}$$
et $$       
S_2  \system{
        &x_1&+&&x_2&+&\ldots&&+&&x_n&=&&1\cr
        &x_1&+&2&x_2&+&\ldots&&+&n&x_n&=&&0\cr
        &x_1&+&2^2&x_2&+&\ldots&&+&n^2&x_n&=&&0\cr
        &   & &   &   & &\vdots&& &   &   & &&\cr
        &x_1&+&2^{n-1}&x_2&+&\ldots&&+&n^{n-1}&x_n&=&&0\cr}   
$$
\finenonce{002478}

\finexercice


\exercice{2731, tumpach, 2009/10/25}
\enonce{002731}{}
D\'ecider, pour chacun des syst\`emes d'\'equations aux inconnues $x_1$, $x_2$, $\dots$, $x_n$ et aux param\`etres $s$, $t$, s'il est lin\'eaire~:
$$
\begin{array}{ll}
a) \left\{ \begin{array}{ccccc}x_1\sin(t) &+ & x_2 & = &3\\x_1 e^t & +&  3 x_2 & = & t^2\end{array}\right.
&
b) \left\{ \begin{array}{cccccc}\frac{1}{x_1} & + \frac{2}{x_2} & +\dots & +\frac{n}{x_n} & =   &n!\\x_1 & + \frac{x_2}{2} & + \dots & + \frac{x_n}{n} & =  &\frac{1}{n!}\end{array}\right.
\\\end{array} $$
$$
c)\left\{ \begin{array}{ccc}\sqrt{(x_1 + sx_2 + t)^2 - 4sx_2(x_1 + t)} & = & 0\\x_1 \ln s - \pi x_2 + e^t x_n & = & 0\end{array}\right.
$$

$$
d)\left\{ \begin{array}{ccc}(1 + sx_1)(3 + t x_2) - (2 + t x_1)(5 + s x_2) & = & 8\\(x_3 + s)^2 - (x_3 - s)^2 + x_2 & = & 0\end{array}\right.
$$
\finenonce{002731}



\finexercice


\exercice{2732, tumpach, 2009/10/25}
\enonce{002732}{}
En appliquant l'algorithme de Gauss, r\'esoudre le syst\`eme lin\'eaire suivant~:
$$
\mathcal{S} ~: \left\{
\begin{array}{ccccccccc} 
2x_1 & + & 4x_2 & - & 6x_3 & - & 2x_4 & = & 2\\
3x_1 & + & 6x_2 & - & 7x_3 & + & 4x_4 &= & 2\\
5x_1 & + & 10 x_2 & - & 11 x_3 & + & 6x_4 &   = & 3
\end{array}
 \right.. 
$$
\finenonce{002732}



\finexercice


\exercice{2733, tumpach, 2009/10/25}
\enonce{002733}{}
R\'esoudre le syst\`eme~:
$$
\left(\begin{array}{cccc}
1 & 0 & 1 & 2\\-1 & 2 & 0 & 1\\ 3 & 2 & 1 & 2\\4 & 3 & 2 & 3\\1 & -1 & 1 & -2
\end{array}
\right)
\left(\begin{array}{c}x \\ y \\ z \\ t\end{array}\right) = \left(\begin{array}{c} 1 \\ 2 \\ -1 \\ 3 \\ 1\end{array}\right).
$$
\finenonce{002733}



\finexercice


\exercice{2734, tumpach, 2009/10/25}
\enonce{002734}{}
R\'esoudre le syst\`eme~:
$$
\left(\begin{array}{ccccc}
1 & 0 & 2 & 1 & 1\\2 & 1 & 3 & -1 & 2\\ -2 & -1 & 1 & -3& 2\\3 & 2 & 0 & 1&-1
\end{array}
\right)
\left(\begin{array}{c}x \\ y \\ z \\ t\\u\end{array}\right) = \left(\begin{array}{c} 2 \\ 0 \\ 1 \\ 1 \end{array}\right).
$$
\finenonce{002734}



\finexercice


\exercice{2735, tumpach, 2009/10/25}
\enonce{002735}{}
Soit $a$ un nombre r\'eel. On \'etudie le syst\`eme lin\'eaire suivant~:
$$
\mathcal{S}_{a}~: \left\{\begin{array}{ccccccc} 
x & - & 2y & + & 3z & = & 2\\
x & + & 3y & - & 2z & = & 5\\
2x & - & y & + & az & = & 1
\end{array}\right.
$$
\begin{enumerate}
\item En fonction des valeurs du param\`etre $a$, d\'eterminer si le syst\`eme $\mathcal{S}_{a}$ peut~:
\begin{enumerate}
\item[(i)] n'admettre aucune solution ;
\item[(ii)] admettre exactement une solution ;
\item[(iii)] admettre une infinit\'e de solutions.
\end{enumerate}
\item R\'esoudre le syst\`eme $\mathcal{S}_a$ lorsque celui-ci admet une (des) solution(s).
\end{enumerate}
\finenonce{002735}



\finexercice


\exercice{2736, tumpach, 2009/10/25}
\enonce{002736}{}
Les vecteurs complexes $(z, w)$ et $(z', w')$ sont li\'es par la formule $(z', w') = (z + iw, (1+i)z + (1-2i)w)$. Un \'etudiant qui n'aime pas les nombres complexes pose $z = x + iy$, $w = u + iv$, $z' = x' + i y'$ et $w' = u'+iv'$.
\begin{enumerate}
\item Exprimer $(x', y', u', v')$ en fonction de $(x, y, u, v)$.
\item R\'esoudre le syst\`eme $(x', y', u', v') = (1, 2, 3, 4)$.
\end{enumerate}
\finenonce{002736}



\finexercice


\exercice{2737, tumpach, 2009/10/25}
\enonce{002737}{}
Un cycliste s'entra\^ine chaque dimanche en faisant l'aller-retour d'Issy \`a Labat. Le trajet Issy-Labat n'est pas horizontal~: il y a des mont\'ees, des descentes et du plat. En mont\'ee, notre cycliste fait du quinze kilom\`etres \`a l'heure, en plat du vingt, en descente du trente. L'aller lui prend deux heures et le retour trois. Sur la portion du trajet qui n'est pas plate, la pente moyenne est de cinq pour cent.
\begin{enumerate}
\item Quelle est la distance d'Issy \`a Labat, quelle est la plus haute de ces deux villes, et quelle est leur diff\'erence d'altitude~?
\item Un autre cycliste, plus sportif, fait du vingt kilom\`etres \`a l'heure en mont\'ee, trente en plat et quarante en descente. Sachant que l'aller-retour Issy-Labat lui prend seulement trois heures quarante, d\'eterminer les trois longueurs~: de la partie du trajet qui monte, de celle qui descend, de celle qui est \`a plat.
\end{enumerate}
\finenonce{002737}



\finexercice


\exercice{2738, tumpach, 2009/10/25}
\enonce{002738}{}
Soient $a$, $b$, et $c$ trois nombres r\'eels. 
\begin{enumerate}
\item Quelle relation doivent satisfaire les param\`etres $a$, $b$ et $c$ pour que le syst\`eme suivant ait au moins une solution~?
$$
\mathcal{S}_{abc}~: \left\{\begin{array}{ccccccc} 
x & + & 2y & - & 3z & = & a\\
2x & + & 6y & - & 11z & = & b\\
x & - & 2y & + & 7z & = & c
\end{array}\right.
$$
\item Est-ce que le syst\`eme $\mathcal{S}_{abc}$ peut avoir une unique solution~?
\end{enumerate}
\finenonce{002738}



\finexercice


\exercice{2739, tumpach, 2009/10/25}
\enonce{002739}{}
R\'esoudre, suivant les valeurs de $m$~:
$$\mathcal{S}_1~:\; \left\{\begin{array}{ccccl}
    x & + & (m+1)y &=& m+2 \\
    mx & + & (m+4)y &=&3 
\end{array}\right.$$
$$\mathcal{S}_2~:\; \left\{\begin{array}{rcccl}
    mx &  + &(m-1)y & =& m+2 \\
    (m+1)x & - &my &=&5m+3
\end{array}\right.$$
\finenonce{002739}



\finexercice


\exercice{2768, tumpach, 2009/10/25}
\video{NxowOl4QnIA}
\enonce{002768}{}
\begin{enumerate}
\item Résoudre de quatre manières différentes le système suivant 
(par substitution, par la méthode du pivot de Gauss, 
en inversant la matrice des coefficients, par la formule de Cramer)~:
$$
\left\{\begin{array}{ccccc}2x & + & y & = & 1\\ 3x & + & 7y & = & -2  \end{array}\right.
$$
\item Choisir la méthode qui vous para\^it la plus rapide pour résoudre, 
selon les valeurs de $a$,  les systèmes suivants~:
$$
\left\{\begin{array}{ccccc} ax & + & y & = & 2\\ (a^2 + 1)x & + & 2ay & = & 1 \end{array} \right.
\qquad 
\left\{\begin{array}{ccccc} (a + 1)x & + & (a - 1)y & = & 1\\ (a - 1)x & + & (a + 1)y & = & 1 \end{array} \right.
$$
\end{enumerate}
\finenonce{002768}



\finexercice


\exercice{2769, tumpach, 2009/10/25}
\enonce{002769}{}
R\'esoudre le syst\`eme suivant de $5$ \'equations \`a $6$ inconnues~:
$$
\left\{\begin{array}{cccccccccccl} 
2x & + & y & + & z & - & 2u & + & 3v & - &w & = 1\\
3x & + & 2y & + & 2z & - & 3u & + & 5v & - & 3w & = 4\\
2x & + &2y & + & 2z & -&2u & +& 4v &- &4w& =  6\\
x & + & y &+& z & - & u &+& 2v & -& 2w & = 3\\
3x & & && & - &3u & + & 3v & + & 3w & = -6
\end{array}\right.
$$
\finenonce{002769}



\finexercice


\exercice{3408, quercia, 2010/03/10}
\enonce{003408}{Système avec paramètre}

\'Etudier l'existence de solutions du système :
 $$\left\{\begin{array}{ccccccc} x      &+&  y &+& (1-m)z  &=& m+2  \cr
           (1+m)x  &-&  y &+& 2z      &=& 0    \cr
           2x      &-& my &+& 3z      &=& m+2. \cr\end{array}\right.$$
\finenonce{003408}



\finexercice


\exercice{3409, quercia, 2010/03/10}
\enonce{003409}{Système avec paramètre}

\'Etudier l'existence de solutions du système :

 $$\left\{\begin{array}{ccccccc} x &-& my   &+& m^2z &=& m   \cr
           mx &-& m^2y &+& mz   &=& 1   \cr
           mx &+& y    &-& m^3z &=& -1. \cr\end{array}\right.$$
\finenonce{003409}



\finexercice


\exercice{3410, quercia, 2010/03/10}
\enonce{003410}{Système avec paramètre}

\'Etudier l'existence de solutions du système :
$$\left\{\begin{array}{ccccccc} x     &+& my &+& z      &=& 1  \cr
           (m+1)x &+& 2y &+& (m-3)z &=& -1 \cr
           (m-1)x & &    &-& 3z     &=& -1.\cr \end{array}\right.$$
\finenonce{003410}



\finexercice


\exercice{3411, quercia, 2010/03/10}
\enonce{003411}{Système avec paramètre}

\'Etudier l'existence de solutions du système :
$$\left\{\begin{array}{ccccccc} 3mx    &+& (3m-7)y &+& (m-5)z  &=& m-1  \cr
                 (2m-1)x &+& (4m-1)y &+& 2mz     &=& m+1  \cr
                 4mx     &+& (5m-7)y &+& (2m-5)z &=& 0.   \cr\end{array}\right.$$
\finenonce{003411}



\finexercice


\exercice{3412, quercia, 2010/03/10}
\enonce{003412}{Système avec paramètre}

\'Etudier l'existence de solutions du système :
$$\left\{\begin{array}{ccccccc}      x &+&       y &+&       z &= &1        \cr
                      ax &+&      by &+&      cz &= &d        \cr
                 a(a-1)x &+& b(b-1)y &+& c(c-1)z &= &d(d-1).  \cr\end{array}\right.$$
\finenonce{003412}



\finexercice


\exercice{3413, quercia, 2010/03/10}
\enonce{003413}{Système avec paramètre}

\'Etudier l'existence de solutions du système :
 $$\left\{\begin{array}{ccccccccc}  x &+& y &+& z &+& t &=& a \cr
             x &-& y &-& z &+& t &=& b \cr
            -x &-& y &+& z &+& t &=& c \cr
           -3x &+& y &-&3z &-&7t &=& d.\cr\end{array}\right.$$
\finenonce{003413}



\finexercice


\exercice{3414, quercia, 2010/03/10}
\enonce{003414}{Système avec paramètre}

\'Etudier l'existence de solutions du système :
$$\left\{\begin{array}{ccccccccc} 4x &+&3y &+&2z &+& t &=& a \cr
             x &+&4y &+&3z &+&2t &=& b \cr
            2x &+& y &+&4z &+&3t &=& c \cr
            3x &+&2y &+& z &+&4t &=& d.\cr\end{array}\right.$$
\finenonce{003414}



\finexercice


\exercice{3415, quercia, 2010/03/10}
\enonce{003415}{Système avec paramètre}

\'Etudier l'existence de solutions du système :
 $$\left\{\begin{array}{ccccccc} x\cos2\alpha &+& y\cos\alpha &+& z &=& a \cr
            x\cos2\beta  &+& y\cos\beta  &+& z &=& b \cr
            x\cos2\gamma &+& y\cos\gamma &+& z &=& c.\cr\end{array}\right.$$
\finenonce{003415}



\finexercice


\exercice{3416, quercia, 2010/03/10}
\enonce{003416}{Système avec paramètre}

\'Etudier l'existence de solutions du système :
$$\left\{\begin{array}{l}ax-by=p\cr by-cz=q\cr cz-ax=r.\cr\end{array}\right.$$
\finenonce{003416}



\finexercice


\exercice{3417, quercia, 2010/03/10}
\video{V12UnKijqcg}
\enonce{003417}{}

\'Etudier l'existence de solutions du système :
 $$\left\{\begin{array}{ccccccc} ax &+&  by &+&  z &=& 1 \cr
             x &+& aby &+&  z &=& b \cr
             x &+&  by &+& az &=& 1.\cr\end{array}\right.$$
\finenonce{003417}



\finexercice


\exercice{3418, quercia, 2010/03/10}
\enonce{003418}{Système avec paramètre}

\'Etudier l'existence de solutions du système :
 $$\left\{\begin{array}{ccccccc}
\frac x{1+a} &+& \frac y{1+2a} &+& \frac z{1+3a} &=& 1 \cr 
\frac x{2+a} &+& \frac y{2+2a} &+& \frac z{2+3a} &=& 1 \cr
\frac x{3+a} &+& \frac y{3+2a} &+& \frac z{3+3a} &=& 1.\cr\end{array}\right.$$
\finenonce{003418}



\finexercice


\exercice{3419, quercia, 2010/03/10}
\enonce{003419}{Système incompatible}

Soit $(S) \iff AX = B$ un système linéaire incompatible. Montrer que les lignes
de $A$ sont liées.

\finenonce{003419}



\finexercice


\exercice{3420, quercia, 2010/03/10}
\enonce{003420}{Combinaison de formes linéaires}

Soient $f,f_1,\dots,f_p$ des formes linéaires sur $ K^n$ linéairement
indépendantes.
Montrer que $f$ est combinaison linéaire de $f_1, \dots, f_p$ si et seulement si
$\mathrm{Ker} f \supset \mathrm{Ker} f_1 \cap \dots \cap \mathrm{Ker} f_p$.

Indication :
\'Etudier le système $$ \left\{\begin{array}{lll} f_1(x) &=& 0 \cr
                               \quad\vdots  \cr
                               f_p(x) &=& 0 \cr
                               f(x)   &=& 1.\cr \end{array}\right.$$


Le résultat est-il encore vrai si on ne suppose pas $(f_1, \dots, f_p)$
libre ?
\finenonce{003420}



\finexercice


\exercice{3421, quercia, 2010/03/10}
\enonce{003421}{Base antiduale}

Soient $f_1, \dots, f_n$, $n$ formes linéaires indépendantes sur un ev $E$
de dimension $n$.
Montrer qu'il existe une base $(\vec e_i)$ de $E$ telle que
$f_i = {\vec e_i}^{\;*}$.
\finenonce{003421}



\finexercice


\exercice{3422, quercia, 2010/03/10}
\enonce{003422}{Orthogonal d'un sev}

Soit $E$ un $ K$-ev de dimension $n$ et $F$ un sev de $E^*$ dimension $p$.

On note $F^\bot = \{ \vec x \in E \text{ tq }
                     \forall\ f \in F \text{ on a } f(\vec x) = 0 \}$.
Chercher $\dim F^\bot$.
\finenonce{003422}



\finexercice


\exercice{3423, quercia, 2010/03/10}
\enonce{003423}{Système de Vandermonde}

Soient $\alpha_1, \dots, \alpha_n$, $n$ scalaires distincts et
$M$ la matrice $\big(\alpha_i^{j-1}\big)$ (matrice de Vandermonde).


Montrer que $M$ est inversible en interprétant le système $MX = 0$ dans
$ K_{n-1}[x]$.
\finenonce{003423}



\finexercice


\exercice{3424, quercia, 2010/03/10}
\video{wKgDXlqxneE}
\enonce{003424}{Formule d'intégration numérique}

Trouver trois réels $\alpha, \beta, \gamma$ tels que pour tout polynôme
de degré $\le 3$ on ait :
$$  \int_{2}^4 P(x)\; dx = \alpha P(2) + \beta P(3) + \gamma P(4).$$
\finenonce{003424}



\finexercice



\exercice{3425, quercia, 2010/03/10}
\enonce{003425}{Système non linéaire}

Résoudre le système : $$\left\{
\begin{array} {lll}
x^3y^2z^6    &=& 1 \cr
                              x^4y^5z^{12} &=& 2 \cr
                              x^2y^2z^5    &=& 3.\cr 
                              \end{array}\right.$$

\begin{enumerate}
  \item Lorsque $x,y,z$ sont réels strictement positifs.
  \item Lorsque $x,y,z \in \C$.



\end{enumerate}
\finenonce{003425}



\finexercice

\exercice{3431, quercia, 2010/03/10}
\enonce{003431}{Comatrice, Ensi P 91}

    Soit $A \in \mathcal{M}_n(\R)$. \'Etudier le rang de $\text{com}(A)$ en fonction du
    rang de $A$.
    
\finenonce{003431}



\finexercice


\exercice{3432, quercia, 2010/03/10}
\enonce{003432}{Comatrice}

Soit $n\ge 2$ et $A \in \mathcal{M}_n(K)$.

\begin{enumerate}
  \item  Calculer $\text{com}\,(\text{com}\, A)$ dans le cas où $A$ est inversible.
  \item  Si $\mathrm{rg} A \le n-2$, démontrer que $\text{com}\, A = 0$.
  \item  Si $\mathrm{rg} A = n-1$, démontrer que $\mathrm{rg}(\text{com}\, A)= 1$.
  \item  Dans le cas général, démontrer que $\text{com}\,(\text{com}\, A) = (\det A)^{n-2}A$.
\end{enumerate}
\finenonce{003432}



\finexercice

\exercice{3461, quercia, 2010/03/10}
\enonce{003461}{Calcul de rang}

Chercher les rangs des matrices suivantes :

\begin{enumerate}
  \item $\begin{pmatrix} 1 &2  &-4 &-2 &-1 \cr
                0 &-2 &4  &2  &0  \cr
                1 &1  &-2 &-1 &1  \cr \end{pmatrix}$.

  \item $\begin{pmatrix} 0  &-1 &2  &-2 \cr
                -7 &-7 &2  &-8 \cr
                0  &4  &-6 &6  \cr
                2  &-2 &0  &-2 \cr \end{pmatrix}$.
		
  \item $\begin{pmatrix} 1  &7  &2  &5  \cr
                -2 &1  &1  &5  \cr
                -1 &2  &1  &4  \cr
                1  &4  &1  &2  \cr \end{pmatrix}$.

  \item $\begin{pmatrix} 1  &4  &-1 &2  &4  \cr
                2  &0  &-3 &-1 &7  \cr
                -2 &3  &2  &1  &4  \cr \end{pmatrix}$.
\end{enumerate}
\finenonce{003461}



\finexercice


\exercice{3462, quercia, 2010/03/10}
\enonce{003462}{Calcul de rang}

Chercher
$\mathrm{rg}\begin{pmatrix}1 &-1  &  0 &  1  \cr
              m &  1 &-1  &-1   \cr
              1 &-m  &  1 &  0  \cr
              1 &-1  &   m&  2  \cr\end{pmatrix}$ en fonction de $m\in\C$.
\finenonce{003462}



\finexercice


\exercice{3463, quercia, 2010/03/10}
\enonce{003463}{Calcul de rang}

Chercher
$\mathrm{rg}\begin{pmatrix} \lambda &1 &\phantom-5 \cr  -1 &4 &\lambda \cr 3 & -1 &5 \end{pmatrix}$ et le
cas échéant, donner une relation de dépendance linéaire entre les lignes.

\finenonce{003463}



\finexercice


\exercice{3464, quercia, 2010/03/10}
\enonce{003464}{Calcul de rang}

Chercher
$\mathrm{rg}\begin{pmatrix} 2&\phantom-4&-5&-7 \cr -1&3&1&2 \cr 1&a&-2&b \end{pmatrix}$ en fonction de
$a$ et $b$.

\finenonce{003464}



\finexercice


\exercice{3465, quercia, 2010/03/10}
\enonce{003465}{Matrice à trou}

Soient
$A \in \mathcal{M}_{3,2}(K),\ B \in \mathcal{M}_{2,2}(K),\ C \in \mathcal{M}_{2,3}(K)$ telles que
$ABC = \begin{pmatrix} 1&1&\phantom-2 \cr -2&x&1 \cr 1&-2&1 \end{pmatrix}$.
Trouver $x$.


\finenonce{003465}



\finexercice


\exercice{3466, quercia, 2010/03/10}
\enonce{003466}{Factorisation, Centrale P' 1996}

Soit la matrice carrée d'ordre~$n$, $I_p$ ($p\le n$), telle que le $i$-ème
terme diagonal vaut 1 si $i$ est compris entre $p$ et~$n$, tous les autres
coefficients étant nuls.
Quelle sont les conditions sur~$A$ (matrice carrée d'ordre~$n$) pour
qu'il existe~$B$ telle que $AB=I_p$~?
\finenonce{003466}



\finexercice


\exercice{3467, quercia, 2010/03/10}
\enonce{003467}{\'Echange de lignes}

Soit $A \in \mathcal{M}_n(K)$ inversible et $B$ la matrice obtenue en échangeant dans $A$
les colonnes $i$ et $j$. Montrer que $B$ est aussi inversible. Comment passe-t-on de
$A^{-1}$ à $B^{-1}$ ?

\finenonce{003467}



\finexercice


\exercice{3468, quercia, 2010/03/10}
\enonce{003468}{Matrices de rang 1}

Soit $M \in \mathcal{M}_n(K)$. Montrer que :
\smallskip
\leavevmode\hfil
$\mathrm{rg}(M) = 1 \iff$ il existe $C$, colonne et $L$, ligne,
non nulles, telles que $M = CL$.
\hfil
\smallskip
\finenonce{003468}



\finexercice


\exercice{3469, quercia, 2010/03/10}
\enonce{003469}{Matrices de projection de rang 1}

Soit $A \in \mathcal{M}_n(K)$ de rang 1. Montrer que $A$ est une matrice de projection si et
seulement si $\mathrm{tr} A = 1$.
\finenonce{003469}



\finexercice


\exercice{3470, quercia, 2010/03/10}
\enonce{003470}{Calcul de rang}

Soient $x_1, \dots, x_p, y_1, \dots y_p \in  K$. On note $A,B$ les matrices
de termes généraux $x_i + y_j$ et $(x_i+y_j)^2$.

Chercher les rangs de A et~B.
\finenonce{003470}



\finexercice


\exercice{3471, quercia, 2010/03/10}
\enonce{003471}{Juxtaposition de matrices}

Soient $A \in \mathcal{M}_{n,p}(K)$ et $B \in \mathcal{M}_{n,q}(K)$. On considère
$C = (A\ B) \in \mathcal{M}_{n,p+q}(K)$.

Montrer que : $\mathrm{rg}(C) = \mathrm{rg}(A) \Leftrightarrow \exists\ P \in \mathcal{M}_{p,q}(K)$ tq $B = AP$.
\finenonce{003471}



\finexercice


\exercice{3472, quercia, 2010/03/10}
\enonce{003472}{$\mathrm{rg}(A) = \mathrm{rg}({}^t\!AA)$}

Soit $A \in \mathcal{M}_n(\R)$, et $B = {}^t\!AA$.

\begin{enumerate}
  \item Montrer que : $\forall\ Y \in \mathcal{M}_{n,1}(\R)$, ${}^tYY = 0 \Leftrightarrow Y = 0$.
  \item Montrer que : $\forall\ X \in \mathcal{M}_{n,1}(\R)$, $BX = 0 \Leftrightarrow AX = 0$.
  \item En déduire que $\mathrm{rg}(A) = \mathrm{rg}(B)$.
  \item Trouver une matrice $A \in \mathcal{M}_2(\C)$ telle que $\mathrm{rg}(A) \ne \mathrm{rg}({}^t\!AA)$.
    

\end{enumerate}
\finenonce{003472}



\finexercice


\exercice{3473, quercia, 2010/03/10}
\enonce{003473}{Rang de $\Re(M)$}

Soit $M \in \mathcal{M}_n(\C)$ de rang 1. On écrit $M = P + iQ$ avec $P,Q \in \mathcal{M}_n(\R)$.
Montrer que rg$(P) \le 2$.
\finenonce{003473}



\finexercice


\exercice{3474, quercia, 2010/03/10}
\enonce{003474}{Calcul de rang}

Soit $M = \Bigl(\cos (i+j-1)\theta \Bigr) \in \mathcal{M}_n(\R)$.
Déterminer $\mathrm{rg} M$ en fonction de $\theta$.
\finenonce{003474}



\finexercice


\exercice{3475, quercia, 2010/03/10}
\enonce{003475}{Décomposition en blocs}

Soit $M = \begin{pmatrix} A & B \cr C & D \cr \end{pmatrix}$
une matrice carrée décomposée en blocs. On suppose que $A$ est inversible.

Montrer que $\mathrm{rg} M = \mathrm{rg} A + \mathrm{rg} (D - CA^{-1}B)$.
\finenonce{003475}



\finexercice


\exercice{3476, quercia, 2010/03/10}
\enonce{003476}{$MA=0$, Chimie P' 90}
    Soit $A \in \mathcal{M}_n(\R)$ et $E = \{ M \in \mathcal{M}_n(\R)$ tq $MA = 0 \}$.
    Quelle est la structure de $E$, sa dimension~?
\finenonce{003476}



\finexercice


\exercice{3477, quercia, 2010/03/10}
\enonce{003477}{Rang des applications $X  \mapsto AX, XB, AXB$}

Soient $A \in \mathcal{M}_{n,p}(K)$ et $B \in \mathcal{M}_{q,r}(K)$. Chercher le rang des
applications :
$$
   f : {\mathcal{M}_{p,q}(K)} \to {\mathcal{M}_{n,q}(K)}, X \mapsto {AX}\qquad
   g : {\mathcal{M}_{p,q}(K)} \to {\mathcal{M}_{p,r}(K)} X \mapsto{XB}  $$
$$  h  : {\mathcal{M}_{p,q}(K)} \to {\mathcal{M}_{n,r}(K)} X \mapsto{AXB}
$$
(on transformera $A$ et $B$ en matrices canoniques équivalentes)

\finenonce{003477}



\finexercice


\exercice{3478, quercia, 2010/03/10}
\enonce{003478}{Rang de  $X  \mapsto AX - XA$}


Soit $A \in \mathcal{M}_{2}(K)$. Chercher le rang de l'application
${\mathcal{M}_{2}(K)}\to {\mathcal{M}_{2}(K)}, M \mapsto{AM-MA}$
\finenonce{003478}



\finexercice


\exercice{3479, quercia, 2010/03/10}
\enonce{003479}{Matrice antisymétrique $3\times3$}

Soit $M \in \mathcal{M}_3(K)$ antisymétrique. Quel est le rang de $M$ ?
\finenonce{003479}



\finexercice


\exercice{3480, quercia, 2010/03/10}
\enonce{003480}{Matrices antisymétriques}

Soit $A=(a_{ij}) \in \mathcal{M}_n(K)$ antisymétrique.

\begin{enumerate}
  \item On suppose $a_{12} \ne 0$, et on décompose $A$ sous la forme :
    $A = \begin{pmatrix} J &U \cr {-{}^tU} &V\cr\end{pmatrix}$
    avec $J = \begin{pmatrix} 0 &a_{12} \cr -a_{12} &0 \cr\end{pmatrix}$.\par
    Soit $P = \begin{pmatrix} I_2 &-J^{-1}U \cr 0 &I_{n-2}\cr\end{pmatrix}$.
   \begin{enumerate}
    \item Montrer que $P$ existe et est inversible.
    \item Calculer $AP$.
    \item En déduire que $\mathrm{rg}(A) = 2 + \mathrm{rg}({}^tUJ^{-1}U + V)$.
   \end{enumerate}
  \item Dans le cas général, montrer que $\mathrm{rg}(A)$ est pair.
\end{enumerate}
\finenonce{003480}



\finexercice


\exercice{3481, quercia, 2010/03/10}
\enonce{003481}{Matrice à diagonale dominante}

Soit $M=(a_{ij}) \in \mathcal{M}_n(\C)$. On dit que $M$ est {\it à diagonale dominante\/}
si : $\forall\ i,\ |a_{ii}| > \sum_{j\ne i}|a_{ij}|$.

\begin{enumerate}
  \item On transforme $M$ en
    $M' = \begin{pmatrix} a_{11} &\ a_{12}&\dots &a_{1n} \cr
                     0                         \cr
                   \vdots &        &M_1        \cr
                     0                         \cr \end{pmatrix}$
    par la méthode du pivot. Montrer que $M_1$ est à diagonale dominante.
  \item En déduire que $M$ est inversible.
\end{enumerate}
\finenonce{003481}



\finexercice


\exercice{3482, quercia, 2010/03/10}
\enonce{003482}{Rang par blocs, Matexo}

Soit une matrice $M$ de rang $r$ telle que~:
$$M=\begin{pmatrix} M_r&M_2\cr M_1&M_3\cr\end{pmatrix},$$
où la matrice $M_r$ est carrée de rang $r$ et de taille $r$.
Montrer que $M_3=M_1M_r^{-1}M_2$.
\finenonce{003482}



\finexercice


\exercice{3483, quercia, 2010/03/10}
\enonce{003483}{$\mathrm{rg}(BC)$, Centrale MP 2006}
\begin{enumerate}
  \item Soient deux matrices $B \in \mathcal{M}_{n,r}(\R)$ et $C \in \mathcal{M}_{r,n}(\R)$ de même rang $r$.
    Montrer que $A = BC$ est de rang $r$.
    

  \item Réciproquement, soit $A \in \mathcal{M}_n(\R)$ de rang $r \ge 1$.
    Montrer qu'il existe des matrices $B$ et $C$ comme
    précédemment telles que $A = BC$.
    

  \item Déterminer explicitement une telle décomposition pour
    $A = \begin{pmatrix}1 & 2 & 3 & 4 \cr 2 & 3 & 4 & 5 \cr
    3 & 4 & 5 & 6 \cr 4 & 5 & 6 & a\cr\end{pmatrix}$, $a \in \R$.

  \item Supposons de plus $A$ symétrique. Montrer que $CB$ est aussi de rang $r$.
    
\end{enumerate}
\finenonce{003483}



\finexercice


\exercice{3484, quercia, 2010/03/10}
\enonce{003484}{$PA$ nilpotente}
Soit $A\in\mathcal{M}_n(K)$. Montrer que $A$ est non inversible si et seulement s'il
existe $P\in GL_n( K)$ telle que $PA$ est nilpotente.
\finenonce{003484}



\finexercice


\exercice{5367, rouget, 2010/07/06}
\enonce{005367}{***}
Résoudre le système $MX=U$ où $M=\left(
\begin{array}{ccccc}
1&1&\ldots&\ldots&1\\
1&2&\ldots&\ldots&n\\
1&2^2&\ldots&\ldots&n^2\\
\vdots&\vdots& & &\vdots\\
1&2^{n-1}&\ldots&\ldots&n^{n-1}
\end{array}
\right)$.
\finenonce{005367}


\finexercice


\exercice{5375, rouget, 2010/07/06}
\enonce{005375}{}
Résoudre (en discutant en fonction des différents paramètres) les systèmes suivants~:

$$
\begin{array}{ccc}
1)\;\left\{
\begin{array}{l}
2x+3y+z=4\\
-x+my+2z=5\\
7x+3y+(m-5)z=7
\end{array}
\right.&
2)\;\left\{
\begin{array}{l}
2x+my+z=3m\\
x-(2m+1)y+2z=4\\
5x-y+4z=3m-2
\end{array}
\right.&
3)\;\left\{
\begin{array}{l}
x+y+z+t=3\\
x+my+z-mt=m+2\\
mx-y-mz-t=-1
\end{array}
\right.\\
4)\;\left\{
\begin{array}{l}
x+2y+3z+mt=m-1\\
2x+y+mz+3t=1\\
3x+my+z+2t=0\\
mx+3y+2z+t=0
\end{array}
\right.&5)\;\left\{
\begin{array}{l}
mx+y+z=m+2\\
-x-y+mz=m-2\\
-mx+y+mz=-m\\
x-y-mz=m-4
\end{array}
\right.&6)\;
\left\{
\begin{array}{l}
x+y+z=1\\
ax+by+cz=m\\
\frac{x}{a}+\frac{y}{b}+\frac{z}{c}=\frac{1}{m}
\end{array}
\right.\\
7)\;
\left\{
\begin{array}{l}
(b+c)^2x+b^2y+c^2z=1\\
a^2x+(c+a)^2y+c^2z=1\\
a^2x+b^2y+(a+b)^2z=1
\end{array}
\right.&8)\;\left\{
\begin{array}{l}
ax+by+cz=p\\
cx+ay+bz=q\\
bx+cy+az=r
\end{array}
\right.& \\
9)\;\left\{
\begin{array}{l}
x+y+z=0\\
ax+by+cz=2\\
a^2x+b^2y+c^2z=3
\end{array}
\right.&(\mbox{où}\;a,\;b,\;\mbox{et}\;c\;\mbox{sont les racines de l'équation}&t^3-t+1=0).
\end{array}$$
\finenonce{005375}


\finexercice


\exercice{5378, rouget, 2010/07/06}
\enonce{005378}{}
Résoudre le système~:~$x_1+x_2=0$, $x_{k-1}+x_k+x_{k+1}=0$ pour $k=2,...,n-1$, $x_{n-1}+x_n=0$.
\finenonce{005378}


\finexercice


\exercice{5381, rouget, 2010/07/06}
\enonce{005381}{}
Soient $a_1$,..., $a_n$, $b_1$,..., $b_n$ $2n$ nombres complexes deux à deux distincts tels que les sommes $a_i+b_j$ soient toutes non nulles. Résoudre le système $\sum_{j=1}^{n}\frac{x_j}{a_i+b_j}=1$, pour tout $i=1,...,n$ (en utilisant la décomposition en éléments simples de $R=\sum_{j=1}^{n}\frac{x_j}{X+b_j}$).
\finenonce{005381}


\finexercice

\exercice{5639, rouget, 2010/10/16}
\enonce{005639}{**}
Résoudre le système $MX = U$ où $M =(j^{i-1})_{1\leqslant i,j\leqslant n}\in M_n(\Rr)$, $U=(\delta_{i,1})_{1\leqslant i\leqslant n}\in M_{n,1}(\Rr)$ et $X$ est un vecteur colonne inconnu.
\finenonce{005639}


\finexercice


\exercice{7026, megy, 2016/10/28}

\enonce{007026}{}
% source : une partie de ces systèmes provient de Lebossé & Hémery, classe de troisième (!), 1964.
Résoudre les systèmes suivants, en précisant auparavant le domaine de résolution.

\[
\begin{cases}
\frac{x+2}{x}=\frac{y+1}{y-2}\\
\frac{5x+1}{5x-2}=\frac{y-1}{y-2}
\end{cases}
\quad
\begin{cases}
\frac{(x+3)^2+(y-1)^2}{x^2+y^2}=1\\
3x+2y=73
\end{cases}
\quad 
\begin{cases}
\frac{(x-1)^2-(x-5)^2}{(y+1)^2-(y-1)^2}=1\\
2y-x=45
\end{cases}
\]

\[
\begin{cases}
\frac{5}{x}+\frac{4}{y}=18\\
\frac{3}{x}-\frac{7}{y}=-55
\end{cases}
\quad
\begin{cases}
\frac{1}{x}+\frac{1}{y}=3\\
\frac{2}{x}+\frac{3}{y}=9
\end{cases}
\quad 
\begin{cases}
\frac{1-2x}{x}+\frac{3}{y}=0\\
\frac{1-x}{x}+\frac{3-y}{y}=0
\end{cases}
\]

\[ \begin{cases}
\frac{3}{4(x-2)}+\frac{7}{3(y-1)}=41\\
\frac{5}{2(x-2)}-\frac{3}{5(y-1)}=11
\end{cases}
\quad
\begin{cases}
\frac{9}{x-y}+\frac{10}{x+y}=75\\
\frac{12}{x-y}+\frac{25}{x+y}=135
\end{cases}
\]
\finenonce{007026}



\finexercice


\exercice{7027, megy, 2016/10/28}

\enonce{007027}{}
La somme de deux nombres $x$ et $y$ est $29$. La différence de leurs carrés est $145$. Quels sont ces nombres ?
\finenonce{007027}



\finexercice


\exercice{7028, megy, 2016/10/28}

\enonce{007028}{}
Trouver les dimension d'un triangle rectangle d'hypoténuse 13 cm et d'aire 30 cm${}^2$.
\finenonce{007028}



\finexercice


\exercice{7029, megy, 2016/10/28}

\enonce{007029}{}
Dans une station de métro, les usagers ont à leur disposition un tapis roulant de 300m de long.
Un piéton marchant à vitesse constante fait l'aller-retour sur le tapis roulant. À l'aller, il met 1 minute et
30 secondes. Au retour, à contre-sens, il met 4 minutes et 30 secondes.
Déterminer la vitesse du piéton et celle du tapis roulant.
\finenonce{007029}



\finexercice


\exercice{7030, megy, 2016/10/28}
\enonce{007030}{}
 Résoudre le système suivant sur $\R_+^*$ puis sur le domaine de définition maximal:
\[
\begin{cases}
xy=1\\
x/y=2\\
\end{cases}
\]
\finenonce{007030}



\finexercice


\exercice{7031, megy, 2016/10/28}

\enonce{007031}{}
 Résoudre dans $\R_+^*$ le système suivant:
\[
\begin{cases}
xyz=1\\
xy^2z^4=2\\
xy^3z^9=3
\end{cases}
\]
\finenonce{007031}



\finexercice


\exercice{7032, megy, 2016/10/28}

\enonce{007032}{}
% vu dans un cours de Term ES

Soient $a, b, c$ trois réels et $f : \R\to \R, \: x\mapsto ae^x+bx+c$. On suppose que le graphe de $f$ contient le point de coordonnées $(0,1)$, et que sa tangente en ce point contient également le point de coordonnées $(2,3)$. On suppose enfin que le graphe admet une tangente horizontale au point d'abscisse $\ln(3)$. Déterminer $a, b, c$.
\finenonce{007032}



\finexercice


\exercice{7033, megy, 2016/10/28}

\enonce{007033}{}
% tag primaire : systèmes linéaires 200.03
% tags secondaires : géométrie plane, isométries planes
Soient $a_1 ,...,a_n $ des points du plan complexe.\\
Déterminer à quelle(s) condition(s) il existe au moins un polygone à $n$ sommets $z_1 ,...,z_n $ tel que~:\\
($a_i $ est le milieu de $\left[ {z_i ,z_{i + 1} } \right]$ et $a_n $ est le milieu de $\left[ {z_n ,z_1 } \right]$.)
\finenonce{007033}



\finexercice


\exercice{7034, megy, 2016/10/28}

\enonce{007034}{}
% tag primaire : systèmes linéaires
% tags secondaires : exponentielle, fonctions usuelles
Résoudre le système suivant:
\[
\begin{cases}
  8^x  = 10y\\
  2^x  = 5y
\end{cases}
\]
\finenonce{007034}



\finexercice
\section{200.04 Applications}


\exercice{1151, legall, 1998/09/01}

\enonce{001151}{}
 Soit $ E $ un espace vectoriel r\'eel de dimension
finie $ n $ et $ \varphi \in \mathcal{L} (E) $ telle que $ \varphi
^2 =-\hbox{id}_E .$
\begin{enumerate}
    \item Donner des exemples de telles applications dans le cas $ n=2 $ ou
$ 4 .$
    \item Montrer que de telles applications existent si et seulement si $ n
$ est pair.
\end{enumerate}

\finenonce{001151}



\finexercice



\exercice{1152, legall, 1998/09/01}

\enonce{001152}{}
Inverser les matrices
$ \begin{pmatrix}
 1 & -1 & 0 & 0 \cr 2 & 1 & 0 & 0 \cr
                    0 & 0 & 1 & 2 \cr 0 & 0 & 2 & 1 \cr \end{pmatrix} $
et $ \begin{pmatrix}
 1 & 0 & 1 & 0 \cr 0 & 1 & 0 & 1 \cr
                    1 & 0 & -1 & 0 \cr 0 & 1 & 0 & -1 \cr \end{pmatrix} $ ainsi que
leurs produits.
\finenonce{001152}



\finexercice



\exercice{3443, quercia, 2010/03/10}
\enonce{003443}{Résultant}

Soient $P = a_0 + a_1X + \dots + a_pX^p$, et $Q = b_0 + b_1X + \dots + b_qX^q$,
avec $a_p \ne 0,\ b_q \ne 0$.


\def\fleche{$\hskip2pt
            \mathord\leftarrow\mkern-6mu
            \cleaders\hbox{$\mkern-2mu\mathord-\mkern-2mu$}\hfill
            \mkern-6mu\mathord\rightarrow
            \hskip2pt$}

Le résultant de $P$ et $Q$ est :
$\text{Res}(P,Q) = \left|\vcenter{
  \halign{&\hskip3pt$#$\hfil\hskip3pt\cr
  a_0    &      &      &       &b_0              &                        \cr
  a_1    &\ddots&      &       &\vdots           &\ddots                  \cr
  \vdots &\ddots&\ddots&       &\vdots           &\ddots&b_0              \cr
  a_p    &\ddots&\ddots&a_0    &b_{q-1}\hidewidth&\ddots&\vdots           \cr
         &\ddots&\ddots&a_1    &b_q              &\ddots&\vdots           \cr
         &      &\ddots&\vdots &                 &\ddots&b_{q-1}\hidewidth\cr
         &      &      &a_p    &                 &      &b_q              \cr
  \multispan4\fleche &\multispan3\fleche \cr
  \noalign{\vskip-0.5\baselineskip}
  \multispan4 $q$    &\multispan3 $p$    \cr}
  \vskip-\baselineskip}\hskip5pt\right|$\quad
  $\vcenter{\hbox{les positions non remplies}
            \hbox{correspondent à des zéros.}}$\vspace{0.2cm}\\




En considérant l'application
${\Phi} : { K_{q-1}[X] \times  K_{p-1}[X]} \to { K_{p+q-1}[X]},
 {(U,V)} \mapsto {UP+VQ,}$

montrer que : $\text{Res}(P,Q) \ne 0 \iff P \wedge Q = 1$.


Application : CNS pour que le polynôme $P = X^4 + aX + b$ ait une racine
multiple ?

\finenonce{003443}



\finexercice


\exercice{3444, quercia, 2010/03/10}
\enonce{003444}{Système unisolvent}

Soient $E$ un ensemble quelconque et
${f_1, \dots, f_n} : E \to K$ des fonctions.

Montrer par récurrence sur $n$ que la famille $(f_1, \dots, f_n)$ est libre
dans $ K^E$ si et seulement s'il existe des éléments $x_1, \dots, x_n$ de $E$
tels que $\det\bigl( f_i(x_j) \bigr) \ne 0$.
\finenonce{003444}



\finexercice

\exercice{3445, quercia, 2010/03/10}
\enonce{003445}{$\prod a_{i\sigma(i)} =$ cste}

Soit $A = (a_{ij}) \in \mathcal{M}_n(K)$ telle qu'il existe $a \ne 0$ tel que :
$\forall\ \sigma \in S_n,\ \prod_{i=1}^n a_{i\sigma(i)} = a$. Montrer que $\mathrm{rg} A = 1$.
\finenonce{003445}



\finexercice


\exercice{3446, quercia, 2010/03/10}
\enonce{003446}{Combinaison linéaire des solutions}

Soit $(S) \iff AX = B$ un système linéaire de $n$ équations à $n$ inconnues
de Cramer.
Montrer que pour tous scalaires $c_1,\dots,c_n$, on a :
$$c_1x_1 + \dots + c_nx_n =
  -\frac1{\det A}\left|\,
                  \vcenter{\offinterlineskip
                           \halign{\strut\hfil$#$\hfil&&\quad\hfil$#$\hfil\cr
                                &      &\hfill\vrule             &b_1     \cr
                                &A     &\hfill\vrule             &\vdots  \cr
                            \multispan3\hrulefill\vrule depth0pt &b_n     \cr
                            c_1 &\dots &c_n                      &0       \cr }}
                  \,\right|.$$
\finenonce{003446}



\finexercice


\exercice{3447, quercia, 2010/03/10}
\enonce{003447}{Problème d'interpolation de Lagrange}
Soit $A$ un {\it anneau\/} commutatif, $x_1,\dots,x_n\in A$.
Démontrer l'équivalence entre les propositions suivantes~:

{\bf a.} Le déterminant de Vandermonde de $x_1,\dots,x_n$ est un élément inversible de~$A$~;

{\bf b.} Pour tous $y_1,\dots,y_n\in A$, il existe un unique polynôme
        $P\in A_{n-1}[X]$ tel que $P(x_i)=y_i$ pour $i=1,\dots,n$.

Donner un exemple d'anneau $A$ et un problème d'interpolation dans $A$ (en des
points $x_i$ distincts) n'ayant pas de solution.

\finenonce{003447}



\finexercice


\exercice{3448, quercia, 2010/03/10}
\enonce{003448}{Polytechnique MP 2002}

Soit $p$ un nombre premier et $a_0,\dots ,a_{p-1}\in\Z$.
Montrer que le déterminant de la matrice 
$A=(a_{j-i\text{ mod }p})\in\mathcal{M}_p(\Z)$ vérifie~:
$\det (A)\equiv a_0+\dots +a_{p-1}$ mod $p$. 

{\it Indication: écrire $A=\sum_{k=0}^{p-1} a_kJ^k$ et calculer $A^p$.}

\finenonce{003448}



\finexercice


\exercice{3449, quercia, 2010/03/10}
\enonce{003449}{Centrale MP 2002}
Soit un déterminant symétrique réel d'ordre impair dont les coefficients
sont entiers, les diagonaux étant de plus pairs.
Montrer que ce déterminant est pair.

\finenonce{003449}



\finexercice


\exercice{3450, quercia, 2010/03/10}
\enonce{003450}{Formule de Cauchy-Binet}

Soit $M=(a_{ij})\in\mathcal{M}_{np}(K)$ et $q\in[[1,\min(n,p)]]$.

Pour $X=\{x_1,\dots,x_q\}$
et $Y=\{y_1,\dots,y_q\}$
avec $1\le x_1< x_2< \dots< x_q\le n$ et $1\le y_1< y_2< \dots< y_q\le p$
on note $\Delta_{X,Y}(M)$ le déterminant de la matrice $q\times q$ de terme général $a_{x_i,y_j}$.

\begin{enumerate}
  \item Soient $M\in\mathcal{M}_{np}(K)$ et $N\in\mathcal{M}_{pn}(K)$ avec $n\le p$.
    Montrer que $\det(MN) = \sum_{X\subset[[1,p]];\mathrm{Card}\, X =
    n}\Delta_{[[1,n]],X}(M)\Delta_{X,[[1,n]]}(N)$
    (considérer les deux membres comme des fonctions des colonnes de~$N$).
  \item Donner une formule pour $\det(MN)$ quand $n>p$.
  \item Soient $M\in\mathcal{M}_{np}(K)$, $N\in\mathcal{M}_{pq}(K)$ et $r\in{[[1,\min(n,q)]]}$.
    Montrer, pour $X\subset{[[1,n]]}$ et $Y\subset{[[1,q]]}$ avec $\mathrm{Card}\,(X) =
    \mathrm{Card}\,(Y) = r$~: $\Delta_{X,Y}(MN) = \sum_{Z\subset{[[1,p]]} ;\mathrm{Card}\,(Z)=r}\Delta_{X,Z}(M)\Delta_{Z,Y}(N)$.
\end{enumerate}
\finenonce{003450}



\finexercice


\exercice{5377, rouget, 2010/07/06}
\enonce{005377}{}
Dans le plan, on donne $n$ points $A_1$, ... , $A_n$. Existe-t-il $n$ points $M_1$,..., $M_n$ tels que $A_1$ soit le
milieu de $[M_1,M_2]$, $A_2$ soit le milieu de $[M_2,M_3]$,..., $A_{n-1}$ soit le milieu de $[M_{n-1},M_n]$ et $A_n$ soit le milieu de $[M_n,M_1]$.

\finenonce{005377}


\finexercice

\section{200.99 Autre}


\exercice{1153, liousse, 2003/10/01}

\enonce{001153}{}
 Montrer que $\left| 
\begin{array}{ccc}
1+a & a & a \\ 
b & 1+b & b \\ 
c & c & 1+c
\end{array}
\right| =1+a+b+c$ sans le d\'{e}velopper.
\finenonce{001153}



\finexercice



\exercice{1154, legall, 1998/09/01}

\enonce{001154}{}
Une matrice carr\'ee $ A=(a_{ij})_{i,j\in \{ 1,\ldots, n
\} }\in M_n({\Rr}) $ est dite triangulaire
 sup\'erieure lorsque pour tout $ i>j $~: $ a_{ij}=0 .$
\begin{enumerate}
    \item Montrer que le produit de deux matrices triangulaires sup\'erieures est
une matrice triangulaire sup\'erieure.
    \item D\'emontrer que $ \hbox{det}(A)=a_{11}\cdots a_{nn} .$
    \item Soit $ E $ un espace vectoriel, $ \epsilon =\{ e_1 , \ldots, e_n\}  $
une base de $ E $ et $ \varphi \in \mathcal{L} (E) .$ On note $ E_i $
l'espace vectoriel engendr\'e par $ \{ e_1, \ldots, e_i\}  $, pour tout $ 1\leq i\leq
n .$ Montrer que $ \hbox{Mat} (\varphi , \epsilon ) $ est triangulaire sup\'erieure
si et seulement si $ \varphi (E_i )\subset E_i  $ pour tout $ 1\leq i\leq n .$
    \item D\'emontrer que l'inverse d'une matrice triangulaire sup\'erieure est
une matrice triangulaire sup\'erieure.
\end{enumerate}
\finenonce{001154}



\finexercice



\exercice{1155, legall, 1998/09/01}

\enonce{001155}{}
On consid\`ere les matrices~:


$$ I=\begin{pmatrix} 1 & 0 & 0  & 0  \cr
                                   0 & 1 & 0 &  0  \cr
                                    0 & 0 & 1 & 0 \cr
                                     0 & 0 & 0 &  1   \cr \end{pmatrix}  \hskip10mm
N=
\begin{pmatrix} 0 & 3 & 1  & 3  \cr
                                   0 & 0 & 0 &  1  \cr
                                    0 & 0 & 0 & -1 \cr
                                     0 & 0 & 0 &  0   \cr \end{pmatrix} \hskip5mm
A=I+N  .$$
\begin{enumerate}
    \item Pour tout $ n \in \N ^* $ calculer $ \hbox{det}(A^n) .$
    \item Calculer $ N^2 $ et $ N^3 .$
    \item Pour tout $ n \in \N ^* $ donner le rang de $ N^n $ et celui de
$ A^n .$
    \item En utilisant 1., donner, en fonction de $ n \in \N ^* ,$
l'expression de la matrice $ M(n)=A^n .$
    \item Pour $ n \in \N ^* ,$ justifier la formule $ (A^n)^{-1}=M(-n)
.$ Expliquer et justifier l'\' ecriture~:
$ A^n=M(n) $ pour tout $ n \in \Z  .$
\end{enumerate}
\finenonce{001155}



\finexercice



\exercice{1156, legall, 1998/09/01}

\enonce{001156}{}
Soit $ S $ la matrice $ 5 \times 5  $ \`a
coefficients r\' eels : $ S=\begin{pmatrix}  0 & 0 & 1 & 0  & 0\cr
                                    1 & 0 & 0 & 0 & 0 \cr
                                   0 & 1 & 0 & 0 & 0  \cr
                                    0 & 0 & 0 & 0 & 1 \cr
                                    0 & 0 & 0 & 1 &  0 \cr \end{pmatrix}  .$
\begin{enumerate}
    \item Calculer $ \hbox{det }(S) .$ D\' eterminer (de pr\' ef\' erence
sans calcul) $ S^{-1} .$
    \item Montrer qu'il existe deux sous espaces vectoriels $ E_1  $ et $
E_2  $ de $ {\Rr }^5 $ de dimension respective $ 2  $ et $ 3 $ tels
que~:
 ${\Rr }^5=E_1 \oplus E_2 \oplus E_3  $ et  $ S(E_1) \subset E_1  $
$ S(E_2) \subset E_2  .$
    \item Montrer qu'il existe $ x \in E_2  $ tels que $ Sx= x .$
En d\' eduire que la d\' ecomposition qui pr\' ec\' ede n'est pas unique.
\end{enumerate}
\finenonce{001156}



\finexercice



\exercice{1157, ridde, 1999/11/01}

\enonce{001157}{}
 Soit $A \in M_3 (\Rr)$ anti-sym\'etrique. Calculer
$\det (A)$. Ce r\'esultat vaut-il encore pour $A \in M_2 (\Rr)$ ?

\finenonce{001157}



\finexercice



\exercice{1158, ridde, 1999/11/01}

\enonce{001158}{}
Soient $n = 2$ ou $3$ et $A \in M_n (\Qq)$.
\begin{enumerate}
\item Montrer que si $\forall X \in M_n (\Qq)$ $\det (A + X) = \det (X)$ alors $A = 0$.
\item Soit $B \in M_n (\Qq)$ telle que $\forall X \in M_n (\Qq)$ $\det (A + X) =
 \det (B + X)$. Montrer que $A = B$.
\end{enumerate}
\finenonce{001158}



\finexercice



\exercice{1159, gourio, 2001/09/01}

\enonce{001159}{}
Soit $\left( A,B\right) \in M_{n}({\Rr})^{2}$ tel que $A^{2}+B^{2}=AB$ et
$AB-BA  $ inversible. Montrer que $3$ divise $n.$
\finenonce{001159}



\finexercice



\exercice{1160, gourio, 2001/09/01}

\enonce{001160}{}
Montrer que si $ n\in \Nn-\{0,1\},A\in M_{n}({\Rr})$, on a :
$$\det(\mathrm{Com}(A))=\det(A)^{n-1}.$$
\finenonce{001160}



\finexercice



\exercice{1161, gourio, 2001/09/01}

\enonce{001161}{}
Montrer que si $ n\in \Nn^{*}, A\in M_{n}({\Rr})$ :
$$\mathrm{rg}(A)=n\Rightarrow \mathrm{rg}(\mathrm{Com}(A))=n ; $$
$$\mathrm{rg}(A)=n-1\Rightarrow \mathrm{rg}(\mathrm{Com}(A))=1 ; $$
$$\mathrm{rg}(A)\leq n-2\Rightarrow \mathrm{rg}(\mathrm{Com}(A))=0. $$
\finenonce{001161}



\finexercice



\exercice{1162, gourio, 2001/09/01}

\enonce{001162}{}
Soit $ A=\left( a_{i,j}\right) _{(i,j)\in \{1,...,n\}^{2}}\in M_{n}({\Rr})$
telle que :
\begin{eqnarray*}
\forall i &\in &\{1,...,n\},\sum_{j=1}^{n}a_{i,j}\leq 1, \\
\forall (i,j) &\in &\{1,...,n\}^{2},a_{i,j}\in [0,1[.
\end{eqnarray*}
Montrer que $ \left| \det(A)\right| <1.$
\finenonce{001162}



\finexercice



\exercice{2451, matexo1, 2002/02/01}

\enonce{002451}{}
Soit $A = (a_{i,j})$ une matrice complexe dont les coefficients
v\'erifient $|a_{i,j}| \le 1$. Montrer que $|\det A| \le 1$.
\finenonce{002451}


\finexercice


\exercice{2457, matexo1, 2002/02/01}

\enonce{002457}{}
Soit $A, B$ deux matrices carr\'ees d'ordre $n$, $A$ inversible.
\begin{enumerate}
\item Montrer que $\det(A +\lambda  B)$ est un polyn\^ome en
$\lambda$ de degr\'e $n$. Quels sont ses termes de plus haut et
de plus bas degr\'e ?
\item En d\'eduire que si $A$ est une matrice inversible,
pour toute matrice $B$, il existe $\varepsilon_0 > 0$ tel que $A
+\varepsilon B$ soit aussi inversible, $\forall \varepsilon \in
[-\varepsilon_0,\varepsilon_0]$.
\end{enumerate}
\finenonce{002457}


\finexercice


\exercice{3426, quercia, 2010/03/10}
\enonce{003426}{$\det(I-AB) = \det(I-BA)$}

Soient $A \in \mathcal{M}_{p,q}(K)$ et $B \in \mathcal{M}_{q,p}(K)$.
Montrer que $\det(I_p-AB) = \det(I_q-BA)$.
(Commencer par le cas où $A$ est la matrice canonique de rang $r$)
\finenonce{003426}



\finexercice

\exercice{3429, quercia, 2010/03/10}
\enonce{003429}{$\det(A^2+B^2)$}

\begin{enumerate}
  \item Soient $A,B \in \mathcal{M}_n(\R)$ telles que $AB = BA$.
    Montrer que $\det(A^2 + B^2) \ge 0$.
  \item Chercher $A,B$ ne commutant pas telles que $\det(A^2 + B^2) < 0$.
    
\end{enumerate}
\finenonce{003429}



\finexercice


\exercice{3430, quercia, 2010/03/10}
\enonce{003430}{Déterminant par blocs}

Soient $A,B,C,D \in \mathcal{M}_n(K)$ avec $A$ inversible et $AC=CA$. On considère
$M = \begin{pmatrix} A & B \cr C & D \cr \end{pmatrix} \in \mathcal{M}_{2n}(K)$.
Montrer que $\det M = \det(AD-CB)$.
\finenonce{003430}



\finexercice

\exercice{3434, quercia, 2010/03/10}
\enonce{003434}{Système linéaire homogène}

On considère un système linéaire homogène : $(S) \iff AX = 0$, avec
$A \in \mathcal{M}_{n,p}(K)$, $n < p$ et $\mathrm{rg} A = n$.

\begin{enumerate}
  \item  Montrer qu'on peut compléter $A$ en une matrice
     $B = \begin{pmatrix} A \cr A' \cr \end{pmatrix}$ inversible.
  \item  Montrer que les colonnes $n+1$ à $p$ de ${}^{t\!}\text{com}\, B$ constituent
     une base des solutions de $(S)$.
  \item  Consid\'erer l'exemple suivant : 
  $$(S) \iff \left\{
  \begin{array} {lllllllll}
  x &+& 2y &+& 3z &+& 4t &=& 0  \cr
 2x &+& 3y &+& 4z &+& 5t &=& 0. \cr
 \end{array}
 \right.$$
     

\end{enumerate}
\finenonce{003434}



\finexercice


\exercice{3435, quercia, 2010/03/10}
\enonce{003435}{Inégalité}

Soit $A = (a_{ij}) \in \mathcal{M}_n(\R)$.
Démontrer que : $|\det A| \le \prod_{i=1}^n \sum_{j=1}^n |a_{ij}|$.


Quand y a-t-il égalité ?
\finenonce{003435}



\finexercice


\exercice{3436, quercia, 2010/03/10}
\enonce{003436}{Déterminants $2\times2$ imposés}

\def \dddd#1#2{\text{det}\bigl(\vec #1,\vec #2\bigr)}\relax
Soient $\vec a, \vec b, \vec c, \vec d$ quatre vecteurs d'un ev $E$ de
dimension 2. On note $\det$ le déterminant dans une base fixée de $E$.

\begin{enumerate}
  \item Démontrer que : $\dddd ab \dddd cd + \dddd ac \dddd db + \dddd ad \dddd bc = 0$.
    \par
    (Commencer par le cas où $(\vec a,\vec b)$ est libre)

  \item On donne six scalaires : $d_{ab}, d_{ac}, d_{ad}, d_{cd}, d_{db}, d_{bc}$
    tels que $d_{ab} d_{cd} + d_{ac} d_{db} + d_{ad} d_{bc} = 0$.
    \par
    Montrer qu'il existe des vecteurs $\vec a, \vec b, \vec c, \vec d$ tels que :
    $\forall\ x,y,\ d_{xy} = \d xy$.
   
\end{enumerate}
\finenonce{003436}



\finexercice

\exercice{3437, quercia, 2010/03/10}
\enonce{003437}{Décomposition d'un vecteur en dimension 3}

\def \dddd #1#2#3{\det(\vec #1,\vec #2, \vec #3)}\relax
Soient $\vec a, \vec b, \vec c, \vec d$ quatre vecteurs d'un ev $E$ de
dimension 3.
On note : det le déterminant dans une base fixée de~$E$.

Démontrer que : $\d abc \vec d = \dddd abd \vec c + \dddd adc \vec b + \dddd dbc \vec a$.


\finenonce{003437}



\finexercice

\exercice{3439, quercia, 2010/03/10}
\enonce{003439}{$\det( u + n )$}

Soient $u,n \in \mathcal{L}(E)$ deux endomorphismes d'un $\C$-ev de dimension finie,
$u$ inversible, $n$ nilpotent, avec $u\circ n = n \circ u$.

\begin{enumerate}
  \item Démontrer que $\det n = 0$.
  \item Chercher le polynôme caractéristique de $n$.
    En déduire que $\det(\text{id}_E + n ) = 1$.
  \item Démontrer que $\det( u + n ) = \det u$.
\end{enumerate}
\finenonce{003439}



\finexercice


\exercice{3440, quercia, 2010/03/10}
\enonce{003440}{Sev stables}

Soit $f \in \mathcal{L}(E)$ tel qu'il existe deux sev $F,G$ supplémentaires et
stables par $f$.

Démontrer que $\det f = (\det f_{|F})(\det f_{|G})$.

\finenonce{003440}



\finexercice


\exercice{3441, quercia, 2010/03/10}
\enonce{003441}{Groupe $SL_n( K)$}

On note $SL_n( K) = \{ M \in \mathcal{M}_n(K) \text{ tq } \det M = 1 \}$.

\begin{enumerate}
  \item
  \begin{enumerate}
     \item Démontrer que $SL_n( K)$ est un groupe pour le produit matriciel.
     \item Démontrer que $SL_n( K)$ est engendré par les matrices :
    $I + \lambda E_{ij}$, $( j \ne i )$
    où $(E_{ij})$ est la base canonique de $\mathcal{M}_n(K)$, et $\lambda \in  K$
    (transformer une matrice $M \in SL_n( K)$ en $I$ par opérations
    élémentaires).
  \end{enumerate}
  \item
  \begin{enumerate}
    \item Soit $M \in \mathcal{M}_n(\Z)$. Démontrer que $M$ a une inverse dans $\mathcal{M}_n(\Z)$ si
     et seulement si $\det M = \pm1$.
    \item  Démontrer que le groupe $SL_n(\Z)$ est engendré par les matrices
     $I + E_{ij}$, $( j \ne i )$.
   \end{enumerate}  
\end{enumerate}
\finenonce{003441}



\finexercice


\exercice{3442, quercia, 2010/03/10}
\enonce{003442}{Déterminant de $X  \mapsto AX$}

Soit $A \in \mathcal{M}_n(K)$ et ${f_A} : {\mathcal{M}_n(K)} \to {\mathcal{M}_n(K)}, X \mapsto {AX.}$
Calculer $\det f_A$.


\finenonce{003442}



\finexercice

\exercice{5368, rouget, 2010/07/06}
\enonce{005368}{***I}
Soit $(A,B)\in(M_n(\Rr))^2$ et $C=\left(
\begin{array}{cc}
A&B\\
-B&A
\end{array}
\right)\in M_{2n}(\Rr)$. Montrer que $\mbox{det}C\geq0$.

\finenonce{005368}


\finexercice


\exercice{5370, rouget, 2010/07/06}
\enonce{005370}{***I}
Déterminer les matrices $A$, carrées d'ordre $n$, telles que pour toute matrice carrée $B$ d'ordre $n$ on a $\mbox{det}(A+B)=\mbox{det}A+\mbox{det}B$.

\finenonce{005370}


\finexercice


\exercice{5379, rouget, 2010/07/06}
\enonce{005379}{}
Soit $E$ un ensemble contenant au moins $n$ éléments et $(f_1,f_2...,f_n)$ un $n$-uplet de fonctions de $E$ dans $\Cc$. Montrer que les propositions suivantes sont équivalentes~:
\begin{enumerate}
\item  la famille $(f_1,...,f_n)$ est libre ;
\item  il existe $n$ éléments $a_1$, $a_2$,..., $a_n$ dans $E$ tels que $\mbox{det}(f_i(a_j))_{1\leq i,j\leq n}\neq0$.
\end{enumerate}
\finenonce{005379}


\finexercice


\exercice{5380, rouget, 2010/07/06}
\enonce{005380}{}
Déteminer l'inverse de $A=(a_{i,j})$ telle que $a_{i,i+1}=a_{i,i-1}=1$ et $a_{i,j}=0$ sinon.
\finenonce{005380}


\finexercice


\exercice{5614, rouget, 2010/10/16}
\enonce{005614}{**I}
\label{ex:rou17}
Soit $A$ une matrice carrée de format $n$. Calculer le déterminant de sa comatrice.
\finenonce{005614}


\finexercice


\exercice{5615, rouget, 2010/10/16}
\enonce{005615}{***I}
\label{ex:rou18bis}
Soit $A$ une matrice carrée de format $n$. Etudier le rang de $\text{com}A$ en fonction du rang de $A$.
\finenonce{005615}


\finexercice


\exercice{5616, rouget, 2010/10/16}
\enonce{005616}{***}
Résoudre dans $\mathcal{M}_n(\Rr)$ l'équation $M=\text{com}M$ ($n\geqslant2$).
\finenonce{005616}


\finexercice


\exercice{5648, rouget, 2010/10/16}
\enonce{005648}{***}
Soit $A$ une matrice carrée complexe de format $n$ ($n\geqslant 2$) telle que pour tout élément $M$ de $M_n(\Cc)$, on ait $\text{det}(A+M)=\text{det}A+\text{det}M$. Montrer que $A = 0$.
\finenonce{005648}


\finexercice

\section{201.01 Valeur propre, vecteur propre}


\exercice{1597, legall, 1998/09/01}

\enonce{001597}{}
Soit $  m \in \R   $ et $  A_m\in M_3 (\R )  $
la matrice $  \begin{pmatrix} m & 1 & 1 \cr
                                      1 & m & 1 \cr
                         1 & 1 & m\cr \end{pmatrix} .$
\begin{enumerate}
    \item Calculer les valeurs propres de $  A_m  $ et une base de vecteurs propres.
    \item D\'eterminer suivant les valeurs de $  m  $ le rang de $  A_m  .$ D\'eterminer lorsque  cela est possible $  A_m^{-1}  .$
    \item Lorsque $  A_m  $ n'est pas inversible d\'eterminer le noyau et l'image de $  A_m  .$
\end{enumerate}
\finenonce{001597}



\finexercice



\exercice{1598, legall, 1998/09/01}

\enonce{001598}{}
Soit $A \in \mathcal{O}_n(\R)$. Montrer que si $-1$ n'est pas valeur
propre de $A$, alors il existe une matrice $Q$ antisym\'etrique (i.e.
${}^tQ = -Q$) telle que $A = (I + Q)^{-1}(I - Q) = (I - Q)(I + Q)^{-1}$ et
qu'on a $A \in \mathcal{SO}_n(\R)$. R\'eciproque ?
\finenonce{001598}



\finexercice



\exercice{1599, ortiz, 1999/04/01}

\enonce{001599}{}
Soient $E$ un K-espace vectoriel de dimension
finie et $f,g\in \mathcal{L}(E)$. Montrer que si
$\lambda$ est valeur propre de $g\circ f$ alors
$\lambda$ est valeur propre de $f\circ g$ (on
distinguera les cas $\lambda=0$ et $\lambda\not=
0$).
\finenonce{001599}



\finexercice



\exercice{1600, roussel, 2001/09/01}

\enonce{001600}{}
\begin{enumerate}
  \item Soient $f$ et $g$ deux endomorphisme s d'un espace vectoriel $E$ de dimension
  $n$ sur \\
  $K=\mathbb{R}$ ou $\mathbb{C}$, ayant chacun $n$ valeurs propres distinctes dans
  $K$. Montrer que
  $$ f \circ g = g \circ f \quad \Longleftrightarrow \quad f \mbox{~et~} g \mbox{~ont les m\^emes vecteurs propres}.$$
  \item Supposons maintenant que $K=\mathbb{C}$ et que $ f \circ g = g \circ f $.
  Si $u$ est un endomorphisme ~on dit qu'un espace vectoriel $F$ est $u$-stable si
  $u(F) \subset F$. Montrer que tout sous-espace propre de $f$ est $g$-stable. 

  \emph{Remarque} : On peut montrer par r\'ecurrence sur $n$ qu'il existe un
  vecteur propre commun \`a $f$et $g$. On admettra ce r\'esultat.
  
  \item Consid\'erons $f$ et $g$ deux endomorphismes de $\mathbb{R}^3$ dont les
  matrices dans la base canonique sont respectivement
  $$ M = \left (
  \begin{array}{cccc}
    5  & -4  & -4   \\
    1  & 0  & -2   \\
    1  & -1  & 1   \\
  \end{array}
  \right )
  ~~~~~\text{et}~~~~~N = \left (
  \begin{array}{ccc}
    -2  & 2  & 2  \\
    1  & -1  & -2  \\
    -2  & 2  & 3  \\
  \end{array}
  \right ) $$
  
  \begin{itemize}
    \item V\'erifier que $f \circ g = g \circ f$ et d\'eterminer les sous-espaces
    propres de $M$ et $N$.
    \item D\'eterminer une base de $\mathbb{R}^3$ dans laquelle les matrices de $f$
    et $g$ sont diagonales.
  \end{itemize}
\end{enumerate}
\finenonce{001600}



\finexercice


\exercice{1601, roussel, 2001/09/01}

\enonce{001601}{}
Soient $A \in \mathcal{M}_4 \left (\mathbb{R} \right )$ et $B \in
\mathcal{M}_3 \left (\mathbb{R} \right )$. Soit $f$ l'endomorphisme ~associ\'e \`a la
matrice $A$.

$$ A = \left (
  \begin{array}{cccc}
  5  & ~3  & -1  & 3 \\
  0  & -1  & ~1  & 2 \\
  0  & ~2  & ~1  & 2 \\
  0  & ~0  & ~0  & 1 \\
  \end{array}
\right )
 ~~~~~~~~~~B = \left (
                  \begin{array}{ccc}
                   5  & ~3  & -1  \\
                   0  & -1  & ~1  \\
                   0  & ~2  & ~1  \\
                  \end{array}
\right ) $$

\begin{enumerate}
\item \emph{Uniquement} en examinant la matrice $A$, trouver deux valeurs propres ~et
un vecteur propre de $A$, puis deux sous-espaces $f-$stables.
\item Que repr\'esente la matrice $B$?
\end{enumerate}
\finenonce{001601}



\finexercice



\exercice{1602, barraud, 2003/09/01}

\enonce{001602}{}
Soit $u\in\mathrm{End}(E)$. On note $\chi_{u}=(-1)^{n}X^{n}+a_{n-1}X^{n-1}+\cdots+a_{0}$. Montrer
que
$$
 a_{0}=\det(u) \qquad \text{ et } \qquad a_{n-1}=(-1)^{n-1}\mathrm{tr}(u)
$$
\finenonce{001602}



\finexercice



\exercice{1603, barraud, 2003/09/01}

\enonce{001603}{}
Soient $u$ et $v$ deux endomorphismes de $E$. Montrer que $u\circ v$ et $v\circ u$ ont les
m\^{e}mes valeurs propres.
\finenonce{001603}



\finexercice



\exercice{1604, barraud, 2003/09/01}

\enonce{001604}{}
Soient $u$ et $v$ deux endomorphismes de $E$ qui commutent, c'est \`{a} dire tels que $u\circ
v=v\circ u$. On suppose que $v$ admet $n$ valeurs propres distinctes. Montrer qu'il existe
une base de $E$, form\'{e}e de vecteurs propres communs \`{a} $u$ et \`{a} $v$.

En d\'{e}duire qu'il existe $(a_{0},\ldots,a_{n-1})\in\mathbb{K}^{n}$ tel que
$
 \quad u=a_{0}\mathrm{id}+a_{1}v+\cdots+a_{n-1}v^{n-1}
$
\finenonce{001604}



\finexercice



\exercice{1605, barraud, 2003/09/01}

\enonce{001605}{}
On considère les matrices suivantes :
$$
  A=\begin{pmatrix}
   1 &  0 &  0 & -1 \\
  -1 & -1 &  0 &  1 \\
  -2 &  0 &  0 &  2 \\
   0 & -1 & 0 & 0
    \end{pmatrix}
\qquad
  B=\begin{pmatrix}
       1 &  0 &  0 & -1 \\
      -1 & -1 &  0 &  1 \\
      -1 & -1 &  1 &  3 \\
      -1 &  0 & -1 & -1
    \end{pmatrix}
\qquad
  C=\begin{pmatrix}
       0 & -1 &  0 &  0 \\
       0 &  0 &  0 &  0 \\
       1 &  1 &  1 &  1 \\
      -1 &  0 & -1 & -1
    \end{pmatrix}
$$
En effectuant le moins de calculs possible,
\begin{enumerate}
\item
montrer que \hspace{1cm}
 $\{0\}\subset\mathrm{Ker} A\subset\mathrm{Ker} A^{2}\subset\mathrm{Ker} A^{3}=\R^{4}$\\
et déterminer les dimensions respectives de $\mathrm{Ker} A$ et $\mathrm{Ker}
A^{2}$,

\item
déterminer un vecteur $e_{1}$ tel que
 $\R^{4}=\mathrm{Ker} A^{2}\oplus\mathrm{Vect}(e_{1})$,

\item
montrer que $(e_{1},Ae_{1},A^{2}e_{1})$ est une famille libre,

\item
montrer que $Ae_{1}\in\mathrm{Ker} A^{2}$, et que
 $\mathrm{Ker} A^{2}=\mathrm{Ker} A\oplus\mathrm{Vect}(Ae_{1})$,

\item
montrer que $A^{2}e_{1}\in\mathrm{Ker} A$ et déterminer un vecteur $e_{2}$
tel que $\mathrm{Ker} A=\mathrm{Vect}(A^{2}e_{1})\oplus\mathrm{Vect}(e_{2})$,

\item
montrer que $(e_{1},Ae_{1},A^{2}e_{1},e_{2})$ est une base de
$\R^{4}$.

\item
Soit $P$ la matrice de passage de la base canonique à la base
$(A^{2}e_{1},Ae_{1},e_{1},e_{2})$. Caluler $P^{-1}AP$.
\end{enumerate}

Adapter ce travail à l'étude de $B$ et $C$
\finenonce{001605}



\finexercice



\exercice{1606, barraud, 2003/09/01}

\enonce{001606}{}
Soit $J$ la matrice $$J=
\begin{pmatrix}
  1 &\cdots &1\\
  \vdots&   &\vdots\\
  1 &\cdots &1
\end{pmatrix}
$$.
\begin{enumerate}
\item 
Trouver une relation entre $J$ et  $J^{2}$.
\item
En déduire les valeurs propres de $J$ et calculer leurs multiplicités.
\item
Donner le polynôme caractéristique de $J$.
\end{enumerate}
\finenonce{001606}



\finexercice



\exercice{1607, barraud, 2003/09/01}

\enonce{001607}{}
  Soient $A$ et $B$ deux matrices de $\mathcal{M}_{n}(\R)$ telles que
  $$
  AB-BA=A
  $$
  Le but de cet exercice est de montrer que $A$ est nilpotente, c'est à
  dire $$\exists k\in\N, A^{k}=0.$$


  On note $E$ l'espace vectoriel $\mathcal{M}_{n}(\R)$ et on considère
  l'application~:\quad
  $$\displaystyle \psi
   \begin{array}{ccc}
     E & \rightarrow  & E \\
     M &\mapsto  & MB-BM
   \end{array}
  $$
  
  \begin{enumerate}
  \item
    Montrer que $\psi$ est linéaire de $E$ dans $E$.
  \item
    Montrer par récurrence que~: $\forall k\in\N\quad
    \psi(A^{k})=kA^{k}$.
  \item 
    On suppose que $\forall k\in\N, A^{k}\neq0$.
    Montrer que $\psi$ a une infinité de valeurs propres.
  \item 
    Conclure.
\end{enumerate}  
\finenonce{001607}



\finexercice



\exercice{1608, barraud, 2003/09/01}

\enonce{001608}{}
  Soit $M$ la matrice suivante~: $M=
  \begin{pmatrix}
    1 &1 &0 \\
    -1&0 &0 \\
    2 &0 &-1
  \end{pmatrix}
$. Calculer le polynôme caractéristique de $M$. En déduire $M^{-1}$.
\finenonce{001608}



\finexercice



\exercice{1609, barraud, 2003/09/01}

\enonce{001609}{}
  Soit $f$ un endomorphisme de $E=\C^{n}$. Soit $\pi_{1},...,\pi_{N}$ des
  endomorphismes tous non nuls de $E$ et $\lambda_{1},...,\lambda_{N}$
  $N$ nombres complexes distincts. On suppose que~:
  $$
  \forall m\in\N \quad f^{m}=\sum_{k=1}^{N}\lambda_{k}^{m}\pi_{k}.
  $$

  \begin{enumerate}
\item
    Montrer que $\forall P\in\C[X],\quad
    P(f)=\sum_{k=1}^{N}P(\lambda_{k})\pi_{k}$

\medskip

  On considère le polynôme $Q=\prod_{1\leq k\leq N}(X-\lambda_{k})$ et
  pour chaque $p\in\{1,...,N\}$ les polynômes suivants~:  
  $$
  Q_{p}=\prod_{\substack{ 1\leq k\leq N \\k\neq p}}
             (X-\lambda_{k})\qquad\text{et}\qquad
  \tilde Q_{p}=\frac{1}{Q_{p}(\lambda_{p})}\ Q_{p}
  $$
\item
 
    Calculer $Q(f)$. Qu'en déduit-on pour $f$~?
\item

    Montrer que $Sp(f)\subset\{\lambda_{1},...,\lambda_{N}\}$
\item

    Montrer que $\tilde Q_{p}(f)=\pi_{p}$. Vérifier alors que
    $\pi_{p}\circ\pi_{q}= \left\{
    \begin{array}{l}
      0       \text{ si } p\neq q\\
      \pi_{p} \text{ si } p=q
    \end{array}
    \right. $
\item

    Calculer $f\circ\pi_{p}$. En déduire que
    $Sp(f)=\{\lambda_{1},...,\lambda_{N}\}$.


\medskip
  On note $E_{p}$ l'espace propre associé à la valeur propre
  $\lambda_{p}$. 
\item
    Montrer que $\mathrm{Im}\pi_{p}\subset E_{p}$. Réciproquement, pour $x\in
    E_{p}$, montrer que $x\in\mathrm{Ker}\pi_{q}$ pour $q\neq p$ (on calculera
    par exemple $\pi_{q}\circ f(x)$ de deux façons différentes) puis que
    $x=\pi_{p}(x)$. En déduire que $E_{p}\subset\mathrm{Im}\pi_{p}$.
\item    En déduire que $\mathrm{Im}\pi_{p}=E_{p}$ et que
    $\mathrm{Ker}\pi_{p}=\bigoplus_{q\neq p}E_{q}$. Décrire géométriquement
    $\pi_{p}$.
  \end{enumerate}
\finenonce{001609}



\finexercice



\exercice{1610, barraud, 2003/09/01}

\enonce{001610}{}
  On considère l'application suivante~:
  $$
  f:
  \begin{array}{ccc}
    \R_{n}[X] &\rightarrow  &\R_{n}[X] \\
    P         &\mapsto &(X^{2}-1)P'-(nX+a)P
  \end{array}
  $$
  Vérifier que cette application est bien définie.

  Déterminer ses valeurs propres, et les espaces propres associés.
\finenonce{001610}



\finexercice



\exercice{1611, barraud, 2003/09/01}

\enonce{001611}{}
  Soit $E$ un espace vectoriel de dimension $n$ et $u$ un endomorphisme
  de $E$ ayant $n$ valeurs propres distinctes
  $\{\lambda_{1},...,\lambda_{n}\}$. 

\begin{enumerate}
\item 
Montrer que l'ensemble $\mathrm{Com}=\{v\in\mathcal{L}(E,E)/ uv=vu\}$ des
endomorphismes de $E$ qui commutent avec $u$ est un espace vectoriel.

\item
  \begin{enumerate}
  \item 
    Soit $v$ un élément de $\mathrm{Com}$. Montrer que $v$ préserve les espaces
    propres de $u$ (c'est à dire que si $E_{\lambda}$ est un espace
    propre de $u$ associé à la valeur propre $\lambda$, on a $\forall
    x\in E_{\lambda}, v(x)\in E_{\lambda}$). 
  \item
    Donner la dimension des espaces propres de $u$ et montrer que si
    $x$ est un vecteur propre de $u$ alors c'est aussi un vecteur propre
    de $v$. 
  \item
    A l'aide d'une base convenablement choisie, décrire tous les éléments
    de $\mathrm{Com}$, et montrer que $\mathrm{Com}$ est de dimension $n$. 
  \end{enumerate}

\item
Montrer que $\mathrm{Vect}(\mathrm{id},u,u^{2},...,u^{n-1})\subset \mathrm{Com}$.

\item
On veut maintenant étudier l'indépendance linéaire de la famille
$\{\mathrm{id},u,u^{2},...,u^{n-1}\}$. Pour cela, on considère $n$ réels
$\alpha_{0},...,\alpha_{n-1}$ tels que $\sum_{i=0}^{n}\alpha_{i}u^{i}=0$.
  \begin{enumerate}
  \item 
    Montrer que les $(\alpha_{i})$ sont solution du système~:
    $$
    (*)\ \left\{
      \begin{array}{ccccccccccc}
         \alpha_{0}
        &+&\alpha_{1}  \lambda_{1}
        &+&\alpha_{2}  \lambda_{1}^{2}
        &+&...
        &+&\alpha_{n-1}\lambda_{1}^{n-1}
        &=&0\\
        \alpha_{0}
        &+&\alpha_{1}  \lambda_{2}
        &+&\alpha_{2}  \lambda_{2}^{2}
        &+&...
        &+&\alpha_{n-1}\lambda_{2}^{n-1}
        &=&0\\
        \vdots&&\vdots&&\vdots&&&&\vdots&&\vdots\\
        \alpha_{0}
        &+&\alpha_{1}  \lambda_{n}
        &+&\alpha_{2}  \lambda_{n}^{2}
        &+&...
        &+&\alpha_{n-1}\lambda_{n}^{n-1}
        &=&0
      \end{array}
    \right.
    $$

  \item
    On rappel que~: $
    \begin{vmatrix}
      1&\lambda_{1}&\lambda_{1}^{2}&...&\lambda_{1}^{n-1}\\
      1&\lambda_{2}&\lambda_{2}^{2}&...&\lambda_{2}^{n-1}\\
      \vdots&\vdots&\vdots&&\vdots\\
      1&\lambda_{n}&\lambda_{n}^{2}&...&\lambda_{n}^{n-1}
    \end{vmatrix}=\prod\limits_{1\leq i<j\leq n}(\lambda_{j}-\lambda_{i})
    $. En déduire l'ensemble des solutions du système $(*)$ et conclure.
  \end{enumerate}

\item  
Montrer que $\mathrm{Com}=\mathrm{Vect}(\mathrm{id},u,u^{2},...,u^{n-1})$.
\end{enumerate}

\finenonce{001611}




\finexercice



\exercice{2467, matexo1, 2002/02/01}

\enonce{002467}{}
Donner les valeurs propres, vecteurs propres et matrice de
diagonalisation \'eventuelle des matrices suivantes dans $\C^2$\,:
$$ \left( \begin{array}{cc} 4&4\\ 1&4 \end{array} \right), \qquad
\left( \begin{array}{cc} 2&5\\ 4&3 \end{array} \right), \qquad
\left( \begin{array}{cc} 5&3\\ -8&6\end{array} \right).$$
\finenonce{002467}


\finexercice


\exercice{2468, matexo1, 2002/02/01}

\enonce{002468}{}
Soit $\mathbb K$ le corps des r\'eels ou des
complexes, et $u$ l'endomorphisme de $\mathbb K^3$ ayant pour matrice
$$A = \left( \begin{array}{ccc} 0&-2&0\\ 1&0&-1\\ 0&2&0 \end{array} \right).$$
\'Etudier, dans les deux cas $\mathbb K = \R$ et $\mathbb K = \C$, si $u$ est
diagonalisable. En donner une forme diagonalis\'ee dans une base
dont on donnera la matrice de passage par rapport \`a la base
canonique.

\finenonce{002468}


\finexercice


\exercice{2470, matexo1, 2002/02/01}

\enonce{002470}{}
Soit $M$ une matrice de $\mathcal M_n(\C)$\,; on suppose qu'il existe un
entier $p$ tel que $M^p = I$. Montrer que si $\omega$ est une
racine $p$-i\`eme de l'unit\'e, c'est une valeur propre de $M$ ou
alors $M$ v\'erifie
$$ M^{p-1} + \omega M^{p-2} +\cdots+\omega^{p-2}M +\omega^{p-1}I=0.$$
\finenonce{002470}


\finexercice


\exercice{2471, matexo1, 2002/02/01}

\enonce{002471}{}
Soit $E$ un espace vectoriel de dimension $n$ sur un corps
$K$, $u$ un endomorphisme de $E$ et $P \in K[X]$. On suppose que
$u$ v\'erifie l'\'equation $P(u) = 0$.

\begin{enumerate}
\item Montrer que si $\lambda$ est une valeur propre de $u$, alors
$P(\lambda) = 0$.
\item On suppose que $P$ est de la forme
$$P(X) = \prod_{i=1}^n (X - a_i), \quad\mbox{\rm avec } a_i \neq a_j
\mbox{ \rm si } i\neq j.$$
Montrer que les seules matrices v\'erifiant
$P(M) = 0$ sont de la forme $Q^{-1}DQ$, avec $Q$ matrice inversible
quelconque et $D$ matrice diagonale que l'on pr\'ecisera.
Combien y a-t-il de matrices diagonales de ce type\,?
\end{enumerate}
\finenonce{002471}


\finexercice


\exercice{2472, matexo1, 2002/02/01}

\enonce{002472}{}
Soient $u, v$ deux endomorphismes d'un espace vectoriel
complexe $E$. On suppose qu'il existe $a, b$ complexes tels que
$u\circ v = a u+ b v$. Montrer que $u$ et $v$ ont un vecteur
propre commun.
\finenonce{002472}


\finexercice


\exercice{2473, matexo1, 2002/02/01}

\enonce{002473}{}
Soit $E$ l'ensemble des polyn\^omes \`a coefficients r\'eels de
degr\'e inf\'erieur ou \'egal \`a $n$. Soit $u$ l'application qui \`a tout
polyn\^ome $P$ de $E$ fait correspondre $u(P) = P(X-1)$.

\begin{enumerate}
\item V\'erifier que $u$ est un endomorphisme de $E$. Est-il
injectif ou surjectif\,?
\item Trouver les valeurs propres et vecteurs propres de
$u$ ainsi qu'une base dans laquelle $u$ est diagonalisable.
\end{enumerate}
\finenonce{002473}


\finexercice


\exercice{2474, matexo1, 2002/02/01}

\enonce{002474}{}
Soit $A, B$ deux matrices carr\'ees d'ordre $n$ \`a coefficients
complexes qui commutent ($AB = BA$). On suppose en outre que
toutes les valeurs propres de $B$ sont distinctes.
\begin{enumerate}
\item Montrer que tout vecteur propre de $B$ est vecteur propre
de $A$.
\item Montrer que $A$ est de la forme
$ A = P(B) $, o\`u $P$ est un polyn\^ome de $\C[X]$ de degr\'e
inf\'erieur ou \'egal \`a $n-1$.
\end{enumerate}
\finenonce{002474}


\finexercice


\exercice{2476, matexo1, 2002/02/01}

\enonce{002476}{}
Soit $E$ l'ensemble des polyn\^omes \`a coefficients r\'eels de
degr\'e inf\'erieur ou \'egal \`a $n$. \'Etant donn\'e deux r\'eels $a,b$, on
note $u$ l'application qui \`a tout polyn\^ome $P$ de $E$ fait
correspondre 
$$ u(P) = (X-a)(X-b)P' -n XP $$
\begin{enumerate}
\item V\'erifier que $u$ est un endomorphisme de $E$.
\item Trouver les valeurs propres et vecteurs propres de $u$.
\item Trouver le noyau de $u$, ainsi que l'ensemble des
polyn\^omes qui v\'erifient $u(P) = 1$.
\end{enumerate}
\finenonce{002476}


\finexercice


\exercice{2479, matexo1, 2002/02/01}

\enonce{002479}{}
On considère la matrice $N\times N$
\[                          
 M=\begin{pmatrix}
b & c & 0   &.. & 0 & 0\cr 
a & b & c & ..& 0 & 0\cr
0   & a & b & ..& 0 & 0\cr
..  &..   &..   &.. &.. & ..\cr
0   & 0   &0    &.. &b&c  \cr
0   & 0   &0    &.. &a    &b  \cr
\end{pmatrix}                       
\]     
où $a,b$ et $c$ sont trois nombres complexes, avec $c\neq 0$.            
On note $V$ un vecteur propre associé à la valeur propre $\lambda$ de $M$.\\
Ecrire les relations reliant les composantes de $V$.\\ 
Déterminer toutes les valeurs propres de $M$.
\finenonce{002479}

\finexercice


\exercice{2570, delaunay, 2009/05/19}
\enonce{002570}{}
 Soit $A$ une matrice carr\'ee d'ordre $n$. On suppose que $A$ est inversible et que $\lambda\in\R$ est une valeur propre de $A$.
\begin{enumerate}
  \item D\'emontrer que $\lambda\neq 0$.
  \item D\'emontrer que si $\vec x$ est un vecteur propre de $A$ pour la valeur propre $\lambda$ alors il est vecteur propre de $A^{-1}$ de valeur propre
 $\lambda^{-1}$ .
\end{enumerate}
\finenonce{002570} 


\finexercice


\exercice{2571, delaunay, 2009/05/19}
\enonce{002571}{}
Soit $f$ un endomorphisme de $E$ v\'erifiant $f^2={\rm mathrm{Id}}_E$.
\begin {enumerate}
  \item  D\'emontrer que les seules valeurs propres possibles de $f$ sont $1$ et $-1$.
  \item V\'erifier que pour tout $\vec x\in E$, on a
$$f(\vec x-f(\vec x))=-(\vec x-f(\vec x))\ \ {\hbox{et}}\ \ f(\vec x+f(\vec x))=(\vec x+f(\vec x))$$
et en d\'eduire que $f$ admet toujours une valeur propre.
  \item  D\'emontrer que si $1$ et $-1$ sont valeurs propres, alors $E$ est somme directe des sous-espaces propres correspondants.

  \item Traduire g\'eom\'etriquement sur un dessin dans le cas $n=2$.
\end {enumerate}
\finenonce{002571} 


\finexercice


\exercice{2579, delaunay, 2009/05/19}
\enonce{002579}{}
Soit $E$ un espace vectoriel sur un corps $K$ ($K=\R$ ou $\C$), et $u$ un endomorphisme de $E$. On suppose $u$ nilpotent, c'est-\`a-dire qu'il existe un entier strictement positif $n$ tel que $u^n=0$.

\begin{enumerate}  
  \item Montrer que $u$ n'est pas inversible. 
  \item D\'eterminer les valeurs propres de $u$ et les sous-espaces propres associ\'es.
 \end{enumerate}
\finenonce{002579} 


\finexercice


\exercice{2580, delaunay, 2009/05/19}
\enonce{002580}{}
Soit $M$ la matrice de $\R^4$ suivante
$$M=\begin{pmatrix}0&1&0&0 \\  2&0&-1&0 \\  0&7&0&6 \\  0&0&3&0\end{pmatrix}$$

\begin{enumerate}  \item D\'eterminer les valeurs propres de $M$ et ses sous-espaces propres.
  \item Montrer que $M$ est diagonalisable.
  \item D\'eterminer une base de vecteurs propres et $P$ la matrice de passage.
  \item On a $D=P^{-1}MP$, pour $k\in\N$ exprimer $M^k$ en fonction de $D^k$, puis calculer $M^k$.
 \end{enumerate}
\finenonce{002580} 


\finexercice


\exercice{2594, delaunay, 2009/05/19}
\enonce{002594}{}
Soit $A$ la matrice suivante 
$$A=\begin{pmatrix}3&0&-1 \\  2&4&2 \\  -1&0&3\end{pmatrix}$$ 
\begin{enumerate}
 \item  D\'eterminer et factoriser le polyn\^ome caract\'eristique de $A$.
 \item  D\'emontrer que $A$ est diagonalisable et d\'eterminer une matrice $D$ diagonale et une matrice $P$ inversible telles  $A=PDP^{-1}$.
 \item  Donner en le justifiant, mais sans calcul, le polyn\^ome minimal de $A$.
 \item Calculer $A^n$ pour $n\in\N$.
\end{enumerate}
\finenonce{002594} 


\finexercice


\exercice{2595, delaunay, 2009/05/19}
\enonce{002595}{}
Soit $A$ la matrice suivante
$$A=\begin{pmatrix}1&1 \\  2&1\end{pmatrix}$$
\begin{enumerate}
 \item  Calculer le polyn\^ome caract\'eristique et d\'eterminer les valeurs propres de $A$.
 \item On note $\lambda_1>\lambda_2$ les valeurs propres de $A$, $E_1$ et $E_2$ les sous-espaces propres associ\'es. D\'eterminer une base $(\vec{\varepsilon_1},\vec{\varepsilon_2})$ de $\R^2$ telle que 
$\vec{\varepsilon_1}\in E_1$, $\vec{\varepsilon_2}\in E_2$, les deux vecteurs ayant des coordonn\'ees de la forme $(1,y)$.
 \item  Soit $\vec x$ un vecteur de $\R^2$, on note $(\alpha,\beta)$ ses coordonn\'ees dans la base
$(\vec{\varepsilon_1},\vec{\varepsilon_2})$. D\'emontrer que, pour $n\in\N$, on a
$$A^n\vec x=\alpha\lambda_1^n\vec{\varepsilon_1}+\beta\lambda_2^n\vec{\varepsilon_2}$$
 \item Notons $A^n\vec x=\begin{pmatrix}a_n \\  b_n\end{pmatrix}$ dans la base canonique de $\R^2$. Exprimer $a_n$ et $b_n$ en fonction de $\alpha$, $\beta$, $\lambda_1$ et $\lambda_2$. En d\'eduire que, si $\alpha\neq 0$, la suite ${\frac{b_n}{a_n}}$ tend vers $\sqrt2$ quand $n$ tend vers $+\infty$.
 \item  Expliquer, sans calcul, comment obtenir \`a partir des questions pr\'ec\'edentes une approximation de $\sqrt2$ par une suite  de nombres rationnels.
\end{enumerate}
\finenonce{002595} 


\finexercice


\exercice{2596, delaunay, 2009/05/19}
\enonce{002596}{}
Soit $P(X)$ un polyn\^ome de $\C[X]$, soit $A$ une matrice de $M_n(\C)$. On note $B$ la matrice : $B=P(A)\in M_n(\C)$.
\begin{enumerate}
 \item  D\'emontrer que si $\vec x$ est un vecteur propre de $A$ de valeur propre $\lambda$, alors $\vec x$ est un vecteur propre de $B$ de valeur propre $P(\lambda)$.

 \item  Le but de cette question est de d\'emontrer que les valeurs propres de $B$ sont toutes de la forme $P(\lambda)$, avec $\lambda$ valeur propre de $A$.

 Soit $\mu\in\C$, on d\'ecompose le polyn\^ome $P(X)-\mu$ en produit de facteurs de degr\'e $1$ :
$$P(X)-\mu=a(X-\alpha_1)\cdots(X-\alpha_r).$$
  \begin{enumerate}
   \item D\'emontrer que 
$$\det(B-\mu I_n)=a^n\det(A-\alpha_1I_n)\cdots\det(A-\alpha_rI_n).$$
   \item En d\'eduire que si $\mu$ est valeur propre de $B$, alors il existe une valeur propre $\lambda$ de $A$ telle que $\mu=P(\lambda)$.
  \end{enumerate}
 \item On note $S_A$ l'ensemble des valeurs propres de $A$, d\'emontrer que 
$$S_B=\{P(\lambda)/\  \lambda\in S_A\}.$$

 \item  Soient $\lambda_1,\dots,\lambda_r$ les valeurs propres de $A$ et soit $Q(X)$ le polyn\^ome :
$$Q(X)=(X-\lambda_1)\cdots(X-\lambda_r),$$
on note $C$ la matrice $C=Q(A)$. 
  \begin{enumerate}
   \item D\'emontrer que $S_C=\{0\}$.
   \item  En d\'eduire que le polyn\^ome caract\'eristique de $C$ est $(-1)^nX^n$ et que $C^n=0$.
  \end{enumerate}
\end{enumerate}
\finenonce{002596} 


\finexercice


\exercice{2761, tumpach, 2009/10/25}
\enonce{002761}{}
\begin{enumerate}
\item
\begin{enumerate}
\item Soit $f~:\mathbb{R}^2\rightarrow \mathbb{R}^2$ l'application lin\'eaire d\'efinie par 
$$
f\left(\begin{array}{c}x\\y\end{array}\right) =\frac{1}{5} \left(\begin{array}{c}3x + 4y\\ 4x - 3y\end{array}\right).
$$
\item \'Ecrire la matrice de $f$ dans la base canonique de $\mathbb{R}^2$. On la notera $A$.
\item Montrer que le vecteur $\vec{v}_1 = \left(\begin{array}{c}2 \\1\end{array}\right)$ est vecteur propre de $f$. Quelle est la valeur propre associ\'ee~?
\item Montrer que le vecteur $\vec{v}_2 = \left(\begin{array}{c}-1 \\2\end{array}\right)$ est \'egalement vecteur propre de $f$. Quelle est la valeur propre associ\'ee~?
\item Calculer graphiquement l'image du vecteur $\vec{v}_3 = \left(\begin{array}{c}1\\ 3\end{array}\right).$ Retrouver ce r\'esultat par le calcul.
\item Montrer que la famille $\{\vec{v}_1, \vec{v}_2\}$ forme une base de $\mathbb{R}^2$.
\item Quelle est la matrice de $f$ dans la base $\{\vec{v}_1, \vec{v}_2\}$~? On la notera $D$.
\item Soit $P$ la matrice  dont la premi\`ere colonne est le vecteur $\vec{v}_1$ et dont la deuxi\`eme colonne est le vecteur $\vec{v}_2$. Calculer $P^{-1}$.
\item Quelle relation y-a-t-il entre $A$, $P$, $P^{-1}$ et $D$~?
\item Calculer $A^n$, pour $n\in\mathbb{N}$.
\end{enumerate}
\item M\^eme exercice avec la matrice $A = \left(\begin{array}{cc}2 & -3\\ -1 & 0\end{array}\right)$ et les vecteurs $\vec{v}_1 = \left(\begin{array}{c}3 \\-1\end{array}\right)$, $\vec{v}_2 = \left(\begin{array}{c}1 \\1\end{array}\right)$ et $\vec{v}_3 = \left(\begin{array}{c}0 \\4\end{array}\right)$.
\end{enumerate}
\finenonce{002761}



\finexercice


\exercice{2762, tumpach, 2009/10/25}
\enonce{002762}{}
D\'eterminer le polyn\^ome caract\'eristique des matrices suivantes
$$
\left(\begin{array}{cc}0 & 1 \\1 & 0\end{array}\right),\quad
\left(\begin{array}{ccc}0 & 1 & 1\\1 & 0 & 1\\1 & 1 & 0\end{array}\right),\quad
\left(\begin{array}{cccc}0 & 1 & 1 & 1\\ 1 & 0 & 1 & 1\\ 1 & 1 & 0 & 1\\1 & 1 & 1 & 0\end{array}\right).
$$
\finenonce{002762}



\finexercice


\exercice{2763, tumpach, 2009/10/25}
\enonce{002763}{}
Rechercher les valeurs propres et vecteurs propres des matrices suivantes~:
$$
\left(\begin{array}{ccc} 1 & 0 & 0\\ 0 & 1 & 1\\ 0 & 1 & -1\end{array}\right), \quad
\left(\begin{array}{ccc} 1 & 0 & 4\\0 & 7 & -2\\ 4 & -2 & 0\end{array}\right),\quad
\left(\begin{array}{ccc} 1 & -1 & -1 \\ -1 & a^2 & 0 \\ -1 & 0 & a^2 \end{array} \right)
\quad(a\neq 0).
$$
\finenonce{002763}



\finexercice


\exercice{2764, tumpach, 2009/10/25}
\enonce{002764}{}
Trouver une matrice carr\'ee inversible $P$ telle que $B = PAP^{-1}$ soit diagonale, et \'ecrire la matrice $B$ obtenue, pour les matrices $A$ suivantes~:
$$
\left(\begin{array}{ccc}2 & 0 & 0\\0 & 1 & -4\\0 &-4 & 1\end{array}\right),\quad
\left(\begin{array}{ccc}1 & 0 & 1\\0 & 1 & 0\\1 & 0 & 1\end{array}\right),\quad
\left(\begin{array}{ccc}1 & 0 & 4\\0 & 7&-2\\4 & -2 & 0\end{array}\right).
$$
\finenonce{002764}



\finexercice


\exercice{2765, tumpach, 2009/10/25}
\enonce{002765}{DS mai 2008}
Soit la matrice 
$$
A = \left(\begin{array}{ccc}7 & 3 & -9\\-2 & -1 & 2\\2 & -1 & -4\end{array}\right)
$$
qui repr\'esente $f$, un endomorphisme de $\mathbb{R}^3$ dans la base canonique $\mathcal{B} = \{\vec{i}, \vec{j}, \vec{k}\}$.
\begin{enumerate}
\item
\begin{enumerate}
\item Montrer que les valeurs propres de $A$ sont $\lambda_1 = -2$, $\lambda_2 = 1$ et $\lambda_3 = 3$.
\item En d\'eduire que l'on peut diagonaliser $A$.
\end{enumerate}
\item 
\begin{enumerate}
\item
D\'eterminer une base $\mathcal{B}' = \{\vec{v}_1, \vec{v}_2, \vec{v}_3\}$ de vecteurs propres tels que la matrice de $f$ dans la base $\mathcal{B}'$ soit
$$
D = \left(\begin{array}{ccc}\lambda_1 & 0 & 0\\ 0 & \lambda_2 & 0\\ 0 & 0 &\lambda_3\end{array}\right).
$$
\item Pr\'eciser la matrice de passage $P$ de la base $\mathcal{B}$ \`a la base $\mathcal{B}'$ ; quelle relation lie les matrices $A$, $P$, $P^{-1}$ et $D$~?
\end{enumerate}
\item Montrer que pour tout entier $n\in\mathbb{N}$, on a $A^n = P D^{n} P^{-1}$.
\item Apr\`es avoir donn\'e $D^n$, calculer $A^n$ pour tout $n\in\mathbb{N}$.
\end{enumerate}
\finenonce{002765}



\finexercice


\exercice{2766, tumpach, 2009/10/25}
\enonce{002766}{DS mai 2008}
Soit la matrice $A = \left(\begin{array}{ccc}-3 & -2 & -2\\2 & 1 & 2\\ 2 & 2 & 1\end{array}\right).$
\begin{enumerate}
\item Calculer les valeurs propres de $A$.
\item \begin{enumerate}
\item Donner une base et la dimension de chaque sous-espace propre de $A$.
\item $A$ est diagonalisable ; justifier cette affirmation et diagonaliser $A$.
\end{enumerate}
\end{enumerate}
\finenonce{002766}



\finexercice


\exercice{2767, tumpach, 2009/10/25}
\enonce{002767}{}
On consid\`ere la matrice 
$$
A = \left(\begin{array}{cccc}1 & 0 & 0 & 0\\a & 1 & 0 & 0 \\a' & b & 2 & 0\\ a'' & b' & c & 2\end{array}\right).
$$
\`A quelles conditions les inconnue doivent-elles satisfaire pour que cette matrice soit diagonalisable~? Ces conditions \'etant remplies, fournir une base de vecteurs propres pour $A$.
\finenonce{002767}



\finexercice


\exercice{2777, tumpach, 2009/10/25}
\enonce{002777}{}
\begin{enumerate}
\item
Calculer les valeurs propres et les vecteurs propres de la matrice 
$$
A = \left(\begin{array}{ccc}0 & 2 & -2\\1 & -1 & 2\\1 & -3 & 4\end{array}\right).
$$
\item Calculer $A^n$ pour tout $n\in\mathbb{N}$.
\end{enumerate}
\finenonce{002777}



\finexercice


\exercice{3379, quercia, 2010/03/09}
\enonce{003379}{Valeurs propres de $AB$ et $BA$}

Soient $A \in \mathcal{M}_{n,p}(K)$ et $B \in \mathcal{M}_{p,n}(K)$. On note $C = I_n - AB$ et
$D = I_p - BA$.\par
($I_n,I_p = $ matrices unité d'ordres $n$ et $p$)

\begin{enumerate}
  \item Montrer que si $C$ est inversible, alors $D$ l'est aussi (résoudre $DX = 0$).
  \item Le cas échéant, exprimer $D^{-1}$ en fonction de $A,B,C^{-1}$.
    
  \item En déduire que $AB$ et $BA$ ont les mêmes valeurs propres non nulles.
    Examiner le cas de la valeur propre $0$ si $n = p$.
\end{enumerate}
\finenonce{003379}



\finexercice

\exercice{3499, quercia, 2010/03/10}
\enonce{003499}{Calcul de valeurs propres}

Chercher les valeurs propres des matrices :
\begin{enumerate}
  \item
$\begin{pmatrix}0      &\dots  &0      &1      \cr
          \vdots &       &\vdots &\vdots \cr
          0      &\dots  &0      &n-1    \cr
          1      &\dots  &n-1    &n      \cr \end{pmatrix}$.
\item
$\begin{pmatrix}0            &\sin\alpha &\sin2\alpha \cr
          \sin\alpha   &0          &\sin2\alpha \cr
          \sin2\alpha  &\sin\alpha &0           \end{pmatrix}$.

\end{enumerate}
\finenonce{003499}


\finexercice


\exercice{3500, quercia, 2010/03/10}
\enonce{003500}{Calcul de valeurs propres}

Soient $a_1,\dots,a_n \in \R$.

Chercher les valeurs et les vecteurs propres de la matrice
$A = \begin{pmatrix}    &        &        & a_1     \cr
                  &(0)     &        & \vdots  \cr
                  &        &        & a_{n-1} \cr
              a_1 &\dots   &a_{n-1} & a_n     \cr \end{pmatrix}$.
On distinguera les cas :

\begin{enumerate}
  \item $(a_1,\dots,a_{n-1}) \ne (0,\dots,0)$.
  \item $(a_1,\dots,a_{n-1}) = (0,\dots,0)$.
    
\end{enumerate}
\finenonce{003500}



\finexercice


\exercice{3501, quercia, 2010/03/10}
\enonce{003501}{Polynômes de Chebychev}

Soit $A = \begin{pmatrix} 0 &1      &       &(0)     \cr
                    1 &\ddots &\ddots          \cr
                      &\ddots &\ddots &1       \cr
                    (0)&      &1      &0       \cr \end{pmatrix} \in \mathcal{M}_n(\R)$.

\begin{enumerate}
  \item Calculer $D_n(\theta) = \det(A+(2\cos\theta) I)$ par récurrence.
    

  \item En déduire les valeurs propres de $A$.
\end{enumerate}
\finenonce{003501}



\finexercice


\exercice{3502, quercia, 2010/03/10}
\enonce{003502}{Matrice tridiagonale}

Déterminer les valeurs propres de la matrice
$A = \begin{pmatrix} 1  &-1     &       &        &(0)     \cr
               -1 &2      &-1     &        &        \cr
                  &\ddots &\ddots &\ddots  &        \cr
                  &       &-1     &2       &-1      \cr
               (0)&       &       &-1      &1       \cr\end{pmatrix} \in \mathcal{M}_n(\R)$.
\finenonce{003502}



\finexercice

\exercice{3506, quercia, 2010/03/10}
\enonce{003506}{Esem 91}

    Soit $C_{pq} = \begin{pmatrix} U_{pq} &(0) &U_{pq} \cr
			     (0)    &(0) &(0)	 \cr
			     U_{pq} &(0) &U_{pq} \cr \end{pmatrix} \in \mathcal{M}_n(\R)$
    où $U_{pq}$ est la matrice $p\times q$ dont tous les coefficients valent~1.
    Chercher les éléments propres de $C_{p,q}$.
    
\finenonce{003506}



\finexercice

\exercice{3508, quercia, 2010/03/10}
\enonce{003508}{Sommes par lignes ou colonnes constantes}

Soit $A \in \mathcal{M}_n(K)$ telle que la somme des coefficients par ligne est constante
($=S$). Montrer que $S$ est une valeur propre de $A$.

Même question avec la somme des coefficients par colonne.
\finenonce{003508}



\finexercice


\exercice{3509, quercia, 2010/03/10}
\enonce{003509}{Matrices stochastiques}

Soit $M=(m_{ij}) \in \mathcal{M}_n(\R)$ telle que :
$\begin{cases} \forall\ i,j,\ m_{ij} \ge 0 \cr
         \forall\ i,\ m_{i,1} + m_{i,2} + \dots + m_{i,n} = 1.\cr \end{cases}$
({\it matrice stochastique\/})

\begin{enumerate}
  \item Montrer que 1 est valeur propre de $M$.
  \item Soit $\lambda$ une valeur propre complexe de $M$.
    Montrer que $|\lambda| \le 1$
    (si $(x_1,\dots,x_n)\in \C^n$ est un vecteur propre associé, considérer le
    coefficient $x_k$ de plus grand module).
    Montrer que si tous les
    coefficients $m_{ij}$ sont strictement positifs alors $|\lambda| = 1  \Rightarrow  \lambda = 1$.
\end{enumerate}
\finenonce{003509}



\finexercice

\exercice{3511, quercia, 2010/03/10}
\enonce{003511}{$(X-a)P'$}

Soit $E =  K_n[X]$ et $u : E \to E,  P \mapsto {(X-a)P'.}$
Chercher les valeurs propres et les vecteurs propres de $u$.
\finenonce{003511}



\finexercice


\exercice{3512, quercia, 2010/03/10}
\enonce{003512}{$X(X-1)P' -2nXP$}

Soit $E =  K_{2n}[X]$ et $u : E \to E,  P \mapsto {X(X-1)P'-2nXP.}$
Chercher les valeurs propres et les vecteurs propres de $u$.

\finenonce{003512}



\finexercice


\exercice{3513, quercia, 2010/03/10}
\enonce{003513}{$X^3P \bmod (X-a)(X-b)(X-c)$}

Soient $\alpha, \beta, \gamma \in  K$ distincts, et
$\varphi : { K_2[X]} \to { K_2[X]}, P\mapsto R$
où $R$ est le reste de la division euclidienne de $X^3P$ par
$(X-\alpha)(X-\beta)(X-\gamma)$.
Chercher les valeurs et les vecteurs propres de $\varphi$.

\finenonce{003513}



\finexercice


\exercice{3514, quercia, 2010/03/10}
\enonce{003514}{$P(2-X)$}

Déterminer les éléments propres de l'endomorphisme
$\theta : { K[X]} \to { K[X]}, P \mapsto {P(2-X).}$
\finenonce{003514}



\finexercice


\exercice{3515, quercia, 2010/03/10}
\enonce{003515}{$P(X+1) - P'$}

Déterminer les éléments propres de l'endomorphisme
$\theta : { K[X]} \to { K[X]}, P \mapsto {P(X+1) - P'.}$

\finenonce{003515}


\finexercice


\exercice{3516, quercia, 2010/03/10}
\enonce{003516}{Eivp 91}

    Soit $f \in \mathcal{L}({\R_n[X]})$ qui à $P$ associe $(X-a)P'+P-P(a)$.
    Donner la matrice de $f$ dans la base $(X^k)_{0\le k\le n}$.
    Chercher $\Im f$, $\mathrm{Ker} f$ et les éléments propres de $f$.
\finenonce{003516}



\finexercice

\exercice{5655, rouget, 2010/10/16}
\enonce{005655}{***}
Soit $E=\Rr_{3}[X]$. Pour $P$ élément de $E$, soit $f(P)$ le reste de la division euclidienne de $AP$ par $B$ où $A= X^4-1$ et $B = X^4 - X$.

Vérifier que $f$ est un endomorphisme de $E$ puis déterminer $\text{Ker}f$, $\text{Im}f$ et les valeurs et vecteurs propres de $f$.
\finenonce{005655}


\finexercice


\exercice{5660, rouget, 2010/10/16}
\enonce{005660}{****}
Soit $E$ un $\Cc$-espace vectoriel de dimension finie non nulle.

Soient $u$ et $v$ deux endomorphismes de $E$ tels que $\exists(\alpha,\beta)\in\Cc^2/\;uv - vu =\alpha u + \beta v$. Montrer que $u$ et $v$ ont un vecteur propre en commun.
\finenonce{005660}


\finexercice


\exercice{5674, rouget, 2010/10/16}
\enonce{005674}{***}
$E =C^0(\Rr,\Rr)$. Pour $f$ élément de $E$, $\varphi(f)$ est l'application définie par :

\begin{center}
$\forall x\in\Rr^*$, $(\varphi(f))(x) =\frac{1}{x}\int_{0}^{x}f(t)\;dt$ si $x\neq 0$ et $(\varphi(f))(0) =f(0)$.
\end{center}

\begin{enumerate}
 \item  Montrer que $\varphi$ est un endomorphisme de $E$. 

\item  Etudier l'injectivité et la surjectivité de $\varphi$.

\item  Déterminer les éléments propres de $\varphi$.
\end{enumerate}
\finenonce{005674}


\finexercice

\section{201.02 Diagonalisation}


\exercice{1612, liousse, 2003/10/01}

\enonce{001612}{}
 Soient trois vecteurs $e_1,$ $e_2,$ 
$e_3$ formant une base de $\Rr^{3}.$ On note $T$ l'application 
lin\'eaire d\'efinie par $T(e_1)=T(e_3)=e_3$ et $T(e_2)=-e_1+e_2+e_3.$
\begin{enumerate}
  \item D\'eterminer le noyau de cette application lin\'eaire. Donner la matrice 
$A$ de $T$ dans la base donn\'ee.\\
 \item On pose $f_1=e_1-e_3,$ $f_2=e_1-e_2,$ $f_3=-e_1+e_2+e_3.$ Calculer 
$e_1,$ $e_2,$ $e_3$ en fonction de $f_1,$ $f_2,$ $f_3.$ Les vecteurs 
$f_1,$ $f_2,$ $f_3$ forment-ils une base de $\Rr^{3}$ ?\\
 \item Calculer $T(f_1),$ $T(f_2),$ $T(f_3)$ en fonction de $f_1,$ $f_2,$ 
$f_3.$ \'Ecrire la matrice $B$ de $T$ dans cette nouvelle base.\\
 \item On pose $P=\left (\begin{array}{rcl}1&1&-1\\0&-1&1\\-1&0&1
\end{array}\right ).$ V\'erifier que $P$ est inversible et calculer 
$P^{-1}.$ Quelle relation relie $A,$ $B,$ $P$ et $P^{-1}$ ? 
\end{enumerate}
\finenonce{001612}



\finexercice





\exercice{1614, liousse, 2003/10/01}

\enonce{001614}{}
 Soit $E$ un espace vectoriel 
de dimension $n$ et $f$ une 
application lin\'eaire de $E$ dans $E$.

\begin{enumerate}
  \item
Montrer que la condition $f^2=0$ est \'equivalente \`a 
$Im f \subset Ker f$. Quelle
condition v\'erifie alors le rang de $f$ ? 
On suppose dans la suite que $f^2=0$.

  \item  
 Soit $F$ un suppl\'ementaire de $Ker f$ dans $E$ et 
soit $(e_1, \ldots , e_r)$ une base de $F$.
Montrer que la famille des vecteurs 
$(e_1, \ldots , e_r,f(e_1), \ldots , f(e_r))$ est libre.
Montrer comment la compl\'eter si n\'ecessaire par 
des vecteurs de $Ker f$ pour obtenir une base de $E$.
Quelle est la matrice de $f$ dans cette base ?

  \item 
 Sous quelle condition n\'ecessaire et suffisante a-t-on $Im f= Ker f$ ?

  \item 
 Exemple. Soit $f$ une 
application lin\'eaire de $\Rr^3$ dans $\Rr^3$ de matrice 
dans la base canonique 
$M(f) = \left( \begin{array}{ccc} 1&0&1 \\ 2&0&2 \\ -1&0&-1 \\ \end{array}\right)$. 
Montrer que $f^2=0$.
D\'eterminer une nouvelle base dans laquelle la matrice de 
$f$ a la forme indiqu\'ee dans la question 2).
\end{enumerate}
\finenonce{001614}



\finexercice



\exercice{1615, liousse, 2003/10/01}

\enonce{001615}{}
 Soit $A=\left (\begin{array}{cc}
1&4\\2&3\end{array}\right ).$ Trouver les valeurs propres de 
$A$ et les sous-espaces propres correspondant. En d\'eduire une matrice 
inversible $P$ telle que $P^{-1}AP$ soit diagonale.
\finenonce{001615}



\finexercice



\exercice{1616, liousse, 2003/10/01}

\enonce{001616}{}
 Soit $A=\left (\begin{array}{ccc}
4&1&-1\\2&5&-2\\1&1&2\end{array}\right ).$ Diagonaliser $A.$ 
\finenonce{001616}



\finexercice



\exercice{1617, liousse, 2003/10/01}

\enonce{001617}{}
Soit $A=\left (\begin{array}{ccc}
1&1&1\\1&1&1\\1&1&1\end{array}\right ).$ Trouver, sans calculer 
le polyn\^ome caract\'eristique, les valeurs propres de $A.$ 
Cette matrice est-elle diagonalisable ?
\finenonce{001617}



\finexercice



\exercice{1618, liousse, 2003/10/01}

\enonce{001618}{}
 On consid\`ere les matrices suivantes
$$A=\left (\begin{array}{ccc}
3&1&1\\2&4&2\\1&1&3\end{array}\right )\hbox{  }
B=\left (\begin{array}{ccc}
1&2&2\\1&2&-1\\-1&1&4\end{array}\right )\hbox{  }
C=\left (\begin{array}{ccc}
1&1&0\\0&1&0\\0&0&1\end{array}\right )$$ 
Ces matrices sont-elles diagonalisables ? Si oui, les r\'eduire.
\finenonce{001618}



\finexercice



\exercice{1619, liousse, 2003/10/01}

\enonce{001619}{}
 Soit $n$ un entier strictement 
sup\'erieur \`a 1. Soit $A$ une matrice $n\times n$ 
telle que $A^{n}=0$ et $A^{n-1}\ne 0.$ Soit $x_0$ un vecteur de 
$\Rr^{n}$ tel que $A^{n-1}x_0\ne 0.$ Montrer que 
$(x_0,Ax_0,A^{2}x_0,\cdots,A^{n-1}x_0)$ est une base de 
$\Rr^{n}.$ Comment s'\'ecrit la matrice $A$ dans cette base ?

Application : on pose $A=\left(\begin{array}{ccc}2&1&2\\-1&-1&-1\\
-1&0&-1\end{array}\right).$ Calculer $A^{3}$ et donner une base 
de $\Rr^{3}$ dans laquelle $A$ a une forme simple.
\finenonce{001619}



\finexercice



\exercice{1620, liousse, 2003/10/01}

\enonce{001620}{}
On consid\`ere la matrice 
$$M=\left (\begin{array}{ccc}
3&2&1\\-1&0&-1\\-1&-1&1\end{array}\right )$$ 
Est-elle diagonalisable ? Justifier. \'Ecrire alors $M$ sous 
une forme plus simple.
\finenonce{001620}



\finexercice



\exercice{1621, liousse, 2003/10/01}

\enonce{001621}{}
Soit $T$ l'application lin\'eaire de $\mathbb R^3$ dans  $\mathbb R^3$ d\'efinie par sa matrice $A$
dans la base canonique $(e_1,e_2,e_3)$ de $\mathbb R^3$ :
 $$ A= {\begin{pmatrix} 1 & 2 & 0 \\ 1 & 3& -1
\\ 1 & -1& 3\end{pmatrix}}.$$
\begin{enumerate}
\item Donner un base de Ker $T$ et Im$T$.

\item
\begin{enumerate}
\item Calculer le polyn{\^o}me caract\'eristique de $T$, puis ses valeurs propres.
\item Justifier, sans calcul, que $T$ soit diagonalisable et
 \'ecrire une matrice diagonale semblable \`a $A$ .
\item Calculer une base de $\mathbb R^3$ form\'ee de vecteurs propres de $T$.
\end{enumerate}

\item Soient $f_1= -2e_1+e_2+e_3$ , $f_2=e_1+e_2+e_3$ et $f_3 = 2e_1+3e_2 -e_3$ 
trois vecteurs de $\mathbb R^3$.
\begin{enumerate}
\item Justifier que $(f_1,f_2,f_3)$ est une base de $\mathbb R^3$ et
 \'ecrire la matrice $P$ de passage de la base $(e_1,e_2,e_3)$ \`a  la base  $(f_1,f_2,f_3)$.
\item Calculer  $P^{-1}$.
\item Ecrire la matrice $D$ de $T$ dans la base $(f_1,f_2,f_3)$ .
\end{enumerate}

\item Quelle relation relie $A^3$, $D^3$, $P$ et $P^{-1}$ ? En d\'eduire  $A^3$.
\end{enumerate}
\finenonce{001621}



\finexercice



\exercice{1622, barraud, 2003/09/01}

\enonce{001622}{}
Lorsque c'est possible, diagonaliser les matrices suivantes :

$$
\begin{pmatrix}
  2 & -1 &  1 \\
  1 &  0 & -1 \\
  2 & -2 &  1
\end{pmatrix}
\qquad
\begin{pmatrix}
  3 & -3 & -4 & -1 \\
  0 &  2 &  0 & -1 \\
  2 & -4 & -3 &  0 \\
  0 &  2 &  0 & -1
\end{pmatrix}
\qquad
\begin{pmatrix}
  1 &  0 &  0 &  0 \\
  1 &  1 & -1 &  1 \\
  2 & -1 &  1 &  1 \\
  3 & -1 & -1 &  3
\end{pmatrix}
\qquad
\begin{pmatrix}
  3 & -1 &  7 & -14\\
  4 & -1 &  7 & -15\\
  0 &  0 &  3 & -4 \\
  0 &  0 &  2 & -3
\end{pmatrix}
$$
\finenonce{001622}



\finexercice



\exercice{1623, barraud, 2003/09/01}

\enonce{001623}{}
  Pour quelles valeurs de $(a,b,c)\in\C^{2}$ la matrice $A=
  \begin{pmatrix}
    1&a&1&0\\
    0&1&b&2\\
    0&0&2&c\\
    0&0&0&2
  \end{pmatrix} 
$ est-elle diagonalisable~? On ne cherchera pas à réduire explicitement
$A$.
\finenonce{001623}



\finexercice



\exercice{1624, barraud, 2003/09/01}

\enonce{001624}{}
Soit $u$ l'application suivante :
$$
  u:
  \begin{array}{rcl}
    \R_{2}[X] & \rightarrow  & \R_{2}[X]  \\[2ex]
    P         & \mapsto & (2X+1)P-(X^{2}-1)P'
  \end{array}
$$
Montrer que $u$ est bien d\'{e}finie et lin\'{e}aire. 
D\'{e}terminer les valeurs propres de $u$, et, si c'est possible, diagonaliser $u$.
\finenonce{001624}



\finexercice



\exercice{1625, barraud, 2003/09/01}

\enonce{001625}{}
Soit $A\in\mathcal{M}_{n}(\R)$. Montrer que si $\lambda$ est une valeur propre
complexe de $A$, alors $\bar{\lambda}$ est aussi une valeur propre de
$A$. De m\^{e}me, montrer que si $x$ est un vecteur propre complexe de
$A$, alors $\bar{x}$ (o\`{u} $\bar{x}$ d\'{e}signe le vecteur dont les
composantes sont les conjugu\'{e}es des composantes de $x$) est aussi un
vecteur propre complexe de $A$.

Diagonaliser
 $
  A=
  \begin{pmatrix}
    -1 &  1 &  0 \\
     0 & -1 &  1 \\
     1 &  0 & -1
  \end{pmatrix}
 $.
\finenonce{001625}



\finexercice



\exercice{1626, barraud, 2003/09/01}

\enonce{001626}{}
Soit $A_{t}$ la matrice
 $
 A_{t}=
 \begin{pmatrix}
     t   &    1   & \cdots &    1   \\
     1   &    t   & \ddots & \vdots \\
  \vdots & \ddots & \ddots &    1   \\
     1   & \cdots &    1   &    t
 \end{pmatrix}
 $.
Sans calculer le polyn\^{o}me caract\'{e}ristique de $A_{t}$, montrer que $(t-1)$ est valeur
propre. D\'{e}terminer l'espace propre associ\'{e}. Que dire de la multiplicit\'{e} de la valeur
propre $(t-1)$ ? En d\'{e}duire le spectre de $A_{t}$. $A_{t}$ est-elle diagonalisable ?
\finenonce{001626}



\finexercice



\exercice{1627, barraud, 2003/09/01}

\enonce{001627}{}
Pour quelles valeurs de $a$, $b$ et $c$ les matrices suivantes sont-elles diagonalisables
?
$$
 \begin{pmatrix}
   1 & a & 1 \\
   0 & 1 & b \\
   0 & 0 & c
 \end{pmatrix}
\hspace{3cm}
 \begin{pmatrix}
   0 & 0 & a \\
   0 & 0 & b \\
   a & b & c
 \end{pmatrix}
$$
\finenonce{001627}



\finexercice



\exercice{1628, barraud, 2003/09/01}

\enonce{001628}{}
Soient $u$ et $v$ deux endomorphismes diagonalisables de $E$, qui commutent (c'est \`{a} dire
tels que $u\circ v=v\circ u$). On note $\lambda_{1},\ldots,\lambda_{p}$ (resp.
$\mu_{1},\ldots,\mu_{q}$) les valeurs propres de $u$ (resp. de $v$), et
$F_{1},\ldots,F_{p}$ les espaces propres associ\'{e}s (resp. $G_{1},\ldots,G_{q}$).

\begin{enumerate}
\item
Montrer que chaque $G_{j}$ (resp. $F_{i}$) est stable par $u$ (resp. $v$) (c'est \`{a} dire
que $u(G_{j})\subset G_{j}$).
\item
On pose $H_{ij}=F_{i}\cap G_{j}$. Soit $i\in\{1,\ldots,p\}$. Montrer que $F_{i}$ est la
somme directe des espaces $(H_{ij})_{1\leq j\leq q}$.

\item
En d\'{e}duire l'\'{e}nonc\'{e} suivant : {\sl Lorsque deux endomorphismes diagonalisables $u$ et $v$
commutent, il existe une base form\'{e}e de vecteurs propres communs \`{a} $u$ et \`{a} $v$ (en
d'autres termes, $u$ et $v$ sont diagonalisables simultan\'{e}ment dans la m\^{e}me base).}
\end{enumerate}
\finenonce{001628}



\finexercice



\exercice{1629, barraud, 2003/09/01}

\enonce{001629}{}
Les matrices suivantes sont-elles diagonalisables,
triangularisables ? Si oui, les réduire.
$$
A_{1}=\begin{pmatrix}
        3 &-1 & 1 \\
        2 & 0 & 1 \\
        1 &-1 & 2
      \end{pmatrix}
\qquad
A_{2}=\begin{pmatrix}
        3 & 2 &-2 \\
       -1 & 0 & 1 \\
        1 & 1 & 0
      \end{pmatrix}
\qquad
A_{3}=\begin{pmatrix}
       13 &-5 &-2 \\
       -2 & 7 &-8 \\
       -5 & 4 & 7
      \end{pmatrix}
\qquad
$$
\finenonce{001629}



\finexercice



\exercice{1630, barraud, 2003/09/01}

\enonce{001630}{}
Soit $f$ un endomorphisme diagonalisable d'un espace vectoriel $E$ et $P$
un polynôme. Montrer que $P(f)$ est diagonalisable.
\finenonce{001630}



\finexercice



\exercice{1631, barraud, 2003/09/01}

\enonce{001631}{}
Soit $P_{0}$ un polynôme de $\R_{n}[X]$, et $f$ l'application
suivante :
$$
 f: \begin{array}{rcl}
      \R_{n}[X] & \rightarrow  & \R_{n}[X] \\[1ex]
      P         & \mapsto & R=
      \text{ reste de la division euclidienne de $P$ par $P_{0}$}
    \end{array}
$$
A l'aide d'un polynôme annulateur de $f$, montrer que $f$ est
diagonalisable.
\finenonce{001631}



\finexercice



\exercice{1632, barraud, 2003/09/01}

\enonce{001632}{}
Soit $\alpha$ et $\beta$ deux réels, et $A$ la matrice suivante :
$$
  A=\begin{pmatrix}
            1 & -\alpha &  -\alpha &       1 \\
      1-\beta &  \alpha & \alpha-1 &  -\beta \\
        \beta & -\alpha & 1-\alpha & 1+\beta \\
            0 &  \alpha &   \alpha &       0
    \end{pmatrix}
$$
A quelle condition sur $\alpha$ et $\beta$, $A$ est-elle diagonalisable ?

On suppose $\alpha=0$ et $\beta=0$. Vérifier que $A(A-I)=0$. En déduire
$A^{n}$ et $(A+I)^{-1}$ .

\finenonce{001632}



\finexercice



\exercice{1633, barraud, 2003/09/01}

\enonce{001633}{}
Les matrices suivantes sont-elles diagonalisables, triangularisables,
sur $\R$ et sur $\C$ ?

Lorsqu'elles sont diagonalisables, donner une matrice diagonale semblable.

$$
 A=
 \begin{pmatrix}
   3 & 2 &-1 &-2 \\
   1 & 3 &-1 &-1 \\
   2 & 2 & 0 &-2 \\
   1 & 2 &-1 & 0
 \end{pmatrix}
\qquad
 B=  
 \begin{pmatrix}
   1 & 2 &-1 & 1 \\
   1 & 1 & 1 &-1 \\
   3 &-4 & 5 &-3 \\
   0 & 0 & 0 & 2
 \end{pmatrix}
\qquad
 C=  
 \begin{pmatrix}
  -1 & 0 & 1 \\
   1 &-1 & 0 \\
  -4 & 2 & 2 
 \end{pmatrix}
$$

Réduire explicitement $A$ et $C$.

\finenonce{001633}



\finexercice



\exercice{1634, barraud, 2003/09/01}

\enonce{001634}{}
On considère un endomorphisme $f$ d'un $\C$ espace vectoriel $E$ de
dimension finie $n$, tel que $f^2$ est diagonalisable. Le but de cet
exercice est de démontrer que :
$$
  f\text{ est diagonalisable }\Leftrightarrow \mathrm{Ker} f=\mathrm{Ker} f^2
$$ 

\begin{enumerate}
\item
  On suppose que $f$ est diagonalisable. On note
  $\alpha_{1},...,\alpha_{r}$ les valeurs propres (distinctes) de
  $A$, et $E_{1},...,E_{r}$ les espaces propres associés.
  \begin{enumerate}
  \item
    Montrer que si $\mathrm{Ker} f=\{0\}$ alors $\mathrm{Ker} f^{2}=\{0\}$.
  \item
    On suppose maintenant que $\mathrm{Ker} f\neq\{0\}$. On note
    $\alpha_{\alpha_{1}},...,\alpha_{\alpha_{r}}$ les autres valeurs
    propres de $f$, et $E_{0},...,E_{r}$ ses espaces propres. En
    utilisant que $E=E_{0}\oplus E_{1}\oplus ...\oplus E_{r}$, montrer
    que si $f^{2}(x)=0$ alors $f(x)=0$. En déduire que $\mathrm{Ker} f=\mathrm{Ker} f^2$.
  \end{enumerate}
\item
  On suppose que $\mathrm{Ker} f=\mathrm{Ker} f^2$.
  \begin{enumerate}
  \item 
    Montrer que si $\mu$ est une valeur propre de $f$, alors $\mu^2$
    est une valeur propre de $f^2$.
    \begin{enumerate}
    \item 
      Soit $\lambda$ une valeur propre non nulle de $f^2$, et $\mu$ et
      $-\mu$ ses deux racines complexes. Montrer que
      $$
        \mathrm{Ker}(f-\mu\mathrm{id})\subset\mathrm{Ker}(f^2-\lambda\mathrm{id})
        \quad\text{ et que }\quad
        \mathrm{Ker}(f+\mu\mathrm{id})\subset\mathrm{Ker}(f^2-\lambda\mathrm{id}).
      $$ 
    \item
      Montrer que
      $$
        \mathrm{Ker}(f^2-\lambda\mathrm{id})=\mathrm{Ker}(f-\mu\mathrm{id})\oplus\mathrm{Ker}(f+\mu\mathrm{id})
      $$
      (remarquer que 
      $\forall y\in \mathrm{Ker} f^{2}\;\;
        y=\frac{1}{2\mu}((f+\mu\mathrm{id})(y)-(f-\mu\mathrm{id})(y))$).
   \end{enumerate}
   \item
     Montrer (avec soin) que $f$ est diagonalisable.
  \end{enumerate}
\end{enumerate}
\finenonce{001634}



\finexercice



\exercice{1635, barraud, 2003/09/01}

\enonce{001635}{}
La matrice suivante est-elle diagonalisable, triangularisable~? Effectuer
explicitement la réduction.
  $$
  A=
  \begin{pmatrix}
     3& 2& 4\\ 
    -1& 3&-1\\
    -2&-1&-3
  \end{pmatrix}
  $$

\finenonce{001635}



\finexercice



\exercice{1636, barraud, 2003/09/01}

\enonce{001636}{}
  Soit $J=\left(
  \begin{smallmatrix}
    \frac{1}{2}&\frac{1}{2}  \\[.5ex]
    \frac{1}{2}&\frac{1}{2}
  \end{smallmatrix}\right)
 $ et $A=
  \begin{pmatrix}
    0 &\vline& J  \\ \hline
    J &\vline&0
  \end{pmatrix}$. 
  Calculer $A^{2}$, puis $A^{3}$. A l'aide d'un polynôme annulateur de
  $A$, montrer que $A$ est diagonalisable.

  Sans chercher à calculer le polynôme caractéristique de $A$, donner un
  ensemble fini contenant toutes les valeurs propres de $A$, puis donner
  les valeurs valeurs propres elles mêmes ainsi que leurs multiplicités. En
  déduire le polynôme caractéristique de $A$.
\finenonce{001636}


\finexercice



\exercice{1637, barraud, 2003/09/01}

\enonce{001637}{}
  On considère une matrice $A\in\mathcal{M}_{nn}(\C)$ et l'application $\phi_{A}$
  définie par~:
  $$
  \phi_{A}:
  \begin{array}{ccc}
    \mathcal{M}_{nn}(\C) & \rightarrow & \mathcal{M}_{nn}(\C)\\
    B           &\mapsto &AB
  \end{array}
  $$
\begin{enumerate}
\item

    Montrer que $\phi_{A}$ est linéaire.

    Le but de l'exercice est de montrer que $\phi_{A}$ est diagonalisable
    si et seulement si $A$ est diagonalisable.
\item

    Calculer $\phi_{A}^{2}(B)$, puis $\phi_{A}^{k}(B)$ pour $k\in\N$. En
    déduire que si $P$ est un polynôme, alors $P(\phi_{A})=\phi_{P(A)}$.
\item

    En déduire que $P$ est un polynôme annulateur de $A$ si et seulement
    si $P$ est un polynôme annulateur de $\phi_{A}$.
\item

    Montrer que $\phi_{A}$ est diagonalisable si et seulement si $A$
    l'est.

\end{enumerate}
\finenonce{001637}



\finexercice



\exercice{1638, barraud, 2003/09/01}

\enonce{001638}{}
A $n$ nombres complexes $(a_{1},...,a_{n})\in\C^{n}$ avec $a_{2}\neq 0$,
on associe la matrice $A_{n}=
\begin{pmatrix}
  a_{1} &a_{2} &\cdots          &a_{n} \\
  a_{2} &      &                &      \\        
  \vdots&      &\text{\huge{0}} &      \\    
  a_{n} &      &                &     
\end{pmatrix}
$.

\begin{enumerate}
\item 
  Quel est le rang de $A_{n}$. Qu'en déduit-on pour le polynôme
  caractéristique $\chi_{n}$ de $A_{n}$~?
  
\item
  Calculer $\chi_{2}$, $\chi_{3}$.
  
\item
  On pose $b_{n}=a_{2}^{2}+\cdots+a_{n}^{2}$. Par récurrence, montrer
  que $\chi_{n}=(-X)^{n-2}(X^{2}-a_{1}X-b_{n})$.
  
\item
  Si $b_{n}=0$, $A_{n}$ est-elle diagonalisable~?
  
\item
  Si $b_{n}\neq0$, à quelle condition $A_{n}$ est-elle
  diagonalisable~?
\end{enumerate}
\finenonce{001638}



\finexercice



\exercice{1639, barraud, 2003/09/01}

\enonce{001639}{}
  Soit $A$ la matrice $A=
  \begin{pmatrix}
    1&-1&-1\\
   -1& 1&-1\\
   -1&-1& 1
  \end{pmatrix}
  $. Calculer $^{t}\!A$. La matrice $A$ est-elle diagonalisable~?

  Trouver une matrice $P$ orthogonale telle que $P^{-1}AP$ soit
  diagonale.
\finenonce{001639}



\finexercice



\exercice{1640, barraud, 2003/09/01}

\enonce{001640}{}
Soit $E$ un espace vectoriel de dimension finie, et $u$ un endomorphisme
de $E$ tel que $u^{p}=0$ pour un certain entier $p$. Quelles sont les
valeurs propres de $u$. A quelle condition $u$ est-il diagonalisable~?
Montrer que $u^{n}=0$.
\finenonce{001640}



\finexercice



\exercice{1641, barraud, 2003/09/01}

\enonce{001641}{}
  Déterminer les valeurs propres des matrices suivantes. Sont-elles
  diagonalisables, triangularisables~?

$$
A=\begin{pmatrix}
  3 &0 &0 \\
  2 &2 &0 \\
  1 &1 &1
  \end{pmatrix}
\qquad
B=\begin{pmatrix}
  2 &-2 &1 \\
  3 &-3 &1 \\
 -1 &2  &0 
  \end{pmatrix}
$$
A l'aide du polynôme caractéristique de $B$, calculer $B^{-1}$.

\finenonce{001641}



\finexercice



\exercice{1642, barraud, 2003/09/01}

\enonce{001642}{}
  Soit $A$ la matrice $A=
  \begin{pmatrix}
    1&-1&-1\\
   -1& 1&-1\\
   -1&-1& 1
  \end{pmatrix}
  $.
  \begin{enumerate}
  \item 
    Calculer $^{t}\!A$. La matrice $A$ est-elle diagonalisable~? 
  \item
    Diagonaliser $A$.
  \item
    Diagonaliser $A$ dans une base orthonormée (pour le produit scalaire
    usuel de $\R^{3}$).
  \end{enumerate}

\finenonce{001642}



\finexercice



\exercice{1643, barraud, 2003/09/01}

\enonce{001643}{}
Dans l'espace vectoriel $\R_{3}[X]$, on considère l'application linéaire
suivante~:
$$
  u~:
  \begin{array}{rcl}
    \R_{3}[X] & \rightarrow  & \R_{3}[X] \\[2ex]
    P         & \mapsto & P (0)X^{3}
                        +P'(0)X^{2}
                        +\frac{1}{2}P''(0)X
                        +\frac{1}{6}P'''(0)
  \end{array}
$$

\begin{enumerate}
\item   Ecrire la matrice $A$ de $u$ dans la base canonique. Calculer  $A^{2}$.
\item  $u$ est-elle diagonalisable~? Si oui, donner une base de $\R^{3}[X]$
  formée de vecteurs propres de $u$.
\end{enumerate}
\finenonce{001643}



\finexercice



\exercice{1644, barraud, 2003/09/01}

\enonce{001644}{}
On considère un réel $\alpha$ et l'application $T_{\alpha}$ suivante :
$$
T_{\alpha} :
\begin{array}{ccl}
  \R[X] & \rightarrow  & \R[X] \\
  P     & \mapsto & X(X-1)P''+(1+\alpha X)P'
\end{array}
$$

\begin{enumerate}
\item 
Montrer que pour tout entier $n>0$, la restriction de $T_{\alpha}$ à
$\R_{n}[X]$ défini un endomorphisme de $\R_{n}[X]$.

\item
On suppose pour cette question que $n=3$.
\begin{enumerate}
\item 
Ecrire la matrice de $T_{\alpha}$ dans la base
$(1,X,X^{2},x^{3})$. 

\item
Déterminer les valeurs propores de $T_{\alpha}$. On les note
$\lambda_{0},\lambda_{1},\lambda_{2},\lambda_{3}$.

\item
Déterminer les valeurs de $\alpha$ pour lesquelles $T_{\alpha}$ a des
valeurs propres multiples.

\item
Donner un vecteur propre de $T_{\alpha}$ pour chaque valeur propre,
lorsque $\alpha=-1$, puis $\alpha=-4$. L'endomorphisme $T_{-4}$ est-il
diagonalisable ?
\end{enumerate}

\item 
On suppose maintenant $n>3$.
\begin{enumerate}
\item 
Ecrire la matrice de $T_{\alpha}$ dans la base
$(1,X,X^{2},...,X^{n})$. 

\item
Déterminer les valeurs propores de $T_{\alpha}$. On les note
$\lambda_{0},\lambda_{1},...,\lambda_{n}$.

\item
Déterminer les valeurs de $\alpha$ pour lesquelles $T_{\alpha}$ a des
valeurs propres multiples. Dans chaque cas, donner la liste des valeurs
propres avec leurs multiplicités

\item
Déterminer la dimension de $\mathrm{Ker} T_{\alpha}$ et de $\mathrm{Im} T_{\alpha}$
lorsque $\alpha\notin\{1-n,...,-1,0\}$. 

\item
Déterminer $\mathrm{Ker} T_{\alpha}$ pour $\alpha=-1$, puis
$\alpha=0$. L'endomorphisme $T_{0}$ est-il diagonalisable ?

\item
Lorsque $\alpha=p-1$ avec $p\in\{1,...,n\}$, donner un polynôme $P$ de
degré inférieur ou égal à $n$ tel que $T_{\alpha}(P)=0$. En déduire
$\mathrm{Ker} T_{\alpha}$. Préciser sa dimension.

\item
Soit $\lambda_{k}$ une valeur propre simple de $T_{\alpha}$. Donner un
vecteur propre de $T_{\alpha}$ associé à $\lambda_{k}$.
\end{enumerate}
\end{enumerate}
\finenonce{001644}



\finexercice



\exercice{1645, legall, 1998/09/01}

\enonce{001645}{}
Soient $\R^n$ euclidien, $f \in \textbf{O}_n(\R)$. Montrer que $f$
est diagonalisable si et seulement si $f$ est une sym\'etrie orthogonale.
\finenonce{001645}



\finexercice



\exercice{1646, legall, 1998/09/01}

\enonce{001646}{}
Diagonaliser en base orthonormale les matrices suivantes :
$$A = \left(\begin{array}{cccc}0&\ldots&0&a_1\\
\vdots&\vdots&\vdots&\vdots\\ 0&\ldots&0&a_{n-1}\\
a_1&\ldots&a_{n-1}&a_n\end{array}
\right), a_i \in \R;
B = \left(\begin{array}{cccc} a&b& & \\ b&\ddots&\ddots&\\
&\ddots&\ddots&b\\ & & b& a\end{array}\right), a,b \in \R.$$
Peut-on d\'eterminer $a, b$ tels que $B$ soit la matrice d'un produit scalaire ?
\finenonce{001646}



\finexercice



\exercice{1647, legall, 1998/09/01}

\enonce{001647}{}
Montrer que si $A$ est une matrice sym\'etrique r\'eelle, alors $A + iI$
est inversible.
\finenonce{001647}



\finexercice



\exercice{1648, bodin, 1999/11/01}

\enonce{001648}{}
Soit $f$ un endomorphisme de $\C^3$ dont la
matrice par rapport \`a la base canonique est :
$$
M=
\left(
\begin{array}{ccc}
2 & -1 & 1 \\
-1 & k & 1 \\
1 &1 & 2
\end{array}
\right),
\; \mbox{o\`u} \; k\in\C.
$$
\emph{(a)} D\'eterminer, suivant les valeurs de
$k$, la dimension du noyau de $f$.  \\
\emph{(b)} Montrer que $M$ admet une valeur propre
r\'eelle enti\`ere ind\'ependante de $k$, et
calculer toutes les valeurs propres de $M$.  \\
\emph{(c)} Indiquer toutes les valeurs de $k$ pour
lesquelles on obtient des valeurs propres
multiples.  Pour quelles valeurs de ces $k$ la
matrice $M$ est-elle semblable \`a une matrice
diagonale?
\finenonce{001648}



\finexercice



\exercice{1649, legall, 1998/09/01}

\enonce{001649}{}
Soit $  A\in M_n(\Rr)  $ telle que $  A^2=-I  .$
\begin{enumerate}
    \item Montrer que $  n   $ est pair, $  n=2p  .$
    \item Calculer $  Sp_{\R}(A)  $ et montrer $  Sp_{\C}(A)= \{ i,-i\}   .$ Pour quelle raison $  A  $ est elle diagonalisable sur $  \C  $~?
    \item Montrer que si $  \{ y_1 , \ldots y_k \}   $ est une base de $  E_i  , $
alors $  \{ \overline{y_1} , \ldots \overline{y_k} \}   $ est une base de $  E_{-i}  .$ Quelle est donc la
valeur de $  k  $~?
    \item D\'emontrer que $  A   $ est semblable (dans $  M_n (\R )  $) \`a une matrice
diagonale par blocs dont chacun des blocs diagonaux est $
\begin{pmatrix}0 & -1\cr 1 & 0 \cr \end{pmatrix}  .$ (on pourra utiliser la question 3.)
\end{enumerate}
\finenonce{001649}



\finexercice



\exercice{1650, legall, 1998/09/01}

\enonce{001650}{}
Soient $\ M  $ et $  N \in M_n (\mathbb{K} )  .$ On note $  \varphi _M \in \mathcal{L}
(M_n (\mathbb{K} ))  $ l'application $  N\mapsto  MN-NM   .$
\begin{enumerate}
    \item Soient $  A= \begin{pmatrix}3 &
-4\cr
                                      2 & -3 \cr \end{pmatrix} $ et $  B= \begin{pmatrix}1 & 2\cr
                                      0 & 1 \cr \end{pmatrix} .$ Diagonaliser $  A  $ et
montrer que $  B   $ n'est pas diagonalisable.
    \item Montrer que si $  N  $ est un vecteur propre associ\'e \`a une
valeur propre non nulle $  \lambda   $ de $  \varphi _M   $ alors $  N   $ est nilpotente. (on
pourra \'etablir que pour tout $  k \in \N   :   MN^k-N^kM=k\lambda N^k  $.)
    \item Montrer que l'identit\'e n'appartient pas \`a l'image de $  \varphi _M   .$ (utiliser la trace.)
    \item Soit $  D=\begin{pmatrix}1 & 0  \cr
                                      0 & -1 \cr \end{pmatrix} .$ Diagonaliser $  \varphi _D
  $ puis $  \varphi _A  .$ Montrer que
$  \varphi _B  $ n'est pas diagonalisable.
    \item Montrer que si $  M  $ est diagonalisable, $  \varphi _M   $ est diagonalisable.
    \item Etablir la r\'eciproque lorsque $  M  $ a au moins une valeur propre.
\end{enumerate}
\finenonce{001650}



\finexercice



\exercice{1651, legall, 1998/09/01}

\enonce{001651}{}
Soit $  E  $ un $  \mathbb{K}$-espace vectoriel. Une application
$  p \in \mathcal{L} (E)  $ est nomm\'ee projecteur lorsque $  p^2=p  .$
\begin{enumerate}
    \item Montrer que si $  p  $ est un projecteur $  1-p  $ est un projecteur. Montrer que $  \hbox{Im}(p)\oplus \hbox{Ker}(p)=E  .$
    \item On suppose que $  \mathbb{K} = \R   .$ Soient $  p  $ et $ q   $ deux projecteurs
tels que $  p+q
$ soit aussi un projecteur. Montrer que~:
    \begin{enumerate}
        \item $  pq=qp=0  .$
        \item $  \hbox{Im}(p+q)=\hbox{Im}(p)+\hbox{Im}(q)  .$
        \item $  \hbox{Ker}(p+q)=\hbox{Ker}(p)\cap \hbox{Ker}(q)  .$
    \end{enumerate}
On suppose d\'esormais $  E   $ de dimension finie et $  \mathbb{K} =\R   .$
    \item Montrer que tout projecteur est diagonalisable et que deux projecteurs sont semblables si et
seulement si ils ont m\^eme trace.
    \item Montrer que toute matrice diagonalisable est combinaison lin\'eaire de projecteurs.
\end{enumerate}
\finenonce{001651}



\finexercice



\exercice{1652, legall, 1998/09/01}

\enonce{001652}{}
Soient $  E  $ un $  \mathbb{K} $-espace vectoriel de dimension finie,
$  P\in {\mathbb{K}}[X]  $ et $  u \in \mathcal{L} (E)  .$ On note $  P (Sp(u))=\{ P(\lambda )  ;   \lambda \in Sp(u)\}   .$
\begin{enumerate}
    \item On suppose que $  u   $ est diagonalisable. Montrer que $  P (Sp(u))= Sp(P(u))  .$
    \item Montrer, dans le cas g\'en\'eral, $  P (Sp(u))\subset Sp(P(u))  .$\vskip1mm
    \item Lorsque $  \mathbb{K} = \C   $ montrer que $   Sp(P(u))\subset P (Sp(u))  .$ Ce r\'esultat est-il vrai lorsque $  \mathbb{K} = \R   $~?
\end{enumerate}
\finenonce{001652}



\finexercice



\exercice{1653, legall, 1998/09/01}

\enonce{001653}{}
Soit $  E  $ un $  \mathbb{K} $-espace vectoriel de dimension finie et $  f \in \mathcal{L} (E)  $ telle
que $  f^2  $ soit diagonalisable. Montrer que $  f  $ est diagonalisable si et seulement si $
\hbox{Ker}(f)=\hbox{Ker}(f^2)  .$
\finenonce{001653}



\finexercice



\exercice{1654, legall, 1998/09/01}

\enonce{001654}{}
Soit $  f \in \mathcal{L} (\R ^3)  $ d\'etermin\'ee par sa matrice $  M=\begin{pmatrix}
1 & 1 & 1 \cr
                                      1 & 1 & 1 \cr
                                      -1 & 1 & 1 \cr \end{pmatrix} $ dans une base $  \{ e_1,
e_2,e_3\}   $ de $  \R ^3  .$
\begin{enumerate}
    \item Montrer que $  M  $ est diagonalisable.
    \item Montrer que la restriction de $  f  $ a tout sous-espace stable est diagonalisable.
    \item En d\'eduire tous les sous-espaces de $  \R ^3  $ stables par $  f  .$
\end{enumerate}
\finenonce{001654}



\finexercice



\exercice{1655, legall, 1998/09/01}

\enonce{001655}{}
Soit $  M \in M_n(\mathbb{K} )  $ et $  \varphi _M\in \mathcal{L} (M_n(\mathbb{K} ))  $ l'application
$  N\mapsto MN  $. Montrer que $  \varphi _M  $ est diagonalisable si et seulement si $  M  $
est diagonalisable. (utiliser le polyn\^ome minimal.)
\finenonce{001655}



\finexercice



\exercice{1656, legall, 1998/09/01}

\enonce{001656}{}
Soit $  E  $ un $  \mathbb{K} $-espace vectoriel de dimension finie et $  f
\in \mathcal{L} (E)  $ diagonalisable. Montrer que les propri\'et\'es suivantes sont \'equivalentes~:
\vskip1mm
(i) La famille $  \{ id, f, f^2, \ldots , f^{n-1}\}   $ est libre.
\vskip1mm
(ii) Il existe $  x \in E   :   \{ x, f(x), f^2(x), \ldots , f^{n-1}(x)\}   $ engendre $  E  .$
\vskip1mm
(iii) Les valeurs propres de $  f  $ sont simples.
\finenonce{001656}



\finexercice



\exercice{1657, legall, 1998/09/01}

\enonce{001657}{}
Soit $  \rho   $ l'application de $  \R _4[X]  $ dans
lui-m\^eme qui \`a un polyn\^ome $  P   $ associe le reste de la division euclidienne de $  P  $ par $  (X^2-1)  .$
\begin{enumerate}
    \item Montrer que $  \rho   $ est lin\'eaire.
    \item Montrer que $  \rho ^2= \rho   .$ En d\'eduire que $  \rho   $ est diagonalisable.
    \item D\'eterminer (de pr\'ef\'erence sans calcul) une base de vecteurs propres pour $  \rho   .$
\end{enumerate}
\finenonce{001657}



\finexercice



\exercice{1658, roussel, 2001/09/01}

\enonce{001658}{}
Soit $f$ l'endomorphisme de $\mathbb{R}^{3}$, dont la matrice dans la base canonique
$\{e_1, e_2, e_3\}$ est
$$A = \left (
        \begin{array}{ccc}
           ~3  & 2  & -2  \\
           -1  & 0  & ~1  \\
           ~1  & 1  & ~0  \\
                  \end{array}
\right ) $$
\begin{enumerate}
\item Calculer les valeurs propores de $A$. L'endomorphisme $f$ est-il diagonalisable ?
\item Calculer $(A-I)^2$. En d\'eduire $A^n$, en utilisant la formule du
bin\^ome de Newton.
\item Soient $P(X)=(X-1)^2$ et $Q \in \mathbb{R}[X]$. Exprimer le reste de la
division euclidienne de $Q$ par $P$ en fonction de $Q(1)$ et $Q'(1)$, o\`u
$Q'$ est le polyn\^ome d\'eriv\'e de $Q$.\\
En remarquant que $P(A)=0$ (on dit alors que $P$ est un polyn\^ome ~annulateur
de $A$) et en utilisant le r\'esultat pr\'ec\'edent avec un choix judicieux du
polyn\^ome ~$Q$, retrouver $A^n$.
\item Montrer que l'image de $\mathbb{R}^{3}$ par l'endomorphisme ~$(A-I)$ est un
sous-espace de dimension $1$, dont on d\'esignera une base par $\epsilon_2$.
D\'eterminer ensuite un vecteur $\epsilon_3$ tel que $f(\epsilon_3)= \epsilon_2 +
\epsilon_3$. Soit enfin $\epsilon_1$, un vecteur propre de $f$, non colin\'eaire \`a
$\epsilon_2$. Ecrire $\tilde{A}$, la matrice de $f$ dans la base $\left \{\epsilon_1,
\epsilon_2, \epsilon_3 \right \}$, ainsi que la matrice de passage $P$ et son inverse
$P^{-1}$. Retrouver $A^n$.
\end{enumerate}
\finenonce{001658}



\finexercice



\exercice{1659, ortiz, 1999/04/01}

\enonce{001659}{}
Soit $f$ un automorphisme d'un $\Cc$-espace
vectoriel $E$ de dimension finie. Montrer que $f$
est diagonalisable si et seulement si $f^2$ est
diagonalisable.
\finenonce{001659}



\finexercice



\exercice{1660, ortiz, 1999/04/01}

\enonce{001660}{}
Les questions sont ind\'ependantes. $K$ d\'esigne
$\Rr$ ou $\Cc$, $E$ est un $K$-espace vectoriel de
dimension finie $n$, $\mathcal{B}=(e_1,...,e_n)$
est une base fix\'ee de $E$ et $f$ un
endomorphisme de $E$.
\begin{enumerate}
     \item Quels sont les valeurs propres de l'endomorphisme nul de $E$ ?
     \item On suppose que la matrice de $f$ dans $\mathcal{B}$ est
     $M=\left(\begin{smallmatrix} 3&2&4\\-1&3&-1\\-2&-1&-3\end{smallmatrix}\right)$.
     \begin{enumerate}
         \item 2 est-il valeur propre de $f$ ?
         \item Le vecteur $2e_1+e_2+e_3$ est-il un vecteur propre de $f$ ?
     \end{enumerate}
     \item Pourquoi un vecteur de $E$ ne peut-il \^etre vecteur propre relativement \`a deux valeurs propres distinctes ?
     \item
     \begin{enumerate}
         \item Est-il vrai que si $\lambda$ est une valeur propre de $f$ et si $P$ est un polyn\^ome annulateur de $f$ alors $\lambda$ est racine de  $P$?
         \item Est-il vrai que si $\lambda$ est une racine d'un polyn\^ome annulateur de $f$ alors $\lambda$ est une valeur propre de $f$ ?
     \end{enumerate}
     \item Montrer que si $f^2-2f+\text{Id}_E=0$ alors $1$ est valeur propre de $f.$
     \item Montrer qu'il existe toujours au moins un scalaire $\alpha$ tel que $f-\alpha \text{Id}_E$ est bijectif.
     \item Donner un exemple d'endomorphisme $f$ de $E$ avec $n=2$ tel que la somme de deux vecteurs propres de $f$ n'est pas un vecteur propre de $f$.
     \item On suppose que $E=E_1\oplus E_2$ et que si $x\in E$ s'\'ecrit $x_1+x_2$ avec $x_1\in E_1$ et $x_2\in E_2$ alors $f(x)=2x_1-3x_2$.
     \begin{enumerate}
         \item Quel r\'esultat assure l'existence d'un tel endomorphisme ?
         \item Montrer que $f$ est diagonalisable.
     \end{enumerate}
     \item La matrice $M=\left(\begin{smallmatrix} 1&0&1\\0&1&0\\0&0&1\end{smallmatrix}\right)$ est-elle diagonalisable ?
     \item Si l' endomorphisme $f$ admet 0 pour valeur propre et est diagonalisable, que peut-on dire de la dimension du noyau de $f$ ?
\end{enumerate}
\finenonce{001660}



\finexercice



\exercice{1661, ortiz, 1999/04/01}

\enonce{001661}{}
\'Etudier le caract\`ere diagonalisable des matrices suivantes et le cas \'ech\'eant,
les diagonaliser :
\begin{enumerate}
     \item $ A=\left(\begin{smallmatrix}
-2&1&1\\
8&1&-5\\
4&3&-3
\end{smallmatrix}\right)
 \in\mathcal{M}_3(\Rr),
$
     \item $ B
=\left(\begin{smallmatrix}
1&-1&1&0&1\\
0&0&1&1&0\\
0&-1&2&0&1\\
0&0&0&1&-2\\
0&0&0&2&-3
\end{smallmatrix}\right) \in \mathcal{M}_5(\Rr),$
    \item $C=\left(\begin{smallmatrix} 0&1&0&0\\
1&k&1&1\\ 0&1&0&0\\ 0&1&0&0
\end{smallmatrix}\right) \in \mathcal{M}_4(\Cc),k\in\Cc .
$
\end{enumerate}
\finenonce{001661}



\finexercice



\exercice{1662, ortiz, 1999/04/01}

\enonce{001662}{}
Soient $A\in \mathcal{M}_n(K)$ telle que
$\text{tr}(A)\not= 0$ et
$$
f: \mathcal{M}_n(K)\to\mathcal{M}_n(K), M\mapsto
\text{tr} (A)M-\text{tr}(M)A.$$
\begin{enumerate}
    \item  Montrer que $f$ est un endomorphisme de $\mathcal{M}_n(K)$.
    \item  Montrer que $\mathcal{T}=\left\{ M\in \mathcal{M}_n(K): \text{tr}(M)=0\right\}$ et
$\text{vect}(A)$ sont des sous-espaces propres de
$f$.
    \item  En d\'eduire que $f$ est diagonalisable et \'ecrire la matrice r\'eduite de $f$.
\end{enumerate}
\finenonce{001662}



\finexercice



\exercice{1663, ortiz, 1999/04/01}

\enonce{001663}{}
Montrer que  si le polyn\^ome minimal d'un
endomorphisme $f$ d'un $K$-espace vectoriel de
dimension finie admet une racine $\lambda\in K$
alors $\lambda$ est valeur propre de $f$.
\finenonce{001663}



\finexercice



\exercice{1664, ortiz, 1999/04/01}

\enonce{001664}{}
\'Etudier le caract\`ere diagonalisable des matrices
suivantes
\begin{enumerate}
    \item $A=\left( \begin{array}{ccc} 3&2&4\\
-1&3&-1\\
-2&-1&-3
\end{array}
\right) \in\mathcal{M}_3(\Rr)$,
    \item  $B=\left( \begin{array}{cccc}
0&\ldots &0&1\\
\vdots&&\vdots&\vdots\\
0&\ldots &0&1\\
1&\ldots&1&0
\end{array}
\right)\in\mathcal{M}_n(\Rr),n\ge 2$,
\end{enumerate}
\finenonce{001664}



\finexercice



\exercice{1665, ortiz, 1999/04/01}

\enonce{001665}{}
Soient $E$ un $K$-espace vectoriel de dimension
$n$ et $f$ un endomorphisme de $E$ de rang $1$.
\begin{enumerate}
    \item  Montrer que si $f$ est diagonalisable alors $\text{tr}(f)\not= 0$.
    \item
Montrer qu'il existe $\lambda\in K$ tel que le
polyn\^ome cararact\'eristique de $f$ s'\'ecrive
$${\chi_f=(-1)^nX^{n-1}(X-\lambda)}.$$
    \item
    \begin{enumerate}
    \item Montrer que $f$ est diagonalisable si et seulement si $\text{tr}(f)\not= 0$.
\item R\'eduire sans calcul la matrice $A=\left( \begin{smallmatrix}
1&1&-1\\
-2&-2&2\\
1&1&-1
\end{smallmatrix}\right) \in \mathcal{M}_3(\Rr)$ et donner sans calcul les sous-espaces
vectoriels propres.
     \end{enumerate}
\end{enumerate}
\finenonce{001665}



\finexercice



\exercice{1666, ortiz, 1999/04/01}

\enonce{001666}{}
Soient les matrices
$A=\left( \begin{smallmatrix}
4&0&0\\
2&3&0\\
-1&0&1
\end{smallmatrix}\right) \in \mathcal{M}_3(\Rr)$,
$D=\left( \begin{smallmatrix}
4&0&0\\
0&3&0\\
0&0&1
\end{smallmatrix}\right) \in \mathcal{M}_3(\Rr)$.
\begin{enumerate}
    \item  Soit $Y\in\mathcal{M}_3(\Rr)$ telle que $Y^2=D$.
    \begin{enumerate}
        \item Montrer que $Y$ et $D$ commutent.
        \item En d\'eduire que $Y$ est diagonale puis d\'eterminer $Y$.
    \end{enumerate}
    \item
    \begin{enumerate}
         \item Montrer que $A$ est diagonalisable.
         \item En d\'eduire les solutions $X\in\mathcal{M}_3(\Rr)$ de l'\'equation $X^2=A$.
    \end{enumerate}
\end{enumerate}
\finenonce{001666}



\finexercice



\exercice{1667, ortiz, 1999/04/01}

\enonce{001667}{}
Soient $E$ un $\Cc$-espace vectoriel de dimension
$n$ et $f$ un endomorphisme de $E.$
\begin{enumerate}
    \item  Montrer que si $f$ est diagonalisable alors $f^2$ est
diagonalisable et $\text{rg}(f)=\text{rg}(f^2).$
    \item  Soit $\mu\in\Cc\setminus\left\{0\right\}$. Montrer que $\ker(f^2-\mu^2\text{Id}_E)=
 \ker(f-\mu\text{Id}_E)\oplus \ker(f+\mu\text{Id}_E).$
    \item  On suppose $\text{rg}(f)=\text{rg}(f^2).$
    \begin{enumerate}
        \item Montrer que $\ker(f)=\ker(f^2).$
        \item On suppose en outre que $f^2$ est diagonalisable. Montrer que $f$ est diagonalisable.
    \end{enumerate}
\end{enumerate}
\finenonce{001667}



\finexercice



\exercice{1668, ortiz, 1999/04/01}

\enonce{001668}{}
On consid\`ere la matrice par blocs  $A=\left( \begin{array}{cc}O&-I_n\\ I_n&0
\end{array}\right )\in\mathcal{M}_{2n}(\Cc).$
\begin{enumerate}
    \item  Calculer $A^2.$
    \item  Rechercher les \'el\'ements propres de $A.$ La matrice $A$ est-elle diagonalisable ?
\end{enumerate}
\finenonce{001668}



\finexercice



\exercice{1669, roussel, 2001/09/01}

\enonce{001669}{}
On d\'esigne par $E$ l'espace vectoriel des polyn\^ome s \`a coefficients r\'eels,
et par $E_n$, le sous-espace des polyn\^ome s de degr\'e au plus $n$.
\begin{enumerate}
\item Montrer que pour tout $x$ dans $\mathbb{R}, ~~
\Delta P(x)=(x+1)P'(x)+2P(x)$ d\'efinit une application lin\'eaire de $E$ dans
$E$. Quel est le degr\'e de $\Delta P$ lorsque $P$ appartient \`a $E_n$?

\item On consid\`ere $\Delta_2$, la restriction de $\Delta $ au sous-espace
$E_2$. D\'eterminer les valeurs propres ~de $\Delta_2$. L'endomorphisme ~$\Delta_2$
est-il diagonalisable ? Est-ce que $\Delta_2$ est un isomorphisme ?
\item En utilisant la d\'efinition des valeurs propres, calculer les valeurs propres ~et les
polyn\^ome s propres de $\Delta $.
\end{enumerate}
\finenonce{001669}



\finexercice



\exercice{1670, roussel, 2001/09/01}

\enonce{001670}{}
Pour tout \'el\'ement non nul $a=(a_1, a_2, \ldots ,a_n)$ de
$\mathbb{R}^{n}$, on consid\`ere l'endomorphisme ~$u$ de
$\mathbb{R}^{n}$ dont la matrice dans la base canonique
$\{e_{ij},i,j=1,2, \ldots, n \}$ est la matrice $A=\left ( \alpha
_{i,j}\right )$ o\`u $\alpha _{i,j}=a_ia_j$.
\begin{enumerate}
\item D\'eterminer le noyau et l'image de $u$.
\item En d\'eduire les sous-espaces propres de $u$. D\'eterminer les valeurs
propres de $u$. L'endomorphisme $u$ est-il diagonalisable ?
\item Quel est le polyn\^ome caract\'eristique de $u$?
\end{enumerate}
\finenonce{001670}



\finexercice



\exercice{1671, roussel, 2001/09/01}

\enonce{001671}{}
Soit $B$ une matrice diagonalisable ~de $\mathcal{M}_n \left (\mathbb{R} \right )$.
On d\'efinit son rayon spectral par
$$\rho(B)= \max \left \{ |\lambda| \mbox{~~avec~~} \lambda \mbox{~~est une
valeur propre ~de~~} B \right \}.$$
\begin{enumerate}
\item Montrer que $\lim _{k \longrightarrow + \infty} B^k= 0$.
\item En d\'eduire que $I-B$ est inversible et que $(I-B)^{-1}=
\displaystyle{\sum _{k=0}^{+ \infty}} B^{k}$.
\end{enumerate}
\finenonce{001671}



\finexercice



\exercice{1672, gineste, 2001/11/01}

\enonce{001672}{Endomorphisme diagonalisable de $\Rr^2$}
 On
considère l'endomorphisme $a$ de $E=\Rr^2$ dont la matrice
représentative $A=[a]_e^e$ dans la base canonique $e$ est $\left[
\begin{array}{ll}7&-10\\5&-8 \end{array} \right].$ Calculer la
trace, le déterminant, le polynôme caractéristique et le spectre
de $a$. Quel théorème du cours garantit l'existence d'une base
$f=(\vec{f}_1,\vec{f}_2)$ de vecteurs propres? Choisir ensuite $f$
telle que $[\mathrm{id}_E]_f^e$ et $[\mathrm{id}_E]_e^f$ soient à
coefficients entiers. Dessiner $\vec{f}_1$ et $\vec{f}_2$, en
prenant des unités d'axes assez petites. Dessiner quelques
vecteurs $\vec{x}$ et leurs images $a(\vec{x})$ à l'aide de $f$.

Trouver deux matrices $P$ et $D$ carrées d'ordre 2 telles que $D$ soit
diagonale, $P$ inversible et $A=PDP^{-1}$. Calculer $[a^{50}]_f^f$,
$[a^{50}]_e^e$ et $A^{50}$. Calculer $\lim_{n\infty}\frac{1}{3^{2n}}a^{2n}.$

\finenonce{001672}



\finexercice



\exercice{1673, gineste, 2001/11/01}

\enonce{001673}{Endomorphisme d'un espace de matrices}
Soit
$K$ un corps commutatif quelconque, et soit $F=\mathcal{M}_n(K)$
l'espace vectoriel sur $K$ des matrices carrées d'ordre $n$ à
coefficients dans $K.$ Si $i$ et $j$ sont des entiers compris
entre 1 et $n$, on note par $F_{ij}$ l'élément de $F$ dont le
coefficient $(i,j)$ est 1 et dont les autres coefficients sont
nuls. Montrer que les $F_{ij}$ forment une base de $F$. Dimension
de $F$? Soit $D$ dans $F$ et diagonale. Soient $\alpha$ et $\beta$
dans $K$ et soit l'endomorphisme $\Phi$ de $F$ qui à la matrice
$X$ fait correspondre la matrice $\Phi(X)=\alpha XD+ \beta DX.$
Calculer $\Phi(F_{ij}).$ $\Phi$ est il un endomorphisme
diagonalisable? Donner son polynôme caractéristique en fonction
des coefficients de $D$ et de $\alpha$ et $\beta$.

\finenonce{001673}



\finexercice



\exercice{1674, gineste, 2001/11/01}

\enonce{001674}{}
 Soit $\theta\in]0,\pi[$. On considère les deux
matrices d'ordre $n$:

$$ A=\left[ \begin{array}{llllll}0&1&0&\cdots&0&0\\
1&0&1&\cdots&0&0\\
0&1&0&\cdots&0&0\\
\cdots&\cdots&\cdots&\cdots&\cdots&\cdots\\
0&0&0&\cdots&0&1\\
0&0&0&\cdots&1&0
 \end{array} \right],
B=\left[ \begin{array}{llllll}
2\cos \theta&1&0&\cdots&0&0\\
1&2\cos \theta&1&\cdots&0&0\\
0&1&2\cos \theta&\cdots&0&0\\
\cdots&\cdots&\cdots&\cdots&\cdots&\cdots\\
0&0&0&\cdots&2\cos \theta&1\\
0&0&0&\cdots&1&2\cos \theta
 \end{array} \right]$$
Montrer par récurrence que
$\det B=\frac{\sin(n+1)\theta}{\sin \theta}$ (Méthode: développer
par rapport à la dernière ligne). Montrer que $\det B$
s'annule pour $n$ valeurs distinctes de $\theta$ de $]0,\pi[$,
et les déterminer. Si $P_A$ est le polynôme caractéristique de $A$, calculer
$P_A(-2\cos \theta)$ et déduire de ce qui précède les valeurs propres de $A.$
Montrer que les valeurs propres des matrices $2I_n+A$ et $2I_n-A$ sont
strictement positives.

\finenonce{001674}



\finexercice



\exercice{2469, matexo1, 2002/02/01}

\enonce{002469}{}
Soit $a,b,c$ trois r\'eels et $u$ l'endomorphisme de $\R^3$
ayant pour matrice
$$M = \left( \begin{array}{ccc} 1&a&b\\ 0&1&c \\ 0&0&-1 \end{array} \right).$$
Discuter la possibilit\'e de le diagonaliser selon les valeurs de
$a,b,c$.
\finenonce{002469}


\finexercice


\exercice{2477, matexo1, 2002/02/01}

\enonce{002477}{}
Soit $A$ une matrice carr\'ee r\'eelle d'ordre $n$ non nulle et
nilpotente.
\begin{enumerate}
\item Montrer que $I-A$ n'est pas diagonalisable.
\item G\'en\'eraliser en montrant que si $B$ est une matrice
diagonalisable dont toutes les valeurs propres sont \'egales, alors
$B+A$ n'est pas diagonalisable.
\item Montrer qu'il existe $A, B$ dans ${\mathcal M}_2(\R)$, $A \neq 0$
nilpotente et $B$ diagonalisable, telles que $A+B$ soit
diagonalisable.
\end{enumerate}
\finenonce{002477}


\finexercice


\exercice{2563, delaunay, 2009/05/19}
\enonce{002563}{}
Soit $M$ la matrice r\'eelle $3\times3$ suivante :
$$M=\begin{pmatrix}0&2&-1\\3&-2&0\\-2&2&1\end{pmatrix}$$
\begin{enumerate}
  \item D\'eterminer les valeurs propres de $M$.

  \item Montrer que $M$ est diagonalisable.
  \item D\'eterminer une base de vecteurs propres et $P$ la matrice de passage.
  \item On a $D=P^{-1}MP$, pour $k\in\N$ exprimer $M^k$ en fonction de $D^k$, puis calculer $M^k$.
\end{enumerate} 
\finenonce{002563} 


\finexercice


\exercice{2566, delaunay, 2009/05/19}
\enonce{002566}{}
Soit $$A=\begin{pmatrix}1&0&0 \\ 0&1&0 \\ 1&-1&2\end{pmatrix}$$
D\'emontrer que $A$ est diagonalisable et trouver une matrice $P$ telle que $P^{-1}AP$ soit diagonale.
\finenonce{002566} 


\finexercice


\exercice{2567, delaunay, 2009/05/19}
\enonce{002567}{}
 Soit $$A=\begin{pmatrix}1&1&-1 \\ 0&1&0 \\ 1&0&1\end{pmatrix}$$
Factoriser le polyn\^ome caract\'eristique de $A$. La matrice $A$ est-elle diagonalisable dans $\R$~? dans $\C$ ?
\finenonce{002567} 


\finexercice


\exercice{2568, delaunay, 2009/05/19}
\enonce{002568}{}
 Soit  $$A=\begin{pmatrix}a&c \\  c&d\end{pmatrix}\in M_2(\R)$$
D\'emontrer que $A$ est diagonalisable dans $\R$.
\finenonce{002568} 


\finexercice


\exercice{2572, delaunay, 2009/05/19}
\enonce{002572}{}
 ({\it 9 points})  Soit $A$ la matrice de $M_3(\R)$ suivante :
$$A=\begin{pmatrix}1&0&1 \\  -1&2&1 \\  1&-1&1\end{pmatrix}$$
\begin {enumerate}
  \item  D\'emontrer que les valeurs propres de $A$ sont $1$ et $2$.
  \item D\'eterminer les sous-espaces propres de $A$. La matrice $A$ est-elle diagonalisable ?

  \item D\'eterminer les sous-espaces caract\'eristiques de $A$.
  \item D\'eterminer une base de $\R^3$ dans laquelle la matrice de l'endomorphisme associ\'e \`a $A$ est 
$$B=\begin{pmatrix}2&0&0 \\  0&1&1 \\  0&0&1\end{pmatrix}$$
 En d\'eduire la d\'ecomposition de 
Dunford de $B$.
  \item  R\'esoudre le syst\`eme diff\'erentiel

$$\left\{\begin{align*}x'&=x+z \\  y'&=-x+2y+z \\  z'&=x-y+z\end{align*}\right.$$
\end {enumerate}
\finenonce{002572} 



\finexercice


\exercice{2573, delaunay, 2009/05/19}
\enonce{002573}{}
 ({\it 7 points})  On consid\`ere la suite $(u_n)_{n\in\N}$ d\'efinie par $u_0=0$, $u_1=1$ et par la relation de r\'ecurrence
 $$u_{n+1}={\frac{1}{2}}(u_n+u_{n-1}). $$
 \begin {enumerate}
   \item  D\'eterminer une matrice $A\in M_2(\R)$ telle que  pour tout $n\geq1$ on ait
$$\begin{pmatrix}u_{n+1} \\  u_n\end{pmatrix}=A^n\begin{pmatrix}u_1 \\  u_0\end{pmatrix}.$$
Justifier.
  \item D\'eterminer le polyn\^ome caract\'eristique $P_A(X)$ de $A$ et calculer ses racines $\lambda_1$ et $\lambda_2$.
  \item  Soit $R_n(X)=a_nX+b_n$ le reste de la division euclidienne de $X^n$ par $P_A(X)$. Calculer $a_n$ et $b_n$ 
(on pourra utiliser les racines $\lambda_1$ et $\lambda_2$).
  \item Montrer que $A^n=a_nA+b_nI_2$, en d\'eduire que la matrice $A^n$ converge lorsque $n$ tend vers $+\infty$ vers une
limite $A_{\infty}$ que l'on d\'eterminera. Calculer $\displaystyle\lim_{n\rightarrow+\infty}u_n$.
\end {enumerate}
\finenonce{002573} 



\finexercice


\exercice{2574, delaunay, 2009/05/19}
\enonce{002574}{}
 ({\it 5 points})  Soit $A$ une matrice carr\'ee, $A\in M_n(K)$ ($K=\R$ ou $\C$). On rappelle que la trace d'une matrice est la somme de ses coefficients diagonaux et que ${\rm tr}(BAB^{-1})={\rm tr}A$.

D\'emontrer que $\det(\exp A)={\rm e^{trA}}$ dans les cas suivants : 
\begin {enumerate}
  \item  $A$ diagonalisable.
  \item $A$ triangulaire sup\'erieure ayant une diagonale de z\'eros.
  \item  $A$ trigonalisable.
  \item $A$ quelconque.
\end {enumerate}
\finenonce{002574} 


\finexercice


\exercice{2575, delaunay, 2009/05/19}
\enonce{002575}{}
 ({\it 4 points}) On suppose qu'une population $x$ de lapins et une population $y$ de loups sont gouvern\'ees par le syst\`eme suivant d'\'equations diff\'erentielles :
$$(S)\ \ \left\{ \begin{align*}x'&= 4x-2y \\  y'&= x+y\end{align*}\right.$$
\begin {enumerate}
  \item  Diagonaliser la matrice 
$$A=\begin{pmatrix}4&-2 \\ 1&1 \\ \end{pmatrix}.$$
  \item Exprimer le syst\`eme $(S)$ et ses solutions dans une base de vecteurs propres de $A$.
  \item  Repr\'esenter graphiquement les trajectoires de $(S)$ dans le rep\`ere $(Oxy)$.
  \item Discuter graphiquement l'\'evolution de la population des lapins en fonction des conditions initiales.
\end {enumerate}
\finenonce{002575} 


\finexercice


\exercice{2576, delaunay, 2009/05/19}
\enonce{002576}{}
({\it 9 points}) Soit $u$ l'endomorphisme de $\R^3$, dont la matrice dans la base canonique est
$$A=\begin{pmatrix}3&2&-2 \\  -1&0&1 \\ 1&1&0 \\ \end{pmatrix}$$ 
\begin {enumerate}
  \item  Calculer les valeurs propres de $A$. L'endomorphisme $u$ est-il diagonalisable ?
  \item Calculer $(A-I)^2$. Montrer que $A^n=nA+(1-n)I$ en utilisant la formule du bin\^ome de Newton.
  \item  Soient $P(X)=(X-1)^2$ et $Q\in\R[X]$. Exprimer le reste de la division euclidienne de $Q$ par $P$ en fonction de $Q(1)$ et $Q'(1)$, o\`u $Q'$ est le polyn\^ome d\'eriv\'e de $Q$. En remarquant que $P(A)=0$ et en utilisant le r\'esultat pr\'ec\'edent avec un choix judicieux du polyn\^ome 
$Q$, retrouver $A^n$.
  \item
  \begin {enumerate}
    \item  Montrer que l'image de $\R^3$ par l'endomorphisme $u-\mathrm{Id}$ est un sous-espace vectoriel de dimension $1$, on notera $\varepsilon_2$ une base. 
    \item D\'eterminer un vecteur $\varepsilon_3$ tel que 
$u(\varepsilon_3)=\varepsilon_2+\varepsilon_3$. D\'eterminer un vecteur propre $\varepsilon_1$  de $u$ non colin\'eaire \`a $\varepsilon_2$. 

    \item Montrer que $(\varepsilon_1,\varepsilon_2,\varepsilon_3)$ est une base de $\R^3$.Ecrire la matrice de $u$ dans cette base, ainsi que les matrices de passage. 
    \item Retrouver $A^n$.
  \end {enumerate}
\end {enumerate}
\finenonce{002576} 


\finexercice


\exercice{2577, delaunay, 2009/05/19}
\enonce{002577}{}
 ({\it 7 points}) Soient $M$ et $A$ deux matrices de ${\cal{M}}_n(\R)$ telles que 
$MA=AM$.
On suppose que $M$ admet $n$ valeurs propres distinctes.
 \begin {enumerate}
  \item  Soit $x$ un vecteur propre de $M$ de valeur propre $\lambda$, montrer que 
$MAx=\lambda Ax,$
en d\'eduire que les vecteurs $x$ et $Ax$ sont colin\'eaires, puis que tout vecteur propre de $M$ est un vecteur propre de $A$.

  \item On note maintenant $\lambda_1,\cdots,\lambda_n$ les valeurs propres de $M$ et 
$\mu_1,\cdots,\mu_n$ celles de $A$.
  \begin {enumerate}
    \item Montrer par r\'ecurrence sur $n$ l'\'egalit\'e suivante :
$$\begin{vmatrix}1 &\lambda_1 &\cdots &\lambda_1^{n-1} \\ 
\vdots& & &\vdots \\  1 &\lambda_n &\cdots & \lambda_n^{n-1}\end{vmatrix}=
\prod_{1\leq i<j\leq n}(\lambda_i-\lambda_j).$$
En d\'eduire que le syst\`eme suivant
$$\left\{\begin{align*}\mu_1&=\alpha_0+\alpha_1\lambda_1+\cdots+\alpha_{n-1}\lambda_1^{n-1} \\ 
\vdots \\  \mu_n &=\alpha_0+\alpha_1\lambda_n+\cdots+\alpha_{n-1}\lambda_n^{n-1} \\ \end{align*}\right.$$
admet une unique solution $(\alpha_0,\cdots,\alpha_{n-1})\in\R^n.$
    \item Soient $M'$ et $A'$ les matrices diagonales suivantes :
$$M'=\begin{pmatrix}\lambda_1 & 0 &\cdots & 0 \\  0 &\ddots  & & 0 \\ \vdots & &  & \vdots \\  0 & \cdots & 0 &\lambda_n \\  \end{pmatrix},\ \
A'=\begin{pmatrix}\mu_1 & 0 &\cdots & 0 \\  0 &\ddots  & & 0 \\ \vdots & &  & \vdots \\  0 & \cdots & 0 &\mu_n \\ \end{pmatrix}.$$
Montrer qu'il existe des r\'eels $\alpha_0,\cdots,\alpha_{n-1}$ tels que
$$A'=\sum_{k=0}^{n-1}\alpha_k M'^k$$ et en d\'eduire qu'il existe des r\'eels $\alpha_0,\cdots,\alpha_{n-1}$ tels que
$$A=\sum_{k=0}^{n-1}\alpha_k M^k.$$   
  \end {enumerate}
\end {enumerate}
\finenonce{002577} 


\finexercice


\exercice{2581, delaunay, 2009/05/19}
\enonce{002581}{}
 Soit $u$ l'endomorphisme de $\R^3$ dont la matrice dans la base canonique est  
$$A=\begin{pmatrix}-3&-2&-2 \\  2&1&2 \\  3&3&2\end{pmatrix}$$ 
\begin{enumerate}
 \item D\'eterminer et factoriser le polyn\^ome caract\'eristique de $A$.
 \item D\'emontrer que les valeurs propres de $A$ sont $-1$ et $2$. D\'eterminer les sous-espaces propres associ\'es. 
 \item D\'emontrer que $A$ est diagonalisable et donner une base de $\R^3$ dans laquelle la matrice de $u$ est diagonale.
 \item Trouver une matrice $P$ telle que $P^{-1}AP$ soit diagonale.
\end{enumerate}
\finenonce{002581} 


\finexercice


\exercice{2583, delaunay, 2009/05/19}
\enonce{002583}{}
 Soit $A=\begin{pmatrix}1&2&0&0 \\ 0&1&2&0 \\ 0&0&1&2 \\ 0&0&0&1\end{pmatrix}$.
Expliquer sans calcul pourquoi la matrice $A$ n'est pas diagonalisable.
\finenonce{002583} 


\finexercice


\exercice{2584, delaunay, 2009/05/19}
\enonce{002584}{}
 Soit $A$ une matrice $2\times 2$ \`a coefficients r\'eels. On suppose que dans chaque colonne de $A$ la somme des coefficients est \'egale \`a $1$. 
\begin{enumerate}
 \item Soient $(x_1,x_2)$, $(y_1,y_2)$ deux vecteurs de $\R^2$, on suppose que  
$$A\begin{pmatrix}x_1 \\  x_2\end{pmatrix}=\begin{pmatrix}y_1 \\  y_2\end{pmatrix}$$ 
montrer qu'alors 
$$y_1+y_2=x_1+x_2.$$
 \item Soit le vecteur $\varepsilon=(1,-1)$, montrer que c'est un vecteur propre de $A$.
On notera $\lambda$ sa valeur propre.
 \item Montrer que si $v$ est un vecteur propre de $A$ non colin\'eaire \`a $\varepsilon$, alors la valeur propre
associ\'ee \`a $v$ est \'egale \`a $1$.
 \item  Soit $e_1=(1,0)$. Montrer que la matrice, dans la base $(e_1,\varepsilon)$, de l'endomorphisme associ\'e \`a $A$  est de la forme
$$\begin{pmatrix}1&0  \\ \alpha&\lambda \end{pmatrix},$$
o\`u $\alpha\in\R$.

En d\'eduire que si $\lambda\neq 1$, alors $A$ est diagonalisable sur $\R$.
\end{enumerate}
\finenonce{002584} 


\finexercice


\exercice{2586, delaunay, 2009/05/19}
\enonce{002586}{}
\centerline{\bf I}

Soit $\alpha\in\R$ et $A_\alpha\in M_3(\R)$ la matrice suivante
$$A_\alpha=\begin{pmatrix}-1 & 0 & \alpha+1 \\  1&-2&0 \\ -1&1&\alpha\end{pmatrix}$$
{\it Premi\`ere partie} :

\begin{enumerate} 
  \item Factoriser le polyn\^ome caract\'eristique $P_{A_\alpha}(X)$ en produit de facteurs du premier degr\'e.

  \item D\'eterminer selon la valeur du param\`etre $\alpha$ les valeurs propres distinctes de $A_\alpha$ et leur multiplicit\'e.

  \item D\'eterminer les valeurs de $\alpha$ pour lesquelles la matrice $A_\alpha$ est diagonalisable.

  \item D\'eterminer selon la valeur de $\alpha$ le polyn\^ome minimal de $A_\alpha$.
 \end{enumerate}
 
{\it Seconde partie} : 

On suppose d\'esormais que $\alpha=0$, on note $A=A_0$ et $f$ l'endomorphisme de $\R^3$ associ\'e \`a la matrice $A$.

\begin{enumerate} 
  \item D\'eterminer les sous-espaces propres et caract\'eristiques de $A$.

  \item D\'emontrer que $f$ admet un plan stable (c'est-\`a-dire $f$-invariant).

  \item D\'emontrer qu'il existe une base de $\R^3$ dans laquelle la matrice de $f$ est 
$$B=\begin{pmatrix}-1 & 1 & 0 \\  0&-1&1 \\ 0&0&-1\end{pmatrix}$$
et trouver une matrice $P$ inversible telle que $A=PBP^{-1}$.

  \item Ecrire la d\'ecomposition de Dunford de $B$ (justifier).

  \item Pour $t\in\R$, calculer $\exp tB$ et exprimer $\exp tA$ \`a l'aide de $P$ et $\exp tB$.

  \item Donner les solutions des syst\`emes diff\'erentiels $Y'=BY$ et $X'=AX$.
 \end{enumerate}
  

\centerline{\bf II}

On rappelle qu'une matrice $N\in M_n(\C)$ est dite nilpotente d'ordre $m$ si $N^m=0$, et si pour tout $k$ dans $\N$, $k<m$, on a $N^k\neq0$. Soient $N\in M_n(\C)$ une matrice nilpotente d'ordre $m$ et $A\in M_n(\C)$ une matrice telle que $AN=NA$.


\begin{enumerate} 
  \item D\'eterminer un polyn\^ome annulateur de $N$. En d\'eduire le polyn\^ome minimal et le polyn\^ome caract\'eristique de $N$.

  \item D\'eterminer les valeurs propres de $N$.

  \item D\'emontrer que $\det (I+N)=1$.

  \item On suppose $A$ inversible. D\'emontrer que les matrices $AN$ et $NA^{-1}$  sont nilpotentes.
En d\'eduire que $$\det(A+N)=\det A.$$

  \item On suppose $A$ non inversible. En exprimant $(A+N)^k$ pour tout $k\in\N$, d\'emontrer que $$\det(A+N)=0.$$ 
 \end{enumerate}

\finenonce{002586} 


\finexercice


\exercice{2587, delaunay, 2009/05/19}
\enonce{002587}{}
Soit $\displaystyle A=\begin{pmatrix}a&c \\  c&d\end{pmatrix}\in M_2(\R)$,
montrer que $A$ est diagonalisable sur $\R$.
\finenonce{002587} 


\finexercice


\exercice{2591, delaunay, 2009/05/19}
\enonce{002591}{}
 Soit $a\in\R$, notons $A$ la matrice suivante
$$A=\begin{pmatrix}0&1 \\  -a&1+a\end{pmatrix}.$$
On d\'efinit une suite $(u_n)_{n\in\N}$, par la donn\'ee de $u_0$ et $u_1$ et la relation de r\'ecurrence suivante, pour $n\in\N$
$$u_{n+2}=(1+a)u_{n+1}-au_n$$
\begin{enumerate}
 \item Pour quelles valeurs de $a$ la matrice $A$ est-elle diagonalisable ?
 \item Lorsque $A$ est diagonalisable, calculer $A^n$ pour $n\in \N$.
 \item On suppose $A$ diagonalisable. On note $U_n$ le vecteur $U_{n}=\begin{pmatrix}u_{n} \\  u_{n+1}\end{pmatrix}$, exprimer $U_{n+1}$ en fonction de $U_n$ et de $A$, puis $U_n$ en fonction de $U_0$ et de $A$.

\end{enumerate}

\finenonce{002591} 


\finexercice


\exercice{2592, delaunay, 2009/05/19}
\enonce{002592}{}
Soit $A$ la matrice de $M_3(\R)$ suivante :
$$A=\begin{pmatrix}0&1&0 \\  -4&4&0 \\ -2&1&2\end{pmatrix}.$$ 
\begin{enumerate}
 \item  La matrice $A$ est-elle diagonalisable ?
 \item  Calculer $(A-2I_3)^2$, puis $(A-2I_3)^n$ pour tout $n\in\N$. En d\'eduire $A^n$.
\end{enumerate}
\finenonce{002592} 


\finexercice


\exercice{2593, delaunay, 2009/05/19}
\enonce{002593}{}
 Soit $f$ l'endomorphisme de $\R^4$ dont la matrice dans la base canonique  est 
$$A=\begin{pmatrix}-8&-3&-3&1 \\  6&3&2&-1 \\ 26&7&10&-2 \\ 0&0&0&2\end{pmatrix}.$$ 
\begin{enumerate}
 \item  D\'emontrer que $1$ et $2$ sont des valeurs propres de $f$.
 \item D\'eterminer les vecteurs propres de $f$.
 \item  Soit $\vec u$ un vecteur propre de $f$ pour la valeur propre $2$. Trouver des vecteurs $\vec v$ et $\vec w$ tels que 
$$f(\vec v)=2\vec v+\vec u\ {\hbox{et}}\ f(\vec w)=2\vec w+\vec v.$$ 
 \item Soit $\vec e$ un vecteur propre de $f$ pour la valeur propre $1$. D\'emontrer que $(\vec e,\vec u,\vec v,\vec w)$ est une base de $\R^4$. Donner la matrice de $f$ dans cette base.
 \item La matrice $A$ est-elle diagonalisable ?
\end{enumerate}
\finenonce{002593} 


\finexercice


\exercice{2599, delaunay, 2009/05/19}
\enonce{002599}{}
Soit $m\in\R$, et $A$ la matrice
$$A=\begin{pmatrix}1+m&1+m&1 \\  -m&-m&-1 \\  m&m-1&0\end{pmatrix}$$
\begin{enumerate}
 \item Factoriser le polyn\^ome caract\'eristique de $A$ et montrer que les valeurs propres de $A$ sont $-1$ et $1$.
 \item Pour quelles valeurs de $m$ la matrice est-elle diagonalisable ? (justifier). D\'eterminer suivant les valeurs de $m$ le polyn\^ome minimal de $A$ (justifier).
\end{enumerate}
\finenonce{002599} 


\finexercice


\exercice{2600, delaunay, 2009/05/19}
\enonce{002600}{}
\begin{enumerate}
 \item Donner un exemple de matrice dans $M_2(\R)$, diagonalisable sur $\C$ mais non diagonalisable sur $\R$ (justifier).
 \item  Donner un exemple de matrice dans $M_2(\R)$ non diagonalisable, ni sur $\C$, ni sur $\R$ (justifier).
\end{enumerate}
\finenonce{002600} 


\finexercice


\exercice{2601, delaunay, 2009/05/19}
\enonce{002601}{}
Soit $A$ la matrice suivante :
$$A=\begin{pmatrix}0&1 \\ 1&0\end{pmatrix}$$
\begin{enumerate}
 \item Diagonaliser la matrice $A$.
 \item  Exprimer les solutions du syst\`eme diff\'erentiel $X'=AX$ dans une base de vecteurs propres et tracer ses trajectoires.  
\end{enumerate}
\finenonce{002601} 


\finexercice


\exercice{2603, delaunay, 2009/05/19}
\enonce{002603}{}
 Soit $a\in\R$ et $A$ la matrice suivante
$$A=\begin{pmatrix}1&a&0 \\  a&0&1 \\  0&1&a\end{pmatrix}.$$
\begin{enumerate}
 \item Calculer le d\'eterminant de $A$ et d\'eterminer 
pour quelles valeurs de $a$ la matrice est inversible.
 \item  Calculer $A^{-1}$ lorsque $A$ est inversible.
\end{enumerate}

\finenonce{002603} 


\finexercice


\exercice{2604, delaunay, 2009/05/19}
\enonce{002604}{}
 Soit $\theta\in \R$, on consid\`ere l'endomorphisme $f$ de $\R^3$ dont la matrice dans la base canonique est la suivante
$$A=\begin{pmatrix}\cos\theta&-\sin\theta&0 \\  \sin\theta&\cos\theta&0 \\ 0&0&1 \\ \end{pmatrix}.$$
\begin{enumerate}
 \item  Quelle est la nature géométrique de cet endomorphisme ?
 \item  Démontrer que, pour tout $\theta\in\R\setminus\pi\Z$, la matrice $A$ admet une unique valeur propre réelle. Quel est le sous-espace propre associé ? Que se passe-t-il si $\theta\in\pi\Z$ ?
\end{enumerate}
\finenonce{002604} 


\finexercice


\exercice{2605, delaunay, 2009/05/19}
\enonce{002605}{}
 Soit $u$ l'endomorphisme de $\R^3$ dont la matrice dans la base canonique est  
$$A=\begin{pmatrix}-4&-2&-2 \\  2&0&2 \\  3&3&1\end{pmatrix}.$$ 
\begin{enumerate}
 \item D\'eterminer et factoriser le polyn\^ome caract\'eristique de $A$.
 \item  D\'emontrer que les valeurs propres de $A$ sont $1$ et $-2$. D\'eterminer les sous-espaces propres associ\'es. 
 \item D\'emontrer que $A$ est diagonalisable et donner une base de $\R^3$ dans laquelle la matrice de $u$ est diagonale.
 \item Trouver une matrice $P$ telle que $P^{-1}AP$ soit diagonale.
\end{enumerate}
\finenonce{002605} 


\finexercice


\exercice{2606, delaunay, 2009/05/19}
\enonce{002606}{}
Soit $u$ l'endomorphisme de $\R^3$, dont la matrice dans la base canonique est
$$A=\begin{pmatrix}3&2&-2 \\  -1&0&1 \\ 1&1&0 \\ \end{pmatrix}.$$ 
\begin{enumerate}
 \item  Calculer les valeurs propres de $A$. L'endomorphisme $u$ est-il diagonalisable ? (Justifier).
 \item  Calculer $(A-I)^2$. Démontrer que $A^n=nA+(1-n)I$. 
\end{enumerate}
\finenonce{002606} 


\finexercice


\exercice{2607, delaunay, 2009/05/19}
\enonce{002607}{}
 Soit $A$ la matrice 
$$A=\begin{pmatrix}1&0&0 \\  -1&2&1 \\  0&0&2\end{pmatrix}$$ 
et $f$ l'endomorphisme de $\R^3$ associ\'e.
\begin{enumerate}
 \item  Déterminer les valeurs propres de $A$.
 \item  Déterminer, sans calculs, des vecteurs $\vec u$ et $\vec v$ tels que $f(\vec u)=2\vec u$ et $f(\vec v)=2\vec v+\vec u$.
 \item  Soit $\vec e$ tel que $f(\vec e)=\vec e$. Démontrer que $(\vec e,\vec u,\vec v)$ est une base de $\R^3$ et écrire la matrice de $f$ dans cette base.
 \item La matrice $A$ est-elle diagonalisable ? (Justifier.)
\end{enumerate}
\finenonce{002607} 


\finexercice


\exercice{2608, delaunay, 2009/05/19}
\enonce{002608}{}
Soit $A$ la matrice 
$$A=\begin{pmatrix}1&0&0 \\  1&-1&0 \\  -1&2&-1\end{pmatrix}$$ 
et $f$ l'endomorphisme de $\R^3$ associ\'e.
\begin{enumerate}
 \item  Factoriser le polyn\^ome caract\'eristique de $A$. 
 \item D\'eterminer les sous-espaces propres et caract\'eristiques de $A$. 
 \item D\'emontrer qu'il existe une base de $\R^3$ dans laquelle la matrice de $f$ est 
$$B=\begin{pmatrix}1 & 0 & 0 \\  0&-1&2 \\ 0&0&-1\end{pmatrix}$$
et trouver une matrice $P$ inversible telle que $ AP=PB$ (ou $A=PBP^{-1}$). 
 \item Ecrire la d\'ecomposition de Dunford de $B$ (justifier). 
 \item  Pour $t\in\R$, calculer $\exp tB$. 
 \item Donner les solutions des syst\`emes diff\'erentiels $y'=By$ et $x'=Ax$, où $x$ et $y$ désignent des fonctions réelles à valeurs dans $\R^3$. 
\end{enumerate}
\finenonce{002608} 


\finexercice


\exercice{2609, delaunay, 2009/05/19}
\enonce{002609}{}
 Soit $a\in\R$ et $A$ la matrice
$$A=\begin{pmatrix}0&1&0 \\  0&a&0 \\  0&a-2&2\end{pmatrix}.$$
\begin{enumerate}
 \item Pour quelles valeurs de $a$ la matrice $A$ est-elle diagonalisable ? 

Lorsque $A$ est diagonalisable, déterminer une base de vecteurs propres de $A$.

 \item Soit $E$ l'espace vectoriel des solutions du système $x'=Ax$, où $x$ est une fonction de la variable réelle $t$ à valeur dans $\R^3$.
   \begin{enumerate}
     \item  Lorsque $A$ est diagonalisable, donner une base de $E$ en fonction des vecteurs propres et des valeurs propres de $A$. Ecrire la solution générale du système.
     \item Lorsque $A$ n'est pas diagonalisable, intégrer directement le système $x'=Ax$.
    \end{enumerate}
 \item Soit $E_0$ l'ensemble des éléments $s$ de $E$ tels que $\lim_{t\rightarrow+\infty}s(t)=\vec0$. Démontrer que $E_0$ est un sous-espace vectoriel de $E$. (hors barème) Déterminer sa dimension en fonction de $a$.
 \item Soit $F$ l'ensemble des éléments $s$ de $E$ bornés sur $[0,+\infty[$. Démontrer que $F$ est un sous-espace vectoriel de $E$. (hors barème) Déterminer sa dimension en fonction de $a$.
\end{enumerate}
\finenonce{002609} 


\finexercice


\exercice{2610, delaunay, 2009/05/19}
\enonce{002610}{}
 Soit $\alpha\in\R$ et $A_\alpha\in M_3(\R)$ la matrice suivante
$$A_\alpha=\begin{pmatrix}-1 & 0 & \alpha+1 \\  1&-2&0 \\ -1&1&\alpha\end{pmatrix}$$
 \centerline{I} 

\begin{enumerate} 
\item Factoriser le polyn\^ome caract\'eristique $P_{A_\alpha}(X)$ en produit de facteurs du premier degr\'e.

  \item D\'eterminer selon la valeur du param\`etre $\alpha$ les valeurs propres distinctes de $A_\alpha$ et leur multiplicit\'e.

  \item D\'eterminer les valeurs de $\alpha$ pour lesquelles la matrice $A_\alpha$ est diagonalisable.

  \item D\'eterminer selon la valeur de $\alpha$ le polyn\^ome minimal de $A_\alpha$.
\end{enumerate}

\centerline{II}  

On suppose, dans cette partie, que $\alpha=0$, on note $A=A_0$ et $f$ l'endomorphisme de $\R^3$ associ\'e \`a la matrice $A$.

\begin{enumerate} 
\item D\'eterminer les sous-espaces propres et caract\'eristiques de $A$.

  \item D\'emontrer que le sous-espace vectoriel $\ker(A+I)^2$ est un plan stable par $f$.

  \item D\'emontrer qu'il existe une base de $\R^3$ dans laquelle la matrice de $f$ est 
$$B=\begin{pmatrix}-1 & 1 & 0 \\  0&-1&1 \\ 0&0&-1\end{pmatrix}$$
et trouver une matrice $P$ inversible telle que $A=PBP^{-1}$ ($AP=PB$).

  \item Ecrire la d\'ecomposition de Dunford de $B$ (justifier).

  \item Pour $t\in\R$, calculer $\exp tB$ et exprimer $\exp tA$ \`a l'aide de $P$ et $\exp tB$.

  \item Donner les solutions des syst\`emes diff\'erentiels $Y'=BY$ et $X'=AX$.
\end{enumerate}

\centerline{III}  

On suppose, dans cette partie, que $\alpha=-1$, on note $A=A_{-1}$.

\begin{enumerate}  \item V\'erifier que la matrice $A$ est diagonalisable.
 
  \item Diagonaliser la matrice $A$.

  \item Donner les solutions du syst\`eme diff\'erentiel $X'=A.X$.
\end{enumerate}

\centerline{IV}

On suppose, dans cette partie, que $\alpha=1$, on note $A=A_{1}$.

\begin{enumerate} 
  \item D\'eterminer les sous-espaces propres et caract\'eristiques de $A$.

  \item Trigonaliser la matrice $A$.
\end{enumerate}
\finenonce{002610} 


\finexercice

\exercice{2611, delaunay, 2009/05/19}
\enonce{002611}{}
 Soit $A$ une matrice $2\times 2$ \`a coefficients r\'eels.
$$A=\begin{pmatrix}a&b \\  c&d\end{pmatrix}$$
On suppose $a+c=b+d=1$ et $a-b\neq 1$. 
\begin{enumerate}
 \item Soient $(x_1,x_2)$, $(y_1,y_2)$ deux vecteurs de $\R^2$, tels que  
$$A\begin{pmatrix}x_1 \\  x_2\end{pmatrix}=\begin{pmatrix}y_1 \\  y_2\end{pmatrix}$$ 
montrer qu'alors 
$$y_1+y_2=x_1+x_2.$$

 \item  Soit le vecteur $\vec x=(1,-1)$, vérifier que $\vec x$ est un vecteur propre de $A$, et déterminer sa valeur propre.
 \item Déterminer le polynôme caractéristique de $A$ et calculer ses racines.

 \item Déterminer un vecteur propre, $\vec y$, de $A$ non colinéaire à $\vec x$ et exprimer la matrice de l'endomorphisme défini par $A$ dans la base $(\vec x,\vec y)$.
\end{enumerate}
\finenonce{002611} 


\finexercice


\exercice{2612, delaunay, 2009/05/19}
\enonce{002612}{}
Soit $E$ un espace vectoriel de dimension $3$. On note ${\cal B}=(\vec e_1,\vec e_2,\vec e_3)$ une base de $E$, si $\vec u$ est un vecteur de $E$ on note $(x,y,z)$  ses coordonnées dans la base $\cal B$. Soit 
$f$ une application linéaire de $E$ dans lui-même, définie par
\begin{align*}
f:E &\longrightarrow E \\  \begin{pmatrix}x \\  y \\  z\end{pmatrix}&\longmapsto \begin{pmatrix}x' \\  y' \\  z'\end{pmatrix}= \begin{pmatrix}-x+y-z \\  2x+2z \\  4x-2y+4z\end{pmatrix}
\end{align*}
\begin{enumerate}
 \item  Donner la matrice $A$ de $f$ dans la base $\cal B$.
 \item Déterminer les sous-espaces $\ker f$ et $\Im f$.
 \item  Soient $\vec u_1=(1,0,-1)$, $\vec u_2=(1,2,0)$ et $\vec u_3=(0,1,1)$. Démontrer que $(\vec u_1,\vec u_2,\vec u_3)$ est une base de $E$.
 \item Calculer $f(\vec u_1)$, $f(\vec u_2)$ et $f(\vec u_3)$ et déterminer la matrice $B$ de $f$ dans la base $(\vec u_1,\vec u_2,\vec u_3)$.
 \item  Déterminer les valeurs propres de $f$ et, pour chacune, un vecteur propre.
\end{enumerate}
\finenonce{002612} 


\finexercice


\exercice{2613, delaunay, 2009/05/19}
\enonce{002613}{}
 Soit $E$ un espace vectoriel de dimension $n$. On cherche à déterminer une matrice $A\in{\cal M}_n(\R)$ telle que $A^2=-I_n$, où $I_n$ désigne la matrice identité d'ordre $n$. On notera $f$ l'endomorphisme de $E$ de matrice $A$ dans la base canonique.
\begin{enumerate}
 \item  Démontrer que l'existence d'une telle matrice implique la parité de $n$.

 \item On suppose maintenant que $n=4$.
   \begin{enumerate}
   \item Démontrer que pour tout $\vec x\in E$, $\vec x\neq 0$, les vecteurs $\vec x$ et $f(\vec x)$ sont linéairement indépendants.

   \item Soit $\vec x_1\neq0$, on note $F$ le sous-espace vectoriel de $E$ engendré par les vecteurs $\vec x_1$ et $f(\vec x_1)$.
             \begin{enumerate}
               \item Démontrer que $F$ est stable par $f$.

               \item Soit $\vec x_2\in E$, on suppose que $\vec x_2\not\notin F$, démontrer que ${\cal B}=(\vec x_1,f(\vec x_1),\vec x_2,f(\vec x_2))$ est une base de $E$.
             \end{enumerate}
   \item  Ecrire la matrice $A$ de $f$ dans la base $\cal B$.
   \item Calculer $\det f$ et $\det(\lambda\mathrm{id}_E-f)$ pour $\lambda\in\R$.
   \item  L'endomorphisme $f$ admet-il des valeurs propres réelles ?
  \end{enumerate}
\end{enumerate}
\finenonce{002613} 


\finexercice


\exercice{2614, delaunay, 2009/05/19}
\enonce{002614}{}

Soient ${a\in \R}$, ${b\in \R}$ et $A$ la matrice 
$$\begin{pmatrix}
1&a&0 \\ 
0&1&b \\ 
0&0&2 \\ 
\end{pmatrix}$$

\begin{enumerate}
 \item  Donner les valeurs de $a$ et de $b$ pour lesquelles la décomposition de Dunford de $A$ est
$$A=\begin{pmatrix}
1 & 0 & 0  \\ 
0 & 1 & 0  \\ 
0 & 0 & 2  \\ 
\end{pmatrix} + 
\begin{pmatrix}
 0 & a & 0  \\ 
0 & 0 & b  \\ 
0 & 0 & 0  \\ 
\end{pmatrix} 
$$
 \item On suppose dans la suite que $b=1$ et $a\ne 0$
   \begin{enumerate}
   \item Déterminer les sous espaces propres et les sous espaces caractéristiques de $A$.
   \item  Trouver $D$ diagonalisable et $N$ nilpotente telles que $D$ commute avec $N$ et 
$$A=D+N.$$
  \end{enumerate}
 \item Soit le système différentiel suivant:

$${\cal E}: \ \ \left\{\begin{align*} x_1 '(t)&= x_1(t)+2x_2(t) \\ 
x_2'(t)&= x_2(t)+x_3(t) \\ 
x_3'(t)&= 2x_3(t) \\ 
\end{align*}\right.$$

Déterminer les solutions de ${\cal E}$.
\end{enumerate}
\finenonce{002614} 


\finexercice


\exercice{2615, delaunay, 2009/05/19}
\enonce{002615}{}
{\bf Questions préliminaires} : 
\begin{enumerate}
 \item[(a)]  Soient  $E$ un espace vectoriel réel de dimension $n$ et  $u$ un endomorphisme de $E$. Soit $P\in\R[X]$ un polynôme. Soit $\lambda$ une valeur propre de $u$ et $\vec x$ un vecteur propre associé à $\lambda$. Démontrer que $\vec x$ est vecteur propre de l'endomorphisme $P(u)$ pour la valeur propre $P(\lambda)$.

 \item[(b)] Enoncer le théorème de Hamilton-Cayley.
\end{enumerate}

 
Soit $$A=\begin{pmatrix} 1&0&0 \\ -9& 1&9 \\  9&0&-8\end{pmatrix}.$$
\begin{enumerate}
 \item  Déterminer les valeurs propres de $A$. Donner une base de vecteurs propres de $A$ et diagonaliser $A$.
 \item On cherche à déterminer une matrice $B$ telle que $B^3=A$.

\begin{enumerate}
 \item Démontrer que si $\lambda$ est une valeur propre de $B$ alors $\lambda^3$ est une valeur propre de $A$.
 \item  Déterminer les valeurs propres de $B$ et leur multiplicité.

 \item Ecrire le polynôme caractéristique de $B$.
 \item Déterminer $B$ telle que $B^3=A$. 
\end{enumerate}
\end{enumerate}
\finenonce{002615} 


\finexercice


\exercice{3496, quercia, 2010/03/10}
% %------------------------------------------------------------------ Dimension 2
% %
% \SV
\enonce{003496}{Diagonalisation en dimension $2$}
Diagonaliser les matrices suivantes :
\begin{enumerate}
\item$A=\begin{pmatrix} 1 & 5 \cr 2 & 4 \cr\end{pmatrix}$
\item $A=\begin{pmatrix} 2 & 5 \cr  4 & 3 \cr\end{pmatrix}$
\item $A=\begin{pmatrix} 5 & 3 \cr -8 &-6 \cr\end{pmatrix}$
\item $A=\begin{pmatrix} 4 & 4 \cr 1 & 4 \cr\end{pmatrix}$
   \end{enumerate}
 \finenonce{003496}
 
 
 \finexercice


\exercice{3497, quercia, 2010/03/10}
% %------------------------------------------------------------------ Dimension 3
% %
% \SV
\enonce{003497}{Diagonalisation en dimension $3$}
Diagonaliser les matrices suivantes :
\begin{enumerate}
\item $ A=\begin{pmatrix}-1 & 2   &-3   \cr  2 & 2   &-6   \cr -2 & 2   &-6   \cr\end{pmatrix}$
\item $ A=\begin{pmatrix} 0 & \phantom-1 & 2 \cr  1 & 1 & 1 \cr 1 & 0 &-1 \cr\end{pmatrix}$
\item $ A=\begin{pmatrix} 1 & 2    & 3     \cr 2 & 3 & 1 \cr 3 & 1     & 2     \cr\end{pmatrix}$
\item $ A=\begin{pmatrix} 1 & 2     & 3     \cr 2  &-1   &-4   \cr 3 & 1     & 2     \cr\end{pmatrix}$
\item $ A=\begin{pmatrix} 1 & 2 & 3 \cr  2 & 1 & 3 \cr 4 & 2 & 0 \cr\end{pmatrix}$ 
\item $  A=\begin{pmatrix} 1  & 1  & 0 \cr -1  & 1  & 2 \cr 0  & 0  & 2 \cr\end{pmatrix}$ 
\item $  A=\begin{pmatrix} 2 &-1 &-1 \cr -1 & 2 &-1 \cr -1 &-1 & 2 \cr\end{pmatrix}$
\item $ A=\begin{pmatrix} 1 &-1 & 2 \cr 3 &-3 & 6 \cr 2 &-2 & 4 \cr\end{pmatrix}$
\item $ A=\begin{pmatrix} 7 &-12 &-2 \cr  3 &-4 & 0 \cr -2 & 0 &-2 \cr\end{pmatrix}$ 
\item$ A=\begin{pmatrix}-2 & 8 & 6 \cr -4 & 10 & 6 \cr 4 &-8 &-4 \cr\end{pmatrix}$
\end{enumerate}
 \finenonce{003497}
 
 
 \finexercice


\exercice{3498, quercia, 2010/03/10}
\enonce{003498}{Diagonalisation en dimension $4$}
Diagonaliser les matrices suivantes :
\begin{enumerate}
\item $\begin{pmatrix} 0 & 1 & 0 & 0 \cr  3 & 0 & 2 & 0 \cr 0 & 2 & 0 & 3 \cr 0 & 0 & 1 & 0 \cr\end{pmatrix}$
\item $\begin{pmatrix} 1 & 1 & 1 & 1 \cr 1 & 1 &-1 &-1 \cr 1 &-1 & 1 &-1 \cr 1 &-1 &-1 & 1 \cr\end{pmatrix}$
\item $\begin{pmatrix} 0 & 2 &-2 & 2 \cr -2 & 0 & 2 & 2 \cr -2 & 2 & 0 & 2 \cr 2 & 2 &-2 & 0 \cr\end{pmatrix}$
\item $\begin{pmatrix}-5  &  2 &  0 &  \phantom-0 \cr 0  &-11 &  5 &  0 \cr 0  &  7 & -9 &  0 \cr 0  &  3 &  1 &  2 \cr\end{pmatrix}$ 
\item $ \begin{pmatrix} 2 & 0 & 3 & 4 \cr 3 & 1 & 2 & 1 \cr 0 & 0 & 3 & 0 \cr 0 & 0 & 4 &-1 \cr\end{pmatrix}$
\item $\begin{pmatrix} 1 &-3 & \phantom-0 & 3 \cr -2 &-6 & 0 & 13 \cr 0 &-3 & 1 & 3 \cr -1 &-4 & 0 & 8 \cr\end{pmatrix}$
\item $\begin{pmatrix} 3 &-1 & 1 &-7 \cr 9 &-3 &-7 &-1 \cr 0 & 0 & 4 &-8 \cr 0 & 0 & 2 &-4 \cr\end{pmatrix}$
\item $\begin{pmatrix} 0 & 0 & 2 & 3 \cr 0 & 0 &-2 &-3 \cr 2 &-2 & 0 &-1 \cr 3 &-3 &-1 &-3 \cr\end{pmatrix}$ 
\end{enumerate}
 \finenonce{003498}
 
 
 \finexercice


\exercice{3503, quercia, 2010/03/10}
\enonce{003503}{Diagonalisation}

Diagonaliser $M = \begin{pmatrix} (0) & &1 \cr &\cdots \cr 1& &(0) \cr \end{pmatrix}
\in \mathcal{M}_n(K)$.

\finenonce{003503}



\finexercice


\exercice{3504, quercia, 2010/03/10}
\enonce{003504}{Diagonalisation}

Diagonaliser $M = \begin{pmatrix} 0 &       &       &1  \cr
                            1 &\ddots &(0)        \cr
                              &\ddots &\ddots     \cr
                            (0) &     &1      &0  \cr \end{pmatrix} \in \mathcal{M}_n(\C)$.


\finenonce{003504}



\finexercice


\exercice{3505, quercia, 2010/03/10}
\enonce{003505}{Engees 93}

Diagonaliser la matrice
$M = \begin{pmatrix} e &a &b &c \cr
               a &e &c &b \cr
               b &c &e &a \cr
               c &b &a &e \cr\end{pmatrix} \in \mathcal{M}_4(\R)$.

\finenonce{003505}



\finexercice

\exercice{3507, quercia, 2010/03/10}
\enonce{003507}{Matrice triangulaire}

Soit $A = \begin{pmatrix} 1 &a &b &c \cr
                    0 &1 &d &e \cr
                    0 &0 &-1&f \cr
                    0 &0 &0 &-1\cr \end{pmatrix}$.
A quelle condition $A$ est-elle diagonalisable ?
\finenonce{003507}



\finexercice

\exercice{3510, quercia, 2010/03/10}
\enonce{003510}{$(X^2-1)P'' + (2X+1)P'$}

Soit $E =  K_n[X]$ et $ u : E \to E, P \mapsto {(X^2-1)P''+(2X+1)P'.}$

\begin{enumerate}
  \item Chercher la matrice de $u$ dans la base canonique de $ K_n[X]$.
    
  \item Montrer que $u$ est diagonalisable.
\end{enumerate}
\finenonce{003510}



\finexercice

\exercice{3517, quercia, 2010/03/10}
\enonce{003517}{$\mathrm{tr}(A)M+\mathrm{tr}(M)A$}

Soit $A\in\mathcal{M}_n(\C)$. L'endomorphisme $f$ de $\mathcal{M}_n(\C)$ défini par
$f(M) = \mathrm{tr}(A)M+\mathrm{tr}(M)A$ est-il diagonalisable~?

\finenonce{003517}



\finexercice


\exercice{5654, rouget, 2010/10/16}
\enonce{005654}{***}
Soit $f$ qui à $P$ élément de $\Rr_{2n}[X]$ associe $f(P)=(X^2-1)P'- 2nXP$.

Vérifier que $f$ est un endomorphisme de $\Rr_{2n}[X]$ puis déterminer les valeurs et vecteurs propres de $f$. $f$ est-il diagonalisable ?
\finenonce{005654}


\finexercice


\exercice{5663, rouget, 2010/10/16}
\enonce{005663}{****}
Soient $A$ un élément de $\mathcal{M}_n(\Cc$) et $M$ l'élément de $\mathcal{M}_{2n}(\Cc)$ défini par blocs par $M=\left(
\begin{array}{cc}
A&4A\\
A&A
\end{array}
\right)$.

Calculer $\text{det}M$. Déterminer les éléments propres de $M$ puis montrer que $M$ est diagonalisable si et seulement si $A$ est diagonalisable.
\finenonce{005663}


\finexercice


\exercice{5667, exo7, 2010/10/16}
\enonce{005667}{**}

\newcommand{\adots}{\mathinner{\mkern2mu\raise 1pt\hbox{.}\mkern 3mu\raise
4pt\hbox{.}\mkern1mu\raise 7pt\hbox{{.}}}}


Soit $A=\left(
\begin{array}{cccc}
0&\ldots&0&1\\
\vdots&\adots&\adots&0\\
0& \adots & \adots &\vdots\\
1&0&\ldots&0
\end{array}
\right)$

 Montrer que $A$ est diagonalisable.
\finenonce{005667}


\finexercice


\exercice{5673, rouget, 2010/10/16}
\enonce{005673}{***}
Soit $A$ une matrice carrée de format $2$ telle que $A^2$ est diagonalisable et $\text{Tr}A\neq 0$. Montrer que $A$ est diagonalisable dans $\Cc$.
\finenonce{005673}


\finexercice

\exercice{5675, rouget, 2010/10/16}
\enonce{005675}{***I}
Sur $E$ un $\Rr$-espace vectoriel. On donne trois endomorphismes $f$, $u$ et $v$ tels qu'il existe deux réels $\lambda$ et $\mu$ tels que pour $k\in\{1,2,3\}$, $f^k =\lambda^ku +\mu^kv$. Montrer que $f$ est diagonalisable.
\finenonce{005675}


\finexercice


\exercice{5682, rouget, 2010/10/16}
\enonce{005682}{**}
Soit $A=\left(
\begin{array}{cccc}
0&\ldots&0&a_1\\
\vdots& &\vdots&\vdots\\
0&\ldots&0&a_{n-1}\\
a_1&\ldots&a_{n-1}&a_n
\end{array}
\right)$ où $a_1$,..., $a_n$ sont $n$ nombres complexes ($n\geqslant2$). $A$ est-elle diagonalisable?
\finenonce{005682}


\finexercice

\section{201.03 Polynôme caractéristique, théorème de Cayley-Hamilton}


\exercice{3524, quercia, 2010/03/10}
\enonce{003524}{Formules pour une matrice $3\times3$}

Soit $A = (a_{ij}) \in \mathcal{M}_{3}(\R)$.

\begin{enumerate}
  \item Vérifier que
    $\chi_A(\lambda) = -\lambda^3 + (\mathrm{tr} A)\lambda^2
       - \left( \begin{vmatrix} a_{11} &a_{12} \cr a_{21} &a_{22} \end{vmatrix}
              + \begin{vmatrix} a_{11} &a_{13} \cr a_{31} &a_{33} \end{vmatrix}
              + \begin{vmatrix} a_{22} &a_{23} \cr a_{32} &a_{33} \end{vmatrix} \right)\lambda
       +\det(A)$.

  \item Soit $\lambda$ une valeur propre de $A$ et $L_1,L_2$ deux lignes non
    proportionnelles de $A-\lambda I$ (s'il en existe).
    On calcule $L = L_1\wedge L_2$ (produit vectoriel) et $X = {}^tL$.
    Montrer que $X$ est vecteur propre de $A$ pour la valeur propre $\lambda$.

\end{enumerate}
\finenonce{003524}



\finexercice


\exercice{3525, quercia, 2010/03/10}
\enonce{003525}{Recherche de vecteurs propres pour une valeur propre simple}

Soit $A \in \mathcal{M}_n(K)$ et $\lambda \in  K$ une valeur propre de $A$ telle que
$\mathrm{rg}(A-\lambda I) = n-1$.

\begin{enumerate}
  \item Quelle est la dimension du sous espace propre $E_\lambda$ ?
  \item Montrer que les colonnes de ${}^t\text{com}(A - \lambda I)$ engendrent
    $E_\lambda$.
  \item Exemple : diagonaliser $A = \begin{pmatrix} 0 & \phantom-1 &  2 \cr
                                       1 & 1          &  1 \cr
                                       1 & 0          & -1 \cr \end{pmatrix}$.
\end{enumerate}
\finenonce{003525}



\finexercice


\exercice{3526, quercia, 2010/03/10}
\enonce{003526}{\'Eléments propres de $C{}^tC$}

Soit $C = \begin{pmatrix} a_1\cr \vdots\cr a_n \end{pmatrix} \in \mathcal{M}_{n,1}(\R)$ et $M = C{}^tC$.

\begin{enumerate}
  \item Chercher le rang de $M$.
    
  \item En déduire le polynôme caractéristique de $M$.
    
  \item $M$ est-elle diagonalisable ?
    
\end{enumerate}
\finenonce{003526}



\finexercice


\exercice{3527, quercia, 2010/03/10}
\enonce{003527}{Ensi Chimie P 94}
    Soit $A=(a_{ij}) \in \mathcal{M}_n(\R)$ telle que $a_{ij} = \frac ij$.
    $A$ est-elle diagonalisable~?
    
\finenonce{003527}



\finexercice


\exercice{3528, quercia, 2010/03/10}
\enonce{003528}{Ensi Chimie P 93}

On considère le polynôme défini par :
$\forall\ n \in \N,\ P_n(x) = x^n - \sum_{i=0}^{n-1} \alpha_ix^i$ avec
$\alpha_0 > 0$ et $\alpha_i \ge 0$ pour $1 \le i \le n-1$.

\begin{enumerate}
  \item Montrer qu'il existe une unique racine dans $\R^{+*}$ pour $P_n$.
    

  \item Soit $A = \begin{pmatrix} 1       &1      &0      &\dots  &0      \cr
                        2       &0      &1      &\ddots &\vdots \cr
                        \vdots  &\vdots &\ddots &\ddots &0      \cr
                        \vdots  &\vdots &       &\ddots &1      \cr
                        n       &0      &\dots  &\dots  &0      \cr\end{pmatrix}$.
    Montrer que $A$ admet une unique valeur propre réelle strictement positive.
    
\end{enumerate}
\finenonce{003528}



\finexercice


\exercice{3529, quercia, 2010/03/10}
\enonce{003529}{Ensi Physique 93}

Soient $x_1,\dots,x_n,y_1,\dots,y_n \in \C$.
Calculer $\Delta_n = \det(I+(x_iy_j))$.

\finenonce{003529}



\finexercice


\exercice{3530, quercia, 2010/03/10}
\enonce{003530}{Centrale MP 2000}

On considère la matrice de $M_n(\C)$, $A=\begin{pmatrix}0      & a      & \dots  & a\cr
                                                    b      & \ddots & \ddots & \vdots \cr
                                                    \vdots & \ddots & \ddots & a\cr
                                                    b      & \dots  & b      & 0\cr\end{pmatrix}$, $a\ne b$.
\begin{enumerate}
  \item Montrer que le polynôme caractéristique de $A$ est $\frac{(-1)^n}{a-b}(a(X+b)^n-b(X+a)^n)$.
    
  \item Montrer qu'en général les valeurs propres de $A$ sont sur un cercle.
    
\end{enumerate}
\finenonce{003530}



\finexercice


\exercice{3531, quercia, 2010/03/10}
\enonce{003531}{Centrale MP 2000}

Soit $a_1,\dots,a_n,b_1,\dots,b_n \in \R$ et
$A_n=\begin{pmatrix}a_1+b_1 & b_2 &\dots &\dots &b_n\cr
b_1    &a_2+b_2 & b_3           &\dots &b_n\cr
\vdots &b_2 &\ddots &\vdots  &\vdots\cr
\vdots &\vdots &     & \ddots& \vdots\cr
b_1  & b_2 & \dots &b_{n-1} & a_n+b_n\cr\end{pmatrix}$
\begin{enumerate}
  \item Calculer $\det A_n$.
    
  \item Calculer $\chi_A$ le polynôme caractéristique de $A$.
  \item On suppose $a_1<a_2<\dots <a_n$ et, pour tout $i$, $b_i>0$. Montrer que $A_n$ est diagonalisable 
    $\Bigl($on pourra utiliser ${\chi_A(t)}/{\prod_{i=1}^n(a_i-t)}\Bigr)$.
    
  \item Le résultat reste-t-il vrai si l'on suppose $a_1\le a_2\le \dots \le a_n$ et, pour tout $i$, 
    $b_i>0$ ?
    
\end{enumerate}
\finenonce{003531}



\finexercice


\exercice{3532, quercia, 2010/03/10}
\enonce{003532}{Polynômes caractéristiques}

Soit $A \in \mathcal{M}_n(K)$ inversible et $B = A^{-1}$, $C = A^2$.
Exprimer les polynômes caractéristiques $\chi_B$ et $\chi_C$ en fonction de $\chi_A$.

\finenonce{003532}



\finexercice


\exercice{3533, quercia, 2010/03/10}
\enonce{003533}{Matrice compagnon}

Soit $P = a_0 + a_1X + \dots + a_{n-1}X^{n-1} - X^n \in  K_n[X]$.
La matrice compagnon de $P$ est
$M = \begin{pmatrix} 0 &\span (0)    &a_0      \cr
               1 &\ddots &     &a_1      \cr
                 &\ddots &0    &\vdots   \cr
              (0)&       &1    &a_{n-1}  \cr \end{pmatrix}$.

Soit $E$ un $ K$-ev de dimension $n$, ${\cal B} = ({\vec e}_1,\dots,{\vec e}_n)$
une base de $E$ et $\varphi$ l'endomorphisme de $E$ de matrice $M$ dans~$\cal B$.

\begin{enumerate}
  \item Déterminer le polynôme caractéristique de $M$.
  \item Calculer $\varphi^k({\vec e}_1)$ pour $0 \le k \le n$.
  \item En déduire que $P(M) = 0$.
\end{enumerate}
\finenonce{003533}



\finexercice


\exercice{3534, quercia, 2010/03/10}
\enonce{003534}{Matexo}

Soient $A,B \in \mathcal{M}_n(\C)$. Montrer que $\text{spec}(A) \cap \text{spec}(B) = \varnothing$
si et seulement si $\chi_A(B)$ est inversible.

Application~: Soient $A,B,P$ trois matrices carrées complexes avec $P\ne 0$
telles que $AP = PB$. Montrer que $A$ et $B$ ont une valeur popre commune.
\finenonce{003534}



\finexercice


\exercice{3535, quercia, 2010/03/10}
\enonce{003535}{Matrices à spectres disjoints}

Soient $A,B\in\mathcal{M}_n(\C)$. Montrer l'équivalence entre~:

(a)~: $\forall\ C\in\mathcal{M}_n(\C)$, il existe un unique $X\in\mathcal{M}_n(\C)$ tel que $AX-XB=C$.

(b)~: $\forall\ X\in\mathcal{M}_n(\C)$ on a $AX = XB  \Rightarrow  X=0$.

(c)~: $\chi_B(A)$ est inversible.

(d)~: $A$ et $B$ n'ont pas de valeur propre en commun.
\finenonce{003535}



\finexercice


\exercice{3536, quercia, 2010/03/10}
\enonce{003536}{$AB$ et $BA$ ont même polynôme caractéristique}

Soient $A,B \in \mathcal{M}_n(K)$.

\begin{enumerate}
  \item Montrer que $AB$ et $BA$ ont les mêmes valeurs propres.

  \item Montrer que si $A$ ou $B$ est inversible, alors $AB$ et $BA$ ont même
    polynôme caractéristique.

  \item Dans le cas général, on note
    $M = \begin{pmatrix} BA &-B\cr 0 &0  \end{pmatrix}$,
    $N = \begin{pmatrix} 0  &-B\cr 0 &AB \end{pmatrix}$,
    $P = \begin{pmatrix} I_n  & 0\cr A &I_n \end{pmatrix}$
    ($M,N,P \in \mathcal{M}_{2n}(K)$).

 Vérifier que $MP$ = $PN$, montrer que $P$ est inversible, et conclure.
\end{enumerate}
\finenonce{003536}



\finexercice


\exercice{3537, quercia, 2010/03/10}
\enonce{003537}{$\det(I+A\overline A) \ge 0$}

Soit $A \in \mathcal{M}_n(\C)$.

\begin{enumerate}
  \item Soit $\lambda \in \C$ une valeur propre non nulle de $A\overline A$, et $X$ un
    vecteur propre associé. Montrer que $A\overline X$ est aussi vecteur propre de
    $A\overline A$.
  \item Lorsque $\lambda \notin \R^+$, montrer que $X$ et $A\overline X$ sont
    linéairement indépendants.
  \item En déduire que $\det(I + A\overline A ) \in \R^+$.
\end{enumerate}
\finenonce{003537}



\finexercice


\exercice{3538, quercia, 2010/03/10}
\enonce{003538}{Centrale PC 1999}

Soit l'application
$\Phi : {\mathcal{M}_n(\R)} \to {\mathcal{M}_n(\R)}, M \mapsto {^tM.}$
Calculer sa trace par un moyen simple.


\finenonce{003538}



\finexercice


\exercice{3539, quercia, 2010/03/10}
\enonce{003539}{Multiplicité d'une valeur propre}

Soit $E$ un $ K$-ev de dimension finie, $u\in \mathcal{L}(E)$ et $\lambda\in\text{Sp}(u)$.
Montrer que la multiplicité de $\lambda$ dans le polynôme minimal de~$u$
est égale au nombre de facteurs invariants de~$u$ ayant $\lambda$ pour racine.
\finenonce{003539}



\finexercice


\exercice{3540, quercia, 2010/03/10}
\enonce{003540}{Fermat pour la trace, ULM-Lyon-Cachan MP$^*$ 2005}

Soit $p$ premier et $A\in\mathcal{M}_n(\Z)$. Montrer que $\mathrm{tr}(A^p)\equiv \mathrm{tr}(A) \pmod p$.

\finenonce{003540}



\finexercice

\exercice{5656, rouget, 2010/10/16}
\enonce{005656}{***}
Soit $A$ une matrice rectangulaire de format $(p,q)$ et $B$ une matrice de format $(q,p)$. Comparer les polynômes caractéristiques de $AB$ et $BA$.
\finenonce{005656}


\finexercice


\exercice{5657, rouget, 2010/10/16}
\enonce{005657}{*** I}
\label{ex:rou7}
Soient $u$ et $v$ deux endomorphismes d'un espace vectoriel de dimension finie. On suppose que $u$ et $v$ commutent et que $v$ est nilpotent. Montrer que $\text{det}(u+v)=\text{det}u$.
\finenonce{005657}


\finexercice


\exercice{5666, rouget, 2010/10/16}
\enonce{005666}{**}
\label{ex:rou16}
Soit $A$ une matrice antisymétrique réelle. Etudier la parité de son polynôme caractéristique.
\finenonce{005666}


\finexercice


\exercice{5678, rouget, 2010/10/16}
\enonce{005678}{**}
Soient $A$ et $B$ deux matrices carrées complexes de format $n$.
Montrer que $A$ et $B$ n'ont pas de valeurs propres communes si et seulement si la matrice $\chi_A(B)$ est inversible.
\finenonce{005678}


\finexercice


\exercice{5679, rouget, 2010/10/16}
\enonce{005679}{**}
Soit $f$ un endomorphisme d'un $\Kk$-espace vectoriel de dimension finie et $P$ un polynôme. Montrer que $P(f)$ est inversible si et seulement si $P$ et $\chi_f$ sont premiers entre eux.
\finenonce{005679}


\finexercice

\section{201.04 Sous-espace stable}


\exercice{1675, ortiz, 1999/04/01}

\enonce{001675}{}
Soit l'endomorphisme $f$ de $\Rr^3$ canoniquement
associ\'e \`a la matrice
$M=\left(\begin{smallmatrix}
1&1&0\\-1&2&1\\1&0&1\end{smallmatrix}\right)$. Le
plan $P$ d'\'equation $y+z=0$ est-il stable par
$f$ ? La droite $\text{vect}\left\{(1,1,1)\right\}$
est-elle stable par $f$ ?
\finenonce{001675}



\finexercice



\exercice{1676, ortiz, 1999/04/01}

\enonce{001676}{}
que $f^3+f^2+f=0$ o\`u $E$ est un $\Rr$-espace
vectoriel de dimension finie et soit
$F=\text{Im\,} f$.
\begin{enumerate}
    \item
    \begin{enumerate}
        \item Montrer que $F$ est un sous-espace vectoriel  stable par $f.$
        \item Montrer que $\ker f\cap \text{Im\,} f=\left\{0\right\}$.
        \item En d\'eduire que la restriction $g$ de $f$ \`a $F$ est un automorphisme de $F$.
    \end{enumerate}
    \item
    \begin{enumerate}
         \item Montrer que si $\lambda \in\text{Sp}_{\Rr}(f)$ alors $\lambda =0$.
         \item En d\'eduire que le rang de $f$ est pair
(raisonner par l'absurde et \'etudier les racines
r\'eelles du polyn\^ome caract\'eristique de $g$).
\end{enumerate}
\end{enumerate}
\finenonce{001676}



\finexercice



\exercice{1677, ortiz, 1999/04/01}

\enonce{001677}{}
Soient $f\in\mathcal{L}(E)$ et $a\in E.$
\begin{enumerate}
    \item  Montrer que le plus petit sous-espace vectoriel de $E$ contenant $a$ et
stable par $f $ est
$F_a=\text{vect}\left\{f^k(a):k\in\Nn\right\}$.
    \item  Montrer que si $\dim(E)=n$ alors $F_a=\text{vect}\left\{f^k(a):k=0,...,n-1\right\}.$
    \item  Soit l'endomorphisme $f$ de $\Rr^3$ canoniquement associ\'e \`a la
matrice $A=\left( \begin{smallmatrix}
1&1&-1\\
-2&-2&2\\
1&1&-1
\end{smallmatrix}\right) \in \mathcal{M}_3(\Rr)$.
Montrer qu'il n'existe pas $a\in \Rr^3$ tel que
$F_a=\Rr^3.$ G\'en\'eraliser \`a un endomorphisme
diagonalisable.
\end{enumerate}
\finenonce{001677}



\finexercice



\exercice{1678, ortiz, 1999/04/01}

\enonce{001678}{}
Soient $f\in\mathcal{L}(E)$, $F$ un sous-espace
vectoriel de $E$ stable par $f$ et $g$
l'endomorphisme de $F$ induit par $f$.
\begin{enumerate}
    \item  Montrer que si $P\in K[X]$ v\'erifie $P(f)=0$ alors $P(g)=0$.
    \item  En d\'eduire que si $f$ est diagonalisable alors $g$ est diagonalisable.
    \item  Application : trouver tous les sous-espaces vectoriels
stables par l'endomorphisme $f$ de $\Rr^3$ canoniquement associ\'e \`a
la matrice $A=\left( \begin{smallmatrix}
1&1&-1\\
-2&-2&2\\
1&1&-1
\end{smallmatrix}\right) \in \mathcal{M}_3(\Rr)$.
\end{enumerate}
\finenonce{001678}



\finexercice



\exercice{1679, ortiz, 1999/04/01}

\enonce{001679}{}
\begin{enumerate}
     \item Montrer que $A=\left( \begin{smallmatrix}
3&2&4\\
-1&3&-1\\
-2&-1&-3
\end{smallmatrix}\right) \in \mathcal{M}_3(\Rr)$ est trigonalisable.
$A$ est-elle diagonalisable ? R\'eduire $A$ et
d\'eterminer son polyn\^ome minimal.
      \item M\^eme question pour $A=\left( \begin{smallmatrix}
2&-1&2\\
5&-3&3\\
-1&0&-2
\end{smallmatrix}\right) \in \mathcal{M}_3(\Rr)$.
\end{enumerate}
\finenonce{001679}



\finexercice



\exercice{1680, ortiz, 1999/04/01}

\enonce{001680}{}
Quel est le polyn\^ome caract\'eristique d'un
endomorphisme nilpotent d'un $\Cc$-espace
vectoriel de dimension finie ?
\finenonce{001680}



\finexercice



\exercice{1681, ortiz, 1999/04/01}

\enonce{001681}{}
Soit $A\in\mathcal{M}_n(\Rr)$ et soient
$\lambda_1,...,\lambda_n$ ses valeurs propres
complexes.
 Exprimer $\text{tr}(A^p)$ o\`u $p\in\Nn$ en fonction des $\lambda_j,$ $j=1,...,n.$
\finenonce{001681}



\finexercice



\exercice{1682, ortiz, 1999/04/01}

\enonce{001682}{}
Soient $f$ et $g$ deux endomorphismes de $E$ tels
que $fg=gf$.
\begin{enumerate}
    \item  Soit $x\in E$. Montrer que si $n\in\Nn$ et $f(x)=g(x)$ alors $f^n(x)=g^n(x)$.
\par
{\textit{Dans toute la suite, on suppose $g$ nilpotent.}}
    \item
    \begin{enumerate}
          \item D\'eduire de {1.} que si $f$ est inversible alors $f+g$ est inversible.
          \item D\'eduire de {(a)} que si  $f+g$ est inversible alors $f$ est inversible.
    \end{enumerate}
    \item
    \begin{enumerate}
        \item Soit $h\in\mathcal{L}(E)$ nilpotent. Montrer que $\det(h+\text{Id}_E)=1$.
        \item Montrer que $\det(f+g)=\det(f)$  (on distinguera selon que $f$ est inversible
ou non et on utilisera les questions
pr\'ec\'edentes.
\end{enumerate}
\end{enumerate}
\finenonce{001682}



\finexercice



\exercice{1683, ortiz, 1999/04/01}

\enonce{001683}{}
Soient $E$ un $K$-espace vectoriel, $f$ et $g$ des
endomorphismes de $E$ tels que $f\circ g=g\circ f$
et $P$ un polyn\^ome de $K[X].$
\begin{enumerate}
    \item  Montrer que $P(g)$ et $f$ commutent.
    \item  Montrer que le noyau et l'image de l'endomorphisme $P(g)$ sont stables par $f$. Donner des cas particuliers de cette situation.
\end{enumerate}
\finenonce{001683}



\finexercice



\exercice{1684, ortiz, 1999/04/01}

\enonce{001684}{}
Soient $E$ un $K$-espace vectoriel de dimension
finie, $f$ un endomorphisme de $E$ et $F$ un
sous-espace vectoriel de $E$ stable par $f.$ On
d\'esigne par $g$ l'endomorphisme de  $F$ induit
par $f$ sur $F$.
\begin{enumerate}
    \item  Montrer que $\text{Sp}(g)\subseteq \text{Sp}(f)$.
    \item  Montrer que si $P(f)=0$ alors $P(g)=0$. En d\'eduire que le polyn\^ome minimal de $g$ divise celui de $f$.
\end{enumerate}
\finenonce{001684}



\finexercice



\exercice{1685, ortiz, 1999/04/01}

\enonce{001685}{}
Soient $E$ un $\Rr$-espace vectoriel de dimension
finie, $f$ un endomorphisme de $E$. Montrer que si
$f$ admet un sous-espace vectoriel propre de
dimension $p\geq 2$ alors il admet une infinit\'e
de sous-espaces vectoriels stables par $f$.
\finenonce{001685}



\finexercice



\exercice{3609, quercia, 2010/03/10}
\enonce{003609}{Droites et hyperplans stables}

Soit $E$ un $\C$-ev de dimension finie et $u \in \mathcal{L}(E)$.

\begin{enumerate}
  \item Montrer qu'il existe une droite vectorielle stable par $u$.
  \item Montrer qu'il existe un hyperplan stable par $u$
    (considérer $\Im(u-\lambda\mathrm{id})$ où $\lambda$ est une valeur propre de $u$).
  \item Donner un exemple où ces propriétés sont en défaut pour un $\R$-ev.
\end{enumerate}
\finenonce{003609}



\finexercice


\exercice{3610, quercia, 2010/03/10}
\enonce{003610}{Plan stable pour une valeur propre non réelle}

Soit $M \in \mathcal{M}_n(\R)$ et $\lambda = a+ib$ une valeur propre non réelle de $M$
($a \in \R$, $b \in \R^*$).
On note X un vecteur propre complexe de $M$.

\begin{enumerate}
  \item Montrer que $\overline X$ est aussi vecteur propre de $M$.
  \item Montrer que $(X,\overline X)$ est libre dans $\C^n$.
  \item Soient $U = \frac12(X+\overline X)$, $V = \frac1{2i}(X-\overline X)$.
    Montrer que $(U,V)$ est libre dans $\R^n$.
\end{enumerate}
\finenonce{003610}



\finexercice


\exercice{3611, quercia, 2010/03/10}
\enonce{003611}{Plans stables}

Soit $E$ un $ K$-ev de dimension finie et $f\in\mathcal{L}(E)$.

\begin{enumerate}
  \item Soit $F$ un plan vectoriel. Montrer que si $F$ est stable par~$f$ alors
    il existe $P\in K_2[x]$ non nul tel que $F \subset \mathrm{Ker} P(f)$.
  \item Réciproquement, si $P\in K_2[x]$ est irréductible, montrer que $\mathrm{Ker} P(f)$
    contient un plan stable par~$f$.
  \item Si $ K = \R$ montrer que $f$ admet toujours une droite ou un plan stable.
\end{enumerate}
\finenonce{003611}



\finexercice


\exercice{3612, quercia, 2010/03/10}
\enonce{003612}{Recherche de sev stables}

Soit $A = \begin{pmatrix} 1 &1 &\phantom-0 \cr -3 &-2 &0 \cr 0 &0 &1 \cr \end{pmatrix}$.

\begin{enumerate}
  \item Déterminer les sev de $\R^3$ stables pour l'endomorphisme associé à $A$.
    

  \item Quelles sont les matrices réelles commutant avec $A$ ?

\end{enumerate}
\finenonce{003612}



\finexercice


\exercice{3613, quercia, 2010/03/10}
\enonce{003613}{$u$ diagonalisable $ \Rightarrow $ diagonalisable dans un sev stable}

Soit $E$ un $ K$-ev de dimension finie et $u \in \mathcal{L}(E)$ diagonalisable.
Montrer que si $F$ est un sev stable par $u$, alors $u_{|F}$ est diagonalisable.

\finenonce{003613}



\finexercice


\exercice{3614, quercia, 2010/03/10}
\enonce{003614}{Plan affine stable}

Soit $E = \R^3$ et $H:x+2y+3z=1$ un plan {\it affine\/} de $E$.
Montrer que si $H$ est stable par $f\in\mathcal{L}(E)$ alors $1$ est
valeur propre de~$f$.


\finenonce{003614}



\finexercice


\exercice{3615, quercia, 2010/03/10}
\enonce{003615}{Endomorphisme cyclique}

Soit $f \in \mathcal{L}(E)$ un endomorphisme cyclique et $F$ un sous-espace
vectoriel stable par $f$. Montrer que $f_{|F}$ est aussi cyclique.
\finenonce{003615}



\finexercice


\exercice{3616, quercia, 2010/03/10}
\enonce{003616}{Endomorphismes semi-simples.}

Un endomorphisme $f$ est dit semi-simple si tout sous-espace stable par~$f$
admet un supplémentaire stable par~$f$.
Montrer qu'un endomorphisme d'un $\C$-ev de dimension finie est semi-simple
si et seulement s'il est diagonalisable.
\finenonce{003616}



\finexercice


\exercice{3617, quercia, 2010/03/10}
\enonce{003617}{$\chi_u$ irréductible}

Soit $u$ un endomorphisme de $E$, espace vectoriel de dimension $n$ sur le corps $ K$.
Montrez que seuls $\{0\}$ et $E$ sont stables par $u$ si et seulement si $\chi_u$ est
irréductible sur $ K$.

\finenonce{003617}



\finexercice


\exercice{3618, quercia, 2010/03/10}
\enonce{003618}{Polytechnique MP$^*$ 2000}
Soit $E$ un espace vectoriel de dimension finie, $f$ un endomorphisme de~$E$
tel que tout sous-espace de~$E$ admette un supplémentaire stable par~$f$.
Que peut-on dire de~$f$~? Réciproque~?

\finenonce{003618}



\finexercice


\exercice{3619, quercia, 2010/03/10}
\enonce{003619}{Sous-espaces stables (Centrale MP 2003)}
Soit $f\in\mathcal{L}(\R^3)$ ayant pour matrice $M=\left(\begin{smallmatrix}1&1&1\cr1&1&1\cr-1&1&1\cr\end{smallmatrix}\right)$
dans la base canonique de~$\R^3$. Déterminer les sous-espaces de~$\R^3$ stables par~$f$.
\finenonce{003619}



\finexercice


\exercice{5683, rouget, 2010/10/16}
\enonce{005683}{***}
Soit $f$ l'endomorphisme de $\Rr^3$ dont la matrice dans la base canonique est $A$. Trouver les sous espaces stables par $f$ dans chacun des cas suivants :

\begin{enumerate}
 \item  $A =\left(
\begin{array}{ccc}
1&1&-1\\
1&1&1\\
1&1&1
\end{array}
\right)$
\item  $A=\left(
\begin{array}{ccc}
2&2&1\\
1&3&1\\
1&2&2
\end{array}
\right)$
\item  $A=\left(
\begin{array}{ccc}
6&-6&5\\
-4&-1&10\\
7&-6&4
\end{array}
\right)$.
\end{enumerate}
\finenonce{005683}


\finexercice


\exercice{5686, rouget, 2010/10/16}
\enonce{005686}{**}
\label{ex:rou36}
Soit $f$ un endomorphisme d'un espace vectoriel $E$ de dimension finie non nulle et $F$ un sous-espace non nul de $E$ stable par $f$. On suppose que $f$ est diagonalisable. Montrer que la restriction de $f$ à $F$ est un endomorphisme diagonalisable de $F$.
\finenonce{005686}


\finexercice

\section{201.05 Trigonalisation}


\exercice{1686, ortiz, 1999/04/01}

\enonce{001686}{}
Trigonaliser les matrices r\'eelles suivantes :
\begin{enumerate}
    \item $A=\left(\begin{smallmatrix}
-2&1&1\\
8&1&-5\\
4&3&-3
\end{smallmatrix}\right),$\quad
    \item $B=\left(\begin{smallmatrix} 3&2&-2\\-1&0&1\\1&1&0\end{smallmatrix}\right) .$
\end{enumerate}
\finenonce{001686}



\finexercice



\exercice{1687, bodin, 1999/11/01}

\enonce{001687}{}
Mettre sous forme triangulaire les matrices suivantes :
\[  \begin{array}{ll}
\left( \begin{array}{ccc} 4 & 2 & -2 \\ 1 & 5 & -1
\\ 1 & 3 & 1 \end{array}\right); \quad & \qquad
\frac{1}{2} \left( \begin{array}{ccc} 0 & 2 & 2 \\
1 & 3 & -1 \\ -1 & 3 & 3
\end{array}\right). \end{array} \]
\finenonce{001687}



\finexercice



\exercice{1688, ortiz, 1999/04/01}

\enonce{001688}{}
Soient les matrices \`a coefficients r\'eels
suivantes
$${ A=\left(\begin{smallmatrix}
-2&-3&2\\
1&2&-2\\
2&4&-3
\end{smallmatrix}\right),\quad
B
=\left(\begin{smallmatrix}
1&-1&1&0&1\\
0&0&1&1&0\\
0&-1&2&0&1\\
0&0&0&1&-2\\
0&0&0&2&-3
\end{smallmatrix}\right)
\quad
C
=\left(\begin{smallmatrix}
0&1&0&2\\
-3&0&4&0\\
0&1&0&3\\
-1&0&1&0\\
\end{smallmatrix}\right).
}$$
\begin{enumerate}
    \item  Trigonaliser les matrices $A$, $B$ et $C.$
    \item  D\'eterminer le polyn\^ome minimal de $A$, $B$ et $C.$
\end{enumerate}
\finenonce{001688}



\finexercice



\exercice{1689, ortiz, 1999/04/01}

\enonce{001689}{}
 Soit $f$ l'endomorphisme de l'espace vectoriel canonique $\Rr^3$
 dont la matrice dans la base canonique $\mathcal{B}$ est
$${A=\left(\begin{smallmatrix}
1&1&-1\\-1&3&-3\\-2&2&-2\end{smallmatrix}\right).}$$
\begin{enumerate}
    \item  Montrer que $\Rr^3=\ker f^2 \oplus \ker(f-2\text{Id}).$
    \item  Trouver une base $\mathcal{B}'$ de $\Rr^3$  telle que
$${\text{mat}(f,\mathcal{B}')=\left(\begin{smallmatrix}
0&1&0\\0&0&0\\0&0&2\end{smallmatrix}\right).}$$
    \item  Soit $g\in\mathcal{L}(\Rr^3)$ tel que $g^2=f$. Montrer que
    $\ker f^2$ est stable par $g$. En d\'eduire qu'un tel endomorphisme $g$ ne peut exister.
\end{enumerate}
\finenonce{001689}



\finexercice



\exercice{1690, legall, 1998/09/01}

\enonce{001690}{}
Soit
$  \displaystyle{ A=\begin{pmatrix} 1 & 1 & 0    \cr
                                   1/2 & 3/2 & -1/2   \cr
                                    -1/2 & 1/2 & 3/2  \cr  \end{pmatrix} \in M_3(\R )}  $ et
$  f   $ l'endomorphisme lin\'eaire de $  \R ^3  $ ayant pour matrice $  A   $ dans la base
canonique $  \epsilon   $ de $  \R ^3  .$
\begin{enumerate}
    \item Calculer le polyn\^ome caract\'eristique de $  A .$
    \item Trouver une base $  \epsilon '=\{ e_1, e_2 ,e_3 \}   $ de $  \R ^3  $ telle
que
$  \hbox{Mat} (f, \epsilon ')=\begin{pmatrix} 2 & 0 & 0    \cr
                                   0 & 1 & 1   \cr
                                    0 & 0 & 1  \cr  \end{pmatrix}  .$
    \item Soit $  g \in \mathcal{L} (\R ^3)  $ un endomorphisme tel que $  f\circ g=g\circ f  .$ Montrer
que $  \hbox{Ker}(f-2Id)  $ et $  \hbox{Ker}(f-Id)^2  $ sont laiss\'es stables par $  g  .$
En d\'eduire que
la matrice de $  g  $ dans $  \epsilon '  $ est de la forme
$  \hbox{Mat} (g, \epsilon ')=\begin{pmatrix} \lambda  & 0 & 0    \cr
                                   0 & a & b   \cr
                                    0 & c & d  \cr  \end{pmatrix}  $ avec
   $  \begin{pmatrix} a & b \cr c & d \cr \end{pmatrix} \begin{pmatrix} 1 & 1 \cr  0 & 1 \cr \end{pmatrix}=
   \begin{pmatrix} 1 & 1 \cr  0 & 1 \cr \end{pmatrix}\begin{pmatrix} a & b \cr c & d \cr \end{pmatrix}  .$
   Pr\'eciser les valeurs possibles de $  a , b , c  $ et $  d  .$
    \item Soit $  F=\{ B\in  M_3(\R )  ;  AB=BA\}   .$ Montrer que $  F  $ est un sous-espace vectoriel de $   M_3(\R )  .$ Calculer sa dimension (on pourra utiliser la question 3.).
\end{enumerate}
\finenonce{001690}



\finexercice



\exercice{1691, ortiz, 1999/04/01}

\enonce{001691}{}
Les questions sont ind\'ependantes. $K$ d\'esigne
$\Rr$ ou $\Cc$, $E$ est un $K$-espace vectoriel de
dimension finie $n$, $\mathcal{B}=(e_1,...,e_n)$
est une base fix\'ee de $E$ et $f$. un
endomorphisme de $E$.
\begin{enumerate}
    \item  Donner un exemple de matrice de $M_2(K)$ non trigonalisable.
    \item  Donner un exemple de matrice de $M_n(K)$ \`a la fois non
diagonalisable et trigonalisable.
    \item  D\'eterminer sans calculs les valeurs propres complexes de $f$ s
i sa matrice dans $\mathcal{B}$ est
$M=\left(\begin{smallmatrix}
1&0&1\\0&1&0\\1&0&1\end{smallmatrix}\right)$.
    \item  On suppose que $n=3$ et que la matrice de $f$ dans la base $\mathcal{B}$
est $M=\left(\begin{smallmatrix}
3&2&4\\-1&3&-1\\-2&-1&-3\end{smallmatrix}\right)$.
Montrer que le plan d'\'equation $x+2z=0$ est
stable par $f.$
    \item  Que peut-on dire d'un vecteur g\'en\'erateur d'une droite stable par $f$ ?
    \item  Montrer que si l'endomorphisme $f$ est trigonalisable alors il admet
au moins un sous-espace vectoriel stable par $f$
et de dimension $k\in[ 0,n]$ fix\'ee.
\end{enumerate}
\finenonce{001691}



\finexercice



\exercice{3578, quercia, 2010/03/10}
\enonce{003578}{Trigonalisation de matrices}

Soit $A = \begin{pmatrix} -1 &\phantom-2 & 0 \cr 2 &2 &-3 \cr -2 &2 &1 \cr \end{pmatrix}$
et $\varphi$ l'endomorphisme de $\R^3$ canoniquement associé à $A$.

\begin{enumerate}
  \item Vérifier que $A$ n'est pas diagonalisable.
    
  \item Chercher deux vecteurs propres de $A$ linéairement indépendants.
    
  \item Compléter ces vecteurs en une base de $\R^3$.
    
  \item \'Ecrire la matrice de $\varphi$ dans cette base.
    
  \item Résoudre le système différentiel : $X' = AX$.
    
\end{enumerate}
\finenonce{003578}



\finexercice

\exercice{3620, quercia, 2010/03/10}
\enonce{003620}{$AB=0$}

Soient $A,B\in\mathcal{M}_n(\C)$ telles que $AB = 0$. Montrer que $A$ et $B$
sont simultanément trigonalisables.
\finenonce{003620}



\finexercice


\exercice{3621, quercia, 2010/03/10}
\enonce{003621}{Produit de matrices nilpotentes}

Soient $A_1,\dots,A_n\in\mathcal{M}_n(K)$ nilpotentes et commutant deux à deux.
Montrer que $A_1\dots A_n = 0$.
\finenonce{003621}



\finexercice


\exercice{3622, quercia, 2010/03/10}
\enonce{003622}{Matrices nilpotentes}

Soit $A\in\mathcal{M}_n(\C)$. Montrer que $A$ est nilpotente si et seulement
si pour tout $k\in\N^*$ on a $\mathrm{tr}(A^k)=0$.

\finenonce{003622}



\finexercice


\exercice{3623, quercia, 2010/03/10}
\enonce{003623}{Mines MP 2003}
Soit~$E$ un ev de dimension finie et $(u_n)$ une suite d'endomorphismes
diagonalisables convergeant vers~$u\in\mathcal{L}(E)$.

$u$ est-il diagonalisable~?

\finenonce{003623}



\finexercice


\exercice{3624, quercia, 2010/03/10}
\enonce{003624}{Mines-Ponts MP 2005}

On donne une matrice carrée réelle $M$ d'ordre $n$. Soient $\alpha,\beta$ les
multiplicités de zéro dans $\chi_M$ et $\mu_M$. Montrer que $\dim(\mathrm{Ker} M) = \alpha$
si et seulement si $\beta = 1$.


\finenonce{003624}



\finexercice


\exercice{5677, rouget, 2010/10/16}
\enonce{005677}{***}
Soient $f$ et $g$ deux endomorphismes d'un $\Cc$-espace vectoriel de dimension finie non nulle qui commutent. Montrer que $f$ et $g$ sont simultanément trigonalisables.
\finenonce{005677}


\finexercice

\section{201.06 Réduction de Jordan}


\exercice{1692, bodin, 1999/11/01}

\enonce{001692}{}
Soit $E$ un espace vectoriel r\'eel de
dimension $4$.  Soit:
\[
U=
\left(
\begin{array}{cccc}
1 & 0 & 0  & 0 \\
-1 & 4 & 1 & -2 \\
2 & 1 & 2 & -1 \\
1 & 2 & 1 & 0
\end{array}
\right)
\]
la matrice d'un endomorphisme $u$ de $E$ dans la base canonique de
$E$.
\begin{enumerate}
\item Calculer le polyn\^ome caract\'eristique de $u$.
D\'eterminer les sous-espaces propres $E_1$ et $E_2$. Pourquoi $u$
est-il non diagonalisable?  Est-il triangularisable ?
\item  D\'eterminer les sous-espaces caract\'eristiques $F_1$
et $F_2$.  Pour $k=1,2$, donner l'ordre $\beta_k$ du nilpotent
$(u-\lambda_k.  \mathrm{id}_E)\vert_{F_k}$ ($\lambda_1=1$,
$\lambda_2=2$).
\item  Si $v\in F_2$ et $v\notin \ker(u-2.
\mathrm{id}_E)^{\beta_2-1}$, montrer que $f_1=(u-2.
\mathrm{id}_E)^{\beta_2-1}(v)$, $f_1=(u-2.
\mathrm{id}_E)^{\beta_2-2}(v)$, \ldots, $f_{\beta_2}=v$ forment
une base de $F_2$.
\item On note \( f=\left\lbrace
f_1,\ldots,f_4\right\rbrace \) la compl\'et\'ee de la base
pr\'ec\'edente par une base de $F_1$.  V\'erifier que $T=[u]_f^f$
est triangulaire.  D\'ecomposer $T$ sous la forme $D+N$, o\`u $D$
est diagonale, $N$ est nilpotente, et $DN=ND$.  Calculer $T^5$.
\end{enumerate}
\finenonce{001692}



\finexercice



\exercice{1693, bodin, 1999/11/01}

\enonce{001693}{}
 Quel est le polyn\^ome caract\'eristique d'un endomorphisme nilpotent d'un
$\Cc$-espace vectoriel de dimension finie ?
\finenonce{001693}



\finexercice



\exercice{1694, bodin, 1999/11/01}

\enonce{001694}{}
Donner toutes les r\'eduites de Jordan de $\mathcal{M}_n(\Cc)$ des endomorphismes nilpotents pour
$1\le n \le 4$.
\finenonce{001694}



\finexercice



\exercice{1695, bodin, 1999/11/01}

\enonce{001695}{}
 Soit $ \rho  $ l'application de $ \R_4[X] $ dans
lui-m\^eme qui \`a un polyn\^ome $ P  $ associe le reste de la
division euclidienne de $ P $ par $ (X^2-1) .$
\begin{enumerate}
\item Montrer que $ \rho  $ est lin\'eaire.
\item  Montrer que $ \rho ^2= \rho .$ En d\'eduire que $
\rho  $ est diagonalisable.
\item  D\'eterminer (de pr\'ef\'erence
sans calcul) une base de vecteurs propres pour $ \rho  .$
\end{enumerate}
\finenonce{001695}



\finexercice



\exercice{1696, bodin, 1999/11/01}

\enonce{001696}{}
 Les matrices $ \begin{pmatrix}0 & 0 & 1 & 0  \cr
                                   0 & 0 & 0 & 1  \cr
                                    0 & 0 & 0 & 0  \cr
                                    0 & 0 & 0 & 0 \cr \end{pmatrix} $
et $ \begin{pmatrix}0 & 1 & 0 & 0  \cr 0 & 0 & 1 & 0  \cr
                                    0 & 0 & 0 & 1  \cr
                                    0 & 0 & 0 & 0 \cr \end{pmatrix} \in M_4(\C )  $
ont-elles une racine carr\'ee~?
\finenonce{001696}



\finexercice



\exercice{1697, bodin, 1999/11/01}

\enonce{001697}{}
R\'eduire sous la forme de Jordan les matrices suivantes :
$$
\begin{pmatrix}  -1 & 1 & 0 \cr
                 1 & 1 & 2 \cr
                 1 & -1 & 0\cr
\end{pmatrix}, \qquad
\begin{pmatrix} 4 & 0 & 0 &0 \cr
                0 & 0 & 1 & 0 \cr
                0 & 1 & 2 & 2 \cr
                0& 1 & -1 & 1\cr
\end{pmatrix},\qquad
\begin{pmatrix} 3 & -1 & 1 &-7 \cr
                9 & -3 & -7 & -1 \cr
                0 & 0 & 4 & -8 \cr
                0& 0 & 2 & -4\cr
\end{pmatrix}.
$$

\finenonce{001697}



\finexercice



\exercice{1698, bodin, 1999/11/01}

\enonce{001698}{}
Soit $E$ un $\Cc$-espace vectoriel de dimension finie $n$. Soit $f\in\mathcal{L}(E)$ un
endomorphisme nilpotent d'indice $N$ (le plus petit entier $p$ tel que $f^p = 0$).
Montrer que
$$N=n \Leftrightarrow \text{rang} f = n-1.$$
\finenonce{001698}



\finexercice



\exercice{2589, delaunay, 2009/05/19}
\enonce{002589}{}
On consid\`ere la matrice suivante
$$A=\begin{pmatrix}1&-1&0 \\  1&0&-1 \\  -1&0&2\end{pmatrix}$$
et $f$ l'endomorphisme de $\R^3$ associ\'e.
\begin{enumerate}
 \item Factoriser le polyn\^ome caract\'eristique de $A$.
 \item D\'eterminer les sous-espaces propres et caract\'eristiques de $A$.
 \item D\'emontrer qu'il existe une base de $\R^3$ dans laquelle la matrice de $f$ s'\'ecrit
$$B=\begin{pmatrix}1&1&0 \\  0&1&1 \\  0&0&1\end{pmatrix}.$$
 \item  Ecrire la d\'ecomposition de Dunford de $B$ (justifier).
\end{enumerate}
\finenonce{002589} 


\finexercice

\exercice{2602, delaunay, 2009/05/19}
\enonce{002602}{}
 Soit $A$ la matrice 
$$A=\begin{pmatrix}3&2&4 \\  -1&3&-1 \\  -2&-1&-3\end{pmatrix}$$ 
et $f$ l'endomorphisme de $\R^3$ associ\'e.
\begin{enumerate}
 \item  Factoriser le polyn\^ome caract\'eristique de $A$.
 \item D\'eterminer les sous-espaces propres et caract\'eristiques de $A$. 
 \item  D\'emontrer qu'il existe une base de $\R^3$ dans laquelle la matrice de $f$ est 
$$B=\begin{pmatrix}-1 &0& 0 \\  0&2&1 \\ 0&0&2\end{pmatrix}$$
et trouver une matrice $P$ inversible telle que $A=PBP^{-1}$.

 \item Ecrire la d\'ecomposition de Dunford de $B$ (justifier).

 \item Calculer $\exp B$.
\end{enumerate}
\finenonce{002602} 


\finexercice

\exercice{3581, quercia, 2010/03/10}
\enonce{003581}{Décomposition de Dunford}

Soit $A\in\mathcal{M}_n(\C)$. Montrer qu'il existe deux matrices $D,N$ telles
que $A = D+N$, $D$ est diagonalisable, $N$ est nilpotente, $DN=ND$.
\finenonce{003581}



\finexercice


\exercice{3582, quercia, 2010/03/10}
\enonce{003582}{$M$ et ${}^tM$ sont semblables}

Montrer qu'une matrice compagnon est semblable à sa transposée.
En déduire que pour toute $M\in\mathcal{M}_n(K)$ les matrices  $M$ et ${}^tM$ sont semblables.
\finenonce{003582}



\finexercice


\exercice{3583, quercia, 2010/03/10}
\enonce{003583}{Réduction de Jordan (Mines MP 2003)}
Soit $f\in\mathcal{L}(\R^3)$ telle que $\mathrm{Spec}(f) = \{\lambda\}$ et $\dim(\mathrm{Ker}(f-\lambda\mathrm{id}))=2$.

Montrer qu'il existe une base~$\cal B$ dans laquelle $\mathrm{Mat}(f) = \begin{pmatrix}\lambda&0&0\cr0&\lambda&1\cr0&0&\lambda\cr\end{pmatrix}$.

\finenonce{003583}



\finexercice

\exercice{5670, rouget, 2010/10/16}
\enonce{005670}{*** Décomposition de \textsc{Dunford}}
\label{ex:rou20}
Soit $E$ un $\Kk$-espace vectoriel de dimension finie non nulle et $f$ un endomorphisme de $E$ dont le polynôme caractéristique est scindé sur $\Kk$.

Montrer qu'il existe un couple d'endomorphismes $(d,n)$ et un seul tel que $d$ est diagonalisable,  $n$ est nilpotent $n$ et $f = d+n$.

\finenonce{005670}


\finexercice

\section{201.07 Applications}


\exercice{1703, barraud, 2003/09/01}

\enonce{001703}{}
  Soit $A\in\mathcal{M}_{3}(\R)$ la matrice 
  $$
  A=
  \begin{pmatrix}
    0&1&1\\
    1&0&1\\
    0&0&1
  \end{pmatrix}
  $$
  Donner un polynôme annulateur de $A$ de degré aussi petit que
  possible. En déduire $A^{-1}$, $A^{3}$, et $A^{5}$.
\finenonce{001703}



\finexercice



\exercice{1704, barraud, 2003/09/01}

\enonce{001704}{}
  Résoudre les systèmes différentiels suivants
  $$
  \left\{
  \begin{array}{ccccccc}
    \frac{dx}{dt} & = &  4x &+& 6y \\[1.5ex]
    \frac{dy}{dt} & = & -3x &-& 5y \\[1.5ex]
    \frac{dz}{dt} & = & -3x &-& 6y &-& 5z
  \end{array}
  \right.
  \qquad
  \left\{
  \begin{array}{ccccccc}
    \frac{dx}{dt} & = &  2x &+&  y &+&  z  \\[1.5ex]
    \frac{dy}{dt} & = &  3x &+& 3y &+& 4z  \\[1.5ex]
    \frac{dz}{dt} & = & -3x &-&  y &-& 2z
  \end{array}
  \right.
  $$
\finenonce{001704}



\finexercice



\exercice{1705, barraud, 2003/09/01}

\enonce{001705}{}
  Déterminer toutes les suites $(u_{n})$ telles que :
  $$
  \begin{cases}
    \forall n\in \N \quad u_{n+3}+u_{n+2}+u_{n+1}+u_{n}=0\\
    u_{0}=1, u_{1}=2, u_{2}=0
  \end{cases}
  $$

  Résoudre l'équation différentielle :
  $$
  \begin{cases}
    f'''+f''+f'+f=0\\
    f(0)=1, f'(0)=0, f''(0)=0
  \end{cases}
  $$
\finenonce{001705}



\finexercice



\exercice{1706, barraud, 2003/09/01}

\enonce{001706}{}
Résoudre le système différentiel suivant :
$$
\left\{
\begin{array}{rcrcrcr}
\frac{dx}{dt} &=& 2x(t) &+& 2y(t) &+& 2z(t) \\
\frac{dy}{dt} &=&  x(t) &+& 3y(t) &+& 2z(t) \\
\frac{dz}{dt} &=& -x(t) &-&  y(t) &-&  z(t)
\end{array}
\right.
$$


Donner toutes les solutions qui satisfont $x(0)=1$, $y(0)=2$, $z(0)=-1$.
\finenonce{001706}



\finexercice



\exercice{1707, barraud, 2003/09/01}

\enonce{001707}{}
Réduire la matrice 
$$
A=
\begin{pmatrix}
  0& 1& 1\\
  1& 1& 0\\
  1&-3& 4
\end{pmatrix}
$$
\textit{(c'est à dire étudier la diagonalisabilité ou la
  triangularisabilité de $A$, et donner une matrice $P$ telle que
  $P^{-1}AP$ soit aussi simple que possible)}

\medskip
\textbf{\it Application~:} Déterminer toutes les fonctions dérivables
$x,y,z$ de $\R$ dans $\R$ satisfaisant les conditions~:
$$
\left\{
  \begin{aligned}
    x'&=y+z\\
    y'&=x+y\\
    z'&=x-3y+4z
  \end{aligned}
\right.\qquad \text{ et }\qquad
\left\{
  \begin{aligned}
    x(0)&=1\\
    y(0)&=0\\
    z(0)&=0
  \end{aligned}
\right.
$$
\textit{(on rappelle qu'il n'est pas utile de calculer $P^{-1}$... )}
\finenonce{001707}



\finexercice



\exercice{1708, barraud, 2003/09/01}

\enonce{001708}{}
  Déterminer toutes les suites $(u_{n})_{n\in\N}$ à valeur complexes
  telles que~:
  $$
  \forall n\in\N, u_{n+3}+2 u_{n+2}+2u_{n+1}+u_{n}=0.
  $$
  Montrer que les suites réelles satisfaisant cette relation sont les
  suites de la forme~:
  $$
  u_{n}=A(-1)^{n}+B\cos(\frac{2n\pi}{3}+\phi)
  $$
  où $A,B$ et $\phi$ sont des réels.
\finenonce{001708}



\finexercice



\exercice{1709, barraud, 2003/09/01}

\enonce{001709}{}

Etant donnés quatre nombres réels $(u_{0},v_{0},w_{0},x_{0})$, on définit quatre nouveaux nombres $(u_{1},v_{1},w_{1},x_{1})$ en calculant les
moyennes suivantes :
 $u_{1}= \frac{2u_{0}+ v_{0}+ w_{0}+ x_{0}}{5}$,
 $v_{1}= \frac{ u_{0}+2v_{0}+ w_{0}+ x_{0}}{5}$,
 $w_{1}= \frac{ u_{0}+ v_{0}+2w_{0}+ x_{0}}{5}$, et
 $x_{1}= \frac{ u_{0}+ v_{0}+ w_{0}+2x_{0}}{5}$.
En itérant ce procédé, on définit quatre suites $(u_{n})$, $(v_{n})$, $(w_{n})$, et $(x_{n})$ telles que pour tout $n\in\N$ on ait :
$$
\left\{
\begin{array}{r@{\,}l} 
  u_{n+1}&= \frac{1}{5} (2u_{n}+ v_{n}+ w_{n}+ x_{n}) \\[1.5ex]
  v_{n+1}&= \frac{1}{5} ( u_{n}+2v_{n}+ w_{n}+ x_{n}) \\[1.5ex]
  w_{n+1}&= \frac{1}{5} ( u_{n}+ v_{n}+2w_{n}+ x_{n}) \\[1.5ex]
  x_{n+1}&= \frac{1}{5} ( u_{n}+ v_{n}+ w_{n}+2x_{n})
\end{array}
\right.
$$ 

\begin{enumerate}
\item
Ecrire la matrice $A$ associée à cette relation de récurrence, et la
matrice $B=5A$. Que dire de la diagonalisabilité de $B$ ?

\item
Sans calculer le polynôme caractéristique de $B$, montrer que $1$ est
valeur propre de $B$. Quelle est la dimension de l'espace propre
associé ? Que dire de la multiplicité de $1$ comme valeur propre de $B$ ?

\item
En utilisant la trace de $B$, déterminer toutes les valeurs propres de $B$.

\item
Donner un polynôme annulateur de $B$ de degré 2.

\item
En déduire l'existence de deux réels $a_{n}$ et $b_{n}$, que l'on
calculera, tels que $B^{n}=a_{n}B+b_{n}I$.

\item
Calculer $\displaystyle \lim_{n\rightarrow\infty}{\frac{a_{n}}{5^{n}}}$ et
$\displaystyle \lim_{n\rightarrow\infty}{\frac{b_{n}}{5^{n}}}$. En déduire que la
suite de matrices $(A^{n})_{n\in\N}$ est convergente et donner sa
limite.

{\it (On rappelle qu'une suite de matrices $M_{n}$ est dite
convergente si chaque suite de coefficient est convergente. On pourra
utiliser sans démonstration la continuité des opérations élémentaires
sur les matrices pour cette notion de limite, c'est à dire que :
\\
\ - si $(\lambda_{n})$ est une suite convergente alors pour toute matrice
$M$, la suite $(\lambda_{n}M)$ est convergente et
$\lim\limits_{n\rightarrow\infty}(\lambda_{n}M) =
(\lim\limits_{n\rightarrow\infty}\lambda_{n})M$  
\\
\ - si $(M_{n})$ est une suite de matrices convergente alors pour tout
vecteur $X$, la suite de vecteurs $(M_{n}X)$ est convergente et
 $\lim\limits_{n\rightarrow\infty}(M_{n}X) = 
 (\lim\limits_{n\rightarrow\infty} M_{n})X$.)}

\item
En déduire que les suites $(u_{n})_{n\in\N}$, $(v_{n})_{n\in\N}$,
$(w_{n})_{n\in\N}$, et $(x_{n})_{n\in\N}$ sont convergentes, et
donner leur limite.

\end{enumerate}
\finenonce{001709}



\finexercice



\exercice{1710, barraud, 2003/09/01}

\enonce{001710}{}
Donner toutes les suites $(x_{n})$, $(y_{n})$ et $(z_{n})$ telles que~:
{\it(on notera $\omega=e^{\frac{i\pi}{3}}$)}
$$
\forall n\in\N,\ 
\left\{
\begin{array}{r}
  x_{n+1}= x_{n}+y_{n}\\
  y_{n+1}= y_{n}+z_{n}\\
  z_{n+1}= z_{n}+x_{n}
\end{array}
\right.
$$

Parmi les solutions de ce système, donner celle qui satisfait $x_{0}=2$
et $y_{0}=z_{0}=1$.

\finenonce{001710}



\finexercice


\exercice{1711, barraud, 2003/09/01}

\enonce{001711}{}
  Soit $a$ un réel. On considère le système à $n$ équations et $n$
  inconnues suivant~:
  $$
  \left\{
    \begin{array}{lccccc}
               a\,x_{1}  - x_{2}  = 0 \\
      -x_{p-1}+a\,x_{p} -x_{p+1}  = 0 &  (2\leq p\leq n-1)\\
      -x_{n-1}+a\,x_{n}           = 0
    \end{array}
  \right.
  $$
  \'Ecrire la matrice $A_{n}$ associée à ce système. On note $D_{n}=\det
  A_{n}$. Calculer $D_{n}$ en fonction de $D_{n-1}$ et $D_{n-2}$
\finenonce{001711}



\finexercice



\exercice{1712, barraud, 2003/09/01}

\enonce{001712}{}
  On considère la matrice $A=
  \begin{pmatrix}
    a&-b&-c&-d\\
    b& a& d&-c\\
    c&-d& a& b\\
    d& c&-b& a  
  \end{pmatrix}
$, avec $(b,c,d)\neq(0,0,0)$.
\begin{enumerate}
\item 
  Calculer $A{}^t{A}$.
  Que vaut $\det A$ au signe près~?

\item
  En étudiant le signe du terme en $a^{4}$ dans le déterminant de $A$,
  montrer que $\det A=(a^{2}+b^{2}+c^{2}+d^{2})^{2}$. Sans calcul
  supplémentaire, en déduire que le polynôme caractéristique de $A$ est
  $\chi_{A}=((a-X)^{2}+b^{2}+c^{2}+d^{2})^{2}$.

\item
  $A$ est-elle diagonalisable sur $\R$~? (justifier)

\item
  On se place maintenant dans le cas où $a=1$, $b=c=d=-1$. Vérifier que
  $(i\sqrt{3},1,1,1)$ et $(-1,i\sqrt{3},-1,1)$ sont des vecteurs propres
  de $A$, puis diagonaliser $A$ sur $\C$.

\item
  Application~: résoudre le système récurent suivant (il n'est pas
  nécessaire de calculer l'inverse de la matrice de passage de la
  question précédente). On notera $\omega=1/2+i\sqrt{3}/2=e^{i\pi/3}$.
$$
\left\{
  \begin{array}{ccr}
    u_{n+1} &=& u_{n}+v_{n}+w_{n}+h_{n}\\
    v_{n+1} &=&-u_{n}+v_{n}-w_{n}+h_{n}\\
    w_{n+1} &=&-u_{n}+v_{n}+w_{n}-h_{n}\\
    h_{n+1} &=&-u_{n}-v_{n}+w_{n}+h_{n}
  \end{array}
\right.
\qquad
\left\{
  \begin{array}{ccr}
    u_{0} &=&1\\
    v_{0} &=&0\\
    w_{0} &=&0\\
    h_{0} &=&0
  \end{array}
\right.
$$
\end{enumerate}

\finenonce{001712}




\finexercice



\exercice{1713, ortiz, 1999/04/01}

\enonce{001713}{}
R\'esoudre le syst\`eme diff\'erentiel $X'=AX$
o\`u $A$ est la matrice :
$$A=\left(
\begin{matrix}
3&2&4\\
-1&3&-1\\
-2&-1&-3
\end{matrix}\right) \in \mathcal{M}_3(\Rr)  .$$
\finenonce{001713}



\finexercice



\exercice{1714, ortiz, 1999/04/01}

\enonce{001714}{}
Soit la matrice
$A=\left(
\begin{smallmatrix}
3&2&4\\
-1&3&-1\\
-2&-1&-3
\end{smallmatrix}\right) \in \mathcal{M}_3(\Rr).$

Par diff\'erentes m\'ethodes, calculer $A^n$, pour
$n\in\Nn$. Montrer que la formule obtenue a un
sens pour $n\in\Zz$ et donner plusieurs m\'ethodes
pour \'etablir sa validit\'e dans ce cas.
\finenonce{001714}



\finexercice



\exercice{1715, ortiz, 1999/04/01}

\enonce{001715}{}
Soit l'endomorphisme $f\in\mathcal{L}(\Rr^3)$ dont
la matrice dans la base canonique de $\Rr^3$ est :
$${ M=\left(
\begin{smallmatrix}
-2&1&1\\
8&1&-5\\
4&3&-3
\end{smallmatrix}\right).
}$$
\begin{enumerate}
    \item  D\'eterminer toutes les droites vectorielles de $\Rr^3$ stables par $f$.
    \item  D\'eterminer toutes les plans vectoriels $P$ de $\Rr^3$
stables par $f$ (on commencera par \'etudier le
polyn\^ome caract\'eristique de la restriction de $f$ \`a $P$).
    \item  Donner la liste de tous les sous-espaces vectoriels de $\Rr^3$ stables par $f$.
\end{enumerate}
\finenonce{001715}



\finexercice



\exercice{1716, bodin, 1999/11/01}

\enonce{001716}{}
Calculer les puissances et l'exponentielle ($e^M=\sum_{k=0}^{+\infty} \frac{M^k}{k!}$) des
 matrices suivantes :
$$
B=
\begin{pmatrix}
4&1&0\\
0&4&1\\
0&0&4
\end{pmatrix},\qquad
A=
\begin{pmatrix}
3&2&4\\
-1&3&-1\\
-2&-1&-3
\end{pmatrix}.
$$
\finenonce{001716}



\finexercice



\exercice{1717, bodin, 1999/11/01}

\enonce{001717}{}

Soit $E$ un espace vectoriel r\'eel de
dimension finie $n$.  Soit $f\in\mathcal{L}(E)$
diagonalisable.  Donner une condition n\'ecessaire
et suffisante pour qu'il existe
$g\in\mathcal{L}(E)$ tel que $g^2=f$.  Dans le cas
d'existence de $g$, donner le nombre exact de $g$
tel que $g^2=f$.

\medskip

\noindent
{\emph{Application}} Soit :
\[
M=
\left(
\begin{array}{ccc}
5 & 1 & -1 \\
2 & 4 & -2 \\
1 & -1 & 3
\end{array}
\right).
\]
Montrer qu'il existe $N\in\mathcal{M}_3(\R)$ telle
que $N^2=M$.  D\'eterminer une $N$.
\finenonce{001717}



\finexercice



\exercice{1718, bodin, 1999/11/01}

\enonce{001718}{}
  Soit $M\in\mathcal{M}_n(\Cc)$. Montrer que $M$ et ${}^tM$ sont semblables.\\
\emph{Indication} : le montrer d'abord pour des blocs de Jordan
n'ayant que des $1$ au-dessus de la diagonale.
\finenonce{001718}



\finexercice



\exercice{1719, bodin, 1999/11/01}

\enonce{001719}{}
 Soit $M\in\mathcal{M}_n(\Cc)$. Donner une condition n\'ecessaire et suffisante sur
$M$ pour que $M$ et $2M$ soient semblables.
\finenonce{001719}



\finexercice



\exercice{1720, ortiz, 1999/04/01}

\enonce{001720}{}
Soit $a\in\mathcal{L}(E)$ un endomorphisme d'un
$K$-espace vectoriel de dimension $n$ ayant $n$
valeurs propres distinctes. On pose
$$
\mathcal{C}=\left\{u\in\mathcal{L}(E):au=ua\right\}.
$$
\begin{enumerate}
    \item  Soit $u\in\mathcal{C}$.
    \begin{enumerate}
        \item Montrer que tout sous-espace vectoriel propre de $a$ est stable par $u$.
        \item En d\'eduire que $u$ est diagonalisable.
    \end{enumerate}
    \item
    \begin{enumerate}
         \item Montrer  que $\mathcal{C}$ est un sous-espace
vectoriel de $\mathcal{L}(E)$ et que
$\dim\mathcal{C}=n$.
         \item Montrer que la famille $(\text{Id}_E,a,...,a^{n-1})$
est une famille libre de $\mathcal{L}(E)$
(raisonner par l'absurde et utiliser le polyn\^ome
minimal de $a$.)
         \item En d\'eduire que $\mathcal{C}=\left\{P(u):P\in K[X]\right\}$.
\end{enumerate}
\end{enumerate}
\finenonce{001720}



\finexercice



\exercice{1721, ortiz, 1999/04/01}

\enonce{001721}{}
Soient $f\in\mathcal{L}(E)$ un endomorphisme et
$a\in E$ tels que la famille
$(a,f(a),...,f^{n-1}(a))$ est une base de $E$.
\begin{enumerate}
    \item  Soit $P\in K[X]\setminus\left\{0\right\}$ un polyn\^ome annulateur de $f$. Montrer que $\deg(P)\ge n$ (raisonner par l'absurde).
    \item  En d\'eduire que le polyn\^ome minimal de $f$ est
(au signe pr\`es) le polyn\^ome caract\'eristique de $f$.
\end{enumerate}
\finenonce{001721}



\finexercice



\exercice{1722, ortiz, 1999/04/01}

\enonce{001722}{}
Donner un exemple de deux matrices de
$\mathcal{M}_4(\Rr)$ ayant m\^eme polyn\^ome caract\'eristique
 et m\^eme polyn\^ome minimal et pourtant non semblables.
Qu'en est-il pour deux
 matrices de $\mathcal{M}_2(\Rr)$ ?
\finenonce{001722}



\finexercice



\exercice{1723, ortiz, 1999/04/01}

\enonce{001723}{}
Soit le $\Rr$-espace vectoriel
$$\mathcal{S}=\left\{
(u_n)_{n\in\Nn}\in
\Rr^{\Nn}:\forall n\ge 3,
u_n=3u_{n-1}-3u_{n-2}+u_{n-3} \right\}.$$
\begin{enumerate}
    \item  Montrer que l'application
$${f: \mathcal{S}\to \Rr^3,
u=(u_n)_{n\in\Nn}\mapsto(u_0,u_1,u_{2})}$$ est un
isomorphisme de $\Rr$-espace vectoriels.
    \item  Soient la matrice
$
A=
\left(\begin{smallmatrix}
0&1&0\\
0&0&1\\
1&-3&3
\end{smallmatrix}\right)\in \mathcal{M}_3(\Rr)
,$ $\sigma\in\mathcal{L}(\Rr^3)$ l'endomorphisme
canoniquement associ\'e \`a $A$ et,  pour $n\ge
2,$ $U_n=(u_{n-2},u_{n-1},u_{n})\in \Rr^3$.
Montrer que $\sigma(U_{n-1})=U_n$ et en d\'eduire
une base de $\mathcal{S}$.
\end{enumerate}
\finenonce{001723}



\finexercice



\exercice{1724, legall, 1998/09/01}

\enonce{001724}{}
Soient $  (x_n)_{n \in \N }  , $ $  (y_n)_{n \in \N }  $ et
$  (z_n)_{n \in \N }  $ trois suites de nombres r\'eels satisfaisant aux relations de r\'ecurrence~:
$$\left\{ \begin{matrix} x_{n+1} &= & y_n & - & x_n & + & z_n  \cr
y_{n+1} &= & x_n & - & y_n & + & z_n \cr
z_{n+1} &= & x_n & + & y_n & - & z_n \cr \end{matrix} \right.$$
Calculer les valeurs de $  x_n  , $ $  y_n  $ et $  z_n  $ en fonction de $  x_0  , $ $
y_0  $ et $  z_0  .$
\finenonce{001724}



\finexercice



\exercice{1725, legall, 1998/09/01}

\enonce{001725}{}
Soit $  E  $ un $  \mathbb{K} $-espace vectoriel de dimension finie et $  f \in \mathcal{L} (E)  $ telle
que $  f^2=f  .$ Pour quelles valeurs de $  t\in \R   $ l'endomorphisme $  f_t=id+tf  $ est
inversible ? Calculer $  f_t^{-1}  .$
\finenonce{001725}



\finexercice



\exercice{1726, legall, 1998/09/01}

\enonce{001726}{}
Etudier les solutions (suivant $  A  $) dans  $  M_2 (\C )  $
de l'\'equation $  X^2=A  .$
\finenonce{001726}



\finexercice



\exercice{1727, legall, 1998/09/01}

\enonce{001727}{}
Soit $  A\in M_n(\mathbb{K} )  .$ On note $  C(A)=\{ B \in M_n(\mathbb{K} )  ;
AB=BA\}   .$
\begin{enumerate}
    \item On suppose que $  A  $ a des valeurs propres simples. Montrer que les propri\'et\'es suivantes sont
\'equivalentes~:
\vskip1mm
\hskip1mm i) $  B\in C(A)  .$
\vskip1mm
\hskip1mm ii) $  B  $ a une base de vecteurs propres en commun avec $  A  .$
\vskip1mm
\hskip1mm iii) Il existe $  P\in \mathbb{K} _{n-1}[X]  $ tel que $  B=P(A)  .$
\vskip1mm
\hskip1mm iv) Il existe $  P\in \mathbb{K} [X]  $ tel que $  B=P(A)  .$
\vskip1mm
    \item On suppose que $  n =3  $ (pour simplifier) et que $  A  $ est diagonalisable
avec une valeur propre double. D\'eterminer $  C(A)  .$
\end{enumerate}
\finenonce{001727}



\finexercice



\exercice{1728, legall, 1998/09/01}

\enonce{001728}{}

{{\em Les parties I, II, III et IV peuvent \^etre trait\'ees
ind\'ependamment les unes des autres.}}


 Soient
$  M_a=\begin{pmatrix} a+1 & 1-a & a-1    \cr
                                   -1 & 3 & 2a-3   \cr
                                    a-2 & 2-a & 3a-2  \cr  \end{pmatrix}\in M_3(\R )  $ une matrice d\'ependant
d'un param\`etre r\'eel $  a  $ et  $  f_a   $ l'endomorphisme lin\'eaire de $  \R
^3  $ ayant pour matrice $  M_a   $ dans la base canonique de $  \R ^3  .$
\vskip1mm
On nomme {\em racine carr\'ee} d'une matrice $  M\in M_n(\R )  $ toute matrice
$  N \in M_n(\R )  $ telle que $  N^2=M  .$
\vskip1mm On d\'esigne par $  I  $ la matrice identit\'e et, pour toute base $  \epsilon   $ de
$  \R ^3  ,$ on
note $  \hbox{Mat}(f_a,\epsilon)  $ la matrice repr\'esentant
l'endomorphisme $  f_a  $ dans la base $  \epsilon  .$
 \vskip2mm
 \centerline{{\bf I}}
\begin{enumerate}
    \item Calculer les valeurs propres de $  M_a  $ en fonction de $  a  .$ Pour quelle raison la matrice $  M_a  $ est-elle triangularisable~?
    \item Pour quelles valeurs du
param\`etre $  a  $ la matrice $  M_a  $ est-elle diagonalisable~?

\vskip2mm
 \centerline{{\bf II}}
\centerline{{\em On pose maintenant (questions 3 et 4) $  a=2  .$}}
    \item  Diagonaliser $  M_2  .$ D\'eterminer une racine carr\'ee $  A $ de
$  M_2  .$
    \item
    \begin{enumerate}
        \item Soit $  g \in \mathcal{L} (\R ^3)  $ telle que $  g^2=f_2  .$ Montrer que
$  g   $ est diagonalisable (on pourra d\'eterminer le polyn\^ome minimal de $  f_2  $).
 Montrer que les sous-espaces propres de $  f_2  $ sont laiss\'es stables par
$  g  .$
        \item D\'emontrer que la matrice $ \begin{pmatrix} 4 & 0     \cr
                                   0 & 4    \cr \end{pmatrix}  $ a une infinit\'e de racines carr\'ees.
En d\'eduire l'existence d'une infinit\'e de racines carr\'ees de $  M_2  .$
    \end{enumerate}

 \vskip2mm
 \centerline{{\bf III}}
    \item On pose $  a=1  .$ Montrer que $  M_1=2I+N  $ avec $  N  $ nilpotente (telle que $
N^2=0  $). En d\'eduire la valeur de $  (M_1)^n  $, pour tout
$  n \in \N   .$ D\'eterminer deux r\'eels $  \alpha   $ et $  \beta   $
tels que $  \alpha I +\beta N  $ soit une racine carr\'ee de $
M_1  .$

\vskip2mm
 \centerline{{\bf IV}}
\centerline{{\em On pose d\'esormais (questions 6 et 7) $  a=0  .$}}
    \item Montrer que $  \R ^3=\hbox{Ker}(f_0^2)\oplus \hbox{Ker}(f_0-2I)  .$
D\'eterminer une base $  \epsilon  $ de $  \R ^3  $ telle que l'on ait~:
$\hbox{Mat}(f_0, \epsilon )=\begin{pmatrix}  0 &
1 & 0 \cr   0 & 0 & 0   \cr  0 & 0 & 2 \cr  \end{pmatrix}  .$
    \item Soit $  g \in \mathcal{L} (\R ^3)  $ un endomorphisme tel que $  g^2=f_0  .$ Montrer que
$ \hbox{Ker}(f_0^2)  $ est laiss\'e stable par $  g  .$ En d\'eduire que $  f_0  $ n'a pas de
racine carr\'ee.

%8) D\'eterminer les sous-espaces de $  \R ^3  $ stables par $  f_0  .$
\end{enumerate}
\finenonce{001728}



\finexercice



\exercice{2590, delaunay, 2009/05/19}
\enonce{002590}{}
 La suite de Fibonacci $0,1,1,2,3,5,8,13,...$ est la suite $(F_n)_{n\geq 0}$ d\'efinie par la relation de r\'ecurrence $F_{n+1}=F_{n}+F_{n-1}$ pour $n\geq1$, avec $F_0=0$ et $F_1=1$.
\begin{enumerate}
 \item D\'eterminer une matrice $A\in M_2(\R)$ telle que, pour tout $n\geq1$,
$$\begin{pmatrix}F_{n+1} \\  F_n\end{pmatrix}=A^n\begin{pmatrix}F_{1} \\  F_0\end{pmatrix}.$$

 \item Montrer que $A$ admet deux valeurs propres r\'eelles distinctes que l'on note $\lambda_1$ et $\lambda_2$ avec $\lambda_1<\lambda_2$.
 
 \item Trouver des vecteurs propres $\varepsilon_1$ et $\varepsilon_2$ associ\'es aux valeurs propres $\lambda_1$ et $\lambda_2$, sous la forme 
$\displaystyle\begin{pmatrix}\alpha \\  1\end{pmatrix}$, avec $\alpha\in\R$.

 \item D\'eterminer les coordonn\'ees du vecteur $\displaystyle\begin{pmatrix}F_1 \\  F_0\end{pmatrix}$ dans la base 
$(\varepsilon_1,\varepsilon_2)$, on les note $x_1$ et $x_2$.

 \item Montrer que $\displaystyle\begin{pmatrix}F_{n+1} \\  F_n\end{pmatrix}
=\lambda_1^nx_1\varepsilon_1+\lambda_2^nx_2\varepsilon_2$. En d\'eduire que 
$$F_n={\frac{\lambda_1^n}{\lambda_1-\lambda_2}}-{\frac{\lambda_2^n}{\lambda_1-\lambda_2}}\ .$$

 \item Donner un \'equivalent de $F_n$ lorsque $n$ tend vers $+\infty$.
\end{enumerate}
\finenonce{002590} 


\finexercice


\exercice{2597, delaunay, 2009/05/19}
\enonce{002597}{}
 Soit $A$ la matrice 
$$A=\begin{pmatrix}1&-1&0 \\  1&0&-1 \\  -1&0&2\end{pmatrix}$$ 
et $f$ l'endomorphisme de $\R^3$ associ\'e.
\begin{enumerate}
 \item Factoriser le polyn\^ome caract\'eristique de $A$. 
 \item D\'eterminer les sous-espaces propres et caract\'eristiques de $A$. 
 \item D\'emontrer qu'il existe une base de $\R^3$ dans laquelle la matrice de $f$ est 
$$B=\begin{pmatrix}1 & 1 & 0 \\  0&1&1 \\ 0&0&1\end{pmatrix}$$
et trouver une matrice $P$ inversible telle que $A=PBP^{-1}$. 
 \item Ecrire la d\'ecomposition de Dunford de $B$ (justifier). 
 \item  Pour $t\in\R$, calculer $\exp tB$.
 \item  Donner les solutions des syst\`emes diff\'erentiels $Y'=BY$ et $X'=AX$. 
\end{enumerate}
\finenonce{002597} 


\finexercice


\exercice{2598, delaunay, 2009/05/19}
\enonce{002598}{}
\begin{enumerate}
 \item  On note $(\vec e_1, \vec e_2, \vec e_3)$ la base canonique de $\R^3$. Soit $A$ la matrice 
$$A=\begin{pmatrix}1&0&0 \\  0&2&0 \\  0&0&3\end{pmatrix}.$$
Donner sans calcul les valeurs propres de $A$ et une base de vecteurs propres. 
 \item On cherche \`a d\'eterminer, s'il en existe, les matrices $B$ telles que $\exp B=A$.
   \begin{enumerate}
   \item  Montrer que si $A=\exp B$, alors $AB=BA$. 
   \item En d\'eduire que la base $(\vec e_1,\vec e_2, \vec e_3)$ est une base de vecteurs propres de  B.
   \item D\'eterminer toutes les matrices  $B\in M_3(\R)$ telles que $\exp B=A$. Justifier. 
  \end{enumerate}
 \item Soit la matrice $C$, 
$$C=\begin{pmatrix}0&1&0 \\  0&0&1 \\  0&0&0\end{pmatrix}.$$
Montrer qu'il n'existe pas de matrice $D\in M_3(\R)$ telle que $C=\exp D$.
 \item  Calculer le polyn\^ome caract\'eristique et le polyn\^ome minimal de $C$.
 \item Supposons qu'il existe une matrice $E\in M_3(\R)$ telle que $E^2=C$. Notons $Q_E(X)$ son polyn\^ome minimal et $Q_C(X)$ le polyn\^ome minimal de $C$.
    \begin{enumerate}
   \item  Montrer que $Q_E(X)$ divise $Q_C(X^2)$. 
   \item En d\'eduire que $E^3=0$ et que $C^2=0$. 
   \item D\'eduire de ce qui pr\'ec\`ede qu'il n'existe pas de matrice $E$ telle que $E^2=C$. 
  \end{enumerate}
 \item Soient $F$ et $G$ des matrices de $M_3(\R)$ telles que $F=\exp G$. D\'emontrer que pour tout $n\in \N^*$, il existe une matrice $H$ telle que $H^n=F$. 
\end{enumerate}
\finenonce{002598} 


\finexercice


\exercice{3585, quercia, 2010/03/10}
\enonce{003585}{Ensi PC 1999}

Soit $A=\begin{pmatrix}-1&2&1\cr2&-1&-1\cr-4&4&3\cr\end{pmatrix}$.
\begin{enumerate}
  \item Calculer $A^n$.
    
  \item Soit $U_0=\begin{pmatrix}-2\cr4\cr1\end{pmatrix}$ et $(U_n)$ défini par la relation~:
    $U_{n+1} = AU_n$. Calculer $U_n$ en fonction de~$n$.
  \item Soit $X(t) = \begin{pmatrix}x(t)\cr y(t)\cr z(t)\cr\end{pmatrix}$.
    Résoudre $\frac{d X}{d t} = AX$.
\end{enumerate}
\finenonce{003585}



\finexercice


\exercice{3586, quercia, 2010/03/10}
\enonce{003586}{Puissances de $A$}

Soit $A \in \mathcal{M}_{3}(\R)$ ayant pour valeurs propres $1,-2,2$, et $n\in \N$.

\begin{enumerate}
  \item Montrer que $A^n$ peut s'écrire sous la forme :
    $A^n = \alpha_n A^2 + \beta_n A + \gamma_n I$
    avec $\alpha_n,\beta_n,\gamma_n \in \R$.

  \item On considère le polynôme $P = \alpha_n X^2 + \beta_n X + \gamma_n$.
    Montrer que : $P(1) = 1$, $P(2) = 2^n$, $P(-2) = (-2)^n$.

  \item En déduire les coefficients $\alpha_n,\beta_n,\gamma_n$.
    
\end{enumerate}
\finenonce{003586}



\finexercice


\exercice{3587, quercia, 2010/03/10}
\enonce{003587}{Suites récurrentes linéaires}

Soit $(u_n)$ une suite réelle vérifiant l'équation de récurrence :
$u_{n+3} = 6u_{n+2} - 11u_{n+1} + 6u_n$.

\begin{enumerate}
  \item On pose $X_n = \begin{pmatrix} u_n\cr u_{n+1} \cr u_{n+2} \cr \end{pmatrix}$.
    Montrer qu'il existe une matrice $A \in \mathcal{M}_{3}(\R)$ telle que $X_{n+1} = AX_n$.

  \item Diagonaliser $A$.
    En déduire $u_n$ en fonction de $u_0$, $u_1$, $u_2$ et $n$.
    
\end{enumerate}
\finenonce{003587}



\finexercice


\exercice{3588, quercia, 2010/03/10}
\enonce{003588}{Centrale P' 1996}

Soit $f$, endomorphisme d'un espace vectoriel $E$ de dimension~$n$.

\begin{enumerate}
  \item On suppose que pour tout sous-ev $D$ de dimension~1 il existe $x\in D$
    tel que $E = \text{vect}(x,f(x),f^2(x),\dots)$. Que dire de $E$ et $f$~?
  \item On suppose qu'il existe  $x\in E$
    tel que $E = \text{vect}(x,f(x),f^2(x),\dots)$.
    Montrer que si $f$ est diagonalisable alors ses valeurs propres sont
    toutes distinctes. Montrer que si $f$ est nilpotente alors $f^{n-1}\ne 0$.
\end{enumerate}
\finenonce{003588}



\finexercice


\exercice{3589, quercia, 2010/03/10}
\enonce{003589}{Chimie P' 1996}

Soit $(M_n)$ une suite de points dans le plan, de coordonnées $(x_n,y_n)$
définies par la relation de récurrence~:
$$\left\{
\begin{array}{lllll}
x_{n+1} &= &-x_n + 2y_n\cr y_{n+1} &= &-3x_n + 4y_n.\cr
\end{array}\right.$$

\begin{enumerate}
  \item Montrer que, quelque soit $M_0$, les points $M_n$ sont alignés.

  \item Étudier la suite $(M_n)$ quand $n$ tend vers l'infini.

  \item Quelle est la limite de $y_n/x_n$ (utiliser une méthode géométrique)~?

\end{enumerate}
\finenonce{003589}



\finexercice


\exercice{3590, quercia, 2010/03/10}
\enonce{003590}{Commutant d'une matrice à valeurs propres distinctes}

\begin{enumerate}
  \item Soit $D = \text{diag}(\lambda_1,\dots,\lambda_n)$ une matrice diagonale
    à valeurs propres distinctes.
  \begin{enumerate}
    \item Montrer qu'une matrice $M$ commute avec $D$ si et seulement si $M$ est
    diagonale.
    \item Montrer que pour toute matrice $M$ diagonale, il existe un polynôme
    $P \in  K_{n-1}[X]$ unique tel que $M = P(D)$.
  \end{enumerate}
  \item Soit $A \in \mathcal{M}_n(K)$ une matrice à valeurs propres distinctes.
    Montrer que les matrices $M$ commutant avec $A$ sont les polynômes en $A$.
\end{enumerate}
\finenonce{003590}



\finexercice


\exercice{3591, quercia, 2010/03/10}
\enonce{003591}{$XY = YX = A$}

Soit $A = \begin{pmatrix} 1 &1 \cr 1 &1 \cr \end{pmatrix}$.

\begin{enumerate}
  \item $A$ est-elle diagonalisable ?
  \item Trouver toutes les matrices $X,Y \in \mathcal{M}_{2}(K)$ telles que $XY = YX = A$.    
\end{enumerate}
\finenonce{003591}



\finexercice


\exercice{3592, quercia, 2010/03/10}
\enonce{003592}{Racine carrée}

Soit $A = \begin{pmatrix} 9 &0 &0 \cr 1 &4 &0 \cr 1 &1 &1 \cr \end{pmatrix}$.
Trouver les matrices $M \in \mathcal{M}_{3}(\R)$ telles que $M^2 = A$.

\finenonce{003592}



\finexercice


\exercice{3593, quercia, 2010/03/10}
\enonce{003593}{ Ensi Physique 93}

Soit $A = \begin{pmatrix}5 &-4 &1 \cr
                    8 &-7 &2 \cr
                   12 &-12&4 \cr\end{pmatrix}$.
Trouver une matrice $B$ différente de $A$ et $-A$ telle que $B^2=A$.

\finenonce{003593}



\finexercice


\exercice{3594, quercia, 2010/03/10}
\enonce{003594}{Esigelec 91}

    Trouver le commutant de $\begin{pmatrix}2&-2&1\\ 2&-3&2\\ 1&2&0\end{pmatrix}$.
    
\finenonce{003594}



\finexercice


\exercice{3595, quercia, 2010/03/10}
\enonce{003595}{Centrale MP 2000}

Si $A\in M_n(\C)$, on note $C(A)$ le commutant de $A$.
\begin{enumerate}
  \item Pour $n=2$, montrer que $C(A)$ est de dimension $2$ ou $4$, en donner une base.
    
  \item Pour $n\in \N^*$, montrer que $C(A)$ est de dimension $\ge n$ (traiter d'abord le cas où $A$
    est diagonalisable).
    
\end{enumerate}
\finenonce{003595}



\finexercice


\exercice{3596, quercia, 2010/03/10}
\enonce{003596}{Ulm MP$^*$ 2001}

En se déplaçant uniquement sur les arêtes d'un cube de côté 1, combien y a-t-il
de chemins de longueur~$n$ pour aller d'un point à un autre~?

\finenonce{003596}



\finexercice


\exercice{5651, rouget, 2010/10/16}
\enonce{005651}{**}
Soit $A =\left(
\begin{array}{ccc}
1&2&2\\
2&1&2\\
2&2&1
\end{array}
\right)$. Pour $n$ entier relatif donné, calculer $A^n$  par trois méthodes différentes.
\finenonce{005651}


\finexercice


\exercice{5652, rouget, 2010/10/16}
\enonce{005652}{**}
Résoudre dans $\mathcal{M}_3(\Rr)$ l'équation $X^2 = A$ où $A =\left(
\begin{array}{ccc}
3&0&0\\
8&4&0\\
5&0&1
\end{array}
\right)$.   
\finenonce{005652}


\finexercice


\exercice{5653, rouget, 2010/10/16}
\enonce{005653}{**}
Soit $A = \left(
\begin{array}{ccc}
3&1&0\\
-4&-1&0\\
4&8&-2
\end{array}
\right)$.

\begin{enumerate}
 \item  Vérifier que $A$ n'est pas diagonalisable.

\item  Déterminer $\text{Ker}(A-I)^2$.

\item  Montrer que $A$ est semblable à une matrice de la forme $\left(
\begin{array}{ccc}
a&0&0\\
0&b&c\\
0&0&b
\end{array}
\right)$

\item  Calculer $A^n$ pour $n$ entier naturel donné.
\end{enumerate}
\finenonce{005653}


\finexercice


\exercice{5664, rouget, 2010/10/16}
\enonce{005664}{***}
Soient $a$ et $b$ deux réels tels que $|a|\neq|b|$. Soit $A=\left(\begin{array}{cccc}
0&b&\ldots&b\\
a&\ddots&\ddots&\vdots\\
\vdots&\ddots&\ddots&b\\
a&\ldots&a&0
\end{array}
\right)$.

Montrer que les images dans le plan complexe des valeurs propres de $A$ sont cocycliques. (Indication : pour calculer $\chi_A$, considérer $f(x) =\left|
\begin{array}{cccc}
-X+x&b+x&\ldots&b+x\\
a+x&\ddots&\ddots&\vdots\\
\vdots&\ddots&\ddots&b+x\\
a+x&\ldots&a+x&-X+x
\end{array}
\right|$.)
\finenonce{005664}


\finexercice


\exercice{5665, rouget, 2010/10/16}
\enonce{005665}{***I Matrices stochastiques}

Soit $A=(a_{i,j})_{1\leqslant i,j\leqslant n}\in\mathcal{M}_n(\Rr)$ telle que $\forall(i,j)\in\llbracket1,n\rrbracket^2,\;a_{i,j}\in[0,1]$ et $\forall i\in\llbracket1,n\rrbracket$, $\sum_{j=1}^{n}a_{i,j}= 1$.

\begin{enumerate}
 \item  Montrer que $1$ est valeur propre de $A$.

\item  Soit $\lambda$ une valeur propre de $A$.
  \begin{enumerate}
   \item Montrer que $|\lambda|\leqslant 1$.
   \item Montrer qu'il existe un réel $\omega$ de $[0,1]$ tel que $|\lambda-\omega|\leqslant 1-\omega$. Conséquence géométrique ?
  \end{enumerate}
\end{enumerate}
\finenonce{005665}


\finexercice


\exercice{5668, rouget, 2010/10/16}
\enonce{005668}{***I Déterminant circulant}
\label{ex:rou18}
\begin{enumerate}
 \item   Soit $J_n=\left(
\begin{array}{ccccc}
0&1&0&\ldots&0\\
\vdots&\ddots&\ddots&\ddots&\vdots\\
\vdots& & &\ddots&0\\
0& & &\ddots&1\\
1&0&\ldots&\ldots&0
\end{array}
\right)$   (de format $n\geqslant 3$). Diagonaliser $J_n$.

\item  En déduire la valeur de $\left|
\begin{array}{ccccc}
a_0&a_1&\ldots&a_{n-2}&a_{n-1}\\
a_{n-1}&a_0&a_1& &a_{n-2}\\
\vdots& &\ddots&\ddots&\vdots\\
a_2& &\ddots&a_0&a_1\\
a_1&a_2&\ldots&a_{n-1}&a_0
\end{array}
\right|$.
\end{enumerate}
\finenonce{005668}


\finexercice


\exercice{5669, rouget, 2010/10/16}
\enonce{005669}{***I Matrices de permutations}

Pour $\sigma\in S_n$, $n\geqslant 2$, on définit la matrice $P_\sigma$ par $P_\sigma=(\delta_{i,\sigma(j)})_{1\leqslant i,j\leqslant n}$.

\begin{enumerate}
 \item  Calculer $\text{det}(P_\sigma)$ pour tout $\sigma\in S_n$.

\item
 \begin{enumerate}
 \item Montrer que $\forall(\sigma,\sigma')\in S_n^2$, $P_\sigma\times P_{\sigma'}=P_{\sigma\circ\sigma'}$.
 \item On pose $G=\{P_\sigma,\;\sigma\in S_n\}$. Montrer que $(G,\times)$ est un groupe isomorphe à $S_n$.
  \end{enumerate}

\item  Soit $A=(a_{i,j})_{1\leqslant i,j\leqslant n}\in \mathcal{M}_n(\Cc)$. Calculer $AP_\sigma$.

\item  Trouver les valeurs propres d'une matrice de pemutation (on pourra utiliser le résultat hors programme : toute permutation se décompose de manière unique à l'ordre près des facteurs en produit de cycles à supports disjoints).
\end{enumerate}
\finenonce{005669}


\finexercice


\exercice{5671, rouget, 2010/10/16}
\enonce{005671}{**}
Trouver une matrice carrée $A$ vérifiant $A^4 - 3A^3 + A^2 - I = 0$.
\finenonce{005671}


\finexercice


\exercice{5672, rouget, 2010/10/16}
\enonce{005672}{**I}
Calculer $\left|
\begin{array}{cccc}
a&b&\ldots&b\\
b&a&\ddots&\vdots\\
\vdots&\ddots&\ddots&b\\
b&\ldots&b&a
\end{array}
\right|$.
\finenonce{005672}


\finexercice


\exercice{5676, rouget, 2010/10/16}
\enonce{005676}{**I}
Résoudre dans $\mathcal{M}_3(\Cc)$ l'équation $X^2=\left(\begin{array}{ccc}
0&1&0\\
0&0&1\\
0&0&0
\end{array}
\right)$.
\finenonce{005676}


\finexercice


\exercice{5680, rouget, 2010/10/16}
\enonce{005680}{** (ESTP1994)}
Soit $M_{a,b}=\left(
\begin{array}{cccc}
1&1&0&0\\
0&1&a&0\\
0&0&1&b\\
0&0&0&1
\end{array}
\right)$. Peut-on trouver deux matrices distinctes semblables parmi
les quatre matrices $M_{0,0}$, $M_{0,1}$, $M_{1,0}$ et $M_{1,1}$ ?
\finenonce{005680}


\finexercice


\exercice{5684, rouget, 2010/10/16}
\enonce{005684}{***}
Résoudre dans $\mathcal{M}_3(\Cc)$ l'équation $X^2= 
\left(
\begin{array}{ccc}
1&3&-7\\
2&6&-14\\
1&3&-7
\end{array}
\right)$.
\finenonce{005684}


\finexercice


\exercice{5685, rouget, 2010/10/16}
\enonce{005685}{}
Commutant de $\left(
\begin{array}{ccc}
1&0&-1\\
1&2&1\\
2&2&3
\end{array}
\right)$.
\finenonce{005685}


\finexercice

\section{201.08 Polynôme annulateur}


\exercice{1574, legall, 1998/09/01}

\enonce{001574}{}
Soit $  A\in M_3 (\R )  $ la matrice $  \begin{pmatrix} 0 & 1 & 1 \cr 
                                      1 & 0 & 1 \cr
                         0 & 0 & 1\cr \end{pmatrix} .$ 
 Calculer le polyn\^ome minimal de $  A  .$ En d\'eduire $  A^{-1}  ,$ $  A^{3}  $ et $  
 A^{5}  .$
\finenonce{001574}



\finexercice



\exercice{1575, legall, 1998/09/01}

\enonce{001575}{}
Soit $  P\in \C [X]  $ tel que $  P(0)=0  $ et $  P'(0)\not = 0  .$
Soit $  E  $ un $  \C $-espace vectoriel de dimension finie et $  f \in \mathcal{L} (E)  $ telle
que $  P(f)=0  .$ 

Montrer que $  \hbox{Ker}(f)=\hbox{Ker}(f^2)  ;$ en d\'eduire $  E=\hbox{Ker}(f)\oplus
\hbox{Im}(f).$
\finenonce{001575}



\finexercice



\exercice{1576, legall, 1998/09/01}

\enonce{001576}{}
 Soit $  E  $ un $  \mathbb{K} $-espace vectoriel de dimension finie
$  n  $ et $  f
\in \mathcal{L} (E)  $ tel que $  \hbox{rg}(f-id)=1  .$ On note $  H = \hbox{Ker}(f-id)  .$
\begin{enumerate}
    \item Soit $  \{ e_1, \cdots , e_{n-1}\}   $ une base de $  H   $ et $  e_n \notin H  .$ Montrer que $  \{ e_1, \ldots , e_{n}\}   $ est une base de $  E   $ et donner l'allure de la matrice de $  f   $ dans cette base. 
    \item Montrer que le polyn\^ome $  (X-1)(X-\hbox{det}(f))  $ annule $  f  .$ Donner une condition n\'ec\'essaire et
suffisante pour que $  f  $ soit diagonalisable.
\end{enumerate}
\finenonce{001576}



\finexercice



\exercice{1577, barraud, 2003/09/01}

\enonce{001577}{}
Soit $E$ un espace vectoriel de dimension $n$, et $u$ un endomorphisme de
$E$ nilpotent, c'est à dire tel que $\exists m\in\N, u^{m}=0$. Montrer que
$u^{n}=0$
\finenonce{001577}



\finexercice



\exercice{1578, barraud, 2003/09/01}

\enonce{001578}{}
Déterminer toutes les matrices $A$ de $\mathcal{M}_{2,2}(\R)$ telles que
$$
  A^{2}-3A+2\mathrm{id}=0
$$
Même question pour
$$
  A^{3}-8A^{2}+21A-18\mathrm{id}=0
$$
\finenonce{001578}



\finexercice



\exercice{1579, barraud, 2003/09/01}

\enonce{001579}{}
 \'Enoncer le théorème de Cayley-Hamilton.  Le démontrer dans le cas
 particulier où le polynôme caractéristique est scindé à racines simples.

\finenonce{001579}



\finexercice



\exercice{1580, barraud, 2003/09/01}

\enonce{001580}{}
  \begin{enumerate}
  \item 
    Réduire la matrice 
    $$
    A=
    \begin{pmatrix}
      2&  0& 0\\
      3& -4& 3\\
      3& -6& 5
    \end{pmatrix}
    $$
  \item
    Donner un polynôme annulateur de $A$ de degré 2.
  \item
    En  déduire qu'il existe des coefficients $a_{n}$ et $b_{n}$ tels que
    $A^{n}=a_{n}A+b_{n}$ et les calculer en fonction de $n$.
  \end{enumerate}
\finenonce{001580}



\finexercice



\exercice{1581, bodin, 1999/11/01}

\enonce{001581}{}
Soit $A\in\mathcal{M}_2(\C)$ de trace non
nulle.  Montrer que toute matrice
$M\in\mathcal{M}_2(\C)$ qui commute avec $A^2$
commute aussi avec $A$. (\emph{Indication :}
utiliser Cayley-Hamilton.)
\finenonce{001581}



\finexercice



\exercice{1582, ortiz, 1999/04/01}

\enonce{001582}{}
Que peut-on dire d'un endomorphisme d'un
$K$-espace vectoriel de dimension finie annul\'e
par les polyn\^omes $P=1-X^3$ et $Q=X^2-2X+1$?
\finenonce{001582}



\finexercice



\exercice{1583, ortiz, 1999/04/01}

\enonce{001583}{}
Soient $E$ un $K$-espace vectoriel de dimension
finie et $f\in\mathcal{L}(E).$ On suppose que le
polyn\^ome minimal de $f$ est $P=(X-2)(X-1)^2.$
Quel est le polyn\^ome minimal de $f+\text{Id}_E$?
\finenonce{001583}



\finexercice



\exercice{1584, ortiz, 1999/04/01}

\enonce{001584}{}
Soit $M\in\mathcal{M}_n(K)$ une matrice diagonale.
Si $P\in K[X]$, calculer $P(M)$ et en d\'eduire le
polyn\^ome minimal de $M.$
\finenonce{001584}



\finexercice



\exercice{1585, ortiz, 1999/04/01}

\enonce{001585}{}
En appliquant la m\'ethode utilis\'ee en cours
pour d\'emontrer l'existence d'un polyn\^ome
annulateur d'un endomorphisme d'un espace
vectoriel de dimension finie, d\'eterminer le
polyn\^ome minimal de la matrice
$B=\left(\begin{smallmatrix}
2&1\\-1&1\end{smallmatrix}\right)$.
\finenonce{001585}



\finexercice



\exercice{1586, ortiz, 1999/04/01}

\enonce{001586}{}
Quel est le polyn\^ome minimal d'un endomorphisme
d'une droite vectorielle ?
\finenonce{001586}



\finexercice



\exercice{1587, ortiz, 1999/04/01}

\enonce{001587}{}
Soient $E$ un espace vectoriel de dimension $n\ge
2$ et $f$ un endomorphisme de $E$ de rang 1.
Montrer que le polyn\^ome minimal de $f$ est de la
forme $X(X-\lambda).$
\finenonce{001587}



\finexercice



\exercice{1588, ortiz, 1999/04/01}

\enonce{001588}{}
D\'eterminer les endomorphismes d'un $K$-espace
vectoriel $E$ de dimension finie $n$ dont le
polyn\^ome minimal est de degr\'e 1.
\finenonce{001588}



\finexercice



\exercice{1589, ortiz, 1999/04/01}

\enonce{001589}{}
\begin{enumerate}
    \item  Montrer que $P=(X-1)^2(X-2)$ est un polyn\^ome annulateur
de la matrice $A=\left(\begin{smallmatrix}
1&0&0\\0&1&0\\0&0&2\end{smallmatrix}\right)$ et en
d\'eduire le polyn\^ome minimal de la matrice $A$.
    \item  Soit $B\in\mathcal{M}_2(\Cc)$. Calculer explicitement
$B^2-\text{tr}(B)\,B+\det(B)I_2$. En d\'eduire le
polyn\^ome minimal de la matrice
$B=\left(\begin{smallmatrix}
3&1\\-1&1\end{smallmatrix}\right)$.
\end{enumerate}
\finenonce{001589}



\finexercice



\exercice{1590, ortiz, 1999/04/01}

\enonce{001590}{}
Soient $E$ un $K$-espace vectoriel de dimension
finie, $f\in\mathcal{L}(E)$ et $P$ son polyn\^ome
minimal. Montrer que $f$ est bijective si et
seulement si $P(0)\not= 0.$
\finenonce{001590}



\finexercice



\exercice{1591, ortiz, 1999/04/01}

\enonce{001591}{}
Soit $f$ un endomorphisme d'un $\Rr$-espace
vectoriel $E$ de dimension 3. Montrer que $f$
admet un plan stable (on discutera en fonction du
caract\`ere trigonalisable de $f$).
\finenonce{001591}



\finexercice



\exercice{1592, ortiz, 1999/04/01}

\enonce{001592}{}
Soit $f$ un endomorphisme d'un $K$-espace
vectoriel $E$ de dimension finie tel que
$${f^4=f^2+f.}$$
\begin{enumerate}
    \item  Montrer que $\ker(f^3-f-\text{Id})\oplus \ker f=E.$
    \item
    \begin{enumerate}
        \item Montrer que $\text{Im\,} f\subseteq \ker(f^3-f-\text{Id}).$
        \item En d\'eduire que $\text{Im\,} f= \ker(f^3-f-\text{Id}).$
    \end{enumerate}
\end{enumerate}
\finenonce{001592}



\finexercice



\exercice{1593, ortiz, 1999/04/01}

\enonce{001593}{}
D\'eterminer le polyn\^ome minimal de la matrice
$A=\left(\begin{smallmatrix}
7&3&-4\\-6&-2&5\\4&2&-1\end{smallmatrix}\right).$
\finenonce{001593}



\finexercice



\exercice{1594, ortiz, 1999/04/01}

\enonce{001594}{}
Soient $J=\left(\begin{smallmatrix}
1&1\\1&1\end{smallmatrix}\right)$ et la matrice
par blocs \`a coefficients r\'eels suivante
$${M=\begin{pmatrix}O&\frac 12 J\\\frac 12
J&O\end{pmatrix}.}$$
\begin{enumerate}
    \item  Calculer $M^2$ et $M^3$ et en d\'eduire que $M$ est diagonalisable.
    \item  D\'eterminer le polyn\^ome caract\;eristique
 et le polyn\^ome minimal de $M$.
\end{enumerate}
\finenonce{001594}



\finexercice



\exercice{1595, gineste, 2001/11/01}

\enonce{001595}{}
 On considére la matrice
$$ A=\left[ \begin{array}{rrr}3&-2&-1\\
2&-1&1\\
6&3&-2
 \end{array} \right].$$
Calculer son polynôme caractéristique, calculer $A^2$ et déduire de ces calculs 
et du théorème de Cayley-Hamilton l'inverse de $A.$
\finenonce{001595}



\finexercice



\exercice{1596, gineste, 2001/11/01}

\enonce{001596}{}
On se place dans $E=\C\,\,^4$ muni de sa base canonique $b=(e_1,e_2,e_3,e_4).$ On d\'esigne par $j$ l'endomorphisme de $E$ dont la matrice dans $b$ est la matrice suivante
 $$J=\left( \begin{array}{cccc}
0&1 &0&0\\
0&0&1&0\\
0&0&0&1\\
1&0&0&0
\end{array}
\right)\in M_4(\C).$$

\begin{enumerate}
\item D\'eterminer l'image de $b$ par $j,$ $j^2,$ $j^3,$ et $j^4.$
\item En d\'eduire $J^2,$ $J^3$ et $J^4.$
\item D\'eterminer un polyn\^ome annulateur  non nul de $J.$
\item Montrer que si $P\in\C[X]$ avec $\mathrm{deg}(P)\leq 3$ v\'erifie $P(J)=0$
alors $P=0$.
\item En d\'eduire le polyn\^ome minimal de $J.$
\item Montrer que $J$ est diagonalisable.
\item  D\'eterminer les valeurs propres de $J.$
\end{enumerate}
\finenonce{001596}



\finexercice



\exercice{2569, delaunay, 2009/05/19}
\enonce{002569}{}
Soit $A$ la matrice suivante
$$A=\begin{pmatrix}0&1&1 \\  1&0&1 \\  1&1&0\end{pmatrix}$$
Calculer $A^2$ et v\'erifier que $A^2=A+2I_3$. En d\'eduire que $A$ est inversible et donner son inverse en fonction de $A$.
\finenonce{002569} 


\finexercice


\exercice{2588, delaunay, 2009/05/19}
\enonce{002588}{}
 Soit $N$ une matrice nilpotente, il existe $q\in\N$ tel que $N^q=0$. Montrer que la matrice $I-N$ est inversible et exprimer son inverse en fonction de $N$.
\finenonce{002588} 


\finexercice


\exercice{3518, quercia, 2010/03/10}
\enonce{003518}{\'Etude d'une matrice}

Soit $A=\begin{pmatrix}a_1 &1  &       &(0)\cr
                 a_2 &   &\ddots &   \cr
              \vdots &   &       &1  \cr
                 a_n &(0)&       &0  \cr\end{pmatrix}$
où les $a_i$ sont des réels positifs ou nuls, avec $a_1a_n > 0$.
\begin{enumerate}
  \item Quel est le polynôme caractérique de~$A$~?

  \item Montrer que $A$ admet une unique valeur propre $r>0$ et que l'on
    a~$r < 1 + \max(a_1,\dots,a_n)$.

  \item Soit $\lambda$ une valeur propre complexe de~$A$.
    Montrer que $|\lambda| \le r$ et $|\lambda|=r  \Rightarrow  \lambda=r$.

  \item Montrer qu'il existe un entier $k$ tel que $A^k$ a tous ses
    coefficients strictement positifs.

\end{enumerate}
\finenonce{003518}



\finexercice

\exercice{3541, quercia, 2010/03/10}
\enonce{003541}{Matrice bitriangulaire}

Donner une CNS sur $a,b\in\C$ pour que la matrice
$M=\begin{pmatrix}0&&(a)\cr&\ddots\cr(b)&&0\cr\end{pmatrix}$ soit diagonalisable.

\finenonce{003541}



\finexercice


\exercice{3542, quercia, 2010/03/10}
\enonce{003542}{$A^2 = A$ et $tr(A) = 0$}

Trouver les matrices $A \in \mathcal{M}_n(\R)$ telles que $A^2 = A$ et $\mathrm{tr}(A) = 0$.
\finenonce{003542}



\finexercice


\exercice{3543, quercia, 2010/03/10}
\enonce{003543}{Matexo}

Soient $E$ un ev de dimension finie sur $\C$ et $u$ un endomorphisme de $E$.

On suppose que $u^3 = u^2, u \ne \mathrm{id}, u^2 \ne 0, u^2 \ne u$.

\begin{enumerate}
  \item Montrer qu'une valeur propre de $u$ ne peut être que $0$ ou $1$.
  \item Montrer que $1$ et $0$ sont effectivement valeurs propres de $u$.
  \item Montrer que u {\it n'est pas} diagonalisable.
  \item Montrer que $E = {\rm Im} (u^2) \oplus {\rm Ker} (u^2)$.
  \item Monter que $u|_F$ avec $F= {\rm Im} (u^2)$ est l'identité.
\end{enumerate}
\finenonce{003543}



\finexercice


\exercice{3544, quercia, 2010/03/10}
\enonce{003544}{INT gestion 94}

    Soit $A=\left(\begin{smallmatrix}1&1&1&1\cr -1&1&-1&1\cr -1&1&1&-1\cr -1&-1&1&1\cr\end{smallmatrix}\right)$.\par
 \begin{enumerate}   
      \item Calculer $\det A$.
      \item Calculer $(A-xI)(\,{}^t\!A-xI)$ et en déduire $\chi_A(x)$.
    
      \item Montrer que $A$ est $\C$-diagonalisable.
    
\end{enumerate}
\finenonce{003544}



\finexercice


\exercice{3545, quercia, 2010/03/10}
\enonce{003545}{$X^nP(1/X)$}

Soit $E =  K_n[X]$ et $ u : E \to  E, P  \mapsto {X^nP(1/X).}$

\begin{enumerate}
  \item Déterminer $u\circ u$. En déduire que $u$ est diagonalisable.
  \item Donner une base de vecteurs propres de $u$.
\end{enumerate}
\finenonce{003545}



\finexercice


\exercice{3546, quercia, 2010/03/10}
\enonce{003546}{TPE 93}

Soit $A \in \mathcal{M}_n(\C)$ telle que $A = A^{-1}$. $A$ est-elle diagonalisable ?
Calculer $e^A$. ($e^A = \sum_{k=0}^\infty \frac{A^k}{k!}$)

\finenonce{003546}



\finexercice


\exercice{3547, quercia, 2010/03/10}
\enonce{003547}{Ensi Physique 93}

Soit $E$ un ev de dimension finie et $u \in \mathcal{L}(E)$ tel que $\mathrm{rg}(u)=1$.
Montrer que : $$\Im u \subset \mathrm{Ker} u \Leftrightarrow u \text{ n'est pas diagonalisable.}$$

\finenonce{003547}



\finexercice


\exercice{3548, quercia, 2010/03/10}
\enonce{003548}{$u^2$ diagonalisable}

Soit $E$ un $\C$-ev de dimension finie et $u \in GL(E)$ tel que $u^2$
est diagonalisable. Montrer que $u$ est diagonalisable.
\finenonce{003548}



\finexercice


\exercice{3549, quercia, 2010/03/10}
\enonce{003549}{Ensi PC 1999}

Soit $A\in\mathcal{M}_n(\C)$ inversible diagonalisable et $B\in \mathcal{M}_n(\C)$, $p\in\N^*$
tels que $B^p=A$.
 
\begin{enumerate}
  \item Montrer que $B$ est diagonalisable.


  \item Si $A$ n'est pas inversible la conclusion subsiste-t-elle~?
\end{enumerate}
\finenonce{003549}



\finexercice


\exercice{3550, quercia, 2010/03/10}
\enonce{003550}{Ensi P 90}

    Soit $E$ un espace vectoriel de dimension $n$ et $p \in \mathcal{L}(E)$ tel que
    $p^2$ est un projecteur. Quelles sont les valeurs propres éventuelles de
    $p$ ? Montrer que $p$ est diagonalisable si et seulement si $p^3=p$.
    
\finenonce{003550}



\finexercice


\exercice{3551, quercia, 2010/03/10}
\enonce{003551}{$A^3 = A + I$}

Soit $A \in \mathcal{M}_n(\R)$ telle que $A^3 = A + I$. Montrer que $\det(A) > 0$.
\finenonce{003551}



\finexercice


\exercice{3552, quercia, 2010/03/10}
\enonce{003552}{Mines-Ponts PC 1999}
Soit $A \in \mathcal{M}_n(\R)$ telle que $A^3 + A^2 + A = 0$. Montrer que $\mathrm{rg} A$ est pair.
\finenonce{003552}



\finexercice


\exercice{3553, quercia, 2010/03/10}
\enonce{003553}{Esem 91}

    Soit $A \in \mathcal{M}_n(\C)$ telle que $A^n = I$ et $(I,A,\dots,A^{n-1})$ est libre. Montrer qu'alors on a $\mathrm{tr}(A) = 0$.
    
\finenonce{003553}



\finexercice


\exercice{3554, quercia, 2010/03/10}
\enonce{003554}{$A^p = I$ et $\mathrm{spec}(A) \subset \R  \Rightarrow  A^2 = I$}

Soit $A \in \mathcal{M}_n(\R)$. On suppose que les valeurs propres de $A$ sont réelles
et qu'il existe $p \ge 1$ tel que $A^p = I$.
Montrer que $A^2 = I$.

\finenonce{003554}



\finexercice


\exercice{3555, quercia, 2010/03/10}
\enonce{003555}{$P(u) = \sum P(\lambda_i) u_i$}

Soit $E$ un $ K$-ev de dimension finie et $u \in \mathcal{L}(E)$.

\begin{enumerate}
  \item On suppose $u$  diagonalisable et on note
    $\lambda_1,\dots,\lambda_p$ ses valeurs propres distinctes.
  \begin{enumerate}
     \item Montrer qu'il existe des endomorphismes $u_1,\dots,u_p$ tels que pour tout
    polynôme $P \in  K[X]$, on ait :\par
    $P(u) = \sum_{i=1}^p P(\lambda_i) u_i$.

     \item Montrer qu'il existe un polynôme $P_i$ tel que $u_i = P_i(u)$.
  \end{enumerate}
  \item Réciproquement, soit $u,u_1,\dots,u_p\in\mathcal{L}(E)$ et $\lambda_1,\dots,\lambda_p\in K$
tels que pour tout $P\in K[X]$, $P(u) = \sum_{i=1}^p P(\lambda_i) u_i$.
Montrer que $u$ est diagonalisable et $\mathrm{Sp}(u) \subset \{\lambda_1,\dots,\lambda_p\}$.
\end{enumerate}
\finenonce{003555}



\finexercice


\exercice{3556, quercia, 2010/03/10}
\enonce{003556}{Mines PSI 1998}

Soit $f$ un endomorphisme diagonalisable d'un ev~$E$ de dimension finie,
$\lambda$ une valeur propre de~$f$ et $p_\lambda$ le projecteur sur le
sous-espace propre associé parallèlement à la somme des autres sous-espaces
propres. Montrer que $p_\lambda$ est un polynôme en~$f$.

\finenonce{003556}



\finexercice


\exercice{3557, quercia, 2010/03/10}
\enonce{003557}{Endomorphismes anticomutant (Centrale MP 2003)}
Soit~$E$ un~$\C$-ev de dimension~$n\in\N^*$ et $u_1,\dots,u_p$ ($p\ge 2$)
des endomorphismes de~$E$ vérifiant~:
$$\forall\ k,\ u_k^2 = -\mathrm{id}_E,
  \quad\forall\ k\ne\ell,\ u_k\circ u_\ell = -u_\ell\circ u_k.$$
  
  \begin{enumerate}
  \item Montrer que les~$u_k$ sont des automorphismes et qu'ils sont diagonalisables.
  \item Montrer que~$n$ est pair.
  \item Donner le spectre de chaque~$u_k$.
    
  \item Donner les ordres de multiplicité des valeurs propres des~$u_k$.
    
  \item Calculer~$\det(u_k)$.
\end{enumerate}
\finenonce{003557}



\finexercice


\exercice{3558, quercia, 2010/03/10}
\enonce{003558}{Ensi PC 1999}

Soit $E$ un ev de dimension finie et $u\in \mathcal{L}(E)$ tel que $u\circ u = 0$.

\begin{enumerate}
  \item Quelle relation y a-t-il entre $\mathrm{Ker} u$ et $\Im u$~?
    Montrer que $2\mathrm{rg} u \le \dim E$.

  \item On suppose ici $\dim E = 4$ et $\mathrm{rg} u = 2$.
    Montrer qu'il existe une base $(\vec{e_1},\vec{e_2},\vec{e_3},\vec{e_4})$
    de~$E$ telle que~:
    $u(\vec{e_1}) = \vec{e_2}$, $u(\vec{e_2}) = \vec{0}$,
    $u(\vec{e_3}) = \vec{e_4}$, $u(\vec{e_4}) = \vec{0}$.

  \item On suppose $\dim E = n$ et $\Im u = \mathrm{Ker} u$.
    Est-ce que $u$ est diagonalisable~?
\end{enumerate}
\finenonce{003558}



\finexercice


\exercice{3559, quercia, 2010/03/10}
\enonce{003559}{Réduction de $M$ tq $M^3 = I$}

Soit $M \in \mathcal{M}_{3}(\R)$ telle que $M \ne I$, et $M^3 = I$.

\begin{enumerate}
  \item Quelles sont les valeurs propres complexes de $M$ ?
    (On vérifiera que ce sont effectivement des valeurs propres de $M$)

  \item Montrer que $M$ est semblable à
    $\begin{pmatrix} 1 &0    &0         \cr 
               0 &-1/2 &-\sqrt3/2 \cr
               0 &\sqrt3/2 &-1/2  \cr \end{pmatrix}$.
\end{enumerate}
\finenonce{003559}



\finexercice


\exercice{3560, quercia, 2010/03/10}
\enonce{003560}{Centrale PSI 1998}

Soient $u,v,h$ trois endomorphismes de $\R^n$ tels que~:
$$u\circ v = v\circ u,\quad
u\circ h - h\circ u = -2u,\quad
v\circ h - h\circ v = -2v.$$

\begin{enumerate}
  \item Cas particulier, $n=3$,
    $\text{Mat}(u) = \begin{pmatrix} 0&1&0\cr 0&0&1\cr 0&0&0\cr\end{pmatrix}$.
    Déterminer si $v$ et $h$ existent et si oui, les donner.
  \item Cas général.
  \begin{enumerate}
     \item Que peut-on dire de $\mathrm{tr}(u)$ et $\mathrm{tr}(v)$~?
     \item Montrer que $u$ et $v$ sont non inversibles.
    Montrer que $\mathrm{Ker} u$ et $\mathrm{Ker} v$ sont stables par~$h$.
     \item Déterminer $u^k\circ h - h\circ u^k$ pour $k\in\N$.
    Déterminer $P(u)\circ h - h \circ P(u)$ pour $P\in\R[X]$.
     \item Quel est le polynôme minimal de~$u$~?
  \end{enumerate}
\end{enumerate}
\finenonce{003560}



\finexercice


\exercice{3561, quercia, 2010/03/10}
\enonce{003561}{Indépendance du polynôme minimal par rapport au corps}

Soient $ K\subset L$ deux corps et $A\in \mathcal{M}_n(K)$.
On note $\mu_ K(A)$ et $\mu_{L}(A)$ les polynômes minimaux
de~$A$ en tant que matrice à coefficients dans $ K$ ou dans $L$.
Montrer que ces polynômes sont égaux.
\finenonce{003561}



\finexercice


\exercice{3562, quercia, 2010/03/10}
\enonce{003562}{Polynôme minimal et caractéristique}

Soit $A\in\mathcal{M}_n(K)$. Montrer que $\chi_A$ et $\mu_A$ ont les mêmes facteurs
irréductibles.
\finenonce{003562}



\finexercice


\exercice{3563, quercia, 2010/03/10}
\enonce{003563}{X MP$^*$ 2004}
Caractériser les polynômes $P$ tels que~: $\forall\ A\in\mathcal{M}_n(\C)$, $(P(A) = 0)  \Rightarrow  (\mathrm{tr}(A)\in\Z)$.


\finenonce{003563}



\finexercice


\exercice{3564, quercia, 2010/03/10}
\enonce{003564}{TPE MP 2005}

Soient $A,B,C\in\mathcal{M}_n(\C)$ telles que $AC = CB$ et $\mathrm{rg}(C) = r$.
Montrer que $A$ et $B$ ont au moins $r$ valeurs propres communes.


\finenonce{003564}



\finexercice


\exercice{3565, quercia, 2010/03/10}
\enonce{003565}{Polynôme minimal imposé, Centrale MP 2005}

Le polyn{\^o}me $X^{4}+X^{3}+2X^{2}+X+1$ peut-il {\^e}tre le polyn{\^o}me minimal d'une matrice de $M_{5}(\R)$ ?

\finenonce{003565}



\finexercice


\exercice{3566, quercia, 2010/03/10}
\enonce{003566}{$\mathrm{Ker} u^p \oplus \Im u^p$, Polytechnique MP$^*$ 2006}
Soit $E$ un $ K$-ev de dimension~$n$.
Soit $u\in\mathcal{L}(E)$, $P$ son polynôme minimal et $p$ le plus petit exposant
de~$X$ dans l'écriture de~$P$.

\begin{enumerate}
  \item Si $p=0$, que dire de~$u$~?
    
  \item Si $p=1$, montrer que $E=\Im u \oplus \mathrm{Ker} u$.
    
  \item Dans le cas général, montrer que $E = \mathrm{Ker} u^p \oplus \Im u^p$.
    
\end{enumerate}
\finenonce{003566}



\finexercice

\exercice{5658, rouget, 2010/10/16}
\enonce{005658}{****}
Soit $A$ une matrice carrée de format $n$.

Montrer que $A$ est nilpotente si et seulement si $\forall k\in\llbracket1,n\rrbracket$, $\text{Tr}(A^k) =0$.
\finenonce{005658}


\finexercice


\exercice{5659, rouget, 2010/10/16}
\enonce{005659}{*** I}
\label{ex:rou9}
Soient $f$ et $g$ deux endomorphismes d'un espace vectoriel de dimension finie vérifiant $fg - gf = f$. Montrer que $f$ est nilpotent.

\finenonce{005659}


\finexercice

\section{201.99 Autre}


\exercice{1699, bodin, 1999/11/01}

\enonce{001699}{}
Soit $u\in\mathcal{L}(\R^4)$ de matrice dans la base canonique :
\[
A=
\left(
\begin{array}{cccc}
1 & -1 & 2  & -2 \\
0 & 0 & 1 & -1 \\
1 & -1 & 1 & 0 \\
1 & -1 & 1 & 0
\end{array}
\right).
\]
\begin{enumerate}
\item  D\'eterminer le polyn\^ome
caract\'eristique $P_u$ de $u$.  Trouver les valeurs propres et
les sous-espaces caract\'eristiques $F_i$.
\item  Donner une base
suivant laquelle la matrice de $u$ se d\'ecompose en deux blocs
diagonaux.
\item  Donner les projections $p_i$ de $\R^4$ sur
$F_i$.
\end{enumerate}
\finenonce{001699}



\finexercice



\exercice{1700, legall, 1998/09/01}

\enonce{001700}{}
Soit $  A\in M_3 (\R )  $ telle que $  A^3=-A  $ et $  A \not = 0  .$ Montrer
que $  A   $ est semblable \`a $  \begin{pmatrix} 0 & 0 & 0 \cr
                                      0 & 0 & -1 \cr
                                      0 & 1 & 0 \cr \end{pmatrix} .$
\finenonce{001700}



\finexercice



\exercice{1701, ortiz, 1999/04/01}

\enonce{001701}{}
Soient $n\in\Nn\setminus\left\{0\right\}$ et $f$
l'endomorphisme de l'espace vectoriel $\Rr^{2n}$
dont la matrice dans la base canonique est la
matrice par blocs $M=\left(\begin{smallmatrix}
I_n&I_n\\O_n&O_n\end{smallmatrix}\right)\in
\mathcal{M}_{2n}(\Rr)$
.
\begin{enumerate}
    \item  D\'eterminer le polyn\^ome caract\'eristique de $M$.
    \item
    \begin{enumerate}
        \item D\'eterminer le noyau de $f$.
        \item Montrer que $f$ est diagonalisable.
    \end{enumerate}
\end{enumerate}
\finenonce{001701}



\finexercice



\exercice{1702, barraud, 2003/09/01}

\enonce{001702}{}
Soit $E$ un $\R$-espace vectoriel de dimension finie $n$, et $u$ un
endomorphisme de $E$.

Soit $x_{0}\in E\setminus\{0\}$. On note $x_{k}=u^{k}(x_{0})$ et $F$ le
sous espace vectoriel engendré par la famille $\{x_{k}, k\in\N\}$, c'est à
dire l'ensemble des combinaisons linéaires finies de vecteurs de ${x_{k},
  k\in\N}$~:
$$
 F=\left\{x\in E\ /\ \exists N\in\N,
   \exists(\alpha_{0}\dots\alpha_{N})\in\R^{N+1},\
   x=\sum_{i=0}^{N}\alpha_{i}x_{i}\right\}
$$

\begin{enumerate}
\item
Montrer que $F$ est stable par $u$, c'est à dire que $\forall x\in F,
u(x)\in F$. 

\item 
Montrer qu'il existe un entier $k\leq n$ tel que
$(x_{0},x_{1},\dots,x_{k})$ soit libre et $(x_{0},x_{1},\dots,x_{k+1})$
soit liée. Montrer alors qu'il existe des scalaires
$(a_{0},a_{1},\dots,a_{k})$ tels que
$$
  x_{k+1}=a_{0}x_{0}+a_{1}x_{1}+\dots+a_{k}x_{k}
$$

\item\label{q:P0(u)x0=0}
En déduire que le polynôme $P_{0}=X^{k+1}-\sum_{i=0}^{k}a_{i}X^{i}$
satisgfait $\big(P_{0}(u)\big)(x_{0})=0$.

\item\label{q:x=P(u)x0}
Montrer que pour tout $x$ de $F$, il existe un polynôme $P\in\R[X]$ tel
que $x=\big(P(u)\big)(x_{0})$.

\item
A l'aide des questions (\ref{q:P0(u)x0=0}) et (\ref{q:x=P(u)x0}), montrer
que $\forall x\in F, \exists R\in\R_{k}[X], x=\big(R(u)\big)(x_{0})$.

{\it(on pourra effectuer la division eulidienne de $P$ par $P_{0}$)}

\item
En déduire que $(x_{0}\dots x_{k})$ est une base de $F$.

\item
Ecrire la matrice de la restriction $u_{|_{F}}$ de $u$ à $F$ dans cette
base. Quel est le polynôme caractéristique de $\tilde{u}$~?

\item
Montrer qu'il existe une base $\mathcal{B}$ de $E$ dans la quelle
$$
 \mathrm{Mat}_{_{\mathcal{B}}}(u)=
 \begin{pmatrix}
   C_{1} &0    &\cdots&0\\
   0     &C_{2}&      &\vdots\\
   \vdots&     &\ddots&0     \\
   0     &\cdots&0    &C_{r}
 \end{pmatrix}
$$
où les matrices $C_{i}$ sont des matrices Compagnon.
\end{enumerate}

\finenonce{001702}



\finexercice



\exercice{3519, quercia, 2010/03/10}
\enonce{003519}{$f  \mapsto f(2x)$}

Soit $E = \{ f : \R \to  \R \text{ continues}$ tq $f(x)\to 0$ lorsque $x\to\pm\infty \}$,
$\varphi : \R \to \R,  x \mapsto {2x}$ et $u : E \to  E,  f \mapsto {f \circ \varphi.}$

Montrer que $u$ n'a pas de valeurs propres
(si $u(f) = \lambda f$, étudier les limites de $f$ en $0$ ou $\pm\infty$).
\finenonce{003519}



\finexercice


\exercice{3520, quercia, 2010/03/10}
\enonce{003520}{Ensi Physique P 94}

    Soit $E$ le sous-espace vectoriel de $\mathcal{C}(\R^+,\R)$ des fonctions ayant
    une limite finie en $+\infty$.
    Soit $T \in \mathcal{L}(E)$ défini par $T(f)(x) = f(x+1).$ Trouver les valeurs
    propres de~$T$.
\finenonce{003520}



\finexercice


\exercice{3521, quercia, 2010/03/10}
\enonce{003521}{\'Equation intégrale}

Soit $E = \mathcal{C}([0,+\infty[ \to\R)$
et $u : E \to  E, f \mapsto {\tilde f}$
avec $\tilde f(x) = \frac1x \int_{t=0}^x f(t)\,d t$.

\begin{enumerate}
  \item Montrer que $\tilde f$ peut être prolongée en une fonction continue sur
    $[0,+\infty[$.
  \item Chercher les valeurs propres et les vecteurs propres de $u$.
    
\end{enumerate}
\finenonce{003521}



\finexercice


\exercice{3522, quercia, 2010/03/10}
\enonce{003522}{Endomorphisme sur les suites}

Soit $E$ l'espace vectoriel des suites réelles $u=(u_n)_{n\ge 1}$
et $f$ l'endomorphisme de~$E$ défini par~:
$$(f(u))_n = \frac{u_1+2u_2+\dots+nu_n}{n^2}.$$
Quelles sont les valeurs propres de~$f$~?

\finenonce{003522}



\finexercice


\exercice{3523, quercia, 2010/03/10}
\enonce{003523}{Opérateur intégral}

Soit $E = \mathcal{C}([0,1],\R)$
et $f : E \to E,  u {\tilde u}$
avec $\tilde u(x) =  \int_{t=0}^1 \min(x,t)u(t)\,d t$.

Chercher les valeurs propres et les vecteurs propres de $f$.

\finenonce{003523}



\finexercice

\exercice{3579, quercia, 2010/03/10}
\enonce{003579}{Somme de projecteurs}

Soit $E$ un $ K$-ev de dimension finie et $u \in \mathcal{L}(E)$.
Montrer que $u$ est diagonalisable si et seulement s'il existe des projecteurs
$p_1,\dots,p_k \in \mathcal{L}(E)$ et des scalaires $\lambda_1,\dots,\lambda_k$
tels que :
$\begin{cases} u = \lambda_1p_1 + \dots + \lambda_kp_k\cr
         \forall\ i\ne j,\ p_i\circ p_j = 0.\cr \end{cases}$
\finenonce{003579}



\finexercice


\exercice{3580, quercia, 2010/03/10}
\enonce{003580}{$A^3$ est semblable à $A^4$}

Quelles sont les matrices $A \in \mathcal{M}_{3}(\C)$ telles que $A^3$ est semblable à $A^4$ ?
On étudiera séparément les cas :

\begin{enumerate}
  \item $A$ a 3 valeurs propres distinctes.
    
  \item $A$ a 2 valeurs propres distinctes
    
  \item $A$ a une seule valeur propre.
    
\end{enumerate}
\finenonce{003580}



\finexercice

\exercice{3584, quercia, 2010/03/10}
\enonce{003584}{$A$ et $2A$ sont semblables}

Soit~$A\in\mathcal{M}_n(\C)$ nilpotente. Montrer que~$A$ et $2A$ sont semblables.


\finenonce{003584}



\finexercice


\exercice{3597, quercia, 2010/03/10}
\enonce{003597}{Matrice bloc}

Soit $A \in \mathcal{M}_n(K)$ et $M = \begin{pmatrix} A &0 \cr A &A \cr\end{pmatrix} \in \mathcal{M}_{2N}(K)$.

\begin{enumerate}
  \item Comparer les valeurs propres de $A$ et $M$.
  \item Soit $P \in  K[X]$ et $Q = XP'$.
    Montrer que $P(M) = \begin{pmatrix} P(A) & 0 \cr Q(A) & P(A) \cr \end{pmatrix}$.
    
  \item A quelle condition sur $A$, $M$ est-elle diagonalisable ?
\end{enumerate}
\finenonce{003597}



\finexercice


\exercice{3598, quercia, 2010/03/10}
\enonce{003598}{Ensi P 90}

    Soit $M \in \mathcal{M}_n(\C)$ diagonalisable.
    Soit $A = \begin{pmatrix}M&M\cr M&M\cr\end{pmatrix} \in \mathcal{M}_{2n}(\C)$.
    La matrice $A$ est-elle diagonalisable ?
    
\finenonce{003598}



\finexercice


\exercice{3599, quercia, 2010/03/10}
\enonce{003599}{Matrice bloc}

Soit $A \in GL_n(\C)$ et $M = \begin{pmatrix} 0 &A \cr I &0 \cr \end{pmatrix} \in \mathcal{M}_{2n}(\C)$.
Montrer que $M$ est diagonalisable si et seulement si $A$ l'est.
(Chercher les sous-espaces propres de $M$ en fonction de ceux de $A$)

\finenonce{003599}



\finexercice


\exercice{3600, quercia, 2010/03/10}
\enonce{003600}{Matrice bloc}

Soit $M = \begin{pmatrix} A &B \cr C &D \end{pmatrix} \in \mathcal{M}_n(K)$ diagonalisable
avec $A$ carrée d'ordre $p$.

Soit $\lambda$ une valeur propre de $M$ de multiplicité $m$.
Montrer que si $p > n-m$, alors $\lambda$ est valeur propre de $A$.
\finenonce{003600}



\finexercice


\exercice{3601, quercia, 2010/03/10}
\enonce{003601}{Réduction par blocs (Centrale MP 2003)}
Soit~$A\in\mathcal{M}_n(\R)$ et $B = \begin{pmatrix}0&A\cr A&2A\cr\end{pmatrix}\in\mathcal{M}_{2n}(\R)$.
Déterminer $\mathrm{Spec}(B)$ et fonction de $\mathrm{Spec}(A)$.

\finenonce{003601}



\finexercice


\exercice{3602, quercia, 2010/03/10}
\enonce{003602}{ $A^m\to 0$ lorsque $m\to\infty$ (Mines MP 2003)}
\smallskip
Soit $A=\begin{pmatrix}a^2&ab&ab&b^2\cr ab&a^2&b^2&ab\cr ab&b^2&a^2&ab\cr b^2&ab&ab&a^2\cr\end{pmatrix}$.
Représenter dans un plan l'ensemble des couples~$(a,b)$ tels que $A^m \to 0$ lorsque $m\to\infty-$.

\finenonce{003602}



\finexercice


\exercice{3603, quercia, 2010/03/10}
\enonce{003603}{Chimie P 1996}

Soit $E$ un espace vectoriel réel de dimension~$n$ et $f$ un endomorphisme
de~$E$.

 Est-il vrai que~:
$f$ est diagonalisable $\Leftrightarrow$ $\mathrm{Ker} f + \Im f = E$~?
\finenonce{003603}



\finexercice


\exercice{3604, quercia, 2010/03/10}
\enonce{003604}{$u$ diagonalisable $ \Rightarrow  Ker(u-\lambda\mathrm{id}) + Im(u-\lambda\mathrm{id})$ est directe}

Soit $E$ un $ K$-ev de dimension finie et $u \in \mathcal{L}(E)$ diagonalisable.
Pour $\lambda \in \text{spec}(u)$, on note $E_\lambda = \mathrm{Ker}(u-\lambda\mathrm{id})$
et $F_\lambda = \Im(u-\lambda\mathrm{id})$.
Montrer que $E_\lambda \oplus F_\lambda = E$.
\finenonce{003604}



\finexercice


\exercice{3605, quercia, 2010/03/10}
\enonce{003605}{$\mathrm{Ker} f \oplus \Im f$}

Soit $E$ un $ K$-ev de dimension finie et $f \in \mathcal{L}(E)$.
On suppose qu'il existe $P\in K[X]$ tel que $P(f) = 0$ et $P'(0)\ne 0$.
Montrer que $\mathrm{Ker} f \oplus \Im f = E$.

\finenonce{003605}



\finexercice


\exercice{3606, quercia, 2010/03/10}
\enonce{003606}{$rg(f-\lambda\mathrm{id})$}

Soit $E$ un $\C$-ev de dimension finie et $f\in\mathcal{L}(E)$.
Montrer que $f$ est diagonalisable si et seulement si
pour tout $\lambda\in\C$ on a $rg(f-\lambda\mathrm{id}) = rg(f-\lambda\mathrm{id})^2$.
\finenonce{003606}



\finexercice


\exercice{3607, quercia, 2010/03/10}
\enonce{003607}{Nombre de noyaux et d'images}

Soit $E$ un $ K$-ev de dimension finie et $u\in \mathcal{L}(E)$.
Montrer que les ensembles ${\cal K} = \{\mathrm{Ker}(P(u)),\ P\in K[X]\}$ et
${\cal I} = \{\Im(P(u)),\ P\in K[X]\}$ sont finis et ont même cardinal.
\finenonce{003607}



\finexercice


\exercice{3608, quercia, 2010/03/10}
\enonce{003608}{$\dim(\mathrm{Ker} f^2) = 2\dim(\mathrm{Ker} f)$, Mines-Ponts MP 2005}

Soit $E$ un espace vectoriel de dimension finie et $f\in\mathcal{L}(E)$ tel que
$\dim(\mathrm{Ker} f^2) = 2\dim(\mathrm{Ker} f) = 2d$. Montrer que s'il existe $g\in\mathcal{L}(E)$
et $k\in\N^*$ tels que $g^k = f$ alors $k$ divise $d$.


\finenonce{003608}



\finexercice

\exercice{5661, rouget, 2010/10/16}
\enonce{005661}{***}
Soit $E =SL_2(\Zz) =\{\text{matrices carrées de format}\;2\;\text{à coefficients dans}\;\Zz\;\text{et de déterminant}\;1\}$.

\begin{enumerate}
 \item  Montrer que $(E,\times)$ est un groupe

\item  Soit $A$ un élément de $E$ tel que $\exists p\in\Nn^*/\;A^p = I_2$. Montrer que $A^{12} = I_2$.
\end{enumerate}
\finenonce{005661}


\finexercice


\exercice{5662, rouget, 2010/10/16}
\enonce{005662}{****}
Montrer que toute matrice de trace nulle est semblable à une matrice de diagonale nulle.
\finenonce{005662}


\finexercice


\exercice{5681, rouget, 2010/10/16}
\enonce{005681}{****}
Trouver $A$ dans $\mathcal{M}_n(\Rr)$ telle que la comatrice de $A$ soit $\left(
\begin{array}{cccc}
1&0&\ldots&0\\
2&\vdots& &\vdots\\
\vdots&\vdots& &\vdots\\
n&0&\ldots&0
\end{array}
\right)$.
\finenonce{005681}


\finexercice


\exercice{5687, rouget, 2010/10/16}
\enonce{005687}{**I}
Soit $A$ une matrice carrée réelle de format $n\geqslant 2$ vérifiant $A^3+A^2+A=0$. Montrer que le rang de $A$ est un entier pair.

\finenonce{005687}


\finexercice

\section{202.01 Endomorphisme du plan}


\exercice{1517, barraud, 2003/09/01}
\enonce{001517}{}
Dessiner l'allure du Shadock ci dessous après qu'il ait subi l'action de
l'endomorphisme de $\R^{2}$ dont la matrice dans la base canonique est

$$
A=
\begin{pmatrix}
 1/2 & 0\\
  0 & 2  
\end{pmatrix}
\quad
B=
\begin{pmatrix}
  1 & 1/2\\
  0 & 1  
\end{pmatrix}
\quad
C=
\begin{pmatrix}
  0 & 1\\
  1 & 0  
\end{pmatrix}
\quad
D=
\begin{pmatrix}
  1/2 &-\sqrt{3}/2\\
  \sqrt{3}/2 & 1/2  
\end{pmatrix}
\quad
E=
\begin{pmatrix}
   1   &-1\\
  -1/2 & 3/2  
\end{pmatrix}
$$

Ecrire la matrice de la dernière transformation dans la base
$((2,1),(-1,1))$.

$$
\includegraphics{../images/img001517-1}
$$
\finenonce{001517}



\finexercice



\exercice{1518, barraud, 2003/09/01}
\enonce{001518}{}
Retrouver la matrice (dans la base indiquée sur le premier dessin) de la
transformation subie par chacun des Shadocks ci-dessous.
$$
\includegraphics[height=8cm, keepaspectratio]{../images/img001518-1}
$$
\finenonce{001518}



\finexercice



\exercice{2564, delaunay, 2009/05/19}
\enonce{002564}{}
Soit $E$ un espace vectoriel sur un corps $K$ ($K=\R$ ou $\C$), on appelle {\it projecteur}
un endomorphisme $p$ de $E$ v\'erifiant $p\circ p=p$. Soit $p$ un projecteur.
\begin{enumerate}
  \item Montrer que $\mathrm{\mathrm{Id}}_E-p$ est un projecteur, calculer $p\circ(\mathrm{Id}_E-p)$ et $(\mathrm{Id}_E-p)\circ p$. 
  \item Montrer que pour tout $\vec x\in \Im p$, on a $p(\vec x)=\vec x$.
  \item En d\'eduire que $\Im p$ et $\ker p$ sont suppl\'ementaires.
  \item Montrer que le rang de $p$ est \'egal \`a la trace de $p$. (On rappelle que la trace de la matrice
 d'un endomorphisme ne d\'epend pas de la base dans laquelle on exprime cette matrice.)
\end{enumerate}
\finenonce{002564} 



\finexercice
\section{202.02 Endomorphisme auto-adjoint}


\exercice{1525, legall, 1998/09/01}

\enonce{001525}{}
Soit $  (E , \langle   ,  \rangle )  $ un espace euclidien et $  p \in
\mathcal{L} (E)  $ un projecteur. Montrer que $  p   $ est orthogonal si 
et seulement si $  p=p^*  .$
\finenonce{001525}



\finexercice



\exercice{1526, legall, 1998/09/01}

\enonce{001526}{}
Soit $  (E , \langle   ,   \rangle )  $ un espace euclidien et
$ \varphi \in \mathcal{L} (E)  .$ Soit $  F  $ un sous-espace vectoriel de $  E  .$
\begin{enumerate}
    \item Soit $  F  $ un sous-espace vectoriel de $  E  .$ Montrer que si $  \varphi = \varphi^*  $
et $  \varphi (F)\subset F  $ alors $  \varphi (F^\perp )\subset F^\perp   .$
    \item Soit $  F  $ un espace propre de $  \varphi   .$ Montrer que si $  \varphi \circ \varphi^*
=\varphi ^*\circ \varphi  $ alors $  \varphi (F^\perp )\subset F^\perp    .$
\end{enumerate}
\finenonce{001526}



\finexercice



\exercice{1527, legall, 1998/09/01}

\enonce{001527}{}
Soient $  A  $ et $  B  $ deux matrices sym\'etriques positives.
Soit $  k \in \N ^* .$
\begin{enumerate}
    \item Montrer que tout vecteur propre de $  A^k  $ est vecteur propre de $  A  .$ 
    \item Si $  A^k = B^k   $ alors $  A = B  .$
    \item  Que se passe-t-il sans l'hypoth\`ese $  A  $ et $  B  $ sym\'etriques positives ?
\end{enumerate}
\finenonce{001527}



\finexercice



\exercice{1528, legall, 1998/09/01}

\enonce{001528}{}
Soit $  (E , \langle   ,   \rangle )  $ un espace euclidien et
$ \varphi \in \mathcal{L} (E)  .$ 
\begin{enumerate}
    \item Montrer que $  \varphi ^*\circ \varphi   $ est sym\'etrique et que $  \hbox{Sp}(
\varphi ^*\circ \varphi )\subset \R _+   . $
    \item On note respectivement $  \lambda   $ et $  \mu   $ la plus grande et la plus petite valeur propre de $  \varphi ^*\circ \varphi   .$ Montrer,
pour tout $  x \in E  ,$ l'in\'egalit\'e~: $$\mu \Vert x\Vert ^2\leq \Vert \varphi (x)\Vert ^2\leq
\lambda \Vert x\Vert ^2   .$$
\end{enumerate}
\finenonce{001528}



\finexercice



\exercice{1529, legall, 1998/09/01}

\enonce{001529}{}
Soit $  (E , \langle   ,   \rangle )  $ un espace euclidien et
$ \varphi \in \mathcal{L} (E)  .$ 
\begin{enumerate}
    \item Montrer que si $  \varphi =\varphi ^*  $ et $  \forall x \in E  :  \langle x , \varphi (x) \rangle =0  $ alors $  \varphi =0  .$
    \item Montrer que les propri\'et\'es suivantes sont \'equivalentes~:
\vskip1mm \hskip2mm {\em i)} $  \varphi \circ \varphi^*
=\varphi ^*\circ \varphi  .$
\vskip1mm \hskip2mm {\em ii)} $  \forall x,y\in E  :  \langle \varphi (x),\varphi (x) \rangle =
\langle \varphi ^*(x),
\varphi ^*(x)\rangle   .$
\vskip1mm \hskip2mm {\em iii)} $  \forall x \in E  :   \Vert \varphi (x) \Vert = \Vert \varphi
^* (x)\Vert   .$
    \item Si $  \hbox{dim}(E)=2  $ et si $  \varphi \circ \varphi^*
=\varphi ^*\circ \varphi  $ alors la matrice de $  \varphi   $ dans une base orthonorm\'ee est soit sym\'etrique, soit
de la forme $  \begin{pmatrix} a & -b   \cr  b & a  \cr \end{pmatrix}  $ avec $  b \not =0  .$
    \item On suppose d\'esormais que $  \hbox{dim}(E)=3  $ et que $  \varphi \circ \varphi^* =\varphi ^*\circ \varphi  .$ 
    \begin{enumerate}
        \item Montrer que $  \varphi   $ a au moins une valeur propre r\'eelle
qu'on notera $  \lambda   .$ Montrer que $  E_{\lambda }  $ et $  E_{\lambda }^\perp   $
sont laiss\'es stables par $  \varphi   $ et $  \varphi ^*  .$
        \item Montrer que si $  \varphi   $ n'est pas sym\'etrique, il existe une base 
orthonorm\'ee $  \epsilon  $ de $  E  $ et deux r\'eels $  a   $ et $  b   $ (avec $ 
b\not =0  $) tels que $  \hbox{Mat}(\varphi ,\epsilon ) =\begin{pmatrix} a & -b & 0  \cr  b & a & 0 \cr 0 & 0 & \lambda \cr \end{pmatrix}  .$
    \end{enumerate}
\end{enumerate}
\finenonce{001529}



\finexercice



\exercice{1530, legall, 1998/09/01}

\enonce{001530}{}
Soit $  E  $ un espace euclidien de dimension $  3  .$ 
\begin{enumerate}
    \item Soit $  \{ e_1, e_2 ,e_3 \}   $ une base orthonorm\'ee de $  E  .$ Soient $ 
x=x_1e_1+x_2e_2+x_3e_3  $ et $ 
y=y_1e_1+y_2e_2+y_3e_3  $ deux vecteurs de $  E  .$ Calculer $  \langle x,y\rangle   $ en
fonction des coefficients $  x_i  $ et $  y_i  $ (pour $  i\in \{ 1,2,3\}  $).
    \item On consid\`ere $  u \in \mathcal{L} (E)  $ un endomorphisme auto-adjoint. On note $  \lambda   $ sa plus petite valeur propre et $  \lambda '  $ sa
plus grande valeur propre. Montrer que l'on a, pour tout $  x   $ appartenant \`a $  E  ,$ les
in\'egalit\'es~:
$$ \lambda  \Vert x\Vert ^2 \leq \langle u(x),x\rangle \leq \lambda ' \Vert x\Vert ^2.$$
(On utilisera une base orthonorm\'ee convenable.)
    \item Soit $  v \in \mathcal{L} (E)  $ un endomorphisme quelconque. Montrer que $  \displaystyle {u=
\frac{1}{ 2} (v+v^*)}  $ est auto-adjoint. Soient $  \mu   $ une valeur propre de $  v  ,$
$  \lambda   $ la
plus petite valeur propre de $  u  $ et $  \lambda '  $ la
plus grande valeur propre de $  u  .$ Montrer que $  \lambda  \leq \mu \leq \lambda '  .$
\end{enumerate}
\finenonce{001530}



\finexercice



\exercice{1531, legall, 1998/09/01}

\enonce{001531}{}
\begin{enumerate}
    \item Soit $A =(a_{ij}) \in \textbf{M}_n(\R)$. Montrer que $S =
{}^t\! A\cdot A$ est une matrice sym\'etrique dont tous les valeurs
propres $\lambda_1, \dots, \lambda_n$ sont positives. D\'emontrer
l'\'egalit\'e : $\sum_{i=1}^n \lambda_i = \sum_{1\leq i,j\leq n}
a^2_{ij}$.
    \item Soit $S \in \textbf{M}_n(\R)$ une matrice sym\'etrique. Existe-t-il
une matrice $A \in \textbf{M}_n(\R)$ telle que $S = {}^t\!
A\cdot A$ ? Donner une condition n\'ecessaire et suffisante sur $S$ pour
que $A$ soit inversible. Application \`a $S = \left(\begin{array}{cc}2&1\\
1&2\end{array}\right)$.
\end{enumerate}
\finenonce{001531}



\finexercice



\exercice{1532, legall, 1998/09/01}

\enonce{001532}{}
Soit $(E,<,>)$ un espace euclidien de dimension $p$. A chaque $n$-uple
$(x_1, \dots, x_n)$ d'\'el\'ements de $E$ on associe le nombre
(d\'eterminant de Gram)
$$G(x_1, \dots, x_n) = \textrm{d\'et}(<x_i,x_j>)_{i,j=1,\dots,n}.$$
\begin{enumerate}
    \item Montrer que $x_1, \dots, x_n$ sont li\'es si et seulement si $G(x_1,
\dots, x_n) = 0$ ;  montrer que si $x_1, \dots, x_n$ sont ind\'ependants,
on a $G(x_1, \dots, x_n) > 0$.
    \item Montrer que, pour toute permutation $\sigma$ de $\{1, \dots, n\}$, on a
$G(x_{\sigma(1)}, \dots, x_{\sigma(n)}) = G(x_1, \dots, x_n)$, et que la
valeur de $G(x_1, \dots, x_n)$ n'est pas modifi\'ee si l'on rajoute \`a un
des vecteurs, soit $x_i$, une combinaison lin\'eaire des autres vecteurs
$x_j (j\not= i)$. Calculer $G({\alpha}x_1, \dots, x_n)$ (${\alpha} \in 
\R$). 
    \item On suppose $x_1, \dots, x_n$ ind\'ependants. Soit $x \in E$, et soit
$d(x,H)$ la distance de $x$ \`a l'hyperplan $H = \textrm{Vect}(x_1, \dots,
x_n)$. Montrer que
$\displaystyle{d(x,H)^2 = \frac{G(x, x_1, \dots, x_n)}{G(x_1, \dots, x_n)}}$.
\end{enumerate}
\finenonce{001532}



\finexercice



\exercice{1533, ortiz, 1999/04/01}

\enonce{001533}{}
Diagonaliser tr\`es rapidement la matrice
$${M=\left(\begin{smallmatrix}
0&1&1\\1&0&1\\1&1&0\end{smallmatrix}\right)\in\mathcal{M}_3(\Rr).}$$
\finenonce{001533}



\finexercice



\exercice{1534, ortiz, 1999/04/01}

\enonce{001534}{}
Montrer que l'endomorphisme de l'espace vectoriel
euclidien canonique $\Rr^3$ de matrice dans la
base canonique de $\Rr^3$
$${ C=-\frac 19\left(\begin{smallmatrix}
4&1&-8\\ 7&4&4\\ 4&-8&1
\end{smallmatrix}\right)
}$$ est un automorphisme orthogonal.
\finenonce{001534}



\finexercice



\exercice{1535, ortiz, 1999/04/01}

\enonce{001535}{}
Soient $E$ un espace vectoriel euclidien et $f$ un
endomorphisme de $E$ tel que
$${\forall x\in E,
{\|f(x)\|}\le {\|x\|}.}$$
\begin{enumerate}
    \item
    \begin{enumerate}
         \item Soit $x\in E$ tel que $f^{\ast}(x)=x$. Montrer que
${\|f(x)-x\|}^2={\|f(x)\|}^2-{\|x\|}^2.$
         \item En d\'eduire que $\ker (f^{\ast}-\text{Id})\subseteq \ker (f-\text{Id}).$
    \end{enumerate}
    \item  Soit $h$ un endomorphisme de $E.$ Montrer que $(\text{Im\,} h)^{\perp}\subseteq \ker h^{\ast}.$
    \item  En d\'eduire que les sous-espace vectoriels $\ker (f-\text{Id})$ et $\text{Im\,} (f-\text{Id})$ sont suppl\'ementaires et orthogonaux.
\end{enumerate}
\finenonce{001535}



\finexercice



\exercice{1536, ortiz, 1999/04/01}

\enonce{001536}{}
Soit $\, E\, $ un espace euclidien de dimension
$\, 3\, .$
\begin{enumerate}
    \item  Soit $\, ( e_1, e_2 ,e_3 ) \, $ une base orthonorm\'ee de $\, E\, .$
Soient $\,
x=x_1e_1+x_2e_2+x_3e_3\, $ et $\,
y=y_1e_1+y_2e_2+y_3e_3\, $ deux vecteurs de $\, E\, .$ Calculer $\, \langle
x,y\rangle \, $ en
fonction des coefficients $\, x_i\, $ et $\, y_i\, $ (pour $\, i\in \{
1,2,3\}\, $).

    \item  On consid\`ere $\, u \in \mathcal{L} (E)\, $ un endomorphisme auto-adjoint. On
note $\, \lambda _1\, $ sa
plus petite valeur propre et $\, \lambda _2\, $ sa
plus grande valeur propre. Montrer que l'on a, pour tout $\, x \, $
appartenant \`a $\, E\, $ les
in\'egalit\'es~:
$$ \lambda _1 \Vert x\Vert ^2 \leq \langle u(x),x\rangle \leq \lambda _2
\Vert x\Vert ^2.$$
(On utilisera une base orthonorm\'ee convenable.)

    \item  Soit $\, v \in \mathcal{L} (E)\, $ un endomorphisme quelconque. Montrer que $\,
\displaystyle {u=
{\frac{1}{2}} (v+v^*)}\, $ est auto-adjoint.
Soient $\, \lambda \, $ une valeur propre de $\,
v\, ,$ $\, \lambda _1\, $ la plus petite valeur
propre de $\, u\, $ et $\, \lambda
_2\, $ la plus grande valeur propre de $\, u\, .$
Montrer que $\,
\lambda _1 \leq
\lambda \leq \lambda _2\, .$
\end{enumerate}
\finenonce{001536}



\finexercice



\exercice{1537, ortiz, 1999/04/01}

\enonce{001537}{}
\begin{enumerate}
    \item  Soient $E$ un espace vectoriel euclidien, $f\in\mathcal{L}(E)$
 un endomorphisme sym\'etrique positif. Montrer que si $x\in E$ alors $(f(x)|x)\geq 0.$
    \item   Soit $M=\left( m_{i,j}\right) _{i,j}\in \mathcal{M}_n(\Rr)$
sym\'etrique positive. Montrer que pour tout $i=1,..,n,\;m_{ii}\geq 0$ et tr$%
(M)\geq 0$
    \item  Soient $A,\;B\in \mathcal{M}_n(\Rr)$ sym\'etriques positives.
    \begin{enumerate}
        \item Montrer qu'il existe $D\in \mathcal{M}_n(\Rr)$ diagonale et $%
M\in \mathcal{M}_n(\Rr)$ sym\'etrique positive
telle que tr$(AB)=$tr$(DM).$
         \item En d\'eduire que tr$(AB)\leq $tr$(A)$tr$(B).$
    \end{enumerate}
\end{enumerate}
\finenonce{001537}



\finexercice



\exercice{1538, legall, 2003/10/01}

\enonce{001538}{}
 Soit $ (E, 
\langle   ,  \rangle )  $ un espace euclidien et $f \in \mathcal{L} (E)
$ un endomorphisme autoadjoint.
Montrer que
les trois propri\'et\'es suivantes sont \'equivalentes~:

\begin{enumerate}
\item $\forall x \in E ,\langle f(x), x\rangle \geq 0$.
\item Il existe $g\in \mathcal{L} (E)$ tel que $f=g^*g.$
\item Il existe $h\in \mathcal{L} (E)$ tel que $h=h^*$ et $f=h^2.$
\end{enumerate}

\finenonce{001538}



\finexercice



\exercice{1539, legall, 2003/10/01}

\enonce{001539}{}
Soit $ (E, \langle   ,  \rangle )  $ un espace euclidien et $
f \in \mathcal{L} (E)
$ un endomorphisme. Montrer que $\Vert f\Vert =\Vert f^*\Vert $.

\finenonce{001539}



\finexercice



\exercice{1540, legall, 2003/10/01}

\enonce{001540}{}
Soit $ (E, 
\langle   ,  \rangle )  $ un espace euclidien (de dimension finie) et 
$
f \in \mathcal{L} (E)
$ un endomorphisme autoadjoint. On note $X=\{ x\in E ; \langle f(x),x 
\rangle \leq 1\} .$
Montrer que $X$ est compacte si et seulement si toutes les valeurs 
propres de $f$ sont strictement positives.

\finenonce{001540}



\finexercice


\section{202.03 Autres endomorphismes normaux}


\exercice{1541, legall, 1998/09/01}

\enonce{001541}{}
Soit $  (E , \langle   ,   \rangle )  $ un espace euclidien. Un
endomorphisme $  \varphi \in \mathcal{L} (E)  $ est dit antisym\'etrique lorsque
 $  \varphi ^*=-\varphi   .$
\begin{enumerate}
    \item Montrer que $  \varphi  $ est antisym\'etrique si et seulement si $  \forall x \in E  ,  
\langle \varphi (x),x \rangle =0  .$ (on pourra remarquer que $  \varphi +\varphi ^*   $ est
autoadjoint.) 
    \item Montrer que si $  \varphi  $ est antisym\'etrique alors
 $  (\hbox{Ker}(\varphi))^\perp =\hbox{Im}
(\varphi )  $ puis que $  \hbox{rg}(\varphi )  $ est pair.
\end{enumerate}
\finenonce{001541}



\finexercice



\exercice{1542, legall, 1998/09/01}

\enonce{001542}{}
Soit $  (E , \langle   ,   \rangle )  $ un espace euclidien. Soit
 $  \varphi \in \mathcal{L} (E)  $ un endomorphisme antisym\'etrique c'est-\`a-dire tel que $  \varphi
^*=-\varphi   .$ 
\begin{enumerate}
    \item Montrer que si $  \lambda \in Sp(\varphi )  $ alors $  \lambda =0  .$ Montrer que
$  (\hbox{Ker}(\varphi ))^\perp   $ est stable par $  \varphi   .$
    \item 
    \begin{enumerate}
        \item  Montrer que $  \varphi ^2  $ est sym\'etrique. 
        \item Montrer que si $  x  $ est un vecteur propre associ\'e \`a une valeur propre $  \mu   $ de $ 
\varphi ^2  $ alors $  E_x= \hbox{vect}\{ x, \varphi (x)\}  $ et $  E_x ^\perp   $ sont laiss\'es
stables par $  \varphi  .$ 
        \item Montrer que $  \mu >0  .$ D\'eterminer une base $  \{ e_1,e_2\}   $ de $  E_x  $ telle que
la matrice de la restriction de $  \varphi   $ ˆ $  E_x  $ dans $  \{ e_1,e_2\}   $ soit
$  \begin{pmatrix} 0 & -\sqrt{\mu } \cr \sqrt{\mu } & 0 \cr \end{pmatrix}. $  
    \end{enumerate}
    \item Montrer que $  E  $ est somme directe orthogonale de $  \hbox{Ker} (\varphi )  $ et de plans stables.
\end{enumerate}
\finenonce{001542}



\finexercice


\section{202.04 Endomorphisme orthogonal}


\exercice{1543, barraud, 2003/09/01}

\enonce{001543}{}
Soit $f$ une transformation orthogonal d'un espace euclidien $E$. Montrer que
$$
 \mathrm{Ker}(f-\mathrm{id})=\mathrm{Im}(f-\mathrm{id})^{\bot}
$$
En déduire que si $(f-\mathrm{id})^{2}=0$, alors $f=\mathrm{id}$.
\finenonce{001543}



\finexercice



\exercice{1544, barraud, 2003/09/01}

\enonce{001544}{}
Déterminer la nature des transformations de $\R^{3}$ dont les matrices
dans la base canonique sont les suivantes :
$$
 A=\frac{1}{3}
  \begin{pmatrix}
    1 & -2 & -2  \\
   -2 &  1 & -2  \\
    2 &  2 & -1
  \end{pmatrix}
 \qquad
 B=\frac{1}{3}
  \begin{pmatrix}
    2 &  2 & -1  \\
   -1 &  2 &  2  \\
    2 & -1 &  2
  \end{pmatrix}
 \qquad
 C=
  \begin{pmatrix}
    0 &  1 &  0  \\
    0 &  0 & -1  \\
   -1 &  0 &  0
  \end{pmatrix}
$$
\finenonce{001544}



\finexercice



\exercice{1545, barraud, 2003/09/01}

\enonce{001545}{}
Diagonaliser dans une base orthonormale (pour le produit scalaire canonique de $\R^{3}$) la
matrice suivante :
$$
 A=
 \begin{pmatrix}
   5 & -1 &  2 \\
  -1 &  5 &  2 \\
   2 &  2 &  2
 \end{pmatrix}
$$
Interpréter géométriquement la transformation de $\R^{3}$ représentée par cette matrice.
\finenonce{001545}



\finexercice



\exercice{1546, barraud, 2003/09/01}

\enonce{001546}{}
Diagonaliser les matrices suivantes dans des bases orthonormées~:
$$
A=
\begin{pmatrix}
  4&i&-i\\
  -i&4&1\\
  i &1&4
\end{pmatrix}
\qquad
B=
\begin{pmatrix}
  1&i\sqrt{2}&0\\
  -i\sqrt{2}&1&-i\sqrt{2}\\
  0 &i\sqrt{2}&1
\end{pmatrix}
\qquad
C=
\begin{pmatrix}
  -1&  0&-3i&  0\\
   0&  1&  0& 3i\\
  3i&  0& -1&  0\\
   0&-3i&  0&  1\\
\end{pmatrix}
$$
\finenonce{001546}



\finexercice



\exercice{1547, barraud, 2003/09/01}

\enonce{001547}{}
Soit $A=(a_{ij})_{\substack{1\leq i\leq n\\1\leq j\leq n}}$ une matrice symétrique réelle.
Montrer que ses valeurs propres $\lambda_{1},\ldots,\lambda_{n}$ vérifient
$$
  \sum_{i=1}^{n}{\lambda_{i}}=\sum_{i=1}^{n}{a_{i,j}^{2}}
$$
\finenonce{001547}



\finexercice



\exercice{1548, barraud, 2003/09/01}

\enonce{001548}{}
Soit $\mathcal{B}=(e_{1},\ldots,e_{n})$ une base orthogonal d'un espace euclidien $E$. On
dit qu'un endomorphisme $f$ de $E$ conserve l'orthogonalité de $\mathcal{B}$ si et
seulement si $(f(e_{1}),\ldots,f(e_{n}))$ est une famille orthogonale.

Montrer que $f$ conserve l'orthogonalité de $\mathcal{B}$ si et seulement si $\mathcal{B}$
est une base de vecteurs propres de ${}^t{f}\,f$.

Montrer que pour tout endomorphisme $f$ de $E$, il existe une base orthogonale dont $f$
conserve l'orthogonalité.
\finenonce{001548}



\finexercice



\exercice{1549, barraud, 2003/09/01}

\enonce{001549}{Décomposition polaire}
\begin{enumerate}
\item
Soit $r$ un endomorphisme symétrique d'un espace euclidien $E$. On dit que $r$ est
positif, si toutes ses valeurs propres sont positives.

Montrer que si $r$ est défini positif, il existe un et un seul endomorphisme symétrique
$s$ positif tel que $s^{2}=r$. On appelle $s$ racine carrée positive de $r$.

On dit que $r$ est défini positif si et seulement si toutes ses racines sont strictement
positives. Montrer que si $r$ est défini positif, alors sa racine positive aussi.

\item
Soit $f$ un endomorphisme de $E$. Montrer que ${}^t{f}f$ est symétrique et positif.
Montrer que si en plus $f$ est bijective, ${}^t{f}f$ est défini positif.

\item
On suppose maintenant que $f$ est une bijection. Soit $s$ la racine carrée positive de
${}^t{f}f$. Montrer que $u=f\circ s^{-1}$ est une transformation orthogonale. En
déduire que tout endomorphisme bijectif de $E$ peut s mettre sous la forme :
$$
 f=u\circ s
$$
où $u$ et une transformation orthogonale, et $s$ est symétrique défini positif.

Montrer que cette décomposition, appelée décomposition polaire de $f$ est unique.

\item
Que se passe-t-il si $f$ n'est pas bijective ?
\end{enumerate}
\finenonce{001549}



\finexercice



\exercice{1550, barraud, 2003/09/01}

\enonce{001550}{}

Dans l'espace vectoriel $\R^{4}$ muni de son produit scalaire
canonique, on considère l'endomorphisme $f$ dont la matrice dans la
base canonique est :
$$
 A =\frac{1}{7} 
 \begin{pmatrix}
   -1 & -4 &  4 & -4 \\
   -4 &  5 &  2 & -2 \\
    4 &  2 &  5 &  2 \\
   -4 & -2 &  2 &  5 
 \end{pmatrix}
\rlap{\qquad {(\it attention au $\frac{1}{7}$...)}}
$$

\begin{enumerate}
\item 
Sans calculs, dire pourquoi $f$ est diagonalisable dans une base
orthonormée.

\item
Montrer que $f$ est orthogonal. En déduire les seules valeurs propres
possibles pour $f$.

\item
Sans calculer le polynôme caractéristique de $f$, déterminer à l'aide
de la trace l'ordre de multiplicité des valeurs propres de $f$. En
déduire le polynôme caractéristique de $f$.

\item
Déterminer l'espace propre $E_{1}$ associé à la valeur propre
1. En donner une base, puis lui appliquer le procédé de Schmidt pour 
obtenir une base orthonormée de $E_{1}.$

\item
Montrer que l'espace propre $E_{-1}$ associé à la valeur propre -1
satisfait $E_{-1}=(E_{1})^{\bot}$. En utilisant l'équation
caractérisant $E_{1}$, en déduire un vecteur générateur de $E_{-1}$.

\item
Donner une base orthonormée dans laquelle la matrice de $f$ est
diagonale. Donner une interprétation géométrique de $f$.

\end{enumerate}
\finenonce{001550}



\finexercice



\exercice{1551, barraud, 2003/09/01}

\enonce{001551}{}
\noindent\textbf{A --- } Soit $E$ un espace vectoriel et $u$ et $v$ deux
endomorphismes de
$E$ diagonalisables qui commutent (c'est à dire qui satisfont $u\circ
v=v\circ u$). On note $\lambda_{1}...\lambda_{k}$les valeurs propres de
$u$ et $E_{1}...E_{k}$ les espaces propres associés.
  \begin{enumerate}
  \item 
    Montrer que $v(E_{i})\subset E_{i}$.
  \item
    On note $v_{i}=v_{|_{E_{i}}}$ la restriction de $v$ à $E_{i}$. Soit
    $P\in\C[X]$, montrer que $P(v_{i})=P(v)_{|_{E_{i}}}$.
  \item 
    En déduire que $v_{i}$ est diagonalisable. Soit $B_{i}$ une base de
    $E_{i}$ formée de vecteurs propres de $v_{i}$. 

    Montrer que $B=\bigcup\limits_{i=1}^{k}B_{i}$ est une base de $E$
    formée de vecteurs propres à la fois pour $u$ et pour $v$.
 \item
    En déduire que $u$ et $v$ sont diagonalisables dans une même base.
    Montrer que $u-v$ est diagonalisable.
  \end{enumerate}

\medskip
\noindent
\textbf{B --- \it Application~:} On considère maintenant une matrice
$A\in M_{n,n}(\R)$, et on lui associe l'endomorphisme
$w_{A}\in\mathrm{End}(M_{n,n}(\R))$ suivant~:
$$ w:
\begin{array}{ccc}
\mathcal{M}_{n,n}(\R) &  \rightarrow & \mathcal{M}_{n,n}(\R)\\
M            &\mapsto  & AM-MA
\end{array}
$$
Le but de l'exercice est de montrer que si $A$ est diagonalisable, $w_{A}$
l'est aussi.

On note 
$$
u_{A}:
\begin{array}{ccc}
\mathcal{M}_{n,n}(\R) &  \rightarrow & \mathcal{M}_{n,n}(\R)\\
M            &\mapsto  & AM
\end{array}
\qquad\text{ et }\qquad
v_{A}:
\begin{array}{ccc}
\mathcal{M}_{n,n}(\R) &  \rightarrow & \mathcal{M}_{n,n}(\R)\\
M            &\mapsto  & MA
\end{array}
$$


\begin{enumerate}
\item 
  Montrer que $\forall k\in\N, (u_{A})^{k}=u_{A^{k}}$. En déduire que
  $\forall P\in\C[X], P(u_{A})=u_{P(A)}$, puis que tout polynôme
  annulateur de $A$ est un polynôme annulateur de $u_{A}$.

\item  
  Montrer que
  $$
  A \text{ diagonalisable }\Rightarrow u_{A}\text{ diagonalisable}
  $$
  
  On admet sans démonstration que le même résultat est vrai pour
  $v_{A}$~:
  $$
  A \text{ diagonalisable }\Rightarrow v_{A}\text{ diagonalisable}
  $$
  
\item
  Montrer que $u_{A}\circ v_{A}=v_{A}\circ u_{A}$.
  
\item
  En déduire que
  $$
  A \text{ diagonalisable }\Rightarrow w_{A}\text{ diagonalisable}
  $$
\end{enumerate}
\finenonce{001551}



\finexercice



\exercice{1552, barraud, 2003/09/01}

\enonce{001552}{}
  Dans un espace euclidien $(E,<\cdot ,\cdot>)$, on considère un vecteur
  $v$ non nul, un scalaire $\lambda$ et l'endomorphisme~:
$$
u:
\begin{array}{ccl}
E &  \rightarrow &E  \\
x &\mapsto  &x+\lambda <x,v> v
\end{array}
$$

\begin{enumerate}
\item 
  Pour $x\in E$, calculer $\Vert u(x)\Vert^{2}$.

\item  
  Donner une condition nécessaire et suffisante sur $\lambda$ et $v$
  pour que $u$ soit une transformation orthogonale.
  
\item
  Lorsque $u$ est orthogonale, dire a priori quelles sont les valeurs
  propres possibles de $u$, puis dire si elles sont effectivement valeur
  propre en étudiant les espaces propres associés.

\item
  Lorsque $u$ est orthogonale, donner une interprétation géométrique de
  $u$.
\end{enumerate}
\finenonce{001552}



\finexercice



\exercice{1553, barraud, 2003/09/01}

\enonce{001553}{}
On considère un espace euclidien $(E,<>)$. On dit qu'un endomorphisme
$u$ de $E$ est une similitude de $E$ si et seulement si il existe un
réel $\lambda>0$ tel que 
$$
 u^{*}u=\lambda\mathrm{id}
$$ 
Montrer que les trois assertions suivantes sont équivalentes :
\begin{enumerate}
\item[\it(i)] 
$u$ est une similitude
\item[\it(ii)]
$u$ est colinéaire à une transformation orthogonale, c'est à dire
$$
 \exists\alpha\in\R\setminus\{0\},\ \exists v\in O(E)\ /\ u=\alpha v
$$
\item[\it(iii)]
$u$ conserve l'orthogonalité, c'est à dire :
$$
 \forall(x,y)\in E^{2}, <x,y>=0\Rightarrow <u(x),u(y)>=0
$$
{\it
  Pour (i)$\Leftrightarrow$(ii), on pourra commencer par montrer que
  (ii)$\Rightarrow$(i).

  Pour (i)$\Rightarrow$(iii), on commencera par
  montrer que $x$ et $u^{*}u(x)$ sont toujours colinéaires, c'est à
  dire que 
  $$\forall x\in E\exists\lambda_{x}/u^{*}u(x)=\lambda_{x}x$$
  puis on montrera que $\lambda_{x}$ est indépendant de $x$.}
\end{enumerate}
\finenonce{001553}



\finexercice



\exercice{1554, barraud, 2003/09/01}

\enonce{001554}{}
  Dans un espace euclidien $E$, on considère un vecteur unitaire $a$, et
  à un réel $k\neq -1$ on associe l'endomorphisme $u_{k}$ de $E$ défini
  par~:
  $$
  u_{k}(x)=k<x,a>a + x
  $$
  
  \begin{enumerate}
  \item
    Montrer que $u_{k}$ est un isomorphisme. Déterminer $u_{k}^{-1}$.
    \emph{(on pourra commencer par calculer $<u_{k}(x),a>$)}

  \item 
    Rappeler la caractérisation de l'adjoint d'un endomorphisme, et 
    montrer que $u$ est auto adjoint.

  \item  
    Pour quelles valeurs de $k$ $u$ est-il orthogonal~? Interpréter alors
    géométriquement cette transformation.

  \item
    Déterminer les valeurs propres et vecteurs propres de $u_{k}$.
  \end{enumerate}

\finenonce{001554}



\finexercice



\exercice{1555, ortiz, 1999/04/01}

\enonce{001555}{}
\begin{enumerate}
    \item  Soient $E$ un espace vectoriel euclidien, $f$ un endomorphisme de $E$ et
$A=\left(a_{ij}\right)_{1\le i,j\le
n}\in\mathcal{M}_n(\Rr)$ la matrice de $f$ dans
une base orthonormale donn\'ee
$\mathcal{B}=(e_1,...,e_n)$ de $E$. Pour
$i,j\in\left\{1,...,n\right\}$, exprimer $a_{ij}$ en
fonction de $f$ et des vecteurs $e_i$ et $e_j$.
    \item  Soient $A=\left(a_{ij}\right)_{1\le i,j\le n}\in\mathcal{O}_n(\Rr)$ et
$S=\sum_{1\le i,j\le n}a_{ij}$.
    \begin{enumerate}
         \item Montrer qu'il existe $u\in E$ tel que $S=(u\vert f(u))$.
         \item En d\'eduire que $|S|\le n$.
    \end{enumerate}
\end{enumerate}
\finenonce{001555}



\finexercice



\exercice{1556, ortiz, 1999/04/01}

\enonce{001556}{}
Soient $A=\left(a_{ij}\right)_{1\le i,j\le
n}\in\mathcal{O}_n(\Rr)$ et $A_{ij}$ le cofacteur
$(i,j)$ de $A$. Montrer que $\det A>0$ si et
seulement si $a_{ij}$ et $A_{ij}$ sont de m\^eme
signe.
\finenonce{001556}



\finexercice



\exercice{1557, ortiz, 1999/04/01}

\enonce{001557}{}
Que peut-on dire d'une matrice carr\'ee r\'eelle
\`a la fois sym\'etrique et orthogonale ?
D\'eterminer la nature et les \'el\'ements
caract\'eristiques de l'endomorphisme de l'espace
vectoriel euclidien canonique $\Rr^3$ de matrice $
A=\frac {1}{7}\left(\begin{smallmatrix}
-2&6&-3\\
6&3&2\\
-3&2&6
\end{smallmatrix}\right)$ dans la base canonique de $\Rr^3.$
\finenonce{001557}



\finexercice



\exercice{1558, ortiz, 1999/04/01}

\enonce{001558}{}
Quelles sont les isom\'etries vectorielles d'un
espace vectoriel euclidien qui sont
diagonalisables.
\finenonce{001558}



\finexercice



\exercice{1559, ortiz, 1999/04/01}

\enonce{001559}{}
Soient $E$ un espace vectoriel euclidien et $f$ un
endomorphisme de $E$ tel que
\begin{center}$\forall x\in E, \|{f(x)}\|\le \|{x}\|.$\end{center}
\begin{enumerate}
    \item
    \begin{enumerate}
         \item Soit $x\in E$ tel que $f^{\ast}(x)=x$. Montrer que
$\|{f(x)-x}\|^2=\|{f(x)}\|^2-\|{x}\|^2.$
         \item En d\'eduire que $\ker (f^{\ast}-\text{Id})\subseteq \ker (f-\text{Id}).$
    \end{enumerate}
    \item  Soit $h$ un endomorphisme de $E.$ Montrer que
$(\text{Im\,} h)^{\perp}\subseteq \ker h^{\ast}.$
    \item  En d\'eduire que les sous-espace vectoriels
$\ker (f-\text{Id})$ et $\text{Im\,}
(f-\text{Id})$ sont suppl\'ementaires et
orthogonaux.
\end{enumerate}
\finenonce{001559}



\finexercice



\exercice{1560, ortiz, 1999/04/01}

\enonce{001560}{}
D\'eterminer une matrice diagonale
$D\in\mathcal{M}_2(\Rr)$ et une matrice
orthogonale $U\in\mathcal{O}_2(\Rr)$ telles que
$UDU^{-1}=\left(
\begin{smallmatrix}
1&\frac 12 \\
\frac 12&1
\end{smallmatrix}\right).$
\finenonce{001560}



\finexercice



\exercice{1561, legall, 1998/09/01}

\enonce{001561}{}
Soit $  (E , \langle   ,   \rangle )  $ un espace euclidien et $  u \in
\mathcal{L} (E)  .$ Montrer que les propri\'et\'es suivantes sont \'equivalentes~:
\vskip1mm \hskip1mm {\em i)} $  u^*=u^{-1}  .$
\vskip1mm \hskip1mm {\em ii)} $  \forall x \in E  ,   \Vert u(x)\Vert=\Vert x\Vert   .$
\vskip1mm \hskip1mm {\em iii)} $  \forall x ,y \in E  ,   \langle u(x), u(y)\rangle =\langle
x,y\rangle    .$ 
\vskip1mm  \hskip1mm {\em iv)} L'image par $  u  $ d'une base orthonorm\'ee de $  E  $ est une base orthonorm\'ee de $  E  .$
\vskip1mm \hskip1mm {\em v)} L'image par $  u  $ de toute base orthonorm\'ee de $  E  $ est une base orthonorm\'ee de $  E  .$
\finenonce{001561}



\finexercice



\exercice{1562, legall, 1998/09/01}

\enonce{001562}{}
Soit $  (E , \langle   ,   \rangle )  $ un espace euclidien et $ \varphi
\in \mathcal{O} (E)  .$  Soit $  F  $ un sous-espace vectoriel de $  E  .$ Montrer que si $  \varphi
(F)\subset F  $ alors $  \varphi (F^\perp )\subset F^\perp    .$ A-t-on \'egalit\'e~?
\finenonce{001562}



\finexercice



\exercice{1563, legall, 1998/09/01}

\enonce{001563}{}
 Soit $  (E , \langle   ,   \rangle )  $ un espace euclidien 
de dimension $  3  $ et $  u\in \mathcal{O} ^-(E)  .$ On pose $  F=\hbox{Ker}(u+id)  .$
\begin{enumerate}
    \item Montrer que $  F\not = \{ 0\}  .$ Montrer que $  F  $ et $  f^\perp  $ sont stables par $  u  .$
Pour quelle raison $  \hbox{dim} (F)\not = 2  ?$
    \item On suppose $  E\not = F  .$ Montrer que la restriction de $  u  $ \`a $F^\perp$ est une rotation.
    \item En d\'eduire qu'il existe $  \theta \in \R   $ et une base $  \epsilon$ de $  E  $ tels que~: 
$$ \hbox{Mat} (u, \epsilon )=\begin{pmatrix}
 \cos (\theta ) & \sin (\theta ) & 0  \cr  -\sin (\theta ) & \cos (\theta )  & 0 \cr 0 & 0 & -1
\cr \end{pmatrix}  .$$
\end{enumerate}
\finenonce{001563}



\finexercice



\exercice{1564, legall, 1998/09/01}

\enonce{001564}{}
Soit $  (E , \langle   ,   \rangle )  $ un espace euclidien 
de dimension $  4  $
et $  \epsilon =\{ e_1, \cdots , e_4\}   $ une base orthonorm\'ee de $  E  .$ Soit $  A  $ la matrice $  \displaystyle{ A=\frac{1}{ 4}  \begin{pmatrix}  0 &

-2 \sqrt{2} & 2 \sqrt2
& 0 \cr   2 \sqrt{2}  & 1 & 1 &  -\sqrt {6} \cr  
-2 \sqrt{2}  & 1 & 1 &  \sqrt {6}  \cr  0  & \sqrt {6} & \sqrt {6} &  2 \cr \end{pmatrix} } $
et $  u \in \mathcal{L} (E)  $ l'endomorphisme d\'etermin\'e par $  \hbox{Mat}(u, \epsilon )=A  .$
\begin{enumerate}
    \item Montrer que $  u \in \mathcal{O} ^+(E)  .$
    \item Montrer que l'espace vectoriel $  F  $ engendr\'e par $  e_1  $ et $  u(e_1)  $ est stable par
$  u  .$ Montrer que la restriction de $  u  $ \`a $  F  $ est une rotation.
    \item Montrer que $  F^\perp   $ est stable par $  u  $ et est engendr\'e par $  e_4  $ et $  u(e_4)  .$ La restriction de $  u  $ \`a $  F^\perp   $ est-elle une rotation~?
\end{enumerate}
\finenonce{001564}



\finexercice



\exercice{1565, legall, 1998/09/01}

\enonce{001565}{}
Soit $  A=(a_{i,j})\in \mathcal{O}(n,\R)  .$ Montrer pour tout $  j \in \{
1, \cdots , n\}  $ l'\'egalit\'e~: $  \displaystyle{ \sum _{i=1}^n a_{i,j}^2=1  .}$ En d\'eduire
que si $  A  $ est triangulaire sup\'erieure elle est diagonale.
\finenonce{001565}



\finexercice



\exercice{1566, legall, 1998/09/01}

\enonce{001566}{}
Soit $  (E , \langle   ,   \rangle )  $ un espace euclidien et 
$  u \in \mathcal{O} (E)  .$ On pose $  v =id-u  .$
\begin{enumerate}
    \item Montrer que $  \hbox{Ker}(v)=\hbox{Im} (v)^\perp   .$
    \item Montrer que $  \displaystyle{ \lim _{n \rightarrow \infty }\frac{1}{ n} \sum _{p=0}^{n-1}u^p (x)} 
$ est la projet\'e orthogonal de $  x  $ sur $  \hbox{Ker}(v)  .$
\end{enumerate}
\finenonce{001566}



\finexercice



\exercice{1567, legall, 1998/09/01}

\enonce{001567}{}
Soit $  (E , \langle   ,   \rangle )  $ un espace euclidien et
$  s \in \mathcal{L} (E)  $ telle que $  s^2=id  .$
\begin{enumerate}
    \item Montrer que $  E=\hbox{Ker}(s-Id)\oplus \hbox{Ker}(s+Id)  .$
    \item Montrer que les propri\'et\'es suivantes sont \'equivalentes~:
\vskip1mm \hskip1mm {\em i)} $  s \in \mathcal{O} (E)  .$
\vskip1mm \hskip1mm {\em ii)} $  \hbox{Ker}(s-Id) \perp  \hbox{Ker}(s+Id)  .$
\vskip1mm \hskip1mm {\em iii)} $  s=s^*  .$
    \item On note d\'esormais $  s_F   $ l'unique sym\'etrie $  s \in \mathcal{O} (E)  $ telle que $ 
F=\hbox{Ker}(s+Id)  .$ Montrer que pour tout $  u\in \mathcal{O} (E)  $ on a~: $  us_Fu^{-1}=s_{u(F)}  .$
    \item Montrer que si $  f   $ est une application de $  E  $ dans lui-m\^eme laissant stables toutes
les droites vectorielles (c'est \`a dire que pour tout $  x\in E  $ il existe $  \lambda _x
\in \R  $ tel que $  f(x)=\lambda _x x  $) alors $  f   $ est lin\'eaire.
    \item En d\'eduire que $  Z(\mathcal{O} (E))=\{ id, -id\}   $ et que si $  n\geq 3  $ alors
$  Z(\mathcal{O} ^+ (E))=\{ id, -id\} \cap \mathcal{O} ^+(E)  .$ (on pourra appliquer 3.) dans le cas o\`u $  F  $
est une droite ou un plan.)
    \item Que se passe-t-il lorsque $  n=1  $ et $  n=2  ?$
\end{enumerate}
\finenonce{001567}



\finexercice



\exercice{1568, ridde, 1999/11/01}

\enonce{001568}{}
Soit $E$ un espace euclidien et $u\in O (E)$ telle que $\ker (u-id) \neq E$.
Soit $x\in E$ tel que $u (x) \neq x$. On pose $y = u (x)$. Alors on sait qu'il
existe une unique r\'eflexion $r$ telle que $r (y) = x$.
\begin{enumerate}
\item Montrer que $\ker (u-id) \subset \ker (r-id)$.
\item Montrer que $\dim \ker (r\circ u-id) >\dim \ker (u-id)$.
\item Montrer par r\'ecurrence que toute isom\'etrie vectorielle est la compos\'ee
de r\'eflexions.
\end{enumerate}
\finenonce{001568}



\finexercice



\exercice{1569, ridde, 1999/11/01}

\enonce{001569}{}
Soit $A = (a_{i, j}) \in O_{n} (\Rr)$. Montrer que $\forall (i, j) \, \left|a_{i, j}\right|
\leq 1$ et que $\left|\sum_{i, j}a_{i, j}\right| \leq n$.
\finenonce{001569}



\finexercice



\exercice{1570, gourio, 2001/09/01}

\enonce{001570}{}
Soit $ E$ euclidien, $ n\in \Nn^{*}, \left(x_{1},...,x_{n},y_{1},
...,y_{n}\right) \in E^{2n}$ tels que :
$$\forall (i,j)\in \{1,...,n\}^{2},\left( x_{i}|x_{j}\right) =\left(y_{i}|y_{j}\right) . $$
Montrer qu'il existe un endomorphisme orthogonal $ f$ de $ E$ tel que :
$$\forall i\in \{1,...,n\},f(x_{i})=y_{i.} $$
\finenonce{001570}



\finexercice


\section{202.99 Autre}


\exercice{1519, barraud, 2003/09/01}

\enonce{001519}{}
On consid\`ere l'application suivante~:
$$
\alpha~:\begin{array}{ccc}
     \R_{n}[X] & \rightarrow & \R_{n}[X] \\
     P         &\mapsto & \int_{0}^{1}P(t)dt
   \end{array}
$$
Montrer que $\alpha$ est une forme lin\'eaire sur $\R_{n}[X]$.

Pour $i\in\{0,...,n\}$, on note $\alpha_{i}$ l'application
$$
\alpha_{i}~:\begin{array}{ccc}
     \R_{n}[X] & \rightarrow & \R_{n}[X] \\
     P         &\mapsto & P(i/n)
   \end{array}
$$
Montrer que $\alpha_{i}$ est une forme lin\'eaire sur $\R_{n}[X]$, et
montrer que la famille $(\alpha_{0},...,\alpha_{n})$ est une base de
${\R_{n}[X]}^{*}$.

En d\'eduire que~:
$$
\exists (\lambda_{0},...,\lambda_{n})\in\R^{n+1},
\forall P\in\R_{n}[X]\quad \int_{0}^{1}P(t)dt = \sum_{i=0}^{n}\lambda_{i}P(i/n)
$$
\finenonce{001519}



\finexercice



\exercice{1520, barraud, 2003/09/01}

\enonce{001520}{}
On consid\`ere l'application $u$ suivante~:
$$
 u~:\begin{array}{ccc}
     \R_{n}[X] & \rightarrow & \R_{n}[X] \\
     P         &\mapsto & P'
   \end{array}
$$
Calculer ${}^{t}u(\alpha)$ lorsque~:
$$
\alpha~: P\mapsto P(0) \qquad \qquad \qquad \alpha~: P\mapsto \int_{0}^{1}P(t)dt
$$
\finenonce{001520}



\finexercice



\exercice{1521, barraud, 2003/09/01}

\enonce{001521}{}
On appelle \emph{trace} d'une matrice $A$, et on note $\mathrm{tr}(A)$, la somme
de ses
coefficients diagonaux.
\begin{enumerate}
\item
Montrer que l'application
 $
 \begin{smallmatrix}
   \mathcal{M}_{n}(\mathbb{K})&\rightarrow  &\mathbb{K}     \\
   A         &\mapsto &\mathrm{tr}(A)
 \end{smallmatrix}
 $
est une forme lin\'{e}aire sur $\mathcal{M}_{n}(\mathbb{K})$.
\item
Montrer que : $\forall(A,B)\in(\mathcal{M}_{n}(\mathbb{K}))^{2},\ \mathrm{tr}(AB)=\mathrm{tr}(BA)$. En d\'{e}duire que deux
matrices semblables ont m\^{e}me trace.
\item
Existe-t-il deux matrices $A$ et $B$ de $\mathcal{M}_{n}(\mathbb{K})$ telles que $AB-BA=I$ ?
\end{enumerate}
\finenonce{001521}



\finexercice



\exercice{1522, barraud, 2003/09/01}

\enonce{001522}{}
Soient $E$ et $F$ deux espaces vectoriels et soit $f\in\mathcal{L}(E,F)$. Montrer que
 $(\mathrm{Im} f)^{\bot}=\mathrm{Ker}{}^t{f}$.

En d\'{e}duire que $f$ est surjective si et seulement si ${}^t{f}$ est
injective.

Lorsque $E$ et $F$ sont de dimension finies, montrer que $\mathrm{rg}(f)=\mathrm{rg}({}^t{f})$. En
d\'{e}duire que pour toute matrice $A\in\mathcal{M}_{m,n}(\mathbb{K})$ on a $\mathrm{rg}(A)=\mathrm{rg}({}^t(A))$.
\finenonce{001522}



\finexercice



\exercice{1523, barraud, 2003/09/01}

\enonce{001523}{}
Soit $f\in\mathrm{End}(\R^{3})$ tel que $f^{2}=0$. Montrer qu'il existe $\alpha\in{(\R^{3})}^{*}$
et $v\in\R^{3}$ tels que
 $$
 \forall x\in \R^{3} \;\; f(x)=\alpha(x)\,v.
 $$
\textit{\small(Indication : commencer par montrer que $\mathrm{rg}(f)\leq 1$)}
\finenonce{001523}



\finexercice



\exercice{1524, barraud, 2003/09/01}

\enonce{001524}{}
On considère un espace euclidien $(E,<>)$. On rappelle que l'adjoint
$u^{*}$ d'un endomorphisme $u$ est l'endomorphisme caractérisé par~:
$$
\forall (x,y)\in E^{2},\ <u(x),y>=<x,u^{*}(y)>
$$

On dit qu'un endomorphisme $u$ de $E$ est une similitude de $E$ si et
seulement si $u$ est la composée d'une homotétie et d'une isométrie,
c'est à dire si et seulement si~:
$$
 \exists\alpha\in\R\setminus\{0\},\ \exists v\in O(E)\ /\ u=\alpha v.
$$

\begin{enumerate}
\item
Redémontrer l'équivalence entre les trois caractérisations suivantes des
isométries~:
\begin{align*}
  v\text{ est une isométrie} 
  &\Leftrightarrow \forall x\in E\quad \Vert v(x)\Vert=\Vert x\Vert\\
  &\Leftrightarrow \forall (x,y)\in E^{2}\quad <v(x),v(y)>=<x,y>\\
  &\Leftrightarrow v^{*}v=\mathrm{id}
\end{align*}
\end{enumerate}

\bigskip

On veut montrer l'équivalence des assertions suivantes~:

\begin{enumerate}
\item[{\it(i)}] 
$u$ est une similitude
\item[{\it(ii)}]
il existe un réel $\lambda>0$ tel que 
$$
 u^{*}u=\lambda\mathrm{id}
$$ 

\item[{\it(iii)}]
$u$ conserve l'orthogonalité, c'est à dire :
$$
 \forall(x,y)\in E^{2}, <x,y>=0\Leftrightarrow <u(x),u(y)>=0
$$
\end{enumerate}

\bigskip


\begin{enumerate}
\setcounter{enumi}{1}
\item
  Montrer que $(i)\Rightarrow (ii)$, puis que $(ii)\Rightarrow(i)$.

\item
  Montrer que $(i)\Rightarrow(iii)$.

\item  
  On suppose $(iii)$.

  \begin{enumerate}
  \item 
    Soit $x\in E$, $x\neq 0$.Montrer que
    $$
    \forall y\in E \quad x\bot y\Leftrightarrow u^{*}u(x)\bot y
    $$
  \item
    En déduire que $u^{*}u(x)$ appartient à la droite engendrée par $x$. On
    note $\lambda_{x}$ le réel tel que $u^{*}u(x)=\lambda_{x} x$.
  \item
    Montrer que~: $\forall t\in\R,\ \lambda_{tx}=\lambda_{x}$
  \item
    Montrer que, pour tout couple $(x,y)$ de vecteurs linéairement
    indépendants de $E$, on a~: $\lambda_{x}=\lambda_{y}$. 
  \item
    En déduire que l'application $x\mapsto\lambda_{x}$ est constante.
    Conclure.
  \end{enumerate}
\end{enumerate}
\finenonce{001524}



\finexercice



\exercice{3567, quercia, 2010/03/10}
\enonce{003567}{\'Equation $AM = \lambda M$}

Soit $A \in \mathcal{M}_n(K)$. Déterminer les scalaires $\lambda$ et les matrices
$M \in \mathcal{M}_n(K)$ telles que $AM = \lambda M$.
\finenonce{003567}



\finexercice


\exercice{3568, quercia, 2010/03/10}
\enonce{003568}{$v \mapsto v\circ u$ (Centrale MP 2003)}
Soit~$E$ un espace vectoriel de dimension finie et $u\in\mathcal{L}(E)$. On considère l'application~$\Phi_u$ qui à $v\in\mathcal{L}(E)$
associe~$v\circ u$.
\begin{enumerate}
  \item Montrer que $\Phi_u\in\mathcal{L}({\mathcal{L}(E)})$.
  \item Montrer l'équivalence~: ($u$ est diagonalisable) $\Leftrightarrow$ ($\Phi_u$ est diagonalisable)\dots
  \begin{enumerate}
    \item en considérant les polynômes annulateurs de~$u$ et de~$\Phi_u$.
    \item en considérant les spectres et sous-espaces propres de~$u$ et de~$\Phi_u$.
  \end{enumerate}
\end{enumerate}
\finenonce{003568}



\finexercice


\exercice{3569, quercia, 2010/03/10}
\enonce{003569}{Matexo}

Soit $E$ un $ K$-espace vectoriel de dimension $n$, $(F,G)$ deux sous-espaces vectoriels
de $E$ tels que 
$E = F \oplus G$. On note $p$ la projection sur $F$ parallèlement à~$G$.
Soit $E_p = \{f\in\mathcal{L}(E) \text{ tq } f\circ p = p\circ f\}$.
Quelle est la dimension de $E_p$ ?
\finenonce{003569}



\finexercice


\exercice{3570, quercia, 2010/03/10}
\enonce{003570}{$f  \mapsto p\circ f\circ p$}

Soit $p \in \mathcal{L}(E)$ une projection et
$\Phi : {\mathcal{L}(E)} \to {\mathcal{L}(E)}, f \mapsto {p\circ f\circ p.}$

Déterminer les éléments propres de $\Phi$.
\finenonce{003570}



\finexercice


\exercice{3571, quercia, 2010/03/10}
\enonce{003571}{$f  \mapsto u\circ f$ et $f  \mapsto f\circ u$}

Soit $E$ un $ K$-ev de dimension finie et $u \in \mathcal{L}(E)$ diagonalisable.

On considère les applications%
${\mathcal{L}(E)} \to  {\mathcal{L}(E)}$ $ {\varphi : f \mapsto u\circ f}$ ${\psi : f \mapsto f\circ u}$

\begin{enumerate}
  \item Montrer que $\varphi$ et $\psi$ sont diagonalisables.
  \item Montrer que $\varphi-\psi$ est diagonalisable.
\end{enumerate}
\finenonce{003571}



\finexercice


\exercice{3572, quercia, 2010/03/10}
\enonce{003572}{$u\circ v - v\circ u = \mathrm{id}$}

Soient $u,v$ deux endomorphisme d'un espace vectoriel $E$ non nul tels que
$u\circ v - v\circ u = \mathrm{id}_E$.

\begin{enumerate}
  \item Simplifier $u^k\circ v - v\circ u^k$ pour $k\in\N$ puis
    $P(u)\circ v - v\circ P(u)$ pour $P\in K[X]$.

  \item Montrer que $u$ et $v$ n'ont pas de polynômes minimaux.
\end{enumerate}
\finenonce{003572}



\finexercice


\exercice{3573, quercia, 2010/03/10}
\enonce{003573}{Ensi PC 1999}

Soit $E$ un ev réel de dimension finie et $f,g\in \mathcal{L}(E)$, $\alpha\in\R^*$
tels que $f\circ g - g\circ f = \alpha f$.

\begin{enumerate}
  \item Montrer pour tout entier naturel~$n$~:
    $f^n\circ g - g\circ f^n = \alpha n f^n$.
  \item Montrer qu'il existe~$n\in\N$ tel que $f^n = 0$
    (raisonner par l'absurde et considérer l'application
     $h  \mapsto h\circ g - g\circ h$ de $\mathcal{L}(E)$ dans $\mathcal{L}(E)$).
\end{enumerate}
\finenonce{003573}



\finexercice


\exercice{3574, quercia, 2010/03/10}
\enonce{003574}{X MP$^*$ 2001}

Soit $f$ un endomorphisme de $E$ (ev de dimension finie sur $ K$) tel que
$\chi_f$ soit irréductible. Montrez que pour aucun endomorphisme $g$
le crochet de Lie $[f,g] = f\circ g - g\circ f$ n'est de rang un.

\finenonce{003574}



\finexercice


\exercice{3575, quercia, 2010/03/10}
\enonce{003575}{$\frac12(p\circ u + u\circ p)$ (Mines MP 2003)}
Soit~$E$ un espace vectoriel de dimension $n$ finie, $p$ un projecteur de rang~$r$
et $\varphi : {\mathcal{L}(E)} \to {\mathcal{L}(E)}, u \mapsto {\frac12(p\circ u + u\circ p).}$
\begin{enumerate}
  \item Est-ce que~$\varphi$ est diagonalisable~?
    
  \item Déterminer les valeurs propres de~$\varphi$ et les dimensions des sous-espaces propres.
    
\end{enumerate}
\finenonce{003575}



\finexercice


\exercice{3576, quercia, 2010/03/10}
\enonce{003576}{Crochet de Lie (Ens Cachan MP$^*$ 2003)}
Soit $\Phi : {\mathcal{M}_n(\C)}\to {\mathcal{M}_n(\C)}$ un automorphisme d'ev tel que~:
$\forall\ A,B\in\mathcal{M}_n(\C),\ \Phi([A,B]) = [\Phi(A),\Phi(B)]$ où $[X,Y] = XY - YX$.
Montrer~: $\forall\ D\in\mathcal{M}_n(\C)$, ($D$ est diagonalisable) $\Leftrightarrow$ ($\Phi(D)$ est diagonalisable).

{\it Indication~: considérer $\phi_D$ : $X  \mapsto[D,X]$ et montrer que
($D$ est diagonalisable) $\Leftrightarrow$ ($\phi_D$ est diagonalisable).}

\finenonce{003576}



\finexercice

\exercice{3625, quercia, 2010/03/10}
\enonce{003625}{Base de $( K^3)^*$}

Dans $ K^3$ on considère les formes linéaires :
$f_1(\vec x) = x+y-z$,
$f_2(\vec x) = x-y+z$,
$f_3(\vec x) = x+y+z$.

\begin{enumerate}
  \item Montrer que $(f_1,f_2,f_3)$ est une base de $( K^3)^*$.
  \item Trouver la base duale.
\end{enumerate}
\finenonce{003625}



\finexercice


\exercice{3626, quercia, 2010/03/10}
\enonce{003626}{Base de $( K^3)^*$}

Dans $ K^3$ on considère les formes linéaires :
$f_1(\vec x) = x+2y+3z$,
$f_2(\vec x) = 2x+3y+4z$,
$f_3(\vec x) = 3x+4y+6z$.

\begin{enumerate}
  \item Montrer que $(f_1,f_2,f_3)$ est une base de $( K^3)^*$.
  \item Trouver la base duale.
    
\end{enumerate}
\finenonce{003626}



\finexercice


\exercice{3627, quercia, 2010/03/10}
\enonce{003627}{Base de $( K^n)^*$}

Pour $\vec x = (x_1,\dots,x_n) \in  K^n$ on pose $f_i(\vec x) = x_i+x_{i+1}$
et $f_n(\vec x) = x_n+x_1$.
Déterminer si ${\cal F} = (f_1,\dots,f_n)$ est une base de $( K^n)^*$ et,
le cas échéant, déterminer la base duale.



\finenonce{003627}



\finexercice


\exercice{3628, quercia, 2010/03/10}
\enonce{003628}{Bases duale des polynômes}

Soit $E =  K_n[X]$. Montrer que la famille ${\cal F} = (f_0,\dots,f_n)$ est
une base de $E^*$ et donner la base duale lorsque~\dots

\begin{enumerate}
  \item $f_i(P) = P(x_i)$ où $x_0,\dots,x_n$ sont des scalaires distincts.
  \item $f_i(P) = P^{(i)}(0)$.
  \item $f_i(P) = P^{(i)}(x_i)$ où $x_0,\dots,x_n$ sont des scalaires quelconques.
    (Ne pas chercher la base duale pour cet exemple)
\end{enumerate}
\finenonce{003628}



\finexercice


\exercice{3629, quercia, 2010/03/10}
\enonce{003629}{Base duale des polynômes}

Soit $E =  \R_{2n-1}[X]$, et $x_1,\dots,x_n \in \R$ distincts.
On note :
$${\phi_i} : E\to \R, P \mapsto {P(x_i)} ;
 \qquad
 {\psi_i} : E \to \R, P \mapsto {P'(x_i)}$$
\begin{enumerate}
  \item Montrer que $(\phi_1,\dots,\phi_n,\psi_1,\dots,\psi_n)$ est une base de
    $E^*$.
  \item Chercher la base duale.
    On notera $P_i = \prod_{j\ne i}\frac{X-x_j}{x_i-x_j}$ et
    $d_i = P_i'(x_i)$.

    
\end{enumerate}
\finenonce{003629}



\finexercice


\exercice{3630, quercia, 2010/03/10}
\enonce{003630}{Intégrale}

Soit $E = \R_3[X]$. On considère les formes linéaires :
$f_i : P  \mapsto  \int_{t=-1}^1 t^iP(t)\,d t$.

\begin{enumerate}
  \item Montrer que $(f_0,f_1,f_2,f_3)$ est une base de $E^*$.
  \item Trouver la base duale.
\end{enumerate}
\finenonce{003630}




\finexercice


\exercice{3631, quercia, 2010/03/10}
\enonce{003631}{\'Evaluations et intégrale}

Soit $E = \R_3[X]$ et $a,b,c \in \R$ distincts.
On considère les formes linéaires sur E :
$$f_a\ :\ P \mapsto P(a),\quad
  f_b\ :\ P \mapsto P(b),\quad
  f_c\ :\ P \mapsto P(c),\quad
  \varphi\ :\ P  \mapsto  \int_{t=a}^b P(t)\,d t.$$
\'Etudier la liberté de $(f_a,f_b,f_c,\varphi)$.

\finenonce{003631}



\finexercice


\exercice{3632, quercia, 2010/03/10}
\enonce{003632}{Base duale de $(1, X, X(X-1), \dots)$}

Soit $E =  K_n[X]$. On note $P_0 = 1$, $P_i = X(X-1)\cdots(X-i+1)$ pour $i\ge1$,
et $f_i : P  \mapsto P(i)$.

\begin{enumerate}
  \item Montrer que $(P_0,\dots,P_n)$ est une base de $E$ et ${\cal B} = (f_0,\dots,f_n)$
    est une base de $E^*$.
  \item Décomposer la forme linéaire $P_n^*$ dans la base $\cal B$.
    (On pourra utiliser les polynômes :
    $Q_i = \prod_{1\le j\le n; j\ne i}(X-j)$)
    
  \item Décomposer de même les autres formes linéaires $P_k^*$.
    
\end{enumerate}
\finenonce{003632}



\finexercice


\exercice{3633, quercia, 2010/03/10}
\enonce{003633}{Base duale de $((X-a)^k(X-b)^{n-k})$}

Soit $E =  K_n[X]$ et $a,b \in  K$ distincts.
On pose $P_k = (X-a)^k(X-b)^{n-k}$.

\begin{enumerate}
  \item Montrer que $(P_0,\dots,P_n)$ est une base de $E$.
  \item On suppose $n = 2$ et on prend comme base de $E^*$ :
    ${\cal B} = (f_a,f_c,f_b)$ où $f_x(P) = P(x)$ et $c = \frac{a+b}2$.
    Exprimer les formes linéaires $(P_0^*,P_1^*,P_2^*)$ dans $\cal B$.
    
\end{enumerate}
\finenonce{003633}



\finexercice


\exercice{3634, quercia, 2010/03/10}
\enonce{003634}{Formes linéaires liées}

Soient $f_1,\dots,f_n$ des formes linéaires sur $ K^n$ telles qu'il existe
$\vec x \in  K^n$ non nul tel que $f_1(\vec x) = \dots = f_n(\vec x) = 0$.
Montrer que $(f_1,\dots,f_n)$ est liée.
\finenonce{003634}



\finexercice


\exercice{3635, quercia, 2010/03/10}
\enonce{003635}{$(P(X),\dots,P(X+n))$ est une base des polynômes}

Soit $E =  K_n[X]$, $Q \in E$ de degré $n$ et $Q_i = Q(X+i)$ ($0 \le i \le n$).

\begin{enumerate}
  \item Montrer que $(Q,Q',Q'',\dots,Q^{(n)})$ est libre.
  \item Montrer que toute forme linéaire sur $E$ peut se mettre sous la forme :\par
    $f : P  \mapsto \alpha_0 P(0) + \alpha_1 P'(0) + \dots + \alpha_nP^{(n)}(0)$.
  \item Soit $f \in E^*$ telle que $f(Q_0) = \dots = f(Q_n) = 0$.
    Montrer que $f=0$.\par
    (considérer le polynôme $P = \alpha_0Q + \dots + \alpha_nQ^{(n)}$)
  \item  Montrer que $(Q_0,\dots,Q_n)$ est une base de $E$.
\end{enumerate}
\finenonce{003635}



\finexercice


\exercice{3636, quercia, 2010/03/10}
\enonce{003636}{$\varphi((X-a)P) = 0$}

Soit $E =  K_n[X]$. Soit $\varphi \in E^*$ telle que :
    $\forall\ P \in  K_{n-1}[X],\ \varphi((X-a)P) = 0$.

    Montrer qu'il existe $\lambda \in  K$ tel que :
    $\forall\ P \in E,\ \varphi(P) = \lambda P(a)$.

\finenonce{003636}



\finexercice


\exercice{3637, quercia, 2010/03/10}
\enonce{003637}{Forme linéaire sur les polynômes}

Montrer l'existence et l'unicité d'une forme linéaire $\Phi$ sur $\R_n[X]$
telle que : $\Phi(1) = 0$, $\Phi(X) = 1$ et $\Phi(P) = 0$ pour tout polynôme
$P \in \R_n[X]$ tel que $P(0) = P(1) = 0$.
\finenonce{003637}



\finexercice


\exercice{3638, quercia, 2010/03/10}
\enonce{003638}{Système unisolvent}

Soient $f_1,\dots,f_n : \R \to \R$ $n$ fonctions linéairement indépendantes.
On pose $E = \text{vect}(f_1,\dots,f_n)$ et pour $x \in \R$,
$ {\delta_x} : E \to \R, f \mapsto {f(x).}$

\begin{enumerate}
  \item Montrer que la famille $(\delta_x)_{x\in\R}$ engendre $E^*$.
  \item En déduire qu'il existe $x_1,\dots,x_n \in \R$ tels que $\det M \ne 0$ où
    $M$ est la matrice de terme général $f_i(x_j)$.
\end{enumerate}
\finenonce{003638}



\finexercice


\exercice{3639, quercia, 2010/03/10}
\enonce{003639}{Polynômes à deux variables}

Soit $f : {\R^2} \to \R$. On dit que $f$ est polynomiale si elle est de la
forme : $f(x,y) = \sum_{i,j} a_{ij}x^iy^j$, la somme portant sur un nombre
fini de termes. Le degré de $f$ est alors $\max(i+j \text{ tq } a_{ij} \ne 0)$.

On note $E_k$ l'ensemble des fonctions $\R^2 \to \R$ polynomiales de degré
inférieur ou égal à $k$.

\begin{enumerate}
  \item Montrer que $E_k$ est un $\R$-ev de dimension finie et donner sa dimension.
    
  \item Soient $A=(0,0)$, $B=(1,0)$, $C=(0,1)$.
    Montrer que les formes linéaires $f  \mapsto f(A)$, $f  \mapsto f(B)$, $f  \mapsto f(C)$
    constituent une base de $E_1^*$.
  \item Chercher de même une base de $E_2^*$.
    
  \item Soit $T$ le triangle plein $ABC$ et $f \in E_1$.
    Montrer que $ \int_{}^{}\!\! \int_T^{} f(x,y)\,d xd y =
    \frac{f(A)+f(B)+f(C)}6$.
  \item Chercher une formule analogue pour $f \in E_2$.
    
\end{enumerate}
\finenonce{003639}



\finexercice


\exercice{3640, quercia, 2010/03/10}
\enonce{003640}{Polynômes trigonométriques}

On note $f_n(x) = \cos nx$ et $g_n(x) = \sin nx$ ($x \in \R$, $n \in \N$).

Soit $E_n$ l'espace engendré par la famille
${\cal F}_n = (f_0,\dots,f_n,g_1,\dots,g_n)$.

\begin{enumerate}
  \item Montrer que pour $k \ge 1$, $(f_k,g_k)$ est libre.
  \item Soit $\varphi : {E_n} \to {E_n}, f \mapsto {f''}$.
    Chercher les sous-espaces propres de $\varphi$.
    En déduire que ${\cal F}_n$ est libre.
  \item On note $a_k = \frac{2k\pi}{2n+1}$ et
    ${\varphi_k} : {E_n} \to \R, f \mapsto {f(a_k).}$

    Montrer que $(\varphi_0,\dots,\varphi_{2n})$ est une base de $E_n^*$.
    (On utilisera la fonction $f : x  \mapsto \prod_{k=1}^n (\cos x - \cos a_k)$)
    

  \item Soit $N \in \N^*$. On note $b_k = \frac{2k\pi}N$ et
    ${\psi_k} : {E_n}\to \R, f \mapsto {f(b_k).}$
    Montrer que $(\psi_0,\dots,\psi_{N-1})$ est libre si et seulement si
    $N \le 2n+1$, et engendre $E_n^*$ si et seulement si $N \ge 2n+1$.
\end{enumerate}
\finenonce{003640}



\finexercice


\exercice{3641, quercia, 2010/03/10}
\enonce{003641}{trace sur $\mathcal{M}_n(K)$}

Soit $E = \mathcal{M}_n(K)$. Pour $A \in \mathcal{M}_n(K)$, on note
$ {\phi_A} : E \to K, M \mapsto {\mathrm{tr}(AM).}$

\begin{enumerate}
  \item Montrer que $E^* = \{ \phi_A \text{ tq } A \in E\}$.
  \item On note $\cal S$ l'ensemble des matrices symétriques
         et $\cal A$ l'ensemble des matrices antisymétriques.
    Montrer que ${\cal S}^\circ = \{ \phi_A \text{ tq } A \in {\cal A}\}$
     et ${\cal A}^\circ = \{ \phi_A \text{ tq } A \in {\cal S}\}$.
\end{enumerate}
\finenonce{003641}



\finexercice


\exercice{3642, quercia, 2010/03/10}
\enonce{003642}{Suites de Fibonacci}

Soit $E$ l'ensemble des suites $u=(u_n)$ à termes réels telles que pour tout
$n$ : $u_{n+2} = u_{n+1} + u_n$.

\begin{enumerate}
  \item Montrer que $E$ est un $\R$-ev de dimension finie.
  \item Soient $ {f_0} : E \to \R, u \mapsto {u_0}$ et
           ${f_1} : E \to \R, u \mapsto {u_1.}$
    Trouver la base duale de $(f_0,f_1)$.
    

\end{enumerate}
\finenonce{003642}



\finexercice


\exercice{3643, quercia, 2010/03/10}
\enonce{003643}{Combinaison de formes linéaires}

Soit $E$ un $ K$-ev de dimension finie et $f$, $f_1,\dots,f_p \in E^*$.
Montrer que $f$ est combinaison linéaire de $f_1,\dots,f_p$ si et seulement si
$\mathrm{Ker} f \supset \mathrm{Ker} f_1 \cap \dots \cap \mathrm{Ker} f_p$.
\finenonce{003643}



\finexercice


\exercice{3644, quercia, 2010/03/10}
\enonce{003644}{Combinaison de formes linéaires}

Soit $E$ un $ K$-ev de dimension finie et $f_1,\dots,f_p \in E^*$.
On considère l'application :
$$\Phi : E \to { K^p}, {\vec x} \mapsto {\Bigl(f_1(\vec x),\dots,f_p(\vec x)\Bigr)}$$
Montrer que $\Phi$ est surjective si et seulement si $(f_1,\dots,f_p)$ est libre.
\finenonce{003644}



\finexercice


\exercice{3645, quercia, 2010/03/10}
\enonce{003645}{Orthogonaux d'une somme directe}

Soit $E$ un $ K$-ev de dimension finie et $F,G$ deux sev de $E$ tels que
$F\oplus G = E$.

\begin{enumerate}
  \item Montrer que $F^\circ \oplus G^\circ = E^*$.
  \item Montrer que $F^\circ$ est naturellement isomorphe à $G^*$ et
 $G^\circ$ à $F^*$.
\end{enumerate}
\finenonce{003645}



\finexercice


\exercice{3646, quercia, 2010/03/10}
\enonce{003646}{$f(u) \ne 0$ et $g(u) \ne 0$}

Soit $E$ un $ K$-ev de dimension finie et $f,g \in E^*$ toutes deux non nulles.
Montrer qu'il existe un vecteur $\vec u \in E$ tel que $f(\vec u) \ne 0$
et $g(\vec u) \ne 0$.
\finenonce{003646}



\finexercice


\exercice{3647, quercia, 2010/03/10}
\enonce{003647}{$p$ formes linéaires}

Soit $E$ un $ K$-ev. On suppose qu'il existe $p$ formes linéaires
$f_1,\dots,f_p$ telles que :
$$\forall\ \vec x \in E,\ (f_1(\vec x) = \dots = f_p(\vec x) = 0)
   \Rightarrow  (\vec x = \vec 0).$$
Montrer que $E$ est de dimension finie inférieure ou égale à $p$.
\finenonce{003647}



\finexercice


\exercice{3648, quercia, 2010/03/10}
\enonce{003648}{$\det(f_i(u_j)) \ne 0$}

Soit $E$ un $ K$-ev de dimension finie $n$,
${\vec u}_1,\dots,{\vec u}_n \in E$ et $f_1,\dots,f_n \in E^*$.

Soit $M$ la matrice de terme général $f_i({\vec u}_j)$.
Montrer que si $\det M \ne 0$, alors $({\vec u}_1,\dots,{\vec u}_n)$
est une base de $E$ et $(f_1,\dots,f_n)$ est une base de $E^*$.
\finenonce{003648}



\finexercice


\exercice{3649, quercia, 2010/03/10}
\enonce{003649}{$e_n^*$ imposé}

Soit $E$ un $ K$-ev de dimension finie $n$,
${\cal F} = ({\vec e}_1,\dots, {\vec e}_{n-1})$ une famille libre de $E$
et $f \in E^*\setminus\{0\}$.
Montrer qu'on peut compléter $\cal F$ en une base
${\cal B} = ({\vec e}_1,\dots,{\vec e}_n)$
telle que $f = e_n^*$ si et seulement si
$f({\vec e}_1) = \dots = f({\vec e}_{n-1}) = 0$.

Y~a-t-il unicité de ${\vec e}_n$ ?
\finenonce{003649}



\finexercice


\exercice{3650, quercia, 2010/03/10}
\enonce{003650}{Modification élémentaire}

Soit $E$ un $ K$-ev de dimension finie $n$,
${\cal B} = ({\vec e}_1,\dots,{\vec e}_n)$ une base de $E$, et
${\cal B}' = ({\vec e}_1{}',\dots,{\vec e}_n{}')$ déduite de $\cal B$ par une
opération élémentaire (échange de deux vecteurs, multiplication d'un vecteur
par un scalaire non nul, addition à un vecteur d'un multiple d'un autre).

\'Etudier comment on passe de la base duale ${\cal B}^*$ à ${\cal B}'^*$
en fonction de l'opération effectuée.


\finenonce{003650}



\finexercice


\exercice{3651, quercia, 2010/03/10}
\enonce{003651}{Séparation}

Soit $E$ un $ K$-ev de dimension finie.

\begin{enumerate}
  \item Soient $\vec x,\vec y \in E$ avec $\vec x \ne \vec y$.
    Montrer qu'il existe $f \in E^*$ telle que $f(\vec x) \ne f(\vec y)$.
  \item Soit $V$ un sev de $E^*$ ayant la propriété :
    $\forall\ \vec x,\vec y \in E,\ \vec x\ne\vec y  \Rightarrow  \exists\ f \in V
    \text{ tq }   f(\vec x) \ne f(\vec y)$.

    Montrer que $V = E^*$.
\end{enumerate}
\finenonce{003651}



\finexercice

\section{203.01 Groupe, sous-groupe}


\exercice{1303, ridde, 1999/11/01}

\enonce{001303}{}
Soit $ABC$ un triangle \'equilat\'eral du plan.
\begin{enumerate}
\item D\'eterminer l'ensemble des rotations qui laissent invariant $\left\{ A, B, C\right\}$.
\item Montrer que c'est un groupe pour la loi $\circ$.
\item Faire de même avec un carr\'e.
\end{enumerate}
\finenonce{001303}



\finexercice



\exercice{1304, cousquer, 2003/10/01}

\enonce{001304}{Entiers modulo $n$}
Étant donné un entier naturel $n$, on appelle classe d'un entier relatif $p$ modulo
$n$ l'ensemble $\overline p = \{p+kn\mid k\in\mathbb{Z}\}$.
L'ensemble des classes modulo
$n$ est noté $\mathbb{Z}_n$.
\begin{enumerate}
    \item Écrire la liste des éléments distincts de $\mathbb{Z}_2$, $\mathbb{Z}_3$, $\mathbb{Z}_4$
et~$\mathbb{Z}_5$.
    \item Montrer que si $x\in\overline p$ et $y\in\overline q$, alors $x+y\in
\overline{p+q}$ et $xy\in\overline{pq}$.
    \item En posant $\overline p+\overline q=\overline{p+q}$ et $\overline
p\cdot\overline q=\overline{pq}$, on définit deux lois de composition,
addition et multiplication sur $\mathbb{Z}_n$.\\
Écrire la table d'addition et de multiplication de $\mathbb{Z}_4$.
\end{enumerate}
Même question pour $\mathbb{Z}_2$, $\mathbb{Z}_3$, et~$\mathbb{Z}_5$.
\finenonce{001304}



\finexercice



\exercice{1305, cousquer, 2003/10/01}

\enonce{001305}{}
\begin{enumerate}
\item Montrer que les transformations géométriques qui conservent
globalement un rectangle forment un groupe. Faire l'étude de ce groupe.
\item Étudier le groupe $\mathbb{Z}/4\mathbb{Z}$.
\item Montrer qu'il n'existe que deux sortes de groupes à quatre éléments.
\end{enumerate}
\finenonce{001305}



\finexercice



\exercice{1306, cousquer, 2003/10/01}

\enonce{001306}{}
\begin{enumerate}
\item Étudier le groupe des isométries du carré.
\item Écrire la liste des éléments du groupe $\mathfrak{S}_4$ des permutations
de quatre lettres. Trouver des sous-groupes de ce groupe isomorphes
aux groupes du rectangle, du triangle équilatéral, du carré.
\end{enumerate}
\finenonce{001306}



\finexercice



\exercice{1307, cousquer, 2003/10/01}

\enonce{001307}{Permutations d'un ensemble de $n$ éléments}
\begin{enumerate}
    \item Une permutation de l'ensemble de $n$ éléments $\{1,2,\ldots,n\}$ est
une bijection de cet ensemble dans lui-même. Il est commode de désigner une
telle permutation $s$ par le tableau de valeurs suivant~:
$s=
\left(\begin{array}{cccc}
  1   &   2&\cdots&n\\
  s(1)&s(2)&\cdots&s(n)
\end{array}\right)$. On note
$\frak{S}_n$ l'ensemble de ces permutations pour $n$ donné.

    \item  Écrire les éléments de $\frak{S}_2$ et de $\frak{S}_3$.

    \item  Établir les tables de composition de ces deux ensembles.

    \item  De la table de $\frak{S}_3$ on peut extraire des parties \emph{stables}
ne faisant intervenir que certains éléments~; lesquelles~? Peut-on les
trouver toutes.

    \item  Voyez-vous des analogies (totales ou partielles) entre ces tables et des
situations rencontrées plus haut~?

    \item  On peut obtenir tous les éléments de $\frak{S}_3$ à partir de la
composition de certains d'entre-eux~; lesquels~?

    \item  Combien d'éléments possède $\frak{S}_n$~? Combien de cases
contient la table de composition de $\frak{S}_4$, $\frak{S}_5,\ldots$~?
Pourrait-on étudier $\frak{S}_4$ et $\frak{S}_5$ à partir de ces tables~?
\end{enumerate}
\finenonce{001307}



\finexercice



\exercice{1308, ridde, 1999/11/01}

\enonce{001308}{}
Soient les quatre fonctions de $\Rr^*$ dans $\Rr^*$
$$ f_1 (x) = x \qquad f_2 (x) = \frac 1x \qquad f_3 (x) = -x \qquad f_4 (x) = -\frac 1x$$
Montre que $G = \left\{ f_1, f_2, f_3, f_4\right\}$ est un groupe pour la loi $\circ$.
\finenonce{001308}



\finexercice



\exercice{1309, ridde, 1999/11/01}

\enonce{001309}{}
Montrer qu'il existe une seule table possible pour un groupe d'ordre $3$.
Est-ce vrai pour $4$ ?
\finenonce{001309}



\finexercice



\exercice{1310, ridde, 1999/11/01}

\enonce{001310}{}
Montrer que si $X$ contient au moins trois \'el\'ements alors $ \sigma (X)$ n'est pas ab\'elien.
\finenonce{001310}



\finexercice



\exercice{1311, ortiz, 1999/04/01}

\enonce{001311}{}
 Les ensembles suivants, pour les lois
consid\'er\'ees, sont-ils des groupes ?
\begin{enumerate}
    \item $]-1,1[$ muni de la loi d\'efinie par $x\star y=\frac{x+y}{1+xy}$ ;
    \item $\{z\in\Cc:{|z|}=2\}$ pour la multiplication usuelle ;
    \item $\Rr_+$ pour la multiplication usuelle;
    \item $\left\{x\in\Rr\mapsto ax+b : a\in\Rr\setminus\left\{0\right\},b\in\Rr\right\}$
pour la loi de composition des applications.
\end{enumerate}

\finenonce{001311}



\finexercice



\exercice{1312, ortiz, 1999/04/01}

\enonce{001312}{}
Soit $K=\left\{\text{Id}, f_1,f_2,f_3\right\}$ o\`u $ f_1,$
$f_2,$ et $f_3$ sont les permutations de
$E=\left\{1,2,3,4\right\}$ d\'efinies par
$${f_1=\left(\begin{smallmatrix}
1&2&3&4\\2&1&4&3\end{smallmatrix}\right)
,f_2=\left(\begin{smallmatrix}
1&2&3&4\\3&4&1&2\end{smallmatrix}\right)
,f_3=\left(\begin{smallmatrix}
1&2&3&4\\4&3&2&1\end{smallmatrix}\right).}$$
 Montrer
que $K$ est un sous-groupe de $\mathcal{S}_4$.
\finenonce{001312}



\finexercice



\exercice{1313, ortiz, 1999/04/01}

\enonce{001313}{}
Soit l'ensemble
$${\mathcal{J}=\left\{
\left (
\begin{array}{ll}
x&x\\
x&x
\end{array}
\right )\in \mathcal{M}_2(\Rr): x\in \Rr\setminus\left\{0\right\}
\right\}.}$$ Montrer que, muni de la multiplication
usuelle des matrices, $\mathcal{J}$ est un groupe
ab\'elien.
\finenonce{001313}



\finexercice



\exercice{1314, ortiz, 1999/04/01}

\enonce{001314}{}
 Pour la multiplication usuelles des matrices
carr\'ees, les ensembles suivants sont-ils des groupes : $${
\text{GL}(2,\Rr)\cap\mathcal{M}_2(\Zz),\quad
\left\{M\in\mathcal{M}_2(\Zz):\det M=1\right\} \text{ ?} }$$

\finenonce{001314}



\finexercice



\exercice{1315, ortiz, 1999/04/01}

\enonce{001315}{}
Soit $G$ un ensemble muni d'une loi de composition
interne associative, admettant un \'el\'ement
neutre \`a droite et tel que chaque \'el\'ement de
$G$ admette un sym\'etrique \`a droite. Montrer
que $G$ est un groupe.
\finenonce{001315}



\finexercice



\exercice{1316, ortiz, 1999/04/01}

\enonce{001316}{}
Soient $(G,.)$ un groupe et $a,b\in G.$ On suppose
que
$${(1): ab^2=b^3a \quad\text{ et }(2):
ba^2=a^3b.}$$
\begin{enumerate}
    \item Montrer, en utilisant seulement (1), que $a^2b^8a^{-2}=b^{18}$
puis que $a^3b^8a^{-3}=b^{27}.$
    \item En d\'eduire, en utilisant (2), que
$a^3b^8a^{-3}=b^{18}$ et enfin que
 $a=b=1.$
\end{enumerate}
\finenonce{001316}



\finexercice



\exercice{1317, ortiz, 1999/04/01}

\enonce{001317}{}
\begin{enumerate}
    \item  L'ensemble $\Rr\setminus\left\{-1\right\}$ muni de la loi $\star$
    d\'efinie par $\forall a,b\in
\Rr,a\star b=a+b+ab$ est-il un groupe~~?
    \item  L'ensemble $E=\left\{-1,1,i,-i \right\}\subseteq \Cc$ muni de la loi usuelle de
multiplication dans $\Cc$ est-il un groupe~?
    \item  L'ensemble $E=\left\{\left(\begin{smallmatrix} a&0\\0&0
\end{smallmatrix}\right) :a\in \Rr\setminus\left\{0\right\}\right\}$ muni
de la loi de multiplication usuelle des matrices
de $\mathcal{M}_2(\Rr)$ est-il un groupe ?
    \item  L'ensemble $\mathcal{S}_2(\Rr) $ des matrices sym\'etriques
r\'eelles d'ordre 2 muni de la loi de
multiplication usuelle des matrices de
$\mathcal{M}_2(\Rr)$ est-il un groupe ?
\end{enumerate}
\finenonce{001317}



\finexercice



\exercice{1318, ridde, 1999/11/01}

\enonce{001318}{}
Soient $ (G, \star)$ et $ (H, \triangle)$ deux groupes. On d\'efinit sur $G \times H$
la loi $\heartsuit$ par $ (x, y) \heartsuit (x', y') = (x \star x', y \triangle y')$.
\begin{enumerate}
\item Montrer que $ (G \times H, \heartsuit)$ est un groupe.
\item Si $G$ est de cardinal 2, dresser la table de $G \times G$ et la reconna\^{\i}tre
parmi les exemples des exercices pr\'ec\'edents.
\end{enumerate}
\finenonce{001318}



\finexercice



\exercice{1319, ridde, 1999/11/01}

\enonce{001319}{}
Montrer que si $H$ et $K$ sont des sous-groupes de $G$ alors $H \cap K$ est un
sous-groupe de $G$. Est-ce vrai pour $H \cup K$ ?
\finenonce{001319}



\finexercice



\exercice{1320, ridde, 1999/11/01}

\enonce{001320}{}
Si $G$ est un groupe, on appelle centre de $G$ et on note $Z (G)$ l'ensemble
$\left\{ x \in G / \forall y \in G, xy = yx\right\}$.
\begin{enumerate}
\item Montrer que $Z (G)$ est un sous-groupe de $G$.
\item Montrer que $G$ est commutatif ssi $Z (G) = G$.
\item Calculer $Z (\sigma_3)$.
\end{enumerate}
\finenonce{001320}



\finexercice



\exercice{1321, legall, 1998/09/01}

\enonce{001321}{}
On nomme  $M_n({\Zz})$  l'ensemble des matrices
de taille  $n\times n$  \`a coefficients entiers relatifs.

- Soit  $M\in M_n({\Zz})$. Montrer que pour que  $M$  admette
un inverse \' el\' ement de  $M_n({\Zz})$  il faut et il suffit que  $\det (M)\in \{ -1,1\} $.

- D\' emontrer que  $Gl_n({\Zz})=\{ M\in M_n({\Zz})\ ; \det (M)\in \{ -1,1\} \} $  est
un sous-groupe de  $Gl_n({\Rr})$.
\finenonce{001321}



\finexercice



\exercice{1322, legall, 1998/09/01}

\enonce{001322}{}
\begin{enumerate}
    \item L'ensemble des matrices  $\begin{pmatrix}
a & c \cr
b & d \cr  \end{pmatrix} $  avec  $a,b,c,d\in {\Rr}$   tels que  $ad-bc\not = 0
$  et  $a^2-b^2-c^2-d^2 \leq 1$  est
il un sous-groupe de  $Gl_2({\Rr})$ ?
    \item L'ensemble des matrices  $\begin{pmatrix}
a & b \cr
0 & a^{-1} \cr \end{pmatrix}$  avec  $a\in {\Rr}^*$  et  $b\in {\Rr}$  est-il un sous groupe
de  $Gl_2({\Rr})$ ?
    \item Existe-t-il une valeur  $M\in {\Rr}$  telle que l'ensemble des matrices  $\begin{pmatrix}
a & c \cr
b & d \cr \end{pmatrix}$  avec  $a,b,c,d\in {\Rr}$  tels que   $ad-bc\not = 0 $  et  $
a\leq M$  forme un sous-groupe de   $Gl_2({\Rr})$ ?
\end{enumerate}

\finenonce{001322}



\finexercice



\exercice{1323, legall, 1998/09/01}

\enonce{001323}{}
 Soit $  G  $ un groupe, $  H  $ et $  K  $ deux
sous-groupes de $ G  .$ Montrer que $  H\cup K  $ est un
sous-groupe de $  G  $ si et seulement si $  H\subset K  $ ou $
K\subset H  .$

\finenonce{001323}



\finexercice



\exercice{1324, ortiz, 1999/04/01}

\enonce{001324}{}
D\'eterminer le sous-groupe de $\Zz$ engendr\'e
par les entiers 24, 36 et $-54.$
\finenonce{001324}



\finexercice



\exercice{1325, ortiz, 1999/04/01}

\enonce{001325}{}
Les questions sont ind\'ependantes. Soit $j$ le
nombre complexe $e^{\frac {2i\pi}{3}}.$
\begin{enumerate}
    \item  D\'eterminer le sous-goupe du groupe additif $\Cc$ engendr\'e par $i$ et $j$.
    \item  D\'eterminer le sous-goupe du groupe multiplicatif
$\Cc^{\ast}$ engendr\'e par $i$ et $j$.
\end{enumerate}
\finenonce{001325}



\finexercice



\exercice{1326, ortiz, 1999/04/01}

\enonce{001326}{}
 Soit $G$ un groupe engendr\'e par $a$ et $b.$
Montrer que $<a>\cap<b>\subseteq Z(G)$ o\`u $Z(G)$ d\'esigne le
centre de $G.$

\finenonce{001326}



\finexercice



\exercice{1327, gourio, 2001/09/01}

\enonce{001327}{}
Soit $G$ un sous-groupe de $(\Rr,+)$ avec  $G \neq \{0\}$.
\begin{enumerate}
\item  Montrer l'existence de $\alpha =\inf (G\cap \Rr^{+*})$.
\item Si $\alpha >0$ montrer que $G=\alpha \Zz.$
\item Si $\alpha =0$ montrer que $G$ est dense dans $\Rr.$
\end{enumerate}
\finenonce{001327}



\finexercice



\exercice{1328, gourio, 2001/09/01}

\enonce{001328}{}
Soit $G$ un groupe. Montrer que l'ensemble $Aut(G)$ des automorphismes de $G$
est un groupe pour la loi de composition. Soit $H $ un sous-groupe de $Aut(G)$,
 et $\pi :G\rightarrow \wp (G)$ d\'{e}finie par :
$\pi (x)=\{f(x)|f\in H\}.$
Montrer que $\pi (G)$ est une partition de $G.$
\finenonce{001328}



\finexercice



\exercice{1329, hilion, 2003/10/01}

\enonce{001329}{}
Soit $E$ un ensemble muni d'une loi interne $\star$. On appelle translation à droite (resp. à gauche) par $a\in E$, l'application $d_{a}$ (resp. $g_{a})$ de $E$ dans $E$ définie par $d_{a}(x)=a\star x$ (resp. $g_{a}(x)=x\star a$).
\begin{enumerate}
\item Montrer que dans un groupe les translations à droite et à gauche sont des bijections.
\item Réciproquement, si la loi $\star$ de $E$ est associative, et que les translations à droite et à gauche sont des bijections, on va montrer que $(E,\star)$ est un groupe.
\begin{enumerate}
\item Montrer que pour tout $x\in E$, il existe un unique élément $e_{x}\in E$ (resp. $f_{x}\in E$) tel que $e_{x}\star x=x$ (resp. $x\star f_{x}=x$).
\item Si $x,y \in E$, montrer que $e_{x}=e_{y}$ (noté $e$ dorénavant) et $f_{x}=f_{y}$ (noté $f$ dorénavant).
\item Montrer que $e=f$ (noté $e$ dorénavant).
\item  Montrer que pour tout $x\in E$, il existe un unique élément $\bar{x}\in E$ (resp. $\bar{\bar{x}}\in E$) tel que $\bar{x}\star x=e$ (resp. $x\star\bar{\bar{x}}=e$).
\item Montrer que $\bar{x}=\bar{\bar{x}}$.
\item Conclure.
\end{enumerate}
\end{enumerate}
\finenonce{001329}



\finexercice



\exercice{1330, hilion, 2003/10/01}

\enonce{001330}{}
Si $K$ est un sous-groupe de $H$ et $H$ un sous-groupe de $G$, montrer que $K$ est un sous-groupe de $G$.
\finenonce{001330}



\finexercice



\exercice{1331, hilion, 2003/10/01}

\enonce{001331}{}
\begin{enumerate}
\item Soit $(G,.)$ un groupe. Montrer l'équivalence de:
\begin{description}
\item{i)} $G$ est abélien.
\item{ii)} Pour tout $a,b\in G$, on a: $(ab)^{2}=a^{2}b^{2}$.
\item{iii)} Pour tout $a,b\in G$, on a: $(ab)^{-1}=a^{-1}b^{-1}$.
\item{iv)} L'application $f$ de $G$ dans $G$ définie par $f(x)=x^{-1}$ est un automorphisme.
\end{description}
\item En déduire que si pour tout $x\in G$, $x^{2}=e$, alors $G$ est abélien.
\end{enumerate}
\finenonce{001331}



\finexercice



\exercice{1332, hilion, 2003/10/01}

\enonce{001332}{}
\begin{enumerate}
    \item Les ensembles $\Nn,\Zz,\Rr, \Rr_{+}, \Rr_{+}^{*}, \Cc, \Cc^{*}$ munis des lois $+$ ou $\times$ sont-ils des groupes ? Quand c'est le cas, chercher des sous-groupes non triviaux.
    \item $\left\{x\in\Rr\mapsto ax+b : a\in\Rr\setminus\left\{0\right\},b\in\Rr\right\}$ muni de la loi de composition des applications est-il un groupe ?
\end{enumerate}
\finenonce{001332}



\finexercice



\exercice{1333, hilion, 2003/10/01}

\enonce{001333}{}
Quel est le plus petit sous-groupe de $(\Rr,+)$ (resp. de $(\Rr^*,\times)$) contenant 1 ? Contenant 2 ?
\finenonce{001333}



\finexercice



\exercice{1334, hilion, 2003/10/01}

\enonce{001334}{}
Soit $\lambda\in\Cc$ fixé. Montrer que $S_{\lambda}=\{ \exp(i\lambda t): t\in \Rr \}$ est un sous-groupe de $(\Cc,\times)$.
Pour quelles valeurs de $\lambda$ retrouve-t-on des sous-groupes bien connus? A quoi ressemblent les courbes $S_{\lambda}$ ?
Que peut-on dire, en terme de morphisme, de l'application $t\mapsto\exp(i\lambda t)$ ?
\finenonce{001334}



\finexercice



\exercice{1343, legall, 1998/09/01}

\enonce{001343}{}
 D\' ecrire tous les homomorphismes de groupes de
${\Zz}$ dans  ${\Zz}$. D\' eterminer ceux qui sont injectifs et
ceux qui sont surjectifs.

\finenonce{001343}




\finexercice



\exercice{1344, ortiz, 1999/04/01}

\enonce{001344}{}
Pour tout couple $(a,b)$ de $\Rr^2$, on pose la matrice
$M_{a,b}=\left(\begin{smallmatrix}
a&-b\\b&a\end{smallmatrix}\right).$
Soit $\mathcal{S}=\left\{M_{a,b}:(a,b)\in
\Rr^2\setminus
\left\{(0,0)\right\}\right\}$. Soit l'application $f: \mathcal{S}
\to\Rr, M_{a,b}\mapsto a^2+b^2$.
\begin{enumerate}
    \item  Montrer que $\mathcal{ S }$ est un groupe pour
la loi usuelle de multiplication des matrices
carr\'ees.
    \item  Montrer que $f$ est un morphisme du groupe
$\left(\mathcal{ S },\times\right)$ dans le groupe
multiplicatif $\Rr\setminus\left\{(0,0)\right\}$.
\end{enumerate}
\finenonce{001344}



\finexercice



\exercice{1345, legall, 1998/09/01}

\enonce{001345}{}
 Soit  $f:{\Rr}\rightarrow {\Cc}^*$  l'application
qui \`a tout $x\in {\Rr}$  associe  $e^{ix}\in {\Cc}^*$. Montrer
que $f$ est un homomorphisme de groupes. Calculer son noyau et son
image. $f$ est-elle injective ?

\finenonce{001345}



\finexercice



\exercice{1346, ortiz, 1999/04/01}

\enonce{001346}{}
Traduire en termes d'homomorphisme de groupes les
propri\'et\'es traditionnelles suivantes :
\begin{enumerate}
    \item $\ln(xy)=\ln x+\ln y$ ;
    \item $\det(MM')= \det(M)\det(M')$ ;
    \item $|{zz'}|=|{z}||{z'}|$ ;
    \item $(xy)^{\frac 12}=x^{\frac 12}y^{\frac 12}$ ;
    \item $e^{z+z'}= e^{z}e^{z'}$ ;
    \item $\overline{z+z'}=\overline{z}+\overline{ z'}$.
\end{enumerate}
\finenonce{001346}



\finexercice



\exercice{1347, ortiz, 1999/04/01}

\enonce{001347}{}
Pour tout couple $(a,b)$ de $\Rr^2$, on pose
$M_{a,b}=\left(\begin{smallmatrix}
a&-b\\b&a\end{smallmatrix}\right),$
$\mathcal{ S }=\left\{M_{a,b}:(a,b)\in
\Rr^2\right\}$ et $\mathcal{ S}^{\ast}=\mathcal{ S}\setminus\left\{M_{0,0}\right\}.$
 Soit l'application $f:
\mathcal{ S }\to\Cc, M_{a,b}\mapsto a+ib$.
\begin{enumerate}
    \item
    \begin{enumerate}
          \item Montrer que $\mathcal{ S }$ est un
sous-groupe du groupe additif  usuel
$\mathcal{M}_2(\Rr).$
          \item Montrer que $\mathcal{ S }^{\ast}$ est un
sous-groupe multiplicatif de ~~$\text{GL}_2(\Rr).$

    \end{enumerate}
    \item  Montrer que $f$ est un isomorphisme du groupe
$\left(\mathcal{ S },+\right)$ sur le groupe
additif ~$\Cc$.
    \item
    \begin{enumerate}
         \item Montrer que $f$ d\'efinit un homomorphisme du groupe
$\left(\mathcal{ S }^{\ast},\times\right)$ sur le
groupe multiplicatif~$\Cc^{\ast}$.
         \item D\'eterminer le noyau et l'image de cet homomorphisme.
    \end{enumerate}
    \item  Montrer que $\Omega=\left\{M_{a,b}:(a,b)\in \Rr^2,a^2+b^2=1\right\}$ est
    un sous-groupe distingu\'e du groupe multiplicatif $\mathcal{ S }^{\ast}$.
\end{enumerate}
\finenonce{001347}



\finexercice



\exercice{1348, gourio, 2001/09/01}

\enonce{001348}{}
Soit $G$ un groupe.\ Montrer que l'application $x\rightarrow x^{-1}$ est un
morphisme si et seulement si $G$ est commutatif.\ On suppose $G$ fini; soit $\phi$
un morphisme involutif de $G $ dont le seul point fixe est $e$,
montrer que :
$$\forall z\in G,\exists t\in G,z=t(\phi (t))^{-1}.$$
En d\'{e}duire $\phi$ puis que $G$ est commutatif.
\finenonce{001348}



\finexercice



\exercice{1349, legall, 1998/09/01}

\enonce{001349}{}
Montrer que les groupes $  (\R ,+)  $ et $  (\R ^*_+ ,\times )  $
sont isomorphes.
\finenonce{001349}



\finexercice



\exercice{1350, ridde, 1999/11/01}

\enonce{001350}{}
Montrer que $\mathbb{U}_2 \times \mathbb{U}_3$ est isomorphe \`a $\mathbb{U}_6$.
Est-ce que $\mathbb{U}_2 \times \mathbb{U}_2 $ est isomorphe est $\mathbb{U}_4$ ?
Pouvez-vous conjecturer \`a quelle condition $\mathbb{U}_n \times \mathbb{U}_m$
est isomorphe \`a $\mathbb{U}_{nm}$ ?
\finenonce{001350}



\finexercice



\exercice{1351, ortiz, 1999/04/01}

\enonce{001351}{}
Soit $G$ un groupe.
\begin{enumerate}
    \item  Montrer que l'ensemble des automorphismes de $G$ muni
de la loi de composition des applications est un
groupe. Ce groupe est not\'e $\text{Aut}\,(G).$
    \item  V\'erifier que l'application $\phi:G\to \text{Aut}\,(G)$
qui associe \`a $g\in G$ l'application
$\phi_g:G\to G, x\mapsto gxg^{-1}$ est un
morphisme de groupes. D\'eterminer son noyau
$Z(G),$ dit \textit{centre} de $G.$
    \item  D\'eterminer $\text{Aut}\,(\Q)$ et $\text{Aut}\,(\Zz).$
\end{enumerate}
\finenonce{001351}



\finexercice



\exercice{1352, hilion, 2003/10/01}

\enonce{001352}{}
Soit $(G,.)$ un groupe. On appelle conjugaison par $a\in G$, l'application $f_{a}$ de $G$ dans $G$ définie par $f_{a}(x)=a.x.a^{-1}$.
\begin{enumerate} 
\item Montrer que $f_{a}$ est un automorphisme de $G$.
\item Soit $\Gamma=\{f_{a}:a\in G\}$. Montrer que $(\Gamma,\circ)$ est un groupe.
\item Soit $\Phi:G\rightarrow\Gamma,a\mapsto f_{a}$. Vérifier que $\Phi$ est un morphisme.
Est-il injectif ? (indication: préciser ce morphisme lorsque $G$ est abélien).
\end{enumerate}
\finenonce{001352}



\finexercice



\exercice{1353, ortiz, 1999/04/01}

\enonce{001353}{}
\begin{enumerate}
    \item  Les sous-groupes $(\Q,+)$ et $(\Zz,+)$ sont-ils isomorphes ?
    \item  Les sous-groupes $(\Q,+)$ et $(\Q\setminus\left\{0\right\},\times)$ sont-ils isomorphes ?
\end{enumerate}
\finenonce{001353}



\finexercice



\exercice{1354, ortiz, 1999/04/01}

\enonce{001354}{}
 Montrer que les groupes multiplicatifs
$\Rr\setminus\left\{0\right\}$ et $\Cc\setminus\left\{0\right\}$ ne sont pas
isomorphes.

\finenonce{001354}



\finexercice



\exercice{1385, cousquer, 2003/10/01}

\enonce{001385}{}
\begin{enumerate}
    \item On suppose que $\varphi$ est un isomorphisme de $(G,\ast)$ sur
$(G',\diamond)$. Si $e$ est l'élément neutre de $G$, que peut-on dire de
$\varphi(e)$~? Si $x'$ est l'inverse de $x$ dans $G$, que peut-on dire de
$\varphi(x')$~? Si $G$ est d'ordre $n$, que peut-on dire de l'ordre de $G'$~?
    \item Pouvez-vous citer des exemples de groupes~? de groupes isomorphes~?
    \item Si $(G,\ast)$ est un groupe fini et si on établit la table de la loi
$\ast$, peut-on rencontrer deux fois le même élément dans la même ligne,
dans la même colonne~? Établir les tables de composition possibles pour des
groupes à $2$, $3$, $4$ éléments. Pouvez-vous donner des exemples de groupes
correspondant à ces tables. Retrouver éventuellement des groupes isomorphes.
\end{enumerate}
\finenonce{001385}



\finexercice



\exercice{1386, legall, 1998/09/01}

\enonce{001386}{}
 Soient $  p  $ un nombre premier et $  G  $ un
groupe d'ordre $  p .$ Montrer que $  G  $ est cyclique et donner
la liste des g\'en\'erateurs de $  G  .$

\finenonce{001386}



\finexercice



\exercice{1387, legall, 1998/09/01}

\enonce{001387}{}
 Soit $  G  $ un groupe d'ordre $  pn  $ avec $  p
$ premier.
\begin{enumerate}
    \item On consid\`ere deux sous-groupes $  H  $ et $  H'  $ de $  G  $ d'ordre $  p  $ avec $  H\not = H'  .$ Montrer que $  H\cap H'=\{ e\}  .$
    \item En d\'eduire que le nombre d'\'el\'ements d'ordre $  p   $ dans $  G  $ est un multiple de $  p-1  .$
\end{enumerate}

\finenonce{001387}



\finexercice



\exercice{1388, legall, 1998/09/01}

\enonce{001388}{}
D\'eterminer (\`a isomorphisme pr\`es) tous les groupes d'ordre $  4  .$
\finenonce{001388}



\finexercice



\exercice{1389, legall, 1998/09/01}

\enonce{001389}{}
\begin{enumerate}
    \item Soit $  G  $ un groupe dans lequel tout \'el\'ement (distinct de
l'\'el\'ement neutre) est d'ordre $  2  .$ Montrer que $  G  $ est commutatif.
    \item Soit $  G  $ un groupe d'ordre pair. Montrer que $  G  $ contient au moins un \'el\'ement d'ordre $  2  .$
\end{enumerate}

\finenonce{001389}



\finexercice



\exercice{1390, ortiz, 1999/04/01}

\enonce{001390}{}
Montrer que tout morphisme de
groupes de $\Q$ dans un  groupe fini $G$ est
trivial.
\finenonce{001390}



\finexercice



\exercice{1391, ortiz, 1999/04/01}

\enonce{001391}{}
Soit $G$ un groupe et $H$ une partie finie non
vide de $G.$ On suppose que $H$ est stable pour la
loi de $G.$ Montrer $H$ est un sous-groupe de $G.$
\finenonce{001391}



\finexercice



\exercice{1392, gourio, 2001/09/01}

\enonce{001392}{}
Soit $G$ un groupe fini de cardinal $2n $ ($n\geq 2$), poss\'{e}dant 2
sous-groupe $H$ et $H^{\prime } $ tels que :
$$\mathrm{Card}(H)=\mathrm{Card}(H^{\prime })=n$$ et
$$H\cap H^{\prime }=\{e\}.$$
\begin{enumerate}
\item Montrer que $G-(H\cup H^{\prime })$ est un singleton, not\'{e} $\{a\}.$
\item Soit $h\in H-\{e\}$, montrer que $hH^{\prime }\subset \{h,a\}$, en
d\'{e}duire que $hH^{\prime }=\{h,a\} $ puis que $n=2.$
\item On \'{e}crit $G=\{a,e,h,h^{\prime }\}, $ donner la table de $G $ (puis un
exemple d'un tel groupe).
\end{enumerate}
\finenonce{001392}



\finexercice



\exercice{2438, matexo1, 2002/02/01}

\enonce{002438}{}
Soit $G$ l'ensemble des matrices d'ordre 2 inversibles.

\begin{enumerate}
\item Montrer que $G$ est un groupe pour le produit matriciel.
Est-il commutatif\,?
\item Montrer que si $A\in G, B \in G$ v\'erifient $A^2 = B^2 =
ABAB = I$, alors $A = A^{-1}, B = B^{-1}$ et $AB = BA$.
\item Trouver deux \'el\'ements de $G$ v\'erifiant $A^2 = B^2 =
I$ et $AB \neq BA$.
\end{enumerate}
\finenonce{002438}


\finexercice


\exercice{2439, matexo1, 2002/02/01}

\enonce{002439}{}
Soit $G$ l'ensemble des matrices de la forme
$$
\left(
\begin{array}{cc}
1 & x \\ 0 & 1
\end{array}
\right), \qquad x \in \R.$$
Montrer que c'est un groupe multiplicatif isomorphe au groupe
additif r\'eel.
\finenonce{002439}


\finexercice


\exercice{2440, matexo1, 2002/02/01}

\enonce{002440}{}
Soit $G$ le groupe multiplicatif des matrices complexes
d'ordre $n$. Parmi les sous-ensembles suivants de $G$, lesquels
sont des sous-groupes ?

\begin{itemize}
\item[---] les matrices \`a coefficients r\'eels\,;
\item[---] les matrices inversibles\,;
\item[---] les matrices r\'eelles inversibles \`a coefficients
positifs\,;
\item[---] les matrices diagonales inversibles\,;
\item[---] les matrices v\'erifiant $a_{i,i} \neq 0, \forall i$,  et
triangulaires sup\'erieures ($a_{i,j} = 0$ si $i > j$).
\item[---] les matrices v\'erifiant $a_{i,j}= \overline{a_{j,i}},
\forall i,j$.
\end{itemize}
\finenonce{002440}


\finexercice


\exercice{2727, tumpach, 2009/10/25}
\enonce{002727}{}
On munit l'ensemble $G = \{a, b, c, d\}$ d'une loi de composition interne dont la table de Pythagore est
$$
\begin{tabular}{r|rrrr}
$\star$ & \textbf{a} & \textbf{b}&\textbf{c}&\textbf{d}\\
\hline
\textbf{a} & c & a & c & a\\
\textbf{b}& a & d& c& b\\
\textbf{c}& c & c& c & c\\
\textbf{d} & a & b & c & d 
\end{tabular}
$$
(La premi\`ere ligne se lit $a \star a = a$, $a \star b = a$, $a \star c = c$,\ldots.) 
\begin{enumerate}
\item Cette loi poss\`ede-t-elle un \'el\'ement neutre~?
\item Cette loi est-elle commutative~?
\item Cette loi est-elle associative~?
\item Est-ce une loi de groupe~?
\end{enumerate}
\finenonce{002727}



\finexercice


\exercice{2728, tumpach, 2009/10/25}
\enonce{002728}{}
On d\'efinit une permutation $\sigma$ de l'ensemble $\{1, 2, \dots, 15\}$ par la suite finie des entiers $\sigma(1)$, $\sigma(2), \dots, \sigma(15)$. Par exemple
$$
\sigma_1 = \left(
\begin{array}{ccccccccccccccc}
1 & 2 & 3 & 4 & 5 & 6 & 7 & 8 & 9 & 10 & 11 & 12 & 13 & 14 & 15\\
2 & 7 & 1 & 14 & 3 & 12 & 8 & 9 & 6 & 15 & 13 & 4 & 10 & 5 & 11
\end{array}
\right)
$$
signifie $\sigma(1) = 2, \sigma(2) = 7$, etc\dots Soient
 $$
\sigma_2 = \left(
\begin{array}{ccccccccccccccc}
1 & 2 & 3 & 4 & 5 & 6 & 7 & 8 & 9 & 10 & 11 & 12 & 13 & 14 & 15\\
7 & 6 & 5 & 8 & 9 & 3 & 2 & 15 & 4 & 11 & 13 & 10 & 12 & 14 & 1
\end{array}
\right)
$$
$$
\sigma_3 = \left(
\begin{array}{ccccccccccccccc}
1 & 2 & 3 & 4 & 5 & 6 & 7 & 8 & 9 & 10 & 11 & 12 & 13 & 14 & 15\\
1 & 15 & 2 & 14 & 3 & 13 & 4 & 12 & 5 & 11 & 6 & 10 & 7 & 9 & 8
\end{array}
\right)
$$
$$
\sigma_4 = \left(
\begin{array}{ccccccccccccccc}
1 & 2 & 3 & 4 & 5 & 6 & 7 & 8 & 9 & 10 & 11 & 12 & 13 & 14 & 15\\
2 & 4 & 6 & 8 & 10 & 12 & 14 & 15 & 13 & 11 & 9 & 7 & 5 & 3 & 1
\end{array}
\right)
$$
\begin{enumerate}
\item Pour  $i=1, \dots, 4$,
\begin{itemize}
\item  d\'ecomposer $\sigma_i$ en cycles \`a supports disjoints.
\item  d\'eterminer l'ordre de $\sigma_i$.
\item  d\'eterminer la signature de $\sigma_i$.
\end{itemize}
\item Calculer les puissances successives du cycle $\sigma = (10\,\,\,15\,\,\,11\,\,\,13)$. Quel est l'inverse de $\sigma_1$~?
\item Calculer $\sigma_2^{2008}$.
\item D\'eterminer, sans fatigue excessive, la signature de 
$$
\sigma_3\circ \sigma_4 \circ \sigma_{3}^{-4} \circ \sigma_{4}^{3}\circ \sigma_3\circ \sigma_4 \circ \sigma_3 \circ \sigma_4 \circ \sigma_3^{-1} \circ \sigma_4^{-6}.
$$
\item Combien y a-t-il de permutations $g$ de $\{1, \dots, 15\}$ telles que $\sigma_1\circ g = g\circ \sigma_1$~?
\end{enumerate}
\finenonce{002728}



\finexercice


\exercice{2729, tumpach, 2009/10/25}
\enonce{002729}{}
\begin{enumerate}
\item
Montrer que les ensembles $G$ suivants, munis des lois $\star$ donn\'ees, sont des groupes. Pr\'eciser quel est l'\'el\'ement neutre de $G$ et quel est  l'inverse d'un \'el\'ement quelconque $x\in G$.
\begin{enumerate}
\item $G = \mathbb{Z}$, $\star = $ l'addition des nombres ;
\item $G = \mathbb{Q}^{*}$ (ensemble des rationnels non nuls), $\star = $la multiplication des nombres  ;
\item $G = \mathbb{Q}^{+*}$ (ensemble des rationnels strictement positifs),  $\star = $ la multiplication des nombres  ;\item  $G = \mathbb{R}$, $\star = $ l'addition des nombres ;
\item $G = \mathbb{R}^{+*}$, $\star = $ la multiplication des nombres ;
\item $G = \mathbb{C}$, $\star = $ l'addition des nombres ;
\item $G = \mathbb{C}^{*}$, $\star = $ la multiplication des nombres ;
\item $G = \{z \in\mathbb{C}, |z| = 1\}$, $\star = $  la multiplication des nombres  ;
\item $G = \{e^{i\frac{2\pi k}{n}}, k = 0, 1, \dots, n-1\}$, $\star = $ la multiplication des nombres ($n$ est un entier fix\'e) ;
\item $G = \textrm{l'ensemble des bijections d'un ensemble non vide} \,\,\,E$, $\star =$ la composition des applications ;
\item $G$ = l'ensemble des isom\'etries de l'espace euclidien $\mathbb{R}^{3}$ (muni du produit scalaire standard), $\star =$ la composition des applications ;
\item $G$ = l'ensemble des isom\'etries du plan euclidien $\mathbb{R}^{2}$ (muni du produit scalaire standard) qui pr\'eservent une figure donn\'ee, $\star =$ la composition des applications.
\end{enumerate}


\item Donner un morphisme de groupes entre $(\mathbb{R}, +)$ et $(\mathbb{R}^{+*}, \cdot)$ ;
\item Donner un morphisme de groupes entre $(\mathbb{R}^{+*}, \cdot)$ et $(\mathbb{R}, +)$ ;
\item Donner un morphisme de groupes surjectif entre $(\mathbb{C}, +)$ et $(\mathbb{C}^*, \cdot)$ ; 
\end{enumerate}
\finenonce{002729}



\finexercice


\exercice{2730, tumpach, 2009/10/25}
\enonce{002730}{}
Dire pour quelle(s) raison(s) les op\'erations $\star$ suivantes ne munissent pas les ensembles $G$ donn\'es d'une structure de groupe~?
\begin{enumerate}
\item[(a)] $G = \mathbb{N}$, $\star = $ l'addition des nombres ;
\item[(b)] $G = \mathbb{N}^{*}$, $\star = $ la multiplication des nombres ;
\item[(c)] $G =\mathbb{R}$, $\star = $ la multiplication des nombres ;

\end{enumerate}
\finenonce{002730}



\finexercice


\exercice{2965, quercia, 2010/03/08}

\enonce{002965}{Groupe produit}
Soient $G,H$ deux groupes multiplicatifs. On munit $G\times H$ de l'op{\'e}ration :
$$\forall\ g,g'\in G,\ \forall\ h,h'\in H,  \qquad (g,h)\cdot(g',h') = (gg',hh').$$
\\
Montrer que $\cdot$ d{\'e}finit une loi de groupe sur $G\times H$.
\finenonce{002965}



\finexercice 


\exercice{2966, quercia, 2010/03/08}

\enonce{002966}{Essai de tables}
Les op{\'e}rations suivantes sont-elles des lois de groupe ?

\begin{enumerate}
  \item
    $$\begin{array}{c|ccc}
      &a &b &c \\
      \hline
      a &a &a &a \\
      b &a &b &b \\
      c &a &b &c 
    \end{array}
    $$
  \item
    $$\begin{array}{c|ccc}
      &a &b &c \\
      \hline
      a &b &c &a \\
      b &c &a &b \\
      c &a &b &c 
    \end{array}
    $$
  \item
    $$\begin{array}{c|cccc}
      &a &b &c &d \\
      \hline
      a &a &b &c &d \\
      b &b &a &d &c \\
      c &d &c &b &a \\
      d &c &d &a &b 
    \end{array}
    $$
\end{enumerate}
\finenonce{002966}


\finexercice 


\exercice{2967, quercia, 2010/03/08}

\enonce{002967}{Translations surjectives}
Soit $G$ un ensemble non vide muni d'une op{\'e}ration interne $\cdot$ associative telle que :
$$\forall\ a,b \in G,\ \exists\ x,y \in G \text{ tq } a = x\cdot b = b\cdot y.$$

Montrer que $(G,\cdot)$ est un groupe.
\finenonce{002967}



\finexercice 


\exercice{2968, quercia, 2010/03/08}

\enonce{002968}{Transport de structure}
Soit $G$ un groupe multiplicatif, $E$ un ensemble, et $\phi : G \to E$ une
bijection.

On d{\'e}finit une op{\'e}ration $\star$ sur $E$ par :
$$\forall\ x,y \in E,\ x \star y = \phi\Bigl( \phi^{-1}(x)\phi^{-1}(y) \Bigr).$$

Montrer que $\star$ est une loi de groupe et que les groupes $G$ et $E$ sont
isomorphes.
\finenonce{002968}



\finexercice 


\exercice{2969, quercia, 2010/03/08}

\enonce{002969}{Transport de structure} Pour tout $x,y \in \R$, on pose $x
\star y = x \sqrt{1+y^2} + y \sqrt{1+x^2}$.
\begin{enumerate}
  \item V{\'e}rifier que $\sqrt{1+(x\star y)^2} = \sqrt{1+x^2}\sqrt{1+y^2} + xy$.
  \item Montrer que $(\R, \star)$ est un groupe.
  \item Montrer que l'application $\sh$ est un isomorphisme entre $(\R,+)$ et
    $(\R, \star)$.
\end{enumerate}
\finenonce{002969}



\finexercice 


\exercice{2970, quercia, 2010/03/08}

\enonce{002970}{Transport de structure} Pour tout $x,y \in \R$, on pose $x
\star y = \root{\scriptstyle 3}\of {x^3+y^3}$.  Montrer que
$(\R,\star)$ est un groupe isomorphe {\`a} $(\R,+)$. 
\finenonce{002970}



\finexercice 


\exercice{2971, quercia, 2010/03/08}

\enonce{002971}{Loi associative r{\'e}guli{\`e}re}
Soit $E$ un ensemble fini muni d'une op{\'e}ration interne $*$ associative pour laquelle
tout {\'e}l{\'e}ment est r{\'e}gulier {\`a} droite et {\`a} gauche.
Montrer que $E$ est un groupe.
\finenonce{002971}



\finexercice


\exercice{2972, quercia, 2010/03/08}

\enonce{002972}{Partie finie stable par produit}
Soit $G$ un groupe multiplicatif et $H$ une partie finie de $G$ non vide,
stable par multiplication.
Montrer que $H$ est un sous-groupe de $G$.
\finenonce{002972}



\finexercice 


\exercice{2973, quercia, 2010/03/08}

\enonce{002973}{Centre d'un groupe et commutant}
Soit $G$ un groupe multiplicatif.
On note $Z(G) = \{ a \in G$ tq $\forall\ b \in G,$ on a $ab = ba\}$
({\it centre de $G$\/}), et
pour ${a \in G\colon}$ $C(a) = \{ b \in G$ tq $ab = ba\}$
({\it commutant de $a$\/}).

Montrer que $Z(G)$ et $C(a)$ sont des sous-groupes de $G$.
\finenonce{002973}



\finexercice 


\exercice{2974, quercia, 2010/03/08}

\enonce{002974}{ Loi $\Delta$}
Soit $E$ un ensemble et $G = {\cal P}(E)$.
\begin{enumerate}
  \item Montrer que $(G,\Delta)$ est un groupe commutatif.
  \item Pour $a \in E$, on note
    ${\phi_a} : G \to {\Z/2\Z}$ d{\'e}finie par $\begin{cases} X &{\dot 0 \text{ si } a \notin X}\cr
                                   X &{\dot 1 \text{ si } a \in X}\cr\end{cases}$.
    \par
    Montrer que $\phi_a$ est un morphisme de groupes.

  \item On prend $E = \{1,\dots, n\}$ et on note
    \begin{align*}
      \Phi : G &\to {\bigl({\Z/2\Z}\bigr)^n}, \\
      X &\mapsto (\phi_1(X),\dots,\phi_n(X)) .
    \end{align*}
    Montrer que $\Phi$ est un isomorphisme de groupes.
\end{enumerate}
\finenonce{002974}



\finexercice


\exercice{2975, quercia, 2010/03/08}

\enonce{002975}{Sous-groupes embo{\^\i}t{\'e}s}
Soit $G$ un groupe additif, et $H,K,L$ trois sous-groupes de $G$ v{\'e}rifiant :
$H \subset K$, $H \cap L = K \cap L$, $H + L = K + L$.
D{\'e}montrer que $H = K$.
\finenonce{002975}



\finexercice 


\exercice{2976, quercia, 2010/03/08}

\enonce{002976}{$\mathrm{Card}\,(HK)$}
Soit $G$ un groupe fini et $H,K$ deux sous-groupes de $G$.
On consid{\`e}re l'application
$\phi : {H\times K} \to  G, {(h,k)} \mapsto {hk}$

\begin{enumerate}
  \item Est-ce que $\phi$ est un morphisme de groupes ?
  \item Soit $z \in HK$, $z=h_0k_0$ avec $h_0 \in H$ et $k_0 \in K$.
    \par
    Montrer que les ant{\'e}c{\'e}dents de $z$ par $\phi$ sont les couples
    $(h_0t,t^{-1}k_0)$ avec $t \in H \cap K$.
  \item En d{\'e}duire que : $\mathrm{Card}\,(HK)\mathrm{Card}\,(H\cap K) = \mathrm{Card}\,(H)\mathrm{Card}\,(K)$.
  \item Montrer que : ($HK$ est un sous-groupe de $G$) $\iff$ ($HK \subset KH$)
    $\iff$ ($HK = KH$).
\end{enumerate}
\finenonce{002976}



\finexercice 

\exercice{2978, quercia, 2010/03/08}

\enonce{002978}{Sous-groupes d'un groupe cyclique}
Soit $n \in \N^*$ et $G = \Z/n\Z$. Soit $k \in \Z$ et
$d = k \wedge n$.
\begin{enumerate}
  \item D{\'e}terminer l'ordre de $\dot k$ dans $G$.
  \item Montrer que $\dot k$ et $\dot d$ engendrent le m{\^e}me sous-groupe de $G$.
  \item Quels sont tous les sous-groupes de $G$ ?
\end{enumerate}
\finenonce{002978}



\finexercice 

\exercice{7290, mourougane, 2021/08/10}
\enonce{007290}{}
\begin{enumerate}
 \item Soit $(G,\star)$ un groupe. Donner la définition du sous-groupe engendré par une partie $P$ de $G$.
 \item Soit $(G,\star)$ un groupe. Décrire le sous-groupe engendré par un élément $g$ de $G$.
 \item Décrire le sous-groupe engendré par un élément $n$ de $(\Z,+)$.
\end{enumerate}
\finenonce{007290}
\finexercice


\exercice{7291, mourougane, 2021/08/10}
\enonce{007291}{}
\begin{enumerate}
 \item Dire si les couples suivants sont des groupes : $(\Z,+)$ ; $(\Z,\times)$ ; $(\C^*,+)$ ; $(\C^*, \times)$.
 \item Déterminer une loi de composition interne pour laquelle $\{-1,1\}$ est un groupe.
Même question avec $\R\times\R$.
 \item Soit $(G,\star)$ un groupe et $E$ un ensemble.
Montrer que l'ensemble $\mathcal{A}(E,G)$ des applications
de $E$ dans $G$ peut être muni d'une structure naturelle de groupe. 
 \item L'ensemble $\{a,b,c\}$ muni de la loi de composition interne définie par la table
$$\begin{array}{|c|c|c|c|}\hline
 \star& a&b&c\\ \hline
 a&a&b&c\\ \hline
b&b&a&b\\ \hline
c& c& b&a\\\hline
\end{array}
$$
admet-il un élément neutre ? est-il commutatif ?
tout élément admet-il un symétrique ?
est-il un groupe ? 
 \item La réunion de deux sous-groupes d'un groupe $G$ est-elle un sous-groupe ? Et l'intersection ?
 \item Donner l'exemple d'un ensemble avec une loi de composition interne sans élément neutre.
\end{enumerate}
\finenonce{007291}
\finexercice


\exercice{7292, mourougane, 2021/08/10}
\enonce{007292}{}
Soit $(E,\star)$ un ensemble muni d'une loi interne
associative avec un élément neutre, noté $\varepsilon$. 
\begin{enumerate}
 \item Démontrer qu'il n'y a pas d'autre élément neutre.
 \item Soit $x$ un élément qui admet un symétrique $x'$.
Démontrer qu'il n'y a pas d'autre symétrique.
 \item Soit $x$ et $y$ dans
$E$ qui admettent un symétrique. Déterminer un symétrique de $x\star y$. 
\end{enumerate}
\finenonce{007292}
\finexercice


\exercice{7293, mourougane, 2021/08/10}
\enonce{007293}{}
\begin{enumerate}
% \item Soit $H$ un sous-groupe de $(\Z,+)$ contenant $\{2,7\}$. Montrer que $1$ est dans $H$.
 \item Montrer que $1$ est dans le sous-groupe de $(\Z,+)$ engendré par $\{2,7\}$.
 \item Quel est le sous-groupe engendré par $\{2,7\}$ ?
\end{enumerate}
\finenonce{007293}
\finexercice


\exercice{7294, mourougane, 2021/08/10}
\enonce{007294}{}
Soit $a$ et $b$ deux entiers non tous nuls et $<a,b>$ le sous-groupe de $(\Z,+)$ engendré par $\{a,b\}$.
\begin{enumerate}
 \item Montrer que $a\Z$ est inclus dans $<a,b>$.
 \item Montrer que $a\Z+b\Z$ est inclus dans $<a,b>$.
 \item Montrer que $a\Z+b\Z$ est un sous-groupe de $(\Z,+)$, contenant $\{a,b\}$.
 \item Montrer que $<a,b>$ est inclus dans $a\Z+b\Z$.
 \item Montrer que $<a,b>$ est $a\Z+b\Z$.
 \item Montrer que $<a,b>=a\Z+b\Z=\pgcd(a,b)\Z$.
\end{enumerate}
\finenonce{007294}
\finexercice


\exercice{7295, mourougane, 2021/08/10}
\enonce{007295}{}
\begin{enumerate}
 \item Déterminer $45\Z \cap 60\Z$.
 \item Déterminer $56\Z+63\Z$.
 \item Trouver les sous-groupes de $\Z$ contenant $48\Z$ et donner les relations d'inclusion existant entre eux.
\end{enumerate}
\finenonce{007295}
\finexercice


\exercice{7296, mourougane, 2021/08/10}
\enonce{007296}{}
\begin{enumerate}
 \item L'application $F~:(\Z,+)\to (\Z,+)$, $n\mapsto 1$ est-elle un morphisme de groupes ?
 \item Déterminer deux morphismes de groupes de $(\R,+)$ vers $(\R_+^\star,\times)$.
\end{enumerate}
\finenonce{007296}
\finexercice


\exercice{7297, mourougane, 2021/08/10}
\enonce{007297}{}
\begin{enumerate}
 \item Soit $F$ un morphisme de groupes de $(\Z,+)$ dans lui-même.
Montrer qu'il suffit de connaître $F(1)$ pour connaître l'image de chaque entier par $F$.
 \item Existe-t-il un morphisme de groupes de $(\Z,+)$ dans lui-même tel que $F(2)=3$ ?
 \item Déterminer tous les morphismes de groupes de $(\Z,+)$ dans lui-même tels que $F(2)=~4$. 
 \item Montrer qu'une condition nécessaire et suffisante sur $(a,b)\in \Z^2$ pour qu'il existe un morphisme de groupes de $(\Z,+)$ dans lui-même tel que $F(2)=a$ et $F(5)=b$ est : $5a=2b$.
\end{enumerate}
\finenonce{007297}
\finexercice


\exercice{7298, mourougane, 2021/08/10}
\enonce{007298}{}
Soit $E=\{a,b\}$ un ensemble à deux éléments.
On cherche à déterminer toutes les structures de groupes sur cet
ensemble. 
\begin{enumerate}
 \item Si $a$ est l'élément neutre, quel doit être le symétrique de $b$ ?
En déduire la table de multiplication dans ce cas.
 \item Combien y a-t-il de structures de groupes différentes sur $E$ ?
\end{enumerate}
\finenonce{007298}
\finexercice


\exercice{7299, mourougane, 2021/08/10}
\enonce{007299}{}
Soit $E = \Q^2 \setminus \{(0,0)\}$, et $\star$ la loi de composition interne définie par :
$$ (x,y) \star (x',y') = (xx'-yy',xy'+yx').$$
$(E,\star)$ est-il un groupe ? commutatif ?
\finenonce{007299}
\finexercice


\exercice{7300, mourougane, 2021/08/10}
\enonce{007300}{}
Soit $E$ un ensemble, on note $\mathcal{P}(E)$ l'ensemble des parties de $E$. Si $A, B \in \mathcal{P}(E)$, on note $A\Delta B = (A \cup B) \setminus (A \cap B)$ la \emph{différence symétrique} de $A$ et $B$. Montrer que pour toute partie $A$ de $E$, $A^2 = A$ et que le triplet $(\mathcal{P}(E),\Delta,\cap)$ forme un anneau commutatif.
\finenonce{007300}
\finexercice


\exercice{7301, mourougane, 2021/08/10}
\enonce{007301}{}
Soit $G$ un groupe tel que $\forall x\in G, x\star x=e_G$.

Montrer que tout élément de $G$ est son propre inverse.
Calculer pour tout $(a,b)\in G^2$, le produit $a\star b\star a^{-1}\star b^{-1}$ et montrer que $G$ est abélien.
\finenonce{007301}
\finexercice


\exercice{7343, mourougane, 2021/08/10}
\enonce{007343}{}
\begin{enumerate}
Voici la table d'un groupe $G$.
$$\begin{array}{|c||c|c|c|c|c|c|c|c|c|c|c|c|}\hline
 * &	 a &	 b 	& c &	 d 	& e 	& f &	 g &	 h &	 i& 	 j& 	 k & 	 l\\ \hline\hline
 a &	 d &	 j &	 g &	 &	 h &	 i & 	 &	 e & 	 f& 	 b & 	 &	 k \\ \hline
 b &	 &	 g &	 j & 	 &	 l &	 e &	 d & 	 i & 	 k& 	 a &	 h &	 \\ \hline
 c &	 b &	 a &	 d &	 	& &	 k &	 j &	 l & 	e &	 g & 	 f &	 h \\ \hline
 d &	  &	 &	 c & 	 d &	 &	 &	 &	 h & 	i 	&  &  &	 l \\ \hline
 e &	 &	 f & 	 &	 e & 	 d & 	 b & 	 l &	 j & 	c &	 h &	 a &	 g \\ \hline
 f &	 i &	 l &	 h &	 	& 	& d &	 e &	 c & 	 a & 	 k & 	 j &	 b \\ \hline
 g &	 j &	 d &	 &	 g & 	 f &	 l &	 	& k &	 h& 	 c & 	 i & 	 e \\ \hline
 h &	 l &	 i &	 &	 h &	 a & 	 j &	 k &	 b & 	& 	 e &	 d &	 c \\ \hline
 i &	 f &	 &	 e &	 i &	 c 	& a &	 h &	 	& d & 	 l &	 b &	 j \\ \hline
 j & 	 &	 c& 	 b &	 j & 	 k &	 &	 a &	 f & 	 l &	 d &	 e &	 i \\ \hline
 k &	 e &	 h &	 l &	 k &	 j &	 c &	 i 	& d &	 &	 &	 g & 	 a \\ \hline
 l & 	 &	 e & 	 k & 	 l &	 b &	 g &	 &	 a & 	 j& 	 i& 	 c &	 d \\ \hline\end{array}
$$
 \item La compléter. 
 \item Le groupe est-il commutatif ?
 \item Déterminer l'ordre de $b$.
\end{enumerate}
\finenonce{007343}
\finexercice


\exercice{7344, mourougane, 2021/08/10}
\enonce{007344}{}
Considérons le groupe
 $\displaystyle\left(\mathbb{F}_{53} \right)^\times$.
Quels sont les ordres possibles d'un élément de ce groupe ?
Combien d'éléments de chaque ordre ce groupe possède -t-il ?
Déterminer en le justifiant un générateur de ce groupe.
\finenonce{007344}
\finexercice


\exercice{7354, mourougane, 2021/08/10}
\enonce{007354}{}
Montrer que tout groupe cyclique d'ordre $n$ est isomorphe à $\Z/n\Z$.
\finenonce{007354}
\finexercice


\exercice{7355, mourougane, 2021/08/10}
\enonce{007355}{}
 Montrer que les deux permutations 
$\tau=(1,2)$ et $c=(1,2,3)$ engendrent le groupe symétrique $\mathcal{S}_3$. On pourra construire un arbre dont $Id$ est la racine et tel qu'en chaque noeud partent deux branches, l'une marquée $\tau$, l'autre $c$.
\finenonce{007355}
\finexercice


\exercice{7356, mourougane, 2021/08/10}
\enonce{007356}{}
\begin{enumerate}
 \item Donner la liste des permutations du groupe symétrique $\mathcal{S}_4$.
 \item Les deux permutations $\tau=(1,3)$ et $c=(1,2,3,4)$ engendrent-elles le groupe symétrique $\mathcal{S}_3$ ?
\end{enumerate}
\finenonce{007356}
\finexercice


\exercice{7357, mourougane, 2021/08/10}
\enonce{007357}{}
\begin{enumerate}
 \item Construire à la règle et au compas un hexagone régulier $E$.
 \item Une translation de vecteur non nul peut-elle conserver globalement l'hexagone $E$ ?
 \item Quels sont les centres et les angles des rotations du plan euclidien orienté qui conservent globalement l'hexagone $E$ ?
 \item Quelles sont les symétries orthogonales du plan euclidien orienté qui conservent globalement l'hexagone $E$ ?
 \item Quelles sont les symétries orthogonales glissées du plan euclidien orienté qui conservent globalement l'hexagone $E$ ?
 \item Soit $r$ une rotation d'angle minimal non nul qui conserve globalement l'hexagone $E$ et $s$ une symétrie orthogonale qui conserve globalement $E$. Quel est l'ordre de $r$ ? et celui de $s$ ? 
 \item \'Ecrire $r$ comme composée de $s$ suivie d'une autre symétrie orthogonale.
Déterminer la nature et les éléments caractéristiques (centre, axe, angle...) de l'isométrie $rs$.
 \item Donner la liste des éléments du groupe $D_{12}$ des isométries du plan euclidien orienté qui conservent globalement l'hexagone $E$. On n'utilisera dans les notations seulement une rotation $r$ et une symétrie $s$.
 \item \`A l'aide de l'écriture de $r$ comme composée de deux symétries, calculer $sr$.
 \item En déduire les produits $rsrs$, $rsr^2s$, $r^3sr^4s$.
\end{enumerate}
\finenonce{007357}
\finexercice


\exercice{7358, mourougane, 2021/08/10}
\enonce{007358}{}
\begin{enumerate}
 \item Montrer que le groupe $GL_2(\mathbb{F}_2)$ des matrices inversibles de taille $2\times 2$ à coefficients dans le corps $\mathbb{F}_2=\Z/2\Z$ est un groupe de cardinal $(2^2-1)(2^2-2)=6$.
 \item On a vu en cours que tout groupe d'ordre $6$ est isomorphe soit au groupe cyclique $\Z/6\Z$, soit au groupe symétrique $\mathcal{S}_3$ et que ces deux groupes ne sont pas isomorphes.
Auquel de ces deux groupes, le groupe $GL_2(\mathbb{F}_2)$ est-il isomorphe ?
\end{enumerate}
\finenonce{007358}
\finexercice


\exercice{7359, mourougane, 2021/08/10}
\enonce{007359}{}
\begin{enumerate}
\item Enoncer le théorème de Lagrange.
\item Donner l'exemple de deux groupes finis de même ordre non isomorphes.
Justifier le fait que les deux groupes choisis ne sont pas isomorphes.
\item Donner la définition d'un sous-groupe \emph{distingué}.
\item L'anneau $(\Z/24\Z,+,\times)$ est-il intègre ? Justifier.
\item Effectuer la division euclidienne de $X^3+2X^2-5X+8$ par $X^2-1$ dans $\Z/5\Z[X]$.
\item Effectuer la division euclidienne de $X^3+2X^2-5X+8$ par $2X^2-1$ dans $\Z/5\Z[X]$.
\end{enumerate}
\finenonce{007359}
\finexercice


\exercice{7360, mourougane, 2021/08/10}
\enonce{007360}{}
\begin{enumerate}
 \item Calculer les produits dans $\mathcal{S}_7$, $(1,2)(1,3)$ et $(1,2)(2,3)(1,2)$.
 \item L'ensemble des transpositions de $\mathcal{S}_7$ est-il un sous-groupe de $\mathcal{S}_7$ ?
\end{enumerate}
\finenonce{007360}
\finexercice


\exercice{7361, mourougane, 2021/08/10}
\enonce{007361}{}
\begin{enumerate}
 \item Montrer que $1$ est dans le sous-groupe de $(\Z,+)$ engendré par $\{3,8\}$.
 \item Quel est le sous-groupe engendré par $\{3,8\}$ ?
 \item Quel est le sous-groupe engendré par $\{3,8,15\}$ ?
\end{enumerate}
\finenonce{007361}
\finexercice


\exercice{7362, mourougane, 2021/08/10}
\enonce{007362}{}
\begin{enumerate}
\item Le sous-ensemble $28\Z+18\Z$ de $\Z$ est-il un sous-groupe de $\Z$ ? Si oui, le déterminer.
\item Le sous-ensemble $28\Z\cup 18\Z$ de $\Z$ est-il un sous-groupe de $\Z$ ? Si oui, le déterminer.
\item Le sous-ensemble $28\Z\cap 18\Z$ de $\Z$ est-il un sous-groupe de $\Z$ ? Si oui, le déterminer.
\item Donner l'exemple d'un groupe commutatif non cyclique.
\item Donner l'exemple d'un groupe non commutatif infini.
\item Donner la définition d'un anneau.
\item Démontrer que tout groupe fini d'ordre premier est cyclique.
\end{enumerate}
\finenonce{007362}
\finexercice


\exercice{7363, mourougane, 2021/08/10}
\enonce{007363}{}
\begin{enumerate}
\item Quels sont les ordres possibles des éléments d'un groupe d'ordre $6$ ?
\item Quels sont les éléments inversibles de $(\Z/7\Z,\times)$ ?
\item Ecrire la table de multiplication de $((\Z/7\Z)^\star,\times)$.
\item Déterminer s'il en existe un générateur de $((\Z/7\Z)^\star,\times)$.
\item Combien $((\Z/7\Z)^\star,\times)$ a-t-il de générateurs ?
\item Le groupe $\Z/2\Z\times \Z/4\Z$ est-il cyclique ?
\item Le groupe des inversibles de $\Z/20\Z$ est-il cyclique ?
\end{enumerate}
\finenonce{007363}

\finexercice


\exercice{7364, mourougane, 2021/08/10}
\enonce{007364}{}
Dans le groupe $\mathcal{S}_7$ des permutations de l'ensemble fini $\{1,2,\cdots,7\}$ on considère les deux permutations 
$$\alpha=(157)(43) \quad \text{ et } \quad \beta=(26).$$
\begin{enumerate}
 \item Montrer que $\alpha\beta=\beta\alpha$.
 \item Déterminer l'ordre de $\alpha$ et de $\beta$.
 \item Calculer $\alpha^{-1}$ et $\beta^{-1}$.
  \item Montrer que $S=\{\alpha^i\beta^j, 0\leq i \leq 5, 0\leq j\leq 1\}$ est un sous-groupe de $\mathcal{S}_7$.
 \item Calculer l'ordre de $S$.
 \item Montrer que tout sous-groupe de $\mathcal{S}_7$ qui contient $\alpha$ et $\beta$ contient $S$.
 \item Que peut-on déduire des questions 4. et 6. précédentes ?
 \item Déterminer l'ordre de $\alpha\beta$.
 \item Le sous-groupe $S$ est-il cyclique ?
\end{enumerate} 
\finenonce{007364}

\finexercice


\exercice{7365, mourougane, 2021/08/10}
\enonce{007365}{}
\begin{enumerate}
 \item Calculer les produits de matrices $ADA^{-1}$ et $BDB^{-1}$ dans $GL(2,\R)$ où
$$A=\begin{pmatrix}0&1\\2&0\end{pmatrix} 
\quad \text{ et } \quad B=\begin{pmatrix}1&4\\0&1\end{pmatrix}\quad \text{ et } \quad D=\begin{pmatrix}2&0\\0&3 \end{pmatrix}$$
 \item Le sous-groupe de $GL(2,\R)$ des matrices diagonales inversibles est-il distingué dans $GL(2,\R)$ ?
 \item Calculer le produit de matrices $AB$ dans $GL(2,\R)$.
 \item L'application trace de $GL(2,\R)$ dans $(\R,+)$ est-elle un morphisme de groupes ?
\end{enumerate} 
\finenonce{007365}

\finexercice


\exercice{7366, mourougane, 2021/08/10}
\enonce{007366}{}
\begin{enumerate}
 \item Décomposer en produit de cycles à supports disjoints la permutation
$$\sigma=\left(\begin{array}{ccccccc}
1&2&3&4&5&6&7 \\
7&6&3&1&4&2&5
    \end{array}
 \right).$$
 \item Déterminer son ordre.
 \item Ecrire sa puissance $\sigma ^6$.
 \item Ecrire son inverse en produit de cycles à supports disjoints.
 \item La décomposer en produit de moins de $7$ tranpositions.
\end{enumerate}
\finenonce{007366}
\finexercice


\exercice{7367, mourougane, 2021/08/10}
\enonce{007367}{}
 \begin{enumerate}
\item Quels sont les éléments inversibles de l'anneau $(\Z/11\Z)$ ?
 \item Déterminer un générateur $g$ du groupe $(\Z/11\Z)^\star$ des inversibles de l'anneau $(\Z/11\Z)$.
\item Ecrire chaque élément de $(\Z/11\Z)^\star$ comme puissance de $g$.
\end{enumerate}
\finenonce{007367}
\finexercice

\section{203.02 Ordre d'un élément}


\exercice{1335, cousquer, 2003/10/01}

\enonce{001335}{}
 On appelle \emph{ordre d'un élément} d'un groupe fini $(G,\ast)$ l'ordre du
sous-groupe engendré dans $G$ par cet élément.
\begin{enumerate}
    \item Montrer que si $x$ est d'ordre $p$, $p$ est le plus petit entier tel que
$x^p=e$.
    \item Déterminer les ordres des éléments des groupes rencontrés au 
\textbf{I}.
    \item Soit $(G,\ast)$ un groupe fini, $a$ un élément de $G$, $H$ un sous-groupe
d'ordre $p$ de $G$~; on note $aH$ l'ensemble $\{a\ast y\mid y\in 
H\}$.\newline
a) Montrer que pour tout $a\in G$, $aH$ a $p$ éléments.\newline
b) Montrer que si $a\in G$ et $b \in G$, $(aH=bH)$ ou $(aH\cap
bH=\emptyset)$.\newline 
c) En déduire que l'ordre de $H$ divise l'ordre de~$G$.
    \item Montrer que si $G$ est un groupe fini d'ordre $n$, les ordres de tous ses
éléments divisent~$n$.
    \item Trouver des sous-groupes de $\mathbb{Z}_2$, $\mathbb{Z}_3$, 
$\mathbb{Z}_4$, $\mathbb{Z}_5$, $\mathbb{Z}_6$, 
$\frak{S}_2$, $\frak{S}_3$.
    \item Si $G$ est un groupe d'ordre~$5$, que peut-on dire de l'ordre de ses
éléments~? En déduire les tables de composition possibles pour un groupe
d'ordre~$5$. Que peut-on dire de deux groupes quelconques d'ordre~$5$~? Mêmes
questions pour les groupes d'ordre~$23$. Généraliser.
\end{enumerate}
\finenonce{001335}



\finexercice



\exercice{1336, legall, 1998/09/01}

\enonce{001336}{}
 Soit  $H$  un groupe ab\' elien. Un \' el\' ement
$x\in H$ est dit d'ordre fini lorsque il existe  $n\in {\Nn}$  tel
que la somme  $x+...+x$  ($n$-fois) soit \' egale \`a  $0$.
Montrer que l'ensemble des \' el\' ements d'ordre fini de  $H$ est
un sous-groupe ab\' elien de  $H$.

\finenonce{001336}



\finexercice



\exercice{1337, legall, 1998/09/01}

\enonce{001337}{}

Soit  $G$  un groupe, $e$  son \' el\' ement neutre. Un
\' el\' ement  $g$  de  $G$  est
dit {\em d'ordre  $n\in {\Nn}$}  si  $g^n=e$  et  $g^k\not = e$  pour tout
entier  $k<n$. $g$  est dit {\em d'ordre fini} si il est d'ordre  $n$  pour
un  $n$  quelconque.
\begin{enumerate}
    \item Montrer que  $Gl_2({\Rr})$  contient des \' el\' ements d'ordre $2$  et
des \' el\' ements qui ne sont pas d'ordre fini.
    \item Soit  $\varphi $  un homomorphisme de  $G$  \`a valeurs
dans  $H$  et  $g$  un \' el\' ement de  $ G$  d'ordre  $n$. Montrer que :

 - $\varphi (g)$  est d'ordre fini
inf\' erieur ou \' egal \`a  $n$.

 - Si  $\varphi $  est injectif, l'ordre de  $\varphi (g)$  est \' egal \`a  $n$.
    \item Montrer que si  $G$  n'a qu'un nombre fini d'\' el\' ements, tous ses \' el\' ements
 ont un ordre fini.
\end{enumerate}

\finenonce{001337}



\finexercice



\exercice{1338, ortiz, 1999/04/01}

\enonce{001338}{}
Soit le groupe $G=\Zz/12\Zz$.
\begin{enumerate}
    \item  D\'eterminer le sous-groupe $H$ de $G$ engendr\'e par
$\overline{6}$ et $\overline{8}$ et d\'eterminer
son ordre.
    \item  Caract\'eriser les g\'en\'erateurs de $G.$
    \item  Quel est l'ordre de l'\'el\'ement $\overline{9}$ ?
\end{enumerate}
\finenonce{001338}



\finexercice



\exercice{1339, ortiz, 1999/04/01}

\enonce{001339}{}
Soient $E$ un espace vectoriel r\'eel de dimension
2 et $(e_1,e_2)$ une base de $E.$ On consid\`ere
les endomorphismes de $E$ d\'efinis par
$${s(e_1)=e_1,\quad s(e_2)=-e_2,}$$
$${r(e_1)=e_2,\quad r(e_2)=-e_1.}$$
\begin{enumerate}
    \item  Montrer que $r$ et $s$ sont des automorphismes du $\Rr$-espace vectoriel $E.$
    \item  D\'eterminer l'ordre de $s$ et l'ordre de $r.$
    \item
    \begin{enumerate}
         \item Montrer que $sr=-rs.$
         \item En d\'eduire que $G:=\left\{\text{Id}_E,s,r,sr,-\text{Id}_E,-s,-r,-s\right\}$
est un sous-groupe du groupe lin\'eaire de $E.$
         \item Montrer que $G$ est le sous-groupe du groupe lin\'eaire
$\text{GL}(E)$ engendr\'e par $s$ et $t.$
     \end{enumerate}
\end{enumerate}
\finenonce{001339}



\finexercice



\exercice{1340, ortiz, 1999/04/01}

\enonce{001340}{}
 Soient  $G$ un groupe et $x\in G$ un \'el\'ement
d'ordre $n.$ Quel est l'ordre de $x^2$ ?

\finenonce{001340}



\finexercice



\exercice{1341, ortiz, 1999/04/01}

\enonce{001341}{}
\begin{enumerate}
    \item  Soient  $G$ un groupe et $x, y\in G$ des \'el\'ements
qui commutent et d'ordres respectifs $m$ et $n$
premiers entre eux. Montrer que $xy$ est d'ordre
$mn.$ Montrer que l'hypoth\`ese $m$ et $n$
\textit{premiers entre eux} est indispensable.
    \item  Montrer que $A:=\left(\begin{smallmatrix} 0&-1\\1&0
\end{smallmatrix}\right)$ et $B:=\left(\begin{smallmatrix}
0&1\\-1&-1\end{smallmatrix}\right)$ sont des
\'el\'ements de $\text{GL}(2,\Rr)$ d'ordres finis
et que $AB$ n'est pas d'ordre fini.
\end{enumerate}

\finenonce{001341}




\finexercice



\exercice{1342, legall, 1998/09/01}

\enonce{001342}{}

 Le groupe $  (\Q ,+)  $ est-il monog\`ene ?

\finenonce{001342}



\finexercice



\exercice{2982, quercia, 2010/03/08}

\enonce{002982}{Sous groupes finis de $\C^*$}
D{\'e}terminer tous les sous-groupes finis de $(\C^*,\times)$.
\finenonce{002982}



\finexercice 


\exercice{2983, quercia, 2010/03/08}

\enonce{002983}{Ordre d'un {\'e}l{\'e}ment}
\begin{enumerate}
  \item Soient $G$ et $G'$ deux groupes et $f$ un morphisme de $G$ dans $G'$.
    Pour $a \in G$, comparer l'ordre de $a$ et celui de $f(a)$.
  \item Soient $a,b \in G$. Comparer les ordres de $a$ et de $bab^{-1}$.
  \item Soient $a,b \in G$. Comparer les ordres de $ab$ et de $ba$.
\end{enumerate}
\finenonce{002983}



\finexercice 


\exercice{2984, quercia, 2010/03/08}

\enonce{002984}{Ordre de $ab$}
Soient $a,b$ deux {\'e}l{\'e}ments d'un groupe multiplicatif $G$ tels que :
$\begin{cases}a \text{ est d'ordre } \alpha\cr
        b \text{ est d'ordre } \beta\cr
        \alpha \wedge \beta = 1 \cr
        ab = ba.\end{cases}$

D{\'e}terminer l'ordre de $ab$.
\finenonce{002984}



\finexercice 


\exercice{2985, quercia, 2010/03/08}

\enonce{002985}{D{\'e}composition d'un {\'e}l{\'e}ment d'ordre fini}
Soit $G$ un groupe multiplicatif et $a \in G$ d'ordre $np$ avec $n \wedge p = 1$.
\par
Montrer qu'il existe $b,c \in G$ uniques tels que
    $b$ est d'ordre $n$,
    $c$ est d'ordre $p$,
    $a = bc = cb$.
\par
\finenonce{002985}



\finexercice 

\section{203.03 Morphisme, isomorphisme}


\exercice{2977, quercia, 2010/03/08}

\enonce{002977}{Groupe des automorphismes}
Soit $G$ un groupe multiplicatif. On note Aut$(G)$ l'ensemble des
isomorphismes $\phi : G \to  G$.
\begin{enumerate}
  \item Montrer que Aut$(G)$ est un groupe pour la loi $\circ$.
  \item D{\'e}terminer Aut$(\Z)$.
  \item Pour $a \in G$ on note ${\phi_a} : G \to G, x \mapsto {axa^{-1}}$
    \par
    Montrer que $\phi_a \in \text{Aut}(G)$, et que l'application
    $a \mapsto \phi_a$ est un morphisme de groupes.
\end{enumerate}
\finenonce{002977}



\finexercice 

\exercice{2979, quercia, 2010/03/08}

\enonce{002979}{Images directes et r{\'e}ciproques}
Soit $G$ un groupe additif et $f : G \to {G'}$ un morphisme de groupes.
\begin{enumerate}
  \item Montrer que pour tout sous-groupe $H$ de $G$ on a : $f^{-1}(f(H)) = H + \mathrm{Ker}\,f$.
  \item Montrer que pour tout sous-groupe $H'$ de $G'$ on a : $f(f^{-1}(H')) = H' \cap \Im\,f$.
\end{enumerate}
\finenonce{002979}



\finexercice 


\exercice{2980, quercia, 2010/03/08}

\enonce{002980}{Morphismes entre deux groupes cycliques}
Soit $G$ un groupe cyclique engendr{\'e} par $a$ d'ordre $n$, $G'$ un deuxi{\`e}me
groupe, et $a' \in G'$.

Montrer qu'il existe un morphisme $\phi : G \to {G'}$ tel que $\phi(a) = a'$
si et seulement si $a'$ est d'ordre fini divisant~$n$.

Application : d{\'e}terminer tous les morphismes :
$\Z/n\Z \to \Z$,\quad $\Z/n\Z \to \C^*$,\quad $\Z/n\Z \to \Z/p\Z$.
\finenonce{002980}



\finexercice 


\exercice{2981, quercia, 2010/03/08}

\enonce{002981}{Morphismes de $\Q$ additif}
D{\'e}terminer tous les morphismes de
\begin{enumerate}
  \item $(\Q,+)$ dans $(\Q,+)$. 
  \item $(\Q,+)$ dans $(\Z,+)$. 
  \item $(\Q,+)$ dans $(\Q^*,\times)$. 
\end{enumerate}
\finenonce{002981}



\finexercice 

\section{203.04 Anneau}


\exercice{1355, ortiz, 1999/04/01}

\enonce{001355}{}
Soient $a,b\in\Cc$. L'application
$f:\Cc\to\Cc,z\mapsto iz-\overline{z}$ est-elle un
(endo)morphisme...
\begin{enumerate}
    \item ...du groupe $\Cc$ ?
    \item ...de l'anneau $\Cc$ ?
    \item ...du $\Rr$-espace vectoriel $\Cc$ ?
\end{enumerate}
\finenonce{001355}



\finexercice



\exercice{1356, ortiz, 1999/04/01}

\enonce{001356}{}
Soient les ensembles
$${L=\left\{
\left (
\begin{array}{ll}
x&0\\
0&0
\end{array}
\right )\in \mathcal{M}_2(\Rr): x\in \Rr
\right\}\text{ et }
M=\left\{
\left (
\begin{array}{ll}
x&x\\
-x&-x
\end{array}
\right )\in \mathcal{M}_2(\Rr): x\in \Rr
\right\}}$$

\'Etudier si, munis des lois usuelles, $L$ et  $M$ sont des anneaux, des corps.
\finenonce{001356}



\finexercice



\exercice{1357, ortiz, 1999/04/01}

\enonce{001357}{}
\begin{enumerate}
    \item  Soit $D=\left\{f\in \Rr[X]: f'(0)=0\right\}.$ Montrer que $D$
n'est pas un id\'eal de l'anneau $\Rr [X]$ et que
c'est un sous-anneau de l'anneau $\Rr [X]$.
    \item  Soit $E=\left\{f\in \Rr[X]: f(0)=f'(0)=0\right\}$. Montrer que
$D$ n'est pas un sous-anneau de l'anneau $\Rr [X]$
et que c'est un id\'eal de l'anneau $\Rr [X]$ dont
on donnera un g\'en\'erateur.
\end{enumerate}
\finenonce{001357}



\finexercice



\exercice{1358, ridde, 1999/11/01}

\enonce{001358}{}
On d\'efinit sur $\Rr$ les deux lois $\oplus$ et $\otimes$ par
$x \oplus y = x + y-1$ et $x \otimes y = x + y-xy$. Montrer que
$ (\Rr, \oplus, \otimes)$ est un corps.
\finenonce{001358}



\finexercice



\exercice{1359, ridde, 1999/11/01}

\enonce{001359}{}
Soit $ (G,  + )$ un groupe commutatif. On note $\text{End} (G)$ l'ensemble des
endomorphismes de $G$ sur lequel on d\'efinit la loi $ + $ par
${f + g} : {G} \to {G}$, ${x}\mapsto {f (x) + g (x)}$.\\
 Montrer que $ (\text{End} (G),  +, \circ)$
est un anneau.
\finenonce{001359}



\finexercice



\exercice{1360, ridde, 1999/11/01}

\enonce{001360}{}
Soit $ (A,  +, \times)$ un anneau. On dit que $x \in A$ est nilpotent ssi il
existe $n \in \Nn$ tel que $x^n = 0$.
\begin{enumerate}
\item Montrer que si $x$ est nilpotent alors $1-x$ est inversible.
\item Montrer que si $x$ et $y$ sont nilpotents et commutent, alors $xy$ et
$x + y$ sont nilpotents.
\item Un corps admet-il des \'el\'ements nilpotents ?
\end{enumerate}
\finenonce{001360}



\finexercice



\exercice{1361, ridde, 1999/11/01}

\enonce{001361}{}
Soit $ (A,  +, \times)$ un anneau.\\ On appelle centre de $A$ l'ensemble
$C = \left\{ x \in A / \forall y \in A, xy = yx\right\}$. \\
Montrer que $C$ est un sous-anneau de $A$.
\finenonce{001361}



\finexercice



\exercice{1362, ridde, 1999/11/01}

\enonce{001362}{}
Soient $A$ et $B$ deux anneaux. On d\'efinit sur $A \times B$ les lois
$$ (x, y) + (x', y') = (x + x', y + y')$$
$$ (x, y) (x', y') = (xx', yy')$$
\begin{enumerate}
\item Montrer que $A \times B$ est alors un anneau.
\item Si $A$ et $B$ sont des corps, en est-il de même pour $A \times B$ ?
\end{enumerate}
\finenonce{001362}



\finexercice



\exercice{1363, ridde, 1999/11/01}

\enonce{001363}{}
Montrer que si $A_1, \ldots, A_n$ sont des sous-anneaux de $A$ alors
$A_1 \cap \ldots \cap A_n$ est un sous-anneau de $A$.
\finenonce{001363}



\finexercice



\exercice{1364, ridde, 1999/11/01}

\enonce{001364}{}
Soit $\Zz[i] = \left\{ a + ib, (a, b) \in \Zz^2\right\}$.
\begin{enumerate}
\item Montrer que $\Zz[i]$ est un anneau commutatif pour les lois usuelles de $\Cc$.
\item D\'eterminer les inversibles de $\Zz[i]$.
\end{enumerate}
\finenonce{001364}



\finexercice



\exercice{1365, ortiz, 1999/04/01}

\enonce{001365}{}
Soit $A$ un anneau commutatif. On dit que $a\in A$
est nilpotent s'il existe $n\in \Nn^{*}$ tel que
$a^n=0$. On pose $\mathcal{N}\left( A\right)
=\left\{ a\in A:a\text{ est nilpotent}\right\} .$
\begin{enumerate}
    \item  Dans cette question, $A=\Zz{/}72\Zz$. Montrer
que $\overline{6}\in \mathcal{N}\left( A\right) $
puis que $\mathcal{N}
\left( A\right) =\left\{ \lambda \overline{6}:\lambda \in \Zz\right\} .$
    \item   Que peut-on dire de $\mathcal{N}\left( A\right) $ si $A$ est
int\`egre?
    \item  Montrer que $\mathcal{N}\left( A\right) $ est un id\'eal de $A$ %(on
%pourra utiliser la formule du bin\^ome de Newton).
\end{enumerate}
\finenonce{001365}



\finexercice



\exercice{1366, ortiz, 1999/04/01}

\enonce{001366}{Extrait de l'examen de juin 1994}
 Sur l'ensemble $\Rr^2,$ on d\'efinit la loi $\star$ par
$${(x_1,x_2)\star
(y_1,y_2)=(x_1y_1,x_1y_2+x_2y_1).}$$
\begin{enumerate}
    \item
    \begin{enumerate}
         \item Montrer que $(\Rr^2,+,\star)$ est
un anneau commutatif not\'e $A.$
         \item Chercher les diviseurs de $0$ de l'anneau $A.$
    \end{enumerate}
    \item  On consid\`ere l'application
$${f:\Rr[X]\to A, P\mapsto (P(0),P'(0)).}$$
    \begin{enumerate}
         \item Montrer que $f$ est un homomorphisme d'anneaux.
         \item $f$ est-il surjectif ?
         \item D\'eterminer le noyau de $f.$
    \end{enumerate}
\end{enumerate}
\finenonce{001366}



\finexercice



\exercice{1367, ortiz, 1999/04/01}

\enonce{001367}{Extrait de l'examen de janvier 1994}
 On d\'efinit
$A=\left\{a+jb:a,b\in \Zz\right\}$  o\`u $j=\exp(\frac
{2i\pi}3)$.
\begin{enumerate}
    \item  Montrer que $A$ est un sous-anneau de $\Cc.$
On d\'esigne par $\mathcal{U}(A)$ le groupe des
\'el\'ements inversibles de $A$ et enfin, on pose,
pour tout $z\in\Cc,$ $N(z)={|z|}^2.$
    \item
    \begin{enumerate}
         \item Montrer que si $z\in A$ alors $N(z)\in\Zz.$
         \item Soit $z\in A.$ Montrer que  $z\in \mathcal{U}(A)$ si et seulement si $N(z)=1.$
         \item Soient $a$ et $b$ des entiers. Montrer que si $N(a+jb)=1$ alors  $a,b\in\left\{-1,0,1\right\}.$
    \end{enumerate}
    \item  D\'ecrire le groupe $\mathcal{U}(A)$ et en d\'eterminer les \'el\'ements d'ordre 3.
    \item  Soit $\Phi: \Qq[X]\to\Cc, P\mapsto P(j).$
    \begin{enumerate}
        \item Montrer que $\Phi$ est un homomorphisme d'anneaux.
        \item D\'eterminer le noyau de $\Phi$ (on pourra
remarquer que $j^2+j+1=0$).
        \item Montrer que $\text{Im\,}\Phi=\left\{a+jb:a,b\in \Qq\right\}$ et que c'est un sous-corps de $\Cc.$
    \end{enumerate}
\end{enumerate}
\finenonce{001367}



\finexercice



\exercice{1369, ortiz, 1999/04/01}

\enonce{001369}{D'apr\`es examen juin 94}
\begin{enumerate}
    \item   Montrer que $\overline{k}$ est inversible dans l'anneau $\Zz
{/}n\Zz$ si et seulement si les entiers $k$ et $n$
sont premiers entre eux.
    \item  On pose $n=10$ et soit $G$ le groupe des \'el\'ements inversibles de $%
\Zz{/}n\Zz$.
    \begin{enumerate}
         \item Donner la liste des \'el\'ements de $G.$
         \item Quel est l'ordre de $\overline{3}$ ? $G$ est-il cyclique ?
    \end{enumerate}
\end{enumerate}
\finenonce{001369}



\finexercice



\exercice{1370, ortiz, 1999/04/01}

\enonce{001370}{Bac 1978}
 Soit l'anneau $A=\Zz{/}91\Zz$.
\begin{enumerate}
    \item  D\'eterminer les diviseurs de z\'ero de l'anneau $A$.
    \item   R\'esoudre dans $A$ l'\'equation $x^2+\overline{2}x-\overline{3}=%
\overline{0}$.
\end{enumerate}
\finenonce{001370}



\finexercice



\exercice{1372, ortiz, 1999/04/01}

\enonce{001372}{}
Montrer que $\Zz{/}n\Zz$ est un anneau principal.
\finenonce{001372}



\finexercice



\exercice{1373, gourio, 2001/09/01}

\enonce{001373}{}
Soit $A$ un anneau fini commutatif int\`{e}gre (i.e.\ $xy=0\Rightarrow x=0 $ ou $y=0$).
Montrer que c'est un corps, i.e. que tout \'{e}l\'{e}ment non nul est
inversible.
\finenonce{001373}



\finexercice



\exercice{1374, gourio, 2001/09/01}

\enonce{001374}{}
Soit $A$ un anneau, on dit que $x\in A$ est nilpotent si $\exists n\in \Nn$
tel que $x^{n}=0.$
\begin{enumerate}
\item Montrer que si $x$ est nilpotent alors $(1-x)$ est inversible.
\item Montrer que si $x$ et $y$ sont nilpotents et commutent alors $xy$ et $x+y$
sont nilpotents.
\item Un corps admet-il des \'{e}l\'{e}m{e}nts nilpotents ?
\end{enumerate}
\finenonce{001374}



\finexercice



\exercice{3005, quercia, 2010/03/08}

\enonce{003005}{Anneau de Boole}
Soit $E$ un ensemble fini et $A = {\cal P}(E)$.
\begin{enumerate}
  \item Montrer que $(A,\Delta,\cap)$ est un anneau commutatif. Est-il int{\`e}gre ?
  \item Soit $I$ un id{\'e}al de $A$. Montrer que :
    $\begin{cases}\forall\ X \in I,\ \forall\ Y \subset X,\text{ on a } Y \in I\cr
            \forall\ X,Y \in I,\text{ on a } X \cup Y \in I.\end{cases}$
  \item En d{\'e}duire que $I = {\cal P}(E')$ avec $E' \subset E$.
  \item {\'E}tudier la r{\'e}ciproque.
  \item Si $E$ est infini, montrer que $I = \{$parties finies de $E \}$ est un
    id{\'e}al qui n'est pas de la forme ${\cal P}(E')$.
\end{enumerate}
\finenonce{003005}



\finexercice  


\exercice{3006, quercia, 2010/03/08}

\enonce{003006}{Id{\'e}aux triviaux}
Soit $A$ un anneau commutatif non nul dont les seuls id{\'e}aux sont $\{0\}$ et $A$.
Montrer que $A$ est un corps.
\finenonce{003006}



\finexercice 


\exercice{3007, quercia, 2010/03/08}

\enonce{003007}{Id{\'e}aux premiers}
Un id{\'e}al $I$ d'un anneau $A$ est dit premier si~: $\forall\ x,y\in A,\ xy\in I \Rightarrow  x\in I$ ou $y\in I$.
\begin{enumerate}
  \item Quels sont les id{\'e}aux premiers de~$\Z$~?
  \item Montrer que si $A$ est commutatif non nul et si tous les id{\'e}aux de~$A$ sont premiers alors $A$ est un corps.
\end{enumerate}
\finenonce{003007}


\finexercice 


\exercice{3008, quercia, 2010/03/08}

\enonce{003008}{Th{\'e}or{\`e}me de Gauss}
Soit $A$ un anneau commutatif et $a,b \in A$. On dit que :
$\begin{cases} a \text{ divise } b \text{ si } b \in aA\cr
         a \text{ est premier {\`a} } b \text{ si } aA + bA =
         A.\end{cases}$

Montrer que si $a$ est premier {\`a} $b$ et $a$ divise $bc$, alors $a$ divise $c$.
\finenonce{003008}



\finexercice   


\exercice{3009, quercia, 2010/03/08}

\enonce{003009}{Caract{\'e}ristique}
Soit $A$ un anneau. On appelle {\it caract{\'e}ristique de\/} $A$ l'ordre de
$1$ dans le groupe additif $(A,+)$.
On suppose $A$ de caract{\'e}ristique finie, $n$.

\begin{enumerate}
  \item Montrer que : $\forall\ x \in A,\ nx = 0$.
  \item Si $A$ est int{\`e}gre, montrer que $n$ est un nombre premier.
  \item Si $A$ est int{\`e}gre et commutatif, montrer que $x  \mapsto x^n$ est un morphisme
    d'anneau.
\end{enumerate}
\finenonce{003009}



\finexercice   


\exercice{3010, quercia, 2010/03/08}

\enonce{003010}{Anneau de caract{\'e}ristique 2}
Soit $A$ un anneau non nul tel que : $\forall\ x \in A,\ x^2 = x$.
\begin{enumerate}
  \item Exemple d'un tel anneau ?
  \item Quels sont les {\'e}l{\'e}ments inversibles de $A$ ? 
  \item Montrer que : $\forall\ x \in A,\ x+x = 0$.
    En d{\'e}duire que $A$ est commutatif.
  \item Pour $x,y \in A$ on pose :
    $x \le y \iff \exists\ a \in A \text{ tel que } x=ay$.
    Montrer que c'est une relation d'ordre.  
\end{enumerate}
\finenonce{003010}



\finexercice   


\exercice{3011, quercia, 2010/03/08}

\enonce{003011}{El{\'e}ments nilpotents}
Soit $A$ un anneau commutatif, et $a \in A$. On dit que $a$ est nilpotent s'il
existe $n \in \N$ tel que $a^n = 0$.
\begin{enumerate}
  \item Exemple : D{\'e}terminer les {\'e}l{\'e}ments nilpotents de $\Z/36\Z$.
  \item Montrer que l'ensemble des {\'e}l{\'e}ments nilpotents est un id{\'e}al de $A$.
  \item Soit $a$ nilpotent. Montrer que $1 - a$ est inversible
    (remarquer que $1 = 1^n - a^n$).
  \item Soient $a$ nilpotent et $b$ inversible. Montrer que $a + b$ est inversible.
\end{enumerate}
\finenonce{003011}



\finexercice   


\exercice{3012, quercia, 2010/03/08}

\enonce{003012}{$1-ab$ et $1-ba$}
Soit $A$ un anneau et $a,b\in A$. Montrer que $1-ab\in A^* \Leftrightarrow 1-ba \in A^*$.
\finenonce{003012}



\finexercice 


\exercice{3013, quercia, 2010/03/08}

\enonce{003013}{Radical d'un id{\'e}al}
Soit $A$ un anneau commutatif et $I$ un id{\'e}al de $A$.

On note $\sqrt{I} = \{ x \in A \text{ tq } \exists\ n \in \N \text{ tq } x^n \in I \}$
(radical de $I$).

\begin{enumerate}
  \item Montrer que $\sqrt{I}$ est un id{\'e}al de $A$.
  \item Montrer que $\sqrt{\sqrt I} = \sqrt I$.
  \item Montrer que $\sqrt{I \cap J} = \sqrt I \cap \sqrt J$ et
    $\sqrt{I + J} \supset \sqrt I + \sqrt J$.
  \item Exemple : $A = \Z$, $I = 3648\Z$. Trouver $\sqrt I$.
\end{enumerate}
\finenonce{003013}



\finexercice   


\exercice{3014, quercia, 2010/03/08}

\enonce{003014}{Produit de deux id{\'e}aux}
Soit $A$ un anneau commutatif et $I,J$ deux id{\'e}aux de $A$.

On note $IJ = \{ a_1b_1 + \dots + a_nb_n \text{ tel que } a_i \in I,\ b_i \in J \}$.

\begin{enumerate}
  \item Montrer que $IJ$ est un id{\'e}al de $A$.
  \item Montrer que $I(J+K) = IJ + IK$.
  \item On suppose $I + J = A$. Montrer que $IJ = I \cap J$.
  \item Pour $A = \Z$, $I = n\Z$, $J = p\Z$, qu'est-ce que $IJ$ ?
\end{enumerate}
\finenonce{003014}



\finexercice   


\exercice{3015, quercia, 2010/03/08}

\enonce{003015}{Relation d'{\'e}quivalence compatible avec les op{\'e}rations d'anneau}
Soit $A$ un anneau commutatif.
\begin{enumerate}
  \item Soit $\cal R$ une relation d'{\'e}quivalence compatible
    avec l'addition et la multiplication dans $A$. On note $I$ la classe de $0$.
    Montrer que $I$ est un id{\'e}al de $A$.
  \item R{\'e}ciproquement, soit $J$ un id{\'e}al de $A$.
    On pose $x \sim y \iff x-y \in J$.
    Montrer que $\sim$ est une relation d'{\'e}quivalence compatible avec $+$ et
    $\times$.
\end{enumerate}
\finenonce{003015}



\finexercice 


\exercice{3016, quercia, 2010/03/08}

\enonce{003016}{{\'E}tude de l'anneau $\Z^2$}
\begin{enumerate}
\item Soit $d \in \N$. On pose
  $$A_d = \{(x,y) \in \Z^2 \text{ tq } x\equiv y (\mathrm{mod}\, d)\}$$
  ($x=y$ pour $d=0$).
    Montrer que $A_d$ est un sous-anneau de $\Z^2$.
  \item Montrer que l'on obtient ainsi tous les sous-anneaux de $\Z^2$.

  \item Soit $I$ un id{\'e}al de $\Z^2$. On note :
    $\begin{cases} I_1 = \{ x \in \Z \text{ tq } (x,0) \in I \}\cr
             I_2 = \{ y \in \Z \text{ tq } (0,y) \in I \}.\end{cases}$

    Montrer que $I_1$ et $I_2$ sont des id{\'e}aux de $\Z$, et que
    $I = I_1 \times I_2$.
  \item En d{\'e}duire que $I$ est un id{\'e}al principal.
\end{enumerate}
\finenonce{003016}



\finexercice 


\exercice{3017, quercia, 2010/03/08}

\enonce{003017}{Id{\'e}aux de $ K^E$}
Soit $ K$ un corps, $E$ un ensemble fini, et $A = K^E$.
Pour $e \in E$, on pose :
$$I_e = \{ f \in A \text{ tq } f(e) = 0 \},  \qquad
  \chi_e\ :\ x  \mapsto \begin{cases} 1 &si x=e\cr 0 &\text{ si } x\ne e.\end{cases}$$

\begin{enumerate}
  \item Montrer que $I_e$ est un id{\'e}al principal de A.
  \item Soit $f \in A$. V{\'e}rifier que $f = \sum_{e\in E} f(e)\chi_e$.
  \item Soit $I$ un id{\'e}al quelconque de $A$, et
    $F = \{ e \in E \text{ tq } \exists\ f \in I \text{ tq } f(e) \ne 0\}$.
    \par
    Montrer que $I$ est un id{\'e}al principal engendr{\'e} par $\sum_{e\in F}\chi_e$.
\end{enumerate}
\finenonce{003017}



\finexercice 


\exercice{3018, quercia, 2010/03/08}

\enonce{003018}{Fonctions trigonom{\'e}triques} On pose\\ $A = \left\{ f : \R \to \R
\text{ de la forme } f(x) = a_0 + \sum_{k=1}^n a_k\cos(kx),\ n \in
\N,\ a_i \in \R \right\}$.
\begin{enumerate}
  \item Montrer que $A$ est un sous-anneau de $\R^\R$.
  \item Soit $f\in A$. Montrer que si $f = 0$, alors les coefficients $a_k$ sont
    tous nuls
    $\left({\text{calculer }  \int_{t=0}^{2\pi} f(t)\cos(nt)\,d t}\right)$.
  \item En d{\'e}duire que $A$ est int{\`e}gre.
\end{enumerate}
\finenonce{003018}



\finexercice   


\exercice{3019, quercia, 2010/03/08}

\enonce{003019}{Suites croissantes d'id{\'e}aux}
Soit $A$ un anneau commutatif et $(I_n)$ une suite croissante d'id{\'e}aux de $A$.
On pose $I = \bigcup_{n \in \N} I_n$.
\begin{enumerate}
  \item Montrer que $I$ est un id{\'e}al de $A$.
  \item On suppose que $A$ est principal. Montrer qu'il existe $n_0 \in \N$ tel que
    $I = I_{n_0}$.
  \item En d{\'e}duire que $\R^\R$ n'est pas principal. 
\end{enumerate}
\finenonce{003019}



\finexercice   


\exercice{3020, quercia, 2010/03/08}

\enonce{003020}{Endomorphismes d'un groupe commutatif}
Soit $G$ un groupe additif et $A = \{ \text{morphismes } f : G \to G\}$.
\begin{enumerate}
  \item Montrer que $(A,+,\circ)$ est un anneau.
  \item On prend $G = \Z/n\Z$, $n \ge 2$.
    Montrer que $A$ est l'ensemble des applications
    $ G \to G,  x \mapsto {kx}$ avec $k \in G$,
    et que $A$ est isomorphe {\`a} l'anneau $\Z/n\Z$.
\end{enumerate}
\finenonce{003020}



\finexercice   


\exercice{3021, quercia, 2010/03/08}

\enonce{003021}{Entiers 2-adiques}
Soit $A = \{ m/n \in \Q \text{ tel que } n \text{ est impair} \}$.
\begin{enumerate}
  \item Montrer que $A$ est un sous-anneau de $\Q$.
  \item Chercher les {\'e}l{\'e}ments inversibles dans $A$.
  \item Montrer que les id{\'e}aux de $A$ sont tous principaux engendr{\'e}s
    par les nom\-bres de la forme $2^k$, $k\in \N$.
\end{enumerate}
\finenonce{003021}



\finexercice   


\exercice{3022, quercia, 2010/03/08}

\enonce{003022}{Morphismes $\Z^n \to \Z$}
\label{morphismes}
Chercher les morphismes d'anneaux : $\Z^n \to \Z$.

\finenonce{003022}



\finexercice 


\exercice{3023, quercia, 2010/03/08}

\enonce{003023}{Suites stationnaires}
Soit $A = \{$suites stationnaires d'entiers re\-la\-tifs$\}$ muni des
op{\'e}rations usuelles.
\begin{enumerate}
  \item Montrer que $A$ est un anneau.
  \item Chercher les morphismes d'anneaux : $A \to \Z$.
  \item Soit $I = \{\text{suites presque nulles}\}$. Montrer que c'est un id{\'e}al non
    principal. 
\end{enumerate}
\finenonce{003023}



\finexercice   


\exercice{3024, quercia, 2010/03/08}

\enonce{003024}{Entiers de Gauss}
Soit $A = \{ a + bi \text{ tq } a,b \in \Z\}$.
\begin{enumerate}
  \item Montrer que $A$ est un sous-anneau de $\C$. Quels sont les {\'e}l{\'e}ments
    inversibles ?

  \item Soient $u,v \in A$ avec $v \ne 0$. Montrer qu'il existe $q,r \in A$
    tels que $u = qv + r$ et $|r| < |v|$.
    A-t-on unicit{\'e} ?

  \item Montrer que $A$ est principal.
\end{enumerate}
\finenonce{003024}



\finexercice   


\exercice{7302, mourougane, 2021/08/10}
\enonce{007302}{Exercice sur le cours}
\begin{enumerate}
 \item Ecrire les tables d'opération dans l'anneau $\Z/6\Z$.
  \item Soit $p$ un nombre premier, et soit $(G,\star)$ un groupe d'ordre $p$. 
Montrer que $G$ est cyclique.
 \item Démontrer que le noyau d'un morphisme d'anneaux commutatifs est un idéal.
 \item L'application $f~: \Z/10\Z\to \Z/3\Z$ construite de la manière suivante est-elle bien définie ?
Pour un élément $c$ de $\Z/10\Z$, on choisit un représentant $x$ et on pose $f(c):=[x]_3$.
\end{enumerate}
\finenonce{007302}
\finexercice


\exercice{7303, mourougane, 2021/08/10}
\enonce{007303}{}
\begin{enumerate}
 \item Montrer que $\Z \times \Z$ peut être muni d'une structure naturelle d'anneau.
 \item Montrer que $\R[X]$ est un anneau, et que pour tout $a \in \R$, l'ensemble 
$$I_a = \{P \in \R[X] ~/~ P(a) = 0\}$$ est un idéal de $\R[X]$.
\end{enumerate}
\finenonce{007303}
\finexercice


\exercice{7304, mourougane, 2021/08/10}
\enonce{007304}{}
\begin{enumerate}
 \item Déterminer l'ordre de $\overline{2}$ dans $(\Z/10\Z,+)$.
 \item Montrer que $\{\overline{0},\overline{3},\overline{6},\overline{9}\}$ est un sous-groupe de $(\Z/12\Z,+)$.
 \item Expliciter un sous-groupe d'ordre $6$ de $(\Z/12\Z,+)$.
 \item Quels sont les ordres possibles des sous-groupes de $(\Z/6\Z,+)$ ?
 \item Soient $n$ un entier naturel non nul. Le but de cette question est de déterminer tous 
les sous-groupes de $\Z/n\Z$. Soit $d$ un diviseur de $n$.
\begin{enumerate}
\item Expliciter un sous-groupe $G_{d}$ de $\Z/n\Z$ d'ordre $d$.
\item Soit $H$ un sous-groupe de $\Z/n\Z$ d'ordre $d$. Montrer que pour tout 
$\overline x\in H$,
$d\cdot \overline x=0.$ Combien d'éléments de $\Z/n\Z$ vérifient cette équation ? 
En déduire que $H=G_{d}$.
\item Conclure.
\end{enumerate}
 \item Donner la liste de tous les sous-groupes de $(\Z/6\Z,+)$.
\end{enumerate}
\finenonce{007304}
\finexercice


\exercice{7305, mourougane, 2021/08/10}
\enonce{007305}{}
Déterminer les puissances de $2$ modulo $9$.
Que dire du groupe $(\Z/9\Z)^\times$ ?
\finenonce{007305}
\finexercice


\exercice{7306, mourougane, 2021/08/10}
\enonce{007306}{}
\begin{enumerate}
 \item Soit $m$ et $n$ deux entiers naturels non nuls premiers entre eux.
On a vu comment trouver une relation de Bézout
$um + vn = 1$. Montrer alors que l'application 
$$\begin{array}{ccc}
  (\Z/m\Z) \times (\Z/n\Z)&\to & \Z/(mn)\Z\\
(x , y) &\mapsto& umy + vnx
 \end{array}$$
est bien définie, est un isomorphisme. Déterminer son isomorphisme réciproque.

 \item Trouver l’entier entre $0$ et $100$ congru à $9$ modulo $11$ et à $3$
modulo $13$.
\end{enumerate}
\finenonce{007306}
\finexercice


\exercice{7307, mourougane, 2021/08/10}
\enonce{007307}{}
 Il est souvent important de calculer $a^t$ modulo $n$ avec $a, t, n$
grands (calculer $a^t$ dans $\Z$ n’est pas envisageable).
Méthode : Ecrire $t$ en binaire : $t =\Sigma t_i2^i$ (où $t^i\in\{0, 1\}$).
Les $a^{2^i}$ se calculent facilement par élévations au carré
successives modulo $n$, et $a^t$ modulo $n$ est le produit (modulo $n$) des $a^{2^i}$ pour lesquels $t^i = 1$.

Calculer $3^{2025}$ modulo $50$.
\finenonce{007307}
\finexercice


\exercice{7308, mourougane, 2021/08/10}
\enonce{007308}{}

Voici la table d'un groupe fini. Déterminer l'ordre de $a$. 

$$\begin{array}{|c|c|c|c|c|}\hline
 * 	& a &	 b & c &	 d \\ \hline
 a 	& c &	 d &	 a &	 b \\ \hline
 b 	& d &	 a &	 b &	 c \\ \hline
 c 	& a &	 b & c &	 d \\ \hline
 d 	& b &	 c &	 d &	 a \\ \hline
 \end{array}$$




\finenonce{007308}
\finexercice


\exercice{7309, mourougane, 2021/08/10}
\enonce{007309}{}
Le but de cet exercice est de déterminer tous les morphismes d'anneaux de $\Z/10\Z$ dans un $\Z/n\Z$.
\begin{enumerate}
 \item Montrer que tout morphisme de groupes de $\Z/10\Z$ dans un $\Z/n\Z$ est déterminé par l'image de $[1]_{10}$.
 \item Quelles sont les valeurs possibles pour l'image de $[1]_{10}$ par un morphisme d'anneaux de $\Z/10\Z$ dans un $\Z/n\Z$ ?
 \item Déterminer tous les morphismes d'anneaux de $\Z/10\Z$ dans un $\Z/n\Z$.
\end{enumerate}
\finenonce{007309}
\finexercice


\exercice{7310, mourougane, 2021/08/10}
\enonce{007310}{}
\begin{enumerate}
 \item L'équation $x^3 + x + 1 = 0$ a-t-elle des solutions dans $\Z/2\Z$.
 \item L'équation $x^3 +x+1 = 0$ a-t-elle des solutions dans $\Z$ ?
\end{enumerate}
\finenonce{007310}
\finexercice


\exercice{7311, mourougane, 2021/08/10}
\enonce{007311}{}
\begin{enumerate}
 \item Quelle opération fait de $\Q/\Z$ un groupe ?
 \item Ce groupe est-il fini ?
 \item Quel est l'ordre de $7/12$ dans ce groupe ?
 \item Montrer que tous les éléments de $\Q/\Z$ sont d'ordre fini.
\end{enumerate}
\finenonce{007311}
\finexercice


\exercice{7312, mourougane, 2021/08/10}
\enonce{007312}{}
\begin{enumerate}
 \item Soit $p$ un nombre premier impair. Quels éléments de $(\Z/p\Z)^\times$ sont leur propre inverse ? 

 \item Démontrer le théorème de Wilson.

 Théorème de Wilson : Soit $n$ un entier naturel supérieur à $3$.
Alors $n$ est premier si, et seulement si, $(n-1)!= -1 \pmod n$.
\end{enumerate}
\finenonce{007312}
\finexercice


\exercice{7313, mourougane, 2021/08/10}
\enonce{007313}{}
\begin{enumerate}
 \item Dans l'anneau $\Z /3 \Z$, calculer $(a+b)^3$ où $a$ et $b$ sont
deux éléments.

 \item Dans l'anneau $\Z /6 \Z$, calculer $(a+b)^6$ où $a$ et $b$ sont
deux éléments.

 \item Dans l'anneau $\Z /7 \Z$, calculer $(a+b)^7$ où $a$ et $b$ sont
deux éléments.
\end{enumerate}
\finenonce{007313}
\finexercice


\exercice{7314, mourougane, 2021/08/10}
\enonce{007314}{}
\begin{enumerate}
 \item $20606$ appartient-il à $14443 + 3079 \Z$ ? 
 \item Calculer l'élément $2169$ dans $\Z/13\Z$. Le résultat doit être représenté par un nombre compris entre $0$ et $12$. 
 \item Considérons une application $f : \Z/8\Z \rightarrow \Z/8\Z$, qui envoie $x$ sur $x^3$. Est-elle injective ?
 \item $[2]_{26}$ est-il un diviseur de $[0]$ dans $\Z/26\Z$ ?
 \item $187489=433^2$, où $433$ est un nombre premier. Combien de diviseurs de zéro y a-t-il dans $\Z/187489\Z$ ? 
 \item Déterminer les puissances de $2$ modulo $9$.
Que dire du groupe $(\Z/9\Z)^\times$ ?
\end{enumerate}
\finenonce{007314}
\finexercice


\exercice{7315, mourougane, 2021/08/10}
\enonce{007315}{}
Soit $D ~:~ \R[X] \rightarrow \R[X]$ l'application de dérivation. 
Est-ce un morphisme de groupes ? Un morphisme d'anneaux ? Une application linéaire ?
\finenonce{007315}
\finexercice


\exercice{7316, mourougane, 2021/08/10}
\enonce{007316}{}
Déterminer le pgcd et le ppcm des polynômes $X^5 - 3X^4 + X^3 + 2X^2 - 6X +2$ et $X^4 - 3X^3 + 3X -1$, éléments de $\Q[X]$.
\finenonce{007316}
\finexercice


\exercice{7317, mourougane, 2021/08/10}
\enonce{007317}{}
Soient $P=X^3-1$ et $Q = X^2-3X+2$ des éléments de $\R[X]$. Quel est l'idéal de $\R[X]$ engendré par $P$ ? L'idéal engendré par $P$ et $Q$ ? L'idéal $(P) \cap (Q)$ ?
\finenonce{007317}
\finexercice


\exercice{7318, mourougane, 2021/08/10}
\enonce{007318}{}
Soit $K$ un corps, et $P \in K[X]$. Les assertions suivantes sont elles vraies ou fausses ?
\begin{enumerate}
 \item Si $P$ n'a pas de racine dans $K$, alors $P$ est irréductible.
 \item Si $P$ est irréductible, alors $P$ n'a pas de racine dans $K$.
\end{enumerate}
\finenonce{007318}
\finexercice


\exercice{7319, mourougane, 2021/08/10}
\enonce{007319}{}
Soit $P = X^4 + 1 \in \C[X]$.
\begin{enumerate}
 \item De quelle forme sont les polynômes irréductibles de $\C[X]$ ? Factoriser $P$ en produit d'irréductibles dans $\C[X]$.
 \item De quelle forme sont les polynômes irréductibles de $\R[X]$ ? Factoriser $P$ en produit d'irréductibles dans $\R[X]$.
 \item Montrer que $P$ ne peut pas se factoriser en produit de deux polynômes de degré 2 à coefficients rationnels (on pourra supposer que c'est possible, et écrire la division euclidienne de $P$ par l'un de ces facteurs). En déduire que $P$ est irréductible dans $\Q[X]$.
\end{enumerate}
\finenonce{007319}
\finexercice


\exercice{7320, mourougane, 2021/08/10}
\enonce{007320}{}
\begin{enumerate}
 \item Soit $n$ un entier naturel non nul et $d$ un diviseur de $n$.
Montrer que $X^d-1$ divise $X^n-1$.
 \item Soit $n$ un entier naturel non nul. Soit $d$ un entier naturel non nul et $r$ le reste de la division euclidienne de $n$ par $d$. Montrer que $X^r-1$ est le reste de la division euclidienne de $X^n-1$ par $X^d-1$. 
 \item Soient $m,n \in \N^*$, et $d = \pgcd(m,n)$. Montrer que $\pgcd(X^m-1,X^n-1) = X^d -1$.
\end{enumerate}
\finenonce{007320}
\finexercice


\exercice{7321, mourougane, 2021/08/10}
\enonce{007321}{}
Déterminer les racines de $X^2 - 1$ dans $\Z/8\Z$. Comparer leur nombre au degré du polynôme. Comment expliquer ce phénomène ?
\finenonce{007321}
\finexercice


\exercice{7322, mourougane, 2021/08/10}
\enonce{007322}{}
Soit $P = a_0 + a_1X + \cdots + a_nX^n \in \Z[X]$. On veut déterminer si $P$ a des racines rationnelles.
\begin{enumerate}
 \item On suppose que $P$ a une racine rationnelle non nulle $x$, avec $x = \frac{p}{q}$ et $\pgcd(p,q) = 1$. Montrer que $p$ divise $a_0$ et $q$ divise $a_n$.
 \item Le polynôme $7X^3 - 5X^2 -9X +4$ a-t-il des racines rationnelles ? et $X^4 - 2X^2 - 3$ ?
 \item Soit $n \in \N^*$. Montrer que $\sqrt{n}$ est soit un entier, soit un irrationnel.
\end{enumerate}
\finenonce{007322}
\finexercice


\exercice{7323, mourougane, 2021/08/10}
\enonce{007323}{}
 Montrer que $X^2 + 4$ est irréductible dans $\Z[X]$ mais réductible dans $\Z/2\Z[X]$.
\finenonce{007323}
\finexercice


\exercice{7333, mourougane, 2021/08/10}
\enonce{007333}{}
\begin{enumerate}
 \item Soit $p=13$. Trouver une racine $c$ de $-1$ modulo $13$.
 \item Représenter le réseau $\mathcal{R}$ de ${\Z}^2$ engendré par $(c,1)$ et $(p,0)$ et montrer que tous ses éléments ont une norme multiple de $p$.  
 \item Trouver deux entiers $x$ et $y$ tels que $p=x^2+y^2$.
\end{enumerate}
\finenonce{007333}
\finexercice


\exercice{7334, mourougane, 2021/08/10}
\enonce{007334}{}
\begin{enumerate}
 \item Le nombre $101$ est-il premier ? Peut-il s'écrire comme somme de deux carrés.
 \item Déterminer une racine $c$ de $-1$ modulo $101$.
 \item Calculer le $\pgcd(101,c+i)$ dans $\Z[i]$.
 \item Ecrire $101$ comme somme de deux carrés.
 \item Mêmes questions avec $2011$.
\end{enumerate}
\finenonce{007334}
\finexercice


\exercice{7335, mourougane, 2021/08/10}
\enonce{007335}{}
\begin{enumerate}
 \item Quelle identité obtient-on quand on écrit que la norme de $(a+ib)(c+id)$ dans $\Z[i]$ est le produit de la norme de $(a+ib)$ et de celle de $(c+id)$ ?
 \item \'Ecrire $2425=5^2\cdot 97$ et $754=2\cdot 13 \cdot 29$ comme sommes de deux carrés.
 \item Tous les entiers naturels sont-ils sommes de trois carrés ?
%1. Montrer que tout entier positif n de la forme 8k + 7 ne peut pas s'écrire comme somme
%de 3 carrés.
%Indication : on pourra raisonner modulo 8...
%2. Montrer que si n est un entier positif tel que 4n est somme de 3 carrés, alors n l'est aussi.
%Indication : après avoir écrit 4n comme somme de 3 carrés, on pourra étudier la parité
%de chacun des termes de la somme et raisonner modulo 4.
%3. A l'aide des deux questions précédentes, deviner une condition qui emp^eche un entier
%positif de s'écrire comme somme de 3 carrés.
\end{enumerate}
\finenonce{007335}
\finexercice


\exercice{7336, mourougane, 2021/08/10}
\enonce{007336}{}
Considérons l'anneau $\Z[i\sqrt{3}]$, le sous-anneaux de $\C$ engendré par $i\sqrt{3}$. 
\begin{enumerate}
 \item Montrer que $\Z[i\sqrt{3}]:=\{a+bi\sqrt{3}, \ (a,b)\in\Z^2\}.$
 \item \`A tout élément $x=a+bi\sqrt{3}$ de $\mathbf{Z}[i\sqrt{3}]$ on associe son conjugué $\overline x=a-bi\sqrt{3}$.
 Montrer que pour tous $x,y \in \mathbf{Z}[i\sqrt{3}]$ on a 
$$\overline{x+y}=\overline x + \overline y \textrm{ ~ et ~ } \overline{xy}=\overline x \overline y.$$
 \item En considérant l'application $N : \mathbf{Z}[i\sqrt{3}]\to \mathbf{Z}, x \mapsto x\overline x$, montrer que les éléments inversibles de $\mathbf{Z}[i\sqrt{3}]$ sont exactement les éléments de norme $1$. Donner la liste des éléments inversibles. 
 \item Quelles sont les normes possibles d'un diviseur de $1+i\sqrt{3}$ dans $\mathbf{Z}[i\sqrt{3}]$ ?
 \item Montrer que éléments $2$, $1+i\sqrt{3}$ sont premiers entre eux dans $\mathbf{Z}[i\sqrt{3}]$.
Peut-on leur écrire un couple de Bézout ?
 \item Donner deux factorisations différentes de $4$ dans $\mathbf{Z}[i\sqrt{3}]$. 
Le lemme d'Euclide est-il valide dans l'anneau $\mathbf{Z}[\sqrt{-3}]$ ?
\end{enumerate}
\finenonce{007336}
\finexercice


\exercice{7337, mourougane, 2021/08/10}
\enonce{007337}{}
Les polynômes symétriques élémentaires en $n$ indéterminées sont par définition
\begin{eqnarray*}
 s_1(t_1,t_2,\cdots,t_n)&:=&t_1+t_2+\cdots+t_n\\
s_2(t_1,t_2,\cdots,t_n)&:=&t_1t_2+t_1t_3+\cdots+t_2t_3+t_2t_4
+\cdots +\cdots t_{n-1}t_n\\
\vdots\\
s_n(t_1,t_2,\cdots,t_n)&:=&t_1t_2\cdots t_n.
\end{eqnarray*}

\begin{enumerate}
 \item Soit $F(X)=a(X-t_1)(X-t_2)\cdots(X-t_n)$. Développer $F(X)$ en puissances de $X$.
 \item Ecrire $t_1^3+t_2^3$ comme polynôme à coefficients entiers de $s_1$, $s_2$. 
 \item Ecrire $t_1^4+t_2^4$ comme polynôme à coefficients entiers de $s_1$, $s_2$. 
 \item Ecrire $t_1^3+t_2^3+t_3^3$ comme polynôme à coefficients entiers de $s_1$, $s_2$, $s_3$. 
 \item Ecrire $f$
$$f(t_1,t_2,t_3)=t_1^2t_2+t_1^2t_3+t_2^2t_1+t_2^2t_3+t_3^2t_1+t_3^2t_2.$$
 comme polynôme à coefficients entiers en les polynômes symétriques élémentaires.
\end{enumerate}
\finenonce{007337}
\finexercice


\exercice{7338, mourougane, 2021/08/10}
\enonce{007338}{}
 \'Ecrire en base 10 le nombre dont l'écriture hexadécimale est 
$DA582$.
\finenonce{007338}
\finexercice


\exercice{7339, mourougane, 2021/08/10}
\enonce{007339}{}
Calculer $P(2)$ où $P=3x^4-x^2-16x-14$ par la méthode de Hörner.
Effectuer la division euclidienne de $P$ par $x-2$.
\finenonce{007339}
\finexercice


\exercice{7340, mourougane, 2021/08/10}
\enonce{007340}{}
Le but de l'exercice est de factoriser le polynôme $x^4 + 4x^3 - 81x^2 -16x + 308$.
\begin{enumerate}
 \item Vérifier que 2 est racine de $P$ et factoriser $P$ par $x-2$ avec la méthode de Hörner.
 \item Vérfier que $-2$ est racine du quotient obtenu
dans la question précédente et factoriser le par $x+2$ par la méthode de Hörner.
 \item Conclure.
%(x - 2)(x + 2)(x - 7)(x + 11)
\end{enumerate}
\finenonce{007340}
\finexercice


\exercice{7341, mourougane, 2021/08/10}
\enonce{007341}{}
Soit $P=3x^5 - 6x^4 + x^3 + 5x^2 - 3x - 4$.
\begin{enumerate}
 \item Effectuer la division euclidienne de $P$ par $x-1$.
 \item Présenter les résultats selon la méthode de Hörner.
 \item \'Ecrire $P$ dans la base $1,(x-1),(x-1)^2,\ldots, (x-1)^5$.
%3(x - 1)5 + 9(x - 1)4 + 7(x - 1)3 + 2(x - 12) + 1(x - 1) - 4
\end{enumerate}
\finenonce{007341}
\finexercice


\exercice{7342, mourougane, 2021/08/10}
\enonce{007342}{}
Le but de l'exercice est de déterminer une valeur approchée de la racine cubique réelle $c$ de $17$.
On considère le polynôme $P(x)=x^3-17$.
\begin{enumerate}
 \item Montrer que $2<c<3$.
 \item Déterminer par la méthode de Hörner le polynôme $R(z):=10^3P(2+\frac{z}{10})$ en la variable $z$.
 \item Déterminer une valeur approchée de $c$ à $10^{-1}$ près.
 \item Déterminer une valeur approchée de $c$ à $10^{-2}$ près.
%(réponse : $2,57<c<2,58$.)
\end{enumerate}
\finenonce{007342}
\finexercice


\exercice{7345, mourougane, 2021/08/10}
\enonce{007345}{}
\begin{enumerate}
    \item L'entier $256$ appartient-il à $115 + 247\Z$ ? 
    \item L'entier $-601$ est-il un représentant de la classe $[-738]_{ 28}$ de $\Z/28\Z$? 
    \item Calculer l'élément $589$ dans $\Z/23\Z$. Le résultat doit être représenté par un nombre compris entre $0$ et $22$. 
    \item Calculer l'élément $13 ^{923}$ dans $\Z/11\Z$. Le résultat doit être représenté par un nombre compris entre $0$ et $10$. 
\end{enumerate}
\finenonce{007345}
\finexercice


\exercice{7346, mourougane, 2021/08/10}
\enonce{007346}{}
\begin{enumerate}
 \item Montrer que le polynôme $X^9-1$ de $\mathbb{F}_3[X]$ vaut $(X-1)^9$. On considère le code ternaire $C$ de longueur $9$ associé au polynôme $g=(X-1)^5$ 
 \item Déterminer l'alphabet, la longueur des mots, la dimension du code, le nombre de mots de code. Le code est-il cyclique ?
 \item Donner une matrice génératrice de $C$.
 \item Déterminer un élément de poids $3$ du code.
 \item Déterminer une matrice de contrôle $H$ de ce code.
 \item Déterminer la distance de ce code. Combien d'erreurs ce code peut-il détecter ? combien d'erreurs peut-il corriger ?
 \item On a reçu le mot $r=121102210$. Calculer son image par $H$. Le mot $r$ est-il un mot du code ?
 \item Corriger le mot $r$ en supposant qu'il n'y a eu au plus qu'une seule erreur de transmission.
\end{enumerate}
\finenonce{007346}
\finexercice


\exercice{7347, mourougane, 2021/08/10}
\enonce{007347}{}
\begin{enumerate}
 \item Le nombre $101$ est-il premier ? Peut-il s'écrire comme somme de deux carrés.
 \item Déterminer une racine $c$ de $-1$ modulo $101$.
 \item Calculer le $pgcd(101,c-i)$ dans $\Z[i]$.
 \item Ecrire $101$ comme somme de deux carrés.
\end{enumerate}
\finenonce{007347}
\finexercice


\exercice{7368, mourougane, 2021/08/10}
\enonce{007368}{Sur le nombre de solutions d'une équation du second degré}

\textit{Le but de ce problème est de montrer que l'équation $X^2+X+1=0$ peut avoir un nombre arbitrairement grand de solutions dans les anneaux $\Z/n\Z$.}


\textit{Pour un nombre premier $p$, on notera $\mathbb{F}_p$ le corps $\Z/p\Z$.
On rappelle qu'alors le groupe $(\mathbb{F}_p^\times,\times)$ des inversibles de $\mathbb{F}_p$ est cyclique.
On pourra admettre le résultat d'une question pour continuer.}
\begin{enumerate}
 \item Si $p$ est un nombre premier, combien au maximum l'équation $X^2+X+1=0$ a-t-elle de solutions dans l'anneau $\Z/p\Z$ ?

  \item Déterminer les solutions de l'équation $X^2+X+1=0$ 
\begin{enumerate}
\item [a] dans $\Z/2\Z$ , dans $\Z/2n\Z$. 
\item [b] dans $\Z/7\Z$.
\item [c] dans $\Z/13\Z$.
\item [d] \`A l'aide des questions précédentes, déterminer les solutions de l'équation\\ $X^2+X+1=0$ dans $\Z/91\Z$.
\end{enumerate}


 \item Soit $p$ est un nombre premier plus grand que $5$.
Montrer que l'équation $X^2+X+1=0$ admet une solution dans le corps $\mathbb{F}_p$ si et seulement si l'équation $X^3=1$ admet une solution \textsl{différente de $1$} dans $\mathbb{F}_p$.

 \item Soit $p$ est un nombre premier plus grand que $5$.
Montrer que l'équation $X^2+X+1=0$ admet une solution dans le corps $\mathbb{F}_p$ si et seulement si $(\mathbb{F}_p)^\times$ a un élément d'ordre $3$.

 \item Soit $p$ est un nombre premier plus grand que $5$.
Montrer que l'équation $X^2+X+1=0$ admet une solution dans le corps $\mathbb{F}_p$ si et seulement si $p$ est congru à $1$ modulo $3$.

 \item Soit $p$ est un nombre premier plus grand que $5$.
Montrer que si l'équation $X^2+X+1=0$ admet une solution $a$ dans le corps $\mathbb{F}_p$ elle admet une autre solution disctincte $-1-a$.

 \item On admet qu'il existe une infinité de nombres premiers congrus à $1$ modulo~$3$.
Montrer que l'équation $X^2+X+1=0$ peut avoir un nombre arbitrairement grand de solutions dans les anneaux $\Z/n\Z$.
%son discriminant réduit $\Delta'=-3$ est un carré dans $\Z/p\Z$,
%si et seulement si $(-3)^\frac{p-1}{2}=1$

\end{enumerate}
\finenonce{007368}
\finexercice


\exercice{7369, mourougane, 2021/08/10}
\enonce{007369}{}
\begin{enumerate}
 \item Soit $G$ un groupe, $e$ son élément neutre et $a$ un élément de $G$ et $m$ un entier naturel tel $a^m=e$. Que peut-on dire de l'ordre de l'élément $a$ ?
 \item Donner la définition d'un idéal dans un anneau commutatif.
 \item Quels sont les idéaux de l'anneau $(\Z,+,\times)$ ?
 \item Le sous-ensemble $28\Z\cap 18\Z$ de $\Z$ est-il un idéal de $\Z$ ? Si oui, le déterminer.
 \item Donner l'exemple de deux groupes finis de même ordre non isomorphes.
Justifier le fait que les deux groupes choisis ne sont pas isomorphes.
 \item Le groupe $\mathcal{S}_4$ des permutations de $\{1,2,3,4\}$ est-il cyclique ? (justifier)
\end{enumerate}
\finenonce{007369}
\finexercice


\exercice{7370, mourougane, 2021/08/10}
\enonce{007370}{}
\begin{enumerate}
 \item L'entier $-1601$ est-il un représentant de la classe $[-7387]_{ 2893}$ de $\Z/2893\Z$? 
 \item Calculer l'élément $11 ^{329}$ dans $\Z/13\Z$. Le résultat doit être représenté par un nombre compris entre $0$ et $12$. 
 \item La classe $[51]$ est-elle inversible dans l'anneau $\Z/131\Z$. Si oui, calculer son inverse dans $\Z/131\Z$. Le résultat doit être représenté par un nombre compris entre $0$ et $130$. 
\end{enumerate}
\finenonce{007370}
\finexercice


\exercice{7371, mourougane, 2021/08/10}
\enonce{007371}{}
\begin{enumerate}
 \item Ecrire une relation de Bezout entre $X^2+X+1$ et $X^3+X+1$ dans $\R[X]$.
 \item La classe du polynôme $X^2+X+1$ est-elle inversible dans l'anneau quotient\\
$\R[X]/(X^3+X+1)$ ? Si oui, donner son inverse.
\end{enumerate}
\finenonce{007371}
\finexercice


\exercice{7372, mourougane, 2021/08/10}
\enonce{007372}{}
\begin{enumerate}
  \item Quels sont les éléments inversibles de l'anneau $(\Z/12\Z)$ ?
  \item \`A quel groupe le groupe $(\Z/12\Z)^\times$ des inversibles de l'anneau $(\Z/12\Z)$ est-il isomorphe ?
\end{enumerate}
\finenonce{007372}
\finexercice


\exercice{7373, mourougane, 2021/08/10}
\enonce{007373}{}
Voici la table d'un groupe $G$.
Quel est l'ordre de $G$ ? Le groupe $G$ est-il nécessairement commutatif ?
 Compléter la table en énonçant précisément les propriétés utilisées.
$$\begin{array}{|c|c|c|c|c|c|}\hline
 * & a & b 	& c &	d & e\\ \hline
 a & 	 & 	& d &	 & c \\ \hline
 b & e &	&a  &	 & d \\ \hline
 c &	 &a 	 &  &  &	 \\ \hline
 d &	 &	 &  &d &e \\ \hline
 e &	 &	 &b &	 & a \\ \hline
\end{array}
$$

\finenonce{007373}
\finexercice


\exercice{7374, mourougane, 2021/08/10}
\enonce{007374}{}
\begin{enumerate}
 \item Calculer les produits dans $\mathcal{S}_7$, $(1,2)(1,3)$ et $(1,2)(2,3)(1,2)$.
 \item L'ensemble des transpositions de $\mathcal{S}_7$ est-il un sous-groupe de $\mathcal{S}_7$ ?
\end{enumerate}
\finenonce{007374}
\finexercice

\section{203.05 Idéal}


\exercice{1368, ortiz, 1999/04/01}

\enonce{001368}{}
\begin{enumerate}
    \item  $\mathcal{J}=\left\{ \left( \alpha ,\alpha \right) :\alpha \in \Zz%
\right\} $ est-il un id\'eal de l'anneau $\Zz^2$?
    \item  $\mathcal{J}=\left\{ P\in \Rr\left[ X\right] :P^{^{\prime
}}\left( 0\right) =0\right\} $ est-il un id\'eal
de $\Rr\left[ X\right] $?
\end{enumerate}
\finenonce{001368}



\finexercice



\exercice{1371, ortiz, 1999/04/01}

\enonce{001371}{}
Soit $\mathcal{J}=\left\{ P\in
\Zz\left[ X\right] :P\left( 0\right) \in 2%
\Zz\right\} .$
\begin{enumerate}
    \item
    \begin{enumerate}
         \item  Montrer que $\mathcal{J}$ est un id\'eal de $\Zz\left[ X\right] .$
         \item  Montrer que $\mathcal{J}$ est engendr\'e par les polyn\^omes $2$ et $X
$.
\end{enumerate}
    \item  En
remarquant que $2\in \mathcal{J}$, montrer que
l'hypoth\`ese ``$\mathcal{J}$ est un id\'eal
principal de $\Zz[X]$'' est absurde.
\end{enumerate}
\finenonce{001371}



\finexercice



\exercice{1375, gourio, 2001/09/01}

\enonce{001375}{}
Soit $(A,+,\times )$ un anneau commutatif, on dit que $I\subset A $ est un
id\'{e}al de $A $ si et seulement si: $I  $ est un sous-groupe de $(A,+)$
et de plus: $\forall a\in A,\forall x\in I,ax\in I.$
\begin{enumerate}
\item Quels sont les id\'{e}aux de $\Zz$ ?
\item On appelle radical de $I$, l'ensemble :
$$\sqrt{I} =\{x\in A|\exists n\in \Nn,x^{n}\in I\}.$$
Montrer que $\sqrt{I\text{ }}$ est un id\'{e}al de $A $contenant $I$.
\'Etudier le cas $A=Z.$
\item Montrer que si  $I$ et $J $ sont deux id\'{e}aux de A tels que $I\subset J$,
alors $\sqrt{I}\subset \sqrt{J}. $En d\'{e}duire $\sqrt{\sqrt{I}}=\sqrt{I}.$
\item Montrer que si $I$ et $J $sont deux id\'{e}aux de A, $\sqrt{I\cap J}=\sqrt{I}\cap \sqrt{J}. $
\end{enumerate}
\finenonce{001375}



\finexercice



\exercice{1376, legall, 2003/10/01}

\enonce{001376}{}
$A$ est nomm\'e id\'eal de $A$
lorsque pour tout $x\in J$ et tout $a\in A$ le produit $ax$ appartient \`a $J$.
\begin{enumerate}
    \item Trouver tous les id\'eaux d'un corps $\mathbb{K}$.
    \item Montrer que tout id\'eal de $\Zz$ est de la forme $a\Zz ,$ 
o\`u $a\in \Zz .$
    \item On note $D $ l'ensemble des rationnels $x$ tels que il 
existe $k\in \Nn $ tel que $x10^{k}\in \Zz .$
Montrer que tout id\'eal de $D $ est de la forme
$aD$ o\`u $a\in D.$\end{enumerate}
\finenonce{001376}



\finexercice



\exercice{1571, ortiz, 1999/04/01}

\enonce{001571}{}
Montrer qu'un id\'eal de $K[X]$ est distinct de
$K[X]$ si et seulement s'il ne contient aucun
polyn\^ome constant non nul.
\finenonce{001571}



\finexercice



\exercice{1572, ortiz, 1999/04/01}

\enonce{001572}{}
Soient les polyn\^omes $P=X^4+X^3-2X+1$ et
$Q=X^2+X+1$ de $\Rr[X]$. D\'eterminer $\mathrm{pgcd}(P,Q)
$ puis la somme et l'intersection des id\'eaux
principaux $(P)$ et $(Q)$ de $\Rr[X]$.
\finenonce{001572}



\finexercice



\exercice{1573, ortiz, 1999/04/01}

\enonce{001573}{}
Les parties $\mathcal{I}=\left\{P\in\Rr[X]:P'(0)=0\right\}$
et $\mathcal{J}=\left\{P\in\Rr[X]:P(0)=P'(0)=0\right\}$
sont-elles des id\'eaux de $\Rr[X]$ ? Dans
l'affirmative, en donner un g\'en\'erateur.
\finenonce{001573}



\finexercice


\section{203.06 Algèbre, corps}


\exercice{1377, ridde, 1999/11/01}

\enonce{001377}{}
D\'eterminer les automorphismes du corps $\Qq$.
\finenonce{001377}



\finexercice



\exercice{1378, ridde, 1999/11/01}

\enonce{001378}{}
Soit $\sigma$ un automorphisme de $\Rr$.
\begin{enumerate}
\item  Montrer que si $x \geq 0$ alors $\sigma (x) \geq 0$.
\item Montrer que $\sigma$ est croissante.
\item D\'eterminer $\sigma$.
\end{enumerate}
\finenonce{001378}



\finexercice



\exercice{1379, ortiz, 1999/04/01}

\enonce{001379}{}
Soient $A=\left(\begin{smallmatrix}
1&1\\0&1\end{smallmatrix}\right)$ et
$C=\left\{M\in\mathcal{M}_2(\Rr): MA=AM\right\}.$
\begin{enumerate}
    \item  Montrer que $C$ est un sous-espace vectoriel de $\mathcal{M}_2(\Rr)$
 et en d\'eterminer une base.
    \item  Montrer que, pour les lois usuelles, $C$ est une $\Rr$-alg\`ebre.
\end{enumerate}
\finenonce{001379}



\finexercice



\exercice{1380, ortiz, 1999/04/01}

\enonce{001380}{}
  Soient $E$ un $\Rr$-espace vectoriel et $u\in \mathcal{L} (E)$ tel que $u^2=u$. On d\'efinit
$${\Rr[u]:=\left\{P(u):P\in\Rr[X]\right\}.}$$
\begin{enumerate}
    \item  Montrer que, muni des lois usuelles sur $\mathcal{L}(E)$, c'est une $\Rr$-alg\`ebre.
    \item  Montrer que cette alg\`ebre est de dimension finie et discuter de sa dimension en fonction de $u$.
    \item  L'anneau $\Rr[u]$ est-il un corps ?
\end{enumerate}
\finenonce{001380}



\finexercice



\exercice{1381, ortiz, 1999/04/01}

\enonce{001381}{}
Soit
$M=\left\{aI_2+bJ\in\mathcal{M}_2(\Rr):a,b\in\Rr\right\}$
o\`u $I_2=
\left (
\begin{array}{ll}
1&0\\
0&1
\end{array}
\right ),
J=
\left (
\begin{array}{ll}
0&2\\
1&0
\end{array}
\right )$.
\begin{enumerate}
    \item  Calculer $J^2$ et montrer que si $a,b\in \Rr$ et $ aI_2+ bJ =O$ alors $a=b=0$.
    \item  Montrer que, muni des lois usuelles sur $\mathcal{M}_2(\Rr)$,
$M$ est un anneau. Cet anneau est-il commutatif,
int\`egre ?
    \item  $M$ est-il un corps, une $\Rr$-alg\`ebre ?
\end{enumerate}
\finenonce{001381}



\finexercice



\exercice{1382, ortiz, 1999/04/01}

\enonce{001382}{}
  Montrer que l'ensemble $S$ des suites r\'eelles convergentes est une $\Rr$-alg\`ebre.
  L'application $S\to \Rr,  u\mapsto \lim{u}$ est-elle un morphisme de $\Rr$-alg\`ebres ?
  L'anneau $S$ est-il int\`egre ?
\finenonce{001382}



\finexercice



\exercice{1383, ortiz, 1999/04/01}

\enonce{001383}{}
 Soient $E$ un $\Rr$-ev et $u\in \mathcal{L} (E)$
 tel que $u^2=u$. On d\'efinit
$${\Rr[u]=\left\{a\text{Id}_E+bu:a,b\in \Rr\right\}.}$$
 Montrer que,
muni des lois usuelles sur $\mathcal{L}(E)$, c'est
une $\Rr$-alg\`ebre. L'anneau $\Rr[u]$ est-il un
corps ?
\finenonce{001383}



\finexercice



\exercice{1384, legall, 2003/10/01}

\enonce{001384}{}
Un automorphisme d'un corps $\mathbb{K} $ est une 
application bijective $\varphi $ de $\mathbb{K} $ dans lui-m\^eme
telle que $\varphi (1)=1,$ $\varphi (0)=0$ et, pour tout $a,b \in \mathbb{K} 
,$ on ait $\varphi (a+b)=\varphi (a)+\varphi (b)$ et $\varphi 
(ab)=\varphi (a)\varphi (b).$

\begin{enumerate}
\item Soit $\varphi $ un automorphisme de $\Rr .$ Montrer que 
l'application $x\mapsto \varphi (x)$ est croissante. En d\'eduire que 
l'identit\'e est le seul automorphisme de $\Rr $.
\item Soit $\psi$ un automorphisme {\em continu} de $\Cc .$ Montrer 
$\psi (x)= x,$ pour tout $x\in \Rr .$ En d\'eduire tous les 
automorphismes {\em continus} de $\Cc$.
\end{enumerate}
\finenonce{001384}



\finexercice



\exercice{3025, quercia, 2010/03/08}

\enonce{003025}{Anneau int{\`e}gre fini}
Soit $A$ un anneau non nul, commutatif et int{\`e}gre.
\begin{enumerate}
  \item Montrer que si $A$ est fini, alors c'est un corps.
  \item Montrer que si $A$ n'a qu'un nombre fini d'id{\'e}aux, alors c'est un corps
    (consid{\'e}rer les id{\'e}aux $I_n = x^nA$ pour $x\in A$ non nul).
\end{enumerate}
\finenonce{003025}



\finexercice   


\exercice{3026, quercia, 2010/03/08}

\enonce{003026}{Corps ${\mathbb F}_4$}
Chercher les structures de corps {\`a} 4 {\'e}l{\'e}ments.
\finenonce{003026}



\finexercice   


\exercice{3027, quercia, 2010/03/08}

\enonce{003027}{Groupe multiplicatif d'un corps fini}
Soit $K$ un corps fini. Pour $x\in K^*$ on note $O(x)$ l'ordre multiplicatif
de~$x$ et $n$ le ppcm des ordres des {\'e}l{\'e}ments de~$K^*$.
\begin{enumerate}
  \item Soient $a,b\in\N^*$. Montrer qu'il existe $a',b'\in\N^*$ tels
    que $a'|a$, $b'|b$, $a'\wedge b' = 1$ et $a'b' = a\vee b$.
  \item Soient $x,y\in K^*$ d'ordres $a$ et $b$. Montrer qu'il existe $u,v$ entiers tels que
    $O(x^uy^v) = a\vee b$. En d{\'e}duire qu'il existe $z\in K^*$ d'ordre~$n$.
  \item Montrer que $n=\mathrm{Card}\,( K^*)$ (ceci prouve que $ K^*$ est cyclique).

\end{enumerate}
\finenonce{003027}



\finexercice 


\exercice{3028, quercia, 2010/03/08}

\enonce{003028}{Groupe multiplicatif d'un corps fini}
Soit $ K$ un corps fini de cardinal $n$.
Si $a,b\in \N$ sont tels que $ab = n-1$, on consid{\`e}re l'application
${f_a} : { K^*} \to { K^*}, x \mapsto {x^a}$ (remarquer que $f_a$ est
un morphisme de groupe). On note $N_a = \mathrm{Card}\,(\mathrm{Ker} f_a)$.
\begin{enumerate}
  \item Expliquer pourquoi $N_a \le a$.
  \item Montrer que $\Im(f_a) \subset \mathrm{Ker} f_b$. En d{\'e}duire que $N_a = a$ et $N_b=b$.
  \item Soit $\varphi$ l'indicateur d'Euler. Montrer par r{\'e}currence sur $a$,
    diviseur de $n-1$, que le nombre d'{\'e}l{\'e}ments de $ K^*$ d'ordre~$a$ est
    {\'e}gal {\`a} $\varphi(a)$ (ceci prouve que $ K^*$ est cyclique).
\end{enumerate}
\finenonce{003028}



\finexercice 


\exercice{3029, quercia, 2010/03/08}

\enonce{003029}{Th{\'e}or{\`e}me de Wedderburn}
On dit que $K$ est un corps gauche si $(K,+,\times)$ est un anneau et si $(K\setminus \{0\},\times)$ 
est un groupe (non n{\'e}cessairement commutatif). On v{\'e}rifiera rapidement que la th{\'e}orie des espaces vectoriels est inchang{\'e}e si on 
remplace le corps de base par un corps gauche.
L'objet de l'exercice est de d{\'e}montrer le th{\'e}or{\`e}me de Wedderburn~:
{\it tout corps gauche fini est commutatif.}

Pour $n\in \N ^*$, soit ${\cal P}_n$ l'ensemble des racines $n$-{\`e}mes primitives de l'unit{\'e} dans $\C$.
On pose $\Phi_1(X)=X-1$ et $\Phi_n(X)=\prod_{\zeta \in {\cal P}_n}(X-\zeta)$.
$\Phi_n$ est appel{\'e} {\it le $n$-{\`e}me polyn{\^o}me cyclotomique}
(son degr{\'e} est $\phi (n)$ o{\`u} $\phi$ est l'indicateur d'Euler).
\begin{enumerate}
  \item D{\'e}montrer : $(\forall\ n\in \N ^*)\ X^n-1 = \prod_{d|n} \Phi_d(X)$. 
    En d{\'e}duire, par r{\'e}currence, que $\Phi_n(X)$ a tous ses coefficients dans $\Z$.
  \item Calculer explicitement $\Phi_n(X)$ pour $n\le 16$.
  \item D{\'e}montrer que, pour $p$ premier et $\alpha \in \N ^*$, $\Phi_{p^\alpha}(X)=
    \sum_{k=0}^{p-1}X^{kp^{\alpha -1}}$.
  \item Calculer le terme constant de chaque $\Phi_n$.
  \item Montrer que, si $d<n$ et $d$ divise $n$, alors $X^d-1$ divise
    $X^n-1$ dans $\Z [X]$, puis que $\Phi_n(X)$ divise $X^n-1$ et
    $\frac{X^n-1}{X^d-1}$ dans $\Z [X]$.
\\
On consid{\`e}re $K$ un corps gauche fini et $Z(K)$ son centre, de cardinal $q$.
  \item Montrer que $Z(K)$ est un corps commutatif. 
  \item Montrer que $K$ est un $Z(K)$-espace vectoriel de dimension finie, not{\'e}e $n$.
    Donner alors le cardinal de $K$ en fonction de $q$ et~$n$.
  \item Soit $a\in K\setminus \{0\}$. On note $C_a=\{x\in K\,|\,ax=xa\}$. 
    \\
    Montrer que $C_a$ est un corps gauche, puis que c'est un
    $Z(K)$-espace vectoriel de dimension finie $d$ divisant $n$ (on
    montrera pour cela que $K$ est un $C_a$-espace vectoriel et l'on
    {\'e}tudiera sa dimension).
  \item On fait op{\'e}rer le groupe multiplicatif $K^*$ sur lui-m{\^e}me par automorphismes int{\'e}rieurs. 
    \\
    En consid{\'e}rant les orbites selon cette op{\'e}ration montrer que l'on a :
$$q^n-1=q-1+\sum_{i=1}^k \frac{q^n-1}{q^{d_i}-1} \text{ avec, pour tout }i,\ d_i|n.$$
  \item En d{\'e}duire que $\Phi_n(q)$ divise $q-1$.
  \item En {\'e}tudiant $|\Phi_n(q)|$ montrer que $n=1$.
\end{enumerate}
\finenonce{003029}



\finexercice 


\exercice{3324, quercia, 2010/03/09}
\enonce{003324}{\'Eléments algébriques}
Soient $ K, \mathbb{L}$ deux corps avec $ K \subset \mathbb{L}$.

Un élément $\alpha \in \mathbb{L}$ est dit algébrique sur $ K$ s'il existe
un polynôme non nul $P \in  K[X]$ tel que $P(\alpha) = 0$.

\begin{enumerate}
  \item Montrer que $\alpha$ est algébrique sur $ K$ si et seulement si $ K[\alpha]$ est
      un $ K$-ev de dimension finie.

  \item On suppose que $\alpha$ et $\beta$ sont algébriques sur $ K$.
      Montrer que $\alpha + \beta$ et $\alpha\beta$ sont algébriques sur $ K$
      (étudier $ K[\alpha,\beta]$).

\end{enumerate}
\finenonce{003324}



\finexercice


\exercice{3325, quercia, 2010/03/09}
\enonce{003325}{Corps emboîtés}
Soient $H \subset  K \subset  L$ trois sous-corps de $\C$.

\begin{enumerate}
  \item Montrer que $ K$ et $ L$ sont des $ H$-ev et $ L$ est un $ K$-ev.
  \item Montrer que $ L$ est de dimension finie sur $ H$ si et seulement si
    $ K$ est de dimension finie sur $ H$ et
    $ L$ est de dimension finie sur $ K$.
  \item Application : Montrer que $\overline\Q$, la cloture algébrique de $\Q$ dans~$\C$, est
    un corps algébriquement clos (si $P\in\overline\Q[X]$, considérer le sous-corps
    de~$\C$ engendré par les coefficients de~$P$).
\end{enumerate}
\finenonce{003325}



\finexercice


\exercice{3326, quercia, 2010/03/09}
\enonce{003326}{Surcorps de $\R$}
Soit $\mathbb{A}$ une $\R$-algèbre commutative, intègre et de dimension finie.

\begin{enumerate}
  \item Montrer que $\mathbb{A}$ est un corps.
  \item Si $\dim{\mathbb{A}} > 1$ montrer que tout élément de $\mathbb{A}$ est
    algébrique de degré 1 ou 2 sur $\R$. En déduire qu'alors $\mathbb{A}$
    est isomorphe à~$\C$.
\end{enumerate}
\finenonce{003326}



\finexercice


\exercice{3352, quercia, 2010/03/09}
\enonce{003352}{Sous algèbres}

Soit $E$ un ev de dimension finie et $\cal A$ une sous-algèbre de $\mathcal{L}(E)$.
Montrer que si $f \in \cal A$ et $f$ est bijective, alors $f^{-1} \in \cal A$.
\par
On pourra étudier l'application
$\phi : {\cal A} \to {\cal A}, g \mapsto {f\circ g.}$
\finenonce{003352}



\finexercice

\exercice{7324, mourougane, 2021/08/10}
\enonce{007324}{}
\begin{enumerate}
 \item Donner la définition d'un groupe.
 \item Dire si les couples suivants sont des groupes : $(\Z,+)$ ; $(\Z,\times)$ ; $(\C^*,+)$ ; $(\C^*, \times)$. Lorsque la réponse est "non", on indiquera une propriété des groupes qui fait défaut (on ne demande pas de justification lorsque la réponse est "oui").
 \item Quels sont les idéaux de l'anneau $(\Z,+,\times)$ ?
 \item Quels sont les idéaux de l'anneau $(\R[X],+,\times)$ ?
\end{enumerate}
\finenonce{007324}
\finexercice


\exercice{7325, mourougane, 2021/08/10}
\enonce{007325}{}
\begin{enumerate}
 \item L'entier $51606$ appartient-il à $2569 + 247\Z$ ? 
 \item L'entier $-1601$ est-il un représentant de la classe $[-7387]_{ 2893}$ de $\Z/2893\Z$? 
 \item Calculer l'élément $2169$ dans $\Z/17\Z$. Le résultat doit être représenté par un nombre compris entre $0$ et $16$. 
 \item Calculer l'élément $11 ^{329}$ dans $\Z/13\Z$. Le résultat doit être représenté par un nombre compris entre $0$ et $12$. 
\end{enumerate}
\finenonce{007325}
\finexercice


\exercice{7326, mourougane, 2021/08/10}
\enonce{007326}{}

On note $\mathbb{F}_{11}$ le corps fini $\Z/11\Z$.
Considérons le groupe
 $\displaystyle\left(\mathbb{F}_{11} \right)^\times$.
Quels sont les ordres possibles d'un élément de ce groupe ?
Combien d'éléments de chaque ordre ce groupe possède-t-il ?
Déterminer tous les générateurs de ce groupe.

\finenonce{007326}
\finexercice


\exercice{7327, mourougane, 2021/08/10}
\enonce{007327}{}
\begin{enumerate}
 \item Déterminer la liste des inversibles de $\Z/15\Z$.
Déterminer l'ordre de chaque élément dans $\left(\Z/15\Z\right)^\times $.
 \item Déterminer la liste des inversibles de $\Z/3\Z\times \Z/5\Z$.
Déterminer l'ordre de chaque élément dans $\left(\Z/3\Z\times \Z/5\Z\right)^\times $.
 \item Donner un isomorphisme de groupes de $\left(\Z/15\Z\right)^\times $ dans $\left(\Z/3\Z\times \Z/5\Z\right)^\times$.
\end{enumerate}
\finenonce{007327}
\finexercice


\exercice{7328, mourougane, 2021/08/10}
\enonce{007328}{}
\begin{enumerate}
 \item Le polynôme $(X^2+X+1)^3$ est-il irréductible dans $\mathbb{F}_2[X]$ ?
 \item Donner la liste des polynômes irréductibles de $\mathbb{F}_2[X]$ de degré $2$ et $3$.
 \item Donner un polynôme irréductible de degré $4$ de $\mathbb{F}_2[X]$.
 \item Ecrire dans $\mathbb{F}_2[X]$, une relation de Bezout pour $X^3+X^2+1$ et $X^2+X+1$.
\end{enumerate}
\finenonce{007328}
\finexercice


\exercice{7329, mourougane, 2021/08/10}
\enonce{007329}{}
\begin{enumerate}
  \item Combien le corps $\mathbb{F}_2[X]/(X^3+X+1)$ a-t-il d'éléments ?
 \item Déterminer la liste des éléments et la table de multiplication de l'anneau quotient
$$\mathbb{F}_2[X]/(X^2+X+1).$$
 \item Multiplier $[X^5+X^4+6X]$ par $[X^4+7X^5+9X^3+4X^2]$ dans $\mathbb{F}_2[X]/(X^2+X+1)$ et donner le résultat avec un représentant de la liste précédente.
 \item Déterminer un inverse de $[X^3+X^2+1]$ dans $\mathbb{F}_2[X]/(X^2+X+1).$
\end{enumerate}
\finenonce{007329}
\finexercice


\exercice{7330, mourougane, 2021/08/10}
\enonce{007330}{}

 Résoudre dans $\R[X]$ le système de congruences
$$\left\{\begin{array}{cccc}
     P&=&X&[X^2+X+1]\cr
P&=&3&[X^2+X]
    \end{array}
\right.$$
\finenonce{007330}
\finexercice


\exercice{7331, mourougane, 2021/08/10}
\enonce{007331}{}
\begin{enumerate}
 \item Le polynôme $X^2+1$ est-il irréductible dans $\mathbb{F}_3[X]$ ?
 \item Quelle est alors la structure de l'ensemble quotient $A=\mathbb{F}_3[X]/(X^2+1)$ ? 
 \item Quelle relation vérifie la classe $\alpha$ du polynôme $X$ dans ce quotient ?
 \item Donner la liste des éléments de $A$.
 \item Déterminer l'ordre multiplicatif de $\alpha$ dans $A^\times$.
 \item Déterminer l'ordre multiplicatif de $a:=\alpha+2$ dans $A^\times$.
 \item Etablir la table des puissances de $a$.
 \item Calculer $(2+a) (2+2a)$.
 \item Calculer $a^3+a^2$ comme puissance de $a$.
 \item Calculer $(1+2a)^{-1}$.
\end{enumerate}
\finenonce{007331}
\finexercice


\exercice{7332, mourougane, 2021/08/10}
\enonce{007332}{}
\begin{enumerate}
 \item Montrer que le polynôme $X^3+X+1$ est irréductible dans $\mathbb{F}_2[X]$.
 \item Quelle est alors la structure de l'ensemble quotient $A=\mathbb{F}_2[X]/(X^3+X+1)$ ? 
 \item Quel est le cardinal de $A$ ? Soit $\alpha$ la classe du polynôme $X$ dans ce quotient $A$. Donner la liste des éléments de $A$.
 \item Sans calculs, mais en justifiant votre réponse, dire ce que valent les quantités suivantes 
$$\alpha+\alpha, \ \ \alpha^3+\alpha+1, \ \ \ \alpha^7.$$
 \item Déterminer l'ordre multiplicatif de $\alpha$ dans $A^\times$.
 \item Etablir la table des puissances de $\alpha$.
 \item Calculer $(1+\alpha) (1+\alpha^2)$.
 \item Calculer $\alpha^3+\alpha^2$ comme puissance de $\alpha$.
 \item Calculer $(1+\alpha)^{-1}$.
\end{enumerate}
\finenonce{007332}
\finexercice


\exercice{7348, mourougane, 2021/08/10}
\enonce{007348}{}
Alice et Bernard décident d'utiliser l'algorithme d'El Gamal.
Il utilise le corps $\mathbb{F}_{19}$ avec l'élément $G=15$.
\begin{enumerate}
    \item Déterminer l'ordre de $15$ dans $\mathbb{F}_{19}^\times$.
    \item Bernard choisit sa clé privée $c=4$. Déterminer sa clé publique $C=G^c$.
    \item Alice choisit une clé temporaire privée $d=5$. Quelle est sa clé publique $D$ ?
    Elle souhaite envoyer le message $m=17$. Elle le chiffre en utilisant la clé publique $C$ de Bernard par $(M_1,M_2)=(D,mC^d)$. Calculer ce message chiffré.
    \item Comment Bernard retrouve-t-il le message $m$ ?
    \item Dans un second envoi, Bernard reçoit $(8,3)$.
    Quel est le message $m$ envoyé cette fois par Alice ?
    Quelle clé privée a-t-elle utilisé cette fois ?
\end{enumerate}
\finenonce{007348}
\finexercice

\section{203.07 Groupe de permutation}


\exercice{1402, hilion, 2003/10/01}

\enonce{001402}{}
\begin{enumerate}
\item Déterminer $\text{card}(S_{3})$ et écrire tous les éléments de $S_{3}$, puis écrire la table de $S_{3}$ et en déduire tous les sous-groupes de $S_{3}$.
\item On considère $T$ un triangle équilatéral du plan, de sommets $A,B,C$. 
\begin{enumerate}
\item Montrer que les isométries du plan qui préservent $T$ forment un groupe pour la loi $\circ$, que l'on note $G$.
\item Montrer qu'un élément de $G$ induit une permutation de l'ensemble $\{A,B,C\}$. On construit ainsi une application $\phi$ de $G$ dans $S_{3}$.
\item Montrer que $\phi$ est un isomorphisme.
\end{enumerate}
\end{enumerate}
\finenonce{001402}



\finexercice



\exercice{1403, hilion, 2003/10/01}

\enonce{001403}{}
On considère le groupe symétrique $S_n$.
\begin{enumerate}
\item Déterminer $\text{card}(S_{n})$.
\item Calculer $(3 4) (4 5) (2 3) (1 2) (5 6) (2 3) (4 5) (3 4) (2 3)$.
\item Rappel: la permutation
$\sigma=
\begin{pmatrix}
a_{1} & a_{2} & \dots & a_{k} \\
a_{2} & a_{3} & \dots & a_{1}
\end{pmatrix}$
est un cycle de longueur $k$, que l'on note $(a_{1}\; a_{2} \dots a_{k})$.

Si $\tau\in S_n$, montrer que $\tau\sigma\tau^{-1}=(\tau(a_{1})\; \tau(a_{2}) \dots \tau(a_{k}))$.
\item Rappel: toute permutation se décompose en produit de cycles à supports disjoints, et cette décomposition est unique à l'ordre près.

Décomposer les permutations suivantes en produits de cycles à supports disjoints:
$\begin{pmatrix}
1&2&3&4&5 \\
3&4&5&1&2
\end{pmatrix}$,
$\begin{pmatrix}
1&2&3&4&5&6&7 \\
7&6&1&2&3&4&5
\end{pmatrix}$,
$\begin{pmatrix}
1&2&3&4&5&6&7&8 \\
6&2&5&7&8&1&3&4
\end{pmatrix}$
\item Rappel: il existe un unique morphisme de $S_n$ dans $(\{-1,1\},\times)$ non trivial, appelé signature, et noté $\epsilon$.
Une manière de calculer $\epsilon(\tau)$ (o\`u $\tau\in S_n$) consiste à décomposer $\tau$ en produit de $p$ transpositions (ie cycles de longueur 2): alors $\epsilon(\tau)=(-1)^p$.

Montrer que la signature d'un cycle de longueur $k$ vaut $(-1)^{k-1}$. En déduire comment se calcule la signature d'une permutation à partir de sa décomposition en produit de cycles disjoints.
\end{enumerate}
\finenonce{001403}



\finexercice



\exercice{1404, barraud, 2003/09/01}

\enonce{001404}{}
  Comment passer de 1234 à 2314 en échangeant seulement deux chiffres à
  chaque étape~? Y a-t-il plusieurs façons d'y parvenir~? Même question
  pour 1234 et 4312.

  Peut-on obtenir n'importe quelle permutation des chiffres 1234 par ce
  procédé~?
\finenonce{001404}



\finexercice



\exercice{1405, barraud, 2003/09/01}

\enonce{001405}{}
Représenter graphiquement les permutations suivantes. Les décomposer
en produit de cycles à supports disjoints, puis en produits de
transpositions.
$$
\sigma_{1}=
 \left(
\begin{array}{@{}c@{}} 1234567\\1425376 \end{array}
 \right)
\qquad
\sigma_{2}=
 \left(
\begin{array}{@{}c@{}} 1234567\\2471635 \end{array}
 \right)
\qquad
\sigma_{3}=
 \left(
\begin{array}{@{}c@{}} 1234567\\3261547 \end{array}
 \right)
\qquad
\sigma_{4}=
 \left(
\begin{array}{@{}c@{}} 1234567\\7146253 \end{array}
 \right)
$$

  Calculer la signature des permutations ci-dessus. Calculer le produit
  $\sigma_{1}\sigma_{2}\sigma_{3}$ et sa signature. Comparer ce 
  résultat aux précédents.
\finenonce{001405}



\finexercice



\exercice{1406, barraud, 2003/09/01}

\enonce{001406}{}
  Soient $a,b,c$ trois éléments distincts de $\{1,...,n\}$. Calculer le
  produit $(ab)(bc)(ab)$.

  En déduire que $\mathcal{S}_{n}$ est engendré par les permutations
  $\{(1,i)\}_{2\leq, i\leq n}$, c'est à dire que toute permutation
  s'écrit comme produit de transpositions de cette forme.

  Montrer que $\mathcal{S}_{n}$ est engendré par $(12)$ et $(123...n)$.
\finenonce{001406}



\finexercice



\exercice{1407, barraud, 2003/09/01}

\enonce{001407}{}
  Décrire tous les morphismes de groupe de
  $(\mathcal{S}_{n},\circ)\rightarrow (\{+1,-1\},\cdot)$, c'est les
  applications $\phi~:\mathcal{S}_{n}\rightarrow\{+1,-1\}$ satisfaisant~:
  $$
  \forall (\sigma,\sigma')\in\mathcal{S}_{n}^{2},\quad
  \phi(\sigma\sigma')=\phi(\sigma)\phi(\sigma')
  $$
  \textbf{Indication}~: Commencer par montrer que toutes les
  transpositions ont même image.
\finenonce{001407}



\finexercice



\exercice{1408, barraud, 2003/09/01}

\enonce{001408}{}
Dans $\R^{n}$, on désigne par $(e_{1},...,e_{n})$ la base canonique.
A une permutation $\sigma\in\mathcal{S}_n$, on associe l'endomorphisme $u_{\sigma}$ de
$\R^{n}$ suivant~:
$$
u_{\sigma}~:
\begin{array}{rcl}
\R^{n} & \rightarrow  & \R^{n} \\
\left(\begin{smallmatrix}
x_{1}\\\vdots\\x_{n}
\end{smallmatrix}\right)
&\mapsto &
\left(\begin{smallmatrix}
x_{\sigma(1)}\\\vdots\\x_{\sigma(n)}
\end{smallmatrix}\right)
\end{array}
$$
\begin{enumerate}
\item
Soit $\tau=(ij)$ une transposition. \'Ecrire la matrice de $u_{\tau}$
dans la base canonique. Montrer que $\det(u_{\tau})=-1$.

\item
Montrer que $\forall\sigma,\sigma'\in\mathcal{S}_n,\ u_{\sigma}\circ
u_{\sigma'}=u_{\sigma\circ\sigma'}$.

\item
En déduire que $\forall\sigma\in\mathcal{S}_n,\ \det u_{\sigma}=\epsilon(\sigma)$ où
$\epsilon$ désigne la signature.
\end{enumerate}
\finenonce{001408}



\finexercice



\exercice{1409, barraud, 2003/09/01}

\enonce{001409}{}
On note $\mathcal{S}_n$ le groupe sym\'{e}trique des permutations sur $n$ \'{e}l\'{e}ments.

Soit $\rho$ un \emph{morphisme de groupes} de $(\mathcal{S}_n,\circ)$ dans $(\{-1,1\},\cdot)$,
c'est \`{a} dire une application de $\mathcal{S}_n$ dans $\{-1,1\}$ satisfaisant
$$
  \forall(\sigma,\tau)\in\mathcal{S}_n\;\; \rho(\sigma\tau)=\rho(\sigma)\rho(\tau)
$$
\begin{enumerate}
\item
Calculer $\rho(\mathrm{id})$. Pour tout cycle $\gamma$ de longueur $p$, calculer $\gamma^{p}$. En
d\'{e}duire que lorsque $p$ est impair, $\rho(\gamma)=1$.
\item
On suppose que pour toute transposition $\tau$, $\rho(\tau)=1$. Montrer que
$\forall\sigma\in\mathcal{S}_n,\;\rho(\sigma)=1$
\item
On suppose maintenant qu'il existe une transposition $\tau_{0}=(a,b)$ pour laquelle
$\rho(\tau_{0})=-1$.


\begin{enumerate}
\item[(a)]
Pour un \'{e}l\'{e}ment $c\in\{1,\ldots,n\}\setminus\{a,b\}$, calculer $(a,b)(a,c)$. En d\'{e}duire
que $\rho(a,c)=-1$.
\item[(b)]
Pour deux \'{e}l\'{e}ments distincts $c$ et $d$ de $\{1,\ldots,n\}$, calculer $(a,c)(a,d)(a,c)$.
En d\'{e}duire que $\rho(c,d)=-1$.
\item[(c)] En déduire que pour toute transposition $\tau$, $\rho(\tau)=-1$ puis 
montrer que pour toute permutation $\sigma\in\mathcal{S}_n$, $\rho(\sigma)$ est la signature de
$\sigma$.
\end{enumerate}
\item Quels sont tous les morphismes de groupes de $(\mathcal{S}_n,\circ)$ dans $(\{-1,1\},\cdot)$ ?
\item

On consid\`{e}re l'application $\varphi$ suivante :
$$
 \varphi :
  \begin{array}{ccc}
    \mathcal{S}_n & \rightarrow  & \{-1,1\} \\[3mm]
    \sigma  & \mapsto & \prod_{i=1}^{n}\frac{\sigma(i)-\sigma(j)}{i-j}
  \end{array}
$$
Montrer que $\forall(\sigma,\tau)\in\mathcal{S}_n,\;
\varphi(\sigma\tau)=\varphi(\sigma)\varphi(\tau)$.

En d\'{e}duire que
$$
 \forall\sigma\in\mathcal{S}_n,\;\;\epsilon(\sigma)=\prod_{i=1}^{n}\frac{\sigma(i)-\sigma(j)}{i-j},
$$
o\`{u} $\epsilon(\sigma)$ d\'{e}signe la signature de $\sigma$.
\end{enumerate}



\finenonce{001409}



\finexercice



\exercice{1410, legall, 1998/09/01}

\enonce{001410}{}
Soit $  G  $ un groupe d'ordre $  2n  $ et $  H  $ un sous-groupe de $  G  $ d'ordre $  n
$ ($  H  $ est donc d'indice deux dans $  G  $).
\begin{enumerate}
    \item Montrer que si $  g \in G  $ et $  g\not \in H  ,$ on a $  H\cap gH=\emptyset   $ puis que $  G=H\cup gH  .$
    \item En d\'eduire que pour tout $  g \in G  ,  g^2\in H  .$
    \item On suppose d\'esormais $  G=\mathcal{A} _4  $ le groupe des permutations paires de l'ensemble
$ \{1,2,3,4 \}   .$ Soit $  \sigma =(a,b,c)  $ un $  3$-cycle. Montrer que $  \sigma  $ peut
s'\'ecrire comme le carr\'e d'une permutation paire c'est \`a dire qu'il existe $  \varphi \in \mathcal{A}_4
$ telle que $  \varphi ^2=\sigma   .$ En d\'eduire que $  \mathcal{A} _4  $ ne poss\`ede pas de sous-groupe d'ordre $  6 .$
\end{enumerate}
\finenonce{001410}



\finexercice



\exercice{1411, legall, 1998/09/01}

\enonce{001411}{}
D\' eterminer tous les \' el\' ements  $\sigma \in S_n$  tels
que  $\sigma ^2=\sigma $.
\finenonce{001411}



\finexercice



\exercice{1412, legall, 1998/09/01}

\enonce{001412}{}
\begin{enumerate}
    \item Rappeler $  \vert S_3\vert   .$ Montrer que $  S_3  $ ne contient pas d'\'el\'ement d'ordre $  6  .$
    \item Montrer que $  S_3  $ contient un unique sous-groupe d'ordre $  3  .$ D\'eterminer tous
les sous-groupes d'ordre $  2  $ de $  S_3  .$
    \item D\'eduire de ce qui pr\'ec\`ede tous
les sous-groupes de $  S_3  .$
\end{enumerate}

\finenonce{001412}



\finexercice



\exercice{1413, bodin, 1999/11/01}

\enonce{001413}{examen juin 1999}
Soit $GL_2(\Rr)$ l'ensemble des matrices inversibles $2\times 2$
\`a c\oe fficients r\'eels. $GL_2(\Rr)$ est naturellement muni d'une structure
de groupe par la multiplication usuelle des matrices.
Soit
$$A=\begin{pmatrix}
     1 & 0 \\
     0 & -1 \\
    \end{pmatrix}
\qquad \text{ et } \qquad
B=\begin{pmatrix}
     0 & -1 \\
     1 & 0 \\
    \end{pmatrix}.
$$

\begin{enumerate}
  \item Montrer que $A$ et $B$ appartiennent \`a $GL_2(\Rr)$.
  \item Quels sont les ordres de $A$ et $B$ ?
  \item Montrer que $AB = -BA$ et en d\'eduire que :
  \begin{enumerate}
    \item $G = \big\{ I, A, B, AB, -I, -A, -B, -AB \big\}$ est un groupe (pour
la loi multiplicative des matrices ; $I$ esl la matrice identit\'e) ;
    \item $G$ est le sous-groupe de $GL_2(\Rr)$ engendr\'e par $\{ A,B\}$.
  \end{enumerate}
  \item On munit $\Rr^2$ de sa structure euclidienne orient\'ee canonique.
  \begin{enumerate}
     \item Montrer que $G$ est inclus dans $O_2(\Rr)$ (le groupe orthogonal).
     \item D\'eterminer l'intersection de $G$ et de $SO_2(\Rr)$ (le groupe sp\'ecial orthogonal).
     \item D\'eterminer la nature g\'eom\'etrique des $8$ \'el\'ements de $G$.
  \end{enumerate}
\end{enumerate}
\finenonce{001413}



\finexercice



\exercice{1414, bodin, 1999/11/01}

\enonce{001414}{examen juin 1999}
\hfil I \\

Soit $(G,\cdot)$ un groupe. On d\'efinit le centre
$\mathcal{Z}(G)$ de $G$ par :
$$\mathcal{Z}(G) = \big\{ x \in G\  /\  \forall a \in G \ \ ax=xa \big\}.$$

Montrer que $\mathcal{Z}(G)$ est un sous-groupe de $G$.

Que peut-on dire de $\mathcal{Z}(G)$ si $G$ est ab\'elien ?


\bigskip

\hfil II \\

On d\'esigne par $\mathcal{A}_n$ le groupe altern\'e d'ordre $n$
(rappel : c'est le sous-groupe de
$(\mathcal{S}_n,\circ)$ form\'e des permutations
de $E_n = \{1,2,\ldots,n\}$
de signature $+1$.)

On se propose de d\'eterminer le centre de
$\mathcal{A}_n$ pour $n\geq 3$.

\begin{enumerate}
  \item Donner la liste des \'el\'ements de $\mathcal{A}_3$ et de  $\mathcal{Z}(\mathcal{A}_3)$.

  \item On suppose d\'esormais $n\geq 4$.
        Dans cette question on fixe $i,j,k$ trois \'el\'ements
distincts de $E_n$.

  \begin{enumerate}
    \item  V\'erifier que le $3$-cycle $(i,j,k)$ est dans $\mathcal{A}_n$.
    \item Soit $s\in \mathcal{S}_n$, montrer que
$s\circ(i,j,k) = (s(i),s(j),s(k)) \circ s$.
    \item En d\'eduire que si $s \in \mathcal{Z}(\mathcal{A}_n)$
alors l'image de  $\{i,j,k\}$ par $s$ est $\{i,j,k\}$.
   \end{enumerate}

   \item Pour $n=4$, on note $E_4= \{ i,j,k,\ell\}$.
Si $s \in \mathcal{Z}(\mathcal{A}_4)$ montrer que $s(\ell) = \ell$.
 En d\'eduire $\mathcal{Z}(\mathcal{A}_4)= \{ \text{id} \}$.

   \item Pour $n \geq 5$, soit $s \in \mathcal{Z}(\mathcal{A}_n)$,
soit $i,j,k,\ell,m$ cinq \'el\'ements distincts de $E_n$.
En consid\'erant les ensembles  $\{i,j,k\}$ et  $\{i,\ell,m\}$ montrer
que  $s = \text{id}$ et d\'eterminer
$\mathcal{Z}(\mathcal{A}_n)$

%  \item Soit $H$ un sous-groupe de $G$. A-t-on $\mathcal{Z}(H) \subset
% \mathcal{Z}(G)$ ?
\end{enumerate}
\finenonce{001414}



\finexercice



\exercice{1415, legall, 1998/09/01}

\enonce{001415}{}
Quel est l'ordre maximal d'un \'el\'ement de $  S_4  ?$ De $  S_5  ?$
De $  A_5  ?$
\finenonce{001415}



\finexercice



\exercice{1416, legall, 1998/09/01}

\enonce{001416}{}
On d\'esigne par  $  K  $ le sous-ensemble $  \{ id, (1,2)(3,4),
(1,3)(2,4),(1,4)(2,3)\}   $ de $  S_4  .$
\begin{enumerate}
    \item Montrer que $  K  $ est un sous-groupe distingu\'e de $  S_4  $ et de $  A_4  .$
    \item Pour quelle raison $  K  $ est-il isomorphe \`a $  \Z /2\Z \times \Z /2\Z   ?$ Calculer le quotient $  A_4/K  .$
    \item Montrer que le quotient $  S_4/K  $
est isomorphe \`a $  S_3  .$
    \item Donner un exemple de sous groupe distingu\'e de $  K  $ et non de $  S_4  .$ Quelle conclusion peut-on en tirer~?
\end{enumerate}
\finenonce{001416}



\finexercice



\exercice{1417, legall, 1998/09/01}

\enonce{001417}{}
Calculer $  Z(S_n)  $ suivant les valeurs de $  n\in \N   .$
\finenonce{001417}



\finexercice



\exercice{1418, ortiz, 1999/04/01}

\enonce{001418}{}
 Trouver la d\'ecomposition en produit de cycles \`a
supports disjoints, la signature, l'ordre et une d\'ecomposition
en produit de transpositions
des permutations suivantes de ${\cal S}_{10}:$%
$${
\sigma =\left(
\begin{array}{cccccccccc}
1 & 2 & 3 & 4 & 5 & 6 & 7 & 8 & 9 & 10 \\
3 & 7 & 1 & 4 & 2 & 6 & 9 & 8 & 5 & 10
\end{array}
\right) ,\
}$$
$${
\varphi =\left( 10,3,4,1\right) \left( 8,7\right) \left( 4,7\right) \left(
5,6\right) \left( 2,6\right) \left( 2,9\right)
.\bigskip }$$ Calculer $\sigma ^{1998}\;$et
$\varphi ^{1998}.$


\finenonce{001418}



\finexercice



\exercice{1419, ortiz, 1999/04/01}

\enonce{001419}{}
${\cal A}_4$ d\'esigne le groupe des permutations
paires sur l'ensemble $E=\left\{ 1,2,3,4\right\}
.$
\begin{enumerate}
    \item  Quels sont les ordres des \'el\'ements de ${\cal A}_4$ ? En d\'eduire
la liste de ces \'el\'ements sous forme d\'ecompos\'ee en produit de cycles
\`a supports disjoints.
    \item  Montrer que les permutations $s=(1\;2)(3\;4)$ et $r=(1\;2\;3)$
engendrent ${\cal A}_4.$
    \item  Montrer que ${\cal A}_4$ admet un unique sous-groupe $H$ d'ordre 4 (on
examinera d'abord les ordres des \'el\'ements d'un
tel sous-groupe) et que ce sous-groupe est un
sous-groupe distingu\'e de ${\cal A}_4$.
\end{enumerate}
\finenonce{001419}



\finexercice



\exercice{1420, ortiz, 1999/04/01}

\enonce{001420}{}
Le groupe $G=\mathcal{S}_3\times \mathcal{S}_3$
est-il ab\'elien ? D\'eterminer tous les
sous-groupes de $G$ d'ordre 4.
\finenonce{001420}



\finexercice



\exercice{1421, ortiz, 1999/04/01}

\enonce{001421}{}
Quel est le nombre de $k$-cycles dans
$\mathcal{S}_k$ puis dans $\mathcal{S}_n$ o\`u
$k\leq n$ ?
\finenonce{001421}



\finexercice



\exercice{1422, ortiz, 1999/04/01}

\enonce{001422}{}
Soit $G$ un sous-groupe de $\mathcal{S}_n.$
\begin{enumerate}
    \item  Montrer que si $G$ est d'ordre impair alors $G$ ne contient aucune permutation impaire.
    \item  Montrer que si $G$ contient au moins une permutation impaire,
alors
$G$ contient autant de permutations paires que de
permutations impaires.
\end{enumerate}
\finenonce{001422}



\finexercice



\exercice{1423, ortiz, 1999/04/01}

\enonce{001423}{}
Soient $a=(1,2)(3,4),b=(1,3)(2,4),c=(1,4)(2,3)\in
\mathcal{A}_4,$ $X=\left\{a,b,c\right\},$
$V=\left\{a,b,c,\text{Id}\right\}$ et
$\Phi:\mathcal{S}_4\to\mathcal{S}(X),g\in G\mapsto
\Phi_g=\left[ x\mapsto gxg^{-1}\right].$
\begin{enumerate}
    \item
    \begin{enumerate}
        \item Montrer que $V$ est un sous-groupe distingu\'e de $\mathcal{A}_4$
(on pourra \'etudier l'ordre des \'elements de
$\mathcal{A}_4$).
        \item Montrer que $<a>$ est un sous-groupe distingu\'e de $V$ et n'est
pas un sous-groupe distingu\'e de $\mathcal{A}_4.$
    \end{enumerate}
    \item  Montrer que $\Phi$ est un homomorphisme de groupes.
    \item
    \begin{enumerate}
         \item Calculer $\Phi(g)$ pour $g=(1,2)$ puis $g=(1,2,3).$
         \item En d\'eduire que $\Phi$ est surjectif.
    \end{enumerate}
    \item  Montrer que $\mathcal{S}_4/V$ est isomorphe \`a $\mathcal{S}_3.$
    \item  Ecrire la d\'ecomposition de $\mathcal{A}_4$ suivant les classes modulo $V.$
\end{enumerate}
\finenonce{001423}



\finexercice



\exercice{1424, ortiz, 1999/04/01}

\enonce{001424}{}
\begin{enumerate}
    \item  D\'eterminer le centre du groupe $\mathcal{S}_n.$
    \item
    \begin{enumerate}
    \item Montrer qu'un groupe $G_1\times G_2$ contient un sous-groupe distingu\'e isomorphe \`a $G_1.$
    \item Montrer que les groupes $\mathcal{S}_n$ et $\Zz/2\Zz\times\mathcal{A}_n$
ne sont pas isomorphes si $n\geq 3$.
    \end{enumerate}
\end{enumerate}
\finenonce{001424}



\finexercice



\exercice{1425, ortiz, 1999/04/01}

\enonce{001425}{}
\begin{enumerate}
    \item  Montrer que dans $\mathcal{S}_n$ on a
$f\circ (a,b)\circ f^{-1}=(f(a),f(b))$.
    \item  Montrer que les permutations $(1,...,n)$ et $(1,2)$ engendrent
 $\mathcal{S}_n$ (on rappelle que les transpositions engendrent $\mathcal{S}_n$).
\end{enumerate}
\finenonce{001425}



\finexercice



\exercice{1426, ortiz, 1999/04/01}

\enonce{001426}{}
\begin{enumerate}
    \item  Montrer que $\mathcal{S}_n$ est isomorphe \`a un sous-groupe de $\mathcal{A}_{n+2}.$
    \item  Montrer que $\mathcal{S}_4$ n'est pas isomorphe \`a un sous-groupe de $\mathcal{A}_5.$
    \item  Montrer que $\mathcal{S}_5$ n'est pas isomorphe \`a un sous-groupe de $\mathcal{A}_6.$
\end{enumerate}

\finenonce{001426}



\finexercice



\exercice{1427, gourio, 2001/09/01}

\enonce{001427}{}
Montrer que tout groupe fini est isomorphe \`{a} un sous-groupe de $S_{n}$
(groupe sym\'{e}trique)  $ $pour un certain $n$.
\finenonce{001427}



\finexercice



\exercice{3074, quercia, 2010/03/08}

\enonce{003074}{G{\'e}n{\'e}rateurs de ${\cal S}_n$}
Soit $n \in \N^*$.
Montrer que ${\cal S}_n$ est engendr{\'e} par les sous-ensembles suivants :
\begin{enumerate}
  \item $A = \{ (i,i+1)$ tq $1 \le i < n    \}$.
  \item $B = \{ (1\ i)$  tq $2 \le i \le n  \}$.
  \item $C = \{ (1\ 2), (1\ 2\ \cdots\ n)  \}$.
\end{enumerate}
\finenonce{003074}



\finexercice 


\exercice{3075, quercia, 2010/03/08}

\enonce{003075}{G{\'e}n{\'e}rateurs de ${\cal S}_n$}
Montrer que toute permutation de ${\cal S}_n$ s'{\'e}crit de mani{\`e}re unique :
$\sigma = c_2^{\alpha_2}\circ c_3^{\alpha_3} \circ \dots \circ c_n^{\alpha_n}$
o{\`u} $c_i = (1\ 2\ \cdots\ i)$ et $0 \le \alpha_i < i$.
\finenonce{003075}



\finexercice 


\exercice{3076, quercia, 2010/03/08}

\enonce{003076}{${\cal A}_n$ est engendr{\'e} par les 3-cycles}
\begin{enumerate}
  \item Calculer $(a\,b\,c) \circ (b\,c\,d)$.
  \item Montrer que le sous-groupe altern{\'e} ${\cal A}_n$ est engendr{\'e}
    par les 3-cycles ($n \ge 3$).
\end{enumerate}
\finenonce{003076}



\finexercice 


\exercice{3077, quercia, 2010/03/08}

\enonce{003077}{${\cal A}_n$ est engendr{\'e} par les 3-cycles}
Soit $n \in \N,n \ge 4$.
\begin{enumerate}
  \item Soit $i,j \in \{3,\dots,n \}\ , i \ne j $.\par
    D{\'e}composer en cycles {\`a} supports disjoints la permutation :
    $\sigma=(1\ i\ 2)\circ(1\ 2\ j)\circ(1\ i\ 2)$.

  \item On note $\cal H$ le sous-groupe de ${\cal A}_n$ engendr{\'e} par les 3-cycles
    $(1\ 2\ k)$,\quad $3 \le k \le n$.
  \begin{enumerate}
    \item Montrer que : $\forall\ i,j \ge 3$, avec $i \ne j$, $\cal H$ contient $(1\ 2)\circ(i\ j)$ et $(i\ j)\circ(1\ 2)$.
    \item Montrer que : $\forall\ j  \ge 3$, $\cal H$ contient $(1\ 2)\circ(1\ j)$ et $(1\ 2)\circ(2\ j)$.
    \item Montrer que : $\forall\ i \ne j$, $\forall\ k \ne l$,\quad $(i\ j)\circ(k\ l) \in \cal H$.
    \item Montrer que ${\cal H}={\cal A}_n$.
   \end{enumerate}   
\end{enumerate}
\finenonce{003077}



\finexercice 


\exercice{3078, quercia, 2010/03/08}

\enonce{003078}{Signature en fonction du nombre d'orbites}
Soit $\sigma \in {\cal S}_n$. On note $c$ le nombre de cycles {\`a} supports disjoints
constituant $\sigma$, et $f$ le nombre de points fixes.

Calculer $\epsilon(\sigma)$ en fonction de $n$, $c$, et $f$.
\finenonce{003078}


\finexercice 


\exercice{3079, quercia, 2010/03/08}

\enonce{003079}{Nombre de transposition pour engendrer un cycle}
Soit $\sigma \in S_n$. On appelle {\it orbite de $\sigma$\/} toute partie
$X$ de $\{1,\dots,n\}$ sur laquelle $\sigma$ induit une permutation circulaire.
(Les orbites sont les supports des cycles de $\sigma$, et les singletons
constitu{\'e}s de points fixes)

On note $N(\sigma)$ le nombre d'orbites de $\sigma$.
\begin{enumerate}
  \item Montrer que si $\tau$ est une transposition, alors $N(\tau\circ\sigma) = N(\sigma) \pm 1$.
  \item Application : Quel est le nombre minimal de transpositions n{\'e}c{\'e}ssaires
    pour obtenir un $n$-cycle ?
\end{enumerate}
\finenonce{003079}



\finexercice 


\exercice{3080, quercia, 2010/03/08}

\enonce{003080}{Conjugaison}
Soient $\sigma, \sigma' \in {\cal S}_n$. On dit que $\sigma$ et  $\sigma'$ sont
conjugu{\'e}es s'il existe $\rho \in {\cal S}_n$ tel que
$\sigma' = \rho \circ \sigma \circ \rho^{-1}$.
\begin{enumerate}
  \item Montrer que tout conjugu{\'e} d'un $k$-cycle est encore un $k$-cycle.
  \item Montrer que $\sigma$ et  $\sigma'$ sont conjugu{\'e}es si et seulement si
    les cycles {\`a} supports disjoints de $\sigma$ et $\sigma'$ ont deux {\`a} deux
    m{\^e}mes longueurs.
\end{enumerate}
\finenonce{003080}



\finexercice 


\exercice{3081, quercia, 2010/03/08}

\enonce{003081}{Caract{\'e}risation de la signature}
Soit $E$ un ensemble fini et $f : {{\cal S}_E} \to {\C^*}$ un morphisme de groupes.
\begin{enumerate}
  \item Si $\sigma$ est une transposition, que peut-on dire de $f(\sigma)$ ?
  \item Montrer que deux permutations conjugu{\'e}es ont m{\^e}me image par $f$.
  \item En d{\'e}duire que $f$ est la fonction constante $1$, ou bien $f$ est la
    signature.
\end{enumerate}
\finenonce{003081}



\finexercice 


\exercice{3082, quercia, 2010/03/08}

\enonce{003082}{Calcul de signature}
Soit $$\sigma = \begin{pmatrix}1 &2 &3 &\dots &n    &n+1 &n+2 &\dots &2n \cr
                        1 &3 &5 &\dots &2n-1 &2   &4   &\dots &2n \end{pmatrix}.$$
Calculer $\varepsilon(\sigma)$.
\finenonce{003082}


\finexercice 


\exercice{3083, quercia, 2010/03/08}

\enonce{003083}{Centre de ${\cal S}_E$}
Soit $E$ un ensemble ayant au moins trois {\'e}l{\'e}ments.
\begin{enumerate}
  \item Pour $a,b \in E$ distincts et $\sigma \in {\cal S}_E$, simplifier
    $\sigma \circ (a\ b) \circ \sigma^{-1}$.
  \item Quelles sont les permutations $\sigma$ qui commutent avec $(a\ b)$ ?
  \item En d{\'e}duire que le centre de ${\cal S}_E$ est r{\'e}duit {\`a} $\{$id$_E\}$.
\end{enumerate}
\finenonce{003083}



\finexercice 


\exercice{3084, quercia, 2010/03/08}

\enonce{003084}{Commutant d'un $n$-cycle}
Soit $\sigma = (1\ 2\ \dots \ n) \in {\cal S}_n$. Trouver toutes les permutations
$\rho \in {\cal S}_n$ commutant avec $\sigma$.
(Reconna{\^\i}tre $\rho \circ \sigma \circ \rho^{-1}$)

\finenonce{003084}



\finexercice 


\exercice{3085, quercia, 2010/03/08}

\enonce{003085}{Commutant d'un produit de 5-cycles}
Dans ${\cal S}_{10}$, quelles sont les permutations qui commutent avec
$\sigma = (1\ 2\ 3\ 4\ 5)\circ(6\ 7\ 8\ 9\ 10)$ ?

\finenonce{003085}


\finexercice 


\exercice{3086, quercia, 2010/03/08}

\enonce{003086}{Puissances d'un $k$-cycle}
Soit $\sigma$ un $k$-cycle de ${\cal S}_n$ et $p \in \Z$.
\begin{enumerate}
  \item Si $p \mid k$, montrer que $\sigma^p$ est le produit de $p$ cycles {\`a} supports
    disjoints de longueur $\frac kp$.
  \item Montrer que pour $p \wedge k = 1$, $\sigma^p$ est un $k$-cycle
    (utiliser l'{\'e}galit{\'e} de B{\'e}zout).
  \item Dans le cas g{\'e}n{\'e}ral, {\'e}tudier la d{\'e}composition en cycles de $\sigma^p$.
\end{enumerate}
\finenonce{003086}



\finexercice 


\exercice{3087, quercia, 2010/03/08}

\enonce{003087}{Ordre maximal}
Trouver l'ordre maximal d'une permutation de ${\cal S}_{10}$.
\finenonce{003087}


\finexercice 


\exercice{3088, quercia, 2010/03/08}

\enonce{003088}{Sous-groupe d'indice 2 dans ${\cal S}_n$}
Soit $H$ un sous-groupe de ${\cal S}_n$ d'ordre $\frac {n!}2$.
On note $K = {\cal S}_n \setminus H$.
\begin{enumerate}
  \item Pour $\sigma \in H$, montrer que $\sigma H = H$ et $\sigma K = K$.
  \item Soit $\sigma \in {\cal S}_n$. D{\'e}terminer les ensembles $\sigma H$,
    $\sigma K$, $H\sigma$, $K\sigma$ suivant que $\sigma \in H$ ou
    $\sigma \in K$.
  \item En d{\'e}duire que si deux permutations sont conjugu{\'e}es, alors elles sont toutes deux
    dans $H$ ou toutes deux dans $K$.
  \item Montrer enfin que $H = {\cal A}_n$.
\end{enumerate}
\finenonce{003088}



\finexercice 


\exercice{3089, quercia, 2010/03/08}

\enonce{003089}{D{\'e}nombrement}
Combien y a-t-il de permutations de ${\cal S}_{26}$ comportant
trois points fixes, deux 3-cycles, un 5-cycle, et deux 6-cycles~?

\finenonce{003089}


\finexercice 


\exercice{5353, rouget, 2010/07/04}
\enonce{005353}{**IT}
Soit $\sigma$ l'élément de $S_{12}$~:~$\sigma=(3\;10\;7\;1\;2\;6\;4\;5\;12\;8\;9\;11)$.
\begin{enumerate}
\item  Combien $\sigma$ possède-t-elle d'inversions~?~Que vaut sa signature~?
\item  Décomposer $\sigma$ en produit de transpositions. Retrouvez sa signature.
\item  Déterminer les orbites de $\sigma$.
\item  Déterminer $\sigma^{2005}$.
\end{enumerate}
\finenonce{005353}


\finexercice


\exercice{5354, rouget, 2010/07/04}
\enonce{005354}{***IT}
\label{exo:suprou2ter}
Démontrer que $S_n$ est engendré par $\tau_{1,2}$, $\tau_{1,3}$,...,$\tau_{1,n}$.
\finenonce{005354}


\finexercice


\exercice{5355, rouget, 2010/07/04}
\enonce{005355}{***IT}
Démontrer que $A_n$ est engendré par les cycles de longueur $3$ (pour $n\geq3$).
\finenonce{005355}


\finexercice


\exercice{5356, rouget, 2010/07/04}
\enonce{005356}{***I}
Démontrer que $S_n$ est engendré par $\tau_{1,2}$ et le cycle $(2\;3\;...\;n\;1)$.
\finenonce{005356}


\finexercice


\exercice{5357, rouget, 2010/07/04}
\enonce{005357}{***I}
Soit $(G,\times)$ un groupe. Montrer que $(G,\times)$ est isomorphe à un sous-groupe de $(S(G),\circ)$ et que, en particulier, tout groupe fini d'ordre $n$ est isomorphe à un sous-groupe de $S_n$ (théorème de \textsc{Cayley}). (Indication~:~montrer que pour chaque $x$ de $G$, l'application $y\mapsto xy$ est une permutation de $G$.)
\finenonce{005357}


\finexercice


\exercice{5358, rouget, 2010/07/04}
\enonce{005358}{***}
Soit $\sigma$ une permutation de $\{1,...,n\}$ et $k$ le nombre d'orbites de $\sigma$. Montrer que $\varepsilon(\sigma)=(-1)^{n-k}$.
\finenonce{005358}


\finexercice


\exercice{5359, rouget, 2010/07/04}
\enonce{005359}{***I}
$\sigma$ étant une permutation de $\{1,...,n\}$ donnée, on définit la matrice notée $P_\sigma$, carrée d'ordre $n$ dont le terme ligne $i$ colonne $j$ est $\delta_{i,\sigma(j)}$ (où $\delta_{i,j}$ est le symbôle de \textsc{Kronecker}. On note $G$ l'ensemble des $P_\sigma$ où $\sigma$ décrit $S_n$.
\begin{enumerate}
\item  
\begin{enumerate}
\item $\sigma$ et $\sigma'$ étant deux éléments de $S_n$, calculer $P_\sigma\times P_{\sigma'}$.
\item En déduire que $(G,\times)$ est un sous-groupe de $(GL_n(\Rr),\times)$, isomorphe à $(S_n,\circ)$ (les matrices $P_\sigma$ sont appelées \og~matrices de permutation~\fg).
\end{enumerate}
\item  (Une utilisation des $P_\sigma$) $A$ étant une matrice carrée donnée, calculer $AP_\sigma$ et $P_\sigma A$. Que constate-t-on~?
\end{enumerate}
\finenonce{005359}


\finexercice


\exercice{5360, rouget, 2010/07/04}
\enonce{005360}{***I}
\label{exo:suprou8ter}
 $A_1$, $A_2$,...,$A_p$ sont $p$ matrices carrées d'ordre $n$, deux à deux distinctes et inversibles. On suppose que $\{A_1,...,A_p\}$ est stable pour $\times$. 
Montrer que $\{A_1,...,A_p\}$ est un sous groupe de $(GL_n(\Rr),\times)$.
\finenonce{005360}


\finexercice


\exercice{5361, rouget, 2010/07/04}
\enonce{005361}{***}
Dans $E=\Rr^n$, on considère l'hyperplan $H$ d'équation $x_1+...+x_n=0$ dans la base canonique $(e_i)_{1\leq i \leq n}$ de $E$. Pour $\sigma\in S_n$ donnée, on considère l'endomorphisme $f_\sigma$ de $E$ défini par~:~$\forall i\in E,\;f_\sigma(e_i)=e_{\sigma(i)}$.

On pose alors $p=\frac{1}{n!}\sum_{\sigma\in S_n}^{}f_\sigma$. Montrer que $p$ est une projection dont on déterminera l'image et la direction.
\finenonce{005361}


\finexercice

\section{203.99 Autre}


\exercice{2960, quercia, 2010/03/08}

\enonce{002960}{$x + y - xy$}
\begin{enumerate}
  \item Sur $E = [0,1]$, on d{\'e}finit l'op{\'e}ration : $x * y = x + y - xy$.
    V{\'e}rifier que $*$ est interne, et
    {\'e}tudier ses propri{\'e}t{\'e}s (commutativit{\'e}, associativit{\'e}, {\'e}l{\'e}ment neutre,
    {\'e}l{\'e}ments sym{\'e}trisables, {\'e}l{\'e}ments r{\'e}guliers).

  \item M{\^e}mes questions avec $E = ]-\infty, 1[$.

\end{enumerate}
\finenonce{002960}



\finexercice 


\exercice{2961, quercia, 2010/03/08}

\enonce{002961}{$x  \mapsto axa$ surjective}
Soit $*$ une op{\'e}ration associative sur $E$, et $a \in E$ tel que l'application%
$ E \to  E, x \mapsto {a*x*a}$ soit surjective.

Montrer qu'il existe un {\'e}l{\'e}ment neutre, et que $a$ est sym{\'e}trisable.
\finenonce{002961}


\finexercice 


\exercice{2962, quercia, 2010/03/08}

\enonce{002962}{Op{\'e}ration induite sur les parties}
Soit $*$ une op{\'e}ration sur $E$. Pour $A,B \subset E$, on pose
$A*B = \{a*b$ tq $a \in A,\ b \in B \}$.
\begin{enumerate}
  \item   {\'E}tudier les propri{\'e}t{\'e}s de $*$ sur ${\cal P}(E)$ en fonction de celles de
      $*$ sur $E$ (commutativit{\'e}, associativit{\'e}, {\'e}l{\'e}ment neutre, {\'e}l{\'e}ments
      sym{\'e}trisables).

  \item   Est-ce que $*$ est distributive par rapport {\`a} $\cup$ ?
\end{enumerate}
\finenonce{002962}


\finexercice 


\exercice{2963, quercia, 2010/03/08}

\enonce{002963}{Loi sur $\Z^2$}
On d{\'e}finit l'op{\'e}ration dans $\Z^2$ : $(a,b) * (a',b') = (aa', ab'+b)$.
\begin{enumerate}
  \item   {\'E}tudier les propri{\'e}t{\'e}s de cette op{\'e}ration.

  \item   Pour $z \in \Z$, on pose $f_{a,b}(z) = az+b$.

      Montrer que $\phi : {\Z^2} \to {\Z^{\Z}},  {(a,b)} \mapsto {f_{a,b}}$
      est un morphisme pour $*$ et $\circ$.

  \item   Est-ce un isomorphisme ?
\end{enumerate}
\finenonce{002963}



\finexercice 


\exercice{2964, quercia, 2010/03/08}

\enonce{002964}{Composition de relations}
Soit $E$ un ensemble, et $\cal F$ l'ensemble des relations binaires sur $E$.
Pour $R,S \in \cal F$, on d{\'e}finit la relation $R\ast S$ par~:
$$x (R\ast S) y \iff \exists\ z \in E\text{ tq }x R z\text{ et }z S y.$$

A toute fonction $f : E \to  E$, on associe la relation :
$y R_f x \iff y = f(x)$.
\begin{enumerate}
  \item Montrer que $\ast$ est associative, mais non commutative en g{\'e}n{\'e}ral.
  \item Simplifier $R_f \ast R_g$.
  \item Est-ce que $\ast$ admet un {\'e}l{\'e}ment neutre ?
\end{enumerate}
\finenonce{002964}



\finexercice 


\exercice{2986, quercia, 2010/03/08}

\enonce{002986}{Groupe sans sous-groupe non trivial}
Soit $G$ un groupe n'ayant pas de sous-groupe non trivial.
Montrer que $G$ est monog{\`e}ne, fini, et que $\mathrm{Card}\, G$ est un nombre premier.
\finenonce{002986}



\finexercice


\exercice{2987, quercia, 2010/03/08}

\enonce{002987}{Groupe di{\'e}dral}
Soit $n \in \N$, $n \ge 3$. On note $\omega = \exp\frac {2i\pi}n$ et :
$${f_k} : {\C} \to {\C},  z \mapsto {\omega^kz}   \qquad
  {g_k} : {\C} \to {\C}, z \mapsto{\omega^k\overline z}   \qquad
  (0 \le k < n)$$
\begin{enumerate}
  \item Montrer que $G = \{f_0,\dots,f_{n-1}, g_0,\dots,g_{n-1}\}$ est un groupe pour
    la composition des applications.

  \item Soit $a > 0$ et $A_k$ le point du plan d'affixe $a\omega^k$.
    Montrer que $G$ repr{\'e}sente le groupe des isom{\'e}tries du polygone $A_0\dots A_{n-1}$.

  \item $G$ est-il cyclique ?

  \item Montrer que $G$ est engendr{\'e} par les applications $f_1$ et $g_0$ et que l'on a :
    $f_1\circ g_0 = g_0\circ f_1^{-1}$.

  \item Soit $H$ un groupe quelconque engendr{\'e} par deux {\'e}l{\'e}ments $\rho$ et $\sigma$
    tels que $\begin{cases}\rho   \text{ est d'ordre } n\cr
                     \sigma \text{ est d'ordre } 2\cr
                     \rho\sigma = \sigma\rho^{-1}.\end{cases}$
    \par
    Montrer que $G$ et $H$ sont isomorphes.
\end{enumerate}
\finenonce{002987}



\finexercice 


\exercice{2988, quercia, 2010/03/08}

\enonce{002988}{Groupe d'ordre pair}
Soit $G$ un groupe fini de cardinal pair. Montrer qu'il existe un {\'e}l{\'e}ment
d'ordre 2.
\finenonce{002988}



\finexercice 


\exercice{2989, quercia, 2010/03/08}

\enonce{002989}{Groupe d'ordre impair}
Soit $G$ un groupe fini de cardinal impair.
Montrer que : $\forall\ x \in G,\ \exists!\ y \in G$ tq $x = y^2$.
\finenonce{002989}



\finexercice 


\exercice{2990, quercia, 2010/03/08}

\enonce{002990}{Groupe d'exposant 2}
Soit $G$ un groupe fini tel que : $\forall\ x \in G,\ x^2 = e$.
\begin{enumerate}
  \item Montrer que $G$ est commutatif (consid{\'e}rer $(xy)(xy)$).
  \item Soit $H$ un sous-groupe de $G$ et $x \in G \setminus H$.
    On note $K$ le sous groupe engendr{\'e} par $H \cup \{x\}$.
    \par
    Montrer que $\mathrm{Card}\, K = 2\mathrm{Card}\, H$.
  \item En d{\'e}duire que $\mathrm{Card}\, G$ est une puissance de 2.
\end{enumerate}
\finenonce{002990}



\finexercice 


\exercice{2991, quercia, 2010/03/08}

\enonce{002991}{Groupes d'ordre 6}
D{\'e}terminer tous les groupes finis de cardinal 6
(on admettra que dans un tel groupe, il existe un {\'e}l{\'e}ment $a$ d'ordre 2,
et un {\'e}l{\'e}ment $b$ d'ordre 3).
\finenonce{002991}



\finexercice 


\exercice{2992, quercia, 2010/03/08}

\enonce{002992}{Groupe d'homographies}
Soit $E = \R \setminus \{0,1\}$, et $ f : E \to E, x \mapsto {\frac 1x}$,
  \qquad $ g : E \to E, x \mapsto {1-x}$

V{\'e}rifier que $f$ et $g$ sont des bijections et d{\'e}terminer le groupe engendr{\'e} par
$f$ et $g$ pour la loi $\circ$.
\finenonce{002992}



\finexercice 


\exercice{2993, quercia, 2010/03/08}

\enonce{002993}{Groupes de similitudes}
Pour $\alpha \in \C^*$ et $\beta \in \C$, on note
${f_{\alpha,\beta}} : {\C} \to {\C}, z \mapsto {\alpha z + \beta}$
\begin{enumerate}
  \item Montrer que l'ensemble des fonctions $f_{\alpha,\beta}$ est un groupe pour la
    loi $\circ$. Est-il commutatif ?
  \item A quelle condition sur $\alpha,\beta$, $f_{\alpha,\beta}$ est-elle d'ordre
    fini ?
\end{enumerate}
\finenonce{002993}



\finexercice 


\exercice{2994, quercia, 2010/03/08}

\enonce{002994}{Th{\'e}or{\`e}me de Lagrange}
Soit $G$ un groupe fini et $H$ un sous-groupe de $G$.
On d{\'e}finit une relation sur $G$ par :
$$\forall\ x,y \in G,\ x \sim y \iff \exists\ h \in H \text{ tel que } x = hy.$$
\begin{enumerate}
  \item Montrer que $\sim$ est une relation d'{\'e}quivalence. Quelle est la classe de $e$ ?
  \item Soit $a \in G$. Montrer que $\dot a$ est {\'e}quipotent {\`a} $H$.
  \item En d{\'e}duire que $\mathrm{Card}\, H$ divise $\mathrm{Card}\, G$ ({\it Th{\'e}or{\`e}me de Lagrange\/}).
\end{enumerate}
\finenonce{002994}



\finexercice 


\exercice{2995, quercia, 2010/03/08}

\enonce{002995}{Relation d'{\'e}quivalence avec deux sous-groupes}
Soient $H,K$ deux sous-groupes d'un groupe $G$.
Pour $x,y \in G$, on pose :
$$x \sim y \iff \exists\ h \in H,\ \exists\ k \in K \text{ tq } y = hxk.$$
\begin{enumerate}
  \item Montrer que c'est une relation d'{\'e}quivalence.
  \item Pour $x \in G$, soit $G_x = \{(h,k) \in H\times K$ tq $hxk^{-1} = x\}$.
    Montrer que $G_x$ est un sous-groupe de $H\times K$.
  \item Si $H$ et $K$ sont finis, montrer que chaque classe d'{\'e}quivalence est finie
    de cardinal divisant $\mathrm{Card}\,(H)\mathrm{Card}\,(K)$.
\end{enumerate}
\finenonce{002995}



\finexercice


\exercice{2996, quercia, 2010/03/08}

\enonce{002996}{Groupe d'ordre $ab$}
Soit $G$ un groupe commutatif fini d'ordre $n = ab$ avec $a \wedge b = 1$.

On pose $A = \{x \in G$ tq $x^a = e\}$ et $B = \{x \in G$ tq $x^b = e\}$.

\begin{enumerate}
  \item Montrer que $A$ et $B$ sont des sous-groupes de $G$.
  \item Montrer que $A \cap B = \{e\}$ et $AB = G$.
\end{enumerate}
\finenonce{002996}


\finexercice 


\exercice{2997, quercia, 2010/03/08}

\enonce{002997}{Sous-groupes de type fini de $\Q$}
\begin{enumerate}
  \item Soit $H$ un sous-groupe additif de $\Q$ engendr{\'e} par un nombre fini
    d'{\'e}l{\'e}ments. Montrer que $H$ est monog{\`e}ne.
  \item Trouver un sous-groupe non trivial de $\Q$ qui n'est pas engendr{\'e} par
    une famille finie.
\end{enumerate}
\finenonce{002997}



\finexercice 


\exercice{2998, quercia, 2010/03/08}

\enonce{002998}{$(\Q,+)$ et ($\Q^{+*},\times)$ ne sont pas isomorphes}
Montrer que les groupes $(\Q, +)$ et $(\Q^{+*}, \times )$ ne sont pas isomorphes
(penser {\`a} $\sqrt 2$).
\finenonce{002998}



\finexercice 


\exercice{2999, quercia, 2010/03/08}

\enonce{002999}{Sous-groupe infini de $\C^*$}
Soit $p$ un entier naturel premier. On appelle $G$ l'ensemble des $z\in \C$
pour lesquels existe $n\in\N$ tel que $z^{p^n}=1$.
\begin{enumerate}
  \item Montrer que $G$ est un groupe multiplicatif infini o{\`u} tout {\'e}l{\'e}ment est d'ordre fini.
  \item Montrer que tout sous-groupe $H$ de $G$, distinct de $G$, est cyclique
    (on pourra consid{\'e}rer un {\'e}l{\'e}ment $z_0$ de $G\setminus H$ et montrer que l'ordre des
    {\'e}l{\'e}ments de $H$ n'exc{\`e}de pas celui de $z_0$).
\end{enumerate}
\finenonce{002999}



\finexercice 


\exercice{3000, quercia, 2010/03/08}

\enonce{003000}{Th{\'e}or{\`e}me du rang}
Soit $f: G \to {G'}$ un morphisme de groupes o{\`u} $G$ est un groupe fini.

Montrer que $\mathrm{Card}\,(\mathrm{Ker} f) \times \mathrm{Card}\,(\Im f) = \mathrm{Card}\,(G)$.
\finenonce{003000}



\finexercice


\exercice{3001, quercia, 2010/03/08}

\enonce{003001}{Centre d'un $p$-groupe}
Soit $G$ un groupe fini de cardinal $p^k$ o{\`u} $p$ est un nombre premier
et $k\in\N^*$. On note $Z$ le centre de $G$.
\begin{enumerate}
  \item En consid{\'e}rant l'action de $G$ sur lui-m{\^e}me par automorphismes int{\'e}rieurs
    montrer que $\mathrm{Card}\,(Z) \equiv 0 \bmod p$.
  \item En d{\'e}duire que tout groupe d'ordre $p^2$, $p$ premier, est commutatif
    et est isomorphe soit {\`a} $\Z/p^2\Z$ soit {\`a} $(\Z/p\Z)^2$.
\end{enumerate}
\finenonce{003001}



\finexercice 


\exercice{3002, quercia, 2010/03/08}

\enonce{003002}{Sous groupes et g{\'e}n{\'e}rateurs de $\Z^2$}
On consid{\`e}re le groupe $G = \Z^2$. Une {\it base\/} de $G$ est une
famille $(\alpha=(a,a'), \beta=(b,b'))$ engendrant $G$.
\begin{enumerate}
  \item
  \begin{enumerate}
\item Montrer que $(\alpha,\beta)$ est une base de $G$ si et seulement
    si $\det(\alpha,\beta) = \pm 1$.
\item Montrer que $\alpha = (a,a')$ appartient {\`a} une base de $G$ si et
    seulement si $a\wedge a' = 1$.
  \end{enumerate}
  \item Soit $H$ un sous-groupe non trivial de $G$.
    On note $H' = \{ux+vy \text{ tq } u\in\Z,\ v\in\Z, (x,y)\in H\}$,
    $n$ le plus petit {\'e}l{\'e}ment de~$H'$ strictement positif
    et $u\in\Z$, $v\in\Z$, $(x,y)\in H$ tels que
    $ux+vy = n$.
    \begin{enumerate}
\item Montrer que $u\wedge v = 1$ et que $x$ et $y$ sont divisibles par $n$.
\item On pose $\alpha = (x/n,y/n)$ et $\beta = (-v,u)$. Montrer que
    $(\alpha,\beta)$ est une base de $G$ et qu'il existe $p\in \N$
    tel que $(n\alpha,np\beta)$ engendre $H$.
    \end{enumerate}
\end{enumerate}
\finenonce{003002}



\finexercice 


\exercice{3003, quercia, 2010/03/08}

\enonce{003003}{Partie g{\'e}n{\'e}ratrice d'un groupe fini}
Soit $G$ un groupe fini de cardinal~$n$. Montrer qu'il existe une
partie g{\'e}n{\'e}ratrice de~$G$ de cardinal inf{\'e}rieur ou {\'e}gal {\`a}~$\log_2(n)$.
\finenonce{003003}



\finexercice 


\exercice{3004, quercia, 2010/03/08}

\enonce{003004}{Groupe fini ?}
Soit $G$ un groupe ayant un nombre fini de sous-groupes. Montrer que $G$ est fini.
\finenonce{003004}



\finexercice 


\exercice{7349, mourougane, 2021/08/10}
\enonce{007349}{Questions de cours}
\begin{enumerate}
 \item Soit $(E,\star)$ un ensemble muni d'une loi de composition interne
associative avec un élément neutre, noté $\varepsilon$. 
 Soit $x$ et $y$ deux éléments de
$E$ qui admettent un symétrique. Déterminer un symétrique de $x\star
y$. 

 \item Retrouver la caractérisation de l'injectivité d'un morphisme à partir
de son noyau.
\end{enumerate}
\finenonce{007349}
\finexercice


\exercice{7350, mourougane, 2021/08/10}
\enonce{007350}{}
Considérons le groupe $(\Zz , +)$.
Le sous-ensemble $4\Zz$ est par définition l'ensemble des multiples
entiers de $4$, autrement dit
$$4\Zz =\{p\in \Zz , \exists m\in \Zz \ \ p=4m\}.$$
\begin{enumerate}
 \item Les sous-ensembles $4\Zz$ et $6\Zz$ sont-ils stables pour la loi $+$ ? Sont-ils alors des sous-groupes ?
 \item Le sous-ensemble $4\Zz \cap 6\Zz$ est-il un sous-groupe ?
(indication~: on pourra énumérer ses premiers éléments positifs)
 \item Le sous-ensemble $2\Zz \cup 3\Zz$ est-il un sous-groupe ?
\end{enumerate}
\finenonce{007350}
\finexercice


\exercice{7351, mourougane, 2021/08/10}
\enonce{007351}{Intersection-Réunion}
Soit $E$ un ensemble et $\mathcal{P}(E)$ l'ensemble de ses parties.
\begin{enumerate}
 \item Déterminer les propriétés de la loi $(\mathcal{P}(E),\cup )$
(commutativité, associativité, existence et unicité des éléments
neutres, éléments symétrisables).

 \item Déterminer les propriétés de la loi $(\mathcal{P}(E),\cap )$.

 \item La loi $\cup$ est-elle distributive par rapport à la loi $\cap$ ?
\end{enumerate}
\finenonce{007351}
\finexercice


\exercice{7352, mourougane, 2021/08/10}
\enonce{007352}{Sur un ensemble à deux éléments}
Soit $E=\{a,b\}$ un ensemble à deux éléments.
On cherche à déterminer toutes les structures de groupe sur cet
ensemble. 
\begin{enumerate}
 \item Si $a$ est l'élément neutre, quel doit être le symétrique de $b$ ?
En déduire la table de multiplication dans ce cas.

 \item Combien y a-t-il de structures de groupes différentes sur $E$ ?
\end{enumerate}
\finenonce{007352}
\finexercice


\exercice{7353, mourougane, 2021/08/10}
\enonce{007353}{Groupes ordonnés}
Soit $G$ un ensemble muni d'une relation d'ordre $\leq$ et d'une
structure de groupe $\star$.
On dit que $(G,\star ,\leq)$ est un {\em groupe ordonné} si
$$\forall a\in G,\ \ \forall (x,y)\in G^2,\ \ \left(x\leq y\Rightarrow
 \left( a\star x \leq a\star y \quad \text{ et } \quad x\star a \leq y\star a\right)\right).$$
\begin{enumerate}
 \item $(\Rr ,+, \leq)$ est-il un groupe ordonné ?
 \item $(\Rr^\star ,\times , \leq)$ est-il un groupe ordonné ?
\end{enumerate}
\finenonce{007353}
\finexercice


\exercice{7375, mourougane, 2021/08/10}
\enonce{007375}{}
\begin{enumerate}
 \item Quelle loi fait de l'ensemble de inversibles de l'anneau $\Z/17\Z$ un groupe ? Ce groupe est-il cyclique ?

 \item Combien un groupe cyclique d'ordre $91$ a-t-il de générateurs ?

 \item Décomposer en produits de cycles à supports disjoints, la permutation $$s=\left(\begin{array}{ccccccccc} 
1&2&3&4&5&6&7&8&9\\
5&7&2&6&1&9&3&8&4\end{array}
\right).$$

 \item Décomposer en produits de transpositions, la permutation
$$\sigma=\left(\begin{array}{ccccccccc} 
1&2&3&4&5&6&7&8&9\\
9&8&6&7&2&3&1&4&5\end{array}
\right).$$

 \item Le polynôme $P(t_1,t_2,t_3,t_4)=t_1t_2+t_2t_3+t_3t_4+t_4t_1$ en quatre variables est-il symétrique ?

\end{enumerate}
\finenonce{007375}
\finexercice


\exercice{7376, mourougane, 2021/08/10}
\enonce{007376}{}
\begin{enumerate}
 \item Combien l'anneau $\Z/24\Z$ a-t-il d'éléments inversibles ? 

 \item Donner un isomorphisme entre $\Z/24\Z$ et un groupe produit.

 \item Le groupe des inversibles de $\Z/24\Z$ est-il cyclique ?
Si oui, combien a-t-il de générateurs ?
\end{enumerate}
\finenonce{007376}
\finexercice


\exercice{7377, mourougane, 2021/08/10}
\enonce{007377}{}
\begin{enumerate}
 \item Factoriser $221$ en produit de nombres premiers.
 \item Déterminer l'ensemble des solutions entières de l'équation
$$(X-3)(X-5)=0\pmod{221}.$$
\end{enumerate}
\finenonce{007377}
\finexercice


\exercice{7378, mourougane, 2021/08/10}
\enonce{007378}{}
\begin{enumerate}
 \item Sur le cercle unité, on a marqué les éléments du groupe des racines de l'unité d'ordre $8$. Soit
$\zeta = \exp(2i \pi/8)$.
Représenter sur un cercle unité les éléments du sous-groupe engendré par~$\zeta^3$.

 \item Soit la racine de l'unité $z = \exp(16 i \pi/11)$. 
Déterminer son ordre dans le groupe $\C^\star$.
Déterminer l'ordre de $z^8$. Déterminer l'argument de $z^8$.

\end{enumerate}
\finenonce{007378}
\finexercice


\exercice{7379, mourougane, 2021/08/10}
\enonce{007379}{}

\begin{enumerate}
 \item Calculer, avec l'algorithme de Hörner appliqué au polynôme $P(x)=x^2-7$, une valeur approchée à $10^{-1}$ près de $\sqrt{7}$.

 \item Calculer avec l'algorithme de Hörner une valeur approchée à $10^{-2}$ près de $\sqrt{7}$.
\end{enumerate}
\finenonce{007379}
\finexercice


\exercice{7380, mourougane, 2021/08/10}
\enonce{007380}{}
\begin{enumerate}
 \item Combien l'anneau $\Z/24\Z$ a-t-il d'éléments inversibles ? 
 \item Quelle loi fait de l'ensemble des inversibles de l'anneau $\Z/24\Z$ un groupe ?
 \item Le groupe des inversibles de $\Z/17\Z$ est-il cyclique ?
Si oui, combien a-t-il de générateurs ?
 \item Combien un groupe cyclique d'ordre $91$ a-t-il de générateurs ?
 \item Décomposer en produits de cycles à supports disjoints la permutation $$\left(\begin{array}{ccccccccc} 
1&2&3&4&5&6&7&8&9\\
5&7&8&6&1&3&9&4&2\end{array}
\right).$$
 \item Décomposer en produits de transpositions la permutation
$$\left(\begin{array}{ccccccccc} 1&2&3&4&5&6&7&8&9\\
9&8&6&7&2&3&1&4&5\end{array}
\right).$$
\end{enumerate}
\finenonce{007380}
\finexercice


\exercice{7381, mourougane, 2021/08/10}
\enonce{007381}{}
\begin{enumerate}
 \item Donner un isomorphisme entre $\Z/24\Z$ et un groupe produit.
 \item Le groupe des inversibles de $\Z/24\Z$ est-il cyclique ?
Si oui, combien a-t-il de générateurs ?
\end{enumerate}
\finenonce{007381}
\finexercice


\exercice{7382, mourougane, 2021/08/10}
\enonce{007382}{}
\begin{enumerate}
 \item Factoriser $221$ en produit de nombres premiers.
 \item Déterminer l'ensemble des solutions de l'équation
$$(X-3)(X-5)=0\pmod{221}.$$
\end{enumerate}
\finenonce{007382}
\finexercice


\exercice{7383, mourougane, 2021/08/10}
\enonce{007383}{}
\begin{enumerate}
 \item Donner la définition d'un idéal dans un anneau commutatif.
 \item Soit $I$ un idéal de l'anneau $\Z[i]$ des entiers de Gauss.
Montrer qu'il existe un entier de Gauss $a$ tel que $I$ soit l'ensemble des multiples de $a$.
\end{enumerate}
\finenonce{007383}
\finexercice


\exercice{7384, mourougane, 2021/08/10}
\enonce{007384}{}
Voici la table d'un groupe $G$.
Quel est l'ordre de $G$ ? Le groupe $G$ est-il nécessairement commutatif ?
 Compléter la table en énonçant précisément les propriétés utilisées.
$$\begin{array}{|c|c|c|c|c|c|c|}\hline
 * & a & b 	 & c &d  & e &f\\ \hline
 a & e & d  &   & b &a  & \\ \hline
 b &   & e  &   & f &  &d \\ \hline
 c &	b &f 	 &   & e & & \\ \hline
 d &f &	 &  e &c  & &b \\ \hline
 e &	 & b	 &  &  & &\\ \hline
 f & d  & c  & b &  & & e\\ \hline
\end{array}
$$
Le groupe est-il cyclique ?
\finenonce{007384}
\finexercice


\exercice{7385, mourougane, 2021/08/10}
\enonce{007385}{}
On rappelle que $\mathbb{F}_5$ désigne le corps $\Z/5\Z$ à cinq éléments.
\begin{enumerate}
 \item Montrer que l'anneau quotient $\mathbb{F}_{25}:=\mathbb{F}_5[X]/(X^2+2X+1)$ est un corps. Combien a-t-il d'éléments ?
 \item Si $a$ est un élément non nul de $\mathbb{F}_{25}$, que valent $5a$, $a^2+2a+1$ et $a^{24}$ ?
 \item Soit $x$ un élément de $\mathbb{F}_{25}$ tel que $x^2+3x+3=0$. Quel est l'ordre multiplicatif de $x$ ? 
\end{enumerate}
\finenonce{007385}
\finexercice


\exercice{7386, mourougane, 2021/08/10}
\enonce{007386}{}
\begin{enumerate}
 \item Quels sont les ordres possibles des éléments de $\left((F_{19})^\times ,\times\right)$ ?
 \item Le groupe $\left((F_{19})^\times ,\times\right)$ est-il cyclique ? si oui, combien a-t-il de générateurs ?
 \item Quel est l'ordre multiplicatif de l'élément $12$ du corps $F_{19}$ ? 
\end{enumerate}
\finenonce{007386}
\finexercice


\exercice{7387, mourougane, 2021/08/10}
\enonce{007387}{}
On admet que $271$ et $281$ sont deux nombres premiers.
\begin{enumerate}
 \item Peut-on écrire $271$ comme somme de deux carrés ? Si oui, faîtes le. On pourra trouver une racine $c$ de $-1$ modulo $271$, puis calculer le $\pgcd(271, c+i)$ dans l'anneau $\Z[i]$ des entiers de Gauss.
 \item Peut-on écrire $281$ comme somme de deux carrés ? Si oui, faîtes le. On pourra trouver une racine $c$ de $-1$ modulo $281$, puis calculer le $\pgcd(281, c+i)$ dans l'anneau $\Z[i]$ des entiers de Gauss.
\end{enumerate}
\finenonce{007387}
\finexercice


\exercice{7388, mourougane, 2021/08/10}
\enonce{007388}{}
\begin{enumerate}
 \item Rappeler la liste à isomorphisme près des groupes d'ordre $4$.
 \item Rappeler la liste à isomorphisme près des groupes d'ordre $5$.
 \item Le nombre premier $173$ peut-il s'écrire comme somme de deux carrés ?
 \item Quelle loi fait de $\Z/135\Z$ un groupe ? Combien le groupe $\Z/135\Z$ a-t-il de générateurs ?
\end{enumerate}
\finenonce{007388}
\finexercice


\exercice{7389, mourougane, 2021/08/10}
\enonce{007389}{}
\begin{enumerate}
 \item Donner la définition d'un idéal dans un anneau commutatif.
 \item Soit $I$ un idéal de l'anneau $\Z[i]$ des entiers de Gauss.
Montrer qu'il existe un entier de Gauss $a$ tel que $I$ soit l'ensemble des multiples de $a$.
\end{enumerate}
\finenonce{007389}
\finexercice


\exercice{7390, mourougane, 2021/08/10}
\enonce{007390}{}
On admet $281$ est un nombre premier.
\begin{enumerate}
 \item Peut-on écrire $281$ comme somme de deux carrés ? Si oui, faîtes le.
 \item Factoriser $280$. On choisit $x=3$. On admet que $3^{35}=60[281]$. Déterminer une racine $c$ de $-1$ modulo $281$.
 \item Calculer le $\pgcd(281, c+i)$ dans l'anneau $\Z[i]$ des entiers de Gauss.
\end{enumerate}
\finenonce{007390}

\finexercice


\exercice{7391, mourougane, 2021/08/10}
\enonce{007391}{}
\begin{enumerate}
 \item Donner la définition d'un groupe.
 \item Dire si les couples suivants sont des groupes : $(\Z,+)$ ; $(\Z,\times)$ ; $(\C^*,+)$ ; $(\C^*, \times)$. Lorsque la réponse est "non", on indiquera une propriété des groupes qui fait défaut (on ne demande pas de justification lorsque la réponse est "oui").
 \item Quels sont les idéaux de l'anneau $(\Z,+,\times)$ ?
 \item Quels sont les idéaux de l'anneau $(\R[X],+,\times)$ ?
\end{enumerate}
\finenonce{007391}
\finexercice


\exercice{7392, mourougane, 2021/08/10}
\enonce{007392}{}
\begin{enumerate}
 \item L'entier $51606$ appartient-il à $2569 + 247\Z$ ? 
 \item L'entier $-1601$ est-il un représentant de la classe $[-7387]_{ 2893}$ de $\Z/2893\Z$? 
 \item Calculer l'élément $2169$ dans $\Z/17\Z$. Le résultat doit être représenté par un nombre compris entre $0$ et $16$. 
 \item Calculer l'élément $11 ^{329}$ dans $\Z/13\Z$. Le résultat doit être représenté par un nombre compris entre $0$ et $12$. 
\end{enumerate}
\finenonce{007392}
\finexercice


\exercice{7393, mourougane, 2021/08/10}
\enonce{007393}{}
Voici la table d'un groupe $G$.
Quel est l'ordre de $G$ ? Le groupe $G$ est-il nécessairement commutatif ?
 Compléter la table en énonçant précisément les propriétés utilisées.
$$\begin{array}{|c|c|c|c|c|c|}\hline
 * & a & b 	& c &	d & e\\ \hline
 a & 	 & 	& d &	 & c \\ \hline
 b & e &	&a  &	 & d \\ \hline
 c &	 &a 	 &  &  &	 \\ \hline
 d &	 &	 &  &d &e \\ \hline
 e &	 &	 &b &	 & a \\ \hline
\end{array}
$$
\finenonce{007393}
\finexercice


\exercice{7394, mourougane, 2021/08/10}
\enonce{007394}{}
Considérons le groupe
 $\displaystyle\left(\mathbb{F}_{41} \right)^\times$.
Quels sont les ordres possibles d'un élément de ce groupe ?
Combien d'éléments de chaque ordre ce groupe possède -t-il ?
Déterminer en le justifiant un générateur de ce groupe.
\finenonce{007394}
\finexercice


\exercice{7395, mourougane, 2021/08/10}
\enonce{007395}{}
\begin{enumerate}
  \item Montrer que le polynôme $X^3+X+1$ est irréductible dans $\mathbb{F}_2[X]$.
 \item Quelle est alors la structure de l'ensemble quotient $A=\mathbb{F}_2[X]/(X^3+X+1)$ ? 
 \item Quel est le cardinal de $A$ ? Soit $\alpha$ la classe du polynôme $X$ dans ce quotient $A$. Donner la liste des éléments de $A$.
 \item  Sans calculs, mais en justifiant votre réponse, dire ce que valent les quantités suivantes 
$$\alpha+\alpha, \ \ \alpha^3+\alpha+1, \ \ \ \alpha^7.$$
 \item Déterminer l'ordre multiplicatif de $\alpha$ dans $A^\times$.
 \item Etablir la table des puissances de $\alpha$.
 \item Calculer $(1+\alpha) (1+\alpha^2)$.
 \item Calculer $\alpha^3+\alpha^2$ comme puissance de $\alpha$.
 \item Calculer $(1+\alpha)^{-1}$.
\end{enumerate}
\finenonce{007395}
\finexercice


\exercice{7396, mourougane, 2021/08/10}
\enonce{007396}{}
\begin{enumerate}
 \item Enoncer le théorème de Lagrange.
 \item Soit $G$ un groupe et $a$ un élément d'ordre $k$ dans $G$. Soit $p$ un entier naturel. Quel est l'ordre de $a^p$ ? 
 \item Le groupe $\mathcal{S}_3$ des permutations de $\{1,2,3\}$ est-il cyclique ? (justifier) 
 \item Enoncer le théorème de la division euclidienne dans $k[X]$.
 \item Enoncer le théorème de la division euclidienne dans $\Z[i]$.
\end{enumerate}
\finenonce{007396}
\finexercice


\exercice{7397, mourougane, 2021/08/10}
\enonce{007397}{}
\begin{enumerate}
 \item La classe $[51]$ est-elle inversible dans l'anneau $\Z/131\Z$. Si oui, calculer $92\times 51^{-1}$ dans $\Z/131\Z$. Le résultat doit être représenté par un nombre compris entre 0 et 130. 
 \item Trouver l'ensemble des diviseurs de zéro dans $\Z/16\Z$. (Dans cet exercice on ne considère pas le 0 comme un diviseur de zéro.) Représenter chaque classe par un nombre compris entre 1 et 15.
\end{enumerate}
\finenonce{007397}
\finexercice


\exercice{7398, mourougane, 2021/08/10}
\enonce{007398}{}
\begin{enumerate}
 \item On rappelle que le seul polynôme irréductible de degré $2$ sur $\mathbb{F}_2$ est $X^2+X+1$.
 Montrer que le polynôme $X^4+X+1$ est irréductible dans $\mathbb{F}_2[X]$.
 \item On note $A:=\mathbb{F}_2[X]/<X^4+X+1>$ l'anneau quotient de $\mathbb{F}_2[X]$ par l'idéal engendré par $P$. La classe de $3X^5+X^2+X+7$ est-elle nulle dans $A$ ?
 L'anneau $A$ est-il un corps ? Combien a-t-il d'éléments ?
 \item On note $\alpha$ la classe du polynôme $X$ dans $A$.
Déterminer $\alpha^4$ et $\alpha^{15}$ comme polynômes de degré au plus $3$ en $\alpha$.
 \item Le polynôme $X^{15}-1$ est-il un multiple de $X^4+X+1$ dans $\mathbb{F}_2[X]$ ? 
\end{enumerate}
\finenonce{007398}
\finexercice


\exercice{7399, mourougane, 2021/08/10}
\enonce{007399}{}
\begin{enumerate}
 \item On considère le code binaire, linéaire engendré par la matrice
$$\left(\begin{array}{ccccccccccccccc}
1&0&0&1&1&0&0&0&0&0&0&0&0&0&0\\
0&1&0&0&1&1&0&0&0&0&0&0&0&0&0\\
0&0&1&0&0&1&1&0&0&0&0&0&0&0&0\\
0&0&0&1&0&0&1&1&0&0&0&0&0&0&0\\
0&0&0&0&1&0&0&1&1&0&0&0&0&0&0\\
0&0&0&0&0&1&0&0&1&1&0&0&0&0&0\\
0&0&0&0&0&0&1&0&0&1&1&0&0&0&0\\
0&0&0&0&0&0&0&1&0&0&1&1&0&0&0\\
0&0&0&0&0&0&0&0&1&0&0&1&1&0&0\\
0&0&0&0&0&0&0&0&0&1&0&0&1&1&0\\
0&0&0&0&0&0&0&0&0&0&1&0&0&1&1\\
\end{array}\right)$$
Quel est son alphabet ? sa longueur ? sa dimension ? un polynôme générateur ? son nombre de mots ?
 \item Le code est-il cyclique ?
 \item Ecrire une matrice de contrôle. Montrer que la distance du code est au moins $3$. Combien d'erreurs peut-on alors détecter ? combien d'erreurs peut-on alors corriger ?
 \item Le mot $(0,0,0,1,0,1,1,1,1,1,1,0,0,1,1)$
est-il un mot de code ? Si non, en supposant qu'il n'a qu'une erreur, écrire le mot de code dont il provient.
\end{enumerate}
\finenonce{007399}
\finexercice


\exercice{7400, mourougane, 2021/08/10}
\enonce{007400}{}
\begin{enumerate}
 \item Le nombre $613$ est-il premier ?
 \item Peut-il s'écrire comme somme de deux carrés ?
 \item Calculer $35^2$ modulo $613$.
 \item Effectuer la division euclidienne de $613$ par $35+i$ dans l'anneau $\Z[i]$ des entiers de Gauss.
 \item Ecrire $613$ comme somme de deux carrés.
\end{enumerate}
\finenonce{007400}
\finexercice


\exercice{7401, mourougane, 2021/08/10}
\enonce{007401}{}
\begin{enumerate}
  \item Enoncer le théorème de la division euclidienne dans $k[X]$.
  \item Enoncer le théorème de la division euclidienne dans $\Z[i]$.
  \item Soit $G$ un groupe et $a$ un élément de $G$ d'ordre $n$. Soit $k$ un entier naturel. Quel est l'ordre de $a^k$ ?
  \item Démontrer que tout groupe d'ordre $13$ est commutatif.
  \item Donner l'exemple d'un nombre premier qui ne peut pas s'écrire comme somme de deux carrés.
\end{enumerate}
\finenonce{007401}
\finexercice


\exercice{7402, mourougane, 2021/08/10}
\enonce{007402}{}
\begin{enumerate}
 \item Les groupes $\Z/9\Z$ et $\Z/3\Z\times \Z/3\Z$ sont-ils isomorphes ?
 \item Les groupes $\mathbb{F}_7$ et $\Z/2\Z\times \Z/3\Z$ sont-ils isomorphes ? Justifier.
 \item Les groupes $(\mathbb{F}_7)^\star$ et $\Z/2\Z\times \Z/3\Z$ sont-ils isomorphes ? Justifier.
\end{enumerate}
\finenonce{007402}
\finexercice


\exercice{7403, mourougane, 2021/08/10}
\enonce{007403}{}
\begin{enumerate}
	\item On rappelle que le seul polynôme irréductible de degré $2$ sur $\mathbb{F}_2$ est $X^2+X+1$.
	Montrer que le polynôme $X^4+X^3+1$ est irréductible dans $\mathbb{F}_2[X]$.
	\item On note $A:=\mathbb{F}_2[X]/(X^4+X^3+1)$ l'anneau quotient de	$\mathbb{F}_2[X]$ par l'idéal engendré par $P$.	L'anneau $A$ est-il un corps ? Combien a-t-il d'éléments ?
	\item On note $\alpha$ la classe du polynôme $X$ dans $A$. Déterminer $\alpha^4$ et $\alpha^{15}$ comme polynômes de degré au
	plus $3$ en $\alpha$.
	\item Déterminer toutes les puissances de $\alpha$, $\alpha$ jusqu'à $\alpha^{15}$, comme polynômes de degré au plus~$3$ en $\alpha$.
	\item Déterminer $\alpha^7+\alpha^8+\alpha^9$ comme polynômes de degré au plus $3$ en $\alpha$.
	\item Ecrire l'inverse de $1+\alpha+\alpha^3$ comme	puissance de $\alpha$.
\end{enumerate}
\finenonce{007403}
\finexercice


\exercice{7404, mourougane, 2021/08/10}
\enonce{007404}{}
Alice et Bernard décident d'utiliser l'algorithme d'El Gamal. Il utilise le corps $\mathbb{F}_{13}$ avec l'élément $G=2$.
\begin{enumerate}
	\item Quels sont les ordres possibles des éléments de $\mathbb{F}_{19}^\times$. Déterminer l'ordre de $2$ dans $\mathbb{F}_{19}^\times$.
	\item Bernard choisit sa clé privée $c=3$. Déterminer sa clé publique $C=G^c$.
	\item Alice choisit une clé temporaire privée
	$d=7$. Quelle est sa clé publique $D$ ? Elle souhaite envoyer le
	message $m=11$. Elle le chiffre en utilisant la clé publique $C$ de
	Bernard par $(M_1,M_2)=(D,mC^d)$. Expliciter ce message chiffré.
	\item Comment Bernard retrouve-t-il le message $m$ ?
	\item Dans un second envoi, Bernard reçoit $(8,3)$. Quel est le message
	$m$ envoyé cette fois par Alice ? Quelle clé privée a-t-elle utilisé
	cette fois ?
\end{enumerate}
\finenonce{007404}
\finexercice


\exercice{7405, mourougane, 2021/08/10}
\enonce{007405}{Questions de cours}
\begin{enumerate}
 \item \'Enoncer le théorème de Lagrange.
 \item Donner l'exemple de deux groupes finis de même ordre non isomorphes. Justifier le fait que les deux groupes choisis ne sont pas isomorphes.
 \item Donner la définition d'un sous-groupe \emph{distingué}.
 \item L'anneau $(\Z/24\Z,+,\times)$ est-il intègre ? Justifier.
 \item Effectuer la division euclidienne de $X^3+2X^2-5X+8$ par $X^2-1$ dans $\Z/5\Z[X]$.
 \item Effectuer la division euclidienne de $X^3+2X^2-5X+8$ par $2X^2-1$ dans $\Z/5\Z[X]$.
\end{enumerate}
\finenonce{007405}
\finexercice


\exercice{7406, mourougane, 2021/08/10}
\enonce{007406}{}
\begin{enumerate}
 \item Calculer les produits dans $\mathcal{S}_7$, $(1,2)(1,3)$ et $(1,2)(2,3)(1,2)$.
 \item L'ensemble des transpositions de $\mathcal{S}_7$ est-il un sous-groupe de $\mathcal{S}_7$ ?
\end{enumerate}
\finenonce{007406}
\finexercice


\exercice{7407, mourougane, 2021/08/10}
\enonce{007407}{}
\begin{enumerate}
 \item Montrer que $1$ est dans le sous-groupe de $(\Z,+)$ engendré par $\{3,8\}$.
 \item Quel est le sous-groupe engendré par $\{3,8\}$ ?
 \item Quel est le sous-groupe engendré par $\{3,8,15\}$ ?
\end{enumerate}
\finenonce{007407}
\finexercice


\exercice{7408, mourougane, 2021/08/10}
\enonce{007408}{}
\begin{enumerate}
 \item Décomposer en produit de cycles à supports disjoints la permutation
$$\sigma=\left(\begin{array}{ccccccc}
1&2&3&4&5&6&7 \\
7&6&3&1&4&2&5
    \end{array}
 \right).$$
 \item Déterminer son ordre.
 \item Ecrire sa puissance $\sigma ^6$.
 \item Ecrire son inverse en produit de cycles à supports disjoints.
 \item La décomposer en produit de moins de $7$ transpositions.
\end{enumerate}
\finenonce{007408}
\finexercice


\exercice{7409, mourougane, 2021/08/10}
\enonce{007409}{}
\begin{enumerate}
 \item Quels sont les éléments inversibles de l'anneau $(\Z/11\Z)$ ?
 \item Déterminer un générateur $g$ du groupe $(\Z/11\Z)^\star$ des inversibles de l'anneau $(\Z/11\Z)$.
 \item Ecrire chaque élément de $(\Z/11\Z)^\star$ comme puissance de $g$.
\end{enumerate}
\finenonce{007409}
\finexercice

\section{204.01 Produit scalaire, norme}


\exercice{1450, barraud, 2003/09/01}

\enonce{001450}{}
  A deux polynômes $P=a_{0}+a_{1}X+a_{2}X^{2}$ et
  $Q=b_{0}+b_{1}X+b_{2}X^{2}$ de $\R_{2}[X]$, on associe
  $$
  <P,Q>=(a_{0}+a_{1})b_{0}+(a_{0}+3a_{1})b_{1}+3a_{2}b_{2}
  $$ Montrer qu'il s'agit d'un produit scalaire.  
\finenonce{001450}



\finexercice



\exercice{1451, barraud, 2003/09/01}

\enonce{001451}{}
  Pour quelles valeurs de $\lambda$ les formes bilinéaires ci-dessous
  définissent-elles un produit scalaire sur $\R^{3}$ ?
  \begin{enumerate}
  \item 
    $f(x,y)= x_{1}y_{1}
           +6x_{2}y_{2}
           +3x_{3}y_{3}
           +2x_{1}y_{2}
           +2x_{2}y_{1}
           +3\lambda x_{1}y_{3}
           +3\lambda x_{3}y_{1}
           $
  \item 
    $g(x,y)=  x_{1}y_{1}
           +10x_{2}y_{2}
           + 6x_{1}y_{2}
           +\lambda x_{3}y_{3}
           - x_{2}y_{3}
           - x_{3}y_{2}
           $
  \item 
    $h(x,y)= 2x_{1}y_{1}
           + 7x_{1}y_{2}
           + 7x_{2}y_{1}
           + 8x_{2}y_{2}
           - 3x_{3}y_{3}
           +\lambda x_{2}y_{3}
           +\lambda x_{3}y_{2}
           $

 \item 
   $i(x,y)=(x_{1}+x_{2})(y_{1}+y_{2})
          +(x_{1}+x_{3})(y_{1}+y_{3})
          +(x_{2}+x_{3})(y_{2}+y_{3})
          -\lambda(x_{1}+x_{2}+x_{3})(y_{1}+y_{2}+y_{3})$
 \end{enumerate}
\finenonce{001451}



\finexercice



\exercice{1452, barraud, 2003/09/01}

\enonce{001452}{}
Vérifier que l'application $\phi~:\R^{3}\times\R^{3}\rightarrow \R$ définie
ci-dessous est une forme bilinéaire symétrique et déterminer la forme
quadratique qui lui est associée~:
$$
 \phi\big((x,y,z),(x',y',z')\big)=
 xx'+ 2yy' + 2yz' + 2y'z + zz'.
$$
S'agit-il d'un produit scalaire~?
\medskip

Vérifier que l'application $q~:\R^{3}\rightarrow \R$ définie ci-dessous est
une forme quadratique et déterminer la forme bilinéaire symétrique qui
lui est associée~:
$$
 q\big((x,y,z)\big)=
 x^{2}+ 3(x+y-z)^{2} + (z-y)^{2}.
$$
S'agit-il d'une norme euclidienne~?

\finenonce{001452}



\finexercice



\exercice{1453, barraud, 2003/09/01}

\enonce{001453}{}
  Sur $\R_{3}[X]$ on considère les formes bilinéaires suivantes. Dire
  lesquelles sont des produits scalaire. 
  \begin{align*}
  \phi(P,Q)&=\int_{-1}^{1}P(t)Q(t)dt\\
  \phi(P,Q)&=\int_{-1}^{1}P'(t)Q(t)+P'(t)Q(t)dt\\
  \phi(P,Q)&=\int_{-1}^{1}P'(t)Q'(t)dt + P(0)Q(0)
\end{align*}
\finenonce{001453}



\finexercice



\exercice{1454, barraud, 2003/09/01}

\enonce{001454}{}
  Pour quelles valeurs de $\lambda$ les formes bilinéaires ci-dessous
  définissent-elles un produit scalaire sur $\R^{3}$ ?
  \begin{enumerate}
  \item 
    $f(x,y)= x_{1}y_{1}
           +6x_{2}y_{2}
           +3x_{3}y_{3}
           +2x_{1}y_{2}
           +2x_{2}y_{1}
           +3\lambda x_{1}y_{3}
           +3\lambda x_{3}y_{1}
           $
  \item 
    $g(x,y)=  x_{1}y_{1}
           +10x_{2}y_{2}
           + 6x_{1}y_{2}
           +\lambda x_{3}y_{3}
           - x_{2}y_{3}
           - x_{3}y_{2}
           $
  \item 
    $h(x,y)= 2x_{1}y_{1}
           + 7x_{1}y_{2}
           + 7x_{2}y_{1}
           + 8x_{2}y_{2}
           - 3x_{3}y_{3}
           +\lambda x_{2}y_{3}
           +\lambda x_{3}y_{2}
           $

 \item 
   $i(x,y)=(x_{1}+x_{2})(y_{1}+y_{2})
          +(x_{1}+x_{3})(y_{1}+y_{3})
          +(x_{2}+x_{3})(y_{2}+y_{3})
          -\lambda(x_{1}+x_{2}+x_{3})(y_{1}+y_{2}+y_{3})$
 \end{enumerate}
\finenonce{001454}



\finexercice



\exercice{1455, legall, 1998/09/01}

\enonce{001455}{}
Soient $  x = (x_1, x_2)  $ et $  y = (y_1, y_2)  $
appartenant \`a $  \R ^2  .$
Pour quelles valeurs de $  a, b, c, d \in \R   $ l'application $  f(x, y) =
ax_1y_1 + bx_1y_2 + cx_2y_1 + dx_2y_2  $ est-elle un produit scalaire sur
$  \R ^2  $?
\finenonce{001455}



\finexercice



\exercice{1456, legall, 1998/09/01}

\enonce{001456}{}
Soient $x,y$ et $z$ trois r\'eels tels que $ x^2+2y^2+3z^2\leq 1$.
Montrer l'in\' egalit\'e :
${(x+y+z)^2
 \leq \frac{11}{6}}.$ (On pourra par exemple appliquer l'in\'egalit\'e de Cauchy-Schwarz \`a
certains vecteurs de $  \R ^3  $ pour un produit scalaire bien choisi.)
\finenonce{001456}



\finexercice



\exercice{1457, legall, 1998/09/01}

\enonce{001457}{}
Soient $  x,y  $ et $  z  $ trois r\'eels tels que $  x^2+y^2+z^2\leq
1  .$ Montrer que $  (x+2y+3z)^2\leq 14   .$
\finenonce{001457}



\finexercice



\exercice{1458, legall, 1998/09/01}

\enonce{001458}{}
 Soient $E$ un $\R$-espace vectoriel non nul, $\varphi$ un produit
scalaire sur $E$, $(a,b,c) \in \R^3$. $\psi : E\times E \to \R$
l'application d\'efinie par $\psi(x,y) = a\varphi(x,x) + b\varphi(x,y) +
c\varphi(y,y)$. Trouver une condition n\'ecessaire et suffisante sur
$(a,b,c)$ pour que $\psi$ soit un produit scalaire sur $E$.

\finenonce{001458}



\finexercice



\exercice{1459, legall, 1998/09/01}

\enonce{001459}{}
\begin{enumerate}
    \item  Soient $  (E, \langle   ,  \rangle )  $ un espace euclidien et $
\Vert   .   \Vert   $ la norme associ\'ee~; $  n \in \N ^*  ,$ et $  v_1, \dots, v_n \in E  .$
Montrer l'in\'egalit\'e~: $  \displaystyle\Vert \sum_{i=1}^n v_i\Vert^2  \leq\ n\sum_{i=1}^n \Vert
v_i\Vert^2  .$
    \item Soient $  n \in \N ^*  ,$ $  x_1, \dots, x_n \in \R _+^*  $ tels que
$  \displaystyle{\sum_{i=1}^n x_i = 1}  .$ Montrer que $  \displaystyle{\sum_{i=1}^n
\frac{1}{ x_i} \geq n^2}  .$ Etudier le cas d'\'egalit\'e.
\end{enumerate}
\finenonce{001459}



\finexercice



\exercice{1460, barraud, 2003/09/01}

\enonce{001460}{}
  Montrer que 
  $$
  \forall(x_{1},...,x_{n})\in\R^{n},\sum_{i=1}^{n}{x_{i}}\leq
  \sqrt{n}\left(\sum_{i=1}^{n}x_{i}^{2}\right)^{\frac{1}{2}}.
  $$
  Etudier le cas d'égalité.

  Soit $f$ et $g$ deux applications continues de $[0,1]$ dans
  $\R$. Montrer que :
  $$
  \forall (f,g)\in C^{0}([0,1],\R)\qquad
  \left(\int_{0}^{1}f(t)g(t)dt\right)^{2}
    \leq \int_{0}^{1}f^{2}(t)dt\int_{0}^{1}g^{2}(t)dt.
   $$
   Etudier le cas d'égalité.

   Soit $f$ une application continue d'un intervalle $[a,b]$ dans
   $\R$. Montrer que :
   $$
  \forall f\in C^{0}([a,b],\R)\qquad
  \left(\int_{a}^{b}f(t)dt\right)^{2}
    \leq (b-a)\int_{a}^{b}f^{2}(t)dt.
  $$
  Etudier le cas d'égalité.
\finenonce{001460}



\finexercice



\exercice{1461, barraud, 2003/09/01}

\enonce{001461}{}
Rappeler l'énoncé de l'inégalité de Cauchy-Schwarz.

Montrer que pour toute fonction continue d'un segment $[a,b]$ dans
$\R$, on a 
$$
  \left(\int_{a}^{b}{f(t)dt}\right)^{2}\leq
  (b-a)\int_{a}^{b}{\Big(f(t)\Big)^{2}dt}
$$

Pour quelles fonctions a-t-on l'égalité ?
\finenonce{001461}



\finexercice



\exercice{1462, ridde, 1999/11/01}

\enonce{001462}{}
Soit $E = \left\{ f : \Rr \rightarrow \Rr \text{ continue $2\pi$-p\'eriodique}\right\}$.
Montrer que $\langle f|g \rangle = \int_{0}^{2\pi}{f (t)g (t)dt}$ est un produit
scalaire sur $E$.
\finenonce{001462}



\finexercice



\exercice{1463, ridde, 1999/11/01}

\enonce{001463}{}
Soit $E = \Rr_{n}[X]$. Montrer que $\langle P|Q \rangle = \sum\limits_{k = 0}^{n}{P (k)Q (k)}$
est un produit scalaire sur $E$.
\finenonce{001463}



\finexercice



\exercice{1464, ridde, 1999/11/01}

\enonce{001464}{}
Soit $E$ un espace euclidien et $f$ et $g$ deux fonctions de $E$ dans $E$ qui v\'erifient :
$\forall (x, y)\in E^{2} \, \langle f (x)|y \rangle  = \langle x|g (y)\rangle$.
Montrer que $f$ et $g$ sont lin\'eaires.
\finenonce{001464}



\finexercice



\exercice{1465, ridde, 1999/11/01}

\enonce{001465}{}
Soient $a_{1}, \ldots, a_{n},b_{1}, \ldots, b_{n}, c_{1}, \ldots, c_{n} $ des
r\'eels positifs.\\
Montrer que $\sum\limits_{k = 1}^{n}{a_{k}b_{k}c_{k}} \leq
 \sum\limits_{k = 1}^{n}{a_{k}^{2}c_{k}} \sum\limits_{k = 1}^{n}{b_{k}^{2}c_{k}}$.
\finenonce{001465}



\finexercice



\exercice{1466, ridde, 1999/11/01}

\enonce{001466}{}
Soit $E$ un espace euclidien de dimension $n$ et $x_{1}, \ldots, x_{p}$ des vecteurs
de $E$ tels que si $i \neq j$ alors $\langle x_{i}|x_{j} \rangle <0$.
Montrer par r\'ecurrence sur $n$ que $p \leq n + 1$.
\finenonce{001466}



\finexercice



\exercice{1467, gourio, 2001/09/01}

\enonce{001467}{}
Soit $ E  $ un espace euclidien, et $ \left( e_{1},...,e_{n}\right)  $ des
vecteurs unitaires v\'{e}rifiant :
$$\forall x\in E,\left\| x\right\| ^{2}=\sum_{i=1}^{n}\left\langle
x,e_{i}\right\rangle ^{2}. $$
Montrer que $ \left( e_{1},...,e_{n}\right) $  est une base orthonormale (i.e.
une base qui est aussi une famille orthonormale).
(NB: on ne suppose pas que la dimension de l'espace est $ n.$)
\finenonce{001467}



\finexercice



\exercice{1468, gourio, 2001/09/01}

\enonce{001468}{}
  \begin{enumerate}
  \item Montrer que sur $ M_{n}({\Rr})$ l'application :
$$\left( A,B\right) \rightarrow tr(^{t}AB) $$
est un produit scalaire.
\item Soit $ N$ la norme associ\'{e}e, montrer que :
$$\forall (A,B)\in M_{n}({\Rr}),  N(AB)\leq N(A)N(B). $$
\item Montrer que :
$$\forall A\in M_{n}({\Rr}),\left| tr(A)\right| \leq \sqrt{n}N(A). $$
\end{enumerate}
\finenonce{001468}



\finexercice



\exercice{1469, gourio, 2001/09/01}

\enonce{001469}{}
Soit $ E$ un espace euclidien et $ f $ et $ g $ deux fonctions de $ E$ dans $ E$
telles que :
$$\forall (x,y)\in E^{2},\left\langle f(x),y\right\rangle =\left\langle
x,g(y)\right\rangle .$$
Montrer que $ f$ et $ g$ sont lin\'{e}aires.
\finenonce{001469}



\finexercice



\exercice{1470, gourio, 2001/09/01}

\enonce{001470}{}
Soit $ E$ un espace euclidien, montrer que :
$$\forall (x,y)\in E^{2},\left\| x+y\right\| ^{2}+1\leq 2\left( 1+\left\|
x\right\| ^{2}\right) \left( 1+\left\| y\right\| ^{2}\right) . $$
\finenonce{001470}



\finexercice



\exercice{1471, gourio, 2001/09/01}

\enonce{001471}{}
Soit $ E$ un espace euclidien et $ f:E\rightarrow E$ tel que $ f(0)=0 $ et :
$$\forall (x,y)\in E^{2},\left\| f(x)-f(y)\right\| =\left\| x-y\right\| . $$
Montrer que $ f$ est lin\'{e}aire.
\finenonce{001471}



\finexercice



\exercice{1472, gourio, 2001/09/01}

\enonce{001472}{}
On munit $ {\Rr}[X]$ du produit scalaire :
$$(P,Q)\rightarrow \int_{0}^{1}P(t)Q(t)dt. $$
Existe t-il $ A\in {\Rr}[X] $ tel que :
$$\forall P\in {\Rr}[X],(P|A)=P(0)\   ? $$
\finenonce{001472}



\finexercice



\exercice{1473, gourio, 2001/09/01}

\enonce{001473}{}
Soit $ E$ un espace euclidien et $ f $ un endomorphisme de $E,$ tel que :
$$\forall (x,y)\in E^{2},(x|y)=0\Rightarrow (f(x)|f(y))=0. $$
Montrer :
$$\exists \alpha \in {\Rr}^{+},\forall (x,y)\in E^{2},(f(x)|f(y))=\alpha(x|y).$$
\finenonce{001473}



\finexercice



\exercice{1474, legall, 2003/10/01}

\enonce{001474}{}
Soit $ (E, \langle   ,  \rangle )  $ un espace euclidien et $
f \in \mathcal{L} (E)
$ un endomorphisme tel que $\forall x, y \in E $ tels que $\langle 
x,y\rangle =0,$ on ait $ \langle f(x), f(y)\rangle =0 $.
Montrer qu'il existe $k\in \Rr _+$ tel que, pour tout $x\in E$~: 
$\Vert f(x)\Vert =k \Vert x\Vert .$

\finenonce{001474}



\finexercice



\exercice{2426, matexo1, 2002/02/01}

\enonce{002426}{}
On consid\`ere l'espace vectoriel $\R^n$ ($n \ge 2$), et $p \ge 1$
un nombre r\'eel. On d\'efinit l'application
$$ N_p: \R^n \to \R, x = (x_1,\ldots,x_n) \mapsto N_p(x) = \left( \sum_{i=1}^n
|x_i|^p \right)^{\frac 1 p}.$$
\begin{enumerate}
\item Montrer que $N_p$ est une norme sur $\R_n$\,; on la note
aussi $\|\cdot\|_p$. Dans les cas o\`u $p \neq 1$, $p \neq 2$, on
pourra s'aider de la relation suivante (admise)\,:
$$ \sum_{i=1}^n |a_i b_i| \leq 
\left(\sum_{i=1}^n |a_i|^p \right)^{\frac 1 p}
\left(\sum_{i=1}^n |b_i|^{\frac p{p-1}} \right)^{\frac{p-1}p}
\qquad\qquad \forall a, b\in \R^n.$$
\item On d\'efinit, pour tout $x\in \R^n$, $N_\infty(x) =
\lim_{p\to\infty} \|x\|_p$. Montrer que cette limite
existe pour tout $x$, et que $N_\infty(x) = \max_{1 \leq i \leq n} |x_i|$.
D\'emontrer
qu'il s'agit aussi d'une norme sur $\R^n$\,; on la note
aussi $\|\cdot\|_\infty$.
\item Dessiner les ``boules'' $\{x\in\R^2,\ N_p(x)\leq 1\}$
dans les cas o\`u $p=1$, $p=2$, $p=\infty$. \`A quoi ressemblent les cas
des valeurs interm\'ediaires\,?
\end{enumerate}
\finenonce{002426}

\finexercice


\exercice{2460, matexo1, 2002/02/01}

\enonce{002460}{}
{\sc Identit\'e du parall\'elogramme}

\bigskip

Soit $E$ un $\R$-espace vectoriel norm\'e. 

\begin{enumerate}
\item On suppose que la norme de $E$ v\'erifie la relation
\begin{equation} \label{eq}
\forall x,y\in E, \qquad
 2\bigl(\|x\|^2 +\|y\|^2\bigr)
=\|x+y\|^2 + \|x-y\|^2.
\end{equation}
On d\'efinit $p : E\times E \to \R$ par
$$p(x,y) =
\frac 1 2 \bigl(\|x+y\|^2 - \|x\|^2 -\|y\|^2\bigr).$$
Montrer que $p$ est un produit scalaire sur $E$.
\item R\'eciproquement, si $E$ est un espace euclidien dont le
produit scalaire est not\'e $\langle x,y \rangle$, montrer que la norme euclidienne
(d\'efinie par $\|x\| =\sqrt{\langle x,x \rangle}$) v\'erifie (\ref{eq}),
et que $\langle x,y \rangle=p(x,y)$.
\item Dans le cas o\`u $E=\R^n$, pour quelles valeurs de $q\geq 1$
les normes $\|\cdot\|_q$ v\'erifient-elles (\ref{eq})\,?

\end{enumerate}
\finenonce{002460}


\finexercice


\exercice{2461, matexo1, 2002/02/01}

\enonce{002461}{}
Soit $E=\R[X]$ l'espace des polyn\^omes \`a coefficients r\'eels. On
d\'efinit $p : E\times E\to \R$ par
$$ p(u, v)  = \int_{-1}^1 u(t) v(t)\,d t.$$

\begin{enumerate}
\item Montrer que $p$ est un produit scalaire sur $E$, et que $u
\mapsto \sqrt{p(u,u)}$ est une norme sur~$E$.
\item Soit $E_2 \subset E$ le sous-espace des polyn\^omes de degr\'e
$\leq 2$, et $u\in E$ d\'efini par $u(t) = t^4$. D\'eterminer la
projection
orthogonale $u_0$ de $u$ sur $E_2$. En d\'eduire la valeur de
$$ \inf_{(a,b,c)\in \R^3} \int_{-1}^1 (t^4 -at^2-bt -c)^2 \,d t.$$
\end{enumerate}

\finenonce{002461}


\finexercice


\exercice{2671, matexo1, 2002/02/01}

\enonce{002671}{}
\begin{enumerate}
\item Sur l'espace $\R^n$, d\'efinir une distance $d$, telle que
$\forall x,y \in \R^n$, on a $d(x,y)\leq 1$. Est-il possible de
choisir $d$ de sorte que, en plus, $x\mapsto \sqrt{d(x,0)}$ soit
une norme\,?
\item On consid\`ere l'ensemble $S$ constitu\'e des points de $\R^3$
de norme euclidienne \'egale \`a 1 (``sph\`ere unit\'e''). Si $x$, $y$,
sont deux \'el\'ements de $S$, il existe au moins un cercle de rayon 1
trac\'e sur $S$ et passant par ces points (dans quels cas en existe-t-il
plusieurs\,?)\,; on note $d(x,y)$ la longueur du plus court arc de ce
cercle joignant $x$ et $y$ (``distance g\'eod\'esique'' --- comparer
avec la distance sur la sph\`ere terrestre). Montrer que $d$ est une
distance sur $S$\,; pourquoi $x\mapsto \sqrt{d(x,0)}$ n'est-elle pas
une norme\,?
\end{enumerate}

\finenonce{002671}
\finexercice


\exercice{3660, quercia, 2010/03/11}
\enonce{003660}{Produits scalaires~?}

Dire si les applications suivantes sont des produits scalaires :

\begin{enumerate}
  \item $E = \R^2$, $(x,x') \mid  (y,y') = axy + bxy' + cx'y + dx'y'$
    (étudier $(1,t) \mid  (1,t),\ t \in \R$).
    

  \item $E = \R^n$, $(x_1,\dots,x_n) \mid  (y_1,\dots,y_n) =
                     a\sum\limits_i x_iy_i + b\sum\limits_{i\ne j} x_iy_j$
    $\Bigl($On montrera que $\left(\sum x_i\right)^2 \le n\sum x_i^2\Bigr)$.
    

  \item $E = \R_n[X]$, $(P \mid  Q) = \sum_{i=0}^{n} P(i)Q(i)$.
\end{enumerate}
\finenonce{003660}



\finexercice


\exercice{3661, quercia, 2010/03/11}
\enonce{003661}{Base de Schmidt}

Trouver une base orthonormée de $\R_3[X]$ pour le produit scalaire :
$(P\mid Q) =  \int_{t=-1}^1 P(t)Q(t)\,d t$.

\finenonce{003661}



\finexercice


\exercice{3662, quercia, 2010/03/11}
\enonce{003662}{Base de Schmidt}
Soit $E  =  \R_2[X]$ muni du produit scalaire :
$(P\mid Q) = \sum_{i=0}^4 P(i)Q(i)$.
Chercher une base orthonormée de $E$.
\finenonce{003662}



\finexercice


\exercice{3663, quercia, 2010/03/11}
\enonce{003663}{Inversion}

Soit $E$ un espace vectoriel euclidien. On pose pour $\vec x \ne \vec 0$ :
$i(\vec x) = \frac{\vec x}{\|\vec x\,\|^2}$.

\begin{enumerate}
  \item Montrer que $i$ est une involution et conserve les angles de vecteurs.
  \item Vérifier que : $\forall\ \vec x,\vec y \in E \setminus\{\vec 0\}$,
    $\|i(\vec x) - i(\vec y)\| = \frac{\|\vec x - \vec y\,\|}{\|\vec x\,\|\,\|\vec y\,\|}$.
    
  \item Déterminer l'image par $i$ :
  \begin{enumerate}
    \item d'un hyperplan affine ne passant pas par $\vec 0$.
        
    \item d'une sphère passant par $\vec 0$.
        
    \item d'une sphère ne passant pas par $\vec 0$.
        
    \end{enumerate}
\end{enumerate}
\finenonce{003663}



\finexercice


\exercice{3664, quercia, 2010/03/11}
\enonce{003664}{Inégalité de Ptolémée}

Soit $E$ un espace euclidien.

\begin{enumerate}
  \item Pour $\vec x \in E\setminus\{\vec 0\}$, on pose
    $f(\vec x) = \frac{\vec x}{\|\vec x\,\|^2}$.
    Montrer que : $\forall\ \vec x,\vec y \in E \setminus\{\vec 0\}$,
    $\|f(\vec x) - f(\vec y)\| = \frac{\|\vec x - \vec y\,\|}{\|\vec x\,\|\,\|\vec y\,\|}$.
    

  \item Soient $\vec a,\vec b,\vec c,\vec d \in E$.
    Montrer que $\|\vec a-\vec c\,\|\,\|\vec b-\vec d\,\| \le
         \|\vec a-\vec b\,\|\,\|\vec c-\vec d\,\| +
         \|\vec b-\vec c\,\|\,\|\vec a-\vec d\,\|$.


Indication : se ramener au cas $\vec a = \vec 0$ et utiliser l'application $f$.
\end{enumerate}
\finenonce{003664}



\finexercice


\exercice{3665, quercia, 2010/03/11}
\enonce{003665}{Calcul de distance}

On munit $E = \R_n[X]$ du produit scalaire :
Pour $P = \sum_i a_iX^i$ et $Q = \sum_i b_iX^i$,
$(P\mid Q) = \sum_i a_ib_i$.

Soit $H = \{ P \in E \text{ tq } P(1) = 0 \}$.

\begin{enumerate}
  \item Trouver une base orthonormale de $H$.
  \item Calculer $d(X,H)$.
    
\end{enumerate}
\finenonce{003665}



\finexercice


\exercice{3666, quercia, 2010/03/11}
\enonce{003666}{Expression analytique}

Soit $E$ un espace euclidien de dimension 4, ${\cal B} = (\vec e_1, \dots, \vec e_4)$
une base orthonormée de $E$, et $F$ le sous-espace vectoriel d'équations dans $\cal B$ :
$$\begin{cases} x+y+z+t = 0\cr x+2y+3z+4t =0 \cr\end{cases}$$

\begin{enumerate}
  \item Trouver une base orthonormée de $F$.
  \item Donner la matrice dans $\cal B$ de la projection orthogonale sur $F$.
  \item Calculer $d(\vec e_1, F)$.
\end{enumerate}
\finenonce{003666}



\finexercice


\exercice{3667, quercia, 2010/03/11}
\enonce{003667}{Projection sur un hyperplan}

On munit $\R^n$ du produit scalaire usuel.
Soit $H = \{ (x_1,\dots,x_n) \in \R^n \text{ tq } a_1x_1 + \dots + a_nx_n = 0 \}$
où $a_1,\dots,a_n$ sont des réels donnés non tous nuls.
Chercher la matrice dans la base canonique de la projection orthogonale
sur $H$.

\finenonce{003667}



\finexercice


\exercice{3668, quercia, 2010/03/11}
\enonce{003668}{Caractérisation des projections orthogonales}

Soit $E$ un espace vectoriel euclidien et $p \in \mathcal{L}(E)$ une projection.
Montrer que :

$p$ est une projection orthogonale 
$\Leftrightarrow \forall\ \vec x,\vec y \in E,\ (\vec x\mid p(\vec y)) = (p(\vec x)\mid \vec y)$\par
$\Leftrightarrow \forall\ \vec x \in E,\ \|p(\vec x)\| \le \|\vec x\,\|$.

(Pour la deuxième caractérisation, considérer $\vec x \in (\mathrm{Ker} p)^\perp$ et
faire un dessin)
\finenonce{003668}



\finexercice


\exercice{3669, quercia, 2010/03/11}
\enonce{003669}{Projection sur un sous-espace vectoriel de dimension finie}

Soit $E$ un espace vectoriel muni d'un produit scalaire (de dimension éventuellement infinie)
et $(\vec u_1,\dots,\vec u_n)$ une famille orthonormée de $E$.
On note $F = \text{vect}(\vec u_1,\dots,\vec u_n)$.

\begin{enumerate}
  \item Démontrer que $F \oplus F^\perp = E$ et $F^{\perp\perp} = F$
    (on utilisera la projection associée aux $\vec u_i$).

  \item Soit $\vec x \in E$. Démontrer que $\sum_{i=1}^n (\vec x\mid \vec u_i)^2 \le \|\vec x\,\|^2$.
    Quand a-t-on égalité ?


  \item Application : Soit $f : {[0,2\pi]} \to {\R}$ continue. On appelle
{\it coefficients de Fourier de $f$\/} les réels :
$$ c_k(f) =  \int_{t=0}^{2\pi} f(t)\cos(kt)\,d t   \qquad \text{et}  \qquad
   s_k(f) =  \int_{t=0}^{2\pi} f(t)\sin(kt)\,d t.$$


Démontrer l'inégalité de Bessel :
$ \int_{t=0}^{2\pi} f^2(t)\,d t \ge \frac {c_0(f)^2}{2\pi}
 + \sum_{k=1}^\infty \frac {c_k(f)^2 + s_k(f)^2}\pi$.

\end{enumerate}
\finenonce{003669}



\finexercice


\exercice{3670, quercia, 2010/03/11}
\enonce{003670}{Composition de projecteurs}

Soient $F$, $G$ deux sous-espaces vectoriels d'un espace vectoriel euclidien $E$ tels que $F^\perp \perp G^\perp$.
On note $p_F$ et $p_G$ les projections orthogonales sur $F$ et sur $G$.
Montrer que $p_F + p_G - p_{F\cap G} = \mathrm{id}_E$ et $p_F \circ p_G = p_G \circ p_F = p_{F\cap G}$.
\finenonce{003670}



\finexercice


\exercice{3671, quercia, 2010/03/11}
\enonce{003671}{Projecteurs commutant}

Soit $E$ un espace vectoriel euclidien et $p$, $q$ deux projections
orthogonales. Montrer que $p$ et $q$ commutent si et seulement si
$(\Im p\cap\Im q)^\bot \cap \Im p$ et
$(\Im p\cap\Im q)^\bot \cap \Im q$ sont orthogonaux.
\finenonce{003671}



\finexercice

\exercice{3672, quercia, 2010/03/11}
\enonce{003672}{Caractérisation des bases orthonormales}

Soit $E$ un espace vectoriel euclidien, et $\vec e_1,\dots,\vec e_n$ des vecteurs unitaires tels que :
$\forall\ \vec x \in E, \quad \|\vec x\,\|^2 = \sum_{i=1}^n (\vec x\mid \vec e_i)^2$.

\begin{enumerate}
  \item Démontrer que $(\vec e_1,\dots,\vec e_n)$ est une base orthonormale de $E$.
  \item On remplace l'hypothèse : $\vec e_i$ unitaire par : $\dim E = n$.
  \begin{enumerate}
    \item Démontrer que $(\vec e_1,\dots,\vec e_n)$ est une base de $E$.
    \item Démontrer que : $\forall\ \vec x,\vec y \in E,\ (\vec x \mid\vec y) = \sum_{i=1}^n (\vec x\mid \vec e_i)(\vec y\mid \vec e_i)$.
    \item On note $G$ la matrice de Gram de $\vec e_1,\dots,\vec e_n$. Démontrer que $G^2 = G$ et conclure.
  \end{enumerate}
\end{enumerate}
\finenonce{003672}



\finexercice


\exercice{3673, quercia, 2010/03/11}
\enonce{003673}{Matrice de Gram}

Soient $\vec x_1,\dots,\vec x_n$ des vecteurs d'un espace vectoriel euclidien $E$, et $G$ leur matrice
de Gram.

\begin{enumerate}
  \item Montrer que $\mathrm{rg} G = \mathrm{rg}(\vec x_1,\dots,\vec x_n)$.
  \item Montrer que $\det G$ est inchangé si on remplace $\vec x_k$ par
    $\vec x_k - \sum_{i\ne k} \lambda_i \vec x_i$.
  \item Soit $F = \text{vect}(\vec x_1,\dots,\vec x_n)$ et $\vec x \in E$.
    On note $d(\vec x,F) = \min(\|\vec x-\vec y\,\|, \vec y \in F)$.
    \par Montrer que $d(\vec x,F)^2 =
    \frac {\text{Gram}(\vec x_1,\dots,\vec x_n, \vec x)}{\text{Gram}(\vec x_1,\dots,\vec x_n)}$.
\end{enumerate}
\finenonce{003673}



\finexercice


\exercice{3674, quercia, 2010/03/11}
\enonce{003674}{Gram$(u(e_i))$}

Soit $E$ un espace vectoriel euclidien, $u \in \mathcal{L}(E)$ et
$(\vec e_1,\dots,\vec e_n)$ une base quelconque de $E$.
On note $G$ le déterminant de Gram.
Montrer que $G(u(\vec e_1),\dots,u(\vec e_n)) = (\det u)^2 G(\vec e_1,\dots,\vec e_n)$.
\finenonce{003674}



\finexercice


\exercice{3675, quercia, 2010/03/11}
\enonce{003675}{\'Equation du second degré}

Soient $E$ espace vectoriel euclidien, $\vec a \in E$ et $\alpha,\beta,\gamma \in \R$.
Résoudre l'équation $\alpha (\vec x\mid \vec x) + \beta (\vec x\mid \vec a) + \gamma = 0$.


\finenonce{003675}



\finexercice


\exercice{3676, quercia, 2010/03/11}
\enonce{003676}{Vecteur défini par ses produits scalaires}

Soient ${f_1,f_2,\dots,f_n} : {[0,1]} \to {\R}$ continues.

Existe-t-il $f : {[0,1]} \to {\R}$ continue telle que :
$\forall\ i,\  \int_{t=0}^1 f(t)f_i(t)\,d t = 1$ ?
\finenonce{003676}



\finexercice


\exercice{3677, quercia, 2010/03/11}
\enonce{003677}{Décomposition $QR$}

\begin{enumerate}
  \item Soit $M \in \mathcal{M}_n(\R)$ inversible. Montrer qu'il existe une matrice orthogonale, $P$,
    et une matrice triangulaire supérieure à coefficients diagonaux positifs, $T$,
    uniques telles que $M = PT$.

  \item Application : inégalité de Hadamard. Soit $E$ un espace vectoriel euclidien,
    $(\vec e_1,\dots,\vec e_n)$ une
    base orthonormée, et $\vec u_1,\dots,\vec u_n$ des vecteurs quelconques.

    Démontrer que $|\det_{(\vec e_i)}(\vec u_j)| \le \prod_j \|\vec u_j\|$.
    \'Etudier les cas d'égalité.

\end{enumerate}
\finenonce{003677}



\finexercice


\exercice{3678, quercia, 2010/03/11}
\enonce{003678}{Coefficients diagonaux dans la méthode de Schmidt}

Soit $E$ un espace euclidien, ${\cal B} = (\vec u_1,\dots,\vec u_n)$ une base de $E$
et ${\cal B}' = (\vec e_1,\dots,\vec e_n)$ la base orthonormée déduite de
$\cal B$ par la méthode de Schmidt.
On note $P$ la matrice de passage de $\cal B$ à ${\cal B}'$.

Montrer que $P_{ii} \times d(\vec u_i, \text{vect}(\vec u_1,\dots,\vec u_{i-1})) = 1$.
\finenonce{003678}



\finexercice


\exercice{3679, quercia, 2010/03/11}
\enonce{003679}{Coordonnées des vecteurs de Schmidt}

Soit $E$ un espace euclidien, ${\cal B} = (\vec u_1,\dots,\vec u_n)$ une base de $E$
et ${\cal B}' = (\vec e_1,\dots,\vec e_n)$ la base orthonormée déduite de
$\cal B$ par la méthode de Schmidt.

On note $G_n$ le déterminant de Gram de $\vec u_1,\dots,\vec u_n$,
et $\Delta_{i,n}$ le cofacteur de $(\vec u_i\mid \vec u_n)$ dans $G_n$.

Montrer que $\vec e_n = \frac1{\sqrt{G_{n-1}G_n}}\sum_{i=1}^n \Delta_{i,n}\vec u_i$.
\finenonce{003679}



\finexercice


\exercice{3680, quercia, 2010/03/11}
\enonce{003680}{$\det({}^t\!AA)$}

Soit $A \in \mathcal{M}_{n,p}(\R)$. Montrer que $\det(^tAA) \ge 0$.
\finenonce{003680}



\finexercice


\exercice{3681, quercia, 2010/03/11}
\enonce{003681}{Angles $> 2\pi/3$}

Soit $E$ un espace euclidien de dimension supérieure ou égale à 3.
Existe-t-il trois vecteurs $\vec u_1,\vec u_2,\vec u_3$ unitaires
faisant entre eux deux à deux des angles strictement supérieurs à
$\frac {2\pi}3$ ?



\finenonce{003681}



\finexercice


\exercice{3682, quercia, 2010/03/11}
\enonce{003682}{Polynômes orthogonaux}

Soit $E$ = $\R[X]$. On pose $(P \mid Q) =  \int_{t=0}^1 P(t)Q(t)\,d t$

\begin{enumerate}
  \item Démontrer que $(\ \mid\ )$ est un produit scalaire sur $E$.
  \item Démontrer qu'il existe une unique famille $(P_0, P_1, \dots, P_n,\dots )$ de polynômes
    vérifiant :
   $$\begin{cases}\deg P_i = i\cr
            \text{le coefficient dominant de $P_i$ est strictement positif}\cr
            \text{la famille $(P_i)$ est orthonormée.}\cr \end{cases}$$

\end{enumerate}
\finenonce{003682}



\finexercice


\exercice{3683, quercia, 2010/03/11}
\enonce{003683}{Centrale PSI 1997}

Soit $E = \R_n[X]$ et $(P\mid Q) =  \int_{t=0}^1 P(t)Q(t)\,d t$.

\begin{enumerate}
  \item Montrer que $E$, muni de $(\ \mid\ )$, est un espace euclidien.
  \item Soit $K = \R_{n-1}[X]^\bot$ et $P \in K\setminus\{0\}$.
    Quel est le degré de~$P$~?
  \item Soit $\Phi\ :\ x  \mapsto  \int_{t=0}^1 P(t)t^x\,d t$.
    Montrer que $\Phi$ est une fonction rationnelle.
    
  \item Trouver $\Phi$ à une constante multiplicative près.
    
  \item En déduire les coefficients de~$P$.
    
  \item En déduire une base orthogonale de~$E$.
\end{enumerate}
\finenonce{003683}



\finexercice


\exercice{3684, quercia, 2010/03/11}
\enonce{003684}{Réduction en carrés d'une forme quadratique}

Soient $f_1, \dots, f_p$ $p$ formes linéaires sur $\R^n$ telles que
$\mathrm{rg}(f_1, \dots, f_p) = n$.

En considérant le produit scalaire :
$ (\vec x \mid \vec y) = \sum_{i=1}^p f_i(\vec x)f_i(\vec y) $, démontrer qu'il existe $n$ formes
linéaires $g_1, \dots, g_n$ telles que :

$$\forall\ \vec x \in \R^n,\ \sum_{i=1}^p f_i(\vec x)^2 = \sum_{i=1}^n g_i(\vec x)^2.$$

Exemple : réduire $x^2 + (x+y)^2 + (x+2y)^2$
\finenonce{003684}



\finexercice


\exercice{3685, quercia, 2010/03/11}
\enonce{003685}{Famille de vecteurs unitaires équidistants}

Soit $E$ un espace vectoriel euclidien, et $(\vec x_1,\dots,\vec x_n)$ une famille libre.
Démontrer qu'il existe une famille $(\vec u_1,\dots,\vec u_n)$ vérifiant~:

$$\begin{cases} \vec u_i \text{ est unitaire} \cr
          \|\vec u_i-\vec u_j\| = 1 \cr
          \text{vect}(\vec u_1,\dots,\vec u_i) = \text{vect}(\vec x_1,\dots,\vec x_i).\cr\end{cases} $$


Démontrer que toute famille $(\vec u_1,\dots,\vec u_n)$ vérifiant les deux premières propriétés
est libre.
\finenonce{003685}



\finexercice


\exercice{3686, quercia, 2010/03/11}
\enonce{003686}{Famille obtusangle}

Soit $E$ un espace vectoriel euclidien et $\vec u_1,\dots,\vec u_n$ une famille de vecteurs vérifiant :
$\forall\ i\ne j,\ (\vec u_i \mid  \vec u_j) < 0$.

\begin{enumerate}
  \item On suppose $(\vec u_1,\dots,\vec u_n)$ libre. Soit $(\vec e_1,\dots,\vec e_n)$ la famille de Schmidt associée et $M$ la matrice de passage de $(\vec u_1,\dots,\vec u_n)$ à $(\vec e_1,\dots,\vec e_n)$. Montrer que $M$ est à coefficients positifs.
  \item Dans le cas général, démontrer par récurrence sur $n$ que $\mathrm{rg}(\vec u_1,\dots,\vec u_n) \ge n-1$.
  \item Si $\mathrm{rg}(\vec u_1,\dots,\vec u_n) = n-1$, démontrer que toute famille de $n-1$ vecteurs extraite
    de $(\vec u_1,\dots,\vec u_n)$ est libre, et que les composantes dans cette famille
    du vecteur retiré sont strictement négatives.
\end{enumerate}
\finenonce{003686}



\finexercice


\exercice{3687, quercia, 2010/03/11}
\enonce{003687}{$F + F^\perp \ne E$}

Soit $E = \mathcal{C}([0,1])$ muni du produit scalaire :
$(f\mid g) =  \int_{t=0}^1 fg(t)\,d t$, et $F = \{ f \in E \text{ tq } f(0) = 0 \}$.

Montrer que $F^\perp = \{0\}$.

\finenonce{003687}



\finexercice


\exercice{3688, quercia, 2010/03/11}
\enonce{003688}{Forme linéaire sur les polynômes}

On munit $\R_2[X]$ du produit scalaire : $(P\mid Q) =  \int_{t=0}^1 PQ(t)\,d t$.

\begin{enumerate}
  \item Vérifier que c'est effectivement un produit scalaire.
  \item Soit $\varphi : {\R_2[X]} \to {\R}, P \mapsto {P(0).}$
    Trouver le polynôme $A$ tel que : $\forall\ P \in \R_2[X],\
    \varphi(P) = (A\mid P)$.
    
\end{enumerate}
\finenonce{003688}



\finexercice


\exercice{3689, quercia, 2010/03/11}
\enonce{003689}{Norme uniforme sur les polynômes}

Soit $P \in \R[X]$ de degré inférieur ou égal à 3 tel que
$ \int_{t=-1}^1 P^2(t)\,d t = 1$.

Montrer que $\sup\{ |P(x)|\text{ tq } -1\le x \le 1\} \le 2\sqrt2$.

Indications : Pour $a\in\R$ montrer qu'il existe $P_a\in\R_3[X]$ tel que :
$\forall\ P\in\R_3[X],\ P(a) =  \int_{t=-1}^1 P(t)P_a(t)\,d t$.
Calculer explicitement $P_a$, et appliquer l'inégalité de Cauchy-Schwarz.

\finenonce{003689}



\finexercice


\exercice{3690, quercia, 2010/03/11}
\enonce{003690}{Centrale MP 2000}
Soit $E=\mathcal{C}^1([0,1],\R)$ et $\varphi (f,g)= \int_{[0,1]} fg+f'g'$.

\begin{enumerate}
  \item Montrer que $\varphi$ est un produit scalaire.
  \item Soit $V=\{f\in E\, |\, f(0)=f(1)=0\}$ et $W=\{f\in E \, |\, f''=f\}$. Montrer que $V$ et $W$
    sont supplémentaires orthogonaux et exprimer la projection orthogonale sur $W$.
    
  \item Soit $E_{\alpha \beta}=\{f\in E\, |\, f(0)=\alpha \text{ et } f(1)=\beta\}$. Déterminer
    $\inf\limits_{f\in E_{\alpha \beta}}  \int_{[0,1]} f^2+f'^2$.
    
\end{enumerate}
\finenonce{003690}



\finexercice


\exercice{3691, quercia, 2010/03/11}
\enonce{003691}{Polytechnique MP$^*$ 2000}
Soit~$H$ un espace euclidien, $(y_j)_{j\in I}$ une famille de vecteurs
de~$H$ telle qu'il existe $A$ et~$B$ strictement positifs vérifiant~:
$$\forall\ x\in H,\ A\|x\|^2 \le \sum_{j\in I}(x\mid y_j)^2 \le B\|x\|^2.$$

\begin{enumerate}
  \item Montrer que $(y_j)_{j\in I}$ engendre~$H$.
    
  \item On choisit $H=\R^2$. Montrer que $y_1=\left(\begin{smallmatrix}0\cr1\end{smallmatrix}\right)$,
    $y_2=\left(\begin{smallmatrix}-\sqrt3/2\cr-1/2\cr\end{smallmatrix}\right)$, $y_3=y_2$ conviennent.
    
  \item Si $A=B=1$ et $\|y_j\|=1$ pour tout~$j$, montrer que $(y_j)_{j\in I}$ est une
    base orthonormale.
    
  \item Si $A=B$, montrer que pour tout~$x\in H$, $x=\frac1A\sum_{j\in I}(x\mid y_j)y_j$.
    
\end{enumerate}
\finenonce{003691}



\finexercice


\exercice{3692, quercia, 2010/03/11}
\enonce{003692}{$\|u(x)\| \le \|x\|$}
Soit $E$ un espace euclidien et $u\in\mathcal{L}(E)$ tel que
$\forall\ x\in E,\ \|u(x)\| \le \|x\|$. Montrer que
$E = \mathrm{Ker}(u-\mathrm{id})\mathop\oplus\limits^\perp\Im(u-\mathrm{id})$.


\finenonce{003692}



\finexercice


\exercice{3693, quercia, 2010/03/11}
\enonce{003693}{X MP$^*$ 2000}

Soit $E$ un espace euclidien de dimension $n>1$. Trouver toutes les fonctions $f$
de $E$ dans $\R$ continues telles que $u\perp v \Rightarrow f(u+v)=f(u)+f(v)$.
\finenonce{003693}



\finexercice


\exercice{5482, rouget, 2010/07/10}
\enonce{005482}{***}
Pour $A=(a_{i,j})_{1\leq i,j\leq n}\in\mathcal{M}_n(\Rr)$, $N(A)=\mbox{Tr}(^tAA)$. Montrer que $N$ est une norme vérifiant de plus $N(AB)\leq N(A)N(B)$ pour toutes matrices carrées $A$ et $B$. $N$ est-elle associée à un produit scalaire~?
\finenonce{005482}


\finexercice


\exercice{5483, rouget, 2010/07/10}
\enonce{005483}{***}
Soit $E$ un $\Rr$ espace vectoriel de dimension finie. Soit $||\;||$ une norme sur $E$ vérifiant l'identité du parallèlogramme, c'est-à-dire~:~$\forall(x,y)\in E^2,\;||x+y||^2+||x-y||^2=2(||x||^2+||y||^2)$. On se propose de démontrer que $||\;||$ est associée à un produit scalaire.
On définit sur $E^2$ une application $f$ par~:~$\forall(x,y)\in E^2,\;f(x,y)=\frac{1}{4}(||x+y||^2-||x-y||^2)$.
\begin{enumerate}
\item  Montrer que pour tout $(x,y,z)$ de $E^3$, on a~:~$f(x+z,y)+f(x-z,y)=2f(x,y)$.
\item  Montrer que pour tout $(x,y)$ de $E^2$, on a~:~$f(2x,y)=2f(x,y)$.
\item  Montrer que pour tout $(x,y)$ de $E^2$ et tout rationnel $r$, on a~:~$f(rx,y)=rf(x,y)$.

On admettra que pour tout réel $\lambda$ et tout $(x,y)$ de $E^2$ on a~:~$f(\lambda x,y)=\lambda f(x,y)$ ( ce résultat provient de la continuité de $f$).
\item  Montrer que pour tout $(u,v,w)$ de $E^3$, $f(u,w)+f(v,w)=f(u+v,w)$.
\item  Montrer que $f$ est bilinéaire.
\item  Montrer que $||\;||$ est une norme euclidienne.
\end{enumerate}

\finenonce{005483}


\finexercice


\exercice{5485, rouget, 2010/07/10}
\enonce{005485}{**}
Sur $\Rr[X]$, on pose $P|Q=\int_{0}^{1}P(t)Q(t)\;dt$. Existe-t-il $A$ élément de $\Rr[X]$ tel que $\forall P\in\Rr[X],\;P|A=P(0)$~?
\finenonce{005485}


\finexercice


\exercice{5495, rouget, 2010/07/10}
\enonce{005495}{***}
Soit $(e_1,...,e_n)$ une base quelconque de $E$ euclidien. Soient $a_1$,..., $a_n$ $n$ réels donnés.
Montrer qu'il existe un unique vecteur $x$ tel que $\forall i\in\{1,...,\},\;x|e_i=a_i$.
\finenonce{005495}


\finexercice


\exercice{5496, rouget, 2010/07/10}
\enonce{005496}{****}
Soit $E$ un espace vectoriel euclidien de dimension $n\geq1$.
Une famille de $p$ vecteurs $(x_1,...,x_p)$ est dite obtusangle si et seulement si pour tout $(i,j)$ tel que $i\neq j$, $x_i|x_j<0$. Montrer que l'on a nécessairement $p\leq n+1$.
\finenonce{005496}


\finexercice


\exercice{5629, rouget, 2010/10/16}
\enonce{005629}{**I}
\begin{enumerate}
 \item  Soient $n\in\Nn^*$ et $E=\Cc_n[X]$. Pour $a\in\Cc$, on définit l'application $\varphi_a$ par : $\forall P\in E,\;\varphi_a(P)=P(a)$. Montrer que pour tout $a\in E$, $\varphi_a\in E^*$.

\item  Soient $a_0$, $a_1$,\ldots, $a_n$ $n+1$ nombres complexes deux à deux distincts. Montrer que la famille $\left(\varphi_{a_0},\ldots,\varphi_{a_n}\right)$ est une base de $E^*$ et déterminer sa préduale.

\item  Montrer qu'il existe $(\lambda_0,\ldots,\lambda_n)\in\Cc^{n+1}$ tel que $\forall P\in\Cc_n[X],\;\int_{0}^{1}P(t)\;dt=\lambda_0P(a_0)+\ldots+\lambda_nP(a_n)$ puis donner la valeur des $\lambda_i$ sous la forme d'une intégrale.
\end{enumerate}
\finenonce{005629}


\finexercice


\exercice{5630, rouget, 2010/10/16}
\enonce{005630}{**}
Sur $E =\Rr_3[X]$, on pose pour tout $P\in E$, $\varphi_1(P)=P(0)$  et $\varphi_2(P) = P(1)$ puis $\psi_1(P) =P'(0)$ et $\psi_2(P) = P'(1)$. Montrer que $(\varphi_1,\varphi_2,\psi_1,\psi_2)$ est une base de $E^*$ et trouver la base dont elle est la duale.
\finenonce{005630}


\finexercice


\exercice{5631, rouget, 2010/10/16}
\enonce{005631}{**}
Soit $E$ un $\Kk$-espace vectoriel et $\varphi$ et $\psi$ deux formes linéaires sur $E$. On suppose que pour tout $x$ de $E$, on a  $\varphi(x)\psi(x) = 0$. Montrer que $\varphi= 0$ ou $\psi= 0$.
\finenonce{005631}


\finexercice


\exercice{5632, rouget, 2010/10/16}
\enonce{005632}{***}
\begin{enumerate}
 \item  Soient $n\in\Nn^*$ puis $\varphi_1$,..., $\varphi_n$ et $\varphi$ $n+1$ formes linéaires sur un $\Kk$-espace vectoriel $E$ de dimension finie.

Montrer que : $\left(\exists(\lambda_1,...,\lambda_n)\in\Kk^n/\;\varphi =\lambda_1\varphi_1+...+\lambda_n\varphi_n\Leftrightarrow\displaystyle\bigcap_{i=1}^{n}\text{Ker}\varphi_i\subset\text{Ker}\varphi\right)$.

\item  Signification du résultat précédent : dans $\Rr^3$, équation d'un plan $P$ contenant $D~:~\left\{
\begin{array}{l}
x+y+z=0\\
2x+3z=0
\end{array}
\right.$ et le vecteur $u=(1,1,1)$ ?
\end{enumerate}
\finenonce{005632}


\finexercice


\exercice{5633, rouget, 2010/10/16}
\enonce{005633}{***}
Soient $n\in\Nn^*$ puis $\varphi_1$,..., $\varphi_n$ $n$ formes linéaires sur un $\Kk$-espace $E$ de dimension $n$.

Montrer que la famille $(\varphi_1,...,\varphi_n)$ est liée si et seulement si il existe un vecteur $x$ non nul tel que $\forall i\in\llbracket1,n\rrbracket,\;\varphi_i(x) = 0$.
\finenonce{005633}


\finexercice


\exercice{5634, rouget, 2010/10/16}
\enonce{005634}{**}
Rang du système de formes linéaires sur $\Rr^4$ 

\begin{center}
$\begin{array}{l}
f_1 = x_1 + 2x_2 - x_3 - 2x_4\\
f_2 = x_1 + x_2 + mx_3 + x_4\\
f_3 = x_1+ x_3 + (m+4)x_4\\
f_4=x_2 - 3x_3 - mx_4
\end{array}
$ ?
\end{center}
\finenonce{005634}


\finexercice


\exercice{5773, rouget, 2010/10/16}
\enonce{005773}{*** I}
Soit $E=\Rr[X]$. Pour $(P,Q)\in E^2$, on pose $\varphi(P,Q)=\int_{0}^{+\infty}P(t)Q(t)e^{-t}\;dt$. Pour $n\in\Nn$, on pose $h_n=(X^ne^{-X})^{(n)}e^X$.

\begin{enumerate}
 \item  Montrer que $\varphi$ est un produit scalaire sur $E$.

\item 
  \begin{enumerate}
   \item  Pour $n\in\Nn$, préciser les coefficients de $h_n$. Montrer que la famille $(h_n)_{n\in\Nn}$ est une base de $E$.
   \item Montrer que la famille $(h_n)_{n\in\Nn}$ est une base orthogonale de l'espace préhilbertien $(E,\varphi)$.
   \item Pour $n\in\Nn$, déterminer $\|h_n\|$. En déduire une base orthonormée de l'espace préhilbertien $(E,\varphi)$.
  \end{enumerate}
\end{enumerate}
\finenonce{005773}


\finexercice


\exercice{5774, rouget, 2010/10/16}
\enonce{005774}{** I Polynômes de \textsc{Tchebychev}}
Soit $E=\Rr[X]$. Pour $(P,Q)\in E^2$, on pose $\varphi(P,Q)=\int_{-1}^{1}\frac{P(t)Q(t)}{\sqrt{1-t^2}}\;dt$. Pour $n\in\Nn$, on note $T_n$ le $n$-ème polynôme de \textsc{Tchebychev} de première espèce c'est-à-dire l'unique polynôme tel que $\forall\theta\in\Rr$, $T_n(\cos\theta)=\cos(n\theta)$.

\begin{enumerate}
 \item  Montrer que $\varphi$ est un produit scalaire sur $E$.

\item 
  \begin{enumerate}
  \item Montrer que $(T_n)_{n\in\Nn}$ est une base orthogonale de l'espace préhilbertien $(E,\varphi)$.
  \item Pour $n\in\Nn$, déterminer $\|T_n\|$.
  \end{enumerate}
\end{enumerate}
\finenonce{005774}


\finexercice


\exercice{5775, rouget, 2010/10/16}
\enonce{005775}{** I}
On note $E$ l'ensemble des suites réelles de carrés sommables c'est-dire les suites réelles $(u_n)_{n\in\Nn}$ telles que

\begin{center}
$\sum_{n=0}^{+\infty}u_n^2<+\infty$.
\end{center}

\begin{enumerate}
 \item  Montrer que $E$ est un $\Rr$-espace vectoriel.

\item  Pour $(u,v)\in E^2$, on pose $\varphi(u,v)=\sum_{n=0}^{+\infty}u_nv_n$. Montrer que $\varphi$ est un produit scalaire sur $E$.
\end{enumerate}
\finenonce{005775}


\finexercice


\exercice{5776, rouget, 2010/10/16}
\enonce{005776}{* I}
 Soit $\Phi$ l'application qui à deux matrices carrées réelles $A$ et $B$ de format $n$ associe $\text{Tr}({^t}A\times B)$. Montrer que $\Phi$ est un produit scalaire sur $\mathcal{M}_n(\Rr)$. Est ce que $\Phi$ est un produit scalaire sur $\mathcal{M}_n(\Cc)$ ?
\finenonce{005776}


\finexercice


\exercice{5777, rouget, 2010/10/16}
\enonce{005777}{****}
 Soit $E$ un $\Rr$-espace vectoriel muni d'une norme, notée $\|\;\|$, vérifiant l'identité du parallélogramme. Montrer que cette norme est hilbertienne.
\finenonce{005777}


\finexercice


\exercice{5778, rouget, 2010/10/16}
\enonce{005778}{**}
Soit $E$ un espace préhilbertien réel et $(e_1,...,e_n)$ une famille de $n$ vecteurs unitaires de $E$ ($n\in\Nn^*$) telle que pour tout vecteur $x$ de $E$, on ait  $\|x\|^2 =\sum_{k=1}^{n}\left(x|e_k\right)^2$. Montrer que la famille $(e_1,...,e_n)$ est une base orthonormée de $E$. 

\finenonce{005778}


\finexercice


\exercice{5788, rouget, 2010/10/16}
\enonce{005788}{**** I}
Soit $E$ un espace euclidien de dimension $n$ non nulle. Soit $(x_1,...,x_p)$ une famille de $p$ vecteurs de $E$ ($p\geqslant 2$) . On dit que la famille $(x_1,...,x_p)$ est une famille obtusangle si et seulement si $\forall(i,j)\in\llbracket1,p\rrbracket^2$ $(i<j\Rightarrow x_i|x_j < 0)$. Montrer que si la famille $(x_1,...,x_p)$ est une famille obtusangle alors $p\leqslant n+1$.
\finenonce{005788}


\finexercice

\section{204.02 Forme quadratique}


\exercice{1509, ortiz, 1999/04/01}

\enonce{001509}{}
Soient $E$ un $K$-espace vectoriel (o\`u $K$ est
$\Rr$ ou $\Cc$) de dimension finie $n>0$ et $q$
une forme quadratique sur $E$.
\begin{enumerate}
    \item  $q$ peut-elle \^etre injective ?
    \item  Trouver une condition n\'ecessaire et suffisante sur $q$ pour qu'elle soit surjective.
\end{enumerate}
\finenonce{001509}



\finexercice



\exercice{1510, bodin, 1999/11/01}

\enonce{001510}{examen juin 1999}
Soit $a$ un nombre r\'eel.
Soit $q$ la forme quadratique d\'efinie sur $\Rr^3$ par
$$q(v)=x^2+(1+a)y^2+(1+a+a^2)z^2+2xy-2ayz$$
pour $v=(x,y,z)$.
Soit $f$ la forme bilin\'eaire sym\'etrique associ\'ee \`a $q$.
\begin{enumerate}
  \item D\'eterminer une d\'ecomposition de $q$ en combinaison lin\'eaire de carr\'es
de formes lin\'eaires ind\'ependantes.
  \item Donner le rang et la signature de $q$ suivant les valeurs
de $a$.
  \item Pour quelles valeurs de $a$, $f$ d\'efinit-elle un produit
scalaire ?
\end{enumerate}
\finenonce{001510}



\finexercice



\exercice{1511, ortiz, 1999/04/01}

\enonce{001511}{}
Soit $q$ la forme quadratique de $\Rr^3$ de
matrice $A=
\left(\begin{smallmatrix}
2 & 1 & 1 \\
1 & 1 & 1 \\
1 & 1 & 2
\end{smallmatrix}\right)
$ dans la base canonique
$\mathcal{B}=(e_1,e_2,e_3)$ de $\Rr^3.$
\begin{enumerate}
    \item  Donner l'expression analytique de $q$ dans $\mathcal{B}$ et
expliciter sa forme polaire $f$.
    \item  V\'erifier que $\mathcal{B}^{^{\prime }}=(e_1,-\frac
12e_1+e_2,-e_2+e_3)$ est une base $\Rr^3$ et
donner la matrice de $q$ dans cette base.
Expliciter $q$ dans cette base.
    \item  Trouver le rang et la signature de $q$.
\end{enumerate}
\finenonce{001511}



\finexercice



\exercice{1512, ortiz, 1999/04/01}

\enonce{001512}{}
Soient $E=\Rr_2\left[ X\right] $ et $q$
l'application de $E$ dans $\Rr$ d\'efinie par
$q(P)=P(0)P(1).$
\begin{enumerate}
    \item
    \begin{enumerate}
         \item Montrer que $q$ est une forme quadratique sur $E$.
         \item D\'eterminer la matrice de $q$ dans la base canonique de $E$.
         \item La forme $q$ est-elle positive, n\'egative ?
    \end{enumerate}
    \item   Soit $P:=X^2+X+1$ et $V=$vect($P$). D\'eterminer $V^{\bot }$ et $V^{\bot
\bot }.$
    \item  D\'eterminer le rang de $q$ puis son noyau.
    \item  D\'eterminer le c\^one isotrope $C(q)$ de $q$ et constuire
une base de $E$ form\'ee de vecteurs isotropes.
$C(q)$ est-il un sous-espace vectoriel de $E$ ?
    \item  D\'eterminer une base $(P_0,P_1,P_2)$ de $E$ telle que
$q(a_0P_0+a_1P_1+a_2P_2)=a_0^2-a_1^2$ et donner la
signature de $q.$
\end{enumerate}
\finenonce{001512}



\finexercice



\exercice{1513, ortiz, 1999/04/01}

\enonce{001513}{}
Soit $q$ une forme quadratique sur un $\Rr$-espace
vectoriel $E$, que l'on suppose d\'efinie (\emph{i.e.} son
c\^one isotrope est $\left\{0\right\}$). Montrer que $q$
garde un signe constant sur $E$ (on pourra
raisonner par l'absurde et consid\'erer $q(a+tb)$
o\`u $a$ et $b$ sont des vecteurs bien choisis et
$t\in\Rr$).
\finenonce{001513}



\finexercice



\exercice{1514, ortiz, 1999/04/01}

\enonce{001514}{}
\begin{enumerate}
    \item  Diagonaliser $A=
\left(\begin{smallmatrix} 11 & -5 & 5 \\
-5 &3 & -3 \\
5 & -3 & 3
\end{smallmatrix}\right)  .$
    \item  Soit $q$ la forme quadratique de $\Rr^3$ de matrice $A$ dans la base
canonique de $\Rr^3.$ Utiliser la question pr\'ec\'edente pour trouver
une base $q$-orthogonale, d\'eterminer la signature de $q$ et une d\'ecomposition
 de $q$ en combinaison lin\'eaire de carr\'es de formes lin\'eaires ind\'ependantes.
\end{enumerate}
\finenonce{001514}



\finexercice



\exercice{1515, ortiz, 1999/04/01}

\enonce{001515}{}
D\'eterminer la signature de la forme quadratique
$${q:\left( x,y,z\right) \in \Rr^3\mapsto
\left( 2x+y-z\right) ^2-\left( 3x-y+2z\right)
^2+\left( 5y-7z\right) ^2.}$$
\finenonce{001515}



\finexercice



\exercice{1516, ortiz, 1999/04/01}

\enonce{001516}{}
Soit la forme quadratique $q$ d\'efinie par
$${
q:\left( x_1,x_2,x_3,x_4\right) \in
\Cc^4\mapsto
x_1x_2+x_2x_4-x_3x_4-2x_1x_4-2x_2x_3-x_1x_3}.$$
\begin{enumerate}
    \item  Montrer, sans r\'eduire $q$, qu'il existe
une base $q$-orthonormale de $\Cc^4.$
    \item  En expliciter une.
\end{enumerate}
\finenonce{001516}



\finexercice



\exercice{3713, quercia, 2010/03/11}
\enonce{003713}{\'Etude de signe}

Déterminer si les formes quadratiques suivantes sont positives :

\begin{enumerate}
  \item $q(x,y) = (1-\lambda)x^2 + 2\mu xy + (1+\lambda)y^2$.
    

  \item $q(x,y,z) = x^2 + y^2 + 2z(x\cos\alpha + y\sin \alpha)$.
    

  \item $q(x,y,z,t) = x^2 + 3y^2 + 4z^2 + t^2 + 2xy + xt$.
    
\end{enumerate}
\finenonce{003713}



\finexercice


\exercice{3714, quercia, 2010/03/11}
\enonce{003714}{Ensi PC 1999}

La matrice $A = \begin{pmatrix}4&1&1\cr1&4&1\cr1&1&4\cr\end{pmatrix}$
est-elle une matrice de produit scalaire~?
\finenonce{003714}



\finexercice


\exercice{3715, quercia, 2010/03/11}
\enonce{003715}{Calcul de signature}

Soit $A\in\mathcal{M}_n(\R)$ à coefficients strictement positifs.
Déterminer la signature de la forme quadratique sur~$\R^n$ définie par~:
$q(x_1,\dots,x_n) = \sum_{i,j} a_{i,j}(x_i-x_j)^2$.

\finenonce{003715}



\finexercice


\exercice{3716, quercia, 2010/03/11}
\enonce{003716}{Signature de ${}^t\!AA$}

Soit $A \in \mathcal{M}_{n,p}(\R)$.

\begin{enumerate}
  \item Montrer que ${}^t\!AA$ est la matrice d'une forme quadratique positive sur
    $\R^p$.
  \item Déterminer sa signature en fonction de~$\mathrm{rg} A$.

\end{enumerate}
\finenonce{003716}



\finexercice


\exercice{3717, quercia, 2010/03/11}
\enonce{003717}{Décomposition en carrés}

Décomposer en carrés la forme quadratique définie sur~$\R^n$ par~:
$q(x_1,\dots,x_n) = \sum_{1\le i\le j\le n} x_ix_j
                  = \frac12\sum_{i\ge 1} x_i^2 + \frac12\Bigl(\sum_{i\ge1} x_i\Bigr)^2$.
On posera $y_i = x_i + (x_{i+1} + \dots + x_n)/(i+1)$.


\finenonce{003717}



\finexercice


\exercice{3718, quercia, 2010/03/11}
\enonce{003718}{Rang d'une décomposition en carrés}

Soit $q$ une forme quadratique sur un espace vectoriel $E$ de dimension finie et
$f_1,\dots,f_p\in E^*$, $\alpha_1,\dots,\alpha_p\in \R$ tels
que ${q = \alpha_1f_1^2 + \dots + \alpha_pf_p^2}$.
Montrer que $\mathrm{rg}(f_1,\dots,f_p) \ge \mathrm{rg}(q)$.
\finenonce{003718}



\finexercice


\exercice{3719, quercia, 2010/03/11}
\enonce{003719}{Différentielle d'une forme quadratique}

Soit $q$ une forme quadratique sur $\R^n$ et $f$ la forme bilinéaire
symétrique associée.

Montrer que :
$\forall\ \vec x,\vec y \in \R^n,\ dq_{\vec x}(\vec y) = 2f(\vec x,\vec y)$.
\finenonce{003719}



\finexercice


\exercice{3720, quercia, 2010/03/11}
\enonce{003720}{$q(a)q(x)-f^2(a,x)$}

Soit $f$ une forme bilinéaire symétrique sur~$E$ et $q$ la forme quadratique
associée.\par On pose pour~$x\in\ E$~: $\varphi(x) = q(a)q(x)-f^2(a,x)$.

\begin{enumerate}
  \item Montrer que $\varphi$ est une forme quadratique sur~$E$.
  \item Si $E$ est de dimension finie comparer les rangs de $\varphi$ et~$q$.
  \item Dans le cas général, déterminer le noyau de la forme polaire de~$\varphi$ en
    fonction de celui de~$f$ et de~$a$.
\end{enumerate}
\finenonce{003720}



\finexercice


\exercice{3721, quercia, 2010/03/11}
\enonce{003721}{$\mathrm{tr}(A^2)$}
Soit pour $A\in\mathcal{M}_n(\R)$~: $q(A) = \mathrm{tr}(A^2)$. Montrer que $q$ est une
forme quadratique sur~$\mathcal{M}_n(\R)$ et déterminer sa signature
(indication~: étudier les restrictions de~$q$ aux sous-espaces vectoriels des matrices
symétriques et antisymétriques).
\finenonce{003721}



\finexercice


\exercice{3722, quercia, 2010/03/11}
\enonce{003722}{Adjoint}

Soit $E$ un $\R$-espace vectoriel de dimension finie et $f$ une forme bilinéaire symétrique
non dégénérée sur~$E$.

\begin{enumerate}
  \item Si $u\in\mathcal{L}(E)$ montrer qu'il existe un unique endomorphisme $v\in\mathcal{L}(E)$
    tel que~: $$\forall\ x,y\in\ E,\ f(u(x),y) = f(x,v(y)).$$
    On note $v=u^*$.

  \item Montrer que l'application $u \mapsto u^*$ est un anti-isomorphisme involutif
    de l'algèbre~$\mathcal{L}(E)$ (c'est-à-dire un isomorphisme linéaire
    tel que $(u\circ v)^* = v^*\circ u^*$ et $u^{**}=u$).
\end{enumerate}
\finenonce{003722}



\finexercice


\exercice{3723, quercia, 2010/03/11}
\enonce{003723}{Restriction d'une forme quadratique à un sous-espace vectoriel}

Soit $E$ un $\R$-espace vectoriel de dimension finie $n\ge 1$ et $q$ une forme
quadratique sur~$E$ de signature $(n-1,1)$.
Soit~$H$ un sous-espace vectoriel de~$E$ de dimension~$d\ge 1$.

\begin{enumerate}
  \item On suppose qu'il
existe $x\in\ H$ tel que~$q(x) < 0$. Montrer que la signature
de $q_{|H}$ est $(d-1,1)$.

  \item On suppose que $q_{|H}$ est positive, quelle est sa signature~?

\end{enumerate}
\finenonce{003723}



\finexercice


\exercice{3724, quercia, 2010/03/11}
\enonce{003724}{Mineurs principaux\label{mineursppx}}

Soit $n\ge 2$ et $A$ une matrice réelle symétrique $n\times n$
représentant une forme quadratique~$q$.
On appelle mineurs principaux de~$A$ les déterminants~:
$$\Delta_k(A) = \det\bigl((a_{i,j})_{i,j\le k}\bigr).$$
On suppose que tous les mineurs principaux de~$A$ sont non nuls,
montrer que la signature de~$q$ est $(r,s)$ où $s$ est le nombre
de changements de signe dans la suite $(1,\Delta_1,\dots,\Delta_n)$
et $r=n-s$.


\finenonce{003724}



\finexercice


\exercice{3725, quercia, 2010/03/11}
\enonce{003725}{Diagonale dominante}

Soit $A = (a_{ij})\in\mathcal{M}_n(\R)$ symétrique. On dit que $A$ est à diagonale
faiblement dominante si pour tout $i$ on a $a_{ii}\ge \sum_{j\ne i}|a_{ij}|$
et que $A$ est à diagonale fortement dominante si pour tout $i$
on a $a_{ii}> \sum_{j\ne i}|a_{ij}|$.

Montrer que si $A$ est à diagonale fortement dominante alors $A$ est définie positive
et si $A$ est à diagonale faiblement dominante alors $A$ est positive.


\finenonce{003725}



\finexercice


\exercice{3726, quercia, 2010/03/11}
\enonce{003726}{Formes quadratiques de signature donnée}

Soit $E$ un $\R$-espace vectoriel de dimension~$n$, on note~:

\halign{\indent\indent#\hfil\quad&#\hfil\cr
Quad$(E)$        &l'ensemble des formes quadratiques sur~$E$~;\cr
Quad$^*(E)$      &l'ensemble des formes quadratiques sur~$E$ de rang~$n$~;\cr
Quad$_{p,q}(E)$  &l'ensemble des formes quadratiques sur~$E$ de signature $(p,q)$.\cr
}

\begin{enumerate}
  \item Montrer que $\mathrm{Quad}^*(E)$ est dense dans $\mathrm{Quad}(E)$.
  \item Montrer que si $p+q=n$ alors $\mathrm{Quad}_{p,q}(E)$ est ouvert dans $\mathrm{Quad}(E)$.
  \item Montrer que $\mathrm{Quad}_{p,q}(E)$ est connexe par arcs.

\end{enumerate}
\finenonce{003726}



\finexercice


\exercice{3727, quercia, 2010/03/11}
\enonce{003727}{$1/(\lambda_i+\lambda_j)$}

\begin{enumerate}
  \item Soit ${f_1,\dots,f_n} : I \to \R$ des fonctions continues de carrés
intégrables sur l'intervalle~$I$.
On pose $a_{i,j} =  \int_I f_if_j$. Montrer que la matrice $(a_{i,j})$
est définie positive ssi la famille $(f_1,\dots,f_n)$ est libre.

  \item En déduire que si $\lambda_1$,\dots,$\lambda_n$ sont des réels
strictement positifs distincts alors la matrice de terme général
$1/(\lambda_i+\lambda_j)$ est définie positive.
\end{enumerate}
\finenonce{003727}



\finexercice


\exercice{3728, quercia, 2010/03/11}
\enonce{003728}{Matrice des inverses}

Soit $A = (a_{i,j})\in \mathcal{M}_n(\R)$ à coefficients tous non nuls.
On note $A'$ la matrice de coefficient général $1/a_{i,j}$.

\begin{enumerate}
  \item Trouver les matrices $A$ telles que $A$ et $A'$ sont symétriques définies positives
    (examiner les cas $n=1$, $n=2$, $n=n$).


  \item Trouver les matrices $A$ telles que $A$ et $A'$ sont symétriques positives
    (examiner les cas $n=2$, $n=3$, $n=n$).

\end{enumerate}
\finenonce{003728}



\finexercice


\exercice{3729, quercia, 2010/03/11}
\enonce{003729}{Centrale MP 2000}
Soit $S$ une matrice carrée d'ordre~$n$, à coefficients réels, symétrique
définie positive. Soit~$X=\begin{pmatrix}x_1\cr\vdots\cr x_n\cr\end{pmatrix}$.
Montrer que $q$ : $X \mapsto\begin{vmatrix}0&x_1&\dots&x_n\cr x_1\cr\vdots&&S\cr x_n\cr\end{vmatrix}$
est définie négative.

\finenonce{003729}



\finexercice


\exercice{3730, quercia, 2010/03/11}
\enonce{003730}{Polytechnique MP$^*$ 2000}
On considère sur $\R^n$ la forme quadratique : $q(x)=\alpha \| x\| ^2 +\beta (x|a)^2$ où
$\alpha >0$, $\beta $ réel et $a\in \R^n$.
Discuter de la signature et du rang de $q$.

\finenonce{003730}



\finexercice


\exercice{3731, quercia, 2010/03/11}
\enonce{003731}{Ensae MP$^*$ 2000}
Soit $q$ une forme quadratique non nulle sur $M_2(\C)$ telle que $q(AB)=q(A) q(B)$. Déterminer $q$.

\finenonce{003731}



\finexercice


\exercice{3732, quercia, 2010/03/11}
\enonce{003732}{Déterminant de Gram, X MP$^*$ 2005}

Soit $E$ un espace vectoriel réel de dimension quelconque, $(x_{1},\ldots,x_{n})$ et $(y_{1},\ldots,y_{n})$ deux familles
de vecteurs de $E$ et $\phi$ une forme bilinéaire symétrique positive.

Montrer que $(\det [\phi(x_{i}, y_{j})])^{2}\le \det [\phi(x_{i}, x_{j})]\times \det [\phi(y_{i}, y_{j})]$.

\finenonce{003732}



\finexercice


\exercice{5807, rouget, 2010/10/16}
\enonce{005807}{**}
Soit $E =\mathcal{L}(\Rr^2)$. Soit $(\lambda,\mu)\in\Rr^2$. On pose

\begin{center}
$\forall f\in\mathcal{L}(\Rr^2)$, $Q(f) =\lambda\text{Tr}(f^2) +\mu\text{det}(f)$.
\end{center}

\begin{enumerate}
 \item  Vérifier que $Q$ est une forme quadratique sur $E$.

\item  Déterminer en fonction de $\lambda$ et $\mu$ le rang et la signature de $Q$.
Analyser en particulier les cas $(\lambda,\mu) = (1,0)$ et $(\lambda,\mu) = (0,1)$.
\end{enumerate}
\finenonce{005807}


\finexercice


\exercice{5808, rouget, 2010/10/16}
\enonce{005808}{**}
Soit $Q$ une forme quadratique sur un $\Rr$-espace vectoriel $E$. On note $\varphi$ sa forme polaire.

On suppose que $\varphi$ est non dégénérée mais non définie. Montrer que $Q$ n'est pas de signe constant.
\finenonce{005808}


\finexercice


\exercice{5809, rouget, 2010/10/16}
\enonce{005809}{*** I}
Soient $f_1$, $f_2$,..., $f_n$ $n$ fonctions continues sur $[a,b]$ à valeurs dans $\Rr$. Pour $(i,j)\in\llbracket1,n\rrbracket^2$, on pose $b_{i,j}=\int_{a}^{b}f_i(t)f_j(t)\;dt$ puis pour $(x_1,...x_n)\in\Rr^n$, $Q((x_1,...,x_n)) =\sum_{1\leqslant i,j\leqslant n}^{}b_{i,j}x_ix_j$.

\begin{enumerate}
 \item  Montrer que $Q$ est une forme quadratique positive.

\item  Montrer que $Q$ est définie positive si et seulement si la famille $(f_1,...,f_n)$ est libre.

\item  Ecrire la matrice de $Q$ dans la base canonique de $\Rr^n$ dans le cas particulier : $\forall i\in\llbracket1,n\rrbracket$, $\forall t\in[a,b]$, $f_i(t)=t^{i-1}$.
\end{enumerate}
\finenonce{005809}


\finexercice


\exercice{5810, rouget, 2010/10/16}
\enonce{005810}{***}
Soit $S$ une matrice symétrique réelle, définie positive. Pour $(x_1,...,x_n)\in\Rr^n$, on pose 

\begin{center}
$Q((x_1,...,x_n)) = -\text{det}\left(\begin{array}{cccc}
0&x_1&\ldots&x_n\\
x_1& & & \\
\vdots& &S& \\
x_n& & & 
\end{array}
\right)$.
\end{center}

Montrer que $Q$ est une forme quadratique définie positive.
\finenonce{005810}


\finexercice

\section{204.03 Espace orthogonal}


\exercice{1475, legall, 1998/09/01}

\enonce{001475}{}
Montrer que l'application $  (A, B)\mapsto tr({}^t\! AB)  $ de $M_2(\R)\times M_2(\R)$ \`a valeurs dans $  \R   $ est un produit scalaire. Calculer l'orthogonal de l'ensemble des
matrices diagonales puis celui des matrices sym\'etriques.
\finenonce{001475}



\finexercice



\exercice{1476, legall, 1998/09/01}

\enonce{001476}{}
Soit $  (E , \langle   ,  \rangle )  $ un espace euclidien, $  F  $
et $  G  $ deux sous-espaces vectoriels de $  E  .$ Montrer que~:
\begin{enumerate}
    \item Si $  F\subset G  $ alors $  G^\perp \subset F^\perp   .$
    \item $  (F+G)^\perp = F^\perp \cap G ^\perp   .$
    \item $  (F\cap G)^\perp = F^\perp +G ^\perp   .$
    \item Si $\hbox{dim}(E)  $ est finie, alors $  (F^\perp )^\perp =F  .$
\end{enumerate}
\finenonce{001476}



\finexercice



\exercice{5792, rouget, 2010/10/16}
\enonce{005792}{**}
On munit $E =\mathcal{M}_3(\Rr)$ muni du produit scalaire usuel.

\begin{enumerate}
 \item  Déterminer l'orthogonal de $\mathcal{A}_3(\Rr)$.

\item  Calculer la distance de la matrice $M =\left(
\begin{array}{ccc}
0&1&0\\
0&0&1\\
0&0&0
\end{array}
\right)$  au sous-espace vectoriel des matrices antisymétriques.
\end{enumerate}
\finenonce{005792}


\finexercice

\section{204.04 Projection, symétrie}


\exercice{1477, barraud, 2003/09/01}

\enonce{001477}{}
  Déterminer la matrice dans la base canonique de $\R^{3}$ de la projection
  orthogonale sur le plan d'équation $x+2y-3z=0$.

  En déduire la matrice de la symétrie orthogonale par rapport à ce plan.

  \medskip

  Dans un espace euclidien de dimension $n$, on considére un sous-espace
  $F$ de dimension $r$ et $(f_{1},...,f_{r})$ une base de orthonormée de
  cet espace. On not $p_{F}$ la projection orthogonale sur $F$, c'est à
  dire la projection sur $F$ associée à la décomposition $E=F\oplus
  F^{\bot}$. Montrer que :
  $$
  \forall v\in F, \qquad p_{F}(v)=<v,f_{1}>f_{1}
                                 +<v,f_{2}>f_{2}
                                 +\cdots
                                 +<v,f_{r}>f_{r}
  $$ 
\finenonce{001477}



\finexercice



\exercice{1478, barraud, 2003/09/01}

\enonce{001478}{}
  Dans $\R^{3}$ muni de son produit scalaire canonique, déterminer la
  projection orthogonale sur le plan d'équation $x+y+z=0$ de $(1,0,0)$,
  et plus généralement d'un vecteur $(x,y,z)$ quelconque.

  Donner la  matrice de cette projection ainsi que celle de la symétrie
  orthogonale par rapport à ce plan.
  \medskip

  Dans un espace euclidien de dimension $n$, on considère un sous-espace
  $F$ de dimension $r$ et $(f_{1},...,f_{r})$ une base de orthonormée de
  cet espace. On not $p_{F}$ la projection orthogonale sur $F$, c'est à
  dire la projection sur $F$ associée à la décomposition $E=F\oplus
  F^{\bot}$. Montrer que :
  $$
  \forall v\in F, \qquad p_{F}(v)=<v,f_{1}>f_{1}
                                 +<v,f_{2}>f_{2}
                                 +\cdots
                                 +<v,f_{r}>f_{r}
  $$ 
\finenonce{001478}



\finexercice



\exercice{1479, legall, 1998/09/01}

\enonce{001479}{}
Soit $  (E , \langle   ,  \rangle )  $ un espace euclidien et $  p \in
\mathcal{L} (E)  $ un projecteur. Montrer que $  p   $ est orthogonal
(c'est-\`a-dire $  \hbox{Ker}(p)\perp
\hbox{Im}(p)  $) si et seulement si~: $  \forall x \in E  :  \Vert p(x)\Vert \leq \Vert x\Vert
.$
\finenonce{001479}



\finexercice



\exercice{1480, legall, 1998/09/01}

\enonce{001480}{}
Soit $  (E , \langle   ,  \rangle )  $ un espace euclidien et $  F
$ un sous-espace vectoriel de $  E  .$ On note $  p  $ la projection orthogonale sur $  F  $
et on pose, pour tout $  \displaystyle{ x\in E  :   d(x, F)= \inf _{y\in F}\Vert x-y\Vert }  .$
Soit $  z \in F  .$
\begin{enumerate}
    \item Montrer que pour tout $  x\in F  , $ les trois conditions sont \'equivalentes~:
\begin{enumerate}
\item[(i)] $  d(x,F)=\Vert x-z\Vert   .$
\item[(ii)] $  z=p(x)  .$
\item[(iii)] $  \forall y \in F  ,   y\perp (x-z)  .$
\end{enumerate}
    \item En d\'eduire $  \displaystyle{ \inf _{a,b\in \R } \int _0 ^1 (x^2-ax-b)^2dx }  .$
\end{enumerate}
\finenonce{001480}



\finexercice



\exercice{1481, legall, 1998/09/01}

\enonce{001481}{}
Soit $  (E , \langle   ,  \rangle )  $ un espace euclidien de dimension
sup\'erieure ou \'egale \`a $  2  .$ Soient $  x  $ et $   y \in E  .$ Montrer que~:
\begin{enumerate}
    \item Si $  \Vert x\Vert = \Vert y\Vert  ,$ alors il existe un hyperplan $  H  $ de
$  E  $ tel que $  y = s(x)  $ o\`u $  s  $ est la sym\'etrie orthogonale par rapport
\`a $  H  .$
    \item Si $  \langle x, y\rangle = \Vert y\Vert ^2  ,$ alors il existe un hyperplan $  H  $ de $  E  $
tel que $  y = p(x)  $ o\`u $  p  $ est la projection orthogonale sur $  H  .$
\end{enumerate}
\finenonce{001481}



\finexercice



\exercice{1482, legall, 1998/09/01}

\enonce{001482}{}
Dans $  \R ^3  $ muni du produit scalaire euclidien
canonique, donner la matrice de
la projection orthogonale sur le plan d'\'equation $  x + 2y - 3z = 0  .$ Donner la matrice de la
sym\'etrie orthogonale par rapport \`a ce m\^eme plan.
\finenonce{001482}



\finexercice



\exercice{1483, legall, 1998/09/01}

\enonce{001483}{}
Soit $  (E , \langle   ,  \rangle )  $ un espace euclidien et $
  F  $ un sous-espace vectoriel muni d'une base
orthonormale $  (e_1, \ldots, e_m)  .$ Soit $  p  $ la projection orthogonale sur
$  F  .$
\begin{enumerate}
    \item Montrer que $  \displaystyle{\forall x \in E  ,   p(x) = \sum_{i=1}^m \langle
x,e_i\rangle e_i}  .$
    \item Donner de m\^eme l'expression de la sym\'etrie orthogonale par rapport
\`a $  F  $ et la projection orthogonale sur $  F^\perp   .$
\end{enumerate}
\finenonce{001483}



\finexercice



\exercice{1484, ridde, 1999/11/01}

\enonce{001484}{}
Quelle est la transformation de $\Rr^{3}$ dont la matrice dans la base canonique
est $\frac{1}{7} \begin{pmatrix} -2 & 6&-3\\ 6&3&2\\-3&2&6\\ \end{pmatrix}$ ?
\finenonce{001484}



\finexercice



\exercice{1485, ridde, 1999/11/01}

\enonce{001485}{}
D\'eterminer la matrice dans la base canonique de $\Rr^{4}$ de la projection orthogonale
sur $\mathrm{Vect} (v_{1}, v_{2})$ où $v_{1} = (1, -1, 0, 0)$ et $v_{2} = (0, 1, 0, 1)$.
\finenonce{001485}



\finexercice



\exercice{1486, ridde, 1999/11/01}

\enonce{001486}{}
Soient $E$ un espace euclidien, $u$ un vecteur non nul et $H = u^{\bot}$.
Soient $p$ la projection orthogonale sur $H$ et $s$ la sym\'etrie orthogonale par
rapport \`a $H$.
\begin{enumerate}
\item Montrer que $\forall x \in E \qquad p (x) = x-\frac{\langle x|u \rangle}
{\left\| u\right\|^{2}}u$.
\item Montrer que $\forall x \in E \qquad s (x) = x-2\frac{\langle x|u \rangle}
{\left\| u\right\|^{2}}u$.
\item On consid\`ere dans $\Rr^{3}$ le plan $ (\Pi : x-y + z = 0)$. D\'eterminer la matrice
dans la base canonique de la sym\'etrie orthogonale par rapport \`a $\Pi$.
\end{enumerate}
\finenonce{001486}



\finexercice



\exercice{1487, ortiz, 1999/04/01}

\enonce{001487}{}
Soit $\left(E,\;|\;\right)$ un espace vectoriel de
dimension $n$.
\begin{enumerate}
    \item  Soient $F$ et $G$ des sous-espace vectoriels de $E$. Montrer que $\left(F\cap G\right)^\perp=F^\perp+G^\perp $.
    \item  Soient $\mathcal{B}=(e_1,...,e_n)$ une base orthonormale  de $E$ et
$(a_1,...,a_n)\in\Rr^n\setminus \left\{(0,...,0)\right\} $
et $H$ le sous-espace vectoriel de $E$
d'\'equation cart\'esienne
$\sum_{k=1}^{n}a_kx_k=0$ dans $\mathcal{B}.$
     \begin{enumerate}
         \item D\'eterminer l'orthogonale de $H$.
         \item D\'eterminer la distance du vecteur $x=\sum_{k=1}^{n}x_ke_k$ de $E$ au sous-espace vectoriel $H$.
    \end{enumerate}
    \item  Soit $P$ le sous-espace vectoriel de l'espace vectoriel $\Rr^4$ d\'efini par
$${u=(x_1,x_2,x_3,x_4)\in P\Leftrightarrow
x_1+x_2+x_3+x_4=x_2+2x_3+3x_4=0.}$$
    \begin{enumerate}
         \item D\'eterminer une base de $P^\perp$ puis une base orthonormale  de $P^\perp$.
         \item En d\'eduire une expression analytique de la projection orthogonale de $\Rr^4$ sur $P$.
    \end{enumerate}
\end{enumerate}
\finenonce{001487}



\finexercice



\exercice{1488, ortiz, 1999/04/01}

\enonce{001488}{}
Soient $E$ un espace vectoriel euclidien, $F$ et
$G$ deux sous-espace vectoriels suppl\'ementaires
de $E$ et $p$ le projecteur de $E$ d'axe $F$ et de
direction $G$.
\begin{enumerate}
    \item  On suppose que $F\perp G$. Montrer que
$\forall x\in E, \|p(x)\|\le\|x\|.$
    \item  On suppose que $\forall x\in E, \|p(x)\|\le\|x\|$.
    \begin{enumerate}
         \item Soient $a\in F$ et $b\in G$. Montrer que $\|a+b\|\ge\|a\|$.
         \item En d\'eduire que $F\perp G$.
    \end{enumerate}
\end{enumerate}
\finenonce{001488}



\finexercice



\exercice{1489, ortiz, 1999/04/01}

\enonce{001489}{}
Soit
$\alpha=\inf\left\{\int_{-1}^{1}\left(ax^2+bx+c-|{x}|\;\right)^2dx:a,b,c\in\Rr\right\}$.
\begin{enumerate}
    \item  D\'eterminer un espace vectoriel euclidien $(E,\,|\,)$, un sous-espace vectoriel $F$ de $E$ et $v\in E$ tel que $\alpha=\text{d}(v,F)^2$.
    \item  D\'eterminer $p\in F$ tel que $\alpha=\text{d}(v,p)^2$ et $\alpha$.
\end{enumerate}
\finenonce{001489}



\finexercice



\exercice{1490, ridde, 1999/11/01}

\enonce{001490}{}
Soit $E$ un espace euclidien (de dimension finie),
 $F$ et $G$ deux sous-espaces vectoriels de $E$.
D\'eterminer $ (F + G)^{\bot}$ et $ (F \cap G)^{\bot}$ en fonction de $F^{\bot}$
 et $G^{\bot}$.
\finenonce{001490}



\finexercice



\exercice{1491, ridde, 1999/11/01}

\enonce{001491}{}
D\'eterminer $\inf\limits_{ (a, b)\in \Rr^{2}} \int_{0}^{1}{ (e^{x}- (ax + b))^{2}dx}$.
\finenonce{001491}



\finexercice



\exercice{1492, gourio, 2001/09/01}

\enonce{001492}{}
Calculer :
$$\inf\limits_{(a,b)\in {\Rr}^{2}}\int_{0}^{1}x^{2}\left| \ln x-ax-b\right|^{2}dx. $$
\finenonce{001492}



\finexercice



\exercice{3733, quercia, 2010/03/11}
\enonce{003733}{Isométries affines}
Soit $E$ un espace euclidien de dimension $n$.
Soit $f : E \to E$ une application non nécéssairement linéaire.

\begin{enumerate}
  \item On suppose que $f$ conserve le produit scalaire.
    Démontrer que $f$ est linéaire.

  \item On suppose que $f$ conserve les distances, c'est à dire :
    $\forall\ \vec x,\vec y \in E,\ \| f(\vec x) - f(\vec y) \| = \|\vec x - \vec y\,\|$.
    Démontrer que $f = f(\vec0) + g$, avec $g \in {\cal O}(E)$.
\end{enumerate}
\finenonce{003733}



\finexercice


\exercice{3734, quercia, 2010/03/11}
\enonce{003734}{Projection sur vect($u$)}

Soit $E$ un espace euclidien de dimension $n$.
Soit $\vec u$ un vecteur unitaire de matrice $U$ dans une base orthonormée
$\cal B$.

\begin{enumerate}
  \item Montrer que $U^tU$ est la matrice dans $\cal B$ de la projection orthogonale
    sur vect$(\vec u)$.
  \item Trouver la matrice de la symétrie associée.
\end{enumerate}
\finenonce{003734}



\finexercice


\exercice{3735, quercia, 2010/03/11}
\enonce{003735}{$\vec x + \lambda (\vec x\mid \vec v)\vec v$}

Soit $\vec v \in E\setminus\{\vec0\}$ et $\lambda \in \R$.
On pose pour $\vec x \in E$ :
$f(\vec x) = \vec x + \lambda (\vec x\mid \vec v)\vec v$.

Déterminer $\lambda$ pour que $f \in {\cal O}(E)$. Reconnaître alors $f$.

\finenonce{003735}



\finexercice


\exercice{3736, quercia, 2010/03/11}
\enonce{003736}{Composition de symétries}
Soit $E$ un espace euclidien de dimension $n$.
Soient $F,G$ deux sous-espaces de $E$ tels que $F \perp G$. On note $s_F$
et $s_G$ les symétries orthogonales de bases $F$ et~$G$.
Montrer que $s_F\circ s_G = s_G \circ s_F = s_{(F\oplus G)^\perp}$.
\finenonce{003736}



\finexercice


\exercice{3737, quercia, 2010/03/11}
\enonce{003737}{Composition de symétries}
Soit $E$ un espace euclidien de dimension $n$.
Soient $F,G$ deux sous-espaces de $E$ tels que $F \subset G$. On note $s_F$
et $s_G$ les symétries orthogonales de bases $F$ et~$G$.
Montrer que $s_F\circ s_G = s_G \circ s_F = s_{F\oplus G^\perp}$.
\finenonce{003737}



\finexercice


\exercice{3738, quercia, 2010/03/11}
\enonce{003738}{Condition pour que deux symétries commutent}
Soit $E$ un espace euclidien de dimension $n$.
Soient $H,K$ deux hyperplans de $E$, et $s_H$, $s_K$ les symétries associées.

Démontrer que $s_H$ et $s_K$ commutent si et seulement si $H = K$ ou
$H^\perp \subset K$.
\finenonce{003738}



\finexercice


\exercice{3739, quercia, 2010/03/11}
\enonce{003739}{Similitudes}
Soit $E$ un espace euclidien de dimension $n$.
Soit $f \in \mathcal{L}(E)$ et $\lambda > 0$.
On dit que $f$ est une similitude de rapport $\lambda$ si :
$\forall\ \vec x \in E,\ \|f(\vec x)\| = \lambda\|\vec x\,\|$.

\begin{enumerate}
  \item Montrer que $f$ est une similitude de rapport $\lambda$ si et seulement si :
    $\forall\ \vec x,\vec y\in E,\
    (f(\vec x)\mid f(\vec y)) = \lambda^2(\vec x\mid\vec y)$.
  \item Caractériser les similtudes par leurs matrices dans une base orthonormée.
  \item Montrer que $f$ est une similitude si et seulement si $f$ est non nulle et
    conserve l'orthogonalité, c'est à dire~:
    $\forall\ \vec x,\vec y\in E,\ \vec x\perp\vec y
       \Rightarrow  f(\vec x)\perp f(\vec y)$.
\end{enumerate}
\finenonce{003739}



\finexercice


\exercice{3740, quercia, 2010/03/11}
\enonce{003740}{Similitudes}

On définit l'application $\varphi : A\in\mathcal{M}_n(\R) \mapsto \sum_{i,j}a_{i,j}^2$.
Trouvez les matrices $P\in GL_n(\R)$ telles que
pour tout $A$ on ait $\varphi(P^{-1}AP)=\varphi(A)$.	

\finenonce{003740}



\finexercice


\exercice{3741, quercia, 2010/03/11}
\enonce{003741}{Sous-espaces stables}
Soit $E$ un espace euclidien de dimension $n$.
Soit $u \in {\cal O}(E)$ et $F$ un sous-espace vectoriel stable par $u$.
Montrer que $F^\perp$ est aussi stable par $u$.
\finenonce{003741}



\finexercice


\exercice{3742, quercia, 2010/03/11}
\enonce{003742}{Projection sur le sous-espace invariant}
Soit $E$ un espace euclidien de dimension $n$.
Soit $u \in {\cal O}(E)$.
On note $v = \text{id}_E - u$.

\begin{enumerate}
  \item Montrer que $\mathrm{Ker} v = (\Im v)^\perp$.
  \item Soit $p$ la projection orthogonale sur $\mathrm{Ker} v$.
    Montrer que :
    $\forall\ \vec x \in E,\ \frac 1m\sum_{k=0}^{m-1}u^k(\vec x)
     \to p(\vec x)$ lorsque $m\to\infty$.

\end{enumerate}
\finenonce{003742}



\finexercice


\exercice{3743, quercia, 2010/03/11}
\enonce{003743}{Endomorphismes orthogonaux diagonalisables}

Quels sont-ils ?
\finenonce{003743}



\finexercice


\exercice{3744, quercia, 2010/03/11}
\enonce{003744}{Valeurs propres d'une isométrie}
Soit $E$ un espace euclidien de dimension $n$.
Soit $f \in {\cal O}(E)$.

\begin{enumerate}
  \item On suppose $n$ impair et $f \in {\cal O}^+(E)$.
    Montrer que 1 est valeur propre de $f$.
    (comparer $\det(f-\mathrm{id})$ et $\det(f^{-1}-\mathrm{id})$)
  \item Que peut-on dire quand $n$ est pair ?
  \item Soit $n$ quelconque, $f \in {\cal O}^{-}(E)$.
    Montrer que $-1$ est valeur propre de $f$.
\end{enumerate}
\finenonce{003744}



\finexercice


\exercice{3745, quercia, 2010/03/11}
\enonce{003745}{Caractérisation des symétries orthogonales}

Soit $M \in {\cal O}(n)$.

\begin{enumerate}
  \item Montrer que $M$ est la matrice d'une symétrie orthogonale si et seulement si
    $M$ est symétrique.
  \item Dans ce cas, déterminer la base et la direction de cette symétrie en
    fonction des matrices $I+M$ et $I-M$.
\end{enumerate}
\finenonce{003745}



\finexercice


\exercice{3746, quercia, 2010/03/11}
\enonce{003746}{$f$ orthogonal antisymétrique}
Soit $E$ un espace euclidien de dimension $n$.
Soit $f\in{\cal O}(E)$. Montrer que les énoncés suivants sont équivalents :
$$f\circ f = -\mathrm{id}_E \iff \forall\ \vec x\in E,\ f(\vec x) \perp \vec x
  \iff \forall\ \vec x,\vec y \in E,\
  (f(\vec x)\mid \vec y) = -(\vec x\mid f(\vec y)).$$
\finenonce{003746}



\finexercice


\exercice{3747, quercia, 2010/03/11}
\enonce{003747}{Centre de O(E)}
Soit $E$ un espace euclidien de dimension $n$.
Soit $f \in {\cal O}(E)$, et $s$ une réflexion par rapport à un hyperplan $H$.
Soit $\vec u \in H^\perp, \vec u \ne \vec0$.

\begin{enumerate}
  \item Montrer que $f\circ s \circ f^{-1}$ est aussi une symétrie et en donner la
    base.
  \item En déduire que $f$ et $s$ commutent si et seulement si $\vec u$ est vecteur
    propre de $f$.
  \item Quel est le centre de $O(E)$ ?
\end{enumerate}
\finenonce{003747}



\finexercice


\exercice{3748, quercia, 2010/03/11}
\enonce{003748}{Nombre de réflexions nécéssaires pour engendrer une application donnée}
Soit $E$ un espace euclidien de dimension $n$.
Soit $f \in {\cal O}(E)$. On note $F = \mathrm{Ker}(f - \text{id}_E)$
et $p = \text{codim}(F)$.

\begin{enumerate}
  \item Montrer qu'on peut décomposer $f$ en produit d'au plus $p$ réflexions.
  \item Inversement, si $f$ est un produit de $k$ réflexions, démontrer que
    $p \le k$.
  \item Application : trouver $f\in {\cal O}(E)$ qui se décompose en $n$ réflexions
    et pas moins.
\end{enumerate}
\finenonce{003748}



\finexercice


\exercice{3749, quercia, 2010/03/11}
\enonce{003749}{Prolongement d'une transformation orthogonale}
Soit $E$ un espace euclidien de dimension $n$.
\begin{enumerate}
  \item Soit $F$ un sous-espace vectoriel de $E$ et $f : F \to E$ une application orthogonale.
    Démontrer qu'on peut prolonger $f$ en une application orthogonale
    de $E$ dans $E$.

  \item Soient $\vec u_1,\dots,\vec u_n$, $\vec v_1,\dots,\vec v_n$ des vecteurs
    de $E$ tels que :
    $\forall\ i,j,\ (\vec u_i \mid \vec u_j) = (\vec v_i \mid \vec v_j)$.
  \begin{enumerate}
    \item Si $\sum \lambda_i \vec u_i = \vec0$, démontrer que
    $\sum \lambda_i \vec v_i = \vec0$.
    \item En déduire qu'il existe $f \in {\cal O}(E)$ telle que :
    $\forall\ i,\ f(\vec u_i) = \vec v_i$.
  \end{enumerate}  
\end{enumerate}
\finenonce{003749}



\finexercice


\exercice{3750, quercia, 2010/03/11}
\enonce{003750}{Action de ${\cal O}(E)$ sur les sous-espaces vectoriels}
Soit $E$ un espace euclidien de dimension $n$.
Soient $F,G$ deux sous-espaces vectoriels de $E$ de même dimension.

\begin{enumerate}
  \item Montrer qu'il existe $u \in {\cal O}(E)$ tel que $u(F) = G$.
  \item $(*)$ Montrer qu'il existe $u \in {\cal O}(E)$ tel que $u(F) = G$ et $u(G) = F$.
\end{enumerate}
\finenonce{003750}



\finexercice


\exercice{3751, quercia, 2010/03/11}
\enonce{003751}{Transformations orthogonales sur $\mathcal{M}_n(\R)$}

$E = \mathcal{M}_n(\R)$ muni du produit scalaire : $(A\mid B) = \mathrm{tr}({}^tAB)$.

\begin{enumerate}
  \item Vérifier que c'est un produit scalaire.

  \item Soit $P \in {\cal O}(n)$. Montrer que les applications
    $\begin{cases} \phi_P : A \longmapsto AP \cr
             \psi_P : A \longmapsto P^{-1}AP \cr\end{cases}$
    sont orthogonales.

  \item Réciproquement, si $\phi_P$ ou $\psi_P \in {\cal O}(\mathcal{M}_n(\R))$,
    est-ce que $P \in {\cal O}(n)$ ?

\end{enumerate}
\finenonce{003751}



\finexercice


\exercice{3752, quercia, 2010/03/11}
\enonce{003752}{$A$ = com$(A)$}

Quelles sont les matrices $A \in \mathcal{M}_n(\R)$ égales à leur comatrice ?
\finenonce{003752}



\finexercice


\exercice{3753, quercia, 2010/03/11}
\enonce{003753}{$A{}^t\!A + A + {}^t\!A = 0$}

\begin{enumerate}
  \item Trouver les matrices $A \in \mathcal{M}_n(\R)$ telles que :
    $A{}^t\!A + A + {}^t\!A = 0$.
    
  \item Montrer que pour une telle matrice, $|\det A| \le 2^n$.
    

\end{enumerate}
\finenonce{003753}



\finexercice


\exercice{3754, quercia, 2010/03/11}
\enonce{003754}{Somme des coefficients d'une matrice orthogonale}

Soit $P \in {\cal O}(n)$. A l'aide de l'inégalité de Cauchy-Schwarz,
démontrer que : $\big|\sum_{i,j} P_{ij}\big| \le n$.

Quand a-t-on égalité ?
\finenonce{003754}



\finexercice


\exercice{3755, quercia, 2010/03/11}
\enonce{003755}{Groupe engendré par les réflexions}

Soit $E$ un espace vectoriel hermitien et $G$ le sous-groupe
de~$U(E)$ engendré par les réflexions.

\begin{enumerate}
  \item On suppose ici $\dim E\ge 2$.
    Soient $u,v\in E$. Montrer qu'il existe une réflexion $\sigma$
    échangeant $u$ et~$v$ si et seulement si $\|u\| = \|v\|$ et $(u\mid v)\in\R$.

  \item Montrer que $G = \{f\in U(E)\ \text{ tq }\det(f)\in\{-1,1\}\}$.

\end{enumerate}
\finenonce{003755}



\finexercice


\exercice{3756, quercia, 2010/03/11}
\enonce{003756}{Orthotrigonalisation}

Montrer que tout endomorphisme d'un espace vectoriel hermitien est
trigonalisable en base orthonormale.
\finenonce{003756}



\finexercice


\exercice{3757, quercia, 2010/03/11}
\enonce{003757}{Réduction des endomorphismes orthogonaux et unitaires}

\begin{enumerate}
  \item Soit $G$ un sous groupe compact de $GL_n(\C)$.
    Si $M\in G$ montrer que $\mathrm{Sp}(M)\subset\mathbb{U}$.
    En déduire que tout élément de~$G$ est diagonalisable.

  \item Soit $M\in U_n(\C)$. Montrer que $M$ est diagonalisable et
    qu'il existe une base orthonormée de~$\C^n$ propre pour~$M$.

  \item Soit $M\in O_n(\R)$. Montrer qu'il existe $P\in O_n(\R)$
    telle que $P^{-1}MP$ est diagonale par blocs~:

    $P^{-1}MP = \mathrm{diag}(\pm1,\dots,\pm1,R_1,\dots,R_k)$ où chaque~$R_i$
    est une matrice de la forme~:
    $R_i = \begin{pmatrix}\cos\theta_i &-\sin\theta_i\cr
                    \sin\theta_i &\cos\theta_i\cr\end{pmatrix}$.
\end{enumerate}
\finenonce{003757}



\finexercice


\exercice{3758, quercia, 2010/03/11}
\enonce{003758}{Groupes orthogonaux égaux}

Soit $E$ un $\R$-espace vectoriel de dimension finie et $\|\ \|_1$, $\|\ \|_2$
deux normes euclidiennes telles que les groupes orthogonaux associés sont
égaux. Que peut-on dire de ces normes~?
\finenonce{003758}



\finexercice


\exercice{3759, quercia, 2010/03/11}
\enonce{003759}{Calcul de norme}

Soit $f : {\R^n} \to {\R^n}, {(x_1,\dots,x_n)} \mapsto {(x_1-x_n,x_2-x_1,\dots,x_n-x_{n-1}).}$
Avec la structure euclidienne canonique de~$\R^n$, calculer la norme de~$f$.
\finenonce{003759}



\finexercice


\exercice{3760, quercia, 2010/03/11}
\enonce{003760}{Densité (Ens Ulm MP$^*$ 2003)}
${\cal O}_n(\Q)$ est-il dense dans~${\cal O}_n(\R)$~?
\finenonce{003760}



\finexercice


\exercice{3761, quercia, 2010/03/11}
\enonce{003761}{Centrale MP 2003}
On considère l'endomorphisme de~$\R^3$ de matrice dans la base canonique~:
$A = \begin{pmatrix}a^2 &ab-c &ac+b\cr ab+c &b^2 &bc-a\cr ac-b &bc+a &c^2\cr\end{pmatrix}$
avec~$a,b,c\in\R$.
Déterminer $a,b,c$ de sorte que~$f$ soit une isométrie, et la préciser.
\finenonce{003761}



\finexercice


\exercice{5486, rouget, 2010/07/10}
\enonce{005486}{***I Matrices et déterminants de \textsc{Gram}}
Soit $E$ un espace vectoriel euclidien de dimension $p$ sur $\Rr$ ($p\geq2$).
Pour $(x_1,...,x_n)$ donné dans $E^n$, on pose $G(x_1,...,x_n)=(x_i|x_j)_{1\leq i,j\leq n}$ (matrice de \textsc{Gram})
et $\gamma(x_1,...,x_n)=\mbox{det}(G(x_1 , ... , x_n))$ (déterminant de \textsc{Gram}).
\begin{enumerate}
\item  Montrer que $\mbox{rg}(G(x_1,...,x_n))=\mbox{rg}(x_1, ... ,x_n)$.
\item  Montrer que $(x_1,...,x_n)$ est liée si et seulement si $\gamma(x_1,...,x_n)=0$ et que $(x_1,...,x_n)$ est libre si et seulement si $\gamma(x_1,...,x_n)>0$.
\item  On suppose que $(x_1,...,x_n)$ est libre dans $E$ (et donc $n\leq p$). On pose $F=\mbox{Vect}(x_1,...,x_n)$.

Pour $x\in E$, on note $p_F(x)$ la projection orthogonale de $x$ sur $F$ puis $d_F(x)$ la distance de $x$ à $F$ (c'est-à-dire $d_F(x)=||x-p_F(x)||$). Montrer que $d_F(x)=\sqrt{\frac{\gamma(x,x_1,...,x_n)}{\gamma(x_1,...,x_n)}}$.
\end{enumerate}
\finenonce{005486}


\finexercice


\exercice{5488, rouget, 2010/07/10}
\enonce{005488}{**I}
Matrice de la projection orthogonale sur la droite d'équations $3x=6y=2z$ dans la base canonique orthonormée de $\Rr^3$ ainsi que de la symétrie orthogonale par rapport à cette même droite.
De manière générale, matrice de la projection orthogonale sur le vecteur unitaire $u=(a,b,c)$ et de la projection orthogonale sur le plan d'équation $ax+by+cz=0$ dans la base canonique orthonormée de $\Rr^3$.
\finenonce{005488}


\finexercice


\exercice{5494, rouget, 2010/07/10}
\enonce{005494}{**I}
Existence, unicité et calcul de $a$ et $b$ tels que $\int_{0}^{1}(x^4-ax-b)^2\;dx$ soit minimum (trouver deux démonstrations, une dans la mentalité du lycée et une dans la mentalité maths sup).
\finenonce{005494}


\finexercice


\exercice{5504, rouget, 2010/07/15}
\enonce{005504}{**T}
Dans $\Rr^3$, espace vectoriel euclidien orienté rapporté à la base orthonormée directe $(i,j,k)$, déterminer l'image du plan d'équation $x+y=0$ par \begin{enumerate}
 \item  la symétrie orthogonale par rapport au plan d'équation $x-y+z=0$,  
 \item  la symétrie orthogonale par rapport au vecteur $(1,1,1)$,  
 \item  par la rotation d'angle $\frac{\pi}{4}$ autour du vecteur $(1,1,1)$.
\end{enumerate}
\finenonce{005504}


\finexercice


\exercice{5509, rouget, 2010/07/15}
\enonce{005509}{**T}
Dans $\Rr^3$, soient $(D)$ $\left\{
\begin{array}{l}
x+y+z=1\\
x-2y-z=0
\end{array}
\right.$ et $(\Delta)~:~6x=2y=3z$ puis $(P)~:~x+3y+2z=6$. Déterminer la projection de $(D)$ sur $(P)$ parallèlement à $(\Delta)$.

\finenonce{005509}


\finexercice


\exercice{5520, rouget, 2010/07/15}
\enonce{005520}{**}
Soient $(D)$ la droite dont un système d'équations cartésiennes est 
$\left\{
\begin{array}{l}
x+y+z=1\\
x-2y-z=0
\end{array}
\right.$ et $(P)$ le plan d'équation cartésienne $x+3y+2z=6$. Déterminer la projetée (orthogonale) de $(D)$ sur $(P)$.
\finenonce{005520}


\finexercice

\section{204.05 Orthonormalisation}


\exercice{1493, legall, 1998/09/01}

\enonce{001493}{}
R\'esoudre l'\'equation $(1-x)^2 + (x-y)^2 + (y-z)^2 + z^2 = \frac{1}{4}$
pour $(x,y,z) \in \R^3$.
\finenonce{001493}



\finexercice



\exercice{1494, legall, 1998/09/01}

\enonce{001494}{}
\begin{enumerate}
    \item Soit $F$ le sous-espace de $\R^5$ engendr\'e par $u = (1, 2, 3, -1,2)$
 et $v = (2, 4, 7, 2, -1)$. Trouver une base de l'orthogonal $F^\perp$
de $F$.
    \item Trouver une base orthonormale du sous-espace $E$ de $\C^3$ engendr\'e
par $v_1 = (1, i, 0)$ et $v_2 = (1, 2, 1-i)$.
\end{enumerate}
\finenonce{001494}



\finexercice



\exercice{1495, legall, 1998/09/01}

\enonce{001495}{}
Soit $F$ un sous-espace d'un espace euclidien $E$. Montrer qu'il existe une
base orthonormale de $F$ qui est inclue dans une base orthonormale de $E$.
\finenonce{001495}



\finexercice



\exercice{1496, legall, 1998/09/01}

\enonce{001496}{}
\begin{enumerate}
    \item Soit $A =
\left(\begin{array}{ccc}2&1&1\\1&1&1\\1&1&2\end{array}\right)$. Montrer que
$A$ d\'efinit un produit scalaire $\varphi$ sur $\R^3$. Construire une
base orthonormale pour $\varphi$.
    \item Consid\'erons une base  $\{v_1 = (1,1,1), v_2 = (0,1,1), v_3 =
(0,0,1)\}$ de l'espace euclidien $\R^3$. Utiliser le proc\'ed\'e
d'orthogonalisation de Schmidt pour transformer $\{v_i\}$ en une base
orthonormale.
\end{enumerate}
\finenonce{001496}



\finexercice



\exercice{1497, legall, 1998/09/01}

\enonce{001497}{}
Soient $E = \R_n[X], I_n = \frac{1}{\sqrt{2\pi}}\int_{-\infty}^{+\infty}
t^ne^{\frac{-t^2}{2}}dt$.
\begin{enumerate}
    \item Montrer que l'int\'egrale $I_n$ est convergente. Que vaut $I_{2p+1}$ ?\par
Soit $\varphi : E \times E \to \R$ d\'efinie par $\varphi(P,Q) =
\frac{1}{\sqrt{2\pi}}\int_{-\infty}^{+\infty}P(t)Q(t)e^{\frac{-t^2}{2}}dt$.
    \item Montrer que $\varphi$ est un produit scalaire.
    \item On suppose $n = 2$. Ecrire la matrice associ\'ee \`a $\varphi$ dans la
base $(1,X,X^2)$. Construire une base orthonormale $(P_0,P_1,P_2)$ par le
proc\'ed\'e d'orthogonalisation de Schmidt appliqu\'e \`a $(1,X,X^2)$.
\end{enumerate}
\finenonce{001497}



\finexercice



\exercice{1498, legall, 1998/09/01}

\enonce{001498}{}
R\'eduire en somme de carr\'es ind\'ependants les formes suivantes :
\begin{enumerate}
    \item $9x^2 - 6y^2 - 8z^2 + 6xy - 14xz  + 18xw + 8yz + 12yw - 4zw$
    \item $x_1^2 + x_2^2 + x_3^2 - 2x_4^2 - 2x_1x_2 + 2x_1x_3 - 2x_1x_4 + 2x_2x_3
- 4x_2x_4$
\end{enumerate}
\finenonce{001498}



\finexercice



\exercice{1499, ortiz, 1999/04/01}

\enonce{001499}{}
$\Rr^3$ est muni de sa structure canonique
d'espace vectoriel euclidien. V\'erifier que les
vecteurs $e_1=(1,0,1),$ $e_2=(1,0,2)$ et
$e_3=(1,1,1)$ forment une base de $\Rr^3$ et  en
d\'eterminer l'orthonormalis\'ee de Gram-Schmidt.
\finenonce{001499}



\finexercice



\exercice{1500, ortiz, 1999/04/01}

\enonce{001500}{}
$\Rr^4$ est muni de sa structure canonique
d'espace vectoriel euclidien. Soient
$e_1=(1,0,1,0)$ et $e_2=(1,-1,1,-1)$ et
$F=\text{vect}(e_1,e_2)$.
\begin{enumerate}
    \item  D\'eterminer une base orthonormale de $F.$
    \item  D\'eterminer la matrice dans la base canonique de $\Rr^4$ du projecteur orthogonal
sur $F.$
    \item  D\'eterminer la distance du vecteur $(1,1,1,1)$ au sous-espace vectoriel  $F.$
\end{enumerate}
\finenonce{001500}



\finexercice



\exercice{1501, ortiz, 1999/04/01}

\enonce{001501}{}
On munit le $\Rr$-espace vectoriel  $\Rr_2[X]$ du
produit scalaire d\'efini par
$${\phi:\Rr_2[X]\to\Rr_2[X],
(P,Q)\mapsto\int_{-1}^{1}P(t)Q(t)dt.}$$
\begin{enumerate}
    \item  D\'eterminer l'orthonormalis\'ee de Gram-Schmidt de la base canonique de
    $\Rr_2[X]$.
    \item  D\'eterminer la distance du polyn\^ome $P=X^2+X+1$ au sous-espace vectoriel
    $F$ de  $\Rr_2[X]$ form\'e des polyn\^omes $f$ tels que $f'(0)=0.$
\end{enumerate}
\finenonce{001501}



\finexercice



\exercice{1502, ortiz, 1999/04/01}

\enonce{001502}{}
Soit $f:\Rr^3\times\Rr^3\to \Rr$ d\'efinie de la
mani\`ere suivante : si $u=(x,y,z)$ et
$u'=(x',y',z')$ alors
$${f(u,u')=2xx'+yy'+2zz'+xy'+yx'+xz'+zx'+yz'+zy'.}$$
\begin{enumerate}
    \item  Montrer que $f$ est un produit scalaire sur l'espace vectoriel canonique $\Rr^3$.
    \item  Soit $P$ le sous-espace vectoriel de $\Rr^3$ d'\'equation cart\'esienne $2x-y+z=0$.
    \begin{enumerate}
        \item D\'eterminer l'orthogonal du sous-espace vectoriel $P$.
        \item D\'eterminer un sous-espace vectoriel de $\Rr^3$ dont l'orthogonal est $P$.
    \end{enumerate}
    \item  D\'eterminer l'orthonormalis\'ee de Gram-Schmidt de la base canonique de $\Rr^3$  pour le produit scalaire $f$.
\end{enumerate}
\finenonce{001502}



\finexercice



\exercice{1503, ridde, 1999/11/01}

\enonce{001503}{}
Orthonormaliser dans $\Rr^{3}$ la famille $x_{1} = (1, -2, 2)$, $x_{2} = (-1, 0, -1)$,
$x_{3} = (5, -3, 7)$.
\finenonce{001503}



\finexercice



\exercice{1504, ridde, 1999/11/01}

\enonce{001504}{}
D\'eterminer une base orthonorm\'ee de $\Rr_{2}[X]$ muni du produit scalaire
$\langle P|Q \rangle = \int_{0}^{1}{P (t)Q (t)dt}$.
\finenonce{001504}



\finexercice



\exercice{1505, barraud, 2003/09/01}

\enonce{001505}{}
  On considère la forme bilinéaire $b$ de $\R^{4}$ définie par :
  $$
  b(x,y)=x_{1}y_{1}+2x_{2}y_{2}+4x_{3}y_{3}+18x_{4}y_{4}
  + x_{1}y_{3}+ x_{3}y_{1}
  +2x_{2}y_{4}+2x_{4}y_{2}
  +6x_{3}y_{4}+6x_{4}y_{3}
  $$
  où $x_{1},x_{2},x_{3}$ et $y_{1},y_{2},y_{3}$ sont les coordonnées
  de $x$ et $y$ dans la base canonique.
  \begin{enumerate}
  \item 
    Montrer qu'il s'agit d'un produit scalaire.
  \item 
    Ecrire la matrice de $b$ dans la base canonique.
  \item 
    Trouver une base orthonormée pour $b$.
  \end{enumerate}
\finenonce{001505}



\finexercice



\exercice{1506, barraud, 2003/09/01}

\enonce{001506}{}
On considère un espace euclidien $(E,<>)$.

\begin{enumerate}
\item
\textbf{Théorème de Pythagore :}

Soient $u$ et $v$ deux vecteurs orthogonaux de $E$. Calculer $||u+v||^{2}$. Illustrer le
résultat obtenu à l'aide d'un dessin.

\item
\textbf{Projection orthogonale et distance à un sous-espace :}

Soit $F$ un sous-espace de $E$. On rappelle que $E=F\oplus F^{\bot}$, et donc que tout
vecteur $x$ de $E$ se décompose de manière unique en une somme $x=x_{1}+_{2}$ avec
$x_{1}\in F$ et $x_{2}\in F^{\bot}$. Le vecteur $x_{1}$ s'appelle alors la projection
orthogonale de $x$ sur $F$.
\begin{enumerate}
  \item
Montrer que l'application $p$ qui à un vecteur asocie sa projection orthogonale sur $E$
est une application linéaire. Vérifier que : $\forall y\in F, <x-p(x),y>=0$.
  \item\label{q: proj-realise-d}
On considère maintenant un vecteur $x$ de $E$. On appelle distance de $x$ à $F$ le nombre
$\mathrm{dist}(x,F)=\inf_{y\in F}\Vert x-y\Vert$.

Pour $y\in F$, vérifier que $x-p(x)$ et $y-p(x)$ sont orthogonaux. Utiliser alors la
question 1 pour montrer que $||x-y||^{2}\geq ||x-p(x)||^{2}$. Illustrer sur un dessin.

En déduire que $\mathrm{dist}(x,F)=||x-p(x)||$. 

\item\label{q: proj-en-coord}
Soit $(e_{1},\ldots,e_{r})$ une base orthonormée de $F$. Montrer que
$p(x)=\sum_{i=1}^{r}<x,e_{i}>\;e_{i}$.
\end{enumerate}

\item
\textbf{Espace de polynômes :}

Sur l'espace $E=\R_{3}[X]$, on considère la forme bilinéaire définie par :
$$
  <P,Q>=\frac{1}{2}\int_{-1}^{1}P(t)Q(t)dt.
$$
\begin{enumerate}
\item
Montrer qu'il s'agit d'un produit scalaire (on admet que l'intégrale sur $[-1,1]$ d'une
fonction $f$ continue et positive est nulle si et seulement si $f$ est nulle sur $[-1,1]$)

\item\label{q: Schmidt}
A l'aide du procédé de Schmidt appliqué à la base $(1,X,X^{2})$, construire une base
orthonormée de $\R_{2}[X]$ pour ce produit scalaire.
\item
On considère le polynôme $P_{0}=X^{3}$. Calculer la projection
orthogonale de $X^{3}$ sur $\R_{2}[X]$. En déduire que pour ce produit
scalaire, on a :
$$
  \mathrm{dist}(X^{3},\R^{2}[X])=\frac{2}{5\sqrt{7}}\ .
$$
\end{enumerate}

\end{enumerate}
\finenonce{001506}



\finexercice



\exercice{1507, barraud, 2003/09/01}

\enonce{001507}{}

Soit $(E,<,>)$ un espace euclidien, $x_{0}$ un point de $E$ et $F$ un
sous espace vectoriel de $E$. On note $\pi$ la projection orthogonale de
$E$ sur $F$. On rappelle que pour $x\in E$, $\pi(x)$ est caractérisé par
les relations~: 
$$
\pi(x)\in F\qquad\text{ et }\qquad x-\pi(x)\in F^{\bot}
$$

Le but de cette partie est de montrer que la projection orthogonale de
$x_{0}$ sur $F$ est le point de $F$ le plus proche de $x_{0}$.

\begin{enumerate}
\item

  En utilisant que $x_{0}-y=(x_{0}-\pi(x_{0}))+(\pi(x_{0})-y)$, montrer
  que 
  $$
  \Vert x_{0}-y\Vert^{2}=\Vert x_{0}-\pi(x_{0})\Vert^{2}+\Vert
  y-\pi(x_{0})\Vert^{2}.
  $$

\item

  En déduire que $\inf\limits_{y\in F}\Vert x_{0}-y\Vert^{2}=\Vert
  x_{0}-\pi(x_{0})\Vert^{2} $, c'est à dire que~:
  $$
  \forall y\in F,\ \Vert
  x_{0}-y\Vert^{2}\geq\Vert x_{0}-\pi(x_{0})\Vert^{2}
  $$
  A quelle condition a-t-on égalité dans la relation ci-dessus~?
\item

  Soit $(e_{1},\dots,e_{k})$ une base orthonormée de $F$. Montrer que
  $\pi(x_{0})=\sum_{i=1}^{k}<e_{i},x_{0}>e_{i}$
\item

  Déduire des deux questions précédentes que 
  $$
  \inf\limits_{y\in F}\Vert
  x_{0}-y\Vert^{2}=\Vert x_{0}-\sum_{i=1}^{k}<e_{i},x_{0}>e_{i}\Vert^{2}
  =\Vert x_{0}\Vert^{2}-\sum_{i=1}^{k}<e_{i},x_{0}>^{2}
  $$

\bigskip


{Application}~: Le but est maintenant de déterminer
$$
\alpha=\inf\limits_{a,b\in \R^{2}}\int_{-1}^{1} (e^{t}-at-b)^{2}dt.
$$ 

On considère à cet effet l'espace $F=\R_{1}[X]$, comme sous espace de
$E=F\oplus\R f_{0}$ où $f_{0}$ est la fonction définie par $f_{0}(t)=e^{t}$.
On admettra sans démonstration que $<f,g>=\int_{-1}^{1}f(t)g(t)dt$ est un
produit scalaire sur $E$.  
\item
Donner une base orthonormée $(P_{1},P_{2})$ de $\R_{1}[X]$ pour ce produit
scalaire.

\item

Calculer $<f_{0},P_{1}>$, $<f_{0},P_{2}>$, et $\Vert f_{0}\Vert^{2}$. En
déduire que 
$$
 \alpha=\frac{e^{2}-e^{-2}}{2}-(2e^{-1})^{2}-\Big(\frac{e-e^{-1}}{2}\Big)^{2}.
$$
\item

  Même question avec le calcul de $\alpha'=\inf\limits_{a,b\in
    \R^{2}}\int_{-1}^{1} (e^{t}-at^{2}-bt-c)^{2}dt. $~: commencer par
  chercher une base orthonormée de $\R_{2}[X]$ pour le même produit
  scalaire, et en déduire $\alpha'$.

\end{enumerate}
\finenonce{001507}



\finexercice



\exercice{1508, barraud, 2003/09/01}

\enonce{001508}{}
  A deux polynômes $P$ et $Q$ de $\R_{n}[X]$, on associe le nombre
  $$
  \phi(P,Q)=\int_{0}^{1}P'(t)Q'(t)dt+P(0)Q(0)
  $$
  \begin{enumerate}
  \item
    Montrer que $\phi$ est un produit scalaire sur $\R_{n}[X]$.
  \item
    Lorsque $n=2$, donner une base orthonormée pour ce produit scalaire.
\end{enumerate}
\finenonce{001508}



\finexercice



\exercice{5484, rouget, 2010/07/10}
\enonce{005484}{**IT}
Dans $\Rr^4$ muni du produit scalaire usuel, on pose~:~$V_1=(1,2,-1,1)$ et $V_2=(0,3,1,-1)$.
On pose $F=\mbox{Vect}(V_1,V_2)$. Déterminer une base orthonormale de $F$ et un système d'équations de $F^\bot$.
\finenonce{005484}


\finexercice


\exercice{5500, rouget, 2010/07/10}
\enonce{005500}{****I}
Sur $E=\Rr_n[X]$, on pose $P|Q=\int_{-1}^{1}P(t)Q(t)\;dt$.
\begin{enumerate} 
\item  Montrer que $(E,|)$ est un espace euclidien.
\item  Pour $p$ entier naturel compris entre $0$ et $n$, on pose $L_p=((X^2-1)^p)^{(p)}$.
Montrer que $\left(\frac{L_p}{||L_p||}\right)_{0\leq p\leq n}$ est l'orthonormalisée de \textsc{Schmidt} de la base canonique de $E$.

Déterminer $||Lp||$.
\end{enumerate}
\finenonce{005500}


\finexercice

\exercice{5772, rouget, 2010/10/16}
\enonce{005772}{*** I Polynômes de \textsc{Legendre}}
Soit $E =\Rr[X]$. On munit $E$ du produit scalaire $P|Q =\int_{-1}^{1}P(t)Q(t)\;dt$.

\begin{enumerate}
 \item  Pour $n\in\Nn$, on pose $L_n=\left((X^2-1)^n\right)^{(n)}$.
   \begin{enumerate}
   \item Montrer que la famille $(L_n)_{n\in\Nn}$ est une base orthogonale de l'espace préhilbertien $(E,\;|\;)$.
   \item Déterminer $\|L_n\|$ pour $n\in\Nn$.
   \end{enumerate}
\item  Déterminer l'orthonormalisée de \textsc{Schmidt} de la base canonique de $E$.
\end{enumerate}
\finenonce{005772}


\finexercice


\exercice{5779, rouget, 2010/10/16}
\enonce{005779}{***}
Soit $f$ une fonction continue sur $[0,1]$, non nulle à valeurs réelles positives. Pour $P$ et $Q$ polynômes donnés, on pose $\Phi(P,Q) = \int_{0}^{1}f(t)P(t)Q(t)\;dt$.

\begin{enumerate}
 \item  Montrer que $\Phi$ est un produit scalaire sur $\Rr[X]$.

\item  Montrer qu'il existe une base orthonormale $(P_n)_{n\in\Nn}$ pour $\Phi$ telle que, pour tout entier naturel $n$, $\text{deg}(P_n)= n$.

\item  Soit $(P_n)_{n\in\Nn}$ une telle base. Montrer que chaque polynôme $P_n$, $n\in\Nn^*$, a $n$ racines réelles simples.
\end{enumerate}
\finenonce{005779}


\finexercice


\exercice{5789, rouget, 2010/10/16}
\enonce{005789}{**I Inégalité de \textsc{Hadamard}}
\label{exo:sperou4}
Soit $E$ un espace euclidien de dimension $n\geqslant1$ et $\mathcal{B}$ une base orthonormée de $E$.

Montrer que pour tout $n$-uplet de vecteurs $(x_1,...x_n)$, on a : $\left|\text{det}_{\mathcal{B}}(x_1,...,x_n)\right|\leqslant\|x_1\|...\|x_n\|$. Cas d'égalité ?
\finenonce{005789}


\finexercice


\exercice{5790, rouget, 2010/10/16}
\enonce{005790}{**}
Montrer que pour toute matrice carrée $A$ réelle de format $n$, on a $|\text{det} A|\leqslant\sqrt{\prod_{j=1}^{n}\left(\sum_{i=1}^{n}a_{i,j}^2\right)}$.
\finenonce{005790}


\finexercice


\exercice{5806, rouget, 2010/10/16}
\enonce{005806}{**}
Rang et signature des formes quadratiques suivantes :

\begin{enumerate}
 \item  $Q((x,y,z)) = 2x^2-2y^2-6z^2+3xy-4xz+7yz$.

\item  $Q((x,y,z)) = 3x^2+3y^2+3z^2-2xy-2xz-2yz$

\item  $Q((x,y,z,t)) = xy + yz+zt+tx$.

\item  $Q((x,y,z,t)) = x^2+(4+\lambda)y^2+(1+4\lambda)z^2+\lambda t^2+4xy+2xz+4(1-\lambda)yz+2\lambda yt+(1-4\lambda)zt$.

\item  $Q((x_1,...,x_5)) =\sum_{1\leqslant i<j\leqslant n}^{}x_ix_j$.

\item  $Q((x_1,...,x_n)) =\sum_{1\leqslant i,j\leqslant n}^{}ijx_ix_j$.

\item  $Q((x_1,...,x_n)) =\sum_{1\leqslant i,j\leqslant n}^{i}x_ix_j$.

\item  $Q((x_1,...,x_n)) =\sum_{1\leqslant i,j\leqslant n}^{}\text{Inf}(i,j)x_ix_j$.
\end{enumerate}
\finenonce{005806}


\finexercice


\exercice{5811, rouget, 2010/10/16}
\enonce{005811}{**}
Sur $E =\Rr^2$ ou $\Rr^3$ muni de sa structure euclidienne usuelle, réduire en base orthonormée les formes quadratiques suivantes :

\begin{enumerate}
 \item  $Q((x,y)) = x^2+10xy+y^2$.

\item  $Q((x,y)) = 6x2+4xy+9y2$.

\item  $Q((x,y,z)) = 4x^2+9y^2-z^2+2\sqrt{6}xy +10\sqrt{2}xz+2\sqrt{3}yz$.
\end{enumerate}
\finenonce{005811}


\finexercice


\exercice{5812, rouget, 2010/10/16}
\enonce{005812}{***}
Soit $E =\Rr_n[X]$. Pour $P\in E$, on pose $Q(P) =\sum_{k=0}^{+\infty}P(k)P(-k)e^{-k}$.

\begin{enumerate}
 \item  Montrer que $Q$ est une forme quadratique sur $E$.

\item  Déterminer sa signature.
\end{enumerate}
\finenonce{005812}


\finexercice

\section{204.06 Espace vectoriel euclidien de dimension 3}


\exercice{3694, quercia, 2010/03/11}
\enonce{003694}{Propriétés du produit vectoriel}

Soient $\vec u,\vec v,\vec w,\vec t$ quatre vecteurs d'un espace vectoriel euclidien orienté de
dimension 3. Démontrer :

\begin{align*}
%
  &(\vec u\wedge\vec v) \mid (\vec w\wedge\vec t)
    = (\vec u\mid\vec w)(\vec v\mid\vec t)
    - (\vec u\mid\vec t)(\vec v\mid\vec w) \\
%
  &(\vec u\wedge\vec v) \wedge (\vec w\wedge\vec t)
    = -[\vec u,\vec v,\vec w]\vec t
    + [\vec u,\vec v,\vec t]\vec w \\
%
  &   [\vec t,\vec v,\vec w]\vec u
    + [\vec u,\vec t,\vec w]\vec v
    + [\vec u,\vec v,\vec t]\vec w
    = [\vec u,\vec v,\vec w]\vec t. \\\end{align*}
\finenonce{003694}



\finexercice


\exercice{3695, quercia, 2010/03/11}
\enonce{003695}{Division vectorielle}

Soit $E$ un espace vectoriel euclidien orienté de dimension 3 et $\vec a,\vec b$ deux vecteurs
donnés, $\vec a \ne \vec 0$.
\'Etudier l'équation : $\vec a\wedge\vec x = \vec b$.
On cherchera une solution particulière de la forme
$\vec x = \vec a\wedge\vec y$.

\finenonce{003695}



\finexercice


\exercice{3696, quercia, 2010/03/11}
\enonce{003696}{$a\wedge b$, $b\wedge c$, $c\wedge a$ donnés}

Soit $E$ un espace vectoriel euclidien orienté de dimension 3.\par
Trouver $\vec a,\vec b,\vec c$ connaissant
$\vec u = \vec a\wedge\vec b$,
$\vec v = \vec b\wedge\vec c$
et $\vec w = \vec c\wedge\vec a$
(calculer $\vec u\wedge\vec v$).

\finenonce{003696}



\finexercice


\exercice{3697, quercia, 2010/03/11}
\enonce{003697}{$f(u)\wedge v + u\wedge f(v) = g(u\wedge v)$}

Soit $E$ un espace vectoriel euclidien orienté de dimension 3\ et $f \in \mathcal{L}(E)$.

\begin{enumerate}
  \item  Prouver que :
     $[f(\vec u),\vec v,\vec w]
     + [\vec u,f(\vec v),\vec w]
     + [\vec u,\vec v,f(\vec w)] = [\vec u,\vec v,\vec w]\mathrm{tr}(f)$.

  \item  Montrer qu'il existe $g \in \mathcal{L}(E)$ telle que :
     $\forall\ \vec u,\vec v$, on a
     $f(\vec u)\wedge\vec v + \vec u\wedge f(\vec v) = g(\vec u\wedge\vec v)$.

  \item  Dans une {\it bond}, exprimer la matrice de $g$ en fonction de celle
     de $f$.
\end{enumerate}
\finenonce{003697}



\finexercice


\exercice{3698, quercia, 2010/03/11}
\enonce{003698}{Applications bilinéaires antisymétriques}

Soit $E$ un espace vectoriel euclidien orienté de dimension 3 et $\phi : {E\times E} \to \R$
une application bilinéaire antisymétrique.

Montrer qu'il existe $f \in E^*$ unique telle que :
$\forall\ \vec x,\vec y,\ \phi(\vec x,\vec y) = f(\vec x\wedge\vec y)$.

\finenonce{003698}



\finexercice


\exercice{3699, quercia, 2010/03/11}
\enonce{003699}{Volume d'un parallélépipède}

Soient $\vec u,\vec v,\vec w$ trois vecteurs d'un espace vectoriel euclidien orienté de
dimension 3.

On donne $\|\vec u\| = a$, $\|\vec v\| = b$, $\|\vec w\| = c$,
$\overline{(\vec u,\vec v)} \equiv \alpha$,
$\overline{(\vec v,\vec w)} \equiv \beta$,
$\overline{(\vec w,\vec u)} \equiv \gamma$.

Quel est le volume du parallélépipède construit sur $\vec u,\vec v,\vec w$ ?

\finenonce{003699}



\finexercice


\exercice{3700, quercia, 2010/03/11}
\enonce{003700}{Applications conservant le produit vectoriel}

Soit $E$ un espace vectoriel euclidien orienté de dimension 3.
Trouver les applications $f \in \mathcal{L}(E)$ vérifiant :

\begin{enumerate}
  \item   $f(\vec u\wedge\vec v) = f(\vec u)\wedge f(\vec v)$.
   \item $f(\vec u\wedge\vec v) =-f(\vec u)\wedge f(\vec v)$.
\end{enumerate}
\finenonce{003700}



\finexercice


\exercice{3701, quercia, 2010/03/11}
\enonce{003701}{Expression d'une rotation}

Soit $E$ un espace vectoriel euclidien orienté de dimension 3, $\vec u \in E$ unitaire,
$\alpha \in \R$ et $f$ la rotation autour de $\vec u$ d'angle de mesure $\alpha$.
\begin{enumerate}
  \item Exprimer $f(\vec x)$ en fonction de $\vec u$, $\vec x$ et $\alpha$.
    
  \item On donne les coordonnées de $\vec u$ dans une base orthonormée : $a,b,c$.
    Calculer la matrice de $f$ dans cette base.
    
\end{enumerate}
\finenonce{003701}



\finexercice


\exercice{3702, quercia, 2010/03/11}
\enonce{003702}{Conjuguée d'une rotation}

Soit $\rho$ une rotation d'un espace vectoriel euclidien orienté de dimension 3,
et $f \in {\cal O}(E)$. Reconnaître $f \circ \rho \circ f^{-1}$.

Application : Déterminer le centre de ${\cal O}^+(E)$.
\finenonce{003702}



\finexercice


\exercice{3703, quercia, 2010/03/11}
\enonce{003703}{Conjugaison dans ${\cal O}(\R^3)$}


Soient $f,g \in {\cal O}(\R^3)$ ayant même polynôme caractéristique.

Montrer qu'il existe $h \in {\cal O}(\R^3)$ tel que $f = h^{-1}\circ g \circ h$.

Si $f$ et $g$ sont positifs, a-t-on $h$ positif ?

\finenonce{003703}



\finexercice


\exercice{3704, quercia, 2010/03/11}
\enonce{003704}{Décomposition des rotations}

Soit $(\vec i,\vec j,\vec k)$ une {\it bond} d'un espace vectoriel euclidien orienté de
dimension 3\ $E$, et $f \in {\cal O}^+(E)$.

\begin{enumerate}
  \item On suppose $f(\vec j) \perp \vec i$.
    Montrer qu'il existe $r,r'$ rotations autour de
    $\vec j$ et $\vec i$ telles que $r' \circ r = f$.

  \item En déduire que tout $f \in {\cal O}^+(E)$ se décompose de deux manières
    sous la forme : $f = r'' \circ r' \circ r$ où $r,r''$ sont des rotations
    autour de $\vec j$ et $r'$ est une rotation autour de $\vec i$.

  \item Décomposer $(x,y,z) \mapsto (y,x,z)$ et
    $(x,y,z) \mapsto (x\cos\alpha-y\sin\alpha, x\sin\alpha+y\cos\alpha, z)$.
\end{enumerate}
\finenonce{003704}



\finexercice


\exercice{3705, quercia, 2010/03/11}
\enonce{003705}{Sous-groupes finis de ${\cal O}^+(3)$}

Déterminer les sous-groupes de ${\cal O}^+(3)$ de cardinal 2,3, ou 4.
\finenonce{003705}



\finexercice


\exercice{3706, quercia, 2010/03/11}
\enonce{003706}{Applications antisymétriques}

Soit $E$ un espace vectoriel euclidien et $f \in \mathcal{L}(E)$ antisymétrique.

\begin{enumerate}
  \item Montrer que $\mathrm{id}_E + f \in GL(E)$.
  \item Montrer que $g = (\mathrm{id} - f) \circ (\mathrm{id} + f)^{-1} \in {\cal O}^+(E)$ et
    $\mathrm{id}+g$ est inversible.
    

  \item Réciproquement, soit $h \in {\cal O}^+(E)$ tq $\mathrm{id} + h$ soit inversible.
    Montrer qu'il existe $f$ antisymétrique tel que
    $h = (\mathrm{id} - f) \circ (\mathrm{id} + f)^{-1}$.
    
\end{enumerate}
\finenonce{003706}



\finexercice


\exercice{3707, quercia, 2010/03/11}
\enonce{003707}{Applications antisymétriques}
Soit $E$ un espace vectoriel euclidien orienté de dimension 3, ${\cal B}$ une base orthonormée
directe de $E$ et $f \in \mathcal{L}(E)$ de matrice dans $\cal B$ :
$M = \begin{pmatrix} 0      &-\gamma &\beta   \cr
               \gamma &0       &-\alpha \cr
               -\beta &\alpha  &0       \cr\end{pmatrix}$.

\begin{enumerate}
  \item Reconaître $f$.
    
  \item Montrer que $\mathrm{id}_E + f$ est une bijection et calculer la bijection
    réciproque.
    
  \item Montrer que $g = (\mathrm{id} - f) \circ (\mathrm{id} + f)^{-1}$ est une rotation et
    préciser son axe et son angle.
    
\end{enumerate}
\finenonce{003707}



\finexercice


\exercice{3708, quercia, 2010/03/11}
\enonce{003708}{Exponentielle d'une application antisymétrique}

Soit $E$ un espace vectoriel euclidien orienté de dimension 3\ et
$\vec a \in E\setminus\{\vec 0\}$. On note $\alpha = \|\vec a\,\|$.
Soit $f$ l'endomorphisme de $E$ défini par : $f(\vec x) = \vec a\wedge\vec x$.

\begin{enumerate}
  \item Vérifier que $f^3 = -\alpha^2f$.
  \item On pose $g(\vec x) = \sum_{k=0}^\infty \frac{f^k(\vec x)}{k!}$.
    Simplifier $g(\vec x)$ et en déduire que $g$ est une rotation.
    
\end{enumerate}
\finenonce{003708}



\finexercice


\exercice{3709, quercia, 2010/03/11}
\enonce{003709}{Matrice à trou}

\begin{enumerate}
  \item Compléter la matrice $A = \frac 17 \begin{pmatrix} 6 &3 &. \cr
                                                 -2 &6 &. \cr
                                                  3 &. &. \cr\end{pmatrix}$
    en une matrice orthogonale positive.
    

  \item Reconnaître l'application de matrice $A$ dans la base canonique
    de $\R^3$.
\end{enumerate}
\finenonce{003709}



\finexercice


\exercice{3710, quercia, 2010/03/11}
% modification le 2016-05-01 par D.M. : 
% ajout de la condition c), et des questions 2) et 3)

\enonce{003710}{Matrice circulante}

Pour tout triplet $(a,b,c)$  de réels on pose 
\[M(a,b,c) = \begin{pmatrix}a &b &c \cr c &a &b \cr b &c &a \cr\end{pmatrix}.\]
\begin{enumerate}
\item Montrer que les assertions suivantes sont équivalentes:
\begin{enumerate}
\item
$M(a,b,c)$ est une matrice de rotation;
\item $(a,b,c)$ appartient au cercle $\mathcal C \subset \R^3$ défini par l'intersection de la sphère unité de $\R^3$ et du plan $x+y+z-1=0$;
\item $a$, $b$ et $c$ sont les racines d'un polynôme de la forme $P = X^3-X^2+\lambda$ avec
$\lambda \in \left[0,\frac4{27}\right]$.
\end{enumerate}
\item Si $M(a,b,c)$ est une rotation, calculer son axe et son angle au signe près.
\item Montrer que l'ensemble des matrices de rotation de la forme $M(a,b,c)$ est un sous-groupe de $SO(3)$. À quel groupe connu est-il isomorphe ?
\end{enumerate}

\finenonce{003710}



\finexercice

\exercice{3711, quercia, 2010/03/11}
\enonce{003711}{Expressions analytiques}

Reconnaître les endomorphismes de $\R^3$ définis par les expressions
analytiques dans la base canonique :

\begin{enumerate}
  \item $\begin{cases} 3x' = 2x+2y+z \cr 3y' = -2x+y+2z \cr 3z' = x-2y+2z \cr\end{cases}$
    
  \item $\begin{cases} 9x' = 8x+y-4z \cr 9y' = -4x+4y-7z \cr 9z'=x+8y+4z \cr\end{cases}$

  \item $\begin{cases} 3x' = -2x+2y-z \cr 3y' = 2x+y-2z \cr 3z' = -x-2y-2z \cr\end{cases}$

  \item $\begin{cases} 4x' = -2x - y\sqrt6 + z\sqrt6 \cr
             4y' = x\sqrt6 + y + 3z \cr
             4z' = -x\sqrt6 + 3y + z \cr \end{cases}$  

  \item $\begin{cases}x' = \frac x{\sqrt3} + \frac y{\sqrt2} - \frac z{\sqrt6} \cr
            y' = \frac x{\sqrt3}  + \frac {2z}{\sqrt6} \cr
            z' = \frac x{\sqrt3} - \frac y{\sqrt2} - \frac z{\sqrt6} \cr\end{cases}$   

  \item $\begin{cases} 3x' = x+2y+2z \cr 3y' = 2x+y-2z \cr 3z' = 2x-2y+z \cr\end{cases}$  

  \item $\begin{cases} 7x' = -2x+6y-3z \cr 7y' = 6x+3y+2z \cr 7z' = -3x+2y+6z \cr\end{cases}$   

  \item $\begin{cases} 3x' = 2x-2y+z \cr 3y' = -2x-y+2z \cr 3z' = x+2y+2z \cr\end{cases}$    

  \item $\begin{cases} 3x' = 2x+y+2z \cr 3y' = 2x-2y-z \cr 3z' = -x-2y+2z \cr\end{cases}$
    
  \item $\begin{cases} 4x' = -x+3y-z\sqrt6 \cr
             4y' = 3x-y-z\sqrt6  \cr
             4z' = x\sqrt6+y\sqrt6+2z \cr\end{cases}$

  \item $\begin{cases} 15x'=5x-10z \cr 15y' = -8x+5y+6z \cr 15z' = 6x-10y+8z \cr\end{cases}$
   \label{rotproj}
\end{enumerate}
\finenonce{003711}


\finexercice


\exercice{3712, quercia, 2010/03/11}
\enonce{003712}{Ensi Physique 92}
Déterminer la matrice de la rotation $\cal R$ de $\R^3$ dans une base
orthonormée $(\vec i,\vec j,\vec k\,)$ telle que ${\cal R}(\vec u)=\vec u$
avec $\vec u\bigl(\frac1{\sqrt3},\frac{-1}{\sqrt3},\frac1{\sqrt3}\bigr)$
et ${\cal R}(\vec i) = \vec k$. Donner son angle de rotation.
\finenonce{003712}


\finexercice


\exercice{5491, rouget, 2010/07/10}
\enonce{005491}{**}
$\mathcal{B}$ est une base orthonormée directe de $\Rr^3$ donnée. Montrer que $\mbox{det}_{\mathcal{B}}(u\wedge v,v\wedge w,w\wedge u)=(\mbox{det}_{\mathcal{B}}(u,v,w))^2$ pour tous vecteurs $u$, $v$ et $w$.
\finenonce{005491}


\finexercice

\exercice{5493, rouget, 2010/07/10}
\enonce{005493}{**}
Montrer que $u\wedge v|w\wedge s=(u|w)(v|s)-(u|s)(v|w)$ et $(u\wedge v)\wedge(w\wedge s)=[u,v,s]w-[u,v,w]s$.
\finenonce{005493}


\finexercice


\exercice{5793, rouget, 2010/10/16}
\enonce{005793}{**}
Soit $(e_1,e_2,e_3)$ une base orthonormée directe d'un espace euclidien orienté $E$ de dimension $3$. Matrice de la rotation d'angle $\frac{\pi}{3}$ autour de $e_1+e_2$.
\finenonce{005793}


\finexercice


\exercice{5797, rouget, 2010/10/16}
\enonce{005797}{***}
Trouver tous les endomorphismes de $\Rr^3$ vérifiant $\forall(x,y)\in(\Rr^3)^2$, $f(x\wedge y) = f(x)\wedge f(y)$.
\finenonce{005797}


\finexercice


\exercice{5801, rouget, 2010/10/16}
\enonce{005801}{**}
Valeurs et vecteurs propres de l'endomorphisme de $\Rr^3$ euclidien orienté défini par

\begin{center}
$\forall x\in\Rr^3$, $f(x) = a\wedge(a\wedge x)$ où $a$ est un vecteur donné.
\end{center}
\finenonce{005801}


\finexercice


\exercice{5803, rouget, 2010/10/16}
\enonce{005803}{** I}
Soit $P$ le plan de $\Rr^4$ d'équations $\left\{
\begin{array}{l}
x+y+z+t=0\\
x+y-2z-t=0
\end{array}
\right.$ dans une base orthonormée $\mathcal{B}$ de $\Rr^4$ muni de sa structure euclidienne canonique.

\begin{enumerate}
 \item Déterminer les matrices dans $\mathcal{B}$ de la projection orthogonale sur $P$ et de la symétrie orthogonale par rapport à $P$.

\item  Calculer la distance d'un vecteur quelconque de $\Rr^4$ à $P$.
\end{enumerate}
\finenonce{005803}


\finexercice

\section{204.07 Endomorphismes auto-adjoints}


\exercice{3762, quercia, 2010/03/11}
\enonce{003762}{Esem 91}
    Soit $A \in \mathcal{M}_n(\R)$. On suppose ${}^tA = A$ et $A^2 = 0$. Montrer que $A=0$.
\finenonce{003762}



\finexercice


\exercice{3763, quercia, 2010/03/11}
\enonce{003763}{Comatrice d'une matrice symétrique}

Soit $M \in \mathcal{M}_n(\R)$ symétrique. Montrer que $\text{com}\,M$ est aussi symétrique.
La réciproque est-elle vraie ?
\finenonce{003763}



\finexercice


\exercice{3764, quercia, 2010/03/11}
\enonce{003764}{Base non orthonormée}

Soit ${\cal B} = (\vec e_1,\dots,\vec e_n)$ une base non orthonormée de $E$,
$G$ la matrice de Gram des $\vec e_i$, $f \in \mathcal{L}(E)$ et $M$ sa matrice dans
$\cal B$.

\begin{enumerate}
  \item Montrer que $f$ est auto-adjoint si et seulement si ${}^tMG = GM$.
  \item Montrer que $f$ est orthogonal si et seulement si ${}^tMGM = G$.
\end{enumerate}
\finenonce{003764}



\finexercice


\exercice{3765, quercia, 2010/03/11}
\enonce{003765}{autoadjoint $ \Rightarrow $ linéaire}

Soit $u : E \to E$ telle que :
$\forall\ \vec x,\vec y \in E,\ (u(\vec x)\mid\vec y) = (\vec x\mid u(\vec y))$.
Montrer que $u$ est linéaire.

\finenonce{003765}



\finexercice


\exercice{3766, quercia, 2010/03/11}
\enonce{003766}{Diagonalisation de matrices symétriques}

Diagonaliser dans une base orthonormée :


\begin{enumerate}
  \item $A = \begin{pmatrix} 6 &-2 &2 \cr -2 &5 &0 \cr 2 &0 &7 \cr \end{pmatrix}$.
    

  \item $A = \frac19\begin{pmatrix} 23 &2 &-4 \cr 2 &26 &2 \cr -4 &2 &23 \cr \end{pmatrix}$.
    
\end{enumerate}
\finenonce{003766}



\finexercice


\exercice{3767, quercia, 2010/03/11}
\enonce{003767}{Diagonalisation de $C{}^tC$}

Soient $a_1,\dots,a_n \in \R$ et $M = (a_ia_j) \in \mathcal{M}_n(\R)$.
Montrer que $M$ est diagonalisable et déterminer ses éléments propres.


\finenonce{003767}



\finexercice


\exercice{3768, quercia, 2010/03/11}
\enonce{003768}{Décomposition en projections orthogonales}

Soit $\varphi$ l'endomorphisme de matrice dans la base canonique de $\R^4$ :
$M = \begin{pmatrix} 2 &0 &0 &3 \cr 0 &2 &3 &0 \cr 0 &3 &2 &0 \cr 3 &0 &0 &2 \cr\end{pmatrix}$.

Montrer qu'il existe des projections orthogonales $p$, $q$ et des réels
$\lambda$, $\mu$ tels que :
$$\varphi = \lambda p + \mu q,\qquad p\circ q = 0,  \qquad p+q = \mathrm{id}_E.$$


\finenonce{003768}

\finexercice


\exercice{3769, quercia, 2010/03/11}
\enonce{003769}{Ensi Physique 93}

Soit $E = \R_n[x]$. On pose pour $P,Q \in E :
{<P,Q>} =  \int_{-1}^1 P(t)Q(t)\,d t$ et on considère
$$u : E \to  {\R[x]}, {P(x)}  \mapsto {2xP'(x) + (x^2-1)P''(x).}$$

\begin{enumerate}
  \item Montrer que l'on définit un produit scalaire et que $u$ est un endomorphisme.
  \item Montrer que $u$ est diagonalisable et que si $P_k,P_\ell$ sont des vecteurs propres
    de valeurs propres distinctes, alors ${{<}P_k,P_\ell{>=} 0}$.
  \item \'Eléments propres de $u$ pour $n=3$ ?
    
\end{enumerate}
\finenonce{003769}



\finexercice


\exercice{3770, quercia, 2010/03/11}
\enonce{003770}{$(X^2-1)P'' + (2X+1)P'$}

Pour $P,Q \in \R_n[X]$ on pose
$(P\mid Q) =  \int_{t=-1}^1 \sqrt{\frac{1-t}{1+t}}P(t)Q(t)\,d t$
et $\Phi(P) = (X^2-1)P'' + (2X+1)P'$.

\begin{enumerate}
  \item Vérifier que $(P\mid Q)$ existe et qu'on définit ainsi un produit scalaire
    sur $\R_n[X]$.

  \item Montrer que pour ce produit scalaire, $\Phi$ est auto-adjoint
    $\Bigl($calculer $ \int_{t=-1}^1 (1-t)^{3/2}(1+t)^{1/2}P''(t)Q(t)\,d t$
    par parties$\Bigr)$.
  \item Déterminer les valeurs propres de $\Phi$ et montrer qu'il existe une
    base propre de degrés étagés.
    


\end{enumerate}
\finenonce{003770}



\finexercice


\exercice{3771, quercia, 2010/03/11}
\enonce{003771}{$\mathrm{Ker} u + \Im u = E$}

Soit $u \in \mathcal{L}(E)$ auto-adjoint.
Montrer que $\mathrm{Ker} u \mathop{\oplus}\limits^\perp \Im u = E$.
\finenonce{003771}



\finexercice


\exercice{3772, quercia, 2010/03/11}
\enonce{003772}{$u\circ v$ autoadjoint ?}

Soient $u,v \in \mathcal{L}(E)$ auto-adjoints. Montrer que $u\circ v$ est auto-adjoint
si et seulement si $u\circ v = v\circ u$.
\finenonce{003772}



\finexercice


\exercice{3773, quercia, 2010/03/11}
\enonce{003773}{Composée de projecteurs}

Soient $p,q$ deux projecteurs orthogonaux.

\begin{enumerate}
  \item Montrer que $p\circ q \circ p$ est auto-adjoint.
  \item Montrer que
    $(\Im p + \mathrm{Ker} q) \mathop{\oplus}\limits^\perp (\mathrm{Ker} p \cap \Im q)  = E$.
  \item En déduire que $p\circ q$ est diagonalisable.
    
\end{enumerate}
\finenonce{003773}



\finexercice


\exercice{3774, quercia, 2010/03/11}
\enonce{003774}{Autoadjoint et orthogonal}

Quels sont les endomorphismes de $E$ à la fois auto-adjoints et orthogonaux ?
\finenonce{003774}



\finexercice


\exercice{3775, quercia, 2010/03/11}
\enonce{003775}{Spectre et rang d'une matrice antisymétrique}

Soit $M \in \mathcal{M}_n(\R)$ antisymétrique et $f$ l'endomorphisme de $\R^n$
canoniquement associé à $M$.

\begin{enumerate}
  \item Montrer que les valeurs propres de $M$ sont imaginaires pures.
  \item Montrer que $\mathrm{Ker} f \perp \Im f$. En déduire que $g = f_{|\Im f}$ est un isomorphisme de $\Im f$.
  \item Montrer que $g^2$ est diagonalisable. En déduire que $\mathrm{rg}(M)$ est pair.
\end{enumerate}
\finenonce{003775}



\finexercice


\exercice{3776, quercia, 2010/03/11}
\enonce{003776}{Racine carrée}

Soit $A \in \mathcal{M}_n(\R)$ symétrique définie positive.
Montrer qu'il existe une unique matrice $B \in \mathcal{M}_n(\R)$ symétrique définie
positive telle que $B^2 = A$.

Calculer $B$ lorsque $A = \begin{pmatrix} 1 &2 \cr 2 &5 \cr\end{pmatrix}$.
\finenonce{003776}



\finexercice


\exercice{3777, quercia, 2010/03/11}
\enonce{003777}{$A = {}^tBB$}

Soit $A \in \mathcal{M}_n(\R)$. Montrer que $A$ est symétrique définie positive si et
seulement s'il existe $B \in GL_n(\R)$ telle que $A = {}^tBB$.
\finenonce{003777}



\finexercice


\exercice{3778, quercia, 2010/03/11}
\enonce{003778}{Mineurs principaux positifs}

Soit $A \in \mathcal{M}_n(\R)$ symétrique.
Pour $1 \le p \le n$, on note $\Delta_p$ le déterminant
de la sous-matrice $(a_{ij})_{1\le i\le p ; 1\le j\le p}$.

\begin{enumerate}
  \item Montrer que si $A$ est définie positive, alors tous les déterminants
    $\Delta_p$ sont strictement positifs.
  \item Réciproque : on suppose $\Delta_1 > 0, \dots, \Delta_n > 0$.
    Montrer qu'il existe
    une matrice $B$ triangulaire supérieure inversible telle $A = {}^tBB$.
    En déduire que $A$ est définie positive.
    
\end{enumerate}
\finenonce{003778}



\finexercice


\exercice{3779, quercia, 2010/03/11}
\enonce{003779}{$q$ positive $ \Rightarrow  q(x) = \|u(x)\|^2$}

Soit $E$ un espace euclidien et $q$ une forme quadratique positive.
Montrer qu'il existe un endomorphisme $u$ auto-adjoint tel que :
$\forall\ \vec x \in E$, $q(\vec x) = \|u(\vec x)\|^2$.
\finenonce{003779}



\finexercice


\exercice{3780, quercia, 2010/03/11}
\enonce{003780}{$A$ symétrique et $A^k = I$}

Soit $A \in \mathcal{M}_n(\R)$ symétrique telle qu'il existe $k \in \N^*$ tel que
$A^k = I$. Montrer que $A^2 = I$.
\finenonce{003780}



\finexercice


\exercice{3781, quercia, 2010/03/11}
\enonce{003781}{$\sum_{i,j} a_{ij}^2$}

Soit $A = (a_{ij}) \in \mathcal{M}_n(\R)$ symétrique de valeurs propres
$\lambda_1,\dots,\lambda_n$.
Montrer que $\sum_{i,j} a_{ij}^2 = \sum_i \lambda_i^2$.
\finenonce{003781}



\finexercice


\exercice{3782, quercia, 2010/03/11}
\enonce{003782}{$u$ autoadjoint et tr$(u) = 0$}

Soit $u \in \mathcal{L}(E)$ auto-adjoint tel que $\mathrm{tr}(u) = 0$.

\begin{enumerate}
  \item Montrer qu'il existe un vecteur $\vec x$ non nul tel que
    $u(\vec x) \perp \vec x$.
    

  \item En déduire qu'il existe une base orthonormée $(\vec e_i)$ telle que :
    $\forall\ i,\ (u(\vec e_i) \mid \vec e_i) = 0$.
    

\end{enumerate}
\finenonce{003782}



\finexercice


\exercice{3783, quercia, 2010/03/11}
\enonce{003783}{Matrices symétriques commutant}

Soit $(A_i)$ une famille de matrices $n\times n$ réelles symétriques
commutant deux à deux.

Montrer qu'il existe une matrice symétrique $A$ et des polynômes $P_i$ tels que :
$\forall\ i,\ A_i = P_i(A)$.
\finenonce{003783}



\finexercice


\exercice{3784, quercia, 2010/03/11}
\enonce{003784}{Valeurs propres de $AB$}

Soient $A,B \in \mathcal{M}_n(\R)$ symétriques, $B$ définie positive.
Montrer que les valeurs propres de $AB$ sont réelles.

\finenonce{003784}



\finexercice


\exercice{3785, quercia, 2010/03/11}
\enonce{003785}{$\mathrm{tr}(AB) \le \mathrm{tr}(A)\mathrm{tr}(B)$}

Soient $A,B \in \mathcal{M}_n(\R)$ symétriques positives.
Montrer que $0 \le \mathrm{tr}(AB) \le \mathrm{tr}(A)\mathrm{tr}(B)$.

\finenonce{003785}



\finexercice


\exercice{3786, quercia, 2010/03/11}
\enonce{003786}{$\det(A+B) \ge \det(A) + \det(B)$}

Soient $A,B \in \mathcal{M}_n(\R)$ symétriques définies positives.
Montrer que $\det(A+B) \ge \det(A) + \det(B)$.

\finenonce{003786}



\finexercice


\exercice{3787, quercia, 2010/03/11}
\enonce{003787}{$f$ quelconque, il existe une BON dont l'image est orthogonale}

Soit $f \in \mathcal{L}(E)$. Montrer qu'il existe une base orthonormée
$(\vec e_1,\dots,\vec e_n)$ dont l'image par $f$ est une famille orthogonale.

\finenonce{003787}



\finexercice


\exercice{3788, quercia, 2010/03/11}
\enonce{003788}{Quotients de Rayleigh}

Soit $f \in \mathcal{L}(E)$ auto-adjoint et
$\lambda_1 \le \lambda_2 \le \dots \le \lambda_n$ ses valeurs propres.

\begin{enumerate}
  \item Montrer que : $\forall\ \vec x \in E,\
                       \lambda_1 \|\vec x\,\|^2
                   \le (f(\vec x)\mid\vec x)
                   \le \lambda_n \|\vec x\,\|^2$.
  \item Montrer que si l'une de ces deux inégalités est une égalité pour un vecteur
    $\vec x\ne\vec 0$, alors $\vec x$ est vecteur propre de~$f$.
  \item Soit $(\vec e_1,\dots,\vec e_n)$ une base orthonormée de $E$ telle que
    pour tout $i$ : $(f(\vec e_i)\mid\vec e_i) = \lambda_i$.

    Montrer que : $\forall\ i,\ f(\vec e_i) = \lambda_i\vec e_i$.
\end{enumerate}
\finenonce{003788}



\finexercice


\exercice{3789, quercia, 2010/03/11}
\enonce{003789}{Spec$(A+B)$}

Soient $A,B \in \mathcal{M}_n(\R)$ symétriques, $\lambda,\lambda'$ leurs plus petites
valeurs propres et $\mu,\mu'$ leurs plus grandes valeurs propres.
Montrer que toute valeur propre de $A+B$ est comprise entre $\lambda+\lambda'$
et $\mu+\mu'$.
\finenonce{003789}



\finexercice


\exercice{3790, quercia, 2010/03/11}
\enonce{003790}{Comparaison de valeurs propres}

Soient $h \in \mathcal{L}(E)$ autoadjoint, $\vec x_0 \in E$ unitaire, $p$ la projection
orthogonale sur $\text{vect}(\vec x_0)$, et $f = h+p$.

On note $\lambda_1 \le \dots \le \lambda_n$ les valeurs propres de $h$ et
$\mu_1 \le \dots \le \mu_n$ celles de $f$.

Montrer que $\lambda_1 \le \mu_1 \le \dots \le \lambda_n \le \mu_n$.


\finenonce{003790}



\finexercice


\exercice{3791, quercia, 2010/03/11}
\enonce{003791}{Mines P' 1996}

Soit $E$ un espace euclidien et $f\in\mathcal{L}(E)$.

\begin{enumerate}
  \item Montrer~: $\mathrm{Ker} f = \Im f  \Rightarrow  f + f^* \in GL(E)$.


  \item Montrer la réciproque lorsque l'on a $f^2 = 0$.

\end{enumerate}
\finenonce{003791}



\finexercice


\exercice{3792, quercia, 2010/03/11}
\enonce{003792}{Rayon spectral}

Soit $f\in\mathcal{L}(E)$.
Montrer que $\|f\|^2 =\max\{|\lambda| \text{ tq } \lambda\in\mathrm{Sp}(f^*\circ f)\}$.
\finenonce{003792}



\finexercice


\exercice{3793, quercia, 2010/03/11}
\enonce{003793}{Décomposition polaire d'un endomorphisme}

Soit $f\in\mathcal{L}(E)$.

\begin{enumerate}
  \item En considérant l'endomorphisme $f^*\circ f$, montrer que si~$f$
    est inversible alors $f$ se décompose de manière unique sous
    la forme $f = u\circ h$ avec $u$ unitaire et $h$ hermitien positif.

  \item Si $f$ est non inversible, montrer qu'une telle décomposition
    existe mais n'est pas unique (on rappelle que $U(E)$ est compact).

  \item Montrer que l'application $f  \mapsto (u,h)$ est continue sur $GL(E)$.
\end{enumerate}
\finenonce{003793}



\finexercice


\exercice{3794, quercia, 2010/03/11}
\enonce{003794}{Endomorphismes normaux}

Soit $E$ un espace vectoriel hermitien.
Un endomorphisme $u\in\mathcal{L}(E)$ est dit normal si $u$ et~$u^*$ commutent.

\begin{enumerate}
  \item Soit $u$ normal, montrer que si $F$ est un sous-espace propre de~$u$ alors $F^\perp$
    est stable par~$u$. En déduire que $u$ est diagonalisable en
    base orthonormale. La réciproque est-elle vraie~?

  \item Soit $u\in\mathcal{L}(E)$. Montrer l'équivalence entre les propriétés suivantes~:

{\obeylines%
(1) $u$ est normal.
(2) $\forall\ x\in E, \|u(x)\|=\|u^*(x)\|$.
(3) Tout sous-espace vectoriel stable par~$u$ est stable par~$u^*$.
(4) Si un sous-espace vectoriel $F$ est stable par~$u$ alors $F^\perp$ est stable par $u$.
(5) Il existe $P\in\C[X]$ tel que $u^* = P(u)$.
\par}
\end{enumerate}
\finenonce{003794}



\finexercice


\exercice{3795, quercia, 2010/03/11}
\enonce{003795}{$\|u(x)\| = \|v(x)\|$}

Soit $E$ un espace hermitien non nul et $u,v\in\mathcal{L}(E)$.
Montrer l'équivalence~:
$$\Bigl(\forall\ x\in E,\ \|u(x)\|=\|v(x)\|\Bigr)\Leftrightarrow
  \Bigl(\exists\ w\in\ U(E)\text{ tq } u = w\circ v\Bigr).$$

\finenonce{003795}



\finexercice


\exercice{3796, quercia, 2010/03/11}
\enonce{003796}{$(u(x)\mid x)$ est réel}

Soit $E$ un espace vectoriel hermitien et $u\in\mathcal{L}(E)$.

\begin{enumerate}
  \item Montrer que $u=u^*$ si et seulement si pour tout~$x\in E$,
    $(u(x)\mid x)$ est réel.


  \item On suppose $u$ autoadjoint positif. Montrer~:
    $\forall\ x\in E,\ \|u(x)\|^4 \le (x\mid u(x))\times(u(x)\mid u^2(x))$.


\end{enumerate}
\finenonce{003796}



\finexercice


\exercice{3797, quercia, 2010/03/11}
\enonce{003797}{Série d'autoadjoints positifs}

Soit $H$ un espace de Hilbert et $(u_n)$ une suite d'endomorphismes
de~$H$ autoadjoints positifs continus telle que la suite
$(u_0+\dots+u_n)$ est bornée dans~${\cal L}_c(H)$.
Montrer que pour tout~$x\in H$ la série $\sum_{n=0}^\infty u_n(x)$ est
convergente.



\finenonce{003797}



\finexercice


\exercice{3798, quercia, 2010/03/11}
\enonce{003798}{Mines MP 2000}
Soit $A\in M_n(\R)$ telle que $A^3={}^t\!A A$. $A$ est-elle diagonalisable dans $M_n(\R)$, dans
$M_n(\C)$ ? 
\finenonce{003798}



\finexercice


\exercice{3799, quercia, 2010/03/11}
\enonce{003799}{Centrale MP 2000 (avec Maple)}
Soit~$E$ un espace euclidien, $u$ et~$v$ deux endomorphismes auto-adjoints
de~$E$, $u$ étant défini positif.

\begin{enumerate}
  \item Montrer qu'il existe un unique endomorphisme~$w$ tel que
    $u\circ w + w\circ u = v$. Que peut-on dire de~$w$~?

  \item On suppose~$E$ de dimension~$3$, rapporté à une base orthonormale
    dans laquelle $u$ et~$v$ ont pour matrices respectives
    $A=\begin{pmatrix}4&1&1\cr 1&4&-1\cr1&-1&4\cr\end{pmatrix}$ et
    $B=\begin{pmatrix}0&\phantom-0&-1\cr 0&0&1\cr -1&1&3\cr\end{pmatrix}$. Déterminer~$w$.

  \item On revient au cas général. Si~$v$ est défini positif, que dire de~$w$~?
    Si~$w$ est défini positif, que dire de~$v$~?   
\end{enumerate}
\finenonce{003799}



\finexercice


\exercice{3800, quercia, 2010/03/11}
\enonce{003800}{Polytechnique MP$^*$ 2000}
Soit $E$ un espace euclidien et $s$ une symétrie de~$E$.

\begin{enumerate}
  \item Que dire de $s^*\circ s$~?
    
  \item Un polynôme $P$ est dit réciproque si $P(X) = X^nP\Bigl(\frac1X\Bigr)$,
    pour $P$ de degré~$n$.\par
    Montrer que~:
    $P(X) = \det(X\mathrm{id}+s^*\circ s)$ est un polynôme réciproque.
    
  \item Montrer que $P(1)\ge 2^n$. A quelle condition y a-til égalité~? Y a-t-il des
    conditions sur~$s$~?
    
  \item Soit la matrice $A = \begin{pmatrix}A_1&A_2\cr A_3&A_4\cr\end{pmatrix}$, carrée, d'ordre~$n$,
    symétrique définie positive, où $A_1$ et $A_4$ sont carrées d'ordres respectifs
    $p$ et~$q$. Montrer que $\det(A)\le \det(A_1)\det(A_4)$.
    
\end{enumerate}
\finenonce{003800}



\finexercice


\exercice{3801, quercia, 2010/03/11}
\enonce{003801}{Cachan MP$^*$ 2000}
On note $P$ l'ensemble des fonctions~$f$ polynomiales par morceaux,
continues sur $[0,1]$ et vérifiant ${f(0)=f(1)=0}$.
Si $f$ et~$g$ sont des fonctions de~$P$, on note
$(f\mid g) =  \int_{t=0}^1f'(t)g'(t)\,d t$.

\begin{enumerate}
  \item Que dire de~$P$ muni de cette application~?
    
  \item Montrer que si $x\in{[0,1]}$, il existe $g_x\in P$ telle que
    $\forall\ f\in P,\ (g_x\mid f) = f(x)$.
    
  \item On considère $n$ réels vérifiant~:
    $0<x_1<x_2<\dots<x_n<1$ et on donne $n$ réels $(\alpha_i)_{i\in{[[1,n]]}}$.
    On pose $\varphi(f) = \|f\|^2 + \sum_{i=1}^n (f(x_i)-\alpha_i)^2$
    et on demande de trouver le minimum de~$\varphi$ sur~$P$.
    

\end{enumerate}
\finenonce{003801}



\finexercice


\exercice{3802, quercia, 2010/03/11}
\enonce{003802}{Centrale MP 2002}

\begin{enumerate}
  \item Que peut-on dire de l'adjoint d'un projecteur orthogonal d'un espace euclidien~? Réciproque~?

  \item Soit $p$ un projecteur d'un espace euclidien tel que $p\circ p^*=p^*\circ p$. Montrer que $p$ est 
    un projecteur orthogonal.

\end{enumerate}
\finenonce{003802}



\finexercice


\exercice{3803, quercia, 2010/03/11}
\enonce{003803}{IIE MP 2004}
Soit $E = \mathcal{C}([0,1],\R)$ muni du produit scalaire défini par $(f\mid g) =  \int_0^1 fg$.

Soient $u,v$ les endomorphismes de~$E$ définis par
$u(f)(x) =  \int_0^x f$ et $v(f)(x) =  \int_x^1 f$.

\begin{enumerate}
  \item Montrer que $(u(f)\mid g) = (f\mid v(g))$.
  \item Déterminer les valeurs propres de~$u\circ v$.

\end{enumerate}
\finenonce{003803}



\finexercice


\exercice{3804, quercia, 2010/03/11}
\enonce{003804}{Centrale MP 2004}
Soit $E$ un espace euclidien de dimension~$n$ et $p$ endomorphismes
autoadjoints $u_1,\dots,u_p$. Soit $q_i$ la forme quadratique associée à~$u_i$
($q_i(x) = (u_i(x)\mid x)$). On suppose~:
$$\forall\ x\in E,\ q_1(x) + \dots + q_p(x) =
\|x\|^2\quad\text{et}\quad\mathrm{rg}(u_1)+\dots+\mathrm{rg}(u_p) = n.$$

\begin{enumerate}
  \item Montrer que $u_1+\dots+u_p = \mathrm{id}_E$.

  \item Montrer que $\Im(u_1)\oplus\dots\oplus\Im(u_p) = E$.

  \item Montrer que les~$u_i$ sont en fait des projecteurs orthogonaux et que la
    somme précédente est orthogonale.

\end{enumerate}
\finenonce{003804}



\finexercice


\exercice{3805, quercia, 2010/03/11}
\enonce{003805}{Mines MP 2005}

Soit $A$ matrice réelle~;
montrer que $A$ est diagonalisable ssi il existe $S$ symétrique réelle définie
positive telle que ${}^tA = SA S^{-1}$.


\finenonce{003805}



\finexercice


\exercice{3806, quercia, 2010/03/11}
\enonce{003806}{Rayon spectral, Centrale MP 2006}
Soient $A,B$ des matrices de $\mathcal{M}_n(\R)$ symétriques et 
$f : \R \to \R, t \mapsto {\max(\mathrm{Sp}(A+tB)).}$
Montrer que $f$ est convexe.


\finenonce{003806}



\finexercice

\section{204.08 Espaces vectoriels hermitiens}


\exercice{3823, quercia, 2010/03/11}
\enonce{003823}{Somme directe orthogonale}

Soit $E$ un espace préhilbertien et $F_1,\dots,F_n$ des sous-espaces vectoriels tels que
pour $i\ne j$, $F_i \perp F_j$.
Montrer que la somme $F_1+\dots+F_n$ est directe.
\finenonce{003823}



\finexercice


\exercice{3824, quercia, 2010/03/11}
\enonce{003824}{Espace $\ell^2$}

Soit $E$ l'ensemble des suites $(u_n)_{n\in\N}$ à termes réels telles que
la série $\sum u_n^2$ converge.

Pour $u,v \in E$, on pose : $(u\mid v) = \sum_{n=0}^{\infty} u_nv_n$.

\begin{enumerate}
  \item Montrer que $E$ est un espace vectoriel sur $\R$.
  \item Montrer que $(u\mid v)$ existe.
  \item Montrer qu'on définit ainsi un produit scalaire sur $E$.
  \item Montrer que $E$, muni de la norme associée, est complet.
\end{enumerate}
\finenonce{003824}



\finexercice


\exercice{3825, quercia, 2010/03/11}
\enonce{003825}{$f(x) \perp x  \Rightarrow  f = 0$}

Soit $E$ un espace vectoriel préhilbertien complexe et $f \in \mathcal{L}(E)$ tel
que pour tout vecteur $\vec x \in E$, on a $f(\vec x) \perp \vec x$.

\begin{enumerate}
  \item Montrer que pour tous vecteurs $\vec x,\vec y \in E$, on a
    $(f(\vec x)\mid \vec y) = 0$.
    
  \item Montrer que $f = 0$.
  \item Comparer avec le cas réel.
\end{enumerate}
\finenonce{003825}



\finexercice


\exercice{3826, quercia, 2010/03/11}
\enonce{003826}{Hanh-Banach pour une boule}

Soit $E$ un espace préhilbertien réel et $B$ une boule ouverte de $E$ ne
contenant pas $\vec 0$.
Montrer qu'il existe une forme linéaire $f \in E^*$ telle que :
$\forall\ \vec x \in B,\ f(\vec x) > 0$.
\finenonce{003826}



\finexercice


\exercice{3827, quercia, 2010/03/11}
\enonce{003827}{Calcul de minimums}

Calculer le minimum sur $\R^2$ de $f : {\R^2} \to \R, {(a,b)} \mapsto
{ \int_{x=0}^\pi (\sin x-ax^2-bx)^2\,d x.}$

\finenonce{003827}



\finexercice


\exercice{3828, quercia, 2010/03/11}
\enonce{003828}{Calcul de minimums}
\begin{enumerate}
  \item Soit $\varphi : {\R^n} \to \R$ définie par
    $\varphi(x_1,\dots,x_n) =  \int_{t=0}^1 (1+tx_1+\dots+t^nx_n)^2\,d t$.
    Montrer que $\varphi$ admet un minimum absolu et le calculer lorsque $n=3$.
    
  \item Même question avec
    $\psi(x_1,\dots,x_n) =  \int_{t=0}^{+\infty} e^{-t}(1+tx_1+\dots+t^nx_n)^2\,d t$.
    
\end{enumerate}
\finenonce{003828}



\finexercice


\exercice{3829, quercia, 2010/03/11}
\enonce{003829}{Trouvez le produit scalaire}

Soit $(P_n)_{n\in \N}$ une suite de polynômes de degrés étagés ($\deg P_n = n$).
Montrer qu'il existe un unique produit scalaire sur $\R[X]$ pour lequel la
famille $P_n$ est orthonormée.

\finenonce{003829}



\finexercice


\exercice{3830, quercia, 2010/03/11}
\enonce{003830}{Décomposition de Cholesky}

Soit $A \in \mathcal{M}_n(\R)$ symétrique définie positive.

\begin{enumerate}
  \item Montrer qu'il existe une matrice $T$ triangulaire supérieure telle que
    $A = {}^tTT$.
    Montrer que $T$ est unique si on impose la condition :
    $\forall\ i,\ T_{ii} > 0$.

  \item Application : Montrer que $\det A \le \prod_{i=1}^n a_{ii}$.
\end{enumerate}
\finenonce{003830}



\finexercice


\exercice{3831, quercia, 2010/03/11}
\enonce{003831}{Changement de base unitaire}

Soit $E$ un espace vectoriel hermitien et $\cal B$, ${\cal B}'$ deux bases
orthonormées de $E$.
On note $P$ la matrice de passage de $\cal B$ à ${\cal B}'$.

Montrer que ${}^t\overline PP = I$. Que peut-on dire de $\det P$ ?
\finenonce{003831}



\finexercice


\exercice{3832, quercia, 2010/03/11}
\enonce{003832}{Déterminant de Gram}

Soit $E$ un espace préhilbertien et $\vec u_1,\dots,\vec u_n \in E$.
On note $G = (g_{ij}) \in \mathcal{M}_n(K)$ la matrice de Gram de ces vecteurs
($g_{ij} = (\vec u_i\mid\vec u_j)$).

\begin{enumerate}
  \item On suppose $E$ de dimension finie, rapporté à une base orthonormée
    ${\cal B} = (\vec e_1,\dots,\vec e_p)$.
    Exprimer $G$ en fonction de
    $M = \text{Mat}_{\cal B}(\vec u_1,\dots,\vec u_n)$.
  \item En déduire que $\det(G)$ est un réel positif ou nul, et nul si et
    seulement si les vecteurs $\vec u_i$ sont liés.
  \item Montrer le même résultat sans supposer que $E$ est de dimension finie.
  \item Examiner le cas particulier $n=2$.
  \item Application : Le tétraèdre $ABCD$ est tel que $AB=AC=AD=1$ et
    $(AB,AC) \equiv \frac \pi4$, $(AB,AD) \equiv \frac \pi3$,
    $(AC,AD) \equiv \frac\pi2$. Calculer son volume.
    
\end{enumerate}
\finenonce{003832}



\finexercice


\exercice{3833, quercia, 2010/03/11}
\enonce{003833}{Congruence des matrices de Gram}

Soit $E$ un espace vectoriel hermitien et $\cal B$, ${\cal B}'$ deux bases quelconques.
On note $P$ la matrice de passage de $\cal B$ à ${\cal B}'$, et $G,G'$ les
matrices de Gram de $\cal B$ et ${\cal B}'$.
Quelle relation y a-t-il entre $P$, $G$ et $G'$ ?

\finenonce{003833}



\finexercice


\exercice{3834, quercia, 2010/03/11}
\enonce{003834}{Normes euclidiennes}

\begin{enumerate}
  \item Montrer que les applications :
    $${N_1} : {\R^2} \to \R, {(x,y)} \mapsto {\sqrt{x^2+xy+y^2}}
      \quad\text{et}\quad
       {N_2} : {\R^2} \mapsto \R, {(x,y)} \mapsto {\sqrt{2x^2-xy+y^2}}$$
    sont des normes.
  \item Montrer qu'il existe $\alpha,\beta > 0$ tels que :
    $\forall\ (x,y)\in\R^2,\ \alpha N_2(x,y) \le N_1(x,y) \le \beta N_2(x,y)$.
  \item Trouver les meilleures constantes $\alpha,\beta$
    (\'etudier si $N_1(x,y)^2 - \lambda N_2(x,y)^2$ est positive, négative).
    
\end{enumerate}
\finenonce{003834}



\finexercice


\exercice{3835, quercia, 2010/03/11}
\enonce{003835}{Famille duale de $1,X,X^2,\dots$}

\begin{enumerate}
  \item Montrer qu'il existe des polynômes $P_0,\dots,P_n \in \R_n[X]$ tels que :
$\forall\ i,j\le n,\  \int_{t=0}^{+\infty} e^{-t}t^iP_j(t)\,d t = \delta_{ij}$.
  \item Montrer qu'il n'existe pas de suite de polynômes $(P_0,\dots,P_n,\dots)$ telle que :
$\forall\ i,j\in\N,\  \int_{t=0}^{+\infty} e^{-t}t^iP_j(t)\,d t = \delta_{ij}$.

\end{enumerate}
\finenonce{003835}



\finexercice


\exercice{3836, quercia, 2010/03/11}
\enonce{003836}{Inégalité de Cauchy-Schwarz}

Soit $E$ l'ensemble des fonctions : $[a,b] \to \R^{+*}$ continues
et $\Phi : E \to \R, f \mapsto { \int_a^b f\times  \int_a^b 1/f.}$

Montrer que $\min\limits_{f\in E} \Phi(f) = (b-a)^2$ et chercher les fonctions
réalisant le minimum.

\finenonce{003836}



\finexercice


\exercice{3837, quercia, 2010/03/11}
\enonce{003837}{Intégrale double}

Soit $D$ le disque unité fermé de $\R^2$.
On considère l'espace $E$ des fonctions $f : D \to \R$ de classe $\mathcal{C}^1$
nulles sur le bord, $C$, de $D$.

Pour $f,g \in E$, on pose $(f\mid g) = \iint_D \Bigl(
\frac{\partial f}{\partial x}
\frac{\partial g}{\partial x} +
\frac{\partial f}{\partial y}
\frac{\partial g}{\partial y}
\Bigr) d x d y$.
Montrer que c'est un produit scalaire.

\finenonce{003837}



\finexercice


\exercice{3838, quercia, 2010/03/11}
\enonce{003838}{Forme quadratique associée à la matrice de Gram}

Soit $E$ un espace euclidien, $(\vec e_1,\dots,\vec e_n)$ une base de $E$,
$G$ sa matrice de Gram et $G^{-1} = (a_{ij})$.

Montrer que :
$\forall\ \vec x\in E,\ \sum_{i,j} a_{ij}
  (\vec e_i\mid \vec x) (\vec e_j\mid \vec x) = \|\vec x\,\|^2$.

\finenonce{003838}



\finexercice


\exercice{3839, quercia, 2010/03/11}
\enonce{003839}{Orthogonal des polynômes}

Soit $E = \mathcal{C}([0,1],\R)$ muni du produit scalaire usuel, $F$ le sous-espace vectoriel des fonctions
polynomiales et $g$ la fonction exponentielle sur $[0,1]$.

\begin{enumerate}
  \item Montrer que $g\notin F$.
  \item Montrer qu'il existe une suite $(f_n)$ de fonctions polynomiales convergeant
    vers $g$ pour la norme eucli\-dienne.
  \item En déduire que $F$ n'a pas de supplémentaire orthogonal.
\end{enumerate}
\finenonce{003839}



\finexercice


\exercice{3840, quercia, 2010/03/11}
\enonce{003840}{Orthogonal d'un hyperplan}

Soit $E = \mathcal{C}([0,1],\R)$ muni du produit scalaire usuel et, pour $f \in E$ :
$\varphi(f) =  \int_{t=0}^{1/2} f(t)\,d t$.

\begin{enumerate}
  \item Montrer que $\varphi$ est continue.
  \item Montrer que $H = \mathrm{Ker}\varphi$ est fermé.
  \item Montrer que $H^\bot = \{0\}$.
    
\end{enumerate}
\finenonce{003840}



\finexercice


\exercice{3841, quercia, 2010/03/11}
\enonce{003841}{Produit scalaire}

Soit $E = \mathcal{C}([a,b],\R)$ et $ u : {[a,b]} \to \R$ une fonction continue
par morceaux. On pose pour $f,g \in E$ :
$(f\mid g) =  \int_{t=a}^b u(t)f(t)g(t)\,d t$.

\begin{enumerate}
  \item A quelle condition sur $u$ définit-on ainsi un produit scalaire ?
    

  \item Soient $u,v$ deux fonctions convenables. A quelle condition les normes
    associées sont-elles équivalentes~?
    
\end{enumerate}
\finenonce{003841}



\finexercice


\exercice{3842, quercia, 2010/03/11}
\enonce{003842}{Produit scalaire~?}

Soit $E = \mathcal{C}([a,b],\R)$ et $(a_n)$ une suite d'éléments de $[a,b]$.

Pour $f,g \in E$,
on pose : $(f\mid g) = \sum_{n=0}^\infty \frac{f(a_n)g(a_n)}{2^n}$.

\begin{enumerate}
  \item A quelle condition sur la suite $(a_n)$ définit-on un produit scalaire ?


  \item Soient $a=(a_n)$ et $b=(b_n)$ deux telles suites telles que les
    ensembles $\{a_n,n\in\N\}$ et $\{b_n,n\in\N\}$ sont distincts.
    Montrer que les normes correspondantes sont non équivalentes.

  \item Question ouverte~: à quelle condition les normes associées à
    deux suites $(a_n)$ et $(b_n)$ sont-elles équivalentes~?

  \item Montrer qu'il n'existe pas de suite $(a_n)$ pour laquelle $E$ soit
    complet.

\end{enumerate}
\finenonce{003842}



\finexercice


\exercice{3843, quercia, 2010/03/11}
\enonce{003843}{Ulm MP$^*$ 2000}
$V$ est un espace hermitien et $u,v,w$ trois vecteurs unitaires.
Montrer que~:
$$\sqrt{1-|(u\mid v)|^2} \le \sqrt{1-|(u\mid w)|^2} + \sqrt{1-|(v\mid w)|^2}.$$

\finenonce{003843}



\finexercice


\exercice{3844, quercia, 2010/03/11}
\enonce{003844}{Ulm  MP$^*$ 2000}
Soit $A\in \mathcal{M}_n(\C)$. Montrer qu'elle admet une décomposition~$A=UT\,^tU$
avec $U$ unitaire et $T$ triangulaire supérieure si et seulement si
le spectre de~$A\overline A$ est inclus dans~$\R^+$.
\finenonce{003844}



\finexercice

\section{204.09 Problèmes matriciels}


\exercice{3807, quercia, 2010/03/11}
\enonce{003807}{$I + a(X^tY - Y^tX)$ inversible}

Soient $X,Y \in \mathcal{M}_{n,1}(\R)$ indépendantes, $a \in \R$ et $M$ la matrice
$n\times n$ telle que $m_{ij} = x_iy_j - x_jy_i$.

A quelle condition $I + aM$ est-elle inversible ?

\finenonce{003807}



\finexercice


\exercice{3808, quercia, 2010/03/11}
\enonce{003808}{Matrice orthogonale pour une forme $(p,q)$}

Soit $J = \begin{pmatrix}I_n&0\cr0&-I_p\cr\end{pmatrix}$, et $M = \begin{pmatrix}A&B\cr C&D\end{pmatrix}$
telle que ${}^tMJM = J$.
Montrer que $A$ et $D$ sont inversibles.
\finenonce{003808}



\finexercice


\exercice{3809, quercia, 2010/03/11}
\enonce{003809}{Calcul d'inverse}

Soit $A = \begin{pmatrix} a &-b &-c &-d \cr
                    b & a & d &-c \cr
                    c &-d & a & b \cr
                    d & c &-b & a \cr \end{pmatrix} \in \mathcal{M}_4(\R)$,
avec $a,b,c,d$ non tous nuls.

Démontrer que $A$ est inversible et calculer $A^{-1}$.
\finenonce{003809}



\finexercice


\exercice{3810, quercia, 2010/03/11}
\enonce{003810}{Matrices normales}

Soit $A \in \mathcal{M}_n(\C)$ de valeurs propres $\lambda_1, \dots, \lambda_n$.
Montrer que $AA^* = A^*A \iff \mathrm{tr}(AA^*) = |\lambda_1|^2+ \dots + |\lambda_n|^2$.
\finenonce{003810}



\finexercice


\exercice{3811, quercia, 2010/03/11}
\enonce{003811}{$f$ normal, $f^2=-\mathrm{id}$}

Soit $f \in \mathcal{L}(E)$ tel que $f\circ f^* = f^*\circ f$ et $f^2 = -$id.
Montrer que $f$ est orthogonal.
\finenonce{003811}



\finexercice


\exercice{3812, quercia, 2010/03/11}
\enonce{003812}{Chimie P 1996}

Soit $S$, matrice orthogonale d'ordre impair, de coefficients fonction
de~$t$ dérivables. Montrer que $\frac{d S}{d t}$ n'est pas inversible.

\finenonce{003812}



\finexercice


\exercice{3813, quercia, 2010/03/11}
\enonce{003813}{Chimie P 1996}

Déterminer $a$, $b$, $c$, réels non nuls pour que la matrice
$M = -\frac23\openup 3pt\begin{pmatrix}-\frac12&\frac ac &\frac ab\cr
                       \frac ca &-\frac12 &\frac ca\cr
                       \frac ba &\frac ac &-\frac12\cr\end{pmatrix}$
soit la matrice d'une isométrie.
\finenonce{003813}



\finexercice


\exercice{3814, quercia, 2010/03/11}
\enonce{003814}{Noyaux de $A$ et ${}^tA$}

Soit~$A\in\mathcal{M}_n(\R)$ telle que pour tout~$X\in\mathcal{M}_{n,1}(\R)$ on a
$\mathrm{tr}({}^tXAX)\ge 0$. Comparer les noyaux de~$A$ et ~${}^tA$.
\finenonce{003814}



\finexercice


\exercice{3815, quercia, 2010/03/11}
\enonce{003815}{Matrice orthogonale ?}

\begin{enumerate}
  \item Peut-on définir sur $\R^2$ une structure euclidienne telle que l'endomorphisme
$f$ dont la matrice dans la base canonique est
$M = \begin{pmatrix}1&1\cr-3&-2\cr\end{pmatrix}$ soit une rotation~?

  \item Généraliser à une matrice $M\in\mathcal{M}_2(\R)$ quelconque.

  \item Généraliser à une matrice $M\in\mathcal{M}_n(\R)$ quelconque.
\end{enumerate}
\finenonce{003815}



\finexercice


\exercice{3816, quercia, 2010/03/11}
\enonce{003816}{Valeurs propres d'une matrice complexe}
Soit~$M\in\mathcal{M}_n(\C)$ et $\lambda\in\mathrm{spec}(M)$. Montrer que $\Re(\lambda)$
est compris entre la plus grande et la plus petite valeur propre de
$\frac12(M+M^*)$.
\finenonce{003816}



\finexercice


\exercice{3817, quercia, 2010/03/11}
\enonce{003817}{Centrale MP 2001}

\begin{enumerate}
  \item Montrer que toute matrice symétrique réelle positive a ses coefficients
diagonaux positifs. Montrer que si l'un des coefficients diagonaux $u_{ii}$
est nul, alors pour tout $j$ on a $u_{ij} = 0$.

  \item $U$ est une matrice symétrique réelle positive de la forme
$U = \begin{pmatrix}A&C\cr{}^tC&B\cr\end{pmatrix}$ avec $A$ et $B$ carrées.
Montrer que la matrice $U' = \begin{pmatrix}A&C\cr0&0\cr\end{pmatrix}$ est diagonalisable.

\end{enumerate}
\finenonce{003817}



\finexercice


\exercice{3818, quercia, 2010/03/11}
\enonce{003818}{Centrale MP 2001}

\begin{enumerate}
  \item Pour $M\in GL_n(\R)$ montrer l'existence de deux matrices orthogonales
    $U$ et $V$ telles que $^tUMV$ soit diagonale.
    
  \item Même question pour $M\in\mathcal{M}_n(\R)$.
    
  \item Déterminer $U$ et $V$ pour $M=\begin{pmatrix}0 &1 &\phantom-1\cr -1 &0 &1\cr -1 &-1 &0\cr\end{pmatrix}$.

\end{enumerate}
\finenonce{003818}



\finexercice


\exercice{3819, quercia, 2010/03/11}
\enonce{003819}{X MP$^*$ 2001}

Soit $A\in\mathcal{M}_n(\R)$ telle que $A^2 = -I_n$.

\begin{enumerate}
  \item Montrer que $n$ est pair.

  \item Montrer que $A$ est semblable à~$A' = \begin{pmatrix}0 &-I_{n/2}\cr I_{n/2}&0\cr\end{pmatrix}$.

  \item On suppose $A\in{\cal O}(n)$. Montrer que $A$ est semblable à la matrice~$A'$
précédente avec une matrice de passage orthogonale.

\end{enumerate}
\finenonce{003819}



\finexercice


\exercice{3820, quercia, 2010/03/11}
\enonce{003820}{X MP$^*$ 2001}

Soient $A$ et $B$ deux matrices hermitiennes et  $C=A+B$. On note
$a_1\ge a_2\ge \dots\ge a_n$ les valeurs propres de la première,
$b_1\ge b_2\ge \dots\ge b_n$ celles de la deuxième, $c_1\ge c_2\ge \dots\ge c_n$ 
celles de la troisième.
Montrez que pour tout $i$ on a $c_i\ge a_i+b_n$.
{\it Indication~:} se ramener au cas $b_n=0$.

\finenonce{003820}



\finexercice


\exercice{3821, quercia, 2010/03/11}
\enonce{003821}{Centrale MP 2002}

Soient $n \in \N^*$, $S_n(\R)$ l'espace des matrices $n\times n$ symétriques
à coefficients réels,
$S_n^{+}(\R)$ le sous-ensemble des matrices positives,
$S_n^{++}(\R)$ le sous-ensemble des matrices définies positives
et $\phi \in \mathcal{L}({S_n(\R)})$.
On suppose que $\phi (S_n^{++}(\R))=S_n^{++}(\R)$.
\smallskip
\begin{enumerate}
  \item Montrer que : $\forall\ M \in S_n(\R),\ \exists\ A \in \R^+ \text{ tq }
    \forall\ \lambda > A,\ M + \lambda I_n \in S_n^{++}(\R)$.


  \item Montrer que $\phi \in GL(S_n(\R))$ et que $\phi (S_n^+(\R))=S_n^+(\R)$.



  \item On suppose $n = 2$ et $\phi (I_2) = I_2$. 
    Montrer que : $\forall\ M \in S_2(\R),\ \chi_{\phi (M)}=\chi_M$.
    Montrer que $\det(\varphi (M))=\det(M)$ (i.e. $\phi$ conserve le déterminant).

    
\end{enumerate}
\finenonce{003821}



\finexercice


\exercice{3822, quercia, 2010/03/11}
\enonce{003822}{Mines MP 2002}

Déterminer $\{M \in {\cal{M}}_n(\R) \text{ tq } M(^tMM)^2=I_n\}$.


\finenonce{003822}



\finexercice


\exercice{5489, rouget, 2010/07/10}
\enonce{005489}{**}
$E=\Rr^3$ euclidien orienté rapporté à une base orthonormée directe $\mathcal{B}$.
Etudier les endomorphismes de matrice $A$ dans $\mathcal{B}$ suivants~:

$$\begin{array}{lll}
1)\;A=-\frac{1}{3}
\left(
\begin{array}{ccc}
-2&1&2\\
2&2&1\\
1&-2&2
\end{array}
\right)
&2/\;A=\frac{1}{4}
\left(
\begin{array}{ccc}
3&1&\sqrt{6}\\
1&3&-\sqrt{6}\\
-\sqrt{6}&\sqrt{6}&2
\end{array}
\right)&3/\;A=\frac{1}{9}
\left(
\begin{array}{ccc}
8&1&4\\
-4&4&7\\
1&8&-4
\end{array}
\right).
\end{array}
$$
\finenonce{005489}


\finexercice


\exercice{5490, rouget, 2010/07/10}
\enonce{005490}{***}
Soit $M=\left(
\begin{array}{ccc}
a&b&c\\
c&a&b\\
b&c&a
\end{array}
\right)$ avec $a$, $b$ et $c$ réels.
Montrer que $M$ est la matrice dans la base canonique orthonormée directe de $R^3$ d'une rotation si et seulement si $a$, $b$ et $c$ sont les solutions d'une équation du type $x^3-x^2+k=0$ où $0\leq k\leq\frac{4}{27}$.
En posant $k=\frac{4\sin^2\varphi}{27}$, déterminer explicitement les matrices $M$ correspondantes ainsi que les axes et les angles des rotations qu'elles représentent.

\finenonce{005490}


\finexercice


\exercice{5780, rouget, 2010/10/16}
\enonce{005780}{*** I Matrices et déterminants de \textsc{Gram}}
Soit $E$ un espace préhilbertien réel.

Pour $n\in\Nn^*$ et $(x_1, ... , x_n)$ dans $E^n$, on pose $G(x_1, ... , x_n) =\left(x_i | x_j\right)_{1\leqslant i,j\leqslant n}$ (matrice de \textsc{Gram}) puis $\gamma(x_1, ... , x_n) = \text{det}(G(x_1 , ... , x_n))$ (déterminant de \textsc{Gram}).

\begin{enumerate}
 \item  Montrer que $\text{rg}(G(x_1, ... , x_n)) =\text{rg}(x_1, ... ,x_n)$.

\item  Montrer que la famille $(x_1, ... ,x_n)$ est liée si et seulement si $\gamma(x_1, ... , x_n) = 0$ et que la famille $(x_1, ... , x_n)$ est libre si et seulement si $\gamma(x_1, ... , x_n) > 0$.

\item  On suppose que la famille $(x_1, ... , x_n)$ est libre dans $E$. On pose $F =\text{Vect} (x_1, ... , x_n)$. Pour $x\in E$, on note $p_F(x)$ la projection orthogonale de $x$ sur $F$ puis $d(x,F)$ la distance de $x$ à $F$ (c'est-à-dire $d(x,F) =\|x-p_F(x)\|^2$). Montrer que $d(x,F) =\sqrt{\frac{\gamma(x,x_1,\ldots,x_n)}{\gamma(x_1,\ldots,x_n)}}$.
\end{enumerate}
\finenonce{005780}


\finexercice


\exercice{5786, rouget, 2010/10/16}
\enonce{005786}{*** I}
Montrer que la matrice de \textsc{Hilbert} $H_n =\left(\frac{1}{i+j-1}\right)_{1\leqslant i,j\leqslant n}$ est définie positive.
\finenonce{005786}


\finexercice


\exercice{5787, rouget, 2010/10/16}
\enonce{005787}{*** I}
\label{ex:rou2}
\begin{enumerate}
 \item  Soit $A$ une matrice carrée réelle de format $n$ et $S ={^t}AA$. Montrer que $S\in\mathcal{S}_n^+(\Rr)$.

\item  Réciproquement, montrer que pour toute matrice $S$ symétrique positive, il existe une matrice $A$ carrée réelle de format $n$ telle que $S ={^t}AA$. A-t-on l'unicité de $A$ ?

\item  Montrer que $S$ est définie positive si et seulement si $A$ est inversible.

\item  Montrer que $\text{rg}(A) =\text{rg}(S)$.

\item  (Racine carrée d'une matrice symétrique positive) Soit $S$ une matrice symétrique positive.

Montrer qu'il existe une et une seule matrice $R$ symétrique positive telle que $R^2 = S$. 
\end{enumerate}
\finenonce{005787}


\finexercice


\exercice{5791, rouget, 2010/10/16}
\enonce{005791}{***}
Soit $A$ une matrice orthogonale. Montrer que la valeur absolue de la somme des coefficients de $A$ est inférieure ou égale à $n$. Cas d'égalité si de plus tous les coefficients de $A$ sont positifs ?
\finenonce{005791}


\finexercice


\exercice{5794, rouget, 2010/10/16}
\enonce{005794}{***}
Soit $A$ une matrice carrée réelle symétrique positive de format $n$.
Montrer que $1 +\sqrt[n]{\text{det}(A)}\leqslant\sqrt[n]{\text{det}(I_n+A)}$.
\finenonce{005794}


\finexercice


\exercice{5795, rouget, 2010/10/16}
\enonce{005795}{**}
Déterminer $\text{card}(O_n(\Rr)\cap\mathcal{M}_n(\Zz))$.
\finenonce{005795}


\finexercice


\exercice{5798, rouget, 2010/10/16}
\enonce{005798}{**}
Soit $A$ une matrice carrée réelle. Montrer que les matrices ${^t}AA$ et $A{^t}A$ sont orthogonalement semblables.
\finenonce{005798}


\finexercice


\exercice{5799, rouget, 2010/10/16}
\enonce{005799}{*** I}
Montrer que le produit de deux matrices symétriques réelles positives est à valeurs propres réelles positives.
\finenonce{005799}


\finexercice


\exercice{5800, rouget, 2010/10/16}
\enonce{005800}{*** I}
Soient $A$ et $B$ deux matrices carrées réelles symétriques positives. 
Montrer que $\text{det}A +\text{det}B\leqslant\text{det}(A+B)$.
\finenonce{005800}


\finexercice


\exercice{5804, rouget, 2010/10/16}
\enonce{005804}{**}
La matrice $\left(
\begin{array}{cccc}
n-1&-1&\ldots&-1\\
-1&\ddots&\ddots&\vdots\\
\vdots&\ddots&\ddots&-1\\
-1&\ldots&-1&n-1
\end{array}
\right)$ est-elle positive ? définie ?
\finenonce{005804}


\finexercice


\exercice{5805, rouget, 2010/10/16}
\enonce{005805}{***}
$O_n(\Rr)$ est-il convexe ?
$O_n(\Rr)$ contient-il trois points alignés?
\finenonce{005805}


\finexercice


\exercice{5813, rouget, 2010/10/16}
\enonce{005813}{** I}
\label{ex:rou8}
Soit $A$ une matrice carrée réelle symétrique définie positive. Montrer qu'il existe une matrice triangulaire inversible $T$ telle que $A ={^t}TT$.

\finenonce{005813}


\finexercice


\exercice{5814, rouget, 2010/10/16}
\enonce{005814}{*** I}
Soit $A$ une matrice carrée réelle symétrique définie positive. Montrer que le déterminant de $A$ est inférieur ou égal au produit de ses coefficients diagonaux (utiliser l'exercice \ref{ex:rou8}).
\finenonce{005814}


\finexercice

\section{204.99 Autre}


\exercice{5487, rouget, 2010/07/10}
\enonce{005487}{**I}
Soit $a$ un vecteur non nul de l'espce euclidien $\Rr^3$. On définit $f$ de $\Rr^3$ dans lui même par~:~$\forall x\in\Rr^3,\;f(x)=a\wedge(a\wedge x)$. Montrer que $f$ est linéaire puis déterminer les vecteurs non nuls colinéaires à leur image par $f$.
\finenonce{005487}


\finexercice


\exercice{5492, rouget, 2010/07/10}
\enonce{005492}{***I Inégalité de \textsc{Hadamard}}
Soit $\mathcal{B}$ une base orthonormée de $E$, espace euclidien de dimension $n$.
Montrer que~:~$\forall(x_1,...,x_n)\in E^n,\;|\mbox{det}_{\mathcal{B}}(x_1,...,x_n)|\leq||x_1||...||x_n||$ en précisant les cas d'égalité.
\finenonce{005492}


\finexercice


\exercice{5497, rouget, 2010/07/10}
\enonce{005497}{***}
Soit $P\in\Rr_3[X]$ tel que $\int_{-1}^{1}P^2(t)\;dt=1$. Montrer que $\mbox{sup}\{|P(x)|,\;|x|\leq1\}\leq2$. Cas d'égalité~?
\finenonce{005497}


\finexercice


\exercice{5498, rouget, 2010/07/10}
\enonce{005498}{**IT}
Soit $r$ la rotation de $\Rr^3$, euclidien orienté, dont l'axe est orienté par $k$ unitaire et dont une mesure de l'angle est $\theta$.
Montrer que pour tout $x$ de $\Rr^3$, $r(x)=(\cos\theta)x+(\sin\theta)(k\wedge x)+2(x.k)\sin^2(\frac{\theta}{2})k$.
Application~:~écrire la matrice dans la base canonique (orthonormée directe de $\Rr^3$) de la rotation autour de $k=\frac{1}{\sqrt{2}}(e_1+e_2)$ et d'angle $\theta=\frac{\pi}{3}$.
\finenonce{005498}


\finexercice


\exercice{5499, rouget, 2010/07/10}
\enonce{005499}{**}
Soit $f$ continue strictement positive sur $[0,1]$. Pour $n\in\Nn$, on pose $I_n=\int_{0}^{1}f^n(t)\;dt$.
Montrer que la suite $u_n=\frac{I_{n+1}}{I_n}$ est définie et croissante.
\finenonce{005499}


\finexercice

\exercice{5796, rouget, 2010/10/16}
\enonce{005796}{}
Soit $f$ une application de $\Cc$ dans $\Cc$, $\Rr$-linéaire.

\begin{enumerate}
 \item  Montrer qu'il existe deux complexes $a$ et $b$ tels que pour tout $z\in\Cc$, $f(z) =az + b\overline{z}$.

\item  Calculer $\text{Tr}(f)$ et $\text{det}(f)$ en fonction de $a$ et $b$.

\item  C.N.S. pour que $f$ soit autoadjoint dans $\Cc$ muni de sa structure euclidienne canonique.
\end{enumerate}
\finenonce{005796}


\finexercice


\exercice{5802, rouget, 2010/10/16}
\enonce{005802}{*** I}
Soit $f$ un endomorphisme d'un espace euclidien de dimension $n$ qui conserve l'orthogonalité. Montrer qu'il existe un réel positif $k$ tel que $\forall x\in E$, $\|f(x)\| = k\|x\|$.
\finenonce{005802}


\finexercice

\section{205.01 Arithmétique de Z}


\exercice{2658, matexo1, 2002/02/01}
\enonce{002658}{}
Le but de cet exercice est de montrer que
$$\forall n\geq3 \quad \pi(2n+1)\geq \ln2\times{2n+1 \over \ln(2n+1)}$$
o{\`u} $\pi(x)$ d{\'e}signe le nombre d'entiers premiers inf{\'e}rieurs ou {\'e}gaux {\`a} $x$.

a) Calculer $I_{p,q}=\int_0^1 x^p(1-x)^qdx$ pour $p$ et $q$ entiers naturels.

b) Soit $D_n$ le ppmc de $n+1,\ n+2,\dots,\ 2n+1$. A l'aide de $I_{n,n}$, {\'e}tablir l'in{\'e}galit{\'e}
$D_n\geq{(2n+1)!\over(n!)^2}$

c) Montrer que $D_n\leq (2n+1)^{\pi(2n+1)}$ et en d{\'e}duire la minoration de $\pi(2n+1)$
annonc{\'e}e au d{\'e}but de l'exercice.
\finenonce{002658}

\finexercice


\exercice{5291, rouget, 2010/07/04}
\enonce{005291}{**}
Montrer que le produit de quatre entiers consécutifs, augmenté de 1, est un carré parfait.
\finenonce{005291}


\finexercice


\exercice{5292, rouget, 2010/07/04}
\enonce{005292}{***T}
\begin{enumerate}
\item  Montrer que $\forall n\in\Zz,\;6|5n^3+n$.
\item  Montrer que $\forall n\in\Nn,\;7|4^{2^n}+2^{2^n}+1$.
\end{enumerate}
\finenonce{005292}


\finexercice


\exercice{5293, rouget, 2010/07/04}
\enonce{005293}{***IT}
Un entier de la forme $8n+7$ ne peut pas être la somme de trois carrés parfaits.
\finenonce{005293}


\finexercice


\exercice{5294, rouget, 2010/07/04}
\enonce{005294}{**IT}
Pour $n\in\Nn^*$, on pose $(1+\sqrt{2})^n=a_n+b_n\sqrt{2}$ où $(a_n,b_n)\in(\Nn^*)^2$. Montrer que $a_n\wedge b_n=1$.
\finenonce{005294}


\finexercice


\exercice{5295, rouget, 2010/07/04}
\enonce{005295}{****}
Montrer que, pour tout entier naturel $n$, $2^{n+1}$ divise $E((1+\sqrt{3})^{2n+1})$.
\finenonce{005295}


\finexercice


\exercice{5296, rouget, 2010/07/04}
\enonce{005296}{***IT}
Soient $A$ la somme des chiffres de $4444^{4444}$ et $B$ la somme des chiffres de $A$. Trouver la somme des chiffres de $B$. (Commencer par majorer la somme des chiffres de $n=a_0+10a_1...+10^pa_p$.)
\finenonce{005296}


\finexercice


\exercice{5297, rouget, 2010/07/04}
\enonce{005297}{**}
Montrer que si $p$ est premier et $8p^2+1$ est premier alors $8p^2-1$ est premier.
\finenonce{005297}


\finexercice


\exercice{5298, rouget, 2010/07/04}
\enonce{005298}{**I}
\begin{enumerate}
\item  Montrer que $\forall(k,n)\in(\Nn^*)^2,\;[k\wedge n=1\Rightarrow n|C_n^k]$.
\item  Montrer que $\forall n\in\Nn^*,\;(n+1)|C_{2n}^n$.
\end{enumerate}
\finenonce{005298}


\finexercice


\exercice{5299, rouget, 2010/07/04}
\enonce{005299}{**T}
Résoudre dans $(\Nn^*)^2$ les équations ou systèmes d'équations suivants~:

$$1)\;\left\{
\begin{array}{l}
x+y=56\\
x\vee y=105
\end{array}
\right.
\quad
2)\;\left\{
\begin{array}{l}
x\wedge y=x-y\\
x\vee y=72
\end{array}
\right.\quad3)\;x\vee y-x\wedge y=243.$$
\finenonce{005299}


\finexercice


\exercice{5300, rouget, 2010/07/04}
\enonce{005300}{***}
Montrer que la somme de cinq carrés parfaits d'entiers consécutifs n'est jamais un carré parfait.
\finenonce{005300}


\finexercice


\exercice{5301, rouget, 2010/07/04}
\enonce{005301}{***IT}
Pour $n\in\Nn$, on pose $F_n=2^{2^n}+1$ (nombres de \textsc{Fermat}). Montrer que les nombres de Fermat sont deux à deux premiers entre eux.
\finenonce{005301}


\finexercice


\exercice{5302, rouget, 2010/07/04}
\enonce{005302}{***}
Soit $(u_n)_{n\in\Nn}$ la suite définie par $u_0=0$, $u_1=1$ et $\forall n\in\Nn,\;u_{n+2}=u_{n+1}+u_n$ (suite de \textsc{Fibonacci}).
\begin{enumerate}
\item  Montrer que $\forall n\in\Nn^*,\;u_{n+1}u_{n-1}-u_n^2=(-1)^n$ et en déduire que $\forall n\in\Nn^*,\;u_n\wedge u_{n+1}=1$.
\item  Montrer que $\forall n\in\Nn,\;\forall m\in\Nn^*,\;u_{m+n}=u_mu_{n+1}+u_{m-1}u_n$ et en déduire que $u_m\wedge u_n=u_{m\wedge n}$ pour $m$ et $n$ non nuls.
\end{enumerate}
\finenonce{005302}


\finexercice


\exercice{5303, rouget, 2010/07/04}
\enonce{005303}{***I}
On veut résoudre dans $\Zz^3$ l'équation $x^2+y^2=z^2$ (de tels triplets d'entiers relatifs sont appelés triplets pythagoriciens, comme par exemple $(3,4,5)$).
\begin{enumerate}
\item  Montrer que l'on peut se ramener au cas où $x\wedge y\wedge z=1$. Montrer alors que dans ce cas, $x$, $y$ et $z$ sont de plus deux à deux premiers entre eux.
\item  On suppose que $x$, $y$ et $z$ sont deux à deux premiers entre eux. Montrer que deux des trois nombres $x$, $y$ et $z$ sont impairs le troisième étant pair puis que $z$ est impair.

On suppose dorénavant que $x$ et $z$ sont impairs et $y$ est pair. On pose $y=2y'$, $X=\frac{z+x}{2}$ et $Z=\frac{z-x}{2}$.
\item  Montrer que $X\wedge Z=1$ et que $X$ et $Z$ sont des carrés parfaits.
\item  En déduire que l'ensemble des triplets pythagoriciens est l'ensemble des triplets de la forme

$$(d(u^2-v^2),\;2duv,\;d(u^2+v^2))$$

où $d\in\Nn$, $(u,v)\in\Zz^2$, à une permutation près des deux premières composantes.
\end{enumerate}
\finenonce{005303}


\finexercice


\exercice{5304, rouget, 2010/07/04}
\enonce{005304}{**}
Résoudre dans $\Nn^2$ l'équation $3x^3+xy+4y^3=349$.
\finenonce{005304}


\finexercice


\exercice{5305, rouget, 2010/07/04}
\enonce{005305}{***}
Résoudre dans $(\Nn^*)^2$ l'équation d'inconnue $(x,y)$~:~$\sum_{k=1}^{x}k!=y^2$.
\finenonce{005305}


\finexercice


\exercice{5306, rouget, 2010/07/04}
\enonce{005306}{***}
Montrer que $n=4...48...89$ ($p$ chiffres $4$ et $p-1$ chiffres $8$ et donc $2p$ chiffres) (en base $10$) est un carré parfait.
\finenonce{005306}


\finexercice


\exercice{5307, rouget, 2010/07/04}
\enonce{005307}{***I}
Montrer que tout nombre impair non divisible par $5$ admet un multiple qui ne s'écrit (en base $10$) qu'avec des $1$ (par exemple, $37.1=37$, $37.2=74$, $37.3=111$).
\finenonce{005307}


\finexercice


\exercice{5308, rouget, 2010/07/04}
\enonce{005308}{***}
Soit $u_n=10...01_2$.($n$ chiffres égaux à $0$). Déterminer l'écriture binaire de~:
\begin{enumerate}
 \item $u_n^2$,
 \item $u_n^3$,
 \item $u_n^3-u_n^2+u_n$.
\end{enumerate}
\finenonce{005308}


\finexercice


\exercice{5309, rouget, 2010/07/04}
\enonce{005309}{**I}
\label{exo:suprou19}
\begin{enumerate}
\item  Déterminer en fonction de $n$ entier non nul, le nombre de chiffres de $n$ en base $10$.
\item  Soit $\sigma(n)$ la somme des chiffres de $n$ en base $10$.
\begin{enumerate}
\item Montrer que la suite $\left(\frac{\sigma(n+1)}{\sigma(n)}\right)_{n\geq 1}$ est bornée. Cette suite converge-t-elle~?
\item Montrer que pour tout naturel non nul $n$, $1\leq\sigma(n)\leq9(1+\log n)$.
\item Montrer que la suite $(\sqrt[n]{\sigma(n)})_{n\geq1}$ converge et préciser sa limite.
\end{enumerate}
\end{enumerate}
\finenonce{005309}


\finexercice


\exercice{5310, rouget, 2010/07/04}
\enonce{005310}{***I}
\begin{enumerate}
\item  (Formule de \textsc{Legendre}) Soit $n$ un entier naturel supérieur ou égal à $2$ et $p$ un nombre premier. Etablir que l'exposant de $p$ dans la décomposition de $n!$ en facteurs premiers est

$$E(\frac{n}{p})+E(\frac{n}{p^2})+E(\frac{n}{p^3})+...$$

\item  Par combien de $0$ se termine l'écriture en base $10$ de $1000!$~?
\end{enumerate}
\finenonce{005310}


\finexercice


\exercice{5311, rouget, 2010/07/04}
\enonce{005311}{***I Petit théorème de \textsc{Fermat}}
Soit $p$ un nombre premier.
\begin{enumerate}
\item  Montrer que, pour tout entier $k$ tel que $1\leq k\leq p-1$, $p$ divise $C_p^k$.
\item  Montrer que $\forall a\in\Nn^*$, $a^p\equiv a$ $(p)$ (par récurrence sur $a$).
\end{enumerate}
\finenonce{005311}


\finexercice


\exercice{5312, rouget, 2010/07/04}
\enonce{005312}{***I Théorème de \textsc{Wilson}}
Soit $p$ un entier supérieur ou égal à $2$. Montrer que~:~$(p-1)!\equiv-1\;(p)\Rightarrow$ $p$ est premier (en fait les deux phrases sont équivalentes mais en Sup, on sait trop peu de choses en arithmétique pour pouvoir fournir une démonstration raisonnablement courte de la réciproque).
\finenonce{005312}


\finexercice

\section{205.02 Anneau Z/nZ, théorème chinois}


\exercice{1393, legall, 1998/09/01}

\enonce{001393}{}
Donner la liste des g\'en\'erateurs de $  (\Z /n\Z ,+)  .$
\finenonce{001393}



\finexercice



\exercice{1394, ortiz, 1999/04/01}

\enonce{001394}{}
Soit le groupe $G$ (additif) $\Zz/40\Zz.$
\begin{enumerate}
    \item  Soit $H$ le sous-groupe de $G$
engendr\'e par $\overline{12}$ et $\overline{20}.$ Montrer
que $H$ est le sous-groupe de $G$ engendr\'e par $\overline{4}
$ et trouver son ordre.
    \item  Caract\'eriser les g\'en\'erateurs de $G.$ Combien en compte-t-on?
    \item   D\'eterminer l'ordre de $\overline{15}.$
\end{enumerate}
\finenonce{001394}



\finexercice



\exercice{1395, ortiz, 1999/04/01}

\enonce{001395}{}
\begin{enumerate}
    \item Montrer qu'il n'existe aucun \'el\'ement d'ordre 3 dans le groupe
$\Zz/2\Zz\times \Zz/4\Zz.$
    \item En d\'eduire les morphismes de groupes de $\Zz/3\Zz$ dans $\Zz/2\Zz\times \Zz/4\Zz.$
\end{enumerate}
\finenonce{001395}



\finexercice



\exercice{1396, ortiz, 1999/04/01}

\enonce{001396}{}
Soit $f$ un morphisme de groupes de $\Zz/15\Zz$
dans $\Zz/18\Zz.$
\begin{enumerate}
     \item Montrer que $f$ est caract\'eris\'e par $f(\overline 1).$
     \item D\'eterminer les ordres possibles de $f(\overline 1).$
     \item En d\'eduire la liste des morphismes de groupes de $\mathbb{Z}/15\mathbb{Z}$ dans
 $\mathbb{Z}/18\mathbb{Z}.$
\end{enumerate}
\finenonce{001396}



\finexercice



\exercice{1397, ortiz, 1999/04/01}

\enonce{001397}{}
Soit $G$ le groupe-produit $\left(
\Zz/2\Zz\right) \times
\left( \Zz/4\Zz\right) .$
\begin{enumerate}
    \item  Donner la liste des \'el\'ements de $G$ et d\'eterminer l'ordre de
chacun d'entre eux. $G$ est-il cyclique ?
    \item   Donner la liste des sous-groupes de $G$ et en constuire le treillis.
\end{enumerate}
\finenonce{001397}



\finexercice



\exercice{1398, ortiz, 1999/04/01}

\enonce{001398}{}
\begin{enumerate}
    \item  Soit $f:\Zz \to  \Zz/3\Zz \times
 \Zz/5\Zz$ d\'efinie par
$f(n)=( \overline{n},\widetilde{n} ).$
    \begin{enumerate}
        \item Montrer que $f$ est un morphisme de groupes.
        \item D\'eterminer le noyau de $f.$
    \end{enumerate}
    \item  En d\'eduire que les groupes $\left( \Zz/3\Zz\right) \times
\left( \Zz/5\Zz\right)$ et $ \Zz/15\Zz$ sont isomorphes.
\end{enumerate}
\finenonce{001398}



\finexercice



\exercice{1399, ortiz, 1999/04/01}

\enonce{001399}{}
Les groupes $\Zz/8\Zz,$
$\left(\Zz/2\Zz\right)\times
\left(\Zz/4\Zz\right)$ et $\left(\Zz/2\Zz\right)^3$
sont-ils isomorphes ?
\finenonce{001399}



\finexercice



\exercice{1400, hilion, 2003/10/01}

\enonce{001400}{}
On note $R_n$ la rotation du plan de centre $O$, d'angle $2\pi/n$, $S$ la symétrie par rapport à l'axe $(Ox)$.
\begin{enumerate}
\item Montrer que $S^2=id$, $(R_n)^n=id$ et $R_nS=SR_n^{-1}$.
\item Montrer que le sous-groupe des isométries du plan engendré par $R_n$ et $S$ (ie le plus petit sous-groupe des isométries du plan qui contient $R_n$ et $S$) est de cardinal $2n$.
On le note $D_n$: c'est le groupe dihédral d'ordre $2n$.
\item Montrer que $D_n$ préserve un polygone régulier à $n$ côtés, centré en $O$.
\item En vous aidant de ce qui précède, construire un isomorphisme entre $D_3$ et $S_3$.
\end{enumerate}
\finenonce{001400}



\finexercice



\exercice{1401, hilion, 2003/10/01}

\enonce{001401}{}
Soit $\mathbb{H}=\biggl\lbrace
\begin{pmatrix}
z & w \\
-\bar{w} & \bar{z}
\end{pmatrix}
: (z,w)\in\Cc^2\biggr\rbrace$ l'ensemble des quaternions.
$\mathbb{H}^*$ désigne $\mathbb{H}$ privé de la matrice nulle.
On note
$\mathbf{1}=
\begin{pmatrix}
1 & 0 \\
0 & 1
\end{pmatrix}$,
$\mathbf{i}=
\begin{pmatrix}
i & 0 \\
0 & i
\end{pmatrix}$,
$\mathbf{j}=
\begin{pmatrix}
0 & 1 \\
-1 & 0
\end{pmatrix}$,
$\mathbf{k}=
\begin{pmatrix}
0 & i \\
-i & 0
\end{pmatrix}$.
\begin{enumerate}
\item Montrer que $\mathbb{H}^*$ est un sous-groupe de $GL_n(\Cc)$.
\item Montrer que $\mathbf{i}^2=\mathbf{j}^2=\mathbf{k}^2=\mathbf{1}$, $\mathbf{i}\mathbf{j}=\mathbf{k}$, $\mathbf{j}\mathbf{k}=\mathbf{i}$, $\mathbf{k}\mathbf{i}=\mathbf{j}$, $\mathbf{j}\mathbf{i}=-\mathbf{k}$, $\mathbf{k}\mathbf{j}==\mathbf{i}$, $\mathbf{i}\mathbf{k}=-\mathbf{j}$.
\item En déduire que le sous-groupe de $\mathbb{H}^*$ engendré par $\mathbf{i}$, $\mathbf{j}$ et $\mathbf{k}$ est d'ordre 8. On le note $H_8$.
\item Ecrire la table de $H_8$.
\item Vérifier que les groupes (tous de cardinal 8) $H_8$, $\Zz/2\Zz\times\Zz/2\Zz\times\Zz/2\Zz$, $\Zz/2\Zz\times\Zz/4\Zz$, $\Zz/8\Zz$ et $D_4$ sont 2 à 2 non isomorphes.
\end{enumerate}
\finenonce{001401}



\finexercice



\exercice{3151, quercia, 2010/03/08}

\enonce{003151}{{\'E}quations lin{\'e}aires}
R{\'e}soudre dans $\Z/37\Z$ :
\begin{enumerate}
  \item $$\left\{ \begin{aligned}\dot3x + \dot7y &= \dot3\cr \dot6x - \dot7y &= \dot0.\end{aligned}\right.$$
  \item $x^2 - \dot{\overline{31}}x + \dot{\overline{18}} = \dot0$.
\end{enumerate}
\finenonce{003151}


\finexercice 


\exercice{3152, quercia, 2010/03/08}

\enonce{003152}{{\'E}quation alg{\'e}brique}
\begin{enumerate}
  \item Dresser la liste des cubes dans $\Z/13\Z$.
  \item Soient $x,y,z \in \Z$ tels que $5x^3 + 11y^3 + 13z^3 = 0$.
    Montrer que 13 divise $x,y,z$.
  \item L'{\'e}quation : $5x^3 + 11y^3 + 13z^3 = 0$ a-t-elle des solutions enti{\`e}res ?
\end{enumerate}
\finenonce{003152}


\finexercice 


\exercice{3153, quercia, 2010/03/08}

\enonce{003153}{Ordre d'un entier modulo $n$}
\begin{enumerate}
  \item Soient $n,p \ge 2$. Montrer que : $n \wedge p = 1 \iff \exists\ k > 0$
    tel que $n^k \equiv 1 (\mathrm{mod}\, p)$.
  \item Soit $n$ un entier impair non divisible par 5. Montrer qu'il existe un multiple
    de $n$ qui s'{\'e}crit $1...1$ en base~10.
\end{enumerate}
\finenonce{003153}



\finexercice 


\exercice{3154, quercia, 2010/03/08}

\enonce{003154}{Th{\'e}or{\`e}me chinois}
Soient $n,p \in \N^*$ tels que $n \wedge p = 1$.
Pour $x\in\Z$ on note $\overline{x}^{\,n}$, $\overline{x}^{\,p}$ et $\overline{x}^{\,np}$
les classes d'{\'e}quivalence de $x$ modulo $n$, $p$ et $np$.
\begin{enumerate}
  \item Montrer que l'application $\phi : 
      {\Z/\bigl(np\Z\bigr)} \to 
      {\bigl(\Z/n\Z\bigr) \times \bigl(\Z/p\Z\bigr)},
      {\overline x^{\,np}} \mapsto {(\overline x^{\,n}, \overline x^{\,p})}$
      est un morphisme d'anneaux.

  \item En d{\'e}duire que $\varphi(np) = \varphi(n)\varphi(p)$ ($\varphi$ = fonction d'Euler).
  \item V{\'e}rifier que l'hypoth{\`e}se $n \wedge p = 1$ est n{\'e}c{\'e}ssaire.
\end{enumerate}
\finenonce{003154}



\finexercice 


\exercice{3155, quercia, 2010/03/08}

\enonce{003155}{Th{\'e}or{\`e}me de Wilson}
Soit $n \ge 2$. Montrer que $n$ est premier si et seulement si
$(n-1)! \equiv -1 (\mathrm{mod}\, n)$.
\finenonce{003155}



\finexercice 


\exercice{3156, quercia, 2010/03/08}

\enonce{003156}{$(\Z/2^n\Z)^*$}
\begin{enumerate}
  \item Montrer que pour tout entier $a$ impair et tout $n\ge 3$ : $a^{2^{n-2}} \equiv 1
    (\mathrm{mod}\, {2^n})$.
  \item Le groupe $\bigl(\Z/2^n\Z\bigr)^*$ est-il cyclique ?
\end{enumerate}
\finenonce{003156}



\finexercice 


\exercice{3157, quercia, 2010/03/08}

\enonce{003157}{{\'E}quation alg{\'e}brique}
\begin{enumerate}
  \item D{\'e}montrer que $f$ est une permutation de $E$.
  \item Chercher l'ordre de $f$ pour $\circ$.
  \item En d{\'e}duire que le nombre de points fixes de $f$ est congru {\`a} $\mathrm{Card}\, E$ modulo 3.
  \item D{\'e}montrer que ce nombre est inf{\'e}rieur ou {\'e}gal {\`a} 2.
  \item Combien l'{\'e}quation $x^2-x+\dot1 = \dot0$ a-t-elle de racines
    dans $\Z/p\Z$ en fonction de $p$ ?
  \item Pour $p=37$, r{\'e}soudre l'{\'e}quation.
\end{enumerate}
\finenonce{003157}



\finexercice 


\exercice{3158, quercia, 2010/03/08}

\enonce{003158}{Carr{\'e}s dans $\Z/p\Z$}
Soit $p$ un nombre premier impair.
Montrer que $\dot k$ est un carr{\'e} dans l'anneau $\Z/p\Z$ si et seulement si
$k^{(p+1)/2} \equiv k (\mathrm{mod}\, p)$.
\finenonce{003158}



\finexercice 


\exercice{3159, quercia, 2010/03/08}

\enonce{003159}{Test de primalit{\'e} de Rabin-Miller}
Soit $n$ un entier premier impair sup{\'e}rieur ou {\'e}gal {\`a}~$3$~: $n=q2^p+1$ avec $p$ impair
et soit $a\in\Z$ premier {\`a}~$n$.
On consid{\`e}re la suite $(b_0,b_1,\dots,b_p)$ d'entiers compris entre $0$ et
$n-1$ d{\'e}finie par~:
$$b_0\equiv a^q(\mathrm{mod}\, n),\quad
b_1\equiv b_0^2(\mathrm{mod}\, n),\quad\dots,\quad
b_p\equiv b_{p-1}^2(\mathrm{mod}\, n).$$
\begin{enumerate}
  \item Montrer que $b_p = 1$.
  \item Si $b_0\ne 1$ montrer qu'il existe un indice $i$ tel que $b_i = n-1$.
\end{enumerate}
\finenonce{003159}



\finexercice 


\exercice{3160, quercia, 2010/03/08}

\enonce{003160}{Coefficients du bin{\^o}me}
Soit $p$ un nombre premier.
Montrer que $\sum_{k=0}^pC_p^kC_{p+k}^k\equiv 2^p+1(\mathrm{mod}\,{p^2})$.
\finenonce{003160}


\finexercice 


\exercice{3161, quercia, 2010/03/08}

\enonce{003161}{Suite r{\'e}currente (Mines MP 2003)}
On consid{\`e}re la suite $(x_n)$ {\`a} valeurs dans~$\Z/11\Z$ telle que pour tout~$n$
on ait $x_{n+3} = 4(x_{n+2}+x_{n+1}+x_n)$.
D{\'e}terminer les diff{\'e}rents comportements possibles de $(x_n)$.
\finenonce{003161}


\finexercice 


\exercice{3162, quercia, 2010/03/08}

\enonce{003162}{$-3$ est-il un carr{\'e} ?}
Soit $p$ un nombre premier impair.
\begin{enumerate}
  \item Montrer qu'une {\'e}quation du second degr{\'e}~: $x^2 + ax + b = \dot 0$ admet
    une solution dans $\Z/p\Z$ si et seulement si son discriminant~: $a^2 - 4b$
    est un carr{\'e} dans $\Z/p\Z$.
  \item On suppose que $p\equiv 1 (\mathrm{mod}\, 3)$~: $p=3q+1$.
  \begin{enumerate}
    \item Montrer qu'il existe $a\in(\Z/p\Z)^*$ tel que $a^q\ne \dot 1$.
    \item En d{\'e}duire que $-\dot 3$ est un carr{\'e}.
  \end{enumerate}
  \item R{\'e}ciproquement, on suppose que $-\dot 3$ est un carr{\'e} dans
    $\Z/p\Z$. Montrer que $p\equiv 1 (\mathrm{mod}\, 3)$.
\end{enumerate}
\finenonce{003162}


\finexercice 


\exercice{3163, quercia, 2010/03/08}

\enonce{003163}{Indicateur d'Euler}
Soit $n\ge 3$. Montrer que $\varphi(n)$ est pair et
que $\sum_{x\wedge n = 1, 1\le x\le n} x = \frac{n\varphi(n)}2$.
\finenonce{003163}


\finexercice 

\section{205.03 Groupe fini commutatif}

\section{205.04 Arithmétique de K[X]}

\section{205.05 Corps fini}

\section{205.06 Applications}

\section{205.99 Autre}

\section{220.01 Convergence normale}

\section{220.02 Critères de Cauchy et d'Alembert}

\section{220.03 Rayon de convergence}


\exercice{4564, quercia, 2010/03/14}
\enonce{004564}{Vrai ou faux ?}

Dire si les affirmations suivantes sont vraies ou fausses. En donner
une démonstration ou un contre-exemple.


\begin{enumerate}
  \item Les séries $\sum a_nz^n$ et $\sum (-1)^na_nz^n$ ont même rayon de
    convergence.
  \item Les séries $\sum a_nz^n$ et $\sum (-1)^na_nz^n$ ont même domaine de
    convergence.
  \item Si la série $\sum a_nz^n$ a un rayon de convergence infini, alors elle
    converge uniformément sur $\R$.
  \item Si $\sum a_nx^n$ a un rayon de convergence fini $R > 0$, alors sa somme
    admet une limite infinie en $(-R)^+$ ou en~$R^-$.
  \item Si $f(x) = \sum a_nx^n$ a un rayon de convergence infini et si les $a_n$
    sont strictement positifs, alors pour tout entier~$p$,
    $\frac{f(x)}{x^p} \to +\infty$ lorsque $x\to+\infty$.
\end{enumerate}
\finenonce{004564}



\finexercice


\exercice{4565, quercia, 2010/03/14}
\enonce{004565}{Calculs de rayons}

Trouver le rayon de convergence de la série entière $\sum a_nz^n$ : 
\begin{enumerate}
  \item $a_n \to \ell \ne 0$ lorsque $n\to\infty$
  \item $(a_n)$ est périodique non nulle.\par
  \item $a_n = \sum_{d\mid n} d^2$.
  \item $a_n = \frac{n^n}{n!}$.
  \item $a_{2n} = a^n$, $a_{2n+1} = b^n$,\par $0<a<b$.
  \item $a_{n^2} = n!$, $a_k = 0$ si $\sqrt k \notin\N$.
  \item $a_n = (\ln n)^{-\ln n}$.
  \item $a_n = e^{\sqrt n}$.
  \item $a_n = \frac{1.4.7\dots(3n-2)}{n!}$.
  \item $a_n = \frac1{\sqrt n^{\sqrt n}}$.
  \item $a_n = \left(1+\frac12+\dots+\frac1n\right)^{\ln n}$.
  \item $a_{n+2} = 2a_{n+1} + a_n$,\par\nobreak $a_0=a_1=1$.
  \item $a_n = C_{kn}^n$.
  \item $a_n = e^{(n+1)^2} - e^{(n-1)^2}$.
  \item $a_n =  \int_{t=0}^1 (1+t^2)^n\,d t$.
  \item $a_n=\sqrt[n]n-\sqrt[{n+1}]{n+1}$.
  \item $a_n=\frac{\cos n\theta}{\sqrt n+(-1)^n}$.
\end{enumerate}
\finenonce{004565}



\finexercice


\exercice{4566, quercia, 2010/03/14}
\enonce{004566}{Centrale P' 1996}

Comment peut-on trouver le rayon de convergence d'une série entière
dont la suite des coefficients admet une infinité de zéros~?
\finenonce{004566}



\finexercice


\exercice{4567, quercia, 2010/03/14}
\enonce{004567}{Mines MP 2003}
Quel est le rayon de convergence de la série entière~:
$\sum_{k=0}^\infty \cos^k\Bigl(\frac{2k\pi}5+\alpha\Bigr)x^k$ où $\alpha\in\R$~?


\finenonce{004567}



\finexercice


\exercice{4568, quercia, 2010/03/14}
\enonce{004568}{Ensi MP 2003}
Rayon de convergence $R$ de la série entière $\sum_{n=1}^\infty \frac{x^n}{\sum_{k=1}^n k^{-\alpha}}$
et étude pour $x=\pm R$.

\finenonce{004568}



\finexercice


\exercice{4569, quercia, 2010/03/14}
\enonce{004569}{Centrale MP 2003}
On considère les suites $(a_n)$ et $(b_n)$ définies par :
$a_n = \frac{\cos(n\pi/3)}{n^{1/3}}$, $b_n = \sin(a_n)$.
\smallskip
\begin{enumerate}
  \item Déterminer les rayons de convergence des séries $\sum a_nx^n$ et $\sum b_nx^n$.
    
  \item Déterminer la nature de $\sum a_nx^n$ et $\sum b_nx^n$ en fonction de $x$.
    
\end{enumerate}
\finenonce{004569}



\finexercice


\exercice{4570, quercia, 2010/03/14}
\enonce{004570}{Transformation de rayons}

Soit $\sum a_nz^n$ une série entière de rayon de convergence $R > 0$.
Déterminer les rayons de convergence des séries :

\begin{enumerate}
  \item $\sum a_n^2 z^n$.               
  \item $\sum \frac {a_n}{n!} z^n$.    
  \item $\sum \frac{n!\,a_n}{n^n}z^n$. 
\end{enumerate}
\finenonce{004570}



\finexercice


\exercice{4571, quercia, 2010/03/14}
\enonce{004571}{Séries paire et impaire}

On suppose que les séries $\sum a_{2n}z^n$ et $\sum a_{2n+1}z^n$ ont
pour rayons de convergence $R$ et $R'$. Déterminer le rayon de convergence de
$\sum a_nz^n$.

\finenonce{004571}



\finexercice


\exercice{4572, quercia, 2010/03/14}
\enonce{004572}{Division par $z-\rho$}

Soit $a(z) = \sum_{n=0}^\infty a_nz^n$ une série entière de rayon de
convergence infini et $\rho > 0$.

On définit la série entière $b(z) = \sum_{n=0}^\infty b_nz^n$ de sorte que
$(z-\rho)b(z) = a(z)$ en cas de convergence de $b(z)$.

\begin{enumerate}
  \item Prouver l'existence et l'unicité des coefficients $b_n$.
    
  \item Quel est le rayon de convergence de $b(z)$ ?
    
\end{enumerate}
\finenonce{004572}



\finexercice


\exercice{4573, quercia, 2010/03/14}
\enonce{004573}{Développer peut être dangereux}

Pour $n\in\N$ et $x\in\R$ on pose $u_n(x) = \Bigl(\frac{x(1-x)}2\Bigr)^{4^n}$.

\begin{enumerate}
  \item Déterminer le domaine de convergence de la série $\sum_{n=0}^\infty u_n(x)$.

  \item On développe $u_n(x)$ par la formule du binôme~:
    $u_n(x) = \sum_{4^n\le k \le 2.4^n}a_kx^k$. Montrer que le
    rayon de convergence de la série entière $\sum_{k\ge 1} a_kx^k$
    est égal à~$1$ (en convenant que les $a_k$ non définis valent zéro).

\end{enumerate}
\finenonce{004573}


\finexercice


\exercice{5745, rouget, 2010/10/16}
\enonce{005745}{**}
Déterminer le rayon de convergence de la série entière proposée dans chacun des cas suivants :

\begin{enumerate}
 \item  $\sum_{n=1}^{+\infty}(\ln n)^n z^n$
\item  $\sum_{n=1}^{+\infty}\left(\sqrt{n}\right)^nz^n$
\item  $\sum_{n=0}^{+\infty}\left(\ln(n!)\right)^2 z^n$
\item  $\sum_{n=1}^{+\infty}\left(\frac{1}{2}\left(\ch\frac{1}{n}+\cos\frac{1}{n}\right)\right)^{n^4}z^n$
\item  $\sum_{n=1}^{+\infty}\frac{C_{2n}^n}{n^n}z^n$
\item  $\sum_{n=1}^{+\infty}\frac{\left(\ln(n!)\right)^a}{n!^b}z^n$
\item  $\sum_{n=0}^{+\infty}\frac{a^n}{1+b^n}z^n$, $(a,b)\in(\Rr_+^*)^2$
\end{enumerate}
\finenonce{005745}


\finexercice


\exercice{7544, mourougane, 2021/08/10}
\enonce{007544}{Calcul de rayon de convergence}
\begin{enumerate}
\item Déterminer le rayon de convergence de la série $\sum 2^{2n}z^{2n}$.
\item Déterminer la limite de $\frac{2^{2(n+1)}}{2^{2n}}$.
\end{enumerate}
\finenonce{007544}
\finexercice


\exercice{7545, mourougane, 2021/08/10}
\enonce{007545}{Opérations sur les rayons de convergence}
Soit $\sum a_n z^n$ et $\sum b_nz^n$ deux séries entières de rayon de convergence respectif $R_1$ et $R_2$.
\begin{enumerate}
 \item Montrer que le rayon de convergence de $\sum(a_n+b_n)z^n$ est supérieur à $\min(R_1,R_2)$ avec égalité si $R_2\not=R_1$.
 \item Montrer que le rayon de convergence de $\sum(a_nb_n)z^n$ est supérieur à $R_1R_2$.
\end{enumerate}
\finenonce{007545}
\finexercice

\section{220.04 Propriétés de la sommme d'une série entière}

\section{220.05 Calcul de la somme d'une série entière}


\exercice{4581, quercia, 2010/03/14}
\enonce{004581}{Sommation de séries entières}

Calculer les sommes des séries suivantes : 
\begin{enumerate}
  \item $\sum_{n=0}^\infty \frac{x^n}{2n-1}$.
  \item $\sum_{n=0}^\infty n^2x^n$.
  \item $\sum_{n=0}^\infty n^3x^n$.
  \item $\sum_{n=0}^\infty \frac{x^n}{(n+1)(n+3)}$.
  \item $\sum_{n=0}^\infty \frac{(-1)^nx^{2n+1}}{4n^2-1}$.
  \item $\sum_{n=0}^\infty \frac{x^n}{4n-1}$, $x\ge0$.
  \item $\sum_{n=0}^\infty \frac{n+3}{2n+1}x^n$.
  \item $\sum_{n=1}^\infty \frac{x^n}n\ch(na)$.
  \item $\sum_{n=0}^\infty \frac{n\sin^2(n\theta)}{2^n}$.
  \item $\sum_{n=0}^\infty \frac{n^2+1}{n+1}x^n$.
  \item $\sum_{n=0}^\infty \frac{x^n}{(2n)!}$.
  \item $\sum_{n=0}^\infty \frac{\sin^2(n\theta)}{n!}x^{2n}$.
  \item $\sum_{n=0}^\infty \frac{n^5x^n}{n!}$.
  \item $\sum_{n=0}^\infty \frac{x^{3n}}{(3n)!}$.
  \item $\sum_{n=1}^\infty C_{2n}^{n+1} x^n$.
  \item $\sum_{n=0}^\infty \frac1{n!} \int_{t=1}^x\ln^nt\,d t$.
  \item $\sum_{n=1}^\infty \left(1+\frac12+\dots+\frac1n\right)x^n$.
\end{enumerate}
\finenonce{004581}



\finexercice


\exercice{4582, quercia, 2010/03/14}
\enonce{004582}{Suite récurrente linéaire}

On définit deux suites $(u_n)$ et $(v_n)$ par :
$\begin{cases} u_0 = 1\cr v_0 = 0\cr\end{cases}$ et
$\begin{cases}u_{n+1} = u_n + 2v_n\cr v_{n+1} = u_n+v_n.\cr\end{cases}$

Déterminer le rayon de convergence et la somme de la série entière
$\sum_{n=0}^\infty u_nx^n$.

\finenonce{004582}



\finexercice


\exercice{4583, quercia, 2010/03/14}
\enonce{004583}{Série matricielle, Centrale MP 2000}

\begin{enumerate}
  \item Montrer l'existence de~$f(z) = \sum_{k=1}^{\infty} kz^k$ pour~$z\in\C$,
    $|z|<1$.
  \item Soit~$A\in\mathcal{M}_n(\C)$. Montrer que $\sum_{k=1}^{\infty} kA^k$ converge
    si et seulement si les valeurs propres de~$A$ sont de module strictement
    inférieur à~$1$.
    
  \item La somme $S = \sum_{k=1}^\infty kA^k$ est-elle inversible~?
    
\end{enumerate}
\finenonce{004583}



\finexercice


\exercice{4584, quercia, 2010/03/14}
\enonce{004584}{Série des traces (Centrale MP 2003)}
\smallskip
Soit~$A=\left(\begin{smallmatrix}1&1&1\cr1&1&0\cr1&0&0\cr\end{smallmatrix}\right)\in\mathcal{M}_{3}(\R)$.
\begin{enumerate}
  \item Montrer que~$A$ est diagonalisable et admet trois valeurs propres réelles
    dont on précisera les parties entières.
    
  \item On pose $t_n = \mathrm{tr}(A^n)$. Exprimer $t_n$ en fonction de $t_{n-1},t_{n-2},t_{n-3}$.
    
  \item Déterminer le rayon de convergence de la série entière $\sum_{n=0}^\infty t_nz^n$ et calculer sa somme.
    
\end{enumerate}
\finenonce{004584}



\finexercice


\exercice{4585, quercia, 2010/03/14}
\enonce{004585}{Centrale MP 2000}
\smallskip
Calculer $\sum_{n=1}^{\infty} \frac{(-1)^n}{3n+1}\cdotp$
\finenonce{004585}



\finexercice


\exercice{4586, quercia, 2010/03/14}
\enonce{004586}{$\sum P(n)x^n$, Ensi P 91}

    Rayon et somme de $\sum P(n)x^n$ où $P$ est un polynôme de degré $p$.
    
\finenonce{004586}



\finexercice


\exercice{4587, quercia, 2010/03/14}
\enonce{004587}{$\sum e^{in\theta}/2^n$, Ensi P 91}

    Calculer $\sum_{n=1}^\infty \frac{\sin n\theta}{2^n}$ et
	     $\sum_{n=1}^\infty \frac{\cos n\theta}{n2^n}$.
    
\finenonce{004587}



\finexercice


\exercice{4588, quercia, 2010/03/14}
\enonce{004588}{Ensae MP$^*$ 2000}

Soit $(u_n)$ définie par, pour tout $n\in \N$, $\sum_{k=0}^n\frac{u_{n-k}}{k!\strut}=1$. Trouver la
limite de $(u_n)$.

\finenonce{004588}



\finexercice

\exercice{5746, rouget, 2010/10/16}
\enonce{005746}{}
Calculer les sommes suivantes dans leur intervalle ouvert de convergence après avoir déterminé le rayon de convergence de la série proposée.

\begin{center}
\begin{tabular}{llll}
\textbf{1) (**)} $\sum_{n=2}^{+\infty}\frac{1}{n(n-1)}x^n$&\textbf{2) (**)} $\sum_{n=0}^{+\infty}\frac{3n}{n+2}x^n$&\textbf{3) (** I)} $\sum_{n=0}^{+\infty}\frac{x^n}{2n+1}$&\textbf{4) (**)} $\sum_{n=0}^{+\infty}\frac{n}{(2n+1)!}x^n$\\
\textbf{5) (*)} $\sum_{n=0}^{+\infty}\frac{x^{4n}}{(4n)!}$&\textbf{6) (**)} $\sum_{n=0}^{+\infty}(\ch n)x^n$&\textbf{7) (** I)} $\sum_{n=1}^{+\infty}\left(\sum_{k=1}^{n}\frac{1}{k}\right)x^n$&\textbf{7)} $\sum_{n=0}^{+\infty}\frac{n^2+4n-1}{n!(n+2)}x^n$\\
\textbf{9) (** I)} $\sum_{n=1}^{+\infty}n^{(-1)^n}x^n$&\textbf{10) (*)} $\sum_{n=1}^{+\infty}(-1)^n\frac{x^{4n-1}}{4n}$&\textbf{11) (**)} $\sum_{n=0}^{+\infty}(n^2+1)2^{n+1}x^n$&\textbf{12) (**)} $\sum_{n=0}^{+\infty}(-1)^{n+1}nx^{2n+1}$
\end{tabular}

\textbf{13) (***)} $\sum_{n=0}^{+\infty}a_nx^n$ où $a_0 = a_1 = 1$ et 
$\forall n\in\Nn$, $a_{n+2}= a_{n+1} + a_n$\rule{7.7cm}{0cm}

\textbf{14) (**)} $\sum_{n=0}^{+\infty}a_nx^n$ où $a_n$ est le nombre de couples $(x,y)$ d'entiers naturels tels que $x+5y = n$.\rule{3.3cm}{0cm}
\end{center}
\finenonce{005746}


\finexercice


\exercice{5752, rouget, 2010/10/16}
\enonce{005752}{***}
Calculer $\sum_{n=1}^{+\infty}\frac{1}{n}\cos\left(\frac{2n\pi}{3}\right)x^n$ pour $x$ dans $]-1,1[$. 
\finenonce{005752}


\finexercice


\exercice{5753, rouget, 2010/10/16}
\enonce{005753}{*** I}
Calculer $\sum_{n=0}^{+\infty}\frac{x^n}{4n^2-1}$ pour $x$ dans $]-1,1[$ et en déduire les sommes  $\sum_{n=0}^{+\infty}\frac{1}{4n^2-1}$ et $\sum_{n=0}^{+\infty}\frac{(-1)^n}{4n^2-1}$.
\finenonce{005753}


\finexercice


\exercice{5754, rouget, 2010/10/16}
\enonce{005754}{****}
 Pour $n$ entier naturel, on pose $u_n=\frac{(-1)^n}{2n+1}\sum_{k=0}^{n}\frac{1}{4k+1}$. Convergence et somme de la série (numérique) de terme général $u_n$.
\finenonce{005754}


\finexercice


\exercice{5757, rouget, 2010/10/16}
\enonce{005757}{**}
On pose $a_0 = 1$ et $b_0 = 0$ puis pour tout entier naturel $n$, $\left\{
\begin{array}{l}
a_{n+1}=-a_n-2b_n\\
b_{n+1}=3a_n+4b_n
\end{array}
\right.$. Rayons et sommes de $\sum_{n=0}^{+\infty}\frac{a_n}{n!}x^n$ et  $\sum_{n=0}^{+\infty}\frac{b_n}{n!}x^n$.
\finenonce{005757}


\finexercice


\exercice{5758, rouget, 2010/10/16}
\enonce{005758}{*** I}
Rayon de convergence et somme de $\sum_{n=1}^{+\infty}\frac{1}{nC_{2n}^n}x^n$.
\finenonce{005758}


\finexercice


\exercice{5763, rouget, 2010/10/16}
\enonce{005763}{***}
\label{ex:rou19x}
Soit $I_n$ le nombre d'involutions de $\llbracket1,n\rrbracket$. Rayon de convergence et somme de la série entière associée à la suite $\left(\frac{I_n}{n!}\right)_{n\in\Nn^*}$.
\finenonce{005763}


\finexercice

\section{220.06 Développement en série entière}


\exercice{4574, quercia, 2010/03/14}
\enonce{004574}{Développements en série entière}

Développer en série entière les fonctions suivantes : 
\begin{enumerate}
  \item $\ln(1+x+x^2)$.
  \item $(x-1)\ln(x^2-5x+6)$.
  \item $x\ln(x+\sqrt{x^2+1}\,)$.
  \item $\frac{x-2}{x^3-x^2-x+1}$.
  \item $\frac1{1+x-2x^3}$.
  \item $\frac{1-x}{(1+2x-x^2)^2}$.
  \item $\sqrt{\frac{1-x}{1+x}}$.
  \item $\Arctan(x+1)$.
  \item $\Arctan(x+\sqrt3\,)$.
  \item $ \int_{t=0}^x \frac{\ln(t^2-5t/2+1)}t\,d t$.
  \item $\Bigl(\frac{(1+x)\sin x}x\Bigr)^2$.
  \item $ \int_{t=x}^{2x} e^{-t^2}\,d t$.
  \item $e^{-2x^2} \int_{t=0}^x e^{2t^2}\,d t$.
  \item $\frac{\Arcsin \sqrt x}{\sqrt{x(1-x)}}$.
  \item $\sin\left(\frac13\Arcsin x\right)$.
\end{enumerate}
\finenonce{004574}



\finexercice


\exercice{4575, quercia, 2010/03/14}
\enonce{004575}{Ensi PC 1999}

Développer en série entière~: $\ln(\sqrt{1-2x\ch a + x^2})$.
\finenonce{004575}



\finexercice


\exercice{4576, quercia, 2010/03/14}
\enonce{004576}{$e^{x^2}/(1-x)$}

Développer en série entière $\frac{e^x}{1-x}$ puis $\frac{e^{x^2}}{1-x}$.


\finenonce{004576}



\finexercice


\exercice{4577, quercia, 2010/03/14}
\enonce{004577}{Mines-Ponts MP 2004}
Développer en série entière $f(x) = \sqrt{x+\sqrt{1+x^2}}$.
\finenonce{004577}



\finexercice


\exercice{4578, quercia, 2010/03/14}
\enonce{004578}{DSE d'une fraction rationnelle par récurrence linéaire}

Développer $f(x) = \frac x{1-x-x^2}$ en série entière en utilisant
la relation : $(1-x-x^2)f(x) = x$.

\finenonce{004578}



\finexercice


\exercice{4579, quercia, 2010/03/14}
\enonce{004579}{Produit de polynômes}

Quel est le coefficient de $x^n$ dans
$\Bigl(1+x+\dots+x^n\Bigr)
\Bigl(1+2x+\dots+(n+1)x^n\Bigr)
\Bigl(1+4x+\dots+(n+1)^2x^n\Bigr)$ ?

\finenonce{004579}



\finexercice


\exercice{4580, quercia, 2010/03/14}
\enonce{004580}{Développement en série entière de $\zeta(1+x)-1/x$}

\begin{enumerate}
  \item Vérifier que pour~$x\in{]0,+\infty[}$ on a~:
    $\zeta(1+x)-\frac1x = \sum_{n=1}^\infty\Bigl( \frac1{n^{1+x}} - \frac1x\bigl(\frac1{n^x}-\frac1{(n+1)^x}\bigr)\Bigr)$.

  \item Pour $p\in\N$ on pose $\gamma_p = \lim_{k\to\infty}\Bigl(\frac{\ln^p(1)}1 + \dots + \frac{\ln^p(k)}k - \frac{\ln^{p+1}(k+1)}{p+1}\Bigr)$.
    Justifier l'existence de $\gamma_p$ et montrer que $|\gamma_p| \le (p/e)^p$.

  \item Montrer alors que pour~$x\in{]0,1[}$ on a~:
    $\zeta(1+x)-\frac1x = \sum_{p=0}^\infty \frac{(-1)^p\gamma_p}{p!}x^p$.
\end{enumerate}
\finenonce{004580}



\finexercice

\exercice{4589, quercia, 2010/03/14}
\enonce{004589}{$\prod_{n=1}^\infty (1-q^nx)$}

Soit $q \in {]-1,1[}$ et $f(x) = \prod_{n=1}^\infty (1-q^nx)$.

\begin{enumerate}
  \item Montrer que $f(x)$ existe pour tout $x \in \R$ et que $f$ est développable
    en série entière au voisinage de 0.
    On admettra que si une fonction $g$ est DSE alors $e^g$ l'est.
    
  \item A l'aide de la relation : $f(x) = (1-qx)f(qx)$, calculer les coefficients
    du développement de $f$ et le rayon de convergence.
    

\end{enumerate}
\finenonce{004589}



\finexercice


\exercice{4590, quercia, 2010/03/14}
\enonce{004590}{Fonction non DSE}

Soit $f(x) = \sum_{n=0}^\infty e^{-n+n^2ix}$.
Montrer que $f$ est de classe $\mathcal{C}^\infty$ sur $\R$ mais n'est pas développable
en série entière autour de $0$.


\finenonce{004590}



\finexercice


\exercice{4591, quercia, 2010/03/14}
\enonce{004591}{Ens Ulm-Lyon-Cachan MP$^*$ 2003}
Soit~$\alpha>0$. On considère la fonction $f_\alpha$ : $x \mapsto\sum_{n=1}^\infty e^{-n^\alpha}e^{inx}$.
Montrer que~$f$ est $\mathcal{C}^\infty$. Donner une CNS sur~$\alpha$ pour que~$f$
soit développable en série entière en tout point de~$\R$.

\finenonce{004591}



\finexercice


\exercice{4592, quercia, 2010/03/14}
\enonce{004592}{Théorème de réalisation de Borel}

Soit $(a_n)$ une suite complexe donnée, on construit dans cet exercice
une fonction $f:\R \to \R$ de classe $\mathcal{C}^\infty$ telle que pour tout
entier~$n$ on ait $f^{(n)}(0) = n!\,a_n$.

Soit $\varphi : \R \to \R$ une fonction de classe $\mathcal{C}^\infty$
vérifiant~: $\forall\ x\in{[-1,1]},\ \varphi(x)=1$ et
$\forall\ x\notin{[-2,2]},\ \varphi(x)=0$
(l'existence de $\varphi$ fait l'objet de la question {2.}).
On pose $\varphi_n(x) = x^n\varphi(x)$,
$M_n = \max(\|\varphi_n'\|_\infty,\dots,\|\varphi_n^{(n)}\|_\infty)$
et $f(x) = \sum_{n=0}^\infty a_nx^n\varphi(\lambda_nx)$
où $(\lambda_n)$ est une suite de réels strictement
positifs, tendant vers $+\infty$ et telle que $\sum |a_n|M_n/\lambda_n$
converge.

\begin{enumerate}
  \item Montrer que $f$ est bien définie, est
    de classe $\mathcal{C}^\infty$ sur $\R$ et vérifie $f^{(n)}(0) = n!\,a_n$.
  \item Construction de~$\varphi$~: à l'aide de fonctions
    du type $x \mapsto\exp(-1/x)$ construire une fonction~$\psi$
    de classe $\mathcal{C}^\infty$ sur~$[0,+\infty[$ nulle sur
    $[0,1]\cup[2,+\infty[$ et strictement positive sur~$]1,2[$.

    Vérifier alors que $\varphi(x) =  \int_{t=|x|}^{+\infty}\psi(t)\,d t\Bigm/ \int_{t=0}^{+\infty}\psi(t)\,d t$
    convient.
    
\end{enumerate}
\finenonce{004592}



\finexercice


\exercice{5747, rouget, 2010/10/16}
\enonce{005747}{}
Développer en série entière les fonctions suivantes :

\begin{center}
\begin{tabular}{lll}
\textbf{1) (*)} $\frac{1}{(x-1)(x-2)}$&\textbf{2) (*** I)} $\frac{1}{x^2-2tx+1}$, $t\in\Rr$&\textbf{3) (*)} $\ln(x^2-5x+6)$\\
\textbf{4) (**)} $\Arctan\left(\frac{x\sin a}{1-x\cos a}\right)$, $a\in]0,\pi[$&
\textbf{5) (**)} $\frac{1}{(x-1)(x-2)\ldots(x-p)}$&\textbf{6) (*** I)} $(\Arcsin x)^2$\\
\textbf{7) (*)} $\int_{0}^{x}\cos(t^2)\;dt$&\textbf{8) (*** I)} $\int_{-\infty}^{x}\frac{dt}{t^4+t^2+1}$&
\textbf{9) (**)} $\cos x\ch x$.
\end{tabular}
\end{center}
\finenonce{005747}


\finexercice


\exercice{5748, rouget, 2010/10/16}
\enonce{005748}{* I}
Pour $x$ réel, on pose $f(x)=\left\{
\begin{array}{l}
\frac{\sin x}{x}\;\text{si}\;x\neq0\\
\rule{0mm}{4mm}1\;\text{si}\;x=0
\end{array}
\right.$. Montrer que $f$ est ce classe $C^\infty$ sur $\Rr$.
\finenonce{005748}


\finexercice


\exercice{5756, rouget, 2010/10/16}
\enonce{005756}{***}
Pour $x$ réel, on pose $F(x) =e^{-x^2}\int_{0}^{x}e^{t^2}\;dt$. En développant $F$ en série entière par deux méthodes différentes, montrer que pour tout entier naturel $n$, 

\begin{center}
$\sum_{k=0}^{n}(-1)^{n-k}\frac{1}{(2k+1)k!(n-k)!}=(-1)^n\frac{2^{2n}n!}{(2n+1)!}$.
\end{center}
\finenonce{005756}


\finexercice


\exercice{5761, rouget, 2010/10/16}
\enonce{005761}{**** I Développement en série entière de la fonction $x\mapsto\tan x$}
Pour $x\in\left]-\frac{\pi}{2},\frac{\pi}{2}\right[$, on pose $f(x) =\tan x$.

\begin{enumerate}
 \item  Montrer qu'il existe une suite de polynômes $(P_n)_{n\in\Nn}$ telle que pour tout entier naturel $n$, $f^{(n)}= P_n\circ f$ et que les $P_n$ sont à coefficients entiers naturels. (Utiliser $\tan' = 1+\tan^2$).

\item  En utilisant la formule de \textsc{Taylor}-\textsc{Laplace}, montrer que la série de \textsc{Taylor} à l'origine de $f$ a un rayon de convergence $R$ supérieur ou égal à $\frac{\pi}{2}$.

\item  On note $a_n$ les coefficients du développement précédent et $g$ la somme de la série entière associée à la suite $(a_n)_{n\in\Nn}$. Montrer que pour tout entier naturel non nul $n$, $(n+1)a_{n+1}= \sum_{k=0}^{n}a_ka_{n-k}$. En déduire que pour tout $x$ de $\left]-\frac{\pi}{2},\frac{\pi}{2}\right[$, $f(x) = g(x)$ et que $R=\frac{\pi}{2}$.

\item  Calculer $a_0$, $a_1$, $a_2$,..., $a_7$.

\item  Vérifier que la fonction $x\mapsto\tanh x$ est développable en série entière. Préciser le rayon et la valeur des coefficients en fonction des $a_n$.
\end{enumerate}
\finenonce{005761}


\finexercice


\exercice{5762, rouget, 2010/10/16}
\enonce{005762}{***}
Développer en série entière $F(x)=\int_{0}^{+\infty}e^{-t^2}\sin(tx)\;dt$ et en déduire que pour tout réel $x$, $F(x)=\frac{e^{-x^2/4}}{2}\int_{0}^{x}e^{t^2/4}\;dt$.
\finenonce{005762}


\finexercice


\exercice{7564, mourougane, 2021/08/10}
\enonce{007564}{Calculs}
Développer si possible en séries entières au voisinage de $0$ les applications 
\begin{enumerate}
 \item $f(z)=\sin (z)^2$.
 \item $g(z)=\cos (z^2-1)$.
 \item $h(z)=\frac{2z+1}{(z^2+1)(z+1)^2}$.
\end{enumerate}
\finenonce{007564}
\finexercice


\exercice{7565, mourougane, 2021/08/10}
\enonce{007565}{Intégrale}
Soit $a$ et $b$ deux nombres complexes tels que $|a|<1<|b|$ et $m,n$ deux entiers naturels.
Déterminer la valeur de $$\int_{\partial\Delta}\frac{d\xi}{(\xi-a)^m(\xi-b)^n}.$$
\finenonce{007565}
\finexercice


\exercice{7566, mourougane, 2021/08/10}
\enonce{007566}{Equation différentielle}
Déterminer toutes les applications holomorphes $f$ sur $\Cc$ telles que $f''+f=0$.
\finenonce{007566}
\finexercice


\exercice{7567, mourougane, 2021/08/10}
\enonce{007567}{Prolongement}
Soit $c$ un point de $\Cc^-$.
\begin{enumerate}
 \item Déterminer un développement en séries entières $\sum_n a_n (z-c)^n$ centré en $c$ de la branche principale du logarithme $\log$.
\item Déterminer le rayon de convergence du développement précédent.
\item La somme $\sum_{n=0}^{+\infty} a_n (z-c)^n$ coïncide-t-elle avec $\log$ sur l'intersection de $\Cc^-$ et de son disque de convergence.
\end{enumerate}
\finenonce{007567}
\finexercice


\exercice{7568, mourougane, 2021/08/10}
\enonce{007568}{Nombres de Bernouilli}
On définit les nombres de Bernouilli comme les nombres complexes $B_n$ tels que sur le domaine de convergence
$$\frac{z}{e^z-1}=\sum_{n=0}^{+\infty}\frac{B_n}{n!}z^n.$$
\begin{enumerate}
 \item Montrer que le rayon de convergence de la série précédente est $2\pi$.
 \item A partir de la relation $\frac{z}{e^z-1}\frac{e^z-1}{z}=1$ déterminer une relation de récurrence entre les $B_n$.
 \item Exprimer $\frac{e^z-1}{z}$ en fonction de $\frac{\sinh(z/2)}{\cosh(z/2)}$
 et en déduire que pour $n$ impair plus grand que $3$, $B_n=0$.
 \item Calculer $B_0, B_1\cdots B_8$.
 \item Montrer que tous les $B_n$ sont rationnels.
 \item La suite des $B_n$ est-elle bornée ?
\end{enumerate}
\finenonce{007568}
\finexercice

\section{220.07 Etude au bord}


\exercice{4593, quercia, 2010/03/14}
\enonce{004593}{\'Etude sur le cercle de convergence}

Pour $x \in \R$ on pose $f(x) = \sum_{n=1}^\infty x^n\sin\frac1{\sqrt n}$.

\begin{enumerate}
  \item Déterminer le rayon de convergence, $R$, de cette série. 
  \item \'Etudier la convergence de $f$ pour $x = \pm R$.
    
  \item Déterminer $\lim_{x\to R^-} f(x)$.
    
\end{enumerate}
\finenonce{004593}



\finexercice


\exercice{4594, quercia, 2010/03/14}
\enonce{004594}{Coefficients équivalents $ \Rightarrow $ séries équivalentes}

Soit $(a_n)$ une suite de réels strictement positifs.
On suppose que le rayon de convergence de la série entière
$A(x) = \sum_{n=0}^\infty a_nx^n$ est 1 et que la série diverge pour $x = 1$.

\begin{enumerate}
  \item Montrer que $A(x) \to +\infty$ lorsque $x\to1^-$.
    

  \item Soit $(b_n)$ une suite telle que $b_n \sim a_n$ et
    $B(x) = \sum_{n=0}^\infty b_nx^n$.
    Montrer que $B(x) \sim A(x)$ pour $x\to1^-$.
    
\end{enumerate}
\finenonce{004594}



\finexercice


\exercice{4595, quercia, 2010/03/14}
\enonce{004595}{Produit de Cauchy}

Soit $(c_n)$ le produit de Cauchy de la suite $(a_n)$ par la suite
$(b_n)$.
Montrer que si les trois séries $\sum a_n$, $\sum b_n$ et $\sum c_n$
convergent vers $A$,$B$,$C$, alors $C = AB$
(considérer les séries entières $\sum a_nz^n$, $\sum b_nz^n$ et $\sum c_nz^n$).

\finenonce{004595}



\finexercice


\exercice{4596, quercia, 2010/03/14}
\enonce{004596}{Produit de Cauchy}

Soit $(c_n)$ le produit de Cauchy de la suite $(a_n)$ par la suite
$(b_n)$. On suppose que la série $A(z) = \sum_{n=0}^\infty a_nz^n$
a un rayon $R > 0$ et que $b_n/b_{n+1} \to \lambda$ lorsque $n\to\infty$
avec $|\lambda| < R$.
Montrer que $c_n/b_n\to A(\lambda)$ lorsque $n\to\infty$.
\finenonce{004596}



\finexercice
\section{220.08 Equations différentielles}


\exercice{4597, quercia, 2010/03/14}
\enonce{004597}{\'Equation différentielle}

Montrer que l'équation $3xy' + (2-5x)y = x$ admet une solution développable en
série entière autour de $0$.
Calculer $y(1)$ à $5.10^{-5}$ près.

\finenonce{004597}



\finexercice


\exercice{4598, quercia, 2010/03/14}
\enonce{004598}{DSE de tan}

\begin{enumerate}
  \item En utilisant la relation : $\tan' = 1 + \tan^2$, exprimer $\tan^{(n)}$
    en fonction de $\tan,\dots,\tan{(n-1)}$.
    En déduire que~: $\forall\ x \in {[0,\pi/2[}$, $\tan^{(n)}(x) \ge 0$.
    
  \item Montrer que la série de Taylor de $\tan$ en 0 converge sur $]-\pi/2,\pi/2[$.
    
  \item Soit $f$ la somme de la série précédente. Montrer que $f' = 1+f^2$ et en
    déduire que $f = \tan$.
  \item Prouver que le rayon de convergence est exactement $\pi/2$.
    
\end{enumerate}
\finenonce{004598}



\finexercice


\exercice{4599, quercia, 2010/03/14}
\enonce{004599}{DSE de $(\Arcsin x)^2$}

On pose $f(x) = \frac{\Arcsin x}{\sqrt{1-x^2}}$.

\begin{enumerate}
  \item Montrer que $f$ admet un développement en série entière au voisinage de $0$
    et préciser le rayon de convergence.
    
  \item Chercher une équation différentielle d'ordre 1 vérifiée par $f$.
    En déduire les coefficients du développement en série entière de $f$.
    
  \item Donner le développement en série entière de $\Arcsin^2 x$.
    
\end{enumerate}
\finenonce{004599}



\finexercice


\exercice{4600, quercia, 2010/03/14}
\enonce{004600}{$\sum_{n=0}^\infty \frac 1{\strut C_{2n}^n}$}

On pose $f(x) = \sum_{n=0}^\infty \frac{x^n}{\strut C_{2n}^n}$.

\begin{enumerate}
  \item Déterminer le rayon de convergence et montrer que $f$ vérifie l'équation :
    $x(4-x)y' - (x+2)y = -2$.
    
  \item Résoudre l'équation précédente pour $x > 0$
    (utiliser le DL de $f$ en 0 à l'ordre 1 pour fixer la constante)
    et en déduire la somme de la série
    $\sum_{n=0}^\infty \frac 1{\strut C_{2n}^n}$.
    
\end{enumerate}
\finenonce{004600}



\finexercice


\exercice{4601, quercia, 2010/03/14}
\enonce{004601}{Calcul de somme}

On pose $f(x) = \sum_{n=0}^\infty \frac{n!\,x^{2n+1}}{1.3.5\dots(2n+1)}$.

\begin{enumerate}
  \item Déterminer le rayon de convergence. 
  \item \'Etudier la convergence aux bornes de l'intervalle de convergence.
    
  \item Calculer $f(x)$.
    
\end{enumerate}
\finenonce{004601}



\finexercice


\exercice{4602, quercia, 2010/03/14}
\enonce{004602}{Fonction génératrice du nombre de partitions}

On note $T_n$ le nombre de partitions d'un ensemble à $n$ éléments.

\begin{enumerate}
  \item Montrer que $T_{n+1} = \sum_{k=0}^n {\binom{n}{k}}T_k$.
  \item Montrer que $\sum_{n=0}^\infty \frac{T_nx^n}{n!} = e^{e^x-1}$.
    
\end{enumerate}
\finenonce{004602}



\finexercice


\exercice{4603, quercia, 2010/03/14}
\enonce{004603}{Suite récurrente}

Soit $(a_n)$ la suite réelle définie par :
$a_0 = 1$, $2a_{n+1} = \sum_{k=0}^n C_n^k a_ka_{n-k}$.
On pose $f(x) = \sum_{n=0}^\infty \frac{a_n}{n!}x^n$.

\begin{enumerate}
  \item Montrer que le rayon de convergence est non nul.
    
  \item Calculer $f(x)$.
    
  \item En déduire $a_n$ en fonction de $n$.
    
\end{enumerate}
\finenonce{004603}



\finexercice


\exercice{4604, quercia, 2010/03/14}
\enonce{004604}{Fonction $\zeta$\label{recurzeta}}

Pour $|x| < 1$ on pose~: $Z(x) = \sum_{n=1}^\infty \zeta(2n)x^n$.

Montrer que $Z$ vérifie l'équation différentielle~:
$2xZ'(x) - 2Z^2(x) + Z(x) = 3x\zeta(2)$ (écrire $Z(x)$ comme somme d'une
série double, intervertir les sommations, remplacer et \dots\ simplifier).

En déduire la relation de récurrence~:
$\forall\ n\ge 2,\ (n+\frac12)\zeta(2n) = \sum_{p=1}^{n-1}\zeta(2p)\zeta(2n-2p)$.
\finenonce{004604}



\finexercice


\exercice{4605, quercia, 2010/03/14}
\enonce{004605}{DSE de tan}

On note $\zeta_i(n) = \sum_{k=0}^\infty \frac1{(2k+1)^n}$
et $Z_i(x) = \sum_{n=1}^\infty \zeta_i(2n)x^n$.
En s'inspirant de l'exercice~\ref{recurzeta}
montrer que $Z_i$ vérifie l'équation différentielle~:
$2xZ_i'(x) - 2Z_i^2(x) - Z_i(x) = x\zeta_i(2)$.

Déterminer alors deux réels $\alpha$ et $\beta$ tels que
$T(x) = Z_i(x^2)/x$ soit égal à~$\alpha\tan\beta x$ sur~$]-1,1[$.

\finenonce{004605}



\finexercice


\exercice{4606, quercia, 2010/03/14}
\enonce{004606}{DSE de $\tan x$.}

\begin{enumerate}
  \item Pour $a,b\in\R$ avec $b\not\equiv 0(\mathrm{mod}\,{\pi})$, vérifier
    l'identité suivante~:
    $\frac{(1+ia) - e^{ib}(1-ia)}{1-e^{ib}} = 1 - \frac a{\tan(b/2)}$.
  \item Pour $a,b\in\C$ et $n\in\N^*$, vérifier
    l'identité suivante~:
    $a^n + b^n = \prod_{k=0}^{n-1}(a-be^{i(2k+1){\frac \pi n}})$.
  \item Pour $x\in\R$ et $p\in\N^*$, vérifier
    l'identité suivante~:
    $\frac{\Bigl(1+\frac{ix}{2p}\Bigr)^{2p} + \Bigl(1-\frac{ix}{2p}\Bigr)^{2p}}{\vrule height 10pt width 0pt 2} =
     \prod_{k=0}^{p-1}\Bigl(1-\frac{x^2}{4p^2\tan^2\frac{(2k+1)\pi}{4p}}\Bigr)$.
  \item Démontrer alors~:
    $\forall\ x\in{]-\frac\pi2,\frac\pi2[},\ \ln(\cos x) = \sum_{k=0}^\infty\ln\Bigl(1-\frac{4x^2}{(2k+1)^2\pi^2}\Bigr)$.

  \item En déduire~:
    $\forall\ x\in{]-\frac\pi2,\frac\pi2[},\
    \tan x = \sum_{k=0}^\infty \frac{8x}{(2k+1)^2\pi^2 - 4x^2}$.

  \item Pour $n\in\N$ avec $n\ge 2$, vérifier l'identité suivante~:
    $\sum_{k=0}^\infty \frac1{(2k+1)^n} = \frac{2^n-1}{\vrule height 10pt width 0pt 2^n}\zeta(n)$.

  \item Démontrer enfin~:
    $\forall\ x\in{]-\frac\pi2,\frac\pi2[},\
    \tan x = \sum_{n=1}^\infty \frac{2(4^n-1)}{\pi^{2n}}\zeta(2n)x^{2n-1}$.
\end{enumerate}
\finenonce{004606}



\finexercice
\section{220.09 Intégrales}


\exercice{4607, quercia, 2010/03/14}
\enonce{004607}{$ \int_{t=0}^1 t^t\,d t$}

\begin{enumerate}
  \item A l'aide d'un développement en série entière, montrer que
    $ \int_{t=0}^1 t^t\,d t = \sum_{n=1}^\infty \frac{(-1)^{n-1}}{n^n}$.
    

  \item Calculer la valeur commune des deux membres à $10^{-5}$ près.
    
\end{enumerate}
\finenonce{004607}



\finexercice


\exercice{4608, quercia, 2010/03/14}
\enonce{004608}{$ \int_{t=0}^1 \ln(t)\ln(1-t)\,d t$}

On admet que $\sum_{n=1}^\infty \frac 1{n^2} = \frac{\pi^2}6$.
Calculer $ \int_{t=0}^1 \ln(t)\ln(1-t)\,d t$.


\finenonce{004608}



\finexercice


\exercice{4609, quercia, 2010/03/14}
\enonce{004609}{Centrale PSI 1997}

Établir la convergence puis calculer la valeur de
$ \int_{t=0}^1 \frac{\ln(t^2)\ln(1-t^2)}{t^2}\,d t$.

\finenonce{004609}



\finexercice


\exercice{4610, quercia, 2010/03/14}
\enonce{004610}{$ \int_{t=0}^x \frac{\ln(1-t)}t\,d t$}

Montrer que pour $x \in {]-1,1[}$ :
$\sum_{n=1}^\infty \frac {x^n}{n^2} = - \int_{t=0}^x \frac{\ln(1-t)}t\,d t$.
En déduire la valeur de $ \int_{t=0}^1 \frac{\ln(1-t)}t\,d t$.
\finenonce{004610}



\finexercice


\exercice{4611, quercia, 2010/03/14}
\enonce{004611}{Intégrale elliptique}

Montrer que la longueur d'une ellipse de demi-axes $a,b$ est :
$L = 2\pi\sqrt{\frac{a^2+b^2}2}\sum_{p=0}^\infty
\left(\frac{a^2-b^2}{a^2+b^2}\right)^{2p}\frac{C_{4p}^{2p} C_{2p}^p}
      {4^{3p}(1-4p)}$.
\finenonce{004611}



\finexercice


\exercice{4612, quercia, 2010/03/14}
\enonce{004612}{Norme $L^2$}

Soit $f(z) = \sum_{n=0}^\infty a_nz^n$ une série de rayon $R>0$.
Montrer, pour~$0\le r< R$~:
$\sum_{n=0}^\infty |a_n|^2r^{2n} =
\frac 1{2\pi}  \int_{\theta=0}^{2\pi}|f(re^{i\theta})|^2\, d\theta$.
\finenonce{004612}



\finexercice


\exercice{5751, rouget, 2010/10/16}
\enonce{005751}{*** I}
Pour $n\in\Nn$, on pose $W_n=\int_{0}^{\pi/2}\cos^nt\;dt$. Rayon de convergence et somme de la série entière associée à la suite $(W_n)_{n\in\Nn}$.
\finenonce{005751}


\finexercice

\section{220.10 Analycité}


\exercice{4613, quercia, 2010/03/14}
\enonce{004613}{Série à valeurs réelles}

Soit $f(z) = \sum a_nz^n$ une série de rayon $R>0$ telle que
pour tout $z\in \mathring D(0,R)$ on a $f(z)\in\R$.
Montrer que $f$ est constante.
\finenonce{004613}



\finexercice


\exercice{4614, quercia, 2010/03/14}
\enonce{004614}{Formules de Cauchy\label{Cauchy}}

Soit $U$ un ouvert de~$\C$ contenant $0$ et $f : U \to \C$ analytique.
On note $\sum_{n=0}^\infty a_nz^n$ le développement en série entière
de $f$ en~$0$, $R$ son rayon et~$d$ la distance de~$0$ à~$\mathrm{fr}(U)$
($d=+\infty$ si $U=\C$).

\begin{enumerate}
  \item Montrer, pour $0<r<\min(R,d)$ et $n\in\N$~:
    $a_n = \frac1{2\pi} \int_{\theta=0}^{2\pi}\frac{f(re^{i\theta})}{r^ne^{in\theta}}\, d\theta$.
    \label{formule-c}

  \item Montrer que l'application $r \mapsto  \int_{\theta=0}^{2\pi}\frac{f(re^{i\theta})}{e^{in\theta}}\, d\theta$
    est analytique sur~$[0,d[$ (minorer le rayon de convergence du
    DSE de~$f$ en $r_0e^{i\theta}$ et majorer en module les coefficients lorsque
    $\theta$ décrit~$[0,2\pi]$ et $r_0$ est fixé dans~$[0,d[$ à l'aide
    d'un recouvrement ouvert de~$[0,2\pi]$). En déduire que l'égalité de la question \ref{formule-c}.\
    a lieu pour tout $r\in{[0,d[}$.
  \item Pour $0<r<d$ et $|z|<r$ on pose $g(z)= \frac1{2\pi} \int_{\theta=0}^{2\pi}\frac{f(re^{i\theta})}{re^{i\theta}-z}re^{i\theta}\, d\theta$.
    Montrer que $g$ est la somme d'une série entière de rayon supérieur
    ou égal à~$r$ et que $g$ coïncide avec $f$ sur $\mathring D(0,r)$.

  \bigskip

Applications~:
  \item $R\ge d$.
  \item Si $U=\C$ et $f$ est bornée alors $f$ est constante (théorème de Liouville).
    
  \item Si $P\in\C[X]$ ne s'annule pas alors $P$ est constant (théorème de d'Alembert-Gauss).
    
  \item Si $(f_n)$ est une suite de fonctions analytiques convergeant uniformément
    sur~$U$ vers une fonction~$f$ alors $f$ est analytique sur~$U$ (théorème de Weierstrass, comparer avec le cas réel).
    
  \item La composée de deux fonctions analytiques est analytique.
    
\end{enumerate}
\finenonce{004614}



\finexercice


\exercice{4615, quercia, 2010/03/14}
\enonce{004615}{Formule des résidus}

Soit $P\in\C[X]$ ayant pour racines $z_1,\dots,z_k$ de multiplicités $m_1,\dots,m_k$
    et $r\in\R^{+*}\setminus\{|z_1|,\dots,|z_k|\}$.

Montrer~:
    $\frac1{2\pi} \int_{\theta=0}^{2\pi}\frac{P'(re^{i\theta})}{P(re^{i\theta})}re^{i\theta}\, d\theta
    = \sum_{|z_j|<r} m_j$.
\finenonce{004615}



\finexercice


\exercice{4616, quercia, 2010/03/14}
\enonce{004616}{Croissance de $f$ en fonction des coefficients}

Soit $f(z) = \sum_{n=0}^\infty a_nz^n$ une série entière de rayon de
convergence infini. Montrer l'équivalence entre les propriétés :

\indent 1: Pour tout $a > 0$, la fonction $z \mapsto f(z)e^{-a|z|}$ est bornée sur
           $\C$.\par
\indent 2: $\sqrt[n]{n!\,|a_n|} \to 0$ lorsque $n\to\infty$.\par

On utilisera les formules de Cauchy (cf. exercice~\ref{Cauchy}).

\finenonce{004616}



\finexercice


\exercice{4617, quercia, 2010/03/14}
\enonce{004617}{Centrale MP 2000}

Soit $(a_n)$ une suite réelle avec $a_0=0$ et $a_1=1$. On note 
$f(z)=\sum_{n=0}^{\infty}a_nz^n$. On suppose que $f$ est injective et que le rayon de convergence
de la série entière vaut 1. On considère $\Omega ^+=\{z\in \C\,|\, \Im z >0\}$ et
$D=\{z\in \C\,|\, |z|<1\}$.

\begin{enumerate}
  \item Montrer que, pour tout $z\in D$, $f(z)\in \R$ si et seulement si $z\in \R$.
    
  \item Montrer que $f(D\cap \Omega ^+)\subset \Omega ^+$.
    
  \item Montrer que, pour tout $|r|<1$, 
    $a_n=\frac{2}{\strut\pi r^n} \int_{\theta=0}^{\pi}\Im (f(re^{i\theta})) \sin (n\theta)\, d \theta$.
    
  \item Montrer que $|\sin (n\theta)|\le n|\sin \theta |$. En déduire que $|a_n|\le n$.
    
\end{enumerate}
\finenonce{004617}



\finexercice

\section{220.99 Autre}


\exercice{4618, quercia, 2010/03/14}
\enonce{004618}{Anneau des séries entières}

Soit $A$ l'ensemble des suites $(a_n)$ de complexes telles que
la série entière $\sum a_nz^n$ a un rayon non nul.
On munit~$A$ de l'addition terme à terme et du produit de Cauchy noté $*$.

\begin{enumerate}
  \item Vérifier que $A$ est un anneau intègre. Quels sont les éléments de~$A$
    inversibles~?
  \item Soit $I_k =\{a=(a_n)\in A\text{ tel que } a_0=\dots=a_k = 0\}$.
    Montrer que les idéaux de~$A$ sont $\{0\}$, $A$ et les $I_k$, $k\in\N$.
  \item Soit $f(x) = 2-\sqrt{\frac{1-2x}{\strut 1-x}}$. Montrer que $f$ est développable
    en série entière sur $]-\frac12,\frac12[$ et que si $f(x) = \sum_{n=0}^\infty u_nx^n$
    alors la suite $(u_n)$ vérifie la relation de récurrence~:
    $2u_{n+1} = 1 + \sum_{k=1}^n u_ku_{n+1-k}$.
  \item Soit $a=(a_n)\in A$ avec $a_0=1$ et $|a_n|\le 1$ pour tout~$n$.
    Montrer qu'il existe une unique
    suite $b=(b_n)\in A$ telle que $b_0 = 1$ et $b*b = a$.
    Pour prouver que le rayon de convergence de~$b$ est non nul
    on établira par récurrence que $|b_n| \le u_n$.
  \item Pour $a\in A$ quelconque, étudier l'équation $b*b = a$ d'inconnue~$b\in A$.
\end{enumerate}
\finenonce{004618}



\finexercice


\exercice{4619, quercia, 2010/03/14}
\enonce{004619}{Ulm MP$^*$ 2000}

Soit $z_1,\dots,z_p\in \C$, $p_1,\dots,p_p \in \R^+$ tels que $\sum_{i=1}^p p_i=1$, et $\omega \in \R$.
Pour $n>p$ on pose $z_n=e^{i\omega}\sum_{j=1}^p z_{n-j}p_j$.
Étudier la suite $(z_n)$.
\finenonce{004619}



\finexercice


\exercice{4620, quercia, 2010/03/14}
\enonce{004620}{X MP$^*$ 2001}

Soit $D$ le disque ouvert de~$\C$ de centre~$0$ et rayon~$1$.

\begin{enumerate}
  \item Soit $\varphi(z) = \sum_{n\in\N}a_nz^n$ une série entière de rayon $R\ge 1$
    et $r\in{]0,1[}$.
    Montrer que $$a_n = \frac1{2\pi r^n} \int_{\theta=0}^{2\pi}\varphi(re^{i\theta})e^{-in\theta}\, d\theta.$$

  \item Soit $E$ l'ensemble des fonctions de~$\overline D$ dans~$\C$ continues
    et dont la restriction à~$D$ est somme d'une série entière. Montrer que
    $f \mapsto\|f\| = \sup\{|f(z)|,\ z\in\overline D\}$ définit une norme sur~$E$ et
    que pour cette norme $E$ est complet.
  \item Montrer que l'ensemble des polynômes à coefficients complexes est dense dans~$E$.
    
\end{enumerate}
\finenonce{004620}



\finexercice


\exercice{4621, quercia, 2010/03/14}
\enonce{004621}{Centrale MP 2002}

\begin{enumerate}
  \item Développer en série entière $f$ : $z \mapsto z(1-z)^{-2}$. Montrer que $f$
    est injective sur $\mathring D(0,1)$.

  \item Soit $f(z)=z+\sum_{n=2}^{+\infty} a_nz^n$ la somme d'une série entière de rayon de 
    convergence au moins $1$ à coefficients réels.
    On suppose $f$ injective sur $\mathring D(0,1)$
    et on veut prouver~: $\forall\ n\ge 1, |a_n|\le n$.
  \begin{enumerate}
  \item Montrer pour $|z|<1$ que $f(z)\in\R\Leftrightarrow z\in\R$ et en déduire~:
    $\Im(z)\ge 0 \Rightarrow  \Im(f(z))\ge 0$.

  \item Pour $0<r<1$ calculer $ \int_{t=0}^{\pi }\Im(f(re^{it}))\sin nt\,d t$.
    En déduire  $|a_n|r^n\le n|a_1|r$ et conclure.

  \end{enumerate}
\end{enumerate}
\finenonce{004621}



\finexercice


\exercice{5749, rouget, 2010/10/16}
\enonce{005749}{*** I}
Soient $P_n =\sum_{k=0}^{n}\frac{X^k}{k!}$ et $R > 0$ donné. Montrer que pour $n$ suffisamment grand, $P_n$ n'a pas de racine dans le disque fermé de centre $0$ et de rayon $R$.
\finenonce{005749}


\finexercice


\exercice{5750, rouget, 2010/10/16}
\enonce{005750}{**** Inverse d'une série entière}
Soit  $\sum_{n=0}^{+\infty}a_nz^n$ une série entière de rayon $R > 0$ et telle que $a_0 = 1$ (ou plus généralement $a_0\neq 0$).

\begin{enumerate}
 \item  Montrer qu'il existe une et une seule suite $(b_n)_{n\in\Nn}$ telle que $\forall n\in\Nn$,  $\sum_{k=0}^{n}a_kb_{n-k}=\delta_{0,n}$.

\item  Montrer que la série entière $\sum_{n=0}^{+\infty}b_nz^n$ a un rayon strictement positif.
\end{enumerate}
\finenonce{005750}


\finexercice


\exercice{5755, rouget, 2010/10/16}
\enonce{005755}{***}
Soit $A$ une matrice carrée complexe de format $p\in\Nn^*$.
Rayon de convergence et somme en fonction de $\chi_A$ de la série entière $\sum_{n=0}^{+\infty}\text{Tr}(A^n)z^n$.
\finenonce{005755}


\finexercice


\exercice{5759, rouget, 2010/10/16}
\enonce{005759}{*** I}
 Soient $(a_n)_{n\in\Nn}$ et $(b_n)_{n\in\Nn}$ deux suites de réels strictement positifs telles que la suite $\left(\frac{a_n}{b_n}\right)_{n\in\Nn}$ ait une limite réelle $k$. (En particulier $a_n\underset{n\rightarrow+\infty}{=}o(b_n)$ si $k = 0$ et $a_n\underset{n\rightarrow+\infty}{\sim}b_n$ si $k = 1$).On suppose de plus que la série entière associée à la suite $(a_n)_{n\in\Nn}$ a un  rayon de convergence égal à $1$ et que la série de terme général $a_n$ diverge.

\begin{enumerate}
 \item  Montrer que $\lim_{x \rightarrow 1}\frac{\sum_{n=0}^{+\infty}a_nx^n}{\sum_{n=0}^{+\infty}b_nx^n}= k$.

\item \textbf{Applications.} 
  \begin{enumerate}
   \item Equivalent simple quand $x$ tend vers $1$ de $\sum_{n=1}^{+\infty}\ln nx^n$.
   \item Déterminer $\lim_{x \rightarrow 1}(1-x)^p\sum_{n=0}^{+\infty}n^{p-1}x^n$ où $p$ est un entier naturel non nul donné.
   \end{enumerate}
\end{enumerate}
\finenonce{005759}


\finexercice


\exercice{5760, rouget, 2010/10/16}
\enonce{005760}{}
Soit $(a_n)_{n\in\Nn}$ une suite à valeurs dans $\{-1,1\}$. Pour $x$ réel on pose $f(x) =\sum_{n=0}^{+\infty}\frac{a_n}{n!}x^n$.

On suppose que pour tout entier naturel $p$ et tout réel positif $x$, $|f^{(p)}(x)|\leqslant 1$. Déterminer $f$.
\finenonce{005760}


\finexercice


\exercice{5764, rouget, 2010/10/16}
\enonce{005764}{*** I Dénombrement de parenthésages}
\label{ex:rou20x}
\begin{enumerate}
 \item  Soit $E$ un ensemble non vide muni d'une loi interne et $a_n$ le nombre de parenthésages possibles d'un produit de $n$ éléménts de $E$ (($a_1 = 1$ conventionnellement), $a_2 = 1$, $a_3 = 2$, $a_4 = 5$, ...). Montrer que pour tout $n\geqslant 2$, $a_n =\sum_{k=1}^{n-1}a_ka_{n-k}$.

\item  Soit $f$ la série entière associée à la suite $(a_n)$. On suppose momentanément le rayon $R$ de cette série strictement positif.  Montrer que pour tout $x$ de $]-R,R[$, $(f(x))^2 -f(x)+ x = 0$.

\item  Calculer $R$ et $f$.

\item  En déduire $a_n$.
\end{enumerate}
\finenonce{005764}


\finexercice


\exercice{7546, mourougane, 2021/08/10}
\enonce{007546}{Estimations de reste}
Montrer que pour tout $N\in\Nn$, et tout $z\in\Delta$, 
$$\left|\exp(z)-\sum_0^N\frac{z^n}{n!} \right|\leq \frac{2}{(N+1)!}.$$
\finenonce{007546}
\finexercice


\exercice{7547, mourougane, 2021/08/10}
\enonce{007547}{Dérivation}
Montrer que pour tout entier naturel $k$ et tout $z\in\Delta$,
$$\frac{1}{(1-z)^{k+1}}=\sum_{n=k}^{+\infty} \binom{n}{k} z^{n-k}.$$
\finenonce{007547}
\finexercice
\section{221.01 Calcul de coefficients}


\exercice{1951, legall, 2003/10/01}

\enonce{001951}{}

  Soit $f$ la fonction $2\pi $-p\'eriodique sur $\Rr$ telle que 
$f(x)=\vert x \vert $ si $\vert x \vert \leq \pi .$

\begin{enumerate}

\item D\'eterminer la s\'erie de Fourier de $f$.
\item Calculer $\displaystyle{ \int _{-\pi}^\pi \vert x\vert ^2dx}$. 
En d\'eduire la valeur de $\displaystyle{ \sum _{p=0} ^{\infty } 
\frac{1}{(2p+1)^4}}.$

\item Calculer $\displaystyle{ \sum _{p=1}^\infty \frac{1}{n^4}}$.

\item Montrer que $\vert x\vert = \displaystyle{\frac{\pi }{2} 
-\frac{4}{\pi} \sum _{p=0}^\infty \frac{ \cos (2p+1)x }{(2p+1)^2}}$.
En d\'eduire les valeurs de
$\displaystyle{ \sum _{p=0}^\infty \frac{1}{(2p+1)^2}}$
puis  $\displaystyle{ \sum _{p=1}^\infty \frac{1}{n^2}}$.


\end{enumerate}
\finenonce{001951}



\finexercice



\exercice{1952, legall, 2003/10/01}

\enonce{001952}{}
  Soit $f$ la fonction $2\pi $-p\'eriodique sur $\Rr$ telle que 
$f(x)=x^2$ si $\vert x \vert \leq \pi .$
\begin{enumerate}
\item D\'eterminer la s\'erie de Fourier de $f$.
\item Calculer $\displaystyle{ \int _{-\pi}^\pi x^4dx}.$ En d\'eduire 
la valeur de $\displaystyle{ \sum _{p=1}^\infty \frac{1}{n^4}}.$
\item Montrer que $x^2= \displaystyle{\frac{\pi ^2}{3} +4 \sum 
_{n=1}^\infty \frac{(-1)^n \cos nx }{n^2}}.$
En d\'eduire $\displaystyle{ \sum _{p=1}^\infty \frac{1}{n^2}}.$
\end{enumerate}
\finenonce{001952}



\finexercice



\exercice{1955, legall, 2003/10/01}

\enonce{001955}{}
Soit $f: \Rr \rightarrow \Rr $ l'application $x\mapsto 
\displaystyle{\int _0^{+\infty}\frac{\cos tx}{1+t^2}dt}.$

\begin{enumerate}
\item Montrer que $\forall x\in \Rr : xf(x)=2\displaystyle{\int 
_0^{+\infty}\frac{t\sin tx}{(1+t^2)^2}dt}.$
\item Montrer que $f$ est de classe $C^2$ puis, en d\'erivant 
l'expression ci-dessus, que $\forall x\in \Rr _+^* : f''(x)=f(x).$

\item Donner une expression de $f$ sur $\Rr _+^*$ puis sur $\Rr$.

\end{enumerate}
\finenonce{001955}



\finexercice



\exercice{4622, quercia, 2010/03/14}
\enonce{004622}{Développements}

Calculer le développement des fonctions $f$ $2\pi$-périodiques telles que :

\begin{enumerate}
  \item $f(x) = \pi-|x|$ sur $]-\pi,\pi[$.
    
  \item $f(x) = \pi-x$ sur $]0,2\pi[$.
    
  \item $f(x) = x^2$ sur $]0,2\pi[$.
    
  \item $f(x) = \max(0,\sin x)$.
    
  \item $f(x) = |\sin x|^3$.
    
\end{enumerate}
\finenonce{004622}



\finexercice


\exercice{4623, quercia, 2010/03/14}
\enonce{004623}{Chimie P' 1996}

Établir la convergence puis calculer
$ \int_{t=0}^{\pi/2} \frac{\sin(nt)}{\sin t}\,d t$.

En déduire les coefficients de Fourier de $f$~: $f(t) = \ln|\tan(t/2)|$.


\finenonce{004623}



\finexercice


\exercice{4624, quercia, 2010/03/14}
\enonce{004624}{Chimie P 1996}

Développer en série de Fourier $f\ :\ t  \mapsto \frac1{1-\cos\alpha\cos t}$
avec $0<\alpha<\pi$.
Indication~: on pourra utiliser une relation de récurrence entre les
coefficients à partir de $(1-\cos\alpha\cos t)f(t) = 1$.
\finenonce{004624}



\finexercice


\exercice{4625, quercia, 2010/03/14}
\enonce{004625}{Mines MP 2002}

Soit $a\in {]-1,1[}$ et $g$ : $x \mapsto\frac{1-a\cos x}{ 1-2a\cos x +a^2}$.

\begin{enumerate}
  \item Prouver~: $\forall\ x\in\R,\ g(x)=\sum_{n=0}^{\infty} a^n\cos nx$.

  \item Quel est le mode de convergence de la série~?

  \item Soit $f : \R \to \C$ continue par morceaux et $2\pi$-périodique.
    Montrer que $h$ : $x \mapsto \int_{t=0}^{2\pi}g(x-t)f(t)\,d t$
    est somme d'une série trigonométrique uniformément convergente.
    Que peut-on déduire pour~$h$~?\vrule height 12pt width 0pt

  \item Soit $\lambda\in\R$. Trouver toutes les fonctions $f : \R \to \C$
    continues par morceaux et $2\pi$-périodiques telles que~:
    $\forall\ x\in\R,\ f(x)=\lambda \int_{t=0}^{2\pi}g(x-t)f(t)\,d t +\sum_{n=1}^{\infty}\frac{\cos nx}{n^2}$.

\end{enumerate}
\finenonce{004625}



\finexercice


\exercice{4626, quercia, 2010/03/14}
\enonce{004626}{Usage d'une série entière}

\begin{enumerate}
  \item Existe-t-il une fonction $f:\R \to \R$ continue telle que les coefficients
    de Fourier soient : $a_n = \frac1{2^n}$ et $b_n = 0$~?
  \item Application : calculer $ \int_{t=0}^\pi \frac{d t}{5-4\cos t}$.
    
\end{enumerate}
\finenonce{004626}



\finexercice


\exercice{4627, quercia, 2010/03/14}
\enonce{004627}{ $1/(\cos x + \ch a)$}

Soit $a > 0$.

\begin{enumerate}
  \item Développer en série entière : $f(x) = \frac1{x+e^a}$.
  \item En déduire le développement en série de Fourier de
    $g(x) = \frac1{\cos x + \ch a}$.
    

\end{enumerate}
\finenonce{004627}



\finexercice


\exercice{4628, quercia, 2010/03/14}
\enonce{004628}{Décomposition en $\sin^2$}

Montrer que : $\forall\ x \in \R,\ |\sin x| =
\frac8\pi\sum_{n=1}^\infty \frac{\sin^2nx}{4n^2-1}$.
\finenonce{004628}



\finexercice


\exercice{4629, quercia, 2010/03/14}
\enonce{004629}{DSF de $f*g$, Mines PSI 1998}

Soient ${f,g} : \R \to \C$ continues $2\pi$-périodiques. On pose pour $x \in \R$ :
$h(x) = \frac1{2\pi} \int_{t=0}^{2\pi} f(x-t)g(t)\,d t$.

\begin{enumerate}
  \item
Montrer que $h$ est $2\pi$-périodique, continue, et calculer les coefficients
de Fourier exponentiels de $h$ en fonction de ceux de $f$ et de $g$.


  \item
Pour $g$ fixée, déterminer les valeurs et vecteurs propres de
$f \mapsto h$.
\end{enumerate}
\finenonce{004629}



\finexercice


\exercice{4630, quercia, 2010/03/14}
\enonce{004630}{DSF d'une série}
On pose $f(x) = \sum_{n=-\infty}^{+\infty} e^{-(x-2n\pi)^2}$.
Montrer que $f$ est définie sur~$\R$, $2\pi$-périodique et de classe~$\mathcal{C}^1$.
Déterminer sa série de Fourier.

\finenonce{004630}



\finexercice


\exercice{4631, quercia, 2010/03/14}
\enonce{004631}{Calcul de séries}

Soit $f$ la fonction $2\pi$-périodique telle que :
$\forall\ x \in {[-\pi,\pi[},\ f(x) = e^x$.

\begin{enumerate}
  \item Chercher le développement en série de Fourier de $f$.
    

  \item En déduire les sommes des séries :
    $S = \sum_{n=1}^\infty \frac1{n^2+1}$ et
    $S' = \sum_{n=1}^\infty \frac{(-1)^n}{n^2+1}$.
    
\end{enumerate}
\finenonce{004631}



\finexercice


\exercice{4632, quercia, 2010/03/14}
\enonce{004632}{$\sum_{n=1}^\infty  \frac{a}{n^2+a^2}$ (Centrale MP 2003)}
\begin{enumerate}
  \item Soit~$a\in\R$. Développer en série de Fourier la fonction $2\pi$-périodique
    valant~$e^{ax}$ sur~$]0,2\pi]$.

Soit~$a\in\R$. On pose~$I(a) =  \int_{u=0}^{+\infty}\frac{e^{-u}}{1-e^{-u}}\sin(au)\,d u$.

  \item Exprimer $I(a)$ sous forme d'une série sans intégrale.
    
  \item Calculer $ \int_{u=0}^{+\infty}e^{-u}\sin(au)\,d u$.
    
  \item Conclure.
    
\end{enumerate}
\finenonce{004632}



\finexercice


\exercice{4633, quercia, 2010/03/14}
\enonce{004633}{$\sum_{n=1}^\infty  \frac{a}{n^2+a^2}$ (Centrale MP 2000)}

\begin{enumerate}
  \item Donner le développement en série de Fourier de la fonction $2\pi$-périodique
    définie sur~$]0,2\pi[$ par $f(x) = e^{ax}$ avec $a\ne 0$.
  \item Calculer $\sum_{n\ge1}\frac{a}{a^2+n^2}$. En déduire $\sum_{n\ge 1}\frac1{n^2}$.
  \item Que vaut $\lim_{a\to+\infty}\,\sum_{n\ge1}\frac{a}{a^2+n^2}$~?
\end{enumerate}
\finenonce{004633}



\finexercice


\exercice{4634, quercia, 2010/03/14}
\enonce{004634}{Calcul de séries, Matexo}

On considère la fonction $2 \pi$-périodique sur $\R$ définie par
$f(x) = x \sin {\frac x 2}\quad \text{ si } 0 \le x < 2\pi$.

\begin{enumerate}
  \item Calculer les coefficients de Fourier de $f$.


  \item Quelle est la nature de la série de Fourier $S_f$ de $f$ ?

  \item En déduire la somme de la série
$\sum^\infty_{n=1}\, \frac{(-1)^{n-1}}{4 n^2 - 1}$.

\end{enumerate}
\finenonce{004634}



\finexercice


\exercice{4635, quercia, 2010/03/14}
\enonce{004635}{$\sin(\pi a)/\pi a$}

Soit $a \in \R\setminus\Z$.

\begin{enumerate}
  \item Développer en série de Fourier la fonction définie sur $[-\pi,\pi]$ par :
    $f(x) = \cos(ax)$.
    
  \item Soit $g(t) = \sum_{n=1}^\infty \ln\left(1-\frac{t^2}{n^2}\right)$.
    Justifier l'existence et la dérivabilité de $g$ et la calculer.
    
\end{enumerate}
\finenonce{004635}



\finexercice


\exercice{4636, quercia, 2010/03/14}
\enonce{004636}{$\sum_{n=1}^\infty \frac{\sin n}n = \sum_{n=1}^\infty \frac{\sin^2 n}{n^2}$}

\begin{enumerate}
  \item Développer en série de Fourier la fonction $f$, $2\pi$-périodique
    telle que $f(x) = \frac{\pi-x}2$ pour $0 \le x < 2\pi$.
    
  \item Donner les développements en série de Fourier de $f(x+1)$ et $f(x-1)$.
  \item Montrer que $\sum_{n=1}^\infty \frac{\sin n}n
    = \sum_{n=1}^\infty \frac{\sin^2 n}{n^2}$.
    
\end{enumerate}
\finenonce{004636}


\finexercice


\exercice{4637, quercia, 2010/03/14}
\enonce{004637}{$f(x+\pi)$}

Soit $f:\R \to \R$ $2\pi$-périodique continue par morceaux.
Que peut-on dire des coefficients de Fourier de $f$ si l'on~a~:

\begin{enumerate}
  \item $\forall\ x \in \R,\ f(x+\pi) = f(x)$ ?
    
  \item $\forall\ x \in \R,\ f(x+\pi) = -f(x)$ ?
    
\end{enumerate}
\finenonce{004637}



\finexercice


\exercice{4638, quercia, 2010/03/14}
\enonce{004638}{$f$ est-elle $\pi$-périodique ?}

Soit $f \in {\cal D}$.
On note $c_k$ les coefficients de Fourier exponentiels de $f$.
Montrer que $f$ est $\pi$-périodique si et seulement si $c_k$ est nul pour tout
$k$ impair (noter que la série de Fourier de $f$ peut ne pas converger vers
$f$).
\finenonce{004638}



\finexercice


\exercice{4639, quercia, 2010/03/14}
\enonce{004639}{DSF de $f'$}

Soit $f:\R \to \R$ continue, $2\pi$-périodique de classe $\mathcal{C}^1$ par morceaux.
On note $a_k$, $b_k$ les coefficients de Fourier de~$f$.
Calculer les coefficients de Fourier de $f'$ en fonction de ceux de $f$.
En déduire que $ka_k \to 0$ et $kb_k \to 0$ lorsque $k\to\infty$.

\finenonce{004639}



\finexercice


\exercice{4640, quercia, 2010/03/14}
\enonce{004640}{DSF de $f'$}

Soit $f : {[0,2\pi]} \to \R$ de classe $\mathcal{C}^1$. On considère la fonction $g$,
$2\pi$-périodique coïncidant avec $f$ sur $[0,2\pi[$.
Soient $a_n,b_n$ les coefficients de Fourier de $g$.

\begin{enumerate}
  \item Montrer que $a_n =  o\left(\frac1n\right)$ et
    $b_n = \frac{f(0)-f(2\pi)}{\pi n} +  o\left(\frac1n\right)$.

  \item Donner le développement en série de Fourier
    de $g'$.
    
\end{enumerate}
\finenonce{004640}



\finexercice


\exercice{4641, quercia, 2010/03/14}
\enonce{004641}{DSF d'une primitive de $f$}

Soit $f$ continue $2\pi$-périodique, $F(x) =  \int_{t=0}^x f(t)\,d t$,
$a_n, b_n$ les coefficients de Fourier trigonométriques de~$f$ et
$C = \frac1{2\pi} \int_{t=0}^{2\pi} (\pi-t)f(t)\,d t$. Montrer~:
$$\forall\ x\in\R,\ F(x) =
\frac{a_0x}2 + C + \sum_{n=1}^\infty \frac{a_n\sin nx - b_n\cos nx}n.
$$
\finenonce{004641}



\finexercice


\exercice{4642, quercia, 2010/03/14}
\enonce{004642}{Concavité, ENS}

Soit $f:\R \to \R$ continue $2\pi$-périodique paire dont
la restriction à $[0,2\pi]$ est concave. Montrer que les
coefficients de Fourier trigonométriques de~$f$ vérifient~: $a_k \le 0$ pour
$k\ge 1$.
\finenonce{004642}



\finexercice


\exercice{5782, rouget, 2010/10/16}
\enonce{005782}{}
Développer en série de \textsc{Fourier} les fonctions suivantes puis déterminer la valeur des sommes indiquées :

\textbf{1) (**)} $f~:~\Rr\rightarrow\Rr$ $2\pi$-périodique paire telle que $\forall x\in[0,\pi]$, $f(x)=1-\frac{2x}{\pi}$. En déduire $\sum_{n=0}^{+\infty}\frac{1}{(2n+1)^2}$, $\sum_{n=1}^{+\infty}\frac{1}{n^2}$ et $\sum_{n=1}^{+\infty}\frac{1}{n^4}$.

\textbf{2) (**)} $f~:~\Rr\rightarrow\Rr$ $2\pi$-périodique impaire telle que $\forall x\in[0,\pi]$, $f(x)=x(\pi-x)$. En déduire $\sum_{n=0}^{+\infty}\frac{(-1)^n}{(2n+1)^3}$, $\sum_{n=0}^{+\infty}\frac{1}{(2n+1)^6}$ et $\sum_{n=1}^{+\infty}\frac{1}{n^6}$.

\textbf{3) (**)} $f~:~\Rr\rightarrow\Rr$ $2\pi$-périodique telle que $\forall x\in]-\pi,\pi]$, $f(x)=\sin\left(\frac{x}{2}\right)$. En déduire $\sum_{n=0}^{+\infty}(-1)^n\frac{2n+1}{16n^2+16n+3}$.

\textbf{4) (***)} $f~:~\Rr\rightarrow\Rr$ $2\pi$-périodique telle que $\forall x\in[-\pi,\pi]$, $f(x)=\ch(\lambda x)$ ($\lambda$ réel strictement positif donné). En déduire $\sum_{n=1}^{+\infty}\frac{(-1)^n}{\lambda^2+n^2}$, $\sum_{n=1}^{+\infty}\frac{1}{\lambda^2+n^2}$ et $\sum_{n=1}^{+\infty}\frac{1}{(\lambda^2+n^2)^2}$.

\textbf{5) (**)} $f~:~\Rr\rightarrow\Rr$ telle que $\forall x\in\Rr$, $f(x)=\text{sup}(0,\sin x)$. En déduire $\sum_{n=1}^{+\infty}\frac{1}{4p^2-1}$.

\finenonce{005782}


\finexercice


\exercice{5783, rouget, 2010/10/16}
\enonce{005783}{***}
Soit $a\in\Cc\setminus[-1,1]$.

\begin{enumerate}
 \item 
   \begin{enumerate}
   \item Développer en série trigonométrique la fonction $f~:~t\mapsto\frac{1}{a-\cos t}$ (utiliser la racine de plus petit module, notée $b$,  de l'équation $z^2-az+1=0$).
   \item La série obtenue est-elle la série de \textsc{Fourier} de $f$ ?
   \end{enumerate}

\item  Déduire de 1) la valeur des intégrales $I_n=\int_{0}^{\pi}\frac{\cos(nt)}{a-\cos t}\;dt$, $n\in\Nn$.
\end{enumerate}
\finenonce{005783}


\finexercice


\exercice{5784, rouget, 2010/10/16}
\enonce{005784}{*** I}
(un développement en série de fonctions de $\frac{\pi}{\sin(\pi z)}$ et $\mathrm{cotan}(\pi z)$).

Soit $\alpha\in\Cc\setminus\Zz$. Soit $f$ l'application de $\Rr$ dans $\Cc$, $2\pi$-périodique telles que $\forall x\in[-\pi,\pi]$, $f(x)=\cos(\alpha x)$.

\begin{enumerate}
 \item  Développer la fonction $f$ en série de \textsc{Fourier}.

\item  En déduire que pour tout $z\in\Cc\setminus\Zz$,

\begin{center}
$\frac{\pi}{\sin(\pi z)}=\frac{1}{z}+\sum_{n=1}^{+\infty}(-1)^n\frac{2z}{z^2-n^2}$ et $\pi\mathrm{cotan}(\pi z)=\frac{1}{z}+\sum_{n=1}^{+\infty}\frac{2z}{z^2-n^2}$.
\end{center}
\end{enumerate}
\finenonce{005784}


\finexercice


\exercice{5785, rouget, 2010/10/16}
\enonce{005785}{**}
Développer en série de \textsc{Fourier} la fonction $f~:~x\mapsto x-E(x)-\frac{1}{2}$. 
\finenonce{005785}


\finexercice

\section{221.02 Convergence, théorème de Dirichlet}


\exercice{4650, quercia, 2010/03/14}
\enonce{004650}{Phénomène de Gibbs pour $\sin kx/k$}

Soit $f_n(x) = \sum_{k=1}^n \frac{\sin kx}k$.

\begin{enumerate}
  \item Calculer l'abscisse, $x_n$, du premier maximum positif de $f_n$.
    

  \item Déterminer $\lim_{n\to\infty} f_n(x_n)$.
    

\end{enumerate}
\finenonce{004650}



\finexercice


\exercice{4651, quercia, 2010/03/14}
\enonce{004651}{Convergence uniforme}

Soit $\sum_{n=1}^\infty (a_n \cos nx + b_n\sin nx)$ une série trigonométrique
convergeant uniformément sur un intervalle $[\alpha,\beta]$.
Montrer que les suites $(a_n)$ et $(b_n)$ tendent vers 0.


\finenonce{004651}



\finexercice


\exercice{4652, quercia, 2010/03/14}
\enonce{004652}{Convergence uniforme}

Soit $(a_n)$ une suite décroissante de limite nulle.
Montrer que la série $\sum_{n=1}^\infty a_n\sin nx$ converge uniformément sur
$\R$ si et seulement si $na_n \to 0$ lorsque $n\to\infty$.

Pour le sens direct : utiliser le critère de convergence uniforme de Cauchy
et l'inégalité : $\sin x \ge \frac{2x}\pi$ sur $[0,\pi/2]$.
\finenonce{004652}



\finexercice


\exercice{4653, quercia, 2010/03/14}
\enonce{004653}{Fonction continue dont la série de Fourier diverge en 0}

\begin{enumerate}
  \item Soit $f : \R \to \C$ paire, $2\pi$-périodique, telle que, pour tout 
$x\in [0,\pi]$, $f(x)=\sum_{p=1}^{\infty}\frac{1}{p^2}\sin \Bigl((2^{p^3}+1)\frac{x}{2}\Bigr)$. 
Vérifier que $f$ est définie et continue sur $\R$.

  \item Soit $A_0 = \frac{1}{\pi} \int_0^{\pi} f(t)\,d t$, pour $n\in \N^*$, $A_n=\frac{2}{\pi} \int_0^{\pi} f(t) \cos (nt)\,d t$.

Pour $\nu \in \N$, on pose
$a_{0,\nu}=\frac12 \int_0^{\pi} \sin \Bigl((2\nu +1)\frac{t}{2}\Bigr)\,d t $ et
$a_{n,\nu}= \int_0^{\pi} \cos (nt)\sin \Bigl((2\nu +1)\frac{t}{2}\Bigr)\,d t $.

Pour $q\in \N$, on note 
$s_{q,\nu}=\sum_{i=0}^q a_{i,\nu}$. Montrer que si $\nu$ est fixé, $s_{n,\nu}\to0$ lorsque $n\to\infty$. Calculer 
explicitement les $a_{n,\nu}$. En déduire que, pour tout $q$, pour tout $\nu$, $s_{q,\nu}>0$, et prouver que
$\mathop{\max}\limits_{q\in \N} (s_{q,\nu})=s_{\nu,\nu}$.

  \item Montrer qu'il existe $B>0$ tel que, pour tout $\nu \ge 1$, $s_{\nu,\nu}\ge B \ln \nu$.

  \item Montrer que, pour tout $n\in \N$, $A_n=\frac{2}{\pi}\sum_{p=1}^{\infty} \frac{1}{p^2}a_{n,2^{p^3-1}}$.

  \item Pour $n\in \N ^*$, on pose $T_n=\frac{\pi}{2} \sum_{k=0}^{n} A_k$. Vérifier que 
$T_n=\sum_{p=1}^{\infty} \frac{1}{p^2} s_{n,2^{p^3-1}}$. Montrer qu'il existe $D>0$ tel que, pour 
tout $p\ge 1$, $T_{2^{p^3-1}}\ge Dp$, et constater que la série de Fourier de $f$ diverge au point $0$.
\end{enumerate}
\finenonce{004653}



\finexercice


\exercice{4654, quercia, 2010/03/14}
\enonce{004654}{$\sum_{n=-\infty}^\infty R(n)e^{inx}$, $R =$ fraction rationnelle}

Soit $R$ une fraction rationnelle à coefficients complexes, de degré strictement
négatif, n'ayant pas de pôle dans $\Z$.
On pose $f(x) = \sum_{n=-\infty}^\infty R(n)e^{inx}$.

\begin{enumerate}
  \item \'Etudier l'existence et la continuité de $f$.
    

  \item Montrer que $f$ est de classe $\mathcal{C}^\infty$ sur $\R\setminus 2\pi\Z$.
    
\end{enumerate}
\finenonce{004654}



\finexercice


\exercice{4655, quercia, 2010/03/14}
\enonce{004655}{$\sum_{n=1}^\infty \frac{\cos nx}{P(n)}$, $P =$ polynôme}

\begin{enumerate}
  \item Donner le développement en série de Fourier de la fonction $f$
    $2\pi$-périodique telle que $f(x) = (\pi-x)^2$ sur $]0,2\pi[$.
    
  \item Soit $P$ un polynôme de degré 2 sans racines dans $\N^*$.
    On pose $g(x) = \sum_{n=1}^\infty \frac{\cos nx}{P(n)}$.
    Montrer que $g$ est de classe $\mathcal{C}^1$ par morceaux.
    
\end{enumerate}
\finenonce{004655}



\finexercice


\exercice{4656, quercia, 2010/03/14}
\enonce{004656}{Noyau de Féjer}

Soit $f : \R \to \C$ $2\pi$-périodique continue, $f_n$ sa
$n$-ème somme de Fourier et $g_n = \frac{f_0 + \dots + f_n}{n+1}$.

\begin{enumerate}
  \item Exprimer $g_n$ à l'aide d'un produit de convolution, $g_n = f*k_n$.
    
  \item Montrer que la suite $(k_n)$ constitue une suite d'approximations de la
    mesure de Dirac sur $]-\pi,\pi[$. Ceci montre que la moyenne des
    sommes partielles de la série de Fourier de~$f$ converge uniformément
    vers~$f$ {\it pour toute $f$ continue}.
\end{enumerate}
\finenonce{004656}



\finexercice

\section{221.03 Formule de Parseval}


\exercice{1953, legall, 2003/10/01}

\enonce{001953}{}
Soit $f : \Rr \rightarrow \Rr $ une fonction $2\pi $-p\'eriodique, 
$C^2$ et telle que
$\displaystyle{ \int _0^{2\pi}f(t)dt=0}$ et, $\forall t\in [0,2\pi ], 
\vert f(t)\vert \geq \vert f''(t)\vert .$On note respectivement 
$(c_n)_{n\in \Zz}$ et $(c''_n)_{n\in \Zz}$ les
coefficients de Fourier (complexes) de $f$ et $f''$.
\begin{enumerate}
\item Calculer $c_0$ puis calculer $c''_n$ en fonction de $c_n$.
\item A l'aide du th\'eor\`eme de Parseval, en d\'eduire que $c_n=0$
pour $\vert n\vert \geq 2.$
\item Montrer qu'il existe $\varphi \in [0,2\pi ]$ et $\rho \in \Rr 
_+$ tels que $f(t)=\rho \cos (t+\varphi )$ pour tout $t\in [0,2\pi ].$


\end{enumerate}
\finenonce{001953}



\finexercice



\exercice{1954, legall, 2003/10/01}

\enonce{001954}{}
Soit $f : \Rr \rightarrow \Cc $ une fonction $2\pi $-p\'eriodique, 
$C^1$ et telle que
$\displaystyle{ \int _0^{2\pi}f(t)dt=0}.$ On note respectivement 
$(c_n)_{n\in \Zz}$ et $(c'_n)_{n\in \Zz}$ les coefficients de Fourier
(complexes) de $f$ et $f'$.
\begin{enumerate}
\item Calculer $c_0$ puis donner une relation entre $c_n$ et $c'_n$.
\item En d\'eduire que que
$\displaystyle{ \int _0^{2\pi}\vert f(t)\vert ^2dt\leq \int 
_0^{2\pi}\vert f'(t)\vert ^2dt}.$
\item Dans quel cas l'\'egalit\'e a-t-elle lieu ?

\end{enumerate}
\finenonce{001954}



\finexercice



\exercice{4643, quercia, 2010/03/14}
\enonce{004643}{ENS MP 2002}

Soit $f: {[0,1]} \to \R$ de classe $\mathcal{C}^2$ telle que $f(0)=f(1)=0$.

\begin{enumerate}
  \item Montrer que l'on peut prolonger $f$ en une fonction impaire et $2$-périodique.
 
  \item En déduire l'existence de $c>0$ indépendant de $f$ tel que $\|f\|_{\infty}\le c\|f''\|_2$.

\end{enumerate}
\finenonce{004643}



\finexercice


\exercice{4644, quercia, 2010/03/14}
\enonce{004644}{Inégalité de Wirtinger}

Soit $f : {[0,2\pi]} \to \R$ de classe $\mathcal{C}^1$
telle que $ \int_{t=0}^{2\pi} f(t)\,d t = 0$ et $f(0) = f(2\pi)$.

Montrer que $ \int_{t=0}^{2\pi} f^2(t)\,d t \le
 \int_{t=0}^{2\pi} f'^2(t)\,d t$ et déterminer les cas d'égalité.

\finenonce{004644}



\finexercice


\exercice{4645, quercia, 2010/03/14}
\enonce{004645}{Inégalité isopérimétrique}

\begin{enumerate}
  \item Soient $f,g$ deux applications $2\pi$-périodiques réelles de classe $\mathcal{C}^1$.
Montrer que~: $2 \int_0^{2\pi}fg' \le  \int_0^{2\pi}f'^2 +  \int_0^{2\pi}g'^2$.


  \item Soit $\Gamma$ un arc $\mathcal{C}^1$, fermé, simple, de longueur $2\pi$.
Montrer que l'aire du domaine limité par~$\Gamma$ est inférieure ou égale
à~$\pi$.
\end{enumerate}
\finenonce{004645}



\finexercice


\exercice{4646, quercia, 2010/03/14}
\enonce{004646}{$|f''| \le |f|$}

Trouver les fonctions $f:\R \to \R$ $2\pi$-périodiques de classe $\mathcal{C}^2$
telles que $ \int_{t=0}^{2\pi}f(t)\,d t = 0$ et $|f''| \le |f|$.
\finenonce{004646}



\finexercice


\exercice{4647, quercia, 2010/03/14}
\enonce{004647}{Calcul de $(f\mid g)$}

Soient ${f,g} : \R \to \C$ $2\pi$-périodiques, continues par morceaux.
On note $c_n(f)$ et $c_n(g)$ les coefficients de Fourier exponentiels de
$f$ et $g$.
Montrer que :  $\sum_{n=-\infty}^{+\infty} \overline{c_n(f)}c_n(g) =
        \frac1{2\pi} \int_{t=0}^{2\pi} \overline{f(t)}g(t)\,d t.$

\finenonce{004647}



\finexercice


\exercice{4648, quercia, 2010/03/14}
\enonce{004648}{Une série trigonométrique qui n'est pas une série de Fourier}

On pose $f(x) = \sum_{n=1}^\infty \frac{\sin nx}{\sqrt n}$.

\begin{enumerate}
  \item Montrer que $f$ est bien définie sur~$\R$, $2\pi$-périodique et
continue sur $\R\setminus2\pi\Z$.

  \item Calculer $\lim_{a\to0^+} \int_{t=a}^{\pi-a}f(t)\sin(pt)\,d t$
    et en déduire que $f$ n'a pas de développement en série de
    Fourier (et donc n'est pas continue en~$0$).
\end{enumerate}
\finenonce{004648}



\finexercice


\exercice{4649, quercia, 2010/03/14}
\enonce{004649}{X MP$^*$ 2001}

Soit $a>0$ et $f$ continue sur $[0,a]$ à valeurs réelles. On suppose que pour tout $x\in\R$
on a $ \int_{t=0}^a f(t)\cos(xt)\,d t=0$. Montrer que $f$ est nulle.

\finenonce{004649}



\finexercice
\section{221.99 Autre}


\exercice{1950, legall, 2003/10/01}

\enonce{001950}{}
Soit $f$ une fonction int\'egrable au sens de Riemann p\'eriodique de 
p\'eriode $2\pi $. On d\'esigne par~:
$\displaystyle{ \frac{a_0}{2}+ \sum _{k=1}^\infty(a_k \cos (kx)+ b_k 
\sin (kx))}$ sa s\'erie de Fourier et on pose, pour tout $n\in \Nn :
\displaystyle{ S_n(x)=\frac{a_0}{2}+ \sum _{k=1}^n(a_k \cos (kx)+ b_k 
\sin (kx))}.$

\begin{enumerate}
\item Soit $\theta \in \Rr-2\pi \Zz .$ Montrer $\displaystyle{ 
\frac{1}{2}+ \sum _{k=1}^n \cos (k\theta)=\frac{\sin 
(n+\frac{1}{2})\theta}
{2\sin \frac{\theta }{2}}}.$

\item
Etablir que $
\displaystyle{ S_n(x)=\frac{1}{\pi}\int _{-\pi}^\pi \frac {\sin 
(n+\frac{1}{2})(x-t)}{2\sin \frac{(x-t)}{2}}f(t)dt}$.

\item En d\'eduire $ \displaystyle{ S_n(x)=\frac{1}{2\pi}\int 
_{-\pi}^\pi f(x+\theta ) \frac{\sin (n+\frac{1}{2})\theta}
{\sin \frac{\theta }{2}}d\theta } $.

\item Calculer $\displaystyle{ \int _{0}^\pi \frac{\sin (n+\frac{1}{2})\theta}
{\sin \frac{\theta }{2}}d\theta } .$

\end{enumerate}
\finenonce{001950}



\finexercice



\exercice{4657, quercia, 2010/03/14}
\enonce{004657}{Formule sommatoire de Poisson}

Soit $f : \R \to \C$ de classe $\mathcal{C}^1$. On suppose qu'il existe $a>1$ tel que
$f(x) =  O(1/|x|^a)$ et $f'(x) =  O(1/|x|^a)$ lorsque $|x|\to\infty$, et on pose
$F(x) = \sum_{n\in \Z}f(x+2n\pi)$.

Montrer que $F$ est bien définie, $\mathcal{C}^1$ et $2\pi$-périodique.
En déduire la formule sommatoire de Poisson~:
$$\sum_{n\in\Z}f(2n\pi) = \frac1{\sqrt{2\pi}} \sum_{n\in\Z} \hat f(n).$$
\finenonce{004657}



\finexercice


\exercice{4658, quercia, 2010/03/14}
\enonce{004658}{Formule d'échange}

Soient ${f,g} : \R \to \C$ de classe $\mathcal{C}^1$ telles que $f,f',g,g'$ sont
intégrables. Montrer~:
$$ \int_{t=-\infty}^{+\infty}\hat f(t)g(t)\,d t =
   \int_{t=-\infty}^{+\infty} f(t)\hat g(t)\,d t.$$
\finenonce{004658}



\finexercice


\exercice{4659, quercia, 2010/03/14}
\enonce{004659}{$ \int_{t=a}^b f(t)|\sin nt|\,d t$}

\begin{enumerate}
  \item Développer en série de Fourier la fonction : $x  \mapsto |\sin x|$.
    
  \item Application : Soit $f : {[a,b]} \to \R$ continue.
    Montrer que $ \int_{t=a}^b f(t)|\sin nt|\,d t \to
    \frac2\pi \int_{t=a}^b f(t)\,d t$ lorsque $n\to\infty$.

\end{enumerate}
\finenonce{004659}



\finexercice


\exercice{4660, quercia, 2010/03/14}
\enonce{004660}{\'Equation différentielle}

Montrer que l'équation : $y^{(4)}+y''+y = |\sin x|$ admet une et une seule
solution $\pi$-périodique.

\finenonce{004660}



\finexercice


\exercice{4661, quercia, 2010/03/14}
\enonce{004661}{\'Equation différentielle}

Soit $k \in \R$.
Résoudre l'équation différentielle $y'' + k^2y = \sum_{n=1}^\infty \frac{\cos nx}{n^2}$.

\finenonce{004661}



\finexercice


\exercice{4662, quercia, 2010/03/14}
\enonce{004662}{\'Equirépartition modulo 1}



\begin{enumerate}
  \item Soit $f:\R \to \R$ de classe $\mathcal{C}^1$ 1-périodique,
    $\alpha \in \R\setminus \Q$ et $x \in \R$.

    Montrer que $\frac{f(x) + f(x+\alpha) + \dots + f(x+n\alpha)}{n+1}
    \to  \int_{t=0}^1 f(t)\,d t$ lorsque $n\to\infty$.
    

  \item Montrer que le résultat est encore vrai en supposant seulement $f$ continue.
    

  \item En déduire la nature de la série $\sum_{n=1}^\infty \frac{\sin^2 n}n$.
    

\end{enumerate}
\finenonce{004662}



\finexercice


\exercice{4663, quercia, 2010/03/14}
\enonce{004663}{Cachan MP$^*$ 2000}

Soit un réel $\beta > 1$ et
$a_k = \iint_{[0,1]^2}e^{-|x-x'|^\beta}e^{2i\pi k(x-x')}\,d x\,d x'$.
Trouver un équivalent quand $n$ tend vers l'infini de
$\sum_{|k|>n\text{ ou }|\ell|>n} a_ka_\ell$,
$k$ et~$\ell$ étant des entiers relatifs.
\finenonce{004663}



\finexercice


\exercice{4664, quercia, 2010/03/14}
\enonce{004664}{Algèbre de séries trigonométriques}
Soit $E$ l'ensemble des fonctions de~$\R$ dans~$\C$ de la forme~:
$f(x)=\sum_{n=-\infty}^{+\infty} c_ne^{2i\pi nx}$ où $\sum |c_n|$ converge.
On pose pour $f\in E$~: $\|f\| = \sum_{n=-\infty}^{+\infty}|c_n|$.

\begin{enumerate}
  \item Justifier la définition de $\|f\|$ et montrer que $E$ est un espace vectoriel normé complet.

  \item Montrer que $E$ est une $\C$-algèbre et que $\|fg\| \le \|f\|\,\|g\|$.
    

  \item Soit $\varphi : E \to \C$ un morphisme d'algèbres.
  \begin{enumerate}
    \item On suppose $\varphi$ continu, montrer qu'il existe $z_0\in\mathbb{U}$ tel
    que $\forall\ f\in E,\ \varphi(f) = \sum_{n=-\infty}^{+\infty}c_nz_0^n$.
    

    \item Vérifier que la formule précédente définit effectivement un morphisme
    continu de~$E$ dans~$\C$.
  \end{enumerate}
\end{enumerate}
\finenonce{004664}



\finexercice


\exercice{4665, quercia, 2010/03/14}
\enonce{004665}{Mines MP 2002}

Déterminer la nature de la série de terme général
$u_n=(-1)^n  \int_{t=0}^1 \cos (nt^2)\,d t$.

\finenonce{004665}



\finexercice


\exercice{4666, quercia, 2010/03/14}
\enonce{004666}{Ens Lyon MP$^*$ 2003}
On note~: 
$E^{\phantom 1} = \{\text{fonctions continues }2\pi\text{-périodiques }f : \R \to \C\}$ ;\par
$E^1 = \{f\in E \text{ de classe }\mathcal{C}^1\}$ ;\par
$E_n = \{f\in E\text{ tel que }\forall\ k\in[[-n,n]],\ {\textstyle \int}_{t=0}^{2\pi}f(t)e^{-ikt}\,d t = 0\}$ ;\par
$E_n^1 = E_n \cap E^1.$

On considère sur~$E$ la norme~$\|\ \|_2$ $\Bigl(\|f\|_2 = \sqrt{\frac1{2\pi}\int |f|^2}\;\Bigr)$.

\begin{enumerate}
  \item Montrer que~$D : {E_0^1} \to {E_0},  f  \mapsto {f'}$ est une bijection.
    
  \item $D$ est-elle continue~?
    
  \item Montrer que~$D^{-1}$ est continue.
    
  \item Montrer que~$D^{-1}(E_n) = E_n^1$ et calculer $\|\kern-1.2pt|D^{-1}_{|E_n}\|\kern-1.2pt|$.
        
\end{enumerate}
\finenonce{004666}



\finexercice


\exercice{4667, quercia, 2010/03/14}
\enonce{004667}{Quatre racines, ENS Cachan MP$^*$ 2005}

Soit $f$ à valeurs réelles, de classe $\mathcal{C}^{2}$, $2\pi$-p{\'e}riodique, de moyenne nulle. Montrer que $g=f+f''$ s'annule au moins quatre fois sur $[0,2\pi[$.

\finenonce{004667}



\finexercice


\exercice{5781, rouget, 2010/10/16}
\enonce{005781}{**}
\begin{enumerate}
 \item  Soit $f$ la fonction définie sur $\Rr$, $2\pi$-périodique et impaire telle que $\forall x\in\left]0,\frac{\pi}{2}\right[$, $f(x)=\sin\left(\frac{x}{2}\right)$. Déterminer $f(x)$ pour tout réel $x$.

\item  Soit $f$ la fonction définie sur $\Rr$, $2\pi$-périodique et paire telle que $\forall x\in\left[0,\frac{\pi}{2}\right]$, $f(x)=\sin\left(\frac{x}{2}\right)$. Déterminer $f(x)$ pour tout réel $x$.
\end{enumerate}
\finenonce{005781}


\finexercice

\section{222.01 Convergence simple, uniforme, normale}


\exercice{2683, matexo1, 2002/02/01}

\enonce{002683}{}
A 

\begin{enumerate}
\item Soient $a$ et $z$ deux réels. Soit $f$ une fonction de
classe $C^{n+1}$ sur le segment d'extrémités
$a$ et $z$ et $\phi$ un polynôme de degré $n$. Prouver que pour tout $t$ compris dans
l'intervalle $[0,1]$,
\[
{d\over dt}\sum_{m=1}^n (-1)^m (z-a)^m
                \phi^{(n-m)}(t) f^{(m)}\left(a+t(z-a)\right)
\]\[
  = -(z-a) \phi^{(n)}(t) f'\left(a+t(z-a)\right)
    + (-1)^n (z-a)^{n+1} \phi(t) f^{(n+1)}\left(a+t(z-a)\right) 
\]

\item 

\begin{enumerate}
\item Montrer que la fonction $t\mapsto {t \over e^t-1}$ est prolongeable par continuité
en zéro, que son prolongement est indéfiniment dérivable et admet des développements
limités en zéro de la forme:
$$ 1 - {t/2} + {b_1 t^2 \over 2!}
                                + {b_2 t^4 \over 4!} + \dots + {b_n t^{2n} \over (2n)!}
+o(t^{2n+1}),$$
où les $b_i$ sont des réels qu'on ne cherchera pas à déterminer.

Montrer que la dérivée $n^{i\grave eme}$ en zéro, notée $\phi_n(z)$, de la fonction
$t\mapsto t {e^{zt}-1 \over e^t-1}$
est un polynôme en $z$ de degré $n$ et que
$$\phi_n(z)= z^n - {1\over 2}nz^{n-1} + C_n^2 b_1 z^{n-2}
                + C_n^4 b_2 z^{n-4} + \dots
                + C_n^{2N} b_N z^{n-2N}$$
où $N=E({n-1\over2})$, $E$ désignant la fonction partie entière.

\item
Prouver que $nz^{n-1}=\phi_n(z+1) - \phi_n(z)$
\end{enumerate}

\item 
Prouver que
$$
\begin{matrix}
(i)     &\phi_n^{(n-k)}(1) = \phi_n^{(n-k)}(0)\ \ (2\leq k\leq n) &\quad
& (ii)  &\phi_n^{(n-2k-1)}(0) = 0 \ \ (1\leq k\leq N)             \cr
(iii)   &\phi_n^{(n-2k)}(0) = {n! b_k \over (2k)!}  \ \ (1\leq k\leq N)  &\quad
& (iv)  &\phi_n^{(n-1)}(0) = -{1\over2} n!         \cr
 (v)    &\phi_n^{(n-1)}(1) = +{1\over2} n! & \quad
&(vi)   &\phi_n^{(n)} = n!  \cr
\end{matrix}
$$

\item 

\begin{enumerate}
\item 
On suppose $f$ de classe $C^{2n+1}$. Prouver que
\begin{eqnarray*}
0 &=& f(z) - f(a) - {z-a \over 2} \left[ f'(z) + f'(a) \right] 
\\
&&+ \sum_{m=1}^{n-1} b_m {(z-a)^{2m} \over (2m)!} \left[ f^{(2m)}(z) - f^{(2m)}(a) \right] 
\\
&& - {(z-a)^{2n+1} \over (2n)!} \int_0^1 \phi_{2n}(t) f^{(2n+1)}\Big( a+(z-a)t \Big) dt
 \end{eqnarray*}

\item
En déduire que si $F$ est de classe $C^{2n}$ sur $[a,a+r\omega]$
où $r\in \Nn$ et $\omega>0$, alors
\begin{eqnarray*}
\int_a^{a+r\omega} F(x)dx &=&
\omega \Big[ {1\over2} F(a) + F(a+\omega) + \dots 
+ F(a+{\scriptstyle (r-1)}\omega) + {1\over2} F(a+r\omega) \Big]
\\
&& - \sum_{m=1}^{n-1}      b_m {\omega^{2m} \over (2m)!} \left[ F^{(2m-1)}(a+r\omega) - F^{(2m-1)}(a) \right] 
+ R_n                                               
 \end{eqnarray*} 
où 
\[
R_n={\omega^{2n+1} \over (2n)!} \int_0^1 \phi_{2n}(t)
\sum_{m=0}^{r-1} F^{(2n)}(a+m\omega+\omega t) dt \; .
\]
\end{enumerate}
\end{enumerate}
                        
B

\begin{enumerate}
\item
Soit $u_k: x>0 \mapsto \ln(x+k)-\ln(k)+x\ln\left({k\over k+1}\right)
                                \qquad(k\in \Nn^*)$\par
Montrer que pour tout $x$ strictement positif, la série
$\sum_{k\geq 1}u_k(x)$ est convergente. On pose pour la suite
$G(x)=\ln(x)+\sum_{k=1}^{\infty}u_k(x)$

\item
Prouver que $G$ vérifie l'équation fonctionnelle 
$$\forall x>0 \quad G(x+1)=G(x)-\ln(x).$$

\item
En déduire que $\forall m\in \Nn \quad \exp(-G(m+1))=m!$

\item
Soit $x$ et $y$ deux réels strictement positifs. Montrer que la
série\par
\centerline{$\sum_{k\geq1}\bigl[ \ln(y+k) - \ln(x+k)
                                + (y-x)\ln\left({k\over k+1}\right) \bigr]$}
est convergente et que sa somme est $G(y) - G(x) - \ln y + \ln x$.

\item
Prouver à l'aide de A que pour tous entiers strictement positifs $n$
et $p$ 
\begin{eqnarray*}
\sum_{k=0}^n \ln(y+k) - \ln(x+k) &=&
\int_0^n f(t)dt + {1\over2} \big( f(0)+f(n) \big)
\\
&&+\sum_{h=1}^{p-1} {b_h\over(2h)!} \big( f^{(2h-1)}(n) - f^{(2h-1)}(0) \big) + T_{p,n}(x,y)
 \end{eqnarray*}
où $f: t\mapsto \ln(y+t) -\ln(x+t)$ et
$T_{p,n}(x,y)$ est une expression que l'on précisera.

\item
Prouver que $R_p(x,y)=\lim_{n \to +\infty}T_{p,n}(x,y)$ existe.

\item
On pose $g(z)=z\ln z - z - {1\over2} \ln z + \sum_{h=1}^{p-1} 
                                {b_h\over(2h)(2h-1)} {1\over z^{2h-1}}$.
Montrer que $G(y)+g(y)=G(x)+g(x)+R_p(x,y)$

\item
Montrer que $ R_p(x,y)  = O\left( {1 \over \bigl[ \inf(x,y) \bigr]^{2p-1}} \right)$
quand $\inf(x,y)\to +\infty$.

\item
Prouver à l'aide de la formule de Stirling que
$G(m)+g(m)\to {1\over2}\ln2\pi$ quand
$m \to +\infty$.

\item
Montrer que
$$G(y)= -y\ln y + y + {1\over2}\ln y - {1\over2}\ln2\pi
- \sum_{h=1}^{p-1}  {b_h\over(2h)(2h-1)} {1\over y^{2h-1}}
+ O({1\over y^{2p-1}})$$

\item
Donner un développement asymptotique de $\ln(m!)$ quand $m$ tend vers $+\infty$ à
un $O({1\over m^7})$ près.
\end{enumerate}

\finenonce{002683}
\finexercice


\exercice{4503, quercia, 2010/03/14}
\enonce{004503}{\'Etude de convergence}

Soit $\alpha \in \R$ et $f_n(x) = n^\alpha x(1-x)^n$ pour $x \in {[0,1]}$.

\begin{enumerate}
  \item Trouver la limite simple des fonctions $f_n$.
  \item Y a-t-il convergence uniforme ?
    

\end{enumerate}
\finenonce{004503}



\finexercice


\exercice{4504, quercia, 2010/03/14}
\enonce{004504}{\'Etude de convergence}

On pose $f_n(x) = x^n(1-x)$ et $g_n(x) = x^n\sin(\pi x)$.

\begin{enumerate}
  \item Montrer que la suite $(f_n)$ converge uniformément vers la fonction nulle sur
    $[0,1]$.
  \item En déduire qu'il en est de même pour la suite $(g_n)$.
    (On utilisera la concavité de sin sur $[0,\pi]$)


\end{enumerate}
\finenonce{004504}



\finexercice


\exercice{4505, quercia, 2010/03/14}
\enonce{004505}{Non interversion limite-intégrale}

Soit $f_n(x) = n\cos^nx\sin x$.

\begin{enumerate}
  \item Chercher la limite simple, $f$, des fonctions $f_n$.
  \item Vérifier que $ \int_{t=0}^{\pi/2} f(t)\,d t\ne
          \lim_{n\to\infty}  \int_{t=0}^{\pi/2} f_n(t)\,d t$.

\end{enumerate}
\finenonce{004505}



\finexercice


\exercice{4506, quercia, 2010/03/14}
\enonce{004506}{Non interversion limite-intégrale}

\begin{enumerate}
  \item Déterminer la limite simple des fonctions
    $f_n : x  \mapsto \frac{x^ne^{-x}}{n!}$
    sur $\R^+$ et montrer qu'il y a convergence uniforme.
    (On admettra la formule de Stirling : $n! \sim n^ne^{-n}\sqrt{2\pi n}$)
  \item Calculer $\lim_{n\to\infty} \int_{t=0}^{+\infty} f_n(t)\,d t$.
    
\end{enumerate}
\finenonce{004506}



\finexercice


\exercice{4507, quercia, 2010/03/14}
\enonce{004507}{\'Etude de convergence}

Soit ${f_n} : {[0,+\infty[} \to   {\R}$, $\begin{cases}
                     x \le n             &(1- x/n)^n \cr
                      x > n               &0.         \cr\end{cases}$

\begin{enumerate}
  \item Déterminer la limite simple, $f$, des fonctions $f_n$. 
  \item Montrer que : $\forall\ x\in \R^+,\ 0 \le f_n(x) \le f(x)$.
  \item Montrer que $(f_n)$ converge uniformément vers $f$ sur tout segment $[0,a]$.
  \item Démontrer que la convergence est uniforme sur $\R^+$.
\end{enumerate}
\finenonce{004507}



\finexercice


\exercice{4508, quercia, 2010/03/14}
\enonce{004508}{\'Etude de convergence}

\'Etudier la convergence simple, uniforme, de la suite de fonctions :
$f_n : x  \mapsto \left(1+\frac xn\right)^{-n}$.


\finenonce{004508}



\finexercice


\exercice{4509, quercia, 2010/03/14}
\enonce{004509}{\'Etude de convergence}

Soit $f_n(x) = \frac {nx}{1+n^2x^2}$.
\'Etudier la convergence simple, puis uniforme des $f_n$
sur $\R^+$ puis sur $[\alpha,+\infty[$, pour $\alpha > 0$.

\finenonce{004509}



\finexercice


\exercice{4510, quercia, 2010/03/14}
\enonce{004510}{$f(nx)$, $f(x/n)$}

Soit $f : {\R^+} \to \R$ continue, non identiquement nulle, telle que $f(0) = 0$
et $f(x) \to 0$ lorsque $x\to+\infty$.

On pose $f_n(x) = f(nx)$ et $g_n(x) = f\left(\frac xn\right)$.

\begin{enumerate}
  \item Donner un exemple de fonction $f$.
  \item Montrer que $f_n$ et $g_n$ convergent simplement vers la fonction nulle,
    et que la convergence n'est pas uniforme sur $\R^+$.
  \item Si $ \int_{t=0}^{+\infty} f(t)\,d t$ converge, chercher
       $\lim_{n\to\infty}  \int_{t=0}^{+\infty} f_n(t)\,d t$ et
       $\lim_{n\to\infty}  \int_{t=0}^{+\infty} g_n(t)\,d t$.
\end{enumerate}
\finenonce{004510}



\finexercice


\exercice{4511, quercia, 2010/03/14}
\enonce{004511}{\'Equation différentielle dépendant d'un paramètre}

Soit $y_n$ la solution de l'équation :
$(*_n) \Leftrightarrow \left(1+\frac1n\right)y'' - \left(2+\frac1n\right)y' + y = 0$
vérifiant les conditions initiales~: $y(0) = 0$, $y'(0) = 1$.

\begin{enumerate}
  \item Calculer explicitement $y_n$.
    
  \item Déterminer la limite simple, $y$, des fonctions $y_n$.
    
  \item Vérifier que $y$ est solution de l'équation limite de $(*_n)$ avec les mêmes
    conditions initiales.
\end{enumerate}
\finenonce{004511}



\finexercice


\exercice{4512, quercia, 2010/03/14}
\enonce{004512}{$f\circ f\circ ...\circ f$}

Soit $f : {[-1,1]} \to {[-1,1]}$ une fonction continue vérifiant :
$\forall\ x \ne 0,\quad |f(x)| < |x|$.

On pose $f_0(x) = x$, puis $f_{n+1}(x) = f(f_n(x))$.
\'Etudier la convergence simple des $f_n$.



\finenonce{004512}



\finexercice


\exercice{4513, quercia, 2010/03/14}
\enonce{004513}{\'Etude de convergence}

On pose $f_0(t) = 0$, $f_{n+1}(t) = \sqrt{t+f_n(t)}$, pour $t \ge 0$.

\begin{enumerate}
  \item Déterminer la limite simple, $\ell$, des fonctions $f_n$.
  \item Y a-t-il convergence uniforme sur $\R^+$ ?
  \item Démontrer que : $\forall\ t > 0,\
    |f_{n+1}(t) - \ell(t)| \le \frac {|f_n(t) - \ell(t)|}{2f_{n+1}(t)}$.
    
  \item En déduire que la suite $(f_n)$ converge uniformément sur tout intervalle
    $[a,+\infty[$, avec $a>0$.
    (Remarquer que $f_n - \ell$ est bornée pour $n \ge 1$)

\end{enumerate}
\finenonce{004513}



\finexercice


\exercice{4514, quercia, 2010/03/14}
\enonce{004514}{Approximation de la racine carrée par la méthode de Newton}

On définit une suite de fonctions ${f_n} : {\R^{+*}} \to {\R^{+*}}$
par récurrence :
$\begin{cases} f_0(x) = x \cr
         f_{n+1}(x) = \frac 12\left({f_n(x) + \frac x{f_n(x)}}\right).\cr \end{cases}$

\'Etudier la convergence simple, puis uniforme des $f_n$
$\left({\text{considérer }
 g_n(x)= \frac {f_n(x) - \sqrt x}{f_n(x) + \sqrt x}}\right)$.

\finenonce{004514}



\finexercice


\exercice{4515, quercia, 2010/03/14}
\enonce{004515}{Approximation polynomiale de la racine carrée}

On considère la suite $(f_n)$ de fonctions sur $[0,1]$ définie par les relations :
$f_0 = 0$, $f_{n+1}(t) = f_n(t) + \frac{t-f_n^2(t)}2$.

\'Etudier la convergence simple, uniforme, des fonctions $f_n$.
\finenonce{004515}



\finexercice


\exercice{4516, quercia, 2010/03/14}
\enonce{004516}{Suite ayant deux limites}

Trouver une suite de polynômes $(P_n)$ convergeant simplement (resp. uniformément) vers la
fonction nulle sur $[0,1]$ et vers la fonction constante égale à 1
sur $[2,3]$.

Remarque~: une telle suite a donc des limites distinctes dans $\R[x]$
pour les normes de la convergence uniforme sur $[0,1]$ et sur $[2,3]$.




\finenonce{004516}



\finexercice

\exercice{4522, quercia, 2010/03/14}
\enonce{004522}{Limite simple de polynômes de degrés bornés}

Soit $p \in \N$ fixé et $(P_n)$ une suite de fonctions polynomiales
de degrés inférieurs ou égaux à $p$ convergeant simplement vers $f$ sur un
intervalle $[a,b]$.

\begin{enumerate}
  \item Démontrer que $f$ est polynomiale de degré inférieur ou égal à $p$,
    et que les coefficients des $P_n$ convergent vers ceux de $f$.
    
  \item Montrer que la convergence est uniforme.

\end{enumerate}
\finenonce{004522}



\finexercice


\exercice{4523, quercia, 2010/03/14}
\enonce{004523}{Polynômes à coefficients entiers, ENS Lyon MP$^*$ 2005}

On consid{\'e}re $f : x\mapsto 2x(1-x)$ d{\'e}finie sur $[0,1]$. 

\begin{enumerate}
  \item \'Etude de la suite de fonction $g_{n}$, avec $g_{n}=f^{n}=f\circ\ldots\circ f$.


  \item Soit $[a,b]\subset {]0,1[}$ et $h$ continue sur $[a,b]$. Montrer que $h$ est limite uniforme 
sur $[a,b]$ d'une suite de polyn{\^o}mes {\`a} coefficients entiers. 

\end{enumerate}
\finenonce{004523}



\finexercice


\exercice{4524, quercia, 2010/03/14}
\enonce{004524}{Théorèmes de Dini}

Soit $(f_n)$ une suite de fonctions continues ${[a,b]}  \to {\R}$
convergeant simplement vers une fonction continue $f$.

\begin{enumerate}
  \item On suppose que chaque fonction $f_n$ est croissante. Montrer qu'il y a convergence uniforme.
  \item On suppose qu'à $x$ fixé la suite $(f_n(x))$ est croissante. Montrer qu'il y a convergence uniforme.
\end{enumerate}
\finenonce{004524}



\finexercice


\exercice{4525, quercia, 2010/03/14}
\enonce{004525}{Théorème d'Ascoli}

Soit $(f_n)$ une suite de fonctions ${[a,b]} \to {\R}$
convergeant simplement vers $f$.
On suppose que toutes les fonctions $f_n$ sont $k$-Lipchitizennes
(avec le même $k$).

\begin{enumerate}
  \item Soit $(a_0,a_1,\dots,a_N)$ une subdivision régulière de $[a,b]$.
    On note $M_n = \max\{|f_n(a_i)-f(a_i)|\text{ tel que } 0\le i \le N\}$.
    Encadrer $\|f_n-f\|_{\infty}$ à l'aide de $M_n$.

  \item Montrer que $f_n$ converge uniformément vers $f$.
\end{enumerate}
\finenonce{004525}



\finexercice

\exercice{5726, rouget, 2010/10/16}
\enonce{005726}{}
Etudier les suites de fonctions suivantes (convergence simple, convergence uniforme, convergence localement uniforme)

\begin{center}
\textbf{1) (**)} $f_n(x)=\frac{nx}{1+n^2x^2}$\quad\textbf{2) (**)} $f_n(x) = e^{-x}\sum_{k=0}^{n}\frac{x^k}{k!}$\quad\textbf{3) (**)} $f_n(x) = n(1-x)^n\sin\left(\frac{\pi x}{2}\right)$. 
\end{center}
\finenonce{005726}


\finexercice


\exercice{5727, rouget, 2010/10/16}
\enonce{005727}{*** I}
 Pour $n\in\Nn^*$, on pose $f_n(x)=\left\{\begin{array}{l}
\left(1-\frac{x}{n}\right)^n\;\text{si}\;x\in[0,n]\\
\rule{0mm}{5mm}0\;\text{si}\;x\geqslant n
\end{array}
\right.$.

\begin{enumerate}
 \item  Montrer que la suite $(f_n)_{n\in\Nn^*}$ converge uniformément sur $\Rr^+$ vers la fonction $f~:~x\mapsto e^{-x}$.

\item  A l'aide de la suite $(f_n)_{n\in\Nn^*}$, calculer l'intégrale de \textsc{Gauss} $\int_{0}^{+\infty}e^{-x^2}\;dx$.
\end{enumerate}
\finenonce{005727}


\finexercice


\exercice{5729, rouget, 2010/10/16}
\enonce{005729}{** I}
 Soit $(P_n)_{n\in\Nn}$ une suite de polynômes convergeant uniformément sur $\Rr$ vers une fonction $f$. Montrer que $f$ est un polynôme.

\finenonce{005729}


\finexercice


\exercice{5732, rouget, 2010/10/16}
\enonce{005732}{**}
Etudier (convergence simple, convergence absolue, convergence uniforme, convergence normale) les séries de fonctions de termes généraux :

\begin{enumerate}
\item  $f_n(x) = nx^2e^{-x\sqrt{n}}$ sur $\Rr^+$
\item  $f_n(x) =\frac{1}{n+n^3x^2}$   sur $\Rr_+^*$
\item  $f_n(x) = (-1)^n\frac{x}{(1+x^2)^n}$.
\end{enumerate}

\finenonce{005732}


\finexercice


\exercice{5733, rouget, 2010/10/16}
\enonce{005733}{** I}
Montrer que pour tout réel $a > 0$, $\int_{0}^{1}\frac{1}{1+x^a}\;dx=\sum_{n=0}^{+\infty}\frac{(-1)^n}{1+na}$.
\finenonce{005733}


\finexercice


\exercice{7542, mourougane, 2021/08/10}
\enonce{007542}{Le cas très particulier des polynômes}
Soit $d$ un entier naturel.
Montrer qu'une suite de polynômes $(P_n(z))_{n\in\Nn}$ de degré au plus $d$ converge localement uniformément sur $\Cc$
si et seulement si les $d+1$ suites de ses coefficients convergent,
si et seulement si, il existe $d+1$ points $c_i$ de $\Cc$ tels que les suites $(P_n(c_i))_{n\in\Nn}$ convergent.
\finenonce{007542}
\finexercice


\exercice{7543, mourougane, 2021/08/10}
\enonce{007543}{Un exemple}
\begin{enumerate}
\item Montrer que la série $\sum_{n\geq 1}\frac{(-1)^n}{z+n}$ converge de façon compacte sur $\Cc-\{-1,-2,-3\cdots\}$.
\item La convergence est-elle normale ?
\end{enumerate}
\finenonce{007543}
\finexercice

\section{222.02 Continuité, dérivabilité}


\exercice{4517, quercia, 2010/03/14}
\enonce{004517}{Fonction orthogonale aux polynômes}

Soit $f : {[a,b]} \to \R$ continue telle que pour tout entier~$k$ on
a $ \int_{t=a}^b f(t)t^k\,d t = 0$. Que peut-on dire de~$f$~?
\finenonce{004517}



\finexercice


\exercice{4518, quercia, 2010/03/14}
\enonce{004518}{Approximation de $f$ et $f'$}

Soit $f : {[a,b]} \to \R$ de classe $\mathcal{C}^1$.

\begin{enumerate}
  \item Montrer qu'il existe
une suite de polynômes $(P_n)$ telle que $P_n$ converge uniformément
vers $f$ et $P_n'$ converge uniformément vers $f'$.

  \item Si $f$ est $\mathcal{C}^\infty$, peut-on trouver une suite de polynômes
$(P_n)$ telle que pour tout $k$ la suite $(P_n^{(k)})$ converge uniformément
vers $f^{(k)}$~?

\end{enumerate}
\finenonce{004518}



\finexercice


\exercice{4519, quercia, 2010/03/14}
\enonce{004519}{Limite de $f_n(x_n)$}

Soient ${f_n} : D \to \R$ des fonctions continues convergeant vers une
fonction continue $f$ et $(x_n)$ une suite d'éléments de $D$ convergeant vers
$x \in D$.

\begin{enumerate}
  \item Si les fonctions $f_n$ convergent uniformément, montrer que
    $f_n(x_n) \to f(x)$ lorsque $n\to\infty$.
  \item Donner un contre-exemple lorsqu'il y a seulement convergence simple.

\end{enumerate}
\finenonce{004519}



\finexercice


\exercice{4520, quercia, 2010/03/14}
\enonce{004520}{Compositon et convergence}

Soit $f_n$ convergeant uniformément vers $f$, et $g$ une fonction continue.
Démontrer que $g \circ f_n \to g \circ f$ uniformément.

\finenonce{004520}



\finexercice


\exercice{4521, quercia, 2010/03/14}
\enonce{004521}{$f_n \circ g_n$}

Soit ${f_n} : {[a,b]} \to {[c,d]}$ et ${g_n} : {[c,d]} \to \R$
des fonctions continues convergeant uniformément vers les fonctions $f$ et $g$.
Montrer que $g_n\circ f_n$ converge uniformément vers $g\circ f$.

\finenonce{004521}



\finexercice


\exercice{4526, quercia, 2010/03/14}
\enonce{004526}{\'Equicontinuité}

Soit $(f_n)$ une suite de fonctions continues sur $D \subset \R$ convergeant
uniformément vers une fonction $f$.
Montrer que les fonctions $f_n$ sont {\it équi-continues\/} c'est à dire :
$$\forall\ x\in D,\ \forall\ \varepsilon > 0,\ \exists\ \delta > 0 \text{ tel que }
\forall\ n\in\N,\ \forall\ y \in {]x-\delta,x+\delta[}\cap D,\
|f_n(x)-f_n(y)| < \varepsilon.$$


\finenonce{004526}



\finexercice


\exercice{4527, quercia, 2010/03/14}
\enonce{004527}{Limite simple de fonctions convexes}

Soit ${f_n} : {[a,b]} \to \R$ des fonctions continues convexes convergeant simplement
vers une fonction continue $f$. Montrer que la convergence est uniforme.


\finenonce{004527}



\finexercice

\exercice{5728, rouget, 2010/10/16}
\enonce{005728}{*** I Polynômes de \textsc{Bernstein}. Théorème de \textsc{Weierstrass}}
Soit $f$ une application continue sur $[0,1]$ à valeurs dans $\Rr$. Pour $n$ entier naturel non nul, on définit le $n$-ème polynôme de \textsc{Bernstein} associé à $f$ par 

\begin{center}
$B_n(f) =\sum_{k=0}^{n}\dbinom{n}{k}f\left(\frac{k}{n}\right)X^k(1-X)^{n-k}$.
\end{center}

\begin{enumerate}
 \item
  \begin{enumerate}
   \item Calculer $B_n(f)$ quand $f$ est la fonction $x\mapsto 1$,  quand $f$ est la fonction $x\mapsto x$, quand $f$ est la fonction $x\mapsto x(x-1)$.
   \item En déduire que $\sum_{k=0}^{n}\dbinom{n}{k}(k-nX)^2X^k(1-X)^{n-k}= nX(1-X)$.
   \end{enumerate}

\item  En séparant les entiers $k$ tels que $\left|x-\frac{k}{n}\right|>\alpha$ et les entiers $k$ tels que $\left|x-\frac{k}{n}\right|\leqslant\alpha$ ($\alpha>0$ donné), montrer que la suite de polynômes $(B_n(f))_{n\in\Nn^*}$ converge uniformément vers $f$ sur $[0,1]$.

\item  Montrer le théorème de \textsc{Weierstrass} : soit $f$ une application continue sur $[a,b]$ à valeurs dans $\Rr$. Montrer que $f$ est limite uniforme sur $[a,b]$ d'une suite de polynômes.
\end{enumerate}
\finenonce{005728}


\finexercice


\exercice{5730, rouget, 2010/10/16}
\enonce{005730}{**}
Soit $f(x) =\sum_{n=1}^{+\infty}\frac{x^n\sin(nx)}{n}$.

\begin{enumerate}
 \item  Montrer que $f$ est de classe $C^1$ sur $]-1,1[$.

\item  Calculer $f'(x)$ et en déduire que $f(x) =\Arctan\left(\frac{x\sin x}{1-x\cos x}\right)$.
\end{enumerate}
\finenonce{005730}


\finexercice


\exercice{5731, rouget, 2010/10/16}
\enonce{005731}{**}
Soit $f(x) =\sum_{n=1}^{+\infty}\frac{(-1)^{n-1}}{\ln(nx)}$.

\begin{enumerate}
 \item  Domaine de définition de $f$. On étudie ensuite $f$ sur $]1,+\infty[$.

\item  Continuité de $f$ et limites de $f$ en $1$ et $+\infty$.

\item  Montrer que $f$ est de classe $C^1$ sur $]1,+\infty[$ et dresser son tableau de variation.
\end{enumerate}
\finenonce{005731}


\finexercice


\exercice{5734, rouget, 2010/10/16}
\enonce{005734}{**}
Pour $n\in\Nn^*$, soit $f_n(t) = (-1)^n\ln\left(1+\frac{t^2}{n(1+t^2)}\right)$.

\begin{enumerate}
 \item  Etudier la convergence simple et uniforme de la série de terme général $f_n$ puis la continuité de la somme $f$.

\item  Montrer que $\lim_{t \rightarrow +\infty}f(t) =\ln\left(\frac{2}{\pi}\right)$ à l'aide de la formule de \textsc{Stirling}.
\end{enumerate}
\finenonce{005734}


\finexercice


\exercice{5735, rouget, 2010/10/16}
\enonce{005735}{**}
Pour $n\in\Nn^*$ et $t\in\Rr$, soit $f_n(t) =\frac{\Arctan(nt)}{n^2}$.

  
Etude complète de $f =\sum_{n=1}^{+\infty}f_n$ : domaine de définition, parité, limites, continuité, dérivabilité (vérifier que $f$ n'est pas dérivable en $0$), allure du graphe.
\finenonce{005735}


\finexercice

\section{222.03 Suites et séries d'intégrales}


\exercice{5738, rouget, 2010/10/16}
\enonce{005738}{** I}
Pour $n\in\Nn^*$, on pose $f_n(x)=\left\{\begin{array}{l}
\left(1-\frac{x^2}{n}\right)^n\;\text{si}\;x\in[0,\sqrt{n}]\\
\rule{0mm}{5mm}0\;\text{si}\;x>\sqrt{n}
\end{array}
\right.$.

\begin{enumerate}
 \item  Montrer que la suite $(f_n)_{n\in\Nn^*}$ converge simplement sur $\Rr^+$ vers la fonction $f~:~x\mapsto e^{-x^2}$.

\item  A l'aide de la suite $(f_n)_{n\in\Nn^*}$, calculer l'intégrale de \textsc{Gauss} $\int_{0}^{+\infty}e^{-x^2}\;dx$.
\end{enumerate}
\finenonce{005738}


\finexercice


\exercice{5739, rouget, 2010/10/16}
\enonce{005739}{**}
Montrer que $\int_{0}^{1}x^{-x}\;dx =\sum_{n=1}^{+\infty}\frac{1}{n^n}$ et $\int_{0}^{1}x^{x}\;dx =\sum_{n=1}^{+\infty}\frac{(-1)^n}{n^n}$.
\finenonce{005739}


\finexercice


\exercice{5740, rouget, 2010/10/16}
\enonce{005740}{**}
\label{ex:rou3bis}
Montrer que $\int_{0}^{+\infty}\frac{x^2}{e^x-1}\;dx=2\sum_{n=1}^{+\infty}\frac{1}{n^3}$.
\finenonce{005740}


\finexercice


\exercice{5741, rouget, 2010/10/16}
\enonce{005741}{**}
Calculer $\int_{0}^{+\infty}\frac{x}{\sh x}\;dx$ en écrivant cette intégrale comme somme d'une série.
\finenonce{005741}


\finexercice


\exercice{5742, rouget, 2010/10/16}
\enonce{005742}{**}
Calculer $\int_{0}^{1}\frac{\ln x}{1+x^2}\;dx$ et $\int_{0}^{+\infty}\frac{\ln x}{1+x^2}\;dx$.
\finenonce{005742}


\finexercice


\exercice{5743, rouget, 2010/10/16}
\enonce{005743}{**}
\begin{enumerate}
 \item  Montrer que pour $x$ réel de $[0,1[$, $-\ln(1-x) =\sum_{n=1}^{+\infty}\frac{x^n}{n}$.

\item  Montrer que $\int_{0}^{1}\frac{\ln(t)\ln(1-t)}{t}\;dt =\sum_{n=1}^{+\infty}\frac{1}{n^3}$.
\end{enumerate}
\finenonce{005743}


\finexercice


\exercice{5744, rouget, 2010/10/16}
\enonce{005744}{*** I}
Montrer que pour tout réel $x$, $\int_{0}^{+\infty}\frac{\cos(xt)}{\ch t}\;dt = 2\sum_{n=0}^{+\infty}(-1)^n\frac{2n+1}{(2n+1)^2+x^2}$

\finenonce{005744}


\finexercice

\section{222.04 Suite et série de matrices}


\exercice{5864, rouget, 2010/10/16}
\enonce{005864}{**}
Déterminer $\lim_{n \rightarrow +\infty}\left(
\begin{array}{cc}
1&- \frac{a}{n}\\
 \frac{a}{n}&1
\end{array}
\right)^n$ ($a$ réel strictement positif donné).
\finenonce{005864}


\finexercice


\exercice{5865, rouget, 2010/10/16}
\enonce{005865}{***}
Soit $A\in\mathcal{M}_p(\Cc)$, $p\geqslant1$. Montrer que les trois propositions suivantes sont équivalentes :

\textbf{(1)} $\text{Sp}(A)\subset B_o(0,1)$ (disque unité ouvert).

\textbf{(2)} $\lim_{n \rightarrow +\infty}A^n=0$

\textbf{(3)} La série de terme général $A^n$, $n\in\Nn$, converge.
\finenonce{005865}


\finexercice


\exercice{5866, rouget, 2010/10/16}
\enonce{005866}{**}
Soit $A=\left(
\begin{array}{cc}
4/3&-5/6\\
5/3&-7/6
\end{array}
\right)$. Convergence et somme de la série de terme général $A^n$, $n\in\Nn$.
\finenonce{005866}


\finexercice


\exercice{5867, rouget, 2010/10/16}
\enonce{005867}{** I}
On munit $\mathcal{M}_p(\Cc)$ d'une norme sous-multiplicative notée $\|\;\|$. Soit $A$ un élément de $\mathcal{M}_p(\Rr)$ tel que $\|A\|<1$. Montrer que la série de terme général $A^n$, $n\in\Nn$, converge puis que $\sum_{n=0}^{+\infty}A^n=(I-A)^{-1}$.

En déduire que $\|(I-A)^{-1}-(I+A)\|\leqslant \frac{\|A\|^2}{1-\|A\|}$.
\finenonce{005867}


\finexercice


\exercice{5868, rouget, 2010/10/16}
\enonce{005868}{** I}
Soit $A\in\mathcal{M}_n(\Cc)$. Montrer qu'il existe $p_0\in\Nn$ tel que $\forall p\geqslant p_0$, $\sum_{k=0}^{p} \frac{A^k}{k!}\in GL_n(\Rr)$.
\finenonce{005868}


\finexercice


\exercice{5869, rouget, 2010/10/16}
\enonce{005869}{** I}
Calculer $\text{exp}(tA)$, $t\in\Rr$, si 
\begin{enumerate}
 \item  $A=\left(
\begin{array}{ccc}
3&2&2\\
1&0&1\\
-1&1&0
\end{array}
\right)$ \item  $A=\left(
\begin{array}{ccc}
4&1&1\\
6&4&2\\
-10&-4&-2
\end{array}
\right)$ 
\end{enumerate}
\finenonce{005869}


\finexercice


\exercice{5870, rouget, 2010/10/16}
\enonce{005870}{**}
Soit $A=\left(
\begin{array}{ccc}
0&1/2&-2\\
1/2&0&0\\
0&0&-1/2
\end{array}
\right)$. Calculer $\ln(I_3+tA)=\sum_{n=1}^{+\infty} \frac{(-1)^{n-1}t^n}{n}A^n$ en précisant les valeurs de $t$ pour lesquelles la série converge.

\finenonce{005870}


\finexercice


\exercice{5871, rouget, 2010/10/16}
\enonce{005871}{** I Exponentielle d'un endomorphisme anti-symétrique de $\Rr^3$}

\begin{enumerate}
 \item 
 \begin{enumerate}
 \item Soit $\overrightarrow{\omega}\in\Rr^3$. Pour $\overrightarrow{x}\in\Rr^3$, on pose $f_{\overrightarrow{\omega}}(\overrightarrow{x})=\overrightarrow{\omega}\wedge\overrightarrow{x}$. Vérifier que $f_{\overrightarrow{\omega}}\in\mathcal{A}(\Rr^3)$.

 \item Réciproquement, soit $f\in\mathcal{A}(\Rr^3)$. Montrer qu'il existe un vecteur $\overrightarrow{\omega}$ unique tel que $f=f_{\overrightarrow{\omega}}$.
  \end{enumerate}

\item  Soit $\overrightarrow{\omega}\in\Rr^3$. Montrer que $\text{exp}(f_{\overrightarrow{\omega}})$ est une rotation dont on déterminera l'axe (quand celui-ci est défini) et l'angle.
\end{enumerate}
\finenonce{005871}


\finexercice


\exercice{5872, rouget, 2010/10/16}
\enonce{005872}{**}
Pour $A\in\mathcal{M}_n(\Cc)$, calculer $\lim_{p \rightarrow +\infty}\left(I_n+ \frac{A}{p}\right)^p$.

\finenonce{005872}


\finexercice


\exercice{5873, rouget, 2010/10/16}
\enonce{005873}{**}
Montrer que $\forall A\in\mathcal{M}_n(\Rr)$, $\text{exp}(A)$ est un polynôme en $A$.
\finenonce{005873}


\finexercice

\section{222.99 Autre}


\exercice{4528, quercia, 2010/03/14}
\enonce{004528}{Fonction définie par une série}

On pose $f(x) = \sum_{n=0}^\infty \frac{\arccos(\cos nx)}{n!}$.

\begin{enumerate}
  \item Montrer que $f$ est définie sur $\R$, continue, paire et $2\pi$-périodique.
  \item Calculer $f(0)$, $f(\pi)$, $f\left(\frac\pi2\right)$.
    
\end{enumerate}
\finenonce{004528}



\finexercice


\exercice{4529, quercia, 2010/03/14}
\enonce{004529}{Fonction définie par une série (Centrale MP 2003)}
Soit~$f(a) = \sum_{n=0}^\infty e^{-a^2n^2}$ sous réserve de convergence $(a\in\R$).

\begin{enumerate}
  \item Domaine de définition de~$f$~?
    
  \item Limite de~$af(a)$ quand $a\to0$~?
    
  \item Limite de~$f(a)$ quand $a\to+\infty$~?
    

\end{enumerate}
\finenonce{004529}



\finexercice


\exercice{4530, quercia, 2010/03/14}
\enonce{004530}{Fonction $\zeta$ de Riemann}

Soit $\zeta(x) = \sum_{n=1}^\infty \frac1{n^x}$.

\begin{enumerate}
  \item Déterminer le domaine de définition de $\zeta$.
    Montrer que $\zeta$ est de classe $\mathcal{C}^\infty$ sur ce domaine.

  \item Prouver que $\zeta(x) \to 1$ lorsque $x\to+\infty$ $\Bigl($majorer
    $\sum_{n=2}^\infty \frac1{n^x}$ par comparaison à une intégrale$\Bigr)$.

  \item Prouver que $\zeta(x) \to +\infty$ lorsque $x\to1^+$.
\end{enumerate}
\finenonce{004530}



\finexercice


\exercice{4531, quercia, 2010/03/14}
\enonce{004531}{Fonction $\zeta$ de Riemann et constante d'Euler}

Soit $\zeta(x) = \sum_{n=1}^\infty \frac1{n^x}$ et
$\gamma = \lim_{n\to\infty}\left(\frac 11 + \dots + \frac 1n -\ln(n) \right)$.

Montrer que $\gamma = 1 + \sum_{n=2}^\infty\left(\frac 1n + \ln\left(1-\frac 1n\right)\right)$
puis que $\gamma = 1 - \sum_{k=2}^\infty\frac{\zeta(k)-1}k$.

\finenonce{004531}



\finexercice


\exercice{4532, quercia, 2010/03/14}
\enonce{004532}{Fonction définie par une série}

\begin{enumerate}
  \item \'Etudier la convergence simple, uniforme, de la série de fonctions :
    $f(x) = \sum_{n=0}^\infty ne^{-nx}$.

  \item Calculer $f(x)$ lorsque la série converge (intégrer terme à terme).
    

\end{enumerate}
\finenonce{004532}



\finexercice


\exercice{4533, quercia, 2010/03/14}
\enonce{004533}{Fonction définie par une série}

\begin{enumerate}
  \item \'Etudier la convergence de la série
    $f(x) = \sum_{n=0}^\infty \frac1{1+x^n}$.
  \item Montrer que $f$ est de classe $\mathcal{C}^1$ sur son domaine de définition.
  \item Tracer la courbe représentative de $f$ sur $]1,+\infty[$.

\end{enumerate}
\finenonce{004533}



\finexercice


\exercice{4534, quercia, 2010/03/14}
\enonce{004534}{Fonction définie par une série}

Soit $g(x) = \sum_{n=0}^\infty \frac{(-1)^n}{n!\,(x+n)}$.

\begin{enumerate}
  \item Déterminer le domaine, $D$ de définition de $g$ et prouver que $g$ est de
    classe $\mathcal{C}^\infty$ sur $D$.
  \item Montrer que la quantité : $xg(x) - g(x+1)$ est constante sur $D$.
    
  \item Tracer la courbe représentative de $g$ sur $]0,+\infty[$.
    
  \item Donner un équivalent de $g(x)$ en $+\infty$ et en $0^+$.
    

\end{enumerate}
\finenonce{004534}



\finexercice


\exercice{4535, quercia, 2010/03/14}
\enonce{004535}{Fonction définie par une série}

\begin{enumerate}
  \item\'Etablir la convergence simple sur $\R$ de la série de fonctions :
     $f(x) = \sum_{n=0}^\infty \frac {(\sin x)^2}{\ch nx}$.

  \item  Montrer que la convergence est uniforme sur toute partie de la forme
     $\R\setminus[-\alpha,\alpha]$, $\alpha > 0$.
     Que pouvez-vous en déduire pour $f$ ?
\end{enumerate}
\finenonce{004535}



\finexercice


\exercice{4536, quercia, 2010/03/14}
\enonce{004536}{Fonction définie par une série}

Soit $u_n(x) = (-1)^n\ln\left(1+\frac x{n(1+x)}\right)$ et
$f(x) = \sum_{n=1}^\infty u_n(x)$.

\begin{enumerate}
  \item Montrer que la série $f(x)$ converge simplement sur $\R^+$.
  \item Majorer convenablement le reste de la série, et montrer qu'il y a
    convergence uniforme sur $\R^+$.
    
  \item Y a-t-il convergence normale ?
    
\end{enumerate}
\finenonce{004536}



\finexercice


\exercice{4537, quercia, 2010/03/14}
\enonce{004537}{Fonction définie par une série}

Soit $f(x) = \sum_{n=0}^\infty \frac1{x(x+1)\dots(x+n)}$.

\begin{enumerate}
  \item \'Etablir l'existence et la continuité de $f$ sur $\R^{+*}$.
  \item Calculer $f(x+1)$ en fonction de $f(x)$.
    
  \item Tracer la courbe de $f$.

\end{enumerate}
\finenonce{004537}



\finexercice


\exercice{4538, quercia, 2010/03/14}
\enonce{004538}{Fonction définie par une série}

\begin{enumerate}
  \item \'Etudier la convergence simple, uniforme, de $f(x) = \sum_{n=0}^\infty
    \bigl(\Arctan(x+n) - \Arctan(n)\bigr)$.
    
  \item Montrer que $f$ est de classe $\mathcal{C}^1$ sur $\R$.
  \item Chercher une relation simple entre $f(x)$ et $f(x+1)$.
    
  \item Trouver $\lim_{x\to+\infty} f(x)$.
    
\end{enumerate}
\finenonce{004538}



\finexercice


\exercice{4539, quercia, 2010/03/14}
\enonce{004539}{Conversion série-intégrale}

Montrer, pour $x > 0$ : $\sum_{n=0}^\infty \frac{(-1)^n}{n+x}
=  \int_{t=0}^1 \frac{t^{x-1}}{t+1}\,d t$.


\finenonce{004539}



\finexercice


\exercice{4540, quercia, 2010/03/14}
\enonce{004540}{Fonction $\Gamma$}

Soit $f_n(x) = \frac{n^x}{(1+x)(1+x/2)\dots(1+x/n)}$.
\begin{enumerate}
  \item \'Etudier la convergence simple des fonctions $f_n$.

  \item On note $f = \lim f_n$. Calculer $f(x)$ en fonction de $f(x-1)$ lorsque
    ces deux quantités existent.
  \item Montrer que $f$ est de classe $\mathcal{C}^1$ sur son domaine de définition
    (on calculera $f_n'(x)/f_n(x)$).

\end{enumerate}
\finenonce{004540}



\finexercice


\exercice{4541, quercia, 2010/03/14}
\enonce{004541}{Ensi Chimie P' 93}

\'Etudier la convergence de la suite de fonctions définies par :
$f_n(x) = \frac{n(n+1)}{x^{n+1}} \int_0^x (x-t)^{n-1}\sin t\,d t$.

\finenonce{004541}



\finexercice


\exercice{4542, quercia, 2010/03/14}
\enonce{004542}{Convergence de $f^{(n)}$}

Soit $f\in \mathcal{C}^{\infty}(\R)$. On définit la suite $(f_n)_{n\in\N^*}$ par
$f_n = {f}^{(n)}$ (dérivée $n$-ème).
On suppose que ${(f_n)}_{n\ge 1}$ converge uniformément vers $\varphi$.
Que peut-on dire de $\varphi$ ?
\finenonce{004542}



\finexercice


\exercice{4543, quercia, 2010/03/14}
\enonce{004543}{Ensi PC 1999}

Soit $f_n(x) = \frac{(-1)^n\cos^nx}{n+1}$.

\begin{enumerate}
  \item \'Etudier la convergence de $f(x) = \sum_{n=0}^\infty f_n(x)$.

  \item Montrer la convergence de la série de terme général
    $u_n =  \int_{x=0}^{\pi/2} f_n(x)\,d x$.

  \item En déduire $\sum_{n=0}^\infty u_n$ sous forme d'une intégrale.

\end{enumerate}
\finenonce{004543}



\finexercice


\exercice{4544, quercia, 2010/03/14}
\enonce{004544}{Développement de $\coth(x)$}

\begin{enumerate}
  \item Décomposer en éléments simples sur $\C$ la fractions rationnelle~:
    $F_n(X) = \frac1{(1+X/n)^n-1}$.

  \item En déduire pour $x\in\R^*$~:
    $\coth x = \frac1{e^{2x}-1} - \frac1{e^{-2x}-1} = \frac1x + \sum_{k=1}^\infty\frac{2x}{x^2+k^2\pi^2}$.

  \item En déduire la valeur de $\zeta(2)$.
    
\end{enumerate}
\finenonce{004544}



\finexercice


\exercice{4545, quercia, 2010/03/14}
\enonce{004545}{$\sum\sin(n)/n$}

Pour $n\in\N^*$ et $x\in[-1,1]$ on pose $u_n(x)=\frac{x^n\sin(nx)}n$.

\begin{enumerate}
  \item Montrer que la série $\sum_{n=1}^\infty u_n(x)$ converge uniformément
    sur $[-1,1]$ vers une fonction continue, $f$.
    
  \item Justifier la dérivabilité de~$f$ sur $]-1,1[$ et calculer $f'(x)$.
    En déduire $f(x)$.
    
  \item En déduire la valeur de $\sum_{n=1}^\infty \frac{\sin n}n$.
    
\end{enumerate}
\finenonce{004545}



\finexercice


\exercice{4546, quercia, 2010/03/14}
\enonce{004546}{Fonctions $\zeta$ et $\eta$}

Pour $x>1$ on pose $\zeta(x) = \sum_{n=1}^\infty \frac1{\strut n^x}$
et pour $x>0$~: $\eta(x) = \sum_{n=1}^\infty \frac{(-1)^{n-1}}{\strut n^x}$.

\begin{enumerate}
  \item \'Etablir pour $x>1$~: $\eta(x) = (1-2^{1-x})\zeta(x)$.
    En déduire $\zeta(x) \sim \frac 1{x-1}$ pour $x\to1^+$.

  \item Montrer que $\zeta(x) = \frac 1{x-1} + \gamma +  o(1)$.
    On remarquera que $\frac 1{x-1} =  \int_{t=1}^{+\infty} \frac{d t}{\strut t^x}$.
    

  \item En déduire la valeur de $\sum_{n=1}^\infty \frac{(-1)^n\ln n}n$.
    
\end{enumerate}
\finenonce{004546}



\finexercice


\exercice{4547, quercia, 2010/03/14}
\enonce{004547}{Centrale MP 2000}

Pour~$y\in\R$ et~$n\in\N^*$, on pose $a_n(y) = \frac{\cos(ny)}{\sqrt n}$.

\begin{enumerate}
  \item Déterminer le rayon de convergence de la série entière $\sum a_n(y)x^n$.
  \item Soit~$D=\{(x,y)\in\R^2,\ |x|<1\}$ et $F(x,y) = \sum_{n=1}^{+\infty} a_n(y)x^n$.
    Montrer que $F$, $\frac{\partial F}{\partial x}$ et $\frac{\partial F}{\partial y}$ existent en tout point de~$D$.
\end{enumerate}
\finenonce{004547}



\finexercice


\exercice{4548, quercia, 2010/03/14}
\enonce{004548}{Série lacunaire}
Soit $(p_n)$ une suite d'entiers naturels, strictement croissante
et telle que $p_n/n\to\infty$  lorsque $n\to\infty$. On pose pour $x\in{]-1,1[}$~:
$f(x)=\sum_{n=0}^\infty x^{p_n}$. Montrer que $(1-x)f(x)\to 0$ lorsque $x\to1^-$.

\finenonce{004548}



\finexercice


\exercice{4549, quercia, 2010/03/14}
\enonce{004549}{Fonctions réciproques (Pugin, MP$^*$-2001)}

Soit $(f_n)$ une suite de fonctions $[a,b]\to[c,d]$ continues, bijectives, strictement
croissantes, convergeant simplement vers une fonction $f$ : $[a,b]\to[c,d]$ elle aussi
continue, bijective strictement croissante.

\begin{enumerate}
  \item Montrer qu'il y a convergence uniforme (deuxième théorème de Dini, considérer une subdivision de~$[a,b]$).
  \item Montrer que les fonctions réciproques $f_n^{-1}$ convergent simplement
    vers une fonction~$g$ et que $g = f^{-1}$.

  \item Montrer que $(f_n^{-1})$ converge uniformément vers~$f^{-1}$.
\end{enumerate}
\finenonce{004549}



\finexercice


\exercice{4550, quercia, 2010/03/14}
\enonce{004550}{Mines MP 2001}

Soit $(f_n)$ une suite de fonctions continues sur le compact $K$, à
valeurs réelles et convergent uniformément sur $K$ vers
la fonction $f$. A-t-on $\sup f_n\to \sup f$  lorsque $n\to\infty$ ?


\finenonce{004550}



\finexercice


\exercice{4551, quercia, 2010/03/14}
\enonce{004551}{Mines MP 2001}

Pour $x\in\R^+$ et $n\in\N$, $n\ge 2$ on pose $f_n(x) = \frac{xe^{-nx}}{\ln n}$
et $S(x) = \sum_{n=2}^\infty f_n(x)$ sous réserve de convergence.

\begin{enumerate}
  \item \'Etudier la convergence simple, normale, uniforme de la série $\sum f_n$
    sur $\R^+$.
    

  \item Montrer que $S$ est de classe $\mathcal{C}^1$ sur $\R^{+*}$.
  \item Montrer que $S$ n'est pas dérivable à droite en $0$.
    
  \item Montrer que $x^kS(x)$ tend vers $0$ en $+\infty$ pour tout $k\in\N$.


\end{enumerate}
\finenonce{004551}



\finexercice


\exercice{4552, quercia, 2010/03/14}
\enonce{004552}{Centrale MP 2001}

Convergence et limite en $1^-$ de $f(x) = \sum_{n=0}^\infty \frac{(1-x)x^n}{1+x^n}$.

\finenonce{004552}



\finexercice


\exercice{4553, quercia, 2010/03/14}
\enonce{004553}{Centrale MP 2001}

Soit $S(t) = \sum_{n=1}^\infty \frac{t^n}{1-t^n}$.

\begin{enumerate}
  \item Pour quelles valeurs de $t$, $S$ est-elle définie~? Est-elle continue~?

  \item Montrer qu'au voisinage de~$1^-$ on a $S(t) = -\frac{\ln(1-t)}{1-t} + O\Bigl(\frac1{1-t}\Bigr)$.
On pourra développer $\ln(1-t)$ en série entière.
\end{enumerate}
\finenonce{004553}



\finexercice


\exercice{4554, quercia, 2010/03/14}
\enonce{004554}{Centrale MP 2002}

On pose $\phi (x)= d(x,\Z)=\inf\{ |x-n| \text{ tel que } n\in\Z\}$.
\begin{enumerate}
  \item Montrer que $f$ : $\R\ni x \mapsto \sum_{n=0}^{+\infty} (\frac34) ^n\phi (4^nx)$
    est définie et continue.

  \item Montrer que $\phi$ est lipschitzienne. Que peut-on en déduire pour $f$~?

  \item Montrer que $f$ n'est dérivable en aucun point.

\end{enumerate}
\finenonce{004554}



\finexercice


\exercice{4555, quercia, 2010/03/14}
\enonce{004555}{ENS Lyon-Cachan MP 2002}

Soin $(a_n)_{n\ge 1}$ une suite complexe telle que la série $\sum a_n$ converge. On pose~:
$f(h) = \sum_{n=1}^\infty a_n\frac{\sin^2(nh)}{(nh)^2}$ si $h\ne 0$ et
$f(0) = \sum_{n=1}^\infty a_n$. Étudier le domaine de définition et la
continuité de~$f$.

\finenonce{004555}



\finexercice


\exercice{4556, quercia, 2010/03/14}
\enonce{004556}{Centrale MP 2002}

Soit $f : \R \to \R$ continue et $2\pi$-périodique. Pour $n\in\N^*$, on pose
$F_n(x)=\frac{1}{n} \int_{t=0}^nf(x+t)f(t)\,d t$.

\begin{enumerate}
  \item Montrer que la suite $(F_n)$ converge vers une fonction $F$ que l'on précisera.

  \item Nature de la convergence~?

  \item Prouver $\|F\|_{\infty}= |F(0)|$.

\end{enumerate}
\finenonce{004556}



\finexercice


\exercice{4557, quercia, 2010/03/14}
\enonce{004557}{Approximation par des fractions rationnelles}

Soit $f:\R \to \R$ continue, ayant même limite finie~$\ell$ en $\pm\infty$.
Montrer que $f$ est limite uniforme sur~$\R$ de fractions rationnelles.
\finenonce{004557}



\finexercice


\exercice{4558, quercia, 2010/03/14}
\enonce{004558}{Fonction définie par une série}

On pose pour $x\in\R$~: $f(x) = \sum_{n=1}^\infty \frac{(-1)^{n-1}}{\sqrt{n^2+x^2}}$.

\begin{enumerate}
  \item Déterminer $\lim_{x\to\infty}f(x)$.

  \item Chercher un équivalent de~$f(x)$ en $+\infty$.
\end{enumerate}
\finenonce{004558}



\finexercice


\exercice{4559, quercia, 2010/03/14}
\enonce{004559}{Recherche d'équivalents, Centrale MP 2006}
Déterminer un équivalent au voisinage de $0$ de
$S_{1}(x) = \sum_{n=1}^{\infty}{\frac{1}{\sh{(nx)}}}$ et
$S_{2}(x) = \sum_{n=1}^{\infty}{\frac{1}{\sh^{2}{(nx)}}}$.



\finenonce{004559}



\finexercice


\exercice{4560, quercia, 2010/03/14}
\enonce{004560}{Étude de $\sum t^{p-1}\sin(px)$ pour~$x\in{}0,\pi[}$, TPE MP 2005]

\begin{enumerate}
  \item Calculer $S_n(t) = \sum_{p=1}^n t^{p-1}\sin(px)$ puis $S(t) = \lim_{n\to\infty} S_n(t)$.
    

  \item Calculer $ \int_{t=0}^1 S_n(t)\,d t$ et  $ \int_{t=0}^1 S(t)\,d t$.
    
  \item En déduire que $\sum_{n=1}^\infty \frac{\sin nx}n$ converge et donner sa valeur.
    

\end{enumerate}
\finenonce{004560}



\finexercice


\exercice{4561, quercia, 2010/03/14}
\enonce{004561}{Fraction rationnelle de meilleure approximation (Ens Ulm-Lyon-Cachan MP$^*$ 2003)}
On note~$R$ l'ensemble des fractions rationnelles continues sur~$[0,1]$
et pour $m,n\in\N$~:

$R_{m,n} = \{f\in R\text{ tel que }\exists\ P,Q\in\R[X]\text{ tel que }\deg(P)\le m,\ \deg(Q)\le n\text{ et }f=P/Q\}$.

\begin{enumerate}
  \item $R$ est-il un espace vectoriel~? Si oui en trouver une base. Même question pour~$R_{m,n}$.
    
  \item Soient $m,n$ fixés. On note $d = \inf\{\|g-f\|,\ f\in R_{m,n}\}$
    où $g$ désigne une fonction continue de~$[0,1]$ dans~$\R$ et ${\|h\| = \sup\{|h(x)|,\ x\in[0,1]\}}$.
    Montrer qu'il existe $r_0\in R_{m,n}$ tel que $\|g-r_0\| = d$.
    
\end{enumerate}
\finenonce{004561}



\finexercice


\exercice{4562, quercia, 2010/03/14}
\enonce{004562}{Dérivation multiple, ULM-Lyon-Cachan MP$^*$ 2005}
\begin{enumerate}
  \item Soit $(f_n)$ une suite de fonctions de classe~$\mathcal{C}^1$ sur $[a,b]$ telle que $(f_n')$ converge
    uniformément vers $g$ et il existe $x_1$ tel que $(f_n(x_1))$ converge.
    Montrer que $(f_n)$ converge uniformément sur~$[a,b]$ vers $f$ telle que $f'=g$.

  \item Soit $(f_n)$ une suite de fonctions de classe~$\mathcal{C}^p$ sur $[a,b]$ telle que $(f_n^{(p)})$ converge
    uniformément vers $g$ et il existe $x_1,\dots,x_p$ distincts tels que $(f_n(x_i))$ converge.
    Montrer que $(f_n)$ converge uniformément sur~$[a,b]$ vers $f$ telle que $f^{(p)}=g$.
\end{enumerate}
\finenonce{004562}



\finexercice


\exercice{4563, quercia, 2010/03/14}
\enonce{004563}{Exponentielle, Polytechnique MP$^*$ 2006}
Soient $A,B\in \mathcal{M}_n(\R)$. Montrer que :
$\exp(A) - \exp(B) =  \int_{s=0}^1\exp(sA) (A-B) \exp((1-s)B)\,d s$.

\finenonce{004563}



\finexercice


\exercice{5736, rouget, 2010/10/16}
\enonce{005736}{**}
Pour $x>0$, on pose $f(x) =\sum_{n=0}^{+\infty}e^{-x\sqrt{n}}$. Trouver un équivalent simple de $f$ en $0$ à droite.
\finenonce{005736}


\finexercice


\exercice{5737, rouget, 2010/10/16}
\enonce{005737}{***}
Pour $x\in]-1,1[$, on pose $f(x)=\sum_{n=1}^{+\infty}x^{n^2}$. Trouver un équivalent simple de $f$ en $1$.
\finenonce{005737}


\finexercice

\section{223.01 Limite}


\exercice{1784, drutu, 2003/10/01}

\enonce{001784}{}
\'Etudier l'existence des limites suivantes :
\begin{enumerate}
\item $\mathrm{lim}_{\begin{smallmatrix}(x,y)\to (0,0)\\ x+y\ne 0\end{smallmatrix}} \frac{x^2y}{x+y}$ 
\item  $\mathrm{lim}_{\begin{smallmatrix}(x,y,z)\to (0,0,0)\\2x^3+yz^2 \ne 0\end{smallmatrix}} \frac{xyz+z^3}{2x^3+yz^2}$
\item $\mathrm{lim}_{\begin{smallmatrix}(x,y)\to (0,0)\\(x,y)\ne (0,0)\end{smallmatrix}} \frac{|x|+|y|}{x^2+y^2}$
\item $\mathrm{lim}_{\begin{smallmatrix}(x,y)\to (0,0)\\ x \ne \pm y\end{smallmatrix}} \frac{x^4y}{x^2-y^2}$
\item  $\mathrm{lim}_{\begin{smallmatrix}(x,y,z)\to (0,0,0)\\(x,y,z) \ne (0,0,0)\end{smallmatrix}} \frac{xy+yz}{x^2+2y^2+3z^2}$
\end{enumerate}

\finenonce{001784}



\finexercice



\exercice{1785, drutu, 2003/10/01}

\enonce{001785}{}
Soit $f \colon \R^2 \setminus \{(0,0)\} \to \R$ la fonction d\'efinie par
\[
f(x,y) = \frac{x^2y^2}{x^2y^2 + (x-y)^2} .
\]
Montrer que
\begin{equation}
\lim_{x\to 0}\lim_{y\to 0}f(x,y) =  \lim_{y\to 0}\lim_{x\to 0}f(x,y) = 0
\label{non}
\end{equation}
et que $\lim_{(x,y)\to (0,0)} f(x,y)$ n'existe pas.
\finenonce{001785}



\finexercice



\exercice{1786, drutu, 2003/10/01}

\enonce{001786}{}
Soit 
$$f : \R^2 \to \R , \quad f(x,y) = \left\{
         \begin{array}{cc}
         (x+y)\sin \frac{1}{x} \sin \frac{1}{y} & \mbox{ si }xy \neq 0 \\
         0 & \mbox{ si } xy = 0 
         \end{array}
         \right .$$
D\'emontrer que  les deux limites it\'er\'ees 
$$\lim_{x\to 0}\lim_{y\to 0}f(x,y) \quad  \mbox{ et }  \lim_{y\to 0}\lim_{x\to 0}f(x,y)$$ 
n'existent pas, et que
$$\lim_{(x,y)\to (0,0)}f(x,y)$$ 
existe et est \'egale \`a $0$.
\finenonce{001786}



\finexercice



\exercice{1787, drutu, 2003/10/01}

\enonce{001787}{}
D\'eterminer les limites
lorsqu'elles existent:
\begin{enumerate}
\item $\lim_{(x,y)\to (0,0)} \frac{x}{x^2+y^2}  $ 
\item$ \lim_{(x,y)\to (0,0)} \frac{(x+2y)^3}{x^2+y^2} $ 
\item$ \lim_{(x,y)\to (1,0)} \frac{\log (x+e^y)}{\sqrt{x^2+y^2}} $ 
\item $\lim_{(x,y)\to (0,0)} \frac{x^4+y^3-xy}{x^4+y^2} $ 
\item $ \lim_{(x,y)\to (0,0)} \frac{x^3y}{x^4+y^4}  $ ;
\item $ \lim_{(x,y)\to (0,0)} \frac{(x^2+y^2)^2}{x^2-y^2} $ ;
\item $ \lim_{(x,y)\to (0,0)} \frac{1-\cos xy}{y^2}  $ ;
\item $ \lim_{(x,y)\to (0,0)} \frac{\sin x}{\cos y-\cosh x} $ 
\end{enumerate}
\finenonce{001787}



\finexercice



\exercice{1788, drutu, 2003/10/01}

\enonce{001788}{}
Pour chacune des fonctions $f$ suivantes, \'etudier l'existence d'une
limite en $(0,0,0)$  :
\begin{enumerate}
\item $f(x,y,z)= \frac{xyz}{x+y+z}$ ;
\item $f(x,y,z)= \frac{x+y}{x^2-y^2+z^2}$.
\end{enumerate}
\finenonce{001788}



\finexercice



\exercice{1789, maillot, 2001/09/01}

\enonce{001789}{}
\'Etudier la continuité des fonctions définies sur $\R^2$ par
\begin{align*}
&f_1(x,y)=\frac{xy}{x^2+y^2} \qquad &\text{si }\quad (x,y)\neq (0,0),\\
&f_1(0,0)=0.\\
&f_2(x,y)=\frac{x^3+y^3}{x^2+y^2} \qquad &\text{si}\quad (x,y)\neq (0,0),\\
&f_2(0,0)=0.
\end{align*}
\finenonce{001789}



\finexercice



\exercice{1790, maillot, 2001/09/01}

\enonce{001790}{partiel 1999}
\begin{enumerate}
\item \'Etudier la continuit\'e de la fonction $f_1:\R^2\rightarrow\R$ d\'efinie par
\[f_1(x,y)=\begin{cases} \frac{(\sin x)\, (\sin y)}{\sqrt{\vert x\vert}+
       \sqrt{\vert y\vert}} &\text{si\ } (x,y)\neq (0,0)\\
               0 &\text{si\ } (x,y)= (0,0).
         \end{cases}\]
\item Soit $a>0$ fix\'e. \'Etudier la continuit\'e de la fonction
$f_2:\R^2\rightarrow\R$ d\'efinie par
\[f_2(x,y)=\begin{cases} \frac{{\vert x\vert}^a\, {\vert y\vert}^a}{x^2+
       y^2} &\text{si\ } (x,y)\neq (0,0)\\
               0 &\text{si\ } (x,y)= (0,0).
         \end{cases}\]
\item \'Etudier la continuit\'e de la fonction $f_3:\R^2\rightarrow\R$ d\'efinie par
\[f_3(x,y)=\begin{cases} y-x^2 &\text{si\ } y>x^2\\
               0 &\text{si\ } y\le x^2.
         \end{cases}\]
\item On d\'efinit une fonction continue de l'ouvert $U=\{(x,y,z)\in
\R^3 \mid xyz\neq 0\}$ dans $\R$ en posant
$$ f_4(x,y,z)=(x^2+y^2+z^2)\, \sin{\frac1x}\, \sin{\frac1y}\,
\cos{\frac1z}.$$
\'Etudier la possibilit\'e de prolonger $f_4$ en une fonction continue
sur $\R^3$.
\end{enumerate}
\finenonce{001790}



\finexercice



\exercice{1791, ridde, 1999/11/01}

\enonce{001791}{}
Prolonger par continuit\'e la fonction ${g}:{(\Rr^{2})^{*}}\to{\Rr}$, ${ (x, y)} \mapsto
{xy\ln (x^{2} + y^{2})}$.
\finenonce{001791}



\finexercice



\exercice{1792, gourio, 2001/09/01}

\enonce{001792}{}
Soit $ f:{\Rr}^{2}\rightarrow {\Rr}$ telle que $ \forall (x,y)\in {\Rr}
^{2},f(x,.)$ et $ f(.,y) $ sont continues. Montrer qu'il existe une suite $
\left( g_{n}\right) _{n\in \Nn } $ d'applications continues sur $ {\Rr}
^{2}$ telles que :
$$\forall (x,y)\in {\Rr}^{2},\lim\limits_{n\rightarrow \infty
}g_{n}(x,y)=f(x,y). $$
\finenonce{001792}



\finexercice



\exercice{2621, debievre, 2009/05/19}
\enonce{002621}{}
Pour chacune des suites $(u_n)_n$ dans le plan $\R^2$ ci-dessous, 
placer quelques-uns des points $u_n$ dans le plan 
et d\'ecrire qualitativement le comportement de la suite lorsque $n$ tend 
vers l'infini. \'Etudier ensuite 
la convergence de chacune des suites et d\'eterminer 
la limite le cas \'ech\'eant.

\begin{enumerate}
  \item $u_n=(\frac{4n^2}{n^2+4n+3}, \cos\frac1n)$
  \item $u_n=(\frac{n^2\arctan n}{n^2+1}, \sin(\frac{\pi}{4}\exp(-{\frac1n})))$
  \item $u_n=(\sinh n, \frac{\ln n}{n})$
  \item $u_n=(a^n\cos (n\alpha), a^n\sin (n\alpha))$, en fonction de 
$a\in \R$, $a>0$ et $\alpha\in\R$.
\end{enumerate}\finenonce{002621}


\finexercice


\exercice{2649, debievre, 2009/05/19}
\enonce{002649}{}
On consid\`ere une suite $(u_n)_n$, de terme g\'en\'eral $u_n\in\R^2$.
\begin{enumerate}
 \item  Donner la d\'efinition de convergence pour une telle suite. (Ceci est une question de cours!)

 \item Soit la suite de terme g\'en\'eral $u_n=(\tanh(n), \cos(n) \exp(-n^2))$. \'Etudier sa convergence.
\end{enumerate}
\finenonce{002649}



\finexercice


\exercice{5553, rouget, 2010/07/15}
\enonce{005553}{**T}
Etudier l'existence et la valeur éventuelle d'une limite en $(0,0)$ des fonctions suivantes :
\begin{enumerate}
 \item  $\frac{xy}{x+y}$
 \item  $\frac{xy}{x^2+y^2}$
 \item  $\frac{x^2y^2}{x^2+y^2}$
 \item   $\frac{1+x^2+y^2}{y}\sin y$
 \item  $\frac{x^3+y^3}{x^2+y^2}$
 \item  $\frac{x^4+y^4}{x^2+y^2}$.
\end{enumerate}

\finenonce{005553}


\finexercice


\exercice{5887, rouget, 2010/10/16}
\enonce{005887}{** I}
Etudier l'existence et la valeur éventuelle des limites suivantes :

\begin{enumerate}
 \item  $ \frac{xy}{x^2+y^2}$ en $(0,0)$
 \item  $ \frac{x^2y^2}{x^2+y^2}$ en $(0,0)$
 \item  $ \frac{x^3+y^3}{x^2+y^4}$  en $(0,0)$
 \item  $ \frac{\sqrt{x^2+y^2}}{|x|\sqrt{|y|}+|y|\sqrt{|x|}}$  en $(0,0)$
 \item  $ \frac{(x^2-y)(y^2-x)}{x+y}$ en $(0,0)$
 \item  $ \frac{1-\cos\sqrt{|xy|}}{|y|}$ en $(0,0)$
 \item  $ \frac{x+y}{x^2-y^2+z^2}$ en $(0,0,0)$
 \item  $ \frac{x+y}{x^2-y^2+z^2}$ en $(2,-2,0)$
\end{enumerate}
\finenonce{005887}


\finexercice

\section{223.02 Continuité}


\exercice{1793, gourio, 2001/09/01}

\enonce{001793}{}
Trouver les fonctions $ f$ continues sur $ {\Rr}^{2}$ telles que :
$$\forall (x,y)\in {\Rr}^{2},f(x,y)=f(x+y,x-y). $$
\finenonce{001793}



\finexercice



\exercice{1794, drutu, 2003/10/01}

\enonce{001794}{}
Etudier la continuit\'e sur $\R^2$ de la fonction suivante :

\begin{enumerate}

\item $$f(x,y) = \left\{
         \begin{array}{cc}
         \frac{x^2y^2}{x^2+y^2} & \mbox{ si }(x,y) \neq (0,0) \\
         0 & \mbox{ sinon. }
         \end{array}
         \right .$$
         

\item $$f(x,y) = \left\{
         \begin{array}{cc}
         \frac{x^2y}{x^2+y^2} & \mbox{ si }(x,y) \neq (0,0) \\
         0 & \mbox{ sinon. }
         \end{array}
         \right .$$
         

\item $$f(x,y) = \left\{
         \begin{array}{cc}
         \frac{x^4y}{x^4+y^6} & \mbox{ si }(x,y) \neq (0,0) \\
         0 & \mbox{ sinon. }
         \end{array}
         \right .$$

\item $$f(x,y) = \left\{
         \begin{array}{cc}
         \frac{xy^4}{x^4+y^6} & \mbox{ si }(x,y) \neq (0,0) \\
         0 & \mbox{ sinon. }
         \end{array}
         \right .$$
         
\item $$f(x,y) = \left\{
         \begin{array}{cc}
         y^2\sin \frac{x}{y} & \mbox{ si }y \neq 0 \\
         0 & \mbox{ sinon. }
         \end{array}
         \right .$$
         
\item $$f(x,y) = \left\{
         \begin{array}{cc}
         xe^{\arctan \frac{y}{x}} & \mbox{ si }x \neq 0 \\
         0 & \mbox{ sinon. }
         \end{array}
         \right .$$
         
\end{enumerate}
\finenonce{001794}



\finexercice



\exercice{1795, drutu, 2003/10/01}

\enonce{001795}{}
On d\'efinit la fonction $f$ sur $\R^2 \setminus \{ (x,x) \; ;\; x\in \R \}$ par 
$$
f(x,y)= \frac{\sin x -\sin y}{x-y}\; .
$$

Peut-on prolonger $f$ en une fonction continue sur $\R^2$ ? 
\finenonce{001795}



\finexercice



\exercice{1796, drutu, 2003/10/01}

\enonce{001796}{}
Etudier la continuit\'e en $(0,0)$ des fonctions suivantes :
\begin{enumerate}

\item $$f(x,y) = \left\{
         \begin{array}{cc}
         \frac{(x+y)^4}{x^4+y^4} & \mbox{ si }(x,y) \neq (0,0) \\
         1 & \mbox{ sinon. }
         \end{array}
         \right .$$
         

\item $$f(x,y) = \left\{
         \begin{array}{cc}
         \frac{|x|^3|y|^5}{(x^2+y^2)^2} & \mbox{ si }(x,y) \neq (0,0) \\
         0 & \mbox{ sinon. }
         \end{array}
         \right .$$
         
\item $$f(x,y) = \left\{
         \begin{array}{cc}
         \frac{e^{xy}-1}{x^2+y^2} & \mbox{ si }(x,y) \neq (0,0) \\
         0 & \mbox{ sinon. }
         \end{array}
         \right .$$        
\end{enumerate}
\finenonce{001796}



\finexercice



\exercice{1797, drutu, 2003/10/01}

\enonce{001797}{}
\begin{enumerate}
\item
$$
f(x,y)=\left\{
         \begin{array}{cc}
         \frac{(x+2y)^3y^3}{x^4+y^4} & \mbox{ si }(x,y) \neq (0,0) \\
         0 & \mbox{ sinon. }
         \end{array}
         \right .
         $$

     \item
$$
f(x,y)=\left\{
        \begin{array}{cc}
        \frac{x^6+x^2y^2}{x^2+y^2} & \mbox{ si }(x,y)\neq (0,0) \\
        0 & \mbox{ si }(x,y)=(0,0)
        \end{array}
        \right .
$$

  \item $$
f(x,y)=\left\{
         \begin{array}{cc}
         e^{\frac{x}{y}} & \mbox{ si }y \neq 0 \\
         0 & \mbox{ sinon. }
         \end{array}
         \right .
        $$

  \item $$
f(x,y)=\left\{
         \begin{array}{cc}
         \frac{\sin xy}{y} & \mbox{ si }y \neq 0 \\
         x & \mbox{ sinon. }
         \end{array}
         \right .
         $$

  \item $$
f(x,y)=\left\{
         \begin{array}{cc}
         \frac{\ln (1+x)-\ln (1+y)}{x-y} & \mbox{ si }x \neq y \\
         \frac{1}{1+x} & \mbox{ sinon. }
         \end{array}
         \right .
         $$ d\'efinie sur $D=\{(x,y) \mid x\geq 0, y\geq 0\}$.
\end{enumerate}
\finenonce{001797}



\finexercice



\exercice{2650, debievre, 2009/05/19}
\enonce{002650}{}
Soit $f:\R^2\to\R$ d\'efinie par
\begin{eqnarray*}
 f(x,y)&=&\frac{x^2y}{x-y},\quad \mathrm{si}\ x\not=y\\
&=& x,\quad \mathrm{si}\ x=y.
\end{eqnarray*}
\begin{enumerate}
 \item  Calculer les d\'eriv\'ees partielles $\frac{\partial f}{\partial x}(1,-2)$ et $\frac{\partial f}{\partial y}(1,-2)$. 
 \item Pour tout $v=(\cos \theta, \sin \theta)$, calculer $D_vf(1,-2)$. Pour quelles valeurs de $\theta\in [0,2\pi[$, $D_vf(1,-2)=0$ ?
 \item \'Etudier la continuit\'e de $f$ au point  $(1,1)\in\R^2$.
 \item \'Etudier la continuit\'e de $f$ au point  $(0,0)\in\R^2$.
 \item Montrer que les d\'eriv\'ees partielles $\frac{\partial f}{\partial x}(0,0)$ et $\frac{\partial f}{\partial y}(0,0)$ existent et les d\'eterminer. 
 \item Montrer que la d\'eriv\'ee directionnelle $D_vf(0,0)$ existe pour $v=(1,1)$, et la d\'eterminer.
On constatera que l'\'egalit\'e $D_vf(0,0)=\partial_xf(0,0) + \partial_yf(0,0)$ n'est pas satisfaite. Expliquer pourquoi cela ne contredit aucun th\'eor\`eme du cours. 
\end{enumerate}
\finenonce{002650}



\finexercice


\exercice{5554, rouget, 2010/07/15}
\enonce{005554}{***}
On pose $\begin{array}[t]{cccc}f_{x,y}~:&[-1,1]&\rightarrow&\Rr\\
 &t&\mapsto&xt^2+yt
\end{array}$ puis $F(x,y) =\underset{t\in[-1,1]}{\text{sup}}f_{x,y}(t)$. Etudier la continuité de $F$ sur $\Rr^2$.
\finenonce{005554}


\finexercice


\exercice{5900, rouget, 2010/10/16}
\enonce{005900}{**}
Soit $(E,\|\;\|)$ un espace vectoriel normé et $B=\{x\in E/\;\|x\| < 1\}$. Montrer que $\begin{array}[t]{cccc}
f~:&E&\rightarrow&B\\
 &x&\mapsto& \frac{x}{1+\|x\|}
\end{array}$ est un homéomorphisme.
\finenonce{005900}


\finexercice


\exercice{5901, rouget, 2010/10/16}
\enonce{005901}{**}
$E =\Rr^n$ est muni de sa structure euclidienne usuelle.
Montrer que $\begin{array}[t]{cccc}
f~:&E&\rightarrow&\Rr\\
 &x&\mapsto&\|x\|_2
\end{array}$ est différentiable sur $E\setminus\{0\}$ et préciser $df$. Montrer que $f$ n'est pas différentiable en $0$.
\finenonce{005901}


\finexercice

\section{223.03 Différentiabilité}


\exercice{1798, ridde, 1999/11/01}

\enonce{001798}{}
Soit $f : \begin{cases} \Rr^{2} \rightarrow \Rr \\ (x, y) \mapsto x \text{ si }
\left|x\right|>\left|y\right|\\ (x, y) \mapsto y \text{ si }
\left|x\right|<\left|y\right|\\ (x, y) \mapsto 0 \text{ si }
\left|x\right| = \left|y\right|\\ \end{cases} $. \\
\'Etudier la continuit\'e de $f$, l'existence des d\'eriv\'ees partielles et leur continuit\'e.
\finenonce{001798}



\finexercice



\exercice{1799, ridde, 1999/11/01}

\enonce{001799}{}
Soit $f : \begin{cases} \Rr^{2} \rightarrow \Rr \\ (x, y) \mapsto \frac{\sin (xy)}
{\left|x\right| + \left|y\right|} \text{ si } (x, y) \neq (0,0)\\
 (0, 0) \mapsto 0\\ \end{cases}$.\\
\'Etudier la continuit\'e de $f$ et l'existence des d\'eriv\'ees partielles.$f$ est-elle
 $C^{1}$ ?
\finenonce{001799}



\finexercice



\exercice{1800, ridde, 1999/11/01}

\enonce{001800}{}
Soit la fonction $f \colon \R^{2} \longrightarrow \R$ d\'efinie par
\[
\begin{array}{lll}
f(x, y)= xy\frac{x^{2}-y^{2}}
{x^{2} + y^{2}}& \mathrm{ si } &(x, y) \neq (0,0),\\
 f(0, 0)= 0
\end{array}
\]
\'Etudier la continuit\'e de $f$. Montrer que $f$ est de classe $C^{1}$.
\finenonce{001800}



\finexercice



\exercice{1801, ridde, 1999/11/01}

\exercice{}
\enonce{001801}{}
Soit $f : \R \rightarrow
\R$ d\'erivable. Calculer les d\'eriv\'ees partielles de :
\[
g (x, y) = f (x + y),\qquad h (x, y) = f (x^{2} + y^{2}),
\qquad k (x, y) = f (xy)
\]
\finenonce{001801}



\finexercice



\exercice{1802, ridde, 1999/11/01}

\enonce{001802}{}
Soit $f : \begin{cases} \Rr^{2} \rightarrow \Rr \\ (x, y) \mapsto \frac{x^{5}}
{ (y-x^{2})^{2} + x^{6}} \text{ si } (x, y) \neq (0,0)\\
 (0, 0) \mapsto 0\\ \end{cases}$.\\
 Montrer que $f$ admet une d\'eriv\'ee en $ (0, 0)$ suivant tout vecteur mais
 n'admet pas de d\'eveloppement limit\'e \`a l'ordre 1 en $ (0, 0)$.
\finenonce{001802}



\finexercice



\exercice{1803, gourio, 2001/09/01}

\enonce{001803}{}
Soit $f :  \R^{2} \rightarrow \R$, d\'efinie par  
\[
\begin{array}{lll}
 f(x, y)=x &\mathrm{si}&
\left|x\right|>\left|y\right|\\ 
f(x, y)=y &\mathrm{ si }&\left|x\right|<\left|y\right|\\ 
f(x, y)= 0 &\mathrm{ si }&
\left|x\right| = \left|y\right|. 
\end{array}
\]
\'Etudier la continuit\'e de $f$, l'existence des d\'eriv\'ees partielles et leur continuit\'e.

\finenonce{001803}



\finexercice



\exercice{1804, gourio, 2001/09/01}

\enonce{001804}{}
Soit $ a\in {\Rr}^{2}$ fix\'{e} ; l'application $ x\rightarrow \left\langle
x,a\right\rangle $  de $ {\Rr}^{2}$ usuel dans $ {\Rr}$ est-elle continue,
admet-elle des d\'{e}riv\'{e}es partielles, celles-ci sont elles continues ?
\finenonce{001804}



\finexercice



\exercice{1805, gourio, 2001/09/01}

\enonce{001805}{}
Soit $ f$ la fonction d\'{e}finie sur $ {\Rr}^{2}$ par :
\begin{itemize}
\item  si $ \left| x\right| \leq y,$ $ f(x,y)=x^{2}.$
\item $ f(x,y)=y^{2}$ sinon.
\end{itemize}
\'Etudier la continuit\'{e} de $ f$ et l'existence de d\'{e}riv\'{e}es
partielles.
\finenonce{001805}



\finexercice



\exercice{1806, gourio, 2001/09/01}

\enonce{001806}{}
Montrer qu'une norme $ N$ sur $ {\Rr}^{2} $ ne peut avoir des
d\'{e}riv\'{e}es partielles qui existent et qui soient continues en $ 0.$
\finenonce{001806}



\finexercice



\exercice{1807, maillot, 2001/09/01}

\enonce{001807}{}

Soient $\alpha>0$ et $f~:\,\R^2 \to \R$ d\'efinie par
$$\left \{\begin{array}{l}
\displaystyle f(x,y) = \frac{|x|^\alpha y}{x^2+y^4} \quad \mbox{si }(x,y)\neq
(0,0)\\
f(0,0)=0
\end{array}\right.$$
\begin{enumerate}
\item \begin{enumerate}
      \item Montrer que
$$\displaystyle \forall (x,y)\neq (0,0)\quad \left|f(x,y)\right|\leq
\left(x^2+y^4\right)^{\frac{2\alpha -3}4}.$$
       \item  Calculer $\displaystyle \lim_{
\begin{array}{c}y\to 0\\y\neq 0\end{array}}
\left|f(y^2,y)\right|$.

       \item \'Etudier la continuit\'e de $f$ en (0,0).
       \end{enumerate}
\item \begin{enumerate}
      \item Montrer que
$$\displaystyle\forall (x,y)\neq (0,0)\quad
\frac{\left|f(x,y)\right|}{\sqrt{x^2+y^2}}\leq
|x|^{\alpha-2}.$$
        \item  Calculer $\displaystyle \lim_{
\begin{array}{c}x\to 0\\x \neq 0\end{array}}
\frac{\left|f(x,x)\right|}{\sqrt 2 |x|}.$

        \item  \'Etudier la diff\'erentiabilit\'e de $f$ en (0,0).
\end{enumerate}
\end{enumerate}
\finenonce{001807}



\finexercice



\exercice{1808, drutu, 2003/10/01}

\enonce{001808}{}
\begin{enumerate}
\item Calculer la d\'eriv\'ee de la fonction $F(x,y)=e^{x^2+y^2}$ au point 
$P(1,0)$ suivant la bissectrice du premier quadrant.

\item Calculer la d\'eriv\'ee de la fonction $F(x,y,z)=x^2-3yz+5$ au point 
$P(1,2,1)$ dans une direction formant des angles \'egaux avec les trois axes de 
coordonn\'ees.

\item Calculer la d\'eriv\'ee de la fonction 
$F(x,y,z)=xy+yz+zx$ au point 
$M(2,1,3)$ dans la direction joignant ce point au point $N(5,5,15)$.
\end{enumerate}
\finenonce{001808}



\finexercice



\exercice{1809, drutu, 2003/10/01}

\enonce{001809}{}
Etudier la continuit\'e, ainsi que l'existence et la continuit\'e des d\'eriv\'ees partielles premi\`eres, des fonctions suivantes :
\begin{enumerate}


\item $$
f(x,y)=\left\{
         \begin{array}{cc}
         \frac{x|y|}{\sqrt{x^2+y^2}} & \mbox{ si }(x,y) \neq (0,0) \\
         0 & \mbox{ sinon. }
         \end{array}
         \right .
         $$      
         
\item $$
f(x,y)=\left\{
         \begin{array}{cc}
         \frac{x\sin y - y\sin x}{x^2+y^2} & \mbox{ si }(x,y) \neq (0,0) \\
         0 & \mbox{ sinon. }
         \end{array}
         \right .
         $$
         
\item $$
f(x,y)=\left\{
         \begin{array}{cc}
         e^{x \ln (x^2+y^2)} & \mbox{ si }(x,y) \neq (0,0) \\
         1 & \mbox{ sinon. }
         \end{array}
         \right .
         $$
\end{enumerate}
\finenonce{001809}



\finexercice



\exercice{1810, drutu, 2003/10/01}

\enonce{001810}{}
 On d\'efinit la fonction
$$
f(x,y)=\left\{
         \begin{array}{cc}
         \frac{x^3-y^3}{x^2+y^2} & \mbox{ si }(x,y) \neq (0,0) \\
         0 & \mbox{ sinon. }
         \end{array}
         \right .
         $$
Montrer que $\frac{\partial f}{\partial x} (x,y)$ et $\frac{\partial f}{\partial y}(x,y)$ existent en tout point de $\R^2$ et 
que $f$ est continue mais pas diff\'erentiable en $(0,0)$.
\finenonce{001810}



\finexercice



\exercice{1811, drutu, 2003/10/01}

\enonce{001811}{}
Soit $f:\rbrack 0,1\lbrack \times \rbrack 0,1\lbrack \to \R 
$, 
$$
f(x,y)=\left\{ 
        \begin{array}{cc}
        x(1-y) & \mbox{ si }x\leq y \\
        y(1-x) & \mbox{ si }x > y
        \end{array}
        \right . 
$$
Etudier la continuit\'e et la diff\'erentiabilit\'e de $f$.
\finenonce{001811}



\finexercice



\exercice{1812, drutu, 2003/10/01}

\enonce{001812}{}
 Soit $f:\R^2 \to \R $, 
$$
f(x,y)=\left\{ 
        \begin{array}{cc}
        \frac{x^2y+xy^2}{x^2+y^2} & \mbox{ si }(x,y)\neq (0,0) \\
        0 & \mbox{ si }(x,y)=(0,0) 
        \end{array}
        \right . 
$$
Montrer que $f$ est continue en $(0,0)$ et admet des d\'eriv\'ees partielles 
dans toutes les directions, mais n'y est pas diff\'erentiable.
\finenonce{001812}



\finexercice



\exercice{1813, drutu, 2003/10/01}

\enonce{001813}{}
Soit $f:\R^2 \to \R $, 
$$
f(x,y)=\left\{ 
        \begin{array}{cc}
        x^2y^2\sin \frac{1}{x} & \mbox{ si }x\neq 0 \\
        0 & \mbox{ si }x=0\; . 
        \end{array}
        \right . 
$$
 Montrer que la fonction $f$ est diff\'erentiable en tout point de $\R^2$ mais 
que $\partial_1f$ et $\partial_2f$ ne sont pas continues en certains points de 
$\R^2$.
\finenonce{001813}



\finexercice



\exercice{1814, drutu, 2003/10/01}

\enonce{001814}{}
 Etudier la diff\'erentiabilit\'e et la continuit\'e des d\'eriv\'ees partielles de la fonction $f:\R^2 \to \R $, 
$$
f(x,y)=\left\{ 
        \begin{array}{cc}
        (x^2+y^2)^{3/2}\sin \frac{1}{x^2+y^2} & \mbox{ si }(x,y)\neq (0,0) \\
        0 & \mbox{ si }(x,y)= (0,0)\; . 
        \end{array}
        \right . 
$$ 
\finenonce{001814}



\finexercice



\exercice{1815, drutu, 2003/10/01}

\enonce{001815}{}
Etudier la diff\'erentiabilit\'e en $(0,0)$ des fonctions d\'efinies par
\begin{enumerate}
\item $$
f(x,y)=\left\{
         \begin{array}{cc}
         \frac{x^3y}{x^4+y^2} & \mbox{ si }(x,y) \neq (0,0) \\
         0 & \mbox{ sinon. }
         \end{array}
         \right .
         $$
         
\item $$
f(x,y)=\left\{
         \begin{array}{cc}
         \frac{xy^3}{x^4+y^2} & \mbox{ si }(x,y) \neq (0,0) \\
         0 & \mbox{ sinon. }
         \end{array}
         \right .
         $$
\end{enumerate}
\finenonce{001815}



\finexercice



\exercice{1816, drutu, 2003/10/01}

\enonce{001816}{}
Calculer les d\'eriv\'ees partielles (d'ordre un) des fonctions 
suivantes en un point arbitraire du domaine de d\'efinition.
\begin{enumerate}
\item $f(x,y)=x^2e^{xy}$ ;

\item $g(x,y,z)=x^2y^3\sqrt{z}$ ;

\item $h(x,y)=\ln (x+\sqrt{x^2+y^2})$.
\end{enumerate}
\finenonce{001816}



\finexercice



\exercice{1817, drutu, 2003/10/01}

\enonce{001817}{}
Calculer les d\'eriv\'ees partielles (d'ordre un) de la fonction 
$f(x,y)=\sqrt{xy+\frac{x}{y}}$ en $(2,1)$.
\finenonce{001817}



\finexercice



\exercice{1818, drutu, 2003/10/01}

\enonce{001818}{}
On d\'efinit la fonction
$$
f(x,y)=\left\{
         \begin{array}{cc}
         \frac{xy}{x^2+y^2} & \mbox{ si }(x,y) \neq (0,0) \\
         0 & \mbox{ sinon. }
         \end{array}
         \right .
         $$
Montrer que $\frac{\partial }{\partial x} f(x,y)$ et $\frac{\partial }{\partial y}f(x,y)$ 
existent en tout point de $\R^2$ bien 
que $f$ ne soit pas continue en $(0,0)$.
\finenonce{001818}



\finexercice



\exercice{1819, drutu, 2003/10/01}

\enonce{001819}{}
\begin{enumerate}
\item  Calculer la d\'eriv\'ee de la fonction 
$F(x,y)=x^2-xy-2y^2$ 
au point $P(1,2)$ dans une direction formant avec l'axe $Ox$ un angle de 
$\frac{\pi}{3}$.

\item  Calculer la d\'eriv\'ee de la fonction $F(x,y)=x^3-2x^2y+xy^2+1$ au point 
$P(1,2)$ dans la direction joignant ce point au point $M(4,6)$.

 
\item  Calculer la d\'eriv\'ee de la fonction $F(x,y)=\ln \sqrt{x^2+y^2}$ au point 
$P(1,1)$ suivant la bissectrice du premier quadrant.
\end{enumerate}
\finenonce{001819}



\finexercice



\exercice{1820, drutu, 2003/10/01}

\enonce{001820}{}
Calculer les diff\'erentielles des fonctions suivantes en un point arbitraire du domaine de d\'efinition :
\begin{enumerate}
\item $f(x,y)=\sin ^2 x+\cos ^2 y$ ;

\item $f(x,y)=\ln \left( 1+\frac{x}{y}\right)$. 
\end{enumerate}
\finenonce{001820}



\finexercice



\exercice{1821, drutu, 2003/10/01}

\enonce{001821}{}
Calculer $df\, (1,1)$, si $f(x,y)=\frac{x}{y^2}$.
\finenonce{001821}



\finexercice



\exercice{1822, drutu, 2003/10/01}

\enonce{001822}{}
 Calculer la d\'eriv\'ee de la fonction $F(x,y,z)=\ln \left( 
e^x+e^y+e^z\right)$ \`a l'origine dans une direction formant avec les axes de 
coordonn\'ees $x,y,z$ les angles $\alpha ,\beta , \gamma$.
\finenonce{001822}



\finexercice



\exercice{2628, debievre, 2009/05/19}
\enonce{002628}{}
Trouver l'\'equation du plan tangent pour chaque surface ci-dessous, au point
$(x_0,y_0,z_0)$ donn\'e:
\begin{enumerate}
 \item  $z=\sqrt{19-x^2-y^2},\quad (x_0,y_0,z_0)=(1,3,3)$; 
 \item $z=\sin(\pi xy)\exp(2x^2y-1), \quad (x_0,y_0,z_0)=(1,1/2,1)$.
\end{enumerate}
\finenonce{002628}


\finexercice


\exercice{2629, debievre, 2009/05/19}
\enonce{002629}{}
On demande \`a un \'etudiant de trouver l'\'equation du plan tangent \`a
la surface d'\'equation
$z=x^4-y^2$ au point $(x_0,y_0,z_0)=(2,3,7)$. Sa r\'eponse est
\[
z=4x^3(x-2)-2y(y-3).
\]
\begin{enumerate}
 \item Expliquer, sans calcul, pourquoi cela ne peut en aucun 
cas \^etre la
bonne r\'eponse.
 \item Quelle est l'erreur commise par l'\'etudiant?
 \item Donner la r\'eponse correcte.
\end{enumerate}
\finenonce{002629}


\finexercice


\exercice{2630, debievre, 2009/05/19}
\enonce{002630}{}
Trouver les points sur 
le parabolo\"\i de $z=4x^2 +y^2$ o\`u le plan tangent est parall\`ele au plan
$x+2y+z=6$. M\^eme question avec le plan $3x+5y-2z=3$. 
\finenonce{002630}


\finexercice


\exercice{2631, debievre, 2009/05/19}
\enonce{002631}{}
Soit $C$ le c\^one d'\'equation $z^2= x^2+y^2$
et $C^+$ le demi-c\^one o\`u $z \geq 0$. Pour un 
point quelconque $M_0$ de  $C\setminus \{(0,0,0)\}$, de
coordonn\'ees $(x_0,y_0,\pm \sqrt{x_0^2+y_0^2})$, 
on note ${\cal P}_{M_0}$ le plan tangent au c\^one $C$ en $M_0$.
\begin{enumerate}
 \item  D\'eterminer un vecteur normal et l'\'equation du plan ${\cal P}_{M_0}$.
 \item Montrer que l'intersection du c\^one $C$ avec 
le plan vertical d'\'equation ${y=ax}$ o\`u $a\in \R$ est constitu\'ee
de deux droites ${\cal D}_1$ et ${\cal D}_2$
et que 
 l'intersection du demi-c\^one $C^+$ avec 
ce plan vertical  est constitu\'ee
de deux demi-droites ${\cal D}^+_1$ et ${\cal D}^+_2$. 
 \item Montrer que le plan tangent au c\^one $C$ 
est le m\^eme en tout point de ${\cal D}_1\setminus \{(0,0,0)\}$ (respectivement  en tout point de
${\cal D}_2\setminus \{(0,0,0)\}$).
\end{enumerate}
\finenonce{002631}


\finexercice


\exercice{2632, debievre, 2009/05/19}
\enonce{002632}{}
 Soit $f$ la fonction d\'efinie sur
$\R^2$ par $f(x,y)= x^2-2y^3$.
\begin{enumerate}
 \item D\'eterminer l'\'equation du plan tangent 
${\cal P}_{M_0}$ 
au graphe $G_f$ de $f$ en un point quelconque 
$M_0$
de $G_f$.
 \item Pour le point 
$M_0$ de coordonn\'ees 
$(2,1,2)$, d\'eterminer tous les points 
$M$ 
tels que le plan tangent 
en $M$  soit
parall\`ele \`a
${\cal P}_{M_0}$.
\end{enumerate}
\finenonce{002632}


\finexercice


\exercice{2633, debievre, 2009/05/19}
\enonce{002633}{}

Soit la fonction $f\colon\R^2\to\R$ d\'efinie par
$$
f(x,y)=\frac{xy^2}{x^2+y^2}, \quad (x,y)\not=(0,0)
$$
et $f(0,0)=0$.
\begin{enumerate}
 \item Montrer que $f$ est continue et que,
quel que soit $v\in\R^2$, la d\'eriv\'ee directionnelle
 $D_vf(x,y)$ existe en chaque
$(x,y)\in\R^2$  mais que $f$ n'est pas diff\'erentiable en
$(0,0)$.
 \item La d\'eriv\'ee directionnelle $D_vf(0,0)$ 
est-elle lin\'eaire en $v$? Les droites
appartenant \`a la famille 
des droites passant par l'origine
et de vecteurs directeurs $(v, D_vf(0,0))\in\R^3$, forment-elles un plan? 
Expliquer comment on peut observer la r\'eponse sur la figure.
 \item Le vecteur $v$ \'etant fix\'e, qu'est-ce qu'on peut dire de
la continuit\'e de $D_vf(x,y)$ en $(x,y)$?

{\center
\includegraphics[height=6cm, keepaspectratio]{../images/img002633-1}
}

\end{enumerate}
\finenonce{002633}


\finexercice


\exercice{2634, debievre, 2009/05/19}
\enonce{002634}{}
Utiliser une approximation affine bien choisie pour calculer une valeur approch\'ee des nombres suivants:
\[
\exp[\sin(3.16)\cos(0.02))], \quad \arctan[\sqrt{4.03}-2\exp(0.01)].
\]
\finenonce{002634}


\finexercice


\exercice{2651, debievre, 2009/05/19}
\enonce{002651}{}
\begin{enumerate}
 \item  Soit $f:D\subset \R^m\to\R^n$ et $a\in D$. Donner la d\'efinition de 
``$f$ est diff\'erentiable en $a$''. 
 \item Montrer que, si $f$ est diff\'erentiable en $a$, alors toutes ses d\'eriv\'ees partielles existent. Exprimer le lien entre la diff\'erentielle $d f_a$ de $f$ en $a$ et les d\'eriv\'ees partielles de $f$ en $a$. 
 \item Les affirmations suivantes, sont-elles vraies ou fausses? On justifiera bri\`evement sa r\'eponse.

(A) Si $f$ est diff\'erentiable en $a$, alors elle y est continue.

(B) Si toutes les d\'eriv\'ees partielles de $f$ en $a$ existent, alors $f$ est diff\'erentiable en $a$.
\end{enumerate}
\finenonce{002651}



\finexercice


\exercice{2652, debievre, 2009/05/19}
\enonce{002652}{}
Utiliser une approximation affine bien choisie pour calculer une valeur approch\'ee de
$$\exp[-0.02\sqrt{4.03}].$$
\finenonce{002652}



\finexercice


\exercice{5555, rouget, 2010/07/15}
\enonce{005555}{***T}
Déterminer la classe de $f$ sur $\Rr^2$ où $f(x,y)=\left\{\begin{array}{l}
0\;\text{si}\;(x,y)=(0,0)\\
\frac{xy(x^2-y^2)}{x^2+y^2}\;\text{si}\;(x,y)\neq(0,0)
\end{array}
\right.$.
\finenonce{005555}


\finexercice


\exercice{5888, rouget, 2010/10/16}
\enonce{005888}{*** I}
Pour  $(x,y)\in\Rr^2$, on pose $f(x,y) =\left\{
\begin{array}{l}
 \frac{xy(x^2-y^2)}{x^2+y^2}\;\text{si}\;(x,y)\neq(0,0)\\
\rule{0mm}{5mm}0\;\text{si}\;(x,y)=(0,0)
\end{array}
\right.$. Montrer que $f$ est de classe $C^1$ (au moins) sur $\Rr^2$.
\finenonce{005888}


\finexercice


\exercice{5894, rouget, 2010/10/16}
\enonce{005894}{** I}
Extremums des fonctions suivantes :

\begin{enumerate}
 \item  $f(x,y) = x^3+3x^2y -15x-12y$

 \item  $f(x,y) = -2(x-y)^2+x^4+y^4$.
\end{enumerate}
\finenonce{005894}


\finexercice


\exercice{5895, rouget, 2010/10/16}
\enonce{005895}{*** I}
Soit $\begin{array}[t]{cccc}
f~:&GL_n(\Rr)&\rightarrow&M_n(\Rr)\\
 &A&\mapsto&A^{-1}
\end{array}$. Montrer que $f$ est différentiable en tout point de $M_n(\Rr)\setminus\{0\}$ et déterminer sa différentielle.
\finenonce{005895}


\finexercice


\exercice{5896, rouget, 2010/10/16}
\enonce{005896}{*}
Déterminer $\text{Max}\{|\sin z|,\;z\in\Cc,\;|z|\leqslant1\}$.
\finenonce{005896}


\finexercice


\exercice{5899, rouget, 2010/10/16}
\enonce{005899}{**}
Déterminer la différentielle en tout point de $\begin{array}[t]{cccc}
f~:&\Rr^3\times\Rr^3&\rightarrow&\Rr\\
 &(x,y)&\mapsto&x.y
\end{array}$ et $\begin{array}[t]{cccc}
g~:&\Rr^3\times\Rr^3&\rightarrow&\Rr\\
 &(x,y)&\mapsto&x\wedge y
\end{array}$.
\finenonce{005899}


\finexercice

\section{223.04 Dérivée partielle}


\exercice{2622, debievre, 2009/05/19}
\enonce{002622}{}
D\'eterminer, pour chacune des fonctions suivantes, le domaine 
de d\'efinition $D_f$. 
Pour chacune des fonctions,
calculer ensuite 
les d\'eriv\'ees partielles en chaque point du
domaine de d\'efinition  lorsqu'elles existent:
\begin{enumerate}
 \item $f(x,y)=x^2\exp(xy)$,
 \item $f(x,y)=\ln(x+\sqrt{x^2+y^2})$,
 \item $f(x,y)=\sin^2 x+ \cos^2y$,
 \item $f(x,y,z)=x^2y^2\sqrt{z}$.
\end{enumerate}
\finenonce{002622}


\finexercice


\exercice{2623, debievre, 2009/05/19}
\enonce{002623}{}

Soit $f$ la fonction sur $\R^2$ d\'efinie par $f(x,y)= x\cos y + y\exp x$. 
\begin{enumerate}
 \item  Calculer ses d\'eriv\'ees partielles.
 \item Soit $v=(\cos \theta, \sin \theta)$, $\theta\in [0,2\pi[$. Calculer $D_vf(0,0)$. Pour quelle(s) valeurs de $\theta$ cette d\'eriv\'ee directionnelle de $f$ est-elle maximale/minimale? Que cela signifie-t-il?
\end{enumerate}
\finenonce{002623}


\finexercice



\exercice{2624, debievre, 2009/05/19}
\enonce{002624}{}
Soit $f:\R^2\rightarrow \R$ la fonction
d\'efinie par $f(x,y)=(x^2+y^2)^x$ pour $(x,y)\not=(0,0)$ et
$f(0,0)= 1$.
\begin{enumerate}
 \item La fonction $f$ est-elle continue en $(0,0)$?
 \item D\'eterminer les d\'eriv\'ees partielles de $f$ en un point
quelconque distinct de l'origine.
 \item  La fonction $f$ admet-elle des d\'eriv\'ees partielles par
rapport \`a $x$, \`a $y$ en $(0,0)$?
\end{enumerate}
\finenonce{002624}


\finexercice


\exercice{2625, debievre, 2009/05/19}
\enonce{002625}{}
Soit $f\colon \R^2\rightarrow \R$ la fonction
d\'efinie par 
\begin{eqnarray*}
f(x,y)&=&\frac{x^2y +3y^3}{x^2+y^2} \ \mathrm{ pour }\ (x,y)\not=(0,0),\\
f(0,0)&=& 0.
\end{eqnarray*}
\begin{enumerate}
 \item La fonction $f$ est-elle continue en $(0,0)$? Justifier la
r\'eponse.
 \item La fonction $f$ admet-elle des d\'eriv\'ees partielles par
rapport \`a $x$, \`a $y$ en $(0,0)$? Donner la ou les valeurs le cas
\'ech\'eant et justifier la r\'eponse.
 \item La fonction $f$ est-elle diff\'erentiable en $(0,0)$? Justifier
la r\'eponse. 
 \item D\'eterminer les d\'eriv\'ees partielles de $f$ en un point
$(x_0,y_0)\not=(0,0)$.
 \item D\'eterminer l'\'equation du plan tangent au graphe de $f$ au 
point $(1,1,2)$.
 \item Soit $F:\R^2\rightarrow \R^2$ la fonction d\'efinie par
$F(x,y)=(f(x,y),f(y,x))$. D\'eterminer la matrice jacobienne de $F$ au
point $(1,1)$. La fonction $F$ admet-elle une r\'eciproque locale au
voisinage du point $(2,2)$?
\end{enumerate}
\finenonce{002625}


\finexercice


\exercice{2626, debievre, 2009/05/19}
\enonce{002626}{}
Soit $f\colon\R^3 \rightarrow \R$ une
fonction de classe $C^1$
et soit  $g\colon \R^3 \rightarrow \R$
la fonction d\'efinie par
\[
g(x,y,z) = f(x-y,y-z,z-x). 
\]
Montrer que
\begin{equation}
\frac{\partial g}{\partial x}  + \frac{\partial g}{\partial y}  + \frac{\partial g}{\partial z} = 0.
\label{ex3}
\end{equation}
\finenonce{002626}


\finexercice


\exercice{2627, debievre, 2009/05/19}
\enonce{002627}{}
On consid\`ere les fonctions $f\colon \R^2\longrightarrow \R^3$ et
$g\colon \R^3\longrightarrow \R$ d\'efinies par
\[
f(x,y)=(\sin (xy), y\cos x, xy\sin(xy)\exp(y^2)),\quad 
g(u,v,w)= uvw .
\]
\begin{enumerate}
 \item  Calculer explicitement $g\circ f$.
 \item En utilisant l'expression trouv\'ee en (1), calculer les d\'eriv\'ees partielles de $g\circ f$.
 \item  D\'eterminer les matrices jacobiennes $J_f(x,y)$ et $J_g(u,v,w)$ de $f$ et de $g$. 
 \item Retrouver le r\'esultat sous (2.) en utilisant un produit appropri\'e de matrices jacobiennes. 
\end{enumerate}
\finenonce{002627}


\finexercice

\exercice{4136, quercia, 2010/03/11}
\enonce{004136}{Calcul de dérivées partielles}

Calculer les dérivées partielles des fonctions :

\begin{enumerate}
  \item $f(x,y,z) = (x+z)^{(y^z)}$\quad
  \item $f(x,y) = \min(x,y^2)$  \qquad
  \item $f(x,y) = \begin{cases}\frac {g(x)-g(y)}{x-y} &\text{ si } x \ne y \cr
                     g'(x) &\text{ si } x = y.\end{cases}$

\end{enumerate}
\finenonce{004136}



\finexercice


\exercice{4137, quercia, 2010/03/11}
\enonce{004137}{DL d'ordre 1}

Soit $f : {\R^3} \to \R$ de classe $\mathcal{C}^1$ telle que
$f(0,1,1) = 0$, $\frac{\partial f}{\partial x}(0,1,1) = 1$, $\frac{\partial f}{\partial y}(0,1,1) = 2$, $\frac{\partial f}{\partial z}(0,1,1) = 3$.

Peut-on déterminer $\lim_{t\to0} \frac{f(t^2,\ch t,e^t)}{f(t,\cos t,\ch t)}$ ?

\finenonce{004137}



\finexercice


\exercice{4138, quercia, 2010/03/11}
\enonce{004138}{Simplification}

Soit $f(x,y) = \arcsin\left(\frac{1+xy}{\sqrt{(1+x^2)(1+y^2)}}\right)$
et $g(x,y) = \arctan x - \arctan y$.

\begin{enumerate}
  \item Vérifier que $f$ est définie sur $\R^2$.
  \item Calculer les dérivées partielles premières de $f$ et de $g$.
    
  \item Simplifier $f$ à l'aide de $g$.
    

\end{enumerate}
\finenonce{004138}



\finexercice


\exercice{4139, quercia, 2010/03/11}
\enonce{004139}{Somme des angles d'un triangle}

Sur quelle partie $D$ de $\R^3$ la fonction
$$f : (x,y,z)  \mapsto \arccos\Bigl(\frac{x^2+y^2-z^2}{2xy}\Bigr)
                 + \arccos\Bigl(\frac{y^2+z^2-x^2}{2yz}\Bigr)
                 + \arccos\Bigl(\frac{z^2+x^2-y^2}{2zx}\Bigr)$$
est-elle définie ? Montrer que $f$ est constante lorsque $x,y,z$ sont strictement
positifs.
\finenonce{004139}



\finexercice


\exercice{4140, quercia, 2010/03/11}
\enonce{004140}{Intégrale fonction de paramètres}

Soit $f : {\R^2} \to \R$ de classe $\mathcal{C}^1$ et
$g : {\R^n} \to \R, {(x_1,\dots,x_n)} \mapsto
{ \int_{t=0}^1 f(t,x_1t+x_2t^2 + \dots + x_nt^n)\,d t.}$

Montrer que $g$ est de classe $\mathcal{C}^1$ et calculer ses dérivées partielles.
\finenonce{004140}



\finexercice


\exercice{4141, quercia, 2010/03/11}
\enonce{004141}{Dérivées secondes composées}

Soient ${u,v,f,g} : {\R^2} \mapsto \R$ des fonctions de classe $\mathcal{C}^2$ liées
par la relation :
$$\forall\ (x,y) \in {\R^2},\ f(x,y) = g(u(x,y),v(x,y)).$$

Calculer les dérivées partielles premières et secondes de $f$
en fonction de celles de $g$.

\finenonce{004141}



\finexercice


\exercice{4142, quercia, 2010/03/11}
\enonce{004142}{Les polynômes complexes sont harmoniques}

Soient $P \in \C[X]$ et $f : {\R^2} \to \C, {(x,y)} \mapsto {P(x+iy).}$
Montrer que $\frac{\partial^2 f}{\partial x^2} + \frac{\partial^2 f}{\partial y^2} = 0$.
\finenonce{004142}



\finexercice


\exercice{4143, quercia, 2010/03/11}
\enonce{004143}{Laplacien en polaires}

Soit $f : {\R^2} \to \R$ de classe $\mathcal{C}^2$, et
$g(\rho,\theta) = f(\rho\cos\theta,\rho\sin\theta)$.
On pose $\Delta f = \frac{\partial^2 f}{\partial x^2} + \frac{\partial^2 f}{\partial y^2}$ (laplacien de $f$).

\begin{enumerate}
  \item Calculer $\frac{\partial g}{\partial \rho}$, $\frac{\partial g}{\partial \theta}$, $\frac{\partial^2 g}{\partial \rho^2}$, $\frac{\partial^2 g}{\partial \theta^2}$ en fonction des dérivées
    partielles de $f$.
  \item Exprimer $\Delta f$ en fonction des dérivées de $g$.
    

\end{enumerate}
\finenonce{004143}



\finexercice


\exercice{4144, quercia, 2010/03/11}
\enonce{004144}{Laplacien en sphériques}

Soient $f : {\R^3} \mapsto \R$ de classe $\mathcal{C}^2$,
$\Phi : {\R^3} \to {\R^3}, {(r,\theta,\varphi)} \mapsto
{(x,y,z)}$ avec $\begin{cases} x = r\cos\theta\cos\varphi\cr
                         y = r\sin\theta\cos\varphi\cr
                         z = r\sin\varphi,\cr \end{cases}$
et $F = f\circ \Phi$.
Vérifier que :
$$(\Delta f)\circ \Phi = \frac{\partial^2 F}{\partial r^2} + \frac2r\frac{\partial F}{\partial r} + \frac1{r^2}\frac{\partial^2 F}{\partial \varphi^2}
  -\frac{\tan\varphi}{r^2}\frac{\partial F}{\partial \varphi} +\frac1{r^2\cos^2\varphi}\frac{\partial^2 F}{\partial \theta^2}.$$

{\it Pour cet exercice, il est conseillé de prendre la feuille dans le sens de
la longueur, et d'y aller calmement, en vérifiant ses calculs.}
\finenonce{004144}



\finexercice


\exercice{4145, quercia, 2010/03/11}
\enonce{004145}{Laplacien en dimension $n$}

Soit $f$ une application de classe $\mathcal{C}^2$ de $\R^{+*}$ dans $\R$.

On définit une application $F$ de $\R^n\setminus\{\vec 0\}$ dans $\R$ par :
$F(x_1,\dots,x_n) = f(\sqrt{x_1^2 + \dots + x_n^2}\,)$.

Calculer le laplacien de $F$ en fonction de $f$.

\finenonce{004145}



\finexercice


\exercice{4146, quercia, 2010/03/11}
\enonce{004146}{Contre-exemple au théorème de Schwarz}

Soit $f:\R \to \R$ une fonction $\pi$-périodique de classe $\mathcal{C}^2$.
On pose pour $(x,y) \in \R^2$ :
$g(x,y) = r^2f(\theta)$ avec $(x,y)=(r\cos\theta,r\sin\theta)$.
Calculer $\frac{\partial g}{\partial x}(0,y)$ et $\frac{\partial g}{\partial y}(x,0)$ en fonction de $f$.
En déduire les valeurs de $\frac{\partial^2 g}{\partial y \partial x}(0,0)$ et $\frac{\partial^2 g}{\partial x \partial y}(0,0)$.
Construire un exemple précis (donner $g(x,y)$ en fonction de $x$ et $y$)
pour lequel ces deux dérivées sont distinctes.

\finenonce{004146}



\finexercice


\exercice{4147, quercia, 2010/03/11}
\enonce{004147}{Contre-exemple au théorème de Schwarz (Centrale MP 2003)}
\smallskip
Soit $f(x,y) = \frac{x^3y}{x^2+y^2}$ si $(x,y)\ne 0$ et $f(0,0)=0$.

\begin{enumerate}
  \item \'Etudier la continuité de~$f$ et de ses dérivées partielles premières sur~$\R^2$.
    
  \item Montrer que~$\frac{\partial^2 f}{\partial x \partial y}(0,0)\ne\frac{\partial^2 f}{\partial y \partial x}(0,0)$.
    

\end{enumerate}
\finenonce{004147}



\finexercice


\exercice{4148, quercia, 2010/03/11}
\enonce{004148}{Dérivées d'ordre $k$ distinctes}

Trouver ${f_k} : {\R^2} \to \R$ de classe $\mathcal{C}^k$ telle que les $k+1$ dérivées d'ordre $k$ en
$(0,0)$ soient distinctes.

\finenonce{004148}



\finexercice


\exercice{4149, quercia, 2010/03/11}
\enonce{004149}{Les isométries conservent le laplacien}

Soit $\varphi : {\R^2} \to {\R^2}$ une isométrie pour la norme $\|\ \|_2$.

\begin{enumerate}
  \item Montrer que la matrice jacobienne de $\varphi$ est constante, égale à la matrice
    dans la base canonique de $\R^2$ de la partie linéaire de $\varphi$.
  \item Soit $f : {\R^2} \to \R$ de classe $\mathcal{C}^2$.
    Montrer que $(\Delta f)\circ \varphi = \Delta(f\circ \varphi)$.

\end{enumerate}
\finenonce{004149}



\finexercice


\exercice{4150, quercia, 2010/03/11}
\enonce{004150}{Changement de variables affine}

Soit $\varphi : {\R^2} \to {\R^2}$ une application affine.

\begin{enumerate}
  \item Montrer que la matrice jacobienne, $J$, de $\varphi$ est constante.
  \item Soit $f : {\R^2} \to \R$ de classe $\mathcal{C}^2$. Pour $A \in \R^2$, on note
    $H_f(A) = \begin{pmatrix} \frac{\partial^2 f}{\partial x^2}(A) &\frac{\partial^2 f}{\partial y \partial x}(A) \cr \frac{\partial^2 f}{\partial x \partial y}(A) &\frac{\partial^2 f}{\partial y^2}(A) \cr\end{pmatrix}$
    ({\it matrice Hessienne de $f$}). Montrer que :
    $\forall\ A \in \R^2,\ H_{f\circ\varphi}(A) = {}^tJH_f(\varphi(A))J$.


\end{enumerate}
\finenonce{004150}



\finexercice


\exercice{4151, quercia, 2010/03/11}
\enonce{004151}{Formule de Leibniz}

Soient ${f,g} : {\R^2} \to \R$ de classe $\mathcal{C}^n$. Calculer
$\frac{\partial^n(fg)}{\partial x^k\partial y^{n-k}}$ en fonction des dérivées
de $f$ et $g$.


\finenonce{004151}



\finexercice


\exercice{4152, quercia, 2010/03/11}
\enonce{004152}{Intégration de formes différentielles}

Déterminer les fonctions $f : {D \subset {\R^2}} \to \R$ vérifiant :

\begin{enumerate}
  \item $\begin{cases} \frac{\partial f}{\partial x} = \frac {2+x}y \cr
             \frac{\partial f}{\partial y} = \frac {2+y}x \cr\end{cases}$
  \item $\begin{cases} \frac{\partial f}{\partial x} = \frac {1-y}{(x+y+1)^2} \cr
             \frac{\partial f}{\partial y} = \frac {2+x}{(x+y+1)^2} \cr\end{cases}$
  \item $\begin{cases} \frac{\partial f}{\partial x} = \frac{y^2}{(x+y)^2} \cr
             \frac{\partial f}{\partial y} = \frac{x^2}{(x+y)^2} \cr \end{cases}$
  \item $\begin{cases} \frac{\partial f}{\partial x} = 2x + \frac 1y \cr
             \frac{\partial f}{\partial y} = 2y - \frac x{y^2} \cr \end{cases}$
\end{enumerate}
\finenonce{004152}



\finexercice


\exercice{4153, quercia, 2010/03/11}
\enonce{004153}{Formes différentielles exactes}

Trouver les fonctions $f:\R \to \R$ de classe $\mathcal{C}^1$ telles que la forme
différentielle
$\omega = f(y)(xe^y\,dx + y\,dy)$ soit exacte. Déterminer alors ses
primitives.

\finenonce{004153}



\finexercice


\exercice{4154, quercia, 2010/03/11}
\enonce{004154}{}
Quelles sont les applications $f : {\R^2} \to \R$ de classe $\mathcal{C}^1$ telles
que la forme différentielle : \par
$\omega = f(x,y)\,d(x^2+y^2)$ soit exacte ?

\finenonce{004154}



\finexercice


\exercice{4155, quercia, 2010/03/11}
\enonce{004155}{}
Trouver les fonctions ${f,g} : \R \to \R$ de classe $\mathcal{C}^1$ telles que la forme
différentielle :\par
$\omega = 2xz\,dx + f(y)g(z)\,dy + \left(x^2+\frac{y^2}2\right)dz$
soit exacte. Déterminer alors ses primitives.


\finenonce{004155}



\finexercice


\exercice{4156, quercia, 2010/03/11}
\enonce{004156}{\'Equation associée à une différentielle exacte}

\begin{enumerate}
  \item Déterminer les fonctions $f : {\R^{+*}\times \R^{+*}} \to  \R$ vérifiant :
    $\begin{cases} \frac{\partial f}{\partial x}(x,y) = \frac {\ln x + y - 1}{x^2y} \cr
             \frac{\partial f}{\partial y}(x,y) = \frac {\ln x }{xy^2}        \cr\end{cases}$
    
  \item Application :
    Résoudre l'équation différentielle : $(x\ln x) y' + (\ln x + y - 1)y = 0$.
    

\end{enumerate}
\finenonce{004156}



\finexercice


\exercice{4157, quercia, 2010/03/11}
\enonce{004157}{\'Equation aux dérivées partielles}

Trouver les fonctions polynomiales $f : {\R^2} \to \R$ vérifiant :
$x\frac{\partial f}{\partial x} + y\frac{\partial f}{\partial y} = 4f$.
\finenonce{004157}



\finexercice


\exercice{4158, quercia, 2010/03/11}
\enonce{004158}{DL d'ordre 2}

Soit $f : {\R^2} \to \R$ de classe $\mathcal{C}^2$.
Démontrer que :
$$f(a+h,b+k) = f(a,b) + \left( h\frac{\partial f}{\partial x} + k\frac{\partial f}{\partial y} \right)(a,b)
  + \frac 12 \left( h^2\frac{\partial^2 f}{\partial x^2} + 2hk\frac{\partial^2 f}{\partial x \partial y} + k^2\frac{\partial^2 f}{\partial y^2} \right)(a,b)
  +  o(h^2+k^2).$$
\finenonce{004158}



\finexercice


\exercice{4159, quercia, 2010/03/11}
\enonce{004159}{Ajustement linéaire}

Problème d'ajustement linéaire : Etant donné $n$ couples de réels
$(x_i,y_i)\ 1 \le i \le n$,
on cherche une droite $D$ d'équation $y = ax+b$ telle que
$\mu(a,b) = \sum_{i=1}^n (y_i-ax_i-b)^2$ soit minimal.

On note $\overline  x    = \frac 1n \sum_{i=1}^n x_i$,
        $\overline  y    = \frac 1n \sum_{i=1}^n y_i$,
        $\overline {x^2} = \frac 1n \sum_{i=1}^n x_i^2$,
        $\overline {xy} = \frac 1n \sum_{i=1}^n x_iy_i$,
et on suppose $\overline {x^2} \ne \overline x^2$.

\begin{enumerate}
  \item  Résoudre le problème.
  \item  Interpréter la relation $\overline {x^2} \ne \overline x^2$
     à l'aide de l'inégalité de Cauchy-Schwarz.
\end{enumerate}
\finenonce{004159}



\finexercice


\exercice{4160, quercia, 2010/03/11}
\enonce{004160}{Jacobien des fonctions symétriques}

Soit $f : {\R^n} \to {\R^n}, {(x_1,\dots,x_n)} \mapsto {(\sigma_1,\dots,\sigma_n)}$

où $\sigma_1,\dots,\sigma_n$ sont les fonctions symétriques élémentaires de
$x_1,\dots,x_n$.
Calculer le déterminant jacobien de~$f$.


\finenonce{004160}



\finexercice


\exercice{4161, quercia, 2010/03/11}
\enonce{004161}{Changement de variables}

On pose $f(x,y) = (x+y,xy) = (u,v)$.
Montrer que $f$ induit un $\mathcal{C}^1$-difféomorphisme de $U$ sur $V$ où $U$ et $V$
sont des ouverts de $\R^2$ à préciser.
Chercher l'expression de $f^{-1}$ et vérifier que le produit des matrices
jacobiennes est égal à $I$.

\finenonce{004161}



\finexercice


\exercice{4162, quercia, 2010/03/11}
\enonce{004162}{Changement de variables}

Soit $f : {\R^3} \to {\R^3}, {(x,y,z)} \mapsto {(e^{2y}+e^{2z}, e^{2x}-e^{2z}, x-y).}$
Montrer que $f$ induit un $\mathcal{C}^1$-difféomorphisme de $\R^3$ sur un ouvert à préciser.

\finenonce{004162}



\finexercice


\exercice{4163, quercia, 2010/03/11}
\enonce{004163}{Inégalité de Taylor-Lagrange}

Soit $U$ un ouvert convexe de $\R^p$ et $f : U \to \R$ de classe $\mathcal{C}^2$
dont les dérivées secondes sont bornées :

$$\forall\ i,j,\ \forall\ A \in U,\
\left|\frac{\partial^2 f}{\partial {x_i}\partial {x_j}}(A)\right| \le M.$$

\begin{enumerate}
  \item Montrer que : $\forall\ A,B \in U,\ \bigl|f(B)-f(A) - df_A(\vec{AB})\bigr|
          \le \frac{M\|\vec{AB}\|_1^2}2$.

  \item Montrer que : $\forall\ A,B \in U,\ \bigl|f(B)-f(A) - df_C(\vec{AB})\bigr|
          \le \frac{M\|\vec{AB}\|_1^2}4$ où $C$ est le milieu de $[A,B]$.

\end{enumerate}
\finenonce{004163}



\finexercice


\exercice{4164, quercia, 2010/03/11}
\enonce{004164}{Application du théorème des fonctions implicites}

On considère la courbe d'équation $e^{x-y} = 1+2x+y$. Donner la tangente à cette
courbe et la position par rapport à la tangente au point $(0,0)$.


\finenonce{004164}



\finexercice


\exercice{4165, quercia, 2010/03/11}
\enonce{004165}{Théorème des fonctions implicites}

\begin{enumerate}
  \item Montrer que l'équation : $x^3+y^3-3xy=1$ définit au voisinage de $0$ une fonction
    implicite : $y=\varphi(x)$ telle que $\varphi(0)=1$.
  \item Donner le DL de $\varphi$ en $0$ à l'ordre $3$.

\end{enumerate}
\finenonce{004165}



\finexercice


\exercice{4166, quercia, 2010/03/11}
\enonce{004166}{Théorème des fonctions implicites, Ensi P 91}

    Montrer que l'égalité $2e^{x+y} + y - x = 0$ définit $y = \varphi(x)$ au voisinage
    de $(1,-1)$. Calculer $\varphi'(1)$ et $\varphi''(1)$.
    

\finenonce{004166}



\finexercice


\exercice{4167, quercia, 2010/03/11}
\enonce{004167}{\'Equation implicite $x\ln x = y\ln y$}

Soit $f(x,y) = x\ln x - y\ln y$ $(x,y > 0)$.

Pour $k \in \R$, on considère la courbe $\mathcal{C}_k$ d'équation $f(x,y) = k$.

\begin{enumerate}
  \item  Suivant la position de $(a,b) \in \mathcal{C}_k$, préciser l'orientation
     de la tangente à $\mathcal{C}_k$ en $(a,b)$.
  \item  Dresser le tableau de variations de $\phi(t) = t\ln t$.
  \item  Dessiner $\mathcal{C}_0$. (\'Etudier en particulier les points
     $(0,1), (1,0)$ et $\Bigl(\frac 1e,\frac 1e\Bigr)$ à l'aide de DL)
  \item  Indiquer l'allure générale des courbes $\mathcal{C}_k$ suivant le signe de $k$.
\end{enumerate}
\finenonce{004167}



\finexercice


\exercice{4168, quercia, 2010/03/11}
\enonce{004168}{Fonction implicite}

Soit $f:\R \to \R$ de classe $\mathcal{C}^1$.

\begin{enumerate}
  \item Montrer que, sous une condition à préciser, l'équation $y-zx=f(z)$ définit
    localement $z$ fonction implicite de $x$ et $y$.
  \item Montrer que l'on a alors : $\frac{\partial z}{\partial x} + z\frac{\partial z}{\partial y} = 0$.
\end{enumerate}
\finenonce{004168}



\finexercice


\exercice{4169, quercia, 2010/03/11}
\enonce{004169}{\'Equation fonction de deux paramètres}

Soit l'équation $(*) \Leftrightarrow x^5 + \lambda x^3 + \mu x^2 -1 = 0$.
Montrer qu'il existe un voisinage, $V$, de $(0,0)$ et
$\varphi : V \to \R$ tels que~:
\begin{align*}&\varphi \text{ est } \mathcal{C}^\infty\\
           &\varphi(0,0) = 1\cr
           &\forall\ (\lambda,\mu)\in V,\ \varphi(\lambda,\mu)
            \text{ est racine simple de } (*).\\ 
\end{align*}
Donner le DL à l'ordre $2$ de $\varphi$ en $(0,0)$.


\finenonce{004169}



\finexercice


\exercice{4170, quercia, 2010/03/11}
\enonce{004170}{Changement de variable singulier, Matexo}

On considère la fonction de $\R^2$ sur lui-même définie par
$f(x,y)=(u,v)$, où
$$u(x,y)=x \sqrt{1+y^2} + y \sqrt{1+x^2}\quad \hbox{et}\quad v(x,y)=(x +
\sqrt{1+x^2})(y+\sqrt{1+y^2}).$$

Calculer sa matrice jacobienne. Est-elle inversible localement ?
Caractériser $f (\R^2)$.

\finenonce{004170}



\finexercice


\exercice{4171, quercia, 2010/03/11}
\enonce{004171}{Longueur d'un arc de courbe}

\def\arc#1{\mathop{#1}\limits^\frown}
Soit $f : {U \subset \R^n} \to {\R^p}$ de classe $\mathcal{C}^1$ dont les dérivées
partielles sont bornées sur~$U$ et $t\in I\  \mapsto M_t$ une courbe paramétrée dans
$U$ de classe $\mathcal{C}^1$. Pour $a,b\in I$ comparer les longueurs des arcs
$\arc{M_aM_b}$ et $\arc{f(M_a)f(M_b)}$.
\finenonce{004171}



\finexercice


\exercice{4172, quercia, 2010/03/11}
\enonce{004172}{Différentielle du déterminant}

Soit $f : {\mathcal{M}_n(\R)} \to  \R, M \mapsto {\det M.}$

Montrer que $f$ est de classe $\mathcal{C}^1$ et que l'on a pour $M,H\in\mathcal{M}_n(\R)$~:
$d f_M(H) = \mathrm{tr}(^t\mathrm{com}(M)H)$.

Application~: soit $M\in\mathcal{M}_n(\R)$ et $P_M(X) = (-1)^nX^n + \dots + a_1X + \det(M)$.
Exprimer $a_1$ en fonction des cofacteurs de~$M$.
\finenonce{004172}



\finexercice


\exercice{4173, quercia, 2010/03/11}
\enonce{004173}{Mise en facteur de $x$ et $y$}

Soit $U$ un ouvert convexe de~$\R^2$ contenant~$(0,0)$
et $f : U \to \R$ une fonction de classe $\mathcal{C}^\infty$ telle que $f(0,0) = 0$.


\begin{enumerate}
  \item Montrer qu'il existe ${g,h} : U \to \R$ de classe $\mathcal{C}^\infty$
    telles que~:
    $$\forall\ (x,y)\in U,\ f(x,y) = xg(x,y) + yh(x,y).$$

  \item Y a-t-il unicité de $g$ et $h$~?
  \item Généraliser au cas où $U$ n'est pas convexe.
\end{enumerate}
\finenonce{004173}



\finexercice


\exercice{4174, quercia, 2010/03/11}
\enonce{004174}{Fonctions convexes}

Soit $U$ un ouvert convexe de~$\R^n$ et $f : U \to \R$. On dit que
$f$ est convexe lorsque~:
$$\forall\ x,y\in U,\ \forall\ t\in{[0,1]},\ f(tx+(1-t)y) \le tf(x) + (1-t)f(y).$$
On dit que $f$ est strictement convexe si l'inégalité précédente est
stricte lorsque $x\ne y$ et $0<t<1$.

\begin{enumerate}
  \item On suppose que $f$ est convexe.
  \begin{enumerate}
    \item Soient $x\in U$, $h\in \R^n$ et $t\in{[0,1]}$ tel que $x-h\in U$ et $x+h\in U$. Montrer~:
$$(1+t)f(x) - tf(x-h) \le f(x+th) \le (1-t)f(x) + tf(x+h).$$
    \item Montrer que $f$ est continue (raisonner sur le cas $n=2$ puis généraliser).
  \end{enumerate}
  \item On suppose que $f$ est de classe $\mathcal{C}^1$.

Montrer que $f$ est convexe si et seulement si pour tous $(x,y)\in U$
on a~: $f(y) \ge f(x) + d f_x(y-x)$. Donner une interprétation
géométrique de cette inégalité lorsque $n=2$.
  \item On suppose que $f$ est de classe $\mathcal{C}^2$.
  \begin{enumerate}
    \item Montrer que $f$ est convexe si et seulement si pour tout~$x\in U$
la forme bilinéaire symétrique $d^2f_x$ est positive.
    \item Si, pour tout $x\in U$, $d^2f_x$ est définie positive, montrer
que $f$ est strictement convexe. Montrer par un exemple que la réciproque est fausse.
  \end{enumerate}
\end{enumerate}
\finenonce{004174}



\finexercice


\exercice{4175, quercia, 2010/03/11}
\enonce{004175}{Les racines d'un polynôme sont des fonctions $\mathcal{C}^\infty$ des coefficients}

Soit $U$ l'ensemble des polynômes à coefficients réels de degré~$n$
et à racines réelles simples.

\begin{enumerate}
  \item Montrer que $U$ est ouvert dans~$\R_n[X]$.
  \item Pour $P\in U$ on note $x_1 < x_2 < \dots < x_n$ les racines de~$P$.
Montrer que l'application $P  \mapsto (x_1,\dots,x_n)$ est de classe $\mathcal{C}^\infty$.
\end{enumerate}
\finenonce{004175}



\finexercice


\exercice{4176, quercia, 2010/03/11}
\enonce{004176}{Non injectivité locale de l'exponentielle}

Soit $f : {\mathcal{M}_n(\C)} \to {\mathcal{M}_n(\C)}, M  \mapsto {\exp(M).}$

\begin{enumerate}
  \item Montrer que $f$ est de classe $\mathcal{C}^1$ sur $\mathcal{M}_n(\C)$ et exprimer,
pour $M,H\in \mathcal{M}_n(\C)$, $d f_M(H)$ sous forme d'une série.


  \item Montrer qu'il existe un voisinage $V$ de $0$ dans $\mathcal{M}_n(\C)$
tel que pour toutes matrices $A,B\in V$ on a~:

$\exp(A) = \exp(B)  \Rightarrow  A=B$.

  \item Trouver une suite $(M_k)$ de matrices de $\mathcal{M}_2(\C)$ distinctes ayant même
exponentielle et convergeant vers une matrice $A$ (donc il n'existe
pas de voisinage de~$A$ sur lequel la restriction de~$f$ est injective).


  \item Donner de même un point de non injectivité locale dans $\mathcal{M}_2(\R)$.

\end{enumerate}
\finenonce{004176}



\finexercice


\exercice{4177, quercia, 2010/03/11}
\enonce{004177}{Caractérisation des isométries}

Soit $E$ un espace vectoriel euclidien et $f : E \to E$ de classe
$\mathcal{C}^1$.

\begin{enumerate}
  \item Montrer que $f$ est une application affine si et seulement
si sa différentielle est constante (c'est-à-dire $d f_x = d f_y$
pour tous $x,y$, égalité dans $\mathcal{L}(E)$).

  \item Soit $X$ un ensemble non vide quelconque et
$\varphi : {X^3}\to \R$ une application vérifiant~:
$$\forall\ x,y,z\in X,\ \varphi(x,y,z) = \varphi(y,x,z) = -\varphi(z,y,x).$$
Montrer que $\varphi=0$ (lemme des tresses).

  \item On suppose $f$ de classe $\mathcal{C}^2$. Montrer que $f$ est une isométrie
de~$E$ pour la distance euclidienne si et seulement si, pour tout
$x\in E$, $d f_x$ est une application orthogonale.
\end{enumerate}
\finenonce{004177}



\finexercice


\exercice{4178, quercia, 2010/03/11}
\enonce{004178}{Différentiabilité de la norme}

Pour chacune des trois normes classiques sur $\R^2$ dire en quels
points elles sont différentiables.
\finenonce{004178}



\finexercice


\exercice{4179, quercia, 2010/03/11}
\enonce{004179}{Difféomorphisme}

Soit $E$ un espace euclidien et $f : E \to E$ de classe $\mathcal{C}^1$, $\alpha > 0$ vérifiant~:
$$\forall\ x\in E,\ \forall\ h\in E,\ (d f_x(h)\mid h) \ge \alpha\|h\|^2.$$

\begin{enumerate}
  \item Montrer pour $x,y\in E$~: $(f(x)-f(y)\mid x-y) \ge \alpha\|x-y\|^2$.
En déduire que $f(E)$ est fermé.
  \item Montrer que $f(E)$ est ouvert puis que $f$ est un $\mathcal{C}^1$-difféomorphisme
de~$E$ sur~$E$.
\end{enumerate}
\finenonce{004179}



\finexercice


\exercice{4180, quercia, 2010/03/11}
\enonce{004180}{Difféomorphisme}

Soit $f:\R \to \R$ de classe $\mathcal{C}^1$ et $k$-lipschitzienne avec $k<1$ et
$\varphi : {\R^2} \to {\R^2}, {(x,y)} \mapsto {(x+f(y),y-f(x)).}$

Montrer que $\varphi$ est un $\mathcal{C}^1$-difféomorphisme de~$\R^2$ sur~$\R^2$.
\finenonce{004180}



\finexercice


\exercice{4181, quercia, 2010/03/11}
\enonce{004181}{Partiellement dérivable $ \Rightarrow $ continue~?}

Soit $U$ un ouvert de~$\R^2$.

\begin{enumerate}
  \item Donner un exemple de fonction $f : U \to \R$ ayant en tout point des
dérivées partielles premières, mais discontinue en au moins un point.


  \item Soit $f : U \to \R$ ayant en tout point des
dérivées partielles premières {\it bornées\/} sur~$U$. Montrer que~$f$
est continue.

\end{enumerate}
\finenonce{004181}



\finexercice


\exercice{4182, quercia, 2010/03/11}
\enonce{004182}{Point non extrémal}

On pose pour $(x,y)\in\R^2$~:
$$f(x,y) = x^2+y^2-2x^2y - \frac{4x^6y^2}{(x^4+y^2)^2}\quad\text{si }(x,y)\ne 0,
  \qquad f(0,0)=0.$$
\begin{enumerate}
  \item Montrer que $f$ est continue sur $\R^2$.

  \item Soit $\theta\in\R$ fixé et $g_\theta(r) = f(r\cos\theta,r\sin\theta)$.
Montrer que $g_\theta$ admet un minimum local strict en~$r=0$.

  \item Calculer $f(x,x^2)$. Conclusion~?

\end{enumerate}
\finenonce{004182}



\finexercice


\exercice{4183, quercia, 2010/03/11}
\enonce{004183}{Contre-exemple au théorème de Leibniz}

On pose~:
$f(x,y) = \begin{cases}
 x            &\text{ si } y\ge0 \text{ et } 0\le x\le \sqrt y~;\cr
 2\sqrt y - x &\text{ si } y\ge0 \text{ et } \sqrt y < x\le 2\sqrt y~;\cr
 0            &\text{ si } y\ge0 \text{ et } 2\sqrt y < x \text{ ou } x\le 0~;\cr
 -f(x,-y)     &\text{ si } y<0.\cr\end{cases}
$
et~: $F(y) =  \int_{x=0}^1 f(x,y)\,d x$.

Faire un dessin, vérifier que $f$ est continue sur~$\R^2$,
calculer $F(y)$ pour $-\frac14 \le y \le \frac 14$, $F'(0)$ et
$ \int_{x=0}^1 \frac{\partial f}{\partial y}(x,0)\,d x$.
\finenonce{004183}



\finexercice


\exercice{4184, quercia, 2010/03/11}
\enonce{004184}{Centrale MP 2000}

Soit ${f} : {\R^n} \to {\R^n}$ de classe $\mathcal{C}^1$ et $c>0$ tels que, pour tous $x,y$, 
$\| f(x)-f(y)\| \ge c \|x-y\|$.

\begin{enumerate}
  \item Montrer que pour tous $x,h$, $\|d f_x(h)\|\ge c\|h\|$.

  \item Montrer que $f$ est un $\mathcal{C}^1$-difféomorphisme sur $\R^n$ (pour la surjectivité on 
    considèrera, si $a\in \R^n$, le minimum de $\|f(x)-a\|^2$).
    
\end{enumerate}
\finenonce{004184}



\finexercice


\exercice{4185, quercia, 2010/03/11}
\enonce{004185}{Centrale MP 2000}
Soit $\Omega$ un ouvert borné de $\R^2$ et ${u} : {\overline{\Omega}} \to {\R}$ continue sur 
$\overline{\Omega}$ et $\mathcal{C}^2$ sur $\Omega$.

\begin{enumerate}
  \item On suppose que $\Delta u>0$. Montrer que $\max\limits_{(x,y)\in \overline{\Omega}}u(x,y)=
    \max\limits_{(x,y)\in \overline{\Omega}\setminus \Omega}u(x,y)$.
    
  \item Même question en supposant seulement $\Delta u \ge 0$.
    
  \item Soit $0<r_1<r_2$, $A=\{(x,y)\in \R^2 \,|\, r_1^2<x^2+y^2<r_2^2\}$. On suppose que $u$ est  
    continue sur $\overline{A}$, $\mathcal{C}^2$ sur $A$ et que $\Delta u \ge 0$ sur $A$. On pose pose
    $M(r)=\max\limits_{x^2+y^2=r^2} (u(x,y))$.

    Montrer que, pour tout $r_1\le r\le r_2$, 
    $M(r)\le \frac{M(r_1)\ln(r_2/r)+M(r_2)\ln (r/r_1)}{\ln (r_2/r_1)}\cdotp$

    Indication~: la fonction $v$ : $(x,y) \mapsto\ln(x^2+y^2)$ vérifie $\Delta v = 0$.
    
\end{enumerate}
\finenonce{004185}



\finexercice


\exercice{4186, quercia, 2010/03/11}
\enonce{004186}{Mines MP 2001}

Soit une fonction $f$ de classe $\mathcal{C}^2$ sur le disque unité du plan, telle que son laplacien
$\frac{\partial^2 f}{\partial x^2} + \frac{\partial^2 f}{\partial y^2}$ soit nul.

\begin{enumerate}
  \item Montrer $ \int_{\theta=0}^{2\pi}f(r\cos\theta,r\sin\theta)\,d\theta$
    ne dépend pas de $r\in[0,1]$.
    

  \item Calculer alors $\iint_{D_r} f(x,y)\,d xd y$ $D_r$ étant le disque fermé
    de centre $0$ et de rayon $r$.
    
 
\end{enumerate}
\finenonce{004186}



\finexercice


\exercice{4187, quercia, 2010/03/11}
\enonce{004187}{Mines MP 2001}

Soient $f$ et $g$ deux fonctions de classe $\mathcal{C}^1$ sur $\R$ à valeurs
dans $\R$ vérifiant~: $\forall\ x\in\R,\ f'(x)\ge 1$ et $|g'(x)| < 1$.
Soit $\varphi$ définie sur $\R^2$ par $\varphi(x,y) = (f(x)+g(y), f(y) + g(x))$.

\begin{enumerate}
  \item Montrer que $\varphi$ est un $\mathcal{C}^1$-difféomorphisme de $\R^2$ sur $\varphi(\R^2)$.
    

  \item On suppose qu'il existe $k\in{]0,1[}$ tel que $\forall\ x\in\R$, $|g'(x)| < k$~;
    montrer que $\varphi(\R^2) = \R^2$.
    
\end{enumerate}
\finenonce{004187}



\finexercice


\exercice{4188, quercia, 2010/03/11}
\enonce{004188}{ENS MP 2002}

Soit $f:{\R^2}\to\R$ de classe $\mathcal{C}^2$ telle que 
$|f(x)|/\|x\| \to +\infty$ lorsque $\|x\|\to\infty$.
Prouver que $\nabla f$ est surjective sur $\R^2$.

\finenonce{004188}



\finexercice


\exercice{4189, quercia, 2010/03/11}
\enonce{004189}{ENS MP 2002}

Soit $n$ un entier $>0$, $\|\ \|$ la norme euclidienne 
sur $\R^n$ et $f : {\R^n} \to \R$ de classe $\mathcal{C}^2$. On suppose
que $f(x)/\|x\| \to +\infty$ lorsque $\|x\|\to\infty$, et 
qu'en tout point la matrice hessienne de $f$ est 
définie positive.

On pose $g(y)= \sup\{ (x\mid y) -f(x),\ x\in \R^n\}$.
Étudier les propriétés  de $g$. 
\finenonce{004189}



\finexercice


\exercice{4190, quercia, 2010/03/11}
\enonce{004190}{$\int \varphi\circ f$, X MP$^*$ 2004}
Soit~$E$ l'ensemble des fonctions continues de~$[0,1]$ dans~$\R$.
On y définit une norme par~: $\|f\|=\sqrt{ \int_{t=0}^1 f^2(t)\,d t}$.
Soit $\varphi : \R \to \R$ de classe $\mathcal{C}^2$ telle que $\varphi''$ est bornée.
Pour~$f\in E$ on pose $T(f) =  \int_{t=0}^1\varphi(f(t))\,d t$.

\begin{enumerate}
  \item Montrer que l'application ainsi définie~$T : E \to \R$ est continue.

  \item Montrer que~$T$ est différentiable en tout point.

\end{enumerate}
\finenonce{004190}



\finexercice


\exercice{5556, rouget, 2010/07/15}
\enonce{005556}{***T}
Soit $\begin{array}[t]{cccc}f~:&\Rr^2&\longrightarrow&\Rr\\
 &(x,y)&\mapsto&\left\{
\begin{array}{l}
0\;\text{si}\;y=0\\
y^2\sin\left(\frac{x}{y}\right)\;\text{si}\;y\neq0
\end{array}
\right.
\end{array}
$.
\begin{enumerate}
 \item  Etudier la continuité de $f$.

 \item  Etudier l'existence et la valeur éventuelle de dérivées partielles d'ordre 1 et 2.
On montrera en particulier que $\frac{\partial^2f}{\partial x\partial y}$ et $\frac{\partial^2f}{\partial y\partial x}$  sont définies en $(0,0)$ mais n'ont pas la même valeur.
\end{enumerate}
\finenonce{005556}


\finexercice


\exercice{5557, rouget, 2010/07/15}
\enonce{005557}{***}
 Le laplacien d'une application $g$ de $\Rr^2$ dans $\Rr$, de classe $C^2$ sur $\Rr^2$  est $\Delta g =\frac{\partial^2g}{\partial x^2}+\frac{\partial^2g}{\partial y^2}$.

 
Déterminer une fontion de classe $C^2$ sur un intervalle $I$ de $\Rr$ à préciser à valeurs dans $\Rr$ telle que la fonction

\begin{center}
$g(x,y) =f\left(\frac{\cos2x}{\ch2y}\right)$
\end{center}
soit non constante et ait un laplacien nul sur un sous-ensemble de $\Rr^2$ le plus grand possible (une fonction de Laplacien nul est dite harmonique).

\finenonce{005557}


\finexercice


\exercice{5893, rouget, 2010/10/16}
\enonce{005893}{*}
Soit $f$ une application de $\Rr^n$ dans $\Rr$ de classe $C^1$. On dit que $f$ est positivement homogène de degré $r$ ($r$ réel donné) si et seulement si $\forall\lambda\in]0,+\infty[$, $\forall x\in\Rr^n$, $f(\lambda x) =\lambda^rf(x)$.

Montrer pour une telle fonction l'identité d'\textsc{Euler} :

\begin{center}
$\forall x = (x_1,...,x_n)\in \Rr^n$ $\sum_{i=1}^{n}x_i \frac{\partial f}{\partial x_i}(x) = rf(x)$.
\end{center}
\finenonce{005893}


\finexercice

\section{223.05 Différentielle de fonctions composées}

\section{223.06 Différentielle seconde}


\exercice{2635, debievre, 2009/05/19}
\enonce{002635}{}
Calculer les diff\'erentielles suivantes, sans calculer des d\'eriv\'ees
 partielles, en utilisant les 
propri\'et\'es des diff\'erentielles de sommes, produits et compos\'ees:\\
\[
\mathrm{(a)}\ \mathrm{d} \left(\ln(xy)\right)
\quad \mathrm{(b)}\ \mathrm{d} \left(xyz(1+\sinh(yz))\right)
\quad \mathrm{(c)}\ \mathrm{d} \left(\sin(x^2y) \mathrm e^{x-y}\right)
\]
\finenonce{002635}


\finexercice


\exercice{2636, debievre, 2009/05/19}
\enonce{002636}{}
\begin{enumerate}
 \item Y a-t-il une fonction $g\colon \R^2\to\R$ telle que
\[
 \mathrm{d} g = x^2y^2 \mathrm{d} x + x^3 y\mathrm{d} y?
\]
 \item Trouver les fonctions $b\colon\R^2\to\R$ telles qu'il existe 
$g\colon \R^2\to\R$ satisfaisant \`a la condition
\[
\mathrm{d} g = x^2y^2 \mathrm{d} x + b(x,y)\mathrm{d} y.
\]
\'Etant donn\'ee alors la fonction $b$, d\'eterminer toutes les fonctions $g$
correspondantes.
\end{enumerate}
\finenonce{002636}


\finexercice


\exercice{2637, debievre, 2009/05/19}
\enonce{002637}{}
Soit $g\colon\R_{>0}\times \R_{>0}\to \R$ une fonction
de classe $C^1$
telle que $g(1,1)=3$ et dont la diff\'erentielle vaille
\begin{equation}\label{eq}
\mathrm{d} g= (2xy+y^2)\mathrm{d} x + (x^2+2xy) \mathrm{d} y.
\end{equation} 
Soit 
\[
h\colon \R_{>0}\times\R_{>0} \longrightarrow \R_{>0}\times\R_{>0}
\]
l'application de classe $C^1$ d\'efinie par 
\[
h(x,y)= (u(x,y), v(x,y))=(x^2y, xy^2)\in \R_{>0}\times \R_{>0}. 
\]
\begin{enumerate}
 \item  Calculer $\mathrm{d} u+\mathrm{d} v$. 
 \item D\'eterminer $g$  \`a partir du calcul 
pr\'ec\'edent et (\ref{eq}), et sans autre calcul.
 \item  Montrer que $h$ est une bijection. (On pourra calculer explicitement $h^{-1}$.)
 \item D\'eterminer explicitement $\mathrm{d} (g\circ h^{-1})$.
 \item  Calculer les matrixes jacobiennes
$J_h(x,y)$ et $J_{h^{-1}}(u,v)$ et v\'erifier par un calcul direct que
\[
J_h(x,y)J_{h^{-1}}(h(x,y))=I_2,
\]
 o\`u $I_2$ est la matrice identit\'e d'ordre 2. 
\end{enumerate}
\finenonce{002637}


\finexercice


\exercice{2638, debievre, 2009/05/19}
\enonce{002638}{}
Calculer les matrices hessiennes des fonctions $f$ d\'efinies par les
 expressions suivantes sur leur domaine de d\'efinition naturel:
\[
\sin(xyz), \quad \sin^2(y/x).
\]
\finenonce{002638}


\finexercice


\exercice{2639, debievre, 2009/05/19}
\enonce{002639}{}
Soit $f\colon \R^2\setminus\{(0,0)\}\longrightarrow \R$ 
une fonction
de classe $C^2$ et soient $r$ et $\theta$ les coordonn\'ees
polaires standard dans le plan de telle sorte que
l'association
\[
]0,+\infty[\times  ]0,2\pi[\longrightarrow  \R^2\setminus\{(0,0)\},
\quad
(r,\theta)\longmapsto (x,y)=(r\cos\theta, r\sin\theta),
\]
soit un changement de variables.
Soit $F$ la fonction d\'efinie par 
\[
F(r,\theta)=f(r\cos\theta, r\sin\theta).
\]
C'est \lq\lq l'expression de $f$ 
en coordonn\'ees polaires\rq\rq. 
Montrer que
\begin{equation}
\frac{\partial^2 f}{\partial x^2}(x,y)+\frac{\partial^2 f}{\partial y^2}(x,y)=
\frac{\partial^2F}{\partial r^2}(r,\theta)+\frac{1}{r}\frac{\partial F}{\partial r}(r,\theta)+\frac{1}{r^2}\frac{\partial^2 F}{\partial \theta^2}(r,\theta).
\label{laplace}
\end{equation}
Cette formule
calcule ``le Laplacien en coordonn\'ees polaires.'' 
L'exercice ne d\'epend pas de la connaissance
du Laplacien cependant.
\finenonce{002639}


\finexercice


\exercice{2640, debievre, 2009/05/19}
\enonce{002640}{}
Les variables \'etant not\'ees $x$ et $t$,
trouver la solution g\'en\'erale 
$f\colon\R^2\to\R$
de ``l'\'equation des ondes'', \`a savoir
\begin{equation}\label{eq:ondes}
\frac{\partial^2 f}{\partial x^2}-\frac{\partial^2 f}{\partial t^2}=0.
\end{equation}
Trouver ensuite la solution unique de l'\'equation des ondes qui
satisfait aux conditions initiales
\begin{equation}
f(x,0)=\sin x, \ \frac{\partial f}{\partial t}(x,0)=-\cos x.
\label{in}
\end{equation}

\finenonce{002640}


\finexercice

\exercice{5889, rouget, 2010/10/16}
\enonce{005889}{*** I}
Soit $f(x,y) =\left\{
\begin{array}{l}
y^2\sin\left( \frac{x}{y}\right)\;\text{si}\;y\neq0\\
\rule{0mm}{5mm}0\;\text{si}\;y=0
\end{array}
\right.$.

Déterminer le plus grand sous-ensemble de $\Rr^2$ sur lequel $f$ est de classe $C^1$. Vérifier que $ \frac{\partial^2f}{\partial x\partial y}(0,0)$ et $ \frac{\partial^2f}{\partial y\partial x}(0,0)$ existent et sont différents.
\finenonce{005889}


\finexercice


\exercice{5904, rouget, 2010/10/16}
\enonce{005904}{***}
Trouver une application non constante $f~:~]-1,1[\rightarrow\Rr$ de classe $C^2$ telle que l'application $g$ définie sur $\Rr^2$ par
$g(x,y) = f\left( \frac{\cos(2x)}{\ch(2y)}\right)$ ait un laplacien nul sur un ensemble à préciser. (On rappelle que le laplacien de $g$ est $\Delta g = \frac{\partial^2g}{\partial x^2}+ \frac{\partial^2g}{\partial y^2}$. Une fonction de laplacien nul est dite harmonique.)
\finenonce{005904}


\finexercice


\exercice{5905, rouget, 2010/10/16}
\enonce{005905}{*** I}
Soit $f~:~\Rr^2\rightarrow\Rr^2$ de classe $C^2$ dont la différentielle en tout point est une rotation. Montrer que $f$ est une rotation affine.
\finenonce{005905}


\finexercice

\section{223.07 Extremums locaux}


\exercice{2641, debievre, 2009/05/19}
\enonce{002641}{}
Pour chacune des fonctions suivantes \'etudier la nature du point
critique donn\'e :
\begin{enumerate}
\item $f(x,y)=x^2-xy+y^2$ au point critique $(0,0)$ ;
\item $f(x,y)=x^2+2xy+y^2+6$ au point critique $(0,0)$ ;
\item $f(x,y)=x^3+2xy^2-y^4+x^2+3xy+y^2+10$ au point critique
$(0,0)$.
\end{enumerate}
\finenonce{002641}


\finexercice


\exercice{2642, debievre, 2009/05/19}
\enonce{002642}{}
Trouver les points critiques de la  fonction $f$  suivante et
d\'eterminer si ce sont des minima locaux, des maxima locaux ou
des points selle.
\[
f(x,y)=\sin x+y^2-2y+1
\] 
\finenonce{002642}


\finexercice


\exercice{2643, debievre, 2009/05/19}
\enonce{002643}{}
\begin{enumerate}
 \item  Soit 
$f$ une fonction r\'eelle d'une variable r\'eelle
de classe ${\rm C}^2$ dans un voisinage de $0\in \mathbb R$ 
telle que $f(0)=0$ et
$f^\prime(0)\not=0$. Montrer que la fonction 
r\'eelle $F$ des deux variables $x$ et $y$
d\'efinie dans un voisinage de $(0,0)$
par 
$F(x,y)= f(x) f(y)$ n'a pas d'extremum relatif en $(0,0)$.
Est-ce que  le point $(0,0)$ est quand m\^eme critique?
Si oui caract\'eriser sa nature.
 \item D\'eterminer les points critiques, puis les minima et les maxima locaux de
\[
f(x,y)=\sin(2\pi x)\sin(2\pi y).
\]
Remarque: en utilisant la 
p\'eriodicit\'e de la fonction, on peut limiter le nombre de cas \`a \'etudier.
\end{enumerate}
\finenonce{002643}


\finexercice


\exercice{2644, debievre, 2009/05/19}
\enonce{002644}{}
 D\'eterminer l'\'equation du plan tangent \`a la surface de niveau
\[\sin(\pi xy)+\sin(\pi yz) =1,\]
au point de coordonn\'ees $(1,\frac{1}{6}, 1)$. Identifier,
 en ce point, un vecteur perpendiculaire \`a la surface. 
Votre r\'esultat est-il compatible avec la figure ci-dessous ? Expliquer. 

{\center \includegraphics[height=8cm, keepaspectratio]{../images/img002644-1}}
\finenonce{002644}


\finexercice



\exercice{2645, debievre, 2009/05/19}
\enonce{002645}{}
Soit $\cal C$ la courbe plane 
d'\'equation $f(x,y)=ye^{x}+e^y\sin(2x)=0$.
\begin{enumerate}
 \item  Appliquer le th\'eor\`eme des fonctions implicites
\`a la courbe ${\cal C}$ au point $(0,0)$.
 \item D\'eterminer la limite de $y/x$ quand $(x,y)$ tend  le long la courbe ${\cal C}$
vers
$(0,0)$.
\end{enumerate}
\finenonce{002645}


\finexercice


\exercice{2646, debievre, 2009/05/19}
\enonce{002646}{}
\begin{enumerate}
 \item  D\'eterminer les points stationnaires de 
la fonction $f$ de deux variables d\'efinie par
$f(x,y)= x\hskip 2pt (x+1)^2 - y^2$ et pr\'eciser
la nature de chacun d'eux.
 \item Tracer la courbe  
constitu\'ee des points tels que $f(x,y)=0$ et $x \geq 0$.
(\emph{Indication}: \'Etudier  la fonction 
$x\mapsto \sqrt{x} \hskip 2pt (x+1)$ pour $x\geq 0$). 
 \item  Montrer que 
le point $(-1,0)$ est un point isol\'e de la partie
\[ {\cal C}=\{(x,y); f(x,y)=0\}
\]
du plan, c'est-\`a-dire, le point $(-1,0)$ appartient \`a cette partie
et il existe un nombre r\'eel 
$\varepsilon >0$ tel que 
$D_{\varepsilon} \cap {\cal C} =\{(-1,0)\}$ o\`u
$D_{\varepsilon}$ est le disque ouvert centr\'e en $(-1,0)$ et de rayon $\varepsilon$.
 \item \'Enoncer le th\'eor\`eme des fonctions implicites.
 \item  Montrer que, quel que soit le point 
$(x_0,y_0)$ de ${\cal C}$ distinct de $(-1,0)$,
au moins une des deux alternatives (i) ou (ii) ci-dessous est v\'erifi\'ee:
\begin{enumerate}
\item[(i)]
Il existe une fonction  $h$ de classe $C^1$ de la variable $x$ 
 d\'efinie dans un intervalle ouvert 
appropri\'e telle que $h(x_0)=y_0$ et telle que,
pour qu'au voisinage de $(x_0,y_0)$
les coordonn\'ees $x$ et $y$ du point $(x,y)$
satisfassent \`a l'\'equation
$f(x,y)=0$ il faut et il suffit que
$y=h(x)$.
\item[(ii)]
Il existe une fonction  $k$ de classe $C^1$ de la variable $y$ 
 d\'efinie dans un intervalle ouvert 
appropri\'e telle que $h(y_0)=x_0$ et telle que,
pour qu'au voisinage de $(x_0,y_0)$
les coordonn\'ees $x$ et $y$ du point $(x,y)$
satisfassent \`a l'\'equation
$f(x,y)=0$ il faut et il suffit que
$x=k(y)$.
\end{enumerate}

\end{enumerate}
\finenonce{002646}


\finexercice


\exercice{2654, debievre, 2009/05/19}
\enonce{002654}{}
On consid\`ere la fonction 
$$f(x,y)=(1+2\cos^2(\pi x))(1-\exp(-y^2))+\sin(\pi x).$$
Son graphe est reproduit dans la figure ci-dessous.

\centerline{\includegraphics[height=8cm, keepaspectratio]{../images/img002654-1}}

\begin{enumerate}
 \item  Trouver tous les points critiques de $f$ et d\'eterminer leur nature. Vos r\'esultats sont-ils compatibles avec le graphe de la fonction, reproduit ci-dessus? 
 \item D\'eterminer l'\'equation du plan tangent au graphe de $f$ au point de coordonn\'ees $(1, 1, f(1,1))$. Tracer la droite d'intersection de ce plan avec le plan $xOy$.
\end{enumerate}
\finenonce{002654}



\finexercice



\exercice{4191, quercia, 2010/03/11}
\enonce{004191}{\'Etude de points critiques}

Chercher les extrémums des fonctions $f(x,y)$ suivantes :


\begin{enumerate}
  \item $3xy - x^3 - y^3$
    

  \item $-2(x-y)^2+x^4+y^4$
    

  \item $x^2y^2(1+3x+2y)$
    

  \item $2x+y-x^4-y^4$
    

  \item $\frac {xy}{(x+y)(1+x)(1+y)},\ x,y> 0$
    

  \item $xe^y + ye^x$
    

  \item $x(\ln^2x + y^2)$, $x>0$
    

  \item $\sqrt{x^2+(1-y)^2} + \sqrt{y^2+(1-x)^2}$
    

  \item $MA+MB-MO$,\quad $O = \text{mil}(A,B)$
    

\end{enumerate}
\finenonce{004191}



\finexercice


\exercice{4192, quercia, 2010/03/11}
\enonce{004192}{Distances aux sommets d'un triangle}

Soit $A \in \R^p$ fixé et $f : {\R^p} \to \R, M \mapsto {AM^2}$
$g : {\R^p} \to \R, M \mapsto {AM}$
(distance euclidienne)

\begin{enumerate}
  \item Calculer les gradients de $f$ et $g$ en un point $M$.
  \item Soient $A,B,C$ trois points non alignés du plan. Trouver les points $M$
    du plan réalisant le minimum de :
  \begin{enumerate}
\item $MA^2 + MB^2 + MC^2$.     
\item $MA   + MB   + MC  $.     
\item $MA\times MB\times MC$.   
  \end{enumerate}
\end{enumerate}
\finenonce{004192}



\finexercice


\exercice{4193, quercia, 2010/03/11}
\enonce{004193}{Aire d'un triangle}

Soit $ABC$ un triangle de cotés $a,b,c$.

\begin{enumerate}
  \item Calculer l'aire, $S$, de $ABC$ en fonction de $a,b,c$.
    
  \item Montrer que $\frac{S}{a^2+b^2+c^2}$ est maximal lorsque $ABC$ est équilatéral.
    

\end{enumerate}
\finenonce{004193}



\finexercice


\exercice{4194, quercia, 2010/03/11}
\enonce{004194}{Centrale MP 2000}

On considère un vrai triangle $ABC$ et $f$ la fonction définie par~:
$f(M) = d(M,AB)\times d(M,AC)\times d(M,BC)$.
Montrer que $f$ admet un maximum à l'intérieur du triangle~$ABC$,
et caractériser géométriquement le point~$M_0$ où~$f$ est maximale.

\finenonce{004194}



\finexercice


\exercice{4195, quercia, 2010/03/11}
\enonce{004195}{Loi de réfraction}

Soient dans $\R^2$ : $A=(0,a)$, $B=(b,-c)$ et $M=(x,0)$ ($a,b,c > 0$).
Un rayon lumineux parcourt la ligne brisée $AMB$ à la vitesse $v_1$ de $A$ à $M$
et $v_2$ de $M$ à $B$.
On note $\alpha_1 = \overline{(\vec j,\vec{MA})}$
        $\alpha_2 = \overline{(-\vec j,\vec{MB})}$.

\begin{enumerate}
  \item  Faire une figure.
  \item Montrer que le temps de parcours est minimal lorsque
    $\frac {\sin\alpha_1}{v_1} = \frac {\sin\alpha_2}{v_2}$.
\end{enumerate}
\finenonce{004195}



\finexercice


\exercice{4196, quercia, 2010/03/11}
\enonce{004196}{Centrale MP 2001}

Soit $f$ une forme linéaire sur $E$ espace euclidien et $g(x) = f(x)e^{-\|x\|^2}$.
Montrer que $g$ admet un minimum et un maximum.

\finenonce{004196}



\finexercice


\exercice{4197, quercia, 2010/03/11}
\enonce{004197}{Centrale MP 2001}

$D_1$, $D_2$, $D_3$ sont trois droites d'un plan portant les côtés
d'un triangle équilatéral de côté~$a$.
On pose $$\varphi : {D_1\times D_2\times D_3} \to \R, {(M,N,P)} \mapsto {MN+NP+PM.}$$
Déterminer $\min\varphi$ et les triplets $(M,N,P)$ où ce minimum est atteint.

\finenonce{004197}



\finexercice


\exercice{4198, quercia, 2010/03/11}
\enonce{004198}{Centrale MP 2006}
$E$ désigne l'espace affine euclidien classique. $D_1$, $D_2$, $D_3$
sont trois droites deux à deux non parallèles.
\smallskip
Soit $f : {D_1\times D_2\times D_3} \to \R, {(M_1,M_2,M_3)} \mapsto
{\|\vec{M_1M_2}\|^2 + \|\vec{M_2M_3}\|^2 + \|\vec{M_3M_1}\|^2.}$

\begin{enumerate}
  \item Montrer que $f$ admet un minimum atteint pour un unique triplet.
    

  \item Dans le cas où $D_1$, $D_2$, $D_3$ sont coplanaires et délimitent
    un triangle équilatéral, trouver ce triplet.
    
\end{enumerate}
\finenonce{004198}



\finexercice


\exercice{4199, quercia, 2010/03/11}
\enonce{004199}{Plus court chemin, ENS Ulm-Lyon-Cachan MP$^*$ 2005}

D{\'e}terminer le plus court chemin entre les p{\^o}les nord et sud d'une sph{\`e}re en dimension $3$.

\finenonce{004199}



\finexercice


\exercice{4200, quercia, 2010/03/11}
\enonce{004200}{Extremums liés, ENS Ulm-Lyon-Cachan MP$^*$ 2005}
Soit $B$ la boule unit{\'e} de $\R^{n}$, $f$ de classe $\mathcal{C}^{1}$ sur $B$ et $x\in B$ tel que $f(x)=\max\{f(y),\, y\in B\}$.

Montrer que $\nabla f(x)=\lambda x$ avec $\lambda \ge 0$.
 
\finenonce{004200}



\finexercice


\exercice{5558, rouget, 2010/07/15}
\enonce{005558}{**T}
Trouver les extrema locaux de 

\begin{enumerate}
 \item  $\begin{array}[t]{cccc}
f~:&\Rr^2&\rightarrow&\Rr\\
 &(x,y)&\mapsto&x^2+xy+y^2+2x+3y
\end{array}$ 
 \item  $\begin{array}[t]{cccc}
f~:&\Rr^2&\rightarrow&\Rr\\
 &(x,y)&\mapsto&x^4+y^4-4xy
\end{array}$ 
\end{enumerate}
\finenonce{005558}


\finexercice


\exercice{5559, rouget, 2010/07/15}
\enonce{005559}{***}
Maximum du produit des distances aux cotés d'un triangle $ABC$ du plan d'un point $M$ intérieur à ce triangle (on admettra que ce maximum existe).
\finenonce{005559}


\finexercice


\exercice{5560, rouget, 2010/07/15}
\enonce{005560}{**}
Soit $a$ un réel strictement positif donné. Trouver le minimum de $f(x,y)=\sqrt{x^2+(y-a)^2}+\sqrt{y^2+(x-a)^2}$.

\finenonce{005560}


\finexercice

\exercice{5902, rouget, 2010/10/16}
\enonce{005902}{***}
Maximum du produit des distances d'un point $M$ intérieur à un triangle $ABC$ aux cotés de ce triangle.
\finenonce{005902}


\finexercice


\exercice{5903, rouget, 2010/10/16}
\enonce{005903}{*}
Minimum de $f(x,y) =\sqrt{x^2+(y-a)^2}+\sqrt{(x-a)^2+y^2}$, $a$ réel donné.
\finenonce{005903}


\finexercice

\section{223.08 Fonctions implicites}


\exercice{2653, debievre, 2009/05/19}
\enonce{002653}{}
 On consid\`ere la courbe $\mathcal C$ d'\'equation
$$y^2(x^2+1)+x^2(y^2+1)=1.$$
\begin{enumerate}
 \item  Montrer qu'il existe un unique $b>0$ tel que le point de coordonn\'ees $(1/2,b)$ se trouve sur $\mathcal C$. D\'eterminer $b$, puis d\'eterminer l'\'equation de la droite tangente  \`a $\mathcal C$, passant par $(1/2,b)$. 
 \item Trouver l'unique fonction $\varphi:x\in ]-1,1[\to\varphi(x)\in\R^+$ telle que $(x,\varphi(x))\in\mathcal C$ pour tout $x\in]-1,1[$. Montrer que $\varphi(-x)=\varphi(x)$ et que $\varphi$ est d\'ecroissante sur $[0,1[$. Tracer $\mathcal C$. 
 \item \'Enoncer le th\'eor\`eme des fonctions implicites et montrer qu'il existe exactement deux points de la courbe $\mathcal C$ o\`u le th\'eor\`eme des fonctions implicites ne  s'applique pas pour \'ecrire, au voisinage de chacun de ces deux points,
$y$ comme fonction de $x$.  
\end{enumerate}
\finenonce{002653}



\finexercice


\exercice{5890, rouget, 2010/10/16}
\enonce{005890}{**}
Montrer que $\begin{array}[t]{cccc}
\varphi~:&\Rr^2&\rightarrow&\Rr^2\\
 &(x,y)&\mapsto&(e^x-e^y , x+y)
 \end{array}$ est un $C^1$-difféomorphisme de $\Rr^2$ sur lui-même.
\finenonce{005890}


\finexercice


\exercice{5891, rouget, 2010/10/16}
\enonce{005891}{***}
Soit $n\in\Nn$. Montrer que l'équation $y^{2n+1}+ y - x = 0$ définit implicitement une fonction $\varphi$ sur $\Rr$ telle que : $(\forall(x,y)\in\Rr^2),\;[y^{2n+1}+ y - x = 0\Leftrightarrow y =\varphi(x)]$.

Montrer que $\varphi$ est de classe $C^\infty$ sur $\Rr$ et calculer $\int_{0}^{2}\varphi(t)\;dt$.
\finenonce{005891}


\finexercice


\exercice{5892, rouget, 2010/10/16}
\enonce{005892}{***}
Donner un développement limité à l'ordre $3$ en $0$ de la fonction implicitement définie sur un voisinage de $0$ par l'égalité $e^{x+y}+y-1 = 0$.
\finenonce{005892}


\finexercice

\section{223.99 Autre}


\exercice{1861, gourio, 2001/09/01}

\enonce{001861}{}
Soit $ f:{\Rr}^{2}\rightarrow {\Rr} $ admettant des d\'{e}riv\'{e}es
partielles continues en $ 0$ et telle que :
$$\forall a\in {\Rr}^{2}-\{0\},\forall t>0,f(ta)=tf(a). $$
Montrer que $ f$ est lin\'{e}aire.
\finenonce{001861}



\finexercice



\exercice{1862, gourio, 2001/09/01}

\enonce{001862}{}
Soit $ f:{\Rr}^{2}\rightarrow {\Rr}$ une application $ C^{1 } $ sur un
ouvert convexe $ O$ telle que :
$$\forall a\in O,\frac{\partial f}{\partial x_{1}}(a)=\frac{\partial f}{
\partial x_{2}}(a)=0 .$$
Montrer que $ f$ est constante sur $ O.$
\finenonce{001862}



\finexercice



\exercice{1863, drutu, 2003/10/01}

\enonce{001863}{}
Soit $f:\R^n\to \R$ une application diff\'erentiable. Montrez que si $\| \nabla f(x)\| 
\leq M,\; \forall x\in \R^n$, alors 
$$
|f(x)-f(y)|\leq M\| x-y\|,\; \forall x,y \in \R^n \; .
$$
\finenonce{001863}



\finexercice



\exercice{1864, gineste, 2001/11/01}

\enonce{001864}{In\'egalit\'e de Cauchy-Schwarz}
\begin{enumerate}
\item Montrer que : $ \forall x_1,x_2, \cdots ,
 x_n \in \R \; \; (x_1 + x_2 + \cdots + x_n)^2 \leq n(x_1^2 + x_2^2 + \cdots + x_n^2) $
\item D\'eterminer : $m = \text{Inf}
\{ (\sum_{i=1}^n x_i)(\sum_{i=1}^n 1/x_i) \mbox{ tels que }  x_1,x_2, \cdots , x_n > 0 \} $
\item D\'eterminer : $M= \sup \{ |x + 2y + 3z + 4t| \mbox{ tels que } (x,y,z,t) \in \R^4 \; , x^2+y^2+z^2+t^2 \leq 1 \} $
\end{enumerate}

\finenonce{001864}



\finexercice



\exercice{2648, debievre, 2009/05/19}
\enonce{002648}{}
On consid\`ere la fonction $f(x,y)=2x^2 + y^2$, $(x,y)\in\R^2$. 
\begin{enumerate}
 \item Tracer les lignes de niveau $f(x,y)=2, f(x,y)=4$. 
 \item Tracer le graphe de la fonction $f$. Expliquer votre dessin en quelques phrases, en identifiant notamment les intersections du graphe de $f$ avec les plans parall\`eles aux trois plans des coordonn\'ees. 
\end{enumerate}
\finenonce{002648}



\finexercice


\exercice{2655, debievre, 2009/05/19}
\addcommand{\usepackage{graphics}}
\enonce{002655}{} 
On consid\`ere les quatre surfaces $\Sigma_1, \Sigma_2, \Sigma_3, \Sigma_4$, d\'efinies par les \'equations suivantes:
\begin{eqnarray*}
z^2-\exp(2x^2+y^2)=0&(\Sigma_1)\\
z=x^2+3y^2+4&(\Sigma_2)\\
z-(x-2y)^2-4=0&(\Sigma_3)\\
\exp(x^2+y^2)+\exp(y^2+z^2)=3&(\Sigma_4)
\end{eqnarray*}
Les quatre surfaces sont trac\'ees dans les parties A, B, C et D de la figure sur la page suivante. Indiquer quelle surface correspond \`a quelle partie de la figure. On justifiera tr\`es bri\`evement ses r\'eponses.


\centerline{
%\hskip-8cm
\begin{tabular}{cc}
\includegraphics[height=8cm, keepaspectratio]{../images/img002655-3}&
\includegraphics[height=8cm, keepaspectratio]{../images/img002655-2}\\
A&B\\
\includegraphics[height=8cm, keepaspectratio]{../images/img002655-4}&
\includegraphics[height=8cm, keepaspectratio]{../images/img002655-1}\\
C&D\end{tabular}
}
\finenonce{002655}



\finexercice


\exercice{2684, matexo1, 2002/02/01}
\enonce{002684}{}
Soit la fonction $f$ d\'efinie sur $\R_+^n$
par $f(x_1,\dots,x_n) = \sqrt[n]{x_1\ldots x_n}$. Chercher ses
extrema quand $x_1+\cdots+x_n = n$. En d\'eduire que la moyenne
g\'eom\'etrique de $n$ nombres positifs est plus petite que leur moyenne
arithm\'etique.
\finenonce{002684}
\finexercice


\exercice{5561, rouget, 2010/07/15}
\enonce{005561}{**}
Trouver toutes les applications $\varphi$ de $\Rr$ dans $\Rr$ de classe $C^2$ telle que l'application $f$ de $U=\{(x,y)\in\Rr^2/\;x\neq0\}$ dans $\Rr$ qui à $(x,y)$ associe $\varphi\left(\frac{y}{x}\right)$ vérifie :

\begin{center}
$\frac{\partial^2f}{\partial x^2}-\frac{\partial^2f}{\partial y^2}=\frac{y}{x^3}$.
\end{center}
\finenonce{005561}


\finexercice


\exercice{5562, rouget, 2010/07/15}
\enonce{005562}{**}
Trouver toutes les applications $f$ de $\Rr^2$ dans $\Rr$ vérifiant

\begin{enumerate}
 \item   $2\frac{\partial f}{\partial x}-\frac{\partial f}{\partial y}=0$   (en utilisant le changement de variables $u=x+y$ et $v=x+2y$)

 \item  $x\frac{\partial f}{\partial x}+ y\frac{\partial f}{\partial y}=\sqrt{x^2+y^2}$ sur $D=\{(x,y)\in\Rr^2/\;x>0\}$ (en passant en polaires).
\end{enumerate}
\finenonce{005562}


\finexercice
\section{224.01 Intégrale multiple}


\exercice{1907, ridde, 1999/11/01}

\enonce{001907}{}
Calculer $I_1 = \iint\limits_{D} (x + y)e^{-x}e^{-y}dxdy$ où
$D  = \left\{ (x, y)\in \Rr^{2}/x, y\geq 0, x + y \leq 1\right\}$.

Calculer $I_2 = \iint\limits_{D} (x^{2} + y^{2})dxdy$ où
$D  = \left\{ (x, y)\in \Rr^{2}/x^{2} + y^{2}<x, x^{2} + y^{2}>y\right\}$.

Calculer $I_3 = \iint\limits_{D} \frac{xy}{1 + x^{2} + y^{2}}dxdy$ où
$D  = \left\{ (x, y)\in [0, 1]^{2}/x^{2} + y^{2} \geq 1\right\}$.

Calculer $I_4 = \iint\limits_{D} \frac{1}{y\cos (x) + 1}dxdy$ où
$D  = [0, \frac{\pi}{2}]\times [0, \frac{1}{2}]$.

Calculer $I_5 = \iiint\limits_{D} zdxdydz$ où
$D  = \left\{ (x, y, z)\in (\Rr^{ + })^{3}/y^{2} + z \leq 1, x^{2} + z \leq 1\right\}$.

Calculer $I_5 = \iint\limits_{D} xydxdy$ où
$D  = \left\{ (x, y)\in \Rr^{2}/x, y>0, \frac{x^{2}}{a^{2}} + \frac{y^{2}}{b^{2}}
\leq 1\right\}$ avec $a, b>0$.
\finenonce{001907}



\finexercice



\exercice{1908, ridde, 1999/11/01}

\enonce{001908}{}
Repr\'esenter et calculer le volume de $\left\{ (x, y, z)\in \Rr^{3}/-1\leq z \leq 1,
x^{2} + y^{2} \leq z^{2} + 1\right\}$.
\finenonce{001908}



\finexercice



\exercice{1909, ridde, 1999/11/01}

\enonce{001909}{}
D\'eterminer le centre de gravit\'e du culbuto (homog\`ene), \emph{i.e.} le cône
$$\left\{ (x, y, z)\in \Rr^{3}/z \in [0, 1], x^{2} + y^{2} \leq z^{2}\right\}$$
auquel on adjoint sur sa base une demi-boule.
\finenonce{001909}



\finexercice



\exercice{1910, maillot, 2001/09/01}

\enonce{001910}{}
Soit $D=[0,1]^2$. Calculer~:
$$\iint_D \frac{dx\,dy}{(x+y+1)^2}.$$
\finenonce{001910}



\finexercice



\exercice{1911, maillot, 2001/09/01}

\enonce{001911}{}
Soit $D$ le disque de centre $(0,1)$ et de rayon $1$ du plan. Calculer~:
$$\iint_D (x^2+y^2) \, dx\, dy.$$
\finenonce{001911}



\finexercice



\exercice{1912, maillot, 2001/09/01}

\enonce{001912}{}
Soit $D=\{x\ge0,y\ge0,x^2+y^2-2y\ge0,x^2+y^2-1\le 0\}$. Calculer~:
$$\iint_D \sqrt{x^2+y^2} \, dx\, dy.$$
\finenonce{001912}



\finexercice



\exercice{1913, maillot, 2001/09/01}

\enonce{001913}{}
Soit $D=\{(x^2+y^2)^2\le xy\}$. Calculer~:
$$\iint_D \sqrt{xy} \, dx\, dy.$$
\finenonce{001913}



\finexercice



\exercice{1914, maillot, 2001/09/01}

\enonce{001914}{}
Soient $a,b>0$. Calculer l'aire de l'ellipse
$E=\{ \frac{x^2}{a^2} + \frac{y^2}{b^2} \le 1 \}$ par deux m\'ethodes
diff\'erentes.

(On rappelle que l'aire d'un domaine $D$ vaut $\iint_D dx\, dy$.)
\finenonce{001914}



\finexercice



\exercice{1915, maillot, 2001/09/01}

\enonce{001915}{}
Soit $a>0$ et $D$ le domaine d\'elimit\'e par la courbe d'\'equation
polaire $\rho=a(1+\cos \theta)$. Calculer l'aire de $D$.
\finenonce{001915}



\finexercice



\exercice{1916, maillot, 2001/09/01}

\enonce{001916}{}
Soient $0<a\le b, 0<c\le d,$ et $D= \{ax^2\le y\le bx^2,
\frac cx\le y\le \frac d x\}$. Calculer l'aire de $D$.

(Indication: poser $u=\frac y{x^2}$ et $v=xy$.)
\finenonce{001916}



\finexercice



\exercice{1917, maillot, 2001/09/01}

\enonce{001917}{}
Soit $p>0$ et $D=\{y^2-2px\le 0,x^2-2py\le 0\}$. Calculer~:
$$\iint_D e^{\frac{x^3+y^3}{xy}} \, dx\, dy.$$

(Indication~: poser $x=u^2v$ et $y=uv^2$.)
\finenonce{001917}



\finexercice



\exercice{1918, maillot, 2001/09/01}

\enonce{001918}{}
Soit $R>0$, $D_R=\{x^2+y^2\le R^2, x>0, y>0\}$ et $K_R=[0,R]^2$.
Montrer que~:
$$\iint_{D_R} e^{-(x^2+y^2)}\, dx\, dy \le
\iint_{K_R} e^{-(x^2+y^2)}\, dx\, dy \le
\iint_{D_{2R}} e^{-(x^2+y^2)}\, dx\, dy.$$

En d\'eduire l'existence et la valeur de
$$\lim_{R\rightarrow +\infty} \int_0^R e^{-t^2}\, dt.$$
\finenonce{001918}



\finexercice



\exercice{1919, maillot, 2001/09/01}

\enonce{001919}{}
Soient $a,R>0$. Dans le plan $(yOz)$,
soit $D$ le disque de centre $(0,a,0)$ et de rayon $R$. En
tournant autour de l'axe $(Oz)$, le disque $D$ engendre un domaine
$T$ (appel\'e un tore plein). Calculer le volume de $T$
(c'est-\`a-dire l'int\'egrale triple $\iiint_T dx\, dy\, dz$).
\finenonce{001919}



\finexercice



\exercice{1920, maillot, 2001/09/01}

\enonce{001920}{}
Soit $D=\{x^2+y^2\le 1, 0\le z\le 1-x^2+y^2\}$. Calculer le
volume de $D$.
\finenonce{001920}



\finexercice



\exercice{1921, maillot, 2001/09/01}

\enonce{001921}{}
Soit $D=\{x\ge 0, y\ge 0, z\ge 0, x+y+z\le 1\}$. Calculer~:
$$\iiint_D \frac{dx\, dy\, dz}{(1+x+y+z)^3}.$$
\finenonce{001921}



\finexercice



\exercice{1922, maillot, 2001/09/01}

\enonce{001922}{}
Quel est le volume d\'elimit\'e par deux cylindres de r\'evolution
d'axes $(Ox)$ et $(Oy)$ et de m\^eme rayon $R>0$~?
\finenonce{001922}



\finexercice



\exercice{1923, drutu, 2003/10/01}

\enonce{001923}{}
En utilisant un changement de variables, calculer l'int\'egrale de
$f$ sur $D$ avec
\begin{enumerate}
\item $D=\{ (x,y)\in \R^2 \mid \pi^2 < x^2+y^2 \leq 4\pi^2 \}\; ;\; f(x,y)=\sin \sqrt{x^2+y^2}$ ;
\item $D=\left\{ (x,y)\in \R^2 \mid \frac{x^2}{a^2}+\frac{y^2}{b^2} \leq 1 \right\}\; \mbox{avec }a\, ,\, b>0\; ;\; f(x,y)=x^2+y^2$ ;
\item $D=\{ (x,y,z)\in \R^3 \mid x^2+y^2 \leq 1\, ,\, 0\leq z\leq h \}\; \mbox{avec }h>0 \; ;\; f(x,y,z)=z$ ;
\item $D=\{ (x,y)\in \R^2 \mid 0 < x^2\leq y\leq 2x^2\, ,\, 1/x\leq y \leq 2/x \}\; ;\; f(x,y)=x+y$  (changement de variable $u=y/x^2\, ,\, v=xy$) ;
\item $D=\{ (x,y,z)\in \R^3 \mid x\geq 0\, ,\, y\geq 0\, ,\, z\geq 0\, ,\, x^2+y^2+z^2 \leq 1 \}\; ;\; f(x,y,z)=xyz $ ;
\item $D=\{ (x,y,z)\in \R^3 \mid 1\leq x^2+y^2+z^2 \leq 4 \}\; ;\; f(x,y,z)=(x^2+y^2+z^2)^\alpha $.
\end{enumerate}
\finenonce{001923}



\finexercice



\exercice{1924, drutu, 2003/10/01}

\enonce{001924}{}
Identifier les ensembles suivants et calculer leur aire s'ils sont
dans $\R^2$, leur volume s'ils sont dans $\R^3$.
\begin{enumerate}
\item $D=\left\{ (x,y)\in \R^2 \mid \frac{x^2}{a^2}+\frac{y^2}{b^2} \leq 1 \right\}$ avec $a,b>0$ ;
\item $D=\left\{ (x,y,z)\in \R^3 \mid \frac{x^2}{a^2}+\frac{y^2}{b^2}+\frac{z^2}{c^2} \leq 1 \right\}$ avec $a,b,c>0$ ; qu'obtient-on dans le cas particulier o\`u $D$ est la boule unit\'e de $\R^3$ ?
\item $D=\left\{ (x,y,z)\in \R^3 \mid x^2 +y^2 \leq R\, ,\, 0\leq z\leq h \right\}$ avec $R,h>0$ ;
\item $D=\left\{ (x,y,z)\in \R^3 \mid x\geq 0\, ,\, y\geq 0\, ,\, z\geq 0\, ,\, x+y+z \leq 1 \right\}$ ;
\item $D=\left\{ (x,y,z)\in \R^3 \mid x^2+y^2\leq z^2/h^2 \, ,\,  0\leq z\leq h \right\}$ avec
$h>0$.
\end{enumerate}
\finenonce{001924}



\finexercice



\exercice{1925, drutu, 2003/10/01}

\enonce{001925}{}
Calculer les coordonn\'ees du centre d'inertie (de gravit\'e) du domaine $D$ :
\begin{enumerate}
\item $D=\left\{ (x,y)\in \R^2 \mid \frac{x^2}{a^2}+\frac{y^2}{b^2} \leq 1\, ,\, x\geq 0\, ,\, y\geq 0 \right\}$ (le quart d'ellipse) ;
\item $D=\left\{ (x,y)\in \R^2 \mid x^2+y^2 \leq 1\, ,\, x\geq 0 \, ,\, |y|\leq ax \right\}$ ;
\item $D=\{ (x,y)\in \R^2 \mid  x^2+y^2\leq 9\, ,\, (x-1)^2+y^2\geq 1 \}$.
\end{enumerate}
\finenonce{001925}



\finexercice



\exercice{1926, drutu, 2003/10/01}

\enonce{001926}{}
\begin{enumerate}
\item {\bf{Th\'eor\`eme de Guldin}}
Soit $D_0$ un domaine trac\'e dans le demi-plan $\{ (x,0,z)\in \R^3 \mid x\geq 0 \}$. Si l'on fait tourner $D_0$ autour de l'axe $Oz$, on obtient un domaine $D$ de $\R^3$. En utilisant les coordonn\'ees cylindriques. montrer que 
$$
Vol\, (D) = 2\pi Aire\, (D_0)\cdot x_G\, ,
$$ o\`u $(x_G,z_G)$ sont les coordonn\'ees du centre d'inertie du domaine $D_0$.
\item Calculer les volumes des domaines suivants :
\begin{enumerate}
\item le tore obtenu en faisant tourner autour de $Oz$ le domaine 
$D_0=\{ (x,0,z)\mid \frac{(x-c)^2}{a^2} +\frac{z^2}{b^2}\leq 1 \}$, o\`u $a<c$ ;
\item $D=\left\{ (x,y,z)\in \R^3 \mid x^2 +y^2 +z^2 \leq 4R^2\, ,\, x^2 +y^2 \leq R^2 \right\}$, o\`u $R>0$.
\end{enumerate}
\end{enumerate}
\finenonce{001926}



\finexercice



\exercice{1927, legall, 2003/10/01}

\enonce{001927}{}
On pose
  $I= \displaystyle{\int _\Rr e^\frac{-t^2}{2}dt  }$ et  $J= 
\displaystyle{\int \! \! \! \int _{\Rr ^2} 
e^\frac{-(x^2+y^2)}{2}dxdy. }$
Calculer $J$ et en d\'eduire la valeur de $I$.
\finenonce{001927}



\finexercice



\exercice{1928, legall, 2003/10/01}

\enonce{001928}{}

On note $D$ le domaine d\'elimit\'e par les droites $x=0,$ $y=x+2$ et $y=-x$.
  \begin{enumerate}
\item Calculer (directement) $I= \displaystyle{\int \! \! \! \int 
_{D} (x-y)dxdy. }$
\item Calculer $I$ au moyen du changement de variable $u=x+y$ et $v=x-y.$
\end{enumerate}
\finenonce{001928}



\finexercice



\exercice{1929, legall, 2003/10/01}

\enonce{001929}{}

Soit $D=\{ (x,y) ; x\geq 0, y\geq 0, x^2+y^2\leq 1\} $.
Calculer $\displaystyle{\int \! \! \! \int _{D} (4-x^2-y^2)dxdy. }$
\finenonce{001929}



\finexercice



\exercice{4381, quercia, 2010/03/12}
\enonce{004381}{Intégrales doubles}

Calculer $\iint_D^{} f(x,y)\,d xd y$ : 

   \begin{enumerate}
  \item $D=\{y\ge0, x+y\le1, y-x\le1\}$,                     \par\nobreak $f(x,y)=x^2y$.                                    
  \item $D=\{x^2+y^2\le R^2\}$,                              \par\nobreak $f(x,y)=x^2y$.                                    
  \item $D=\{\frac{x^2}{a^2} + \frac{y^2}{b^2} \le 1\}$,   \par\nobreak $f(x,y)=x^2 + y^2$.                               
  \item $D=\{0 \le x \le 1-\frac{y^2}4\}$,                  \par\nobreak $f(x,y)=x^2 + y^2$.                               
  \item $D=\{x^2+y^2 \le 1\}$,                               \par\nobreak $f(x,y)=(x+y)^2$.                                 
  \item $D=\{x^2+y^2 \le 1\}$,                               \par\nobreak $f(x,y)=\frac{(x+y)^2}{x^2+y^2+1}$.              
  \item $D=\{x\ge0, y\ge0, x+y\le1\}$,                       \par\nobreak $f(x,y)=x+y+1$.                                   
  \item $D=\{|x+y|\le1,|x-y|\le1\}$,                         \par\nobreak $f(x,y)=\ln(x+y+1)$.                              
  \item $D=\{x\ge0, y\ge0, x+y\le\pi\}$,                     \par\nobreak $f(x,y)=(x+y)\sin x\sin y$.                       
  \item $D=\{|x|\le x^2+y^2\le 1\}$,                         \par\nobreak $f(x,y)=(1+x^2+y^2)^2$.                           
  \item $D=\{x\ge0, y\ge0, x+y\le a\}$,                      \par\nobreak $f(x,y)=x+y+\sqrt{a^2+(x+y)^2}$.                  
  \item $D=\{x\ge0, y\ge0, x^2+y^2\le 1\}$,                  \par\nobreak $f(x,y)=xy\sqrt{x^2+4y^2}$.                       
  \item $D=\{x^2+y^2-2y \le 0\}$,                            \par\nobreak $f(x,y)=y\exp(x^2+y^2-2y)$.                       
  \item $D=\{y^2 \le 2px, x^2 \le 2py\}$,                    \par\nobreak $f(x,y)=\exp\left(\frac{x^3+y^3}{xy}\right)$.    
\end{enumerate}
\finenonce{004381}



\finexercice


\exercice{4382, quercia, 2010/03/12}
\enonce{004382}{ESEM 94}
    Calculer $I = \iint_\Delta xy\,d xd y$ où
    $\Delta = \{(x,y) \text{ tel que } y\ge 0 \text{ et } (x+y)^2 \le 2x/3 \}$.
\finenonce{004382}



\finexercice


\exercice{4383, quercia, 2010/03/12}
\enonce{004383}{Ensi PC 1999}
    Calculer $I = \iint_\Delta (x^2 + xy + y^2)\,d xd y$ où
    $\Delta = \{(x,y) \text{ tel que } y\ge 0 \text{ et } x^2+y^2-2x\le 0 \text{ et } x^2+y^2-2y\le 0 \}$.
\finenonce{004383}



\finexercice


\exercice{4384, quercia, 2010/03/12}
\enonce{004384}{Intégrales triples}

Calculer $\iiint_D^{} f(x,y,z)\,d xd yd z$ : 
   \begin{enumerate}
  \item $D=\{0\le x\le1, 0\le y\le1, 0\le z\le1\}$,   \par\nobreak $f(x,y,z)=\frac1{(x+y+z+1)^3}$.                         
  \item $D=\{x^2+y^2+z^2 \le R^2\}$,                  \par\nobreak $f(x,y,z)=\frac1{\sqrt{a^2-x^2-y^2-z^2}} \quad(a>R>0)$. 
  \item $D=\{x\ge0, y\ge0, z\ge0, x+y+z\le1\}$,       \par\nobreak $f(x,y,z)=xyz$.                                          
  \item $D=\{x\ge0, y\ge0, z\ge0, x+y+z\le1\}$,       \par\nobreak $f(x,y,z)=\frac1{(x+y+z+1)^2}$.                         
  \item $D=\{x^2+y^2\le R^2, 0\le z\le a\}$,          \par\nobreak $f(x,y,z)=x^3+y^3+z^3-3z(x^2+y^2)$.                      
  \item $D=\{x^2+y^2\le z^2, 0\le z\le 1\}$,          \par\nobreak $f(x,y,z)=\frac z{(x^2+y^2+1)^2}$.   
  \item $D=\left\{\frac{x^2}{a^2} + \frac{y^2}{b^2} + \frac{z^2}{c^2} \le 1\right\}$, \par\nobreak  $f(x,y,z)=x^2+y^2$.
\end{enumerate}
\finenonce{004384}



\finexercice


\exercice{4385, quercia, 2010/03/12}
\enonce{004385}{Ensi Chimie P 93}
\begin{enumerate}
  \item Calculer $\iiint_D \frac{d x\,d y\,d z}{(1+x^2z^2)(1+y^2z^2)}$ avec
    $D = \{(x,y,z) \text{ tel que } 0\le x \le 1,\ 0\le y\le 1,\ 0\le z\}$.
    
  \item En déduire $ \int_{t=0}^{+\infty} \Bigl(\frac{\Arctan t}t\Bigr)^2d t$.
\end{enumerate}
\finenonce{004385}



\finexercice


\exercice{4386, quercia, 2010/03/12}
\enonce{004386}{Ensi Chimie P 93}

Soit $I =  \int_0^1\frac{\ln(1+x)}{1+x^2}\,d x$.\par
En calculant $J=\iint_D\frac{x\,d x\,d y}{(1+x^2)(1+xy)}$ avec
$D = \{(x,y) \text{ tel que } 0\le x \le 1,\ 0\le y\le 1\}$ de deux façons différentes,
trouver $I$.

\finenonce{004386}



\finexercice


\exercice{4387, quercia, 2010/03/12}
\enonce{004387}{Ensi Chimie P 93}

Soit $T$ un tore plein d'axe $(Oz)$ et de rayons $R$, $r$ ($R > r$).
Calculer $\iiint_T(x^2+y^2)\,d x\,d y\,d z$.

\finenonce{004387}



\finexercice


\exercice{4388, quercia, 2010/03/12}
\enonce{004388}{$MF + MF'$}

Soit $\cal E$ l'ellipse d'équation $\frac{x^2}{a^2} + \frac{y^2}{b^2} = 1$
($0<b<a$),
$E$ le domaine limité par $\cal E$ et $F,F'$ les foyers de $\cal E$.
Calculer $I = \iint_{M\in E}^{} (MF+MF')\,d xd y$.

On effectuera le changement de variable :
$x = \sqrt{ u^2+c^2 }\cos v$, $y = u\sin v$ où $c = \sqrt{a^2-b^2}$.
\finenonce{004388}



\finexercice


\exercice{4389, quercia, 2010/03/12}
\enonce{004389}{$ \int_{x=0}^{\pi/2} \frac{\ln(1+\cos x)}{\cos x}d x$}

\begin{enumerate}
  \item Montrer l'existence de $I =  \int_{x=0}^{\pi/2} \frac{\ln(1+\cos x)}{\cos x}d x$.
  \item Montrer que $I = \iint_D^{} \frac{\sin y}{1+\cos x\cos y} d xd y$ où $D = [0,\frac\pi2]^2$.
  \item En déduire la valeur de $I$.
    
\end{enumerate}
\finenonce{004389}



\finexercice


\exercice{4390, quercia, 2010/03/12}
\enonce{004390}{Intégrale de Gauss}

Calcul de $I =  \int_{t=0}^{+\infty} e^{-t^2}\,d t$.

\begin{enumerate}
  \item Justifier la convergence de cette intégrale.
  \item Pour $a > 0$ on note $\Delta_a = [0,a] \times [0,a]$ et
    $C_a$ le quart de disque d'équations : $x^2+y^2 \le a^2$, $x\ge0$, $y\ge0$.
 \begin{enumerate} 
    \item Encadrer l'intégrale sur $\Delta_a$ de $f(x,y) = e^{-x^2-y^2}$ par
        les intégrales de $f$ sur des domaines du type $C_b$.

    \item Calculer $\iint_{C_b}^{} f(x,y)\,d xd y$ en polaires et en déduire la
        valeur de $I$.
  \end{enumerate}
\end{enumerate}
\finenonce{004390}



\finexercice


\exercice{4391, quercia, 2010/03/12}
\enonce{004391}{$ \int_{t=0}^1 \frac{\ln t}{1-t^2}\,d t$}

\begin{enumerate}
  \item Calculer $A = \iint_{0\le y\le x\le 1}^{} \frac{d xd y}{(1+x^2)(1+y^2)}$.
    
  \item Démontrer la convergence des intégrales :\par
    $B =  \int_{\theta=0}^{\pi/4} \frac{\ln(2\cos^2\theta)}{2\cos2\theta}  d\theta$,
    $C =  \int_{\theta=0}^{\pi/4} \frac{\ln(2\sin^2\theta)}{2\cos2\theta}  d\theta$,
    et $D =  \int_{t=0}^1 \frac{\ln t}{1-t^2}\,d t$.
  \item Démontrer que $A = B$ (passer en coordonnées polaires dans $A$).
  \item Calculer $B+C$ et $B-C$ en fonction de $D$.
    
  \item En déduire les valeurs de $C$ et $D$.
    
\end{enumerate}
\finenonce{004391}



\finexercice


\exercice{4392, quercia, 2010/03/12}
\enonce{004392}{Aires}

Calculer l'aire des domaines suivants :

\begin{enumerate}
  \item $D$ est la partie du disque unité située dans la concavité de
    l'hyperbole d'équation $xy = \frac{\sqrt3}4$.
    


  \item $D$ est l'intersection des domaines limités par les ellipses d'équation
    $\frac{x^2}{a^2} + \frac{y^2}{b^2} = 1$ et $\frac{x^2}{b^2} + \frac{y^2}{a^2} = 1$.
\end{enumerate}
\finenonce{004392}



\finexercice


\exercice{4393, quercia, 2010/03/12}
\enonce{004393}{Ensi P 90}
    Soit $\cal P$ le plan rapporté au repère $(O,\vec i,\vec j)$.
    Calculer l'aire du domaine délimité par la courbe d'équation
    $x^{2/3}+y^{2/3} = a^{2/3}$.
\finenonce{004393}



\finexercice


\exercice{4394, quercia, 2010/03/12}
\enonce{004394}{Chimie P 91}
    On considère les courbes planes : ${\cal Q}_i$ : $(x^2 = 2q_iy)$ et
				      ${\cal P}_i$ : $(y^2 = 2p_ix)$.
    On suppose $0 < q_1 < q_2$ et $0 < p_1 < p_2$. Calculer l'aire du
    ``quadrilatère'' limité par ${\cal P}_1$,${\cal P}_2$,${\cal Q}_1$ et
    ${\cal Q}_2$.
    
\finenonce{004394}



\finexercice


\exercice{4395, quercia, 2010/03/12}
\enonce{004395}{Chimie P 1996}
    Calculer l'aire délimitée par la courbe d'équation $(y-x)^2 = a^2-x^2$.
\finenonce{004395}



\finexercice


\exercice{4396, quercia, 2010/03/12}
\enonce{004396}{Volumes}
Calculer le volume des domaines suivants :

\begin{enumerate}
  \item $D$ est l'intersection du cylindre de révolution d'axe $Oz$ de rayon $a$
    et de la boule de centre $O$ de rayon 1 ($0 < a < 1$).
    

  \item $D$ est l'intersection de la boule de centre $O$ de rayon 1 et du cône de
    révolution d'axe $Oz$ et de demi-angle~$\frac\pi4$.
    


  \item $D$ est le volume engendré par la rotation d'un disque de rayon $r$ autour
    d'une droite coplanaire avec le disque, située à la distance $R > r$
    du centre du disque (tore de révolution ou chambre à air).
    
\end{enumerate}
\finenonce{004396}



\finexercice


\exercice{4397, quercia, 2010/03/12}
\enonce{004397}{Ensi Physique P 94}
    Calculer le volume intérieur au paraboloïde d'équation $x^2+y^2 = 2pz$
    et extérieur au cône d'équation $x^2+y^2 = \lambda^2z^2$
    ($p>0$, $\lambda>0$).
\finenonce{004397}



\finexercice


\exercice{4398, quercia, 2010/03/12}
\enonce{004398}{Volume}
Dans le plan $Oxy$  on considère la courbe $\mathcal{C}$ d'équation polaire
$\rho = a\sqrt{\cos 2\theta}$ $(a>0,\ -\frac\pi4 \le \theta \le \frac\pi4 )$.
En tournant autour de $Ox$, $\mathcal{C}$ engendre une surface dont on calculera le
volume qu'elle limite
(on posera $x = \rho\cos\theta$, $y=\rho\sin\theta\cos\phi$,
$z=\rho\sin\theta\sin\phi$).
\finenonce{004398}



\finexercice


\exercice{4399, quercia, 2010/03/12}
\enonce{004399}{Volume}
On coupe une demi-boule par un plan $P$ parallèle à sa base. Quelle doit
être la position de $P$ pour que les deux morceaux aient même volume ?
(Donner un résultat approché)

\finenonce{004399}



\finexercice


\exercice{4400, quercia, 2010/03/12}
\enonce{004400}{Somme double}

Soit $f : {[0,1]} \to \R$ continue.

Chercher $\lim_{n\to\infty} \frac1{n^2} \sum_{0\le i\le j \le n} f\left(\frac in\right)f\left(\frac jn\right)$.
\finenonce{004400}



\finexercice


\exercice{4401, quercia, 2010/03/12}
\enonce{004401}{Nombre de couples $(a,b)$ tel que $a^2+b^2 \le n$}

Pour $n \in \N$ on pose
$E_n = \{ (p,q)\in \N^2 \text{ tel que } p^2+q^2 \le n \}$ et $C_n = \mathrm{Card}\,(E_n)$.

Interpréter $C_n$ comme une aire et donner un équivalent de $C_n$ lorsque
$n\to\infty$.
\finenonce{004401}



\finexercice


\exercice{4402, quercia, 2010/03/12}
\enonce{004402}{Ens MP 2002}

Soit $f\in\mathcal{C}([0,1],\R^+)$ telle que $ \int_0^1f = 1$. Pour $\psi\in\mathcal{C}([0,1],\R)$
on pose $$\Lambda_n(\psi) =  \int_0^1\!\dots \int_0^1\psi\Bigl(\frac{x_1+\dots+x_n}n\Bigr)f(x_1)\dots f(x_n)\,d x_1\dots d x_n.$$
Montrer que $\Lambda_n(\psi)\to\psi\Bigl( \int_{x=0}^1xf(x)\,d x\Bigr)$  lorsque $n\to\infty$.


\finenonce{004402}



\finexercice


\exercice{5908, rouget, 2010/10/16}
\enonce{005908}{**}
Calculer les intégrales multiples suivantes

\begin{enumerate}
 \item  $I=\displaystyle\iint_{D}(x+y)\;dxdy$ où $D=\{(x,y)\in\Rr^2/\;x\leqslant1,\;y\leqslant1,\;x+y\geqslant1\}$.

\item  $I=\displaystyle\iint_{[-1,1]^2}|x+y|\;dxdy$.

\item  $I=\displaystyle\iint_{D}xy\;dxdy$ où $D$ est la partie du plan limitée par les paraboles d'équations respectives $y=x^2$ et $x=y^2$.

\item  $I=\displaystyle\iint_{x^2+y^2\leqslant1} \frac{1}{1+x^2+y^2}\;dxdy$.

\item  $I=\displaystyle\iint_{x\leqslant x^2+y^2\leqslant1} \frac{dxdy}{(1+x^2+y^2)^2}$.

\item  $I=\displaystyle\iiint_{0\leqslant x\leqslant y\leqslant z\leqslant1}xyzdxdydz$.

\item  $I=\displaystyle\iiint_{\sqrt{x}+\sqrt{y}+\sqrt{z}\leqslant1}zdxdydz$.
\end{enumerate}
\finenonce{005908}


\finexercice


\exercice{5910, rouget, 2010/10/16}
\enonce{005910}{***}
Soient $(p_1,p_2,q_1,q_2)\in]0,+\infty[^4$ tel que $p_1<p_2$ et $q_1<q_2$.

Calculer l'aire du domaine $D=\{(x,y)\in\Rr^2/\;2p_1x\leqslant y^2\leqslant2p_2x\;\text{et}\;2q_2y\leqslant x^2\leqslant2q_2y\}$.
\finenonce{005910}


\finexercice


\exercice{5911, rouget, 2010/10/16}
\enonce{005911}{*** I}
Calculer le volume de $B=\{(x_1,\ldots,x_n)\in\Rr^n/\;x_1^2+\ldots+x_n^2\leqslant1\}$ (boule unité fermée de $\Rr^n$ pour $\|\;\|_2$).
\finenonce{005911}


\finexercice


\exercice{5912, rouget, 2010/10/16}
\enonce{005912}{**}
Calculer le volume de l'intérieur de l'ellipsoïde d'équation $x^2+ \frac{1}{2}y^2+ \frac{3}{4}z^2+xz=1$.
\finenonce{005912}


\finexercice


\exercice{5914, rouget, 2010/10/16}
\enonce{005914}{***}
Calculer $I=\displaystyle\iint_{\frac{x^2}{a^2}+\frac{y^2}{b^2}\leqslant1}(x^2-y^2)\;dxdy$.
\finenonce{005914}


\finexercice
\section{224.02 Calcul approché d'intégrale}

\section{224.03 Intégrale de Riemann dépendant d'un paramètre}


\exercice{4313, quercia, 2010/03/12}
\enonce{004313}{Calcul de limite}

Chercher $\lim_{x\to0}  \int_{t=x}^{2x} \frac{\cos t\ln(1+t^2)}{\sin^2t\sh t}\,d t$.

\finenonce{004313}



\finexercice


\exercice{4314, quercia, 2010/03/12}
\enonce{004314}{Calcul de limite, Ensi P 90}

    Calculer les limites : $\lim_{x\to0} \int_x^{3x} \frac t{\tan^2t}\,dt$
    et $\lim_{x\to0} \frac1{x^3} \int_0^x \frac {t^2}{t+e^{3t}}\,dt$.
    

\finenonce{004314}


$\frac t{\tan^2t} = \frac1t + \varphi(t)$ avec $\varphi$
	     prolongeable par continuité en $0$, donc
	     $\lim_{x\to0} \int_x^{3x} \frac t{\tan^2t}\,dt = \ln 3$.\par
	     $\frac{t^2}{t+e^{3t}} = t^2 + \text{o}(t^2)$ donc
	     $\lim_{x\to0} \frac1{x^3} \int_0^x \frac {t^2}{t+e^{3t}}\,dt = \frac13$.
\finexercice


\exercice{4315, quercia, 2010/03/12}
\enonce{004315}{Calcul de limite}

Chercher $\lim_{x\to+\infty} \frac1{x^2} \int_{t=3}^{x^2+x} \frac{\sin t\,d t}{3+\ln(\ln t)}$.

\finenonce{004315}



\finexercice


\exercice{4316, quercia, 2010/03/12}
\enonce{004316}{Calcul de limite}

Chercher $\lim_{x\to0^+}  \int_{t=x}^{x^2} \frac{e^{-t}\,d t}{\sin t\ln t}$.

\finenonce{004316}



\finexercice


\exercice{4317, quercia, 2010/03/12}
\enonce{004317}{Série d'intégrales, Esem 91}

    \'Etablir la convergence et calculer la somme de
    $\sum_{n=1}^\infty (-1)^n \int_0^{+\infty} \frac{dt}{(1+t^3)^n}$.
    

\finenonce{004317}



\finexercice


\exercice{4318, quercia, 2010/03/12}
\enonce{004318}{$\sin(t)/(t+x)$}

\begin{enumerate}
  \item Prouver l'existence pour $x > 0$ de
    $I(x) =  \int_{t=0}^{+\infty} \frac{\sin t}{t+x}\,d t$.

  \item Déterminer $\lim_{x\to+\infty} I(x)$.
    

\end{enumerate}
\finenonce{004318}



\finexercice


\exercice{4319, quercia, 2010/03/12}
\enonce{004319}{Calcul de limite}

Soit $f : {[0,1]} \to \R$ continue.
Chercher $\lim_{x\to0^+}  \int_{t=0}^1 \frac{xf(t)}{x^2+t^2}\,d t$.

\finenonce{004319}



\finexercice


\exercice{4320, quercia, 2010/03/12}
\enonce{004320}{Calcul d'équivalent, Mines 1999}

Donner un équivalent pour $x\to+\infty$ de
$ \int_{t=0}^{+\infty}\frac{\sin t}{x^2+t^2}\,d t$.

\finenonce{004320}



\finexercice


\exercice{4321, quercia, 2010/03/12}
\enonce{004321}{ Calcul de limite}

Soit $a > 0$. Donner le DL en $x=1$ à l'ordre 3 de
$f(x) =  \int_{t=a/x}^{ax} \frac{\ln t}{a^2+t^2}\,d t$.

\finenonce{004321}



\finexercice


\exercice{4322, quercia, 2010/03/12}
\enonce{004322}{$(\int f^x)^{1/x}$}

Soit $f : {[a,b]} \to {\R^+}$ continue. On pose
$\varphi(x) = \left({ \int_{t=a}^b (f(t))^x\,d t}\right)^{1/x}$.

\begin{enumerate}
  \item Montrer que $\varphi(x) \to \max(f)$ lorsque $x\to+\infty$.
  \item On suppose $f > 0$ et $b-a = 1$.
    Montrer que $\varphi(x) \to \exp\left({ \int_{t=a}^b\ln(f(t))\,d t}\right)$ lorsque $x\to0^+$.
    

\end{enumerate}
\finenonce{004322}



\finexercice


\exercice{4323, quercia, 2010/03/12}
\enonce{004323}{$t^nf(t)$}

Soit $I_n =  \int_{t=0}^1 t^n\ln(1+t^2)\,d t$.
Montrer que $I_n \to 0$ lorsque $n\to\infty$.

\finenonce{004323}



\finexercice


\exercice{4324, quercia, 2010/03/12}
\enonce{004324}{$t^nf(t)$}

Soit $f : {[0,1]} \to \R$ continue.
Montrer que $ \int_{t=0}^1 t^nf(t)\,d t = \frac{f(1)}n +  o\Bigl(\frac1n\Bigr)$.

\finenonce{004324}



\finexercice


\exercice{4325, quercia, 2010/03/12}
\enonce{004325}{$f(t^n)$}

\begin{enumerate}
  \item Soit $f : {[0,1]} \to \R$ continue.
    Montrer que $ \int_{t=0}^1 f(t^n)\,d t \to f(0)$  lorsque $n\to\infty$.
    
  \item Chercher un équivalent pour $n\to\infty$ de $ \int_{t=0}^1 \frac{t^n\,d t}{1+t^n}$.
    
  \item Chercher un équivalent pour $n\to\infty$ de $-1 +  \int_{t=0}^1 \sqrt{1+t^n}\,d t$.
    
\end{enumerate}
\finenonce{004325}



\finexercice


\exercice{4326, quercia, 2010/03/12}
\enonce{004326}{$f(t^n)$}

Donner les deux premiers termes du DL pour $n\to\infty$ de
$I_n =  \int_{t=0}^1 \frac{d t}{1+t^n}$.

\finenonce{004326}



\finexercice


\exercice{4327, quercia, 2010/03/12}
\enonce{004327}{$f(t^n)$}
Donner les deux premiers termes du DL pour $n\to\infty$ de
$I_n =  \int_{t=0}^1 \sqrt{1+t^n}\,d t$.


\finenonce{004327}



\finexercice


\exercice{4328, quercia, 2010/03/12}
\enonce{004328}{$f(t^n)$}

Chercher un équivalent pour $n\to\infty$ de $ \int_{t=1}^{1+1/n} \sqrt{1+t^n}\,d t$.


\finenonce{004328}



\finexercice


\exercice{4329, quercia, 2010/03/12}
\enonce{004329}{Calcul de limite}

Déterminer $\lim_{n\to\infty}  \int_{x=0}^{\pi/2} \frac{\sin^nx}{x+2}\,d x$.

\finenonce{004329}



\finexercice


\exercice{4330, quercia, 2010/03/12}
\enonce{004330}{Calcul de limite}

Soit $f : {[0,1]} \to \R$ continue.
Déterminer $\lim_{n\to\infty} \int_{t=0}^1 nf(t)e^{-nt}\,d t$.

\finenonce{004330}



\finexercice


\exercice{4331, quercia, 2010/03/12}
\enonce{004331}{$(1-x/n)^n$, Ensi PSI 1998}

Soit $x\in{[0,n]}$. Montrer que $(1-x/n)^n \le e^{-x}$.
En déduire $\lim_{n\to\infty}  \int_{x=0}^n(1-x/n)^n\,d x$.

\finenonce{004331}



\finexercice


\exercice{4332, quercia, 2010/03/12}
\enonce{004332}{\'Equation intégrale, Ensi P 91}

Déterminer les fonctions $f \in \mathcal{C}^0(\R,\R)$ telles que :
$\forall\ x\in \R,\ f(x) +  \int_0^x (x-t)f(t)\,dt = 1$.


\finenonce{004332}



\finexercice


\exercice{4333, quercia, 2010/03/12}
\enonce{004333}{$\tan^nt$, Ensi Physique P 94}

On pose $I_n =  \int_0^{\pi/4} \tan^nt\,d t$.\par

\begin{enumerate}
  \item Montrer que $I_n \to 0$  lorsque $n\to\infty$.

  \item Calculer $I_n$ en fonction de~$n$.

  \item Que peut-on en déduire~?

\end{enumerate}
\finenonce{004333}



\finexercice


\exercice{4334, quercia, 2010/03/12}
\enonce{004334}{Calcul de limite, École de l'air 94}

    Chercher $\lim_{n\to\infty}  \int_0^1\frac{t^n-t^{2n}}{1-t}\,d t$.
    
\finenonce{004334}



\finexercice


\exercice{4335, quercia, 2010/03/12}
\enonce{004335}{Approximation de la mesure de Dirac}

Soit $f : {[a,b]} \to {\R^+}$ continue atteignant son maximum en un unique
point $c \in {]a,b[}$.

\begin{enumerate}
  \item Soit $\mu > 0$ tel que $[c-\mu,c+\mu] \subset [a,b]$.
    Chercher $\lim_{n\to\infty} \biggl( \int_{t=a}^b f^n(t)\,d t\biggm/ \int_{t=c-\mu}^{c+\mu} f^n(t)\,d t\biggr)$.
    
  \item Soit $g : {[a,b]} \to {\R^+}$ continue.
    Chercher $\lim_{n\to\infty} \biggl( \int_{t=a}^b f^n(t)g(t)\,d t\biggm/ \int_{t=a}^b f^n(t)\,d t\biggr)$.
    

\end{enumerate}
\finenonce{004335}



\finexercice


\exercice{4336, quercia, 2010/03/12}
\enonce{004336}{\'Equation intégrale}

Soit $f : {[0,+\infty[} \to {\R^+}$ continue telle que
$f(x) \int_{t=0}^x f^2(t)\,d t \to \ell\ne 0$ (lorsque $x\to+\infty$).
Trouver un équivalent de $f$ en~$+\infty$.

\finenonce{004336}



\finexercice


\exercice{4337, quercia, 2010/03/12}
\enonce{004337}{Convolution}

Soient ${f,g} : \R \to \R$ continues et $a,b \in \R$.
On pose $\varphi(x) =  \int_{t=a}^b f(t)g(x-t)\,d t$.

\begin{enumerate}
  \item Montrer que $\varphi$ est continue et que si $g$ est de classe $\mathcal{C}^k$, alors
    $\varphi$ l'est aussi.

  \item Montrer que si $f$ est de classe $\mathcal{C}^1$ (et $g$ continue), alors $\varphi$ est
    aussi de classe $\mathcal{C}^1$.
\end{enumerate}
\finenonce{004337}



\finexercice


\exercice{4338, quercia, 2010/03/12}
\enonce{004338}{Convolution (Mines MP 2003)}
Soient $f,g\in\mathcal{C}([0,+\infty[,\R)$. On pose $h(x) =  \int_{t=0}^x f(x-t)g(t)\,d t$.

\begin{enumerate}
  \item Existence et continuité de~$h$.
  \item Peut-on inverser $f$ et $g$~?
  \item On suppose $f$ intégrable sur~$[0,+\infty[$ et $g$ bornée. Montrer que~$h$ est bornée.
  \item On prend $f(x) = \frac{\sin x}x$ et $g(x) = \cos(\alpha x)$ avec $0\le\alpha\le 1$.
    $h$ est-elle bornée (on pourra étudier les cas $\alpha=0$ et $\alpha=1$)~?
    
\end{enumerate}
\finenonce{004338}



\finexercice


\exercice{4339, quercia, 2010/03/12}
\enonce{004339}{Calcul d'intégrale}

\begin{enumerate}
  \item Calculer $\varphi(a) =  \int_{t=0}^1 \frac{d t}{1+at}$.
    

  \item En déduire la valeur de $ \int_{t=0}^1 \frac{td t}{(1+at)^2}$.
    

\end{enumerate}
\finenonce{004339}



\finexercice


\exercice{4340, quercia, 2010/03/12}
\enonce{004340}{Fonction définie par une intégrale}

On pose $\varphi(x) =  \int_{t=0}^1 e^{-x/t}\,d t$.

\begin{enumerate}
  \item Montrer que $\varphi$ est de classe $\mathcal{C}^\infty$ sur $\R^{+*}$.
  \item Vérifier que $\varphi''(x) = \frac{e^{-x}}x$.

\end{enumerate}
\finenonce{004340}



\finexercice


\exercice{4341, quercia, 2010/03/12}
\enonce{004341}{Fonction définie par une intégrale, Mines 1999}

Soit $I(\alpha) =  \int_{x=0}^{+\infty} \frac{x^{\alpha-1}}{1+x}\,d x$.
Montrer que $I(\alpha)$ existe et définit une fonction de classe
$\mathcal{C}^1$ sur $]0,1[$. \'Ecrire $I(\alpha)$ comme somme d'une série.
\finenonce{004341}



\finexercice


\exercice{4342, quercia, 2010/03/12}
\enonce{004342}{Fonction définie par une intégrale}

On pose pour $x \ge 0$ :
$f(x) =  \int_{t=0}^{+\infty} \frac{\ln(x^2+t^2)}{1+t^2}\,d t$.

Calculer explicitement $f'(x)$ et en déduire $f(x)$
(on calculera $f(0)$ à
l'aide du changement de variable $u = 1/t$).


\finenonce{004342}



\finexercice


\exercice{4343, quercia, 2010/03/12}
\enonce{004343}{Fonction définie par une intégrale}

On pose $I(x) =  \int_{t=0}^{\pi/2} \ln(\cos^2t + x^2\sin^2t)\,d t$.

\begin{enumerate}
  \item Montrer que $I$ est de classe $\mathcal{C}^1$ sur $\R^{+*}$.
  \item Calculer $I'(x)$ et en déduire $I(x)$.


\end{enumerate}
\finenonce{004343}



\finexercice


\exercice{4344, quercia, 2010/03/12}
\enonce{004344}{Intégrale de Gauss}

On considère les fonctions définies par :
$f(x) = \left({ \int_{t=0}^x e^{-t^2}\,d t}\right)^2$ et
$g(x) =  \int_{t=0}^1 \frac{e^{-x^2(1+t^2)}}{1+t^2}\,d t$.

\begin{enumerate}
  \item Montrer que $f$ et $g$ sont dérivables et calculer $f'$ et $g'$.
    
  \item Montrer que $f(x) + g(x) = \frac\pi4$ pour tout $x \in \R^+$.
  \item En déduire la valeur de $ \int_{t=0}^{+\infty} e^{-t^2}\,d t$.
\end{enumerate}
\finenonce{004344}



\finexercice


\exercice{4345, quercia, 2010/03/12}
\enonce{004345}{Intégrale de Gauss, Ensi PC 1999}

On donne~: $ \int_{t=0}^{+\infty} e^{-t^2}\,d t =\frac{\sqrt\pi}2$.
Existence et valeur de
$ \int_{t=0}^{+\infty} e^{-(t^2+a^2/t^2)}\,d t$.
\finenonce{004345}



\finexercice


\exercice{4346, quercia, 2010/03/12}
\enonce{004346}{Fonction définie par une intégrale}

\begin{enumerate}
  \item Soit $I(x) =  \int_{t=0}^{+\infty} e^{-t^2}\cos(2xt)\,d t$.
    Prouver que $I$ est de classe $\mathcal{C}^1$ sur $\R$.

  \item Chercher une relation simple entre $I$ et $I'$.
    

  \item En déduire la valeur de $I(x)$
    $\Bigl($on admet que $I(0) = \frac{\sqrt\pi}2\Bigr)$.
    

\end{enumerate}
\finenonce{004346}



\finexercice


\exercice{4347, quercia, 2010/03/12}
\enonce{004347}{Fonctions définies par des intégrales}

On pose, pour $x$ réel,
$F(x) =  \int^{+\infty}_{t=0} \frac{1 - \cos tx}{t^2}\, d t$
et $G(x) =  \int^{+\infty}_{t=0} \frac{1 - \cos tx}{t\sqrt{1+t^2}}\,d t$.

\begin{enumerate}
  \item Montrer que les intégrales $F(x)$ et $G(x)$
convergent absolument pour tout $x$ réel et que $F(x) = |x| F(1)$.

  \item Montrer que la fonction $F-G$
est de classe $\mathcal{C}^1$ sur~$\R$. En déduire que $G$ est $\mathcal{C}^1$ sur $\R^*$ et n'est
pas dérivable en~$0$.

\end{enumerate}
\finenonce{004347}



\finexercice


\exercice{4348, quercia, 2010/03/12}
\enonce{004348}{Théorème de division des fonctions $\mathcal{C}^\infty$}

Soit $f:\R \to \R$ de classe $\mathcal{C}^\infty$ et
$g(x) = \begin{cases}\frac{f(x)-f(0)}x & \text{ si } x \ne 0\cr
               f'(0) & \text{ si } x = 0.\cr\end{cases}$

Vérifier que $g(x) =  \int_{t=0}^1 f'(tx)\,d t$. En déduire que $g$ est de classe $\mathcal{C}^\infty$.

Montrer de même que la fonction $g_k : x  \mapsto \frac1{x^k}\left({f(x) - f(0) - xf'(0) - \dots - \frac{x^{k-1}}{(k-1)!}f^{(k-1)}(0)}\right)$
se prolonge en une fonction de classe $\mathcal{C}^\infty$ en 0.

\finenonce{004348}



\finexercice


\exercice{4349, quercia, 2010/03/12}
\enonce{004349}{$y''+y = f$}

Soit $f:\R \to \R$ continue.
On pose $g(x) =  \int_{t=0}^x f(t)\sin(x-t)\,d t$.
Montrer que $g$ est l'unique solution de l'équation différentielle : $y''+y = f(x)$
telle que $y(0) = y'(0) = 0$.

\finenonce{004349}



\finexercice


\exercice{4350, quercia, 2010/03/12}
\enonce{004350}{Fonction définie par une intégrale}

Soit $f:\R \to \R$ continue. On définit pour $x \in \R^*$ et $y \in \R$ :
$g(x,y) = \frac 1x  \int_{t=x}^{xy} f(t)\,d t$.

\begin{enumerate}
  \item Montrer que $g$ peut être prolongée en une fonction continue sur $\R^2$.
    
  \item On suppose de plus $f$ dérivable en 0. Montrer que $g$ est de classe $\mathcal{C}^1$.
    

\end{enumerate}
\finenonce{004350}



\finexercice


\exercice{4351, quercia, 2010/03/12}
\enonce{004351}{Fonctions définies par des intégrales}

Construire les courbes représentatives des fonctions suivantes :
\begin{enumerate}
  \item $f(x) =  \int_{t=-\pi/2}^{\pi/2} |x+t|\sin t\,d t$.\par
  \item $f(x) =  \int_{t=x}^{x^2} \frac{d t}{\ln t}$.
  \item $f(x) =  \int_{t=x}^{x^2} \frac{d t}{\sqrt{t^4+t^2+1}}$.
  \item $f(x) =  \int_{t=0}^1 \frac{e^t\,d t}{t+x}$.
  \item $f(x) =  \int_{t=0}^{\pi/2} x\exp\Bigl(\frac{\sin t}x\Bigr)d t$.
  \item $f(x) =  \int_{t=0}^x \frac{\ln(1+xt)}{1+t^2}\,d t$.
\end{enumerate}
\finenonce{004351}



\finexercice


\exercice{4352, quercia, 2010/03/12}
\enonce{004352}{Fonction définie par une intégrale}

Montrer qu'il existe un unique réel $x \in {[0,\pi]}$ tel que
$ \int_{\theta=0}^\pi \cos(x\sin\theta)\, d\theta = 0$.
Calculer une valeur approchée de $x$ à $10^{-2}$ près.


\finenonce{004352}



\finexercice


\exercice{4353, quercia, 2010/03/12}
\enonce{004353}{Développement en série, Ensam PSI 1998, Mines MP 1999}

Soit $I(\alpha) =  \int_{x=0}^{+\infty} \frac{\sin\alpha x}{e^x-1}\,d x$.

\begin{enumerate}
  \item Justifier l'existence de $I(\alpha)$.
  \item Déterminer les réels $a$ et~$b$ tels que~:
    $I(\alpha) = \sum_{n=1}^\infty \frac a{b+n^2}$.
    
  \item Donner un équivalent de $I(\alpha)$ quand $\alpha\to+\infty$.
    
\end{enumerate}
\finenonce{004353}



\finexercice


\exercice{4354, quercia, 2010/03/12}
\enonce{004354}{Formule de Stirling}
Montrer que $\Gamma(x+1) \sim x^xe^{-x}\sqrt{2\pi x}$ pour $x$ réel
tendant vers $+\infty$.

\finenonce{004354}



\finexercice


\exercice{4355, quercia, 2010/03/12}
\enonce{004355}{Développement en série, Mines 1999}

Soit $\theta\in{]0,\pi[}$.
Montrer que $ \int_{t=0}^1 \frac{d t}{e^{-i\theta}-t} = \sum_{n=1}^\infty \frac{\exp(in\theta)}n$.
\finenonce{004355}



\finexercice


\exercice{4356, quercia, 2010/03/12}
\enonce{004356}{Fonction définie par une intégrale, X 1999}

\begin{enumerate}
  \item Calculer $f(a) =  \int_{t=0}^{+\infty} e^{-t^2}\cos(at)\,d t$.
    
  \item Soit $g(a) =  \int_{t=0}^{+\infty} e^{-t^2}\frac{\sin(at)}t\,d t$~;
    calculer $\lim_{a\to+\infty} g(a)$.
    
\end{enumerate}
\finenonce{004356}



\finexercice


\exercice{4357, quercia, 2010/03/12}
\enonce{004357}{Développement asymptotique}

Soient $J(x) =  \int_{t=0}^{\pi/2}\frac{d t}{\sqrt{\sin^2t + x^2\cos^2t}}$
et $K(x) =  \int_{t=0}^{\pi/2}\frac{\cos t\,d t}{\sqrt{\sin^2t + x^2\cos^2t}}$.

Calculer $\lim_{x\to0^+}(J(x)-K(x))$ et montrer que $J(x) = -\ln x + 2\ln2 + o_{x\to0^+}(1)$.
\finenonce{004357}



\finexercice


\exercice{4358, quercia, 2010/03/12}
\enonce{004358}{Transformée de Laplace}

Soit $f : {[0,+\infty[} \to \R$ continue telle que $ \int_{t=0}^{+\infty}f(t)\,d t$ converge (pas forcément absolument).

On pose $\varphi(a) =  \int_{t=0}^{+\infty}e^{-at}f(t)\,d t$.

\begin{enumerate}
  \item Montrer que $\varphi$ est de classe $\mathcal{C}^\infty$ sur $]0,+\infty[$.
    
  \item Montrer que $\varphi$ est continue en $0$.
    
\end{enumerate}
\finenonce{004358}



\finexercice


\exercice{4359, quercia, 2010/03/12}
\enonce{004359}{} 

On pose pour $n\ge2$~: $v_n= \int_{x=0}^1 \frac 1{1+x^n}\, d x$.
Montrer que la suite $(v_n)$ converge. Nature de la série $\sum(v_n-1)$ ?

\finenonce{004359}



\finexercice


\exercice{4360, quercia, 2010/03/12}
\enonce{004360}{} 

On pose pour $n\ge2$ $u_n= \int_{x=0}^{+\infty} \frac1{1+x^n}\, d x$.
Montrer que la suite $(u_n)$ converge, puis que la série $\sum(u_n-1)$ converge
également.

\finenonce{004360}



\finexercice


\exercice{4361, quercia, 2010/03/12}
\enonce{004361}{Centrale MP 2000}
Domaine de définition de~$I(\alpha) =  \int_{x=0}^{+\infty}\frac{x\ln x}{(1+x^2)^\alpha}\,d x$.
Calculer $I(2)$ et~$I(3)$.
Déterminer la limite de $I(\alpha)$ en~$+\infty$.

\finenonce{004361}



\finexercice


\exercice{4362, quercia, 2010/03/12}
\enonce{004362}{Centrale MP 2000}
On considère $f(x)= \int_{t=0}^{+\infty}\frac{d t}{t^x(1+t)}$.

\begin{enumerate}
  \item Domaine de définition, monotonie, convexité de $f$ (sans dériver $f$).
    
  \item Continuité, dérivabilité, calcul de $f^{(k)}(x)$.
  \item Donner un équivalent de $f(x)$ en $0$ et en $1$.
    
  \item Calculer $f(1/n)$ pour $n\in \N$, $n\ge 2$.
    
\end{enumerate}
\finenonce{004362}



\finexercice


\exercice{4363, quercia, 2010/03/12}
\enonce{004363}{Ensae MP$^*$ 2000}

Soit $\alpha \in \R$. Trouver la limite de $u_n=\sum_{k=1}^n \frac{\sin (k\alpha)}{n+k}\cdotp$

\finenonce{004363}



\finexercice


\exercice{4364, quercia, 2010/03/12}
\enonce{004364}{Polytechnique MP$^*$ 2000}

Existence et continuité de $f(x) = \int_{t=-\infty}^{+\infty} \frac{e^{-|x+t|}\cos(x+t)}{\sqrt{|t|}(1+|t|)}\,d t$.
Montrer que $f$ est intégrable.
\finenonce{004364}



\finexercice


\exercice{4365, quercia, 2010/03/12}
\enonce{004365}{Centrale MP 2001}

\begin{enumerate}
  \item Développer, pour tout $x>0$, $s(x) =  \int_{t=0}^{+\infty}\frac{\sin t}{e^{xt}-1}\,d t$
    en série de fractions rationnelles.
  \item Montrer qu'en $0^+$, $s(x)$ est équivalente à $\frac\pi{2x}$.
\end{enumerate}
\finenonce{004365}



\finexercice


\exercice{4366, quercia, 2010/03/12}
\enonce{004366}{X MP$^*$ 2001}
\'Etudier $ \int_{t=0}^{+\infty} \frac{e^{-tx}}{1+t^2}\,d t$.

\finenonce{004366}



\finexercice


\exercice{4367, quercia, 2010/03/12}
\enonce{004367}{Ensi MP 2004}
\smallskip
Soit $f(x) =  \int_{t=0}^{+\infty}\frac{e^{-t^2x}}{1+t^2}\,d t$.

\begin{enumerate}
  \item Trouver le domaine de définition de~$f$.
  \item Montrer que $f$ est dérivable sur $\R^{+*}$.
  \item Calculer $f-f'$.
  \item Donner un équivalent simple de~$f'(x)$ pour $x\to+\infty$.
  \item Montrer que $f(x) = \frac{\sqrt\pi}{2\sqrt x} - \frac{\sqrt\pi}{4x\sqrt x} +  o\Bigl(\frac1{x\sqrt x}\Bigr)$.
  \item Tracer la courbe de~$f$.
\end{enumerate}
\finenonce{004367}



\finexercice


\exercice{4368, quercia, 2010/03/12}
\enonce{004368}{Ensea MP 2004}
Soit $\alpha>0$.

\begin{enumerate}
  \item Montrer que $f : x \mapsto e^{-\alpha x} \int_{\theta=0}^\pi \cos(x\sin\theta)\, d\theta$ est intégrable
    sur $\R^+$.
  \item Calculer $I =  \int_{x=0}^{+\infty}f(x)\,d x$.
    Indication~: écrire $I = \lim_{a\to+\infty} \int_{x=0}^a f(x)\,d x$.
    
\end{enumerate}
\finenonce{004368}



\finexercice


\exercice{4369, quercia, 2010/03/12}
\enonce{004369}{X MP$^*$ 2000}

Étudier la limite en $0+$ de $I(x) =  \int_{t=0}^{+\infty}
\frac{e^{-t}-\cos t}{t}e^{-xt}\, d t$.

\finenonce{004369}



\finexercice


\exercice{4370, quercia, 2010/03/12}
\enonce{004370}{$\zeta$ et $\Gamma$}
Montrer, pour $x>1$~: $\zeta(x)\Gamma(x) =  \int_{t=0}^{+\infty} \frac{t^{x-1}}{e^t-1}\,d t$.
\finenonce{004370}



\finexercice


\exercice{4371, quercia, 2010/03/12}
\enonce{004371}{Centrale MP 2002}

Soit $f$ : $x \mapsto  \int_{t=0}^{+\infty}\frac{d t}{t^{x+1}+t+1}$.
Déterminer son domaine de définition~; étudier sa continuité et sa monotonie.
Calculer $ \int_{t=1}^{+\infty}\frac{d t }{t^{x+1}+t}$ et en déduire des
équivalents et les limites de~$f$ en~$0$ et en~$+\infty$.
\finenonce{004371}



\finexercice


\exercice{4372, quercia, 2010/03/12}
\enonce{004372}{Polytechnique MP 2002}

Soit $\alpha\in{]0,\frac\pi2[}$ et $\lambda\in\R$.
Chercher un équivalent pour $n\to\infty$ de
$I_n =  \int_{x=0}^{\alpha}\sin(x)\exp(\lambda n\sin^2(x))\,d x$.

\finenonce{004372}



\finexercice


\exercice{4373, quercia, 2010/03/12}
\enonce{004373}{Centrale MP 2004}
Soit~$(a_n)_{n\in\N}$ la suite définie par $a_0=1$ et $a_n =
\frac1{n!} \int_{t=0}^1t(t-1)\dots(t-n)\,d t$.

\begin{enumerate}
  \item Quel est le rayon de convergence de la série entière $\sum_{n=0}^\infty
    a_nx^n$~?
    
  \item Donner un équivalent de~$a_n$.
    
\end{enumerate}
\finenonce{004373}



\finexercice


\exercice{4374, quercia, 2010/03/12}
\enonce{004374}{$\sum_{n=1}^\infty \frac{\sin(nx)}n$, Mines-Ponts MP 2004}

Soit $u_n(t)$ le terme général d'une série~: $u_n(t)= t^{n-1}\sin(nx)$
avec $0< x< \pi$.

\begin{enumerate}
  \item \'Etudier la convergence de la série.


  \item Calculer $\sum_{p=0}^n u_p(t) = S_n(t)$.
    Mettre $S_n(t)$ sous la forme $\frac{P_n(t)}{Q(t)}$ avec $Q(t)>\alpha$, $\alpha>0$.


  \item Calculer $\lim_{n\to\infty} S_n(t)$ et $\lim_{n\to\infty} \int_{t=0}^1 S_n(t)\,d t$.


  \item En déduire $\sum_{n=1}^\infty \frac{\sin(nx)}n$.

\end{enumerate}
\finenonce{004374}



\finexercice


\exercice{4375, quercia, 2010/03/12}
\enonce{004375}{Lemme de Lebesgue, Centrale MP 2004}
Soit $f$ continue par morceaux définie sur $\R$, à valeurs dans $\C$.

\begin{enumerate}
  \item Soient $a,b\in\R$.
Montrer que $ \int_{t=a}^b f(t) \cos(nt)\,d t\to0$ lorsque $n\to\infty$.

  \item On suppose que $f$ est intégrable sur $]0,+\infty[$.
Soit $u_n=  \int_{t=0}^{n\pi}  \sin^2(nt) f(t)\,d t $.
Montrer que $(u_n)_{n \in \N}$ admet une limite quand $n\to\infty$ et la préciser.

\end{enumerate}
\finenonce{004375}



\finexercice


\exercice{4376, quercia, 2010/03/12}
\enonce{004376}{Suite d'intégrales, Centrale MP 2004}
Soit $(f_n)_{n\in\N^*}$ une suite de fonctions définie par~:
$\forall\ n\in\N^*,\ \forall\ x\in{[0,1]},\ f_n(x) = \Bigl(\frac{x+x^n}2\Bigr)^n$.

\begin{enumerate}
  \item Montrer que~$(f_n)$ converge simplement vers une fonction~$\varphi$.
    
  \item
  \begin{enumerate}  
    \item La convergence est-elle uniforme~?
   
   \item La convergence est-elle monotone~?
   
  \end{enumerate} 
  \item Soit, pour~$n\in\N^*$, $J_n =  \int_{x=0}^1f_n(x)\,d x$.
    Montrer que~$J_n\sim \frac2{n^2}$.
     
\end{enumerate}
\finenonce{004376}



\finexercice


\exercice{4377, quercia, 2010/03/12}
\enonce{004377}{Mines-Ponts MP 2004}
\smallskip
Soit $f(x) =  \int_{t=0}^1\frac{1-t}{\ln t}t^x\,d t$.
Étudier le domaine de définition de~$f$, sa dérivabilité, puis calculer $f(x)$.


\finenonce{004377}



\finexercice


\exercice{4378, quercia, 2010/03/12}
\enonce{004378}{Mines-Ponts MP 2004}
Soit $I(a) =  \int_{x=0}^{+\infty}\frac{\sh x}xe^{-ax}\,d x$.

\begin{enumerate}
  \item Quel est le domaine de définition de~$I$~?

  \item Étudier la continuité et la dérivabilité de~$I$.
  \item Calculer $I(a)$.

\end{enumerate}
\finenonce{004378}



\finexercice


\exercice{4379, quercia, 2010/03/12}
\enonce{004379}{ENS Lyon MP${}^*$ 2004}
\begin{enumerate}
  \item Soit $f : {[a,b]} \to \R$ de classe $\mathcal{C}^2$ et $F(x) =  \int_{t=a}^b
    f(t)e^{-itx}\,d t$ avec $a<0<b$. Montrer que $F(x)\to0$ lorsque $x\to+\infty$.
    
  \item Montrer que $F(x) = \frac{f(a)e^{-iax} - f(b)e^{-ibx}}{ix} +  o\Bigl(\frac1x\Bigr)$.
    
  \item Montrer la convergence de l'intégrale
    $I= \int_{t=-\infty}^{+\infty}e^{-it^2/2}\,d t$.
    
  \item Soit $g(x) =  \int_{t=a}^bf(t)e^{-ixt^2/2}\,d t$.
    Montrer que $g(x) = \frac{I.f(0)}{\sqrt x} +  o\Bigl(\frac1{\sqrt x}\Bigr)$.
    
\end{enumerate}
\finenonce{004379}



\finexercice


\exercice{4380, quercia, 2010/03/12}
\enonce{004380}{Théorème de d'Alembert-Gauss}

Soit $P\in \C[X]$ de degré $n\ge 1$.
Le but de cet exercice est de prouver que $P$ admet une racine dans~$\C$.
On suppose au contraire que $P$ ne s'annule pas et on considère pour
$r \ge 0$, $\theta\in{[0,2\pi]}$~:
$f(r,\theta) = \frac{r^ne^{in\theta}}{P(re^{i\theta})}$
et $F(r) =  \int_{\theta=0}^{2\pi} f(r,\theta)\, d\theta$.

\begin{enumerate}
  \item Montrer que $F$ est de classe $\mathcal{C}^1$ sur $[0,+\infty[$.
  \item Vérifier que $ir \frac{\partial f}{\partial r} = \frac{\partial f}{\partial \theta}$. En déduire que $F$ est constante.
  \item Obtenir une contradiction.
\end{enumerate}
\finenonce{004380}



\finexercice


\exercice{5453, rouget, 2010/07/10}
\enonce{005453}{****}
Soit $f$ une fonction continue sur $[a,b]$. Pour $x$ réel, on pose $F(x)=\int_{a}^{b}|t-x|f(t)\;dt$. Etudier la dérivabilité de $F$ sur $\Rr$.
\finenonce{005453}


\finexercice


\exercice{5460, rouget, 2010/07/10}
\enonce{005460}{***}
Etude complète de $F(x)=\int_{x}^{2x}\frac{dt}{\sqrt{t^4+t^2+1}}$. 
\finenonce{005460}


\finexercice

\exercice{5464, rouget, 2010/07/10}
\enonce{005464}{***I}
\begin{itemize}
\item  Déterminer $\lim_{x\rightarrow 1}\int_{x}^{x^2}\frac{dt}{\ln t}$. 
\item  Etude complète de $F(x)=\int_{x}^{x^2}\frac{dt}{\ln t}$.
\end{itemize}
\finenonce{005464}


\finexercice

\exercice{5472, rouget, 2010/07/10}
\enonce{005472}{}
Etude de $f(x)=\int_{-1}^{1}\frac{\sin x}{1-2t\cos x+t^2}\;dt$.  
\finenonce{005472}


\finexercice


\exercice{5473, rouget, 2010/07/10}
\enonce{005473}{}
Etude de $f(x)=\int_{0}^{1}\mbox{Max}(x,t)\;dt$. 
\finenonce{005473}


\finexercice


\exercice{5765, rouget, 2010/10/16}
\enonce{005765}{**}
Pour $n\in\Nn^*$ et $x\in]0,+\infty[$, on pose $I_n(x)=\int_{0}^{+\infty}\frac{1}{(t^2+x^2)^n}$.

\begin{enumerate}
 \item  Calculer la dérivée de la fonction $I_n$ sur $]0,+\infty[$.

\item  En déduire la valeur de $\int_{0}^{+\infty}\frac{1}{(t^2+1)^3}\;dt$.
\end{enumerate}
\finenonce{005765}


\finexercice


\exercice{5766, rouget, 2010/10/16}
\enonce{005766}{*** I (très long) Intégrale de \textsc{Poisson}}
Pour $x\in\Rr$, on pose $F(x)=\int_{-\pi}^{\pi}\ln(1-2x\cos\theta+x^2)\;d\theta$.

\begin{enumerate}
 \item  
  \begin{enumerate}
   \item Montrer que $F$ est paire, définie et continue sur $\Rr$, dérivable sur $\Rr\setminus\{-1,1\}$.
Préciser une expression de $F'(x)$ sous forme intégrale.
   \item Calculer $F'(x)$.
   \item Déterminer $\lim_{x \rightarrow +\infty}(F(x)-4\pi\ln x)$.
   \item En déduire $F(x)$ pour tout réel $x$.
   \end{enumerate}

\item
  \begin{enumerate}
   \item Quand $x\in]-1,1[$, retrouver ce résultat en écrivant d'abord $\ln(x^2-2x\cos\theta+1)$ comme somme d'une série (commencer par dériver la fonction de $\theta$).

   \item Déterminer une relation entre $F(x)$ et $F\left(\frac{1}{x}\right)$ et en déduire $F(x)$ pour tout réel $x$.
  \end{enumerate}
\end{enumerate}
\finenonce{005766}


\finexercice


\exercice{5767, rouget, 2010/10/16}
\enonce{005767}{** I Un calcul de l'intégrale de \textsc{Gauss} $I=\int_{0}^{+\infty}e^{-t^2}\;dt$}

Pour $x\in\Rr$, on pose $F(x)=\int_{0}^{1}\frac{e^{-x^2(1+t^2)}}{1+t^2}\;dt$ et $G(x)=\left(\int_{0}^{x}e^{-t^2}\;dt\right)^2$.

\begin{enumerate}
 \item  Montrer que $F$ est de classe $C^1$ sur $\Rr$ et préciser $F'$.

\item  Montrer que $G$ est de classe $C^1$ sur $\Rr$ et préciser $G'$.

\item  Montrer que la fonction $F+G$  est constante sur $\Rr$.

\item  Déterminer $\lim_{x \rightarrow +\infty}F(x)$.

\item  En déduire $I$.
\end{enumerate}
\finenonce{005767}


\finexercice


\exercice{5768, rouget, 2010/10/16}
\enonce{005768}{***}
Existence et calcul de $\int_{0}^{+\infty}e^{-t^2}\ch(tx)\;dt$ (on admettra que $\int_{0}^{+\infty}e^{-t^2}\;dt=\frac{\sqrt{\pi}}{2}$).

\finenonce{005768}


\finexercice


\exercice{5769, rouget, 2010/10/16}
\enonce{005769}{***}
Existence et calcul de $\int_{0}^{1}\frac{t-1}{\ln t}t^x\;dt$.
\finenonce{005769}


\finexercice


\exercice{5770, rouget, 2010/10/16}
\enonce{005770}{**** I (très long)}
Montrer que pour tout réel $x$ strictement positif, $\int_{0}^{+\infty}\frac{e^{-xt}}{1+t^2}\;dt =\int_{0}^{+\infty}\frac{\sin t}{x+t}\;dt$ et en déduire $\int_{0}^{+\infty}\frac{\sin t}{t}\;dt$
(indication : trouver une équation différentielle du second ordre vérifiée par ces deux fonctions).
\finenonce{005770}


\finexercice


\exercice{5771, rouget, 2010/10/16}
\enonce{005771}{** I (Produit de convolution)}
\begin{enumerate}
 \item  Soient $f$ et $g$ deux fonctions définies sur $\Rr$, continues et $T$-périodiques ($T$ réel strictement positif). Pour $x\in\Rr$, on pose

\begin{center}
$f*g(x)=\int_{0}^{T}f(x-t)g(t)\;dt$.
\end{center}

Montrer que la fonction $f*g$ est définie sur $\Rr$, continue et $T$-périodique.

\item  $*$ est donc une loi interne sur $E$, l'espace vectoriel des fonctions définies et continues sur $\Rr$ et $T$-périodiques. Montrer que cette loi est commutative.

\end{enumerate}
\finenonce{005771}


\finexercice

\section{224.04 Tranformée de Laplace et transformée de Fourier}

\section{224.99 Autre}

\section{225.01 Résolution d'équation différentielle du premier ordre}


\exercice{845, bodin, 1998/09/01}

\enonce{000845}{}
R\'esoudre les \'equations diff\'erentielles suivantes :
\begin{enumerate}
    \item  $y'=y+x$ avec $y(0)=1$,
    \item  $y'= \cos x + y$,
    \item  $y'+2y = (x-2)^2$.
\end{enumerate}
\finenonce{000845}



\finexercice



\exercice{846, bodin, 1998/09/01}

\enonce{000846}{}
Pour chacune des \'equations diff\'erentielles qui suit : \'ecrire la solution
passant par le point M(.,.) et tracer sommairement le graphe de la solution.
\begin{enumerate}
    \item  $y'+2xy=0,\ \ \ M=(0,1) $,
    \item  $y'+y \tan x = \sin x \cos x \ \ \ M=(\frac{\pi}{4},0)$,
    \item  $x(x^2-1)y'+2y=x^2$, On d\'eterminera $a,b,c \in \Rr$ tels que
$\frac{1}{x(x^2-1)}=\frac{a}{x}+\frac{b}{x-1}+\frac{c}{x+1}$.
\end{enumerate}
\finenonce{000846}



\finexercice



\exercice{847, bodin, 1998/09/01}
\enonce{000847}{}
On se propose d'int\'egrer sur l'intervalle le plus grand possible contenu dans $]0,
\infty[$ l'\'equation diff\'erentielle :

$$
(E) \ \ \ \ \ \ \ y'(x)- \frac{y(x)}{x}-y(x)^2=-9x^2.
$$
\begin{enumerate}
    \item D\'eterminer $ a \in ]0, \infty[ $ tel que $y(x)=ax$ soit une solution
particuli\`ere
$y_0$ de $(E)$.
    \item  Montrer que le changement de fonction inconnue : $y(x) = y_0(x)-\frac{1}{z(x)}$
transforme l'\'equation (E) en l'\'equation diff\'erentielle
$$
(E_1) \ \ \ \ \ \ \ z'(x) + (6x+ \frac{1}{x})z(x) =1.
$$
    \item Int\'egrer $(E_1)$ sur $]0,\infty[$.
    \item Donner toutes les solutions de $(E)$ d\'efinies sur $]0,\infty[$.
\end{enumerate}
\finenonce{000847} 
\finexercice


\exercice{848, bodin, 1998/09/01}

\enonce{000848}{}
Trouver les solutions r\'eelles des \'equations diff\'erentielles
suivantes :
\begin{enumerate}
    \item $y'(t)+2y(t) = 0$ ;
    \item $\displaystyle{\frac{dx}{dt}-x=0}$ ;
    \item $y'(x)+2y(x) = 0 \text{ avec } (y-y')(0)=0$.
\end{enumerate}
\finenonce{000848}



\finexercice



\exercice{849, bodin, 1998/09/01}

\enonce{000849}{}
Trouver les solutions r\'eelles des \'equations  diff\'erentielles
suivantes :
\begin{enumerate}
    \item $(1+x^2)y'-xy = 0$ ;
    \item $y'+y\tan x = 0$, pour $x$ dans
$]\frac{\pi}{2},\frac{3\pi}{2}[$.
\end{enumerate}
\finenonce{000849}



\finexercice



\exercice{850, bodin, 1998/09/01}

\enonce{000850}{}
Trouver les solutions r\'eelles sur l'intervalle maximal
de l'\'equation diff\'erentielle :
$$t^2y'+y = 1.$$
\finenonce{000850}



\finexercice



\exercice{851, bodin, 1998/09/01}

\enonce{000851}{}
Soit l'\'equation diff\'erentielle
$$(E)\qquad y'+2xy = x.$$
\begin{enumerate}
    \item R\'esoudre l'\'equation homog\`ene asoci\'ee.
    \item Calculer la solution de $(E)$ v\'erifiant $y(0)=1$.
\end{enumerate}

\finenonce{000851}



\finexercice



\exercice{852, ridde, 1999/11/01}

\enonce{000852}{}
R\'esoudre et raccorder \'eventuellement :
\begin{enumerate}
\item $xy'-2y = x^4$.
\item $x (1 + x^2)y' = y$.
\item $ (x^2 + 1)y' + (x-1)^2y = x^3-x^2 + x + 1$.
\item $ (e^x-1)y' + (e^x + 1)y = 3 + 2e^x$.
\end{enumerate}
\finenonce{000852}



\finexercice



\exercice{853, ridde, 1999/11/01}

\enonce{000853}{}
R\'esoudre le syst\`eme diff\'erentiel :
$\begin{cases} \dot{x} (t) = x (t) + y (t) \\ \dot{y} (t) = 3x (t)-y (t) \end{cases}$
et $ \begin{cases} x (0) = 2 \\ y (0) = -2 \end{cases}$.
\finenonce{000853}




\finexercice


\exercice{854, ridde, 1999/11/01}

\enonce{000854}{}
R\'esoudre l'\'equation diff\'erentielle de Ricatti $x^2y' = x^2y^2 + xy + 1$ en trouvant
une solution particuli\`ere $y_0$ et en posant $z = \frac 1{y-y_0}$.
\finenonce{000854}



\finexercice



\exercice{855, gourio, 2001/09/01}

\enonce{000855}{}
Soit l'\'{e}quation diff\'{e}rentielle :
$$(E):\frac{dy(x)}{dx}+y(x)=x^{2}+2x $$
Int\'{e}grer $(E)$ et montrer que par un point donn\'{e} il passe une et une
seule courbe int\'{e}grale. Soit $H  $ l'ensemble des points $M  $ tels
que la courbe int\'{e}grale passant par $M$ a une tangente horizontale en ce
point, et $I   $ l'ensemble des points $M$ tels que la courbe
int\'{e}grale passant par ce point a un point d'inflexion en ce point.\
Tracer $H, I $ et la courbe int\'{e}grale passant par $O(0,0). $ En
d\'{e}duire un trac\'{e} g\'{e}om\'{e}trique des courbes int\'{e}grales.
\finenonce{000855}



\finexercice



\exercice{856, ruette, 2013/02/01}

\enonce{000856}{}
On consid\`ere l'\'equation diff\'erentielle suivante :
$$
(E)\qquad y'-(1-\tan(x))y=\cos(x).
$$
On s'int\'eresse \`a cette \'equation sur l'intervalle $]-\frac{\pi}{2},\frac{\pi}{2}[$.


\begin{enumerate}
\item
\begin{enumerate}
\item \'Ecrire l'\'equation homog\`ene associ\'ee \`a (E), puis  r\'esoudre cette \'equation 
sur l'intervalle  $]-\frac{\pi}{2},\frac{\pi}{2}[$.
\item  Quelle est  la solution de l'\'equation homog\`ene v\'erifiant la condition 
initiale $y(0)=1$ ? On la note $g_{0}$. Donner le tableau de variations de $g_{0}$ sur l'intervalle 
$]-\frac{\pi}{2},\frac{\pi}{2}[$. D\'eterminer les limites aux bornes de l'intervalle. 
Tracer  le graphe de $g_{0}$.
\item
Tracer sur le m\^eme dessin, rapidement, les graphes des trois  solutions de l'\'equation homog\`ene  v\'erifiant respectivement  $y(0) = 2$, $y(0)=3$ , $y(0) = -1$.
\end{enumerate}
%

\item 
R\'esoudre compl\`etement l'\'equation $(E)$. 

%%%%%%%%%%%%%%%
\item %variante
 R\'esoudre  de m\^eme, sur l'intervalle $]-\frac{\pi}{2},\frac{\pi}{2}[$,  
l'\'equation diff\'erentielle suivante :
$$
y'-(1-\tan(x))y=\cos^2(x).
$$
\textit{Indication : on pourra calculer la primitive en  utilisant l'exponentielle complexe, ou des int\'egrations par parties.}


\item %superposition
  Quelles sont les solutions de l'\'equation diff\'erentielle suivante  sur l'intervalle $]-\frac{\pi}{2},\frac{\pi}{2}[$ ?
$$
y'-(1-\tan(x))y=\cos(x) + \cos^2(x).
$$
\end{enumerate}
\finenonce{000856}



\finexercice

\exercice{857, gourio, 2001/09/01}

\enonce{000857}{}
Soit $f\in C^{1}({\Rr},\Cc),\alpha \in {\Rr}^{+*}. $ Montrer que
si :
$$\lim\limits_{x\rightarrow \infty }\left( f^{\prime }(x)+\alpha f(x)\right)=0 $$
alors :
$$\lim_{x\rightarrow \infty }f^{\prime }(x)=\lim_{x\rightarrow \infty }f(x)=0. $$
\finenonce{000857}



\finexercice



\exercice{858, gourio, 2001/09/01}

\enonce{000858}{}
Soit $f \in C^{1}({\Rr},{\Rr})$ telle que $f(0)=1 $ et $f\leq f^{\prime
}\leq 2f. $Encadrer $f(-1)$ et $f(1).$
\finenonce{000858}



\finexercice



\exercice{5476, rouget, 2010/07/10}
\enonce{005476}{**IT}
Résoudre sur l'intervalle $I$ de $\Rr$ proposé les équations différentielles suivantes~:

$$\begin{array}{ll}
1)\;x\ln xy'+y=x,\;I=]1,+\infty[&2)\;x(xy'+y-x)=1,\;I=]-\infty,0[\\
3)\;2xy'+y=x^4,\;I=]-\infty,0[&4)\;y'+2y=x^2-3x,\;I=\Rr\\
5)\;y'+y=\frac{1}{1+2e^x},\;I=\Rr&6)\;y'\sin x-y\cos x+1=0,\;I=]0,\pi[
\end{array}$$
\finenonce{005476}


\finexercice


\exercice{5477, rouget, 2010/07/10}
\enonce{005477}{***I}
Résoudre l'équation différentielle $(1-x^2)y'-2xy=x^2$ sur chacun des intervalles $I$
suivants~:~$I=]1,+\infty[$, $I=]-1,1[$, $I=]-1,+\infty[$, $I=\Rr$.
\finenonce{005477}


\finexercice


\exercice{5478, rouget, 2010/07/10}
\enonce{005478}{***}
Résoudre sur $\}-\infty,0[$ et $]0,+\infty[$ l'équation différentielle~:~$|x|y'+(x-1)y=x^3$.
\finenonce{005478}


\finexercice


\exercice{5481, rouget, 2010/07/10}
\enonce{005481}{**}
Soit $a$ un réel non nul. Soit $f$ continue sur $\Rr$ et pèriodique de pèriode $T\neq0$. Montrer que l'équation
différentielle $y'+ay=f$ admet une et une seule solution périodique sur $\Rr$, de période $T$.
\finenonce{005481}


\finexercice


\exercice{6991, blanc-centi, 2015/07/04}
\video{aFS_Z30s_RY}
\enonce{006991}{}
Résoudre sur $\R$ les équations différentielles suivantes: 
\begin{enumerate}
\item $y'+2y=x^2$ $(E_1)$
\item $y'+y=2\sin x$ $(E_2)$
\item $y'-y=(x+1)e^x$ $(E_3)$
\item $y'+y=x-e^x+\cos x$ $(E_4)$
\end{enumerate}
\finenonce{006991}
 

\finexercice


\exercice{6992, blanc-centi, 2015/07/04}
\video{hVq0LPY0kV0}
\enonce{006992}{}
Déterminer toutes les fonctions $f:[0;1]\to\R$, dérivables, telles que 
$$\forall x\in[0;1],\ f'(x)+f(x)=f(0)+f(1)$$
\finenonce{006992}
 

\finexercice


\exercice{6993, blanc-centi, 2015/07/04}
\video{Fg3TtiARvxI}
\enonce{006993}{}

\begin{enumerate}
\item Résoudre l'équation différentielle $(x^2+1)y'+2xy=3x^2+1$ sur $\R$.
Tracer des courbes intégrales. Trouver la solution vérifiant $y(0) = 3$.


\item Résoudre l'équation différentielle $y'\sin x-y\cos x+1=0$ sur $]0;\pi[$.
Tracer des courbes intégrales. Trouver la solution vérifiant $y(\frac\pi4) = 1$.

\end{enumerate}
\finenonce{006993}
 

\finexercice


\exercice{6994, blanc-centi, 2015/07/04}
\video{4M0txOmYb1I}
\enonce{006994}{Variation de la constante}
Résoudre les équations différentielles suivantes en trouvant 
une solution particulière par la méthode de variation de la constante :
\begin{enumerate}
\item $y' - (2x - \frac{1}{x})y = 1$ sur $]0;+\infty[$
\item $y'-y = x^k \exp(x)$ sur $\R$, avec $k \in \Nn$
\item $x(1+\ln^2(x))y'+2\ln(x)y=1$ sur $]0;+\infty[$
\end{enumerate}
\finenonce{006994}
 

\finexercice


\exercice{6995, blanc-centi, 2015/07/04}
\video{UW3uUaXMS34}
\enonce{006995}{}
On considère l'équation différentielle
$$y'-e^xe^y=a$$
Déterminer ses solutions, en précisant soigneusement leurs intervalles de définition, pour
\begin{enumerate}
\item $a=0$
\item $a=-1$ (faire le changement de fonction inconnue $z(x)=x+y(x)$)
\end{enumerate}
Dans chacun des cas, construire la courbe intégrale qui passe par l'origine.
\finenonce{006995}
 

\finexercice


\exercice{6996, blanc-centi, 2015/07/04}
\video{1-v-EhSJHlk}
\enonce{006996}{}
Pour les équations différentielles suivantes,
trouver les solutions définies sur $\R$ tout entier :
\begin{enumerate}
\item $x^2y'-y = 0$ $(E_1)$
\item $xy'+y-1 = 0$ $(E_2)$
\end{enumerate}
\finenonce{006996}
 

\finexercice


\exercice{7001, blanc-centi, 2015/07/04}
\video{k9R8_U4-Pyk}
% extrait de 4112, 4113
\enonce{007001}{\'Equations de Bernoulli et Riccatti}
\ 
\begin{enumerate}
  \item \textbf{\'Equation de Bernoulli}
  \begin{enumerate}
    \item Montrer que l'équation de Bernoulli
    $$y'+a(x)y+b(x)y^n = 0 \qquad n \in \Zz \quad n \neq 0, n\neq 1$$
    se ramène à une équation linéaire par le changement de fonction 
    $z(x) = 1/ y(x)^{n-1}$.
    
    
    \item Trouver les solutions de l'équation $xy'+y-xy^3 = 0$.
    
  \end{enumerate}
    \item \textbf{\'Equation de Riccati}
  \begin{enumerate}
    \item Montrer que si $y_0$ est une solution particulière de l'équation de Riccati 
    $$y'+ a(x)y +b(x)y^2 = c(x)$$
    alors la fonction définie par $u(x) = y(x)-y_0(x)$ vérifie une équation de Bernoulli
    (avec $n=2$).
    
    \item Résoudre $x^2(y'+y^2) = xy-1$ en vérifiant d'abord que $y_0(x) = \frac1x$ est une solution.
  \end{enumerate}

\end{enumerate}
\finenonce{007001}


\finexercice


\exercice{7002, blanc-centi, 2015/07/04}
\video{xrbDXpFKgfE}
\enonce{007002}{}
\ 
\begin{enumerate}
\item Montrer que toute solution sur $\R$ de $y'+e^{x^2}y=0$ tend vers 0 en $+\infty$.
\item Montrer que toute solution sur $\R$ de $y''+e^{x^2}y=0$ est bornée.
(\emph{Indication :} étudier la fonction auxiliaire $u(x)=y(x)^2+e^{-x^2}y'(x)^2$.)
\end{enumerate}
\finenonce{007002}


\finexercice

\section{225.02 Résolution d'équation différentielle du deuxième ordre}


\exercice{859, bodin, 1998/09/01}

\enonce{000859}{}
R\'esoudre les \'equations diff\'erentielles du second ordre suivantes :
\begin{enumerate}
    \item  $y''+4y'+3y=0$,
    \item  $y''-6y'+9y=0$,
    \item  $y''-2y'+2y=0$.
\end{enumerate}
\finenonce{000859}



\finexercice



\exercice{860, bodin, 1998/09/01}

\enonce{000860}{}
R\'esoudre les \'equations diff\'erentielles du second ordre suivantes :
\begin{enumerate}
    \item  $y''-y=x^3+x^2$,
    \item  $y''-2y'+y=e^x$,
    \item  $y''-2y'+y= \cos (mx)$ où $m \in \Rr$,
    \item  $y''-2y'+y=x^3e^x+2\cos x+ (x^3+3)$ (utiliser le principe de
superposition).
\end{enumerate}
\finenonce{000860}



\finexercice



\exercice{861, bodin, 1998/09/01}

\enonce{000861}{}
On consid\`ere l'\'equation homog\`ene $(E) \ \ ay''+by'+cy=0$, avec $a \neq 0$. Donner des
conditions necessaires et suffisantes liant les coefficients $a,b$ et $c$ dans les deux cas
suivants :\\
(i) toutes les solutions de $(E)$ tendent vers $0$ lorsque $x$ tend vers
l'infini ;\\
(ii) toutes les solutions sont p\'eriodiques.
\finenonce{000861}



\finexercice



\exercice{862, bodin, 1998/09/01}

\enonce{000862}{}
 R\'esoudre l'\'equation :
$$ y^{\prime\prime} +k^2y = \cos{mx}, \quad k,m \in \R.$$
On discutera suivant les valeurs de $k$ et $m$.
\finenonce{000862}



\finexercice



\exercice{863, bodin, 1998/09/01}
\enonce{000863}{}
R\'esoudre l'\'equation suivante :
$$ y^{\prime\prime}- 3y^\prime + 2 y = e^x. $$
\finenonce{000863} 
\finexercice


\exercice{864, bodin, 1998/09/01}
\enonce{000864}{}
R\'esoudre l'\'equation suivante :
 $$ y^{\prime\prime}- y =-6\cos x + 2x\sin x. $$
\finenonce{000864} 
\finexercice


\exercice{865, bodin, 1998/09/01}
\enonce{000865}{}
R\'esoudre l'\'equation suivante :
$$4y^{\prime\prime}+4y^\prime + 5y = \sin x e^{-x/2} .$$
\finenonce{000865} 
\finexercice


\exercice{866, bodin, 1998/09/01}
\enonce{000866}{}
 On consid\`ere l'\'equation  :
$$ y^{\prime\prime} + 2y^\prime + 4y = xe^x \qquad (E) $$
\begin{enumerate}
    \item  R\'esoudre l'\'equation diff\'erentielle homog\`ene associ\'ee \`a $(E)$.
    \item  Trouver une solution particuli\`ere de $(E)$ (expliquer votre
d\'emarche), puis donner l'ensemble de toutes les solutions de
$(E)$.
    \item  D\'eterminer l'unique solution $h$ de $(E)$ v\'erifiant
$h(0)=1$ et $h(1)=0$.
    \item  Soit $f : ]0,\infty [ \longrightarrow \R$
une fonction deux fois d\'erivable sur $]0,\infty [ $ et qui v\'erifie :
$$ t^2f^{\prime \prime}(t) +3tf^\prime (t) + 4f(t) = t\log{t}. $$
    \begin{enumerate}
        \item  On pose $g(x)=f(e^x)$, v\'erifier que $g$ est solution de $(E)$.
        \item  En d\'eduire une expression de $f$.
    \end{enumerate}
\end{enumerate}

\finenonce{000866} 


\finexercice


\exercice{867, bodin, 1998/09/01}

\enonce{000867}{}
 Soit $m \in \R$. D\'eterminer la solution de l'\'equation :
$$ (E_m) \qquad y^{\prime\prime} - 2y^\prime + (1+m^2)y = (1+4m^2)\cos{mx}$$
qui v\'erifie $y(0)=1$ et $y^\prime (0)=0$ (Indication : On traitera
s\'eparement les cas $ m=0$ et $ m \not = 0$).
\finenonce{000867}



\finexercice



\exercice{868, bodin, 1998/09/01}

\enonce{000868}{}
On consid\`ere l'\'equation diff\'erentielle :
$$y''+6y'+9y=d(x) \quad (E)$$
\begin{enumerate}
    \item R\'esoudre l'\'equation diff\'erentielle
homog\`ene associ\'ee \`a $(E)$.
    \item Trouver une solution particuli\`ere de $(E)$
lorque, respectivement, on pose :
$$ d(x) = (x^2+1)e^{-3x} \ \ \text{et}\ \ d(x)=\cos x.$$
    \item Donner la forme g\'en\'erale des solutions de $(E)$
lorsque :
$$d(x) = 2(x^2+1)e^{-3x}+50\cos x.$$
\end{enumerate}
\finenonce{000868}



\finexercice



\exercice{869, bodin, 1998/09/01}

\enonce{000869}{}
D\'eterminer une \'equation diff\'erentielle v\'erifi\'ee
par la famille de fonctions
$$y(x) = C_1e^{2x}+C_2e^{-x} \quad C_1,C_2 \in \Rr.$$
\finenonce{000869}



\finexercice



\exercice{870, bodin, 1998/09/01}

\enonce{000870}{}
D\'eterminer une \'equation diff\'erentielle admettant $(r-2)^2=0$
comme \'equation caract\'eristique et $e^x+(x^3/6)e^{2x}$
comme solution particuli\`ere.
\finenonce{000870}



\finexercice



\exercice{871, bodin, 1998/09/01}

\enonce{000871}{}
 D\'eterminer l'ensemble des solutions r\'eelles des  \'equations :\\
a) $y^{\prime\prime}+ y^\prime -6y=e^{3x}$, \hspace{1.4cm}
b) $y^{\prime\prime}+ y^\prime -6y=e^x(2x+1)$,\\
c) $y^{\prime\prime}-4y^\prime +13y=\cos{x}$,\hspace{1cm}
d) $y^{\prime\prime}+3y^\prime +2y=e^{-2x}(x+1)$ avec $y(0)=1$, $y(1)=0$.
\finenonce{000871}



\finexercice



\exercice{872, bodin, 1998/09/01}
\enonce{000872}{}
On consid\`ere l'\'equation diff\'erentielle suivante :
$$
(E.D.) \quad y''-4y'+4y = d(x),
$$
o\`u $d$ est une fonction qui sera pr\'ecis\'ee plus loin.
\begin{enumerate}
    \item R\'esoudre l'\'equation diff\'erentielle homog\`ene (ou sans second membre)
associ\'ee \`a $(E.D.)$.
    \item Trouver une solution particuli\`ere de $(E.D.)$ lorsque $d(x)=e^{-2x}$
et lorsque $d(x)=e^{2x}$ respectivement.
    \item Donner la forme g\'en\'erale des solutions de $(E.D)$ lorsque
$$d(x) = \frac{e^{-2x}+e^{2x}}{4}.$$
\end{enumerate}
\finenonce{000872} 
\finexercice


\exercice{873, ridde, 1999/11/01}

\enonce{000873}{}
R\'esoudre sur $\Rr$ :
\begin{enumerate}
\item $y''-4y = 4e^{-2x}$.
\item $y''-3y' + 2y = (x^2 + 1)e^x$.
\item $y''-2y' + y = e^x \sin x$.
\item $y'' + y = e^{-\left|x\right|}$.
\end{enumerate}
\finenonce{000873}



\finexercice



\exercice{874, ridde, 1999/11/01}

\enonce{000874}{}
Trouver les $f : \Rr \rightarrow \Rr$ deux fois d\'erivables telles que $\forall
x \in \Rr$ $f'' (x) + f (-x) = x$.
\finenonce{000874}



\finexercice



\exercice{875, ridde, 1999/11/01}

\enonce{000875}{}
R\'esoudre sur $]0,  + \infty[$ $xy''-y'-x^3y = 0$ en posant $z (t) = y (\sqrt t)$.
\finenonce{000875}



\finexercice



\exercice{876, ridde, 1999/11/01}

\enonce{000876}{}
R\'esoudre en posant $z (t) = y (e^t)$ ou $y (-e^t)$ suivant le signe de $x$, les
\'equations diff\'erentielles (d'Euler) suivantes :
\begin{enumerate}
\item $x^2y''-2y = x$.
\item $x^2y'' + xy' + y = x \ln \left|x\right|$.
\end{enumerate}
\finenonce{000876}



\finexercice



\exercice{877, ridde, 1999/11/01}

\enonce{000877}{}
R\'esoudre l'\'equation diff\'erentielle de Bernouilli $x^2y^2-xy'-3y = 0$ en supposant
que $y$ ne s'annule pas et en posant $z = \frac 1y$.
\finenonce{000877}



\finexercice



\exercice{878, gourio, 2001/09/01}

\enonce{000878}{}
R\'{e}soudre sur ${\Rr }$ :
\begin{eqnarray*}
x\frac{dy(x)}{dx}-2y(x) &=&x^{4}, \\
y"(x)-4y(x) &=&4e^{-2x}, \\
y"(x)-2y^{\prime }(x)+y(x) &=&e^{x}\sin x.
\end{eqnarray*}
\finenonce{000878}



\finexercice



\exercice{879, gourio, 2001/09/01}

\enonce{000879}{}
En posant $z=\frac{1}{y}$ et en supposant que $y  $ ne s'annulle pas,
r\'{e}soudre l'\'{e}quation (de Bernoulli) :
$$x^{2}\frac{d^{2}y(x)}{dx^{2}}-x\frac{dy(x)}{dx}-3y(x)=0. $$
\finenonce{000879}



\finexercice



\exercice{880, gourio, 2001/09/01}
\enonce{000880}{}
 R\'{e}soudre : $y''(x)+2y^{\prime }(x)+y(x)=2x\cos
x\cosh x$.

\finenonce{000880} 
\finexercice


\exercice{881, gourio, 2001/09/01}
\enonce{000881}{}
 D\'{e}terminer les $f\in C^{2}({\Rr},{\Rr}) $
telles que :
$$\forall x\in \Rr,f''(x)+f(-x)=x\cos x. $$
\finenonce{000881} 
\finexercice


\exercice{882, gourio, 2001/09/01}

\enonce{000882}{}
Soit $p$ continue positive non nulle ; montrer que toute solution
de $y''(x)+p(x)y(x)=0$ s'annule au moins une fois sur ${\Rr}.$
\finenonce{000882}



\finexercice



\exercice{883, gourio, 2001/09/01}

\enonce{000883}{}
Montrer que toute solution de $y''(x)e^{-x^{2}}+y(x)=0 $ est
born\'{e}e sur ${\Rr}.$
\finenonce{000883}



\finexercice



\exercice{884, gourio, 2001/09/01}
\enonce{000884}{}
 En posant $t=\arctan x,$ r\'{e}soudre :
$$y''(x)+\frac{2x}{1+x^{2}}y^{\prime }(x)+\frac{y(x)}{(1+x^{2})^{2}}=0. $$
\finenonce{000884} 
\finexercice


\exercice{885, gourio, 2001/09/01}
\enonce{000885}{}
 R\'{e}soudre par le changement de fonction
$z=\frac{y}{x} $ l'\'{e}quation diff\'{e}rentielle :
$${x}^{2}y''(x)-2xy^{\prime }(x)+(2-x^{2})y(x)=0. $$
\finenonce{000885} 
\finexercice


\exercice{5479, rouget, 2010/07/10}
\enonce{005479}{**}
Résoudre sur $\Rr$ les équations différentielles~:~

$$
\begin{array}{lll}
1)\;y''-2y'+2y=x\cos x\ch x&2)\;y''+6y'+9y=x^2e^{2x}&3)\;y''-2y'+y=\ch x\\
4)\;y''-2ky'+(1+k^2)y=e^x\sin x,\;k\in\Rr\setminus\{1\}
\end{array}
$$
\finenonce{005479}


\finexercice


\exercice{5480, rouget, 2010/07/10}
\enonce{005480}{**}
On considère l'équation différentielle $(E)$~:~$ax^2y''+bxy'+cy= 0$ ($a$, $b$, $c$ réels, $a\neq0$)
pour $x\in]0,+\infty[$.

\begin{enumerate}
\item  Soit $y$ une fonction deux fois dérivable sur $]0,+\infty[$. Pour $t\in\Rr$, on pose $z(t)=y(e^t)$. Vérifier
que $y$ est deux fois dérivable sur $]0,+\infty[$ si et seulement si $z$ est deux fois dérivable sur $\Rr$.

\item  Effectuer le changement d'inconnue précédent dans l'équation différentielle $(E)$ et vérifier que la
résolution de $(E)$ se ramène à la résolution d'une équation linéaire du second ordre à coefficients constants.
\item  Résoudre sur $]0,+\infty[$, l'équation différentielle $x^2y''-xy'+y=0$.
\end{enumerate}

\finenonce{005480}


\finexercice


\exercice{6997, blanc-centi, 2015/07/04}
\video{P_IrLQoPqYs}
\enonce{006997}{}
Résoudre
\begin{enumerate}
\item $y''-3y'+2y=0$
\item $y''+2y'+2y=0$
\item $y''-2y'+y=0$
\item $y''+y=2\cos^2x$
\end{enumerate}
\finenonce{006997}


\finexercice


\exercice{6998, blanc-centi, 2015/07/04}
\video{iRTCx-ikFP0}
\enonce{006998}{}
On considère $y''-4y'+4y=d(x)$. 
Résoudre l'équation homogène, puis trouver une solution particulière 
lorsque $d(x)=e^{-2x}$, puis $d(x)=e^{2x}$. 
Donner la forme générale des solutions quand $d(x)=\frac{1}{2}\ch(2x)$.
\finenonce{006998}


\finexercice


\exercice{6999, blanc-centi, 2015/07/04}
\video{cNVcKnV6F2E}
\enonce{006999}{}
Résoudre sur $]0;\pi[$ l'équation différentielle
$y''+y=\mathrm{cotan}\, x$, où $\mathrm{cotan}\, x = \frac{\cos x}{\sin x}$.
\finenonce{006999}


\finexercice


\exercice{7000, blanc-centi, 2015/07/04}
\video{nFoOhAmqUtQ}
\enonce{007000}{}
Résoudre les équations différentielles suivantes à l'aide du changement de variable suggéré.
\begin{enumerate}
\item $x^2y''+xy'+y=0$, sur $]0;+\infty[$, en posant $x=e^t$;
\item $(1+x^2)^2y''+2x(1+x^2)y'+my=0$, sur $\R$, en posant $x=\tan t$ (en fonction de $m\in\R$).
\end{enumerate}
\finenonce{007000}


\finexercice

\exercice{7003, blanc-centi, 2015/07/04}
\video{lBNXBvuCSD0}
\enonce{007003}{}
\ \begin{enumerate}
\item Résoudre sur $]0;+\infty[$ l'équation différentielle $x^2y''+y=0$ (utiliser le changement de variable $x=e^t$).
\item Trouver toutes les fonctions de classe $\mathcal{C}^1$ sur $\R$ vérifiant
$$\forall x\not=0,\ f'(x)=f\left(\frac{1}{x}\right)$$ 
\end{enumerate}
\finenonce{007003}


\finexercice

\section{225.03 Raccordement de solutions}

\section{225.04 Equations différentielles linéaires}


\exercice{4054, quercia, 2010/03/11}
\enonce{004054}{\'Equations linéaires d'ordre 1}

Intégrer les équations suivantes :


\begin{enumerate}
  \item   $(2+x)y' = 2-y$.
      

  \item   $xy'+y = \cos x$.
      

  \item   $(1+x)y' + y = (1+x)\sin x$.
      

  \item   $x^3y' -x^2y = 1$.
      

  \item   $3xy' - 4y = x$.
      

  \item   $y'+y = \sin x + 3\sin 2x$.
      

  \item   $2x(1-x)y' + (1-2x)y = 1$.
      


  \item   $x(x+1)y' + y = \arctan x$. \label{arctan}

  \item   $x(x^2-1)y' + 2y = x\ln x - x^2$.
      

pour \ref{arctan} : \'Etudier les problèmes de raccordement.

\end{enumerate}
\finenonce{004054}



\finexercice


\exercice{4055, quercia, 2010/03/11}
\enonce{004055}{\'Equations d'ordre 2 à coefficients constants}

Intégrer :


\begin{enumerate}
  \item   $y'' -2y' + 2y = xe^x$.
      

  \item   $y'' - 4y' + 4y = 2(x-2)e^x$.
      

  \item   $y'' - 4y' + 13y = 10\cos 2x + 25\sin 2x$.
      

  \item   $y'' + y = \mathrm{cotan} x$.
      

  \item   $y'' + 3y' + 2y = \frac{x-1}{x^2} e^{-x}$.
      

  \item   $y'' + y = P(x)$ où $P$ est un polynôme.
      

  \item   $y'^2 + y^2 = 1$ (dériver).
      
\end{enumerate}
\finenonce{004055}



\finexercice


\exercice{4056, quercia, 2010/03/11}
\enonce{004056}{\'Equations d'ordre 2 à coefficients non constants}

Intégrer les équations suivantes :

\begin{enumerate}
  \item $y''-y'-e^{2x}y = e^{3x}$ (poser $u = e^x$).
    

  \item $y'' - \left(6x+\frac1x\right)y' + 8x^2y = x^4$ (poser $u=x^2$).
    

  \item $x(1-2\ln x)y'' + (1+2\ln x)y' - \frac4xy = 0$
    (chercher une solution de la forme $y = x^\alpha$).
    

  \item $x^2y''-2xy'+2y = 2+2x^3\sin x$ (poser $u = \ln x$).
    

  \item $x(x+1)y'' - y' - 2y = 3x^2$
    (chercher une solution de l'équation homogène de la forme $y = x^\alpha$).
    

  \item $x^2y'' + 4xy' + (2-x^2)y = 1$ $\Bigl($poser $y = \frac u{x^2}\Bigr)$.
    

  \item $(x^2+3)y'' + xy' - y = 1$ (chercher les solutions polynomiales).
    

  \item $xy''-2y'-xy = 0$ (dériver deux fois).
    
\end{enumerate}
\finenonce{004056}



\finexercice


\exercice{4057, quercia, 2010/03/11}
\enonce{004057}{Résolution par DSE}

Chercher les solutions développables en série entière des équations suivantes
et résoudre complètement ces équations.

\begin{enumerate}
  \item $4xy'' - 2y' + 9x^2y = 0$.
    


  \item $xy'' + 2y' - xy = 0$.
    


  \item $4xy'' + 2y' - y = 0$.
    

  \item $y'' + xy' + 3y = 0$.\label{y2+xy1+3y}
    


  \item $x^2y'' + 6xy' + (6-x^2)y = -1$.
    


  \item $x(x-1)y'' + 3xy' + y = 0$.
    


\end{enumerate}
\finenonce{004057}



\finexercice


\exercice{4058, quercia, 2010/03/11}
\enonce{004058}{$y^{(4)}+y''+y = |\sin x|$}

Montrer que l'équation : $y^{(4)}+y''+y = |\sin x|$ admet une et une seule solution
$\pi$-périodique.

\finenonce{004058}



\finexercice


\exercice{4059, quercia, 2010/03/11}
\enonce{004059}{$y'' + k^2y = \sum_{n=1}^\infty \frac{\cos nx}{n^2}$}

Soit $k \in \R$.
Résoudre $y'' + k^2y = \sum_{n=1}^\infty \frac{\cos nx}{n^2}$.


\finenonce{004059}



\finexercice


\exercice{4060, quercia, 2010/03/11}
\enonce{004060}{$y' = |x-y|$}

Résoudre l'équation : $y' = |x-y|$. \'Etudier les problèmes de raccordement.


\finenonce{004060}



\finexercice


\exercice{4061, quercia, 2010/03/11}
\enonce{004061}{$y'' + |y| = 1$}

Résoudre l'équation $y'' + |y| = 1$, avec $y(0) = 0,\ y'(0) = 1$.

\finenonce{004061}



\finexercice


\exercice{4062, quercia, 2010/03/11}
\enonce{004062}{Mines MP 2000 }

Résoudre $(E) : 4xy''+2y'+y=0$ sachant que $(E)$ admet deux solutions $y$ et $z$ telles que
$yz=1$. Comment résoudre cette équation sans l'indication ?

\finenonce{004062}



\finexercice


\exercice{4063, quercia, 2010/03/11}
\enonce{004063}{Mines MP 2000}

Soit $E=\mathcal{C}(\R^+,\R)$, $b\in \R$ et $a>0$.
\begin{enumerate}
  \item Montrer que, pour tout $f\in E$, il existe un unique $g$ de $\mathcal{C}^1(\R^+,\R)$ tel que 
    $\begin{cases}g'+ag=f\cr g(0)=b.\cr\end{cases}$
    
  \item Montrer que si $f$ est intégrable sur $\R^+$, $g$ l'est également. Relation entre
    $ \int_{t=0}^{+\infty}f(t)\,d t$ et $ \int_{t=0}^{+\infty}g(t)\,d t$.
    
\end{enumerate}
\finenonce{004063}



\finexercice


\exercice{4064, quercia, 2010/03/11}
\enonce{004064}{Systèmes différentiels à coefficients constants}

\def\qpar{}%
$x,y,z$ sont des fonctions de $t$. Résoudre les systèmes :

\begin{enumerate}
  \item $\begin{cases} x' = 2y + 2z\cr y' = -x + 2y + 2z\cr z' = -x + y  + 3z.\cr\end{cases}$
    
%
  \item $\begin{cases} y'+y = z \cr z'+2z = y-1. \cr\end{cases}$
    
%
  \item $\begin{cases} y' = y+z+ \sin t\cr z'=-y+3z.\cr\end{cases}$
    
%
  \item $\begin{cases} x' = x+y-z \cr y'=2x+y-2z\cr z'=-2x+2y+z.\cr\end{cases}$
    


  \item $\begin{cases} x' = 2x + y + z   \cr
             y' = x - y - z    \cr
             z' = -x + 2y + 2z.\cr\end{cases}$
    
%
  \item $\begin{cases} x' = 2x + z + \sh t\cr
             y' = x - y - z + \ch t\cr
             z' = -x + 2y + 2z - \ch t.\cr\end{cases}$
    


\end{enumerate}
\finenonce{004064}



\finexercice


\exercice{4065, quercia, 2010/03/11}
\enonce{004065}{Système différentiel à coefficients non constants}
Résoudre le système différentiel suivant :
$\begin{cases}(t^2+1)x' = tx + y + 2t^2-1\cr
        (t^2+1)y' = x - ty + 3t.\cr\end{cases}$

\finenonce{004065}



\finexercice


\exercice{4066, quercia, 2010/03/11}
\enonce{004066}{Lemme des noyaux, Matexo}

Soit $(p,q)\in\R^2$ et $(E)$ l'équation~: $y'' + p y' + q  y = 0$.
On note $S$ l'ensemble des solutions de $(E)$ et $D$ l'application de $S$ dans 
$S$ définie par $D(f) = f'$.

\begin{enumerate}
  \item $D$ peut elle être une homothétie ?
  \item Déterminer les valeurs de $p$ et $q$ pour lequelles $D$ n'est pas un 
    isomorphisme de~$S$.
  \item Vérifiez que $D\circ D + p D + q\mathrm{id}_S=0$ et montrer qu'il existe des 
    nombres complexes $r_1$ et $r_2$ tels que :
    $(D - r_1 \mathrm{id}_S) \circ (D - r_2 \mathrm{id}_S) = 0$.
  \item Les applications $D - r_1 \mathrm{id}_S$, $D - r_2 \mathrm{id}_S$
    peuvent-elles être inversibles ?

\end{enumerate}
\finenonce{004066}



\finexercice


\exercice{4067, quercia, 2010/03/11}
\enonce{004067}{$y'' + x\ y' + y = 0$, Matexo}

On désigne par $y$ la solution de l'équation différentielle
$y'' + x\ y' + y = 0$,
avec les conditions de Cauchy $y(0) = 0, y'(0) = 1$.

\begin{enumerate}
  \item Montrer que les dérivées de $y$ vérifient
$y^{(n)} + x\ y^{(n-1)} + (n-1)\ y^{(n-2)} = 0, \quad \forall n \geq 2$.
  \item Calculer par récurrence les dérivées successives de $y$ en zéro.
  \item Montrer que $y$ admet le développement limité à l'origine
$(x \to 0)$:

$$y(x) = x - \frac{2 x^3}{3!} + \frac{8x^5}{5!} +\dots 
       + \frac{(-2)^kk!\,x^{2k+1}}{(2 k+1)!} +  o(x^{2k+2}).$$

\end{enumerate}
\finenonce{004067}



\finexercice


\exercice{4068, quercia, 2010/03/11}
\enonce{004068}{$f''(x) + f(-x) = x\cos x$}

Trouver les fonctions $f:\R \to \R$ de classe $\mathcal{C}^2$ telles que :
$\forall\ x \in R,\ f''(x) + f(-x) = x\cos x$.



\finenonce{004068}



\finexercice


\exercice{4069, quercia, 2010/03/11}
\enonce{004069}{$y'' + \frac {2y'}{\tanh x} + y = 0$}

On considère l'équation différentielle :
$(*) \iff y'' + \frac {2y'}{\tanh x} + y = 0$.

\begin{enumerate}
  \item  On pose $z(x) = y'(x) + \frac {y(x)}{\tanh x}$.
     \'Ecrire l'équation différentielle (d'ordre 1) sur $z$ déduite de $(*)$.

  \item  Résoudre sur $]-\infty,0[$ et $]0,+\infty[$ l'équation en $z$, puis
     $(*)$.
     

  \item  Parmi les solutions trouvées, quelles sont celles prolongeables en 0 ?

     On note $y_0$ la solution de $(*)$ telle que $y_0(x) \to 1$ lorsque $x\to0$.

  \item  Démontrer que $y_0$ est de classe $\mathcal{C}^1$ et que $\frac {y_0'(x)}{\tanh x}$
     admet une limite finie en $0$.

     En déduire que $y_0$ est de classe $\mathcal{C}^2$ sur $\R$.

  \item  Est-ce que l'aire comprise entre la courbe de $y_0$
     et l'axe des abscisses est finie~?

\end{enumerate}
\finenonce{004069}



\finexercice


\exercice{4070, quercia, 2010/03/11}
\enonce{004070}{$y'' + 2xy' + (x^2-1)y = 0$}

Soit $E = \mathcal{C}^\infty(\R,\C)$ et
$\Phi : E \to E$, définie par $\Phi(f) = g$ où $g$ est l'application ${g:t \mapsto f'(t)+tf(t).}$

\begin{enumerate}
  \item Trouver les valeurs propres et les vecteurs propres de $\Phi$.
    

  \item Trouver les valeurs propres et les vecteurs propres de $\Phi^2$.
    

  \item Résoudre l'équation : $y'' + 2xy' + (x^2-1)y = 0$.
    

\end{enumerate}
\finenonce{004070}



\finexercice


\exercice{4071, quercia, 2010/03/11}
\enonce{004071}{$x^2f''(x) + xf'(x) = \lambda f(x)$}

Déterminer les éléments propres des endomorphismes suivants :
\begin{enumerate}
  \item $E = \R[X]$
   $\Phi(P)(X) = X^2P''(X) + XP'(X)$.
   

  \item $E = \mathcal{C}^\infty(]0,+\infty[)$
   $\Phi(f)(x) = x^2f''(x) + xf'(x)$.
   

  \item $E = \mathcal{C}^\infty(]0,1[)$
  $\Phi(f)(x) = \sqrt{\frac {1-x}x}f'(x)$.
  

\end{enumerate}
\finenonce{004071}



\finexercice


\exercice{4072, quercia, 2010/03/11}
\enonce{004072}{$AP' - nA'P = \lambda P$}

Soit $A$ un polynôme à coefficients réels de degré 2 donné.
Au polynôme $P$ de degré inférieur ou égal à $2n$ on fait correspondre le
polynôme $Q = AP' - nA'P$.

\begin{enumerate}
  \item Montrer qu'on définit ainsi un endomorphisme $\Phi$ de $\R_{2n}[X]$.
  \item Chercher les valeurs propres et les vecteurs propres de $\Phi$ dans les
    cas particuliers :\par
   \begin{enumerate} 
    \item $A = X^2-1$, 
    \item $A = X^2$,   
    \item $A = X^2+1$.
  \end{enumerate}  
\end{enumerate}
\finenonce{004072}



\finexercice


\exercice{4073, quercia, 2010/03/11}
\enonce{004073}{\'Equation intégrale}

Trouver les applications $g : {\R^+} \to {\R^+}$ continues vérifiant pour tout
$x > 0$ :
$$\frac12 \int_{t=0}^x g^2(t)\,dt = \frac1x\left( \int_{t=0}^x g(t)\,d t\right)^2.$$



\finenonce{004073}



\finexercice


\exercice{4074, quercia, 2010/03/11}
\enonce{004074}{Inéquations différentielles}

Soient ${a,b} : \R \to \R$ continues, et $y,z$ solutions de
$\begin{cases}y(0) = z(0)\cr y' = a(t)y + b(t)\cr z'\le a(t)z + b(t).\end{cases}$

Démontrer que : $\forall\ t \ge 0$, on a $y(t) \ge z(t)$.



\finenonce{004074}



\finexercice


\exercice{4075, quercia, 2010/03/11}
\enonce{004075}{Tangentes parallèles ou concourantes}

Soit l'équation $(*) \Leftrightarrow y' = a(x)y + b(x)$ et $x_0$ un réel fixé.
Montrer que les tangentes aux courbes intégrales au point d'abscisse $x_0$ sont
parallèles ou concourantes.


\finenonce{004075}



\finexercice


\exercice{4076, quercia, 2010/03/11}
\enonce{004076}{$y' + ay =$ fct périodique}

Soit $\lambda \in \C$ et $\varphi : \R \to \C$ $T$-périodique.
On considère l'équation : $(*) \Leftrightarrow y' + \lambda y = \varphi(x)$.

\begin{enumerate}
  \item Montrer que si $y$ est solution de $(*)$, alors $y(x+T)$ est aussi solution.
  \item En déduire que $y$, solution de $(*)$, est $T$-périodique si et seulement si $y(0) = y(T)$.
  \item Montrer que, sauf pour des valeurs exceptionnelles de $\lambda$, l'équation
    $(*)$ admet une et une seule solution $T$-périodique.
\end{enumerate}
\finenonce{004076}



\finexercice


\exercice{4077, quercia, 2010/03/11}
\enonce{004077}{Coefficients périodiques}

Soit l'équation $(*) \Leftrightarrow y' + a(x)y = b(x)$ où $a,b$ sont des fonctions
continues, $T$-périodiques.

\begin{enumerate}
  \item Montrer que si $y$ est solution de $(*)$, alors la fonction définie par
    $z(x) = y(x+T)$ est aussi solution.
  \item En déduire que si $ \int_{t=0}^T a(t)\,d t \ne 0$, alors $(*)$ admet une unique
    solution $T$-périodique.
    

\end{enumerate}
\finenonce{004077}



\finexercice


\exercice{4078, quercia, 2010/03/11}
\enonce{004078}{\'Equation intégrale}

Soit $E = \{ \text{fcts : } [0,1] \to \R \text{ continues } \}$
et $\Phi : E \to E, f \mapsto g$ avec $g(x) =  \int_{t=0}^1 \inf(t,x)f(t)\,d t$.

Chercher les valeurs propres et les fonctions propres de $\Phi$.
\finenonce{004078}



\finexercice


\exercice{4079, quercia, 2010/03/11}
\enonce{004079}{Matexo}

Soit $k\in\R^*$ fixé. On considère :
$E=\{f\in\mathcal{C}^2([0,1],\R)\text{ tq } f(0) = 0 \text{ et } f(1)=3\}$.

Déterminer $\displaystyle \inf_{f\in E}  \int_{t=0}^1 (f'(t)+kf(t))^2\,d t$.
{\it Indication~: poser $f'+kf = g$ et calculer $f(1)$ en fonction de~$g$.}

\finenonce{004079}



\finexercice


\exercice{4080, quercia, 2010/03/11}
\enonce{004080}{Ulm-Lyon-Cachan MP$^*$ 2000}

Soient $u,v,w$ trois applications bornées et de classe~$\mathcal{C}^1$ sur~$\R$,
à valeurs dans~$\R^3$, vérifiant~:
$u'+v' = w$; $w'=-v$; $ \int_{0}^{\infty}\|u'\|^2<+\infty$.
On suppose qu'il existe une suite de terme général $t_n$ tendant vers~$+\infty$
telle que $u(t_n)$ tend vers~$a\in\R^3$.

\begin{enumerate}
  \item Montrer que la suite de terme général $u_n=\frac1{2\pi} \int_{t=t_n}^{t_n+2\pi} u(t)\,d t$
    tend vers~$a$.


  \item Montrer que les suites de termes généraux $v_n=\frac1{2\pi} \int_{t=t_n}^{t_n+2\pi} v(t)\,d t$
    et $w_n=\frac1{2\pi} \int_{t=t_n}^{t_n+2\pi} w(t)\,d t$ tendent vers~$0$.
\end{enumerate}
\finenonce{004080}



\finexercice


\exercice{4081, quercia, 2010/03/11}
\enonce{004081}{Centrale MP 2001}

Soit $f$ continue et intégrable sur~$\R$. On considère l'équation
différentielle $(E)$ : $y' - y + f = 0$.

\begin{enumerate}
  \item Montrer que $(E)$ admet une unique solution $F$ bornée sur~$\R$.
    
    
  \item Montrer que $F$ est intégrable sur $\R$ et comparer $ \int_{-\infty}^{+\infty} F$ et  $ \int_{-\infty}^{+\infty} f$.
    
\end{enumerate}
\finenonce{004081}



\finexercice


\exercice{4082, quercia, 2010/03/11}
\enonce{004082}{Centrale MP 2001}

Trouver les fonctions $f,g$ continues vérifiant~:
$ \int_{t=0}^x f(t)\,d t = x - 1 + g(x)$ et
$ \int_{t=0}^x g(t)\,d t = x - 1 + f(x)$.

\finenonce{004082}



\finexercice


\exercice{4083, quercia, 2010/03/11}
\enonce{004083}{X MP$^*$ 2000}

On considère l'équation différentielle $y'=\sin(x+y)$. Montrer que
pour toute condition initiale l'intervalle maximal est
$\R$. Ensemble des points d'inflexion des courbes solutions ?


\finenonce{004083}



\finexercice


\exercice{4084, quercia, 2010/03/11}
\enonce{004084}{Polytechnique PC 2002}
Soit l'équation différentielle~:
$(E)\Leftrightarrow u''(x) + (k-2d\cos(x))u(x)= 0$.

\begin{enumerate}
  \item Existence et domaine de définition des solutions
    maximales $A$ et $B$ telles que
    $A(0)=1$, $A'(0)=0$ et $B(0)=0$, $B'(0)=1$.


  \item Montrer que $A$ est paire et $B$ est impaire.


  \item Montrer que $A(k,d,x) = A(k,0,x) +
    2d \int_{t=0}^xB(d,0,x-t)A(k,d,x)\cos(t)\,d t$.

\end{enumerate}
\finenonce{004084}



\finexercice


\exercice{4085, quercia, 2010/03/11}
\enonce{004085}{$f \mapsto f+f'$ (Mines MP 2003)}
Soit~$E$ l'ensemble des fonctions $f : \R \to \R$ de classe $\mathcal{C}^\infty$
telle que pour tout~$k\in\N$, $f^{(k)}$ est bornée.

Soit $u : E \to E, f \mapsto {f+f'.}$

\begin{enumerate}
  \item Montrer que $u\in\mathcal{L}(E)$.
  \item Est-ce que~$u$ est injectif~?
    
  \item Est-ce que~$u$ est surjectif~?
    

\end{enumerate}
\finenonce{004085}



\finexercice


\exercice{4086, quercia, 2010/03/11}
\enonce{004086}{Centrale MP 2004}
On considère l'équation différentielle~: $-y'' + \frac y{p^2} = f$
où $p\in\N^*$ et $f$ est une fonction continue donnée.

\begin{enumerate}
  \item Donner les solutions de cette équation.
    Montrer que $x \mapsto - \int_{t=0}^x pf(t)\sh\Bigl(\frac{x-t}p\Bigr)\,d t$
    est solution.
  \item Montrer qu'il existe une unique solution telle que $y(0) = y(1) = 0$.
    On la note $u_p$.
    
  \item Montrer que $(u_p)$ converge simplement vers une fonction que l'on déterminera.
    
\end{enumerate}
\finenonce{004086}



\finexercice


\exercice{4087, quercia, 2010/03/11}
\enonce{004087}{Centrale MP 2004}

Soit $\lambda\in\R$ et $({\cal E})\Leftrightarrow \frac{d^2 u(x)}{d x^2} + (1-\lambda)x^2u(x) = 0$.

\begin{enumerate}
  \item Montrer que les solutions de~$({\cal E})$ sont de la forme $H(x)e^{-x^2/2}$
    où $H$ est une fonction développable en série entière.
    
  \item Déterminer les valeurs de~$\lambda$ telles que $H$ soit une fonction
    polynomiale non nulle.
    

\end{enumerate}
\finenonce{004087}



\finexercice


\exercice{4088, quercia, 2010/03/11}
\enonce{004088}{Lemme de Gronwall (X MP$^*$ 2003)}

Soient~$f,g$ deux fonctions continues et $a\in\R$ vérifiant~:
$\forall\ t\ge 0,\ g(t)\ge 0\quad\text{et}\quad f(t)\le a +  \int_{u=0}^t f(u)g(u)\,d u.$
Montrer~: $\forall\ t\ge0,\ f(t)\le a\exp\Bigl( \int_{u=0}^t g(u)\,d u\Bigr)$.

\finenonce{004088}



\finexercice


\exercice{4089, quercia, 2010/03/11}
\enonce{004089}{$y''+y \ge 0$}

Soit $f : \R \to \R$ de classe $\mathcal{C}^2$ telle que :
$\forall\ x \in \R,\ f(x) + f''(x) \ge 0$.

Montrer que : $\forall\ x \in \R,\ f(x) + f(x+\pi) \ge 0$.

\finenonce{004089}



\finexercice


\exercice{4090, quercia, 2010/03/11}
\enonce{004090}{$f'' + f' + f \to 0$}

Soit $f : \R \to\R$ de classe $\mathcal{C}^2$ telle que
$f''(t) + f'(t) + f(t) \to 0$ lorsque $t\to+\infty$. Démontrer que $f(t) \to 0$ lorsque $t\to\infty$.
\finenonce{004090}



\finexercice


\exercice{4091, quercia, 2010/03/11}
\enonce{004091}{$f'' \ge f + 2/ch(x)^3$, Centrale PC 1997}

Soit $f$ de classe $\mathcal{C}^2$ de $\R$ dans~$\R$ telle que $f(0) = f'(0) = 0$
et pour tout $x$~: $f''(x) \ge f(x) + \frac 2{\ch(x)^3}$.
Montrer pour tout $x$~: $f(x) \ge \frac{\sh(x)^2}{\ch(x)}$.
\finenonce{004091}



\finexercice


\exercice{4092, quercia, 2010/03/11}
\enonce{004092}{$y'+ay = b$, $y(-\infty) = 0$}

Soit ${a,b} : \R \to \R$ continues telles que : $\forall\ x\in\R,\ a(x)\ge 1$
et $b(x)\to 0$ lorsque $x\to+\infty$.

\begin{enumerate}
  \item Montrer que toute solution de l'équation : $y' + ay = b$ tend vers $0$ en $+\infty$.
    

  \item On suppose $b(x)\to 0$ lorsque $x\to-\infty$. Montrer qu'il y a une unique solution $y$
    qui tend vers $0$ en $-\infty$.
    
\end{enumerate}
\finenonce{004092}



\finexercice


\exercice{4093, quercia, 2010/03/11}
\enonce{004093}{ $y'' + ay = 0$, $a > 0  \Rightarrow  y$ s'annule}

Soit $a : \R \to {\R^{+*}}$ une fonction continue.

\begin{enumerate}
  \item Soit $y$ une solution de l'équation $ y'' + a(x)y = 0$.
    Montrer que $y$ s'annule au moins une fois sur $\R$.
    

  \item Soit $z$ une solution de l'équation $ z'' - a(x)z = 0$.
    Montrer que $z=0$ ou bien $z$ s'annule au plus une fois sur $\R$.
    

\end{enumerate}
\finenonce{004093}



\finexercice


\exercice{4094, quercia, 2010/03/11}
\enonce{004094}{$y'' + ay' = 0$ $a$ croissante positive $=> y$ est bornée}

Soit $a : \R \to \R$ une fonction de classe $\mathcal{C}^1$ croissante strictement positive et
$y$ une solution de l'équation : $y'' + a(t)y = 0$.
Montrer que $y$ est bornée au voisinage de $+\infty$
(on étudiera $z = y^2 + y'^2/a$).
\finenonce{004094}



\finexercice


\exercice{4095, quercia, 2010/03/11}
\enonce{004095}{$y'' + ay =0$, $a$ intégrable}

Soit $a : {\R^+} \to \R$ continue intégrable. Montrer l'équation
$y'' + a(t)y =0$ admet des solutions non bornées sur $[0,+\infty[$
(on commencera par prouver que si $y_1,y_2$ sont deux solutions alors le déterminant
wronskien de $y_1$ et $y_2$ est constant).



\finenonce{004095}



\finexercice


\exercice{4096, quercia, 2010/03/11}
\enonce{004096}{Zéros entrelacés (Centrale MP 2003)}
Soient $r$ et $q$ deux fonctions continues définies sur~$I={[a,b]}$ telles que~:
$\forall\ x\in I,\ r(x)\ge q(x)$.
On considère les équations différentielles~:
$$(E_1)\Leftrightarrow y'' + qy = 0,
  \qquad (E_2)\Leftrightarrow z'' + rz = 0.$$
\begin{enumerate}
  \item Soit~$y$ une solution de~$(E_1)$ et $x_0,x_1$ deux zéros consécutifs
    de~$y$. $y'(x_0)$ et $y'(x_1)$ peuvent-ils être nuls~? Que dire de leurs signes~?
    
  \item Soit~$z$ une solution de~$(E_2)$. On considère
    $W(x) = \begin{vmatrix}y(x)&z(x)\cr y'(x)&z'(x)\cr\end{vmatrix}$. Calculer $W'(x)$
    et $W(x_1)-W(x_0)$.
    
  \item Montrer que $z$ a un zéro dans~$]x_0,x_1[$ ou $z(x_0)=z(x_1)=0$.
    
  \item Soit~$u$ une solution de~$(E_1)$. Montrer que $u$ est soit proportionnelle
    à~$y$, soit admet un unique zéro dans~$]x_0,x_1[$.
    
\end{enumerate}
\finenonce{004096}



\finexercice


\exercice{4097, quercia, 2010/03/11}
\enonce{004097}{Zéros des solutions de $y''+ay'+by = 0$}

On considère l'équation $(*) \iff y'' + a(t)y' + b(t)y = 0$, avec $a,b$ continues.

\begin{enumerate}
  \item Soit $y$ une solution non nulle de $(*)$. Montrer que les zéros de $y$ sont isolés.
  \item Soient $y,z$ deux solutions de $(*)$ non proportionelles.
  \begin{enumerate}
     \item Montrer que $y$ et $z$ n'ont pas de zéros commun.
    \item Montrer que si $u,v$ sont deux zéros consécutifs de $y$, alors $z$ possède un
    unique zéro dans l'intervalle $]u,v[$
    (étudier $\frac zy$).
  \end{enumerate}
\end{enumerate}
\finenonce{004097}



\finexercice


\exercice{4098, quercia, 2010/03/11}
\enonce{004098}{$y'' + \left(1+\frac\lambda{t^2}\right)y = 0$}

Soit $\lambda > 0$ et $y$ une solution de
$y'' + \left(1+\frac\lambda{t^2}\right)y = 0$.
Montrer que pour tout $a \in \R$, $y$ a un zéro dans l'intervalle $]a,a+\pi[$.
(\'Etudier $z = y'\varphi - y\varphi'$ où $\varphi(t) = \sin(t-a)$)

\finenonce{004098}



\finexercice


\exercice{4099, quercia, 2010/03/11}
\enonce{004099}{$y'' + e^t y = 0$}

Soit $y : \R \to \R$ une solution non identiquement nulle de
$y'' + e^t y = 0$.

\begin{enumerate}
  \item Montrer que l'ensemble des zéros de $y$ est infini dénombrable.
    

  \item On note $a_n$ le $n$ème zéro positif de $y$.
    En utilisant les fonctions
    $\begin{cases} \varphi(t) = \sin\Bigl(e^{a_n/2}(t-a_n)\Bigr) \cr
             \psi(t) = \sin\Bigl(e^{a_{n+1}/2}(t-a_n)\Bigr),\cr\end{cases}$ montrer que
    $\frac\pi{e^{a_{n+1}/2}} \le a_{n+1}-a_n \le \frac\pi{e^{a_n/2}}$.

  \item Donner un équivalent de $a_n$ quand $n\to\infty$.
    
\end{enumerate}
\finenonce{004099}



\finexercice


\exercice{4100, quercia, 2010/03/11}
\enonce{004100}{Conditions aux limites}

Soient ${f,g} : {[a,b]} \to \R$ continues avec $f$ positive.
Montrer qu'il existe une unique solution pour le problème aux limites~:
$y'' = f(t)y + g(t)$, $y(a) = y(b) = 0$.

\finenonce{004100}



\finexercice


\exercice{4101, quercia, 2010/03/11}
\enonce{004101}{Comparaison de solutions}

Soient ${p,q} : \R  \to \R$ continues avec : $\forall\ x \in \R,\ q(x) < 0$.
Soient ${y,z} : \R \to \R$ de classe $\mathcal{C}^2$ vérifiant :
$$\begin{cases}\forall\ x \in \R,\ y''+p(x)y'+q(x)y \le z''+p(x)z'+q(x)z\cr
         y(a) \le z(a),\ y'(a)<z'(a).\cr\end{cases}$$
Montrer que : $\forall\ x > a,\ y(x) < z(x)$.

\finenonce{004101}



\finexercice


\exercice{4102, quercia, 2010/03/11}
\enonce{004102}{Système différentiel à coefficients positifs}
Soit $A : {\R^+} \to {\mathcal{M}_n(\R)}, t \mapsto {(a_{ij}(t))}$ continue avec~:
$\forall\ t\ge 0,\ \forall\ i,j,\ a_{ij}(t)\ge 0$ et $X$ une solution
du système différentiel $X'(t) = A(t)X(t)$.
Montrer que si toutes les coordonnées de $X(0)$ sont positives ou
nulles il en est de même pour $X(t)$ pour tout~$t$ (commencer par
le cas strictement positif).

\finenonce{004102}



\finexercice


\exercice{4103, quercia, 2010/03/11}
\enonce{004103}{Intégrale fonction d'un paramètre}

On pose $f(x) =  \int_{t=0}^{+\infty} e^{-t}e^{itx} \frac {d t}{\sqrt t}$.
Former une équation différentielle satisfaite par $f$. En déduire $f$.
\finenonce{004103}



\finexercice


\exercice{4104, quercia, 2010/03/11}
\enonce{004104}{Ulm MP$^*$ 2000}

Soit $I=[a,b]\subset \R$. On suppose que la fonction $\Delta$ est strictement
positive sur $I$.

On pose $E=\{f\in \mathcal{C}^2(I) \, |\, f(a)=f(b)=0\}$. On
considère enfin l'opérateur $K:f\mapsto \frac{f''}{\Delta}\cdotp$

\begin{enumerate}
  \item Montrer que Sp$(K)\subset {]-\infty,0[}$.
    
  \item Trouver un produit scalaire $(\ \mid\ )$ pour lequel deux vecteurs propres associés à des
    valeurs propres distinctes sont orthogonaux. 
    
  \item On suppose que $I=\R^+$ et que $\Delta (x)\ge 1$ pour $x\ge 2$. Soit $\lambda <0$.
  \begin{enumerate}
     \item Montrer qu'il existe une unique $f_{\lambda}\in \mathcal{C}^1(\R^+,\R)$ telle 
    $\begin{cases}f''_{\lambda}=\lambda \Delta f_\lambda\cr f_{\lambda}(0)=0\cr f'_{\lambda}(0)=1.\cr\end{cases}$
     \item Montrer $f_{\lambda}$ a une infinité dénombrable de zéros $(x_0<x_1<\dots <x_n<\dots)$ et 
    que la suite $(x_n)$ tend vers~$+\infty$.
    
  \end{enumerate}

\end{enumerate}
\finenonce{004104}



\finexercice


\exercice{4105, quercia, 2010/03/11}
\enonce{004105}{Centrale MP 2001}

\begin{enumerate}
  \item Soit $f$ de classe $\mathcal{C}^1$ de $[0,\pi]$ dans $\R$ telle que $f(0) = f(\pi) = 0$.
Montrer que $ \int_0^\pi f^2 \le  \int_0^\pi f'^2$.

{\it Indication~: prolonger $f$ en une fonction impaire $2\pi$-périodique.}


  \item Soit une fonction $q$ de classe $\mathcal{C}^1$ sur $[0,\pi]$, à valeurs
dans $]-\infty,1[$. Montrer que l'unique fonction $x$ de classe $\mathcal{C}^2$
s'annulant en $0$ et en $\pi$ et vérifiant l'équation différentielle
$x''(t) + q(t)x(t) = 0$ est la fonction nulle.

  \item Soit $f$ une fonction de classe $\mathcal{C}^1$ sur $[0,\pi]$ et deux réels $a,b$ fixés.
Montrer qu'il existe une unique solution $x$ de classe $\mathcal{C}^2$ vérifiant $x(0) = a$,
$x(\pi) = b$ et $x''(t) + q(t)x(t) = f(t)$.

\end{enumerate}
\finenonce{004105}



\finexercice


\exercice{4106, quercia, 2010/03/11}
\enonce{004106}{X MP$^*$ 2000}

On considère l'équation différentielle à coefficients continus sur
$\R$ : $x''+p(t)x'+q(t)x=0$. Trouver une condtion nécessaire
portant sur $p$ et $q$ pour qu'il existe deux solutions sur $\R$ dont
le produit vaut constamment~un.



\finenonce{004106}



\finexercice


\exercice{4107, quercia, 2010/03/11}
\enonce{004107}{X MP$^*$ 2000}

Soit $A$ coninue de $\mathcal{M}_n(\R)$ dans lui-même. On suppose
que les $a_{ij}(t)$ restent positifs quand $t$ décrit $\R^+$, et l'on se donne un
vecteur $X_0$ dont toutes les composantes sont positives. Montrer
qu'en désignant par $X(t)$ la valeur en $t$ du système $Y'=AY$ valant
$X_0$ en $t=0$, on a pour tout $t\ge0$ et pour tout $i$ l'inégalité
$x_i(t)\ge 0$. 


\finenonce{004107}



\finexercice


\exercice{4108, quercia, 2010/03/11}
\enonce{004108}{ENS MP 2002}

Soit une application $A$ de classe $\mathcal{C}^\infty$ sur~$\R$ à valeurs dans~$\mathcal{M}_n(\C)$,
telle que les valeurs propres de $A(0)$ aient toutes une partie réelle
strictement positive. Soit~$F$ de classe $\mathcal{C}^\infty$ sur~$\R$, à valeurs
dans~$\C^n$. Montrer qu'il existe $X$ de classe $\mathcal{C}^\infty$ sur~$\R$, à valeurs
dans $\C^n$, solution de $tX'(t) + A(t)X(t) = F(t)$.

{\it Indication~: commencer par le cas $n=1$, $A$ constante.}

\finenonce{004108}



\finexercice


\exercice{4109, quercia, 2010/03/11}
\enonce{004109}{X MP$^*$ 2003}
\begin{enumerate}
  \item Soient ${f,g}: \R \to \R$ continues et $k>0$, $t_0\in\R$
    tels que~: $\forall\ t\ge t_0,\ f(t)\le g(t) + k \int_{u=t_0}^t f(u)\,d u$.

    Montrer que~: $\forall\ t\ge t_0,\ f(t) \le g(t) + k \int_{u=t_0}^t e^{k(t-u)}g(u)\,d u$.
    

  \item Soient ${A,B} : \R\to {\mathcal{M}_n(\R)}$ continues, $T>t_0$, $K>0$ et $\eta>0$
    tels que~: $\forall\ t\in[t_0,T],\ \|\kern-1.2pt|A(t)\|\kern-1.2pt|\le K$
    et $\|\kern-1.2pt|A(t)-B(t)\|\kern-1.2pt|\le\eta$.
    On note~$M_0$ (resp. $N_0$) la solution du problème de Cauchy~:
    $M(t_0)=I,\ M'(t) = A(t)M(t)$ (resp. $N(t_0)=I,\ N'(t) = B(t)N(t)$).
    Montrer que~: $\forall\ t\in[t_0,T],\ \|\kern-1.2pt|M_0(t)-N_0(t)\|\kern-1.2pt|
    \le e^{K(t-t_0)}(e^{\eta(t-t_0)}-1)$.
    

  \item On note $X_0$ (resp. $Y_0$) la solution du problème de Cauchy dans~$\R^n$~:
    $X(t_0) = \alpha,\ X'(t) = A(t)X(t)$ (resp. $Y(t_0) = \alpha,\ Y'(t) = B(t)Y(t)$)
    où $\alpha\in\R^n$. Quelle majoration a-t-on sur $\|X_0(t)-Y_0(t)\|$~?
    
\end{enumerate}
\finenonce{004109}



\finexercice


\exercice{5874, rouget, 2010/10/16}
\enonce{005874}{**}
Résoudre sur $\Rr$ l'équation différentielle proposée :

\begin{enumerate}
 \item  $y'+y = 1$
 \item  $2y'-y =\cos x$
 \item  $y'-2y = xe^{2x}$  
 \item  $y''-4y'+4y = e^{2x}$
 \item  $y''+4y =\cos(2x)$
 \item  $y''+2y'+2y=\cos x\ch x$.
\end{enumerate}
\finenonce{005874}


\finexercice


\exercice{5875, rouget, 2010/10/16}
\enonce{005875}{*** I}
\begin{enumerate}
 \item  Soit $\alpha\in\Cc$ tel que $\text{Re}(\alpha) > 0$. Soit $f~:~\Rr\rightarrow\Rr$ de classe $C^1$ sur $\Rr$.

On suppose que quand $x$ tend vers $+\infty$, $f'+\alpha f$ tend vers $\ell\in\Cc$. Montrer que $f(x)$ tend vers $ \frac{\ell}{\alpha}$ quand $x$ tend vers $+\infty$.

\item  Soit $f~:~\Rr\rightarrow\Cc$ de classe $C^2$ sur $\Rr$ telle que $\lim_{x \rightarrow +\infty}(f+f'+f'')(x) = 0$. Montrer que $\lim_{x \rightarrow +\infty}f(x) = 0$.

\item  Soient $n\in\Nn^*$ et $f~:~\Rr\rightarrow\Cc$ de classe $C^n$ sur $\Rr$.

On note $D$ l'opérateur de dérivation. Soit $P$ un polynôme de degré $n$ unitaire dont tous les zéros ont des parties réelles strictement négatives. Montrer que $\lim_{x \rightarrow +\infty}(P(D))(f)(x)= 0\Rightarrow\lim_{x \rightarrow +\infty}f(x) = 0$.
\end{enumerate}
\finenonce{005875}


\finexercice


\exercice{5876, rouget, 2010/10/16}
\enonce{005876}{*** I}
Soit $f$ une application de classe $C^2$ sur $\Rr$ à valeurs dans $\Rr$ telle que $\forall x\in\Rr$, $f(x)+f''(x)\geqslant 0$. Montrer que $\forall x\in\Rr$, $f(x)+f(x+\pi)\geqslant 0$.
\finenonce{005876}


\finexercice


\exercice{5877, rouget, 2010/10/16}
\enonce{005877}{*** I}
Résoudre sur l'intervalle $I$ proposé :

\begin{enumerate}
 \item  $xy'-2y = 0$ ($I=\Rr$)
 \item  $xy' -y = 0$ ($I=\Rr$)
 \item  $xy'+y = 0$ ($I=\Rr$) 
 \item  $xy'-2y = x^3$ ($I=]0,+\infty[$)
 \item  $x^2y'+2xy = 1$ ($I=\Rr$)
 \item  $2x(1-x)y'+(1-x)y = 1$ ($I=]-\infty,0[$, $]0,1[$, $]1,+\infty[$, $]-\infty,1[$, $]0,+\infty[$, $\Rr$)
 \item  $|x|y' + (x-1)y = x^3$ ($I=\Rr$).
\end{enumerate}
\finenonce{005877}


\finexercice


\exercice{5878, rouget, 2010/10/16}
\enonce{005878}{*** I}
Déterminer le rayon de convergence puis calculer $\sum_{n=0}^{+\infty}(-1)^{n-1} \frac{C_{2n}^n}{2n-1}x^n$ quand $x$ appartient à l'intervalle ouvert de convergence. En déduire la valeur de $\sum_{n=0}^{+\infty} \frac{(-1)^{n-1}C_{2n}^n}{(2n-1)4^n}$.
\finenonce{005878}


\finexercice


\exercice{5879, rouget, 2010/10/16}
\enonce{005879}{**}
Résoudre les systèmes suivants :

\begin{enumerate}
 \item  $\left\{
\begin{array}{l}
x'=4x-2y\\
y'=x+y
\end{array}
\right.$\quad\item  $\left\{
\begin{array}{l}
x'=x-y+ \frac{1}{\cos t}\\
y'=2x-y
\end{array}
\right.$  sur $\left]- \frac{\pi}{2}, \frac{\pi}{2}\right[$\quad\item  $\left\{
\begin{array}{l}
x'=5x-2y+e^t\\
y'=-x+6y+t
\end{array}
\right.$  

\item   $\left\{
\begin{array}{l}
x'=5x+y-z\\
y'=2x+4y-2z\\
z'=x-y+z
\end{array}
\right.$\quad\item  $\left\{
\begin{array}{l}
x'=2x+y\\
y'=-x\\
z'=x+y+z
\end{array}
\right.$ (trouver la solution telle que $x(0)=0$, $y(0)=1$ et $z(0)=-1$).
\end{enumerate}
\finenonce{005879}


\finexercice


\exercice{5880, rouget, 2010/10/16}
\enonce{005880}{**}
Soit $A\in\mathcal{S}_n^+(\Rr)$. Montrer que pour toute solution de $X'=AX$, la fonction $t\mapsto\|X(t)\|_2$ est croissante sur $\Rr$.
\finenonce{005880}


\finexercice


\exercice{5881, rouget, 2010/10/16}
\enonce{005881}{**}
Résoudre les systèmes :

\begin{enumerate}
 \item  $\left\{
\begin{array}{l}
x'=- \frac{1}{2t}x+ \frac{1}{2t^2}y+2t\\
y'= \frac{1}{2}x+ \frac{1}{2t}y+t^2
\end{array}
\right.$ sur $]0,+\infty[$\quad\item  $\left\{
\begin{array}{l}
(t^2+1)x'=tx-y+2t^2-1\\
(t^2+1)y'=x+ty+3t
\end{array}
\right.$

\item  $\left\{
\begin{array}{l}
\sh(2t)x'=\ch(2t)x-y\\
\sh(2t)y'=-x+\ch(2t)y
\end{array}
\right.$ sur $]0,+\infty[$ sachant qu'il existe une solution vérifiant $xy = 1$.
\end{enumerate}
\finenonce{005881}


\finexercice


\exercice{5882, rouget, 2010/10/16}
\enonce{005882}{*** I}
Résoudre les équations différentielles suivantes :

\begin{enumerate}
 \item  $(2x+1)y''+ (4x-2)y' - 8y = 0$ sur $\left]- \frac{1}{2},+\infty\right[$ puis sur $\Rr$.

\item  $(x^2+x)y''- 2xy'+ 2y = 0$ sur $]0,+\infty[$.

\item  $4xy''- 2y'+ 9x^2y = 0$ sur $]0,+\infty[$.

\item  $(1+x)y''- 2y'+ (1-x)y = xe^{-x}$ sur $]-1,+\infty[$.

\item  $y''+4y'+4y = \frac{e^{-2x}}{\sqrt{x^2+1}}$.

\item  $4xy''+2y'-y = 0$ sur $]0,+\infty[$.
\end{enumerate}
\finenonce{005882}


\finexercice


\exercice{5883, rouget, 2010/10/16}
\enonce{005883}{**}
 Trouver les fonctions $f$ dérivables sur $\Rr$ vérifiant $\forall x\in\Rr$, $f'(x) + f(-x) = e^x$.
\finenonce{005883}


\finexercice


\exercice{5884, rouget, 2010/10/16}
\enonce{005884}{***}
Trouver toutes les fonctions $f$ dérivables sur $]0,+\infty[$ vérifiant $\forall x >0$, $f'(x) = f\left( \frac{3}{16x}\right)$.
\finenonce{005884}


\finexercice


\exercice{5885, rouget, 2010/10/16}
\enonce{005885}{*** I}
Trouver toutes les applications $f~:~\Rr\rightarrow\Rr$, continues sur $\Rr$ telles que $\forall(x,y)\in\Rr^2$, $f(x)f(y) =\int_{x-y}^{x+y}f(t)\;dt$.
\finenonce{005885}


\finexercice


\exercice{5886, rouget, 2010/10/16}
\enonce{005886}{*** I}
 Montrer que $\forall x>0$, $\int_{0}^{+\infty} \frac{e^{-tx}}{1+t^2}\;dt=\int_{0}^{+\infty} \frac{\sin t}{t+x}\;dt$.
\finenonce{005886}


\finexercice

\section{225.05 Equations différentielles non linéaires}


\exercice{4110, quercia, 2010/03/11}
\enonce{004110}{\'Equations à variables séparables}

\begin{enumerate}
  \item $y' = y(1+y)$.
    

  \item $y' = \sin x\cos y$.
    

  \item $2yy'\sqrt x = \sqrt{y^2-1}$.
    


  \item $1+xy' = e^y$, condition initiale : $y(1)=1$.
    

  \item $y' = \sqrt{|y|}$ : étudier les problèmes de raccordements.
    

\end{enumerate}
\finenonce{004110}



\finexercice


\exercice{4111, quercia, 2010/03/11}
\enonce{004111}{\'Equations homogènes}

\begin{enumerate}
  \item $y-xy' = \sqrt{x^2+y^2}$.
    

  \item $y' = \frac{x-y}{x+y}$.
    

  \item   $(x^2+y^2)y' = 2xy$.
      

  \item $(x+y)y' = 2x-y$.
    

\end{enumerate}
\finenonce{004111}



\finexercice


\exercice{4112, quercia, 2010/03/11}
\enonce{004112}{\'Equations de Bernouilli}

\begin{enumerate}
  \item $xy'+y = xy^3$.
    

  \item $2xy'+y = \frac{2x^2}{y^3}$.
    

  \item $\sqrt x y' - y + (x+2\sqrt x)\sqrt y = 0$.
    

  \item $xy'+y = (xy)^{3/2}$.
    

  \item $x^3y' = y(3x^2+y^2)$.
    

\end{enumerate}
\finenonce{004112}



\finexercice


\exercice{4113, quercia, 2010/03/11}
\enonce{004113}{\'Equations de Riccati}

\begin{enumerate}
  \item $x^2(y'+y^2) = xy-1$.
    

\end{enumerate}
\finenonce{004113}



\finexercice


\exercice{4114, quercia, 2010/03/11}
\enonce{004114}{Divers ordre 1}
 $(x^2-y^2-1)y' = 2xy$ : poser $z = \frac{y}{x^2+y^2-1}$.
\finenonce{004114}



\finexercice


\exercice{4115, quercia, 2010/03/11}
\enonce{004115}{Centrale MP 2004}

Soit $n>0$ et $(S)\Leftrightarrow\begin{cases}x'(t) = \frac2nx(t)y(t)\cr y'(t) = -x^2(t)+y^2(t).\cr\end{cases}$

\begin{enumerate}
  \item Soit $\gamma$ : $t \mapsto(x(t),y(t))$ une solution de~$(S)$. Trouver
une autre solution présentant une symétrie avec $\gamma$. Peut-on
avoir comme solution $\sigma(t) = \lambda\gamma(\mu t)$~? En déduire
une propriété géométriques des solutions maximales de~$(S)$.

  \item Déterminer les courbes du plan formées des points $(x_0,y_0)$ où
les solutions de~$(S)$ ont des tangentes parallèles aux axes $(Ox)$ et
$(Oy)$. En déduire quelques solutions particulières.

  \item A supposer qu'il existe $\Phi : {I\subset \R} \to \R$ telle que
$\gamma(t) = (x(t),y(t))$ vérifie $y(t) = \Phi(x(t))$, déterminer
$\Phi$ et en déduire toutes les courbes intégrales.

\end{enumerate}
\finenonce{004115}



\finexercice


\exercice{4116, quercia, 2010/03/11}
\enonce{004116}{Chimie P 91}

Résoudre numériquement le système 
$$\begin{cases} y'=-y\cr z'=y-z\cr y(0)=1 \text{ et } z(0)=0. \end{cases}$$
Prendre $h = 0.1$ et faire un tableau avec $10$ valeurs.
Faire la résolution analytique.
\finenonce{004116}



\finexercice


\exercice{4117, quercia, 2010/03/11}
\enonce{004117}{\'Equations d'ordre 2}


\begin{enumerate}
  \item $y'' = \sin y$, $y(0) = \frac\pi2$, $y'(0) = \sqrt 2$.
    

  \item $2(2a-y)y'' = y'^2$.
    

  \item $yy''=y'^2-y^2$ : poser $z=y'/y$.
    
\end{enumerate}
\finenonce{004117}



\finexercice


\exercice{4118, quercia, 2010/03/11}
\enonce{004118}{Centrale MP 2000}

Existe-t-il des solutions de classe $\mathcal{C}^1$ sur $\R$ de l'équation différentielle : 
$y'+2\sqrt y =0$ ? Que peut-on dire de l'équation : $y'^2=4y$ ?
\finenonce{004118}



\finexercice


\exercice{4119, quercia, 2010/03/11}
\enonce{004119}{\'Etude qualitative : $y'=x^3+y^3$}

Soit $y$ la solution maximale de l'équation $y' =x^3+y^3$ telle que
$y(0) = a \ge 0$, et $I = {]\alpha,\beta[}$ son intervalle de définition.
Montrer que $y$ est strictement croissante sur $[0,\beta[$, que $\beta$ est borné,
et que $y \to +\infty$ lorsque $x\to\beta^-$.
\finenonce{004119}



\finexercice


\exercice{4120, quercia, 2010/03/11}
\enonce{004120}{\'Etude qualitative : $y'=x-e^y$}

Soit $y$ une solution maximale de l'équation $y'=x-e^y$.

\begin{enumerate}
  \item Montrer que $y$ est décroissante puis croissante.
    
  \item Montrer que $y$ est définie jusqu'en $+\infty$ et que sa courbe représentative
    admet une branche parabolique horizontale.
    
  \item Montrer que $\alpha > -\infty$ et que $y \to +\infty$ lorsque $x\to\alpha^-$.
    
\end{enumerate}
\finenonce{004120}



\finexercice


\exercice{4121, quercia, 2010/03/11}
\enonce{004121}{\'Etude qualitative : $x'=\cos(t) + \cos(x)$}

Soit $x$ la solution maximale du problème de Cauchy~:
$x'=\cos(t) + \cos(x)$, $x(0) = x_0 \in{]0,\pi[}$.

Montrer que $x$ est définie sur $\R$ et~: $\forall\ t>0,\ 0<x(t)<\pi$.
\finenonce{004121}



\finexercice


\exercice{4122, quercia, 2010/03/11}
\enonce{004122}{\'Etude qualitative : $x'=x^{2}-t$, ENS Cachan MP$^*$ 2005}

On consid{\`e}re l'{\'e}quationn diff{\'e}rentielle $(E) : x'=x^{2}-t$ et l'ensemble $D_{0}=\{(t,x)\,|\, x^{2}-t<0\}$.

Montrer que si $x$ est une solution de $(E)$ v{\'e}rifiant $(t_{0},x(t_{0}))\in D_{0}$, alors $x$ est d{\'e}finie sur
$[t_{0},+\infty[$ et la courbe int{\'e}grale reste dans $D_{0}$. En d{\'e}duire que $x(t)\mathop{\sim}\limits_{t\to +\infty} -\sqrt t$. 

\finenonce{004122}



\finexercice


\exercice{4123, quercia, 2010/03/11}
\enonce{004123}{Intervalle maximal pour $y' = f(y)$}

Soit $f:\R \to \R$ de classe $\mathcal{C}^1$ strictement positive et $y$ la solution
maximale définie sur $]\alpha,\beta[$ du problème de Cauchy : $y' = f(y)$,
$y(x_0) = y_0$.
Montrer que $\beta = x_0 +  \int_{t=y_0}^{+\infty} \frac{d t}{f(t)}$ et que
    $y \to +\infty$ lorsque $x\to\beta^-$.
\finenonce{004123}



\finexercice


\exercice{4124, quercia, 2010/03/11}
\enonce{004124}{\'Etude qualitative de $y' = 2ty + y^2$}

On considère l'équation : $y' = 2ty + y^2,\ y(t_0) = y_0$.
Soit $y$ une solution maximale.

\begin{enumerate}
  \item Montrer que $y = 0$ ou bien $y$ ne s'annulle pas.
  \item On choisit $y_0 > 0, t_0 < 0$. Soit $]t_1,t_2[$ le domaine d'existence de $y$.
   \begin{enumerate} 
    \item Montrer que si $y_0 \ge -2t_0$, alors $y$ est strictement croissante sur
      $[t_0,t_2[$.
    \item Montrer que $t_1 = -\infty$. (sinon, $y$ et $y'$ seraient bornées sur $]t_1,t_0]$.)
    \item Donner l'allure générale de la courbe de $y$.
  \end{enumerate}
  \item Résoudre l'équation en posant $z(t) = \frac {\exp(t^2)}{y(t)}$.

\end{enumerate}
\finenonce{004124}



\finexercice


\exercice{4125, quercia, 2010/03/11}
\enonce{004125}{\'Equation admettant simultanément $t$ et $\sin t$ comme solution}

Existe-t-il une fonction $f : (y,t) \to f(y,t)$ de classe $\mathcal{C}^1$ et
$n \in \N^*$ tels que l'équation : $y^{(n)} = f(y,t)$ admette les
deux solutions $y(t) = t$ et $y(t) = \sin t$ sur $\R$ ?


Même question avec l'équation $y^{(n)} = f(y^{(n-1)}, \dots, y, t )$.
\finenonce{004125}



\finexercice


\exercice{4126, quercia, 2010/03/11}
\enonce{004126}{$y'' = F(x,y)$, $y(a) = \alpha$, $y(b) = \beta$}


Soit $F : {[a,b]\times \R} \to \R$ de classe $\mathcal{C}^1$ telle que pour tout
$x \in {[a,b]}$, l'application $y \mapsto F(x,y)$ est strictement croissante.
Montrer que pour tous $\alpha,\beta \in \R$, il existe au plus une solution à
l'équation :
$y'' = F(x,y)$ avec les conditions aux limites $y(a) = \alpha$, $y(b) = \beta$.



\finenonce{004126}



\finexercice


\exercice{4127, quercia, 2010/03/11}
\enonce{004127}{Comparaison d'équations}

Soient $y,z$ solutions de $\begin{cases} y' = f(y,t)\cr z'=g(z,t)\cr y(0) = z(0) \cr \end{cases}$
où $f$,$g$ sont deux fonctions localement lipschitziennes telles que :
$$\begin{cases}
\forall\ u,t,\ f(u,t) \le g(u,t) \cr
\forall\ u,v,t,\ u \le v  \Rightarrow  f(u,t) \le f(v,t).\cr \end{cases}$$

\begin{enumerate}
  \item   Pour $\varepsilon > 0$, on note $z_\varepsilon$ la solution de
      $\begin{cases} z'_\varepsilon = g(z_\varepsilon,t) + \varepsilon \cr
               z_\varepsilon(0) = y(0).\cr\end{cases}$

      Montrer que $z_\varepsilon \ge y$ (sur leur domaine commun de définition).

  \item   Démontrer que $z_\varepsilon \to z$\ lorsque $varepsilon\to0^+$ uniformément sur tout intervalle
      borné. Conclusion ?


\end{enumerate}
\finenonce{004127}



\finexercice


\exercice{4128, quercia, 2010/03/11}
\enonce{004128}{\'Etude de l'équation}
    $\begin{cases} y'' + \sin y = 0 \cr  y(0) = 0,\ y'(0) = \alpha \ge 0.\cr\end{cases}$

Soit $y$ la solution maximale.
On a l'intégrale première : $\frac {y'^2}2 - \cos y = C = \alpha^2 - 1$.

\begin{enumerate}
  \item
  \begin{enumerate}
    \item Montrer que $y$ est définie sur $\R$.
    \item Montrer que $y$ est impaire.
  \end{enumerate}
  \item On suppose ici que $C > 1$.
  \begin{enumerate}
    \item Montrer qu'il existe un plus petit $T > 0$ tel que $y(T) = 2\pi$.
    \item Montrer que : $\forall\ t \in \R,\ y(t+T) = y(t) + 2\pi$.
  \end{enumerate}
  \item On suppose ici que $-1 < C < 1$: On pose $C = -\cos \theta$,
    et $F(x) =  \int_{u=0}^x \frac {d u}{\sqrt{2(\cos u - \cos \theta)}}$.
  \begin{enumerate}
    \item Soit $a$ maximal tel que $y'(t) > 0$ sur $[0,a[$.
      Montrer que $y(a) = \theta$ et $F(\theta) = a$.
    \item Montrer que $y$ est $4a$-périodique.
  \end{enumerate}
  \item \'Etudier les cas $C = 1$, $C = -1$.

\end{enumerate}
\finenonce{004128}



\finexercice


\exercice{4129, quercia, 2010/03/11}
\enonce{004129}{Résolution approchée de $y' = f(y,t),\ y(a) = y_0$  par la méthode d'Euler}

On suppose que $f$ est bornée par $M$ et $|f(y,s)-f(z,t)| \le K(|y-z|+|s-t|)$.
On divise $[a,b]$ en $n$ intervalles $[a_k, a_{k+1}]$,
$a_k = a + k\frac {b-a}n$ et on approche
la solution $y$ par la fonction $z$, continue affine par morceaux définie
par :
$$\begin{cases}
z(a_0) = y_0 \cr
\text{sur } ]a_k, a_{k+1}[,\ z' = f(z(a_k),a_k). \cr\end{cases}$$

\begin{enumerate}
  \item  Soit $\varepsilon_k = |z(a_k) - y(a_k)|$.
     Montrer que : $\forall\ t \in {[a_k, a_{k+1}]},\quad |y(t)-z(t)| \le kh^2(M+1) + (1+Kh)\varepsilon_k$
     $\Bigl(h = \frac {b-a}n \Bigr)$.
  \item  En déduire que $\sup|y-z| \le (M+1)(e^{K(b-a)}-1) \frac {b-a}n$.
\end{enumerate}
\finenonce{004129}



\finexercice


\exercice{4130, quercia, 2010/03/11}
\enonce{004130}{Lyon MP$^*$ 2000}

\begin{enumerate}
  \item Soit~$f$ une application minorée et de classe~$\mathcal{C}^1$ sur~$\R$,
à valeurs dans~$\R$. Montrer qu'il existe une suite $(a_n)$ telle
que la suite $(f'(a_n))$ tende vers~$0$.


  \item Soit~$f$ une application minorée et de classe~$\mathcal{C}^2$ sur~$\R^p$,
à valeurs dans~$\R$. Montrer qu'il existe une suite $(a_n)$ de~$\R^p$ telle
que la suite $(d f(a_n))$ tende vers~$0$, c'est à dire
$\nabla f(a_n)$ tend vers~$0$.

\end{enumerate}
\finenonce{004130}



\finexercice


\exercice{4131, quercia, 2010/03/11}
\enonce{004131}{ENS MP$^*$ 2001}

Soit un vecteur $v=(v_1,v_2,v_3)$ de~$\R^3$ muni de sa base canonique $(e_1,e_2,e_3)$.
Montrer qu'il existe une unique fonction $u = (u_1,u_2,u_3)$ de classe~$\mathcal{C}^1$
de~$\R$ dans~$\R^3$ telle que $u' + u\wedge u' = -u\wedge(u_3e_3)$ et $u(0) = v$.

{\it Indication : étudier la fonction $p \mapsto p + u\wedge p$ avant de pouvoir
évoquer le théorème de Cauchy-Lipschitz.}

\finenonce{004131}



\finexercice


\exercice{4132, quercia, 2010/03/11}
\enonce{004132}{Centrale MP 2001}

On définit une suite de fonctions sur $[0,1]$ de la manière suivante~:
$f_0$ est la fonction constante $1$ et pour tout $x\in{[0,1]}$ et $n\in\N$,
$f_{n+1}(x) = 1 +  \int_{t=0}^x f_n(t-t^2)\,d t$.

\begin{enumerate}
  \item En étudiant $f_{n+1}-f_n$ montrer que la suite $(f_n)$ converge
    uniformément sur~$[0,1]$. On note $f$ sa limite.
    

  \item Montrer que $f$ est de classe $\mathcal{C}^\infty$ sur $[0,1]$.
    Que valent $f'(0)$ et $f'(1)$~?
    

  \item \'Etudier la concavité de~$f$.
    

  \item Montrer que pour tout $x\in{[0,1]}$ on a $1+x \le f(x)\le \exp(x)$.
    

\end{enumerate}
\finenonce{004132}



\finexercice


\exercice{4133, quercia, 2010/03/11}
\enonce{004133}{ENS MP 2002}
Soient $f : {\R^2} \to \R, {(t,x)} \mapsto {f(t,x)}$ de classe $\mathcal{C}^1$,
et $a,b$ des réels tels que $a<b$. On suppose que $f$ 
est $T$-périodique par rapport à $t$ et que l'on a~:
$\forall\ t\in\R^+,\ f(t,a)>0\text{ et }f(t,b)<0$.

\begin{enumerate}
  \item Que peut-on dire des solutions du problème de Cauchy $E_y$ : $( x'(t)=f(t,x(t)),\ x(0)=y\in [a,b])$~?

  \item Montrer que toute solution maximale est définie sur $\R^+$ et à valeurs dans $[a,b]$.

  \item Montrer qu'il existe une solution de $E_y$ qui est $T$-périodique.

\end{enumerate}
\finenonce{004133}



\finexercice


\exercice{4134, quercia, 2010/03/11}
\enonce{004134}{X MP$^*$ 2005}

Soit $J$ un intervalle de~$\R$ et $f : {J\times \R^n} \to {\R^n}$ continue.
On suppose qu'il existe $a,b$ continues de $J$ dans $\R^+$ telles que,
pour tous $t,y$~: $(f(t,y)\mid y) \le a(t)\|y\|^2 + b(t)$.
Montrer que toute solution maximale de $y' = f(t,y)$ est définie sur $J$ entier.


\finenonce{004134}



\finexercice


\exercice{4135, quercia, 2010/03/11}
\enonce{004135}{Système autonome, ENS Ulm-Lyon-Cachan MP$^*$ 2006}
On considère le système différentiel : $$(V) \Leftrightarrow\left\{\begin{array}{l} x' = x(1-y) \\ y' = y(x-1).\\ \end{array}\right.$$
dont on cherche les solutions $(x, y)$ définies sur $\R$ à valeurs dans $(\R^{+*})^{2}$.

\begin{enumerate}
  \item Trouver une fonction $f \in \mathcal{C}^{2}((\R^{+*})^{2}, \R)$ telle que pour toute solution
$(x, y)$ de $V$, $f(x, y)$ soit constante.

  \item Montrer que les solutions de $(V)$ sont périodiques.

\end{enumerate}
\finenonce{004135}



\finexercice

\section{225.06 Equations aux dérivées partielles}


\exercice{4201, quercia, 2010/03/11}
\enonce{004201}{$2\frac{\partial f}{\partial x} + 3\frac{\partial f}{\partial y} = 4f$}

Résoudre l'équation $2\frac{\partial f}{\partial x} + 3\frac{\partial f}{\partial y} = 4f$ avec la {\it condition aux limites\/}
: $f(t,t) = t\ (t \in \R)$.

(\'Etudier $\varphi : t  \mapsto f(a+bt,a+ct)$ avec $a,b,c$ bien choisis)

\finenonce{004201}



\finexercice


\exercice{4202, quercia, 2010/03/11}
\enonce{004202}{$\frac{\partial f}{\partial x} - \frac{\partial f}{\partial y} =$ cste}

Déterminer les applications $f$ de classe $\mathcal{C}^1$ de ${\R^2}$ dans $\R$ vérifiant :
$\frac{\partial f}{\partial x} - \frac{\partial f}{\partial y} = a$ où $a$ est une constante réelle donnée.
On utilisera le changement de variable : $u = x+y$, $v = x-y$.

\finenonce{004202}



\finexercice


\exercice{4203, quercia, 2010/03/11}
\enonce{004203}{$x\frac{\partial f}{\partial x} = y\frac{\partial f}{\partial y}$}

Résoudre sur $(\R^{+*})^2$ : $x\frac{\partial f}{\partial x} = y\frac{\partial f}{\partial y}$,
en posant $\begin{cases} u = xy \cr v = \frac xy. \end{cases}$

\finenonce{004203}



\finexercice


\exercice{4204, quercia, 2010/03/11}
\enonce{004204}{$x\frac{\partial f}{\partial x} + y\frac{\partial f}{\partial y} = 2$}
Soit $U$ l'ouvert de ${\R^2}$ : $U = \{(x,y) \text{ tq } x > 0,\ y > 0\}$.
Trouver les applications $f : U \to \R$ de classe $\mathcal{C}^1$
vérifiant~: $x\frac{\partial f}{\partial x} + y\frac{\partial f}{\partial y} = 2$.
On utilisera le changement de variable : $u = xy$, $v = \frac yx$.

\finenonce{004204}



\finexercice


\exercice{4205, quercia, 2010/03/11}
\enonce{004205}{$x\frac{\partial f}{\partial x} = -y\frac{\partial f}{\partial y}$}
Résoudre sur ${\R^2} \setminus \{(0,0)\}$ : $x\frac{\partial f}{\partial x} = -y\frac{\partial f}{\partial y}$,
en posant $\begin{cases}x = \rho\cos\theta \cr y = \rho\sin\theta. \end{cases}$

\finenonce{004205}



\finexercice


\exercice{4206, quercia, 2010/03/11}
\enonce{004206}{$y\frac{\partial f}{\partial x} - x\frac{\partial f}{\partial y} = 2f$}

Soit $f : U \to \R$ de classe $\mathcal{C}^1$ vérifiant :
$y\frac{\partial f}{\partial x} - x\frac{\partial f}{\partial y} = 2f$. où $U$ est un ouvert de ${\R^2}$.

On pose $g(\rho,\theta) = f(\rho\cos\theta,\rho\sin\theta)$.
Calculer $\frac{\partial g}{\partial \rho}$, $\frac{\partial g}{\partial \theta}$, puis trouver $f$ \dots

\begin{enumerate}
  \item Si $U = \{(x,y) \text{ tq } x > 0\}$.
    
  \item Si $U = {\R^2}$.
\end{enumerate}
\finenonce{004206}



\finexercice


\exercice{4207, quercia, 2010/03/11}
\enonce{004207}{Ensi Physique P 94}

    Résoudre l'équation aux dérivées partielles suivante~:
    $2xy\frac{\partial f}{\partial x} + (1+y^2)\frac{\partial f}{\partial y} = 0$
    en utilisant, par exemple, le changement de variable~:
    $x=\frac{u^2+v^2}2$ et $y=\frac uv$.
\finenonce{004207}



\finexercice


\exercice{4208, quercia, 2010/03/11}
\enonce{004208}{Fonctions homogènes}

Soit $\Omega = \{ (x,y) \in {\R^2} \text{ tq } (x,y) \ne (0,0) \}$, et
$f : \Omega \to \R$ une fonction de classe $\mathcal{C}^1$.

Montrer que $f$ est positivement homogène de degré $\alpha$ si et seulement si :
$$\forall\ (x,y)\in\Omega,\ x\frac{\partial f}{\partial x}(x,y) + y\frac{\partial f}{\partial y}(x,y) = \alpha f(x,y).$$
(On étudiera $g(\rho,\theta) = f(\rho\cos\theta,\rho\sin\theta)$)
\finenonce{004208}



\finexercice


\exercice{4209, quercia, 2010/03/11}
\enonce{004209}{}
Résoudre l'équation : $x^2\frac{\partial^2 f}{\partial x^2} + 2xy\frac{\partial^2 f}{\partial x \partial y} + y^2\frac{\partial^2 f}{\partial y^2} = \alpha(\alpha-1)f$
où $\alpha$ est un réel fixé, $\alpha \ne \frac 12$.
On posera $x=\rho\cos\theta$, $y=\rho\sin\theta$.

\finenonce{004209}



\finexercice


\exercice{4210, quercia, 2010/03/11}
\enonce{004210}{\'Equation d'ordre 2 à coefficients constants}

Soient $a,b,c \in \R$ non tous nuls. On considère l'équation aux dérivées
partielles : $$(*) \Leftrightarrow a\frac{\partial^2 f}{\partial x^2} + b\frac{\partial^2 f}{\partial x \partial y} + c\frac{\partial^2 f}{\partial y^2} = 0$$ où $f$ est une
fonction inconnue : ${\R^2} \to \R$ de classe $\mathcal{C}^2$.
Soient $\alpha,\beta\in \R$ distincts, fixés. On fait le changement de
variable : $u = x+\alpha y$, $v = x+\beta y$.

\begin{enumerate}
  \item \'Ecrire l'équation déduite de $(*)$ par ce changement de variable.
    
  \item En déduire que l'on peut ramener $(*)$ à l'une des trois formes réduites :\par
    $(1) : \frac{\partial^2 g}{\partial u \partial v}  = 0$,  \qquad
    $(2) : \frac{\partial^2 g}{\partial u^2}  = 0$,  \qquad
    $(3) : \frac{\partial^2 g}{\partial u^2} + \frac{\partial^2 g}{\partial v^2} = 0$.
\end{enumerate}
\finenonce{004210}



\finexercice


\exercice{4211, quercia, 2010/03/11}
\enonce{004211}{$x^2\frac{\partial^2 f}{\partial x^2} + 2xy\frac{\partial^2 f}{\partial x \partial y} + y^2\frac{\partial^2 f}{\partial y^2} = 0$}
Trouver les applications $f : {(\R^{+*})^2} \to \R$ de classe $\mathcal{C}^2$
vérifiant : $x^2\frac{\partial^2 f}{\partial x^2} + 2xy\frac{\partial^2 f}{\partial x \partial y} + y^2\frac{\partial^2 f}{\partial y^2} = 0$.
On utilisera le changement de variables : $u = xy$, $v = \frac xy$.


\finenonce{004211}



\finexercice


\exercice{4212, quercia, 2010/03/11}
\enonce{004212}{$\frac{\partial^2 f}{\partial x^2} + \frac{\partial^2 f}{\partial y^2} = \frac y{x^3}$}

Soit $g : \R \to \R$ de classe $\mathcal{C}^2$ et
$f : {\R^*\times \R} \to \R, {(x,y)} \mapsto {g(y/x).}$
Trouver $g$ telle que $\frac{\partial^2 f}{\partial x^2} + \frac{\partial^2 f}{\partial y^2} = \frac y{x^3}$.


\finenonce{004212}



\finexercice


\exercice{4213, quercia, 2010/03/11}
\enonce{004213}{$\frac{\partial^2 g}{\partial x^2} - 4\frac{\partial^2 g}{\partial y^2} = 1$}

Soit $f : {\R^2} \to \R$ de classe $\mathcal{C}^2$.
On pose $g(x,y) = f(2x+y,2x-y)$.

\begin{enumerate}
  \item Calculer les dérivées partielles secondes de $g$ en fonction de celles
    de $f$.
    

  \item Trouver $f$ telle que $\frac{\partial^2 g}{\partial x^2} - 4\frac{\partial^2 g}{\partial y^2} = 1$.
    
\end{enumerate}
\finenonce{004213}



\finexercice


\exercice{4214, quercia, 2010/03/11}
\enonce{004214}{$x^2\frac{\partial^2 f}{\partial x^2} - y^2 \frac{\partial^2 f}{\partial y^2} = 0$}

On considère l'équation aux dérivées partielles sur $\Omega = (\R^{+*})^2$ :
$x^2\frac{\partial^2 f}{\partial x^2} - y^2 \frac{\partial^2 f}{\partial y^2} = 0$.

Résoudre cette équation en posant $\begin{cases} u = xy \cr v = \frac xy. \end{cases}$
\finenonce{004214}



\finexercice


\exercice{4215, quercia, 2010/03/11}
\enonce{004215}{$f(\cos x/\ch y)$ harmonique}


Soit $f : {]-1,1[} \to \R$ de classe $\mathcal{C}^2$. On considère
$g : {D \subset {\R^2}} \to \R, {(x,y)} \mapsto {f\Bigl(\frac {\cos x}{\ch y}\Bigr).}$

Déterminer $f$ pour que $g$ vérifie : $\frac{\partial^2 g}{\partial x^2} + \frac{\partial^2 g}{\partial y^2} = 0$.

\finenonce{004215}



\finexercice


\exercice{4216, quercia, 2010/03/11}
\enonce{004216}{$(x,y)  \mapsto (x^2-y^2,2xy)$ préserve les fonctions harmoniques}

Pour $(x,y) \in {\R^2}$, on pose $u = x^2-y^2$, $v = 2xy$.

Soit $F : {\R^2} \to \R, {(u,v)} \mapsto {F(u,v)}$ et $f$ définie par :
$f(x,y) = F(u,v)$.

Montrer que $\frac{\partial^2 F}{\partial u^2} + \frac{\partial^2 F}{\partial v^2} = 0$ entraîne $\frac{\partial^2 f}{\partial x^2} + \frac{\partial^2 f}{\partial y^2} = 0$.

\finenonce{004216}



\finexercice


\exercice{4217, quercia, 2010/03/11}
\enonce{004217}{$x\frac{\partial^2 f}{\partial x^2} + y\frac{\partial^2 f}{\partial x \partial y} + \frac{\partial^2 f}{\partial y^2} + \frac{\partial f}{\partial x} = y$}
Soit $f : {\R^2} \to \R$ de classe $\mathcal{C}^2$ et $g(u,v) = f(uv,u+v)$.

\begin{enumerate}
  \item Calculer $\frac{\partial^2 g}{\partial u \partial v}$.
    
  \item Résoudre l'équation :
    $x\frac{\partial^2 f}{\partial x^2} + y\frac{\partial^2 f}{\partial x \partial y} + \frac{\partial^2 f}{\partial y^2} + \frac{\partial f}{\partial x} = y$.
    
\end{enumerate}
\finenonce{004217}



\finexercice


\exercice{4218, quercia, 2010/03/11}
\enonce{004218}{$f(\sqrt{x^2+y^2+z^2})$ tq $\Delta f = -f$}

Soit $f:\R \to \R$ de classe $\mathcal{C}^2$ et $g(x,y,z) = \frac {f(r)}r$
avec $r = \sqrt{x^2+y^2+z^2}$.
Déterminer $f$ de sorte que $\frac{\partial^2 g}{\partial x^2} + \frac{\partial^2 g}{\partial y^2} + \frac{\partial^2 g}{\partial z^2} = -g$.

\finenonce{004218}



\finexercice


\exercice{4219, quercia, 2010/03/11}
\enonce{004219}{$x^4 \frac{\partial^2 f}{\partial x^2} - \frac{\partial^2 f}{\partial y^2} = 0$}

\begin{enumerate}
  \item Trouver les fonctions $g : {\{(u,v)\in\R^2\text{ tq } u > v\}} \to \R$ de classe $\mathcal{C}^2$ vérifiant :
    $\frac{\partial }{\partial u}\Bigl(g+v\frac{\partial g}{\partial v}\Bigr) = \frac{\partial}{\partial  v}\Bigl(g+u\frac{\partial g}{\partial u}\Bigr)$

    (penser au théorème de Poincaré).
    

  \item Résoudre sur $\R^{+*}\times\R$ l'équation : $x^4 \frac{\partial^2 f}{\partial x^2} - \frac{\partial^2 f}{\partial y^2} = 0$ en posant
    $u = y+\frac 1x$, $v = y-\frac 1x$.
    
\end{enumerate}
\finenonce{004219}



\finexercice


\exercice{5898, rouget, 2010/10/16}
\enonce{005898}{*** I}
Résoudre les équations aux dérivées partielles suivantes :

\begin{enumerate}
 \item  $2 \frac{\partial f}{\partial x}- \frac{\partial f}{\partial y}= 0$ en posant $u = x+y$ et $v = x+2y$.

\item  $x \frac{\partial f}{\partial y}-y \frac{\partial f}{\partial x}=0$ sur $\Rr^2\setminus\{(0,0)\}$ en passant en polaires.

\item  $x^2 \frac{\partial^2f}{\partial x^2}+2xy \frac{\partial^2f}{\partial x\partial y}+y^2 \frac{\partial^2f}{\partial y^2}= 0$ sur $]0, +\infty[\times\Rr$ en posant $x = u$ et $y = uv$.
\end{enumerate}
\finenonce{005898}


\finexercice

\section{225.99 Autre}

\section{229.01 Ouvert, fermé, intérieur, adhérence}


\exercice{1741, drutu, 2003/10/01}

\enonce{001741}{}
Repr\'esenter graphiquement et d\'eterminer 
si les ensembles suivants sont des ouverts.

$A=\{(x,y)\in \R^2 \mid 0<\vert x-1\vert <1 \}$ ;   $B=\{(x,y)\in \R^2 \mid 0<x\leq 1 \}$ ;

$C=\{(x,y)\in \R^2 \mid \vert x\vert <1,\; \vert y\vert \leq 
1 \}$ ;   $D=\{(x,y)\in \R^2 \mid x\in \Q ,y\in \Q \}$ ;

$E=\{(x,y)\in \R^2 \mid x\not\in \Q ,y\not\in \Q  \}$ ;   $F=\{(x,y)\in \R^2 \mid x^2+y^2 <4 \}$ .
\finenonce{001741}



\finexercice



\exercice{1742, drutu, 2003/10/01}

\enonce{001742}{}
Montrer que toute reunion et toute intersection finie 
d'ensembles ouverts est un ensemble ouvert. 
Que peut-on dire des intersections infinies d'ensembles ouverts ~?
\finenonce{001742}



\finexercice



\exercice{1743, maillot, 2001/09/01}

\enonce{001743}{partiel 1999}
On d\'efinit un sous-ensemble $A$ de $\R^2$ en posant
$$A=\{(x,y)\in\R^2\mid x^2+y^2\le 2\} \setminus \{(x,y)\in \R^2
\mid (x-1)^2+y^2<1\}.$$
D\'eterminer l'int\'erieur, l'adh\'erence et la fronti\`ere de $A$.
L'ensemble $A$ est-il connexe~?
\finenonce{001743}



\finexercice



\exercice{1754, maillot, 2001/09/01}

\enonce{001754}{}
Les sous-ensembles de $\R^2$ suivants sont-ils ouverts~? Fermés~?
Compacts~?
\begin{gather*}
A=\{(x,y)\in \R^2 \mid x^2-\sin(y)\le 4\} \\
B=\{(x,y)\in \R^2 \mid x^3-4e^y> 4\} \\
C=\{(x,y)\in {[0,1]\times [0,1]} \mid \cos(x)\ge 0\}
\end{gather*}
\finenonce{001754}



\finexercice



\exercice{1755, maillot, 2001/09/01}

\enonce{001755}{}
On se propose de montrer que tout ouvert de $\R$ est une réunion
d'intervalles ouverts disjoints. On considère donc un ouvert
$U\subset\R$ et pour tout $x\in U$ on pose $$C(x)=\{y\in [x,+\infty[\ \mid
[x,y]\subset U\} \cup \{y\in ]-\infty,x[\ \mid [y,x]\subset U\}.$$
\begin{enumerate}
\item Montrer que $C(x)$ est un intervalle ouvert pour tout $x$.
(Considérer $\inf_{y\in C(x)}y$ et $\sup_{y\in C(x)}y$.)
\item Pour tous $x,y$ dans $U$, montrer qu'on a $C(x)=C(y)$ ou
$C(x)\cap C(y)=\emptyset$.
\item Conclure.
\end{enumerate}
\finenonce{001755}



\finexercice



\exercice{1756, maillot, 2001/09/01}

\enonce{001756}{}
Soit $E$ un espace vectoriel norm\'e. Soient $A$ et $B$ deux parties de $E$.
Montrer~:
\begin{enumerate}
\item $C_{\stackrel{\circ}{A}}= \overline{C_A}$, $C_{\bar A}= \stackrel{\circ}
{C_A}$
\item $\overline{A \cup B} = \bar A\cup \bar B$\par
En \textbf{d\'eduire} $\stackrel{\circ}{\overbrace{A\cap
B}}=\stackrel{\circ}{A}
\cap \stackrel{\circ}{B}$.
\item $\overline{A \cap B} \subset \bar A\cap \bar B$\par
En \textbf{d\'eduire} $\stackrel{\circ}{A}
\cup \stackrel{\circ}{B} \subset \stackrel{\circ}{\overbrace{A\cup B}}$.\par
Donner un exemple pour lequel l'inclusion r\'eciproque n'est pas
r\'ealis\'ee.
\end{enumerate}
\finenonce{001756}



\finexercice



\exercice{1757, maillot, 2001/09/01}

\enonce{001757}{}
Soit $A$ une partie d'un espace vectoriel norm\'e $E$. On
rappelle que la fronti\`ere de $A$ est l'ensemble $\mathrm{Fr}(A)=\bar A-
\stackrel{\circ}{A}$. Montrer que~:
\begin{enumerate}
\item $ \mathrm{Fr}(A)=\{x\in E \mid \forall \epsilon>0, B(x,\epsilon)\cap A
\neq\emptyset \text{\ et\ } B(x,\epsilon)\cap C_A\neq\emptyset\}$
\item $\mathrm{Fr}(A)=\mathrm{Fr}(C_A)$
\item $A$ est ferm\'e si et seulement si $\mathrm{Fr}(A)$ est inclus dans $A$.
\item $A$ est ouvert si et seulement si $\mathrm{Fr}(A)\cap A=\emptyset$.
\end{enumerate}
\finenonce{001757}



\finexercice



\exercice{1758, maillot, 2001/09/01}

\enonce{001758}{}
Soit $A$ une partie d'un espace vectoriel norm\'e $E$.
\begin{enumerate}
\item Montrer que $\bar A$ est l'ensemble des limites de suites
convergentes d'\'el\'ements de $A$.
\item On suppose maintenant que $E=\R$. D\'eduire de la question
pr\'ec\'edente que si $A$ est born\'ee, alors $\sup A\in \bar A$.
(Construire une suite de points appropri\'ee.)
\end{enumerate}
\finenonce{001758}



\finexercice



\exercice{1759, maillot, 2001/09/01}

\enonce{001759}{}
Montrer que l'adh\'erence d'une boule ouverte est la boule ferm\'ee
de m\^eme centre et m\^eme rayon.
\finenonce{001759}



\finexercice



\exercice{1760, maillot, 2001/09/01}

\enonce{001760}{}
Soit $E$ un espace vectoriel norm\'e. Soient $A$ et $B$ deux parties de $E$.
On pose $A+B= \{z\in E \mid \exists x\in A, \exists y\in B, z=x+y\}$.

Montrer que si $A$ est ouvert, $A+B$ est ouvert. (Commencer par le
cas o\`u $B$ est un singleton.)
\finenonce{001760}



\finexercice



\exercice{1761, maillot, 2001/09/01}

\enonce{001761}{}
Soit $E$ un espace vectoriel norm\'e de dimension finie.

Montrer que tout sous-espace vectoriel de $E$ est ferm\'e.
\finenonce{001761}



\finexercice



\exercice{1762, maillot, 2001/09/01}

\enonce{001762}{}
Soit $E$ un espace vectoriel norm\'e. Soit $A$ une partie non vide et
born\'ee de $E$. On d\'efinit $\mathrm{diam}(A)=\sup \{\|y-x\|, x,y\in A\}$.
\begin{enumerate}
\item Montrer que si $A$ est born\'ee, alors $\bar A$ et $\mathrm{Fr}(A)$
sont born\'es.
\item Comparer $\mathrm{diam}(A)$, $\mathrm{diam}(\stackrel{\circ}{A})$ et $\mathrm{diam}(\bar A)$
lorsque $\stackrel{\circ}{A}$ est non vide.
\item \begin{enumerate}
   \item Montrer que $\mathrm{diam}(\mathrm{Fr}(A)) \le \mathrm{diam}(A)$.
   \item Soit $x$ et $u$ des \'el\'ements de $A$ avec $u\neq 0$. On
consid\`ere l'ensemble $X=\{t\ge 0 \mid x+tu\in A\}$. Montrer que
$\sup X$ existe.
   \item En d\'eduire que toute demi-droite issue d'un point $x$ de
$A$ coupe $\mathrm{Fr}(A)$.
   \item En d\'eduire que $\mathrm{diam}(\mathrm{Fr}(A)) = \mathrm{diam} (A)$.
  \end{enumerate}
\end{enumerate}
\finenonce{001762}



\finexercice



\exercice{1764, maillot, 2001/09/01}

\enonce{001764}{}
Soit $E=\R^d$ muni d'une norme $\|\cdot\|$. On d\'efinit la
\emph{distance} d'un \'el\'ement $x_0$ de $E$ \`a une partie $A$ de $E$,
not\'ee $d(x_0,A)$, par la formule $$d(x_0,A)=\inf_{x\in A}\|x-x_0\|.$$
\begin{enumerate}
\item Supposons $A$ compact. Montrer que pour tout $x_0\in E$
il existe $y\in A$ tel que $d(x_0,A)=\|y-x_0\|$.
\item Montrer que le r\'esultat est encore vrai si on suppose seulement
que $A$ est ferm\'e. (On remarquera que pour toute partie $B$ de $A$ on
a $d(x_0,B)\ge d(x_0,A)$.)
\item Montrer que l'application qui \`a $x_0$ associe $d(x_0,A)$ est
continue sur $E$ (sans aucune hypoth\`ese sur $A$).
\item En d\'eduire que si $A$ est un ferm\'e de $E$ et $B$ un compact de
$E$ tels que $A$ et $B$ sont disjoints, alors il existe une
constante $\delta>0$ telle que
$$\|a-b\| \ge \delta \qquad \forall (a,b)\in A\times B.$$
\item Montrer par un contre-exemple que le r\'esultat est faux si on
suppose seulement que $A$ et $B$ sont deux ferm\'es disjoints.
\end{enumerate}
\finenonce{001764}



\finexercice



\exercice{1766, maillot, 2001/09/01}

\enonce{001766}{}
Soit $(E,\|\cdot\|)$ un espace vectoriel norm\'e. Pour toutes parties
$A$ et $B$ de $E$ on note $$A+B=\{z\in E \mid \exists(x,y)\in A\times
B, z=x+y\}.$$

Montrer que si $A$ est compact et $B$ ferm\'e, alors $A+B$ est ferm\'e.
\finenonce{001766}



\finexercice



\exercice{1769, gourio, 2001/09/01}

\enonce{001769}{}
Soit $X$  une partie de ${\Rr}^{2}$; montrer qu'elle est
ferm\'{e}e si et seulement si pour toute partie ferm\'{e}e born\'{e}e
$K, K\cap X$ est ferm\'{e}e born\'{e}e.
\finenonce{001769}



\finexercice



\exercice{1770, gourio, 2001/09/01}

\enonce{001770}{}
Soient $k\in {\Rr}^{+*},$
$$ \omega _{n}=\left\{ \left( x,y\right) \in {\Rr}^{2}|\left( x-\frac{1}{n}%
\right) ^{2}+\left( y-\frac{1}{n}\right) ^{2}\leq \frac{k^{2}}{n^{2}}%
\right\} , $$
et
$$\Omega =\bigcup\limits_{n\in \Nn^{*}}\omega _{n}. $$
$\Omega $ est-il ouvert? ferm\'{e}? ...
\finenonce{001770}



\finexercice



\exercice{1771, gourio, 2001/09/01}

\enonce{001771}{}
Soit $\left( K_{n}\right) _{n\in \Nn^{*}}$ une suite d'ensembles
ferm\'{e}s born\'{e}s de ${\Rr}^{2}$ telle que $\forall n\in \Nn%
,K_{n+1}\subset K_{n},$ et $K_{n}\neq \emptyset .$

Montrer que :
$$\bigcap\limits_{n\in \Nn^{*}}K_{n}\neq \emptyset . $$
\finenonce{001771}



\finexercice



\exercice{1772, gourio, 2001/09/01}

\enonce{001772}{}
Montrer que l'intersection de deux ensembles ouvert est ouvert, que l'union
de deux ensembles ferm\'{e}s est ferm\'{e}e, que cela reste vrai pour un
nombre fini d'ensembles, mais que cela peut devenir faux si l'on
consid\`{e}re des suites infinies.
\finenonce{001772}



\finexercice



\exercice{1773, gourio, 2001/09/01}

\enonce{001773}{}
Soit $E\subset {\Rr}^{2} $ un ensemble ; on pose
$$\mathrm{Int}(E)=^{c}\overline{^{c}E}. $$
Montrer que $\mathrm{Int}(E)$ est le plus grand ouvert contenu dans $E.$
\finenonce{001773}



\finexercice



\exercice{1774, gourio, 2001/09/01}

\enonce{001774}{}
Soit $ A$ une partie born\'{e}e de $ {\Rr}^{2}, $  montrer que $ \overline{A}$
est aussi born\'{e}e et que
$$\sup\limits_{x\in A}\left\| x\right\| =\sup\limits_{x\in \overline{A}%
}\left\| x\right\| . $$
\finenonce{001774}



\finexercice



\exercice{1776, gourio, 2001/09/01}

\enonce{001776}{}
Classer (pour l'inclusion) les parties : $ \overline{A\cap B}, \overline{A}
\cap \overline{B}  $ et $ \overline{A\cup B}, \overline{A}\cup \overline{B}. $
\finenonce{001776}



\finexercice



\exercice{1777, gineste, 2001/11/01}

\enonce{001777}{}
Dans l'espace vectoriel norm\'e $\R$, chacune des parties suivantes est-elle ouverte ? ferm\'ee ? \
\par
$ \N, \Z, \Q, \R, [0,1[, [0,+\infty[, ]0,1[ \cup \{2\}, \{ 1/n, \, n \in \N^* \}, \bigcap_{n \geq 1} ]-1/n,1/n[ .$
\finenonce{001777}



\finexercice



\exercice{1778, gineste, 2001/11/01}

\enonce{001778}{}
Soit $E$ un evn (espace vectoriel norm\'e). Soit $A$ une partie de $E$. Montrer l'\'egalit\'e
$$ E \backslash \overline{A} = \stackrel{\circ}{\stackrel{\frown} {E\backslash A}} \mbox{ et  } E \backslash \stackrel{\circ}{A}= \overline{E \backslash A } $$
\finenonce{001778}



\finexercice



\exercice{1779, gineste, 2001/11/01}

\enonce{001779}{}
Soit $E$ un evn, $V$ un  sous-espace vectoriel de $E$.
\begin{enumerate}
\item Montrer que $\overline{V}$ est un sous-espace vectoriel de $E$.
\item Montrer que si $\stackrel{\circ}{V} \not = \varnothing$  alors $V=E.$
\end{enumerate}
\finenonce{001779}



\finexercice



\exercice{1780, gineste, 2001/11/01}

\enonce{001780}{}
Représenter graphiquement les parties suivantes de $\R^2$ et dire pour chacune d'elle si c'est un ouvert, un fermé, ou ni l'un ni l'autre. Déterminer leurs adhérences et intérieurs.
\begin{enumerate}
\item $$ \{ (x,y) \in \R^2 \; , \;|x| \not = 1 \mbox{ et } |y| \not = 1 \} $$
\item $$ \{ (x,y) \in \R^2 \; , \; |x|  = 1 \mbox{ et } |y| \not = 1 \} $$
\item $$ \{ (x,y) \in \R^2 \; , \; |x| \not = 1 \mbox{ ou } |y| \not = 1 \} $$
\item $$ \{ (x,y) \in \R^2 \; , \; 1 - xy >  0 \} $$
\item $$ \{ (x,y) \in \R^2 \; , \; 3x+4y=2  \} $$
\item $$ \{ (x,y) \in \R^2 \; , x^2 + y^2 = 1 \} $$
\item $$ \{ (x,y) \in \R^2 \; , xy=1 \} $$
\item $$ \bigcup_{n \in \N^*} \{ 1/n \} \times [0,1] $$
\end{enumerate}
\finenonce{001780}



\finexercice



\exercice{1781, gineste, 2001/11/01}

\enonce{001781}{}
Déterminer l'adhérence de chacune des parties de $\R$ suivantes :
\begin{enumerate}
\item $\N, \Z, \Q $
\item $ \{ 1/n, \, n \in \N^* \} $
\item $ \{ \frac{ (-1)^n }{ 1 + 1/n } , \, n \in \N^* \} $
\end{enumerate}
\finenonce{001781}



\finexercice



\exercice{1782, gineste, 2001/11/01}

\enonce{001782}{}
Soient $A$ et $B$, deux parties d'un evn $E$.
\begin{enumerate}
\item Montrer que si $O$ est un ouvert de $E$, alors $A+O$ est ouvert. (Indication : Prendre d'abord $A=\{ a \}$ puis $A$ quelconque .... )
\item Etablir que $\overline{A \cup B}=\overline{A} \cup \overline{B}$ et que  $\overline{A \cap B} \subset \overline{A} \cap \overline{B}$. (Trouver un exemple o\`u l'inclusion est stricte)
\end{enumerate}
\finenonce{001782}



\finexercice



\exercice{2616, debievre, 2009/05/19}
\enonce{002616}{}
\begin{enumerate}
  \item Tracer le graphe de la fonction 
$f\colon \R^2 \longrightarrow\R$ d\'efinie par
$f(x,y)= x^2+y^2$ et tracer les lignes de niveau de cette fonction.
  \item Tracer les graphes des fonctions $f$ et $g$ d\'efinies 
par $f(x,y)=25-(x^2+y^2)$  et $g(x,y)=5-\sqrt{x^2+y^2}$ sur 
$D=\{(x,y)\in\R^2\mid x^2+y^2\leq 25\}$. 
  \item  Tracer le graphe de la courbe  param\'etr\'ee
$f\colon \R\longrightarrow \R^2$ d\'efinie par
$f(x)= (x\cos x,  x\sin x)$.

  \item Peut-on repr\'esenter graphiquement 
%et  de fa\c con  diff\'erente 
l'application de la question (3.)? Comment?  
  \item  D\'ecrire les surfaces de niveau de la fonction 
$f\colon \R^3\longrightarrow  \R$ d\'efinie par
$f(x,y,z)= \exp (x+y^2-z^2)$.
  \item Pourquoi ne peut-t-on pas 
na\"\i vement
repr\'esenter le graphe
de l'application 
\[f\colon\R^2\longrightarrow \R^2,\  
f(x,y)= (-y, x),
\] 
sur une feuille de
papier. Comment peut-on graphiquement repr\'esenter cette application?
\end{enumerate}
\finenonce{002616}


\finexercice


\exercice{2617, debievre, 2009/05/19}
\enonce{002617}{}

D\'eterminer si chacune des parties suivantes du plan
sont ouvertes ou
ferm\'ees, ou ni l'un ni l'autre. D\'eterminer chaque fois l'int\'erieur 
et l'adh\'erence. 
\begin{enumerate}
  \item $ A_1=\{ (x, y)\in\R^2 | x^2y^2 > 1\}$,
  \item $ A_2=\{ (x,y)\in\R^2 | x^2+y^2=1, y>0\}$.
\end{enumerate}
\finenonce{002617}


\finexercice


\exercice{2618, debievre, 2009/05/19}
\enonce{002618}{}
\begin{enumerate}
  \item  Soient $B_1\subset \R^n$ et $B_2\subset \R^m$ des boules 
ouvertes. Montrer
que $B_1\times B_2\subset \R^{n+m}$ est un ouvert.
  \item Soit $A$ un ouvert de $\R^2$ et $B$ un ouvert de $\R$. Montrer 
que $A\times B$
est un ouvert de $\R^3$. 
\end{enumerate}
\finenonce{002618}


\finexercice


\exercice{2619, debievre, 2009/05/19}
\enonce{002619}{}
\begin{enumerate}
  \item Soit $(A_n)$ ($n\in\N$) 
une suite de parties ouvertes de $\R^2$. Est-ce que la
r\'eunion des $A_n$ est encore une partie ouverte? Et leur intersection?
  \item M\^eme question pour une famille de parties ferm\'ees. 
\end{enumerate}
\finenonce{002619}


\finexercice


\exercice{2620, debievre, 2009/05/19}
\enonce{002620}{}
Soit $A=\{(t,\sin\frac1t)\in\R^2 ; t>0 \}$. Montrer que $A$
n'est ni ouvert ni ferm\'e. D\'eterminer l'adh\'erence $\overline A$
de $A$. 
\finenonce{002620}


\finexercice

\exercice{2647, debievre, 2009/05/19}
\enonce{002647}{}
Soit $D$ un sous-ensemble de $\R^n$.
\begin{enumerate}
 \item Donner la d\'efinition de ``$D$ est ouvert.'' (Ceci est une question de cours!)
 \item Donner la d\'efinition de ``$a\in\R^2$ est un point adh\'erent de $D$.'' (Ceci est une question de cours!)

On consid\`ere dans la suite de l'exercice l'ensemble
$$D=\{(x,y)\in\R^2 \mid |x|\leq |y|, x^2+y^2<1\}\subset \R^2.$$
 \item Dessiner $D$.
 \item Montrer que $D$ n'est pas ouvert.
 \item  D\'eterminer $\overline D$, l'adh\'erence de $D$. On justifiera bri\`evement sa r\'eponse, en s'aidant d'un dessin.
\end{enumerate}
\finenonce{002647}



\finexercice


\exercice{5844, rouget, 2010/10/16}
\enonce{005844}{**}
\label{ex:rou6}
Soient $A$ et $B$ des parties d'un espace vectoriel normé $E$. Montrer que

\begin{enumerate}
 \item  $\overline{(\overline{A})}=\overline{A}$ et $\overset{\circ}{\overset{\circ}{A}}=\overset{\circ}{A}$.

\item   $A\subset B\Rightarrow \overline{A}\subset\overline{B}$ et$A\subset B\Rightarrow \overset{\circ}{A}\subset\overset{\circ}{B}$.

\item  $\overline{A\cup B}=\overline{A}\cup\overline{B}$  et $\overset{\circ}{A\cap B}=\overset{\circ}{A}\cap\overset{\circ}{B}$.

 
\item    $\overline{A\cap B}\subset\overline{A}\cap\overline{B}$  et $\overset{\circ}{A\cup B}\subset\overset{\circ}{A}\cup\overset{\circ}{B}$. Trouver un exemple où l'inclusion est stricte.
 

\item  $\overset{\circ}{A\setminus B}=\overset{\circ}{A}\setminus\overline{B}$.

\item  $\overline{\overset{\circ}{\overline{\overset{\circ}{A}}}}=\overline{\overset{\circ}{A}}$ et $\overset{\circ}{\overline{\overset{\circ}{\overline{A}}}}=\overset{\circ}{\overline{A}}$.
\end{enumerate}
\finenonce{005844}


\finexercice


\exercice{5845, rouget, 2010/10/16}
\enonce{005845}{**}
Trouver une partie $A$ de $\Rr$ telle que les sept ensembles $A$,  $\overline{A}$, $\overset{\circ}{A}$, $\overline{\overset{\circ}{A}}$, $\overset{\circ}{\overline{A}}$, $\overline{\overset{\circ}{\overline{A}}}$ et $\overset{\circ}{\overline{\overset{\circ}{A}}}$  soient deux à deux distincts.
\finenonce{005845}


\finexercice


\exercice{5846, rouget, 2010/10/16}
\enonce{005846}{**}
Soit $E$ le $\Rr$-espace vectoriel des fonctions continues sur $[0,1]$ à valeurs dans $\Rr$. On munit $E$ de $\|\;\|_\infty$.

$D$ est la partie de $E$ constituée des applications dérivables et $P$ est la partie de $E$ constituée des fonctions polynomiales. Déterminer l'intérieur de $D$ et l'intérieur de $P$.
\finenonce{005846}


\finexercice


\exercice{5848, rouget, 2010/10/16}
\enonce{005848}{**}
\begin{enumerate}
 \item  Soient $(E,N_E)$ et $(F,N_F)$ deux espaces vectoriels normés. Soient $f$ et $g$ deux applications continues sur $E$ à valeurs dans $F$. Soit $D$ une partie de $E$ dense dans $E$. Montrer que si $f_{/D}= g_{/D}$ alors $f = g$.

\item   Déterminer tous les morphismes continus de $(\Rr,+)$ dans lui-même.
\end{enumerate}
\finenonce{005848}


\finexercice


\exercice{5849, rouget, 2010/10/16}
\enonce{005849}{***}
Soit $u$ une suite bornée d'un espace vectoriel normé de dimension finie ayant une unique valeur d'adhérence. Montrer que la suite $u$ converge.
\finenonce{005849}


\finexercice


\exercice{7522, mourougane, 2021/08/10}
\enonce{007522}{Intérieur, adhérence}
Soit $A$ une partie de l'espace métrique $(\Cc,|~~|)$ et $z$ un nombre complexe.
\begin{enumerate}
 \item Montrer que $z$ appartient $\mathring{A}$ si et seulement si il existe un nombre réel $\epsilon>0$ 
 tel que $B(z,\epsilon)$ soit incluse dans $A$.
 \item Montrer que $z$ appartient à $\overline{A}$ si et seulement si pour tout réel $\epsilon>0$,
 la boule $B(z,\epsilon)$ rencontre~$A$.
\end{enumerate}
\finenonce{007522}
\finexercice

\section{229.02 Compacité}


\exercice{1744, maillot, 2001/09/01}

\enonce{001744}{partiel 1999}
Soit $f:\R^n\rightarrow \R$ une application continue. Montrer que les trois
conditions suivantes sont \'equivalentes~:

(1) $\forall M>0$, $\exists R>0$ tel que $\|x\|>R \Rightarrow \vert
f(x)\vert >M$.

(2) Pour toute partie born\'ee $B$ de $\R$, $f^{-1} (B)$ est une partie
born\'ee de $\R^n$.

(3) Pour toute partie compacte $K$ de $\R$, $f^{-1} (K)$ est une partie
compacte de $\R^n$.
\finenonce{001744}



\finexercice



\exercice{1763, maillot, 2001/09/01}

\enonce{001763}{}
Dans $\R^2$ euclidien, les ensembles suivants sont-ils compacts~?
\begin{itemize}
\item $A=\{(x,y)\in\R^2 \mid \frac12 \le \| (x,y)\| \le 2\ \text{et}
\ xy=1\}$.
\item $B=\{(x,y)\in\R^2 \mid \frac12 < \| (x,y)\| \le 2\ \text{et}
\ xy=1\}$.
\item $C= \{(x,\cos n)\in\R^2 \mid 0\le x \le 18 \ \text{et}\ n\in\N\}$.
\end{itemize}
\finenonce{001763}



\finexercice



\exercice{1764, maillot, 2001/09/01}

\enonce{001764}{}
Soit $E=\R^d$ muni d'une norme $\|\cdot\|$. On d\'efinit la
\emph{distance} d'un \'el\'ement $x_0$ de $E$ \`a une partie $A$ de $E$,
not\'ee $d(x_0,A)$, par la formule $$d(x_0,A)=\inf_{x\in A}\|x-x_0\|.$$
\begin{enumerate}
\item Supposons $A$ compact. Montrer que pour tout $x_0\in E$
il existe $y\in A$ tel que $d(x_0,A)=\|y-x_0\|$.
\item Montrer que le r\'esultat est encore vrai si on suppose seulement
que $A$ est ferm\'e. (On remarquera que pour toute partie $B$ de $A$ on
a $d(x_0,B)\ge d(x_0,A)$.)
\item Montrer que l'application qui \`a $x_0$ associe $d(x_0,A)$ est
continue sur $E$ (sans aucune hypoth\`ese sur $A$).
\item En d\'eduire que si $A$ est un ferm\'e de $E$ et $B$ un compact de
$E$ tels que $A$ et $B$ sont disjoints, alors il existe une
constante $\delta>0$ telle que
$$\|a-b\| \ge \delta \qquad \forall (a,b)\in A\times B.$$
\item Montrer par un contre-exemple que le r\'esultat est faux si on
suppose seulement que $A$ et $B$ sont deux ferm\'es disjoints.
\end{enumerate}
\finenonce{001764}



\finexercice



\exercice{1767, maillot, 2001/09/01}

\enonce{001767}{}
Soit $(E,\|\cdot\|)$ un espace vectoriel norm\'e. Soit $(x_n)$ une suite
convergente de $E$ et $x$ sa limite. Montrer que l'ensemble $\{x\}\cup
\{x_n,n\in\N\}$ est compact.
\finenonce{001767}



\finexercice



\exercice{1783, gineste, 2001/11/01}

\enonce{001783}{}
Soit $(u_n)_{n \geq 1}$ une suite r\'eelle. $\forall n \geq 1$, on pose $A_n=\{ u_p / \; p\geq n \}.$ D\'emontrer que l'ensemble des valeurs d'adh\'erence de la suite  $(u_n)_{n \geq 1}$ est $V= \bigcap_{n \geq 1} \overline{A_n}$, et qu'ainsi $V$ est ferm\'e. En d\'eduire que si la suite est born\'ee, alors l'ensemble $V$ est un compact non vide.
\finenonce{001783}



\finexercice



\exercice{4827, quercia, 2010/03/16}
\enonce{004827}{Graphe ferm{\'e}}
Soient $E,F$ deux espaces vectoriels norm{\'e}s et $f : E \to  {F.}$
On note $Gr(f) = \{(x,y) \in E\times F\text{ tq } y = f(x)\}$.

\begin{enumerate}
  \item Montrer que si $f$ est continue, alors $Gr(f)$ est ferm{\'e} dans $E\times F$.
  \item Prouver la r{\'e}ciproque lorsque $f(E)$ est inclus dans un compact de~$F$.
  \item Donner un contre-exemple si $f(E)$ n'est pas inclus dans un compact.

\end{enumerate}
\finenonce{004827}



\finexercice


\exercice{4828, quercia, 2010/03/16}
\enonce{004828}{Application presque contractante}
Soit $A$ une partie compacte d'un evn $E$ et $f : A \to  A$ telle que :
$\forall\ x,y \in A,\ x\ne y  \Rightarrow  d(f(x),f(y)) < d(x,y)$.

\begin{enumerate}
  \item Montrer que $f$ admet un point fixe unique, $a$.
  \item Soit $(x_n)$ une suite d'{\'e}l{\'e}ments de $A$ telle que $x_{n+1} = f(x_n)$.
    Montrer qu'elle converge vers $a$.
    

\end{enumerate}
\finenonce{004828}



\finexercice


\exercice{4829, quercia, 2010/03/16}
\enonce{004829}{Application presque contractante (Mines MP 2003)}
Soit $C$ un compact convexe d'un evn $E$. Soit~$f : C \to {C,}$
$1$-lipschitzienne.  Montrer que~$f$ admet un point fixe. On pourra
utiliser la fonction $f_n$~:~$x \mapsto\frac an +
\Bigl(1-\frac1n\Bigr)f(x)$ avec $a\in C$.

\finenonce{004829}



\finexercice


\exercice{4830, quercia, 2010/03/16}
\enonce{004830}{Fonction bicontinue sur un compact}
Soit $A$ une partie compacte d'un evn $E$ et $f : A \to F$ une fonction
continue et injective ($F$ = evn).

\begin{enumerate}
  \item Montrer que ${f^{-1}} : {f(A)} \to A$ est aussi continue.
  \item Donner un exemple o{\`u} $A$ n'est pas compact et $f^{-1}$ n'est pas continue.
\end{enumerate}
\finenonce{004830}



\finexercice


\exercice{4831, quercia, 2010/03/16}
\enonce{004831}{Isom{\'e}tries d'un compact}
Soit $A$ une partie compacte d'un evn $E$ et $f : A \to A$ telle que :
$\forall\ x,y \in A,\ d(f(x),f(y)) \ge d(x,y)$.

\begin{enumerate}
  \item Soit $a \in A$ et $(a_n)$ la suite d{\'e}finie par :
    $a_0 = a$, $a_{n+1} = f(a_n)$. Montrer que $a$ est valeur d'adh{\'e}rence de la
    suite $(a_n)$.
  \item Soient $a,b \in A$. Montrer que $d(f(a),f(b)) = d(a,b)$.
  \item Montrer que $f(A) = A$.

\end{enumerate}
\finenonce{004831}



\finexercice


\exercice{4832, quercia, 2010/03/16}
\enonce{004832}{Partie dense dans un compact}
Soit $A$ une partie compacte d'un evn $E$.
Montrer qu'il existe une suite $(a_k)$ d'{\'e}l{\'e}ments de~$A$
qui est dense dans $A$.
\finenonce{004832}



\finexercice


\exercice{4833, quercia, 2010/03/16}
\enonce{004833}{Intersection emboit{\'e}e}
Soit $E$ un espace vectoriel norm{\'e},
$(K_n)$ une suite d{\'e}croissante de compacts non vides de $E$ et $K = \bigcap_n K_n$.

\begin{enumerate}
  \item Montrer que $K \ne \varnothing$.
  \item Soit $U$ un ouvert contenant $K$. Montrer qu'il existe $n$ tel que $K_n \subset U$.
  \item Montrer que $\delta(K) = \lim_{n\to\infty} \delta(K_n)$ ($\delta$ est le diam{\`e}tre).



\end{enumerate}
\finenonce{004833}



\finexercice


\exercice{4834, quercia, 2010/03/16}
\enonce{004834}{Image d'une intersection}
Soient $E,F$ deux espaces vectoriels norm{\'e}s et $f : E \to F$ continue.
Soit $(K_n)$ une suite d{\'e}croissante de compacts de $E$.
Montrer que $f(\mathop{\cap}\limits_nK_n) = \bigcap_nf(K_n)$.

\finenonce{004834}



\finexercice


\exercice{4835, quercia, 2010/03/16}
\enonce{004835}{Recouvrement ouvert}
Soit $A$ une partie compacte d'un evn $E$
et $(O_i)_{i\in I}$ un recouvrement ouvert
de $A$. Montrer qu'il existe $r > 0$ tel que toute partie de $A$ de diam{\`e}tre inf{\'e}rieur
ou {\'e}gal {\`a} $r$ soit incluse dans l'un des $O_i$.


\finenonce{004835}



\finexercice


\exercice{4836, quercia, 2010/03/16}
\enonce{004836}{Ensemble compact de suites}
Soit $E = {\cal B}(\N,\R) = \{ \text{suites } u = (u_n) \text{ born{\'e}es}\}$.
On munit $E$ de la norme :
$\|u\| = \sum_{n=0}^\infty \frac{|u_n|}{2^n}$.
Montrer que $A = \{u\in E \text{ tq } \forall\ n\in\N,\ 0\le u_n\le 1\}$
est compact.


\finenonce{004836}



\finexercice


\exercice{4837, quercia, 2010/03/16}
\enonce{004837}{Boule unit{\'e} non compacte}
Soit $E = {\cal C}([0,2\pi])$ muni de la norme $\|.\|_2$.
Pour $n \in \N$, on pose $f_n(x) = \cos(nx)$.

\begin{enumerate}
  \item Calculer $\|f_n-f_p\|_2$ pour $n,p \in \N$.
  \item En d{\'e}duire que $\overline B(0,1)$ n'est pas compacte.

\end{enumerate}
\finenonce{004837}



\finexercice


\exercice{4838, quercia, 2010/03/16}
\enonce{004838}{Plus petite boule contenant une partie}
Soit $\|.\|$ une norme sur $\R^2$, $A \subset \R^2$ une partie born{\'e}e contenant au moins deux points.

\begin{enumerate}
  \item Montrer qu'il existe une boule ferm{\'e}e de rayon minimum contenant $A$.
  \item Montrer que cette boule n'est pas n{\'e}cessairement unique (on prendra $\|.\| = \|.\|_\infty$).
  \item Montrer que si $\|.\|$ est une norme euclidienne, alors la boule pr{\'e}c{\'e}dente est unique.

\end{enumerate}
\finenonce{004838}



\finexercice


\exercice{4839, quercia, 2010/03/16}
\enonce{004839}{Polytechnique MP$^*$ 2000}
Soit~$E$ un espace vectoriel norm{\'e}, $K$ un compact convexe de~$E$,
$f$ une application de~$K$ dans~$K$, $1$-lipchitzienne.
Montrer que $f$ a un point fixe.

\finenonce{004839}



\finexercice


\exercice{4840, quercia, 2010/03/16}
\enonce{004840}{Thm. de Riesz, Stival 2003}

Soit $E$ un evn de dimension infinie.

\begin{enumerate}
  \item Soit $F$ un sev de dimension finie et $a\in E\setminus F$.
  \begin{enumerate}
    \item Montrer qu'il existe $b\in F$ tel que $\|a-b\| = d(a,F)$.
    \item En d{\'e}duire qu'il existe $c\in E$ tel que $\|c\| = 1 = d(c,F)$.
  \end{enumerate}
  \item Montrer que la boule unit{\'e} de~$E$ n'est pas compacte.
    
\end{enumerate}
\finenonce{004840}



\finexercice

\section{229.03 Borne supérieure}


\exercice{1750, maillot, 2001/09/01}

\enonce{001750}{}
Soient $A$ et $B$ deux parties non vides et majorées de $\R$. Montrer
les implications suivants :
\begin{itemize}
\item
$\exists M \in\R\;\forall x\in A,\; x<M\Rightarrow\sup A\leq M$
\item
$A\subset B\Rightarrow\sup A\leq\sup B$.
\end{itemize}

\finenonce{001750}



\finexercice



\exercice{1751, maillot, 2001/09/01}

\enonce{001751}{}
Soient $A$ et $B$ deux parties non vides et majorées de $\R$. On définit
:
$$
A+B=\{c\in\R\ | \ \exists a\in A,\exists b\in B, c=a+B\}.
$$
\begin{enumerate}
\item Montrer que $A+B$ admet une borne supérieure, puis que
$\sup(A+B)=\sup A+\sup B$.

\item Montrer l'implication :
$$
\exists M\in\R\;\forall x\in A,\forall y\in B,\;x+y<M\Rightarrow\sup A+\sup B\leq M.
$$
\end{enumerate}
\finenonce{001751}



\finexercice



\exercice{1752, maillot, 2001/09/01}

\enonce{001752}{}
Soit $\epsilon\in\R^+$ tel que $\forall
x\in\R_+^*,\;x\geq\epsilon$. Montrer que $\epsilon=0$.
\finenonce{001752}



\finexercice



\exercice{1753, maillot, 2001/09/01}

\enonce{001753}{}
Soit $A$ une partie non vide et bornée de $\R$. Montrer que :
$$
\sup\{|x-y|:(x,y)\in A^2\}=\sup A-\inf A.
$$
\finenonce{001753}



\finexercice



\exercice{5850, rouget, 2010/10/16}
\enonce{005850}{***}
Calculer $\underset{\alpha\in]0,\pi[}{\text{Inf}}\left\{\underset{n\in\Zz}{\text{Sup}}|\sin(n\alpha)|\right\}$.
\finenonce{005850}


\finexercice

\section{229.04 Topologie de la droite réelle}


\exercice{4717, quercia, 2010/03/16}
\enonce{004717}{Partie {\`a} un seul point d'accumulation}
Soit $A$ une partie born{\'e}e de $\R$ ayant un seul point d'accumulation, $a$.

\begin{enumerate}
  \item Montrer que $A$ est d{\'e}nombrable.
    
  \item On num{\'e}rote les {\'e}l{\'e}ments de $A$ d'une mani{\`e}re quelconque :
    $A = \{ x_1, x_2, \dots, x_n, \dots \}$.

    Montrer que $x_n \xrightarrow[n\to\infty]{} a$.

\end{enumerate}
\finenonce{004717}



\finexercice


\exercice{4718, quercia, 2010/03/16}
\enonce{004718}{$(\sin(n))$ est dense}
Soit $a \in \R\setminus \Q$ et $A = \{ ma + n \text{ tq } m\in \Z,\ n\in \N\}$.
Montrer que $A$ est dense dans $\R$.

Application : Montrer que tout r{\'e}el de $[-1,1]$ est valeur d'adh{\'e}rence de la suite
$(\sin n)$.
\finenonce{004718}



\finexercice


\exercice{4719, quercia, 2010/03/16}
\enonce{004719}{$\sqrt m - \sqrt n$}
Montrer que l'ensemble $A = \{ \sqrt m - \sqrt n \text{ tq } m,n \in \N \}$ est dense dans
$\R$.




\finenonce{004719}



\finexercice


\exercice{4720, quercia, 2010/03/16}
\enonce{004720}{Unit{\'e}s quadratiques}
Soit $A =\{ n+p\sqrt2 \text{ tq } n,p \in \N,\ n+p\sqrt2 > 0,\ n^2-2p^2 = 1\}$.
Montrer que $A$ est un sous-groupe discret de $\R^{+*}$.


\finenonce{004720}



\finexercice


\exercice{4721, quercia, 2010/03/16}
\enonce{004721}{Olympiades 1991}
Soit $a > 1$. Montrer qu'il existe une suite r{\'e}elle born{\'e}e, $(x_n)$,
telle que : $\forall\ i\ne j,\ |x_i-x_j| \ge \frac 1{|i-j|^a}$.


\finenonce{004721}



\finexercice


\exercice{4722, quercia, 2010/03/16}
\enonce{004722}{$u_{n+1} - u_n \to 0$}
Soit $(u_n)$ une suite r{\'e}elle born{\'e}e telle que $u_{n+1} - u_n \xrightarrow[n\to\infty]{} 0$.
Montrer que l'ensemble des valeurs d'adh{\'e}rence de $(u_n)$ est un intervalle.
\finenonce{004722}



\finexercice


\exercice{4723, quercia, 2010/03/16}
\enonce{004723}{$u_{n+1} - u_n \to 0$}
Soit $f : {[0,1]} \to {[0,1]}$ continue, $u_0 \in {[0,1]}$ et $(u_n)$
la suite des it{\'e}r{\'e}es de $f$ en $u_0$.

On suppose que $u_{n+1} - u_n \xrightarrow[n\to\infty]{} 0$.
Montrer que la suite $(u_n)$ converge vers un point fixe de $f$.


\finenonce{004723}



\finexercice


\exercice{4724, quercia, 2010/03/16}
\enonce{004724}{$\exp(iu_n)$}
Soit $(u_n)$ une suite r{\'e}elle telle que la suite $(\exp(iu_n))$ converge
et la suite $(|u_{n+1} - u_n|)$ est major{\'e}e par $\alpha < \pi$.
Montrer que $(u_n)$ converge.

\finenonce{004724}



\finexercice


\exercice{4725, quercia, 2010/03/16}
\enonce{004725}{$\exp(iu_n)$}
Soit $(u_n)$ une suite r{\'e}elle telle que $u_{n+1} - u_n \xrightarrow[n\to\infty]{} 0$
et $u_n \xrightarrow[n\to\infty]{}+\infty$.
D{\'e}montrer que la suite $(\exp(iu_n))$ est dense dans $\mathbb{U}$.
\finenonce{004725}



\finexercice


\exercice{4726, quercia, 2010/03/16}
\enonce{004726}{$\exp(iu_n)$}
Soit $(x_n)$ une suite r{\'e}elle born{\'e}e et $u>0,\,v>0$. On suppose que
$\frac{u}{v} \notin \Q$ et que les suites $(e^{iux_n})$ et $(e^{ivx_n})$
convergent.
Montrer que la suite $(x_n)$ converge.
\finenonce{004726}



\finexercice


\exercice{4727, quercia, 2010/03/16}
\enonce{004727}{$u_{n+p} \le u_n + u_p$}
Soit $(u_n)$ une suite r{\'e}elle positive telle que :
$\forall\ n,p\in\N,\ u_{n+p} \le u_n + u_p$.
Montrer que la suite $\left(\frac {u_n}n\right)$ est convergente.


\finenonce{004727}



\finexercice


\exercice{4728, quercia, 2010/03/16}
\enonce{004728}{Fonctions p{\'e}riodiques (Ens Ulm-Lyon-Cachan MP$^*$ 2003)}
\begin{enumerate}
  \item D{\'e}terminer toutes les fonctions~$f : \R \to \R$ continues, p{\'e}riodiques de
    p{\'e}riodes~$1$ et~$\sqrt2$.
    
  \item D{\'e}terminer les fonctions~$f : {\R^2} \to {\R^2}$ continues telles que~:\par
    pour tout~$X\in\R^2$,
    $f(X) = f(X+(1,0)) = f(X+(0,1)) = f(AX)$ o{\`u} $A=\left(\begin{smallmatrix}\ 1&1\cr0&1\cr\end{smallmatrix}\right)$.
    
 \end{enumerate}
\finenonce{004728}



\finexercice   


\exercice{5843, rouget, 2010/10/16}
\enonce{005843}{**}
Montrer qu'entre deux réels distincts, il existe un rationnel (ou encore montrer que $\Qq$ est dense dans $\Rr$).
\finenonce{005843}


\finexercice


\exercice{5851, rouget, 2010/10/16}
\enonce{005851}{*** I}
Soit $f~:~\Rr\rightarrow\Rr$ une application uniformément continue sur $\Rr$. Montrer qu'il existe deux réels $a$ et $b$ tels que $\forall x\in\Rr$, $|f(x)|\leqslant a|x| +b$.
\finenonce{005851}


\finexercice

\section{229.05 Topologie des espaces métriques}


\exercice{4729, quercia, 2010/03/16}
\enonce{004729}{Ouverts disjoints}
Soient $U,V$ deux ouverts disjoints d'un espace vectoriel norm{\'e}.
Montrer que $\mathring{\overline U}$ et $\mathring{\overline V}$ sont disjoints.


Donner un contre-exemple lorsque $U$ et $V$ ne sont pas ouverts.

\finenonce{004729}



\finexercice


\exercice{4730, quercia, 2010/03/16}
\enonce{004730}{$A$ ouvert disjoint de $B$}
Soient $A,B$ deux parties d'un espace vectoriel norm{\'e} disjointes. Si $A$ est ouvert, montrer que
$A$ et $\overline B$ sont disjoints.

\finenonce{004730}



\finexercice


\exercice{4731, quercia, 2010/03/16}
\enonce{004731}{$\overline{\mathring{\overline U}} = \overline U$.}
Soit $U$ un ouvert d'un espace vectoriel norm{\'e}. Montrer que $\overline{\mathring{\overline U}} = \overline U$.

\finenonce{004731}



\finexercice


\exercice{4732, quercia, 2010/03/16}
\enonce{004732}{Fronti{\`e}re d'un ouvert}
Soit $U$ un ouvert d'un espace vectoriel norm{\'e}. Montrer que la fronti{\`e}re de $U$ est d'int{\'e}rieur vide.

\finenonce{004732}



\finexercice


\exercice{4733, quercia, 2010/03/16}
\enonce{004733}{La distance est 1-lipchitzienne}
Soit $A$ une partie non vide d'un espace vectoriel norm{\'e} $E$.
Pour $x \in E$, on pose $d(x,A) = \inf\{d(x,a)\text{ tq } a\in A\}$.

\begin{enumerate}
  \item Montrer que : $\forall\ x,y \in E,\ |d(x,A) - d(y,A)| \le d(x,y)$.
  \item Montrer que l'application $x  \mapsto d(x,A)$ est continue.

\end{enumerate}
\finenonce{004733}



\finexercice


\exercice{4734, quercia, 2010/03/16}
\enonce{004734}{Diam{\`e}tre de la fronti{\`e}re}
Soit $A$ une partie non vide et born{\'e}e d'un evn $E$.
On note $\delta(A) = \sup\{d(x,y) \text{ tq } x,y\in A\}$
({\it diam{\`e}tre de $A$}).

Montrer que $\delta(A) = \delta(\text{Fr}(A))$.

\finenonce{004734}



\finexercice


\exercice{4735, quercia, 2010/03/16}
\enonce{004735}{Ensemble d{\'e}riv{\'e}}
Soit $A$ une partie d'un espace vectoriel norm{\'e} $E$. Un point $x \in E$ est dit
{\it point d'accumulation de $A$} si toute boule de centre $x$
contient une infinit{\'e} de points de $A$.
On note $A'$ l'ensemble des points d'accumulation de $A$
({\it ensemble d{\'e}riv{\'e} de $A$}).
Montrer que $A'$ est ferm{\'e}, et comparer $A'$ et $\overline A$.

\finenonce{004735}



\finexercice


\exercice{4736, quercia, 2010/03/16}
\enonce{004736}{Caract{\'e}risation des fonctions continues}
Soient $E$, $F$ deux espaces vectoriels norm{\'e}s et $f : E \to {F.}$

Montrer que $f$ est continue 
si et seulement si : $\forall\ A \subset E,\ f(\overline A) \subset \overline{f(A)}$\par
si et seulement si : $\forall\ B \subset F,\ f^{-1}(\mathring B) \subset f^{-1}(B)^\circ$.\par

\finenonce{004736}



\finexercice

\section{229.06 Topologie des espaces vectoriels normés}


\exercice{4737, quercia, 2010/03/16}
\enonce{004737}{Unicit{\'e} du centre et du rayon d'une boule}
Soit $E$ un evn non nul et $\vec a,\vec a' \in E$, $r,r' > 0$ tels que
$B_f(\vec a,r) = B_f(\vec a',r')$.
Montrer que $\vec a = \vec a'$ et $r = r'$.
\finenonce{004737}



\finexercice


\exercice{4738, quercia, 2010/03/16}
\enonce{004738}{$x+y+z = 0$}
\def \xxxx{\vec x}\def\yyyy{\vec y}\def\zzzz{\vec z}%
Soient $\xxxx,\yyyy,\zzzz$ trois vecteurs d'un evn $E$ tels que $\xxxx+\yyyy+\zzzz = \vec 0$.

Montrer que : $\|\xxxx-\yyyy\,\| + \|\yyyy-\zzzz\,\| + \|\zzzz-\xxxx\,\| \ge
       \frac32(\|\xxxx\,\|+\|\yyyy\,\|+\|\zzzz\,\|)$.
\finenonce{004738}



\finexercice

\exercice{4739, quercia, 2010/03/16}
\enonce{004739}{Une boule est convexe}
Soit $E$ un evn, et $\vec a \in E$, $r > 0$.
On note $\overline B = \overline B(\vec a,r)$ et $\mathring B = \mathring B(\vec a,r)$.

\begin{enumerate}
  \item Montrer que $\overline B$ et $\mathring B$ sont convexes.
  \item Si la norme est euclidienne, montrer que si $\vec u,\vec v \in \overline B$ avec
    $\vec u \ne \vec v$, alors $]\vec u,\vec v\,[ \subset \mathring B$.

    $(]\vec u,\vec v\,[ = \{ (1-t)\vec u + t\vec v \text{ tq } t \in {]0,1[}\,\})$
  \item En d{\'e}duire que si la norme est euclidienne, toute partie $A$ telle que
    $\mathring B\subset A \subset \overline B$ est convexe.
  \item Donner un contre-exemple avec une norme non euclidienne.
\end{enumerate}
\finenonce{004739}



\finexercice


\exercice{4740, quercia, 2010/03/16}
\enonce{004740}{Distance {\`a} un ensemble}
Soit $E$ un evn et $A \subset E$ une partie non vide.
Pour $\vec x \in E$ on pose :
$d(\vec x,A) = \inf\{\|\vec x-\vec a\,\| \text{ tq } \vec a \in A \}$.

\begin{enumerate}
  \item Montrer que : $\forall\ \vec x,\vec y \in E,\
    \Bigl|d(\vec x,A) - d(\vec y,A)\Bigr| \le \|\vec x-\vec y\,\|$.
    (Enlever la valeur absolue et d{\'e}montrer s{\'e}par{\'e}ment chaque in{\'e}galit{\'e})
  \item Montrer que l'application $\vec x  \mapsto d(\vec x,A)$ est continue.
\end{enumerate}
\finenonce{004740}



\finexercice


\exercice{4741, quercia, 2010/03/16}
\enonce{004741}{Distance {\`a} un ensemble (ENS Cachan MP 2002)}

Soit $A$ une partie de~$\R^n$ non vide. On note pour $x\in\R^n$~:
$d_A(x) = \inf\{\|x-y\|\text{ tq }y\in A\}$.

\begin{enumerate}
  \item Montrer que $d_A$ est continue.

  \item Soient deux parties de $\R^n$ non vides $A,B$. Donner une condition {\'e}quivalente
    {\`a} $d_A=d_B$.

  \item On note $\rho(A,B) = \sup\{|d_A(y)-d_B(y)|,\ y\in\R^n\}$, valant {\'e}ventuellement~$+\infty$.

    Montrer que l'on a $\rho(A,B) = \max\Bigl(\sup\limits_{x\in A} d_B(x),\ \sup\limits_{x\in B} d_A(x)\Bigr)$.

\end{enumerate}
\finenonce{004741}



\finexercice


\exercice{4742, quercia, 2010/03/16}
\enonce{004742}{Distance entre un ferm{\'e} et un compact}
Soient $A,B$ deux parties compactes non vides de $\R^n$.

Montrer qu'il existe $a\in A$ et $b \in B$ tels que
$\|a-b\| = \min\{\|x-y\| \text{ tq } x\in A,\ y\in B\}$.

Montrer que ceci est encore vrai si on suppose $A$ compact et $B$ ferm{\'e}.
\finenonce{004742}



\finexercice


\exercice{4743, quercia, 2010/03/16}
\enonce{004743}{Diam{\`e}tre}
Soit $E$ un evn de dimension finie et $A \subset E$ born{\'e}, ferm{\'e}, non vide.
Montrer qu'il existe $\vec a,\vec b \in A$ tels que
$\|\vec a-\vec b\| = \max( \|\vec x-\vec y\,\| \text{ tq } \vec x,\vec y \in A )$.
(Consid{\'e}rer l'ensemble $A\times A$ dans l'evn $E\times E$)
\finenonce{004743}



\finexercice


\exercice{4744, quercia, 2010/03/16}
\enonce{004744}{Diam{\`e}tres concourants (Ens Ulm MP$^*$ 2003)}
\begin{enumerate}
  \item Soit~$K$ un compact convexe de~$\R^2$ d'int{\'e}rieur non vide.
    Soit~$O\in\mathring K$. Montrer qu'il existe une fonction
    $f : \R \to {\R^+}$ continue $2\pi$-p{\'e}riodique telle qu'en coordonn{\'e}es polaires
    de centre~$O$, $K$ est d{\'e}fini par~$\rho\le f(\theta)$.
    

  \item Soit~$g : {[0,1]} \to \R$ continue telle que
    $ \int_{x=0}^{\pi}g(x)\cos(x)\,d x =  \int_{x=0}^{\pi}g(x)\sin(x)\,d x = 0$.
    Montrer que~$g$ s'annule au moins deux fois sur~$]0,\pi[$.
    

  \item Soit~$G$ le centre de gravit{\'e} de~$K$. Montrer que~$G$ est le milieu
    d'au moins trois "diam{\`e}tres" de~$K$ (trois segments joignant deux points de
    la fronti{\`e}re).
    
\end{enumerate}
\finenonce{004744}




\finexercice


\exercice{4745, quercia, 2010/03/16}
\enonce{004745}{$x/\max(1,\|x\|)$, Centrale MP 2005}

Soit $f$ d{\'e}finie par $f(x)=\frac{x}{\max(1,\|x\|)}$. Montrer que $f$ est $2$-lipschitzienne.

\finenonce{004745}



\finexercice


\exercice{4746, quercia, 2010/03/16}
\enonce{004746}{$u_n$ colin {\`a} $v_n  \Rightarrow  \lim u_n$ colin {\`a} $\lim v_n$}
Soit $E$ un evn de dimension finie et $({\vec u}_n)$, $({\vec v}_n)$
deux suites de vecteurs telles que :
$$\forall\ n \in \N,\ {\vec u}_n \text{ est colin{\'e}aire {\`a} } {\vec v}_n,
  \qquad
  {\vec u}_n \xrightarrow[n\to\infty]{} \vec u,  \qquad
  {\vec v}_n \xrightarrow[n\to\infty]{} \vec v.$$
Montrer que $\vec u$ et $\vec v$ sont colin{\'e}aires
(raisonner par l'absurde et compl{\'e}ter $(\vec u,\vec v)$ en une base de $E$).
\finenonce{004746}



\finexercice


\exercice{4747, quercia, 2010/03/16}
\enonce{004747}{Suites de Cauchy}
Soient $(u_n)$, $(v_n)$ deux suites d'un evn $E$ telles que
$u_n - v_n \xrightarrow[n\to\infty]{} 0$ et $(u_n)$ est de Cauchy.
Montrer que $(v_n)$ est de Cauchy.
\finenonce{004747}



\finexercice


\exercice{4748, quercia, 2010/03/16}
\enonce{004748}{Suite de Cauchy non convergente}
Soit $E = \R[X]$ muni de la norme : $\|\sum a_kX^k \| = \max(|a_k|,\ k\in \N)$.
On note $P_n = 1 + X + \frac{X^2}2 + \dots + \frac{X^n}n$.
Montrer que la suite $(P_n)$ est de Cauchy, mais ne converge pas.
\finenonce{004748}



\finexercice


\exercice{4749, quercia, 2010/03/16}
\enonce{004749}{Mines PC 1998}
Soit $B$ une matrice antisym{\'e}trique.
On suppose que la suite $(B^n)$ converge vers une matrice~$C$.
Que peut-on dire de $C$~?
\finenonce{004749}



\finexercice


\exercice{4750, quercia, 2010/03/16}
\enonce{004750}{Suite de matrices inversibles}
Soit $(A_n)$ une suite de matrices de $\mathcal{M}_p(\R)$ v{\'e}rifiant les propri{\'e}t{\'e}s
suivantes :
$$\begin{cases} 1:\ A_n\xrightarrow[n\to\infty]{} A \in \mathcal{M}_p(\R) \cr
          2:\ \text{pour tout $n$, $A_n$ est inversible} \cr
          3:\ A_n^{-1}\xrightarrow[n\to\infty]{} B \in \mathcal{M}_p(\R). \cr \end{cases}$$

\begin{enumerate}
  \item Montrer que $A$ est inversible et $A^{-1} = B$.
  \item Peut-on retirer la propri{\'e}t{\'e} 3 ?
\end{enumerate}
\finenonce{004750}



\finexercice


\exercice{4751, quercia, 2010/03/16}
\enonce{004751}{Suite de matrices inversibles}
Soit $A \in \mathcal{M}_p(\R)$ quelconque. Montrer qu'il existe une suite de matrices inversibles
convergeant vers $A$.
\finenonce{004751}



\finexercice


\exercice{4752, quercia, 2010/03/16}
\enonce{004752}{DSE de $I-A$}
Soit $A \in \mathcal{M}_p(\R)$. On suppose que la suite de matrices :
$A_n = I + A + A^2 + \dots + A^n$ converge vers une matrice $B$.
Montrer que $I-A$ est inversible, et $B = (I-A)^{-1}$.

Remarque : La r{\'e}ciproque est fausse, c'est {\`a} dire que la suite $(A_n)$
peut diverger m{\^e}me si $I-A$ est inversible. Chercher un contre-exemple.
\finenonce{004752}



\finexercice


\exercice{4753, quercia, 2010/03/16}
\enonce{004753}{Ensam PSI 1998}
Soit $A\in \mathcal{M}_n(\C)$ telle que la suite $(A^k)$ converge vers une matrice~$P$.
Montrer que $P$ est une matrice de projection.
\finenonce{004753}



\finexercice


\exercice{4754, quercia, 2010/03/16}
\enonce{004754}{Suites de fonctions}
Soient $E = \mathcal{C}([a,b]\to\R)$, $(f_n)$ une suite de fonctions de $E$
et $f \in E$.
Comparer les {\'e}nonc{\'e}s :
$$1 : \|f_n-f\|_1 \xrightarrow[n\to\infty]{} 0
  \qquad
  2 : \|f_n-f\|_2 \xrightarrow[n\to\infty]{}0  \qquad
  3 : \|f_n-f\|_\infty \xrightarrow[n\to\infty]{} 0.$$
\finenonce{004754}



\finexercice


\exercice{4755, quercia, 2010/03/16}
\enonce{004755}{} 
On note $E$ l'espace vectoriel des suites r{\'e}elles $(x_n)$ telles que la s{\'e}rie $\sum x_n^2$
converge. On le munit du produit scalaire $(x\mid y)=\sum_{n=0}^\infty x_ny_n$.
Soit $(y^s)$ une suite born{\'e}e d'{\'e}l{\'e}ments de $E$. Montrer qu'on peut en
extraire une sous-suite convergent faiblement, c'est-{\`a}-dire qu'il existe $z$
telle que pour tout $x$ de $E$ on ait $(x\mid y^{s_k})\xrightarrow[k\to\infty]{}(x\mid z)$.

\finenonce{004755}



\finexercice


\exercice{4756, quercia, 2010/03/16}
\enonce{004756}{ENS Lyon MP 2002}
Soit $E$ un espace vectoriel norm{\'e} sur $\R$ ou $\C$ de dimension finie,
et $u\in\mathcal{L}(E)$ tel que pour tout $x\in E$ la suite $(u^n(x))_{n\in\N}$ est
born{\'e}e.

\begin{enumerate}
  \item Montrer que la suite $(\|\hskip-1pt|u^n\|\hskip-1pt|)_{n\in\N}$ est born{\'e}e.

  \item D{\'e}terminer la limite quand $n\to\infty$ de $\frac1{n+1}\sum_{i=0}^n u^i(x)$.

\end{enumerate}
\finenonce{004756}



\finexercice


\exercice{4757, quercia, 2010/03/16}
\enonce{004757}{Norme bizarre}
Montrer que $(x,y)\mapsto \sup\limits_{t\in \R} \frac{|x+ty|}{1+t+t^2}$ est
une norme sur $\R ^2$ ; dessiner la boule unit{\'e}.
\finenonce{004757}



\finexercice


\exercice{4758, quercia, 2010/03/16}
\enonce{004758}{Normes de polyn{\^o}mes}
Soit $E = \R[X]$. Pour $P = \sum_{k=0}^n a_kX^k$, on pose :
$$
\begin{aligned}
  \|P\|_1 &= \sum_{k=0}^n|a_k|, \\
  \|P\|_\infty &=
  \max\{|a_0|,\dots,|a_n|\}, \\ \|P\|_* &= \max\{|P(t)| \text{ tq }
  0 \le t \le 1\}.
\end{aligned}
  $$

Montrer que ce sont des normes, et qu'elles sont deux {\`a} deux non {\'e}quivalentes.
(On consid{\`e}rera $P_n(t) = (t-1)^n$ et $Q_n(t) = 1 + t + t^2 + \dots + t^n$)
\finenonce{004758}



\finexercice


\exercice{4759, quercia, 2010/03/16}
\enonce{004759}{Norme de polyn{\^o}mes}
Soit $E = \R[X]$. Pour $P\in E$ on pose
$\|P\| = \sup(|P(t)-P'(t)| \text{ tq } t \in {[0,1]})$.

Montrer qu'on d{\'e}finit ainsi une norme sur $E$.
\finenonce{004759}



\finexercice


\exercice{4760, quercia, 2010/03/16}
\enonce{004760}{Normes de polyn{\^o}mes}

Soit $a\in\R$. On pose pour $P\in\R[X]$~: $N_a(P) = |P(a)| +  \int_{t=0}^1 |P'(t)|\,d t$.
Montrer que\dots

\begin{enumerate}
  \item $N_a$ est une norme.
  \item $N_0$ et $N_1$ sont {\'e}quivalentes.
  \item Si $a,b\in{[0,1]}$, alors $N_a$ et $N_b$ sont {\'e}quivalentes.
  \item Soit $P_n = (X/2)^n$. D{\'e}terminer pour quelles normes $N_a$ la suite
    $(P_n)$ est convergente et quelle est sa limite.
    

  \item Si $0 \le a < b$ et $b > 1$
    alors aucune des normes $N_a$, $N_b$ n'est plus fine que l'autre.
     
\end{enumerate}
\finenonce{004760}



\finexercice


\exercice{4761, quercia, 2010/03/16}
\enonce{004761}{Normes sur les polyn{\^o}mes}
Soit $(\lambda_n)$ une suite de r{\'e}els strictement positifs.
On lui associe la norme sur $\R[x]$ :
$N(\sum\limits_i a_ix^i) = \sum\limits_i \lambda_i|a_i|$.

Soient $(\lambda_n)$ et $(\lambda'_n)$ deux suites et $N$, $N'$ les normes
associ{\'e}es.
Montrer que $N$ et $N'$ sont {\'e}quivalentes si et seulement si les suites
$\Bigl(\frac{\lambda_n}{\lambda'_n}\Bigr)$ et
$\Bigl(\frac{\lambda'_n}{\lambda_n}\Bigr)$ sont born{\'e}es.
\finenonce{004761}



\finexercice


\exercice{4762, quercia, 2010/03/16}
\enonce{004762}{Centrale MP 2006}
$E$ est l'ensemble des fonctions $f$ de classe $\mathcal{C}^2$ sur $[0,1]$ telles que $f(0) = f'(0) =0$.
Pour $f \in E$, on pose :

$$
\begin{aligned}
  N_{\infty}(f) &= \sup_{x\in[0,1]}{|f(x)|}, \\ N(f) &=
  \sup_{x\in[0,1]}{|f(x) + f''(x)|}, \\ N_{1}(f) &=
  \sup_{x\in[0,1]}{|f''(x)|} + \sup_{x\in[0,1]}{|f(x)|}.
\end{aligned}
 $$
\begin{enumerate}
  \item Montrer que $N_{\infty}$, $N$ et $N_{1}$ sont des normes sur $E$.

  \item Montrer que $N_{\infty}$ n'est {\'e}quivalente ni {\`a} $N_{1}$ ni {\`a} $N$.

  \item Montrer que $N$ et $N_{1}$ sont {\'e}quivalentes (introduire l'{\'e}quation diff{\'e}rentielle $y'' + y = g$).
    
\end{enumerate}
\finenonce{004762}



\finexercice


\exercice{4763, quercia, 2010/03/16}
\enonce{004763}{Norme de Frobenius}
Pour $A \in \mathcal{M}_n(\R)$, on pose $\|A\| = \sqrt{\mathrm{tr}({}^tAA)}$.

Montrer que c'est une norme et que : $\forall\ A,B \in \mathcal{M}_n(\R),\ \|AB\| \le \|A\|\times\|B\|$.

\finenonce{004763}



\finexercice


\exercice{4764, quercia, 2010/03/16}
\enonce{004764}{Semi-norme}
Soit $p$ une semi-norme sur $\mathcal{M}_n(\C)$ (ie. il manque juste l'axiome
$p(A)=0 \Rightarrow  A=0$). On suppose de plus que $\forall\ (A,B)\in(\mathcal{M}_n(\C))^2,\
p(AB)\le p(A)p(B)$. Montrer que $p=0$ ou $p$ est en fait une norme.

\finenonce{004764}



\finexercice


\exercice{4765, quercia, 2010/03/16}
\enonce{004765}{Normes produit}
Soient $E,F$ deux evn et $G = E \times F$.
On pose pour $u = (\vec x,\vec y) \in G$ :
$$\|u\|_1 = \|\vec x\,\|_E + \|\vec y\,\|_F,\quad
  \|u\|_2 = \sqrt{\|\vec x\,\|_E^2 + \|\vec y\,\|_F^2},\quad
  \|u\|_\infty = \max(\|\vec x\,\|_E , \|\vec y\,\|_F).$$

\begin{enumerate}
  \item Montrer que ce sont des normes sur $G$ et qu'elles sont deux {\`a} deux {\'e}quivalentes
    (sans hypoth{\`e}se de dimension finie).
  \item On prend $E=F$. Montrer que pour chacune de ces normes, l'application
    $G \to E, {(\vec x,\vec y\,)} \mapsto {\vec x+\vec y}$ est continue.
\end{enumerate}
\finenonce{004765}



\finexercice


\exercice{4766, quercia, 2010/03/16}
\enonce{004766}{Normes sur les suites}
Soit $E$ l'ensemble des suites $u=(u_n)$ r{\'e}elles born{\'e}es.
On pose $\begin{cases} \|u\| = \sup(|u_n| \text{ tq } n\in\N) \cr
                 N(u) = \sup(|u_n|+|u_{2n}| \text{ tq } n\in \N). \cr \end{cases}$

Montrer que ce sont des normes sur $E$ et qu'elles sont {\'e}quivalentes.
\finenonce{004766}



\finexercice


\exercice{4767, quercia, 2010/03/16}
\enonce{004767}{Norme sur les suites}
Soit $E$ l'ensemble des suites r{\'e}elles $u = (u_n)_{n\ge1}$ telles que la suite
$(\sqrt[n]{|u_n|}\,)$ est born{\'e}e.
Pour $u\in E$, on pose $\|u\| = \sup(\sqrt[n]{|u_n|}\text{ tq } n\in \N^*)$.
Montrer que $E$ est un $\R$-ev et que $\|\;.\;\|$ {\it n'est pas\/} une norme sur $E$.
\finenonce{004767}



\finexercice


\exercice{4768, quercia, 2010/03/16}
\enonce{004768}{Fonctions lipschitziennes}
Soit $E$ l'ensemble des fonctions $\R \to \R$ lipschitziennes.
Pour $f \in E$, on pose :
$$
\begin{aligned}
  \|f\| &= |f(0)| + \sup\left( \left|\frac{f(x)-f(y)}{x-y}\right|
    \text{ tq } x\ne y \right),
  \\
  N(f) &= |f(0)| + \sup\left( \left|\frac{f(x)-f(0)}{x}\right| \text{
      tq } x\ne 0 \right).
\end{aligned}
  $$
\begin{enumerate}
  \item Montrer que $E$ est un $\R$-ev.
  \item Montrer que $\|\;.\;\|$ et $N$ sont des normes sur $E$.
  \item Sont-elles {\'e}quivalentes ?
\end{enumerate}
\finenonce{004768}



\finexercice


\exercice{4769, quercia, 2010/03/16}
\enonce{004769}{Fonctions $\mathcal{C}^1$}
Soit $E$ l'ensemble des fonctions $f : {[0,1]} \to \R$ de classe $\mathcal{C}^1$.
Pour $f \in E$, on pose :
$N_1(f) = \sup(|f| + |f'|)$,
$N_2(f) = \sup|f| + \sup|f'|$.
Montrer que $N_1$ et $N_2$ sont deux normes {\'e}quivalentes sur $E$.
\finenonce{004769}



\finexercice


\exercice{4770, quercia, 2010/03/16}
\enonce{004770}{Norme sur les fonctions continues}
$E=C([0,1],\R)$. Soit $g\in E$. Pour tout $f\in E$ on pose
$N(f)=\sup\limits_{x\in [0,1]}\{|f(x)g(x)|\}$.

\begin{enumerate}
  \item Donner une condition n{\'e}cessaire et suffisante sur $g$ pour que $N$ soit une
norme sur $E$.
  \item Si pour tout $x\in [0,1]$, $g(x)\ne 0$, montrer qu'alors $N$ et $\|\ \|_{\infty}$
sont des normes sur $E$ {\'e}quivalentes.
  \item D{\'e}montrer la r{\'e}ciproque de la proposition pr{\'e}c{\'e}dente.
\end{enumerate}
\finenonce{004770}



\finexercice


\exercice{4771, quercia, 2010/03/16}
\enonce{004771}{Comparaison de normes (ENS MP 2002)}

\begin{enumerate}
  \item Soit $E$ un espace pr{\'e}hilbertien r{\'e}el et $u_1,\dots,u_n$ des {\'e}l{\'e}ments de~$E$.
    Calculer $\sum_\sigma\bigl\|\sum_{i=1}^n\sigma(i)u_i\bigr\|^2$ o{\`u} $\sigma$
    parcourt l'ensemble des fonctions de~$[[1,n]]$ dans $\{-1,1\}$.

  \item On se place dans l'ensemble des fonctions continues de~$[0,1]$ dans~$\R$.
    Montrer que la norme infinie n'est {\'e}quivalente {\`a} aucune norme euclidienne.

  \item M{\^e}me question avec la norme $\|\ \|_p$, $p\in{[1,+\infty[}\setminus\{2\}$.

\end{enumerate}
\finenonce{004771}



\finexercice


\exercice{4772, quercia, 2010/03/16}
\enonce{004772}{Jauge}
Soit $(E,\|\ \|)$ un $\R$-evn et $K\subset E$ une partie convexe, born{\'e}e,
sym{\'e}trique par rapport {\`a} l'origine et telle que
$0 \in \mathring K$.

Pour $x \in E$, on pose $n(x) = \inf\{|\lambda| \text{ tq } x \in \lambda K\}$.
Montrer que $n$ est une norme {\'e}quivalente {\`a} $\|\ \|$.
\finenonce{004772}



\finexercice


\exercice{4773, quercia, 2010/03/16}
\enonce{004773}{Polytechnique MP$^*$ 2006}
Soit $E$ un espace vectoriel r{\'e}el.
On consid{\`e}re une application $N : E \to {\R^+}$ telle que~:
    $$\begin{aligned}(i)
  \qquad&\forall\ \lambda,x,\ \ N(\lambda x) = |\lambda|N(x)\ ;\cr
               (ii)  \qquad&\forall\ x,\ N(x)=0 \Leftrightarrow x=0.\cr\end{aligned}$$

\begin{enumerate}
  \item Montrer que $N$ est une norme si et seulement $B = \{x$ tq. $N(x)\le 1\}$
    est convexe.
  \item Montrer que si $N$ v{\'e}rifie aussi
    $$\begin{aligned} (iii)  \qquad&\forall\ x,y,\ N(x+y)^2 \le 2N(x)^2 + 2N(y)^2\cr\end{aligned}$$
    alors c'est une norme.
    
\end{enumerate}
\finenonce{004773}



\finexercice


\exercice{4774, quercia, 2010/03/16}
\enonce{004774}{Parties de $\R^n$}
Les parties suivantes sont-elles ouvertes ? ferm{\'e}es ? born{\'e}es ?

\begin{enumerate}
  \item $A = \{ (x,y) \in \R^2 \text{ tq } xy=1 \}$.
  \item $B = \{ (x,y) \in \R^2 \text{ tq } x^2+xy+y^2 < 1 \}$.
  \item $C = \{ z \in \C \text{ tq } \Re(z^2) \le 1 \}$.
\end{enumerate}
\finenonce{004774}



\finexercice


\exercice{4775, quercia, 2010/03/16}
\enonce{004775}{Addition de parties}
Soient $A,B$ deux parties non vides d'un evn $E$.
On note $A+B = \{\vec a + \vec b \text{ tq } \vec a \in A,\ \vec b \in B \}$.
Montrer que \dots

\begin{enumerate}
  \item Si $A$ ou $B$ est ouvert, alors $A+B$ est ouvert.
  \item Si $A$ et $B$ sont ferm{\'e}s, alors $A+B$ n'est pas n{\'e}c{\'e}ssairement ferm{\'e}.
    (Prendre $A = \{(x,y) \in \R^2 \text{ tq } xy=1\}$ et $B = \{(x,0) \text{ tq } x\in \R\}$)
  \item Si $A$ et $B$ sont compacts, alors $A+B$ est compact.
\end{enumerate}
\finenonce{004775}



\finexercice


\exercice{4776, quercia, 2010/03/16}
\enonce{004776}{Voisinage ferm{\'e} d'un ferm{\'e}}
Soit $F$ un ferm{\'e} de $\R^n$ et $r > 0$.
On pose $F' = \bigcup_{\vec x \in F} \overline B(\vec x,r)$. Montrer que $F'$ est ferm{\'e}.
\finenonce{004776}



\finexercice


\exercice{4777, quercia, 2010/03/16}
\enonce{004777}{Ev engendr{\'e} par un ouvert}
Soit $\cal O$ un ouvert non vide d'un ev norm{\'e} $E$.
Montrer que $\text{vect}({\cal O}) = E$.
\finenonce{004777}



\finexercice


\exercice{4778, quercia, 2010/03/16}
\enonce{004778}{Adh{\'e}rence et int{\'e}rieur d'un sev}
Soit $E$ un evn et $F$ un sev de $E$.

\begin{enumerate}
  \item Montrer que $\overline{F}$ est un sev de $E$.
  \item Si $E$ est de dimension finie, montrer que $F = \overline F$.
  \item Dans le cas g{\'e}n{\'e}ral, montrer que $\mathring F = \varnothing$ ou $F = E$.
\end{enumerate}
\finenonce{004778}



\finexercice


\exercice{4779, quercia, 2010/03/16}
\enonce{004779}{C{\^o}ne convexe engendr{\'e} par un ensemble fini, ENS ULM-Lyon-Cachan MP$^*$ 2005}

$E$ est un $\R$-espace vectoriel norm{\'e} et $a_{1},\ldots,a_{n}\in E$. On pose $C=\Bigl\{\sum_{i=1}^{n}\lambda_{i}a_{i},\lambda_{i}\ge 0\Bigr\}$.

\begin{enumerate}
  \item Montrer que, pour tout $x\in E$, il existe $c\in C$ tel que $\|x-c\|=\inf \{\|x-a\|,a\in C\}$.


  \item En d{\'e}duire que $C$ est ferm{\'e}.
\end{enumerate}
\finenonce{004779}



\finexercice


\exercice{4780, quercia, 2010/03/16}
\enonce{004780}{Partie convexe dense}
Soit $E$ un evn de dimension finie et $C \subset E$ convexe et dense.
Montrer que $C=E$.
\finenonce{004780}



\finexercice


\exercice{4781, quercia, 2010/03/16}
\enonce{004781}{L'ensemble des projecteurs est ferm{\'e}}
Soit $E$ un evn de dimension finie et $\cal P$ l'ensemble des projecteurs
de $E$. Montrer que $\cal P$ est ferm{\'e} dans $\mathcal{L}(E)$.
\finenonce{004781}



\finexercice


\exercice{4782, quercia, 2010/03/16}
\enonce{004782}{Adh{\'e}rence et int{\'e}rieur dans les fonctions continues}

\begin{enumerate}
  \item Soit $E = \mathcal{C}([0,1],\R)$ muni de la norme de la convergence uniforme.
    Soit $P$ l'ensemble des fonctions de $E$ positives ou nulles.
    Chercher $\overline P$ et $\mathring P$.
    
  \item M{\^e}mes questions avec la norme : $\|f\| =  \int_{t=0}^1 |f(t)|\,dt$.
    
\end{enumerate}
\finenonce{004782}



\finexercice


\exercice{4783, quercia, 2010/03/16}
\enonce{004783}{Mines MP 2000}
On pose $E=C([0,1],\R)$ et on le munit de la norme $N_{\infty}$. Soit 
$F=\{f\in E \, | \, f(0)=f(1)\}$.
D{\'e}terminer l'adh{\'e}rence et l'int{\'e}rieur de $F$.

\finenonce{004783}



\finexercice


\exercice{4784, quercia, 2010/03/16}
\enonce{004784}{Points isol{\'e}s (Ens Ulm MP$^*$ 2003)}

Les solutions de l'{\'e}quation $u^2=\mathrm{id}_{\R^n}$ pour~$u\in\mathcal{L}({\R^n})$
sont-elles isol{\'e}es~?

\finenonce{004784}



\finexercice


\exercice{4785, quercia, 2010/03/16}
\enonce{004785}{Adh{\'e}rence et int{\'e}rieur d'un convexe}
Soit $A$ une partie convexe d'un evn $E$.

\begin{enumerate}
  \item D{\'e}montrer que $\overline A$ et $\mathring A$ sont aussi convexes
    (pour $\mathring A$ : faire un dessin).
  \item Montrer que l'application $x \mapsto d(x,A)$ est convexe (c.a.d.
    $d(tx+(1-t)y,A) \le td(x,A) + (1-t)d(y,A)$).
\end{enumerate}
\finenonce{004785}



\finexercice


\exercice{4786, quercia, 2010/03/16}
\enonce{004786}{Th{\'e}or{\`e}me des ferm{\'e}s emboit{\'e}s}
Soit $E$ un evn de dimension finie, et $(B_n = B(\vec a_n, r_n))$ une suite de
boules ferm{\'e}es, d{\'e}croissante pour l'inclusion, tq $r_n\xrightarrow[n\to\infty]{} 0$.

\begin{enumerate}
  \item Montrer que la suite $(\vec a_n)$ admet une sous-suite convergeant vers $\vec a \in E$.
  \item Montrer que $\vec a_n \xrightarrow[n\to\infty]{} \vec a$.
  \item Montrer que $\bigcap_{n\in \N} B_n = \{\vec a \}$.
\end{enumerate}
\finenonce{004786}



\finexercice


\exercice{4787, quercia, 2010/03/16}
\enonce{004787}{X MP$^*$ 2001}

On consid{\`e}re l'espace $\mathcal{M}_n(\C)$ muni d'une norme quelconque.

\begin{enumerate}
  \item Montrer que $GL_n(\C)$ est ouvert dense de $\mathcal{M}_n(\C)$.
    

  \item Soit $D_n(\C)$ l'ensemble des matrices diagonalisables de $\mathcal{M}_n(\C)$.
    Montrer que $D_n(\C)$ est dense dans $\mathcal{M}_n(\C)$.
    

  \item Quel est l'int{\'e}rieur de $D_n(\C)$~?

\end{enumerate}
\finenonce{004787}



\finexercice


\exercice{4788, quercia, 2010/03/16}
\enonce{004788}{Matrices nilpotentes, ENS Ulm-Lyon-Cachan MP$^*$ 2006}
Soit $N\in\mathcal{M}_n(\C)$. Montrer que $N$ est nilpotente si et seulement si
la matrice nulle est adh{\'e}rente {\`a} l'ensemble
$\{P^{-1}NP,\ P\in GL_n(\C)\}$.


\finenonce{004788}



\finexercice


\exercice{4789, quercia, 2010/03/16}
\enonce{004789}{Polyn{\^o}mes scind{\'e}s (Ens Ulm-Lyon-Cachan MP$^*$ 2003)}
Soit~$n\in\N^*$ et~$\sigma\in\R^n$.
On note~$P_\sigma = X^n - \sigma_1 X^{n-1} + \dots + (-1)^{n-1}\sigma_{n-1}X + (-1)^n\sigma_n$.

Soit~$\Omega = \{\sigma\in\R^n$ tq $P_\sigma$ est {\`a} racines r{\'e}elles, distinctes$\}$.

\begin{enumerate}
  \item $\Omega$ est-il ouvert~? ferm{\'e}~?
    

  \item Notons~$f$ : $\sigma  \mapsto P_\sigma$. D{\'e}terminer $f(\overline\Omega)$.
    
\end{enumerate}
\finenonce{004789}



\finexercice


\exercice{4790, quercia, 2010/03/16}
\enonce{004790}{$(f(x)-f(y))/(x-y)$}
Soit $f:\R \to \R$ une fonction de classe $\mathcal{C}^1$ et
$g : {\R^2} \to \R$ $\begin{cases}{(x,y)} &{\frac {f(x)-f(y)}{x-y} \text{ si } x \ne y}\cr
 {(x,x)} &{f'(x)}\cr\end{cases}$

Montrer que $g$ est continue.
(Attention : pour une fonction d{\'e}finie par cas, se placer au voisinage d'un
point $(x_0,y_0)$ et d{\'e}terminer si un seul ou plusieurs cas sont {\`a} consid{\'e}rer
dans ce voisinage)
\finenonce{004790}



\finexercice


\exercice{4791, quercia, 2010/03/16}
\enonce{004791}{$\sup(f(x,y))$}
Soit $f : {\R^2} \to \R$ continue.
On pose $g(x) = \sup(f(x,y) \text{ tq } y \in {[0,1]})$.
Montrer que $g$ est continue.

\finenonce{004791}



\finexercice


\exercice{4792, quercia, 2010/03/16}
\enonce{004792}{Fonction tendant vers $+\infty$ {\`a} l'infini}
Soit $E$ un evn de dimension finie et $f : E \to\R$ continue.
On suppose que $f(\vec x\,) \xrightarrow[\|\vec x\,\|\to\infty]{} +\infty$, c'est {\`a} dire :
$$\forall\ A\in\R,\ \exists\ B\in\R\text{ tq } \forall\ \vec x\in E,\
  \|\vec x\,\| \ge B  \Rightarrow  f(\vec x\,)\ge A.$$

\begin{enumerate}
  \item On prend $A = f(\vec 0)$ et $B$ le nombre correspondant.\par
    Montrer que $\inf\{f(\vec x\,) \text{ tq } \vec x\in E\} =
         \inf\{f(\vec x\,) \text{ tq } \|\vec x\,\| \le B \}$.
  \item En d{\'e}duire que $f$ admet un minimum.
  \item Exemple : soit $E = \R_n[X]$ et $f : {[a,b]} \to \R$ born{\'e}e.\par
    Montrer qu'il existe $P\in E \text{ tq } \|f-P\|_\infty = \sup\{|f(t)-P(t)| \text{ tq } t\in{[a,b]} \}$
    soit minimal
    ($P$ est appel{\'e} : \emph{un} polyn{\^o}me de meilleure approximation de $f$ sur $[a,b]$).
\end{enumerate}
\finenonce{004792}



\finexercice


\exercice{4793, quercia, 2010/03/16}
\enonce{004793}{Fonctions homog{\`e}nes}
On note $\Omega = \R^2\setminus \{(0,0)\}$ et $D = \{(x,y)\in\R^2 \text{ tq } 0 < x^2+y^2 \le 1 \}$.
Soit $f : \Omega \to \R$. On dit que $f$ est
{\it positivement homog{\`e}ne de degr{\'e} $\alpha$} si :

$$ \forall\ (x,y) \in \Omega, \forall\ t > 0,\ f(tx,ty) = t^\alpha f(x,y). $$

\begin{enumerate}
  \item Donner des exemples de telles fonctions pour $\alpha = 1,0,-2, \frac 12$.
  \item Soit $f$ continue, positivement homog{\`e}ne de degr{\'e} $\alpha$.
    Montrer que si $\alpha \ge 0$, $f$ est born{\'e}e sur $D$ et que si $\alpha > 0$,
    $f$ admet une limite finie en $(0,0)$.
    Examiner les cas $\alpha = 0$, $\alpha < 0$.
\end{enumerate}
\finenonce{004793}



\finexercice


\exercice{4794, quercia, 2010/03/16}
\enonce{004794}{Fonction partiellement continue dans toutes les directions}
Trouver une fonction $f : {\R^2} \to \R$ discontinue en $(0,0)$
mais telle que pour tous $\alpha,\beta \in \R$,
$f(\alpha t,\beta t) \xrightarrow[t\to0]{} 0$.
\finenonce{004794}



\finexercice


\exercice{4795, quercia, 2010/03/16}
\enonce{004795}{Continuit{\'e} du polyn{\^o}me caract{\'e}ristique}
Soit $E = \mathcal{M}_n(\R)$, $F = \R_n[X]$, et
$\varphi : E \to F, A \mapsto {P_A}$ (polyn{\^o}me caract{\'e}ristique).
Montrer que $\varphi$ est continue.
\finenonce{004795}



\finexercice


\exercice{4796, quercia, 2010/03/16}
\enonce{004796}{Ouverts et non ouverts}
Soit $E = \R[X]$. Pour $P \in E$, on pose :
$\begin{cases} N_1(P) = \sup(|P(t)| \text{ tq } 0 \le t \le 1 )\cr
         N_2(P) = \sup(|P(t)| \text{ tq } 1 \le t \le 2 )\cr
         \varphi(P) = P(0). \cr \end{cases}$

\begin{enumerate}
  \item V{\'e}rifier que $N_1$ et $N_2$ sont des normes.
  \item Montrer que $\varphi$ est continue pour $N_1$.
  \item Montrer que $\varphi$ est discontinue pour $N_2$. (Consid{\'e}rer $P_n(t) = (1-t/2)^n$)
  \item $N_1$ et $N_2$ sont-elles {\'e}quivalentes ?
  \item Soit ${\cal O} = \{ P \in E \text{ tq } P(0) \ne 0 \}$.
    Montrer que $\cal O$ est ouvert pour $N_1$ mais pas pour $N_2$.
\end{enumerate}
\finenonce{004796}



\finexercice


\exercice{4797, quercia, 2010/03/16}
\enonce{004797}{Thm du point fixe}
Soit $E$ un evn de dimension finie et $f : E \to E$ une fonction
$k$-lipchitzienne avec $k < 1$.
On choisit ${\vec u}_0 \in E$ arbitrairement, et on consid{\`e}re la suite
$({\vec u}_n)$ telle que pour tout $n$ :
${\vec u}_{n+1} = f({\vec u}_n)$.

\begin{enumerate}
  \item Montrer que $\|{\vec u}_{n+1} - {\vec u}_n\| \le k^n \|{\vec u}_1 - {\vec u}_0\|$.
  \item En d{\'e}duire que la suite $({\vec u}_n)$ est de Cauchy.
  \item Soit $\vec \ell = \lim({\vec u}_n)$. Montrer que $\vec \ell$ est l'unique solution
    dans $E$ de l'{\'e}quation $f(\vec x\,) = \vec x$.
\end{enumerate}  
\finenonce{004797}



\finexercice


\exercice{4798, quercia, 2010/03/16}
\enonce{004798}{Points fixes, ULM-Lyon-Cachan MP$^*$ 2005}
\begin{enumerate}
  \item Montrer que les points fixes de $f$, continue sur $[0,1]$, {\`a} valeurs dans $[0,1]$,
    forment un ensemble ferm{\'e} non vide.

  \item Montrer que tout ferm{\'e} de $[0,1]$ non vide est l'ensemble des points fixes
    d'une fonction continue de $[0,1]$ dans~$[0,1]$.

\end{enumerate}
\finenonce{004798}



\finexercice


\exercice{4799, quercia, 2010/03/16}
\enonce{004799}{Racines de polyn{\^o}mes X MP$^*$ 2004}
Soit~$E=\C_d[X]$ norm{\'e} par $\|P\| = \sum|a_i|$,
$P\in E$ de degr{\'e}~$d$ {\`a} racines simples et $P_n$ une suite de polyn{\^o}mes de~$E$ convergeant vers~$P$.

Soit~$z\in\C$ tel que $P(z) = 0$ et $\delta>0$.

\begin{enumerate}
  \item Montrer que pour~$n$ assez grand, $P_n$ a au moins un z{\'e}ro dans $\overline{B(z,\delta)}$.
  \item Montrer qu'il existe $\delta_0>0$ tel que pour tout $\delta\in{]0,\delta_0]}$
    $P_n$ a exactement une racine dans $\overline{B(z,\delta)}$ si $n$ est assez grand.
    
  \item Que peut-on dire si les z{\'e}ros de~$P$ ne sont plus suppos{\'e}s simples~?
    
\end{enumerate}
\finenonce{004799}



\finexercice


\exercice{4800, quercia, 2010/03/16}
\enonce{004800}{Une application polynomiale est ferm{\'e}e, ULM-Lyon-Cachan MP$^*$ 2005}

Soit $f$ une fonction polynomiale sur~$\C$. Montrer que l'image par $f$ de tout
ferm{\'e} est un ferm{\'e}.



\finenonce{004800}



\finexercice


\exercice{4801, quercia, 2010/03/16}
\enonce{004801}{Principe du maximum, ULM-Lyon-Cachan MP$^*$ 2005}

Soit $P\in\C[X]$ et $U$ un ouvert de~$\C$ born{\'e}.
Montrer que $\sup(|P(x)|,\ x\in U) = \sup(|P(x)|,\ x\in \mathrm{Fr}(U))$.


\finenonce{004801}



\finexercice


\exercice{4802, quercia, 2010/03/16}
\enonce{004802}{Fonction presque additive, Centrale MP 2001}

Soient $E,F$ deux $\R$-espaces vectoriels norm{\'e}s, $F$ {\'e}tant complet.
Soit $f$ une application continue de $E$ dans $F$ telle qu'il existe $M\in\R^+$
v{\'e}rifiant~:
$$\forall\ x,y\in E,\ \|f(x+y)-f(x)-f(y)\| \le M.$$

\begin{enumerate}
  \item Dans le cas $M=0$ montrer que $f$ est lin{\'e}aire. Ce r{\'e}sultat subsiste-t-il
    si $E$ et $F$ sont des $\C$-ev~?
    

  \item On suppose $M>0$. Soit pour $x\in E$ et $n\in\N$~: $f_n(x) = 2^{-n}f(2^nx)$.
    Montrer que la suite $(f_n)$ converge simplement sur~$E$.
    

  \item On note $g = \lim_{n\to\infty}f_n$. Montrer que $g$ est une application
    lin{\'e}aire continue et que c'est l'unique application lin{\'e}aire telle que
    $f-g$ soit born{\'e}e.
    
\end{enumerate}
\finenonce{004802}



\finexercice


\exercice{4803, quercia, 2010/03/16}
\enonce{004803}{$f$ uc $ \Rightarrow  f$(born{\'e}) est born{\'e}}

Soit $A\subset E$ une partie non vide born{\'e}e et $f : A  \to F$ uniform{\'e}ment continue.
Montrer que $f$ est born{\'e}e dans les cas suivants~:

\begin{enumerate}
  \item $A$ est convexe.
  \item $A$ est connexe par arcs.
  \item $A$ est quelconque et $E$ est de dimension finie.
\end{enumerate}
\finenonce{004803}



\finexercice


\exercice{4804, quercia, 2010/03/16}
\enonce{004804}{Applications lin{\'e}aires continues}
\def\N{|||}%
Soit $E$ un $\R$-ev de dimension finie et $u \in \mathcal{L}(E)$.
On pose $\N u \N = \sup(\|u(\vec x)\| \text{ tq } \|\vec x\,\| = 1)$.

\begin{enumerate}
  \item Montrer que $\N u \N$ existe et que c'est un maximum.
  \item Montrer que $\N\;.\;\N$ est une norme sur $\mathcal{L}(E)$
    (appel{\'e}e : {\it norme lin{\'e}aire associ{\'e}e {\`a}\/} $\|\;.\;\|$).
  \item Montrer que : $\forall\ \vec x \in E,\ \|u(\vec x)\| \le \N u \N \times \|\vec x\,\|$.
  \item En d{\'e}duire que : $\forall\ u,v \in \mathcal{L}(E),\ \N u\circ v \N \le \N u \N \times \N v \N$.
\end{enumerate}
\finenonce{004804}



\finexercice


\exercice{4805, quercia, 2010/03/16}
\enonce{004805}{Centrale MP 2006}
$E$ est l'ensemble des fonctions de $\R$ dans $\R$ continues et born{\'e}es sur $\R$.

Pour $p \in \N$ et $f \in E$ on pose : 
$N_{p}(f) = \sup\{|t^{p} e^{-|t|} f(t)|,\ t\in\R\}$.

\begin{enumerate}
  \item Montrer que $N_{p}$ est une norme sur $E$.

  \item Soit $c \in \R$ et ${P_{c}} : E \to \R, f \mapsto {f(c).}$
    {\'E}tudier la continuit{\'e} de $P_c$ sur $(E,N_p)$.
    

  \item Montrer que, pour $p$ et $q$ distincts dans $\N$, les normes $N_p$ et $N_q$ ne sont pas {\'e}quivalentes.
\end{enumerate}
\finenonce{004805}



\finexercice


\exercice{4806, quercia, 2010/03/16}
\enonce{004806}{Applications lin{\'e}aires continues}
Soient $E,F$ deux evn de dimensions finies et $\varphi : E \to F$ lin{\'e}aire.
Montrer que $\varphi$ est continue.

En d{\'e}duire que tout sev de $E$ est ferm{\'e}.
\finenonce{004806}



\finexercice


\exercice{4807, quercia, 2010/03/16}
\enonce{004807}{Continuit{\'e} du polyn{\^o}me de Lagrange}
Soient $x_1,\dots,x_n \in \R$ distincts.
{\`A} $(y_1,\dots,y_n) \in \R^n$ on fait correspondre le polyn{\^o}me
$P \in \R_{n-1}[X]$ tel que pour tout $i$ : $P(x_i) = y_i$.

\begin{enumerate}
  \item Montrer que l'application $(y_1,\dots,y_n)  \mapsto P$ est continue.
  \item Montrer que l'application r{\'e}ciproque est aussi continue.
\end{enumerate}
\finenonce{004807}



\finexercice


\exercice{4808, quercia, 2010/03/16}
\enonce{004808}{Ensi PSI 1998}

Soit $\mathcal{C}(\R)$ l'ensemble des suites r{\'e}elles convergentes muni de la
norme~$\|u\| = \sup\{|u_n|, n\in\N\}$.

Soit $L : {\mathcal{C}(\R)} \to \R, u \mapsto {\ell = \lim_{n\to\infty} u_n.}$
Montrer que $L$ est une application lin{\'e}aire continue et calculer sa norme.
\finenonce{004808}



\finexercice


\exercice{4809, quercia, 2010/03/16}
\enonce{004809}{It{\'e}ration d'un endomorphisme}
Soit $E$ un evn de dimension finie et $u \in \mathcal{L}(E)$.
On choisit ${\vec x}_0 \in E$, et on consid{\`e}re la suite $({\vec x}_n)$
d{\'e}finie par la relation de r{\'e}currence :
${\vec x}_{n+1} = \frac{u({\vec x}_n)}{\|u({\vec x}_n)\|}$.

On suppose que la suite $({\vec x}_n)$ converge. Montrer que la limite est un vecteur
propre de $u$.

Exemples :
{\bf 1)} $E = \R^3$, $\text{mat}(u) = \left(\begin{smallmatrix} 1 &2 &3 \cr 4 &5 &6 \cr 7 &8 &9 \cr\end{smallmatrix}\right)$.
{\bf 2)} $E = \R^3$, $\text{mat}(u) = \left(\begin{smallmatrix} 0 &0 &1 \cr 1 &0 &0 \cr 0 &1 &0 \cr\end{smallmatrix}\right)$.
\finenonce{004809}



\finexercice


\exercice{4810, quercia, 2010/03/16}
\enonce{004810}{Puissances de $u$}
Soit $E$ un evn de dimension finie et $u \in \mathcal{L}(E)$ tel que $\|u\| \le 1$.
Montrer que la suite $(u^n)$ contient une sous-suite simplement convergente.

\finenonce{004810}



\finexercice


\exercice{4811, quercia, 2010/03/16}
\enonce{004811}{$\mathrm{id} - u$ est bicontinu}
Soit $E$ un $\C$-evn et $u \in {\cal L}_c(E)$ tel que $\mathrm{id}_E - u$ est
bicontinu.
Montrer que pour tout entier $n$, $\mathrm{id}_E - u^n$ est bicontinu et comparer son inverse
{\`a} $\sum_{k=0}^{n-1} \bigl(\mathrm{id}_E - e^{2ik\pi/n}u\bigr)^{-1}$.


\finenonce{004811}



\finexercice


\exercice{4812, quercia, 2010/03/16}
\enonce{004812}{Norme lin{\'e}aire sur $\R$ = norme lin{\'e}aire sur $\C$}
Soit $A \in \mathcal{M}_n(\R)$ et $f \in \mathcal{L}(\R^n)$, $g\in \mathcal{L}({\C^n})$ les endomorphismes
canoniquement associ{\'e}s {\`a} $A$.
Montrer que si on munit $\R^n$ et $\C^n$ des normes euclidiennes usuelles, alors
$\|f\| = \|g\|$.

\finenonce{004812}



\finexercice


\exercice{4813, quercia, 2010/03/16}
\enonce{004813}{Applications lin{\'e}aires sur les polyn{\^o}mes}
Soit $E = \R[x]$ muni de la norme :
$\bigl\|\sum\limits_i a_ix^i\bigr\| = \sum\limits_i|a_i|$.

\begin{enumerate}
  \item Est-ce que $\varphi : P  \mapsto P(x+1)$ est continue ?
    
  \item Est-ce que $\psi : P  \mapsto AP$ est continue ? ($A \in E$ fix{\'e})
    

  \item Reprendre les questions pr{\'e}c{\'e}dentes avec la norme :
    $\|P\| = \sup\{e^{-|t|}|P(t)|,\ t\in \R\}$.
    
\end{enumerate}
\finenonce{004813}



\finexercice


\exercice{4814, quercia, 2010/03/16}
\enonce{004814}{$uv - vu = \mathrm{id}$}
Soit $E$ un espace vectoriel norm{\'e} et $u,v \in \mathcal{L}(E)$ tels que
$u\circ v - v \circ u = \mathrm{id}_E$.

\begin{enumerate}
  \item Calculer $u\circ v^n - v^n \circ u$ pour $n \in \N^*$.
    
  \item Montrer que $u$ ou $v$ est discontinu.
    
\end{enumerate}
\finenonce{004814}



\finexercice


\exercice{4815, quercia, 2010/03/16}
\enonce{004815}{La d{\'e}rivation peut-elle {\^e}tre continue~?}
On note $E=\mathcal{C}^\infty([0,+\infty[,\R)$ et $D$ l'endomorphisme de $E$ de d{\'e}rivation~:
$D(f) = f'$.

\begin{enumerate}
  \item Montrer qu'il n'existe aucune norme sur~$E$ pour laquelle $D$ soit continu
    (consid{\'e}rer $x \mapsto e^{\alpha x}$).
  \item Soit $F$ le sous-ev de~$E$ constitu{\'e} des fonctions polynomiales.
    Trouver une norme sur~$F$ pour laquelle $D_{|F}$ est continu.

\end{enumerate}
\finenonce{004815}



\finexercice


\exercice{4816, quercia, 2010/03/16}
\enonce{004816}{Mines MP 2002}
On munit $E_k = \R_k[X]$ de la norme $\|P\|_k = \sum_{i=0}^k |P(i)|$.
Calculer $\|\hskip-1pt|\varphi\|\hskip-1pt|$ avec
$\varphi : {E_2} \to {E_3}, P \mapsto {X^2P'.}$
\finenonce{004816}



\finexercice


\exercice{4817, quercia, 2010/03/16}
\enonce{004817}{Normes sur les suites born{\'e}es}
Soit $E$ l'ensemble des suites r{\'e}elles $u = (u_n)$ born{\'e}es et $F$
le sev des suites telles que la s{\'e}rie de terme g{\'e}n{\'e}ral $|u_n|$ converge.
Pour $u \in E$, on pose $\|u\|_\infty = \sup\limits_n |u_n|$ et
pour $u \in F$ : $\|u\|_1 = \sum\limits_n |u_n|$.

Soit $a \in E$ et $f : E \to E, u  \mapsto {au = (a_nu_n).}$

\begin{enumerate}
  \item Montrer que $f$ est une application lin{\'e}aire continue de $E$ dans $E$ et calculer
    sa norme.
  \item Montrer que $F$ est stable par $f$ et calculer la norme de $f_{|F}$ quand on prend
    la norme $\|\ \|_1$ sur $F$.
\end{enumerate}
\finenonce{004817}



\finexercice


\exercice{4818, quercia, 2010/03/16}
\enonce{004818}{Thm de l'hyperplan ferm{\'e}}
Soit $E$ un $\R$-evn et $f \in E^*$.

\begin{enumerate}
  \item Montrer que $f$ est continue si et seulement si $\mathrm{Ker} f$ est ferm{\'e}
    (pour la r{\'e}ciproque : supposer $\mathrm{Ker} f$ ferm{\'e}, montrer que $\{x \text{ tq } f(x) > 0\}$
     est ouvert, puis {\'e}tudier $\{ x \text{ tq } -1 < f(x) < 1\}$).
  \item On suppose $f$ continue. Soit $x \in E$. Montrer que $|f(x)| = \|f\|d(x,\mathrm{Ker} f)$.
    
\end{enumerate}
\finenonce{004818}



\finexercice


\exercice{4819, quercia, 2010/03/16}
\enonce{004819}{Th{\'e}or{\`e}me de Hahn-Banach (Polytechnique MP$^*$ 2003)}
Soit~$E$ un espace vectoriel norm{\'e} de dimension finie et~$F$ un hyperplan de~$E$.
Soit~$\varepsilon\in E$ tel que~$\R\varepsilon$ soit suppl{\'e}mentaire de~$F$.
Soit~$f$ une forme lin{\'e}aire sur~$F$.

\begin{enumerate}
  \item Montrer que~: 
    $\forall\ x_1,x_2\in F,\ f(x_1)-\|\kern-1.2pt|f\|\kern-1.2pt|\;\|x_1-\varepsilon\|
    \le \|\kern-1.2pt|f\|\kern-1.2pt|\;\|x_2+\varepsilon\| - f(x_2)$.
    
  \item Montrer qu'il existe~$\alpha\in\R$ tel que~:
    $\forall\ x_1,x_2\in F,\ f(x_1)-\|\kern-1.2pt|f\|\kern-1.2pt|\;\|x_1-\varepsilon\|
    \le \alpha \le \|\kern-1.2pt|f\|\kern-1.2pt|\;\|x_2+\varepsilon\| - f(x_2)$.

  \item On d{\'e}finit $\varphi : E \to \R$ par~$\varphi_{|F} = f$ et $\varphi(\varepsilon) = \alpha$.
    Montrer que~$\|\kern-1.2pt|\varphi\|\kern-1.2pt| = \|\kern-1.2pt|f\|\kern-1.2pt|$.
    

  \item On consid{\`e}re~$E = \{u = (u_n)_{n\in\N}\in\R^\N\text{ tq }\sum_{n\in\N}|u_n|<+\infty\}$
    avec la norme d{\'e}finie par~: $\|u\| = \sum_{n\in\N}|u_n|$.
    Montrer que~$E$ est complet pour cette norme.
    

  \item Donner une famille d{\'e}nombrable de sev de~$E$ de dimensions finies dont la r{\'e}union
    est dense dans~$E$.
    

  \item Soit~$F$ un sev de~$E$ de dimension finie et~$f$ une forme lin{\'e}aire sur~$F$.
    Montrer qu'il existe une forme lin{\'e}aire~$\varphi$ sur~$E$ telle que
    $\varphi_{|F} = f$ et $\|\kern-1.2pt|\varphi\|\kern-1.2pt| = \|\kern-1.2pt|f\|\kern-1.2pt|$.
    
\end{enumerate}
\finenonce{004819}



\finexercice


\exercice{4820, quercia, 2010/03/16}
\enonce{004820}{Rayon spectral}
Soit $E$ un evn de dimension finie et $u \in \mathcal{L}(E)$. On pose $x_n = \|u^n\|$.

\begin{enumerate}
  \item Montrer que $\rho = \inf\{\sqrt[n]{x_n},\ n\in \N\}$ est ind{\'e}pendant de la norme
    choisie sur $E$.
  \item En utilisant l'in{\'e}galit{\'e} : $x_{p+q} \le x_p x_q$, montrer que
    la suite $(\sqrt[n] {x_n}\,)$ converge vers $\rho$.
\end{enumerate}
\finenonce{004820}



\finexercice


\exercice{4821, quercia, 2010/03/16}
\enonce{004821}{Rayon spectral (Centrale MP 2001)}

Soit $E$ un espace vectoriel complexe de dimension finie.
On consid{\`e}re un endomorphisme $f$ de~$E$ et on note
$\rho(f) = \sup\{|\lambda|\text{ tq }\lambda\in\mathrm{Sp}(f)\}$
(rayon spectral de~$f$).
Soit $\nu$ une norme sur $\mathcal{L}(E)$.

\begin{enumerate}
  \item Montrer que $\rho(f)\le \lim_{p\to\infty}(\nu(f^p)^{1/p})$.
    On pourra pour commencer supposer que $\nu$ est la norme subordonn{\'e}e
    {\`a} une norme sur~$E$.
    

  \item Montrer que si $f$ est diagonalisable l'in{\'e}galit{\'e} pr{\'e}c{\'e}dente est une {\'e}galit{\'e}.
    

  \item {\'E}tudier le cas g{\'e}n{\'e}ral.
    
\end{enumerate}
\finenonce{004821}



\finexercice


\exercice{4822, quercia, 2010/03/16}
\enonce{004822}{Polytechnique MP$^*$ 2000}
Soit~$u$ une application lin{\'e}aire de~$\R^n$ dans~$\R^m$. Prouver que
$u$ est surjective si et seulement si elle transforme tout ouvert
de~$\R^n$ en ouvert de~$\R^m$.
\finenonce{004822}



\finexercice


\exercice{4823, quercia, 2010/03/16}
\enonce{004823}{Connexit{\'e} d'un evn}
Soit $E$ un evn et $A \subset E$. On suppose que $A$ est {\`a} la fois ouvert
et ferm{\'e}.

\begin{enumerate}
  \item Exemples de telles parties ?
  \item On d{\'e}finit la fonction $f : E \to {\R}$ par
  $\begin{cases}
    \vec x &1 \text{ si } \vec x \in A,\cr
     \vec x &0 \text{ si } \vec x \notin A\cr \end{cases}$
  \begin{enumerate}
    \item  Montrer que $f$ est continue.
    \item On prend $\vec a \in A$ et $\vec b \notin A$.
        Montrer que $\varphi : {[0,1]} \to {\R}, t \mapsto {f(t\vec a + (1-t)\vec b)}$
        est continue.
    \item Conclure.
  \end{enumerate}   
\end{enumerate}
\finenonce{004823}



\finexercice


\exercice{4824, quercia, 2010/03/16}
\enonce{004824}{$A \ne E$ et $A \ne \varnothing => \mathrm{Fr}(A) \ne\varnothing$}
Soit $E$ un evn et $A$ une partie de $E$ ni vide, ni {\'e}gale {\`a} $E$.
Montrer que $\mathrm{Fr}(A) \ne \varnothing$.
\finenonce{004824}



\finexercice


\exercice{4825, quercia, 2010/03/16}
\enonce{004825}{$A\cup B$ ferm{\'e} $ \Rightarrow $ $A\cup B =E$.}
Soit $E$ un evn de dimension sup{\'e}rieure ou {\'e}gale {\`a} $2$
et $A,B$ deux parties de $E$ telles que
$A$ est ouvert non vide, $B$ est fini et $A\cup B$ est ferm{\'e}.
Montrer que $A\cup B = E$.
\finenonce{004825}



\finexercice


\exercice{4826, quercia, 2010/03/16}
\enonce{004826}{Compl{\'e}mentaire d'un hyperplan (Ens Ulm MP$^*$ 2005)}

Soit $E$ un evn r{\'e}el et $H$ un hyperplan de~$E$.
Montrer que $E\setminus H$ est connexe par arcs si et seulement si
$H$ {\it n'est pas\/} ferm{\'e}.
\finenonce{004826}



\finexercice


\exercice{5839, rouget, 2010/10/16}
\enonce{005839}{**}
\label{ex:rou1bis}
Montrer que la boule unité d'un espace vectoriel normé est un convexe de cet espace.
\finenonce{005839}


\finexercice


\exercice{5840, rouget, 2010/10/16}
\enonce{005840}{*** I}
\begin{enumerate}
 \item  Inégalités de \textsc{Hölder} et de \textsc{Minkowski}. Soit $(p,q)\in]0,+\infty[^2$  tel que $ \frac{1}{p}+ \frac{1}{q}=1$. 
   \begin{enumerate}
   \item Montrer que pour $(x,y)\in[0,+\infty[^2$, $xy\leqslant \frac{x^p}{p}+ \frac{x^q}{q}$.  

   \item En déduire que $\forall((a_1,...,a_n),(b_1,...,b_n))\in(\Rr^n)^2$, $\left|\sum_{k=1}^{n}a_kb_k\right|\leqslant\left(\sum_{k=1}^{n}|a_k|^p\right)^{1/p}\left(\sum_{k=1}^{n}|b_k|^q\right)^{1/q}$.

   \item En déduire que $\forall((a_1,...,a_n),(b_1,...,b_n))\in(\Rr^n)^2$,  $\left(\sum_{k=1}^{n}\left|a_k+b_k\right|^p\right)^{1/p}\leqslant\left(\sum_{k=1}^{n}\left|a_k\right|^p\right)^{1/p}+\left(\sum_{k=1}^{n}\left|b_k\right|^p\right)^{1/p}$.
   \end{enumerate}

 

\item  Soit $\alpha$ un réel strictement positif. Pour $x =(x_1,...,x_n)\in\Rr^n$, on définit $N_\alpha(x)=\left(\sum_{k=1}^{n}|x_k|^\alpha\right)^{1/\alpha}$.
   \begin{enumerate}
   \item Montrer que $\forall\alpha\geqslant1$, $N_\alpha$ est une norme sur $\Rr^n$.

   \item Dessiner les \og boules unités \fg~de $\Rr^2$ dans le cas où $\alpha\in\left\{ \frac{2}{3},1, \frac{3}{2},2,+\infty\right\}$.

   \item Montrer que, pour $x=(x_k)_{1\leqslant k\leqslant n}$ fixé, $\lim_{\alpha \rightarrow +\infty}N_\alpha(x) =\text{Max}\{|x_k|,\;1\leqslant k\leqslant n\}= N_\infty(x)$.

   \item Montrer que si $0 <\alpha< 1$, $N_\alpha$ n'est pas une norme sur $\Rr^n$ (si $n\geqslant2)$.
   \end{enumerate}
\end{enumerate}
\finenonce{005840}


\finexercice


\exercice{5841, rouget, 2010/10/16}
\enonce{005841}{** I}
Soit $E = C^2([0,1],\Rr)$. Pour $f$ élément de $E$, on pose $N(f) =\int_{0}^{1}|f(t)|\;dt$, $N'(f) = |f(0)| +\int_{0}^{1}|f'(t)|\;dt$ et

$N''(f) =|f(0)|+|f'(0)|+\int_{0}^{1}|f''(t)|\;dt$. Montrer que $N$, $N'$ et $N''$ sont des normes et les comparer.
\finenonce{005841}


\finexercice


\exercice{5847, rouget, 2010/10/16}
\enonce{005847}{** I Distance d'un point à une partie}

Soit $A$ une partie non vide d'un espace vectoriel normé $(E,\|\;\|)$.

Pour $x\in E$, on pose $d_A(x) =d(x,A)$ où $d(x,A) =\text{Inf}\left\{\|x-a\|,\;a\in A\right\}$.

\begin{enumerate}
 \item  Justifier l'existence de $d_A(x)$ pour chaque $x$ de $E$.

\item 
  \begin{enumerate}
  \item Montrer que si $A$ est fermée, $\forall x\in E$, $d_A(x) = 0\Leftrightarrow x\in A$.
  \item Montrer que si $A$ est fermée et $E$ est de dimension finie, $\forall x\in E$, $\exists a\in A/\;d_A(x)=\|x-a\|$.
  \end{enumerate}

\item  Si $A$ est quelconque, comparer $d_A(x)$ et $d_{\overline{A}}(x)$.

\item  Montrer $d_A$ est continue sur $E$.

\item  A chaque partie fermée non vide $A$, on associe l'application $d_A$ définie ci-dessus. Montrer que l'application $A\mapsto d_A$  est injective.

\item   Dans l'espace des applications continues sur $[0,1]$ à valeurs dans $\Rr$ muni de la norme de la convergence uniforme, on considère $A =\left\{f\in E/\;f(0)=0\;\text{et}\;\int_{0}^{1}f(t)\;dt\geqslant 1\right\}$. Calculer $d_A(0)$.
\end{enumerate}
\finenonce{005847}


\finexercice

\section{229.07 Connexité}


\exercice{4841, quercia, 2010/03/16}
\enonce{004841}{Fronti{\`e}re connexe}
Soit $E$ un espace vectoriel norm{\'e} et $A \subset E$ ferm{\'e}.
Montrer que si $\text{Fr}(A)$ est connexe, alors $A$ est connexe.


\finenonce{004841}



\finexercice


\exercice{4842, quercia, 2010/03/16}
\enonce{004842}{$\mathbb{U}$ et $\R$ ne sont pas hom{\'e}omorphes}
Soit $\mathbb{U}$ le cercle unit{\'e} de $\C$ et $f : \mathbb{U} \to \R$ continue.
Montrer que $f$ n'est pas injective.
\finenonce{004842}



\finexercice


\exercice{4843, quercia, 2010/03/16}
\enonce{004843}{$u_{n+1} - u_n \to 0$}
Soit $E$ un evn de dimension finie et
$(u_n)$ une suite born{\'e}e d'{\'e}l{\'e}ments de $E$ telle que $u_{n+1} - u_n \xrightarrow[n\to\infty]{} 0$.
Montrer que l'ensemble des valeurs d'adh{\'e}rence de la suite est connexe.



\finenonce{004843}



\finexercice


\exercice{4844, quercia, 2010/03/16}
\enonce{004844}{Compl{\'e}mentaire d'une partie {\'e}toil{\'e}e}
Soit $\Omega$ une partie born{\'e}e {\'e}toil{\'e}e d'un evn r{\'e}el de dimension sup{\'e}rieure
ou {\'e}gale {\`a} 2. Montrer que le com\-pl{\'e}\-men\-taire de $\Omega$ est connexe.

\finenonce{004844}



\finexercice


\exercice{4845, quercia, 2010/03/16}
\enonce{004845}{Compl{\'e}mentaire d'un sev}
Soit $E$ un $\R$-evn de dimension finie et $F$ un sev propre de $E$. Montrer que $E\setminus F$
est connexe si et seulement si codim$(F) \ge 2$.
Que peut-on dire dans un $\C$-ev ?
\finenonce{004845}



\finexercice


\exercice{7523, mourougane, 2021/08/10}
\enonce{007523}{Connexité}
Le but de cet exercice est de montrer qu'une partie $P$ ouverte de $\Cc$ est connexe
si et seulement si pour tout couple $(a,b)$ de points de $P$ il existe une ligne polygonale incluse dans $P$
qui joint $a$ et $b$.
\begin{enumerate}
 \item On suppose d'abord que pour tout couple $(a,b)$ de points de $P$ il existe une ligne polygonale incluse dans $P$
qui joint $a$ et $b$. Supposons que $P$ n'est pas connexe. Montrer qu'il existe deux ouverts disjoints $A$ et $B$,
un point $a$ de $A$, un point $b$ de $B$ tel que le segment $[ab]$ soit inclus dans~$P$. 
 \item Noter \begin{eqnarray*}
       \alpha&:=&\{t\in[0,1]\ \ /\ \ ta+(1-t)b\in A\}\\
       \beta&:=&\{t\in[0,1]\ \ /\ \ ta+(1-t)b\in B\}\\
       \end{eqnarray*}
Montrer que $\alpha$ et $\beta$ sont non vides ouverts et disjoints
et utiliser la connexité de $[0,1]$ pour obtenir une contradiction.
 \item On suppose réciproquement que $P$ est connexe et on fixe un point $o$ de $P$.
 On considère $$A:=\{b\in P \text{ il y a une ligne polygonale incluse dans } P \text{ qui joint } o \text{ et } b.\}.$$
 Montrer que $A$ est ouvert dans $P$, en montrant que pour tout point $b$ de $A$, il existe une boule $B(b,\epsilon)$ incluse dans $A$.
 \item Avec les notations précédentes, montrer que $P-A$ est ouvert dans $P$, en montrant que pour tout point $z$ de $P-A$
 et toute boule $B(z,\epsilon)$ incluse dans $P$, $B(z,\epsilon)$ est incluse dans $P-A$.
 \item Conclure.
\end{enumerate}
\finenonce{007523}
\finexercice

\section{229.08 Espaces complets}


\exercice{4846, quercia, 2010/03/16}
\enonce{004846}{Norme pour les fonctions lipschitziennes}
Soit $E = \{ \text{fonctions lipchitziennes } f:\R \to \R \}$.
Pour $f \in E$, on pose
$\|f\| = |f(0)| + \sup\limits_{x\ne y} \left|\frac{f(x)-f(y)}{x-y}\right|$.

Montrer que $E$ est complet.

\finenonce{004846}



\finexercice


\exercice{4847, quercia, 2010/03/16}
\enonce{004847}{Image d'une intersection de ferm{\'e}s}
Soit $E$ un espace vectoriel norm{\'e} complet, $F$ un espace vectoriel norm{\'e} quelconque,
$f : E \to F$ une application continue
et $(E_n)$ une suite d{\'e}croissante de ferm{\'e}s de $E$ dont le diam{\`e}tre tend
vers~0.

Montrer que $f(\mathop{\cap}\limits_n E_n) = \bigcap_n f(E_n)$.
\finenonce{004847}



\finexercice


\exercice{4848, quercia, 2010/03/16}
\enonce{004848}{Intersection de boules}
Soit $E$ un evn complet et $(B_n(a_n,r_n))$ une suite d{\'e}croissante de boules
ferm{\'e}es dont le rayon ne tend pas vers 0.
Montrer que $\bigcap_n B_n$ est une boule ferm{\'e}e.

\finenonce{004848}



\finexercice


\exercice{4849, quercia, 2010/03/16}
\enonce{004849}{Intersection vide}
Soit $E = {\cal B}(\N,\R) = \{ \text{suites } u = (u_n) \text{ born{\'e}es}\}$.
On munit $E$ de la norme : $\|u\| = \sup\{|u_n|,\ n\in\N\}$.

\begin{enumerate}
  \item Montrer que $E$ est complet.
  \item Soit $F_k = \{ u\in E\text{ tq } \|u\| = 1\text{ et }u_0=\dots=u_k=0\}$.
    V{\'e}rifier que les $F_k$ forment une suite de ferm{\'e}s born{\'e}s
    embo{\^\i}t{\'e}s dont l'intersection est vide.

\end{enumerate}
\finenonce{004849}



\finexercice


\exercice{4850, quercia, 2010/03/16}
\enonce{004850}{Th{\'e}or{\`e}me de Baire}
Soit $E$ un espace vectoriel norm{\'e} complet et $(F_n)$
une suite de ferm{\'e}s de $E$ d'int{\'e}rieurs vides.
On pose $\smash{F = \bigcup_n F_n}$. Montrer que $\mathring F = \varnothing$.

\finenonce{004850}



\finexercice


\exercice{4851, quercia, 2010/03/16}
\enonce{004851}{$f\circ f$ est contractante}
Soit $E$ un espace vectoriel norm{\'e} complet et $f : E \to E$ telle que $f\circ f$
est contractante. Montrer que $f$ admet un unique point fixe.

\finenonce{004851}



\finexercice


\exercice{4852, quercia, 2010/03/16}
\enonce{004852}{Centrale MP 2001}

Montrer qu'un plan euclidien n'est pas r{\'e}union de cercles disjoints non
r{\'e}duits {\`a} un point.
\finenonce{004852}



\finexercice

\section{229.09 Fonctions vectorielles}


\exercice{4853, quercia, 2010/03/16}
\enonce{004853}{Centre de gravit{\'e} d'une courbe param{\'e}tr{\'e}e}

Soit $\varphi : {[a,b]} \to {\R^2}, t \mapsto {M_t}$ une courbe param{\'e}tr{\'e}e de classe
$\mathcal{C}^1$ de longueur non nulle.
Le centre de gravit{\'e} de la courbe est le point $G$ tel que
$ \int_{t=a}^b \vec{GM_t}\|\vec M\,'(t)\|\,dt = \vec 0$.

\begin{enumerate}
  \item Montrer l'existence et l'unicit{\'e} de $G$.
  \item D{\'e}terminer le centre de gravit{\'e} d'un demi-cercle.
    (On admet que $G$ est ind{\'e}pendant du param{\'e}trage)
    
  \item Montrer que $G$ appartient {\`a} l'enveloppe convexe de la courbe.
    
  \item Montrer que si la courbe admet un axe de sym{\'e}trie, $\Delta$, alors $G \in \Delta$.
    (Si $\sigma$ est la sym{\'e}trie associ{\'e}e, consid{\'e}rer la courbe d{\'e}crite par
    $N_t = \sigma(M_t)$)
    
  \item Soit $\Phi : {\R^2} \to {\R^2}$ une isom{\'e}trie affine.
    Montrer que si $G$ est le centre de gravit{\'e} de $\mathcal{C}$, alors $\Phi(G)$ est le centre
    de gravit{\'e} de $\Phi(\mathcal{C})$.
\end{enumerate}
\finenonce{004853}



\finexercice


\exercice{4854, quercia, 2010/03/16}
\enonce{004854}{D{\'e}riv{\'e}e d'une base orthonorm{\'e}e}

Soient ${\vec e_1,\vec e_2,\vec e_3} : {I\subset \R} \to {\R^3}$
de classe $\mathcal{C}^1$ telles que pour tout $t\in I$,
${\cal B}_t = (\vec e_1(t),\vec e_2(t), \vec e_3(t))$ est une base orthonorm{\'e}e
de $\R^3$ (base orthonorm{\'e}e mobile).

\begin{enumerate}
  \item Soit $M_t$ la matrice dans ${\cal B}_t$ des vecteurs d{\'e}riv{\'e}s
    $\vec e_1\,'(t), \vec e_2\,'(t), \vec e_3\,'(t)$.
    Montrer que $M_t$ est antisym{\'e}trique.
  \item En d{\'e}duire qu'il existe un vecteur $\vec \Omega(t)$ tel que
    $\vec e_i\,'(t) = \vec \Omega(t)\wedge\vec e_i(t)$, $i=1,2,3$.
  \item Si $\vec e_1,\vec e_2,\vec e_3$ sont de classe $\mathcal{C}^2$, montrer que $\vec \Omega$
    est de classe $\mathcal{C}^1$ et calculer $\vec e_i\,''$ en fonction de
    $\vec \Omega$, $\vec \Omega\,'$ et $\vec e_i$.
    
\end{enumerate}
\finenonce{004854}



\finexercice


\exercice{4855, quercia, 2010/03/16}
\enonce{004855}{$f'$ est colin{\'e}aire {\`a} $f$}

Soit ${\vec f} : {I\subset \R} \to {\R^3}$ une fonction de classe $\mathcal{C}^1$
telle que :
$$\forall\ t\in I,\ \vec f(t) \ne \vec 0 \text{ et la famille }
  (\vec f(t),\vec f\,'(t)) \text{ est li{\'e}e}.$$
On pose $\vec g(t) = \frac{\vec f(t)}{\|\vec f(t)\|}$.

\begin{enumerate}
  \item Montrer que $g$ est de classe $\mathcal{C}^1$ et que $\vec g\,'(t)$ est {\`a} la fois orthogonal
    et colin{\'e}aire {\`a} $\vec g(t)$.
  \item En d{\'e}duire que $\vec f(t)$ garde une direction constante.
  \item Chercher un contre-exemple lorsqu'on retire la propri{\'e}t{\'e} :
    $\forall\ t\in I,\ \vec f(t) \ne \vec 0$.

\end{enumerate}
\finenonce{004855}



\finexercice


\exercice{4856, quercia, 2010/03/16}
\enonce{004856}{$f''$ est colin{\'e}aire {\`a} $f$}

Soit ${\vec f} : {I\subset \R} \to {\R^3}$ une fonction de classe $\mathcal{C}^2$
telle que pour tout $t\in I$, $\vec f\,''(t)$ est colin{\'e}aire {\`a} $\vec f(t)$.
({\it mouvement {\`a} acc{\'e}l{\'e}ration centrale})
On note $\vec \sigma(t) = \vec f(t)\wedge\vec f\,'(t)$.

\begin{enumerate}
  \item Montrer que $\vec \sigma(t)$ est un vecteur constant.
  \item S'il existe $t_0\in I$ tel que $(\vec f(t_0),\vec f\,'(t_0))$ est libre,
    montrer que $\vec f(I)$ est inclus dans un plan.
\end{enumerate}
\finenonce{004856}



\finexercice


\exercice{4857, quercia, 2010/03/16}
\enonce{004857}{$f$ et $g$ colin{\'e}aires}

Soient ${f,g} : {I\subset \R} \to E$ deux fonctions vectorielles de classe $\mathcal{C}^1$.

\begin{enumerate}
  \item On suppose : $\forall\ t \in I,\ f(t)$ et $g(t)$ sont colin{\'e}aires.
    Est-ce que $f'(t)$ et $g'(t)$ sont colin{\'e}aires ?
    

  \item On suppose : $\forall\ t \in I,\ f'(t)$ et $g'(t)$ sont colin{\'e}aires.
    Existe-t-il $\vec c \in E$ tel que $f-\vec c$ et $g$ soient colin{\'e}aires~?
    
\end{enumerate}
\finenonce{004857}



\finexercice


\exercice{4858, quercia, 2010/03/16}
\enonce{004858}{$\R^2\setminus$ une droite n'est pas connexe}

Soit $f : {I\subset \R}  \to {\R^2}$ une fonction continue,
$D$ une droite de $\R^2$ et $P^+$, $P^-$ les demi-plans d{\'e}limit{\'e}s par $D$.
Montrer que s'il existe $a,b\in I$ tels que $f(a)\in P^+$ et $f(b)\in P^-$, alors
il existe $c$ compris entre $a$ et $b$ tel que $f(c)\in D$.
G{\'e}n{\'e}raliser en dimension $n$.
\finenonce{004858}



\finexercice

\section{229.10 Application linéaire continue, norme matricielle}


\exercice{5854, rouget, 2010/10/16}
\enonce{005854}{*}
On munit $E=\Rr[X]$ de la norme $\|\;\|_\infty$ définie par : $\forall P\in E$, $\|P\|_\infty=\text{Sup}\left\{\left| \frac{P^{(n)}(0)}{n!}\right|,\;n\in\Nn\right\}$.

\begin{enumerate}
 \item  Vérifier brièvement que $\|\;\|_\infty$ est une norme sur $E$.

\item  Soit $f$ l'endomorphisme de $E$ défini par $\forall P\in E$, $f(P)=XP$. Démontrer que l'application $f$ est continue sur $(E,\|\;\|_\infty)$ et déterminer $|||f|||$.
\end{enumerate}
\finenonce{005854}


\finexercice


\exercice{5855, rouget, 2010/10/16}
\enonce{005855}{**}
On munit $E =\ell^\infty(\Cc)$ le $\Cc$-espace vectoriel des suites bornées de la norme $\|u\|_\infty=\underset{n\in\Nn}{\text{sup}}\;|u_n|$.

On considère les endomorphismes $\Delta$ et $C$ de $\ell^\infty(\Cc)$ définis par :

\begin{center}
$\forall u\in E$, $\Delta(u) = v$ où $\forall n\in\Nn$, $v_n = u_{n+1}-u_n$ et $\forall u\in E$, $C(u)=w$ où $\forall n\in\Nn$, $w_n= \frac{1}{n+1}\sum_{k=0}^{n}u_k$.
\end{center}

Montrer que $\Delta$ et $C$ sont continus sur $(E,\|\;\|_\infty)$ et calculer leur norme.
\finenonce{005855}


\finexercice


\exercice{5856, rouget, 2010/10/16}
\enonce{005856}{*** I}
On munit $E=C^0([0,1],\Rr)$ de la norme $1$ définie par $\forall f\in E$, $\|f\|_1=\int_{0}^{1}|f(t)|\;dt$.

On pose $\begin{array}[t]{cccl}
T~:&E&\rightarrow&\rule{1.5mm}{0mm}E\\
 &f&\mapsto&\begin{array}[t]{cccc}Tf~:&[0,1]&\rightarrow&\Rr\\
  &x&\mapsto&\int_{0}^{x}f(t)\;dt
  \end{array}
 \end{array}$ et on admet que $T$ est un endomorphisme de $E$.
 
 
 \begin{enumerate}
 \item  Démontrer que $T$ est continu sur $(E,\|\;\|_1)$ et déterminer $|||T|||$.
 
 
 \item  Vérifier que la borne supérieure n'est pas atteinte.
\end{enumerate}
\finenonce{005856}


\finexercice


\exercice{5857, rouget, 2010/10/16}
\enonce{005857}{**}
On munit $E=\mathcal{M}_n(\Rr)$ de la norme $N$ définie par $\forall A\in E$, $N(A)=\underset{1\leqslant i\leqslant n}{\text{Sup}}\left\{\sum_{j=1}^{n}|a_{i,j}|\right\}$ (on admet que $N$ est une norme sur $E$).

Soit $f$ l'application de $E$ dans $\Rr$ définie par $\forall A\in E$, $f(A)=\text{Tr}(A)$. Démontrer que l'application $f$ est continue sur $(E,N)$ et déterminer $|||f|||$.
\finenonce{005857}


\finexercice


\exercice{5858, rouget, 2010/10/16}
\enonce{005858}{***}
\label{ex:rou5}
Déterminer $s=\text{Sup}\left\{ \frac{\|AB\|}{\|A\|\|B\|},\;(A,B)\in(\mathcal{M}_n(\Cc)\setminus\{0\})^2\right\}$ quand $\|\;\|$ est 
\begin{enumerate}
 \item  $\|\;\|_1$, 
 \item  $\|\;\|_2$, 
 \item  $\|\;\|_\infty$.
\end{enumerate}
\finenonce{005858}


\finexercice


\exercice{5859, rouget, 2010/10/16}
\enonce{005859}{*}
Une norme sur $\mathcal{M}_n(\Rr)$ ($n\geqslant2$), est-elle nécessairement une \og norme trois barres \fg~?
\finenonce{005859}


\finexercice


\exercice{5860, rouget, 2010/10/16}
\enonce{005860}{**}
Soit $N$ une norme sur $\mathcal{M}_n(\Rr)$.

Montrer qu'il existe $k>0$ tel que $\forall(A,B)\in(\mathcal{M}_n(\Rr))^2$, $N(AB)\leqslant k(A)N(B)$.
\finenonce{005860}


\finexercice


\exercice{5861, rouget, 2010/10/16}
\enonce{005861}{**}
Existe-t-il une norme $N$ sur $\mathcal{M}_n(\Rr)$ ($n\geqslant2$) telle que $\forall(A,B)\in(\mathcal{M}_n(\Rr))^2$, $N(AB)=N(A)N(B)$.
\finenonce{005861}


\finexercice


\exercice{5862, rouget, 2010/10/16}
\enonce{005862}{***}
On pose $\forall X=(x_i)_{1\leqslant i\leqslant n}\mathcal{M}_{n,1}(\Rr)$, $\|X\|_1=\sum_{i=1}^{n}|x_i|$ et $\|X\|_\infty=\underset{1\leqslant i\leqslant n}{\text{Max}}|x_i|$.

Déterminer les normes sur $\mathcal{M}_n(\Rr)$ respectivement associées aux normes $\|\;\|_1$ et $\|\;\|_\infty$ de $\mathcal{M}_{n,1}(\Rr)$. On notera $|||\;|||_1$ et $|||\;|||_\infty$ ces normes.
\finenonce{005862}


\finexercice


\exercice{5863, rouget, 2010/10/16}
\enonce{005863}{**I}
Pour $X=(x_i)_{1\leqslant i\leqslant n}\in\mathcal{M}_{n,1}(\Rr)$, on pose $\|X\|_2=\sqrt{\sum_{i=1}^{n}x_i^2}$. Pour $A\in\mathcal{S}_n(\Rr)$, on note $\rho(A)$ le rayon spectral de $A$ c'est-à-dire $\rho(A)=\text{Max}\{|\lambda|,\;\lambda\in\text{Sp}(A)\}$.

Montrer que $\forall A\in\mathcal{S}_n(\Rr)$, $|||A|||_2=\rho(A)$ où $|||A|||_2=\text{Sup}\left\{ \frac{\|AX\|_2}{\|X\|_2},\;X\in\mathcal{M}_{n,1}(\Rr)\setminus\{0\}\right\}$.

\finenonce{005863}


\finexercice

\section{229.99 Autre}


\exercice{1740, drutu, 2003/10/01}

\enonce{001740}{}
Soit $x = (x_1, \cdots x_n) \in \R^n$. On pose 
$$ \|x\|_{1} = \sum_{i=1}^n |x_i| \quad ; \quad  \|x\|_{2} = \left(\sum_{i=1}^n |x_i|^2\right)^{1/2} $$
et
$$ \|x\|_{\infty} = \sup \{ |x_i| : 1 \leq i \leq n\} . $$
\begin{enumerate}
\item D\'emontrer que $\|\cdot\|_{1}$ est une norme sur $\R^n$.

\item D\'emontrer que  
$$ \|x\|_{\infty} \leq \|x\|_{2} \leq \|x\|_{1} \leq n \|x\|_{\infty}$$
et 
$$ \|x\|_{2} \leq \sqrt{n}\|x\|_{\infty},$$
pour tout $x \in \R^n$. Discuter le cas $n = 1$. 

\item Repr\'esenter dans $\R^2$ la boule unit\'e ferm\'ee 
$$B_{\|\cdot\|} = \{ x = (x_1,x_2) \in \R^2 \; ;\;  \|x\| \leq 1\} $$
pour chacune des normes $\|\cdot\|_{1}, \|\cdot\|_{2}$ et $\|\cdot\|_{\infty}$.
\end{enumerate}
\finenonce{001740}



\finexercice



\exercice{1745, maillot, 2001/09/01}

\enonce{001745}{}
\begin{enumerate}
\item
 Dans $\R^2$ ou $\R^3$ euclidien muni d'une b.o.n., représenter les
ensembles suivants :
  \begin{itemize}
  \item
$\mathcal{A}=\{(x,y)\in\R^2 \ | \  x^2-y^2>1 \mbox{ et } x^2+y^2<4\}$
  \item
$\mathcal{B}=\{(x,y)\in\R^2 \ | \  (x-1)^2-y^2>1 \mbox{ et } x^2+\frac{y^2}{4}<4\}$
  \item
$\mathcal{C}=\{(x,y,z)\in\R^3 \ | \  1<x+y+z<3 \mbox{ et } x>0 \mbox{ et }
y>0 \mbox{ et } z>0\}$
  \item
$\mathcal{D}=\left\{
\begin{array}{cc}
(x,y,z)\in\R^3&\left|
\begin{array}{cc}
&x+y+z<1\\
\mbox{et }&x-y+z<1\\
\mbox{et }&-x-y+z<1
\end{array}
\right.
\end{array}
\right\}$
\item
$\mathcal{E}=\{(x,y,z)\in\R^3 \ | \  x^2+y^2-z^2<0 \mbox{ et } 2<z<4\}$
\item
$\mathcal{F}=\{(x,y,z)\in\R^3\ | \  x^2+y^2+z^2<1 \mbox{ et } x^2+y^2<z^2
  \mbox{ et } z>0\}$
\item
$\mathcal{G}=\{(x,y,z)\in\R^3\ | \  x^2+y^2=4 \mbox{ et } z=x-1\}$.
\end{itemize}
\item Déterminer les projections de $\mathcal{E}$ et $\mathcal{G}$ sur le
plan $(xOy)$.
\end{enumerate}
\finenonce{001745}



\finexercice



\exercice{1746, maillot, 2001/09/01}

\enonce{001746}{Images directes et réciproques}
\begin{enumerate}
\item Soit $f$ l'application affine par morceaux, de $\R$ dans $\R$, définie par :
$$
f(x)=\left\{
\begin{array}{ccc}
0&\mbox{ si }&x\leq -2\\
1+x&\mbox{ si }&-2<x<0\\
x&\mbox{ si }&0\leq x\leq 1\\
1&\mbox{ si }&x>1.
\end{array}
\right.
$$
Soient $A=[-1,0[$ et $B=[0,2[$.
Déterminer $f(A)$, $f^{-1}(B)$, $f(\R\setminus A)$, $f^{-1}(f(A))$,
$f(f^{-1}(B))$, $f(A\cap B)$, et $f(A)\cap f(B)$.

\item Soient deux ensembles $E$ et $F$, et $f:E\rightarrow F$ une
application. Comparer les ensembles $f(A\cap B)$ et $f(A)\cap f(B)$,
$f^{-1}(f(A))$ et $A$, $f(f^{-1}(B))$ et $B$, $f(E\setminus A)$ et
$F\setminus f(A)$.
\end{enumerate}
\finenonce{001746}



\finexercice



\exercice{1747, maillot, 2001/09/01}

\enonce{001747}{}
Soit l'application
$\left(\begin{array}{cccl}
G:&\R^2&\longrightarrow&\R^2\\
&(u,v)&\longmapsto&(\frac{u}{u+v},\frac{\sqrt{v(v+2u)}}{u+v})
\end{array}\right)$. On note
$\mathcal{D}$ l'ensemble de définition de $G$. Déterminer $G(\mathcal{D})$.

\finenonce{001747}



\finexercice



\exercice{1748, maillot, 2001/09/01}

\enonce{001748}{}
Soient les applications $f$ et $g$ de $\R^2$ dans $\R^2$ définies par
:
$$
f(x,y)=(\frac{x+y}{2},\frac{\sqrt{3}}{2}y) \mbox { et }
g(x,y)=(2x,\frac{y}{\sqrt{2}}).
$$
Soient les ensembles
$$
\mathcal{D}_1=\{(x,y)\in\R^2\ | \ \frac{x^2}{4}+y^2+\frac{xy}{2}=1\},
$$
$$
\mbox{et }
\mathcal{D}_2=\{(x,y)\in\R^2\ | \ \frac{x^2}{4}+2y^2=1\}.
$$
Déterminer $f(\mathcal{D}_1)$ et $g^{-1}(\mathcal{D}_2)$.

\finenonce{001748}



\finexercice



\exercice{1749, maillot, 2001/09/01}

\enonce{001749}{}
Simplifier l'écriture des ensembles suivants :
$$
I=\bigcup_{n>1}[\frac{1}{n},1-\frac{1}{n}] \mbox{ et }
J=\bigcap_{i>0,j>0}]-\frac{1}{i},1+\frac{1}{j}[.
$$

\finenonce{001749}



\finexercice



\exercice{1765, maillot, 2001/09/01}

\enonce{001765}{}
Soit $f:\R^d\rightarrow \R$ une fonction continue telle que $\lim_{x\rightarrow -\infty}
f(x)=\lim_{x\rightarrow +\infty} f(x) = +\infty$. Montrer que $f$ admet un
minimum.
\finenonce{001765}



\finexercice



\exercice{1768, maillot, 2001/09/01}

\enonce{001768}{}
Soit $(E,\|\cdot\|)$ un espace vectoriel norm\'e et $(x_n)_{n\in\N}$ une
suite d'\'el\'ements de $E$. On suppose que $(x_n)$ est de Cauchy.
Montrer qu'elle converge si et seulement si elle admet une sous-suite
convergente.
\finenonce{001768}



\finexercice



\exercice{1775, gourio, 2001/09/01}

\enonce{001775}{}
Soit $ C$ une partie convexe de $ {\Rr}^{2}, $ montrer que $ \overline{C}$
est aussi convexe.
\finenonce{001775}



\finexercice



\exercice{5842, rouget, 2010/10/16}
\enonce{005842}{*** I Topologie dans $\mathcal{M}_n(\Kk)$}

\begin{enumerate}
 \item   Montrer que $GL_n(\Rr)$ est un ouvert de $\mathcal{M}_n(\Rr)$, dense dans $\mathcal{M}_n(\Rr)$.

\item  Montrer que $\mathcal{M}_n(\Rr)\setminus GL_n(\Rr)$ est fermé mais non compact (pour $n\geqslant2$).

\item  Montrer que $O_n(\Rr)$ est compact. $O_n(\Rr)$ est-il convexe ?

\item  Montrer que $S_n(\Rr)$ est fermé.

\item  Soit $p\in\llbracket0,n\rrbracket$. Montrer que l'ensemble des matrices de rang inférieur ou égal à $p$ est un fermé de $\mathcal{M}_n(\Rr)$.

\item   Montrer que l'ensemble des matrices diagonalisables dans $\mathcal{M}_n(\Cc)$ est dense dans $\mathcal{M}_n(\Cc)$. Peut-on remplacer $\mathcal{M}_n(\Cc)$ par $\mathcal{M}_n(\Rr)$ ?

\item  Propriétés topologiques de l'ensemble des triplets de réels $(a,b,c)$ tels que la forme quadratique $(x,y)\mapsto ax^2+2bxy+cy^2$ soit définie positive ?

\item   Montrer que l'ensemble des matrices stochastiques (matrices $(a_{i,j})_{1\leqslant i,j\leqslant n}\in\mathcal{M}_n(\Rr)$ telles que $\forall (i,j)\in\llbracket1,n\rrbracket^2$, $a_{i,j}\geqslant0$ et $\forall i\in\llbracket1,n\rrbracket$, $\sum_{j=1}^{n}a_{i,j}= 1$) est un compact convexe de $\mathcal{M}_n(\Rr)$.

\item  Montrer que l'ensemble des matrices diagonalisables de $\mathcal{M}_n(\Rr)$ est connexe par arcs.
\end{enumerate}
\finenonce{005842}


\finexercice

\section{240.00 Géométrie affine dans le plan et dans l'espace}


\exercice{1956, liousse, 2003/10/01}
\video{zrMBUrw-HzU}
\enonce{001956}{}
Soit $P$ un  plan muni d'un repère  $\mathcal{R}(O,\vec{i},\vec{j})$, les points et les vecteurs sont exprimés par leurs
coordonnées dans $\mathcal{R}$.
\begin{enumerate}
\item
 Donner un vecteur directeur, la pente 
 une équation paramétrique et une équation cartésienne des droites $(AB)$ suivantes :

 \begin{enumerate}
   \item $A(2,3)$ et $B(-1,4)$    
  \item  $A(-7,-2)$ et $B(-2,-5)$  
  \item  $A(3,3)$ et $B(3,6)$
  \end{enumerate} 

 \item  Donner des équations paramétriques et cartésiennes des droites passant par $A$ et 
dirigées par $\vec{v}$ avec : 
  \begin{enumerate}
    \item  $A(2,1)$ et $\vec{v}(-3,-1)$ 
    \item  $A(0,1)$ et $\vec{v}(1,2)$
    \item  $A(-1,1)$ et $\vec{v}(1,0)$
  \end{enumerate} 

\item   Donner des équations paramétriques et cartésiennes
 des droites définies comme suit  :
  \begin{enumerate}
    \item  passant par le point $(0,4)$ et  de pente $3$,
     \item  passant par le point $(2,-3)$ et  parallèle à l'axe des $x$,
    \item  passant par le point $(-2,5)$ et  parallèle à la droite $D : 8x+4y=3$.
\end{enumerate}
\end{enumerate}
\finenonce{001956}



\finexercice

\exercice{1957, liousse, 2003/10/01}

\enonce{001957}{}
\begin{enumerate}
\item
 Les trois points $A$, $B$ et $C$ de $P$ sont-ils align\'es ?
 Si oui donner une \'equation cart\'esienne de la droite qui les
contient.

 
\begin{enumerate}
   \item $A(-3,3)$, $B(5,2)$ et $C(2,1)$,
   \item $A(1,1)$, $B(-2,2)$ et $C(2,1)$,
   \item $A(4,-3)$, $B(0,-1)$ et $C(2,-2)$,
   \item $A(2,-1)$, $B(1,-2)$ et $C(-3,4)$.

 \end{enumerate}


 \item  Dans les cas suivant, donner un vecteur directeur de $D$ et d\'eterminer 
si le point $C$ appartient  ou non \`a $D$ 
\begin{enumerate}
\item $(D) : 3x + 5y +1= 0$, \hfil $C(3, -2)$.

\item $(D) : \left\{ \begin{array}{l}
                             x = 3 + t \\
                             y = 2 - t
                           \end{array} \right.$, \hfil $C(5,3)$.
 
                            
\end{enumerate}
\end{enumerate}
\finenonce{001957}



\finexercice



\exercice{1958, liousse, 2003/10/01}

\enonce{001958}{}
Dans l'exercice suivant, 
on consid\`ere des couples de deux droites $D_1$ et $D_2$ : on doit 
d\'eterminer si
elles sont s\'ecantes, parall\`eles ou confondues. Si elles sont s\'ecantes, 
on d\'eterminera les coordonn\'ees
du point d'intersection, et si elles 
sont parall\`eles ou confondues on d\'eterminera un vecteur directeur.
\begin{enumerate}

\item $(D_1) : 3x + 5y - 2 = 0$ et $(D_2) : x - 2y + 3 = 0$
\item $(D_1) : 2x - 4y + 1 = 0$ et $(D_2) : -5x + 10y + 3 = 0$
\item $(D_1) : \left\{ \begin{array}{l}
                             x = 3 + 4t \\
                             y = 2 - t
                           \end{array} \right.$
 et $(D_2) : \left\{ \begin{array}{l}
                       x = 5 - s \\
                       y = 2 + 3s 
                     \end{array} \right.$
\item $(D_1) : \left\{ \begin{array}{l}
                             x = 1 + 2t \\
                             y = 2 - 3t
                           \end{array} \right.$
 et $(D_2) : \left\{ \begin{array}{l}
                       x = 3 - 4s \\
                       y = -1 + 6s 
                     \end{array} \right.$
\item $(D_1) : x - 2y + 3 = 0$ et $D_2 :\left\{ \begin{array}{l}
                                                     x = 2 + t  \\
                                                     y = 3 - 2t 
                                                    \end{array} \right.$ 
\item $(D_1) : 3x - 2y + 1 = 0$ et $(D_2) : \left\{ \begin{array}{l}
                                                          x = 1 - 4t \\
                                                          y = 2 - 6t 
                                                        \end{array} \right.$
\end{enumerate}
\finenonce{001958}



\finexercice



\exercice{1959, liousse, 2003/10/01}

\enonce{001959}{}
 On consid\`ere les deux droites du 
plan $D : 2x-3y+4 = 0$ et $D' : x+3y+1 = 0$. On consid\`ere le 
point $A$, intersection des deux droites et le point $B$ de coordonn\'ees $(3,8)$.
Donner une \'equation de (AB).
\finenonce{001959}



\finexercice



\exercice{1960, liousse, 2003/10/01}
\video{tdFJc9r8NSM}
\enonce{001960}{}
On considère le triangle 
$ABC$ dont les c\^otés ont pour équations
$(AB) : x+2y=3 , (AC) : x+y=2 , (BC) : 2x+3y=4$.
\begin{enumerate}
\item Donner les coordonnées des points $A, B, C$.
\item Donner les coordonnées des milieux $A', B', C'$ des segments $[BC]$, $[AC]$ et $[AB]$ respectivement.
\item Donner une équation de chaque médiane et vérifier qu'elles sont concourantes.
\end{enumerate}
\finenonce{001960}



\finexercice


\exercice{1961, liousse, 2003/10/01}

\enonce{001961}{M\'edianes}
 
On consid\`ere dans $P$ trois points $A$, $B$ et $C$.
 
\begin{enumerate}
\item 
 D\'eterminer dans le rep\`ere $(A,\vec{AB},\vec{AC})$ des \'equations pour les 
m\'edianes du triangle ABC.

\item En d\'eduire que les m\'edianes d'un triangle sont concourantes.
\end{enumerate}
\finenonce{001961}



\finexercice



\exercice{1962, liousse, 2003/10/01}

\enonce{001962}{Th\'eor\`eme de Menela\"us}
Dans le triangle $ABC$, on consid\`ere trois points $P, Q, R$, sur les c\^ot\'es
$(BC), (AC)$ et $(AB)$ respectivement, ces points n'\'etant pas les points $A,B$ ou $C$.
Montrer que $P,Q$ et $R$ sont align\'es si et seulement si
$$ {\overline {PB} \over \overline {PC}} = {\overline {QC}\over \overline {QA}} = 
{\overline {RA} \over \overline {RB}}=1$$
\finenonce{001962}



\finexercice



\exercice{1963, liousse, 2003/10/01}

\enonce{001963}{Th\'eor\`eme de Pappus}
 Soient $(A_1,A_2,A_3)$ et $(B_1,B_2,B_3)$ deux 
syst\`emes de trois points align\'es. Montrer que les points 
$C_1,$ $C_2$ et $C_3,$ intersections des droites $(A_2B_3)$ et $(A_3B_2),$ 
$(A_3B_1)$ et $(A_1B_3),$ $(A_1B_2)$ et $(A_2B_1)$ (que l'on 
suppose exister) sont align\'es.
\finenonce{001963}



\finexercice



\exercice{1964, liousse, 2003/10/01}

\enonce{001964}{Th\'eor\`eme de Ceva}
 Dans le triangle $ABC$, on consid\`ere trois points $P, Q, R$, sur 
les droites $(BC), (AC)$ et $(AB)$ respectivement, ces points n'\'etant pas 
les points $A,B$ ou $C$.
Montrer que les droites $(AP),$ $(BQ)$ et $(CR)$ sont 
concourantes ou parall\`eles si et seulement si
$$ {\overline {PB} \over \overline {PC}}.{\overline {QC}\over \overline {QA}}.
{\overline {RA} \over \overline {RB}}=-1$$
\finenonce{001964}



\finexercice



\exercice{1965, ridde, 1999/11/01}

\enonce{001965}{}
Montrer que l'intersection de deux parties convexes est convexe. Est-ce vrai
pour l'union ?
\finenonce{001965}



\finexercice



\exercice{1966, gourio, 2001/09/01}

\enonce{001966}{}
Soient $C$ et $C^{\prime }$ deux ensembles convexes d'un espace affine,
montrer que
$$D=\left\{ \frac{M+M^{\prime }}{2}|\left( M,M^{\prime }\right)
\in C\times C^{\prime }\right\} $$ est convexe.
\finenonce{001966}



\finexercice



\exercice{1967, gourio, 2001/09/01}

\enonce{001967}{}
On appelle enveloppe convexe $co(A) $ d'une partie non vide $A$ d'un espace
affine $E $  l'intersection des ensembles convexes contenant $A$ ;
c'est le plus petit ensemble convexe contenant $A.$ Montrer que c'est aussi
l'ensemble des barycentres \`{a} coefficients positifs de points de $A.$ Que
sont $co(\{A,B\}), co(\{A,B,C\}) $ ?
\finenonce{001967}



\finexercice



\exercice{1968, gourio, 2001/09/01}

\enonce{001968}{}
Un c\^{o}ne d'un espace vectoriel est une partie $K $ telle que :
$$\forall x\in K,\forall t\geq 0,tx\in K. $$
Montrer qu'un c\^{o}ne est convexe si et seulement si il est stable par
addition.
\finenonce{001968}



\finexercice



\exercice{1969, gourio, 2001/09/01}

\enonce{001969}{}
Trouver les parties $C$ convexes de ${\Rr}^{2}$ telles que le
compl\'{e}mentaire $^{c}C$ soit aussi convexe.
\finenonce{001969}



\finexercice



\exercice{1970, gourio, 2001/09/01}

\enonce{001970}{}
Soit $E  $ un espace affine de dimension $n,$ et $\left(
x_{1},...,x_{n}\right)$ des points de $E. $On consid\`{e}re une
combinaison convexe de points de $A$, sous ensemble de $E:$
$$x=\sum_{i=1}^{m}t_{i}x_{i}  \text{ avec }
\forall i\in \{1,...,m\}:t_{i}\geq 0  \text{ et }\sum_{j=1}^{m}t_{j}=1. $$
Montrer qu'on peut \'{e}crire :
$$x=\sum_{k=1}^{n+1}g_{k}x_{k} \text{ avec }
\forall k\in \{1,...,n+1\}:g_{k}\geq 0  \text{ et  }
\sum_{k=1}^{n+1}g_{k}=1. $$
Ainsi il suffit de $n+1$ points dans un espace de dimension $n$ pour
\'{e}crire une combinaison convexe.
\finenonce{001970}



\finexercice



\exercice{1971, ridde, 1999/11/01}

\enonce{001971}{}
Une bim\'ediane d'un t\'etra\`edre est une droite qui passe par les milieux de deux
arêtes oppos\'ees. Montrer que les trois bim\'edianes sont concourantes.
\finenonce{001971}



\finexercice



\exercice{1972, ridde, 1999/11/01}

\enonce{001972}{}
Soient $A, B, C$ trois points non align\'es d'un plan affine. D\'eterminer l'ensemble
des points ayant mêmes coordonn\'ees dans les rep\`eres $ (A, \overrightarrow{AB},
\overrightarrow{AC})$ et $ (B, \overrightarrow{BA}, \overrightarrow{BC})$.
\finenonce{001972}



\finexercice



\exercice{1973, ridde, 1999/11/01}

\enonce{001973}{}
Soit $R_{1} = (0, e_{1}, e_{2}, e_{3})$ un rep\`ere cart\'esien d'un espace affine.
Soient $O^{'} = (1, 0, 0)$, $e_{1}^{'} = e_{1} + e_{2}$, $e_{2}^{'} = e_{1} - e_{2}$,
$e_{3}^{'} = e_{3}$ et $R_{2} = (0^{'}, e_{1}^{'}, e_{2}^{'}, e_{3}^{'})$.
D\'eterminer les coordonn\'ees d'un point dans $R_{2}$ en fonction de ses coordonn\'ees dans $R_{1}$.
\finenonce{001973}



\finexercice



\exercice{1974, ridde, 1999/11/01}

\enonce{001974}{}
Soient $ (D_{i})_{i = 1 \ldots 4}$ quatre droites du plan affine s\'ecantes deux
\`a deux en six points distincts. Si deux d'entre elles se coupent en $A$ et les
deux autres en $B$, on dit que $[AB]$ est une diagonale. Montrer que les milieux des
trois diagonales sont align\'es (on \'etudiera le probl\`eme analytiquement en choisissant
un bon rep\`ere).
\finenonce{001974}



\finexercice



\exercice{1975, ridde, 1999/11/01}

\enonce{001975}{}
\begin{enumerate}
\item Soient $ (D_{i} : u_{i}x + v_{i}y + h_{i} = 0)_{i = 1 \ldots 3}$ trois droites du plan affine.
 Montrer qu'elles sont parall\`eles ou concourantes ssi $\begin{vmatrix}
u_{1} & v_{1} & h_{1}\\u_{2} & v_{2} & h_{2}\\u_{3} & v_{3} & h_{3}\\
\end{vmatrix} = 0$.
\item Soient $ (D_{1} : x + 2y = 1)$, $ (D_{2} : x + y = 2)$, $ (D_{3} : 2x + y = 3)$,
$ (D_{4} : 3x + 2y = 1)$. D\'eterminer une \'equation de la droite $D$ qui passe par
$D_{1} \cap D_{2}$ et $D_{3} \cap D_{4}$ sans calculer ces points d'intersection.
\end{enumerate}
\finenonce{001975}



\finexercice



\exercice{1976, ridde, 1999/11/01}

\enonce{001976}{}
Soient $A, B, C$ trois points non align\'es d'un plan affine.
\begin{enumerate}
\item Soit $f$ une application
affine telle que $f (A) = A$, $f (B) = B$ et $f (C) = C$. Montrer que $f = id$.
\item
Soient $f$ et $g$ affines telles que $f (A) = g (A)$, $f (B) = g (B)$ et $f (C) = g (C)$.
Que peut-on dire ?
\item
Soit $f$ affine telle que $f (A) = B$, $f (B) = C$ et $f (C) = A$. Que peut-on dire ?
\end{enumerate}
\finenonce{001976}



\finexercice



\exercice{1977, ridde, 1999/11/01}

\enonce{001977}{}
Soit $E$ un espace affine et f une application affine de $E$ dans $E$.
\begin{enumerate}
\item Montrer que $f$ est une translation ssi $\vec {f} = id$.

\item Montrer que si $\vec {f} = \lambda id$ où $\lambda \neq 1$ alors $f$ est
une homoth\'etie (on montrera que $f$ admet un point fixe).

\item On note $\mathcal{T}$ l'ensemble des translations. Montrer que $\mathcal{T}$
est un sous-groupe du groupe affine.

\item  On note $\mathcal{H}$ l'ensemble des homoth\'eties bijectives.
Montrer que $\mathcal{T} \cup \mathcal{H}$ est un sous-groupe du groupe affine.
\end{enumerate}
\finenonce{001977}



\finexercice



\exercice{1978, ridde, 1999/11/01}

\enonce{001978}{}
Soient $f$ et $g$ deux applications affines de $E$ dans $E$ telles que $\vec{f}
 = \vec{g}$. Montrer qu'il existe $u \in \vec{E}$ tel que $f = t_{u}\circ g$ où
 $t_{u}$ est la translation de vecteur $u$.
 Que peut-on dire si de plus il existe $M\in E$ tel que $f (M) = g (M)$ ?
\finenonce{001978}



\finexercice



\exercice{1979, ridde, 1999/11/01}

\enonce{001979}{}
Reconnaître les application affines de $\Rr^{3}$ suivantes :
$$\begin{pmatrix} x \\ y\\ z\\ \end{pmatrix} \mapsto \begin{pmatrix}
-x + 2y-2z-2\\ -3y + 2z + 6\\ -4y + 3z + 6\\ \end{pmatrix} \text{ et }\begin{pmatrix}
x \\ y\\ z\\
\end{pmatrix} \mapsto \begin{pmatrix} \frac{y}{2}-\frac{z}{2} + \frac{2}{3}\\
-x + \frac{3y}{2}-\frac{z}{2} + \frac{2}{3}\\ -x + \frac{y}{2} + \frac{z}{2} + \frac{2}{3}\\
\end{pmatrix}$$
\finenonce{001979}



\finexercice



\exercice{1980, ridde, 1999/11/01}

\enonce{001980}{}
Soit $E$ un espace affine, $f$ une application affine de $E$ dans $E$ et
 $$F = \left\{M\in E /f (M) = M \right\}.$$
 On suppose que $F \neq \emptyset$.
\begin{enumerate}
\item Montrer que $\vec{F} = \ker (\vec{f}-id)$.

\item On suppose que $\vec{f}\circ \vec{f} = \vec{f}$. Soit $s$ la projection
affine sur $F$ parall\`element \`a $\ker (\vec{f})$. Montrer que
$f = s$.

\item Faire la même chose si $\vec{f}\circ \vec{f} = id$.
\end{enumerate}
\finenonce{001980}



\finexercice



\exercice{1981, ridde, 1999/11/01}

\enonce{001981}{}
Soit $E$ un espace affine et $f$ une application affine de $E$ dans $E$.
\begin{enumerate}
\item Montrer que si $f\circ f = f$ alors $f$ est une projection affine.

\item Montrer que si $f\circ f = id$ alors $f$ est une sym\'etrie affine.
\end{enumerate}
\finenonce{001981}



\finexercice



\exercice{1991, liousse, 2003/10/01}

\enonce{001991}{}
  On consid\`ere les droites $D : x+2y=5$ et $D' : 3x- y = 1$ et on note 
$A$ l'intersection des deux droites et $B$ le point de coordonn\'ees $(5,2)$.
\begin{enumerate}
\item Donner une \'equation cart\'esienne de la droite $(AB)$.

\item Donner une \'equation cart\'esienne de la perpendiculaire \`a $D$ 
passant par $B$.

\item Donner une \'equation cart\'esienne de la parall\`ele \`a $D'$ 
passant par $B$.

\item Soit $C$ le point de coordonn\'ees $(2,-7)$). Donner une \'equation
cart\'esienne de la m\'ediatrice $\Delta$ du segment $[B,C]$. $\Delta$ est-elle parall\`ele \`a $D$ ? Et \`a $D'$ ? 
\end{enumerate}
\finenonce{001991}



\finexercice



\exercice{1992, liousse, 2003/10/01}

\enonce{001992}{}
\begin{enumerate}
\item
   On consid\`ere la famille des 
droites $D_{\lambda} : x + \lambda y + 1 =0$, o\`u $\lambda \in \Rr$.
\begin{enumerate}
\item V\'erifier que ces droites passent toutes par un m\^eme point $A$ dont 
on donnera les coordonn\'ees.

\item Parmi toutes ces droites, y en a-t-il une qui est verticale ? Si oui donner une \'equation de
cette droite.

\item Parmi toutes ces droites, y en a-t-il une qui est horizontale ? Si oui donner une \'equation de
cette droite.

\item Parmi toutes ces droites, y en a-t-il qui sont  parall\`eles, confondues  ou perpendiculaires \`a la droite 
$\Delta$  d'\'equation $2x-3y+1=0$ ? Si oui donner des \'equations
 de ces droites.
\end{enumerate}

\item   On consid\`ere la famille de droites $D_m$  : 
$(2m-1)x+(3-m)y+m+1=0$, $m \in \Rr$.

Parmi toutes ces droites y en a-t-il une perpendiculaire 
\`a $(\Delta ): x+y-1=0$? Si oui, laquelle?
\end{enumerate}
\finenonce{001992}



\finexercice



\exercice{1993, liousse, 2003/10/01}

\enonce{001993}{}
On consid\`ere les trois points de $ P$ : $A(2,-3)$, $B(0,-1)$ et $C(-2,-5)$.
\begin{enumerate}
\item Dessiner le triangle $ABC$ puis calculer son aire.
\item Calculer les coordonn\'ees de l'orthocentre $H$, du centre du cercle
 circonscrit $\Omega$  et du centre de gravit\'e $G$ de $ABC$.
\item V\'erifier que $H$, $\Omega$ et $G$ sont align\'es et qu'en particulier $
\overrightarrow{\Omega G} =  {1\over 3}\overrightarrow{\Omega H}$.
\end{enumerate}
\finenonce{001993}



\finexercice



\exercice{1994, liousse, 2003/10/01}

\enonce{001994}{}
\begin{enumerate}
\item Calculer les  angles : 
\begin{enumerate}
\item entre les vecteurs ${\buildrel\rightarrow \over u_1}(\sqrt 3,2)$ et 
${\buildrel\rightarrow \over v_1}(1,3\sqrt 3),$

\item entre les vecteurs ${\buildrel\rightarrow \over u_2}(1,\sqrt 2)$ et 
${\buildrel\rightarrow \over v_2}(\sqrt 2-2,\sqrt 2+2),$

\item du triangle de sommets $A(-1,0),$ $B(\frac{1}{2},\frac{\sqrt{3}}{2})$ et 
$C(\frac{1}{2},-\frac{\sqrt{3}}{2}).$
\end{enumerate}

\item Calculer la distance du point $A$ \`a la droite $D$ : 
\begin{enumerate}
\item $A(1,1)$ et $D : 2x+y-1 =0$ 

\item $A(2,-1)$ et $D : 3x-2y+4 =0$ 

\item $A(3,3)$ et $D : -x+3y+2=0 $. 
\end{enumerate}

\item   Trouver les bissectrices de :
\begin{enumerate}
\item
 $D : 5x-12y+7 =0$ et $D' : 3x+4y-7 =0$,


\item  $D : x-3y+5 =0$ et$D' : 3x-y-1 =0$.
\end{enumerate}
\end{enumerate}
\finenonce{001994}



\finexercice



\exercice{1995, ridde, 1999/11/01}

\enonce{001995}{}
 Soit $ (0, \vec{i}, \vec{j})$ un rep\`ere du plan. D\'eterminer
l'expression analytique dans ce rep\`ere de la r\'eflexion d'axe
$x + y = 1$.
\finenonce{001995}



\finexercice



\exercice{1996, ridde, 1999/11/01}

\enonce{001996}{}
 Soit $G$ un sous-groupe fini de l'ensemble des isom\'etries
du plan. Montrer que $G$ ne peut pas contenir de translation non triviale.
\finenonce{001996}



\finexercice



\exercice{1997, ridde, 1999/11/01}

\enonce{001997}{}
On consid\`ere dans le plan les deux droites
$ (D : 3x + y = 5)$ et \\
$ (D' : x-2y + 3 = 0)$. Quel est l'angle entre ces deux droites ?
\finenonce{001997}



\finexercice



\exercice{1998, ridde, 1999/11/01}

\enonce{001998}{}
Soit $C$ un cercle de centre $ I = (x_{0}, y_{0})$ et de rayon
$R$ et $ (D : ax + by + c = 0)$. En param\'etrant $D$, montrer que $D$ est tangente \`a
$C$ (\emph{i.e.} $D\cap C$ est un singleton) ssi $d (I, D) = R$.
\finenonce{001998}



\finexercice



\exercice{1999, ridde, 1999/11/01}

\enonce{001999}{}
Soient $A$ et $B$ deux points du plan et $\alpha$ un r\'eel. D\'eterminer l'ensemble
des points $M$ qui v\'erifient $\overrightarrow{MA}.\overrightarrow{MB}  = \alpha$.
\finenonce{001999}



\finexercice



\exercice{2000, ridde, 1999/11/01}

\enonce{002000}{}
Soient $A, B, C$ les sommets d'un triangle \'equilat\'eral de cot\'e $1$. D\'eterminer
l'ensemble des points $M$ qui v\'erifient $MA^{2} + MB^{2} + MC^{2} = 2$.
\finenonce{002000}



\finexercice



\exercice{2001, ridde, 1999/11/01}

\enonce{002001}{}
Soient $A$ et $B$ deux points du plan et $k$ un r\'eel strictement positif.
D\'eterminer l'ensemble des points $M$ qui v\'erifient $MA = kMB$.
\finenonce{002001}



\finexercice



\exercice{2002, ridde, 1999/11/01}

\enonce{002002}{}
 Quelle est l'application $f : \begin{cases}
\Rr^{2} \rightarrow \Rr^{2}\\ \begin{pmatrix} x\\y \end{pmatrix} \mapsto
\frac{1}{5} \begin{pmatrix} -3x-4y\\-4x + 3y-2\end{pmatrix}\end{cases}$ ?
\finenonce{002002}



\finexercice



\exercice{2003, ridde, 1999/11/01}

\enonce{002003}{}
Soit $X = \left\{ A, B, C, D\right\}$ les sommets d'un carr\'e
du plan et $G = \left\{ f\in I_{2}/f (X) = X\right\}$. Montrer que $G$ est un sous-groupe
 de $I_{2}$. Montrer que si $f\in G$ alors $f (O) = O$ où $O$ est l'isobarycentre
 de $A, B, C, D$. En d\'eduire les \'el\'ements de $G$.
\finenonce{002003}



\finexercice



\exercice{2004, ridde, 1999/11/01}

\enonce{002004}{}
 D\'eterminer les $z\in \Cc$ tels que $z, z^{2}, z^{4}$ soient
align\'es.
\finenonce{002004}



\finexercice



\exercice{2005, ridde, 1999/11/01}

\enonce{002005}{}
Si $a$ et $b$ sont les affixes de deux sommets oppos\'es
d'un carr\'e, calculer les affixes des deux autres.
\finenonce{002005}



\finexercice



\exercice{2006, ridde, 1999/11/01}

\enonce{002006}{}
Soit $O, A, B$ un triangle rectangle en $O$. A toute droite
$D$ issue de $O$ on associe le cercle de diam\`etre $A^{'}B^{'}$ où $A^{'}$ et $B^{'}$
sont les projet\'es orthogonaux de $A$ et $B$ sur $D$. Montrer que tous les cercles
passent par un même point fixe (on pourra utiliser une similitude... ).
\finenonce{002006}



\finexercice



\exercice{2007, ridde, 1999/11/01}

\enonce{002007}{}
Pour $a, b, c$ trois nombres complexes tels que $b\neq c$, on note
$V (a, b, c) = \frac{c-a}{c-b}$. Soient $ z_{1}, z_{2}, z_{3}, z_{4}$
quatre nombres complexes distincts. Montrer que les images
de ces nombres complexes sont align\'ees ou cocycliques ssi
$\frac{V (z_{1}, z_{2}, z_{3})}{V (z_{1}, z_{2}, z_{4})} \in \Rr$.
\finenonce{002007}



\finexercice



\exercice{2008, ridde, 1999/11/01}

\enonce{002008}{}
Soit $ABCD$ un carr\'e direct et $M$ un point de la droite $ (DC)$. La perpendiculaire
\`a $ (AM)$ passant par $A$ coupe $ (BC)$ en $N$. On note I le milieu de $[MN]$.
D\'eterminer le lieu des points $I$ lorsque $M$ d\'ecrit la droite $ (DC)$.
\finenonce{002008}



\finexercice



\exercice{2009, ridde, 1999/11/01}

\enonce{002009}{}
Soient $A, B, C, D$ quatre points distincts du plan tels que $\overrightarrow{AB}
 \neq \overrightarrow{CD}$. Montrer que le centre de la similitude directe transformant $A$ en $C$
et $B$ en $D$ est aussi le centre de celle transformant $A$ en $B$ et $C$ en $D$.
\finenonce{002009}



\finexercice



\exercice{2010, liousse, 2003/10/01}

\enonce{002010}{}
Les quatre points $A$, $B$, $C$ et $D$ de l'espace 
sont-ils coplanaires ? Si oui, donner une \'equation cart\'esienne du plan qui les contient :
\begin{enumerate}
\item $A(1,2,2)$,  $B(-1,-2,-1)$, $C(3,4,4)$  et $D(-2,3,1)$.
\item  $A(0,1,3)$,  $B(1,2,-1)$, $C(1,1,-1)$  et $D(1,2,2)$.
\item  $A(-1,2,4)$,  $B(3,-3,0)$, $C(1,3,4)$  et $D(5,1,-6)$.
\item  $A(2,-1,0)$,  $B(0,-4,5)$, $C(4,-13,13)$  et $D(-4,5,-3)$.
\end{enumerate}
\finenonce{002010}



\finexercice



\exercice{2011, liousse, 2003/10/01}
\video{UmHg_4CL8x4}
\enonce{002011}{}
\begin{enumerate}
\item   Trouver une équation du plan $(P)$ 
défini par les éléments suivants.


\begin{enumerate}
\item  $A$, $B$ et $C$ sont des points de $(P)$

\begin{enumerate}
\item $A(0,0,1)$, $B(1,0,0)$ et $C(0,1,0)$.
\item $A(1,1,1)$, $B(2,0,1)$ et $C(-1,2,4)$.
\end{enumerate}

\item $A$ est un point de $(P)$,  $\vec{u}$ et $\vec{v} $ sont des vecteurs directeurs de $(P)$
\begin{enumerate}
\item $A(1,2,1)$, $\vec{u} (4,0,3)$ et $\vec{v} (1,3,-1)$.
\item $A(1,0,2)$, $\vec u(2,-1,3)$ et  $\vec v(-1,4,5)$.
\end{enumerate}

 \item  $A$ est un point de $(P)$,  $D$ est une droite contenue dans $(P)$

\begin{enumerate}
\item $A(0,0,0)$ et $ (D) : \left\{ \begin{array}{l}
                                           x + y - z + 3 = 0 \\
                                           4x - y + 2z = 0 
                                         \end{array} \right.$
\item $A(1,1,0)$ et $(D) : \left\{ \begin{array}{l}
                                         x = t \\
                                         y = -1 + 2t \\
                                         z = 1 - 3t
                     \end{array} \right.$
\end{enumerate}             

\item   $D$ et $D'$ sont des droites contenues dans $(P)$           
             
\begin{enumerate}
\item  $(D) :\left\{ \begin{array}{l}
                           x + y - z + 3 = 0\\
                           x - y - 2 = 0 
                     \end{array} \right.$
 et $(D') : \left\{ \begin{array}{l}
                       3x - y - z + 5 = 0 \\
                        x + y - z + 1 =0 
                     \end{array} \right.$
             
             

\item $(D) :\left\{ \begin{array}{l}
                           x + 2y - z + 1 = 0\\
                           x + 3y + z - 4 = 0 
                     \end{array} \right.$
 et $(D') : \left\{ \begin{array}{l}
                       2x + y - 3z + 7 = 0 \\
                        3x + 2y + z - 1 =0 
                     \end{array} \right.$
\end{enumerate}              
\end{enumerate}  
 \item  Montrer que les représentations paramétriques suivantes 
définissent le m\^eme plan :


$$\left\{ \begin{array}{l}x=2+s+2t \\ y=2+2s+t \\ z=1-s-t \end{array} \right.
\quad \text{ et } \quad  \left\{ \begin{array}{l} x=1+3s'-t'\\ y=3+3s'+t'\\ z=1-2s' \end{array} \right.$$
\end{enumerate} 
\finenonce{002011}



\finexercice



\exercice{2012, liousse, 2003/10/01}

\enonce{002012}{}
Les plans suivants 
sont-ils parall\`eles ou s\'ecants? Dans ce dernier cas, donner un vecteur
directeur de la droite $(D)=(P) \cap (P')$.

\begin{enumerate}
\item  $(P): 5x-y-1=0$ et $(P'): z=3$.
\item  $(P): x+y+z+1=0$ et $(P'): 2x-y+3z+2=0$.
\item  $(P): 2x-z+1=0$ et $(P'): 4x-3y+2z+5=0$.
\item  $(P): 4x-6y+8z-1=0$ et $(P'): -6x+12y-9z+11=0$.
\end{enumerate} 
\finenonce{002012}



\finexercice



\exercice{2013, liousse, 2003/10/01}

\enonce{002013}{}
 Quelle est 
la nature de l'intersection des trois plans suivants? 
Si c'est un point en donner
les coordonn\'ees, si c'est une droite 
en donner un vecteur directeur.
\begin{enumerate}
\item $(P): z=1$, $(P'): x-y-2=0$ et $(P"): 4x-2y+z+2=0$.
\item $(P): 4x-2y+3z+5=0$, $(P'): 3x+y-z+2=0$ et $(P"): x-y+z+1=0$.
\item $(P): 4x-2y+10z-4=0$, $(P'): -10x+5y-25z+13=0$ et $(P"): x+y-z+1=0$.
\item $(P): 3x-y+2z-5=0$, $(P'): x-y+3z-7=0$ et $(P"): 4x+2y-z+1=0$.
\item $(P): x-y+2z-1=0$, $(P'): 2x+y+z+3=0$ et $(P"): x-4y+5z-6=0$.
\item $(P): x-y+2z-1=0$, $(P'): 2x+y-z+1=0$ et $(P"): x+5y-8z+2=0$.
\end{enumerate} 
\finenonce{002013}



\finexercice



\exercice{2014, liousse, 2003/10/01}

\enonce{002014}{}
Les droites suivantes sont-elles s\'ecantes, parall\`eles 
ou non coplanaires ?
Si elles sont s\'ecantes donner leur point d'intersection 
et si elles sont parall\`eles
donner un vecteur directeur.
\begin{enumerate}
\item $(D) : \left\{ \begin{array}{l}
                            x+y-z+2=0\\
                           x+y+z+1=0
                     \end{array} \right.$ 
et $(D') : \left\{ \begin{array}{l}
                                3x-y+2z-7=0 \\
                 x-y=0 
                 \end{array} \right.$
        
\item $(D) :\left\{ \begin{array}{l}
                          x=1-2t \\
              y=t+2 \\
                          z=3t+1  
                     \end{array} \right.$
 et $(D') : \left\{ \begin{array}{l}
                       x=3t-1  \\
                       y=-t+2  \\
               z=2t
                     \end{array} \right.$
             
\end{enumerate} 
\finenonce{002014}



\finexercice



\exercice{2015, liousse, 2003/10/01}

\enonce{002015}{}
Dans chacun des cas suivants 
dire si la droite $(D)$ et le plan $(P)$ sont parall\`eles ou s\'ecants.
Donner alors leur point d'intersection.
\begin{enumerate}
\item $(D): \left\{ \begin{array}{l}
                                5x-3y+2z-5=0 \\
                  2x-y-z-1=0  \end{array} \right.$ et $(P): 4x-3y+7z-7=0$.
\item $(D):  \left\{ \begin{array}{l}
                                  x=3+2t\\
                   y=5-3t\\
                   z=2-2t \end{array} \right.$ et $(P): -3x+2y+3z-5=0$.
\end{enumerate} 
\finenonce{002015}



\finexercice



\exercice{2016, liousse, 2003/10/01}

\enonce{002016}{}
On consid\`ere les cinq points suivants:
 $A(1,2,-1)$, $B(3,2,0)$, $C(2,1,-1)$, $D(1,0,4)$ et $E(-1,1,1)$.
\begin{enumerate}
\item Ces quatre points sont-ils coplanaires ?
\item D\'eterminer la nature du triangle $ABC$. $A$, $B$ et $C$ sont-ils align\'es, si non
 donner une  \'equation cat\'esienne du plan $P$ qui les contient.
\item D\'eterminer les coordonn\'ees
du barycentre $G$ des points $A$, $B$, $C$ et $D$.
\item Montrer que $O$, $D$ et $G$ sont align\'es et que la droite $OD$ est perpendiculaire \`a $P$.
\end{enumerate} 
\finenonce{002016}



\finexercice



\exercice{2017, liousse, 2003/10/01}

\enonce{002017}{}
  Soient $D_1$, $D_2$ et $D_3$ trois droites concourrantes en $\Omega$ et soient  $P$, $P'$ et $P''$ trois plans tels que aucun ne contient aucune des 3 droites ci dessus.
   On peut alors d\'efinir les 9 points d'intersections   :
$P$ coupe $D_1$, $D_2$, $D_3$ en    $A$, $B$, $C$ ;
$P'$ coupe $D_1$, $D_2$, $D_3$ en    $A'$, $B'$, $C'$ ;
$P'$ coupe $D_1$  $D_2$, $D_3$ en    $A''$, $B''$, $C''$ ;

  On consid\`ere aussi les intersections suivantes :
$I=(AB')\cap (A'B)$ , $J=(AC')\cap (A'C)$ ,$K=(BC')\cap (B'C)$.

  Montrer que les droites $(A''K)$, $B''J)$ et $(C''I)$ sont parall\`elles ou concourrantes.
  ({\it Indication : utiliser  un bon rep\`ere affine}).
\finenonce{002017}



\finexercice



\exercice{2018, liousse, 2003/10/01}

\enonce{002018}{}
On consid\`ere les quatre points suivants:
$A(2,0,0)$, $B(-1,\sqrt{3},0)$, $C(-1,-\sqrt{3},0)$, $D(0,0,4)$.
D\'eterminer un vecteur directeur de la droite $(ABC) \cap (ADE)$.
\finenonce{002018}



\finexercice



\exercice{2019, liousse, 2003/10/01}

\enonce{002019}{}
Donner une condition sur $m$ pour que les trois plans 
suivants se coupent sur une m\^eme droite.
$(P): x+my-z+1=0$, $(P'): (m+1)x+3y+4z-2=0$ et $(P"): y+(2m+4)z-(2m+2)=0$.
\finenonce{002019}



\finexercice



\exercice{2020, liousse, 2003/10/01}
\video{0i30Mt7oXdU}
\enonce{002020}{}
On considère la famille de  plans $(P_m)_{m\in \Rr}$ définis par les équations cartésiennes :
$$ m^2x+(2m-1)y+mz=3$$
\begin{enumerate}
  \item Déterminer les plans $P_m$ dans chacun des cas suivants :
  \begin{enumerate}
    \item $A(1,1,1)\in P_m$ 
    \item $\vec{n}(2,-\frac 52,-1)$ est normal à $P_m$.
    \item  $\vec{v}(1,1,1)$ est un vecteur directeur de $P_m$ 
  \end{enumerate} 
  \item Montrer qu'il existe un unique point $Q$ appartenant à tous les plans $P_m$.
\end{enumerate} 
\finenonce{002020}



\finexercice


\exercice{2021, liousse, 2003/10/01}
\video{6HAwYKkwS8Y}
\enonce{002021}{}
\begin{enumerate}
\item Déterminer la distance du point $A$ au plan $(P)$

\begin{enumerate}
\item $A(1,0,2)$ et $(P): 2x+y+z+4=0$.
\item $A(3,2,1)$ et $(P): -x+5y-4z=5$.
\end{enumerate} 
\item Calculer la distance du point $A(1,2,3)$ à la droite 
$(D):  \left\{ \begin{array}{l}
-2x+y-3z=1\\ x+z=1  \end{array} \right.$
\end{enumerate} 
\finenonce{002021}



\finexercice


\exercice{2022, liousse, 2003/10/01}

\enonce{002022}{}

 On consid\`ere les deux 
droites  $(D): \left\{ \begin{array}{l} y-z=3\\ -x-y+2=0 \end{array} \right. $ et
$(\Delta):\left\{ \begin{array}{l}  -x+3z=1 \\ -x-3y=2 \end{array} \right. .$
\begin{enumerate}
\item Donner un vecteur directeur de $D$ et de $\Delta$.

\item Donner une \'equation param\'etrique de $\Delta$.

\item On fixe un point $M_{\alpha}$ de $\Delta$ d\'ependant du param\`etre $\alpha$ o\`u
$\alpha$ est l'abscisse de point $M_{\alpha}$. Donner une \'equation du plan 
$P_{\alpha}$ passant par $M_{\alpha}$ et contenant $D$.

\item Parmi tous ces plans, y en a-t-il un qui est perpendiculaire \`a $\Delta$ ?
Pour quelle valeur $\alpha_0$ de $\alpha $ est il obtenu ? Donner une \'equation
de ce plan. Donner les coordonn\'ees de $M_{\alpha_0}$.
\end{enumerate} 
\finenonce{002022}



\finexercice



\exercice{2023, liousse, 2003/10/01}

\enonce{002023}{}

On se donne $2$ droites $D_1$ et $D_2$ ayant comme vecteurs directeurs respectifs
 $ \vec{u_1}$ et $\vec{u_2}$.
\begin{enumerate}
\item 
  Perpendiculaire commune \`a ces deux droites.

\begin{enumerate}
\item  On suppose que $ \vec{u_1}$ et $\vec{u_2}$ ne sont pas colin\'eaires et 
  on note $  \vec{n}:=\vec{u_1} \wedge \vec{u_2} $.



\begin{enumerate}
\item
 Montrer que le plan $P_1$ contenant $D_1$ et admettant  $\vec{n}$ comme vecteur directeur et
le plan $P_2$ contenant $D_2$ et admettant  $\vec{n}$ comme vecteur directeur se coupent en une droite
$\Delta$.

\item  Montrer que $\Delta$ est une perpendiculaire commune \`a $D_1$ et $D_2$
 (c'est \`a dire $\Delta$ coupe  $D_1$  et $D_2$, et 
 est  orthogonale \`a $D_1$ et \`a  $D_2$).


\item  Montrer que $\Delta$ est la seule  perpendiculaire commune \`a $D_1$ et $D_2$.
\end{enumerate} 


\item   Comment  construire $\Delta$ dans le cas o\`u 
$D_1$ et $D_2$ sont parall\`elles ?

\end{enumerate} 

\item   Distance entre  ces deux droites.

Soit $H_1:= D_1\cap \Delta$ et $H_2:= D_2\cap \Delta$.


 Montrer que pour tout
$A_1\in D_1$ et tout $A_2\in D_2$, on  a $d(A_1,A_2)\geq d(H_1,H_2)$.

 $ d(H_1,H_2)$ est appel\'ee {\it distance entre les deux droites $D_1$ et $D_2$}.


\item Donner des \'equations cart\'esiennes pour $\Delta$ et calculer
la distance entre les  deux droites $D_1$ et $D_2$ dans le cas suivant :


\begin{enumerate}
\item $(D_1): \left\{\begin{array}{l} 
                           x-y-z+4=0 \\
                         -x-2y-3z+9=0\end{array} \right.$ et 
$(D_2): \left\{\begin{array}{l} -x+2y+z+2=0 \\-2x+4y-z+1=0\end{array} \right. $




\item  $(D_1): \left\{\begin{array}{l}
                   x+y-z+2=0 \\
            x+y+z+1=0 \end{array} \right.$ et 
$(D_2):\left\{\begin{array}{l} 3x-y+2z-7=0 \\ x-y=0 \end{array} \right.$

\item  $(D_1):\left\{\begin{array}{l}
                              x=1-2t\\ 
                  y=t+2\\  
                   z=3t+1  \end{array} \right.$ et 
$(D_2):\left\{\begin{array}{l}
                      x=3t-1\\
              y=-t+2\\   
              z=2t \end{array} \right. $

\item  $(D_1): \left\{ \begin{array}{l}
                                   x-y-z-2=0 \\
                   x-2y-3z+1=0 \end{array} \right. $ et 
$(D_2): \left\{\begin{array}{l}
                        x+y+2z-1=0 \\
             2x+y+z+2=0  \end{array} \right.$
\end{enumerate} 
 \end{enumerate} 
\finenonce{002023}



\finexercice



\exercice{2024, liousse, 2003/10/01}

\enonce{002024}{}
\begin{enumerate}
\item D\'eterminer les plans bissecteurs de :

$P: x+y+z+3=0$ et $P' : 2x + y+2z=1$

$Q : 5x + 3y -4z =8$ et $Q':4x-5y-3z=2$.


\item D\'eterminer l'ensemble des points de l'espace \'equidistants des trois axes de coordonn\'ees.

\item On consid\`ere la droite 
$D$ d'\'equation param\'etrique
$\left\{\begin{array}{l}
                   x=3t-1 \\
            y=1 \\
             z=-t-1\end{array} \right.$ 
             
             
Donner une \'equation des deux plans $P$ et $P'$ contenant $D$ \`a une 
distance de 1 de l'origine (point $O$ de coordonn\'ees $(0,0,0)$).
\end{enumerate} 
\finenonce{002024}



\finexercice



\exercice{2025, ridde, 1999/11/01}

\enonce{002025}{}
D\'eterminer l'expression analytique de la r\'eflexion
$s$ de plan $x + y-z = 1$. Quelle est l'image par $s$ du plan
$x + 2y -3z + 1 = 0$ ?
\finenonce{002025}



\finexercice



\exercice{2026, ridde, 1999/11/01}

\enonce{002026}{}
D\'eterminer la distance du point $M = (1, 2, 3)$ aux droites
$$D\begin{cases}x + y-2z = 1\\2x-y + z + 1 = 0\end{cases} \text{ et }
\Delta \begin{cases}x = 1 + 2t\\y = 2-t\\z = 2 + 2t\end{cases}$$
\finenonce{002026}



\finexercice



\exercice{2027, cousquer, 2003/10/01}

\enonce{002027}{}
 Soit deux plans 
$\left\{
\begin{array}{lrcl}
    \pi: & ux+vy+wz+h & = & 0  \\
    \pi': & u'x+v'y+w'z+h' & = & 0
\end{array}
\right.$.
\begin{enumerate}
\item  Montrer que si $\pi$ et $\pi'$ sont sécants, tout plan passant
par leur droite d'intersection~$D$ a une équation du type
$$\lambda(ux+vy+wz+h)+\mu(u'x+v'y+w'z+h')=0$$
et réciproquement, tout plan ayant une équation de ce type, 
(pour un couple $(\lambda, \mu$) donné)
passe par~$D$.
\item  Si $\pi$ et $\pi'$ sont parallèles, que représente
l'ensemble des plans d'équation~:
$$\lambda(ux+vy+wz+h)+\mu(u'x+v'y+w'z+h')=0$$
\end{enumerate}

\finenonce{002027}



\finexercice



\exercice{2028, cousquer, 2003/10/01}

\enonce{002028}{}
 Écrire l'équation du plan passant par la droite
$\left\{
\begin{array}{rcl}
    3x+2y+5z+6 & = & 0  \\
    x+4y+3z+4 & = & 0
\end{array}\right.$
et parallèle à la droite $\displaystyle\frac{x-1}{3}=
\frac{y-5}{2}=\frac{z+1}{-3}$.
\finenonce{002028}



\finexercice



\exercice{2029, cousquer, 2003/10/01}

\enonce{002029}{}
 Soit la droite d'équations
$\left\{
\begin{array}{rcl}
    3x-2y-z+4 & =& 0  \\
    x-4y-3z-2 & = & 0
\end{array}\right.$.
Trouver sa projection sur le plan $5x+2y+2z-7=0$.
\finenonce{002029}



\finexercice



\exercice{2030, cousquer, 2003/10/01}

\enonce{002030}{}
Soit les droites $D$ et $D'$ non coplanaires~:
$$(D) \;\left\{
\begin{array}{rcl}
    x-y+z+1 & = & 0  \\
    2x+y-z  & = & 0
\end{array}\right.\quad
\mbox{et}\quad (D') \; \left\{
\begin{array}{rcl}
    x+2y+z & = & 0  \\
    2x-2y-2z-1 & = & 0
\end{array}\right.$$
Trouver des équations de leur perpendiculaire commune.
\finenonce{002030}



\finexercice



\exercice{2698, matexo1, 2002/02/01}
\enonce{002698}{}
Soit $ABC$ un triangle \'equilat\'eral de
c\^ot\'e unit\'e et $T_0$ son int\'erieur. On consid\`ere les figures
g\'eom\'etriques $T_n$ obtenues par r\'ecurrence de la mani\`ere
suivante\,:
sur chaque c\^ot\'e $MN$ de $T_{n-1}$, on ajoute l'int\'erieur d'un
triangle \'equilat\'eral
$PQR$, o\`u $P$ et $Q$ sont sur le segment $[MN]$, aux tiers de sa
longueur, et $R$ est
ext\'erieur \`a $T_{n-1}$. Finalement on d\'efinit le sous-ensemble du
plan $K$ par
$$ K = \bigcup_{n\geq 0} T_n.$$

\begin{enumerate}
\item Faire un dessin respr\'esentant $T_0$, $T_1$, $T_2$...
\item Donner l'aire de $T_n$ sous forme de s\'erie. Quelle est l'aire de
$K$\,? 
\item M\^emes questions avec le p\'erim\`etre, puis le diam\`etre de
$T_n$ et $K$.
\end{enumerate}
\finenonce{002698}
\finexercice


\exercice{5195, rouget, 2010/06/30}
\enonce{005195}{**I}
$(ABC)$ est un vrai triangle.
\begin{enumerate}
\item  Montrer que ses médianes sont concourantes en $G$ l'isobarycentre de $(ABC)$.
\item  Montrer que ses médiatrices sont concourantes en $O$ le centre du cercle circonscrit à $(ABC)$.
\item  Montrer que ses hauteurs sont concourantes en $H$ l'orthocentre de $(ABC)$ puis montrer la relation d'\textsc{Euler}~:

$\overrightarrow{OH}=3\overrightarrow{OG}$ (considérer l'homothétie de centre $G$ et de rapport $-2$).
\item  Montrer que ses  bissectrices (intérieures) sont concourantes en $I$ le centre du cercle inscrit.
\end{enumerate}
\finenonce{005195}


\finexercice


\exercice{5196, rouget, 2010/06/30}
\enonce{005196}{**IT}
On donne les points $A(1,2)$, $B(-2,1)$ et $C(0,4)$.
\begin{enumerate}
\item  Déterminer $\widehat{BAC}$ au degré près.
\item  Déterminer l'aire du triangle $(ABC)$.
\item  Déterminer son isobarycentre, son orthocentre, le centre de son cercle circonscrit puis une équation de ce cercle.
\item  Déterminer une équation des bissectrices de l'angle $\widehat{BAC}$ puis de la bissectrice intérieure à l'angle $\widehat{A}$.
\end{enumerate}
\finenonce{005196}


\finexercice


\exercice{5199, rouget, 2010/06/30}
\enonce{005199}{**}
Soit $(E)$ l'ensemble d'équation cartésienne $2x^2+5xy+3y^2-3x-2y-5=0$. Montrer que $(E)$ est une réunion de deux droites. Déterminer l'aire du parallélogramme formé par ces deux droites et les parallèles à ces deux droites passant par $O$.
\finenonce{005199}


\finexercice


\exercice{5200, rouget, 2010/06/30}
\enonce{005200}{**}
Déterminer un cercle tangent aux trois droites d'équations respectives $y=2x+1$, $y=2x+7$ et $y=-\frac{1}{2}x$.
\finenonce{005200}


\finexercice


\exercice{5202, rouget, 2010/06/30}
\enonce{005202}{***I}
\label{exo:suprou8bis}
Soient $n$ un entier supérieur ou égal à $2$, puis $A_1$, $A_2$,..., $A_n$ $n$ points du plan. Existe-t-il $n$ points $B_1$, $B_2$,..., $B_n$ tels que, pour $i\in\{1,...,n\}$, $A_i$ soit le milieu de $[B_i,B_{i+1}]$ (avec la convention $B_{n+1}=B_1$)~? (Utiliser l'exercice précédent.)
\finenonce{005202}


\finexercice


\exercice{5203, rouget, 2010/06/30}
\enonce{005203}{*T}
Soit $\mathcal{C}$ la courbe d'équation $x^2+y^2-2x+4y+1=0$.
\begin{enumerate}
\item  Déterminer une équation de la tangente au point de $\mathcal{C}$ de coordonnées $(2,-2+\sqrt{3})$.
\item  Déterminer l'intersection de $\mathcal{C}$ et du cercle de centre $(1,0)$ et de rayon $2$.
\end{enumerate}
\finenonce{005203}


\finexercice


\exercice{5204, rouget, 2010/06/30}
\enonce{005204}{*** Théorème de \textsc{Ménélaüs}}
Soient $A$, $B$ et $C$ trois points non alignés. Soient $M$, $N$ et $P$ trois points appartenant respectivement aux droites $(BC)$, $(CA)$ et $(AB)$ et distincts de $A$, $B$ et $C$. Montrer que~:~

$$(M,\;N,\;\mbox{et}\;P\;\mbox{sont alignés})\Leftrightarrow(\frac{\overline{MB}}{\overline{MC}}.\frac{\overline{NC}}{\overline{NA}}\frac{\overline{PA}}
{\overline{PB}}=1).$$

(Trouver une démonstration utilisant le théorème de \textsc{Thalès}, une utilisant la composée de deux homothéties et une utilisant des coordonnées.)
\finenonce{005204}


\finexercice

\exercice{5208, rouget, 2010/06/30}
\enonce{005208}{** Faisceaux de droites}
\begin{enumerate}
\item  Soient $(D)$ et $(D')$ deux droites sécantes d'équation respectives $ax+by+c=0$ et $a'x+b'y+c'=0$, $(a,b)\neq(0,0)$, $(a',b')\neq(0,0)$. Soit $(\Delta)$ une droite. Montrer que $(D)$, $(D')$ et $(\Delta)$ sont concourantes si et seulement si il existe $(\Delta)$ a une équation cartésienne de la forme $\lambda(ax+by+c)+\mu(a'x+b'y+c')=0$, $(\lambda,\mu)\neq(0,0)$.
\item  Equation cartésienne de la droite passant par le point $(1,0)$ et par le point d'intersection des droites d'équations respectives $5x+7y+1=0$ et $-3x+2y+1=0$
\item  Pour $m\in\Rr$, on considère $(D_m)$ la droite d'équation $(2m-1)x+(m+1)y-4m-1=0$. Montrer que les droites $(D_m)$ sont concourantes en un point $A$ que l'on précisera. Toute droite passant par $A$ est-elle une droite $(D_m)$~?
\end{enumerate}
\finenonce{005208}


\finexercice

\exercice{5510, rouget, 2010/07/15}
\enonce{005510}{**}
Dans $\Rr^3$, soient $(D)$ $\left\{
\begin{array}{l}
x-z-a=0\\
y+3z+1=0
\end{array}
\right.$ et $(D')$ $\left\{
\begin{array}{l}
x+2y+z-2b=0\\
3x+3y+2z-7=0
\end{array}
\right.$.
Vérifier que $(D)$ et $(D')$ ne sont pas parallèles puis trouver $a$ et $b$ pour que $(D)$ et $(D')$ soient sécantes. Former alors une équation cartésienne de leur plan.
\finenonce{005510}


\finexercice


\exercice{5511, rouget, 2010/07/15}
\enonce{005511}{**}
Système d'équations cartésiennes de la droite $(\Delta)$ parallèle à la droite $(D)$~:~$2x=3y=6z$ et sécante aux droites $(D_1)$~:~$x=z-4=0$ et $(D_2)$~:~$y=z+4=0$.
\finenonce{005511}


\finexercice


\exercice{5512, rouget, 2010/07/15}
\enonce{005512}{***}
Trouver toutes les droites sécantes aux quatre droites $(D_1)$~$x-1=y=0$, $(D_2)$~:~$y-1=z=0$, $(D_3)$~:~$z-1=x=0$ et $(D_4)$~:~$x=y=-6z$.

\finenonce{005512}


\finexercice


\exercice{5513, rouget, 2010/07/15}
\enonce{005513}{**T}
Dans $\Rr^3$ euclidien rapporté à un repère orthonormé, on donne $A(2,-2,0)$, $B(4,2,6)$ et 
$C(-1,-3,0)$. Déterminer l'orthocentre, le centre de gravité, les centres des cercles circonscrits et inscrits au triangle $(A,B,C)$.
\finenonce{005513}


\finexercice

\exercice{5514, rouget, 2010/07/15}
\enonce{005514}{**T}
Soit $M(x,y,z)$ un point de $\Rr^3$ rapporté à un repère orthonormé. Déterminer la distance de $M$ à la droite 
$(D)$ $\left\{
\begin{array}{l}
x+y+z+1=0\\
2x+y+5z=2
\end{array}
\right.$. En déduire une équation du cylindre de révolution d'axe $(D)$ et de rayon $2$. 
\finenonce{005514}


\finexercice


\exercice{5515, rouget, 2010/07/15}
\enonce{005515}{**T}
Dans $\Rr^3$ rapporté à un repère orthonormé, soient $(D)$ $\left\{
\begin{array}{l}
x+y+z+1=0\\
2x+y+5z=2
\end{array}
\right.$ et $(D')$ $\left\{
\begin{array}{l}
x+y+z=2\\
2x+y-5z=3
\end{array}
\right.$. Déterminer la distance de $(D)$ à $(D')$ puis la perpendiculaire commune à ces deux droites.
\finenonce{005515}


\finexercice


\exercice{5516, rouget, 2010/07/15}
\enonce{005516}{**}
Montrer que les plans $(P_1)$~:~$z-2y=5$, $(P_2)$~:~$2x-3z=0$ et $(P_3)~:~3y-x=0$ admettent une parallèle commune. Ils définissent ainsi un prisme. Déterminer l'aire d'une section perpendiculaire.
\finenonce{005516}


\finexercice


\exercice{5517, rouget, 2010/07/15}
\enonce{005517}{*T}
Angle des plans $x+2y+2z=3$ et $x+y=0$.
\finenonce{005517}


\finexercice


\exercice{5518, rouget, 2010/07/15}
\enonce{005518}{**T}
Soient $(P_1)~:~4x+4y-7z-1=0$ et $(P_2)~:~8x-4y+z+7=0$. Trouver une équation cartésienne des plans bissecteurs de $(P_1)$ et $(P_2)$.
\finenonce{005518}


\finexercice


\exercice{5519, rouget, 2010/07/15}
\enonce{005519}{**T}
Déterminer la perpendiculaire commune aux droites $(D)$ et $(D')$~:~$(D)$ $\left\{
\begin{array}{l}
x+y-3z+4=0\\
2x-z+1=0
\end{array}
\right.$ et $(D')$ $\left\{
\begin{array}{l}
x=z-1\\
y=z-1
\end{array}
\right.$. 
\finenonce{005519}


\finexercice


\exercice{5521, rouget, 2010/07/15}
\enonce{005521}{**I}
Déterminer les différents angles d'un tétraèdre régulier (entre deux faces, entre deux arêtes et entre une arête et une face).
\finenonce{005521}


\finexercice


\exercice{5522, rouget, 2010/07/15}
\enonce{005522}{**T}
Déterminer la distance de l'origine $O$ à la droite $(D)$ dont un système d'équations cartésiennes est $\left\{
\begin{array}{l}
x-y-z=0\\
x+2y-z=10
\end{array}
\right.$.
\finenonce{005522}


\finexercice


\exercice{6884, rouget, 2012/09/05}
\video{CGIhcrD5Jtg}
\enonce{006884}{}
% Exos de J.-L. Rouget #5197 légèrement modifié
Déterminer le projeté orthogonal du point $M_0(x_0,y_0)$ sur la droite $(D)$ 
d'équation $2x-3y=5$ ainsi que son symétrique orthogonal.
\finenonce{006884}


\finexercice

\exercice{7415, mourougane, 2021/08/10}
\enonce{007415}{Questions de cours}
\begin{enumerate}
\item Rappeler les définitions du parallélisme et du parallélisme faible.
\item Soit $(\lambda_i)$ $p$ nombre réels de somme égale à $1$ et $M_i$ 
$p$ points d'un espace affine.
Rappeler le sens de la notation $\sum_i \lambda_i M_i$.
Est-il nécessaire que les points $M_i$ soient distincts ?
\end{enumerate}
\finenonce{007415}
\finexercice


\exercice{7416, mourougane, 2021/08/10}
\enonce{007416}{}
Donner la liste des sous-espaces affines du plan affine $\Rr^2$ et de l'espace affine $\Rr^3$.
\finenonce{007416}
\finexercice


\exercice{7417, mourougane, 2021/08/10}
\enonce{007417}{Avec le cours}
\begin{itemize}
\item Deux plans de l'espace affine $\Rr^3$ peuvent-ils avoir un unique point d'intersection ?
\item Dans l'espace affine $\Rr^n$, quelle est la dimension 
d'un sous-espace affine défini par deux équations affines
indépendantes ? 
\end{itemize}
\finenonce{007417}
\finexercice


\exercice{7418, mourougane, 2021/08/10}
\enonce{007418}{Avec le cours}
Soit $A,B,C,D$ quatre points d'un plan affine. Soit $I$ le milieu de $[C,D]$.
L'isobarycentre des points $A,B,C,D$ est-il le centre de gravité du triangle $A,B,I$?
\finenonce{007418}
\finexercice


\exercice{7419, mourougane, 2021/08/10}
\enonce{007419}{Dans $\Rr^3$, à partir d'équations}
Soit l'espace affine $\Rr^3$ muni d'un repère affine
$A_0,A_1,A_2,A_3$.
\begin{enumerate}
\item
Montrer que le sous-ensemble $A$ de l'espace affine $\Rr^3$ d'équation 
$$M\left(\begin{array}{c}
x\\y\\z\end{array}\right)\in A\iff \left\{ 
\begin{array}{c}
2x-y-z-3=0\\x+y-2z=3
\end{array}
\right.$$
est un sous-espace affine. Préciser sa dimension, l'espace vectoriel
directeur et un repère affine.
\item
Même question avec $B$ d'équation 
$$M\left(\begin{array}{c}
x\\y\\z\end{array}\right)\in B\iff \left\{ 
\begin{array}{c}
x+y=2\\2x+2y=3z+1\\5x+5y=10z
\end{array}
\right.$$
\end{enumerate}
\finenonce{007419}
\finexercice


\exercice{7420, mourougane, 2021/08/10}
\enonce{007420}{Dans $\Rr^3$, trouver des équations}
Soit l'espace affine $\Rr^3$ muni d'un repère affine
$A_0,A_1,A_2,A_3$.

\begin{enumerate}
\item
Trouver un système d'équations pour le sous-espace affine
$\mathcal{A}$ passant par le point
 $\displaystyle A\left(\begin{array}{c}
1\\2\\3\end{array}\right)$ et parallèle au plan d'équation 
$(2x-y-z=5)$.
\item
Trouver un système d'équations pour le sous-espace affine
$\mathcal{B}$ passant par
le point
 $\displaystyle A\left(\begin{array}{c}
1\\2\\3\end{array}\right)$ de direction $\Rr \vec{u}\oplus \Rr\vec{v}$
 o{ù} $\vec{u}$
est le vecteur de coordonées $\left(\begin{array}{c}
-1\\0\\3\end{array}\right)$ et $\vec{v}$ le vecteur de coordonées
$\left(\begin{array}{c} 
2\\1\\3\end{array}\right)$
dans la base
$(\vec{A_0A_1},\vec{A_0A_2}\vec{A_0A_3})$.
\item
Trouver un système d'équations pour le sous-espace affine
$\mathcal{C}$ engendré par les points
 $\displaystyle A\left(\begin{array}{c}
1\\2\\3\end{array}\right)$, $\displaystyle B\left(\begin{array}{c}
4\\0\\-1\end{array}\right)$ et $\displaystyle C\left(\begin{array}{c}
1\\1\\0\end{array}\right)$.
\end{enumerate}
\finenonce{007420}
\finexercice


\exercice{7421, mourougane, 2021/08/10}
\enonce{007421}{Dans $\Rr^3$, trouver encore des équations}
Soit l'espace affine $\Rr^3$ muni d'un repère affine
$A_0,A_1,A_2,A_3$.

\begin{enumerate}
\item
Trouver un système d'équations pour le sous-espace affine
$\mathcal{D}$ passant par le point
 $\displaystyle A\left(\begin{array}{c}
1\\2\\3\end{array}\right)$ et parallèle à la droite d'équation 
$\left\{\begin{array}{c}2x-y-z=5\\x+y+z=3\end{array}\right.$.
\item
Trouver un système d'équations pour le sous-espace affine
$\mathcal{E}$ passant par
le point
 $\displaystyle A\left(\begin{array}{c}
1\\2\\3\end{array}\right)$ de direction $\Rr \vec{u}$ o{ù} $\vec{u}$
est le vecteur de coordonées $\left(\begin{array}{c}
-1\\0\\3\end{array}\right)$ 
dans la base
$(\vec{A_0A_1},\vec{A_0A_2}\vec{A_0A_3})$.
\item
Trouver un système d'équations pour le sous-espace affine
$\mathcal{F}$ engendré par
les points
 $\displaystyle A\left(\begin{array}{c}
1\\2\\3\end{array}\right)$ et $\displaystyle B\left(\begin{array}{c}
4\\0\\-1\end{array}\right)$.
\end{enumerate}
\finenonce{007421}
\finexercice


\exercice{7422, mourougane, 2021/08/10}
\enonce{007422}{Calcul en coordonnées cartésiennes dans le plan}
Soit $a$ et $b$ deux réels non nuls distincts.
Dans un plan affine muni d'un repère, calculer une équation
cartésienne de la droite $\Delta$ passant par le point de coordonnées
$(a,b)$ et par le point d'intersection des deux droites $D$ d'équation
$x/a+y/b=1$ et $D'$ d'équation $x/b+y/a=1$.
\finenonce{007422}
\finexercice


\exercice{7423, mourougane, 2021/08/10}
\enonce{007423}{Deux droites sécantes dans l'espace}
Dans un espace affine de dimension $3$ muni d'un repère affine.
\begin{enumerate}
\item Montrer que les droites $D$ et $D'$ d'équations cartésiennes
$$D ~:~x+y-z-2=0\quad \text{ et } \quad 2x-y+3z-1=0$$
et $$D' ~:~x-2y-3\quad \text{ et } \quad 3x+6y-1$$
sont concourantes.
\item Trouver une équation cartésienne du plan qu'elles déterminent.
\end{enumerate}
\finenonce{007423}
\finexercice


\exercice{7424, mourougane, 2021/08/10}
\enonce{007424}{Sur les médianes}
Dans l'espace affine $\Rr^2$, on considère le triangle $ABC$.
Choisir un repère affine $A_0,A_1,A_2$ adapté pour montrer simplement que les
médianes du triangle $ABC$ sont concourantes. On devra calculer une
équation pour chaque médiane et montrer qu'elles sont concourantes.
\finenonce{007424}
\finexercice


\exercice{7425, mourougane, 2021/08/10}
\enonce{007425}{Dessiner}
Dans le plan affine $\Rr^2$ muni d'un repère $A_0,A_1,A_2$,
représenter l'enveloppe convexe des points donnés en coordonnées cartésiennes
$\displaystyle A\left(\begin{array}{c}
-2\\ -3\end{array}\right)$, $\displaystyle B\left(\begin{array}{c}
0\\0\end{array}\right)$, $\displaystyle C\left(\begin{array}{c}
0\\4\end{array}\right)$, $\displaystyle D\left(\begin{array}{c}
4\\3\end{array}\right)$, $\displaystyle E\left(\begin{array}{c}
5\\0\end{array}\right)$.
\finenonce{007425}
\finexercice


\exercice{7426, mourougane, 2021/08/10}
\enonce{007426}{Une propriété des tétraèdres}
 Soient $\mathcal{E}$ un espace affine de dimension 3, et $A,B,C,D$ un
tétraèdre de $\mathcal{E}$. Montrer que les droites joignant les milieux des cotés
opposés du tétraèdre sont concourantes.
\finenonce{007426}
\finexercice


\exercice{7427, mourougane, 2021/08/10}
\enonce{007427}{Centre de gravité}
Dans un plan affine, soit $ABC$ et $A'B'C'$ deux triangles de centre
de gravité $G$ et $G'$ respectivement.
\begin{enumerate}
\item Calculer $\vec{AB'}+\vec{BC'}+\vec{CA'}$ en fonction de $G$ et $G'$.
\item Montrer que $ABC$ et $A'B'C'$ ont le même centre de gravité si
 et seulement s'il existe un point $D$ tel que $DBA'C$ et $DB'AC'$
 soient des parallélogrammes.
\end{enumerate}
\finenonce{007427}
\finexercice


\exercice{7428, mourougane, 2021/08/10}
\enonce{007428}{Localiser des points en coordonnées barycentriques}
\begin{enumerate}
\item Soit $(AB)$ une droite dans un espace affine déterminée par deux points distincts $A$ et $B$. Décrire à l'aide des coordonnées barycentriques dans le repère $AB$ les trois régions de la droite découpées par les points $A$ et $B$.
\item Soit $ABC$ un triangle non plat dans un plan affine $E$.
Décrire à l'aide des coordonnées barycentriques dans le repère $ABC$ les sept régions découpées par les droites qui portent les cotés du triangle $ABC$.
\end{enumerate}
\finenonce{007428}
\finexercice


\exercice{7429, mourougane, 2021/08/10}
\enonce{007429}{Exercice de construction}
On suppose savoir tracer la parallèle à une droite donnée passant par un point donné.
On peut utiliser un compas mais seulement pour reporter des longueurs égales.
Partager en sept parties de même longueur un segment donné.
\finenonce{007429}
\finexercice


\exercice{7430, mourougane, 2021/08/10}
\enonce{007430}{Dans un plan affine}
On considère dans l'espace affine $\Rr^2$ un triangle non aplati
$ABC$. Soit $(a,b)$, $(c,d)$ et $(e,f)$ trois couples de nombre réels
de somme non nulle. On désigne par $C'$ le barycentre de $(A,a ;B,b)$,
$A'$ le barycentre de $(B,c ;C,d)$ et $B'$ le barycentre de $(C,e
;A,f)$. Montrer que $A', B'$ et $C'$ sont alignés si et seulement si
$ace=-bdf$.
\finenonce{007430}
\finexercice


\exercice{7431, mourougane, 2021/08/10}
\enonce{007431}{}
Une application affine peut-elle avoir exactement deux points fixes distincts ?

Donner un exemple d'application affine sans point fixe, qui n'est pas une translation.
\finenonce{007431}
\finexercice


\exercice{7432, mourougane, 2021/08/10}
\enonce{007432}{}
Soit le plan affine $E$ muni du repère cartésien
$(A,\vec{u},\vec{v})$. On considère l'application $f$ qui à tout
$\displaystyle M\left(\begin{array}{c}
x\\y\end{array}\right)$ associe le point $M'\left(\begin{array}{c}
2x-5y+3\\-4x+10y-1\end{array}\right)$.
\begin{enumerate}
\item  Montrer que $f$ est affine et écrire la matrice de son
 application linéaire associée $\vec{f}$ dans la base $(\vec{u},\vec{v})$.
\item  Déterminer les points fixes de $f$.
\item  Montrer que $Im\vec{f}$ et $Ker\vec{f}$ sont
 supplémentaires dans $\vec{E}$. Donner une base de chacun de ces
 sous-espaces.
\end{enumerate}
\finenonce{007432}
\finexercice


\exercice{7433, mourougane, 2021/08/10}
\enonce{007433}{Les translations}
Montrer à l'aide de la règle du parallélogramme
que les seules applications affines dont la partie linéaire est l'identité sont les translations. 
\finenonce{007433}
\finexercice


\exercice{7434, mourougane, 2021/08/10}
\enonce{007434}{Les homothéties}
Soit $E$ un espace affine et $f$ une application affine dont la partie linéaire est multiple de l'identité de $\vec{E}$.

On suppose que $\vec{f}=\lambda Id_{\vec{E}}$ avec $\lambda\not =1$.
Décrire les valeurs propres de $\vec{f}$. Montrer que
$f$ a un unique point fixe. Montrer alors que $f$ est une homothétie.
\finenonce{007434}
\finexercice


\exercice{7435, mourougane, 2021/08/10}
\enonce{007435}{Les projections}
Soit $E$ un espace affine et $F$ et $G$ deux sous espaces affines de direction supplémentaire dans $\vec{E}$. Décrire $F\cap G$.
\begin{enumerate}
\item Soit $M$ un point de $E$. 
Montrer que le sous-espace affine passant par $M$ et parallèle à $G$ rencontre $F$ en un unique point noté $p(M)$.

\item Montrer que l'application $p$ est affine.

\item Déterminer les points fixes de $p$.

\item Montrer que $p\circ p=p$.
\end{enumerate}
\finenonce{007435}
\finexercice


\exercice{7436, mourougane, 2021/08/10}
\enonce{007436}{Les symétries}
Soit $E$ un espace affine et $F$ et $G$ deux sous espaces affines de direction supplémentaire dans $\vec{E}$. 
\begin{enumerate}
\item Décrire la symétrie $s$ par rapport à $F$ parallèlement à $G$. Est-elle bijective ?

\item Soit $M$ un point de $E$. Quel est le milieu du segment $[M,s(M)] $?

\item Déterminer $ker(\vec{s}-Id_{\vec{E}})$ et $ker(\vec{s}+Id_{\vec{E}})$.

\item Soit $\vec{u}$ un vecteur de $G$. Déterminer l'application $s\circ t_{\vec{u}}\circ s^{-1}$.
\end{enumerate}
\finenonce{007436}
\finexercice


\exercice{7437, mourougane, 2021/08/10}
\enonce{007437}{Les affinités et transvections}

Dans l'espace affine $E=\Rr^3_{aff}$, on considère un endomorphisme affine $f$ qui admet un hyperplan de points fixes $F$. Soit $A$ un point de $E$ qui n'est pas dans $F$.
\begin{enumerate}
\item Que dire de $f$ si $f(A)=A$ ?

\item On suppose maintenant que $f(A)$ est différent de $A$. On suppose de plus que la droite $(Af(A))$ coupe le plan $F$ en un point $\Omega$.
Expliquer comment construire $f(M)$? 
On dit alors que l'application $f$ est une affinité.

\item On suppose maintenant que $f(A)$ est différent de $A$. On suppose de plus que la droite $(Af(A))$ ne rencontre pas le plan $F$.
Expliquer comment construire $f(M)$? 
On dit alors que l'application $f$ est une transvection.
\end{enumerate}
\finenonce{007437}
\finexercice


\exercice{7438, mourougane, 2021/08/10}
\enonce{007438}{Écrire des expressions analytiques}
Dans l'espace affine $\Rr^3$ muni d'un repère affine
$A_0,A_1,A_2,A_3$, donner l'expression analytique 
\begin{enumerate}
 \item \begin{enumerate}
\item de l'homothétie $h$ de centre 
$\displaystyle C\left(\begin{array}{c}
1\\2\\3\end{array}\right)$ et de rapport $4$.
\item de la symétrie $s$ d'axe $(x+y+z=1)$ parallèlement à
 $\vec{A_0A_1}$. 
\item de l'affinité $a$ de base $(x+y+z=1)$ de rapport $3$
 parallèlement à  $\vec{A_0A_1}$. 
\item de la transvection $t$ de base $(x+y+z=1)$ qui envoie $A_0$ sur 
$\displaystyle \left(\begin{array}{c}
1\\1\\ -2\end{array}\right)$
\end{enumerate}

\item Donner des expressions analytiques des applications affines
 précédentes dans des repères mieux adaptés à
 déterminer. 
\end{enumerate}
\finenonce{007438}
\finexercice


\exercice{7439, mourougane, 2021/08/10}
\enonce{007439}{Constructions}
\begin{enumerate}
\item Deux droites se coupent hors de la feuille en un point $I$.
Soit $A$ un point de la feuille.
Tracer la droite $(AI)$.

\item Soit $d$ et $d'$ deux droites parallèles et $M$ un point du plan.
Tracer avec une règle (non graduée) la parallèle à $d$ passant par $M$.
\end{enumerate}
\finenonce{007439}
\finexercice


\exercice{7440, mourougane, 2021/08/10}
\enonce{007440}{}
Dans un plan affine, quel est l'ensemble des milieux des segments dont les extrémités
appartiennent respectivement à deux segments donnés.
\finenonce{007440}
\finexercice 


\exercice{7441, mourougane, 2021/08/10}
\enonce{007441}{}
Soit $x_i$ $n$ nombres réels. Montrer que 
$\displaystyle\left(\sum_1^n x_i\right)^2\leq n \sum_1^n x_i ^2$.
\finenonce{007441}
\finexercice


\exercice{7442, mourougane, 2021/08/10}
\enonce{007442}{}
On suppose que $F$ et $G$ sont deux sous-espaces vectoriels de
$E$. Montrer que  $(F+G)^{\perp}=F^{\perp}\cap G^{\perp}$ et que
$(F\cap G)^{\perp}=F^{\perp}+G^{\perp}$.
\finenonce{007442}
\finexercice


\exercice{7443, mourougane, 2021/08/10}
\enonce{007443}{Procédé d'orthonormalisation de Gram-Schmidt}
Dans l'espace vectoriel euclidien $\Rr^3$ muni du produit scalaire
standard  et  de  la  base  canonique,  appliquer  le  procédé
d'orthonormalisation de Gram-Schmidt à la base
$$v_1=\left( \begin{array}{c}  1\\0\\1\end{array}\right)\ ; \ 
v_2=\left(
\begin{array}{c}   2\\ 1\\ 0 \end{array}\right)\   ;   \  
 v_3=\left(
\begin{array}{c} -1\\ -1\\ -1 \end{array}\right)
$$ pour obtenir une base orthonormée.
\finenonce{007443}
\finexercice


\exercice{7444, mourougane, 2021/08/10}
\enonce{007444}{}
 On  considère $\Rr^{4}$  muni de  son  produit scalaire
usuel. Soit $F$ le sous espace engendré par $v_{1}=(1,1,0,0)$,
$v_{2}=(0,1,-1,1)$. Déterminer une base orthonormée de $F$ et la
compléter pour obtenir une base orthonormée de $\Rr^{4}$.
\finenonce{007444}
\finexercice


\exercice{7445, mourougane, 2021/08/10}
\enonce{007445}{Vrai ou faux, Préciser}
Soit $(E,<.>)$ un espace vectoriel euclidien. Répondre par vrai ou
faux, puis dire sous quelles hypothèses supplémentaires sur la base
l'affirmation est vraie.
\begin{enumerate}
 \item La matrice $A$ d'une application linéaire orthogonale $u$ dans une base $\mathcal B$ de $E$ vérifie $\mbox{}^tAA=Id$.

 \item Dans toute base orthonormée de $\Rr^3$, la matrice d'une
rotation  est  de   la  forme  $  \left(\begin{array}{ccc}
1&0&0\\0&\cos\theta&-\sin\theta\\0&\sin\theta&\cos\theta
\end{array}\right)
$
\end{enumerate}
\finenonce{007445}
\finexercice


\exercice{7446, mourougane, 2021/08/10}
\enonce{007446}{}
Soit $E$ l'espace vectoriel $\Rr^3$ muni du produit scalaire canonique
et de la base canonique. Soit $H$ le plan d'équation $x+2y+2z=0$. Soit
$\pi$ la projection orthogonale sur $H$ et $s$ la symétrie orthogonale
par rapport à $H$.

 \begin{enumerate}
 \item Déterminer un vecteur $\epsilon_1$ normal à $H$ et unitaire.

 \item Pour tout vecteur $V$ de $E$, écrire $V-\pi (V)$, puis
  $\pi(V)$ à l'aide de $V$ et $\epsilon_1$ seulement. (On pourra
  utiliser des produits scalaires comme coefficients).

 \item Déterminer les matrices de $\pi$ et de $s$ dans la base
 canonique. Sont-elles orthogonales, symétriques ?
\end{enumerate}
\finenonce{007446}
\finexercice


\exercice{7447, mourougane, 2021/08/10}
\enonce{007447}{Cours}
Démontrer que dans un espace affine
euclidien $E$ la fonction $E\times E\to\Rr, (x,y)\mapsto \| x-y\| $
est une fonction distance.
\finenonce{007447}
\finexercice


\exercice{7448, mourougane, 2021/08/10}
\enonce{007448}{Distance entre deux droites de l'espace}
Soit $\mathcal{E}$ un espace affine euclidien de dimension $3$. Pour
tout couple de droites $(D_1,D_2)$ on appelle $$d(D_1,D_2)=\inf\{ \|
y_1-y_2 \| , y_1 \in D_1, y_2\in D_2\}.$$ 
\begin{enumerate}
 \item Calculer $d(D_1,D_2)$ quand $D_1$ et $D_2$ sont concourantes.

 \item Soit $a_1\in D_1$ et $a_2\in D_2$. En décomposant $a_1-a_2$ dans
 $\vec{D_1} +\vec{D_2}+(\vec{D_1} +\vec{D_2})^\perp$ montrer qu'il
 existe $x_1\in D_1$ et $x_2\in D_2$ tels que
 $d(D_1,D_2)=d(x_1,x_2)$. 

 \item Montrer que pour $z_1\in D_1$ et $ z_2\in D_2$, 
$$d(D_1,D_2)=d(z_1,z_2) _iff z_1-z_2\in \vec{D_1}^\perp \cap
\vec{D_2}^\perp.$$ 

 \item Montrer que si $e_i$ est un vecteur directeur de $D_i$
 $$d(D_1,D_2)^2=\frac{Gram(a_1-a_2,e_1,e_2)}{Gram(e_1,e_2)}$$ 
o{ù} $Gram(u_1,u_2,\cdots ,u_r):=det(\langle u_i,u_j\rangle
)_{1\leq i,j\leq r}$.

 \item Calculer la distance entre les deux droites données par les équations
cartésiennes dans un repère orthonormé de $\mathcal{E}$~:
\begin{eqnarray*}
M\left( \begin{array}{c}x\\y\\z\end{array}\right)\in D_1
\iff\left\{\begin{array}{c}x+y=1\\x+y+2z=1\end{array}\right.\ \ \ 
M\left( \begin{array}{c}x\\y\\z\end{array}\right)\in D_2
\iff\left\{\begin{array}{c}y+z=1\\x+y-2z=3\end{array}\right.
\end{eqnarray*}
\end{enumerate}
\finenonce{007448}
\finexercice


\exercice{7449, mourougane, 2021/08/10}
\enonce{007449}{Angle de deux demi-droites de même origine}
Dans le plan euclidien orienté $\mathcal{P}$,
 on considère deux demi-droites $d_1$ et $d_2$ de même origine $\Omega$.

\begin{enumerate}
 \item Rappeler la forme générale d'une matrice de rotation vectorielle dans une
 base orthonormée directe de $\vec{\mathcal{P}}$. Comment change cette
 matrice quand on change de base orthonormée ?

 \item Montrer qu'il existe une unique rotation de centre $\Omega$ qui
 envoie $d_1$ sur $d_2$. On pourra donc définir l'angle orienté de
 deux demi-droites comme l'angle de la rotation qui envoie la
 première sur la seconde.
\end{enumerate}
\finenonce{007449}
\finexercice


\exercice{7450, mourougane, 2021/08/10}
\enonce{007450}{}
Dans le plan euclidien muni d'un repère orthonormé on considère les
points $A\left( \begin{array}{c} 1\\2 \end{array} \right)$ 
et $B\left( \begin{array}{c} -1\\ -1 \end{array}\right)$.

\begin{enumerate}
 \item Déterminer selon la valeur de $r$ l'ensemble des
points $M$ du plan tels que $MA^2-2MB^2=r$.

 \item Déterminer selon la valeur de $r$ l'ensemble des
points $M$ du plan tels que $MA^2-MB^2=r$.
\end{enumerate}
\finenonce{007450}
\finexercice


\exercice{7451, mourougane, 2021/08/10}
\enonce{007451}{Sur les homothéties translations}
On travaille dans un espace affine euclidien.
On rappelle que les homothéties-translations sont caractérisées par le
fait qu'elles transforment toute droite en une droite parallèle.
On appelle dilatation une application affine dont la partie linéaire
est une homothétie vectorielle.

\begin{enumerate}
 \item Montrer que l'ensemble des dilatations coïncide avec
 l'ensemble des homothéties-translations.

 \item Montrer que l'ensemble des dilatations est un sous-groupe
 distingué du groupe des applications affines.

 \item On travaille maintenant dans le plan affine euclidien $\mathcal{P }$.
 Montrer qu'il existe exactement deux dilatations qui
 transforment un cercle donné en un cercle donné.
\end{enumerate}
\finenonce{007451}
\finexercice


\exercice{7452, mourougane, 2021/08/10}
\enonce{007452}{Inégalité isopérimétrique}
On considère dans le plan euclidien muni d'un repère orthonormé la
courbe $C$ d'équation polaire $r=f(\theta )$ où $f$ est une fonction
continue positive sur $[-\frac{\pi}{2},\frac{\pi}{2}]$ avec
$f(-\frac{\pi}{2})=f(\frac{\pi}{2})=0$.
\begin{enumerate} 
 \item Représenter $C$ dans le cas où $f=cos$.

 \item On rappelle que l'aire de la surface $S$ délimitée par la courbe $C$
est donnée par $$a(S)=\int_{-\frac{\pi}{2}}^{\frac{\pi}{2}}
\frac{1}{2}f^2(\theta )d\theta. $$
Montrer que $$a(S)\leq \frac{\pi}{4}\left(diam (S)\right)^2.$$
\end{enumerate}
\finenonce{007452}
\finexercice


\exercice{7453, mourougane, 2021/08/10}
\enonce{007453}{Cours}
Dans un espace vectoriel euclidien, deux sous-espaces vectoriels non nuls orthogonaux sont-ils toujours en somme directe ?
\begin{enumerate}
 \item Montrer en utilisant la forme réduite que le déterminant d'une isométrie
est $(-1)^{codim E_1}$ où $E_1$ est l'espace propre de valeur propre $1$.

 \item Rappeler toutes les isométries du plan euclidien et de l'espace
euclidien de dimension 3, en précisant leur
partie linéaire, leur point fixe, leur axe, leur composante à point fixe et leur composante de glissement.

 \item Rappeler la construction de l'axe radical de deux cercles non sécants.

 \item Déterminer le groupe d'isométries d'un segment dans le plan euclidien.
\end{enumerate}
\finenonce{007453}
\finexercice


\exercice{7454, mourougane, 2021/08/10}
\enonce{007454}{}
Soit $A$ et $B$ deux points fixés dans un plan affine euclidien. Déterminer le lieu des points $M$ du plan où $MA\perp MB$.
\finenonce{007454}
\finexercice


\exercice{7455, mourougane, 2021/08/10}
\enonce{007455}{}
Le but de cet exercice est de 
démontrer que le symétrique de l'orthocentre $H$ d'un triangle non plat $ABC$ par rapport à un des cotés (par exemple $(AC)$) est sur le cercle circonscrit.

Soit $ABC$ un triangle non plat. La hauteur issue de $A$ coupe $(BC)$ en $A'$ et la hauteur issue de $C$ coupe $(AB)$ en $C'$. Démontrer que les points $A',B,C'$ et $H$ sont cocycliques. Démontrer que les angles de droites $((BC'),(BA'))$ et $((HC'),(HA'))$ sont égaux.
Conclure.
\finenonce{007455}
\finexercice


\exercice{7456, mourougane, 2021/08/10}
\enonce{007456}{}
Deux cercles sont dits orthogonaux si les tangentes aux points d'intersection sont orthogonales.
Montrer que les cercles $C$et $C'$ sont orthogonaux si et seulement si la puissance du centre de $C$
par rapport à $C'$ est égale au carré du rayon de $C$.
\finenonce{007456}
\finexercice


\exercice{7457, mourougane, 2021/08/10}
\enonce{007457}{}
Soit $A$ et $B$ deux points d'un plan affine euclidien.
Déterminer suivant la valeur de la constante~$k$,
\begin{enumerate}
 \item l'ensemble des points $M$ du plan tels que $MA^2+MB^2=k$
 \item l'ensemble des points $M$ du plan tels que $MA^2-MB^2=k$
 \item et l'ensemble des points $M$ du plan tels que $MA/MB=k$
 \item l'ensemble des points $M$ du plan tels que $MA/MB<k$
\end{enumerate}
\finenonce{007457}
\finexercice


\exercice{7458, mourougane, 2021/08/10}
\enonce{007458}{}
En décomposant les rotations en produits de réflexions, déterminer
le centre d'une composée de deux rotations dont la somme des angles n'est pas nulle modulo $2\pi$.
\finenonce{007458}
\finexercice


\exercice{7459, mourougane, 2021/08/10}
\enonce{007459}{Décomposition des isométries planes}
\begin{enumerate}
 \item Que peut-on dire d'une isométrie plane qui a trois points fixes non alignés ?

 \item Soit $\phi$ une isométrie plane qui a deux points fixes distincts $A$ et $B$.
Soit $C$ un point hors de la droite $(AB)$ et $C'$ son image par $\phi$.
Si $C'$ est différent de $C$, déterminer la médiatrice du segment $[CC']$
et montrer que $s_{(AB)}\circ \phi$ est l'identité.
En déduire la nature de $\phi$.

 \item Soit $\phi$ une isométrie différente de l'identité qui a un point fixe $A$.
Soit $B$ un autre point du plan et $B'$ son image.
Si $B'$ est différent de $B$, soit $d$ la médiatrice du segment $[BB']$
Montrer que $s_d\circ \phi$ a deux points fixes.
Montrer que $\phi$ est composée de réflexions.

 \item Montrer qu'une isométrie plane qui n'a pas de point fixe
peut s'écrire composée de moins de trois réflexions.
\end{enumerate}
\finenonce{007459}
\finexercice


\exercice{7460, mourougane, 2021/08/10}
\enonce{007460}{}
On munit le plan affine euclidien d'un repère orthonormé $(O,i,j)$ et on l'identifie au plan complexe.
\'Ecrire à l'aide des affixes complexes, la symétrie glissée d'axe d'équation $x+y=2$ et de vecteur $(3,3)$.
\finenonce{007460}
\finexercice


\exercice{7461, mourougane, 2021/08/10}
\enonce{007461}{Déterminer une rotation à partir d'images}
 Soit $E$ un plan affine euclidien muni d'un repère cartésien 
orthonormé. Soient
$A$, $B$, $C$ et $D$ les points de $E$ dont les coordonnées sont
$$A : (0,3), \quad B : (2,1), \quad C : (2,3) \text{ et } D 
: (0,1).$$
\begin{enumerate}
 \item Montrer que les droites $(A,B)$ et $(C,D)$ sont 
orthogonales et expliciter les coordonnées de leur
point d'intersection.
 \item Prouver l'existence d'une {\it {rotation}} qui envoie $A$ 
sur $C$, $C$ sur $B$, $B$ sur $D$ et $D$
sur $A$. Expliciter une représentation matricielle de cette rotation.
\end{enumerate}
\finenonce{007461}
\finexercice


\exercice{7462, mourougane, 2021/08/10}
\enonce{007462}{Trouver l'isométrie}
 Soit $E$ un espace affine euclidien de dimension 3 muni d'un repère 
cartésien orthonormé. On note $v$ la transformation de $E$ dans $E$
qui  envoie le point de coordonnées $(x,y,z)$ sur le point de 
 coordonnées $(x',y',z')$ définies par:
 $$ x' = \frac{2x-2y+z+1}{3}; y' =\frac{2x+y-2z+2}{3} ; z'
 =\frac{x+2y+2z+5}{3}.$$  
  
 Montrer que $v$ est une isométrie de $E$. Préciser de quel type
 d'isométrie il s'agit. Expliciter son axe et son vecteur de
 glissement. 
\finenonce{007462}
\finexercice


\exercice{7463, mourougane, 2021/08/10}
\enonce{007463}{Composée d'isométries}
Soit $E$ un espace affine euclidien de dimension $3$ muni d'un repère cartésien
orthonormé $(O,i,j,k)$. On désigne par
$D$ la droite d'équation $(x=0,z=1)$ et par $D'$ la droite d'équation 
$(y=0,z=0)$. On note $S_D$ la
symétrie par rapport à la droite $D$ et $R_{\theta}$ la rotation 
d'axe $D'$ et d'angle $\theta$ (en
considérant la base $(j,k)$ comme directe). On pose $\varphi=S_D\circ 
R_{\theta}$.
\begin{enumerate}
 \item {\'E}crire dans la base $(i,j,k)$ la matrice de 
$\overrightarrow {S_D}$, celle de
$\overrightarrow {R_{\theta}}$ et celle de $\overrightarrow{\varphi}$.
{\'E}crire les expressions analytiques de $S_D$ et de $R_{\theta}$ dans
le repère $(O,i,j,k)$.
 \item Montrer que $\varphi$ est une symétrie éventuellement 
glissée d'axe une droite $\Delta$.
 \item Pour tout point $M$ de $E$, prouver que les milieux de
$\bigl(M,s_\Delta (M)\bigr)$ et de $\bigl(M,\varphi(M)\bigr)$ 
sont sur $\Delta$.
 \item En utilisant le point $O$, montrer que $\Delta$ passe par le
point de coordonnées $(0,0,1)$.
et est contenue dans le plan affine d'équation $x=0$.
 \item Donner les composantes du vecteur de glissement de 
$\varphi$ en fonction de $\theta$.
\end{enumerate}
\finenonce{007463}
\finexercice


\exercice{7464, mourougane, 2021/08/10}
\enonce{007464}{}
\begin{enumerate}
 \item Soit $C$ un cercle et $\vec{u}$ un vecteur.
Construire une corde $[AB]$ du cercle $C$ telle que $\vec{AB}=\vec{u}$.

 \item Construire un segment $[AB]$ connaissant son milieu $I$ et sachant que $A$ appartient à une droite donnée $d$ et $B$ à un cercle donné $C$.

 \item Construire un carré $ABCD$ sachant que $A$ et $C$ sont sur une droite donnée $d_1$ que $B$ est sur une droite donnée $d_2$ et que $D$ est sur une droite donnée $d_3$.
\end{enumerate}
\finenonce{007464}
\finexercice


\exercice{7465, mourougane, 2021/08/10}
\enonce{007465}{}
 Soit $d$ une droite.
Soit $\mathcal{C}$ un cercle et $A$ un
point de $\mathcal{C}$.
Construire un cercle tangent à la droite $d$ et 
tangent en $A$ au cercle $\mathcal{C}$.
\finenonce{007465}
\finexercice


\exercice{7466, mourougane, 2021/08/10}
\enonce{007466}{}
On considére le plan muni d'un un repère orthonormé ($O, \overrightarrow {\imath},\overrightarrow{\jmath}$) et la courbe $(C)$ d'équation 
\begin{center}$x^{2} + xy + y^{2} + x - y =0 $ \end{center}
\begin{enumerate}
 \item Montrer que cette courbe possède un centre de symétrie $\Omega$ et donner son équation dans le repère ($\Omega, \overrightarrow {\imath},\overrightarrow{\jmath}$)


 \item Montrer que dans le repère ($\Omega, \overrightarrow {\imath },\overrightarrow{\jmath}$) la première bissectrice $\Delta$ est axe de symétrie.

 \item On considère le repère orthonormé direct ($\Omega, \overrightarrow {I},\overrightarrow{J}$) où $\overrightarrow {I}$ est un vecteur unitaire de $\Delta$. Donner l'équation de $(C)$ dans ce repère.

 \item Montrer que dans le repère ($\Omega, \overrightarrow {\imath },\overrightarrow{\jmath}$) la seconde bissectrice $\Delta^{'}$ est axe de symétrie.

Donner les équations de $\Delta$ et $\Delta^{'}$ dans le repère $(O, \overrightarrow {\imath},\overrightarrow{\jmath})$.

Sachant que $(C)$ est une ellipse, tracer $(C)$ dans le repère $(O, \overrightarrow {\imath},\overrightarrow{\jmath})$.
\end{enumerate}

Reprendre cet exercice mais en utilisant la théorie des formes quadratiques et leur application aux coniques. Notamment retrouver les axes de $(C)$.
\finenonce{007466}
\finexercice


\exercice{7467, mourougane, 2021/08/10}
\enonce{007467}{}
On considére le plan muni d'un un repère orthonormé ($O, \overrightarrow {\imath},\overrightarrow{\jmath}$) et la courbe $(C)$ d'équation 

\begin{center}$x^{2} - 3xy +2y^{2} + 2x - 3y + 1=0 $ \end{center}

\begin{enumerate}
 \item Montrer que cette courbe possède un centre de symétrie $\Omega$ et donner son équation dans le repère ($\Omega, \overrightarrow {\imath},\overrightarrow{\jmath}$)

 \item En déduire que $(C)$ est la réunion de deux droites dont on donnera les équations dans le rèpere ($O, \overrightarrow {\imath},\overrightarrow{\jmath}$) 
\end{enumerate}
\finenonce{007467}
\finexercice


\exercice{7468, mourougane, 2021/08/10}
\enonce{007468}{}
On considère le plan muni d'un un repère orthonormé ($O, \overrightarrow {\imath},\overrightarrow{\jmath}$) et la courbe $(C)$ d'équation 

\begin{center}$2x^{2} + 3xy -2y^{2} - 10 = 0$ \end{center}
\begin{enumerate}
 \item Vérifier que $O$ est un centre de symétrie. Trouver une base orthonormale tel que l'équation de $(C)$ soit de la forme $\displaystyle \frac {x^{2}} {a^{2}} - \frac {y^{2}} {b^{2}} = 1$ $(a,b$ réels).

 \item Donner l'équation des asymptotes de $(C)$ et tracer $(C)$ dans le repère orthonormé ($O, \overrightarrow {\imath},\overrightarrow{\jmath}$).
\end{enumerate}
\finenonce{007468}
\finexercice


\exercice{7469, mourougane, 2021/08/10}
\enonce{007469}{}
On considère le plan euclidien muni d'un un repère orthonormé ($O, \overrightarrow {\imath},\overrightarrow{\jmath}$) et la courbe $(C)$ d'équation 

\begin{center}$4x^{2} - 4xy +y^{2} -3x -y - 1= 0$ \end{center}
\begin{enumerate}
 \item Montrer que $(C)$ est une parabole.

 \item Trouver un repère orthonormé ($S, \overrightarrow {u_{1}},\overrightarrow{u_{2}}$) tel que $(C)$ ait une équation de la forme $ x^{2} = 2py$ dans ce repère.
\end{enumerate}
\finenonce{007469}
\finexercice


\exercice{7470, mourougane, 2021/08/10}
\enonce{007470}{}
On considère le plan euclidien muni d'un repère orthonormé ($O, \overrightarrow {\imath},\overrightarrow{\jmath}$)  et la courbe $(C)$ d'équation

\begin{center}$ x^{2} - 2xy + y^{2} -6x - 10y + 9 = 0$ \end{center}

\begin{enumerate}
 \item Montrer que $(C)$ est une parabole.

 \item Trouver un repère orthonormé ($S, \overrightarrow {u_{1}},\overrightarrow{u_{2}}$) tel que $(C)$ ait une équation de la forme $ x^{2} = 2py$ dans ce repère. 

 \emph{Indication.} On devra trouver que dans le repère orthonormé ($O, \overrightarrow {\imath},\overrightarrow{\jmath}$) le sommet a pour coordonnées $(0,1)$ et pour axe la droite $ y = x + 1$.
\end{enumerate}


\finenonce{007470}
\finexercice


\exercice{7471, mourougane, 2021/08/10}
\enonce{007471}{}
Etudier les coniques suivantes dont les équations sont données dans un repère orthonormé ($O, \overrightarrow {\imath},\overrightarrow{\jmath}$) 

On précisera suivant la nature le centre, les axes, le sommet et les asymptotes.
\begin{enumerate}
 \item $x^{2} + 8xy -5y^{2} -28x + 14y +3 = 0$.

\emph{Indication.}  On trouvera une hyperbole de centre $(2,3)$ dans le repère ($O, \overrightarrow {\imath},\overrightarrow{\jmath}$) et d'équation réduite (donc rapportée à ses axes) $ 3x^{2} - 7y^{2} = 4$. 


 \item $ 2x^{2} + xy + y^{2} +4x - y - 2 =0$

\emph{Indication.}  On trouvera une ellipse de centre $(\frac{-9}{7}, \frac {8} {7}$) dans le repère orthonormé ($O, \overrightarrow {\imath},\overrightarrow{\jmath}$) et d'équation réduite (donc rapportée à ses axes) $ \frac{3 + \sqrt{2}}{2} x^{2} +\frac{3 - \sqrt{2}} {2} y^{2} - \frac {36}{7} = 0$ 
\end{enumerate}
\finenonce{007471}
\finexercice


\exercice{7472, mourougane, 2021/08/10}
\enonce{007472}{}
On considére le plan euclidien muni d'un un repère orthonormé ($O, \overrightarrow {\imath},\overrightarrow{\jmath}$) et la courbe $(C)$ d'équation 

\begin{center}$sqrt{3}x^{2} +6xy + \sqrt{3}y^{2} +2(\sqrt{3} - 6)x -2(3+\sqrt{3})y + 1 =0 $ \end{center}

\begin{enumerate}
 \item  Montrer que cette courbe posséde un centre de symétrie $\Omega$ et donner son équation dans le repère ($\Omega, \overrightarrow {\imath},\overrightarrow{\jmath}$)

 \item Montrer que dans le repère ($\Omega, \overrightarrow {\imath},\overrightarrow{\jmath}$) la droite $\Delta$ d'équation $y = \sqrt{3}x$ est axe de symétrie.

  \item On considère le repère orthonormé direct ($\Omega, \overrightarrow {I},\overrightarrow{J}$) où $\overrightarrow {I}$ est un vecteur unitaire de $\Delta$. Donner l'équation de $(C)$ dans ce repère.

 \item Montrer que dans le repère ($\Omega, \overrightarrow {\imath},\overrightarrow{\jmath}$) la perpendiculaire $\Delta^{'}$ à $\Delta$ est axe de symétrie.

 \item Sachant que $(C)$ est une hyperbole dessiner graphiquement $(C)$ dans le repère ($O, \overrightarrow {\imath},\overrightarrow{\jmath})$.
\end{enumerate}
\finenonce{007472}
\finexercice

\section{240.01 Sous-espaces affines}


\exercice{4859, quercia, 2010/03/17}
\enonce{004859}{Ensi Physique P 94}

    Soient $I,J,K$ trois points du plan. Montrer l'équivalence entre les trois
    propriétés~:
    
       a) $I$, $J$, $K$ sont alignés.\par
       b) Il existe $M$ tel que $\det(\vec{MI},\vec{MJ}) +  \det(\vec{MJ},\vec{MK}) +  \det(\vec{MK},\vec{MI}) = 0$.\par
       c) Pour tout point $M$, on a $\det(\vec{MI},\vec{MJ}) +  \det(\vec{MJ},\vec{MK}) +  \det(\vec{MK},\vec{MI}) = 0$.\par

\finenonce{004859}



\finexercice


\exercice{4860, quercia, 2010/03/17}
\enonce{004860}{Faisceau de plans}

On considère deux plans non parallèles de ${\cal E}_3$ ayant pour équation dans un repère
${\cal R} = (O,\vec i,\vec j,\vec k\,)$ :

$$\begin{cases} P :& ax+by+cz+d = 0\cr P':&a'x+b'y+c'z+d' = 0.\cr\end{cases}$$

Soit $D = P \cap P'$. Montrer qu'un plan $Q$ contient $D$ si et seulement
s'il a pour équation dans ${\cal R}$ :
$$\alpha(ax+by+cz+d) + \beta(a'x+b'y+c'z+d') = 0$$

avec $\alpha, \beta \in \R$ non tous deux nuls.


\finenonce{004860}



\finexercice


\exercice{4861, quercia, 2010/03/17}
\enonce{004861}{\'Equation d'un plan}

Dans ${\cal E}_3$ muni d'un repère $(O,\vec i,\vec j,\vec k\,)$, on donne :
$A : \left(\begin{smallmatrix} 1 \cr -1 \cr 1 \cr\end{smallmatrix}\right)$
et $D : \begin{cases} x - 3y + 2z = 1 \cr 2x + y - 3z = -1.\cr\end{cases}$

Donner l'équation cartésienne du plan passant par $A$ et $D$.

\finenonce{004861}



\finexercice


\exercice{4862, quercia, 2010/03/17}
\enonce{004862}{Droites coplanaires}


Dans ${\cal E}_3$ muni d'un repère $(O,\vec i,\vec j,\vec k\,)$, on donne :
$D : \begin{cases} x - 2z = 1 \cr y - z = 2 \cr\end{cases}$
et $D' : \begin{cases} x + y + z = 1 \cr x - 2y + 2z = a.\cr\end{cases}$

\begin{enumerate}
  \item  Pour quelles valeurs de $a$, $D$ et $D'$ sont-elles coplanaires ?
     

  \item  Donner alors l'équation du plan contenant $D$ et $D'$.
     
\end{enumerate}
\finenonce{004862}



\finexercice


\exercice{4863, quercia, 2010/03/17}
\enonce{004863}{Droites non coplanaires}

Soit $\cal E$ un espace affine de dimension 3, et $D,D',D''$ trois droites
parallèles à un même plan $\cal P$, mais deux à deux non coplanaires.

\begin{enumerate}
  \item Montrer que par tout point $A$ de $D$, il passe une unique droite $\Delta_A$
    rencontrant $D'$ et $D''$.
    
  \item Montrer que les droites $\Delta_A$ sont toutes parallèles à un même plan
    $\cal Q$.
    
\end{enumerate}
\finenonce{004863}



\finexercice


\exercice{4864, quercia, 2010/03/17}
\enonce{004864}{Droites concourantes}


Dans ${\cal E}_2$ muni d'un repère $(O,\vec i,\vec j)$, on considère les trois
droites : $\begin{cases}D   :&$ax   + by    = c   $\cr
                  D'  :&$a'x  + b'y   = c'  $\cr
                  D'' :&$a''x + b''y  = c''.$\cr\end{cases}$


Montrer que $D,D',D''$ sont parallèles ou concourantes si et seulement si
$\begin{vmatrix}a & b & c \cr a' & b' & c' \cr a'' & b'' & c'' \cr\end{vmatrix} = 0$.
\finenonce{004864}



\finexercice


\exercice{4865, quercia, 2010/03/17}
\enonce{004865}{Droites concourantes}
Soit $ABCD$ un parallélogramme, et $M \in (ABC)$.
On note $I,J$ les projections de $M$ sur $(AB)$ et $(CD)$ parallèlement
à $(AD)$, et $K,L$ les projections de $M$ sur $(AD)$ et $(BC)$ parallèlement
à $(AB)$.

Montrer que les droites $(IK)$, $(JL)$, $(BD)$ sont parallèles ou concourantes.


\finenonce{004865}



\finexercice


\exercice{4866, quercia, 2010/03/17}
\enonce{004866}{\'Equation d'une droite variable}

Soit $(O,\vec i,\vec j)$ un repère de ${\cal E}_2$, et
$A : \left(\begin{smallmatrix} 1 \cr 0 \cr\end{smallmatrix}\right)$,
$B : \left(\begin{smallmatrix} 0 \cr 1 \cr\end{smallmatrix}\right)$,
$C : \left(\begin{smallmatrix} 0 \cr 2 \cr\end{smallmatrix}\right)$.

Pour $m \in \R$, on construit les droites $D : y = mx$ et $D' : y = -mx$,
puis $M \in D \cap (AB)$, et $M' \in D' \cap (AC)$
(si possible).

Montrer que la droite $(MM')$ passe par un point fixe (= indépendant de $m$).


\finenonce{004866}



\finexercice


\exercice{4867, quercia, 2010/03/17}
\enonce{004867}{Dimensions}

Soient $\cal F$, $\cal G$, deux sous-espaces affines de dimension finie
d'un espace affine $\cal E$.
On note $\cal H$ le sous-espace affine engendré par ${\cal F} \cup {\cal G}$.
Déterminer $\dim({\cal H})$.

\finenonce{004867}



\finexercice


\exercice{4868, quercia, 2010/03/17}
\enonce{004868}{Dimensions}
Soient $\cal F$, $\cal G$, deux sous-espaces affines disjoints
de dimensions $f,g$ d'un espace affine $\cal E$ avec $f \le g$.

Montrer que ${\cal F} \ /\!/\ {\cal G}$ si et seulement s'il existe un
sous-espace affine ${\cal H}$ de dimension $g+1$ contenant $\cal F$ et $\cal G$.
\finenonce{004868}



\finexercice


\exercice{4869, quercia, 2010/03/17}
\enonce{004869}{Centrale PSI 1997}

Soit la famille de droites~:
$$(D_\lambda)\quad\begin{cases}x=\lambda+\lambda^2z\cr
                         y=\lambda^2+\lambda z.\cr\end{cases}$$
\begin{enumerate}
  \item En écrivant leurs équations sous la forme
    $\begin{cases}z=a\cr ux+vy+h=0\end{cases}$ montrer qu'il existe deux droites
    $\Delta_1$ et $\Delta_2$ horizontales coupant toutes les droites
    $D_\lambda$.
    
  \item Trouver les équations des plans passant par $M(\lambda,\lambda^2,0)$
    et contenant respectivement $\Delta_1$ et $\Delta_2$.
    
  \item Retrouver l'ensemble $(D_\lambda)$.
\end{enumerate}
\finenonce{004869}



\finexercice


\exercice{5505, rouget, 2010/07/15}
\enonce{005505}{*T}
Dans $\Rr^3$ affine, déterminer un repère de la droite $(D)$ $\left\{
\begin{array}{l}
x-y+2z+7=0\\
2x+2y+3z-5=0
\end{array}
\right.$.
\finenonce{005505}


\finexercice


\exercice{5506, rouget, 2010/07/15}
\enonce{005506}{*T}
Dans $\Rr^3$, déterminer l'intersection de $(D)$ $\left\{
\begin{array}{l}
x=2+\lambda\\
y=3-\lambda\\
z=7
\end{array}
\right.$ et $(P)~:~x+3y-5z+2=0$.
\finenonce{005506}


\finexercice


\exercice{5507, rouget, 2010/07/15}
\enonce{005507}{**}
Dans $\Rr^3$ affine, déterminer le réel $a$ pour que les droites $\left\{
\begin{array}{l}
x+2=-2z\\
y=3x+z
\end{array}
\right.$ et $\left\{
\begin{array}{l}
x+y+z=1\\
2x+y-z=a
\end{array}
\right.$ soient coplanaires, puis déterminer une équation du plan les contenant.
\finenonce{005507}


\finexercice


\exercice{5508, rouget, 2010/07/15}
\enonce{005508}{**T}
Dans $\Rr^3$, équation du plan $P$ parallèle à la droite $(Oy)$ et passant par $A(0,-1,2)$ et $B(-1,2,3)$.
\finenonce{005508}


\finexercice

\section{240.02 Applications affines}


\exercice{4870, quercia, 2010/03/17}
\enonce{004870}{$f^p = \mathrm{id} => f$ a un point fixe}

Soit $f : {\cal E} \to {\cal E}$ affine telle qu'il existe $p \in \N^*$ tel que
$f^p = \mathrm{id}_{\cal E}$. Montrer que $f$ admet au moins un point fixe.
\finenonce{004870}



\finexercice


\exercice{4871, quercia, 2010/03/17}
\enonce{004871}{$1$ non valeur propre $ \Rightarrow $ un pt fixe unique}

Soit ${\cal E}$ un espace affine de dimension finie et
$f : {\cal E} \to {\cal E}$ affine.
Montrer que $f$ admet un unique point fixe si et seulement si $1$ n'est pas
valeur propre de $\vec f$.
\finenonce{004871}



\finexercice


\exercice{4872, quercia, 2010/03/17}
\enonce{004872}{Expressions analytiques}

On fixe un repère ${\cal R} = (O,\vec e_1, \vec e_2, \vec e_3)$ d'un espace
affine de dimension 3.
Déterminer les expressions analytiques des applications suivantes :

\begin{enumerate}
  \item Symétrie de base le plan d'équation $x+2y+z = 1$ et de direction
    vect$(\vec e_1 + \vec e_2 + \vec e_3)$.
    


  \item Symétrie de base la droite d'équations
    $\begin{cases} x+y+1 = 0 \cr 2y+z+2 = 0,\end{cases}$ \par
    de direction le plan vectoriel
    d'équation $3x + 3y - 2z = 0$.
    
\end{enumerate}
\finenonce{004872}



\finexercice


\exercice{4873, quercia, 2010/03/17}
\enonce{004873}{Expression analytique}
On fixe un repère ${\cal R} = (O,\vec e_1, \vec e_2, \vec e_3)$ d'un espace
affine de dimension 3.
Reconaître l' application ayant l'expression  analytique  suivante :

   $$\begin{cases}x' =  3x + 4y + 2z - 4 \cr
            y' = -2x - 3y - 2z + 4 \cr
            z' =  4x + 8y + 5z - 8.\cr\end{cases}$$

   (chercher les points fixes de $f$ et étudier $\overrightarrow{MM'}$)
   
\finenonce{004873}



\finexercice


\exercice{4874, quercia, 2010/03/17}
\enonce{004874}{Permutation circulaire de 4 points}

Dans un espace affine ${\cal E}$, on considère quatre points $A,B,C,D$.
\'Etudier l'existence d'une application affine $f$ telle que
$f(A) = B$, $f(B) = C$, $f(C) = D$, $f(D) = A$.
\finenonce{004874}



\finexercice


\exercice{4875, quercia, 2010/03/17}
\enonce{004875}{$f^3 = \mathrm{id}$}

Soit ${\cal P}$ un plan, et $f : {\cal P} \to {\cal P}$ une application
affine telle que $f^3 = \mathrm{id}$, avec $f \ne \mathrm{id}$.

\begin{enumerate}
  \item Montrer que si $A \ne f(A)$, alors $A,f(A), f^2(A)$ sont non alignés.
  \item En déduire que $f$ est le produit de deux symétries.
\end{enumerate}
\finenonce{004875}



\finexercice


\exercice{4876, quercia, 2010/03/17}
\enonce{004876}{Produit d'affinités}

Soit $\cal P$ un plan, $\cal D$ une droite de $\cal P$, et $f,g$ deux affinités
de base $\cal D$, de directions $\vec{\Delta}$, $\vec{\Delta'}$ et de rapports
$\lambda$, $\mu$.

\'Etudier la nature de $f\circ g$.
\finenonce{004876}



\finexercice


\exercice{4877, quercia, 2010/03/17}
\enonce{004877}{Barycentre de projections}

Soient $\pi,\pi'$ deux projections dans un espace affine ${\cal E}$ ayant même
direction $\vec{\cal F}$.

Pour $\lambda \in \R$, on note $\pi_\lambda$ l'application~:
$M \longmapsto \text{Bar}(\pi(M):\lambda, \pi'(M):1-\lambda)$.

Montrer que $\pi_\lambda$ est encore une projection de direction $\vec{\cal F}$.
\finenonce{004877}



\finexercice


\exercice{4878, quercia, 2010/03/17}
\enonce{004878}{Symétrie-translation}

Soit $f : {\cal E} \to {\cal E}$ affine.
On dit que $f$ est une {\it symétrie-translation\/}
s'il existe une symétrie $s$ et une translation $t$ telles que
$f = s\circ t = t\circ s$.

\begin{enumerate}
  \item Soient $s$ une symétrie de base $\cal B$ de direction $\vec {\cal F}$,
    et $t$ une translation de vecteur $\vec u$.

    Montrer que $s\circ t = t\circ s \iff \vec u \in \vec {\cal B}$.

  \item Soit $f$ une symétrie-translation. Montrer que le couple $(s,t)$ tel que
    $f = s\circ t = t\circ s$ est unique.

  \item Soit $f$ affine quelconque. Montrer que $f$ est une symétrie-translation si
    et seulement si $f\circ f$ est une translation.

  \item En déduire que le produit d'une symétrie par une translation quelconques
    est une symétrie-translation.

  \item AN : décomposer l'application $f$ d'expression analytique dans un
    repère ${\cal R} = (O,\vec e_1, \vec e_2, \vec e_3)$ :
    $$\begin{cases}x' = (x-2y-2z+1)/3 \cr y' = (-2x+y-2z+2)/3 \cr z' = (-2x-2y+z-1)/3.\cr\end{cases}$$

\end{enumerate}
\finenonce{004878}



\finexercice


\exercice{4879, quercia, 2010/03/17}
\enonce{004879}{Transitivité des homothéties-translations}

Dans un espace affine ${\cal E}$ on donne quatre points $P,Q,P',Q'$ avec
$P \ne Q$.
Existe-t-il une homothétie-translation $f$ telle que $f(P) = P'$ et
$f(Q) = Q'$ ?
\finenonce{004879}



\finexercice


\exercice{4880, quercia, 2010/03/17}
\enonce{004880}{Usage d'applications affines}

On considère dans l'espace deux plans parallèles distincts $\cal P$, ${\cal P}'$,
$A,B,C \in \cal P$, $O \notin \cal P$,
et on construit les points suivants :

-- $A',B',C'$ : les intersections avec ${\cal P}'$ des droites $(OA)$, $(OB)$,
   $(OC)$.

-- $\alpha,\beta,\gamma$ : les milieux des segments $[B,C]$, $[C,A]$, $[A,B]$.

Montrer que les droites $(A'\alpha)$, $(B'\beta)$, $(C'\gamma)$ sont parallèles
ou concourantes.
\finenonce{004880}



\finexercice


\exercice{4881, quercia, 2010/03/17}
\enonce{004881}{Usage d'applications affines}
Soient $A_1,\dots,A_n$ $n$ points de ${\cal E}$.

\'Etudier l'existence de points $B_1,\dots,B_n$ tels que
$A_i = \text{mil}(B_i,B_{i+1})$ ($A_n = \text{mil}(B_n,B_{1}))$.
\finenonce{004881}



\finexercice


\exercice{4882, quercia, 2010/03/17}
\enonce{004882}{Projection stéréographique}

Dans l'espace, on considère un point $O$ et un plan $\cal P$ ne passant pas
par $O$. On définit l'application $f$: $M \longmapsto M'$ où $M'$ est le
point intersection de $\cal P$ et $(OM)$.
(Projection stéréographique sur $\cal P$ de pôle $O$)

\begin{enumerate}
  \item Est-ce que $f$ est affine ?
  \item Etudier l'image par $f$ d'une droite, d'un plan, d'une partie convexe.

\end{enumerate}
\finenonce{004882}



\finexercice


\exercice{4883, quercia, 2010/03/17}
\enonce{004883}{ Caractérisation des produits de symétries}

Soit ${\cal E}$ un espace affine de dimension finie et
$f : {\cal E} \to {\cal E}$ affine.

Montrer que $f$ est un produit de symétries si et seulement si
$\det(\vec f) = \pm1$.
\finenonce{004883}



\finexercice


\exercice{4885, quercia, 2010/03/17}
\enonce{004885}{Points dans l'espace}

Dans l'espace, les droites $(AA')$, $(BB')$, $(CC')$ sont concourantes en $O$,
$O \notin (ABC)$ et $A,B,C$ non alignés.
Soient $G,G'$ les isobarycentres des triangles $ABC$, $A'B'C'$.
CNS pour que $O,G,G'$ soient alignés ?

\finenonce{004885}



\finexercice


\exercice{4886, quercia, 2010/03/17}
\enonce{004886}{Polygone des milieux}

Soit $P = A_1A_2\dots A_n$ un polygone à $n$ sommets :
on lui associe le polygone $P' = A'_1A'_2\dots A'_{n-1}A'_n$
où $A'_i$ est le milieu de $A_i$ et $A_{i+1}$ ($A_{n+1} = A_1$).

On définit alors une suite de polygones par récurrence :
$\begin{cases} P_0 = P\cr P_{k+1} = (P_k)'.\cr \end{cases}$

Montrer que chaque sommet de $P_k$ converge vers le centre de gravité de $P_0$
lorsque $k$ tend vers l'infini.

(\'Ecrire un sommet de $P_k$ comme barycentre de $A_1,\dots,A_n$)

\finenonce{004886}



\finexercice


\exercice{4887, quercia, 2010/03/17}
\enonce{004887}{Isobarycentre de tous les points sauf un}

Soit $P = A_1A_2\dots A_n$ un polygone à $n$ sommets :
on lui associe le polygone $P' = A'_1A'_2\dots A'_n$
où $A'_i$ est l'isobarycentre de tous les sommets sauf $A_i$.

On définit alors une suite de polygones par récurrence :
$\begin{cases} P_0 = P\cr P_{k+1} = (P_k)'.\cr \end{cases}$

Montrer que chaque sommet de $P_k$ converge vers le centre de gravité de $P_0$
lorsque $k$ tend vers l'infini.

\finenonce{004887}



\finexercice


\exercice{4888, quercia, 2010/03/17}
\enonce{004888}{Suite récurente}

Soient $A_0,A_1,A_2$ trois points donnés. On considère la suite $(A_k)$
de points vérifiant la relation de récurrence~:
$$\forall\ k\ge3,\ A_k = \text{Bar}( A_{k-1}:1, A_{k-2},1, A_{k-3}:2 ).$$
\'Etudier la convergence de cette suite.



\finenonce{004888}



\finexercice


\exercice{4889, quercia, 2010/03/17}
\enonce{004889}{Théorème de Ménélaüs}

Soit $ABC$ un triangle et trois points
$P \in (AB)$, $Q \in (BC)$, $R \in (CA)$, distincts de $A,B,C$.

\begin{enumerate}
  \item Montrer que $P,Q,R$ sont alignés si et seulement si
    $\frac {\overline{PA}}{\overline{PB}} \times
     \frac {\overline{QB}}{\overline{QC}} \times
     \frac {\overline{RC}}{\overline{RA}} = 1$.

  \item Dans ce cas, montrer que $P' = \text{mil}(P,C)$, $Q' = \text{mil}(Q,A)$,
    et $R' = \text{mil}(R,B)$ sont aussi alignés.
\end{enumerate}
\finenonce{004889}



\finexercice


\exercice{4890, quercia, 2010/03/17}
\enonce{004890}{Théorème de Céva}

Soit $ABC$ un triangle et trois points
$P \in (AB)$, $Q \in (BC)$, $R \in (CA)$, distincts de $A,B,C$.
Démontrer que les droites $(AQ)$, $(BR)$, $(CP)$ sont parallèles ou
concourantes si et seulement
si ${\frac {\overline{PA}}{\overline{PB}} \times
    \frac {\overline{QB}}{\overline{QC}} \times
    \frac {\overline{RC}}{\overline{RA}} = -1}$.
\finenonce{004890}



\finexercice


\exercice{4891, quercia, 2010/03/17}
\enonce{004891}{Droites parallèles}

Soient $ABC$ et $A'B'C'$ deux triangles tels que les parallèles à $(AB)$,
$(BC)$, $(CA)$ passant respectivement par $C',A',B'$ soient concourantes.


Montrer qu'il en est de même pour les parallèles à $(A'B')$, $(B'C')$, $(C'A')$
passant par $C,A,B$.

\finenonce{004891}



\finexercice


\exercice{4892, quercia, 2010/03/17}
\enonce{004892}{Points aux tiers des côtés}

Soit $ABC$ un triangle, $A_1 = \text{Bar}(B:2,C:1)$, $B_1 = \text{Bar}(C:2,A:1)$ et
$C_1 = \text{Bar}(A:2,B:1)$.

On note $A_2,B_2,C_2$ les points d'intersection des droites $(AA_1)$, $(BB_1)$,
et $(CC_1)$.

\begin{enumerate}
  \item  Montrer que $A_2$ est le milieu de $[B,B_2]$.
  \item  Comparer les surfaces des triangles $ABC$ et $A_2B_2C_2$.
\end{enumerate}
\finenonce{004892}



\finexercice


\exercice{4893, quercia, 2010/03/17}
\enonce{004893}{Symétriques d'un point par rapport aux milieux des cotés}

Soit un triangle $ABC$, $A',B',C'$, les milieux des côtés, et $M$ un point
du plan $(ABC)$ de coordonnées barycentriques $(\alpha,\beta,\gamma)$.

\begin{enumerate}
  \item  Chercher les coordonnées barycentriques de $P,Q,R$ symétriques de $M$
     par rapport aux points $A',B',C'$.
  \item  Montrer que les droites $(AP)$, $(BQ)$, $(CR)$ sont concourantes en un
     point $N$.
     
  \item  Montrer que $N$ est le milieu de $[A,P]$, $[B,Q]$, $[C,R]$.
  \item  Reconnaître l'application $M \mapsto N$.
     
\end{enumerate}
\finenonce{004893}



\finexercice


\exercice{4894, quercia, 2010/03/17}
\enonce{004894}{Parallélogrammes}

Dans le plan, on considère :

--  trois points non alignés $A,B,C$.\par
--  trois points alignés $P,Q,R$ avec $P \in (AB)$, $Q \in (AC)$, $R \in (BC)$.

On construit les points $I,J,K$ de sorte que $BPIR$, $APJQ$, $CQKR$ soient
des parallélogrammes.

Montrer que $I,J,K$ sont alignés.

\finenonce{004894}



\finexercice


\exercice{4895, quercia, 2010/03/17}
\enonce{004895}{Projections en cascade}

Soient $A,B,C$ trois points non alignés et $M_1 \in (AB)$.
On construit les points $M_2,M_3,M_4$ de la manière suivante~:

-- $M_2$ est le projeté de $M_1$ sur $(BC)$ parallèlement à $(AC)$.\par
-- $M_3$ est le projeté de $M_2$ sur $(AC)$ parallèlement à $(AB)$.\par
-- $M_4$ est le projeté de $M_3$ sur $(AB)$ parallèlement à $(BC)$.\par

On recommence ensuite les mêmes constructions à partir de $M_4$, ce qui
donne les points $M_5, M_6, M_7$.

Montrer que $M_7 = M_1$.
\finenonce{004895}



\finexercice


\exercice{4896, quercia, 2010/03/17}
\enonce{004896}{Caractérisation du barycentre par les surfaces}

Soit $ABC$ un triangle, et $M \in (ABC)$. On note $\alpha,\beta,\gamma$ les
aires des triangles $MBC$, $MCA$, $MAB$.

\begin{enumerate}
  \item On suppose que $M$ est dans l'enveloppe convexe de $\{A,B,C\}$.\par
    Montrer que : $(\alpha = \beta = \gamma) \iff
                   M = \text{Bar}(A:1,\ B:1,\ C:1)$.
  \item Quels sont tous les points du plan $(ABC)$ tels que
    $\alpha = \beta = \gamma$ ?
\end{enumerate}
\finenonce{004896}



\finexercice


\exercice{4897, quercia, 2010/03/17}
\enonce{004897}{Coord. barycentriques du centre du cercle circonscrit}

Soit $ABC$ un triangle.
On note : $a = BC$, $b = CA$, $c = AB$,
$\alpha \equiv (\overline{ \vec{AB}, \vec{AC}})$,
$\beta  \equiv (\overline{ \vec{BC}, \vec{BA}})$,
$\gamma \equiv (\overline{ \vec{CA}, \vec{CB}})$.


\begin{enumerate}
  \item   Montrer que pour tout point $M$ du cercle ($ABC$), on a :

      $$a\cos\alpha MA^2 + b\cos\beta MB^2 + c\cos\gamma MC^2 = abc.$$

  \item   En déduire les coordonnées barycentriques du centre du cercle $(ABC)$.
\end{enumerate}
\finenonce{004897}



\finexercice


\exercice{4898, quercia, 2010/03/17}
\enonce{004898}{Cercle inscrit}

Soit $ABC$ un triangle.
On note : $a = BC$, $b = CA$, $c = AB$,

\begin{enumerate}
  \item Soit $A'$ le pied de la bissectrice intérieure issue de $A$.
    Montrer que $\frac{A'B}{A'C} = \frac cb$.
    

  \item En déduire les coordonnées barycentriques de $I$, centre du cercle inscrit.
    
\end{enumerate}
\finenonce{004898}



\finexercice

\section{240.03 Barycentre}


\exercice{4884, quercia, 2010/03/17}
\enonce{004884}{\'Equation barycentrique d'une droite}


Soit $(A,B,C)$ une base affine de ${\cal E}_2$, et $M,M',M''$ trois points de
coordonnées barycentriques
$(\alpha, \beta, \gamma)$,
$(\alpha', \beta', \gamma')$,
$(\alpha'', \beta'', \gamma'')$.

Montrer que $M,M',M''$ sont alignés si et seulement si
$\begin{vmatrix}\alpha& \beta& \gamma \cr
          \alpha'& \beta'& \gamma' \cr
          \alpha''& \beta''& \gamma''\cr\end{vmatrix} = 0$.
\finenonce{004884}



\finexercice


\exercice{4885, quercia, 2010/03/17}
\enonce{004885}{Points dans l'espace}

Dans l'espace, les droites $(AA')$, $(BB')$, $(CC')$ sont concourantes en $O$,
$O \notin (ABC)$ et $A,B,C$ non alignés.
Soient $G,G'$ les isobarycentres des triangles $ABC$, $A'B'C'$.
CNS pour que $O,G,G'$ soient alignés ?

\finenonce{004885}



\finexercice


\exercice{4886, quercia, 2010/03/17}
\enonce{004886}{Polygone des milieux}

Soit $P = A_1A_2\dots A_n$ un polygone à $n$ sommets :
on lui associe le polygone $P' = A'_1A'_2\dots A'_{n-1}A'_n$
où $A'_i$ est le milieu de $A_i$ et $A_{i+1}$ ($A_{n+1} = A_1$).

On définit alors une suite de polygones par récurrence :
$\begin{cases} P_0 = P\cr P_{k+1} = (P_k)'.\cr \end{cases}$

Montrer que chaque sommet de $P_k$ converge vers le centre de gravité de $P_0$
lorsque $k$ tend vers l'infini.

(\'Ecrire un sommet de $P_k$ comme barycentre de $A_1,\dots,A_n$)

\finenonce{004886}



\finexercice


\exercice{4887, quercia, 2010/03/17}
\enonce{004887}{Isobarycentre de tous les points sauf un}

Soit $P = A_1A_2\dots A_n$ un polygone à $n$ sommets :
on lui associe le polygone $P' = A'_1A'_2\dots A'_n$
où $A'_i$ est l'isobarycentre de tous les sommets sauf $A_i$.

On définit alors une suite de polygones par récurrence :
$\begin{cases} P_0 = P\cr P_{k+1} = (P_k)'.\cr \end{cases}$

Montrer que chaque sommet de $P_k$ converge vers le centre de gravité de $P_0$
lorsque $k$ tend vers l'infini.

\finenonce{004887}



\finexercice


\exercice{4888, quercia, 2010/03/17}
\enonce{004888}{Suite récurente}

Soient $A_0,A_1,A_2$ trois points donnés. On considère la suite $(A_k)$
de points vérifiant la relation de récurrence~:
$$\forall\ k\ge3,\ A_k = \text{Bar}( A_{k-1}:1, A_{k-2},1, A_{k-3}:2 ).$$
\'Etudier la convergence de cette suite.



\finenonce{004888}



\finexercice

\section{240.04 Propriétés des triangles}


\exercice{4889, quercia, 2010/03/17}
\enonce{004889}{Théorème de Ménélaüs}

Soit $ABC$ un triangle et trois points
$P \in (AB)$, $Q \in (BC)$, $R \in (CA)$, distincts de $A,B,C$.

\begin{enumerate}
  \item Montrer que $P,Q,R$ sont alignés si et seulement si
    $\frac {\overline{PA}}{\overline{PB}} \times
     \frac {\overline{QB}}{\overline{QC}} \times
     \frac {\overline{RC}}{\overline{RA}} = 1$.

  \item Dans ce cas, montrer que $P' = \text{mil}(P,C)$, $Q' = \text{mil}(Q,A)$,
    et $R' = \text{mil}(R,B)$ sont aussi alignés.
\end{enumerate}
\finenonce{004889}



\finexercice


\exercice{4890, quercia, 2010/03/17}
\enonce{004890}{Théorème de Céva}

Soit $ABC$ un triangle et trois points
$P \in (AB)$, $Q \in (BC)$, $R \in (CA)$, distincts de $A,B,C$.
Démontrer que les droites $(AQ)$, $(BR)$, $(CP)$ sont parallèles ou
concourantes si et seulement
si ${\frac {\overline{PA}}{\overline{PB}} \times
    \frac {\overline{QB}}{\overline{QC}} \times
    \frac {\overline{RC}}{\overline{RA}} = -1}$.
\finenonce{004890}



\finexercice


\exercice{4891, quercia, 2010/03/17}
\enonce{004891}{Droites parallèles}

Soient $ABC$ et $A'B'C'$ deux triangles tels que les parallèles à $(AB)$,
$(BC)$, $(CA)$ passant respectivement par $C',A',B'$ soient concourantes.


Montrer qu'il en est de même pour les parallèles à $(A'B')$, $(B'C')$, $(C'A')$
passant par $C,A,B$.

\finenonce{004891}



\finexercice


\exercice{4892, quercia, 2010/03/17}
\enonce{004892}{Points aux tiers des côtés}

Soit $ABC$ un triangle, $A_1 = \text{Bar}(B:2,C:1)$, $B_1 = \text{Bar}(C:2,A:1)$ et
$C_1 = \text{Bar}(A:2,B:1)$.

On note $A_2,B_2,C_2$ les points d'intersection des droites $(AA_1)$, $(BB_1)$,
et $(CC_1)$.

\begin{enumerate}
  \item  Montrer que $A_2$ est le milieu de $[B,B_2]$.
  \item  Comparer les surfaces des triangles $ABC$ et $A_2B_2C_2$.
\end{enumerate}
\finenonce{004892}



\finexercice


\exercice{4893, quercia, 2010/03/17}
\enonce{004893}{Symétriques d'un point par rapport aux milieux des cotés}

Soit un triangle $ABC$, $A',B',C'$, les milieux des côtés, et $M$ un point
du plan $(ABC)$ de coordonnées barycentriques $(\alpha,\beta,\gamma)$.

\begin{enumerate}
  \item  Chercher les coordonnées barycentriques de $P,Q,R$ symétriques de $M$
     par rapport aux points $A',B',C'$.
  \item  Montrer que les droites $(AP)$, $(BQ)$, $(CR)$ sont concourantes en un
     point $N$.
     
  \item  Montrer que $N$ est le milieu de $[A,P]$, $[B,Q]$, $[C,R]$.
  \item  Reconnaître l'application $M \mapsto N$.
     
\end{enumerate}
\finenonce{004893}



\finexercice


\exercice{4894, quercia, 2010/03/17}
\enonce{004894}{Parallélogrammes}

Dans le plan, on considère :

--  trois points non alignés $A,B,C$.\par
--  trois points alignés $P,Q,R$ avec $P \in (AB)$, $Q \in (AC)$, $R \in (BC)$.

On construit les points $I,J,K$ de sorte que $BPIR$, $APJQ$, $CQKR$ soient
des parallélogrammes.

Montrer que $I,J,K$ sont alignés.

\finenonce{004894}



\finexercice


\exercice{4895, quercia, 2010/03/17}
\enonce{004895}{Projections en cascade}

Soient $A,B,C$ trois points non alignés et $M_1 \in (AB)$.
On construit les points $M_2,M_3,M_4$ de la manière suivante~:

-- $M_2$ est le projeté de $M_1$ sur $(BC)$ parallèlement à $(AC)$.\par
-- $M_3$ est le projeté de $M_2$ sur $(AC)$ parallèlement à $(AB)$.\par
-- $M_4$ est le projeté de $M_3$ sur $(AB)$ parallèlement à $(BC)$.\par

On recommence ensuite les mêmes constructions à partir de $M_4$, ce qui
donne les points $M_5, M_6, M_7$.

Montrer que $M_7 = M_1$.
\finenonce{004895}



\finexercice


\exercice{4896, quercia, 2010/03/17}
\enonce{004896}{Caractérisation du barycentre par les surfaces}

Soit $ABC$ un triangle, et $M \in (ABC)$. On note $\alpha,\beta,\gamma$ les
aires des triangles $MBC$, $MCA$, $MAB$.

\begin{enumerate}
  \item On suppose que $M$ est dans l'enveloppe convexe de $\{A,B,C\}$.\par
    Montrer que : $(\alpha = \beta = \gamma) \iff
                   M = \text{Bar}(A:1,\ B:1,\ C:1)$.
  \item Quels sont tous les points du plan $(ABC)$ tels que
    $\alpha = \beta = \gamma$ ?
\end{enumerate}
\finenonce{004896}



\finexercice


\exercice{4897, quercia, 2010/03/17}
\enonce{004897}{Coord. barycentriques du centre du cercle circonscrit}

Soit $ABC$ un triangle.
On note : $a = BC$, $b = CA$, $c = AB$,
$\alpha \equiv (\overline{ \vec{AB}, \vec{AC}})$,
$\beta  \equiv (\overline{ \vec{BC}, \vec{BA}})$,
$\gamma \equiv (\overline{ \vec{CA}, \vec{CB}})$.


\begin{enumerate}
  \item   Montrer que pour tout point $M$ du cercle ($ABC$), on a :

      $$a\cos\alpha MA^2 + b\cos\beta MB^2 + c\cos\gamma MC^2 = abc.$$

  \item   En déduire les coordonnées barycentriques du centre du cercle $(ABC)$.
\end{enumerate}
\finenonce{004897}



\finexercice


\exercice{4898, quercia, 2010/03/17}
\enonce{004898}{Cercle inscrit}

Soit $ABC$ un triangle.
On note : $a = BC$, $b = CA$, $c = AB$,

\begin{enumerate}
  \item Soit $A'$ le pied de la bissectrice intérieure issue de $A$.
    Montrer que $\frac{A'B}{A'C} = \frac cb$.
    

  \item En déduire les coordonnées barycentriques de $I$, centre du cercle inscrit.
    
\end{enumerate}
\finenonce{004898}



\finexercice


\exercice{4899, quercia, 2010/03/17}
\enonce{004899}{Orthocentre}

Soit $ABC$ un triangle.
On note :
$\alpha \equiv (\overline{ \vec{AB}, \vec{AC}})$,
$\beta  \equiv (\overline{ \vec{BC}, \vec{BA}})$,
$\gamma \equiv (\overline{ \vec{CA}, \vec{CB}})$.

\begin{enumerate}
  \item Soit $A'$ le pied de la hauteur issue de $A$.
    Calculer $\frac{\overline{A'B}}{\overline{A'C}}$.
    
  \item En déduire les coordonnées barycentriques de l'orthocentre $H$.
    

\end{enumerate}
\finenonce{004899}



\finexercice


\exercice{5206, rouget, 2010/06/30}
\enonce{005206}{**}
Montrer qu'il n'existe pas de triangle équilatéral dont les sommets appartiennent aux points d'intersection des lignes d'une feuille blanche quadrillée usuelle.
\finenonce{005206}


\finexercice


\exercice{7059, megy, 2017/01/11}

\enonce{007059}{}
% collège
% triangle rectangle, cercle circonscrit
Soit $ABC$ un triangle et $\mathcal C$ son cercle circonscrit, de centre $O$.
Soit $A’$ le point diamétralement opposé à $A$ sur le cercle $\mathcal C$.
La hauteur $(AH)$ issue de $A$ du triangle $ABC$ recoupe le cercle $\mathcal C$ au point $D$.

Montrer que la droite $(DA’)$ est parallèle à $(BC)$.
\finenonce{007059}
\finexercice  


\exercice{7060, megy, 2017/01/11}

\enonce{007060}{} 
% collège
% orthocentre, triangle rectangle inscrit 
% cercle circonscrit, hauteurs

Soit $[AB]$ un segment et $M$, $N$ deux points appartenant au cercle  $\mathcal C$  de diamètre $[AB]$. On suppose que les droites $(MB)$ et $(AN)$ (respectivement $(NB)$ et $(AM)$ ) s'intersectent en $P$ (respectivement en $Q$). Déterminer l'angle formé par  les droites $(AB)$ et $(PQ)$.
\finenonce{007060}
\finexercice  


\exercice{7061, megy, 2017/01/11}

\enonce{007061}{} 
% collège
% orthocentre, triangle rectangle inscrit 
% cercle circonscrit, hauteurs
Soit $ABC$ un triangle. Le cercle $\mathcal C$ (resp. $\mathcal C'$) de diamètre $[BC]$  (resp. $[CA]$) coupe la droite $(CA)$ (resp. la droite $(BC)$) en $P$ (resp. $Q$). Les cercles $\mathcal C$  et $\mathcal C'$ se recoupent en un second point $R$. Montrer que $(CR)$, $(BP)$ et $(AQ)$ sont concourantes.

\begin{center}
\includegraphics{../images/img007061-1}
\end{center}

\finenonce{007061}
\finexercice  


\exercice{7062, megy, 2017/01/11}

\enonce{007062}{} 
% orthocentre, triangle rectangle inscrit 
% cercle circonscrit, hauteurs
On donne  un cercle $\mathcal C$, un diamètre $[AB]$ et un troisième point $M$ du cercle. L'objectif est de construire le projeté orthogonal de $M$ sur $(AB)$ à la règle seule.
\begin{enumerate}
\item Montrer qu'il suffit de construire une droite orthogonale à $(AB)$ coupant le cercle en deux points.
\item Construire une telle droite.
\end{enumerate}
\finenonce{007062}
\finexercice  


\exercice{7063, megy, 2017/01/11}

\enonce{007063}{}
%  triangle rectangle, cercle
On donne deux cercles $\mathcal C$ et $\mathcal C'$ de rayons distincts, de centres $O$ et $O'$, tangents extérieurement en  un point $A$. On admet qu'il existe trois tangentes communes  à $\mathcal C$ et $\mathcal C'$ : la tangente commune en $A$, qui est directement constructible, et deux autres droites. L'objectif de l'exercice est de tracer ces deux dernières tangentes. 
\begin{enumerate}

\item Considérons donc une droite tangente à $\mathcal C$ en $B$ et à $\mathcal C'$ en $C$, avec $B\neq C$. La tangente commune en $A$ aux deux cercles coupe $(BC)$ en $I$. Montrer que $I$ est le milieu de $[BC]$ et que $ABC$ est rectangle en $A$.
\item Finir l'exercice (c'est-à-dire construire $B$ et $C$) de l'une des deux façons suivantes:
\begin{enumerate}
\item Soit $D$ tel que $ABDC$ soit un rectangle. Quels sont les points d'intersection entre $(DB)$, $(DC)$ et $(OO')$ ? En déduire une construction du point $D$.
\item Montrer que $OIO'$ est rectangle en $I$ et en déduire une construction du point $I$.
\end{enumerate}
\end{enumerate}
\finenonce{007063}
\finexercice  


\exercice{7064, megy, 2017/01/11}

\enonce{007064}{Théorème de Varignon}
Soit $ABCD$ un quadrilatère convexe, et $I$, $J$, $K$, $L$ les milieux de ses côtés. Montrer que $IJKL$ est un parallélogramme soit en utilisant des barycentres, soit le théorème de Thalès. Montrer que l'aire de $ABCD$ est le double de celle de $IJKL$ de deux façons différentes.
\finenonce{007064}


\finexercice  


\exercice{7065, megy, 2017/01/11}

\enonce{007065}{Quadrilatères orthodiagonaux}
%  aire, hauteur,
Un quadrilatère est dit \emph{orthodiagonal} si ses diagonales sont perpendiculaires.

Soit $ABCD$ un quadrilatère orthodiagonal non croisé. Montrer que son aire vaut $\frac12 AC\cdot BD$. 

\finenonce{007065}
     
\finexercice  


\exercice{7066, megy, 2017/01/11}

\enonce{007066}{Somme de distances}

Soit $ABC$ un triangle équilatéral. Pour tout point $M$ à l'intérieur du triangle, on note 
\[d = dist(M,[AB]) + dist(M,[BC]) + dist(M,[AC])\]
 la somme des distances de $M$ aux trois côtés. Montrer que $d$ ne dépend en fait pas du point $M$.
\finenonce{007066} 

\finexercice   


\exercice{7067, megy, 2017/01/11}

\enonce{007067}{}
% collège
% triangles isocèles, médiatrices, cercle circonscrit
Trois cercles sont tangents extérieurement deux à deux. Montrer que les tangentes communes sont concourantes.
\finenonce{007067}
\finexercice  


\exercice{7068, megy, 2017/01/11}

\enonce{007068}{*}
% joli
%  médiane, centre de gravité, bissectrice
Soit $ABC$ un triangle avec $AB = 2 BC$ et $M$ un point de $[AC]$ tel que $AM = 2 MC$. Comparer les angles $\widehat{ABM}$ et $\widehat{MBC}$.

\begin{center}
\includegraphics{../images/img007068-1}
\end{center}

\finenonce{007068}
\finexercice  


\exercice{7069, megy, 2017/01/11}

\enonce{007069}{}
% triangles rectangles
À l'extérieur d'un triangle $BOA$ on construit deux triangles rectangles :
\begin{itemize}
\item Le triangle $OAC$, ayant pour hypoténuse le côté $[OA]$, tel que le sommet $C$ de l'angle droit soit situé sur la bissectrice extérieure de $OAB$.
\item Le triangle $OBE$, ayant pour hypoténuse le côté $[OB]$, tel que le sommet $E$ de l'angle droit soit situé sur la bissectrice extérieure de $OBA$.
\end{itemize}
Que dire de $[EC]$ et de sa longueur ?
\finenonce{007069}
    
\finexercice  


\exercice{7070, megy, 2017/01/11}

\enonce{007070}{}
% pythagore, trinôme
On considère un carré $ABCD$, et un cercle $\mathcal C$ passant par $A$ et $B$ et tangent à $[CD]$.
\begin{enumerate}
\item Montrer que le point de tangence est le milieu de $[CD]$.
\item Montrer que si le rayon vaut $r=10$ alors $AB=16$, et réciproquement.
%\item Montrer que $AB = \frac{8}{5} r$.
\end{enumerate}
\finenonce{007070}
    
\finexercice  


\exercice{7071, megy, 2017/01/11}

\enonce{007071}{}
% triangle isocèle, cercles
Soit $A$  un point quelconque du diamètre d'un cercle $\mathcal C$ et $B$ l'extrémité d'un rayon perpendiculaire à ce diamètre. On mène une droite $(BA)$ qui coupe le cercle en $P$, puis la tangente au point $P$ qui coupe en $C$ le diamètre prolongé. Démontrer que $CA = CP$.
\finenonce{007071}
\finexercice  


\exercice{7072, megy, 2017/01/11}

\enonce{007072}{}
% cercle inscrit, somme des angles d'un triangle
Soit $ABC$ un triangle et $I$ le centre de son cercle inscrit, dont on note $r$ le rayon. Montrer qu'un des sommets du triangle est à distance $\geq 2r$ de $I$, et qu'un autre est à distance $\leq 2r$.
\finenonce{007072}

\finexercice  


\exercice{7073, megy, 2017/01/11}

\enonce{007073}{Quadrilatères tangentiels}
%  cercle inscrit, bissectrice, triangles isocèles
Un quadrilatère convexe est dit \emph{tangentiel} ou \emph{circonscriptible} s'il possède un cercle inscrit, c'est-à-dire si ses quatre côtés sont tangents à un même cercle.

\begin{enumerate}
\item Montrer qu'un quadrilatère est tangentiel ssi ses bissectrices intérieures sont concourantes.
\item Montrer le théorème de Pitot (1725) : dans un quadrilatère tangentiel, la somme des longueurs de deux côtés opposés est égale à la somme des deux autres. Réciproque ?
\item Montrer qu'un cerf-volant isocèle (ou rhomboïde) est tangentiel.
\end{enumerate}
\finenonce{007073}
      
\finexercice  


\exercice{7074, megy, 2017/01/11}

\enonce{007074}{Théorème des trois tangentes}
% triangle isocèle, tangentes
% facile, intersection de deux tangentes => équidistance
Soit $ABC$ un triangle. Le cercle exinscrit dans l'angle en $A$ touche les côtés $[BC]$, $[AC)$ et $[AB)$ en $P$, $Q$ et $R$. Montrer que la somme $AR+AQ$ est égale au périmètre du triangle $ABC$.
\finenonce{007074}


\finexercice 


\exercice{7075, megy, 2017/01/11}

\enonce{007075}{Aire, périmètre et cercle inscrit}

Soit $ABC$ un triangle dont on note $a$, $b$ et $c$ les longueurs des côtés.
\begin{enumerate}
\item Exprimer l'aire $S$ du triangle en fonction du périmètre $a+b+c=2p$ et du rayon $r$ du cercle inscrit.
\item Exprimer également $S$ en fonction de $a$ et du rayon $r_A$ du cercle exinscrit en $A$.
\item En déduire $\frac{1}{r} = \frac{1}{r_A}+\frac{1}{r_B}+\frac{1}{r_C}$.
\end{enumerate}
\finenonce{007075}
   
\finexercice  


\exercice{7076, megy, 2017/01/11}

\enonce{007076}{}
% Euclide, I prop. 17
% cercles inscrits et exinscrits
 On donne un cercle  $\mathcal C$ (de centre $O$), un point $M$ à l'extérieur du cercle, les deux tangentes $\mathcal D$ et $\mathcal D'$ à $\mathcal C$ passant par $M$. On notera $A$ et $B$ les points de tangence. 

Le cercle $\mathcal C$ coupe $(MO)$ en deux points $P$ et $Q$. D'autre part, soit $H$ l'intersection de la corde $[AB]$ avec $(OM)$. Montrer que les cercles de centres $P$ et $Q$ et passant par $H$ sont tangents à $\mathcal D$ et $\mathcal D'$. 


\finenonce{007076}
\finexercice  

\section{240.99 Autres}

\section{241.00 Isométrie vectorielle}


\exercice{1982, ridde, 1999/11/01}

\enonce{001982}{}
Compl\'eter $x_{1} = (1, 2, 1)$ en base orthogonale directe
de $\Rr^{3}$ euclidien canonique.
\finenonce{001982}



\finexercice



\exercice{1983, ridde, 1999/11/01}

\enonce{001983}{}
Montrer que $\forall (x, y, z)\in (\Rr^{3})^{3} \; \;
x\wedge (y\wedge z) + y\wedge (z\wedge x) + z\wedge (x\wedge y) = 0$.
\finenonce{001983}



\finexercice



\exercice{1984, ridde, 1999/11/01}

\enonce{001984}{}
Soit $E$ euclidien orient\'e de dimension 3 et $a\in E$. \\
Soit $f : \begin{cases} E \rightarrow E \\ x \mapsto x\wedge a \end{cases}$.
$f$ est-elle lin\'eaire, bijective ? Comparer $f^{3}$ et $f$.
\finenonce{001984}



\finexercice



\exercice{1985, ridde, 1999/11/01}

\enonce{001985}{}
Soient $a$ et $b$ deux vecteurs de $\Rr^{3}$. Discuter
et r\'esoudre l'\'equation $a \wedge x  = b$.
\finenonce{001985}



\finexercice



\exercice{1986, ridde, 1999/11/01}

\enonce{001986}{}
 Soit $R$ le rotation vectorielle d'angle $\theta$ et
 d'axe orient\'e par le vecteur unitaire $k$. Montrer que $\forall x\in \Rr^{3}$
 $ R (x) = (\cos\theta)x + (\sin\theta)k\wedge x + 2 (x|k)\sin^{2} (\tfrac{\theta}{2})k$.
\finenonce{001986}



\finexercice



\exercice{1987, ridde, 1999/11/01}

\enonce{001987}{}
 Determiner la matrice dans la base canonique de $\Rr^{3}$
du retournement d'axe $\Rr (1, 2, 1)$.
\finenonce{001987}



\finexercice



\exercice{1988, ridde, 1999/11/01}

\enonce{001988}{}
 Reconnaître les transformations g\'eom\'etriques dont les
matrices respectives dans la base canonique de $\Rr^{3}$ sont : \\
$$\frac{1}{3} \begin{pmatrix} 3&1&\sqrt{6}\\1&3&-\sqrt{6}\\-\sqrt{6}&\sqrt{6}&2
\end{pmatrix} \qquad \frac{1}{3} \begin{pmatrix} -2&2&1\\2&1&2\\-1&-2&2 \end{pmatrix}$$
\finenonce{001988}



\finexercice



\exercice{1989, ridde, 1999/11/01}

\enonce{001989}{}
 Soit $R$ une rotation de $\Rr^{3}$ d'axe $\Rr u$ et d'angle
 $\theta$ et $r$ une rotation quelconque. D\'eterminer $rRr^{-1}$.
 En d\'eduire que le centre de $SO_{3} (\Rr)$ est bîîîîîîp (le centre est l'ensemble
 des rotations qui commutent avec toutes les autres).
\finenonce{001989}



\finexercice



\exercice{1990, ortiz, 1999/04/01}

\enonce{001990}{}
On consid\`ere l'espace vectoriel euclidien
canonique et orient\'e $\Rr^3$. Soient
$a,b,c\in\Rr^3$ et $p=\left[a,b,c\right]$ le
produit mixte de $a$, $b$ et $c$. Exprimer \`a
l'aide de $p$ les quantit\'es  suivantes
\begin{enumerate}
   \item $s=\left[a+b,b+c,c+a\right]$,
   \item $t=\left[a\wedge b,b\wedge c,c\wedge a\right].$
\end{enumerate}
\finenonce{001990}



\finexercice



\exercice{5207, rouget, 2010/06/30}
\enonce{005207}{*T}
Nature et éléments caractéristiques de la transformation d'expression complexe~:
\begin{enumerate}
\item  $z'=z+3-i$
\item  $z'=2z+3$
\item  $z'=iz+1$
\item  $z'=(1-i)z+2+i$
\end{enumerate}
\finenonce{005207}


\finexercice

\exercice{7480, mourougane, 2021/08/10}
\enonce{007480}{}
Construire un segment $[A,B]$ connaissant son milieu $I$ et sachant
que $A$ appartient à une droite donnée $d$ et $B$ à un cercle donné
$\mathcal{C}$. 
\finenonce{007480}
\finexercice


\exercice{7481, mourougane, 2021/08/10}
\enonce{007481}{}
Construire un triangle équilatéral $ABC$ connaissant $A$ et sachant
que $B$ appartient à un cercle donné $\mathcal{C}$ et $C$ à un autre 
cercle donné $\mathcal{C}'$. 
\finenonce{007481}
\finexercice


\exercice{7482, mourougane, 2021/08/10}
\enonce{007482}{}
Soit $\mathcal{C}$ un cercle et $\vec{u}$ un vecteur. Construire une
corde $AB$ de $\mathcal{C}$ telle que $\overrightarrow{AB}=\vec{u}$.
\finenonce{007482}
\finexercice


\exercice{7483, mourougane, 2021/08/10}
\enonce{007483}{}
Construire un carré $ABCD$ sachant que $A$ et $C$ sont sur une droite
$D_1$ donnée, $B$ sur une droite $D_2$ donnée 
et $D$ sur une droite $D_3$ donnée.
\finenonce{007483}
\finexercice


\exercice{7484, mourougane, 2021/08/10}
\enonce{007484}{}
Soit $d$ une droite.
Soit $\mathcal{C}$ un cercle et $A$ un
point de $\mathcal{C}$.
Construire un cercle tangent à la droite $d$ et 
tangent en $A$ au cercle $\mathcal{C}$. 
\finenonce{007484}
\finexercice 


\exercice{7485, mourougane, 2021/08/10}
\enonce{007485}{}
Dans un triangle $ABC$ rectangle en $A$ trouver la relation entre la
longueur de la hauteur $[AH]$ issue de $A$ et la longueur des segments
$[BH]$ et $[CH]$. {\'E}tant donné un segment de longueur $a$ construire
un segment de longueur $\sqrt{a}$. 
\finenonce{007485}
\finexercice


\exercice{7486, mourougane, 2021/08/10}
\enonce{007486}{}
Étant donné un segment de longueur $a$ contruire
un segment de longueur $a\sqrt{5}$.
\finenonce{007486}
\finexercice

\section{242.00 Géométrie affine euclidienne}

\section{242.01 Géométrie affine euclidienne du plan}


\exercice{2035, liousse, 2003/10/01}

\enonce{002035}{Transformations affines et Isom\'etries}

Soit 
$P$ un plan muni d'un rep\`ere $(O,{\buildrel\rightarrow \over i},
{\buildrel\rightarrow \over j})$ quelconque.

\begin{enumerate}
\item  On consid\`ere $D$ une droite d'\'equation cart\'esienne 
$2x-y+3=0$ et ${\buildrel\rightarrow \over u}(3,-2).$
\begin{enumerate}
\item  Soit $A(4,2).$ Donner une \'equation param\'etrique de 
$D_A$ droite passant par $A$ de direction ${\buildrel\rightarrow \over u}.$ 
En d\'eduire les coordonn\'ees de $A'=D_A\cap D$ projet\'e de 
$A$ sur $D$ selon ${\buildrel\rightarrow \over u}.$

\item  D\'efinir plus g\'en\'eralement analytiquement la projection sur 
$D$ selon ${\buildrel\rightarrow \over u}$ en exprimant 
les coordonn\'ees $x',$ $y'$ de $M'$ projet\'e de $M(x,y)$ 
en fonction de $x$ et $y.$\\
\end{enumerate}
\item D\'efinir analytiquement les projections sur $D$ selon $\Delta$ 
dans les cas suivants :
\begin{enumerate}
\item $\Delta$ d'\'equation $x-2y+1=0.$
\item $\Delta$ d'\'equation $3x+2y+2=0.$
\item  $\Delta$ d'\'equation $x+y-1=0.$
\item  $\Delta$ d'\'equation $2x-2y+4=0.$
\end{enumerate}
\end{enumerate}
\finenonce{002035}



\finexercice



\exercice{2036, liousse, 2003/10/01}

\enonce{002036}{}
  Soit  $P$ un plan muni d'un rep\`ere $(O,{\buildrel\rightarrow \over i},
{\buildrel\rightarrow \over j})$ quelconque.
\begin{enumerate}
\item  Donner l'expression 
analytique de la translation $t_1$ de vecteur $(1,2).$\\
\item  Donner l'expression analytique de la translation $t_2$ de vecteur $(-1,2).$\\
\item  Donner l'expression analytique de l'homoth\'etie $h_1$ de centre l'origine 
du rep\`ere et de rapport 2 et de l'homoth\'etie $h_2$ de centre $A(2,-1)$ de 
rapport 3.\\
\item Donner l'expression analytique de $t_1\circ h_1,$ $t_2\circ h_2,$ 
$h_1\circ t_1,$ $h_2\circ t_2.$
\item  Soit  $M(x,y)$ un point de $P$. Donner les coordonn\'ees du sym\'etrique de 
$M$ par rapport \`a la droite d'\'equation $y=ax+b.$
\end{enumerate}
\finenonce{002036}



\finexercice



\exercice{2037, liousse, 2003/10/01}

\enonce{002037}{}
\begin{enumerate}
\item  On consid\`ere $S_1$ la transformation du plan 
d\'efinie par le syst\`eme d'\'equations suivant :

$x'= {\sqrt 3 \over 2} x + {1 \over 2} y -1,\ 
y'=-{1\over 2}x + {\sqrt 3 \over 2}y +2.$ 
Reconna\^itre cette transformation.

\item  De m\^eme avec la transformation $S_2$ d\'efinie par
$x'= 5{\sqrt 2} x + 5{\sqrt 2} y,\ 
y'=-5{\sqrt 2}x + 5{\sqrt 2}y.$ 

\item  On compose $S_1$ avec $S_2$. Donner l'expression de $S_1 \circ S_2$, et trouver la nature de 
cette transformation.
\end{enumerate}
\finenonce{002037}



\finexercice



\exercice{2040, liousse, 2003/10/01}

\enonce{002040}{}

Dans le plan muni d'un rep\`ere orthonorm\'e direct $(O,\vec{OI},\vec{OJ})$.
\begin{enumerate}
\item  Soit $f$ la transformation 
du plan d\'efinie analytiquement par
$$\left\{\begin{array}{ll}x'=&{1\over \sqrt{5}}(x+2y-1)\\y'=&{1\over \sqrt{5}}
(-2x+y+2)\end{array}\right.$$
\begin{enumerate}
\item  Calculer les coordonn\'ees de $O'$, $I'$, $J'$ les images par $f$ des points $O$, $I$, $J$.

\item  Montrer que le rep\`ere  $(O',\vec{O'I'},\vec{O'J'})$ est orthonorm\'e, est-il direct ?

\item  En d\'eduire  que $f$ est une isom\'etrie, est-elle directe ?

\item  D\'eterminer l'ensemble  des points invariants par $f$ et reconnaitre $f$.

\item   Donner l'expression analytique de la transformation inverse de $f$.

\item  Calculer l'image par $f$ la droite d'\'equation $2x-y-1=0$.
\end{enumerate}

\item  Donner l'expression analytique de la rotation de centre $A(1,1)$ et
d'angle $\pi\over 3$, calculer l'image de $0$ par cette transformation.


\item  M\^eme question pour la sym\'etrie d'axe la droite d'\'equation $ x+y+1=0$

\item  Donner l'expression analytique  de la compos\'ee des deux applications pr\'ec\'edentes.
\end{enumerate}
\finenonce{002040}



\finexercice



\exercice{2041, liousse, 2003/10/01}

\enonce{002041}{}
  Dans le plan cart\'esien identifi\'e \`a $\Cc$, un point $M$ est repr\'esent\'e par son affixe $z$.
\begin{enumerate}
\item  Dessiner les ensembles suivants puis les exprimer en fonction de $(x,y)$ ($(z=x+iy)$) :

(i) $z+\overline{z}=1$ \hfill  (ii) $z-\overline{z}=i$ \hfill  (iii) $iz-i\overline{z}=1$ 

\item  Donner l'expression analytique en complexe des transformations suivantes, puis 
calculer l'image de $i$ par ces transformations :
\begin{enumerate}
\item  la rotation de centre $1+i$ et
d'angle $\pi\over 3$,


\item  la sym\'etrie d'axe la droite d'\'equation  $iz-i\overline{z}=1$,

\item   la compos\'ee des deux applications pr\'ec\'edentes.
\end{enumerate}
\smallskip

\item  Soit $f$ la transformation  du plan d\'efinie analytiquement par $z'= (1+i)z+ 1$.

\begin{enumerate}
\item  D\'eterminer l'ensemble  des points invariants par $f$.

\item  Donner l'expression analytique de la transformation inverse de $f$.

\item  Calculer l'image par $f$ de l'ensemble $z+\overline{z}=1$.

\item  Ecrire $f$ comme la compos\'ee d'une homoth\'etie et d'une isom\'etrie.
\end{enumerate}
\end{enumerate}
\finenonce{002041}



\finexercice



\exercice{4949, quercia, 2010/03/17}
\enonce{004949}{Fonction numérique de Leibniz}

Soit $ABC$ un triangle équilatéral de côté $a$.

Quels sont les points $M$ du plan $(ABC)$ tels que
$MA^2 + a^2 = 2(MB^2 + MC^2)$ ?

\finenonce{004949}



\finexercice


\exercice{4950, quercia, 2010/03/17}
\enonce{004950}{Cercle circonscrit à un triangle}

Soit $ABC$ un triangle et $\mathcal{C}$ son cercle circonscrit.
Soit $M$ un point du plan de coordonnées barycentriques $(x,y,z)$ dans
le repère affine $(ABC)$.

Montrer que : $M\in\mathcal{C} \Leftrightarrow xAM^2 + yBM^2 + zCM^2 = 0
              \Leftrightarrow xyAB^2 + xzAC^2 + yzBC^2 = 0$.

\finenonce{004950}



\finexercice


\exercice{4951, quercia, 2010/03/17}
\enonce{004951}{Cercle stable par une application affine}

Soit ${\cal C} = {\cal C}(O,r)$ un cercle du plan et $f$ une application affine
telle que $f({\cal C}) = {\cal C}$.
Montrer que $f$ est une isométrie de point fixe $O$.
\finenonce{004951}



\finexercice


\exercice{4952, quercia, 2010/03/17}
\video{zldT8ImfOX8}
\enonce{004952}{Point équidistant d'une famille de droites}

Pour $\lambda \in \R$ on considère la droite $D_\lambda$ d'équation
cartésienne : $(1-\lambda^2)x + 2\lambda y = 4\lambda + 2$.

Montrer qu'il existe un point $M_0$ équidistant de toutes les droites
$D_\lambda$.
\finenonce{004952}


\finexercice

\exercice{4953, quercia, 2010/03/17}
\enonce{004953}{Bissectrice de deux droites}

Soient $D,D'$ deux droites distinctes sécantes en $O$.

On note ${\cal H} = \{ M$ tq $d(M,D) = d(M,D') \}$.

\begin{enumerate}
  \item Montrer que $\cal H$ est la réunion de deux droites perpendiculaires.
    (appelées bissectrices de $(D,D')$)

  \item Soit $s$ une symétrie orthogonale telle que $s(D) = D'$.
    Montrer que l'axe de $s$ est l'une des droites de $\cal H$

  \item Soit $\cal C$ un cercle du plan tangent à $D$.
    Montrer que $\cal C$ est tangent à $D$ et à $D'$
    si et seulement si son centre appartient à $\cal H$.
\end{enumerate}
\finenonce{004953}



\finexercice


\exercice{4954, quercia, 2010/03/17}
\enonce{004954}{Trois figures isométriques}


Trois figures $F_1,F_2,F_3$ se déduisent l'une de l'autre par rotations.
Montrer qu'il existe une figure $F$ dont $F_1,F_2,F_3$ se déduisent par
symétries axiales.


\finenonce{004954}



\finexercice


\exercice{4955, quercia, 2010/03/17}
\enonce{004955}{Produit de 3 rotations}

Soit $ABC$ un triangle d'angles $\alpha,\beta,\gamma$.

On note $\rho,\rho',\rho''$ les rotations autour de $A,B,C$ d'angles
$\alpha,\beta,\gamma$, orientés suivant le dessin :
$$
\includegraphics[height=3cm]{../images/img004955-1}
$$

Qu'est-ce que $\rho \circ \rho' \circ \rho''$ ?
\finenonce{004955}



\finexercice


\exercice{4956, quercia, 2010/03/17}
\enonce{004956}{Sous-groupes finis de déplacements}

\begin{enumerate}
  \item Soit $G$ un sous-groupe fini de déplacements du plan.
  \begin{enumerate}
    \item   Montrer que $G$ est constitué uniquement de rotations.
    \item   Soient $f,g \in G$. Montrer que $f$ et $g$ ont même centre
          (étudier $f \circ g \circ f^{-1} \circ g^{-1}$).
    \item   Prouver enfin que $G$ est cyclique.
  \end{enumerate}
  \item Soit $G$ un sous groupe fini d'ordre $p$ d'isométries du plan, non toutes
    positives.
 \begin{enumerate} 
    \item   Montrer que $G$ contient autant d'isométries positives que négatives.
    \item   Montrer que $G$ est un groupe diédral (groupe d'isotropie d'un polygone
          régulier).
  \end{enumerate}
\end{enumerate}
\finenonce{004956}



\finexercice


\exercice{4957, quercia, 2010/03/17}
\enonce{004957}{Centrale MP 2000}

Soit~$E$ un plan affine euclidien muni d'un repère orthonormé
d'origine~$O$. Soit~$A$ le point de coordonnées~$(a,0)$.
Pour tout point~$M$, on définit $M'=f(M)$ de la manière suivante~:
$A,M,M'$ sont alignés et $(MO)$ est orthogonale à~$(M'O)$.
Expliciter~$f$ en fonction des coordonnées~$(x,y)$ de~$M$.
Donner son domaine de définition. Montrer que~$f$ réalise une bijection
entre le demi-disque supérieur de diamètre~$[AO]$ et le quart de plan d'équations $x<0, y>0$.

\finenonce{004957}



\finexercice


\exercice{5197, rouget, 2010/06/30}
\enonce{005197}{*IT}
Déterminer le projeté orthogonal du point $M(x_0,y_0)$ sur la droite $(D)$ d'équation $x+3y-5=0$ ainsi que son symétrique orthogonal.
\finenonce{005197}


\finexercice


\exercice{5198, rouget, 2010/06/30}
\enonce{005198}{*}
Soit $(ABDC)$ un parallélogramme. Déterminer les coordonnées de $D$ dans le repère $(A,\overrightarrow{AB},\overrightarrow{AC})$.
\finenonce{005198}


\finexercice


\exercice{5201, rouget, 2010/06/30}
\enonce{005201}{**I}
\label{exo:suprou7bis}
\begin{enumerate}
\item  $h$ (resp. $h'$) est l'homothétie de centre $\Omega$ et de rapport $k$ (resp. $k'$) non nul. Déterminer la nature et les éléments caractéristiques de $h'\circ h$.
\item  $s$ (resp. $s'$) est la symétrie centrale de centre $\Omega$ (resp. $\Omega'$). Déterminer la nature et les éléments caractéristiques de $s'\circ s$.
\item  $s$ est la symétrie centrale de centre $\Omega$ et $t$ est la translation de vecteur $\overrightarrow{u}$. Déterminer la nature et les éléments caractéristiques de $t\circ s$.
\end{enumerate}
\finenonce{005201}


\finexercice


\exercice{7047, megy, 2017/01/08}

\enonce{007047}{Construction de parallèles et de perpendiculaires}
% Source :
% Tags :  
On donne une droite $\mathcal D$ et un point $P \not\in \mathcal D$.  Tracer la parallèle ainsi que la perpendiculaire à $\mathcal D$ passant par $P$.
\finenonce{007047}

\finexercice  


\exercice{7048, megy, 2017/01/08}

\enonce{007048}{Construction du centre et de  tangentes}
% tags : corde, médiatrice, symétrique, milieu, triangle rectangle, cercle circonscrit
On donne un cercle $\mathcal C$ (sans son centre).
\begin{enumerate}
\item  Tracer son centre, si possible de plusieurs façons. 
\item On donne un point $P$ à l'extérieur du cercle. Tracer les tangentes à $\mathcal C$ passant par $P$.
\item Même question si $P$ est sur le cercle.
\end{enumerate}
\finenonce{007048}
\finexercice  


\exercice{7049, megy, 2017/01/08}

\enonce{007049}{Triangle de l'écolier, 1}
%tags : triangle rectangle, cercle circonscrit
On donne deux points $A$ et $B$. Construire un triangle $ABC$ rectangle en $C$ tel que $AB = 2AC$.
\finenonce{007049}
\finexercice  


\exercice{7050, megy, 2017/01/08}
\enonce{007050}{Triangle de l'écolier, 2}
Soit $ADC$ un triangle équilatéral et $B$ le symétrique de $A$ par rapport à $D$. Montrer que $ABC$ est rectangle en $C$ et déterminer également $\widehat A$ et $\widehat B$.
\begin{center}
\begin{tikzpicture}[line cap=round,line join=round,x=1.0cm,y=1.0cm]
\clip(-1.38,-0.88) rectangle (9.2,5.2);
\fill[fill opacity=0.1] (-0.22,0.26) -- (3.88,0.26) -- (1.83,3.81) -- cycle;
\draw  (-0.22,0.26)-- (3.88,0.26);
\draw  (3.88,0.26)-- (1.83,3.81);
\draw  (1.83,3.81)-- (-0.22,0.26);
\draw (7.98,0.26)-- (1.83,3.81);
\draw (7.98,0.26)-- (3.88,0.26);
\begin{scriptsize}
\draw (-0.5,0.62) node {$A$};
\draw (4.02,0.62) node {$D$};
\draw (1.98,4.18) node {$C$};
\draw (8.12,0.62) node {$B$};
\end{scriptsize}
\end{tikzpicture}
\end{center}
\finenonce{007050}
\finexercice  


\exercice{7051, megy, 2017/01/08}

\enonce{007051}{Subdivision et barycentres}
On donne un segment $[AB]$.
\begin{enumerate}
\item Diviser le segment en quatre.
\item  Diviser le segment en trois parties égales sans utiliser le théorème de Thalès.
\item Diviser le segment en sept parties égales.
\item Construire le barycentre de $(A,3)$ et $(B,4)$.
\end{enumerate}
\finenonce{007051}
\finexercice  


\exercice{7052, megy, 2017/01/08}

\enonce{007052}{Isobarycentre}
\begin{enumerate}
\item On donne un quadrilatère $ABCD$. Construire l'isobarycentre de ses sommets. Construire le barycentre de $(A,1)$, $(B,1)$, $(C,3)$ et $(D,3)$.
\item On donne un pentagone quelconque. Construire l'isobarycentre de ses sommets.
\end{enumerate}
\finenonce{007052}

\finexercice  


\exercice{7053, megy, 2017/01/08}

\enonce{007053}{Introduction aux nombres constructibles}
% source : wikipedia, Carrega
On donne deux points $O$ et $I$, avec $OI=1$. Un réel $r$ est constructible si on peut construire à la règle et au compas un point $M$ tel que $\overrightarrow{OM}=r\overrightarrow{OI}$. Le but de l'exercice est de montrer que l'ensemble des nombres constructibles est un sous-corps de $\R$ stable par racine carrée.
\begin{enumerate}
\item  Construire sur la droite $(OI)$ des points $A$, $B$ et  $C$  tels que $OA = \frac{1}{\sqrt{2}}$,  $OB =\sqrt{2}$ et $OC =\sqrt{3}$.
\item (Construction du produit et de l'inverse de deux nombres constructibles.) On donne deux points $A$ et $B$ alignés avec $O$. Construire sur la droite $(AB)$ des points $C$ et $D$ tel que $OC = OA\times OB$ et $OD = \frac{OA}{OB}$. 
\item (Construction de la racine carrée.) Soit $A$ un point sur la demi-droite $[OI)$. Soit $I'$ le symétrique de $I$ par rapport à $O$, soit $\mathcal C$ le cercle de diamètre $I'A$, et soit $F$ l'une des intersections du cercle $\mathcal C$ avec la perpendiculaire à $(OA)$ passant par $O$. Montrer que $OF = \sqrt{OA}$.
\end{enumerate}
\finenonce{007053}
      
\finexercice  


\exercice{7054, megy, 2017/01/08}

\enonce{007054}{Tangentes communes à deux cercles}
% tags triangle rectangle, cercle circonscrit
On donne deux cercles $\mathcal C$ et $\mathcal C'$ de rayons $r < r'$, de centres $O$ et $O'$, disjoints et extérieurs l'un à l'autre. On admet qu'il existe quatre tangentes communes  à $\mathcal C$ et $\mathcal C'$. L'objectif est de les construire.
\begin{enumerate}
\item (Analyse) Soit $\mathcal D$ une tangente commune. On note $A$ et $A'$ les points de contact de $\mathcal D$ avec les deux cercles. Que peut-on dire de la parallèle à $(AA')$ passant par $O$ et de son intersection avec $(O'A')$ ? 
\item (Synthèse) En déduire une construction du point d'intersection de ces deux droites, puis ces deux droites et enfin de $\mathcal D$. Tracer les quatre tangentes communes de cette façon.
\end{enumerate}

\finenonce{007054}

\finexercice  


\exercice{7055, megy, 2017/01/08}

\enonce{007055}{Construction de cercles, 1}
% tags : construction de cercles, CDP
\begin{enumerate}
\item On donne une droite $\mathcal D$ et un point $O$ hors de la droite. Tracer le cercle de centre $O$ et tangent à $\mathcal D$.

\item On donne deux droites $\mathcal D$ et $\mathcal D'$ sécantes et un point $O$ d'une bissectrice (et hors des droites). Construire le cercle de centre $O$ et tangent aux deux droites. 
\end{enumerate}

\finenonce{007055}
\finexercice  


\exercice{7056, megy, 2017/01/08}

\enonce{007056}{Construction de cercles, 2}
\begin{enumerate}
\item On donne une droite $\mathcal D$, un point $H$ sur $\mathcal D$ et un point $A$ en-dehors. Tracer le cercle passant par $A$ et tangent à la droite en $H$.% cas particulier de PPP ou PPD



\item On donne trois droites dont deux parallèles. Dénombrer et construire les cercles tangents aux trois droites. 


\end{enumerate}

\finenonce{007056}

\finexercice  


\exercice{7057, megy, 2017/01/08}

\enonce{007057}{Construction de cercles}
% tags : construction de cercles, cas particulier de CCC
On donne trois cercles distincts $\mathcal C_1$, $\mathcal C_2$ et $\mathcal C_3$ de même rayon et dont les centres ne sont pas alignés. Construire deux cercles tangents à $\mathcal C_1$, $\mathcal C_2$ et $\mathcal C_3$.

\finenonce{007057}
\finexercice  


\exercice{7058, megy, 2017/01/08}

\enonce{007058}{Construction de cercles de rayon donné}
% tags : construction de cercles, cas particulier de CCC
Dans tout l'exercice, on fixe $R>0$. Dénombrer et construire les cercles de rayon $R$ tangents à :
\begin{enumerate}
\item deux cercles distincts $\mathcal C_1$ et $\mathcal C_2$;
\item deux droites sécantes $\mathcal D_1$ et $\mathcal D_2$;
\item un cercle $\mathcal C$ et une droite $\mathcal D$.
\end{enumerate}
\finenonce{007058}
\finexercice  


\exercice{7077, megy, 2017/01/21}

\enonce{007077}{}
% Source :
% Tags :  
Soit $\mathcal D$ une droite et $\sigma$ la réflexion orthogonale suivant cette droite. Construire l'image par $\sigma$ d'un point, d'un segment, d'un droite, d'un cercle (que ces objets intersectent l'axe  de symétrie ou pas).\finenonce{007077}


\finexercice  


\exercice{7078, megy, 2017/01/21}

\enonce{007078}{Constructions}
\begin{enumerate}
\item On donne deux points $A$ et $B$. Soit $\tau$ la translation de vecteur $\overrightarrow{AB}$ et $M$ un point du plan. Construire $\tau(M)$.
\item Soit $\mathcal D$ une droite du plan, et $\sigma$ la réflexion orthogonale d'axe $\mathcal D$. Soit $M$ un point du plan. Construire $\sigma(M)$. Réciproquement, soit $\sigma$ une réflexion orthogonale, et supposons donnés un point $A$ et son image $\sigma(A)$. Construire $\mathcal D$.
\item On donne un point $O$, et $\phi$ l'homothétie de centre $O$ et de rapport $5/8$. Soit $M$ un point du plan. Construire $\phi(M)$. Remarque : ceci marche pour tout rapport rationnel.
\item On donne un triangle auxiliaire $ABC$ et on considère l'angle $\alpha=(\overrightarrow{AB},\overrightarrow{AC})$. Soit $\rho$ la rotation de centre $O$ et d'angle $\alpha$. Si $M$ est un point du plan, construire $\rho(M)$.
\end{enumerate}\finenonce{007078}


\finexercice  


\exercice{7079, megy, 2017/01/21}

\enonce{007079}{}
\begin{enumerate}
\item À quelle condition sur quatre points $P_1, ... P_4$ existe-t-il une homothétie $h$ telle que $h(P_1)=P_2$ et $h(P_3)=P_4$ ?
\item Soit $\phi$ une homothétie. On donne deux points $A$ et $B$, ainsi que leurs images $\phi(A)$ et $\phi(B)$. Le centre de l'homothétie n'est pas donné. Le construire, y compris si les quatre points donnés sont alignés.
% dans le deuxième cas, tracer un triangle ABC, et tracer son image par l'homthétie en tracant les parallèles. On trouve un point C', et CC' intersecte la droite (AB) en O.
\item  Soit $\rho$ une rotation du plan. On donne deux points $A$ et $B$, ainsi que leurs images $\rho(A)$ et $\rho(B)$. Construire le centre de la rotation, en distinguant les cas.
\end{enumerate}\finenonce{007079}


\finexercice  


\exercice{7080, megy, 2017/01/21}

\enonce{007080}{Polygone des milieux}
% Source : Dehornoy
% tags : symétries centrales
On donne un nombre impair de points du plan $M_1$, ...$M_n$. Existe-t-il un polygone $P_1$, $P_2$, ... $P_n$ tel que les $M_i$ soient les milieux des côtés du polygone ? Commencer par $n=3$. Et si $n$ est pair  ? En particulier, si $n=4$, trouver une condition nécessaire et suffisante pour que le problème admette une solution et  déterminer l'ensemble des solutions.
\finenonce{007080}
\finexercice  


\exercice{7081, megy, 2017/01/21}

\enonce{007081}{Trapèze rectangle}
%tags : isocèle, trapèze, symétrie centrale, projection
Soit $\mathcal D$ une droite, $A$ et $B$ deux points hors de cette droite, et $A'$, $B'$ leurs projetés orthogonaux sur $\mathcal D$, supposés distincts. Soit enfin $I$ le milieu de $[AB]$. Montrer que $A'IB'$ est isocèle en $I$.
\begin{center}
\includegraphics{../images/img007081-1}
\end{center}
\finenonce{007081}
\finexercice  


\exercice{7082, megy, 2017/01/21}

\enonce{007082}{Trapèze isocèle}
% tags : symétrie centrale, réflexion orthogonale
Soit $ABC$ un triangle,  $P$ le pied de la hauteur issue de $A$ et $I$, $J$, $K$  les milieux des côtés $[AB]$, $[BC]$ et $[CA]$. Montrer que le quadrilatère $IPJK$ est un trapèze isocèle.
\begin{center}
\includegraphics{../images/img007082-1}
\end{center}
\finenonce{007082}


Par le théorème des milieux, on a $(IK) // (PJ)$ donc $IPJK$ est un trapèze.

Soit $Q$ le milieu de $[IK]$. La perpendiculaire à $(IK)$ en $Q$ est alors un axe de symétrie de $IPJK$, par le théorème des milieux appliqué à $APJ$.
\finexercice  


\exercice{7083, megy, 2017/01/21}

\enonce{007083}{Deux cercles sont homothétiques}
%tags : homothéties
% source : http://debart.pagesperso-orange.fr/geoplan/construc_probleme_classique.html#ch13
Soient $\mathcal C$ et $\mathcal C'$ deux cercles de rayons distincts. Montrer qu'il existe des homothéties transformant l'un en l'autre. Suivant la position et la taille des cercles, combien y a-t-il de telles homothéties ? Tracer leurs centres.
\finenonce{007083}
\finexercice  


\exercice{7084, megy, 2017/01/21}

\enonce{007084}{Tangentes communes à deux cercles}
% tags: homothéties et translations
On donne deux cercles. Déterminer le nombre de tangentes communes aux deux cercles et tracer ces tangentes.
\finenonce{007084}
\finexercice  


\exercice{7085, megy, 2017/01/21}

\enonce{007085}{Carré inscrit dans un triangle}
% homothétie, Thalès
Soit $ABC$ un triangle. Construire un carré dont un sommet appartient à $[AB]$, un à $[AC]$ et deux sommets adjacents appartiennent à $[BC]$.
\finenonce{007085}
\finexercice  


\exercice{7086, megy, 2017/01/21}

\enonce{007086}{Cercle d'Euler}
% tags : homothéties, cercle circonscrit, cercle des neuf points
% cercle de Feuerbach, cercle de Terquem, cercle médian
Soit $ABC$ un triangle. On note $G$, $\Omega$ et $H$ le centre de gravité, le centre du cercle circonscrit $\mathcal C$ et l'orthocentre. Soit $\mathcal C'$ le cercle passant par les milieux $I_A$, $I_B$ et $I_C$ des côtés de $ABC$.
\begin{enumerate}
\item Montrer que le centre de $\mathcal C'$ appartient à la droite $(G\Omega)$. Calculer son rayon.
\item Montrer que $\mathcal C'$ coupe les segments reliant les sommets à l'orthocentre $H$ en leur milieu.
\end{enumerate}
% il y a beaucoup d'autres choses à montrer, mais c'est plus pratique avec les angles inscrits.
\begin{center}
\includegraphics{../images/img007086-1}
\end{center}

\finenonce{007086}
\finexercice  


\exercice{7087, megy, 2017/01/21}

\enonce{007087}{Théorème du trapèze}
% homothéties
Soient $[AB]$ et $[CD]$ deux segments parallèles de longueurs différentes. Montrer qu'il existe des homothéties transformant l'un en l'autre. Combien y a-t-il de telles homothéties ? Tracer leurs centres. Montrer que la droite reliant ces centres coupe les segments en leur moitié.

Application : construire à la règle seule le symétrique de $A$ par rapport à $B$. Expliquer comment construire à la règle seule n'importe quel barycentre à coefficients rationnels de $A$ et de $B$.
\finenonce{007087}
\finexercice  


\exercice{7088, megy, 2017/01/21}

\enonce{007088}{Application du théorème du trapèze}
Soit $\mathcal D$ une droite, $A$ et $B$ deux points distincts  n'appartenant pas à cette droite, et $A'$ le symétrique de $A$ par rapport à $\mathcal D$. On suppose que $A'B$ n'est pas parallèle à $\mathcal D$. Construire à la règle seule le symétrique $B'$ de $B$ par rapport à $\mathcal D$.
\finenonce{007088}
\finexercice  


\exercice{7089, megy, 2017/01/21}

\enonce{007089}{Application du théorème du trapèze, bis}
Soient deux droites parallèles distinctes $\mathcal D$ et $\mathcal D'$, et $P$ un point non situé sur ces droites. Construire à la règle seule la droite passant par $P$ et parallèle aux deux autres.\finenonce{007089}
\finexercice  


\exercice{7090, megy, 2017/01/21}

\enonce{007090}{Où est le centre?}
% source : Leitchnam
% se fait facilement en coisissant des coordonnées affines, en fait.

Soit $ABCD$ un parallélogramme, et $M$ un point sur la diagonale $(BD)$. Soit $I$ le symétrique de $C$ par rapport à $M$. Soit $E$ la projection de $I$ sur $(AB)$ parallèlement à $(AD)$, et $F$ la projection de $I$ sur $(AD)$ parallèlement à $(AB)$. Montrer que $E$, $M$ et $F$ sont alignés.
\finenonce{007090}
\finexercice  


\exercice{7091, megy, 2017/01/21}

\enonce{007091}{Tangentes communes à plusieurs cercles}
% homothéties
%(\cite[I, prop. 17]{Euclide})
 On donne un cercle  $\mathcal C$ (de centre $O$), un point $M$ à l'extérieur du cercle, les deux tangentes $\mathcal D$ et $\mathcal D'$ à $\mathcal C$ passant par $M$. On notera $A$ et $B$ les points de tangence. 

Le cercle $\mathcal C$ coupe $(MO)$ en deux points $P$ et $Q$. D'autre part, soit $H$ l'intersection de la corde $[AB]$ avec $(OM)$. À l'aide d'homothéties, montrer que les cercles de centres $P$ et $Q$ et passant par $H$ sont tangents à $\mathcal D$ et $\mathcal D'$. 
\begin{center}
\includegraphics{../images/img007091-1}
\end{center}
\finenonce{007091}
\finexercice  


\exercice{7092, megy, 2017/01/21}

\enonce{007092}{}
% homothéties
On donne trois cercles disjoints, de rayons distincts et à l'extérieur les uns des autres. Chacune des trois paires de cercles fournit deux tangentes communes extérieures qui se croisent en un point.  Montrer que ces trois points sont alignés.
\finenonce{007092}
\finexercice  


\exercice{7093, megy, 2017/01/21}

\enonce{007093}{}
% tags : translation

Soit $ABCD$ un rectangle et $M$ un point du plan.

On note $C’$ le projeté orthogonal de $C$ sur $(AM)$,
$D’$  le projeté orthogonal de $D$ sur $(BM)$ et 
$M’$  le projeté orthogonal de $M$ sur $(AB)$. 
Enfin, on note $I$ le point d'intersection des droites $(CC’)$ et $(DD’)$.

Montrer que les points $M$, $M’$ et $I$ sont alignés.
\finenonce{007093}
\finexercice  


\exercice{7094, megy, 2017/01/21}

\enonce{007094}{Théorème de Pappus, version affine}
% Source :  Audin par exemple
Soient $D$ et $D'$ deux droites. Soient $A, B, C$ trois points sur $D$, et $A'$, $B'$ et $C'$ trois points sur $D'$. Si $(AB') // (BC')$ et $(BA') // (CB')$, alors $(AA') // (CC')$.
\finenonce{007094}
\finexercice  


\exercice{7095, megy, 2017/01/21}

\enonce{007095}{Théorème de Desargues, version  affine}
% tags : homothéties et translations.
\begin{enumerate}
\item Soient $ABC$ et $A'B'C'$ deux triangles (non aplatis) sans sommet commun. Montrer qu'ils se déduisent l'un de l'autre par homothétie ou translation ssi leurs côtés sont parallèles.

\item (Application) On donne deux droites se coupant en un point $O$ hors de la feuille, ainsi qu'un point $M$ hors de ces droites. Tracer la droite $(OM)$.
\end{enumerate}
\finenonce{007095}
\finexercice  


\exercice{7096, megy, 2017/01/21}

\enonce{007096}{}
% tags : homothéties, translations
Soient $f$ et $g$ deux homothéties de même rapport et de centres distincts. Déterminer la nature de $f\circ g^{-1}$.
\finenonce{007096}
\finexercice  


\exercice{7097, megy, 2017/01/21}

\enonce{007097}{Corde de longueur fixée}
% translations
Soit $\mathcal C$ un cercle et $D$ une droite. Construire une droite parallèle à $D$ coupant le cercle $\mathcal C$ en deux points situés à une distance $a$ donnée (inférieure au diamètre).
\finenonce{007097}
\finexercice  


\exercice{7098, megy, 2017/01/21}
\enonce{007098}{Cercle tangent à deux droites et passant par un point}
On donne deux droites parallèles et un point $A$ entre les deux droites. Dénombrer et tracer les cercles passant par $A$ et tangents aux deux droites. %translations sous-jacentes.
\finenonce{007098}
\finexercice


\exercice{7099, megy, 2017/01/21}

\enonce{007099}{}
% tags : translation, homothétie
\begin{enumerate}
\item On donne deux droites parallèles et un point $A$ entre les deux droites. Dénombrer et tracer les cercles passant par $A$ et tangents aux deux droites.
\item On donne deux droites sécantes et un point $A$ n'appartenant pas aux deux droites. Dénombrer et tracer les cercles passant par $A$ et tangents aux deux droites.

\end{enumerate}

\finenonce{007099}
\finexercice  


\exercice{7100, megy, 2017/01/21}

\enonce{007100}{Permutation cyclique}
% variations avec d'autres types permutations
Soient $A$, $B$, $C$ et $D$ quatre points du plan tels que trois d'entre eux ne soient jamais alignés. Donner une condition nécessaire et suffisante sur $ABCD$ pour qu'il existe une transformation affine $f$ du plan telle que $f(A)=B$, $f(B)=C$, $f(C)=D$ et $f(D)=A$. Montrer qu'une telle transformation, si elle existe, est nécessairement d'ordre quatre dans le groupe affine.
%
\finenonce{007100}
\finexercice  


\exercice{7101, megy, 2017/01/21}

\enonce{007101}{Minimisation}
Soient $C$ et $C'$ deux cercles sécants de centres $O$ et $O'$, et $A$ un de leurs points d'intersection. Une droite $D$ passant par $A$ recoupe les deux cercles en $M$ et $M'$. Déterminer la position de la droite qui maximise la distance $MM'$ et calculer le maximum.
\finenonce{007101}
\finexercice  


\exercice{7102, megy, 2017/01/21}

\enonce{007102}{Cloître}
% Tags :  carré, rotation
On se donne un carré, et on cherche à construire un carré de même centre, aux côtés parallèles, et d'aire deux fois plus petite, comme ci-dessous:

\begin{center}
\includegraphics{../images/img007102-1}
\end{center}



\begin{enumerate}
\item (Question préliminaire) Soit $\mathcal C$ un cercle. Montrer qu'un carré circonscrit au cercle a une aire deux fois plus grande qu'un carré inscrit dans le cercle.
\item En déduire une solution au problème initial.
\end{enumerate}
\finenonce{007102}
\finexercice  


\exercice{7103, megy, 2017/01/21}

\enonce{007103}{Construction de l'octogone régulier}
%rotation, octogone, analyse-synthèse
Construire  un octogone régulier inscrit dans un carré donné (c'est-à-dire un octogone dont quatre des huit cotés s'appuient sur les cotés du carré).
\finenonce{007103}
\finexercice  


\exercice{7104, megy, 2017/01/21}

\enonce{007104}{Carré invisible}
% tags : rotation, translation, difficile
On considère un carré $ABCD$ et on place quatre points $E$, $F$, $G$, et $H$ sur les côtés de ce carré, en-dehors des sommets. Puis, on efface le carré. En considérant une rotation et une translation, reconstruire le carré.
\finenonce{007104}
\finexercice  


\exercice{7105, megy, 2017/01/21}

\enonce{007105}{}
% source : Debart "rotation au collège"
% tags: collège, rotation

Sur les côtés $[AB]$ et $[BC]$ d'un carré direct $ABCD$, on place des points $M$ et $N$vérifiant $AM = BN$. Soit $H$ le point d'intersection des droites $(AN)$ et $(CM)$. Montrer que $H$ est l'orthocentre du triangle $DMN$.\finenonce{007105}
\finexercice  


\exercice{7106, megy, 2017/01/21}

\enonce{007106}{}
% source : Debart rotations au collège
% tags : rotations, collège
Soit $ABCD$ un carré de centre $O$, et $OPQR$ un second carré de même taille. Calculer l'aire de l'intersection de ces deux carrés en fonction de l'aire de $ABCD$.
\finenonce{007106}
\finexercice  


\exercice{7107, megy, 2017/01/21}

\enonce{007107}{}
% source : Debart rotations aux collège
% rotations, collège
Soit $ABCD$ un carré de centre $O$, et $M$ un point sur $(AB)$. À partir de $M$, on construit le triangle isocèle $OMN$, rectangle en $O$. Montrer que les points $B$, $C$ et $N$ sont alignés.
\finenonce{007107}
\finexercice  


\exercice{7108, megy, 2017/01/21}

\enonce{007108}{}
% tags : rotation, complexes

Soit $BOA$ un triangle indirect quelconque, $OAC$ et $OEB$  deux triangles rectangles isocèles en $O$ directs.

Montrer que $ACEB$ est un pseudo-carré, c'est-à-dire que les droites $(AE)$ et $(BC)$ sont perpendiculaires et que $BC = AE$.
\finenonce{007108}
\finexercice  


\exercice{7109, megy, 2017/01/21}

\enonce{007109}{}
% tags : rotation, complexes

Soit $BOA$ un triangle indirect isocèle en $O$, $OAC$ et $OEB$  deux triangles équilatéraux directs.

Montrer que $BC = AE$ et vérifier que l'angle des droites $(BC)$ et $(AE)$ est de $\pi/3$.
\finenonce{007109}
\finexercice  


\exercice{7110, megy, 2017/01/21}

\enonce{007110}{}
% pas génial
Soit $ABCD$ un carré, $E$ le symétrique de $C$ par rapport à $D$, $I$  le milieu de $[BC]$ et $J$ le milieu de $[DE]$.

Montrer que le triangle $AIJ$ est rectangle isocèle en $A$.
\finenonce{007110}
\finexercice  


\exercice{7111, megy, 2017/01/21}

\enonce{007111}{}
% tags : rotation, dur, utile
On considère trois droites parallèles $D_1$, $D_2$ et $D_3$. Construire un triangle équilatéral dont les sommets appartiennent respectivement à $D_1$, $D_2$ et $D_3$. 
\finenonce{007111}
\finexercice  


\exercice{7112, megy, 2017/01/21}

\enonce{007112}{}
% tags : similitude
Soit $ABCD$ et $A'B'C'D'$ deux carrés du plan inclus l'un dans l'autre. On suppose que ce sont des cartes routières de la même région, tracées à différentes échelles, posées l'une sur l'autre. Montrer qu'il existe un unique point dont les représentations sur les deux cartes coïncident. (Bonus : construire ce point.)
\finenonce{007112}
\finexercice  


\exercice{7113, megy, 2017/01/21}

\enonce{007113}{}
% similitudes, cercles circonscrit, triangle rectangle
Soit $ABC$ un triangle et $D$ le projeté orthogonal de $A$ sur $(BC)$. On considère des points $E$ et $F$ appartenant à une droite passant par $D$ tels que $(AE)$ et $(BE)$ sont perpendiculaires ainsi que $(AF)$ et $(CF)$. Enfin, on note $M$ et $N$ sont les milieux respectifs de $[BC]$ et $[EF]$.
Montrer que $(AN)$ et $(NM)$ sont perpendiculaires.
\finenonce{007113}
\finexercice  


\exercice{7114, megy, 2017/01/21}

\enonce{007114}{}
% tags : similitude
Soit $ABC$ un triangle isocèle en $A$, $A'$ et $B'$ les pieds des hauteurs issues de $A$ et $B$, $I$ le milieu de $[CB']$ et $J$ le milieu de $[A'I]$. Montrer que $(BI)$ et $(AJ)$ sont orthogonales.

\begin{center}
\includegraphics{../images/img007114-1}
\end{center}

\finenonce{007114}
\finexercice  


\exercice{7115, megy, 2017/01/21}

\enonce{007115}{}
% rotations, similitudes, complexes
% Von Aubel déguisé (pas très bien)
Soit $ABC$ un triangle direct non isocèle rectangle. Soient $ARB$, $BPC$, $CQA$ les triangles isocèles rectangles en $P$, $Q$ et $R$,  intérieurs à $ABC$.
\begin{enumerate}
\item Montrer que $\overrightarrow{AP}$ et $\overrightarrow{QR}$ ont même norme et sont orthogonaux.
\item Montrer que $(AP)$, $(BQ)$ et $(CE)$ sont concourantes.
\end{enumerate}
\finenonce{007115}
\finexercice  


\exercice{7116, megy, 2017/01/21}

\enonce{007116}{}
% similitudes, translations, ou bien complexes
% Ressemble à Van Aubel
Soit un quadrilatère convexe $ABCD$. Les points $E$, $F$, $G$, $H$ sont tels que $AEB$, $BFC$, $CGD$, $DHA$ sont rectangles isocèles en respectivement $E$, $F$, $G$, $H$. Les triangles $AEB$ et $CGD$ sont vers l'extérieur de $ABC$ , les triangles $BFC$ et $DHA$ vers l'intérieur.
Montrer que $EFGH$ est un parallélogramme.
\finenonce{007116}
\finexercice  


\exercice{7117, megy, 2017/01/21}

\enonce{007117}{}
% similitudes, ou complexes
Soient $ABCD$ et $AB'C'D'$  des carrés de même orientation et de centres respectifs $O$ et $O'$. On note $R$ et $S$ sont les milieux respectifs de $[BD']$ et $[B'D]$. Montrer que $OSO'R$ est un carré.\finenonce{007117}
\finexercice  


\exercice{7118, megy, 2017/02/08}
\enonce{007118}{Angle inscrit dans le cas limite}
% n'utilise que des triangles isocèles et rectangles
Soit $\mathcal C$ un cercle de centre $O$, $[AB]$ une corde et $\mathcal T$ la tangente de $A$. Alors $(\mathcal T,AB)=\frac12(OA,OB)$.

\begin{center}
\includegraphics{../images/img007118-1}
\end{center}
\finenonce{007118}

\finexercice  


\exercice{7119, megy, 2017/02/08}

\enonce{007119}{Construction de l'arc capable}
% application directe
On donne un segment $[AB]$ et un réel $\alpha \in ]-\pi,\pi[$. On suppose que l'on dispose également d'un triangle auxiliaire $XYZ$ avec $\widehat{(\overrightarrow{XY}, \overrightarrow{XZ})}=\alpha$, de sorte que les angles de mesure $\alpha$ sont constructibles.

 Construire le lieu des points $M$ tels que $\widehat{(\overrightarrow{MA}, \overrightarrow{MB})}=\alpha$.\finenonce{007119}
\finexercice  


\exercice{7120, megy, 2017/02/08}

\enonce{007120}{Octogone appuyé sur un segment}
% angle au centre, inscrit. application directe
 Construire un octogone convexe régulier dont un des côtés est un segment $[AB]$ donné.
\finenonce{007120}
\finexercice  


\exercice{7121, megy, 2017/02/08}

\enonce{007121}{Trapèzes inscriptibles}
% angles inscrits, facile, application directe
Montrer qu'un trapèze est isocèle si et seulement s'il est inscriptible.
\finenonce{007121}
\finexercice  


\exercice{7122, megy, 2017/02/08}

\enonce{007122}{antiparallélogramme}
% angles inscrits
Un antiparallélogramme est un quadrilatère croisé dont les cotés opposés sont deux à deux de même longueur. Soit $ABCD$ un antiparallélogramme. Montrer les assertion suivantes.

\begin{center}
\includegraphics{../images/img007122-1}
\end{center}

\begin{enumerate}
\item Les angles opposés ont la même mesure.
\item Les diagonales $(AC)$ et $(BD)$ sont parallèles.
\item La médiatrice des diagonales est un axe de symétrie de $ABCD$.
\item Deux côtés opposés ont leur point d'intersection situé sur cette médiatrice.
\item Le quadrilatère convexe $ADBC$ formé par les deux côtés non croisés et les diagonales est un trapèze isocèle.
\item $ABCD$ est inscriptible.
\end{enumerate}\finenonce{007122}
\finexercice  


\exercice{7123, megy, 2017/02/08}

\enonce{007123}{Théorème de Reim}
% application directe
Soient $\mathcal C_1$ et $\mathcal C_2$ deux cercles sécants en $A$ et $B$, et $\mathcal D_A$ (respectivement $\mathcal D_B$) une droite passant par $A$  (resp. $B$). On note $C$ et $E$ (resp. $D$ et $F$) l'intersection de $\mathcal D_A$ (resp. $\mathcal D_B$) avec 
 les deux cercles. Montrer que $(CD) // (EF)$.
 \finenonce{007123}
\finexercice  


\exercice{7124, megy, 2017/02/08}

\enonce{007124}{Bissectrices et cercle circonscrit}
% exercice simple sur les angles inscrits
Les bissectrices intérieure et extérieure en $A$ d'un triangle $ABC$ non isocèle en $A$ recoupent le cercle $\Gamma$ circonscrit à ce triangle respectivement en $I$ et $J$. Montrer que $I$ et $J$ appartiennent à la médiatrice de $[BC]$.
\finenonce{007124}

\finexercice  


\exercice{7125, megy, 2017/02/08}
\enonce{007125}{Cas limite du théorème de Reim}
% application directe
% angle inscrit et angle au centre, cas limite
% ou homothéties
Soient $\mathcal C$ et $\mathcal C'$ deux cercles tangents en un point $T$, et $\mathcal D_1$, $\mathcal D_2$ deux droites sécantes en $T$. On note $A$ et $A'$ (resp. $B$ et $B'$) les points d'intersection de $\mathcal D_1$ (resp. $\mathcal D_2$) avec $\mathcal C$ et $\mathcal C'$. Montrer que les droites $(AB)$ et $(A'B')$ sont parallèles.


\begin{center}
\includegraphics{../images/img007125-1}
\end{center}

\finenonce{007125}
\finexercice


\exercice{7126, megy, 2017/02/08}

\enonce{007126}{Théorème des trois cercles de Miquel}
% Mettre l'autre sens
% points cocycliques
Soit $ABC$ un triangle direct, et $P$, $Q$ $R$ trois points situés sur $[BC]$, $[CA]$ et $[AB]$ respectivement, et distincts des sommets. Montrer que les cercles circonscrits à $ARQ$, $BPR$ et $CQP$ sont concourants.

\begin{center}
\includegraphics{../images/img007126-1}
\end{center}

\finenonce{007126}

\finexercice  


\exercice{7127, megy, 2017/02/08}

\enonce{007127}{Ligne brisée}
% angle inscrit, aires
% plus difficile
Une ligne brisée $ABCDE$ formée d'angles de 45 degrés est inscrite dans un cercle. Elle partage le disque en deux régions, chacune d'un côté différent de la ligne. Calculer l'aire de ces deux régions.
% la moitié. 
\begin{center}
\includegraphics{../images/img007127-1}
\end{center}
\finenonce{007127}
\finexercice  


\exercice{7128, megy, 2017/02/08}

\enonce{007128}{Pentagramme}
% tags : chasse aux angles, un peu plus difficile

Soit $\mathcal C$ un cercle, $[BC]$ une corde, et $A \in \mathcal C$ tels que les arcs $AB$ et $AC$ soient égaux. Soient $[AD]$ et $[AE]$ deux autres cordes d'extrémités $A$, qui coupent $[BC]$ en $F$ et en $G$, respectivement. Montrer que $DEFG$ est inscriptible.\finenonce{007128}

\finexercice  


\exercice{7129, megy, 2017/02/08}

\enonce{007129}{}

Soient $\mathcal C_1$ et $\mathcal C_2$ deux cercles se coupant en $P$ et $Q$, et considérons une droite $\mathcal D$ coupant $\mathcal C_1$ en $A$ et $B$, et coupant $\mathcal C_2$ en  $C$ et $D$. Montrer que $(PA,PC)=(DQ,BQ)$.


Plus précisément, montrer que les angles $\widehat{APC}$ et $\widehat{DQB}$ sont égaux si $\mathcal D$ coupe le segment $[PQ]$ et que $A$, $C$, $B$ et $D$ sont alignés dans cet ordre. Que peut-on dire dans les autres cas ?
\finenonce{007129}
\finexercice  


\exercice{7130, megy, 2017/02/08}

\enonce{007130}{}
% nom :  centre similitude
% source : ? voir ausi Audin, p.91-93
On donne deux segments $[AB]$ et $[CD]$ non parallèles et de longueur différente. On admet qu'il existe une similitude directe $\phi$ envoyant $A$ sur $C$ et $B$ sur $D$. Le but de l'exercice est de construire le centre $O$ de cette similitude.
\begin{enumerate}
\item Montrer que l'angle de la similitude est $\widehat{(\overrightarrow{AB},\overrightarrow{CD}}$.
\item On note $Q = (AB) \cap (CD)$. Montrer que $AQCO$ est inscriptible.
\item Terminer le raisonnement et construire $O$.
\item Que faire si les segments sont parallèles ? De même longueur ? Réfléchir à d'autres méthodes pour construire le centre, sans utiliser de cercles.
\end{enumerate}\finenonce{007130}
\finexercice  


\exercice{7131, megy, 2017/02/08}

\enonce{007131}{Triangle orthique}
% joli, points cocycliques
Soit $ABC$ un triangle non rectangle et $A'$, $B'$, $C'$ les pieds des hauteurs. Montrer que les hauteurs de $ABC$ sont des bissectrices du triangle $A'B'C'$, dit \emph{triangle orthique}. 
\finenonce{007131}
\finexercice  


\exercice{7132, megy, 2017/02/08}

\enonce{007132}{Ptolémée, preuve géométrique du sens direct}
% similitudes

On considère l'énoncé suivant :

\begin{theoreme}[Ptolémée] Soit $ABCD$ un quadrilatère convexe, direct. Alors $A, B, C, D$ sont cocycliques si et seulement si $ AC\cdot BD = AB\cdot CD + BC\cdot AD$.
\end{theoreme}

L'objectif de cet exercice est de prouver le sens direct du théorème.

On suppose $A, B, C, D$ cocycliques.
\begin{enumerate}
\item Faire une figure au brouillon et montrer que $\widehat{BAC} = \widehat{BDC}$ et trois relations similaires sur d'autres angles.
\item Soit $K$ le point de la diagonale $[AC]$ tel que $\widehat{ABK} = \widehat{DBC}$. Faire une figure et construire $K$ en expliquant (faire la figure de telle sorte que $K$ soit lisible).
\item Montrer que les triangles $ABK$ et $DBC$ sont semblables, de même que $ABD$ et $KBC$, par des similitudes dont on précisera les centres et les rapports. Note : il suffit pour cela de montrer qu'ils ont mêmes angles. En déduire des relations sur les côtés de ces triangles.
\item Conclure.

\end{enumerate}
\finenonce{007132}


\finexercice  


\exercice{7133, megy, 2017/02/08}

\enonce{007133}{Une application de Ptolémée}
% source : Deschamps chap 2 ex 8
Soit $ABC$ un triangle équilatéral et $M$ un point du cercle circonscrit appartenant à l'arc $BC$ ne contenant pas $A$. Montrer que $MA= MB+MC$.\finenonce{007133}
\finexercice  


\exercice{7134, megy, 2017/02/08}

\enonce{007134}{Trisection}
% cocyclicité
\begin{enumerate}
\item Soit $\Gamma$ un cercle, $O$ son centre et $M$ un point n'appartenant pas à $\Gamma$. Deux sécantes issues de $M$ coupent $\Gamma$ respectivement en $A$ et $B$, et en $C$ et $D$. Démontrer l'égalité:
\[2 (\overrightarrow{MA},\overrightarrow{MC}) = (\overrightarrow{OB},\overrightarrow{OD}) -  (\overrightarrow{OC},\overrightarrow{OA}).\]
\item Soient $A$ et $B$ deux points d'un cercle $\Gamma$ de centre $O$ et de rayon $r$. Sur la droite $(OA)$, soit $C$ le point extérieur au cercle tel que la droite $(CB)$ recoupe le cercle en un point $M$ vérifiant $MC=r$. (On ne demande pas de construire ce point, le placer approximativement sur la figure.) Montrer que $\widehat{ACB} = \frac13 \widehat{AOB}$.
\end{enumerate}
\finenonce{007134}
\finexercice  


\exercice{7135, megy, 2017/02/08}

\enonce{007135}{Problème \og DPP\fg }

Soit $\mathcal D$ une droite et $A$ et $B$ deux points situés d'un seul côté de $\mathcal D$. L'objectif est de construire un cercle passant par les deux points et tangent à la droite. 

\begin{enumerate}
\item Construire un tel cercle si les droites $(AB)$ et $\mathcal D$ sont parallèles. Dans la suite, on supposera qu'elles sont sécantes.
\item (Analyse) Soit $\mathcal C$ un tel cercle et $T$ son point de tangence avec $\mathcal D$. Montrer que $(AB,AT) = (TB,\mathcal D)$. % angle inscrit avec le cas limite.
\item (Synthèse) Soit $I$ le point d'intersection de $(AB)$ avec $\mathcal D$, $B'$ le symétrique de $B$ par rapport à $I$, et $B''$ le symétrique de $B$ par rapport à $\mathcal D$. Montrer que le cercle circonscrit à $AB'B''$ (de diamètre $[AB']$ dans le cas où $B'=B''$) coupe $\mathcal D$ en deux points qui conviennent pour le choix de $T$.
\end{enumerate}

\finenonce{007135}
\finexercice  


\exercice{7136, megy, 2017/02/08}

\enonce{007136}{Bissectrices d'un quadrilatère convexe}
Soit $ABCD$ un quadrilatère convexe. On note $\mathcal B_A$ (resp. $\mathcal B_B$, $\mathcal B_C$, $\mathcal B_D$) la bissectrice intérieure en $A$ (resp. en $B$, $C$, $D$). Soient $I=\mathcal B_A\cap \mathcal B_B$, $J=\mathcal B_B\cap \mathcal B_C$, $K=\mathcal B_C\cap \mathcal B_D$ et $L=\mathcal B_D\cap \mathcal B_A$. Montrer que $IJKL$ est inscriptible.

\begin{center}
\includegraphics{../images/img007136-1}
\end{center}

\finenonce{007136}
\finexercice  


\exercice{7137, megy, 2017/02/08}

\enonce{007137}{Médiatrices d'un quadrilatère convexe}
Soit $ABCD$ un quadrilatère convexe. Que dire du quadrilatère formé par les intersections des médiatrices des quatre côtés ? En particulier, que dire si $ABCD$ est un parallélogramme ? Un losange ?
\finenonce{007137}
\finexercice  


\exercice{7138, megy, 2017/02/08}

\enonce{007138}{Carré à retrouver}
On considère un carré $ABCD$ et on place quatre points $E$, $F$, $G$, et $H$ sur les côtés de ce carré (en-dehors des sommets). Puis, on efface le carré. L'objectif est de reconstruire le carré en utilisant le théorème de l'angle inscrit.

Si $A$ est le sommet entre $E$ et $F$, montrer que la diagonale du carré partant de $A$ passe par l'intersection du cercle de diamètre $[EF]$ avec la médiatrice de $[EF]$. % angle inscrit
En déduire une construction des diagonales du carré, puis du carré.\finenonce{007138}
\finexercice  


\exercice{7139, megy, 2017/04/05}
\enonce{007139}{Symétrique de l'orthocentre}

Soit $ABC$ un triangle, $H$ son orthocentre et $\mathcal C$ son cercle circonscrit. La hauteur issue de $B$ recoupe $\mathcal C$ en $H'$. Montrer que $H'$ est le symétrique de $H$ par rapport à la droite $(AC)$. 

\begin{center}
\includegraphics{../images/img007139-1}
\end{center}

\finenonce{007139}


\finexercice


\exercice{7140, megy, 2017/04/05}
\enonce{007140}{Un théorème de Brahmagupta}
Soit $ABCD$ un quadrilatère convexe inscriptible dont les diagonales sont perpendiculaires, et soit $O$ leur point d'intersection. Soit $H$ le projeté orthogonal de $O$ sur $[CD]$, et $I$ l'intersection de $(OH)$ avec $[AB]$. L'objectif est de montrer que $I$ est le milieu de $[AB]$. 


\begin{center}
\includegraphics{../images/img007140-1}
\end{center}

\begin{enumerate}
\item Montrer qu'il est suffisant d'établir que $IO=IA$.
\item Conclure.
\end{enumerate}
\finenonce{007140}

\finexercice


\exercice{7141, megy, 2017/04/05}

\enonce{007141}{Puissance d'un point par rapport à un cercle}
Soit $\mathcal C$ un cercle et $P$ un point du plan.

\begin{enumerate}
\item Soit $\mathcal D$  une droite passant par $P$ et intersectant le cercle en deux points $A$ et $B$. Montrer que la quantité $\overline{PA}\cdot\overline{PB}:= \overrightarrow{PA} \cdot \overrightarrow{PB}$ ne dépend pas de la droite choisie. On l'appelle la \emph{puissance du point $P$ par rapport au cercle $\mathcal C$}, et on la note $p_{\mathcal C}(P)$.
%Elle est positive si $P$ est à l'extérieur du cercle, négative s'il est à l'intérieur, et nulle si $P$ appartient au cercle. 
\item Avec les mêmes notations, si $P$ est à l'extérieur du cercle et $\mathcal D$ est tangente au cercle en $T$, montrer que $p_{\mathcal C}(P) = PT^2$.
\item Montrer que $p_{\mathcal C}(P)$ vaut $OP^2 - r^2$, où $r$ est le rayon du cercle.
\item Soit $\lambda \in \R$. Déterminer l'ensemble des points dont la puissance par rapport au cercle vaut $\lambda$.
\item Réciproquement, soient $(AB)$ et $(CD)$ deux droites se coupant en $P$. On suppose que $\overline{PA}\cdot \overline{PB} = \overline{PC}\cdot \overline{PD}$. Montrer que $ABCD$ est inscriptible.

\end{enumerate}
\finenonce{007141}
\finexercice


\exercice{7142, megy, 2017/04/05}

\enonce{007142}{}
Soient deux cercles se coupant en deux points distincts $A$ et $B$, et $\mathcal D$ une tangente commune, touchant les cercles en $C$ et $D$. Montrer que $(AB)$ coupe $[CD]$ en son milieu.
\begin{center}
\includegraphics{../images/img007142-1}
\end{center}
\finenonce{007142}

     
\finexercice  


\exercice{7143, megy, 2017/04/05}
\enonce{007143}{}
% puissance, triangles rectangles, semblables
Soit $ABC$ un triangle rectangle en $A$, et $H$ le pied de la hauteur issue de $A$. Montrer $AH^2 = HB\cdot HC$.
\finenonce{007143}
\finexercice


\exercice{7144, megy, 2017/04/05}
\enonce{007144}{}
Soit $\mathcal C$ un cercle et $P$ un point du plan. On considère une droite  $\mathcal D$ passant par $P$ et intersectant le cercle en au moins un point $E$. Soit $F$ le point diamétralement opposé à $E$. Montrer que $p_\mathcal C(P) = \overrightarrow{PE}\cdot \overrightarrow{PF}$.
\begin{center}
\includegraphics{../images/img007144-1}
\end{center}
\finenonce{007144}
\finexercice  


\exercice{7152, megy, 2017/05/13}
\enonce{007152}{Nombre d'or}
Soit $ABCD$ un carré direct de côté $1$, et soient $E$  et $F$ deux points tels que $AEFD$ soit un rectangle direct, avec $AE > 1$. Dans la suite on note $l=AE$.


Montrer que qu'il existe une similitude directe $s$ envoyant $A$ (resp. $E$, $F$ et $D$) sur $C$ (resp.  $B$, $E$ et $F$) ssi $l$ est égal au nombre d'or $(1+\sqrt 5)/2$.
\finenonce{007152}

\finexercice 


\exercice{7153, megy, 2017/05/13}
\enonce{007153}{}
Soit $ABC$ un triangle rectangle en $A$ et non isocèle. La médiatrice de $[BC]$ recoupe le demi-cercle circonscrit en $I$. On considère deux points $D\in [AB]$ et $E \in [AC]$ tels que $BD=CE$. Montrer que $IDE$ est rectangle isocèle en $I$.

\begin{center}
\includegraphics{../images/img007153-1}
\end{center}

\finenonce{007153}

\finexercice 


\exercice{7154, megy, 2017/05/13}
\enonce{007154}{}
%[d'après bac Amérique du sud 2004]
Soient $A$ et $B_0$ deux points et soit $s$ la similitude directe de centre $A$, de rapport $\frac12$ et  d'angle $\frac{3\pi}{4}$. On considère la suite de points $(B_n)_{n\in\N}$ définie par la relation de récurrence $\forall n\in \N,\: B_{n+1}=s(B_n)$.


\begin{enumerate}
\item Faire une figure avec $AB_0=8$ et placer les points $B_n$ jusqu'à $n=4$.
\item Montrer que pour tout  $n\in\N$, les triangles $AB_nB_{n+1}$ et $AB_{n+1}B_{n+2}$ sont semblables.
\item Dans la suite, on considère le sous-ensemble du plan $S = \bigcup_{n\in \N} [B_nB_{n+1}]$. C'est une spirale polyédrale. Sa longueur est-elle finie ou infinie ? Dans le premier cas, calculer sa longueur.
\end{enumerate}
\finenonce{007154}

\finexercice 


\exercice{7155, megy, 2017/05/13}
\enonce{007155}{Point intérieur de Vecten}
% composition de similitudes
% hauteurs, orthocentre
Soit $ABC$ un triangle direct non rectangle isocèle, et soit $P$ (resp. $Q$, $R$) tel que $BCP$ (resp. $CAQ$, $ABR$) soit direct isocèle rectangle en $P$ (resp. $Q$ et $R$). 

\begin{center}
\includegraphics{../images/img007155-1}
\end{center}

\begin{enumerate}
\item Montrer que $\overrightarrow{AP}$ et $\overrightarrow{QR}$ sont orthogonaux et de même norme.
\item Montrer que les droites $(AP)$, $(BQ)$ et $(CE)$ sont concourantes.
\end{enumerate}

\finenonce{007155}

\finexercice 


\exercice{7156, megy, 2017/05/13}
\enonce{007156}{}
% concours général 88, cocyclicité, angle d'une similitude directe
Soient $M$, $M_1$, $M_2$, $M_3$ et $M_4$ cinq points distincts sur un cercle $\mathcal C$. Montrer que le produit des distances de $M$ aux droites $(M_1M_2)$ et $(M_3M_4)$ est égal au produit des distances de $M$ aux droites $(M_1M_3)$ et $(M_2M_4)$.
\begin{center}
\includegraphics{../images/img007156-1}
\end{center}

\finenonce{007156}

\finexercice 


\exercice{7157, megy, 2017/05/13}
\enonce{007157}{}
% groupes de frises
\begin{enumerate}
\item Soit $\mathcal S \subset \R^2$ le graphe de la fonction sinus. Décrire les isométries de $\mathcal S$. 
\item Même question pour le graphe de la fonction tangente.
\end{enumerate}
\finenonce{007157}
\finexercice  


\exercice{7158, megy, 2017/05/13}
\enonce{007158}{Isométries d'un triangle}
Soit $\mathcal T=ABC$ un triangle et soit $g$ une isométrie de $\mathcal T$, c'est-à-dire une isométrie du plan telle que $g(\mathcal T)=\mathcal T$.
\begin{enumerate}
\item Montrer que $g$ envoie les sommets sur les sommets.
\item Montrer que $g$ admet au moins un point fixe.
\item Montrer que $g$ n'est ni une translation ni une réflexion glissée.
\end{enumerate}
\finenonce{007158}

\finexercice  


\exercice{7159, megy, 2017/05/13}
\enonce{007159}{Isométries d'un triangle équilatéral}
Soit $\mathcal T=ABC$ un triangle équilatéral et $G=Isom(\mathcal T)$ son groupe d'isométries.
\begin{enumerate}
\item Trouver six isométries laissant $\mathcal T$ invariant.
\item Montrer qu'une isométrie de $\mathcal T$ doit envoyer un sommet sur un sommet.
\item \'Ecrire un morphisme injectif $\phi$ entre $G$ et $\mathfrak S_3$.
\item Montrer qu'il est bijectif.
\item Décrire le groupe $H=Isom^+(\mathcal T)$ et écrire un isomorphisme  entre ce groupe et $\Z/3\Z$. À quel sous-groupe de $\mathfrak S_3$ correspond $H$ ?
\end{enumerate}
\finenonce{007159}

\finexercice  


\exercice{7160, megy, 2017/05/13}
\enonce{007160}{Isométries d'un triangle isocèle}
Soit $\mathcal T=ABC$ un triangle isocèle non équilatéral. Déterminer son groupe d'isométries ainsi qu'un isomorphisme entre ce groupe et un groupe classique.
\finenonce{007160}

\finexercice  


\exercice{7161, megy, 2017/05/13}
\enonce{007161}{Isométries d'un carré}
Soit $\mathcal C=ABCD$ un carré plein du plan (c'est-à-dire l'enveloppe convexe de ses sommets) et $O$ son centre.
\begin{enumerate}
\item Soit $g\in \operatorname{Isom}(\mathcal C)$ une isométrie du carré, c'est-à-dire une isométrie du plan préservant le carré. Montrer qu'elle permute les sommets, puis que c'est une rotation de centre $O$ ou bien une réflexion d'axe contenant $O$.
\item Déterminer entièrement  le groupe $G=\operatorname{Isom}(\mathcal C)$.
\item  Décrire le groupe $H=\operatorname{Isom}^+(\mathcal C)$ des isométries directes du carré et écrire un isomorphisme entre ce groupe et $\Z/4\Z$.
\end{enumerate}
\finenonce{007161}

\finexercice  


\exercice{7162, megy, 2017/05/13}
\enonce{007162}{Isométries d'un rectangle}
Soit $\mathcal C=ABCD$ un rectangle plein, non carré. Déterminer son groupe d'isométries, et préciser un isomorphisme entre ce groupe et un groupe usuel abstrait.
\finenonce{007162}

\finexercice  


\exercice{7165, megy, 2017/06/11}
\enonce{007165}{}

Soit $ABC$ un triangle équilatéral direct, et $P$ un point n'appartenant pas aux droites $(AB)$, $(BC)$ et $CA)$. La parallèle à $(BC)$ (resp. à $(CA)$ et $(AB)$) passant par $P$ coupe $(AB)$ (resp. $(BC)$ et $(CA)$) en $U$ (resp. $V$, $W$).

\begin{center}
\includegraphics{../images/img007165-1}
\end{center}

\begin{enumerate}
\item Montrer que $AUPW$ et $PCVW$ sont inscriptibles.
\item Dresser une liste d'égalités d'angles (de droites) que l'on peut  déduire.
\item Faire une (grande) figure dans le cas où $P$ est sur le cercle circonscrit à $ABC$.
\item Montrer que les points $U$, $V$ et $W$ sont alignés ssi $ABCP$ est inscriptible.
\end{enumerate}
\finenonce{007165}

\finexercice  

\exercice{7167, megy, 2017/07/08}
\enonce{007167}{}
% rotations
Soient $AOB$ et $COD$ deux triangles directs, isocèles rectangles en $O$. Soient $I$, $J$, $K$ et $L$ les milieux des segments $[AB]$, $[BC]$, $[CD]$ et $[DA]$.

\begin{center}
\includegraphics{../images/img007167-1}
\end{center}

\begin{enumerate}
\item Montrer que $(AC)\bot (BD)$ et que $AC=BD$.
\item Montrer que $IJKL$ est un carré. 
\end{enumerate}
\finenonce{007167}
\finexercice


\exercice{7168, megy, 2017/07/08}
\enonce{007168}{}
% cocyclicité, bissectrice, puissance
Soit $ABC$ un triangle non isocèle en $A$. % acutangle? Pas besoin
La bissectrice intérieure $\Delta$ de $\widehat{BAC}$ coupe $[BC]$ en $A_1$ et le cercle circonscrit à $ABC$ en $A_2$.
\begin{enumerate}
\item Soit $D$ le symétrique de $B$ par rapport à $\Delta$. Justifier que $D \in (AC)$.
\item Montrer que $A_1A_2CD$ est inscriptible.
\item Montrer que $AA_1\cdot AA_2 = AB\cdot AC$. % puissance
\end{enumerate}
\begin{center}
\includegraphics{../images/img007168-1}
\end{center}
\finenonce{007168}
\finexercice


\exercice{7249, mourougane, 2021/08/10}
\enonce{007249}{}
Soient \(P\), \(Q\), \(R\) trois points du plan. Dans cet exercice, 
on notera \(\vec{u}\cdot\vec{v}\) le produit scalaire de deux vecteurs 
\(\vec{u}\) et \(\vec{v}\).
\begin{enumerate}
\item Montrer que, pour tout \( \lambda \in \Rr\), on~a
\[
(\vec{QP} + \lambda \vec{QR})^2
= \vec{QP}^2 + 2 \lambda \vec{QP}\cdot\vec{QR} + \lambda^2 \vec{QR}^2.
\]
\item En considérant le discriminant du polynôme (en la variable \( \lambda\)) 
de droite dans l'égalité précédente, montrer que
\[
\left\lvert\vec{QP}\cdot\vec{QR}\right\rvert \leq QP \times  QR.
\]
\item Montrer que
\[
\vec{PR}^2 = \vec{QP}^2 - 2 \vec{QP}\cdot\vec{QR} + \vec{QR}^2.
\]
\item En déduire que
\[
PR \leq PQ + QR.
\]
\item Montrer que \(PR = PQ + QR\) si et seulement si \(Q \in [PR]\).
\item On considère maintenant quatre points \(P\), \(Q\), \(R\) 
et \(S\). Montrer que
\[
PS \leq PQ + QR + RS
\]
et caractériser les configurations de quatre points \(P\), \(Q\), 
\(R\) et \(S\) qui vérifient l'égalité \(PS = PQ + QR + RS\).
\end{enumerate}
\finenonce{007249}
\finexercice


\exercice{7250, mourougane, 2021/08/10}
\enonce{007250}{Diagramme de Voronoï. Médiatrice, diagramme de 2 points}

% cf. 978-0-387-69216-6.pdf page 529

Un \emph{diagramme de Voronoï} est une famille de parties du plan 
(ou de l'espace, mais dans cet exercice on se limitera au plan) et 
de points associés telle que:
\begin{itemize}
\item chaque partie du plan a un unique point associé, qui est 
contenu dedans;
\item chaque partie est exactement égale à l'ensemble des points 
du plan qui sont plus proches du point associé à cette partie que 
des points associés aux autres parties.
\end{itemize}
Autrement dit, c'est une famille \((A_i, P_i)_{i \in I}\), où:
\begin{itemize}
\item \(I\) est un ensemble;
\item pour tout \(i \in I\), \(A_i\) est une partie (i.e. un 
sous-ensemble) du plan et \(P_i \in A_i\);
\item pour tout \(i \in I\), on~a (en notant \( \mathcal{P}\) le plan):
\[
A_i = \left\{ \; Q \in \mathcal{P} \; \middle/ \; \forall j \in I \setminus \{i\} \; P_iQ \leq P_jQ \; \right\}.
\]
\end{itemize}
Les parties \(A_i\) sont appelées les \emph{cellules} du diagramme 
de Voronoï. Le point \(P_i\) associé à la cellule \(A_i\) est appelé 
le \emph{germe} de la cellule.

Les diagrammes de Voronoï sont un outil utile pour représenter les 
zones de couverture d'antennes radio, ou pour étudier l'implantation 
d'écoles, d'hôpitaux, de bureaux de poste, etc, dans une région.


\begin{enumerate}
\item Soient \(A\) et \(B\) deux points du plan. Montrer que 
l'ensemble des points équidistants de \(A\) et \(B\) (autrement 
dit, l'ensemble des points \(P\) du plan tels que \(PA = PB\)) 
est une droite (qu'on notera \( \Delta\), et qui est appelée la 
\emph{médiatrice} du segment \([AB]\)).
\item Montrer que la droite \( \Delta\) est orthogonale à la 
droite \((AB)\).
\item Montrer que l'ensemble des points \(P\) du plan tels 
que \(PA \leq PB\) est le demi-plan de frontière \( \Delta\) contenant \(A\).
\item Quel est le diagramme de Voronoï d'un ensemble de deux 
points distincts?
\end{enumerate}
\finenonce{007250}
\finexercice


\exercice{7251, mourougane, 2021/08/10}
\enonce{007251}{Diagramme de Voronoï}
Soit \((A_i, P_i)_{i \in I}\) un diagramme de Voronoï. Si \(i, j \in I\) 
avec \(i \neq j\), notons \( \Pi_{i,j}\) le demi-plan de frontière la 
médiatrice du segment \([P_iPj]\) et qui contient le point \(P_i\). 
\begin{enumerate}
\item Montrer que, pour tout \(i \in I\), on~a:
\[
A_i = \bigcap_{j \in I \setminus \{i\}} \Pi_{i,j}.
\]
\item Dessiner le diagramme de Voronoï de trois points formant un 
triangle équilatéral.
\end{enumerate}
\finenonce{007251}
\finexercice


\exercice{7252, mourougane, 2021/08/10}
\enonce{007252}{Diagramme de Voronoï}
\begin{enumerate}
\item Que dire de la frontière entre deux cellules de Voronoï ?
\item Que dire du point commun à trois cellules de Voronoï, appelé 
sommet de Voronoï?
\item Que dire du cercle centré en un sommet de Voronoï et passant 
par un germe d'une des trois cellules?
\item Ajouter un point au triangle équilatéral de l'exercice précédent, 
et tracer le nouveau diagramme de Voronoï.
\item Ajouter un cinquième point très proche d'un sommet de Voronoï 
et tracer le nouveau diagramme de Voronoï. Toutes les cellules 
ont-elles changé ?
\end{enumerate}
\finenonce{007252}
\finexercice


\exercice{7253, mourougane, 2021/08/10}
\enonce{007253}{Convexité des cellules de Voronoï}
On dit qu'une partie \(X\) du plan (ou de l'espace) est \emph{convexe} 
si elle vérifie:
$$
 \forall (P,Q) \in X^2 \quad [PQ] \subset X,
$$
autrement dit, pour tout couple \((P,Q)\) de points de \(X\), 
le segment \([PQ]\) tout entier est contenu dans \(X\).

\begin{enumerate}
\item Montrer qu'une intersection de parties convexes du plan est 
convexe.
\item En déduire que les cellules d'un diagramme de Voronoï sont 
convexes.
\end{enumerate}
\finenonce{007253}
\finexercice


\exercice{7254, mourougane, 2021/08/10}
\enonce{007254}{}
On rappelle qu'un quadrilatère d'un espace euclidien \(E\) est un 
\emph{parallélogramme} si ses diagonales se coupent en leur milieu, 
appelé centre du parallélogramme.

Cette définition est aussi valable en dimension 1 et pour les cas où 
deux sommets coïncident. (Dans ces cas, le parallélogramme est plat).

On dit que deux bipoints \((A,B)\) et \((C,D)\) sont \emph{équipollents} 
si le quadrilatère \((ABDC)\) est un parallélogramme.

\begin{enumerate}
\item Vérifier que pour tout couple de points \((A,B)\), les bipoints 
\((A,B)\) et \((A,B)\) sont équipollents. On dit alors que la relation 
d'équipollence est \emph{réflexive}.
\item Montrer que pour tous bipoints \((A,B)\) et \((C,D)\), si les 
bipoints \((A,B)\) et \((C,D)\) sont équipollents alors les bipoints 
\((C,D)\) et \((A,B)\) le sont aussi. On dit alors que la relation 
d'équipollence est \emph{symétrique}.
\item Démontrer que la relation d'équipollence est \emph{transitive}, 
c'est à dire que pour tous triplets \((A,B)\), \((C,D)\) et \((F,G)\) 
de bipoints, si les bipoints \((A,B)\) et \((C,D)\) sont équipollents 
et si les bipoints \((C,D)\) et \((F,G)\) sont équipollents alors les 
bipoints \((A,B)\) et \((F,G)\) le sont aussi. (Indication : dans le 
cas où le quadrilatère \((ABGF)\) n'est pas plat, on pourra considérer 
la droite joignant les centres des parallélogrammes \((ABDC)\) et 
\((CDGF)\); dans le cas où le quadrilatère \((ABGF)\) est plat, on 
pourra utiliser le théorème de Thalès.)
\item On résume les trois propriétés précédentes en disant que la 
relation d'équipollence est \emph{une relation d'équivalence}.
La \emph{classe d'équipollence} du bipoint \((A,B)\) est par 
définition l'ensemble des bipoints équipollents à \((A,B)\).
Elle est appelée \emph{vecteur} et notée \(\vec{AB}\).
Si \((C,D)\) est équipollent à \((A,B)\), on dit que \((C,D)\) est 
un \emph{représentant} de \(\vec{AB}\). Montrer qu'étant donné un 
point \(A\) et un vecteur \(\vec{u}\), il existe un unique point \(B\) 
tel que \(\vec{AB} = \vec{u}\). On notera \(B = t_{\vec{u}}(A)\).
\item Étant donnés un point \(A\) et un représentant \((F, G)\) 
du vecteur \(\vec{u}\), construire à la règle et au compas le 
point \(t_{\vec{u}}(A)\).
\item Montrer que si deux bipoints \((A,B)\) et \((C,D)\) sont 
équipollents, alors les bipoints \((A,C)\) et \((B,D)\) le sont aussi.
\item On définit la somme de deux vecteurs \(\vec{u}\) et \(\vec{v}\) 
par le procédé suivant:
\begin{itemize}
\item On choisit un point \(A\).
\item On détermine le point \(B\) tel que \(\vec{AB} = \vec{u}\).
\item On détermine le point \(C\) tel que \(\vec{BC} = \vec{v}\).
\item On définit \(\vec{u} + \vec{v} := \vec{AC}\).
\end{itemize}
Montrer que la \emph{somme} ainsi définie est indépendante du choix du 
point de base \(A\), c'est à dire, montrer que si on choisit un autre 
point \(A'\) comme point de base, le bipoint \((A',C')\) construit 
alors est équipollent au bipoint \((A,C)\) construit en partant du 
point \(A\). (Indication : On pourra montrer que \((A,A')\) est 
équipollent à \((C,C')\).)
\end{enumerate}
\finenonce{007254}
\finexercice


\exercice{7258, mourougane, 2021/08/10}
\enonce{007258}{}
Construire à la règle et au compas un angle de mesure $\pi/12$.
\finenonce{007258}
\finexercice


\exercice{7259, mourougane, 2021/08/10}
\enonce{007259}{}
\begin{enumerate}

\item On dit qu'un angle est inscrit dans un cercle si son sommet appartient à ce cercle. Démontrer le théorème des angles inscrits :\\

Deux angles de vecteurs inscrits dans un cercle interceptant le même arc de cercle sont de même mesure.

\item Soient $\mathcal{C}_1$ et $\mathcal{C}_2$ deux cercles ayant deux points d'intersection $I$ et $J$. Soient $A$ et $M$ deux points distincts de $\mathcal{C}_1$ (et différents de $I$ et $J$). 
On note $B$ le point d'intersection de la droite $(AJ)$ avec $\mathcal{C}_2$ et $N$ le point d'intersection de la droite $(MJ)$ avec $\mathcal{C}_2$.\\

En considérant la somme des mesures des angles des triangles $AIB$ et $MIN$, montrer que $\text{Mes}\widehat{AIB}=\text{Mes}\widehat{MIN}$.

\end{enumerate}
\finenonce{007259}
\finexercice


\exercice{7260, mourougane, 2021/08/10}
\enonce{007260}{}
Soit $E$ un plan euclidien orienté, muni d'un repère $(O,\vec{\imath},\vec{\jmath})$ orthonormé direct.
\begin{enumerate}
\item Soit $n$ un entier naturel supérieur à $3$.
Exprimer à l'aide des fonctions trigonométriques $\cos$ et $\sin$,
 le périmètre $p_n$ d'un polygone régulier à $n$ côtés inscrit dans le cercle trigonométrique
(c'est à dire le cercle de centre $0$ et de rayon $1$.)
\item On rappelle que pour tout $\theta\in ]0,\pi/2[$, $$\theta\cos\theta\leq \sin\theta\leq \theta.$$
Montrer que la suite $(p_n)_{n\in\mathbb{N}}$ admet une limite et déterminer cette limite.
\end{enumerate}
\finenonce{007260}
\finexercice


\exercice{7261, mourougane, 2021/08/10}
\enonce{007261}{}
\begin{enumerate}

\item Soit $\mathcal{C}$ un cercle de centre $O$, $P$ et $Q$ deux points de $\mathcal{C}$ non diamétralement opposés. Calculer $\text{Mes}\widehat{OPQ}$ en fonction de $\text{Mes}\widehat{POQ}$. 
\item Soit $d$ la droite perpendiculaire à $(OP)$ passant par $P$. 
En calculant la distance entre $O$ et tout point $M$ de la droite $d$,
montrer que $P$ est l'unique point d'intersection entre $d$ et $\mathcal{C}$. 
La droite $d$ est appelée \emph{tangente au cercle $\mathcal{C}$ en $P$}. \\



\exercice{7262, mourougane, 2021/08/10}
\enonce{007262}{}
Soit $E$ un plan et $(0,\vec{\imath},\vec{\jmath})$ un repère.
On considère l'application 
$$
\phi : (\vec{u}\begin{pmatrix}x\\y\end{pmatrix} , \vec{u'}\begin{pmatrix}x'\\y'\end{pmatrix})\mapsto xx'+3yy'.
$$

\begin{enumerate}

\item Montrer que $\phi$ est un produit scalaire.
\item Déterminer un repère $(O,\vec{I},\vec{J})$ du plan orthonormé pour ce produit scalaire.
\item Calculer $\phi (X\vec{I}+Y\vec{J}, X'\vec{I}+Y'\vec{J})$.

\end{enumerate}
\finenonce{007262}
\finexercice


\exercice{7263, mourougane, 2021/08/10}
\enonce{007263}{}
Soit $E$ un plan et $(0,\vec{\imath},\vec{\jmath})$ un repère.
Soient $\alpha$ et $\beta$ des nombres réels. On considère l'application 
$$
\varphi : (\vec{u}\begin{pmatrix}x,y\end{pmatrix} , \vec{u'}\begin{pmatrix}x',y'\end{pmatrix})\mapsto \alpha xx'+\beta yy'
$$

\begin{enumerate}

\item Montrer que $\varphi$ est un bilinéaire et symétrique.
\item Montrer que $\varphi$ est un produit scalaire si, et seulement si, $\alpha$ et $\beta$ sont strictement positifs.
\item Montrer que l'application 
$$
\psi : (\vec{u}\begin{pmatrix}x\\y\end{pmatrix} , \vec{u'}\begin{pmatrix}x'\\y'\end{pmatrix})\mapsto 2xx'+yy' +xy'+x'y 
$$
est un produit scalaire. \emph{Indication :} Chercher une identité remarquable.
\end{enumerate}
\finenonce{007263}
\finexercice


\exercice{7264, mourougane, 2021/08/10}
\enonce{007264}{}
Soit $E$ un plan euclidien orienté muni d'un repère orthonormé direct $(0,\vec{I},\vec{J})$.
Soit $\vec{u}$ un vecteur de norme $1$ et de coordonnées $\begin{pmatrix}X\\Y\end{pmatrix}$ dans le repère 
$(0,\vec{I},\vec{J})$.
\begin{enumerate}
\item Calculer le produit scalaire $\vec{I}\cdot\vec{u}$.
 \item Calculer, à l'aide de la définition de la fonction cosinus, la quantité
$||\vec{I}||\ ||\vec{u}||\cos\widehat{(\vec{I},\vec{u})}$.
\item Reprendre les calculs précédents, sans l'hypothèse que $\vec{u}$ est de norme $1$.

\end{enumerate}
\finenonce{007264}
\finexercice


\exercice{7265, mourougane, 2021/08/10}
\enonce{007265}{}
Montrer que si $A$ est inclus dans une partie $B$ d'un plan euclidien $E$,
 alors $$\mathcal{A}(A)\leq\mathcal{A}(B).$$
\finenonce{007265}
\finexercice


\exercice{7266, mourougane, 2021/08/10}
\enonce{007266}{}
Soit $ABC$ un triangle dans un plan euclidien orienté $E$.
On notera $a=BC, b=CA$ et $c=AB$ et $\hat{A}=\widehat{(BAC)}$.
Notons $p$ la moitié de son périmètre. 
\begin{enumerate}
\item Montrer que \begin{eqnarray*}
2p(p-a)=bc+\vec{AB}\cdot\vec{AC}\\
          2(p-b)(p-c)=bc-\vec{AB}\cdot\vec{AC}
         \end{eqnarray*}
\item En déduire que l'aire du triangle $ABC$ est $\sqrt{p(p-a)(p-b)(p-c)}$.
\end{enumerate}
\finenonce{007266}
\finexercice


\exercice{7267, mourougane, 2021/08/10}
\enonce{007267}{}
Soit \(ABC\) un triangle. Notons \(\Delta_A\), \(\Delta_B\) et 
\(\Delta_C\) les bissectrices des angles en \(A\), \(B\), \(C\) 
respectivement. On note aussi \(a\), \(b\), \(c\) les longueurs 
\(BC\), \(CA\) et \(AB\) respectivement, et \(\alpha\), \(\beta\), 
\(\gamma\) les mesures des angles (non orientés) \(\widehat{CAB}\), 
\(\widehat{ABC}\) et \(\widehat{BCA}\) respectivement.
\begin{enumerate}
\item Montrer que si \(\Delta_A\) est parallèle à \(\Delta_B\), 
alors le triangle \(ABC\) est plat. On supposera dans la suite de 
l'exercice que ce triangle n'est pas plat.
\item Montrer que pour tout point \(P\) de \(\Delta_A\), la distance 
de \(P\) à la droite \((AB)\) est égale à la distance de \(P\) à 
la droite \((AC)\).
\item En déduire que le point d'intersection de \(\Delta_A\) et de 
\(\Delta_B\), que l'on notera \(\Omega\), est équidistant de 
\((AB)\), \((BC)\) et \((CA)\). On admettra que cela permet de démontrer 
que \(\Omega \in \Delta_C\).
\item On note \(A'\) le projeté orthogonal de \(\Omega\) sur la droite 
\((BC)\), c'est-à-dire l'unique point \(A' \in (BC)\) tel que 
\((\Omega A') \perp (BC)\). De même, on note \(B'\) le projeté orthogonal 
de \(\Omega\) sur \((CA)\) et \(C'\) le projeté orthogonal de \(\Omega\) 
sur \((AB)\). Montrer que \(AB' = AC'\).
\item On admet que \(A' \in [BC]\), \(B' \in [CA]\), \(C' \in [AB]\). 
Montrer que
\[
2 AB' + a = 2 AB' + A'B + A'C = 2 AB' + BC' + CB' = b + c.
\]
\item On note \(p = \frac{a + b + c}{2}\). Montrer que \(AB' = p - a\).
\item On note \(r = \Omega A'\). Montrer que
\[
r = (p - a) \tan\frac{\alpha}{2}
 = (p - b) \tan\frac{\beta}{2}
 = (p - c) \tan\frac{\gamma}{2}.
\]
\item En utilisant les formules d'addition pour \(\sin\) et \(\cos\), montrer que
\[
\tan\left(\frac{\alpha}{2} + \frac{\beta}{2}\right)
= \frac{\tan\frac{\alpha}{2} + \tan\frac{\beta}{2}}{1 - \tan\frac{\alpha}{2} \tan\frac{\beta}{2}}
\quad\text{et}\quad
\tan\left(\frac{\pi}{2} - \frac{\gamma}{2}\right)
= \frac{1}{\tan\frac{\gamma}{2}}.
\]
\item Montrer que
\(\frac{\alpha}{2} + \frac{\beta}{2} + \frac{\gamma}{2} = \frac{\pi}{2}\).
En déduire que
\[
\frac{1}{\tan\frac{\alpha}{2}} + \frac{1}{\tan\frac{\beta}{2}} + \frac{1}{\tan\frac{\gamma}{2}}
= \frac{1}{\tan\frac{\alpha}{2} \tan\frac{\beta}{2} \tan\frac{\gamma}{2}}.
\]
\item En déduire que \((p - a) (p - b) (p - c) = r^2 p\), puis que l'aire 
du triangle \(ABC\) est
\[
\sqrt{p (p - a) (p - b) (p - c)}.
\]
\end{enumerate}
\finenonce{007267}
\finexercice


\exercice{7268, mourougane, 2021/08/10}
\enonce{007268}{}
Soit $\mathcal{T}$ un triangle rectangle dont les côtés adjacents à l'angle droit mesurent $a$ et $b$,
et l'hypoténuse $c$. On construit dans un carré $ABCD$ de côté $a+b$ les points $A'B'C'D'$
sur les côtés aux distances indiquées des extrémités.

\begin{center}
\setlength{\unitlength}{1657sp}%
%
\begingroup\makeatletter\ifx\SetFigFont\undefined%
\gdef\SetFigFont#1#2#3#4#5{%
  \reset@font\fontsize{#1}{#2pt}%
  \fontfamily{#3}\fontseries{#4}\fontshape{#5}%
  \selectfont}%
\fi\endgroup%
\begin{picture}(7095,7050)(1021,-7186)
\thinlines
{\color[rgb]{0,0,0}\put(1351,-6811){\framebox(6300,6300){}}
}%
{\color[rgb]{0,0,0}\put(1351,-3211){\line( 4, 3){3600}}
}%
{\color[rgb]{0,0,0}\put(1351,-3211){\line( 3,-4){2700}}
}%
{\color[rgb]{0,0,0}\put(4051,-6811){\line( 4, 3){3600}}
}%
{\color[rgb]{0,0,0}\put(4951,-511){\line( 3,-4){2700}}
}%
\put(1081,-1591){\makebox(0,0)[lb]{\smash{{\SetFigFont{5}{6.0}{\rmdefault}{\mddefault}{\updefault}{\color[rgb]{0,0,0}b}%
}}}}
\put(6346,-286){\makebox(0,0)[lb]{\smash{{\SetFigFont{5}{6.0}{\rmdefault}{\mddefault}{\updefault}{\color[rgb]{0,0,0}b}%
}}}}
\put(7921,-5866){\makebox(0,0)[lb]{\smash{{\SetFigFont{5}{6.0}{\rmdefault}{\mddefault}{\updefault}{\color[rgb]{0,0,0}b}%
}}}}
\put(2431,-7126){\makebox(0,0)[lb]{\smash{{\SetFigFont{5}{6.0}{\rmdefault}{\mddefault}{\updefault}{\color[rgb]{0,0,0}b}%
}}}}
\put(2926,-331){\makebox(0,0)[lb]{\smash{{\SetFigFont{5}{6.0}{\rmdefault}{\mddefault}{\updefault}{\color[rgb]{0,0,0}a}%
}}}}
\put(7966,-1996){\makebox(0,0)[lb]{\smash{{\SetFigFont{5}{6.0}{\rmdefault}{\mddefault}{\updefault}{\color[rgb]{0,0,0}a}%
}}}}
\put(6391,-7171){\makebox(0,0)[lb]{\smash{{\SetFigFont{5}{6.0}{\rmdefault}{\mddefault}{\updefault}{\color[rgb]{0,0,0}a}%
}}}}
\put(1036,-5056){\makebox(0,0)[lb]{\smash{{\SetFigFont{5}{6.0}{\rmdefault}{\mddefault}{\updefault}{\color[rgb]{0,0,0}a}%
}}}}
\put(3331,-2086){\makebox(0,0)[lb]{\smash{{\SetFigFont{5}{6.0}{\rmdefault}{\mddefault}{\updefault}{\color[rgb]{0,0,0}c}%
}}}}
\put(6256,-2626){\makebox(0,0)[lb]{\smash{{\SetFigFont{5}{6.0}{\rmdefault}{\mddefault}{\updefault}{\color[rgb]{0,0,0}c}%
}}}}
\put(5671,-5326){\makebox(0,0)[lb]{\smash{{\SetFigFont{5}{6.0}{\rmdefault}{\mddefault}{\updefault}{\color[rgb]{0,0,0}c}%
}}}}
\put(2926,-4831){\makebox(0,0)[lb]{\smash{{\SetFigFont{5}{6.0}{\rmdefault}{\mddefault}{\updefault}{\color[rgb]{0,0,0}c}%
}}}}
\put(1216,-286){\makebox(0,0)[lb]{\smash{{\SetFigFont{5}{6.0}{\rmdefault}{\mddefault}{\updefault}{\color[rgb]{0,0,0}A}%
}}}}
\put(7876,-286){\makebox(0,0)[lb]{\smash{{\SetFigFont{5}{6.0}{\rmdefault}{\mddefault}{\updefault}{\color[rgb]{0,0,0}B}%
}}}}
\put(7741,-7081){\makebox(0,0)[lb]{\smash{{\SetFigFont{5}{6.0}{\rmdefault}{\mddefault}{\updefault}{\color[rgb]{0,0,0}C}%
}}}}
\put(1171,-7081){\makebox(0,0)[lb]{\smash{{\SetFigFont{5}{6.0}{\rmdefault}{\mddefault}{\updefault}{\color[rgb]{0,0,0}D}%
}}}}
\put(1081,-3211){\makebox(0,0)[lb]{\smash{{\SetFigFont{5}{6.0}{\rmdefault}{\mddefault}{\updefault}{\color[rgb]{0,0,0}A'}%
}}}}
\put(4906,-331){\makebox(0,0)[lb]{\smash{{\SetFigFont{5}{6.0}{\rmdefault}{\mddefault}{\updefault}{\color[rgb]{0,0,0}B'}%
}}}}
\put(8101,-4201){\makebox(0,0)[lb]{\smash{{\SetFigFont{5}{6.0}{\rmdefault}{\mddefault}{\updefault}{\color[rgb]{0,0,0}C'}%
}}}}
\put(4096,-7081){\makebox(0,0)[lb]{\smash{{\SetFigFont{5}{6.0}{\rmdefault}{\mddefault}{\updefault}{\color[rgb]{0,0,0}D'}%
}}}}
\end{picture}%
\end{center}

Montrer que $A'B'C'D'$ est un carré.
En déduire une démonstration du théorème de Pythagore.
\finenonce{007268}
\finexercice


\exercice{7269, mourougane, 2021/08/10}
\enonce{007269}{}
Soit $ABC$ un triangle équilatéral de côté $a$. Soit $B'$ le milieu de $[AC]$.
 Soit $L$ le point de $[BB']$ tel que $B'L=a/2$. Soit $d$ la droite orthogonale à $(BB')$ qui passe par $L$.
 Soit $M$ un point d'intersection de $d$ avec le cercle de diamètre $[BB']$.
 Calculer $BM$.
\finenonce{007269}
\finexercice


\exercice{7270, mourougane, 2021/08/10}
\enonce{007270}{}
Le \emph{théorème de Bolyai} affirme que deux polygones de même aire peuvent toujours être obtenus 
l'un à partir de l'autre par découpage et recollement.
On va étudier le découpage d'un triangle équilatéral pour obtenir un carré de même aire.


\textit{On considère un triangle équilatéral $ABC$ que l'on souhaite découper de façon à pouvoir le transformer en un rectangle, voire un carré.}

\includegraphics[scale=0.6]{images/img-mour-022-1}

\begin{enumerate}
\item On part du triangle équilatéral $ABC$ et on place un point $D$ sur $[BC]$ tel que $0<BD<\frac{1}{2}BC$. On construit ensuite $E$ sur le même segment tel que $DE=\frac{1}{2}BC$. 
On joint $D$ au milieu $B'$ de $[AC]$, et on appelle $H$ et $K$ les projetés orthogonaux de $C'$ et $E$ sur $[B'D]$.\\
On suppose le découpage effectué selon le modèle des figures a et b donne un rectangle sur la figure $b$. 
Refaire des figures en choisissant $D$ proche de $B$.
Montrer en utilisant le fait que la figure b est un rectangle avec coïncidence de certains points que :
\begin{enumerate}
\item $B'$ et $C'$ sont les milieux de $[AB]$ et $[AC]$,
\item $DE=\frac{1}{2}BC$,
\item $KE=HC'$,
\item $DH=KB'$.
\end{enumerate}
\item Montrer que $DEB'C'$ est un parallélogramme. 
\item En déduire que les triangles $BC'D$ et $CB'E$ ont deux côtés et un angle égaux. 
Peut-on en déduire qu'ils sont isométriques ?
% Montrer que dans ce cas, $BD=EC$, et donc les points $D$ et $E$ sont au quart et trois quart de $[BC]$, et les droites $(C'D)$ et $(B'E)$ sont perpendiculaires à $(BC)$. \\ \\

%\textit{En fait, les triangles $BC'D$ et $CB'E$ ne sont pas isométriques (on n'est pas dans le premier cas d'isométrie puisque l'angle en commun n'est pas celui compris entre les deux côtés de même longueur).}

\item On part du triangle équilatéral $ABC$ et on place un point $D$ sur $[BC]$ tel que $0<BD<\frac{1}{2}BC$. On construit ensuite $E$ sur le même segment tel que $DE=\frac{1}{2}BC$. 
On joint $D$ au milieu $B'$ de $[AC]$, et on appelle $H$ et $K$ les projetés orthogonaux de $C'$ et $E$ sur $[B'D]$.\\
Montrer que l'on obtient bien un rectangle en découpant la figure comme indiqué sur les figures c et d .

\includegraphics[scale=0.6]{images/img-mour-022-2}

\item Calculer, en fonction du côté $a$ du triangle, les longueurs des côtés du rectangle obtenu dans le cas où les points $D$ et $E$ sont au quart et trois quart du segment $[BC]$. Ce rectangle est-il un carré ?
\item Déterminer la position de $D$ pour laquelle on obtient effectivement un carré. On précisera la longueur $B'D$ à l'aide d'un calcul d'aire.
\end{enumerate}
\finenonce{007270}

\finexercice


\exercice{7271, mourougane, 2021/08/10}
\enonce{007271}{}
 \begin{enumerate}
  \item Une translation transforme-t-elle une droite en une droite qui lui est parallèle ?
  \item Montrer qu'une translation conserve le parallélisme, c'est à dire qu'elle transforme deux droites parallèles $d_1 // d_2$
  en deux droites parallèles $d'_1 // d'_2$.
  \item Reprendre les questions pour une rotation.
 \end{enumerate}
\finenonce{007271}
\finexercice


\exercice{7272, mourougane, 2021/08/10}
\enonce{007272}{}
Soit $E$ un plan euclidien et $A$ et $B$ deux points de $E$
 \begin{enumerate}
\item Déterminer la composée de l'homothétie de centre $A$ et de rapport $2$
avec l'homothétie de centre $B$ et de rapport $1/2$.
On pourra construire l'image de quelques points.
\item Déterminer la composée de l'homothétie de centre $A$ et de rapport $2$
avec l'homothétie de centre $B$ et de rapport $3$.
 \item Que dire en général de la composé de deux homothéties ?
 \item Que dire de la composée d'une homothétie et d'une translation ?
\end{enumerate}
\finenonce{007272}
\finexercice


\exercice{7273, mourougane, 2021/08/10}
\enonce{007273}{}
 \begin{enumerate}
\item Soit $d$ une droite d'un plan euclidien orienté $E$ de vecteur directeur $\vec{u}$.
 Soit $\vec{v}$ un vecteur orthogonal à $\vec{u}$.
 Déterminer la composée de la symétrie axiale $s_d$ d'axe $d$ avec la translation $t$ de vecteur $\vec{u}$.
 \textit{Indication : on pourra décomposer la translation comme composée de deux symétries bien choisies.}
 \item Soit $d$ une droite d'un plan euclidien orienté $E$ de vecteur directeur $\vec{u}$.
 Soit $A$ un point de $d$.
 Déterminer la composée de la symétrie axiale $s_d$ d'axe $d$ avec la rotation $r$ de centre $A$
 et d'angle $\pi/4$.
 \end{enumerate}
\finenonce{007273}
\finexercice


\exercice{7274, mourougane, 2021/08/10}
\enonce{007274}{}
On rappelle que l'\emph{homothétie} de centre \(O\) et de 
rapport \(\lambda\), pour \(O\) un point du plan \(\mathcal{P}\) et 
\(\lambda \in \Rr \smallsetminus \{0\}\), est l'application 
de \(\mathcal{P}\) dans lui-même qui à un point \(P \in \mathcal{P}\) 
associe l'unique point \(Q\) vérifiant
\(\overrightarrow{OQ} = \lambda \overrightarrow{OP}\).

Dans cet exercice, on considère une homothétie \(h\) de centre \(O\) 
et de rapport \(\lambda\).

\begin{enumerate}
\item Montrer que si \(\lambda \neq 1\), alors \(O\) est l'unique 
point du plan dont l'image par \(h\) est lui-même.
\item La propriété précédente est-elle vraie si \(\lambda = 1\)? 
(Rappel: votre réponse doit être accompagnée d'une démonstration.)
\item Si \(P\) est un point du plan autre que \(O\), montrer que 
\(h(P)\) est un point de la droite \((OP)\) (\emph{indication :} on pourra 
considérer les vecteurs \(\overrightarrow{OP}\) et 
\(\overrightarrow{O \; h(P)}\)).
\item En déduire que si \(\mathcal{D}\) est une droite passant par 
le point \(O\), alors \(h(\mathcal{D}) \subseteq \mathcal{D}\). 
(Rappel: on note \(h(\mathcal{D})\) l'ensemble des points \(h(P)\) 
pour \(P \in \mathcal{D}\).)
\item En considérant l'homothétie de centre \(O\) et de rapport 
\(\frac{1}{\lambda}\), en déduire que \(h(\mathcal{D}) = \mathcal{D}\) 
(avec les mêmes notations qu'à la question précédente).
\item Soint \(\Delta\) une droite. On admet dans cet exercice que 
\(h(\Delta)\) est aussi une droite. Montrer que si \(\Delta\) et 
\(h(\Delta)\) ont un point commun et si \(\lambda \neq 1\) alors 
ce point commun est \(O\).
\item En déduire que si \(\Delta\) est une droite ne passant pas 
par \(O\) et si \(\lambda \neq 1\) alors les droites \(\Delta\) et 
\(h(\Delta)\) sont parallèles et sont deux droites distinctes.
\item Montrer que l'image d'une droite par une homothétie est une 
droite parallèle. (Attention, dans cette question, il n'y a plus 
d'hypothèse particulière sur la droite considérée, ni sur le rapport 
de l'homothétie.)
\end{enumerate}
\finenonce{007274}
\finexercice


\exercice{7275, mourougane, 2021/08/10}
\enonce{007275}{}
\label{centre-regle-compas}
\begin{enumerate}
\item Étant donnés deux points \(A\) et \(B\) du plan, construire 
à la règle et au compas la médiatrice du segment \([AB]\). 
(Indication: on pourra, à l'aide du compas, construire deux points 
équidistants de \(A\) et \(B\).)
\item Étant donné un cercle \(\mathcal{C}\) dont on ne connaît pas 
le centre, déterminer son centre par une construction à la règle et 
au compas. (\emph{Indication :} on pourra utiliser la question précédente.)
\end{enumerate}
\finenonce{007275}
\finexercice


\exercice{7276, mourougane, 2021/08/10}
\enonce{007276}{}
Étant données deux demi-droites \([OA)\) et \([OB)\), on veut construire 
la bissectrice de l'angle géométrique \(\widehat{AOB}\).
\begin{enumerate}
\item Comment construire au compas deux points \(P \in ]OA)\) et \(Q \in ]OB)\) 
de sorte que le triangle \(OPQ\) soit isocèle?
\item Montrer que la médiatrice du segment \([PQ]\) est aussi la bissectrice 
de l'angle \(\widehat{POQ}\). (Indication: utiliser une symétrie axiale bien 
choisie.)
\item Expliquer comment construire à la règle et au compas la bissectrice de 
l'angle \(\widehat{AOB}\).
\end{enumerate}
\finenonce{007276}
\finexercice


\exercice{7277, mourougane, 2021/08/10}
\enonce{007277}{}
\label{sym-compas-seul}
Comment construire, au compas seul, le symétrique d'un point par rapport 
à une droite? Comment construire à la règle et au compas le projeté orthogonal 
d'un point sur une droite?
\finenonce{007277}
\finexercice


\exercice{7278, mourougane, 2021/08/10}
\enonce{007278}{}
\label{para-compas-seul}
Étant donnés trois points \(A\), \(B\), \(C\) non alignés, construire au compas 
seul le point \(D\) tel que \(ABCD\) soit un parallélogramme.
\finenonce{007278}
\finexercice


\exercice{7279, mourougane, 2021/08/10}
\enonce{007279}{}
Étant données trois droites concourantes, construire un triangle dont ce 
sont les médiatrices. (Indication: on pourra considérer l'application 
composée des trois symétries axiales correspondantes, et ses points fixes.)
\finenonce{007279}
\finexercice


\exercice{7280, mourougane, 2021/08/10}
\enonce{007280}{Construction du pentagone régulier}
\begin{enumerate}
\item Montrer que pour tout \( \theta \in \Rr\), on~a
\[
\cos(5 \theta) = 16 \cos( \theta)^5 - 20 \cos( \theta)^3 + 5 \cos( \theta).
\]
\item En déduire que
$$
\cos\frac{\pi}{5} = \frac{1 + \sqrt{5}}{4} \quad \text{ et } \quad 
\cos\frac{3 \pi}{5} = \frac{1 - \sqrt{5}}{4} $$

$$\sin\frac{ \pi}{10} = \frac{-1 + \sqrt{5}}{4} \quad \text{ et } \quad 
\sin\frac{3 \pi}{10} = \frac{1 + \sqrt{5}}{4}.
$$

On vérifiera que
\[
16 X^5 - 20 X^3 + 5 X + 1 = (X + 1) (4 X^2 - 2 X - 1)^2.
\]
\item Quelle est la longueur des côtés d'un pentagone régulier inscrit 
dans un cercle de rayon \(1\)? celle d'un décagone régulier?
\item On considère la construction suivante:

\begin{center}
\includegraphics[scale=1]{images/img-mour-032}
\end{center}

% [[figure asymptote]]

%\begin{center}
%\medskip
%\begin{asy}
%size(7cm, 0);
%point O = (0, 0); label("\(O\)", O, NE);
%circle C0 = circle(O, 1); draw(C0);
%point A = (-1, 0); label("\(A\)", A, W);
%point B = (1, 0); label("\(B\)", B, E);
%draw(A--B);
%circle C1 = circle(A, O); draw(C1);
%point C = (0, -1); label("\(C\)", C, S);
%point D = (0, 1); label("\(D\)", D, N);
%draw(C--D);
%perpendicularmark(line(O, A), line(O, C));
%line D0 = line(0.5*A + 0.5*O, C);
%draw(D0);
%point[] PQ = intersectionpoints(C1, D0);
%point P = PQ[1]; label("\(P\)", P, S);
%point Q = PQ[0]; label("\(Q\)", Q, S);
%circle C2 = circle(C, abs(P-C));
%clipdraw(C2);
%circle C3 = circle(C, abs(Q-C));
%clipdraw(C3);
%point[] R = intersectionpoints(C0, C2);
%point[] T = intersectionpoints(C0, C3);
%label("\(R_0\)", R[0], W);
%label("\(R_1\)", R[1], E);
%label("\(S_0\)", T[0], W);
%label("\(S_1\)", T[1], E);
%draw(D -- R[0] -- T[0] -- T[1] -- R[1] -- cycle, gray);
%\end{asy}
%\end{center}


Les segments \([AB]\) et \([CD]\) sont deux diamètres orthogonaux d'un 
cercle \(\mathcal{C}\) de centre \(O\). Notons \(\mathcal{C}_1\) le cercle 
de diamètre \([AO]\), et \(P\) et \(Q\) les points d'intersection du 
cercle \(\mathcal{C}_1\) et de la droite passant par \(C\) et par le 
centre de \(\mathcal{C}_1\). Soient \(R_0\) et \(R_1\) les points 
d'intersection de \(\mathcal{C}\) et du cercle de centre \(C\) passant 
par \(P\), et soient \(S_0\) et \(S_1\) les points d'intersection 
de \(\mathcal{C}\) et du cercle de centre \(C\) passant par \(Q\). 
Calculer les longueurs \(CP\) et \(CQ\).
\item En déduire la mesure des angles \(\widehat{COR_0}\) 
et \(\widehat{COS_0}\).
\item En déduire que les points \(D R_0 S_0 S_1 R_1\) forment un pentagone 
régulier.
\end{enumerate}
\finenonce{007280}



\finexercice


\exercice{7281, mourougane, 2021/08/10}
\enonce{007281}{}
Étant donnés quatre points \(A\), \(B\), \(C\), \(D\), construire à 
la règle et au compas deux points dont la distance soit \(AB + CD\). 
Même question pour \(AB - CD\).
\finenonce{007281}
\finexercice


\exercice{7282, mourougane, 2021/08/10}
\enonce{007282}{}
Étant données trois longueurs \(a\), \(b\) et \(c\), construire à la 
règle et au compas la longueur \(\frac{ab}{c}\). (Indication: on pourra 
considérer la figure ci-dessous.)


\begin{center}
    \includegraphics[scale=1]{images/img-mour-034}
\end{center}

% [[figure asymptote]]

%\begin{center}
%\begin{asy}
%size(7cm, 0);
%point A = (0, 0); label("\(A\)", A, S);
%point B = (1, 0); label("\(B\)", B, S);
%point C = (1.7, 0); label("\(C\)", C, S);
%line L0 = line(A, B); draw(L0);
%distance("\(b\)", A, C, 5mm, gray);
%distance("\(c\)", A, B, 1cm, gray);
%point D = (0.3, 0.8); label("\(D\)", D, E);
%distance("\(a\)", A, D, -5mm, gray);
%line L1 = line(A, D); draw(L1);
%line L2 = line(B, D); draw(L2);
%line L3 = parallel(C, L2); draw(L3);
%point E0 = intersectionpoint(L1, L3); label("\(E\)", E0, E);
%distance("?", A, E0, -1cm, gray);
%\end{asy}
%\end{center}

\finenonce{007282}
\finexercice


\exercice{7283, mourougane, 2021/08/10}
\enonce{007283}{}
Étant données deux longueurs \(a\) et \(b\), construire à la 
règle et au compas la longueur \(\sqrt{ab}\). (Indication: on pourra 
considérer la figure ci-dessous.)

Pourquoi, dans cet exercice et dans le précédent, a-t-on considéré 
les longueurs \(\frac{ab}{c}\) et \(\sqrt{ab}\) plutôt que \(ab\), 
\(\frac{a}{c}\) et \(\sqrt{a}\)?

\begin{center}
    \includegraphics[scale=1]{images/img-mour-035}
\end{center}

% [[figure asymptote]]

%\begin{center}
%\begin{asy}
%size(7cm, 0);
%point A = (0, 0); label("\(A\)", A, S);
%point B = (1, 0); label("\(B\)", B, S);
%point C = (3, 0); label("\(C\)", C, S);
%line L0 = line(A, C);
%draw(A -- C);
%distance("\(a\)", A, B, 8mm, gray);
%distance("\(b\)", B, C, 8mm, gray);
%circle C0 = circle(A, C);
%line L1 = perpendicular(B, L0);
%point[] D = intersectionpoints(C0, L1);
%label("\(D\)", D[1], NW);
%draw(L1);
%perpendicularmark(L0, L1);
%distance("?", B, D[1], 5mm, gray);
%clipdraw(C0);
%\end{asy}
%\end{center}

\finenonce{007283}
\finexercice


\exercice{7284, mourougane, 2021/08/10}
\enonce{007284}{Puissance d'un point par rapport à un cercle}
Soit \(\mathcal{C}\) un cercle de centre \(O\) et de rayon \(r\), et 
soit \(P\) un point. Soit \(\mathcal{D}\) une droite passant par \(P\) et 
qui coupe \(\mathcal{C}\) en deux points \(A\) et \(B\). La \emph{puissance} 
du point \(P\) par rapport au cercle \(\mathcal{C}\) est le produit scalaire 
\(\vec{PA} \cdot \vec{PB}\).

\begin{center}
    \includegraphics[scale=1]{images/img-mour-036}
\end{center}

% [[figure asymptote]]

%\begin{center}
%\begin{asy}
%size(7cm, 0);
%point O = (0, 0); dot("\(O\)", O, S);
%circle C0 = circle(O, 1.3); draw(C0);
%point P = (-3, 0); dot("\(P\)", P, S);
%line L0 = line(P, 0); draw(L0);
%line D0 = rotate(20, P) * L0; draw(D0);
%point[] AB = intersectionpoints(C0, D0);
%label("\(A\)", AB[0], N); label("\(B\)", AB[1], N);
%point Q = projection(D0) * O;
%draw(O--Q, gray); perpendicularmark(line(O, Q), D0, gray, quarter=3);
%draw(O--AB[0], gray); draw(O--AB[1], gray);
%\end{asy}
%\end{center}


Montrer que la puissance de \(P\) par rapport à \(\mathcal{C}\) est 
égale à \(OP^2 - r^2\). (Indication: utiliser le théorème de Pythagore.) 
Dépend-elle de la droite \(\mathcal{D}\) choisie?
\finenonce{007284}
\finexercice


\exercice{7285, mourougane, 2021/08/10}
\enonce{007285}{Inversions}
\begin{enumerate}
\item Si \( \Omega\) est un point et si \( \lambda\) est un réel non nul, 
l'\emph{inversion} de centre \( \Omega\) et de rapport \( \lambda\) est l'application 
du plan privé de \(P\) dans lui-même qui au point \(P\) associe l'unique 
point \(P' \in ( \Omega P)\) tel que \(\vec{ \Omega P} \cdot \vec{ \Omega P'} = \lambda\). Montrer que 
les inversions sont involutives, c'est-à-dire que ce sont des bijections 
qui sont leur propre bijection réciproque.
\item Montrer que l'ensemble des points fixes d'une inversion est un 
cercle.
\item\label{q-inv-puissance} Montrer que si \(\mathcal{C}\) est un cercle 
tel que la puissance de \( \Omega\) par rapport à \(\mathcal{C}\) est \( \lambda\), 
et si \( \iota\)~est l'inversion de centre \( \Omega\) et de rapport \( \lambda\), alors 
\( \iota(\mathcal{C}) = \mathcal{C}\).
\item Si \( \Delta\) est une droite qui ne passe pas par \( \Omega\), montrer que son 
image par l'inversion de centre \( \Omega\) et de rapport \( \lambda\) est un cercle 
passant par \( \Omega\) privé du point \( \Omega\). (Indication: noter \(H\) le 
projeté orthogonal de \( \Omega\) sur \( \Delta\), \( \iota\)~l'inversion; en appliquant 
une homothétie de centre \( \Omega\) bien choisie on peut supposer que 
\( \iota(H) = H\); montrer alors que \(P \in \Delta\) si et seulement si 
\(\vec{ \Omega \iota(P)} \cdot \vec{H \iota(P)} = 0\).)
\item Montrer que l'image par une inversion d'un cercle ne passant pas 
par le centre de l'inversion est un cercle. (Indication: utiliser une 
homothétie de même centre que l'inversion pour se ramener à la 
question~\ref{q-inv-puissance}.)
\end{enumerate}
\finenonce{007285}
\finexercice 


\exercice{7286, mourougane, 2021/08/10}
\enonce{007286}{Intersection d'un cercle et d'une droite}
On considère un cercle \(\mathcal{C}\) de centre \(O\) et deux points 
distincts \(A\) et \(B\). On veut construire au compas seul les points 
d'intersection du cercle \(\mathcal{C}\) et de la droite \((AB)\).
\begin{enumerate}
\item Si \(O \notin (AB)\), expliquer comment construire ces deux points. 
(\emph{Indication :} on pourra construire le symétrique du cercle \(\mathcal{C}\) 
par rapport à la droite \((AB)\), en utilisant 
l'exercice~\ref{sym-compas-seul}.)
\item Si \(O \in (AB)\), on se ramène à une intersection de cercles en 
utilisant une inversion. Plus précisément, on considère la construction 
suivante.

\begin{center}
    \includegraphics[scale=1]{images/img-mour-038}
\end{center}

% [[figure asymptote]]

%\begin{center}
%\begin{asy}
%size(7cm, 0);
%point O = (0, 0); dot("\(O\)", O, S);
%point A = (2, 0); dot("\(A\)", A, S);
%point B = (3, 0); dot("\(B\)", B, S);
%circle C0 = circle(O, 1.3); draw(C0);
%draw(line(A, B));
%point O1 = (1.5, -0.5);
%circle C1 = circle(O1, abs(A - O1)); draw(C1);
%point[] DE = intersectionpoints(C0, C1);
%dot("\(D\)", DE[0], NE); dot("\(E\)", DE[1], NE);
%line L0 = line(DE[0], DE[1]); draw(L0);
%point F = reflect(L0) * O; dot("\(F\)", F, S);
%point G = rotate(60, F) * O; dot("\(G\)", G, E);
%line L1 = line(A, G); draw(L1);
%point[] AH = intersectionpoints(C1, L1);
%point H = (AH[0] == A) ? AH[1] : AH[0]; dot("\(H\)", H, S);
%point I = intersectionpoint(bisector(segment(G, H)),
%      perpendicular(G, line(O, A)));
%dot("\(I\)", I, S);
%circle C2 = circle(I, abs(G - I)); draw(C2);
%point[] M = intersectionpoints(C0, C2);
%dot("\(M\)", M[0], W); dot("\(M'\)", M[1], W);
%draw(line(G, M[0])); draw(line(G, M[1]));
%\end{asy}
%\end{center}

Soit \(\mathcal{C}'\) un cercle passant par \(A\), qui coupe \(\mathcal{C}\) 
en deux points \(D\) et \(E\). Soit \(F\) le symétrique de \(O\) par rapport 
à la droite \((DE)\), et soit \(G\) un point tel que le triangle \(OFG\) soit 
équilatéral. Montrer que \(G \in (DE)\) et expliquer comment le construire au 
compas seul.
\item Montrer que \(G\) a la même puissance par rapport à \(\mathcal{C}\) 
et à \(\mathcal{C}'\). En déduire qu'il existe une inversion \( \iota\) de 
centre \(G\) telle que \( \iota(\mathcal{C}) = \mathcal{C}\) et 
\( \iota(\mathcal{C}') = \mathcal{C}'\).
\item On note \(\mathcal{C}''\) l'image de la droite \((OA)\) par 
l'inversion \( \iota\) (en ajoutant son centre \(G\) pour avoir un cercle), 
et on note \(I\) le centre de \(\mathcal{C}''\). Montrer que les droites 
\((GI)\) et \((OA)\) sont orthogonales.
\item Soit \(H\) l'intersection du cercle \(\mathcal{C}'\) et de la 
droite \((AG)\) qui n'est pas \(A\). Montrer que \( \iota(A) = H\).
\item Montrer que \(I\)~est sur la médiatrice du segment \([GH]\). Comment 
construire au compas seul deux points de cette médiatrice?
\item Comment construire deux points de la perpendiculaire à \((OA)\) 
passant par \(G\)? (\emph{Indication :} construire le symétrique de \(G\) par rapport 
à \((OA)\).)
\item Montrer que \(I\) est l'intersection de la médiatrice de \([GH]\) et 
de la perpendiculaire à \((OA)\) passant par \(G\) (voir 
l'exercice~\ref{inter-droites-compas-seul} pour comment construire cette 
intersection au compas seul).
\item En déduire comment construire \(\mathcal{C}''\). (\emph{Indication :} 
\(G \in \mathcal{C}''\).)
\item On note \(M\) et \(M'\) les points d'intersection de \(\mathcal{C}\) 
et \(\mathcal{C}'\). Montrer que \( \iota(M)\) et \( \iota(M')\) sont les points 
d'intersection de \(\mathcal{C}\) et \((AB)\).
\item Montrer que \( \iota(M)\) est l'intersection de \(\mathcal{C}\) et de 
la droite \((GM)\), autre que \(M\).
\end{enumerate}
\finenonce{007286}
\finexercice 


\exercice{7287, mourougane, 2021/08/10}
\enonce{007287}{Intersection de deux droites}
\label{inter-droites-compas-seul}
On considère quatre points distincts \(A\), \(B\), \(C\), \(D\), 
tels que les droites \((AB)\) et \((CD)\) soient sécantes en un point \(X\). 
On veut construire le point \(X\) au compas seul.

\begin{center}
    \includegraphics[scale=1]{images/img-mour-039}
\end{center}

% [[figure asymptote]]

%\begin{center}
%\begin{asy}
%size(7cm, 0);
%point A = (0, 0); dot("\(A\)", A, S);
%point B = (1, 0); dot("\(B\)", B, S);
%line L0 = line(A, B); draw(L0);
%point C = (1.5, 1.5); dot("\(C\)", C, N);
%point D = (0.8, 3); dot("\(D\)", D, N);
%line L1 = line(C, D); draw(L1);
%point X = intersectionpoint(L0, L1); dot("\(X\)", X, S);
%line L2 = parallel(D, L0); draw(L2);
%real r = 3;
%point O = intersectionpoints(circle(C, r), circle(D, r))[1];
%dot("\(O\)", O, N);
%circle C0 = circle(O, r); clipdraw(C0);
%point[] DF = intersectionpoints(C0, L2);
%point F = (DF[0] == D) ? DF[1] : DF[0]; dot("\(F\)", F, N);
%circle C1 = circle(C, abs(F - C)); clipdraw(C1);
%point[] H = intersectionpoints(C1, L0);
%dot("\(H\)", H[0], S); dot("\(H'\)", H[1], S);
%draw(C--F, gray); draw(C--H[0], gray); draw(C--H[1], gray);
%point K0 = intersectionpoints(circle(C, r), circle(H[0], r))[0];
%point K1 = intersectionpoints(circle(C, r), circle(H[1], r))[1];
%dot("\(K\)", K0, S); dot("\(K'\)", K1, N);
%circle C2 = circle(K0, r); draw(C2);
%circle C3 = circle(K1, r); draw(C3);
%\end{asy}
%\end{center}

\begin{enumerate}
\item À l'aide de l'exercice~\ref{para-compas-seul}, expliquer comment 
construire la parallèle \( \Delta\) à \((AB)\) passant par \(D\).
\item Construire un point \(O\) équidistant de \(C\) et \(D\), et \(F\) 
le point d'intersection du cercle de centre \(O\) passant par \(C\) et \(D\) 
avec ma droite \( \Delta\), autre que \(D\). On note \(H\) et \(H'\) les points 
d'intersection du cercle de centre \(C\) passant par \(F\) avec la 
droite \((AB)\). Comment construire ces points au compas seul?
\item Construire le centre \(K\) de \(\mathcal{C}\), un des deux cercles 
de rayon \(OC\) passant par \(C\) et \(H\). De même, construire le centre \(K'\) 
de \(\mathcal{C}'\), un des deux cercles de rayon \(OC\) passant par \(C\) 
et \(H'\). Montrer que le point \(C\) est commun à \(\mathcal{C}\) 
et \(\mathcal{C}'\).
\item Notons \( \alpha\) la mesure de l'angle \(\widehat{AXD}\). En supposant que 
les points sont disposés comme sur la figure ci-dessus, montrer que:
\begin{enumerate}
\item la mesure de \(\widehat{CDF}\) est \( \alpha\);
\item la mesure de \(\widehat{COF}\) est \(2 \alpha\);
\item la mesure de \(\widehat{CKH}\) est \(2 \alpha\);
\item la mesure de \(\widehat{CK'H'}\) est \(2 \alpha\).
\end{enumerate}
\item Notons \( \zeta \) le point d'intersection de \(\mathcal{C}\) et \((AB)\) 
autre que \(C\). Montrer que la mesure de l'angle \(\widehat{C \zeta H}\) 
est \( \alpha\). En déduire que \( \zeta = X\).
\item Montrer que les points d'intersection de \(\mathcal{C}\) 
et \(\mathcal{C}'\) sont \(C\) et \(X\).
\end{enumerate}
\finenonce{007287}
\finexercice 


\exercice{7288, mourougane, 2021/08/10}
\enonce{007288}{}
La notion de «construction au compas» est un peu ambiguë. On distingue:
\begin{itemize}
\item le \emph{compas traçant}, qui permet, à partir de deux points 
\(A\) et \(B\), de construire le cercle de centre \(A\) passant 
par \(B\);
\item le \emph{compas transporteur}, qui permet, à partir de trois 
points \(A\), \(B\) et \(C\), de construire le cercle de centre \(A\) 
et de rayon \(BC\) (il permet de «transporter» la distance \(BC\), 
d'où son nom).
\end{itemize}
Le compas transporteur permet évidemment toutes les constructions 
possibles au compas traçant. Le but de cet exercice est de montrer 
que toute construction au compas transporteur peut être transformée 
en une construction au compas traçant.

On considère trois points \(A\), \(B\) et \(C\), deux à deux 
distincts.
\begin{enumerate}
\item Construire, au compas traçant, deux points de la médiatrice 
du segment \([AB]\).
\item Construire, au compas traçant, le symétrique \(C'\) de \(C\) 
par rapport à la médiatrice du segment \([AB]\).
\item Montrer que \(AC' = BC\).
\item Construire, au compas traçant, le cercle de centre \(A\) et 
de rayon \(BC\).
\end{enumerate}


\emph{Remarques.} La question de savoir quelles constructions peuvent être faites au compas 
seul (sans règle) a été étudiée dans l'espoir d'obtenir des constructions 
plus précises. (Il est plus facile de fabriquer un bon compas qu'une règle 
vraiment droite.)

D'après le théorème de Mohr et Mascheroni, les constructions réalisables 
au compas seul sont exactement les mêmes que celles réalisables à la règle 
et au compas, à condition de considérer qu'une droite est construite à partir 
du moment où on en connaît deux points. On peut trouver plus de 
renseignements sur ce théorème dans le livre \emph{Leçons sur les 
constructions géométriques} d'Henri Lebesgue.

\finenonce{007288}
\finexercice


\exercice{7289, mourougane, 2021/08/10}
\enonce{007289}{}
% cf. https://en.wikipedia.org/wiki/Napoleon%27s_problem
La question de trouver le centre d'un cercle à la règle et au compas a été 
étudiée à l'exercice~\ref{centre-regle-compas}. On étudie maintenant une 
construction au compas seul.

\begin{center}
    \includegraphics[scale=1]{images/img-mour-041}
\end{center}

% [[figure asymptote]]

%\begin{center}
%\begin{asy}
%size(7cm, 0);
%point O = (0, 0); dot("\(O\)", O, N, blue);
%circle C0 = circle(O, 1); draw(C0, blue);
%point A = (0, -1); dot("\(A\)", A, S);
%circle C1 = circle(A, 1.3); draw(C1);
%point[] B = intersectionpoints(C0, C1);
%dot("\(B\)", B[0], W); dot("\(B'\)", B[1], E);
%circle C2 = circle(B[0], abs(A - B[0])); clipdraw(C2);
%circle C3 = circle(B[1], abs(A - B[1])); clipdraw(C3);
%point[] AD = intersectionpoints(C2, C3);
%point D = (AD[0] == A) ? AD[1] : AD[0];
%dot("\(D\)", D, N);
%circle C4 = circle(D, abs(A - D)); clipdraw(C4);
%point[] F = intersectionpoints(C1, C4);
%dot("\(F\)", F[0], W); dot("\(F'\)", F[1], E);
%clipdraw(circle(F[0], abs(A - F[0])));
%clipdraw(circle(F[1], abs(A - F[1])));
%\end{asy}
%\end{center}

Soit \(\mathcal{C}\) un cercle (dont on ne connaît pas le centre). 
Soit \(A\) un point de \(\mathcal{C}\), et soit \(\mathcal{C}'\) 
un cercle de centre \(A\), qui coupe \(\mathcal{C}\) en deux points 
\(B\) et \(B'\). Les cercles de centres \(B\) et \(B'\) passant 
par \(A\) se coupent en un autre point \(D\). Notons \(\mathcal{C}''\) 
le cercle de centre \(D\) passant par \(A\), et notons \(F\) et \(F'\) 
les points d'intersection de \(\mathcal{C}'\) et \(\mathcal{C}''\). 
Finalement, notons \( \Omega\) l'autre point d'intersection des cercles 
de centres \(F\) et \(F'\) passant par \(A\).
\begin{enumerate}
\item Notons \(O\) le centre du cercle \(\mathcal{C}\), \(r\)~son rayon, 
\(a\)~le rayon du cercle \(\mathcal{C}'\). Montrer que l'on~a 
\(\vec{OA} \cdot \vec{OD} = r^2 - a^2\). (Indication: considérer la 
puissance du point \(O\) par rapport au cercle de centre \(B\) 
passant par \(A\).)
\item En déduire que \(\vec{DA} = \frac{a^2}{r^2} \vec{OA}\) 
et \(DA = \frac{a^2}{r}\).
\item En considérant la puissance de \(D\) par rapport au cercle 
de centre \(F\) passant par \(A\), montrer de même que 
\(\vec{ \Omega A} = \frac{a^2}{DA^2} \vec{DA}\).
\item En déduire que \( \Omega = O\).
\end{enumerate}
\finenonce{007289}
\finexercice


\exercice{7474, exo7, 2021/08/10}
\enonce{007474}{Déterminer une rotation à partir d'images}
 Soit $E$ un plan affine euclidien muni d'un repère cartésien 
orthonormé. Soient
$A$, $B$, $C$ et $D$ les points de $E$ dont les coordonnées sont
$$A : (0,3), \quad B : (2,1), \quad C : (2,3) \quad \text{ et } \quad D 
: (0,1).$$
\begin{enumerate}
\item Montrer que les droites $(A,B)$ et $(C,D)$ sont 
orthogonales et expliciter les coordonnées de leur
point d'intersection.
\item Prouver l'existence d'une {\it {rotation}} qui envoie $A$ 
sur $C$, $C$ sur $B$, $B$ sur $D$ et $D$
sur $A$. Expliciter une représentation matricielle de cette rotation.
\end{enumerate}
\finenonce{007474}
\finexercice


\exercice{7475, exo7, 2021/08/10}
\enonce{007475}{Reconnaître une application affine}
Soit $E$ un plan affine euclidien. Soient $\mathcal{R}$ un repère 
cartésien orthonormé de $E$,
et $f : E\to E$ l'application affine
définie dans $\mathcal{R}$ par l'égalité
$$f \bigl((x,y) \bigr)=\biggl({3x+4y+8\over 5},{4x-3y-1\over 5}
\biggr).$$
\begin{enumerate}
\item L'application $f$ possède-t-elle des points fixes ?
\item Démontrer que lorsque $M$ décrit $E$, le milieu de 
$(M,f(M))$ décrit une droite dont on précisera
l'équation.
\item Démontrer que $f$ est une isométrie dont on précisera la 
nature, l'axe et la composante translation.
\end{enumerate}
\finenonce{007475}
\finexercice


\exercice{7478, mourougane, 2021/08/10}
\enonce{007478}{Sur un groupe d'isométries}
 Dans un plan affine euclidien orienté on considère deux points 
distincts $O$ et $A$. On note $r$ la rotation de centre $O$ et d'angle $2\pi/3$
et $\rho$ la rotation de centre $A$ et d'angle $2\pi/3$. On pose $B=r(A)$ et
$C=r(B)$. Enfin on note $G$ le groupe d'isométries engendré par $r$
et $\rho$. 

\begin{enumerate}
\item Montrer que $G$ ne contient que des translations et des
 rotations d'angle $2\pi/3$ et $-2\pi/3$. 


\item Expliciter une relation de dépendance entre les vecteurs $\vec{OA},
 \vec{OB}$  et $\vec{OC}$
 (on pourra remarquer que la somme de ces vecteurs est
 invariante par $r$). 


\item Montrer que $r\circ\rho ^{-1}$  et $r^{-1} \circ \rho$
 sont des translations dont on
 précisera le vecteur (on pourra étudier l'image de $A$).


\item Montrer que $G$ contient toutes les translations de vecteur
 $p \vec{OA} + q \vec{OB}$ avec $p\in \Zz$, $q\in\Zz$ et $p+q\in 3\Zz$.
\end{enumerate}
\finenonce{007478}
\finexercice


\exercice{7487, mourougane, 2021/08/10}
\enonce{007487}{}
Dans le plan affine $P$ réel, muni d'un repère affine $(A_0,A_1,A_2)$, à quelle condition le point $M$ de coordonnées barycentriques $(\alpha,\beta,\gamma)$ et les points $A(1,5,-2)$ et $B(3,2,0)$ sont-ils alignés ?
\finenonce{007487}
\finexercice


\exercice{7488, mourougane, 2021/08/10}
\enonce{007488}{}
\begin{enumerate}
    \item Donner la définition d'un sous-espace affine d'un espace affine $E$.
    \item Donner l'exemple d'une application $f~:\R^2\to\R^2$ qui n'est pas affine.
    \item Enoncer et démontrer le sens direct du théorème de Thalès.
\end{enumerate}
\finenonce{007488}
\finexercice


\exercice{7489, mourougane, 2021/08/10}
\enonce{007489}{}
Dans le plan affine $P$ réel, muni d'un repère affine $(A_0,A_1,A_2)$, 
on considère les points donnés en coordonnées barycentriques par $A(2,-1,5)$ et $B(1,1,2)$ et $C(2,3,0)$.
Déterminer les coordonnées barycentriques normalisées, du barycentre $G$ des points massiques $(A,1)$, $(B,2)$, $(C,-1)$.
\finenonce{007489}
\finexercice


\exercice{7495, mourougane, 2021/08/10}
\enonce{007495}{}
Soit $A(0,0)$ et $B(3,0)$ deux points d'un plan affine euclidien muni d'un repère $\mathcal{R}=(O,i,j)$ 
cartésien orthonormé.
Déterminer l'ensemble des points $M$ du plan tels que $$2MA^2-MB^2=-2.$$
\finenonce{007495}
\finexercice


\exercice{7496, mourougane, 2021/08/10}
\enonce{007496}{Reconnaître une application affine}
Soit $E$ un plan affine euclidien. Soit $\mathcal{R}=(O,i,j)$ un repère 
cartésien orthonormé de $E$.
\begin{enumerate}
    \item Déterminer l'expression analytique de la translation $t$ de vecteur $\vec{u}=3i+j$.
    \item Déterminer l'expression analytique de la symétrie orthogonale $s$ d'axe la droite $d$ d'équation $(x+y=1)$.
    \item Déterminer l'expression analytique de la composée $f=t\circ s$. 
    \item Démontrer que $f$ est une isométrie. Préciser le déterminant de sa partie linéaire. Que peut-on en déduire ?
    \item Déterminer l'ensemble des points fixes de $f$.
    Déterminer son axe (on pourra montrer que pour tout point $M$ du plan le milieu du segment $[M,f(M)]$ est sur une droite dont on précisera l'équation).
    \item Déterminer la composante de translation de $f$.
\end{enumerate}
\finenonce{007496}
\finexercice


\exercice{7497, mourougane, 2021/08/10}
\enonce{007497}{}
On reprend les notations de l'exercice précédent. 
Soit $E$ un plan affine euclidien. Soit $\mathcal{R}=(O,i,j)$ un repère 
cartésien orthonormé de $E$. On considère la translation $t$ de vecteur $\vec{u}=3i+j$ et la symétrie orthogonale $s$ d'axe la droite $d$ d'équation $(x+y=1)$.
\begin{enumerate} 
    \item Décomposer le vecteur $\vec{u}$ dans la somme directe orthogonale
    $\vec{E}=\vec{d}^\perp\oplus\vec{d}$ en $\vec{u}=\vec{v}+\vec{w}$.
    \item Déterminer géométriquement la nature et les éléments caractéristiques de la composée 
    $t_{\vec{v}}\circ s$. On pourra décomposer chaque isométrie en produits de réflexions.
    \item Déterminer géométriquement la nature et les éléments caractéristiques de la composée 
    $t\circ s$.
\end{enumerate}
\finenonce{007497}
\finexercice


\exercice{7498, mourougane, 2021/08/10}
\enonce{007498}{}
Rappeler la construction des centres d'homothéties qui envoient un cercle sur un autre de rayon différent.

Construire un cercle tangent à deux droites données et passant par un point donné.
\finenonce{007498}
\finexercice


\exercice{7501, mourougane, 2021/08/10}
\enonce{007501}{Questions de cours}
\begin{enumerate}
\item Soit $P$ un plan affine et $(A,B,C)$ un repère affine.
Donner la condition d'alignement de trois points
 en coordonnées barycentriques.

\item Démontrer que si $f$ est une application linéaire orthogonale 
d'un espace vectoriel euclidien $\vec{E}$ de dimension finie, alors
$Ker (f-Id_{\vec{E}})$ et $Im (f-Id_{\vec{E}})$ sont supplémentaires orthogonaux.

\item Donner l'exemple d'une application affine qui n'est pas une isométrie.

\item Les isométries de déterminant $-1$
dans un espace affine euclidien de dimension trois
ont-elles toutes un point fixe ?

\item Peut-on écrire une rotation dans le plan euclidien comme composée de cinq
réflexions ?

\item Dans l'espace affine euclidien muni d'un repère orthonormé, 
on considère le cube de sommets $A(-1,-1,-1)$, $B(-1,-1,1)$, 
$C(-1,1,1)$, $D(-1,1,-1)$, $E(1,-1,-1)$, $F(1,-1,1)$, $G(1,1,1)$, $H(1,1,-1)$.
 Existe-t-il une isométrie de l'espace qui envoie $A$ sur $B$ et $G$ sur $F$ ?
Si oui, la déterminer.
\end{enumerate}
\finenonce{007501}
\finexercice


\exercice{7502, mourougane, 2021/08/10}
\enonce{007502}{Ligne de niveau d'une fonction de Leibniz}
On munit le plan affine euclidien $(P,<,>)$ d'un repère orthonormé 
$(O,\vec{\imath},\vec{\jmath})$.
On considère les points $A(-1,2)$ et $B(5,4)$.
\begin{enumerate}
    \item Déterminer les coordonnées du barycentre $G$ des points massiques 
    $A(-3)$ et $B(1)$.
    \item Calculer $-3GA^ 2+GB^2$.
    \item Démontrer que pour tout point $M$ du plan, on a $$-3MA^2+MB^2=-2MG^2 -3GA^ 2+GB^2.$$
    \item Déterminer l'ensemble $\mathcal L$ des points $M$ du plan tels que $-3MA^2+MB^2=50$.
\end{enumerate}
\finenonce{007502}
\finexercice


\exercice{7503, mourougane, 2021/08/10}
\enonce{007503}{Isométries et constructions}
\begin{enumerate}
    \item On considère dans le plan euclidien orienté un point $A$ et 
    la rotation $r$ de centre $A$ d'angle $+\pi/2$.
    Soit $M$ un point et $M'=r(M)$ son image par $r$. Soit $d$ une droite passant par $M$.
    Décrire un point et la direction de l'image $r(d)$ de la droite $d$.
    
    
    \item
    Soit $d_1$ et $d_2$ deux droites et $B\in d_1$ et $C\in d_2$ tel que 
    $ABC$ soit un triangle rectangle isocèle en $A$ (avec $mes \widehat{(\vec{AB},\vec{AC})}=+\pi/2$).
    Démontrer que $C$ appartient à l'image de la droite $d_1$ par $r$.
    
    \item
    Soit $\delta_1$ et $\delta_2$ deux droites non perpendiculaires et $E$ un point du plan. 
    Construire un triangle $EFG$ rectangle isocèle en $E$ et 
    tel que $F\in \delta_1$ et $G\in \delta_2$.
\end{enumerate}
\finenonce{007503}
\finexercice


\exercice{7504, mourougane, 2021/08/10}
\enonce{007504}{Produits de réflexions et composition}
Dans le plan euclidien orienté muni d'un repère orthonormé direct $(O,\vec{\imath},\vec{\jmath})$, 
on considère la translation $t$ de vecteur $\vec{u}= \vec{\imath}+2\vec{\jmath}$
et la rotation $r$ de centre $A(-3,-1)$ et d'angle $+\pi/2$.
\begin{enumerate}
    \item Montrer que $t\circ r$ a un unique point fixe.
    \item
    Décomposer la translation $t$ en produit $s_{d_2}\circ s_{d_1}$ de réflexions
    par rapport à des droites $d_1$ et $d_2$ que l'on décrira.
    \item Décomposer la rotation $r$ en produit $s_{d_4}\circ s_{d_3}$ de réflexions
    par rapport à des droites $d_3$ et $d_4$ que l'on décrira.
    \item Est-il possible de choisir $d_1=d_4$ dans les questions précédentes ?
    \item Déterminer par une méthode géométrique la nature et 
    les éléments caractéristiques de la composée $t\circ r$.
\end{enumerate}
\finenonce{007504}
\finexercice


\exercice{7505, mourougane, 2021/08/10}
\enonce{007505}{}
On munit le plan affine euclidien $(P,<,>)$ d'un repère orthonormé
$(O,\vec{\imath},\vec{\jmath})$.
On considère les points $A$ de coordonnées $(-2,1)$ et $B$ de coordonnées $(4,4)$.

\begin{enumerate}
    \item Déterminer les coordonnées du barycentre $G$ des points massiques $A(2)$ et $B(1)$.
    \item Calculer $2GA^ 2+GB^2$.
    \item Démontrer que pour tout point $M$ du plan, on a $$2MA^2+MB^2=3MG^2 +2GA^ 2+GB^2.$$
    \item Déterminer l'ensemble $\mathcal L$ des points $M$ du plan tels que $2MA^2+MB^2=42$.
\end{enumerate}
\finenonce{007505}
\finexercice


\exercice{7506, mourougane, 2021/08/10}
\enonce{007506}{}
\begin{enumerate}
    \item On identifie le plan affine $\R^2$ muni du produit scalaire standard à l'espace vectoriel $\C$ muni du produit scalaire $(z,z')=re(z\overline{z'})$ par l'application linéaire
    $(1,0)\mapsto 1$, $(0,1)\mapsto i$.
    Parmi les transformations suivantes, lesquelles sont des isométries du plan euclidien 
    \begin{enumerate}
        \item[a.] $\ z\mapsto 3z+4$
        \item[b.] $\ z\mapsto 3\overline{z}+4$
        \item[c.] $\ z\mapsto \overline{z}+4$
        \item[d.] $\ z\mapsto i\overline{z}+4$
    \end{enumerate}
    
    \item Pour chacune des isométries, dire s'il s'agit d'une translation, d'une rotation ... (sans préciser les éléments caractéristiques).
\end{enumerate}
\finenonce{007506}
\finexercice


\exercice{7507, mourougane, 2021/08/10}
\enonce{007507}{}
\begin{enumerate}
    \item Démontrer que si $A, B, C$ sont trois points distincts d'un plan affine euclidien $\mathcal{P}$
    la somme des angles de vecteurs 
    $$((\vec{AB}, \vec{AC}))+((\vec{BC},\vec{BA}))+((\vec{CA},\vec{CB}))$$
    est un angle plat.
    \item Soit $\mathcal{C}$ un cercle de $\mathcal{P}$ et $A$ un point de $\mathcal{C}$.
    Soit $\mathcal{C}'$ l'image par une rotation $r$ de centre $A$ du cercle $\mathcal{C}$.
    Soit $B$ l'autre point d'intersection de $\mathcal{C}$ et $\mathcal{C}'$.
    Soit $D$ le point de $\mathcal{C}$ diamétralement opposé à $A$ sur $\mathcal{C}$.
    Soit $D'=r(D)$ son image par $r$. Montrer que $D$, $D'$ et $B$ sont alignés.
    \item Soit $M$ un point quelconque de $\mathcal{C}$.
    Montrer que $M$, $M'=r(M)$ et $B$ sont alignés.
\end{enumerate}
\finenonce{007507}
\finexercice


\exercice{7508, mourougane, 2021/08/10}
\enonce{007508}{}
Soit $\R^2$ le plan affine euclidien muni du produit scalaire standard et de la base canonique.
\begin{enumerate}
    \item Ecrire la matrice $A$ de la forme bilinéaire symétrique donnée en coordonnées par
    $$f(\begin{pmatrix}x\\y\end{pmatrix},\begin{pmatrix}x'\\y'\end{pmatrix})
    =xx'+19yy'+12xy'+12x'y.$$
    \item Diagonaliser $A$ dans une base orthonormée.
    \item Déterminer la nature de la conique d'équation
    $$x^2+19y^2+24xy+5y=0.$$
\end{enumerate}
\finenonce{007508}
\finexercice


\exercice{7512, mourougane, 2021/08/10}
\enonce{007512}{}
Soit $A, B ,C$ un repère affine d'un plan affine $E$.
\begin{enumerate}
    \item Déterminer les équations barycentriques des médianes du triangle $ABC$.
    \item En utilisant la question précédente, montrer que les médianes du triangle $ABC$ sont concourantes.
\end{enumerate}
\finenonce{007512}
\finexercice


\exercice{7513, mourougane, 2021/08/10}
\enonce{007513}{}
Soit $P$ un plan euclidien orienté.
\begin{enumerate}
    \item Donner la liste des éléments du groupe des isométries du plan qui conservent un carré.
    \item Ce groupe est-il commutatif ?
    \item Le groupe des déplacements du plan qui conservent un carré est-il commutatif ?
\end{enumerate}

\finenonce{007513}
\finexercice


\exercice{7514, mourougane, 2021/08/10}
\enonce{007514}{}
Construire un point $A$ sur le cercle $\mathcal{C}$ dont le symétrique par rapport à $O$ est sur la droite $d$.

\includegraphics[scale=0.5]{images/img-mour-269-1}

Justifier votre construction.
\finenonce{007514}
\finexercice

\section{242.02 Géométrie affine euclidienne de l'espace}


\exercice{2031, liousse, 2003/10/01}

\enonce{002031}{}
On consid\`ere les 4 points $A$, $B$, $C$, $D$ donn\'es. 
$(A,\vec{AB},\vec{AC},\vec{AD})$ d\'efinit-il bien un nouveau rep\`ere ? 
Dans ce cas, trouver les formules de changements de rep\`ere exprimant
 les coordonn\'ees $(x,y,z)$ dans $(O,\vec i, \vec j,\vec k)$
  en fonction de celles $(x',y',z')$ dans $(A, \vec{AB},\vec{AC},\vec{AD})$.

\begin{enumerate}
\item $A(2,-1,0)$, $B(7,-1,-1)$, $C(-3,0,-2)$, $D(3,-6,-3)$
\item $A(4,1,4)$, $B(7,3,1)$, $C(9,0,0)$, $D(5,2,3)$
\item $A(0,-1,3)$, $B(5,-6,4)$, $C(-4,1,-2)$, $D(-3,3,6)$
\item $A(1,1,0)$, $B(1,5,2)$, $C(0,-1,1)$, $D(3,4,-1)$
\item $A(2,-1,4)$, $B(0,0,1)$, $C(3,2,-1)$, $D(1,3,4)$
\item $A(4,4,2)$, $B(5,3,2)$, $C(4,3,3)$, $D(3,5,2)$  
\item $A (1,3,1)$ ,$B (1,2,2)$ ,$C (2,-1,-4)$, $D(0,8,6)$.
\end{enumerate} 
\finenonce{002031}



\finexercice



\exercice{2032, liousse, 2003/10/01}

\enonce{002032}{}
Les formules suivantes d\'efinissent-elles bien un changement de rep\`ere ? Dans ce cas, 
donner le changement de rep\`ere inverse.

\begin{enumerate}
\item $\left\{ 
         \begin{array}{l}
           x' =  y - z + 1 \\
           y' = -x - 4y + 5z + 2 \\
           z' = x - 5y + 5z + 1
         \end{array}
      \right. $

\item $\left\{  
         \begin{array}{l}
           x' = 5x + 4y + 3z - 2 \\
           y' = 2x + 3y + z + 2 \\
           z' = 4x - y + 3z + 2
         \end{array}
      \right. $  

\item $\left\{  
         \begin{array}{l}
           x' = -2x - 4y + 2z - 2 \\
           y' = x + y - 5z + 1 \\
           z' = -3x - 4y + 4z - 2
         \end{array}
      \right. $  

\item  $\left\{  
         \begin{array}{l}
           x' = 3x - 5y + z + 2 \\
           y' = 2x - y + z - 1 \\
           z' = -3x - 4y - z - 5
         \end{array}
      \right. $ 

\item $\left\{  
         \begin{array}{l}
           x' = 2x - z + 1 \\
           y' = -2x + 2y + 2z - 2 \\
           z' = -2x + y - z 
         \end{array}
      \right. $  

\item $\left\{  
         \begin{array}{l}
           x' = x - 2y - 3z + 5 \\
           y' = -3x + 4y + z - 2 \\
           z' = 2x - y + 6z + 3
         \end{array}
      \right. $             
\end{enumerate}
\finenonce{002032}



\finexercice



\exercice{2033, liousse, 2003/10/01}

\enonce{002033}{}

On consid\`ere les droites et les plans suivants dont les \'equations sont donn\'ees dans le
rep\`ere $(O,{\buildrel\rightarrow \over i},{\buildrel\rightarrow \over j}
,{\buildrel\rightarrow \over k})$. 
Donner leurs \'equations dans le nouveau rep\`ere $(A,
{\buildrel\rightarrow \over {AB}}, {\buildrel\rightarrow \over {AC}}, 
{\buildrel\rightarrow \over {AD}})$, sachant que dans $(O,{\buildrel\rightarrow \over i},
{\buildrel\rightarrow \over j},{\buildrel\rightarrow \over k})$ les points $A, B, C$ et
$D$ ont pour coordonn\'ees respectives $A (4,-1,2)$, $B (2, -5, 4)$, $C(5, 0, -3)$, $D(1,-5,6)$.
\begin{enumerate}
\item  $ P : x+y=1 $ 
\item  $ P : 2x -3y +4z -1=0$
\item  $ P : x-y+z+3=0  $ 
\item  $P : \left\{  
         \begin{array}{l}
           x = 2t +3s + 1 \\
           y = t -s +2 \\
           z = 4t - 2s - 3
         \end{array}
     \right. $  

\item  $(D):\left\{ \begin{array}{l} x+y+z=1 \\ 2x-y+4z=3 \end{array}\right. $ 
\item  $(D):\left\{ \begin{array}{l} 3x-y-z=-1 \\ 4x-3y-z=-2 \end{array}\right. $  
\end{enumerate}
\finenonce{002033}



\finexercice




\exercice{2034, liousse, 2003/10/01}
\video{W0w1WraqsVc}
\enonce{002034}{}

\begin{enumerate}
\item  On considère le point $A (-2,4,1) $, les vecteurs ${\buildrel\rightarrow \over {u}}
 (1,1,1), {\buildrel\rightarrow \over {v}} (2,2,-4) $, 
${\buildrel\rightarrow \over {w}} (3,-1,1)$
et le repère $(A, {\buildrel\rightarrow \over {u}} , 
{\buildrel\rightarrow \over {v}} , {\buildrel\rightarrow \over {w}}) $.
On note $x',y'$et $z'$ les coordonnées dans ce repère. Donner les formules analytiques du changement de 
repère exprimant $x,y,z$ en fonction de $x',y',z'$.

\item On considère la droite $(D):\left\{\begin{array}{l} y-z=3 \\ x+y=2 \end{array}\right.$. 
Utiliser le changement de repère pour donner une équation de $D$ dans le repère
$(A, {\buildrel\rightarrow \over {u}} , 
{\buildrel\rightarrow \over {v}} , {\buildrel\rightarrow \over {w}}) $.

\item  Donner les formules analytiques du changement de repère inverse.
\end{enumerate}
\finenonce{002034}



\finexercice


\exercice{2038, liousse, 2003/10/01}

\enonce{002038}{}
\begin{enumerate}
\item  Soit $f$ la transformation 
de l'espace d\'efinie analytiquement par
$$\left\{\begin{array}{ll}x'=&-3x+2y-2z+4\\y'=&-8x+5y-4z+8\\z'=&-4x+2y-z+4
\end{array}\right.$$
\begin{enumerate}
\item D\'eterminer l'ensemble $P$ des points invariants par $f.$

\item  Montrer que pour $M$ d'image $M',$ le milieu de $[MM']$ est dans $P,$ 
(MM') est parall\`ele \`a une direction fixe.

\item  En d\'eduire une description simple de $f.$
\end{enumerate}
\vskip1mm

\item  Soit $f$ la transformation 
de l'espace d\'efinie analytiquement par
$$\left\{\begin{array}{ll}x'={1\over 3}(&2x-y-z+1)\\y'={1\over 3}(&-x+2y-z+1)\\z'={1\over 3}(&-x-y+2z+1)
\end{array}\right.$$
\begin{enumerate}
\item  D\'eterminer l'ensemble $P$ des points invariants par $f.$

\item  Montrer que pour $M$ d'image $M'$ le vecteur $\vec{MM'}$ est colin\'eaire \`a un vecteur fixe.

\item En d\'eduire une description simple de $f.$
\end{enumerate}
\end{enumerate}
\finenonce{002038}



\finexercice



\exercice{2039, liousse, 2003/10/01}
\video{csVlnEEMcVw}
\enonce{002039}{}
\begin{enumerate}
\item Définir analytiquement la projection orthogonale sur le plan  d'équation $2x+2y-z=1$.

\item Définir analytiquement la projection orthogonale sur la droite 
d'équation $\left\{\begin{array}{l}x+y+z=1 \\2x-z=2 \end{array}\right.$.

\item   Donner l'expression analytique de la projection sur le plan $(P)$ contenant 
le point $C(2,-1,1)$ et ayant pour vecteurs directeurs $\vec {u}(0,-1,1)$ et $\vec {u'}(-2,0,1)$,  
selon la droite $(AB)$, où $A(1,-1,0)$ et $B(0,-1,3)$.
\end{enumerate}
\finenonce{002039}



\finexercice


\exercice{2042, liousse, 2003/10/01}

\enonce{002042}{}
Tout ce probl\`eme se situe dans l'espace euclidien tridimensionnel muni d'un rep\`ere
 orthonorm\'e direct $\mathcal R 
= (0,\overrightarrow{i}, \overrightarrow{j},\overrightarrow{k})$.

\begin{enumerate}
\item
On consid\`ere les
deux droites $d$ et $D$ donn\'ees par les syst\`emes d'\'equations 
cart\'esiennes suivant :

$d \left\lbrace 
\begin{array}{ll} x +y-3z& = 0\cr y +z& =0 \end{array}\right.$\hfil et \hfil
$D\left\lbrace \begin{array}{ll}x -1& = 0 \cr y - z -1& =0 \end{array}\right.$  

\begin{enumerate}
\item
\begin{enumerate}
\item Donner un point et un vecteur directeur de $d$.
 Donner un point et un vecteur directeur  de $D$.

\item Dire si les droites $d$ et $D$ sont parall\`eles, s\'ecantes 
 ou non coplanaires.

\item Justifier l'existence de deux plans parall\`eles  (en donnant pour chacun de ces 
deux plans un point et deux vecteurs directeurs) tels que $d$ est contenue dans l'un 
 et $D$ est contenue dans l'autre.
\end{enumerate} 

\item
\begin{enumerate}
\item
 Soient $\overrightarrow{u}$ le vecteur de coordonn\'ees $(4,-1,1)$
 dans $\mathcal R$, $\overrightarrow{v}$ le vecteur de coordonn\'ees $(0,1,1)$
 dans $\mathcal R$ et $\Omega$ le point de coordonn\'ees $(1,1,0)$
 dans $\mathcal R$. 

D\'eterminer une \'equation cart\'esienne pour le  
 plan  $P$ de rep\`ere cart\'esien $(O,\overrightarrow{u},\overrightarrow{v})$, en d\'eduire 
 une \'equation cart\'esienne pour le  
 plan $Q$  de rep\`ere cart\'esien $(\Omega,\overrightarrow{u},\overrightarrow{v})$.

\item Donner des  \'equations param\'etriques pour
la  droite $\Delta$  normale \`a $P$
 passant par $O$. D\'eterminer les deux points 
$\Delta\cap P$ et $\Delta\cap Q$ puis  calculer la distance entre eux.
\end{enumerate}
\end{enumerate}

Interpr\'eter cette distance.



\item
 On consid\`ere les vecteurs de l'espace $
\overrightarrow{a} = ({1\over 3}, {2\over 3}, -{2\over 3})$,  
$\overrightarrow{b} = ({2\over 3}, {1\over 3}, {2\over 3})$, 
$\overrightarrow{c} = ({-2\over 3}, {2\over 3}, {1\over 3}).$

\begin{enumerate}
\item Montrer que $(0, \overrightarrow{a},\overrightarrow{b},
\overrightarrow{c})$ est un rep\`ere orthonorm\'e. Est-il direct ?

\item  Ecrire les formules de changement de rep\`eres
 de $\mathcal R$ \`a  $(0, \overrightarrow{a},\overrightarrow{b},
\overrightarrow{c})$.

\item Quelle est l'\'equation dans le rep\`ere
 $(0, \overrightarrow {a},\overrightarrow {b},\overrightarrow {c})$ 
du plan  d'\'equation $x+2y-2z = 0$ dans $\mathcal R$ ? M\^eme question avec le
 plan d'\'equation $x+2y-2z= 3$ dans $\mathcal R$.
\end{enumerate}
\end{enumerate}
\finenonce{002042}



\finexercice



\exercice{2043, liousse, 2003/10/01}

\enonce{002043}{}
Tout ce probl\`eme se situe dans l'espace euclidien tridimensionnel muni d'un rep\`ere
 orthonorm\'e direct $\mathcal R 
= (O,\overrightarrow{i}, \overrightarrow{j},\overrightarrow{k})$.
On d\'efinit les trois  points : $A=(3, \sqrt{6},3)$, $B=(3, -\sqrt{6},3)$  et $C=(4,0,0)$.
\begin{enumerate}
\item
\begin{enumerate}
\item Montrer que les points $O$, $A$ et $B$ ne sont pas align\'es et donner
une \'equation cart\'esienne du plan $P$ contenant $O$, $A$ et $B$.
\item Calculer les distances $OA$, $OB$ et $AB$. En d\'eduire la nature du triangle
$OAB$.
\item Les points $O,A,B$ et $C$ sont-ils coplanaires ?
\end{enumerate}


\item Soit $G$ le barycentre des points $O$, $A$, $B$ et $C$, c'est \`a dire, par 
d\'efinition l'unique point $G$ de l'espace tel que :  $\overrightarrow{GO}+\overrightarrow{GA}+\overrightarrow{GB}+\overrightarrow{GC}=
\overrightarrow{0}$.
\begin{enumerate}
\item Montrer que   $\overrightarrow{OG}= {1\over 4}(
\overrightarrow{OA}+\overrightarrow{OB}+\overrightarrow{OC}).$
\item En d\'eduire les coordonn\'ees de $G$ dans $\mathcal R$. 
\end{enumerate}

\item
\begin{enumerate}
\item Montrer que la droite $(GC)$ est perpendiculaire au plan $P$.
\item Calculer les coordonn\'ees du
 point d'intersection de la droite $(GC)$ avec le plan $P$.
\end{enumerate}

\item Montrer que la transformation de l'espace d\'efinie par les formules :
$(x'=x, y'=-y, z'=z)$ est une isom\'etrie.
Quels sont ses points fixes ? 
D\'eterminer les images des points $O,A,B,C$ par cette isom\'etrie. Que remarque-t-on~?
\end{enumerate}
\finenonce{002043}



\finexercice



\exercice{2044, liousse, 2003/10/01}

\enonce{002044}{}
L'espace est rapport\'e \`a un rep\`ere orthonorm\'e direct 
$(0,\vec{\strut\imath},\vec{\strut\jmath},\vec{\strut k})$. On d\'efinit 
les points
$$A :\, (1,2,3) \; ; \quad B :\, (2,3,1) \; ; \quad C :\, (3,1,2) \; ; 
\quad D :\; (1,1,1)$$
et le plan
$$\Pi : 2x-3y+4z=0.$$
\begin{enumerate}
\item Montrer que les points $A,B,C$ ne sont pas align\'es.
\item Montrer que les points $A,B,C,D$ ne sont pas coplanaires.
\item Donner une \'equation cart\'esienne du plan $P$ passant par $A,B,C$.
\item Calculer la distance de $D$ au plan $P$.
\item Donner une repr\'esentation param\'etrique de la droite $d=P\cap\Pi$.
\end{enumerate}
\finenonce{002044}



\finexercice



\exercice{2045, liousse, 2003/10/01}

\enonce{002045}{}
Soient $A$, $B$ et $C$ trois points distincts et non align\'es de l'espace affine 
tridimensionnel $\mathcal E$.
On note $P$ le plan qui contient $A$, $B$ et $C$.
Soit $O$ un point de $\mathcal E$ n'appartenant pas \`a $P$.
\begin{enumerate}
\item

\begin{enumerate}
\item
Expliquer rapidement pourquoi $\mathcal R = (O,\overrightarrow{OA},\overrightarrow{OB},
\overrightarrow{OC})$ est un rep\`ere  cart\'esien de $\mathcal E$.
\item Dans ce rep\`ere $\mathcal R$, \'ecrire les coordonn\'ees des points $O$, $A$, $B$ et $C$, et d\'eterminer une
\'equation cart\'esienne du plan $P$.
\end{enumerate}


\item Soit $A'$ le point de la droite $(OA)$ tel que 
$\overrightarrow{OA'}= 2 \overrightarrow{OA}$.
On note $P'$ le plan parall\`ele \`a $P$ passant par $A'$.
$P'$ coupe $(OB)$ en $B'$ et $(OC)$ en $C'$.

Dans $\mathcal R$, \'ecrire les  coordonn\'ees des points $A'$, $B'$ et $C'$  et d\'eterminer des
\'equations param\'etriques pour les droites $(BC')$ et $(B'C)$, en d\'eduire des \'equations 
cart\'esiennes de ces droites.

Calculer les coordonn\'ees des points $I=(BC')\cap(B'C)$, $J=(AC')\cap(A'C)$ et 
$K=(AB')\cap(A'B)$.



\item Soit $A''$ le point de la droite $(OA)$ tel que 
$\overrightarrow{OA''}=-{ 2\over 3} \overrightarrow{OA}$.
On note $P''$ le plan parall\`ele \`a $P$ passant par $A''$.
$P''$ coupe $(OB)$ en $B''$ et $(OC)$ en $C''$.

Montrer que les droites $(IA'')$, $(JB'')$, $(KC'')$ sont parall\`eles.
\end{enumerate}
\finenonce{002045}



\finexercice



\exercice{4958, quercia, 2010/03/17}
\enonce{004958}{\'Equation au produit vectoriel}

Soient $A$, $B$, $C$ trois points distincts de l'espace.

Déterminer le lieu des points $M$ tels que
$\vec{MA}\wedge\vec{MB} + \vec{MB}\wedge\vec{MC} = 2\vec{MC}\wedge\vec{MA}$.

\finenonce{004958}



\finexercice


\exercice{4959, quercia, 2010/03/17}
\enonce{004959}{$\| \vec{MA} + 2\vec{MB} + k\vec{MC} \| = \| \vec{MD} + \vec{ME} \|$}

Soient $A,B,C,D,E$ cinq points de l'espace et $k \in \R$.

Déterminer le lieu des points $M$ de l'espace tels que
$\| \vec{MA} + 2\vec{MB} + k\vec{MC} \| = \| \vec{MD} + \vec{ME} \|$.
\finenonce{004959}



\finexercice


\exercice{4960, quercia, 2010/03/17}
\enonce{004960}{Ensi Chimie P 93}

Trouver les coordonnées des projetés du point $C(3,4,-2)$ sur les droites
définies par les équations~:

$D_1 : \frac{x-5}{13} = \frac{y-1}6 = \frac{z+3}{-4}$.

$D_2 : \frac{x-2}{13} = \frac{y-3}1 = \frac{z-3}{-4}$.
\finenonce{004960}



\finexercice


\exercice{4961, quercia, 2010/03/17}
\enonce{004961}{Projections sur 4 plans}

Dans un \emph{rond} on donne les plans
$\begin{cases}
P&x+y=1\cr Q&y+z=1\cr R&x+z=1\cr S&x+3y+z=0 \cr\end{cases}$
et le point $A : (1,1,\lambda)$.

Donner une CNS sur $\lambda$ pour que les projections de $A$ sur les
quatre plans soient coplanaires.

\finenonce{004961}



\finexercice


\exercice{4962, quercia, 2010/03/17}
\enonce{004962}{Calculs de points et plans}

Dans un \emph{rond} on donne les points
$A:(1,2,3)$, $B:(2,3,1)$, $C:(3,1,2)$, $D:(1,0,-1)$.

\begin{enumerate}
  \item Chercher le centre et le rayon de la sphère circonscrite à $ABCD$.
    
  \item Chercher les équations cartésiennes des plans $(ABC)$, $(ABD)$, $(ACD)$,
    $(BCD)$.
    
  \item Chercher le centre et le rayon de la sphère inscrite dans le tétraèdre
    $ABCD$.
    
\end{enumerate}
\finenonce{004962}



\finexercice


\exercice{4963, quercia, 2010/03/17}
\enonce{004963}{Perpendiculaire commune à deux droites}

Dans un \emph{rond} on donne les droites
$D  : \begin{cases}x - y + z = -1 \cr 2x + y - z = 0 \cr\end{cases}$ et
$D' : \begin{cases}x +2y + z = 0  \cr  x - y - z = \frac 12.\cr\end{cases}$

Chercher la perpendiculaire commune, $\Delta$, à $D$ et $D'$
(On donnera les points $H\in D\cap\Delta$ et $K\in D'\cap \Delta$).

\finenonce{004963}



\finexercice


\exercice{4964, quercia, 2010/03/17}
\enonce{004964}{Perpendiculaire commune à deux droites}
Dans un \emph{rond} on donne les droites
$D  : \begin{cases}x +2y - z = 1  \cr 2x - y +2z = 2 \cr\end{cases}$ et
$D' : \begin{cases}x + y + z = 3  \cr  x - y +2z = 0.\cr\end{cases}$

Calculer $d(D,D')$.
\finenonce{004964}



\finexercice


\exercice{4965, quercia, 2010/03/17}
\enonce{004965}{Tétraèdre dont les faces ont même aire}

Soit $ABCD$ un tétraèdre dont les quatre faces ont même aire.
Montrer que les côtés non coplanaires ont deux à deux mêmes longueurs.

\finenonce{004965}



\finexercice


\exercice{4966, quercia, 2010/03/17}
\enonce{004966}{Distance entre les côtés d'un tétraèdre}

Soit $ABCD$ un tétraèdre régulier de côté $a$.
Chercher la distance entre deux côtés non coplanaires.

\finenonce{004966}



\finexercice


\exercice{4967, quercia, 2010/03/17}
\enonce{004967}{Distance d'un point à une droite}

Dans un \emph{rond} on donne la droite
$D  : \begin{cases} x +2y - z = -3 \cr  x - y +2z = -4\cr\end{cases}$ et $M(x,y,z)$.
Calculer $d(M,D)$.

\finenonce{004967}



\finexercice


\exercice{4968, quercia, 2010/03/17}
\enonce{004968}{Projection orthogonale}

Dans un \emph{rond} on donne le plan  $P : x + 2y + 3z = 4$.
Déterminer l'expression analytique de la projection orthogonale sur $P$.
\finenonce{004968}



\finexercice


\exercice{4969, quercia, 2010/03/17}
\enonce{004969}{Projection orthogonale}
Dans un \emph{rond} on donne les points
$A:(1, 0, -1)$, $B:(-1, 1, 1)$, $C:(2, -1, 1)$, $D:(1, 2, -2)$, $E:(-2, -2, 0)$.

Déterminer, par un point et un vecteur directeur, la projection de $(DE)$
sur le plan $(ABC)$.
\finenonce{004969}



\finexercice


\exercice{4970, quercia, 2010/03/17}
\enonce{004970}{Symétrique d'un plan}

Dans un \emph{rond} on donne les plans
$P : x+y+z = 1$ et $Q : 2x-y+z=1$.
Chercher une équation cartésienne du plan $Q'$ symétrique de $Q$ par rapport
à $P$.

\finenonce{004970}



\finexercice


\exercice{4971, quercia, 2010/03/17}
\enonce{004971}{Repères orthonormés}

Soient $(O,\vec{OA}, \vec{OB}, \vec{OC})$ et $(O,\vec{OA'}, \vec{OB'}, \vec{OC'})$
deux repères orthonormés directs de l'espace.

Montrer que $\vec{AA'}$, $\vec{BB'}$, $\vec{CC'}$ sont coplanaires.
\finenonce{004971}



\finexercice


\exercice{4972, quercia, 2010/03/17}
\enonce{004972}{Angle d'un plan et d'une droite}

Soient $P$ un plan, $D$ une droite tels que
$\overline{(P,D)} \equiv \theta (\mathrm{mod}\,\pi)$.

Montrer que pour toute droite $\Delta \subset P$, on a
$\cos(D,\Delta) \ge \cos \theta$.
Quand y a-t-il égalité ?
\finenonce{004972}



\finexercice


\exercice{4973, quercia, 2010/03/17}
\enonce{004973}{Angle entre deux faces d'un dodécaèdre}

Quel est l'angle entre deux faces d'un dodécaèdre régulier ?
(on donne : $4\sin\frac\pi5 = \sqrt{10-2\sqrt5}$)
\finenonce{004973}



\finexercice


\exercice{4974, quercia, 2010/03/17}
\enonce{004974}{Ensi P 90}

    Déterminer l'équation de la sphère contenant les cercles d'équations
    $\begin{cases} x^2+y^2 = 9\cr z=0\cr \end{cases}$ et
    $\begin{cases} x^2+y^2 = 25\cr z=2.\cr \end{cases}$
    
\finenonce{004974}



\finexercice


\exercice{4975, quercia, 2010/03/17}
\enonce{004975}{Sphère définie par ses intersections}

Soit $S$ une partie de l'espace contenant au moins deux points et telle
que pour tout plan $P$, $P\cap S$ est un cercle, un singleton ou vide.
Montrer que $S$ est une sphère.
\finenonce{004975}



\finexercice


\exercice{4976, quercia, 2010/03/17}
\enonce{004976}{CNS pour que deux vissages commutent}

Soient $f,g$ deux vissages d'angles $\ne \pi$. Trouver une CNS pour que
$f\circ g = g \circ f$.
(On étudiera $f \circ g \circ f^{-1}$)
\finenonce{004976}



\finexercice


\exercice{4977, quercia, 2010/03/17}
\enonce{004977}{Composée de 3 demi-tours}

Soient $D_1,D_2,D_3$, trois droites, et $\sigma_1,\sigma_2,\sigma_3$ les
$\frac12$-tours correspondants.

Démontrer que $\sigma_1 \circ \sigma_2 \circ \sigma_3$ est un $\frac12$-tour
si et seulement si $D_1,D_2,D_3$ ont une perpendiculaire commune ou sont
parallèles.
\finenonce{004977}



\finexercice


\exercice{4978, quercia, 2010/03/17}
\enonce{004978}{Composée de demi-tours par rapport aux arêtes d'un tétraèdre}

Soit $ABCD$ un tétraèdre régulier, et $d_{AB},d_{AC},d_{AD}$ les $\frac12$-tours
autour des droites $(AB),(AC),(AD)$.
Simplifier $f = d_{AB} \circ d_{AC} \circ d_{AD}$.
\finenonce{004978}



\finexercice


\exercice{4979, quercia, 2010/03/17}
\enonce{004979}{Isométries transformant un triangle en un triangle donné}

Soient $ABC$ et $A'B'C'$ deux triangles tels que
$AB = A'B'$, $AC = A'C'$, $BC = B'C'$.
Combien y a-t-il d'isométries transformant $ABC$ en $A'B'C'$ ?

Indication : si $f$ et $g$ sont deux telles isométries, alors $f\circ g^{-1}$
est une isométrie conservant $ABC$.
\finenonce{004979}



\finexercice


\exercice{4980, quercia, 2010/03/17}
\enonce{004980}{Groupes d'isotropie}

Déterminer {\it toutes\/} les isométries

\begin{enumerate}
  \item d'un tétraèdre régulier.
  \item d'un cube.
  \item de deux droites non coplanaires.
\end{enumerate}
\finenonce{004980}



\finexercice


\exercice{4981, quercia, 2010/03/17}
\enonce{004981}{Composée de projections}

Soient $D_1$, $D_2$, $D_3$ trois droites de l'espace non toutes parallèles.
Pour $M_1 \in D_1$ on construit : $M_2$, projeté de $M_1$ sur $D_2$,
$M_3$, projeté de $M_2$ sur $D_3$,
$M_4$, projeté de $M_3$ sur $D_1$.

Montrer qu'il existe un unique point $M_1 \in D_1$ tel que $M_4 = M_1$.

\finenonce{004981}



\finexercice


\exercice{5501, rouget, 2010/07/10}
\enonce{005501}{**T}
Dans $E_3$ rapporté à un repère $(O,i,j,k)$, on donne les points $A(1,2,-1)$, $B(3,2,0)$, $C(2,1,-1)$ et $D(1,0,4)$. Déterminer l'intersection des plans $(OAB)$ et $(OCD)$.
\finenonce{005501}


\finexercice


\exercice{5502, rouget, 2010/07/10}
\enonce{005502}{**T}
Dans $E_3$ rapporté à un repère $(O,i,j,k)$, on donne~:
la droite $(D)$ dont un système d'équations paramétriques est $\left\{
\begin{array}{l}
x=2+3t\\
y=-t\\
z=1+t
\end{array}
\right.
$,
le plan $P$ dont un système d'équations paramétriques est $\left\{
\begin{array}{l}
x=1+2\lambda+\mu\\
y=-1-3\lambda+2\mu\\
z=1+\lambda
\end{array}
\right.
$,
le plan $P'$ dont un système d'équations paramétriques est $\left\{
\begin{array}{l}
x=-5-\nu\\
y=3+\nu+3\eta\\
z=\nu+\eta
\end{array}
\right.
$,
Etudier $D\cap P$ et $P\cap P'$
\finenonce{005502}


\finexercice


\exercice{5503, rouget, 2010/07/10}
\enonce{005503}{**T}
Matrice dans la base canonique orthonormée directe de la rotation vectorielle de $\Rr^3$ autour de $(1,2,2)$ qui transforme $j$ en $k$.
\finenonce{005503}


\finexercice


\exercice{7473, mourougane, 2021/08/10}
\enonce{007473}{Questions de cours}
Rappeler toutes les isométries du plan euclidien et de l'espace
euclidien de dimension 3, en précisant leur
partie linéaire, leur point fixe, leur axe, leur composante à point
fixe et leur composante de glissement.
\finenonce{007473}
\finexercice


\exercice{7476, mourougane, 2021/08/10}
\enonce{007476}{Trouver l'isométrie}
 Soit $E$ un espace affine euclidien de dimension 3 muni d'un repère 
cartésien orthonormé. On note $v$ la transformation de $E$ dans $E$
qui  envoie le point de coordonnées $(x,y,z)$ sur le point de 
 coordonnées $(x',y',z')$ définies par:
 $$ x' = \frac{2x-2y+z+1}{3}; y' =\frac{2x+y-2z+2}{3} ; z'
 =\frac{x+2y+2z+5}{3}.$$  
  
 Montrer que $v$ est une isométrie de $E$. Préciser de quel type
 d'isométrie il s'agit. Expliciter son axe et son vecteur de
 glissement. 
\finenonce{007476}
\finexercice


\exercice{7477, mourougane, 2021/08/10}
\enonce{007477}{Composée d'isométries}
Soit $E$ un espace affine euclidien de dimension $3$ muni d'un repère cartésien
orthonormé $(O,i,j,k)$. On désigne par
$D$ la droite d'équation $(x=0,z=1)$ et par $D'$ la droite d'équation 
$(y=0,z=0)$. On note $S_D$ la
symétrie par rapport à la droite $D$ et $R_{\theta}$ la rotation 
d'axe $D'$ et d'angle $\theta$ (en
considérant la base $(j,k)$ comme directe). On pose $\varphi=S_D\circ 
R_{\theta}$.
\begin{enumerate}
\item {\'E}crire dans la base $(i,j,k)$ la matrice de 
$\overrightarrow {S_D}$, celle de
$\overrightarrow {R_{\theta}}$ et celle de $\overrightarrow{\varphi}$.
{\'E}crire les expressions analytiques de $S_D$ et de $R_{\theta}$ dans
le repère $(O,i,j,k)$.
\item Montrer que $\varphi$ est une symétrie éventuellement 
glissée d'axe une droite $\Delta$.
\item Pour tout point $M$ de $E$, prouver que les milieux de
$\bigl(M,s_\Delta (M)\bigr)$ et de $\bigl(M,\varphi(M)\bigr)$ 
sont sur $\Delta$.
\item En utilisant le point $O$, montrer que $\Delta$ passe par le
point de coordonnées $(0,0,1)$.
et est contenue dans le plan affine d'équation $x=0$.
\item Donner les composantes du vecteur de glissement de 
$\varphi$ en fonction de $\theta$.
\end{enumerate}
\finenonce{007477}
\finexercice


\exercice{7479, mourougane, 2021/08/10}
\enonce{007479}{Groupe laissant stable une partie}
Soient $E$ un espace affine euclidien de dimension 3 et $\mathcal{R}$ 
un repère cartésien orthonormé
de $E$. Soit
$n$ un entier $\geq 3$. On considère l'ensemble $X$ des points de $E$ 
dont les coordonnées $(x,y,z)$ dans
$\mathcal{R}$ satisfont aux deux conditions suivantes :
\begin{enumerate}
\item Les nombres $x$, $y$ et $z$ sont dans $\Zz$ ;
\item Le nombre $x+y+z$ est divisible par $n$.
\end{enumerate}

On note $O$ l'origine de $\mathcal{R}$ et $A$, $B$, $C$ et $D$ les points 
de coordonnées dans $\mathcal{R}$ :
$$A=(n,0,0), \quad B=(0,n,0), \quad C=(0,0,n) \quad \text{ et } \quad 
D=(1,-1,0).$$
Soient $G$ le sous-groupe des isométries affines de $E$ qui 
conservent globalement l'ensemble $X$ et $G_0$ le
sous-groupe de $G$ formé des éléments qui fixent $O$.
\begin{enumerate}
\item Soient $P$ et $Q$ deux éléments de $X$. Montrer que la 
translation de vecteur $\overrightarrow{PQ}$
appartient à $G$. En déduire une caractérisation des translations qui 
appartiennent à $G$.
\item Quels sont les centres des symétries par rapport à un point 
qui appartiennent à $G$ ?
\item Décrire l'ensemble des éléments $P$ de $X$ tels que 
$\|\overrightarrow{OP}\|^2=2$.
\item On considère l'ensemble $\mathcal{T}$ des transformations affines 
qui laissent fixe $O$ et conservent globalement l'ensemble
$\lbrace A,B,C\rbrace$. Montrer que les éléments de $\mathcal{T}$ sont des
isométries et qu'ils conservent $X$. 
 Montrer que $\mathcal{T}$ est un sous-groupe de $G_0$ isomorphe au groupe de
 permutations $S_3$.
Expliciter la représentation matricielle de 
des transformations de $\mathcal{T}$.
\item Déduire des questions précédentes l'orbite de $D$ sous 
l'action de $G_0$. Montrer que tout élément de
$G_0$ laisse globalement invariant le plan $H$ d'équation $x+y+z=0$.
\item Expliciter une représentation matricielle de la symétrie 
orthogonale par rapport à $H$. Cette
symétrie appartient-elle à $G$ ?
\item On considère l'ensemble $G'_0$ des restrictions à $H$ des 
éléments de $G_0$. En étudiant l'action de $G'_0$ sur $H$ montrer que
l'ordre de $G'_0$ est 12.
Montrer que $G_0$ est d'ordre 12 ou 24 suivant la valeur de $n$.
\end{enumerate}
\finenonce{007479}
\finexercice


\exercice{7490, mourougane, 2021/08/10}
\enonce{007490}{}
Dans l'espace affine $E$ réel, on considère une homothétie $h$ de centre $\Omega$ et de rapport $1/2$ et une translation $t$ de vecteur $\vec{u}$.
Quelle est la nature de $t\circ h$ ?
\finenonce{007490}
\finexercice


\exercice{7491, mourougane, 2021/08/10}
\enonce{007491}{}
 Dans l'espace affine $E$ réel, muni d'un repère cartésien $(O, \vec{\imath},\vec{\jmath},\vec{k})$,
donner un système d'équations pour la droite passant par le point $A(1,2,3)$ et de vecteur directeur $\vec{u}=2\vec{\imath}-\vec{k}$.
\finenonce{007491}
\finexercice


\exercice{7492, mourougane, 2021/08/10}
\enonce{007492}{}
 Dans l'espace affine $E$ réel, muni d'un repère cartésien $(O, \vec{\imath},\vec{\jmath},\vec{k})$, soit $f$ une application affine dont la partie linéaire a pour matrice 
$$\begin{pmatrix}-1&0&0\cr
0&2&3\cr 0&2&3\end{pmatrix}.$$
Que peut-on dire de l'ensemble des points fixes de $f$ ?
\finenonce{007492}
\finexercice


\exercice{7493, mourougane, 2021/08/10}
\enonce{007493}{}
 Dans l'espace affine $E$ réel, muni d'un repère cartésien $(O, \vec{\imath},\vec{\jmath},\vec{k})$,
donner un système d'équations pour la droite passant par le point $A(3,2,1)$ et de vecteur directeur $\vec{u}=2\vec{\imath}-\vec{k}$.
\finenonce{007493}
\finexercice


\exercice{7494, mourougane, 2021/08/10}
\enonce{007494}{}
 Dans l'espace affine $E$ réel, muni d'un repère cartésien $(O, \vec{\imath},\vec{\jmath},\vec{k})$, on considère une application affine $f$ dont la partie linéaire a pour matrice 
$$\begin{pmatrix}-1&0&0\cr
0&1&2\cr 0&0&3
\end{pmatrix}.$$
Que peut-on dire de l'ensemble des points fixes de $f$ ?
\finenonce{007494}
\finexercice


\exercice{7499, mourougane, 2021/08/10}
\enonce{007499}{Questions de cours}
\begin{enumerate}
\item Dans un espace affine, quel est le sens de la notation $1/3\ A+1/3\ B+1/3\ C$
où $A,B,C$ sont trois points ?

\item Soit $A,B,C,D$ quatre points d'un plan affine. Soit $I$ le milieu de $[C,D]$.
L'isobarycentre des points $A,B,C,D$ est-il le centre de gravité du triangle $A,B,I$?

\item Une isométrie est-elle caractérisée par sa partie linéaire ?

\item Deux isométries d'un espace affine euclidien de dimension trois qui ont exactement
le même plan $P$ de points fixes sont-elles égales ?
\end{enumerate}
\finenonce{007499}
\finexercice


\exercice{7500, mourougane, 2021/08/10}
\enonce{007500}{Procédé d'orthonormalisation de Gram-Schmidt}
Dans l'espace vectoriel euclidien $\Rr^3$ muni du produit scalaire
standard  et  de  la  base  canonique $\mathcal{C}$,  appliquer  le  procédé
d'orthonormalisation de Gram-Schmidt à la base $\mathcal{B}$
$$v_1=\left( \begin{array}{c}  3\\1\\1\end{array}\right)\ ; \ 
v_2=\left(
\begin{array}{c}   2\\ 1\\ 0\end{array}\right)\   ;   \  
 v_3=\left(
\begin{array}{c} -1\\ -1\\ -1 \end{array}\right)
$$ pour obtenir une base orthonormée $\mathcal{B'}$.
\finenonce{007500}
\finexercice


\exercice{7509, mourougane, 2021/08/10}
\enonce{007509}{Une symétrie}
Soit $E$ un espace affine euclidien de dimension trois et $(O,\vec{\imath},\vec{\jmath}, \vec{k})$ un repère.
Soit $A(1,2,3)$ et $B(3,2,1)$ deux points de $E$.
\begin{enumerate}
\item
Déterminer l'expression analytique dans le repère $(O,\vec{\imath},\vec{\jmath}, \vec{k})$
de la symétrie ortho\-gonale par rapport au plan média\-teur 
du segment $[AB]$.
\item Vérifier votre résultat en déterminant l'image du point $A$.
\end{enumerate}
\finenonce{007509}
\finexercice


\exercice{7510, mourougane, 2021/08/10}
\enonce{007510}{Conique}
Dans l'espace affine $\mathcal{E}$ muni d'un repère orthonormé 
$(O,\vec{\imath},\vec{\jmath},\vec{k})$, on considère le cône $\mathcal{C}$ d'équation
$y^2+z^2=3(x-2)^2$.
\begin{enumerate}
\item Déterminer un plan dont l'intersection avec le cône $\mathcal{C}$ soit un cercle.

\item Déterminer la nature de l'intersection de $\mathcal{C}$ avec le plan d'équation $z=1$.
On précisera (s'ils existent) le centre, les axes de symétrie et les asymptotes .
\end{enumerate}
\finenonce{007510}
\finexercice


\exercice{7511, mourougane, 2021/08/10}
\enonce{007511}{}
\'Ecrire la table de toutes les isométries d'un espace affine euclidien de dimension $3$. 
\finenonce{007511}
\finexercice


\exercice{7515, mourougane, 2021/08/10}
\enonce{007515}{}
Démontrer que si $f$ est une application linéaire orthogonale 
d'un espace vectoriel euclidien $\vec{E}$ de dimension finie, alors
$Ker (f-Id_{\vec{E}})$ et $Im (f-Id_{\vec{E}})$ sont supplémentaires orthogonaux.
\finenonce{007515}
\finexercice


\exercice{7516, mourougane, 2021/08/10}
\enonce{007516}{}
 Soient $E$ un espace affine de dimension 3, et $A,B,C,D$ un
tétraèdre de $E$. Montrer que les droites joignant les milieux des cotés
opposés du tétraèdre sont concourantes.
\finenonce{007516}
\finexercice


\exercice{7517, mourougane, 2021/08/10}
\enonce{007517}{}
Soit $(\R_{ev}^3, standard)$ l'espace euclidien standard muni de la base canonique.

Déterminer une base orthonormée du sous-espace $vect(e_1, e_1+e_2+e_3)$.
La compléter en une base orthonormée de $\R^3$. 
\finenonce{007517}
\finexercice

\section{243.00 Conique}


\exercice{5540, rouget, 2010/07/15}
\enonce{005540}{*IT}
Le plan est rapporté à un repère orthonormé $\mathcal{R}=(0,\overrightarrow{i},\overrightarrow{j})$.
Eléménts caractéristiques de la conique dont une équation cartésienne dans $\mathcal{R}$ est

\begin{itemize}
\item  
       \begin{enumerate}
        \item  $y^2=x$,
        \item $y^2=-x$,
        \item $y=x^2$,
        \item $y=-x^2$.
       \end{enumerate}

\item
        \begin{enumerate}
        \item $\frac{x^2}{25}+\frac{y^2}{9}=1$,
        \item $\frac{x^2}{9}+\frac{y^2}{25}=1$,
        \item $x^2+2y^2=1$.
       \end{enumerate} 
\item
        \begin{enumerate}
        \item $\frac{x^2}{16}-\frac{y^2}{9}=1$,
        \item $-\frac{x^2}{16}+\frac{y^2}{9}=1$,
        \item $x^2-y^2=1$.
       \end{enumerate}
\end{itemize}
\finenonce{005540}


\finexercice


\exercice{5541, rouget, 2010/07/15}
\enonce{005541}{*IT}
Le plan est rapporté à un repère orthonormé $\mathcal{R}=\left(0,\overrightarrow{i},\overrightarrow{j}\right)$.
Eléménts caractéristiques de la courbe dont une équation dans $\mathcal{R}$ est

\begin{itemize}
\item  
  \begin{enumerate}
   \item $y=x^2+x+1$,
   \item $y^2+y-2x=0$,
   \item $y=\sqrt{2x+3}$.
  \end{enumerate}
\item
   \begin{enumerate}
   \item $x^2+x+2y^2+y=0$,
   \item $y=-2\sqrt{-x^2+x}$.
  \end{enumerate}
\item  $x^2-y^2+x+y+1=0$.
\end{itemize}
\finenonce{005541}


\finexercice


\exercice{5542, rouget, 2010/07/15}
\enonce{005542}{**IT}
Le plan est rapporté à un repère orthonormé $\mathcal{R}=(0,\overrightarrow{i},\overrightarrow{j})$.
Nature et éléments caractéristiques de la courbe dont une équation en repère orthonormé est
\begin{enumerate}
 \item $y=\frac{1}{x}$,
 \item $41x^2-24xy+34y^2-106x+92y+74=0$,
 \item $x^2+2xy+y^2+3x-2y+1=0$,
 \item $(x-y+1)^2+(x+y-1)^2=0$,
 \item $x^2+y^2-3x-y+3=0$,
 \item $x(x-1)+(y-2)(y-3)=0$,
 \item $(x+y+1)(x-y+3)=3$,
 \item $(2x+y-1)^2-3(x+y)=0$.
\end{enumerate}
\finenonce{005542}


\finexercice


\exercice{5543, rouget, 2010/07/15}
\enonce{005543}{*IT}
\label{exo:routhe4}
Etudier les courbes dont une équation polaire (en repère orthonormé direct) est

\begin{enumerate}
 \item $r=\frac{1}{1+2\cos\theta}$,
 \item $r=\frac{1}{1+\cos\theta}$,
 \item $r=\frac{1}{2+\cos\theta}$,
 \item $r=\frac{1}{1-\sin\theta}$,
 \item $r=\frac{1}{2-\cos\theta}$.
\end{enumerate}
\finenonce{005543}


\finexercice


\exercice{5544, rouget, 2010/07/15}
\enonce{005544}{***}
Déterminer l'image du cercle trigonométrique par la fonction
$\begin{array}[t]{cccc}f~:&\Cc&\rightarrow&\Cc\\
 &z&\mapsto&\frac{1}{1+z+z^2}
\end{array}$.
\finenonce{005544}


\finexercice


\exercice{5546, rouget, 2010/07/15}
\enonce{005546}{***}
\begin{enumerate}
 \item  \textbf{Droite de } \textsc{Simson.} Soit $(A,B,C)$ un triangle et $M$ un point du plan. Montrer que les
projetés orthogonaux $P$, $Q$ et $R$ de $M$ sur les cotés $(BC)$, $(CA)$ et $(AB)$ du triangle $(ABC)$ sont alignés si
et seulement si $M$ est sur le cercle circonscrit à $(ABC)$. La droite passant par $P$, $Q$ et $R$ s'appelle la droite
de \textsc{Simson} du point $M$ relativement au triangle $ABC$ (ou au cercle $(ABC)$).
 \item  \textbf{Parabole tangente aux trois côtés d'un triangle.} Lieu des foyers des paraboles tangentes à trois
droites deux à deux non parallèles. En particulier, fournir la construction des points de contacts.
\end{enumerate}
\finenonce{005546}


\finexercice


\exercice{5547, rouget, 2010/07/15}
\enonce{005547}{**}
$(\mathcal{C})$ est le cercle de diamètre $[A,B]$. $(D)$ est la tangente en $A$ à $(\mathcal{C})$. $P$ est
un point variable sur $(\mathcal{C})$ et $(T)$ la tangente en $P$ à $(\mathcal{C})$. $(T)$ recoupe $(D)$ en $S$. La perpendiculaire à $(AB)$
passant par $P$ coupe $(BS)$ en $M$. Ensemble des points $M$~?
\finenonce{005547}


\finexercice


\exercice{5815, rouget, 2010/10/16}
\enonce{005815}{**}
Etudier les courbes dont une équation en repère orthonormé est :

\begin{enumerate}
 \item  $2x^2+6xy+5y^2+4x+6y+1 = 0$.

\item  $x^2+2xy+y^2 +3x-2y+1 = 0$.

\item  $2x^2-4xy-3x+3y+1 = 0$.

\item  $-5x^2+6\sqrt{3}xy+y^2-4 = 0$.

\item  $4x^2+12xy+9y^2-2x+1 = 0$.

\item  $(x-y+1)^2+(x+y-1)^2 = 0$.

\item  $x^2+y^2-3x-y+2=0$.

\item  $x(x-1)+(y-2)(y-3) = 0$.

\item  $(x+2y-4)(x-y-1)=3$.

\item  $(2x+y-1)^2-3(x+y) = 0$.
\end{enumerate}
\finenonce{005815}


\finexercice


\exercice{5816, rouget, 2010/10/16}
\enonce{005816}{**}
Etudier les courbes dont une équation polaire (en repère orthonormé direct) est
\begin{enumerate}
\item  $r =\frac{2}{1-2\cos\theta}$,
\item  $r =\frac{6}{2+\cos\theta}$,
\item  $r =\frac{2}{1-\sin\theta}$.
\end{enumerate}
\finenonce{005816}


\finexercice


\exercice{5817, rouget, 2010/10/16}
\enonce{005817}{***}
\begin{enumerate}
 \item  Montrer que toute courbe de degré inférieur ou égal à $2$ admet une représentation paramétrique de la forme

\begin{center}
$\left\{
\begin{array}{l}
x(t)=\frac{P(t)}{R(t)}\\
y(t)=\frac{Q(t)}{R(t)}
\end{array}
\right.$
\end{center}

où $P$, $Q$ et $R$ sont des polynômes de degré inférieur ou égal à $2$ et montrer réciproquement que toute courbe paramétrée du type précédent est une courbe de degré inférieur ou égal à $2$.

\item  Etudier la courbe $\left\{
\begin{array}{l}
x=\frac{2t+1}{t^2+2t-1}\\
y=\frac{t^2-1}{t^2+2t-1}
\end{array}
\right.$.
\end{enumerate}
\finenonce{005817}


\finexercice

\section{243.01 Ellipse}


\exercice{2065, ridde, 1999/11/01}

\enonce{002065}{}
Soit $\mathcal{E}$ une ellipse de foyers $F$ et $F'$, $M$ un point fixé de
$\mathcal {E}$ et $M'$ un point qui se prom\`ene sur $\mathcal{E}$. Soient
$\mathcal{C}$ et  $\mathcal{C}'$ les cercles de centres $M$ et $M'$ de rayons
$MF'$ et $M'F'$. Soient $I$ le point de $(FM) \cap \mathcal{C}$ tel que $M \in
[FI]$ et $J$ le deuxi\`eme point d'intersection de $\mathcal{C}$ et  $\mathcal{C}'$.
\begin{enumerate}
\item Montrer que $ (MM')$ est bissectrice de l'angle $F'MJ$.
\item Que devient $J$ si $M'$ tend vers $M$ (on ne demande pas de preuve) ?
\item Montrer que la tangente \`a $\mathcal{E}$ en $M$ est bissectrice ext\'erieure
 de l'angle $FMF'$.
\end{enumerate}
\finenonce{002065}



\finexercice



\exercice{2071, ridde, 1999/11/01}

\enonce{002071}{}
Montrer que la courbe param\'etr\'ee $x (t) =\dfrac 1{ t^2 + t + 1}$ et $y (t) =
\dfrac t{ t^2 + t + 1}$ est une ellipse et la tracer.
\finenonce{002071}



\finexercice



\exercice{4912, quercia, 2010/03/17}
\enonce{004912}{Orthoptique d'une ellipse}

Soit ${\cal E}$ une ellipse de foyers $F,F'$, de centre $O$, de dimensions
$a$ et $b$.

Soient $M,M' \in {\cal E}$ tels que les tangentes à ${\cal E}$ sont
perpendiculaires en un point $T$.

Montrer que $TF^2 + TF'^2 = 4a^2$. Quel est le lieu de $T$ quand $M$ et $M'$
varient ?


\finenonce{004912}



\finexercice


\exercice{4913, quercia, 2010/03/17}
\enonce{004913}{Tangentes à une ellipse}

Soient ${\cal E} : \frac {x^2}{4a^2} + \frac {y^2}{4b^2} = 1$,
et ${\cal E}' : \frac {x^2}{a^2} + \frac {y^2}{b^2} = 1$.


\begin{enumerate}
  \item CNS sur $u,v,w$ pour que la droite d'équation $ux+vy+w=0$ soit
    tangente à ${\cal E}'$ ?
    

  \item Soient $(MP)$, $(MQ)$ deux tangentes à ${\cal E}'$ avec $M,P,Q \in {\cal E}$.
    Montrer que $(PQ)$ est aussi tangente à ${\cal E}$.
    
\end{enumerate}
\finenonce{004913}



\finexercice


\exercice{4914, quercia, 2010/03/17}
\enonce{004914}{Points mobiles avec $PQ =$ constante}

Soient $P$ un point mobile sur $Ox$, et $Q$ un point mobile sur $Oy$ tels
que $PQ$ reste constante.

\begin{enumerate}
  \item Pour $\alpha\in \R$, déterminer le lieu, ${\cal C}_\alpha$,
    de Bar$(P:1-\alpha,Q:\alpha)$.
    

  \item Soit $R$ le quatrième point du rectangle $OPQR$. Démontrer que la tangente à
    ${\cal C}_\alpha$ en un point $M$ est perpendiculaire à $(RM)$.
\end{enumerate}
\finenonce{004914}



\finexercice


\exercice{4915, quercia, 2010/03/17}
\enonce{004915}{$FMT$ est rectangle en $F$}

Soit ${\cal E}$ une ellipse de foyer $F$, directrice $D$.
Soit $M \in {\cal E}$ hors de l'axe focal, et $T$ le point d'intersection
de la tangente en $M$ et de la directrice $D$.
Montrer que $FMT$ est rectangle en $F$.

\finenonce{004915}



\finexercice


\exercice{4916, quercia, 2010/03/17}
\enonce{004916}{$1/OM^2 + 1/OP^2$}

Soit ${\cal E}$ un ellipse de centre $O$ et de dimensions $a$, $b$.
Soient $M,P \in {\cal E}$ tels que $OMP$ soit rectangle en $O$.

\begin{enumerate}
  \item Montrer que $\frac{1}{OM^2} + \frac1{OP^2} = \frac1{a^2} + \frac1{b^2}$.
  \item En déduire que $(MP)$ reste tangente à un cercle fixe de centre $O$.
\end{enumerate}
\finenonce{004916}



\finexercice


\exercice{4917, quercia, 2010/03/17}
\enonce{004917}{Cercle sur une tangente}

Soit ${\cal E}$ une ellipse de sommets $A,A'$, et $M \in {\cal E}$.
La tangente en $M$ coupe les tangentes en $A,A'$ en $P,P'$.
Montrer que le cercle de diamètre $[P,P']$ passe par les foyers de ${\cal E}$.
\finenonce{004917}



\finexercice


\exercice{5550, rouget, 2010/07/15}
\enonce{005550}{***}
Soit $P$ un polynôme de degré $3$ à coefficients réels. Montrer que la courbe d'équation
$P(x)=P(y)$ dans un certain repère orthonormé, est en général la réunion d'une droite et d'une ellipse d'excentricité
fixe.
\finenonce{005550}


\finexercice

\exercice{5820, rouget, 2010/10/16}
\enonce{005820}{***}
 $(\mathcal{C})$ est le cercle de diamètre $[AB]$. $(\mathcal{D})$ est la tangente en $A$ au cercle $(\mathcal{C})$. $P$ est un point variable sur le cercle $(\mathcal{C})$ et $(\mathcal{T})$ la tangente en $P$ au cercle $(\mathcal{C})$. $(\mathcal{T})$ recoupe $(\mathcal{D})$ en $S$. La perpendiculaire à la droite $(AB)$ passant par $P$ coupe la droite $(BS)$ en $M$. Ensemble des points $M$ ?
\finenonce{005820}


\finexercice


\exercice{5822, rouget, 2010/10/16}
\enonce{005822}{**}
Soit $P$ un polynôme de degré $3$ à coefficients réels. Montrer que la courbe d'équation $P(x) = P(y)$ dans un certain repère orthonormé, est en général la réunion d'une droite et d'une ellipse d'excentricité fixe.
\finenonce{005822}


\finexercice

\section{243.02 Parabole}


\exercice{2066, ridde, 1999/11/01}

\enonce{002066}{}
Soit $\mathcal{P}$ une parabole de foyer $F$,  de directrice $\mathcal{D}$,
$M$ un point de $\mathcal{P}$ et $H$ son projet\'e orthogonal sur $\mathcal{D}$.
Montrer que la tangente \`a $\mathcal{P}$ en $M$ est la m\'ediatrice de $[FH]$.
En d\'eduire un proc\'ed\'e de construction d'une parabole.
\finenonce{002066}



\finexercice



\exercice{2067, ridde, 1999/11/01}

\enonce{002067}{}
D\'eterminer l'ensemble des points d'où l'on peut mener deux tangentes
orthogonales \`a une parabole.
\finenonce{002067}



\finexercice



\exercice{2070, ridde, 1999/11/01}

\enonce{002070}{}
D\'eterminer astucieusement le sommet et l'axe de la parabole
$x (t) = t^2 + t + 1$ et $y (t) = t^2-2t + 2$.
\finenonce{002070}



\finexercice



\exercice{4903, quercia, 2010/03/17}
\enonce{004903}{Orthoptique d'une parabole}

Soit $P$ une parabole de foyer $F$ et de directrice $D$.
Soit $M \in P$, et $M'$ le point de $P$ tel que les tangentes en $M$ et $M'$
sont orthogonales.

\begin{enumerate}
  \item Montrer que ces tangentes se coupent au milieu de $[H,H']$.
    
  \item Montrer que $M,F,M'$ sont alignés.

En déduire dans un repère $(O,\vec i,\vec j)$ donné, toutes les paraboles
tangentes aux axes de coordonnées.
\end{enumerate}
\finenonce{004903}



\finexercice


\exercice{4904, quercia, 2010/03/17}
\enonce{004904}{Cercle circonscrit}

Soit ${\cal C}$ un cercle de centre $O$, et $A,B$ deux points distincts
de ${\cal C}$.
Soit $\Delta$ le diamètre parallèle à $(AB)$.

Pour $M \in {\cal C}$, on note $P,Q$ les intersections de $(MA)$ et $(MB)$ avec
$\Delta$.
Chercher le lieu du centre du cercle circonscrit à $MPQ$.

\finenonce{004904}



\finexercice


\exercice{4905, quercia, 2010/03/17}
\enonce{004905}{Projection sur le diamètre d'un cercle}

On donne un cercle ${\cal C}$ de centre $O$ et $A \in {\cal C}$.
Pour $M \in {\cal C}$, on construit le projeté $N$ sur le diamètre
perpendiculaire à $(OA)$, et $I$, le point d'intersection de $(OM)$ et $(AN)$.
Quel est le lieu de $I$ ?
\finenonce{004905}



\finexercice


\exercice{4906, quercia, 2010/03/17}
\enonce{004906}{$MF + MH = 2a$}

Soit $F$ un point, $D$ une droite ne passant pas par $F$, et
$a > \frac12d(F,D)$.

Trouver l'ensemble des points $M$ tels que $MF + d(M,D) = 2a$.


\finenonce{004906}



\finexercice


\exercice{4907, quercia, 2010/03/17}
\enonce{004907}{Paraboles passant par un point}

Soient $D$ une droite et $F \in D$.

Montrer que pour tout point $M \notin D$, il passe exactement deux paraboles
de foyer $F$ et d'axe $D$.

Montrer que les tangentes à ces paraboles en $M$ sont orthogonales.
\finenonce{004907}



\finexercice


\exercice{4908, quercia, 2010/03/17}
\enonce{004908}{Longueur minimale d'une corde normale, Ensi Physique 93}

Soit $\cal P$ une parabole de paramètre $p$ et $A \in \cal P$. Soit $B$ le
point où la normale à $\cal P$ en $A$ recoupe $\cal P$. Déterminer la longueur
minimale de $AB$.


\finenonce{004908}



\finexercice


\exercice{4909, quercia, 2010/03/17}
\enonce{004909}{Cordes perpendiculaires, Centrale P' 1996}

On considère une parabole dans le plan euclidien.

\begin{enumerate}
  \item Exprimer l'équation d'une droite passant par deux points $A$ et~$B$
    de la parabole à l'aide d'un déterminant d'ordre~3.


  \item $A,B,C$ étant trois points sur la parabole, exprimer le fait que
    $(AB)$ et $(AC)$ sont perpendiculaires.


  \item On fixe $A$ sur la parabole, $B$ et $C$ sont deux points de la parabole
    variables tels que $(AB)$ et $(AC)$ sont perpendiculaires.
    Montrer que $(BC)$ passe par un point fixe~$M$.


  \item Quel est le lieu de $M$ quand $A$ varie~?
\end{enumerate}
\finenonce{004909}



\finexercice


\exercice{4910, quercia, 2010/03/17}
\enonce{004910}{Normales concourantes, Centrale P' 1996}

Soit~${\cal P}$ la parabole d'équation $y^2 = 2px$ et $M_0=(x_0,y_0)\in{\cal P}$.

\begin{enumerate}
  \item Discuter l'existence et le nombre de points $M\in{\cal P}$ distincts de~$M_0$
    tels que la normale à~${\cal P}$ en~$M$ passe par~$M_0$.


  \item Dans le cas où il y a deux solutions, $M_1$ et $M_2$, trouver le lieu
    géométrique du centre de gravité du triangle $M_0M_1M_2$.

\end{enumerate}
\finenonce{004910}



\finexercice


\exercice{4911, quercia, 2010/03/17}
\enonce{004911}{Croisillons sur une parabole, Centrale MP 2000}

Pour $p>0$ on donne la courbe~$\Gamma$ d'équation $y^2=2px$.
Soit un carré $ABCD$ tel que $B,D\in \Gamma$ et $A,C$ appartiennent
à l'axe de symétrie de~$\Gamma$.

\begin{enumerate}
  \item Quelle relation lie les abscisses de $A$ et~$C$~?
    
  \item On construit une suite $(M_n)$ de points de~$Ox$, $M_n$ d'abscisse~$x_n$,
    telle que $x_{n+1}>x_n$ et $M_nM_{n+1}$ est la diagonale
    d'un carré dont les deux autres sommets appartiennent à~$\Gamma$.
    Déterminer un équivalent de~$x_n$ quand~$n\to\infty$.
    
\end{enumerate}
\finenonce{004911}



\finexercice


\exercice{5545, rouget, 2010/07/15}
\enonce{005545}{**}
Déterminer l'orthoptique d'une parabole , c'est-à-dire l'ensemble des points du plan par
lesquels il passe deux tangentes à la parabole, perpendiculaires l'une à l'autre.
\finenonce{005545}


\finexercice

\exercice{5548, rouget, 2010/07/15}
\enonce{005548}{***}
Soit, dans $\Rr^3$ rapporté à un repère orthonormé $(O,\overrightarrow{i},\overrightarrow{j},\overrightarrow{k})$, la courbe
$(\Gamma)$ d'équations $\left\{
\begin{array}{l}
y=x^2+x+1\\
x+y+z=1
\end{array}\right.$. Montrer que $(\Gamma)$ est une parabole dont on déterminera le sommet, l'axe, le foyer et la 
directrice.
\finenonce{005548}


\finexercice


\exercice{5552, rouget, 2010/07/15}
\enonce{005552}{***}
Equation cartésienne de la parabole tangente à $(0x)$ en $(1,0)$ et à $(0y)$ en $(0,2)$.
\finenonce{005552}


\finexercice


\exercice{5818, rouget, 2010/10/16}
\enonce{005818}{**}
Déterminer l'orthoptique d'une parabole, c'est-à-dire l'ensemble des points du plan par lesquels il passe deux tangentes à la parabole qui soient perpendiculaires.
\finenonce{005818}


\finexercice


\exercice{5819, rouget, 2010/10/16}
\enonce{005819}{***}
\begin{enumerate}
 \item  (Droite de \textsc{Simson}) Soient $ABC$ un triangle et $M$ un point du plan. 

Montrer que les projetés orthogonaux $P$, $Q$ et $R$ du point $M$ sur les côtés $(BC)$, $(CA)$ et $(AB)$ du triangle $ABC$ sont alignés si et seulement si $M$ est sur le cercle circonscrit au triangle $ABC$. La droite passant par les points $P$, $Q$ et $R$ s'appelle la droite de \textsc{Simson} du point $M$ relativement au cercle $(ABC)$.

\item  (Parabole tangente aux trois côtés d'un triangle) Lieu des foyers des paraboles tangentes à trois droites deux à deux non parallèles. Fournir en particulier la construction des points de contacts.
\end{enumerate}
\finenonce{005819}


\finexercice


\exercice{5821, rouget, 2010/10/16}
\enonce{005821}{**}
L'espace de dimension $3$ est rapporté à un repère orthonormé $(O,i,j,k)$. On note $(\Gamma)$ la courbe d'équations

\begin{center}
$\left\{
\begin{array}{l}
y=x^2+x+1\\
x+y+z-1=0
\end{array}
\right.$.
\end{center}

Montrer que $(\Gamma)$ est une parabole dont on déterminera le sommet, l'axe, le foyer et la directrice.

\finenonce{005821}


\finexercice


\exercice{5824, rouget, 2010/10/16}
\enonce{005824}{**}
Equation cartésienne de la parabole $(\mathcal{P})$ tangente à $(Ox)$ en $(1,0)$ et à $(Oy)$ en $(0,2)$.
\finenonce{005824}


\finexercice

\section{243.03 Hyperbole}


\exercice{4918, quercia, 2010/03/17}
\enonce{004918}{Projection non orthogonale}

Soient $F$ un point, $D$ une droite ne passant pas par $F$, et $\vec{\Delta}$
une direction ni égale ni perpendiculaire à $\vec{D}$.
Pour $M \in {\cal P}$, on note $H$ le projeté de $M$ sur $D$ parallèlement
à $\vec{\Delta}$.
Quel est l'ensemble des points $M$ tels que $MF = MH$ ?

\finenonce{004918}



\finexercice


\exercice{4919, quercia, 2010/03/17}
\enonce{004919}{Triangle rectangle sur une hyperbole}

Soit ${\cal H}$ une hyperbole équilatère de dimension $a$.
On se place dans un ROND $(O,\vec i,\vec j)$ construit sur les
asymptotes de ${\cal H}$.

\begin{enumerate}
  \item Déterminer l'équation de ${\cal H}$ dans ce repère.
    

  \item Soit $ABC$ un triangle rectangle en $A$ dont les trois sommets sont
    sur ${\cal H}$. Montrer que la tangente en $A$ est orthogonale à $(BC)$.
    

  \item Soit $ABC$ un triangle quelconque dont les sommets sont sur ${\cal H}$.
    Montrer que l'orthocentre y est aussi.
    
\end{enumerate}
\finenonce{004919}



\finexercice


\exercice{4920, quercia, 2010/03/17}
\enonce{004920}{Cercle sur une tangente}

Soit ${\cal H}$ une hyperbole de sommets $A,A'$, et $M \in {\cal H}$.
La tangente en $M$ coupe les tangentes en $A,A'$ en $P,P'$.
Montrer que le cercle de diamètre $[P,P']$ passe par les foyers de ${\cal H}$.
\finenonce{004920}



\finexercice


\exercice{4921, quercia, 2010/03/17}
\enonce{004921}{Triangle équilatéral}

Soient $A,F$ deux points distincts, $D$ leur médiatrice, ${\cal H}$
l'hyperbole de foyer $F$, directrice $D$, excentricité 2, et ${\cal C}$
un cercle passant par $A$ et $F$, de centre $I$.

\begin{enumerate}
  \item Pour $M \in {\cal C}$, montrer que $M \in {\cal H}
    \Leftrightarrow 3\overline{(\vec{IM},D)} \equiv \overline{(\vec{IF},D)}\ [2\pi]$.
    

  \item En déduire que si $I \notin (AF)$, ${\cal C}$ coupe ${\cal H}$ aux sommets
    d'un triangle équilatéral.
\end{enumerate}
\finenonce{004921}



\finexercice


\exercice{4922, quercia, 2010/03/17}
\enonce{004922}{$\overline{(\vec{OA},\vec{OM})} \equiv 2\overline{(\vec{AM},\vec{AO})}$}

Soient $O,A$ deux points distincts du plan.
Trouver les points $M$ tels que
$\overline{(\vec{OA},\vec{OM})} \equiv 2\overline{(\vec{AM},\vec{AO})}$.

\finenonce{004922}



\finexercice


\exercice{4923, quercia, 2010/03/17}
\enonce{004923}{Lieu géométrique}

Soient $A,A'$ deux points distincts et ${\cal C}$ le cercle de diamètre $[A,A']$.
Pour $P \in {\cal C}$, on construit : $P'$ le symétrique de $P$ par rapport à
$(AA')$, et $M$ le point d'intersection de $(AP)$ et $(A',P')$.
Quel est le lieu de $M$ ?
\finenonce{004923}



\finexercice


\exercice{4924, quercia, 2010/03/17}
\enonce{004924}{Triangle sur une hyperbole, Ensi P 91}

    Soit $\cal H$ une hyperbole équilatère et $ABC$ un triangle dont les
    sommets appartiennent à $\cal H$. Montrer que l'orthocentre, $H$, du triangle
    appartient aussi à $\cal H$. Comparer $H$ et le point $Q$ où le cercle
    circonscrit à $ABC$ recoupe $\cal H$.
    
\finenonce{004924}



\finexercice


\exercice{5549, rouget, 2010/07/15}
\enonce{005549}{*}
Que vaut l'excentricité de l'hyperbole équilatère (une hyperbole est équilatère si et seulement
si ses asymptotes sont perpendiculaires)~?
\finenonce{005549}


\finexercice


\exercice{5551, rouget, 2010/07/15}
\enonce{005551}{***}
Soit $(\mathcal{H})$ une hyperbole équilatère de centre $O$ et $P$ et $Q$ deux points de
$(\mathcal{H})$ symétriques par rapport à $O$. Montrer que le cercle de centre $P$ et de rayon $PQ$ recoupe
$(\mathcal{H})$ en trois points formant un triangle équilatéral de centre $P$.
\finenonce{005551}


\finexercice


\exercice{5823, rouget, 2010/10/16}
\enonce{005823}{***}
Soit $(\mathcal{H})$ une hyperbole équilatère de centre $O$ et $P$ et $Q$ deux points de $(\mathcal{H})$ symétriques par rapport à $O$. Montrer que le cercle de centre $P$ et de rayon $PQ$ recoupe $(\mathcal{H})$ en trois points formant un triangle équilatéral de centre $P$.
\finenonce{005823}


\finexercice

\section{243.04 Quadrique}


\exercice{4925, quercia, 2010/03/17}
\enonce{004925}{\'Etude d'équations}

Déterminer les natures des surfaces d'équation :

\begin{enumerate}
  \item $x^2 + y^2 + z^2 -2xy + 2xz + 3x - y + z + 1 = 0$.
  \item $(x-y)(y-z) + (y-z)(z-x) + (z-x)(x-y) + (x-y) = 0$.
  \item $x^2 + 9y^2 + 4z^2 - 6xy - 12yz + 4zx + 4 = 0$.
  \item $x^2 - 2y^2 - z^2 + 2xz - 4yz + 3 = 0$.
  \item $2x^2 + 2y^2 + z^2 + 2xz - 2yz + 4x - 2y - z + 3 = 0$.
  \item $xy + xz + yz + 1 = 0$.
  \item $2x^2 + 2y^2 - z^2 + 5xy - yz + xz = 0$.
  \item $xy + yz = 1$.
  \item $x^2 + 4y^2 + 5z^2 - 4xy - 2x + 4y = 0$.

On fera le minimum de calculs nécéssaires pour pouvoir conclure.
\end{enumerate}
\finenonce{004925}



\finexercice


\exercice{4926, quercia, 2010/03/17}
\enonce{004926}{Repère non orthonormé}

Soit ${\cal S}$ une surface d'équation
$ax^2 + by^2 + cz^2 + 2dxy + 2exz + 2fyz + 2gx + 2hy + 2iz + j = 0$
dans un repère non orthonormé. Montrer que c'est quand même une quadrique.
\finenonce{004926}



\finexercice


\exercice{4927, quercia, 2010/03/17}
\enonce{004927}{Centre de symétrie}

Soit ${\cal S}$ une quadrique d'équation $f(x,y,z) = 0$. On note $q$ la forme
quadratique associée à $f$.

\begin{enumerate}
  \item Montrer que, pour tout point $A$ et tout vecteur $\vec h$, on a :
    $f(A+\vec h) = f(A) + (\vec{\nabla f}(A)\mid \vec h) + q(\vec h)$.
  \item On suppose que ${\cal S}$ n'est pas incluse dans un plan.
    Montrer qu'un point $\Omega$ est centre de symétrie de ${\cal S}$
    si et seulement si $\vec{\nabla f}(\Omega) = \vec 0$.
  \item En déduire que si $0$ n'est pas valeur propre de la matrice de $q$,
    alors ${\cal S}$ admet un centre unique.
\end{enumerate}
\finenonce{004927}



\finexercice


\exercice{4928, quercia, 2010/03/17}
\enonce{004928}{Cône s'appuyant sur une ellipse}

Soit ${\cal E}$ l'ellipse d'équations :
$\begin{cases}\frac{x^2}{a^2} + \frac{y^2}{b^2} = 1 \cr
        z = 0 \cr\end{cases}$
et $\Omega = (x_0,y_0,z_0)$ avec $z_0 \ne 0$.
On note $\mathcal{C}$ le cône de sommet $\Omega$ engendré par ${\cal E}$.

\begin{enumerate}
  \item Chercher une équation cartésienne de $\mathcal{C}$.
    

  \item Quels sont les points $\Omega$ tels que $\mathcal{C} \cap Oyz$ soit un cercle ?
    
\end{enumerate}
\finenonce{004928}



\finexercice


\exercice{4929, quercia, 2010/03/17}
\enonce{004929}{Sections circulaires}


\begin{enumerate}
  \item On considère la forme quadratique $q(x,y,z) = ax^2 + by^2 + cz^2$
    avec $a \in {[b,c]}$.
  \begin{enumerate}
    \item Montrer qu'il existe $y,z \in \R$ tels que $y^2 + z^2 = 1$ et
        $by^2 + cz^2 = a$.
    \item En déduire qu'il existe une base orthonormée de $\R^3$ dans laquelle la
        matrice de $q$ est de la forme :
        $M = \begin{pmatrix} a &0 &* \cr 0 &a &* \cr * &* &* \cr\end{pmatrix}$.
  \end{enumerate}
  \item Soit ${\cal E}$ un ellipsoïde de centre $O$.
    Montrer qu'il existe un plan $P$ qui coupe ${\cal E}$ selon un cercle de
    centre $O$.
    Montrer que les sections de ${\cal E}$ par des plans parallèles à $P$ sont
    des cercles.

  \item Peut-on généraliser à une quadrique quelconque ?
\end{enumerate}
\finenonce{004929}



\finexercice


\exercice{4930, quercia, 2010/03/17}
\enonce{004930}{Rotation d'une droite}

\begin{enumerate}
  \item Soit $D$ la droite d'équations $\begin{cases} y = 1\cr x = \lambda z\cr\end{cases}$ où
    $\lambda$ est un réel non nul fixé.
    Déterminer une équation cartésienne et la nature de la surface ${\cal S}$
    engendrée par la rotation de $D$ autour de $Oz$.
    

  \item En déduire que tout hyperboloïde de révolution à une nappe est réunion
    d'une famille de droites ({\it surface réglée\/}).

  \item Généraliser à un hyperboloïde à une nappe quelconque.

\end{enumerate}
\finenonce{004930}



\finexercice


\exercice{4931, quercia, 2010/03/17}
\enonce{004931}{Droites sur un paraboloïde hyperbolique}

Soit ${\cal P}$ le paraboloïde d'équation $z= xy$.
Montrer que par tout point $M \in {\cal P}$, il passe deux droites et deux
seulement incluses dans ${\cal P}$.
\finenonce{004931}



\finexercice


\exercice{4932, quercia, 2010/03/17}
\enonce{004932}{Hyperbole en rotation}

Soit $\mathcal{C}$ la courbe d'équations :
$\begin{cases} x^2 - y^2 - 4x + 2 = 0 \cr x + z = 1.\cr\end{cases}$

\begin{enumerate}
  \item Déterminer la nature et les éléments remarquables de $\mathcal{C}$.
    

  \item Chercher une équation cartésienne de la surface ${\cal S}$ engendrée par la
    rotation de $\mathcal{C}$ autour de $Oz$ et reconnaître~${\cal S}$.
    

\end{enumerate}
\finenonce{004932}



\finexercice


\exercice{4933, quercia, 2010/03/17}
\enonce{004933}{Volume d'un ellipsoïde}

Soit ${\cal S}$ la surface d'équation
$x^2 + \frac{y^2}2 + \frac{3z^2}4 + xz = 1$.
Montrer que ${\cal S}$ est un ellipsoïde et en calculer le volume intérieur.
\finenonce{004933}



\finexercice


\exercice{4934, quercia, 2010/03/17}
\enonce{004934}{}\'Equation d'un cône

Déterminer les réels $\lambda$ tels que la surface d'équation :
$x(\lambda-y) + y(\lambda-z) + z(\lambda-x) = \lambda$ soit un cône.
Préciser alors le sommet et la nature du cône.
\finenonce{004934}



\finexercice


\exercice{4935, quercia, 2010/03/17}
\enonce{004935}{Plan tangent à un ellipsoïde}

Soit ${\cal E}$ un ellipsoïde d'équation
$\frac{x^2}{a^2} + \frac{y^2}{b^2} + \frac{z^2}{c^2} = 1$
et $P$ un plan d'équation $ux + vy + wz = 1$.
Montrer que $P$ est tangent à $E$ si et seulement si
$a^2u^2 + b^2v^2 + c^2w^2 = 1$.
\finenonce{004935}



\finexercice


\exercice{4936, quercia, 2010/03/17}
\enonce{004936}{Normale à un ellipsoïde}

Soit ${\cal E}$ un ellipsoïde d'équation
$\frac{x^2}{a^2} + \frac{y^2}{b^2} + \frac{z^2}{c^2} = 1$,
$M$ un point de ${\cal E}$, et $P,Q,R$ les intersections de la normale en $M$ à
${\cal E}$ avec les plans $Oyz$, $Oxz$, $Oxy$.
Montrer que $\overline{MP}$, $\overline{MQ}$, $\overline{MR}$ sont dans un
rapport constant (indépendant de $M$).
\finenonce{004936}



\finexercice


\exercice{4937, quercia, 2010/03/17}
\enonce{004937}{Points équidistants de deux droites}

Soient $D,D'$ deux droites non coplanaires et ${\cal S}$ l'ensemble des points
équidistants de $D$ et $D'$.
Montrer que ${\cal S}$ est un paraboloïde hyperbolique.
(Utiliser un repère judicieux)
\finenonce{004937}



\finexercice


\exercice{4938, quercia, 2010/03/17}
\enonce{004938}{$MF = eMH$}

On considère un point $F$, un plan $P$ ne passant pas par $F$ et un réel
$e > 0$. Montrer que l'ensemble, ${\cal S}$, des points $M$ tels que
$MF = ed(M,P)$ est une quadrique de révolution.
Préciser les différents cas possibles.
\finenonce{004938}



\finexercice


\exercice{4939, quercia, 2010/03/17}
\enonce{004939}{$MF = ed(M,D)$}

Dans l'espace, on considère un point $F$, une droite $D$ ne passant pas par $F$
et un réel $e > 0$. Montrer que l'ensemble, ${\cal S}$, des points $M$ tels que
$MF = ed(M,D)$ est une quadrique. Préciser les différents cas possibles.

\finenonce{004939}



\finexercice


\exercice{4940, quercia, 2010/03/17}
\enonce{004940}{$d(M,P)^2 + d(M,D)^2 =$ cste}

On considère un plan $P$ et une droite $D$ sécants. Déterminer le lieu des
points $M$ tels que $d(M,P)^2 + d(M,D)^2 = a$ (constante fixée).


\finenonce{004940}



\finexercice


\exercice{4941, quercia, 2010/03/17}
\enonce{004941}{Points équidistants d'un plan et d'une droite}

Dans l'espace, soit $P$ le plan d'équation $z=0$ et $D$ la droite d'équations :
$\begin{cases} y = 0\cr x\cos\theta - z\sin\theta = 0.\cr\end{cases}$
($0\le \theta < \frac\pi2$)

Quel est le lieu des points $M$ tels que $d(M,P) = d(M,D)$ ?
\finenonce{004941}



\finexercice


\exercice{4942, quercia, 2010/03/17}
\enonce{004942}{Sphères équidistantes d'une sphère et d'un plan}

Dans l'espace, on considère un plan $P$ et une sphère $S$.
Quel est le lieu des centres des sphères tangentes à $S$ et à $P$ ?
\finenonce{004942}



\finexercice


\exercice{4943, quercia, 2010/03/17}
\enonce{004943}{Appellations incontrolées}

La liste des quadriques semble comporter des oublis : paraboloïde parabolique,
cône hyperbolique,\dots
Dresser la liste de toutes les surfaces oubliées et constater qu'elles sont
connues sous d'autres appellations.
\finenonce{004943}



\finexercice


\exercice{5825, rouget, 2010/10/16}
\enonce{005825}{** I}
Nature et \og éléments caractéristiques \fg~de la quadrique $(\mathcal{S})$ dont une équation dans un repère orthonormé donné $\mathcal{R}=(O,i,j,k)$ de l'espace de dimension $3$ est :

\begin{enumerate}
 \item  $x^2+y^2+z^2-2yz-4x+4y-1=0$.

\item  $x^2+y^2+z^2+2xy-1=0$.

\item  $x^2+y^2+z^2-2xy+2xz+3x-y+z+1=0$.

\item  $x^2+4y^2+5z^2-4xy-2x+4y=0$.

\item  $x^2-4x-3y-2=0$.

\item  $7x^2-2y^2+4z^2+4xy+20xz+16yz-36x+72y-108z+36=0$.

\item  $(x-y)(y-z)+(y-z)(z-x)+(z-x)(x-y)+(x-y) = 0$.

\item  $xy+yz =1$.

\item  $xy+yz+zx+2y+1=0$.
\end{enumerate}
\finenonce{005825}


\finexercice


\exercice{5826, rouget, 2010/10/16}
\enonce{005826}{**}
Déterminer la quadrique contenant le point $A(2,3,2)$ et les deux paraboles $(\mathcal{P})$ d'équations $\left\{
\begin{array}{l}
z=0\\
y^2=2x
\end{array}
\right.$ et $(\mathcal{P}')$ d'équations  $\left\{
\begin{array}{l}
x=0\\
y^2=2z
\end{array}
\right.$.
\finenonce{005826}


\finexercice


\exercice{5827, rouget, 2010/10/16}
\enonce{005827}{***}
Démontrer que toute équation du second degré symétrique en $x$, $y$ et $z$ est l'équation d'une surface de révolution (une surface $(\mathcal{S})$ est dite de révolution d'axe $(\mathcal{D})$ si et seulement si $(\mathcal{S})$ est invariante par toute rotation d'axe $(\mathcal{D})$).

\finenonce{005827}


\finexercice


\exercice{5828, rouget, 2010/10/16}
\enonce{005828}{***}
Former l'équation de la surface de révolution $(\mathcal{S})$ engendrée par la rotation de la droite $(\mathcal{D})$ $\left\{
\begin{array}{l}
x=z+2\\
y=2z+1
\end{array}
\right.$ autour de la droite $(\Delta)$ d'équations $x = y = z$. Quelle surface obtient-on ?
\finenonce{005828}


\finexercice


\exercice{5829, rouget, 2010/10/16}
\enonce{005829}{***}
Equation du cône de sommet $S$ et de directrice $(\mathcal{C})$ dans les cas suivants :

\begin{enumerate}
 \item  $S(0,0,0)$ et $(\mathcal{C})$ : $x = t$, $y = t^2$, $z = t^3$, $t\in\Rr^*$.

\item  $S(1,-1,0)$ et $(\mathcal{C})$ : $\left\{
\begin{array}{l}
y+z=1\\
x^2+y^2=z
\end{array}
\right.$. 
\end{enumerate}
\finenonce{005829}


\finexercice


\exercice{5830, rouget, 2010/10/16}
\enonce{005830}{***}
Trouver une équation du cône de sommet $S$ circonscrit à la surface $(\mathcal{S})$ quand

\begin{enumerate}
 \item  $S(0,5,0)$ et $(\mathcal{S})$ : $x^2+y^2+z^2=9$,

\item  $S(0,0,0)$ et $(\mathcal{S})$ : $x^2+xy+z-1=0$.  (Préciser la courbe de contact.)

\end{enumerate}

(\textbf{Définitions.} Le cône $(\mathcal{C})$ de sommet $S$ circonscrit à la surface $(\mathcal{S})$ est la réunion des tangentes à $(\mathcal{S})$ passant par $S$. D'autre part, une droite est tangente à la surface $(\mathcal{S})$ en un point $M$ si et seulement si elle passe par $M$ et est contenue dans le plan tangent à $(\mathcal{S})$ en $M$).
\finenonce{005830}


\finexercice


\exercice{5831, rouget, 2010/10/16}
\enonce{005831}{***}
Pour quelles valeurs de $\lambda$ la surface $(\mathcal{S})$ d'équation $x(\lambda-y)+y(\lambda-z)+z(\lambda-x)-\lambda=0$ est-elle un cône du second degré ? En préciser alors le sommet et une directrice.
\finenonce{005831}


\finexercice


\exercice{5832, rouget, 2010/10/16}
\enonce{005832}{*}
Montrer que l'arc paramétré $\left\{
\begin{array}{l}
x= \frac{1}{2}e^t(\cos t-\sin t)\\
\rule[-4mm]{0mm}{10mm}y= \frac{1}{2}e^t(\cos t+\sin t)\\
z=e^t
\end{array}
\right.$  est tracé sur un cône du second degré de sommet $O$.
\finenonce{005832}


\finexercice


\exercice{5833, rouget, 2010/10/16}
\enonce{005833}{***}
Equation cartésienne du cylindre $(\mathcal{C})$ de direction $\overrightarrow{u}$ et de directrice $(C)$ dans les cas suivants :

\begin{enumerate}
 \item   $\overrightarrow{u}(1,0,1)$ et $(C)$ : $x = a\cos t$,  $y = b\sin t$,  $z =a\sin t\cos t$ ($a$ et $b$ tous deux non nuls).

\item   $\overrightarrow{u}(0,1,1)$ et $(C)$ :  $\left\{
\begin{array}{l}
y+z=1\\
x^2+y^2=1
\end{array}
\right.$.
\end{enumerate}
\finenonce{005833}


\finexercice


\exercice{5834, rouget, 2010/10/16}
\enonce{005834}{**}
Equation du cylindre $(\mathcal{C})$ de section droite la courbe $(C)$ d'équations $\left\{
\begin{array}{l}
z=x\\
2x^2+y^2=1
\end{array}
\right.$
\finenonce{005834}


\finexercice


\exercice{5835, rouget, 2010/10/16}
\enonce{005835}{** I}
Equation cartésienne du cylindre de révolution $(\mathcal{C})$ de rayon $R$ et d'axe $(\mathcal{D})$ d'équations $\left\{
\begin{array}{l}
x=z+2\\
y=z+1
\end{array}
\right.$. Déterminer $R$ pour que la droite $(Oz)$ soit tangente au cylindre.
\finenonce{005835}


\finexercice


\exercice{5836, rouget, 2010/10/16}
\enonce{005836}{*}
Trouver les plans tangents à l'ellipsoïde d'équation $x^2+2y^2+3z^2 = 21$ qui sont  parallèles au plan d'équation $x+4y+6z=0$.
\finenonce{005836}


\finexercice


\exercice{5837, rouget, 2010/10/16}
\enonce{005837}{**}
Trouver les plans tangents à la surface $(\mathcal{S})$ d'équation $x-8yz=0$ et contenant la droite $(\mathcal{D})$

d'équations $\left\{
\begin{array}{l}
y=1\\
x+4z+2=0
\end{array}
\right.$.

\finenonce{005837}


\finexercice


\exercice{5838, rouget, 2010/10/16}
\enonce{005838}{** I}
\begin{enumerate}
 \item  Equation du cylindre de révolution $(\mathcal{C})$ d'axe la droite d'équations $x=y+1=3z-6$ et de rayon $3$.

\item  Equation du cône de révolution $(\mathcal{C})$ d'axe la droite d'équations $x=y+1=3z-6$, de sommet $S(0,-1,2)$ et de demi-angle au sommet $ \frac{\pi}{3}$.
\end{enumerate}
\finenonce{005838}


\finexercice

\section{243.99 Autre}


\exercice{2068, ridde, 1999/11/01}

\enonce{002068}{}
Soit $\mathcal{E} = \left\{ M (z)/2\left| z\right|^2-\frac i2 (z^2-\bar{z}^2) = 1\right\}$,
$R$ la rotation de centre $O$ et d'angle $\frac{\pi}4$ et $\mathcal{E'} = R (
\mathcal{E})$. D\'eterminer une \'equation cart\'esienne de $\mathcal{E'}$ et en d\'eduire
le trac\'e de $\mathcal{E}$.
\finenonce{002068}



\finexercice



\exercice{2069, ridde, 1999/11/01}

\enonce{002069}{}
\begin{enumerate}
\item $13x^2 -32xy + 37y^2-2x + 14y-5 = 0$
\item $xy + 3x + 5y-4 = 0$
\item $ (2x + 3y)^2 + 4x + 6y-5 = 0$
\end{enumerate}
\finenonce{002069}



\finexercice



\exercice{4900, quercia, 2010/03/17}
\enonce{004900}{\'Equations du second degré}

Déterminer la nature et les éléments de la courbe d'équation dans un
repère $(O,\vec i,\vec j)$ orthonormé :

\begin{enumerate}
  \item $16x^2 -24xy + 9y^2 + 35x -20y = 0$.
    

  \item $5x^2 + 7y^2 + 2xy\sqrt3 -(10+2\sqrt3\,)x - (14+2\sqrt3\,)y - 4+2\sqrt3 = 0$.
    

  \item $x^2+xy+y^2 = 1$.
    


  \item $x^2+2y^2+4xy\sqrt3 + x + y\sqrt3 + 1 = 0$.

  \item $mx^2 + 4mx + (m-1)y^2 + 2 = 0$ ($m \in \R$).
    
\end{enumerate}
\finenonce{004900}



\finexercice


\exercice{4901, quercia, 2010/03/17}
\enonce{004901}{Courbe paramétrée}

Montrer que le support de la courbe paramétrée :
$\begin{cases} x = \cos t\cr y = \cos t+\sin t\cr\end{cases}$
est une ellipse, et en préciser les éléments.
\finenonce{004901}



\finexercice


\exercice{4902, quercia, 2010/03/17}
\enonce{004902}{Points alignés avec le foyer}

Soit $\cal C$ une conique de foyer $F$, directrice $D$, excentricité $e$.
On considère deux points de $\cal C$, $M\ne M'$ alignés avec $F$.
Montrer que les tangentes à $\cal C$ en $M$ et $M'$ se coupent sur $D$ ou sont
parallèles.

\finenonce{004902}



\finexercice
\section{244.01 Courbes paramétrées}


\exercice{2046, ridde, 1999/11/01}

\enonce{002046}{}
Tracer les courbes param\'etr\'ees suivantes\\

$$x (t) = \cos^{2} (t) \qquad y (t) = \cos^{3} (t)\sin (t)$$
$$x (t) = \frac{t}{1 + t^{4}} \qquad y (t) = \frac{t^{3}}{1 + t^{4}}$$
$$x (t) =t^{2} + \frac{2}{t}  \qquad y (t) = t + \frac{1}{t}$$
$$x (t) = \frac{1-t^{2}}{1 + t^{2}} \qquad y (t) = t\frac{1-t^{2}}{1 + t^{2}}$$
$$x (t) =\tan (t) + \sin (t)  \qquad y (t) = \frac{1}{\cos (t)}$$
$$x (t) = \sin (2t) \qquad y (t) = \sin (3t)$$
\finenonce{002046}



\finexercice



\exercice{2047, ridde, 1999/11/01}

\enonce{002047}{}
On fait rouler sans glissement un cercle de rayon $1$ sur l'axe $ (Ox)$.
D\'eterminer et tracer la courbe d\'ecrite par un point du cercle.
\finenonce{002047}



\finexercice



\exercice{2048, ridde, 1999/11/01}

\enonce{002048}{}
Tracer la courbe d'\'equation $x^{3} + y^{3} = 3xy$ en la coupant par les droites
$y = tx$ où $t \in \Rr$.
\finenonce{002048}



\finexercice



\exercice{2049, gourio, 2001/09/01}

\enonce{002049}{}
Tracer la courbe param\'{e}tr\'{e}e d\'{e}finie par :
$$x(t)=\int_{0}^{t}\cos (2u)\sin (u)du,\quad  y(t)=\int_{0}^{t}\sin (2u)\cos(u)du. $$
\finenonce{002049}



\finexercice



\exercice{2050, gourio, 2001/09/01}

\enonce{002050}{}
Tracer la courbe param\'{e}tr\'{e}e d\'{e}finie par :
$$x(t)=t^{2}+2t,  y(t)=\frac{1+2t}{t^{2}}. $$
\finenonce{002050}



\finexercice



\exercice{2699, matexo1, 2002/02/01}
\enonce{002699}{}
Dans le plan euclidien $\R^2$, donner des \'equations
param\'etriques pour une droite\,; un cercle\,; une ellipse\,; une
hyperbole\,; une parabole.
\finenonce{002699}
\finexercice


\exercice{2700, matexo1, 2002/02/01}
\enonce{002700}{}
Pour chacune des courbes suivantes, d\'eterminer la
tangente en tout point, les points d'inflexion et de
rebroussement, les branches infinies, les points doubles\,;
construire et tracer la courbe.

\begin{enumerate}
\item $\displaystyle{ x = \sin 4t, \quad y = \cos 3t}$\,;
\item $\displaystyle{ x = t - \sin t, \quad y = 1-\cos t}$\,;
\item $\displaystyle{ x = \cos^3 t, \quad y = \sin ^3 t}$\,;
\item $\displaystyle{ x = \cos t, \quad y = {\sin^2 t \over 2+\sin t}}$\,;
\item $\displaystyle{ x = {t^3 \over 1-t^2}, \quad y = {1+t\over
(1-t)^2}}$\,;
\item $\displaystyle{ x = \cos^3 t+\sin t, \quad y = \sin^3t+\cos t}$\,;
\item $\displaystyle{ x = 3 \cos t -2\sin^3 t, \quad y = \cos4t}$.
\end{enumerate}
\finenonce{002700}
\finexercice


\exercice{2701, matexo1, 2002/02/01}
\enonce{002701}{}

On consid\`ere l'ensemble $\Gamma$ des points du plan

$(x,y)$ qui v\'erifient $0<x<2$ et $x = 2\sin(y/x)$. 

Montrer

que c'est un arc dont on trouvera une repr\'esentation

param\'etrique. Construire $\Gamma$.

\finenonce{002701}



\finexercice



\exercice{2702, matexo1, 2002/02/01}
\enonce{002702}{}
 On consid\`ere l'arc param\'etr\'e du plan d\'efini par
$$ x= {t^2+1 \over t^3-1}, \quad y = {2t \over t^3-1}.$$
\'Etudier ses branches infinies. Trouver ses points d'inflexion
et montrer qu'ils sont align\'es. Tracer l'arc.
\finenonce{002702}
\finexercice


\exercice{2703, matexo1, 2002/02/01}
\enonce{002703}{}
Soit l'arc param\'etr\'e d\'efini par
$$ x = {t-\sin t \over t^2}, \quad y = {1-\cos t \over t^2}. $$
Montrer qu'il peut \^etre prolong\'e contin\^ument pour tout $t
\in\R$ et qu'il poss\`ede un axe de sym\'etrie. Montrer qu'il
poss\`ede une infinit\'e de points de rebroussement situ\'es
sur un m\^eme cercle, et que les tangentes en ces points sont
concourantes. Tracer l'arc.
\finenonce{002703}
\finexercice


\exercice{2704, matexo1, 2002/02/01}
\enonce{002704}{}
{\bf \'Epicyclo\"{\i}des, hypocyclo\"{\i}des - } \\ 
Soit $R >0$ un r\'eel, et $C_R$ le cercle 
de centre $O$ et de rayon $R$ dans le plan.
\begin{enumerate}
\item On consid\`ere le cercle $\gamma$ de rayon 1 tangent
ext\'erieurement \`a $C_R$ en $A = (R, 0)$\,; on fait rouler
$\gamma$ le long de $C_R$ sans glisser. Trouver des \'equations
param\'etriques de l'ensemble $\Gamma_R$ des points occup\'es
par $A$. Dans quels cas cet ensemble est-il un arc\,? Sinon quel
est-il\,?
\item On suppose \`a pr\'esent que $R >1$ et que $\gamma$
est tangent int\'erieurement \`a $C_R$\,; m\^emes questions sur
l'ensemble $\Gamma'_R$ ainsi construit.
\item Tracer $\Gamma_R$ et $\Gamma'_R$ pour $R = 6$ et $R =
8/3$.
\end{enumerate}
\finenonce{002704}
\finexercice


\exercice{2705, matexo1, 2002/02/01}
\enonce{002705}{}
Montrer que les deux droites de l'espace d'\'equations
param\'etriques ${x = 2+2t}, {y =2+4t}, {z=2-4t}$ et 
${x=4+t}, {y=6+2t}, {z=-2-2t}$ sont identiques.
\finenonce{002705}
\finexercice


\exercice{2706, matexo1, 2002/02/01}
\enonce{002706}{}
Montrer que la courbe de l'espace d'\'equations
param\'etriques
$$ x=4 \sqrt2 \cos t,\quad y = t+2\sin t,\quad z
=-2\cos t$$
est plane.
\finenonce{002706}
\finexercice


\exercice{2707, matexo1, 2002/02/01}
\enonce{002707}{}
 {\bf H\'elice - }  \\ \'Etudier la courbe param\'etr\'ee de l'espace d\'efinie
par 
$$ x = \cos t,\quad y = \sin t,\quad z = t. $$
Tracer ses projections orthogonales sur les trois plans $xOy$,
$yOz$, $xOz$. Montrer que la projection de cette courbe sur le
plan $xOy$ parall\`element \`a la direction d'une de ses
tangentes est une cyclo\"\i de.

\finenonce{002707}
\finexercice


\exercice{2708, matexo1, 2002/02/01}
\enonce{002708}{}
Trouver en tout point l'\'equation du plan osculateur \`a
la courbe 
$$ x = {t^3 \over t^2+1},\quad y={t^2 \over t^2+1},\quad z = {t \over t^2+1}. $$
\finenonce{002708}
\finexercice


\exercice{2709, matexo1, 2002/02/01}
\enonce{002709}{}
On consid\`ere l'arc de l'espace d\'efini en
coordonn\'ees param\'etriques par
$$ M(t):\qquad x = t^3,\quad y = t^2,\quad z = t. $$
D\'eterminer l'intersection $\mu(t)$ de sa tangente en $M(t)$
avec le plan osculateur en $O = M(0)$, et montrer que la
tangente \`a l'arc $t \mapsto \mu(t)$ n'est autre que
l'intersection des plans osculateurs en $O$ et en $M(t)$.
\finenonce{002709}
\finexercice


\exercice{2710, matexo1, 2002/02/01}
\enonce{002710}{}
{\bf Loxodromie de sph\`ere - } \\ On consid\`ere la courbe
param\'etrique de l'espace d\'efinie par
$$ x = \cos(k\log\sin t) \sin t,\quad
   y = \sin(k\log\sin t) \sin t,\quad
   z = \cos t, $$
pour $ 0 < t < \pi $, o\`u $k > 0$ est un r\'eel fix\'e.
\begin{enumerate}
\item Montrer qu'elle est trac\'ee sur une
sph\`ere de centre $O$, et qu'elle est sym\'etrique par
rapport \`a $O$. Montrer qu'elle poss\`ede deux points limites
que l'on pr\'ecisera.
\item Calculer sa tangente en tout point. Montrer qu'elle fait
un angle constant avec les m\'eridiens de la sph\`ere, angle que
l'on d\'eterminera en fonction de $k$. 
\item Tracer les projections de la courbe sur les trois
plans $xOy$, $yOz$ et $xOz$. Quelle est l'allure de cette
courbe dans l'espace\,?
\end{enumerate}
\finenonce{002710}
\finexercice


\exercice{2716, matexo1, 2002/02/01}
\enonce{002716}{}
\'Etudier la courbe d\'efinie par 
$$ \theta = t-2\sin t, \qquad \rho =\tan t $$
Trouver asymptotes et points doubles.
\finenonce{002716}
\finexercice


\exercice{2717, matexo1, 2002/02/01}
\enonce{002717}{}
Construire la courbe ayant pour \'equation implicite
$(x^2+y^2)^2-ax(x^2+2y^2) = 0, (a>0)$.

\finenonce{002717}
\finexercice


\exercice{4982, quercia, 2010/03/17}
\enonce{004982}{Rebroussements}

\'Etudier les points stationnaires des courbes paramétrées suivantes :

\begin{enumerate}
  \item $x=\sin t$,\quad             $y=\frac{\cos^2t}{2-\cos t}$. (Bicorne)
  \item $x=(1+\cos^2t)\sin t$,\quad  $y=\sin^2t\cos t$.
  \item $x = (1+\cos t)\sin2t$,\quad $y=\cos2t$.
  \item $x=2t^3+3t^2$,\quad          $y=3t^2+6t$.
  \item $x=t^3-3t$,\quad             $y=t^3-t^2-t+1$.
\end{enumerate}
\finenonce{004982}



\finexercice


\exercice{4983, quercia, 2010/03/17}
\enonce{004983}{Branches infinies}

\'Etudier les branches infinies des courbes paramétrées suivantes :

\begin{enumerate}
  \item $x=t^5-t^3+\frac t4$,\quad           $y=\frac{3t}{3t^2+1}$.
  \item $x=2\cos^2t+\ln|\sin t|$,\quad        $y=\sin2t$.
  \item $x=\sqrt{\frac{t^2-2}{t^4-1}}$,\quad $y=tx$.
    L'aire comprise entre la courbe et ses asymptotes est-elle finie ?
  \item $x=\frac{t^3-t}{2t-1}$,\quad         $y=tx$.
  \item $x=\frac1t+\frac1{t+1}$,\quad       $y=\frac1t+\frac1{(t+1)^2}$.
  \item $x=\frac1t+\frac1{t+1}$,\quad       $y=\frac1t-\frac1{t+1}$.
  \item $x=\frac{3t}{1+t^3}$,\quad           $y=tx$.
  \item $x=\frac{te^t}{t+1}$,\quad           $y=\frac{e^t}{t+1}$.
  \item $x=2t^3+3t^2$,\quad                   $y=3t^2+6t$.
  \item $x=t^3-3t$,\quad                      $y=t^3-t^2-t+1$.
  \item $x=\frac t{t^2-1}$,\quad             $y=\frac{t^2}{t-1}$.
\end{enumerate}
\finenonce{004983}



\finexercice


\exercice{4984, quercia, 2010/03/17}
\enonce{004984}{Inflexions}

Déterminer les points d'inflexion des courbes paramétrées suivantes :

\begin{enumerate}
  \item $x=\sin t$,\quad             $y=\frac{\cos^2t}{2-\cos t}$. (Bicorne)
  \item $x=\sin\frac t2$,\quad      $y=\tan t$.
  \item $x=\frac{e^t}t$,\quad       $y=te^t$.
    
  \item $x=\sin t\cos2t$,\quad       $y=\cos t\sin2t$.
    
\end{enumerate}
\finenonce{004984}



\finexercice


\exercice{4985, quercia, 2010/03/17}
\enonce{004985}{Matexo}

Soit $\mathcal{C}$ la courbe d'équations paramétriques~:
$x(t) = \frac{t^2+1}{t^3-1},
 y(t) = \frac{2t}{t^3-1}$.

\begin{enumerate}
  \item Montrer que les points de paramètres $t,u,v$ (distincts) sont alignés si et
    seulement si $tuv = t + u + v + 1$.
  \item Prouver que $\mathcal{C}$ admet exactement trois points d'inflexion et qu'ils
    sont alignés.
\end{enumerate}
\finenonce{004985}



\finexercice


\exercice{4986, quercia, 2010/03/17}
\enonce{004986}{Construction}

Construire la courbe d'équations paramétriques :
$x=\frac t{t^2-1}$, $y=\frac{t^2}{t-1}$.

Déterminer les coordonnées du point double et vérifier que les tangentes en ce
\finenonce{004986}



\finexercice


\exercice{4987, quercia, 2010/03/17}
\enonce{004987}{Construction}
Dessiner la courbe d'équation cartésienne : $x^3+y^3 = 3xy$
(folium de Descartes)
On prendra $t = \frac yx$ comme paramètre.
\finenonce{004987}



\finexercice


\exercice{4988, quercia, 2010/03/17}
\enonce{004988}{Construction}

Construire les courbes d'équation polaire :

\begin{enumerate}
  \item $\rho = \frac{\cos(\theta/2)}{1+\sin\theta}$.
  \item $\rho = \frac{\cos2\theta}{\cos\theta}$.
    (Strophoïde, calculer l'aire limitée par la boucle)
  \item $\rho = \frac{\sin\theta}{2\cos\theta-1}$.
    Vérifier que la courbe traverse ses asymptotes au point double.
  \item $\rho = \frac1{\cos\theta+\sin2\theta}$.
  \item $\rho = \cos\theta+ \frac1{\cos\theta}$.
  \item $\rho = \frac{\cos2\theta}{2\cos\theta-1}$.
  \item $\rho = \cos\frac\theta3$.
  \item $\rho = 1 + \sin3\theta$.
  \item $\rho = \frac1{\sqrt\theta}$.
  \item $\rho = \ln\theta$.
\end{enumerate}
\finenonce{004988}



\finexercice


\exercice{4989, quercia, 2010/03/17}
\enonce{004989}{Strophoïde}

Soit $\Gamma$ un cercle de centre $O$ et de rayon 1, $A \in \Gamma$, et $D$ le
diamètre de $\Gamma$ perpendiculaire à $(OA)$.

Pour $M \in \Gamma \setminus\{A\}$, on construit le point $N$ intersection de
$D$ et $(AM)$, puis le point $P$ tel que $\vec {AP} = \vec {MN}$.

\finenonce{004989}



\finexercice


\exercice{4990, quercia, 2010/03/17}
\enonce{004990}{Cochléoïde}

\begin{enumerate}
  \item Tracer la courbe $\mathcal{C}$ d'équation polaire $\rho = \frac{\sin\theta}\theta$
    (cochléoïde)

  \item Une droite passant par $O$ coupe $\mathcal{C}$ en un certain nombre de points.
    Montrer que les tangentes à $\mathcal{C}$ en ces points sont concourantes.
    
\end{enumerate}
\finenonce{004990}



\finexercice


\exercice{4991, quercia, 2010/03/17}
\enonce{004991}{Chimie P 91}

    Soient $O$ et $A$ deux points distincts dans un plan $\cal P$.
    Déterminer le lieu des points $M \in \cal P$ tels que
    $\overline{(\vec{OA},\vec{AM})} \equiv 3\overline{(\vec{OA},\vec{OM})}
    (\mathrm{mod}\,{\pi})$.
    
\finenonce{004991}



\finexercice


\exercice{4992, quercia, 2010/03/17}
\enonce{004992}{Ensi Chimie P' 93}
Déterminer les points doubles de la courbe d'équation polaire
$\rho = \frac{\theta}{\theta^2-1}$.
\finenonce{004992}



\finexercice


\exercice{5523, rouget, 2010/07/15}
\enonce{005523}{Quelques grands classiques}
\label{exo:routhe1}
\begin{enumerate}
\item (**) \textbf{L'astroïde.}
\begin{enumerate}
\item $a$ est un réel strictement positif donné. Etudier et construire la courbe de paramétrisation~:
$\left\{
\begin{array}{l}
x=a\cos^3t\\
y=a\sin^3t
\end{array}
\right.$.
\item Pour $t\in]0,\frac{\pi}{2}[$, on note $A(t)$ et $B(t)$ les points d'intersection de la tangente au point
courant $M(t)$ avec respectivement $(Ox)$ et $(Oy)$. Calculer la longueur $A(t)B(t)$.
\end{enumerate}

\item (**) \textbf{La cycloïde.}

\begin{enumerate}
\item Un cercle $(\mathcal{C})$, de rayon $R>0$, roule sans glisser sur l'axe $(Ox)$. On note $I$ le point de
contact entre $(\mathcal{C})$ et $(Ox)$ et on note $\Omega$ le centre de $(\mathcal{C})$ ($\Omega$ et $I$ sont mobiles).
$M$ est un point donné de $(\mathcal{C})$ ($M$ est mobile, mais solidaire de $(\mathcal{C})$). On pose
$t=(\widehat{(\overrightarrow{\Omega M},\overrightarrow{\Omega I})}$.

$$\includegraphics{../images/img005523-3}$$


Déterminer une paramétrisation de la courbe décrite par le point $M$ (on prendra $t$ pour paramètre).

\item Etudier et construire l'arc paramétré~:~$\left\{
\begin{array}{l}
x=R(t-\sin t)\\
y=R(1-\cos t)
\end{array}
\right.$ où $R$ est un réel strictement positif donné.
\end{enumerate}
\item (**) \textbf{Une courbe de} \textsc{Lissajous}. Etudier et construire l'arc
paramétré~:~$\left\{\begin{array}{l}x=\sin(2t)\\
y=\sin(3t)
\end{array}\right.$
\item (**) \textbf{La lemniscate de} \textsc{Bernoulli}. Etudier et construire l'arc
paramétré~:~$\left\{\begin{array}{l}x=\frac{t}{1+t^4}\\
y=\frac{t^3}{1+t^4}
\end{array}\right.$
\item (***) \textbf{Les tractrices.}
\begin{enumerate} 
\item Trouver les trajectoires orthogonales à la famille des cercles de rayon $R$ ($R>0$ donné) et centrés sur $(Ox)$.
\item Etudier et construire l'arc paramétré~:~$\left\{
\begin{array}{l}
x=R(\ln|\tan\frac{t}{2}|+\cos t)\\
y=R\sin t
\end{array}
\right.$ où $R$ est un réel strictement positif donné.
\end{enumerate}
\end{enumerate}
\finenonce{005523}


\finexercice


\exercice{5524, rouget, 2010/07/15}
\enonce{005524}{}
Construire les courbes de paramétrisations :
\begin{enumerate}
 \item $\left\{\begin{array}{l}
x=\frac{t^3}{(t+1)^2(t-1)}\\
y=\frac{t^2}{t^2-1}
\end{array}\right.$

 \item $\left\{\begin{array}{l}
x=(t+2)e^{1/t}\\
y=(t-2)e^{1/t}
\end{array}\right.$

 \item $\left\{\begin{array}{l}
x=(t-1)\ln(|t|)\\
y=(t+1)\ln(|t|)
\end{array}\right.$

 \item $\left\{\begin{array}{l}
x=\frac{2t}{1+t^2}\\
y=\frac{t+2}{1-t^2}
\end{array}\right.$

 \item $\left\{\begin{array}{l}
x=\frac{t}{t^2-1}\\
y=\frac{t+2}{(t-1)^2}
\end{array}\right.$

 \item $\left\{\begin{array}{l}
x=\frac{t^3}{t^2-9}\\
y=\frac{t(t-2)}{t-3}
\end{array}\right.$

 \item $\left\{\begin{array}{l}
x=\frac{t^3}{1+3t}\\
y=\frac{3t^2}{1+3t}
\end{array}\right.$

 \item $\left\{\begin{array}{l}
x=t^2+t^3\\
y=t^2+t^3-2t^4-2t^5
\end{array}\right.$
\end{enumerate}
\finenonce{005524}


\finexercice


\exercice{5525, rouget, 2010/07/15}
\enonce{005525}{}
La courbe orthoptique d'une courbe $(\mathcal{C})$ est le lieu  des points du plan d'où l'on peut mener (au moins) deux
tangentes à $(\mathcal{C})$, orthogonales. Déterminer l'orthoptique de $(\mathcal{C})$ dans chacun des cas suivants~:

\begin{enumerate}
\item  $(\mathcal{C})$ est un astroïde de paramétrisation $\left\{\begin{array}{l}
x=a\cos^3t\\
y=a\sin^3t
\end{array}\right.$, $a>0$ donné.

\item  $(\mathcal{C})$ est l'arc paramétré~:~$\left\{
\begin{array}{l}
x=t^2-2t\\
y=2t^3-3t^2
\end{array}
\right.$.

\item  $(\mathcal{C})$ est l'ellipse d'équation $\frac{x^2}{a^2}+\frac{y^2}{b^2}=1$, $(a,b)\in]0,+\infty[^2$.
\end{enumerate}
\finenonce{005525}


\finexercice


\exercice{5526, rouget, 2010/07/15}
\enonce{005526}{}
Trouver les droites à la fois tangentes et normales à l'arc paramétré~:~$\left\{\begin{array}{l}
x=3t^2\\
y=4t^3
\end{array}\right.$
\finenonce{005526}



\finexercice


\exercice{5527, rouget, 2010/07/15}
\enonce{005527}{}
Dans chacun des cas suivants, trouver une paramétrisation rationnelle de la courbe proposée puis construire

$$\begin{array}{cc}
1)\;x(y^2-x^2)=2y^2-x^2&2)\;x^3-y^3+xy-2x+2y+3=0
\end{array}$$
\finenonce{005527}



\finexercice


\exercice{5528, rouget, 2010/07/15}
\enonce{005528}{}
Trouver une équation cartésienne des supports des arcs suivants~:

\begin{enumerate}
 \item $\left\{\begin{array}{l}
x=t^2\\
y=-t^2
\end{array}\right.$
 \item $\left\{\begin{array}{l}
x=t^2\\
y=t^3
\end{array}\right.$
  \item $\left\{\begin{array}{l}
x=\frac{t}{1+t^4}\\
y=\frac{t^3}{1+t^4}
\end{array}\right.$
\end{enumerate}
\finenonce{005528}



\finexercice

\exercice{6981, blanc-centi, 2015/07/04}
\video{Lq1AdnlleAw}
\enonce{006981}{}
Représenter les courbes d'équation cartésienne $y=f(x)$, donner l'équation de 
leur tangente au point d'abscisse $x=0$ et la position de la courbe par rapport à cette tangente, pour :
\begin{enumerate}
\item $f(x)=\sin^2x+\cos x$
\item $f(x)=x+\ln(1+e^x)$
\end{enumerate}
\finenonce{006981}
 

\finexercice


\exercice{6982, blanc-centi, 2015/07/04}
\video{tNHUCWICKS8}
\enonce{006982}{}\
\begin{enumerate}
\item Donner une paramétrisation $(x(t),y(t))$ de la courbe d'équation $$y=\sqrt{-x^2-3x+4}$$ en précisant le domaine de variation du paramètre $t$.
\item Montrer que le support de la courbe paramétrée par
$$\left\{\begin{array}{l}x(t)=\cos t +3\\ y(t)=\sin t\end{array}\quad(t\in\R)\right.$$
ne peut pas \^etre décrit par une équation de la forme $y=f(x)$.
\item Montrer que le support de la courbe paramétrée par
$$\left\{\begin{array}{l}x(t)=\cos^2t-2\\ y(t)=\sin^4t+4\sin^2t+4\end{array}\quad(t\in\R)\right.$$
est le graphe d'une fonction $f$ que l'on précisera, ainsi que son domaine de définition.
\end{enumerate}
\finenonce{006982}
 

\finexercice


\exercice{6983, blanc-centi, 2015/07/04}
\video{cue85AzsR6Q}
\enonce{006983}{}
\'Etudier et tracer les courbes paramétrées suivantes:
\begin{enumerate}
\item $\displaystyle\left\{\begin{array}{l}x(t)=\cos^3t\\ y(t)=\sin^3t\end{array}\right.$ {\it (L'astro\"ide)}
\item $\displaystyle\left\{\begin{array}{l}x(t)=t-\tanh t\\ y(t)=\frac{1}{\ch t}\end{array}\right.$
\item $\displaystyle\left\{\begin{array}{l}x(t)=t-\sin t\\ y(t)=1-\cos t\end{array}\right.$ {\it(La cyclo\"ide)}
\end{enumerate}
\finenonce{006983}
 

\finexercice


\exercice{6984, blanc-centi, 2015/07/04}
\video{Eq0ELGOAv88}
\enonce{006984}{}
Soit $\mathcal{C}$ la courbe plane paramétrée par 
$$\left\{\begin{array}{l}x(t)=t\ln t\\ y(t)=\frac{\ln t}{t}\end{array}\right.\ (t\in]0;+\infty[)$$
\begin{enumerate}
\item Comparer les points de paramètres $t$ et $1/t$, en déduire un domaine d'étude de $\mathcal{C}$.
\item Représenter $\mathcal{C}$.
\end{enumerate}
\finenonce{006984}
 

\finexercice


\exercice{6985, blanc-centi, 2015/07/04}
\video{DO8B1OcGOTE}
\enonce{006985}{}
Montrer que la courbe paramétrée 
$$\left\{\begin{array}{l}x(t)=\frac{1}{t^2-t}\\ \ \\ y(t)=\frac{t}{t^2-1}\end{array}\right.$$
possède un point double et que les tangentes en ce point sont perpendiculaires.
\finenonce{006985}
 

\finexercice


\exercice{6986, blanc-centi, 2015/07/04}
\video{-Jx6iFJgAgU}
\enonce{006986}{}
Montrer que la courbe paramétrée 
$$\left\{\begin{array}{l}x(t)=\frac{4t-3}{t^2+1}\\ \ \\ y(t)=\frac{2t-1}{t^2+2}\end{array}\right.$$
admet un unique point singulier, et tracer l'allure de la courbe au voisinage de ce point.
\finenonce{006986}
 

\finexercice


\exercice{6987, blanc-centi, 2015/07/04}
\video{xJr2AjXWZto}
\enonce{006987}{}
On considère la courbe paramétrée définie par 
$$\left\{\begin{array}{l}x(t)=t+\frac{4}{t}\\ \ \\ y(t)=\frac{t}{3}+2+\frac{3}{t+1}\end{array}\right.$$
\begin{enumerate}
\item Dresser le tableau de variations conjointes de $x$ et $y$.
\item Calculer les tangentes horizontales, verticales et les asymptotes.
\item Trouver le point singulier de la courbe, étudier son type et 
écrire l'équation de la tangente à la courbe en ces points.
\item Tracer la courbe.
\end{enumerate}
\finenonce{006987}
 

\finexercice


\exercice{6988, blanc-centi, 2015/07/04}
\video{nXE8OreNJv8}
\enonce{006988}{}
Trouver les droites à la fois tangentes et orthogonales à la courbe 
$$\left\{\begin{array}{l}x(t)=3t^2\\ y(t)=4t^3\end{array}\right.$$
\finenonce{006988}
 

\finexercice

\section{244.02 Coordonnées polaires}


\exercice{2051, ridde, 1999/11/01}

\enonce{002051}{}
Tracer les courbes en polaires suivantes\\

$$\rho (\theta) = \sin (2\theta)$$
$$\rho (\theta) = \frac{\sin (\theta)}{\theta}$$
$$\rho (\theta) = \frac{\theta-1}{\theta + 1}$$
$$\rho (\theta) = \cos (\theta)-\cos (2\theta)$$
$$\rho (\theta) = \frac{\cos (\theta)}{1 + \sin (\theta)}$$
\finenonce{002051}



\finexercice



\exercice{2052, ridde, 1999/11/01}

\enonce{002052}{}
Soit $C$ un cercle du plan de centre $ (1, 0)$ et de rayon $a$. D\'eterminer
et tracer le lieu des projet\'es orthogonaux de $O$ sur les tangentes de $C$.
\finenonce{002052}



\finexercice



\exercice{2053, ridde, 1999/11/01}

\enonce{002053}{}
D\'eterminer et tracer les courbes dont la tangente en tout point $M$ fait un
angle de $\dfrac{\pi}{4}$ avec $\overrightarrow{OM}$.
\finenonce{002053}



\finexercice



\exercice{2054, ridde, 1999/11/01}

\enonce{002054}{}
Grâce aux coordonn\'ees polaires, tracer la courbe d\'efinie implicitement par la relation
$2xy (x^{2} + y^{2}) = x^{2}-y^{2}$.
\finenonce{002054}



\finexercice



\exercice{2055, gourio, 2001/09/01}

\enonce{002055}{}
Tracer la courbe d'\'{e}quation polaire :
$$r=1+\cos \theta .$$
\finenonce{002055}



\finexercice



\exercice{2056, gourio, 2001/09/01}

\enonce{002056}{}
Tracer les courbes d'\'{e}quations polaires :
$$r=\frac{\tan \theta }{\cos \theta } \ ;\
r^{2}=\frac{1}{\sin (2\theta )} .$$
\finenonce{002056}



\finexercice



\exercice{2711, matexo1, 2002/02/01}
\enonce{002711}{}
\begin{enumerate}
\item 
Montrer qu'un cercle $C$ de
diam\`etre $a$ et passant par le p\^ole $O$ peut \^etre
repr\'esent\'e en coordonn\'ees polaires par l'\'equation $\rho = a
\cos(\theta-\theta_0)$.
On consid\`ere un r\'eel $b > 0$ et la concho\"\i
de de $C$ de valeur $b$, c'est-\`a-dire la courbe
$\Gamma$ d\'efinie comme suit: \`a tout point $P$ de $C$, on
associe le point $M$ situ\'e sur la demi-droite $OP$, du
c\^ot\'e oppos\'e \`a $O$ par rapport \`a $P$ et tel que
$PM = b$; $\Gamma$ est le lieu des points $M$. Donner une
\'equation polaire de $\Gamma$. Construire $\Gamma$ en distinguant
quatre cas: $a>b$, $a=b$, $a<b<2a$ et $b \ge 2a$. D\'eterminer en
particulier les points d'inflexion dans chaque cas.

\item 
En s'inspirant de la question pr\'ec\'edente, tracer les
concho\"\i des d'une droite $\rho = {a / \cos(\theta-\theta_0)}$.
\end{enumerate}
\finenonce{002711}
\finexercice


\exercice{2712, matexo1, 2002/02/01}
\enonce{002712}{}
\'Etudier en fonction des param\`etres $a, b >0$ les courbes
d'\'equation $\rho = a/(1+b\cos\theta)$ (en particulier branches
infinies, position par rapport aux asymptotes). Montrer que ce
sont des coniques, et en d\'eterminer les foyers.
\finenonce{002712}
\finexercice


\exercice{2713, matexo1, 2002/02/01}
\enonce{002713}{}

\'Etudier et tracer les courbes d\'efinies en coordonn\'ees

polaires ci-apr\`es; s'il y a des branches infinies, les

pr\'eciser, et pr\'eciser la position de la courbe par rapport

aux \'eventuelles asymptotes; trouver aussi les points doubles:


$$
\begin{array}{cccc}
\text{rosace \`a quatre branches} &\rho &=& a \sin2\theta\cr
&&&\cr
& \rho &=&\sin\theta+\cos{\theta\over2} \cr
&&&\cr
\text{stropho\"\i de droite} & \rho &=& a{\cos2\theta \over \cos\theta} \cr
&&&\cr
& \rho &=& 1+2\cos{3\theta\over2}\cr
&&&\cr
\text{scarab\'ee} & \rho &=& 5\cos2\theta -3\cos\theta \cr
&&&\cr
\text{courbe du diable} & \rho^2 &=& 49 + {1\over \cos2\theta} \cr
&&&\cr
\text{spirale d'Archim\`ede} & \rho &=& a \theta \cr
&&&\cr
& \rho &=& \theta + 1/\theta  \text{ (asymptote ?)} \cr
&&&\cr
\text{spirale parabolique} & \theta &=& (\rho-1)^2 \cr
&&&\cr
\text{cochl\'eo\"\i de} & \rho &=& a {\sin\theta \over \theta} \cr
&&&\cr
\text{courbe du spiral} & \rho &=& {a \over 1+e^{\theta/5}}\cr
&&&\cr
& \rho&=& {1-2\cos\theta \over 1+\sin\theta}  \text{ (parabole asymptote)}\cr
&&&\cr
\text{\'epi} & \rho &=& {a \over \sin(5\theta/3)} \cr
\end{array}
$$


\finenonce{002713}



\finexercice



\exercice{2714, matexo1, 2002/02/01}
\enonce{002714}{}
Une {\em spirale logarithmique} est une courbe
d'\'equation en coordonn\'ees polaires $\rho = a e^{k\theta}$. Montrer
qu'elle coupe ses rayons vecteurs suivant un angle constant,
qu'on d\'eterminera en fonction de $k$.
\finenonce{002714}
\finexercice


\exercice{2715, matexo1, 2002/02/01}
\enonce{002715}{}
Montrer que la courbe d\'efinie par
$$\rho = 1+\tan{\theta\over2}$$ 
admet une asymptote; pr\'eciser la position de la courbe
par rapport \`a elle.
\finenonce{002715}
\finexercice


\exercice{4993, quercia, 2010/03/17}
\enonce{004993}{Courbes en polaires}
Construire les courbes en polaires suivantes :
\begin{enumerate}
\item  $\rho=\dfrac{\cos\theta/2}{1+\sin\theta}$                     \par
% \vtop to 7cm{\mapleplot{% 				img14
% r:=cos(t/2)/(1+sin(t));
% print(plot([r,t,t=0..4*Pi],view=[-2.8..2.8,-2..2],coords=polar));}   \par
% \vss}


\item $\rho=\dfrac{\cos2\theta}\cos\theta$                          \par
% \vtop to 7cm{\mapleplot{% 				img15
% r:=cos(2*t)/cos(t);
% print(plot([r,t,t=-Pi/2..Pi/2],view=[-2..2,-2..2],coords=polar));}   \par
% \vss}


\item $\rho=\dfrac{\sin\theta}{2\cos\theta-1}$                      \par
% \vtop to 7cm{\mapleplot{%				img16
% r:=sin(t)/(2*cos(t)-1);
% print(plot([r,t,t=-Pi..Pi],view=[-2..2,-3..1],coords=polar));}       \par
% \vss}


\item $\rho=\dfrac1{\cos\theta+\sin2\theta}$                        \par
% \vtop to 7cm{\mapleplot{%				img17
% r:=1/(cos(t)+sin(2*t));
% print(plot([r,t,t=-Pi..Pi],view=[-4..4,-3..3],coords=polar));}       \par
% \vss}


\item  $\rho=\cos\theta+\dfrac1\cos\theta$                           \par
% \vtop to 7cm{\mapleplot{%				img18
% r:=cos(t)+1/cos(t);
% print(plot([r,t,t=-Pi/2..Pi/2],view=[0..3,-5..5],coords=polar));}    \par
%
% \vss}


\item $\rho=\dfrac{\cos2\theta}{2\cos\theta-1}$                     \par

% \vtop to 7cm{\mapleplot{%				img19
% r:=cos(2*t)/(2*cos(t)-1);
% print(plot([r,t,t=-Pi..Pi],view=[-2..2,-2..2],coords=polar));}       \par
% \vss}
% 


\item  $\rho=\cos\dfrac\theta3$                                                  \par
% \vtop to 7cm{\mapleplot{%				img20
% r:=cos(t/3);
% print(plot([r,t,t=0..6*Pi],scaling=constrained,coords=polar,axes=frame));}\par
% \vss}

\item  $\rho=1+\sin3\theta$                                                      \par
% \vtop to 7cm{\mapleplot{%				img21
% r:=1+sin(3*t);
% print(plot([r,t,t=0..2*Pi],scaling=constrained,coords=polar,axes=frame));}\par
% \vss}

\item $\rho=\dfrac1{\sqrt\theta}$                                   \par
% \vtop to 7cm{\mapleplot{%				img22
% r:=1/sqrt(t);
% print(plot([r,t,t=0..6*Pi],view=[-2..2,-1.4..1.4],coords=polar));}   \par
% \vss}

\item $\rho=\ln\theta$                                              \par
% \vtop to 7cm{\mapleplot{%				img23
% r:=ln(t);
% print(plot([r,t,t=0..20],view=[-4..4,-3..3],coords=polar));}         \par
% \vss}

\item $\rho=\frac{\sin\theta}\theta$                                               \par
% \vtop to 7cm{\mapleplot{%				img24
% r:=sin(t)/t;
% print(plot([r,t,t=-6*Pi..6*Pi],scaling=constrained,coords=polar,axes=frame));}\par
% \vss}
% 
% 
\end{enumerate}
\finenonce{004993}
 
 
\finexercice


\exercice{5205, rouget, 2010/06/30}
\enonce{005205}{***}
Construire l'ensemble des points $M$ de coordonnées polaires $(r,\theta)\in\Rr^2$ vérifiant 

$r=\frac{1}{\sqrt{1+\sin(2\theta)}+\sqrt{1-\sin(2\theta)}}$ (commencer par étudier toutes les symétries de l'ensemble considéré).
\finenonce{005205}


\finexercice


\exercice{5530, rouget, 2010/07/15}
\enonce{005530}{}
Construire les courbes suivantes~:
\begin{enumerate}
 \item $r=\sqrt{\cos(2\theta)}$,
 \item $r=\sin\left(\frac{2\theta}{3}\right)$,
 \item $r=ae^{b\theta},\;(a,b)\in]0,+\infty[^2$,
 \item $r=2\cos(2\theta)+1$,
 \item $r=\tan\left(\frac{2\theta}{3}\right)$.
\end{enumerate}
\finenonce{005530}


\finexercice


\exercice{5531, rouget, 2010/07/15}
\enonce{005531}{}
Etude complète de la courbe d'équation polaire $r=\frac{2\cos\theta+1}{2\sin\theta+1}$.
\finenonce{005531}


\finexercice


\exercice{5532, rouget, 2010/07/15}
\enonce{005532}{La cardioïde}
Soit la courbe d'équation polaire $r=a(1+\cos\theta)$, $a>0$. 
\begin{enumerate}
 \item  Construire la courbe.
 \item  Longueur et développée. 
\end{enumerate}
\finenonce{005532}


\finexercice


\exercice{5533, rouget, 2010/07/15}
\enonce{005533}{}
Construire la courbe d'équation cartésienne $x^2(x^2+y^2)-(y-x)^2=0$ après être passé en polaires .
\finenonce{005533}


\finexercice


\exercice{5534, rouget, 2010/07/15}
\enonce{005534}{}
Développée de la spirale logarithmique d'équation polaire $r=ae^{\theta}$ ($a>0$).
\finenonce{005534}


\finexercice


\exercice{6989, blanc-centi, 2015/07/04}
\video{QZhe-OF_d10}
\enonce{006989}{}
\'Etudier les courbes d'équations polaires suivantes:
\begin{enumerate}
\item $\displaystyle r(\theta)=\frac{1}{\sqrt{\tan(2\theta)}}$ \quad pour $\theta \in ]0,\frac\pi4[$
\item $\displaystyle r(\theta)=\frac{\sin^2\theta}{\cos\theta}$ \quad pour $\theta \in ]-\frac\pi2,\frac\pi2[$\quad {\it (La cisso\"ide droite)}
\item $\displaystyle r(\theta)=\sqrt{\cos(2\theta)}$ \quad {\it (La lemniscate de Bernoulli)}
\end{enumerate}
\finenonce{006989}
 

\finexercice


\exercice{6990, blanc-centi, 2015/07/04}
\video{jlbJn2n89PY}
\enonce{006990}{}
On considère les courbes $\mathcal{C}_1$ et $\mathcal{C}_2$ {\it (des limaçons de Pascal)} 
respectivement données en polaires par 
$$r_1(\theta)=1+\cos\theta\quad\quad r_2(\theta)=3+\cos\theta$$
Pour $i=1,2$, on note $N_i(\theta)$ la droite orthogonale au point $M_i(\theta)\in\mathcal{C}_i$. 
Vérifier que pour tout $\theta\not\equiv 0\ [2\pi]$, les droites $N_1(\theta)$ et $N_2(\theta)$ sont sécantes, en un point $P(\theta)$. Déterminer le lieu du point $P$ quand $\theta$ varie.
\finenonce{006990}
 

\finexercice

\section{244.03 Courbes définies par une condition}


\exercice{4994, quercia, 2010/03/17}
\enonce{004994}{Sous-tangente, sous-normale}

Soit $\mathcal{C}$ une courbe du plan. A un point $M$ un point de $\mathcal{C}$, on associe les
points $H$, $T$ et $N$ selon le dessin :
$$
\includegraphics[height=3cm]{../images/img004994-1}
$$

Déterminer les courbes d'équation $y = f(x)$ vérifiant la condition suivante :

\begin{enumerate}
  \item $\overline{HT} ={}$cste.
    
  \item $\overline{HN} ={}$cste.
    
  \item $MN = {}$cste.
    
  \item $MT = {}$cste.
    
  \item $AN=MN$ où $A$ est le point de coordonnées $(0,a)$.
    

\end{enumerate}
\finenonce{004994}



\finexercice


\exercice{4995, quercia, 2010/03/17}
\enonce{004995}{Sous-tangente, sous-normale}
Soit $\mathcal{C}$ une courbe du plan. A un point $M$ un point de $\mathcal{C}$, on associe les
points $T$ et $N$ selon le dessin :
$$
\includegraphics[height=4cm]{../images/img004995-1}
$$

Déterminer les courbes vérifiant la condition suivante :

\begin{enumerate}
  \item $\overline{OT} = {}$cste.
    
  \item $\overline{ON} = {}$cste.
    
\end{enumerate}
\finenonce{004995}



\finexercice


\exercice{4996, quercia, 2010/03/17}
\enonce{004996}{Milieu fixe}

Soit $D$ une droite du plan et $\mathcal{C}$ une courbe paramétrée.
Pour $M \in \mathcal{C}$ on note $T$ et $N$ les points d'intersection de $D$ avec la
tangente et la normale à $\mathcal{C}$ en $M$.
Déterminer $\mathcal{C}$ telle que le milieu de $[T,N]$ reste fixe.

$\Bigl($On paramètrera $\mathcal{C}$ par $t = \frac{y'}{x'}\Bigr)$
\finenonce{004996}



\finexercice


\exercice{4997, quercia, 2010/03/17}
\enonce{004997}{Distance $TN$ constante}

Soit $D$ une droite du plan et $\mathcal{C}$ une courbe paramétrée.
Pour $M \in \mathcal{C}$ on note $T$ et $N$ les points d'intersection de $D$ avec la
tangente et la normale à $\mathcal{C}$ en $M$.
Déterminer $\mathcal{C}$ telle que la distance $TN$ reste constante.

$\Bigl($On paramètrera $\mathcal{C}$ par $t = \frac{y'}{x'}\Bigr)$
\finenonce{004997}



\finexercice


\exercice{4998, quercia, 2010/03/17}
\enonce{004998}{Ensi Chimie P' 93}
$$
\includegraphics[width=4cm]{../images/img004998-1}
$$
 Trouver les courbes $\mathcal{C}$ telles que $MN = ON$.
\finenonce{004998}



\finexercice


\exercice{4999, quercia, 2010/03/17}
\enonce{004999}{Ensi Physique P 94}

    Trouver les arcs biréguliers du plan dont le cercle osculateur est en
    tout point tangent à une droite fixe.
    
\finenonce{004999}



\finexercice


\exercice{5000, quercia, 2010/03/17}
\enonce{005000}{L'homothétique du cercle osculateur reste tangent à $Ox$}

Déterminer les courbes planes telles que l'image du cercle osculateur en un
point $M$ par l'homothétie de centre $M$ et de rapport 2 reste tangente à $Ox$.

On prendra $\varphi$ comme paramètre et on cherchera une équation différentielle
sur le rayon de courbure $R$.

\finenonce{005000}



\finexercice


\exercice{5001, quercia, 2010/03/17}
\enonce{005001}{Ensi P 91}

    On se place dans un plan affine euclidien rapporté à un repère orthonormé.
    Donner l'ensemble des trajectoires orthogonales de la famille des cercles
    de rayon constant $a$ $(a > 0)$ centrés sur $Ox$.
    
\finenonce{005001}



\finexercice


\exercice{5002, quercia, 2010/03/17}
\enonce{005002}{\'Equations intrinsèques}

Soit $f:\R \to \R$ une fonction continue. On étudie les
courbes planes paramétrées par une abscisse curviligne, $s$, telles que la
courbure au point $M_s$ soit $c = f(s)$.

\begin{enumerate}
  \item Montrer que si l'on impose la position de $M_0$ et la tangente en ce point,
    le problème admet une solution unique.

  \item Dans le cas général, démontrer que les courbes solutions se déduisent
    d'une courbe particulière en appliquant un déplacement du plan arbitraire.

  \item \'Etudier les équations : $c = {}$cste, $c = \frac 1s$
    (spirale logarithmique).
\end{enumerate}
\finenonce{005002}



\finexercice


\exercice{5003, quercia, 2010/03/17}
\enonce{005003}{\'Equations intrinsèques}
Chercher les courbes planes vérifiant l'équation intrinsèque :

\begin{enumerate}
  \item $R=s$.
    

  \item $Rs = 1$.
    

  \item $R^2 = 2as$, $a > 0$ donné.
    

  \item $R = 1+s^2$.
    

  \item $R^2+s^2 = a^2$.
    
\end{enumerate}
\finenonce{005003}



\finexercice


\exercice{5004, quercia, 2010/03/17}
\enonce{005004}{$I$ reste sur un cercle}

Trouver les courbes planes $\mathcal{C}$ telles que le centre de courbure reste
sur un cercle $\mathcal{C}(O,r)$ fixe.
(On prendra $\varphi$ comme paramètre)


\finenonce{005004}



\finexercice


\exercice{5005, quercia, 2010/03/17}
\enonce{005005}{$M - s/2M'$ reste sur $Ox$}

Soit $\mathcal{C}$ une courbe plane et $s$ une abscisse curviligne sur $\mathcal{C}$.
A chaque point $M \in \mathcal{C}$ d'abscisse curviligne $s$, on associe le point
$N = M - \frac s2\vec T$.
Trouver $\mathcal{C}$ telle que $N$ reste sur $Ox$.

\finenonce{005005}



\finexercice


\exercice{5006, quercia, 2010/03/17}
\enonce{005006}{$MC = kMN$}

Trouver les courbes $\Gamma$ du plan ayant la propriété suivante :
Soit $M \in \Gamma$, $C$ le centre de courbure de $\Gamma$ en $M$ et $N$
le projeté de $O$ sur la normale à $\Gamma$ en $M$.
Alors $\vec{MC} = k\vec{MN}$ où $k$ est un réel fixé.

\'Etudier les cas particuliers :
$k=1$, $k=\frac23$, $k=2$, $k=\frac13$ et $k=-1$.

\finenonce{005006}



\finexercice


\exercice{5529, rouget, 2010/07/15}
\enonce{005529}{}
Soit $T$ l'intersection de $(Ox)$ et de la tangente en $M$ et $H$ le projeté orthogonal de $M$ sur $(Ox)$. Trouver les
courbes telles que
\begin{enumerate}
\item  $MT=a$ ($a>0$ donné)
\item  $HT=a$ (sans rapport avec 1))
\end{enumerate}
\finenonce{005529}



\finexercice

\section{244.04 Branches infinies}


\exercice{5007, quercia, 2010/03/17}



\enonce{005007}{Branches infinies}
Déterminer les branches infinies pour les courbes paramétrées suivantes :
\begin{enumerate}
\item$x=4t^5-4t^3+t$,
$y=\dfrac{t}{3t^4+1}$                                 \par


% %-----------------------------------------------------------------------------
% \vtop to 7cm{\mapleplot{%					img1
% x:=4*t^5-4*t^3+t;
% y:=t/(3*t^4+1);
% print(plot([x,y,t=-2..2],-3..3,numpoints=200));}      \par
% \vss}


\item $x=2\cos^2t+\ln|\sin t|$,
 $y=\sin 2t$        \par

% %-----------------------------------------------------------------------------
% \vtop to 7cm{\mapleplot{%					img2
% x:=2*cos(t)^2+ln(abs(sin(t)));
% y:=sin(2*t);
% print(plot([x,y,t=0..Pi],-2..2,scaling=constrained));}  \par
                                     
                                     % \vss}

\item$x=\sqrt{\dfrac{t^2-2}{t^4-1}}$,
$y=tx$                                                  \par

% %-----------------------------------------------------------------------------
% \vtop to 7cm{\mapleplot{%					img3
% x:=sqrt((t^2-2)/(t^4-1));
% u:=sqrt((1-2*t^2)/(1-t^4));
% p1:=plot([x,t*x,t=-1..1],-2..3,-2..2):
% p2:=plot([t*u,u,t=0..sqrt(1/2)],-2..3,-2..2):
% p3:=plot([t*u,-u,t=0..sqrt(1/2)],-2..3,-2..2):
% p4:=plot([abs(t),t,t=-2..2],-2..3,-2..2,color=yellow):
% with(plots):
% print(display({p1,p2,p3,p4}));}                         \par
% \vss}


\item $x=\dfrac{t^3-t}{2t-1}$,
 $y=tx$                                                   \par

% %-----------------------------------------------------------------------------
% \vtop to 7cm{\mapleplot{%					img4
% x:=(t^3-t)/(2*t-1);
% y:=t*x;
% print(plot([x,y,t=-2..2],-1..1,-1..1));}                 \par
% \vss}

\item$x=\dfrac1t+\dfrac1{t+1}$, $y=\dfrac1t+\dfrac1{(t+1)^2}$                            \par
% %-----------------------------------------------------------------------------
% \vtop to 7cm{\mapleplot{%					img5
% x:=1/t+1/(t+1);
% y:=1/t+1/(t+1)^2;
% print(plot([x,y,t=-10..10],-4..4,-3..3,numpoints=200));} \par
% \vss}

\item $x=\sin \dfrac t2$,
 $y=\tan t$                                               \par
% %-----------------------------------------------------------------------------
% \vtop to 7cm{\mapleplot{%					img6
% x:=sin(t/2);
% y:=tan(t);
% print(plot([x,y,t=0..4*Pi],-1.5..1.5,-2..2));}           \par
% \vss}\endmulticolonne
% 
% \vfill\eject
% \vphantom{\Gros\bf Branches infinies}\SV\multicolonne 2

 \item$x=\dfrac1t+\dfrac1{t+1}$,
 $y=\dfrac1t-\dfrac1{t+1}$                                \par
% %-----------------------------------------------------------------------------
% \vtop to 7cm{\mapleplot{%					img7
% x:=1/t+1/(t+1);
% y:=1/t-1/(t+1);
% print(plot([x,y,t=-5..5],-7..7,-6..4,numpoints=200));}   \par
% \vss}

 \item$x=\dfrac{3t}{1+t^3}$,
 $y=tx$                                                   \par
% %-----------------------------------------------------------------------------
% \vtop to 7cm{\mapleplot{%					img8
% x:=(3*t)/(1+t^3);
% y:=t*x;
% print(plot([x,y,t=-10..10],-4..4,-3..3));}               \par
% \vss}

 \item$x=\dfrac{te^t}{t+1}$,
 $y=\dfrac{e^t}{t+1}$                                     \par
% %-----------------------------------------------------------------------------
% \vtop to 7cm{\mapleplot{%					img9
% y:=exp(t)/(t+1);
% x:=t*y;
% print(plot([x,y,t=-5..3],-4..4,-3..3));}                 \par
% \vss}

\item$x=2t^3+3t^2$,
$y=3t^2+6t$                                              \par
% %-----------------------------------------------------------------------------
% \vtop to 7cm{\mapleplot{%					img10
% x:=2*t^3+3*t^2;
% y:=3*t^2+6*t;
% print(plot([x,y,t=-5..5],-50..50,-2..50));}              \par
% \vss}

\item$x=t^3-3t$,               
$y=t^3-t^2-t+1$                                          \par
% %-----------------------------------------------------------------------------
% \vtop to 7cm{\mapleplot{%					img11
% x:=t^3-3*t;
% y:=t^3-t^2-t+1;
% print(plot([x,y,t=-3..3],-10..10,-7..7));}               \par
% \vss}

\item$x=\dfrac t{t^2-1}$, 
$y=\dfrac{t^2}{t-1}$                                     \par
% %-----------------------------------------------------------------------------
% \vtop to 7cm{\mapleplot{%					img12
% x:=t/(t^2-1);
% y:=t^2/(t-1);
% print(plot([x,y,t=-5..5],-4..4,-3..5));}                 \par
% \vss}
\end{enumerate}
\finenonce{005007}



\finexercice

\section{244.05 Points de rebroussement}


\exercice{5008, quercia, 2010/03/17}
\enonce{005008}{Rebroussements}
\begin{enumerate}
\item $x=2t^3+3t^2$, $y=3t^2+6t$                  

% 
% %-----------------------------------------------------------------------------
% \vtop to 7cm{\mapleplot{% img25
% x:=2*t^3+3*t^2;
% y:=3*t^2+6*t;
% print(plot([x,y,t=-3..3],-6..4,-3..3));}  \par
% \vss}


\item $x=t^3-3t$, $y=t^3-t^2-t+1$ 
% %-----------------------------------------------------------------------------
% \vtop to 7cm{\mapleplot{% img26
% x:=t^3-3*t;
% y:=t^3-t^2-t+1;
% print(plot([x,y,t=-3..3],-10..10,-7..7));}   \par
% \vss}


\item  $x=\sin t$, $y=\dfrac{\cos^2t}{2-\cos t}$           

% %-----------------------------------------------------------------------------
% \vtop to 7cm{\mapleplot{% img27
% x:=sin(t);
% y:=cos(t)^2/(2-cos(t));
% print(plot([x,y,t=0..2*Pi],scaling=constrained));} \par
% \vss}

\item  $x=(1+\cos^2t)\sin t$, $y=\sin^2t\cos t$               
% %-----------------------------------------------------------------------------
% \vtop to 7cm{\mapleplot{% img28
% x:=(1+cos(t)^2)*sin(t);
% y:=sin(t)^2*cos(t);
% print(plot([x,y,t=0..2*Pi],scaling=constrained));} \par
% \vss}

\item  $x=(1+\cos t)\sin 2t$, $y=\cos 2t$

% %-----------------------------------------------------------------------------
% \vtop to 7cm{\mapleplot{% img29
% x:=(1+cos(t))*sin(2*t);
% y:=cos(2*t);
% print(plot([x,y,t=0..2*Pi],scaling=constrained));} \par
% \vss}
% 
% 

\end{enumerate}
\finenonce{005008}


\finexercice

\section{244.06 Enveloppes}


\exercice{5009, quercia, 2010/03/17}
\enonce{005009}{Esem 91}

    Soit $\mathcal{C}$ le cercle : $x^2+y^2=1$. Soit $M$ un point de $\mathcal{C}$ d'angle
    polaire $\theta$ et $D_\theta$ la droite passant par $M$ d'angle polaire
    $2\theta$. Trouver l'enveloppe des droites $D_\theta$.
    
\finenonce{005009}



\finexercice


\exercice{5010, quercia, 2010/03/17}
\enonce{005010}{Ensi Physique 93}

Soit $\mathcal{C}$ un cercle de centre $O$ et de rayon $R$, et $S$ un point du
plan différent de $O$. Donner l'enveloppe des normales en $M$ à $(SM)$
lorsque $M$ décrit $\mathcal{C}$.
\finenonce{005010}



\finexercice


\exercice{5011, quercia, 2010/03/17}
\enonce{005011}{Cordes sur une parabole}

Soit $\cal P$ la parabole d'équation $y^2 = 2px$.
Chercher l'enveloppe des cordes $[A,B]$ de $\cal P$ de hauteur $h > 0$
donnée.

\finenonce{005011}



\finexercice


\exercice{5012, quercia, 2010/03/17}
\enonce{005012}{Cordes sur une parabole}
Soit $\cal P$ la parabole d'équation $y^2 = 2px$.
Pour $A,B \in \cal P$ distincts, on note $C$ le point d'intersection des
tangentes en $A$ et $B$.
Trouver l'enveloppe des droites $(AB)$ lorsque l'aire du triangle $ABC$ reste
constante.

\finenonce{005012}



\finexercice


\exercice{5013, quercia, 2010/03/17}
\enonce{005013}{Cordes sur une parabole}
Soient $M$, $M'$ deux points d'une parabole $\cal P$ tels que
$(MM')$ passe par le foyer $F$. Quels sont :

\begin{enumerate}
  \item L'enveloppe des droites $(MM')$ ?
    
  \item Le lieu des milieux des segments $[M,M']$ ?
    
  \item L'enveloppe des médiatrices de $[M,M']$ ?
    
\end{enumerate}
\finenonce{005013}



\finexercice


\exercice{5014, quercia, 2010/03/17}
\enonce{005014}{Rayons réfléchis sur une parabole}

Soit $\cal P$ la parabole d'équation $y^2 = 2px$.

\begin{enumerate}
  \item Un rayon incident arrive suivant une parallèle à $Ox$ et se réfléchit à
    ``l'intérieur'' de $\cal P$ avec le même angle.
    Trouver l'enveloppe des rayons réfléchis.
    

  \item Même question, mais le rayon incident est parallèle à $Oy$.
    
\end{enumerate}
\finenonce{005014}



\finexercice


\exercice{5015, quercia, 2010/03/17}
\enonce{005015}{Cercle osculateur à une parabole}

Soit $\cal P$ une parabole, $M \in \cal P$ et $\cal C$ le cercle osculateur à
$\cal P$ en $M$.
Montrer que, sauf cas particulier, $\cal C$ recoupe $\cal P$ en un deuxième
point $P$.
Déterminer l'enveloppe des droites $(MP)$.

\finenonce{005015}



\finexercice


\exercice{5016, quercia, 2010/03/17}
\enonce{005016}{Cordes d'une hyperbole}

Soit $\cal H$ une hyperbole de foyer $F$. Trouver l'enveloppe des cordes
$[P,Q]$ de $\cal H$ vues depuis $F$ sous un angle droit.
\finenonce{005016}



\finexercice


\exercice{5017, quercia, 2010/03/17}
\enonce{005017}{Cardioïde}

Pour $\theta \in \R$, on note $A_\theta = (\cos\theta,\sin\theta)$.
Chercher l'enveloppe des droites $D_\theta = (A_\theta\ A_{2\theta})$.

\finenonce{005017}



\finexercice


\exercice{5018, quercia, 2010/03/17}
\enonce{005018}{Cycloïde}

Chercher l'enveloppe d'un diamètre $\Delta$ d'un cercle $\cal C$ roulant sans
glisser sur une droite $D$.
Comparer le point caractéristique à la projection orthogonale du point de
contact $I$ sur $\Delta$.

\finenonce{005018}



\finexercice


\exercice{5019, quercia, 2010/03/17}
\enonce{005019}{Hypocycloïde}

Soit $\cal C$ un cercle passant par $O$ centré sur $Ox$. Pour $M \in \cal C$, on
note $D_M$ la droite symétrique de $(OM)$ par rapport à l'horizontale passant
par $M$.
Déterminer l'enveloppe des droites $D_M$ et la construire.

\finenonce{005019}



\finexercice


\exercice{5020, quercia, 2010/03/17}
\enonce{005020}{Cordes de $\rho = a/\cos(3\theta)$}

Tracer la courbe d'équation polaire $\rho = \frac{a}{\cos 3\theta}$, $a>0$.
Chercher l'enveloppe des cordes vues de $O$ sous un angle droit.

\finenonce{005020}



\finexercice


\exercice{5021, quercia, 2010/03/17}
\enonce{005021}{Perpendiculaire à $OM$ sur une ellipse}

Soit $\cal E$ une ellipse de centre $O$, de paramètres $a$ et $b$.
Pour $M \in \cal E$, soit $D$ la perpendiculaire en $M$ à $(OM)$.

\begin{enumerate}
  \item Donner les équations paramétriques de l'enveloppe des droites $D$.
    
  \item Tracer les enveloppes sur ordinateur pour différentes valeurs de $a/b$.
  \item \'Etudier les points stationnaires de l'enveloppe quand il y en a.
\end{enumerate}
\finenonce{005021}



\finexercice


\exercice{5022, quercia, 2010/03/17}
\enonce{005022}{$AM \perp D$}

Soit $D$ une droite du plan et $A$ un point non élément de $D$.
Soit $M$ un point variable sur $D$. Trouver l'enveloppe de la normale en $M$
à $(AM)$.
\finenonce{005022}



\finexercice


\exercice{5023, quercia, 2010/03/17}
\enonce{005023}{Concavité}

Soient $u,v,w$ de classe ${\cal C}^2$,
$D_t$ la droite d'équation : $u(t)x + v(t)y + w(t) = 0$,
et $\Gamma$ l'enveloppe des droites $D_t$.

On note : $\delta = \begin{vmatrix} u &v \cr u' &v' \cr\end{vmatrix}$,
          $\Delta = \begin{vmatrix} u &v &w \cr u' &v' &w' \cr u'' &v'' &w'' \cr\end{vmatrix}$,
et on suppose pour tout $t$ : $\delta\Delta w(t) \ne 0$.

Montrer que $\Gamma$ tourne sa concavité vers $O$ si et seulement si
pour tout $t$ : $\delta\Delta w(t) > 0$.
\finenonce{005023}



\finexercice

\section{244.07 Propriétés métriques : longueur, courbure,...}


\exercice{2059, ridde, 1999/11/01}

\enonce{002059}{}
D\'eterminer la longueur de la courbe $y = \sqrt{x} (1-\dfrac x3)$ pour $0 \leq
x \leq 3$.
\finenonce{002059}



\finexercice



\exercice{2060, ridde, 1999/11/01}

\enonce{002060}{}
D\'eterminer une abscisse curviligne, la longueur et la d\'evelopp\'ee de l'astroïde.
\finenonce{002060}



\finexercice



\exercice{2061, ridde, 1999/11/01}

\enonce{002061}{}
Calculer le rayon de courbure de $\rho (\theta) = \cos (\dfrac{\theta}3)$ en
fonction de $\rho$.
\finenonce{002061}



\finexercice



\exercice{2062, ridde, 1999/11/01}

\enonce{002062}{}
Soit $\mathcal{P}$ la parabole $y^2 = x$. D\'eterminer une \'equation param\'etr\'ee
et une \'equation cart\'esienne de $\Gamma$ la d\'evelopp\'ee de $\mathcal{P}$.
Tracer $\Gamma$.
\finenonce{002062}



\finexercice



\exercice{2063, ridde, 1999/11/01}

\enonce{002063}{}
Soit $\Gamma$ la courbe $\rho (\theta) = \sqrt{\sin (2 \theta)}$.
\begin{enumerate}
\item Tracer cette courbe.
\item Calculer le rayon de courbure.
\item Soient $I$ le centre de courbure en $M$ et $H$ le projet\'e orthogonal de $I$
sur $ (OM)$. D\'eterminer $\overrightarrow{MH}$.
\item En d\'eduire une construction g\'eom\'etrique de la d\'evelopp\'ee de $\Gamma$.
\end{enumerate}
\finenonce{002063}



\finexercice



\exercice{2064, ridde, 1999/11/01}

\enonce{002064}{}
Soit $M (s)$ un arc $C^2$ bir\'egulier param\'etr\'e par une abscisse curviligne.
Soit $\mathcal{R}$ le rep\`ere de Fr\'enet $ (M (0), \vec{t} (0), \vec{n} (0))$.
On note $ (X (s), Y (s))$ les coordonn\'ees dans ce rep\`ere d'un point $M (s)$ de la courbe.
\begin{enumerate}
\item Montrer que si $R_{0}$ est le rayon de courbure en $M (0)$ alors
$R_0 = \lim\limits_{s \rightarrow 0} \dfrac{X^2 (s)}{2Y (s)}$.
\item En d\'eduire le rayon de courbure au point $\theta  = 0$ de la courbe
$\rho (\theta) = 1 + 2\cos (\frac{\theta}2)$. 
\end{enumerate}
\finenonce{002064}



\finexercice



\exercice{5024, quercia, 2010/03/17}
\enonce{005024}{Calcul de longueur}

\def\arc#1{\mathop{#1}\limits^\frown}
Déterminer la longueur d'un arc $\arc{M_0M_t}$
ou $\arc{M_0M_\theta}$ pour les courbes :

\begin{enumerate}
  \item $x = t - \ch t\sh t$,\quad $y = 2\ch t$ 
  \item $\rho = \tanh \frac\theta2$.             
\end{enumerate}
\finenonce{005024}



\finexercice


\exercice{5025, quercia, 2010/03/17}
\enonce{005025}{Calcul de longueur}
\def\arc#1{\mathop{#1}\limits^\frown}
Soit la courbe paramétrée par : $x = 2t^3+3t^2$,\quad $y = 3t^2+6t$.
Calculer la longueur de l'arc $\arc{AO}$ où $A$ est le point de
rebroussement.
\finenonce{005025}



\finexercice


\exercice{5026, quercia, 2010/03/17}
\enonce{005026}{Calcul de longueur}
\def\arc#1{\mathop{#1}\limits^\frown}
Calculer la longueur totale des courbes suivantes :

\begin{enumerate}
  \item $x = (1+\cos^2t)\sin t$,\quad $y=\sin^2t\cos t$.
    
  \item $\rho = \sin^2\frac\theta2$.
    
\end{enumerate}
\finenonce{005026}



\finexercice


\exercice{5027, quercia, 2010/03/17}
\enonce{005027}{TPE MP 2003}
Nature, construction et longueur de la courbe d'équation $\sqrt x + \sqrt y = 1$.

\finenonce{005027}



\finexercice


\exercice{5028, quercia, 2010/03/17}
\enonce{005028}{Comparaison de longueurs (ENS MP 2002)}

Soit $f : {[a,b]} \to \R$ continue concave, $\mathcal{C}^1$ par morceaux,
$L_1$ la courbe paramétrée $x  \mapsto(x,f(x))$ et
$L_2$ un chemin continu $\mathcal{C}^1$ par morceaux joignant
les extrémités de $L_1$ et situé au-dessus de $L_1$.
Montrer que la longueur de $L_2$ est supérieure ou égale à celle de $L_1$.


\finenonce{005028}



\finexercice


\exercice{5029, quercia, 2010/03/17}
\enonce{005029}{Centre de courbure}

Déterminer les coordonnées du centre de courbure au point $M$ pour les courbes
suivantes :

\begin{enumerate}
  \item $x = 3t-t^3$,\quad $y=3t^2$.
    

  \item $x=2\cos t+\cos 2t$,\quad $y=2\sin t-\sin2t$.
    

  \item $x=t-\sin t$,\quad $y=1-\cos t$.
    (Cycloïde, indiquer une relation géométrique simple entre la courbe décrite
    par $M$ et celle décrite par $I$)
    

  \item $x = a\cos^3t$,\quad $y=a\sin^3t$.
    (Astroïde) Construire le courbe $\mathcal{C}$ et sa développée,
    puis prouver par le calcul qu'elles sont semblables.
 
  \item Hyperbole d'équation $xy=1$.
    

  \item Ellipse d'équation $\frac{x^2}{a^2} + \frac{y^2}{b^2} = 1$.
    

  \item $\rho = e^\theta$. (Spirale logarithmique)
    

  \item $\rho = 1+\cos\theta$. (Cardioïde)
    
\end{enumerate}
\finenonce{005029}



\finexercice


\exercice{5030, quercia, 2010/03/17}
\enonce{005030}{Points sur une hyperbole (Ensi P 91)}

    Soit la courbe $\Gamma$ définie par : $xy = a^2$, $(a>0)$.
    Pour chaque point $M$ on définit le point $\Omega$ par :
    $2\vec{\Omega M} = \vec{MN}$,
    où $N$ est le point où $\Gamma$ recoupe sa normale en $M$.
    Montrer que $\Omega$ est le centre de courbure de $\Gamma$ en $M$.
    
\finenonce{005030}



\finexercice


\exercice{5031, quercia, 2010/03/17}
\enonce{005031}{Cercle circonscrit à trois points}

Soit $\mathcal{C}$ une courbe plane paramétrée par une abscisse curviligne $s$.
Soit $s_0$ fixé.

\begin{enumerate}
  \item   Donner le DL à l'ordre 2 de $M_s$ pour $s \to s_0$ dans le repère
      de Frenet en $M_{s_0}$.

  \item   On suppose $c(s_0) \ne 0$. Montrer que pour $h$ assez petit, les points
      $M_{s_0-h}$, $M_{s_0}$, $M_{s_0+h}$ ne sont pas alignés.

  \item   Soit $\Gamma_h$ le cercle circonscrit à ces trois points, et $R_h$
      son rayon. Chercher $\lim_{h\to 0} R_h$.
\end{enumerate}
\finenonce{005031}



\finexercice


\exercice{5032, quercia, 2010/03/17}
\enonce{005032}{Propriétés de la cycloïde}

Soit $\mathcal{C}$ la courbe d'équations paramétriques
$\begin{cases} x = a(t-\sin t) \cr y = a(1-\cos t) \cr \end{cases}$
pour $t \in {]0,2\pi[}$ (arche de cycloïde).
On note $S$ le point de paramètre $\pi$, et $D$ la tangente à $\mathcal{C}$ en $S$.

Soit $M \in \mathcal{C}\setminus\{S\}$, $I$ le point d'intersection de la normale à $\mathcal{C}$ en $M$ et
de $Ox$, et $J$ le point d'intersection de la tangente en $M$ avec $D$.

\begin{enumerate}
  \item Faire un dessin.
  \item Montrer que $I$ et $J$ ont même abscisse.
  \item On prend $S$ comme origine des abscisses curvilignes. Trouver une relation
    entre $s$ et $\vec{MJ}$.
\end{enumerate}
\finenonce{005032}



\finexercice


\exercice{5033, quercia, 2010/03/17}
\enonce{005033}{Normales à une cardioïde}

On considère la courbe $\mathcal{C}$ d'équation polaire $\rho = 1 + \cos\theta$
(cardioïde).

\begin{enumerate}
  \item Dessiner $\mathcal{C}$.
  \item Une droite $D$ passant par $O$ coupe $\mathcal{C}$ en deux points $M_1$ et $M_2$.
    Soient $\Delta_1$, $\Delta_2$ les normales à $\mathcal{C}$ en ces points et $P$
    le point d'intersection de $\Delta_1$ et $\Delta_2$. Quelle est la courbe
    décrite par $P$ lorsque $D$ tourne autour de $O$~?
    
\end{enumerate}
\finenonce{005033}



\finexercice


\exercice{5034, quercia, 2010/03/17}
\enonce{005034}{Calcul de courbure par TFI}

Déterminer le rayon de courbure de la courbe $\mathcal{C}$ d'équation : $2x^2+y^2=1$
aux points intersection de $\mathcal{C}$ et des axes $Ox$ et $Oy$.

\finenonce{005034}



\finexercice


\exercice{5035, quercia, 2010/03/17}
\enonce{005035}{Calcul de courbure par TFI}
Soit $\mathcal{C}$ la courbe d'équation cartésienne $ x^4+y^4+x^3+y^3=2$.
En utilisant le théorème des fonctions implicites, calculer la courbure de $\mathcal{C}$
en $A = (-1,1)$.

\finenonce{005035}



\finexercice


\exercice{5036, quercia, 2010/03/17}
\enonce{005036}{Calcul de courbure (Chimie P' 90)}

Déterminer l'ensemble des centres de courbure en $O$ aux courbes intégrales
de l'équation dif\-fé\-ren\-tielle $(1-x^2)y''-xy'-2y=1$ telles que $y(0)=0$.


\finenonce{005036}



\finexercice


\exercice{5037, quercia, 2010/03/17}
\enonce{005037}{Courbe parallèle à une parabole}

Soit $\mathcal{C} : t \mapsto M_t$ une courbe plane paramétrée sans point stationnaire.
Les courbes parallèles à $\mathcal{C}$ sont les courbes de la forme :
$t \mapsto M_t + \lambda \vec N$,
ou $\vec N$ est le vecteur normal en $M_t$ et $\lambda$ est constant.

\begin{enumerate}
  \item Montrer que le parallélisme est une relation d'équivalence entre arcs
    sans points stationnaires.
  \item Construire les parallèles à la parabole d'équation $y = x^2$ pour
    $\lambda = \pm 2$.
\end{enumerate}
\finenonce{005037}



\finexercice


\exercice{5038, quercia, 2010/03/17}
\enonce{005038}{Points équidistants sur la tangente}

Soit $\mathcal{C}$ une courbe paramétrée, $(M,\vec t,\vec n)$ le repère de Frenet en un
point $M$ de $\mathcal{C}$. Soit $a > 0$ fixé et $P_1 = M + a\vec t$, $P_2 = M - a\vec t$.
On note $\mathcal{C}_1$ et $\mathcal{C}_2$ les courbes décrites par $P_1$ et $P_2$ quand $M$
décrit $\mathcal{C}$ et $c_1$, $c_2$ les courbures correspondantes.
Soit $C$ le centre de courbure à $\mathcal{C}$ en $M$.

Montrer que $c_1 + c_2 = \frac2{CP_1}$ et que les trois normales sont
concourantes.



\finenonce{005038}



\finexercice


\exercice{5039, quercia, 2010/03/17}
\enonce{005039}{Paraboles de cercle osculateur donné}

Soit $\mathcal{C}$ le cercle d'équation $x^2+y^2-2Rx = 0$ et $\Delta$ une droite
variable passant par $O$.

\begin{enumerate}
  \item Chercher l'équation de la parabole $\cal P$ d'axe parallèle à $\Delta$,
    passant par $O$, dont $\mathcal{C}$ est le cercle osculateur en~$O$.
    

  \item Quelle est l'enveloppe des paraboles précédentes ?
\end{enumerate}
\finenonce{005039}



\finexercice


\exercice{5040, quercia, 2010/03/17}
\enonce{005040}{Développante}

\begin{enumerate}
  \item Construire la courbe $\mathcal{C}$ d'équations paramétriques :
    $x=3t-t^3$,\quad $y=3t^2$.
  \item Chercher les équations paramétriques des développantes de $\mathcal{C}$.
  \item  Tracer la développante qui rencontre $\mathcal{C}$ à l'origine.
\end{enumerate}
\finenonce{005040}



\finexercice


\exercice{5041, quercia, 2010/03/17}
\enonce{005041}{Développante}
Déterminer la développante de la chainette $\mathcal{C}$ d'équation $y = a\ch(x/a)$ qui
rencontre $\mathcal{C}$ pour $x=0$. (Tractrice) Dessiner les deux courbes.

\finenonce{005041}



\finexercice


\exercice{5535, rouget, 2010/07/15}
\enonce{005535}{}
Longueur $L$ de $(\Gamma)$ dans chacun des cas suivants~:
\begin{enumerate}
\item  $\Gamma$ est l'astroïde de représentation paramétrique $\left\{
\begin{array}{l}
x=a\cos^3t\\
y=a\sin^3t
\end{array}
\right.$ ($a>0$ donné).
\item  $\Gamma$ est l'arche de cycloïde de représentation paramétrique $\left\{
\begin{array}{l}
x=R(t-\sin t)\\
y=R(1-\cos t)
\end{array}
\right.$, $0\leqslant t\leqslant2\pi$.
\item  $\Gamma$ est l'arc de parabole d'équation cartésienne $x^2=2py$, $0\leq x\leq a$ ($p>0$ et $a>0$ donnés).
\item  $\Gamma$ est la cardioïde d'équation polaire $r=a(1+\cos\theta)$ ( $a>0$ donné).
\end{enumerate}
\finenonce{005535}


\finexercice


\exercice{5536, rouget, 2010/07/15}
\enonce{005536}{}
Déterminer et construire la développée
\begin{enumerate}
\item  $\left\{
\begin{array}{l}
x=R\left(\cos t+\ln\left|\tan\frac{t}{2}\right|\right)\\
y=R\sin t
\end{array}
\right.$
\item  $\left\{
\begin{array}{l}
x=R(t-\sin t)\\
y=R(1-\cos t)
\end{array}
\right.$.
\item  $y=x^3$
\end{enumerate}
\finenonce{005536}


\finexercice


\exercice{5537, rouget, 2010/07/15}
\enonce{005537}{}
Trouver le point de la courbe d'équation $y=\ln x$ en lequel la valeur absolue du rayon de courbure est minimum.
\finenonce{005537}


\finexercice


\exercice{5538, rouget, 2010/07/15}
\enonce{005538}{}
Soit $(\Gamma)$ la courbe d'équation $y=\ln(\cos x)$, pour $-\frac{\pi}{2}<x<\frac{\pi}{2}$. Calculer l'abscisse curviligne $s$ quand $O$ est l'origine des abscisses curvilignes et l'orientation est celle des $x$ croissants. Trouver une relation entre $R$ et $s$. Tracer $(\Gamma)$ et sa développée.
\finenonce{005538}


\finexercice


\exercice{5539, rouget, 2010/07/15}
\enonce{005539}{}
Pour $\lambda\in\Rr$, on note $(\Gamma_\lambda)$ la courbe d'équation $y=\lambda xe^{-x}$. Quel est le lieu des centres de courbure $C_\lambda$ en $O$ à $(\Gamma_\lambda)$ quand $\lambda$ décrit $\Rr$.
\finenonce{005539}


\finexercice
\section{244.08 Courbes dans l'espace}


\exercice{5042, quercia, 2010/03/17}
\enonce{005042}{Ensi P 90}

    On considère la courbe $\mathcal{C}$ définie par :
    $x(t) = \frac{t^4}{1+t^2}$,
    $y(t) = \frac{t^3}{1+t^2}$,
    $z(t) = \frac{t^2}{1+t^2}$.\par
    A quelle condition $M_1,M_2,M_3,M_4$ quatre points de $\mathcal{C}$ de paramètres
    respectifs $t_1,t_2,t_3,t_4$ sont-ils coplanaires ?
    
\finenonce{005042}



\finexercice


\exercice{5043, quercia, 2010/03/17}
\enonce{005043}{Courbure de $M$ cste $ \Rightarrow $ courbure de $I$ cste}

Soit $\mathcal{C}$ une courbe de l'espace, et $\Gamma$ la courbe décrite par le centre
de courbure, $I$, en un point $M$ de $\mathcal{C}$.
On suppose que la courbure de $\mathcal{C}$ est constante et sa torsion non nulle.

\begin{enumerate}
  \item Montrer que la courbure de $\Gamma$ est aussi constante.
    
  \item Chercher la torsion de $\Gamma$ en $I$ en fonction de la courbure et la
    torsion de $\mathcal{C}$ en $M$.
    
\end{enumerate}
\finenonce{005043}



\finexercice


\exercice{5044, quercia, 2010/03/17}
\enonce{005044}{\'Eléments de courbure de $T$}

Soit $s \mapsto M_s$ une courbe de l'espace de classe $\mathcal{C}^3$ paramétrée par une
abscisse curviligne,
et $P$ le point tel que $\vec{OP} = \frac {\vec{dM}}{ds}$.
Chercher les éléments de courbure de la trajectoire de $P$.

\finenonce{005044}



\finexercice


\exercice{5045, quercia, 2010/03/17}
\enonce{005045}{Enveloppe de normales}

Soit $s \mapsto M_s$ une courbe de l'espace de classe $\mathcal{C}^3$ paramétrée par une
abscisse curviligne.
Pour tout $s$ on choisit une normale à la courbe en $M_s$ : $\Delta_s$.
A quelle condition les droites $\Delta_s$ admettent-elles une enveloppe~?
\finenonce{005045}



\finexercice


\exercice{5046, quercia, 2010/03/17}
\enonce{005046}{\'Equations intrinsèques en dimension 3}

Trouver les courbes de l'espace vérifiant les équations intrinsèques :
$c = \tau = \frac 1{s\sqrt2}$.

\finenonce{005046}



\finexercice

\section{244.99 Autre}


\exercice{2057, gourio, 2001/09/01}

\enonce{002057}{}
Soit $f:[0,1]\rightarrow [0,1]^{2} $ de classe $C^{1}, $ montrer que $f  $
ne peut \^{e}tre bijective.
\finenonce{002057}



\finexercice



\exercice{2058, gourio, 2001/09/01}

\enonce{002058}{}
Soit $\gamma :[0,1]\rightarrow \Cc $ continue, et $z\in \Cc$
quelconque. Montrer :
\begin{eqnarray*}
\forall \epsilon &>&0,\exists \gamma ^{\prime }\in C([0,1],\Cc)
\text{ tel que:} \\
1 &:&\forall t\in [0,1],\left| \gamma (t)-\gamma ^{\prime }(t)\right|
<\epsilon , \\
2 &:&\forall t\in [0,1],\gamma (t)\neq z.
\end{eqnarray*}
\finenonce{002058}



\finexercice



\exercice{2691, matexo1, 2002/02/01}
\enonce{002691}{}

Un cercle de rayon $R$ roule sans glisser sur l'axe ${\rm O}x$ dans le sens des
$x$ croissants. Soit {\cal C} la courbe d{\'e}crite par le point M li{\'e} {\`a} la
circonf{\'e}rence qui, dans la position initiale, co\"{\i}ncide avec l'origine
O\,(cyclo\"{\i}de). Soit ${\rm M}(\theta)$ la position du point M quand le cercle
a tourn{\'e} d'un angle $\theta$ {\`a} partir de la position initiale
et $\Omega(\theta)$ le point de contact correspondant entre la circonf{\'e}rence et
l'axe ${\rm O}x$.
\begin{itemize}
\item D{\'e}terminer l'abscisse de $\Omega(\theta)$ et les coordonn{\'e}es
$x(\theta)$ et $y(\theta)$ du point ${\rm M}(\theta)$. Montrer que la courbe
{\cal C} est p{\'e}riodique et repr{\'e}senter graphiquement la premi{\`e}re
p{\'e}riode. 
\item D{\'e}terminer, en fonction de $\theta$, le vecteur tangent
$d\overrightarrow{\rm OM}/d\theta$, le vecteur tangent unitaire
$\overrightarrow{\rm T}$, et l'{\'e}l{\'e}ment de longueur $ds$. 
\item D{\'e}terminer
le vecteur normal unitaire $\overrightarrow{\rm N}$ et le rayon de courbure
$\rho$ au point param{\'e}tr{\'e} par $\theta$. Montrer que le centre de courbure est
situ{\'e} sur la droite d{\'e}finie par ${\rm M}(\theta)$ et $\Omega(\theta)$, et
pr{\'e}ciser sa position sur cette droite. 

\item(facultatif) L'angle de rotation
est d{\'e}fini en fonction du temps par la fonction $\theta(t)$. Calculer le
vecteur vitesse $\overrightarrow{v}$ du point M {\`a} l'instant $t$. Montrer que
$\overrightarrow{v}$ s'exprime en fonction de $\Omega{\rm M}$, $d\theta/dt$ et
$\overrightarrow{\rm T}$, et donner un vecteur $\overrightarrow{\omega}$
orthogonal au plan du mouvement tel que
$\overrightarrow{v}=\overrightarrow{\omega}\land \overrightarrow{\Omega {\rm
M}}$. Obtenir g{\'e}om{\'e}triquement {\`a} un instant donn{\'e} quelconque le vecteur
vitesse d'un point P quelconque de la circonf{\'e}rence.
\end{itemize}
\finenonce{002691}



\finexercice



\exercice{2694, matexo1, 2002/02/01}
\enonce{002694}{}
Un segment AB de longueur $l$ se d{\'e}place dans le plan de fa\c{c}on que le point
A reste constamment sur l'axe ${\rm O}x$ et le point B sur l'axe ${\rm O}y$,
l'angle $\theta=(\vec{Ox},\vec {AB})$ variant de 0 {\`a} $ 2 \pi$. Soit M le point de AB tel que AM=$\alpha l$, $\alpha=C^{te}$, $0 \leq \alpha \leq 1$.
Calculer les coordonn{\'e}es de M en fonction de $\theta$ et d{\'e}terminer la courbe
qu'il d{\'e}crit. 
\begin{itemize}
\item A l'instant z{\'e}ro un oiseau s'envole d'un point A d'un
mouvement rectiligne uniforme de vitesse $\overrightarrow{v}$. Au m{\^e}me instant un
chasseur situ{\'e} au point B tire un coup de fusil en vue d'abattre l'oiseau. La
vitesse de la balle de fusil est en valeur absolue {\'e}gale {\`a} $u$. On suppose
{\'e}videmment que l'on a $u>\|\overrightarrow{v}\|$.
D{\'e}terminer la direction dans laquelle le chasseur doit tirer pour abattre
l'oiseau et l'instant $t_0$ de l'impact: on {\'e}crira deux {\'e}quations
d{\'e}terminant la vitesse vectorielle $\overrightarrow{u}$ de la balle de fusil
et l'instant $t_0$, et on en donnera les solutions. Donner l'expression de la
distance $d$ parcourue par l'oiseau entre les instants 0 et $t_0$.

Appliquer num{\'e}riquement les r{\'e}sultats pr{\'e}c{\'e}dents aux deux cas d{\'e}finis
par $A(0,0,a)$, $B(b,0,0)$, $\overrightarrow{v}=(0,v,0)$ avec~:
\item Envol {\`a} partir du repos~:
$a$=15m, $b$=10m, $v$=5m/s, $u$=300m/s 
\item Passage en plein vol~: 
$a$=20m, $b$=0, $v$=90km/h, $u$=300m/s
\end{itemize}
\finenonce{002694}
\finexercice


\exercice{2695, matexo1, 2002/02/01}
\enonce{002695}{}
Soit $\cal P$ la parabole d'{\'e}quation $y^2=2px, p>0$.
\begin{itemize}
\item Montrer que la tangente {\`a} $\cal P$ au point ${\rm M}_{0}= (x_{0},y_{0})$
a pour {\'e}quation $yy_{0}=p(x+x_{0})$. 
\item Un rayon lumineux, port{\'e} par la
droite d'{\'e}quation $y=y_{0}$ et se propageant en sens inverse de l'axe des $x$,
se r{\'e}fl{\'e}chit au point ${\rm M}_{0}$ sur la tangente {\`a} $\cal P$ selon la loi
de Descartes. D{\'e}terminer l'{\'e}quation du rayon r{\'e}fl{\'e}chi. 
\item V{\'e}rifier
que les rayons r{\'e}fl{\'e}chis correspondant aux diverses valeurs de $y_{0}$
passent tous par un m{\^e}me point F situ{\'e} sur l'axe des $x$\,(foyer de la
parabole).

Citer des applications pratiques de cette propri{\'e}t{\'e}. 
\end{itemize}
\finenonce{002695}
\finexercice


\exercice{2696, matexo1, 2002/02/01}
\enonce{002696}{}
On dispose d'un oscilloscope {\`a} deux voies. On applique sur la voie X une
tension sinuso\"{\i}dale de pulsation $\omega$, et sur la voie Y une tension de
m{\^e}me amplitude et de pulsation $2\omega$. En pla\c{c}ant l'oscilloscope en mode
X--Y et pour un choix appropri{\'e} du gain de chaque voie, on observe sur
l'{\'e}cran une courbe param{\'e}tr{\'e}e d{\'e}finie en coordonn{\'e}es cart{\'e}siennes par
les {\'e}quations: \[\hspace*{-2cm}\left\{\begin{array}{l} x(t)=a\sin \omega t \\
y(t)=a\sin 2\omega t \end{array} \right.\]
\begin{itemize}
\item D{\'e}terminer la p{\'e}riode du mouvement T. 
\item Donner le tableau des
variations de $x(t)$ et $y(t)$ sur l'intervalle [0,T], et en d{\'e}duire l'allure
de la courbe. 
\item D{\'e}terminer les coordonn{\'e}es des points de la courbe
d'abscisse ou d'ordonn{\'e}e maximum. 
\item D{\'e}terminer les sym{\'e}tries de la
courbe et donner les transformations correspondantes du param{\`e}tre $t$.
\end{itemize}
\finenonce{002696}
\finexercice


\exercice{2697, matexo1, 2002/02/01}
\enonce{002697}{}
Sur l'{\'e}cran d'un oscilloscope on observe la courbe dont les {\'e}quations
param{\'e}triques sont les suivantes: \[\hspace*{-3cm}\left\{\begin{array}{l}
x(t)=a\sin \omega t \\ y(t)=a\sin(\omega t-\varphi) \end{array} \right.\]
\begin{itemize}
\item Exprimer puis factoriser la somme et la diff{\'e}rence $x+y$ et $x-y$. 
\item
Soient $X$ et $Y$ les coordonn{\'e}es par rapport aux axes d{\'e}duits des axes ${\rm
O}x$ et ${\rm O}y$ par une rotation de $\pi/4$. Donner les {\'e}quations
param{\'e}triques de la courbe dans ce syst{\`e}me de coordonn{\'e}es. 
\item Tracer la
courbe et discuter de sa forme et du sens de parcours sur celle-ci en fonction du
param{\`e}tre $\varphi \in [0,2\pi]\, $ (consid{\'e}rer les valeurs multiples de
$\pi/2$ et les r{\'e}gions qu'elles d{\'e}limitent). 
\item La courbe {\'e}tant
suppos{\'e}e donn{\'e}e, en d{\'e}duire g{\'e}om{\'e}triquement la valeur de $\varphi$.
\end{itemize}
\finenonce{002697}
\finexercice

\section{245.00 Analyse vectorielle : forme différentielle, champ de vecteurs, circulation}

\section{245.01 Forme différentielle, champ de vecteurs, circulation}


\exercice{2072, maillot, 2001/09/01}

\enonce{002072}{}
On consid\`ere le champ de vecteurs $P:\R^2\rightarrow \R^2$ d\'efini par
$$P(x,y)=(2x {\mathrm e}^{x^2-2y};-2 {\mathrm e}^{x^2-2y}).$$
\begin{enumerate}
\item V\'erifier que la forme diff\'erentielle associ\'ee \`a $P$ est ferm\'ee.
\item En d\'eduire que $P$ est un champ de gradients et en d\'eterminer un
potentiel.
\item Calculer la circulation de $P$ le long du chemin $$\gamma:[0,1]\rightarrow
\R^2;\quad t\mapsto (\ln (1+t);{\mathrm e}^t +1).$$
\end{enumerate}
\finenonce{002072}



\finexercice



\exercice{2073, maillot, 2001/09/01}

\enonce{002073}{}
Soient $a,b$ des nombres tels que $0<a<b$ et soit
 \[D=\{(x,y)\in{(\R_+)}^2\mid a\le xy \le b,\ y\ge x,\ y^2-x^2\le 1\}.\]

En effectuant le changement de variable $u=xy, v=y^2-x^2$, calculer
$$I=\iint_D (y^2-x^2)(x^2+y^2)\, dx\, dy.$$
\finenonce{002073}



\finexercice



\exercice{2074, ridde, 1999/11/01}

\enonce{002074}{}
Soit le champ de vecteurs ${\vec{V}} : {\Rr^2} \to {\Rr^2}$, ${ (x, y)} \mapsto { (2xy + e^y, x^2 + xe^y)}$.
Calculer la circulation de $\vec{V}$ le long de la parabole $x = y^2$ entre les points
$ (0, 0)$ et $ (1, 1)$.
\finenonce{002074}



\finexercice



\exercice{2075, ridde, 1999/11/01}

\enonce{002075}{}

Soit le champ de vecteurs ${\vec{V}}: {\Rr^3} \to {\Rr^3}$, ${ (x, y, z)} \mapsto { (xy, -z, xz)}$.
$\vec{V}$ est-il un champ de gradient ? Calculer la circulation de $\vec{V}$ le
long de l'h\'elice $x = \cos t \,,  \, y = \sin t \,, \, z = t$ pour $t \in [0, 2\pi]$.
\finenonce{002075}



\finexercice



\exercice{2076, ridde, 1999/11/01}

\enonce{002076}{}
Montrer que $\omega (x, y ) = \dfrac{ (1-x^2 + y^2)y}{ (1 + x^2 + y^2)^2}dx
 + \dfrac{ (1 + x^2 - y^2)x}{ (1 + x^2 + y^2)^2}dy$ est une forme diff\'erentielle
 exacte sur $\Rr^2$ et l'int\'egrer.
\finenonce{002076}



\finexercice



\exercice{2077, ridde, 1999/11/01}

\enonce{002077}{}
Sur $D = ]0,  + \infty[^2$ on d\'efinit $\omega (x, y) = (\dfrac x{x + y} + \ln (x^2 + xy))dx
 + \dfrac{\varphi (y)}{x + y}dy$.

\begin{enumerate}
\item Trouver une CNS sur $\varphi$ pour que $\omega$ soit ferm\'ee.
\item Montrer qu'alors $\omega$ est exacte et l'int\'egrer.
\end{enumerate}

\finenonce{002077}



\finexercice



\exercice{2078, ridde, 1999/11/01}

\enonce{002078}{}
\label{exo:diff}
Soit $\omega (x, y, z) = P (x, y, z)dx  + Q (x, y, z)dy + R (x, y, z)dz$ une forme
diff\'erentielle $C^1$ sur un ouvert \'etoil\'e $U$ de $\Rr^3$.

\begin{enumerate}
\item A quelle condition $\omega$ est-elle exacte ?
\item On suppose qu'elle n'est pas exacte et on cherche alors $\lambda : \Rr^3
\rightarrow \Rr^*$ de classe $C^1$ telle que $\lambda \omega$ soit exacte.
On dit alors que $\lambda$ est un facteur int\'egrant.
En \'eliminant $\lambda$ dans la condition trouv\'ee \`a la question pr\'ec\'edente, trouver
une condition n\'ecessaire sur $P, Q, R$ pour qu'il existe un facteur int\'egrant.
\end{enumerate}

\finenonce{002078}



\finexercice



\exercice{2079, ridde, 1999/11/01}

\enonce{002079}{}
Soit $U = \left\{ (x, y, z) \in \Rr^3 / z>0\right\}$ et $\omega (x, y, z) =
2xzdx-2yzdy- (x^2-y^2)dz$.
\begin{enumerate}
\item En utilisant l'exercice pr\'ec\'edent (exercice \ref{exo:diff}), montrer que $\omega$ admet un
facteur int\'egrant.
\item Chercher un facteur int\'egrant ne d\'ependant que de $z$.
\item On suppose qu'un mouvement dans $U$ v\'erifie l'\'equation diff\'erentielle
 $2x (t)z (t)\dot{x} (t)-2y (t)z (t)\dot{y} (t)- (x^2 (t)-y^2 (t))\dot{z} (t)$.
Trouver une int\'egrale premi\`ere du mouvement.
\end{enumerate}
\finenonce{002079}



\finexercice



\exercice{2080, ridde, 1999/11/01}

\enonce{002080}{}
Calculer l'aire d'une astroïde.
\finenonce{002080}



\finexercice



\exercice{2480, matexo1, 2002/02/01}

\enonce{002480}{}
On rappelle que la formule de Stokes g\'en\'erale affirme que si $\omega
$ est une forme diff\'erentielle de degr\'e $p-1$, $\Omega$ une vari\'et\'e de
$\R^N$ de dimension $p$ et de bord  $\partial  \Omega$, alors
$$\int_\Omega d\omega  = \int_{\partial  \Omega} \omega.$$ 
Dans le cas o\`u $p=1$, $\Omega=[a,b]$ un segnment et $\omega  = f$ une
fonction r\'eelle, que donne cette formule\,? Et si $\Omega$ est la r\'eunion
de plusieurs intervalles\,? Plus g\'en\'eralement, si $\Omega$ est une courbe
de $\R^3$, et $g$ une fonction d\'efinie sur $\R^3$, quel est le travail de $g$
le long de $\Omega$\,? Montrer qu'il ne d\'epend pas du chemin parcouru.
\finenonce{002480}



\finexercice


\exercice{2481, matexo1, 2002/02/01}

\enonce{002481}{}
Soit $\Sigma$ une surface de $\R^3$ de bord $\Gamma  = \partial  \Sigma
$, et $\bf U$ un champ de
vecteurs de $\R^3$. En consid\'erant la forme diff\'erentielle $\omega  = {\bf
U}_1 \,dx +
{\bf U}_2 \,dy + {\bf U}_3 \,dz$, montrer la forme vectorielle de la formule de
Stokes\,:
$$\int_\Sigma  \operatorname{rot}{\bf U}\cdot {\bf n} \,d\sigma  = 
\oint_{\Gamma } {\bf U}\cdot d {\bf r},$$
o\`u ${\bf n}$ d\'esigne le vecteur normal \`a $\Sigma$.
Expliciter cette formule dans les cas o\`u $\Sigma $ est donn\'ee\,: a) sous la
forme directe
$z = f(x,y)$, avec $(x,y) \in D \subset \R^2$\,; b) sous la forme intrins\`eque
$f(x,y,z) =
0$\,; c) sous la forme param\'etrique $x= x(u,v), y=y(u,v), z=z(u,v)$, avec
$(u,v)\in
D\subset \R^2$.
\finenonce{002481}



\finexercice


\exercice{2482, matexo1, 2002/02/01}

\enonce{002482}{}
Calculer
$$\oint_\Gamma  (2xy^2+\sin z)\,dx + 2x^2y\,dy + x\cos z \,dz$$
le long de la courbe $\Gamma$ donn\'ee par $x=\cos t, y = z = \sin t$,
$0\leq t <2\pi$.

\finenonce{002482}



\finexercice


\exercice{2483, matexo1, 2002/02/01}

\enonce{002483}{}
Montrer que si $\Sigma $ est une surface ferm\'ee de $\R^3$ et ${\bf U}$
un champ de vecteurs $C^1$ sur $\Sigma $, alors $\int_\Sigma  \operatorname{rot}{\bf
U}\cdot{\bf n} \,d\sigma  = 0$. En d\'eduire la valeur de $\int_S \operatorname{rot}{\bf
U}\cdot{\bf n} \,ds$, o\`u ${\bf U} = (-y^3, x^3+z, z^3)$ et $S$
l'h\'emisph\`ere $z>0$ de la sph\`ere unit\'e.
\finenonce{002483}



\finexercice


\exercice{2484, matexo1, 2002/02/01}

\enonce{002484}{}
Soit ${\bf U} =  (e^x + y^2, -ye^x, x^2+y^2)$\,; calculer $\div
{\bf U}$ et en d\'eduire
$\int_\Sigma  {\bf U}\cdot{\bf n} \,d\sigma $, o\`u $\Sigma $ est une surface
ferm\'ee de $\R^3$.
\finenonce{002484}



\finexercice


\exercice{2485, matexo1, 2002/02/01}

\enonce{002485}{}
Sous les conditions du th\'eor\`eme de Stokes, montrer les identit\'es
suivantes, o\`u $\bf V$ est un champ de vecteur arbitraire,  $\Sigma$ une
surface de bord $\Gamma$, et $\phi $ et $\psi $ des fonctions $C^1$\,:
\begin{eqnarray*}
\oint \phi  \,d{\bf r} 
&=& \int_\Sigma  d{{\bf\sigma }}\wedge\nabla  \varphi,  \\
\oint_{\Gamma } \,d{\bf r}\wedge{\bf U} 
&=& \int_\Sigma  \left({\bf n}\wedge\nabla \right)\wedge {\bf U} \,d\sigma, \\
\oint_{\Gamma } \varphi  \nabla \psi  \cdot d{\bf r} 
&=& \int_\Sigma  \nabla \phi \wedge\nabla \psi  \cdot{\bf n} \,d\sigma.
\end{eqnarray*}

\finenonce{002485}



\finexercice


\exercice{2486, matexo1, 2002/02/01}

\enonce{002486}{}
D\'eduire de la formule g\'en\'erale de Stokes la formule de Green en deux
dimensions\,: si $\Omega\subset \R^2$ de bord $\partial  \Omega$ de classe $C^1$
par morceaux, et si
$P$, $Q$ sont deux fonctions de $\R^2$ dans $\R$ de classe $C^1$ dans
$\Omega$, alors
$$\oint_{\partial \Omega} P(x,y)\,dx+Q(x,y)\,dy = \int_\Omega 
\left({\partial  Q\over \partial  x} - {\partial  P\over \partial  y}\right)
\,dxdy.$$

\finenonce{002486}



\finexercice


\exercice{2487, matexo1, 2002/02/01}

\enonce{002487}{}
Soit $C$ une courbe ferm\'ee du plan, et $P$ et $Q$  deux polyn\^omes de
degr\'e 1 en $x$, $y$\,; montrer que la valeur de $\oint_C 
P(x,y)\,dx+Q(x,y)\,dy$ ne change pas si l'on effectue une
translation sur $C$. En d\'eduire la valeur de $\oint_C
(3x+4y)\,dx+(x-3y)\,dy$, o\`u $C$ est un cercle quelconque de rayon
$a>0$.
\finenonce{002487}



\finexercice


\exercice{2488, matexo1, 2002/02/01}

\enonce{002488}{}
Soit $C$ le bord du carr\'e $[-1,1]\times[-1,1]$. Calculer
$$\oint_C (e^{x^3}+3y^2)\,dx + \sqrt{\cos y} \,dy.$$
\finenonce{002488}



\finexercice


\exercice{2489, matexo1, 2002/02/01}


\enonce{002489}{}

Soit $C$ la cyclo\¨{\i}de d'\'equations $x = t-\sin t, y = 1-\cos t$, et

$C_0$ l'arc de $C$ joignant 

$O=(0,0)$ \`a $A=(\pi ,2)$; calculer $\int_{C_0}

(2x^2+3y^2)\,dx+(6xy+4y^2)\,dy$.

\finenonce{002489}









\finexercice



\exercice{2490, matexo1, 2002/02/01}

\enonce{002490}{}
Soit $C$ une courbe ferm\'ee du plan enclosant une aire $S$, et $a$,
$b$ deux r\'eels.
\begin{enumerate}
\item Calculer $\oint_C ay\,dx+bx\,dy$.
En d\'eduire que 
$$S = \oint_C x\,dy = -\oint_C y\, dx = {1\over2}\oint_C x\,dy-y\,dx. $$
\item En utilisant la formule pr\'ec\'edente, calculer l'aire comprise
entre l'axe $Ox$ et l'arche de la cyclo\"{\i}de d'\'equations 
$x = t-\sin t, y = 1-\cos t$, avec $0\leq t\leq 2\pi$.
\item De m\^eme, trouver l'aire
int\'erieure \`a la boucle du folium de Descartes d'\'equation
$x^3+y^3=3xy$, qui est comprise dans le quadrant $x>0$, $y>0$ (on
pourra chercher une repr\'esentation param\'etrique de la boucle du
folium en posant $y = tx$).
\end{enumerate}
\finenonce{002490}



\finexercice


\exercice{2491, matexo1, 2002/02/01}


\enonce{002491}{}

 D\'eduire de la formule g\'en\'erale de Stokes la formule de Green en trois

dimensions\,: si $V$ est un volume de $\R^3$ de bord $\Sigma  = \partial  V$, et

${\bf U}$ un champ de

vecteur $C^1$ dans V, alors

$$\int_V div{\bf U} \,dv = \oint_\Sigma  {\bf U}\cdot{\bf n} \,d\sigma.$$

En d\'eduire que le volume de $V$ est donn\'e par

$$ |V| = \oint_\Sigma  x \,dy dz = \oint_\Sigma  y \,dz dx = \oint_\Sigma  z

\,dx dy

       = {1\over 3} \oint_\Sigma  x \,dy dz + y \,dz dx + z \,dx dy. $$

\finenonce{002491}









\finexercice



\exercice{2492, matexo1, 2002/02/01}

\enonce{002492}{}
Sous les conditions de la formule de Green, montrer les identit\'es
suivantes, o\`u ${\bf U}$ est un champ de vecteur arbitraire, et $\varphi$ une
fonction $C^1$\,:
\begin{eqnarray*}
\int_V \nabla  \varphi  \,dv &=& \oint_\Sigma  \varphi \,{\bf
n}\,d\sigma, \\ 
\int_V \operatorname{rot}{\bf U} \,dv &=& \oint_\Sigma  {\bf n}\wedge{\bf U}
\,d\sigma.
\end{eqnarray*}
\finenonce{002492}



\finexercice


\exercice{2493, matexo1, 2002/02/01}

\enonce{002493}{}
Soit ${\bf U} = (x/r^3, y/r^3, z/r^3)$\,; calculer directement $I_S =
\oint_S
{\bf U}\cdot{\bf n} \,ds$ dans le cas o\`u $S$ est une boule de rayon $r$ et de
centre $O = (0,0,0)$. Calculer \'egalement $\div{\bf U}$. Que
constatez-vous\,? Expliquer pourquoi la formule de Green ne
s'applique pas. Calculer $I_S$ respectivement dans le cas o\`u $S$ est
une surface ferm\'ee dont l'int\'erieur ne contient pas $O$ (resp.
contient $O$). Que se passe-t-il si $O$ se trouve sur $S$\,?
\finenonce{002493}



\finexercice


\exercice{2685, matexo1, 2002/02/01}
\enonce{002685}{}
 Dans $\R^3$, soit ${\bf U}$ un champ vectoriel arbitraire, ${\bf a}$
un vecteur constant, ${\bf r}$ le champ vectoriel de coordonn\'ees $(x, y, z)$
et $r = \sqrt{x^2+y^2+z^2}$. Calculer $\displaystyle\Delta{1\over r}$\,;
$\operatorname{div}({\bf r}/r^3)$\,; $\operatorname{div}({\bf U}\wedge{\bf r})$\,; 
$\operatorname{rot}(\operatorname{rot}{\bf U})$\,; $\operatorname{rot}({\bf a}\wedge{\bf r})$.
\finenonce{002685}
\finexercice


\exercice{2686, matexo1, 2002/02/01}
\enonce{002686}{}
Soit $a$, $b$, $c$ des constantes et ${\bf U}$ le champ de vecteurs de
$\R^3$ de coordonn\'ees 
$${(x+2y+az,\ bx-3y-z,\ 4x+cy+2z)}.$$
D\'eterminer pour quelles valeurs de $a$, $b$, $c$, ${\bf U}$ est
irrotationnel. Dans ce cas, de quel potentiel d\'erive-t-il\,? M\^emes
questions avec le champ de coordonn\'ees
$$\left({xz r^2\over (ax^2+by^2+c)^{3\over2}},\ {yz r^2\over
(ax^2+by^2+c)^{3\over2}},\ {-r^2\over (ax^2+by^2+c)^{3\over2}} \right),$$
o\`u $r^2 = x^2+y^2+z^2$.
\finenonce{002686}
\finexercice


\exercice{2687, matexo1, 2002/02/01}
\enonce{002687}{}
Un {\em champ \'electromagn\'etique} est caract\'eris\'e par deux champs
de vecteurs ${\bf E}$ et ${\bf H}$, \'egalement fonctions du temps
$t$, et satisfaisant les {\em \'equations de Maxwell}
$$\operatorname{div} {\bf H} = 0,\quad \operatorname{div} {\bf E} = 4\pi \rho ,\quad 
\operatorname{rot} {\bf H} ={1\over c} {\partial {\bf E}\over\partial  t},\quad \operatorname{rot}
{\bf E} =-{1\over c} {\partial {\bf H}\over\partial  t}.
$$ 
o\`u $c$ est une constante et $\rho$ une fonction scalaire de
$(x,y,z,t)$. Montrer qu'il existe un champ de vecteurs ${\bf A}$ tel
que $\operatorname{rot} {\bf A} = {\bf H}$, et un champ scalaire $\phi$  tel que
${\bf E}$ s'exprime en fonction de ${\bf A}$ et $\phi$. Calculer
$\operatorname{div} {\bf A}$ \`a l'aide de $\phi$, et montrer que 
${\bf A}$ et $\phi$ satisfont une \'equation des ondes.
\finenonce{002687}
\finexercice


\exercice{5897, rouget, 2010/10/16}
\enonce{005897}{**}
Les formes différentielles suivantes sont elles exactes ? Si oui, intégrer et si non chercher un facteur intégrant.

\begin{enumerate}
 \item  $\omega= (2x+2y+e^{x+y})(dx+dy)$ sur $\Rr^2$.

\item  $\omega= \frac{xdy-ydx}{(x-y)^2}$ sur $\Omega=\{(x,y)\in\Rr^2/\;y > x\}$

\item  $\omega= \frac{xdx+ydy}{x^2+y^2}- ydy$

\item  $\omega= \frac{1}{x^2y}dx - \frac{1}{xy^2}dy$ sur $(]0,+\infty[)^2$ (trouver un facteur intégrant non nul ne dépendant que de $x^2+y^2$).
\end{enumerate}
\finenonce{005897}


\finexercice


\exercice{5906, rouget, 2010/10/16}
\enonce{005906}{**}
Calculer l' intégrale de la forme différentielle $\omega$ le long du contour orienté $C$ dans les cas suivants :

\begin{enumerate}
 \item  $\omega= \frac{x}{x^2+y^2}dx+ \frac{y}{x^2+y^2}dy$ et $C$ est l'arc de la parabole d'équation $y^2=2x+1$ joignant les points $(0,-1)$ et $(0,1)$
parcouru une fois dans le sens des $y$ croissants.

\item  $\omega=(x-y^3)dx+x^3dy$ et $C$ est le cercle de centre $O$ et de rayon $1$ parcouru une fois dans le sens direct.

\item  $\omega=xyzdx$ et $C$ est l'arc $x=\cos t$, $y=\sin t$, $z=\cos t\sin t$, $t$ variant en croissant de $0$ à $ \frac{\pi}{2}$.
\end{enumerate}
\finenonce{005906}


\finexercice


\exercice{5907, rouget, 2010/10/16}
\enonce{005907}{**}
Soit $\omega=x^2dx+y^2dy$. Calculer l'intégrale de $\omega$ le long de tout cercle du plan parcouru une fois dans le sens trigonométrique. Même question avec $\omega=y^2dx+x^2dy$.

\finenonce{005907}


\finexercice


\exercice{5909, rouget, 2010/10/16}
\enonce{005909}{*** I}
 (Un calcul de $\int_{0}^{+\infty} \frac{\sin x}{x}\;dx$).

\begin{enumerate}
 \item  $r$ et $R$ sont deux réels strictement positifs tels que $r<R$. On considère le contour $\Gamma$ orienté suivant

$$\includegraphics{../images/img005909-1}$$



Calculer l'intégrale de la forme différentielle 

\begin{center}
$\omega= \frac{e^{-y}}{x^2+y^2}((x\sin x-y\cos x)dx+(x\cos x+y\sin x)dy)$
\end{center}

le long de ce contour orienté.

\item  En déduire $\int_{r}^{R} \frac{\sin x}{x}\;dx$ en fonction d'une autre intégrale.

\item  En faisant tendre $r$ vers $0$ et $R$ vers $+\infty$, déterminer la valeur de $\int_{0}^{+\infty} \frac{\sin x}{x}\;dx$.
\end{enumerate}
\finenonce{005909}


\finexercice


\exercice{5913, rouget, 2010/10/16}
\enonce{005913}{**** Inégalité isopérimétrique}
Une courbe fermée $(C)$ est le support d'un arc paramétré $\gamma$ de classe $C^1$ régulier et simple. On note $\mathcal{L}$ sa longueur et $\mathcal{A}$ l'aire délimitée par la courbe fermée $(C)$. Montrer que

\begin{center}
$\mathcal{A}\leqslant \frac{\mathcal{L}^2}{4\pi}$.
\end{center}

Pour cela, on supposera tout d'abord $\mathcal{L}=2\pi$ et on choisira une paramétrisation normale de l'arc. On appliquera ensuite la formule de \textsc{Parseval} aux intégrales permettant de calculer $\mathcal{L}$ et $\mathcal{A}$ et on comparera les sommes des séries obtenues.
\finenonce{005913}


\finexercice


\exercice{6873, gammella, 2012/05/29}

\enonce{006873}{}
Déterminer si les formes différentielles suivantes sont exactes et dans ce cas, les intégrer :
\begin{enumerate}
\item $\omega_1=2xy dx +x^2dy $
\item $\omega_2=xy dx - z dy +xz dz$
\item $\omega_3=2xe^{x^2-y} dx -2e^{x^2-y}dy$
\item $\omega_4=yz^2 dx + (xz^2+z) dy + (2xyz+2z+y) dz.$
\end{enumerate}
\finenonce{006873}


\finexercice


\exercice{6874, gammella, 2012/05/29}

\enonce{006874}{}
On considère le changement de variables en coordonnées sphériques suivant :
$$  \left\{ \begin{array}{lll}
x& = & r\cos \varphi\cos \theta\\
y&= & r\cos \varphi\sin \theta \\
z & = & r \sin \varphi \\
\end{array} \right .$$
\begin{enumerate}
\item Calculer $dx$, $dy$, $dz$.

\item  Vérifier que $x dx+ydy+zdz=rdr.$ En déduire 
$ \frac{\partial r}{\partial x}$, $ \frac{\partial r}{\partial y}$ 
et $ \frac{\partial r}{\partial z}$.
\end{enumerate}
\finenonce{006874}


\finexercice


\exercice{6875, gammella, 2012/05/29}

\enonce{006875}{}
On considère la forme différentielle
$\omega=(x^2+y^2+2x) dx+2y dy$.
\begin{enumerate}
\item
Montrer que $\omega$ n'est pas exacte.
\item
Trouver une fonction $\psi(x)$ telle que $\psi(x) \omega=df$. Préciser
alors $f$. (On dit que $\psi$ est un facteur intégrant.)
\end{enumerate}
\finenonce{006875}


\finexercice


\exercice{6876, gammella, 2012/05/29}

\enonce{006876}{}
On considère le champ vectoriel $\vec{V}(x,y)=(1+2xy,x^3-3)$. 
Ce champ est-il un champ de gradient ?
\finenonce{006876}


\finexercice


\exercice{6877, gammella, 2012/05/29}

\enonce{006877}{}
Quel est le champ vectoriel qui dérive du potentiel
$$U(x,y,z)=1+x+xy+xyz ?$$ 
\finenonce{006877}


\finexercice


\exercice{6878, gammella, 2012/05/29}

\enonce{006878}{}
Calculer la circulation du champ vectoriel
$\vec{V}(x,y)=(3x,x+y)$ le long du cercle $C$ de centre $O$ et de rayon $1$, 
parcouru dans le sens direct.
\finenonce{006878}


\finexercice


\exercice{6879, gammella, 2012/05/29}

\enonce{006879}{}
Calculer le travail $W$ de la force $\vec{F}(x,y,z)=(yz,zx,xy)$ le long de l'hélice $H$ 
paramétrée par $x=\cos t$, $y=\sin t$ et $z=t$ o\`u $t$ varie de $0$ à $ \frac{\pi}{4}$.
\finenonce{006879}


\finexercice


\exercice{6880, gammella, 2012/05/29}

\enonce{006880}{}
On donne le champ vectoriel
$$\vec{V}(x,y,z)=(y^2\cos x, 2y\sin x+e^{2z},2y e^{2z}).$$
\begin{enumerate}
\item Montrer que ce champ est un champ de gradient.
\item Déterminer le potentiel $U(x,y,z)$ dont dérive ce champ sachant qu'il vaut $1$ à l'origine.
\item Quelle est la circulation de ce champ
de $A(0,1,0)$ à $B( \frac{\pi}{2},3 ,0)$ ?
\end{enumerate}
\finenonce{006880}


\finexercice


\exercice{6881, gammella, 2012/05/29}

\enonce{006881}{}
En utilisant la formule de Green-Riemann, calculer
$I= \iint_{\mathcal{D}} xydxdy$
o\`u $$\mathcal{D}=\{(x,y)\in \Rr^2\, |\, x\geq0; y\geq 0;x+y\leq 1\}.$$
\finenonce{006881}


\finexercice


\exercice{6882, gammella, 2012/05/29}

\enonce{006882}{}
 On considère la forme différentielle $$\omega= \frac{-y}{x^2+y^2}dx+\frac{x}{x^2+y^2} dy.$$
 \begin{enumerate}
 \item  Dans quel domaine cette forme différentielle est-elle définie ?
 \item  Calculer l'intégrale curviligne $ \int_{C} \omega$ o\`u $C$ est le cercle de centre $O$ et
 de rayon $1$, parcouru dans le sens direct.
 \item  La forme $\omega$ est-elle exacte ?
  \end{enumerate}
\finenonce{006882}


\finexercice

\section{245.02 Torseurs}


\exercice{4944, quercia, 2010/03/17}
\enonce{004944}{Moment parallèle à un plan}

Soient ${\cal T}$ un torseur et ${\cal P}$ un plan.
Déterminer le lieu des points $M \in {\cal P}$ tels que
${\cal T}(M) \in \vec{{\cal P}}$.

\finenonce{004944}



\finexercice


\exercice{4945, quercia, 2010/03/17}
\enonce{004945}{Torseurs de sommes orthogonales}

Soient ${\cal T},{\cal T}'$ deux torseurs de sommes non nulles, orthogonales.
Montrer que le comoment de ${\cal T}$ et ${\cal T}'$ est nul si et seulement si
les axes centraux sont concourants.
\finenonce{004945}



\finexercice


\exercice{4946, quercia, 2010/03/17}
\enonce{004946}{Somme de glisseurs}

Soit ${\cal R} = (O,\vec i,\vec j,\vec k)$ un repère orthonormé direct de
l'espace. On considère les glisseurs :

\indent ${\cal G}_1$ d'axe $\begin{cases}y=mx\cr z=1\end{cases}$
        et de vecteur $\vec u = \vec i + m\vec j$.

\indent ${\cal G}_2$ d'axe $\begin{cases}y=-mx\cr z=-1\end{cases}$
        et de vecteur $\vec v = \vec i - m\vec j$.


Déterminer l'axe central de ${\cal G}_1 + {\cal G}_2$.
\finenonce{004946}



\finexercice


\exercice{4947, quercia, 2010/03/17}
\enonce{004947}{Glisseurs associés à un tétraèdre}

Soit $ABCD$ un tétraèdre non aplati de l'espace.
Pour $X,Y \in \{A,B,C,D\}$ distincts, on note
${\cal G}_{XY}$ le glisseur d'axe la droite $(XY)$ et de vecteur $\vec{XY}$.

Montrer que $( {\cal G}_{AB},
               {\cal G}_{AC},
               {\cal G}_{AD},
               {\cal G}_{BC},
               {\cal G}_{BD},
               {\cal G}_{CD} )$ est une base de l'espace des torseurs.
\finenonce{004947}



\finexercice


\exercice{4948, quercia, 2010/03/17}
\enonce{004948}{Produit vectoriel de torseurs}

Soient ${\cal T}_1$, ${\cal T}_2$ deux torseurs de sommes
$\vec{R}_1$, $\vec{R}_2$.
On définit le champ $\cal T$ par :
$${\cal T}(M) = \vec{R}_1 \wedge {\cal T}_2(M) + {\cal T}_1(M)\wedge \vec{R}_2.$$

\begin{enumerate}
  \item Montrer que $\cal T$ est un torseur de somme $\vec R_1 \wedge \vec R_2$
    (produit vectoriel de ${\cal T}_1$ et ${\cal T}_2$).
  \item Si $\vec R_1 \wedge \vec R_2 \ne \vec 0$, montrer que l'axe central de
    $\cal T$ est la perpendiculaire commune des axes centraux de ${\cal T}_1$ et
    ${\cal T}_2$.
\end{enumerate}
\finenonce{004948}



\finexercice

\section{246.00 Autre}

\section{246.01 Plan tangent, vecteur normal}


\exercice{5915, rouget, 2010/10/16}
\enonce{005915}{}
 Soit $\mathcal{S}$ la surface d'équation $x^4-x^3+xy-y^2-z=0$.

\begin{enumerate}
 \item   Déterminer les plans tangents à la surface $\mathcal{S}$ parallèle au plan $\left(O,\overrightarrow{i},\overrightarrow{j}\right)$.

\item  Etudier localement la position relative de la surface $\mathcal{S}$ et de son plan tangent en chacun des points ainsi obtenu.

\item  Etudier la position relative globale de la surface $\mathcal{S}$ et du plan $\left(O,\overrightarrow{i},\overrightarrow{j}\right)$.
\end{enumerate}
\finenonce{005915}



\finexercice


\exercice{5916, rouget, 2010/10/16}
\enonce{005916}{}
Trouver toutes les droites tracées sur la surface d'équations $x^3+y^3+z^3=1$ puis vérifier que ces droites sont coplanaires.
\finenonce{005916}



\finexercice

\section{246.02 Surfaces paramétrées}


\exercice{5047, quercia, 2010/03/17}
\enonce{005047}{Chimie P 91}

    \'Equation de la surface de révolution engendrée par la rotation de $\Gamma$
    autour de $Oz$ où $\Gamma$ est la courbe d'équations paramétriques :
    $\begin{cases}x=a\cos^3u\cr y=a\sin^3u\cr z=a\cos2u.\cr\end{cases}$ ($a > 0$)
    
\finenonce{005047}



\finexercice


\exercice{5048, quercia, 2010/03/17}
\enonce{005048}{Ensi Physique 93}
Soit la courbe d'équations dans $\R^3$ :
$$(\Gamma)\quad \left\{\begin{array}{cc} x^2-y^2-4x+2&=0 \cr x+z&=1.\cr\end{array}\right.$$\par
Déterminer la surface engendrée par la rotation de $(\Gamma)$ autour de $Oz$.
\finenonce{005048}



\finexercice


\exercice{5049, quercia, 2010/03/17}
\enonce{005049}{Le plan tangent coupe $Oz$ en un point fixe}

On considère la surface ${\cal S}$ d'équations paramétriques :
$\begin{cases} x = \rho\cos\theta\cr
         y = \rho\sin\theta\cr
         z = f(\rho,\theta)\cr \end{cases}$
où $f$ est une fonction de classe $\mathcal{C}^1$.

\begin{enumerate}
  \item Donner l'équation du plan tangent à ${\cal S}$ en un point $M(\rho,\theta)$.
    
  \item Déterminer $f$ de sorte que, le long d'une ligne $\theta = \text{cste}$,
    le plan tangent coupe $Oz$ en un point fixe.
    

  \item Exemple : $f(\rho,\theta) = \theta$. Dessiner la surface ${\cal S}$.

\end{enumerate}
\finenonce{005049}



\finexercice


\exercice{5050, quercia, 2010/03/17}

\enonce{005050}{Pseudo-sphère}

Dessiner la surface ${\cal S}$ d'équations paramétriques :
$\begin{cases} x = a\cos u/\ch v\cr
         y = a\sin u/\ch v\cr
         z = a(v-\tanh v)\cr \end{cases}$
où $a$ est un réel strictement positif (pseudo-sphère).
\finenonce{005050}



\finexercice


\exercice{5051, quercia, 2010/03/17}
\enonce{005051}{Les normales coupent $Oz \Leftrightarrow$ révolution}

\label{revolution}
Soit ${\cal S}$ une surface d'équation $z = f(x,y)$.
Montrer que ${\cal S}$ est de révolution si et seulement si en tout point $M$,
la normale à ${\cal S}$ en $M$ est parallèle ou sécante à $Oz$.

\finenonce{005051}



\finexercice


\exercice{5052, quercia, 2010/03/17}
\enonce{005052}{}
Que dire d'une surface ${\cal S}$ telle que toutes les normales sont
concourantes ? (cf ex.\ref{revolution})
\finenonce{005052}



\finexercice


\exercice{5053, quercia, 2010/03/17}
\enonce{005053}{Contour apparent}

Soit ${\cal S}$ la surface d'équation cartésienne $z^2-x^2-y^2 = 1$.

\begin{enumerate}
  \item Reconnaître ${\cal S}$.
    

  \item Soit $D$ la droite d'équations : $2x+y = 0$, $z=0$. Déterminer les points
    $M$ de ${\cal S}$ tels que le plan tangent à ${\cal S}$ en $M$ est
    parallèle à $D$.
    (Contour apparent de ${\cal S}$ dans la direction de $D$)
    
\end{enumerate}
\finenonce{005053}



\finexercice


\exercice{5054, quercia, 2010/03/17}
\enonce{005054}{Cylindre circonscrit}

Soit ${\cal S}$ la surface d'équations paramétriques :
$\begin{cases} x = u/(u^2+v^2)\cr
         y = v/(u^2+v^2)\cr
         z = 1/(u^2+v^2).\cr \end{cases}$

\begin{enumerate}
  \item Donner une équation cartésienne de ${\cal S}$.
    

  \item Déterminer l'ensemble $\mathcal{C}$ des points de ${\cal S}$ où le plan tangent est
    parallèle à la droite $D$ d'équations : ${x=y=z}$.
    

  \item Déterminer l'équation cartésienne du cylindre de génératrices parallèles à
    $D$ s'appuyant sur $\mathcal{C}$. (Cylindre circonscrit à ${\cal S}$)
    
\end{enumerate}
\finenonce{005054}



\finexercice


\exercice{5055, quercia, 2010/03/17}
\enonce{005055}{\'Equation de cône}

Soit $\mathcal{C}$ le cercle intersection de la sphère d'équation $x^2+y^2+z^2 = 1$
et du plan d'équation $x+y=1$, et $S=(1,1,1)$.
Déterminer l'équation cartésienne du cône de sommet $S$ s'appuyant sur $\mathcal{C}$.

\finenonce{005055}



\finexercice


\exercice{5056, quercia, 2010/03/17}
\enonce{005056}{Cône $=$ cylindre ?}

Soit ${\cal S}$ la surface d'équation cartésienne :
$\frac1{(x-y)^2} + \frac1{(y-z)^2} = \frac1{(x-z)^2}$.

\begin{enumerate}
  \item Montrer que ${\cal S}$ est à la fois un cylindre et un cône.
  \item Comment est-ce possible ?
\end{enumerate}
\finenonce{005056}



\finexercice


\exercice{5057, quercia, 2010/03/17}
\enonce{005057}{Position d'une surface de révolution  par rapport au plan tangent}

Soit ${\cal S}$ une surface d'équation cartésienne $z=f(\rho)$ où
$\rho = \sqrt{x^2+y^2}$ et $f$ est une fonction de classe $\mathcal{C}^2$.
Montrer que la position de ${\cal S}$ par rapport à son plan tangent est donnée
par le signe de $f'(\rho)f''(\rho)$. Interpréter géométriquement ce fait.
\finenonce{005057}



\finexercice


\exercice{5058, quercia, 2010/03/17}
\enonce{005058}{Intersection de deux cylindres}

Soient ${\cal S}_1$, ${\cal S}_2$ les surfaces d'équations
$x^2 + y^2 + xy = 1$ et $y^2 + z^2 + yz = 1$,
et $\mathcal{C} = {\cal S}_1 \cap {\cal S}_2$.

\begin{enumerate}
  \item Donner en tout point de $\mathcal{C}$ le vecteur tangent à $\mathcal{C}$.
    

  \item Montrer que $\mathcal{C}$ est la réunion de deux courbes planes.
    

  \item Quelle est la projection de $\mathcal{C}$ sur $Oxz$ ?
    
\end{enumerate}
\finenonce{005058}



\finexercice


\exercice{5059, quercia, 2010/03/17}
\enonce{005059}{Conoïde}

Soit ${\cal S}$ la sphère de centre $A=(a,0,0)$ et de rayon $r$ ($0 < r < a$) et
${\cal S}'$ la surface constituée des droites horizontales tangentes à ${\cal S}$
et sécantes à $Oz$. Déterminer l'équation cartésienne de ${\cal S}'$.

\finenonce{005059}



\finexercice


\exercice{5060, quercia, 2010/03/17}
\enonce{005060}{Surface cerclée}

Soit $A=(0,1,0)$ et ${\cal S}$ la surface constituée des cercles verticaux de
diamètre $[A,B]$ où $B$ est un point variable sur $Ox$. Chercher une
équation cartésienne de ${\cal S}$.
\finenonce{005060}



\finexercice


\exercice{5061, quercia, 2010/03/17}
\enonce{005061}{Chimie P' 91}
    On considère la droite $\Delta$ d'équations : $x=a$, $z=0$.
    $P$ est un point décrivant $\Delta$ et $\mathcal{C}_P$ le cercle tangent à $Oz$ en
    $O$ et passant par $P$.
    Faire un schéma et paramétrer la surface engendrée par les cercles $\mathcal{C}_P$
    quand $P$ décrit $\Delta$.
    
\finenonce{005061}



\finexercice


\exercice{5062, quercia, 2010/03/17}

\enonce{005062}{Ensi Chimie P' 93}
Soit $(\Gamma)$ : $\begin{cases} x(t) = a\cos(t)/\ch(mt)\cr
                          y(t) = a\sin(t)/\ch(mt)\cr
                          z(t) = a\tanh(mt).\cr\end{cases}$
\begin{enumerate}
  \item Montrer que $(\Gamma)$ est tracée sur une surface $(\Sigma)$ simple. Montrer que $(\Sigma)$
    est de révolution autour de $Oz$ et donner son équation.
  \item Montrer que $(\Gamma)$ coupe les méridiennes de $(\Sigma)$ suivant un angle constant (loxodromie).
  \item Réciproquement, déterminer toutes les loxodromies de $(\Sigma)$.
  \item Dessiner la projection de $(\Gamma)$ sur $xOy$.
\end{enumerate}
\finenonce{005062}





\finexercice


\section{260.01 Probabilité et dénombrement}


\exercice{5983, quinio, 2011/05/18}

\enonce{005983}{}
Une entreprise décide de classer $20$ personnes
susceptibles d'être embauchées; leurs CV étant très proches,
le patron décide de recourir au hasard : combien y-a-il de classements
possibles : sans ex-aequo; avec exactement $2$ ex-aequo ?
\finenonce{005983}



\finexercice


\exercice{5984, quinio, 2011/05/18}

\enonce{005984}{}
Un étudiant s'habille très vite le matin et prend, au hasard
dans la pile d'habits, un pantalon, un tee-shirt, une paire de chaussettes;
il y a ce jour-là dans l'armoire $5$ pantalons dont $2$ noirs, 
$6$ tee-shirt dont $4$ noirs, $8$ paires de chaussettes, dont $5$ paires
noires. Combien y-a-t-il de façons de s'habiller? Quelles sont les
probabilités des événements suivants : il est tout en noir; une
seule pièce est noire sur les trois.
\finenonce{005984}


\finexercice


\exercice{5985, quinio, 2011/05/18}

\enonce{005985}{}
Si $30$ personnes sont présentes à un réveillon et si, à minuit, chaque personne fait $2$ bises à toutes les autres,
combien de bises se sont-elles échangées en tout ? (On appelle bise un contact entre deux joues...)
\finenonce{005985}


\finexercice


\exercice{5986, quinio, 2011/05/18}

\enonce{005986}{}
Un QCM comporte $10$ questions, pour chacune desquelles $4$ réponses sont proposées, une seule est exacte. Combien y-a-t-il de
grilles-réponses possibles? Quelle est la probabilité de répondre au hasard au moins $6$ fois correctement?
\finenonce{005986}


\finexercice


\exercice{5987, quinio, 2011/05/18}

\enonce{005987}{}
 Amédée, Barnabé, Charles tirent sur un oiseau; si les
probabilités de succès sont pour Amédée : $70$\%, Barnabé : $50$\%, Charles : $90$\%, quelle est la probabilité que l'oiseau soit
touché?
\finenonce{005987}


\finexercice


\exercice{5988, quinio, 2011/05/18}

\enonce{005988}{}
Lors d'une loterie de Noël, $300$ billets sont
vendus aux enfants de l'école ; $4$ billets sont gagnants.
J'achète $10$ billets, quelle est la probabilité pour que je gagne au moins un lot?
\finenonce{005988}


\finexercice


\exercice{5989, quinio, 2011/05/18}

\enonce{005989}{}
La probabilité pour une population d'être atteinte
d'une maladie $A$ est $p$ donné; dans cette même population, un individu
peut être atteint par une maladie $B$ avec une probabilité $q$ donnée aussi; 
on suppose que les maladies sont indépendantes : quelle est la
probabilité d'être atteint par l'une et l'autre de ces maladies?
Quelle est la probabilité d'être atteint par l'une ou l'autre de ces
maladies?
\finenonce{005989}


\finexercice


\exercice{5990, quinio, 2011/05/18}
\enonce{005990}{}
Dans un jeu de $52$ cartes, on prend une carte au hasard : les événements <<tirer un roi>> et 
<<tirer un pique>> sont-ils indépendants? quelle est la probabilité de <<tirer un roi ou
un pique>> ? 
\finenonce{005990}


\finexercice


\exercice{5991, quinio, 2011/05/18}

\enonce{005991}{}
La famille Potter comporte $2$ enfants; les événements $A$ : 
<<il y a deux enfants de sexes différents chez les Potter>> 
et $B$ : <<la famille Potter a au
plus une fille>> sont-ils indépendants? 
Même question si la famille Potter comporte $3$ enfants. Généraliser.
\finenonce{005991}


\finexercice


\exercice{6889, ruette, 2013/01/24}
\enonce{006889}{}
On lance deux dés à 6 faces. Décrire l'ensemble $\Omega$ des résultats
possibles et la probabilité $P$ associée à cette expérience. Donner la
probabilité d'obtenir:
\begin{enumerate}
\item un double,
\item au plus un nombre pair,
\item exactement un nombre pair,
\item deux nombres qui se suivent.
\end{enumerate}
\finenonce{006889}

\finexercice


\exercice{6890, ruette, 2013/01/24}
\enonce{006890}{}
Au loto, on choisit 6 numéros principaux, qui sont 6 nombres différents 
entre 1 et 49, et un numéro 
complémentaire, qui est un nombre entre 1 et 49 différent des 6 précédents. 
Quelle est la probabilité d'avoir :
\begin{enumerate}
\item
les six bons numéros principaux ?
\item cinq bons numéros parmi les 6 principaux?
\item cinq bons numéros parmi les principaux 
et le bon numéro complémentaire ?
\end{enumerate}
\finenonce{006890}
\finexercice


\exercice{6891, ruette, 2013/01/24}
\enonce{006891}{}
Une urne contient une boule rouge, trois boules vertes et seize boules blanches. 
La boule rouge permet de gagner 10 euros, chaque boule verte permet de gagner 5 euros 
et les boules blanches ne rapportent rien. Un joueur tire simultanément cinq boules. 
Quelle est la probabilité pour que ce joueur gagne exactement 10 euros ?
\finenonce{006891}

\finexercice


\exercice{6892, ruette, 2013/01/24}
\enonce{006892}{}
On lance trois dés non pipés. On note le nombre de points (1, 2, 3, 4, 5 ou 6) qui 
apparaît sur la face supérieure de chaque dé. Calculer la probabilité d'avoir :
\begin{enumerate}
\item trois 3,
\item deux 2 et un 1,
\item un 1, un 3, un 5,
\item la somme des points égale à 9,
\item la somme des points égale à 10.
\end{enumerate}
\textit{Remarque : Ces calculs ont été effectués à l'origine par Galilée pour expliquer 
la différence entre 4. et 5.}
\finenonce{006892}

\finexercice


\exercice{6893, ruette, 2013/01/24}
\enonce{006893}{}
Les trois mousquetaires (donc quatre personnes) ont mélangé leurs bottes dans 
le couloir de l'auberge. D'Artagnan se lève le premier et prend deux bottes 
au hasard. Calculer la probabilité pour que :
\begin{enumerate}
\item 
Les deux bottes soient les siennes.
\item
Les deux bottes forment une paire (une paire est la réunion d'un pied droit et d'un pied gauche).
\item
Les deux bottes soient deux pieds droits.
\item
Les deux bottes appartiennent à deux personnes différentes.
\end{enumerate}
\finenonce{006893}
\finexercice


\exercice{6894, ruette, 2013/01/24}
\enonce{006894}{} 
On tire simultanément 6 cartes dans un jeu de 32 cartes.
Quelle est la probabilité d'avoir exactement 2 dames et 3 trèfles ?
\finenonce{006894}
\finexercice


\exercice{6895, ruette, 2013/01/24}
\enonce{006895}{}
On veut transmettre un message électronique composé des digits 0 et 1. Les 
conditions imparfaites de transmission font en sorte qu'il y a une probabilité 
égale à  0,1 qu'un 0 soit changé en un  1  et un  1  en un  0  lors de la 
réception, et ce de façon indépendante pour chaque digit. Pour améliorer 
la qualité de la transmission, on propose d'émettre le bloc  00000 au lieu de  
0  et le bloc  11111  au lieu de  1  et de traduire une majorité de  0  dans 
un bloc lors de la réception par  0  et une majorité de  1  par  1.
\begin{enumerate}
\item 
Quelle est la probabilité de recevoir une majorité de 1 si 00000 est émis ?
\item
Quelle est la probabilité de recevoir une majorité de 1 si 11111 est émis ?
\end{enumerate}
\finenonce{006895}

\finexercice

\section{260.02 Probabilité conditionnelle}


\exercice{5992, quinio, 2011/05/11}

\enonce{005992}{}
Dans la salle des profs $60$\% sont des femmes ; une femme sur trois
porte des lunettes et un homme sur deux porte des lunettes : quelle est la
probabilité pour qu'un porteur de lunettes pris au hasard soit une femme?
\finenonce{005992}


\finexercice

\exercice{5993, quinio, 2011/05/20}

\enonce{005993}{}
Une fête réunit $35$ hommes, $40$ femmes, $25$ enfants ; sur
une table, il y a $3$ urnes $H$, $F$, $E$ contenant des boules de couleurs dont
respectivement $10$\%, $40$\%, $80$\% de boules noires. Un présentateur
aux yeux bandés désigne une personne au hasard et lui demande de
tirer une boule dans l'urne $H$ si cette personne est un homme, dans l'urne $F$
si cette personne est une femme, dans l'urne $E$ si cette personne est un
enfant. La boule tirée est noire : quelle est la probabilité pour que
la boule ait été tirée par un homme? une femme? un enfant? Le présentateur n'est pas plus magicien que vous et moi et pronostique le
genre de la personne au hasard : que doit-il dire pour avoir le moins de
risque d'erreur?
\finenonce{005993}


\finexercice


\exercice{5994, quinio, 2011/05/20}

\enonce{005994}{}
Un fumeur, après avoir lu une série de statistiques
effrayantes sur les risques de cancer, problèmes cardio-vasculaires 
liés au tabac, décide d'arrêter de fumer; toujours d'après des
statistiques, on estime les probabilités suivantes : si cette personne
n'a pas fumé un jour $J_{n}$, alors la probabilité
pour qu'elle ne fume pas le jour suivant $J_{n+1}$ est $0.3$; 
mais si elle a fumé un jour $J_{n}$, alors la probabilité 
pour qu'elle ne fume pas le jour suivant $J_{n+1}$ est $0.9$; 
quelle est la probabilité $P_{n+1}$ pour qu'elle
fume le jour $J_{n+1}$ en fonction de la probabilité 
$P_{n}$ pour qu'elle fume le jour $J_{n}$ ? Quelle est la
limite de $P_{n}$ ? Va-t-il finir par s'arrêter?
\finenonce{005994}


\finexercice


\exercice{5995, quinio, 2011/05/20}

\enonce{005995}{}
Un professeur oublie fréquemment ses clés. Pour tout $n$, on note :
$E_n$ l'événement <<le jour $n$, le professeur oublie ses clés>>, 
$P_{n}=P(E_n)$, $Q_n=P(\overline{E_n})$.

On suppose que : $P_{1}=a$ est donné et que si le jour $n$ il oublie ses clés, 
le jour suivant il les oublie avec la probabilité $\frac{1}{10}$ ; 
si le jour $n$ il n'oublie pas ses clés, le jour suivant il les oublie
avec la probabilité $\frac{4}{10}$.

Montrer que $P_{n+1}=\frac{1}{10}P_{n}+\frac{4}{10}Q_{n}$.
En déduire une relation entre $P_{n+1}$ et $P_{n}$

Quelle est la probabilité de l'événement <<le jour $n$, le professeur oublie ses clés>> ?
\finenonce{005995}


\finexercice


\exercice{5996, quinio, 2011/05/20}

\enonce{005996}{}
Dans les barres de chocolat N., on trouve des images équitablement
réparties des cinq personnages du dernier Walt Disney, une image par
tablette. Ma fille veut avoir le héros Princecharmant : combien dois-je
acheter de barres pour que la probabilité d'avoir la figurine attendue dépasse $80$\%? 
Même question pour être sûr à $90$\%.
\finenonce{005996}


\finexercice


\exercice{5997, quinio, 2011/05/20}

\enonce{005997}{}
En cas de migraine trois patients sur cinq prennent de l'aspirine
(ou équivalent), deux sur cinq prennent un médicament M présentant des effets secondaires :

Avec l'aspirine, 75\% des patients sont soulagés.

Avec le médicament M, 90\% des patients sont soulagés.
\begin{enumerate}
\item Quel est le taux global de personnes soulagées?

\item Quel est la probabilité pour un patient d'avoir pris de l'aspirine
sachant qu'il est soulagé?
\end{enumerate}
\finenonce{005997}


\finexercice


\exercice{5998, quinio, 2011/05/20}

\enonce{005998}{}
Dans une population 40\% des individus ont les yeux bruns, 25\% des
individus ont les cheveux blonds, 15\% des individus ont les yeux bruns
et les cheveux blonds.

On choisit un individu au hasard. Calculez :
\begin{enumerate}
\item La probabilité de l'événement : si un individu a les yeux
bruns d'avoir les cheveux blonds.

\item La probabilité de l'événement : si un individu a les cheveux
blonds d'avoir les yeux bruns.

\item La probabilité de l'événement : si un individu a les cheveux
blonds, de ne pas avoir les yeux bruns.
\end{enumerate}
\finenonce{005998}


\finexercice


\exercice{5999, quinio, 2011/05/20}

\enonce{005999}{}
Un constructeur aéronautique équipe ses avions trimoteurs d'un
moteur central de type A et de deux moteurs, un par aile, de type B; chaque
moteur tombe en panne indépendamment d'un autre, et on estime à $p$ la
probabilité pour un moteur de type A de tomber en panne et à $q$ la
probabilité pour un moteur de type B de tomber en panne.

Le trimoteur peut voler si le moteur central \emph{ou} les deux moteurs d'ailes
fonctionnent : quelle est la probabilité pour l'avion de voler?
Application numérique : $p = 0.001\%$, $q = 0.02\%$.
\finenonce{005999}


\finexercice


\exercice{6000, quinio, 2011/05/20}

\enonce{006000}{}
On sait qu'à une date donnée, 3\% d'une population
est atteinte d'hépatite
On dispose de tests de dépistage de la maladie :
\begin{itemize}
\item Si la personne est malade, alors le test est positif avec une probabilité de 95\%.
\item Si la personne est saine, alors le test est positif avec une probabilité de 10\%.
\end{itemize}

\begin{enumerate}
\item Quelle est la probabilité pour une personne d'être malade
si son test est positif ?
\item Quelle est la probabilité pour une personne d'être saine si
son test est positif ?
\item Quelle est la probabilité pour une personne d'être malade
si son test est négatif ?
\item Quelle est la probabilité pour une personne d'être saine si
son test est négatif ?
\end{enumerate}
\finenonce{006000}


\finexercice


\exercice{6001, quinio, 2011/05/20}

\enonce{006001}{}
Dans mon trousseau de clés il y a $8$ clés; elles sont toutes
semblables. Pour rentrer chez moi je mets une clé au hasard; je
fais ainsi des essais jusqu'à ce que je trouve la bonne; j'écarte au
fur et à mesure les mauvaises clés. Quelle est la probabilité
pour que j'ouvre la porte : 
\begin{enumerate}
\item du premier coup ? 
\item au troisième essai ? 
\item au cinquième essai ? 
\item au huitième essai?
\end{enumerate}
\finenonce{006001}


\finexercice


\exercice{6002, quinio, 2011/05/20}

\enonce{006002}{}
Six couples sont réunis dans une soirée de réveillon. Une
fois les bises de bonne année échangées, on danse, de façon
conventionnelle: un homme avec une femme, mais pas forcément la sienne.

\begin{enumerate}
\item Quelle est la probabilité $P(A)$ pour que chacun des 6 hommes danse
avec son épouse légitime ?
\item Quelle est la probabilité $P(B)$ pour que André danse avec son 
épouse ?
\item Quelle est la probabilité $P(C)$ pour que André et René
dansent avec leur épouse ?
\item Quelle est la probabilité $P(D)$ pour que André ou René
danse(nt) avec leur épouse ?
\end{enumerate}
\finenonce{006002}


\finexercice


\exercice{6003, quinio, 2011/05/20}

\enonce{006003}{}
Dans l'ancienne formule du Loto il fallait choisir 6 numéros parmi 49.
\begin{enumerate}
  \item Combien y-a-t-il de grilles possibles ? En déduire la probabilité de
gagner en jouant une grille.

  \item Quelle est la probabilité que la grille gagnante comporte 2
nombres consécutifs?
\end{enumerate}
\finenonce{006003}


\finexercice


\exercice{6004, quinio, 2011/05/20}

\enonce{006004}{}
Un débutant à un jeu effectue plusieurs parties successives. Pour la
première partie, les probabilités de gagner ou perdre sont les mêmes; puis, on suppose que:
\begin{itemize}
  \item Si une partie est gagnée, la probabilité de gagner la suivante est $0.6$.

  \item Si une partie est perdue, la probabilité de perdre la suivante est $0.7$.
\end{itemize}
Soit $G_n$ l'événement <<Gagner la partie $n$>>, et $u_n=P(G_n)$.
On note $v_n = P(\overline{G_n})$.
\begin{enumerate}
  \item Ecrire 2 relations entre $u_n$, $u_{n+1}$, $v_n$, $v_{n+1}$.

  \item A l'aide de la matrice mise en évidence en déduire $u_n$ et $v_n$.
Faire un calcul direct à l'aide de $u_n+v_n$.
\end{enumerate}
\finenonce{006004}


\finexercice


\exercice{6896, ruette, 2013/01/24}
\enonce{006896}{} 
Aurore arrive en retard en cours avec une probabilité $1/2$. Elle ne va
pas en cours avec une probabilité $1/6$. Aujourd'hui, le cours commence sans
elle. Quelle est la probabilité qu'elle vienne aujourd'hui ?
\finenonce{006896}
\finexercice


\exercice{6897, ruette, 2013/01/24}
\enonce{006897}{}
On a décelé dans une certaine population une probabilité de 0,01 pour
qu'un enfant soit atteint par une maladie M. La 
probabilité qu'un enfant qui n'est pas atteint par M ait une réaction 
négative à un test T est de 0,9. S'il est atteint par M, la probabilité 
qu'il ait une réaction positive au test est de 0,95.\\ 
Quelle est la probabilité  qu'un enfant pris
au hasard ait une réaction positive au test~?
Quelle est la probabilité
qu'un enfant pris au hasard et ayant une réaction 
positive soit atteint par M~?
\finenonce{006897}

\finexercice


\exercice{6898, ruette, 2013/01/24}
\enonce{006898}{}
Les ampoules de la marque X sont fabriquées dans deux usines, A et
B. 20$\%$ des ampoules de l'usine A et 5$\%$ de l'usine B sont
défectueuses. Chaque semaine l'usine A produit $2n$ ampoules et
l'usine B produit $n$ ampoules (où $n\ge 1$ est un entier).
On tire une une ampoule au hasard dans la production d'une semaine.
\begin{enumerate}
\item
Quelle est la probabilité que l'ampoule tirée ne soit pas 
défectueuse~?
\item  Si l'ampoule tirée est défectueuse, quelle est la probabilité 
 qu'elle vienne de l'usine A~? 
\end{enumerate}
\finenonce{006898}
\finexercice


\exercice{6899, ruette, 2013/01/24}
\enonce{006899}{}
Une personne lance deux dés à 6 faces et dit qu'elle a obtenu au moins
un nombre pair. Quelle est la probabilité que les deux nombres obtenus
soient pairs ?
\finenonce{006899}
\finexercice


\exercice{6900, ruette, 2013/01/24}
\enonce{006900}{}
Il y a $5\%$ de daltoniens chez les hommes et $0,25\%$
chez les femmes. Il y a $48\%$ d'hommes et
$52\%$ de femmes dans la population. Quelle est la probabilité pour
qu'un daltonien soit un homme ?

\medskip
\textit{Remarque : la forme la plus courante du daltonisme est génétique, due
à un gène récessif porté par le chromosome X. Un homme (XY) est daltonien
dès que le chromosome X porte ce gène. Une femme (XX) n'est daltoniene
que si les 2 chromosomes X portent ce gène. Ceci explique les taux
très différents chez les hommes et les femmes.}
\finenonce{006900}
\finexercice


\exercice{6901, ruette, 2013/01/24}
\enonce{006901}{}
Deux urnes sont remplies de boules. La première contient
10 boules noires et 30 boules blanches. La seconde contient 20
boules noires et 20 boules blanches.
On tire une des urnes au hasard, de façon équiprobable, et dans cette urne, 
on tire une boule au hasard. La boule est blanche. Quelle est la probabilité 
qu'on ait tiré cette boule dans la première urne sachant qu'elle est blanche ?
\finenonce{006901}
\finexercice
\section{260.03 Variable aléatoire discrète}


\exercice{6005, quinio, 2011/05/20}

\enonce{006005}{}
Une entreprise pharmaceutique décide de faire des économies sur les
tarifs d'affranchissements des courriers publicitaires à envoyer aux
clients.
Pour cela, elle décide d'affranchir, au hasard, une proportion de 3
lettres sur 5 au tarif urgent, les autres au tarif normal.
\begin{enumerate}
\item Quatre lettres sont envoyées dans un cabinet médical de quatre médecins:
quelle est la probabilité des événements:

A : <<Au moins l'un d'entre eux reçoit une lettre au tarif urgent>>.

B : <<Exactement 2 médecins sur les quatre reçoivent une lettre au tarif urgent>>.

\item Soit $X$ la variable aléatoire: <<nombre de lettres affranchies au tarif
urgent parmi 10 lettres>>:
Quelle est la loi de probabilité de $X$, quelle est son espérance,
quelle est sa variance?
\end{enumerate}
\finenonce{006005}


\finexercice


\exercice{6006, quinio, 2011/05/20}

\enonce{006006}{}
On prend au hasard, en même temps, trois ampoules dans un lot de 15 dont
5 sont défectueuses. Calculer la probabilité des événements:

$A$ : au moins une ampoule est défectueuse;

$B$ : les 3 ampoules sont défectueuses;

$C$ : exactement une ampoule est défectueuse.
\finenonce{006006}


\finexercice


\exercice{6007, quinio, 2011/05/20}

\enonce{006007}{}
Un avion peut accueillir 20 personnes; des statistiques montrent 
que 25\% clients ayant réservé ne viennent pas.
Soit $X$ la variable aléatoire: <<nombre de clients qui viennent après
réservation parmi 20>>.
Quelle est la loi de $X$ ? (on ne donnera que la forme générale)
quelle est son espérance, son écart-type ?
Quelle est la probabilité pour que $X$ soit égal à 15 ?
\finenonce{006007}


\finexercice


\exercice{6008, quinio, 2011/05/20}

\enonce{006008}{}
L'oral d'un concours comporte au total 100 sujets; les candidats
tirent au sort trois sujets et choisissent alors le sujet traité parmi
ces trois sujets. Un candidat se présente en ayant révisé 60
sujets sur les 100.
\begin{enumerate}
\item Quelle est la probabilité pour que le candidat ait révisé:
\begin{enumerate}
\item les trois sujets tirés;
\item exactement deux sujets sur les trois sujets;
\item aucun des trois sujets.
\end{enumerate}
\item Définir une variable aléatoire associée à ce problème 
et donner sa loi de probabilité, son espérance.
\end{enumerate}
\finenonce{006008}


\finexercice


\exercice{6009, quinio, 2011/05/20}
\enonce{006009}{}
Un candidat se présente à un concours où, cette fois, les
20 questions sont données sous forme de QCM. A chaque question, sont
proposées 4 réponses, une seule étant exacte. L'examinateur fait
le compte des réponses exactes données par les candidats.
Certains candidats répondent au hasard à chaque question; pour
ceux-la, définir une variable aléatoire associée à ce problème 
et donner sa loi de probabilité, son espérance.
\finenonce{006009}


\finexercice


\exercice{6010, quinio, 2011/05/20}

\enonce{006010}{}
Dans une poste d'un petit village, on remarque qu'entre 10 heures
et 11 heures, la probabilité pour que deux personnes entrent durant 
la même minute est considérée comme nulle et que l'arrivée des
personnes est indépendante de la minute considérée. 
On a observé que la probabilité pour qu'une personne se présente entre la
minute $n$ et la minute $n+1$ est: $p = 0.1$. On veut calculer la probabilité 
pour que : 3,4,5,6,7,8... personnes se présentent au guichet entre 10h et 11h.
\begin{enumerate}
\item Définir une variable aléatoire adaptée, puis répondre au problème considéré.

\item Quelle est la probabilité pour que au moins 10 personnes
se présentent au guichet entre 10h et 11h?
\end{enumerate}
\finenonce{006010}


\finexercice


\exercice{6011, quinio, 2011/05/20}

\enonce{006011}{}
Si dans une population une personne sur cent est un centenaire,
quelle est la probabilité de trouver au moins un centenaire parmi $100$
personnes choisies au hasard ? Et parmi $200$ personnes ?
\finenonce{006011}


\finexercice


\exercice{6012, quinio, 2011/05/20}

\enonce{006012}{}
Un industriel doit vérifier l'état de marche de ses machines et en
remplacer certaines le cas échéant. D'après des statistiques précédentes, 
il évalue à 30\% la probabilité pour une
machine de tomber en panne en 5 ans; parmi ces dernières, 
la probabilité de devenir hors d'usage suite à une panne plus grave est 
évaluée à 75\%; cette probabilité est de 40\% pour une machine
n'ayant jamais eu de panne.
\begin{enumerate}
\item Quelle est la probabilité pour une machine donnée de plus de cinq
ans d'être hors d'usage ?

\item Quelle est la probabilité pour une machine hors d'usage de n'avoir
jamais eu de panne auparavant ?

\item Soit $X$ la variable aléatoire <<nombre de machines
qui tombent en panne au bout de 5 ans, parmi 10 machines choisies au
hasard>>. Quelle est la loi de probabilité de $X$, (on
donnera le type de loi et les formules de calcul), son espérance, sa
variance et son écart-type ?

\item Calculer $P[X=5]$.
\end{enumerate}
\finenonce{006012}


\finexercice


\exercice{6013, quinio, 2011/05/20}

\enonce{006013}{}
Une population comporte en moyenne une personne mesurant plus de 1m90 sur
80 personnes.
Sur 100 personnes, calculer la probabilité qu'il y ait au moins une
personne mesurant plus de 1.90m (utiliser une loi de Poisson).
Sur 300 personnes, calculer la probabilité qu'il y ait au moins une
personne mesurant plus de 1.90m.
\finenonce{006013}


\finexercice


\exercice{6902, ruette, 2013/01/24}
\enonce{006902}{}
Soit $n$ un entier strictement positif.
Quelle est la loi du nombre de garçons dans une famille de $n$
enfants ? (Préciser les hypothèses que vous faites)
\finenonce{006902}
\finexercice


\exercice{6903, ruette, 2013/01/24}
\enonce{006903}{}
Une usine fabrique des transistors. Chaque transistor a une probabilité de 3\%
d'être défectueux. Quelle est la loi du nombre de transistors défectueux dans 
un lot de 100 transistors ? Que vaut son espérance et que représente-t-elle ?
\finenonce{006903}
\finexercice


\exercice{6904, ruette, 2013/01/24}
\enonce{006904}{}
L'entreprise Luminex fabrique des lampes, dont 80\% durent plus de 3000 heures. 
Des tests sont effectués sur des échantillons de taille $n = 15$. 
\begin{enumerate}
\item
Quelle est le nombre moyen de lampes qui ont une durée de vie inférieure à  3000 heures dans un échantillon de taille 15 ?
\item
Quelle est la probabilité que toutes les lampes de l'échantillon durent plus de 3000 heures ?
\item
Quelle est la probabilité que 13 lampes ou plus, dans un échantillon de taille 15, durent plus de 3000 heures ?
\end{enumerate}
\finenonce{006904}

\finexercice


\exercice{6905, ruette, 2013/01/24}
\enonce{006905}{}
Un transporteur aérien a observé que  25\% en moyenne des personnes ayant réservé 
un siège pour un vol ne se présentent pas au départ. Il décide d'accepter jusqu'à  
23  réservations alors qu'il ne dispose que de  20  sièges pour ce vol.
\begin{enumerate}
\item
Soit $X$ la variable aléatoire ``nombre de clients qui viennent après
réservation quand 23 places ont été réservées''.
Quelle est la loi de $X$ (précisez les hypothèses que vous faites pour modéliser la situation) ?
Quelle est son espérance ?
\item
Si 23 personnes ont réservé, quelle est la probabilité que toutes les personnes qui se présentent au départ aient un siège ?
\end{enumerate}
\finenonce{006905}

\finexercice


\exercice{6906, ruette, 2013/01/24}
\enonce{006906}{}
 On lance $10$ fois une pièce supposée bien équilibrée. 
 On désigne par $X$ la fréquence du nombre de fois où pile a 
 été obtenu (c'est-à-dire le nombre de pile divisé
par 10).   
\begin{enumerate}
\item Quelle est  la loi de $X$ ?
\item Avec quelle probabilité $X$  est-elle strictement au dessus de 0,5 ?
\item Avec quelle probabilité $X$ est-elle comprise entre 0,4 et 0,6 (bornes
incluses) ?
\item Déterminer le plus petit entier $a>0$ telle que la probabilité que $X$ 
soit dans l'intervalle $[0,5-\frac{a}{10},0,5+\frac{a}{10}]$ soit supérieure à $95\%$. 
\item On lance la pièce $10$ fois. Elle tombe $3$ fois sur pile et
$7$ fois sur face. D'après vous la pièce est-elle bien
équilibrée (on justifiera sa réponse  en utilisant la question
3~? Même question si on obtient 1 fois pile et 9 fois face.
\end{enumerate}
\finenonce{006906}

\finexercice


\exercice{6907, ruette, 2013/01/24}
\enonce{006907}{}
Un standard téléphonique reçoit en moyenne 2 appels par minute. Les appels sont répartis au hasard dans le temps.
\begin{enumerate}
\item
Quelle est la loi de probabilité régissant le nombre d'appels reçus 
en 3 minutes ? Quelle est la probabilité qu'il n'y ait aucun appel en 3 minutes ?
\item
Quelle est la probabilité que le nombre d'appels en 2 minutes
soit supérieur ou égal à 5 ?
\end{enumerate}
\finenonce{006907}

\finexercice


\exercice{6908, ruette, 2013/01/24}
\enonce{006908}{}
Dans une dictature militaire, le dictateur veut augmenter le nombre de
naissances de garçons. Il impose la règle suivante : 
si une femme donne naissance à une fille, elle doit continuer à faire
des enfants ; si elle donne naissance à un garçon, elle doit arrêter
d'avoir des enfants. On suppose que chaque femme a au moins un enfant et 
pas plus de 5 enfants.
\begin{enumerate}
\item
Soit $X$ le nombre de filles par femme. Quelle est la loi de $X$ ?
\item
Quel est le nombre moyen de filles d'une femme ? Le nombre moyen 
de garçons ?
Cette règle est-elle efficace pour augmenter le nombre de garçons ?
\end{enumerate}
\finenonce{006908}

\finexercice


\exercice{6909, ruette, 2013/01/24}
\enonce{006909}{Loi multinomiale}
On lance quatre dés. Quelle est la probabilité d'obtenir deux 6 et deux 3 ?
\finenonce{006909}
\finexercice


\exercice{6910, ruette, 2013/01/24}
\enonce{006910}{} 
On lance deux dés à 6 faces.
Déterminer la loi de la variable aléatoire donnant le maximum des
deux chiffres obtenus.
\finenonce{006910}
\finexercice


\exercice{6911, ruette, 2013/01/24}
\enonce{006911}{}
On considère $\Omega=\{1,2,3,4\}$. Parmi les choix suivants, quels sont ceux qui donnent une probabilité $P$ sur $\Omega$ ?
\begin{eqnarray*}
a)&&\quad P(1) = 1/4,\quad P(2) = 3/8, \quad P(3) = 1/16,\quad P(4) = 3/16.\\
b)&&\quad P(1) = 0,\quad P(2) = 1/3,\quad P(3) = 1/6,\quad P(4) = 1/2.\\
c)&& \quad P(1) = 1/5,\quad P(2) = 1/4, \quad P(3) = 1/3, \quad P(4) = 1/2.\\
d)&& \quad P(1) =1/4, \quad P(2) = 1/2, \quad P(3) = -1/4, \quad P(4) = 1/2.
\end{eqnarray*}
\finenonce{006911}

\finexercice


\exercice{6912, ruette, 2013/01/24}
\enonce{006912}{} 
Le tableau ci-dessous donne la loi d'un couple de variables aléatoires
$Z=(X,Y)$, avec $X$ prenant ses valeurs dans $\{-1,1\}$
et $Y$ prenant ses valeurs dans $\{-1,0,1\}$.
$$
\begin{array}{|c|c|c|c|}
\hline X\ \backslash\ Y & -1 & 0 & 1 \\ \hline
-1&1/12 &1/3&1/12\\ \hline
1&1/6&1/6&1/6\\
\hline\end{array}
$$
\begin{enumerate}
\item
Déterminer la probabilité que $X$ et $Y$ soient égales.
\item
Déterminer les lois de $X$ et de $Y$.
\item
Déterminer les lois de $X+Y$ et de $XY$.
\end{enumerate}
\finenonce{006912}
\finexercice


\exercice{6913, ruette, 2013/01/24}
\enonce{006913}{}
Le tableau ci-dessous donne la loi d'un couple de variables aléatoires
$(X,Y)$, avec $X$ et $Y$ prenant chacune leurs valeurs dans $\{1,2,3\}$.
$$
\begin{array}{|c|c|c|c|}
\hline X\ \backslash\ Y & 1 & 2 & 3 \\ \hline
1&0&1/9&2/9\\ \hline
2&2/9&0&1/9\\ \hline
3&1/9&2/9&0\\
\hline\end{array}
$$
\begin{enumerate}
\item
Calculer la probabilité que $X$ et $Y$ soit égales.
\item
Déterminer les lois de $X$ et de $Y$.
\end{enumerate}
\finenonce{006913}
\finexercice


\exercice{6914, ruette, 2013/01/24}
\enonce{006914}{}
Charles ne supporte pas les chats et Sophie déteste les chiens. 
Charles n'élève pas plus d'un chien et Sophie pas plus d'un chat. 
La probabilité pour que Charles ait un chien est de $0,2$. Si 
Charles n'a pas de chien, la probabilité pour que Sophie ait un chat est de $0,1$.
On note $X$ le nombre de chiens de Charles, $Y$ le nombre de chats de Sophie
et $Z$ le nombre d'animaux du couple. 
\begin{enumerate}
\item 
Calculer la probabilité pour qu'ils n'aient pas d'animaux.
\item
On suppose de plus que la probabilité que $Z$ soit égal à 1 est de $0,1$.
\begin{enumerate}
\item
Calculer la probabilité pour que $Z$ soit égal à 2.
\item
Déterminer l'espérance et l'écart-type de $Z$.
\item
\'Etablir la loi de probabilité du couple $(X,Y)$.
Quelle est la loi de probabilité de $Y$ ?
\item
Les variables $X$ et $Y$  sont-elles indépendantes ?
\end{enumerate}\end{enumerate}
\finenonce{006914}

\finexercice


\exercice{6915, ruette, 2013/01/24}
\enonce{006915}{}
Dans une pile de $n$ ($n \geq 2$) feuilles dactylographiées, se trouvent les
deux lettres que l'on doit envoyer. On enlève une par une les
feuilles du paquet jusqu'à ce que l'une des lettres à envoyer se
trouve sur le dessus du paquet. On note $X_1$ la variable
aléatoire donnant le nombre de feuilles enlevées. On recommence
l'opération jusqu'à trouver la deuxième lettre et on note  $X_2$
la variable aléatoire donnant le nombre de feuilles
qu'il a fallu retirer du paquet après avoir trouvé la première lettre et
avant que la deuxième lettre soit
sur le dessus du paquet. Sans information supplémentaire, on peut
supposer que toutes les positions possibles pour les deux lettres sont
équiprobables. 
\begin{enumerate}
\item
Décrire l'ensemble $\Omega$ des résultats possibles pour cette
expérience aléatoire et la probabilité $P$ mise sur $\Omega$.
\item
Déterminer la loi du couple $(X_1,X_2)$ puis la loi de $X_1$ et de
$X_2$.
\item
Calculer la probabilité de l'événement ``$X_1=X_2$''.
\item
On note $Z=X_1+X_2+2$.  Que représente la variable aléatoire $Z$~?
Déterminer sa loi.
\end{enumerate} 
\finenonce{006915}
\finexercice

\section{260.04 Lois de distributions}


\exercice{6014, quinio, 2011/05/20}

\enonce{006014}{}
Une usine fabrique des billes de diamètre 8mm. Les erreurs d'usinage
provoquent des variations de diamètre.
On estime, sur les données antérieures, que l'erreur est une
variable aléatoire qui obeit à une loi normale les paramètres étant :
moyenne: $0$mm, écart-type: $0.02$mm. 
On rejette les pièces dont le diamètre n'est pas compris entre $7.97$mm et $8.03$mm.
Quelle est la proportion de billes rejetées?
\finenonce{006014}


\finexercice


\exercice{6015, quinio, 2011/05/20}

\enonce{006015}{}
Des machines fabriquent des plaques de tôle destinées à être
empilées. 
\begin{enumerate}
\item Soit $X$ la variable aléatoire <<épaisseur de la plaque en mm>> ; 
on suppose que $X$ suit une loi normale de paramètres $m=0.3$ et $\sigma =0.1$. 
Calculez la probabilité pour que $X$ soit inférieur à 0.36mm et la
probabilité pour que $X$ soit compris entre 0.25 et 0.35mm.

\item L'utilisation de ces plaques consiste à en empiler $n$, numérotées de $1$ à $n$ 
en les prenant au hasard : soit $X_{i}$ la variable aléatoire <<épaisseur de la plaque numéro $i$ en mm>> 
et $Z$ la variable aléatoire <<épaisseur des $n$ plaques en mm>>.
Pour $n = 20$, quelle est la loi de $Z$, son espérance et sa variance?
\end{enumerate}
\finenonce{006015}


\finexercice


\exercice{6016, quinio, 2011/05/20}

\enonce{006016}{}
Des machines fabriquent des plaques de tôle destinées à être
empilées; on estime à 0.1\% la proportion de plaques inutilisables.
L'utilisation de ces plaques consiste à en empiler $n$, numérotées
de $1$ à $n$ en les prenant au hasard.
Pour $n = 2000$, quelle est la loi suivie par la variable aléatoire $N$
<<nombre de plaques inutilisables parmi les 2000>> ? (on utilisera une loi de probabilité adaptée); 
quelle est la probabilité pour que $N$ soit inférieure ou égal à 3 ? 
Quelle est la probabilité pour que $N$ soit strictement inférieure à 3?
\finenonce{006016}


\finexercice


\exercice{6017, quinio, 2011/05/20}

\enonce{006017}{}
Des machines fabriquent des crêpes destinées à être empilées dans des paquets de 10.
Chaque crêpe a une épaisseur qui suit une loi normale de paramètres $m=0.6$mm et $\sigma =0.1$.
Soit $X$ la variable aléatoire <<épaisseur du paquet en mm>>.
Calculez la probabilité pour que $X$ soit compris entre 6.3mm et 6.6mm.
\finenonce{006017}


\finexercice


\exercice{6018, quinio, 2011/05/20}

\enonce{006018}{}
Sur un grand nombre de personnes on a constaté que la
répartition du taux de cholestérol suit une loi normale 
avec les résultats suivants:

- $56$\% ont un taux inférieur à $165$ cg;

- $34$\% ont un taux compris entre $165$ cg et $180$ cg;

- $10$\% ont un taux supérieur à $180$ cg.

Quelle est le nombre de personnes qu'il faut prévoir de soigner dans une
population de $10\,000$ personnes, si le taux maximum toléré sans
traitement est de $182$ cg?
\finenonce{006018}


\finexercice


\exercice{6019, quinio, 2011/05/20}

\enonce{006019}{}
Pour chacune des variables aléatoires qui sont décrites ci-dessous, indiquez quelle est la loi exacte 
avec les paramètres éventuels (espérance, variance) et indiquez éventuellement une loi approchée.
\begin{enumerate}
\item Nombre annuel d'accidents à un carrefour donné où la
probabilité d'accident par jour est estimée à $\frac{4}{365}$.

\item Nombre de garçons dans une famille de 6 enfants; nombre de filles par
jour dans une maternité où naissent en moyenne 30 enfants par jour.

\item Dans un groupe de 21 personnes dont 7 femmes, le nombre de femmes dans
une délégation de 6 personnes tirées au hasard.
\end{enumerate}
\finenonce{006019}


\finexercice


\exercice{6020, quinio, 2011/05/20}

\enonce{006020}{}
On effectue un contrôle de fabrication sur des pièces dont une
proportion $p=0.02$ est défectueuse.
\begin{enumerate}
\item On contrôle un lot de 1000 pièces :

Soit $X$ la variable aléatoire: <<nombre de pièces défectueuses
parmi 1000>>.
Quelle est la vraie loi de $X$ ? (on ne donnera que la forme générale);
quel est son espérance, son écart-type ?

\item En approchant cette loi par celle d'une loi normale adaptée, calculez
la probabilité pour que $X$ soit compris entre 18 et 22 ($P[18 \leq X\leq 22]$) ; 
on fera les calculs avec et sans correction de continuité.
On fera également les calculs avec la vraie loi pour comparer.
\end{enumerate}
\finenonce{006020}


\finexercice


\exercice{6021, quinio, 2011/05/20}

\enonce{006021}{}
 On effectue un contrôle sur des pièces de un euro
dont une proportion $p=0,05$ est fausse et sur des pièces 
de 2 euros dont une proportion $p'=0,02$ est fausse.
Il y a dans un lot $500$ pièces dont $150$ pièces de 
un euro et $350$ pièces de 2 euros.
\begin{enumerate}
\item On prend une pièce au hasard dans ce lot: quelle est la probabilité qu'elle soit fausse?

\item Sachant que cette pièce est fausse, quelle est la probabilité
qu'elle soit de un euro?

\item On contrôle à présent un lot de 1000 pièces de un euro.
Soit $X$ la variable aléatoire: <<nombre de pièces fausses parmi 1000>>.
Quelle est la vraie loi de $X$ ? (on ne donnera que la forme générale);
quelle est son espérance, son écart-type?

\item En approchant cette loi par celle d'une loi normale adaptée, calculez
la probabilité pour que $X$ soit compris entre 48 et 52.
\end{enumerate}
\finenonce{006021}


\finexercice


\exercice{6022, quinio, 2011/05/20}

\enonce{006022}{}
On jette un dé 180 fois.
On note $X$ la variable aléatoire : <<nombre de sorties du 4>>.
\begin{enumerate}
\item Quelle est la loi de $X$ ?

\item Calculez la probabilité pour que $X$ soit compris entre 29 et 32.
\end{enumerate}
\finenonce{006022}


\finexercice


\exercice{6023, quinio, 2011/05/20}

\enonce{006023}{}
Aux dernières élections présidentielles en France, le candidat A
a obtenu $20$\% des voix. On prend au hasard dans des bureaux de vote de
grandes villes des lots de $200$ bulletins: on note $X$ la variable aléatoire 
<<nombre de voix pour A dans les différents bureaux>>.
\begin{enumerate}
\item Quelle est la loi de probabilité de $X$?

\item Comment peut-on l'approcher?

\item Quelle est alors la probabilité pour que : $X$ soit supérieur 
à $45$? $X$ compris entre $30$ et $50$?

\item Pour un autre candidat B moins heureux le pourcentage des voix est de $2$\%.
En notant $Y$ le nombre de voix pour B dans les différents bureaux, sur $100$ bulletins, 
reprendre les questions 1 et 2. Quelle est alors la probabilité pour que : $Y$ soit supérieur à
$5$? $Y$ compris entre $1$ et $4$ ?
\end{enumerate}
\finenonce{006023}


\finexercice


\exercice{6024, quinio, 2011/05/20}

\enonce{006024}{}
On suppose qu'il y a une probabilité égale à $p$ d'être contrôlé lorsqu'on prend le tram. 
Monsieur A fait $n$ voyages par an sur cette ligne.
\begin{enumerate}
\item On suppose que $p=0.10$, $n=700$.
\begin{enumerate}
\item Quelle est la probabilité que Monsieur A soit contrôlé entre
60 et 80 fois dans l'année ?

\item Monsieur A voyage en fait toujours sans ticket. Afin de
prendre en compte la possibilité de faire plusieurs passages avec le même ticket, on suppose que le prix d'un ticket est de 1,12 euros.
Quelle amende minimale la compagnie doit-elle fixer pour que le fraudeur
ait, sur une période d'une année, une probabilité supérieure 
à $0.75$ d'être perdant ?
\end{enumerate}

\item On suppose que $p=0.50$, $n=300$.
Monsieur A voyage toujours sans ticket. Sachant que le prix d'un ticket est
de 1,12 euros, quelle amende minimale la compagnie doit-elle fixer pour que
le fraudeur ait, sur une période d'une année, une probabilité supérieure à $0.75$ d'être perdant ?
\end{enumerate}
\finenonce{006024}


\finexercice

\section{260.05 Espérance, variance}


\exercice{6916, ruette, 2013/01/24}
\enonce{006916}{} 
On lance deux dés à 6 faces. Soit $X_1$ le résultat du premier dé,
$X_2$ le résultat du deuxième dé, et $S=X_1+X_2$.
\begin{enumerate}
\item
Calculer $E(X_1)$ et $Var(X_1)$.
\item
En déduire $E(S)$ et $\text{Var}(S)$.
\end{enumerate} 
\finenonce{006916}

\finexercice


\exercice{6917, ruette, 2013/01/24}
\enonce{006917}{}
Une urne contient 10 boules rouges et 5 boules blanches.
\begin{enumerate}
\item
On effectue des tirages successifs en remettant après chaque tirage 
la boule sortie dans l'urne. Sur trois tirages, combien de boules
rouges va-t-on tirer en moyenne ? En moyenne, combien faudra-t-il effectuer 
de tirages avec remise avant de tirer une boule rouge ?
\item
On effectue maintenant 3 tirages sans remise. Combien de boules
rouges va-t-on tirer en moyenne ?
\end{enumerate} 
\finenonce{006917}
\finexercice


\exercice{6918, ruette, 2013/01/24}
\enonce{006918}{}
Pour une élection, une population de $N$ individus a eu à choisir
entre voter pour le candidat $A$ ou le candidat $B$. On note $m$ le
nombre de personnes ayant voté pour $A$.  On interroge au
hasard $k$ individus différents dans cette population ($1\leq k\leq N$). 
\begin{enumerate}
\item
On désigne par $a_1,\ldots, a_m$ les $m$ personnes qui ont
voté pour $A$. Pour $i\in\{1,\ldots,m\}$, on note $X_i$ l'indicatrice de l'événement ``la
personne $a_i$ est interrogée''. Quelle est la loi de $X_i$~?
\item
Déterminer l'espérance ainsi que la matrice de covariance 
de $(X_{1},\ldots,X_{m})$. 
\item
On pose $S=\sum_{i=1}^{m}X_i$. Quelle est la loi de $S$~?
\item
Déterminer le nombre moyen de personnes votant pour $A$ sur 
$k$ personnes tirées au hasard.
\item
Déterminer la variance de $S$.
\end{enumerate} 
\finenonce{006918}
\finexercice


\exercice{6919, ruette, 2013/01/24}
\enonce{006919}{} 
Le tableau ci-dessous donne la loi d'un couple de variables aléatoires
$(X,Y)$, avec $X$ prenant ses valeurs dans $\{0,3\}$ et
$Y$ prenant ses valeurs dans $\{0,1\}$.
Calculer la covariance entre $X$ et $Y$, puis la corrélation entre $X$ et $Y$.
$$
\begin{array}{|c|c|c|}
\hline X\ \backslash\ Y & 0 & 1 \\ \hline
0&1/3&0\\ \hline
3&1/6&1/2\\
\hline\end{array}
$$
\finenonce{006919}

\finexercice


\exercice{6920, ruette, 2013/01/24}
\enonce{006920}{}
Soit $U$ et $V$ deux variables aléatoires de même loi, à
valeurs dans $\{1,\ldots, N\}$. On pose $X=U-V$ et $Y=U+V$. Déterminer la
covariance entre $X$ et $Y$.
\finenonce{006920}

\finexercice

\section{260.06 Droite de régression}

\section{260.07 Fonctions génératrices}

\section{260.99 Autre}

\section{261.01 Densité de probabilité}


\exercice{6924, ruette, 2013/01/24}
\enonce{006924}{}
Madame Michel et Monsieur Lustucru vont chaque semaine au marché hebdomadaire 
de Kerplou. Madame Michel arrive à une heure aléatoire 
entre 8h et 12h et elle reste 30 minutes ; on suppose que son heure d'arrivée 
suit une loi uniforme. Monsieur Lustucru arrive à 10h pile et reste également
30 minutes. Quelle est la probabilité pour qu'ils se trouvent ensemble 
au marché à un certain moment ?
\finenonce{006924}

\finexercice


\exercice{6925, ruette, 2013/01/24}
\enonce{006925}{}
\def\I1{{ \rm 1\:\!\!\! l}}
Soit $X$ et $Y$ des variables aléatoires indépendantes de loi uniforme 
de densité $\I1_{[0,1]}$
et $Z=X+Y$. Calculer la loi de $Z$.
\finenonce{006925}

\finexercice


\exercice{6926, ruette, 2013/01/24}
\enonce{006926}{Fonction de répartition, indépendance}
Soit $X_1,\ldots, X_n$ des variables aléatoires indépendantes de même loi 
exponentielle
$\mathcal{E}(1)$ et $Z=\min(X_1,\ldots,X_n)$. Déterminer la loi de $Z$.
\finenonce{006926}

\finexercice

\section{261.02 Loi faible des grands nombres}


\exercice{6927, ruette, 2013/01/24}
\enonce{006927}{}
On suppose que le nombre de pièces sortant d'une
usine donnée en une journée est une variable aléatoire
d'espérance $50$.
\begin{enumerate}
\item
Peut-on estimer la probabilité que la production de demain
dépasse 75 pièces~?
\item
Que peut-on dire de plus sur cette probabilité si on sait que l'écart-type de la production quotidienne est
de $5$ pièces~?
\end{enumerate} 
\finenonce{006927}

\finexercice


\exercice{6928, ruette, 2013/01/24}
\enonce{006928}{}
Pour étudier les particules émises par une substance radioactive, 
on dispose d'un détecteur. On note $X$ la variable aléatoire 
représentant le nombre de particules qui atteignent le détecteur pendant 
un intervalle de temps $\Delta t$. Le nombre maximal de particules que le 
détecteur peut compter pendant un intervalle de temps $\Delta t$ est de 
$10^3$. 
On suppose que  $X$ suit une loi de Poisson de paramètre 
$\lambda=10^2$. Donner  une majoration de la probabilité que $X$ 
dépasse~$10^3$.\\ \textit{(On rappelle que l'espérance et la variance d'une loi
de Poisson $\mathcal{P}(\lambda)$ sont égales à $\lambda$.)}
\finenonce{006928}

\finexercice

\section{261.03 Convergence en loi}

\section{261.04 Loi normale}


\exercice{6921, ruette, 2013/01/24}
\enonce{006921}{}
\begin{enumerate}
\item
Soit $X$ une variable aléatoire suivant une loi normale $N(m,\sigma^2)$.
Quelle est la probabilité que $X$ soit supérieure à $m$ ? 

\item
Soit  $Y$ une variable aléatoire suivant une loi normale $N(0,1)$.
Quelle est la loi de $-Y$ ? la loi de $\sigma Y+m$ ?
\item
En utilisant la table de $\mathcal{N}(0,1)$, déterminer 
$P(0\le Y\le 0,8)$,
$P(-0,6\le Y\le 0)$ et $P(Y\le 0,8)$.
\end{enumerate}
\finenonce{006921}
\finexercice


\exercice{6922, ruette, 2013/01/24}
\enonce{006922}{}
Soient $X_1$, $X_2$, $X_3$, trois variables aléatoires de loi normale 
indépendantes telles que $E(X_1 ) = 100$, $\text{Var}(X_1 ) = 100$, $E(X_2 ) = 20$, 
$\text{Var}(X_2 ) = 4$, $E(X_3 ) = 50$, $\text{Var}(X_3 ) = 25$. On forme la combinaison linéaire $Y=X_1 +2X_2 -X_3$.

Déterminer $E(Y)$ et $\text{Var}(Y)$. Quelle est la loi de $Y$ ? 
\finenonce{006922}
\finexercice


\exercice{6923, ruette, 2013/01/24}
\enonce{006923}{}  
Une machine est conçue pour confectionner 
des paquets d'un poids de 500g, mais ils n'ont pas exactement tous le même poids. 
On a constaté que la distribution des poids autour de la valeur moyenne de 500g avait un écart-type de 25g.
\begin{enumerate}
\item
Par quelle loi est-il raisonnable de modéliser le poids des paquets ?
\item
Sur 1000 paquets, quel est le nombre moyen de paquets pesant entre 480g et 520g ? 
(\textit{utiliser la table  de $\mathcal{N}(0,1)$})
\item
Combien de paquets pèsent entre 480g et 490g ?
\item
Sur 1000 paquets, quel est le nombre moyen de paquets pesant plus de 450g ?
\item
Trouver $a$ tel que les 9/10 de cette production aient un poids compris
entre $500-a$ et $500+a$.
\end{enumerate}
\finenonce{006923}

\finexercice
\section{261.99 Autre}


\exercice{6954, ruette, 2013/01/24}
\enonce{006954}{Le paradoxe des cordes de Bertrand}

On considère un cercle. On choisit une corde au hasard et on cherche la
probabilité pour que cette corde passe à l'intérieur d'un triangle 
équilatéral inscrit dans le cercle dont un des sommets
est une des extrémités de la corde (voir la figure) ou, de façon équivalente,
pour que cette corde soit plus longue que le côté d'un triangle
équilatéral inscrit dans le cercle.

\bigskip
%\bigskip
\centerline{\includegraphics{../images/img006954-1}}


\bigskip
{\it Cet exercice est volontairement formulé de façon imprécise. Essayez
de trouver la solution de façon intuitive (pas de théorème à utiliser ici).
Il est fort probable que vous ne trouviez pas tous la même probabilité.
Historiquement, ce genre de paradoxe a motivé l'introduction d'une théorie
rigoureuse du hasard.
}
\finenonce{006954}
\finexercice

\section{262.01 Estimation}


\exercice{6028, quinio, 2011/05/20}

\enonce{006028}{}
Le staff médical d'une grande entreprise fait ses
petites statistiques sur le taux de cholestérol de ses employés; les
observations sur 100 employés tirés au sort sont les suivantes.

\begin{tabular}{cc}
taux de cholestérol en cg:(centre classe) & effectif d'employés: \\
$120$ & $9$ \\ 
$160$ & $22$ \\ 
$200$ & $25$ \\ 
$240$ & $21$ \\ 
$280$ & $16$ \\ 
$320$ & $7$
\end{tabular}

\begin{enumerate}
\item Calculer la moyenne $m_{e}$ et l'écart-type $\sigma_{e}$ sur l'échantillon.
\item Estimer la moyenne et l'écart-type pour le taux de cholestérol 
dans toute l'entreprise.
\item Déterminer un intervalle de confiance pour la moyenne.
\item Déterminer la taille minimum d'échantillon pour que
l'amplitude de l'intervalle de confiance soit inférieure à 10.
\end{enumerate}
\finenonce{006028}


\finexercice


\exercice{6029, quinio, 2011/05/20}

\enonce{006029}{}
Sur $12\,000$ individus d'une espèce, on a dénombré $13$
albinos. Estimer la proportion d'albinos dans l'espèce. On comparera les
méthodes d'approximation des lois réelles par d'autres lois
classiques.
\finenonce{006029}


\finexercice

\section{262.02 Tests d'hypothèses, intervalle de confiance}


\exercice{6025, quinio, 2011/05/20}

\enonce{006025}{}
Un échantillon de $10\,000$ personnes sur une population étant donné, 
on sait que le taux moyen de personnes à soigner pour un problème de 
cholestérol élevé est de $7,5$\%. Donner un
intervalle dans lequel on soit <<sûr>>  à $95$\%, 
de trouver le nombre exact de personnes à soigner sur les $10\,000$.
\finenonce{006025}


\finexercice


\exercice{6026, quinio, 2011/05/20}

\enonce{006026}{}
Un vol Marseille - Paris est assuré par un Airbus de $150$
places ; pour ce vol des estimations ont montré que la probabilité
pour qu'une personne confirme son billet est $p=0.75$. La
compagnie vend $n$ billets, $n>150$. Soit $X$ la variable aléatoire <<nombre de personnes
parmi les $n$ possibles, ayant confirmé leur réservation pour ce vol>>.
\begin{enumerate}
\item Quelle est la loi exacte suivie par $X$ ?
\item Quel est le nombre maximum de places que la compagnie peut vendre pour
que, à au moins $95$\%, elle soit sûre que tout le monde puisse
monter dans l'avion, c'est-à-dire $n$ tel que : $P[X>150] \leq 0.05$ ?
\item Reprendre le même exercice avec un avion de capacité de $300$
places; faites varier le paramètre $p = 0.5$ ; $p=0.8$.
\end{enumerate}
\finenonce{006026}


\finexercice


\exercice{6027, quinio, 2011/05/20}

\enonce{006027}{}
Un petit avion (liaison Saint Brieuc-Jersey) peut accueillir chaque jour 30
personnes; des statistiques montrent que 20\% des clients ayant réservé ne viennent pas.
Soit $X$ la variable aléatoire: <<nombre de clients qui
se présentent au comptoir parmi 30 personnes qui ont réservé>>.
\begin{enumerate}
\item Quelle est la loi de $X$ ? (on ne donnera que la forme générale);
quelle est son espérance, son écart-type ?

\item Donner un intervalle de confiance au seuil 95\%, permettant d'estimer le
nombre de clients à prévoir.
\end{enumerate}
\finenonce{006027}


\finexercice


\exercice{6030, quinio, 2011/05/20}

\enonce{006030}{}
Une compagnie aérienne a demandé des statistiques afin d'améliorer la sûreté au décollage 
et définir un poids
limite de bagages. Pour l'estimation du poids des voyageurs et du poids des
bagages, un échantillon est constitué de 300 passagers qui ont accepté d'être pesés : on a obtenu une moyenne $m_{e}$ de 68kg, 
avec un écart-type $\sigma_{e}$ de 7 kg.
\begin{enumerate}
\item Définir un intervalle de confiance pour la moyenne des
passagers. (On admet que le poids des passagers suit une loi normale de moyenne
$m$, d'écart-type $\sigma $.)
\item Montrer que l'on peut considérer que le poids des passagers est
une variable aléatoire $X$ de moyenne 70 kg, d'écart-type 8 kg.
\item En procédant de même pour le poids des bagages, on admet les résultats :
\begin{itemize}
\item Si le poids maximum autorisé est de 20 kg, le poids des bagages
peut être considéré comme une variable aléatoire $Y$ de
moyenne 15 kg, d'écart-type 5 kg.
\item La capacité de l'avion est de 300 passagers; l'avion pèse, 
à vide, 250 tonnes. Le décollage est interdit si le poids total dépasse 276.2 tonnes.
Quelle est la probabilité pour que le décollage soit interdit ? 
\end{itemize}
\end{enumerate}
\finenonce{006030}


\finexercice


\exercice{6031, quinio, 2011/05/20}

\enonce{006031}{}
Afin de mieux satisfaire leurs clients, une grande société
fournisseur d'accès internet fait ses statistiques sur le nombre
d'appels reçus en \emph{hotline}, elle pourra ainsi évaluer le temps
d'attente pour le client et le nombre d'employés à mettre au
standard; les résultats de l'enquête portent sur 200 séquences \
consécutives de une minute, durant lesquelles le nombre d'appels moyen a 
été de 3 appels par minute. On suppose que les appels sont répartis également dans le
temps: on partage un intervalle de temps en unités de une seconde;
alors dans chaque unité de temps, il y a au plus un appel.
\begin{enumerate}
\item Quelle est la loi de probabilité du nombre d'appels reçus
en 4 minutes?
\item Montrer que l'on peut approcher cette loi par une loi de Poisson.
\item Donner un intervalle de confiance pour le nombre moyen d'appels en 4 minutes.
\end{enumerate}
\finenonce{006031}


\finexercice


\exercice{6032, quinio, 2011/05/20}

\enonce{006032}{}
On s'intéresse au problème des algues toxiques qui
atteignent certaines plages de France; après étude on constate que
10\% des plages sont atteintes par ce type d'algues et on veut tester
l'influence de rejets chimiques nouveaux sur l'apparition de ces algues.
Pour cela 50 plages proches de zones de rejet chimiques, sont observées; 
on compte alors le nombre de plages atteintes par l'algue nocive :
on constate que 10 plages sont atteintes par l'algue.
Pouvez-vous répondre à la question 
<<Les rejets chimiques ont-t-il modifié, de façon significative, 
avec le risque $\alpha =0.05$, le nombre de plages atteintes ?>> 
\finenonce{006032}


\finexercice


\exercice{6033, quinio, 2011/05/20}

\enonce{006033}{}
On veut étudier la liaison entre les caractères : 
<<être fumeur>> (plus de 20 cigarettes par
jour, pendant 10 ans) et <<avoir un cancer de la gorge>>, 
sur une population de 1000 personnes, dont 500 sont
atteintes d'un cancer de la gorge. Voici les résultats observés:

Tableau observé

\begin{center}
\begin{tabular}{cccc}
\emph{Observé} & cancer & non cancer & marge \\ 
\text{fumeur} & 342 & 258 & 600 \\ 
\text{non fumeur} & 158 & 242 & 400 \\ 
\text{marge} & 500 & 500 & 1000
\end{tabular}
\end{center}
Faire un test d'indépendance pour établir la liaison entre ces caractères.
\finenonce{006033}


\finexercice

\exercice{6929, ruette, 2013/01/24}
\enonce{006929}{}  
On mesure une certaine grandeur physique $G$ avec un appareil dont la
précision est caractérisée par l'écart-type $\sigma$. On fait l'hypothèse que 
les mesures suivent une loi normale.
\begin{enumerate}
\item
On effectue une seule mesure, on trouve $g_{1} = 1,364$.
On suppose connue la précision de l'appareil de mesure : 
$\sigma=4,3.10^{-3}$ (unité arbitraire).
%Donner la meilleure estimation possible de $G$.
Déterminer un intervalle de confiance contenant, avec une probabilité de 
90\%, la valeur $G$.
\item
On ignore la précision de l'appareil de mesure.
On effectue 5 mesures. On trouve : 

\begin{center}
\begin{tabular}{|c|c|c|c|c|}
\hline
1,365&1,371&1,368&1,359&1,362\\\hline
%1.366	&1.365&1.367&1.363&1.364\\\hline
\end{tabular}
\end{center}
 
Donner des estimations  de $G$ et de $\sigma$.
Déterminer un intervalle de confiance contenant, avec une probabilité de 
90\%, la valeur $G$.
\end{enumerate}
\finenonce{006929}

\finexercice


\exercice{6930, ruette, 2013/01/24}
\enonce{006930}{}
On interroge 1000 électeurs, 521 d'entre eux ont déclaré avoir l'intention 
de voter pour le candidat A.
Indiquer avec une probabilité de  0,95  entre quelles limites se situe la 
proportion du corps électoral favorable à A au moment du sondage.
\finenonce{006930}

\finexercice


\exercice{6931, ruette, 2013/01/24}
\enonce{006931}{}  
La firme Comtec vient de développer un 
nouveau dispositif électronique. Avant de le mettre en production, on 
veut en estimer la fiabilité en termes de durée de vie. D'après le 
bureau de Recherche et Développement de l'entreprise, l'écart-type 
de la durée de vie de ce dispositif serait de l'ordre de  100 heures.

Déterminer le nombre d'essais requis pour estimer, avec un niveau de 
confiance de 95\%, la durée de vie moyenne d'une grande production 
de sorte que la marge d'erreur dans l'estimation n'excède pas  $\pm$ 50 heures.
 Même question pour une marge d'erreur n'excédant pas $\pm$ 20 heures.
\finenonce{006931}

\finexercice


\exercice{6932, ruette, 2013/01/24}
\enonce{006932}{} 
On effectue un sondage sur un échantillon de 10\,000 personnes à la veille d'un
référendum : 4903 d'entre elles s'apprêtent à voter oui, et 5097 à voter non. 
Quel risque d'erreur court-on en prédisant la victoire du non ?
\finenonce{006932}
\finexercice


\exercice{6933, ruette, 2013/01/24}
\enonce{006933}{} 
Chaque jour, un train subit un retard aléatoire au départ, évalué en minutes,
indépendant des retards des autres jours et dont la loi est approximativement 
$\mathcal{E}(\lambda)$ (\textit{on rappelle qu'une loi exponentielle $\mathcal{E}(\lambda)$
a pour espérance $1/\lambda$ et pour écart-type $1/\lambda$}).
Sur 400 jours, le retard moyen est de 10 minutes. 
Donner un intervalle de confiance de niveau approximativement 95\% pour 
$\lambda$. 
\finenonce{006933}

\finexercice


\exercice{6934, ruette, 2013/01/24}
\enonce{006934}{} Sur un lot de  700 boulons soumis à des 
essais de rupture, 300  ont résisté. Sur un second lot de  225, 125  
ont résisté. Peut-on admettre, au seuil de  5\%, que ces deux lots 
appartiennent à la même population ?
\finenonce{006934}

\finexercice


\exercice{6935, ruette, 2013/01/24}
\enonce{006935}{}
A l'issue d'une expérience de 1000 tirages, un générateur
de chiffres aléatoires a donné les résultats suivants :
\begin{center}
\begin{tabular}
[c]{|c|llllllllll|}%
\hline
chiffre & 1 & 2 & 3 & 4 & 5 & 6 & 7 & 8 & 9 & 0\\
nombre d'apparitions & 87 & 103 & 90 & 110 & 81 & 108 & 85 & 123 & 90
& 123\\
\hline
\end{tabular}
\end{center}
 Tester au niveau 5\% l'hypothèse selon laquelle le générateur 
simule de façon satisfaisante un tirage uniforme sur 
les entiers $\{0,\ldots, 9\}$.
\finenonce{006935}

\finexercice


\exercice{6936, ruette, 2013/01/24}
\enonce{006936}{}
Le couvert végétal du domaine vital d'un élan d'Amérique se
composait de  feuillus ($25,8\%$ de la superficie), de forêts mixtes 
($38\%$  de la superficie), de résineux ($25,8\%$ de la superficie) et d'un marécage ($10,4\%$ de la
superficie). Dans ce domaine, l'élan fut localisé à $511$ reprises
au cours de l'année : $118$ fois dans les feuillus, $201$ fois dans les forêts mixtes, $110$ fois dans les
résineux et $82$ fois dans le marécage. [Source~: B. Scherrer
``Biostatistique'', éditeur Gaetan Morin, 1984, page 556]

L'élan fréquente-t-il indifféremment les quatre types de
végétation de son domaine vital~?
\finenonce{006936}
\finexercice


\exercice{6937, ruette, 2013/01/24}
\enonce{006937}{} 
\label{exBateson} A la suite de la formulation des lois de Mendel, Bateson a
effectué des croisements avec des pois de senteur afin d'étudier
la couleur (pourpre ou rouge) et la forme du pollen (allongée ou
ronde). Ces croisements ont été fait sur des hybrides qui sont de
couleur pourpre et ont un pollen de forme allongé ; en
supposant que la couleur et la forme du pollen sont chacun controlés par un gène qui a deux
formes différentes (appelées allèles) notées $S$ et $s$ pour
la couleur et $T$, $t$  pour  la forme du pollen, le génotype d'un
pois hybride pour ces deux caractères est $(Ss,Tt)$. Les majuscules
désignent les allèles dominants (ici la couleur pourpre et la
forme allongée du pollen).  Les résultats
des croisements sont donnés dans le tableau ci-dessous~:
\begin{center}
\begin{tabular}{|c|cc|}
\hline
& couleur pourpre & couleur rouge \\
\hline
pollen allongé & 1528 & 117 \\
pollen rond & 106 &381\\
\hline
\end{tabular}
\end{center}
En faisant l'hypothèse que les quatre possibilités de
transmission des allèles pour chacun de ces deux gènes sont
équiprobables, peut-on affirmer que les  gènes contrôlant  ces deux caractères
sont transmis indépendamment l'un de l'autre~?\\
Comment s'assurer que l'hypothèse sur la transmission
équiprobable des quatre possibilités de
transmission des allèles est valide pour chacun de ces deux
gènes~?
\finenonce{006937}
\finexercice

\section{262.99 Autre}

\section{300.00 Groupe quotient, théorème de Lagrange}


\exercice{1434, ortiz, 1999/04/01}

\enonce{001434}{}
Soit $G$ un groupe non r\'eduit \`a un
\'el\'ement. Un sous-groupe $M$ de $G$ est dit
\textit{maximal} si le seul sous-groupe de $G$,
distinct de $G$ et contenant $M,$ est $M$
lui-m\^eme. Les questions sont ind\'ependantes.
\begin{enumerate}
    \item
    \begin{enumerate}
        \item Montrer que $6\Zz$ n'est pas un sous-groupe maximal de $\Zz$.
        \item Montrer que $5\Zz$ est un sous-groupe maximal de $\Zz$.
    \end{enumerate}
    \item  On pose $G:=\Zz/8\Zz.$ Soit $H_1$ le sous-groupe de $G$
engendr\'e par $\overline{4}$ et $H_2$ le sous-groupe de $G$ engendr\'e par $%
\overline{2}$.
    \begin{enumerate}
         \item Expliciter les \'el\'ements de $H_1$ et $H_2$.
         \item Montrer que $H_1$ n'est pas un sous-groupe maximal de $G$ et que $%
H_2$ est un sous- groupe maximal de $G$.
    \end{enumerate}
\end{enumerate}
\finenonce{001434}



\finexercice



\exercice{1435, ortiz, 1999/04/01}

\enonce{001435}{}
 D\'eterminer tous les sous-groupes de $\Zz/8\Zz.$

\finenonce{001435}



\finexercice



\exercice{1436, ortiz, 1999/04/01}

\enonce{001436}{}
Montrer que le groupe-quotient $\Cc/\Rr$ est
isomorphe \`a $\Rr.$
\finenonce{001436}



\finexercice



\exercice{1437, ortiz, 1999/04/01}

\enonce{001437}{}
Soit $G$ le groupe $\Qq/\Zz$. Si $q\in\Q$, on note
$\text{cl}(q)$ la classe de $q$ modulo $~\Zz.$
\begin{enumerate}
    \item  Montrer que $\text{cl}(\frac {35}{6})=\text{cl}(\frac 56)$ et
d\'eterminer l'ordre de $\text{cl}(\frac
{35}{6})$.
    \item  Montrer que si $x\in G$ il existe un unique $\alpha \in \Qq\cap
\left[ 0,1\right[ $ tel que $x=\text{cl}(\alpha ).$
    \item   Montrer que tout \'el\'ement de $G$ est d'ordre fini et qu'il existe
des \'el\'ements d'ordre arbitraire.
\end{enumerate}
\finenonce{001437}



\finexercice



\exercice{1438, ortiz, 1999/04/01}

\enonce{001438}{}
 D\'ecrire le groupe-quotient $\Rr^{*}/\Rr_{+}^{*}$
et montrer qu'il est isomorphe \`a $\Zz/2\Zz.$

\finenonce{001438}



\finexercice



\exercice{1439, ortiz, 1999/04/01}

\enonce{001439}{}
Montrer que tout quotient d'un groupe monog\`ene
est monog\`ene.
\finenonce{001439}



\finexercice



\exercice{1440, ortiz, 1999/04/01}

\enonce{001440}{}
Soient $G$ le groupe-produit $\left(
\Zz/4\Zz\right) \times
\left( \Zz/4\Zz\right) $ et $H$ le sous-groupe de $G$
engendr\'e \nolinebreak par
$(\overline{3},\overline{2}).$
\'Ecrire la d\'ecomposition de $G$ suivant les classes \`a
gauche modulo $H.$ D\'ecrire le groupe-quotient
\mbox{$G/H.$}
\finenonce{001440}



\finexercice



\exercice{1441, legall, 1998/09/01}

\enonce{001441}{}
Soit $  G  $ un groupe $  Z(G)=\{ h\in G   ;   \forall g \in g   ,
  gh=hg\}  .$
\begin{enumerate}
    \item Montrer que $  Z(G)  $ est un sous-groupe distingu\'e de $  G  .$
    \item Montrer que si $  G/Z(G)  $ est monog\`ene $  G  $ est cyclique.
\end{enumerate}
\finenonce{001441}



\finexercice



\exercice{1442, legall, 1998/09/01}

\enonce{001442}{}
 Soit $  G  $ un groupe ; on note $  D(G)  $ le
groupe engendr\'e par les \'el\'ements de la forme
$ghg^{-1}h^{-1}$  ; $g,h\in G.$
\begin{enumerate}
    \item Montrer que $  D(G)  $ est distingu\'e dans $  G  .$
    \item Montrer que $  G/D(G)  $ est commutatif~; plus g\'en\'eralement montrer qu'un sous-groupe distingu\'e
$  H  $ de $  G  $ contient $  D(G)  $ si et seulement si $  G/H  $ est commutatif.
\end{enumerate}

\finenonce{001442}



\finexercice



\exercice{1443, legall, 1998/09/01}

\enonce{001443}{}
Soit $  G  $ un groupe~; on note, pour tout $  g \in G  $
$  \varphi _g   $ l'application $  x\mapsto gxg^{-1}  $ de $  G  $ dans lui-m\^eme et
$  \hbox{Int}(G) =\{ \varphi _g   ;   g \in G \}   .$
\begin{enumerate}
    \item Montrer que $  \hbox{Int}(G)  $ est un sous-groupe distingu\'e de $  \hbox{Aut}(G)  .$
    \item Soit $  f : G\rightarrow \hbox{Int}(G)  $ l'application $  g \mapsto \varphi _g  .$
Montrer que $  f   $ est un homomorphisme de groupe. Calculer $  \hbox{Ker}(f)  .$
    \item En d\'eduire que $  G/Z(G)  $ est isomorphe \`a $  \hbox{Int}(G)  .$
\end{enumerate}
\finenonce{001443}



\finexercice



\exercice{1444, legall, 1998/09/01}

\enonce{001444}{}
Soit $  G  $ un groupe, $  H  $ et $  K  $ deux sous-groupes de $
G  .$ On note $  HK=\{ hk  ;   h\in H   ,   k\in K\}  .$ On suppose
que $  K  $ est distingu\'e dans $ G  .$
\begin{enumerate}
    \item Montrer que  $  HK=KH  $ et que $  HK  $ est un sous-groupe de $  G  .$
    \item Montrer que $  H  $ et $  K  $ sont des sous-groupes de $  KH  $ et que
$  K\cap H  $ est un sous-groupe distingu\'e de $  H  $ et que $  K  $ est distingu\'e
dans $  KH  .$
    \item Soit $  \varphi : H\rightarrow (HK)/K  $ la restriction \`a $  H  $ de l'application quotient.
Calculer le noyau et l'image de $  \varphi  .$ En d\'eduire que les groupes $  H/(K\cap H)  $ et $  (HK)/K  $ sont isomorphes.
\end{enumerate}
\finenonce{001444}



\finexercice



\exercice{1445, legall, 1998/09/01}

\enonce{001445}{}
Soit $  G  $ un groupe, $  H  $ et $  K  $ deux sous-groupes
distingu\'es de $  G  $ avec $  H\subset K  .$
\begin{enumerate}
    \item Montrer que $  K/H  $ est un sous-groupe distingu\'e de $  G/H  .$
    \item Montrer que le quotient $  (G/H)/(K/H)  $ est isomorphe \`a $  G/K  .$
\end{enumerate}
\finenonce{001445}



\finexercice



\exercice{1446, legall, 1998/09/01}

\enonce{001446}{}
 Soit $  G  $ le sous-groupe de $  Gl(2, \R )  $
engendr\'e par les matrices $  \displaystyle{  A= \frac{1 }{\sqrt
2}
\begin{pmatrix}-1 & 1\cr 1 & 1 \cr \end{pmatrix} }  $ et $  B= \begin{pmatrix}-1 & 0\cr 0 & 1 \cr \end{pmatrix}  .$
\begin{enumerate}
    \item Soit $  H  $ le sous-groupe de $  G  $ engendr\'e par $  AB  .$ Calculer $  \vert H\vert  $
    \item Montrer que $  H  $ est distingu\'e dans $  G  .$ Calculer le quotient $  G/H  ;$ en d\'eduire $  \vert G\vert   .$
\end{enumerate}

\finenonce{001446}



\finexercice



\exercice{1447, ortiz, 1999/04/01}

\enonce{001447}{}
 Les questions sont ind\'ependantes.
\begin{enumerate}
    \item
    \begin{enumerate}
         \item Montrer que l'application $f:\Zz^2\to \Zz$, $(x,y)\mapsto3x+6y$
est un morphisme de groupes.
         \item D\'eterminer le noyau $\ker f$ de $f$ et montrer qu'il n'existe pas
$(p,q)\in\Zz^2$ tel que $\ker f=p\Zz\times q\Zz.$
         \item Montrer que le groupe-quotient $\Zz^2/\Zz(-2,1)$ est isomorphe au groupe $3\Zz.$
    \end{enumerate}
    \item  Soit $G$ le sous-groupe de $\Zz^2$ engendr\'e par $(2,0)$ et $(0,2).$
Montrer que le groupe-quotient $\Zz^2/G$ est
isomorphe \`a $\Zz/2\Zz\times\Zz/2\Zz.$
\end{enumerate}

\finenonce{001447}



\finexercice



\exercice{1448, hilion, 2003/10/01}

\enonce{001448}{}
\begin{enumerate}
\item Montrer que les sous-groupes de $\Zz$ sont de la forme $n\Zz$ o\`u $n\in\Nn$.
(indication: utiliser la division euclidienne).
\item Rappeler pourquoi ces sous-groupes sont distingués. On peut donc considérer les groupes quotients $\Zz/n\Zz$.
\item Montrer que $\Zz/n\Zz$ est isomorphe au groupe des racines $n^{\text{i\`eme}}$ de l'unité.
\item Montrer que $\Zz/n\Zz$ est isomorphe au groupe engendré par un cycle de longueur $n$ dans $S_N$ ($N\geq n)$.
\item Plus généralement, montrer qu'il existe, à isomorphisme près, un seul groupe monogène (ie engendré par un seul élément) d'ordre $n$, appelé groupe cyclique d'ordre $n$.
\end{enumerate}
\finenonce{001448}



\finexercice



\exercice{1449, hilion, 2003/10/01}

\enonce{001449}{}
Rappel: si $A$ est un anneau (en particulier, si $A$ est un corps), on note $GL_n(A)$ l'ensemble des matrices carrées de dimension $n$ à coefficient dans $A$, qui sont inversibles.
$GL_n(A)$ forme un groupe pour la loi $\times$ de multiplication des matrices, appelé groupe linéaire.
Une matrice carrée de dimension $n$ est dans $GL_n(A)$ ssi son déterminant est un inversible de l'anneau $A$ (ce qui revient à dire, lorsque $A$ est un corps, que son déterminant est non nul).

Pour simplifier, on suppose dans l'exercice que $A$ est un corps, noté $\mathbb{K}$.
\begin{enumerate}
\item Montrer que $\det:GL_n(\mathbb{K})\rightarrow\mathbb{K}^*$ est un morphisme de groupes.
\item On note $SL_n(\mathbb{K})=\ker(\det)$.
Dire pourquoi $SL_n(\mathbb{K})$ est un sous-groupe distingué de $GL_n(\mathbb{K})$ et montrer que $GL_n(\mathbb{K})/SL_n(\mathbb{K})\cong \mathbb{K}^*$.
\item Reconnaître $GL_1(\mathbb{K})$ et $SL_1(\mathbb{K})$.
\item Montrer que les matrices diagonales (resp. triangulaires supérieures) de $GL_n(\mathbb{K})$ forment un sous-groupe. Sont-ils distingués?
\item Montrer que $Z(GL_n(\mathbb{K}))$ est le sous-groupe formé par les homothéties.
\end{enumerate}
\finenonce{001449}



\finexercice


\section{301.00 Ordre d'un élément}


\exercice{2101, debes, 2008/02/12}
\enonce{002101}{}
On dispose d'un \'echiquier et de dominos. Les dominos sont pos\'es sur
l'\'echiquier soit horizontalement, soit verticalement de fa\c con \`a couvrir deux cases
contigu\"es. Est-il possible de couvrir ainsi enti\`erement l'\'echiquier \`a l'exception
des deux cases extr\`emes, en haut \`a gauche et en bas \`a droite?
Reprendre cette question dans le cas o\`u l'on exclut deux cases quelconques \`a la place des deux cases
extr\`emes ci-dessus.
\finenonce{002101} 



\finexercice


\exercice{2102, debes, 2008/02/12}
\enonce{002102}{}
(I) Soit $X$ un ensemble et ${\cal P} (X)$ l'ensemble des parties de $X$ ordonn\'e par l'inclusion. Soit $\varphi$ une application croissante de
${\cal P}(X)$ dans lui-m\^eme.
\smallskip 

(a) Montrer que l'ensemble $E$ des parties $A$ de $X$ qui v\'erifient $\varphi (A)\subset A$
est non vide et admet un plus petit \'el\'ement $A_0$.

\smallskip
(b) Montrer que $\varphi (A_0)=A_0$.

\medskip
(II) Soit deux ensembles $X$ et $Y$ munis de deux injections $g$ de $X$ dans $Y$ et $h$ de $Y$
dans $X$. 
\smallskip

(a) Montrer que l'application de ${\cal P}(X)$ dans lui-m\^eme d\'efini par 
$$\varphi (A)=X-h(Y-g(A))$$
est croissante.

\smallskip
(b) D\'eduire de ce qui pr\'ec\`ede qu'il existe une bijection de $X$ sur $Y$.
\finenonce{002102} 


\finexercice


\exercice{2103, debes, 2008/02/12}
\enonce{002103}{}
Soit $X$ un ensemble non vide et ordonn\'e. Montrer qu'il existe une partie
$Y$ totalement ordonn\'ee de $X$ qui v\'erifie la propri\'et\'e
$$\forall x \notin Y \quad \exists y\in X \quad x\  \hbox {\rm et} \ y \ \hbox
{non comparables}$$

L'ensemble $Y$ est-il unique?
\finenonce{002103} 



\finexercice


\exercice{2104, debes, 2008/02/12}
\enonce{002104}{}
Un jardinier doit planter $10$ arbres en $5$ rang\'ees de $4$ arbres. Donner une disposition
possible. Quel est le nombre minimal d'arbres dont il doit disposer pour planter $6$ rang\'ees de $5$
arbres? G\'en\'eraliser.
\finenonce{002104} 


\finexercice


\exercice{2105, debes, 2008/02/12}
\enonce{002105}{}
Soit $n$ et $p$ deux entiers, $p\leq n$. D\'emontrer, gr\^ace \`a un d\'enombrement, la
formule suivante:
$$\sum _{ 0\leq k\leq p} C_n ^k C_ {n-k}^{p-k} =2^p C_n^p$$
\finenonce{002105} 


\finexercice


\exercice{2106, debes, 2008/02/12}
\enonce{002106}{}
Soit $n$ un entier impair non divisible par $3$. Montrer que $24$ divise $n^2-1$.
\finenonce{002106} 


\finexercice


\exercice{2107, debes, 2008/02/12}
\enonce{002107}{}
On consid\`ere sur $\R$ la loi de composition d\'efinie par $x\star y=
x+y-xy$. Cette loi est-elle associative, commutative? Admet-elle un \'el\'ement neutre?
 Un r\'eel $x$ admet-il un inverse pour cette loi? Donner une formule pour la puissance
$n$-i\`eme d'un \'el\'ement $x$ pour cette loi.
\finenonce{002107} 


\finexercice


\exercice{2108, debes, 2008/02/12}
\enonce{002108}{}
Soit $E$ un mono\" \i de unitaire. On dit qu'un \'el\'ement $a$
de $E$ admet un {\it inverse \`a gauche}  (resp. {\it inverse \`a droite}) s'il existe $b \in
E$ tel que $ba=e$ (resp. $ab=e$). \smallskip

(a) Supposons qu'un \'el\'ement $a$ admette un inverse \`a gauche $b$  qui lui-m\^eme admet un
inverse \`a gauche. Montrer que $a$ est inversible. \smallskip

(b) Supposons que tout \'el\'ement de $E$ admette un inverse \`a gauche. Montrer que $E$ est
un groupe.
\finenonce{002108} 

\finexercice


\exercice{2109, debes, 2008/02/12}
\enonce{002109}{}
Soit $E$ un ensemble muni d'une loi $\star $ associative 

(i) admettant un \'el\'ement neutre \`a gauche $e$ (i.e. $\forall x \in E \quad e\star x=x$)
et

(ii) tel que tout \'el\'ement poss\`ede un inverse \`a gauche (i.e. $\forall x \in E \quad \exists y\in E \quad y\star x =e$).

Montrer que $E$ est un groupe pour la loi $\star$.
\finenonce{002109} 


\finexercice


\exercice{2110, debes, 2008/02/12}
\enonce{002110}{}
Les rationnels non nuls forment-ils un sous-groupe
multiplicatif de $\R^{\times}$?
\finenonce{002110} 


\finexercice


\exercice{2111, debes, 2008/02/12}
\enonce{002111}{}
Montrer que l'ensemble $\{2^n \hskip 2pt | \hskip 2pt n\in \Z \} $ est un sous-groupe
multiplicatif de $\Q^{*}$, ainsi que l'ensemble 
$\{\frac{1+2m}{1+2n}  \hskip 2pt | \hskip 2pt  n,m\in \Z \}$. 
\finenonce{002111}


\finexercice


\exercice{2112, debes, 2008/02/12}
\enonce{002112}{}
Montrer que l'ensemble des matrices carr\'ees \`a $n$ lignes et $n$
colonnes de d\'eterminant non nul est un groupe pour la multiplication.
\finenonce{002112} 


\finexercice


\exercice{2113, debes, 2008/02/12}
\enonce{002113}{}
On consid\`ere l'ensemble $E$ des matrices carr\'ees \`a coefficients
r\'eels de la forme 
%$$\pmatrix { a &0 \cr b&  0 \cr },\quad  a \in \R ^\times ,\quad b
%\in\R$$ 
$$\left[
\begin{array}{cc}
a&0\\
b&0\\
\end{array}
\right] ,\quad  a \in \R ^\times ,\quad b
\in\R$$
muni du produit des matrices. \smallskip
 
(a) Montrer que $E$ est ainsi muni d'une loi de composition interne associative.
\smallskip

(b) D\'eterminer tous les \'el\'ements neutres \`a droite de $E$.
\smallskip

(c) Montrer que $E$ n'admet pas d'\'el\'ement neutre \`a gauche.
\smallskip

(d) Soit $e$ un \'el\'ement
neutre \`a droite. Montrer que tout \'el\'ement de $E$ 
poss\`ede un inverse \`a gauche pour cet \'el\'ement neutre, i.e. 
$$\forall g\in E\quad \exists h \in E \quad hg=e$$

\finenonce{002113} 

%correction
%(1) Pour $a,a^\prime \in \R^\times$ et $b,b^\prime \in \R$, 
%on a $\displaystyle\left[\matrix{
%a & 0 \cr 
%b & 0 \cr
%}\right]\ \left[\matrix{
%a^\prime & 0 \cr 
%b^\prime & 0 \cr
%}\right] = \left[\matrix{
%aa^\prime & 0 \cr 
%ba^\prime & 0 \cr
%}\right]$.
%fincorrection

\finexercice


\exercice{2114, debes, 2008/02/12}
\enonce{002114}{}
\label{ex:deb14}
 Soit $G$ un groupe v\'erifiant 
$$\forall x \in G  \quad x^2=e$$
Montrer que $G$ est commutatif. D\'eduire que si $G$ est fini, alors l'ordre de $G$ est une
puissance de $2$.
\finenonce{002114} 
\finexercice


\exercice{2115, debes, 2008/02/12}
\enonce{002115}{}
Soit $G$ un groupe d'ordre pair. Montrer qu'il existe un \'el\'ement $x\in
G$, $x \not= e$ tel que $x^2=e$.
\finenonce{002115} 


\finexercice


\exercice{2116, debes, 2008/02/12}
\enonce{002116}{}
Soit $G$ un groupe d'ordre impair. Montrer que l'application $f$ de $G$ sur lui-m\^eme
donn\'ee par $f(x)=x^{2}$ est une bijection. En d\'eduire que l'\'equation $x^2=e$ a une
unique solution, \`a savoir $x=e$.
\finenonce{002116} 


\finexercice


\exercice{2117, debes, 2008/02/12}
\enonce{002117}{}
 Soient $G$ un groupe fini et $m$ un entier premier \`a l'ordre de $G$. 
Montrer que pour tout $a\in G$ l'\'equation $x^m=a$ admet une unique solution.
\finenonce{002117} 


\finexercice


\exercice{2118, debes, 2008/02/12}
\enonce{002118}{}
Soit $G$ un groupe et $H<G$, $K<G$ deux sous-groupes de $G$. On suppose
qu'il existe deux \'el\'ements $a,b \in G$ tels que $Ha \subset Kb$. Montrer que $H<K$.
\finenonce{002118} 

\finexercice


\exercice{2119, debes, 2008/02/12}
\enonce{002119}{}
Soit $H$ une partie non vide d'un groupe $G$. On pose $H^{-1} = \{ x^{-1} ;
x \in H \} $. Montrer les \'equivalences suivantes:\smallskip

(a) $H<G \Leftrightarrow HH^{-1} \subset H$
\smallskip

(b) $H<G \Leftrightarrow \forall a \in H \quad Ha=H$.
\finenonce{002119} 


\finexercice


\exercice{2120, debes, 2008/02/12}
\enonce{002120}{}
Soit $G$ un groupe et $H,K$ deux sous-groupes de $G$.\smallskip

(a) Montrer que $H\cup K$ est un sous-groupe de $G$ si et seulement si $H<K$ ou $K<H$.
\smallskip

(b) Montrer qu'un groupe ne peut \^etre la r\'eunion de deux sous-groupes propres.
\finenonce{002120} 

\finexercice


\exercice{2121, debes, 2008/02/12}
\enonce{002121}{}
Montrer que dans un groupe $G$, toute partie non vide finie stable par la
loi de composition est un sous-groupe. Donner un contre-exemple \`a la propri\'et\'e
pr\'ec\'edente dans le cas d'une partie infinie.
\finenonce{002121} 

\finexercice


\exercice{2122, debes, 2008/02/12}
\enonce{002122}{}
(a) Montrer que les seuls sous-groupes de $\Z $ sont de la forme $n\Z $ o\`u $n$
est un entier.
\smallskip

(b) Un \'el\'ement $x$ d'un groupe est dit d'ordre fini s'il existe un entier $k$ tel que
$x^k=e_G$. Montrer que $\{ k\in \Z \hskip 2pt |\hskip 2pt  x^k=e_G \} $ est alors un
sous-groupe non nul de $\Z$. On appelle ordre de $x$ le g\'en\'erateur positif de ce
sous-groupe.

\smallskip
(c) Soit $x$ un \'el\'ement d'un groupe $G$. Montrer que $x$ est d'ordre $d$ si et seulement
si le sous-groupe $< x >$ de $G$ engendr\'e par $x$ est d'ordre $d$.
\finenonce{002122}


\finexercice


\exercice{2123, debes, 2008/02/12}
\enonce{002123}{}
On pose $SL_2 (\Z )= \{\left[\begin{array}{cc}
a&b\\
c&d\\
\end{array}
\right] \hskip 2pt | \hskip 2pt a,b,c,d \in \Z, ad-bc=1 \}$. \smallskip

(a) Montrer que $SL_2 ( \Z )$ est un sous-groupe du groupe des
matrices inversibles \`a coefficients dans $\Z$. \smallskip

(b) On consid\`ere les deux matrices 
$$\left[
\begin{array}{cc}
0&-1\\
1&0\\
\end{array}
\right] \quad  \left[
\begin{array}{cc}
0&1\\
-1&-1\\
\end{array}
\right]$$
%$$A= \pmatrix { 0&-1 \cr 1&0\cr } \quad B= \pmatrix { 0&1 \cr -1&-1 \cr } $$
D\'emontrer que $A$ et $B$ sont d'ordres finis mais que $AB$ est d'ordre infini.
\finenonce{002123} 


\finexercice


\exercice{2124, debes, 2008/02/12}
\enonce{002124}{}
\label{ex:deb24}
Soit $G$ un groupe ab\'elien et $a$ et $b$ deux \'el\'ements d'ordres finis.
Montrer que $ab$ est d'ordre fini et que l'ordre de $ab$ divise le ppcm des ordres de $a$
et $b$. Montrer que si les ordres de $a$ et $b$ sont premiers entre eux, l'ordre de $ab$
est \'egal au ppcm des ordres de $a$ et de $b$.
\finenonce{002124} 

\finexercice


\exercice{2125, debes, 2008/02/12}
\enonce{002125}{}
Soit $G$ un groupe commutatif. Montrer que l'ensemble des \'el\'ements
d'ordre fini de $G$ forme un sous-groupe de $G$.
\finenonce{002125} 


\finexercice


\exercice{2126, debes, 2008/02/12}
\enonce{002126}{}
D\' eterminer tous les sous-groupes de $\mu_2 \times \mu_2$.
\finenonce{002126} 


\finexercice


\exercice{2127, debes, 2008/02/12}
\enonce{002127}{}
 Soient $G$ un groupe fini et commutatif et $\{ G_i\} _{i\in I }$ la famille des sous-groupes
propres maximaux de $G$. On pose $F=\bigcap _{i\in I} G_i$. Montrer que $F$ est
l'ensemble des \'el\'ements $a$ de $G$ qui sont tels que, pour toute partie $S$ de
$G$ contenant $a$ et engendrant $G$, $S-\{ a \}$ engendre encore $G$.
\finenonce{002127} 

\finexercice


\exercice{2128, debes, 2008/02/12}
\enonce{002128}{}
D\'eterminer tous les groupes d'ordre $\leq 5$. En d\'eduire qu'un groupe non 
commutatif poss\`ede au moins 6 \'el\'ements. Montrer que le groupe sym\'etrique 
$S_3$ est non commutatif.
\finenonce{002128} 


\finexercice


\exercice{2129, debes, 2008/02/12}
\enonce{002129}{} 
Le centre d'un groupe $G$ est l'ensemble $Z(G)$ des \'el\'ements de $G$ qui commutents \`a tous les \'el\'ements de $G$. V\'erifier que $Z(G)$ est un sous-groupe ab\'elien de $G$. Montrer que 
si $G$ poss\`ede un unique \'el\'ement d'ordre $2$, alors cet \'el\'ement est dans le centre $Z(G)$.
\finenonce{002129} 


\finexercice


\exercice{2130, debes, 2008/02/12}
\enonce{002130}{}
Soient $G$ un groupe et $H$ et $K$ deux sous-groupes de $G$.

(a) Montrer que l'ensemble $HK=\{xy \hskip 2pt | \hskip 2pt x\in H, y\in K\}$ est un sous-groupe de $G$ si et seulement si $HK=KH$.

(b) Montrer que si $H$ et $K$ sont finis alors $\displaystyle |HK| = \frac{|H|\cdot |K|}{|H\cap K|}$.
\finenonce{002130} 

\finexercice


\exercice{2131, debes, 2008/02/12}
\enonce{002131}{}
D\' eterminer tous les sous-groupes du groupe sym\'etrique $S_3$.
\finenonce{002131} 

\finexercice


\exercice{2132, debes, 2008/02/12}
\enonce{002132}{}
\label{ex:deb32}
Montrer que dans un groupe d'ordre $35$, il existe un \'el\'ement d'ordre $5$ et un \'el\'ement d'ordre $7$. 
\finenonce{002132} 


\finexercice


\exercice{2133, debes, 2008/02/12}
\enonce{002133}{}
Soit $G$ un groupe d'ordre $2p$ avec $p$ un nombre premier. Montrer qu'il existe un \'el\'ement d'ordre $2$ et un \'el\'ement d'ordre $p$. 
\finenonce{002133} 


\finexercice


\exercice{2134, debes, 2008/02/12}
\enonce{002134}{}
Soient $n\geq 0$ un entier et $p$ un nombre premier tels que $p$ divise $2^{2^n}+1$. Montrer que $p$ est de la forme $p= k2^{n+1}+1$ o\`u $k$ est un entier.
\finenonce{002134} 


\finexercice


\exercice{2135, debes, 2008/02/12}
\enonce{002135}{}
Montrer que tout entier $n>0$ divise toujours $\varphi(2^n-1)$ (o\`u $\varphi$ est la fonction indicatrice d'Euler).
\finenonce{002135} 


\finexercice


\exercice{2661, matexo1, 2002/02/01}
\enonce{002661}{}
Soit $G$ un groupe multiplicatif (c'est-{\`a}-dire dont la loi est
not{\'e}e multiplicativement). Soient $a$ et $b$ deux {\'e}l{\'e}ments de $G$. Montrer
que si $ab$ est d'ordre fini, alors $ba$ l'est {\'e}galement et son ordre est
celui de $ab$.
\finenonce{002661}

\finexercice


\exercice{2662, matexo1, 2002/02/01}
\enonce{002662}{}
Montrer que les {\'e}l{\'e}ments d'ordre fini d'un groupe commutatif $G$ forment
un sous-groupe de $G$. En est-il de m{\^e}me si $G$ n'est pas commutatif?
\finenonce{002662}

\finexercice


\exercice{2663, matexo1, 2002/02/01}
\enonce{002663}{}
Soit $G$ un groupe commutatif multiplicatif, $a$ et $b$ deux {\'e}l{\'e}ments de $G$ d'ordres
$n$ et $m$. Que peut-on dire de l'ordre de $ab$? Que peut-on dire de plus si l'intersection
des sous-groupes $G_a$ et $G_b$ engendr{\'e}s par $a$ et $b$ est r{\'e}duite {\`a} $\{1_G\}$?
\finenonce{002663}

\finexercice


\exercice{2664, matexo1, 2002/02/01}
\enonce{002664}{}
Montrer qu'un groupe fini dont l'ordre est un nombre premier est
cyclique.
\finenonce{002664}

\finexercice


\exercice{2665, matexo1, 2002/02/01}
\enonce{002665}{}
Soient $\sigma$ et $\tau$ deux transpositions de $\{1,\dots,n\}$.
Montrer que $\sigma \circ \tau$ est d'ordre 1, 2 ou 3.
\finenonce{002665}

\finexercice

\section{302.00 Groupe symétrique, décomposition en cycles disjoints, signature}


\exercice{2666, matexo1, 2002/02/01}
\enonce{002666}{}
Pour tout $n\in\N^*$, calculer la signature de la permutation
$[n\ n-1\ n-2 \dots\ 3\ 2\ 1]\in S_n$.
\finenonce{002666}

\finexercice


\exercice{2667, matexo1, 2002/02/01}
\enonce{002667}{}
Soit $n\in\N^*$. Calculer la somme des nombres d'inversions de toutes
les permutations de $\{1,2,\ldots,n\}$.
\finenonce{002667}

\finexercice


\exercice{2668, matexo1, 2002/02/01}
\enonce{002668}{}
Matrices de permutation\par
Soit $n\in\N^*$ et $K$ un corps. Montrer que l'application
$$\begin{aligned}
\phi: S_n&\to M_n(K)\cr
\sigma&\mapsto(\delta_{i,\sigma(j)})_{1\leq i,j\leq n}
\end{aligned}
$$
o{\`u} $\delta_{k,l}=\left\{\begin{matrix}1\hbox{ si }k=l\cr 0\hbox{ si }k\neq l\end{matrix}\right.$ (symbole de Kronecker)
induit un morphisme de groupes de $S_n$ dans $GL_n(K)$.
\finenonce{002668}

\finexercice


\exercice{2669, matexo1, 2002/02/01}
\enonce{002669}{}
Soit $\sigma\in S_n$ et $c=(a_1\ a_2\ \dots\ a_k)$ un cycle. Quelle est la nature de la
permutation $\sigma\circ c \circ\sigma^{-1}$?
\finenonce{002669}

\finexercice


\exercice{2670, matexo1, 2002/02/01}
\enonce{002670}{}
Expliciter les 24 rotations de l'espace laissant un cube de sommets
$A_1, A_2, \dots, A_8$ invariant.\par
D{\'e}composer en cycles les permutations de $S_8$ correspondantes.\par
Ecrire les produits "typiques" de 2 quelconques de ces permutations.
\finenonce{002670}

\finexercice


\exercice{7872, mourougane, 2021/08/11}
\enonce{007872}{Les cycles d'ordre $3$ engendrent $\mathcal{A}_n$}
Soit $n$ un entier naturel supérieur à $3$, $\mathcal{S}_n$
le groupe symétrique de $n$ lettres et $\mathcal{A}_n$ le groupe alterné. 
\begin{enumerate}
 \item Calculer les produits de transpositions $(a,b)(b,c)$, puis $(a,b)(c,d)$.
 \item Montrer que les cycles d'ordre $3$ engendrent $\mathcal{A}_n$.
\end{enumerate}
\finenonce{007872}
\finexercice


\exercice{7873, mourougane, 2021/08/11}
\enonce{007873}{Groupes symétriques}
Soint $n$ un entier naturel supérieur à $3$.
\begin{enumerate}
 \item Montrer que les permutations $(i,j)(j,k)$ et $(i,j)(k,l)$ s'écrivent comme produit de $3$-cycles.
 \item En déduire que le groupe alterné $\mathcal{A}_n$ est engendré par les $3$-cycles.
\item Montrer que si $n\geq 5$, tous les $3$-cycles sont conjugués dans $\mathcal{A}_n$.
\end{enumerate} 
\finenonce{007873}
\finexercice


\exercice{7874, mourougane, 2021/08/11}
\enonce{007874}{}
\begin{enumerate}
 \item Le groupe $\mathcal{S}_n$ est-il simple ?
 \item Le groupe $\Z/89\Z$ est-il simple ?
 \item Le groupe $\Z/221\Z$ est-il simple ?
\end{enumerate}
\finenonce{007874}
\finexercice


\exercice{7875, mourougane, 2021/08/11}
\enonce{007875}{Simplicité de $\mathcal{A}_5$}
\begin{enumerate}
 \item Faire la liste des classes de conjugaison de $\mathcal{S}_n$ dans $\mathcal{A}_n$ en les dénombrant.
\item Montrer que les $3$-cycles sont conjugués dans $\mathcal{A}_n$.
\item Montrer que les éléments d'ordre $2$ sont conjugués dans $\mathcal{A}_n$.
\item Montrer que tout sous-groupe distingué $H$ de $\mathcal{A}_n$ qui contient un élément d'ordre $5$ les contient tous. (On remarquera que le groupe engendré par un élément d'ordre $5$ est un Sylow.)
\item Montrer que tout sous-groupe distingué $H$ de $\mathcal{A}_n$ non réduit à $\{id\}$ contient au moins deux types d'éléments en plus de l'identité. Montrer alors que $H=\mathcal{A}_n$.
\end{enumerate}
\finenonce{007875}
\finexercice

\section{303.00 Sous-groupe distingué}


\exercice{1428, legall, 1998/09/01}

\enonce{001428}{}
 Soit $  G  $ un groupe, $  H  $ et $  K  $ deux
sous-groupes d'ordre fini de $  G  $ tels que $  H\cap K= \{ e_G\}
.$
\begin{enumerate}
    \item Montrer que le cardinal de $  HK  $ est \'egal  $  \vert H\vert \vert K\vert   .$
    \item En d\'eduire que si $  \vert G\vert =pq  $ o\`u $  p   $ est premier
et $  p>q  $ alors $  G  $ a au plus un sous-groupe d'ordre $  p  .$ Montrer que si ce
sous-groupe existe il est distingu\'e dans $  G  .$
\end{enumerate}

\finenonce{001428}



\finexercice



\exercice{1429, legall, 1998/09/01}

\enonce{001429}{}
Soit $  G  $ un groupe, $  A  $ une partie non vide de $  G  .$
On note $  N(A)=\{ g\in G  ;   gAg^{-1}=A\}  $ et $  C(A)=\{ g\in G  ; \forall a \in A  ;
gag^{-1}=a\}  .$ Montrer que $  N(A)   $ et $  C(A)  $ sont des sous-groupes de $  G  $ et que $  C(A)  $ est un sous-groupe distingu\'e de $  N(A)  .$
\finenonce{001429}



\finexercice



\exercice{1430, legall, 1998/09/01}

\enonce{001430}{}
Soit $  G  $ un groupe, $  H  $ et $  K  $ deux sous-groupes de $
G  .$ On note $  HK=\{ hk  ;   h\in H   ,   k\in K\}  .$
\begin{enumerate}
    \item Montrer que $  HK  $ est un sous-groupe de $  G  $ si et seulement si $  HK=KH  .$ En d\'eduire que si $  H  $ est distingu\'e dans $ G  $ alors $  HK  $ est un sous-groupe de $
G  .$
    \item On suppose d\'esormais que $  \forall  h\in H   ,   k\in K   :  hk=kh  .$
Montrer que l'application $  f : H\times K\rightarrow G   $ d\'efinie par $
\forall  h\in H   ,   k\in K   :  f(h,k)=hk  $ est un homomorphisme de groupes.
    \item Calculer le noyau et l'image de $  f  .$ Donner une condition n\'ec\'essaire et suffisante pour que $  f  $ soit un isomorphisme de groupes.
\end{enumerate}
\finenonce{001430}



\finexercice



\exercice{1431, legall, 1998/09/01}

\enonce{001431}{}
\begin{enumerate}
    \item Soit $  G  $ un groupe, $  H  $ un sous-groupe de $  G  .$
Montrer que les propri\'et\'es suivantes sont \'equivalentes~:
\vskip1mm \hskip2mm {\em i)} $  \forall g \in G  :  gHg^{-1}\subset H  .$
\vskip1mm \hskip2mm {\em ii)} $  \forall g \in G  :  gHg^{-1}= H  .$
\vskip1mm \hskip2mm {\em iii)} $  \forall g \in G  :  gH=Hg  .$
    \item En d\'eduire que tout sous-groupe d'indice $  2  $ est distingu\'e.
\end{enumerate}
\finenonce{001431}



\finexercice



\exercice{1432, ortiz, 1999/04/01}

\enonce{001432}{}
Soient $T=\left\{\left(\begin{smallmatrix}
a&b\\0&c\end{smallmatrix}\right) :
a,c\in\Rr\setminus\left\{0\right\},b\in
\Rr\right\}$ et $U=\left\{\left(\begin{smallmatrix} 1&b\\0&1\end{smallmatrix}\right) : b\in
\Rr\right\}.$\\
\begin{enumerate}
    \item  Montrer que $T$ est un sous-groupe de $\text{GL}_2(\Rr).$
    \item  Montrer que $U$ est un sous-groupe distingu\'e de $T.$
\end{enumerate}
\finenonce{001432}



\finexercice



\exercice{1433, hilion, 2003/10/01}

\enonce{001433}{}
Soit $G$ un groupe.
\begin{enumerate}
\item Un sous-groupe $H$ de $G$ est distingué si: $\forall x\in G, xH=Hx$, ce qui est équivalent à dire que $H$ est le noyau d'un morphisme de $G$ dans un groupe.
Rappeler la démonstration de cette équivalence.
\item Si $H$ est un sous-groupe d'indice 2 de $G$, montrer que $H$ est distingué.
\item Si $G$ est abélien, montrer que tout sous-groupe de $G$  est distingué.
\item Le centre de $G$ est l'ensemble $Z(G)=\{z\in G: \forall x\in G, xz=zx\}$. Montrer que $Z(G)$ est un sous-groupe distingué.
\end{enumerate}
\finenonce{001433}



\finexercice



\exercice{2136, debes, 2008/02/12}
\enonce{002136}{}
Soit $G$ un groupe tel que l'application $x\rightarrow x^{-1}$ soit un morphisme.
Montrer que $G$ est commutatif. 
\finenonce{002136} 


\finexercice


\exercice{2137, debes, 2008/02/12}
\enonce{002137}{} 
Soient $G$ un groupe et $n\geq 1$ un entier tels que
l'application $x\rightarrow x^n$ soit un automorphisme de $G$.
Montrer que pour tout \'el\'ement $x$ de $G$, $x^{n-1}$
appartient au centre de $G$.
\finenonce{002137} 

\finexercice


\exercice{2138, debes, 2008/02/12}
\enonce{002138}{}
Montrer que le groupe des automorphismes du groupe $\Z /2\Z \times \Z /2 \Z$ est
isomorphe au groupe sym\'etrique $S_3$.
\finenonce{002138} 

\finexercice


\exercice{2139, debes, 2008/02/12}
\enonce{002139}{}
Montrer qu'un sous-groupe d'indice $2$ dans un groupe
$G$ est distingu\'e dans $G$.
\finenonce{002139} 


\finexercice


\exercice{2140, debes, 2008/02/12}
\enonce{002140}{} 
Soit $G$ un groupe et $H$ un sous-groupe. On suppose que le produit de deux classes 
\`a gauche modulo $H$ est une classe \`a gauche modulo $H$. Montrer que $H$ est
distingu\'e  dans $G$.
\finenonce{002140} 


\finexercice


\exercice{2141, debes, 2008/02/12}
\enonce{002141}{}
Soit $G$ un groupe et $\simeq $ une relation d'\'equivalence sur $G$. On suppose que
cette relation est compatible avec la loi de groupe, c'est-\`a-dire que
$$\forall x,y \in G \quad \forall x',y' \in G \quad  x\simeq x' \quad {\rm et}\quad 
y\simeq y' \quad {\rm alors} \quad  xy\simeq x'y' $$
Montrer que la classe $H$ de l'\'el\'ement neutre $1$ est un sous-groupe distingu\'e de
$G$ et que  $$\forall x, x' \in G \quad x\simeq x' \quad \hbox{\rm est\quad
\'equivalent \quad  \`a} \quad x'x^{-1} \in H$$
\finenonce{002141} 


\finexercice


\exercice{2142, debes, 2008/02/12}
\enonce{002142}{}
\label{ex:le7}
Soit $G$ un groupe et $K\subset H\subset G$ deux sous-groupes. On
suppose que $H$ est distingu\'e  dans $G$ et que $K$ est caract\'eristique dans $H$ (i.e.
stable par tout automorphisme de $H$). Montrer qu'alors $K$ est distingu\'e dans
$G$.

Donner un exemple de groupe $G$ et de deux sous-groupes $K\subset H \subset G$, $H$ \'etant
distingu\'e dans $G$ et $K$ \'etant distingu\'e dans $H$, mais $K$ n'\'etant pas distingu\'e dans
$G$.
\finenonce{002142} 


\finexercice


\exercice{2143, debes, 2008/02/12}
\enonce{002143}{}
(a) Montrer que pour tous entiers premiers entre eux $m,n>0$, les deux groupes $(\Zz/mn\Zz)^\times$ et $(\Zz/m\Zz)^\times \times (\Zz/n\Zz)^\times$ sont isomorphes. En d\'eduire que $\varphi (mn) = \varphi (m) \varphi (n)$, o\`u $\varphi$ est la fonction indicatrice d'Euler.
\smallskip

(b) Le groupe multiplicatif $(\Zz/15\Zz)^\times$ est-il cyclique? Montrer que $(\Zz/8\Zz)^\times \simeq (\Zz/2\Zz) \times (\Zz/2\Zz)$, que $(\Zz/16\Zz)^\times \simeq (\Zz/4\Zz)\times (\Zz/2\Zz)$. Etudier le groupe multiplicatif $(\Zz/24\Zz)^\times$.
\finenonce{002143} 


\finexercice


\exercice{2144, debes, 2008/02/12}
\enonce{002144}{}
(a) Montrer que si $m$ et $n$ sont des entiers premiers entre eux et qu'un
\'el\'ement $z$ d'un groupe $G$ v\'erifie $z^m=z^n=e$ o\`u $e$ d\'esigne l'\'el\'ement
neutre de $G$, alors
$z=e$.
\smallskip

(b) Montrer que si $m$ et $n$ sont deux entiers premiers entre eux, l'application
$$\phi: \mu_m \times \mu_n \rightarrow \mu_{mn} $$ qui au couple $(s,t)$ fait correspondre le
produit
$st$ est un isomorphisme de groupes
\finenonce{002144} 


\finexercice


\exercice{2145, debes, 2008/02/12}
\enonce{002145}{}
Montrer que les groupes $\mu_4$ et $\mu_2 \times \mu_2 $ ne sont pas
isomorphes. De fa\c con g\' en\' erale montrer que si $m$ et $n$ sont des entiers qui ne
sont pas premiers entre eux, les groupes $\mu_{mn}$ et $\mu_m \times \mu_n$ ne sont pas
isomorphes.
\finenonce{002145} 

\finexercice


\exercice{2146, debes, 2008/02/12}
\enonce{002146}{}
Soit $n$ et $d$ deux entiers tels que $d$ divise $n$. On d\' efinit une
application $f : \mu_n \rightarrow  \mu_{d} $
qui \`a $s$ associe $s^{n/d}$. Montrer que $f$ est un morphisme surjectif de groupes dont
le noyau est $\mu_{n/d} $.
\finenonce{002146} 


\finexercice


\exercice{2147, debes, 2008/02/12}
\enonce{002147}{}
\label{ex:le12}
Soit $f:G \rightarrow H$  un morphisme de groupes finis. Soit $G^\prime$
un sous-groupe de $G$. Montrer que l'ordre de $f(G^\prime)$ divise les ordres de
$G^\prime$ et de $H$.
\finenonce{002147} 


\finexercice


\exercice{2148, debes, 2008/02/12}
\enonce{002148}{}
Soit $f:G  \rightarrow H$  un morphisme de groupes finis. Soit $G^\prime$ un
sous-groupe de $G$ d'ordre premier \`a l'ordre de $H$. Montrer que $G^\prime\subset
\hbox{\rm ker}(f)
$.
\finenonce{002148} 




\finexercice


\exercice{2149, debes, 2008/02/12}
\enonce{002149}{}
Soit $G$ un groupe fini et $H$ et $K$ deux sous-groupes de $G$. On suppose que $H$ est
distingu\'e dans $G$, que $|H|$ et $|G/H|$ sont premiers entre eux et 
$|H|=|K|$. Montrer que $H=K$.
\finenonce{002149} 



\finexercice


\exercice{2150, debes, 2008/02/12}
\enonce{002150}{}
Soit $f$ un morphisme de groupes $f : \Q  \rightarrow \Q ^{\times}
_{>0}$, $\Q$ \'etant muni de l'addition et $\Q ^{\times} _{>0}$ muni de la
multiplication. Calculer $f(n)$ en fonction de $f(1)$ pour tout entier $n>0$. Montrer
que les deux groupes pr\'ec\'edents ne sont pas isomorphes.
\finenonce{002150} 

\finexercice


\exercice{2151, debes, 2008/02/12}
\enonce{002151}{}
Trouver tous les morphismes du groupe additif $\Q$ dans lui m\^eme. 

M\^eme question de $\Q $ dans $\Z$.

M\^eme question de $\Z /m \Z$ dans $\Z$.
\finenonce{002151} 


\finexercice


\exercice{2152, debes, 2008/02/12}
\enonce{002152}{}
\label{ex:le17}
Etant donn\'es deux entiers $m, n>0$, d\' eterminer tous les morphismes de groupe de $\Zz/m\Zz$ dans $\Zz/n\Zz$, puis tous les automorphismes de $\Zz/n\Zz$.
\finenonce{002152} 


\finexercice


\exercice{2153, debes, 2008/02/12}
\enonce{002153}{}
\label{ex:le18}
Soit $G$ un groupe et $H$ un sous groupe distingu\'e de $G$ d'indice
$n$. Montrer que pour tout $a\in G$, $a^n \in H$. Donner un exemple de sous-groupe $H$
non distingu\'e de $G$ pour lequel la conclusion pr\'ec\'edente est fausse.
\finenonce{002153} 

\finexercice


\exercice{2154, debes, 2008/02/12}
\enonce{002154}{}
\label{ex:le19}
Soit $G$ un groupe fini et $H$ un sous-groupe distingu\'e d'ordre $n$ et d'indice $m$.
On suppose que $m$ et $n$ sont premiers entre eux. Montrer que $H$ est   
l'unique sous-groupe de $G$ d'ordre $n$. 
\finenonce{002154} 


\finexercice


\exercice{2155, debes, 2008/02/12}
\enonce{002155}{}
Montrer que $\textrm{SL}_n(\Rr)$ est un sous-groupe distingu\'e du groupe $\textrm{GL}_n(\Rr)$ et que le groupe quotient est isomorphe \`a $\Rr^\times$.
\finenonce{002155} 


\finexercice


\exercice{2156, debes, 2008/02/12}
\enonce{002156}{}
On consid\`ere les groupes suivants :
$$T= \{z\in\C \hskip 2pt | \hskip 2pt |z|=1\} \quad \mu _n = \{ z \in \C \hskip 2pt |
\hskip 2pt z ^n =1 \}
\quad 
\mu _\infty =
\{ z
\in \C \hskip 2pt | \hskip 2pt \exists n \quad z ^n =1 \}$$

(a) Montrer les isomorphismes suivants :
$${\rm } \R /\Z \simeq T \quad {\rm } \C ^\times  / \R ^{\times } _{>0}\simeq T \quad {\rm } \C ^{\times } /\R ^{\times } \simeq T  \quad {\rm } T/ \mu _n \simeq T \quad {\rm } \C ^{\times } /\mu _n \simeq \C ^{\times } $$

(b) Montrer que $\mu _ \infty \simeq \Q /\Z $. Quels sont les sous-groupes finis de $\mu
_\infty $? 

\smallskip
(c) Montrer qu'un sous-groupe de type fini de $\Q $ contenant $\Z$ est de la forme
$\frac{1}{q} \Z$. En d\'eduire la forme des sous-groupes de type fini de $\Q /\Z $ et de $\mu
_\infty$.

\smallskip
(d) Soit $p$ un nombre premier. Montrer que $\mu _{p^\infty }= \{ z \in \C \hskip 2pt |
\hskip 2pt \exists n \in \N \quad z^{p ^n }=1 \}$ est un sous-groupe de $\mu _\infty $.
Est-il de type fini?
\finenonce{002156} 

\finexercice


\exercice{2157, debes, 2008/02/12}
\enonce{002157}{}
Soit $G$ un sous-groupe d'indice fini du groupe multiplicatif $\C ^\times $. Montrer que $G=
\C ^\times$.
\finenonce{002157} 

\finexercice


\exercice{2158, debes, 2008/02/12}
\enonce{002158}{}
\label{ex:le23}
Soit  $G$  un groupe et  $H$  un sous-groupe contenu dans le centre  $Z(G)$ de $G$. Montrer
que $H$ est distingu\'e dans $G$ et que, si le groupe quotient $G/H$ est cyclique, 
$G=Z(G)$.
\finenonce{002158} 


\finexercice


\exercice{2159, debes, 2008/02/12}
\enonce{002159}{}
\label{ex:le24}
Montrer qu'un groupe d'ordre $p^2$ o\`u $p$ est un nombre premier est
ab\'elien. (On utilisera que le centre d'un $p$-groupe est non trivial, ce qui est une cons\'equence classique de la ``formule des classes'' (voir chapitre suivant)).
\finenonce{002159} 


\finexercice


\exercice{2160, debes, 2008/02/12}
\enonce{002160}{}
\label{ex:le25}
(a) Soit $p$ un nombre premier. Montrer que tout morphisme de groupes entre $\mathbb{F}_p ^n$
et $\mathbb{F}_p^m$ est une application $\mathbb{F}_p$ -lin\'eaire.

\smallskip
(b) Montrer que le groupe des automorphismes de $\Z/p\Z$ est isomorphe au groupe
multiplicatif $\mathbb{F}_p^{\ast}$.

\smallskip
(c) D\'eterminer le nombre d'automorphismes de $\mathbb{F}_p^n$.
\finenonce{002160} 

\finexercice


\exercice{2161, debes, 2008/02/12}
\enonce{002161}{}
D\'eterminer le centre du groupe $GL_n (\mathbb{F}_p)$ des automorphismes de $(\mathbb{F}_p)^n$.
\finenonce{002161} 


\finexercice


\exercice{2162, debes, 2008/02/12}
\enonce{002162}{}
Soit $p$ un nombre premier. Montrer qu'un groupe ab\'elien fini, dont
tous les \'el\'ements diff\'erents de l'\'el\'ement neutre sont d'ordre $p$, est
isomorphe \`a $(\Z /p
\Z)^n $.
\finenonce{002162} 

\finexercice


\exercice{2163, debes, 2008/02/12}
\enonce{002163}{}
(a) Soit $G$ un groupe et $H$ un sous-groupe distingu\'e de $G$. On note $\varphi $ la
surjection canonique $\varphi : G \rightarrow G/H$. Montrer que l'ordre d'un \'el\'ement $x$ de $G$
est un multiple de l'ordre de $\varphi (x)$.

\smallskip
(b) Pour tout $x\in G$ on pose $\tau _x$ l'application de $G$ dans $G$ d\'efinie par $\tau
_x(y)=xyx^{-1}$. Montrer que $\tau _x$ est un automorphisme de $G$ et que l'application 
$$x\rightarrow \tau _x$$  est un morphisme de groupes de $G$ dans $\hbox{\rm Aut}(G)$. Quel
est le noyau de ce morphisme?

\smallskip

(c) On suppose que $G$ est fini et que $H$ est un sous-groupe distingu\'e dont l'ordre est le
plus petit nombre premier $p$ divisant l'ordre de $G$. Montrer que pour tout $x\in G$ l'ordre
de la restriction \`a $H$ de $\tau _x$ est un diviseur de $p-1$ et de l'ordre de $G$. En
d\'eduire que $\tau _x$ restreint \`a $H$ est l'identit\'e pour tout $x$ et donc que $H$ est
contenu dans le centre de $G$.
\finenonce{002163} 


\finexercice


\exercice{2164, debes, 2008/02/12}
\enonce{002164}{}
\label{ex:le29}
Soit $G$ un groupe. On appelle groupe des commutateurs de $G$ et l'on note
$D(G)$ le sous-groupe de $G$ engendr\'e par les \'el\'ements de la forme
$xyx^{-1}y^{-1}$. Montrer que $D(G)$ est distingu\'e dans $G$ et que le quotient
$G/D(G)$ est ab\'elien. Montrer que $D(G)$ est le plus petit sous-groupe
distingu\'e de $G$ tel que le quotient de $G$ par ce sous-groupe soit ab\'elien.
\finenonce{002164} 


\finexercice


\exercice{2165, debes, 2008/02/12}
\enonce{002165}{}
Soit $G$ un groupe d'ordre $p^3$ o\`u $p$ est un
nombre premier. Montrer que si $G$ n'est pas commutatif, $Z(G)=D(G)$
et que ce sous-groupe est d'ordre $p$.
\finenonce{002165} 

\finexercice

\exercice{7776, mourougane, 2021/08/11}
\enonce{007776}{}
\begin{enumerate}
 \item Soit $f~: \mathcal{S}_n \to A$ un homomorphisme de groupes de $\mathcal{S}_n$ vers un groupe abélien.
   Démontrer que les transpositions ont toutes même image.
   Démontrer que si $A=\{1,-1\}$ $f=signature$ ou $f = \textrm{application constante} 1$.
 \item Soit $G$ d'indice 2 dans $\mathcal{S}_n$.
   Démontrer que $G$ est distingué (reprendre la méthode de la feuille 1) puis que $G=\mathcal{A}_n$.
\end{enumerate}
\finenonce{007776}
\finexercice


\exercice{7777, mourougane, 2021/08/11}
\enonce{007777}{Les sous-groupes distingués de $\mathcal{S}_n$}
Le but de l'exercice est de déterminer les sous-groupes distingués de
$\mathcal{S}_n$ (pour $n\geq 5$).
\begin{enumerate}
\item Soit $H$ un sous-groupe distingué de $\mathcal{S}_n$.
Montrer que $H\cap \mathcal{A}_n$ est un sous-groupe distingué de $\mathcal{A}_n$.
En déduire que $H$ contient $\mathcal{A}_n$ ou que $H\cap\mathcal{A}_n=\{id\}$ ? 
\item On suppose que $H\cap \mathcal{A}_n=\{id\}$. Montrer que la restriction à $H$ du morphisme signature est injective.
Montrer que dans ce cas que tous les éléments de $H$ sont dans le centre de $\mathcal{S}_n$ et en déduire que $H=\{id\}$.
\item On suppose que $H$ contient $\mathcal{A}_n$. Montrer alors que $H=\mathcal{S}_n$ ou $H=\mathcal{A}_n$ suivant l'indice de $H$ dans $\mathcal{S}_n$.
\item Conclure~: si $n\geq 5$, les seuls sous-groupes distingués de $\mathcal{S}_n$ sont $\{id\}$, $\mathcal{A}_n$ et $\mathcal{S}_n$..
\end{enumerate}
\finenonce{007777}
\finexercice


\exercice{7778, mourougane, 2021/08/11}
\enonce{007778}{Groupe triangulaire supérieur}
\begin{enumerate}
\item Montrez que le groupe de Heisenberg $H$ des matrices $3*3$ triangulaires supérieures dont les coefficients diagonaux sont égaux à $1$, est résoluble.
\item Montrez que le groupe $B$ des matrices triangulaires supérieures $3*3$ inversibles (coefficients diagonaux non nuls) est résoluble. 
\end{enumerate}
\finenonce{007778}
\finexercice


\exercice{7779, mourougane, 2021/08/11}
\enonce{007779}{}
Le but de l'exercice est de déterminer les groupes dérivés successifs de $\mathcal{S}_4$.
On notera $V_4$ le sous-groupe des permutations de profil $(\cdot,\cdot)(\cdot,\cdot)$ (avec l'identité).
\begin{enumerate}
\item Montrer que $D(\mathcal{S}_4) \subset \mathcal{A}_4$.
\item Calculer les commutateurs $(1,2)(1,3)(1,2)^{-1}(1,3)^{-1}$ et $(1,2,3)(1,2,4)(1,2,3)^{-1}(1,2,4)^{-1}$.
\item Montrer que $D(\mathcal{S} _4)=A_4$. 
\item Montrer que $V_4\subset D(\mathcal{A}_4)$.
\item Vérifier que $V_4$ est distingué dans $\mathcal{A}_4$ et que le quotient $\mathcal{A}_4/V_4$ est un groupe abélien. En déduire que $D(\mathcal{A}_4)\subset V_4$.
\item En déduire $D^2(\mathcal{S}_4)$.
\item Calculer les autres groupes dérivés de $\mathcal{S}_4$.
\end{enumerate}
\finenonce{007779}
\finexercice


\exercice{7780, mourougane, 2021/08/11}
\enonce{007780}{Groupe dérivé de $GL(3,\mathbb{F}_2)$ et de $SL(3,\mathbb{F}_2)$}
\begin{enumerate}
\item Soit $$t=\begin{pmatrix}1&0&1\cr 0&1&0\cr 0&0&1\end{pmatrix}
\quad \text{ et } \quad s=\begin{pmatrix}0&-1&0\cr 1&0&0\cr 0&0&1\end{pmatrix}.$$
Montrer que le commutateur $tst^{-1}s^{-1}$ est une transvection.
\item Rappeler la démonstration du fait que deux transvections de $SL(3,\mathbb{F}_2)$ sont conjuguées dans $SL(3,\mathbb{F}_2)$. 
\item Déterminer $D(SL(3,\mathbb{F}_2))$ 
\item Déterminer $D(GL(3,\mathbb{F}_2))$.
\end{enumerate}
\finenonce{007780}
\finexercice


\exercice{7781, mourougane, 2021/08/11}
\enonce{007781}{Groupe dérivé de $GL(2,k)$}
On travaille dans $GL(2,k)$ pour un corps $k$ qui a au moins $4$
éléments. Soit $g\in GL(2,k)$. On notera $i_g$ l'automorphisme intérieur donné par $g$.
\begin{enumerate}
\item Démontrer qu'il existe un scalaire non nul $a\in k$ tel que $a^2\not=1$. Que se passe-t-il dans $\mathbb{F}_2$ et dans $\mathbb{F}_3$ ?
\item Soit $$t=\begin{pmatrix}1&1\cr 0&1\end{pmatrix}
\quad \text{ et } \quad s=\begin{pmatrix}a&0\cr 0& a^{-1}\end{pmatrix}.$$
Montrer que $T=tst^{-1}s^{-1}$ est une transvection.
\item Soit $\tau$ une transvection de $SL(2,k)$. Il existe $g\in GL(2,k)$ tel que $\tau = i_g(T):=gTg^{-1}$. Calculer $\tau$ à l'aide de $g$, $s$ et $t$ et montrer que $D(SL(2,k))$ contient toutes les transvections. 
\item Déterminer $D(SL(2,k))$.
\item Déterminer $D(GL(2,k))$.
\end{enumerate}
\finenonce{007781}
\finexercice


\exercice{7782, mourougane, 2021/08/11}
\enonce{007782}{Groupe dérivé de $GL(2,\mathbb{F}_3)$}
On travaille dans $GL(2,\mathbb{F}_3)$.

\begin{enumerate}
\item Soit $$t=\begin{pmatrix}1&1\cr 0&1\end{pmatrix} \quad \text{ et } \quad s=\begin{pmatrix}1&0\cr 0&-1\end{pmatrix}.$$
Calculer $tst^{-1}s^{-1}$.
\item Déterminer $D(GL(2,k))$.
Noter que ce calcul ne suffit pas pour déterminer $D(SL(2,k))$.
\end{enumerate}
\finenonce{007782}
\finexercice


\exercice{7783, mourougane, 2021/08/11}
\enonce{007783}{}
Montrer qu'un groupe d'ordre 63 n'est pas simple.
\finenonce{007783}
\finexercice


\exercice{7784, mourougane, 2021/08/11}
\enonce{007784}{Simplicité de $\mathcal{A}_5$}
\begin{enumerate}
 \item Faire la liste des classes de conjugaison de $\mathcal{S}_n$ dans $\mathcal{A}_n$ en les dénombrant.
\item Montrer que les $3$-cycles sont conjugués dans $\mathcal{A}_n$.
\item Montrer que les éléments d'ordre $2$ sont conjugués dans $\mathcal{A}_n$.
\item Montrer que tout sous-groupe distingué $H$ de $\mathcal{A}_n$ qui contient un élément d'ordre $5$ les contient tous. (On remarquera que le groupe engendré par un élément d'ordre $5$ est un Sylow.)
\item Montrer que tout sous-groupe distingué $H$ de $\mathcal{A}_n$ non réduit à $\{id\}$ contient au moins deux types d'éléments en plus de l'identité. Montrer alors que $H=\mathcal{A}_n$.
\end{enumerate}
\finenonce{007784}
\finexercice


\exercice{7786, mourougane, 2021/08/11}
\enonce{007786}{Groupe finis résolubles}
\begin{enumerate}
\item On rappelle que le centre d'un $p$-groupe n'est jamais réduit à $\{e\}$. Montrer par récurrence qu'un $p$-groupe est toujours résoluble.
\item Soient $p,q$ deux nombres premiers distincts. Montrer qu'un groupe de cardinal $pq$ est toujours résoluble. (Supposer $p>q$ et considérer un $p$-Sylow)
\item Soit $G$ un groupe d'ordre $12$. Montrez que $G$ est résoluble (En supposant les $3$-Sylow non distingués, comptez le nombre d'éléments d'ordre $3$).
\item Soient $p,q$ deux nombres premiers distincts. Montrez qu'un groupe de cardinal $p^2q$ est toujours résoluble.
\end{enumerate}
\finenonce{007786}
\finexercice


\exercice{7787, mourougane, 2021/08/11}
\enonce{007787}{Les sous-groupes distingués de $GL(E)$}
Le but de l'exercice est de déterminer les sous-groupes distingués de
 $SL(E)$.
 On supposera que le corps $k$ est de
caractéristique nulle ou que sa caractéristique est différente de $2$
et première avec la dimension $n$ de $E$. On supposera aussi $n\geq
 3$.

\begin{enumerate}
\item Donner l'exemple d'un sous-groupe non distingué de $SL(E)$.
\item Soit $\phi~:G\to H$ un morphisme surjectif de groupes.
 Montrer que l'image par $\phi$ d'un
 sous-groupe distingué de $G$ est un sous-groupe distingué de $H$.
\item Soit $H$ un sous-groupe de $SL(E)$. Déterminer les possibilités
 pour son image dans $PSL(E)$ pour la projection canonique.
\item Montrer que tout sous-groupe du centre de $SL(E)$ est distingué dans $SL(E)$.
\item On supposera désormais que $H$ n'est pas un sous-groupe du
 centre de $SL(E)$.
Soit $\tau$ une transvection. Montrer que $\tau^n$ est une
 transvection de $H$. 
\item Montrer que $H$ contient toutes les transvections de $E$.
\item Conclure.
\end{enumerate}
\finenonce{007787}
\finexercice


\exercice{7844, mourougane, 2021/08/11}
\enonce{007844}{}
Soit $G$ un groupe et $f~:G\to A$ un morphisme de groupes de $G$ dans un groupe abélien $A$. On suppose de plus que le noyau $N(f)$ de $f$ est résoluble.
\begin{enumerate}
    \item Montrez que le groupe dérivé $D(G)$ de $G$ est inclus dans le noyau de $F$. 
    \item Montrez que $G$ est résoluble.
\end{enumerate}
\finenonce{007844}
\finexercice


\exercice{7845, mourougane, 2021/08/11}
\enonce{007845}{Groupe dérivé}
Soit $G$ un groupe. On appelle groupe des commutateurs de $G$ et l'on note
$D(G)$ le sous-groupe de $G$ engendré par les éléments de la forme
$xyx^{-1}y^{-1}$. Montrer que $D(G)$ est distingué dans $G$ et que le quotient
$G/D(G)$ est abélien. Montrer que $D(G)$ est le plus petit sous-groupe
distingué de $G$ tel que le quotient de $G$ par ce sous-groupe soit abélien.
\finenonce{007845}
\finexercice


\exercice{7857, mourougane, 2021/08/11}
\enonce{007857}{Exemples de sous-groupes caractéristiques}
\begin{enumerate}
\item Montrer qu'un $p$-Sylow distingué est caractéristique.
\item Soit $H$ un sous-groupe distingué d'un groupe fini $G$ tel que son ordre est premier avec son indice.
 Montrer alors que $H$ est le seul sous-groupe d'ordre $|H|$ et donc que $H$ est caractéristique.
\end{enumerate}
\finenonce{007857}
\finexercice


\exercice{7858, mourougane, 2021/08/11}
\enonce{007858}{Produit semi-direct interne}
Soit $N$ un sous-groupe distingué d'un groupe $G$ ($N\lhd G$)
et $H$ un sous groupe de $G$ tel que $H\cap N=\{e_G\}$.
\begin{enumerate}
\item Montrer que $NH$ est un sous-groupe de $G$.
\item On suppose désormais que $|G|=|N||H|$.
Montrer que $\phi : N\times H\to G$, $(n,h)\mapsto nh$ est une bijection.
\item Montrer que si on munit $N\times H$ de la loi 
$$(n,h)\star(n',h')=(n(hn'h^{-1}), hh'),$$
alors $N\times H$ est un groupe et $\phi$ un isomorphisme de groupes.
\end{enumerate}
\finenonce{007858}
\finexercice


\exercice{7859, mourougane, 2021/08/11}
\enonce{007859}{Groupes d'automorphismes}
\begin{enumerate}
 \item Montrer que les automorphismes du groupe $(\Z/n\Z,+)$ sont obtenus par multiplication par un inversible de 
 $(\Z/n\Z,\times)$.
 \item Décrire un isomorphisme de $\Z/10\Z$ sur $(\Z/11\Z)^\star$.
 \item Montrer que si $G$ et $H$ sont deux groupes d'ordre premiers entre eux, alors
 $$Aut(G\times H)=Aut(G)\times Aut (H).$$
 \item En déduire le groupe des automorphismes de $\Z/133\Z$.
 \item Soit $p$ un nombre premier et $n$ un entier naturel non nul.
 Montrer que $$Aut((\Z/p\Z)^n)=GL(n,\mathbb{F}_p).$$
 \item Montrer que $Aut(\Z/2\Z\times\Z/2\Z)\to Bij((1,0),(1,1),(0,1))$
 est un isomorphisme.
\end{enumerate}
\finenonce{007859}
\finexercice


\exercice{7860, mourougane, 2021/08/11}
\enonce{007860}{Exemple de produits semi-directs}
\begin{enumerate}
 \item Montrer que, après avoir fixé un générateur de $(\Z/11\Z)^\star$,
 la donnée d'un morphisme de groupes de $\Z/5\Z$ dans $Aut(\Z/11\Z)$
 revient à la donnée d'un morphisme de groupes de $\Z/5\Z$ dans $\Z/10\Z$.
 \item En déduire une structure de produit semi-direct sur $\Z/11\Z\rtimes\Z/5\Z$.
 \item Montrer que toutes les structures de produit semi-direct donnent des groupes isomorphes.
 \textit{On pourra montrer que si $\phi, \psi\in Hom(\Z/5\Z,Aut(\Z/11\Z)$ 
 alors il existe $\gamma\in Aut(\Z/11\Z)$ tel que $\psi (h) = \gamma\circ\phi (h)\circ \gamma^{-1}$.}
 \item Montrer que tous les morphismes de $\Z/p\Z\to Aut (\Z/q\Z)$ sont de la forme 
 $t\mapsto \{x\mapsto k^t x\}$ où $k$ est un élément de $(\Z/q\Z)^\star$ d'ordre $p$.
\end{enumerate}
\finenonce{007860}
\finexercice


\exercice{7871, mourougane, 2021/08/11}
\enonce{007871}{}
Montrer qu'un groupe d'ordre 63 n'est pas simple.
\finenonce{007871}
\finexercice


\exercice{7876, mourougane, 2021/08/11}
\enonce{007876}{Généralités}
\begin{enumerate}
\item Montrer que le groupe des matrices triangulaires supérieures de diagonale identité est résoluble.
\item Montrer que si $H$ est un sous-groupe distingué d'un groupe $G$, alors $G$ est résoluble si et seulement si $H$ et $G/H$ le sont.
(On pourra commencer par le cas où $G/H$ est abélien).
\item Montrer qu'un $p$-groupe est résoluble.
\end{enumerate}
\finenonce{007876}
\finexercice


\exercice{7877, mourougane, 2021/08/11}
\enonce{007877}{}
\begin{enumerate}
\item L'ensemble des permutations de profil $(\cdot, \cdot)(\cdot, \cdot)$ avec l'identité est-il un sous-groupe distingué de $\mathcal{A}_6$.
(Justifier)
\item
Donner l'exemple d'un groupe résoluble. 
\item
Donner si possible l'exemple d'un groupe simple résoluble.
\item
Donner si possible l'exemple d'un groupe simple résoluble non abélien.
\end{enumerate}
\finenonce{007877}
\finexercice


\exercice{7878, mourougane, 2021/08/11}
\enonce{007878}{Groupe triangulaire supérieur}
\begin{enumerate}
\item Montrez que le groupe de Heisenberg $H$ des matrices $3*3$ triangulaires supérieures dont les coefficients diagonaux sont égaux à $1$, est résoluble.
\item Montrez que le groupe $B$ des matrices triangulaires supérieures $3*3$ inversibles (coefficients diagonaux non nuls) est résoluble. 
\end{enumerate}
\finenonce{007878}
\finexercice


\exercice{7879, mourougane, 2021/08/11}
\enonce{007879}{Groupes dérivés de $\mathcal{S}_4$}
Le but de l'exercice est de déterminer les groupes dérivés successifs de $\mathcal{S}_4$.
On notera $V_4$ le sous-groupe des permutations de profil $(\cdot,\cdot)(\cdot,\cdot)$ (avec l'identité).
\begin{enumerate}
\item Montrer que $D(\mathcal{S}_4) \subset \mathcal{A}_4$.
\item Calculer les commutateurs $(1,2)(1,3)(1,2)^{-1}(1,3)^{-1}$ et $(1,2,3)(1,2,4)(1,2,3)^{-1}(1,2,4)^{-1}$.
\item Montrer que $D(\S _4)=A_4$. 
\item Montrer que $V_4\subset D(\mathcal{A}_4)$.
\item Vérifier que $V_4$ est distingué dans $\mathcal{A}_4$ et que le quotient $\mathcal{A}_4/V_4$ est un groupe abélien. En déduire que $D(\mathcal{A}_4)\subset V_4$.
\item En déduire $D^2(\mathcal{S}_4)$.
\item Calculer les autres groupes dérivés de $\mathcal{S}_4$.
\end{enumerate}
\finenonce{007879}
\finexercice


\exercice{7880, mourougane, 2021/08/11}
\enonce{007880}{Résolubilité}
\begin{enumerate}
\item Montrer que $\mathcal{S}_3$ est résoluble.
\item Montrer que le groupe $D_4$ des permutations de profil $(.,.)(.,.)$ est un groupe abélien d'ordre $4$ distingué dans $A_4$.
En déduire que $\mathcal{A}_4$ et donc $\mathcal{S}_4$ sont résolubles.
\item On suppose désormais $n\geq 3$. Soit $c$ un $3$ cycle. En considérant $c^2$, montrer que $c$ est un commutateur dans $\mathcal{S}_n$.
En déduire le sous groupe dérivé $D(\mathcal{S}_n)$.
\item Montrer que pour $n\geq 5$, $\mathcal{A}_n$ et $\mathcal{S}_n$ ne sont pas résolubles.
\end{enumerate}
\finenonce{007880}
\finexercice


\exercice{7881, mourougane, 2021/08/11}
\enonce{007881}{Groupe d'ordre $pqr$}
Soit $p>q>r$ trois nombres premiers.
Le but de l'exercice est de montrer qu'un groupe d'ordre $pqr$ est résoluble.
\begin{enumerate}
 \item Montrer qu'un groupe d'ordre $pq$ est résoluble.
\item Soit $G$ un groupe d'ordre $pqr$. Supposons qu'il n'admette pas de sous-groupe distingué.
On note $N_p$ (resp. $N_q$, $N_r$ le nombre de sous-groupes de Sylow d'ordre $p$ (resp. $q$, $r$).
Montrer que $m_p=qr$, $m_q\geq p$ et $m_r\geq q$.
\item Conclure. 
\end{enumerate}
\finenonce{007881}
\finexercice


\exercice{7882, mourougane, 2021/08/11}
\enonce{007882}{Groupes linéaires}
Soit $E$ un $k$-espace vectoriel. Soit $f$ une forme linéaire sur $E$ et $a$ un élément non nul de $H=ker(f)$.
On appelle transvection associée à $f$ et $a$ l'application $u : E\to E$, $x\mapsto x+f(x)a$.
On rappelle que les transvections de $E$ engendrent $SL(E)$.
\begin{enumerate}
\item Soit $u$ une transvection. En considérant une base $(e_i)$ de $E$ avec 
$e_{n-1}=a$, $(e_j)_{1\leq j\leq n-1}$ base de $H$, et $e_n$ tel que $f(e_n)=1$, écrire la matrice de $u$.
 \item Montrer que si $u$ est une transvection, $Ker(u-Id)=H$, $\det u=1$, $u$ n'est pas diagonalisable.
 \item Montrer que si $\dim E\geq 3$, les transvections de $E$ sont conjuguées dans $SL(E)$.
 \item On suppose $k$ de caractéristique différente de $2$ et $\dim E\geq 3$.
 Montrer que $$D(GL(E)=D(SL(E))=SL(E)$$ et donc que ni $GL(E)$, ni $SL(E)$ ne sont résolubles.
\end{enumerate}
\finenonce{007882}
\finexercice


\exercice{7883, mourougane, 2021/08/11}
\enonce{007883}{Groupe dérivé de $GL(3,\mathbb{F}_2)$ et de $SL(3,\mathbb{F}_2)$}
\begin{enumerate}
\item Soit $$t=\begin{pmatrix}1&0&1\cr 0&1&0\cr 0&0&1\end{pmatrix}
\quad \text{ et } \quad s=\begin{pmatrix}0&-1&0\cr 1&0&0\cr 0&0&1\end{pmatrix}.$$
Montrer que le commutateur $tst^{-1}s^{-1}$ est une transvection.
\item Rappeler la démonstration du fait que deux transvections de $SL(3,\mathbb{F}_2)$ sont conjuguées dans $SL(3,\mathbb{F}_2)$. 
\item Déterminer $D(SL(3,\mathbb{F}_2))$.
\item Déterminer $D(GL(3,\mathbb{F}_2))$.
\end{enumerate}
\finenonce{007883}
\finexercice


\exercice{7884, mourougane, 2021/08/11}
\enonce{007884}{Groupe dérivé de $GL(2,k)$}
On travaille dans $GL(2,k)$ pour un corps $k$ qui a au moins $4$
éléments. Soit $g\in GL(2,k)$. On notera $i_g$ l'automorphisme intérieur donné par $g$.
\begin{enumerate}
\item Démontrer qu'il existe un scalaire non nul $a\in k$ tel que $a^2\not=1$. Que se passe-t-il dans $\mathbb{F}_2$ et dans $\mathbb{F}_3$ ?
\item Soit $$t=\begin{pmatrix}1&1\cr 0&1\end{pmatrix}
\quad \text{ et } \quad s=\begin{pmatrix}a&0\cr 0& a^{-1}\end{pmatrix}.$$
Montrer que $T=tst^{-1}s^{-1}$ est une transvection.
\item Soit $\tau$ une transvection de $SL(2,k)$. Il existe $g\in GL(2,k)$ tel que $\tau = i_g(T):=gTg^{-1}$. Calculer $\tau$ à l'aide de $g$, $s$ et $t$ et montrer que $D(SL(2,k))$ contient toutes les transvections. 
\item Déterminer $D(SL(2,k))$.
\item Déterminer $D(GL(2,k))$.
\end{enumerate}
\finenonce{007884}
\finexercice


\exercice{7885, mourougane, 2021/08/11}
\enonce{007885}{Groupe dérivé de $GL(2,\mathbb{F}_3)$}
On travaille dans $GL(2,\mathbb{F}_3)$.

\begin{enumerate}
\item Soit $$t=\begin{pmatrix}1&1\cr 0&1\end{pmatrix} \quad \text{ et } \quad s=\begin{pmatrix}1&0\cr 0&-1\end{pmatrix}.$$
Calculer $tst^{-1}s^{-1}$.
\item Déterminer $D(GL(2,k))$.
Noter que ce calcul ne suffit pas pour déterminer $D(SL(2,k))$.
\end{enumerate}
\finenonce{007885}
\finexercice


\exercice{7886, mourougane, 2021/08/11}
\enonce{007886}{Des petites questions}
\begin{enumerate}
 \item L'entier $374$ divise-t-il l'ordre du groupe $\Z/374\Z$ ? Le groupe $\Z/374\Z$ a-t-il un élément d'ordre $374$ ? a-t-il un élément d'ordre $187$ ?
 \item Déterminer tous les diviseurs de $374$ ? Le groupe $\Z/374\Z$ a-t-il un sous-groupe d'ordre chacun des diviseurs de $374$ ?
\end{enumerate}
\finenonce{007886}
\finexercice

\section{304.00 Action de groupe}


\exercice{2166, debes, 2008/02/12}
\enonce{002166}{}
Soit $\sigma \in S_5$ d\'efini par 
$$\sigma = 
\begin{array}{ccccc}
1&2&3&4&5\\
5&4&1&2&3\\
\end{array}$$

(a) Ecrire la d\'ecomposition de $\sigma $ en produit de cycles de
supports disjoints. Quelle est la signature de $\sigma $?

\smallskip
(b) Donner la liste des \'el\'ements de $<\sigma >$. D\'eterminer $<\sigma
> \cap A_5.$
\finenonce{002166} 


\finexercice


\exercice{2167, debes, 2008/02/12}
\enonce{002167}{}
\label{ex:deb67}
(a) Montrer que le produit de deux transpositions distinctes
est un $3$-cycle ou un produit de deux $3$-cycles. En d\'eduire que $A_n$
est engendr\'e par les $3$-cycles.
\smallskip

(b) Montrer que $A_n = <(123), (124), \dots , (12n)>$.
\finenonce{002167} 

\finexercice


\exercice{2168, debes, 2008/02/12}
\enonce{002168}{}
\label{ex:deb68}
On appelle cycle une permutation $\sigma $ v\'erifiant la propri\'et\'e
suivante:  il existe une partition de $\{ 1, \dots , n \} $ en deux sous-ensembles
$I$ et $J$ tels que la restriction de $\sigma $ \`a $I$ est l'identit\'e
de $I$
et il existe $a\in J$ tel que $J= \{ a, \sigma (a) , \dots , \sigma ^{r-1}
(a)\} $ o\`u $r$ est le cardinal de $J$. Le sous-ensemble $J$ est appel\'e le
support du cycle $\sigma $.

Un tel cycle sera not\'e $(a, \sigma (a), \dots , \sigma ^{r-1} (a) )$
\smallskip

(a) Soit
$\sigma
\in S_n$ une permutation. On consid\`ere le sous-groupe
$C$ engendr\'e par
$\sigma$ dans
$S_n$. Montrer que la restriction de $\sigma$ \`a chacune des orbites de
$\{1, \dots ,n\} $ sous
l'action de $C $ est un cycle, que ces diff\'erents cycles commutent entre
eux, et que $\sigma $
est le produit de ces cycles.
\smallskip

(b) D\'ecomposer en cycles les permutations suivantes de $\{ 1, \dots , 7\} $ :

$\hskip 20mm 1\hskip 2mm 2\hskip 2mm 3\hskip 2mm 4\hskip 2mm 5\hskip 2mm 6\hskip 2mm 7 \hskip
20mm  1\hskip 2mm 2\hskip 2mm 3\hskip 2mm 4\hskip 2mm 5\hskip 2mm 6\hskip 2mm 7\hskip 20mm   
1\hskip 2mm 2\hskip 2mm 3\hskip 2mm 4\hskip 2mm 5\hskip 2mm 6\hskip 2mm 7$

$\hskip 20mm 3\hskip 2mm 6\hskip 2mm 7\hskip 2mm 2\hskip 2mm 1\hskip 2mm 4\hskip 2mm 5 \hskip
20mm  7\hskip 2mm 4\hskip 2mm 2\hskip 2mm 3\hskip 2mm 5\hskip 2mm 6\hskip 2mm 1\hskip 20mm    
1\hskip 2mm 3\hskip 2mm 7\hskip 2mm 2\hskip 2mm 4\hskip 2mm 5\hskip 2mm 6$

\smallskip
(c) Montrer que si $\sigma $ est un cycle, $\sigma =(a, \sigma (a), \dots ,
\sigma^{r-1}(a)) $, la
conjugu\'ee $\tau \sigma \tau ^{-1} $ est un cycle et que
$\tau \sigma \tau ^{-1} =(\tau (a), \tau (\sigma  (a)), \dots , \tau
(\sigma ^{r-1} (a)))$.
\smallskip

(d) D\'eterminer toutes les classes de conjugaison des permutations dans $S_5$
(on consid\'erera leur d\'ecomposition en cycles). D\'eterminer tous les
sous-groupes distingu\'es de $S_5$.
\finenonce{002168} 


\finexercice


\exercice{2169, debes, 2008/02/12}
\enonce{002169}{}
\label{ex:deb69}
Montrer que  les permutations circulaires engendrent $S_n$ si
$n$ est pair, et $A_n$ si $n$ est impair.
\finenonce{002169} 

\finexercice


\exercice{2170, debes, 2008/02/12}
\enonce{002170}{}
Soit $I$ un sous-ensemble de $\{ 1, \dots , n\} $ et
$\sigma $ un cycle de support $I$. Soit $\tau $ une autre permutation.
Montrer que $\tau $ commute avec $\sigma $ si et seulement si $\tau $
laisse invariant $I$ et la restriction de $\tau $ \`a $I$ est \'egale \`a une
puissance de la restriction de $\sigma $ \`a $I$.
\finenonce{002170} 

\finexercice


\exercice{2171, debes, 2008/02/12}
\enonce{002171}{}
Soit $H$ un sous-groupe
distingu\'e de $S_n$ contenant une
transposition. Montrer que $H=S_n$.
\finenonce{002171} 

\finexercice


\exercice{2172, debes, 2008/02/12}
\enonce{002172}{}
Dans le groupe sym\'etrique $S_4$ on consid\`ere les
sous-ensembles suivants :
$$H = \{ \sigma \in S_4 \hskip 2pt |\hskip 2pt  \sigma ( \{ 1,2\} ) = \{ 1,2\} \} $$
$$K= \{ \sigma \in S_4 \hskip 2pt |\hskip 2pt  \forall a,b \quad a\equiv b\ [\hbox{\rm mod}\
2] \Rightarrow
\sigma (a) \equiv \sigma (b)\ [\hbox{\rm mod}\ 2] \} $$
Montrer que $H$ et $K$ sont des sous-groupes de $S_4$. Les d\'ecrire.
\finenonce{002172} 

\finexercice


\exercice{2173, debes, 2008/02/12}
\enonce{002173}{}
Montrer que l'ordre d'une permutation impaire est un nombre pair.
\finenonce{002173} 



\finexercice


\exercice{2174, debes, 2008/02/12}
\enonce{002174}{}
Montrer que toute permutation d'ordre $10$ dans $S_8$ est
impaire.
\finenonce{002174} 

\finexercice


\exercice{2175, debes, 2008/02/12}
\enonce{002175}{}
\label{ex:deb75}
(a) Montrer que tout $3$-cycle est un carr\'e. En d\'eduire
que le groupe altern\'e $A_n$ est engendr\'e par les carr\'es de
permutations.
\smallskip

(b) Montrer que $A_n$ est le seul sous-groupe de $S_n$ d'indice $2$.
\finenonce{002175} 

\finexercice


\exercice{2176, debes, 2008/02/12}
\enonce{002176}{}
Trouver toutes les classes de conjugaison de $S_4$. Donner la
liste des sous-groupes distingu\'es de $S_4$.
\finenonce{002176} 

\finexercice


\exercice{2177, debes, 2008/02/12}
\enonce{002177}{} 
Etant donn\'es un groupe $G$ et un sous-groupe $H$, on d\'efinit le normalisateur $\textrm{Nor}_G(H)$ de $H$ dans $G$ comme l'ensemble des \'el\'ements $g\in G$ tels que $gHg^{-1} = H$. 
\smallskip

(a) Montrer que $\textrm{Nor}_G(H)$ est le plus grand sous-groupe de $G$ contenant $H$ comme sous-groupe distingu\'e. 
\smallskip

(b) Montrer que le nombre de sous-groupes distincts conjugu\'es de $H$ dans $G$ est \'egal \`a l'indice $[G:\textrm{Nor}_G(H)]$ et qu'en particulier c'est un diviseur de l'ordre de $G$.
\finenonce{002177} 


\finexercice


\exercice{2178, debes, 2008/02/12}
\enonce{002178}{}
Montrer que pour $m\geq 3$, un groupe simple d'ordre $\geq m!$ ne peut avoir de sous-groupe d'indice $m$.
\finenonce{002178} 


\finexercice


\exercice{2179, debes, 2008/02/12}
\enonce{002179}{}
Soit $G$ un groupe et $H$ un sous-groupe d'indice fini $n$. Montrer que
l'intersection $H^\prime$ des conjugu\'es de $H$ par les \'el\'ements de $G$ est un
sous-groupe distingu\'e de $G$ et d'indice fini dans $G$. Montrer que c'est le plus
grand sous-groupe distingu\'e de $G$ contenu dans $H$.
\finenonce{002179} 


\finexercice


\exercice{2180, debes, 2008/02/12}
\enonce{002180}{}
a) Montrer qu'un groupe $G$ v\'erifiant
$$\forall a,b \in G \quad a^2b^2=(ab)^2$$
est commutatif.
\smallskip

(b) Le but de cette question est de donner un exemple de groupe $G$ 
v\'erifiant la propri\'et\'e 
$$\forall a,b\in G \quad a^3b^3=(ab)^3$$
et qui n'est pas commutatif.

\hskip 5mm (i) montrer qu'il existe un automorphisme $\sigma$ de
$\mathbb{F}_3^2$ d'ordre $3$.

\hskip 5mm (ii) montrer que le groupe $G$ d\'efini comme le produit semi-direct
de $\mathbb{F}_3^2$ par $\Z_3$, $\Z_3$ agissant sur $\mathbb{F}_3^2$ via
$\sigma$ r\'epond \`a la question. 
\finenonce{002180}


\finexercice


\exercice{2181, debes, 2008/02/12}
\enonce{002181}{}
Soient $G$ un groupe et $H$ un sous-groupe d'indice fini dans $G$.  
On d\'efinit sur $G$ la relation $xRy$ si et seulement si $x\in HyH$.
\smallskip

(a) Montrer que $R$ est une relation d'\'equivalence et que toute classe
d'\'equivalence pour la relation $R$ est une union finie disjointe de classes 
\`a gauche modulo $H$. 
\smallskip

Soit $HxH=\bigcup _{1\leq i\leq d(x)} x_i H$  la partition de la classe $HxH$ en
classes \`a gauche distinctes.
\smallskip

(b) Soit $h\in H$ et $i$ un entier compris entre $1$ et $d(x)$; posons  $h\hskip
2pt \ast \hskip 2pt x_iH= hx_iH$. Montrer que cette formule d\'efinit une action
transitive de $H$ sur l'ensemble  des classes $x_1H, \dots, x_{d(x)} H$ et que le
fixateur de $x_i H$ dans cette action est  $H\cap x_i Hx_i ^{-1}$. En d\'eduire que 

$$d(x)=[H:H\cap x Hx ^{-1}]$$ 

et qu'en particulier $d(x)$ divise l'ordre de $G$.
\smallskip

(c) Montrer que $H$ est distingu\'e dans  $G$  si et seulement si 
$d(x)=1$  pour tout $x\in G$.
\smallskip

(d) On suppose que $G$ est fini et que $[G:H]=p$, o\`u $p$ est le plus petit 
nombre premier divisant l'ordre de $G$. Le but de cette question est de montrer 
que $H$ est distingu\'e dans $G$.

\hskip 5mm (i) Montrer que pour tout $x\in G$ , $d(x)\leq p$. En d\'eduire que
$d(x)=1$ ou $d(x)=p$. 

\hskip 5mm (ii) Montrer que si $H$ n'est pas distingu\'e dans $G$, il existe une
unique classe d'\'equivalence pour la relation $R$ et que $G=H$, ce qui contredit
l'hypoth\`ese  $[G:H ]=p$.
\finenonce{002181} 

\finexercice


\exercice{2182, debes, 2008/02/12}
\enonce{002182}{}
Soit $G$ un groupe fini agissant sur un ensemble fini $X$.
\smallskip

(a) On suppose que toute orbite contient au moins deux \'el\'ements, que
$|G|=15$ et que $\hbox{\rm card}(X)=17$. D\'eterminer le nombre d'orbites et le
cardinal de chacune.
\smallskip

(b) On suppose que $|G|=33$ et $\hbox{\rm card}(X)=19$. Montrer qu'il existe au
moins une orbite r\'eduite \`a un \'el\'ement.
\finenonce{002182} 

\finexercice


\exercice{2183, debes, 2008/02/12}
\enonce{002183}{}
\label{ex:deb83}
(a) Soit $G$ un groupe et $H$ un sous-groupe. Montrer que la
formule
$$g. g'H=gg'H$$
d\'efinit une action de $G$ sur l'ensemble quotient $G/H$. D\'eterminer le
fixateur d'une classe $gH$.
\medskip

(b) Soit $G$ un groupe et $X$ et $Y$ deux ensembles sur lesquels $G$ agit
(on parlera de
$G$-ensembles). Soit
$f$ une application de $X$ dans $Y$. On dira que $f$ est compatible \`a
l'action de $G$ (ou que $f$
est un morphisme de $G$-ensembles) si pour tout \'el\'ement
$x$ de $X$ et tout $g$ dans $G$, $f(g.x)=g.f(x)$. Montrer que si $f$ est
bijective et compatible \`a
l'action de $G$ il en est de m\^eme de $f^{-1}$. On dira dans ce cas que
$f$ est un isomorphisme de
$G$-ensembles.
\medskip

(c) Soit $G$ un groupe agissant transitivement sur un ensemble $X$ (i.e.
pour tout couple d'\'el\'ements
$x$ et $y$ de $X$ il existe au moins un \'el\'ement $g$ du groupe tel que
$g.x=y$). Montrer qu'il
existe un sous-groupe $H$ de $G$ tel que $X$ soit isomorphe en tant que
$G$-ensemble \`a $G/H$ (on
prendra pour $H$ le fixateur d'un point quelconque de $X$).
\medskip

(d) i) Soit $H$ et $K$ deux sous-groupes de $G$. Montrer qu'il existe une
application $f$ de $G/H$
vers
$G/K$ compatible avec l'action de
$G$ si et seulement si $H$ est contenu dans un conjugu\'e de $K$. Montrer
que dans ce cas $f$ est
surjective.  Montrer que
$G/H$ et
$G/K$ sont isomorphes en tant que $G$-ensembles si et seulement si $H$ et
$K$ sont conjugu\'es dans
$G$.
\smallskip

ii) Soit $X$ et $Y$ deux $G$-ensembles transitifs. Montrer qu'il existe une
application de $X$ vers
$Y$ compatible avec l'action de $G$ si et seulement si il existe deux
\'el\'ements $x$ et $y$ de $X$ et
$Y$ tels que le fixateur de $x$ soit contenu dans un conjugu\'e du
fixateur de $y$. Montrer que $X$ et
$Y$ sont isomorphes si et seulement si les fixateurs de $x$ et de $y$ sont
conjugu\'es dans $G$.
\finenonce{002183} 

\finexercice


\exercice{2184, debes, 2008/02/12}
\enonce{002184}{}
\label{ex:deb84}
Soit $G$ un groupe fini et $X$ un $G$-ensemble transitif. On
dira que $X$ est
{\it imprimitif } si $X$ admet une partition $X=\bigcup _{1\leq i\leq r}
X_i$ telle que tout \'el\'ement
$g$ de $G$ respecte cette partition, i.e. envoie un sous-ensemble $X_i$
sur un sous-ensemble $X_k$
(\'eventuellement $k=i$) et telle que $2\leq r$ et les parties $X_i$ ne
sont pas r\'eduites \`a un
\'el\'ement. Dans le cas contraire on dit que $X$ est {\it primitif}.
\smallskip

(a) Montrer que dans la d\'ecomposition pr\'ec\'edente, si elle existe,
tous les sous-ensembles $X_i$ ont
m\^eme nombre $m$ d'\'el\'ements.
\smallskip

(b) Soit $H$ un sous-groupe de $G$. Montrer que $G/H$ est imprimitif si et
seulement s'il existe un
sous-groupe propre $K$ de $G$ diff\'erent de $H$ tel que $H\subset K
\subset G$ (on regardera la
partition de $G/H$ en classes modulo $K$).
\smallskip

(c) D\'eduire de ce qui pr\'ec\`ede que $X$ est primitif si et seulement si
le fixateur d'un \'el\'ement $x$
de $X$ est maximal parmi les sous-groupes propres de $G$.
\smallskip

(d) On suppose ici que $X$ est primitif et que $H$ est un sous-groupe
distingu\'e de $G$ dont l'action
n'est pas triviale sur $X$. Montrer qu'alors $H$ agit transitivement sur $X$.
\finenonce{002184} 


\finexercice


\exercice{2185, debes, 2008/02/12}
\enonce{002185}{}
\label{ex:deb85}
Montrer qu'un sous-groupe primitif de $S_n$ qui contient une
transposition est $S_n$
tout entier.
\finenonce{002185} 


\finexercice


\exercice{2186, debes, 2008/02/12}
\enonce{002186}{}
\label{ex:deb86}
Soit $G$ un groupe fini et $X$ un $G$-ensemble. Si $k$ est
un entier ($1\leq k$),
on dit que $X$ est $k$-transitif, si pour tout couple de $k$-uplets 
$(x_1, \dots ,x_k)$ et
$(y_1, \dots , y_k)$ d'\'el\'ements de $X$ distincts deux \`a deux,
il existe au moins un \'el\'ement $g$
de
$G$ tel que pour tout $i$, $1\leq i \leq k$, $g.x_i =y_i$. Un $G$-ensemble
$1$-transitif est donc
simplement un
$G$-ensemble transitif.
\smallskip

(a) Montrer que si $X$ est $k$-transitif, il est aussi $l$-transitif pour
tout $l$, $1\leq l \leq k$.
\smallskip

(b) Montrer que $X$ est $2$-transitif si et seulement si le fixateur d'un \'el\'ement $x$
de $X$ agit transitivement sur $X\setminus\{x\}$.

(c) Montrer que si $X$ est imprimitif, il n'est pas $2$-transitif.
\smallskip

(d) Montrer qu'un groupe cyclique $C$ d'ordre premier consid\'er\'e comme
$C$-ensemble par l'action de
translation de $C$ sur lui-m\^eme, est primitif mais n'est pas $2$-transitif.
\smallskip

(e) Montrer que l'ensemble $\{ 1, \dots ,n \} $ muni de l'action du groupe
$S_n$ est $k$-transitif
pour tout
$k$,
$1\leq k\leq n$. En d\'eduire que
l'ensemble $\{ 1, \dots ,n \} $ muni de l'action du groupe $S_n$ est
primitif.
\smallskip

(f) Montrer que le fixateur de $1$ dans $S_n$ est isomorphe \`a $S_{n-1}$.
Dans la suite on
identifie $S_{n-1}$ \`a ce fixateur. D\'eduire de l'exercice
19 que $S_{n-1}$ est un
sous-groupe propre maximal de $S_n$.
\finenonce{002186} 


\finexercice


\exercice{2187, debes, 2008/02/12}
\enonce{002187}{}
D\'ecrire le groupe $D_n$ des isom\'etries du plan affine
euclidien qui laissent invariant un
polygone r\'egulier \`a $n$ c\^ot\'es. Montrer que $D_n$ est engendr\'e
par deux \'el\'ements $\sigma $ et $\tau
$ qui v\'erifient les relations: $\sigma ^ n =1$, $\tau ^2=1$ et $\tau
\sigma \tau ^{-1} =\sigma
^{-1}$.  Quel est l'ordre de $D_n$? D\'eterminer le centre de $D_n$.
Montrer que $D_3 \simeq S_3$.
\finenonce{002187} 


\finexercice


\exercice{2188, debes, 2008/02/12}
\enonce{002188}{}
Montrer que le groupe des isom\'etries de l'espace affine
euclidien de dimension $3$
qui laissent invariant un t\'etra\`edre r\'egulier de sommets $a_1,
a_2,a_3,a_4$ est isomorphe \`a $S_4$
et que le sous-groupe des isom\'etries directes qui laissent invariant le
t\'etra\`edre est isomorphe \`a
$A_4$.
\finenonce{002188} 

\finexercice


\exercice{2189, debes, 2008/02/12}
\enonce{002189}{}
D\'eterminer le groupe des isom\'etries de l'espace affine
euclidien de dimension $3$
qui laissent invariant un cube.
\finenonce{002189}

\finexercice


\exercice{6379, potyag, 2011/10/16}
\enonce{006379}{}
Soit $G$ un groupe, $H$ un sous-groupe d'indice $n$ dans $G$.
\begin{enumerate}
\item A l'aide des classes à gauche modulo $H$ dans $G$, construire un
homomorphisme $\phi: G \mapsto{\mathcal S}_n$.

\item Montrer que si $N \subset H$ et $N$ est normal dans $G$, on a
 $N < {\rm Ker}\ \phi < H$.

\item En déduire que tout groupe fini $G$ est isomorphe à un
sous-groupe de ${\mathcal S}_n$.
\end{enumerate}
\finenonce{006379}



\finexercice


\exercice{6380, potyag, 2011/10/16}
\enonce{006380}{}
Soit $K$ un corps fini à $q$ éléments, $GL(n,K)$ l'ensemble
des matrices $(n,n)$ inversibles à coefficients dans $K$.
Montrer par récurrence sur $n$ que $\vert GL(n,K)\vert= (q^n-1)
(q^n-q)...(q^n-q^{n-1})$ en considérant l'action de $Aut(K^n)$
sur l'espace vectoriel $K^n$   (de base $\{v_1,...,v_n\}$),
l'orbite et le stabilisateur d'un vecteur de base ($v_1$ par
exemple).
\finenonce{006380}



\finexercice


\exercice{6381, potyag, 2011/10/16}
\enonce{006381}{}
Soit $\varphi$ une action d'un  groupe $G$ opèrant dans $X$
(notée $G\stackrel{\varphi} {\curvearrowright} X$).

\begin{enumerate}


  \item  Montrer que $G_{g(x)}=gG_xg^{-1}$, où $g\in G$ et
  $G_x$ désigne le stabilisateur du point $x$.


  \item  Si l'action $\varphi$ est transitive et fidèle et
  $G$ est abélien alors montrer que $\varphi$ est simplement transitive.
\end{enumerate}
\finenonce{006381}



\finexercice


\exercice{6382, potyag, 2011/10/16}
\enonce{006382}{}
Soit $G=<\gamma_1,\gamma_2>$ opère sur le plan complexe $\Cc$ où
$\gamma_1 : z\mapsto z+1$ et $\gamma_2 : z\mapsto z+i$.
\begin{enumerate}
  \item Montrer que $G\cong\Z^2$ et $G$ agit isométriquement sur
  $\Cc$.

  \item Trouver un ensemble fondamental $F$ pour cette
  action et l'ensemble d'orbites\quad $\Cc/\!\!\thicksim\ = \Cc/G =
  \overline{F}/\!\!\thicksim$\quad en identifiant les points équivalents
  sur le bord de $\overline{F}$.

\end{enumerate}
\finenonce{006382}



\finexercice


\exercice{6383, potyag, 2011/10/16}
\enonce{006383}{}
\begin{enumerate}
  \item  Montrer que la quantité suivante (appelée forme de
  Killing) est un produit scalaire sur le groupe matriciel
  $M_n(\R)$.
  $$<X,Y> = {\rm tr}\ (X^{T}Y).\quad X, Y\in M_n(\R)$$

  \item  Montrer que la forme de Killing reste invariante par
  rapport à l'action de $O(n)$ par conjugaison :

  $$<gXg^{-1}, gYg^{-1}>\ =\ <X, Y>,\ \ X,Y\in M_n(\R),\ g\in O(n).$$

\end{enumerate}
\finenonce{006383}



\finexercice


\exercice{6384, potyag, 2011/10/16}
\enonce{006384}{}
 Soient $G$ un groupe et $S$ un système de
générateurs de $G$ contenant avec chaque élément $s$ son
 inverse $s^{-1}$. Rappelons la construction du graphe de
Cayley $C(G,S)$. L'ensemble $V$ des sommets de $C(G,S)$
est en bijection  avec l'ensemble des éléments de $G$.
Deux sommets $g_1$ et $g_2$ sont joints par une arête si
$g^{-1}_1\cdot g_2 = s\in S$.
%
%(je ne crois pas  que ce soit un bonne idée de "noter" $s$,
%car cela donne plusieurs arêtes  avec le même nom)
%
%
La longueur de cette arête est déclarée par définition égale à $1$. Un
chemin $l\subset C(G,S)$ entre deux sommets $g$ et $h$ est une
succession finie d'arêtes $\{e_1,...,e_n\}$ joignant $g$ et $h$.
La longueur $|l|$ de $l$ vaut par définition $n$  : le nombre des
arêtes qui le constituent.

\begin{enumerate}
  \item  Montrer que la fonction $d:G \times G\mapsto \N$
  donnée par $$\displaystyle d(g,h)={\rm inf}\{{\rm longueurs\
  des\ chemins\ joignant\ } g {\rm\ et\ } h\}$$
  est une distance et qu'il existe un chemin $l\subset C(G,S)$ qui la
  réalise  c.-à.-d. $|l|=d(g,h)$.

  \item Pour chaque $g\in G$ posons $|g|=d(0,g)$. Montrer que
  $\displaystyle  |g|={\rm inf}\{k\ |\  g=s_{i_1}\cdot ... \cdot
  s_{i_k},\ s_{i_j}\in S\}$.

  \item Montrer que $G$ agit isométriquement sur les sommets de $C(G,S)$,
  c.-à.-d. $\ \forall g\in G\ d(g\gamma_1, g\gamma_2)=d(\gamma_1,\gamma_2)\
  {\rm o\grave u}\ \gamma_i\in V\ (i=1,2).$ En déduire que
  $d(f,h)=|f^{-1}\cdot h|\ (f,h\in V)$.



  \item   Soit $F_2=<a,b>$ un groupe libre sur les générateurs
  $a$ et $b$ (voir l'exercice \ref{pot:exo9}). Donner un fragment
  (initial) de son graphe de Cayley $C(F_2, \{a,b\})$.


  \item Démontrer que
  le graphe de Cayley d'un groupe libre est toujours un arbre (un
  graphe sans lacet s'appelle arbre).

\end{enumerate}
\finenonce{006384}



\finexercice


\exercice{7729, mourougane, 2021/08/11}
\enonce{007729}{Coloriages}
\begin{enumerate}
  \item Rappeler la formule de Burnside qui calcule le nombre d'orbites de l'action d'un groupe fini sur un ensemble fini.

 \item Rappeler la liste des éléments du groupe d'isométries directes (déplacements) d'un tétraédre régulier.
On fera une figure pour chaque type d'axe de rotation, en indiquant l'ordre des rotations.

 \item De combien de façons différentes peut-on peindre les faces d'un tétraèdre régulier avec $c$~couleurs ?
Chaque face n'est peinte que d'une couleur.
On ne distingue pas deux résultats qui se déduisent l'un de l'autre par un déplacement du tétraèdre.
\end{enumerate}
\finenonce{007729}
\finexercice


\exercice{7758, mourougane, 2021/08/11}
\enonce{007758}{}
On fixe une action d'un groupe $G$ sur un ensemble fini $E$.
On suppose 
que l'ordre de $G$ est 15,
 que le cardinal de $E$ est 17 
et que $E$ n'a pas de point fixé par tous les éléments du groupe $G$.
Déterminer le nombre d'orbites et le cardinal de chacune d'elles.
\finenonce{007758}
\finexercice


\exercice{7759, mourougane, 2021/08/11}
\enonce{007759}{Des petites questions}
On considère l'action d'un groupe $G$ sur un ensemble $E$.
\begin{enumerate}
\item Montrer qu'un sous-ensemble de $E$ est globalement stable par $G$ si
et seulement s'il est réunion d'orbites.
\item Montrer que deux éléments dans la même orbite ont des stabilisateurs
conjugués. 
\item Montrer que deux éléments conjugués dans le groupe $G$ fixent le même nombre d'éléments.
\end{enumerate}
\finenonce{007759}
\finexercice


\exercice{7760, mourougane, 2021/08/11}
\enonce{007760}{Le théorème de Cayley}
\begin{enumerate}
\item Pour tout élément $a$ d'un groupe fini $G$ d'ordre $n$, 
on définit l'application
$$\begin{array}{ccccc}
l_a&:&G&\to & G\\
  & &g&\mapsto &ag
\end{array}$$
Montrer que $l_a$ est une bijection de $G$, produit de $\frac{n}{ordre
 (a)}$ cycles à support disjoints tous de longueur $ordre (a)$.
\item Montrer alors que l'application
$$\begin{array}{ccccc}
l&:&G&\to & \S(G)\\
  & &a&\mapsto &l_a
\end{array}$$
est un morphisme de groupes, injectif.
Tout groupe fini est donc isomorphe à un sous-groupe du groupe des
permutations de ses éléments.
\end{enumerate}
\finenonce{007760}
\finexercice


\exercice{7761, mourougane, 2021/08/11}
\enonce{007761}{}
\begin{enumerate}
\item Soit $p$ un nombre premier. Montrer que le centre d'un $p$-groupe 
$G$,( i.e. un groupe fini d'ordre une
puissance non nulle de $p$), n'est pas réduit à l'élément neutre.
\item Montrer que si le quotient d'un groupe par son centre est
 cyclique alors le groupe est abélien, donc égal à son centre.
\item Montrer qu'un groupe d'ordre $p^2$ est abélien.
\item Montrer que le centre d'un groupe non abélien d'ordre $p^3$ est d'ordre $p$. En déduire que le nombre de classes de conjugaison est $p^2+p-1$. (On pourra étudier l'action de $G$ sur lui-même par
conjugaison~: ses points fixes, l'orbite des éléments, le
stabilisateur des éléments...) 
\end{enumerate}
\finenonce{007761}
\finexercice


\exercice{7762, mourougane, 2021/08/11}
\enonce{007762}{}
Soit $G$ un groupe fini. Soit $p$ le plus petit 
facteur premier de l'ordre de $G$.
Soit $H$ un sous-groupe de $G$ d'indice $p>1$.
\begin{enumerate}
\item Montrer que les orbites de l'action de $H$ sur $G/H$ (l'ensemble quotient
$G/H$ des classes à gauche de $G$ modulo $H$) par translation à gauche
sont réduites à des points.
\item Montrer que $H$ est distingué.
\end{enumerate}
\finenonce{007762}
\finexercice


\exercice{7771, mourougane, 2021/08/11}
\enonce{007771}{}
Soit $G$ un groupe agissant sur un ensemble $E$.
Soit $E_1$ et $E_2$ deux parties non vides et disjointes de $E$.
Soit $g_+$ et $g_-$ deux éléments de $G$ tels que toute puissance (positive ou négative) de $g_+$ envoie tout élément de $E_1$ dans $E_2$ et toute puissance $g_-$ envoie tout élément de $E_2$ dans $E_1$.
\begin{enumerate}
\item Montrer que les mots de la forme $g_+^{k_1}g_-^{l_1}g_+^{k_2}g_-^{l_2}\cdots g_+^{k_d}g_-^{l_d}g_+^{k_{d+1}}$ ne sont pas égaux à l'élément neutre $e_G$.
\item En déduire en utilisant une conjugaison qu'aucun mot du groupe engendré par $g_+$ et $g_-$ autre que le mot vide n'est égal à l'élément neutre. On dit alors que le groupe engendré par $g_+$ et $g_-$ est un groupe libre.
\item Que dire si on suppose seulement que $E_1$ n'est pas inclus dans $E_2$ ?
\item Montrer le sous groupe de $SL(2,\Zz)$ engendré par 
$$A:=\begin{pmatrix}1&2\\ 0&1\end{pmatrix}\quad \text{ et } \quad B:=\begin{pmatrix}1&0\\ 2&1\end{pmatrix}$$
est libre, en considérant l'action naturelle sur $\Rr^2$ et les domaines $\{(x,y)\in\Rr^2, |x|<|y|\}$ et $\{(x,y)\in\Rr^2, |x|>|y|\}$ délimités par les diagonales.
\end{enumerate}
\finenonce{007771}
\finexercice


\exercice{7772, mourougane, 2021/08/11}
\enonce{007772}{Sous-groupes finis de $GL_n(\Zz)$}
Si $p$ et $q$ sont deux entiers naturels, on note $p\wedge q$ le plus grand commun diviseur de $p$ et $q$, on note également $p|q$ si $p$ divise $q$. Si $m$ est un entier supérieur ou égal à 1, on note $\Phi_m(X)$ le polynôme cyclotomique d’ordre $m$, $$\Phi_m(X) =\Pi_{\{k\in\{1,..,m\}/k\wedge m=1\}}(X-e^{2ik\pi/m}).$$
On rappelle que $\Phi_m(X)$ est un polynôme unitaire à coefficients entiers, irréductible dans $\Qq[X]$. Le degré de $\Phi_m(X)$ est $\varphi(m)$ où $\varphi$ est la fonction indicatrice d’Euler.
On rappelle enfin que $X^m-1 =\Pi_{d|m}\Phi_d(X$).
\begin{enumerate}
 \item Soit $M\in GL_2(\Zz)$, d’ordre fini $m$.
\begin{itemize}
 \item Montrer que si $z$ est une racine complexe du polynôme caractéristique $\chi_M(X)$ alors $z$ est racine du polynôme 
$X^m-1$.
\item Montrer, en résolvant l’équation $\phi(k)=1$, qu’il y a exactement deux polynômes cyclotomiques de degré un.
\item Montrer de même qu’il y a exactement trois polynômes cyclotomiques de degré deux dont on donnera les expressions développées.
\item En déduire que le polynôme $\chi_M(X)$ appartient à l’ensemble
$\{X^2 + X + 1,X^2 + 1,X^2-X + 1,X^2-1, (X-1)^2, (X + 1)^2\}$.
\item En déduire que $m\in \{1, 2, 3, 4, 6\}$.
\item Donner une matrice compagnon de $GL_2(\Zz)$ d’ordre $6$.
\end{itemize}

 \item Soit $p$ un nombre premier, $p\geq 3$. On note $\mathbb{F}_p$ un corps de cardinal $p$. On rappelle que la surjection naturelle $\Zz\to \mathbb{F}_p$ induit un morphisme de groupes
$R_p~:~GL_n(Z)\to GL_n(\mathbb{F}_p)$.


Soit $M\in GL_n(\Zz)$ d’ordre $m\geq 2$ et dans le noyau de $R_p$. On suppose que $M$ n'est pas l'identité.
La matrice $M$ peut donc s'écrire $M = I_n +p^rN$
avec $r\in\Nn^\star$ et $N\in M_n(\Zz)-pM_n(\Zz)$.
\begin{itemize} 
 \item Montrer que $mp^rN\in p^{2r}M_n(\Zz)$. En déduire que $p$ divise $m$.
On pose alors $m = pm'$ et $M'= M^p$.
\item Montrer que $p$ divise $m'$.
\item En déduire une contradiction.
\end{itemize}

 \item Soit $G$ un sous-groupe fini de $GL_n(\Zz)$.
Montrer que $G$ est isomorphe à un sous-groupe de $GL_n(\mathbb{F}_p)$.

 \item Soit $G$ un sous-groupe fini de $GL_2(\Zz)$.
\begin{itemize}
 \item Montrer que le cardinal de $G$ est un diviseur de $48$.
\item Montrer que le cardinal de $G$ ne peut pas être égal à $48$.
(On pourra, éventuellement, étudier $\Phi_8(X)$ considéré comme un polynôme à coefficients dans $\mathbb{F}_3$.)
\end{itemize}
\end{enumerate}
\finenonce{007772}
\finexercice


\exercice{7773, mourougane, 2021/08/11}
\enonce{007773}{Classification des groupes d'ordre $8$}
Soit $G$ un groupe d'ordre $8$.
\begin{enumerate}
 \item Enumérer quatre groupes d'ordre $8$, deux à deux non isomorphes, et même $5$ si possible.

 \item On suppose que tous les éléments de $G$ sont d'ordre $2$.
Montrer que $G$ est abélien. Soit $a$ et $b$ deux éléments non neutres distincts de $G$. Montrer que $\{e,a,b,ab\}$ est un sous-groupe d'ordre $4$ de $G$. 
Déterminer un isomorphisme de $G$ avec un groupe connu.
 \item On suppose que $G$ admet un élément $a$ d'ordre $4$. Soit $b$ un élément hors du sous-groupe engendré par $a$. Montrer que $\langle a\rangle$ est distingué et que $b^2$ appartient à $\langle a\rangle$.
\begin{enumerate}
 \item Quel est l'ordre de $b$ si $b^2=a$ ou si $b^2=a^3$ ? Conclure dans ce cas.
\item Si $b^2=e$, montrer que $G$ est un produit semi-direct et en déduire un isomorphisme avec un groupe connu.
\item Si tous les éléments hors de $<a>$ ont un carré égal à $a^2$, établir la liste des éléments et la table de multiplication de $G$ à l'aide seulement de $a$ et $b$.
\end{enumerate}
\end{enumerate}
\finenonce{007773}
\finexercice


\exercice{7785, mourougane, 2021/08/11}
\enonce{007785}{}
On considère l'action du groupe $G:=\{-1,+1\}$ sur l'algèbre des polynômes $k[X,Y]$ par $(-1)\cdot X=-X$ et $(-1)\cdot Y=-Y$.
Déterminer l'algèbre des polynômes invariants.
Est-ce un anneau factoriel ?
Est-il une algèbre de polynômes ?
\finenonce{007785}
\finexercice

\section{305.00 Groupe cyclique}


\exercice{2659, matexo1, 2002/02/01}
\enonce{002659}{}
Montrer que tout sous-groupe d'un groupe cyclique est cyclique.
\finenonce{002659}

\finexercice


\exercice{2660, matexo1, 2002/02/01}
\enonce{002660}{}
Montrer qu'un groupe $G$ dont les seuls sous-groupes sont $G$ et $\{e_G\}$
est cyclique et que son ordre est premier.
\finenonce{002660}

\finexercice


\exercice{7730, mourougane, 2021/08/11}
\enonce{007730}{La réciproque du théorème de Lagrange}
\begin{enumerate}
  \item Le groupe symétrique $\mathcal{S}_5$ a-t-il un élément d'ordre $6$ ?
 \item Chercher dans un groupe symétrique un contre-exemple à la réciproque du théorème de Lagrange sur l'ordre d'un élément dans un groupe.
\end{enumerate}
\finenonce{007730}
\finexercice


\exercice{7764, mourougane, 2021/08/11}
\enonce{007764}{}
\begin{enumerate}
\item Montrer qu'un groupe d'ordre $p^3q$ (avec $p$ premier et $q$ premier avec $p$) admet un sous-groupe d'ordre $p$, un d'ordre $p^2$ et un d'ordre $p^3$.
\item Donner la liste des éléments du groupe alterné $\mathcal{A}_4$. Soit $H$ un sous-groupe d'ordre $3$ de $\mathcal{A}_4$ et $\sigma$ un élément de $\mathcal{A}_4$ qui n'est pas dans $H$.
Montrer que le sous-groupe engendré par $H$ et $\sigma$ est le groupe $\mathcal{A}_4$. En déduire que $\mathcal{A}_4$ n'a pas de sous-groupe d'ordre $6$.
\end{enumerate}
\finenonce{007764}
\finexercice


\exercice{7765, mourougane, 2021/08/11}
\enonce{007765}{Décompositions explicites}
\setcounter{MaxMatrixCols}{20}
\begin{enumerate}
\item On considère l'élément de $\mathcal{S}_8$ :
$$\sigma=\begin{pmatrix}1&2&3&4&5&6&7&8 \\ 5&6&4&1&2&8&3&7\end{pmatrix}.$$
Décomposer $\sigma$ en un produit de transpositions et calculer
sa signature. Peut-on écrire $\sigma$ comme produit de douze
transpositions ?
\item Soit $\sigma$ l'élément de $\mathcal{S}_{11}$ :
$$\sigma=\begin{pmatrix}
1&2&3&4&5&6&7&8&9&10&11\\
10&7&9&11&2&1&3&5&8&4&6
\end{pmatrix}.$$
Décomposer $\sigma$ en un produit de cycles à support disjoints. Préciser 
l'ordre de $\sigma$, et la signature de $\sigma$. Calculer $\sigma^2$
et $\sigma^3$. \'Ecrire $\sigma^{-1}$ 
en un produit de cycles à support disjoints.
\end{enumerate}
\finenonce{007765}
\finexercice


\exercice{7766, mourougane, 2021/08/11}
\enonce{007766}{}
\begin{enumerate}
\item Si $c$ est le cycle $(1,2,3,4,5,6)$, $c^2$ est-il un cycle ?
\item
Si $c$ est un cycle de $\mathcal{S}_n$ d'ordre $l$ et $k$ un entier naturel,
calculer l'ordre de $c^k$.
\end{enumerate}
\finenonce{007766}
\finexercice


\exercice{7767, mourougane, 2021/08/11}
\enonce{007767}{{\'E}tude de $\mathcal{S}_3$}
Donner les structures de
cycles possibles dans $\mathcal{S}_3$, le nombre d'éléments de $\mathcal{S}_3$ ayant 
cette structure, et leur signature. Décrire les
sous-groupes de $\mathcal{S}_3$, et ceux qui sont distingués dans $\mathcal{S}_3$.

Déterminer les sous-groupes de Sylow de $\mathcal{S}_3$.
\finenonce{007767}
\finexercice


\exercice{7768, mourougane, 2021/08/11}
\enonce{007768}{{\'E}tude de $\mathcal{S}_4$}
Donner les structures de
cycles possibles dans $\mathcal{S}_4$, le nombre d'éléments de $\mathcal{S}_4$ ayant 
cette structure, et leur signature.
 Déterminer les sous-groupes distingués 
de $\mathcal{S}_4$ (utiliser l'exer\-cice sur les classes de conjugaison). En déduire que $A_4$
n'est pas un groupe simple, i.e. qu'il 
possède des sous-groupes distingués autres que $\lbrace e\rbrace$ et 
$A_4$.

Déterminer les sous-groupes de Sylow de $\mathcal{S}_4$.
\finenonce{007768}
\finexercice


\exercice{7769, mourougane, 2021/08/11}
\enonce{007769}{}
On appelle groupe dihédral $D_{2n}$ le groupe des isométries d'un polygone régulier à $n$ côtés.
\begin{enumerate}
 \item Déterminer (par exemple à l'aide d'une action de groupe) le cardinal puis la liste des éléments de $D_{2n}$.
\item On suppose $n$ impair. Déterminer les $2$-Sylow de $D_n$ et vérifier qu'ils sont conjugués.
\item On suppose $n=6$. Déterminer un $2$-Sylow de $D_6$.
Déterminer deux sous-groupes d'ordre $2$ de $D_6$ non conjugués dans $D_6$.
\end{enumerate}
\finenonce{007769}
\finexercice


\exercice{7770, mourougane, 2021/08/11}
\enonce{007770}{}
Soit $G$ le sous-groupe de $\mathcal{S}_7$ engendré par 
$\alpha=(2,4,6)(5,7,1)$ et $\beta=(3,4)(5,6)$. On se
propose de déterminer l'ordre de $G$. On considère pour cela les 
ensembles suivants :
$$G_1=\lbrace \varphi\in G \mid \varphi(1)=1\rbrace \quad \quad 
G_2=\lbrace \varphi\in G_1 \mid
\varphi(2)=2\rbrace \quad \quad G_3=\lbrace \varphi\in G_2 \mid 
\varphi(3)=3\rbrace$$
$$X_1=\lbrace \varphi(1) \mid \varphi\in G\rbrace \quad \quad 
X_2=\lbrace \varphi(2) \mid
\varphi\in G_1\rbrace \quad \quad X_3=\lbrace \varphi(3) \mid 
\varphi\in G_2\rbrace.$$
{\'E}tant donné un ensemble $Y$, on note $|Y|$ le cardinal de $Y$.

\begin{enumerate}
\item Montrer que 6 divise $\mid G\mid$.
\item Quelle relation existe-t-il entre $|G|$ et $|X_1| \ |X_2| 
\ |X_3| \ |G_3|$ ?
\item Expliciter $X_1$.
\item Expliciter $\gamma=\alpha \beta \alpha^{-1}$ et 
$\delta=\gamma \beta \gamma^{-1}$. En déduire
$X_3=\{3,4,5,6\}$ ou $X_3=\{3,4,5,6,7\}$ et $X_2=\{2,7\}$ ou
$X_2=\{2,3, 4, 5,6,7\}$. 
\item On fait agir $G$ sur l'ensemble des parties de $\lbrace 
1,2,3,4,5,6,7\rbrace$. Déterminer
l'orbite de la partie $\lbrace 1,2,7\rbrace$. En déduire que 7 est 
fixé par les éléments de
$G_2$ et que $G_3$ est réduit à l'identité.
\item En déduire $|G|$.
\end{enumerate}
\finenonce{007770}
\finexercice


\exercice{7775, mourougane, 2021/08/11}
\enonce{007775}{}
\begin{enumerate}
 \item Soit $G$ un groupe, $a$ et $b$ deux éléments d'ordre fini qui commutent.
On suppose que les sous-groupes engendrés $\langle a\rangle$ et $\langle b\rangle$ ont une intersection réduite au singleton élément neutre $\{e\}$.
\item Montrer qu'une égalité $(ab)^m=e$ implique $a^m=e$ et $b^m=e$.
\item Calculer l'ordre de $ab$.
\item Montrer que si deux éléments d'un groupe ont des ordres premiers entre eux, l'intersection des sous-groupes qu'ils engendrent est réduite au singleton élément neutre. 
\item Montrer que tout groupe abélien d'ordre $77$ est cyclique. 
\end{enumerate}
\finenonce{007775}
\finexercice

\section{306.00 Théorème de Sylow}


\exercice{2190, debes, 2008/02/12}
\enonce{002190}{} 
Soient $G$ un groupe fini et $H$ un sous-groupe distingu\'e de $G$. Montrer
que si $H$ et $G/H$ sont des $p$-groupes, il en est de m\^eme de $G$. 
\finenonce{002190} 


\finexercice


\exercice{2191, debes, 2008/02/12}
\enonce{002191}{} 
Soit $G$ un $p$-groupe et $H$ un sous-groupe distingu\'e de $G$. 
Montrer que $H\cap Z(G)$ n'est pas r\'eduit \`a l'\'el\'ement neutre.
\finenonce{002191} 

\finexercice


\exercice{2192, debes, 2008/02/12}
\enonce{002192}{} 
Soit $G$ un $p$-groupe d'ordre $p^r$. 
\smallskip

(a) Montrer que pour tout entier $k\leq r$, $G$ poss\`ede un sous-groupe
distingu\'e d'ordre $p^k$. 
\smallskip

(b) Montrer qu'il existe une suite $G_0=\{1\} \subset G_1 \subset \ldots \subset
G_r=G$ de sous-groupes $G_i$ distingu\'es d'ordre $p^i$ ($i=1,\ldots,r$).
\smallskip

(c) Montrer que pour tout sous-groupe $H$ de $G$ d'ordre $p^s$ avec $s<r$, il
existe un sous-groupe d'ordre $p^{s+1}$ de $G$ qui contient $H$.
\finenonce{002192} 


\finexercice


\exercice{2193, debes, 2008/02/12}
\enonce{002193}{} 
\label{ex:deb93}
Soit $G$ un groupe d'ordre $2p$, o\`u $p$ est un nombre premier  
sup\'erieur ou \'egal \`a $3$. Montrer que $G$ contient un unique sous-groupe $H$
d'ordre $p$ et que ce sous-groupe est distingu\'e. V\'erifier que les seuls
automorphismes d'ordre $2$ d'un groupe cyclique d'ordre $p$ sont l'identit\'e et le
passage \`a l'inverse. En d\'eduire que le groupe $G$ est soit cyclique, soit
non commutatif, auquel cas il poss\`ede deux g\'en\'erateurs $s$ et $t$
v\'erifiant les relations $s^p=1$, $t^2=1$ et $tst^{-1} =s^{-1}$.
\finenonce{002193} 

\finexercice


\exercice{2194, debes, 2008/02/12}
\enonce{002194}{} 

Soit $G$ un groupe non commutatif d'ordre $8$.\smallskip 

(a) Montrer que $G$ contient un \'el\'ement $a$ d'ordre $4$ et que le
sous-groupe $H$ de $G$ engendr\'e par $a$ est distingu\'e dans $G$.
\smallskip

(b) On suppose ici qu'il existe un \'el\'ement $b$ de $G\setminus  H$ qui est
d'ordre $2$. Soit $K$ le sous-groupe engendr\'e par $b$. Montrer que dans ce cas $G$
est isomorphe au produit semi-direct de $H$ par $K$, le g\'en\'erateur $b$ de $K$
agissant sur $H$ via l'automorphisme $x
\rightarrow x^{-1}$. Le groupe est alors isomorphe au groupe di\'edral $D_4$.
\smallskip

(c) Dans le cas contraire, soit $b$ un \'el\'ement d'ordre $4$ de $G$
n'appartenant pas \`a $H$. Montrer que $a^2$ est le seul \'el\'ement d'ordre $2$ de
$G$, que le centre $Z(G)$ de $G$ est \'egal \`a $\{ 1, a^2 \}$. On pose
$-1=a^2$. Montrer que $a$ et $b$ v\'erifient les relations suivantes: $a^2=b^2=-1$,
$bab^{-1} =a^{-1}$. Enfin on pose $ab=c$. V\'erifier les relations suivantes:
$$a^2=b^2=c^2=-1 \quad ab=-ba=c \quad bc=-cb=a \quad ca=-ac=b$$
(l'\'ecriture $-x$ signifiant ici $(-1)x$). Ce dernier groupe est le groupe des
quaternions.
\finenonce{002194} 

\finexercice


\exercice{2195, debes, 2008/02/12}
\enonce{002195}{} 
Montrer que le groupe di\'edral $D_6$ est isomorphe au
produit direct $\mu _2 \times S_3$.
\finenonce{002195} 



\finexercice


\exercice{2196, debes, 2008/02/12}
\enonce{002196}{}
 (a) Soit $G$ un groupe non ab\'elien d'ordre $12$. Soit $H$ un $3$-Sylow
de $G$. On consid\`ere le morphisme $\theta: G \rightarrow S_{G/H}$ correspondant
\`a l'action de $G$ par translation de $G$ sur $G/H$. Montrer que ce morphisme n'est
pas injectif si et seulement si $H$ est distingu\'e dans $G$. En d\'eduire que si
$H$ n'est pas distingu\'e dans $G$, le groupe $G$ est isomorphe \`a $A_4$.
\smallskip

(b) On suppose que $G$ n'est pas isomorphe \`a $A_4$. Montrer qu'alors $G$ admet
un unique $3$-Sylow $H=\{ 1, a,a^2 \}$. Montrer ensuite que si $G$ contient un
\'el\'ement $b$ d'ordre $4$, $a$ et $b$ v\'erifient les relations:
$$a^3=b^4=1 \quad bab^{-1}= a^2=a^{-1}$$

Montrer que dans le cas contraire $G\simeq D_6$.
\smallskip

(c) Donner la liste des classes d'isomorphisme de groupes d'ordre $12$.
\finenonce{002196} 

\finexercice


\exercice{2197, debes, 2008/02/12}
\enonce{002197}{}
Soient $G$ un groupe et $H$ un sous-groupe distingu\'e de
$G$. On se donne un nombre premier $p$ et l'on suppose que $H$ admet
un unique $p$-Sylow $S$. Montrer que $S$ est distingu\'e dans $G$.
\finenonce{002197} 



\finexercice


\exercice{2198, debes, 2008/02/12}
\enonce{002198}{} 
Soient $G$ un groupe et $H$ un sous-groupe distingu\'e
de $G$. On se donne un nombre premier $p$ et un $p$-Sylow $P$ de
$G$. Montrer que $H\cap P$ est un $p$-Sylow de $H$ et que
$HP/H$ est un $p$-Sylow de $G/H$.
\finenonce{002198} 

\finexercice


\exercice{2199, debes, 2008/02/12}
\enonce{002199}{} 
Montrer qu'un groupe d'ordre $200$ n'est pas simple.
\finenonce{002199} 

\finexercice


\exercice{2200, debes, 2008/02/12}
\enonce{002200}{}
Pour $p$ un nombre premier, d\'eterminer le nombre de $p$-sous-groupes de Sylow du groupe sym\'etrique $S_p$.
\finenonce{002200} 

\finexercice


\exercice{2201, debes, 2008/02/12}
\enonce{002201}{} 
(a) Donner l'ensemble ${\cal D}$ des
ordres possibles des \'el\'ements du groupe altern\'e $A_5$ et pour
chaque $d\in {\cal D}$, indiquer le nombre d'\'el\'ements de $A_5$
d'ordre $d$.

(b) Montrer que, pour $d=2$ et $d=3$, les \'el\'ements
d'ordre $d$ sont conjugu\'es, et que les sous-groupes d'ordre $5$ sont
conjugu\'es.

(c) D\'eduire une preuve de la simplicit\'e de
$A_5$. 

\finenonce{002201} 

\finexercice


\exercice{2202, debes, 2008/02/12}
\enonce{002202}{}
D\'eterminer les sous-groupes de Sylow du groupe altern\'e $A_5$.
\finenonce{002202} 


\finexercice


\exercice{2203, debes, 2008/02/12}
\enonce{002203}{} 
Soit $G$ un groupe simple d'ordre $60$. \smallskip

(a) Montrer que $G$ admet $6$ $5$-Sylow, et que l'action de conjugaison sur ses
$5$-Sylow d\'efinit un morphisme injectif $\alpha : G \rightarrow S_6$, une fois
une num\'erotation des $5$-Sylow de $G$ choisie. Montrer que l'image $\alpha (G)=H$
est contenue dans $A_6$.
\smallskip

(b) On consid\`ere l'action de $A_6$ par translation \`a gauche sur l'ensemble
$A_6/.H$ des classes \`a gauche. Montrer qu'elle d\'efinit un isomorphisme
$\varphi : A_6 \rightarrow A_6$, une fois une num\'erotation des \'el\'ements de
$A_6/.H$ choisie. 
\smallskip

(c) Montrer que $\varphi (H) $ est le fixateur de la classe de l'\'el\'ement
neutre $H$, et en conclure que $G \simeq A_5$.
\finenonce{002203} 

\finexercice


\exercice{2204, debes, 2008/02/12}
\enonce{002204}{} 
\label{ex:deb104}
Soient $p<q$ deux nombres premiers distincts et $G$ un
groupe d'ordre $pq$. Montrer que $G$ admet un unique $q$-Sylow $Q$
qui est distingu\'e et que $G=QP$, o\`u $P$ est un $p$-Sylow de $G$.
Montrer que $G$ est isomorphe au produit semi-direct d'un groupe
cyclique d'ordre $q$ par un groupe cyclique d'ordre $p$. Montrer que
si $q-1$ n'est pas divisible par $p$, ce produit semi-direct est en
fait un produit direct.
\finenonce{002204} 


\finexercice


\exercice{2205, debes, 2008/02/12}
\enonce{002205}{} 
Montrer qu'un groupe d'ordre $35$ est cyclique.
\finenonce{002205} 


\finexercice


\exercice{2206, debes, 2008/02/12}
\enonce{002206}{}
 Soient $p$ et $q$ deux nombres premiers et $G$ un groupe
d'ordre $p^2q$. On suppose que $p^2-1$ n'est pas divisible par $q$ et
que $q-1$ n'est pas divisible par $p$. Montrer que $G$ est ab\'elien.
\finenonce{002206} 

\finexercice


\exercice{2207, debes, 2008/02/12}
\enonce{002207}{}
 Soient $p$ et $q$ deux nombres premiers. Montrer qu'il n'existe pas de
groupe simple d'ordre $p^2q$.
\finenonce{002207} 


\finexercice


\exercice{2208, debes, 2008/02/12}
\enonce{002208}{}
 Soit $G$ un groupe d'ordre $399$. \smallskip

(a) Montrer que $G$ admet un unique $19$-Sylow $P$ qui est distingu\'e dans $G$.
\smallskip

(b) Soit $Q$ un $7$-Sylow. Montrer que $N=PQ$ est un sous-groupe d'ordre $133$    
de $G$ et que ce groupe est cyclique.
\smallskip

(c) On suppose que $Q$ n'est pas distingu\'e dans $G$. Montrer que $G$ admet $57$
sous-groupes cycliques d'ordre $133$ distincts deux \`a deux. Quel serait le nombre
d'\'el\'ements d'ordre $133$ dans $G$? Aboutir \`a  une contradiction. En d\'eduire
que $Q$ est distingu\'e dans $G$ et que $N$ est distingu\'e dans $G$.
\smallskip

(d) Montrer que $G=NR$, o\`u $R$ est un $3$-Sylow. En d\'eduire que $G$ est
isomorphe au produit semi-direct d'un groupe cyclique d'ordre $133$ par un groupe
cyclique d'ordre $3$.
\finenonce{002208} 

\finexercice


\exercice{2209, debes, 2008/02/12}
\enonce{002209}{}
 Soit $G$ un groupe simple d'ordre $60$. 
\smallskip

(a) Montrer que $G$ n'admet pas de sous-groupe d'ordre $20$. 
\smallskip

(b) Montrer que si $G$ admet un sous-groupe $K$ d'ordre $12$, alors $K$ admet $4$
$3$-Sylow. 
\smallskip

(c) Montrer que si $H$ et $K$ sont deux sous-groupes distinct d'ordre $4$ de $G$
alors $H\cap K=\{1\}$. 
\smallskip

(d) Montrer que si $H$ est un $2$-Sylow, alors $H\not= \hbox{\rm Nor}_G(H)$.
\smallskip

(e) Montrer que $G$ poss\`ede $5$ $2$-Sylow.
\smallskip

(f) Conclure en consid\'erant l'action de $G$ par conjugaison sur les $5$-Sylow.
\finenonce{002209} 


\finexercice


\exercice{7739, mourougane, 2021/08/11}
\enonce{007739}{Sylow des groupes diédraux}
Soit $\mathcal{P}_n$ un polygone régulier à $n$ côtés dans le plan euclidien orienté.
On appelle groupe diédral $D_{n}$ le groupe des isométries de $\mathcal{P}_n$.
\begin{enumerate}
    \item Rappeler sans démonstration la liste complète des éléments de $D_{n}$.
    \item Décrire un sous-groupe cyclique d'ordre $n$ dans $D_{n}$.
    \item On écrit $n=2^kn'$ où $n'$ est impair. En considérant un polygone régulier inscrit, décrire une application injective de $D_{2^{k}}$ dans $D_{n}$. Faire une figure dans le cas $n=12$.
    \item Décrire pour chaque diviseur premier $p$ de $2n$, un $p$-Sylow de $D_{n}$.
\end{enumerate}
\finenonce{007739}
\finexercice


\exercice{7740, mourougane, 2021/08/11}
\enonce{007740}{Sylow des groupes diédraux}
Soit $\mathcal{P}_n$ un polygone régulier à $n$ côtés dans le plan euclidien orienté.
On appelle groupe diédral $D_{n}$ le groupe des isométries de $\mathcal{P}_n$.
\begin{enumerate}
 \item Parmi les translations, les rotations, les symétries orthogonales, et les symétries glissées, décrire des isométries du plan qui conservent le polygone régulier $\mathcal{P}_n$.

 \item Déterminer, à l'aide de l'action naturelle de $D_{n}$ sur l'ensemble des sommets de $\mathcal{P}_n$, le cardinal de $D_{n}$. En déduire la liste complète des éléments de $D_{n}$.

 \item On suppose $n$ impair. Déterminer les $2$-Sylow de $D_{n}$ et vérifier (sans référence au cours) qu'ils sont conjugués.

 \item On suppose $n=6$. Déterminer un $2$-Sylow de $D_{6}$.
Ce $2$-Sylow est-il distingué ?
Déterminer deux sous-groupes d'ordre $2$ de $D_6$ non conjugués dans $D_6$.
Donner un $3$-Sylow de $D_{6}$.
\end{enumerate}
\finenonce{007740}
\finexercice


\exercice{7741, mourougane, 2021/08/11}
\enonce{007741}{Les groupes d'ordre 33}
Déterminer à isomorphisme près tous les groupes d'ordre $33$.
\finenonce{007741}
\finexercice


\exercice{7763, mourougane, 2021/08/11}
\enonce{007763}{}
Déterminer les sous-groupes de Sylow de $\Zz/24\Zz$.
\finenonce{007763}
\finexercice


\exercice{7774, mourougane, 2021/08/11}
\enonce{007774}{Théorème de Sylow pour les groupes abéliens finis}
Soit $G=\{a_1,a_2,\cdots a_n\}$ un groupe fini. On note $\alpha_i$ l'ordre de l'élément $a_i$. Soit $p$ un diviseur premier de $\mid G\mid$ l'ordre de $G$.
\begin{enumerate}
 \item Montrer que l'application
$$\begin{array}{cccc}
 f~:&\Zz/\alpha_1\Zz\times\Zz/\alpha_2\Zz\times\cdots\times\Zz/\alpha_n\Zz
&\to &G\\
& (\overline{h_1},\overline{h_2},\cdots,\overline{h_n})&\mapsto & a_1^{h_1}a_2^{h_2}\cdots a_n^{h_n}
\end{array}$$
est une application bien définie. Démontrer que c'est un homomorphisme de groupes puis qu'il est surjectif.
\item En déduire que $p$ divise $\alpha_1\alpha_2\cdots\alpha_n$.
\item Montrer qu'il y a dans $G$ un élément d'ordre $p$.
\item En raisonnant par récurrence sur l'ordre du groupe et en considérant l'ensemble $G/<x>$ où $x$ est un élément d'ordre $p$ dans $G$, montrer que $G$ admet un $p$-Sylow.
\item Montrer qu'un groupe abélien est simple si et seulement s'il est cyclique, d'ordre un nombre premier.
\end{enumerate}
\finenonce{007774}
\finexercice


\exercice{7846, mourougane, 2021/08/11}
\enonce{007846}{Autour de cours}
\begin{enumerate}
    \item Montrez que les morphismes d'un groupe simple vers un groupe quelconque sont constants ou injectifs.
    \item Quel est le nombre moyen de points fixes d'une permutation de $\mathcal{S}_n$ ?
    \item Démontrez que si $G$ est un groupe fini, $S$ un $p$-Sylow de $G$, et $H$ un sous-groupe de $G$, il existe un conjugué de $S$ qui rencontre $H$ en un $p$-Sylow de $H$. 
\end{enumerate}
\finenonce{007846}
\finexercice


\exercice{7847, mourougane, 2021/08/11}
\enonce{007847}{}
Soit $G$ un groupe fini. Soit $p$ le plus petit facteur premier de l'ordre de $G$. Soit $H$ un sous-groupe de $G$ d'indice $p>1$.
\begin{enumerate}
    \item Montrez que les orbites de l'action de $H$ sur $G/H$ par translation à gauche sont réduites à des points.
    \item Montrez que $H$ est distingué.
\end{enumerate}
\finenonce{007847}
\finexercice


\exercice{7848, mourougane, 2021/08/11}
\enonce{007848}{}
\begin{enumerate}
    \item Explicitez un $7$-Sylow du groupe symétrique $\mathcal{S}_{9}$
    \item Déterminez le nombre d'élements d'ordre $7$ dans $\mathcal{S}_{9}$.
    \item Déterminez le nombre de $7$-Sylows.
    \item Vérifiez les congruences données par le théorème de Sylow sur le nombre de $7$-Sylows.
\end{enumerate}
\finenonce{007848}
\finexercice


\exercice{7849, mourougane, 2021/08/11}
\enonce{007849}{}
On fixe une action d'un groupe $G$ sur un ensemble fini $E$.
On suppose 
que l'ordre de $G$ est 15,
 que le cardinal de $E$ est 17 
et que $E$ n'a pas de point fixé par tous les éléments du groupe $G$.
Déterminer le nombre d'orbites et le cardinal de chacune d'elles.
\finenonce{007849}
\finexercice


\exercice{7850, mourougane, 2021/08/11}
\enonce{007850}{}
\begin{enumerate}
 \item Soit $G$ un $p$-groupe agissant sur un ensemble fini $E$.
 Montrer que le cardinal de l'ensemble des points fixes de l'action
 est congru, modulo $p$, au cardinal de $E$.
 \item En considérant une action de $G$ sur lui-même, montrer que le théorème de Burnside : 
 le centre d'un $p$-groupe non réduit à l'élément neutre
 n'est pas réduit à l'élément neutre.
\end{enumerate}
\finenonce{007850}
\finexercice


\exercice{7851, mourougane, 2021/08/11}
\enonce{007851}{Groupe d'ordre $p^3$}
Soit $G$ un groupe non abélien d'ordre $p^3$ où $p$ est un
nombre premier. 
\begin{enumerate}
\item Montrer que le centre de $G$ est d'ordre $p$ et égal à son sous-groupe dérivé $Z(G)=D(G)$.
\item En déduire que le nombre de classes de conjugaison est $p^2+p-1$. (On pourra étudier l'action de $G$ sur lui-même par
conjugaison~: ses points fixes, l'orbite des éléments, le
stabilisateur des éléments et appliquer la formule de Burnside...) 
\item Montrer que $G/Z(G)$ est isomorphe à $\Z/p\Z\times \Z/p\Z$.
\item Montrer que tout sous-groupe de $G$ d'ordre $p^2$ contient le centre $Z(G)$ de $G$,
et que donc $G$ n'est pas un produit semi-direct de son centre par son abélianisé.
% \href{https://groupprops.subwiki.org/wiki/Semidirect_product_of_cyclic_group_of_prime-square_order_and_cyclic_group_of_prime_order}{Voir sur les groupes d'ordre $p^3$}
\end{enumerate}
\finenonce{007851}
\finexercice


\exercice{7852, mourougane, 2021/08/11}
\enonce{007852}{}
\begin{enumerate}
 \item En comptant le nombre de base de $\mathbb{F}_p^n$ déterminer le cardinal de $GL(n,\mathbb{F}_p)$.
 \item Montrer que l'ensemble des matrices triangulaires supérieures strictes est un $p$-sous-groupe de Sylow de $GL(n,\mathbb{F}_p)$.
\end{enumerate}
\finenonce{007852}
\finexercice


\exercice{7853, mourougane, 2021/08/11}
\enonce{007853}{}
\begin{enumerate}
\item Vérifier que les $p$-Sylow de $GL(2,\mathbb{F}_p)$ sont monogènes.
\item Soit $A$ et $B$ deux matrices de $GL(2,\mathbb{F}_p)$ d'ordre $p$.
Montrer que $A$ est conjuguée à une puissance de $B$.
\end{enumerate}
\finenonce{007853}
\finexercice


\exercice{7854, mourougane, 2021/08/11}
\enonce{007854}{}
Soit $p$ un nombre premier et $m$ un entier non multiple de $p$.
Soit $G$ un groupe de cardinal $|G|=p^dm$.
\begin{enumerate}
 \item Montrer que le nombre de $p$-Sylow de $G$ divise $m$.
 \item Montrer que pour tout $0\leq i\leq d$, $G$ possède un sous groupe d'ordre $p^i$.
\end{enumerate}
\finenonce{007854}
\finexercice


\exercice{7855, mourougane, 2021/08/11}
\enonce{007855}{}
Déterminer les sous-groupes de Sylow de $\Zz/24\Zz$.
\finenonce{007855}
\finexercice


\exercice{7856, mourougane, 2021/08/11}
\enonce{007856}{Groupes de matrices sur $\mathbb{F}_2$}
\begin{enumerate}
\item Décrire un $2$-Sylow de $GL_3(\mathbb{F}_2)$.
\item Soit $A=\begin{pmatrix}1&1&0\\1&0&1\\0&1&0\end{pmatrix}$.
Montrer que le polynôme minimal de $A$ est irréductible de degré 3.
En déduire que $GL_3(\mathbb{F}_2)\cap\mathbb{F}_2[A]$ est un $7$-Sylow de $GL_3(\mathbb{F}_2)$.
\item Déterminer un $3$-Sylow de $GL_3(\mathbb{F}_2)$ à l'aide de la matrice 
$B=\begin{pmatrix}1&0&0\\0&1&1\\0&1&0\end{pmatrix}$.
\end{enumerate}
\finenonce{007856}
\finexercice


\exercice{7861, mourougane, 2021/08/11}
\enonce{007861}{Des petites questions}
\begin{enumerate}
\item Soit $p$ un nombre premier. Déterminer à isomorphisme près, tous les groupes d'ordre $p$.
\item Soit $p$ un nombre premier. Déterminer à isomorphisme près, tous les groupes d'ordre $p^2$.
\item Donner des exemples de groupes d'ordre $6$ non abéliens.
\item Déterminer l'ordre des groupes diédraux $D_n$.
\item Déterminer l'ordre des groupes alternés $\mathcal{A}_n$.
\end{enumerate}
\finenonce{007861}
\finexercice


\exercice{7862, mourougane, 2021/08/11}
\enonce{007862}{{\'E}tude de $\mathcal{S}_3$}
Donner les structures de
cycles possibles dans $\mathcal{S}_3$, le nombre d'éléments de $\mathcal{S}_3$ ayant 
cette structure, et leur signature. Décrire les
sous-groupes de $\mathcal{S}_3$, et ceux qui sont distingués dans $\mathcal{S}_3$.

Déterminer les sous-groupes de Sylow de $\mathcal{S}_3$.
\finenonce{007862}
\finexercice


\exercice{7863, mourougane, 2021/08/11}
\enonce{007863}{Groupes d'ordre 6}
Soit $G$ un groupe d'ordre $6$.
\begin{enumerate}
 \item Montrer que $G$ admet un élément $\tau$ d'ordre $2$ et un élément $\sigma$ d'ordre $3$.
 \item Quelles sont les valeurs possibles de $\tau\sigma\tau$ ?
 \item Déterminer, dans chacun des cas précédents, la structure de $G$ à isomorphisme près. 
\end{enumerate}
\finenonce{007863}
\finexercice


\exercice{7864, mourougane, 2021/08/11}
\enonce{007864}{{\'E}tude de $\mathcal{S}_4$}
Donner les structures de
cycles possibles dans $\mathcal{S}_4$, le nombre d'éléments de $\mathcal{S}_4$ ayant 
cette structure, et leur signature.
 Déterminer les sous-groupes distingués 
de $\mathcal{S}_4$. En déduire que $\mathcal{A}_4$
n'est pas un groupe simple, i.e. qu'il 
possède des sous-groupes distingués autres que $\lbrace e\rbrace$ et 
$\mathcal{A}_4$.

Déterminer les sous-groupes de Sylow de $\mathcal{S}_4$.
\finenonce{007864}
\finexercice


\exercice{7865, mourougane, 2021/08/11}
\enonce{007865}{}
Soit $p$ un nombre premier impair, on se propose de décrire les groupes d'ordre $p^2$ à isomorphisme près.
\begin{enumerate}
\item Soit $G$ un groupe de cardinal $p^2$, montrer que, ou bien $G$ est cyclique ou bien
tous les éléments différents de l’élément neutre sont d’ordre $p$.
\item Soit $G$ un groupe non cyclique d’ordre $p^2$, soit $K$ un sous-groupe d’ordre $p$, montrer que $K$ est distingué dans $G$ et qu’il existe $H$ sous-groupe d’ordre $p$ tel que $K\cap H = \{e\}$. En déduire que $G$ est ismorphe à un produit semi-direct de $\Z/p\Z$
par $\Z/p\Z$.
\item Montrer que tout groupe de cardinal $p^2$ est isomorphe à $\Z/p^2\Z$ ou $\Z/p\Z \times \Z/p\Z$.
\end{enumerate}
\finenonce{007865}
\finexercice


\exercice{7866, mourougane, 2021/08/11}
\enonce{007866}{Groupes d'ordre $pq$}
Soit $G$ un groupe d'ordre $pq$, où $p$ et $q$ sont deux nombres premiers distincts.
 On suppose que $p<q$.
\begin{enumerate}
 \item Montrer qu'il n'y a qu'un $q$-Sylow $Q$ et qu'il est distingué.
 \item Montrer que $G$ est produit semi-direct $Q\rtimes P$ où $P$ est un $p$-Sylow de $G$.
 \item Si $p$ ne divise pas $q-1$, déterminer la structure de $G$.
 \item Si $p=2$, déterminer le morphisme structurel $\phi : \Z/p\Z\to Aut(\Z/q\Z)\equiv \Z/(q-1)\Z$.
 Déterminer alors la structure de $G$.
 \item Si $p$ divise $q-1$, montrer qu'il n'y a qu'un seul produit semi-direct non abélien, à isomorphisme près.
\end{enumerate}
\finenonce{007866}
\finexercice


\exercice{7867, mourougane, 2021/08/11}
\enonce{007867}{Groupe non abélien d'ordre $8$}
Soit $G$ un groupe d'ordre $8$.
\begin{enumerate}
\item Enumérer quatre groupes d'ordre $8$, deux à deux non isomorphes, et même $5$ si possible.
\item
On suppose que tous les éléments de $G$ sont d'ordre $2$.
Montrer que $G$ est abélien. Soit $a$ et $b$ deux éléments non neutres distincts de $G$. Montrer que $\{e,a,b,ab\}$ est un sous-groupe d'ordre $4$ de $G$. 
Déterminer un isomorphisme de $G$ avec un groupe connu.
\item
On suppose que $G$ admet un élément $a$ d'ordre $4$. Soit $b$ un élément hors du sous-groupe engendré par $a$. Montrer que $<a>$ est distingué et que $b^2$ appartient à $<a>$.
\begin{enumerate}
 \item Quel est l'ordre de $b$ si $b^2=a$ ou si $b^2=a^3$ ? Conclure dans ce cas.
\item Si $b^2=e$, montrer que $G$ est un produit semi-direct et en déduire un isomorphisme avec un groupe connu.
\item Si tous les éléments hors de $<a>$ ont un carré égal à $a^2$, établir la liste des éléments et la table de multiplication de $G$ à l'aide seulement de $a$ et $b$.
\end{enumerate}
\end{enumerate}
\finenonce{007867}
\finexercice


\exercice{7868, mourougane, 2021/08/11}
\enonce{007868}{Groupe non abélien d'ordre $8$}
Soit $G$ un groupe non abélien d'ordre $8$.
\begin{enumerate}
\item Montrer que $G$ contient un élément d'ordre $4$. Soit $H$ le sous-groupe qu'il engendre.
\item Montrer que si $G-H$ contient un élément d'ordre $2$, $G$ est un produit semi-direct.
Après avoir vérifié que $Aut(\Z/4\Z)=\Z/2\Z$, montrer qu'il existe une unique structure de tel produit semi-direct
non abélien.
\item Montrer que si $G-H$ n'a pas d'élément d'ordre $2$, on retrouve la table de $H_8$
en choisissant $i$ l'élément d'ordre $4$ qui engendre $H$ et $j$ un élément d'ordre $4$ dans $G-H$.
\textit{On pourra montrer que $i^2$ est le seul élément d'ordre $2$ est qu'il est donc central.
On le notera $-1$}.
\end{enumerate}
\finenonce{007868}
\finexercice


\exercice{7869, mourougane, 2021/08/11}
\enonce{007869}{Les groupes d'ordre 10}
Soit $G$ un groupe d'ordre $10$.
\begin{enumerate}
\item Montrer que $G$ est un produit semi-direct.
\item Déterminer les automorphismes d'ordre $2$ de $\Z/5\Z$.
\item En déduire les deux possibilités pour les classes d'isomorphismes de $G$.
\end{enumerate}
\finenonce{007869}
\finexercice


\exercice{7870, mourougane, 2021/08/11}
\enonce{007870}{Les groupes d'ordre 33}
Déterminer à isomorphisme près tous les groupes d'ordre $33$.(On pourra déterminer le nombre de sous-groupe d'ordre $11$ et le nombre de sous-groupe d'ordre $3$.)
\finenonce{007870}
\finexercice

\section{307.00 Autre}

\section{310.00 Isométrie euclidienne}


\exercice{6385, potyag, 2011/10/16}
\enonce{006385}{}
Soit $E$ un espace vectoriel euclidien de dimension $n$.
\begin{enumerate}
\item Montrer que $A\in GL(E)$ appartient à $\mathcal{O}(n)$ si et seulement si
${ }^T AA =I$.

\item Montrer que si $A\in \mathcal{O}(n)$ alors $\det A = \pm 1$.

\item Montrer que $A \in U(n)$ si et seulement si ${}^T A\bar A =I$.
\end{enumerate}
\finenonce{006385}



\finexercice


\exercice{6386, potyag, 2011/10/16}
\enonce{006386}{}
Soit $L$  un espace hermitien. Est-il vrai que $A\in Iso L$  implique $Ax = Ux+b$
avec $U \in U(n)$.
\finenonce{006386}



\finexercice


\exercice{6387, potyag, 2011/10/16}
\enonce{006387}{}
Soit $E$ est un espace  euclidien de dimension $n$. Montrer que $Iso(E) \not\simeq
\mathcal{O}(n) \times T(E)$.
\finenonce{006387}



\finexercice


\exercice{6388, potyag, 2011/10/16}
\enonce{006388}{}
Déterminer la nature des applications suivantes : 
$\Rr^2 \rightarrow \Rr^2$, $x  \mapsto Ax$ et  
$\Rr^2 \rightarrow \Rr^2$, $x \mapsto Bx$
où  $A = 
\begin{pmatrix}
  \cos \varphi  & -\sin \varphi \\
\sin \varphi  & \cos \varphi  
\end{pmatrix} $
 et $B = \begin{pmatrix}
  \cos \varphi  & \sin \varphi \\
\sin \varphi  & -\cos \varphi  
\end{pmatrix}$.
\finenonce{006388}



\finexercice


\exercice{6389, potyag, 2011/10/16}
\enonce{006389}{}
\begin{enumerate}
\item Notons $l \subset \Rr^2$ une droite affine et $ \tau_l$ la réflexion
par rapport à $l$. Montrer que si $f \in Iso (\Rr^2)$ vérifie $f_{\vert
l}\equiv id$ alors soit $f = id $ soit $ f= \tau _l$.

\item Soient $l$ et $m$ deux droites affines dans $\Rr^2$.
  
  \begin{enumerate}
  \item Montrer qu'il existe $\alpha \in Iso (\Rr^2)$ telles que
  $\alpha \tau_l \alpha^{-1}= \tau_m. $
 
  \item Montrer que $\tau _l .\tau _m $ est une translation si et seulement si $l$ et
  $m$ sont parallèles.
  \end{enumerate}
\end{enumerate}
\finenonce{006389}



\finexercice


\exercice{6390, potyag, 2011/10/16}
\enonce{006390}{} %exo 6
\label{exo:avant}
Notons $R(a,\alpha)$ la rotation d'angle $\alpha$ autour du point
$a\in\Rr^2$ et $t_b$ la translation $t_b:x\mapsto x+b$. Montrer
que
\begin{enumerate}
\item
$\exists\ \beta\in Iso (\Rr^2) :\ \beta
R(a,\alpha)\beta^{-1}\in SO(2)$.

\item $R(a,\alpha)=\tau_l\cdot\tau_m$ où $m$ est une droite quelconque
passant par $a$ et $l$ est une droite passant par $a$ fixée.

\item $t_b$ et $t_c$ sont conjuguée ssi $\vert\vert
b\vert\vert=\vert\vert c\vert\vert$.
\end{enumerate}
\finenonce{006390}



\finexercice


\exercice{6391, potyag, 2011/10/16}
\enonce{006391}{} 
Une application du type $G(l,a)= \tau _l. t_a$ s'appelle {\it réflexion
glissée} si le vecteur $a$ est parallèle à la droite $l$ $\subset \Rr^2$.
\begin{enumerate}
\item Si $G = \tau _l. t_a$ est une réflexion glissée alors
montrer que $\tau_l . t_a = t_a  .\tau _l$ et  $G^2 = t_{2a}$.

\item Montrer que $G = \tau _l t_a$ est une réflexion si $ l$ et
$a$ sont perpendiculaires et est une réflexion glissée si $l $
et $a$ ne sont pas perpendiculaires.


\item En regardant l'ensemble des points fixes $fix(f) := \{ x \in \Rr^2 \vert
f(x) =x\}$ d'une isométrie $f \in Iso(\Rr^2)$ montrer que :
  \begin{enumerate}
  \item si $fix (f) \not= \emptyset$ alors $f=R(a,\alpha )$ ou $f = \tau_l$

  \item si $fix (f) = \emptyset$ alors $f=t_a$ ou $f=G(l,a)$ ({\it indication : utiliser
  la question 2. et l'exercice \ref{exo:avant}, question 2}).
  \end{enumerate}
\end{enumerate}
\finenonce{006391}



\finexercice


\exercice{6392, potyag, 2011/10/16}
\enonce{006392}{}
 Notons $l \subset \Rr^2$ une droite affine de $\Rr^2$.
\begin{enumerate}
\item Montrer que l'ensemble $I_l$  des $g \in Iso (\Rr^2)$ telles que $g(l)=l$
 est un sous-groupe de $Iso (\Rr)$.

\item Déterminer les translations qui appartiennent à $I_l$.

\item Montrer que si $g\in I_l$ possède un point fixe alors $g$ a
un point fixe sur $l$.

\item Soit $g\in I_l$ , montrer qu'il existe une translation
$t$ de $I_l$ telle que $g .t$  possède  un point fixe.

\item Décrire $I_l$.
\end{enumerate}
\finenonce{006392}



\finexercice


\exercice{6393, potyag, 2011/10/16}
\enonce{006393}{}
Soit $E$ un espace euclidien de dimension $n$ et $X=\{ x_1, x_2,...,x_k\}$ un
sous-ensemble de $E$.
\begin{enumerate}
\item Montrer que $\displaystyle Aff(X):= \{ \sum_{i=1}^{k} \lambda_i x_i \vert x_i \in
X ,
\sum_{i=1}^k \lambda _i=1\}$ est le plus petit sous-espace affine de $E$ contenant
$X$.

\item Soit $S= \{x_0, ..., x_n \}$  un
sous-ensemble de $E$, montrer que $S$ est un repère affine de $E$ si et
seulement si $S$ n'est contenu dans aucun hyperplan.
\end{enumerate}
\finenonce{006393}



\finexercice


\exercice{6394, potyag, 2011/10/16}
\enonce{006394}{}
Etudier la composée de deux rotations, puis la composée de
deux réflexions glissées et finalement la composée d'une
rotation et d'une réflexion glissée.
\finenonce{006394}



\finexercice


\exercice{6395, potyag, 2011/10/16}
\enonce{006395}{}
Soit $f$ une  application  qui préserve les rapports de longueur
: $ \forall x,y,z,t \in \Rr^2$,on a : $\frac{d(f(x),f(y))}{d(f(z), f(t))} = \frac{d(x,y)}{d(z,t)}$. (Par définition une
telle application est une {\it similitude})
\begin{enumerate}
\item Monter $\exists k \in \Rr^{+*}$ tel que $d(f(x),f(y)) =
kd(x,y)$.

\item Montrer qu'un similitude s'écrit comme composée d'une
homothétie et d'une isométrie.
\end{enumerate}
\finenonce{006395}



\finexercice


\exercice{6396, potyag, 2011/10/16}
\enonce{006396}{}
Soit $A B C$ un triangle isocèle en $A$ non equilatéral, le
but de cet exercice est  d'étudier l'ensemble des isométries
de $\mathcal{P}$ qui préservent globalement $A B C$.
\begin{enumerate}
\item  Montrer que cet ensemble est groupe.

\item Montrer que si $f$ préserve $A B C $ alors $f$ fixe le
barycentre $G$ de $A B C $.


\item En étudiant les distances $GA$, $GB$, $GC$ montrer que
$f(A)=A$

\item En déduire (en utilisant la classification des isométries
de $\Rr^2$) le groupe de symétries de $ABC$.
\end{enumerate}
\finenonce{006396}



\finexercice


\exercice{6397, potyag, 2011/10/16}
\enonce{006397}{\'Etude des sous-groupes commutatifs d'isométries}

Soient $f$, $g$ deux isométries du plan affine euclidien $\Rr^2$.
\begin{enumerate}
\item Montrer que si $g$ et $f$ commutent alors
$g=f\circ g\circ f^{-1}$  et $f$ conserve (globalement) l'ensemble
des points fixes de $g.$

\item Décrire les cas dans lesquels $f$ et $g$
commutent ($f\circ g =g\circ f$).

\item En déduire une description des sous-groupes
commutatifs d'isométries.
\end{enumerate}
\finenonce{006397}



\finexercice

\section{311.00 Géométrie différentielle élémentaire de Rn}


\exercice{6398, potyag, 2011/10/16}
\enonce{006398}{}

\begin{enumerate}
\item  Soient $\xi(t)$ et $\eta(t)$ deux courbes paramétrées de $\Rr^3$ de classe
$C^1$, montrer que $ \frac{d}{dt} [\xi, \eta]= [\frac{d\xi}{dt}, \eta] +[
{\xi, \frac{d\eta}{dt}}]$, où $[, ]$ désigne le produit vectoriel.

\item Montrer que $\kappa = - \frac{ <r', [r '', r ''']>}{k^2}$, où $ s\mapsto
r(s)$ une courbe de classe $C^2$ paramétrée par sa longueur, $\kappa$ et
$k$ sont respectivement sa torsion et sa courbure.
\end{enumerate}
\finenonce{006398}



\finexercice


\exercice{6399, potyag, 2011/10/16}
\enonce{006399}{}
Pour la courbe $r = (a \cos t, a\sin t, bt)$, $a>0$, $b>0$ trouver  courbure,
 torsion et  repère de Frenet.
\finenonce{006399}



\finexercice


\exercice{6400, potyag, 2011/10/16}
\enonce{006400}{}
\begin{enumerate}
\item Montrer qu'une  courbe $s \mapsto r (s)$ est plane ssi $<r ', [r'',
r''']>=0$.

\item Soit  $s \mapsto r (s)$ une courbe de classe $C^2$ paramétrée par sa
longueur. On considère la nouvelle courbe $s \mapsto n(s)$ où $n$ est le
vecteur normal unitaire. Notons $s^*$ le paramètre naturel de cette courbe.
Montrer que $\frac{ds^{*}}{ds}=
 \sqrt {k^2+\kappa ^2}$.
\end{enumerate}
\finenonce{006400}



\finexercice


\exercice{6401, potyag, 2011/10/16}
\enonce{006401}{}
\begin{enumerate}
\item Montrer que si une courbe $s\mapsto r(s)$ est tracée sur une sphère de
rayon $R$ et si $ \kappa(s)\not=0, k(s)\not=0\ (\forall s)$ alors $R^2 = <r ,r
> = \frac{1}{k^2 }(1+ \frac{(k')^2}{(\kappa k)^2})).$

\item Soit $s \mapsto r (s) $ une courbe à courbure constante qui est tracée
sur la sphère $S^2$. Montrer que son image $r $ est un arc de cercle.
La propriété d'être à courbure constante dépend-t-elle de
la paramétrisation de la courbe ?
\end{enumerate}
\finenonce{006401}



\finexercice


\exercice{6402, potyag, 2011/10/16}
\enonce{006402}{}
\begin{enumerate}
\item Montrer que $3$ points $x, y$ et $z$ sont colinéraires
dans $\Rr^n$ avec $y$ entre $x$ et $z$ (rappelons que \c{c}a
signifie que $y=x+t(z-x)\ t\in[0,1])$ ssi

$$\vert\vert x-y\vert\vert + \vert\vert y-z\vert\vert = \vert\vert
x-z\vert\vert.$$

\item Montrer que si $\gamma :[a,b]\mapsto {\Rr}^n$ est une courbe alors
$\vert\gamma([a,b])\vert\geq\vert\vert\gamma(a)-\gamma(b)\vert\vert$ et que
l'égalité a lieu ssi $\gamma$ est une géodésique (où $\vert\cdot\vert$
désigne la longueur d'une courbe).
\end{enumerate}
\finenonce{006402}



\finexercice

\section{312.00 Géométrie et trigonométrie sphérique}


\exercice{6403, potyag, 2011/10/16}
\enonce{006403}{}
Soit $\alpha:[a,b]\to S^n$ une courbe  avec $b-a < \pi.$ Montrer l'équivalence des
conditions suivantes :
\begin{enumerate}

  \item  $\alpha$ est une courbe géodésique.


  \item  Il existe deux vecteurs orthogonaux $A, B\in S^n$ tels que
$\alpha(t)= A\cos (t-a) + B \sin(t-a).$
  \item  La courbe $\alpha$ vérifie l'équation $\alpha''+\alpha =0.$
\end{enumerate}
\finenonce{006403}



\finexercice


\exercice{6404, potyag, 2011/10/16}
\enonce{006404}{}
 Soit $ S^n$ est la sphère unité dans l'espace linéaire $E$
de dimension $n+1.$
\begin{enumerate}

  \item  Montrer que la distance sphérique induit sur $ S^n$ une topologie équivalente à
celle induite de l'espace ambiant $E.$

  \item  Montrer que l'intersection d'un sous-espace  linéaire $L$ de $E$ de dimension $k$
  avec $S^n$ est une sphère de dimension $k-1$ (si $k=2$ cette intersection est
  un cercle appelée grand cercle de $ S^n$).
\end{enumerate}
\finenonce{006404}



\finexercice


\exercice{6405, potyag, 2011/10/16}
\enonce{006405}{}
Nous noterons $[X, Y]$ le produit vectoriel de deux vecteurs
$X$ et $Y$ dans $\R^3$.
Montrer que trois vecteurs $X,\ Y,\ {\rm et}\ Z$ dans
$\R^3$ sont libres ssi les vecteurs $[X, Y]$, $[Y, Z]$ et
$[Z, X]$ sont libres.


{\it Indication :} Démontrer d'abord l'identité suivante :
$$[[X, Y], Z] = <X, Z>Y - <Y, Z>X$$
\finenonce{006405}



\finexercice


\exercice{6406, potyag, 2011/10/16}
\enonce{006406}{}
\label{exo:prec}
Soit $T=\triangle ABC\subset S^2$ un triangle
sphérique d'angles intérieurs $\alpha=\angle A,\ \beta=\angle B,\
\gamma=\angle C$
\begin{enumerate}
  \item  Montrer qu'il existe un triangle $T'=\triangle
  A'B'C'$ dit polaire tel que $$a'=\pi-\alpha,\ b'=\pi-\beta,\
  c'=\pi-\gamma,$$\ où comme d'habitude on note $a',\ b',\ c'$ les longueurs
  des
  c\^otés opposés aux sommets $A',\ B',\ C'.$
{\it Indication.} Poser : 
$C'=\frac{[A, B]}{\vert\vert [A, B]\vert\vert},
A'=\frac{[B, C]}{\vert\vert [B, C]\vert\vert},
B'= \frac{[C, A]}{\vert\vert [C,A]\vert\vert}$

  \item  Montrer que $(T')' = {\rm sign} (<A, [B,C]>)\cdot T.$ En déduire que
   $$\angle A'=\pi-a,\ \angle B'=\pi-b,\
  \angle C'=\pi-c,$$
\end{enumerate}
\finenonce{006406}



\finexercice


\exercice{6407, potyag, 2011/10/16}
\enonce{006407}{}
\label{exo:apres}
En utilisant le résultat et les notations de l'exercice
\ref{exo:prec} montrer que :
\begin{enumerate}
  \item  $\ \cos\gamma=\cos c\cdot\sin\alpha\cdot\sin\beta -\cos\alpha\cdot\cos\beta$
  \item  $\ \cos\beta=\cos b\cdot\sin\alpha\cdot\sin\gamma -\cos\alpha\cdot\cos\gamma$
  \item  $\ \cos\alpha=\cos a\cdot\sin\beta\cdot\sin\gamma -\cos\beta\cdot\cos\gamma$
\end{enumerate}
\finenonce{006407}



\finexercice


\exercice{6408, potyag, 2011/10/16}
\enonce{006408}{}
 Le but de cet exercice est démontrer que pour
chaque triplet $\alpha,\ \beta,\ \gamma\in\   ]0,\frac{\pi}{2}[$ tel que $\alpha+\beta+\gamma>\pi$ il
existe un triangle sphérique d'angles intérieurs égaux à $\alpha,\ \beta,\
\gamma$.

\begin{enumerate}
  \item  Montrer que $\forall\ \alpha\in\  ]0,\frac{\pi}{2}[$\  $\forall\ d\in\  ]0,
  \alpha]$
   il existe un triangle
  sphérique $\triangle ABC$ tel que $\angle C=\frac{\pi}{2},\ \angle
  A=\alpha$ et $a=\vert BC\vert=d$ (dans les notations précédentes).

  \item  En utilisant 1. et les formules de l'exercice \ref{exo:apres}
 démontrer le résultat.
\end{enumerate}
\finenonce{006408}



\finexercice

\section{313.00 Groupe orthogonal et quaternions}


\exercice{6409, potyag, 2011/10/16}
\enonce{006409}{}
 Le but de cet exercice est d'introduire une topologie sur le
groupe linéaire $GL(E).$ Soit $A\in GL(E),$ alors on introduit la norme:
$$\displaystyle \vert\vert A\vert\vert =\sup_{\vert\vert
x\vert\vert=1}\vert\vert Ax\vert\vert\hfill\eqno(1).$$ Montrer que les
topologies suivantes sur $GL(E)$ sont équivalentes:

\begin{enumerate}
  \item  $A_n\to A$ si la suite $a_{ij}^n$ des coefficients  de $A_n$
  converge vers $a_{ij}.$

  \item  $A_n\to A$ si $\vert\vert A_n -A\vert\vert \to 0.$

 \item  $A_n\to A$ si l'application $A_n:x\to A_n(x)$ converge vers l'application $A(x)$ uniformément sur chaque compact $K\subset
  E.$
\end{enumerate}
\finenonce{006409}



\finexercice


\exercice{6410, potyag, 2011/10/16}
\enonce{006410}{}
\begin{enumerate}
  \item  En utilisant la norme (1) démontrer que $O(n)$ est compact dans
  $GL(E).$

  \item  En utilisant le résultat du cours qu'une matrice orthogonale est une
matrice en blocs démontrer que $O(n)$ contient deux composantes
 connexes :
$SO(n)=\{A\in O(n)\ \vert\ {\rm det} A=+1\}$ et $O^{-}(n)=\{A\in
O(n)\ \vert\ {\rm det} A=-1\}.$
$O^{-}(n)$ est-il un sous-groupe de $O(n)$ ?

 \end{enumerate}
\finenonce{006410}



\finexercice


\exercice{6411, potyag, 2011/10/16}
\enonce{006411}{}
Rappelons qu'une application $g\in SO(3)$ est dite \emph{retournement}
 si $g$ est une rotation d'angle $\pi$ autour d'une droite fixée 
$L\subset \Rr^3\ (g\vert_{L}\equiv {\rm id})$.
Soit $g\in SO(3)$ un retournement d'axe la droite $L\in \R^3\ (0\in L).$
Montrer que $g_1\in SO(3)$ est un retournement ssi il existe $f\in SO(3)$
tel que $g_1=fgf^{-1}$ et l'axe de $g_1$ est $f(L).$
\finenonce{006411}



\finexercice


\exercice{6412, potyag, 2011/10/16}
\enonce{006412}{}
\begin{enumerate}
  \item Soit $S=\{z+tj\ \vert\ z\in\C,\ t\in\R\}$ un sous-ensemble $\mathbb{H}$
  des quaternions. Montrer que $S$ est invariant par l'application :
$$\rho_j\ :\ \to jqj^{-1},\ q\in\mathbb{H}.$$

  \item En identifiant $S$ avec $\R^3$ décrire $\rho_j:S \to S.$

  \item Montrer que l'application $\rho_k:q\to kqk^{-1}$ laisse $S$
  invariant (i.e. $\rho_k(S)=S$). Décrire $\rho_k.$
\end{enumerate}
\finenonce{006412}



\finexercice


\exercice{6413, potyag, 2011/10/16}
\enonce{006413}{}
\begin{enumerate}
  \item Montrer que $ S^3$ est un sous-groupe de $\mathbb{H}^*$ considéré
comme groupe multiplicatif. Est-il normal ?
 
  \item Montrer que si $\exists\ q\in  S^3\ :\ \forall\ q_1\in\R^3=\{y\cdot i+u\cdot j +
v\cdot k\ \vert\ y,v,u\in\R\}\ : \ qq_1q^{-1}=q_1$ alors $q=\pm {\rm id_{\mathbb{H}}}.$
\end{enumerate}
\finenonce{006413}



\finexercice


\exercice{6414, potyag, 2011/10/16}
\enonce{006414}{}
Soit $q,r\in(\mathbb{H}\cong \R^4)$ alors montrer que $<q,r>=\frac{1}{2}({\overline
q}r+{\overline r}q)$ où $<\ ,\ >$ est le produit scalaire euclidien de $\R^4,$
et on note $\overline{q}=\overline{z}-wj$ si $q=z+wj.$
\finenonce{006414}



\finexercice


\exercice{6415, potyag, 2011/10/16}
\enonce{006415}{}
\begin{enumerate}
  \item On considère l'application $\xi_s(q)=-sqs^{-1}\ (s\in(\R^3)^*,\
  q\in\R^3).$ Montrer que  $\xi_s$ est la reflection dans
  $\R^3$  par rapport au plan $s^\bot.$

  \item Démontrer la surjectivité de l'application $\varphi :  S^3\to SO(3)$ où
$\varphi(s)=\rho_s,$ et $\rho_s(q)=sqs^{-1}$ ($s\in S^3,\ q\in\R^3$).
\end{enumerate}
\finenonce{006415}



\finexercice


\exercice{6416, potyag, 2011/10/16}
\enonce{006416}{}
Si $q\in  S^3$ tel que $q=\cos\theta + I\sin\theta,\ I\in S^2,$ alors montrer
que $\rho_q\in SO(3)$ est la rotation d'angle $2\theta$ autour de l'axe $OI$.
\finenonce{006416}



\finexercice


\exercice{7801, mourougane, 2021/08/11}
\enonce{007801}{}
Montrer sans calcul que la matrice suivante $A$ à coefficients complexes est inversible
$$A=\left[
\begin{array}{cccc}
1-i&2&1&3\\ 2&-i&-1&0\\ 1&-1&1-i&-1\\ 3&0&-1&-i
\end{array}\right]
$$
\finenonce{007801}
\finexercice


\exercice{7802, mourougane, 2021/08/11}
\enonce{007802}{}
Déterminer l'orbite et le stabilisateur d'un vecteur de norme 1 sous l'action du groupe orthogonal du produit scalaire standard dans
$\Rr^n$. Montrer que groupe $O(n-1)$ est isomorphe à un sous-groupe $O_{n-1}$ de $O(n)$. Ce sous-groupe est-il distingué ?
\finenonce{007802}
\finexercice


\exercice{7803, mourougane, 2021/08/11}
\enonce{007803}{Décomposition polaire d'un endomorphisme}
Soit $H$ un espace hermitien et $a$ un endomorphisme inversible de $H$.
 Montrer que $a$ s'écrit de fa{\c c}on unique sous la forme $a=hu$ o{ù} $h$ est un
 endomorphisme auto-adjoint positif et $u$ unitaire.
Déterminer $h$ et $u$ pour l'endomorphisme $a$ dont la matrice dans la base canonique de $\Cc^2$ est $\left[
  \begin{array}{cc}
1&i\\ i&1
  \end{array}
\right]$.
\finenonce{007803}
\finexercice


\exercice{7804, mourougane, 2021/08/11}
\enonce{007804}{Involutions}
Soit $f$ une forme bilinéaire symétrique non dégénérée sur un espace vectoriel $E$. 
\begin{enumerate}
\item
Soit $u\in GL(E)$ une involution. Montrer que $u$ est orthogonale si
et seulement si $E_+(u)=E_+:=Ker (u-Id)$ et $E_-:=Ker(u+Id)$ sont
orthogonaux.
 Montrer alors que $(E_+)^\perp = E_-$ et que $E_+$ n'est pas isotrope. 
\item Soit $F\subset E$ un sous-espace non isotrope. Montrer qu'il
 existe une unique involution orthogonale telle que $E_+(u)=F$.
\item Montrer que $O(f)\equiv SO(f)\ltimes \{-1,1\}$.
\end{enumerate}
\finenonce{007804}
\finexercice


\exercice{7805, mourougane, 2021/08/11}
\enonce{007805}{Dilatations}
Soit $f$ une forme sesquilinéaire non dégénérée. Déterminer les
dilatations orthogonales (resp. unitaires. resp. symplectiques.
\finenonce{007805}
\finexercice


\exercice{7806, mourougane, 2021/08/11}
\enonce{007806}{}
Soit $k$ un corps de caractéristique différente
de $2$.
Soit $E$ un $k$-espace vectoriel de dimension $n\geq 2$ et $\phi$ une
forme sesquilinéaire non dégénérée sur $E$ symétrique, hermitienne ou alternée. Soit
$\tau$ une transvection de $E$ donnée à l'aide d'une forme linéaire non nulle 
$f$ sur $E$ et un vecteur $a$ de $Ker f$ par $\forall x\in E, \tau
(x)=x+f(x)a$. 
\begin{enumerate}
\item On suppose désormais que $\tau$ est une isométrie relativement à $\phi$. Montrer que $a$ est isotrope.
\item Montrer que $f$ et $\phi (\cdot , a)$ sont proportionnelles. On
 notera $\lambda\in k^\star$ tel que $f=\lambda\phi (\cdot , a)$.
\item Montrer que si $\sigma\not =Id$ et $\phi$ est hermitienne ou
 symétrique, alors $\lambda +\sigma (\lambda)=0$.
\item Montrer qu'il n'existe pas de transvections orthogonales, qu'il
 existe des transvections unitaires si et seulement si l'indice est
 plus grand que $1$ et qu'il existe toujours des transvections
 symplectiques. 
\end{enumerate}
\finenonce{007806}
\finexercice 


\exercice{7807, mourougane, 2021/08/11}
\enonce{007807}{Sur les similitudes}
Soit  $f$  une forme  sesquilinéaire  non dégénérée  symétrique
(resp. hermitienne, alternée) sur un $K$-espace vectoriel $E$. On note
$GO(f)$ (resp. $GU(f)$, $GSp(f)$) le groupe des similitudes de $f$. On
note  $\mu$ l'application  qui  à une  similitude associe  son
multiplicateur dans $K^\star$.
\begin{enumerate}
\item Déterminer les similitudes de la forme symplectique standard sur
 $K^2$.
\item Montrer que $u$ est une similitude si et seulement si elle
 conserve l'orthogonalité.
\item On suppose $f$ symétrique. Montrer que $Im (\mu)=\{\lambda\in
K^\star / q\equiv \lambda q\}\supset (K^\star)^2$.
\end{enumerate}
\finenonce{007807}
\finexercice


\exercice{7808, mourougane, 2021/08/11}
\enonce{007808}{Etude du groupe orthogonal $O(1,1)$}
\begin{enumerate}
\item Montrer que dans un plan d'Artin il y a exactement deux droites
 isotropes $I$ et $J$.
\item Soit $u\in O(1,1)$. Montrer que $u$ envoie $I\cup J$ sur lui-même.
\item Soit $u\in O(1,1)$. Montrer que $u$ est directe si et seulement
 si $u$ laisse fixes chaque droite isotrope.
\item En déduire la forme des éléments de $O(1,1)$.
\end{enumerate}
\finenonce{007808}
\finexercice


\exercice{7809, mourougane, 2021/08/11}
\enonce{007809}{}
On considère la forme symplectique sur $\Rr^{2n}$ donnée par 
$$f\left(\left({x\atop y}\right),\left({x'\atop y'}\right)\right)
=x^t\ y'-y^t\ x'$$
o{ù} $x,y,x',y'$ sont dans $\Rr^n$.
\begin{enumerate}
\item En décomposant par bloc $n\times n$ une matrique quelconque $g$ de
 $M_{2n}(\Rr)$ sous la forme $\left({AB\atop CD}\right)$,
 caractériser les matrices symplectiques en termes de systèmes
 d'équation pour $A,B,C,D$.
\item Déterminer toutes les matrices symplectiques vérifiant de plus
 $B=C=0$.
\item Déterminer toutes les matrices symplectiques vérifiant
 $A=D=I_n$ et $C=0$.
\item Déterminer toutes les matrices symplectiques vérifiant $C=0$.
\item Déterminer toutes les matrices $Q$ de $M_{n}(\Rr)$ telles que
 l'espace $W_Q:=\left\{ \left({Qy\atop y}\right)/ y\in
 \Rr^n\right\}$ 
soit totalement isotrope. Quel est le lien avec les questions
 précédentes ?
\end{enumerate}
\finenonce{007809}
\finexercice


\exercice{7893, mourougane, 2021/08/11}
\enonce{007893}{Le groupe $SO(2)$}
\begin{enumerate}
\item Rappeler un isomorphisme de groupes entre $SO(2)$ et $S^1$. 
\item L'application $SO(2)\to SO(2), A\mapsto A^2$ est-elle un morphisme de groupes ? Déterminer son image et son noyau.
\end{enumerate}
\finenonce{007893}
\finexercice


\exercice{7894, mourougane, 2021/08/11}
\enonce{007894}{L'espace $V$ des matrices anti-hermitiennes de trace nulle}
\begin{enumerate}
\item Déterminer la nature de l'espace vectoriel $V$ des matrices anti-hermitiennes (i. e. $M^\star:=-M$) de $M(2,\Cc)$ de trace nulle.
\item Écrire la forme générale d'une matrice de $V$ à l'aide de trois nombres réels. En déduire une base de $V$.
\item Montrer que $\ll P,P'\gg:=-\frac{1}{2}trace(PP')$ définit un produit scalaire sur $V$.
\end{enumerate}
\finenonce{007894}
\finexercice


\exercice{7895, mourougane, 2021/08/11}
\enonce{007895}{Le groupe $SU(2)$}
\begin{enumerate}
\item Soit $P=\begin{bmatrix}a&b\\c&d\end{bmatrix}\in SU(2)$. Montrer que \( c=-\overline{b}, d=\overline{a} \quad \text{ et } \quad \overline{a}a+\overline{b}b=1\)
et écrire la forme générale d'une matrice $P$ de $SU(2)$ à l'aide de deux nombres complexes, puis de quatre nombres réels.
\item En déduire un homéomorphisme de $SU(2)$ sur la sphère unité $S^3$ de $\Cc^2$. 
(On munit ici $S^3$ de la topologie induite par celle de $\Cc^2$ et $SU(2)$ de la topologie induite par la topologie d'une norme sur l'espace vectoriel $M(2,\Cc)$.)
\item Soit $-1<c<1$. Décrire topologiquement le sous-espace de $SU(2)$ des matrices de trace $c$, appelé ``latitude $c$''. 
\item Montrer que les latitudes sont des classes de conjugaison dans $SU(2)$. 
(On pourra remarquer que les éléments de $SU(2)$ sont associés à des endomorphismes normaux (i.e. qui commutent avec leur adjoint)).
\item Quelles sont les autres classes de conjugaison ?
\item Décrire topologiquement le sous-groupe $D$ des matrices diagonales de $SU(2)$.
\end{enumerate}
\finenonce{007895}
\finexercice


\exercice{7896, mourougane, 2021/08/11}
\enonce{007896}{Les groupes $SO(3)$ et $SU(2)$}
\begin{enumerate}
\item Montrer que la classe de conjugaison $C$ de $SU(2)$ des matrices de trace nulle (i.e. la latitude $0$) 
est la sphère unité de l'espace euclidien $(V,\ll~~,\gg)$ des matrices anti-hermitienne de trace nulle.
\item Montrer que $SU(2)$ agit par conjugaison sur l'espace $V$.
\item Montrer que cette action est transitive.
\item En déduire un morphisme de groupes $\phi$ de $SU(2)$ dans le groupe orthogonal de $(V,\ll~~,\gg)$.
\item Déterminer le noyau de $\phi$.
\item En utilisant la connexité de $SU(2)$ montrer que l'image de $\phi$ est incluse dans $SO(V)$.
\item Montrer que l'image par $\phi$ du sous-groupe $D$ des matrices diagonales de $SU(2)$
est le sous-groupe des rotations de $V$ qui fixent $\begin{bmatrix}i&0\\0&-i\end{bmatrix}$.
\item En déduire l'image de $\phi$.
\end{enumerate}
\finenonce{007896}
\finexercice


\exercice{7901, mourougane, 2021/08/11}
\enonce{007901}{Représentation matricielle des nombres complexes et des quaternions}
\begin{enumerate}
 \item Déterminer dans $M(2,\Rr)$ une matrice $I$ telle que $I^2=-Id$.
 \item En déduire un morphisme non nul d'anneaux de $\Cc$ dans $M(2,\Rr)$.
 \item Soit $$I =\begin{bmatrix}1&0\\0&1\end{bmatrix} \ A=\begin{bmatrix}i&0\\0&-i\end{bmatrix} 
 \ B=\begin{bmatrix}0&1\\ -1&0\end{bmatrix} \quad \text{ et } \quad C=\begin{bmatrix}0&i\\i&0\end{bmatrix}$$
 Établir la table de multiplication de $G:=\{I, -I, A, -A, B, -B, C, -C\}$. Est-ce un groupe ? Est-il cyclique ?
 \item La sous-algèbre $H$ de $M(2,\Cc)$ engendré par $G$ est-elle commutative ? un corps gauche ? 
\end{enumerate}
\finenonce{007901}
\finexercice

\section{314.00 Géométrie projective}


\exercice{6417, potyag, 2011/10/16}
\enonce{006417}{}
Trouver la formule explicite suivante pour la projection stéréographique $\pi
: \R^n\cup\{\infty\}\to  S^n:$

$$\displaystyle \pi(x)=\Big(\frac{2x_1}{1+\vert\vert x\vert\vert^2},\ ... ,\
\frac{2x_n}{1+\vert\vert x\vert\vert^2},\ \frac{\vert\vert
x\vert\vert^2 -1}{\vert\vert x\vert\vert^2+1}\Big),$$
où $x=(x_1,...,x_n,0)\in \R^n\subset \Rr^{n+1}.$

{\it Indication:} \'Ecrire $\pi(x)-e_{n+1}=t(x-e_{n+1}),\ t\in\R.$
\finenonce{006417}



\finexercice


\exercice{6418, potyag, 2011/10/16}
\enonce{006418}{}
Soit $L$ un espace vectoriel de dimension $n+1.$
\begin{enumerate}
  \item  Montrer que si $M_i\ (i\in I)$ sont des sous-espaces vectoriels de
  $L$ alors  $\displaystyle P(\bigcap_{i\in I} M_i)=\bigcap_{i\in
  I} P(M_i).$

  \item  Soient $M_i\ (i\in \{1,...,k\}$ des sous-espaces linéaires de $L,$
  montrer que $$<P(M_1),..., P(M_k)>
 =P(M_1+...+M_k).$$

  \item  Soient $p:L^*\to P(L)=L^{*} / _{\bf\tilde{}}\ $ l'application de la
  projectivisation (qui associe à chaque $x\in L^*$ sa classe $[x]\in P(L)$)
  et $S$ un sous-ensemble de l'espace $P(L).$ Alors montrer que $<S>= P(D),$
  où $D$ est le sous-espace de $L$ engendré par $p^{-1}(S).$
 \end{enumerate}
\finenonce{006418}



\finexercice


\exercice{6419, potyag, 2011/10/16}
\enonce{006419}{}
\begin{enumerate}
  \item  Montrer que le plan projectif $P_1$
et la droite $P_2$ dans $\mathbb{P}^3$ soit se coupent en un point soit $P_2\subset
P_1.$

  \item Soient $P_i=P(M_i)$ deux sous-espaces projectifs $(i=1,2)$. Montrer
  que si $P_1\cap P_2=\emptyset$ alors la somme $M_1+M_2$ est directe.
\end{enumerate}

\finenonce{006419}



\finexercice


\exercice{6420, potyag, 2011/10/16}
\enonce{006420}{}
Montrer que la fibration de Hopf $FH:S^3\to S^2$ ( $FH^{-1}(x)$ est un  grand
cercles de $S^3\ (\forall\ x\in S^2)$) s'écrit comme ceci :

$$FH(z,z')=(2z'\overline{z}, \vert z'\vert^2-\vert z\vert^2).$$
\finenonce{006420}



\finexercice


\exercice{6421, potyag, 2011/10/16}
\enonce{006421}{}
\begin{enumerate}
  \item  Montrer que tout sous-espace projectif de dimension $k-1$ dans $\mathbb{P}^n$
  peut être recouvert par au moins $k$ cartes affines.

  \item Trouver le nombre d'éléments d'un espace projectif de dimension $n$
  sur un corps avec $q$ éléments.
\end{enumerate}
\finenonce{006421}



\finexercice


\exercice{6422, potyag, 2011/10/16}
\enonce{006422}{}
\begin{enumerate}
  \item  Montrer que  le groupe projectif $PGL(L)$ agit transitivement  sur
  l'ensemble de sous-espaces projectifs de la dimension $k$ fixée.

\item Montrer que $PGL(L)$ agit transitivement sur l'ensemble de couples
ordonnés

$\{(P_1, P_2)\ \vert\  {\rm dim} P_1= k_1, \  {\rm dim} P_2= k_2,\
{\rm dim} (P_1\cap P_2) = k_3\},$ où $k_1, k_2, k_3$ sont fixés.

\item Montrer que $PGL(L)$ agit transitivement sur l'ensemble des drapeaux
projectifs

${\cal D}=\{ (P_1,...,P_k)\ \vert P_1\subset ... \subset P_k, \}$ où la
longueur  $k$ est fixée et  $P_i$ est un sous-espace projectif de $P(L)$ de
dimension $i$ fixée $(i=1,...,k).$
\end{enumerate}
\finenonce{006422}



\finexercice


\exercice{6423, potyag, 2011/10/16}
\enonce{006423}{}
\begin{enumerate}
  \item  Montrer que toute homographie $\gamma\in PGL_n\C$ possède au moins un
  point fixe

  \item Montrer que toute homographie $\gamma:\R\mathbb{P}^{2n}\to \R\mathbb{P}^{2n}$ a toujours
au moins un point fixe.

\item Soit $\gamma\in PGL(L)$ tel que ${\rm card}({\rm fix}(\gamma)) < \infty$ et ${\rm dim} P(L)=n$ alors montrer que
${\rm card}({\rm fix}(\gamma))\leq n+1$ où ${\rm fix}(\gamma)=\{x\ \vert \gamma(x)=x\}.$

\end{enumerate}
\finenonce{006423}



\finexercice


\exercice{6424, potyag, 2011/10/16}
\enonce{006424}{}
 Cet exercice ne concerne pas directement la géométrie
projective mais sera utilisé par la suite.

\begin{enumerate}
\item  Montrer que chaque réflexion par rapport à un hyperplan dans
$\R^n$
  est une application conforme dans  $\R^n.$ En déduire que chaque
  isométrie euclidienne et chaque isométrie sphérique sont
  conformes dans $\R^n$ et sur $ S^n$ respectivement.

  \item Montrer qu'une application linéaire $A:E\mapsto E$
  conserve les angles non-orientés  entre les vecteurs non-nuls ssi $A$
  est une matrice conforme.

  \item  Montrer qu'une application $f:D\mapsto \R^n$ d'un
  ouvert $D\subset \R^n$ est conforme dans $D$ ssi elle conserve les
  angles dans $D$.
\end{enumerate}
\finenonce{006424}



\finexercice


\exercice{6425, potyag, 2011/10/16}
\enonce{006425}{}
 Démontrer que toute application de M\"obius $\gamma\in M(2)$
est conforme sur $\overline\C.$
\finenonce{006425}



\finexercice


\exercice{6426, potyag, 2011/10/16}
\enonce{006426}{}
 Rappelons qu'un cercle généralisé est soit un cercle
euclidien $\Sigma(z_o,r)=\{z\in \C\ \vert\ \vert z-z_0\vert=r\}$ soit une
droite à laquelle on ajoute le point $\{\infty\}$ (à l'aide la projection
stéréographique). On note $M=\{\frac{az+b}{cz+d}\ \vert\ a, b, c, d\in\C,\
ad-bc\not=0\}.$

\begin{enumerate}

 \item Montrer que le groupe
  $PGL_2\C$ agit trois fois transitivement sur $\overline{\Cc}.$

\item Vérifier que chaque cercle généralisé dans $\overline{\Cc}$ s'écrit sous la
  forme :
  $$A(x^2+y^2)+Bx+Cy+D=0,\ B^2+C^2 > 4 A D$$

  \item Soit $C_1\in \C$ un cercle généralisé, alors montrer que un
  sous-espace $C_2\subset \C$ est un cercle généralisé ssi il existe
  $\gamma\in M$ telle que $\gamma(C_1)=C_2.$

  \item Soit $K\subset \C$ un cercle généralisé et $f\in M(2)$ tel que $\displaystyle
f\vert_K\equiv {\rm id_K}$ alors montrer que soit $f\equiv {\rm id}$ soit $f$
est la réflexion par rapport à $K.$
\end{enumerate}
\finenonce{006426}



\finexercice


\exercice{6427, potyag, 2011/10/16}
\enonce{006427}{}
 Montrer que $$\displaystyle M(2)=\Big\{\frac{a_2 \overline
z+b_2}{c_2 \overline z+d_2}\ ;\ \frac{a_1 z+b_1}{c_1 z+d_1}\ \Big\vert\ a_i,
b_i, c_i, d_i \in\C\ ;\ a_id_i-b_ic_i\not=0\Big\}$$

et en déduire que $\vert M(2) : M\vert=2$ où $M=M_+(2)$ est le groupe  des
transformations de M\"obius paires.
\finenonce{006427}



\finexercice


\exercice{6428, potyag, 2011/10/16}
\enonce{006428}{}
Soit $\tau_{\Sigma}$ la réflexion par rapport au cercle
euclidien $\Sigma=\{z\in\C\ \vert \vert z-z_0\vert=r\}$ alors montrer que

$$\displaystyle \vert \tau_{\Sigma}(z)-\tau_{\Sigma}(w)\vert =  r^2\frac{\vert z-w\vert}
{\vert z -z_0\vert\vert w-z_0\vert}$$
\finenonce{006428}



\finexercice


\exercice{6429, potyag, 2011/10/16}
\enonce{006429}{}
\begin{enumerate}

\item Montrer que chaque application $g\in M$ possède soit un point fixe
dans $\overline{\Cc}$ soit deux points fixes. Cette affirmation reste-t-elle vraie pour
les éléménts de $M(2)$ ?

\item Notons ${\rm fix}(g)$ l'ensemble $\{x\in\overline{\Cc}\ \vert\ g(x)=x\}$ des points fixes
de $g.$ Montrer que si $\gamma=fgf^{-1}$ alors ${\rm fix} (\gamma)=f({\rm fix} (g).$

\item Soient $C_i\ (i=1,2)$ deux cercles généralisés. Montrer que
  $\exists\ \gamma\in M\ :\ \tau_{C_1}=\gamma\tau_{C_2}\gamma^{-1},$ où $\tau_{C_i}$ désigne
  l'inversion par rapport à $C_i.$
\end{enumerate}
\finenonce{006429}



\finexercice


\exercice{6430, potyag, 2011/10/16}
\enonce{006430}{}
Soient $C_i\ (i=1,2)$ deux cercles généralisés. Montrer que $\tau_{C_1}$ et
$\tau_{C_2}$ commutent si et seulement si le cercle $C_1$ est orthogonal à $C_2$
(e.g. si $C_i$ sont deux cercles euclidiens alors ils sont orthogonaux ssi les
angles entre deux rayons aux points de l'intersection $C_1$ et $C_2$  sont
égaux à $\frac{\pi}{2}$).
\finenonce{006430}



\finexercice


\exercice{7733, mourougane, 2021/08/11}
\enonce{007733}{Homographie}
Soit $d$ et $d'$ deux droites d'un plan projectif.
Soit $h$ une homographie de $d$ dans $d'$ telle que $A'=h(A)$, $B'=h(B)$, $C'=h(C)$. Construire $h(M)$.

\includegraphics[scale=0.5]{images/img-mour-486}

\finenonce{007733}
\finexercice


\exercice{7738, mourougane, 2021/08/11}
\enonce{007738}{Questions de cours}
\begin{enumerate}
 \item L'ensemble des permutations de profil $(\cdot, \cdot)(\cdot, \cdot)$ avec l'identité est-il un sous-groupe distingué de $\mathcal{A}_6$.

 \item Décrire les différentes possibilités pour la dimension de l'intersection de deux plans projectifs de $\mathbb{P}^3$.
 Décrire les différentes possibilités pour la dimension de l'intersection de deux plans projectifs de $\mathbb{P}^4$.
 
 \item Donner l'exemple de deux quintuplets de points deux à deux distincts d'une droite projective qui ne peuvent pas être l'image l'un de l'autre par une homographie.

\end{enumerate}
\finenonce{007738}
\finexercice


\exercice{7742, mourougane, 2021/08/11}
\enonce{007742}{Une homographie}
Soit $9$ points $A,B,\ldots, I$ du plan $\R^2$ tels que $ABED$,
$BCFE$, $DEHG$ et $EFIH$ soient des carrés.
Etant données les images $A'=h(A)$, $B'=h(B)$, $E'=h(E)$ et $D'=h(D)$ par une homographie de $P^2(\R)$ dans lui-même, construire à la règle les images des autres points.
\finenonce{007742}
\finexercice


\exercice{7743, mourougane, 2021/08/11}
\enonce{007743}{}
~

\includegraphics[scale=0.5]{images/img-mour-495}

\begin{enumerate}
    \item Dessiner la configuration correspondante dans le plan affine $P(V)-(AM)$.
    \item Choisir des coordonnées homogènes telles que l’on ait $A = (1 : 0 : 0), B = (0 : 1 : 0), L = (0 : 0 : 1)
    $ et $ C = (1 : 1 : 0)$. Déterminer les équations des droites projectives $(BL), (DM), (NA), (LE), (MP)$ et les
    coordonnées homogènes des points $N = (BL) \cap (DM), P = (NA) \cap (LE)$ et $F = (AB) \cap (MP)$ en termes
    des coordonnées homogènes des points $D,E$ et $M$.
    \item Exprimer le birapport $[A,B,C,F]$ en termes de $x = [A,B,C,D]$ et $y = [A,B,C,E]$.
\end{enumerate}
\finenonce{007743}
\finexercice


\exercice{7744, mourougane, 2021/08/11}
\enonce{007744}{Homographie plane}
\begin{enumerate}
    \item Etant données les images $A'=h(A)$, $B'=h(B)$, $C'=h(C)$, et $D'=h(D)$ par une homographie $h$ de $P^2(\Rr)$ dans lui-même, construire à la règle les images des autres points.
    On indiquera l'ordre dans lequel les constructions sont effectuées.
    
    \item Suffirait-il de connaître les images $A'=h(A)$, $B'=h(B)$ et $C'=h(C)$ ?
\end{enumerate}
\finenonce{007744}
\finexercice


\exercice{7745, mourougane, 2021/08/11}
\enonce{007745}{}
Soit $F$ une homographie du plan projectif $P^2$ qui admet une droite $d$ de points fixes.
\begin{enumerate}
    \item Montrer qu'on peut choisir $f\in GL(3,k)$ telle que $F=P(f)$ et $f$ admet un plan de points fixes. 
    
    \item Montrer alors qu'il existe un point $O$ de $P^2$ (appelé centre de $F$) tel que pour tout point $M$ de $P^2$ non fixé par $F$, la droite $(MF(M))$ passe par $O$.
    
    \item Soit $d$ une droite et $O$ un point hors de $d$. Soit $A$ et $A'$ deux points hors de $d$ et $A\not=O$ et $O,A,A'$ alignés. Montrer en choisissant un repère convenable qu'il existe une unique homographie $F$ telle que $d$ soit une droite de points fixes et $O$ le centre et qui envoie $A$ sur $A'$.
    
    \item Soit $F$ une homographie du plan projectif $P^2$ qui admet une droite $d$ de points fixes et de centre $O$. Sachant que $F(A)=A'$, construire l'image du point $M$ par $F$ dans les cas suivants.
    
    \begin{figure}[!ht] 
        \begin{center}\includegraphics[scale=0.5]{images/img-mour-497-1}\end{center}
        \caption{}
    \end{figure}
    
    \begin{figure}[!ht] 
        \begin{center}\includegraphics[scale=0.5]{images/img-mour-497-2}
            \caption{le point $O$ est à l'infini}\end{center}
    \end{figure}
    
    \begin{figure}[!ht] 
        \begin{center}\includegraphics[scale=0.4]{images/img-mour-497-3}
            \caption{le point $0$ est à l'infini dans la direction de $d$}\end{center}
    \end{figure}
    
    
    \item Soit $H$ une involution. On considère deux points $P$ et $Q$ tels que avec leur image $P'$ et $Q'$ ils forment un repère projectif (aucun triplet n'est formé de points alignés).
    On définit $O:=(PP')\cap (QQ')$ et $d$ la droite reliant $(PQ')\cap(QP')$ et $(PQ)\cap(P'Q')$. Montrer que $H$ est l'homographie de droite fixe $d$ de centre $O$ qui envoie $P$ sur $P'$. 
    
\end{enumerate}
\finenonce{007745}
\finexercice


\exercice{7746, mourougane, 2021/08/11}
\enonce{007746}{Géométrie projective}
Soit $\Delta = P(E)$ une droite projective. Soient $F = P(f)$ et $F' = P(f')$ deux homographies de $\Delta$ dans elle-même telles que $F^2\neq Id_\Delta$, $F'^2\neq Id_\Delta$ et qui possèdent chacune exactement deux points fixes distincts.\newline On se propose de montrer que $F$ et $F'$ commutent si et seulement elles ont les mêmes points fixes. On note $A$ et $B$ les points fixes de $F$ et on note $A'$ et $B'$ les points fixes de $F'$. 
\begin{enumerate}
 \item On suppose que $F$ et $F'$ ont les mêmes points fixes.
Comment traduire cette hypothèse à l'aide des applications linéaires associées $f$ et $f'$ ?
 Montrer que $F$ et $F'$ commutent (on pourra considérer un repère projectif de $\Delta$).

Dans la suite de l'exercice, on montre l'implication réciproque : on suppose donc que $F$ et $F'$ commutent.

 \item Rappeler la démonstration du fait qu'une homographie d'une droite projective dans elle-même possédant trois points fixes deux à deux distincts est l'identité.

 \item En considérant l'image par $F \circ F'$ des points $A,B,A',B'$ montrer que $\{F'(A),F'(B)\} = \{ A,B\}$ et que $\{F(A'),F(B')\} = \{A',B'\}$.

 \item Supposons que $F(A') = A'$ et $F(B') = B'$, montrer que $\{A',B'\} = \{A,B\}$.

 \item  Supposons que $F(A') = B'$ et $F(B') = A'$, montrer que $\{A',B'\} = \{A,B\}$, et en déduire que ce second cas ne peut pas se produire.

 \item Conclure.
\end{enumerate}
\finenonce{007746}
\finexercice


\exercice{7747, mourougane, 2021/08/11}
\enonce{007747}{Le plan $P(\mathbb{F}^3_2)$}
\begin{enumerate}
 \item Déterminer le nombre de points et de droites du plan projectif $P(\mathbb{F}^3_2)$.
Représenter les relations d'incidence.
\item
 Déterminer le nombre de points de l'espace projectif $\mathbb{P}^n_{\mathbb{F}_q}$ de dimension $n$ sur le corps $\mathbb{F}_q$.
\end{enumerate}
\finenonce{007747}
\finexercice


\exercice{7748, mourougane, 2021/08/11}
\enonce{007748}{}
On rappelle que le nombre de racine $n$ième de l'unité dans $\mathbb{F}_q$ est $pgcd(n,q-1)$.
En considérant, le morphisme $SL(E)\to \mathcal{S}(P(E))$ associé à l'action de $SL(E)$ sur les droites de $E$, démontrer l'existence des isomorphismes suivants
\begin{itemize}
 \item $GL(2,\mathbb{F}_2)=SL(2,\mathbb{F}_2)=PSL(2,\mathbb{F}_2)=PGL(2,\mathbb{F}_2)=\mathcal{S}_3$
\item $PGL(2,\mathbb{F}_3)=\mathcal{S}_4$ et $PSL(2,\mathbb{F}_3)=\mathcal{A}_4$.
\item $PGL(2,\mathbb{F}_4)=PSL(2,\mathbb{F}_4)=\mathcal{A}_5$.
\end{itemize}
\finenonce{007748}
\finexercice


\exercice{7749, mourougane, 2021/08/11}
\enonce{007749}{}
Le but de l'exercice est de 
démontrer le théorème de Pappus affine~:
\textit{Soit $d$ et $d'$ deux droites d'un plan affine $E$. Soit $A, B, C$ (resp. $A', B', C'$) trois points sur $d$ (resp. sur $d'$.)
Si les droites $(AB')$ et $(BA')$ sont parallèles ainsi que les droites $(BC')$ et $(CB')$, alors les droites $(CA')$ et $(AC')$ le sont aussi.}

Dans le cas où $d$ et $d'$ sont sécantes en $I$
\begin{enumerate}
 \item On considère l'homothétie $h$ de centre $I$ qui envoie $A$ sur $B$. Déterminer l'image de $B'$ par $h$. 
\item On considère l'homothétie $H$ de centre $I$ qui envoie $B$ sur $C$. Déterminer l'image de $C'$ par $H$.
\item Déterminer l'image de $A$ et celle de $C'$ par $H\circ h$.
\item Conclure.
\end{enumerate}

Comment raisonner dans le cas où $d$ et $d'$ sont parallèles ?
\finenonce{007749}
\finexercice


\exercice{7750, mourougane, 2021/08/11}
\enonce{007750}{Avec un repère projectif}
Dans un plan projectif réel, on considère le repère projectif $(A, B, C; I)$.
Soit $A', B', C'$ respectivement sur $(BC), (CA)$ et $(AB)$ tels que
$(AA'), (BB')$ et $(CC')$ soient concourantes en $I$.
Montrer analytiquement que les points $P:=(BC)\cap (B'C')$,
$Q:=(CA)\cap(C'A')$ et $R:=(AB)\cap(A'B')$ sont alignés.
\finenonce{007750}
\finexercice


\exercice{7751, mourougane, 2021/08/11}
\enonce{007751}{}
Soit $F=P(f)$ une homographie d'une droite projective dans elle-même. \`A quoi correspondent en terme de $f$ les points fixes de $F$ ?
Montrer que si $F$ admet trois points fixes deux à deux distincts, $F$ est l'identité.
\finenonce{007751}
\finexercice


\exercice{7752, mourougane, 2021/08/11}
\enonce{007752}{Perspective}
Soit $\vec{V}$ un espace vectoriel et $P:=P(\vec{V})$.
\begin{enumerate}
\item Soit $F$ et $G$ deux sous-espaces projectifs disjoints de $P$.
Montrer qu'il existe un unique sous-espace projectif $<F,G>$ de $P$ de
dimension $\dim F+\dim G+1$ contenant $F$ et $G$.
\item Soit $F$ et $G$ deux sous-espaces projectifs disjoints de $P$
 tels que $\dim F+\dim G= \dim P-1 $. Quel est le domaine de
 définition de l'application (appelée perspective) ?
$$\begin{array}{ccccc}
P(\vec{V})&\to& G &\subset& P(\vec{V})\\
M&\mapsto & G\cap <M,F>&=&P (\vec{G}\cap (\vec{M}\oplus\vec{F}))
\end{array}$$
Montrer que c'est une application projective.
\end{enumerate}
\finenonce{007752}
\finexercice


\exercice{7753, mourougane, 2021/08/11}
\enonce{007753}{Quadrique}
Soit $\vec{V}$ un espace vectoriel et $q$ une forme quadratique sur
$\vec{V}$. On appelle quadrique projective associée à $q$ le sous
ensemble de $P:=P(\vec{V})$ défini par ($p$ est la projection
canonique $\vec{V}-\{0\}\to P$)
$$Q:=p\left( \{x\in\vec{V}-\{0\}/ q(x)=0\}\right).$$
\begin{enumerate}
\item On suppose $\dim P=1$. Montrer que si $Q$ contient trois points
 distincts, $Q=P$.
\item On suppose $\dim P=2$. Montrer que si $Q$ contient une droite $d$,
 soit $Q=P$, soit il existe une droite $d'$ telle que $Q=d\cup d'$.
\item Soit $d$ une droite de $P$. Montrer que si $d$ rencontre $Q$ en
 au moins trois points, $d$ est incluse dans $Q$. Montrer que si
 $k=\mathbb{C}$ alors $d$ rencontre $Q$. 
\end{enumerate}
\finenonce{007753}
\finexercice


\exercice{7754, mourougane, 2021/08/11}
\enonce{007754}{Coordonnées}
\begin{enumerate}
\item Soit $\vec{V}$ un espace vectoriel de dimension $n+1$ muni d'une
 base $\mathcal{B}:=(\vec{v_i})$ et $(x_i)$ les coordonnées cartésiennes
 associées. Soit $\mathcal{R}$ le repère projectif associé de
 $P(\vec{V})$. Soit $\vec{v}$ un vecteur de $\vec{V}$.
 Déterminer un système de coordonnées homogènes dans $\mathcal{R}$ pour
 $vect(\vec{v})$ en fonction des coordonnées cartésiennes de
 $\vec{v}$ dans $\mathcal{B}$. 
\item Soit $E$ un espace affine de dimension $n$ et $\widehat{E}$ son
 complété vectoriel. On identifie $E$ à un ouvert affine de
 $P(\widehat{E})$ par l'application naturelle $M\mapsto vect\left(
 (\!( 1,M )\!)\right)$. Soit 
$\mathcal{A}:=(A_i)_{0\leq i\leq n}$ un repère affine de
 $E$.
 \begin{enumerate}
 \item Considérons d'abord la base 
$\mathcal{B}_1:=\left( (\!( 1,A_i)\!)\right) $ de
$\widehat{E}$ et $\mathcal{R}_1$ le repère projectif associé. 
Déterminer un sytème de coordonnées homogènes dans
$\mathcal{R}_1$ de $M\in
E$ considéré dans $P(\widehat{E})$ en fonction de ses coordonnées
barycentriques dans $\mathcal{A}$. Donner une équation de l'hyperplan
à l'infini.
\item Considérons maintenant la base 
$\mathcal{B}_2:=\left( (\!( 1,A_0)\!) , (\!( 0,\vec{A_0A_i})\!)\right) $ de
$\widehat{E}$ et $\mathcal{R}_2$ le repère projectif associé. 
Déterminer un sytème de coordonnées homogènes dans
$\mathcal{R}_2$ de $M\in
E$ considéré dans $P(\widehat{E})$ en fonction de ses coordonnées
cartésiennes dans $\mathcal{A}$. Donner une équation de l'hyperplan
à l'infini.
 \end{enumerate}
\end{enumerate}
\finenonce{007754}
\finexercice


\exercice{7755, mourougane, 2021/08/11}
\enonce{007755}{Cas particulier d'un grand théorème}
Soit $E$ un plan affine muni d'un repère affine
$\mathcal{A}':=(A_0,A_1, A_3)$ et $\mathcal{C}$ la conique d'équation
cartésienne $x^2+y^2=1$. Soit $\mathcal{A}:=
(A_1,A_2=s_{A_0}(A_1),A_3)$ un nouveau repère affine de $E$.
\begin{enumerate}
\item Déterminer une équation barycentrique homogène dans $\mathcal{A}$ de
 $\mathcal{C}$. 
\item Soit $B_1, B_2, B_3$ trois points de $\mathcal{C}$ 
distincts de $A_1, A_2, A_3$. Montrer à l'aide d'un calcul en
coordonnées barycentriques que les points d'intersection 
$P=(A_1B_2)\cap (A_2B_1)$, $Q=(A_2B_3)\cap (A_3B_2)$ et
$R=(A_3B_1)\cap (A_1B_3)$ sont alignés.
\end{enumerate}
\finenonce{007755}
\finexercice


\exercice{7756, mourougane, 2021/08/11}
\enonce{007756}{}
\begin{enumerate}
    \item Soit $6$ points $A,B,\ldots, F$ du plan $\R^2$ tels que $ABCDEF$ soit un hexagone régulier,
    Etant données les images $A'=h(A)$, $B'=h(B)$, $D'=h(D)$ et $E'=h(E)$ par une homographie $h$ de $P^2(\R)$ dans lui-même, construire à la règle les images des autres points.
    \includegraphics[scale=0.5]{images/img-mour-508-1}
    \item Même question en supposant données cette fois, les images $A'=h(A)$, $B'=h(B)$, $C'=h(C)$ et $D'=h(D)$. 
\end{enumerate}
\finenonce{007756}
\finexercice


\exercice{7757, mourougane, 2021/08/11}
\enonce{007757}{}
\begin{enumerate}
    \item Soit $d$ et $d'$ deux droites du plan projectif $P^2(\R)$ et $O$ un point hors de $d\cup d'$ (figure~1).
    Construire l'axe de la projection de $d$ sur $d'$ depuis $O$. 
    
    \item Deux droites se coupent en dehors de la feuille en un point $I$.
    Soit $A$ un point de la feuille. Construire la droite $(AI)$. (voir figure 2)
    
    \includegraphics[scale=0.5]{images/img-mour-509}
    %\textit{Justifier votre construction par des arguments de géométrie projective}
\end{enumerate}
\finenonce{007757}
\finexercice


\exercice{7799, mourougane, 2021/08/11}
\enonce{007799}{Droites et quadriques}
Une quadrique d'un espace projectif $P(V)$ est le lieu des zéros d'une forme quadratique $f$ sur $V$. 
\begin{enumerate}
 \item Montrer que toute quadrique qui contient trois points distincts d'une droite $d$ contient toute la droite $d$.
\item Déterminer la dimension de l'espace des quadriques de $P^3(K)$.
\item Soit $d_1, d_2, d_3$ trois droites de $P^3(K)$. Montrer qu'il existe
une quadrique qui les contient.
\end{enumerate}
\finenonce{007799}
\finexercice


\exercice{7800, mourougane, 2021/08/11}
\enonce{007800}{Polarité}
Soit $\mathcal{Q}$  une quadrique de $P(E)$  (muni d'un repère
projectif) d'équation $q(x)=0$ o{ù} $q$ est la forme quadratique d'une
forme bilinéaire symétrique non dégénérée $f$ sur un espace vectoriel
$E$. Si $\vec{A}$ et $\vec{B}$ sont deux sous-espaces orthogonaux dans
$E$, on note $P(\vec{A})\perp P(\vec{B})$. On appelle hyperplan
polaire d'un point $A=P(\vec{A})$ de $P(E)$ l'hyperplan projectif
$A^\perp:=P(\vec{A}^\perp)$.

\begin{enumerate}
\item On munit de plan projectif d'un repère. Déterminer une équation
de la droite polaire du point $M(x_0,y_0,1)$ par rapport à la
quadrique d'équation $x^2+y^2-z^2=0$. La représenter dans l'espace
affine d'équation $z\not=0$.
\item Soit $F$ un sous-espace non-isotrope de $E$. Soit $A$ et $B$
 deux points de $P(F)$. Montrer si $A\perp B$ pour $f$ alors $A\perp
 B$ pour $f_{|F}$.
\item Soit $\mathcal{Q}$ une quadrique de $P^1(K)$ dont l'image est
 composée des deux points $A$ et $B$. Montrer en utilisant un bon
 repère que pour tout $M$ in $P^1(K)$,
$$M\perp N\iff M \textrm{ et } N \textrm{ sont conjugués harmoniques
 par rapport à } M \textrm{ et } N.$$
\item En déduire une construction géométrique de la polaire d'un point
 par rapport à une conique.
\end{enumerate}
\finenonce{007800}
\finexercice

\section{315.00 Géométrie et trigonométrie hyperbolique}


\exercice{6431, potyag, 2011/10/16}
\enonce{006431}{}
{\it On notera $\mathbb{H}^2$  le plan de Poincaré dans
l'un de deux modèles du disque ou du demi-plan, muni de la
distance hyperbolique $\rho$.}
\begin{enumerate}
  \item Montrer que $\mathbb{H}^2$ est un espace métrique complet
  mais pas compact.

    \item Dans le modèle du demi-plan on suppose que  $z, w$
  sont deux points distincts dans $\mathbb{H}^2$, montrer que
  $\rho(z,w)=\vert\ln([z^*, z, w, w^*])\vert$, où
$[z^*, z, w, w^*]$ désigne le birapport de quatre points, où
$z^*, w^*$ sont les extrémités de la géodésique passant
par $z$ et $w$ dans l'ordre indiqué sur le Figure \ref{fig:pot1} :
\end{enumerate}
\begin{figure}[tbh]
\begin{center}
  \includegraphics[scale=0.5]{../images/img006431-1}
\end{center}
\caption{\label{fig:pot1}Une géodésique}
\end{figure}
\finenonce{006431}



\finexercice


\exercice{6432, potyag, 2011/10/16}
\enonce{006432}{}
Soit $\triangle abc$ un triangle dans $\mathbb{H}^2$ (c-à-d le
sous-ensemble de $\mathbb{H}^2$ bordé par trois géodésiques dont les
points de l'intersection sont $a, b, c$) d'angles intérieurs
$\alpha, \beta, \gamma$. On suppose que $\gamma=\frac{\pi}{2}$ ; en utilisant
les notations indiquées sur le Figure \ref{fig:pot2} démontrer les
identités suivantes :

\begin{enumerate}
  \item $ \cosh c = \cosh a \cdot \cosh b$

  \item $\tanh b = \sinh a\cdot \tan\beta$

 \item $\cosh a\cdot\sin\beta = \cos\alpha$
\end{enumerate}

\begin{figure}[tbh]
\begin{center}
  \includegraphics[scale=0.5]{../images/img006432-1}
\end{center}
\caption{\label{fig:pot2}Triangles géodésiques} 
\end{figure}
\finenonce{006432}



\finexercice


\exercice{6433, potyag, 2011/10/16}
\enonce{006433}{}
Soient $\alpha, \beta, \gamma$ trois nombres positifs tels que $0\leq
\alpha,\beta,\gamma < \pi ; \ \alpha + \beta + \gamma < \pi$ alors montrer qu'il
existe un triangle hyperbolique d'angles intérieurs $\alpha, \beta,
\gamma$.
\finenonce{006433}



\finexercice


\exercice{6434, potyag, 2011/10/16}
\enonce{006434}{}
En utilisant le théorème du cours  que la somme d'angles intérieurs d'un
polygone convexe à $n$ sommets dans $\mathbb{H}^2$ est inférieure à
$(n-2)\pi$ démontrer que :
\begin{enumerate}
  \item  Soient $\theta_1, \theta_2, ..., \theta_n$ une collection
  ordonnée de nombres tels que $0\leq\theta_i < \pi\ (i=1,...,n)$. Pour
  qu'il existe un polygone convexe d'angles intérieurs $\theta_i$ il
  faut et il suffit que $\theta_1 +\theta_2 +...+ \theta_n < (n-2)\pi$.

  \item  Un polygone à $n$ c\^otés et d'angles droits existe
  ssi $n\geq 5$.
\end{enumerate}
\finenonce{006434}



\finexercice


\exercice{6435, potyag, 2011/10/16}
\enonce{006435}{}
\begin{enumerate}
  \item  Démontrer que chaque homographie non-triviale  $g\in M^+(2)\setminus\{{\rm id}\}$ possède soit un point fixe soit deux.
   Un élément $g\in M^+(2)$ est dit {\it parabolique} si son ensemble
   des points fixes : ${\rm fix}(g)=\{x\in\overline{\C}\ \vert\ f(x)=x\}$
   est un singleton. Montrer qu'un élément est parabolique {ssi}
   il est conjugué dans $M^+(2)$ à la translation $z\mapsto z+1$.


  \item  Un élément $f\in M^+(2)$ est dit {\it elliptique}
  s'il est conjugué dans $M^+(2)$ à une rotation $z\mapsto k\cdot z$
  où $\vert k\vert = 1$ , $k\in\C\setminus \{1\}$.

  Un élément $g\in M^+(2)$ est dit {\it loxodromique} s'il est conjugué
  dans $M^+(2)$ à $z\mapsto k\cdot z$ où $\vert k\vert\not=1$.
  De plus, un élément loxodromique est dit {\it hyperbolique} si
  $k\in\R^*_+\setminus\{1\}$ ; un élément loxodromique est dit
  srtictement loxodromique si $k=\lambda\cdot e^{i\theta},\ \lambda > 0\ ,\
  \theta\not=2\pi m$.
  Pour un élément $g\in M^+(2)$ on utilise la même lettre pour une de deux
  matrices dans $SL_2\C$ qui le représentent (en fait c'est $g$ ou $-g$).
  Montrer que pour tout élément $g\in M^+(2)$ l'une des
  possibilités suivantes peut avoir lieu :


\begin{enumerate}
  \item  $g$ est parabolique ssi ${\rm tr}^2(g)=4$, où $tr^2$
  est la trace carré de la matrice $g$.
  \item  $g$ est elliptique ssi ${\rm
  tr}^2(g)\in [0,4[.$
  \item  $g$ est hyperbolique ssi ${\rm tr}^2(g)\in ]4, \infty[.$

\item  $g$ est strictement loxodromique {\bf ssi} ${\rm
tr}^2(g)\not\in[0,\infty[$.
\end{enumerate}



\item En utilisant le modèle du demi-plan démontrer que les
points fixes d'un élément parabolique ou hyperbolique se
trouvent toujours sur le bord $\overline{\R}$ de $\mathbb{H}^2$.

En utilisant le modèle du disque démontrer  qu'un élément
elliptique a exactement  un des deux  points fixes dans $\mathbb{H}^2$.

\end{enumerate}
\finenonce{006435}



\finexercice


\exercice{6436, potyag, 2011/10/16}
\enonce{006436}{}
Soit $g\in M$ une homographie. Montrer que
\begin{enumerate}
  \item Si $g$ est parabolique ${\rm fix}(g)=\{x\}$ alors
  $$\displaystyle\forall z\ \in \overline{\Cc}\ :\ \lim_{n\to \pm\infty} g^n(z)=x.$$

  \item Si $g$ est loxodromique et    ${\rm fix}(g)=\{x, y\}$ alors
 $$\displaystyle\forall z\ \in \overline{\Cc}\setminus\{y\}\  :\  \lim_{n\to +\infty} g^n(z)=x,$$
et le point $x$ est dit {\it point fixe attractif} de $g.$

 $$\displaystyle\forall z\ \in \overline{\Cc}\setminus\{x\}\ :\  \lim_{n\to +\infty} g^{-n}(z)=y,$$
et le point $y$ est dit {\it point fixe répulsif} de $g.$
\end{enumerate}
\finenonce{006436}



\finexercice


\exercice{6437, potyag, 2011/10/16}
\enonce{006437}{}
 Montrer que si un élément $g\in M$ n'est pas strictement
loxodromique conjugué à

$z\to ke^{i\theta}\ (\theta\not= \pi+\pi m,\ k>0)$ alors il existe deux
familles ${\cal C}_i\ (i=1,2)$ de cercles généralisés vérifiant les conditions
suivantes :

\begin{enumerate}
  \item $\forall\ C\in {\cal C}_1\ :\ g(C)=C$

\item $\forall\ C\in {\cal C}_2\ :\ g(C)= {\cal C}_2\setminus \{C\}$



\item $\forall\ C_1\in {\cal C}_1,\ \forall\ C_2\in {\cal C}_2\ :\ C_2\ \bot\ C_1.$
\end{enumerate}

De plus, si $g\in M$ est  strictement loxodromique conjugué à $z\to
ke^{i\theta}\ (\theta\not= \pi+\pi m,\ k>0)$ alors montrer que $g$ n'a pas de
cercle invariant.
\finenonce{006437}



\finexercice

\section{316.00 Autre}

\section{320.00 Groupe}


\exercice{6369, potyag, 2011/10/16}
\enonce{006369}{}
Soit $G$ un groupe et $S$ une partie de $G$.


\begin{enumerate}
\item Montrer que $H:= \{ a_{i_1}^{\epsilon_1}...
a_{i_n}^{\epsilon _n}, a_i \in S, \epsilon_i \in \{-1,+1\} \}$
 est le sous-groupe engendré par $S$ (i.e le plus petit sous-groupe de $G$
contenant $S$), noté $<S>$.

\item Soit $A$ une partie de $G$, on appelle {\it centralisateur} de $A$,
l'ensemble : $C_A:= \{g\in G : \forall a \in A \ ga=ag\}$.
  \begin{enumerate}
  \item Montrer que $C_A$ est un sous-groupe de $G$.
  \item Montrer que $C_A= C_{<A>}$
  \end{enumerate}


\item Donner une condition nécessaire et suffisante sur $S$ pour que $<S>$
soit abélien, $<S>$ soit  normal dans $G$.
\end{enumerate}
\finenonce{006369}



\finexercice


\exercice{6370, potyag, 2011/10/16}
\enonce{006370}{}
Soit $G$ un groupe et $A$, $B$ deux sous-groupes   de $G$, on note $AB:= \{
g=ab : a\in A, b\in B\}$.
\begin{enumerate}
\item Montrer que $AB$ est un sous-groupe de $G$ si et seulement si $AB=BA$.

\item Montrer que si $AB$ est un sous-groupe de $G$ alors $AB= <A,B>$.
\end{enumerate}
\finenonce{006370}



\finexercice


\exercice{6371, potyag, 2011/10/16}
\enonce{006371}{}
Dans $GL(2,\Rr)$ : le groupe des matrices $(2,2)$
inversibles à coefficients réels.
\begin{enumerate}
\item
  \begin{enumerate}
  \item Montrer que $H:=\left(  
\begin{smallmatrix} 
1  & p \cr 
0  & 1 \cr
\end{smallmatrix}\right), 
p\in \Zz $ est un sous-groupe abélien.


  \item Montrer qu'il est cyclique, est-il normal ?
  \end{enumerate}

\item Soient  $ A:= 
\begin{pmatrix} 
0  & -1 \cr 
1 & 1 \cr
\end{pmatrix}
$ et 
 $B:= \begin{pmatrix} 
0  & 1 \cr 
-1 \ & 0 \cr
\end{pmatrix}$ deux matrices.
  
  \begin{enumerate}
  \item  Montrer que $A$ et $B$ appartiennent à $SL(2,\Zz)$, calculer
  leur ordre et montrer que $H$ est contenu dans $<A,B>$.

  Que pensez-vous des  assertions suivantes ?
  \begin{itemize}
  \item ``un groupe engendré par des éléments d'ordre fini est fini."
  \item ``tous les éléments d'un groupe
  engendré par des éléments d'ordre fini sont d'ordre fini."
  \end{itemize}

  \item Le groupe engendré par $A$ et $B$  est-il abélien ?

  \item Calculer l'intersection  du groupe cyclique engendré par $A$ et
  du groupe cyclique engendré par $B$.
  \end{enumerate}
\end{enumerate}
\finenonce{006371}



\finexercice


\exercice{6372, potyag, 2011/10/16}
\enonce{006372}{}
\begin{enumerate}
\item Montrer que tout sous-groupe d'un groupe cyclique (monogène)
est cyclique.

\item Rappelons qu'un groupe s'appelle localement cyclique si chaque
 sous-ensemble fini engendre un sous-groupe cyclique . Montrer
que $\Qq$ est localement cyclique, mais pas cyclique et en
déduire que $\Qq$ n'est pas de type fini.
\end{enumerate}
\finenonce{006372}



\finexercice


\exercice{6373, potyag, 2011/10/16}
\enonce{006373}{}
 Soit $G$ un groupe.

\begin{enumerate}
\item  Soient $A$, $B$ deux sous-groupes   de $G$.
  \begin{enumerate}
  \item On suppose que  $A$ est d'indice fini dans $G$, montrer alors que
  $A\cap B$ est d'indice fini dans $B$.

  \item On suppose que  $A$ et $B$  sont  d'indice fini dans $G$, montrer alors
  que $A\cap B$ est d'indice fini dans $G$, généraliser au cas d'un  nombre fini
  de sous-groupes.
  \end{enumerate}

\item Montrer que $\cap \{A : A$ est d'indice fini dans $\Z\}$ = $\{$id$\}$.
(Comparer avec 1.b).
\end{enumerate}
\finenonce{006373}



\finexercice


\exercice{6374, potyag, 2011/10/16}
\enonce{006374}{}
\begin{enumerate}
\item Supposons que  $H$ est d'indice fini dans  $G$ montrer qu'il
existe $K<\infty$ tel que $\forall\ g\in G\ \exists n_g\in {\Nn}^*\ :\ g^{n_g}\in H$ et  $n_g\leq K$ .


\item Montrer que $\Qq$ ne possède pas de sous-groupe d'indice fini (autre que lui-même).
\end{enumerate}
\finenonce{006374}



\finexercice


\exercice{6375, potyag, 2011/10/16}
\enonce{006375}{}
\begin{enumerate}
\item Soit $G$ un groupe et  $A$ un sous-groupe   de $G$ d'indice
fini.
 Montrer  qu'il existe  un sous-groupe  $B$ de $A$ normal dans $G$ et
d'indice fini
dans $G$.

({\it Indication} : poser $\displaystyle B= \bigcap _{g\in G} gA
g^{-1}$.)

\item Montrer qu'un groupe infini simple ne contient pas de
sous-groupe propre d'indice fini.
\end{enumerate}
\finenonce{006375}



\finexercice


\exercice{6376, potyag, 2011/10/16}
\enonce{006376}{}
Soit $G$ un groupe, $A$ et $B$ deux sous-groupes de $G$ tels que $A\subset
B$. On suppose que $A$ est d'indice fini dans $G$.
Montrer que $ \vert G : A\vert =  \vert G :  B\vert   \vert B : A\vert $.
\finenonce{006376}



\finexercice


\exercice{6377, potyag, 2011/10/16}
\enonce{006377}{}
\label{pot:exo9}
 Le but de cet exercice
est de donner la construction d'un groupe libre et d'introduire la notion de
 présentation d'un groupe.

Soit $S=\{s_i\}_{i\in I}$ un ensemble quelconque qu'on appellera
alphabet. Un mot dans l'alphabet $S$ est par définition une
succession finie (ou vide) :

$$w=s^{\epsilon_1}_{i_1}\cdot\cdot\cdot s^{\epsilon_k}_{i_k},\
o\grave u\ \epsilon_i=\pm 1\ et\ s_{i_j}\in S,\
k\in\N\hfill\eqno(1)$$


Notons $W$ l'ensemble de tous les mots. Un mot $w\in W$ est dit
réduit si son écriture (1)  ne contient pas deux lettres
consécutives du type $ s^{\epsilon}_{i}$ et $ s^{-\epsilon}_{i}$.
Les mots $w_1$ et $w_2$ sont dits voisins si $w_2=g s_i^{\epsilon}
s_i^{-\epsilon} h$ et $w_1=g h$. Deux mots $f$ et $g$ s'appellent
équivalents (on note $f\sim g$) s'il existe une succession finie
de mots : $f=w_0, w_1, ..., w_n=g$ où les mots $w_i$ et
$w_{i-1}$ sont voisins ($i\in\{1,...,n\}$).

\begin{enumerate}
\item Montrer que $\sim$ est une relation
d'équivalence.



Etant donné un mot $f=a_1 ... a_t$   (où $a_j=
s_{i_j}^{\epsilon_j}$) définissons une suite de transformations
appelée $R$-procédé :
$R_0=e$ (le mot vide), $R_1=a_1$
et 
$$R_{i+1}= 
\begin{cases}
  R_i a_{i+1} & \text{ si } R_i \text{ n'est  pas un  mot réduit du type } Xa^{-1}_{i+1} \cr
  X           & \text{ si } R_i \text{ est un mot réduit du type } Xa^{-1}_{i+1} \cr
\end{cases}.$$


Autrement dit un $R$-procédé consiste à faire
toutes les simplifications de droite à gauche.

\item On suppose que $w_1=a_1 ... a_r a_{r+1} ... a_t$ et $w_1=a_1
... a_r s_j^{\epsilon} s_j^{-\epsilon} a_{r+1} ... a_t$ sont deux
mots et que $R^i$ désigne le $R$-procédé appliqué au mot
$w_i$. Montrer que $R^1_t=R^2_{t+2}$ c.-à.-d.
$R^1(w_1)=R^2(w_2)$.
En déduire que chaque classe de $W/\sim$
contient un mot réduit et un seul.

Pour deux classes $[w_i]\in W/\sim\ (i=1,2)$ définissons
maintenant leur produit  comme suit (de gauche à droite) :

$$[w_1] [w_2] = [w_1 w_2]\hfill\eqno (2).$$

\item Démontrer que (2) ne dépend pas du choix des
représentants des classes $[w_i]$. Montrer que l'ensemble
$\displaystyle F=W/\sim$ muni de l'opération (2) est un
groupe.


Ce groupe s'appelle groupe libre engendré par $S$, on appelle
les éléments de $S$ géné\-ra\-teurs libres de $F$.

Soit maintenant $G$
un groupe quelconque engendré par un système $X$ où
$X=\{x_i\}_{i\in I}$ et $F$ est le groupe libre engendré par
$S$. Supposons qu'il existe une bijection $f:S\mapsto X$ telle que
$f(s_i)=x_i\ (i\in I)$.



\item 
  \begin{enumerate}
  \item Montrer que $f$ se prolonge en un homomorphisme $f:F\mapsto
  G$.


  \item En particulier, en déduire que si ${\text { Card } (S)}= \text { Card } (S') $, alors
   le groupe libre engendré par $S$ est isomorphe au le groupe libre engendré par
  $S'$. 
  \end{enumerate}

Si de plus $\text { Card } (S) =n$ ce groupe est  noté $F_n$.


Notons $H= {\text Ker}\ f$, et appelons un sous-ensemble
$H'\subset H$ ensemble des relations de $G$ si le plus petit
sous-groupe normal de $G$ contenant $H'$ co\"\i ncide avec $H$. La donnée
du couple $S$, $H'$ définit le groupe $G$ à un
isomorphisme près ($G$ est isomorphe au groupe quotient  $X/H$).
La donnée d'un tel couple est notée $< S\ \vert\ H'> $
 et s'appelle présentation de $G$.




\item Montrer que le groupe libre ne contient pas d'élément non-trivial d'ordre
fini.
\end{enumerate}
\finenonce{006377}



\finexercice


\exercice{6378, potyag, 2011/10/16}
\enonce{006378}{}
 Montrer que le groupe $G$ donné par sa présentation:

$$<x_1,...,x_n\ \vert\ [x_i, x_j],\ \forall\ i,j\in\{1,..., n\}\
>\hfill\eqno(3)$$

\noindent est isomorphe à $\Z^n$.
\finenonce{006378}



\finexercice


\exercice{6438, drutu, 2011/10/16}
\enonce{006438}{}
\begin{enumerate}
\item Soit l'ensemble $E=\Z \times \Z$
muni de l'opération binaire $(x_1,y_1)*(x_2,y_2)=(x_1x_2,\, x_2y_1+y_2)$.
Est-ce que $(E,*)$ est un mono\"{\i}de ? Un groupe ? Est-ce que l'opération est commutative ?

\item  Mêmes questions pour l'ensemble de matrices 
$$
M:=\left\{ \left( \begin{array}{cc}
                    a & b \\
                    c & d 
                   \end{array} \right) ; a,b,c,d\in \Z, a+b=c+d \right\}\; , $$
 muni de la multiplication des matrices. 
\end{enumerate}
\finenonce{006438}



\finexercice


\exercice{6439, drutu, 2011/10/16}
\enonce{006439}{}
 Soit l'ensemble $G=\lbrack 0,1\lbrack $
muni de la
 loi de composition interne $x*y=\lbrace x+y\rbrace$, où
$\lbrace x\rbrace$ represente la partie fractionelle du nombre
réel $x$. Montrer que $(G,*)$ est un groupe. Montrer que
$\Q\cap G$ est un sous-groupe.
\finenonce{006439}



\finexercice


\exercice{6440, drutu, 2011/10/16}
\enonce{006440}{}
 Soit l'ensemble de matrices 
$$
K:=\left\{ A(x)= \left( \begin{array}{ccc}
                    1-x & 0 & x \\
                    0 & 0 & 0 \\
                    x & 0 & 1-x 
                   \end{array} \right) ; x\in \R \setminus \{ 1/2 \}\right\}\; , $$
 muni de la multiplication des matrices. Montrer que c'est un groupe commutatif (abélien).
\finenonce{006440}



\finexercice


\exercice{6441, drutu, 2011/10/16}
\enonce{006441}{}
 Soient $a,b,c$ les éléments d'un groupe. Montrer que
l'équation $xaxba=xbc$ admet une solution et une seule.
\finenonce{006441}



\finexercice


\exercice{6442, drutu, 2011/10/16}
\enonce{006442}{}
Dans le groupe des entiers modulo $11$ muni de l'opération de
multiplication, lesquels parmi les ensembles suivants forment des sous-groupes ?


$$
(a) \{ 1,3,4,5,9 \},\hspace{0.5cm} (b) \{ 1,3,5,7,8\},\hspace{0.5cm} (c) \{ 1,8\},\hspace{0.5cm} (d) \{ 1,10\},\hspace{0.5cm} (e) \{ 1,3,10\}.
$$
\finenonce{006442}



\finexercice


\exercice{6443, drutu, 2011/10/16}
\enonce{006443}{}
 Montrer qu'un groupe ayant au plus $4$ éléments  est abélien. 
({\it{Indication}} : Comparer les éléments $e,a,b,ab,ba$.)
\finenonce{006443}



\finexercice


\exercice{6444, drutu, 2011/10/16}
\enonce{006444}{}
Lesquels des ensembles de nombres suivants sont des groupes ?
\begin{enumerate}
\item Les nombres irrationnels munis de l'addition ; de la multiplication ;

\item Les nombres complexes de valeur absolue $1$ munis de l'addition ; de la multiplication ; 

\item Les nombres complexes munis de l'opération binaire $z*z'=|z|\cdot z'$.
\end{enumerate} 
\finenonce{006444}



\finexercice


\exercice{6445, drutu, 2011/10/16}
\enonce{006445}{}
Lesquelles parmi les tables de Cayley suivantes décrivent un groupe ?

\begin{tabular}{c|c|c|c|c}
* & a & b & c & d \\
\hline
a &  d & c & b & a \\
\hline
b & c & d & a & b \\
\hline 
c & b & a & d & c \\
\hline
d & a & b & c & d 
\end{tabular}\hspace{1cm} 
\begin{tabular}{r|r|r|r|r|r}
* & a & b & c & d & e \\
\hline
a &  e & d & b & c & a\\
\hline
b & c & e & d & a & b \\
\hline 
c & d & a & e & b & c \\
\hline
d & b & c & a & e & d \\
\hline
e & a & b & c & d & e 
\end{tabular}
\finenonce{006445}



\finexercice


\exercice{6446, drutu, 2011/10/16}
\enonce{006446}{}
 Montrer que si un ensemble $E$ est muni d'une opération binaire qui vérifie les propriétés 
\begin{itemize}
\item $(P_1)$ $(ab)c=a(cb)$ ;
\item $(P_2)$ il existe $e\in E$ tel que $ea=a,\; \forall a$;
\item $(P_3)$ pour tout $a\in E$ il existe $b\in E$ tel que $ba=e$;
\end{itemize}
alors c'est un groupe abélien.
\finenonce{006446}



\finexercice


\exercice{6447, drutu, 2011/10/16}
\enonce{006447}{}
Montrer que, dans un groupe $G$, pour tout élément $a\in G$, l'ensemble des $x\in G$ tels que $ax=xa$ est un sous-groupe (apellé le {\it{centralisateur de $a$ dans $G$}}).

Montrer que l'ensemble des $x\in G$ tels que $ax=xa,\; \forall a\in G$, est un sous-groupe (apellé le {\it{centre de $G$}}).  
\finenonce{006447}



\finexercice


\exercice{6448, drutu, 2011/10/16}
\enonce{006448}{}
Combien de générateurs différents a un groupe cyclique d'ordre $6$?
\finenonce{006448}



\finexercice


\exercice{6449, drutu, 2011/10/16}
\enonce{006449}{}
Soit $G$ un un groupe engendré par deux éléments $x$ et $y$ qui vérifient les relations $x^2=y^3=e,\, xy=yx$. Ecriver tous les éléments de $G$ et sa table de multiplication. 
\finenonce{006449}



\finexercice


\exercice{6450, drutu, 2011/10/16}
\enonce{006450}{}
Trouver une décomposition en produit de
transpositions des permutations suivantes :
\begin{enumerate}
\item $$\sigma_1 =\left(
                       \begin{array}{ccc}
                         1 & 2 & 3 \\
                         2 & 3 & 1
                       \end{array} 
                                         \right) \in {\bf{S}}_3\; ;$$
\item   $$\sigma_2 =\left(
                       \begin{array}{cccc}
                         1 & 2 & 3 & 4 \\
                         4 & 2 & 3 & 1
                       \end{array} 
                                         \right) \in {\bf{S}}_4\; ;$$
\item   $$\sigma_3 =\left(
                       \begin{array}{ccccc}
                         1 & 2 & 3 & 4 & 5\\
                         5 & 1 & 4 & 3 & 2
                       \end{array} 
                                          \right) \in {\bf{S}}_5\; .$$
\end{enumerate} 
 Quelle est leur signature ? 
\finenonce{006450}



\finexercice


\exercice{6451, drutu, 2011/10/16}
\enonce{006451}{}
Déterminer la signature des permutations :
\begin{enumerate}
\item $$\sigma =\left(
                       \begin{array}{cccccc}
                         1 & 2 & \ldots & k & \ldots & n \\
                         n & n-1 & \ldots & n-k+1 & \ldots & 1
                       \end{array} 
                                            \right) \in {\bf{S}}_n\; ;$$


\item   $$\sigma =\left(
                       \begin{array}{ccccccc}
                         1 & 2 & \ldots & n-p & n-p+1 & \ldots & n \\
                         p+1 &p+2 & \ldots & n & 1 & \ldots & p
                       \end{array} 
                                             \right) \in {\bf{S}}_n\; ,$$ où $p$ est
fixé, $1\leq p\leq n-1$ et $n\geq 2$.
\end{enumerate} 
\finenonce{006451}



\finexercice


\exercice{6452, drutu, 2011/10/16}
\enonce{006452}{}
 Décomposer en produit de cycles à supports deux à deux disjoints
 les permutations suivantes
$$\sigma_1 =\left(
                       \begin{array}{ccccccc}
                         1 & 2 & 3 & 4 & 5 & 6 & 7 \\
                         6 & 4 & 5 & 7 & 3 & 1 & 2
                       \end{array} 
                                             \right) \in {\bf{S}}_7\; ,$$
$$\sigma_2 =\left(
                       \begin{array}{ccccccccc}
                         1 & 2 & 3 & 4 & 5 & 6 & 7 & 8 & 9\\
                         3 & 7 & 8 & 9 & 4 & 5 & 2 & 1 & 6
                       \end{array} 
                                             \right) \in {\bf{S}}_9\; .$$


En déduire la signature.
\finenonce{006452}



\finexercice


\exercice{6453, drutu, 2011/10/16}
\enonce{006453}{}
 Ecriver les produits suivants comme produits de cycles disjoints.
$$
(a) (1234)(567)(261)(47){\hspace{0.5cm}}(b) (12345)(67)(1357)(163){\hspace{0.5cm}}(c) (14)(123)(45)(14)
$$


Trouver la signature de chaque produit. 
\finenonce{006453}



\finexercice


\exercice{6454, drutu, 2011/10/16}
\enonce{006454}{}
\begin{enumerate}
\item  Vérifier que pour tout cycle $(i_1\, i_2\dots i_p),\;
p\geq 2$, dans $S_n$ et toute permutation $\sigma \in S_n$, 
$$
\sigma (i_1\, i_2\dots i_p)\sigma^{-1}=(\sigma(i_1)\, \sigma(i_2)\dots
\sigma(i_p))\, .
$$

\item Vérifier que pour tous entiers distincts $i,j\in \lbrace 1,2,\dots
n\rbrace$ on a $(i\, j)=(i\, k)(j\, k)(i\, k)$.

\item En déduire que les familles de transpositions $\lbrace (1\, i)\mid i\in
\lbrace 1,2,\dots n\rbrace \rbrace$ et $\lbrace (i\, i+1)\mid i\in
\lbrace 1,2,\dots n-1\rbrace \rbrace$ engendrent $S_n$.

\item Soient $\tau =(1\, 2)$ et $c=(1\, 2\, \dots n)$. Calculer $c^{i-1}\tau
c^{1-i}$ pour tout $i\in \lbrace 1,2,\dots n-1\rbrace$. Montrer que toute
permutation de $S_n$ s'écrit comme un produit de puissances de $\tau$ et
$c$.
\end{enumerate} 
\finenonce{006454}



\finexercice


\exercice{6455, drutu, 2011/10/16}
\enonce{006455}{}
 Dans le groupe symétrique $S_4$ trouver les sous-groupes suivants 
\begin{enumerate}
\item le sous-groupe des éléments $\sigma$ tels que l'image par $\sigma$ de l'ensemble $\{ 1,2 \}$ est $\{ 1,2 \}$.

\item le sous-groupe des éléments $\tau $ tels que si $a\equiv b$ mod $2$ alors $\tau(a)\equiv \tau(b)$ mod $2$. ({\it{Indication}} : $(13)(24)$ fait partie de ce sous-groupe).  
\end{enumerate} 
\finenonce{006455}



\finexercice


\exercice{7887, mourougane, 2021/08/11}
\enonce{007887}{Groupe unipotent}
\begin{enumerate}
 \item \`A quel groupe le groupe $U$ des matrices de la forme $\begin{bmatrix}1&x\\0&1\end{bmatrix}$ avec $x\in\Rr$ est-il isomorphe ?
\item Est-ce un sous-groupe normal de $SL(2,\Rr)$ ?
\item Déterminer l'inverse de $\begin{bmatrix}1&3\\0&1\end{bmatrix}$.
\end{enumerate}
\finenonce{007887}
\finexercice

\section{321.00 Sous-groupe, morphisme}


\exercice{6456, drutu, 2011/10/16}
\enonce{006456}{}
\begin{enumerate}
\item Trouver l'ordre des permutations suivantes 
$$
(abcdef)(ghij)(klm)\; ;\; \; \;  \; (123456)(1234)(123)\, .
$$

\item Montrer que toute permutation d'ordre $10$ dans $S_8$ est impaire.
\end{enumerate} 
\finenonce{006456}



\finexercice


\exercice{6457, drutu, 2011/10/16}
\enonce{006457}{}
 Soit l'ensemble $E=\{ a,b,c,d,e \}$ qui, muni de l'opération
binaire $*$, devient un groupe. Trouver la table de multiplication de ce groupe
sachant que $a*b=d,\; c*a=e,\; d*c=b$.
\finenonce{006457}



\finexercice


\exercice{6458, drutu, 2011/10/16}
\enonce{006458}{}
\begin{enumerate}
\item Soit $n\geq 2$ un nombre entier. Montrer que $n$ est premier si et seulement si tout groupe d'ordre $n$ a seulement deux sous-groupes.

\item Montrer qu'un groupe d'ordre $p^m,\; m\geq 1,$ où $p\in \N$ est nombre premier, contient un sous-groupe d'ordre $p$. 
\end{enumerate} 
\finenonce{006458}



\finexercice


\exercice{6459, drutu, 2011/10/16}
\enonce{006459}{}
Montrer que les groupes suivants ne sont pas isomorphes :
\begin{enumerate}
\item $( \Z ,\; +)$ et $(\Z[X],\; +)$ ;

\item $(\Q ,\; +)$ et $(\Q[i] ,\; +)$, où $\Q[i]:=\{ a+ib\in \C \mid a,b\in \Q \}$.
\end{enumerate} 
\finenonce{006459}



\finexercice


\exercice{6460, drutu, 2011/10/16}
\enonce{006460}{}
Soit $K=\{ e,a,b,c \}$ un groupe d'ordre $4$ tel que $x^2=e, \forall x$.
\begin{enumerate}
\item Ecrire la table de multiplication de $K$.

\item Montrer que $K$ est isomorphe à $(\Z/2 \Z \times \Z/2 \Z ,\; +)$ et n'est pas isomorphe à $\Z /4 \Z$.

\item Montrer que tout groupe d'ordre 4 est ou bien isomorphe à $K$ ou bien à $\Z /4 \Z$.
\end{enumerate} 
\finenonce{006460}



\finexercice


\exercice{6461, drutu, 2011/10/16}
\enonce{006461}{}
Soit $(G,\cdot)$ un groupe fini d'ordre $n$. Trouver toutes les morphismes de groupe $\phi :(\Q,\; +)\to (G,\cdot)$.
\finenonce{006461}



\finexercice


\exercice{6462, drutu, 2011/10/16}
\enonce{006462}{}
 Soit $(G,\cdot)$ un groupe. On appelle sous-groupe {\it{caractéristique}} de $G$ un sous-groupe invariant par tout automorphisme de $G$. Montrer que tout sous-groupe caractéristique est distingué.
\finenonce{006462}



\finexercice


\exercice{6463, drutu, 2011/10/16}
\enonce{006463}{}
\begin{enumerate}
\item Montrer que $S_3$ contient un sous-groupe $H$ distingué d'ordre $3$ et que $G/H$ est isomorphe à $\Z /2\Z $.

\item Montrer que les seuls groupes d'ordre $6$ sont, à isomorphisme près, le groupe cyclique et $S_3$.
\end{enumerate} 
\finenonce{006463}



\finexercice


\exercice{6464, drutu, 2011/10/16}
\enonce{006464}{}
Combien d'automorphismes a un groupe cyclique d'ordre $p$, où $p$ est un nombre premier ? Et un groupe cyclique d'ordre $pq$, où $p$ et $q$ sont des nombres premiers distincts ?
\finenonce{006464}



\finexercice


\exercice{6465, drutu, 2011/10/16}
\enonce{006465}{} 
Montrer que les seuls groupes distingués de $S_5$ sont $\{ e\} ,\; S_5,\; A_5$.
\finenonce{006465}



\finexercice


\exercice{6466, drutu, 2011/10/16}
\enonce{006466}{} 
 Soit $(G,\cdot )$ un groupe. On appelle {\it commutateur} un élément de la forme $xyx^{-1}y^{-1}$. On note $[x,y]:=xyx^{-1}y^{-1}$. 

\begin{enumerate}
\item Montrer que l'ensemble $G'=[G,G]$ des produits de commutateurs est un sous-groupe distingué de $G$. Montrer que $G/G'$ est un groupe abélien. En particulier cela implique que $G'=\{e \}\Leftrightarrow G$ abélien.

\item Soit $N \vartriangleleft G$ tel que $G/N$ est abélien. Montrer que $G'\subset N$.

\item Trouver $G'$ pour $G=S_3,\; A_5,\; S_5$.

\item Montrer que si $\phi :G\to A$ est un morphisme de $G$ à un groupe abélien $A$, alors il existe un morphisme $\psi :G/G' \to A$ tel que $\phi = \psi \circ \pi $, où $\pi : G \to G/G'$ est la projection canonique.  
\end{enumerate} 
\finenonce{006466}



\finexercice


\exercice{6467, drutu, 2011/10/16}
\enonce{006467}{}
 Soit $V_4:=\{ e,\; (12)(34),\; (13)(24),\; (23)(14) \}\subset S_4$.
\begin{enumerate}
\item Montrer que c'est un sous-groupe distingué de $S_4$.

\item Montrer que $V_4\simeq \Z/2 \Z \times \Z/2 \Z$.

\item Soit $H=\{ e,\; (12)(34) \}$. Montrer que $H \vartriangleleft V_4$ mais que $H$ n'est pas un sous-groupe distingué de $S_4$.

\item Montrer que $S_4/V_4$ est isomorphe à $S_3$.
\end{enumerate} 
\finenonce{006467}



\finexercice


\exercice{6468, drutu, 2011/10/16}
\enonce{006468}{}
\begin{enumerate}
\item Montrer que tout sous-groupe d'un groupe cyclique est cyclique.

\item Soit $d\in \N^*$ un diviseur de $n\in \N^*$. Montrer qu'il existe un unique sous-groupe d'ordre $d$ dans $\Z /n\Z $. 
\end{enumerate} 
\finenonce{006468}



\finexercice


\exercice{6469, drutu, 2011/10/16}
\enonce{006469}{}
\begin{enumerate}
\item Soit $G$ un groupe fini et $H$ un sous-groupe d'indice $p$, où $p$ est le plus petit facteur premier de $|G|$. Montrer que $H \vartriangleleft G$.

\item Soit $H<S_4$, $H$ de cardinal 12. Montrer que $H=A_4$.

\item Montrer que $A_4$ n'est pas simple.

\item Montrer que $A_4$ n'a pas de sous-groupe d'ordre $6$. 
\end{enumerate}  
\finenonce{006469}



\finexercice


\exercice{6470, drutu, 2011/10/16}
\enonce{006470}{}
Démontrer les propriétés suivantes 

\begin{enumerate}
\item $(\C / \R ,+) \simeq (\R,+)$ ;
\item $\C^*/\R_+^*\simeq \mathbb{S}^1$ ;
\item $\C^*/\mathbb{S}^1\simeq \R_+^*\simeq \R^*/\{ \pm 1 \}$ ;
\item $\R^*/ \R_+^*\simeq \{ \pm 1 \}$ ;
\item $(\C / \Z ,+)\simeq \C^*$ ;
\item $\C^* /\R^* \simeq \mathbb{S}^1/\{ \pm 1 \}$.
\end{enumerate} 
\finenonce{006470}



\finexercice


\exercice{6471, drutu, 2011/10/16}
\enonce{006471}{}
Prouver que les permutations 
$$\sigma =\left(
                       \begin{array}{cccccc}
                         1 & 2 & 3 & 4 & 5 & 6\\
                         2 & 5 & 3 & 6 & 1 & 4
                       \end{array} 
                                          \right) \hspace{0.5cm}\mbox{ et }\hspace{0.5cm} \tau =\left(
                       \begin{array}{rrrrrr}
                         1 & 2 & 3 & 4 & 5 & 6\\
                         5 & 3 & 4 & 2 & 1 & 6
                       \end{array} 
                                          \right)$$ 
 sont conjuguées dans $S_6$.  
\finenonce{006471}



\finexercice


\exercice{6472, drutu, 2011/10/16}
\enonce{006472}{}
Soit $G$ un sous-groupe du groupe symétrique $S_n,\; n\geq 2$, contenant une permutation impaire. Montrer que $GA_n=S_n$ et en déduire que $G$ contient au moins un sous-groupe distingué d'indice 2.
\finenonce{006472}



\finexercice


\exercice{6473, drutu, 2011/10/16}
\enonce{006473}{}
 Soit $SL(n,\R)=\{ M\in M_n(\R )\mid det\, M=\pm 1 \}$ et soit l'action du groupe $SL(n,\R)$ sur $\R^n$ donnée par 
$$
SL(n,\R)\times \R^n \to \R^n\, ,
$$
$$
(M,\, X)\to MX\; .
$$

\begin{enumerate}
\item  Trouver les orbites de cette action. Montrer que le centre de $SL(n,\R)$ est $\{ \pm I \}$. 

\item Montrer que, pour tout vecteur $X\in \R^n$, son stabilisateur est conjugué au sous-groupe 
$$
P= \left(
                       \begin{array}{cc}
                         1 & *\\
                         0 & SL(n-1,\R)
                       \end{array} 
                                          \right)\; .$$
                                          
\item Soit $SL(2,\Z )$ muni de son action sur $\Z^2$ définie par
$$
SL(2,\Z)\times \Z^2 \to \Z^2\, ,
$$
$$
(M,\, X)\to MX\; .
$$

Trouver ses orbites. Montrer que l'ensemble $\Gamma(2)$ des matrices $M$ tels que $M=I$ (mod 2) est un sous-groupe de $SL(2,\Z )$. Trouver ses orbites dans $\Z^2$. 
 \end{enumerate} 
\finenonce{006473}



\finexercice


\exercice{6474, drutu, 2011/10/16}
\enonce{006474}{}
Soit $D_\infty $ le groupe des isométries de la droite affine $\R $ formé par l'ensemble des éléments de la forme $\tau_1^n$ et $\tau_1^n \circ \sigma$, où $n\in \Z$, $\tau_1(x)=x+1$ et $\sigma (x)=-x$. On appelle $D_\infty$ {\it le groupe diédral infini}. 
\begin{enumerate}
\item Montrer que $H=<\tau_1>$ est le seul sous-groupe cyclique infini d'indice 2 de $D_\infty $. Montrer que $H$ est un sous-groupe distingué de $D_\infty $.

\item Montrer que, pour tout sous-groupe $S$ d'ordre 2 de $D_\infty$, on a $D_\infty = SH$.

\item Soit $K< D_\infty$ tel que $K \not\subset H$. Montrer que $D_\infty =HK$. En déduire que $K\cap H$ est d'indice 2 dans $K$. Montrer que $K\cap H \neq (e)$ implique $K\simeq D_\infty$.

\item Montrer que tout sous-groupe propre de $D_\infty $ est isomorphe soit à $\Z$, soit à $(\pm 1)$, soit à $D_\infty$.

\item On note ${\mathcal S}_n$ l'ensemble des sous-groupes $K<D_\infty$ tels que $K \not\subset H$ et $K\cap H=H_n$, où $H_n:=<\tau_1^n>$. Prouver que ${\mathcal S}_n$ contient $n$ éléments. 
\end{enumerate}   
\finenonce{006474}



\finexercice

\section{322.00 Groupe fini}


\exercice{6475, drutu, 2011/10/16}
\enonce{006475}{}
\begin{enumerate}
\item Soit les applications $T:\R \cup\{ \infty \} \to \R \cup\{ \infty \} ,\; T(x)=x+4,\; T(\infty )=\infty $, et $g:\R \cup\{ \infty \} \to \R \cup\{ \infty \}$, 
$$
g(x)=\left\{ \begin{array}{ccc} 
          \frac{x}{1-2x} & \mbox{ si } & x\in \R\setminus \{ \frac{1}{2}\} \; ; \\
          \infty & \mbox{ si } & x= \frac{1}{2} \; ;\\ 
          -\frac{1}{2} & \mbox{ si } & x= \infty \;   
          \end{array}\right.
$$ 

Montrer que $T^k([-2,2])\subset ]-\infty ,\, -2]\cup [2,\, \infty [,\; \forall k\in \Z^* $ et que $g^m(]-\infty ,\, -1]\cup [1,\, \infty [)\subset [-1,\, 1],\; \forall m\in \Z^* $.


\item Soit $G$ le groupe des applications bijectives $h:\R\cup \{ \infty \} \to \R\cup \{ \infty \}$, muni de l'opération de composition. Montrer que $G$ a une action naturelle sur $\R \cup\{ \infty \}$.

\item Montrer que le sous-groupe $\Gamma $ de $G$ engendré par $T$ et $g$ est un groupe libre. {\it Indication }: Regarder les orbites des nombres dans l'intervalle $]1,2[$.
\end{enumerate}
\finenonce{006475}



\finexercice


\exercice{6476, drutu, 2011/10/16}
\enonce{006476}{}
\begin{enumerate}
\item Soit $G$ le groupe engendré par deux éléments $a,b$, défini par les relations $a^m=b^2=e,\; (ab)^2=e$, où $m\in \N^*\cup \{\infty  \}$. Montrer que $G\simeq D_m$, où $D_m$ est le groupe diédral de degré $m$.

\item Montrer que pour tout groupe $H$ engendré par deux éléments $\alpha,\; \beta$, qui vérifient les relations $\alpha^m=\beta^2=e,\; (\alpha\beta)^2=e$ il existe un epimorphisme de $D_m$ dans $H$.
\end{enumerate}
\finenonce{006476}



\finexercice


\exercice{6477, drutu, 2011/10/16}
\enonce{006477}{}
 Soit $H_i$ et $N_i$ deux paires de groupes et $\phi_i :H_i \to Aut\, (N_i)$ deux morphismes, $i=1,2$. Montrer que s'il existe deux isomorphismes $\alpha :H_1 \to H_2$ et $\beta :N_1 \to N_2$ tels que $\phi_2(\alpha(h_1))=\beta \circ \phi_1(h_1)\circ \beta^{-1},\, \forall h_1\in H_1$, alors $N_1 \rtimes_{\phi_1} H_1 \simeq N_2 \rtimes_{\phi_2} H_2$. 
\finenonce{006477}



\finexercice


\exercice{6478, drutu, 2011/10/16}
\enonce{006478}{}
Soit $H$ et $N$ deux groupes et $\phi ,\; \psi :H \to Aut\, (N)$ deux morphismes. Montrer que 
\begin{enumerate}
\item s'il existe $\alpha \in Aut\, (H)$ tel que $\psi =\phi \circ \alpha $, alors $N \rtimes_\phi H \simeq N \rtimes_\psi H$.

\item s'il existe $u \in Aut\, (N)$ tel que $\psi(h) =u\phi(h)u^{-1}$, alors $N \rtimes_\phi H \simeq N \rtimes_\psi H$.
\end{enumerate}
\finenonce{006478}



\finexercice


\exercice{6479, drutu, 2011/10/16}
\enonce{006479}{}
\begin{enumerate}
\item (Théorème d'Euler) Soit $n\in \N^*$ et $a\in \Z^*$ premier avec $n$. Démontrer que $a^{\phi(n)}\equiv 1$ (mod $n$), où $\phi$ est la fonction d'Euler.

\item (Théorème de Fermat) Soit $p$ un nombre premier et $a\in \Z$. Montrer que $a^p\equiv a $(mod $p$). 
\end{enumerate}
\finenonce{006479}



\finexercice


\exercice{6480, drutu, 2011/10/16}
\enonce{006480}{}
Soit $\phi $ un automorphisme de $S_n$. Montrer que si $\phi $ transforme toute transposition en une transposition alors $\phi $ est un automorphisme intérieur. {\it Indication} : Utiliser le fait que $S_n$ est engendré par $\{ (12),\; (13),\dots ,\; (1n) \}$. 
\finenonce{006480}



\finexercice


\exercice{6481, drutu, 2011/10/16}
\enonce{006481}{}
\begin{enumerate}
\item Montrer que tout groupe d'ordre $pq$ avec $p,q$ premiers
distincts, est un produit semi-direct de deux sous-groupes cycliques.

\item Déterminer à isomorphisme près tous les groupes d'ordre $pq$.
\end{enumerate}
\finenonce{006481}



\finexercice


\exercice{6482, drutu, 2011/10/16}
\enonce{006482}{}
Soit $G$ un groupe fini opérant sur un ensemble fini $E$. 
\begin{enumerate}
\item On suppose que l'action est telle que toute orbite de $G$ contient au moins 2 points. Si $|G|=15$ et card $E$=17, trouver le nombre d'orbites de $G$ dans $E$ et le cardinal de chacune.

\item Montrer que si $|G|=33$ et card $E$=19, il existe au moins une orbite qui contient un unique point.
\end{enumerate}
\finenonce{006482}



\finexercice


\exercice{6483, drutu, 2011/10/16}
\enonce{006483}{}
Soit $G$ un $p$-groupe opérant sur un ensemble fini $X$ et soit $Fix(G):=\{x\in X \mid gx=x,\; \forall g\in G \}$. Montrer que card $X=$ card $Fix(G)$ (mod $p$). En déduire que le centre d'un $p$-groupe est toujours non-trivial.
\finenonce{006483}



\finexercice


\exercice{6484, drutu, 2011/10/16}
\enonce{006484}{}
Montrer qu'un groupe infini $G$ qui a un sous-groupe propre d'indice fini $H$ n'est pas simple. {\it Indication} : Etudier l'action de $G$ sur $G/H$.
\finenonce{006484}



\finexercice


\exercice{6485, drutu, 2011/10/16}
\enonce{006485}{}
Soit $V$ un $\R$-espace vectoriel de dimension finie et $GL(V)$ le groupe des automorphismes linéaires. Pour $f\in GL(V)$ et $a\in V$ on considère l'application $A_{f,a}:V\to V,\; A_{f,a}(v)=f(v)+a$.
\begin{enumerate}
\item Montrer que ${\mathcal{A}}:=\{A_{f,a}\mid f\in GL(V),\; a\in V \}$ muni de l'opération de composition est un groupe. Ce groupe s'appelle {\it le groupe des transformations affines}.

\item \label{q:2} Trouver deux sous-groupes de ${\mathcal{A}}$ isomorphes à $(GL(V),\cdot)$ et à $(V,+)$, respectivement.

\item Montrer que ${\mathcal{A}}$ est produit semi-direct des deux sous-groupes obtenus dans \ref{q:2}.
\end{enumerate}
\finenonce{006485}



\finexercice


\exercice{6486, drutu, 2011/10/16}
\enonce{006486}{}
Prouver que tout sous-groupe d'ordre 35 est cyclique.
\finenonce{006486}



\finexercice


\exercice{6487, drutu, 2011/10/16}
\enonce{006487}{}
 Soit $G$ un groupe fini avec $|G|=p^2q$, où $p,q$ sont deux nombres premiers distincts tels que $p^2\not\equiv 1$ (mod $q$) et $q\not\equiv 1$ (mod $p$). Montrer que $G$ est abélien.
\finenonce{006487}



\finexercice


\exercice{6488, drutu, 2011/10/16}
\enonce{006488}{}
\begin{enumerate}
\item Prouver qu'un groupe d'ordre $p^2q$ ne peut pas être simple.

\item Montrer qu'un groupe d'ordre 63 n'est pas simple.
\end{enumerate}
\finenonce{006488}



\finexercice


\exercice{6489, drutu, 2011/10/16}
\enonce{006489}{}
 Déterminer à isomorphisme près tous les groupes d'ordre 12. Reconna\^{\i}tre parmi eux $D_6$ et $A_4$. Même question pour les groupes d'ordre 18. 
\finenonce{006489}



\finexercice


\exercice{6490, drutu, 2011/10/16}
\enonce{006490}{}
Soit $G$ un groupe d'ordre 399. Vérifier que $G$ a un sous-groupe distingué d'ordre 19 et un sous-groupe distingué d'ordre 133. En déduire que $G$ est un produit semi-direct de deux groupes cycliques.
\finenonce{006490}



\finexercice


\exercice{6491, drutu, 2011/10/16}
\enonce{006491}{}
 Soit $G$ un groupe de cardinal 24. Montrer que, si aucun de ses sous-groupes de Sylow n'est distingué, $G\simeq S_4$. {\it Indication} : Faire opérer $G$ sur ses 3-sous-groupes de Sylow.
\finenonce{006491}



\finexercice


\exercice{6492, drutu, 2011/10/16}
\enonce{006492}{}
Soit $G$ un groupe et $K$ un sous-groupe fini distingué. Montrer que tout $p$-sous-groupe de Sylow distingué de $K$ est distingué dans $G$.
\finenonce{006492}



\finexercice


\exercice{6493, drutu, 2011/10/16}
\enonce{006493}{}
 Soit $G$ un groupe fini, $H$ un sous-groupe distingué et $p$ un nombre premier divisant $[G:H]$. Montrer que $\Sigma $ est un $p$-sous-groupe de Sylow de $G/H$ si et seulement si il existe un $p$-sous-groupe de Sylow $S$ de $G$ tel que $\Sigma = SH/H$.
\finenonce{006493}



\finexercice


\exercice{6494, drutu, 2011/10/16}
\enonce{006494}{}
Soit $G$ un groupe abélien de type fini. Montrer que si tout
élément de $G$ est d'ordre fini, alors $G$ est fini.
\finenonce{006494}



\finexercice


\exercice{6495, drutu, 2011/10/16}
\enonce{006495}{}
Montrer que $\Q /\Z$ est un groupe abélien infini dont tout élément est d'ordre fini. En déduire qu'il ne peut pas avoir une famille finie de générateurs.
\finenonce{006495}



\finexercice


\exercice{6496, drutu, 2011/10/16}
\enonce{006496}{}
 Soit $G$ un groupe fini abélien. Montrer que pour tout diviseur $d$ de $| G|$ il existe un sous-groupe de $G$ d'ordre $d$.
\finenonce{006496}



\finexercice


\exercice{6497, drutu, 2011/10/16}
\enonce{006497}{}
\begin{enumerate}
\item Trouver les invariants et la décomposition canonique du groupe abélien fini dont les diviseurs élémentaires (appellées aussi {\it invariants primaires}) sont $2^3,\; 2,\; 3^2,\; 3,\; 3$. 

\item Trouver les diviseurs élémentaires/invariants primaires, les invariants et la décomposition canonique de $G={\mathcal{C}}_{30}\times {\mathcal{C}}_{18}$.

\item Trouver les diviseurs élémentaires/invariants primaires, les invariants et la décomposition canonique des groupes abéliens suivants :
  \begin{enumerate}
  \item $G_1$ engendré par $a$ et $b$ tels que $10a=9b=0$;
  \item  $G_2$ engendré par $a,\; b$ et $c$ tels que $15a=6b=4c=0$.
  \end{enumerate}
\end{enumerate}
\finenonce{006497}



\finexercice


\exercice{6498, drutu, 2011/10/16}
\enonce{006498}{}
Soit $\C $ muni des opérations binaires 

$$
z_1*z_2=z_1+z_2,\; z_1\perp z_2=z_1z_2+Im\; z_1 Im\; z_2 \; .
$$
\begin{enumerate}
\item Montrer que $(\C , *,\perp )$ est un anneau. Trouver ses éléments inversibles.

\item Montrer que l'ensemble de matrices 
$$
D:=\left\{ \left(
         \begin{array}{cc}
         a & b \\
         0 & a
         \end{array} \right) \;  ;\;  a\in \R,\; b\in \R \right\}\,  .  
$$
 muni de l'addition et la multiplication des matrices est un anneau.
 
\item Montrer que la fonction $f:\C \to D,$ 
$$
f(x+iy):=\left(
         \begin{array}{cc}
         x & y \\
         0 & x
         \end{array} \right)
$$
 est un isomorphisme d'anneaux.
\end{enumerate}
\finenonce{006498}



\finexercice


\exercice{6499, drutu, 2011/10/16}
\enonce{006499}{}  
Soit $A$ un anneau non nécessairement commutatif. Soit $a,b\in A$ tels que $a,\; b,\; ab-1$ sont inversibles. Montrer que $a-b^{-1}$ et $(a-b^{-1})^{-1}-a^{-1}$ sont inversibles. Montrer qu'on a l'égalité\\ \centerline{$[(a-b^{-1})^{-1}-a^{-1}]^{-1}=aba-a$.}
\finenonce{006499}



\finexercice


\exercice{6500, drutu, 2011/10/16}
\enonce{006500}{} 
Montrer que les anneaux de polyn\^omes $\R[X]$ et $\C[X]$ ne sont pas isomorphes.
\finenonce{006500}



\finexercice


\exercice{6501, drutu, 2011/10/16}
\enonce{006501}{} 
Soit $A$ un anneau non nécessairement commutatif, $A^*$ le groupe des éléments inversibles et $I$ un idéal bilatère de $A$. Soit $U=\{ a\in A^*\mid a\equiv 1\; mod\; I \}$. Montrer que $I$ est un sous-groupe distingué de $A^*$.
\finenonce{006501}



\finexercice


\exercice{7718, mourougane, 2021/08/11}
\enonce{007718}{}
\begin{enumerate}
 \item Montrez que les morphismes d'un groupe simple vers un groupe quelconque sont constants ou injectifs.
 \item Quel est le nombre moyen de points fixes d'une permutation de $\mathcal{S}_n$ ?
 \item Démontrez que si $G$ est un groupe fini, $S$ un $p$-Sylow de $G$, et $H$ un sous-groupe de $G$, il existe un conjugué de $S$ qui rencontre $H$ en un $p$-Sylow de $H$. 
\end{enumerate}
\finenonce{007718}
\finexercice


\exercice{7719, mourougane, 2021/08/11}
\enonce{007719}{}
 Soit $G$ un groupe fini. Soit $p$ le plus petit facteur premier de l'ordre de $G$. Soit $H$ un sous-groupe de $G$ d'indice $p>1$.
\begin{enumerate}
 \item Montrez que les orbites de l'action de $H$ sur $G/H$ par translation à gauche sont réduites à des points.
 \item Montrez que $H$ est distingué.
\end{enumerate}
\finenonce{007719}
\finexercice


\exercice{7720, mourougane, 2021/08/11}
\enonce{007720}{}
\begin{enumerate}
 \item Explicitez un $7$-Sylow du groupe symétrique $\mathcal{S}_{9}$
 \item Déterminez le nombre d'éléments d'ordre $7$ dans $\mathcal{S}_{9}$.
 \item Déterminez le nombre de $7$-Sylow.
 \item Vérifiez les congruences données par le théorème de Sylow sur le nombre de $7$-Sylow.
\end{enumerate}
\finenonce{007720}
\finexercice


\exercice{7721, mourougane, 2021/08/11}
\enonce{007721}{}
Soit $G$ un groupe et $f~:G\to A$ un morphisme de groupes de $G$ dans un groupe abélien $A$. On suppose de plus que le noyau $N(f)$ de $f$ est résoluble.
\begin{enumerate}
 \item Montrez que le groupe dérivé $D(G)$ de $G$ est inclus dans le noyau de $F$. 
 \item Montrez que $G$ est résoluble.
\end{enumerate}
\finenonce{007721}
\finexercice


\exercice{7722, mourougane, 2021/08/11}
\enonce{007722}{} 
Un groupe $G$ de cardinal $169$ agit sur un espace $X$ à $207$ éléments. On suppose qu'il y a exactement $15$ orbites distinctes. Déterminer le nombre d'orbites de chaque cardinal. 
\finenonce{007722}
\finexercice


\exercice{7723, mourougane, 2021/08/11}
\enonce{007723}{}
Combien le groupe $\mathcal{S}_5$ contient-il de $5$-Sylow ?
\finenonce{007723}
\finexercice


\exercice{7724, mourougane, 2021/08/11}
\enonce{007724}{Le plan $P(\mathbb{F}^3_2)$}
Déterminer le nombre de points est de droites du plan projectif $P(\mathbb{F}^3_2)$.
Représenter les relations d'incidence.
\finenonce{007724}
\finexercice


\exercice{7725, mourougane, 2021/08/11}
\enonce{007725}{}
\begin{enumerate}
 \item Déterminer le centre de $SL(4,\mathbb{F}_2)$ et calculer l'ordre de $PSL(4,\mathbb{F}_2)$.

 \item Dans $SL(4,\mathbb{F}_2)$ déterminer l'ordre de 
$$A=\left[\begin{array}{cccc}1&0&0&0\\0&1&0&0\\0&0&1&1\\0&0&0&1
    \end{array}\right] 
\textrm{ et de }
B=\left[\begin{array}{cccc}1&1&0&0\\0&1&0&0\\0&0&1&1\\0&0&0&1
    \end{array}\right]$$
Ces deux matrices sont-elles conjuguées dans $SL(4,\mathbb{F}_2)$ ?

 \item Déterminer le centre de $SL(3,\mathbb{F}_4)$ et calculer l'ordre de $PSL(3,\mathbb{F}_4)$.

 \item Le but de cette question est de montrer que toutes les involutions de $PSL(3,\mathbb{F}_4)$ sont des applications projectives associées à des transvections.
Soit $F$ une homographie de $PSL(3,\mathbb{F}_4)$ involutive (i.e. $F^2=Id\in PSL(3,\mathbb{F}_4)$) et $f\in SL(3,\mathbb{F}_4)$ telle que $P(f)=F$. 
\begin{enumerate}
 \item Montrer que $P(f^3)=F$ et que $f^3$ est d'ordre $2$ dans $SL(3,\mathbb{F}_4)$. On notera $g=f^3$.

\item Quelle est la caractéristique de $\mathbb{F}_4$ ?
Montrer que $(g-Id)^2=0$. Déterminer $\dim Ker(g-Id)$ et $\dim Im(g-Id)$. 

\item Conclure.
\end{enumerate}
 \item En étudiant les classes de conjugaison des éléments d'ordre $2$, déterminer si $PSL(4,\mathbb{F}_2)$ et $PSL(3,\mathbb{F}_4)$ sont isomorphes.
\end{enumerate}
\finenonce{007725}
\finexercice

\section{323.00 Anneau, corps}


\exercice{2240, barraud, 2008/04/24}
\enonce{002240}{}
\begin{enumerate}
  \item Trouver 
$$
999\cdot 1998\mod 1999, \qquad
136^7 \mod 137, \qquad
1997\cdot 1998\cdot 1999\cdot 2000 \mod 2001.
$$
  \item Trouver $\quad 2792^{217}\mod 5\quad$ et $\quad 10^{1000} \mod 13$.
\end{enumerate}
\finenonce{002240} 



\finexercice


\exercice{2241, barraud, 2008/04/24}
\enonce{002241}{}
\begin{enumerate}
  \item Examiner  les carr\'es 
$\quad a^2 \mod n$ pour $n=3, \ 4,\ 8$.
  \item  Examiner $\quad a^3 \mod 9\quad$ et $\quad b^4 \mod 16$.
\end{enumerate}
\finenonce{002241} 


\finexercice


\exercice{2242, barraud, 2008/04/24}
\enonce{002242}{}
Passer $\mod n\ $ avec un  module 
appropri\'e et montrer que chacune des \'equations  suivantes
n'a aucune   solution dans $\Zz$ :
\begin{enumerate}
  \item $\quad 3x^2+2=y^2$;
  \item $\quad x^2+y^2=n\quad$ pour $\quad n=2003$, $2004$;
  \item $\quad x^2+y^2+z^2=1999$;
  \item $\quad x^3+y^3+z^3=5$;
  \item $\quad x_1^4+x_2^4+\dots + x_{15}^4=7936$. 
\end{enumerate}
\finenonce{002242} 


\finexercice


\exercice{2243, barraud, 2008/04/24}
\enonce{002243}{}
On dit que $\quad a \mod n$ est inversible si 
il existe $b \mod n$ tel que $ab \equiv 1 \mod n$. 
\begin{enumerate}
  \item Trouver  tous les \'el\'ements inversibles modulo $5$, $6$, $9$, $11$.
  \item   Trouver $\hbox{pgcd} (107, 281)$  et sa representation lin\'eaire
en utilisant {\it l'algorithme d'Euclide}.
  \item Trouver l'inverse de $107\mod 281$ et l'inverse de $281 \mod 107$.
  \item Montrer que $\quad a \mod n$ est inversible ssi $a$ et $n$
sont premiers entre eux.
\end{enumerate}
\finenonce{002243} 


\finexercice


\exercice{2244, barraud, 2008/04/24}
\enonce{002244}{}
Trouver toutes les solutions dans $\Zz$ :
\begin{enumerate}
  \item $\quad 2x+3 \equiv 10 \mod 13$;
  \item $ \quad \begin{cases}
2x+3y\equiv 5 \mod 7\\
5x+2y \equiv 2 \mod 7;
\end{cases}
$
  \item $\quad x^2 +2x+14 \equiv 0 \mod 17$.
\end{enumerate}
\finenonce{002244} 


\finexercice


\exercice{2245, barraud, 2008/04/24}
\enonce{002245}{Le petit th\'eor\`eme de Fermat}
\label{ptfermat}
Soit $p$ un nombre premier et $a$ un nombre premier \`a $p$.
Montrer que :
\begin{enumerate}
  \item $am\equiv an \mod p$ ssi  $m\equiv n \mod p$;
  \item La suite $a,\ 2a, \ 3a,\dots , (p-1)a \mod p$
est une permutation de la suite $1,\ 2, \ 3,\dots , (p-1) \mod p$;
  \item $a^{p-1}\equiv 1 \mod p$.
\end{enumerate}
\finenonce{002245} 


\finexercice


\exercice{2246, barraud, 2008/04/24}
\enonce{002246}{}
\begin{enumerate}
  \item Examiner  $7^n+11^n \mod 19$.
  \item Trouver $2792^{217}\mod 5\quad$ et $\quad 10^{1000} \mod 13$.
  \item Montrer que $13$ divise $2^{70}+3^{70}$ et  
$11$ divise $2^{129}+3^{118}$.
\end{enumerate}
\finenonce{002246} 


\finexercice


\exercice{2247, barraud, 2008/04/24}
\enonce{002247}{Th\'eor\`eme de Wilson}
Soit $p=2m+1$ un nombre premier. Montrer que :
\begin{enumerate}
  \item $\quad (p-1)!\equiv -1 \mod p$;
  \item $\quad (m!)^2\equiv (-1)^{m+1} \mod p$.
\end{enumerate}
\finenonce{002247} 


\finexercice


\exercice{2248, barraud, 2008/04/24}
\enonce{002248}{}
Soit $p>2$ un nombre premier.
\begin{enumerate}
  \item Soit $a$ premier \`a $p$. Supposons que la congruence
$x^2\equiv a\mod p$ poss\`ede une solution.
Montrer que 
$a^{(p-1)/2}\equiv 1 \mod p$.
  \item La congruence $x^2\equiv -1 \mod p$ a une solution 
ssi $p\equiv 1 \mod 4$.
\end{enumerate}
\finenonce{002248} 


\finexercice


\exercice{2249, barraud, 2008/04/24}
\enonce{002249}{}
Donner la d\'efinition  d'un corps.
Les op\'erations binaires $+$ et $\cdot$, sont-elles \'equivalentes dans 
la d\'efinition ?
\finenonce{002249} 


\finexercice


\exercice{2250, barraud, 2008/04/24}
\enonce{002250}{}
Trouver  toutes les solutions des \'equations :
\begin{enumerate}
  \item $ax+b=c\ $ ($a,b,c\in K$, $K$ est un corps);
  \item $2x\equiv 3 \mod 10\ \ $  et  $\ 2x\equiv 6 \mod 10$ dans l'anneau 
$\Zz_{10}=\Zz/10\Zz$.
\end{enumerate}
\finenonce{002250} 


\finexercice


\exercice{2251, barraud, 2008/04/24}
\enonce{002251}{}
Soit $A$ un anneau.  D\'emontrer que :
\begin{enumerate}
  \item $\forall a\in A\ \ 0_A
\cdot a=0_A$;
  \item $(-1_A)\cdot a=-a$;
  \item $|A|\geq 2\ \Longleftrightarrow\ 1_A\neq 0_A$ dans $A$.
\end{enumerate}
\finenonce{002251} 


\finexercice


\exercice{2252, barraud, 2008/04/24}
\enonce{002252}{}
\begin{enumerate}
  \item Si $x\cdot y$ est inversible dans un anneau $A$,
alors $x$ et $y$ sont inversibles.
  \item Dans un anneau, un \'el\'ement inversible n'est pas diviseur de z\'ero
et un diviseur de z\'ero n'est pas inversible.
\end{enumerate}
\finenonce{002252} 



\finexercice


\exercice{2253, barraud, 2008/04/24}
\enonce{002253}{}
\label{ex:bar5}
D\'emontrer que tout anneau int\`egre fini est un corps.
\finenonce{002253} 


\finexercice


\exercice{2254, barraud, 2008/04/24}
\enonce{002254}{}
\label{exodessus}
Lesquels de ces sous-ensembles donn\'es de $\Cc$
sont des anneaux ? Lesquels sont des corps ?
\begin{enumerate}
  \item $\bigcup\limits_{n\in\Nn}10^{-n}\Zz$ ;
  \item $\{\frac{m}{n}\mid m\in\Zz,n\in\Nn^*, (m,n)=1, p\nmid n\}$ ($p$ est
un nombre premier fix\'e) ;
  \item $\Zz[\sqrt{-1}]=\Zz +\Zz\sqrt{-1}$, 
$\ \ \Zz[\sqrt{2}]=\Zz +\Zz\sqrt{2}$; 
  \item $\Qq[\sqrt{-1}]=\Qq +\Qq\sqrt{-1}$, 
$\ \ \Qq[\sqrt{2}]=\Qq +\Qq\sqrt{2}$. 
\end{enumerate}
\finenonce{002254} 


\finexercice


\exercice{2255, barraud, 2008/04/24}
\enonce{002255}{}
Les \'el\'ements inversibles d'un anneau $A$
forment le groupe multiplicatif $(A^{\times}, \cdot)$.
\begin{enumerate}
  \item Trouver $A^{\times}$ pour les anneaux 1. et 2. 
de l'exercice \ref{exodessus}.
  \item Trouver le groupe $\Zz[\sqrt{-1}]^{\times}$ en utilisant 
la norme complexe.
  \item Montrer que le groupe $\Zz[\sqrt{2}]^{\times}$ est infini.
\end{enumerate}
\finenonce{002255} 


\finexercice


\exercice{2256, barraud, 2008/04/24}
\enonce{002256}{}
Un \'el\'ement $a$ d'un anneau $A$ s'appelle nilpotent, s'il existe $n\in \Nn$ tel que $a^n=0$. Trouver tous les \'el\'ements inversibles, les diviseurs de z\'ero, les nilpotents des anneaux suivants :
\begin{enumerate}
  \item $\Zz /360\Zz$;
  \item $\Zz /n\Zz$;
  \item D\'emontrer que, pour tout nilpotent $x$ de $A$, l'\'el\'ement $1+x$ est inversible.
\end{enumerate}
\finenonce{002256} 


\finexercice


\exercice{2257, barraud, 2008/04/24}
\enonce{002257}{}
Soit $I$ un id\'eal d'un anneau $A$. On note par $(a)=a\cdot A$ l'id\'eal principal engendr\'e par $a$. Montrer que :
\begin{enumerate}
  \item $I=A$ si et seulement si $I$ contient une unit\'e;
  \item $(a)=A$ ssi $a$ est inversible;
  \item Un anneau $A$ est  un corps ssi  $(0)$ est le seul id\'eal propre de $A$.
\end{enumerate}
\finenonce{002257} 


\finexercice


\exercice{2258, barraud, 2008/04/24}
\enonce{002258}{}
 Montrer  que les \'el\'ements nilpotents d'un anneau
forment un id\'eal.
\finenonce{002258} 


\finexercice


\exercice{2259, barraud, 2008/04/24}
\enonce{002259}{Sommes et produits d'id\'eaux}
\begin{enumerate}
  \item Soient $I$, $J$ deux id\'eaux d'un anneau $A$. Montrer que
$$I\cap J,\qquad  I+J=\{x+y\,|\,x\in I, y\in J\}$$
sont des id\'eaux de $A$.
  \item Montrer que $I+J$ est le plus petit id\'eal de $A$ contenant $I$ et $J$.
  \item Soit $n,m\in \Zz$, $I=(n)=n\Zz$, $J=(m)=m\Zz$. Trouver $I\cap J$ et  $I+J$. 
  \item Montrer que 
$$
I\cdot J=\{x_1y_1+x_2y_2+\dots x_ny_n\,|\ n\in \Nn,\  x_k\in I,\ 
y_k\in J\ \text{pour\ }\  1\le k\le n\}
$$ 
est un id\'eal. Il s'appelle {\it produit des id\'eaux} $I$ et $J$.
  \item  On consid\`ere les id\'eaux $I=(x_1,\dots x_n)=
Ax_1+\dots+ Ax_n$ et $J=(y_1,\dots y_m)= Ay_1+\dots+ Ay_m$.
D\'ecrire les id\'eaux $I+J$, $I\cdot J$, $I^2$ en fonction de  
$x_k$, $y_l$.
\end{enumerate}
\finenonce{002259} 


\finexercice


\exercice{2260, barraud, 2008/04/24}
\enonce{002260}{Id\'eaux \'etrangers}
\begin{enumerate}
  \item Montrer que $I\cdot J\subset I\cap J$ et $(I+J)\cdot (I\cap J)\subset I\cdot J$
  \item On dit que deux id\'eaux $I$ et $J$ de $A$ sont {\it \'etrangers}
si $I+J=A$. Montrer que $I\cap J$=$I\cdot J$ si $I$, $J$ sont \'etrangers.
\end{enumerate}
\finenonce{002260} 


\finexercice


\exercice{2300, barraud, 2008/04/24}
\enonce{002300}{}
Soient $A$ un anneau et $I$ et $J$ les  idéaux de $A$ tels que $I+J=(1)$. Démontrer que $I^n+J^m=(1)$ quels que soient entiers positifs non-nuls $n$ et $m$.
\finenonce{002300} 


\finexercice


\exercice{2301, barraud, 2008/04/24}
\enonce{002301}{}
Trouver toutes les solutions des syst\`emes suivantes :
\begin{enumerate}
  \item $ \quad \begin{cases}
x\equiv 1 \mod 3\\
x \equiv 3 \mod 5\\
x \equiv 4 \mod 7\\
x\equiv 2 \mod 11
\end{cases}
$
  \item $ \quad \begin{cases}
x\equiv 997\mod 2001\\
x \equiv 998 \mod 2002\\
x \equiv 999 \mod 2003
\end{cases}.
$
\end{enumerate}
\finenonce{002301} 


\finexercice


\exercice{2302, barraud, 2008/04/24}
\enonce{002302}{}
D\'emontrer que les anneaux suivants sont isomorphes
$$
\Zz/72\Zz\times \Zz/84\Zz\cong \Zz/36\Zz\times \Zz/168\Zz.
$$
\finenonce{002302} 


\finexercice


\exercice{2303, barraud, 2008/04/24}
\enonce{002303}{}
\begin{enumerate}
  \item Montrer que $20^{15}-1$ est divisible par $11\times 31\times 61$.
  \item Trouver le reste de la division de $2^{6754}$  par $1155$.
\end{enumerate}
\finenonce{002303} 


\finexercice


\exercice{2304, barraud, 2008/04/24}
\enonce{002304}{}
\begin{enumerate}
  \item Quels sont les restes des division de $10^{100}$ par $13$ et par $19$ ?
  \item Quel est le reste de la  division de $10^{100}$ par $247=13\cdot 19$ ? 
En d\'eduire que  $10^{99}+1$ est multiple de $247$.
\end{enumerate}
\finenonce{002304} 


\finexercice


\exercice{2305, barraud, 2008/04/24}
\enonce{002305}{}
 Soit $C=A\times B$ le produit direct de deux anneaux.
D\'ecrire les ensembles des \'el\'ements inversibles, des diviseurs
de z\'ero et des \'el\'ements nilpotents de l'anneau $C$.
\finenonce{002305} 


\finexercice


\exercice{2306, barraud, 2008/04/24}
\enonce{002306}{}
\begin{enumerate}
  \item D\'etermin\'er la structure des anneaux 
quotients suivants:
$$ \Zz_2[x]/(x^3+x^2+x+1),\quad\Zz[x]/(x^2-1),\quad  \Qq[x]/(x^8-1). $$
  \item Consid\'erons l'anneau quotient $K[x]/(f^ng^m)$
o\`u $f$ et  $g$ sont deux polyn\^omes  distincts irr\'eductibles sur
le corps $K$. D\'ecrirer  les diviseurs  de z\'ero et 
les \'el\'ements nilpotents de l'anneau  $K[x]/(f^ng^m)$.
  \item Quels id\'eaux a-t-il cet anneau ?
  \item Soit $K$ le corps fini \`a $p$ \'el\'ements.
Trouver le nombre des \'el\'ements du groupe multiplicatif
de l'anneau $K[x]/(f^mg^l)$.
  \item Donner une g\'en\'eralisation de la question 4) dans le cas
du produit de  $n$ polyn\^omes irr\'eductibles sur un corps fini $K$
\`a $q$ \'el\'ements.
\end{enumerate}
\finenonce{002306} 


\finexercice


\exercice{2307, barraud, 2008/04/24}
\enonce{002307}{}
 Trouver les facteurs multiples des polyn\^omes suivants :
\begin{enumerate}
  \item $x^6-15x^4+8x^3+51x^2-72x+27$;
  \item $x^6-2x^5-x^4-2x^3+5x^2+4x+4$.
\end{enumerate}
\finenonce{002307} 


\finexercice


\exercice{2308, barraud, 2008/04/24}
\enonce{002308}{}
Trouver le polyn\^ome   $f\in \Zz[x]$ du derg\'e le plus petit tel que  
$$
\begin{cases}
f\equiv 2x \mod (x-1)^2\\
f \equiv 3x\mod (x-2)^3
\end{cases}.
$$
\finenonce{002308} 


\finexercice


\exercice{2309, barraud, 2008/04/24}
\enonce{002309}{}
Soit $\sqrt{d}$ non rationel. Dans l'anneau 
$$
{\mathbb{Z}}[\sqrt{d}]=\{n+m\sqrt{d}\,|\,n,m\in \mathbb{Z}\}
$$
on definit la ``conjugaison" $\bar{z}$ :
\medskip

\centerline{si $z=n+m\sqrt{d}$, alors $\bar{z}=n-m\sqrt{d}$.}
\medskip

On peut aussi d\'efinir la norme  
$N_d:{\mathbb Z}[\sqrt{d}] \rightarrow{\mathbb Z}$ 
par $N_d(z) = z\bar{z}=(n+m\sqrt{d})(n-m\sqrt{d})$.
\medskip

0.  Montrer que les aplications $\bar{z}$ et $N(z)$ sont multiplicatives :
$$
\overline{z_1\cdot z_2}=\bar{z_1}\cdot \bar{z_2},
\qquad N_d(z_1\cdot z_2)=N_d(z_1)\cdot N_d(z_2).
$$
\finenonce{002309} 


\finexercice


\exercice{2310, barraud, 2008/04/24}
\enonce{002310}{}
\begin{enumerate}
  \item Montrer que $z \in \mathbb{Z}[\sqrt{d}]$ est inversible ssi $N_d(z)
   = \pm 1$.  D\'eterminer les \'el\'ements inversibles de $\mathbb{Z}[\sqrt{-5}]$.
  \item  Montrer que si $N_d(z) = \pm p$, o\`u $p$ est un premier, alors
   $z$ est irr\'eductible dans $\mathbb Z[\sqrt{d}]$.  Donner quelques
   exemples d'\'el\'ements irreductibles dans $\mathbb{Z}[\sqrt{d}]$ pour
   $d=-1$, $2$, $-6$, $p$, o\`u $p$ un premier.
  \item On note $A=\mathbb{Z}[\sqrt{-5}]$.  Montrer que $3$ et $2+\sqrt{-5}$
   sont irr\'eductibles dans $A$.
  \item Trouver tous les irr\'eductibles de $A$ de norme $9$.
  \item Trouver tous les diviseurs de $9$ et de $3(2+\sqrt{-5})$ dans
  l'anneau $A$ \`a association pr\`es.
  \item Trouver un $pgcd\,(3, 2+\sqrt{-5})$, et montrer que $3$ et
   $2+\sqrt{-5}$ n'ont pas de $ppcm$ dans l'anneau $A$.
  \item Montrer que l'id\'eal $I = (3, 2+\sqrt{-5}) \subset A$ n'est pas
  principal.  Donc l'anneau $A$ n'est pas principal. Est-il factoriel
  ?
  \item Montrer que $9$ et $3(2+\sqrt{-5})$ n'ont pas de $pgcd$ dans $A$.
  Poss\`edent-ils un $ppcm$ ?
\end{enumerate}
\finenonce{002310} 


\finexercice


\exercice{2311, barraud, 2008/04/24}
\enonce{002311}{}
Soit $\Zz_{36}=\Zz/36\Zz\ $ l'anneau des entiers modulo $36$.
\begin{enumerate}
  \item D\'ecrire tous les \'el\'ements inversibles, tous les diviseurs
de z\'ero et tous les \'el\'ements nilpotents de l'anneau $\Zz_{36}$.
({\it Un \'el\'ement $a$ d'un anneau $A$ est dit nilpotent
si il existe $n$ tel que $a^n=0$.})
  \item Trouver tous les id\'eaux de l'anneau $\Zz_{36}$.
  \item Soit $A$ un anneau arbitraire. Montrer que
  $$
  (a\in A^{\times} \hbox{ et } b\in A^{\times}) \Longleftrightarrow
  (a\cdot b)\in A^{\times}.
  $$
  \item Donner un exemple d'un polyn\^ome inversible de degr\'e $1$ sur
  $\Zz_{36}$.
  \item D\'ecrire tous les \'el\'ements inversibles de l'anneau
  $\Zz_{36}[x]$.
\end{enumerate}
\finenonce{002311} 


\finexercice


\exercice{2312, barraud, 2008/04/24}
\enonce{002312}{}
 Montrer que les polyn\^omes suivantes sont irr\'eductibles dans
  $\Zz[x]$ :
\begin{enumerate}
  \item $P=x^{2004}+4x^{2002}+2000x^4+2002$;
  \item $Q=x^6+6x^5+12x^4+12x^3+3x^2+6x+25$.
\end{enumerate}
\finenonce{002312} 


\finexercice


\exercice{2313, barraud, 2008/04/24}
\enonce{002313}{}
Soit $p$ un nombre premier impair. Montrer que la congruence
  \newline\noindent $x^2\equiv-1 \mod p\ $ a une solution si et
  seulement si $p\equiv 1 \mod 4$.
\finenonce{002313} 


\finexercice


\exercice{2314, barraud, 2008/04/24}
\enonce{002314}{}
 Soient $f=x^6+x^5+x^4+x^3+1\in \Zz_2[x]\ $, $g=x^3+x^2+1 \in
  \Zz_2[x]$ deux polyn\^omes sur le corps $\Zz_2$.
\begin{enumerate}
  \item En utilisant l'algorithme d'Euclide trouver le p.g.c.d. de $f$ et
  $g$ et sa repr\'esentation lin\'eaire.
  \item Les polyn\^omes $f$ et $g$ sont-ils irr\'eductibles ?
  \item Soit $(g)$ l'id\'eal principal engendr\'e par $g$.  Combien
  d'\'el\'ements contient l'anneau quotient $A=\Zz_2[x]/(g)$ ?
  \item Soit $\pi\,:\,\Zz_2[x]\to A$ la projection canonique.  On pose
  $\pi(x)=\bar x\in A$. Trouver l'inverse de l'\'el\'ement $\pi(f)$
  dans l'anneau quotient $A$.
  \item L'anneau quotient $B=\Zz_2[x]/(f)$ est-il un corps ?  Justifier
  la r\'eponse, i.e. donner une d\'emonstration si $B$ est un corps ou
  trouver un \'el\'ement non-inversible dans $B$ dans le cas
  contraire.
\end{enumerate}
\finenonce{002314} 


\finexercice


\exercice{6502, drutu, 2011/10/16}
\enonce{006502}{}
 Les ensembles suivants sont-ils des anneaux vis-à-vis des opérations usuelles d'addition et de multiplication? Sont-ils des corps ?
$$
(a)\; \left\{ \frac{p}{q}\in \Q \mid q\in \{1,\, 2,\, 4 \} \right\}\; ;\hspace{0.5cm}(b)\; \left\{ \frac{p}{q}\in \Q \mid q\in \{ 2^n\; ; \; n\in \N \} \right\}\; ; \hspace{0.5cm}(c)\; \Z +\sqrt{5}\Z\; ; 
$$
$$
(d)\; \Z +\sqrt[3]{5}\Z \; ;\hspace{0.5cm}(e)\; \Z +\sqrt[3]{2}\Z +\sqrt[3]{4}\Z; 
$$
\finenonce{006502}



\finexercice


\exercice{6503, drutu, 2011/10/16}
\enonce{006503}{}
\begin{enumerate}
\item Montrer qu'il existe un morphisme d'anneaux et un seul de $\Z /4\Z $ à $\Z /2\Z $, de $\Z /12\Z $ à $\Z /3\Z $, de $\Z /12\Z $ à $\Z /4\Z $.

\item Montrer qu'il n'existe pas de morphisme d'anneaux de $\Z /3\Z $ à $\Z /4\Z $.

\item $m$ et $n$ étant des entiers positifs, trouver des conditions pour qu'il existe un morphisme d'anneaux de $\Z /m\Z $ à $\Z /n\Z $. 
\end{enumerate}
\finenonce{006503}



\finexercice


\exercice{6504, drutu, 2011/10/16}
\enonce{006504}{}
Montrer que $\Z /4\Z \times \Z /4\Z $ a exactement trois sous-anneaux.
\finenonce{006504}



\finexercice


\exercice{6505, drutu, 2011/10/16}
\enonce{006505}{}
 Dans $\Z /10\Z $ montrer que les multiples de 5 forment un anneau isomorphe à $\Z /2\Z $ qui n'est pas un sous-anneau de $\Z /10\Z $. 
\finenonce{006505}



\finexercice


\exercice{6506, drutu, 2011/10/16}
\enonce{006506}{}
Soit $M_n(F)$ l'anneau des matrices sur un corps $F$.
\begin{enumerate}
\item Montrer que si $X$ est une matrice non-dégénérée, alors $XY\neq 0$ et $YX\neq 0,\; \forall Y\neq 0$.

\item Soit $X\in M_n(F)$ et soit l'application linéaire $f_X:F^n\to F^n,\; f_X(v):=Xv$. Trouver une écriture de la relation $XY=0,\; X,\; Y\in M_n(F)$ en termes de noyaux et images des applications $f_X,\; f_Y$. Même question pour la relation $YX=0$.
En déduire que toute matrice dégénérée non-nulle est un diviseur de 0.

\item Montrer que tout idéal bilatère de $M_n(F)$ contient, avec une matrice de rang $r$, toutes les matrices diagonales de rang $r$.

\item Montrer qu'un idéal qui contient une matrice non-dégénérée co\"{\i}ncide avec $M_n(F)$. 

\item Montrer que les seuls idéaux de $M_n(F)$ sont $0$ et $M_n(F)$. 
\end{enumerate}
\finenonce{006506}



\finexercice


\exercice{6507, drutu, 2011/10/16}
\enonce{006507}{}  
\begin{enumerate}
\item Montrer que l'image réciproque d'un idéal par un morphisme d'anneaux est un idéal.

\item Montrer par un contre-exemple que l'image d'un idéal par un morphisme d'anneaux n'est pas toujours un idéal. Montrer que l'affirmation précédente est toutefois vraie si le morphisme est surjectif.

\item Qu'est-ce qu'on peut dire sur l'image réciproque et l'image d'un idéal premier/maximal par un morphisme d'anneaux eventuellement surjectif ?  
\end{enumerate}
\finenonce{006507}



\finexercice


\exercice{6508, drutu, 2011/10/16}
\enonce{006508}{}
Trouver les idéaux maximaux de $\Z /n\Z$. 
\finenonce{006508}



\finexercice


\exercice{6509, drutu, 2011/10/16}
\enonce{006509}{}
Montrer que tout anneau intègre fini est un corps.
\finenonce{006509}



\finexercice


\exercice{6510, drutu, 2011/10/16}
\enonce{006510}{}
Soit $A$ un anneau principal. Montrer que $(a_1,\dots ,a_n)=(d)$ si et seulement si 
$d\sim \pgcd (a_1,\dots ,a_n)$.
\finenonce{006510}



\finexercice


\exercice{6511, drutu, 2011/10/16}
\enonce{006511}{}
Montrer que tout morphisme de corps est injectif.
\finenonce{006511}



\finexercice


\exercice{6512, drutu, 2011/10/16}
\enonce{006512}{}
 Soit $p$ un nombre premier et $A$ un anneau intègre de caractéristique $p$. 
\begin{enumerate}
\item Montrer que $p\cdot a=0$ pour tout $a\in A$.

\item Montrer que $p| C^k_p,\; \forall k=1,\dots p-1$, et en déduire que l'application $F_p :A\to A,\; F_p(a)=a^p$, est un endomorphisme d'anneaux. On appelle $F_p$ {\it l'endomorphisme de Frobenius}.

\item Montrer que $F_p(a)=a$ pour tout $a$ dans le plus petit sous-anneau $A_0$ de $A$ contenant $1$.

\item Montrer que $F_p$ est un automorphisme si $A$ est fini.

\item Montrer que $(\Sigma a_ib_i)^{p^k}=\Sigma a_i^{p^k}b_i^{p^k},\; \forall a_i,\; b_i\in A,\; k\in \N^*$.
\end{enumerate}
\finenonce{006512}



\finexercice


\exercice{6513, drutu, 2011/10/16}
\enonce{006513}{}
 Soit $A$ un anneau intègre et $a,b,c \in A\setminus \{0 \}$.
 Montrer que, chaque fois que les pgcd suivants existent, on a les égalités :
\begin{enumerate}
\item  $\pgcd (ca,cb)\sim c \pgcd (a,b)$

\item $\pgcd ( \pgcd (a,b),c)\sim \pgcd (a, \pgcd (b,c))$.
  
\medskip

 Si $A$ est en plus factoriel, montrer que 
\item $\pgcd (a,b)\sim 1$ et $\pgcd(a,c)\sim 1$ implique $\pgcd(a,bc)\sim 1$.
 
\item Si $a|bc$ et $\pgcd (a,b)\sim 1$ alors $a|c$.
  
\item Si $b|a$ et $c|a$ et $\pgcd (b,c)\sim 1$ alors $(bc)|a$.  
\end{enumerate}
\finenonce{006513}



\finexercice


\exercice{6514, drutu, 2011/10/16}
\enonce{006514}{}
Montrer que dans $\Z[i]$ $3$ est premier et 2 ne l'est pas. 
\finenonce{006514}



\finexercice


\exercice{6515, drutu, 2011/10/16}
\enonce{006515}{}
 Montrer que $\Z[i\sqrt{3}]$ et $\Z[i\sqrt{5}]$ ne sont pas des anneaux factoriels.
\finenonce{006515}



\finexercice


\exercice{6516, drutu, 2011/10/16}
\enonce{006516}{}
\begin{enumerate}
\item Montrer que $\Z[X]$ est un anneau factoriel.

\item  Dans $\Z [X]$ montrer que l'ensemble des polyn\^omes de terme constant un nombre pair est un idéal mais non-principal. En déduire que $\Z[X]$ est un anneau non-principal.

\item Montrer que l'idéal $(X)$ est premier mais non maximal dans $\Z[X]$.
\end{enumerate}
\finenonce{006516}



\finexercice


\exercice{6517, drutu, 2011/10/16}
\enonce{006517}{}
 Montrer que l'anneau des entiers $\mathfrak{o}_{19}$ du corps quadratique $\Q[i\sqrt{19}]$ est principal mais non-euclidien.
\finenonce{006517}



\finexercice


\exercice{6518, drutu, 2011/10/16}
\enonce{006518}{}
Soit $A$ un anneau commutatif et soit 
$$
Nil\; (A):=\{a\in A \mid \exists n\in N,\; a^n=0 \}\; .
$$
\begin{enumerate}
\item Montrer que $Nil\; (A)$ est un idéal de $A$.

\item Montrer que si $P$ est un idéal premier de $A$, alors $Nil\; (A)\subset P$.

\item Montrer que pour tout $x\not\in Nil\; (A)$ il existe un idéal premier $P$ 
tel que $x\not\in P$ ({\it Indication}: Utiliser le théorème de Zorn).

\item En déduire que 
$$
Nil\; (A)=\bigcap_{P\mbox{ premier }}P\; .
$$  
\end{enumerate}
\finenonce{006518}



\finexercice


\exercice{6519, drutu, 2011/10/16}
\enonce{006519}{}
\begin{enumerate}
\item Montrer que $\Q (\sqrt{d}),\; d=7,\; 11$, est un corps.

\item Montrer que $a+b\sqrt{7}\to a+b\sqrt{11}$ n'est pas un isomorphisme de corps entre $\Q (\sqrt{7})$ et $\Q (\sqrt{11})$.

\item Montrer que les corps $\Q (\sqrt{7})$ et $\Q (\sqrt{11})$ ne sont pas isomorphes. 
\end{enumerate}
\finenonce{006519}



\finexercice


\exercice{6520, drutu, 2011/10/16}
\enonce{006520}{}
 $F$ et $F'$ étant deux corps, montrer que $F\times F'$ a seulement deux idéaux non-triviaux, et que c'est un anneau principal.
\finenonce{006520}



\finexercice


\exercice{6521, drutu, 2011/10/16}
\enonce{006521}{}
 Montrer que $\{\frac{p}{q}\in \Q \mid q \mbox { impair }\}$ est un anneau principal.
\finenonce{006521}



\finexercice


\exercice{6522, drutu, 2011/10/16}
\enonce{006522}{}
 Résoudre les congruences simultanées suivantes :
\begin{enumerate}
\item $x\equiv 2$ (mod 7), $x\equiv 3$ (mod 6) ;

\item $3x\equiv 2$ (mod 5), $x\equiv 6$ (mod 7), $x\equiv 1$ (mod 6). 
\end{enumerate}
\finenonce{006522}



\finexercice


\exercice{6523, drutu, 2011/10/16}
\enonce{006523}{}
Montrer que deux congruences dans $\Z $ de la forme 
$$
mx\equiv c\; {\rm{(mod }}\; a{\rm{)}},\; nx\equiv d\; {\rm{(mod }}\; b{\rm{)}}
$$
 ont une solution commune $x\in \Z$ quand les coefficients vérifient les conditions $\pgcd (a,b)=1$, 
$\pgcd(a,m)=1$,  $\pgcd(n,b)=1$.   
\finenonce{006523}



\finexercice


\exercice{6524, drutu, 2011/10/16}
\enonce{006524}{}
\begin{enumerate}
\item Pour les anneaux de polyn\^omes sur un corps $F$ montrer que $F[x]/(x^2-1)\simeq F[x]/(x^2-4),\; \; F[x]/(x^2+1)\simeq F[x]/(x^2+2x+2)\, ,$ $F[x,y]/(x+y)\simeq F[x],\; F[x]/(x+1)\simeq F$.


\item Montrer que $\Q[x]/(x^2-2)\simeq  \Q[\sqrt{2}],\; \Z[x]/(x^2+1)\simeq  \Z[i].$ Déduire du dernier isomorphisme que $(x^2+1)$ est un idéal premier non-maximal. 
\end{enumerate}
\finenonce{006524}



\finexercice


\exercice{6525, drutu, 2011/10/16}
\enonce{006525}{}
 Si $M$ est un idéal maximal d'un anneau intègre $A$ montrer que l'anneau local $A_M$ a exactement un idéal maximal.
\finenonce{006525}



\finexercice


\exercice{7731, mourougane, 2021/08/11}
\enonce{007731}{Polarisation} 
\begin{enumerate}
 \item Montrer que l'anneau quotient $\mathbb{F}_3[X]/X^2+1$ est un corps que l'on notera $\mathbb{F}_{9}$. On notera $\xi:=[X]$ la classe du polynôme $X$. 
Montrer que $(1,\xi)$ est une base de $\mathbb{F}_{9}$ comme $\mathbb{F}_{3}$-espace vectoriel. On définit les applications $\ell_i : \mathbb{F}_{9} \rightarrow \mathbb{F}_{3}$ pour $i=1,2$ par la condition $$\forall \lambda \in \mathbb{F}_{9},\ \ \ \lambda=\ell_1(\lambda)+\ell_2(\lambda)\xi.$$

  \item Soit $\sigma : \mathbb{F}_{9}\mapsto\mathbb{F}_{9}$ défini par $\sigma(x)=x^3$.
Montrer que $\sigma$ est un automorphisme de corps de $\mathbb{F}_{9}$ d'ordre $2$. Calculer $\sigma(\xi)$.

 \item Montrer que, pour tout $\lambda \in  \mathbb{F}_{9} $, 
 \begin{eqnarray*}
\ell_2(\lambda)&=&-\ell_1(\xi\lambda)\\
2\ell_1(\lambda)&=&\lambda+\sigma(\lambda) \\
2\xi \ell_2(\lambda)&=&\lambda-\sigma(\lambda).
\end{eqnarray*}


 \item {\itshape (Polarisation)}
Soit $E$ un $\mathbb{F}_{9}$-espace vectoriel.
Soit $h$ une forme $\sigma$-hermitienne sur $E$. 
Soit $q$ la forme quadratique associée (i.e. pour $x\in E$, $q(x)=h(x,x)$). Soit $(x,y)\in E^2$.
Calculer $\ell_1(f(x,y))$ en fonction de $q$
et en déduire $f(x,y)$ en fonction de $q$.

  \item En déduire que les sous-espaces totalement isotropes pour $h$ sont exactement les sous-espaces vectoriels de $E$ contenus dans le cône isotrope de $h$. 
\end{enumerate}
\finenonce{007731}
\finexercice


\exercice{7736, mourougane, 2021/08/11}
\enonce{007736}{}
\begin{enumerate}
 \item On considère l'application $\sigma~:~\mathbb{F}_{25}\to \mathbb{F}_{25}, \lambda\mapsto\lambda^5$. Montrer que c'est un automorphisme involutif du corps $\mathbb{F}_{25}$.
 \item Les formes suivantes sur $\mathbb{F}_{25}^3$ sont-elles $\sigma$-sesquilinéaires ?
\begin{eqnarray*}
 f_1( 
\left(\begin{array}{c}x\\ y\\ z\end{array}\right),
\left(\begin{array}{c}x'\\ y'\\ z'\end{array}\right)
)
&=&x(x')^5+3z(y')^5+3y(z')^5.\\
f_2( 
\left(\begin{array}{c}x\\ y\\ z\end{array}\right),
\left(\begin{array}{c}x'\\ y'\\ z'\end{array}\right)
)
&=&x^5(x')^5+x^5(y')^5+y^5(x')^5.\\
f_3(
\left(\begin{array}{c}x\\ y\\ z\end{array}\right),
\left(\begin{array}{c}x'\\ y'\\ z'\end{array}\right)
)
&=&3x(x')^5+z(y')^5+y(z')^5.
\end{eqnarray*}
 \item Parmi les formes $\sigma$-sesquilinéaires précédentes, lesquelles sont équivalentes ?
\end{enumerate}
\finenonce{007736}
\finexercice


\exercice{7737, mourougane, 2021/08/11}
\enonce{007737}{}
Soit $E$ un $\mathbb{F}_{3^2}$ espace vectoriel. Soit $\sigma : \mathbb{F}_{3^2} \mapsto \mathbb{F}_{3^2}$ défini par $\sigma(x)=x^3$.
\begin{enumerate} 


 \item Montrer que $\sigma$ est un automorphisme de corps de $\mathbb{F}_{3^2}$ d'ordre $2$.

Soit $h$ une forme $\sigma$-hermitienne sur $E$. 
 \item Montrer qu'il existe un élément $\xi \in \mathbb{F}_{3^2}$ tel que $\xi^2+1=0$ et $\mathbb{F}_{3^2}=\mathbb{F}_{3}[\xi]$.
 \item Montrer $\xi^3=-\xi$.

On définit les applications $\ell_i : \mathbb{F}_{3^2} \rightarrow \mathbb{F}_{3}$ pour $i=1,2$ par la condition $$\forall x \in \mathbb{F}_{3^2}, x=\ell_1(x)+\ell_2(x)\xi.$$
 \item Montrer que, pour tout $x \in  \mathbb{F}_{3^2} $, $x+\sigma(x)=\ell_1(x)$ et $x-\sigma(x)=\xi \ell_2(x)$.


Soit $W$ un sous-espace vectoriel de $E$. Soient $x, y \in W$.
  \item {\itshape (Polarisation)} Décrire $h(x,y)$ comme polynôme à coefficients dans $\mathbb{F}_{3^2}$ en les $(h(u,u))_{u \in W}$. 
  \item En déduire que les sous-espaces totalement isotropes pour $h$ sont exactement les sous-espaces vectoriels de $E$ contenus dans le cône isotrope de $h$. 

\end{enumerate}
\finenonce{007737}
\finexercice

\section{324.00 Polynôme}


\exercice{2261, barraud, 2008/04/24}
\enonce{002261}{}
\begin{enumerate}
  \item Soit $A$ un anneau quelconque. Alors l'anneau de polyn\^omes $A[x]$ n'est pas un corps.
  \item Montrer que pour un anneau int\`egre $A$, les polyn\^omes unitaires
lin\'eaires de $A[x]$ sont irr\'eductibles.
  \item D\'ecrire tous les polyn\^omes irr\'eductibles de $\Cc[x]$ et 
de $\Rr[x]$.
  \item D\'emontrer que pour tout  corps $K$, l'anneau de polyn\^omes 
$K[x]$ a une infinit\'e de polyno\^mes unitaires irr\'eductibles.
\end{enumerate}
\finenonce{002261} 


\finexercice


\exercice{2262, barraud, 2008/04/24}
\enonce{002262}{}
\begin{enumerate}
  \item Montrer que l'id\'eal 
$(x, n)$ o\`u $n\in \Zz$, $n>1$ de l'anneau $\Zz[x]$ n'est pas principal.
  \item Soit $A$ un anneau int\`egre. Montrer que
$A[x]$  est principal ssi $A$ est un corps.
\end{enumerate}
\finenonce{002262} 


\finexercice


\exercice{2263, barraud, 2008/04/24}
\enonce{002263}{}
Soit $f(x)\in A[x]$ un polyn\^ome
sur un anneau $A$. Supposons que $(x-1)\,|\,f(x^n)$. Montrer que
$(x^n-1)\,|\,f(x^n)$.
\finenonce{002263} 


\finexercice


\exercice{2264, barraud, 2008/04/24}
\enonce{002264}{}
Pour $n,m\ge 2$, d\'eterminer le reste de la division
euclidienne du polynôme \mbox{$(x-2)^m+(x-1)^n-1$} par $(x-1)(x-2)$ 
dans $\Zz[x]$.
\finenonce{002264} 


\finexercice


\exercice{2265, barraud, 2008/04/24}
\enonce{002265}{}
\label{exoprec}
\begin{enumerate}
  \item Si $K$ est un corps, montrer qu'un polyn\^ome $P$ de degr\'e 2 ou 3 
dans $K [x]$ est irr\'eductible si et seulement si il n'a pas 
de z\'ero dans $K$. 
  \item Trouver tous les polyn\^omes irr\'eductibles
de degr\'e $2$, $3$ \`a coefficients dans  $\Zz/2\Zz$.
  \item En utilisant la partie pr\'ec\'edente, montrer que 
les polyn\^omes $5x^3+8x^2+ 3x+ 15$ et  $x^5+2x^3+3x^2-6x-5$ sont
irr\'eductibles dans $\Zz[x]$.
  \item D\'ecrire tous les polyn\^omes irr\'eductibles de degr\'e $4$
et $5$ sur $\Zz/2\Zz$.
\end{enumerate}
\finenonce{002265} 


\finexercice


\exercice{2266, barraud, 2008/04/24}
\enonce{002266}{}
\begin{enumerate}
  \item Trouver tous les polyn\^omes irr\'eductibles
de degr\'e $2$, $3$ \`a coefficients dans le corps $\mathbb{F}_3=\Zz/3\Zz$.
  \item D\'ecomposer les polyn\^omes suivants en facteurs irr\'eductibles dans 
$\mathbb{F}_3[x]$.
$$
x^2+x+1,\hspace{1cm}x^3+x+2,\hspace{1cm}x^4+x^3+x+1\; . 
$$
\end{enumerate}
\finenonce{002266} 


\finexercice


\exercice{2267, barraud, 2008/04/24}
\enonce{002267}{}
En utilisant les r\'eductions $\!\mod 2$ ou $\!\mod 3$
montrer que les polyn\^omes
\mbox{$x^5-6x^3+2x^2-4x+5$},
\ \ \mbox{$7x^4+8x^3+11x^2-24x-455$}
sont irr\'eductibles dans $\Zz[x]$.
\finenonce{002267} 


\finexercice


\exercice{2268, barraud, 2008/04/24}
\enonce{002268}{}
 Soient
$$
f(x)=(x-a_1)(x-a_2)\dots (x-a_n)-1,\quad
g(x)=(x-a_1)^2(x-a_2)^2\dots (x-a_n)^2+1
$$
o\`u $a_1,\dots a_n\in \Zz$ soient deux \`a deux distincts.
Montrer que  $f$ et $g$ sont irr\'eductibles dans $\Qq[x]$.
\finenonce{002268} 


\finexercice


\exercice{2269, barraud, 2008/04/24}
\enonce{002269}{}
\label{ex:bar9}
Soient $f,g\in \Qq[x]$. Supposons que
$f$ soit  irr\'eductible et qu'il existe $\alpha \in \Cc$ tel que
$f(\alpha)=g(\alpha)=0$. Alors $f$ divise $g$.
\finenonce{002269} 


\finexercice


\exercice{2270, barraud, 2008/04/24}
\enonce{002270}{}
Pour quel $n$, $m$ dans $\Zz$ la fraction
$$
\frac{11n+2m}{18n+5m}
$$
est r\'eductible ?
\finenonce{002270} 


\finexercice


\exercice{2271, barraud, 2008/04/24}
\enonce{002271}{}
Trouver le $\pgcd(x^n-1,\ x^m-1)$ dans
$\Zz[x]$.
\finenonce{002271} 


\finexercice


\exercice{2272, barraud, 2008/04/24}
\enonce{002272}{}
 Trouver le $\pgcd(f,g)$ dans $\Zz_2[x]$ 
et sa repr\'esentation lin\'eaire 
$fu+gv$ o\`u $d, u,v\in \Zz_2[x]$ :

\begin{enumerate}
  \item $$ f=x^5+x^4+1,\qquad g=x^4+x^2+1;
$$
  \item $$\quad f=x^5+x^3+x+1,\qquad g=x^4+1.
$$
\end{enumerate}
\finenonce{002272} 


\finexercice


\exercice{2273, barraud, 2008/04/24}
\enonce{002273}{}
 Trouver le $\pgcd(f,g)$ dans $\Zz_3[x]$ et 
$\Zz_5[x]$
de $f=x^4+1$,  $g=x^3+x+1$.
\finenonce{002273} 


\finexercice


\exercice{2274, barraud, 2008/04/24}
\enonce{002274}{}
Trouver le $\pgcd(f,g)$ dans $\Zz[x]$ de 
$f=x^4+x^3-3x^2-4x-1$ et  $g=x^3+x^2-x-1$.
\finenonce{002274} 


\finexercice


\exercice{2275, barraud, 2008/04/24}
\enonce{002275}{}
Montrer que $f$ est irr\'eductible dans 
$\Qq[x]$ :
\begin{enumerate}
  \item $f=x^4-8x^3+12x^2-6x+2$;
  \item $f=x^5-12x^3+36x-12$;
  \item $f=x^4-x^3+2x+1$;
  \item $f=x^{p-1}+\dots+x+1$, o\`u $p$ est premier.
\end{enumerate}
\finenonce{002275} 


\finexercice


\exercice{2276, barraud, 2008/04/24}
\enonce{002276}{}
Soient $A=\Zz [\sqrt{-3}]$ et $K$ son corps de fractions. 
Montrer que $x^2-x+1$ est irr\'eductible dans $A[x]$ sans 
pour autant \^etre irr\'eductible dans $K[x]$. Expliquer la contradiction
apparente avec le corollaire du lemme de Gauss.
\finenonce{002276} 


\finexercice


\exercice{2277, barraud, 2008/04/24}
\enonce{002277}{}
 Soit $P\in \Zz[x]$.
\begin{enumerate}
  \item Supposons que $P(0)$, $P(1)$ soient impairs. Montrer que
$P$ n'a pas de racine dans $\Zz$.
(\emph{Indication :} Utiliser la r\'eduction modulo $2$.)
  \item  Soit $n\in \Nn$ tel qu'aucun des entiers $P(0),\dots,P(n-1)$
ne soit divisible par $n$. Montrer que
$P$ n'a pas de racine dans $\Zz$.
\end{enumerate}
\finenonce{002277} 


\finexercice


\exercice{2278, barraud, 2008/04/24}
\enonce{002278}{}
\begin{enumerate}
  \item Soit $P\in \Zz[x]$. Soit $\dfrac{a}{b}$
 sa racine rationnelle : $P(\dfrac{a}{b})=0$, $\pgcd(a,b)=1$.
Montrer que $\forall\,k\in \Zz$ $\ (a-bk)$ divise $P(k)$. 
  \item Quelles racines rationnelles ont les polyn\^omes
$f(x)=x^3-6x^2+15x-14$ et $g(x)=2x^3+3x^2+6x-4$ ?
\end{enumerate}
\finenonce{002278} 


\finexercice


\exercice{2279, barraud, 2008/04/24}
\enonce{002279}{}
\begin{enumerate}
  \item Soient $P\in \Zz[x]$, $n\in \Nn$, 
$m=P(n)$. Montrer que $\forall\, k\in \Zz\ $ $\ m\,|\, P(n+km)$.
  \item En d\'eduire qu'il n'existe aucun polyn\^ome $P\in \Zz[x]$,
non constant, tel que, pour tout $n\in \Zz$,
$P(n)$ soit un nombre premier.
\end{enumerate}
\finenonce{002279} 


\finexercice


\exercice{2280, barraud, 2008/04/24}
\enonce{002280}{}
Dans le cours nous avons d\'ej\`a montr\'e que le produit de polyn\^omes
primitifs est  aussi primitif et   que 
$$
c(f\cdot g)=c(f)\cdot c(g)\qquad \forall\ f,g\in \Zz[x].
$$
\begin{enumerate}
  \item Etant donn\'e $f\in \Qq[x]$, alors $f=\alpha \cdot f_0$ o\`u
$f_0\in \Zz[x]$ est un polyn\^ome primitif et $\alpha\in \Qq$.
  \item Soit $g\in \Zz[x]$ un polyn\^ome primitif, $\alpha\in \Qq$
tel que $\alpha \cdot g\in \Zz[x]$. Alors $\alpha\in \Zz$.
  \item Consid\`erons deux polyn\^omes  $d$,   $f$ sur $\Zz$.
Si $d$ est primitif et $d$ divise
$f$ dans $\Qq[x]$ alors $d$ divise $f$ dans $\Zz[x]$.
  \item Supposons que $d=\pgcd_{\Qq[x]}(f,g)$  soit le p.g.c.d.  dans l'anneau
$\Qq[x]$ de deux polyn\^omes primitifs $f$ et $g$ de $\Zz[x]$.
Soit  $d=\alpha\cdot d_0$  sa repr\'esentation de type  1).
Montrer que : $d_0=\pgcd_{\Zz[x]}(f,g)$ dans l'anneau $\Zz[x]$.
  \item Soient $f$, $g\in \Zz[x]$, $f=c(f)f_0$, $g=c(g)g_0$. Alors
$$
\pgcd_{\Zz[x]}(f,g)=\pgcd_{\Zz}(c(f),c(g))\cdot\pgcd_{\Zz[x]}(f_0,g_0).
$$
\end{enumerate}
\finenonce{002280} 


\finexercice


\exercice{2281, barraud, 2008/04/24}
\enonce{002281}{}
D\'emontrer que tout morphisme d'un corps dans
un anneau non-trivial est injectif.
\finenonce{002281} 


\finexercice


\exercice{2282, barraud, 2008/04/24}
\enonce{002282}{}
Soit $R$ un anneau int\`egre
dans lequel toute cha\^{\i}ne d\'ecroissante d'id\'eaux est finie.
D\'emontrer que $R$ est un corps.
\finenonce{002282} 


\finexercice


\exercice{2283, barraud, 2008/04/24}
\enonce{002283}{}
Montrer que dans un anneau fini tout id\'eal premier est maximal.
\finenonce{002283} 


\finexercice


\exercice{2284, barraud, 2008/04/24}
\enonce{002284}{}
Montrer que un id\'eal propre  $I$ de l'anneau
$A$ est premier ssi quand le produit de deux id\'eaux est 
contenue dans $I$, alors l'un de deux est contenu dans $I$.
En d\'eduire que si $M$ est un id\'eal maximal de $A$, alors
le seul id\'eal premier de $A$ qui contient $M^n$ est $M$.
\finenonce{002284} 


\finexercice


\exercice{2285, barraud, 2008/04/24}
\enonce{002285}{}
 Soit $A$ un anneau. Trouver les anneaux quotients
$$
A[x]/(x),\quad A[x,y]/(x),\quad A[x,y]/(x,y),\quad
A[x_1,x_2,\dots,x_n]/(x_1,x_2,\dots,x_n)
$$
o\`u $(x)$, $(x,y)$, $(x_1,x_2,\dots,x_n)$ sont les id\'eaux engendr\'es
r\'espectivement par $x$, $x$ et $y$, $x_1$, $x_2$, ... ,$x_n$.
Sous quelle condition sur l'anneau $A$ ces id\'eaux sont-ils premiers
(maximaux) ?
\finenonce{002285} 


\finexercice


\exercice{2286, barraud, 2008/04/24}
\enonce{002286}{}
\begin{enumerate}
  \item Trouver le nombre d'\'el\'ements de l'anneau
quotient $\Zz[\sqrt{d}]/(m)$ o\`u $m\in \Zz$ et $m\ne 0$.
  \item L'id\'eal principal endendr\'e par $2$ est-il premier
dans l'anneau $\Zz[\sqrt{d}]$ ?
\end{enumerate}
\finenonce{002286} 


\finexercice


\exercice{2287, barraud, 2008/04/24}
\enonce{002287}{}
Soit $A$ un anneau int\`egre. On appelle
{\it \'el\'ement premier} de $A$ un \'el\'ement qui engendre un id\'eal
principal premier. 
\begin{enumerate}
  \item Montrer que  un \'el\'ement premier est irr\'eductible.
  \item D'apr\`es le cours  tout  \'el\'ement irr\'eductible dans un anneau  
factoriel  est  premier. 
Montrer que dans un anneau factoriel, tout id\'eal premier
non nul contient un \'el\'ement irr\'eductible.

  \item Nous avons vu que l'\'el\'ement $3\in \Zz[\sqrt{-5}]$ 
est irr\'eductible. Montrer que $3$ n'est pas premier 
dans $\Zz[\sqrt{-5}]$.
  \item L'\'el\'ement  $2$ est-il irr\'eductible dans l'anneau 
$\Zz[\sqrt{-5}]$ ?
\end{enumerate}
\finenonce{002287} 


\finexercice


\exercice{2288, barraud, 2008/04/24}
\addcommand{\usepackage[all]{xy}}
\enonce{002288}{}
\label{ex:bar49}
\begin{enumerate}
  \item Soit $A$ un anneau principal, $I$ un id\'eal
de $A$. Montrer que  tous les id\'eaux de l'anneau quotient $A/I$ 
sont  principaux.
  \item Trouver tous les id\'eaux des anneaux suivants:
$\Zz/n\Zz$, 
$\Qq[x]/(f)$ o\`u $(f)$ est l'id\'eal principal engendr\'e par 
un polyn\^ome $f$. 
  \item Trouver les id\'eaux maximaux de $\Zz /n\Zz$ et de $\Qq[x]/(f)$.
\end{enumerate}
\finenonce{002288} 


\finexercice


\exercice{2289, barraud, 2008/04/24}
 \addcommand{\usepackage[all]{xy}}
\enonce{002289}{}
Soit $I$ et $J$ deux id\'eaux de l'anneau $A$.
Consid\'erons la projection canonique 
\newline
$\pi_I : A\to A/I$ 
et l'image $\bar J=\pi_I(J)$ de l'id\'eal $J$.
\begin{enumerate}
  \item Montrer que  $\bar J$ est un id\'eal de l'anneau quotient $A/I$.
  \item  D\'emontrer qu'on a l'isomorphisme suivant
: $(A/I)/\bar J\cong A/(I+J)$.

({\it Indication :.} Consid\'erer le morphisme $a+I\mapsto a+(I+J)$
de l'anneau $A/I$ vers l'anneau $A/(I+J)$.)
\end{enumerate}
\finenonce{002289} 


\finexercice


\exercice{2290, barraud, 2008/04/24}
\enonce{002290}{}
Soit $f$ un morphisme de l'anneau $A$ vers l'anneau $B$.
\begin{enumerate}
  \item  Montrer que l'image r\'eciproque d'un id\'eal premier  est
aussi  un id\'eal premier.  Cette proposition est-elle vraie
pour id\'eaux maximaux ?
  \item  Montrer par un exemple, que l'image $f(I)$ d'un id\'eal $I$ de $A$
n'est pas forc\'ement un id\'eal de $B$. D\'emontrer cependant
que si $f$ est surjectif, alors $f(I)$ est un id\'eal pour tout
id\'eal $I$ de $A$. (Voir le cours.)
  \item Toujours sous l'hypoth\`ese que $f$ est surjective, montrer que
l'image d'un id\'eal maximal par $f$ est soit $B$ tout entier, 
soit un id\'eal maximal de $B$.
  \item Consid\'erons la reduction de polyn\^omes sur $\Zz$ modulo $m$ :
$r_m\,:\,\Zz[x]\to \Zz_m[x]$ et deux id\'eaux premiers principaux
$(x)$ et $(x^2+1)$. Les id\'eaux $r_6((x))$ et $r_2((x^2+1))$
sont-ils premiers ?
\end{enumerate}
\finenonce{002290} 


\finexercice


\exercice{2291, barraud, 2008/04/24}
\addcommand{\usepackage[all]{xy}}
\enonce{002291}{}
 Soit $A$ un anneau, $B$ un sous-anneau de $A$,
$I$ un id\'eal de $A$. 
\begin{enumerate}
  \item Montrer que $B\cap I$ est un id\'eal de $B$,
$B+I=\{b+i\,|\, b\in B,\ i\in I\}$ est un sous-anneau
de l'anneau $A$ et $I$ est un id\'eal de ce sous-anneau.
  \item Montrer que l'anneau quotient $B/(B\cap I)$  est isomorphe \`a
l'anneau quotient $(B+I)/I$. ({\it Indication :} Consid\'erer  le compos\'e de l'inclusion
$B\to B+I$ avec la projection canonique  $B+I \to (B+I)/I$.)
\end{enumerate}


\finenonce{002291} 


\finexercice


\exercice{2292, barraud, 2008/04/24}
\enonce{002292}{}
Soit $(x^3-x+2)$ l'id\'eal principal engendr\'e
par $x^3-x+2$ dans l'anneau  $\Qq[x]$.
\begin{enumerate}
  \item Montrer que l'anneau quotient $\Qq[x]/(x^3-x+2)$ est un corps.
  \item Soit $y$ l'image de $x$ dans $\Qq[x]/(x^3-x+2)$ par la surjection
canonique. Calculer son inverse.
  \item  Montrer que $1+y+y^2$ est non nul et calculer son inverse.
\end{enumerate}
\finenonce{002292} 


\finexercice


\exercice{2293, barraud, 2008/04/24}
\enonce{002293}{}
Soit $f\in A[x]$ un polyn\^ome primitif de degr\'e positif sur l'anneau
factoriel $A$. Soit $\pi\in A$ un \'el\'ement irr\'eductible.
Supposons que le coefficient dominant de $f$ ne soit  pas divisible
par $\pi$ et que  $f\mod \pi$ soit  irr\'eductible dans l'anneau quotient
$A/(\pi)$. Montrer que $f$ est irr\'eductible dans $A[x]$.
\finenonce{002293} 


\finexercice


\exercice{2294, barraud, 2008/04/24}
\enonce{002294}{}
Les polynômes suivants sont-ils irréductibles ?
\begin{enumerate}
  \item $X^5+121X^4+1221X^3+12221X^2+122221X+222222$ dans $\Qq [X]$.
  \item  $f(X,Y)=X^2Y^3+X^2Y^2+Y^3-2XY^2+Y^2+X-1$ dans $\Cc [X,Y]$ et $\mathbb{F}_2[X,Y]$.
  \item $f(X,Y)=Y^7+Y^6+7Y^4+XY^3+3X^2Y^2-5Y+X^2+X+1$ dans $\Qq[X,Y]$.
\end{enumerate}
\finenonce{002294} 


\finexercice


\exercice{2295, barraud, 2008/04/24}
\enonce{002295}{}
 L'id\'eal  principal $(x^2+y^2+1)$
est-il maximal dans les anneaux $\Cc[x,y]$, $\Rr[x,y]$, $\Qq[x,y]$,
$\Zz[x]$, $\Zz_2[x,y]$ ?
\finenonce{002295} 


\finexercice


\exercice{2296, barraud, 2008/04/24}
\enonce{002296}{}
\begin{enumerate}
  \item  Soit $f\in \Zz[x]$. 
Consid\'erons la reduction du polyn\^ome $f$ modulo $m$ :
$f\mod m\in \Zz_m[x]$. 
Montrer que
$$
\Zz[x]/(m,f)\cong \Zz_m[x]/(f\hbox{\ mod } m)
$$
o\`u $(m,f)$ est 
l'id\'eal engendr\'e par $m$ et $f$ dans $\Zz[x]$ et 
$(f\hbox{\ mod } m)$ est l'id\'eal engendr\'e par $f\mod m$  dans 
$\Zz_m[x]$.
({\it Indication :} Utiliser l'exercice 10 de fiche 4.)
  \item  Si $p$ est un nombre premier et $f$ est un polyn\^ome tel que
$f\mod p$ est irr\'eductible sur le corps $\Zz_p$, alors
l'id\'eal $(p,f)$ est maximal dans $\Zz[x]$.
\end{enumerate}
\finenonce{002296} 


\finexercice


\exercice{2297, barraud, 2008/04/24}
\enonce{002297}{}
Soit $A$ un anneau factoriel.
\begin{enumerate}
  \item Pour $a,b\ne 0$ on a $(a)\cdot (b)= (a)\cap(b)\ \ $ ssi $\ \ \pgcd(a,b)\thicksim 1$.
  \item  Si $(a,b)$ est principal, alors $(a,b)=(\pgcd(a,b))$.
\end{enumerate}
\finenonce{002297} 


\finexercice


\exercice{2298, barraud, 2008/04/24}
\enonce{002298}{}
\begin{enumerate}
  \item Montrer que  les id\'eaux $(5, x^2+3)$, $(x^2+1,x+2)$, $(x^3-1,x^4-1)$
ne sont pas principaux dans $\Zz[x]$. 
  \item  Les id\'eaux $(x,x+1)$, $(5, x^2+4)$ et  $(x^2+1,x+2)$  
sont-ils  premiers  ou maximaux dans $\Zz[x]$ ?
\end{enumerate}
\finenonce{002298} 


\finexercice


\exercice{2299, barraud, 2008/04/24}
\enonce{002299}{}
D\'emontrer que si $J$ est un id\'eal premier
de l'anneau $\Zz[x]$, alors
$$
J=(0), \quad (p), \quad (f)\ \ \text{\ ou\ } \ \  (p,g),
$$
o\`u $p$ est premier, $f\in \Zz[x]$ est un polyn\^ome irr\'eductible
de degr\'e positif et 
$g$ est un polyn\^ome, tel que  sa r\'eduction modulo $p$ 
est irr\'eductible sur  $\Zz_p$. 
Le dernier cas, $J=(p,g)$ , nous donne la forme g\'en\'erale d'un id\'eal
maximal dans $\Zz[x]$.
{\it Le plan de la d\'emonstration est le suivant.}
\begin{enumerate}
  \item Soit $B$ un sous-anneau de l'anneau $A$, $I$ un id\'eal premier de
$A$. Montrer que $B\cap I$ est  soit un id\'eal premier de $B$, soit
l'anneau $B$ lui-m\^eme.
  \item  Soit $J$ un id'eal premier de $\Zz[x]$. Montrer que 
$\Zz\cap J=(0)$ ou $(p)$ o\`u $p$ est premier.
  \item Supposons que $\Zz\cap J=(0)$. Montrer que si $J\ne (0)$,
alors $J$ est engendr\'e par 
un polyn\^ome primitif  de $J$ de degr\'e minimal.
  \item Supposons que $\Zz\cap J=(p)$. Soit $r_p\,:\,\Zz[x]\to \Zz_p[x]$
la r\'eduction modulo $p$. Montrer que l'id\'eal $r_p(J)$ est premier
et que $J=(p,g)$.
  \item Montrer que $J$ est maximal ssi $J=(p,g)$ o\`u
$p$ est premier et $r_p(g)$ est irr\'eductible dans $\Zz_p[x]$.
\end{enumerate}
\finenonce{002299} 


\finexercice


\exercice{6526, drutu, 2011/10/16}
\enonce{006526}{}
Soit $\mathbb{F} $ un corps et $P,\; Q\in \mathbb{F}[X]$ tels que $\pgcd(P,\, Q)=1$. Montrer qu'il existe $U,\; V\in \mathbb{F}[X]$ tels que $UP+VQ=1,\; d^0\; U<d^0\; Q,\; d^0\; V<d^0\; P$. 
\finenonce{006526}



\finexercice


\exercice{6527, drutu, 2011/10/16}
\enonce{006527}{}
Trouver tous les automorphismes de l'anneau $\mathbb{F} [X]$, où $\mathbb{F}$ est un corps.      
\finenonce{006527}



\finexercice


\exercice{6528, drutu, 2011/10/16}
\enonce{006528}{}
 Utiliser l'algorithme d'Euclide dans $\Q[X]$ pour exprimer le pgcd demandé sous forme de combinaison linéaire des deux polyn\^omes donnés :
\begin{enumerate}
\item $ (X^4+2X^3-X^2-4X-2,X^4+X^3-X^2-2X-2)$,
\item $(X^5+3X^4+X^3+X^2+3X+1,X^4+2X^3+X+2)$,
\item $(3X^3-2X^2+X+2,X^2-X+1)$.
\end{enumerate}
\finenonce{006528}



\finexercice


\exercice{6529, drutu, 2011/10/16}
\enonce{006529}{}
 Montrer que tout monomorphisme $A\to A'$ d'anneaux intègres engendre un monomorphisme des corps des fractions. 
\finenonce{006529}



\finexercice


\exercice{6530, drutu, 2011/10/16}
\enonce{006530}{}
 Si $Q$ est le corps des fractions de l'anneaux intègre $A$, démontrer que le corps des fractions $Q(X)$ est isomorphe au corps des fractions de $A[X]$.
\finenonce{006530}



\finexercice


\exercice{6531, drutu, 2011/10/16}
\enonce{006531}{}
Soit $P(X)=a_nX^n+a_{n-1}X^{n-1}+\cdots +a_0\in \Z[X]$ et soit $c\in \Q$ une racine de $P$. Montrer que $c=\frac{p}{q}$, où $p|a_0$ et $q|a_n$.
\finenonce{006531}



\finexercice


\exercice{6532, drutu, 2011/10/16}
\enonce{006532}{}
\begin{enumerate}
\item Soit $a\in \Z$ un nombre sans facteurs carrés, $a\not\in \{ -1,\; 0,\; 1\}$. Montrer que $X^n-a$ est irréductible sur $\Q $ pour tout $n\in \N^*$.

\item Montrer que $2X^{10}-21$, $3X^5-35,$ $X^5+1000X^2+6$ sont irréductibles sur $\Q $.

\item Démontrer l'irréducibilité dans $\Q[X]$ des polyn\^omes $X^4-8X^3+12X^2-6X+2$ et $X^5-12X^3+36X-12$.
\end{enumerate}
\finenonce{006532}



\finexercice


\exercice{6533, drutu, 2011/10/16}
\enonce{006533}{}
Montrer que le polyn\^ome $X^2-1$ a quatre zéros dans $\Z /15\Z$. Est-ce que ceci contredit un théorème du cours ?
\finenonce{006533}



\finexercice


\exercice{6534, drutu, 2011/10/16}
\enonce{006534}{}
 Montrer que $X^2+1$ est irréductible en tant qu'élément de $\Q[X]$ mais réductible en tant qu'élément de $\mathbb{F}_5[X]$.
\finenonce{006534}



\finexercice


\exercice{6535, drutu, 2011/10/16}
\enonce{006535}{}
\begin{enumerate}
\item Si $\mathbb{F} $ est un corps, montrer qu'un polyn\^ome $P$ de degré 2 ou 3 dans $\mathbb{F} [X]$ est irréductible si et seulement si il n'a pas de zéro dans $\mathbb{F}$. 

\item Décomposer les polyn\^omes suivants en facteurs irréductibles dans $\mathbb{F}_3[X]$.
$$
X^2+X+1,\hspace{1cm}X^3+X+2,\hspace{1cm}X^4+X^3+X+1\; . 
$$
\end{enumerate}
\finenonce{006535}



\finexercice


\exercice{6536, drutu, 2011/10/16}
\enonce{006536}{}
Lesquels parmi les polyn\^omes suivants sont irréductibles dans $\Q[X]$ ?

$$
X^3+4X^2-5X+7\; ;\hspace{0.5cm}X^3-6X^2-4X-13\; ;\hspace{0.5cm}X^3+4X^2-4X+25\; ;
$$
$$
X^3-X^2-X-1\; ;\hspace{0.5cm}X^4+7X^2+4X+1\; ;\hspace{0.5cm}X^6+X^3+1\; ;\hspace{0.5cm}X^6+X^2+1\; .
$$
\finenonce{006536}



\finexercice


\exercice{6537, drutu, 2011/10/16}
\enonce{006537}{}
\begin{enumerate}
\item Soit $A$ un anneau commutatif et $I\subset A$ un idéal.
Montrer que $I[X] $ est un idéal de $A[X]$. Montrer que $A[X]/I[X]\simeq A/I[X]$. Montrer que si $I$ est premier, $I[X]$ est premier.

\item Soit $p$ un nombre premier. Montrer que $\Z[X]/(p)\simeq \mathbb{F}_p[X]$.

\item Montrer que les idéaux premiers non-nuls de $\Z[X]$ sont :
  \begin{enumerate}
  \item  $(p)$, où $p\in \N$ est un nombre premier ;
  \item $(R(X))$, où $R\in \Z[X]$ est un polyn\^ome irréductible ;
  \item $(p,R(X))$, où $p\in \N$ est un nombre premier et $R\in \Z[X]$ est un polyn\^ome irréductible dont la réduction modulo $p$, $\overline{R}(X)\in \mathbb{F}_p[X]$, est un polyn\^ome irréductible. 
  \end{enumerate}
\end{enumerate}
\finenonce{006537}



\finexercice


\exercice{6538, drutu, 2011/10/16}
\enonce{006538}{}
Soient $A=\Z[i\sqrt{3}]$ et $K$ son corps de fractions. Montrer que $X^2-X+1$ est primitif et irréductible dans $A[X]$ sans pour autant être irréductible dans $K[X]$. Est-ce que ceci contredit un théorème du cours ?
\finenonce{006538}



\finexercice


\exercice{6539, drutu, 2011/10/16}
\enonce{006539}{}
 Soit $p\in \N^*$ un nombre premier. Montrer que les polyn\^omes
 $P=\Sigma_{i=1}^pC_p^iX^{i-1}$ et $Q=1+\Sigma_{i=1}^{p-1}X^{i}$ sont irréductibles dans $\Q[X]$. 
\finenonce{006539}



\finexercice


\exercice{6540, drutu, 2011/10/16}
\enonce{006540}{}
Soit $\mathbb{F}$ un corps de caractéristique $p>0$. 

\begin{enumerate}
\item Montrer que 

  \begin{enumerate}
  \item $X^n-a^n,\; a\in \mathbb{F}^*$, n'a pas de racines multiples ssi $p|n$ ;
  \item $a$ est l'unique racine de $X^p-a^p$ et sa multiplicité est $p$. 
  \end{enumerate}

\item Soit $a\in \mathbb{F}_p$. Ecriver la décomposition en facteurs irréductibles du polyn\^ome $X^p+a\in \mathbb{F}_p[X]$. 
\end{enumerate}
\finenonce{006540}



\finexercice


\exercice{6541, drutu, 2011/10/16}
\enonce{006541}{}
Soit $D$ un anneau intègre fini contenant $n$ éléments distincts $c_1,c_2,\dots ,c_n$. Soit le polyn\^ome $P_0:=\Pi_{i=1}^n(X-c_i)$.
\begin{enumerate}
\item Démontrer que deux polyn\^omes $Q, R$ dans $D[X]$ ont la même fonction polynomiale associée si et seulement si $P_0 | Q-R$.

\item Calculer le polyn\^ome $P_0$ pour $D=\mathbb{F}_3,\; \mathbb{F}_5 $.
\end{enumerate}
\finenonce{006541}



\finexercice


\exercice{6542, drutu, 2011/10/16}
\enonce{006542}{}
Soit $A$ un anneau commutatif. Montrer l'équivalence des deux affirmations suivantes :
\begin{enumerate}
\item[(i)] $A $ est un corps fini ;

\item[(ii)] tout polyn\^ome $P\in A[X]$ de degré $n\geq 1$ a au plus $n$ racines dans $A$ et toute fonction $f:A \to A$ est polynomiale.
\end{enumerate}
\finenonce{006542}



\finexercice


\exercice{6543, drutu, 2011/10/16}
\enonce{006543}{}
 Calculer la dérivée du  polyn\^ome $(3X^2+2X-4)(4X^3-2X+3)\in \mathbb{F}_5 [X] $. 
\finenonce{006543}



\finexercice


\exercice{6544, drutu, 2011/10/16}
\enonce{006544}{}
\begin{enumerate}
\item Montrer que si $A[X]$ est principal, alors $A$ est principal.

\item Montrer que si $A[X]$ est principal est si $A$ n'est pas un corps, alors il existe un corps $K$ tel que $K[X]$ soit un corps.

\item En déduire que $A[X]$ est principal si et seulement si $A$ est un corps.
{\it Remarque}: Ceci implique que si $A[X]$ est principal alors $A[X]$ est euclidien.   
\end{enumerate}
\finenonce{006544}



\finexercice


\exercice{6545, drutu, 2011/10/16}
\enonce{006545}{}
Soit $K$ un corps et soit $K(X)$ le corps des fractions de l'anneau $K[X]$. Montrer que si $P$ est un élément de $K[X]$ qui admet au moins une racine simple alors, pour tout $n\in \N^*$, $Y^n-P$ est un élément irréductible de $K[X,Y]$ et de $K(X)[Y]$.   
\finenonce{006545}



\finexercice


\exercice{6546, drutu, 2011/10/16}
\enonce{006546}{}
Montrer que si $K$ est un corps de caractéristique différente de $2$ alors $X^2+Y^2-1$ est irréductible dans $K[X,Y]$. 
\finenonce{006546}



\finexercice


\exercice{6547, drutu, 2011/10/16}
\enonce{006547}{}
\begin{enumerate}
\item Calculer la somme des carrés des racines de l'équation $x^3+2x-3=0$.

\item Calculer $x_1^3x_2+x_1x_2^3+x_2^3x_3+x_2x_3^3+x_3^3x_1+x_3x_1^3$, où $x_1,\; x_2,\; x_3$ sont les racines de l'équation $x^3-x^2-4x+1=0$.
\end{enumerate}
\finenonce{006547}



\finexercice


\exercice{6548, drutu, 2011/10/16}
\enonce{006548}{}
 Exprimer à l'aide des polyn\^omes symétriques fondamentaux :
\begin{enumerate}
\item $(x_1^2+x_2^2)(x_1^2+x_3^2)(x_2^2+x_3^2)$ ;

\item $(x_1+x_2)(x_1+x_3)(x_1+x_4)(x_2+x_3)(x_2+x_4)(x_3+x_4)$ ;

\item $(x_1+x_2-x_3-x_4)(x_1-x_2+x_3-x_4)(x_1-x_2-x_3+x_4)$.
\end{enumerate}
\finenonce{006548}



\finexercice


\exercice{6549, drutu, 2011/10/16}
\enonce{006549}{}
Soient $x_1,x_2,\dots ,x_n$ les zéros du polyn\^ome $X^n+a_1x_{n-1}+\dots +a_n$. Démontrer que tout polyn\^ome symétrique en $x_2,x_3,\dots ,x_n$ peut être représenté sous forme de polyn\^ome en $x_1$. 
\finenonce{006549}



\finexercice


\exercice{6550, drutu, 2011/10/16}
\enonce{006550}{}
 Calculer le résultant des polyn\^omes :
\begin{enumerate}
\item $2X^3-3X^2-X+2$ et $X^4-2X^2-3X+4$ ;

\item $3X^3+2X^2+X+1$ et $2X^3+X^2-X-1$ ;

\item $a_0X^2+a_1X+a_2$  et $b_0X^2+b_1X+b_2$.
\end{enumerate}
\finenonce{006550}



\finexercice


\exercice{6551, drutu, 2011/10/16}
\enonce{006551}{}
Pour quelle valeur de $\lambda$ les polyn\^omes suivants ont-ils un zéro en commun ? 
\begin{enumerate}
\item $X^3-\lambda X+2$ et $X^2+\lambda X+2$ ;

\item $X^3+\lambda X^2-9$ et $X^3+\lambda X-3$. 
\end{enumerate}
\finenonce{006551}



\finexercice

\section{325.00 Extension de corps}


\exercice{6552, drutu, 2011/10/16}
\enonce{006552}{}
 Soit $K$ un corps et $k,\; A,\; B$ des sous-corps tels que $k\subset A$ et $k\subset B$, $[A:k]=m,\; [B:k]=n$. Soit $L$ le plus petit sous-corps de $K$ qui contient $A\cup B$. 
\begin{enumerate}
\item Montrer que $[L:A]\leq n,\; [L:B]\leq m,\; [L:k]\leq mn$. Caractériser le cas $[L:k]= mn$ à l'aide d'une propriété de $A$ par rapport à $B$.

\item Si $[K:k]=4,\, m=n=2$ montrer l'équivalence des propriétés suivantes
  \begin{enumerate}
  \item[(b$_1$)] $A\neq B$ ;
  \item[(b$_2$)] $L=K$ ;
  \item[(b$_3$)] il existe $a\in A$ et $b\in B$ tels que $\{ 1,\, a,\, b,\, ab \}$ soit une base de $L$ sur $k$.
  \end{enumerate}
\end{enumerate}
\finenonce{006552}



\finexercice


\exercice{6553, drutu, 2011/10/16}
\enonce{006553}{} 
 Est-ce qu'un extension algébrique est toujours finie ?
\finenonce{006553}



\finexercice


\exercice{6554, drutu, 2011/10/16}
\enonce{006554}{}
 Comparer les corps de décomposition des polyn\^omes $X^2+X+1,\; X^2+3\in \Q [X]$.   
\finenonce{006554}



\finexercice


\exercice{6555, drutu, 2011/10/16}
\enonce{006555}{}
 Soit $E$ une extension du corps $k$. On dit que $E$ est une {\it extension quadratique} de $k$ si $[E:k]=2$. 
\begin{enumerate}
\item On suppose $k$ de caractéristique $\neq 2$. Montrer que les extensions quadratiques de $k$ (s'il en existe) sont des corps de rupture des polyn\^omes irréductibles de la forme $X^2-a,\; a\in k$. En déduire toutes les extensions quadratiques de $\Q$ à isomorphisme près.

\item On suppose $k$ de caractéristique $2$. Montrer que les extensions quadratiques de $k$ (s'il en existe) sont des corps de rupture des polyn\^omes irréductibles de l'une des deux formes : $X^2-a$ ou $X^2-X-a,\; a\in k$. Une extension quadratique du premier type et une extension quadratique du deuxième type peuvent-elles être isomorphes au dessus de $k$ ? 
\end{enumerate}
\finenonce{006555}



\finexercice


\exercice{6556, drutu, 2011/10/16}
\enonce{006556}{}
\begin{enumerate}
\item Montrer qu'un corps fini n'est jamais algébriquement clos.

\item Quelle est la cl\^oture algébrique du corps fini $\mathbb{F}_{p^n}$ ?
\end{enumerate}
\finenonce{006556}



\finexercice


\exercice{6557, drutu, 2011/10/16}
\enonce{006557}{}
 Soit $K$ un corps et $\overline{K}$ une cl\^oture algébrique de $K$. Soient $a,b\in \overline{K}\setminus K$. Etablir l'équivalence des conditions :
\begin{enumerate}
\item il existe un automorphisme $\phi $ de $\overline{K}$ au dessus de $K$ tel que $\phi (a)=b$ ;

\item $a$ et $b$ ont le même polyn\^ome minimal $f(X)\in K[X]$.
\end{enumerate}
\finenonce{006557}



\finexercice


\exercice{6558, drutu, 2011/10/16}
\enonce{006558}{}
Pour quels nombres premiers $p,\; q$ a-t-on $\Q (\sqrt{p})
\subset \Q (\sqrt[3]{q})$ ? 
\finenonce{006558}



\finexercice


\exercice{6559, drutu, 2011/10/16}
\enonce{006559}{}
 Soient $K\subset L\subset M$ des extensions. On suppose $K\subset L $ et $L\subset M$ algébriques. Montrer que $K\subset M$ est algébrique. 
\finenonce{006559}



\finexercice


\exercice{6560, drutu, 2011/10/16}
\enonce{006560}{}
Donner les polyn\^omes minimaux sur $\Q $ des éléments suivants de $\C$ : $j\sqrt{2},\; \sqrt{2}+\sqrt[3]{7},\; i+j,\; i+\sqrt{2},\, j+\sqrt{3}$. 

\finenonce{006560}



\finexercice


\exercice{6561, drutu, 2011/10/16}
\enonce{006561}{}    
Soit $P\in K[X]$ un polyn\^ome de degré $n$ et $E \supset K$ son corps de décomposition. Montrer que $[E:K]$ divise $n!$.
\finenonce{006561}



\finexercice


\exercice{6562, drutu, 2011/10/16}
\enonce{006562}{}
\begin{enumerate}
\item Soit $P\in K[X]$ un polyn\^ome irréductible de degré $n$ et $E \supset K$ une extension de degré $m$. Montrer que si $(m,n)=1$ alors $P$ est irréductible dans $E[X]$.

\item Montrer que $X^3+X+1$ est irréductible sur $\Q [i]$.
\end{enumerate}
\finenonce{006562}



\finexercice


\exercice{6563, drutu, 2011/10/16}
\enonce{006563}{}
 Soit $p,\; q\in \N^*$ deux nombres premiers distincts. Montrer que $\Q (\sqrt{p}+\sqrt{q})=\Q (\sqrt{p},\; \sqrt{q})$. Trouver le polyn\^ome minimal de $\sqrt{p}+\sqrt{q}$ sur $\Q$. Déterminer tous les plongements de $\Q (\sqrt{p}+\sqrt{q})$ dans $\C$. 
\finenonce{006563}



\finexercice


\exercice{6564, drutu, 2011/10/16}
\enonce{006564}{}
 Soit $p\neq 2$ un nombre premier. Montrer que $-1$ est un carré modulo $p$ ssi $p$ est de la forme $4k+1$.   
\finenonce{006564}



\finexercice


\exercice{6565, drutu, 2011/10/16}
\enonce{006565}{}
\begin{enumerate}
\item Est-ce que $2$ et $3$ sont des carrés modulo $7$ ?
 
\item Est-ce que $23$ est un carré modulo 59 ? 
\end{enumerate}
\finenonce{006565}



\finexercice


\exercice{6566, drutu, 2011/10/16}
\enonce{006566}{}
 Soient $k,\; K$ deux corps finis, $k\subset K$, $[K:k]=m$. Montrer que pour tout $d$ diviseur de $m$ il existe un unique corps intermédiaire $k\subset L\subset K$ tel que $[K:L]=d$.
\finenonce{006566}



\finexercice

\section{326.00 Extension d'anneau}


\exercice{6567, drutu, 2011/10/16}
\enonce{006567}{}
 Donner un exemple d'extension infinie et d'extension finie de
$\C$. 
\finenonce{006567}



\finexercice


\exercice{6568, drutu, 2011/10/16}
\enonce{006568}{}   
 Soit $p>2$ un nombre premier et soit $\left( \frac{x}{p}\right)$ le symbole de Legendre pour $x\in \mathbb{F}^*_p$. 
Montrer qu'on a $\sum_{x\in \mathbb{F}^*_p}\left( \frac{x}{p}\right) =0$.
\finenonce{006568}



\finexercice


\exercice{6569, drutu, 2011/10/16}
\enonce{006569}{} 
 Soit $p$ un nombre premier. 
\begin{enumerate}
\item Soit $a$ une racine de $X^p-X-1\in \mathbb{F}_p[X]$, $a\in \overline{\mathbb{F} }_p$. Montrer que $a+k$ est aussi une racine de $X^p-X-1$, pour tout $k\in \mathbb{F}_p$.  

\item Montrer que $X^p-X-1$ est irréductible sur $\mathbb{F}_p$ et sur $\Z $. 
\end{enumerate}
\finenonce{006569}



\finexercice


\exercice{6570, drutu, 2011/10/16}
\enonce{006570}{}
\begin{enumerate}
\item Soit $K$ un corps et $P\in K[X]$ un polyn\^ome de degré $n$. Montrer que $P$ est irréductible sur $K$ ssi $P$ n'a pas de racines dans les extensions $E$ de $K$ de degré au plus $\frac{n}{2}$.

\item Montrer que $X^4+X+1\in \mathbb{F}_2[X]$ est irréductible sur $\mathbb{F}_2$.

\item Montrer que $X^4-4X^3+2X^2-13X+1\in \Z[X]$ est irréductible sur $\Z$.

\item Montrer que $X^4+1$ est irréductible sur $\Z$ et sur $\Q$ et est réductible sur $\mathbb{F}_p$ pour tout nombre premier $p$.  
\end{enumerate}  
\finenonce{006570}



\finexercice


\exercice{6571, drutu, 2011/10/16}
\enonce{006571}{}
 Soit $K\subset E$ une extension finie.
\begin{enumerate}
\item Montrer que la trace $Tr^E_K$ est une forme linéaire sur $E$ en tant que $K$-espace vectoriel.

\item Calculer la norme $N^E_K(\alpha \beta )$ en fonction de $N^E_K(\alpha )$ et de  $N^E_K(\beta )$. Calculer  $N^E_K(a)$, où $a\in K$.
\end{enumerate}
\finenonce{006571}



\finexercice


\exercice{6572, drutu, 2011/10/16}
\enonce{006572}{}
\begin{enumerate}
\item \'Ecrire une base de $\Q [i,\, \sqrt{2}]$ sur $\Q$. 

\item Dans cette base, écrire la matrice de l'application $m_\alpha^E$, où $\alpha =i+\sqrt{2}$. Déduire le polyn\^ome minimal de $\alpha $ dans $\Q[X]$ à partir de cette matrice.

\item Retrouver le polyn\^ome minimal de $\alpha $ dans $\Q[X]$ par calcul direct.
\end{enumerate}
\finenonce{006572}



\finexercice


\exercice{6573, drutu, 2011/10/16}
\enonce{006573}{}
 Soient $A$ un anneau intégralement clos, $K$ son corps de fractions et $P\in A[X]$ un polyn\^ome unitaire. Si $P$ est réductible dans $K[X]$, montrer qu'il est réductible dans $A[X]$.
{\it Indication}: Considérer les racines de $P$ dans une extension de $K$.
\finenonce{006573}



\finexercice


\exercice{6574, drutu, 2011/10/16}
\enonce{006574}{}
Donner un exemple d'anneau intègre qui n'est pas intégralement clos.
\finenonce{006574}



\finexercice


\exercice{6575, drutu, 2011/10/16}
\enonce{006575}{}
 Pour tout $n\geq 1$, on désigne par $f(n)$ le nombre de polyn\^omes unitaires irréductibles de degré $n$ sur $\mathbb{F}_q$. Montrer qu'on a $$
\sum_{d|n}df(d) = q^n\; .
$$

En déduire les valeurs de $f(1),\; f(2),\; f(3),\; f(4)$ et $f(p)$ pour $p$ premier.
\finenonce{006575}



\finexercice


\exercice{6576, drutu, 2011/10/16}
\enonce{006576}{} 
 Soient $R$ un anneau, $A$ un sous-anneau de $R$ et $x$ un élément inversible de $R$. Montrer que tout $y\in A[x]\cap A[x^{-1}]$ est entier sur $A$.
{\it Indication}: On montrera qu'il existe un entier $n$ tel que le $A$-module $M=A+Ax+\cdots +Ax^n$ vérifie $yM\subset M$.
\finenonce{006576}



\finexercice


\exercice{6577, drutu, 2011/10/16}
\enonce{006577}{}
 Combien de solutions rationnelles ont les équations suivantes ?
$$ a) \; x^2+2y^2=5\; ;\hspace{1cm} b)\; 19x^2-12xy+2y^2=4\; ;\hspace{1cm} c) \; x^2+y^2=3.$$
\finenonce{006577}



\finexercice


\exercice{6578, drutu, 2011/10/16}
\enonce{006578}{} 
\'Eliminer $x$ du système d'équations 
$$
\left\{ \begin{array}{c}
        x^3-xy-y^3+y=0\\
        x^2+x-y^2=1
        \end{array}\right.
$$
\finenonce{006578}



\finexercice


\exercice{6579, drutu, 2011/10/16}
\enonce{006579}{}
 Démontrer les formules 
$$
R(fg,h)=R(f,h)R(g,h),\; \; D(fg)=D(f)D(g)(R(f,g))^2\; .
$$
\finenonce{006579}



\finexercice

\section{327.00 Autre}

\section{328.00 Forme bilinéaire}


\exercice{7715, mourougane, 2021/08/11}
\enonce{007715}{}
Montrer que toute forme quadratique non-dégénérée d'indice $1$ en dimension $3$ est à un scalaire près, équivalente à la forme de Lorentz
$x^2+y^2-y^2$.
\finenonce{007715}
\finexercice


\exercice{7716, mourougane, 2021/08/11}
\enonce{007716}{}
Donner l'exemple de deux formes bilinéaires avec même rang, même indice et même discriminant mais non-équivalentes.
\finenonce{007716}
\finexercice


\exercice{7717, mourougane, 2021/08/11}
\enonce{007717}{}
Soit $P$ un plan muni d'une forme quadratique.
Montrer que les conditions suivantes sont équivalentes.
\begin{itemize}
 \item $P$ est un plan hyperbolique.
\item $q$ est non-dégénérée d'indice $1$.
\item le discriminant de $q$ est égal à $-1$, modulo un carré
(et donc dans une base orthonormée, l'expression de $q$ est $q(x)=ax_1^2+bx_2^2$, avec des scalaires $a$ et $b$ tels que $-ab$ est un carré). 
\end{itemize}
\finenonce{007717}
\finexercice


\exercice{7726, mourougane, 2021/08/11}
\enonce{007726}{}
\begin{enumerate}
 \item Montrer que deux formes quadratiques équivalentes sur un espace $E$ prennent les mêmes valeurs dans $k$.

 \item Soit $(E,f)$ un espace muni d'une forme bilinéaire symétrique non dégénérée de forme quadratique associée $q$.
Montrer que si $f$ admet un vecteur non nul isotrope, la forme $q$ prend toutes les valeurs de $k$.

 \item Montrer que $E$ se décompose comme somme directe orthogonale de plans hyperboliques et d'un sous-espace sur lequel la forme quadratique n'a pas de vecteur isotrope non nul.

 \item Montrer que le nombre de plans hyperboliques dans une telle décomposition est indépendant de la décomposition.
Le décrire à l'aide d'un invariant défini en cours.
\end{enumerate}
\finenonce{007726}
\finexercice


\exercice{7727, mourougane, 2021/08/11}
\enonce{007727}{Formes $\sigma$-sesquilinéaires sur $\mathbb{F}_{25}$}
\begin{enumerate}
 \item On considère l'application $\sigma~:~\mathbb{F}_{25}\to \mathbb{F}_{25}, \lambda\mapsto\lambda^5$. Montrer que c'est un automorphisme involutif du corps $\mathbb{F}_{25}$.

 \item Les formes suivantes sur $E:=\mathbb{F}_{25}^3$ sont-elles $\sigma$-sesquilinéaires ?
\begin{eqnarray*}
 f_1( 
\left(\begin{array}{c}x\\ y\\ z\end{array}\right),
\left(\begin{array}{c}x'\\ y'\\ z'\end{array}\right)
)
&=&x(x')^5+3z(y')^5+3y(z')^5.\\
f_2( 
\left(\begin{array}{c}x\\ y\\ z\end{array}\right),
\left(\begin{array}{c}x'\\ y'\\ z'\end{array}\right)
)
&=&x^5(x')^5+x^5(y')^5+y^5(x')^5.\\
f_3(
\left(\begin{array}{c}x\\ y\\ z\end{array}\right),
\left(\begin{array}{c}x'\\ y'\\ z'\end{array}\right)
)
&=&3x(x')^5+z(y')^5+y(z')^5.
\end{eqnarray*}

 \item Parmi les formes $\sigma$-sesquilinéaires précédentes, lesquelles sont équivalentes ?
\end{enumerate}
\finenonce{007727}
\finexercice


\exercice{7728, mourougane, 2021/08/11}
\enonce{007728}{Action symplectique sur les droites et les plans}
\begin{enumerate}
 \item Soit $E$ un espace muni d'une forme alternée non-dégénérée $f$.
Soit $G$ le groupe (dit symplectique) des isométries de $(E,f)$.
Combien y a-t-il d'orbites dans l'action du groupe $G$ sur l'ensemble $P(E)$ des droites de $E$ ?

 \item Soit $E$ un espace muni d'une forme alternée non-dégénérée $f$. Quelles sont les restrictions possibles à équivalence près de $f$ sur les plans de $E$ ?

 \item Combien y a-t-il d'orbites dans l'action du groupe symplectique d'un espace vectoriel $E$ de dimension~$6$ muni d'une forme alternée non-dégénérée sur l'ensemble des plans de $E$ ?

\end{enumerate}
\finenonce{007728}
\finexercice


\exercice{7732, mourougane, 2021/08/11}
\enonce{007732}{Signature d'une restriction}
Soit $E$ un espace vectoriel réel de dimension $n$ muni d'une forme bilinéaire symétrique $f$ de signature $(p,q)$.
Soit $V$ un sous-espace vectoriel de $E$ de dimension $r$.
Montrer la signature $(p',q')$ de $f$ restreinte à $V$ vérifie
$p'\leq p$ et $p'\geq p+r-n$.
\finenonce{007732}
\finexercice


\exercice{7734, mourougane, 2021/08/11}
\enonce{007734}{}
 \item Soit $E$ un espace muni d'une forme sesquilinéaire non dégénérée réflexive.
Soit $V$ un sous-espace de $E$. Supposons que le radical $\text{rad}(V)=V\cap V^\perp$ de $V$ est de dimension $2$, $\text{rad}(V)=\vec (N_1,N_2)$.
Soit $W$ un supplémentaire de $\text{rad}(V)$ dans $V$.
Calculer en détails le radical de $V'=\vec (N_1)\oplus W$.
\finenonce{007734}
\finexercice


\exercice{7735, mourougane, 2021/08/11}
\enonce{007735}{}
\begin{enumerate}
 \item Soit $E$ un espace muni d'une forme sesquilinéaire non dégénérée réflexive.
Soit $V$ un sous-espace de $E$. Supposons que le radical $\text{rad}(V)=V\cap V^\perp$ de $V$ est de dimension $2$, $\text{rad}(V)=\vec (N_1,N_2)$.
Soit $W$ un supplémentaire de $\text{rad}(V)$ dans $V$.
Calculer en détails le radical de $V'=\vec (N_1)\oplus W$.
 \item Soit $E$ un espace muni d'une forme alternée non-dégénérée $f$.
Soit $G$ le groupe (dit symplectique) des isométries de $(E,f)$.
Combien y a-t-il d'orbites dans l'action du groupe $G$ sur l'ensemble $P(E)$ des droites de $E$ ?
 \item Soit $E$ un espace muni d'une forme alternée non-dégénérée $f$. Quelles sont les restrictions possibles à équivalence près de $f$ sur les plans de $E$ ?
 \item Combien y a-t-il d'orbites dans l'action du groupe symplectique d'un espace vectoriel $E$ de dimension $6$ muni d'une forme alternée non-dégénérée sur l'ensemble des plans de $E$ ?
\end{enumerate}
\finenonce{007735}
\finexercice


\exercice{7788, mourougane, 2021/08/11}
\enonce{007788}{Formes très dégénérées}
Déterminer toutes les formes sesquilinéaires symétriques, anti-symétriques et hermitiennes $f$ sur un espace
$E$ avec $dim E=dim Ker f+1$.
\finenonce{007788}
\finexercice


\exercice{7789, mourougane, 2021/08/11}
\enonce{007789}{Isotropes}
\begin{enumerate}
\item Montrer que si un sous-espace d'un espace non-singulier $(E,f)$ est totalement isotrope, il est de dimension inférieure à $\dim E/2$.
\item Déterminer des vecteurs isotropes et un sous-espace
totalement isotrope maximal pour les formes non dégénérées
suivantes (données par leur forme quadratique). Dans chaque cas, on donnera l'indice, c'est à dire la
dimension des sous-espaces totalement isotropes maximaux. 
\begin{enumerate}
\item $k=\Rr$ : $q=x^2+y^2, x^2-y^2, x^2+y^2-z^2, x^2+y^2+z^2-t^2,
 x^2+y^2-z^2-t^2$.
\item $k=\Cc$ : $q= x^2+y^2, x^2+y^2+z^2$.
\end{enumerate}
\end{enumerate}
\finenonce{007789}
\finexercice


\exercice{7790, mourougane, 2021/08/11}
\enonce{007790}{}
\begin{enumerate}
\item Déterminer si possible un plan totalement isotrope pour chacune des formes bilinéaires symétriques données par les formes quadratiques suivantes sur $\mathbb{C}^4$.
$$P(x,y,z,t)=xy+zt \textrm{ et } Q(x,y,z)=x^2+y^2+z^2+t^2.$$
\item L'entier $-3$ est-il un carré modulo 7 ?
\item Déterminer si possible un vecteur non-nul isotrope pour chacune des formes bilinéaires symétriques données par les formes quadratiques suivantes sur $\mathbb{F}_7^4$.
$$R(x,y,z,t)=xy+zt\textrm{ et } S(x,y,z)=x^2+y^2+z^2+t^2.$$
\end{enumerate}
\finenonce{007790}
\finexercice


\exercice{7791, mourougane, 2021/08/11}
\enonce{007791}{}
Donner l'exemple de deux formes bilinéaires avec même rang, même indice et même discriminant mais non-équivalentes.

\finenonce{007791}
\finexercice


\exercice{7792, mourougane, 2021/08/11}
\enonce{007792}{}
Les formes bilinéaires symétriques données par les formes quadratiques suivantes sur $\mathbb{F}_7^3$ sont-elles équivalentes ?
$$q(x,y,z)=x^2+6y^2+2z^2 \textrm{ et } Q(x,y,z)=xy+3z^2.$$
\finenonce{007792}
\finexercice


\exercice{7793, mourougane, 2021/08/11}
\enonce{007793}{}
\begin{enumerate}
\item Quels sont les sous-espace propres de la matrice $J$ de taille $n\times n$ dont tous les coefficients sont égaux à $1$ ?
\item Diagonaliser dans une base orthonormée pour le produit scalaire standard, la forme bilinéaire symétrique sur $\Rr^n$. Discuter son rang et sa signature.
$$f(x,y)=\sum_{{1\leq i,j\leq  n\atop i\not = j}}x_iy_j+x_jy_i.$$
\end{enumerate}
\finenonce{007793}
\finexercice


\exercice{7794, mourougane, 2021/08/11}
\enonce{007794}{En caractéristique 2}
Soit $g$ la forme bilinéaire symétrique sur $(\Zz/2\Zz)^3$ dont la
matrice dans la base canonique est 
$\left(
 \begin{array}{ccc}
1&0&0\\ 0&0&1\\0&1&0
 \end{array}
\right)$.
Montrer que $g$ est non dégénérée. Peut-on compléter $e_1$ en une base orthogonale ?
\finenonce{007794}
\finexercice


\exercice{7795, mourougane, 2021/08/11}
\enonce{007795}{Espace de matrices}
Soit $E:=M_n(\Rr)$. Déterminer la signature des formes quadratiques 
\begin{enumerate}
\item $q_1 (A)=\operatorname{Tr}(A ^2)$. On cherchera des «grands» sous-espaces où $q_1$ est définie positive ou définie négative.
\item $q_2 (A)=\operatorname{Tr}(A ^2)-(\operatorname{Tr}A)^2$.
\end{enumerate}
\finenonce{007795}
\finexercice


\exercice{7796, mourougane, 2021/08/11}
\enonce{007796}{}
Soit $(E, q(xe_1+ye_2)=2xy)$ un plan hyperbolique orienté. Peut-on prolonger
l'application $u$ définie sur $vect (e_1)$ par $u(e_1)=e_2$ en une
isométrie de $E$ ? En une isométrie directe de $E$ ? 
\finenonce{007796}
\finexercice


\exercice{7797, mourougane, 2021/08/11}
\enonce{007797}{}
Démontrer le théorème de Witt dans le cas particulier suivant :

Soit $E$ et $E'$ deux espaces symplectiques non-singuliers de dimension $4$.
Soit $d\subset E$ et $d'\subset E'$ deux droites 
et $f$ une application linéaire bijective de $d$ sur $d'$.
Montrer qu'il existe une isométrie de $E$ sur $E'$ qui prolonge $f$. 
\finenonce{007797}
\finexercice


\exercice{7798, mourougane, 2021/08/11}
\enonce{007798}{Action des groupes orthogonaux}
Soit $E=\Rr^2$ muni de la forme quadratique $q(x_1,x_2)=x_1^2-x_2^2$ non dégénérée de signature $(1,1)$. Le groupe $O(1,1)$ est par définition le groupe des isométries de $(E,q)$.
\begin{enumerate}
\item Le groupe $O(1,1)$ agit-il transitivement sur les droites de
 $\Rr^2$ ?
\item Soit $q$ une forme quadratique réelle non dégénérée.
 Montrer que si $q$ a la
 même signature en restriction à $F$ et à $F'$ alors $F$ et $F'$ sont
 dans la même orbite sous l'action de $O(q)$.
\item Décrire les orbites de l'action de $O(2,1)$ sur les droites de
 $\Rr^3$.
\end{enumerate}
\finenonce{007798}
\finexercice


\exercice{7810, mourougane, 2021/08/11}
\enonce{007810}{}
Diagonaliser dans une base orthonormée pour le produit scalaire
standard, la forme bilinéaire symétrique sur $\R^n$
$$f(x,y)=\sum_{{1\leq i,j\leq  n\atop i\not = j}}x_iy_j+x_jy_i.$$
Discuter son rang et sa signature.
\finenonce{007810}
\finexercice


\exercice{7811, mourougane, 2021/08/11}
\enonce{007811}{}
On considère dans $\C ^4$ la matrice
$$A=\left[
\begin{array}{cccc}
1-i&2&1&3\\ 2&-i&-1&0\\ 1&-1&1-i&-1\\ 3&0&-1&-i
\end{array}\right]
$$ Montrer sans calcul que $A$ est inversible ?
\finenonce{007811}
\finexercice


\exercice{7812, mourougane, 2021/08/11}
\enonce{007812}{Décomposition polaire d'un endomorphisme}
Soit $H$ un espace hermitien et $a$ un endomorphisme de $H$.
\begin{enumerate}
\item Montrer que $b=a a ^*$ est un endomorphisme auto-adjoint
positif.
\item Montrer en utilisant une diagonalisation qu'il existe un unique
 endomorphisme auto-adjoint positif $c$ de $H$ tel que $c^2=b$.
\item On suppose désormais que $a$ est inversible. Montrer que $a$
 s'écrit de fa{\c c}on unique sous la forme $a=hu$ o{ù} $h$ est un
 endomorphisme auto-adjoint positif et $u$ un endomorphisme unitaire.
\item Déterminer $h$ et $u$ pour l'endomorphisme $a$ dont la matrice
 dans la base canonique de $\Cc^2$ est $$\left[
  \begin{array}{cc}
1&i\\ i&1
  \end{array}
\right]$$
\end{enumerate}
\finenonce{007812}
\finexercice


\exercice{7813, mourougane, 2021/08/11}
\enonce{007813}{Isotropes}
Déterminer l'ensemble des vecteurs isotropes et les sous-espaces
totalement isotropes maximaux pour les formes non dégénérées
suivantes. Dans chaque cas, on donnera l'indice, c'est à dire la
dimension des sous-espaces totalement isotropes maximaux. 
\begin{enumerate}
\item $k=\Rr$ : $q=x^2+y^2, x^2-y^2, x^2+y^2-z^2, x^2+y^2+z^2-t^2,
 x^2+y^2-z^2-t^2$.
\item $k=\Cc$ : $q= x^2+y^2, x^2+y^2+z^2$.
\end{enumerate}
\finenonce{007813}
\finexercice


\exercice{7814, mourougane, 2021/08/11}
\enonce{007814}{Formes très dégénérées}
Déterminer toutes les formes sesquilinéaires réflexives $f$ sur un espace
$E$ avec $dim E=dim Kerf +1$.
\finenonce{007814}
\finexercice


\exercice{7815, mourougane, 2021/08/11}
\enonce{007815}{En caractéristique 2}
Soit $g$ la forme bilinéaire symétrique sur $(\Zz/2\Zz)^3$ dont la
matrice dans la base canonique $(e_1, e_2, e_3)$ est 
$$\left(
 \begin{array}{ccc}
1&0&0\\ 0&0&1\\0&1&0
 \end{array}
\right)$$
Montrer que $g$ est non dégénérée. Peut-on compléter $e_1$ en une
base orthogonale ?
\finenonce{007815}
\finexercice


\exercice{7816, mourougane, 2021/08/11}
\enonce{007816}{Involutions}
Soit $f$ une forme bilinéaire symétrique non dégénérée
sur un espace vectoriel $E$. 
\begin{enumerate}
\item
Soit $u\in GL(E)$ une involution. Montrer que $u$ est orthogonale si
et seulement si $E_+(u)=E_+:=Ker (u-Id)$ et $E_-:=Ker(u+Id)$ sont
orthogonaux.
 Montrer alors que $(E_+)^\perp = E_-$ et que $E_+$ n'est pas isotrope. 
\item Soit $F\subset E$ un sous-espace non isotrope. Montrer qu'il
 existe une unique involution orthogonale telle que $E_+(u)=F$.
\item Montrer que $O(f)\equiv SO(f)\rtimes \{-1,1\}$.
\end{enumerate}
\finenonce{007816}
\finexercice


\exercice{7817, mourougane, 2021/08/11}
\enonce{007817}{Dilatations}
Soit $f$ une forme sesquilinéaire non dégénérée. Déterminer les
dilatations 
\begin{enumerate}
\item orthogonales.
\item unitaires.
\item symplectiques.
\end{enumerate}
\finenonce{007817}
\finexercice


\exercice{7818, mourougane, 2021/08/11}
\enonce{007818}{}
Soit $k$ un corps de caractéristique différente
de $2$.
Soit $E$ un $k$-espace vectoriel de dimension $n\geq 2$ et $\phi$ une
forme sesquilinéaire
non dégénérée sur $E$ symétrique, hermitienne ou alternée. Soit
$\tau$ une transformation de $E$ donnée à l'aide d'une forme linéaire
non nulle 
$f$ sur $E$ et un vecteur $a$ de $Ker f$ par $\forall x\in E, \tau
(x)=x+f(x)a$. 
\begin{enumerate}
\item Déterminer la nature de $\tau$.
\item On suppose désormais que $\tau$ est une isométrie relativement à
 $\phi$. Montrer que $a$ est isotrope.
\item Montrer que $f$ et $\phi (\cdot , a)$ sont proportionnelles. On
 notera $\lambda\in k^\star$ tel que $f=\lambda\phi (\cdot , a)$.
\item Montrer que si $\sigma\not =Id$ et $\phi$ est hermitienne ou
 symétrique, alors $\lambda +\sigma (\lambda)=0$.
\item Montrer qu'il n'existe pas de transvections orthogonales, qu'il
 existe des transvections unitaires si et seulement si l'indice est
 plus grand que $1$ et qu'il existe toujours des transvections
 symplectiques. 
\end{enumerate}
\finenonce{007818}
\finexercice 


\exercice{7819, mourougane, 2021/08/11}
\enonce{007819}{Espace de matrices}
Soit $E:=M_n(\Rr)$. Déterminer la signature des formes quadratiques 
\begin{enumerate}
\item $q_1 (A)=trace(A ^2)$.
\item $q_2 (A)=trace(A ^2)-(trace A)^2$.
\end{enumerate}
\finenonce{007819}
\finexercice


\exercice{7820, mourougane, 2021/08/11}
\enonce{007820}{Produit de formes linéaires}
\begin{enumerate}
\item Montrer que le polyn{ô}me quadratique à coefficients complexes avec $a\not =0$
$$ax^2+a'y^2+a''z^2+2byz+2b'zx+2b''xy$$ est le produit
de deux polyn{ô}mes linéaires si et seulement si
$$aa'-{b''}^2=aa''-{b'}^2=ab-b'b''=0.$$ 
\item L'analogue à coefficients réels est-il vrai ?
\end{enumerate}
\finenonce{007820}
\finexercice


\exercice{7821, mourougane, 2021/08/11}
\enonce{007821}{Espace d'Artin}
On appelle espace d'Artin (ou espace hyperbolique) tout espace
vectoriel $E$ muni d'une forme quadratique $q$ équivalente à 
$$x=(x_i)_{1\leq i\leq 2p}\ \ \ \ \ \ \ 
 q(x)=2\sum_{i=1}^p x_ix_{i+p}.$$

\begin{enumerate}
\item Montrer que sur $\Cc^{2p}$, toute forme quadratique non dégénérée
 définit un espace d'Artin. Caractériser à l'aide de la signature les
 espaces d'Artin sur $\Rr^{2p}$. Les caractériser à l'aide de
 l'indice en supposant la forme non dégénérée.
\item Montrer que tout espace d'Artin est somme directe orthogonale de
 plans d'Artin orthogonaux.
\item Soit $(E,q)$ quelconque, $x$ un vecteur isotrope mais pas dans
 le noyau de $q$. Montrer qu'il existe un plan $P$ qui contient $x$
 et tel que $(P,q_{|P})$ soit un plan d'Artin.
\end{enumerate}
\finenonce{007821}
\finexercice


\exercice{7822, mourougane, 2021/08/11}
\enonce{007822}{Complété non singulier}
Soit $(E,q)$ un espace vectoriel muni d'une forme quadratique non
dégénérée $q$. Un sous-espace vectoriel $F$ de $E$ est dit isotrope (ou
singulier) si $F^\perp\cap F\not= \{0\}$.

Vérifier que si $F$ et $G$ sont deux sous-espaces orthogonaux avec
$G$ non isotropes alors la somme $F+G$ est directe.


 Soit $F$ un sous espace singulier de $E$. On note $s$ la
 dimension de $F^\perp\cap F$, $G$ un supplémentaire de 
$F^\perp\cap F$ dans $F$ et $(x_i)_{1\leq i\leq s}$ une base 
de $F^\perp\cap F$. 
\begin{enumerate}
\item Montrer que $G$ n'est pas isotrope.
\item Montrer que $G^\perp$ contient strictement $F^\perp$.
\item On suppose que $s=1$. Montrer qu'il existe un plan d'Artin $P_1$
 contenant $x_1$ et tel que $G\oplus ^\perp P_1$ soit non singulier. 
\item Montrer qu'il existe des plans d'Artin
 $(P_i,q_{|P_i})$ contenant $x_i$ et tels que $\widehat{F}:=G\oplus^\perp
 P_1\oplus^\perp P_2\cdots \oplus^\perp P_s$ soit non singulier.
(On pourra considérer $G':=G\oplus^\perp Vect(x_1,x_2,\cdots ,x_{s-1}$)
\item Soit $F$ un sous espace singulier de $E$.
 Montrer que si $H$ est un sous espace
 non singulier qui contient $F$ alors $\dim H\geq \dim F+\dim
 F^\perp\cap F$. 
\item Montrer que tout espace quadratique non dégénéré 
qui contient un sous-espace totalement isotrope de dimension moitié
est un espace d'Artin.
\item Montrer que toute isométrie $u~:F\to F'$ se prolonge en une
 isométrie $\widehat{u}~:\widehat{F}\to \widehat{F'}$ pour un bon
 choix de base de $F^\perp\cap F$ et ${F'}^\perp\cap F'$
\end{enumerate}
\finenonce{007822}
\finexercice


\exercice{7823, mourougane, 2021/08/11}
\enonce{007823}{}
Soit $(E, q(xe_1+ye_2)=2xy)$ un plan d'Artin. Peut-on prolonger
l'application $u$ définie sur $vect (e_1)$ par $u(e_1)=e_2$ en une
isométrie de $E$ ? En une isométrie directe de $E$ ? 
\finenonce{007823}
\finexercice


\exercice{7824, mourougane, 2021/08/11}
\enonce{007824}{Le théorème de Witt}
Soit $(E,q)$ un espace vectoriel muni d'une forme quadratique $q$ non
dégénérée. 
\begin{enumerate}
\item Montrer que si $q(x)=q(y)\not=0$ alors l'un des deux vecteurs
 $x+y$ et $x-y$ est non isotrope.
\item Montrer que si $F$ est non singulier de dimension au moins $2$,
 on peut écrire $F=F_1\oplus^\perp F_2$ avec $F_i$ non singulier de
 dimension $\dim F_i<\dim F$.
\item Soit $F=F_1\oplus^\perp F_2$ ($F_i$ non singulier) et soit $u~:F\to F'$
 une isométrie. Soit $v~: E\to E$ une isométrie qui coïncide avec $u$
 sur $F_1$.. Soit $F'_1=u(F_1)$. Montrer que ${F'_1}^\perp $
 contient $u(F_2)$ et $v(F_2)$ et que $u\circ v^{-1}~: v(F_2)\to
 u(F_2)$ est une isométrie.
\item Démontrer le théorème de Witt

Soit $F$ et $F'$ deux sous espaces de $(E,q)$ ($q$ non dégénérée)
 et $u~: (F,q_{|F})\to (F',q_{|F'})$ une isométrie. 
Montrer qu'il existe une isométrie de $E$
qui prolonge $u$. 
\end{enumerate}
\finenonce{007824}
\finexercice


\exercice{7825, mourougane, 2021/08/11}
\enonce{007825}{Action des groupes orthogonaux}
\begin{enumerate}
\item Le groupe $O(1,1)$ agit-il transitivement sur les droites de
 $\Rr^2$ ?
\item Soit $q$ une forme quadratique réelle non dégénérée.
 Montrer que si $q$ a la
 même signature en restriction à $F$ et à $F'$ alors $F$ et $F'$ sont
 dans la même orbite sous l'action de $O(q)$.
\item Décrire les orbites de l'action de $O(2,1)$ sur les droites de
 $\Rr^3$.
\end{enumerate}
\finenonce{007825}
\finexercice


\exercice{7826, mourougane, 2021/08/11}
\enonce{007826}{Etude du groupe orthogonal $O(1,1)$}
\begin{enumerate}
\item Montrer que dans un plan d'Artin il y a exactement deux droites
 isotropes $I$ et $J$.
\item Soit $u\in O(1,1)$. Montrer que $u$ envoie $I\cup J$ sur lui-même.
\item Soit $u\in O(1,1)$. Montrer que $u$ est directe si et seulement
 si $u$ laisse fixes chaque droite isotrope.
\item En déduire la forme des éléments de $O(1,1)$.
\end{enumerate}
\finenonce{007826}
\finexercice


\exercice{7827, mourougane, 2021/08/11}
\enonce{007827}{Sur les similitudes}
Soit  $f$  une forme  sesquilinéaire  non dégénérée  symétrique
(resp. hermitienne, alternée) sur un $K$-espace vectoriel $E$. On note
$GO(f)$ (resp. $GU(f)$, $GSp(f)$) le groupe des similitudes de $f$. On
note  $\mu$ l'application  qui  à une  similitude associe  son
multiplicateur dans $K^\star$.
\begin{enumerate}
\item Déterminer les similitudes de la forme symplectique standard sur
 $K^2$.
\item Montrer que $u$ est une similitude si et seulement si elle
 conserve l'orthogonalité.
\item On suppose $f$ symétrique. Montrer que $Im (\mu)=\{\lambda\in
K^\star / q\equiv \lambda q\}\supset (K^\star)^2$.
\end{enumerate}
\finenonce{007827}
\finexercice


\exercice{7828, mourougane, 2021/08/11}
\enonce{007828}{}
Déterminer l'orbite et le stabilisateur d'un vecteur de norme 1 sous
l'action du groupe orthogonal du produit scalaire standard dans
$\Rr^n$. Montrer que groupe $O(n-1)$ est isomorphe à un sous-groupe
$O_{n-1}$ de $O(n)$. Ce sous-groupe est-il distingué ?
\finenonce{007828}
\finexercice


\exercice{7829, mourougane, 2021/08/11}
\enonce{007829}{}
\begin{enumerate}
\item Montrer que l'application 
$$\begin{array}{cccc}h : &\Rr^3&\to& H_0\\
&(x_1,x_2, x_3)&\mapsto&\left(
\begin{array}{cc}
x_3&x_1+ix_2\\x_1-ix_2&-x_3
\end{array}\right)\end{array}$$
est un isomrphisme entre $\Rr^3$ et le $\Rr$-espace vectoriel
$H_0$ des matrices hermitiennes de trace nulle.
\item Montrer que le groupe $SU(2)$ agit sur $H_0$ par conjuguaison.
\item Par l'isomorphisme $h$, cette action permet de définir une action 
de $SU(2)$sur $\Rr^3$. Montrer que cette action est par isométrie de
déterminant $1$. (On pourra utiliser la connexité de $SU(2)$
homéomorphe à $S^3$ sphère unité du corps des quaternions.).
En déduire un homomorphisme $\phi$ de $SU(2)$ dans $SO(3)$.
\item Montrer que les seules matrices de $SU(2)$ qui commutent à tous
 les éléments de $H_0$ sont $Id$ et $-Id$. En déduire le noyau de
 $\phi$.
\item En utilisant les formes réduites des matrices de $SO(3)$,
 montrer que l'application exponentielle de l'espace $so(3)$
des matrices anti-symétriques réelles $3\times 3$ (de trace nulle) 
sur $SO(3)$ est surjective.
\item En utilisant la diagonalisation des matrices unitaires dans une
 base orthonormée pour le produit scalaire hermitien standard sur
 $\Cc^2$, montrer que l'application exponentielle de l'espace $su(2)$
des matrices anti-hermitiennes de trace nulle sur $SU(2)$ est surjective.
\item Déterminer l'image par $\phi$ des matrices ( avec $a,b,c\in\Rr$)
$$\exp \left(
\begin{array}{cc}
ia&0\\
0&-ia
\end{array}\right)
\ ;\
\exp \left(
\begin{array}{cc}
0&b\\
-b&0
\end{array}\right)
\ ;\
\exp \left(
\begin{array}{cc}
0&ic\\
ic&0
\end{array}\right)$$
En déduire que $\phi$ est surjective. 
\end{enumerate}
\finenonce{007829}
\finexercice


\exercice{7830, mourougane, 2021/08/11}
\enonce{007830}{}
Soit $d_1, d_2, d_3$ trois droites de $P^3(K)$. Montrer qu'il existe
une quadrique qui les contient.
\finenonce{007830}
\finexercice


\exercice{7831, mourougane, 2021/08/11}
\enonce{007831}{Polarité}
Soit $\mathcal{Q}$  une quadrique de $P(E)$  (muni d'un repère
projectif) d'équation $q(x)=0$ o{ù} $q$ est la forme quadratique d'une
forme bilinéaire symétrique non dégénérée $f$ sur un espace vectoriel
$E$. Si $\vec{A}$ et $\vec{B}$ sont deux sous-espaces orthogonaux dans
$E$, on note $P(\vec{A})\perp P(\vec{B})$. On appelle hyperplan
polaire d'un point $A=P(\vec{A})$ de $P(E)$ l'hyperplan projectif
$A^\perp:=P(\vec{A}^\perp)$.

\begin{enumerate}
\item On munit de plan projectif d'un repère. Déterminer une équation
de la droite polaire du point $M(x_0,y_0,1)$ par rapport à la
quadrique d'équation $x^2+y^2-z^2=0$.  La représenter dans l'espace
affine d'équation $z\not=0$.
\item Soit $F$ un sous-espace non-isotrope de $E$. Soit $A$ et $B$
 deux points de $P(F)$. Montrer si $A\perp B$ pour $f$ alors $A\perp
 B$ pour $f_{|F}$.
\item Soit $\mathcal{Q}$ une quadrique de $P^1(K)$ dont l'image est
 composée des deux points $A$ et $B$. Montrer en utilisant un bon
 repère que pour tout $M$ in $P^1(K)$,
$$M\perp N\iff M \textrm{ et } N \textrm{ sont conjugués harmoniques
 par rapport à } M \textrm{ et } N.$$
\item En déduire une construction géométrique de la polaire d'un point
 par rapport à une conique.
\end{enumerate}
\finenonce{007831}
\finexercice


\exercice{7832, mourougane, 2021/08/11}
\enonce{007832}{Classification euclidienne des quadriques}
On considère dans l'espace euclidien $E$ de dimension $3$ une
quadrique $\mathcal{Q}$ d'équation $q(x,y,z)=0$ o{ù} $q$ est un polyn{ô}me
de degré 2. On note $h$ sa partie homogène de degré $2$. C'est une
forme quadratique sur $\vec{E}$. On l'appelle forme quadratique à
l'infini.  On note $Q$ l'homogénéisée de $q$. C'est une forme
quadratique sur $\widehat{E}$ qui définit la complétion projective de
$\mathcal{Q}$.

\begin{enumerate}
\item Montrer que chaque argument de la signature de $Q$ est plus
 grand que l'argument correspondant de la signature de $h$.
\item On suppose que la signature de $h$ est $(3,0)$ ou $(0,3)$.
Déterminer, suivant la signature de $Q$ une forme réduite (dans un bon
repère) pour $q$ et  représenter dans chaque cas la quadrique
$\mathcal{Q}$.
\item Indiquer le résultat pour les autres signatures.
\item Affecter aux différents cas les noms suivants :

plan double réel, couple de plans imaginaires conjugués parallèles
distincts, couple de plans réels parallèles distincts,

cylindre à base parabolique, cylindre à base hyperbolique, cylindre à
base elliptique, cylindre imaginaire,

c{ô}ne imaginaire de sommet réel, c{ô}ne de base une conique propre
réelle,

ellipsoïde imaginaire, ellipsoïde réel,

paraboloïde hyperbolique, paraboloïde elliptique,

hyperboloïde à une nappe, hyperboloïde à deux nappes.
\end{enumerate}
\finenonce{007832}
\finexercice


\exercice{7833, mourougane, 2021/08/11}
\enonce{007833}{Isométries}
\begin{enumerate}
\item Soit $(V,f)$ un espace vectoriel muni d'une forme bilinéaire symétrique
 non dégénérée. Montrer que toute application $f~:V\to V$ qui
 conserve $f$ est une bijection linéaire.
\item Soit $(E,d)$ un espace affine euclidien. Montrer que toute
 application $f~:E\to E$ qui conserve la distance est une bijection
 affine dont la partie linéaire est orthogonale.
\end{enumerate}
\finenonce{007833}
\finexercice


\exercice{7834, mourougane, 2021/08/11}
\enonce{007834}{Droite et conique plane}
\begin{enumerate}
\item Soit $(A,B,C;D)$ un repère projectif du plan projectif $P^2$.
{\'E}crire en coordonnées homogènes la projection $p$ de centre $A$
 sur la droite $(BC)$.
\item Soit $\mathcal{Q}$ une quadrique projective non singulière
de $P^2$ qui passe par le point $A$. Montrer que la projection $p$
se prolonge en une bijection entre $\mathcal{Q}$ et la droite $(BC)$.
\end{enumerate}
\finenonce{007834}
\finexercice


\exercice{7835, mourougane, 2021/08/11}
\enonce{007835}{Système de droites dans une quadrique}
Soit l'espace $P^3$ muni d'un système de coordonnées homogènes
et $\mathcal{Q}$ la quadrique projective d'équation
$X_0^2+X_1^2-X_2^2-X_3^2=0$.
\begin{enumerate}
\item Déterminer un système de coordonnées homogènes dans lequel
 $\mathcal{Q}$ a pour équation $Y_0Y_3-Y_1Y_2=0$.
\item Montrer alors que l'application $v~:\left((a:b),(c:d)\right)\to 
(ac:ad:bc:bd)$ réalise une bijection de $P^1\times P^1$ sur $\mathcal{Q}$.
\end{enumerate}
\finenonce{007835}
\finexercice


\exercice{7836, mourougane, 2021/08/11}
\enonce{007836}{Avec un théorème connu}
Résoudre dans $\Zz^3$ l'équation $a^2+b^2=c^2$.
\finenonce{007836}
\finexercice


\exercice{7837, mourougane, 2021/08/11}
\enonce{007837}{}
On considère la forme symplectique sur $\Rr^{2n}$ donnée par 
$$f\left(\left({x\atop y}\right),\left({x'\atop y'}\right)\right)
=x^t\ y'-y^t\ x'$$
o{ù} $x,y,x',y'$ sont dans $\Rr^n$.
\begin{enumerate}
\item En décomposant par bloc $n\times n$ une matrique quelconque $g$ de
 $M_{2n}(\Rr)$ sous la forme $\left({AB\atop CD}\right)$,
 caractériser les matrices symplectiques en termes de systèmes
 d'équation pour $A,B,C,D$.
\item Déterminer toutes les matrices symplectiques vérifiant de plus
 $B=C=0$.
\item Déterminer toutes les matrices symplectiques vérifiant
 $A=D=I_n$ et $C=0$.
\item Déterminer toutes les matrices symplectiques vérifiant $C=0$.
\item Déterminer toutes les matrices $Q$ de $M_{n}(\Rr)$ telles que
 l'espace $W_Q:=\left\{ \left({Qy\atop y}\right)/ y\in
 \Rr^n\right\}$ 
soit totalement isotrope. Quel est le lien avec les questions
 précédentes ?
\end{enumerate}
\finenonce{007837}
\finexercice


\exercice{7838, mourougane, 2021/08/11}
\enonce{007838}{}
Soit $n$ un entier supérieur à $3$.
Le groupe $D_{2n}$ des isométries d'un polygone régulier à $n$ côtés d'un plan euclidien réel
est-il résoluble ?
\finenonce{007838}
\finexercice


\exercice{7839, mourougane, 2021/08/11}
\enonce{007839}{}
Les formes bilinéaires symétriques données par les formes quadratiques suivantes sur $\mathbb{F}_7^3$ sont-elles équivalentes ?
$$q(x,y,z)=x^2+6y^2+2z^2$$
et $$Q(x,y,z)=xy+3z^2.$$

\finenonce{007839}
\finexercice


\exercice{7840, mourougane, 2021/08/11}
\enonce{007840}{}
Démontrer le théorème de Witt dans le cas particulier suivant :

Soit $E$ et $E'$ deux espaces symplectiques non-singuliers de dimension $4$.
Soit $d\subset E$ et $d'\subset E'$ deux droites 
et $f$ une application linéaire bijective de $d$ sur $d'$.
Montrer qu'il existe une isométrie de $E$ sur $E'$ qui prolonge $f$.
\finenonce{007840}
\finexercice


\exercice{7841, mourougane, 2021/08/11}
\enonce{007841}{}
\begin{enumerate}
 \item Donner la liste des éléments du groupe $D$ des isométries qui conservent un hexagone régulier.
 \item Donner un $3$-Sylow de ce groupe.
 \item Quels sont les ordres possibles d'éléments d'un $2$-Sylow de $D$ ? 
 \item Expliciter un $2$-Sylow de $D$.
 \item Est-il distingué dans $D$ ?
 \item Combien $D$ a-t-il de $2$-Sylow ?
\end{enumerate}
\finenonce{007841}
\finexercice


\exercice{7842, mourougane, 2021/08/11}
\enonce{007842}{}
\begin{enumerate}
 \item Combien le groupe $\mathcal{A}_5$ a-t-il d'éléments d'ordre $3$ ?

 \item Montrer que les déplacements qui conservent 
un nombre fini $A_1,A_2,\ldots, A_n$ de points du plan euclidien sont des rotations d'ordre fini 
dont on précisera le centre.

 \item Le groupe $\mathcal{A}_5$ peut-il être le groupe des déplacements qui conservent 
un nombre fini de points du plan euclidien ?

 \item Le groupe $\mathcal{A}_5$ peut-il être le groupe des isométries qui conservent 
un nombre fini de points du plan euclidien ?
(Indication : Quel serait alors le sous-groupe des déplacements ?)
\end{enumerate}
\finenonce{007842}
\finexercice


\exercice{7843, mourougane, 2021/08/11}
\enonce{007843}{}
\begin{enumerate}
 \item Déterminer si possible un plan totalement isotrope pour chacune des formes bilinéaires symétriques données par les formes quadratiques suivantes sur $\mathbb{C}^4$.
$$P(x,y,z,t)=xy+zt$$
et $$Q(x,y,z)=x^2+y^2+z^2+t^2.$$
 \item L'entier $-3$ est-il un carré modulo 7 ?
 \item Déterminer si possible un vecteur non-nul isotrope pour chacune des formes bilinéaires symétriques données par les formes quadratiques suivantes sur $\mathbb{F}_7^4$.
$$R(x,y,z,t)=xy+zt$$
et $$S(x,y,z)=x^2+y^2+z^2+t^2.$$
\end{enumerate}
\finenonce{007843}
\finexercice


\exercice{7888, mourougane, 2021/08/11}
\enonce{007888}{Décomposition polaire d'un endomorphisme}
Soit $H$ un espace hermitien et $a$ un endomorphisme inversible de $H$.
Montrer que $a$ s'écrit de fa{\c c}on unique sous la forme $a=hu$ o{ù} $h$ est un
endomorphisme auto-adjoint positif et $u$ unitaire.
Déterminer $h$ et $u$ pour l'endomorphisme $a$ dont la matrice dans la base canonique de $\Cc^2$ est $\left[
\begin{array}{cc}
1&i\\ i&1
\end{array}
\right]$.
\finenonce{007888}
\finexercice


\exercice{7889, mourougane, 2021/08/11}
\enonce{007889}{Forme alternée}
Une forme bilinéaire $f$ sur un $k$-espace vectoriel $E$ est dite alternée, si tout vecteur de $E$ est isotrope.
Soit $(E,f)$ un $k$-espace vectoriel de dimension finie muni d'une forme bilinéaire alternée.
\begin{enumerate}
\item Soit $(V,f)$ un $k$-espace vectoriel de dimension $2$, muni d'une forme alternée non-dégénérée. Soit $x$ un vecteur non nul.
Montrer qu'il existe un vecteur isotrope $y$ tel que $f(x,y)=1$. On dit alors que $(V,f)$ est un plan symplectique.
\item Soit $V$ un sous-espace vectoriel de $E$.
Montrer que si $(V,f_{\mid V})$ est un espace non singulier (i.e. $f_{\mid V}$ non dégénérée)
alors $E=V\oplus^\perp V^\perp$.
\item Montrer que $E $ est somme directe orthogonale de droites isotropes et de plans symplectiques.
\item Montrer que tous les sous-espaces isotropes maximaux de $E$ ont même dimension.
Déterminer cette dimension en fonction de la dimension de $E$ et du rang de $f$.
\item Retrouver ce résultat en utilisant le théorème de Witt symplectique :
Soit $(E,f)$ et $(E',f')$ deux $k$-espaces vectoriels de dimension finie muni d'une forme symplectique (i.e. bilinéaire alternée non-dégénérée).
On suppose $(E,f)$ et $(E',f')$ isométriques.
Alors, toute isométrie d'un sous-espace de $(E,f)$ sur un sous-espace de $(E',f')$ se prolonge en une isométrie de $(E,f)$ sur $(E',f')$.
\end{enumerate}
\finenonce{007889}
\finexercice


\exercice{7890, mourougane, 2021/08/11}
\enonce{007890}{}
 Montrer qu'en dimension $2$ le groupe symplectique d'une forme alternée non dégénérée est isomorphe au groupe spécial linéaire $SL(2,k)$.
\finenonce{007890}
\finexercice


\exercice{7891, mourougane, 2021/08/11}
\enonce{007891}{}
\begin{enumerate}
\item Les formes bilinéaires symétriques données par les formes quadratiques suivantes sur $\mathbb{F}_7^3$ sont-elles équivalentes ?
$$q(x,y,z)=x^2+6y^2+2z^2$$
et $$Q(x,y,z)=xy+4z^2.$$
\item Montrer que deux formes quadratiques équivalentes sur un espace $E$ prennent les mêmes valeurs dans $k$.
\item La forme $q$ prend-elle toutes les valeurs de $\mathbb{F}_7$ ? Vérifier qu'elle prend les valeurs $3$ et $5$.
\end{enumerate}
\finenonce{007891}
\finexercice


\exercice{7892, mourougane, 2021/08/11}
\enonce{007892}{}
\begin{enumerate}
\item Montrer que deux formes quadratiques équivalentes sur un espace $E$ prennent les mêmes valeurs dans $k$.
\item Soit dans tout l'exercice $(E,f)$ un espace muni d'une forme bilinéaire symétrique non dégénérée de forme quadratique associée $q$. Montrer que si $f$ admet un vecteur non nul isotrope, la forme $q$ prend toutes les valeurs de $k$.
\item Montrer que $E$ se décompose comme somme directe orthogonale de plans hyperboliques et d'un sous-espace sur lequel la forme quadratique n'a pas de vecteur isotrope non nul.
\item Montrer que le nombre de plans hyperboliques dans une telle décomposition est indépendant de la décomposition.
\end{enumerate}
\finenonce{007892}
\finexercice


\exercice{7897, mourougane, 2021/08/11}
\enonce{007897}{}
On cherche à décrire le groupe $G$ des isométries d'un plan affine euclidien qui conservent globalement un triangle $\mathcal{T}=ABC$ isocèle en $A$ non équilatéral non aplati. 
\begin{enumerate}
 \item Déterminer deux éléments différents du groupe $G$.
\item Soit $f$ un élément de $G$. 
\begin{enumerate}
\item Montrer que $f$ a un point fixe.
\item Montrer que $f(A)=A$ et que 
($f(B)=B$ ou $f(B)=C$).
\item Montrer que $f$ est soit l'identité soit une réflexion.
\end{enumerate}
\item Écrire la table de multiplication du groupe $G$.
\end{enumerate}
\finenonce{007897}
\finexercice


\exercice{7898, mourougane, 2021/08/11}
\enonce{007898}{}
Soit $D_8$ le groupe des isométries du carré.
Déterminer un morphisme injectif de groupes de $D_8$ dans $\mathcal{S}_4$.
Les éléments $(1,3)$ et $(1,2,3,4)$ engendrent-ils le groupe symétrique $\mathcal{S}_4$ ?
\finenonce{007898}
\finexercice


\exercice{7899, mourougane, 2021/08/11}
\enonce{007899}{Sylow des groupes diédraux}
Soit $\mathcal{P}_n$ un polygone régulier à $n$ côtés dans le plan euclidien orienté.
On appelle groupe diédral $D_{n}$ le groupe des isométries de $\mathcal{P}_n$.
\begin{enumerate}
\item Parmi les translations, les rotations, les symétries orthogonales, et les symétries glissées (composées d'une symétrie orthogonale et d'une translation dans l'axe de la symétrie), décrire des isométries du plan qui conservent le polygone régulier $\mathcal{P}_n$.
\item Déterminer, à l'aide de l'action naturelle de $D_{n}$ sur l'ensemble des sommets de $\mathcal{P}_n$, le cardinal de $D_{n}$. En déduire la liste complète des éléments de $D_{n}$.
\item On suppose $n$ impair. Déterminer les $2$-Sylow de $D_{n}$ et vérifier (sans référence au cours) qu'ils sont conjugués.
\item On suppose $n=6$. Déterminer un $2$-Sylow de $D_{6}$.
Déterminer le nombre de $2$-Sylow de $D_{6}$.
Déterminer deux sous-groupes d'ordre $2$ de $D_{6}$ non conjugués dans $D_{6}$.
Donner un $3$-Sylow de $D_{6}$.
\end{enumerate}
\finenonce{007899}
\finexercice


\exercice{7900, mourougane, 2021/08/11}
\enonce{007900}{Sous-groupe fini de $SO(3)$}
\begin{enumerate}
\item Montrer le \emph{théorème de Burnside} : soit $G$ un groupe fini agissant sur un ensemble fini $E$.
Alors le nombre~$N$ d'orbites est la moyenne des cardinaux des points fixes des éléments de $G$
et aussi $$N=\frac{1}{|G|}\sum_{g\in G}\text{Card} Fix(\phi(g))=\frac{1}{|G|}\sum_{x\in E}\text{Card} stabl(x).$$
On pourra considérer $\{(x,g)\in E\times G / g\cdot x=x\}$.
\item Soit $G$ un sous-groupe fini de $SO(3)$. On considère son action sur la sphère unité.
Soit $X$ l'ensemble des points fixés par un des éléments de $G$ différents de l'identité.
Montrer que $X$ est stable par l'action de $G$.
Montrer que le stabilisateur d'un élément de $X$ est un groupe cyclique.
On notera $N$ le nombre d'orbites de l'action de $G$ sur $X$
et $n_j$ le cardinal du stabilisateur d'un élément de l'orbite $O_j$.
\item Montrer que 
$$N|G|=2(|G|-1)+\text{Card} X.$$
\item Montrer que 
$$2-\frac{2}{|G|}=\sum_{j=1}^N(1-\frac{1}{n_j}).$$
\item En déduire que $N=2$ ou $N=3$.
\item Montrer que si $N=2$, $G$ est un sous-groupe cyclique de rotations.
\item Si $N=3$, déterminer les possibilités pour les $n_j$.
\end{enumerate}
\finenonce{007900}
\finexercice


\exercice{7902, mourougane, 2021/08/11}
\enonce{007902}{Droites et quadriques}
Une quadrique d'un espace projectif $P(V)$ est le lieu des zéros d'une forme quadratique $f$ sur $V$. 
\begin{enumerate}
\item Montrer que toute quadrique qui contient trois points distincts d'une droite $d$ contient toute la droite $d$.
\item Déterminer la dimension de l'espace des quadriques de $P^3(K)$.
\item Soit $d_1, d_2, d_3$ trois droites de $P^3(K)$. Montrer qu'il existe
une quadrique qui les contient.
\end{enumerate}
\finenonce{007902}
\finexercice


\exercice{7903, mourougane, 2021/08/11}
\enonce{007903}{}
On considère le plan euclidien muni d'un un repère orthonormé ($O, \overrightarrow {\imath},\overrightarrow{\jmath}$) et la courbe $(C)$ d'équation 

\begin{center}$4x^{2} - 4xy +y^{2} -3x -y - 1= 0$\end{center}

\begin{enumerate}
 \item Montrer que $(C)$ est une parabole.
\item Trouver un repère orthonormé ($S, \overrightarrow {u_{1}},\overrightarrow{u_{2}}$) tel que $(C)$ ait une équation de la forme $ x^{2} = 2py$ dans ce repère.
\end{enumerate}
\finenonce{007903}
\finexercice


\exercice{7904, mourougane, 2021/08/11}
\enonce{007904}{}
\begin{enumerate}
  \item Construire à la règle et au compas la tangente à un cercle $\mathcal{C}$ passant par un point $A$ hors du disque délimité par $\mathcal{C}$.
  \item Soit maintenant $\mathcal{C}$ et $\mathcal{C}'$ deux cercles non concentriques. Construire les centres $\Omega$ et $O$ des homothéties qui envoient $\mathcal{C}$ sur $\mathcal{C}'$.
  \item Construire deux tangentes communes à $\mathcal{C}$ et à $\mathcal{C}'$.
\end{enumerate}
\finenonce{007904}
\finexercice

\section{350.00 Variété}

\section{351.00 Immersion, submersion, plongement}

\section{352.00 Sous-variété}


\exercice{2547, tahani, 2009/04/01}
\enonce{002547}{}
Pour $\lambda \in
\mathbb{R}$, soit $S_\lambda=\{(x_1,x_2,x_3)\in \mathbb{R}^3;
x_1^2+x_2^2-x_3^2=\lambda\}$.
\begin{enumerate}
\item D\'eterminez les $\lambda \in \mathbb{R}$ pour lesquels
$S_\lambda$ est une sous-vari\'et\'e de $\mathbb{R}^3$. Dessiner
$S_\lambda$ en fonction de $\lambda$. 
\item Pour $x,y \in \mathbb{R}^3$, soit $B(x,y)=x_1y_1+x_2y_2-x_3y_3$. Soit $x\in
S_\lambda$, exprimer $T_x S_\lambda$ \`a l'aide de $B$.
\end{enumerate}
\finenonce{002547} 


\finexercice


\exercice{2548, tahani, 2009/04/01}
\enonce{002548}{}
 On muni $\mathbb{R}^n$ de la norme $||x||=<x,x>=\sum_{i=1}^n
 x_i^2$ o\`u $x=(x_1,...,x_n)$ et $<x,y>=\sum_{i=1}^n x_iy_i$.
 Soit $u: \mathbb{R}^n \rightarrow \mathbb{R}^n$ lin\'eaire telle
 que $<u(x),y>=<x,u(y)>$ et soit $Q=\{x \in \mathbb{R}^n;
 <u(x),x>=1\}$ Montrez que $Q$ est une sous-vari\'et\'e et
 d\'eterminez le plan tangent.
\finenonce{002548} 


\finexercice


\exercice{2549, tahani, 2009/04/01}
\enonce{002549}{}
Soit $f:\mathbb{R} \rightarrow \mathbb{R}^3$ d\'efinie par
$f(\theta,\varphi)=(\cos\theta (1+1/2 \cos\varphi),\sin\theta
(1+1/2\cos \varphi), 1/2 \sin \varphi)$ et soit
$T=f(\mathbb{R}^2)$.
\begin{enumerate}
\item Soit $R_\theta$ la rotation d'angle $\theta$ autour de
$(0z)$, et soit $C=\{(1+1/2 \cos \varphi, 0, 1/2 \sin \varphi);
\varphi \in \mathbb{R}\}$. Montrer que
$f(\mathbb{R}^2)=\Cup_{\theta \in \mathbb{R}}R_\theta(C)$.
Dessiner $T$. 
\item Montrer que $f(\theta,
\varphi)=f(\theta_0,\varphi_0)$ si et seulement si il existe
$(k,l) \in \mathbb{Z}^2$ tels que $\theta=\theta_0+2k\pi$ et
$\varphi=\varphi_0+2l\pi$. 
\item Montrer que pour tout ouvert $U
\subset \mathbb{R}^2$, $f(U)$ est un ouvert de $T$. 
\item Montrer
que $T$ est une sous-vari\'et\'e de $\mathbb{R}^3$.
\end{enumerate}
\finenonce{002549} 


\finexercice


\exercice{2550, tahani, 2009/04/01}
\enonce{002550}{}
Soit $f:{\cal M}_n(\mathbb{R}) \rightarrow \mathbb{R}$ de classe
$C^\infty$ d\'efinie par $f(A)=det(A)$.
\begin{enumerate}
\item Montrer que $\lim_{\lambda \rightarrow 0}
\frac{det(I+\lambda X)-1}{\lambda}=tr(X)$ (penser au polyn\^ome
caract\'eristique). En d\'eduire $D_{Id_n}f(X)$. 
\item En remarquant que $\frac{det(A+\lambda X)-det(A)}{\lambda}$ est
\'egal \`a $det(A)\frac{det(I+\lambda A^{-1}X)-1}{\lambda}$, pour
$A$ inversible, calculer $D_Af(X)$ lorsque $A$ est inversible.
\item Montrer que $Sl_n(\mathbb{R})$ est une sous-vari\'et\'e de
${\cal M}_n(\mathbb{R})$, de dimension $n^2-1$, dont l'espace
tangent en $Id$ est $\{X \in M_n(\mathbb{R}); tr(X)=0\}$.
\end{enumerate}
\finenonce{002550} 


\finexercice


\exercice{2551, tahani, 2009/04/01}
\enonce{002551}{}
 Soit $E$ l'espace vectoriel des matrices sym\'etriques r\'eelles
 d'ordre $n$. Soit $f: {\cal M}(\mathbb{R}) \rightarrow E$
 d\'efinie par $f(A)={\ }^tAA$.
 \begin{enumerate}
 \item Montrer que $D_Af(X)={\ }^tAX+{\ }^tXA$.
 \item Soit $A \in {\cal O}_n(\mathbb{R})$, $S \in E$ et $X=1/2
 AS$. Montrer que $D_Af(X)=S$. En d\'eduire que ${\cal
 O}_n(\mathbb{R})$ est une sous-vari\'et\'e de ${\cal
 M}_n(\mathbb{R})$ de dimension $n(n-1)/2$, dont l'espace tangent
 en $Id$ est $\{X \in {\cal M}_n(\mathbb{R}); {\ }^tX=-X\}$.
 \end{enumerate}
\finenonce{002551} 


\finexercice


\exercice{2552, tahani, 2009/04/01}
\enonce{002552}{}
Soit $E$ un espace vectoriel de dimension finie, $a \in E$ et $f:
E \rightarrow E$ un diff\'eomorphisme de classe $C^1$. On suppose
que $f^n=Id$ et $f(a)=a$. On pose $A=D_af$ et $u(x)=\sum_{p=1}^n
A^{-p}f^p(x)$ pour $x \in E$.
\begin{enumerate}
\item Montrer que $u$ est un diff\'eomorphisme local en $a$ tel
que $u \circ f=A \circ u$. 
\item Soit $F$ l'ensemble des points
fixes de $f$. Montrer que $F$ est une sous-vari\'et\'e de $E$.
\item Soit $g: \mathbb{R}^2 \rightarrow \mathbb{R}^2$,
$g(x,y)=(x,y+y^3-x^2)$. Montrer que $g$ est un diff\'eomorphisme
de $\mathbb{R}^2$. En d\'eduire que $2/$ n'est plus
n\'ecessairement vrai si on supprime l'hypoth\`ese $f^n=Id$.
\end{enumerate}
\finenonce{002552} 


\finexercice


\exercice{6280, mayer, 2011/10/16}
\enonce{006280}{} 
Déterminer, parmi les sous-ensembles définis ci-dessous, ceux
qui sont des sous-variétés: \begin{enumerate}
\item $\{(x,y,z)\in \Rr^3 \; ; \; x^3+y^3+ z^3 -3xyz =1\}$;
\item $\{(x,y) \in \Rr^2 \; ; \; xy=0\}$;
\item $\{(x,y,z)\in \Rr^3 \; ; \; x^2+y^2+z^2 =1 \; et \; x^2+y^2-x=0\}$;
\item $\{(x,y) \in \Rr^2 \; ; \; y^2 = x^3\}$;
\item $\{(x,y,z)\in \Rr^3 \; ; \; x^2 + y^2 = \tan (\alpha ) z^2\}$;
\end{enumerate}
\finenonce{006280}



\finexercice


\exercice{6281, mayer, 2011/10/16}
\enonce{006281}{}  Soient $\alpha $ et $\beta $ des fonctions de
${\cal C}^1 (\Rr , \Rr )$.

 \begin{enumerate}
\item On considère l'application $\varphi : \Rr \to \Rr^3$ donnée par $\varphi (x) =
(\alpha (x) , 0, \beta (x))$. Donner des conditions à $\alpha ,
\beta$ pour que ${\cal C} = \varphi (\Rr )$ soit une
sous-variété de $\Rr^3$.
\item Soit maintenant $f:\Rr^2 \to \Rr^3$ l'application
$f(x,y) = (\alpha (x) \cos (y) , \alpha (x) \sin (y) , \beta (x)
)$. On cherche encore des conditions pour $\alpha , \beta$ sous
lesquelles  ${\cal S} =f(\Rr^2 )$ soit une sous-variété de
$\Rr^3$.
\item Notons $p=f(x,0)$, $x\in \Rr$. Quel est le lien entre les espaces
tangents $T_p {\cal S}$ et $T_p {\cal C}$.
 \end{enumerate}
\finenonce{006281}



\finexercice


\exercice{6282, mayer, 2011/10/16}
\enonce{006282}{} 
\begin{enumerate}
\item Montrer que l'équation $xy+xz+yz+2x+2y-z=0$ définit au
voisinage de $(0,0,0)$ une surface. Donner l'équation du plan
tangent de cette surface à l'origine.
\item Montrer que les équations $4xy+2xz+4y-z=0$ et $
xy+xz+yz+2x+2y-z=0$ définissent au voisinage de l'origine une
courbe. Déterminer l'espace tangent de cette courbe à
l'origine. \end{enumerate}
\finenonce{006282}



\finexercice


\exercice{6283, mayer, 2011/10/16}
\enonce{006283}{}  Soit $F=(F_1,...,F_k)$ une application $C^1$ d'un
ouvert $U$ de $\Rr^m$ dans $\Rr^k$. Notons $M=\{x\in U\; ; \;
F(x) =0 \}$ et soit $a\in M$. 
\begin{enumerate}
\item \'Etablir l'équivalence des propriétés suivantes:
\begin{itemize}
\item $DF(a)$ est surjective.
\item Les formes linéaires $DF_1(a), ..., DF_k(a)$ sont
linéairement indépendantes.
\item $\mathrm{Ker}\, DF(a) = \bigcap_{i=1}^k \mathrm{Ker}\, DF_i(a)$ est de dimension
$m-k$.
\end{itemize}
\item Un point $a\in M$ est dit {\it point régulier}  si
$DF(a)$ est surjective. Montrer que l'ensemble des points
réguliers de $M$ est un ouvert de $M$. 
\end{enumerate}
\finenonce{006283}



\finexercice


\exercice{6284, mayer, 2011/10/16}
\enonce{006284}{} 
\label{exn14}
Soit $f: \Rr^n \to \Rr$ un polyn\^ome homogène de degré
$\alpha >0$ à $n$ variables. 
\begin{enumerate}
\item En calculant la dérivée de $\lambda \mapsto f(\lambda x)$ de
deux manières différentes, établir l'identité d'Euler:
$$ \sum_{i=1}^n x_i \frac{\partial f}{\partial x_i} = \alpha\, f(x)
\quad  \text{pour tout} \;\; x\in \Rr^n \; .$$
\item Soit $a$ un réel non nul. Montrer que $X_a =
f^{-1}(\{a\})$ est une sous-variété de dimension $n-1$ de
$\Rr^n$. \'Etablir ensuite que, pour $a_1>a_2>0$, $X_{a_1}$ et
$X_{a_2}$ sont difféomorphes.
\item Supposons que $\varphi$ est un difféomorphisme de $\Rr^n$ avec $\varphi ( X_{a_1} )
=X_{a_2}$ et soit $p\in X_{a_1}$. Exprimer l'espace tangent
$T_{\varphi (p)}X_{a_2}$ en fonction de $T_pX_{a_1}$. 
\end{enumerate}
\finenonce{006284}



\finexercice


\exercice{6285, mayer, 2011/10/16}
\enonce{006285}{} 
Soit $f:M_n(\Rr ) \to \Rr$ l'application $C^\infty$ donnée
par $f(A) = \det(A)$. 
\begin{enumerate}
\item Montrer que
$$\lim_{\lambda \to 0} \frac{\det(I+\lambda X) -1}{\lambda} = \mathrm{tr}(X) \quad ,
\;\; X\in M_n(\Rr )\; .$$ En déduire $Df(I)(X)$.
\item En remarquant que
$$ \frac{\det(A+\lambda X)-\det(A)}{\lambda} = \det(A) \frac{\det(I+\lambda A^{-1}X
-1)}{\lambda}\; ,$$ pour $A$ une matrice inversible, calculer
$Df(A)(X)$ lorsque $A$ est inversible.
\item Montrer que $Sl_n(\Rr) =\{A\in M_n(\Rr ) \; ; \; \det(A) =1\} $ est une sous-variété de
$M_n(\Rr )$ de dimension $n^2-1$ (on pourra faire le lien avec
l'exercice \ref{exn14}) dont l'espace tangent en $I$ est
$$T_I Sl_n(\Rr) = \{X\in M_n(\Rr ) \; ; \; \mathrm{tr}(X) =0\}\; .$$ 
\end{enumerate}
\finenonce{006285}



\finexercice


\exercice{6286, mayer, 2011/10/16}
\enonce{006286}{} 
Soit $E$ un espace vectoriel de dimension finie, $a\in E$ et
$f:E\to E$ un difféomorphisme de classe $C^1$. On suppose que
$f^n =id$ et $f(a)=a$. On pose $A=Df(a)$ et $u(x) = \sum_{p=1}^n
A^{-p} f^p(x)$ pour $x\in E$. 
\begin{enumerate}
\item Montrer que $u$ est un difféomorphisme local en $a$ tel
que $u\circ f = A\circ u$.
\item Soit $F$ l'ensemble des points fixes de $f$. Montrer que $F$
est une sous-variété de $E$.
\item Soit $g:\Rr^2\to \Rr^2$, $g(x,y) = (x,y+y^3-x^2)$.
Montrer que $g$ est un difféomorphisme de $\Rr^2$. En
déduire que $2)$ n'est plus nécessairement vrai si on supprime
l'hypothèse $f^n =id$.
\end{enumerate}
\finenonce{006286}



\finexercice


\exercice{7642, mourougane, 2021/08/11}
\enonce{007642}{Hélice}
\begin{enumerate}
 \item Soient $r$ et $h$ deux nombres réels strictement positifs.
La courbe paramétrée suivante, appelée hélice, est-elle régulière ?
$$\begin{array}{cccc}
        c~:~&\Rr&\to&\Rr^3\\
&t&\mapsto&\left(\begin{array}{c}r\cos t\\r\sin t\\ht\end{array}\right)
        \end{array}
$$

 \item Les courbes paramétrées suivantes sont-elle régulières ?
$$\begin{array}{cccc}
        d~:~&\Rr&\to&\Rr^2\\
&t&\mapsto&\left(\begin{array}{c}t^2\\ t^3\end{array}\right)
        \end{array} \quad \text{ et } \quad 
\begin{array}{cccc}
        e~:~&]0,+\infty[&\to&\Rr^2\\
&t&\mapsto&\left(\begin{array}{c}t^2\\ t^3\end{array}\right)
        \end{array}
$$

\end{enumerate}
\finenonce{007642}
\finexercice


\exercice{7643, mourougane, 2021/08/11}
\enonce{007643}{Reparamétrage d'une courbe régulière}
\begin{enumerate}
 \item Montrer que le reparamétrage d'une courbe régulière est encore régulier.
 \item Peut-on reparamétrer toute courbe paramétrée par sa longueur d'arc ?
 \item Déterminer un reparamétrage de la courbe
$\left\{\begin{array}{cccc}
        c~:~&[0,2\pi[&\to&\Rr^2\\
&t&\mapsto&\left(\begin{array}{c}\cos t\\ \sin t\end{array}\right)
        \end{array}\right.$
qui en change l'orientation.
\end{enumerate}
\finenonce{007643}
\finexercice


\exercice{7644, mourougane, 2021/08/11}
\enonce{007644}{Equivalence}
\begin{enumerate}
 \item Les courbes paramétrées suivantes sont-elles équivalentes ?
$$\begin{array}{cccc}
        c~:~&[0,2\pi[&\to&\Rr^2\\
&t&\mapsto&\left(\begin{array}{c}\cos t\\ \sin t\end{array}\right)
        \end{array}
\quad \text{ et } \quad 
\begin{array}{cccc}
        d~:~&[0,2\pi[&\to&\Rr^2\\
&t&\mapsto&\left(\begin{array}{c}\cos 2t\\ \sin 2t\end{array}\right)
        \end{array}
$$

 \item Les courbes paramétrées suivantes sont-elles équivalentes ?
$$\begin{array}{cccc}
        c~:~&\Rr&\to&\Rr^2\\
&t&\mapsto&\left(\begin{array}{c}\cos t\\ \sin t\end{array}\right)
        \end{array}
\quad \text{ et } \quad 
\begin{array}{cccc}
        d~:~&\Rr&\to&\Rr^2\\
&t&\mapsto&\left(\begin{array}{c}\cos 2t\\ \sin 2t\end{array}\right)
        \end{array}
$$

 \item Deux courbes paramétrées équivalentes ont-elles même image ?
La réciproque est-elle vraie ?

 \item Peut-on retrouver une courbe à partir de son image ?
\end{enumerate}
\finenonce{007644}
\finexercice


\exercice{7645, mourougane, 2021/08/11}
\enonce{007645}{Calcul de longueur}
\begin{enumerate}
 \item Calculer la longueur de la portion d'hélice 
$$\begin{array}{cccc}
        c~:~&[0,1]&\to&\Rr^3\\
&t&\mapsto&\left(\begin{array}{c}r\cos t\\r\sin t\\ht\end{array}\right)
        \end{array}
$$
 \item Déterminer un paramétrage de cette portion d'hélice par sa longueur d'arc.
\end{enumerate}
\finenonce{007645}
\finexercice


\exercice{7646, mourougane, 2021/08/11}
\enonce{007646}{Différents paramétrages par longueur d'arc}
Montrer que deux paramétrages par longueur d'arc d'une même courbe régulière orientée
diffèrent au plus par un changement de paramétrage de la forme 
$t\mapsto t_0+t.$
\finenonce{007646}
\finexercice


\exercice{7647, mourougane, 2021/08/11}
\enonce{007647}{Approximation par lignes polygonales}
 Vérifier l'approximation polygonale pour un cercle de rayon $r$
approché par des polygones réguliers à $n$ côtés.
\finenonce{007647}
\finexercice


\exercice{7648, mourougane, 2021/08/11}
\enonce{007648}{Les cycloïdes}
 On considère un cercle $\mathcal{C}_1$ de rayon $1$ qui glisse sur l'axe des $x$ du plan euclidien orienté $\Rr ^2$,
et pour tout nombre réel $r$ strictement positif le cercle $\mathcal{C}_r$ de rayon $r$, concentrique avec $\mathcal{C}_1$
et solidement attaché à $\mathcal{C}_1$.
\begin{enumerate}
 \item Déterminer par un paramétrage, la trajectoire $T_r$, appelée ``cycloïde`` 
du point $M$ de coordonnées $(0,1-r)$ de $\mathcal{C}_r$ lorsque $\mathcal{C}_1$ glisse sur l'axe des $x$.
 \item Faire une ébauche de $T_r$ suivant la position de $r$ par rapport à $1$.
 \item $T_r$ est-elle une courbe régulière ?
 \item Calculer la longueur de $T_1$ quand $\mathcal{C}_1$ fait un tour.
\end{enumerate}
\finenonce{007648}
\finexercice


\exercice{7649, mourougane, 2021/08/11}
\enonce{007649}{Approximation par lignes polygonales}
\begin{enumerate}
 \item On rappelle qu'une ligne polygonale $P$ de $\Rr^n$ est la donnée d'un uplet $P=(a_0,a_1,\cdots,a_k)$
de points de $\Rr^n$. On supposera aussi que deux points consécutifs sont distincts.
Rappeler la formule pour la longueur d'une ligne polygonale $P$.
 \item On cherche à montrer le théorème suivant.
 
\textbf{Théorème.}
\emph{
Soit $c~:~[a,b]\to \Rr^n$ une courbe paramétrée.
Alors, pour tout $\epsilon>0$, il existe $\delta>$ tel que 
pour toute partition $(t_0=a<t_1<t_2<\cdots<t_m=b)$ de $[a,b]$ de pas inférieur à $\delta$,
$$L[P]\leq L[c]\leq L[P]+\epsilon$$
où $P=(c(t_0),c(t_1),\cdots,c(t_m))$ est la ligne polygonale simplement inscrite dans la courbe $c$
associée à la partition $(t_0=a<t_1<t_2<\cdots<t_m=b)$ de $[a,b]$.
}

\begin{enumerate}
\item Écrire en termes de quantificateurs, en partant d'un $\epsilon_1>0$,
le théorème des sommes de Riemann pour la fonction $[a,b]\to\Rr$, 
$t\mapsto \| \dot{c} (t)\|$.
\item \'Ecrire en termes de quantificateurs en partant d'un $\epsilon_2>0$
la propriété de continuité uniforme des fonctions composantes $[a,b]\to\Rr$, 
$t\mapsto \dot{c}_j (t)$.
\item Soit une partition $(t_0=a<t_1<t_2<\cdots<t_m=b)$ de $[a,b]$.
\'Ecrire en termes de quantificateurs le théorème des accroissements finis
pour la fonction $[t_i,t_{i+1}]\to\Rr$, $t\mapsto {c}_j (t)$.
\end{enumerate}

 \item Majorer la quantité 
$\left| \| c(t_{i+1})-c(t_i)\| -\|\dot{c} (t_{i+1})\| (t_{i+1}-t_i)\right|$
par $\sqrt{n}\epsilon_2(t_{i+1}-t_i)$.
 \item Majorer la quantité $\left|L[c]-L[P]\right|$ par $\epsilon_1+\sqrt{n}\epsilon_2(b-a)$.
 \item Montrer que pour tout $\epsilon>0$, il existe $\delta>$ tel que 
pour toute partition $(t_0=a<t_1<t_2<\cdots<t_m=b)$ de $[a,b]$ de pas inférieur à $\delta$,
$$L[P]-\epsilon\leq L[c]\leq L[P]+\epsilon.$$
 \item Conclure à l'aide de l'inégalité triangulaire sur les polygones.
\end{enumerate}
\finenonce{007649}
\finexercice


\exercice{7650, mourougane, 2021/08/11}
\enonce{007650}{Calculs de courbure}
On considère l'espace affine euclidien orienté $\Rr^2$ muni d'un repère $(O,\vec{\imath},\vec{\jmath})$.
\begin{enumerate}
 \item Calculer la fonction courbure d'un cercle de rayon $r>0$.
 \item Soit $c~:I\to \Rr^2$ une courbe plane paramétrée par la longueur d'arc. 
Soit $R$ la rotation de centre $0$ et d'angle $\alpha$,
et $S$ la symétrie d'axe $x=y$.
Déterminer la fonction courbure de $R\circ c$ et celle de $S\circ c$.
 \item Calculer la fonction courbure de la courbe de Lissajous
$\left\{\begin{array}{cccc}
        c~:~&\Rr&\to&\Rr^2\\
&t&\mapsto&\left(\begin{array}{c}\cos 3t\\ \sin 2t\end{array}\right)
        \end{array}
\right.$
 \item Déterminer une courbe fermée à courbure partout strictement négative.
 \item Montrer qu'une courbe plane régulière de courbure nulle est un segment de droite. 
\end{enumerate}
\finenonce{007650}
\finexercice


\exercice{7651, mourougane, 2021/08/11}
\enonce{007651}{Courbure en un point extrémal}
Soit $c~:I\to \Rr^2$ une courbe plane paramétrée par la longueur d'arc. 
On suppose que $c$ reste dans le disque de rayon $r$ et qu'au point de paramètre $\tau$,
$\| c(\tau)\|=r$. 
\begin{enumerate}
 \item Montrer en dérivant une fois la fonction $t\mapsto \| c(t)\|^2$ 
que $\ddot{c}(\tau)$ est colinéaire à $c(\tau)$ ?
 \item Montrer en dérivant à nouveau la fonction $t\mapsto \| c(t)\|^2$ 
que la courbure en $\tau$ vérifie $|\kappa(\tau)|\geq 1/r$.
\end{enumerate}
\finenonce{007651}
\finexercice


\exercice{7652, mourougane, 2021/08/11}
\enonce{007652}{Construction d'une courbe plane à courbure prescrite}
Le but de l'exercice est de montrer le théorème suivant.

\textbf{Théorème.} 
\emph{
Soit $I$ un intervalle et $\kappa~:~I\to\Rr$ une fonction de classe $\mathcal{C}^\infty$.
Alors, il existe une courbe plane $c~:~I\to\Rr^2$ paramétrée par la longueur d'arc
et de fonction courbure $\kappa$. 
Cette courbe est unique à composition au but par un déplacement près.
}

\begin{enumerate}
 \item On considère le système d'équations différentielles linéaire du premier ordre
$$\frac{d}{dt}\left(\begin{array}{c}c\\v\\n\end{array}\right)
=\left(\begin{array}{ccc}0&1&0\\0&0&\kappa\\0&-\kappa&0\end{array}\right)
\left(\begin{array}{c}c\\v\\n\end{array}\right).$$
où les fonctions $c,v,n~:~I\to\Rr^2$ sont les inconnues.
 Montrer que ce système admet une unique solution $(c(t),v(t),n(t))$ 
avec pour valeur initiale un vecteur fixé $(c_0,v_0,n_0)$ 
avec $(v_0,n_0)$ base orthonormée directe.
 \item \'Ecrire un système d'équations différentielles linéaire du premier ordre
satisfait par le vecteur $(<v,v>,<n,n>,<v,n>)$.
 \item Montrer que pour les solutions obtenues précédemment, $(v(t),n(t))$
reste un repère orthonormé direct.
 \item En déduire que la courbe $c$ obtenue est paramétrée par la longueur d'arc
et que sa fonction courbure est $\kappa$.
 \item Conclure.
\end{enumerate}
\finenonce{007652}
\finexercice


\exercice{7653, mourougane, 2021/08/11}
\enonce{007653}{Nombre de rotation}
 Calculer le nombre de rotation de la courbe régulière fermée suivante
\begin{center}
\includegraphics[scale=0.5]{images/img-mour-405}
\end{center}
\finenonce{007653}
\finexercice


\exercice{7654, mourougane, 2021/08/11}
\enonce{007654}{Courbure totale}
\begin{enumerate}
 \item Calculer la somme des mesures des angles extérieurs d'un triangle et d'un carré.
 \item Déterminer la fonction courbure et la courbure totale de l'ovale
paramétré par $c~:~[0,2\pi]\to\Rr^2$, $t\mapsto(a\cos(t),b\sin(t))$.
\end{enumerate}
\finenonce{007654}
\finexercice


\exercice{7655, mourougane, 2021/08/11}
\enonce{007655}{Formules de Frenet}
 Démontrer les formules de Frenet pour les courbes paramétrées par la longueur d'arc 
dans $\Rr^3$.
\finenonce{007655}
\finexercice


\exercice{7656, mourougane, 2021/08/11}
\enonce{007656}{Torsion}
\begin{enumerate}
 \item Calculer la torsion de l'hélice
$\left\{\begin{array}{cccc}
        c~:~&\Rr&\to&\Rr^3\\
&t&\mapsto&\left(\begin{array}{c}r\cos t\\r\sin t\\ht\end{array}\right)
        \end{array}
\right.$.
 \item Calculer la torsion de la courbe de $\Rr^3$ paramétrée 
par $t\mapsto (4\cos t,5-5\sin t, -3\cos t)$.
 \item La courbe précédente est-elle plane ?
\end{enumerate}
\finenonce{007656}
\finexercice


\exercice{7657, mourougane, 2021/08/11}
\enonce{007657}{Courbes sur une surface}
Soit $C$ la courbe tracée sur la surface d'équation $3z=xy+x^3$
et dont la projection orthogonale sur le plan d'équation $z=0$ 
est la courbe paramétrée $C'$ définie par $x=t$, $y=t^2$ pour $t$ parcourant $[0,1]$.
\begin{enumerate}
 \item Donner une expression intégrale pour la longueur de $C'$ puis celle de $C$, puis les comparer.
 \item Calculer la longueur de $C'$ puis celle de $C$.
\end{enumerate}
\finenonce{007657}
\finexercice


\exercice{7658, mourougane, 2021/08/11}
\enonce{007658}{Construction des courbes de $\Rr^3$}
 Énoncer pour les courbes de l'espace $\Rr^3$, le théorème analogue à la caractérisation
des courbes planes à déplacements près par leur courbure.
\finenonce{007658}
\finexercice


\exercice{7659, mourougane, 2021/08/11}
\enonce{007659}{Difféomorphisme local, global}
\begin{enumerate}
  \item Soit $U$ le plan $\Rr^2$ privé de l'origine.
Montrer que l'application 
\begin{eqnarray*}
  f : U&\to& U\\
(x,y)&\mapsto& (x^2-y^2,2xy)
 \end{eqnarray*}
est un difféomorphisme au voisinage de chacun des points de $U$.
 \item Est-ce un difféomorphisme de $U$ ?
\end{enumerate}
\finenonce{007659}
\finexercice


\exercice{7660, mourougane, 2021/08/11}
\enonce{007660}{Sous-variété de $\Rr^2$}
 Soit $$\mathcal{C}:=\{(x,y)\in\Rr^2, x^3+y^3-3xy=0\}.$$
 \begin{enumerate}
 \item Au voisinage de quels points de $\mathcal{C}$, cette équation définit-elle $y$ comme fonction de $x$ ?
Quelle est alors la dérivée de cette fonction ?
 \item Paramétrer $\mathcal{C}$ à l'aide de $t$ tel que $y=tx$.
 \item Représenter $\mathcal{C}$ avec ses asymptotes.
 \item Le sous-ensemble $\mathcal{C}$ est-il une sous-variété différentiable de $\Rr^2$ ?
\end{enumerate}
\finenonce{007660}
\finexercice


\exercice{7661, mourougane, 2021/08/11}
\enonce{007661}{Surfaces paramétrées}
La surface paramétrée par $(u,v)\mapsto (u+v^2,u^2-v^2,v)$ est-elle régulière
au voisinage de $A(1,-1,-1)$ ?
\finenonce{007661}
\finexercice


\exercice{7662, mourougane, 2021/08/11}
\enonce{007662}{Surfaces implicites}
\begin{enumerate}
 \item L'ensemble d'équation $x^2+y^2+z^2=4$ est-il une surface régulière de $\Rr^3$.
Le décrire.
 \item L'ensemble d'équation $(x^2+y^2+z^2-4)^2=0$ est-il une surface régulière de $\Rr^3$.
Le décrire.
 \item L'ensemble d'équation $x^2+y^2-z^2=0$ est-il une surface régulière de $\Rr^3$.
Le décrire. Montrer que tous les chemins du point $A(1,0,1)$ au point $B(1,0,-1)$ 
passent par un même point à déterminer.
 \item Les ensembles d'équations respectives
\begin{eqnarray*}
 a)x^2+3y^2+z^2+4xy+2z=5 &&b)2x^2+3y^2=1\\ c)\frac{x^2}{2}+\frac{y^2}{3}-\frac{z^2}{4}=1&&d)\frac{x^2}{2}+\frac{y^2}{3}-\frac{z^2}{4}=-1
\end{eqnarray*}
sont-ils des surfaces régulières de $\Rr^3$ ?
Les décrire.
\end{enumerate}
\finenonce{007662}
\finexercice


\exercice{7663, mourougane, 2021/08/11}
\enonce{007663}{Paramétrages de la sphère}
 On considère la sphère $S$ unité dans $\Rr^3$.
 \begin{enumerate}
 \item Déterminer par une expression en coordonnées de
la projection stéréographique d'un ouvert $U$ de $S$ (à déterminer)
depuis le pôle nord sur le plan d'équation $z=0$.
Vérifier qu'il s'agit bien de paramétrage de $U$.
 \item Fa\^ites de même depuis le pôle sud.
 \item Montrer que le changement de paramétrage est de classe $\mathcal{C}^\infty$.
 \item Interpréter les coordonnées sphériques comme un paramétrage de la sphère $S$.
\end{enumerate}
\finenonce{007663}
\finexercice


\exercice{7670, mourougane, 2021/08/11}
\enonce{007670}{Paramétrage}
On considère l'application
$$\begin{array}{ccc}
 F: \Rr^2&\to&\Rr^3\\ 
\left(\begin{array}{c}t \\ s\end{array}\right)
&\mapsto& 
\left(\begin{array}{c}1-t^2\\t(1-t^2)\\s\end{array}\right).
\end{array}$$
\begin{enumerate}
 \item Déterminer des déplacements de l'espace $\Rr^3$ qui conservent l'image de $F$.
 \item Montrer que l'image de $F$ est l'ensemble d'équation $y^2=x^2(1-x)$.
 \item L'image de $F$ est-elle une sous-surface différentiable de $\Rr^3$ ?
\end{enumerate}
\finenonce{007670}
\finexercice


\exercice{7671, mourougane, 2021/08/11}
\enonce{007671}{Applications différentiables entre surfaces}
Soit $S_1$ et $S_2$ deux sous-surfaces différentiables de $\Rr^3$.
\begin{enumerate}
 \item Soit $g~:\Rr^3\to\Rr^3$ une application de classe $\mathcal{C}^\infty$.
On suppose que $g(S_1)\subset S_2$. On note $f : S_1\to S_2$ la restriction de $g$ à $S_1$.
Montrer $f$ est une application différentiable.
 \item Soit $f : S_1\to S_2$ une application différentiable et $p$ un point de $S_1$.
Montrer qu'il existe un voisinage de $p$ dans $\Rr^3$ sur lequel $f$ se prolonge en une application $\mathcal{C}^\infty$ 
à valeurs dans $\Rr^3$.
 \item Montrer que la composée de deux applications différentiables entre sous-surfaces différentiables de $\Rr^3$
est différentiable et expliciter la différentielle de la composée. 
\end{enumerate}
\finenonce{007671}
\finexercice


\exercice{7672, mourougane, 2021/08/11}
\enonce{007672}{Différentielle}
Soit $S_1$ d'équation $x^6+y^6+z^6=1$ et $S_ 2$ d'équation $x^2+y^2+z^2=1$.
\begin{enumerate}
 \item Montrer que $S_1$ et $S_2$ sont deux sous-surfaces différentiables de $\Rr^3$.
 \item On considère l'application $f$ de $S_1$ vers $S_2$, qui à $p$ de coordonnées $(x,y,z)$
associe le point de coordonnées $(x^3, y^3, z^3)$. Montrer que $f$ est bijective de $S_1$ sur $S_2$. 
 \item Montrer que $f$ est différentiable.
 \item Déterminer la différentielle de $f$. Est-elle inversible ?
 \item La bijection réciproque $f^{-1}$ est-elle différentiable ?
 \item Reprenez l'exercice pour l'application $g$ de $S_1$ vers le plan d'équation $z=0$,
qui à $p$ de coordonnées $(x,y,z)$
associe le point de coordonnées $(x,x,0)$.
\end{enumerate}
\finenonce{007672}
\finexercice


\exercice{7673, mourougane, 2021/08/11}
\enonce{007673}{Théorème d'inversion locale}
Énoncer et démontrer le théorème d'inversion locale pour une application entre deux sous-surfaces différentiables de $\Rr^3$.
\finenonce{007673}
\finexercice


\exercice{7674, mourougane, 2021/08/11}
\enonce{007674}{Extrema}
\begin{enumerate}
 \item Trouver les extrema de la fonction $f: (x,y,z)\mapsto xy$ sur la sphère unité.
 \item Trouver les extrema de la fonction $f: (x,y,z)\mapsto xy$ sur l'ellipsoïde de $\Rr^3$ d'équation $x^2+y^2+3z^2=1$.
\end{enumerate}
\finenonce{007674}
\finexercice


\exercice{7675, mourougane, 2021/08/11}
\enonce{007675}{Première forme fondamentale}
Calculer la première forme fondamentale du graphe $S$ d'une application $\varphi : \Rr^2\to \Rr$ différentiable. On pourra déterminer d'abord un paramétrage global de $S$.
\finenonce{007675}
\finexercice


\exercice{7676, mourougane, 2021/08/11}
\enonce{007676}{Dépendance en le paramétrage}
\begin{enumerate}
 \item On considère le plan d'équation $z=0$ avec le paramétrage 
$\left\{\begin{array}{ccc}
 F: \Rr^2&\to&\Rr^3\\ (x,y)&\mapsto& (x,y,0).
\end{array}\right.$\\
Déterminer la matrice de la première forme fondamentale dans la base $\mathcal{B}_F$ correspondante.
  \item On considère le plan d'équation $z=0$ avec le paramétrage local 
\begin{eqnarray*} G: ]0,+\infty[\times ]-\pi,\pi[&\to&\Rr^3\\ (r,\theta)&\mapsto& (r\cos\theta,r\sin\theta,0).\end{eqnarray*}
Déterminer la matrice de la première forme fondamentale dans la base $\mathcal{B}_G$ correspondante.
 \item Déterminer la matrice de changement de base au point de coordonnées $(1,0,0)$ de la base $\mathcal{B}_F$
dans la base $\mathcal{B}_G$. Relier les deux matrices obtenues pour la première forme fondamentale.
\end{enumerate}
\finenonce{007676}
\finexercice


\exercice{7677, mourougane, 2021/08/11}
\enonce{007677}{Sphère}
\begin{enumerate}
 \item Calculer l'aire de la sphère en coordonnées sphériques.
 \item Calculer l'aire de l'ellipsoïde $\mathcal{E}$ d'équation $x^2+y^2+5z^2=1$.
\end{enumerate}
\finenonce{007677}
\finexercice


\exercice{7678, mourougane, 2021/08/11}
\enonce{007678}{Calcul d'aire}
On considère l'application $f$ de la sphère unité $\mathcal{S}$ de $\Rr^3$ 
à valeurs dans le cylindre 
$\mathcal{C}$ de $\Rr^3$ d'équation $x^2+y^2=1$ donnée en coordonnées cartésiennes par
$$f~: \begin{pmatrix}x\\y\\z\end{pmatrix}\mapsto \begin{pmatrix}\frac{x}{\sqrt{x^2+y^2}}\\ \frac{y}{\sqrt{x^2+y^2}}\\ 
z\end{pmatrix}.$$
\begin{enumerate}
 \item Faire une figure pour décrire géométriquement cette application.
 \item On considère un paramétrage local de la sphère unité $\mathcal{S}$ 
en coordonnées polaires
$$\begin{array}{ccc}
 F : (\theta,\phi)\in ]0,2\pi[\times]0,\pi[&\mapsto&\begin{pmatrix}\cos\theta\sin\phi\\ 
\sin\theta\sin\phi\\ \cos\phi\end{pmatrix}.
\end{array}$$
Déterminer la première forme fondamentale de la sphère dans cette paramétrisation.
 \item On considère l'application $G=f\circ F$. Montrer que c'est un paramétrage local 
du cylindre $\mathcal{C}$
et calculer la première forme fondamentale du cylindre dans cette paramétrisation.
 \item L'application $f$ est-elle un difféomorphisme local ? Un difféomorphisme de $\mathcal{S}$
sur $\mathcal{C}$ ?
 \item L'application $f$ est-elle une isométrie locale ? Conserve-t-elle les aires ?
 \item En déduire l'aire de la portion de sphère entre deux grands cercles passant 
par les poles nord
et sud et faisant entre eux un angle de mesure $\alpha$. 
Vérifier en déterminant l'aire de la sphère.
\end{enumerate}
\finenonce{007678}
\finexercice


\exercice{7679, mourougane, 2021/08/11}
\enonce{007679}{Applications conformes}
 Soit $g$ et $g'$ deux produits scalaires euclidiens sur $\Rr^n$.
Montrer en utilisant la règle des sinus dans un triangle $a/\sin \hat{A}=b/\sin \hat{B}$
que $g$ et $g'$ définissent les mêmes mesures d'angle non orientés
si et seulement s'il existe une constante $c>0$ telle que $g=cg'$.
\finenonce{007679}
\finexercice


\exercice{7680, mourougane, 2021/08/11}
\enonce{007680}{Cartes}
On cherche des applications d'un ouvert du plan $\Rr^2$ dans un ouvert de la sphère $S^2$ de $\Rr^3$
qui conservent les longueurs, ou les mesures d'angles, ou les aires.
\begin{enumerate}
 \item Traduire chacune de ces trois propriétés à l'aide des formes fondamentales.
 \item Les coordonnées sphériques donnent-elles une application 
qui conserve les longueurs, ou les mesures d'angles, ou les aires ?
 \item La projection stéréographique depuis le pôle nord 
conserve-t-elle les longueurs, ou les mesures d'angles, ou les aires ?
 \item L'application du cylindre dans la sphère
$$\begin{array}{ccc}
  ]0,2\pi[\times ]-1,1[&\to& S^2\\ (\theta, h)&\mapsto&
\left( \begin{array}{c}\sqrt{1-h^2} \cos\theta\\ \sqrt{1-h^2} \sin\theta\\ h\end{array}\right)
 \end{array}$$
conserve-t-elle les longueurs, ou les mesures d'angles, ou les aires ?
 \item L'application du cylindre dans la sphère
$$\begin{array}{ccc}
  ]0,2\pi[\times ]-\infty,+\infty[&\to& S^2\\ (\theta, x)&\mapsto&
\left( \begin{array}{c}\sqrt{1-h^2(x)} \cos\theta\\ \sqrt{1-h^2(x)} \sin\theta\\ h(x)\end{array}\right)
 \end{array}$$
avec $h(x)=\tanh (x)$ conserve-t-elle les longueurs, ou les mesures d'angles, ou les aires ?
\end{enumerate}
\finenonce{007680}
\finexercice


\exercice{7681, mourougane, 2021/08/11}
\enonce{007681}{Sphère}
Calculer la seconde forme fondamentale de la sphère. 
En déduire sa courbure de Gauss et ses directions principales.
\finenonce{007681}
\finexercice


\exercice{7682, mourougane, 2021/08/11}
\enonce{007682}{Cylindre}
\begin{enumerate}
 \item Déterminer les plans tangents et un champs de vecteurs normaux au cylindre $S^1\times \Rr$.
 \item Déterminer l'application de Weingarten en chaque point $p$ du cylindre.
 \item En déduire sa courbure de Gauss et ses directions principales.
 \item Reprendre les questions précédentes pour le paraboloïde hyperbolique d'équation $z=y^2-x^2$
au point $p(0,0,0)$.
\end{enumerate}
\finenonce{007682}
\finexercice


\exercice{7683, mourougane, 2021/08/11}
\enonce{007683}{Point hyperbolique/elliptique}
Donner l'exemple d'une sous-surface différentiable de $\Rr^3$ avec un point hyperbolique et un point elliptique.
\finenonce{007683}
\finexercice


\exercice{7684, mourougane, 2021/08/11}
\enonce{007684}{Etude jusqu'à l'ordre 2}
Soit $\Sigma$ le graphe de la fonction $f(x,y) = xy$ de $\Rr^2$ dans $\Rr$. 
\begin{enumerate}
 \item Est-ce une sous-surface différentiable de $\Rr^3$ ?
 \item Déterminer une paramétrisation de $\Sigma$.
 \item Dans cette paramétrisation, calculer la première et la seconde forme fondamentale.
 \item Déterminer la courbure de Gauss et la courbure moyenne.
\end{enumerate}
\finenonce{007684}
\finexercice


\exercice{7685, mourougane, 2021/08/11}
\enonce{007685}{Etude jusqu'à l'ordre 2}
Soit l'ouvert $W =]0; 1[\times \Rr$ dans $\Rr^2$. On considère la surface $M$ de $\Rr^3$
paramétrée par 
$$\begin{array}{ccc}
 F: W&\to&\Rr ^3\\(u, v)&\mapsto& (u \cos v , u \sin v , v)
 \end{array}$$
\begin{enumerate} 
 \item Faire un dessin donnant l'allure de $M$.
 \item Calculer l'aire de la partie de $M$ comprise entre les plans d'équations $z = 0$ et $z = 2\pi$.
 \item Calculer, dans le paramétrage $F$, la courbure de Gauss $K$, la courbure
moyenne $H$, les courbures principales de la surface $M$.
\end{enumerate}
\finenonce{007685}
\finexercice


\exercice{7686, mourougane, 2021/08/11}
\enonce{007686}{Définition}
\begin{enumerate}
 \item Une surface régulière de $\Rr^3$ est dite réglée si elle admet des paramétrages
 locaux de la forme $F(t,s)=c(t)+sv(t)$ pour $t,s)\in I\times J$
où $c$ est une courbe régulière de $\Rr^3$ paramétrée sur l'intervalle $I$ de $\Rr$,
et $v~:I\to\Rr^3_{ev}$ une application de classe $\mathcal{C}^\infty$ avec $v(t)$ et $\dot{c}(t)$ 
partout indépendant. 
Montrer que si $J$ est un petit intervalle autour de $0$, $F$ est alors un paramétrage régulier.
 \item Montrer que la courbure de Gauss d'une surface réglée est négative en tout point.
\end{enumerate}
\finenonce{007686}
\finexercice


\exercice{7687, mourougane, 2021/08/11}
\enonce{007687}{Exemples}
\begin{enumerate}
 \item Montrer qu'un cylindre est une surface réglée.
 \item Montrer qu'un paraboloïde hyperbolique d'équation $z=xy$ est réglé
sur la droite d'équation $y=z=0$.
 \item Montrer qu'un hyperboloïde à une nappe d'équation $z^2=x^2+y^2-1$
est réglé sur le cercle à l'altitude $0$.
\end{enumerate}
\finenonce{007687}
\finexercice


\exercice{7688, mourougane, 2021/08/11}
\enonce{007688}{}
Calculer la torsion de la courbe de $\Rr^3$ paramétrée par $c~:~t\mapsto (4\cos t,5-5\sin t, -3\cos t)$.
La courbe  est-elle plane ?
\finenonce{007688}
\finexercice


\exercice{7689, mourougane, 2021/08/11}
\enonce{007689}{}
Soit $C$ la courbe tracée sur la surface d'équation $3z=xy+x^3$
et dont la projection orthogonale sur le plan d'équation $z=0$ 
est la courbe paramétrée $C'$ définie par $x=t$, $y=t^2$ pour $t$ parcourant $[0,1]$.
\begin{enumerate}
 \item Donner une expression intégrale pour la longueur de $C'$ 
puis celle de $C$, puis les comparer.

 \item Calculer la longueur de $C'$ puis celle de $C$.
\end{enumerate}

\finenonce{007689}
\finexercice


\exercice{7690, mourougane, 2021/08/11}
\enonce{007690}{}
Soit $S$ la surface de $\Rr^3$ d'équation $2(2z^2+y^2)+x=0$.
\begin{enumerate}
 \item La surface $S$ est-elle régulière ?
 \item Paramétrer la surface $S$ (de manière polynomiale) en prenant les paramètres $u$ et $v$
 parmi les variables $x$, $y$ et $z$.
 \item Déterminer une base de l'espace tangent à la surface $S$ en $A(-6,1,-1)$.
 \item Calculer un vecteur normal à la surface $S$ en $A(-6,1,-1)$.
 \item Le vecteur $V = (27,-29,-1)$ appartient-il au plan tangent à $S$ en $A (-6,1,1)$ ?
\end{enumerate}
\finenonce{007690}
\finexercice


\exercice{7691, mourougane, 2021/08/11}
\enonce{007691}{}
Calculer l'aire de l'ellipsoïde $\mathcal{E}$ d'équation $x^2+y^2+5z^2=1$.
\finenonce{007691}
\finexercice


\exercice{7692, mourougane, 2021/08/11}
\enonce{007692}{Paraboloïde hyperbolique}
Calculer une application de Weingarten du paraboloïde hyperbolique $S$ d'équation $z=x^2-y^2$
au point $p(0,0,0)$. En déduire sa courbure de Gauss et ses directions principales.
\finenonce{007692}
\finexercice


\exercice{7693, mourougane, 2021/08/11}
\enonce{007693}{}
Soit $C$ une courbe régulière plane fermée simple convexe paramétrée 
par la longueur d'arc par l'application
$[0,\ell]\to\Rr^2, t\mapsto c(t)$. Un coin est un point de la courbe 
où la fonction courbure a une dérivée nulle. Le but de l'exercice est de montrer 
que la courbe $C$ a au moins trois coins. On suppose que $C$ n'est pas un cercle.
 On notera $v(t)=\dot{c}(t)$ et $\gamma(t)=\ddot{c}(t)$ 
et $\kappa(t)$ la courbure au point de paramètre $t$.
\begin{enumerate}
 \item Montrer que $C$ a au moins deux coins distincts $P$ et $Q$. Faire une figure.
 \item Montrer à l'aide d'une intégration par partie que $\int_0^\ell \dot{\kappa}(t)c(t)=0$.
 \item On suppose que $P$ et $Q$ sont sur l'axe des $x$ et qu'il n'y a pas d'autres coins.
Aboutir à une contradiction.
 \item Montrer que la courbe $C$ a au moins quatre coins.
\end{enumerate}
\finenonce{007693}
\finexercice


\exercice{7694, mourougane, 2021/08/11}
\enonce{007694}{}
On considère dans $\Rr^3$ euclidien le cylindre $\mathcal{C}$ d'équation $x^2+y^2=1$
muni de la métrique riemannienne restriction du produit scalaire de $\Rr^3$.
\begin{enumerate}
 \item Montrer que $\mathcal{C}$ est une surface régulière.
 \item Montrer que l'application $F: ]-\pi,\pi[\times\Rr\to \Rr^3$, $(\theta,h)\mapsto (\cos\theta,\sin\theta,h)$
donne un paramétrage de $\mathcal{C}$ au voisinage du point $p$ de coordonnées $(1,0,0)$.
 \item Calculer la matrice $G(u)$ de la première forme fondamentale $I$
dans la base $$(X_\theta(\theta,h):=\frac{\partial F}{\partial \theta}(\theta,h),
X_h:=\frac{\partial F}{\partial h}(\theta,h))$$ de $T_{F(\theta,h)}\mathcal{C}$
correspondant à ce paramétrage $F$. 
 \item Déterminer un champs de vecteurs normaux unitaires $N(\theta,h)$ au point $F(\theta,h)$.
 \item Calculer les symboles de Christoffel de la base 
$(X_\theta(\theta,h),X_h(\theta,h))$ de $T_{F(\theta,h)}\mathcal{C}$.
 \item Soit $a$ un nombre réel fixé et $c :\Rr\to S$, $t\mapsto (\cos t,\sin t,at)$
la courbe paramétrée tracée sur le cylindre $\mathcal{C}$.
Exprimer le vecteur vitesse au point de paramètre $t$ dans la base 
$(X_\theta(u), X_h)$ de $T_{F(\theta,h)}\mathcal{C})$.
 \item Les courbes paramétrées précédentes sont-elles des géodésiques ?
 \item Les sections planes du cylindre paramétrées par la longueur d'arc
sont-elles des géodésiques ?
\end{enumerate}
\finenonce{007694}
\finexercice


\exercice{7695, mourougane, 2021/08/11}
\enonce{007695}{}
Montrer que l'image d'une géodésique par une isométrie entre deux surfaces munies de métriques riemanniennes est une géodésique. Retrouver les géodésiques du cylindre $\mathcal{C}$ d'équation $x^2+y^2=1$ de l'exercice précédent.
\finenonce{007695}
\finexercice


\exercice{7696, mourougane, 2021/08/11}
\enonce{007696}{Métrique riemannienne}
Soit $\kappa\in\Rr^+$. On considère sur le plan affine $\Rr^2$, la métrique riemannienne
donnée par 
$$g_{ij}(x_1,x_2)=\frac{1}{(1+\kappa(x^2+y^2))^2}\begin{pmatrix}1&0\\ 0&1\end{pmatrix}$$
dans la paramétrisation $Id$ de $\Rr^2$.
\begin{enumerate}
 \item Calculer les symboles de Christoffel, par la formule $$\Gamma_{ij}^k=\frac{1}{2}\sum_{m}
\left(\frac{\partial g_{im}}{\partial u^j}
+\frac{\partial g_{jm}}{\partial u^i}-\frac{\partial g_{ij}}{\partial u^m}\right)g^{mk}.$$
 \item On rappelle la formule des coefficients de l'endomorphisme de courbure
$$R(X_i,X_j)X_k=\sum_l \left(\frac{\partial\Gamma_{kj}^l}{\partial u^i}
-\frac{\partial\Gamma_{ki}^l}{\partial u^j}
+\sum_m\left(\Gamma_{mi}^l\Gamma_{kj}^m-\Gamma_{mj}^l\Gamma_{ki}^m\right)
\right)X_l$$
Calculer la courbure de Gauss de $(\Rr^2,g)$.
\end{enumerate}
\finenonce{007696}
\finexercice


\exercice{7697, mourougane, 2021/08/11}
\enonce{007697}{}
\begin{enumerate}
 \item La courbe paramétrée suivante est-elle régulière ?
$$\begin{array}{ccc}
 c: \Rr &\to&\Rr ^3\\ t&\mapsto& \left(\begin{array}{c}t\cos t \\ t\sin t\\ t\end{array}\right).
\end{array}$$
 \item On fixe deux nombres réels $a<b$. Comparer la longueur de la portion de la courbe $c$ 
paramétrée par $[a,b]$ et $\sqrt{2}(b-a)$.
 \item Déterminer une surface quadrique qui contient l'image $\mathcal{C}$ de la courbe $c$.
 \item Déterminer un déplacement (non égal à l'identité) de l'espace $\Rr ^3$ qui conserve l'image $\mathcal{C}$ de la courbe $c$.
 \item Décrire l'allure de la courbe $\mathcal{C}$.
 \item Calculer la fonction torsion de la courbe $c$.
\end{enumerate}
\finenonce{007697}
\finexercice


\exercice{7698, mourougane, 2021/08/11}
\enonce{007698}{}
La courbure moyenne de la surface paramétrée suivante est-elle partout nulle ? 
$$\begin{array}{ccc}
 F: \Rr \times ]-\pi,\pi[&\to&\Rr ^3\\ 
\left(\begin{array}{c}u \\ v\end{array}\right)
&\mapsto& 
\left(\begin{array}{c}\cosh(u)\cos(v)\\ \cosh(u)\sin(v)\\ u\end{array}\right).
\end{array}$$
\finenonce{007698}
\finexercice


\exercice{7699, mourougane, 2021/08/11}
\enonce{007699}{}
On considère l'application
$$\begin{array}{ccc}
 F: \Rr ^2&\to&\Rr ^3\\ 
\left(\begin{array}{c}t \\ s\end{array}\right)
&\mapsto& 
\left(\begin{array}{c}1-t^2\\t(1-t^2)\\s\end{array}\right).
\end{array}$$
\begin{enumerate}
 \item Déterminer des déplacements de l'espace $\Rr ^3$ qui conservent l'image de $F$.
 \item Montrer que l'image de $F$ est l'ensemble d'équation $y^2=x^2(1-x)$.
 \item L'image de $F$ est-elle une surface régulière de $\Rr ^3$ ?
\end{enumerate}
\finenonce{007699}
\finexercice


\exercice{7700, mourougane, 2021/08/11}
\enonce{007700}{}
Trouver les extrema de la fonction $f: (x,y,z)\mapsto xy$ sur la sphère unité.
\finenonce{007700}
\finexercice


\exercice{7701, mourougane, 2021/08/11}
\enonce{007701}{}
On considère la courbe paramétrée
$$\begin{array}{cccc}
        c~:~&]0,1[&\to&\Rr ^3\\
&t&\mapsto&\left(\begin{array}{c}3\cos t\\3\sin t\\4t\end{array}\right)
        \end{array}
$$
\begin{enumerate}
 \item La courbe $c$ est-elle régulière ?
 \item Paramétrer l'image de $c$ par la longueur d'arc.
 \item Déterminer en tout point de $c$ le repère de Frenet.
\end{enumerate}
\finenonce{007701}
\finexercice


\exercice{7702, mourougane, 2021/08/11}
\enonce{007702}{}
 On considère l'application 
$$\begin{array}{ccc}
 F: ]-\pi,\pi[\times ]-\pi,\pi[&\to&\Rr^3\\ 
\left(\begin{array}{c}\varphi\\ \theta\end{array}\right)
&\mapsto& 
\begin{pmatrix}(2+\cos(\varphi))\cos(\theta)\\ (2+\cos(\varphi))\sin(\theta)\\ \sin (\varphi)\end{pmatrix}.
\end{array}$$
\begin{enumerate}
 \item  Représenter l'image $T$ de $F$. 
 \item Montrer que $T$ est une surface régulière de $\Rr ^3$. 
 \item Déterminer la courbure de Gauss $K$ de $T$ avec la métrique induite par le produit scalaire sur $\Rr ^3$,
c'est à dire la première forme fondamentale.(Expliquer d'abord votre démarche globale. Tout résultat
intermédiaire sera pris en compte). 
 \item Calculer $\int_T K(m)d\sigma(m)$.
\end{enumerate}
\finenonce{007702}
\finexercice


\exercice{7703, mourougane, 2021/08/11}
\enonce{007703}{}
 On considère l'application 
$$\begin{array}{ccc}
 F:]-\pi,\pi[\times ]-\pi,\pi[ &\to&\Rr ^3\\ 
\left(\begin{array}{c}\varphi\\ \theta\end{array}\right)
&\mapsto& 
\begin{pmatrix}(2+\cos(\varphi))\cos(\theta)\\ (2+\cos(\varphi))\sin(\theta)\\ r\sin (\varphi)\end{pmatrix}.
\end{array}$$
\begin{enumerate}
 \item Montrer qu'on définit une métrique riemannienne sur $Im(F)=T$
en posant $g(\frac{\partial F}{\partial \varphi})=g(\frac{\partial F}{\partial \theta})=1$
et $g(\frac{\partial F}{\partial \varphi},\frac{\partial F}{\partial \theta})=0$.
 \item Déterminer la courbure de Gauss $K$ de $T$ avec la métrique $g$.
 \item Calculer $\int_T K(m)d\sigma(m)$.
\end{enumerate}
\finenonce{007703}
\finexercice


\exercice{7704, mourougane, 2021/08/11}
\enonce{007704}{}
Soit $S$ une surface différentiable munie d'une métrique riemannienne $g$.
Soit $p$ et $q$ deux points fixés sur $S$.
Soit $\epsilon$ un nombre réel strictement positif. Soit $a\leq b$ deux nombres réels.
Soit $C :]-\epsilon,\epsilon [\times [a,b]\to S$ une application différentiable vérifiant 
pour tout $s\in ]-\epsilon,\epsilon [$,
$C(s,a)=p$ et $C(s,b)=q$.
On notera $c_s(t):=C(s,t)$, $V(t):=\frac{\partial C}{\partial s}_{(0, t)}$.
On admettra l'identité des dérivées covariantes
$$\nabla_{V(t)}\dot{c_0}(t)=\nabla_{\dot{c_0}(t)}V(t).$$
\begin{enumerate}
 \item Représenter sur un dessin $S$, l'application $C$ et le champs de vecteurs $V$,
en particulier $V(a)$ et $V(b)$.
 \item Rappeler la définition de l'énergie $E[c_s]$ de la courbe $c_s : [a,b]\to S$.
 \item Exprimer à l'aide de la dérivée covariante $\nabla_{\dot{c_0}(t)}\dot{c_0}(t)$, 
la dérivée $\frac{dE[c_s]}{ds}_{s=0}$.
\end{enumerate}
\finenonce{007704}
\finexercice


\exercice{7705, mourougane, 2021/08/11}
\enonce{007705}{}
Soit $\Sigma$ le graphe de la fonction $f(x,y) = xy$ de $\Rr^2$ dans $\Rr$. 
\begin{enumerate}
 \item Est-ce une surface régulière ?
 \item Déterminer une paramétrisation de $\Sigma$.
 \item Dans cette paramétrisation, calculer la première et la seconde forme fondamentale.
 \item Déterminer la courbure de Gauss et la courbure moyenne.
\end{enumerate}
\finenonce{007705}
\finexercice


\exercice{7706, mourougane, 2021/08/11}
\enonce{007706}{}
Soit l'ouvert $W =]0; 1[\times \Rr$ dans $\Rr^2$. On considère la surface $M$ de $\Rr^3$
paramétrée par 
$$\begin{array}{ccc}
 F: W&\to&\Rr ^3\\(u, v)&\mapsto& (u \cos v , u \sin v , v)
 \end{array}$$
\begin{enumerate} 
 \item Faire un dessin donnant l'allure de $M$.
 \item Calculer l'aire de la partie de $M$ comprise entre les plans d'équations $z = 0$ et $z = 2\pi$.
 \item Calculer, dans le paramétrage $F$, la courbure de Gauss $K$, la courbure
moyenne $H$, les courbures principales de la surface $M$.
\end{enumerate}
\finenonce{007706}
\finexercice


\exercice{7707, mourougane, 2021/08/11}
\enonce{007707}{}
Trouver les extrema de la fonction 
$\phi(x,y,z)=xy$ définie sur l'ellipsoïde de $\Rr^3$ d'équation $x^2+y^2+3z^2=1$.
\finenonce{007707}
\finexercice


\exercice{7708, mourougane, 2021/08/11}
\enonce{007708}{}
On considère dans $\Rr^4$ l'intersection des ensembles d'équation 
$x^2+y^2+z^2+t^2=8$ et $x^2+y^2=4$.
\begin{enumerate}
 \item Est-ce une surface différentiable ?
 \item Déterminer un système d'équation pour ses espaces tangents.
\end{enumerate}
\finenonce{007708}
\finexercice


\exercice{7709, mourougane, 2021/08/11}
\enonce{007709}{}
Démontrer les formules de Frenet pour les courbes paramétrées par la longueur d'arc 
dans $\Rr^3$.
\finenonce{007709}
\finexercice


\exercice{7710, mourougane, 2021/08/11}
\enonce{007710}{}
Soit $I$ un intervalle de $\Rr$ et $c~:I\to \Rr^2$ une courbe plane paramétrée par la longueur d'arc. 
On suppose que $c$ reste dans le disque de rayon $r>0$ et qu'au point de paramètre $\tau$,
$\| c(\tau)\| = r$.
\begin{enumerate} 
 \item Rappeler la valeur absolue de la courbure d'un cercle de rayon $r$.

 \item Montrer en dérivant une fois la fonction $\phi : t\mapsto \| c(t)\|^2$ 
que $\ddot{c}(\tau)$ est colinéaire à $c(\tau)$.

 \item Montrer en dérivant à nouveau la fonction $\phi$ 
que la courbure en $\tau$ vérifie $|\kappa(\tau)|\geq 1/r$.

 \item Interpréter graphiquement ce résultat.
\end{enumerate}
\finenonce{007710}
\finexercice


\exercice{7711, mourougane, 2021/08/11}
\enonce{007711}{}
\begin{enumerate}
 \item Déterminer une équation cartésienne du plan tangent 
à la surface $\mathcal{S}$ de $\Rr^3$ paramétrée par $F : (u,v)\mapsto (u+v^2,u^2-v^2,v)$
au point $M(u_0,v_0)$ de paramètres $(u_0,v_0)$.

 \item Déterminer une équation cartésienne du plan tangent 
à la surface $\Sigma$ de $\Rr^3$ d'équation $x^5+y^5+z^5=1$
au point $M(x_0,y_0,z_0)$ de coordonnées $(x_0,y_0,z_0)$.

 \item Soit $U$ un ouvert de $\Rr^2$ et $f~:U\to \Rr$ une application de classe $\mathcal{C}^\infty$.
Déterminer une équation cartésienne du plan tangent au graphe de $f$ en chacun de ses points.
\end{enumerate}
\finenonce{007711}
\finexercice


\exercice{7712, mourougane, 2021/08/11}
\enonce{007712}{}
\begin{enumerate}
 \item Déterminer les plans tangents et un champs de vecteurs normaux unitaires
au paraboloïde hyperbolique $\mathcal{P}$ d'équation $z=y^2-x^2$ au voisinage du point $A(0,0,0)$
de coordonnées $(0,0,0)$.
 \item Déterminer l'application de Weingarten au point $A$ du paraboloïde hyperbolique $\mathcal{P}$.
 \item En déduire la courbure de Gauss et les directions principales du paraboloïde hyperbolique $\mathcal{P}$ au point $A$.
\end{enumerate}
\finenonce{007712}
\finexercice


\exercice{7713, mourougane, 2021/08/11}
\enonce{007713}{}
On considère la sphère de centre $0$ et de rayon $1$ de $\Rr^3$.
La projection stéréographique de centre le pôle Nord sur le plan de hauteur nulle (d'équation $z=0$)
conserve-t-elle les longueurs ? les angles ? les aires ?
\finenonce{007713}
\finexercice


\exercice{7714, mourougane, 2021/08/11}
\enonce{007714}{}
L'intersection de deux surfaces régulières de $\Rr^3$ est-elle une courbe régulière de $\Rr^3$ ?
\finenonce{007714}
\finexercice

\section{353.00 Espace tangent, application linéaire tangente}


\exercice{6793, gijs, 2011/10/16}
\enonce{006793}{}
Soit $f:\Rr^n \to \Rr^p$ une fonction dérivable
et soit $M = f^{-1}(\mathbf{0})$. On suppose que $M$ est une
sous-variété de $\Rr^n$ de dimension $k$.
\begin{enumerate}

\item
Donner la définition d'un  vecteur
tangent à $M$ au point $x \in M$. Donner la définition
d'un champ de vecteurs sur $M$.

\item Soit $f$ une fonction sur $M$ et $\alpha$ une
1-forme sur $M$. Donner la définition de $df$ sans
utiliser une carte de $M$. Donner la définition de
$d\alpha$.

\item 
Soit $X$ un vecteur tangent à $M$ au point $x\in M$.
Montrer que $\bigl(\ df|_x(X) \equiv\ \bigr)\ Xf =
\mathbf{0}$. En déduire que $X \in \operatorname{ker}(J(x))$
(où $J(x)$ est la matrice jacobienne de $f$ en $x$). 

\item
Soit $X$ un vecteur tangent $\underline{ \text{à\ }
\Rr^n}$ au point $x\in M$, et supposons que $Xf  =
\mathbf{0}$, et que $\operatorname{rang}(J(x)) = p$.
Montrer que $k=n-p$ et que  $X$ est un vecteur tangent
$\underline{ \text{à }M}$. 
\end{enumerate}
\finenonce{006793}



\finexercice


\exercice{7664, mourougane, 2021/08/11}
\enonce{007664}{Applications régulières ?}
\begin{enumerate}
 \item La restriction à la surface d'équation $x^2+y^2+z^2-6x=-5$ de la fonction
$(x,y,z)\mapsto 5x^6+7xy^4-2$ est-elle différentiable ?
 \item La restriction à la surface d'équation $x^2+y^2+z^2-6x=-5$ de la fonction
$(x,y,z)\mapsto \sqrt{|x|}$ est-elle différentiable ?
\end{enumerate}
\finenonce{007664}
\finexercice


\exercice{7665, mourougane, 2021/08/11}
\enonce{007665}{Difféomorphismes}
\begin{enumerate}
 \item Soit $U$ un ouvert de $\Rr^2$ et $f~:U\to \Rr$ une application de classe $\mathcal{C}^\infty$.
Montrer que le graphe de $f$ est difféomorphe à $U$.
 \item Les surfaces $\mathcal E$ d'équation $2x^2+3y^2+4z^2=1$ et $\mathcal S$ d'équation
$x^2+y^2+z^2=4$ sont-elles difféomorphes ?
 \item Les surfaces $\mathcal C$ d'équation $2x^2+3y^2=1$ et $\mathcal S$ d'équation
$x^2+y^2+z^2=4$ sont-elles difféomorphes ?
\end{enumerate}
\finenonce{007665}
\finexercice


\exercice{7666, mourougane, 2021/08/11}
\enonce{007666}{Détermination de plan tangent}
\begin{enumerate}
 \item Déterminer le plan tangent en $A(1,-1,-1)$
à la surface paramétrée par $(u,v)\mapsto (u+v^2,u^2-v^2,v)$.
 \item Déterminer le plan tangent en $A(1,1,-1)$
à la surface d'équation $x^5+y^5+z^5=1$.
 \item Soit $U$ un ouvert de $\Rr^2$ et $f~:U\to \Rr$ une application de classe $\mathcal{C}^\infty$.
Déterminer le plan tangent au graphe de $f$ en chacun de ses points.
\end{enumerate}
\finenonce{007666}
\finexercice


\exercice{7667, mourougane, 2021/08/11}
\enonce{007667}{Plan tangent}
Soit $S$ la surface de $\Rr^3$ d'équation $2(2z^2+y^2)+x=0$. 
\begin{enumerate}
 \item La surface $S$ est-elle régulière ?
  \item Paramétrer la surface $S$ (de manière polynomiale) en prenant les paramètres $u$ et $v$
 parmi les variables $x$, $y$ et $z$.
 \item Déterminer une base de l'espace tangent à la surface $S$ en $A(-6,1,-1)$.
 \item Calculer un vecteur normal à la surface $S$ en $A(-6,1,-1)$.
 \item Le vecteur $V = (27,-29,-1)$ appartient-il au plan tangent à $S$ en $A (-6,1,-1)$ ? 
\end{enumerate}
\finenonce{007667}
\finexercice


\exercice{7668, mourougane, 2021/08/11}
\enonce{007668}{Plan tangent}
\begin{enumerate}
 \item Déterminer une équation cartésienne du plan tangent 
à la surface $\mathcal{S}$ de $\Rr^3$ paramétrée par $F : (u,v)\mapsto (u+v^2,u^2-v^2,v)$
au point $M(u_0,v_0)$ de paramètres $(u_0,v_0)$.
 \item Déterminer une équation cartésienne du plan tangent 
à la surface $\Sigma$ de $\Rr^3$ d'équation $x^5+y^5+z^5=1$
au point $M(x_0,y_0,z_0)$ de coordonnées $(x_0,y_0,z_0)$.
 \item Soit $U$ un ouvert de $\Rr^2$ et $f~:U\to \Rr$ une application de classe $\mathcal{C}^\infty$.
Déterminer une équation cartésienne du plan tangent au graphe de $f$ en chacun de ses points.
\end{enumerate}
\finenonce{007668}
\finexercice


\exercice{7669, mourougane, 2021/08/11}
\enonce{007669}{Détermination de minima}
Dans tout l'exercice, $r$ et $h$ parcourent $]0,+\infty[$.
\begin{enumerate}
 \item Déterminer le volume $V(r,h)$ d'un cylindre plein de hauteur $h$ et de rayon $r$. 
 \item Déterminer l'aire $A(r,h)$ d'une casserole de hauteur $h$ et de rayon $r$.
 \item Montrer que le sous-ensemble de $\Rr^3$ avec coordonnées $(r,h,v)$ 
d'équation $1=V(r,h)$ est une surface régulière.
 \item Déterminer le gradient de la fonction $\alpha : (r,h,v)\mapsto A(r,h)$ 
et celui de la fonction $\nu :(r,h,v)\mapsto V(r,h)-1$
 \item Soit $(r_0,h_0,v_0)$ un minimum de la fonction $\alpha$ sur la surface d'équation $1=V(r,h)$.
Comparer $grad _{(r_0,h_0,v_0)}\alpha$ et $grad _{(r_0,h_0,v_0)}\nu$.
 \item Déterminer le minimum de $A(r,h)$ sur la surface d'équation
$V(r,h)=1$. Interpréter ce résultat.
\end{enumerate}
\finenonce{007669}
\finexercice

\section{354.00 Champ de vecteurs}


\exercice{6776, gijs, 2011/10/16}
\enonce{006776}{Questions de cours}
\begin{enumerate}
\item Donner les définitions de : 
  \begin{enumerate}
  \item une norme sur un espace vectoriel,
  \item un espace vectoriel normé complet, 
  \item une sous-variété $M$ de $\Rr^n$ de dimension $k$, 
  \item un champ de vecteurs sur $M$, 
  \item un champ de vecteurs  complet sur $M$.
  \end{enumerate}
\item Enoncer le théorème des fonctions implicites.
\item Soit $A\subset \Rr^n$ compact, soit $\mathcal{C}^0(A) 
= \{\,f:A \to \Rr \mid f \text{ continue }\,\}$,
et $  ||f||_\infty = \sup_{x\in A} |f(x)|$.
Montrer que $(\mathcal{C}^0(A), ||{\ }||_\infty)$ est un espace
de Banach.
\end{enumerate}
\finenonce{006776}



\finexercice


\exercice{6777, gijs, 2011/10/16}
\enonce{006777}{}
Soit $M\subset \Rr^n$ une sous-variété
de dimension $k$, et soit $X$ un champ de vecteurs sur
$M$. On définit $D = \{\,m\in M\mid X(m) \neq 0\,\}$ et
$S = \overline D = \text{ fermeture } D$. On vous demande
de démontrer l'énoncé : ``si $S$ est compact, alors
$X$ est complet''. Les questions suivantes peuvent vous
guider.
\begin{enumerate}
\item Pour $x \notin S$ trouver la courbe intégrale
maximale $\gamma : J_x \to M$ passant par $x$.
\item Montrer que si $x\in S$, et si $\gamma : J \to M$
est une courbe intégrale passant par $x$, alors
$\forall\ t\in J\ :\ \gamma(t) \in S$.
\item En utilisant la compacité de $S$, montrer que
$X$ est complet (sur $M$ !).
\end{enumerate}
\finenonce{006777}



\finexercice


\exercice{6778, gijs, 2011/10/16}
\enonce{006778}{}
 Soient $M\subset \Rr^n$ et $N\subset \Rr^p$ deux sous-variétés de dimension $k$ et $\ell$
respectivement. Soit $F : M \to N$ une application
différentiable et soit $X$ un champ de vecteurs sur
$M$. Trouver un contre exemple pour l'énoncé : $$F(m)
= F(\widehat m) \Longrightarrow  TF(m)(X(m)) =
TF(\widehat m)(X(\widehat m))\ .$$ 
Rappel : $TF(m) \equiv
F'(m)$ est la ``dérivée'' de $F$ au point $m$.
\finenonce{006778}



\finexercice


\exercice{6779, gijs, 2011/10/16}
\enonce{006779}{}
 Soit $\phi_2 : \Rr^2 \to S^2 \subset \Rr^3$ la carte de la sphère $S^2$
donnée par la projection stéréographique du p\^ole
nord. En identifiant $\Rr^2$ avec le plan complex
$\Cc$, on définit l'application $F : \Rr^4 \to
S^2$ par : $$F(x,y,z,t) = \phi_2(\frac{z+it}{x+iy})\ .
$$
\begin{enumerate}
\item Calculer l'expression explicite de $F$ et
montrer que la restriction de $F$ à la sphère $S^3
\subset \Rr^4$ est une application $F : S^3 \to S^2$
qui est bien définie.
\item Sur $\Rr^4$ on définit le champ de
vecteurs $$X(x,y,z,t) = x \frac\partial{\partial y} - 
y \frac\partial{\partial x} + z \frac\partial{\partial t}
- t \frac\partial{\partial z}\ .
$$
Calculer le flot de $X$; est ce que $X$ est complet ?
\item En utilisant le résultat de 2., montrer que
si $m\in S^3$, alors $X(m) \in T_m S^3$.
\item Pour tout $m\in S^3$ calculer $TF(m)(X(m)) \in
T_{F(m)} S^2$.
\end{enumerate}
\finenonce{006779}



\finexercice


\exercice{6794, gijs, 2011/10/16}
\enonce{006794}{}
Soit $f:\Rr^4 \to \Rr^2$ défini par
$f(x_1,x_2,x_3,x_4) = (x_1^2 + x_2^2 -1, x_3^2 + x_4^2
-1)$,  soit $M = f^{-1}(0,0)$, et soit $X$ le champ de
vecteurs sur $\Rr^4$ défini par~:
$$
X|_x = x_1 \frac\partial{\partial x_2} - x_2
\frac\partial{\partial x_1} + \alpha ( x_3
\frac\partial{\partial x_4} - x_4 \frac\partial{\partial
x_3})\ ,
$$
où $\alpha\in \Rr$ est un paramètre.
\begin{enumerate}
\item Montrer que pour tout $(x_1,x_2,x_3,x_4) \equiv x\in M$ le
vecteur $X|_x$ est tangent à $M$.

\item Exprimer les vecteurs tangents $X|_x$, $x\in M$
dans la carte $$(\theta,\psi) \mapsto (\cos(\theta),
\sin(\theta), \cos(\psi), \sin(\psi))\ .$$

\item Calculer le flot $\phi_t$ du champ de vecteurs
$X|_x$, $x\in M$ sur $M$ (par exemple en utilisant la
carte $(\theta,\psi)$). Est ce que ce champ est complet\,?

\item Déterminer les 6-uplets $(\alpha,t,x_1,x_2,x_3,x_4)$,
$(x_1,x_2,x_3,x_4) \in M$, tels que
$\phi_t(x_1,x_2,x_3,x_4) = (x_1,x_2,x_3,x_4)$.
\end{enumerate}
\finenonce{006794}



\finexercice


\exercice{6800, gijs, 2011/10/16}
\enonce{006800}{}
Soit $M \subset \Rr^3$ le cylindre défini par
l'équation $x^2 + y^2 = 1$. Dans la carte
$(\theta,z) \mapsto (\cos(\theta), \sin(\theta), z)$ de
$M$ on  donne le champ de vecteurs $X$ défini par~:
$$
X(\theta,z) =
\frac{\partial}{\partial\theta}{\Big\vert_{(\theta,z)}} +
\tfrac12 z(z^2-1) 
\frac{\partial}{\partial z}{\Big\vert_{(\theta,z)}}\ . $$
\begin{enumerate}
\item Dire pourquoi $X$ définit un champ de vecteurs sur
$M$.

\item Esquisser le champ $X$ dans la carte et calculer son
flot $\phi_t$  sur $M$.

\item Le champ $X$ est-il complet~?

\item Calculer les points $(t,\theta,z)$ tel que
$\phi(t,\theta,z) = (\theta,z)$ {\bf sur le cylindre $M$}.
\end{enumerate}
\finenonce{006800}



\finexercice

\exercice{6802, gijs, 2011/10/16}
\enonce{006802}{}
\begin{enumerate}
\item Donner la définition d'une sous-variété $M$ de
$\Rr^n$ de dimension $k$.

\item Enoncer le théorème des fonctions implicites.

\item  Soit $M = \{\,(x,y,z)\in \Rr^3  \mid x^2 + y^2 -
z^2 +1 = 0\ ,\ z>0\,\}$. Montrer que $M$ est une
sous-variété de $\Rr^3$ de dimension 2.

\medskip
On définit la projection stéréographique $s$ de
$\Rr^3 \setminus\{\,(x,y,z)\mid z=-1\,\}$ sur
$\Rr^2 \cong \{\,(x,y,z) \mid z=0\,\}$ par la
procédure suivante. Pour un point $P=(x,y,z)$ on trace
la droite $d=\overline{PS}$ où $S=(0,0,-1)$. L'image
$s(P)$ est l'intersection de la droite $d$ avec le plan
$z=0$.

\item Calculer explicitement l'application $s$.

\item Calculer l'image $D= s(M)$.

\item Calculer ``l'inverse'' de l'application $s:M \to D$.

\medskip
Soit $X$ le champ de vecteurs sur $\Rr^3$ donné
par $$X_{|(x,y,z)} = x
\frac{\partial}{\partial z}_{|(x,y,z)} + z
\frac{\partial}{\partial x}_{|(x,y,z)}
\ \cong\  (z,0,x)\ .
$$

\item Calculer le flot du champ $X$.

\item  Montrer que $X$ est tangent à $M$,
c'est-à-dire que pour tout $(x,y,z)\in M$ le vecteur
$X_{|(x,y,z)}$ appartient à l'espace
tangent $T_{(x,y,z)}M$.

\item Calculer l'expression de $X$ dans la carte $D$ de
$M$.

\item Calculer le flot du champ sur $D$ obtenu en i).
\end{enumerate}
\finenonce{006802}



\finexercice


\exercice{6803, gijs, 2011/10/16}
\enonce{006803}{}
Pour la sphère $S^2$ on considère la carte $U=
\ ]0,\infty[\ \times\ ]0,2\pi[\ $
donnée par~:
$$
\varphi : \ U \to S^2 \quad,\quad
(r,\theta) \mapsto (\frac{2r}{r^2+1} \cos(\theta),
\frac{2r}{r^2+1} \sin(\theta), 
\frac{r^2-1}{r^2+1} )\ .
$$
Dans la carte $U$ on donne le champ de vecteurs $X$
par~:
$$
X_{|(r,\theta)} = f(r)
\frac{\partial}{\partial r}_{|(r,\theta)}
 \ \cong\ (f(r),0)\ .
$$

\begin{enumerate}
\item Calculer le champ sur $S^2$, c'est-à-dire les
vecteurs $\varphi'(r,\theta) X_{|(r,\theta)}$.

\item Soit $f(r) = r^2$. Existe-t-il un
champ de vecteurs continue $Y$ sur la sphère $S^2$
{\bf entière} telle que
$Y_{|\varphi(r,\theta)} =
\varphi'(r,\theta) X_{|(r,\theta)}$ pour tout
$(r,\theta)$ ?

\item Même question qu'en 2. dans le cas $f(r) =
\dfrac{r^2-1}{r^2+1}$.

\item Quelle condition nécessaire et suffisante
(la plus simple possible) doit vérifier la fonction
continue
$f: \ ]0,\infty[\ \to \Rr$ pour qu'il existe un
champ de vecteurs continue $Y$ sur la sphère $S^2$
{\bf entière} telle que
$Y_{|\varphi(r,\theta)} =
\varphi'(r,\theta) X_{|(r,\theta)}$ pour tout
$(r,\theta)$ ?
\end{enumerate}
\finenonce{006803}



\finexercice


\exercice{6823, gijs, 2011/10/16}
\enonce{006823}{}

Soit $Y:\Rr^2 \to \Rr^2$ le champ de vecteurs $Y(x,y) =
(1,1+y^2)$.
\begin{enumerate}
\item Trouver les solutions maximales du champ $Y$,
y compris leur domaine de définition.

\item Esquisser le portrait de phase du champ $Y$.

\item Déterminer le flot du champ $Y$, y compris
son domaine de définition.
\end{enumerate}
\finenonce{006823}



\finexercice


\exercice{6829, gijs, 2011/10/16}
\enonce{006829}{}
Soit $Y:\Rr^2 \to \Rr^2$ le champ de vecteurs $Y(x,y) =
(x^2,xy)$.
\begin{enumerate}
\item Trouver les solutions maximales du champ $Y$,
y compris leur domaine de définition. 

\item Esquisser le portrait de phase du champ $Y$.

\item Déterminer le flot du champ $Y$, y compris
son domaine de définition.
\end{enumerate}
\finenonce{006829}



\finexercice


\exercice{6840, gijs, 2011/10/16}
\enonce{006840}{}
Soit $Y:\Rr^2 \to \Rr^2$ le champ de vecteurs $Y(x,y) =
(x^2y,xy^2)$.
\begin{enumerate}
\item Si $\gamma(t) = (x(t), y(t))$ est solution de
l'équation $\gamma'(t) = Y(\gamma(t))$, que peut-on dire
de la dérivée de
$y/x$ par rapport
à
$t$~?

\item Trouver les solutions maximales du champ $Y$
(y compris leur domaine de définition). 

\item Esquisser le portrait de phase du champ $Y$.

\item Déterminer le flot du champ $Y$ (y compris
son domaine de définition).
\end{enumerate}
\finenonce{006840}



\finexercice


\exercice{6846, gijs, 2011/10/16}
\enonce{006846}{}
Soit $U = \{(x,y) \in \Rr^2 \mid y\neq0\,\}$ et
soit $Y:U \to \Rr^2$ le champ de vecteurs $Y(x,y) =
(1, \dfrac{2x}{3y^2})$.
\begin{enumerate}
\item Trouver les solutions maximales du champ $Y$,
y compris leur domaine de définition (n'oubliez pas
qu'on est dans $U$). 

\item Esquisser le portrait de phase du champ $Y$.

\item Déterminer le flot du champ $Y$, y compris
son domaine de définition.
\end{enumerate}
\finenonce{006846}



\finexercice

\section{355.00 Forme différentielle}


\exercice{6780, gijs, 2011/10/16}
\enonce{006780}{}
 Soit $\alpha(x,y,z,t) = x\,dy - (1+t^2)y\,dx -
z^3\,dt$ une 1-forme sur $\Rr^4$, soit $D =
\{\,(x,y,z,t) \in S^3 \mid t=0, z\geq 0\,\}$, et soit $C
= \partial D$ le bord de $D$.
\begin{enumerate}
\item Calculer la 2-forme $d\alpha$ sur $\Rr^4$.
\item Montrer que $D$ est une demi-sphère de
dimension 2 et que $C$ est un cercle.
\item En utilisant des coordonnées sphériques,
donner une carte $\phi_D : I_1 \times I_2 \to D \subset
S^3 \subset \Rr^4$, et une carte $\phi_C : I_3 \to C
\subset D \subset \Rr^4$, où les $I_j \subset \Rr$ sont des intervalles ouverts.
\item Calculer explicitement $\int_D d\alpha$ et
$\int_{\partial D} \alpha$. Est ce que votre résultat
confirme le théorème de Stokes ?
\end{enumerate}
\finenonce{006780}



\finexercice

\exercice{6785, gijs, 2011/10/16}
\enonce{006785}{} 
Dans tout ce qui suit, on note
$(x,y)$ les coordonnées sur $\Rr^2 = \Cc$, $z=x+
iy$, et on sépare systématiquement les parties
réelle et imaginaire de tous les objets. 
Soit  $z_o \in \Cc$, et soit  $C_r$, $r\in \Rr^+$  la
courbe  donnée par l'équation $|z-z_o|=r$.


\begin{enumerate}
\item
Soit $f:\Cc\setminus\{z_o\} \to \Cc$ l'application
$f(z) = (z-z_o)^n$. Calculer pour tout $n\in \Zz$
l'intégrale $\int_{C_r} f(z)\,dz$ (indication~:
utiliser une variante des coordonnées polaires).


\medskip
Soit $g,h : \Cc\setminus\{z_o\} \to \Rr$
deux applications de classe $C^1$ verifiant les
équations~:
$$
\frac{\partial g}{\partial x} = 
\frac{\partial h}{\partial y}\qquad\&\qquad
\frac{\partial g}{\partial y} = 
-\frac{\partial h}{\partial x}\ ,
$$
et soit $f:\Cc\setminus\{z_o\} \to \Cc$
l'application $f(x,y) = g(x,y) + i h(x,y)$.

\item
 Montrer que la 1-forme
$ \frac{f(z)}{z-z_o}\,dz$ sur $\Rr^2\setminus\{z_o\}$ est fermée (ne pas oublier de
séparer la partie réelle et imaginaire). En déduire
(Stokes!) que $ \int_{C_r}  \frac{f(z)}{z-z_o}\,dz$
est indépendant de $r\in \Rr^+$.

\item En faisant un développement limité de $g$
et de $h$ d'ordre 1 autour $z_o = (x_o,y_o)$, montrer que 
$$
\int_{C_r} \frac{f(z)}{z-z_o}\,dz = 2\pi i\,\, f(z_o)\ .
$$
Indication : utiliser le résultat de 1.; vous avez le
droit d'être un petit  peu vague en ce qui concerne les
$\varepsilon$ dans le développement limité.
\end{enumerate}
\finenonce{006785}



\finexercice


\exercice{6795, gijs, 2011/10/16}
\enonce{006795}{}
Soit $M \subset \Rr^3$ le cylindre défini par
l'équation $x^2 + y^2 = 1$, soit $V\subset \Rr^3$
l'ensemble défini par $V = \{\,(x,y,z)\mid z\ge 0\ \&
\ (x/2)^2 + (2y)^2 + z^2 \le 1\,\}$, soit $K = M \cap V$,
et soit $\partial K$ le bord de $K$. Soit finalement
$\alpha$ la 1-forme sur $\Rr^3$ définie par $\alpha
= z^2x\,dy - z^2y\,dx + (xz-yz^3)\,dz$.

\begin{enumerate}
\item Calculer $d\alpha$.

\item Exprimer $\alpha$ et $d\alpha$ dans la carte
$(\theta,z) \mapsto (\cos(\theta), \sin(\theta), z)$ de
$M$.

\item Exprimer $K$ et $\partial K$ dans cette carte.

\item Calculer séparément $\int_K d\alpha$ et
$\int_{\partial K} \alpha$ sans utiliser le théorème
de Stokes. Est ce que votre résultat confirme ce
théorème\,?
\end{enumerate}
\finenonce{006795}



\finexercice


\exercice{6801, gijs, 2011/10/16}
\enonce{006801}{}
Soit $\alpha = f(x,y)\,dx + g(x,y)\,dy$ une 1-forme
fermée ($d\alpha=0$) sur $\Rr^2$, et soient $(x_0,
y_0)$, $(x_1, y_1)$, $(x, y)$ trois points. On définit
les courbes $\gamma_0, \gamma_1 : [0,1] \to \Rr^2$
par~: $$
\gamma_i(t) = (1-t) (x_i,y_i) + t(x,y)\ ,
$$
et les fonctions $h_0, h_1 : \Rr^2 \to
\Rr$ par~:
$  h_i(x,y) = \int_{\gamma_i} \alpha$.


\begin{enumerate}
\item Montrer que $dh_i = \alpha$ (indication~:
calculer $\tfrac{d}{dt} f(\gamma_i(t))$ et $\tfrac{d}{dt}
g(\gamma_i(t))$ et utiliser  $d\alpha = 0$).

\item Montrer que $h_1 - h_0$ est constante et donner une expression
explicite en terme de $\alpha$ pour cette constante
(indication~: utiliser le théorème de Stokes).

\item 
Sur $\Rr^2 \setminus \{\mathbf{0}\}$ on donne la 1-forme
$\alpha = \dfrac{x\,dy - y\,dx}{\sqrt{x^2+y^2}}$.
Montrer que $d\alpha=0$, et dire pourquoi il n'existe
pas de fonction $h: \Rr^2 \setminus \{\mathbf{0}\} \to
\Rr$ telle que $\alpha = dh$.

\item Pourquoi la construction donnée en 1. ne
marche-t-elle pas dans le cas de $\Rr^2 \setminus
\{\mathbf{0}\}$ (voir 3.)~?
\end{enumerate}
\finenonce{006801}



\finexercice


\exercice{6804, gijs, 2011/10/16}
\enonce{006804}{}
Dans $\Rr^3$ on se donne $M = \{\,(x,y,z) \mid 
x^2 + y^2 - z^2 -1 = 0\,\}$, avec la carte $\varphi$
définie par~:
$$
\varphi\ :\ \ ]0,2\pi[\ \times \Rr \quad ,\quad
(\theta,z) \mapsto  (\left(\sqrt{z^2+1}\,\right)
\cos(\theta), \left(\sqrt{z^2+1}\,\right) \sin(\theta), 
z)\ .
$$
On considère aussi la 2-forme $\alpha = z\,dx\wedge dy$
sur
$\Rr^3$.
\begin{enumerate}
\item Calculer la 3-forme $d\alpha$.

\item Calculer la 2-forme $\alpha$ sur $M$ dans la carte
$\varphi$.

\smallskip
Soit $a<b$, et soit $\widehat M$ la partie de $M$
comprise entre $z=a$ et $z=b$, c'est-à-dire 
$\widehat M = \{\,(x,y,z) \mid 
x^2 + y^2 - z^2 -1 = 0\ ,\ a\le z\le b\,\}$.
\item Calculer 
$\int_{\kern1pt{\widehat{\kern-2pt\vrule width0pt
height6pt M\kern2pt}}} \alpha$. 
(Nota Bene: pour
l'orientation ne pas oublier que
$\theta$ est la première coordonnée, et que $z$ est
la deuxième dans la carte $\varphi$.)

\smallskip
Soit $V$ la partie de $\Rr^3$ donné par les
inégalités $a\le z \le b$ et $x^2 + y^2 - z^2 -1\le
0$. On vous demande d'utiliser le théorème de
Stokes pour calculer le volume de
$V$, qui est donné par la formule $\int_V d\alpha$.
Les questions suivantes vous amènent à ce but.

\item \'Enoncer le théorème de Stokes.

\item Décrire le bord $\partial V$ de $V$.

\item Calculer $\int_{\partial V} \alpha$.
\end{enumerate}
\finenonce{006804}



\finexercice

\section{356.00 Orientation}

\section{357.00 Intégration sur les variétés}

\section{358.00 Autre}

\section{370.00 Différentiabilité, calcul de différentielles}


\exercice{2494, sarkis, 2009/04/01}
\enonce{002494}{}
\begin{enumerate}
\item Montrez que $d(x,y)=|x-y|$ est bien une distance sur
l'ensemble des r\'eels. 
\item Pour tout couple d'\'el\'ements
$X=(x_1,...,x_n)$ et $Y=(y_1,...,y_n)$ de $\mathbb{R}^n$, on
d\'efinit $d(X,Y)=\sup_{i=1..n}|x_i-y_i|$. Montrez que $d$ est
bien une distance sur $\mathbb{R}^n$. 
\item Faire de m\^eme avec
$d(X,Y)=|x_1-y_1|+|x_2-y_2|+...+|x_n-y_n|$.
\end{enumerate}
\finenonce{002494} 


\finexercice


\exercice{2495, sarkis, 2009/04/01}
\enonce{002495}{}
 D\'ecrire la boule de centre
l'origine et de rayon $1$ dans les espaces suivants:
\begin{enumerate}
\item $\mathbb{R}$ muni de la distance $d(x,y)=|x-y|.$ 
\item $\mathbb{R}^2$ muni de la distance
$d_1((x_1,x_2),(y_1,y_2))=\sqrt{(x_1-y_1)^2+(x_2-y_2)^2}.$ 
\item $\mathbb{R}^2$ muni de la distance $d_2((x_1,x_2),(y_1,y_2))=\sup
(|x_1-y_1|,|x_2-y_2|).$ 
\item $\mathbb{R}^2$ muni de la distance
$d_3((x_1,x_2),(y_1,y_2))=|x_1-y_1|+|x_2-y_2|.$
\end{enumerate}
Montrez que les 3 dernières distances sont équivalentes.
\finenonce{002495} 


\finexercice


\exercice{2496, sarkis, 2009/04/01}
\enonce{002496}{}
 Soit $E$ l'ensemble des
fonctions continues de l'intervalle $[0,1]$ dans $\mathbb{R}$ qui
sont continues. Montrez que l'application $\|f\|_1=\int_a^b
|f(t)|dt$ est une norme sur $E$. Montrez que $E$ n'est pas
complet.
\finenonce{002496} 


\finexercice


\exercice{2497, sarkis, 2009/04/01}
\enonce{002497}{}
Etudiez la continuit\'e des
applications suivantes:
\begin{enumerate}
\item $f(x)=\frac{xy^2}{x^2+y^2}.$ \item
$f(x)=\frac{xy}{x^2+y^2}$. \item
$f(x)=\frac{exp(\frac{-1}{x^2+y^2})}{|x|+|y|}.$
\end{enumerate}
\finenonce{002497} 


\finexercice


\exercice{2498, sarkis, 2009/04/01}
\enonce{002498}{}
Soient $E$ et $F$ deux espaces
norm\'es r\'eels et $f:E \rightarrow F$ une application born\'ee
sur la boule unit\'e de $E$ et v\'erifiant
$$f(x+y)=f(x)+f(y) \mbox{ pour tout } x,y \in E.$$
Montrez que $f$ est lin\'eaire continue.
\finenonce{002498} 


\finexercice


\exercice{2499, sarkis, 2009/04/01}
\enonce{002499}{}
Soient $||.||_1$ et $||.||_2$
deux normes sur $\mathbb{R}^2$ et $M=\left (
\begin{array}{cc}
a & b \\
c & d
\end{array}
\right )$ une matrice de ${\cal M}_{n,n}(\mathbb{R} \mbox{ ou }
\mathbb{C})$. On d\'efinit la norme de $M$ (ou de l'application
lin\'eaire associ\'ee) de la mani\`ere suivante:
$$||M||=\sup_{X \in S_1(0,1)}||M.X||_2$$
o\`u $S_1(0,1)$ est la sph\`ere unit\'e pour la norme $||.||_1$.
Dans chacun des cas suivant, calculez la norme de $M$.
\begin{enumerate}
\item $||(x,y)||_1=||(x,y)||_2=sup(|x|,|y|).$ 
\item $||(x,y)||_1=||(x,y)||_2=\sqrt{x^2+y^2}.$ 
\item $||(x,y)||_1=\sqrt{x^2+y^2}$ et $||(x,y)||_2=sup(|x|,|y|).$
\end{enumerate}
\finenonce{002499} 


\finexercice


\exercice{2500, sarkis, 2009/04/01}
\enonce{002500}{}
Continuit\'e sur
$\mathbb{R}^2$ des fonctions suivantes:
\begin{enumerate}
\item $f(x,y)=\frac{xy}{x^2+y^2}$ 
\item $f(x,y)=\frac{x+y}{x^2+y^2}$ 
\item $f(x,y)=\frac{x^3y^2}{x^2+y^2}$
\end{enumerate}
\finenonce{002500}


\finexercice


\exercice{2501, sarkis, 2009/04/01}
\enonce{002501}{}
Calculez la norme des op\'erateurs suivants:
\begin{enumerate}
\item Le shift sur $l^\infty$ d\'efini par $S(x)_{n+1}=x_n,
S(x)_0=0$ (sur $l^\infty$ on d\'efinit $||(x_n)||_\infty=\sup_{n
\in \mathbb{N}}|x_n|)$. 
\item $X={\mathcal C}([0,1])$ avec la
norme sup et l'op\'erateur $Tf(x)=f(x)g(x)$ o\`u $g \in X$. 
\item $X={\mathcal C}([0,1])$ muni de la norme sup et
$u(f)=\int_0^1f(x)g(x)dx$ o\`u $g \in X$ est une fonction qui
s'annule qu'en $x=1/2$. 
\item $X=l^2$ et $u(x)=\sum a_nx_n$ o\`u $(a_n)$ est dans $X$. 
\item $X$ l'espace des suites convergentes
muni de la norme sup et $u: X \rightarrow \mathbb{R}$
l'application $u(x)=\lim_{j \rightarrow \infty} x_j$.
\end{enumerate}
\finenonce{002501} 


\finexercice


\exercice{2502, sarkis, 2009/04/01}
\enonce{002502}{} 
Soit $X={\mathcal C}([0,1])$
avec la norme $||f||=\int_0^1 |f(t)|dt$. Montrez que la forme
lin\'eaire $T: X \rightarrow \mathbb{R}$ d\'efinie par $T(f)=f(0)$
n'est pas continue en $0$. Que peut-on en d\'eduire pour le
sous-espace des fonctions de $X$ nulles en $0$ ?
\finenonce{002502} 


\finexercice

\exercice{2503, queffelec, 2009/04/01}
\enonce{002503}{} 
Soit $f$ une application
$f$ de $E$ dans $F$ espaces vectoriels norm\'es de dimension
finie.


 On rappelle les implica\-tions suivantes : si $x_0\in E$,  ``$f$ de
classe $C^1$ en $x_0$'' $\Rightarrow$ ``$f$ diff\'erentiable en
$x_0$'' $\Rightarrow$ ``$f$ continue en $x_0$''.
 On sait de m\^eme que ``$f$ diff\'erentiable en
$x_0$'' $\Rightarrow$ ``$f$ admet des d\'eriv\'ees partielles en
$x_0$'' montrer que les r\'eciproques sont fausses en g\'en\'eral
en s'inspirant de :

$$f(x)= \left
\{\begin{array}{ccc}
        & x^2\sin{\frac 1 x}+y^2\sin{\frac 1 y} & {\rm si}\ \ xy\neq0\\
        & x^2\sin{\frac 1 1 x}     &{\rm si}\ \ y=0\\

        & y^2\sin{\frac 1 y}     &{\rm si}\ \ x=0\\
        &0      &{\rm en}\ \ (0,0)\\

 \end{array}\right.$$



ou de
$$f(x)= \left \{\begin{array}{ccc}
        & {\frac{xy^2}{x^2+y^2}}& {\rm si}\ \ (x,y)\neq(0,0)\\
        &0    &{\rm si}\ \ (x,y)=(0,0)\\
  \end{array}\right.$$
\finenonce{002503} 


\finexercice


\exercice{2504, queffelec, 2009/04/01}
\enonce{002504}{}
\begin{enumerate}
\item Soit $f$ une application de $E$ dans $F$ espaces vectoriels
norm\'es et supposons $f$ diff\'eren\-tiable en $a$; montrer que
pour tout vecteur $u\in E^*$, la d\'eriv\'ee de $f$ en $a$ dans la
direction $u$ existe , i.e. $\lim_{h\to 0} {\frac 1 h}\big(f(a+hu)-f(a)\big)$ et l'exprimer \`a l'aide de $f'(a)$.


\item On consid\`ere $f:{\Rr^2}\to{\Rr}$ d\'efinie par $f(0,0)=0$
et, si $(x,y)\neq(0,0)$,\ $f(x,y)={\frac{x^3y}{x^4+y^2}}$. Montrer
que $f$ est d\'erivable en $(0,0)$ dans toutes les directions,
mais que $f$ n'est pas diff\'erentiable en $(0,0)$.
\end{enumerate}
\finenonce{002504} 


\finexercice


\exercice{2505, queffelec, 2009/04/01}
\enonce{002505}{} 
Soit $g:{\Rr}\to{\Rr}$ une
application de classe $C^2$ et $F:{\Rr^2}\to{\Rr}$ d\'efinie par
$$F(x,y)={\frac{g(x)-g(y)}{x-y}}\ \hbox{si}\ x\neq y,\ \ F(x,x)=g'(x).$$
Montrer que $F$ est de classe $C^1$ en tout point de ${\Rr^2}$ et
calculer sa diff\'erentielle.
\finenonce{002505} 


\finexercice


\exercice{2506, queffelec, 2009/04/01}
\enonce{002506}{}
Soit $E^n$ l'espace des
polyn\^omes de degr\'e $\leq n$. Etudier la diff\'erentiabilit\'e
des applications $P\mapsto\int_0^1(P^3(t)-P^2(t))\ dt$ et \
$P\mapsto P'-P^2$. 
\finenonce{002506} 


\finexercice


\exercice{2507, queffelec, 2009/04/01}
\enonce{002507}{} 
Soit $f$ une application
diff\'erentiable de $\Rr^2$ dans lui-m\^eme, propre (i.e.
$||f(x)||$ tend vers $\infty$ quand $||x||\to\infty$), telle que
pour tout $x\in\Rr^2$ $Df(x)$ soit injective. On va montrer que
$f$ est surjective. Soit $a\in\Rr^2$ et $g(x)=||f(x)-a||^2$;
\begin{enumerate}
\item Calculer $Dg(x)$.

\item Montrer que $g$ atteint sa borne inf\'erieure en un point
$x_0$ de $\Rr^2$, et que $Dg(x_0)=0$; en d\'eduire le r\'esultat.

\end{enumerate}
\finenonce{002507} 


\finexercice


\exercice{2508, queffelec, 2009/04/01}
\enonce{002508}{}
Soit, dans ${\Rr^n}$, $F$
un sous-espace ferm\'e, et soit $f:{\Rr^n}\to\Rr$ d\'efinie par
$f(x)=d(x,F)$. On rappelle que $f$ est $1$-lipschitzienne, et que
pour chaque $x$ il existe $y\in F$ tel que $f(x)=d(x,y)$.
\begin{enumerate}
\item On suppose que $f$ est diff\'erentiable en $x\notin F$.
Montrer que $||Df(x)||_{{\cal L}(\Rr^n,\Rr)}\leq 1$.

\item On consid\`ere la fonction $\varphi : t\in[0,1]\to
f((1-t)x+ty)$; en calculant $\varphi'(0)$ de deux fa\c cons,
montrer que $Df(x).{\frac{x-y}{||x-y||}}=1$ et $||Df(x)||_{{\cal
L}(\Rr^n,\Rr)}= 1$.

\item En d\'eduire que $y$ est unique.
\end{enumerate}
\finenonce{002508} 


\finexercice


\exercice{2509, queffelec, 2009/04/01}
\enonce{002509}{}
Soit $E$ un espace de
Banach et ${\cal L}(E)$ l'espace des endomorphismes lin\'eaires
continus de $E$.
\begin{enumerate}
\item Soit $A\in{\cal L}(E)$; montrer que l'application
$\varphi:t\in{\Rr}\to e^{tA}$ est d\'erivable et calculer sa
d\'eriv\'ee.

\item On suppose que la norme de $E$ est associ\'ee au produit
scalaire $\langle \cdot,\cdot\rangle$. Soit $x\in E$. Montrer que
l'application $\Phi:t\to \langle e^{tA}x,e^{tA}x\rangle$ est
d\'erivable et calculer sa d\'eriv\'ee.

\item On suppose que $A$ est antisym\'etrique. Montrer que pour
tout $t$, $e^{tA}$ est unitaire.
\end{enumerate}
\finenonce{002509} 


\finexercice

\exercice{2510, mayer, 2009/04/01}
\enonce{002510}{}
Soit $\alpha >0$. \'Etudier la
diff\'erientiabilit\'e \`a l'origine de l'application $f:\Rr^2 \to
\Rr$ qui est d\'efinie par $f(0,0) =0$ et par
$$ f(x,y) = \frac{|xy|^\alpha}{\sqrt{x^2+3y^2}} \quad \text{si} \;\;
(x,y) \neq (0,0)\; .$$ 
\finenonce{002510} 


\finexercice


\exercice{2511, mayer, 2009/04/01}
\enonce{002511}{} 
Soit $f:\Rr ^2 \to \Rr$
d\'efinie par
$$f(x,y)=\frac{xy^2}{x^2+y^2} \quad \text{si} \quad  (x,y)\neq (0,0)$$
 et $f(0,0)=0$. Montrer que $f$ est continue sur $\Rr^2$,
que pour tout $u\in \Rr^2 \setminus \{0\}$ $\frac{\partial
f}{\partial u}(0,0)$ existe, mais que $f$ n'est pas
diff\'erentiable en $(0,0)$. 
\finenonce{002511} 


\finexercice


\exercice{2512, mayer, 2009/04/01}
\enonce{002512}{} 
Soit $X= {\cal C} ([0,1])$
muni de la norme uniforme et soit $f$ une application de ${\cal
C}^1 (\Rr , \Rr )$. On note $F$ l'application $\varphi \mapsto
f\circ \varphi $ de $X$ dans $X$. Montrer que pour chaque $\varphi
\in X$, $DF(\varphi ) $ est l'op\'erateur lin\'eaire de
multiplication par $f'\circ \varphi $ dans $X$:
$$ DF(\varphi )\cdot ( h ) = h \, f'\circ \varphi \; ,$$
et que $DF$ est continue.
\finenonce{002512} 


\finexercice


\exercice{2513, mayer, 2009/04/01}
\enonce{002513}{}
Soit ${\cal F}$ l'alg\`ebre
des matrices carr\'es $p\times p$ munie d'une norme.
\begin{enumerate}

\item Soit $f:\mathcal{F} \to \Rr$ l'application qui associe \`a une matrice $A$
son determinant 
$f(A) = \det (A)$.  
Montrer qu'elle est diff\'erentiable et d\'eterminer $Df$. 

\item Pour
$n\geq 1$, on consid\`ere l'application $\varphi _n (A) =A^n $ de
${\cal F} $ dans ${\cal F}$. Montrer qu'elle est diff\'erentiable
en toute matrice $A\in {\cal F}$. 
\item On d\'esigne par $U$
l'ensemble des matrices inversibles de ${\cal F}$. Montrer que $U$
est un ouvert de ${\cal F}$ et calculer la diff\'erentielle de
l'application $A\mapsto A^{-1}$ de $U$ dans $U$.
\end{enumerate}
\finenonce{002513} 


\finexercice


\exercice{2514, mayer, 2009/04/01}
\enonce{002514}{}
\begin{enumerate}
\item Que peut-on dire de la diff\'erentiabilit\'e de
l'application $f:\Rr^2 \to \Rr$ d\'efinie par $f(x_1,x_2) = \|x\|
_\infty = \max{(|x_1|, |x_2|)}$? \item G\'en\'eraliser ceci \`a
$f:{\cal F} \to \Rr$, $f(x) = \|x\| _\infty$, avec ${\cal F} =
\Rr^n$ ou ${\cal F}$ l'ensemble des suites convergentes vers zero.
\end{enumerate}
\finenonce{002514} 


\finexercice


\exercice{2515, mayer, 2009/04/01}
\enonce{002515}{}
Soit $f:\Rr^2 \to \Rr$
l'application $x=(x_1,x_2) \mapsto \|x\|_1=|x_1|+ |x_2|$. Est-ce
qu'elle est diff\'erentiable?

Consid\'erons maintenant $l^1$ l'espace des suites r\'eelles muni
de la norme $\|x\|_1=\sum_{j=1}^\infty |x_j|$.
\begin{enumerate}
\item Montrer que pour toute forme lin\'eaire continue $L$ sur
$l^1$ il existe une suite born\'ee $\alpha =
(\alpha_1,\alpha_2,....)$ telle que
$$ L(x) =\sum _{j=1}^\infty \alpha _j x_j \;\; .$$
\item Montrer que la norme $\|.\|_1 : l^1 \to \Rr$ n'est pas
 diff\'erentiable en aucun point de $l^1$ (raisonner par l'absurde en utilisant (1.)).
\end{enumerate}
\finenonce{002515} 


\finexercice


\exercice{2516, mayer, 2009/04/01}
\enonce{002516}{} 
Dans un espace norm\'e $({\cal
F} , N)$, on consid\'ere l'application
 $x \mapsto  N(x)$.
Rappeler que, lorsque cette application $N$ est diff\'erentiable
en $x\in {\cal F}$, alors
$$ DN(x) \cdot (h) = \lim _{t\rightarrow 0} \frac{1}{t}\left( N (x +th) -N (x) \right)\; .$$
En d\'eduire que $N$ n'est pas diff\'erentiable en $0\in {\cal
F}$. Supposons $N$ diff\'erentiable en $x\in {\cal F}$, alors
justifier que $N$ l'est aussi en $\lambda x$, o\`u $\lambda >0$,
et que $DN(x) =DN(\lambda x )$. En consid\'erant la d\'eriv\'ee en
$\lambda =1$ de l'application $\lambda \mapsto N (\lambda x)$,
montrer que $DN(x)  \cdot (x)=N(x)$ et en d\'eduire $\|| DN(x) \||
=1$. 
\finenonce{002516} 


\finexercice


\exercice{2517, mayer, 2009/04/01}
\enonce{002517}{} 
Soit ${\cal E}$ un espace
vectoriel r\'eel muni d'un produit scalaire $(x,y) \mapsto \langle
x,y \rangle$ et de la norme associ\'ee $\|x\| = \langle x,x
\rangle ^\frac{1}{2}$.
 Soit $u$ un endomorphisme continu de ${\cal E}$ que l'on suppose sym\'etrique, i.e.
$$\langle u(x),y\rangle =\langle x,u(y)\rangle  \quad  \text{pour tout} \;\; x,y\in {\cal E} \; .$$
\begin{enumerate}
\item Montrer que l'application $x\in {\cal E} \mapsto \langle
u(x),x\rangle $ est diff\'erentiable sur ${\cal E}$ et calculer sa
diff\'erentielle. L'application $x\mapsto \|x\|^2$ est donc
diff\'erentiable. \item On d\'efinit une application $\varphi :
{\cal E} \setminus \{0\} \to \Rr$ en posant $\varphi (x) =
\frac{\langle u(x),x\rangle }{\langle x,x\rangle }$. \'Etablir
qu'il s'agit d'une application diff\'erentiable. Calculer ensuite
$D\varphi$. Montrer que, pour un \'el\'ement non nul $a\in {\cal
E}$, on a $D\varphi (a) =0$ si et seulement si $a$ est vecteur
propre de $u$.
\end{enumerate}
\finenonce{002517} 


\finexercice


\exercice{2518, queffelec, 2009/04/01}
\enonce{002518}{}
\begin{enumerate}
\item Soit $f$ une application r\'eelle continue et d\'erivable
sur $]a,b[$ telle que $f'(x)$ ait une limite quand
$x\buildrel{<}\over{\to}b$; alors $f$ se prolonge en une fonction
continue et d\'erivable \`a gauche au point $b$.

\item Soit $f$ une application continue et d\'erivable sur un
intervalle $I\subset\Rr$, et de d\'eriv\'ee croissante; montrer
que $f$ est convexe sur $I$ i.e. $f((1-t)x+ty)\leq
(1-t)f(x)+tf(y)$ pour tous $x<y$ de $I$ et $t\in[0,1]$. (Poser
$z=(1-t)x+ty$ et appliquer les AF \`a $[x,z]$ puis $[z,y]$.)
\end{enumerate}
\finenonce{002518} 


\finexercice


\exercice{2519, queffelec, 2009/04/01}
\enonce{002519}{}
Montrer que l'identit\'e des accroissements finis n'est pas
vraie pour les fonctions vectorielles en consid\'erant
$f(x)=e^{ix}$.
\finenonce{002519} 


\finexercice


\exercice{2520, queffelec, 2009/04/01}
\enonce{002520}{partiel du 5 d\'ecembre 1999} 
Soit $f:{\Rr^2}\to{\Rr^2}$ d\'efinie par
$f(x,y)=(x^2-y,x^2+y^2)$ et $g=f\circ f$.
\begin{enumerate}
\item Montrer que $f$ et $g$ sont de classe $C^1$.

\item Calculer en tout point $(x,y)\in{\Rr^2}$ la matrice
jacobienne de $f$ not\'ee $Df{(x,y)}$; calculer la matrice
jacobienne de $g$ au point $(0,0)$ not\'ee $Dg{(0,0)}$.

\item Montrer qu'il existe $\rho>0$ tel que pour tout $(x,y)\in
\overline{B_\rho((0,0))}$ (la boule ferm\'ee de centre $(0,0)$ et
de rayon $\rho$) on a $||Dg{(x,y)}||\leq {\frac 1 2}$.

\item Montrer que la fonction $g$ admet un unique point fixe dans
$\overline{B_\rho((0,0))}$.
\end{enumerate}
\finenonce{002520} 


\finexercice


\exercice{2521, queffelec, 2009/04/01}
\enonce{002521}{}
On consid\`ere
l'application $F:{\Rr^2}\to{\Rr^2}$ d\'efinie par $F(x,y)=(\cos
x-\sin y,\ \sin x-\cos y)$; on note $F^{(k)}$ l'application $F$
compos\'ee $k$-fois
\begin{enumerate}
\item Montrer que $||DF(x,y)||\leq \sqrt2$ pour tout $(x,y)$.

\item En d\'eduire que la suite r\'ecurrente d\'efinie par
$x_0,y_0$ et pour $n\geq 1$
$$x_{n+1}=
{\frac 1 2}(\cos x_n-\sin y_n),\quad y_{n+1}= {\frac 1 2}(\sin
x_n-\cos y_n)$$ converge pour tout $(x_0,y_0)$. Donnez
l'\'equation que v\'erifie sa limite ?
\end{enumerate}
\finenonce{002521} 


\finexercice


\exercice{2522, queffelec, 2009/04/01}
\enonce{002522}{}
Soit $f$ une application
diff\'erentiable de $]a,b[\subset\Rr$ dans $\Rr^n$; on suppose
qu'il existe $k>0$ tel que
$$||f'(x)||\leq k||f(x)||, \hskip2mm \forall x\in]a,b[.$$
Montrer que si $f$ s'annule en un point $x_0\in ]a,b[$, $f$ est
identiquement nulle dans $]a,b[$ (montrer que $E=\{x\in ]a,b[\ ;\
f(x)=0\}$ est ouvert). 
\finenonce{002522} 


\finexercice


\exercice{2523, queffelec, 2009/04/01}
\enonce{002523}{}
Soit $E$ un espace de
Banach, $U$ un ouvert de $E$ et $f$ une application
diff\'eren\-tiable de $U$ dans $\Rr$ telle que l'on ait
$||f'(x)||\leq k |f(x)|,\hskip2mm \forall x\in U$. Montrer que
pour $x$ assez voisin de $a\in U$,
$$|f(x)|\leq e^{k||x-a||}\ |f(a)|.$$

\emph{Indication :} consid\'erer l'application $t\in[0,1]\to
f(a+t(x-a))$). 
\finenonce{002523} 


\finexercice


\exercice{2524, queffelec, 2009/04/01}
\enonce{002524}{}
On consid\`ere
l'application $F:{\Rr^2}\to{\Rr^2}$ d\'efinie par
$F(x,y)=(x^2+y^2,y^2)$; on note $F^{(k)}$ l'application $F$
compos\'ee $k$-fois avec elle-m\^eme. On consid\`ere
$\Omega=\{(x,y)\in {\Rr^2}\ /\
\lim_{k\to\infty}F^{(k)}(x,y)=(0,0)\}$.
\begin{enumerate}
\item V\'erifier que $(x,y)\in\Omega\Longleftrightarrow
F(x,y)\in\Omega$.

\item Montrer qu'il existe $\varepsilon>0$ tel que
$||(x,y)||<\varepsilon\Longrightarrow||F'(x,y)||\leq{\frac 1 2}$; en
d\'eduire que $0$ est int\'erieur \`a $\Omega$ puis que $\Omega$
est ouvert.

\item Montrer que  $\Omega$ est connexe.
\end{enumerate}
\finenonce{002524} 


\finexercice


\exercice{2525, mayer, 2009/04/01}
\enonce{002525}{}
On consid\`ere l'application
$F:\Rr ^2 \to \Rr ^2$ d\'efinie par $$F(x,y) = (x^2+y^2, y^2)\;
.$$ Soit $\Omega = \{p\in \Rr^2 \, ; \; \lim _{k\rightarrow \infty
} F^k (p) = (0,0)\}$.
\begin{enumerate}
\item V\'erifier que $p\in \Omega$ si et seulement si $F(p)\in
\Omega$. \item Montrer qu'il existe $\delta >0$ tel que $\| | D
F(p) \| | <\frac{1}{2}$ si $\|p\|<\delta$. En d\'eduire que $(0,0)
$ est dans l'int\'erieur de $\Omega$ puis que $\Omega$ est un
ouvert. \item Utiliser l'homog\'en\'eit\'e de $F$ pour montrer que
$\Omega $ est connexe.
\end{enumerate}
\finenonce{002525} 


\finexercice


\exercice{2526, mayer, 2009/04/01}
\enonce{002526}{}
Soit $\Omega$ un ouvert
convexe de $\Rr^n$ et $f:\Omega \to \Rr^n$ une application de
classe $C^1$ qui est injective sur $\Omega$ et telle que $Df(x)$
soit injective pour tout $x\in \Omega$.
Montrer que, pour tous $a,b \in \Omega$,
$$ \| f(b) -f(a)- Df(a) (b-a)\| \leq \| b -a \| \sup _{c\in [a,b]} \| Df(c) -Df(a)\| \; .$$
\finenonce{002526} 


\finexercice


\exercice{6255, queffelec, 2011/10/16}
\enonce{006255}{} 
Soit $H$ un espace préhilbertien sur $\Rr$, et $f(x)=||x||$ de $H$ dans $\Rr$; 
montrer que $f$ est différentiable en tout point de $H\backslash\{0\}$,
et calculer sa différentielle.
(indic. étudier directement $||x+h||$ ou considérer la fonction composée
$x\to||x||^2\to\sqrt{||x||^2}$.)
Décrire le noyau Ker$f'(x)$ 
en tout $x\neq0$.
\finenonce{006255}



\finexercice


\exercice{6256, queffelec, 2011/10/16}
\enonce{006256}{}
Soit $a\in \Rr^n$ et $f:{\Rr^n}\backslash\{a\}\to {\Rr^n}$ définie par
$f(x)={a-x\over ||x-a||^2}$.
\begin{enumerate}
\item Calculer $f'(x)$ pour tout $x\in{\Rr^n}\backslash\{a\}$.

\item Montrer que $f'(x).h={Sh\over ||x-a||^2}$ où $S$ est la symétrie orthogonale
d'axe
$x-a$. Que peut-on dire de la transformation $f'(x)$ de ${\Rr^n}$ ?
\end{enumerate}
\finenonce{006256}



\finexercice


\exercice{6257, queffelec, 2011/10/16}
\enonce{006257}{}
Soit $B$ une application bilinéaire de $E\times F$ dans $G$, où $E,F,G$ sont
des evn de dimension finie.
\begin{enumerate}
\item Calculer $B'(a)$ sa différentielle en un point $a=(a_1,a_2)$ de $E\times F$.

\item En déduire, pour $f$ et $g$ deux applications différentiables de $I$
intervalle de $\Rr$ dans $\Rr^3$, la différentielle de $t\to f(t)\wedge
g(t)$ et de $t\to \langle f(t),g(t)\rangle$ en tout $t\in I$.

\item Application : Soit $A$ un opérateur de $\Rr^n$ tel que $Ax\bot x$ pour
tout $x$; montrer que $e^{tA}$ est une isométrie pour tout réel $t$. (Dériver
$t\to ||e^{tA}x||^2$.)
\end{enumerate}
\finenonce{006257}



\finexercice


\exercice{6258, queffelec, 2011/10/16}
\enonce{006258}{}

Soit $E$ et $F$ deux evn sur $\Cc$. Une application de $E$ dans
$F$ $\Cc$-linéaire est $\Rr$-linéaire, mais la réciproque est fausse.
\begin{enumerate}
\item Soit $\varphi:E\to F$ une application $\Rr$-linéaire. Montrer que
$\varphi$ est  $\Cc$-linéaire si et seulement si $\varphi(ix)=i\varphi(x)$
pour tout $x\in E$. En déduire les applications de $\Rr^2$ dans $\Rr^2$ qui
sont
$\Cc$-linéaires.

Soit $U$ un ouvert de $E$ et $f:U\to F$. On suppose $f$ 
$\Rr$-différentiable en $a\in U$. Il est clair que $f$ est 
$\Cc$-différentiable en $a$ si et seulement si $f'(a)$ est 
$\Cc$-linéaire.

\item Si $f:{\Cc}\to\Cc$ s'écrit $f(z)=u(z)+iv(z)=f(x+iy)$ avec $u$ et $v$
réelles, qu'on identifie à $f(x,y)=(u(x,y),v(x,y))$, traduire à l'aide de a)
``$f$ est
$\Cc$-différentiable en $a=\alpha+i\beta$''.
En quels points les applications de
$\Cc$ dans
$\Cc$ sont-elles
$\Cc$-différentiables :
$f_1(z)=e^x$;\  $f_2(z)=|z|^2$;\ $f_3(z)=e^{x-iy}$ ?

\item (extrait de septembre 99)   
 Soit $U$ un ouvert de ${\Cc}$ et soit $f:U\to\Cc$ 
$\Cc$-différentiable en $a=\alpha+i\beta\in U$, telle que $f(a)\not=0$.
Montrer que si $g=|f|$ est $\Cc$-différentiable en $a=\alpha+i\beta\in U$,
alors $f'(a)=0$. 
\end{enumerate}


\finenonce{006258}


\finexercice


\exercice{6259, queffelec, 2011/10/16}
\enonce{006259}{}
\begin{enumerate}
\item Montrer que l'identité des accroissements finis n'est pas vraie pour les
fonctions vectorielles en considérant $f(x)=e^{ix}$.

\item Si $f$ est une fonction continue sur $[a,b]$ et dérivable sur $]a,b[$, on a
vu que $f'(]a,b[)$ est connexe. Montrer que ceci est faux  pour les
fonctions vectorielles en considérant $f(x)=(x^2\cos({1\over
x}),x^2\sin({1\over x}))$.
\end{enumerate}
\finenonce{006259}



\finexercice


\exercice{6260, queffelec, 2011/10/16}
\enonce{006260}{}
Soit $E=C([0,1],\Rr^n)$ et soit $F$ un sous-espace vectoriel de $E$ constitué
de fonctions différentiables, telles que
$$||f'(x)||\leq M, \hskip2mm \forall x\in[0,1],\hskip2mm \forall f\in F,\
||f||\leq1$$ où $M$ est une constante fixée à l'avance.
Montrer que la boule unité de $F$ est compacte; que peut-on dire de $F$ ?
\finenonce{006260}



\finexercice


\exercice{6261, queffelec, 2011/10/16}
\enonce{006261}{} 
Soient $E,F$ des espaces normés, $\Omega$ un ouvert de $E$
 et $f:\Omega \to F$ une application continue.
 \begin{enumerate}
 \item Soit $a$ un point de $\Omega$.
Si $f$ est différentiable dans $\Omega \setminus \{a\}$ et si
l'application $x\in \Omega \setminus \{a\} \mapsto Df(x)$ admet
une limite $T\in {\cal L} (E,F)$ quand $x$ tend vers $a$ dans $\Omega$,
montrer que $f$ est différentiable au point $a$ et que $Df(a)=T$
(appliquer le théorème des accroissements finis à la
fonction $g: x\mapsto f(x) -T(x)$).
\item Supposons $f$ différentiable dans $\Omega$. Montrer que
$Df:\Omega \to {\cal L}(E,F) $ est continue en $a\in \Omega$ si et
seulement si, pour tout $\epsilon >0$, il existe $\delta
>0$ tel que
$$ \|f(a+h)-f(a+k) -Df(a) (h-k) \| \leq \epsilon \| h-k \| \quad \text{si} \;\; \|h \| <\delta \; et \; \|k\| <\delta \; .$$
\item Supposons maintenant qu'il existe une application continue
$x\in \Omega \mapsto T_x \in {\cal L} (E,F)$ telle que pour tout $x\in \Omega$ et tout
$h\in E$
$$ \lim_{t\to 0 , t\neq 0} \frac{f(x+th)-f(x)}{t} =T_x (h) \; .$$
Montrer que $f$ est de classe ${\cal C}^1$ et que $Df(x)=T_x$ pour tout
$x\in \Omega$. (On pourra considérer la fonction $g(t) =
f(x+th)-tT_x (h)$.) 
\end{enumerate}
\finenonce{006261}



\finexercice


\exercice{6262, queffelec, 2011/10/16}
\enonce{006262}{} 
Soient $E,F$ des espaces de Banach, $\Omega$ un ouvert connexe de $E$
et $f_n: \Omega \to F$ une suite d'applications différentiables.
On suppose que cette suite vérifie :
\begin{enumerate}
\item[(i)] Il existe $x_0\in \Omega$ tel que $(f_n (x_0))$ converge dans $F$.
\item[(ii)] La suite $(Df_n)$ converge uniformément sur toute boule fermée
$B_F (a,r) \subset \Omega$. 
\end{enumerate} 
Alors, montrer que $(f_n)$
converge uniformément sur toute boule fermée de $\Omega$ et
que, si $f(x) = \lim _{n\to \infty } f_n(x)$ et $L_x = \lim _{n\to
\infty } Df_n(x)$, alors $f$ est différentiable avec $Df(a) =
L_a$, $a\in \Omega$.
\finenonce{006262}



\finexercice


\exercice{6263, queffelec, 2011/10/16}
\enonce{006263}{} 
Soit $\Omega$ un ouvert convexe de $\Rr^n$ et $f:\Omega \to \Rr^n$ une
application de classe $C^1$ qui est injective sur $\Omega$ et
telle que $Df(x)$ soit injective pour tout $x\in \Omega$. 
\begin{enumerate}
\item Montrer que, pour tous $a,b \in \Omega$,
$$ \| f(b) -f(a)- Df(a) (b-a)\| \leq \| b -a \| \sup _{c\in [a,b]} \| Df(c) -Df(a)\| \; .$$
\item Soit $(f_n)$ une suite de fonctions de classe $C^1$ telle que $f_n \to f$ et $Df_n \to Df$
uniformément sur tout compact de $\Omega$. On va montrer:
{\it pour tout compact $K$ de $\Omega$ il existe $n_0$ tel que $f_n$ soit injective
sur $K$ pour $n\geq n_0$.} 
\begin{itemize}
\item En raisonnant par l'absurde, montrer qu'il existerait $K$ compact et, pour une infinité d'entiers
$n$, des points $a_n, b_n \in K$ tels que $f_n(a_n) = f_n (b_n)$.
\item Quitte à extraire, montrer qu'alors $b_n - a_n \to 0$.
\item Utiliser (1.) pour en déduire une contradiction.
\end{itemize}
\end{enumerate}
\finenonce{006263}



\finexercice


\exercice{6796, gijs, 2011/10/16}
\enonce{006796}{}
\begin{enumerate}
\item Enoncer le théorème des accroissements finis pour
une fonction $f:\Rr^n \to \Rr^p$.

\item Soit $B_\epsilon(x_o)$ la boule de rayon
$\epsilon$ dans $\Rr^n$, et soit $f:
\overline{B_\epsilon(x_o)} \to \Rr$ une fonction de
classe $C^1$. Montrer que $f$ est Lipschitzienne et donner
une expression de son rapport.
\end{enumerate}
\finenonce{006796}



\finexercice


\exercice{6832, gijs, 2011/10/16}
\enonce{006832}{}
Soit $U$ un ouvert de $\Rr^n$ et $f: U \to \Rr^p$ une
application dérivable.
\begin{enumerate}
\item \'Enoncer l'inégalité des accroissements
finis. (Pour cette question on peut supposer que $U$ est
con{\bf v}exe.)

\item Démontrer, à l'aide de 1., la proposition
suivante :

Si pour tout $x\in U$ la dérivée
de $f$ en $x$ est nulle~: $Df(x) = 0$, alors pour tout $x$
dans $U$ il existe un voisinage $V$ de $x$ dans $U$ (par
exemple une boule centrée en $x$) tel que $f$ est
constante sur $V$.

\item \`A l'aide de 2., démontrer que si en plus
$U$ est co{\bf nn}exe, alors $f$ est constante sur $U$.
\end{enumerate}
\finenonce{006832}



\finexercice


\exercice{6839, gijs, 2011/10/16}
\enonce{006839}{}
Soit $O \subset \Rr^n$ un ouvert et $f:O \to \Rr^p$ une
fonction. Supposons qu'il existe une fonction continue
$L:O \times O \to \mathcal{L}(\Rr^n; \Rr^p)$ (l'ensemble
des applications linéaires de $\Rr^n$ dans $\Rr^p$ muni
de la norme d'opérateurs) telle que pour tout $x, y\in O$
on a
$$
f(x) -f(y) = L(x,y)\,(x-y)
\ .
$$
Démontrer que $f$ est de classe $C^1$ sur $O$ et que
$Df(x) = L(x,x)$.
\finenonce{006839}



\finexercice

\section{371.00 Différentielle d'ordre supérieur, formule de Taylor}


\exercice{2553, tahani, 2009/04/01}
\enonce{002553}{} 
Calculez $D^2f(x)$ dans les
cas suivants:
\begin{enumerate}
\item $f\in L(E,G)$ continue 
\item $f: E \times F \rightarrow G$, bilinéaire continue. 
\item $f:M_n(\mathbb{R}) \rightarrow M_n(\mathbb{R})$, $f(A)=A^2$
\end{enumerate}
\finenonce{002553} 


\finexercice


\exercice{2554, tahani, 2009/04/01}
\enonce{002554}{}
Etudier les extrémas locaux
et globaux des fonctions suivantes:
\begin{enumerate}
\item $f(x,y)=x^2+xy+y^2+\frac{1}{4}x^3$ 
\item $f(x,y)=x^2y-x^2/2-y^2$ 
\item $f(x,y)=x^4+y^4-2(x-y)^2$ 
\item $f(x,y)=\sin^2x-\sh^2y$ 
\item $f(x,y)=x^3+y^3$ 
\item $f(x,y)=y^2-3x^2y+2x^4$
\end{enumerate}
\finenonce{002554} 


\finexercice


\exercice{2555, tahani, 2009/04/01}
\enonce{002555}{}
Trouver le volume maximum
d'une boite rectangulaire inscrite dans la sph\`ere d'\'equation
$x^2+y^2+z^2=R^2$.
\finenonce{002555} 


\finexercice


\exercice{2556, tahani, 2009/04/01}
\enonce{002556}{}
D\'eterminez le
parall\'elipipède rectangle de volume $V$ donn\'e dont la surface
totale est minimale.
\finenonce{002556} 


\finexercice

\exercice{6287, mayer, 2011/10/16}
\enonce{006287}{Rappel du Cours} 
Soient $E_1,E_2$ et $F$ des espaces normés et $B: E_1\times E_2
\to F$ une application bilinéaire continue. Montrer que $B$ est
de classe $C^\infty$ et déterminer les différentielles $D^k
B$.
\finenonce{006287}



\finexercice


\exercice{6288, mayer, 2011/10/16}
\enonce{006288}{} 
Soient $E$ et $F$ des espaces de Banach et $f: E\to F$ une
application de classe $C^2$. \begin{enumerate}
\item Soit $h\in E$ et $\varphi _h : E\to F$ l'application définie
par $\varphi_h (x) = Df(x) (h)$. Justifier que
$$D^2 f(a) (k,h) = D\varphi _h(a) (k) \quad \text{pour tout} \;\; k\in E\;
.$$
\item Supposons que, pour tous $t\in \Rr$ et $x\in E$, $f(tx) =
t^2 f(x)$. Montrer que $D^2f(0) (x,x) =2f(x)$ pour tout $x\in E$.
\item Soit $a,h,k\in E$ et soit $\Psi : \Rr ^2 \to F$ définie
par $\Psi (t,s) = f(a+th+sk)$. Calculer $\frac{\partial ^2
\Psi}{\partial t \partial s}(0,0)$.
  \end{enumerate}
\finenonce{006288}



\finexercice


\exercice{6289, mayer, 2011/10/16}
\enonce{006289}{} 
Soit $f:\Rr^n\to\Rr^n$ de classe $C^2$ telle que, pour tout $x\in
\Rr^n$, l'application $Df(x)$ est un automorphisme orthogonal, i.e.
$Df(x)$ est linéaire bijective et conserve le produit scalaire:
$$\langle Df(x)(h), Df(x)(k)\rangle=\langle h,k\rangle \quad \text{pour tout} \;\; h,k\in
\Rr^n\; .$$ Montrer que l'application $f$ est elle même un
automorphisme orthogonal.

Indications: 
\begin{enumerate}
\item Déterminer la différentielle de $x\mapsto \langle Df(x)(h),
Df(x)(k)\rangle$.
\item Vérifier que $A(h,k,l) = \langle Df(x)(h), D^2f(x)(k,l)\rangle$ est
antisymétrique par rapport aux deux premières variables et
symétrique par rapport aux deux dernières variables.
\item En déduire que $A(h,k,l)=0$ pour tous $h,k,l\in \Rr^n$ puis conclure.
\end{enumerate}
\finenonce{006289}



\finexercice


\exercice{6290, mayer, 2011/10/16}
\enonce{006290}{} 
\begin{enumerate}
\item Trouver les applications $G: \Rr^2 \to \Rr$ de classe
$C^2$ telles que  $\frac{\partial ^2 G}{\partial x \partial y}=0$.
\item Trouver les applications $F:\Rr^2 \to \Rr$ de classe
$C^2$ solutions de
$$\frac{\partial ^2 F}{\partial x^2} - \frac{\partial ^2 F}{\partial
y^2}=0 \;\; .$$ (Indication: poser $\varphi (u,v) =
\frac{1}{2}(u+v,u-v)$ et $G=F\circ \varphi $ ). \end{enumerate}
\finenonce{006290}



\finexercice


\exercice{6291, mayer, 2011/10/16}
\enonce{006291}{} 
Soient $E,F,G$ des Banach et $u:E\to F$, $v:F\to G$ deux
applications $C^2$. Calculer, à l'aide de la définition, la
différentielle seconde de $w=v\circ u$.
\finenonce{006291}



\finexercice


\exercice{6301, mayer, 2011/10/16}
\enonce{006301}{} 
\begin{enumerate}
\item Soit $f:\Rr \to \Rr$ une application $C^\infty$ et soit
$n\geq 1$. \'Etablir l'équivalence des propriétés suivantes:
\begin{itemize}
\item $f(0)=f'(0)=...=f^{(n-1)}(0) =0$.
\item $f(x) = x^n g(x)$ avec $g\in C^\infty$.
\end{itemize}
\item Soit $\Omega$ un ouvert convexe de $\Rr^n$ contenant $0$ et
soit $f\in C^\infty (\Omega , \Rr )$. On suppose $f(0)=Df(0)=0$.
Montrer qu'il existe $g_{i,j} \in C^\infty (\Omega , \Rr )$
telles que $f(x) = \sum_{i,j=1}^n x_ix_j g_{i,j}(x)$.
 \end{enumerate}
\finenonce{006301}



\finexercice


\exercice{6302, mayer, 2011/10/16}
\enonce{006302}{} 
Déterminer approximativement la valeur de $1,05^{1,02}$ avec une
erreur d'au plus $\epsilon < 10^{-2}$ (Indication: Appliquer Taylor à
la fonction $f(x,y) = x^y$).
\finenonce{006302}



\finexercice


\exercice{6303, mayer, 2011/10/16}
\enonce{006303}{} 
Montrer que si $x = 1,32 \pm  10^{-2}$ et $y = 0,45
\pm  10^{-2}$, alors $f(x,y) = e^{-(x^2 + y^2)} = 0,14
\pm  10^{-2}$.
\finenonce{006303}



\finexercice     


\exercice{6304, mayer, 2011/10/16}
\enonce{006304}{} 
Ecrire le développement de Taylor-Young à
l'ordre 2 au voisinage de (0,0) pour la fonction $f(x,y) =
{e^x\over \cos y}$. En déduire la limite ${e^x - (1 + x)
\cos y\over (x^2 + y^2) \cos y}$ quand $(x,y)$ tend vers
(0,0).
\finenonce{006304}



\finexercice     


\exercice{6305, mayer, 2011/10/16}
\enonce{006305}{} 
Soit $f(x,y)$ une fonction de classe $C^2$ au
voisinage du cercle $x^2 + y^2 = 1$. On pose ${\partial
f\over \partial x} (1,0) = a$ et ${\partial^2 f\over
\partial y^2} (1,0) = b$. Pour tout nombre réel $\theta
$, soit $F(\theta ) = f(\cos \theta, \sin \theta )$.
Calculez $F''(0)$ en fonction de $a$ et $b$.
\finenonce{006305}



\finexercice

\section{372.00 Difféomorphisme, théorème d'inversion locale et des fonctions implicites}


\exercice{1858, maillot, 2001/09/01}

\enonce{001858}{}
Soit $f:\R^2\rightarrow \R$ la fonction d\'efinie par $$f(x,y)=
((x-2)^2+y^2-4)((x-1)^2+\frac{y^2}4 -1).$$
\begin{enumerate}
\item Tracer rapidement la courbe $C$ d'\'equation $f(x,y)=0$.
\item En quels points de $C$ la relation $f(x,y)=0$ permet-elle de
d\'efinir une fonction implicite de la forme $y=\phi(x)$~?
\end{enumerate}
\finenonce{001858}



\finexercice



\exercice{1859, maillot, 2001/09/01}

\enonce{001859}{}
Montrer que les relations propos\'ees d\'efinissent au voisinage du
couple $(a,b)$ indiqu\'e une fonction implicite $y=\phi(x)$.

Donner un d\'eveloppement limit\'e \`a l'ordre 3 de $\phi$ en $a$.

\begin{enumerate}
\item $f(x,y)=x^3+y^3-3xy-1=0 \qquad \qquad \qquad (a,b)=(0,1)$.
\item $f(x,y)=2 e^{x+y-1}+ \ln(x-y) -2x+y^3 \qquad \qquad (a,b)=(1,0)$.
\end{enumerate}
\finenonce{001859}



\finexercice



\exercice{1860, maillot, 2001/09/01}

\enonce{001860}{}
Montrer que la relation $$f(x,y,z)=x^3+y^3+z^3 -2z(x+y)-2x+y-2z-1=0$$
d\'efinit au voisinage de $(0,0,-1)$ une fonction implicite $z=
\phi(x,y)$. Donner un d\'eve\-lop\-pe\-ment limit\'e de $\phi$ \`a
l'ordre 2 en $(0,0)$.
\finenonce{001860}



\finexercice



\exercice{2527, queffelec, 2009/04/01}
\enonce{002527}{}
\begin{enumerate}
\item Soit $f$ une application de ${\Rr}$ dans ${\Rr}$,
d\'erivable en tout point de ${\Rr}$ et telle que, pour tout $x$
de ${\Rr}$, $f'(x)\neq 0$. Montrer que $f$ est un hom\'eomorphisme
de ${\Rr}$ sur $f({\Rr})$ et que $f^{-1}$ est diff\'erentiable en
tout point de $f({\Rr})$.

\item Soit $f$ d\'efinie par $f(x)=x+x^2\sin{\frac \pi x}$ si
$x\not=0$ et $f(0)=0$.

Montrer que $f'(0)$ existe et est $\not=0$, mais que $f$ n'est
inversible sur aucun voisinage de $0$. Expliquer.
\end{enumerate}
\finenonce{002527} 


\finexercice


\exercice{2528, queffelec, 2009/04/01}
\enonce{002528}{}
\begin{enumerate}
\item Montrer que l'application $\varphi: (r,\theta)\to
(x,y)=(r\cos\theta, r\sin\theta)$ est un $C^1$-diff\'eomorphisme
de l'ouvert $]0,\infty[\times]-\pi,\pi[$ sur le plan priv\'e de la
demi-droite $\Rr^-$. Si $f(x,y)=g(r,\theta)$ donner les formules
de passage entre les d\'eriv\'ees partielles de $f$ et celles de
$g$.

\item Soit $U$ le plan priv\'e de l'origine, et $f(x,y)=(x^2-y^2,
2xy)$.

Montrer que $f$ est un diff\'eomorphisme local au voisinage de
tout point de $U$ mais n'est pas un diff\'eomorphisme global.

\item Soit $g$ l'application de ${\Rr}^2$ dans ${\Rr}^2$ d\'efinie
par $g(x,y)= (x+y,xy)$. Trouver un ouvert connexe maximal
$U\subset\Rr^2$ tel que $g$ soit un diff\'eomorphisme de $U$ sur
$g(U)$.

\item Soit $h$ l'application de ${\Rr}^2$ dans ${\Rr}^2$ d\'efinie
par $(x,y)\to (e^x\cos y,e^x\sin y)$.

 Montrer que $h$ est de classe $C^1$ dans ${\Rr}^2$;
que $h'(x,y)$ est un \'el\'ement de Isom(${\Rr}^2,{\Rr}^2$) pour
tout $(x,y)$ de ${\Rr}^2$; mais que $h$ n'est pas un
hom\'eomorphisme de ${\Rr}^2$ sur $h({\Rr}^2)$.
\end{enumerate}
\finenonce{002528} 


\finexercice


\exercice{2529, tahani, 2009/04/01}
\enonce{002529}{}
Soit $\varphi$
l'application de $\mathbb{R}^2$ dans $\mathbb{R}^2$ d\'efinie pas
$$\varphi(x,y)=(\sin(y/2)-x,\sin(x/2)-y).$$
\begin{enumerate}
\item Justifier que $\varphi$ est de classe $C^1$, calculer sa
diff\'erentielle et voir que $D\varphi(x,y)$ est inversible pour
tout $(x,y) \in \mathbb{R}^2$. \item Montrer que $\varphi$ est un
$C^1$-diff\'eomorphisme de $\mathbb{R}^2$ sur
$\varphi(\mathbb{R}^2)$ et justifier que $\varphi(\mathbb{R}^2)$
est un ouvert. 
\item Montrer que $\varphi^{-1}$ est lipschitzienne
(on prendra comme norme sur $\mathbb{R}^2$: $\|(x,y)\|=|x|+|y|$).
\item En d\'eduire que $\varphi$ est un diff\'eomorphisme de
$\mathbb{R}^2$ sur $\mathbb{R}^2$ 
\item Calculer
$D\varphi^{-1}(p)$ o\`u $p=(1-\pi/2, \sqrt{2}/2-\pi)$.
\end{enumerate}
\finenonce{002529} 


\finexercice


\exercice{2530, tahani, 2009/04/01}
\enonce{002530}{}
Soit $f:
\mathbb{R}^n \rightarrow \mathbb{R}^n$ une application $C^1$. On
suppose qu'il existe $\alpha >0$ tel que pour tout $h,x \in
\mathbb{R}^n$,
\[
\langle Df(x)(h),h \rangle \geq \alpha\langle h,h\rangle.
\]

\begin{enumerate}

\item En consid\'erant la fonction 
$t \rightarrow \varphi(t)= \langle f(a+t(b-a)),b_a) \rangle$, montrez que
$$ \langle f(b)-f(a),b-a \rangle \geq \alpha \langle b-a,b-a \rangle 
\text{ pour tout } a,b \in \mathbb{R}^n.$$ 
En déduire que $f$ est une application ferm\'ee.

\item D\'emontrer que, pour tout $x \in E, Df(x)$ est un
isomorphisme de $\mathbb{R}^n$. En déduire que $f$ est une
application ouverte. 

\item Conclure que $f$ est un diff\'eomorphisme de classe $C^1$ de $\mathbb{R}^n$ sur lui même.
\end{enumerate}
\finenonce{002530} 


\finexercice

\exercice{2531, queffelec, 2009/04/01}
\label{gijsexodiff} 
\enonce{002531}{}
Soit $U$ l'ouvert $\Rr^3 \setminus \{0\}$.
Soit $(x,y,z) \to (X,Y,Z)$ l'application inversion de p\^ole $0$, de
puisssance 1, d\'efinie dans $U$, \`a valeurs dans $\Rr^3$, par
les formules
$$X = {\frac x{x^2 + y^2 +
z^2}}\quad ; \quad Y = {\frac y {x^2 + y^2 + z^2}} \quad ; \quad Z =
{\frac z {x^2 + y^2 + z^2}}
$$
Calculer la matrice jacobienne de cette transformation (on posera
$\rho = \sqrt{x^2 + y^2 + z^2}$) et v\'erifier que cette matrice
est \'egale \`a son inverse. 
\finenonce{002531} 


\finexercice


\exercice{2532, queffelec, 2009/04/01}
\enonce{002532}{}
Reconsid\'erez
l'exercice \ref{gijsexodiff} dans l'esprit suivant : ``si $f$ est
un diff\'eomorphisme, la matrice inverse de la matrice jacobienne
de $f$ est la matrice jacobienne de $f^{-1}$.'' 
\finenonce{002532} 


\finexercice


\exercice{2533, queffelec, 2009/04/01}
\enonce{002533}{}
 Soit $E=M_n({\Rr})$ et $I$ la matrice unit\'e dans $E$. En consid\`erant
$\varphi:E\to E$ telle que $\varphi(A)=A^2$, montrer qu'il existe
$\alpha>0$ tel que toute matrice $A$ v\'erifiant $||A-I||<\alpha$
admette une racine carr\'ee. 
\finenonce{002533} 


\finexercice


\exercice{2534, queffelec, 2009/04/01}
\enonce{002534}{}
\begin{enumerate}
\item Montrer que si $a,b$ sont voisins de $1$, on peut trouver
$x,y\in{\Rr}$ tels que $y+e^{xy}=a,\ x+e^{-xy}=b$.

\item Soit $f$ l'application de ${\Rr}^2$ dans lui-m\^eme
d\'efinie par $f(x,y)=(x\sin(xy)+y, y\cos(xy)+x)$, et soit
$(a_n,b_n)$ une suite tendant vers $(0,0)$. Montrer que si
$f(a_n,b_n)=0$ pour tout $n$, la suite $(a_n,b_n)$ stationne.
\end{enumerate}
\finenonce{002534} 


\finexercice


\exercice{2535, queffelec, 2009/04/01}
\enonce{002535}{}
Soit $U$ un ouvert de
$\Rr^2$ et $\varphi : U\to \Rr^2$ une application de classe $C^1$
$\varphi=(f, g)$. On consid\`ere $u,v$ r\'eels et on cherche $x,y$
tels que
$$(*)\quad f(x,y)=u,\ g(x,y)=v.$$
\begin{enumerate}
\item On suppose que la diff\'erentielle de $\varphi$ est de rang
$2$ en tout point de $U$. Montrer que pour tout $(u,v)$ le
syst\`eme $(*)$ admet une solution, unique localement. Que peut-on
dire si la diff\'erentielle est de rang $2$ en un point de $U$
seulement ?

\item A-t-on des solutions si la diff\'erentielle est de rang $0$
?

\item On suppose maintenant que la diff\'erentielle de $\varphi$
est de rang $1$ en tout point de $U$. Si $f'_x$ ne s'annule pas
sur $U$, montrer que $\psi:(x,y)\to (f(x,y),y)$ d\'efinit un
diff\'eomorphisme d'un ouvert $V\subset U$ sur $\psi(V)$. En
d\'eduire $G$ telle que $g(x,y)=G(f(x,y))$ sur $V$. Que peut-on
dire des solutions du syst\`eme $(*)$ ?
\end{enumerate}
\finenonce{002535} 


\finexercice


\exercice{2536, queffelec, 2009/04/01}
\enonce{002536}{}
Soit $E=\Rr^n$ muni d'une
norme quelconque, et $B_r$ la boule ferm\'ee $||x||\leq r$. Soit
$f$ un $C^1$-diff\'eomorphisme entre deux ouverts $U$ et $V$ de
$E$, contenant $0$, tel que $f(0)=0$. On pose $A=f'(0)\in{\cal
L}(E)$. Soit $0<\varepsilon<1$.

\begin{enumerate}
\item Montrer qu'il existe $R>0$ tel que pour tout $x\in B_R$,
$$||A^{-1}(f(x))-x||\leq \varepsilon\ ||x||.$$

\item Montrer qu'il existe $R'>0$ tel que pour $0\leq r\leq R'$,
$$(1-\varepsilon)\ A(B_r)\subset f(B_r)\subset (1+\varepsilon)\ A(B_r).$$

\item En d\'eduire que $\lim_{r\to 0}\frac{\hbox{vol}\ f (B_r)}{
\hbox{vol}\  (B_r)}= |\det A|$.
\end{enumerate}
\finenonce{002536} 


\finexercice


\exercice{2537, queffelec, 2009/04/01}
\enonce{002537}{}
Soit $a,b\in{\Rr}$ et
$f:{\Rr}^2\to{\Rr}^2$ d\'efinie par $f(x,y)=(x+a\sin y, y+b\sin
x)$.
\begin{enumerate}
\item Montrer que si $|ab|<1$, $f$ est un diff\'eomorphisme de
${\Rr}^2$ sur lui-m\^eme.

\item Montrer que si $|ab|=1$, $f$ n'est plus un diff\'eomorphisme
mais reste un hom\'eomorphisme de ${\Rr}^2$ sur lui-m\^eme.
\end{enumerate}
\finenonce{002537} 


\finexercice


\exercice{2538, queffelec, 2009/04/01}
\enonce{002538}{}
Soit $f:{\Rr}^n\to
{\Rr}^n$ une application de classe $C^1$ telle que
$$||f(x)-f(y)||\geq k||x-y||$$
pour tous $x,h\in {\Rr}^n$, $k$ \'etant une constante $>0$. On va
montrer que $f$ est un $C^1$-diff\'eomorphisme de ${\Rr}^n$ sur
lui-m\^eme.
\begin{enumerate}
\item Montrer que $f$ est injective et que $f({\Rr}^n)$ est
ferm\'ee dans ${\Rr}^n$.

\item Montrer que $f'(x)$ est inversible pour tout $x\in{\Rr}^n$.

\item  En d\'eduire que $f({\Rr}^n)$ est un ouvert-ferm\'e de
${\Rr}^n$.
\end{enumerate}
\finenonce{002538} 


\finexercice

\exercice{2539, mayer, 2009/04/01}
\enonce{002539}{}
Soit $G$ un ouvert born\'e de
$\Rr^n$ et soit $f:\overline{G} \to \Rr ^n$ une application
continue dans $\overline{G}$ et $C^1$ dans $G$. Pour tout $x\in
G$,  on suppose $Df(x)$ inversible. D\'emontrer que, sous ces
conditions, l'application $x\mapsto \|f(x)\|$ atteint son maximum
en un point du bord $\partial G = \overline{G}\setminus G$.
\finenonce{002539} 


\finexercice


\exercice{2540, mayer, 2009/04/01}
\enonce{002540}{}
Soit $E$ un espace vectoriel
de dimension finie, $\Omega$ un ouvert connexe de $E$ et soit
$f:\Omega \to E$ une application de classe $C^1$ telle que $\| |
Df(x)\|| \leq c$, pour tout $x\in \Omega$, o\`u $0\leq c <1$.
Montrer que $Id_{|E} -f$ est un diff\'eomorphisme $C^1$ de $\Omega
$ sur son image. 
\finenonce{002540} 


\finexercice

\exercice{2541, tahani, 2009/04/01}
\enonce{002541}{}
Soit $f:\mathbb{R}^3 \rightarrow \mathbb{R}^2$ d\'efinie par
$f(x,y,z)=(x^2-y^2+z^2-1, xyz-1)$. Soit $(x_0,y_0,z_0) \in
\mathbb{R}^3$ tel que $f(x_0,y_0,z_0)=(0,0)$. Montrez qu'il existe
un intervalle $I$ contenant $x_0$ et une application $\varphi: I
\rightarrow \mathbb{R}^2$ tels que $\varphi(x_0)=(y_0,z_0)$ et
$f(x,\varphi(x))=0$ pour tout $x\in I$.
\finenonce{002541} 


\finexercice


\exercice{2542, tahani, 2009/04/01}
\enonce{002542}{}
Soit $F:\mathbb{R}^2
\rightarrow \mathbb{R}$ l'application $F(x,y)=x^2+y^2-1$.
D\'emontrer que, pour $x$ suffisamment proche de $0$, il existe un
nuique $y=y(x)>0$ tel que $F(x,y)=0$. V\'erifier, sans
r\'esolution explicite, que $y'(x)=-x/y$. 
\finenonce{002542} 


\finexercice


\exercice{2543, tahani, 2009/04/01}
\enonce{002543}{}
On consid\`ere le syst\`eme
d'\'equations:
$$\left( \begin{array}{c}
x^2+y^2-2z^2=0 \\ x^2+2y^2+z^2=4
\end{array} \right) $$
Montrer que, pour $x$ proche de l'origine, il existe des fonctions
positives $y(x)$ et $z(x)$ telles que $(x,y(x),z(x))$ soit
solution du syst\`eme. On d\'eterminera $y'$ en fonctionde $x,y$
et $z'$ en fonction de $x,z$.
\finenonce{002543} 


\finexercice


\exercice{2544, tahani, 2009/04/01}
\enonce{002544}{}
Consid\'erons
$F(x,y)=y^n+a_{n-1}(x)y^{n-1}+...+a_1(x)y+a_0(x)$ un polyn\^ome
\`a coefficients variables. On suppose :
\begin{enumerate}
\item Les fonctions $x \rightarrow a_j(x)$ sont $C^1$,
$j=0,1,...,n-1$. 
\item pour un certain $x_0 \in \mathbb{R}$, le
polyn\^ome $y \rightarrow F(x_0,y)$ a un z\'ero simple $y_0 \in
\mathbb{R}$.
\end{enumerate}
D\'emontrer que, dans ces conditions, $F(x,y)$ poss\`ede, pour $x$
voisin de $x_0$, un z\'ero $y(x)$ qui lui est proche de $y_0$ et
que la d\'ependance $x \rightarrow y(x)$ est $C^1$.
\finenonce{002544} 


\finexercice


\exercice{2545, tahani, 2009/04/01}
\enonce{002545}{}
Donner l'allure de $C=\{(x,y)
\in \mathbb{R}^2; x^4+y^3-y^2+x-y=0\}$ au voisinage des points
$(0,0)$ et $(1,1)$.
\finenonce{002545} 


\finexercice


\exercice{2546, tahani, 2009/04/01}
\enonce{002546}{}
Montrer que l'\'equation
$e^x+e^y+x+y-2=0$ d\'efinit, au voisinage de l'origine, une
fonction implicite $\varphi$ de $x$ dont on calculera le
d\'eveloppement limit\'e d'ordre trois en $0$. 
\finenonce{002546} 


\finexercice

\exercice{6264, queffelec, 2011/10/16}
\enonce{006264}{} 
Soit $i = \sqrt{-1}$. Calculer la matrice jacobienne
de l'application $(x,y) \to (X,Y)$ de $\Rr^2$ dans
$\Rr^2$ définie par $X + iY = (x + iy)^3.$
\finenonce{006264}



\finexercice     


\exercice{6265, queffelec, 2011/10/16}
\enonce{006265}{}
Démontrer le résultat suivant (théorème d'inversion globale) :

Soit $E$, $F$ deux Banach, $U$ un ouvert de $E$ et $f:U\to F$ une application
de classe $C^1$ sur $U$. Alors $f$ est un $C^1$-difféomorphisme
 de $U$ sur $f(U)$ si et seulement si :

(i) $f$ est injective;

(ii) $f'(x)\in$ Isom($E,F$) pour tout $x\in U$.
\finenonce{006265}



\finexercice


\exercice{6266, queffelec, 2011/10/16}
\enonce{006266}{}
\begin{enumerate}
\item On considère l'application $\varphi$ de $\Rr^3$ dans lui-même définie par
$(x,y,z)\to (e^{2y}+e^{2z}, e^{2x}-e^{2z}, x-y)$. Montrer que $\varphi$ est un
$C^1$-difféomorphisme de $\Rr^3$ sur son image que l'on précisera.

\item Soit $\lambda\in \Rr$ et $F$ l'application de $\Rr^3$ dans lui-même
définie par
$(x,y,z)\to (e^{x-y+2z}+e^{-x+y+2z}, e^{2x}+e^{2y}-2\lambda e^{x-y},
e^{2x}+e^{2y}-2e^{-x+y})$. Montrer que
$F$ s'écrit $G\circ \varphi$, $G$ à préciser, et que c'est un
$C^1$-difféomorphisme de $\Rr^3$ sur son image si et seulement si
$\lambda\geq0$.
\end{enumerate}
 \finenonce{006266}



\finexercice


\exercice{6267, queffelec, 2011/10/16}
\enonce{006267}{}
On va proposer trois démonstrations possibles de l'exercice classique suivant :
soit $f$ une application de classe $C^1$ de ${\Rr}$ dans lui-même, telle que
$|f'(x)|\leq k$ pour tout $x$ réel, où $k\in ]0,1[$. Alors $F$ définie par
$F(x,y)=(x+f(y),y+f(x))$ est un difféomorphisme de classe $C^1$ de ${\Rr}^2$
{\bf sur} lui-même. 
\begin{enumerate}
\item Remarquer que $F$ est injective et $F'(x,y)\in$ Isom(${\Rr}^2,{\Rr}^2$)
pour tout $(x,y)$.

Reste à établir la surjection.

\item 1ère méthode : Montrer que $F$ est propre ($\lim ||F(x,y)||=+\infty$ quand
$||(x,y)||\to+\infty$) et que si $(a,b)\in{\Rr}^2$, la fonction
$g(x,y)=||F(x,y)-(a,b)||^2$ est différentiable et atteint sa borne inférieure
en un point annulant $g'(x,y)$; conclure.

\item 2ème méthode : Montrer que $F({\Rr}^2)$ est à la fois ouverte et fermée.
Conclure.

\item 3ème méthode : Si $(a,b)\in{\Rr}^2$, appliquer le théorème du point fixe à
l'application $\phi(x,y)=(a-f(y),b-f(x))$; conclure.
\end{enumerate}
\finenonce{006267}



\finexercice


\exercice{6268, queffelec, 2011/10/16}
\enonce{006268}{}
Soit $P$ un polyn\^ome de degré $3$ normé, de racines $x_1<x_2<x_3$ :
$$P(t,x_1,x_2,x_3)=\Pi_{l=1}^3(t-x_l)=t^3+\sum_{k=1}^3a_k t^{k-1}.$$
Les coefficients $a_k$ sont des fonctions polyn\^omiales, donc de classe
$C^1$, des racines. On pose $\Omega=\{x_1<x_2<x_3\}$ et on définit
$f:x\in\Omega\to (a_1,a_2,a_3)\in {\Rr}^3$. On va montrer que $f$ est un
$C^1$-difféomorphisme de $\Omega$ sur
$f(\Omega)$.
\begin{enumerate}
\item Vérifier que $f$ est injective sur $\Omega$.

\item On appelle $J$ la matrice jacobienne de $f$, et $V$ la matrice de coefficients
$\displaystyle v_{ij}=x_i^{j-1}$. En calculant $\displaystyle{\partial
P\over\partial x_k}(t,x_1,x_2,x_3)$ de deux fa\c cons, montrer que $VJ$ est une
matrice diagonale inversible si
$x\in\Omega$. Conclure.

\item En déduire la dérivée de $f^{-1}$ en tout point de $f(\Omega)$.
\end{enumerate}
\finenonce{006268}



\finexercice


\exercice{6269, queffelec, 2011/10/16}
\enonce{006269}{}
\begin{enumerate}
\item Soit $f$ la fonction de ${\Rr}^2$ dans ${\Rr}$ définie par
$f(x,y)=x^3+y^3-3xy$ et $C$ l'ensemble des $(x,y)\in\Rr^2$ tels que
$f(x,y)=0$.

En quels points $(a,b)$ peut-on appliquer le théorème des fonctions implicites
? Calculer la dérivée de la fonction implicite lorsqu'elle existe et
écrire l'équation de la tangente à $C$.

\item Montrer que l'équation $e^x+e^y+x+y-2=0$ définit au voisinage de $0$ une
fonction implicite $\varphi$ de $x$ dont on calculera le développement limité à
l'ordre $3$ en $0$. 

\item Montrer que les équations $x+y-zt=0,\ xy-z+t=0$ définissent au voisinage de
$(0,1)$ deux fonctions implicites $x=\varphi_1(z,t),\ y=\varphi_2(z,t)$ avec $\varphi_1(0,1)=1$,
dont on calculera les différentielles en ce point.
\end{enumerate}
 \finenonce{006269}



\finexercice


\exercice{6270, queffelec, 2011/10/16}
\enonce{006270}{}
Soit $(u,v)\in{\Rr}^2\to F(u,v)\in{\Rr}$ une application de classe $C^1$, telle que $F(0,0)=0$ et
${\partial F\over\partial v}(0,0)\neq 0$. On considère $\varphi:{\Rr}^3\to {\Rr}^2$
telle que $\varphi(x,y,z)=(xy, x^2-y^2-z)$ et l'application $f=F\circ\varphi$.
Montrer que l'équation $f(x,y,z)=0$ définit au voisinage de $(0,0)$ une
applica\-tion $z=\psi(x,y)$ vérifiant 
$$x{\partial \psi\over\partial x}-y{\partial \psi\over\partial y}=2(x^2+y^2).$$
\finenonce{006270}



\finexercice


\exercice{6271, queffelec, 2011/10/16}
\enonce{006271}{novembre 1999}
Soit $O$ un ouvert de $\Rr^3$, $f:O\to \Rr$ une application de classe $C^1$
et $(a,b,c)\in O$ tel que $f(a,b,c)=0$.
\begin{enumerate}
\item Donner une condition suffisante pour qu'on puisse résoudre : et $x$ en
fonction de $(y,z)$ et $y$ en
fonction de $(x,z)$ et $z$ en
fonction de $(x,y)$. Plus précisément, donner une condition suffisante pour
qu'il existe $U$ voisinage de $a$, $V$ voisinage de $b$, $W$ voisinage de $c$
avec $U\times V\times W\subset O$, et $\varphi:V\times W\to U$,
$\chi:U\times W\to V$, $\psi:U\times V\to W$ des fonctions de classe $C^1$
tels que pour $(x,y,z)\in U\times V\times W$ :

$$f(x,y,z)=0\Longleftrightarrow x=\varphi(y,z)\Longleftrightarrow
y=\chi(x,z)\Longleftrightarrow z=\psi(x,y).$$

\item Si la condition donnée en 1. est satisfaite, démontrer que pour $(x,y,z)\in
U\times V\times W$ tel que $f(x,y,z)=0$ on a

$$\partial_1\varphi(y,z)
\partial_2\chi(x,z)\partial_1\psi(x,y)=-1.$$ 
\end{enumerate}
\finenonce{006271}



\finexercice


\exercice{6272, queffelec, 2011/10/16}
\enonce{006272}{}
On considère $E=M_n({\Rr})$, $F=GL(n,{\Rr})$ et l'application $\Psi$ de
$F\times E$ dans $E$ définie par $\Psi(A,B)=AB-I$. Montrer à l'aide du théorème
des fonctions implicites que $\varphi:A\in F\to A^{-1}$ est différentiable en
tout point de $F$ et retrouver sa différentielle.
\finenonce{006272}



\finexercice


\exercice{6273, queffelec, 2011/10/16}
\enonce{006273}{}
On considère le système d'équations d'inconnues $x$ et $y$ :
$$x={1\over2}\sin(x+y)+t-1,\quad y={1\over2}\cos(x-y)-t+{1\over2}.$$
  
\begin{enumerate}
\item Montrer que pour chaque $t_0\in\Rr$, il existe une unique solution
$(x_0,y_0)$, et que la fonction ainsi définie est continue..

\item Montrer en considérant la fonction
$F(x,y,t)=(x-{1\over2}\sin(x+y)+t-1,\quad y-{1\over2}\cos(x-y)-t+{1\over2})$,
 que le système admet une unique solution $x=x(t),\ y=y(t)$ constituée
de fonctions $C^\infty$.

\item Donner un développement limité à l'ordre $2$ de $x(t),y(t)$ au point
$(0,0)$.
\end{enumerate}
\finenonce{006273}



\finexercice


\exercice{6274, mayer, 2011/10/16}
\enonce{006274}{} 
Soit $f:\Rr^3 \to \Rr^2$ définie par
$f(x,y,z)=(x^2-y^2+z^2-1,xyz-1)$. Soit $(x_0,y_0,z_0) \in \Rr^3$
tel que $f(x_0,y_0,z_0)=(0,0)$. Montrer qu'il existe un intervalle
$I$ contenant $x_0$ et une application $\varphi : I \to \Rr^2$ tels
que $\varphi (x_0) = (y_0,z_0)$ et $f(x,\varphi (x))=0$ pour tout $x\in
I$.
\finenonce{006274}



\finexercice


\exercice{6275, mayer, 2011/10/16}
\enonce{006275}{} 
 Soit $F:\Rr^2\to \Rr$ l'application $F(x,y)= x^2+y^2-1$.
Démontrer que, pour $x$ suffisamment proche de $0$, il existe un
unique $y=y(x)>0$ tel que $F(x,y)=0$. Vérifier, sans
résolution explicite, que $y'(x) = -\frac{x}{y}$.
\finenonce{006275}



\finexercice


\exercice{6276, mayer, 2011/10/16}
\enonce{006276}{} 
On considère le système d'équations:
$$ \left\{
\begin{array}{ccc}
  x^2 + y^2 -2z^2 &=&0 \\
  x^2 +2 y^2 +z^2 & = & 4
\end{array}\right. \;\; .$$
Montrer que, pour $x$ proche de l'origine, il existe des fonctions
positives $y(x)$ et $z(x)$ telles que $(x,y(x),z(x))$ soit
solution du système. On déterminera $y'$ en fonction de $x,y$
et $z'$ en fonction de $x,z$.
\finenonce{006276}



\finexercice


\exercice{6277, mayer, 2011/10/16}
\enonce{006277}{} 
Considérons $F(x,y)= y^n + a_{n-1}(x)y^{n-1}+...+a_1(x)y+a_0(x)$
un polyn\^ome à coefficients variables. On suppose: 
\begin{enumerate}
\item les fonctions $x\mapsto a_j(x)$ sont $C^1$, $j=0,1,...,
n-1$,
\item pour un certain $x_0 \in \Rr$, le polyn\^ome $y\mapsto
F(x_0,y)$ a un zéro simple $y_0 \in \Rr$.
 \end{enumerate}
Démontrer que, dans ces conditions, $F(x,y)$ possède, pour $x$
voisin de $x_0$, un zéro $y(x)$ qui lui est proche de $y_0$ et
que la dépendance $x\mapsto y(x)$ est $C^1$.
\finenonce{006277}



\finexercice


\exercice{6278, mayer, 2011/10/16}
\enonce{006278}{} 
Donner l'allure de ${\cal C} =\{(x,y)\in \Rr^2 \; ; \;\; x^4+y^3
-x^2-y^2 +x-y=0\}$ au voisinage des points $(0,0) $ et $(1,1)$.
\finenonce{006278}



\finexercice


\exercice{6279, mayer, 2011/10/16}
\enonce{006279}{} 
Montrer que l'équation $e^x+e^y+x+y-2=0$ définit, au voisinage
de l'origine, une fonction implicite $\varphi$ de $x$ dont on
calculera le développement limité d'ordre trois en $0$.
\finenonce{006279}



\finexercice


\exercice{6792, gijs, 2011/10/16}
\enonce{006792}{}
\label{gijsexo2}
\begin{enumerate}
\item Donner la définition d'une sous-variété $M$ de
$\Rr^n$ de dimension $k$.

\item Enoncer le théorème des fonctions implicites.

\smallskip
 Soit $f : \Rr^n \to \Rr^p$ une fonction
de classe $C^1$ et soit $M = \{\,x\in \Rr^n \mid f(x) =
\mathbf{0}\,\}$. Supposons en plus que $f$ vérifie la
condition~:
$\forall x\in M : \operatorname{rang}(J(x)) = p\ ,$
où $J(x) = \Bigl(\frac{\partial f_i}{\partial
x_j}(x)\Bigr){}_{j=1,\dots, n}^{i=1,\dots, p}$ est la
matrice Jacobienne de $f$ en $x$ de taille $n\times p$.

\item Montrer, en utilisant le théorème des
fonctions implicites, que $M$ est une sous-variété de
$\Rr^n$ de dimension $k = n-p$.
\end{enumerate}
\finenonce{006792}



\finexercice

\exercice{6820, gijs, 2011/10/16}
\enonce{006820}{}
Soit $E=\Rr^n$ muni de la norme euclidienne et $F =
M(n,\Rr)$ l'ensemble des matrices $n\times n$ à
coefficients réels~; on identifie $F$ avec l'ensemble
d'applications linéaires de $E$ dans $E$. Soit
$\|\cdot\|_{op}$ la norme d'opérateurs sur $F$.
\begin{enumerate}
\item Démontrer que $\forall A,B \in F : \|AB\|_{op} \le
\|A\|_{op} \cdot \|B\|_{op}$.

\item Démontrer que pour $A\in F$, la série
$\sum_{n=0}^\infty \frac1{n!} A^n$ est une série
normalement convergente sur tout compact de $F$. (Nota
Bene~: $A^0 = Id \in F$ désigne la matrice identité.)
On note $\exp : F \to F$ l'application $A \mapsto
\sum_{n=0}^\infty\frac1{n!} A^n$. 

\item Calculer $D_{\mathbf{0}} \exp$, où $\mathbf{0}$ est
l'application nulle de $E$ dans $E$. Quelle est la taille
de la matrice $D_{\mathbf{0}} \exp$~?  Déduire qu'il existe
un voisinage $U$ de $\mathbf{0} \in F$ et un voisinage $V$ de
$Id = \exp(\mathbf{0})$ tels que $\exp$ est un
($C^1$-)difféomorphisme de $U$ sur $V$. Justifier votre
réponse~!
\end{enumerate}
\finenonce{006820}



\finexercice


\exercice{6821, gijs, 2011/10/16}
\enonce{006821}{}
Soit $f:O \subset \Rr^3 \to \Rr$ une fonction de classe
$C^1$ sur l'ouvert $O$ et $(x_o,y_o,z_o) \in O$ tel que
$f(x_o,y_o,z_o) = 0$. 
\begin{enumerate}
\item 
Donner une condition suffisante pour
qu'on puisse résoudre~: et $x$ en fonction de $(y,z)$,
et $y$ en fonction de $(x,z)$, et $z$ en fonction de
$(x,y)$. Plus précisément, donner une condition
suffisante pour qu'il existe $U$ un voisinage de $x_o$,
$V$ un voisinage de $y_o$, $W$ un voisinage de $z_o$, $U
\times V \times W \subset O$, et $\phi: V \times W \to
U$, $\chi : U \times W \to V$ et $\psi : U \times V \to W$
des fonctions de classe $C^1$ telles que $\forall (x,y,z)
\in U \times V \times W$~:
$$
f(x,y,z) = 0 \Longleftrightarrow x=\phi(y,z) 
\Longleftrightarrow y = \chi(x,z)
\Longleftrightarrow z = \psi(x,y)
\ .
$$
Justifier votre réponse.

\item
Si la condition donnée en 1. est satisfaite,
démontrer que $\forall (x,y,z) \in U \times V \times W$
tels que $f(x,y,z) = 0$ on a
$$
(\partial_1\phi)(y,z) \cdot (\partial_2\chi)(x,z) \cdot
(\partial_1\psi)(x,y) = -1
\ .
$$
(Remarque, cette relation est beaucoup utilisée en
thermodynamique, où elle est écrite comme $(\frac
xy)_z \cdot (\frac yz)_x \cdot (\frac zx)_y = -1$, avec
l'interprétation que $(\frac xy)_z$ est la dérivée
de $x$ par rapport à $y$ en gardant $z$ constant.)
\end{enumerate}
\finenonce{006821}



\finexercice


\exercice{6833, gijs, 2011/10/16}
\enonce{006833}{} 
Dans cet exercice, $\mathcal{O}$ désigne un ouvert de
$\Rr^n$.
\begin{enumerate}
\item Quand dit-on qu'une fonction $f:\mathcal{O} \to \Rr$
est de classe $C^k$ ($k\ge1$)~?

\item Démontrer par récurrence sur $k$ que si $f
: \mathcal{O} \to \Rr$ et $g: \mathcal{O} \to \Rr\setminus\{0\}$ sont de classe
$C^k$, alors  $h:\mathcal{O} \to \Rr$, $h(x) =
f(x)/g(x)$ est aussi de classe $C^k$.

\item Démontrer par récurrence sur $k$ que la
composée de deux fonctions de classe $C^k$ est aussi de
classe $C^k$.

\item Soit $Gl(2,\Rr)$ l'ensemble des matrices à
coefficients réels $2\times2$ inversibles.
Dé\-mon\-trer que l'applica\-tion $I: Gl(2,\Rr) \to
Gl(2,\Rr)$, $I(\left(\smallmatrix a&b \\ c&d
\endsmallmatrix\right)) = \left(\smallmatrix a&b \\ c&d
\endsmallmatrix\right)^{-1}$ est de classe $C^k$ pour
tout $k\ge1$.

\item \'Enoncer le théorème de l'inversion locale
pour une fonction $f:\mathcal{O} \to \Rr^n$.

\item On se restreint au cas $n=2$ et on se place
dans les conditions du théorème de l'inversion
locale. Démontrer que si $f$ est de classe $C^k$,
$k>1$, alors la réciproque $f^{-1}$ (donnée par le
théorème de l'inversion locale) est  de classe
$C^k$.
\end{enumerate}
\finenonce{006833}



\finexercice

\section{373.00 Extremum, extremum lié}


\exercice{1823, drutu, 2003/10/01}

\enonce{001823}{}
Soit $f$ la fonction d\'efinie sur $\R^2$ par $f(x,y)=x^2-xy^2$.
Montrer que $(0,0)$ est le seul  point critique de $f$,  qu'il
n'est pas un extremum local, mais que pourtant la restriction de
$f$ \`a toute droite passant par $(0,0)$ admet en ce point un
minimum local.
\finenonce{001823}



\finexercice



\exercice{1824, drutu, 2003/10/01}

\enonce{001824}{}
Ecriver la formule de Taylor de second ordre pour chacune des
fonctions suivantes au point $(x_0,y_0)$ donn\'e.
\begin{enumerate}
\item $f(x,y)=\sin (x+2y),\; (x_0,y_0)=(0,0)$ ;
\item $f(x,y)=\frac{1}{x^2+y^2+1},\; (x_0,y_0)=(0,0)$ ;
\item $f(x,y)=e^{-x^2-y^2}\cos xy,\; (x_0,y_0)=(0,0)$ ;
\item $f(x,y)=\sin (xy)+\cos (xy),\; (x_0,y_0)=(0,0)$ ;
\item $f(x,y)=e^{(x-1)^2}\cos y,\; (x_0,y_0)=(1,0)$.
\end{enumerate}
\finenonce{001824}



\finexercice



\exercice{1825, drutu, 2003/10/01}

\enonce{001825}{}
Pour chacune des fonctions suivantes etudiez la nature du point
critique donn\'e :
\begin{enumerate}
\item $f(x,y)=x^2-xy+y^2$ au point critique $(0,0)$ ;
\item $f(x,y)=x^2+2xy+y^2+6$ au point critique $(0,0)$ ;
\item $f(x,y,z)=x^2+y^2+2z^2+xyz$ au point critique $(0,0,0)$ ;
\item $f(x,y)=x^3+2xy^2-y^4+x^2+3xy+y^2+10$ au point critique
$(0,0)$.
\end{enumerate}
\finenonce{001825}



\finexercice



\exercice{1826, drutu, 2003/10/01}

\enonce{001826}{}
Trouvez les points critiques des fonctions suivantes et
d\'eterminez si ce sont des minima locaux, des maxima locaux ou
des points selle.
\begin{enumerate}
\item $f(x,y)=x^3+6x^2+3y^2-12xy+9x$ ;
\item $f(x,y)=\sin x+y^2-2y+1$ ;
\item $f(x,y,z)=\cos 2x\cdot \sin y+z^2$ ;
\item  $f(x,y,z)=(x+y+z)^2$.
\end{enumerate}
\finenonce{001826}



\finexercice



\exercice{1827, maillot, 2001/09/01}

\enonce{001827}{}
Soit $f:\R^2\rightarrow \R$ la fonction d\'efinie par $f(x,y)=x^3-3x(1+y^2)$.
\begin{enumerate}
\item \'Etudier les extremums locaux de $f$.
\item Soit $D=\{(x,y)\in\R^2 \mid x^2+y^2 \le 1\}$. Montrer que $f$ a un
maximum $M$ et un minimum $m$ sur $D$.
\item Soit $(x,y)\in D$. Montrer que si $f(x,y)=M$ ou $f(x,y)=m$, alors
$x^2+y^2=1$.
\item \'Etudier la fonction $t\mapsto f(\cos t, \sin t)$. En d\'eduire les
valeurs de $M$ et $m$.
\end{enumerate}
\finenonce{001827}



\finexercice



\exercice{1828, ridde, 1999/11/01}

\enonce{001828}{}
Trouver le point du plan $ (2x-y + z = 16)$ le plus proche de l'origine.
\finenonce{001828}



\finexercice



\exercice{1829, ridde, 1999/11/01}

\enonce{001829}{}
D\'eterminer les extremums de $f (x, y) = xy (1-x^{2}-y{2})$ sur $[0, 1]^{2}$.
\finenonce{001829}



\finexercice



\exercice{1830, ridde, 1999/11/01}

\enonce{001830}{}
Soit $f (x, y) = x^{2} + xy + y^{2}-3x-6y$. Montrer que $f$ admet au plus un extremum.
Ecrire $f (x, y) + 9$ comme la somme de deux carr\'es et en d\'eduire que $f$ admet $-9$
comme valeur minimale.
\finenonce{001830}



\finexercice



\exercice{1831, ridde, 1999/11/01}

\enonce{001831}{}
D\'eterminer un triangle d'aire maximale inscrit dans un cercle donn\'e.
\finenonce{001831}



\finexercice



\exercice{1832, ridde, 1999/11/01}

\enonce{001832}{}
Soit $f (x, y) = (x^{2}-y) (3x^{2}-y)$.

 Montrer que $f$ admet un minimum local en $0$ suivant tout vecteur de $\Rr^{2}$
 mais n'admet pas de minimum local en $ (0, 0)$.
\finenonce{001832}



\finexercice



\exercice{1833, ridde, 1999/11/01}

\enonce{001833}{}
Soit ${f}: {\Rr^{2}} \to {\Rr}$, ${ (x, y)}\mapsto {xe^{y} + ye^{x}}$.\\
 Montrer que $ (-1, -1)$ est le seul extremum possible. A l'aide d'un
 d\'eveloppement limit\'e de $\varphi (h) = f (-1 + h, -1 + h)$ et de
 $\psi (h) = f (-1 + h, -1-h)$, montrer que $f$ n'a pas d'extremum.
\finenonce{001833}



\finexercice



\exercice{1834, gourio, 2001/09/01}

\enonce{001834}{}
D\'{e}terminer les extr\'{e}mums de $ f:(x,y,z)\rightarrow
x^{2}+y^{2}+z^{2}+2xyz.$
\finenonce{001834}



\finexercice



\exercice{1835, gourio, 2001/09/01}

\enonce{001835}{}
D\'{e}terminer $ \max\limits_{\left| z\right| \leq 1}\left| \sin z\right| .$
On rappelle que : $\sin z=\frac{e^{iz}-e^{-iz}}{2i}. $
\finenonce{001835}



\finexercice



\exercice{1836, gourio, 2001/09/01}

\enonce{001836}{}
Si $ f$ est concave sur un ouvert convexe $ U\subset {\Rr}^{2}$ et si :
$$\exists  a\in U,\frac{\partial f}{\partial x_{1}}(a)=\frac{\partial f}{%
\partial x_{2}}(a)=0, $$
alors $ f$ admet un maximum local en $ a.$
\finenonce{001836}



\finexercice



\exercice{1837, gourio, 2001/09/01}

\enonce{001837}{}
Soit $ A\subset {\Rr}^{2}, $ on d\'{e}finit $\int(A)$ comme
l'ensemble $ \{x\in A |\exists \rho >0,B(x,\rho )\subset A\}. $ On
supposera $ A$ ferm\'{e}e born\'{e}e et $ \int(A)\neq \emptyset .$
On suppose que $ f$ est une fonction $ C^{1 }$ sur $ A$ telle que
$ f$ est constante sur $ A\setminus \mathrm{Int}(A).  $ Montrer qu'il
existe $ z\in \mathrm{Int}(A)$ tel que :
$$\frac{\partial f}{\partial x_{1}}(z)=\frac{\partial f}{\partial x_{2}}(z)=0. $$
\finenonce{001837}



\finexercice



\exercice{1838, gourio, 2001/09/01}

\enonce{001838}{}
Chercher les extr\'{e}mums sur $ {\Rr}^{2}$ des applications :
$$(x,y)\rightarrow x^{4}+y^{4}-4xy; $$
$$(x,y)\rightarrow (x-y)e^{xy}; $$
$$(x,y)\rightarrow xe^{y}+ye^{x}; $$
$$(x,y)\rightarrow e^{x\sin y}; $$
$$(x,y)\rightarrow x^{3}+y^{3}. $$
\finenonce{001838}



\finexercice



\exercice{1839, maillot, 2001/09/01}

\enonce{001839}{}
Soit $f$ une fonction  r\'eelle de classe ${\cal C}^2$ sur un ouvert
$\Omega$ de $\R^2$.
  \begin{enumerate}
   \item
 Rappeler une condition n\'ecessaire pour que $f$ pr\'esente un
   extremum local en $(x_0,y_0)$.

Dans la suite de l'exercice, ${\bf a}=(x_0, y_0)$ v\'erifie cette
   condition, c'est-\`a-dire est un {\em point critique} de $f$. On
   pose
$$A = \frac{\partial^2 f}{\partial x^2}({\bf a}),\quad
B = \frac{\partial^2 f}{\partial x \partial y}({\bf a}),\quad
C = \frac{\partial^2 f}{\partial y^2}({\bf a}),$$
$$Q(x,y)=A x^2+2Bxy+Cy^2, \quad \Delta=B^2-AC,$$
$$R(t)=At^2+2Bt+C,\quad S(t)=Ct^2+2Bt+A.$$
  \item On suppose $\Delta<0$ et $A (\mbox{ou }C) >0$.
  \begin{enumerate}
     \item Montrer que $\forall t\in \R,$ $R(t)\geq \delta$ et $S(t)\geq
\delta$ pour un certain $\delta>0$.

       \item On pose $x=r \cos \theta$, $y=r \sin \theta$, avec $\displaystyle
r=\sqrt{x^2+y^2}$, et on suppose que $\sin \theta. \cos \theta \neq
0$.
Montrer successivement~:
$$\begin {array}{l}
\displaystyle Q(x,y) \geq r^2 \delta \sin^2 \theta,\\
Q(x,y) \geq r^2 \delta \cos^2 \theta,\\
Q(x,y) \geq \frac{r^2}{2} \delta.\end{array}$$
En d\'eduire que
$$\forall (x,y) \quad Q(x,y) \geq \frac{r^2}{2} \mbox{Inf}(\delta, 2A,
2C).$$

  \item  Montrer que ${\bf a}$ est un point de minimum local strict de
$f$. On \'ecrira pour cela la formule de Taylor-Young pour $f$ en ce point.
  \end{enumerate}
\item On suppose $\Delta<0$ et  $A (\mbox{ou }C) <0$.

Montrer que $(x_0,y_0)$ est un point de maximum local strict de $f$.

\item On suppose maintenant $\Delta>0$.
  \begin{enumerate}
     \item
      Montrer qu'il existe $t_1, t_2\in \R$ tels que $S(t_1)>0$
et $S(t_2)<0$.
     \item Soient $\theta_1, \theta_2 \in \R$ tels que
$\mbox{tan}\theta_1=
t_1$ et $\mbox{tan}\theta_2=t_2$. En examinant les fonctions
$$g(t)~:=f(x_0+t \cos \theta_1, y_0+t \sin \theta_1), \quad
h(t)~:=f(x_0+t \cos \theta_2, y_0+t \sin \theta_2)$$
pour $t\in \R$ assez petit, montrer que ${\bf a}$ n'est ni un point
de maximum local, ni un point de minimum local de $f$.
      \end{enumerate}
\item  Dessiner l'allure du graphe de $f$ au voisinage du point
$({\bf a},f({\bf a}))$
dans les trois cas \'etudi\'es ci-dessus (questions 1, 3 et 4).
 \item  Que peut-on dire en g\'en\'eral quand $\Delta=0$~?
Pour r\'epondre \`a cette question, on pourra s'appuyer sur l'\'etude
des deux
cas suivant au voisinage de $(0,0)$~:
$$f_1(x,y)=x^2+x^4+y^4 \quad \mbox{et} \quad f_2(x,y)=x^2-y^4.$$
\end{enumerate}
\finenonce{001839}



\finexercice



\exercice{1840, gourio, 2001/09/01}

\enonce{001840}{}
Existe-t-il un triangle d'aire maximale inscrit dans un cercle
donn\'{e}? Le d\'{e}terminer par une m\'{e}thode
g\'{e}om\'{e}trique.
\finenonce{001840}



\finexercice



\exercice{1841, gourio, 2001/09/01}

\enonce{001841}{}
Soit $ f:{\Rr}^{2}\rightarrow {\Rr}$ continue telle que :
$$\lim\limits_{\left\| x\right\| \rightarrow \infty }\left| f(x)\right|
=+\infty . $$
Montrer que $ f$ est minor\'{e}e et atteint sa borne inf\'{e}rieure.
\finenonce{001841}



\finexercice



\exercice{1842, legall, 2003/10/01}

\enonce{001842}{}
 Soit $f : \Rr 
^2\rightarrow \Rr $ l'application $(x,y) \mapsto 6xy+(y-x)^3.$ On note
$\Delta =\{ (x,y)\in \Rr^2 , -1\leq x\leq y \leq 1\} .$
\begin{enumerate}
\item Dessiner $\Delta $. Montrer que $f$ est born\'ee et atteint ses 
bornes sur $\Delta .$
\item Calculer les extrema de $f $ sur le bord de $\Delta $ puis dans 
l'int\'erieur de $\Delta $.
\item En d\'eduire les bornes de $f$ sur $\Delta $.

\end{enumerate}
\finenonce{001842}



\finexercice



\exercice{1843, legall, 2003/10/01}

\enonce{001843}{}
$S=\{ z\in \Cc ; \vert z\vert = 1\} .$ Soit $f $ l'application
de $D$ dans $\Rr$ d\'efinie par $f(z)=\vert \sin z\vert .$
\begin{enumerate}
\item Pour quelle raison $f$ est-elle born\'ee sur $D$~? On note 
$\displaystyle{ M=\sup_{z\in D}f(z)}$ et $\displaystyle{m=\inf_{z\in 
D}f(z)}.$
Est-ce que $M$ et $m$ sont atteints~? Donner la valeur de $m$.

\item Soit $z= x+iy\in \Cc , x, y \in \Rr.$ Montrer que $\vert \sin z\vert ^2=
\frac{1}{2}(\ch 2y-\cos 2x).$ (On rappelle que $\sin z 
=\frac{e^{i(x+iy)}-e^{-i(x+iy)}}{2i}$ et $\ch 
y=\frac{e^{y}-e^{-y}}{2}.$)

\item En d\'eduire que $M$ est atteint en un point de $S$.

\item Montrer que $\displaystyle{M=\frac{e^2-1}{2e}}$.
\end{enumerate}
\finenonce{001843}



\finexercice



\exercice{1844, maillot, 2001/09/01}

\enonce{001844}{}
On pose $\Omega=\R^2\setminus \{(0,0)\}$.

Soit $f:\R^2\rightarrow\R$ la fonction d\'efinie par
\[f(x,y)=\begin{cases} xy\frac{x^2-y^2}{x^2+y^2} &\text{si\ } (x,y)\in\Omega\\
               0 &\text{si\ } (x,y)= (0,0).
         \end{cases}\]

\begin{enumerate}
\item Montrer que $f$ est diff\'erentiable sur $\Omega$ et calculer sa
diff\'erentielle.
\item Montrer que $f$ est diff\'erentiable en $(0,0)$ et que sa
diff\'erentielle est nulle.
\item Montrer que $f$ admet en tout point des d\'eriv\'ees partielles secondes
$\frac{\partial^2 f}{\partial x\partial y}$ et
$\frac{\partial^2 f}{\partial y\partial x}$ et calculer la valeur de ces
 d\'eriv\'ees en $(0,0)$. Que peut-on en d\'eduire pour la continuit\'e de ces
d\'eriv\'ees partielles en $(0,0)$~?
\end{enumerate}
\finenonce{001844}



\finexercice



\exercice{6306, mayer, 2011/10/16}
\enonce{006306}{} 
Déterminer les extremums (locaux et/ou globaux) de: 
\begin{enumerate}
\item $f(x,y) = x^2 -y^2$, $(x,y) \in \Rr^2$.
\item $f(x,y) = x^3 -y^3$, $(x,y) \in \Rr^2$.
\item $f(x,y) = x^3 +y^3-3xy$, $(x,y) \in \Rr^2$.
\item $f(x,y) = x^2 +y^2-2xy+1$, $(x,y) \in \Rr^2$.
\end{enumerate}
\finenonce{006306}



\finexercice


\exercice{6307, mayer, 2011/10/16}
\enonce{006307}{} 
Discuter, suivant les valeurs du paramètre réel $\lambda $, la
nature des extremums de la fonction $f(x,y) = y(x^2+y^2 -2\lambda y)$.
\finenonce{006307}



\finexercice


\exercice{6308, mayer, 2011/10/16}
\enonce{006308}{} 
Soit $f: (t,x,y) \in \Rr^3 \mapsto f(t,x,y) \in \Rr $ une
application de classe $C^2$ telle que
\begin{itemize}
\item $\frac{\partial ^2 f}{\partial (x,y)^2}$ soit une matrice
définie positive en tout point et
\item $(x,y) \mapsto f(0,x,y)$ atteint son minimum en $(x_0,y_0)$.
\end{itemize}
Montrer que, si $t$ est voisin de $0$, l'application $(x,y)
\mapsto f(t,x,y)$ atteint son minimum en $(x(t),y(t))$, où
$t\mapsto (x(t),y(t))$ est une application de classe $C^1$ sur ce
voisinage de $0$.
\finenonce{006308}



\finexercice


\exercice{6309, mayer, 2011/10/16}
\enonce{006309}{} 
Soit $g(x,y,z)=xyz -32$, ${\cal S} =\{(x,y,z)\in \Rr^3 \; ;\;
g(x,y,z) =0 \}$ et soit $f: \Rr^3 \to \Rr$ l'application
$f(x,y,z) = xy+2yz + 2xz$. Déterminer $\min \{f(x,y,z) \; ;
(x,y,z) \in S\}$.
\finenonce{006309}



\finexercice


\exercice{6310, mayer, 2011/10/16}
\enonce{006310}{} 
Déterminer le point $p$ du plan $\Sigma = \{(x,y,x+y) \; ; \;
x,y\in \Rr\}$ qui réalise la distance $\mathrm{dist} \left( \Sigma \; ,
\; (1,0,0) \right)$.
\finenonce{006310}



\finexercice


\exercice{6311, mayer, 2011/10/16}
\enonce{006311}{} 
\begin{enumerate}
\item Déterminer les extremums de la fonction $f(x,y) =xy$ sur le
cercle unité $S= \{(,x,y) \in \Rr^2 \; ; x^2+y^2=1\}$.
\item Même question pour la fonction $f(x,y) = xy^2$.
\end{enumerate}
\finenonce{006311}



\finexercice


\exercice{6312, mayer, 2011/10/16}
\enonce{006312}{} 
Déterminer le minimum et maximum de la fonction $f(x,y,z) =
5x+y-3z$ sur l'intersection du plan $\Sigma = \{x+y+z=0\}$ avec la
sphère unité de $\Rr^3$.
\finenonce{006312}



\finexercice


\exercice{6313, mayer, 2011/10/16}
\enonce{006313}{} 
Déterminer les extremums de la fonction $f(x,y,z) = 2x+3y+2z$ sur
l'intersection du plan d'équation $x+z=1$ avec le cylindre
${\cal Z}= \{x^2+y^2 = 2\}\subset \Rr^3$.
\finenonce{006313}



\finexercice


\exercice{6883, romon, 2012/08/30}
\enonce{006883}{}
\begin{enumerate}
  \item Soit $E$ un espace vectoriel de dimension finie, $F$ un fermé non 
  vide de $E$ et $x \in E$ un point, appartenant ou non à $F$.
  Montrer qu'il existe un point $\bar{x} \in F$ tel que
  \[ 
  \| x - \bar{x} \| = \inf_{z \in F} \| x - z \| . 
  \]
  (Question subsidiaire: ce point est-il unique?)
  
  \item On se place désormais dans l'espace vectoriel $\mathcal{M}_2 (
  \mathbb{R})$ des matrices $2 \times 2$ à coefficients réels,
  muni de la norme
  \[ \left\| \left(\begin{array}{cc}
       a & c\\
       b & d
     \end{array} \right) \right\| = \sqrt{a^2 + b^2 + c^2 + d^2} \]
  et on considère l'ensemble $\text{SL}_2 (\mathbb{R})$
  des matrices $\left(\begin{array}{cc}
    a & c\\
    b & d
  \end{array} \right)$ de déterminant égal à $1$.
  \begin{enumerate}
    \item Montrer que $\text{SL}_2(\mathbb{R})$ est fermé dans $\mathcal{M}_2 (\mathbb{R})$.
    
    \item Montrer que $\text{SL}_2(\mathbb{R})$ n'est \emph{pas} bornée.
    
    \item Soit $f : \mathcal{M}_2 (\mathbb{R}) \rightarrow
    \mathbb{R}$ la fonction qui à une matrice $M$ associe $f (M
) = \| M \|$. On cherche la ou les matrices $M$ réalisant l'infimum de la fonction $f$
    sur $\text{SL}_2(\mathbb{R})$, autrement dit la ou les matrices les plus proches de la matrice nulle.
    \begin{enumerate}
      \item Montrer que $f_{|\text{SL}_2(\mathbb{R})}$ est minorée et atteint son infimum.
      
      \item Calculer le gradient de $f$. Est-il toujours défini?
      
      \item Trouver le ou les extrema de $f_{|\text{SL}_2(\mathbb{R})}$. Montrer qu'il s'agit du
      minimum et en déduire
      \[
      \inf_{M \in \text{SL}_2 (\mathbb{R})}  \| M \|.
      \]
    \end{enumerate}
  \end{enumerate}
\end{enumerate}
\finenonce{006883}


\finexercice

\section{374.00 Autre}


\exercice{6292, queffelec, 2011/10/16}
\enonce{006292}{} 
Une fonction $f:U \subset \Rr^n \to \Rr$ est
dite \emph{harmonique} si $\sum_{i=1}^n \frac{\partial^2 f}{\partial x_i^2}(x) =
0$ pour tout $x\in U$.
Une fonction $f(x,y)$ est dite \emph{radiale} si ses
valeurs au point $(x,y)$ ne dépendent que de la distance
$r = \sqrt{x^2 + y^2}$ à l'origine, c'est à dire si
$f(x,y) = F(r) = F(\sqrt{x^2 + y^2})$, où $F = F(r)$ est
une fonction d'une seule variable. 

Montrez que les seules
fonctions radiales et harmoniques, dans $\Rr^2$ privé
de l'origine, sont les fonctions $C \ln(r) + D = C \ln(
\sqrt{x^2 + y^2}) + D$, où $C$ et $D$ sont des
constantes.
\finenonce{006292}



\finexercice     


\exercice{6293, queffelec, 2011/10/16}
\enonce{006293}{} 
Vérifiez que les fonctions suivantes sont
harmoniques dans $\Rr^2$ :
\item{}
\begin{enumerate}  
\item $e^x \cos y$~; 
\item $x^3 - 3xy^2$~; 
\item pour tout entier
$k\geq 0$, la fonction $f(x,y) = r^k \cos (k\theta)$, où
$r$ et $\theta $ sont les coordonnées polaires de
$(x,y)$.
\end{enumerate}
\finenonce{006293}



\finexercice     


\exercice{6294, queffelec, 2011/10/16}
\enonce{006294}{} 
Exprimez en coordonnées polaires :
$\displaystyle y^2 {\partial^n f\over \partial x^2} - 2xy
{\partial^n f\over \partial x  \partial y} + x^2 
{\partial^n f\over \partial y^2} - x {\partial f\over
\partial x} - y {\partial f\over \partial x}.$
\finenonce{006294}



\finexercice     


\exercice{6295, queffelec, 2011/10/16}
\enonce{006295}{} 
Soit $U$ l'ouvert $\Rr^3$ privé de l'axe des
$z$.
\begin{enumerate}
\item Vérifiez que la fonction $f(x,y,z)$, qui vaut $e^z
\cos {\theta\over 2}{\sin r\over \sqrt r}$ en
coordonnées cylindriques, est harmonique sur $U$.

\item Soit $\lambda $ une constante réelle. Montrer qu'une
fonction du type $\displaystyle f(x,y,z) = e^z \cos
(\lambda \theta) u(r)$ est harmonique dans $U$ si et
seulement si $u = u(r)$ est solution de l'équation
différen\-tielle (dite de \emph{Bessel}) :
$$r^2 u''(r) + r u'(r) + [r^2 - 
\lambda^2] u (r) = 0.\eqno (E_\lambda)$$
\item
Vérifiez, que pour $\lambda = 3/2$, la fonction
$\displaystyle u(r) = {\sin r - r \cos r\over r \sqrt r}$
convient.
\end{enumerate}
\finenonce{006295}



\finexercice     


\exercice{6296, queffelec, 2011/10/16}
\enonce{006296}{} 
Dans $\Rr^3$ privé de l'origine, montrez que
les seules fonctions harmoniques et \emph{radiales}
(c'est-à-dire ne dépendant que de la distance $\rho$ de
$(x,y,z)$ à l'origine) sont les fonctions $f(x,y,z) =
{C\over \rho} + D = {C\over \sqrt {x^2 + y^2 + z^2}} + D$,
où $C$ et $D$ sont des constantes.
\finenonce{006296}



\finexercice     


\exercice{6297, queffelec, 2011/10/16}
\enonce{006297}{} 
Soient $\rho, \theta, \varphi$ les coordonnées
sphériques dans $\Rr^3$. On pose $\sin \varphi = t$.
Montrer que, pour qu'une fonction de la forme $f(x,y,z) =
\rho^n P(t)$, où $n$ est un entier $\geq 0$, soit
harmonique, il faut et il suffit que la fonction $t \mapsto
P(t)$ soit solution de l'équation différentielle (dite
\emph{de Legendre}) : $$(1 - t^2) P''(t) - 2t
P'(t) + n(n + 1) P(t) = 0.\eqno (D_n)$$
Pour $n = 0,1,2,3,4,5$, vérifiez, en le calculant par la
méthode des coefficients indéter\-minés, qu'il y a un
polyn\^ome $P_n(t)$, et un seul, de degré $n$, solution
de $(D_n)$, et tel que $P_n(1) = 1$. [Remarque : ce fait
vaut pour tout $n$ ; les polyn\^omes $P_n$ s'appellent
polyn\^omes de Legendre].
\finenonce{006297}



\finexercice     


\exercice{6298, queffelec, 2011/10/16}
\enonce{006298}{} 
Dans $\Rr^n$, on pose $\Delta f =
{\partial^f\over \partial x^2_1} + \cdots + 
{\partial^f\over \partial x^2_n}$, et $\rho = 
 \sqrt{x^2_1 +\cdots + x^2_n}$.
Soit une fonction \emph{radiale} $f(x_1, x_2, \ldots ,x_n)
= F(\rho).$
Montrer que $\Delta f = F''(\rho) + (n-1) F'(\rho)$. Si $n
\geq 3$, en déduire que les seules fonctions radiales et
harmoniques dans $\Rr^n$ privé de l'origine sont les
$f(x,y,z) = {C\over \rho ^{n-2}} + D$, où $C$ et $D$
sont des constantes.
\finenonce{006298}



\finexercice     


\exercice{6299, queffelec, 2011/10/16}
\enonce{006299}{} 
$\Delta(fg) = f\Delta g + g\Delta f + \langle
\nabla f \mid \nabla g \rangle$.
\finenonce{006299}



\finexercice     


\exercice{6300, queffelec, 2011/10/16}
\enonce{006300}{} 
Une fonction $f$ de classe $C^4$ (par exemple à 2
variables) est dite \emph{biharmonique} si $$\Delta (\Delta f)
= {\partial^4 f\over \partial x^4} + 2 {\partial^4 f\over
\partial x^2 \partial y^2} + {\partial^4 f\over \partial
y^4} \equiv 0.$$
Ces fonctions interviennent en théorie de
l'Elasticité. Bien entendu toute fonction harmonique est
biharmonique. Montrez que, si $f$ et $g$ sont deux
fonctions harmoniques, alors la fonction $xf + (x^2 +
y^2)g$ est biharmonique.
\finenonce{006300}



\finexercice  

\section{380.00 Solution maximale}


\exercice{2557, tahani, 2009/04/01}
\enonce{002557}{}
\begin{enumerate}
\item Pour chacune des équations suivantes o\`u $y=y(x)$ est
r\'elle de variable r\'eelle, d\'ecrire les solutions en
pr\'ecisant leur intervalle maximal de d\'efinition et dessiner
les trajectoires:
$$\mbox{ (i) } y'=e-y \mbox{ (ii) } y'-y=e^x \mbox{ (iii) }
xy'-2y=0.$$ 
\item Quelles sont les courbes isoclines de
l'\'equation $y'=y^2-x$; en d\'eduire l'allure des trajectoires.
\end{enumerate}
\finenonce{002557} 


\finexercice


\exercice{2558, tahani, 2009/04/01}
\enonce{002558}{} 
On consid\`ere l'\'equation
$$x'=3x^{2/3} \mbox{                          :} (1)$$
avec condition initiale $x(0)=0$.
\begin{enumerate}
\item Soit $\varphi$ une solution de $(1)$ d\'efinie sur
$\mathbb{R}$ telle que $\varphi(0)=0$; on pose $\lambda=\inf \{ t
\leq 0; \varphi(t)=0\} \leq +\infty$. Montrez que $\varphi$ est
identiquement nulle sur $(\lambda,\mu)$. 
\item Montrer que
$\varphi$ vaut $(t-\lambda)^3$ si $t\leq \lambda$, $0$ sur
$[\lambda, \mu]$ et $(t-\mu)^3$ si $t \geq \mu$; en d\'eduire
toutes les solutions maximales de $(1)$ d\'efinies sur
$\mathbb{R}$ avec $x(0)=0$.
\end{enumerate}
\finenonce{002558} 


\finexercice


\exercice{2559, tahani, 2009/04/01}
\enonce{002559}{}
On considère l'\'equation diff\'erentielle $x'=|x|+|t|$. 
\begin {enumerate} 
\item Montrez
que pour tout r\'eel $x_0$, il existe une solution maximale
$(\varphi,J)$ telle que $\varphi(0)=x_0$. 
\item D\'etermner la
solution maximale correspondant à $x_0=1$, en distinguant les cas
$t \geq 0$ et $t<0$, et v\'erifiez qu'elle est d\'efinie sur
$\mathbb{R}$ tout entier. Combien de fois est-elle d\'erivable?
\end{enumerate}
\finenonce{002559} 


\finexercice


\exercice{2560, tahani, 2009/04/01}
\enonce{002560}{} 
Soit $f:\mathbb{R}^2
\rightarrow \mathbb{R}$ donn\'ee par
$f(t,x)=4\frac{t^3x}{t^4+x^2}$ si $(t,x)\neq (0,0)$ et $f(0,0)=0$.
On s'interesse \`a l'\'equation diff\'erentielle
$$x'(t)=f(t,x(t)).$$
\begin{enumerate}
\item L'application $f$, est-elle continue ? est-elle localement
lipschitzienne par rapport à sa seconde variable ? Que peut-on en
d\'eduire pour l'\'equation $(2)$ ? 
\item Soit $\varphi$ une
solution de $(2)$ qui est d\'efinie sur un intervalle $I$ ne
contenant pas $0$. On d\'efinit une application $psi$ par
$\varphi(t)=t^2\psi(t), t\in I$. D\'eterminer une \'equation
diff\'erentielle $(E)$ telle que $\psi$ soit solution de cette
\'equation, puis r\'esoudre cette \'equation $(E)$. 
\item Que
peut-on en d\'eduire pour l'existence et l'unicit\'e de
l'\'equation diff\'erentielle $(2)$ avec donn\'ee initiale
$(t_0,x_0)=(0,0)$
\end{enumerate}
\finenonce{002560} 


\finexercice


\exercice{2561, tahani, 2009/04/01}
\enonce{002561}{}
Soit l'\'equation
diff\'erentielle $$x'''-xx''=0$$ o\`u $x$ est une application
trois fois d\'erivable, d\'efinie sur un intervalle ouvert de
$\mathbb{R}$ et \`a valeurs dans $\mathbb{R}$.
\begin{enumerate}
\item Mettre cette \'equation diff\'erentielle sous la forme
canonique $y'(t)=f(t,y(t))$, o\`u $f$ est une application que l'on
d\'eterminera. 
\item Soient $t_0, a, b, c \in \mathbb{R}$. Montrer
qu'il existe une unique solution maximale $\varphi$ de
l'\'equation $(3)$ qui satisfasse aux conditions initiales
$$\varphi(t_0)=a, \varphi'(t_0) \mbox{ et }
\varphi''(t_0)=c.$$ 
\item Soit $\varphi$ une telle solution
maximale. Calculer la d\'eriv\'ee de la fonction $$t \rightarrow
\varphi''(t)exp\left(- \int_{t_0}^t \varphi(u)du \right)$$ En
d\'eduire que la fonction $\varphi$ est soit convexe, soit concave
sur son intervalle de d\'efinition. D\'eterminer $\varphi$ dans le
cas o\`u $\varphi''(t_0)=0$.
\end{enumerate}
\finenonce{002561} 


\finexercice


\exercice{2562, tahani, 2009/04/01}
\enonce{002562}{}
On consid\`ere l'\'equation
$xx''=(x')^2+1$ sur $\mathbb{R}$. 
\begin{enumerate} 
\item Montrer
que, $x_0\neq 0$ et $x_0'$ \'etant donn\'es dans $\mathbb{R}$, il
existe une unique solution $\varphi$ d\'efinie au voisinage de
$0$, telle que $\varphi(0)=x_0$ et $\varphi'(0)=x_0'$. 
\item Si de
plus $x_0' \neq 0$, on peut supposer que $\varphi$ est un
$C^1$-diff\'eomorphisme d'un voisinage de $0$ sur un voisinage de
$x_0$ (pourquoi ?); on note $\psi$ l'application r\'eciproque et
on pose $z(x)=\varphi'(\psi(x))$. Calculez $z'(x)$, trouver
l'\'equationb satisfaite par $z$ et expliciter $z$; en d\'eduire
une expression de $\varphi$. 
\item Quelle est la solution
$\varphi$ de l\'equation telle que $\varphi(0)=x_0\neq 0$ et
$\varphi'(0)=0$.
\end{enumerate} 
\finenonce{002562} 


\finexercice


\exercice{6314, queffelec, 2011/10/16}
\enonce{006314}{}
\begin{enumerate}
\item Pour chacune des équations suivantes où $y=y(x)$ est réelle de variable
réelle, décrire les solutions en précisant leur intervalle maximal de définition
et dessiner les trajectoires :
$$(i)\ y'=e^{-y};\quad (ii)\ y'-y=e^x ;\quad (iii)\ xy'-2y=0; \quad (iv)\ 
x^2y'-y=x^2-x.$$

\item Quelles sont les courbes isoclines de l'équation
$ y'=y^2-x$; en déduire l'allure des trajectoires.
\end{enumerate}
\finenonce{006314}



\finexercice


\exercice{6315, queffelec, 2011/10/16}
\enonce{006315}{}
 Soit $f,g$ deux fonctions réelles continues sur un intervalle $I$ de $\Rr$,
$x_0\in I$, et on suppose que $g^{-1}(0)$ est fini ou discret; montrer que la
solution de l'équation
$y'=f(x)g(y),\ y(x_0)=y_0$ s'obtient sous forme implicite, et préciser son
intervalle de définition.

Exemples : $f(x)=\displaystyle{1\over\sqrt{1-x^2}},\ g(y)=\sqrt{1-y^2}$;
$f(x)=1,\  g(y)=y-y^2$.
\finenonce{006315}



\finexercice


\exercice{6316, queffelec, 2011/10/16}
\enonce{006316}{}
On considère l'équation différentielle 
$$(E)\qquad y''+y'+y=0$$
\begin{enumerate}
\item Résoudre l'équation différentielle $(E)$ et étudier le comportement des
solu\-tions en
$+\infty$ et $-\infty$.

\item Soit $f$ une fonction de classe ${\cal C}^2$ sur
${\Rr}$ telle que
$f''+f'+f$ soit
$T$-périodique. Montrer que $f(x+T)-f(x)$ est solution de $(E)$.
En déduire que si $f$ est bornée sur ${\Rr}$, $f$ est elle-même
$T$-périodique.
\end{enumerate}
\finenonce{006316}



\finexercice


\exercice{6317, queffelec, 2011/10/16}
\enonce{006317}{}
On considère les deux équations différentielles du second ordre
$$({\cal E}_1)\quad\quad y''=\sin x\qquad({\cal E}_2)\quad\quad y''+ \omega^2y=
\sin x$$ où $\omega$ est un nombre réel de module strictement inférieur à
$1$.

\begin{enumerate}
\item Trouver la solution $y$ de l'équation ${\cal E}_1$ vérifiant
$y(0)=0,\   y(\pi)=4\pi.$


\item Décrire la solution générale de l'équation ${\cal E}_2$, et prouver
ainsi que la solution
$y_\omega$ vérifiant $y_\omega(0)=0,\   y_\omega(\pi)=4\pi$, a pour expression
$$y_\omega(x)=4\pi{{\sin\omega x}\over{\sin\pi\omega}}- {{\sin x}\over{1-\omega^2}}.$$ 


\item Trouver, à $x$ fixé, la limite de $y_\omega(x)$ quand $\omega$ tend
vers $0$. Interprétation.


\item On  restreint $x$ à parcourir l'inter\-valle
$[0,\pi]$, et on suppose $0\leq\omega\leq{1\over2}$. A l'aide
 de la formule de Taylor, 
%à la fonction $f(\omega)=\pi\sin\omega x-x\sin\pi\omega$, 
montrer que 
$\vert
\pi\sin\omega x-x\sin\pi\omega\vert\leq\pi^3\omega^3$.
 En déduire :  $\vert y_\omega(x)-y(x)\vert \leq A\omega^2$,
où $A$ est une constante.
\end{enumerate}
\finenonce{006317}



\finexercice


\exercice{6318, queffelec, 2011/10/16}
\enonce{006318}{}
On considère l'équation différentielle du second ordre
$$({\cal E})\quad\quad y''+ y=f(x)$$
où $f$ est une fonction continue sur ${\Rr}$ tout entier.
\begin{enumerate}
\item On suppose que $f$ est un polyn\^ome de degré $n$. 
On note $E$ l'espace vectoriel des polyn\^omes réels de degré $\leq n$. Montrer
que l'application
$u$, définie sur $E$ par $u(P) = P''+ P$, est une application injective de $E$
 dans lui-même. 

En déduire que l'équation ($\cal E$) a une et une seule solution polyn\^omiale $g$ qui est de même
degré que $f$.


\item On suppose maintenant que $f$ est une fonction continue quelconque sur ${\Rr}$. 
Montrer que la fonction
$$\int_0^x f(t) \sin(x-t) \ dt $$ est une solution particulière de ($\cal E$); en déduire la
solution générale de l'équation.


\item Montrer que si $f$ vérifie l'inégalité $f(t) \leq {1\over{1+t^2}}$ pour tout
$t\in {\Rr}$, toutes les solutions de ($\cal E$) sont des fonctions bornées.
A-t-on la même conclusion si la fonction $f$ est seulement bornée sur ${\Rr}$?  
\end{enumerate}
\finenonce{006318}



\finexercice


\exercice{6319, queffelec, 2011/10/16}
\enonce{006319}{Méthode de Picard}
 On construit de proche en proche
la suite de fonctions réelles $(y_n)$ par la relation
$$y_0(x)=1,\quad y_n(x)=1+\int_0^x (y_{n-1}(t))^2 dt.$$
\begin{enumerate}
\item On suppose d'abord $x\in [0,1[$. Montrer que la suite $(y_n)$ est croissante
 et majorée par
${1\over 1-x}$. En déduire que 
$y_n$ tend, quand $n\to\infty$, vers une limite qui est la solution
de l'équation différentielle $y'=y^2$ valant $1$ en $0$.

\item Si $-1<x\leq0$,  
montrer que $(y_{2n})$ et $(y_{2n+1})$ sont des suites adjacentes, telles que,
pour
$n\geq1$,
${1\over 1-x}\leq y_{2n}\leq1$ et $1+x\leq y_{2n+1}\leq {1\over 1-x}$; retrouver
la solution de l'équation sur
$]-1,1[$ valant $1$ en $0$. 
\end{enumerate}
\finenonce{006319}



\finexercice


\exercice{6320, queffelec, 2011/10/16}
\enonce{006320}{}
\begin{enumerate}
\item En suivant la méthode d'itération de Picard, trouver la solution des
équa\-tions avec condition initiale :

\quad (i) $x'(t)=ax(t)+b;\ x(0)=0.$

\quad (ii) $x'(t)=\sin x(t);\ x(0)=0.$

\item Soit $A$ une matrice $n\times n$ constante. Trouver par la méthode de Picard
la solution de $X'(t)=A\  X(t); X(0)=X_0$; retrouver ainsi la solution de

$x''(t)=- x(t);\ x(0)=0, x'(0)=1$.

\item Soit cette fois $A(t)$ une famille de matrices $n\times n$ de fonctions
continues, telle que pour $s,t$, on ait $A(s)A(t)=A(t)A(s)$. Trouver la
solution de l'équation $X'(t)=A\  X(t); X(0)=X_0$ (on montrera que
$B(s)B(t)=B(t)B(s)$ où
$B(t)=\int_0^tA(u)\ du$).
\end{enumerate}
\finenonce{006320}



\finexercice


\exercice{6321, queffelec, 2011/10/16}
\enonce{006321}{}
On considère l'équation $(1)\quad x'=3x^{2/3}$ avec condition initiale
$x(0)=0$.
\begin{enumerate}
\item Soit $\varphi$ une solution de $(1)$ définie sur
$\Rr$ telle que
$\varphi(0)=0$; on pose $\lambda=\inf\{t\leq0;\ \varphi(t)=0\}\geq -\infty$ et
$\mu=\sup\{t\geq0;\
\varphi(t)=0\}\leq +\infty$. 
Mon\-trer que $\varphi$ est identiquement nulle sur
$(\lambda,\mu)$.

\item Montrer que $\varphi$ vaut $(t-\lambda)^3$ si
$t\leq\lambda$, $0$ sur $[\lambda,\mu]$ et $(t-\mu)^3$ si
$t\geq\mu$; 
 en déduire toutes les solu\-tions maximales de $(1)$ définies sur $\Rr$ avec $x(0)=0$.
\end{enumerate}
\finenonce{006321}



\finexercice


\exercice{6322, queffelec, 2011/10/16}
\enonce{006322}{}
On considère l'équation différentielle $x'=|x|+|t|$;
\begin{enumerate}
\item Montrer que pour tout $x_0$ réel, il existe une solution maximale
$(\varphi,J)$ telle que $\varphi(0)=x_0$.

\item Déterminer la solution maximale correspondant à $x_0=1$, en distinguant les
cas $t\geq0$ et $t<0$, et vérifier qu'elle est définie sur $\Rr$ tout
entier. Combien de fois est-elle dérivable ? 
\end{enumerate}
\finenonce{006322}



\finexercice


\exercice{6323, queffelec, 2011/10/16}
\enonce{006323}{} 
On considère l'équation différentielle $x'= |x|+|t|$.
\begin{enumerate}
\item Montrer que pour tout réel $x_0$, il existe une solution
maximale $(\varphi , J)$ telle que $\varphi (0) = x_0$.
\item Déterminer la solution maximale correspondant  à $x_0
=1$, en distinguant les cas $t\geq 0$ et $t<0$, et vérifier
qu'elle est définie sur $\Rr$ tout entier. Combien de fois
est-elle dérivable?
 \end{enumerate}
\finenonce{006323}



\finexercice


\exercice{6324, queffelec, 2011/10/16}
\enonce{006324}{} 
Soit $f: \Rr ^2 \to \Rr$ donnée par $f(t,x)=
4\frac{t^3x}{t^4+x^2}$ si $(t,x) \neq (0,0)$ et $f(0,0)=0$. On
s'intéresse à l'équation différentielle
\begin{equation} \label{eq 1}
x'(t) = f(t,x(t)) \;\; .
\end{equation}
\begin{enumerate}
  \item L'application $f$, est-elle continue et/ou localement
lipschitzienne par rapport à sa seconde variable? Que peut-on en
déduire pour l'équation (\ref{eq 1})?
  \item Soit $\varphi $ une solution de (\ref{eq 1}) qui est définie sur un
intervalle $I$ ne contenant pas $0$. On définit une application
$\psi$  par $\varphi (t) = t^2 \psi (t)$, $t\in I$. Déterminer une
équation différentielle (\ref{eq 1}') telle que $\psi$ soit solution de
cette équation, puis résoudre cette équation (\ref{eq 1}').
  \item  Que peut-on en déduire pour l'existence et
  l'unicité des solutions de l'équation différentielle
  (\ref{eq 1}) avec donnée initiale $(t_0,x_0)) = (0,0)$?
\end{enumerate}
\finenonce{006324}



\finexercice


\exercice{6325, queffelec, 2011/10/16}
\enonce{006325}{} 
Soit l'équation différentielle
\begin{equation} \label{eq 10}
x''' - x x'' =0 \; \; .
\end{equation}
où $x$ est une application trois fois dérivable, définie sur
un intervalle ouvert de $\Rr$ et à valeurs dans $\Rr$.
\begin{enumerate}
  \item Mettre cette équation différentielle sous la forme
  canonique
  $$ y'(t) = f(t,y(t)) \; , $$
  où $f$ est une application que l'on déterminera.
  \item Soient $t_0, a,b,c\in \Rr$. Montrer qu'il existe une
  unique solution maximale $\varphi$ de l'équation (\ref{eq 10}) qui
  satisfasse aux conditions initiales
  $$ \varphi (t_0) =a \; , \; \varphi ' (t_0)= b \;\; et \;\; \varphi '' (t_0) =c\;
  . $$
  \item Soit $\varphi $ une telle solution maximale. Calculer la
  dérivée de la fonction
  $$t\mapsto \varphi ''(t) \exp\left( -\int _a ^t \varphi (u) \, du
  \right)\; .$$
  En déduire que la fonction $\varphi $ est soit convexe, soit
  concave sur son intervalle de définition. Déterminer $\varphi $
  dans le cas où $\varphi ''(a)=0$.
\end{enumerate}

\finenonce{006325}



\finexercice


\exercice{6809, gijs, 2011/10/16}
\enonce{006809}{}
{\it Introduction.\/}
Si $A$ est une matrice $n\times n$ à coefficients
réels, on sait que les solutions de l'équation 
\begin{equation}
\varphi'(t) = A\varphi(t)
\tag{*}
\end{equation}
sont définies (au moins) sur l'intervalle $[0,\infty[$
(et à valeurs dans $\Rr^n$). Dans la question 1.
on vous demande de démontrer que, sous certaines
hypothèses sur $A$, on a $ \lim_{t\to\infty}
\varphi(t) = 0$.

Le but des questions 2. et 3. est de
démontrer les mêmes résultats pour l'équation
perturbée
\begin{equation}
\varphi'(t) = A\varphi(t) + g(\varphi(t))
\tag {**}
\end{equation}
où $g:\Rr^n \to \Rr^n$ est une fonction de classe $C^1$
vérifiant $g(0) = 0$ et $g'(0) = 0$.

\medskip

Les questions 1. et 2. sont
indépendantes. La question 3. utilise des
résultats des questions 1. et 2.

\begin{enumerate}
\item
Soit $M(n,\Rr)$ l'ensemble des matrices carrées 
$n\times n$ à coefficients réels et $M(n,\Cc)$
l'ensemble des matrices  $n\times n$ à
coefficients complexes. Soit $\| \quad \| : \Cc^n \to
\Rr$ la norme définie par $\|(x_1, \ldots, x_n)\|
=  \max_{i} |x_i|$. 

  \begin{enumerate}
  \item
Démontrer
que pour tout $x\in\Cc^n$, $B\in M(n,\Cc)$ on a~:
$$  \max_{i=1,\ldots,n} \left|{\sum_{j=1}^n B_{ij}x_j}\right|
\le (\max_i \sum_j |B_{ij}|) \cdot (\max_j
|x_j|)\ ,$$ 
c'est-à-dire $\|Bx\| \le (\max_i
\sum_j |B_{ij}) \cdot \|x\|$.


  \item Soit $D\in M(n,\Cc)$ une matrice diagonale avec
$(\lambda_1, \ldots,\lambda_n)$ sur la diagonale. Montrer
que $\forall x\in \Cc^n$~: $\|Dx\| \le (\max_i 
|\lambda_i|) \cdot \|x\|$.


  \item
Soit $A\in M(n,\Cc)$ diagonalisable et  $\lambda_1,
\ldots,\lambda_n $ ses valeurs propres. Soit $\alpha =
\max_i \Re \lambda_i$. Montrer qu'il existe une constante
$K>0$ telle que $\forall x\in \Cc^n$, $\forall t\in
[0,\infty[$~: $\|e^{tA}x\| \le K e^{\alpha t}
\|x\|$.

  \item
Soit $A\in M(n,\Rr)$ diagonalisable sur $\Cc$ et
$\lambda_1,
\ldots,\lambda_n $ ses valeurs propres  (attention, ici
$A$ est  à coefficients réels). Soit en plus $\alpha =
\max_i \Re \lambda_i$  strictement inférieur à 0.
Déduire de ce qui précède que si $\varphi(t)$ est
solution de l'équation (*), alors $ \lim_{t\to
\infty} \varphi(t) = 0$.
  \end{enumerate}

\item
Soit $g:\Rr^n \to \Rr^n$ comme dans l'introduction.
  \begin{enumerate}
  \item
 Démontrer que $\forall \epsilon>0$, $\exists
\eta>0$, $\forall x,y\in \Rr^n$~: 
$\|x\| < \eta$,
$\|y\| < \eta$ $\Longrightarrow$ $\|g(x) - g(y)\|
\le \epsilon \cdot \|x-y\|$.


  \item
Quelle est la solution de domaine de définition
maximal (à préciser) de l'équa\-tion (**)
vérifiant la condition initiale $\varphi(0) = 0$~?
  \item
Soit $I\subset \Rr$ un intervalle contenant 0. Démontrer
qu'une application continue $\varphi:I \to \Rr^n$ est
solution de l'équation (**) avec condition initiale
$\varphi(0) = {x_0\in \Rr^n}$ si et seulement si on a sur
$I$~: $$
\varphi(t) = e^{tA} x_0 + \int_0^t e^{(t-s)A}
g(\varphi(s))\,ds
\ .
$$
Indication~: pour $\Longrightarrow$ on pourra considérer
la fonction $\psi(t) = e^{-tA} \varphi(t)$.
  \end{enumerate}

\item
Soit $g:\Rr^n \to \Rr^n$ comme dans l'introduction~;  soit
$A$ et $\alpha$ comme dans 1.(d) et $K$ comme dans
1.(c) Soit $\beta\in \Rr$ tel que
$\alpha < \beta < 0$. En posant $\epsilon =
(\beta-\alpha)/K$, la question 2.(a) nous donne un
$\eta>0$. Soit finalement $x_0\in \Rr^n$ tel que $\|x_0\| < \eta/K$.

On définit une suite de fonction $\varphi_p : [0, \infty[
\ \to \Rr^n$ par les formules ${\varphi_0(t) =
e^{tA} x_0}$ et 
$$
\varphi_{p+1}(t) = e^{tA} x_0 + \int_0^t e^{(t-s)A}
g(\varphi_p(s)) \, ds
\ .
$$
  \begin{enumerate}
  \item
Démontrer par récurrence que $\forall p\in \Nn$,
$\forall t\in [0,\infty[$~: $\|\varphi_p(t)\| \le K
e^{t\beta} \|x_0\|$.
  \item
Démontrer que $\chi(t) = \|\varphi_1(t) -
\varphi_0(t)\|$ est une fonction bornée sur
$[0,\infty[\,$.
  \item
Démontrer par récurrence que $\forall p\in \Nn$,
$\forall t\in [0,\infty[$~: 
$$
\|\varphi_{p+1}(t) - \varphi_p(t)\|\le \left(
\frac{\beta-\alpha}{-\alpha } \right)^p \cdot
\sup_{s\in [0,\infty[} \|\varphi_1(s) - \varphi_0(s)\|
\ .
$$
  \item
En déduire que $\forall p,q\in \Nn$, $\forall t\in
[0,\infty[$~: 
$$
\|\varphi_{p+q}(t) - \varphi_p(t)\| \le
\left(\frac{\beta-\alpha}{-\alpha } \right)^p \cdot
\left(1 - \frac{\beta-\alpha}{-\alpha } \right)^{-1}
\cdot \sup_{s\in [0,\infty[} \| \varphi_1(s) -
\varphi_0(s)\| \ .
$$
  \item
Démontrer que la suite $(\varphi_p)$ converge vers une
fonction $\varphi:[0,\infty[\ \to \Rr^n$ et que la
convergence est uniforme.
  \item
Démontrer que $\varphi$ est continue et solution de
l'équation (**) vérifiant $\varphi(0) = x_0$.
  \item
Démontrer que $\forall t\in [0,\infty[$~: $\|
\varphi(t)\| \le K e^{t\beta} \|x_0\|$.
  \item
Déduire de ce qui précède qu'il existe un voisinage
$U$ de $0\in \Rr^n$ (à préciser) tel que $\forall
x_0\in U$ il existe un intervalle $I\subset \Rr$ et une
fonction $\varphi:I \to \Rr^n$ vérifiant
    \begin{enumerate} 
    \item $\varphi$ est l'unique solution de (**) vérifiant $\varphi(0) = x_0$,
    \item $[0,\infty[\ \subset I$,
    \item $\lim_{t\to \infty} \varphi(t) = 0$.
    \end{enumerate}
  \end{enumerate}
\end{enumerate}
\finenonce{006809}



\finexercice


\exercice{6822, gijs, 2011/10/16}
\enonce{006822}{}
Pour une fonction réelle  $x$ de la variable $t$ on
considère le système 
\begin{equation}
x'''' + 2x'' + x = 0\qquad \& \qquad x(0) = 2\ .
\tag{*}
\end{equation}
\begin{enumerate}
\item Le système \thetag* a-t-il une solution
unique~? Pourquoi~?

\medskip
Si votre réponse à la question 1. est oui,
répondre à la question 2., si votre réponse à la
question 1. est non, répondre à la question 3..

\item Trouver la solution unique de \thetag*.
Est-elle périodique~?

\item Trouver toutes les solutions de \thetag*.
Y-a-t-il des solutions périodiques~? Si oui, les
expliciter.
\end{enumerate}
\finenonce{006822}



\finexercice


\exercice{6828, gijs, 2011/10/16}
\enonce{006828}{}
Pour une fonction {\bf réelle}  $x$ de la variable $t$ on
considère le système 
\begin{equation}
\begin{cases}
x^{(4)} - 2x^{(3)} + 2x'' - 2x' + x = 0 
\\ 
x(0) = 0 \quad,\quad x''(0) = -2 \quad ,\quad  x^{(3)}(0) =
0 
\end{cases}
\tag{*}
\end{equation}
\begin{enumerate}
\item Le système \thetag* a-t-il une solution
unique~? Pourquoi~?

\medskip
\noindent
Si votre réponse à la question 1. est oui,
répondre à la question 2., si votre réponse à la
question 1. est non, répondre à la question 3..

\item Trouver la solution unique de \thetag*.
Est-elle périodique~?

\item Trouver toutes les solutions de \thetag*.
Y-a-t-il des solutions périodiques~? Si oui, les
expliciter.
\end{enumerate}
\finenonce{006828}



\finexercice


\exercice{6837, gijs, 2011/10/16}
\enonce{006837}{}
Soit $f: \ ]0,+\infty[\ \times \ ]0,+\infty[ \ \times \Rr
\to \Rr$ la fonction $f(x,t,c) = x- \ln(x) - t + \ln(t) -
c$. 
\begin{enumerate}
\item
 Démonter que si une fonction dérivable de $t$,
$x=\gamma_c(t)$ est solution de l'équation
$f(\gamma_c(t),t,c)=0$, alors elle est solution de
l'équa\-tion différentielle ${t(x-1)x' = (t-1)x}$.

\item
 Existe-t-il une fonction $g(t,c)$ définie
dans un voisinage de $(2,0)$ vérifiant $g(2,0)=2$ et
$f(g(t,c),t,c)=0$ (justifier)~? Si oui, calculer le
développement limité (le polyn\^ome de Taylor) d'ordre
2 de cette solution.
\end{enumerate}
\finenonce{006837}



\finexercice


\exercice{6845, gijs, 2011/10/16}
\enonce{006845}{}
Pour une fonction {\bf réelle}  $x$ de la variable $t$ on
considère le système 
\begin{equation}
\begin{cases}
x^{(4)} - 2x^{(3)} + x'' + 2x' - 2x = 0 
\\ 
x(0) = 0 \qquad,\qquad x''(0) =
0 
\end{cases}
\tag{*}
\end{equation}
\begin{enumerate}
\item Le système \thetag* a-t-il une solution
unique~? Pourquoi~?

\noindent
Si votre réponse à la question 1. est oui,
répondre à la question 2., si votre réponse à la
question 1. est non, répondre à la question 3..

\item Trouver la solution unique de \thetag*.
Est-elle périodique~?

\item Trouver toutes les solutions de \thetag*.
Y-a-t-il des solutions périodiques~? Si oui, les
expliciter.
\end{enumerate}
\finenonce{006845}



\finexercice

\section{381.00 Théorème de Cauchy-Lipschitz}


\exercice{6326, queffelec, 2011/10/16}
\enonce{006326}{}
On considère  l'équation différentielle (de Ricatti) sur $\Rr$ :
$$({\cal E})\quad\quad x'(t)=a(t)x^2(t)+b(t)x(t)+c(t),$$
où $a,b,c\in C(\Rr,R)$. Soit $x_i, 1\leq i\leq4$, quatre solutions distinctes
définies sur $I$. On pose $B=\displaystyle{x_3-x_1\over x_3-x_2}
{x_4-x_2\over x_4-x_1}$.
\begin{enumerate}
\item Montrer que $B$ est bien défini sur $I$.

\item Montrer que $B$ est une fonction constante sur $I$ (utiliser la dérivée
logarith\-mique).
\end{enumerate}
\finenonce{006326}



\finexercice


\exercice{6327, queffelec, 2011/10/16}
\enonce{006327}{}
On considère l'équation différentielle (de Bernoulli) sur $\Rr$ :
$$({\cal E})\quad\quad y'+y+xy^2=0.$$
\begin{enumerate}
\item Recherche des solutions qui ne s'annulent jamais. Transformer l'équation
par le difféomorphisme $(x,y)\in{\Rr}\times{\Rr^*}\to (x,{1\over y})$ en une
équation
$({\cal E'})$ qu'on résoudra. En déduire une famille $(\varphi_\lambda)$ de
solutions de
$({\cal E})$ avec leur intervalle maximal de définition.

\item Montrer que par tout point $(x_0,y_0)$ du plan avec $y_0\not=0$, il passe
une solution $\varphi_\lambda$. En déduire toutes les solutions de $({\cal
E})$. 
\end{enumerate}
\finenonce{006327}



\finexercice


\exercice{6328, queffelec, 2011/10/16}
\enonce{006328}{}
Soit $f$ un champ de vecteurs de classe $C^1$ sur un ouvert $U$ de $\Rr^n$,
et,\break $ (1)\quad x'=f(x)$, l'équation associée.
\begin{enumerate}
\item Soit $x_0\in U$ tel que $f(x_0)=0$. Si $\varphi : J\to U$ est une solution de
$(1)$ telle que $\varphi(t_0)=x_0$ pour un $t_0\in J$, alors $\varphi(t)=x_0$
pour tout $t\in J$.

\item Si $f$ est bornée sur $U$ et $\varphi : J\to U$ est une solution de
$(1)$ où $J=]a,b[$, $b\in \Rr$, alors $\lim_{t\to b} \varphi(t)$ existe.


\item Soit $\varphi : J\to U$ une solution de $(1)$ où $J\supset ]0,+\infty[$, et
supposons en outre que $\lim_{t\to\infty}\varphi(t)=a\in U$. Montrer que
$f(a)=0$.
\end{enumerate}
\finenonce{006328}



\finexercice


\exercice{6329, queffelec, 2011/10/16}
\enonce{006329}{}
Soit $f$ un champ de vecteurs de classe $C^1$ de $\Rr^n$ dans $\Rr^n$, et
on suppose (pour simplifier) que, pour toute donnée initiale $x$ de $\Rr^n$, il
existe une unique solution passant par $x$ au temps $t=0$, définie sur $\Rr$
tout entier. On note $\phi(t, x)$ cette solution (ou
$\phi$ le flot du champ).

\begin{enumerate}
\item Montrer qu'on a $\phi(s,\phi(t,x))=\phi(s+t,x)$ pour tous $s,t\in
\Rr$.
et la vérifier sur l'équation  $x'=x^2,
x(0)=\alpha\geq0$ aprés avoir précisé le domaine de définition du
flot.

\item On fait $n=2$; décrire le flot lorsque $f(x,y)=(-x,y);\quad
(y,x);\quad (-y,x)$, et vérifier la relation précédente.
\end{enumerate}
\finenonce{006329}



\finexercice


\exercice{6330, queffelec, 2011/10/16}
\enonce{006330}{Fonctions implicites et Cauchy-Lipschitz}
Soit $f:\Rr^2\to R$, de classe $C^2$, et $(a,b)$ un point tel que $f(a,b)=0$,
et $f'_y(a,b)\not=0$. Montrer que si $y=\varphi(x)$ est la fonction implicite
associée à $f(x,y)=0$, $\varphi$ est solution de l'équation différentielle :
$$y'=-{f'_x(x,y)\over f'_y(x,y)};\ y(a)=b.$$
Réciproque.
\finenonce{006330}



\finexercice


\exercice{6331, queffelec, 2011/10/16}
\enonce{006331}{Inégalité de Gronwall}
Soit $u,v$ deux applications de $[0,\beta]$ dans $\Rr$ continues et
positives telles que 
$$u(t)\leq C+\int_0^t u(s)v(s)\ ds$$
où $C$ est une constante positive ou nulle.
Montrer que $u(t)\leq C e^{\int_0^t v(s)\ ds}$. 
\finenonce{006331}



\finexercice


\exercice{6332, queffelec, 2011/10/16}
\enonce{006332}{}
Soit $f$ une application
$K$-lipschitzienne sur un ouvert $U\subset \Rr^n$. On va démon\-trer que
le flot de solutions de $x'=f(x)$, supposé défini sur un
intervalle $[t_0,t_1]$, dépend continument de la condition initiale $
x(t_0)=x_0$.
\begin{enumerate}
\item Soit $x_1,\ x_2$ deux telles solutions; montrer que si $t\in[t_0,t_1]$, 
$$||x_1(t)-x_2(t)||\leq ||x_1(t_0)-x_2(t_0)||e^{K(t-t_0)}$$


\item En déduire le résultat et le vérifier sur l'exemple : $x'=x^2$ sur $\Rr$.
\end{enumerate}
\finenonce{006332}



\finexercice


\exercice{6333, queffelec, 2011/10/16}
\enonce{006333}{}
Soit $I$ un intervalle ouvert de $\Rr$, $E=\Rr^n$ , $f(t,x)$ une fonction
continue de $I\times E$ dans $E$ telle que  $||f(t,x_1)-f(t,x_2)||\leq k(t)\
||x_1-x_2||$, où $k$ est une fonction continue $\geq0$ définie sur $I$.
\begin{enumerate}
\item  On considère $J$ intervalle compact $\subset I$ et l'opérateur $T$ défini
sur $C(J,\Rr^n)$ par 
$$Tx(t)=x_0+\int_{t_0}^tf(s,x(s))\ ds;$$ montrer que pour $p$ assez grand, $T^p$
est contractante;  en déduire que l'équation
$x'=f(t,x)$ admet une unique solution définie sur
$J$ tout entier telle que
$x(t_0)=x_0$. 
 
\item Montrer que l'équation $x'=f(t,x)$ admet une unique solution telle que
$x(t_0)=x_0$, définie sur
$I$ tout entier (on pourra écrire $I$ comme union d'intervalles compacts). 

\item Exemples : Montrer que les solutions maximales des équations 

$y''=-\sin y,y(0)=a, y'(0)=b$ (qu'on mettra
sous forme canonique), et $x'=A(t).x,
x(0)=x_0$ où $A(t)\in {\cal L}({\Rr^n})$ est constituée de fonctions continues
sur
$\Rr$, sont définies sur
$\Rr$ tout entier. 
\end{enumerate}
\finenonce{006333}



\finexercice


\exercice{6334, queffelec, 2011/10/16}
\enonce{006334}{} 
On considère l'équation  $x x'' = (x')^2 +1$ sur $\Rr$.
\begin{enumerate}
\item Montrer que, $x_0 \neq 0 $ et $x_0'$ étant donnés dans
$\Rr$, il existe une unique solution $\varphi$ définie au
voisinage de $0$, telle que $\varphi (0) =x_0$ et $\varphi '(0) =x'_0$.
\item Si de plus $x'_0 \neq 0$, on peut supposer que $\varphi $ est
un $C^1$--difféomorphisme d'un voisinage  de $0$ sur un
voisinage de $x_0$ (pourquoi?); on note $\psi$ l'application
réciproque et on pose $z(x) = \varphi '( \psi (x))$. Calculer $z'(x)
$, trouver l'équation satisfaite par $z$ et expliciter $z$; en
déduire une expression de $\varphi$.
\item Quelle est la solution  $\varphi$ de l'équation telle que
$\varphi (0) = x_0 \neq 0$ et $\varphi ' (0) =0$.
 \end{enumerate}
\finenonce{006334}



\finexercice


\exercice{6335, queffelec, 2011/10/16}
\enonce{006335}{} 
\begin{enumerate}
\item
On cherche à résoudre le problème
$$
 x'= t^2 +t x\quad , \quad  x(0) = 0 \; .
 $$
  \'Ecrire l'équation intégrale associée et utiliser
les cylindres de sécurités pour justifier que le procédé
itératif de Picard donne une suite de fonctions $(x_n)$
convergent uniformément sur $[-1/2, 1/2]$ vers une solution du
problème. Partant de $x_0 \equiv 0$, déterminer ensuite cette
suite $(x_n)$ et la solution du problème donné.
\item Résoudre avec ce procédé itératif le problème
$$ x' = t x \quad , \quad x(0) =1 \; , $$
puis aussi
$$\begin{array}{ll}
  x'_1= x_2 x_3 \; ,& x_1 (0) = 0 \; , \\
  x'_2 = - x_1 x_3\; , & x_2(0) =1 \; ,\\
  x'_3 = 2\; , & x_3(0) =0\; ,
\end{array}$$
en commen\c{c}ant avec $x_0 (t) = (0,1,0)$.
 \end{enumerate}
\finenonce{006335}



\finexercice


\exercice{6336, queffelec, 2011/10/16}
\enonce{006336}{} 
Calculer les premiers termes de l'itération de
Picard avec les conditions initiales données. Si
possible trouver des solutions explicites, y compris
leurs domaines de définition.
\begin{enumerate}
\item $x' = x+2$; $x(0) = 2$.
\item $x' = x^{4/3}$; $x(0) = 0$.
\item $x' = x^{4/3}$; $x(0) = 1$.
\item $x' = 1/(2x)$; $x(1) = 1$.
\end{enumerate}
\finenonce{006336}



\finexercice    

\section{382.00 Système linéaire à  coefficients constants}


\exercice{6337, queffelec, 2011/10/16}
\enonce{006337}{}
On rappelle les différentes méthodes pour résoudre un système différentiel
linéaire à coefficients constants $X'(t)= A.X(t)$ sur $E$ de dimension finie:
\begin{enumerate}
\item On met $A$ sous forme triangulaire et on résout de proche en proche
le nou\-veau système obtenu par changement de base avant de revenir au système
initial.

\item On met $A$ sous forme de Dunford, $P^{-1}AP=D+N$, où $D$ semi-simple
et
$N$ nilpotente, qui commutent. On calcule ainsi $e^{tA}.X_0$, la solution
valant
$X_0$ au temps $t=0$.

\item On utilise le théorème de Cayley-Hamilton pour établir des
relations entre les puissances de $A$ et calculer ainsi $e^{tA}$.


\item (cf. Cartan) On décompose $E=\oplus_i E_i$ en sous-espaces
caractéris\-tiques, on calcule $e^{tA_i}$ où $A_i=A_{|E_i}$, puis
$X(t)=\sum_ie^{tA_i}v_i$, si $X_0=\sum_iv_i$.

\item On cherche une
base de solutions par identification sous la forme de polyn\^omes-exponentielles,
suivant le résultat du cours.
\end{enumerate}    

Résoudre les systèmes différentiels de matrice :

$$
\begin{pmatrix}1 & 1\\ 2 & 0\end{pmatrix};
\quad 
\begin{pmatrix}
0 & 1 & -12 \\  -1 & 2  &-20\\ 1
& 0 & -5 \end{pmatrix};\quad 
\begin{pmatrix}1 & 0 & -1&1 \\  0&1 & 1  &0\\ 0&0&1 & 0  
\\0&0&1 & 0   \end{pmatrix}
$$ 
\finenonce{006337}



\finexercice


\exercice{6338, queffelec, 2011/10/16}
\enonce{006338}{}
Soit $A$ un opérateur de $\Rr^n$ et $x'=Ax$ le système associé.
\begin{enumerate}
\item On suppose que $A$ laisse un sous-espace $E$ invariant; montrer que si
$\varphi$ est une solution de condition initiale $\varphi(t_0)\in E$ alors
$\varphi(t)\in E$ pour tout $t\in \Rr$.

\item  Que peut-on dire des solutions du
système si $A$ est  nilpotente;?

\item On suppose que $A$ a une valeur propre de partie réelle $<0$; montrer qu'il
existe au moins une solution $\varphi$ telle que
$\lim_{t\to+\infty}\varphi(t)=0$.

\item A quelles conditions le système n'a-t-il que des solutions bornées ?
\end{enumerate}
\finenonce{006338}



\finexercice


\exercice{6339, queffelec, 2011/10/16}
\enonce{006339}{}
Montrer que $\Big(I+{A\over n}\Big)^n$ converge vers $e^A$ quand $n\to\infty$
 en majorant la différence. Retrouver ainsi la
valeur de $\det e^A$.   
\finenonce{006339}



\finexercice


\exercice{6340, queffelec, 2011/10/16}
\enonce{006340}{}
Soit $E=\Cc^n$, $A\in{\cal L}(E,E)$ et le système différentiel $X'=AX$.
\begin{enumerate}
\item On suppose que $X=e^{r_1t} u_1 + e^{r_2t} u_2$ est solution, où $u_i\in E$
et $r_i\in \Cc$ distincts. Montrer que $e^{r_1t} u_1$ et $e^{r_2t} u_2$ sont
solutions.

\item On suppose que $e^{rt}(u+tv)$ est solution où $u,v\in E$ et
$v\not=0$. Montrer que $u$ n'est pas proportionnel à $v$ et que $\dim\hbox{
Ker} (A-rI_E)^2\geq 2$. 
\end{enumerate}
\finenonce{006340}



\finexercice


\exercice{6341, queffelec, 2011/10/16}
\enonce{006341}{}
Trouver la solution générale de l'équation $y^{(4)}+y=0$ sous forme réelle.
On admet que la fonction $\displaystyle f(a)=\int_{\Rr} {\cos(at)\over1+t^4}\
dt$ vérifie cette équation. Sachant que $\int_{\Rr}
{dt\over1+t^4}={\pi\over\sqrt2}$, montrer que pour $a\geq0$
$$f(a)={\pi\over\sqrt2}e^{-{a\over\sqrt2}}\Big(\cos({a\over\sqrt2})+
\sin({a\over\sqrt2})\Big).$$
\finenonce{006341}



\finexercice


\exercice{6342, queffelec, 2011/10/16}
\enonce{006342}{} 
Résoudre le système
\begin{equation*}
(S) \quad \left\{
 \begin{array}{ccc}
  x'_1 & = & x_2 \\
  x'_2 & = & x_1
\end{array}
\right.
\end{equation*}
en utilisant d'abord les valeurs propres de la matrice $A$
définissant ce système, puis en calculant $A^n$ et $e^{tA}$.
\finenonce{006342}



\finexercice


\exercice{6343, queffelec, 2011/10/16}
\enonce{006343}{} 
\begin{enumerate} \item Soit le système
\begin{equation*}
x' = Ax \quad \text{avec} \quad A= \begin{pmatrix}
  1 & -2 & 0 \\
  2 & 0 & -1\\
  4 & -2 & -1
\end{pmatrix} \; .
\end{equation*}
Déterminer les valeurs propres de $A$ puis un système de
solutions de $x' = Ax$.
\item Même exercice avec
\begin{equation*}
A= \begin{pmatrix}
  1 & -1 \\
  4 & -3
\end{pmatrix} \; \;
\text{ou encore} \;\;\; A= \begin{pmatrix}
  2 & 1 & 0 & 0\\
  0 & 0 & -1 & 0\\
  1 & 2 & 1& 0\\
  -1 & 2 & 2 & 1
\end{pmatrix} \; .
\end{equation*}
\end{enumerate}
\finenonce{006343}



\finexercice


\exercice{6344, queffelec, 2011/10/16}
\enonce{006344}{} 
On considère le système linéaire $x'(t) = A(t) x(t)$, où
$A\in {\cal C} ([0, \infty ))$. Soit $\varphi $ une solution
non-triviale de ce système et soit
$$ \gamma = \limsup_{t\to \infty } \frac{1}{t} \log \|\varphi (t) \|
\;\; , -\infty \leq \gamma \leq \infty \; .$$
\begin{enumerate}
\item Montrer que $\gamma$ ne dépend pas du choix de la norme
$\| . \|$ de $\Rr^n$.
\item Montrer que $\gamma$ est une valeur finie si on suppose
 que les coefficients de la matrice $A(t)$ sont des fonctions bornées (on utilisera
 l'inégalité de Gronwall).
\item Dans le cas où $A$ est une matrice constante diagonalisable, montrer que $\gamma$
est forcément la partie réelle d'une valeur propre de $A$.
 \end{enumerate}
\finenonce{006344}



\finexercice


\exercice{6345, queffelec, 2011/10/16}
\enonce{006345}{} 
Résoudre le système $x' = 2x - y$, $y' = x+ 2y$.
Quelle est la solution vérifiant $x(0) = 1$, $y(0) =
-2$~?
\finenonce{006345}



\finexercice     


\exercice{6346, queffelec, 2011/10/16}
\enonce{006346}{} 
Soit $E$ un sous-espace vectoriel de $\Rr^n$ et
$A\in M(n,\Rr)$ une matrice préservant $E$. Si $x(t)$
est une solution de l'équation $x' = Ax$ telle que
$x(t_0) \in E$, montrer que $\forall t\in \Rr : x(t) \in
E$.
\finenonce{006346}



\finexercice     


\exercice{6347, queffelec, 2011/10/16}
\enonce{006347}{} 
Classifier et esquisser les portraits de phase des
équations $x' = Ax$ pour $A\in M(2,\Rr)$ ayant zéro
comme valeur propre.
\finenonce{006347}



\finexercice     


\exercice{6348, queffelec, 2011/10/16}
\enonce{006348}{}
 Pour quelle(s) valeur(s) de $k$ l'origine est-elle un
puits pour l'équation $x' = Ax$~?
\begin{enumerate}
\item $\begin{pmatrix} a & -k \cr k &2 \cr \end{pmatrix}$,
\item $\begin{pmatrix} 3 & 0 \cr k &-4 \cr \end{pmatrix}$,
\item $\begin{pmatrix} k^2 & 1 \cr 0 &k \cr \end{pmatrix}$,
\item $\begin{pmatrix} 0 & -1 &0 \cr 1&0&0 \cr -1 &0 &k
\cr \end{pmatrix}$.
\end{enumerate}
\finenonce{006348}



\finexercice     


\exercice{6349, queffelec, 2011/10/16}
\enonce{006349}{} 
Trouver les solutions du système $x' = -y$, $y''
= -x - y + y'$.
\finenonce{006349}



\finexercice     


\exercice{6350, queffelec, 2011/10/16}
\enonce{006350}{} 
Soit $A\in M(n,\Rr)$. Montrer que si toutes les
solutions de l'équation $x' = Ax$ sont périodiques de
même période, alors $A$ est semi-simple et le
polyn\^ome caractéristique est une puissance de
$\lambda^2 + b^2$ pour un certain $b\in \Rr$.
\finenonce{006350}



\finexercice     


\exercice{6351, queffelec, 2011/10/16}
\enonce{006351}{} 
Soit $A\in M(4,\Rr)$ semi-simple, et soient $\pm
ai$, $\pm bi$, $a>0$, $b>0$ les valeurs propres.
\begin{enumerate}
\item Montrer que si $a/b$ est rationnelle, alors
toutes les solutions de $x' = Ax$ sont périodiques.
\item Montrer que si $a/b$ est irrationnelle, alors
il existe une solution non-périodique $x(t)$ telle que
$M < \vert x(t) \vert < N$ pour certaines constantes
$M,N>0$.
\end{enumerate}
\finenonce{006351}



\finexercice     


\exercice{6352, queffelec, 2011/10/16}
\enonce{006352}{} 
Si $A$ est nilpotente, quelle est la forme des
solutions de l'équation $x'=Ax$~?

\finenonce{006352}



\finexercice     


\exercice{6353, queffelec, 2011/10/16}
\enonce{006353}{} 
Trouver des conditions nécessaires et suffisantes
pour que la matrice $\begin{pmatrix} a&b\cr
c&d\cr\end{pmatrix}$ soit
\begin{enumerate}
\item diagonalisable,
\quad semi-simple,
\quad nilpotente.
\end{enumerate}
\finenonce{006353}



\finexercice     


\exercice{6354, queffelec, 2011/10/16}
\enonce{006354}{} 
Trouver toutes les solutions périodiques de
l'équation $x^{(4)} + 2x'' + x = 0$.
\finenonce{006354}



\finexercice

\section{383.00 Etude qualititative : équilibre, stabilité}


\exercice{6841, gijs, 2011/10/16}
\enonce{006841}{}
Soit $A\in M(n,\Rr)$ une matrice de taille $n\times n$
à coefficients dans $\Rr$.
\begin{enumerate}
\item Quand dit-on que $0\in \Rr^n$ est une source
pour l'équation différentielle linéaire $x'=Ax$~?

\item Démontrer que $0\in \Rr^n$ est une source
pour  l'équation %différentielle linéaire 
$x'=Ax$ si
et seulement si pour toute solution $x(t)$ de l'équation
différentielle on a $  \lim_{t\to-\infty} x(t) =
0$.
\end{enumerate}
\finenonce{006841}



\finexercice

\section{384.00 Equation aux dérivées partielles}


\exercice{1845, ridde, 1999/11/01}

\enonce{001845}{}
R\'esoudre \`a l'aide des coordonn\'ees polaires l'\'equation aux d\'eriv\'ees partielles :
$$x\frac{\partial f}{\partial x} (x, y) + y\frac{\partial f}{\partial y} (x, y) =
\sqrt{x^{2} + y^{2}} $$
\finenonce{001845}



\finexercice



\exercice{1846, ridde, 1999/11/01}

\enonce{001846}{}
R\'esoudre l'\'equation des cordes vibrantes : $\dfrac{\partial^{2}f}{\partial x^{2}}
 =\dfrac{\partial^{2}f}{\partial y^{2}}$ \`a l'aide du changement de variables
 $u = \frac{x + y}{2}$ et $v = \frac{x-y}{2}$ (on suppose que $f$ est $C^2$).
\finenonce{001846}



\finexercice



\exercice{1847, gourio, 2001/09/01}

\enonce{001847}{}
R\'{e}soudre l'\'{e}quation aux d\'{e}riv\'{e}es partielles :
$$x\frac{\partial f}{\partial y}-y\frac{\partial f}{\partial x}=f $$
en passant en coordonn\'{e}es polaires.
\finenonce{001847}



\finexercice



\exercice{1848, gourio, 2001/09/01}

\enonce{001848}{}
R\'{e}soudre en utilisant le changement de variable $ x=u,y=uv$
l'\'{e}quation aux d\'{e}riv\'{e}es partielles suivante :
$$x^{2}\frac{\partial ^{2}f}{\partial x^{2}}+2xy\frac{\partial ^{2}f}{\partial
x\partial y}+y^{2}\frac{\partial ^{2}f}{\partial y^{2}}=0. $$
\finenonce{001848}



\finexercice



\exercice{1849, gourio, 2001/09/01}

\enonce{001849}{}
Soit $ f:{\Rr}^{2}\rightarrow {\Rr}$ une application $ C^{1 }$
homog\`{e}ne de degr\'{e} $ s>0$, i.e. telle que :
$$\forall \lambda \in {\Rr}^{+*},\forall x\in {\Rr}^{2},f\left( \lambda
x\right) =\lambda ^{s}f(x). $$
Montrer que les d\'{e}riv\'{e}es partielles de $ f$ sont homog\`{e}nes de
degr\'{e} $ s-1$ et :
$$sf(x)=x_{1}\frac{\partial f}{\partial x_{1}}(x)+x_{2}\frac{\partial f}{%
\partial x_{2}}(x). $$
\finenonce{001849}



\finexercice



\exercice{1850, ridde, 1999/11/01}

\enonce{001850}{}
Soit $f : \Rr^3 \rightarrow \Rr$ d\'erivable. On pose $F (x, y, z) = f (x-y, y-z, z-x)$.\\
Calculer $\frac{\partial F}{\partial x} + \frac{\partial F}{\partial y} +
\frac{\partial F}{\partial z}$.
\finenonce{001850}



\finexercice



\exercice{1851, ridde, 1999/11/01}

\enonce{001851}{}
Soit $f : \Rr^{2} \rightarrow \Rr$ une fonction $C^{2}$. On pose $g (x, y) = f (x^{2}
-y^{2}, 2xy)$.\\
 Calculer $\triangle (g)$ en fonction de $\triangle (f)$.
\finenonce{001851}



\finexercice



\exercice{1852, maillot, 2001/09/01}

\enonce{001852}{}
On cherche les fonctions $f:\R^2\rightarrow \R$ telles que:
\begin{equation}\label{truc}
\frac{\partial f}{\partial u} (u,v) +2u \frac{\partial f}{\partial v} (u,v)=0 \qquad \text{pour tout $(u,v)\in\R^2$}.
\end{equation}
Soit $\phi:\R^2\rightarrow\R^2$ l'application d\'efinie par $\phi(x,y)=(x,y+x^2)$.

\begin{enumerate}
\item En calculant l'application r\'eciproque, montrer que $\phi$ est
bijective. V\'erifier que $\phi$ et $\phi^{-1}$ sont de classe $\mathcal C^1$.
\item Soit $f:\R^2\rightarrow \R$ une fonction de classe $\mathcal C^1$. Posons
$g=f\circ \phi$.
  \begin{enumerate}
  \item Montrer que $g$ est de classe $\mathcal C^1$.
  \item Montrer que $f$ est solution de (\ref{truc}) si et seulement si
  $\frac{\partial g}{\partial x}=0$.
  \end{enumerate}
\item Soit $f:\R^2\rightarrow \R$ une fonction de classe $\mathcal C^1$. Montrer
que $f$ v\'erifie~(\ref{truc}) si et seulement s'il existe une fonction
$h:\R\rightarrow\R$ de classe $\mathcal C^1$ telle que $f(u,v)=h(v-u^2)$ pour tout
$(x,y)\in\R^2$.
\end{enumerate}
\finenonce{001852}



\finexercice



\exercice{1853, maillot, 2001/09/01}

\enonce{001853}{}
Soient $f~:\,\R^2 \to \R$ diff\'erentiable et $g~:\,\R \to \R$
d\'efinie par $\displaystyle g(x)=f\left(e^x \sin x, \ln
(1+x^2)\right)$.

Montrer que $g$ est d\'erivable sur $\R$ et calculer sa d\'eriv\'ee en
fonction des d\'eriv\'ees partielles de $f$.

\finenonce{001853}



\finexercice



\exercice{1854, maillot, 2001/09/01}

\enonce{001854}{}
Soient $U=\{(x,y)\in \R^2,\, x>0\}$ et $V=]0,+\infty[ \times
]-\frac \pi 2, \frac \pi 2[$. On d\'efinit la fonction
$$\begin{array}{lrll}
\Psi ~: &V &\to& \R^2\\
&(r,\theta) &\mapsto& (r \cos \theta, r\sin \theta)\end{array}$$
\begin{enumerate}
\item Montrer que $U$ et $V$ sont des ouverts de $\R^2$ et que $\Psi$ est
   de classe ${\cal C}^1$ et bijective de $V$ sur $U$. D\'eterminer
$\Psi^{-1}$.
\item Soit $f~:\, U \to \R$ de classe ${\cal C}^1$ sur $U$. On pose
$$F(r,\theta)=f \circ \Psi (r,\theta) = f(r \cos \theta, r \sin
\theta).$$
  \begin{enumerate}
   \item Montrer que $f$ est de classe ${\cal C}^1$ sur $U$ et
calculer $\frac{\partial F}{\partial r}$ et  $\frac{\partial
  F}{\partial \theta}$ en fonction de   $\frac{\partial f}{\partial
  x}$ et  $\frac{\partial f}{\partial y}$.
   \item Montrer que $f$ v\'erifie l'\'equation
$$(E)\qquad
a \frac{\partial f}{\partial x} (a,b) + b \frac{\partial f}{\partial
  y}(a,b) = \sqrt{a^2+b^2} \,\mbox{arctan} \left(\frac b a\right) \quad
\forall (a,b)\in U$$
si et seulement si $F$ v\'erifie l'\'equation

$$(E')\qquad
\frac{\partial F}{\partial r} (r_0,\theta_0)=\theta_0 \quad
\forall (r_0,\theta_0)\in V.$$
  \item D\'eterminer toutes les fonctions $f~:\, U\to \R$
 de classe ${\cal C}^1$ sur $U$ qui v\'erifient l'\'equation $(E)$.
\end{enumerate}
\end{enumerate}

\finenonce{001854}



\finexercice



\exercice{1855, maillot, 2001/09/01}

\enonce{001855}{}
Soit $D=\{(x,y)\in \R^2,\, x>0\}$. On cherche les fonctions $f \in {\cal
C}^1(D,\R)$ qui v\'erifient
$$(E) \qquad x \frac{\partial f}{\partial x} + y  \frac{\partial
f}{\partial y} =0 \quad \forall (x,y) \in D.$$
\begin{enumerate}
\item
 V\'erifier que $\varphi(x,y)=y/x$ est solution de (E).

\item Soit $g \in {\cal C}^1(\R, \R)$. Montrer que $g \circ \varphi$ est
   solution de $(E)$.

\item Soit $f$ une solution de $(E)$. Montrer que $f(u, uv)$ ne d\'epend
   que de $v$.

\item Donner l'ensemble des solutions de $(E)$.
\end{enumerate}
\finenonce{001855}



\finexercice



\exercice{1856, maillot, 2001/09/01}

\enonce{001856}{}
D\'eterminer les fonctions $f\in {\cal C}^1(\R^2,\R)$ v\'erifiant
$$\frac{\partial f}{\partial x} - \frac{\partial f}{\partial y} =0
\quad \forall (x,y) \in \R^2.$$
On pourra effectuer le changement de variables $u=x+y, \, v=x-y$.

\finenonce{001856}



\finexercice



\exercice{1857, drutu, 2003/10/01}

\enonce{001857}{}
Soient $f:\R^n \to \R$ et $g:\R^n \to \R$ deux
fonctions diff\'erentiables. En utilisant des propri\'et\'es de la diff\'erentielle, 
montrer que  $\nabla (fg)=f\cdot \nabla g +g\cdot \nabla f$.
\finenonce{001857}



\finexercice


\section{385.00 Autre}


\exercice{6355, queffelec, 2011/10/16}
\enonce{006355}{}
On considère l'équation différentielle ${d^2x(t)\over dt^2}={dx(t)\over
dt}+2x(t)$.

Ecrire le système différentiel du premier ordre associé et déterminer le noyau
résolvant $R(t,t_0)$ de ce système. En déduire $e^{tA}$ pour la matrice
$A=\begin{pmatrix}0 & 1\cr 2 & 1\cr\end{pmatrix}$.
\finenonce{006355}



\finexercice


\exercice{6356, queffelec, 2011/10/16}
\enonce{006356}{}
\begin{enumerate}
\item Montrer que la matrice $\begin{pmatrix}e^{-2t}\cos t & -\sin t\cr e^{-2t}\sin t &
\cos t\cr\end{pmatrix}$ est la résolvante du système linéaire $x'(t)=A(t).x(t)$ où 
$$A(t)=\begin{pmatrix}-2\cos^2 t & -1-\sin 2t\cr 1-\sin 2t &
-2\sin^2 t\cr\end{pmatrix}.$$ En déduire la solution du système $x'(t)=A(t).x(t)+ b(t)$ avec
$b(t)=\begin{pmatrix}e^{-t}\cos t   \cr e^{-t}\sin t 
\cr\end{pmatrix}$

\item On considère maintenant l'équation différentielle ${d^2x(t)\over
dt^2}+x(t)=f(t)$ où $f$ est une application continue sur $\Rr$. En
appliquant la méthode de variation de Lagrange, trouver la solution du système
telle que
$x(0)=x_0$.
\end{enumerate}
\finenonce{006356}



\finexercice


\exercice{6357, queffelec, 2011/10/16}
\enonce{006357}{}
On considère les équations $$\quad (1)\quad y''+p(x)y'+q(x)y=0,\quad (2)\quad
y''+p(x)y'+q(x)y=r(x)$$ où $p$, $q$ et
$r$ sont des fonctions continues d'un intervalle $I\subset\Rr$ dans $\Rr$.
Etablir ce qui suit :
\begin{enumerate}
\item Pour tout $x_0\in I$, et tout $(a,b)\in \Rr^2$, $(1)$ admet une solution
maximale définie sur $I$ tout entier, telle que $y(x_0)=a, y'(x_0)=b$.

\item Soit $x_0\in I$; les solutions de $(1)$ forment un espace vectoriel $V$ de
dimension 2 dont une base est $(y_1,y_2)$ avec $y_1(x_0)=1, y_1'(x_0)=0$,  
$y_2(x_0)=0, y_2'(x_0)=1$.

\item Soit $u$ et $v$ deux solutions de $(1)$ et $W=u'v-uv'$ leur wronskien;
trouver une équation différentielle satisfaite par $W$; en déduire que $W$ est
soit identiquement nul, soit jamais nul, et que $W\not=0\Longleftrightarrow
(u,v)$ est une base de $V$. Quel est le rapport entre $W$ et la résolvante du
système associé ?

\item La solution $y$ de $(2)$ vérifiant $y(x_0)=0,y'(x_0)=1$, où $x_0$ fixé dans
$I$, est
$$ y(x)=y(x_0)y_1(x)+y'(x_0)y_2(x)+\int_{x_0}^x
{r(t)\Big(y_1(t)y_2(x)-y_1(x)y_2(t)\Big)\over W(t)}\ dt.
$$

\item Exemple : Résoudre  $y''+4y=\tan x$.
\end{enumerate}
\finenonce{006357}



\finexercice


\exercice{6358, queffelec, 2011/10/16}
\enonce{006358}{}
On considère l'équation différentielle linéaire sur $\Rr^n$\quad
$$(1)\quad y'=A(x).y$$
où $A(x)$ est continue sur un intervalle $I$.
\begin{enumerate}
\item Montrer que si l'on suppose $A(x)A(x')=A(x')A(x)$ pour tous $x,x'\in I$, la
résolvante de $(1)$ est 
$$R(x,x_0)=\exp\Big(\int_{x_0}^x A(s)\ ds\Big)=:\exp B(x).$$

(Indic. : remarquer que $B(x)B(x')=B(x')B(x)$.)

\item Montrer que si $A(x)=\begin{pmatrix}a(x) & b(x)\cr -b(x) & a(x)\cr\end{pmatrix}$, $A$ vérifie
l'hypothèse de a) et $B(x)$ est de la forme $\begin{pmatrix}\alpha(x) & \beta(x)\cr
-\beta(x) & \alpha(x)\cr\end{pmatrix}$.

\item Résoudre l'équation $y'=A(x).y$ lorsque $a(x)=-{x\over2(1+x^2)}$ et $b(x)=
{1\over2(1+x^2)}$.

\item Résoudre l'équation $y'=A(x).y + C(x)$ lorsque $A(x)=\begin{pmatrix}\sh(x) &
1\cr -1 & \sh(x)\cr\end{pmatrix}$ et $C(x)=\begin{pmatrix}\sin x\ \sh(x) 
\cr \cos x\ \sh(x) \cr\end{pmatrix}$.
\end{enumerate}
\finenonce{006358}



\finexercice


\exercice{6359, queffelec, 2011/10/16}
\enonce{006359}{}
Soit $E$ un espace de Banach et $t\to A(t)$ une application continue de
$\Rr$ dans ${\cal L}(E,E)$. On suppose
que $A$ est périodique de période $\omega$. Cela n'implique pas nécessairement
que les solutions de $\quad (1)\quad x'=A(t).x$ soient également
$\omega$-périodiques.
\begin{enumerate}
\item Dans le cas où $E$ est un espace de dimension $2$ et $A$ est une matrice
constante, donner une condition nécessaire et suffisante pour que les solutions
de $(1)$ soient $\omega$-périodiques.


\item Dans le cas général, soit $R(t,a)$ le noyau résolvant associé à $(1)$.
\begin{enumerate}
\item Montrer que $R(t+\omega,a+\omega)=R(t,a)$ pour tout $t$.

\item Montrer que la  solution $x(t)$ de $(1)$  telle que $x(0)=x_0$ est
$\omega$-périodique si et seulement si 
$ R(\omega,0)x_0=x_0$.

\item  A quelle condition l'équation
$x'=A(t).x$ a-t-elle une solution $\omega$-périodique ?.
\end{enumerate}

\item On considère l'équation $(1)\quad x''+f(t)x=0$ où $f$ est une fonction
continue,
$\omega$-périodique.
Calculer $\det R(a+\omega,a)$; $(1)$ a-t-elle
toujours une solution $\omega$-périodique ? 
\end{enumerate}
\finenonce{006359}



\finexercice


\exercice{6360, queffelec, 2011/10/16}
\enonce{006360}{}
On considère l'équation du pendule $x''+\sin x=0$.

On sait que les solutions maximales sont définies sur $\Rr$ tout entier.
\begin{enumerate}
\item Soit $\varphi$ la solution maximale de condition initiale
$\varphi(0)=a,\varphi(0)=0$; montrer que $\varphi'(t)^2=2(\cos x(t)- \cos
a)$ et en déduire que $|x(t)|\leq a$ pour tout $t$.

\item Soit $y''=-y,\ y(0)=a,\ y'(0)=0$ le problème linéarisé correspondant.
Montrer que $Z$ définie par $Z=(x-y,\ x'-y')$ vérifie un système différentiel
du premier ordre de la forme $Z'(t)=AZ(t)+B(t)$, où $A$ est antisymétrique.
En déduire, pour tout $t$, $|x(t)-y(t)|\leq {a^3\over6}|t|$. 
 \end{enumerate}
\finenonce{006360}



\finexercice


\exercice{6361, queffelec, 2011/10/16}
\enonce{006361}{}
Soit $V$ un champ de vecteurs défini sur $\Omega\subset\Rr^n$. On dit qu'une
application $h:\Omega\to\Rr$ de classe $C^1$ est une {\it intégrale première}
de $V$, si $h\circ\varphi(t)$ est constante sur $J$ pour toute solution
$(\varphi, J)$ de l'équation autonome associée. On suppose le champ de classe
$C^1$ sur $\Omega$. 
\begin{enumerate}
\item  Montrer que 
$h$ est une intégrale première de
$V$ si et seulement si 

$h'(x).V(x)=0$ pour tout $x\in\Omega$.

\item Donner une intégrale première sur $\Rr^n$ du système différentiel $X'=AX$
où $A$ est une matrice antisymétrique $n\times n$ (commencer avec $n=2$).

\item Soit $f$ une application de classe $C^\infty$ de $\Rr$ dans $\Rr$ telle que $f(0)=0$, et on note $F(x)=\int_0^xf(u)\ du$.

Montrer que la fonction $(x,y)\to y^2+2F(x)$ est une intégrale première sur
$\Rr^2$ du champ de vecteurs 
$V(x,y)=(y,\ -f(x))$ défini sur $\Rr^2$.
On suppose que $F$ tend vers $+\infty$ lorsque $x$ tend vers
$\pm\infty$. Montrer que si une solution  $(x(t),y(t))$ de $X'=V(X)$ est définie
sur un intervalle quelconque $I$, les fonctions $x$ et $y$ sont bornées sur $I$
(remarquer que $F$ est bornée inférieurement).
\end{enumerate}
\finenonce{006361}



\finexercice


\exercice{6362, queffelec, 2011/10/16}
\enonce{006362}{Extrait de l'épreuve de septembre 97}
Soit $f$ une application de classe $C^1$ d'un ouvert $\Omega$ de ${\Rr}^n$
dans ${\Rr}^n$ telle que $f(0)=0$. On considère l'équation différentielle

$$(1)\qquad  {dx\over dt}= f(x).$$
\begin{enumerate}
\item Soit $F$ un difféomorphisme de classe $C^1$ de $\Omega$ sur un ouvert
$\Omega_1$ de ${\Rr}^n$ tel que $F(0)=0$, et on note $G$ le difféomorphisme
inverse. 
Montrer que si $\varphi$ est solution de $(1)$, $\psi=F\circ\varphi$ est
solution de l'équation 

$$(2)\qquad  {dy\over dt}= g(y),$$

où $g$ est une application de classe $C^1$ de $\Omega_1$ dans ${\Rr}^n$ que
l'on déterminera.

On suppose maintenant $n=3$.

\item Montrer que l'application $F$ de ${\Rr}^3$ dans ${\Rr}^3$ définie par
$F(x_1,x_2,x_3)=(2x_2-x_3, x_1-x_2^2, x_3)$ est un difféomorphisme de ${\Rr}^3$
 de classe $C^\infty$ tel que $F(0)=0$.

\item Déduire à l'aide de a) et b) les solutions de l'équation différentielle
$$ \begin{array}{ccc}
        {dx_1\over dt} & =&2(x_1-x_2^2)-2x_2+x_3+2x_2(x_1-x_2^2+5x_2-x_3)\\ 
{dx_2\over dt} & =&x_1-x_2^2+5x_2-x_3\\ 
{dx_3\over dt} &
=&2(x_1-x_2^2)+4x_2+x_3 
\end{array}$$
\end{enumerate}
\finenonce{006362}



\finexercice


\exercice{6363, queffelec, 2011/10/16}
\enonce{006363}{Calcul fonctionnel holomorphe}
Soit $A\in {\cal L}(\Rr^n)$ et $0<\rho=\sup\{|\lambda|;\ \lambda$ valeur 
propre de $A\}$. On va montrer sur un exemple que l'on peut calculer $f(A)$
pour toute $f$ somme d'une série entière de rayon $>\rho$.

Soit donc $A$ un opérateur de $\Rr^n$ tel que $(A-I)^2(A-2I)=0$.
\begin{enumerate}
\item On note $E_1=\ker (A-I)^2$, $E_2=\ker (A-2I)$, $p_i$ le projecteur sur $E_i$
(parallèlement à l'autre). Calculer $p_1$ et $p_2$ en fonction de $A$
(Solution : $p_1=-A(A-2I)$ et $p_2=(A_I)^2$.)

\item Calculer $A^nx$ pour $x\in E_1$, puis $x\in E_2$. Déduire de a) l'expression
de $A^n$ pour tout
$n\geq 0$ (Solution : $A^n=(I+n(A-I))A(2I-A)+2^n(A_I)^2$). 

\item Soit $f$ un polyn\^ome de degré $>2$ et $P$ le polyn\^ome minimal de
$A$. Montrer que ${f(x)\over P(x)}=g(x)+{f(2)\over x-2}-{xf(1)\over
(x-1)^2}-{f'(1)\over x-1}$ où $g$ est lui-même un polyn\^ome. En déduire $f(A)$
pour $f$ polyn\^ome puis $f$ somme d'une série entière de rayon $>2$.

\item Trouver ainsi $e^{tA}$ si $t\in\Rr$ et résoudre le système $x'=A.x$ où $A=${\bf ??}
\end{enumerate}
\finenonce{006363}



\finexercice


\exercice{6364, queffelec, 2011/10/16}
\enonce{006364}{}
On considère $A$ la matrice $ \begin{pmatrix}0 & 1 & 0  \cr  0 & 0 & 1 \cr 1 & 0 &0 
\cr\end{pmatrix}$.
\begin{enumerate}
\item Calculer $A^3$ et montrer que $e^{tA}=\begin{pmatrix}f & g & h  \cr  h & f & g
\cr g & h &f 
\cr\end{pmatrix}$, où $f(t)=\sum_0^\infty{t^{3n}\over 3n!},\
g(t)=\sum_0^\infty{t^{3n+1}\over 3n+1!},\ h(t)=\sum_0^\infty{t^{3n+2}\over
3n+2!}$. Montrer que $f(t)={1\over3}(e^t+e^{jt}+e^{j^2t})
$ et donner l'expression
de
$h$.

\item On considère $\varphi:{\Rr}\to {\Rr}$ une application de classe
$C^\infty$. Montrer qu'une solution particulière de l'équation\quad $({\cal
E})\quad y'''-y=\varphi(t)$\quad est 
$$y(t)=\int_0^th(t-s)\varphi(s)\ ds.$$

\item On suppose $\varphi$ $1$-périodique (ie $\varphi(t+1)=\varphi(t)\ \forall
t\in\Rr$).
Soit $y$ une solution de $({\cal E})$ telle que $y(0)=y(1),\ y'(0)=y'(1),\
y''(0)=y''(1)$. Montrer que $y$ est $1$-périodique.

Montrer que $({\cal E})$ possède une et une seule solution $1$-périodique.
\end{enumerate}
 \finenonce{006364}



\finexercice


\exercice{6365, queffelec, 2011/10/16}
\enonce{006365}{}
Soit $E=\Rr^n$ et $t\to A(t)$, $t\to B(t)$ deux applications de $J$ dans
${\cal L}(E)$ où $J=]\alpha,+\infty[$.

On considère les deux équations
$$(1)\ x'=A(t).x\qquad (2)\ x'=\big(A(t)+B(t)\big).x,$$
et $a\in J$. On note $R(t,a)$ la résolvante de $(1)$ telle que
$R(a,a)=I_E$.
\begin{enumerate}
\item Si $y$ est une solution de $(2)$, montrer que la fonction $z$ définie par
$y(t)=R(t,a).z(t)$ est solution d'une équation de la forme $(3)\ z'=C(t).z$,
où  $C(t)=R(a,t)B(t)R(t,a)$.

\item On suppose que $||R(t,s)||\leq k$ pour tous $t,s\in J$ où $k$ est une
constante et que $||B(t)||\leq \epsilon(t)$ où $\epsilon$ est continue sur $J$.

Montrer que $||C(t)||\leq k^2\epsilon(t)$.

\item On suppose de plus que $\int_a^\infty \epsilon(t)\ dt$ converge. Montrer (à
l'aide de Gronwall) que si $z$ est telle que $z(a)\neq0$, $||z(t)||$ est
uniformément bornée sur
$[a,+\infty[$, puis que $z$ a une limite lorsque $t\to+\infty$.
\end{enumerate}
\finenonce{006365}



\finexercice


\exercice{6366, queffelec, 2011/10/16}
\enonce{006366}{}
Soit $f\in C^1 (\Rr \times \Rr^n , \Rr^n )$ une application
bornée et soit $\varphi$ la solution maximale du problème
$$ x'(t) = f(t,x(t)) \;\; , \;\; x(t_0) =x_0 \;\; ,$$
que l'on suppose définie sur l'intervalle $I\subset \Rr$.
(Rappeler pourquoi une telle solution existe). Montrer que $\varphi$
est définie sur $I=\Rr$ tout entier. (Indication: supposer
$\beta = \sup \{t\; ; \; t\in I\} <\infty$. \'Etablir que $\varphi$
est bornée sur $[t_0, \beta[$ et que $\lim_{t\to \beta} \varphi (t)$
existe. Conclure).
\finenonce{006366}



\finexercice


\exercice{6367, queffelec, 2011/10/16}
\enonce{006367}{} 
Soit $f$ une application $ C^1$ et bornée de $ \Rr^n $ dans
$\Rr^n$ et soit $x_0, x'_0 \in \Rr^n$. Montrer que le
problème
$$ x''(t) = f(x(t) ) \;\; , \;\; x(0 ) = x_0 \; , x'(0) = x'_0 $$
admet une unique solution maximale définie sur $\Rr$.

Exemple: $x''  + \sin (x) =0$, l'équation du pendule simple .
\finenonce{006367}



\finexercice


\exercice{6368, queffelec, 2011/10/16}
\enonce{006368}{} 
Soit $a>0$ et soit $f: \Rr \times \Rr^n \to \Rr^n$ une
application de classe $C^1$ vérifiant
$$ |\langle x  ,  f(t,x)\rangle| \leq a \langle x  ,  x \rangle \quad \text{pour tout} \;\;
 (t,x)\in \Rr\times \Rr^n \; .$$ 
Soit $\varphi$ une solution
de l'équation différentielle $x'=f(t,x)$ que l'on suppose
définie sur l'intervalle $I$. 
\begin{enumerate}
\item On pose $N(t) = \langle \varphi (t)  ,  \varphi (t) \rangle$. Montrer que
l'application $N$ est dérivable sur $I$, calculer sa dérivée
et montrer qu'elle vérifie $|N'(t)|\leq 2 a N(t)$.
\item Soient $t$ et $t_0$ deux points de $I$. Comparer $N(t) $ et
$N(t_0)$.
\item Montrer que les solutions maximales de l'équation
différentielle en considération sont définies sur $\Rr$.
\item Montrer que les solutions maximales du système
\begin{equation*}
(S) \quad \left\{ \begin{array}{ccc}
  x_1'(t) & = & 2x_1(t) +t x_2(t) +x_2^2(t) \\
  x'_2(t) & = & -tx_1(t) +x_2(t) -x_1(t) x_2(t)
\end{array}\right.
\end{equation*}
sont définies sur $\Rr$.

 \end{enumerate}
\finenonce{006368}



\finexercice


\exercice{6810, gijs, 2011/10/16}
\enonce{006810}{}
Soit $W = \{\,(x,y) \in \Rr^2 \mid x>0 \,\}$ et
$\Omega = \{\,(x,y) \in \Rr^2 \mid x\ne 0 \,\}$. On
veut étudier le système d'équations
différentielles
\begin{equation}
%
\left\{ \aligned x' &= 1 \\ y' &= \frac{2y}{x} \endaligned
%
\right.
%x' = 1 \qquad \& \qquad y' = \frac{2y}{x}\ .
\tag{*}
\end{equation}

\begin{enumerate}
\item
Pour chaque condition initiale $(x_0,y_0) \in W$,
trouver la solution $(x(t),y(t))$ sur un intervalle
maximal $I$ du système $(*)$ (préciser $I$). 
Tracer quelques courbes intégrales dans $W$ pour
des conditions initiales variées dans $W$.

\item
Pour chaque condition initiale $(x_0,y_0) \in \Omega$,
trouver la solution $(x(t),y(t))$ sur un intervalle
maximal $I$ du système $(*)$ (préciser $I$). 
Tracer quelques courbes intégrales dans $\Omega$ pour
des conditions initiales variées dans $\Omega$.

\item
Trouver toutes les courbes $\gamma: \Rr \to \Rr^2$ de
classe $C^1$ vérifiant $\gamma(0) = (-1,1)$ et telles que
$\gamma$ est solution de $(*)$ partout où $\gamma(t)$
appartient à $\Omega$.

\item
Même question pour des courbes $\gamma$ de classe $C^2$.

\item
Même question pour des courbes $\gamma$ de classe $C^3$.
\end{enumerate}
\finenonce{006810}



\finexercice

\section{400.00 Tribu, fonction mesurable}


\exercice{5933, tumpach, 2010/11/11}
\enonce{005933}{}
Montrer les \'egalit\'es ensemblistes suivantes~:
$$
[a, b] =\bigcap_{n=1}^{\infty}\, ]a-\frac{1}{n}, b+\frac{1}{n}[
\quad\quad\text{et}\quad\quad]a, b[ = \bigcup_{n=1}^{\infty}\, [a+
\frac{1}{n}, b-\frac{1}{n}]
$$
\finenonce{005933}



\finexercice


\exercice{5934, tumpach, 2010/11/11}
\enonce{005934}{}
Soit $(\Omega, \Sigma, \mu)$ un espace mesur\'e et $f~:\Omega
\rightarrow \mathbb{R}$ une fonction
($\Sigma$-$\mathcal{B}(\mathbb{R})$)-mesurable. Montrer que la
troncature $f_A$ de $f$ d\'efinie par~:
$$
f_A(x) = \left\{\begin{array}{ll}-A & \text{si}\quad f(x) < -A\\
                                  f(x) & \text{si}\quad |f(x)| \leq A\\
                                  A & \text{si}\quad f(x)> A
\end{array} \right.
$$
est ($\Sigma$-$\mathcal{B}(\mathbb{R})$)-mesurable. 
\finenonce{005934}



\finexercice


\exercice{5935, tumpach, 2010/11/11}
\enonce{005935}{}
Soit $\Omega = \mathbb{N}$, $\Sigma = \mathcal{P}(\mathbb{N})$ et
$\mu$ la mesure de comptage sur $\mathbb{N}$ d\'efinie par~:
$$
\mu(E) = \sharp E = \sum_{k \in E} 1,
$$
o\`u $E\in \Sigma$. Soit $f~:\mathbb{N} \rightarrow \mathbb{R}$
une fonction positive ou nulle. Montrer que $f$ est
($\Sigma$-$\mathcal{B}(\mathbb{R})$)-mesurable et que~:
$$
\int_{\Omega} f d\mu = \sum_{n=1}^{\infty} f(n).
$$
\finenonce{005935}



\finexercice


\exercice{5936, tumpach, 2010/11/11}
\enonce{005936}{}
\label{ex:barb20}
Soit $(\Omega, \Sigma)$ un espace mesurable. On dit que
$\varphi~:\Omega \rightarrow \mathbb{R}$ est une \emph{fonction
simple} ou \emph{\'etag\'ee} si $\varphi$ est mesurable et ne
prend qu'un nombre fini de valeurs, i.e. si $\varphi$ s'\'ecrit~:
$$
\varphi = \sum_{j \in J} c_{j} \mathbf{1}_{E_{j}},
$$
o\`u $J$ est un ensemble fini, les ensembles $E_{j}$ sont
mesurables et o\`u, pour $i\neq j$, $c_{i} \neq c_{j}$ et
$E_{i}\cap E_{j} = \emptyset$. Soit $\varphi$ une fonction simple
positive. On rappelle que l'int\'egrale de $\varphi$ par rapport
\`a une mesure $\mu$ est d\'efinie par~:
$$
\int_{\Omega} \varphi \,d\mu = \int_{0}^{\infty}
\mu\left(S_{\varphi}(t) \right)\,dt,
$$
o\`u $S_{\varphi}(t) = \{x\in \Omega, \varphi(x)>t\}$.
\begin{enumerate}
\item Montrer que $$\int_{\Omega} \varphi \,d\mu = \sum_{j \in J}
c_{j} \mu(E_{j}).$$ 
\item Montrer que pour toute fonction r\'eelle
mesurable positive, $f \in \mathcal{M}^{+}(\Omega, \Sigma)$, il
existe une suite $\{\varphi_{n}\}_{n\in\mathbb{N}}$ de fonctions
simples positives telle que~:
\begin{enumerate}
\item[(a)] $0 \leq \varphi_{n}(x) \leq \varphi_{n+1}(x)$ pour tout
$x\in\Omega$ et pour tout $n\in\mathbb{N}$~; 
\item[(b)]
$\lim_{n\rightarrow +\infty} \varphi_{n}(x) = f(x)$ pour tout
$x\in\Omega$.
\end{enumerate}
\end{enumerate}
\finenonce{005936}



\finexercice


\exercice{5937, tumpach, 2010/11/11}
\enonce{005937}{}
Soit $(\Omega, \Sigma, \mu)$ un espace mesur\'e et
$f\in\mathcal{M}^+(\Omega, \Sigma)$ (i.e $f$ est une fonction
r\'eelle mesurable positive). Pour tout $E\in\Sigma$, on pose~:
$$
\lambda(E) = \int_{E} f\,d\mu = \int_{\Omega} \mathbf{1}_{A}\cdot f
\,d\mu.
$$
Monter que $\lambda$ d\'efinit une mesure sur $(\Omega, \Sigma)$.
\finenonce{005937}



\finexercice


\exercice{5938, tumpach, 2010/11/11}
\enonce{005938}{}
Soit $p>0$. Soit $f~:\mathbb{R}^n \rightarrow\mathbb{R}^+$ la fonction
d\'efinie par $$f(x) = |x|^{-p} \mathbf{1}_{\{|x|<1\}}(x).$$ Calculer
l'int\'egrale de $f$ par rapport \`a la mesure de Lebesgue de
$\mathbb{R}^n$ de deux mani\`eres diff\'erentes~:
\begin{enumerate}
\item [(i)] En utilisant les coordonn\'ees polaires et les
m\'ethodes standard de calcul d'int\'egrales~; 
\item[(ii)] En
calculant la mesure des ensembles $S_{f}(a) = \{x\in\Omega,
f(x)>a\}$ et la d\'efinition de l'int\'egrale de Lebesgue.
\end{enumerate}
\finenonce{005938}



\finexercice

\section{401.00 Mesure}


\exercice{5926, tumpach, 2010/11/11}
\enonce{005926}{}
Montrer que $\mathcal{A}$ est une $\sigma$-alg\`ebre si et
seulement si $\mathcal{A}$ est une alg\`ebre et une classe
monotone.
\finenonce{005926}



\finexercice


\exercice{5927, tumpach, 2010/11/11}
\enonce{005927}{}
Soit $(\Omega, \Sigma)$ un espace mesurable (i.e. un ensemble
$\Omega$ muni d'une tribu $\Sigma \subset \mathcal{P}(\Omega)$).
Soit $\mu$ une mesure finie sur $(\Omega, \Sigma)$. Montrer les
propri\'et\'es suivantes~: ($A,B,A_i$ sont des élements de de $\Sigma$)

\begin{enumerate}
 \item Si $A_{1}, A_{2}, \dots, A_{k}$ sont deux \`a deux
disjoints, alors
$$
\mu\left(\cup_{i=1}^{k} A_{i}\right) = \sum_{i=1}^{k} \mu(A_{i}).
$$

  \item Si $B \subset A$ alors $\mu(A\setminus B) =
\mu(A) -\mu(B)$.

  \item \emph{Monotonie~:} Si $B \subset A$ alors $\mu(B) \leq \mu(A)$.

  \item \emph{Principe inclusion-exclusion~:} $\mu(A\cup
B) = \mu(A) + \mu(B) - \mu(A\cap B)$.

  \item $\mu\left(\cup_{i=1}^{+\infty} A_{i}\right) \leq
\sum_{i=1}^{+\infty}\mu(A_{i}).$ \textit{(Rappelons que l'on a
\'egalit\'e si l'union est disjointe.)}
\end{enumerate}
\finenonce{005927}



\finexercice

\section{402.00 Lemme de Fatou, convergence monotone}


\exercice{5939, tumpach, 2010/11/11}
\enonce{005939}{}
\begin{enumerate}
\item  Soit $\{g_{n}\}_{n\in\mathbb{N}}$ une suite dans
$\mathcal{M}^{+}(\Omega, \Sigma)$. Montrer que
$$
\int_{\Omega} \left(\sum_{n=1}^{+\infty} g_{n}\right)\,d\mu =
\sum_{n=1}^{+\infty}\int_{\Omega} g_{n} \,d\mu.
$$

\item Montrer que
$$
\int_{0}^{+\infty} \frac{x^{s-1}}{e^x -1}\,dx = \Gamma(s)
\zeta(s),
$$
o\`u $\Gamma$ est la fonction d'Euler et o\`u $\zeta(s) =
\sum_{n=1}^{+\infty} n^{-s}$. (On pourra consid\'erer les
fonctions $g_{n}(x) = x^{s-1} e^{-nx} \mathbf{1}_{[0, +\infty)}$.)
\end{enumerate}
\finenonce{005939}



\finexercice


\exercice{5940, tumpach, 2010/11/11}
\enonce{005940}{}
Soit $\Omega = \mathbb{R}$, $\Sigma = \mathcal{B}(\mathbb{R})$ et
$\mu$ la mesure de Lebesgue sur $\mathbb{R}$. Si on pose $f_{n} =
\mathbf{1}_{[0, n]}$, $n\in\mathbb{N}$, alors la suite
$\{f_{n}\}_{n\in\mathbb{N}}$ est monotone croissante vers $f =
\mathbf{1}_{[0, +\infty)}$. Bien que les fonctions $f_{n}$ soient
uniform\'ement born\'ees par $1$ et que les int\'egrales des
$f_{n}$ sont finies, on a~:
$$
\int_{\Omega} f\,d\mu = +\infty.
$$
Est-ce que le th\'eor\`eme de convergence monotone s'applique dans
ce cas ?
\finenonce{005940}



\finexercice


\exercice{5941, tumpach, 2010/11/11}
\enonce{005941}{}
Soit $\Omega = \mathbb{R}$, $\Sigma = \mathcal{B}(\mathbb{R})$ et
$\mu$ la mesure de Lebesgue sur $\mathbb{R}$. Si on pose $f_{n} =
\frac{1}{n} \mathbf{1}_{[n, +\infty)}$, $n\in\mathbb{N}$, alors la suite
$\{f_{n}\}_{n\in\mathbb{N}}$ est monotone d\'ecroissante et
converge uniform\'ement vers $0$, mais
$$
0 = ~\int_{\Omega} f\,d\mu ~\neq~ \lim \int_{\Omega} f_{n}\,d\mu
~= +\infty.
$$
Est-ce que cela contredit le th\'eor\`eme de convergence
monotone ? 
\finenonce{005941}



\finexercice


\exercice{5942, tumpach, 2010/11/11}
\enonce{005942}{}
Soit $f_{n} = \frac{1}{n} \mathbf{1}_{[0, n]}$, $n\in\mathbb{N}$, et $f =
0$. Montrer que $f_{n}$ converge uniform\'ement vers $f$, mais que
$$
\int_{\Omega} f\,d\mu ~\neq ~\lim \int_{\Omega} f_{n}\,d\mu
$$
Est-ce que cela contredit le th\'eor\`eme de convergence
monotone ?
\finenonce{005942}



\finexercice


\exercice{5943, tumpach, 2010/11/11}
\enonce{005943}{}
Soit $\Omega = \mathbb{R}$, $\Sigma = \mathcal{B}(\mathbb{R})$ et
$\mu$ la mesure de Lebesgue sur $\mathbb{R}$. Soit $f_{n} =
-\frac{1}{n} \mathbf{1}_{[0, n]}$, $n\in\mathbb{N}$, et $f = 0$. Montrer
que $f_{n}$ converge uniform\'ement vers $f$ sur $\mathbb{R}$ mais
que
$$\lim \inf_{n\rightarrow+\infty} \int_{\Omega} f_{n}\,d\mu  ~<~
\int_{\Omega} f\,d\mu.
$$
Est-ce que cela contredit le lemme de Fatou ?
\finenonce{005943}



\finexercice


\exercice{5950, tumpach, 2010/11/11}
\enonce{005950}{}
\begin{enumerate}
\item Montrer que pour tout $x\in \mathbb{R}_+$,
$\displaystyle\left\{\left(1+\frac{x}{n}\right)^n\right\}$ est une
suite croissante et
$$\lim\limits_{n\rightarrow\infty}\left(1+\frac{x}{n}\right)^n
= e^x = \sum\limits_{k=0}^{\infty} \frac{x^k}{k!}.$$

\item Calculer la limite
$$ \lim_{n\to\infty} \int_{\mathbb{R}_+} \left(1+\frac{x}{n}\right)^n\mathrm{e}^{-bx} d\lambda (x)$$
o\`{u} $b>1$.
\end{enumerate}
\finenonce{005950}



\finexercice

\section{403.00 Théorème de convergence dominée}


\exercice{5944, tumpach, 2010/11/11}
\enonce{005944}{}
Soit $f\in\mathcal{M}^{+}(\Omega, \Sigma)$ tel que $\int_{\Omega}
f\,d\mu ~<~+\infty$. Montrer que $$\mu\{ x\in\Omega, ~f(x) =
+\infty\} = 0. $$ 
On pourra consid\'erer les fonctions $f_{n} = n\mathbf{1}_{\{f\geq n\}}$.
\finenonce{005944}



\finexercice


\exercice{5945, tumpach, 2010/11/11}
\enonce{005945}{}
Soit $(\Omega, \Sigma, \mu)$ un espace mesur\'e avec
$\mu(\Omega)~<~+\infty$. Soit $\{ f_{n} \}_{n\in\mathbb{N}}$ une
suite de fonctions mesurables convergeant presque partout vers une
fonction mesurable $f$. On suppose qu'il existe une constante
$C>0$ telle que $|f_{n}| \leq C$ pour tout $n\geq 1$. Montrer que
$$
\lim_{n\rightarrow+\infty}\int_{\Omega} f_{n}\,d\mu =
\int_{\Omega} f\,d\mu.
$$
\finenonce{005945}



\finexercice


\exercice{5946, tumpach, 2010/11/11}
\enonce{005946}{}
 Soit $f\in
\mathcal{L}^1(\mathbb{R})$. Que vaut la limite
$$ \lim_{n\to\infty} \int_{\mathbb{R}} f(x)\cos^n(\pi x) d\lambda (x) \
?$$
\finenonce{005946}



\finexercice


\exercice{5947, tumpach, 2010/11/11}
\enonce{005947}{}
On rappelle qu'une fonction $f~:\Omega \rightarrow \mathbb{R}$ est
dite int\'egrable si $f_{+} := \max\{f, 0\}$ et $f_{-} = \max\{-f,
0\}$ v\'erifient $\int_{\Omega} f_{+} \,d\mu ~<~+\infty$ et
$\int_{\Omega} f_{-} \,d\mu ~<~+\infty$. On note
$\mathcal{L}^{1}(\Omega, \Sigma, \mu)$ l'ensemble des fonctions
r\'eelles int\'egrables. Pour $f\in \mathcal{L}^{1}(\Omega,
\Sigma, \mu)$, on pose
$$
\int_{\Omega} f\,d\mu = \int_{\Omega} f_{+}\,d\mu - \int_{\Omega}
f_{-}\,d\mu.
$$
\begin{enumerate}
\item Montrer l'\'equivalence
$$
f\in \mathcal{L}^{1}(\Omega, \Sigma, \mu) \Leftrightarrow |f| \in
\mathcal{L}^{1}(\Omega, \Sigma, \mu)
$$
et \begin{equation}\label{leq} \left|\int_{\Omega} f\,d\mu
\right|~\leq~ \int_{\Omega} |f|\,d\mu.
\end{equation}

\item  Montrer que si $f$ est mesurable, $g$ int\'egrable et
$|f|~\leq |g|$, alors $f$ est int\'egrable et
$$
\int_{\Omega} |f|\,d\mu ~\leq~\int_{\Omega} |g|\,d\mu.
$$

\item On rappelle qu'une fonction $f~:\Omega \rightarrow
\mathbb{C}$ est dite int\'egrable si la partie r\'eelle $\text{Re}
f$ et la partie imaginaire $\text{Im} f$ de $f$ sont
int\'egrables. On pose alors
$$
\int_{\Omega} f\,d\mu =  \int_{\Omega}\text{Re} f\,d\mu + i
\int_{\Omega}\text{Im} f\,d\mu.
$$
Montrer que l'in\'egalit\'e  \eqref{leq} est v\'erifi\'ee.
\end{enumerate}
\finenonce{005947}



\finexercice


\exercice{5948, tumpach, 2010/11/11}
\enonce{005948}{}
Soit $(\Omega, \Sigma, \mu)$ un espace mesur\'e.  On dit que
$f_{n}$ converge vers $f$ \emph{en mesure} si pour tout
$\varepsilon$,
$$\lim_{n\rightarrow+\infty} \mu\{x\in\Omega, ~|f_{n}(x) - f(x)| ~> \varepsilon\} ~ =
0.
$$
Montrer que si $f_{n}\rightarrow f$ en mesure, alors il existe une
sous-suite $\{f_{n_{k}}\}_{k\in\mathbb{N}}$ de
$\{f_{n}\}_{n\in\mathbb{N}}$ qui converge vers $f$ $\mu$-presque
partout. 
\finenonce{005948}



\finexercice


\exercice{5949, tumpach, 2010/11/11}
\enonce{005949}{}
 Donner un exemple de fonction
$f:\mathbb{R}\rightarrow \mathbb{R}$ qui est int\'{e}grable au
sens de Lebesgue mais pas au sens de Riemann.
\finenonce{005949}



\finexercice


\exercice{5951, tumpach, 2010/11/11}
\enonce{005951}{}
Montrer que

\begin{enumerate}
 \item  $\displaystyle \lim_{n \to \infty} \int_0^n \left(1-\frac{x}{n}\right)^n
x^{m}dx = m !$ \ (pour tout $m\in \mathbf{N}$).

  \item  $\displaystyle \lim_{n \to \infty} \int_0^n \left(1+\frac{x}{n}\right)^n
e^{-2x}dx = 1$.
\end{enumerate}

\finenonce{005951}



\finexercice


\exercice{5952, tumpach, 2010/11/11}
\enonce{005952}{}
\label{derivation}
Montrer le th\'eor\`eme suivant, $\Omega$ étant un espace mesurable.
(On pourra utiliser le th\'eor\`eme des accroissements finis.) \\
\textbf{Th\'eor\`eme.}(D\'erivation sous le signe $\int$) \\
Soit $f~:\Omega \times \mathbb{R} \rightarrow \mathbb{C}$ une
fonction telle que
\begin{enumerate}
\item[(i)] Pour tout $s\in[s_{1}, s_{2}]$, la fonction $x\mapsto
f(x, s)$ est int\'egrable~;

\item[(ii)] pour presque tout $x$, la fonction $s\mapsto f(x, s) $
est d\'erivable sur $(s_1, s_2)$~;

\item[(ii)] il existe $g\in \mathcal{L}^1(\Omega, \mathbb{R}^{+})$
tel que pour tout $s \in [s_1, s_2]$ et pour presque tout $x\in\Omega$ on ait $|\frac{\partial f(x,s)}{\partial s}| \leq g(x)$ .
\end{enumerate}
Alors la fonction $I(s):= \int_{\Omega} f(x, s)\,d\mu(x)$ est
d\'erivable sur $(s_{1}, s_{2})$ et
$$
\frac{dI}{ds} = \int_{\Omega} \frac{\partial f(x,s)}{\partial s}
\,d\mu(x).
$$

\finenonce{005952}



\finexercice


\exercice{5953, tumpach, 2010/11/11}
\enonce{005953}{}
    Soit \ $f\in
\mathcal{L}^{1}(\Rr )$. Sa \emph{transform\'{e}e de Fourier}
est la fonction \ $\hat{f}:\Rr \rightarrow \Cc$ \ \ d\'{e}finie
par
\[
\hat{f}(y):=\int_{\Rr }\text{e}^{-ixy}f(x)dx,
\]
montrer que
\begin{enumerate}
 \item $\hat{f}$ \ est continue,

  \item $\hat{f}$ \ est born\'{e}e et \ $\sup |\hat{f}|\leq
\left\| f\right\| _{L^{1}}$ \ (=$\int_{\Rr }\left| f(x)\right|
dx$),

  \item Si $x \to xf(x)$ est int\'{e}grable, alors $\hat{f}$ \
est d\'{e}rivable et on a
\[
\frac{d}{dy}\ \hat{f}\ =\widehat{-ixf(x)}.
\] 
\end{enumerate}
\finenonce{005953}



\finexercice

\section{404.00 Intégrales multiples, théorème de Fubini}


\exercice{5957, tumpach, 2010/11/11}
\enonce{005957}{}
 Soit $f(x,y)=
\frac{x^2-y^2}{(x^2+y^2)^2}$. Montrer que
$$
\int_{-1}^{1} \left(\int_{-1}^{1}f(x,y)dx\right) dy \neq
\int_{-1}^{1} \left(\int_{-1}^{1} f(x,y)dy\right) dx.
$$
Y a-t-il contradiction avec le th\'{e}or\`{e}me de Fubini ? (on
pourra calculer l'int\'egrale de $|f|$ sur l'anneau
$S_\varepsilon=\{(x,y)\in\mathbb{R}^2| \varepsilon\leq x^2+y^2\leq
1\}$.)
\finenonce{005957}



\finexercice


\exercice{5958, tumpach, 2010/11/11}
\enonce{005958}{}
Montrer que la fonction $(x, y) \mapsto  e^{-y} \sin 2xy$ est
int\'egrable pour la mesure de Lebesgue sur $[0, 1]\times (0,
+\infty)$ ; en d\'eduire la valeur de
$$
\int_{0}^{+\infty} \frac{1}{y} (\sin y)^2 e^{-y}\,dy.
$$
\finenonce{005958}



\finexercice


\exercice{5959, tumpach, 2010/11/11}
\enonce{005959}{}
Soient $f\in L^1(\mathbb{R}^n)$ et $g\in L^p(\mathbb{R}^n)$ avec
$1\leq p\leq +\infty$, o\`u $\mathbb{R}^n$ est muni de la mesure
de Lebesgue. Montrer que, pour presque tout $x\in \mathbb{R}^n$,
la fonction $y\mapsto f(x-y)\, g(y)$ est int\'egrable sur
$\mathbb{R}^n$ et que le \emph{ produit de convolution} de $f$ et
$g$ d\'efini par
$$
f*g(x) = \int_{\mathbb{R}^n} f(x-y) \,g(y)\,dy
$$
v\'erifie $f*g(x) = g*f(x)$ et
$$
\|f*g\|_{p} \leq \|f\|_{1}\,\|g\|_{p}.
$$
\finenonce{005959}



\finexercice


\exercice{5960, tumpach, 2010/11/11}
\enonce{005960}{}
Soient $a, b>0$, et $f$ et $g$ les fonctions d\'efinies sur
$\mathbb{R}^n$ par $f(x) = e^{-\frac{a |x|^2}{2}}$ et $g(x) =
e^{-\frac{b |x|^2}{2}}$. Calculer $f*g(x)$.
\finenonce{005960}



\finexercice


\exercice{5961, tumpach, 2010/11/11}
\enonce{005961}{}
\begin{enumerate}
\item Pour tout $t>0$, on pose~:
$$
f_{t}(x) = \left(4\pi t \right)^{-\frac{n}{2}}
e^{-\frac{|x|^2}{4t}}.
$$
\begin{enumerate}
  \item Montrer que, pour tout $t>0$, $\int_{\mathbb{R}^n}
f_t(x)\,dx = 1$. 
  \item
 Montrer  que, pour tout $\delta>0$, $\lim_{t\rightarrow
0}\int_{\{|x|>\delta\}} f_t(x)\,dx = 0$.
\end{enumerate}
(On dit que $f_t$ est une \emph{approximation de la distribution de Dirac.})
\item Soit $g$ une fonction continue born\'ee. Montrer que $f_t*g$
est bien d\'efinie et que
$$\lim_{t\rightarrow 0} f_t*g(x) = g(x).$$
\end{enumerate}
\finenonce{005961}



\finexercice


\exercice{5962, tumpach, 2010/11/11}
\enonce{005962}{}
Soient $f, g\in L^1(\mu)$ o\`u $\mu$ est la mesure de Lebesgue sur
$\mathbb{R}^n$. On note $\hat{f}$ la transform\'ee de Fourier
d\'efinie par
$$
\hat{f}(y) = \int_{\mathbb{R}^n} f(x)\,e^{-2\pi i (y, x)} \,dx,
$$
o\`u $(\cdot, \cdot)$ d\'esigne le produit scalaire de
$\mathbb{R}^n.$ Montrer que
\begin{enumerate}
\item $\quad \int_{\mathbb{R}^n} f(x)\hat{g}(x)\,dx =
\int_{\mathbb{R}^n} \hat{f}(x) g(x)\,dx.$ 
\item $\quad
\widehat{f*g} = \hat{f}\,\hat{g}.$
\end{enumerate}
\finenonce{005962}



\finexercice


\exercice{5963, tumpach, 2010/11/11}
\enonce{005963}{}
Calculer la transform\'ee de Fourier de la gaussienne d\'efinie,
pour $x\in\mathbb{R}^n$, par $f(x) = e^{-\frac{a |x|^2}{2}}$, o\`u
$a>0$.
\finenonce{005963}



\finexercice

\section{405.00 Intégrale dépendant d'un paramètre}

\section{406.00 Espace Lp}


\exercice{2692, matexo1, 2002/02/01}
\enonce{002692}{}
Soit $k\in\N$ et $\alpha\in\left]0,1\right]$.
On rappelle qu'on note $C^{k,\alpha}(\R)$ l'ensemble des
fonctions $g$ de classe $C^k$ sur~$\R$, dont la
$k$-i{\`e}me d{\'e}riv{\'e}e est h{\"o}ld{\'e}rienne, c'est-{\`a}-dire v{\'e}rifie
$$\exists C>0,\ \forall x,y\in \R,\qquad
|g^{(k)}(x)-g^{(k)}(y)|\leq  C |x-y|^\alpha.$$
(En particulier, si $\alpha =1$, ce sont les fonctions
de $k$-i{\`e}me d{\'e}riv{\'e}e lipschitzienne.)
\begin{itemize}
\item Soit $f\in L^1(\R)$ {\`a} support compact, et $g
\in C^{0,\alpha}(\R)$. Montrer que $f*g \in
C^{0,\alpha}(\R)$.
 En d{\'e}duire que si $g\in C^{k,\alpha}(\R)$,
alors $f*g \in C^{k,\alpha}(\R)$.
\item Soit $f\in L^1(\R)$ {\`a} support quelconque, et $g
\in C^{0,\alpha}(\R)$ {\it born{\'e}e}. Montrer que $f*g \in
C^{0,\alpha}(\R)$ et est born{\'e}e.
 En d{\'e}duire que si $g\in C^{k,\alpha}(\R)$, born{\'e}e ainsi
que toutes ses  d{\'e}riv{\'e}es, alors $f*g$ aussi.
\end{itemize}
\finenonce{002692}



\finexercice


\exercice{5954, tumpach, 2010/11/11}
\enonce{005954}{}
\begin{enumerate}
\item Soit $a, b \geq 0$ et soit $p,q \in (1, +\infty)$ tel que
$\frac{1}{p} + \frac{1}{q} = 1$ \emph{(on dit que $p$ et $q$ sont
conjugu\'es au sens de Young)}.  Montrer l'in\'egalit\'e de
Young~:
$$
ab \leq \frac{1}{p} a^p + \frac{1}{q} b^q.
$$
On pourra consid\'erer la fonction
$\theta~:\mathbb{R}^{+}\rightarrow \mathbb{R}$ d\'efinie par
$\theta(a) = \frac{1}{p} a^{p} + \frac{1}{q} b^{q} - ab$.

\item Soit de nouveau $p,q \in (1, +\infty)$ tel que $\frac{1}{p}
+ \frac{1}{q} = 1$ et $f\in L^{p}(\mu)$, $g\in L^{q}(\mu)$. En
utilisant la question pr\'ec\'edente, montrer que pour tout
$\lambda>0$
$$
\int_{\Omega} |fg|\,d\mu \leq \frac{\lambda^{p}}{p}\int_{\Omega}
|f|^{p}\,d\mu + \frac{\lambda^{-q}}{q} \int_{\Omega}
|g|^{q}\,d\mu.
$$
Optimiser cette in\'egalit\'e par rapport \`a $\lambda$ et montrer
l'in\'egalit\'e de H\"older~:
$$
\| fg\|_{1} \leq \|f\|_{p} \,\|g\|_{q}.
$$
 Cette
in\'egalit\'e est-elle vraie pour $p=1$ et $q=+\infty$ ?

\item Soient $p$ et  $p'$ dans $[1, +\infty[$ (pas
n\'ecessairement conjugu\'es). Montrer que si $f$ appartient \`a $
L^{p}(\mu) \cap L^{p'}(\mu)$, alors $f$ appartient \`a $
L^{r}(\mu)$ pour tout $r$ compris entre $p$ et $p'$.

\item Montrer que si $\mu$ est une mesure finie alors
$$L^{\infty}(\mu) \subset \bigcap_{p\geq 1} L^{p}(\mu),
$$
et, pour tout $f$,
$$
\lim_{p\rightarrow+\infty}\|f\|_{p} = \|f\|_{\infty}.
$$

 \item  Montrer que si  $f \in L^p(\mu)$ et
$g\in L^q(\mu)$ avec $\frac 1p + \frac 1q = \frac 1r$, alors
$f\cdot g\in L^r(\mu)$ et
$$
 \| fg\|_{r} \leq \|f\|_{p}\|g\|_{q}.
$$
\end{enumerate}
\finenonce{005954}



\finexercice


\exercice{5955, tumpach, 2010/11/11}
\enonce{005955}{}
\ \\

\textbf{Th\'eor\`eme 1.}(Th\'eor\`eme de Riesz)
Pour tout $1\leq p\leq +\infty$, l'espace $L^{p}(\mu)$ est
complet.

\bigskip

\textbf{Th\'eor\`eme 2.}
Soit $p$  tel que $1\leq p \leq +\infty$ et soit
$\{f_{n}\}_{n\in\mathbb{N}}$ une suite de Cauchy dans $L^{p}(\mu)$
convergeant vers une fonction $f\in L^{p}(\mu)$. Alors il existe
une sous-suite de $\{f_{n}\}_{n\in\mathbb{N}}$ qui converge
ponctuellement presque-partout vers $f$.

\bigskip

\emph{Le but de cet exercice est de d\'emontrer les
th\'eor\`emes~1 et 2.}


\begin{enumerate}
\item \emph{Cas de $L^{\infty}(\mu)$.}
  \begin{enumerate}
  \item Soit $\{f_{n}\}_{n\in\mathbb{N}}$ une suite de Cauchy
de $L^{\infty}(\mu)$. Pour $k, m, n\geq 1$,  consid\'erons les
ensembles
\begin{eqnarray*}
A_{k} := \{ x\in\Omega, |f_{k}(x)| > \|f_{k}\|_{\infty}\}~; &
B_{m,n} := \{ x\in\Omega, |f_{m}(x) - f_{n}(x)| > \|f_{m} -
f_n\|_{\infty} \}.
\end{eqnarray*}
Montrer que $E := \bigcup_{k} A_{k}  \bigcup_{n,m} B_{m,n}$ est de
mesure nulle. 
  \item Montrer que sur le compl\'ementaire de
$E$, la suite $\{f_n\}_{n\in\mathbb{N}}$ converge uniform\'ement
vers une fonction $f$. 
  \item En d\'eduire que
$L^{\infty}(\mu)$ est complet.
  \end{enumerate}
\item \emph{Cas de $L^{p}(\mu)$.}
  \begin{enumerate}
  \item Soit $1\leq p<+\infty$ et $\{f_{n}\}_{n\in\mathbb{N}}$
une suite de Cauchy dans $L^{p}(\mu)$. Montrer qu'il existe une
sous-suite $\{f_{n_{k}}\}_{k\in\mathbb{N}}$ de
$\{f_{n}\}_{n\in\mathbb{N}}$ telle que
$\|f_{n_{k+1}}-f_{n_{k}}\|_{p}\leq 2^{-k}$. 
  \item Posons
\begin{eqnarray*}
g_{k} = \sum_{i=1}^{k}|f_{n_{i+1}} - f_{n_{i}}|\quad
\text{et}\quad g= \sum_{i=1}^{+\infty}|f_{n_{i+1}} - f_{n_{i}}|,
\end{eqnarray*}
o\`u $g$ est \`a valeurs dans $\mathbb{R}\cup\{+\infty\}$. Montrer
que pour tout $k\geq1$, on a $\|g_{k}\|_{p} < 1$, puis que
$\|g\|_{p} \leq 1$. 
  \item En d\'eduire que la s\'erie
$$f_{n_{1}} + \sum_{i=1}^{\infty} \left(f_{n_{i+1}} -
f_{n_{i}}\right)$$ est absolument convergente pour presque tout
$x\in\Omega$. Notons $f(x)$ sa somme lorsque celle-ci est finie et
posons $f(x)=0$ sinon. V\'erifier que $f$ est la limite ponctuelle
des $\{f_{n_{k}}\}_{k\in\mathbb{N}}$ pour presque tout
$x\in\Omega$. 
  \item Montrer que $f - f_{m} \in L^{p}(\mu)$,
$f\in L^{p}(\mu)$ et  que $\|f - f_{m}\|_{p}\rightarrow 0$ quand
$m\rightarrow +\infty$. Conclure.
\end{enumerate}
\end{enumerate}
\finenonce{005955}



\finexercice


\exercice{5956, tumpach, 2010/11/11}
\enonce{005956}{}
Soient $f$ et $g$ deux fonctions de $L^{p}(\mu)$ avec
$1<p<+\infty$. Montrer que la fonction $N~:\mathbb{R} \rightarrow
\mathbb{R}$ d\'efinie par
$$
N(t) = \int_{\Omega} |f(x) + t \cdot g(x)|^{p}\,d\mu
$$
est diff\'erentiable et que sa d\'eriv\'ee en $t=0$ est donn\'ee
par~:
$$
\frac{dN}{dt}_{|t=0} = p \int_{\Omega} |f(x)|^{p-2}f(x)
g(x)\,d\mu,
$$
o\`u par convention $|f(x)|^{p-2}f(x) = 0$ lorsque $f=0$.
\finenonce{005956}



\finexercice


\exercice{5964, tumpach, 2010/11/11}
\enonce{005964}{}
Soit $\Omega$ un sous-ensemble de $\mathbb{R}^n$ dont la mesure de
Lebesgue   est \emph{finie}~: $\mu(\Omega)< +\infty$. Pour tout
$1\leq p <+\infty$, on note $L^p(\Omega)$ l'espace des fonctions
$f~:\Omega \rightarrow \mathbb{C}$ telles que $\|f\|_{p} :=
\left(\int_{\Omega} |f|^{p}(x)\,dx\right)^{\frac{1}{p}}<+\infty$
modulo l'\'equivalence $f\sim g \Leftrightarrow f-g = 0 ~\mu-p.p$.
L'espace des fonctions essentiellement born\'ees sera not\'e
$L^{\infty}(\Omega).$
\begin{enumerate}
\item Montrer que si $q\leq p$, alors $L^{p}(\Omega) \subset
L^{q}(\Omega)$. En particulier, pour $1<q<2<p$, on a~:
$$
L^{\infty}(\Omega) \subset L^{p}(\Omega) \subset L^{2}(\Omega)
\subset L^{q}(\Omega) \subset L^{1}(\Omega).
$$
\item Soit $\mathcal{B}^{n}(0, 1)$ la boule unit\'e centr\'ee en
$0$ de $\mathbb{R}^{n}$. En consid\'erant les fonctions
$$
f_{\alpha}(x)= |x|^{-\alpha}
$$
montrer que pour $q<p$, l'inclusion $L^{p}(\mathcal{B}^{n}(0,1))
\subset L^{q}(\mathcal{B}^{n}(0,1))$ est stricte.
\end{enumerate}
\finenonce{005964}



\finexercice


\exercice{5965, tumpach, 2010/11/11}
\enonce{005965}{}
Soit $\Omega = \mathbb{N}$ muni de la mesure de comptage. Pour
tout $1\leq p <+\infty$, on note $\ell^p$ l'espace des suites
complexes $(u_{n})_{n\in\mathbb{N}}$ telles que $\|u\|_{p} :=
\left(\sum_{i=0}^{+\infty}
|u_{n}|^{p}\right)^{\frac{1}{p}}<+\infty$. L'espace des suites
born\'ees sera not\'e $\ell^{\infty}.$
\begin{enumerate}
\item Montrer que si $q\leq p$, alors $\ell^{q}\subset \ell^{p}$.
En particulier, pour $1<q<2<p$, on a~:
$$
\ell^{1} \subset \ell^{q} \subset \ell^{2} \subset \ell^{p}
\subset \ell^{\infty}.
$$
\item En consid\'erant les suites $u^{(\alpha)}_{n} =
n^{-\alpha}$, montrer que pour $q<p$, l'inclusion $\ell^{q}
\subset \ell^{p}$ est stricte.
\end{enumerate}
\finenonce{005965}



\finexercice


\exercice{5966, tumpach, 2010/11/11}
\enonce{005966}{}
Soit $\Omega = \mathbb{R}^{n}$ muni de la mesure de Lebesgue. Pour
tout $1\leq p <+\infty$, on note $L^p(\mathbb{R}^{n})$ l'espace
des fonctions $f~:\mathbb{R}^{n} \rightarrow \mathbb{C}$ telles
que $\|f\|_{p} := \left(\int_{\mathbb{R}^{n}}
|f|^{p}(x)\,dx\right)^{\frac{1}{p}}<+\infty$ modulo
l'\'equivalence $f\sim g \Leftrightarrow f-g = 0 ~\mu-p.p$.
L'espace des fonctions essentiellement born\'ees sera not\'e
$L^{\infty}(\mathbb{R}^n).$
\begin{enumerate}
\item 
  \begin{itemize} 
  \item Pour quelle valeur de $\alpha$ la
fonction $x\mapsto \frac{1}{\left(1 + |x|^2 \right)^{\alpha}}$
appartient-elle \`a $L^{p}(\mathbb{R}^{n})$ ?
  \item Pour
quelle valeur de $\beta$ la fonction $x\mapsto
\frac{1}{|x|^{\beta}} e^{-\frac{|x|^2}{2}}$ appartient-elle \`a
$L^{p}(\mathbb{R}^{n})$ ? 
  \item Soit $1\leq q < p \leq
+\infty$. En utilisant $(a)$ et $(b)$, trouver une fonction $f$
qui appartienne \`a $L^{q}(\mathbb{R}^{n})$ mais pas \`a
$L^{p}(\mathbb{R}^{n})$ et une fonction $g$ qui appartienne \`a
$L^{p}(\mathbb{R}^{n})$ mais pas \`a $L^{q}(\mathbb{R}^{n})$.
  \end{itemize}
\item 
  \begin{itemize} 
  \item Soit $1\leq q< p<+\infty$.
Montrer que l'espace $L^{p}(\mathbb{R}^{n}) \cap
L^{q}(\mathbb{R}^{n})$ est un espace de Banach pour la norme
$\|\cdot \|_{p,q} = \| \cdot \|_{p} + \| \cdot \|_{q}$.  
  \item
Soit $r$ tel que $q< r<p$. Montrer que
$$
\|f\|_{r} \leq \|f\|_{p}^{\alpha}\|f\|_{q}^{1-\alpha}
$$
o\`u $\frac{1}{r} = \frac{\alpha}{p} + \frac{1-\alpha}{q}$,
$\alpha\in[0, 1]$. On pourra \'ecrire $r = r\alpha + r(1-\alpha)$
et utiliser l'in\'egalit\'e de H\"older pour un couple de r\'eels
conjugu\'es bien choisi.
 \item En d\'eduire que si $f_{n}$
converge vers $f$ dans $L^{p}(\mathbb{R}^{n}) \cap
L^{q}(\mathbb{R}^{n})$ alors $f_{n}$ converge vers $f$ dans
$L^{r}(\mathbb{R}^{n})$, i.e $L^{p}(\mathbb{R}^{n}) \cap
L^{q}(\mathbb{R}^{n})$ est un sous-espace de Banach de
$L^{r}(\mathbb{R}^{n})$.
\end{itemize}
 \item  Soit $f\in L^{p}([0, +\infty[)\cap
L^{q}([0, +\infty[)$ avec $1\leq q<2<p$. Montrer que la fonction
$h$ d\'efinie par $h(r) = \frac{1}{\sqrt{r}}f(r)$ appartient \`a
$L^{1}([0, +\infty[)$ et trouver des constantes $C_{p,q}$ et
$\gamma$ telles que $\|h\|_{1} \leq C_{p, q}
\|f\|_{q}^{\gamma}\|f\|_{p}^{(1-\gamma)}$.
\end{enumerate}
\finenonce{005966}



\finexercice


\exercice{5967, tumpach, 2010/11/11}
\enonce{005967}{}
Soit $\{f_{n}\}_{n\in\mathbb{N}}$ la suite de fonctions d\'efinies
par~:
$$
f_{n}(x) = \frac{1}{\sqrt{n}} \mathbf{1}_{[n, 2n]}(x).
$$
\begin{enumerate}
\item Montrer que $f_{n}$ converge faiblement vers $0$ dans
$L^{2}([0, +\infty[)$ mais ne converge pas fortement dans $L^2([0,
+\infty[)$.

\item Montrer que $f_{n}$ converge fortement vers $0$ dans
$L^{p}([0, +\infty[)$ pour $p>2$.
\end{enumerate}
\finenonce{005967}



\finexercice


\exercice{5968, tumpach, 2010/11/11}
\enonce{005968}{}
Soit $\{f_{n}\}_{n\in\mathbb{N}}$ la suite de fonctions d\'efinies
par~:
$$
f_{n}(x) = \sqrt{n} \mathbf{1}_{[n, n+\frac{1}{n}]}(x).
$$
\begin{enumerate}
\item Montrer que $f_{n}$ converge faiblement vers $0$ dans
$L^{2}([0, +\infty[)$ mais ne converge pas fortement dans $L^2([0,
+\infty[)$.

\item Montrer que $f_{n}$ converge fortement vers $0$ dans
$L^{p}([0, +\infty[)$ pour $p<2$.
\end{enumerate}
\finenonce{005968}



\finexercice

\section{407.00 Transformée de Fourier}


\exercice{5977, tumpach, 2010/11/11}
\enonce{005977}{} 


\emph{Le but de cet exercice est de d\'emontrer le th\'eor\`eme
de Plancherel.} 

\textbf{D\'efinition.}
Soient $f, g\in L^1(\mathbb{R}^n)$. On note $\hat{f}$ la
transform\'ee de Fourier d\'efinie par
$$
\hat{f}(y) = \int_{\mathbb{R}^n} f(x)\,e^{-2\pi i (y, x)} \,dx,
$$
o\`u $(\cdot, \cdot)$ d\'esigne le produit scalaire de
$\mathbb{R}^n.$ 

\bigskip

\textbf{Th\'eor\`eme de Plancherel.}
Si $f \in L^{1}(\mathbb{R}^n)\cap L^2(\mathbb{R}^{n})$, alors
$\|\hat{f}\|_{2} = \| f\|_{2}$.

\bigskip


Soit $f \in L^{1}(\mathbb{R}^n)\cap
L^2(\mathbb{R}^{n})$.
\begin{enumerate}
\item Montrer que $\|\hat{f}\|_{\infty}\leq \|f\|_{1}$.

\item Montrer que la fonction $g_{\varepsilon}(k) = |\hat{f}(k)|^2
\,e^{-\varepsilon\pi|k|^2}$ appartient \`a
$L^{1}(\mathbb{R}^{n})$.

\item Montrer que $$\int_{\mathbb{R}^{n}} g_{\varepsilon}(k)\,dk =
\int_{\mathbb{R}^{3n}} \bar{f}(x) f(y) e^{2\pi i (k, x-y)}
e^{-\varepsilon \pi |k|^2}\,dx dy dk.$$

 \item Sachant que la transform\'ee de Fourier de la
gaussienne $h_{\varepsilon}(x) = e^{-\pi \varepsilon |x|^2}$
($\varepsilon>0$, $x\in\mathbb{R}^n$) est donn\'ee par
$\hat{h}_{\varepsilon}(k) = {\varepsilon}^{-\frac{n}{2}}
e^{-\frac{\pi|k|^2}{\varepsilon}}$, montrer que
$$\int_{\mathbb{R}^{n}} g_{\varepsilon}(k)\,dk =
\int_{\mathbb{R}^{2n}}{\varepsilon}^{-\frac{n}{2}}
e^{-\frac{\pi|x-y|^2}{\varepsilon}} \bar{f}(x) f(y) \,dx dy.$$

\item Soit $\{s_{\varepsilon}\}_{\varepsilon>0}$ la famille de
fonctions d\'efinies par~:
$$
s_{\varepsilon} = \int_{\mathbb{R}^n} {\varepsilon}^{-\frac{n}{2}}
e^{-\frac{\pi|x-y|^2}{\varepsilon}} \bar{f}(x)\,dx.
$$
Quelle est la limite dans $L^{2}(\mathbb{R}^{n})$ de
$s_{\varepsilon}$ lorsque $\varepsilon$ tend vers $0$?


\item Montrer que $$\lim_{\varepsilon\rightarrow 0}
\int_{\mathbb{R}^{n}} g_{\varepsilon}(k)\,dk =
\int_{\mathbb{R}^{n}} (\lim_{\varepsilon\rightarrow
0}s_{\varepsilon}) f(y) \,dy.$$


\item Montrer que $\lim_{\varepsilon\rightarrow 0}
\int_{\mathbb{R}^{n}} g_{\varepsilon}(k)\,dk =
\|\hat{f}\|_{2}^2$.

\item En d\'eduire que $\|\hat{f}\|_{2} = \| f\|_{2}$.

\end{enumerate}
\finenonce{005977}



\finexercice


\exercice{5978, tumpach, 2010/11/11}
\enonce{005978}{}
Soit $f~:\mathbb{R}^{3}\rightarrow \mathbb{R}$ une fonction
radiale, i.e. telle que $f(x)= h({r})$ o\`u $x=(x_1,x_2,x_3)$ et $r = |x|$ et
$h~:\mathbb{R}^{+}\rightarrow \mathbb{R}$. Montrer que la
transform\'ee de Fourier $\hat{f}$ de $f$ s'\'ecrit~: $$
\hat{f}(k) = \frac{2}{|k|}\int_{0}^{+\infty} h(r) r
\sin(2\pi|k|r)\,dr.
$$
\finenonce{005978}



\finexercice

\section{408.00 Autre}


\exercice{5928, tumpach, 2010/11/11}
\enonce{005928}{}

\emph{Le but de cet exercice est de prouver le Th\'eor\`eme de
Carath\'eodory.} 

\textbf{D\'efinition.}
Une mesure ext\'erieure sur un ensemble $\Omega$ est une
application~$m_*~:\mathcal{P}(\Omega) \rightarrow [0,+\infty]$
telle que
\begin{enumerate}
\item[(i)] $m_*(\emptyset) = 0~;$ \item[(ii)] (monotonie) $A
\subset B \Rightarrow m_*(A) \leq m_*(B)~;$ \item[(iii)]
($\sigma$-sous-additivit\'e) Pour toute suite d'ensembles $\{
A_{i} \}_{i\in\mathbb{N}^*} \subset \mathcal{P}(\Omega)$ on a
$$
m_{*}\left(\bigcup_{i=1}^{\infty} A_{i} \right) \leq
\sum_{i=1}^{\infty} m_{*}(A_{i}).
$$
\end{enumerate}


\textbf{Th\'eor\`eme de Carath\'eodory}
Soit $m_*$ une mesure ext\'erieure sur $\Omega$. Un ensemble
$A\subset\Omega$ est dit $m_*$-mesurable si pour tout $Q\subset
\Omega$ on a
$$
m_*(Q) \geq m_*(Q \cap A) + m_*(Q\cap A^c).
$$
Notons $\mathcal{M}_{m_*} \subset \mathcal{P}(\Omega)$ l'ensemble
des parties $m_*$-mesurables. Alors
\begin{enumerate}
\item $\mathcal{M}_{m_*}$ est une $\sigma$-alg\`ebre. \item $m =
m_{*}|_{\mathcal{M}_{m_*}}$ est une mesure sur $(\Omega,
\mathcal{M}_{m_*}).$ \item L'espace mesur\'e $(\Omega,
\mathcal{M}_{m_*}, m)$ est complet, i.e. si $E \in
\mathcal{M}_{m_*}$ et $m(E) = 0$, alors tout sous-ensemble
$A\subset E$ appartient \`a $\mathcal{M}_{m_*}$.
\end{enumerate}


Début de l'exercice :
\begin{enumerate} 
\item 
  \begin{enumerate} 
  \item Rappeler la d\'efinition d'une $\sigma$-alg\`ebre.

  \item V\'erifier que $\emptyset$ et $\Omega \in \mathcal{M}_{m_*}$, et
$A\in\mathcal{M}_{m_*} \Rightarrow A^c\in\mathcal{M}_{m_*}$.
  
  \item Soit $\{A_{i}\}_{i\in \mathbb{N}^*}$ un suite quelconque
d'ensembles $m_*$-mesurables. On pose $B_1 = \emptyset$,
$B_{2}=A_1$ et $B_{j} = \cup_{i=1}^{j-1} A_{i}$, pour $j\geq 2$.
Soit $Q$ un sous-ensemble de $\Omega$. Montrer par r\'ecurrence
que l'assertion $(P_k)$ suivante est v\'erifi\'ee pour tout $k\geq
1$~:
$$
(P_k)\quad \quad \quad m_*(Q) = m_*(Q \cap B_{k+1}^c) + \sum
_{j=1}^{k} m_*(Q \cap B_{j}^c \cap A_j).
$$

  \item Soit $A = \cup_{j=1}^{\infty} A_{j}$. D\'eduire de la
question pr\'ec\'edente que $$m_{*}(Q) \geq m_*(Q\cap A^c) +
\sum_{j=1}^{\infty} m_*(Q \cap B_{j}^c\cap A_j).$$ $e)$ En
remarquant que $Q\cap A = \bigcup_{j=1}^{\infty}(Q\cap B_{j}^c
\cap A_{j})$, montrer~:
$$
m_*(Q\cap A^c) + m_*(Q\cap A) \leq m_*(Q),
$$
et conclure. 
  \end{enumerate}

\item 
  \begin{enumerate} 
  \item  Rappeler la d\'efinition d'une mesure.
  
  \item En utilisant la question 1.d), montrer la
$\sigma$-additivit\'e de $m$. 
  \end{enumerate}

\item Montrer que $m$ est
compl\`ete.
\end{enumerate}
\finenonce{005928}



\finexercice


\exercice{5929, tumpach, 2010/11/11}
\enonce{005929}{}
On d\'efinit la mesure ext\'erieure de Lebesgue sur $\mathbb{R}$,
$~m_{*}: \mathcal{P}(\mathbb{R}) \rightarrow \mathbb{R},~$ par la
formule
$$
m_{*}(A) = \inf \left\{\sum_{i-1}^{\infty} (b_{i} - a_{i}) ~|~ A
\subset \bigcup_{i=1}^{\infty} ]a_{i}, b_{i}[ \right\}.
$$
Montrer qu'il s'agit bien d'une mesure ext\'erieure. 
\finenonce{005929}



\finexercice


\exercice{5930, tumpach, 2010/11/11}
\enonce{005930}{}
On d\'efinit $m_{*}~:\mathcal{P}(\Omega) \rightarrow \mathbb{R}$
par
$$
m_{*}(A) = \left\{\begin{array}{ll} 0 & \text{si}~A = \emptyset\\
1 & \text{sinon}. \end{array}\right.
$$
\begin{enumerate}
\item Montrer que $m_*$ est une mesure ext\'erieure. 
\item Quels sont les ensembles $m_*$-mesurables ? 
\item V\'erifier le th\'eor\`eme de Carath\'eodory sur cet exemple.
\end{enumerate}
\finenonce{005930}



\finexercice


\exercice{5931, tumpach, 2010/11/11}
\enonce{005931}{}
\begin{enumerate}
\item Calculer $\int_{-\infty}^{+\infty} e^{-x^2}\,dx$.\\
\emph{Indication~:} On pourra d'abord calculer
$\int_{\mathbb{R}^{2}} e^{-(x^2 + y^2)}\, dx dy$ en passant en
coordonn\'ees polaires. 

\item  \emph{Calcul de l'aire de la
sph\`ere unit\'e de $\mathbb{R}^{n}$.} Soit $ \mathcal{S}_{n-1} =
\{(x_1, \dots, x_n) \in \mathbb{R}^n ,\, \sum_{i=1}^{n} x_{i}^2 =
1 \} $ la sph\`ere unit\'e de $\mathbb{R}^{n}$. On note
$\mathcal{A}_{n-1}$ son aire. Calculer
$$
\int_{\mathbb{R}^n} e^{-\sum_{i=1}^{n} x_{i}^2} \,dx_{1}\dots
dx_{n}
$$
en fonction de $\mathcal{A}_{n-1}$. En d\'eduire l'expression de
$\mathcal{A}_{n-1}$ en fonction de la fonction $\Gamma$~:
$$
\Gamma(s) := \int_{0}^{+\infty} x^{s-1} e^{-x}\,dx.
$$

\item \emph{Calcul du volume de la boule unit\'e de
$\mathbb{R}^{n}$.} Soit $ \mathcal{B}_{n} = \{(x_1, \dots, x_n)
\in \mathbb{R}^n ,\, \sum_{i=1}^{n} x_{i}^2 \leq 1 \} $ la boule
ferm\'ee de rayon 1 dans $\mathbb{R}^n$. On note $\mathcal{V}_{n}$
son volume. Montrer que $\mathcal{V}_{n} =
\frac{\mathcal{A}_{n-1}}{n}$. En d\'eduire que~:
$$
\mathcal{V}_{n} =
\frac{{\pi}^{\frac{n}{2}}}{\Gamma\left(\frac{n}{2} + 1 \right)}.
$$

\item \emph{Application~:} Que vaut l'aire de la sph\`ere de rayon
$R$ dans $\mathbb{R}^2$? $\mathbb{R}^3$? Que vaut le volume de la
boule de rayon $R$ dans $\mathbb{R}$? $\mathbb{R}^{2}$?
$\mathbb{R}^{3}$?
\end{enumerate}
\finenonce{005931}



\finexercice


\exercice{5932, tumpach, 2010/11/11}
\enonce{005932}{}

\emph{Cet exercice fournit une autre m\'ethode de calcul du volume
de la boule unit\'e $\mathcal{B}_{n}$ de $\mathbb{R}^{n}$ et de
l'aire de la sph\`ere $\mathcal{S}_{n-1} \subset \mathbb{R}^{n}$.}
On conserve les notations de l'exercice pr\'ec\'edent.
\begin{enumerate}
\item Montrer que $\mathcal{V}_{n} = I_{n} \cdot
\mathcal{V}_{n-1}$, o\`u $I_{n} = \int_{0}^{\pi} \left(\sin \theta
\right)^{n} \,d\theta.$ \item V\'erifier que $I_{n} =
\frac{n-1}{n} I_{n-2}.$ \item Calculer $\mathcal{V}_{n}$ pour $n =
 1, 2, \dots, 7$. \item Calculer $\mathcal{A}_{n-1}$ pour $n = 
1, 2, \dots, 6$.
\end{enumerate}

\finenonce{005932}



\finexercice


\exercice{5969, tumpach, 2010/11/11}
\enonce{005969}{}


\textbf{D\'efinition.}

On dit qu'un espace m\'etrique $E$ est \emph{s\'eparable} s'il
existe un sous-ensemble $\mathcal{F} \subset E$ d\'enombrable et
dense. 


\bigskip

\textbf{Th\'eor\`eme}
L'espace $L^p(\mathbb{R}^{n})$ est s\'eparable pour $1\leq p
<+\infty$.

\bigskip

Le but de cet exercice  est de d\'emontrer ce th\'eor\`eme.


\begin{enumerate}
\item Pour $j = 1, 2, 3, \dots$ et $m = (m_1,\ldots,m_n) \in\mathbb{Z}^{n}$, on
consid\`ere les cubes
$$
\Gamma_{j, m}~:= \{ x\in\mathbb{R}^n,~ 2^{-j}m_{i} < x_i \leq
2^{-j}(m_i + 1),~ i = 1, \dots, n\}.
$$
Montrer que pour tout $j \in \mathbb{N}^*$,
$\bigcup_{m\in\mathbb{Z}^n} \Gamma_{j, m} = \mathbb{R}^{n}$. 

\item Pour $j\in\mathbb{N}^*$, on consid\`ere l'ensemble
$\mathcal{F}_{j}$ de fonctions $\varphi$ de la forme~:
$$
\varphi(x) = \sum_{m\in\mathbb{Z}^n} c_{j, m}\, \mathbf{1}_{\Gamma_{j,
m}},
$$
o\`u les constantes $c_{j, m}\in\mathbb{Q}$ et sont nulles sauf un
nombre fini. Montrer que l'ensemble
$$
\mathcal{F} = \bigcup_{j = 1}^{\infty} \mathcal{F}_{j}
$$
est d\'enombrable.

\item Le but de cette question est de montrer que toute fonction
continue \`a support compact peut \^etre approch\'ee \`a
$\varepsilon$ pr\`es en norme $L^{p}$ par un \'el\'ement de la
famille $\mathcal{F}$. Soit $\tilde{f}$ une fonction continue \`a
support compact et soit $\varepsilon>0$ fix\'e.

\begin{itemize}
\item  Montrer que pour tout $\varepsilon'>0$, il existe $j
\in\mathbb{N}^*$, tel que $\forall m\in\mathbb{Z}^n$,
$$
x, y \in \Gamma_{j, m} \Rightarrow |\tilde{f}(x) - \tilde{f}(y)|<
\varepsilon'.
$$
\item Soit $\varepsilon'>0$ fix\'e et $j$ comme dans la
question pr\'ec\'edente. On consid\`ere la fonction
$\tilde{f}_{j}$ d\'efinie par~:
$$
\tilde{f}_{j}(x) = 2^{nj} \int_{\Gamma_{j, m}}\tilde{f}(y)\,dy
\quad\text{lorsque}\quad x\in\Gamma_{j, m},
$$
i.e. la valeur de $\tilde{f}_{j}$ en un point $x\in\mathbb{R}^{n}$
est la valeur moyenne de la fonction $\tilde{f}$ sur le cube
$\Gamma_{j, m}$ de cot\'e $2^{-j}$ qui contient $x$. Montrer que
$\forall m\in\mathbb{Z}^n$,
$$
x \in \Gamma_{j, m} \Rightarrow |\tilde{f}(x) - \tilde{f}_{j}(x)|<
\varepsilon',
$$
et en d\'eduire que
$$
\|\tilde{f} - \tilde{f}_{j}\|_p <
\text{Volume}(\gamma)^{\frac{1}{p}}\cdot\varepsilon'
$$
o\`u $\gamma$ est un cube de la forme $\{x\in\mathbb{R}^{n},
-2^{J} \leq x_i \leq 2^J\}$ en dehors duquel $\tilde{f}$ est
nulle.


\item En d\'eduire qu'il existe $f_{j}\in\mathcal{F}_{j}$
telle que $\|\tilde{f} -f_{j}\|_{p} < \varepsilon.$ (On rappelle
que les \'el\'ements de $\mathcal{F}_{j}$ ne prennent que des
valeurs rationnelles.)
\end{itemize}


\item Montrer que toute fonction $f\in L^{p}(\mathbb{R}^{n})$,
$1\leq p <+\infty$, peut \^etre approch\'ee \`a $\varepsilon$
pr\`es en norme $L^{p}$ par un \'el\'ement de la famille
$\mathcal{F}$. Conclure.

\end{enumerate}
\finenonce{005969}



\finexercice


\exercice{5970, tumpach, 2010/11/11}
\enonce{005970}{}

\textbf{Th\'eor\`eme.}
L'espace $L^{\infty}(\mathbb{R}^{n})$ n'est pas s\'eparable.

\bigskip

Le but de cet exercice  est de d\'emontrer ce th\'eor\`eme.
\begin{enumerate}
\item Soit $E$ un espace de Banach. On suppose qu'il existe une
famille $(O_{i})_{i\in I}$ telle que
\begin{itemize}
\item Pour tout $i\in I$, $O_i$ est un ouvert non vide de
$E$. 
\item $O_i \cap O_j = \emptyset $ si $i\neq j$.
\item $I$ n'est pas d\'enombrable.
\end{itemize}
Montrer que $E$ n'est pas s\'eparable. (On pourra raisonner par
l'absurde).

\item Pour tout $a\in \mathbb{R}^{n}$,  on pose $f_{a} =
\mathbf{1}_{\mathcal{B}(a, 1)}$ o\`u $\mathcal{B}(a, 1)$ est la boule de
$\mathbb{R}^n$ de rayon $1$ centr\'ee en $a$. Montrer que la
famille
$$
O_a = \{ f\in L^{\infty}(\mathbb{R}^n), ~ \|f - f_a \|_{\infty} <
\frac{1}{2}\},
$$
o\`u $a$ parcourt les points de $\mathbb{R}^n$, satisfait (a), (b)
et (c). Conclure.

\end{enumerate}
\finenonce{005970}



\finexercice


\exercice{5971, tumpach, 2010/11/11}
\enonce{005971}{}



\textbf{D\'efinition.}
Soient $\mu$ et $\nu$ deux mesures  sur un espace mesur\'e
$(\Omega, \Sigma)$. On dit que $\nu$ est absolument continue par
rapport \`a  $\mu$ et on \'ecrit $\nu<<\mu$ si
$$
\mu(S) = 0 \Rightarrow \nu(S) = 0
$$
pour tout $S\in\Sigma$. 

\bigskip

\textbf{Th\'eor\`eme de Radon-Nikodym.}
%\label{t}
Soient $\mu$ et $\nu$ deux mesures finies  sur un espace mesur\'e
$(\Omega, \Sigma)$. Si $\nu$ est absolument continue par rapport
\`a $\mu$, alors il existe une fonction positive $h\in L^1(\Omega,
\mu)$ telle que pour toute fonction positive mesurable $F$ on a~:
\begin{equation}\label{radon}
\int_{\Omega} F(x) \,d\nu(x) = \int_{\Omega} F(x) h(x)\,d\mu(x).
\end{equation}

\bigskip


Le but de cet exercice est de d\'emontrer ce th\'eor\`eme de Radon-Nikodym.
\begin{enumerate}
\item Posons
$$
\alpha = \mu + 2\nu, \quad\quad \omega = 2\mu + \nu.
$$
On consid\`ere l'espace de Hilbert $L^2(\Omega, \alpha)$ des
fonctions de carr\'e int\'egrable par rapport \`a la mesure
$\alpha$ et l'application lin\'eaire $\varphi~: L^2(\Omega,
\alpha) \rightarrow \mathbb{C}$ donn\'ee par~:
$$
\varphi(f) = \int_{\Omega} f(x)\,d\omega(x).
$$
Montrer que $\varphi~: L^2(\Omega, \alpha) \rightarrow \mathbb{C}$
est une application lin\'eaire continue.

\item En d\'eduire qu'il existe $g\in L^{2}(\Omega, \alpha)$ tel
que pour tout $f \in L^2(\Omega, \alpha)$~:
$$
\int_{\Omega} f(2g - 1) \,d\nu =  \int_{\Omega} f(2 - g) \,d\mu.
$$
\item Montrer que les ensembles $S_{1l} := \{x\in\Omega, g(x) <
\frac{1}{2} - \frac{1}{l}\}$ et  $S_{2l} := \{x\in \Omega, g(x)
> 2 + \frac{1}{l}\}$ o\`u $l\in\mathbb{N}^*$ v\'erifient $\mu(S_{jl}) = \nu(S_{jl})  = 0$.
En d\'eduire que l'on peut choisir la fonction $g$ de telle
mani\`ere que $\frac{1}{2} \leq g \leq 2$.  Montrer que l'ensemble
$Z = \{ x\in \Omega~: g(x) = \frac{1}{2}\}$ est de $\mu$-mesure
$0$.

\item Montrer que la fonction $$h(x) = \frac{2 - g(x)}{2g(x) -
1}$$ est bien d\'efinie, positive, appartient \`a $L^{1}(\Omega,
\mu)$ et satisfait \eqref{radon}.

\end{enumerate}
\finenonce{005971}



\finexercice


\exercice{5972, tumpach, 2010/11/11}
\enonce{005972}{}
\label{exo:beta}
\begin{enumerate}
\item On d\'{e}finit  la fonction B\^{e}ta  par $B(a,b):=\int_{0}^{1}s^{a-1}(1-s)^{b-1}ds$, montrer que \\
$$
B\left(1+\frac{d}{2},\frac{m}{2}\right) = 2\int_{0}^{1}\left(
1-r^{2}\right) ^{d/2}r^{m-1}dr
$$

\item D\'{e}montrer que $\displaystyle B(a,b)=\frac{\Gamma
(a)\Gamma (b)}{\Gamma (a+b)}$.

\item Calculer $\int_{\mathbb{R}^{n}} \frac{1}{\left( 1 +
|x|^2\right)^{\alpha}}\,dx $ en fonction de la fonction B\^{e}ta.
\end{enumerate}
\finenonce{005972}



\finexercice


\exercice{5973, tumpach, 2010/11/11}
\enonce{005973}{Coordonn\'ees sph\'eriques dans $\mathbb{R}^{n}$}
Soit $\Omega'$ l'ouvert de $\mathbb{R}^{n}$ d\'efini par
$$
\Omega' = \{(r, \theta_1, \dots, \theta_{n-1})\in\mathbb{R}^{n}~|~
0< r, 0 <\theta_1, \dots, \theta_{n-2} < \pi, 0 < \theta_{n-1} <
2\pi\}.
$$
Soit l'application $S$ de $\Omega'$ dans $\mathbb{R}^{n}$
d\'efinie par
$$
\begin{array}{lcl}
x_1 &= &r \cos \theta_1,\\
x_2 &=& r \sin \theta_1 \cos \theta_2,\\
\vdots & & \\
x_{n-2} &=&  r \sin \theta_1 \sin \theta_2 \dots \cos \theta_{n-2}\\
x_{n-1} &=& r \sin \theta_1 \sin\theta_2 \dots \sin \theta_{n-2}
\cos \theta_{n-1}\\
x_{n} &=& r \sin \theta_1 \sin \theta_2 \dots \sin\theta_{n-2}
\sin \theta_{n-1},
\end{array}
$$
o\`u $(x_1, \dots, x_n)$ sont les coordonn\'ees cart\'esiennes de
$x\in\mathbb{R}^n$.
\begin{enumerate}
\item Soit $\Omega = \mathbb{R}^n\setminus\{
x\in\mathbb{R}^n~|~x_{n} = 0~\text{et}~x_{n-1} \geq 0~\}.$ Montrer
que $\Omega$ est une partie ouverte de $\mathbb{R}^n$ dont le
compl\'ementaire est de mesure nulle, et que $S$ est un
diff\'eomorphisme de $\Omega'$ sur $\Omega$.

 \item Soit $f$
une fonction bor\'elienne positive sur $\mathbb{R}^{n}$. Montrer
que
$$
\begin{array}{lcl}
\int_{\mathbb{R}^n} f(x)\,dx &=& \int_{\Omega'} (f\circ S)(r,
\theta_1, \dots, \theta_{n-1})\,r^{n-1}
\sin^{n-2}\theta_{1}\,\sin^{n-3}\theta_2 \dots \sin
\theta_{n-2}\,dr\,d\theta_1 \dots d\theta_{n-1}\\ & & \\ & =&
\int_{\Omega'} (f\circ S)(r, \theta_1, \dots,
\theta_{n-1})\,r^{n-1}\,dr\,d\sigma,
\end{array}
$$
o\`u $d\sigma$ est la mesure uniforme sur la sph\`ere unit\'e de
$\mathbb{R}^{n}$. 

\item En utilisant les coordonn\'ees
sph\'eriques, calculer le volume $\mathcal{V}_4$ de la boule
unit\'e de $\mathbb{R}^4$ et l'aire $\mathcal{A}_{3}$ de la
sph\`ere unit\'e $\mathcal{S}^3$ de $\mathbb{R}^4$.
\end{enumerate}
\finenonce{005973}



\finexercice


\exercice{5974, tumpach, 2010/11/11}
\enonce{005974}{Th\'eor\`eme de Newton}
Soit $g$ une fonction sur $\mathbb{R}^+$ et
$f~:\mathbb{R}^3\rightarrow \mathbb{R}$ telle que $f(x) = g(|x|)$,
o\`u $|x|$ d\'esigne la norme de $x$ dans $\mathbb{R}^3$.
\begin{enumerate}
\item Montrer que pour $r = |x|$, on a
$$
\int_{\mathbb{R}^3} \frac{f(y)}{|x - y|}\,dy = 4\pi
\frac{1}{r}\int_{0}^{r} g(s) s^2\,ds + 4\pi\int_{r}^{+\infty} g(s)
s\,ds.
$$
\item Que peut-on en d\'eduire pour une distribution de masse
$f(x) = g(|x|)$ lorsque $g$ est \`a support dans $[0, R]$ ?
\end{enumerate}
\finenonce{005974}



\finexercice


\exercice{5975, tumpach, 2010/11/11}
\enonce{005975}{}
Soit $x\in\mathbb{R}^{d}$, $d = 1, 2$ et $r = |x|$. On consid\`ere
$f~: \mathbb{R}^d\rightarrow \mathbb{R}$ donn\'ee par $$f(x) =
h(r) = r^2 (1+r^2)^{-2}.$$ 
\begin{enumerate} 
\item Calculer
pour $d = 1$  le r\'earrangement \`a sym\'etrie sph\'erique
d\'ecroissant $f^*$ de $f$. 
\item M\^eme question pour $d =
2$. 
\item Calculer $\|f\|_{2}^{2}$ pour $d = 1$ puis $d = 2$.
\end{enumerate}
\finenonce{005975}



\finexercice


\exercice{5976, tumpach, 2010/11/11}
\enonce{005976}{}
Soit $f~:\mathbb{R}^d \rightarrow \mathbb{R}$ la fonction donn\'ee
par $f(x) = e^{-x^2 + ax}$, o\`u $a\in\mathbb{R}$. Calculer le
r\'earrangement \`a sym\'etrie sph\'erique d\'ecroissant $f^*$ de
$f$. 
\finenonce{005976}



\finexercice


\exercice{5979, tumpach, 2010/11/11}
\enonce{005979}{}

\textbf{D\'efinition.}
Soit $h\in \mathbb{R}^{n}$. On d\'efinit l'op\'erateur de
translation par $h$, not\'e $\tau_{h}$, agissant sur une fonction
$f~:\mathbb{R}^n \rightarrow \mathbb{R}$ par
$
\tau_{h}f(x) := f(x-h), \quad \forall x\in\mathbb{R}^n.
$

\bigskip

\textbf{Th\'eor\`eme.}
%\label{th:t1}
Si $f\in L^{p}(\mathbb{R}^n)$ avec $1\leq p < +\infty$, alors
$
\lim_{h\rightarrow 0}\|\tau_{h}f - f\|_{p} = 0,
$
i.e. $\tau_{h}f$ tend vers $f$ dans $L^{p}(\mathbb{R}^n)$ lorsque
$h$ tend vers $0$.

\bigskip

Le but de cet exercice est de d\'emontrer ce th\'eor\`eme.
Soit $1 \leq p < +\infty$.
\begin{enumerate}
\item Montrer que si $f$ est continue \`a support compact dans la
boule $\mathcal{B}(0, M)$ centr\'ee en $0$ et de rayon $M$, et si
$|h|\leq 1$, alors
$$
 |f(x - h) - f(x)|^p \leq \mathbf{1}_{B(0, M+1)}
2^p \|f\|_{\infty}^p.
$$
o\`u $\mathcal{B}(0, M+1)$ est la boule centr\'ee en $0$ de rayon
$M+1$.

\item En d\'eduire que pour $f$ continue \`a support compact, on a
$$
\lim_{h\rightarrow 0}\|\tau_{h}f - f\|_{p} = 0.
$$

\item D\'emontrer le th\'eor\`eme pour une fonction quelconque
dans $L^p(\mathbb{R}^n)$, $1\leq p <+\infty$.

\item Que se passe-t-il pour $p = \infty$ ?
\end{enumerate}
\finenonce{005979}



\finexercice


\exercice{5980, tumpach, 2010/11/11}
\enonce{005980}{}


\textbf{Th\'eor\`eme}
Soit $\{\varphi_{n}\}_{n\in\mathbb{N}}$ une suite de fonctions de
$\mathbb{R}^n$ dans $\mathbb{R}$  telles que~:
\begin{itemize}
\item[(i)] $\int_{\mathbb{R}^n} \varphi_{n} = 1$ \item[(ii)] il
existe une constante $K
> 0$ telle que $~\sup_{n\in\mathbb{N}} \int_{\mathbb{R}^n} |\varphi_{n}|(x) \,dx \leq
K$ \item[(iii)] Pour tout $\varepsilon>0$, on a
$\lim_{n\rightarrow +\infty}\int_{\|x\|>\varepsilon}
|\varphi_{n}(x)|\,dx = 0$.
\end{itemize}
Alors pour tout $f\in L^p(\mathbb{R}^n)$, $1\leq p <+\infty$,
$\lim_{n\rightarrow+\infty}\|\varphi_{n}*f - f\|_p = 0$.


\bigskip

Le but de cet exercice est de d\'emontrer ce th\'eor\`eme.


Soit $\{\varphi_{n}\}_{n\in\mathbb{N}}$ une suite de fonctions
v\'erifiant les hypoth\`eses (i), (ii) et (iii) du
th\'eor\`eme, et soit $1\leq p <+\infty$.
\begin{enumerate}
\item En notant $q$ l'exposant conjugu\'e de $p$ ($\frac{1}{p} +
\frac{1}{q} = 1$), et en utilisant l'in\'egalit\'e de H\"older
pour la mesure $d\nu(x) = |\varphi_n|(x)\,dx$, montrer que
$$
|\varphi_n * f - f|^p(x) \leq \left(\int_{\mathbb{R}^n}
|\varphi_n|(x)\,dx \right)^{\frac{p}{q}} \left(\int_{\mathbb{R}^n}
|f(x - y) - f(x)|^p |\varphi_n|(y)\,dy\right).
$$
\item En d\'eduire que
$$
\|\varphi_n * f - f\|_{p}^{p} \leq K^{\frac{p}{q}}
\int_{\mathbb{R}^{n}} \|\tau_{y}f - f\|_{p}^p |\varphi_n(y)|\,dy.
$$

\item Soit $\delta>0$, montrer que
$$
\|\varphi_n * f - f\|_{p}^{p} \leq
K^{\frac{p}{q}}\left(\sup_{|y|\leq\delta}\|\tau_{y}f - f\|_{p}^p +
2^p\|f\|_{p}^{p}\int_{|y|>\delta} |\varphi_{n}(y)|\,dy \right).
$$

\item En d\'eduire le th\'eor\`eme cherch\'e.
\end{enumerate}
\finenonce{005980}



\finexercice


\exercice{5981, tumpach, 2010/11/11}
\enonce{005981}{} Soit $f$ une fonction dans
$\mathcal{C}^{\infty}_c(\mathbb{R}^{n})$ et $0< \alpha< n$. Posons
$c_{\alpha}~:= \pi^{-\alpha/2} \Gamma(\alpha/2)$. En utilisant
l'identit\'e
$$
c_{\alpha} |k|^{-\alpha} = \int_{0}^{\infty} e^{-\pi
|k|^2\lambda}\,\lambda^{\frac{\alpha}{2}-1}\,d\lambda,
$$
montrer que
$$
c_{\alpha}\left(|k|^{-\alpha}\hat{f}(k)\right)^{\vee}(x) =
c_{n-\alpha}\int_{\mathbb{R}^{n}} |x- y|^{\alpha - n}f(y)dy,
$$
o\`u la notation $h^{\vee}$ d\'esigne la transform\'ee de Fourier
inverse d'une fonction $h$ donn\'ee par $h^{\vee}(x) :=
\hat{h}(-x)$.
\finenonce{005981}



\finexercice

\section{420.00 Espace topologique, espace métrique}


\exercice{1867, roussel, 2001/09/01}

\enonce{001867}{}
 Soit $(E,d)$ un espace
m\'etrique.
%Soit $f : \mathbb{R}_{+} \longrightarrow \mathbb{R}_{+}$, strictement
%\hspace*{5cm}croissante, v\'erifiant : $$f(0)=0 \mbox{~et~} \forall ~(u,v)
%\in ~\mathbb{R}^{2}_{+}, ~f(u+v) \leq f(u)+f(v).$$
\begin{enumerate}
\item Montrer que $d'(x,y)=\sqrt{d(x,y)}$ est une distance sur  $E$. Enoncer
des conditions suffisantes sur une fonction $f$, d\'efinie de
$\mathbb{R}_{+}$ dans $\mathbb{R}_{+}$ pour que $(x,y) \longrightarrow
f(d(x,y))$ soit une distance sur $E$.
\item Montrer que l'application $d''$ d\'efinie sur $E\times E$ par
$d''(x,y)={\displaystyle d(x,y)\over \displaystyle1+d(x,y)}$ est une distance
sur $E$. \emph{Indication} : On utilisera la croissance de la fonction
$u \longrightarrow \displaystyle{\frac{u}{1+u}}.$
\item Comparer les distances $d$ et $d''$ .
\item Dans le cas o\`u $E$ est l'ensemble des nombres r\'eels et o\`u $d$ est
la distance valeur absolue, construire $B_{d''}(0,a)$ o\`u $a$ est un r\'eel.
\end{enumerate}
\finenonce{001867}



\finexercice



\exercice{1868, roussel, 2001/09/01}

\enonce{001868}{}

Soit $(E,d)$ un espace m\'etrique complet, et $f$ une application
de $E$ dans $E$ telle qu'il existe \hspace*{18pt}$k\in \mathbb{R},\ 0< k< 1$
tel que $d(f(x),f(y))\leq k\ d(x,y)\ \ \forall x\in E,\ \forall y\in E$.
\begin{enumerate}
\item Montrer que $f$ est continue sur $(E,d)$.
\item Soient $x_0\in E \mbox{ et pour } n\geq 0,\ x_{n+1}=f(x_n)$.
Montrer que la suite $(x_n)_{n\geq 0}$ est de Cauchy dans $(E,d)$.
\item Montrer que cette suite converge vers un point fixe de $f$,
c'est-\`a-dire une solution de $f(l)=l$. Montrer que ce point fixe est unique.
\item Application: montrer que le syst\`eme
$\displaystyle
  \left\{
 \begin{array}{ll}
x_1 &=\frac{1}{5}(2\sin x_1 + \cos x_2 )\\
x_2 &=\frac{1}{5}( \cos x_1 +3 \sin x_2)
\end{array}
\right.
$ admet une solution unique $(x_1,x_2)\in \mathbb{R}^2$.
\end{enumerate}
\finenonce{001868}



\finexercice



\exercice{2340, queffelec, 2003/10/01}
\enonce{002340}{}
\begin{enumerate}
\item  Rappeler les d\'efinitions d'une borne sup\'erieure (inf\'erieure) d'un ensemble de
nombres r\'eels. Si $A$ et $B$ sont deux ensembles born\'es non vides de $\Rr$, comparer 
avec $\sup A$, $\inf A$, $\sup B$ et $\inf B$ les nombres suivants : 

\quad (i) $\sup(A+B)$, \quad (ii) $\sup(A\cup B)$, \quad (iii) $\sup(A\cap
B)$, \quad (iv) $\inf(A\cup B)$, \quad (v) $\inf(A\cap B)$.

\item Pour $x\in\Rr^n$ et $A\subset \Rr^n$ on d\'efinit $d(x,A)=\inf_{a\in A}
||x-a||$. Trouver $d(0,\Rr-\Qq)$, $d(\sqrt2,\Qq)$, $d(M,{\cal D})$ o\`u
$M=(x,y,z)\in\Rr^3$ et $\cal D$ est la droite de vecteur unitaire $(a,b,c)$.

\item Pour  $A, B\subset \Rr^n$ on d\'efinit $d(A,B)=\inf_{a\in A,b\in B}
||a-b||$. Trouver $d(A,B)$ lorsque $A$ est une branche de l'hyperbole
$\{(x,y)\in{\Rr^2}\ ;\  xy=1\}$ et $B$ une asymptote.

\item On d\'efinit $\hbox {diam} A=\sup_{a,b\in A}||a-b||$. Quel est $\hbox
{diam} (]0,1[\cap\Qq)$ ? $\hbox{diam} ([0,1]\cap\Rr-\Qq)$ ?
\end{enumerate}


\finenonce{002340}


\finexercice


\exercice{2341, queffelec, 2003/10/01}

\enonce{002341}{}
Montrer que tout ouvert de $\Rr$ est union d\'enombrable d'intervalles ouverts
deux \`a deux disjoints. 
(\emph{Indication :} si $x\in O$ ouvert, consid\'erer $J_x$ qui est l'union des
intervalles ouverts inclus dans $O$ et  contenant $x$). \'Enoncer un résultat similaire pour
 les ouverts de $\Rr^n$.
\finenonce{002341}


\finexercice


\exercice{2342, queffelec, 2003/10/01}
\enonce{002342}{}
On va montrer que l'ensemble $D$ des r\'eels de la forme $p+q\sqrt 2$ o\`u $p$ et
$q$ d\'ecrivent ${\Zz}$, est dense dans ${\Rr}$.
\begin{enumerate}
\item Remarquer que $D$ est stable par addition et multiplication.

\item Posons $u=\sqrt 2 -1$; montrer que pour tous $a<b$, on peut trouver $n\geq 1$
tel que $0<u^n<b-a$, puis $m\in \Zz$ v\'erifiant $a<mu^n<b$.

 En d\'eduire le r\'esultat.
\end{enumerate}

\finenonce{002342}


\finexercice


\exercice{2343, queffelec, 2003/10/01}
\enonce{002343}{}
 Montrer que dans tout espace m\'etrique $(E,d)$ une boule ferm\'ee est un ferm\'e,
mais que l'adh\'erence d'une boule ouverte $B(a,r)$ ne coincide pas
n\'ecessai\-rement avec la boule ferm\'ee $B'(a,r)$ (on pourra consid\'erer dans 
$(\Rr^2,||.||_\infty)$, $E=[0,1]\times\{0\}\cup\{0\}\times[0,1]$ et la boule
centr\'ee en
$({1\over2},0)$ de rayon $1/2$).
\finenonce{002343}


\finexercice


\exercice{2344, queffelec, 2003/10/01}
\enonce{002344}{}
 $(E,||.||)$ un espace vectoriel norm\'e.
\begin{enumerate}
\item Montrer que dans ce cas la boule ferm\'ee $B'(a,r)$  est l'adh\'erence
de la boule ouverte $B(a,r)$.

\item Montrer que $\overline B(a,r)\subset \overline B(b,R) \Longleftrightarrow$
$r\leq R$ et
$||a-b||\leq R-r$.
\end{enumerate}
\finenonce{002344}


\finexercice


\exercice{2345, queffelec, 2003/10/01}
\enonce{002345}{}
\begin{enumerate}
\item Si $(x,y)\in\Rr^2$, on pose $||(x,y)||=\max(|x+y|,|x-2y|)$. Montrer qu'il
s'agit d'une norme sur $\Rr^2$ et dessiner sa boule unit\'e ferm\'ee.

\item Soit $p,q$ deux normes sur $\Rr^n$, $B_p$ et $B_q$ leurs boules unit\'es
ferm\'ees. Montrer que 
$$ B_q\subset B_p\Longleftrightarrow p\leq q.$$

Que signifie $ {1\over2}B_p\subset B_q\subset 2B_p$ ? Exemples.
\end{enumerate}
\finenonce{002345}


\finexercice


\exercice{2346, queffelec, 2003/10/01}
\enonce{002346}{}
On note $X=l^{\infty}$ l'espace des suites r\'eelles born\'ees, et $Y=c_0$
l'espace des suites r\'eelles tendant vers $0$, tous deux munis de la
m\'etrique (\`a v\'erifier) $d(x,y)=\sup_n\vert x(n)-y(n)\vert$. Montrer que $Y$
est ferm\'e dans $X$.  Montrer que l'ensemble des suites nulles \`a
partir d'un certain rang est dense dans
$Y$ mais pas dans $X$.
\finenonce{002346}


\finexercice


\exercice{2347, queffelec, 2003/10/01}
\enonce{002347}{}
Soit $E=\{f\in C^1([0,1],{\Rr})\ ;\ f(0)=0\}$. On pose 
$$||f||=\sup_{0\leq x\leq 1}|f(x)+f'(x)|,\ \hbox{et}\ N(f)=\sup_{0\leq x\leq
1}|f(x)|+\sup_{0\leq x\leq 1}|f'(x)|.$$
Montrer que ce sont deux normes \'equivalentes sur $E$. 
\finenonce{002347}


\finexercice


\exercice{2348, queffelec, 2003/10/01}
\enonce{002348}{}
 On d\'esigne par $d(a,b)$ la distance euclidienne usuelle de $a,b\in\Rr^2$ et
on pose

$$ \delta(a,b)= \left \{\begin{array}{ccc}
&d(a,b)\enskip &\hbox{si}\enskip
a,b\enskip
\hbox{sont align\'es avec l'origine}\enskip O\\
&d(0,a)+d(0,b)\enskip & \hbox{sinon}
\end{array}\right.$$
\begin{enumerate}
\item Montrer que $\delta$ est une distance sur $\Rr^2$ (``distance SNCF") plus
fine que la distance usuelle. 

 Dans
la suite, on suppose $\Rr^2$ muni de la topologie associ\'ee \`a $\delta$.

\item Soit $H$ le demi-plan $\{(x,y)\ ;\ y>0\}$~; montrer que $H$ est 
un ouvert ; d\'eterminer $\overline H$.

\item Quelle est la topologie induite sur une droite vectorielle; sur
le cercle unit\'e
$\Gamma
$~? 

\item Lesquelles des transformations suivantes sont continues~: homoth\'eties de centre $O$~;
rotations de centre $O$~; translations~?
\end{enumerate}
\finenonce{002348}



\finexercice


\exercice{2349, queffelec, 2003/10/01}
\enonce{002349}{}
\begin{enumerate}
\item Montrer que 
$||f||_\infty= \sup_{0\leq x\leq
1}|f(x)|$ et  $||f||_1=\int_0^1|f(t)|\ dt$ sont deux normes 
sur $C([0,1],{\Rr})$. Sont-elles \'equivalentes ?

\item Les deux m\'etriques associ\'ees sont-elles topologiquement \'equivalentes ?
\end{enumerate}
\finenonce{002349}


\finexercice


\exercice{2350, queffelec, 2003/10/01}
\enonce{002350}{}
Soit $E=C^1([0,1],{\Rr})$. Comparer les normes 
$N_1(f)=||f||_\infty,\  N_2(f)=||f||_\infty+||f||_1,\
N_3(f)=||f'||_\infty+||f||_\infty,\  N_4(f)=||f'||_1+||f||_\infty.$
\finenonce{002350}


\finexercice


\exercice{2351, queffelec, 2003/10/01}
\enonce{002351}{}
Soit $(x_n)$ une suite d'un espace topologique $X$ séparé; on note $A$
l'ensemble $\{x_1,x_2,\ldots\}$.
\begin{enumerate}
\item Toute valeur d'adhérence $a$ de la suite est un point de $\overline
A$ : donner un exemple où $a$ est un point isolé de $A$; un exemple où $a$ est
un point d'accumulation dans $A$; un exemple où $a$ est
un point d'accumulation dans $\overline A\backslash A$.

\item Montrer que tout point d'accumulation de $A$ est valeur d'adhérence de la
suite.
\end{enumerate}
\finenonce{002351}

\finexercice


\exercice{2352, queffelec, 2003/10/01}
\enonce{002352}{}
Soit $\R^n$ considéré comme groupe additif muni de sa topologie usuelle.
Soit $G$ un sous-groupe de $\R^n$.
\begin{enumerate}
\item On suppose que $0$ est isolé dans $G$. Montrer que tout point est isolé,
que $G$ est discret et fermé dans $\R^n$.

 On se restreint maintenant au cas $n=1$.

\item Montrer qu'alors, $G$ est soit $\{0\}$, soit de la forme
$a\Zz$, $a>0$.

\item Montrer que si $0$ est point d'accumulation, $G$ est partout dense dans $\R$. 
En déduire ainsi les sous-groupes fermés de $\R$.

\item On considère $\alpha\notin \Qq$; montrer que ${\Zz}+\alpha{\Zz}$ est un
sous-groupe dense de $\Rr$. En déduire les valeurs d'adhérence de la suite
$(e^{2i\pi n\alpha})_{n\in \Zz}$.
\end{enumerate}
\finenonce{002352}

\finexercice


\exercice{2418, drutu, 2007/10/01}
\enonce{002418}{} 
Soit $X=\{ a,b,c,d \}$. Lesquelles parmi les collections de
sous-ensembles suivants d\'eterminent une topologie sur $X$ ?
Justifier.

\begin{enumerate}
  \item $\emptyset$, $X$, $\{a\}$, $\{ b\}$, $\{a,c\}$, $\{a,b,c\}$,
  $\{a,b\}$;

  \item $\emptyset$, $X$, $\{a\}$, $\{ b\}$, $\{a,b\}$, $\{
  b,d\}$;

   \item $\emptyset$, $X$, $\{a,c,d\}$, $\{b,c,d\}$.
\end{enumerate}
\finenonce{002418}


\finexercice


\exercice{2419, drutu, 2007/10/01}
\enonce{002419}{} 
Soit $\R$ et soit $\mathcal{T}$ une collection de sous-ensembles de $\R$
contenant $\emptyset$, $\R$ et tous les complementaires
d'ensembles finis. Est-ce une topologie sur $\R$ ? Est-ce une
topologie s\'epar\'ee ?
\finenonce{002419}


\finexercice


\exercice{2420, drutu, 2007/10/01}
\enonce{002420}{} 
On appelle \textit{base} d'une topologie $\mathcal{T}$ un sous-ensemble
$\mathcal{B}$ de $\mathcal{T}$ tel que tout ouvert $\mathcal{O} \in \mathcal{T}$ s'\'ecrit comme
$\mathcal{O}=\bigcup_{i\in I} B_i$, o\`u $B_i \in \mathcal{B}$ pour tout $i\in I$.

\begin{enumerate}
  \item Montrer que $\mathcal{B}$ est une base de $\mathcal{T}$ si et
  seulement si pour tout ouvert $\mathcal{O}$ et tout point $x\in \mathcal{O}$ il
  existe un $B\in \mathcal{B}$ tel que $x\in B\subset \mathcal{O}$.

  \item Soit $\mathcal{T}_n$ la topologie sur $\R^n$ induite par la
  m\'etrique euclidienne
  $$\mathrm{dist} (\bar{x}, \bar{y})=\sqrt{(x_1-y_1)^2+...+(x_n-y_n)^2}\, .
  $$ Montrer que l'ensemble $\mathcal{B}$ de boules ouvertes ayant leur centre dans
  $\Q^n$et leur rayon dans $\Q$ est une base de $\mathcal{T}_n$.

  \item Soit $\mathcal{B}'$ l'ensemble de parallelipip\`edes ouverts dans
  $\R^n$ dont les ar\^etes sont parall\`eles aux axes de
  coordonn\'ees. Est-ce que $\mathcal{B}'$ est une base de $\mathcal{T}_n$ ?
 
  \item Est-ce que $\{ ]-\infty ,a[\; ;\; a\in \R  \}\cup \{ ]b, +\infty[\; ;\; b\in \R
  \}$ est une base pour $\mathcal{T}_1$ ?

  \item Pour tout $a\in \Q$ on note par $\delta_a$ la droite d'\'equation $y=ax$
   dans $\R^2$, et on note par $Y$ la r\'eunion des droites $\delta_a$. Soit $\mathcal{T}$ la
  topologie sur $Y$ induite par la topologie sur $\R^2$ et soit
  $\mathcal{T}'$ la topologie de base $\mathcal{B}'$ compos\'ee par tous les
  segments ouverts
   $]M,N[\subset \delta_a$, $O\not\in ]M,N[$, et par toutes les reunions $\bigcup_{a\in \Q, O\in ]M_a,N_a[} ]M_a,N_a[$. Les
  deux topologies $\mathcal{T}$ et $\mathcal{T}'$ sont-elles \'equivalentes ?
\end{enumerate}
\finenonce{002420}


\finexercice


\exercice{2421, drutu, 2007/10/01}
\enonce{002421}{} 
Soit $X$ un espace muni d'une m\'etrique $\mathrm{dist} : X\times X \to
\R_+$.
\begin{enumerate}
  \item Montrer que si $f:\R_+ \to R_+$ est une fonction
  croissante telle que $f(0)=0$ et $f(x+y)\leq f(x) +f(y)$ alors
  $\mathrm{dist}_f (x,y)=f(\mathrm{dist} (x,y))$ est une m\'etrique sur $X$.
  \item Montrer que
$$
\mathrm{dist}'(x,y)=\frac{\mathrm{dist}(x,y)}{1+\mathrm{dist} (x,y)}\, ,\; \forall x,y,
$$ est une m\'etrique sur $X$.
  \item Montrer que les m\'etriques $\mathrm{dist}$ et $\mathrm{dist}'$ sont
  topologiquement \'equivalentes.
\end{enumerate}
\finenonce{002421}


\finexercice


\exercice{2422, bodin, 2007/10/01}
\enonce{002422}{} 
%Abou-Jaoudé, Chevalier p76
Soit $(E,d)$ un espace métrique. On dit que $d$ est \emph{ultramétrique} si elle vérifie:
$$\forall (x,y,z)\in E^3 \qquad d(x,z) \le \sup \left( d(x,y), d(y,z)\right).$$
Cette inégalité entraine évidemment l'inégalité triangulaire.

\begin{enumerate}
\item Montrer que $E$ muni de la distance $d$ définie par 
$$d(x,y) = 1 \ \text{ si }\  x\neq y, \qquad d(x,x)=0$$
est un espace ultramétrique.

\medskip
On suppose maintenant que $(E,d)$ est ultramétrique.
  \item Montrer que si $d(x,y)\neq d(y,z)$, on a $ d(x,z) = \sup \left( d(x,y), d(y,z)\right)$.
  
  \item Montrer qu'une boule ouverte (resp. fermée) est une partie à la fois ouverte et fermée.

  \item Montrer que si deux boules ont un point commun l'une est contenue dans l'autre. Montrer de plus que si ces boules ont même rayon et sont toutes les deux des boules ouvertes (resp. fermées) elles sont confondues.

  \item Montrer que si deux boules ouvertes distinctes $B_1, B_2$  de rayon $r$ sont contenues dans une boule fermée de même rayon, alors leur distance est égale à $r$: 
$$d(B_1,B_2) := \inf_{(a,b)\in B_1\times B_2}d(a,b) = r.$$
\end{enumerate}
\finenonce{002422}


\finexercice


\exercice{2423, bodin, 2007/10/01}
\enonce{002423}{} 
%Abou-Jaoudé, Chevalier p75
Soit $p$ un nombre premier. Pour $n\in \Nn$ on définit $\nu(n)$ comme étant l'exposant de $p$ dans la décomposition de $n$ en facteurs premiers.
Pour $x = \pm \frac ab$, ($a,b\in \Nn^*$), on définit $\nu(x) = \nu(a)-\nu(b)$.
\begin{enumerate}
\item Montrer que $\nu(x)$ est indépendant du choix de la représentation $\pm \frac ab$.
  
  \item Montrer que $\nu(xy) = \nu(x)+\nu(y)$, $x,y\in \Qq$.

  \item Montrer que $\nu(x+y) \ge \min (\nu(x),\nu(y))$ pour $x,y\in \Zz$, puis pour $x,y\in \Qq$.

  \item Montrer que sur $\Qq$, $d$ définie par :
$$d(x,y) = p^{-\nu(x-y)} \ \text{ si } \ x\neq y, \qquad d(x,x)=0$$
est une distance ultramétrique.
\end{enumerate}
\finenonce{002423}


\finexercice

\exercice{2424, lescure, 2007/10/01}

\enonce{002424}{} 
\begin{enumerate}
\item  Soit $E$ un espace m\'etrique et $A \subset E$ une de ses parties. 
On désigne par $\overline A$ l'adhérence de $A$ et par $\mathrm{Fr} (A)$ la frontière  de $A$ dans $E.$ 
On a $\mathrm{Fr} (A) = \overline{A} \cap \overline{A^c}$.

\begin{enumerate}
\item Montrez que $x \in\mathrm{Fr} (A),$ si et seulement si il existe une 
suite $(x_n)$ d'éléments de $A$ et une suite $(y_n)$ d'éléments du complémentaire
$E \setminus A$ de $A$ dans $E,$ qui convergent l'une et l'autre vers $x.$

\item Soit $E=]-\infty,-1]\cup [0,1[ \cup [2,+\infty[$ muni de la topologie induite
par $\Rr.$ Avec $A=[0,\frac 12],$ qu'elle est la frontière de $A$ dans $E.$ Considérée
 comme sous-partie de $\Rr,$ qu'elle serait la frontière de $A$ dans $\Rr$?
\end{enumerate}

\item 
Soient $E$ et $F$ deux espaces métriques respectivement au moyen des distances $d$ et
$d'$.
\begin{enumerate}
\item  Précisez ce que l'on entend par la distance $\sup(d,d')$ sur $E \times F.$ Dîtes
rapidement pourquoi cette distance définit sur $E \times F$ le produit des topologies
métriques sur $E$ et $F.$ 
\item Soient $A\subset E$ et $B\subset F$. Montrez que l'intérieur $A\times B \setminus \mathrm{Fr} (A\times B)$ 
de $A\times B$ dans $E \times F$ est le produit cartésien de l'intérieur 
$A  \setminus \mathrm{Fr} (A)$ de $A$ dans $E$ avec l'intérieur $ B \setminus \mathrm{Fr} ( B)$ de $B$
dans $F.$
\end{enumerate}

\item 
 $E$ et $F$ sont toujours comme dans la deuxième question \c{c}i dessus.
\begin{enumerate}
\item Si $( \xi_n, \xi'_n)$ est une suite de points dans le complémentaire $E\times F \setminus A\times B$ de
$A\times B$ dans $E \times F,$ montrez qu'au moins une des deux alternatives suivantes (i) ou (ii) est vérifiée:

(i) il existe une suite extraite $\xi_{n_k}$ dont tous les termes sont dans $E \setminus A.$
  
(ii) il existe une suite extraite $\xi'_{n_k}$ dont tous les termes sont dans $F \setminus B.$

\item Déduire, de tout ce qui précède, que la frontière $\mathrm{Fr} (A \times B)$ de $A \times B$ dans
$E \times F$ est donnée par la formule:
$$\mathrm{Fr} (A \times B)=\bigl( \mathrm{Fr} (A) \times \overline B \bigr) \cup 
\bigl( \overline A \times\mathrm{Fr} ( B) \bigr)$$
\end{enumerate}

\item 
 Supposons $E$ et $F$ comme \c{c}i dessus mais avec l'hypothèse supplémentaire d'être
connexes, et avec des inclusions strictes $A \subset E$ et $B \subset F.$
\begin{enumerate}
\item  Soient, dans $E\times F,$ les points  $(x,x')\notin A \times B$ et $(y,y') \notin A \times B.$
Supposons que $x \in A$ et $y \notin A;$  Montrez qu'il existe une partie connexe entièrement contenue
dans le complémentaire de $A \times B$ qui contient $(x,x')$ et $(y,y')$. 
\item En déduire, sous les présentes hypothèses de cette quatrième question, que le complémen\-taire
de  $A \times B$ dans $E \times F$ est connexe.
\end{enumerate}
\end{enumerate}
\finenonce{002424}


\finexercice


\exercice{6034, queffelec, 2011/10/16}
\enonce{006034}{}
\begin{enumerate}
\item Soit $X=\{0,1\}$ muni de la famille d'ouverts $\{\emptyset, \{0\}, X\}$.
Cette topologie est-elle séparée?

\item Soit $X$ un ensemble non vide. Décrire la topologie dont les singletons
forment une base d'ouverts.

\item Décrire la topologie sur ${\Rr}$ dont la famille des
intervalles fermés forme une base d'ouverts; même question avec les intervalles
ouverts symétriques.

\item Soit $X$ un ensemble infini. Montrer que la famille d'ensembles constituée
de l'ensemble vide et des parties de $X$ de complémentaire fini définit une
topologie sur $X$.
\end{enumerate}
\finenonce{006034}



\finexercice

\exercice{6035, queffelec, 2011/10/16}
\enonce{006035}{}
Soit $X$ un espace topologique, et $f$ une application quelconque de $X$ dans
un ensemble $Y$. On dit qu'une partie $A$ de $Y$ est ouverte, si $f^{-1}(A)$
est un ouvert de $X$. Vérifier qu'on a défini ainsi une topologie sur $Y$.
\finenonce{006035}



\finexercice


\exercice{6036, queffelec, 2011/10/16}
\enonce{006036}{}
Montrer qu'on peut construire sur ${\Rr}\cup \{\infty\}$ une topologie
séparée en prenant comme ouverts,
les ouverts de  ${\Rr}$ et les ensembles de la forme
$\{x/\vert x\vert>a\}\cup \{\infty\}$ où $a$ est réel. Comment construire une
topologie séparée sur ${\Rr}\cup \{+\infty\}\cup \{-\infty\}$?
\finenonce{006036}



\finexercice


\exercice{6037, queffelec, 2011/10/16}
\enonce{006037}{}
\label{exn7}
Soit $X$ un ensemble non vide et $\Sigma$ une famille de parties de $X$ stable
par intersection finie et contenant $X$. Montrer que la plus petite topologie
$\cal T$ conte\-nant $\Sigma$ (la topologie engendrée par $\Sigma$) est
constituée des unions d'ensembles de $\Sigma$, ou, de fa\c con équivalente, 

$$A\in{\cal T}\Longleftrightarrow \forall x\in A\ \exists S\in \Sigma\ ;\ x\in
S\subset A.$$

Montrer que l'on peut affaiblir l'hypothèse de stabilité par intersection finie
en :

$(*)\qquad\forall S_1, S_2\in\Sigma,\ \forall x\in S_1\cap S_2,\ \exists
S_3\in\Sigma\ ;\ x\in S_3\subset S_1\cap S_2.$
\finenonce{006037}



\finexercice


\exercice{6038, queffelec, 2011/10/16}
\enonce{006038}{}
Soit $C$ l'ensemble des fonctions continues réelles sur $[0,1]$. Pour toute
$f\in C$ et $\varepsilon>0$ on définit
$$M(f,\varepsilon)=\{g/\int_0^1\vert f-g\vert<\varepsilon\}.$$
Montrer que la famille {\cal M} des ensembles $M(f,\varepsilon)$ lorsque $f\in
C$ et  $\varepsilon>0$ est une base de topologie. Même question avec la famille
$$U(f,\varepsilon)=\{g/\sup_x\vert f(x)-g(x)\vert<\varepsilon\}.$$
\finenonce{006038}



\finexercice       


\exercice{6039, queffelec, 2011/10/16}
\enonce{006039}{}
$U$ dans ${\Nn}$ est dit ouvert s'il est stable par divisibilité, c.a.d.
tout diviseur de $n\in U$ est encore dans $U$.
Montrer qu'on a défini ainsi une topologie sur ${\Nn}$ qui n'est pas la
topologie discrète.
\finenonce{006039}



\finexercice       


\exercice{6040, queffelec, 2011/10/16}
\enonce{006040}{}
On considère dans ${\Nn}^*$, la famille de progressions arithmétiques
$$P_{a,b}=\{a+bn/n\in {\Nn}^*\},$$ 
où $a$ et $b$ sont deux entiers premiers entre eux.

\begin{enumerate}
\item Montrer que l'intersection de deux telles progressions est soit vide,
soit une progression arithmétique de même nature, plus
précisément,
$$P_{a,b}\cap P_{a',b'}=P_{\alpha,\beta}$$
où $\alpha$ est le minimum de l'ensemble $P_{a,b}\cap P_{a',b'}$, et
$\beta=\mathrm{ppcm}(b,b')$.

\item En déduire que cette famille d'ensembles (en y adjoignant $\emptyset$) 
forme une base de topologie sur ${\Nn}^*$ dont on décrira les ouverts.

\item Montrer que cette topologie est séparée. 
\end{enumerate}
\finenonce{006040}



\finexercice 


\exercice{6041, queffelec, 2011/10/16}
\enonce{006041}{}
\begin{enumerate}
\item Montrer que si $B$ est un ouvert de l'espace topologique $X$ et $A\cap
B=\emptyset$, alors $\overline A\cap B=\emptyset$, mais que $\overline A\cap
\overline B$ n'est pas nécessairement vide.

\item Montrer à l'aide d'exemples que l'égalité $\overline{\cup_i
A_i}=\cup_i\overline{A_i}$ n'a pas lieu en général pour une infinité d'indices.
\end{enumerate}
\finenonce{006041}



\finexercice


\exercice{6042, queffelec, 2011/10/16}
\enonce{006042}{}
Déterminer l'adhérence et l'intérieur des ensembles suivants :

 ${\Qq}$;\ \ ${\Rr}\backslash{\Qq}$;\ \ $\{(x,y)\in {\Rr^2}\ /\  
0<x<1, y=0\}$;\ \ $\{(x,y,z)\in {\Rr^3}\ /\  x=0\}$
$\{{1\over n}, n \geq 1\}$;\ \ le cercle unité de $\Rr^2$.
\finenonce{006042}



\finexercice


\exercice{6043, queffelec, 2011/10/16}
\enonce{006043}{}
Si $A$ est une partie de l'espace topologique $X$, on pose
$\alpha(A)= \stackrel{\circ}{\overline{A}}$ et
$\beta(A)=\overline{ \stackrel{\circ}{A}}$.
\begin{enumerate}
\item Montrer que $\alpha$ et $\beta$ sont des applications croissantes pour
l'inclusion de ${\cal P}(X)$ dans ${\cal P}(X)$.

\item Montrer que si $A$ est ouvert, $A\subset \alpha(A)$ et si $A$ est fermé,
$\beta(A)\subset A$. En déduire que $\alpha^2=\alpha$ et $\beta^2=\beta$.

\item Construire $A\subset\Rr$ tel que les cinq ensembles  :

$A$, $\overline{A}$, $\stackrel{\circ}{A}$, $\alpha(A)$, $\beta(A)$
soient tous distincts.  
\end{enumerate}
\finenonce{006043}



\finexercice


\exercice{6044, queffelec, 2011/10/16}
\enonce{006044}{}
Déterminer l'adhérence dans ${\Rr}^2$ du graphe 
$$G=\{(x,y)/ y=\sin {1\over x}, 0<x\leq 1\}.$$
\finenonce{006044}



\finexercice


\exercice{6045, queffelec, 2011/10/16}
\enonce{006045}{}
Dans un espace topologique, on définit la frontière d'une partie $A$
comme étant $\partial A =\overline A \ \backslash \stackrel{\circ}{A}$.
\begin{enumerate}
\item Montrer que $\partial A =\partial (A^c) $ et que $A=\partial A
\Longleftrightarrow$ $A$ fermé d'intérieur vide. 

\item Montrer que $\partial (\overline A) $ et $\partial (\stackrel{\circ}{A}) $
sont toutes deux incluses dans
$\partial A $, et donner un exemple où ces inclusions sont strictes. 

\item Montrer que $\partial(A\cup B)\subset \partial A \cup \partial B $, et que
l'inclusion peut
être stricte; montrer qu'il y a égalité lorsque $\overline A\cap \overline B
=\emptyset$ (établir $\stackrel{\circ}{A\cup B}\subset\stackrel{\circ}{A} \cup
\stackrel{\circ}{B}$).

Montrer que $\stackrel{\circ}{A\cup B}=\stackrel{\circ}{A} \cup
\stackrel{\circ}{B}$ reste vrai lorsque $\partial A \cap \partial B =\emptyset$
(raisonner par l'absurde).
\end{enumerate}
\finenonce{006045}



\finexercice


\exercice{6046, queffelec, 2011/10/16}
\enonce{006046}{}
\begin{enumerate}
\item Soit $X$ un espace topologique, et $D$ un sous-ensemble (partout) dense dans
$X$. Montrer qu'il est aussi équivalent de dire 

(i) Le complémentaire de $D$ est d'intérieur vide.

(ii) Si $F$ est un fermé contenant $D$, alors $F=X$.

(iii) $D$ rencontre tout ouvert non vide de $X$.

Montrer qu'un ensemble $A\subset X$ rencontre toute partie dense dans $X$ si et
seulement si il est d'intérieur non vide.

\item Soit $E$ et $G$ deux ouverts denses dans $X$; montrer que $E\cap G$ est encore
dense dans $X$. En déduire que toute intersection dénombrable d'ouverts denses
est une intersection décroissante d'ouverts denses.
\end{enumerate}
\finenonce{006046}



\finexercice


\exercice{6047, queffelec, 2011/10/16}
\enonce{006047}{}
\label{gijsexopr}
Etablir les propriétés suivantes de l'adhérence d'un ensemble dans un
espace topologique :

\begin{enumerate}  
\item $\overline{\overline A}=\overline A$

\item Si $A\subset B$ alors $\overline A\subset \overline B$.

\item $\overline{A\cup B}=\overline A\cup \overline B$
\end{enumerate}
Montrer que la formule $\overline{A\cap B}=\overline A\cap \overline B$ n'est
pas vraie en général; montrer que 3. n'est pas vrai en général
pour une infinité d'ensembles.
\finenonce{006047}



\finexercice       


\exercice{6048, queffelec, 2011/10/16}
\enonce{006048}{}
Etablir l'équivalence entre les propriétés suivantes :
\begin{enumerate}  
\item $\stackrel{\circ}{A}$ est le plus grand ouvert contenu dans $A$.

\item $a\in\stackrel{\circ} A$ si et seulement si il existe un voisinage de $a$
entièrement contenu dans $A$.
\end{enumerate}
Etablir pour l'intérieur d'un ensemble des propriétés analogues à
celles de l'exercice \ref{gijsexopr}.
\finenonce{006048}



\finexercice       


\exercice{6049, queffelec, 2011/10/16}
\enonce{006049}{}
On rappelle la construction de l'ensemble triadique de Cantor : on part du
segment $[0,1]$ dont on supprime l'intervalle médian $]{1\over 3}, {2\over
3}[$; à la deuxième étape, on supprime les intervalles  $]{1\over 9},
{2\over 9}[$ et $]{7\over 9}, {8\over 9}[$ etc. On note $K_n$ la réunion des
intervalles restants à la $n$-ième étape, et $K=\bigcap K_n$. Quelle est
l'adhérence et l'intérieur de $K$?
\finenonce{006049}



\finexercice       


\exercice{6050, queffelec, 2011/10/16}
\enonce{006050}{}
 Soit $X$ un espace topologique, et $D$ un sous-ensemble dense dans $X$.
Montrer qu'il est aussi équivalent de dire 

\begin{enumerate}  
\item Le complémentaire de $D$ est d'intérieur vide.

\item Si $F$ est un fermé contenant $D$, alors $F=X$.

\item  $D$ rencontre tout ouvert de $X$.
\end{enumerate}
Montrer qu'un ensemble $A\in X$ rencontre toute partie dense dans $X$ si et
seulement si il est d'intérieur non vide.

\finenonce{006050}



\finexercice       


\exercice{6051, queffelec, 2011/10/16}
\enonce{006051}{}
Soit $E$ et $G$ deux ouverts denses dans $X$ ; montrer que $E\cap G$ est encore
dense dans $X$.
\finenonce{006051}



\finexercice       


\exercice{6052, queffelec, 2011/10/16}
\enonce{006052}{}
Soit $f$ une application de ${\Rr}$ dans ${\Rr}$ telle que pour tout $a>0$,
l'ensemble des $x$ vérifiant $\vert f(x)\vert>a$ est fini.
Montrer que $\{x/f(x)=0\}$ est dense dans ${\Rr}$. Le vérifier sur
l'exemple suivant : on énumère les rationnels $r_1, r_2, r_3,\cdots,
r_n,\cdots$ et on pose $f(r_n)={1\over n}$ si $n\geq 1$, $f(x)=0$ ailleurs.
\finenonce{006052}



\finexercice       


\exercice{6053, queffelec, 2011/10/16}
\enonce{006053}{}
Montrer que $\{\sqrt n - E(\sqrt n), n\geq 1\}$ est dense dans $[0,1]$, où
$E(x)$ désigne la partie entière de $x$.
\finenonce{006053}



\finexercice  


\exercice{6054, queffelec, 2011/10/16}
\enonce{006054}{}
Soit $E$ un ensemble non vide, et $X=E^{\Nn}$ l'ensemble des suites $x=(x_n)$
d'éléments de $E$. Pour $x,y\in X$, on pose $p(x,y)=\min\{n/x_n\neq y_n\}$ si
$x\neq y$, et $\infty$ si $x=y$.
\begin{enumerate}
\item Montrer que $d(x,y)={1\over p(x,y)}$ (avec ${1\over\infty}=0$) est une
distance sur $X$ qui vérifie l'inégalité ultramétrique
$$d(x,z)\leq\max(d(x,y),d(y,z)).$$

\item Quelles sont les boules ouvertes et les boules fermées pour cette métrique ?
\end{enumerate}
\finenonce{006054}



\finexercice

\exercice{6060, queffelec, 2011/10/16}
\enonce{006060}{Espace quasi-séparé}

Soit $(X,\cal T)$ un espace topologique. 
\begin{enumerate}
\item Montrer que les conditions suivantes
sont équivalentes : 

\quad (i) $\forall x,y\in X,\ x\not=y,\ \exists V\ \hbox{voisinage de}\ x\ ; \
y\notin V$.

\quad (ii) $\forall x\in X,\ \{x\}$ est fermé.

\quad (iii) $\forall x\in X,\ \cap\ \{V\ ; \ V\ \hbox{voisinage de}\
x\}=\{x\}$.

\item Soit $(X,\cal T)$ ainsi et $A\subset X$ tel que $\overline A\not=A$.
Montrer que si $x\in\overline A\backslash A$, tout voisinage de $x$ coupe
$A$ en une infinité de points.
\end{enumerate}
\finenonce{006060}



\finexercice


\exercice{6061, queffelec, 2011/10/16}
\enonce{006061}{Exemple de topologie non séparée}
Dans $\Cc$, on note $[z_0\rightarrow[$ la demi-droite $\{\rho e^{i\theta_0}\
;\ \rho\geq\rho_0\}$, si $z_0=\rho_0 e^{i\theta_0}$. On déclare ouvert toute
réunion (éventuellement vide) de telles demi-droites. 
\begin{enumerate}
\item Montrer qu'on a ainsi défini sur $\Cc$ une topologie $\cal T$ non séparée.

\item Montrer que l'adhérence du point $\{z_0\}$ pour cette topologie est
$[0,z_0]$.

\item En déduire que les fermés de $\cal T$ sont les ensembles étoilés par rapport
à $0$ ($A$ est dit ``étoilé par rapport à $0$" si, pour tout $z\in A$, le
segment $[0,z]$ est encore dans $A$).
\end{enumerate}



\finenonce{006061}


\finexercice


\exercice{6062, queffelec, 2011/10/16}
\enonce{006062}{}
Soit $(X,\cal T)$ un espace topologique séparé. Montrer que la diagonale
$\Delta$ de
$X\times X$ est fermée dans $X\times X$.
\finenonce{006062}



\finexercice


\exercice{6063, queffelec, 2011/10/16}
\enonce{006063}{}
\begin{enumerate}
\item Quels sont les ouverts de $[1,2]\cup\{3\}$ induits par ceux de ${\Rr}$?

\item Quelle est la topologie induite sur ${\Zz}$ par celle de ${\Rr}$?

\item Quels sont les ouverts du cercle $\Gamma=\{z/\vert z\vert=1\}$? du demi-plan
$\{z/\Im z >0\}$? du demi-plan $\{z/\Im z \geq0\}$ dans ${\Cc}$?
\end{enumerate}
\finenonce{006063}



\finexercice


\exercice{6064, queffelec, 2011/10/16}
\enonce{006064}{}
Soit $Y$ un sous-ensemble de l'espace topologique $X$, muni de la topologie
induite. Décrire les ouverts (fermés) induits de $Y$ lorsque $Y$ est ouvert
(fermé). 

Soit
$A\subset Y$. Montrer que l'adhérence de
$A$ dans
$Y$,
$\overline A^Y=Y\cap \overline A$; a-t-on pour l'intérieur de $A$ dans $Y$,
$\stackrel\circ {A^ Y} = Y\cap \stackrel\circ A$?
\finenonce{006064}



\finexercice


\exercice{6065, queffelec, 2011/10/16}
\enonce{006065}{}
On dit qu'un espace topologique $X$ a la propriété (P) si la famille de
parties de $X$ qui sont à la fois ouvertes et fermées est une base pour les
ouverts de $X$.
\begin{enumerate}
\item Montrer qu'un espace topologique discret a cette propriété.

\item Montrer que la topologie induite sur
${\Qq}$ par la topologie usuelle de
${\Rr}$ n'est pas la topologie discrète, mais qu'elle possède aussi la
propriété (P).

\item Autre exemple ?
\end{enumerate}
\finenonce{006065}



\finexercice


\exercice{6082, queffelec, 2011/10/16}
\enonce{006082}{}
Si $A$ est une partie bornée d'un espace métrique $(E,d)$, on pose
$\hbox{diam} A=\sup_{a,b\in A}d(a,b)$.
\begin{enumerate}
\item Montrer que  $\hbox{diam} A=\hbox{diam} \overline A$.

\item Trouver le diamètre de $\{f\in C([0,1])\ ;\  0\leq f\leq 1\}$; de $\{f\in
C([0,1])\ ;\ 0\leq f\leq 1,\ f(0)=0\}$, $C$ étant muni de la métrique $d_1$.
\end{enumerate}
\finenonce{006082}



\finexercice


\exercice{6090, queffelec, 2011/10/16}
\enonce{006090}{}
Peut-on construire dans ${\Rr}$ un ensemble infini, fermé, constitué
uniquement d'irrationnels?
\finenonce{006090}



\finexercice       


\exercice{6091, queffelec, 2011/10/16}
\enonce{006091}{}
Montrer que sur ${\Rr}^n$, les distances $d$ euclidienne, $d_{\infty}$ et
$d_1$ définissent la même topologie.
\finenonce{006091}



\finexercice       


\exercice{6092, queffelec, 2011/10/16}
\enonce{006092}{}
\begin{enumerate}  
\item Dans ${\Rr}^2$, on considère $U={\Rr}^2\backslash\{(0,y)\in {\Rr}^2
/ y\geq 0\}$. Vérifier qu'il est ouvert et qu'il peut s'écrire comme une
union dénombrable de fermés (un tel ensemble est dit de type $F_\sigma$).

\item Dans ${\Rr}^n$, on considère le sous-ensemble des points à
coordonnées entières, et le sous-ensemble des points à
coordonnées rationnelles. Vérifier que le premier est fermé mais que le
second n'est ni ouvert ni fermé.
\end{enumerate}
\finenonce{006092}



\finexercice       


\exercice{6093, queffelec, 2011/10/16}
\enonce{006093}{}
Soit $M_n({\Rr})$ l'ensemble des matrices carrées d'ordre $n$, muni de la
distance $d(A,B)=\max_{i,j}\vert a_{i,j}-b_{i,j}\vert$ où $A=(a_{i,j})$ et
$B=(b_{i,j})$.

\begin{enumerate}  
\item Montrer que l'ensemble des matrices inversibles est un ouvert dense de
$M_n({\Rr})$.

\item Dans le cas $n=2$, décider si les ensembles suivants sont ouverts,
fermés, ni ouverts ni fermés :
 
 ${\cal A}$ = matrices ayant deux valeurs propres distinctes et $>0$.

 ${\cal B}$ = matrices ayant deux valeurs propres $>0$.
\end{enumerate}
\finenonce{006093}



\finexercice       


\exercice{6094, queffelec, 2011/10/16}
\enonce{006094}{}
On note $X$ l'espace des suites réelles $x=(x(n))$ et on le munit de
la topologie dont les ouverts élémentaires sont
$$V(x;n_1,n_2,\cdots n_k;\varepsilon)=\{y\in X /\vert
x(n_i)-y(n_i)\vert<\varepsilon,\  i=1\cdots k\}.$$
Vérifier qu'on a bien défini une base de topologie.

Comparer la topologie qu'elle engendre sur $l^{\infty}$ et $c_0$ avec la
topologie métrique de l'exercice précédent. 
\finenonce{006094}



\finexercice       


\exercice{6095, queffelec, 2011/10/16}
\enonce{006095}{}
\label{gijsexosep}
Soit $X$ un espace topologique. On considère les propriétés suivantes :
\begin{itemize}
  \item[(i)] $X$ contient un dénombrable dense.
  \item[(ii)] la topologie sur $X$ possède une base dénombrable d'ouverts.
\end{itemize}
Montrer que (ii) implique (i) et que la réciproque a lieu si $X$ est
métrisable. Un espace vérifiant (i) est dit séparable.
\finenonce{006095}



\finexercice       


\exercice{6096, queffelec, 2011/10/16}
\enonce{006096}{}
Soit $X$ un espace métrique séparable (cf exercice \ref{gijsexosep}), et $A$ une
partie quelconque de $X$. Montrer que $A$ est encore séparable.
\finenonce{006096}



\finexercice    


\exercice{6097, queffelec, 2011/10/16}
\enonce{006097}{}
On considère dans $\Rr$ les trois topologies 
${\cal T}_1, \  {\cal T}_2, \  {\cal T}_3$, engendrées respectivement par les
intervalles de la forme
$]a,b[, \  [a,b[, \  [a,b]$,
$a$ et $b$ décrivant $\Rr$. Comparer les topologies, et décrire
les fonctions continues de $({\Rr}, {\cal T}_1)$ dans $({\Rr},{\cal T}_2)$;
de $({\Rr}, {\cal T}_1)$ dans $({\Rr},{\cal T}_3)$.
\finenonce{006097}



\finexercice


\exercice{6098, queffelec, 2011/10/16}
\enonce{006098}{}
Soit $\cal T$ et $\cal T'$ deux topologies sur $X$. Montrer que $\cal T'$ est
plus fine que $\cal T$ ssi $(X,{\cal T'})\buildrel \hbox{id}\over{\to} (X,\cal
T)$ est continue. Montrer qu'alors $\overline A^{\cal T'}\subset
\overline A^{\cal T}$; quelle inclusion a-t-on entre $\stackrel\circ{ A^{\cal
T'}}$ et 
$\stackrel\circ{A^{\cal T}}$ ?
\finenonce{006098}



\finexercice


\exercice{6099, queffelec, 2011/10/16}
\enonce{006099}{}
Comparer sur $X=\{0,1\}^{\Nn^*}$, l'espace des suites de $0-1$, les
topologies définies par les distances 

\quad $d(x,y)=\displaystyle{1\over \min\{n/x_n\neq y_n\}}$ si
$x\neq y$,\quad $0$ sinon,

et

\quad $\delta(x,y)=\sup_n|x(n)-y(n)|$.
\finenonce{006099}



\finexercice


\exercice{6100, queffelec, 2011/10/16}
\enonce{006100}{}
On se donne une application $f:\Rr\to\Rr^n$, et on note $d$ la distance
euclidienne sur $\Rr^n$. A quelles conditions sur
$f$, $\delta(x,y)=d(f(x),f(y))$ définit-elle une distance sur $\Rr$
équivalente topologiquement à la distance usuelle (ie définissant la même
topologie.)?   
\finenonce{006100}



\finexercice       


\exercice{6101, queffelec, 2011/10/16}
\enonce{006101}{}
Soit $E$ un ensemble non vide, et $X=E^{\Nn}$ l'ensemble des suites $x=(x_n)$
d'éléments de $E$. Pour $x,y\in X$, on pose $p(x,y)=\min\{n/x_n\neq y_n\}$ si
$x\neq y$, et $\infty$ si $x=y$.

Montrer que $d(x,y)={1\over p(x,y)}$ (avec ${1\over\infty}=0$) est une distance
sur $X$ qui vérifie l'inégalité ultramétrique
$$d(x,z)\leq\max(d(x,y),d(y,z)).$$
\finenonce{006101}



\finexercice       


\exercice{6102, queffelec, 2011/10/16}
\enonce{006102}{}
On dit qu'une distance est \emph{ultramétrique} si elle vérifie l'inégalité
triangulaire renforcée :
$$d(x,z)\leq\max(d(x,y),d(y,z)).$$
Etablir les assertions suivantes :
\begin{enumerate}  
\item Si $d(x,y)\neq d(y,z)$, alors $d(x,z)=\max(d(x,y),d(y,z)).$ En déduire que
tout triangle dans $E$ est isocèle.

\item Toute boule ouverte $B(x,r)$ est un ensemble à la fois ouvert et fermé, et 
$$B(x,r)=B(y,r)\ \ \forall y\in B(x,r).$$

\item Toute boule fermée $B'(x,r)$ est un ensemble à la fois ouvert et fermé, et 
$$B'(x,r)=B'(y,r)\ \ \forall y\in B'(x,r).$$

\item Si deux boules ont un point commun, elles sont embo\^\i tées.
\end{enumerate}
\finenonce{006102}



\finexercice       


\exercice{6103, queffelec, 2011/10/16}
\enonce{006103}{}
Soit $(X,d)$ un espace métrique, et soit $\varphi$ une fonction réelle définie
pour $x\geq0$, vérifiant  (i) $\varphi(0)=0$, (ii) $\varphi$ croissante, (iii)
$\varphi(u)>0$ si $u>0$, (iv) $\varphi(u+v)\leq\varphi(u)+\varphi(v)$.


\begin{enumerate}  
\item Montrer que $\delta(x,y)=\varphi(d(x,y))$ définit une distance sur $X$.

\item Vérifier que les fonctions
$\varphi_1(u)=\inf(u,1)$, $\varphi_2(u)={u\over{1+u}}$,
$\varphi_3(u)=\log(1+u)$, et $\varphi_4(u)=u^\alpha$ où $0<\alpha<1$ remplissent
les conditions (i) (ii) et (iii); plus généralement, montrer que toute fonction
$f$ strictement croissante, concave, telle que
$f(0)=0$ remplit ces conditions.

\item  On suppose de plus que la fonction $\varphi$ est continue en $0$. Montrer
que les métriques $d$ et $\delta$ sont topologiquement équivalentes.

\item  Montrer que $\delta_1=\varphi_1(d)$ et $\delta_2=\varphi_2(d)$ sont
lipschitz-équivalentes.
\end{enumerate}
\finenonce{006103}



\finexercice       


\exercice{6104, queffelec, 2011/10/16}
\enonce{006104}{}
Soit $(X,d)$ un espace métrique avec métrique bornée. On note $\cal F$ l'ensemble
des fermés non vides de $X$, et on définit pour $A$ et $B$ dans $\cal F$,
$$\delta (A,B)=\Vert d_A-d_B\Vert_\infty$$
où $d_A$ est la fonction bornée $x\to d(x,A)$.
Montrer qu'on a défini ainsi une métrique sur $\cal F$, et que l'application
$a\to\{a\}$ est une isométrie de $X$ dans $\cal F$.
\finenonce{006104}



\finexercice 


\exercice{6105, queffelec, 2011/10/16}
\enonce{006105}{}
Trouver les valeurs d'adhérence de la  suite : 

\quad
$0,1,0,{1\over2},1,0,{1\over4},{1\over2},{3\over4},1,\cdots,0,{1\over2^k},
{2\over2^k},\cdots,{2^k-1\over2^k},1,0,...$
\finenonce{006105}



\finexercice


\exercice{6106, queffelec, 2011/10/16}
\enonce{006106}{}
\begin{enumerate}
\item Soit $(u_n)$ une suite réelle telle que $e^{iu_n}$ et $e^{i\sqrt2
  u_n}$ convergent. Montrer que $(u_n)$ a au plus une valeur d'adhérence.

\item Soit $(u_n)$ une suite réelle telle que $e^{itu_n}$ 
converge pour $t\in T$ où $T$ est non dénombrable. Même conclusion.
\end{enumerate}
\finenonce{006106}



\finexercice


\exercice{6107, queffelec, 2011/10/16}
\enonce{006107}{}
Soit $(u_n)$ une série positive divergente telle que $u_n$ décroit vers
$0$ et on pose $A=\{\pm u_1\pm u_2\pm\cdots\pm u_n,\ n\geq1\}$. Montrer
que $\overline A=\R$.
\finenonce{006107}



\finexercice


\exercice{6108, queffelec, 2011/10/16}
\enonce{006108}{}
Soit dans un espace métrique $(X,d)$ une suite $(x_n)$ telle que les trois
sous-suites $(x_{2n})$, $(x_{2n+1})$, et $(x_{3n})$ convergent. Montrer que la
suite elle-même converge.
\finenonce{006108}



\finexercice       


\exercice{6109, queffelec, 2011/10/16}
\enonce{006109}{}
Soit $(a_{m,n})_{(m,n)\in{\Nn}^2}$ une suite d'un espace métrique $(X,d)$. On
suppose que $\lim_{n\to\infty}a_{m,n}=a_m$, et que $\lim_{m\to\infty}a_m=a$.
Montrer qu'il existe une sous-suite de la suite initiale $(a_{p,n_p})$ telle que
$\lim_{p\to\infty}a_{p,n_p}=a$.
\finenonce{006109}



\finexercice       


\exercice{6110, queffelec, 2011/10/16}
\enonce{006110}{}
Soit $(F_n)$ une suite décroissante de fermés dans un espace topologique
$X$, et soit $(x_n)$ une suite convergente dans $X$ telle que pour chaque $n$,
$x_n\in F_n$. Vérifier que $\lim_{x\to\infty} x_n \in \bigcap F_n$.

Que peut-on dire si la suite de fermés n'est plus décroissante?
\finenonce{006110}



\finexercice       


\exercice{6111, queffelec, 2011/10/16}
\enonce{006111}{}
\label{gijsexopol}
On va montrer que les polyn\^omes sont denses dans les fonctions continues sur
$[-1,1]$. Pour commencer, on approche la fonction $\vert t\vert$.

\begin{enumerate}  
\item Montrer que la suite de polyn\^omes définis par récurrence : 
$$p_{n+1}(t) = p_n(t) + {1\over 2} (t^2-p_n^2(t)),\ \ p_0(t)=0,$$
converge vers $\vert t\vert$.

\item En déduire que toute fonction affine par morceaux sur
$[-1,1]$ est limite d'une suite de polyn\^omes.

\item Montrer que les polyn\^omes sont denses dans les fonctions continues sur
$[-1,1]$.
\end{enumerate}
\finenonce{006111}



\finexercice       


\exercice{6112, queffelec, 2011/10/16}
\enonce{006112}{}
\begin{enumerate}  
\item Déterminer l'ensemble des valeurs d'adhérence de la suite de réels
$x_n=(1+{1\over n})\sin(n{\pi\over6})$; de la suite $({1\over m}+{1\over
n})_{m\geq1,n\geq1}$.

\item Montrer que l'ensemble ${\Zz}\alpha+{\Zz}$ est dense dans ${\Rr}$ si
$\alpha$ est irrationnel. En déduire l'ensemble des valeurs d'adhérence de la
suite $x_n=\cos(2\pi n\alpha)$.

(\emph{Indication :} on pourra montrer que tout sous-groupe fermé de
$\Rr$ est soit $\Rr$, soit
discret, de la forme $a\Zz$.)
\end{enumerate}
\finenonce{006112}



\finexercice       


\exercice{6113, queffelec, 2011/10/16}
\enonce{006113}{}
On sait que l'ensemble des valeurs d'adhérence d'une suite réelle est un fermé
de $\Rr$. Montrer que tout fermé de $\Rr$ est l'ensemble des valeurs
d'adhérence d'une suite réelle : si $F$ est fini, trouver une suite qui
prend une infinité de fois chaque valeur de $F$; si
$F$ est infini, montrer que $F$ contient un dénombrable dense $D$ et trouver
une suite qui prend une infinité de fois chaque valeur de $D$.
\finenonce{006113}



\finexercice       


\exercice{6114, queffelec, 2011/10/16}
\enonce{006114}{}
Soit $(\epsilon_k)$ une suite à valeurs dans $\{-1,1\}$ et
$S_n=\sum_{k=0}^n\epsilon_k$. Montrer que l'ensemble des valeurs d'adhérence
de la suite $(S_n)$ est un intervalle de $\Zz$. 
\finenonce{006114}



\finexercice       


\exercice{6115, queffelec, 2011/10/16}
\enonce{006115}{}
On considère une suite $(x_n)$ de $[0,1]$
telle que $x_{n+1}-x_n$ tend vers $0$. 

\begin{enumerate}  
\item Montrer que l'ensemble $A$ de ses valeurs d'adhérence est un intervalle
fermé de $[0,1]$.

\item On suppose de plus que cette suite est une suite récurrente i.e. définie par $
x_{n+1}=f(x_n)$ où $f$ est continue de $[0,1]$ dans lui-même, et un point initial
$x_0\in [0,1]$. Montrer alors que la suite converge (on commencera par remarquer
que si
$x\in A$, alors $x=f(x)$, et que si $x_m\in A$ pour un indice $m$, alors la
suite converge.)

\item Soit $x=(x_n)$ une suite de $l^{\infty}$; montrer que l'ensemble des valeurs
d'adhé\-rence de la suite $y$ de terme général $y_n={{x_1+x_2+\cdots+x_n}\over
n}$ est un intervalle. En déduire que l'application $f$ de $l^{\infty}$ dans
lui-même qui associe $y$ à $x$, n'est pas bijective.
\end{enumerate}
\finenonce{006115}



\finexercice  


\exercice{6136, queffelec, 2011/10/16}
\enonce{006136}{}
On considère l'espace métrique $E=C([0,1])$ muni de $d_\infty$, et pour $f\in
E$, on note $M(f)$ le maximum de $f$ sur $[0,1]$.
Montrer que l'application $f\to M(f)$ est $1$-lipschitzienne.
\finenonce{006136}



\finexercice       


\exercice{6137, queffelec, 2011/10/16}
\enonce{006137}{}
Soit $(f_n)$ une suite de polyn\^omes qui converge uniformément sur $[0,1]$ vers
une fonction qui n'est pas un polyn\^ome. Montrer que la suite des degrés tend
vers l'infini.
\finenonce{006137}



\finexercice       


\exercice{6138, queffelec, 2011/10/16}
\enonce{006138}{}
On considère la suite de polyn\^omes sur $[-1,1]$ 
$$f_n(x)={{\int_0^x(1-t^2)^n\ dt}\over {\int_0^1(1-t^2)^n\ dt}}.$$

\begin{enumerate}  
\item Montrer que
pour tout $\varepsilon$, cette suite converge uniformément vers
$1$ sur l'intervalle $[\varepsilon,1]$, et  vers
$-1$ sur l'intervalle $[-1,-\varepsilon]$.

\emph{Indication :} Comparer $\int_0^1(1-t^2)^n\ dt$ à $\int_0^1(1-t)^n\ dt$.

\item En déduire que la suite $g_n(x)=\int_0^x f_n(t)\ dt$ converge uniformément
vers $\vert x\vert$ sur $[-1,1]$.

\item Montrer que dans l'exercice \ref{gijsexopol} la convergence est aussi
uniforme sur $[-1,1]$, en établissant une relation de récurrence satisfaite par
l'erreur  $\epsilon_n(t)=\vert t\vert-p_n(t).$
\end{enumerate}
\finenonce{006138}



\finexercice       


\exercice{6139, queffelec, 2011/10/16}
\enonce{006139}{}
Une fonction $f:[0,1]\to\Rr$ est dite réglée, si elle a en tout point une
limite à droite et une limite à gauche (et bien s\^ur, une limite à droite en $0$,
une limite à gauche en $1$.)
Montrer qu'une limite uniforme de fonctions en escalier est une fonction réglée
(la réciproque sera établie ultérieurement).
\finenonce{006139}



\finexercice       


\exercice{6140, queffelec, 2011/10/16}
\enonce{006140}{}
Soit $E$ l'espace $C_b(\Rr,C)$ des fonctions continues bornées sur $\Rr$
muni de la métrique de la convergence uniforme $d$.

\begin{enumerate}  
\item On rappelle qu'un espace topologique est séparable s'il contient un
dénom\-brable dense. Montrer que dans un espace métrique séparable, toute
collection d'ouverts deux à deux disjoints est au plus dénombrable. 

\item Soit $\lambda$ et $\mu$ deux réels distincts. Montrer que
$d(e^{i\lambda x}, e^{i\mu x})\geq 2$. En déduire que $E$ n'est pas séparable.
\end{enumerate}
\finenonce{006140}



\finexercice  

\exercice{6772, gijs, 2011/10/16}
\enonce{006772}{Questions de cours}

\begin{enumerate}
\item Donner les définitions
d'une topologie, d'un espace topologique Hausdorff, d'un
espace topologique quasi compact, d'un espace topologique
compact, d'un espace topologique connexe, et d'un espace
topologique connexe par arcs.
\item Donner un exemple d'un espace topologique connexe
mais pas connexe par arcs (une démonstration n'est pas
demandée !).
\item Soit $A\subset \Rr^n$. Montrer que $A$ est
quasi compact si et seulement si $A$ est fermé et borné
(énoncer clairement les théorèmes utilisés!).
\item Soit $X$ un espace topologique. Montrer que si
$X$ ne contient qu'un seul élément, alors $X$ est
compact.
\item Soit $(X,d)$ un espace métrique. Donner la
définition d'une suite de Cauchy dans $X$. Sous quelle
condition dit-on qu'une suite $x_n$ converge vers $a\in
X$ (noté $x_n \to a$)~?
\item Soit $(X,d)$ un espace métrique et $x_n$ une
suite dans $X$. Montrer que (1) si $x_n$ est une suite de
Cauchy, alors $x_n$ est bornée, et (2) si $x_n \to a$,
alors $x_n$ est une suite de Cauchy. 
\end{enumerate}
\finenonce{006772}



\finexercice

\exercice{6774, gijs, 2011/10/16}
\enonce{006774}{}
 Soient $X$ et $Y$ des espaces topologiques, et
soit $Y^X = \{\,f:X \to Y \,\}$ l'espace de toutes les
fonctions de $X$ dans $Y$. Pour $A\subset X$ et $B
\subset Y$ on définit $V(A,B) \subset Y^X$ par $V(A,B)
= \{\,f:X \to Y \mid f(A) \subset B\,\}$. La topologie
compacte-ouverte sur $Y^X$ a une sous-base constituée
des ensembles $V(A,B)$ où $A\subset X$ est compact et
$B\subset Y$ est ouvert.

 Montrer que $Y^X$ avec la topologie compacte-ouverte est
Hausdorff si et seulement si $Y$ est Hausdorff.
(Indication : pour le ``seulement si'' penser à une
fonction constante.)
\finenonce{006774}



\finexercice


\exercice{6786, gijs, 2011/10/16}
\enonce{006786}{}
\begin{enumerate}
\item Donner la définition d'un espace topologique $T_1$.

\item Montrer qu'un espace topologique $X$  est $T_1$
si et seulement si~: $\forall x\in X\ :\ \{x\}$ est
fermé.

\item Soit $X$ un espace topologique contenant un
nombre fini de points. Montrer que si $X$ est $T_1$,
alors sa toplogie est la topologie discrète.

\item Soit $X$ un espace topoogique $T_1$ ayant la
propriété~: $\forall x\in X, \forall A\subset X\ :
\text{ $A$ fermé et $x\notin A$ } \Longrightarrow
\exists U,V\text{ ouverts } : x\in U,\  A\subset V,\  U\cap
V = \emptyset$. Montrer que $X$ est $T_2$.
\end{enumerate}
\finenonce{006786}



\finexercice

\exercice{6787, gijs, 2011/10/16}
\enonce{006787}{}
Soit $K \subset \Rr^n$ un fermé tel que $K \subset
B(x_0,R)$ (où $B(x_0,R)$ est la boule ouverte de centre
$x_0$ et rayon $R$). Montrer qu'il existe un $R' < R$ tel
que $K\subset B(x_0,R')$. (indication~: regarder $\sup
d(x,x_0)$\,)
\finenonce{006787}



\finexercice


\exercice{6798, gijs, 2011/10/16}
\enonce{006798}{}
Soit $f:\Rr^2 \to \Rr$ une fonction continue,
posons  $M = \Rr^2 \setminus \{\, (x,y) \mid f(x,y) = 0
\,\}$ et $V= \{\, g:M \to \Rr \mid g \text{
dérivable, } \forall x\in M\,:\, g'(x) = 0 \,\}$. 
\begin{enumerate}
\item Montrer que $M$ est un ouvert de $\Rr^2$ et que
$V$ est un espace vectoriel. 

\item Calculer la
dimension de l'espace $V$ dans les cas suivants~:
\begin{itemize}
\item $f(x,y) = x^2 + y^2 - 1$,
\item $f(x,y) = \sin(x)$,
\item $f(x,y) = y^2 - x(x-1)(x-t)$.
\end{itemize}

Dans le dernier cas on vous demande de calculer la
dimension de $V$ en fonction de $t\in [0,1]$
(indication~: esquisser l'ensemble $f(x,y)=0$ et
distinguer les cas $t=0$, $t=1$, $0<t<1$).
\end{enumerate}
\finenonce{006798}



\finexercice


\exercice{6799, gijs, 2011/10/16}
\enonce{006799}{}
\begin{enumerate}
\item
Soit $X$ un espace $T_2$ et $K$ un sous ensemble quasi
compact. Montrer que $K$ est fermé.

\item Donner la définition d'un espace $T_4$.
Montrer qu'un espace $X$ est $T_4$ si et seulement si~:
pour tout fermé $A$ contenu dans un ouvert $U$ il
existe un ouvert $V$ tel que $A \subset V \subset
\overline{V} \subset U$.

\smallskip
Soient $f,g : X \to Y$ deux fonctions continues et $Y$ un
espace Hausdorff. 
\item Montrer que $\{\, x\in X \mid f(x) = g(x) \,\}$
est fermé.

\item Montrer que si $f$ et $g$ co\"\i ncident
sur un ensemble dense dans $X$, alors $f=g$.
\end{enumerate}
\finenonce{006799}



\finexercice


\exercice{6843, gijs, 2011/10/16}
\enonce{006843}{}
Soit $X$ un espace topologique, $Y$ un espace
topologique séparé, $U\subset X$ un sous ensemble dense de
$X$,  et soit $f,g:X
\to Y$ deux applications continues. Démontrer que 
$\{x\in X \mid f(x) = g(x) \,\}$ est un fermé. En déduire
que si $f$
et $g$ coincident sur $U$,
alors $f$ et $g$ coincident sur $X$.
\finenonce{006843}



\finexercice

\section{421.00 Compacité}


\exercice{2370, mayer, 2003/10/01}
\enonce{002370}{} 
Soit $X$ un espace m\'etrique.
\begin{enumerate}
\item Soit $A$ et $B$ deux compacts disjoints dans $X$. Montrer qu'ils
poss\`edent des voisinages ouverts disjoints (commencer par le cas
o\`u $B$ est r\'eduit \`a un point).
\item Soit $K$ un compact non vide de $X$ et $U$ un ouvert de $X$ contenant $K$.
Montrer qu'il existe $r>0$ tel que pour tout $x\in X$, on ait
l'implication: $$d(x,K)<r \Rightarrow x\in U\; .$$
\end{enumerate}
\finenonce{002370}


\finexercice


\exercice{2371, mayer, 2003/10/01}
\enonce{002371}{} 
Montrer qu'une suite convergente et sa limite forment un ensemble compact.
\finenonce{002371}


\finexercice


\exercice{2372, mayer, 2003/10/01}
\enonce{002372}{} 
Soient $K,F\subset \Rr^n$ des parties non vides, $K$ compact et $F$ ferm\'e.
Montrer qu'il existe $a\in K$ et $b \in F$ tel que $\|a-b\| = \mathrm{dist}(K,F)$.
\finenonce{002372}


\finexercice


\exercice{2373, mayer, 2003/10/01}
\enonce{002373}{} 
Soit $E$ un espace compact et soit $(F,d)$ un espace m\'etrique. Soit
$f:E\to F$ une application localement born\'ee, ce qui signifie que, pour tout $y\in E$,
il existe un voisinage $V_y$ de $y$ sur lequel $f$ est born\'ee. Montrer que $f$ est born\'ee
sur $E$.
\finenonce{002373}


\finexercice


\exercice{2374, mayer, 2003/10/01}
\label{exo 2}
\enonce{002374}{} 
Soit $X$ un espace m\'etrique.
\begin{enumerate}
\item Soit $(F_n)_n$ une suite d\'ecroissante de ferm\'es de $X$
et soit $(x_n)_n$ une suite convergente telle que $x_n\in F_n$
pour tout $n\geq 0$. Montrer que
$$\lim _{n\rightarrow \infty} x_n
\in \bigcap _{n\geq 0} F_n\;\; .$$
 Donner un exemple pour lequel
$\bigcap _{n\geq 0} F_n =\emptyset$.
\item Soit maintenant $(K_n)_n$ une suite d\'ecroissante de
{\it compacts}  non vides de $X$. V\'erifier que $K=\bigcap_{n\geq 0} K_n $
est non vide et que tout ouvert $\Omega$ qui contient $K$ contient
tous les $K_n$ \`a partir d'un certain rang.
\end{enumerate}
\finenonce{002374}


\finexercice


\exercice{2375, mayer, 2003/10/01}
\enonce{002375}{} 
Soit $X$ un espace topologique et $f:X\times [0,1] \to \Rr$
continue. Montrer que l'application $g: x\in X \to \int_0^ 1
f(x,y)\, dy$ est continue.
\finenonce{002375}


\finexercice


\exercice{2376, mayer, 2003/10/01}
\enonce{002376}{} 
Soit $E$ un espace norm\'e. Si $A$ et $B$ sont deux parties de $E$, on note
$A+B$ l'ensemble $\{a+b \; ; \; a\in A \; \text{et} \; b\in B\}$.
\begin{enumerate}
\item Montrer que si $A$ est compact et $B$ est ferm\'e, alors $A+B$ est ferm\'e.
\item Donner un exemple de deux ferm\'es de $\Rr^2$ dont la somme n'est pas ferm\'e.
\end{enumerate}
\finenonce{002376}


\finexercice


\exercice{2377, mayer, 2003/10/01}
\enonce{002377}{} 
Soit $f:\Rr^n \to \Rr^n$ une application continue. Elle est dite {\it propre}  si pour tout compact
$K\subset \Rr^n$, l'image r\'eciproque $f^{-1}(K)$ est compact.
\begin{enumerate}
\item Montrer que, si $f$ est propre, alors
l'image par $f$ de tout ferm\'e de $\Rr^n$ est un ferm\'e.
\item \'Etablir l'\'equivalence suivante: l'application $f$ est propre si et seulement si elle a
la propri\'et\'e:
$$ \|f(x)\|\to \infty  \quad \text{quand} \quad \|x\|\to \infty \; .$$
\end{enumerate}
\finenonce{002377}


\finexercice


\exercice{2378, mayer, 2003/10/01}
\enonce{002378}{} 
Soit $E =\{ f:[0,1] \rightarrow \Rr \;\; \text{continue} \}$. On munit
$E$ de la m\'etrique  $d_\infty (f,g) = \sup_{t\in [0,1]}
|f(t)-g(t)|$. Montrer que la boule unit\'e ferm\'ee de $E$ n'est
pas compact (on pourra construire une suite dont aucune sous suite
n'est de Cauchy).

Que peut-on dire de la boule unit\'e ferm\'ee de $l^\infty$ (l'espace des
suites born\'ees muni de la norme sup)?
\finenonce{002378}


\finexercice


\exercice{2379, mayer, 2003/10/01}
\enonce{002379}{} 
Soit $(X,d)$ un espace m\'etrique, soit $(Y,\delta )$ un espace m\'etrique compact
et soit $f:X\to Y$ une application dont le graphe
$$G=\{(x,f(x)) \; x\in X\} \subset X\times Y $$
est ferm\'e dans $X\times Y$. Notons $p:G\to X$ et $q:G\to Y$ les restrictions des deux
projections $p(x,y) = x$ et $q(x,y)=y$. Montrer que $p$ est un hom\'eomorphisme de $G$ sur $X$.
En d\'eduire que $f$ est continue.
\finenonce{002379}



\finexercice


\exercice{2380, mayer, 2003/10/01}
\enonce{002380}{} 
Soit $(X,d)$ un espace m\'etrique compact et $f:X\to X$ une application
v\'erifiant
$$d(f(x),f(y))<d(x,y) \quad \text{pour tout } x,y \in X \; , x\neq y\; .$$
Le but ici est de montrer que $f$ a un unique point fixe $p\in X$.
\begin{enumerate}
\item Justifier que $f$ peut avoir au plus un point fixe.
\item Montrer que les ensembles $X_n=f^n(X)$, $n\in \Nn$, forment une suite
d\'ecroissante de compacts et que $Y=\bigcap _{n\geq 0} X_n$ n'est pas vide.
\item Montrer que $Y$ est un ensemble invariant, i.e. $f(Y)=Y$, et
en d\'eduire que le diam\`etre de cet ensemble est zero.
\item Conclure que $f$ a un unique point fixe $p\in X$ et que pour
tout $x_0\in X$ la suite $x_n=f^n(x_0)\to p$, lorsque $n\to \infty$.
\end{enumerate}
\finenonce{002380}



\finexercice


\exercice{2381, mayer, 2003/10/01}
\enonce{002381}{} 
Soient $(E,d)$ un espace m\'etrique compact et $f:E\rightarrow E$
une application v\'erifiant
$$ d(f(x),f(y)) \geq d(x,y) \quad \text{pour tout} \;\; x,y \in E \;
.$$ On se propose de montrer que $f$ est une isom\'etrie
surjective. Soient $a,b\in E$ et posons, pour $n\geq 1$, $a_n
=f^n(a)=f\circ f^{n-1}(a)$ et $b_n = f^n(b)$.
\begin{enumerate}
  \item Montrer que pour tout $\epsilon >0$, il existe $k \geq 1$ tel
  que $d(a,a_k)<\epsilon$ et $d(b,b_k)< \epsilon$ (Consid\'erer une valeur
  d'adh\'erence de la suite $z_n=(a_n,b_n)$).
  \item En d\'eduire que $f(E)$ est dense dans $E$ et que
  $d(f(a),f(b)) =d(a,b)$ (Consid\'erer la suite $u_n
  =d(a_n,b_n)$).
\end{enumerate}
\finenonce{002381}


\finexercice


\exercice{2382, mayer, 2003/10/01}
\enonce{002382}{} 
On se donne une m\'etrique $d$ sur $X=[0,1]$ telle que
l'identit\'e $ i: (X, |.|) \to (X,d)$ soit continue (i.e. la
topologie d\'efinie par $d$ est moins fine que la topologie
usuelle de $X$).
\begin{enumerate}
\item Montrer que tout sous-ensemble de $X$ compact pour la
topologie usuelle est aussi compact pour la topologie d\'efinie
par $d$; puis montrer cette propri\'et\'e pour les ferm\'es.
\item En d\'eduire que la topologie d\'efinie par $d$ est la topologie
usuelle.
\end{enumerate}
\finenonce{002382}


\finexercice


\exercice{6164, queffelec, 2011/10/16}
\enonce{006164}{}
\begin{enumerate}
\item Soit $X$ un espace topologique séparé. Montrer qu'il est compact et discret
si et seulement si il est fini.

\item Montrer que dans un espace topologique séparé, l'ensemble constitué d'une
suite convergente et de sa limite est compact.
\end{enumerate}
\finenonce{006164}



\finexercice


\exercice{6165, queffelec, 2011/10/16}
\enonce{006165}{}
Soit $X$ un espace topologique compact et $f_1,f_2,\ldots,f_n$, $n$ fonctions
continues réelles qui séparent les points de $X$. Montrer que $X$ est homéomorphe
à une partie de $\Rr^n$. 
\finenonce{006165}



\finexercice


\exercice{6166, queffelec, 2011/10/16}
\enonce{006166}{}
Soit $X, Y$ deux espaces topologiques séparés et $(K_n)$ une suite décroissante
de com\-pacts non vides de $X$. 
Soit $f:X\to Y$ une application continue. Montrer
que $f(\cap_n K_n)=\cap_n f(K_n)$.
\finenonce{006166}



\finexercice


\exercice{6167, queffelec, 2011/10/16}
\enonce{006167}{}
\label{exn11}
Soit $X$ un espace topologique séparé et $A$ et $B$ deux compacts disjoints dans
$X$. Montrer qu'ils possèdent des voisinages ouverts disjoints. (Commencer par le
cas où $B$ est réduit à un point).
\finenonce{006167}



\finexercice


\exercice{6168, queffelec, 2011/10/16}
\enonce{006168}{}
Soit $(f_n)$ une suite croissante de fonctions réelles définies sur un espace
topolo\-gique compact $X$, convergeant simplement vers une fonction $f$; on
suppose que les fonctions $f_n$ et $f$ sont continues.
Montrer que la convergence est uniforme sur $X$.

Application : montrer que la suite de fonctions $f_n$ définies sur $[0,1]$
par $f_n(x)=\sum_1^{n-1} x^k(1-x)^{n-k}$ converge vers $0$ uniformément sur
$[0,1]$.  
\finenonce{006168}



\finexercice


\exercice{6169, queffelec, 2011/10/16}
\enonce{006169}{}
Soit $X$ un espace topologique compact et $C(X)$ l'espace des fonctions réelles
conti\-nues sur $X$ avec la norme uniforme.

Soit $J$ un idéal propre de $C(X)$; on va montrer par l'absurde que toutes les
fonctions de
$J$ s'annulent en un même point de $X$.

\begin{enumerate}
\item Sinon, montrer qu'on peut trouver $n$ points de $X$, $x_1,\cdots,x_n$,
$V_1,\cdots,V_n$ où $V_i$ voisinage de $x_i$ et $n$ fonctions de $J$,
$f_1,\cdots,f_n$ tels que
 $$X=\cup_i\ V_i, \  \  f_i{}_{|V_i}\not=0.$$

\item Construire alors une fonction $g$ dans $J$ ne s'annulant jamais et en déduire
que ${\bf 1}\in J$, d'où la contradiction. 
\end{enumerate}
\finenonce{006169}



\finexercice


\exercice{6170, queffelec, 2011/10/16}
\enonce{006170}{}
Soit $X$ un espace topologique et $f:X\times[0,1]\to\Rr$ continue. Montrer
que l'application $g:x\in X\to\int_0^1f(x,y)\, dy$ est continue.
\finenonce{006170}



\finexercice       


\exercice{6171, queffelec, 2011/10/16}
\enonce{006171}{}
Soit $X=[a,b]$ et on se donne une métrique $d$ sur $X$ telle que la topologie
définie par $d$ est moins fine sur $X$ que la topologie usuelle. Montrer que
tout sous-ensemble de $X$ compact pour la topologie usuelle est aussi compact
pour la topologie définie par $d$; puis montrer cette propriété pour les fermés.

En déduire que la topologie définie par $d$ est la topologie usuelle. 
\finenonce{006171}



\finexercice       


\exercice{6172, queffelec, 2011/10/16}
\enonce{006172}{}
Soit $X$ un espace topologique séparé et $(K_n)$ une suite décroissante de
com\-pacts non vides de $X$. Montrer que $K=\cap K_n$ est non vide et que si
$\Omega$ est un ouvert contenant $K$, il contient tous les $K_n$ à partir d'un
certain rang.
\finenonce{006172}



\finexercice       


\exercice{6173, queffelec, 2011/10/16}
\enonce{006173}{}
Soit $f$ et $g$ deux fonctions réelles continues sur un espace topologique
compact $X$, telles que $f\geq0$, et $f(x)>0$ si $g(x)\leq0$.
Montrer qu'il existe une constante $A>0$ telle que $$Af(x)+g(x)>0,\;  \forall
x\in X.$$
(\emph{Indication :} raisonner par l'absurde, et considérer les ensembles $A_n=\{x\in X
/nf(x)+g(x)\leq 0\}$).
\finenonce{006173}



\finexercice       


\exercice{6174, queffelec, 2011/10/16}
\enonce{006174}{}
Soit $X$ un espace topologique compact et $f_1,f_2,\ldots,f_n$, $n$ fonctions
continues réelles qui séparent les points de $X$. Montrer que $X$ est homéomorphe
à une partie de $\Rr^n$. 
\finenonce{006174}



\finexercice       


\exercice{6175, queffelec, 2011/10/16}
\enonce{006175}{}
Montrer que toute fonction réglée sur $[0,1]$ s'approche uniformément par des
fonctions en escalier.
\finenonce{006175}



\finexercice       


\exercice{6176, queffelec, 2011/10/16}
\enonce{006176}{}
Soit $(f_n)$ une suite croissante de fonctions réelles définies sur un espace
topolo\-gique compact $X$, convergeant simplement vers une fonction $f$; on
suppose que les fonctions $f_n$ et $f$ sont continues.

Montrer que la convergence est uniforme sur $X$.  
\finenonce{006176}



\finexercice       


\exercice{6177, queffelec, 2011/10/16}
\enonce{006177}{}
Soit $X$ un espace topologique compact et $C(X)$ l'espace des fonctions
conti\-nues sur $X$ avec la norme uniforme.

\begin{enumerate}  
\item Soit $J$ un idéal propre de $C(X)$ ; montrer que toutes les fonctions de $J$
s'annulent en un même point de $X$.
\emph{Indication :} raisonner par l'absurde, utiliser le fait qu'une fonction continue
$\neq0$ en $x$, est $\neq0$ sur un voisinage de $x$ et recouvrir $X$ avec de
tels voisinages. 

\medskip

Pour $f\in J$, on note $Z_f=f^{-1}(\{0\})$, l'ensemble des zéros de $f$.

\item Soit $J$ un idéal de $C(X)$ et $Z=\cap_{f\in J}Z_f$; $Z$ est fermé.
  \begin{enumerate}  
  \item Soit $K$ un fermé de $X$ disjoint de $Z$. Par un raisonnement
analogue à celui du 1., construire $f\in J$, $f\geq 0$ et ne s'annulant pas sur
$K$.

Etudier la limite $F$ de ${nf\over{1+nf}}$ dans $C(X)$.

  \item  Montrer que si $g\in C(X)$ s'annule sur un ouvert contenant $Z$,
alors $g\in J$ et $Z\neq\emptyset$.
  \item Soit $g\in C(X)$ nulle sur $Z$; par un bon choix de $K$, montrer
que $g\in \overline J$.
  \end{enumerate}
En déduire la description des idéaux fermés de $C(X)$. 
\end{enumerate}
\finenonce{006177}



\finexercice       


\exercice{6178, queffelec, 2011/10/16}
\enonce{006178}{}
Montrer que les sous-groupes compacts du groupe multiplicatif $\Cc^*$ sont
con\-tenus dans $\mathbb{U}$ le sous-groupe des nombres complexes de module $1$.
\finenonce{006178}



\finexercice


\exercice{6179, queffelec, 2011/10/16}
\enonce{006179}{}
On rappelle la construction de l'ensemble triadique de Cantor : on part du
segment $[0,1]$ dont on supprime l'intervalle médian $]{1\over 3}, {2\over
3}[$; à la deuxième étape, on supprime les intervalles  $]{1\over 9},
{2\over 9}[$ et $]{7\over 9}, {8\over 9}[$ etc. On note $K_n$ la réunion des
intervalles restants à la $n$-ième étape, et $K=\bigcap K_n$. Montrer
que $K$ est un compact d'intérieur vide, sans point isolé.
\finenonce{006179}



\finexercice


\exercice{6180, queffelec, 2011/10/16}
\enonce{006180}{}
On considère dans $M_n(\Rr)$ le sous-ensemble des matrices de déterminant
égal à $1$. Est-il compact ? On note $O(n)$ le sous-ensemble des matrices
orthogonales ($^tA.A=I$); montrer que $O(n)$ est compact.
\finenonce{006180}



\finexercice


\exercice{6181, queffelec, 2011/10/16}
\enonce{006181}{}
Montrer que dans un evn, la boule unité fermée est compacte si et seulement si
la sphère unité est compacte.
\finenonce{006181}



\finexercice


\exercice{6182, queffelec, 2011/10/16}
\enonce{006182}{}
Soit $A$ une partie d'un espace normé $E$. On note co$(A)$, l'enveloppe convexe
de $A$ ie l'ensemble $\{\sum_{\text{finie}}\lambda_j a_j,\,
\lambda_j\geq0,\, \sum\lambda_j =1\}$ des combinaisons convexes de points de $A$.
\begin{enumerate}
\item Montrer que si $A$ est fini, co$(A)$ est compacte.

\item Montrer que si $E$ est de dimension finie $n$ et $A$ compact, co$(A)$ est
compacte (on admettra que tout point de co$(A)$ est combinaison convexe d'au plus
$n+1$ points de $A$).
\end{enumerate}
\finenonce{006182}



\finexercice


\exercice{6183, queffelec, 2011/10/16}
\enonce{006183}{}
Soit $(X,d)$ un espace métrique.
\begin{enumerate}
\item Soit $A$ une partie compacte de $X$; montrer qu'il
existe $x,y\in A$ tels que  diam$A=d(x,y)$.

\item Soit $A$ et $B$ deux parties compactes
disjointes de
$X$. Montrer qu'il existe $\delta>0$ tel que  
$d(a,b)\geq \delta\ \forall {a\in A, b\in B} $. En déduire une démonstration
simple de l'exercice 10 dans le cadre métrique.

\item Montrer que le résultat est encore
vrai si l'une est compacte et l'autre fermée, mais devient faux si 
les deux parties
sont seulement fermées.
\end{enumerate}
\finenonce{006183}



\finexercice


\exercice{6184, queffelec, 2011/10/16}
\enonce{006184}{}
Soit $f$ une surjection continue de $\Rr^2$ sur $\Rr$. On va montrer que
l'image réciproque de tout point est non bornée. On raisonne par l'absurde : 

 Sinon, il existe $a\in \Rr$ et un disque fermé $ D$ du plan tel que
$f^{-1}(\{a\})\subset  D$; en étudiant $f(D^c)$ et $f(D)$ montrer que
$f(\Rr^2)$ ne peut être égal à $\Rr$ tout entier.
\finenonce{006184}



\finexercice


\exercice{6185, queffelec, 2011/10/16}
\enonce{006185}{}
Soit $F_1,F_2,...,F_p$, $p$ fermés d'un espace métrique compact $E$, tels que
$F_1\cap ...\cap F_p= \emptyset$. Montrer qu'il existe $\varepsilon>0$ tel que
toute partie $A$ de $E$ rencontrant tous les $F_i$ ait un diamètre
$\geq\varepsilon$ (raisonner par l'absurde).
\finenonce{006185}



\finexercice


\exercice{6186, queffelec, 2011/10/16}
\enonce{006186}{}
Soit $X=X_1\times X_2\times\cdots\times X_n$ où les $X_i$ sont $n$ espaces
métriques, et on note $p_i$ la projection de $X$ sur $X_i$. Montrer que
$A\subset X$ est compact si et seulement si $A$ est fermé dans $X$ et les
$p_i(A)$ sont tous compacts.
\finenonce{006186}



\finexercice


\exercice{6187, queffelec, 2011/10/16}
\enonce{006187}{}
\begin{enumerate}
\item Montrer que la boule unité fermée d'un evn de dimension finie est compacte.

\item Soit $E$ un espace vectoriel normé et $F$ un sous-espace de $E$ de dimension
finie. Montrer que $d(x,F)=\inf\{d(x,y),\ y\in F,\ ||y||\leq2||x||\}$; en
déduire que $F$ est fermé dans $E$.  


\item Soit $(f_n)$ une suite de polyn\^omes qui converge uniformément sur $[0,1]$ vers
une fonction qui n'est pas un polyn\^ome. Montrer que la suite des degrés tend
vers l'infini (raisonner par l'absurde).
\end{enumerate}
\finenonce{006187}



\finexercice


\exercice{6188, queffelec, 2011/10/16}
\enonce{006188}{Partiel de décembre 1998}
 
Soit $E$ un espace vectoriel normé sur $\Cc$ de boule unité fermée
$\overline B$ et
$F$ un sous-espace vectoriel fermé de $E$. On a montré dans le liste précédente
que si
$F\not=E$,\quad $\sup_{x\in\overline B} d(x,F)=1.$

On va montrer qu'un evn dont la boule unité fermée est compacte est
nécessaire\-ment de dimension finie.
On suppose donc que $\overline B$ est compacte.
\begin{enumerate}
\item Montrer que pour tout $\varepsilon >0$ on peut trouver un nombre fini de
points $x_1,\cdots,x_k\in\overline B$ tels que $\overline B\subset\cup_{j=1}^k
B(x_j,\varepsilon)$.

\item Montrer que $E$ est de dimension finie : pour cela, considérer le
sous-espace vectoriel engendré par $x_1,\cdots,x_k$.
\end{enumerate}
\finenonce{006188}



\finexercice


\exercice{6189, queffelec, 2011/10/16}
\enonce{006189}{}
Voici quelques applications du fait important suivant : dans un espace
mé\-trique compact, toute suite ayant une seule valeur d'adhérence converge.
\begin{enumerate}
\item Soit $(a_n)$ une suite bornée de réels, telle que $(e^{ita_n})$ converge pour
un ensemble non dénombrable de $t\in \Rr$; montrer que la suite $(a_n)$
converge.

\item Soit $f$ une application de ${\Rr}$ dans ${\Rr}$ et $G$ son graphe.
Montrer que  si $G$ est connexe par arcs, $f$ est continue.

\item Soit $f$ une application de $X$ dans $Y$, espaces métriques et $G$ le graphe
de
$f$. Montrer que $G$ est fermé dans $X\times Y$ si $f$ est continue.
Montrer que la réciproque est vraie lorsque
$Y$ est compact.

\item Soit $X$ un espace métrique, $Y$ un espace métrique compact et
$f:X\times Y\to \Rr$ une application continue telle que, pour tout $x\in X$,
l'équation $f(x,y)=0$ ait une unique solution $y\in Y$. Montrer que
l'application $u:x\in X\to y\in Y$ ainsi définie est continue.
\end{enumerate}
\finenonce{006189}



\finexercice


\exercice{6190, queffelec, 2011/10/16}
\enonce{006190}{}
On considère une suite $(x_n)$ de $[0,1]$
telle que $x_{n+1}-x_n$ tend vers $0$. Soit $A$ l'ensemble de ses valeurs
d'adhérence.
\begin{enumerate}
\item Justifier le fait que $A$ est non vide. Si
$\alpha\notin A$, montrer qu'il existe $\varepsilon>0$ et
$n_0$ tels que les points $x_n,\ n\geq n_0$, soient en dehors de
$[\alpha-\varepsilon,\alpha+\varepsilon]$. Montrer ainsi que $A$ est un
intervalle (si $\alpha$ et $\beta \in A$, ${\alpha+\beta\over2}\in A$).

\item On suppose de plus que cette suite est une suite récurrente i.e. définie par $
x_{n+1}=f(x_n)$ où $f$ est continue de $[0,1]$ dans lui-même, et un point initial
$x_0\in [0,1]$. Montrer alors que la suite converge (on commencera par remarquer
que si
$x\in A$, alors $x=f(x)$, et que si $x_m\in A$ pour un indice $m$, alors la
suite converge.)

\item Soit $x=(x_n)$ une suite de $l^{\infty}$; montrer que l'ensemble des valeurs
d'adhé\-rence de la suite $y$ de terme général $y_n={{x_1+x_2+\cdots+x_n}\over
n}$ est un intervalle. En déduire que l'application $f$ de $l^{\infty}$ dans
lui-même qui associe $y$ à $x$, n'est pas bijective.
\end{enumerate}
\finenonce{006190}



\finexercice


\exercice{6191, queffelec, 2011/10/16}
\enonce{006191}{}
On note ${\mathbb{S}^1}$ le cercle unité dans ${\Rr}^2$, et $h$ l'application de
${\Rr}$ dans ${\mathbb{S}^1}$: 
$t\to(\cos 2\pi t, \sin 2\pi t)$. 
\begin{enumerate}
\item Montrer que le cercle privé d'un point, ${\mathbb{S}^1}\backslash
\{a\}$, est homéomorphe à l'inter\-valle $]0,1[$. 

\item Montrer que $h$ est une bijection continue de $[0,1[$ sur ${\mathbb{S}^1}$, mais
n'est pas un homéomorphisme.
\end{enumerate}
\finenonce{006191}



\finexercice


\exercice{6192, queffelec, 2011/10/16}
\enonce{006192}{}
Démontrer de plusieurs fa\c cons que le cercle unité
$\mathbb{S}^1\subset\Rr^2$ est compact.
\finenonce{006192}



\finexercice       


\exercice{6193, queffelec, 2011/10/16}
\enonce{006193}{}
Soit $(X,d)$ un espace métrique, $A$ et $B$ deux parties de $X$. On pose
$d(A,B)=\inf_{a\in A, b\in B} d(a,b)$.

\begin{enumerate}  \item Si $A$ et $B$ sont disjointes, l'une compacte et l'autre fermée, montrer que
$d(A,B)>0$.
\item Montrer, par un contre-exemple, que ceci peut être faux si les deux parties
sont seulement fermées. 
\end{enumerate}
\finenonce{006193}



\finexercice       


\exercice{6194, queffelec, 2011/10/16}
\enonce{006194}{}
Soit $E$ un espace normé, $X$ et $Y$ deux sous-ensembles de $E$. Montrer que
\begin{enumerate}  
\item $X+Y$ est ouvert si $X$ est ouvert;

\item $X+Y$ est compact si $X$ et $Y$ sont compacts;

\item $X+Y$ est fermé si $X$ est compact et $Y$ fermé. 
\end{enumerate}
Que peut-on dire de $X+Y$ si $X$ et $Y$ sont seulement fermés ?
\finenonce{006194}



\finexercice       


\exercice{6195, queffelec, 2011/10/16}
\enonce{006195}{}
Soit $E$ un espace normé, $X$ et $Y$ deux parties compactes de $E$.
Montrer que la réunion des segments joignant un point $x\in X$ à un point $y\in
Y$ est encore compacte.
\finenonce{006195}



\finexercice       


\exercice{6196, queffelec, 2011/10/16}
\enonce{006196}{}
Soit $K$ un convexe compact symétrique de $\Rr^n$ contenant $0$ comme point
intérieur. Alors
$K$ est la boule unité fermée associée à une norme de $\Rr^n$ : considérer pour
cela $$p(x)= \inf\{t>0 /{x\over t}\in K\}$$
\finenonce{006196}



\finexercice       


\exercice{6197, queffelec, 2011/10/16}
\enonce{006197}{}
Trouver l'ensemble des valeurs d'adhérence quand $x\to 0$ de $f(x)=\sin{1\over
x}$; $g(x)={1\over x}\sin{1\over x}$.
\finenonce{006197}



\finexercice       


\exercice{6198, queffelec, 2011/10/16}
\enonce{006198}{}
\begin{enumerate}  
\item Soit $X$ un espace métrique compact et $(f_n)$ une suite d'applications
continues à valeurs dans un espace métrique $Y$, convergeant vers $f$
uniformément sur $X$. Montrer que si
$(x_n)$ est une suite de points de $X$ convergeant vers $x\in X$, alors
$f_n(x_n)$ tend vers
$f(x)$.

\item Application : Soit $X$ un espace métrique compact, et soit $(f_n)$ une suite
d'applications de $X$ dans $X$, ayant chacune un point
fixe; on suppose que la suite $(f_n)$ converge vers une fonction $f$
uniformément sur $X$. Montrer que $f$ a aussi un point fixe.

\item Soit $K$ un convexe compact de $\Rr^n$ et $f$ une application continue de
$K$ dans $K$ vérifiant
$$\Vert f(x)-f(y)\Vert \leq \Vert x-y\Vert;$$
En considérant les fonctions $f_n$ définies sur $K$ par $f_n(x)={1\over
n}f(x_0)+(1-{1\over n})f(x)$, où $x_0\in K$, montrer que
$f$ a un point fixe. Est-il unique ? Que se passe-t-il si $K$ n'est plus
convexe ?
\end{enumerate}
\finenonce{006198}



\finexercice       


\exercice{6199, queffelec, 2011/10/16}
\enonce{006199}{}
Soit $A$ une partie d'un espace normé $E$. On note co$(A)$, l'enveloppe convexe
de $A$ ie l'ensemble $\{\sum_{finie}\lambda_j a_j,\,
\lambda_j\geq0,\, \sum\lambda_j =1\}$ des combinaisons convexes de points de $A$.

\begin{enumerate}  
\item Montrer que si $A$ est fini, co$(A)$ est compacte.

\item Montrer que si $E$ est de dimension finie et $A$ compact, co$(A)$ est
compacte. 
\end{enumerate}
\finenonce{006199}



\finexercice       


\exercice{6200, queffelec, 2011/10/16}
\enonce{006200}{}
Soit $E=C_b(\Rr)$ muni de la norme uniforme; pour $f\in E$, on note
$f_a$ la translatée de $f$ par $a$, ie la fonction $x\to f(x-a)$, et $O_f$
l'ensemble des translatées de $f$.

\begin{enumerate}  
\item Montrer que si $f$ est périodique, $O_f$ est compact (considérer
l'application $a\to f_a$).

\item Soit $f$ une limite uniforme sur $\Rr$ de fonctions périodiques; montrer
que $O_f$ est précompact.

\item On suppose cette fois $O_f$ précompact; on va montrer que $f$ est
uniformé\-ment continue.
  \begin{enumerate}  
  \item De toute suite $(f_{a_n})$ de $O_f$ on peut extraire une sous-suite
convergente dans $E$.
  \item Si $x_n-y_n$ tend vers $0$, montrer que $(f_{x_n-y_n})$ n'a qu'une
valeur d'adhé\-rence $f$; en déduire que $f(x_n)-f(y_n)$ tend vers $0$.
  \item  Montrer que $f$ est uniformément continue.
  \end{enumerate}
\end{enumerate}
\finenonce{006200}



\finexercice  


\exercice{6201, queffelec, 2011/10/16}
\enonce{006201}{}
Soit $E$ l'ensemble des suites infinies de nombres réels $x=(x_1,x_2,\cdots)$
à valeurs $0$ ou $1$.
Si $x$ et $y$ sont deux éléments de $E$, on pose
$$d(x,y)=\sup_{k\geq1}({1\over k}\vert x_k-y_k\vert)$$
\begin{enumerate}
\item Montrer que $d$ est une distance sur $E$.

\item Soit $\varepsilon>0$; montrer qu'il existe une partie finie $E_\varepsilon$
de $E$ qui possède la propriété suivante : les boules fermées de rayon 
$\varepsilon$ centrées en un point de $E_\varepsilon$ recouvrent $E$.

\item Montrer que $E$ est compact.
\end{enumerate}
\finenonce{006201}



\finexercice


\exercice{6202, queffelec, 2011/10/16}
\enonce{006202}{}
Soit $K$ un convexe compact de $\Rr^2$.
\begin{enumerate}
\item Si $K$ est d'intérieur vide, montrer que $K$ est homéomorphe au segment
$[0,1]$.

\item Si $K$ n'est pas d'intérieur vide, montrer que $K$ est homéomorphe au disque
unité fermé en considérant l'application $p(x)=\inf\{a>0\ ;\ {x\over a}\in
K\}$; on montrera que $0$ est un point intérieur, que $\delta ||x||\leq
p(x)\leq C||x||$ puis que $p$ est continue.
\end{enumerate}
\finenonce{006202}



\finexercice


\exercice{6203, queffelec, 2011/10/16}
\enonce{006203}{}
Soit $(A_n)$ une suite décroissante de compacts connexes non vides dans un
espace topologique séparé. Montrer que $\cap_n A_n$ est encore un compact
connexe non vide. (Pour la connexité on pourra raisonner avec de fermés et
utiliser l'exercice \ref{exn11}.) 
\finenonce{006203}



\finexercice


\exercice{6783, gijs, 2011/10/16}
\enonce{006783}{}
Soit $X$ un espace topologique et $A \subset X$
un sous-ensemble de $X$.
\begin{enumerate}
\item Montere que si $X$ est quasi compact et
$A$ est fermé dans $X$, alors $A$ est quasi compact.

\item Montrer que si $X$ est Hausdorff et $A$ est
quasi compact, alors $A$ est fermé dans $X$.
\end{enumerate}
\finenonce{006783}



\finexercice

\exercice{6788, gijs, 2011/10/16}
\enonce{006788}{}
\label{gijsexo5}
Soit $(X,d)$ un espace métrique  compact, et
soit $f:X  \to X$ une application continue verifiant
$\forall x\neq y\ :\ d(f(x),f(y)) < d(x,y)$. Montrer qu'il
existe un point fixe unique pour $f$. (indication~:
regarder $\inf d(x,f(x))$\,)
\finenonce{006788}



\finexercice

\exercice{6797, gijs, 2011/10/16}
\enonce{006797}{}
Soit $K$ un compact contenu dans $U$ un ouvert de
$\Rr^n$, et soit $K_\epsilon$ défini par
 $K_\epsilon = \{\,x\in \Rr^n \mid d(x,K)
\le \epsilon \,\}$.

\begin{enumerate}
\item Montrer que pour tout $\epsilon>0$ l'ensemble
$K_\epsilon$ est compact.

\item  Montrer qu'il existe un $\epsilon>0$ tel que
$K_\epsilon$ soit contenu dans $U$ (indication : regarder
la fonction $x \mapsto d(x, \Rr^n \setminus U)$
définie sur $K$).
\end{enumerate}
\finenonce{006797}



\finexercice


\exercice{6817, gijs, 2011/10/16}
\enonce{006817}{}
Soit $(X,d)$ un espace métrique. On suppose que {\it
toutes\/} les boules fermées (c'est-à-dire les
ensembles de la forme $\{y\in X \mid d(x,y) \le r\}$) sont
compactes.
\begin{enumerate}
\item Démontrer que $X$ est complet.

\item Démontrer que si $A \subset X$ est fermé
et borné, alors $A$ est compact.
\end{enumerate}
\finenonce{006817}



\finexercice


\exercice{6818, gijs, 2011/10/16}
\enonce{006818}{}
Soit $(X,d)$ un espace métrique compact et $(a_n)_{n\in
\Nn}$ une suite dans $X$. Démontrer que si $(a_n)_{n\in
\Nn}$ n'a qu'une seule valeur d'adhérence $x\in X$, alors
$ \lim_{n\to \infty} a_n = x$. (Indication~: on
suppose que la suite ne converge pas vers $x$.)
\finenonce{006818}



\finexercice


\exercice{6827, gijs, 2011/10/16}
\enonce{006827}{}
 Soit $X$ un espace topologique compact. Soit
$\mathcal{F}$ une collection de fonctions continues sur $X$ à
valeurs dans $\Rr$. La collection $\mathcal{F}$ a les deux
propriétés suivantes~:
  \begin{enumerate}
  \item
Si $f$ et $g$ appartiennent à $\mathcal{F}$, alors leur
produit $f\cdot g$ appartient à $\mathcal{F}$.

  \item
Pour tout $x\in X$ il existe un voisinage $U$ de $x$ et
une fonction $f\in \mathcal{F}$ tel que $f$ est identiquement
nulle sur $U$.

  \end{enumerate}
Démontrer que la fonction qui est identiquement nulle
sur $X$ appartient à $\mathcal{F}$.
\finenonce{006827}



\finexercice


\exercice{6835, gijs, 2011/10/16}
\enonce{006835}{}
Soit $X$ un espace topologique séparé et soit $\infty$
un point qui n'appartient pas à $X$. On pose
$X_\infty = X \cup \{\infty\}$ et  on définit $\mathcal{T}_\infty 
\subset \mathcal{P}(X_\infty) $  par~: $A\in \mathcal{T}_\infty$ si et seulement si 
\begin{align*}
\infty \notin A &\Rightarrow \text{ $A$ est
un ouvert de $X$~;}
\\
\infty\in A &\Rightarrow \text{ $X\setminus A$ est
un compact de $X$.}
\end{align*}

\begin{enumerate}
\item
Démontrer que si $A$ appartient à $\mathcal{T}_\infty$,
alors
$A \cap X$ est un ouvert de $X$.

\item 
Démontrer que $\mathcal{T}_\infty$
est  une topologie sur $X_\infty$.

\item
Démontrer que $X$ est dense dans $X_\infty$.

\item
Démontrer que $X_\infty$ est compact.
\end{enumerate}
\finenonce{006835}



\finexercice

\section{422.00 Continuité, uniforme continuité}


\exercice{2353, queffelec, 2003/10/01}
\enonce{002353}{}
Soit $X$ un espace topologique et $f:X\to \Rr$.
\begin{enumerate}
\item Montrer que $f$ est continue si et seulement si pour tout $\lambda\in \Rr$, 
les ensembles $\{x\ ;\ f(x)<\lambda\}$ et $\{x\ ;\ f(x)>\lambda\}$ sont des
ouverts de $X$.

\item Montrer que si $f$ est continue, pour tout $\omega$ ouvert de $\Rr$,
$f^{-1}(\omega)$ est un $F_\sigma$ ouvert de $X$ ($F_\sigma$= r\'eunion
d\'enombrable de ferm\'es).

\end{enumerate}
\finenonce{002353}


\finexercice


\exercice{2354, queffelec, 2003/10/01}
\enonce{002354}{}
\begin{enumerate}
\item Soit $C$ l'espace des fonctions continues r\'eelles sur $[0,1]$ muni de la
m\'etrique
$d_1(f,g)=\int_0^1\vert f-g\vert\ dx$, puis de la m\'etrique
$d_\infty(f,g)=\sup_x\vert f(x)-g(x)\vert$. V\'erifier que l'applica\-tion
$f\to \int_0^1 |f| dx$ de
$C$ dans
$\Rr$ est $1$-lipschitzienne dans les deux cas.   

\item Soit $c$ l'espace des suites r\'eelles convergentes, muni de la m\'etrique
$d(x,y)=\sup_n\vert x(n)-y(n)\vert$. Si on d\'esigne par $\ell(x)$ la limite de la
suite $x$, montrer que $\ell$ est une application continue de $c$ dans $\Rr$. En
d\'eduire que $c_0$ est ferm\'e dans $c$.
\end{enumerate}
\finenonce{002354}


\finexercice


\exercice{2355, queffelec, 2003/10/01}
\enonce{002355}{}
Soit $f,g$ deux applications continues de $X$ dans $Y$, espaces topologiques, $Y$ \'etant s\'epar\'e.
Montrer que $\{f=g\}$ est ferm\'e dans $X$; en d\'eduire que si $f$ et $g$
coïncident sur une partie dense de $X$, alors $f=g$.
\finenonce{002355}


\finexercice


\exercice{2356, queffelec, 2003/10/01}
\enonce{002356}{}
Une application de $X$ dans $Y$ est dite {\it ouverte} si l'image de tout ouvert
de $X$ est un ouvert de $Y$; {\it ferm\'ee} si l'image de tout ferm\'e de $X$ est un
ferm\'e de $Y$. 

\begin{enumerate}
\item Montrer qu'une fonction polynomiale de ${\Rr}$ dans ${\Rr}$ est une
application ferm\'ee.

\item Montrer que l'application $(x,y)\in X\times Y \to x\in X$ est ouverte mais pas
n\'ecessairement ferm\'ee (consid\'erer l'hyperbole \'equilat\`ere de  ${\Rr}^2$).

\item Montrer que la fonction indicatrice de l'intervalle $[0, {1\over2}]$, comme
application de $\Rr$ dans $\{0,1\}$, est surjective, ou\-verte, ferm\'ee, mais
pas continue.

\item Montrer que toute application ouverte de $\Rr$ dans $\Rr$ et continue est monotone.
\end{enumerate}
\finenonce{002356}


\finexercice


\exercice{2357, queffelec, 2003/10/01}
\enonce{002357}{}
\begin{enumerate}
\item Montrer que $f$ est continue si et seulement si $f(\overline A)\subset
\overline{f(A)}$ pour tout $A$ dans $X$. Que peut-on dire alors de l'image par
$f$ d'un ensemble dense dans $X$?  

\item Montrer que $f$ est ferm\'ee si et seulement si $\overline{f(A)}\subset
f(\overline A)$, et que $f$ est ouverte si et seulement si ${f(\stackrel\circ
A)}\subset \stackrel\circ{f( A)}$.
\end{enumerate}
\finenonce{002357}


\finexercice


\exercice{2358, queffelec, 2003/10/01}
\enonce{002358}{} 
\begin{enumerate}
\item Soit $f$ une fonction r\'eelle continue sur $[0,1]$; montrer que $f$ est
``presque lipschitzienne'' au sens :

$\forall\varepsilon>0 \ \exists C_\varepsilon\ ;\ \forall x,y \in [0,1]\quad  \
|f(x)-f(y)|\leq C_\varepsilon|x-y| + \varepsilon.$

\item Montrer qu'une fonction $f$ uniform\'ement continue de $\Rr$ dans $\Rr$
v\'erifie pour tout $x\in\Rr$,
$\vert f(x)\vert\leq a\vert x\vert+b$ o\`u $a$ et $b$ sont des constantes.
\end{enumerate}
\finenonce{002358}



\finexercice


\exercice{2359, queffelec, 2003/10/01}
\enonce{002359}{}
Soit $f$ une fonction continue de $]0,1[$ dans $\Rr$. Montrer que, si $f$ est
unifor\-m\'ement continue, elle est born\'ee. R\'eciproque ?
\finenonce{002359}



\finexercice


\exercice{2360, queffelec, 2003/10/01}
\enonce{002360}{}
Soit
$f$ une fonction uniform\'ement continue sur
$\Rr$ telle que
$\int_0^\infty f(t) dt$ con\-verge.
Montrer que $f$ tend vers $0$ quand $x\to+\infty$. Retrouver ainsi le fait que
la fonction $\sin(x^2)$ n'est pas uniform\'ement continue.
\finenonce{002360}


\finexercice

\exercice{6066, queffelec, 2011/10/16}
\enonce{006066}{}
 Soit $f$ une isométrie de $\R$ dans $\R$. Montrer qu'on a soit
 $f(x)=a-x$, soit $f(x)=a+x$, où $a=f(0)$. (Se ramener à $a=0$.)
\finenonce{006066}



\finexercice


\exercice{6067, queffelec, 2011/10/16}
\enonce{006067}{}
Soit $f$ une application de $\R$ dans $\R$, telle que $f(x+y)=f(x)+f(y)$
et $f(xy)=f(x)f(y)$ pour tous $x,y\in \R$. On va montrer que $f$ est soit
nulle, soit la fonction identité. 
\begin{enumerate}
\item Remarquer que $f(x)\geq0$ si
$x\geq0$ et ainsi, que $f$ est croissante.

\item Montrer que pour tout $x$ réel on peut construire une suite $(r_k)$ et une
suite $(s_k)$ de rationnels telles que $r_k\uparrow x$ et $s_k\downarrow x$.
En déduire le résultat.
\end{enumerate}
\finenonce{006067}



\finexercice


\exercice{6068, queffelec, 2011/10/16}
\enonce{006068}{}
Soit $f$ une application continue de $\R$ dans $\R$. On rappelle que $t$
est une période de $f$ si $f(x+t)=f(x)$ pour tout $x$ réel. Soit $E$ le
groupe des périodes de $f$, supposé non vide et $T=\inf \{t\in E\ ;\
t>0,\}$.
\begin{enumerate}
\item Montrer que si $T=0$ alors $f$ est constante.

\item Si $T>0$, $f$ est $T$-périodique et $E={\Zz}.T$.
\end{enumerate}
 \finenonce{006068}



\finexercice


\exercice{6069, queffelec, 2011/10/16}
\enonce{006069}{}
Soit $f$ une application de $\R$ dans $\R$ et $\omega$ sa fonction
oscillation définie pour $x_0\in \R$ et $\delta>0$ par 
$$\omega(x_0,\delta)=\sup_{\{|x_0-y|=\delta,|x_0-z|=\delta\}}|f(y)-f(z)|.$$

\begin{enumerate}
\item Remarquer que $f$ est continue en $x_0$ si et seulement si
$$\omega(x_0)=\inf_{\delta>0}\omega(x_0,\delta)=0.$$ 

\item Montrer que pour tout $\varepsilon >0$, $O_\varepsilon=\{x\ ;\
\omega(x)<\varepsilon\}$ est un ouvert.

En déduire que $C(f)$, l'ensemble des points de continuité de $f$, est un
$G_\delta$. 
\end{enumerate}
\finenonce{006069}



\finexercice


\exercice{6070, queffelec, 2011/10/16}
\enonce{006070}{}
Existe-t-il une application continue $f$ de $[0,1]$ dans $\Rr$, telle que
$f(x)$ soit rationnel si $x$ est irrationnel, et  $f(x)$ irrationnel si $x$
est rationnel ?
\finenonce{006070}



\finexercice       


\exercice{6071, queffelec, 2011/10/16}
\enonce{006071}{}
On note pour tout $x\in \Rr, \ \varphi(x)=$ dist$(x,\Zz)$. 

\begin{enumerate}  
\item Montrer que la fonction $\varphi$ est continue, $1$-périodique, et étudier
la fonction  $f$ telle que
$$f(x)=\sum_n {\varphi(2^n x)\over{2^n}}.$$

\item On fixe $x_0\in \Rr$, et on considère les deux suites de terme 
$$z_k={1\over 2^k}E(2^kx_0),\ \ \ y_k=z_k+{1\over 2^k}.$$
Montrer que la suite $(z_k)$ cro\^\i t vers $x_0$ et que la suite $(y_k)$
décro\^\i t vers $x_0$.
Calculer ${f(z_k)-f(y_k)\over z_k-y_k}$ et en déduire que $f$ n'est pas dérivable
en $x_0$. 

 On a ainsi construit une fonction continue, nulle part dérivable.
\end{enumerate}
\finenonce{006071}



\finexercice       


\exercice{6072, queffelec, 2011/10/16}
\enonce{006072}{} 
Soit $X$ un ensemble infini muni de la topologie dont les seuls ouverts sont :
l'ensemble vide, et les parties de complémentaire fini. Montrer que si $Y$ est
un espace séparé, toute application continue de $X$ dans $Y$ est constante.
\finenonce{006072}



\finexercice


\exercice{6073, queffelec, 2011/10/16}
\enonce{006073}{}
Soit $E$ et $F$ deux espaces vectoriels normés et on note $B_E$ la boule unité
fermée de $E$. Soit $u$ une application de $E$ dans $F$ telle que 

\quad (i) $u(x+y)=u(x)+u(y),\ \forall x,y\in E$.

\quad (ii) $u(B_E)$ est bornée dans $F$.
\begin{enumerate}
\item Calculer $u(rx)$, $x\in E$, $r$ rationnel.

\item Montrer que $u$ est continue en $0$, plus précisément : 

$$\exists M>0\ ;\ \forall x\not=0\ \ ||u(x)||\leq M||x||.$$

\item Montrer que $u$ est continue et linéaire.
\end{enumerate}
\finenonce{006073}



\finexercice


\exercice{6074, queffelec, 2011/10/16}
\enonce{006074}{}
Soit $O$ un ouvert de l'espace topologique produit $X\times Y$. Montrer que 
pour tout $x\in X$, l'ensemble $A_x=\{y\in Y/(x,y)\in O\}$ est un ouvert de $Y$.
Le vérifier sur $\{(x,y)\in {\Rr}^2/ xy>1,\  x+y<4\}$.
\finenonce{006074}



\finexercice


\exercice{6075, queffelec, 2011/10/16}
\enonce{006075}{}
Montrer que si $f$ est continue de $X$ dans $Y$, espaces
topologiques, $Y$ étant séparé, son graphe $G$ est fermé dans $X\times Y$.
Etudier la réciproque en considérant l'hyperbole équilatère.
\finenonce{006075}



\finexercice


\exercice{6076, queffelec, 2011/10/16}
\enonce{006076}{} 
Soit $f:X\to Y$, espaces topologiques.
Montrer que les conditions suivantes sont équivalentes : 

\hskip 3mm (i) $f$ est continue.  

\hskip 3mm (ii) $\overline{f^{-1}(B)}\subset f^{-1}(\overline B)$ pour toute partie $B$
de $Y$. 

\hskip 3mm (iii) ${f^{-1}(\stackrel\circ B)}\subset \stackrel\circ{f^{-1}( B)}$
pour toute partie $B$
de $Y$.

En déduire $\partial f^{-1}(B)\subset f^{-1}(\partial B)$ pour toute partie $B$
de $Y$.
\finenonce{006076}



\finexercice


\exercice{6077, queffelec, 2011/10/16}
\enonce{006077}{}
Soit $C$ l'espace des fonctions continues réelles sur $[0,1]$ muni de la métrique
$d(f,g)=\int_0^1\vert f-g\vert\ dx$, puis de la métrique
$d(f,g)=\sup_x\vert f(x)-g(x)\vert$. Vérifier que l'applica\-tion
$f\to \int_0^1 f dx$ de
$C$ dans
$\Rr$ est continue dans les deux cas.   
\finenonce{006077}



\finexercice       


\exercice{6078, queffelec, 2011/10/16}
\enonce{006078}{}
Soit $c$ l'espace des suites réelles convergentes, muni de la métrique
$d(x,y)=\sup_n\vert x(n)-y(n)\vert$. Si on désigne par $l(x)$ la limite de la
suite $x$, montrer que $l$ est une application continue de $c$ dans $\Rr$.
\finenonce{006078}



\finexercice       


\exercice{6079, queffelec, 2011/10/16}
\enonce{006079}{}
Soit $X$ un ensemble infini muni de la topologie dont les seuls ouverts sont :
l'ensemble vide, et les parties de complémentaire fini. Montrer que si $Y$ est
un espace séparé, toute application continue de $X$ dans $Y$ est constante.
\finenonce{006079}



\finexercice       


\exercice{6080, queffelec, 2011/10/16}
\enonce{006080}{}
 Soit $X$ un espace métrique et $Y$ un sous-ensemble de $X$. Montrer que $Y$ est
fermé si et seulement si il existe une application continue $f:X\to \Rr$
telle que  $Y=\{x/f(x)=0\}$.
\finenonce{006080}



\finexercice       


\exercice{6081, queffelec, 2011/10/16}
\enonce{006081}{}
Soit $f$ une application ouverte de $X$ dans ${\Rr}^n$, et $A$ une partie de
$X$. Montrer que pour tout $a$ dans l'intérieur de $A$,
$$\Vert f(a)\Vert< \sup_{x\in A}\Vert f(x)\Vert.$$
\finenonce{006081}



\finexercice       


\exercice{6083, queffelec, 2011/10/16}
\enonce{006083}{}
Soit $(X,d)$ un espace métrique; montrer que l'application $(x,y)\to
d(x,y)$ est continue sur le produit $X\times X$.
\finenonce{006083}



\finexercice


\exercice{6084, queffelec, 2011/10/16}
\enonce{006084}{}
Soit $(E,d)$ un espace métrique et $A$ une partie de $E$; retrouver les
propriétés de la fonction $d_A : x\to d(x,A)$ : 
\begin{enumerate}
\item $d_A$ est $1$-lipschitzienne;
$d(x,A)=d(x,\overline A)$ et $d_A(x)=0$ si et seulement si
$x\in\overline A$.

\item Montrer que $\{x\in E\ ; \
d(x, A)< \varepsilon\}$ est un ouvert contenant $A$.

\item Montrer que tout fermé de $E$ est un $G_\delta$ et que tout
ouvert est un $F_\sigma$.
\end{enumerate}
\finenonce{006084}



\finexercice


\exercice{6085, queffelec, 2011/10/16}
\enonce{006085}{Support d'une fonction continue}

Soit $f:E\to\R$ une fonction continue
définie sur un espace topologique $E$. On appelle support (fermé) de $f$,
$S=S(f)=\overline{\{x\in E\ ;\ f(x)\not=0\}}$.
\begin{enumerate}
\item Montrer que $S=\overline{\stackrel\circ{S}}$.

\item Réciproque. On suppose $E$ métrique et $A\subset E$ fermé vérifiant
$A=\overline{\stackrel\circ{A}}$. Montrer qu'il existe $f:E\to\R$ une
fonction continue telle que $A=S(f)$.
\end{enumerate}
\finenonce{006085}



\finexercice


\exercice{6086, queffelec, 2011/10/16}
\enonce{006086}{}
\begin{enumerate}
\item Montrer qu'un espace métrique possède une propriété forte de
séparation, à savoir : deux fermés disjoints $F_1$ et $F_2$ peuvent être
séparés par deux ouverts disjoints, en considérant $\{x/
d(x,F_1)>d(x,F_2) \}$.

\item Montrer que la propriété précédente est équivalente à l'existence d'une
fonction continue $f$ valant $0$ sur $F_1$ et $1$ sur $F_2$ (considérer
$f(x)={d(x,F_1)\over d(x,F_1)+d(x,F_2)}$).
\end{enumerate}
\finenonce{006086}



\finexercice


\exercice{6087, queffelec, 2011/10/16}
\enonce{006087}{}
Soit $(X,d)$ un espace métrique avec métrique bornée. On note $\cal F$
l'ensemble des fermés non vides de $X$, et on définit pour $A$ et $B$
dans $\cal F$,
$$\delta (A,B)=\Vert d_A-d_B\Vert_\infty$$
où $d_A$ est la fonction bornée $x\to d(x,A)$.

Montrer qu'on a défini ainsi une métrique sur $\cal F$, et que l'application
$a\to\{a\}$ est une isométrie de $X$ dans $\cal F$.

\finenonce{006087}



\finexercice


\exercice{6116, queffelec, 2011/10/16}
\enonce{006116}{}
\begin{enumerate}  
\item Montrer que $\Zz$ et $\Qq$ (munis de la topologie induite par celle de
$\Rr$) ne sont pas homéomorphes. On peut par ailleurs montrer
que deux sous-ensembles dénombrables denses de
$\Rr$ sont toujours homéomorphes.

\item Trouver un homéomorphisme de $]-1, 1[$ sur $\Rr$;
de $]-1, 1[$ sur $]a, b[$.

\item Montrer que si $I$ est un intervalle ouvert de $\Rr$, et $c$ un point
n'appartenant pas à $I$, les ensembles $I$ et $I\cup \{c\}$ ne sont pas
homéomorphes bien qu'en bijection.
\end{enumerate}
\finenonce{006116}



\finexercice       


\exercice{6117, queffelec, 2011/10/16}
\enonce{006117}{}
Soit $f$ une injection continue de $\Rr$ dans $\Rr$.
 
\begin{enumerate}  
\item Montrer à l'aide
du théorème des valeurs intermédiaires que
$f$ est stric\-tement monotone.

\item Montrer que l'image par $f$ d'un intervalle ouvert est encore un
intervalle ouvert; en déduire que $f$ est ouverte et donc un homéomorphisme
de $\Rr$ sur
$f(\Rr)$.
\end{enumerate}
\finenonce{006117}



\finexercice       


\exercice{6118, queffelec, 2011/10/16}
\enonce{006118}{}
Soit $f$ une application de $X$ dans $Y$ séparé. Montrer que si $f$ est
continue, son graphe $G$ est fermé dans $X\times Y$, et l'application
$x\to (x,f(x))$ est un homéomorphisme de $X$ sur le graphe $G$ de $f$.

Montrer sur un exemple que la réciproque est fausse en général (mais vraie si $Y$
est compact).
\finenonce{006118}



\finexercice       


\exercice{6119, queffelec, 2011/10/16}
\enonce{006119}{}
Montrer que le carré unité fermé et le disque fermé dans ${\Rr}^2$ sont
homéo\-morphes.
\finenonce{006119}



\finexercice       


\exercice{6120, queffelec, 2011/10/16}
\enonce{006120}{}
Montrer que la boule unité ouverte de ${\Rr}^n$ est homéomorphe à ${\Rr}^n$ tout entier, et que deux boules ouvertes sont homéomorphes entre elles.
\finenonce{006120}



\finexercice       


\exercice{6121, queffelec, 2011/10/16}
\enonce{006121}{}
On note $\mathbb{S}^1$ le cercle unité dans ${\Rr}^2$, et $h$ l'application de
${\Rr}$ dans $\mathbb{S}^1$: 
$t\to(\cos 2\pi t, \sin 2\pi t)$. 

\begin{enumerate}  
\item Montrer que le cercle privé d'un point, $\mathbb{S}^1\backslash
\{a\}$, est homéomorphe à l'inter\-valle $]0,1[$. 

\item Montrer que $h$ est une bijection continue de $[0,1[$ sur $\mathbb{S}^1$, mais
n'est pas un homéomorphisme.

\item Soit $f$ une application continue de $\Rr$ dans $\mathbb{S}^1\backslash
\{a\}$, cette fois plongé dans $\Cc$. Montrer que $f$ admet un ``logarithme
continu", c'est-à-dire qu'il existe $g$ continue de $\Rr$ dans $\Rr$ telle
que
$f=e^{ig}$.
\end{enumerate}
\finenonce{006121}



\finexercice       


\exercice{6122, queffelec, 2011/10/16}
\enonce{006122}{}
Soit $F$ l'application de ${\Rr}^+$ dans ${\Cc}^2$ qui à $x$ associe
$(\exp(2i\pi x), \exp(2i\pi x\sqrt2))$ dont l'image est la courbe $\gamma$.

\begin{enumerate}  
\item Montrer que $F$ est continue injective.

\item Montrer que l'adhérence de $\gamma$ dans ${\Cc}^2$ est
$\mathbb{S}^1\times\mathbb{S}^1$.

\item Montrer que $F^{-1}$ n'est continue en aucun point de $\gamma$.
\end{enumerate}
\finenonce{006122}



\finexercice       


\exercice{6123, queffelec, 2011/10/16}
\enonce{006123}{Projection stéréographique}
 Soit $S^{n-1}=\{x=(x_1,\cdots,x_n)\in {\Rr}^n/ \ 
\Vert x\Vert^2=\sum_1^nx_i^2=1\}$, la sphère unité de ${\Rr}^n$, $p$ son p\^ole
nord i.e. le point $p=(0,\cdots,0,1)$, et $A=S^{n-1}\backslash \{p\}$.

\begin{enumerate}  
\item Montrer que le ``plan" de l'équateur $E$ est homéomorphe à ${\Rr}^{n-1}$.

\item A tout point $x$ de $A$ on associe $h(x)$ le point d'intersection de la droite
issue de
$p$ passant par ce point, avec le plan $E$. Expliciter $h$, puis $h^{-1}$ et
montrer ainsi que la sphère est homéomorphe à ${\Rr}^{n-1}$.

(On établira $h(x)=p +{{x-p}\over{1-x_n}}$ et $h^{-1}(y)={{2y}\over{1+\Vert
y\Vert^2}} + p\ {{1-\Vert
y\Vert^2}\over{1+\Vert
y\Vert^2}}$).

\item En déduire un homéomorphisme de $\mathbb{S}^1$ sur $\overline{\Rr}$.
\end{enumerate}
\finenonce{006123}



\finexercice  


\exercice{6131, queffelec, 2011/10/16}
\enonce{006131}{}
\begin{enumerate}  
\item Montrer que si deux fonctions continues sur un espace topologique $X$
co\"\i ncident sur un ensemble dense dans $X$, elles sont égales.

\item Soit $f$ une fonction réelle définie continue sur $[-1,1]$. Montrer que si
pour tout $n$, $\int_{-1}^1f(x)\ x^n\ dx$ est nulle, alors $f$ est nulle.

(\emph{Indication :} Considérer l'application $g\to \int_{-1}^1f(x)\ g(x)\ dx$.)
\end{enumerate}
\finenonce{006131}



\finexercice       


\exercice{6132, queffelec, 2011/10/16}
\enonce{006132}{}
Soit $F$ un fermé de $\Rr$, et $f$ une application continue de $F$ dans $\Rr$. Montrer que $f$ se prolonge en une fonction continue sur $\Rr$ tout
entier. Peut-on remplacer ``fermé" par ``ouvert"?
\finenonce{006132}



\finexercice       


\exercice{6133, queffelec, 2011/10/16}
\enonce{006133}{}
Soit $n\to r_n$ une bijection de $\Nn$ sur $\Qq\cap [0,1]$, et $f$ la
fonction définie sur $\Qq\cap [0,1]$ par 
$$f(x) = \sum_{r_n<x} 2^{-n}.$$
Montrer que $f$ est continue, mais qu'elle ne peut être prolongée en aucune
fonction continue sur $[0,1]$.
\finenonce{006133}



\finexercice       


\exercice{6134, queffelec, 2011/10/16}
\enonce{006134}{}
Soit $(X,d)$ un espace métrique; on rappelle tout d'abord les propriétés de la
fonction
$d_A : x\to d(x,A)$ où $A$ est une partie de $X$ : 

\begin{enumerate}  
\item $d_A$ est $1$-lipschitzienne, et $d_A(x)=0$ si et seulement si
$x\in\overline A$. On en déduit que tout fermé est un $G_\delta$ et que tout
ouvert est un $F_\sigma$. 

\item Montrer qu'un espace métrique possède une propriété forte de séparation, à
savoir : deux fermés disjoints $F_1$ et $F_2$ peuvent être séparés par deux
ouverts disjoints, en considérant $\{x/ d(x,F_1)>d(x,F_2) \}$.

\item Montrer que la propriété précédente est équivalente à l'existence d'une
fonction continue $f$ valant $0$ sur $F_1$ et $1$ sur $F_2$.

\item Soit $F_1$, $F_2$,...,$F_n$, $n$ fermés disjoints dans $X$, et $c_1$,
$c_2$,...$c_n$, $n$ nombres réels. Montrer que la fonction $f$ valant $c_i$
sur $F_i$ peut se prolonger en une fonction continue à $X$ tout entier. 
\end{enumerate}
\finenonce{006134}



\finexercice       


\exercice{6135, queffelec, 2011/10/16}
\enonce{006135}{}
Soit $(X,d)$ un espace métrique, et $Y$ un sous-espace non vide de $X$. On va
montrer que toute fonction $f:Y\to \Rr$, $k$-lipschitzienne, admet un
prolongement $g:X\to\Rr$ qui est aussi $k$-lipschitzien. Soit donc $f$ ainsi;
pour tout $x\in X$ et $y\in Y$ , on pose
$$f_y(x)=f(y)+kd(x,y).$$

\begin{enumerate}  
\item Montrer que pour $x$ fixé, l'ensemble $\{f_y(x)\}$ lorsque $y$ parcourt $Y$
est minoré. On pose $g(x) = \inf_{y\in Y}\{f_y(x)\}$.

\item Montrer que l'application $g$ ainsi définie sur $X$, réalise un prolongement
$k$-lipschitzien de $f$.

\item Donner une condition suffisante pour que ce prolongement soit unique.
\end{enumerate}
\finenonce{006135}



\finexercice  


\exercice{6241, queffelec, 2011/10/16}
\enonce{006241}{}
\begin{enumerate}
\item Montrer qu'une fonction de $(X,d)$ dans $(Y,\delta)$ n'est pas uniformément
continue, si et seulement si on peut trouver $\varepsilon>0$ et deux suites de
points de
$X$,
$(x_n)$ et $(y_n)$ vérifiant

(i) $d(x_n,y_n)$ tend vers $0$.

(ii) $\delta(f(x_n),f(y_n))\geq\varepsilon$.

\item Parmi les fonctions de variable réelle suivantes, lesquelles sont
unifor\-mément continues : $\sin(x^2)$, $x\sin x$, $\sin {1\over
x},\  x\neq0$ ?
\end{enumerate}
\finenonce{006241}



\finexercice


\exercice{6242, queffelec, 2011/10/16}
\enonce{006242}{}
Soit $E=C_b(\Rr)$ muni de la norme uniforme; pour $f\in E$, on note
$f_a$ la translatée de $f$ par $a$, ie la fonction $x\to f(x-a)$, et $O_f$
l'ensemble des translatées de $f$. Soit
$f$ une fonction continue périodique de $\Rr$ dans $\Rr$; 

\begin{enumerate}
\item  Montrer que $f$
est uniformément continue sur
$\Rr$.

\item Montrer que $O_f$ est compact et connexe (considérer
l'application $a\to f_a$).
\end{enumerate}
\finenonce{006242}



\finexercice


\exercice{6243, queffelec, 2011/10/16}
\enonce{006243}{}
Soit $(f_n)$ une suite d'applications croissantes de $[0,1]$
dans $\Rr$, qui converge simplement vers une fonction $f$ continue.
Montrer que la convergence est uniforme sur $[0,1]$. 
\emph{Indication :} $\varepsilon >0$
étant fixé, montrer qu'il existe $0=x_0<x_1<\cdots<x_k=1$ tels que
$f(x_{j+1})-f(x_j)\leq\varepsilon, \ 1\leq j\leq k-1$ et établir
$|f(x)-f_n(x)|\leq\sup_j|f_n(x_j)-f(x_j)|+\varepsilon$.
\finenonce{006243}



\finexercice


\exercice{6244, queffelec, 2011/10/16}
\enonce{006244}{} 
Parmi les métriques suivantes définies sur $\Rr$, lesquelles sont
uniformément équivalentes à la métrique usuelle ?
\begin{enumerate}  
\item $\vert x^3-y^3\vert$

\item $\vert \arctan x-\arctan y\vert$

\item ${{\vert x-y\vert}\over{1+\vert x-y\vert}}$
\end{enumerate}
\finenonce{006244}



\finexercice       


\exercice{6245, queffelec, 2011/10/16}
\enonce{006245}{}
Soit $d_1$ et $d_2$ deux distances sur un espace $X$. On considère les
quatre assertions suivantes :

   (i) Les métriques sont topologiquement équivalentes.

  (ii) Les métriques sont uniformément équivalentes.
 
  (iii) Les métriques sont Lipschitz-équivalentes (il existe $A$ et $B$
constantes telles que $A\, d_1\leq d_2\leq B\, d_1$).

  (iv) $(X,d_1)$ et $(X,d_2)$ sont simultanément complets.

  
Etablir les implications entre ces propriétés et donner des contre-exemples
lorsque les implications  n'ont pas lieu.
\finenonce{006245}



\finexercice       


\exercice{6246, queffelec, 2011/10/16}
\enonce{006246}{}
Soit $d_1$ et $d_2$ deux distances sur un espace $X$. Montrer qu'elles sont
unifor\-mément équivalentes si et seulement si $(X,d_1)$ et $(X,d_2)$ ont les
mêmes applica\-tions réelles uniformément continues.

\emph{Indication :} Raisonner par contraposition et considérer pour $(x_n)$ et $(y_n)$
vérifiant
$\lim d_1(x_n,y_n)=0$ et
$d_2(x_n,y_n)\geq\varepsilon$,
$A=\overline{\{x_1,x_2,...\}}$, $B=\overline{\{y_1,y_2,...\}}$ dans $(X,d_2)$,
et 
$$f(x)={{d_2(x,A)}\over{d_2(x,A)+d_2(x,B)}}$$
\finenonce{006246}



\finexercice       


\exercice{6247, queffelec, 2011/10/16}
\enonce{006247}{}
Soit $(f_n)$ une suite de fonctions réelles convergeant uniformément vers $f$
sur $\Rr$ et soit $g$ une fonction uniformément continue sur $\Rr$. Montrer
que la suite $(g\circ f_n)$ converge uniformément vers $g\circ f$ sur
$\Rr$.
\finenonce{006247}



\finexercice       


\exercice{6248, queffelec, 2011/10/16}
\enonce{006248}{}
Soit $X$ un espace métrique.

\begin{enumerate}  
\item Montrer que si $X$ n'est pas complet, 
il existe une suite de Cauchy $(a_n)$,
non convergente, et telle que $a_p\neq a_q$ pour $p\neq q$.

\item Soit $(b_n)$ une suite de Cauchy non convergente; montrer que l'ensemble
$B=\{b_n,\ n\in\Nn\}$ est fermé dans $X$.

\item Déduire des questions précédentes que si $X$ n'est pas complet, on peut
trouver une fonction continue $f:X\to [0,1]$ qui n'est pas uniformément
continue.

\emph{Indication :} Si $(a_n)$ est définie par 1., construire $f:X\to [0,1]$ telle que
$f(a_{2n})=0$ et $f(a_{2n+1})=1$.
\end{enumerate}
\finenonce{006248}



\finexercice       


\exercice{6249, queffelec, 2011/10/16}
\enonce{006249}{}
Soit $f$ une application bijective d'espaces métriques $f:X\to Y$ uniformément
continue et d'inverse continue. Montrer que si $Y$ est complet, $X$ l'est aussi.
\finenonce{006249}



\finexercice       


\exercice{6250, queffelec, 2011/10/16}
\enonce{006250}{}
Soit $(X,d)$ et $(Y,\delta)$ deux espaces métriques; soit $f$ une application
surjective de $X$ sur $Y$ telle que  $\delta(f(x),f(x'))=d(x,x')$ pour tous
$x,x'$ dans $X$. Vérifier que $f$ est un homéomorphisme uniformément continu
ainsi que $f^{-1}$. Donner des exemples sur $\Rr^n$ et décrire les isométries
de $\Rr$.
\finenonce{006250}



\finexercice       


\exercice{6251, queffelec, 2011/10/16}
\enonce{006251}{}
On considère $l^1$ et $l^2$ les espaces de suites réelles
absolument et de carré sommables, et l'application
$F$ (non linéaire) de  $l^1$ dans $l^2$ définie par
$F(a)=b$ si $a=(a_n)$, $b=(b_n)$ avec $b_n= {\rm sign}\,(a_n)\sqrt{\vert
a_n\vert}$. Vérifier que $F$ est un homéomor\-phisme de $l^1$ sur $l^2$,
uniformément continu mais d'inverse non uniformément continu.
\finenonce{006251}



\finexercice   


\exercice{6775, gijs, 2011/10/16}
\enonce{006775}{}
\begin{enumerate}
\item
(D'abord un cas particulier) Dans $\Rr^3$ on
considère les objets suivants~: le demi-espace $D =
\{\,(x,y,z)\in \Rr^3 \mid z<1\,\}$, le plan $P =
\{\,(x,y,z)\in \Rr^3 \mid z=0\,\}$, la sphère
unité $S^2 = \{\, (x,y,z) \in \Rr^3 \mid x^2 + y^2 +
z^2 - 1 = 0\,\}$, et le pole nord $N = (0,0,1)$. 

\begin{enumerate}
\item On définit une application $p: D \to P$ par la
procédure suivante~: pour un point $A \in D$,
la droite dans $\Rr^3$ qui passe par $A$ et $N$
coupe le plan $P$ en $p(A)$. Trouver
l'expression explicite de l'application $p$ et en
déduire qu'elle est continue.
\item On définit une application $i : P \to S^2$ par
la procédure suivante~: pour $B\in P$, la droite
dans $\Rr^3$ qui passe par $B$ et $N$
coupe la sphère unité $S^2$ en $i(B)$. Trouver
l'expression explicite de l'application $i$ et en
déduire qu'elle est continue.
\item En utilisant les applications $p$ et $i$,
montrer que $P$ est homéomorphe à $S^2 \setminus
\{N\}$.
\item Montrer que $S^2$ est compacte.
\end{enumerate}

\item
(Le cas général) Soit $X$ un espace topologique
Hausdorff et $\infty$ un élément qui n'appartient pas
à $X$. On définit $\widehat X = X \cup \{\infty\}$ et
on dit qu'un sous-ensemble $U \subset \widehat X$ est
ouvert si et seulement si~: ou bien $\infty \notin U$ et
$U$ ouvert dans $X$, ou bien $\infty \in U$ et $X\setminus
U$ est quasi compact dans $X$. 

\begin{enumerate}
\item Montrer que si $U$ est un ouvert de $\widehat X$
contenant $\infty$, alors $X\setminus U$ est fermé dans
$X$.
\item Montrer que les ouverts dans $\widehat X$
forment bien une topologie.
\item Montrer que $\widehat X$ avec la topologie
décrite ci-dessus est quasi compact. 
\item En considérant $X \subset \widehat X$, on
donne $X$ la topologie induite par $\widehat X$. Montrer
que cette topologie co\"\i ncide avec la topologie de
départ de $X$.
\end{enumerate}
\end{enumerate}
\finenonce{006775}



\finexercice

\exercice{6816, gijs, 2011/10/16}
\enonce{006816}{}
Soit $X$ un espace topologique et $Y$ un espace
topologique séparé. Soit $f,g : X \to Y$  deux
applications continues.
\begin{enumerate}
\item Démontrer que $U = \{x\in X \mid f(x) \neq
g(x)\}$ est un ouvert de $X$.

\item Soit $D\subset X$ une partie dense.
Démontrer que si $f\vert_D = g\vert_D$, alors $f=g$.
\end{enumerate}
\finenonce{006816}



\finexercice

\section{423.00 Application linéaire bornée}


\exercice{2361, mayer, 2003/10/01}
\enonce{002361}{} 
Soient $E_1, E_2$ et $F$ des espaces norm\'es sur $\Rr$
et soit $B: E_1 \times E_2 \to F$ une application bilin\'eaire.
Montrer que $B$ est continue si et seulement s'il existe
$M>0$ tel que
$$ \|B(x) \|\leq M \| x_1\| \| x_2\|\quad  \text{pour tout } x=(x_1,x_2) \in E_1\times E_2 \; .$$
\finenonce{002361}


\finexercice


\exercice{2362, mayer, 2003/10/01}
\enonce{002362}{} 
Soient $E$ et $F$ deux espaces norm\'es et $L:E\to F$ une application lin\'eaire
v\'erifiant:
\emph{$(L(x_n))_n$ est born\'ee dans $F$ pour toute suite $(x_n)_n$ de $E$
tendant vers $0\in E$.}
Montrer que $L$ est continue.
\finenonce{002362}


\finexercice


\exercice{2363, mayer, 2003/10/01}
\enonce{002363}{} 
Soient $E$ et $F$ deux espaces norm\'es r\'eels et $f:E\to F$ une application born\'ee sur
la boule unit\'e de $E$ et v\'erifiant
$$f(x+y)=f(x) +f(y)\quad  \text{pour tout } x, y \in E \; .$$
Montrer que $f$ est lin\'eaire continue.

\finenonce{002363}



\finexercice


\exercice{2364, mayer, 2003/10/01}
\enonce{002364}{} 
Calculer la norme des op\'erateurs suivants:
\begin{itemize}
\item Le shift sur $l^\infty$ d\'efini par $S(x)_{n+1}= x_n$,
$S(x)_0=0$.
\item $X={\cal C}([0,1])$ muni de la norme $\|.\|_\infty$ et $Tf(x) = f(x)g(x)$ o\`u $g\in X$.
\end{itemize}
Calculer la norme des formes lin\'eaires suivantes:
\begin{itemize}
\item $X={\cal C}([0,1])$  muni de la norme $\|.\|_\infty$ et $u(f) = \int _0^1 f(x)g(x) \; dx$ o\`u $g\in X$ est
une fonction qui ne s'annule qu'en $x=1/2$.
\item $X= l^2$ et $u(x) = \sum a_nx_n$ o\`u $(a_n)$ est dans $X$.
\item $X= l^1$ et $u(x) = \sum a_nx_n$ o\`u $(a_n)$ est dans $l^\infty$.
\item $X$ l'espace des suites convergentes muni de la norme sup et $u:X\to \Rr$
l'application $u(x) = \lim_{j\to \infty} x_j$.
\end{itemize} 
\finenonce{002364}


\finexercice


\exercice{2365, mayer, 2003/10/01}
\enonce{002365}{} 
\label{exopol}
Soit $X = \Rr[x]$ l'ensemble des polyn\^omes. Pour $P(x) = \sum _{k=0}^p a_k x^k$
on pose $\|P\|= \sup_k |a_k|$, $U(P)(x) = \sum _{k=1}^n \frac{1}{k} a_k x^k$ et
$V(P)(x) = \sum _{k=1}^n k a_k x^k$.
\begin{enumerate}
\item Montrer que $\|.\|$ d\'efinit une norme et que $U$ et $V$ d\'efinissent des applications
lin\'eaires de $X$ dans $X$.
\item Examiner si $U$ et $V$ sont continues?
\end{enumerate}
\finenonce{002365}


\finexercice


\exercice{2366, mayer, 2003/10/01}
\enonce{002366}{} 
Soit $l^\infty$ l'espace des suites r\'eelles muni avec la norme
uniforme, i.e. $\|x\| _\infty = \sup _n |x_n|$. On consid\'ere
l'application $A:l^\infty \to l^\infty$ d\'efinie par
$$A(x_1,x_2,...,x_n,...) = (x_1, x_2/2,..., x_n/n,...) \; .$$ Montrer
que :
\begin{enumerate}
\item $A$ est injective et continue avec $\|A\| =1$. Par contre, $A$
n'est pas surjective.
\item $A$ admet un inverse \`a gauche mais qu'il n'est pas continu.
\end{enumerate}
\finenonce{002366}


\finexercice


\exercice{2367, mayer, 2003/10/01}
\enonce{002367}{} 
\label{exoferm}
Soit $X$ un espace norm\'e, $L:X\to \Rr$ une forme lin\'eaire non nulle
et $H = L^{-1}(\{0\})$ son noyau.
\begin{enumerate}
\item Montrer que, si $L$ est continue, alors $H$ est un sous-espace ferm\'e dans $X$.
\'Etablir la relation
$$ \mathrm{dist}(a,H) = \frac{|L(a)|}{\|L\|} \quad \text{pour tout } a\in X \; .$$
\item R\'eciproquement, supposons que le noyau $H$ est un ferm\'e. D\'emontrer alors
que $\mathrm{dist}(a,H)>0$ d\`es que $a\in X\setminus H$ et en d\'eduire que $L$ est continue
de norme au plus $|L(a)|/\mathrm{dist}(a,H)$.
\item Peut-on g\'en\'eraliser ceci a des applications lin\'eaires entre espaces
norm\'es?
\end{enumerate}
\finenonce{002367}


\finexercice


\exercice{2368, mayer, 2003/10/01}
\enonce{002368}{} 
Soit $X={\cal C}([0,1])$ avec la norme $\| f\| = \int _0^ 1 |f(t)|\, dt$.
Montrer que la forme lin\'eaire $f\in X \mapsto f(0)\in \Rr$ n'est pas continue.
Que peut-on en d\'eduire pour le sous-espace des fonctions de $X$ nulles en 0?
\finenonce{002368}


\finexercice


\exercice{2369, mayer, 2003/10/01}
\enonce{002369}{} 
Soit $X = \{f\in {\cal C}(\Rr ) \; ; \; (1+x^ 2 ) |f(x)| \text{ soit born\'ee}\}$. On pose
$N(f) = \sup_{x\in \Rr}(1+x^ 2 ) |f(x)|$. V\'erifier que $N$ est une norme, puis
montrer que la forme lin\'eaire suivante $L$ est continue et calculer sa norme:
$$ L:X\to \Rr \quad \mbox{d\'efinie par} \quad L(f) = \int _\Rr f(x) \; dx \; .$$

\finenonce{002369}


\finexercice


\exercice{6204, queffelec, 2011/10/16}
\enonce{006204}{}
On désigne par $E$ l'espace $C([-1,1])$ muni de la norme uniforme et
par $T$ la forme linéaire
 définie par $$Tf=\int_{-1}^1 \sin(\pi t)\ f(t)\ dt $$
pour $f\in E$. Vérifier que $T$ est continue et calculer la norme de $T$.
\finenonce{006204}



\finexercice


\exercice{6205, queffelec, 2011/10/16}
\enonce{006205}{}
Soit $E=C([0,1])$, $\mu(x)=\int_0^1 x(t)\ dt$, $\mu_n(x)={1\over n}\sum_{k=1}^n
x({k\over n})$.
\begin{enumerate}
\item Calculer $\Vert\mu\Vert$ et $\Vert\mu_n\Vert$.

\item Montrer que $\mu_n(x)$ converge vers $\mu(x)$ pour toute $x$ dans $E$, mais
que $\Vert\mu-\mu_n\Vert=2$.
\end{enumerate}
\finenonce{006205}



\finexercice


\exercice{6206, queffelec, 2011/10/16}
\enonce{006206}{}
On désigne par $E$ l'espace $C([0,1])$ muni de la norme uniforme et l'opérateur
$A$ défini par $$Af(x)=\int_0^x tf(t)\ dt + x\int_x^1 f(t)\ dt$$
pour $f\in E$ et $x\in [0,1]$.
\begin{enumerate}
\item Vérifier que $A$ est continu et calculer sa norme opérateur.

\item L'équation $Af=f$ a-t-elle dans $E$ des solutions $f$ non nulles ?
\end{enumerate}
\finenonce{006206}



\finexercice


\exercice{6207, queffelec, 2011/10/16}
\enonce{006207}{}
Soit $K$ un compact convexe d'un evn $E$. Soit $u$ une application linéaire
continue de $E$ dans $E$ telle que $u(K)\subset K$. On va montrer que $u$ a un
point fixe dans $K$.
\begin{enumerate}
\item On peut supposer que $0\notin K$. Pour chaque $n\geq1$, on désigne par $S_n$
l'application définie sur $E$ par 

$S_n(x)={1\over n}(x+u(x)+\cdots+u^{n-1}(x))$.
Montrer que $S_n(K)\subset K$.

\item Montrer que pour tous entiers $n_1,n_2,\ldots,n_k$ en nombre fini,

$S_{n_1}\circ\cdots\circ S_{n_k}(K)\subset S_{n_1}(K)\cap
S_{n_2}(K)\cap\cdots\cap S_{n_k}(K)$.

En déduire que $A=\cap_{n\geq1} S_{n}(K)$ est non vide.

\item Montrer que tout $x\in A$ est  point fixe de $u$.
\end{enumerate}




\finenonce{006207}


\finexercice


\exercice{6208, queffelec, 2011/10/16}
\enonce{006208}{}
\begin{enumerate}
\item Montrer que pour tout $x\in\Rr^n$, $||x||=\sup\{ |\langle x,y \rangle|\ ; \
||y||\leq1\}$.

\item Montrer que l'espace des formes linéaires ${\cal L}(\Rr^n,\Rr)$ sur 
$\Rr^n$ (plus généra\-lement ${\cal L}(E,\Rr)$ où $\dim E$ finie) est
isométriquement isomorphe à $\Rr^n$ (ou $E$).
\end{enumerate}
\finenonce{006208}



\finexercice


\exercice{6209, queffelec, 2011/10/16}
\enonce{006209}{}
Sur $M_n(\Rr)$ on note $|A|=\sup_{||x||\leq1}||Ax||$ la norme opérateur de la
matrice $A$, où $||x||$ désigne la morme euclidienne de $x$. Montrer que
$$|A|=\sup_{||x||\leq1,||y||\leq1}|\langle Ax,y \rangle |$$
et en déduire que $|A|\leq \Big(\sum|a_{ij}|^2\Big)^{1\over2}$.
\finenonce{006209}



\finexercice


\exercice{6210, queffelec, 2011/10/16}
\enonce{006210}{}
Soit $E=\Rr^n$ muni de sa norme euclidienne et $f$ une application continue de
$[0,1]$ dans
$\Rr^n$.

Montrer que  $||\int_{[0,1]}f(t)\ dt||\leq \int_{[0,1]}||f(t)||\ dt$. 
\finenonce{006210}



\finexercice

\section{424.00 Espace vectoriel normé}


\exercice{1865, gineste, 2001/11/01}

\enonce{001865}{Normes sur $\R^2$}
Pour tout $(x,y) \in
\R^2$, on pose $N_1(x,y)=\text{Max} (\sqrt{x^2+y^2},|x-y|)$ et
$N_2(x,y)=\sqrt{ x^2/9 + y^2/4 }.$
\begin{enumerate}
\item Montrer que $N_1$ et $N_2$ sont des
normes sur $\R^2$ et repr\'esenter les boules unit\'es ferm\'ees associ\'ees \`a ces normes.
\item Montrer que $N_2 \leq \| . \|_\infty \leq \| . \|_2 \leq N_1 \leq \|.\|_1 \leq 4N_2.$
\item D\'eterminer le plus petit r\'eel $k>0$, tel que $\|.\|_1 \leq kN_2.$ (utiliser Cauchy-Schwarz)
\end{enumerate}
\finenonce{001865}



\finexercice



\exercice{1869, roussel, 2001/09/01}

\enonce{001869}{}
On consid\`ere les trois normes d\'efinies sur $\mathbb{R}^2$ par:
$$%\begin{array}{llll}
%\mbox{a)}
\| {X}\|_{1} = |x_1|+|x_2|~~,~~ %\vspace*{0.3cm} \\ \mbox{b)}
\| {X}\|_{2} = (x_1^2+x_2^2)^{1\over 2}~~,~~ %\vspace*{0.3cm}\\ \mbox{c)}
\| {X}\|_\infty =  \max \{|x_1|,|x_2| \}. %\\ \end{array}
$$
Repr\'esenter graphiquement les boules unit\'es de chacune d'entre
elles. Peut-on ``comparer" ces trois normes? Ecriver les
d\'efinitions des distances $d_1$,$d_2$ et $d_{\infty}$
associ\'ees \`a chacune d'entre elles.
\finenonce{001869}



\finexercice



\exercice{1870, roussel, 2001/09/01}

\enonce{001870}{}
Soit $E$ l'espace vectoriel
des fonctions \`a valeurs dans $\mathbb{R}$, d\'efinies et
\hspace*{5cm}continues sur $[ \mbox{-1,1}]$.
\begin{enumerate}
\item Montrer que les trois applications suivantes sont des normes sur $E$:
$$%\begin{array}{lcl}
f\longrightarrow \| {f}\|_{1} = \int_{-1}^{+1}|f(x)|dx ,~~~%\vspace*{0.3cm} \\
f\longrightarrow \| {f}\|_{2} = (\int_{-1}^{+1}f^2(x)dx) ^{1\over 2}$$
%\vspace*{0.3cm}\\
$$f\longrightarrow \| {f}\|_\infty = \sup_{x\in [-1,+1]}\{ |f(x)|\}
%\end{array}
$$
\item On consid\`ere la suite $(f_n)_{n\in N^*}$ de fonctions d\'efinies par
{$f_n(x)=
  \left\{
 \begin{array}{lll}
-1 &\mbox{ si }x\in [-1,-{1\over n}]\vspace*{0.3cm} \\
nx &\mbox{ si }x\in ]-{1\over n},{1\over n}] \vspace*{0.3cm} \\
1 &\mbox{ si }x\in ]{1\over n},1]  \\
\end{array}
\right.
$} \\

La suite $f_n$ est-elle de Cauchy dans $(E,\|{.}\|_{1})$, $(E,\|{.}\|_{2})$ et
dans $(E,\| {.}\|_\infty)$? Conclusions?
\end{enumerate}
\finenonce{001870}



\finexercice



\exercice{1871, roussel, 2001/09/01}

\enonce{001871}{}
Soit $E$ l'espace vectoriel des fonctions \`a valeurs dans $\mathbb{R}$,
d\'efinies, continues et d\'erivables sur  $[ \mbox{0,1}]$ et
v\'erifiant $f(0)=0$. On d\'efinit sur cet espace les deux normes suivantes :
$$N_1(f)= \| {f}\|_\infty \mbox{~et~} N_2(f)= \| {f'}\|_\infty.$$
\begin{enumerate}
\item Montrer que $N_1(f) \leq N_2(f)$. En d\'eduire que l'application
identique de $(E, N_2)$ vers $(E, N_1)$ est continue.
\item A l'aide de la fonction $f_n(x)={x^n \over n}$, montrer que l'application
identique de $(E, N_1)$ vers $(E, N_2)$ n'est pas continue.
\end{enumerate}
\finenonce{001871}



\finexercice



\exercice{1872, roussel, 2001/09/01}

\enonce{001872}{}
Lorsqu'un espace vectoriel $E$ est en outre muni d'une multiplication,
 l'application
 $\nolinebreak {N:\ E \rightarrow \mathbb{R}}$  est dite norme
multiplicative si:
\begin{itemize}
\item $N$ est une norme,
\item pour tous $A$ et $B$ dans $E$, $N(A.B)\leq N(A).N(B)$.
\end{itemize}
Soit $E=M_n(\mathbb{R})$, l'espace vectoriel des matrices carr\'ees \`a $n$
lignes et $n$ colonnes. $A\in E$ se  note
$A=(a_{i,j})_{1\leq i,j\leq n}$
\begin{enumerate}
\item Montrer que $\displaystyle N_{\infty }(A)=
\max_{1\leq i\leq n}\{\  \sum_{j=1}^n|a_{i,j}|\ \}$ d\'efinit une norme
multiplicative sur $E$.
\item Montrer que $\displaystyle N_{\infty }(A)=\max_{\{ X\in \mathbb{R}^n,\ \|{X}\|_\infty=1\} }
\{\  \| {A.X}\|_\infty \}$.
\item Soit $A\in M_n(\mathbb{R})$ telle que $\displaystyle \forall \ 1\leq i \leq n,\ |a_{i,i}|>
\sum_{j=1,j\neq i}^n|a_{i,j}|$ et $D$ la matrice diagonale form\'ee avec les
\'el\'ements diagonaux de $A$. Soit aussi $F$ un vecteur de $\mathbb{R}^n$.
On consid\`ere la suite des $X^{(p)}\in \mathbb{R}^n$ d\'efinie pour $p\geq 0$ par:
$$
  \left\{
 \begin{array}{lll}
X^{(0)} &=X_0\in \mathbb{R}^n &\\
X^{(p+1)} &=(I-D^{-1}A)X^{(p)}+D^{-1}F &\mbox{ pour p}\geq 0
\end{array}
\right.
$$
Montrer qu'elle est convergente et calculer sa limite.
\end{enumerate}
\finenonce{001872}



\finexercice



\exercice{1873, maillot, 2001/09/01}

\enonce{001873}{partiel 1999}
Soit $(E,\|\cdot\|)$ un espace vectoriel norm\'e, $x$ un \'el\'ement
de $E$ et $A$ un compact de $E$.
\begin{enumerate}
\item Montrer que l'application de $E$ dans $\R$ qui \`a $y$ associe
$\|y\|$ est continue.
\item Montrer que l'application de $E$ dans $\R$ qui \`a $y$ associe
$\|y-x\|$ est continue.
\item Montrer que la distance de $x$ \`a $A$ est atteinte,
c'est-\`a-dire qu'il existe $a\in A$ tel que
$$\inf_{y\in A}\|y-x\| = \|a-x\|.$$
\end{enumerate}
\finenonce{001873}



\finexercice



\exercice{1874, maillot, 2001/09/01}

\enonce{001874}{}
Soient $(E,\|\cdot\|_E)$ et $(F,\|\cdot\|_F)$ deux espaces vectoriels
normés. Soit $L$ une application linéaire de $E$ dans $F$.
\begin{enumerate}
\item Montrer que $L$ est continue en $0$ si et seulement si elle est
continue en tout point de $E$.
\item On suppose qu'il existe une constante $K>0$ telle que
$$\|L(x)\|_F\le K\|x\|_E \qquad \forall x\in E.$$
Montrer que $L$ est continue.
\item Dans la suite, on suppose que $L$ est continue et on pose
$$K=\sup_{\|x\|_E=1} \|L(x)\|_F.$$
\begin{enumerate}
\item Supposons que $K=+\infty$. Montrer qu'alors il existe une
suite $(x_n)$ dans $E$ telle que $\|x_n\|=1$ pour tout $n$ et telle
que $\|L(x_n)\|_F$ tend vers $+\infty$. En déduire qu'il existe une
suite $y_n$ tendant vers $0$ et telle que $\|L(y_n)\|_F=1$.
\item En déduire que $K\in \R_+$ et que pour tout $x\in E$ on a
$$\|L(x)\|_F\le K\|x\|_E.$$
\end{enumerate}
\end{enumerate}
\finenonce{001874}



\finexercice



\exercice{1875, maillot, 2001/09/01}

\enonce{001875}{}
Soit $E$ l'espace vectoriel des fonctions continues de $[-1,1]$ à
valeurs dans $\R$ muni de la norme $$\|f\|_1=\int_0^1 |f(x)| \,dx.$$
On considère l'application $L:E\rightarrow \R$ définie par $L(f)=f(1)$.
\begin{enumerate}
\item Montrer que $L$ est une application linéaire.
\item En considérant les fonctions $f_n:x\mapsto \sqrt n x^n$,
montrer que $L$ n'est pas continue.
\end{enumerate}
\finenonce{001875}



\finexercice



\exercice{1876, maillot, 2001/09/01}

\enonce{001876}{}
Soit $(E,\|\cdot\|)$ un espace vectoriel normé et $(x_n)_{n\in\N}$ une
suite d'éléments de $E$. On suppose que $(x_n)$ est de Cauchy.
Montrer qu'elle converge si et seulement si elle admet une sous-suite
convergente.
\finenonce{001876}



\finexercice



\exercice{1877, maillot, 2001/09/01}

\enonce{001877}{}
Soit $E$ l'espace vectoriel des fonctions continues de $[-1,1]$ à
valeurs dans $\R$. On définit une norme sur $E$ en posant
$$\|f\|_1=\int_{-1}^1 |f(t)| \,dt.$$ On va montrer que $E$ muni de
cette norme n'est pas complet. Pour cela, on définit une suite
$(f_n)_{n\in\N^*}$ par
\[f_n(t)=\begin{cases} -1 &\text{si\ } -1\le t \le -\frac1n\\
               nt &\text{si\ } -\frac1n\le t \le \frac1n\\
               1  &\text{si\ \ \ \ \ } \frac1n \le t\le 1.
          \end{cases}\]

\begin{enumerate}
\item Vérifier que $f_n\in E$ pour tout $n\ge 1$.
\item Montrer que $$\|f_n-f_p\|\le \sup(\frac2n,\frac2p)$$ et en
déduire que $(f_n)$ est de Cauchy.
\item Supposons qu'il existe une fonction $f\in E$ telle que $(f_n)$
converge vers $f$ dans $(E,\|\cdot\|_1)$. Montrer qu'alors on a
$$\lim_{n\rightarrow+\infty}\int_{-1}^{-\alpha} |f_n(t)-f(t)|\, dt=0
\qquad \text{et} \qquad \lim_{n\rightarrow+\infty}\int_{\alpha}^1
|f_n(t)-f(t)|\, dt=0$$ pour tout $0<\alpha<1$.
\item Montrer qu'on a $$\lim_{n\rightarrow+\infty}\int_{-1}^{-\alpha}
|f_n(t)+1|\, dt=0
\qquad \text{et} \qquad \lim_{n\rightarrow+\infty}\int_{\alpha}^1
|f_n(t)-1|\, dt=0$$ pour tout $0<\alpha<1$.
En déduire que
\begin{align*}
&f(t)=-1\qquad &\forall t\in[-1,0[\\
&f(t)=1\qquad &\forall t\in ]0,1].
\end{align*}
Conclure.
\end{enumerate}
\finenonce{001877}



\finexercice



\exercice{1878, maillot, 2001/09/01}

\enonce{001878}{}
Soit $E=\R^d$ muni d'une norme $\|\cdot\|$. On rappelle qu'une application
continue $g$ de $E$ dans $E$ est dite \emph{contractante} s'il existe
$K\in ]0,1[$ tel que
$$ \|g(x)-g(y)\|\le K \|x-y\| \qquad \forall x,y\in E.$$ On rappelle
aussi que toute application contractante admet un unique point fixe.

Soit $f$ une application continue de $E$ dans $E$ telle qu'il existe un
entier $n$ tel que $f^n$ soit contractante. On note $x_0$ le point
fixe de $f^n$.
\begin{enumerate}
\item Montrer que tout point fixe de $f$ est un point fixe de $f^n$.
\item Montrer que si $x$ est un point fixe de $f^n$, il en est de même
pour $f(x)$.
\item En déduire que $x_0$ est l'unique point fixe de $f$.
\end{enumerate}
\finenonce{001878}



\finexercice



\exercice{1879, maillot, 2001/09/01}

\enonce{001879}{}
Soit $E=\R^d$ muni d'une norme $\|\cdot\|$. On définit la
\emph{distance} d'un élément $x_0$ de $E$ à une partie $A$ de $E$,
notée $d(x_0,A)$, par la formule $$d(x_0,A)=\inf_{x\in A}\|x-x_0\|.$$
\begin{enumerate}
\item Supposons $A$ compact. Montrer que pour tout $x_0\in E$
il existe $y\in A$ tel que $d(x_0,A)=\|y-x_0\|$.
\item Montrer que le résultat est encore vrai si on suppose seulement
que $A$ est fermé. (On remarquera que pour toute partie $B$ de $A$ on
a $d(x_0,B)\ge d(x_0,A)$.)
\item Montrer que l'application qui à $x_0$ associe $d(x_0,A)$ est
continue sur $E$ (sans aucune hypothèse sur $A$).
\item En déduire que si $A$ est un fermé de $E$ et $B$ un compact de
$E$ tels que $A$ et $B$ sont disjoints, alors il existe une
constante $\delta>0$ telle que
$$\|a-b\| \ge \delta \qquad \forall (a,b)\in A\times B.$$
\item Montrer par un contre-exemple que le résultat est faux si on
suppose seulement que $A$ et $B$ sont deux fermés disjoints.
\end{enumerate}
\finenonce{001879}



\finexercice





\exercice{1881, maillot, 2001/09/01}

\enonce{001881}{}
\begin{enumerate}
\item
 Montrer que $\forall p\geq 1$, l'application $\left(\begin{array}{cccl}
N_p:&\R^n&\longrightarrow&\R\\
&x&\longmapsto &(\sum_{k=1}^n |x_i|^p)^\frac{1}{p}
\end{array}\right)$ est une norme (on utilisera la convexité de $x^p$).
\item Pour $x\in\R^n$ fixé, montrer que $\lim_{p\rightarrow +\infty}
N_p(x)=\max(x_i,1\leq i\leq n)$, et que cela définit une norme, appelée
{\bf norme infinie}, et notée $N_\infty$.

\item \'Etablir les inégalités suivantes :
$$
\forall x \in \R^n,\; N_\infty(x)\leq N_1(x)\leq\sqrt{n}N_2(x)\leq
nN_\infty(x).
$$
Que peut-on en déduire ?

\item Dessiner les boules unités des normes 1,2, et $\infty$ dans $\R^2$.
\end{enumerate}
\finenonce{001881}



\finexercice



\exercice{1882, maillot, 2001/09/01}

\enonce{001882}{}
Soit $\left(\begin{array}{cccl}
N:&\R^n&\longrightarrow&\R\\
&x&\longmapsto &\sum_{k=1}^n |\sum_{i=1}^k x_i|
\end{array}\right)$. Montrer que $N$ est une norme.
\finenonce{001882}



\finexercice



\exercice{1883, maillot, 2001/09/01}

\enonce{001883}{}
$A$ est dit {\it convexe} s'il contient tout segment reliant deux
quelconques de ses points :
$$
\forall (x,y)\in A^2,\;[x,y]=\{x+t(y-x),\,t\in [0,1]\}\subset A.
$$

Soit $E$ un espace vectoriel muni d'une norme $N$. Montrer que toute
boule fermée (ou ouverte) est convexe et symétrique par rapport à son
centre.
\finenonce{001883}



\finexercice



\exercice{1884, maillot, 2001/09/01}

\enonce{001884}{}
Soit $(E,N)$ un espace vectoriel normé. Montrer :
$$
\forall
(x,y)\in(E\setminus\{0\})^2,\;N(x-y)\geq\frac{1}{2}\sup(N(x),N(y))\cdot
N\left(\frac{x}{N(x)}-\frac{y}{N(y)}\right).
$$

\finenonce{001884}



\finexercice



\exercice{1885, maillot, 2001/09/01}

\enonce{001885}{}
Soit $E$ un espace vectoriel normé, et $(a,a')\in E^2$,
$(r,r')\in(\R_+^*)^2$. Montrer :
\begin{enumerate}
\item
$B(a,r)=\{a\}+B(0,r)$
\item
$B(a,r)=B(a',r') \Leftrightarrow a=a' \text{ et } r=r'$
\item
$B(a+a',r+r')=B(a,r)+B(a',r')$
\item
$B(a,r)\cap B(a',r')\neq\emptyset\Leftrightarrow\|{a'-a}\|{}<r+r'$.
\end{enumerate}
\finenonce{001885}



\finexercice



\exercice{1886, maillot, 2001/09/01}

\enonce{001886}{}
Soit $(E,N)$ une espace vectoriel. Montrer les équivalences :
$$
\begin{array}{ccl}
A\subset E \text{ est borné} &\Leftrightarrow&\exists(a,r)\in E\times\R^+ \, : \,
A\subset B(a,r)\\
&\Leftrightarrow&\exists R\geq 0\, : \, A\subset B(0,R)\\
&\Leftrightarrow&\exists R\geq 0\, : \, A\subset B_f(0,R)\\
&\Leftrightarrow&A\text{ est inclus dans une boule de }E.
\end{array}
$$
\finenonce{001886}



\finexercice



\exercice{1887, maillot, 2001/09/01}

\enonce{001887}{Topologie du $\R$-espace vectoriel $\R$}
\begin{enumerate}
\item Quelles sont toutes les normes sur le $\R$-espace vectoriel $\R${} ?

On se place désormais dans $(\R,|\, .\, |)$.

\item Quelles sont les boules ouvertes ? fermées ?

\item Ouverts et fermés de $\R${} :
\begin{enumerate}
\item
soit $(I_a)_{a\in A}$ une famille d'intervalles ouverts non vides de
$\R$, deux à deux disjoints. Montrer que A est au plus dénombrable.
\item
soit $O$ un ouvert de $\R$, et $a\in O$. On pose $A=\{x\in\R\ | \  x\not\in
O\text{ et } x>a\}$ et $B=\{x\in\R\ | \  x\not\in O \text{ et }
x<a\}$. Etudier l'existence de $\inf A$ et $\sup B$.
\item
en déduire que :
\begin{itemize}
\item
tout ouvert de $\R${} est réunion d'une famille au plus dénombrable
d'intervalles ouverts
\item
tout fermé de $\R${} est réunion d'une famille au plus dénombrable
d'intervalles fermés.
\end{itemize}
\end{enumerate}
\end{enumerate}
\finenonce{001887}



\finexercice



\exercice{1888, legall, 2003/10/01}

\enonce{001888}{}
Soit $E$ l'espace vectoriel des fonctions de classe $C^1$ sur $[0,1]$ 
telles que $f(0)=0.$
\begin{enumerate}
\item On pose pour tout $f\in E$, $N(f)= \Vert f\Vert _{\infty}$ et 
$N'(f)=\Vert f'\Vert _{\infty}$.
Montrer que $N$ et $N'$ sont des normes.
\item Montrer que $N$ et $N'$ ne sont pas \'equivalentes.
\end{enumerate}
\finenonce{001888}



\finexercice



\exercice{1889, legall, 2003/10/01}

\enonce{001889}{}
Soit $E$ l'espace vectoriel des fonctions de classe $C^1$ sur $[0,1]$ 
telles que $f(0)=0.$
\begin{enumerate}
\item On pose pour tout $f\in E$, $N(f)= \Vert f\Vert _{\infty}+\Vert 
f'\Vert _{\infty}$. Montrer que $N$ est une norme sur $E$
\item Montrer que, si $f\in E$ alors, pour tout $x\in [0,1]$~:
$  \displaystyle{ f(x) = e^{-x}\int _0 ^x e^t(f(t)+f'(t))dt}$.
\item On pose, pour tout $f\in E$, $N'(f)= \Vert f+f'\Vert 
_{\infty}$. Montrer que $N'$ est une norme sur $E$,
\'equivalente \`a $N$.
\end{enumerate}
\finenonce{001889}



\finexercice



\exercice{1890, legall, 2003/10/01}

\enonce{001890}{}
 Soit $E$ un espace vectoriel norm\'e, $A$ une partie 
de $E$ et $x$ un \'el\'ement de $E$. Comparer les deux assertions~:
\vskip1mm {\em i)} Pour tout $\epsilon > 0 $ l'ensemble $A\cap 
B(x, \epsilon )$ est infini.
\vskip1mm {\em ii)} Pour tout $\epsilon > 0 $ il existe un 
\'el\'ement $y$ distinct de $x$ dans $A\cap B(x, \epsilon )$.


\finenonce{001890}



\finexercice



\exercice{1891, legall, 2003/10/01}

\enonce{001891}{}
 Soit $A$ l'ensemble des fonctions continues sur 
$[0,1]$ telles que $f(x)\geq 0$ pour tout $x\in [0,1].$
\begin{enumerate}
\item On munit $C[0,1]$ de la norme $\displaystyle{\Vert f \Vert 
_{\infty}=\sup _{x\in [0,1]}\vert f(x)\vert }.$ Montrer que $A$ est 
ferm\'e et calculer son int\'erieur.
\item On munit $C[0,1]$ de la norme $\displaystyle{\Vert f \Vert 
_{1}=\int _{0}^1\vert f(x)\vert dx}.$ Montrer que l'int\'erieur de 
$A$ est vide et que $A$ est ferm\'e.
\end{enumerate}
\finenonce{001891}



\finexercice



\exercice{1892, legall, 2003/10/01}

\enonce{001892}{}
 Soit $ 
(E,\Vert \quad \Vert ) $ un espace vectoriel
norm\'e sur $\Rr$. On pose $$\mu (E) = \sup _{x,y \in 
E-(0,0)}\frac{\Vert x+y\Vert ^2+\Vert x-y\Vert ^2}{2(\Vert x\Vert 
^2+\Vert y\Vert ^2)}.$$

\begin{enumerate}
\item Montrer que $1\leq \mu (E)\leq 2$.
\item Calculer $\mu (\Rr ^2)$ lorsque $\R^2$ est muni de la norme 
euclidienne puis de la norme $\Vert (x,y)\Vert _{\infty}=\max \{\vert 
x\vert, \vert y\vert \}$.
\end{enumerate}
\finenonce{001892}



\finexercice



\exercice{1893, legall, 2003/10/01}

\enonce{001893}{}
Soit $\Vert \quad \Vert $ une norme sur $\Rr ^n$ et 
$A=(a_{i,j})_{i,j\in {1,\cdots n}}\in M_n(\Rr ).$ On pose~:
$$\Vert A\Vert=\sup_{x\in \Rr^n ;\Vert x\Vert =1}\Vert Ax\Vert .$$
\begin{enumerate}
\item Montrer qu'on d\'efinit ainsi une norme sur $M_n(\Rr ).$
\item On munit $\Rr ^n$ de la norme $\Vert \quad \Vert _1.$ Montrer 
que $\displaystyle{\Vert A\Vert _1= \max _{ 1\leq j \leq n}(\sum 
_{j=1}^n \vert a_{i,j}\vert ).}$
\end{enumerate}


\finenonce{001893}



\finexercice



\exercice{1895, legall, 2003/10/01}

\enonce{001895}{}
On munit 
$C[0,1]$, l'espace vectoriel des fonctions continues sur $[0,1]$ \`a 
valeurs r\'eelles
de la norme $\Vert f\Vert _{\infty }= \displaystyle{ \sup _{x\in 
[0,1]}\vert f(x)\vert }.$
\begin{enumerate}
\item Soit $\varphi : C[0,1]\rightarrow \Rr $
une application lin\'eaire. On pose $N(\varphi ) =\displaystyle{ \sup 
_{f\in C[0,1]; \Vert f \Vert _{\infty}=1}\vert \varphi (f)\vert } .$
Montrer que $\varphi $ est continue si et seulement si $N(\varphi ) $ est fini.
\item Calculer $N(\psi )$ lorsque  $\psi (f) = \displaystyle{\int _0 
^1f(t)dt }.$
\item Posons, pour toute fonction $f\in C[0,1]$ : $\varphi (f) = 
\displaystyle{\int _0 ^\frac{1}{2}f(t)dt -\int _\frac{1}{2}^1f(t)dt}.$
Montrer que $N(\varphi )=1$.
\end{enumerate}
\finenonce{001895}



\finexercice



\exercice{1896, legall, 2003/10/01}

\enonce{001896}{}
On munit $E$, 
l'espace vectoriel des fonctions continues sur $[0,1]$ \`a valeurs 
r\'eelles
telles que $f(0)=0$ de la norme $\Vert f\Vert _{\infty }= 
\displaystyle{ \sup _{x\in [0,1]}\vert f(x)\vert }.$
\begin{enumerate}
\item Soit $\varphi : E\rightarrow \Rr $
une application lin\'eaire. On pose $N(\varphi ) =\displaystyle{ \sup 
_{f\in E; \Vert f \Vert _{\infty}=1}\vert \varphi (f)\vert } .$
Montrer que $\varphi $ est continue si et seulement si $N(\varphi ) $ 
est fini. Montrer que $\varphi \mapsto N(\varphi )$ est une norme
sur l'espace vectoriel des formes lin\'eaires continues sur $E$.
\item Calculer $\mu =N(\psi )$ lorsque $\psi $ est d\'efinie en 
posant, pour toute fonction $f\in E$ : $\psi (f) = \displaystyle{\int 
_0 ^1f(t)dt }.$

\item Peut-on trouver une fonction $f\in E$ telle que  $\vert \psi 
(f)\vert =\mu $ et $ \Vert f \Vert _{\infty}=1$~?
\end{enumerate}
\finenonce{001896}



\finexercice



\exercice{1897, legall, 2003/10/01}

\enonce{001897}{}
On munit $E =C^1[0,1]$
   et $F=C [0,1]$ de la norme $\Vert f\Vert _{\infty }= \displaystyle{ 
\sup _{x\in [0,1]}\vert f(x)\vert }$.
\begin{enumerate}
\item Soit $\varphi : E\rightarrow F $
une application lin\'eaire. On pose $N(\varphi ) =\displaystyle{ \sup 
_{f\in E; \Vert f \Vert _{\infty}=1}\vert \varphi (f)\vert } .$
Montrer que $\varphi $ est continue si et seulement si $N(\varphi ) $ est fini.
\item Montrer que l'application $f\mapsto f' $ n'est pas continue.
\end{enumerate}
\finenonce{001897}



\finexercice



\exercice{1898, legall, 2003/10/01}

\enonce{001898}{}
 Soit $(E, 
\langle , \rangle )$ un espace euclidien et $S=\{ x \in E ; \Vert x 
\Vert =1  \} .$
\begin{enumerate}
\item Soient $x, y\in E $ et $I$ le segment $[x,y]$. Calculer $S\cap I .$

\item Les normes $\Vert \quad \Vert _{1} $ et $\Vert \quad \Vert 
_{\infty } $ de $\Rr ^n$ sont-elles euclidiennes ?
\end{enumerate}
\finenonce{001898}



\finexercice



\exercice{1899, legall, 2003/10/01}

\enonce{001899}{}
\begin{enumerate}
\item Soit $A\in M_n(\Cc )$. Montrer quil existe une suite de 
matrices $(A_n)_{n\in \Nn }$
inversibles convergeant vers $A$ (en un sens que l'on pr\'ecisera).
\item Soit $N\in M_n(\Cc )$ une matrice nilpotente. Calculer les 
valeurs propres de $N$. Montrer
que $\hbox{det}(I+N)=1.$
\item Soit $A\in M_n(\Cc )$ telle que $AN=NA$. Calculer $\hbox{det}(A+N).$

\end{enumerate}

\finenonce{001899}



\finexercice



\exercice{1900, legall, 2003/10/01}

\enonce{001900}{}
\ \par
\begin{center}{\large \bf I Pr\'eliminaires}
\end{center}


\begin{enumerate}
\item Soit $\mathcal{P} $ l'espace vectoriel des fonctions polynomiales de 
$[0,1]$ \`a valeurs dans $\Rr .$ Montrer que $\mathcal{P} $ est de dimension 
infinie.
\item Soit $X$ une partie born\'ee de $\Rr .$ Montrer que $\sup 
(X)=\sup \bar{X}.$
\end{enumerate}


\vskip5mm
\begin{center}{\large \bf II}
\end{center}
\vskip5mm

On note $\mathcal{L}$ l'ensemble des fonctions {\em lipschitziennes} de 
$[0,1]$ \`a valeurs dans $\Rr$, c'est \`a dire telles qu'il existe
$k\in \R _+$ tel que, pour tout $x,y\in [0,1], $ $ \vert 
f(x)-f(y)\vert \leq k\vert x-y\vert $. On note $\C ^1$ l'ensemble des 
fonctions de $[0,1]$ \`a valeurs dans $\Rr$
de classe $C ^1,$ c'est \`a dire d\'erivables \`a d\'eriv\'ee continue.
\begin{enumerate}
\item Montrer que $\mathcal{L} $ est un sous espace vectoriel de l'espace 
vectoriel des fonctions de $[0,1]$ \`a valeurs dans $\Rr$, que $\mathcal{L} $ 
contient $\C ^1$ et est de dimension infinie.





\item On pose, pour tout $f\in \mathcal{L} :$
$$\displaystyle{ N_1(f)=\vert f(0)\vert +\sup _{(x,y)\in [0,1]^2, 
x\not = y}\dfrac {\vert f(x)-f(y)\vert }{\vert x-y\vert }}$$
$$\displaystyle{ N_2(f)=\vert f(0)\vert +\sup _{x\in ]0,1]}\dfrac 
{\vert f(x)-f(0)\vert }{\vert x\vert }}$$
$$ \Vert f \Vert _{\infty} =\sup _{x\in [0,1]}\vert f(x)\vert $$
$$\displaystyle{ \lambda (f)=\Vert f\Vert _{\infty} +\sup _{(x,y)\in 
[0,1]^2,  x\not = y}\dfrac {\vert f(x)-f(y)\vert }{\vert x-y\vert 
}}.$$

\begin{enumerate}
\item Montrer que $N_1$, $N_2$,  $
\Vert \ \Vert _{\infty}$ et $\lambda $ sont des normes sur $\mathcal{L} .$
\item En consid\'erant la suite $f_n(x)=\sin (2\pi n x),$ montrer que 
$N_2$ n'est pas \'equivalente \`a $
\Vert \ \Vert _{\infty}.$
\item Montrer que $N_1$ n'est \'equivalente ni \`a  $
\Vert \ \Vert _{\infty}$, ni \`a $N_2.$
\item Construire une suite $(g_n)_{n\in \Nn }$ d'\'el\'ements de $\mathcal{L} 
$ qui converge vers $0$ pour $
\Vert \ \Vert _{\infty}$ mais pas pour $N_2$. En d\'eduire (de 
nouveau) que $N_2$ n'est pas \'equivalente \`a $
\Vert \ \Vert _{\infty}.$
\item Montrer que $\lambda $ et $N_1$ sont \'equivalentes.
\end{enumerate}

\item On pose, pour tout $f\in \C ^1 :$
$\displaystyle{ \nu _1(f)=\vert f(0)\vert +\Vert f'\Vert _{\infty}}$
et $\displaystyle{ \nu (f)=\Vert f\Vert _{\infty} +\Vert f'\Vert _{\infty}}.$
\begin{enumerate}
\item Montrer que $\nu _1$ et $\nu $ sont des normes sur $\C ^1$.
\item Montrer que $\nu _1(f)= N_1(f) ,$ pour tout $f\in \C ^1.$
\item Les normes $\nu $ et $\nu _1$ sont-elles \'equivalentes~?
\end{enumerate}


\item  Soit $(E, \Vert \ \Vert )$ un espace vectoriel norm\'e.
Une suite $(x_n)_{n\in \Nn }$ d'\'el\'ements de $E$ est dite de {\em 
Cauchy} si, pour tout $\epsilon >0, $ il existe $N $ tel que, si 
$m, n\geq N$
alors $\Vert x_n-x_m\Vert \leq \epsilon .$ On dit que $(E, \Vert \ 
\Vert )$ est {\em complet} si toute suite de Cauchy y est 
convergente. On rappelle
que $\Rr $ muni de la norme $x\mapsto \vert x \vert $ est complet.
\begin{enumerate}
\item Soit $C^0$ l'espace vectoriel des fonctions continues de 
$[0,1]$ \`a valeurs dans $\Rr$. Montrer que $(C^0, \Vert \ \Vert 
_{\infty})$ est complet.
\item L'espace vectoriel norm\'e $(\C^1, \nu )$ est-il complet~? 
Qu'en est-il de $(\C^1, \nu _1)$~?
\item Soit $(f_n)_{n\in \Nn }$ une suite de Cauchy dans $(\mathcal{L} , 
\lambda ).$ Montrer que $(f_n)_{n\in \Nn }$ converge uniform\'ement 
vers une fonction continue $f.$
\item D\'emontrer que pour $n $ assez grand $f-f_n$ est lipschitzienne.
\item En d\'eduire que $(\mathcal{L} , \lambda )$ est complet.
\end{enumerate}
\end{enumerate}

\vskip5mm
\begin{center}{\large \bf III}
\end{center}
\vskip5mm


On munit $\C ^1$ d'une norme $N$ et
   et $C^0$ de la norme $\Vert \ \Vert _{\infty }$. On note $d$ 
l'application $f\mapsto f'$ de $\C ^1$ \`a valeurs dans $\C ^0.$
\begin{enumerate}
\item Soit $\varphi : \C ^1\rightarrow C^0 $
une application lin\'eaire. On pose $N(\varphi ) =\displaystyle{ \sup 
_{f ; N( f )\leq 1}\Vert \varphi (f)\Vert _{\infty}} .$
D\'emontrer que $\varphi $ est continue si et seulement si $N(\varphi 
) $ est fini. V\'erifier que $N$ est une norme sur l'espace vectoriel 
des applications lin\'eaires
continues de $\C ^1$ \`a valeurs dans $\C^0.$
\item Montrer que l'application $d $ n'est pas continue si $N= \Vert 
\ \Vert _{\infty}$.
\item On munit $\C ^1$ de la norme $\nu .$ Montrer que $d$ est 
continue et calculer $N(d).$
\end{enumerate}




\finenonce{001900}


\finexercice



\exercice{6055, queffelec, 2011/10/16}
\enonce{006055}{}
\begin{enumerate}
\item Soit $||.||$ une norme sur $\Rr^n$ et $K$ sa boule unité fermée. Montrer
que   

(i) $K$ est symétrique,

(ii) $K$ est convexe, fermé, borné,

(iii) $0$ est un point intérieur à $K$.

\item Réciproquement, montrer que si $K$ possède les trois propriétés ci-dessus,
il existe une norme dont $K$ soit la boule unité fermée, en considérant

\qquad $p(x)=\inf\{a>0\ ;\ {x\over a}\in K\}$.
\end{enumerate}



\finenonce{006055}


\finexercice


\exercice{6056, queffelec, 2011/10/16}
\enonce{006056}{}
Montrer que dans un espace normé, la boule unité est convexe.

Réciproquement, supposons que l'espace vectoriel soit muni d'une
applica\-tion $N$ de $E$ dans $\Rr^+$ telle que $N(\lambda x)=\vert\lambda\vert
N(x)$, et telle que $\{y/N(y)\leq1\}$ soit convexe. Montrer que 
$$N(x+y)\leq 2\sup(N(x), N(y)), \ \ x,y\in E.$$
\finenonce{006056}



\finexercice       


\exercice{6057, queffelec, 2011/10/16}
\enonce{006057}{}
 On considère dans ${\Rr}^2$, les deux applications
$$n((x,y))=\sup_{t\in [0,1]}\vert x+ty\vert,$$
$$m((x,y))=\int_0^1\vert x+ty\vert\ dt.$$

\begin{enumerate}  
\item Montrer que $n$ et $m$ définissent deux normes sur ${\Rr}^2$.

\item Dessiner les boules unités fermées associées, et trouver des constantes
effec\-tives $A$,$B$, telles que $A\ n((x,y))\leq m((x,y))\leq B\ n((x,y))$ pour
tout $(x,y)\in {\Rr}^2$.
\end{enumerate}
\finenonce{006057}



\finexercice       


\exercice{6058, queffelec, 2011/10/16}
\enonce{006058}{}
\begin{enumerate}  
\item On considère dans ${\Rr}^2$ les 4 boules euclidiennes fermées de rayon $1$
centrées  aux points $(1,0), (-1,0), (0,1), (0,-1)$; $A$ leur réunion contient
$0$ comme point intérieur. Trouver le rayon de la plus grande boule ouverte
centrée en $0$ et contenue dans $A$.

\item On se pose plus généralement le problème dans ${\Rr}^n$ : $A$ désigne
l'union $\cup_j\overline B(e_j,1)\cup_j\overline B(-e_j,1)$ où $(e_j)$ est la
base canonique de ${\Rr}^n$. Montrer que $x\in A$ si et seulement si $\Vert
x\Vert_2^2\leq 2\Vert x\Vert_\infty$. En déduire que le rayon de la plus grande
boule ouverte centrée en $0$ et contenue dans $A$ est ${2\over\sqrt n}$. 
\end{enumerate}
\finenonce{006058}



\finexercice       


\exercice{6059, queffelec, 2011/10/16}
\enonce{006059}{}
Soit $N$ un entier $\geq1$, et $E$, l'espace des polyn\^omes trigonométriques
$p$ de degré $\leq N$, $p(t)=\sum_{-N}^N c_k\exp(ikt).$

On pose, pour $p\in E$, $\Vert p\Vert_{\infty}=\sup_{t\in[0,2\pi]} \vert
p(t)\vert$, et $\Vert p\Vert=\sum_{-N}^N \vert c_k\vert$. Montrer, à l'aide de
l'identité de Parseval, que ces deux normes vérifient
$$\Vert p\Vert_{\infty}\leq \Vert p\Vert\leq\sqrt{2N+1}\Vert p\Vert_{\infty}.$$
\finenonce{006059}



\finexercice       


\exercice{6088, queffelec, 2011/10/16}
\enonce{006088}{}
Soit $E$ un espace vectoriel normé sur $\R$ ou $\C$.
\begin{enumerate}
\item Vérifier que l'application $(\lambda,x)\to \lambda x$ est continue; que
$(x,y)\to x+y$ est lipschitzienne ainsi que l'applica\-tion $x\to\Vert
x\Vert$; et que les translations et les homothéties sont des
homéomor\-phismes de
$E$.

\item Montrer que la boule unité ouverte est
homéomorphe à
$E$ tout entier (considérer l'application $x\to {x\over 1-||x||}$).

\item Montrer que deux
boules ouvertes de $(E,||.||)$ sont homéomorphes entre elles.

\item Montrer que le seul sous-espace ouvert de $E$ est $E$ lui-même, et que tout
sous-espace propre est d'intérieur vide dans $E$.

\item Montrer que l'adhérence d'un sous-espace vectoriel est encore un sous-espace
vectoriel; en déduire qu'un hyperplan de $E$ est fermé ou partout dense dans $E$.
\end{enumerate}
\finenonce{006088}



\finexercice


\exercice{6089, queffelec, 2011/10/16}
\enonce{006089}{extrait du partiel de décembre 98}

Soit $E$ un espace vectoriel normé sur $\C$ de boule unité fermée
$\overline B$ et
$F$ un sous-espace vectoriel fermé de $E$. On va montrer que si $F\not=E$,

$$\sup_{x\in\overline B} d(x,F)=1.$$
\begin{enumerate}
\item Etablir les propriétés  pour $x,x'\in E, y\in F, \lambda\in\C$ :   

\quad (i) $d(x,F)\leq ||x||$.

\quad (ii) $d(\lambda x,F)=|\lambda| d(x,F)$.

\quad (iii) $d(x-y,F)=d(x,F)$

\quad (iv) $d(x+x',F)\leq d(x,F)+d(x',F)$.

\item Soit $x\in\overline B$ tel que $\alpha=d(x,F)>0$. Montrer que
pour tout $\varepsilon>0$ il existe $y\in F$ tel que : 
$$\alpha\leq ||x-y||<\alpha(1+\varepsilon).$$

\item Montrer qu'il existe $x'\in\overline B$ tel que : \qquad
${1\over{1+\varepsilon}}=d(x',F)<1$.

\item En déduire le résultat.
\end{enumerate}
\finenonce{006089}



\finexercice 

\section{425.00 Espace métrique complet, espace de Banach}


\exercice{2395, mayer, 2003/10/01}
\enonce{002395}{} 
L'espace $(\Rr , d)$ est-il complet si $d$ est l'une des
m\'etriques suivantes? \begin{enumerate}
\item $d(x,y) = |x^3-y^3|$.
\item $d(x,y) = |\exp (x) -\exp (y)|$.
\item $d(x,y) = \log (1+ |x-y|)$.
\end{enumerate}
\finenonce{002395}


\finexercice


\exercice{2396, mayer, 2003/10/01}
\enonce{002396}{} 
On consid\`ere pour $x,y \in \Rr$, $d(x,y) =\|f(x) -f(y)\|$,
o\`u $f$ est une application injective de $\Rr$ dans
$\Rr ^2$. Montrer que cette distance est compl\`ete si et
seulement si $f$ est d'image ferm\'ee dans $\Rr^2$.
\finenonce{002396}


\finexercice


\exercice{2397, mayer, 2003/10/01}
\enonce{002397}{} 
\label{exocomp}
On consid\`ere l'espace des fonctions continues $X={\cal
C}([a,b])$. \begin{enumerate}
\item Soit $\omega \in X$ une fonction qui ne s'annule pas sur
$[a,b]$. Posons $$d_\omega (f,g) = \sup_{t\in[a,b]} |\omega (t)
(f(t)-g(t))| \; .$$
 L'espace $(X,d_\omega)$ est-il complet?
 \item Montrer que l'espace $(X, \|.\|_1)$ n'est pas complet (o\`u
 $\|f\|_1 = \int _0^1 |f(t)| \, dt$).
 \end{enumerate}
\finenonce{002397}


\finexercice


\exercice{2398, mayer, 2003/10/01}
\enonce{002398}{} 
Soit $X={\cal C}^1([a,b])$.
\begin{enumerate}
\item Est-ce un espace complet si on le muni de la norme
uniforme $\| . \| _\infty $?
\item Consid\'erons maintenant, pour $f\in X$, la norme
$$ N(f) = \sup _{t\in [a,b]} \|f(t)\| + \sup _{t\in [a,b]}
\|f'(t)\|\; .$$ L'espace $(X,N)$ est-il complet?
\end{enumerate}
\finenonce{002398}


\finexercice


\exercice{2399, mayer, 2003/10/01}
\enonce{002399}{} 
Soit $X$ l'espace des suites r\'eelles nulles \`a partir d'un certain rang, et soit
$$\rho (x,y) =\sum _{k=1}^\infty 2^{-k}
\frac{|x_k-y_k|}{1+|x_k-y_k|}\quad  \text{pour} \; x,y\in X\; .$$
\begin{enumerate}
\item Montrer que $X$ n'est pas complet pour la m\'etrique $\rho $.
\item Trouver un espace de suites $Y$ tel que $(Y, \rho )$ soit complet et
tel que $X$ soit dense dans $Y$.
\item Que donne l'exercice si on remplace $\rho$ par la norme uniforme?
\end{enumerate}
\finenonce{002399}


\finexercice


\exercice{2400, mayer, 2003/10/01}
\enonce{002400}{} 
Soit $E$ un espace vectoriel norm\'e. On dit qu'une s\'erie $\sum u_k$ est normalement convergente
si la s\'erie $\sum \| u_k \|$ est convergente. On veut d\'emontrer que $E$ est complet si et seulement
si toute s\'erie normalement convergente est convergente.
\begin{enumerate}
\item Soit $(x_n)$ une suite de Cauchy de $E$; montrer qu'on peut en extraire une sous-suite
$(x_{n_k})$ telle que la s\'erie de terme g\'en\'eral $u_k = x_{n_{k+1}} - x_{n_k}$ 
soit normalement convergente.
En d\'eduire que si toute s\'erie normalement convergente est convergente, alors $E$ est complet.

\item Soit $\sum u_k$ une s\'erie normalement convergente. On note $S_n = \sum _{k=0}^nu_k$. 
Montrer que $S_n$ est une suite de Cauchy.
En d\'eduire que si $E$ est complet, alors toute série normalement convergente est convergente.
\end{enumerate}
\finenonce{002400}


\finexercice


\exercice{2401, mayer, 2003/10/01}
\enonce{002401}{} 
Soient $E,F$ des espaces norm\'es et $A_n,A \in {\cal L}(E,F)$.
Montrer l'\'equivalence entre:
\begin{enumerate}
\item $A_n \to A$ dans ${\cal L}(E,F)$.
\item Pour toute partie born\'ee $M\subset E$, la suite $A_n x$
converge uniform\'ement vers $Ax$, $x\in M$.
\end{enumerate}
\finenonce{002401}


\finexercice


\exercice{2402, mayer, 2003/10/01}
\enonce{002402}{Cours} 
Soit $E$ un espace norm\'e et $F$ un espace de Banach. Alors
${\cal L}(E,F)$ est aussi un espace de Banach.
\finenonce{002402}


\finexercice


\exercice{2403, mayer, 2003/10/01}
\enonce{002403}{}
Soit $\delta$ la m\'etrique sur $\Rr$ d\'efinie par
$\delta(x,y)=\vert{x\over{1+\vert x\vert}}-{y\over{1+\vert y\vert}}\vert$.
Montrer, \`a l'aide du th\'eor\`eme de prolongement de fonction uniform\'ement
continue, que l'identit\'e $i:({\Rr},\delta)\to ({\Rr},\vert.\vert)$ n'est
pas uniform\'ement continue.
\finenonce{002403}



\finexercice  


\exercice{6211, queffelec, 2011/10/16}
\enonce{006211}{}
\begin{enumerate}
\item Pour $n,m\in \Nn^*$, on pose $d(n,m)=|{1\over n}-{1\over m}|$. 
Montrer que $d$ est une distance sur $\Nn^*$ qui induit la topologie discrète
sur $\Nn^*$; est-elle complète ? 

\item Montrer que $\overline{\Rr}$ est un espace métrique complet pour la
distance $d(x,y)=\vert \arctan x-\arctan y\vert$.
\end{enumerate}
\finenonce{006211}



\finexercice


\exercice{6212, queffelec, 2011/10/16}
\enonce{006212}{}
Soit $E$ un espace normé. Montrer qu'il est complet si et seulement si la
sphère unité $S=\{x/\Vert x\Vert=1\}$ est complète.
\finenonce{006212}



\finexercice


\exercice{6213, queffelec, 2011/10/16}
\enonce{006213}{}
\begin{enumerate}
\item Pour $x,y\in\Rr^*$ on pose  $d(x,y)=|x-y|+|{1\over x}-{1\over y}|$.

 Montrer que $d$ définit une distance sur $\Rr^*$
qui induit la topologie usuelle et que $({\Rr^*},d)$ est complet. 

\item Plus généralement soit $U$ un ouvert d'un espace complet $(X,d)$; comment
peut-on définir une métrique $\delta$ sur $U$, équivalente à la métrique
initiale, qui fasse de $U$ un espace complet ?
\end{enumerate}
\finenonce{006213}



\finexercice


\exercice{6214, queffelec, 2011/10/16}
\enonce{006214}{}
Soit $E$ un espace vectoriel normé.
\begin{enumerate}
\item Soit $(x_n)$ une suite de Cauchy de $E$; montrer qu'on peut en extraire une
sous-suite $(x_{n_k})$ telle que la série de terme général
$u_k=x_{n_{k+1}}-x_{n_k}$ soit normalement convergente.

\item En déduire que $E$ est complet si et seulement si toute série normalement
convergente est convergente.
\end{enumerate}
\finenonce{006214}



\finexercice


\exercice{6215, queffelec, 2011/10/16}
\enonce{006215}{}
\begin{enumerate}
\item Montrer que l'espace $C([0,1])$ est complet pour la norme uniforme mais pas
pour la norme $||.||_1$.

\item Montrer que l'espace $S$ des suites réelles nulles à partir d'un certain rang,
muni de la norme uniforme, n'est pas complet. Trouver un espace métrique complet
contenant $S$ comme sous-espace dense.
\end{enumerate}
\finenonce{006215}



\finexercice


\exercice{6216, queffelec, 2011/10/16}
\enonce{006216}{}
\begin{enumerate}
\item Soit $f$ un homéomorphisme d'espaces métriques $f:X\to Y$; montrer que $X$
peut être complet sans que $Y$ le soit.

\item On suppose de plus que $f$ est uniformément
continue. Montrer que si $Y$ est complet, $X$ l'est aussi.

\item On considère $E=\{f\in C^1([0,1])\ ;\ f(0)=0\}$, muni de la métrique
$d(f,g)=\inf(1,\sup|f'(t)-g'(t)|)$. Montrer que $E$ est complet pour cette
métrique.
\end{enumerate}
\finenonce{006216}



\finexercice


\exercice{6217, queffelec, 2011/10/16}
\enonce{006217}{}
Soit $E$ un Banach, $A$,$B$ deux sous-espaces de $E$ tels que $A\cap B=\{0\}$,
$A$ étant fermé et $B$ de dimension finie.
\begin{enumerate}
\item Pour $b\in B$, on définit $[b]=d(b,A)=\inf_{a\in A}||a+b||$. Vérifier que
$[.]$ est une norme sur $B$.

\item En déduire qu'il existe $C>0$ telle que $||a+b||\geq C||b||$ pour tous $a\in
A$ et $b\in B$.

\item Montrer que $A\oplus B$ est encore un sous-espace fermé de $E$.
\end{enumerate}
\finenonce{006217}



\finexercice


\exercice{6218, queffelec, 2011/10/16}
\enonce{006218}{}
Soit $(X,d)$ un espace métrique, et $(x_n)$ une suite de Cauchy dans $X$.
Vérifier :

\begin{enumerate}  
\item La suite $(x_n)$ est bornée même si la métrique est non
bornée, mais il existe des suites bornées dont aucune sous-suite n'est de Cauchy.

\item Si $(x_n)$ contient une sous-suite convergente, elle est convergente.

\item  Soit $(\epsilon_k)$ une suite quelconque de réels $>0$; il existe une
sous-suite $(x_{n_k})$ de $(x_n)$ telle que $d(x_{n_k},x_{n_{k+1}})\leq
\epsilon_k.$

\item  Soit $(y_n)$ une suite quelconque de $X$. Si $\sum_1^\infty
d(y_n,y_{n+1})<\infty$, la suite $(y_n)$ est de Cauchy. Réciproque ?

\item  On suppose cette fois la distance $d$ ultramétrique. Dans ce cas $(y_n)$
est de Cauchy si et seulement si $d(y_n,y_{n+1})$ tend vers $0$. 
\end{enumerate}
\finenonce{006218}



\finexercice       


\exercice{6219, queffelec, 2011/10/16}
\enonce{006219}{}
Vérifier que $\overline{\Rr}$ est un espace métrique complet pour la
distance $d(x,y)=\vert \arctan x-\arctan y\vert$.
\finenonce{006219}



\finexercice       


\exercice{6220, queffelec, 2011/10/16}
\enonce{006220}{}
Sur l'ensemble $\Nn$ des entiers naturels, définissons 
\begin{eqnarray*}
d(n,m)&=& 0 \ \ {\rm pour}\  m=n\\
      &=& 1+{1\over{n+m}}\ \  {\rm pour}\  m\neq n
\end{eqnarray*}

\begin{enumerate}  
\item Montrer que $d$ est une métrique sur $\Nn$ pour laquelle il est complet.

\item Construire dans $({\Nn},d)$ une suite de boules fermées non vides embo\^\i
tées dont les rayons ne tendent pas vers $0$, et d'intersection vide.
\end{enumerate}
\finenonce{006220}



\finexercice       


\exercice{6221, queffelec, 2011/10/16}
\enonce{006221}{}
Soit $U$ un ouvert d'un espace complet $(X,d)$; on note $F=U^c$ et
$f(x,y)=\vert{1\over{d(x,F)}}-{1\over{d(y,F)}}\vert$ pour $x,y \in U$.

Montrer que $\delta(x,y)=\max(d(x,y),f(x,y))$ définit une distance sur $U$
équiva\-lente (topologiquement) à $d$ et que $(U,\delta)$ est complet. 
\finenonce{006221}



\finexercice       


\exercice{6222, queffelec, 2011/10/16}
\enonce{006222}{}
Soit $X$ un espace métrique et $(a_n)$ une suite de Cauchy dans $X$.

\begin{enumerate}  
\item Montrer que pour tout $x\in X$, la suite de réels $(d(a_n,x))$ a une limite.
On note $f(x)$ cette limite; montrer que l'application $x\to f(x)$ est continue
de $X$ dans $\Rr$.

\item Calculer $\inf_{x\in X}  f(x)$. Quand cette borne inférieure est-elle
atteinte?

\item Déduire de ce qui précède que si $X$ n'est pas complet, il existe une
applica\-tion $\phi:X\to\Rr$ continue et non bornée.
\end{enumerate}
\finenonce{006222}



\finexercice       


\exercice{6223, queffelec, 2011/10/16}
\enonce{006223}{}
On considère pour $f$ et $g$ dans $E=C(\Rr,C)$, $$d(f,g)=\sum_n{1\over
{2^n}}\min(1,\sup_{\vert x\vert\leq n}\vert f(x)-g(x)\vert).$$
Vérifier que $d$ est une métrique sur $E$ pour laquelle il est complet. 

Montrer que la convergence pour $d$ n'est autre que la convergence uniforme
sur tout compact de $\Rr$.
\finenonce{006223}



\finexercice       


\exercice{6224, queffelec, 2011/10/16}
\enonce{006224}{}
Soit $(X,d)$ un espace métrique et $Y$ une partie de $X$; on considère $f$ une
application surjective de $X$ sur $Y$, et on pose pour $u$ et $v$ dans $Y$
$$D(u,v)=d(f^{-1}(\{u\}),f^{-1}(\{v\}))=\inf_{x\in f^{-1}(\{u\}),y\in
f^{-1}(\{v\})} d(x,y).$$

\begin{enumerate}  
\item Montrer que pour $u$ et $v$ dans $Y$, $D(u,v)\geq0$; que $D(u,u)=0$ et
$D(u,v)=D(v,u)$; $D$ vérifie-t-elle l'inégalité triangulaire ?

\item On suppose que pour $u$ dans $Y$, $f^{-1}(\{u\})$ est un fermé de $Y$, et
que pour $u$ et $v$ dans $Y$
$$d(x,f^{-1}(\{v\}))=d(x',f^{-1}(\{v\}))$$
pour tous $x$ et $x'$ dans $f^{-1}(\{u\})$. Montrer alors que $D$ est une
distance.

\item  On suppose les conditions de 2. vérifiées. Montrer que $Y$ est complet si
$X$ est complet. 
\end{enumerate}
\finenonce{006224}



\finexercice


\exercice{6225, queffelec, 2011/10/16}
\enonce{006225}{}
On considère sur $C^1([0,1],\Rr)$ les normes suivantes :
\begin{enumerate}  
\item $\Vert f\Vert=\sup_{[0,1]} \vert f(x)\vert$

\item $\Vert f\Vert=\sup_{[0,1]} \vert f'(x)\vert + \vert f(0)\vert$

\item $\Vert f\Vert=\sup_{[0,1]} \vert f'(x)+f(x)\vert + \vert f(0)\vert$
\end{enumerate}
Lesquelles sont complètes sur $C^1([0,1],\Rr)$ ?
\finenonce{006225}



\finexercice       


\exercice{6226, queffelec, 2011/10/16}
\enonce{006226}{}
Soit $(B,\Vert.\Vert)$ un espace de Banach et $M$, $N$ deux sous-espaces de $B$
tels que
$B=M\oplus N$. On met sur $B$ une nouvelle norme $\Vert
z\Vert'=\Vert x\Vert+\Vert y\Vert$ si $z=x+y$.

\begin{enumerate}  
\item Vérifier que $\Vert.\Vert'$ est bien une norme sur $B$ et que
$(B,\Vert.\Vert')$ est complet si et seulement si $M$ et $N$ sont fermés.

\item Montrer que si les projections $P_M$ et $P_N$ sur
$M$ et $N$ sont continues, $(B,\Vert.\Vert')$ est encore un Banach.
\end{enumerate}
\finenonce{006226}



\finexercice       


\exercice{6227, queffelec, 2011/10/16}
\enonce{006227}{}
On considère $E=c$, l'espace des suites réelles convergentes; montrer que, muni
de la norme uniforme, $E$ est complet et décrire son dual topologique. 
\finenonce{006227}



\finexercice       


\exercice{6228, queffelec, 2011/10/16}
\enonce{006228}{}
On considère $E$ l'espace des séries convergentes, et on pose
$$\Vert\xi\Vert=\sup_n\vert\sum_{k=1}^n\xi_k\vert$$

\begin{enumerate}  
\item Vérifier que ceci définit une norme sur $E$ pour laquelle il est complet.

\item L'espace $l^1$ des séries absolument convergentes est un sous-espace de
$E$; montrer que les normes $\Vert.\Vert$ et $\Vert.\Vert_1$
 ne sont pas équivalentes sur $l^1$ (en considérant une série de terme général
$\xi_k={{(-1)^k}\over {k^\alpha}}$.)

\item Montrer que $l^1$ est dense dans $(E,\Vert.\Vert)$.
\end{enumerate}
\finenonce{006228}



\finexercice       


\exercice{6229, queffelec, 2011/10/16}
\enonce{006229}{}
Pour tout $k>0$ on note $H_k$ le sous-espace de $C([0,1])$ constitué des
fonctions lipschitziennes de constante $k$ ie des fonctions $f$ vérifiant
$\vert f(x)-f(y)\vert\leq k\vert x-y\vert$ pour tous $x$ et $y$ dans $[0,1]$.
On pose aussi $H=\bigcup_{k>0}H_k$.

\begin{enumerate}  
\item Montrer que $H$ contient les fonctions de classe $C^1$ sur $[0,1]$, mais
que la fonction $\sqrt x $ n'est pas dans $H$.

\item Montrer que pour tout $k$, $H_k$ est un espace de Banach pour la norme
uniforme.

\item Montrer qu'il existe une suite de fonction de $H$ qui converge uniformément
sur $[0,1]$ vers $\sqrt x$. En déduire que $H$ n'est pas complet pour la
norme uniforme.

\item Montrer que si on pose $$\Vert f\Vert=\sup_{x\neq y}{{\vert
f(x)-f(y)\vert}\over{\vert x-y\vert}} +\vert f(0)\vert,$$ on définit ainsi une
norme sur l'espace $E$, pour laquelle l'espace est complet.  
\end{enumerate}
\finenonce{006229}



\finexercice       


\exercice{6230, queffelec, 2011/10/16}
\enonce{006230}{}
Soit $E$ un espace de Banach, $A\in{\cal L}(E)$, et $s,t\in\Rr$.

\begin{enumerate}  
\item On rappelle que $e^{tA}=\sum_0^\infty {{t^nA^n}\over{n!}}\in{\cal L}(E)$.
Montrer que $\Vert e^{tA}\Vert\leq e^{\vert t\vert\Vert A\Vert}$ et que
$e^{(t+s)A}=e^{tA} e^{sA}$.

\item Soit $u_0\in E$ et $u$ la fonction vectorielle de variable réelle définie
par $u(t)=e^{tA}u_0$. Montrer que $u$ est dérivable sur $\Rr$ et calculer sa
dérivée.
\end{enumerate}
\finenonce{006230}



\finexercice       


\exercice{6231, queffelec, 2011/10/16}
\enonce{006231}{}
Soit $E$ un espace de Banach et $F$ un sous-espace fermé de $E$.

\begin{enumerate}  
\item Montrer que $N(\bar x)=\inf_{y\in F} \Vert x+y\Vert=d(x,F)$ définit une norme
sur l'espace vectoriel quotient $E/F$.

\item Montrer à l'aide du critère sur les séries que $E/F$ muni de $N$ est un
espace de Banach. 
\end{enumerate}
\finenonce{006231}



\finexercice  


\exercice{6781, gijs, 2011/10/16}
\enonce{006781}{}
 Soit $A\subset \Rr^n$ compact, soit $\mathcal{C}^0(A) = \{\,f:A \to \Rr \mid f \text{ continue }\,\}$,
et $  ||f||_\infty = \sup_{x\in A} |f(x)|$.
Montrer que $(\mathcal{C}^0(A), ||{\ }||_\infty)$ est un espace
de Banach.
\finenonce{006781}



\finexercice

\exercice{6784, gijs, 2011/10/16}
\enonce{006784}{}
Soit $M\subset \Rr^n$ une sous-variété
de dimension $k$, et soit $X$ un champ de vecteurs sur
$M$. On définit $D = \{\,m\in M\mid X(m) \neq 0\,\}$ et
$S = \overline D = \text{ fermeture } D$. On vous demande
de démontrer l'énoncé : ``si $S$ est compact, alors
$X$ est complet''. Les questions suivantes peuvent vous
guider.


\begin{enumerate}
\item Pour $x \notin S$ trouver la courbe intégrale
maximale $\gamma : J_x \to M$ passant par $x$.
\item Montrer que si $x\in S$, et si $\gamma : J \to M$
est une courbe intégrale passant par $x$, alors
$\forall\ t\in J\ :\ \gamma(t) \in S$.
\item En utilisant la compacité de $S$, montrer que
$X$ est complet (sur $M$ !).
\end{enumerate}
Nota Bene  : $S$ n'est pas un sous-variété d'un $\Rr^m$; montrer le résultat (du cours!) ``si $M$
est compact, alors $X$ est complet'' rapporte moins de
points.
\finenonce{006784}



\finexercice


\exercice{6789, gijs, 2011/10/16}
\enonce{006789}{}
Soit $(X,d)$ un espace métrique et soit $a:\Nn \to
X$ une suite dans $X$.

\begin{enumerate}
\item Quand est-ce que $a$ est une suite de Cauchy,
une suite convergente ?

\item Donner les définitions d'un point
d'accumulation de la suite $a$, et de ``$(X,d)$ est
complet''.

\item Enoncer le théorème de Bolzano-Weierstrass.

\item Montrer que si $a$ est de Cauchy et si $x$ est un
point d'accumulation de $a$, alors $a$ converge vers $x$.

\item Montrer que si $X$ est compact, alors $X$ est complet.
\end{enumerate}
\finenonce{006789}



\finexercice

\exercice{6791, gijs, 2011/10/16}
\enonce{006791}{}
Soit $B$ et $C$ deux  sous-ensembles compacts de $\Rr^n$, et soit $E = \{\, f : B\to C \mid f \text{
continue}\,\}$. On définit $  d(f,g) = \sup_{x\in
B} |f(x) - g(x)|$.
\begin{enumerate}
\item Montrer que $d : E \times E \to \Rr$ est
une métrique sur $E$.

\item Montrer que $(E,d)$ est un espace métrique complet.
\end{enumerate}
\finenonce{006791}



\finexercice


\exercice{6819, gijs, 2011/10/16}
\enonce{006819}{}
Soit $(E, \|\cdot\|)$ un espace vectoriel normé. Le
but de cet exercice est de montrer que $E$ est complet si
et seulement si toute série absolument convergente
converge.
\begin{enumerate}
\item
Soit $E$ complet et $(a_n)_{n\in \Nn}$ une suite dans $E$
telle que $\sum_{n=0}^\infty \|a_n\|$ converge.
Démon\-trer que la suite $s_N = \sum_{n=0}^N a_n$ est
une suite de Cauchy~; en déduire que $
\sum_{n=0}^\infty a_n$ converge.

\item
On suppose que toute série absolument convergente
converge, c'est-à-dire si $\sum_{n=0}^\infty
\|b_n\|$ converge, alors $\sum_{n=0}^\infty b_n$
converge. Soit $(a_n)_{n\in \Nn}$ une suite de Cauchy dans
$E$. Trouver une suite strictement croissante $i\mapsto n_i
\in \Nn$ telle que $\forall i : \|a_{n_{i+1}} -
a_{n_i}\| < 2^{-i}$. En déduire que
$\sum_{i=0}^\infty (a_{n_{i+1}} - a_{n_i})$
converge. Déduire de ce résultat que $E$ est complet.
\end{enumerate}
\finenonce{006819}



\finexercice

\section{426.00 Théorème du point fixe}


\exercice{2404, mayer, 2003/10/01}
\enonce{002404}{} 
\label{exn13}
Soit $\alpha _n >0$ tel que la s\'erie $\sum_{n=1}^\infty \alpha _n $
converge. Soit $(X,d)$ un espace m\'etrique complet et $f:X\to X$
une application pour laquelle
$$ d(f^n(x),f^n(y))\leq \alpha_n d(x,y) \quad \text{pour tout } x,y\in X \;
\text {et }\; n\in \Nn\; .$$ Montrer que, sous ces conditions, $f$
poss\`ede un unique point fixe $p\in X$, que pour tout point
initial $x_0 \in X$, la suite  des it\'er\'ees $(x_n=f^n
(x_0))_{n\geq 0}$ converge vers $p$ et que la vitesse de
convergence d'une telle suite est contr\^ol\'ee par
$$ d(p,x_n) \leq \left ( \sum_{\nu =n}^\infty \alpha _\nu \right ) d(x_1,x_0) \; .$$
\finenonce{002404}


\finexercice


\exercice{2405, mayer, 2003/10/01}
\enonce{002405}{} 
\label{exoiter}
Soit $(X,d)$ un espace m\'etrique complet et soit $f: X\to X$ une
application telle que l'une de ces it\'er\'ees $f^n$ est
strictement contractante, i.e. il existe $\rho <1$ tel que
$$ d(f^n(x),f^n(y)) \leq \rho d(x,y) \quad \mbox{pour tout} \quad
x,y\in X \; .$$ Montrer que $f$ poss\`ede un unique point fixe.
Faire le rapprochement avec l'exercice \ref{exn13}.
\finenonce{002405}


\finexercice


\exercice{2406, mayer, 2003/10/01}
\enonce{002406}{} 
Soit $X = ({\cal C}^1([0,1]),N)$ avec $N(f) = \|f\|_\infty+\|f'\|_\infty$.
Montrer qu'il existe une fonction $f\in X$ qui est
point fixe de l'op\'erateur $T$ donn\'e par
$$ Tf(x)= 1 + \int _0^x f(t-t^2)\, dt \; .$$
On pourra commencer par \'etablir que $T\circ T$ est une
contraction. Utiliser ceci pour \'etablir l'existence d'une
fonction unique $f\in X$ qui v\'erifie $f(0)=1 $
et $f'(x) = f(x-x^2) $.
\finenonce{002406}


\finexercice


\exercice{2407, mayer, 2003/10/01}
\enonce{002407}{} 
Soient $y\in {\cal C}([a,b])$ et $k\in {\cal C}([a,b]\times
[a,b])$ des fonctions continues. On se propose de r\'esoudre
l'\'equation (int\'egrale de Fredholm) suivante:
\begin{equation}\label{eq 11}
 x(s) -\int _a^b
k(s,t) x(t) \, dt =y(s) \quad \text{pour} \; s\in [a,b] \;
\end{equation}
 d'inconnue $x\in{\cal C}([a,b])$. Pour ce faire on
suppose que le "noyau" $k$ satisfait l'hypoth\`ese suivante:
$$ \lambda := \max _{a\leq s\leq b} \int _a^b |k(s,t)|\, dt <1 \quad \left( \mbox{ou
m\^eme} \quad \max_{a\leq s,t\leq b} |k(s,t)| <\frac{1}{b-a}
\;\right) .$$
\begin{enumerate}
  \item Rappeler que $({\cal C}([a,b]),\| .\| _\infty)$ est
un espace complet.
\item Soit $x\in {\cal C}([a,b]) \mapsto Ax \in {\cal C}([a,b])$ l'application
donn\'ee par
$$(Ax)(s) :=\int _a^b k(s,t) x(t) \, dt +y(s)  \; .$$
Noter que (\ref{eq 11}) \'equivaut \`a $Ax=x$ et qu'on cherche donc
un point fixe de $x \mapsto Ax$. D\'eduire des hypoth\`eses faites
sur $k$ qu'un tel point fixe $x\in {\cal C}([a,b])$ existe et que
toute suite $A^n x_0$, $x_0\in{\cal C}([a,b])$, converge
uniform\'ement vers ce point fixe $x$.
\item {\it D\'ependance continue de la solution $x =
x(y)$}.

Soient $y_1,y_2\in{\cal C}([a,b])$ deux fonctions et $x_1,
x_2\in{\cal C}([a,b])$ les deux solutions associ\'ees de (\ref{eq
11}) ou, de fa\c{c}on \'equivalente, les points fixes des
applications associ\'ees $x\mapsto A_i x$. Montrer que
$$ \| x_1 - x_2\|_\infty = \| A_1 x_1 -A_2
x_2\|_\infty\leq \|y_1 -y_2\|_\infty+ \lambda \|x_1 -x_2 \|_\infty \;
.$$ En d\'eduire que
$$\|x_1 -x_2\|_\infty    \leq \frac{1}{1-\lambda}\| y_1
-y_2\|_\infty$$ et donc que la solution $x$ de (\ref{eq 11})
d\'epend continuement de la fonction $y$.
\end{enumerate}
\finenonce{002407}


\finexercice

\exercice{6232, queffelec, 2011/10/16}
\enonce{006232}{} 
\begin{enumerate}
\item Soit $X$ un espace métrique et $(f_n)$ une suite d'applications
conti\-nues à valeurs dans un espace métrique $Y$, convergeant vers $f$
uniformément sur $X$. Montrer que si
$(x_n)$ est une suite de points de $X$ convergeant vers $x\in X$, alors
$f_n(x_n)$ tend vers
$f(x)$.

\item Application : Soit $X$ un espace métrique compact, et soit $(f_n)$ une suite
d'applications continues de $X$ dans $X$, ayant chacune un point
fixe; on suppose que la suite $(f_n)$ converge vers une fonction $f$
uniformément sur $X$. Montrer que $f$ a aussi un point fixe.

\item Soit $K$ un convexe compact de $\Rr^n$ et $f$ une application continue de
$K$ dans $K$ vérifiant
$$\Vert f(x)-f(y)\Vert \leq \Vert x-y\Vert;$$
En considérant les fonctions $f_n$ définies sur $K$ par $f_n(x)={1\over
n}f(x_0)+(1-{1\over n})f(x)$, où $x_0\in K$, montrer que
$f$ a un point fixe. Est-il unique ? Que se passe-t-il si $K$ n'est plus
convexe ?
\end{enumerate}
\finenonce{006232}



\finexercice


\exercice{6233, queffelec, 2011/10/16}
\enonce{006233}{}
Soit $E$ un espace métrique compact, $f$ une application continue de $E$ dans
$E$ et on note $\Omega$ l'ensemble de ses points fixes.
\begin{enumerate}
\item Montrer que $\Omega$ est un compact, qui est non vide dans le cas où
$E=[a,b]$.

\item Si $\Omega=\emptyset$, montrer qu'il existe $r>0$ tel que $d(x,f(x))\geq r$
pour tout $x\in E$.

\item On suppose que $d(f(x),f(y))<d(x,y)$ pour tous $x\not=y$ de $E$. Montrer que 
$\Omega$ est réduit à un point $a$ et que pour tout choix initial de $x_0\in
E$, la suite récurrente $x_{n+1}=f(x_n)$ converge vers $a$.
\end{enumerate}
\finenonce{006233}



\finexercice


\exercice{6234, queffelec, 2011/10/16}
\enonce{006234}{}
Pour $x,y\in X=]0,+\infty[$ on pose $\delta(x,y)=\vert \log x-\log y\vert$
\begin{enumerate}
\item Montrer que $X$ muni de $\delta$ est complet alors qu'il ne l'est pas pour
la métrique usuelle de $\Rr$.

\item Soit $f$ une application de classe $C^1$ de $X$ dans $X$ vérifiant pour
tout $x\in X$  
$$x\vert f'(x)\vert\leq k f(x)$$
 où $k$ est un réel de $]0,1[$ fixé.
Montrer que $f$ a un seul point fixe dans $X$.
\end{enumerate}
\finenonce{006234}



\finexercice


\exercice{6235, queffelec, 2011/10/16}
\enonce{006235}{}
\begin{enumerate}
\item On considère une matrice $A=(a_{ij})$ à coefficients réels telle que
$\sum_{i,j=1}^n  a_{ij}^2 <1$.
En utilisant le théorème du point fixe, montrer que quels que soient les réels
$b_1,b_2,\cdots b_n$, le système d'équations linéaires
$$x_i-\sum_{j=1}^n  a_{ij} x_j = b_i,\  \  \  1\leq i\leq n$$
admet toujours une solution unique. En déduire $\det (I-A)\neq 0$.

\item Montrer sous les mêmes hypothèses que le système non linéaire
$$x_i-\sum_{j=1}^n  \sin(a_{ij} x_j) = b_i,\  \  \  1\leq i\leq n$$
admet une unique solution.
\end{enumerate}
\finenonce{006235}



\finexercice


\exercice{6236, queffelec, 2011/10/16}
\enonce{006236}{}
On va montrer qu'il existe une et une seule $h$ continue sur $[0,1]$ vérifiant
$h(0)=0$ et $h'(t)=\cos (th(t))$ pour tout $t\in[0,1]$. On note $E$ l'espace des
fonctions continues sur $[0,1]$ muni de la métrique uniforme.
\begin{enumerate}
\item $h$ est solution si et seulement si $h$ est continue et $h(s)=\int_0^s \cos
(th(t))\ dt$.

\item L'opérateur $T:E\to E$ défini par $Tf(s)=\int_0^s \cos
(tf(t))\ dt$ est $1/2$-contrac\-tant. Conclure.
\end{enumerate}
\finenonce{006236}



\finexercice


\exercice{6237, queffelec, 2011/10/16}
\enonce{006237}{}
Soit $a,b\in E$ evn, $B_1=\{x\in E/ \Vert x-a\Vert=\Vert x-b\Vert={1\over2}\Vert
a-b\Vert\}$, et pour $n>1$,  $B_n=\{x\in B_{n-1}/ \Vert x-y\Vert\leq
{1\over2}\delta(B_{n-1}),\
\forall y\in B_{n-1}\}$, où $\delta(B)$ désigne le diamètre de l'ensemble $B$.
\begin{enumerate}
\item Montrer que $\delta(B_{n})\leq {1\over2}\delta(B_{n-1})$, et que $\bigcap_n
B_n=\{{{a+b}\over2}\}$.

\item Soit $f$ une isométrie de $E$ {\it sur} $F$ evn, telle que $f(0)=0$.
Montrer en considérant la suite $(f(B_n))$ que pour tous $a,b\in E$,
$$f\Big({{a+b}\over2}\Big)={{f(a)+f(b)}\over2}.$$
En déduire que $f$ est une isométrie linéaire. Que peut-on dire plus
généralement d'une isométrie $f$ de $E$ sur $F$ ?
\end{enumerate}
\finenonce{006237}



\finexercice


\exercice{6238, queffelec, 2011/10/16}
\enonce{006238}{}
On va montrer qu'il existe une et une seule $h$ continue sur $[0,1]$ vérifiant
$h(0)=0$ et $h'(t)=\cos (th(t))$ pour tout $t\in[0,1]$. On note $E$ l'espace des
fonctions continues sur $[0,1]$ muni de la métrique uniforme.

\begin{enumerate}  
\item $h$ est solution si et seulement si $h$ est continue et $h(s)=\int_0^s \cos
(th(t))\ dt$.

\item L'opérateur $T:E\to E$ défini par $Tf(s)=\int_0^s \cos
(tf(t))\ dt$ est $1/2$-contrac\-tant. Conclure.
\end{enumerate}
\finenonce{006238}



\finexercice       


\exercice{6239, queffelec, 2011/10/16}
\enonce{006239}{}
On désigne par $E$ l'espace $C([0,1])$ muni de la norme uniforme et l'opérateur
$A$ défini par $$Af(x)=\int_0^x tf(t)\ dt + x\int_x^1 f(t)\ dt$$
pour $f\in E$ et $x\in [0,1]$.

\begin{enumerate}  
\item Vérifier que $A$ est continu et calculer sa norme opérateur.

\item L'équation $Af=f$ a-t-elle dans $E$ des solutions $f$ non nulles ?
\end{enumerate}
\finenonce{006239}



\finexercice       


\exercice{6240, queffelec, 2011/10/16}
\enonce{006240}{}
On considère $T:C([0,1])\to C([0,1])$ qui à $f$ associe $F$ définie par 

$$F(t)= \left \{\begin{array}{ccc}
    & 3/4\  f(3t) & {\rm si}\ \ 0\leq t\leq 1/3\\
    & 1/4 + 1/2\  f(2-3t)   &{\rm si}\ \ 1/3\leq t\leq 2/3\\
&1/4 +3/4\ f(3t-2)&{\rm si}\ \ 2/3\leq t\leq 1 

\end{array}\right.$$

\begin{enumerate}  
\item Vérifier que $F$ est bien continue et que $T$ est $3/4$-contractante.

\item On note $h$ le point fixe de $T$. Montrer par récurrence $\vert
h({{k-1}\over3^n})-h({k\over3^n})\vert\geq 2^{-n}$. 

Soit $a\in[0,1]$; montrer qu'il existe une suite $(t_n)$ telle que $\lim t_n
=a$ et $\lim \vert {{h(t_n)-h(a)}\over{t_n-a}}\vert = +\infty$.

\item En déduire l'existence d'une fonction
continue nulle part dérivable.
\end{enumerate}
\finenonce{006240}



\finexercice       


\exercice{6790, gijs, 2011/10/16}
\enonce{006790}{}
Soit $I\subset \Rr$ un intervalle fermé borné et
soit $f: I \to I$ une application dérivable verifiant
$\forall x\in I\ :\ |f'(x)| < 1$. Pour un $x_0 \in I$ on
définit la suite récurrente $a : \Nn \to I$ par
$a_0 = x_0$, $a_{n+1} = f(a_n)$. On vous demande de
montrer qu'il existe un unique $\ell \in I$, indépendant
de $x_0$, tel que $ \lim_{n\to \infty} a_n = \ell$.
Les questions suivantes peuvent vous guider dans la
démonstration.

\begin{enumerate}

\item Montrer que $f$ admet un point fixe unique $\ell$. 

\item Montrer que $d(a_n,\ell)$ converge.

\item Montrer que la suite $a$ admet une sous-suite
convergente $b$.

\item Notons $ \lim_{k\to \infty} b_k = \beta$
et  $ \lim_{n\to \infty} d(a_n,\ell) = r$.
Montrer~: $d(\beta,\ell) = r = d(f(\beta),\ell)$. 
\end{enumerate}


Est-ce-que le résultat est vrai si $I$ n'est pas
fermé ?  \ Justifier votre réponse.
\finenonce{006790}



\finexercice


\exercice{6805, gijs, 2011/10/16}
\enonce{006805}{}
Soit $f:\Rr^2 \to \Rr^2$ définie par $f(x,y) =
(x^2-y, x^2 + y^2)$. On définit la fonction $g:\Rr^2 \to \Rr^2$ par $g = f \circ f$.
\begin{enumerate}
\item Montrer que $f$ et $g$ sont de classe $C^1$.

\item Calculer pour tout $(x,y) \in \Rr^2$ la matrice
Jacobienne de $f$ en $(x,y)$ notée $D_{(x,y)}f$; 
calculer la matrice Jacobienne de $g$ en $(0,0)$ notée
$D_{(0,0)}g$.

\item Montrer qu'il existe $\rho>0$ tel que pour tout $(x,y) \in
\overline{B_\rho((0,0))}$ (la boule fermée de centre
$(0,0)$ et de rayon $\rho$) on a $\Vert D_{(x,y)}g \Vert
\le \frac12$. 

\item
Montrer que la fonction $g$ admet un unique point fixe
dans $\overline{B_\rho((0,0))}$ avec $\rho$ comme dans 3).
\end{enumerate}
\finenonce{006805}



\finexercice

\section{427.00 Espace de Hilbert, théorème de projection}


\exercice{1894, legall, 2003/10/01}

\enonce{001894}{}
\begin{enumerate}
\item Montrer que l'application $(f,g)\mapsto \langle f,g\rangle 
=\displaystyle{\int _0 ^1f(t)g(t)dt}$ est un produit scalaire 
euclidien
sur $C[0,1]$, l'espace vectoriel des fonctions continues sur $[0,1]$ 
\`a valeurs r\'eelles.
\item On note $C=\{ f\in C[0,1] ; \displaystyle{\int _0 ^1f(t)dt=1}$. 
Montrer que $\displaystyle{ \inf _{f\in C} \int _0 ^1f^2(t)dt=1}$
et que cette borne inf\'erieure est atteinte.
\end{enumerate}
\finenonce{001894}



\finexercice


\section{428.00 Théorème de Baire}


\exercice{2392, queffelec, 2003/10/01}
\enonce{002392}{}
 \`A l'aide du th\'eor\`eme de Baire, montrer qu'un ferm\'e d\'enombrable non vide $X$ de
$\Rr$ a au moins un point isol\'e.
\emph{Indication :} on pourra considérer $\omega_x = X \setminus \{x\}$.

Que peut-on dire de l'ensemble de Cantor ?
\finenonce{002392}


\finexercice       


\exercice{2393, queffelec, 2003/10/01}
\enonce{002393}{}
Soit $f$ une application d\'efinie sur un espace m\'etrique complet $(X,d)$,
\`a valeurs r\'eelles et semi-continue inf\'erieurement. Montrer qu'il existe un
ouvert non vide $O$ sur lequel $f$ est major\'ee.

Application : soit $(f_n)$ une suite de formes lin\'eaires continues sur un Banach $B$,
v\'erifiant
$$\forall {x\in B},\;  \sup_n\vert f_n(x)\vert<\infty.$$

En utilisant ce qui pr\'ec\`ede, montrer que $ \sup_n\Vert f_n\Vert<\infty$.
\finenonce{002393}


\finexercice       


\exercice{2394, queffelec, 2003/10/01}
\enonce{002394}{}
On sait que $l^1$ est inclus dans $l^2$ (au fait pourquoi ?) mais n'est pas
ferm\'e dans $l^2$ (re-pourquoi ?); on va montrer qu'il est de premi\`ere cat\'egorie
dans
$l^2$ c.a.d. r\'eunion d\'enombrable de ferm\'es d'int\'erieur vide (dans $l^2$).

\begin{enumerate}  
\item On consid\`ere pour chaque $p\geq1$, 
$$F_p=\{(a_n)\in l^2\, / \sum\vert a_n\vert \leq p\}$$
Montrer que $F_p$ est ferm\'e dans $l^2$ et d'int\'erieur vide.

\item En d\'eduire le r\'esultat.
\end{enumerate}
\finenonce{002394}



\finexercice 

\exercice{6141, queffelec, 2011/10/16}
\enonce{006141}{}
Montrer que $\Qq$ n'est pas un $G_\delta$ c'est-à-dire n'est pas intersection
dénombrable d'ouverts de $\Rr$.

\emph{Indication :} on pourra raisonner par l'absurde et considérer $\omega_n={\Rr}\backslash\{q_n\}$ si ${\Qq}=\{q_1,\ldots,q_n,\ldots\}$.
\finenonce{006141}



\finexercice       


\exercice{6142, queffelec, 2011/10/16}
\enonce{006142}{}
Soit $B$ un espace de Banach; on rappelle que tout sous-espace propre de $B$
est d'intérieur vide dans $B$.
Montrer que si $B$ est de dimension infinie, $B$ ne possède pas de base
algébrique dénombrable.

En déduire que l'espace des polyn\^omes n'est complet pour aucune norme.
\finenonce{006142}



\finexercice       


\exercice{6143, queffelec, 2011/10/16}
\enonce{006143}{}
Soit $(X,{\cal T})$ un espace topologique de Baire, c'est-à-dire pour lequel le
théo\-rème de Baire est valide. On va montrer que tout ouvert de $X$ muni de la
topologie induite est encore un espace de Baire.

\begin{enumerate}  
\item Soit $(O_n)$ une suite d'ouverts denses dans $O$; montrer que chaque
$\omega_n=O_n\cup\overline{O}^c$ est un ouvert dense dans $X$ (on rappelle
qu'un ensemble est dense dans $X$ s'il rencontre tout ouvert de $X$).

\item Montrer que pour tout ouvert $\omega$ de $O$,
$$(\cap_nO_n)\cap\omega\neq\emptyset$$
En déduire le résultat.
\end{enumerate}
\finenonce{006143}



\finexercice       

\section{429.00 Dualité, topologie faible}


\exercice{6124, queffelec, 2011/10/16}
\enonce{006124}{} 
Soit $E$ un evn, $f$ un élément non nul du dual de $E$, et $L$
l'hyperplan affine $\{x\in E/f(x)=1\}$.

\begin{enumerate}  
\item Montrer que $$\inf_{x\in L}\Vert x\Vert\geq {1\over{\Vert f\Vert}}.$$
\item On peut trouver dans la sphère unité une suite $(x_n)$ telle que $\vert
f(x_n)\vert\geq {n\over{n+1}}\Vert f\Vert$ (justifier) et, à l'aide de cette
suite, montrer que l'on a finalement 
$$\inf_{x\in L}\Vert x\Vert= {1\over{\Vert f\Vert}}.$$
\end{enumerate}
\finenonce{006124}



\finexercice       


\exercice{6125, queffelec, 2011/10/16}
\enonce{006125}{}
Soit $E=C([0,1])$, $\mu(x)=\int_0^1 x(t)\ dt$, $\mu_n(x)={1\over n}\sum_{k=1}^n
x({k\over n})$.

\begin{enumerate}  
\item Calculer $\Vert\mu\Vert$ et $\Vert\mu_n\Vert$.

\item Montrer que $\mu_n(x)$ converge vers $\mu(x)$ pour tout $x$ dans $E$, mais
que $\Vert\mu-\mu_n\Vert=2$.
\end{enumerate}
\finenonce{006125}



\finexercice       


\exercice{6126, queffelec, 2011/10/16}
\enonce{006126}{}
Soit $E=C([0,1])$ et $(t_n)$ une suite de points distincts, convergente dans
$[0,1]$. Montrer que
$f$ définie par $f(x)=\sum_1^\infty {{(-1)^n}\over{2^n}} x(t_n)$ est un élément
de
$E'$ de norme $1$ qui n'atteint sa norme en aucun point de la boule unité de $E$.
\finenonce{006126}



\finexercice       


\exercice{6127, queffelec, 2011/10/16}
\enonce{006127}{}
Soit $a,b\in E$ evn, $B_1=\{x\in E/ \Vert x-a\Vert=\Vert x-b\Vert={1\over2}\Vert
a-b\Vert\}$, et pour $n>1$,  $B_n=\{x\in B_{n-1}/ \Vert x-y\Vert\leq
{1\over2}\delta(B_{n-1}),\
\forall y\in B_{n-1}\}$, où $\delta(B)$ désigne le diamètre de l'ensemble $B$.

\begin{enumerate}  
\item Montrer que $\delta(B_{n})\leq {1\over2}\delta(B_{n-1})$, et que $\bigcap_n
B_n=\{{{a+b}\over2}\}$.

\item Soit $f$ une isométrie de $E$ {\it sur} $F$ evn, telle que $f(0)=0$.
Montrer en considérant la suite $(f(B_n))$ que pour tous $a,b\in E$,
$$f({{a+b}\over2})={{f(a)+f(b)}\over2}.$$
En déduire que $f$ est une isométrie linéaire. Que peut-on dire plus
généralement d'une isométrie $f$ de $E$ sur $F$.

\item On note $l^\infty_n$ l'espace $\Rr^n$ muni de la norme
$\sup_{1\leq i\leq n}\vert x_i\vert$, et on considère l'application $f:
l^\infty_n\to l^\infty_{n+1}$ définie par
$f(x_1,\cdots,x_n)=(x_1,\cdots,x_n,\sin x_1)$. Vérifier que $f$ est une
isométrie non linéaire entre evn; pourquoi n'a-t-on pas de contra\-diction avec
ce qui précède ?
\end{enumerate}
\finenonce{006127}



\finexercice       


\exercice{6128, queffelec, 2011/10/16}
\enonce{006128}{}
\begin{enumerate}  
\item Soit $(u_n)$ une suite de nombres complexes. On suppose que, pour toute suite
bornée de complexes $(v_n)$, la série $\sum u_nv_n$ converge. Montrer que
$(u_n)$ est dans l'espace $l^1$.

\item Soit $(u_n)$ une suite de nombres complexes. On suppose que, pour toute
suite $(v_n)$ dans $l^2$, la série $\sum u_nv_n$ converge. Montrer
que
$(u_n)$ est dans l'espace $l^2$.

\emph{Indication :} Soit $(a_n)$ une série positive divergente. Montrer que la série de terme
général ${{a_n}\over{S_n^\alpha}}$, où $S_n=\sum_{k=0}^n a_k$ converge si
$\alpha>1$ et diverge sinon. Utiliser ensuite cette remarque pour conduire un
raisonnement par l'absurde.
\end{enumerate}
\finenonce{006128}



\finexercice       


\exercice{6129, queffelec, 2011/10/16}
\enonce{006129}{}
\label{gijsexodual}
On va montrer que le dual de $l^2$ est isométriquement isomorphe à $l^2$.
On note comme d'habitude $e_n$ l'élément de $l^2$ dans toutes les composantes
sont nulles, sauf la $n$-ième qui vaut $1$.  
\begin{enumerate}  
\item Soit $x\in l^2$. Montrer que la suite d'éléments de $l^2$ $x_n=\sum_1^n
x(k)e_k$ converge vers $x$ dans $l^2$ (autrement dit, les suites nulles 
à partir d'un certain rang sont denses dans $l^2$.)
En déduire que si $f\in(l^2)'$, $f(x)=\sum_1^\infty x(n)f(e_n).$

\item Montrer que $\Vert f\Vert \geq (\sum_1^n\vert f(e_k)\vert^2)^{1\over2}$, et
que $(f(e_n))_n$ est un élément de $l^2$.

\item Montrer alors que pour tout $x\in l^2$, $\vert f(x)\vert\leq \Vert
x\Vert_2\Vert (f(e_n))\Vert_2$, et que  $\Vert f\Vert=\Vert (f(e_n))\Vert_2$.

En déduire que l'application $f\to(f(e_n))$ est un isomorphisme isométrique du
dual de $l^2$ sur $l^2$.

\end{enumerate}
\finenonce{006129}



\finexercice       


\exercice{6130, queffelec, 2011/10/16}
\enonce{006130}{}
En suivant la même démarche que l'exercice \ref{gijsexodual}, montrer que le dual topologique de $c_0$ est
isométriquement isomorphe à $l^1$. 
\finenonce{006130}



\finexercice  

\section{430.00 Connexité}


\exercice{2383, mayer, 2003/10/01}
\enonce{002383}{}
\label{exocon}
Soit $X$ un espace m\'etrique. 
\begin{enumerate}
\item Montrer que $X$ est connexe si et seulement si toute application
continue $f:X\to \{0,1\}$ est constante.
\item Soit $A$ une partie de $X$ connexe. Montrer que toute partie
$B \subset E$ v\'erifiant $A \subset B\subset \overline{A}$ est
connexe.
\item Si $(A_n)_{n\geq 0}$ est une suite de parties connexes de
$X$ telle que $A_n \cap A_{n+1}\neq \emptyset$ pour tout $n\geq
0$. Prouver que $\bigcup _{n\geq 0} A_n $ est connexe.
\end{enumerate}
\finenonce{002383}


\finexercice


\exercice{2384, mayer, 2003/10/01}
\enonce{002384}{} 
D\'eterminer les parties connexes de
$$\{(x,y)\in \Rr^2 \; ; x\neq y\} \quad  \text{et  de} \quad
\{(z,w) \in \Cc^2 \; ; z\neq w\}\; .$$
\finenonce{002384}


\finexercice


\exercice{2385, mayer, 2003/10/01}
\enonce{002385}{} 
Soit $A$ et $B$ des parties de $X$. On suppose $B$ connexe et que
$B\cap A$ et $B \cap \complement A$ sont non vides. Montrer que $B$
coupe la fronti\`ere de $A$.
\finenonce{002385}


\finexercice


\exercice{2386, mayer, 2003/10/01}
\enonce{002386}{} 
Notons $T=\{0\}\times [-1,1]\cup [-1,1]\times \{0\}$ muni de la
topologie induite par celle de $\Rr^2$.
\begin{enumerate}
\item Montrer que $T$ est compact et connexe et que $f(T)$ est un
segment si $f:T\to \Rr$ est une fonction continue.
\item D\'eterminer les points $x\in T$ pour lesquels $T\setminus \{x\}$ est connexe.
\item Montrer que $T$ n'est hom\'eomorphe \`a aucune partie de
$\Rr$.
\end{enumerate}
\finenonce{002386}


\finexercice


\exercice{2387, mayer, 2003/10/01}
\enonce{002387}{}  
\begin{enumerate}
\item Montrer qu'il existe une surjection continue de $\Rr $ sur
$\mathbb{S}^1 =\{z\in \Cc \; ; \;\; |z|=1\}$ et qu'il n'existe pas
d'injection continue de $\mathbb{S}^1$ dans $\Rr$.
\item Montrer qu'il n'existe pas d'injection continue de $\Rr^2$
dans $\Rr$.
\end{enumerate}
\finenonce{002387}


\finexercice


\exercice{2388, mayer, 2003/10/01}
\enonce{002388}{} 
Dans $\Rr^2$, soit $B_a$ l'ensemble $\{a\}\times ]0,1]$ si $a$
est rationnel et $B_a=\{a\}\times [-1,0]$ si $a$ est irrationnel.
Montrer que $B = \bigcup _{a\in \Rr } B_a$ est une partie
connexe de $\Rr^2$.
\finenonce{002388}


\finexercice


\exercice{2389, mayer, 2003/10/01}
\enonce{002389}{} 
Soit $I$ un intervalle ouvert de $\Rr$ et soit $f:I\to \Rr$
une application d\'erivable. Notons $A=\{(x,y)\in I\times I \; ;\;
\; x<y \}$.
\begin{enumerate}
\item Montrer que $A$ est une partie connexe de $\Rr^2$.
\item Pour $(x,y)\in A$, posons $g(x,y) =\frac{f(y)-f(x)}{y-x}$.
Montrer que $g(A)\subset f'(I) \subset \overline{g(A)}$.
\item Montrer que $f'(I)$ est un intervalle.
\end{enumerate}
Ce r\'esultat signifie que {\it la d\'eriv\'ee de toute fonction
d\'erivable poss\`ede la propri\'et\'e de la valeur interm\'ediaire}
 (un th\'eor\`eme de Darboux).
\finenonce{002389}


\finexercice


\exercice{2390, mayer, 2003/10/01}
\enonce{002390}{} 
Soit $X$ un espace m\'etrique et $(A_i)_{i\in I}$ une famille de
parties connexes par arcs de $X$ telle que $\bigcap _{i\in I}A_i
\neq \emptyset$. Montrer que $\bigcup _{i\in I}A_i$ est connexe
par arcs.
\finenonce{002390}


\finexercice


\exercice{2391, mayer, 2003/10/01}
\enonce{002391}{}
Dans $\Rr^2$ on consid\`ere l'ensemble $A= \{ (x, \sin(
\frac{1}{x}))\; ; \; x >0 \}$.
\begin{enumerate}
\item Montrer que $A$ est une partie connexe et connexe par arcs
de $\Rr^2$.
\item D\'eterminer $\overline{A}$ et justifier que
$\overline{A}$ est connexe.
\item Montrer que $\overline{A} $ n'est pas connexe par arcs.
\end{enumerate}
\finenonce{002391}



\finexercice

\exercice{6144, queffelec, 2011/10/16}
\enonce{006144}{}
Montrer que $\Zz$ et $\Qq$ (munis de la topologie induite par celle
de $\R$) ne sont pas homéomorphes, mais sont tous les deux ``totalement
discontinus" au sens  suivant : leurs seuls connexes sont les points.
(Remarquer que A connexe dans $Y\Rightarrow A$ connexe dans $X$ si
$Y\subset X$).
\finenonce{006144}



\finexercice


\exercice{6145, queffelec, 2011/10/16}
\enonce{006145}{}
Soit $A$ une partie du cercle unité ${\mathbb{S}^1}=\partial D$~; montrer que
$D\cup A$ est connexe.
\finenonce{006145}



\finexercice


\exercice{6146, queffelec, 2011/10/16}
\enonce{006146}{}
\begin{enumerate}
\item Montrer qu'il y a équivalence pour $X$ espace topologique entre :

\hskip3mm i)  Toute application continue $\varphi  : X\to {\Zz}$ est constante.

\hskip3mm ii) $X$ est connexe.

\item   Retrouver ainsi différents résultats du cours ($f(C)$ connexe si $C$ connexe
et $f$ continue; $B$ connexe si $A$ connexe et $A\subset B\subset \overline A$;
un produit de deux connexes est encore connexe; etc)

\item Soit $A,B$ connexes de $X$ tels que $\overline A\cap B\not =\emptyset $~;
montrer à l'aide de a) que $A\cup B$ est connexe.
\end{enumerate}
 \finenonce{006146}



\finexercice


\exercice{6147, queffelec, 2011/10/16}
\enonce{006147}{}
Existe-t-il une application continue $f:{\R}\to{\R}$ telle que
$f(x)\in\Qq$ si $x\notin \Qq$ et $f(x)\notin\Qq$ si $x\in \Qq$ ?
(Regarder l'image de $f$.)
\finenonce{006147}



\finexercice


\exercice{6148, queffelec, 2011/10/16}
\enonce{006148}{}
Soit $X=\Qq$ muni de la topologie induite par celle de $\Rr$. Montrer que
les seuls connexes de $X$ sont les points. ($A$ connexe dans $X\Rightarrow A$
connexe dans $\Rr$)
\finenonce{006148}



\finexercice


\exercice{6149, queffelec, 2011/10/16}
\enonce{006149}{}
Soit $X$ un ouvert d'un espace vectoriel normé $E$ ; montrer que $X$ est connexe
si et seulement si il est connexe par arcs. 
(\emph{Indication :} fixer $a\in X$ et
considérer $A=\{x\in X, \hbox{relié à}\ a\ \hbox{par un chemin dans}\
X\}$.) 
\finenonce{006149}



\finexercice


\exercice{6150, queffelec, 2011/10/16}
\enonce{006150}{}
Soit $f$ une surjection continue de $\Rr^2$ sur $\Rr$. Montrer que l'image
réciproque de tout point est non bornée (raisonner par l'absurde et utiliser
que le complémen\-taire d'un disque dans $\Rr^2$ est connexe).
\finenonce{006150}



\finexercice


\exercice{6151, queffelec, 2011/10/16}
\enonce{006151}{}
 Soit $P\in{\Cc}[X]$ un polyn\^ome de racines $z_1,\ldots ,z_n$  distinctes ou
non, situées dans un convexe $K$ de ${\Cc}$.

\begin{enumerate}  \item On suppose que $P'(z)=0$ et $z\notin \{z_1,\ldots ,z_n\}$~; montrer qu'il
existe des réels
$\lambda _1(z),\ldots ,\lambda _n(z)$, inconnus mais $>0$, tels que l'on ait :
$\sum^n_{k=1} \lambda _k(z)(z-z_k)=0$.  (Indication~: considérer
${P'(z)\over P(z)}$ et son conjugué).

\item Montrer que $P'$ a aussi toutes ses racines dans $K$ (théorème de
Gauss-Lucas).
\end{enumerate}
\finenonce{006151}



\finexercice


\exercice{6152, queffelec, 2011/10/16}
\enonce{006152}{}
On dit qu'un espace topologique possède la propriété du point fixe si
toute fonction continue de $X$ dans $X$ admet un point fixe.
\begin{enumerate}
\item Montrer qu'un espace topologique possédant cette propriété est
nécessai\-rement connexe.

\item Montrer que si
$X$ a cette propriété, tout
$Y$ homéomorphe à
$X$ la possède aussi. 

\item Montrer ainsi que ${\mathbb{S}^1}$ n'est pas homéomorphe à un segment.
\end{enumerate}
\finenonce{006152}



\finexercice


\exercice{6153, queffelec, 2011/10/16}
\enonce{006153}{}
Soit $I=[a,b]$ et $f : I\to\R$ dérivable~; soit $A=\big\{(x,y)\in
I\times I; y>x\big\}$  et $g : A\to \R$ définie par
$g(x,y)={f(y)-f(x)\over y-x}$. 
\begin{enumerate}
\item Montrer que $g(A)\subset f'(I)\subset\overline {g(A)}$.

\item Montrer que $f'$ a la propriété de la valeur intermédiaire~: si elle
prend les valeurs
$\alpha $ et $\beta $, elle prend toute valeur $\gamma \in[\alpha ,\beta ]$.
\end{enumerate}
\finenonce{006153}



\finexercice


\exercice{6154, queffelec, 2011/10/16}
\enonce{006154}{}
On va démontrer à l'aide de la connexité, le résultat classique:

``$f:{\R}\to{\R}$ continue injective $\Longrightarrow$ $f$ strictement
monoton''.

\noindent Pour cela, considérons l'application $F$ définie sur $\R^2$ par
$F(x,y)=f(x)-f(y)$ et $C=\{(x,y)\in{\R^2}\ /\ x>y\}$
\begin{enumerate}
\item Montrer que $F(C)$ est un connexe de $\R$.

\item En déduire le résultat.
\end{enumerate}
\finenonce{006154}



\finexercice


\exercice{6155, queffelec, 2011/10/16}
\enonce{006155}{}
On définit la projection stéréographique $h$ de ${\mathbb{S}^1}$ sur ${\Rr}
\cup\{\infty\}$, $h(x,y)$ étant le point d'intersection avec l'axe
réel de la droite issue de $(0,1)$  passant par $(x,y)$ si
$(x,y)\not=(0,1)$  et $h(0,1)=\infty$. Vérifier qu'il s'agit d'un
homéomorphisme. En déduire que ${\R}\cup\{\infty\}$ et
${\R}\cup\{+\infty\}\cup\{-\infty\}$ ne sont pas homéomorphes.
\finenonce{006155}



\finexercice


\exercice{6156, queffelec, 2011/10/16}
\enonce{006156}{} 
Soit $X$ un espace métrique. \'Etablir l'équivalence des assertions suivantes:
\begin{enumerate}
\item $X$ est compact connexe.
\item Pour tout recouvrement ouvert $(U_i)_{i\in I}$, il existe $n\in \Nn$ et
$i_1,...,i_n \in I$ tels que
$$\bigcup_{k=1}^n U_{i_k} =X \quad et \quad U_{i_k} \cap U_{i_{k+1}}\neq \emptyset \; \text{pour} \; k=1,...,n-1\; .$$
\end{enumerate}
\finenonce{006156}



\finexercice


\exercice{6157, queffelec, 2011/10/16}
\enonce{006157}{}
$A$ et $B$ sont des parties d'un espace topologique $X$. Vrai ou faux ? 
\begin{enumerate}
\item Si $A$ est connexe, $\partial A$ est connexe ?

\item Si $\overline A$ est connexe, $A$ est connexe ?

\item Si $A$ et $B$ sont connexes et $A\cap B\not=\emptyset$, $A\cap B$
est connexe ?

\item Si $X$ est un evn et $A$ et $B$ convexes avec $A\cap B\not=\emptyset$,
$A\cap B$ est connexe ?

\item Si $A$ et $B$ sont connexes, $A\cup B$
est connexe ?

\item Soit $f$ continue de $X$ dans $Y$ espace topologique. Si $A$ est connexe
par arcs , $f(A)$ est connexe par arcs ?

\item Soit $f$ continue de $X$ dans $Y$ evn. Si $A$ est convexe, $f(A)$ est
convexe ?
\end{enumerate}
\finenonce{006157}



\finexercice


\exercice{6158, queffelec, 2011/10/16}
\enonce{006158}{}

Dans $\Rr^2$ on considère l'ensemble $A$ des points dont une coordonnée au
moins est irrationnelle.
\begin{enumerate}
\item Soit  $\alpha\in {\Rr\backslash \Qq}$; décrire l'ensemble $A\cap\{(x,y)\in
{\Rr^2}, x=\alpha\}.$

\item Montrer que $A$ est connexe par arcs (plus précisément deux points de $A$
peuvent être reliés par une ligne polygonale).
\end{enumerate}
\finenonce{006158}



\finexercice


\exercice{6159, queffelec, 2011/10/16}
\enonce{006159}{}
\begin{enumerate}
\item Montrer que dans ${\Rr^n}, n\geq 2$, les sous-ensembles suivants sont
 connexes : 
$$ B(0,r);\  {\Rr^n}\backslash B(0,r);\ S^{n-1}(0,r)=\{x\in{\Rr^n}\ /\
||x||=r\}.$$

\item Montrer que $\Rr$ et $\Rr^2$ ne sont pas homéomorphes ( sinon enlever un
point à 
$\Rr$).
\end{enumerate}
\finenonce{006159}



\finexercice


\exercice{6160, queffelec, 2011/10/16}
\enonce{006160}{}
On rappelle que si $X$ est réunion disjointe de parties non vides $\omega_i$
ouvertes et connexes, les $\omega_i$ sont les composantes connexes de $X$.

Trouver les composantes connexes du complémentaire des ensembles sui\-vants :
$$\quad \{(x,y)\in{\Rr^2} / y^2-x=0\};\quad \{(x,y,z)\in{\Rr^3} /
0<x^2+y^2+z^2\leq1\};$$
$$ S^{n-1}=\{x\in{\Rr^n} / ||x||=1\};\quad {\Qq}\times {\Qq}\subset{\Rr^2}.$$
\finenonce{006160}



\finexercice


\exercice{6161, queffelec, 2011/10/16}
\enonce{006161}{}
Soit $H$ un sous-espace vectoriel de $\Rr^n$, $n\geq2$. Montrer que
\begin{enumerate}
\item si dim $H = n-1$, ${\Rr^n}\backslash H$ a deux composantes connexes;

\item si dim $H \leq n-2$, ${\Rr^n}\backslash H$ est connexe. 
\end{enumerate}
\finenonce{006161}



\finexercice


\exercice{6162, queffelec, 2011/10/16}
\enonce{006162}{}
On considère le sous-ensemble suivant du plan complexe :

$$C=\cup_{n\geq1} [0,1+{i\over n}]\cup\ [{1\over2},1]= A\cup\ [{1\over2},1]$$

\begin{enumerate}
\item Montrer que $C$ est connexe.

\noindent Soit $\gamma$ un chemin reliant un point de $A$ à un point de
$[{1\over2},1]$ et d'image dans $C$.

\item Si $\gamma$ ne passe pas par $0$, montrer que $\gamma(t)=r(t)e^{i\theta(t)}$
où $r(t)>0$ et $0\leq\theta(t)<{\pi\over2}$, et $r,\theta$ continues.

\item Montrer que $\theta$ ne prend qu'un nombre dénombrable de valeurs et
aboutir à une contradiction.

\item Dans tous les cas, montrer qu'il existe $t_0\in]0,1[$ tel que $\gamma(t)$ ne
passe pas par $0$ pour $t\geq t_0$. En déduire que $C$ n'est pas connexe par
arcs.  
\end{enumerate}
\finenonce{006162}



\finexercice


\exercice{6163, queffelec, 2011/10/16}
\enonce{006163}{}
Soit $f$ une application continue de $[a,b]$ dans $\Rr$ vérifiant
$$f({x+y\over2})\leq {1\over2}\Big(f(x)+f(y)\Big)\quad \forall x,y\in [a,b].$$
\begin{enumerate}
\item On suppose $f(a)=f(b)=0$. On considère $E=\{x\in ]a,b[\ /\
f(x)=\displaystyle\sup_{t\in[a,b]} f(t)\}$. Montrer que $E$ est ouvert et
fermé dans
$]a,b[$. En déduire que $f$ est $<0$ ou identiquement nulle sur $]a,b[$.

\item Montrer dans tous les cas que $f$ est convexe ie $f$ vérifie
$f((1-t)x+ty)\leq (1-t)f(x)+tf(y)$ pour tous $x,y\in [a,b]$ et $t\in[0,1]$ (On se
ramènera au cas a) en considérant
 $f$ privée de sa corde sur $[a,b]$).
\end{enumerate}
\finenonce{006163}



\finexercice


\exercice{6773, gijs, 2011/10/16}
\enonce{006773}{} 
Soit $A$ un ouvert connexe non-vide de $\Rr^n$ et $a\in A$. Soit $G_a$ l'ensemble des points de $A$
pouvant être reliés à $a$ par un chemin contenu
dans $A$. 
\begin{enumerate}
\item Montrer que $G_a$ et $A\setminus G_a$ sont
ouverts. 
\item En déduire que $A$ est connexe par arcs.
\end{enumerate}
\finenonce{006773}



\finexercice

\exercice{6782, gijs, 2011/10/16}
\enonce{006782}{}
\begin{enumerate}
\item Donner la {\it définition\/} d'un intervalle
dans $\Rr$ et montrer que si $A\subset \Rr$ est
connexe, alors $A$ est un intervalle.

\item En utilisant la notion de connexité, montrer
qu'il n'existe pas un homéo\-mor\-phisme $f: \Rr^2
\to \Rr$.
\end{enumerate}
\finenonce{006782}



\finexercice


\exercice{6814, gijs, 2011/10/16}
\enonce{006814}{}
Soit $X$ un espace topologique, $I$ un ensemble d'indices
et pour chaque $\alpha \in I$, soit $A_\alpha \subset X$
une partie connexe de $X$. On suppose que $\forall
\alpha,\beta \in I : A_\alpha \cap A_\beta \neq
\emptyset$. Démontrer que $ {\cup}_{\alpha\in
I} A_\alpha$ est connexe.
\finenonce{006814}



\finexercice


\exercice{6815, gijs, 2011/10/16}
\enonce{006815}{}
Soit $X$ un espace topologique, $A\subset X$ une partie
et $C\subset X$ une partie connexe telle que $C \cap A
\neq \emptyset$ et $C \cap (X\setminus A) \neq \emptyset$.
\begin{enumerate}
\item Démontrer que $C$ contient des points de la
frontière de $A$.

\item Pourquoi $A = \{(x,y,0) \mid x^2 + y^2 \le 1\}
\subset \Rr^3$ et $C = \{(0,0,z) \mid |z|\le1\}$
n'est-il pas un contre-exemple pour la propriété
énoncée dans 1. ?
\end{enumerate}
\finenonce{006815}



\finexercice


\exercice{6834, gijs, 2011/10/16}
\enonce{006834}{}
Soit $X$ un espace topologique connexe, soit $\mathcal{U}$ un
recouvrement de $X$ par ouverts et soit $x_0 \in X$ un
point de $X$. 
On dit que $x$ est $n$-éloigné de $x_0$ s'il existe
$U_0, \dots, U_n \in \mathcal{U}$ tels  que $x_0 \in U_0$,
$U_{i-1} \cap U_i \neq \emptyset$, $i=1, \dots, n$ et
$x\in U_n$. Prenez le temps de dessiner ce que veut dire
$n$-éloigné. On définit l'ensemble $A\subset X$ par
$$ 
A = \{ \,x\in X \mid \exists n \in \Nn : \text{$x$ est
$n$-éloigné de $x_0$}\,\}
\ .
$$
Démontrer que $A$ est ouvert et fermé dans $X$. En
déduire que $A=X$.
\finenonce{006834}



\finexercice

\section{431.00 Autre}


\exercice{1864, gineste, 2001/11/01}

\enonce{001864}{In\'egalit\'e de Cauchy-Schwarz}
\begin{enumerate}
\item Montrer que : $ \forall x_1,x_2, \cdots ,
 x_n \in \R \; \; (x_1 + x_2 + \cdots + x_n)^2 \leq n(x_1^2 + x_2^2 + \cdots + x_n^2) $
\item D\'eterminer : $m = \text{Inf}
\{ (\sum_{i=1}^n x_i)(\sum_{i=1}^n 1/x_i) \mbox{ tels que }  x_1,x_2, \cdots , x_n > 0 \} $
\item D\'eterminer : $M= \sup \{ |x + 2y + 3z + 4t| \mbox{ tels que } (x,y,z,t) \in \R^4 \; , x^2+y^2+z^2+t^2 \leq 1 \} $
\end{enumerate}

\finenonce{001864}



\finexercice



\exercice{1866, roussel, 2001/09/01}

\enonce{001866}{}
Soient $(a_i)_{1\leq i\leq n}$ et $(b_i)_{1\leq i\leq n}$ deux
familles de $n$ nombres r\'eels. Montrer, en \'etudiant le signe
du trin\^ome $\displaystyle \lambda \longrightarrow  \sum_{i=1}^n
(a_i+\lambda b_i)^2$ que $\displaystyle{\sum_{i=1}^n} a_ib_i\leq
\left (\sum_{i=1}^n a_i^2 \right )^{1\over 2} \left (\sum_{i=1}^n
b_i^2 \right )^{1\over 2}.$
\finenonce{001866}



\finexercice



\exercice{6807, gijs, 2011/10/16}
\enonce{006807}{}
Soit $E$ un espace vectoriel normé sur $\Cc$. On
désigne par $\overline B = \{x\in E \mid \Vert x\Vert
\le1\}$ la boule fermée de centre $0$ et de rayon $1$.

Soit $F\subset E$ un sous-espace vectoriel {\bf fermé}
de $E$. Pour $x\in E$ on pose~:
$$
d(x,F) = \inf_{y\in F} \Vert\, x-y\,\Vert \equiv \inf \{\,
\Vert\, x-y\, \Vert \mid y\in F\,\}
\ .
$$
\begin{enumerate}
\item Montrer que pour tout $x\in E$ on a $0\le d(x,F) \le
\Vert x\Vert$.

\item Montrer que les deux conditions suivantes sont
équivalentes. 
  \begin{enumerate}
  \item $d(x,F) = 0$;
  \item $x\in F$.
  \end{enumerate}

\item
  \begin{enumerate}
  \item Montrer que, quels que soient $x,x'\in E$,
$\lambda\in \Cc$, $y\in F$, on a~:

\begin{align*}
%\forall x\in E, \lambda\in \Cc : &
d(\lambda x,F) &= \vert\lambda\vert\,d(x,F)
\\
%\forall x\in E, y\in F : &
d(x-y,F) &= d(x,F)
\\
%\forall x,x'\in E : &
d(x+x',F) &\le d(x,F) + d(x',F)
\end{align*}

  \item Montrer que l'application $x\mapsto d(x,F)$ est
uniformément continue dans $E$.
  \end{enumerate}

\item Soit $x\in \overline B$. On pose $\alpha =
d(x,F)$ et on suppose $\alpha>0$. Soit de plus
$\epsilon>0$.
  \begin{enumerate}
  \item Montrer qu'il existe $y\in F$ tel que 
$$
\alpha \le \Vert\, x-y\, \Vert < \alpha(1+\epsilon)
\ .
$$
  \item Montrer qu'il existe $x' \in \overline B$
tel que
$$
\tfrac{1}{1+\epsilon} = d(x',F)  <1
\ .
$$
  \item Montrer que, si $F\not= E$, 
${ \sup_{x\in \overline B}}\ d(x,F) = 1$.
  \end{enumerate}

\item Sachant que tout espace vectoriel normé de
dimension finie est complet, dé\-mon\-trer que tout
sous-espace vectoriel de dimension finie d'un espace
vectoriel normé est fermé.

\item On suppose maintenant que $\overline B$ est
compact. 
  \begin{enumerate}
  \item Montrer que, pour tout $\epsilon>0$, il existe un nombre
fini $k$ de points de $\overline B$~: $x_1, \dots, x_k \in
\overline B$ tels que
$$
\overline B \subset \bigcup_{j=1}^k B_{\epsilon}(x_j)
\ ,
$$
où $B_\epsilon(x_j)$ désigne la boule ouverte de
centre $x_j$ et de rayon $\epsilon$.

  \item Déduire de ce qui précède que $E$
est de dimension finie (on pourra considérer le
sous-espace vectoriel engendré par les vecteurs $x_1,
\dots, x_k$).
  \end{enumerate}

\item On revient au cas général (c'est-à-dire~:
on ne suppose plus que $\overline B$ est compact). Soit
$u:E \to E$ une application linéaire continue. On
suppose que $\overline{u(\overline B)}$ (l'adhérence de
l'image de $\overline B$ par $u$) est compacte. Soit
$\lambda \not= 0$ une valeur propre de $u$
(c'est-à-dire~: il existe
$x_0
\in E$, $x_0\not= 0$ tel que
$u(x_0) = \lambda\, x_0$). On pose
$V_\lambda = \{\,x\in E \mid u(x) = \lambda x\,\}$.

  \begin{enumerate}
  \item Montrer que $V_\lambda$ est un sous-espace
vectoriel fermé de $E$.
  \item Montrer que $\overline B \cap V_\lambda
\subset \frac1\lambda \, u(\overline B)$.
  \item Montrer que $V_\lambda$ est de dimension
finie.
  \end{enumerate}
\end{enumerate}
\finenonce{006807}



\finexercice

\section{432.00 Théorème de Stone-Weirstrass, théorème d'Ascoli}


\exercice{2408, bodin, 2003/10/01}
\enonce{002408}{} 
Soit $f \in \mathcal{C}([a,b],\Rr)$ telle que
$$\forall n \in \Nn \quad \int_a^b f(t) t^n \, dt = 0.$$
Montrer que $f$ est la fonction nulle.

\finenonce{002408}



\finexercice


\exercice{2409, bodin, 2003/10/01}
\enonce{002409}{} 
Montrer qu'une fonction de $\mathcal{C}(\Rr,\Rr)$ admettant une limite finie en $+\infty$
n'est pas limite uniforme de polynômes de $\Rr[x]$.
\finenonce{002409}



\finexercice


\exercice{2410, bodin, 2003/10/01}
\enonce{002410}{} 
Soit $E$ un espace compact. Soit $f_i$, $i=1,\ldots,n$ une famille de $n$ élements
de $\mathcal{C}(E,\Rr)$ qui sépare les points de $E$. Montrer que $E$ est homéomorphe 
à une partie de $\Rr^n$.
\finenonce{002410}


\finexercice


\exercice{2411, bodin, 2003/10/01}
\enonce{002411}{} 
Soient $X$ et $Y$ deux espaces métriques compacts. Soit $\mathcal{A}$ l'ensembles des combinaisons linéaires finies
$f \in \mathcal{C}(X\times Y,\Rr)$ de la forme :
$$f (x,y) = \sum_{i\in I} \lambda_i u_i(x) \cdot v_i(y), \quad \text {avec } u_i \in \mathcal{C}(X,\Rr), v_i \in \mathcal{C}(Y,\Rr), \lambda_i \in \Rr,  I \text{ fini}.$$

Montrer que toute fonction de $\mathcal{C}(X\times Y,\Rr)$ est limite uniforme de suites d'éléments de $\mathcal{A}$.
\finenonce{002411}


\finexercice


\exercice{2412, mayer, 2003/10/01}
\enonce{002412}{} 
\begin{enumerate}
\item Soit $k>0$ et ${\cal F}$ l'ensemble des fonctions diff\'erentiables 
$f:[a,b]\to \Rr$ telles que $|f'(t)|\leq k$
pour tout $t \in ]a,b[$. Montrer que ${\cal F}$ est une famille \'equicontinue.
\item Si $L>0$ et $f_n: \Rr^n \to \Rr^n $ est une suite d'applications $L$-lipschitziennes
avec $\|f_n(0)\| = \sqrt 2$, alors montrer que l'on peut extraire une sous-suite convergente
de $(f_n)$.
\end{enumerate}
\finenonce{002412}


\finexercice


\exercice{2413, mayer, 2003/10/01}
\enonce{002413}{} 
Soient $E,F$ des espaces norm\'es et $(f_n)$ une suite d'applications  de $E$ dans $F$
\'equicontinue en $a\in E$. Montrer que, si la suite $(f_n(a))$ converge vers $b$,
alors $(f_n(x_n))$ converge \'egalement vers $b$, si $(x_n)$ est une suite de $E$
telle que $\lim_{n\to \infty} x_n =a$.

L'\'equicontinuit\'e est-elle n\'ecessaire ici?
\finenonce{002413}


\finexercice


\exercice{2414, mayer, 2003/10/01}
\enonce{002414}{} 
Soient $E,F$ des espaces norm\'es et $(f_n)$ une suite d'applications \'equi\-continues
de $E$ dans $F$. Montrer que l'ensemble des $x\in E$, pour lesquels $(f_n(x))$
est une suite de Cauchy dans $F$, est un ferm\'e.
\finenonce{002414}


\finexercice


\exercice{2415, mayer, 2003/10/01}
\enonce{002415}{} 
Soient $(E,d)$ un espace m\'etrique et ${\cal H}$ une famille \'equicontinue d'appli\-cations
de $E$ dans $\Rr$. \'Etablir:
\begin{enumerate}
\item L'ensemble $A$ des $x\in E$ pour lesquels ${\cal H} (x)$ est born\'e est ouvert et ferm\'e.
\item Si $E$ est compact et connexe et si ${\cal H} (x_0)$ est born\'e pour un point quelconque $x_0 \in E$,
alors ${\cal H}$ est relativement compact dans ${\cal C} (E, \Rr )$.
\end{enumerate}
\finenonce{002415}


\finexercice


\exercice{2416, mayer, 2003/10/01}
\enonce{002416}{} 
On consid\`ere la suite de fonctions $f_n(t) = \sin (\sqrt{t+ 4(n\pi )^2})$,
$t\in [0, \infty [$.
\begin{enumerate}
\item Montrer qu'il s'agit d'une suite de fonctions \'equicontinues convergent
simplement vers $f\equiv 0$.
\item  La suite $(f_n)$ est elle relativement compacte
dans $({\cal C}_b ([0, \infty [), \|.\|_\infty )$, l'ensemble des fonctions continues et bornées ?
Que dit le th\'eor\`eme d'Ascoli?
\end{enumerate}
\finenonce{002416}


\finexercice

\exercice{2417, mayer, 2003/10/01}
\enonce{002417}{} 
Soit $K: {\cal C} ([a,b]) \to {\cal C} ([a,b])$ donn\'e par $(Kf)(s)=\int _a^b k(s,t) f(t) \, dt$,
$k \in {\cal C} ([a,b]\times [a,b])$, et soit $(f_n)$ une suite born\'ee de $X= ({\cal C} ([a,b]), \|.\|_\infty )$.
\begin{enumerate}
\item Rappeler pourquoi $k$ est uniform\'ement continue.
\item En d\'eduire l'\'equicontinuit\'e de $(Kf_n)$.
\item Montrer que $(Kf_n)$ contient une sous-suite convergente dans $X$.
\end{enumerate}
\finenonce{002417}


\finexercice


\exercice{6252, queffelec, 2011/10/16}
\enonce{006252}{}
Soit $X=[0,1]$, $Y=[-1,1]$, $f_n(x)={1\over n}\sin(n^2x)$; on note $A$
le sous-ensemble de $C(X,Y)$ constitué des $(f_n)$, $n\geq1$.

\begin{enumerate}  
\item Montrer sans calculs que $A$ est équicontinu.

\item Montrer que l'on a plus précisément
$|f_n(x)-f_n(y)|\leq\sqrt2|x-y|^{1\over2}$ pour tous n, x, y. En déduire que le
module d'équicontinuité de $A$ est $\geq \varepsilon^2/2$.
\end{enumerate}
\finenonce{006252}



\finexercice


\exercice{6253, queffelec, 2011/10/16}
\enonce{006253}{}
Pour quelles valeurs de $\alpha\geq0$, la fonction $f(x)=\cos x^\alpha$ est -elle uniformément
continue ?

On suppose cette condition remplie et on définit $f_n$ par $f_n(x)=f(x+n)$. Montrer que l'on peut
extraire de $(f_n)$ une suite convergeant uniformément sur tout compact de $\Rr^+$. Peut-on avoir
convergence uniforme sur tout $\Rr^+$ ?
\finenonce{006253}



\finexercice


\exercice{6254, queffelec, 2011/10/16}
\enonce{006254}{}
Soit $H=\{f\in C^1({\Rr})\ /\ \int_{\Rr}|f(x)|^2 dx + \int_{\Rr}|f'(x)|^2 dx \leq1\}$.

\begin{enumerate}  
\item Montrer que si $f\in H$, $\lim_{|x|\to\infty} f(x)=0$.

\item Montrer que $||f||_\infty\leq {1\over\sqrt2}$, pout toute $f\in H$.

\item Montrer que $H$ est une partie équicontinue et bornée de $C_0({\Rr})$.

\item Soit $\varphi(x)=(x^2-1)^2\ {\bf 1}_{|x|\leq1}$, et $f_n(x)=\varphi(x-n)$; montrer que
$\lambda f_n\in H$ pour une constante $\lambda$ bien choisie et que cette suite n'a aucune valeur
d'adhérence dans $C(\Rr)$. Conclusion ?
\end{enumerate}
\finenonce{006254}



\finexercice

\section{440.00 Fonction holomorphe}


\exercice{2681, matexo1, 2002/02/01}

\enonce{002681}{}
Tout complexe $z$ qui n'est pas un r{\'e}el positif ou nul peut s'{\'e}crire sous la
forme $z = r e^{i \theta}$, avec $r > 0$ et~$\theta\in\left]0,2\pi \right[$. 
On d{\'e}finit une fonction $f$
sur $\Omega = \C\setminus[0,+\infty [$ par
$$ f(z) = P(r, \theta)+iQ(r,\theta) $$
o{\`u} $P$ et $Q$ sont des fonctions r{\'e}elles donn{\'e}es. Donner des conditions sur les
d{\'e}riv{\'e}es partielles de $P$, $Q$ pour que $f$ soit holomorphe sur $\Omega $.

En d{\'e}duire que la fonction
$$ \log z = \log r + i\theta $$
est holomorphe sur $\Omega $. Quelle est sa d{\'e}riv{\'e}e ?
\finenonce{002681}
\finexercice


\exercice{2682, matexo1, 2002/02/01}

\enonce{002682}{}
Soit $a$ un r{\'e}el dans $]0,1[$. On souhaite calculer l'int{\'e}grale
$$ I_a = \int_0^{+\infty } {x^{a-1} \over 1+x}\,dx $$

Pour cela, on d{\'e}finit une fonction $f$ sur
$\Omega =\C\setminus [0,+\infty [$ de la mani{\`e}re
suivante:
$$ f(z) = {\exp\left((a-1)\log z\right) \over 1+z} $$
o{\`u} la fonction $\log$ est d{\'e}finie comme dans l'exercice pr{\'e}c{\'e}dent.
\begin{itemize}
\item Montrer que $f$ est m{\'e}romorphe sur $\Omega $. Quels en sont les p{\^o}les ?

\item Soit $\varepsilon\in\left]0,1\right[$ et $R>1$ deux r{\'e}els.
On consid{\`e}re {\`a} pr{\'e}sent le chemin~$C$ 
r{\'e}union du segment $[\varepsilon, R]+0i$, du
cercle de rayon~$R$, parcouru positivement, du segment
$[R, \varepsilon]-0i$ et du cercle de rayon~$\varepsilon$
parcouru n{\'e}gativement.
Calculer $ \int_C f(z)\,dz $; en d{\'e}duire
la valeur de $I_a$.
\end{itemize}
\finenonce{002682}
\finexercice


\exercice{2783, burnol, 2009/12/15}
\enonce{002783}{}
\label{ex:burnol1.1.1}
Montrer que la fonction $f(z) = \frac1z$ est holomorphe sur
$\Cc\setminus\{0\}$ et vérifie $f'(z) = -\frac1{z^2}$.
\finenonce{002783}
\finexercice


\exercice{2784, burnol, 2009/12/15}
\enonce{002784}{}
\label{ex:burnol1.1.2}
Si $f$ et $g$ sont deux fonctions dérivables au sens
complexe au point $z_0$ ; montrer que $f+g$, $f-g$ et $fg$ le
sont et donner la valeur de leurs dérivées au point
$z_0$. 
\finenonce{002784}
\finexercice


\exercice{2785, burnol, 2009/12/15}
\enonce{002785}{}
\label{ex:burnol1.1.3}
Si $f$ et $g$ sont deux fonctions dérivables au sens
complexe au point $z_0$ montrer que $\frac fg$ est dérivable au sens complexe
et donner la valeur de la dérivée lorsque $g(z_0)\neq 0$.
\finenonce{002785}
\finexercice


\exercice{2786, burnol, 2009/12/15}
\enonce{002786}{}
\label{ex:burnol1.1.4}
Montrer la formule pour la dérivée d'une composition $g\circ f$.
\finenonce{002786}
\finexercice


\exercice{2787, burnol, 2009/12/15}
\enonce{002787}{}
\label{ex:burnol1.1.5}
Soit $f$ et $g$ deux fonctions $n$-fois dérivables au sens
complexe sur un ouvert non vide $U$ (remarque: d'après le
cours il suffit qu'elles soient dérivables une fois sur $U$
pour qu'elles le soient un nombre quelconque de
     fois). Montrer la formule de Leibniz généralisée: 
\[ \forall z\in U\qquad (fg)^{(n)}(z) = \sum_{j=0}^n \binom{n}{j}
f^{(j)}(z)g^{(n-j)}(z)\]
\finenonce{002787}
\finexercice


\exercice{2788, burnol, 2009/12/15}
\enonce{002788}{}
\label{ex:burnol1.1.6}
  On se donne deux séries entières $f(z) = \sum_{n=0}^\infty
  a_n z^n$ et $g(z) = \sum_{n=0}^\infty
  b_n z^n$ de rayons de convergences $R_1$ et $R_2$ non
  nuls. En utilisant le théorème sur les séries doubles
  prouver $f(z)g(z) = \sum_{n=0}^\infty c_n z^n$ pour $|z|<R
  = \min(R_1,R_2)$ avec (formules dites de Cauchy):
\[ \forall n\in\Nn\qquad c_n = \sum_{j=0}^n a_j b_{n-j}\]
Le rayon de convergence de la série $\sum_{n=0}^\infty c_n
     z^n$ est-il toujours égal à 
$\min(R_1,R_2)$ ou peut-il être plus grand?
\finenonce{002788}
\finexercice


\exercice{2789, burnol, 2009/12/15}
\enonce{002789}{}
  Retrouver le résultat de l'exercice précédent (l'exercice \ref{ex:burnol1.1.6}) de manière
  plus indirecte en montrant que les coefficients $c_n =
     \sum_{j=0}^n a_j b_{n-j}$ sont 
  ceux de la série de Taylor à l'origine de la fonction
  holomorphe $k(z)=f(z)g(z)$.
\finenonce{002789}
\finexercice


\exercice{2790, burnol, 2009/12/15}
\enonce{002790}{}
  En quels points la fonction $z\mapsto \overline{z}$
  est-elle dérivable au sens complexe, et/ou holomorphe?
  Même question pour les fonctions $z\mapsto x$ et $z\mapsto y$.
\finenonce{002790}
\finexercice


\exercice{2791, burnol, 2009/12/15}
\enonce{002791}{}
\label{ex:burnol1.1.9}
  Prouver qu'une fonction holomorphe sur un ouvert connexe,
  de dérivée  identiquement nulle, est constante. Et si
  l'ouvert n'est pas connexe?
\finenonce{002791}
\finexercice


\exercice{2792, burnol, 2009/12/15}
\enonce{002792}{}
\label{ex:burnol1.1.10}
  Sur un ouvert connexe $U$ on se donne une fonction
  holomorphe $f$ qui a la propriété de ne prendre que des
  valeurs réelles. En utilisant les équations de
  Cauchy-Riemann, montrer que $f$ est constante.
\finenonce{002792}
\finexercice


\exercice{2793, burnol, 2009/12/15}
\enonce{002793}{}
Cet exercice propose une variante pour développer la
théorie de la fonction exponentielle.\medskip
\begin{enumerate}
  \item On se donne une fonction $f$ qui est $n+1$-fois dérivable au sens
  complexe  sur le disque ouvert $D(0,R)$ (on
  sait qu'une fois suffit mais on ne va pas utiliser ce
  théorème difficile ici). Soit
  $z\in D(0,R)$. En appliquant la formule de Taylor avec
  reste intégral de Lagrange à la fonction de la variable
  réelle $t\mapsto g(t) = f(tz)$ pour $0\leq t\leq 1$,
  prouver:
\[ f(z) = f(0) + f'(0) z + \frac{f^{(2)}(0)}2 z^2 + \dots +
  \frac{f^{(n)}(0)}{n!} z^n + z^{n+1}\int_0^1
  \frac{(1-t)^{n}}{n!} f^{(n+1)}(tz)\,dt\]  
  \item On suppose que $f$ est dérivable au sens complexe une fois sur
  $D(0,R)$ et vérifie $f'=f$ et $f(0) = 1$. Montrer que $f$
  est infiniment dérivable au sens complexe. En utilisant
  la question précédente montrer :
\[ \left| f(z) - \sum_{k=0}^n \frac{z^k}{k!}\right| \leq
   (\sup_{|w|\leq |z|} |f(w)|)
   \frac{|z|^{n+1}}{(n+1)!}\]
et en déduire que, pour tout $z\in \Cc$ on a : $f(z) = \sum_{k=0}^\infty \frac{z^k}{k!}$.
  \item Réciproquement on considère la fonction $F(z) =
  \sum_{k=0}^\infty \frac{z^k}{k!}$. Vérifier que le rayon
  de convergence est infini. Établir par un calcul direct
  que $F'(0)$ existe et vaut $1$. En utilisant le théorème sur
  les séries doubles, montrer $F(z+w) = F(z)F(w)$. En
  déduire ensuite que $F$ est holomorphe sur $\Cc$ et vérifie
  $F' = F$.
\end{enumerate}
\finenonce{002793}
\finexercice


\exercice{2794, burnol, 2009/12/15}
\enonce{002794}{}
  Soit $\sum_{n=0}^\infty a_n z^n$ une série entière de
   rayon de convergence $R$. Est-il exact que pour $|z|>R$
   on a $\lim |a_nz^n| = +\infty$?
\finenonce{002794}
\finexercice


\exercice{2795, burnol, 2009/12/15}
\enonce{002795}{}
  Déterminer les séries de Taylor à l'origine de
$\frac1{1-z}$, $\frac1{(1-z)^2}$, $\frac1{(1-z)^3}$,
$\frac1{(1-z)^4}$. 
\finenonce{002795}
\finexercice


\exercice{2796, burnol, 2009/12/15}
\enonce{002796}{}
Déterminer en tout $z_0\neq1$ la série de Taylor et son rayon
de convergence pour la fonction analytique $\frac1{z-1}$.
\finenonce{002796}
\finexercice


\exercice{2797, burnol, 2009/12/15}
\enonce{002797}{}
Déterminer en  tout $z_0\neq 1,\;2$ la série de Taylor et son rayon
de convergence pour la fonction analytique
$\frac1{(z-1)(z-2)}$. On aura intérêt à réduire en éléments
simples. De plus on demande d'indiquer le rayon de
convergence \emph{avant} de déterminer explicitement la série
de Taylor.
\finenonce{002797}
\finexercice


\exercice{2798, burnol, 2009/12/15}
\enonce{002798}{}
Déterminer en tout point $z_0$ où elle est définie la série de
Taylor de la fonction $\frac1{z^3-1}$. On déterminera son
rayon de convergence en fonction de $z_0$.
\finenonce{002798}
\finexercice


\exercice{2799, burnol, 2009/12/15}
\enonce{002799}{}
  On considère la série entière $\sum_{k=0}^\infty
z^{2^k}$. Quel est son rayon de convergence? On note $f(z)$
sa somme. Que vaut $\lim_{t\to1} f(t)$? (on prend $0<t<1$; minorer $f$ par
ses sommes partielles). Plus généralement que vaut $\lim_{t\to1}
f(tw)$ (ici
   encore $t$ est pris dans $]0,1[$), lorsque $w$ vérifie
   une équation $w^{2^N} = 1$?
En déduire qu'il est impossible de trouver un ouvert $U$
   connexe intersectant $D(0,1)$ mais non inclus entièrement
   dans $D(0,1)$ et une fonction holomorphe
   $g(z)$ sur $U$ tels que $g = f$ sur $U\cap D(0,1)$. Pour
   tout $z_0\in D(0,1)$ déterminer alors le rayon de
convergence de la série de Taylor de $f$ au point $z_0$.
\finenonce{002799}
\finexercice


\exercice{2800, burnol, 2009/12/15}
\enonce{002800}{}
Montrer que le rayon de convergence de chacune des séries
   concernées est $1$ et prouver:  
\begin{enumerate}
\item $\sum_{n=1}^\infty nz^n$ ne converge en aucun
  point du cercle $|z|=1$.
\item $\sum_{n=1}^\infty \frac{z^n}{n^2}$ converge en tout point du cercle
  $|z|=1$.
\item $\sum_{n=1}^\infty \frac{z^n}{n}$ converge en tout point du cercle
  $|z|=1$ \textbf{sauf} en $z=1$.
\end{enumerate}
Pour ce dernier cas on définit $S_0 = 1$, $S_1 = 1+z$,
   $S_2 = 1+ z+z^2$, \dots 
(on pose aussi $S_{-1} = 0$). En écrivant $z^n = S_n -
   S_{n-1}$  exprimer $\sum_{n=1}^N \frac{z^n}{n}$ en
   fonction des $S_n$. Montrer que les $S_n$ sont bornées
   lorsque $|z|=1$, $z\neq 1$. Conclure.
\finenonce{002800}
\finexercice


\exercice{2801, burnol, 2009/12/15}
\enonce{002801}{}
Montrer qu'un entier
$k\geq1$ s'écrit de manière unique sous la forme
$2^n(2m+1)$, $n\geq0$, $m\geq0$. Puis prouver pour $|z|<1$
:
\[ \frac z{1-z^2} + \frac {z^2}{1-z^4} + \dots +
\frac{z^{2^n}}{1 - z^{2^{n+1}}} + \dots = \frac z{1-z}
\;. \]
On justifiera les interversions de séries. Prouver aussi:
\[ \frac z{1+z} + \frac {2z^2}{1+z^2} + \dots +
\frac{2^n z^{2^n}}{1 + z^{2^{n}}} + \dots = \frac z{1-z} \;.
\]
\finenonce{002801}
\finexercice


\exercice{2802, burnol, 2009/12/15}
\enonce{002802}{}
Lorsque $z$ est complexe les fonctions $\sin(z)$, $\cos(z)$, $\sh(z)$ et
   $\ch(z)$ sont définies par les formules:
\begin{align*}
\sin(z) &= \frac{e^{iz} - e^{-iz}}{2i}&\qquad \sh(z) &=
   \frac{e^{z} - e^{-z}}{2}  \\
 \cos(z) &=
   \frac{e^{iz} + e^{-iz}}{2}&\qquad \ch(z) &=
   \frac{e^{z} + e^{-z}}{2}
\end{align*}
\begin{enumerate}
  \item Montrer que $\cos$ et $\ch$ sont des fonctions paires et
   $\sin$ et $\sh$ des fonctions impaires et donner leurs
   représentations comme
   séries entières. Prouver $e^{iz} = \cos(z) + i \sin(z)$,
   $\sin(iz) = i\sh(z)$, $\cos(iz) = \ch(z)$, $\sh(iz) =
   i\sin(z)$, $\ch(iz) = \cos(z)$.
  \item Établir les formules:
\[ \cos(z+w) = \cos(z)\cos(w) - \sin(z)\sin(w)\]
\[ \sin(z+w) = \sin(z)\cos(w) + \cos(z)\sin(w)\]
en écrivant de deux manières différentes $e^{\pm
   i(z+w)}$. Donner une autre preuve en utilisant le
   principe du prolongement analytique et la validité
   (admise) des formules pour $z$ et $w$ réels. 
  \item   Prouver pour tout $z$ complexe $\cos(\pi + z) = -\cos(z)$,
  $\sin(\pi+ z) =
     -\sin(z)$. Prouver $\cos(\frac\pi2 - z) = \sin(z)$.
  \item   Prouver les formules $\cos^2 z + \sin^2 z = 1$ et $\ch^2 z
   - \sh^2 z = 1$ pour tout $z\in\Cc$.
\end{enumerate}
\finenonce{002802}
\finexercice


\exercice{2803, burnol, 2009/12/15}
\enonce{002803}{}
  Montrer $\sin(a+ib) =
   \sin(a)\ch(b) + i\cos(a)\sh(b)$. Puis en prenant
   dorénavant $a$ et $b$ réels, prouver:
\[ a,b\in\Rr\implies\quad |\sin(a+ib)|^2 = \sin^2(a) + \sh^2(b)\]
Déterminer alors les nombres complexes $z = a+ib$ tels que
   $\sin(z) = 0$. Donner une autre preuve.
\finenonce{002803}
\finexercice


\exercice{2804, burnol, 2009/12/15}
\enonce{002804}{}
  Montrer :
\[ a,b\in\Rr\implies\quad |\cos(a+ib)|^2 = \cos^2(a) +
   \sh^2(b) = \ch^2(b) - \sin^2(a) \]
Déterminer les nombres complexes $z$ avec $\cos(z) = 0$.
\finenonce{002804}
\finexercice


\exercice{2805, burnol, 2009/12/15}
\enonce{002805}{}
Les fonctions de Bessel sont très importantes en
     Analyse. Elles apparaissent très souvent dans des
     problèmes de physique mathématique. L'analyse complexe
     permet d'étudier de manière approfondie ces
     fonctions. Ici nous nous contentons des tout débuts de
     la théorie. Nous ne considérons que les
     fonctions\footnote{dites ``fonctions de Bessel
     de
     première espèce (et d'indices entiers)''.} $J_0$, $J_1$,
     $J_2$, \dots, qui sont définies par les formules:
\footnote{Autrement dit:
\[ J_\nu(z) = \frac{z^{\nu}}{2.4.\dots.(2\nu)}\left(1 -
     \frac{z^2}{2.(2\nu+2)} 
     + \frac{z^4}{2.4.(2\nu+2).(2\nu+4)}
     - \frac{z^6}{2.4.6.(2\nu+2).(2\nu+4).(2\nu+6)} +\dots\right)\]
Remarquez que seule la constante
     $2.4.\dots.(2\nu) = 2^\nu \nu!$ nous restreint (pour le
     moment) à des valeurs entières de $\nu$. Si on en fait
     abstraction on obtient
     avec $\nu = -\frac12$ la fonction ``multiforme''
     $z^{-1/2}\cos(z)$ ; tandis qu'avec
     $\nu=+\frac12$ on obtient $z^{-1/2}\sin(z)$. Les
     définitions exactes sont $J_{-1/2}(z) =
     \sqrt{\frac2{\pi z}}\cos(z)$ et $J_{1/2}(z) =
     \sqrt{\frac2{\pi z}}\sin(z)$.}
\[ \nu\in\Nn,z\in\Cc\qquad J_\nu(z) = \sum_{n=0}^\infty 
     (-1)^n\frac{(\frac z2)^{2n+\nu}}{n!(n+\nu)!}\]
\medskip
\begin{enumerate}
\item   Montrer que le rayon de convergence de la série
     définissant $J_\nu$ est $+\infty$.
 \item
  En dérivant terme à terme prouver les formules:
\begin{align*}
 (z^\nu J_\nu)' &= z^\nu J_{\nu-1}\qquad\quad\ (\nu\geq1)\\
 (z^{-\nu} J_\nu)' &= - z^{-\nu} J_{\nu+1}\qquad(\nu\geq0)
\end{align*}
En particulier on a $(zJ_1)' = zJ_0$ et $J_0^{\,'} = - J_1$.
\item  Réécrire les équations précédentes sous la forme 
$(z\frac d{dz} + \nu) J_\nu = z J_{\nu-1}$
($\nu\geq1$) et $(z\frac d{dz} - \nu) J_\nu = - z
     J_{\nu+1}$ ($\nu\geq0$)
et en déduire $-(\frac d{dz} + \frac{\nu + 1}z)(\frac d{dz}
     - \frac\nu z)J_\nu = J_\nu$, puis, après simplification,
     l'équation différentielle de Bessel:
\[ z^2 J_\nu'' + zJ_\nu' + (z^2 - \nu^2)J_\nu = 0\]
  \item Montrer, pour tout $\nu\in\Nn$, que la série
     entière définissant $J_\nu$ est la seule (à une
     constante multiplicative près) qui donne une solution
     de l'équation différentielle de Bessel.\footnote{les
     autres solutions de l'équation différentielle sont
     singulières en $z=0$, avec une composante logarithmique
     ($\nu\in\Zz$). Pour $\nu\notin\Zz$ il y a une solution
     en $z^\nu(\sum_{k\geq0} c_k z^k)$ et une autre en
     $z^{-\nu}(\sum_{k\geq0} d_k z^k)$.}
\end{enumerate}
\finenonce{002805}
\finexercice

\exercice{6580, hueb, 2011/10/16}
\enonce{006580}{}
\begin{enumerate}
\item Montrer que $u(x,y) = e^{-x}(x \sin y - y \cos y)$ est
  harmonique.
\item Déterminer $v$ telle que $f = u +iv$ soit holomorphe.
\item Ecrire la fonction $f$ trouvée ci-dessus comme fonction d'une
  variable complexe.
\end{enumerate}
\finenonce{006580}



\finexercice


\exercice{6581, hueb, 2011/10/16}
\enonce{006581}{} Montrer que
$$f(x+iy) = x^2 + i y^3$$
n'est holomorphe en aucun point bien que les
equations de Cauchy-Riemann soient vérifiées à l'origine, même
sur une parabole que l'on précisera.  \finenonce{006581}



\finexercice


\exercice{6582, hueb, 2011/10/16}
\enonce{006582}{} Etude de l'exponentielle
complexe $f(z) = e^z$ et du logarithme complexe.
\begin{enumerate}
\item Décrire l'image d'une droite $y =c$, $c$ étant une
  constante, par rapport à $f$.
\item Décrire l'image d'une droite $x =c$, $c$ étant une
  constante, par rapport à $f$.
\item Vérifier que la restriction de $f$ au domaine
  $$W = \{ z = x+iy; |y| < \pi\}$$
  est une bijection de $W$ sur
  $$D=\Bbb C \setminus \{z; z = -x,\ x \in \Bbb R,\ x \geq 0\}.$$
\item En déduire l'existence d'une fonction complexe unique $\Phi$,
  avec domaine de définition $D_{\Phi} = D$, de sorte que
  $$
  e^{\Phi (z)} = z, \quad |\textrm{Im} \Phi (z)| < \pi.  $$
  Cette
  fonction est appelée {\it détermination principale\/} du
  logarithme, notée Log; en utilisant un peu plus de théorie on
  montre qu'elle est holomorphe, avec $\textrm{Log}'(z) = \frac 1z$.
\end{enumerate}
\finenonce{006582}



\finexercice


\exercice{6583, hueb, 2011/10/16}
\enonce{006583}{} Vérifier que, pour tout
nombre complexe $z$ avec $|z| <1$, on a
$$\textrm{Log}(1+z) = \sum_1^{\infty} (-1)^{n-1}\frac {z^n}n =z -
\frac{z^2}2 + \frac{z^3}3 - \dots.$$
\finenonce{006583}



\finexercice


\exercice{6584, hueb, 2011/10/16}
\enonce{006584}{} Développer en série
entière $\sin z$ et $\cos z$.  \finenonce{006584}



\finexercice


\exercice{6585, hueb, 2011/10/16}
\enonce{006585}{} Vérifier que
$$\textrm {Arctg} w = w - \frac {w^3}3 + \frac {w^5}5 - \frac {w^7}7 +
\dots$$
dans un domaine de convergence que l'on précisera.  \finenonce{006585}



\finexercice


\exercice{6586, hueb, 2011/10/16}
\enonce{006586}{} Evaluer
$$\int_{1+i}^{2+4i} z^2 dz$$
\begin{enumerate}
\item le long de la parabole $y = x^2,\ 1 \leq x \leq 2$;
\item le long du segment de droite $(1+i), (2+4i)$;
\item le long des segments $(1+i), (2+i)$ et $(2+i),(2+4i)$.
\end{enumerate}
\finenonce{006586}



\finexercice


\exercice{6620, queffelec, 2011/10/16}
\enonce{006620}{}
Pour tout $z\in\Cc$, on définit $\exp(z)=\sum_{n=0}^\infty {1\over n!}z^n$.
\begin{enumerate}
\item Montrer que la fonction $\exp$ est continue et vérifie :

$\exp(0)=1;\quad \forall a,b\in{\Cc},\ \exp(a+b)=\exp(a).\exp(b).
$

\item On note $e=\exp(1)$ et $f(z)=\exp(z)=e^z$. Etablir les résultats suivants :
  \begin{enumerate}
  \item $\forall z\in{\Cc},\ e^z\not=0$.
  \item $f'(z)=f(z)$.
  \item $f_{|\Rr}$ est croissante, positive, tend vers $+\infty$ quand $x\to+\infty$,
tend vers $0$ quand $x\to-\infty$.
  \item Il existe un nombre  positif $\pi$ tel que
$e^{i{\pi\over2}}=i$ et tel que $e^z=1$ ssi $z\in2i\pi\Zz$.
  \item $f$ est périodique de période $2i\pi$.
  \item  $t\to e^{it}$ est une surjection de $\Rr$ sur le cercle unité.
  \item  $f({\Cc})= {\Cc}^*$.
  \end{enumerate}
\end{enumerate}
\finenonce{006620}



\finexercice


\exercice{6621, queffelec, 2011/10/16}
\enonce{006621}{}
Trouver le domaine maximal de convergence des séries entières suivantes :
$$(1)\ \sum_{n=1}^\infty {z^n\over n},\quad (2)\ \sum_{n=0}^\infty
(2z)^n\cos n\theta,\quad  (3)\ \sum_{n=0}^\infty
z^{n^2}.$$
\finenonce{006621}



\finexercice


\exercice{6622, queffelec, 2011/10/16}
\enonce{006622}{}
Développer en série entière les fonctions suivantes et préciser le domaine
maximal de convergence ($a,b\not=0$):
$$
f(z)={1\over 1+z+z^2},\quad g(z)={1\over (a-z)(b-z)},\quad 
h(z)={z\sin a\over z^2-2(\cos a)z+1}.$$
\finenonce{006622}



\finexercice


\exercice{6623, queffelec, 2011/10/16}
\enonce{006623}{}
Développer en série entière sur $\Rr$ : $f(x)=e^{-{x^2\over2}}\displaystyle
\int_0^x e^{{t^2\over2}}\ dt$.
\finenonce{006623}



\finexercice


\exercice{6624, queffelec, 2011/10/16}
\enonce{006624}{}
Soit $f:[-1,1]\to \Rr$ définie par $f(x)=\sum_{n\geq2}{(-1)^n\over x+n}$.
Montrer que $f$ est de classe $C^\infty$ et développable en série entière.
\finenonce{006624}



\finexercice


\exercice{6625, queffelec, 2011/10/16}
\enonce{006625}{}
Trouver les solutions de l'équation
différentielle :$(1+x^2)y''-2y=0$ en commen\c cant par chercher des solutions
développables en série entière.
\finenonce{006625}



\finexercice


\exercice{6626, queffelec, 2011/10/16}
\enonce{006626}{}
Soit $\sum a_nz^n$ une série entière de rayon $1$ .
Montrer les équivalences (i)$\Longleftrightarrow$(ii)

$\Longleftrightarrow$(iii)
où

(i) La série converge uniformément sur $D$.

(ii) La série converge uniformément sur $\overline D$.

(iii) La série converge uniformément sur $\partial D$

(\emph{Indication} : pour l'implication (iii)$\Longrightarrow$(ii), on posera
$r_N=\sum_N^\infty a_n e^{in\theta}$ et on fera une tranformation d'Abel dans
la somme $\sum_M^N a_n \rho^n e^{in\theta}$ où $0\leq\rho\leq1$.)
\finenonce{006626}



\finexercice


\exercice{6627, queffelec, 2011/10/16}
\enonce{006627}{}
Soit $\sum a_nz^n$ une série entière de rayon $1$. On pose
$s_n=a_0+a_1+\cdots+a_n$, $t_n={1\over n+1}(s_0+s_1+\cdots+s_n)$, $u(z)=\sum
s_nz^n$ et $v(z)=\sum t_nz^n$.
\begin{enumerate}
\item Montrer que les rayons de convergence de $u$ et de $v$ sont égaux à $1$.

\item Etablir pour tout $|z|<1$, $u(z)={1\over 1-z}\sum a_nz^n$. Retrouver ainsi le
théorème d'Abel : soit $f(z)=\sum a_nz^n$ une série entière de rayon $1$, telle
que
$\sum a_n$ converge vers $A$. Alors $f(z)$ tend vers $A$ quand $z\to 1$ non
tangentiellement.
\end{enumerate}
\finenonce{006627}



\finexercice


\exercice{6628, queffelec, 2011/10/16}
\enonce{006628}{} 
Soit $(u_n)_{n\ge 0}$ une suite de réels. 
On dit que $a\in \overline{\Rr}$ est une valeur d'adhérence de la suite
$(u_n)_{n\ge 0}$ s'il existe une sous-suite de la suite $(u_n)_{n\ge 0}$
qui tend vers $a$. Montrer que $L=\limsup_{n\to\infty}u_n=\inf_n\sup_{k\ge n}u_k$ 
et $l=\liminf_{n\to\infty}u_n=\sup_n\inf_{k\ge n}u_k$ sont des valeurs 
d'adhérence de la suite $(u_n)_{n\ge 0}$. Vérifier que ce sont 
respectivement la plus grande et la plus petite des valeurs d'adhérence.
\finenonce{006628}



\finexercice       


\exercice{6629, queffelec, 2011/10/16}
\enonce{006629}{} 
Soit $(u_n)_{n\ge 0}$ et $(v_n)_{n\ge 0}$ deux suites de réels.
Etablir les propositions suivantes :
$$\limsup_{n\to\infty} \left( u_n+v_n\right)\le\limsup_{n\to\infty}u_n+
\limsup_{n\to\infty}v_n\ ;$$
$$\left(\forall n\ge n_0, u_n\le v_n\right)\Rightarrow\limsup_{n\to\infty}u_n\le
\limsup_{n\to\infty}v_n.$$
Donner l'exemple de deux suites pour lesquelles la première inégalité
est stricte.
\finenonce{006629}



\finexercice       


\exercice{6630, queffelec, 2011/10/16}
\enonce{006630}{} 
Soit $\sum_{n\ge 0}a_nz^n$ une série
entière de rayon de convergence $R$. On suppose que $\lim_{n\to\infty}
\left\vert a_{n+1}/a_n\right\vert$ existe dans $\overline{\Rr}$.
Montrer que 
$$\frac{1}{R} =\lim_{n\to\infty}\left\vert\frac{a_{n+1}}{a_n}\right\vert.$$
\finenonce{006630}



\finexercice       


\exercice{6631, queffelec, 2011/10/16}
\enonce{006631}{} 
Déterminer le rayon de convergence des séries
$\sum_{n\ge 0}a^nz^n$ ($a\in \C$) ; $\sum_{n\ge 0}e^{\sqrt n}z^n$ ; 
$\sum_{n\ge 0}z^n/n!$ ; $\sum_{n\ge 0}z^{2n}$ ; 
$\sum_{n\ge 0}z^{n!}$.

Comparer le rayon de convergence des séries $\sum_{n\ge 0}c_nz^n$ et
$\sum_{n\ge 0}c_nz^{2n}$.

Si le rayon de convergence de $\sum_{n\ge 0}c_nz^n$ est strictement positif
et fini, quel est celui de $\sum_{n\ge 0}c_nz^{n^2}$ ?
\finenonce{006631}



\finexercice       


\exercice{6632, queffelec, 2011/10/16}
\enonce{006632}{} 
Montrer que le rayon de convergence de la série 
$\sum_{n\ge 0}\left( a_n+b_n\right) z^n$ est supérieur ou égal au
plus petit des rayons de convergence des séries $\sum_{n\ge 0}a_nz^n$ 
et $\sum_{n\ge 0}b_nz^n$. Montrer par un exemple que l'inégalité peut être
stricte.
\finenonce{006632}



\finexercice       


\exercice{6633, queffelec, 2011/10/16}
\enonce{006633}{} 
On considère une série entière 
$\sum_{n\ge 0}c_nz^n$ de rayon de convergence $R>0$. Sa somme est
notée $f(z)$.
\begin{enumerate}
\item Montrer que la série $\sum_{n\ge
0}\frac{c_n}{n!}z^n$ a un rayon de convergence infini. Sa somme notée
$F(z)$ est appelée transformée de Borel.

\item Soit $r$ un réel vérifiant $0<r<R$. Montrer qu'il existe un polynôme
$P$ tel que
$$\forall z\in \C,\sup_{N\ge0}\left\vert \sum_{n=0}^N\frac{c_n}{n!}
z^n\right\vert \le P\left( \vert z\vert\right) +\exp \left( \frac{\vert z\vert}{r}\right).$$
On pourra considérer un entier $n_0$ tel que, pour tout $n>n_0$, on ait
$\vert c_n\vert \le r^{-n}$.

\item Montrer que, pour tout $z$ de $\C$ tel que $\vert z\vert <R$,
on a
$$f(z)=\int_0^{+\infty}F(tz)e^{-t}\, dt.$$
\end{enumerate}
\finenonce{006633}



\finexercice


\exercice{6634, queffelec, 2011/10/16}
\enonce{006634}{}
\begin{enumerate}
\item Montrer que $L=\limsup u_n\in{\Rr}$ est caractérisée par la
condition        suivante        :

Pour tout $\varepsilon>0$, tous les $u_n$ sauf un nombre fini d'entre eux sont
$\leq L+\varepsilon$, et une infinité d'entre eux est $\geq L-\varepsilon$.

\item Ecire l'analogue pour la $\liminf$.
\end{enumerate}
\finenonce{006634}



\finexercice


\exercice{6635, queffelec, 2011/10/16}
\enonce{006635}{}
\begin{enumerate}
\item Déterminer à l'aide de la formule de Hadamard le rayon de convergence des
séries suivantes :

$$(1)\ \sum_{n=0}^\infty e^nz^{n^2}\quad (2)\ \sum_{n=0}^\infty
n!z^{n^2}\quad  (3)\ \sum_{n=1}^\infty \Big( 1+{(-1)^n\over n}\Big)^{n^2}
z^n.$$

\item Trouver le rayon de convergence de la série  $\sum a_nz^n$ où
$$a_n=q^{n^2} (|q|<1);\quad a_n=C_{kn}^n;\quad a_n=e^{n^\alpha} (0<\alpha<1).$$
\end{enumerate}
\finenonce{006635}



\finexercice


\exercice{6636, queffelec, 2011/10/16}
\enonce{006636}{}
\begin{enumerate}
\item Trouver le rayon de convergence de la série  $\sum a_nz^n$ où
$a_{2n+1}=a^{2n+1}$ et $a_{2n}=b^{2n}$ avec $0<a,b<1$.

\item Pour la série $\sum a_nz^n$, où
$a_{2n}=a^{n}b^{n}$ et $a_{2n+1}=a^{n}b^{n+1}$ avec $a,b>0$, comparer l'inverse
du rayon,
$R^{-1}$, avec
$\limsup {a_{n+1}\over a_n}$.
\end{enumerate}
\finenonce{006636}



\finexercice


\exercice{6637, queffelec, 2011/10/16}
\enonce{006637}{}
Montrer que si la série $\sum a_nz^n$ a un rayon $0<R<+\infty$, la série
lacunaire $\sum a_nz^{\lambda(n)}$, où $\lim{\lambda(n)\over n}=+\infty$, a
comme rayon $R'=1$. Montrer sur un exemple que la réciproque peut être fausse.
\finenonce{006637}



\finexercice


\exercice{6638, queffelec, 2011/10/16}
\enonce{006638}{}
\begin{enumerate}
\item Soit $f=P+iQ$ une fonction holomorphe dans un ouvert connexe non vide $\Omega$
de
$\Cc$. Montrer que les assertions suivantes sont équivalentes :
 
(i) $f$ est constante,\quad (ii)  $P$ est constante, \quad (iii)  $Q$ est
constante,\quad (iv)  $\bar f$ est holomorphe dans $U$, \quad (v)  $|f|$ est
constant.

\item Soit $f,g\in H(U)$. On suppose que $g$ ne s'annule pas dans $U$ et
$f(z)\bar g(z)\in\Rr$ pour $z\in U$. Montrer qu'il existe $c\in\Rr$ telle
que $f=cg$.
\end{enumerate}
\finenonce{006638}



\finexercice


\exercice{6639, queffelec, 2011/10/16}
\enonce{006639}{}
\begin{enumerate}
\item Pour $z=x+iy,\ x,y\in \Rr$, on pose $f(z)=x+iy^2$. Montrer que $f$ est 
$\Rr$-différentiable sur $\Cc$ et calculer sa différentielle. Existe-t-il un
ouvert $U$ de $\Cc$ telle que $f_{|U}\in H(U)$ ? 

\item Même question avec $f(z)=|\sin z|$.
\end{enumerate}
\finenonce{006639}



\finexercice


\exercice{6640, queffelec, 2011/10/16}
\enonce{006640}{}
\begin{enumerate}
\item Soit $U=\{z=x+iy\in{\Cc}\ ;\ -\pi<x<\pi,\ y\in \Rr\}$. Soit
$P(x,y)={\sin x\over \cos x+\cosh  y}$ pour $z\in U$. Montrer qu'il existe
$f\in H(U)$ unique telle que $f(0)=0$ et $P=\Re f$.

\item Soit $a,b,c\in\Rr$. On pose $P(x,y)=ax^2+2bxy+cy^2$ pour $x,y\in\Rr$.

Donner une CNS pour qu'il existe $f\in H({\Cc})$ telle que 
$P=\Re f$. Sous cette condition trouver alors toutes les applications $f\in
H({\Cc})$ telles que $P=\Re f$.
\end{enumerate}
\finenonce{006640}



\finexercice


\exercice{6641, queffelec, 2011/10/16}
\enonce{006641}{}
\begin{enumerate}
\item Montrer que les équations de Cauchy-Riemann en polaires s'écrivent
$\frac{\partial f}{\partial r}+\frac{i}{r}\frac{\partial f}{\partial\theta}=0$.
\item Écrire les opérateurs $\frac{\partial}{\partial z}$ et $\frac{\partial}{\partial \bar z}$ en coordonnées polaires.
\item Si $f=u+iv$ est holomorphe sur $U$, montrer qu'en tout $z=re^{i\theta}$ de $U$, on a  $f'(z) = \frac{e^{-i\theta}}{r}\left(\frac{\partial v}{\partial \theta}(z) -i\frac{\partial u}{\partial \theta}(z)\right)$.
\item Soit $P$ un polyn\^ome non constant, supposé sans zéros. On pose
alors $I(r)=\int_0^{2\pi}{d\theta\over P(re^{i\theta})}$. Montrer que $I'(r)=0$.
En calculant $\lim_{r\to\infty}I(r)$ et $I(0)$, aboutir à une contradiction.
Conclusion ?
\end{enumerate}
\finenonce{006641}



\finexercice


\exercice{6642, queffelec, 2011/10/16}
\enonce{006642}{}
 En quels points la fonction $f(z)=z\Re z$ est-elle
$\C$-différentiable ? Même question pour $f(z) =\exp \overline z$.
\finenonce{006642}



\finexercice       


\exercice{6643, queffelec, 2011/10/16}
\enonce{006643}{} 
Ecrire les équations de Cauchy-Riemann en coordonnées polaires. En déduire
que chaque détermination du logarithme est holomorphe.
\finenonce{006643}



\finexercice       


\exercice{6644, queffelec, 2011/10/16}
\enonce{006644}{} 
Soit $f$ la fonction définie sur $\C$ par
$$\left\{ \begin{array}{ll}
f(z)=e^{-1/z^4}&\mbox{ si }z\ne 0\\
f(0)=0&\\
\end{array}\right.$$
Montrer que $f$ vérifie les équations de Cauchy-Riemann en tout point de
$\C$ mais n'est pas holomorphe dans $\C$.
\finenonce{006644}



\finexercice       


\exercice{6645, queffelec, 2011/10/16}
\enonce{006645}{} Soit $\Omega$ un ouvert connexe de $\C$ et $f$ une fonction
holomorphe sur $\Omega$. On désigne par $P$ et $Q$ respectivement ses
parties réelle et imaginaire. On suppose qu'il existe des constantes réelles
non toutes nulles $a$, $b$ et $c$ telles que la fonction $aP+bQ+c$ soit
identiquement nulle sur $\Omega$. Montrer que $f$ est constante sur
$\Omega$.
\finenonce{006645}



\finexercice       


\exercice{6646, queffelec, 2011/10/16}
\enonce{006646}{} Soit $f$ une fonction holomorphe sur un ouvert $\Omega $, $u$
sa partie réelle et $v$ sa partie imaginaire. On suppose que les dérivées
partielles secondes de $u$ et $v$ existent et sont continues sur $\Omega $.
Montrer que $u$ (resp. $v$) est  harmonique (c'est-à-dire $\displaystyle
{\partial ^2u\over \partial x^2}+ {\partial ^2u\over \partial y^2}=0$).
\finenonce{006646}



\finexercice       


\exercice{6647, queffelec, 2011/10/16}
\enonce{006647}{} On dit que deux fonctions réelles $u(x,y)$ et $v(x,y)$ sont conjuguées 
har\-mo\-ni\-ques si elles vérifient les équations de Cauchy-Riemann. 
\begin{enumerate}
\item Montrer que si $u$ et $v$ sont conjuguées harmoniques, alors $u$ et $v$
sont harmoniques.

\item Trouver
les conjuguées har\-mo\-ni\-ques des fonctions har\-mo\-ni\-ques suivantes
dans les ouverts indiqués :
\begin{enumerate}
\item $u(x,y)=x^2-y^2+x$ sur $\C$
\item $u(x,y)={x\over x^2+y^2}$ sur $\C\setminus \{0\}$
\item $u(x,y)={1\over 2}\ln{(x^2+y^2)}$ sur
\begin{enumerate}
\item $\C\setminus \{x+iy\vert\ y=0,x\le 0\}$
\item $\C\setminus \{0\}$
\end{enumerate}
\end{enumerate}
\end{enumerate}
\finenonce{006647}



\finexercice       


\exercice{6648, queffelec, 2011/10/16}
\enonce{006648}{} 
Trouver des domaines de définition $\Omega $ (le plus grand
possible) et des fonctions holomorphes $f=u+iv$ sur $\Omega $  étant donné
la partie réelle $u$ ou la partie imaginaire $v$. 
\begin{enumerate}
\item $u=x^2-y^2+5x+y-{y\over x^2+y^2}$
\item $u=e^x(x\cos y-y\sin y)+2\sin x\sh y+x^3-3xy^2+y$
\item $v=\ln{(x^2+y^2)}+x-2y$
\end{enumerate}
\finenonce{006648}



\finexercice       


\exercice{6649, queffelec, 2011/10/16}
\enonce{006649}{} 
Soit $f(z)=u+iv$ une fonction holomorphe dans un ouvert connexe
$\Omega $. Montrer que les familles de courbes $u(x,y)=c_1$ et
$v(x,y)=c_2$ sont or\-tho\-go\-na\-les ; plus précisément, montrer qu'en
tout point d'intersection $z_0=x_0+iy_0$ de deux de ces courbes tel que
$f'(z_0)\ne 0$, leurs tangentes respectives sont perpendiculaires.
\finenonce{006649}



\finexercice       


\exercice{6650, queffelec, 2011/10/16}
\enonce{006650}{} 
Montrer que si $f(z)$ est holomorphe dans un ouvert connexe 
$\Omega $ et si $f'(z)\ne 0$ en tout point de $\Omega $, alors la 
transformation $w=f(z)$ conserve les angles.
\finenonce{006650}



\finexercice      


\exercice{6651, queffelec, 2011/10/16}
\enonce{006651}{} 
\begin{enumerate}
\item Montrer les inégalités suivantes, pour tout $z\in \C$,
$$\vert e^z-1\vert\le\vert z\vert e^{\vert z\vert}\mbox{ et }
\vert e^z-1-z\vert\le{\vert z\vert^2\over 2}e^{\vert z\vert}$$

\item Soit $K$ un compact inclus dans ${\Cc}$, $f$ une fonction
continue sur $K$ et $(f_n)_{n\ge 1}$ une suite de fonctions définies sur
$K$ convergeant uniformément vers $f$ sur $K$. Montrer que
$(\exp{f_n})_{n\ge 1}$ tend vers $\exp f$ uniformément sur $K$ (on
pourra appliquer la première inégalité de a) à $f_n(z)-f(z)$ ).

\item Montrer que $\left(\left(1+{z\over n}\right)^n\right)_{n\ge 1}$ tend
vers $e^z$ uniformément sur tout compact de ${\Cc}$.
\end{enumerate}
\finenonce{006651}



\finexercice      


\exercice{6652, queffelec, 2011/10/16}
\enonce{006652}{} 
\begin{enumerate}
\item Soient $\rho $ un réel strictement positif, $z$ et $w$ des
nombres complexes tels que $\vert z\vert >\rho $ et $\vert w\vert
>\rho $, et $n$ un entier naturel.

Montrer que 
$$\left\vert {1\over w^n}-{1\over z^n}\right\vert\le \vert z-w\vert
{n\over \rho ^{n+1}}$$
et que
$$\left\vert {1\over z-w}\left({1\over w^n}-{1\over
z^n}\right)-{n\over z^{n+1}}\right\vert =
\left\vert\sum_{k=1}^n\left({1\over w^n}-{1\over z^n}\right){1\over
z^{n-k+1}}\right\vert \le \vert z-w\vert {n^2\over \rho ^{n+2}}$$

Soient maintenant $\sigma $ et $\phi $ deux fonctions continues à
valeurs complexes définies sur un intervalle $I=[a,b]$. On fixe un point
$z\in \C\setminus \sigma (I)$ et on pose $\rho ={1\over 2}\inf_{a\le
t\le b}\vert \sigma (t)-z\vert $.

\item Soit $\displaystyle g(z)=\int_a^b{\phi (t)\over (\sigma (t)-z)^n}dt$.
En remplaçant dans a) $z$ par $\sigma (t)-z$ et $w$ par $\sigma
(t)-z-h$ avec $\vert h\vert <\rho $, montrer que
$$\left\vert {g(z+h)-g(z)\over h}-n\int_a^b{\phi (t)\over (\sigma
(t)-z)^{n+1}}dt\right\vert\le {\vert h\vert n^2\over \rho
^{n+2}}\int_a^b\vert \phi (t)\vert dt$$

En déduire que $g(z)$ est holomorphe sur $\C\setminus \sigma (I)$ et
que $g'$ est donnée par
$$g'(z)=n\int_a^b{\phi (t)\over (\sigma (t)-z)^{n+1}}dt$$
\end{enumerate}
\finenonce{006652}



\finexercice


\exercice{6813, gijs, 2011/10/16}
\enonce{006813}{}
Soit $f$ une application holomorphe (c'est-à-dire
$\Cc$-différentiable) d'un ouvert $\Omega \subset \Cc$
dans $\Cc$.
\begin{enumerate}
\item
Soit $z_0 \in \Omega$ tel que $f'(z_0) = 0$. Démontrer
que $z \mapsto \vert f(z) \vert$ est
$\Cc$-différentiable en $z_0$. (On pourra utiliser le
développement en série entière de $f$ au voisinage
de $z_0$.)

\item
Soit $z_1 = x_1 + iy_1 \in \Omega$ tel que $f(z_1) \neq
0$. Démontrer que $(x,y) \mapsto \vert f(x+iy) \vert$
est $\Rr$-différentiable en $(x_1,y_1)$.

\item
Soit $z_2 = x_2 + iy_2 \in \Omega$ tel que $f(z_2) \neq
0$ et tel que $z \mapsto \vert f(z) \vert$ soit
$\Cc$-différentiable en $z_2$. Démontrer que $f'(z_2)
= 0$. (On pourra utiliser les conditions de
Cauchy-Riemann.)

\item
Soit $z_3 = x_3 + iy_3 \in \Omega$ tel que $f(z_3) = 0$
et $f'(z_3) \neq 0$. 

L'application $(x,y) \mapsto \vert
f(x+iy) \vert$ est-elle $\Rr$-différentiable en
$(x_3,y_3)$~? 
L'application $z \mapsto \vert f(z) \vert$ est-elle
$\Cc$-différentiable en $z_3$~?

\item
Donner le domaine où $z \mapsto \vert f(z) \vert$ est
continue, puis celui où $z \mapsto \vert f(z) \vert$
est $\Cc$-différentiable, et enfin celui où
$(x,y) \mapsto \vert f(x+iy) \vert$ est
$\Rr$-différentiable.
\end{enumerate}
\finenonce{006813}



\finexercice


\exercice{7214, megy, 2021/02/22}

\enonce{007214}{}
On considère la fonction $f : \C\to \C, z\mapsto \sqrt{|xy|}$, où l'on a noté $z=x+iy$ la forme algébrique de $z$. Montrer qu'elle vérifie les conditions de Cauchy-Riemann en l'origine, mais qu'elle n'est pas $\C$-dérivable en l'origine.
\finenonce{007214}



\finexercice


\exercice{7215, megy, 2021/02/22}

\enonce{007215}{}
On considère la fonction $f(z)=|x^2-y^2|+2i|xy|$ où $z=x+iy.$ Déterminer l'ensemble où $f$ est $\C$-dérivable.
\finenonce{007215}



\finexercice


\exercice{7216, megy, 2021/02/22}

\enonce{007216}{}
Soit $U\subset \C$  un ouvert connexe. Soit \(f:U\to \C\) une fonction holomorphe. On note \(f=u+iv\) l'écriture sous forme algébrique des valeurs de $f$. Montrer que les assertions suivantes sont équivalentes.
\begin{enumerate}
\item \(f\) est constante.
\item \(2u+3v=5\).
\item \(v=u^2\).
\item \(u^2+3v^2=1\).
\end{enumerate}
\finenonce{007216}



\finexercice


\exercice{7217, megy, 2021/02/22}

\enonce{007217}{}
Soit $U\subset \C$  un ouvert connexe. Soit \(f:U\to \C\) une fonction holomorphe. On note \(f=u+iv\) l'écriture sous forme algébrique des valeurs de $f$. Montrer que les assertions suivantes sont équivalentes.
\begin{enumerate}
\item \(f\) est constante.
\item Il existe une fonction $\phi \in \mathcal C^1(\R,\R)$ telle que $v=\phi(u)$.
\item Il existe une fonction $\psi \in \mathcal C^1(\R^2,\R)$ dont la différentielle ne s'annule qu'en un nombre fini de points, et telle que $\psi(u,v)$ est constante.
\end{enumerate}
\finenonce{007217}



\finexercice


\exercice{7218, megy, 2021/02/22}

\enonce{007218}{}
Soit \(U\subset \C\) un ouvert. Soit \(f:U\to \C\) une fonction de classe \(\mathcal{C}^1\). 
 Montrer que \[\overline{\frac{\partial f}{\partial z}}=\frac{\partial \overline{f}}{\partial \bar z}
 \quad \text{et} \quad 
 \overline{\frac{\partial f}{\partial\bar z}}=\frac{\partial \overline{f}}{\partial z}.\]
\finenonce{007218}



\finexercice


\exercice{7219, megy, 2021/02/22}

\enonce{007219}{}
Pour chacune des fonctions \(f\), déterminer l'ensemble de définition de \(f\) et calculer \(\frac{\partial f}{\partial z}\) et \(\frac{\partial f}{\partial \bar z}\).
\begin{enumerate}
\item \(f(z)=z^2\bar z\)
\item \(f(z)=e^{\bar z}\)
\item \(f(z)=\frac{\bar z-i}{z^2+iz-2}\)
\item \(f(z)=\frac{\Re(z)}{\Im(z)}\)
\end{enumerate}
\finenonce{007219}



\finexercice


\exercice{7220, megy, 2021/02/22}

\enonce{007220}{}
Montrer que si $f$ est holomorphe et ne s'annule jamais, alors $\log |f|$ est harmonique.
\finenonce{007220}



\finexercice


\exercice{7221, megy, 2021/02/22}

\enonce{007221}{}
Soit \(V\subset\C \) un ouvert et soit \(g:V\to \C\) une fonction harmonique. On suppose que  \(f\) est  holomorphe et que \(f(U)\subset V\). Montrer que \(g\circ f\) est harmonique.
\finenonce{007221}



\finexercice


\exercice{7222, megy, 2021/03/01}

\enonce{007222}{}
Soit $U\subseteq \C$ un ouvert et $f:U\to \C$ une fonction analytique. Soit $F : U\to \C$ une fonction holomorphe telle que $F'=f$. Montrer que $F$ est analytique.  En déduire que le logarithme principal $\mathrm{Log}$ est analytique sur $\C\setminus \R_-$.
\finenonce{007222}



\finexercice


\exercice{7223, megy, 2021/03/01}

\enonce{007223}{}
Montrer que la fonction \(f:\C\setminus \{1,2\}\to \C\) définie par 
\[f(z)=\frac{1}{(z-1)(z-2)}\]
est analytique, puis pour tout \(z_0\in \C\setminus \{1,2\}\) calculer le développement en série entière de \(f\) centré en \(z_0\).
\finenonce{007223}



\finexercice


\exercice{7224, megy, 2021/03/01}

\enonce{007224}{}
Soit \(U\subset \C\) un ouvert connexe. On note \(\mathcal{A}(U)\) l'ensemble des fonctions analytiques sur \(U\). Montrer que, muni de l'addition et de la multiplication de fonctions, \(\mathcal{A}(U)\) est un anneau intègre. Que se passe-t-il si l'on retire l'hypothèse de connexité ?
\finenonce{007224}



\finexercice


\exercice{7225, megy, 2021/03/01}

\enonce{007225}{}
Soit \(\displaystyle{f(z)=\sum_{n=0}^{+\infty}}a_nz^n\) une série entière de rayon de convergence \(R>0\). Un point \(z_0\in \partial B(0,R)\) est un \emph{point régulier} de \(f\) si il existe une  extension analytique de \(f\) dans un voisinage de \(B(0,R)\cup\{z_0\}\). Si \(z_0\in \partial B(0,R)\) n'est pas régulier pour \(f\), on dit que c'est un \emph{point singulier} de \(f\). On note \({\rm Sing}(f)\subset \partial B(0,R)\) l'ensemble des points singuliers de \(f\).
\begin{enumerate}
\item Déterminer les points réguliers et singuliers pour \(\displaystyle{f(z)=\sum_{n=0}^{+\infty}z^n}\) et \(\displaystyle{f(z)=\sum_{n=0}^{+\infty}(-1)^n\frac{z^n}{n}}\). 
\item Quel est le lien entre la régularité d'un point \(z_0\in \partial B(0,R)\) et la convergence de \(\displaystyle{\sum_{n=0}^{+\infty}a_nz_0^n}\) ?
\item Montrer que \({\rm Sing}(f)\) est fermé.
\item On considère \(\displaystyle{f(z)=\sum_{n=0}^{+\infty}z^{2^n}}\).
\begin{enumerate}
\item Montrer que le rayon de convergence de \(f\) est \(1\).
\item Montrer que \(1\) est un point singulier de \(f\). (Indication, montrer que \(\lim_{t\to 1^-}\sum_{n=0}^{+\infty}t^{2^n}=+\infty\)).
\item Montrer que pour tout \(m\in \N\), toute racine \(2^m\)-ième de \(1\) est un point singulier de \(f\). (Indication, observer que pour tout \(m\in \N\) et pour tout \(z\in B(0,1)\),  \(f\big(z^{2^{m+1}}\big)=f\big(z^{2^m}\big)-z^{2^m}\)).
\item En déduire que \({\rm Sing}(f)=\partial B(0,1)\).
\end{enumerate}
\end{enumerate}
\finenonce{007225}



\finexercice


\exercice{7226, megy, 2021/03/01}

\enonce{007226}{}
Soit \(\displaystyle{f(z)=\sum_{n\geqslant 0}a_nz^n}\) une série entière de rayon de convergence \(R>0\) et telle que \(a_0\neq 0\). L'objectif de cet exercice est de démontrer que la fonction \(\frac{1}{f}\) est développable en série entière en \(0\).
\begin{enumerate}
\item On suppose que ceci est le cas et que \(\displaystyle{\frac{1}{f(z)}=\sum_{n=0}^{+\infty}b_nz^n}\). Quelle relation de récurrence vérifie la suite \((b_n)_{n\in \N}\)?
\item Soit \((b_n)_{n\in \N}\) une suite vérifiant le relation de récurrence précédente. Montrer qu'il existe \(C>0\) tel que pour tout \(n\geqslant 0\), on a 
\[|b_n|\leqslant \frac{C^n}{|a_0|}.\]
\item En déduire que \(\frac{1}{f}\) est développable en série entière en \(0\).
\end{enumerate}
\finenonce{007226}



\finexercice


\exercice{7524, mourougane, 2021/08/10}
\enonce{007524}{$\Cc$-dérivabilité}
\begin{enumerate}
 \item \'Ecrire en termes d'$\varepsilon$ et $\delta$ la définition de la $\Cc$-dérivabilité d'une fonction $f : D\to \Cc$
 d'un ouvert $D$ de $\Cc$ en un point $a$ de $D$.
 \item Soit $n$ un entier naturel non nul. Soit $a$ un point de $\Cc$.
 Déterminer une fonction $f_1 : \Cc\to\Cc$ telle que sur $\Cc$,
 $$z^n=a^n+(z-a)f_1(z).$$
 En déduire que l'application $\Cc\to\Cc$, $z\mapsto z^n$ est $\Cc$-dérivable en $a$ et déterminer sa dérivée en $a$.
\item Démontrer que si $f$ et $g$ sont deux applications d'un ouvert $D$ de $\Cc$ dans $\Cc$ qui sont $\Cc$-dérivables en $a$,
alors leur produit est $\Cc$-dérivable en $a$. Déterminer alors la dérivée du produit en $a$ à l'aide des dérivées en $a$ de $f$ et de $g$.
\item Donner un exemple d'application $f :\Cc\to\Cc$ continue mais nulle part $\Cc$-dérivable.
\end{enumerate}
\finenonce{007524}
\finexercice


\exercice{7525, mourougane, 2021/08/10}
\enonce{007525}{\'Equations de Cauchy-Riemann}
\begin{enumerate}
\item Soit $f : D\to\Cc$ une application $\Cc$-dérivable en un point $a$ de $D$.
\'Enoncer les équation de Cauchy-Riemann.
\item \'Ecrire les équations de Cauchy-Riemann pour l'application $\Cc\to\Cc$, $z\mapsto z^3$.
\item Donner l'exemple d'une application $f :\Rr^2\to\Rr^2$ différentiable en tout point et
qui ne vérifie en aucun point les équations de Cauchy-Riemann. 
\end{enumerate}
\finenonce{007525}
\finexercice


\exercice{7526, mourougane, 2021/08/10}
\enonce{007526}{Holomorphie}
\begin{enumerate}
 \item Donner l'exemple d'une fonction holomorphe $f$ sur $\Cc$.
 \item Vérifier en calculant $\Delta u$, que sa partie réelle $u$ est harmonique.
\end{enumerate}
\finenonce{007526}
\finexercice


\exercice{7527, mourougane, 2021/08/10}
\enonce{007527}{Un exemple}
Soit $f : \Cc\to\Cc$, $z=x+iy\mapsto x^3y^2+ix^2y^3$ et $g=f_{\Rr} : \Rr^2\to\Rr^2$ l'application réelle associée.
\begin{enumerate}
\item En quels points de $\Rr^2$, $g$ est-elle différentiable ?
\item En quels points de $\Rr^2$, $g$ vérifie-t-elle les équations de Cauchy-Riemann ?
\item En quels points de $\Cc$, $f$ est-elle $\Cc$-dérivable ?
\item En quels points de $\Cc$, $f$ est-elle holomorphe ?
\end{enumerate}
\finenonce{007527}
\finexercice


\exercice{7528, mourougane, 2021/08/10}
\enonce{007528}{Exponentielle}
\begin{enumerate}
\item Reprendre l'exercice précédent pour la fonction $\exp : \Cc\to\Cc$, $z=x+iy\mapsto e^x\cos y+ie^x\sin y$.
\item Déterminer la dérivée de cette application, là où elle est $\Cc$-dérivable.
\item Démontrer ou réfuter~: la fonction $f(z)=\exp(\bar z)$ est holo\-morphe dans $\Cc$.
\end{enumerate}
\finenonce{007528}
\finexercice


\exercice{7529, mourougane, 2021/08/10}
\enonce{007529}{Logarithme}
On considère $D:=\{z\in\Cc\ /\ \Re(z)\not=0.\}$ et l'application
$$\begin{array}{cccc}
\ell : &D&\to&\Cc\cr &z=x+iy&\mapsto& \frac{1}{2}\log (x^2+y^2)+i\arctan\frac{y}{x}.
          \end{array}$$

\begin{enumerate}
\item Montrer que $\ell$ est $\Cc$-dérivable sur $D$.
\item Déterminer sa dérivée sur $D$.
\end{enumerate}
\finenonce{007529}
\finexercice


\exercice{7530, mourougane, 2021/08/10}
\enonce{007530}{En coordonnées polaires}
Déterminer les équations de Cauchy Riemann en coordonnées polaires. 
\finenonce{007530}
\finexercice


\exercice{7531, mourougane, 2021/08/10}
\enonce{007531}{Exemples}
Montrer qu'un polynôme $P(z, \overline{z})$ est holomorphe si et seulement aucun monôme ne contient le facteur~$\overline{z}$. 
\finenonce{007531}
\finexercice


\exercice{7532, mourougane, 2021/08/10}
\enonce{007532}{Fonctions localement constantes}
\begin{enumerate}
\item Montrer qu'une fonction holomorphe sur un ouvert $D$ de $\Cc$ qui ne prend que des valeurs réelles
est localement constante.
\item Que dire d'une fonction holomorphe sur un ouvert de $\Cc$ dont la partie réelle est constante ?
\item Que dire d'une fonction holomorphe $f=u+iv$ sur un ouvert de $\Cc$ dont la conjuguée $\overline{f}:=u-iv$ 
est aussi holomorphe ?
\item Montrer qu'une fonction holomorphe qui ne prend que des valeurs de module $1$ est localement constante.
\end{enumerate}
\finenonce{007532}
\finexercice


\exercice{7533, mourougane, 2021/08/10}
\enonce{007533}{Fonctions harmoniques}
\begin{enumerate}
\item Déterminer, si possible, une fonction holomorphe sur $\Cc$ dont la partie réelle est $u(x,y)=x^2+y^2$.
\item Déterminer, si possible, une fonction holomorphe sur $\Cc$ dont la partie réelle est $u(x,y)=x^2-y^2$.
 \item Montrer sans calcul que la fonction $f : \Rr^2\to\Rr^2$, $(x,y)\mapsto 2xy$ est harmonique sur $\Rr^2$.
\end{enumerate}
\finenonce{007533}
\finexercice


\exercice{7534, mourougane, 2021/08/10}
\enonce{007534}{Fonctions holomorphes}
\begin{enumerate}
 \item Déterminer, si possible, une fonction holomorphe sur $\Cc$ dont la partie réelle est $u(x,y)=e^{-y}\sin(x) -e^{y}\cos(x)$.
 \item Déterminer toutes les fonctions holomorphes sur $\Cc$ dont la partie réelle est $u(x,y)=e^{-y}\sin(x) -e^{y}\cos(x)$.
\end{enumerate}
\finenonce{007534}
\finexercice


\exercice{7535, mourougane, 2021/08/10}
\enonce{007535}{Polynômes harmoniques}
 Soit $u(x,y)\in\Rr[x,y]$ un polynôme harmonique.
 Montrer que la fonction $$p(z):=2u(\frac{z}{2},\frac{z}{2i})-u(0,0)$$
 est holomorphe de partie réelle $u$.
\finenonce{007535}
\finexercice


\exercice{7536, mourougane, 2021/08/10}
\enonce{007536}{Laplacien}
Soit une fonction holomorphe $ f $ définie sur un ouvert de $\Cc$. Montrer que 
$$
\Delta( \arrowvert f \arrowvert ^2 ) = 4 \arrowvert f' \arrowvert ^2 \;\;\; ; \;\;\; 
\Delta \ln ( 1 + \arrowvert f \arrowvert ^2 ) = \frac{4 \arrowvert f' \arrowvert ^2}{(1 + \arrowvert f' \arrowvert ^2)^2}
$$ 
\finenonce{007536}
\finexercice


\exercice{7537, mourougane, 2021/08/10}
\enonce{007537}{Conservation des angles}
On considère l'application $f :\Cc^\times\to\Cc^\times$, $z\mapsto z^2$.
\begin{enumerate}
 \item Montrer que $f$ conserve les angles.
 \item Soit $a$ un nombre réel non nul. Montrer que l'image de la droite d'équation $x=a$ est incluse dans la parabole d'équation
$v^2=4a^2(a^2-u)$. Représenter cette parabole pour $a=1$ et $a=2$.
\item Soit $b$ un nombre réel non nul. Déterminer une parabole contenant l'image de la droite d'équation $y=b$.
Représenter cette parabole pour $b=1$ et $b=2$.
\item Vérifier l'orthogonalité aux points d'intersection des quatre paraboles.
\end{enumerate}
\finenonce{007537}
\finexercice


\exercice{7538, mourougane, 2021/08/10}
\enonce{007538}{Conservation des angles}
On considère l'application $f :\Cc^\times\to\Cc^\times$, $z\mapsto \frac{1}{2}(z+\frac{1}{z})$.
\begin{enumerate}
\item Déterminer le lieu où $f$ préserve les angles.
\item Montrer que si on note $r=|z|$, $u=re(f)$ et $v=im(f)$, alors 
$$u=\frac{1}{2}(r+\frac{1}{r})\frac{x}{r}\quad \text{ et } \quad  v=\frac{1}{2}(r-\frac{1}{r})\frac{x}{r}.$$
puis
$$\frac{u^2}{\frac{1}{4}(r+\frac{1}{r})^2}+\frac{v^2}{\frac{1}{4}(r-\frac{1}{r})^2}=1
\quad \text{ et } \quad \frac{u^2}{\frac{x^2}{r^2}}-\frac{v^2}{\frac{y^2}{r^2}}=1.$$
\item Déterminer l'image par $f$ des cercles de centre $o$ et de rayon $1$ et $2$.
\item Déterminer l'image par $f$ des segments radiaux $z=\frac{1+i}{\sqrt{2}}t$ et $z=\frac{1-i}{\sqrt{2}}t$, quand $t$ varie dans $]0,1[$.
\item Montrer que $f$ est surjective.
\item Montrer que l'image réciproque d'un point de $\Cc-[-1,1]$ est composée de deux points,
l'un dans $\Delta^\times$ l'autre dans $\Cc-\Delta$.
\item Montrer que $f$ est une bijection holomorphe de $\Delta^\times$ sur $\Cc-[-1,1]$.
\end{enumerate}
\finenonce{007538}
\finexercice


\exercice{7539, mourougane, 2021/08/10}
\enonce{007539}{Applications linéaires fractionnaires}
On rappelle qu'à toute matrice $A=\begin{pmatrix} a&b\\c&d\end{pmatrix}$ de $GL(2,\Cc)$, on associe l'application linéaire fractionnaire
$$\begin{array}{cccc}
 f_A :& \Cc-\{-\frac{-d}{c}\}&\longrightarrow&\Cc-\{\frac{a}{c}\}\\&z&\longmapsto& \frac{az+b}{cz+d}
\end{array}$$

\begin{enumerate}
\item Vérifier que $f_{A^{-1}}=f_A^{-1}$. 
\item On choisit désormais la matrice $C=\begin{pmatrix} 1&-i\\1&i\end{pmatrix}$. 
Expliciter le biholomorphisme $h=f_C$ et son inverse~$h^{-1}$.
\item Montrer que 
\begin{eqnarray*}
 1-|h(z)|^2=\frac{4 Im(z)}{|z+i|^2}\\
 Im(h^{-1}(z))=\frac{1-|z|^2}{|1-z|^2}
\end{eqnarray*}

\item En déduire l'image du demi-plan de Poincaré $\mathbb{H}$ par $h$.
\end{enumerate}
\finenonce{007539}
\finexercice


\exercice{7540, mourougane, 2021/08/10}
\enonce{007540}{Biholomorphismes entre domaines}
On rappelle que $\Delta:=\{z\in\Cc / |z|<1\}$, $\mathbb{H}:=\{z\in\Cc / Im(z)>0\}$ et $\Cc^-:=\Cc-\{z\in\Cc / Im(z)=0 \quad \text{ et } \quad Re(z)\leq 0\}$. 
\begin{enumerate}
\item Montrer que l'application 
$$\begin{array}{cccc}
 q:& \mathbb{H}&\longrightarrow&\Cc^{-}\\&z&\longmapsto& -z^2
\end{array}$$
est holomorphe et bijective.
\item En déduire une application holomorphe et bijective de $\Delta$ sur $\Cc^-$.
\end{enumerate}
\finenonce{007540}
\finexercice


\exercice{7541, mourougane, 2021/08/10}
\enonce{007541}{Biholomorphisme de $\mathbb{H}$}
On rappelle qu'à toute matrice $A=\begin{pmatrix} a&b\\c&d\end{pmatrix}$ de $SL(2,\Rr)$,
on associe l'application linéaire fractionnaire
$$\begin{array}{cccc}
 h_A :& \Cc-\{-\frac{-d}{c}\}&\longrightarrow&\Cc-\{\frac{a}{c}\}\\&z&\longmapsto& \frac{az+b}{cz+d}
\end{array}$$
\begin{enumerate}
\item Montrer que $h_A$ envoie $\mathbb{H}$ sur $\mathbb{H}$.
\item Montrer que pour tout élément $z$ de $\mathbb{H}$, il existe $A\in SL(2,\Rr)$ 
tel que $h_A(i)=z$.
\end{enumerate}
\finenonce{007541}
\finexercice


\exercice{7569, mourougane, 2021/08/10}
\enonce{007569}{Polynômes}
Soit $D$ un ouvert de $\Cc$ et $f : D\to\Cc$ une application holomorphe.
Montrer qu'il y a équivalence entre 
\begin{enumerate}
\item $f$ est polynomiale.
\item Il existe $c\in D$ tel que en dehors d'un nombre fini de $n\in\Nn$, $f^{(n)}(c)=0$.
\end{enumerate}
\finenonce{007569}
\finexercice


\exercice{7570, mourougane, 2021/08/10}
\enonce{007570}{Propriétés}
Pour chacune des propriétés suivantes, donner un exemple d'application holomorphe au voisinage de $0$
qui la satisfait, ou bien démontrer qu'il n'existe pas de telles applications.
\begin{enumerate}
 \item Pour presque tout $n\in\Nn$, $f(\frac{1}{n})=(-1)^n\frac{1}{n}$.
 \item Pour presque tout $n\in\Nn$,$f(\frac{1}{n})=\frac{1}{n^2-1}$.
 \item Pour presque tout $k\in\Nn$,$|f^{(k)}(0)|\geq (n!)^2$.
 \item Pour presque tout $n\in\Nn$,$|f(\frac{1}{n})|\leq e^{-n}$ et $f$ est non nulle.
\end{enumerate}
\finenonce{007570}
\finexercice


\exercice{7571, mourougane, 2021/08/10}
\enonce{007571}{Prolongement}
Soit $D$ un ouvert de $\Cc$ $a\in D$ et $f : D-\{a\}\to \Cc$ holomorphe.
On suppose que $f'$ admet un prolongement holomorphe à $D$. 
Est-ce aussi le cas pour $f$ ?
\finenonce{007571}
\finexercice


\exercice{7572, mourougane, 2021/08/10}
\enonce{007572}{Contraintes}
\begin{enumerate}
 \item Soit une fonction holomorphe $f$ entière définie sur le plan complexe tout entier, on suppose que $ Re(f) \leq 0 $. 
 Montrer que $f$ est une fonction constante. On pourrait considérer la fonction $e^{f(z)}$.
 \item Soit une fonction holomorphe $f$ au voisinage de $0$.
 Montrer que si $ f \left( \frac{1}{n} \right) = f \left( \frac{1}{2n} \right)$ pour tout $n$ assez grand alors $f$ est une constante. 
\item Soit $D$ un ouvert de $\Cc$ et $f=u+iv : D\to\Cc$ une application holomorphe. On suppose que sur $D$,
$v=u^2$. Montrer que $f$ est une fonction constante.
\item Soit une fonction entière $f$ telle que $\arrowvert f \arrowvert $ tend vers l'infini si $\arrowvert z \arrowvert $ tend vers l'infini. Montrer que :
\begin{enumerate}
\item $f$ n'admet qu'un nombre fini de zéros.
\item $f$ est un polynôme.
\end{enumerate}
\item Soit une fonction entière $f$ non constante, montrer que l'image du plan complexe par $f$ est dense dans $\Cc$.
\item Soit une fonction $f$ holomorphe sur le plan complexe en entier, vérifiant $ f(x+1)=f(x)$ pour $ x \in \Rr $. Montrer que $f$ est une fonction périodique
de période $1$. 
\end{enumerate}
\finenonce{007572}
\finexercice

\section{441.00 Fonction logarithme et fonction puissance}


\exercice{2851, burnol, 2009/12/15}
\enonce{002851}{} Montrer qu'il existe une (unique) fonction analytique
sur $\Cc\setminus[-1,1]$ qui vaut $\sqrt{a^2-1}$ pour
$a>1$. 
\emph{Indication :} montrer pour commencer que la formule
$f(a) = \exp(\frac12 \mathrm{Log}(a-1) + \frac12\mathrm{Log}(a+1))$ donne
une solution sur l'ouvert
$\Cc\setminus]-\infty,1]$. Puis montrer que  $g(a)
=  -\exp(\frac12 \mathrm{Log}(-a-1) + \frac12\mathrm{Log}(-a+1)) = -f(-a)$
est analytique sur
$\Cc\setminus[-1,+\infty[$. Enfin montrer que $g(a) = f(a)$
dans le demi-plan supérieur et aussi dans le demi-plan
inférieur en calculant $f(\pm i)$ et donc $g(\pm i)$ et en
expliquant pourquoi a priori le quotient $g(a)/f(a)$ est
constant dans ces deux demi-plans.
\finenonce{002851}
\finexercice


\exercice{2852, burnol, 2009/12/15}
\enonce{002852}{}
\begin{enumerate}
  \item   On considère la fonction analytique $\phi(a) = \mathrm{Log}(a-1) -
  \mathrm{Log}(a+1)$ dans le demi-plan supérieur et la fonction
  analytique $\psi(a) =  \mathrm{Log}(a-1) -
  \mathrm{Log}(a+1)$ dans le demi-plan inférieur. Montrer que $\phi$
  et $\psi$ sont la restriction à leurs demi-plans
  respectifs d'une fonction analytique sur
  $\Cc\setminus[-1,+1]$. 
  \emph{Indication :} il y a plusieurs
  raisonnements possibles et plusieurs indications
  possibles. Donc, débrouillez vous.
  \item On considère la fonction $a\mapsto
\frac{a-1}{a+1}$. Quelle est l'image par cette fonction de
l'intervalle $]-1,1[$?  Quelle est l'image par cette
fonction de $\Cc\setminus[-1,+1]$? En déduire que la
fonction composée $\Phi(a) = \mathrm{Log}\frac{a-1}{a+1}$ existe et
est analytique sur $\Cc\setminus[-1,+1]$. Retrouver le
résultat de la question précédente (et montrer que $\phi$,
$\psi$ et $\Phi$ coïncident dans les intersections
deux-à-deux de leurs ouverts de définitions).
  \item Quel est le développement en série de Laurent de la
  fonction analytique $\Phi$ dans la couronne
  $1<|a|<\infty$? Que vaut par exemple $\int_{|a|=2}
  \Phi(a)a^{18} da$?
\end{enumerate}
\finenonce{002852}
\finexercice

\exercice{6653, queffelec, 2011/10/16}
\enonce{006653}{}
Montrer qu'il n'existe pas de détermination holomorphe du logarithme de $z$ sur
$\Cc\backslash\{0\}$ tout entier. (On raisonnera par l'absurde et on exhibera
ainsi une application continue injective du cercle unité dans $\Rr$).
\finenonce{006653}



\finexercice


\exercice{6654, queffelec, 2011/10/16}
\enonce{006654}{}
Soit $\hbox{Log} z$ la détermination principale du logarithme dans 
$\Cc\backslash \Rr^-$, ie $\hbox{Log} z=\hbox{Ln}|z|+i\hbox{Arg} z$ où
$|\hbox{Arg}z|<\pi$, et on définit $z^\alpha=e^{\alpha{\rm Log}z}$
\begin{enumerate}
\item On considère $z=e^{2i\pi\over3}$; comparer $\hbox{Log} (z^2)$ et $2\hbox{Log}
z$. 

\item On considère $z=e^{3i\pi\over4}$; comparer $z^{2i}$, $(z^2)^i$ et $(z^i)^2$.
\end{enumerate}
\finenonce{006654}



\finexercice


\exercice{6655, queffelec, 2011/10/16}
\enonce{006655}{}
On se propose de calculer les sommes de séries convergentes pour $0<t<2\pi$ 

$$\sum_1^\infty {\cos nt\over n},\quad  \sum_1^\infty {\sin nt\over n}.$$
\begin{enumerate}
\item Rappeler pourquoi $S(z)=-\sum_{n\geq 1}{z^n\over n}$ coincide
sur $D$ avec la détermi\-nation principale $\hbox{Log}(1-z)$.

\item Soit $r<1$; calculer $\sum_{n\geq 1} {r^n\cos nt\over n}$ et
$\sum_{n\geq 1} {r^n\sin nt\over n}$.

\item En déduire la valeur de ces sommes (on pourra utiliser le théorème d'Abel).
\end{enumerate}
\finenonce{006655}



\finexercice


\exercice{6656, queffelec, 2011/10/16}
\enonce{006656}{}
Soit $\Omega$ un ouvert connexe de $\Cc$ et $f$ une fonction complexe sans
zéro sur $\Omega$. 
On rappelle que
$f$ admet un logarithme continu (resp. holomorphe) sur $\Omega$ s'il existe une
fonction
$g$
 continue (resp. holomorphe) sur $\Omega$ telle que $e^{g(z)}=f(z)$. 

Montrer
que deux déterminations continues du logarithme de $f$ sur $\Omega$ diffèrent
d'une constante $2ki\pi$.

 En reproduisant la démonstration du théorème d'inversion locale, montrer que
si
$f$ admet sur $\Omega$ un logarithme continu, elle y admet un logarithme
holomorphe.
\finenonce{006656}



\finexercice


\exercice{6657, queffelec, 2011/10/16}
\enonce{006657}{}
 On rappelle qu'une fonction complexe $f$ a une
racine
$n$ième holomorphe dans un ouvert connexe $\Omega$ s'il existe $g\in H(\Omega)$
telle que $g^n(z)=f(z)$.
\begin{enumerate}
\item Montrer que si $f$ admet un logarithme holomorphe dans $\Omega$, elle y
admet des racines de tous ordres; montrer, sur un exemple, qu'une fonction
holomorphe $f$ peut admettre une racine sans admettre de logarithme
(holomorphe).

\item Si $g_1,g_2$ sont deux fonctions continues de 
$\Omega$ connexe dans ${\Cc}\backslash\{0\}$, telles que $g_1^n=g_2^n$ pour un entier $n\geq1$,
montrer que  $g_1=e^{2i\pi k\over n}g_2$ où $k$ est un entier et $g_1=g_2$ dès
que les fonctions coincident en un point.
\end{enumerate}
\finenonce{006657}



\finexercice


\exercice{6658, queffelec, 2011/10/16}
\enonce{006658}{}
\begin{enumerate}
\item Montrer que $\Re(\cos z)>0$ si $|\Re z|<{\pi\over2}$. En déduire une
détermination holo\-morphe du logarithme de  $\cos z$ dans
$\{|\Re z|<{\pi\over2}\}$.

\item Montrer que l'on peut définir une fonction holomorphe $f(z)=\hbox{Log}
{1+iz\over 1-iz}$ sur l'ouvert $U={\Cc}\backslash S$ où $S=\{ix\ ;\
|x|\geq1\}$.

\item Montrer que l'on peut définir une fonction holomorphe $f(z)=\hbox{Log}
\sqrt{z^3-1}$ sur un ouvert $U$ à déterminer (où $\sqrt{}$ désigne la
détermination principale de la racine).
\end{enumerate}
\finenonce{006658}



\finexercice


\exercice{6659, queffelec, 2011/10/16}
\enonce{006659}{}
\begin{enumerate}
\item Soit $\Omega={\Cc}\backslash [-1,1]$; montrer que la fonction
$$g(z)=\hbox{exp}({1\over2}\hbox{Log}(z+1)+{1\over2}\hbox{Log}(z-1))$$
se prolonge en une fonction continue sur $\Omega$ 
et fournit ainsi une racine carrée holomorphe de $z^2-1$ dans
$\Omega$, bien que $\hbox{Log}(z^2-1)$ n'ait pas de prolongement continu à
$\Omega$.

\item Construire de même une racine carrée holomorphe de $z^2+1$ sur l'ouvert
connexe $\Omega'{\Cc}\backslash [-i,i]$.
\end{enumerate}
\finenonce{006659}



\finexercice


\exercice{6660, queffelec, 2011/10/16}
\enonce{006660}{} 
Montrer que
$$\mathrm{Arg} {(x+iy)} =\left\{ \begin{array}{ll}
\arctan{\left({y\over x}\right)}&\mbox{ si }x>0\\
{\pi\over 2}-\arctan{\left({x\over y}\right)}&\mbox{ si }y>0\\
-{\pi\over 2}-\arctan{\left({x\over y}\right)}&\mbox{ si }y<0
\end{array}\right.$$
En déduire que la fonction $\mathrm{Arg}$ est ${\Rr}$-différentiable
sur l'ouvert $\Omega$
$$\Omega ={\Rr}^2\setminus \{ (x,y)\vert\ y=0,x\le 0\}$$
et calculer sa différentielle.
\finenonce{006660}



\finexercice       


\exercice{6661, queffelec, 2011/10/16}
\enonce{006661}{}
Soit $w$ dans $\C$. Déterminer les nombres 
complexes $z$ tels que   $\cos z=\cos w$. Même question avec le sinus.
\finenonce{006661}



\finexercice       


\exercice{6662, queffelec, 2011/10/16}
\enonce{006662}{}
 Si $z=x+iy$ avec $x$ et $y$ réels,  montrer que l'on a
$$\begin{array}{lll}
\vert \sin z\vert^2&=&\sin^2x+\sh^2y\\
\vert \cos z\vert^2&=&\cos^2x+\sh^2y
\end{array}$$
En déduire les zéros de $\sin z$ et $\cos z$ dans $\C$.
\finenonce{006662}



\finexercice       


\exercice{6663, queffelec, 2011/10/16}
\enonce{006663}{}
 Montrer que la fonction $f(z)=\tan z$ définie
par $\tan z={\sin z\over \cos z}$ réalise une bijection de
$T=\{ z\in \C\vert -\pi/2<\Re z\le \pi/2,\ z\ne \pi/2\}$ sur 
$\C\setminus \{-i,i\}$. On pourra écrire $f=f_4\circ f_3\circ
f_2\circ f_1$ avec $f_1(z)=2iz$, $f_2(z)=e^z$,
$f_3(z)={1-z\over 1+z}$, $f_4(z)=iz$.
\finenonce{006663}



\finexercice      


\exercice{6664, queffelec, 2011/10/16}
\enonce{006664}{} Résoudre les équations $e^z=-3$, $\cos z=2$, $\sin z=2$,
$\tan z=2i$, $\ch z=1/2$ de la manière suivante :
\begin{enumerate}
\item en identifiant les parties réelles et imaginaires
\item en utilisant le logarithme
\end{enumerate}
\finenonce{006664}



\finexercice       


\exercice{6665, queffelec, 2011/10/16}
\enonce{006665}{} 
\begin{enumerate}
\item En utilisant la détermination principale du
logarithme, on définit les fonctions $z\mapsto
z^{1/2}$, $z\mapsto (1-z)^{1/3}$, $z\mapsto \left(
(1-2i)z\right)^{2i/5}$. Donner leur domaine de définition.

\item Soit $n\in{\Nn}$. Soit $f$ une détermination continue de $z^{1/n}$.
Montrer que $f(z)^n=z$ sur son domaine de définition.
\end{enumerate}
\finenonce{006665}



\finexercice      


\exercice{6666, queffelec, 2011/10/16}
\enonce{006666}{} On considère $f(z)=\sqrt{z^2}=(z^2)^{1/2}$ définie à l'aide de la
détermination principale du logarithme.
\begin{enumerate}
\item Trouver le domaine de définition et donner une expression
explicite de $f$.

\item Peut-on prolonger $f$ par continuité sur un ouvert plus grand ?

\item Mêmes questions en prenant la détermination ${\rm log}z=\mathrm{Log}
z+2i\pi$.
\end{enumerate}
\finenonce{006666}



\finexercice       


\exercice{6667, queffelec, 2011/10/16}
\enonce{006667}{} On définit à l'aide de la détermination principale du logarithme les
fonctions $f_1(z)=(z^3-1)^{1/2}$,
$f_2(z)=(z-1)^{1/2}(z-j)^{1/2}(z-j^2)^{1/2}$, et 
$f_3(z)=(1-z)^{1/2}(iz-ij)^{1/2}(iz-ij^2)^{1/2}$, où $j=e^{i\pi/3}$.
\begin{enumerate}
\item Trouver les domaines de définition des $f_k$ et montrer que l'on a
toujours $f_k(z)^2=z^3-1$.

\item Montrer que l'on peut prolonger $f_3$ par continuité sur un ouvert
plus grand.
\end{enumerate}
\finenonce{006667}



\finexercice       


\exercice{6668, queffelec, 2011/10/16}
\enonce{006668}{}
On veut démontrer qu'il existe une détermination continue $f$ de 
$\sqrt{1-z^2}$ sur $U=\C\setminus [-1,1]$ telle que $f(i)=\sqrt 2$.
\begin{enumerate}
\item Définir $f$ sur $\C\setminus \mathopen]-\infty,1]$ au moyen des fonctions
$\mathrm{Arg}{(z+1)}$ et $\mathrm{Arg}{(z-1)}$.

\item Soit $x$ un réel strictement inférieur à 1. Etudier $\lim_{y\to 0}
f(x+iy)$ quand $y$ tend vers 0 par valeurs positives puis négatives.
Conclure.

\item Montrer qu'on obtient ainsi une application $f$ telle que
$f(U)\subset U$ et $f\circ f=-{\rm Id_{\vert U}}$. En déduire que $f$ est
une bijection de $U$ sur lui-même.
\end{enumerate}
\finenonce{006668}



\finexercice  


\exercice{6806, gijs, 2011/10/16}
\enonce{006806}{}
On note $\mathrm{Log} : \Cc\setminus\,]-\infty,0] \to \Cc$
la détermination principale du logarithme, qui est
réelle pour $z$ réel positif.
On pose 

\begin{align*}
f(z) &= \exp\bigl(\tfrac13\mathrm{Log}(z) + \tfrac13\mathrm{Log}(z-1) +
\tfrac13\mathrm{Log}(z+1)\bigr)
\\
g(z) &= \exp\bigl(\tfrac13\mathrm{Log}(-z) +
\tfrac13\mathrm{Log}(1-z^2)\bigr)
\ .
\end{align*}


\begin{enumerate}
\item
Déterminer les domaines de définition $\Omega_f
\subset \Cc$ de $f$ et $\Omega_g \subset \Cc$ de
$g$ et montrer que $f$ et $g$ sont des branches continues
de $\root3\of {z^3-z}$.

\item 
Est-ce-qu'on peut élargir le domaine de définition de
$f$, c'est à dire~: existe-t-il un ouvert
$\widehat\Omega_f \subset \Cc$ et une fonction
continue $\widehat f : \widehat\Omega_f \to \Cc$
telle que $\Omega_f \subset \widehat\Omega_f$ et pour
tout $z\in \Omega_f$ on a $\widehat f(z) = f(z)$~?
N'oubliez pas de justifier votre réponse~!

\item
Est-ce-qu'on peut élargir le domaine de définition de
$g$~?
\end{enumerate}
\finenonce{006806}



\finexercice


\exercice{6824, gijs, 2011/10/16}
\enonce{006824}{}
On définit les fonctions $f_1$, $f_2$ et $f_3$ par les
formules
\begin{align*}
f_1(z) &= \exp(\tfrac13[ \mathrm{Log}(z+1) + \mathrm{Log}(z) + \mathrm{Log}(z-1)
+ \mathrm{Log}(z - \sqrt3)])
\\
f_2(z) &= \exp(\tfrac13[ \mathrm{Log}(z+1) + \mathrm{Log}(-z) + \mathrm{Log}(1-z)
+ \mathrm{Log}(\sqrt3 - z) +i\pi])
\\
f_3(z) &= \exp(\tfrac13[\mathrm{Log}(-1-z) + \mathrm{Log}(-z) + \mathrm{Log}(1-z)
+ \mathrm{Log}(z-\sqrt3) + i\pi])
\ ,
\end{align*}
où $\mathrm{Log}$ désigne le logarithme principal.
\begin{enumerate}
\item Calculer $f_1(\pm i)$, $f_2(\pm i)$ et
$f_3(\pm i)$.

\item Déterminer les domaines de définition de
$f_1$, $f_2$ et $f_3$.

\item Démontrer que $f_1$, $f_2$ et $f_3$ sont des
déterminations continues de $\root3 \of {z^4
-\sqrt3 z^3 -z^2 + \sqrt3 z}$.

\item Peut-on prolonger $f_1$ sur un ouvert plus
grand~? Si oui, lequel~?

\item Peut-on prolonger $f_2$  sur un ouvert plus
grand~?  Si oui, lequel~?

\item Peut-on prolonger $f_3$  sur un ouvert plus
grand~?  Si oui, lequel~?

\item Y-a-t-il un lien entre $f_1$, $f_2$ et $f_3$~?
\end{enumerate}
\finenonce{006824}



\finexercice


\exercice{7548, mourougane, 2021/08/10}
\enonce{007548}{trigonométrie hyperbolique}
On définit $\cosh z$ comme la somme de la série $\sum \frac{z^{2n}}{(2n)!}$ et
$\sinh z$ comme la somme de la série $\sum \frac{z^{2n+1}}{(2n+1)!}$.
\begin{enumerate}
 \item Déterminer leur rayon de convergence.
 \item Exprimer $\cosh$ et $\sinh$ à l'aide de la fonction exponentielle.
 \item Démontrer les formules d'addition pour $\cosh(z+w)$ et $\sinh(z+w)$ pour $z$ et $w$ dans $\Cc$.
 \item Démontrer que pour tous $x$ et $y$ dans $\Rr$,
 \begin{eqnarray*}
 \cos(x+iy)&=&\cos(x)\cosh(y)-i\sin(x)\sinh(y)\\
 \sin(x+iy)&=&\sin(x)\cosh(y)+i\cos(x)\sinh(y).
 \end{eqnarray*}
\end{enumerate}
\finenonce{007548}
\finexercice


\exercice{7549, mourougane, 2021/08/10}
\enonce{007549}{Exponentielle}
\begin{enumerate}
 \item Trouver la valeur minimale de $\arrowvert f(z) \arrowvert $ o\` u $ f(z) = e^{z^{2}}$ sur le disque unité.
 \item Montrer que pour tout $z\in\Cc$, $2i\sin(z)=e^{-iz}(e^{2iz}-1)$. En déduire les zéros de la fonction sinus sur $\Cc$
 et en particulier que la fonction sinus ne s'annule que pour des valeurs réelles.
 \item Résoudre dans $\Cc$, les équations suivantes : $ e^{z} = -5 $ ; $ \sin(z) = 2 $.
\end{enumerate}
\finenonce{007549}
\finexercice


\exercice{7550, mourougane, 2021/08/10}
\enonce{007550}{Logarithme}
On considère la branche principale du logarithme, toujours notée $\log$.
\begin{enumerate}
 \item Calculer $\log(i)$.
 \item Calculer $i^i$.
 \item Démontrer que pour tout $\alpha$ et $\beta$ fixés dans $\Cc$ et $z\in\Cc^-$, 
 $$(z^\alpha)'=\alpha z^{\alpha-1}\quad \text{ et } \quad z^{\alpha+\beta}=z^\alpha z^\beta.$$
\item En utilisant la détermination principale du logarithme, on définit les fonctions suivantes : 
\begin{enumerate}
\item [a.]$ z \longmapsto z^{1/2} $
\item [b.]$ z \longmapsto( 1-z) ^{1/3} $.
\end{enumerate}
Déterminer leurs domaines de définition. 
\end{enumerate}
\finenonce{007550}
\finexercice


\exercice{7551, mourougane, 2021/08/10}
\enonce{007551}{Une autre application logarithme ?}
On considère l'application
$$\begin{array}{cccc}
 L :& \Cc-\{z\in\Cc, re(z)=0\}&\longrightarrow&\Cc\\
 &z=x+iy&\longmapsto& \frac{1}{2}\log_\Rr(x^2+y^2)+i\arctan (\frac{y}{x}).
\end{array}$$
\begin{enumerate}
 \item Montrer que $L$ est holomorphe.
 \item L'application $L$ coïncide-t-elle avec la branche principale $\log$ du logarithme sur\\ $\{z\in\Cc, re(z)>0\}$ ?
 \item L'application $L$ est-elle un logarithme sur $\{z\in\Cc, re(z)<0\}$ ? 
\end{enumerate}
\finenonce{007551}
\finexercice


\exercice{7552, mourougane, 2021/08/10}
\enonce{007552}{Propriétés de la branche principale du logarithme}
\begin{enumerate}
\item Pour quelles valeurs de $z\in \Cc$ le nombre complexe $\exp (z)$ est-il dans $\Cc-\Cc^-$ ?
\item Vérifier que $\exp : B\to \Cc^-$ et $\log : \Cc^-\to B$ sont deux biholomorphismes réciproques.
 \item A-t-on pour tous $z,w\in\Cc^-$, $$\log (zw)=\log z +\log w ?$$ 
 \item A-t-on pour tous $z,w\in\Cc$ tels que $re (z)>0$ et $re (w)>0$, $$\log (zw)=\log z +\log w ?$$ 
\end{enumerate}
\finenonce{007552}
\finexercice


\exercice{7553, mourougane, 2021/08/10}
\enonce{007553}{La fonction zeta de Riemann}
On rappelle que pour tout entier naturel $n$, et tout nombre complexe $z$, $n^z$ désigne $p_z(n)=\exp(z\log n)$.
\begin{enumerate}
\item Déterminer le module de $n^z$.
\item Montrer que la série $\sum_{n\geq 1} \frac{1}{n^z}$ converge normalement sur l'ouvert $\{z\in\Cc / re(z)>1\}$.
\end{enumerate}
\finenonce{007553}
\finexercice

\section{442.00 Formule de Cauchy}


\exercice{2806, burnol, 2009/12/15}
\enonce{002806}{}
Le Laplacien $\Delta = \frac{\partial^2}{\partial x^2}
+\frac{\partial^2}{\partial y^2}$ est un opérateur
différentiel qui joue un rôle important en analyse
complexe. Soit $f(z) = u(x,y) + i\, v(x,y)$ une fonction holomorphe sur
un ouvert du plan complexe. On sait que $f$, donc $u$ et
$v$, admettent des dérivées partielles de tous les
ordres. En
utilisant les équations de 
Cauchy-Riemann, montrer que $u$ et $v$ vérifient l'équation de Laplace:
\[ \frac{\partial^2 u}{\partial x^2}
+\frac{\partial^2 u}{\partial y^2} =0 \qquad
\frac{\partial^2 v}{\partial x^2} +\frac{\partial^2
v}{\partial y^2} =0\] On dit d'une fonction vérifiant
l'équation de Laplace qu'elle est harmonique. La fonction
holomorphe $f = u + i v$ est aussi une fonction harmonique
puisque $\Delta(f) = \Delta(u) + i \Delta(v) = 0$.
\finenonce{002806}
\finexercice


\exercice{2807, burnol, 2009/12/15}
\enonce{002807}{}
\label{exo:CRpolaires}
  On veut exprimer les équations de Cauchy-Riemann avec les coordonnées
  polaires $r$ et $\theta$. Les équations de Cauchy-Riemann
 peuvent s'écrire sous la forme:
\[  \left(\frac\partial{\partial x} +
 i\frac\partial{\partial y}\right) F = 0\]
donc il s'agit d'exprimer $\frac\partial{\partial x}$ et
$\frac\partial{\partial y}$ en fonction de
$\frac\partial{\partial r}$ et de $\frac\partial{\partial
\theta}$. Lorsque l'on travaille sur un ouvert (ne
 contenant pas l'origine) sur lequel une détermination
 continue de l'argument $\theta$ est possible (par exemple
 sur $\Omega = \Cc\setminus]-\infty,0]$). Montrer:
\begin{align*}
 \frac\partial{\partial r} &=
 \cos(\theta)\frac\partial{\partial x} +
 \sin(\theta)\frac\partial{\partial y}\\
 \frac\partial{\partial\theta} &= -r\sin(\theta)\frac\partial{\partial x} +
 r\cos(\theta)\frac\partial{\partial y}
\end{align*}
En déduire $\frac\partial{\partial x} =
 \cos(\theta)\frac\partial{\partial r} - \sin(\theta)\frac1r
 \frac\partial{\partial\theta}$ et $\frac\partial{\partial y} =
 \sin(\theta)\frac\partial{\partial r} + \cos(\theta)\frac1r
 \frac\partial{\partial\theta}$. Montrer alors:
\[ \frac\partial{\partial x} +
 i\frac\partial{\partial y} = e^{i\theta}\left(\frac\partial{\partial r} +
 i\frac1 r\frac\partial{\partial \theta}\right) =
 e^{i\theta}\frac1 r\left(r\frac\partial{\partial r} + 
 i\frac\partial{\partial \theta}\right)\]
En déduire qu'en coordonnées polaires les équations de
 Cauchy-Riemann peuvent s'écrire (en particulier) sous la forme:
\[ \frac{\partial F}{\partial\theta} = i\, r\frac{\partial
 F}{\partial r}\]
\finenonce{002807}
\finexercice


\exercice{2808, burnol, 2009/12/15}
\enonce{002808}{}
  Il est intéressant que l'équation de l'exercice précédent
 $\frac{\partial F}{\partial\theta} = i\, r\frac{\partial 
 F}{\partial r}$, peut se réécrire dans le système de coordonnées $(a,b) =
 (\log(r),\theta)$ sous la forme: 
\[ \frac{\partial F}{\partial b} = i\, \frac{\partial
 F}{\partial a}\;,\]
autrement dit exactement sous la même forme qu'ont les
 équations de Cauchy-Riemann originelles dans les
 coordonnées cartésiennes $(x,y)$.\footnote{$\frac{\partial
F}{\partial y} = i \frac{\partial
 F}{\partial x}$, ou, plus mnémotechnique: $\frac{\phantom{i}\,\partial
F}{i\,\partial y} = \frac{\partial
 F}{\partial x}$ qui dit ``holomorphe $\Leftrightarrow$ $iy$
 est comme $x$''.} 
 Or $a$ et $b$ sont les parties réelles et imaginaires de la
 combinaison $a+ib$ qui est holomorphe comme fonction de
 $x+iy$: $a+ib = \log(x+iy)$. Montrer que cela est général:
 dans un système de coordonnées $(a,b)$ telles que $w= a+ib$
 est une fonction holomorphe de $z = x+iy$ les équations de
 Cauchy-Riemann pour l'holomorphie (par rapport à $(x,y)$)
 d'une fonction $F$ sont $\frac{\partial F}{\partial
 b} = i \frac{\partial F}{\partial a}$ (ce qui équivaut à
 l'holomorphie de $F$ comme fonction ``sur le plan de
 $w=a+ib$''\footnote{autrement dit pour qu'une fonction
 soit holomorphe comme fonction de $x+iy$ il est nécessaire
 et suffisant qu'elle soit holomorphe comme fonction de
 $a+ib$. En particulier $x+iy$ est une fonction holomorphe
 de $a+ib$: on a donc prouvé que la réciproque d'une
 bijection holomorphe est aussi holomorphe. Nous reviendrons
 là-dessus avec d'autres méthodes (dont celle très concrète de
 l'``inversion'' d'une série entière).}). Indication: prouver
 l'identité :
\[ \frac\partial{\partial x} +
 i\frac\partial{\partial y} = \left(\frac{\partial
 a}{\partial x} - i \frac{\partial
 b}{\partial x}\right)\left(\frac\partial{\partial a} +
 i\frac\partial{\partial b}\right)\;,\]
en exploitant les équations de Cauchy-Riemann
$\frac{\partial
 b}{\partial x} = -\frac{\partial a}{\partial y}$,
 $\frac{\partial a}{\partial x} = +\frac{\partial
 b}{\partial y}$ pour $a+ib = g(x+iy)$.
\finenonce{002808}
\finexercice


\exercice{2809, burnol, 2009/12/15}
\enonce{002809}{}
  On veut exprimer le Laplacien avec les coordonnées
  polaires $r$ et $\theta$: autrement dit pour toute
  fonction deux fois différentiable $\Phi$ on veut calculer
  la fonction $\Delta(\Phi)$ à l'aide des opérateurs de
  dérivées partielles $\frac\partial{\partial r}$ et
  $\frac\partial{\partial
 \theta}$, lorsque l'on travaille sur un ouvert (ne
 contenant pas l'origine) dans lequel une détermination
 continue de l'argument $\theta$ est possible (par exemple
 sur $\Omega = \Cc\setminus]-\infty,0]$). Une méthode
 possible est d'utiliser les expressions obtenues dans
 l'exercice \ref{exo:CRpolaires} :
\[ \frac\partial{\partial x} =
 \cos(\theta)\frac\partial{\partial r} - \sin(\theta)\frac1r
 \frac\partial{\partial\theta}\qquad \frac\partial{\partial y} =
 \sin(\theta)\frac\partial{\partial r} + \cos(\theta)\frac1r
 \frac\partial{\partial\theta}\;,\]
et de calculer ensuite $(\frac\partial{\partial x})^2$ et
 $(\frac\partial{\partial y})^2$ puis de faire la
 somme. Mais cela donne des calculs un peu longs. Voici une
 ruse: en reprenant une formule déjà établie dans
 l'exercice \ref{exo:CRpolaires} montrer
\[ (x-iy)\left(\frac\partial{\partial x} +
 i\frac\partial{\partial y}\right) = r\frac\partial{\partial r} + 
 i\frac\partial{\partial \theta}\]
\[ (x+iy)\left(\frac\partial{\partial x} -
 i\frac\partial{\partial y}\right) = r\frac\partial{\partial r} - 
 i\frac\partial{\partial \theta}\]
On remarquera maintenant que l'opérateur différentiel
 $\frac\partial{\partial x} + 
 i\frac\partial{\partial y}$ appliqué à la fonction $x+iy$
 donne zéro. Donc (expliquer!):
\[ (x-iy)\left(\frac\partial{\partial x} +
 i\frac\partial{\partial y}\right)(x+iy)\left(\frac\partial{\partial x} -
 i\frac\partial{\partial y}\right) =
 (x-iy)(x+iy)\left(\frac\partial{\partial x} + 
 i\frac\partial{\partial y}\right)\left(\frac\partial{\partial x} -
 i\frac\partial{\partial y}\right)\]
Prouver alors en conclusion:
\[ \frac{\partial^2\Phi}{\partial x^2}
+\frac{\partial^2\Phi}{\partial y^2} = \frac{1}{r^2}\left(
 (r\frac\partial{\partial r})^2 + \frac{\partial^2}{\partial
 \theta^2}\right)\Phi
= \frac{\partial^2\Phi}{\partial r^2} + \frac1r
 \frac{\partial\Phi}{\partial
 r} +  \frac1{r^2}\frac{\partial^2\Phi}{\partial 
 \theta^2}.\]
\finenonce{002809}
\finexercice


\exercice{2810, burnol, 2009/12/15}
\enonce{002810}{}
\label{ex:burnol2.2.1}
  Soit $\gamma = [A,B]+[B,C]+[C,D]+[D,A]$ le bord (parcouru
 dans le sens direct) du carré de sommets $A = 1-i$, $B
 = 1+i$, $C=-1+i$, $D=-1-i$. Déterminer les intégrales suivantes:
\begin{enumerate}
\item $\int_\gamma \,dx$, $\int_\gamma x\,dx$, $\int_\gamma
  x^2\,dx$, $\int_\gamma y\,dx$, $\int_\gamma
  y^2\,dx$, $\int_\gamma
  y^3 \,dx$,
\item $\int_\gamma x\,dx+y\,dy$,  $\int_\gamma x\,dy+y\,dx$,
 $\int_\gamma x\,dy-y\,dx$, 
\item $\int_\gamma \,dz$, $\int_\gamma z\,dz$, 
$\int_\gamma
 x\,dz$, $\int_\gamma z\,dx$, 
\item  $\int_\gamma
  z^{-1} \,dz$, $\int_\gamma
  z^{-2} \,dz$, $\int_\gamma
  z^{n} \,dz$, pour $n\in\Zz$. 
\end{enumerate}
\finenonce{002810}
\finexercice


\exercice{2811, burnol, 2009/12/15}
\enonce{002811}{}
Avec les mêmes notations on veut évaluer $\int_\gamma
  \overline{z^{n}} \,dz$, $n\in\Zz$. Justifier les étapes
  suivantes:
\begin{align*}
 \int_\gamma
  \overline{z^{n}} \,dz &= \overline{\int_\gamma z^n \,\overline{dz}}\\
\int_\gamma z^n \,\overline{dz}&= \int_{[B,C]}  z^n \,dz - \int_{[C,D]}
  z^n  \,dz + \int_{[D,A]} z^n \,dz - \int_{[A,B]}z^n \,dz\;,
\end{align*}
et compléter le calcul, pour tout $n\in\Zz$.
\finenonce{002811}
\finexercice


\exercice{2812, burnol, 2009/12/15}
\enonce{002812}{}
\label{ex:burnol2.2.3}
On note $C$ le cercle de rayon $1$ parcouru dans le sens
direct. Calculer $ \int_C z^n \,dz$ et  $\int_\gamma z^n
\,dz$ pour tout $n\in\Zz$, et vérifier qu'il y a toujours
égalité (ici $\gamma = \partial\mathcal{R}$ est à nouveau le bord du
carré qui a été 
 utilisé dans les exercices précédents). Calculer $\int_C
 \overline{z}^n \,dz$ et
$\int_\gamma \overline{z}^n \,dz$ et trouver les cas
 d'égalités et
d'inégalités.
\finenonce{002812}
\finexercice


\exercice{2813, burnol, 2009/12/15}
\enonce{002813}{}
Soit $C$ un cercle de centre quelconque, parcouru dans le
sens direct, et ne passant pas par l'origine. Calculer
$\int_C z^n \,dz$ pour tout $n\in\Zz$ dans le cas où $C$
encercle l'origine, et dans le cas où $C$ n'encercle pas
l'origine. 
\emph{Indication} pour $n=-1$: soit $w$ l'affixe du
 centre du cercle, et $R$ son rayon. Paramétrer le cercle
 par $z = w(1 + \frac{R}{|w|}e^{i\theta})$,
 $-\pi<\theta\leq+\pi$, puis utiliser un développement en
 série en distinguant les cas $R>|w|$ et $R<|w|$. Ou encore
 invoquer la fonction $\mathrm{Log}(z/w)$.
\finenonce{002813}
\finexercice


\exercice{2814, burnol, 2009/12/15}
\enonce{002814}{}
Soit $0<a<b$ sur l'axe réel positif et soit $C=\{|z|=r\}$ le
cercle de rayon $r$ centré en l'origine, parcouru dans le
sens direct. Montrer:
\[ \int_C \frac1{(z-a)(z-b)}\,dz = 2\pi i \begin{cases}
0& r<a\\
\frac1{a-b}& a<r<b\\
0& r>b
\end{cases}\]
On pourra réduire la fraction en élément simples, puis se
 ramener au résultat de l'exercice précédent. Ou encore, on
 pourra envisager des développements en
séries, pour se ramener par étapes aux intégrales $\int_C
z^n dz$, $n\in\Zz$.
\finenonce{002814}
\finexercice


\exercice{2815, burnol, 2009/12/15}
\enonce{002815}{}
Soit $C$ le cercle unité parcouru dans le sens
direct. Calculer 
\[ \int_C \left( z + \frac1z \right)^{n}\frac{dz}z\quad(
n\in\Nn)\] en développant par la formule du binôme et en
utilisant les valeurs connues de $\int_C z^k dz$,
$k\in\Zz$. En déduire $\int_{-\pi}^{+\pi} \cos^{n}
t\,dt$. En déduire la valeur de $\int_{0}^{\pi/2} \cos^{n}
t\,dt$ pour $n$ pair:
\[ I_m = \int_0^{\frac\pi2} \cos^{2m} t\,dt =
\frac{1.3.\cdots.(2m-1)}{2.4.\cdots.(2m)}\frac\pi2\] 
\finenonce{002815}
\finexercice


\exercice{2816, burnol, 2009/12/15}
\enonce{002816}{}
  On pose $J_m = \int_0^{\pi/2} \cos^{2m+1} t\,dt$, pour
  $m\in\Nn$. En intégrant par parties $J_{m+1}$ obtenir la relation de
  récurrence $J_{m+1} = \frac{2m+2}{2m+3}J_m$ et
  prouver:\footnote{par convention lorsque qu'un produit porte sur un
  ensemble vide il vaut $1$. Donc la formule est bien
  compatible avec $J_0 = 1$.}
\[ J_m = \frac{2.4.\cdots.(2m)}{3.5.\cdots.(2m+1)}\]
\finenonce{002816}
\finexercice


\exercice{2817, burnol, 2009/12/15}
\enonce{002817}{}
  En utilisant $I_{m+1}\leq J_m\leq I_{m}$, obtenir:
\[ \frac{2m+1}{2m+2}
  \frac{(1.3.\cdots.(2m-1))(3.5.\cdots.(2m+1))}{(2.4.\cdots.(2m))^2}\leq
  \frac2\pi \leq
  \frac{(1.3.\cdots.(2m-1))(3.5.\cdots.(2m+1))}{(2.4.\cdots.(2m))^2}\] 
En déduire la formule de Wallis:
\[ \frac2\pi = \lim_{m\to\infty}
  \frac{1.3}{2.2}\frac{3.5}{4.4}\cdots\frac{(2m-1).(2m+1)}{(2m).(2m)}
  = \prod_{m=1}^\infty \left( 1 - \frac1{4m^2}\right)\]
\finenonce{002817}
\finexercice


\exercice{2818, burnol, 2009/12/15}
\enonce{002818}{}
  Justifier le réarrangement suivant (qui découle aussi du
  terme de gauche dans l'inégalité de l'exercice précédent):
$ \frac2\pi = \frac12 \frac{3.3.5.5.\dots}{2.4.4.6.\dots} =
\frac12\prod_{m=1}^\infty \left( 1 -
\frac1{(2m+1)^2}\right)^{-1}$, soit
encore:
\[ \frac\pi4 = \prod_{m=1}^\infty \left( 1 -
\frac1{(2m+1)^2}\right)\;.\]
\finenonce{002818}
\finexercice


\exercice{2819, burnol, 2009/12/15}
\enonce{002819}{}
Justifier également sur la base des formules précédentes les
équivalents asymptotiques: 
$$\binom{2m}{m} \sim_{m\to\infty} \frac{2^{2m}}{\sqrt{\pi m}}$$
$$\frac{(1+\frac12)(2+\frac12)\cdots(m+\frac12)}{1.2.\cdots.m}\sim 2\sqrt{\frac{m}{\pi}}$$
$$\frac{(\frac12)_m}{m!}\sim \frac1{\sqrt{\pi m}}$$
\finenonce{002819}
\finexercice

\exercice{6669, queffelec, 2011/10/16}
\enonce{006669}{}
Soit $f(z)=\sum_{n\geq0} a_n z^n$ la somme d'une série entière de rayon infini.
Montrer que pour
$r>0$ et $n\geq 0$  
$$a_n={1\over 2\pi r^n}\int_0^{2\pi}f(re^{it})\ e^{-int}\ dt$$.

En déduire que si $|f(z)|\leq A+B|z|^k$ pour tout $z$ de module $\geq R$, $f$
est un polyn\^ome.
\finenonce{006669}



\finexercice


\exercice{6670, queffelec, 2011/10/16}
\enonce{006670}{}
On se propose de calculer $I=\int_{-\pi}^\pi \ln|re^{i\theta}-a|\ d\theta$, où
$0<r<|a|$.
\begin{enumerate}
\item Vérifier que $I$ est bien définie.

\item On considère $f$ la fonction définie sur $\Omega={\Cc}\backslash {\Rr^-}$
par $f(z)={1\over z}\hbox{ Log}(1-{z\over a})$, où $\hbox {Log}$ est la
détermination principale du logarithme sur $\Omega$. Montrer que $f$ est
holomorphe sur un ouvert contenant le cercle $\{z=re^{i\theta}\}$.

\item En déduire $I=2\pi \ln|a|$.
\end{enumerate}
\finenonce{006670}



\finexercice


\exercice{6671, queffelec, 2011/10/16}
\enonce{006671}{}
Soit $\lambda\in\Cc$ de module $\not=0,1$, et $I=\displaystyle\int_0^{2\pi}
{\cos(n\theta)\ d\theta\over \lambda^2-2\lambda\cos \theta+1}$.
\begin{enumerate}
\item Vérifier que $I$ est bien définie.

\item On considère $f$ définie par
$f(z)=\displaystyle{z^n\over (z-\lambda)(z-\lambda^{-1})}$; en calculant
l'intégrale de
$f$ sur le cercle unité, trouver la valeur de $I$ (distinguer les deux cas
$|\lambda|>1, |\lambda|<1$).
\end{enumerate}
\finenonce{006671}



\finexercice


\exercice{6672, queffelec, 2011/10/16}
\enonce{006672}{}
Soit $U$ un ouvert de $\Cc$ contenant $\bar D$ et $f\in H(U)$. On note
 $\gamma$ le paramétrage de $\partial D$ par $t\in [0,2\pi]\to e^{it}$.
\begin{enumerate}
\item Calculer $I_1=\displaystyle\int_{\gamma} (2+z+{1\over z}){f(z)\over z}\ dz$ et
$I_2=\displaystyle\int_{\gamma} (2-z-{1\over z}){f(z)\over z}\ dz$.

\item En déduire la valeur de $\displaystyle\int_0^{2\pi}
f(e^{i\theta})\cos^2({\theta\over2})\ d\theta$ et $\displaystyle\int_0^{2\pi}
f(e^{i\theta})\sin^2({\theta\over2})\ d\theta$.

\item Pour $|a|\not=1$, évaluer $I(a)=\displaystyle{1\over 2i\pi}\int_{\gamma}
{\overline{f(z)}\over z-a}\ dz.$
\end{enumerate}
\finenonce{006672}



\finexercice


\exercice{6673, queffelec, 2011/10/16}
\enonce{006673}{}
\begin{enumerate}
\item Soit $\Omega $ un ouvert de $\C$, $F$ et $G$ deux fonctions
holomorphes dans $\Omega $ et $\gamma \colon [a,b]\to\C$ un chemin tel
 que $\gamma ^*\subset \Omega $. Montrer que 
$$\int_\gamma F(z)G'(z)\ dz=F(\gamma (b))G(\gamma (b))-
F(\gamma (a))G(\gamma (a))-\int_\gamma F'(z)G(z)\ dz$$

\item Calculer $\int_\gamma (z+2)e^{iz}\ dz$, où $\gamma $ est l'arc de
parabole $\displaystyle\gamma (t)=t+i{t^2\over\pi^2}$ joignant (0,0) à $(\pi,1)$.
\end{enumerate}
\finenonce{006673}



\finexercice       


\exercice{6674, queffelec, 2011/10/16}
\enonce{006674}{} 
En évaluant $\int_Ce^z\ dz$ sur le cercle unité, montrer que
$$\int_0^{2\pi}e^{\cos\theta }\cos{(\theta +\sin\theta )}\ d\theta =
\int_0^{2\pi}e^{\cos\theta }\sin{(\theta +\sin\theta )}\ d\theta =
0$$
\finenonce{006674}



\finexercice       


\exercice{6675, queffelec, 2011/10/16}
\enonce{006675}{} 
Calculer 
$$\int_\gamma \left( z+{1\over z}\right)^n \ {dz\over z},$$
où $\gamma (t)=e^{it}$ ($t\in [0,2\pi]$) et $n\in {\Nn}$. 
En déduire la valeur de $\int_0^{2\pi} \cos^n t\ dt$.
\finenonce{006675}



\finexercice       


\exercice{6676, queffelec, 2011/10/16}
\enonce{006676}{} 
Calculer les intégrales suivantes.
\begin{enumerate}
\item $\displaystyle\int_\gamma {\cos z\over z}\ dz$ où $\gamma (t)=e^{it}$ ($t\in [0,2\pi]$)
\item $\displaystyle\int_\gamma {\sin z\over z}\ dz$ où $\gamma (t)=e^{it}$ ($t\in [0,2\pi]$)
\item $\displaystyle\int_\gamma {\cos {z^2}\over z}\ dz$ où $\gamma (t)=e^{it}$ ($t\in [0,2\pi]$)
\item $\displaystyle\int_\gamma {e^{\pi z}\over z^3+z}\ dz$ où
$\gamma (t) =2e^{it}$ ($t\in [0,2\pi]$)
\item $\displaystyle\int_\gamma {dz\over(z-a)^n}$ ($n\in{\Zz}$), où $\gamma $
est un chemin fermé ne passant pas par $a$
\item $\displaystyle\int_{\gamma _r}{3z^2-12z+11\over z^3-6z^2+11z-6}\ dz$,
où $\gamma _r(t)=re^{it}$ ($t\in [0,2\pi]$)
\item $\displaystyle\int_{\gamma }\left({z\over z-1}\right)^n\ dz$, 
$n\in{\Nn}^*$,
où $\gamma (t)=1+e^{it}$ ($t\in [0,2\pi]$)
\end{enumerate}
\finenonce{006676}



\finexercice       


\exercice{6677, queffelec, 2011/10/16}
\enonce{006677}{} 
Soit $I(z)=\int_0^{+\infty}{\ln t\over t^2+z^2}dt$.
\begin{enumerate}
\item Pour quelles valeurs de $z$ $I(z)$ est-elle définie ?

\item Montrer que pour $\Re z>0$, on a
$$I(z)={1\over z}\left( {\pi\over 2}\mathrm{Log} z+
\int_0^{+\infty}{\ln t\over 1+t^2}dt\right) ,$$
où $\mathrm{Log}$ est la détermination principale du logarithme. On pourra
considérer le chemin fermé $\Gamma_{\varepsilon, R}=[\varepsilon,
R]+\gamma_R+[Re^{i\varphi},\varepsilon
e^{i\varphi}]-\gamma_{\varepsilon}$, où $\varphi=\mathrm{Arg} z$ et
$\gamma_r\colon t\mapsto re^{it}$, $t\in [0,\varphi]$.

\item Qu'obtient-on pour $\Re z<0$ ?
\end{enumerate}
\finenonce{006677}



\finexercice       


\exercice{6678, queffelec, 2011/10/16}
\enonce{006678}{} 
Soit $a>0$ et $\gamma\colon t\mapsto a+it$, $t\in {\Rr}$. Pour
$\alpha\in {\Rr}$, on considère
$$J(\alpha)={1\over 2i\pi}\int_\gamma {e^{\alpha z}\over z^2}dz.$$
\begin{enumerate}
\item Montrer que cette intégrale converge.

\item  En considérant les chemins fermés $[a-iR,a+iR]+\gamma_R$, où
$\gamma_R^*$ est un demi-cercle de diamètre $[a-iR,a+iR]$, montrer
que $J(\alpha)=0$ si $\alpha\le 0$ et $j(\alpha)=\alpha$ si $\alpha\ge
0$.
\end{enumerate}
\finenonce{006678}



\finexercice       


\exercice{6679, queffelec, 2011/10/16}
\enonce{006679}{} Pour $n\ge 1$ et $0\le k\le n$, on désigne par $C_n^k$ le
coefficient binomial. Pour $r>0$, soit $c_r\colon t\mapsto re^{it}$, $t\in
[0,2\pi]$.

\begin{enumerate}  
\item Montrer que 
$$C_n^k={1\over 2i\pi}\int_{c_r}(1+z)^n{dz\over z^{k+1}}.$$
En déduire que $C_{2n}^n\le 4^n$.

\item Montrer que 
$$C_{2n}^n={1\over 2i\pi}\int_{c_r}\left( {1\over z}
+2+z\right) ^n{dz\over z}$$
(on pourra utiliser 1.). En déduire que 
$$\sum_{n=0}^{+\infty}C_{2n}^n{1\over 5^n}={1\over
2i\pi}\int_{c_r}{dz\over 3z-1-z^2)}=\sqrt 5 ,$$
à condition que $r_1<r<r_2$, où $r_1<r_2$ sont les deux racines de
$3z-1-z^2=0$.

\item Montrer que
$${1\over 2i\pi}\int_{c_r}(1+z)^n\left( 1+{1\over z}\right)^n{dz\over
z}=\sum_{k=0}^n\left( C_n^k\right) ^2.$$
En déduire que 
$$\sum_{k=0}^n\left( C_n^k\right) ^2=C_{2n}^n.$$

\item Montrer que
$${1\over 2i\pi}\int_{c_1}{(z-1)^{2n}(z+1)^n\over z^{2n+1}}dz=
\sum_{k=0}^n(-1)^kC_n^kC_{2n}^k.$$
Montrer que si $z\in c_1^*$, $\vert z-1\vert ^2\vert z+1\vert \le
{16\over 9}\sqrt 3$. En déduire que
$$\sum_{k=0}^n(-1)^kC_n^kC_{2n}^k\le \left( {16\over 9}\sqrt 3\right)
^n.$$
\end{enumerate}
\finenonce{006679}



\finexercice       


\exercice{6680, queffelec, 2011/10/16}
\enonce{006680}{} 
\begin{enumerate}  
\item Montrer que $\int_0^{2\pi} e^{re^{it}}\ dt$ est indépendant de
$r>0$.

\item Montrer que $\int_0^{2\pi} \mathrm{Log}{(1-re^{it})}\ dt$ est indépendant de
$r\in \mathopen]0,1\mathclose[$ ($\mathrm{Log}$ désigne la branche principale
du logarithme sur $\C\setminus \mathopen] -\infty,0]$).
\end{enumerate}
\finenonce{006680}



\finexercice       


\exercice{6681, queffelec, 2011/10/16}
\enonce{006681}{}
Calculer de deux manières $\int_\gamma {dz\over z}$ où
$\gamma (t)=a\cos t+ib\sin t$ ($t\in [0,2\pi]$, $a,b\in {\Rr}_+^*$) et en 
déduire la valeur de l'intégrale
$$\int_0^{2\pi}{dt\over a^2\cos ^2t+b^2\sin ^2t}$$
\finenonce{006681}



\finexercice       


\exercice{6682, queffelec, 2011/10/16}
\enonce{006682}{} 
Soient $f$ et $g$ deux fonctions holomorphes sur un ouvert connexe
$\Omega $ contenant le disque unité fermé et $\gamma (t) =e^{it}$ ($t\in
[0,2\pi]$). Montrer que 
$${1\over 2i\pi}\int_\gamma \left({f(z)\over z-a}+{ag(z)\over az-1}\right)\ dz=
\left\{\begin{array}{ll}
f(a)\ &\ \mbox{ si }\vert a\vert <1\\
g(1/a)\ &\ \mbox{ si }\vert a\vert >1\\
\end{array}\right.$$
\finenonce{006682}



\finexercice       


\exercice{6683, queffelec, 2011/10/16}
\enonce{006683}{} 
Soit $f$ une fonction holomorphe sur un ouvert $\Omega $ contenant
$\overline{D(a,r)}$. Montrer que
$$f(a)={1\over 2\pi}\int_0^{2\pi}f(a+re^{i\theta })\ d\theta $$
(propriété de la moyenne). 
\finenonce{006683}



\finexercice       


\exercice{6684, queffelec, 2011/10/16}
\enonce{006684}{}
Soit $f$ une fonction continue dans le secteur
$$\{ z\in {\Cc}\vert\ -\alpha  \le \mathrm{Arg} z\le \alpha \}$$
On suppose que $zf(z)$ tend vers $A\in{\Cc}$ quand $\vert z\vert$ 
tend vers l'infini, $z$ restant dans ce secteur. Notons $C_R$ la partie du 
cercle de centre 0 et de rayon $R$ contenue dans ce secteur.
Montrer que
$$\lim_{R\to +\infty}\int_{C_R}f(z)dz=2i\alpha A$$
\finenonce{006684}



\finexercice      


\exercice{6685, queffelec, 2011/10/16}
\enonce{006685}{}
\begin{enumerate}  
\item Soit $C\subset{\Cc}$ une courbe orientée $C^1$ et fermée. Soit
$\gamma $ un chemin $C^1$ d'origine $a$ et d'extrémité $b$, tels que $a$ et $b$ ne
soient pas des points de $C$. On suppose que l'intersection de $C$ et
$\gamma $ est constituée d'un nombre fini de points $m_1,\dots,m_n$
et que les tangentes à $C$ et à $\gamma $ sont distinctes en ces points.
Soit $\varepsilon _i=1$ si l'angle de la tangente à $\gamma $ avec la
tangente à $C$ en $m_i$ est entre 0 et $\pi$, $\varepsilon _i=-1$
sinon. Montrer que
$$\sum_i \varepsilon _i=\mathrm{Ind}_C(a)-\mathrm{Ind}_C(b)$$
où $\mathrm{Ind}_C(z)$ désigne l'indice de $z$ par rapport à $C$.

\item Calculer l'indice du point $z={3\over 4}$ par rapport à la courbe dont l'équation en
coordonnées polaires est $r=\cos {\theta \over 3}$ avec $0\le \theta 
\le 3\pi$, parcourue dans le sens des $\theta $ croissants.
\end{enumerate}
\finenonce{006685}



\finexercice       


\exercice{6686, queffelec, 2011/10/16}
\enonce{006686}{}
On considère la série entière
$$L(z)=\sum_{n\ge 1}{z^n\over n^2}$$
Soit $f(z)={-\mathrm{Log}{(1-z)}\over z}$ où $\mathrm{Log}$ désigne la détermination
principale du logarithme complexe.

\begin{enumerate}  
\item On note $U={\Cc}\setminus [1,+\infty\mathclose[$. Montrer que $f$ est définie 
dans $U$.

\item Vérifier que si $D=\{ z\in {\Cc}\vert\ \vert z\vert<1\}$, 
on a $L'(z)=f(z)$.

\item En déduire qu'il existe une primitive de $f$, définie dans $U$ tout
entier, dont la restriction à $D$ est égale à $L$. On note cette primitive 
$L$ par abus de langage.

\item Soit $x\in {\Rr}$, $x>1$. Calculer $\lim_{y\to 0}L(x+iy)-L(x-iy)$
comme fonction de $x$.
\end{enumerate}
\finenonce{006686}



\finexercice       


\exercice{6687, queffelec, 2011/10/16}
\enonce{006687}{}
\begin{enumerate}  
\item Soit $P$ un polynôme qui ne s'annule pas sur le cercle $\vert z\vert=1$.
Montrer que le nombre de zéros de $P$ à l'intérieur du cercle unité est
$${1\over 2\pi}\left[ \mathrm{Arg}{P(e^{i\theta })}\right] _0^{2\pi}$$
(variation d'une détermination continue de l'argument sur le cercle
unité). On utilisera le théorème de d'Alembert pour factoriser
$P$ en facteurs de degré 1, puis on considèrera l'indice de chacune des
racines par rapport au cercle.

\item Soit $P$ un polynôme n'ayant aucun zéro sur le cercle $\vert z\vert=1$
et ayant exatement $k$ racines (comptées avec multiplicité) à l'intérieur du
cercle unité. Montrer que la fonction
$$\theta \mapsto\Re P(e^{i\theta })$$
s'annule au moins $2k$ fois pour $\theta \in [0,2\pi]$ (indication : étudier
les zéros de la fonction $\cos{\mathrm{Arg}{P(e^{i\theta })}}$, où $\mathrm{Arg}$ est
une détermination continue de l'argument).
\end{enumerate}
\finenonce{006687}



\finexercice       


\exercice{6688, queffelec, 2011/10/16}
\enonce{006688}{} 
\begin{enumerate}  
\item Montrer qu'il existe sur $\C\setminus [-1,1]$ une détermination
holomorphe de $(z^2-1)^{-1/2}$ qui prend la valeur 1 pour $z=\sqrt 2$.
Unicité ? On désignera par $f$ cette détermination.

\item Pour $z\in \C\setminus{\Rr}$, on désigne par $\gamma _z$ le
segment $\mathopen] 0,z]$ orienté de 0 vers $z$ et on pose
$$F(z)=\int_{\gamma _z}f(\zeta )\ d\zeta .$$

\begin{enumerate}  
\item Montrer que $F$ est holomorphe sur $\C\setminus {\Rr}$ (quelle en
est la dérivée $F'$ ?).

\item Etudier $\lim_{z'\to z}F(z')$ lorsque $z\in {\Rr}\setminus [-1,1]$.

\item En déduire que $f$ n'a pas de primitive sur $\C\setminus [-1,1]$.
\end{enumerate}
\end{enumerate}
\finenonce{006688}



\finexercice       


\exercice{6689, queffelec, 2011/10/16}
\enonce{006689}{}
\begin{enumerate}
\item Soit $f$ une fonction entière telle que $|f(z)|\leq 1+e^{|z|}\sin|z|$ pour
tout 
$z$. Montrer que $f$ est une constante.

\item Soit $f\in H(D)$ telle que $|f(z)|(1-|z|)\leq 1$ pour $z\in D$. Montrer que
pour tout $n$ $|a_n|< e(n+1)$.
\end{enumerate}
\finenonce{006689}



\finexercice


\exercice{6690, queffelec, 2011/10/16}
\enonce{006690}{}
Soit $f$ une fonction holomorphe dans $\Cc$ tout entier et soit
$\gamma_R$ un paramétrage du cercle $C(0,R)$, $R>0$. Calculer pour $|z|<R$ :
${z\over 2i\pi}\int_{\gamma_R}{f(w)\over w(w-z)}\ dw$; en déduire que si
$\sup_R\int_0^{2\pi}|f(Re^{it})|\ dt<\infty$,
$f$  est constante. Quel théorème retrouve-t-on ?
\finenonce{006690}



\finexercice


\exercice{6691, queffelec, 2011/10/16}
\enonce{006691}{}
 Soit $f$ un
polyn\^ome de degré $n>0$, supposé sans zéros dans $\Cc$, et $\gamma_R$ le
cercle centré en
$0$, de rayon $R$. Calculer $I_R={1\over 2i\pi}\int_{\gamma_R}
{nf(z)-zf'(z)\over zf(z)}\ dz$, et trouver sa limite lorsque $R\to+\infty$.
Quel théorème retrouve-t-on ?
\finenonce{006691}



\finexercice


\exercice{6692, queffelec, 2011/10/16}
\enonce{006692}{}
 Soit $f$ une fonction entière telle que $|f(z)|\to+\infty$ quand
$|z|\to+\infty$. 
\begin{enumerate}
\item Montrer que $f$ n'a qu'un nombre fini de zéros dans $\Cc$ notés $z_1,\ldots
z_k$.

\item En déduire que $f$ est un polyn\^ome.
(Pour cela considérer  $f(z)/P(z)=g(z)$ où $P(z)=(z-z_1)\cdots (z-z_k)$ et
montrer que
$1/g$ est entière) . 
\end{enumerate}
\finenonce{006692}



\finexercice


\exercice{6693, queffelec, 2011/10/16}
\enonce{006693}{}
Soit $f$ une fonction holomorphe dans un ouvert connexe $\Omega$ contenant
$0$. 
Montrer que  :
\begin{enumerate}
\item Si $f({1\over n})={1\over n+1}$ pour $n$ assez grand alors $f(z)={z\over
z+1}$ sur $\Omega\cap D(0,1)$.

\item Si $f({1\over n})=f({1\over 2n})$ pour $n$ assez grand alors $f$ est 
constante sur $\Omega$.

\item Si $f({1\over n})=f({1\over n+1})$ pour $n$ assez grand alors $f$ est
constante sur $\Omega$.

\item $f({1\over n})=2^{-n}$ pour $n$ assez grand est impossible.
\end{enumerate}
\finenonce{006693}



\finexercice


\exercice{6694, queffelec, 2011/10/16}
\enonce{006694}{}
On considère la série entière $\sum_0^\infty 2^nz^n$. Calculer sa somme $f$ dans
son disque de convergence.
Trouver le plus grand ouvert connexe de $\Cc$ sur lequel $f$ se prolonge en
une fonction holomorphe. 
Donner le développement en série de $f$ au point $z=-1/4$ et le rayon de
convergence de cette série.
\finenonce{006694}



\finexercice


\exercice{6695, queffelec, 2011/10/16}
\enonce{006695}{}
Soit $f$ une fonction holomorphe dans un ouvert connexe $\Omega$ sur lequel elle
ne s'annule pas. Alors sont équivalentes :

(i) Il existe une détermination holomorphe du logarithme de $f$ sur $\Omega$.

ii) $\int_\gamma
{f'(z)\over f(z)}\ dz=0$ pour toute $\gamma$ une courbe fermée dans $\Omega$
de classe $C^1$ par morceaux. 

(iii)  ${f'\over f}$ admet une primitive sur $\Omega$.
\finenonce{006695}



\finexercice


\exercice{6696, queffelec, 2011/10/16}
\enonce{006696}{}
Soit $f$ une fonction holomorphe dans un ouvert connexe $\Omega$. On va montrer
l'équivalence entre 

(i) $f$ admet un logarithme holomorphe dans $\Omega$.

(ii)  $f$ admet des racines {\bf de tous ordres} holomorphes dans $\Omega$.

On a vu que (i) implique (ii).
Supposons maintenant que (ii) est vérifié : pour chaque $n$ on note $f_n$ la
fonction de $H(\Omega)$ telle que
$f_n^n(z)=f(z)$ si $z\in\Omega$.
\begin{enumerate}
\item Soit $a$ un zéro de $f$; que peut-on dire de la multiplicité de $a$ ? En
déduire que
$f$ ne s'annule pas sur
$\Omega$.

\item Soit $\gamma$ une courbe fermée dans $\Omega$
de classe $C^1$ par morceaux. On pose $I={1\over 2\pi i}\int_\gamma
{f'(z)\over f(z)}\ dz$, et $I_n={1\over 2\pi i}\int_\gamma
{f_n'(z)\over f_n(z)}\ dz$. Montrer que $I$ et $I_n$ sont des entiers, $I=nI_n$,
puis que
$I=0$. Conclure.
\end{enumerate}
\finenonce{006696}



\finexercice


\exercice{6697, queffelec, 2011/10/16}
\enonce{006697}{} 
Soit $f$ une fonction entière ; on pose, pour $r\in {\Rr}_+^*$,
$$M(r)=\sup_{\vert z\vert =r}\vert f(z)\vert$$

\begin{enumerate}  
\item On suppose qu'il existe $p\in {\Nn}$ tel que 
$$\lim_{r\to +\infty}{M(r)\over r^{p+1}}=0$$
Montrer qu'alors $f$ est un polynôme de degré au plus $p$.

\item On suppose qu'il existe $R\ge 0$, $K>0$ et $p\in {\Nn}$ tels que 
$$\vert z\vert  >R\Longrightarrow \vert f(z)\vert \le K\vert z\vert ^p$$
Montrer qu'alors $f$ est un polynôme de degré au plus $p$. Montrer que si
de plus $R=0$ alors $f$ est un monôme de degré $p$.

\item En déduire que si $f$ vérifie 
$$\forall z\in \C,\ \vert f'(z)\vert \le \vert z\vert$$
alors $f$ est de la forme $f(z)=a+bz^2$ avec $\vert b\vert \le 1/2$.
\end{enumerate}
\finenonce{006697}



\finexercice       


\exercice{6698, queffelec, 2011/10/16}
\enonce{006698}{} 
Soit $f(z)=\sum_{n\ge 0}a_nz^n$ une série entière de rayon de
convergence $R>0$. Pour $r<R$, soit $\gamma_r\colon t\mapsto re^{it}$,
$t\in [0,2\pi]$ et
$$I(r)={1\over 2i\pi}\int_{\gamma_r}\vert f(z)\vert ^2{dz\over z}.$$

\begin{enumerate}  
\item Montrer que $I(r)=\sum_{n=0}^{+\infty}\vert a_n\vert ^2 r^{2n}$.

\item En déduire une nouvelle démonstration des inégalités de Cauchy et
montrer que si $f$ n'est pas un monôme, ces inégalités sont strictes.

\item En considérant $f(z)=1/(1-z)^2$, montrer que pour $r<1$,
$${1\over 2\pi}\int_0^{2\pi}{dt\over (1-2r\cos
t+r^2)^2}=\sum_{n=1}^{+\infty}n^2r^{2n-2}.$$
\end{enumerate}
\finenonce{006698}



\finexercice       


\exercice{6699, queffelec, 2011/10/16}
\enonce{006699}{} 
Soit $f$ une fonction holomorphe dans le disque unité ouvert $D$
vérifiant 
$$\forall z\in D,\ \vert f(z)\vert <{1\over 1-\vert z\vert}.$$
Montrer que les coefficients $a_n$ du développement $f(z)=\sum_{n\ge
0}a_nz^n$ vérifient
$$\forall n\ge 1,\ \vert a_n\vert <(n+1)\left( 1+{1\over n}\right)
^n<(n+1)e
{\rm \ et \ } \ a_0\vert <e.$$
\finenonce{006699}



\finexercice       


\exercice{6700, queffelec, 2011/10/16}
\enonce{006700}{} 
Montrer qu'une fonction entière admettant 1 et $i$ comme périodes
est constante.
\finenonce{006700}



\finexercice       


\exercice{6701, queffelec, 2011/10/16}
\enonce{006701}{} 
Soit $f$ une fonction entière non constante ne s'annulant pas. 
Mon\-trer que 
$$\forall \varepsilon >0,\ \forall r>0,\ \exists z\in\C, \vert z\vert >r
\mbox{ et }\vert f(z)\vert <\varepsilon$$
Application : tout polynôme non constant admet un zéro dans $\C$ 
(théorème de d'Alembert).
\finenonce{006701}



\finexercice       


\exercice{6702, queffelec, 2011/10/16}
\enonce{006702}{} 
Déterminer toutes les fonctions $f$ holomorphes dans le disque
$D(0,R)=\{ z\in \C\vert \ \vert z\vert <R\}$ telles que 
$$\forall z\in D(0,{R\over 2}),\ f(z)=f(2z).$$
\finenonce{006702}



\finexercice       


\exercice{6703, queffelec, 2011/10/16}
\enonce{006703}{} 
On considère la série entière $f(z)=\sum_{n\ge 1} {z^n\over n}$.

\begin{enumerate}  
\item Montrer que $f$ admet un prolongement analytique sur un ouvert que
l'on déterminera.

\item Déterminer le développement en série entière de ce prolongement en
$i$. Quel est son rayon de convergence ?
\end{enumerate}
\finenonce{006703}



\finexercice       


\exercice{6704, queffelec, 2011/10/16}
\enonce{006704}{} 
Soit $F_1(z)=\int_0^{+\infty}e^{-(z+1)^2t}dt$.

\begin{enumerate}  
\item Déterminer l'ouvert $E$ sur lequel $F_1$ définit une fonction
holomorphe.

\item Montrer que $F_1$ se prolonge analytiquement en une fonction $F$
sur un ouvert que l'on déterminera. Calculer $F(2-4i)$. 
\end{enumerate}
\finenonce{006704}



\finexercice       


\exercice{6705, queffelec, 2011/10/16}
\enonce{006705}{} 
Existe-t-il des fonctions $f$ holomorphes sur $D(0,1)$ telles que
pour tout entier $n>0$

\begin{enumerate}  
\item $f\left( {1\over 2n}\right) = f\left( {1\over 2n+1}\right)={1\over n}$
?

\item $f\left( {1\over n}\right) = f\left( -{1\over n}\right)={1\over n^2}$
?

\item $f\left( {1\over n}\right) = f\left( -{1\over n}\right)=-{1\over n^3}$
?
\end{enumerate}
\finenonce{006705}



\finexercice       


\exercice{6706, queffelec, 2011/10/16}
\enonce{006706}{} 
Déterminer explicitement pour $a\in {\Rr}$
$$\max_{\vert z\vert \le 1}\vert z^2+2aiz+1\vert.$$
\finenonce{006706}



\finexercice       


\exercice{6707, queffelec, 2011/10/16}
\enonce{006707}{}
 Soient $P_1$, $P_2, \dots ,P_k$ $k$ points du plan. On désigne par
$d(A,B)$ la distance entre deux points $A$ et $B$ du plan. Soit $D$ un
ouvert connexe borné du plan. Montrer que le sup de $\prod_{1\le j\le
k}d(M,P_j)$ sur $D$ est atteint sur la frontière de $D$.
\finenonce{006707}



\finexercice       


\exercice{6708, queffelec, 2011/10/16}
\enonce{006708}{} 
Soit $f$ une fonction holomorphe sur un ouvert connexe $\Omega $ et ne
s'annulant pas sur $\Omega $. On suppose qu'il existe $a\in \Omega $ et
$\varepsilon >0$ tels que $\vert f(z)\vert\ge\vert f(a)\vert$ pour
$\vert z-a\vert <\varepsilon$. Montrer que $f$ est constante.
\finenonce{006708}



\finexercice      


\exercice{6709, queffelec, 2011/10/16}
\enonce{006709}{} 
Soit $\Omega $ un ouvert connexe borné de $\C$, $f$ une fonction
holomorphe sur $\Omega $, continue sur $\overline{\Omega }$, non
constante, telle que $\vert f\vert$ est constant sur la frontière de 
$\Omega $. Montrer que $f$ admet un zéro dans $\Omega $.
\finenonce{006709}



\finexercice       


\exercice{6710, queffelec, 2011/10/16}
\enonce{006710}{} 
Soit $\Omega $ un ouvert connexe de $\C$ contenant le disque
fermé  $\{ z\in\C\vert \ \vert z-z_0\vert \le r\}$, et $f$ une fonction
holomorphe sur $\Omega $ telle que $f(z)\in {\Rr}$ si
$\vert z-z_0\vert =r$. Montrer que $f$ est constante (considérer
$e^{if}$).
\finenonce{006710}



\finexercice       


\exercice{6711, queffelec, 2011/10/16}
\enonce{006711}{} 
Soient $f$ et $g$ deux fonctions holomorphes et ne s'annulant pas
dans un ouvert connexe $\Omega $ contenant le disque unité fermé.
On suppose que pour $\vert z\vert =1$, $\vert f(z)\vert=\vert g(z)\vert$
et que $f(0)$ et $g(0)$ appartiennent à ${\Rr}_+^*$. Montrer que $f=g$
sur $\Omega $.
\finenonce{006711}



\finexercice       


\exercice{6712, queffelec, 2011/10/16}
\enonce{006712}{} 
Soit $f$ une fonction holomorphe sur $\{z\in\C\vert\
\vert z\vert <R\}$. Si $f(z)=\sum_{n=0}^\infty a_nz^n$ pour $\vert
z\vert<R$, on pose pour $0\le r< R$ :
\begin{eqnarray*}\lefteqn{M(r,f)=\max_{\vert z\vert =r}\vert f(z)\vert
,\  M_1(r,f)=\sum_{n=0}^\infty \vert a_n\vert r^n,}\\ 
M_2(r,f)&=&\left({1\over 2\pi}\int_0^{2\pi}\vert f(re^{i\theta })\vert
^2d\theta \right)^{1/2}.\end{eqnarray*}

\begin{enumerate}  
\item 
\begin{enumerate}  
\item Montrer que pour $0\le r<R$,
$$M_2(r,f)=\left( \sum_{n=0}^\infty\vert
a_n\vert^2r^{2n}\right)^{1/2}.$$

 \item Déduire de (a) que $r\mapsto M_2(r,f)$ est une fonction continue
croissante.

  \item Déduire de (a) une autre démonstration des inégalités de Cauchy.
\end{enumerate}
\item
\begin{enumerate}  
\item Montrer que pour $0\le r<r\alpha <R$, on a
$$M(r,f)\le M_1(r,f)\le {\alpha \over \alpha -1}M(\alpha r,f)$$
(pour démontrer la seconde inégalité, on pourra utiliser les inégalités
de Cauchy).

\item Montrer que la fonction $r\mapsto M(r,f)$ est continue et
croissante, et même strictement croissante si $f$ n'est pas constante.
\end{enumerate}
\item On rappelle que si les deux séries $\sum_{n\ge 0}\vert
\alpha _n\vert^2$ et $\sum_{n\ge 0}\vert \beta
_n\vert^2$ convergent (où $\alpha _n\in\C$ et $\beta _n\in\C$), la série
$\sum_{n\ge 0}\alpha _n\overline{\beta _n}$ converge et on a
l'inégalité de Cauchy-Schwarz :
$$\left\vert \sum_{n\ge 0}\alpha _n\overline{\beta _n}\right\vert^2\le
\left(\sum_{n\ge 0}\vert \alpha _n\vert^2\right)
\left(\sum_{n\ge 0}\vert \beta _n\vert^2\right) .$$
Montrer que pour $\alpha \in \mathopen]0,1\mathclose[$, on a
$$\sqrt{1-\alpha ^2}M_1(r\alpha ,f)\le M_2(r,f)\le M(r,f).$$

\item En utilisant les questions 2. et 3., montrer que si $0\le
r<R$, $$\lim_{n\to \infty}M_1(r,f^n)^{1/n}=\lim_{n\to
\infty}M_2(r,f^n)^{1/n}= M(r,f)$$
(on pourra remarquer que $M(r,f^n)=M(r,f)^n$).
\end{enumerate}
\finenonce{006712}



\finexercice       


\exercice{6713, queffelec, 2011/10/16}
\enonce{006713}{} 
Soit $f$ une fonction holomorphe non identiquement nulle dans un
ouvert $\Omega $ connexe con\-te\-nant le disque fermé
$\overline{D(z_0,r)}$. Soit $\gamma $ le cercle $\{z\vert\ \vert
z-z_0\vert =r\}$ orienté positivement.

\begin{enumerate}  
\item Montrer que $f$  a un nombre fini
de zéros dans $D(z_0,r)$.

\item On suppose que $f$ n'a pas de zéros sur $\gamma ^*$. Calculer
$\displaystyle {1\over 2i\pi}\int_\gamma {f'(z)\over f(z)}z^pdz$ pour
$p=0,1,\dots$ (utiliser les zéros de $f$ dans $D(z_0,r)$). 

\item On suppose toujours que $f$ n'a pas de zéros sur $\gamma ^*$. On prend
$p=0$. Montrer que l'intégrale précédente est égale au nombre total de
zéros de $f$ dans $D(z_0,r)$. Montrer que ce nombre est aussi l'indice du
point $0$ par rapport à la courbe fermée $\Gamma =f(\gamma )$.

\item Soit $w_0=f(z_0)$ et $n$ la multiplicité du zéro $z_0$ pour la
fonction $f(z)-w_0$. Montrer que l'on peut choisir $\varepsilon >0$ tel
que $f'(z)$ ne s'annule pas pour $0<\vert z-z_0\vert <\varepsilon $, et
qu'il existe un $\delta >0$ tel que pour tout $a$ avec $\vert a-w_0\vert
<\delta $, l'équation $f(z)=a$ a exactement $n$ racines dans le disque
$\vert z-z_0\vert <\varepsilon $.

\item En déduire le principe du maximum : si $f$ est holomorphe et non
constante dans $\Omega $, alors $\vert f\vert$ n'a pas de maximum
dans $\Omega $.
\end{enumerate}
\finenonce{006713}



\finexercice       


\exercice{6714, queffelec, 2011/10/16}
\enonce{006714}{} 
\label{gijsexophi}
Soit $D$ le disque unité ouvert : 
$$D=\{ z\in \C\vert\ \vert z\vert
<1\}$$
 et $C$ le cercle unité : 
$$C=\{ z\in \C\vert\ \vert
z\vert =1\} .$$
Pour $a\in D$, on considère l'application
homographique
$\Phi _a$ :
$$\Phi _a(z)={z-a\over 1-\overline{a}z}.$$

\begin{enumerate}  
\item Montrer que $\Phi _a(C)\subset C$, puis que $\Phi_a$ est une bijection
de $D$ sur lui-même, de réciproque $\Phi _{-a}$.

\item Soit $f$ un biholomorphisme de $D$, c'est-à-dire une fonction
holomorphe de
$D$ sur lui-même, bijective, telle que $f^{-1}$ soit aussi holomorphe.
Soit $a=f^{-1}(0)$. En considérant $f\circ (\Phi _a)^{-1}$ et sa
réciproque, montrer qu'il existe $\varphi\in {\Rr}$ tel que $$\forall
z\in D,\ f(z)=e^{i\varphi}\Phi _a(z),$$ c'est-à-dire qu'à rotation près, les
seuls biholomorphismes du disque sont les $\Phi _a$.
\end{enumerate}
\finenonce{006714}



\finexercice       


\exercice{6715, queffelec, 2011/10/16}
\enonce{006715}{} 
Soit $D$ le disque unité ouvert et  $f$ une fonction holomorphe de
$D$ dans lui-même. Soit $\Phi _a$ la fonction définie dans l'exercice
précédent. Quelle est l'image de 0 par $h=\Phi _{f(a)}\circ f\circ (\Phi
_a)^{-1}$ ? En déduire que pour tout $z$ de $D$, $$\left\vert
{f(z)-f(a)\over 1-\overline{f(a)}f(z)}\right\vert \le
\left\vert {z-a\over 1-\overline{a}z}\right\vert$$
puis 
$$\vert f'(a)\vert\le {1-\vert f(a)\vert^2\over 1-\vert a\vert^2}$$
(lemme de Schwarz-Pick).
\finenonce{006715}



\finexercice       


\exercice{6716, queffelec, 2011/10/16}
\enonce{006716}{} 
Soit $D$ le disque unité ouvert et $f$ une fonction holomorphe de
$D$ dans $D$. On suppose que $f$ admet au moins deux points fixes,
c'est-à-dire qu'il existe $a$ et $b$ dans $D$ , $a\ne b $, tels que $f(a)=a$
et $f(b)=b$. Montrer que $f$ est l'identité de $D$. On pourra utiliser
l'application  $\Phi _a$ définie dans l'exercice \ref{gijsexophi} pour se
ramener au cas où l'un des points fixes est $0$.
\finenonce{006716}



\finexercice       


\exercice{6717, queffelec, 2011/10/16}
\enonce{006717}{} 
Soit $D$ le disque unité ouvert. On dira qu'une fonction $E$ est
unitaire si elle est holomorphe dans $D$, continue sur $\overline{D}$ et
si $\vert f(z)\vert =1$ si $\vert z\vert=1$.

\begin{enumerate}  
\item Montrer qu'une fonction unitaire dans $D$ n'a qu'un nombre fini de
zéros.

Montrer qu'une fonction unitaire sans zéro est une constante.

Montrer qu'une fonction unitaire ayant les points $a_1,a_2,\dots,a_n$
pour zéros (chacun étant compté avec son ordre de multiplicité) s'écrit
$$E(z)=c\prod_{j=1}^n{z-a_j\over 1-\overline{a_j}z}.$$

\item Soit $f$ holomorphe sur $D$ et non identiquement nulle et supposons
qu'il existe $M>0$ tel que $\vert f(z)\vert \le M$ sur $D$. Soit $E$ une
fonction unitaire dans $D$ et telle que $f(z)/E(z)$ soit holomorphe dans
$D$. Montrer que l'on a
$$\forall z\in D,\ \vert f(z)\vert \le M\vert E(z)\vert .$$
Soit $a_1,a_2,\dots, a_n,\dots $ la suite des zéros de $f$ dans $D$,
chacun étant compté avec son ordre de multiplicité. Montrer que
$$\forall n\ge 1,\ \vert f(0)\vert \le M\vert a_1\vert\vert
a_2\vert\cdots\vert a_n\vert .$$
En déduire que si $f(0)\ne 0$, la série $\sum_{n\ge 1}(1-\vert
a_n\vert)$ converge.
\end{enumerate}
\finenonce{006717}



\finexercice       


\exercice{6718, queffelec, 2011/10/16}
\enonce{006718}{} 
Soit $f$ une fonction entière telle que $\vert f(z)\vert =1$ si
$\vert z\vert =1$. Soit $D$ le disque unité fermé.

\begin{enumerate}  
\item Supposons que $f$ n'a pas de zéros dans $D$. Montrer qu'il existe
$k\in\C$ tel que $f(z)=k$ pour tout $z$ de $\C$.

\item Soit $a_1,\dots,a_n$ (pourquoi un nombre fini ?) les zéros de $f$
dans $D$ , chacun étant compté avec son ordre de multiplicité. En
étudiant la fonction
$$g(z)=f(z)\prod_{j=1}^n{1-\overline{a_j}z\over z-a_j}$$
montrer qu'il existe $k\in\C$, $\vert k\vert =1$, et $n\in {\Nn}$ tels
que $f(z)=kz^n$.
\end{enumerate}
\finenonce{006718}



\finexercice       


\exercice{6719, queffelec, 2011/10/16}
\enonce{006719}{Théorème des trois cercles d'Hadamard}

Soit $f$ une fonction holomorphe dans un domaine contenant la couronne
fermée constituée par les $z\in\C$ tels que $r_1\le\vert z\vert \le
r_2$ (où $0<r_1<r_2$). On pose $M(r)=\max_{\vert z\vert =r}\vert
f(z)\vert$ pour $r_1\le r\le r_2$.

\begin{enumerate}  
\item Montrer qu'il existe un nombre réel $\alpha $ tel que
$r_1^\alpha M(r_1)=r_2^\alpha M(r_2)$.

\item Montrer que
$$M(r)\le M(r_1)^{\ln{r_2}-\ln{r}\over \ln{r_2}-\ln{r_1}}
M(r_2)^{\ln{r}-\ln{r_1}\over \ln{r_2}-\ln{r_1}}$$
(on appliquera le principe du maximum à la fonction $z^pf(z)^q$ où
$p\in{\Zz}$ et $q\in {\Nn}^*$, puis on considèrera une suite
$(p_n,q_n)$, $p_n\in{\Zz}$ et $q_n\in {\Nn}^*$, telle que
$\lim_{n\to\infty}p_n/q_n=\alpha $).
\end{enumerate}
\finenonce{006719}



\finexercice


\exercice{7227, megy, 2021/03/06}

\enonce{007227}{}
Représenter graphiquement les chemins suivants de \(\C\) :
\begin{enumerate}
\item \(\gamma(t)=t+it\), \(t\in [0,1]\),
\item \(\gamma(t)=t^2-it\), \(t\in [0,1]\),
\item \(\gamma(t)=|t|+it\), \(t\in [-1,1]\),
\item \(\gamma=\gamma_1\vee\gamma_2\vee\gamma_3\) où : \(\gamma_1(t)= it^2\), \(t\in [0,1]\); \(\gamma_2(t)=t+(1-t)i\), \(t\in [0,1]\); \(\gamma_3(t)=e^{-it}\), \(t\in [0,\pi]\).
\end{enumerate}
\finenonce{007227}



\finexercice


\exercice{7228, megy, 2021/03/06}

\enonce{007228}{}
Calculer les intégrales suivantes :
\begin{enumerate}
\item $\int_\gamma z^2dz$ pour $\gamma(t)=1+t(1+i),$  $t\in [0,1].$
\item $\int_\gamma (z^3+1)dz$ pour $\gamma(t)=|t|-it$, $t\in [-1,1].$
\item $\int_{\gamma}\Im(z)dz$ pour $\gamma(t)=t+i\mathrm{inf}\{t,1\}$, $t\in [0,2].$
\item $\int_{\gamma}\Re(z^2)dz$ pour $\gamma(t)=1-it+t^2,$ $t\in [0,1].$
\end{enumerate}
\finenonce{007228}



\finexercice


\exercice{7229, megy, 2021/03/06}

\enonce{007229}{}
On considère la fonction $F$ définie sur le plan complexe  par 
\[F(z)=z^2e^{\cos\left(\frac{\pi}{4}(z-1)\right)}.\]
On considère de plus le chemin suivant \(\gamma(t)=\frac{1}{2}+\frac{1}{2}e^{-i\pi t}\) pour \(t\in [0,1]\).
\begin{enumerate}
\item Justifier que $F$ est holomorphe et calculer $F'.$
\item En déduire 
$$\int_\gamma z\left(2-z\frac{\pi}{4}\sin\left(\frac{\pi}{4}(z-1)\right)\right)e^{\cos\left(\frac{\pi}{4}(z-1)\right)}dz.$$
\end{enumerate}
\finenonce{007229}



\finexercice


\exercice{7230, megy, 2021/03/06}

\enonce{007230}{}
Soit \(\gamma\) un chemin \(\mathcal{C}^1\) par morceaux de \(\C\) allant de \(0\) à \(i\). Calculer \(\int_\gamma f(z)dz\) avec :
\begin{enumerate}
\item \(f(z)=z^2\sin z\)
\item \(f(z)=ze^{iz}\).
\end{enumerate}
\finenonce{007230}



\finexercice


\exercice{7231, megy, 2021/03/06}

\enonce{007231}{}
(Variations sur l'aire.)
\begin{enumerate}
\item Pour tout \(r>0\), montrer que 
\[\int_{\partial B(0,r)}\bar z dz=2i\mathrm{Aire}\big(B(0,r)\big).\]
\item Soient \(z_1,z_2,z_3\in \C\) des points non-alignés. On note \(\Delta\) le triangle de sommets \(z_1,z_2,z_3\). Montrer que 
\[\int_{\partial \Delta}\bar z dz=2i \mathrm{Aire}(\Delta).\]
\item De manière générale, soit \(K\) un compact à bord régulier. Montrer que 
\[\int_{\partial K}\bar z dz=2i\mathrm{Aire}(K).\]
\end{enumerate}
\finenonce{007231}



\finexercice


\exercice{7232, megy, 2021/03/06}

\enonce{007232}{}
Soit \(U\subset \C\) un ouvert contenant \(0\). Soit \(f:U\to \C\) une fonction de classe \(\mathcal{C}^1\).
\begin{enumerate}
\item Montrer que pour tout \(r>0\) tel que \(\overline{B}(0,r)\subset U\), on a
\[\int_{\partial B(0,r)}f(z)dz=2i\int_{B(0,r)} \frac{\partial f}{\partial \bar z}dxdy.\]
\item En déduire que 
\[\lim_{r\to 0^+}\frac{1}{2i\pi r^2}\int_{\partial B(0,r)}f(z)dz=\frac{\partial f}{\partial \bar z}(0).\]
\end{enumerate}
\finenonce{007232}



\finexercice


\exercice{7233, megy, 2021/03/06}

\enonce{007233}{}
Calculer les intégrales suivantes:
\begin{enumerate}
\item \(\displaystyle{I_1=\int_{C(0,2)}\frac{z^3-iz+1}{z-i}dz}\),
\item \(\displaystyle{ I_2=\int_{C(2,2)}\frac{z^3-iz+1}{z-i}dz}\),
\item \(\displaystyle{I_3=\int_{C(0,2)}\frac{\sin(z)\cos(z)}{3z-\pi}dz}\),
\item \(\displaystyle{ I_4=\int_{C(2i,1)}\frac{e^{z^2}}{z^3(z-2i)}dz}\).
\end{enumerate}
\finenonce{007233}



\finexercice


\exercice{7234, megy, 2021/03/06}

\enonce{007234}{}
Calculer 
\(\displaystyle{\int_{C(1,\frac{1}{2})}\frac{e^z}{(z-1)(z+1)}dz}\)
et 
 \(\displaystyle{ \int_{C(-1,\frac{1}{2})}\frac{e^z}{(z-1)(z+1)}dz}\).
 En déduire 
\(\displaystyle{\int_{C(0,2)}\frac{e^z}{z^2-1}dz}\).
\finenonce{007234}



\finexercice


\exercice{7235, megy, 2021/03/06}

\enonce{007235}{}
Déterminer les racines du polynôme $P(z)=z^2+(1-i)z-i$.
En déduire l'intégrale \(\int_{C(0,2)}\frac{z-1}{z^2+(1-i)z-i}dz\).
\finenonce{007235}



\finexercice


\exercice{7236, megy, 2021/03/06}

\enonce{007236}{}
Pour $r>0$, calculer l'intégrale suivante :
$$\int_{C(0,r)}(|z|-e^{\cos(z)}\sin(z)+\bar{z})dz.$$

\finenonce{007236}



\finexercice


\exercice{7237, megy, 2021/03/06}

\enonce{007237}{}
Calculer les intégrales suivantes :
\begin{enumerate}
\item \(\displaystyle{I_1=\int_{C(0,2)}\frac{z^3-i}{(z-1)^2}dz}\),
\item \(\displaystyle{ I_2=\int_{C(0,1)}\frac{e^z}{z^3}dz}\),
\item \(\displaystyle{ I_3=\int_{C(i,5)}\frac{ze^{iz}}{(1+z)^3}dz}\),
\item \(\displaystyle{ I_4=\int_{C(0,1)}\frac{\cos(z)}{z^3(z-2)}dz}\).
\end{enumerate}
\finenonce{007237}



\finexercice


\exercice{7238, megy, 2021/03/06}

\enonce{007238}{}
Calculer les intégrales 
\(\displaystyle{\int_{C(i,\frac{1}{2})}\frac{iz^3-3}{(z-i)^2(z+i)^2}dz}\)
 et 
 \(\displaystyle{ \int_{C(-i,\frac{1}{2})}\frac{iz^3-3}{(z-i)^2(z+i)^2}dz}\).
En déduire 
\(\displaystyle{\int_{C(0,2)}\frac{iz^3-3}{(z-i)^2(z+i)^2}dz}\).
\finenonce{007238}



\finexercice


\exercice{7239, megy, 2021/03/06}

\enonce{007239}{}
Soit \(U\subset \C\) un ouvert contenant \(B(0,1)\). Soit \(f:U\to \C\) une fonction holomorphe. 
\begin{enumerate}
\item Calculer 
\[\int_{C(0,1)}\left(2+z+\frac{1}{z}\right)\frac{f(z)}{z}dz.\]
\item En déduire que
\[\frac{2}{\pi}\int_0^{2\pi}f(e^{i\theta})\cos^2\big(\theta/2\big)d\theta=2f(0)+f'(0).\]
\item Utiliser une stratégie similaire pour calculer 
\(\displaystyle{\frac{2}{\pi}\int_0^{2\pi}f(e^{i\theta})\sin^2\big(\theta/2\big)d\theta}.\)
\end{enumerate}
\finenonce{007239}



\finexercice


\exercice{7240, megy, 2021/03/06}

\enonce{007240}{}
(Généralisation de Liouville.) 
Soit \(f:\C\to \C\) une fonction entière non-constante. Montrer que l'image  \(f(\C)\) de \(\C\) par \(f\) est dense dans \(\C\). 
\finenonce{007240}



\finexercice


\exercice{7241, megy, 2021/03/06}

\enonce{007241}{}
Soit \(f:\C\to \C\) une fonction holomorphe telle que \(\Re(f)\) est bornée.
Montrer que \(f\) est constante.
\finenonce{007241}



\finexercice


\exercice{7242, megy, 2021/03/06}

\enonce{007242}{}
Soit $f:\C\to\C$ une fonction holomorphe.
Montrer que si elle vérifie 
\[ \forall z\in \C, \left( |f(z)|\leq |z| \text{ et } |f(z)|\leq |z|^2 \right),\]
alors elle est nulle.
\finenonce{007242}



\finexercice


\exercice{7243, megy, 2021/03/06}

\enonce{007243}{}
(Raffinement de Goursat.)
L'objectif de cet exercice est de démontrer la version suivante du lemme de Goursat :\\
\textbf{Lemme. (Raffinement de Goursat)}\\
Soit \(U\) un ouvert. Soit  \(\Delta\) un triangle dans \(U\). Soit \(z_0\in \Delta\). Soit \(f:U\to \C\) une fonction continue sur \(U\) holomorphe sur \(U\setminus\{z_0\}\). Alors,
\[\int_{\partial \Delta}f(z)dz=0.\]

\begin{enumerate}
\item On traite d'abord le cas où \(z_0\) est un sommet du triangle. On dénote les deux autres sommets par \(z_1\) et \(z_2\).
% AJOUTER FIGURE ?
\begin{enumerate}
\item Montrer que pour tout \(z_1'\) sur le segment \([z_0,z_1]\) et pour tout \(z_2'\) sur le segment \([z_0,z_2]\), on a 
\[\int_{\partial \Delta_{z_0z_1z_2}}f(z)dz=\int_{\partial \Delta_{z_0z'_1z'_2}}f(z)dz.\]
\item Montrer que quand \(z_1'\to z_0\) et \(z_2'\to z_0\), alors \(\displaystyle{\int_{\partial \Delta_{z_0z'_1z'_2}}f(z)dz\to 0}\).
\item Démontrer le lemme dans ce cas.
\end{enumerate}
\item On suppose dans cette question que \(z_0\in \partial \Delta\). À l'aide de la question précédente, démontrer le lemme  dans ce cas.
\item On suppose maintenant que \(z_0\in \mathring{\Delta}\). Démontrer le lemme  dans ce cas. Conclure.
\item Déduire de ce résultat que si \(f:U\to \C\) est une fonction continue, et holomorphe sauf peut-être en un nombre fini de points, alors \(f\) est holomorphe sur \(U\)
\end{enumerate}
\finenonce{007243}



\finexercice


\exercice{7244, megy, 2021/03/06}

\enonce{007244}{}
(Théorème d'extension de Riemann.) 
Le but de cet exercice est de démontrer le théorème d'extension de Riemann à l'aide du lemme de Goursat raffiné. \\
\textbf{Théorème. (d'extension de Riemann)}\\
Soit \(U\subset \C\) un ouvert. Soit \(z_0\in U\) soit \(f:U\setminus\{z_0\}\to \C\) une fonction holomorphe bornée dans un voisinage de \(z_0\). Alors \(f\) s'étend en une fonction holomorphe sur \(U\).

%Ici, «\(f\) bornée dans un voisinage de \(z_0\)» signifie qu'il existe \(\epsilon>0\) tel que \(B(z_0,\epsilon)\subset U\) et tel que \(f\) est bornée sur \(B(z_0,\epsilon)\setminus \{z_0\}\). L'expression «\(f\) s'étend en une fonction holomorphe sur \(U\)» signifie qu'il existe une fonction holomorphe \(\hat{f}:U\to \C\) telle que \(\hat{f}|_{U\setminus \{z_0\}}=f\).\\

Pour démontrer ce théorème, on considère la fonction \(F:U\to \C\) définie par 
\[F(z)=\left\{\begin{array}{ccc}0&\text{si}& z=z_0\\ (z-z_0)f(z)&\text{si}&z\neq z_0.\end{array}\right.\]

\begin{enumerate}
\item À l'aide du lemme de Goursat raffiné, démontrer que \(F\) est holomorphe sur \(U\).
\item À l'aide de la \(\C\)-dérivabilité de \(F\) en \(z_0\), montrer que \(f\) s'étend en une fonction continue sur \(U\).
\item Conclure.
\end{enumerate}
\finenonce{007244}



\finexercice


\exercice{7245, megy, 2021/03/06}

\enonce{007245}{}
Soit \(U\subset \C\) un ouvert. Soit \(L\subset \C\) une droite. Soit \(f:U\to \C\) une fonction continue telle que \(f|_{U\setminus L}\) est holomorphe. Montrer que \(f\) est holomorphe. 
\finenonce{007245}



\finexercice


\exercice{7246, megy, 2021/03/06}

\enonce{007246}{}
[Principe de réflexion de Schwarz]
Soit \(U\subset \C\) un ouvert  symétrique par rapport à l'axe réel.
Notons \(U_+:=U\cap\{\Im(z)>0\}\) et \(U_-:=U\cap \{\Im(z)<0\}\).
Soit \(f:\overline{U_+}\to \C\) une fonction continue telle que \(f|_{U_+}\) est holomorphe et telle que \(f(x)\in\R\) pour tout \(x\in \overline{U_+}\cap \R\).
\begin{enumerate}
\item Montrer que la fonction \(g:\overline{U_-}\to \C\) définie par \(g(z):=\overline{f(\bar z)}\) est continue sur \(\overline{U_-}\) et holomorphe sur \(U_-\).
\item Montrer que l'on peut étendre la fonction \(f\) en une fonction continue \(h:\overline{U}\to \C\) en posant \(h(z)=f(z)\) si \(z\in \overline{U_+}\) et \(h(z)=g(z)\) si \(z\in \overline{U_-}\).
\item Montrer que \(h|_{U}\) est holomorphe.
\end{enumerate}
\finenonce{007246}



\finexercice


\exercice{7247, megy, 2021/03/06}

\enonce{007247}{}
[Extension de \(\Gamma\)]
Notons \(U_{\Re>0}:=\left\{z\in \C\ ; \ \Re(z)>0\right\}\).
\begin{enumerate}
\item Montrer que pour tout \(z\in U_{\Re>0}\), on a 
\[\Gamma(z+1)=z\Gamma(z).\]
\item En déduire que \(\Gamma(n)=(n-1)!\) pour tout \(n\in \N^*\).
\item Utiliser la première question pour montrer que la fonction \(\Gamma\) s'étend en une fonction holomorphe sur \(\C\setminus \mathbb{Z}_-\).
\end{enumerate}
\finenonce{007247}



\finexercice


\exercice{7248, megy, 2021/03/06}

\enonce{007248}{}
Cet exercice a pour but de donner une introduction à la géométrie du disque. Cela donnera aussi une interprétation géométrique du lemme de Schwarz. On note \(\mathbb{D}\) le disque unité. Étant donné un vecteur tangent \(\xi\) en un point \(z\in \mathbb{D}\), on note sa norme euclidienne par \(\|\xi\|_{2}\) et on défini sa norme \emph{hyperbolique} ou sa norme de \emph{Poincaré} 
par \[\|\xi\|_{\rm hyp}:=\frac{\|\xi\|_2}{1-|z|^2}.\]
La \emph{longueur hyperbolique} d'un chemin \(\gamma:[a,b]\to\mathbb D\) est définie par
\[\ell_{\rm hyp}(\gamma)=\int_{a}^b\|\gamma'(t)\|_{\rm hyp}dt=\int_a^b\frac{\|\gamma'(t)\|_2}{1-|\gamma(t)|^2}dt.\]
La \emph{distance hyperbolique} ou \emph{distance de Poincaré} entre deux points \(z_0,z_1\in \mathbb D\) est définie par
\[d_{\rm hyp}(z_0,z_1):=\inf_{\gamma}\ell_{\rm hyp}(\gamma)\]
où le \(\inf\) est pris sur tous les chemins \(\mathcal{C}^1\) par morceaux de \(\mathbb D\) allant de \(z_0\) à \(z_1\). L'espace métrique \((\mathbb D,d_{\rm hyp})\) est appelé \emph{disque de Poincaré}. Les courbes de longueur minimale sont appelées \emph{géodésiques} ou \emph{droites hyperboliques}.
\begin{enumerate}
\item \begin{enumerate}
\item
À l'aide du lemme de Schwarz-Pick, montrer que toute application holomorphe \(f:\mathbb D\to \mathbb D\) est décroissante par rapport à la distance de Poincaré, c'est à dire que pour tout \(z_0,z_1\in \mathbb D\),
\[d_{\rm hyp}(f(z_1),f(z_1))\leqslant d_{\rm hyp}(z_0,z_1).\]
\item En déduire que pour tout automorphisme du disque \(f\in \rm{Aut}(\mathbb D)\) est une isométrie de \((\mathbb D,d_{\rm hyp})\).
\end{enumerate}
\item On veut maintenant obtenir une expression explicite pour \(d_{\rm hyp}\).
\begin{enumerate}
\item Soit \(w\in ]0,1[\). Montrer que la géodésique allant de \(0\) à \(w\) est le segment \([0,w]\) et montrer  \[d_{\rm hyp}(0,w)=\frac{1}{2}\ln\left(\frac{1+w}{1-w}\right).\]
\item Soit \(w\in\mathbb D^*\). À l'aide d'un automorphisme judicieusement choisi, montrer que la géodésique de allant de \(0\) à \(w\) est le segment \([0,w]\) et montrer que 
\[d_{\rm hyp}(0,w)=\frac{1}{2}\ln\left(\frac{1+|w|}{1-|w|}\right).\]
\item Soit \(z_1,z_2\in \mathbb D^*\). À l'aide d'un automorphisme judicieusement choisi, montrer que la géodésique de allant de \(z_1\) à \(z_2\) est une partie du cercle orthogonal à \(\partial \mathbb D\) passant par \(z_1\) et \(z_2\) et que 
\[d_{\rm hyp}(z_1,z_2)=\frac{1}{2}\ln\left(\frac{|1-z_1\bar z_2|+|z_1-z_2|}{|1-z_1\bar z_2|-|z_1-z_2|}\right)=2{\rm argtanh}\left|\frac{z_1-z_2}{1-z_1\bar z_2}\right|.\]
\item Montrer que \(d_{\rm {\rm hyp}}\) est bien une distance sur \(\mathbb D\). Montrer aussi que les \emph{boules hyperboliques}  sont des boules euclidiennes (pas nécessairement avec le même centre ou le même rayon).
\end{enumerate}
\item La géométrie du disque hyperbolique définie ci-dessus est appelée \emph{géométrie hyperbolique}. Vérifier qu'en géométrie hyperbolique, les axiomes d'Euclide sont vérifiés, à l'exception du cinquième postulat.
\item Rappelons que \(\mathbb D\) est biholomorphe au demi-plan \(\mathbb{H}:=\left\{z\in \C;\ \Im(z)>0\right\}\)  (dit \emph{demi-plan de Poincaré}), via l'application \(\varphi:z\mapsto i\frac{1+z}{1-z}\). Décrire les géodésiques de \((\mathbb{H},d_{\mathbb{H}})\) où \(d_{\mathbb{H}}\) est la métrique induite par \(d_{\rm hyp}\) via \(\varphi\).
\item Nous concluons maintenant par une preuve plus <<géométrique>> du théorème de Liouville. Soit \(f:\C\to \C\) une application holomorphe bornée.
\begin{enumerate}
\item Montrer que l'on peut supposer que \(f(z)\in \mathbb{D}\) pour tout \(z\in \C\) (ce que l'on supposera dans la suite).
\item Pour tout \(R\in \R_+^*\) on considère l'application \(f_R:\mathbb D\to \mathbb D\) définie par \(f_R(z)=f(Rz)\) \(\forall z\in \mathbb{D}\). À l'aide la propriété de décroissance de la distance de Poincaré, montrer que pour tout  \(R\in \mathbb{R}_+^*\) et pour tout \(z_1,z_2\in B(0,R)\), 
\[d_{\rm hyp}(f(z_1),f(z_2))=d_{\rm hyp}\big(f_R\big(\frac{z_1}{R}\big),\big(\frac{z_2}{R}\big)\big)\leqslant d_{\rm hyp}\big(\frac{z_1}{R},\frac{z_2}{R}\big)\]
\item En déduire  que \(f\) est constante.
\end{enumerate}
\end{enumerate}
\finenonce{007248}



\finexercice

\section{443.00 Singularité}


\exercice{2680, matexo1, 2002/02/01}

\enonce{002680}{}
{\`A} l'aide de la formule
$$ f^{(n)}(a) = {n!\over 2i\pi } \int_\Gamma {f(z)\over (z-a)^{n+1}} \,dz$$
o{\`u} $f$ est m{\'e}romorphe dans un domaine contenant le contour simple 
$\Gamma$, et $a$ est
un point int{\'e}rieur {\`a} $\Gamma$, montrer que l'on a
$$ \left(x^n\over n!\right)^2 = {1\over 2i\pi } 
\int_C {x^n e^{xz}\over n! z^{n+1}} \,dz$$
o{\`u} $C$ est le cercle unit{\'e} de $\C$. En d{\'e}duire que l'on a
$$ \sum_{n=0}^{+\infty } {x^{2n}\over \left(n!\right)^2} = {1\over 2\pi }
 \int_0^{2\pi } e^{2x\cos \theta}\,d\theta. $$
\finenonce{002680}
\finexercice


\exercice{6720, queffelec, 2011/10/16}
\enonce{006720}{}
Montrer que la fonction $f$ définie sur $\Cc^*$ par $f(z)=\displaystyle{z\over
\sin z+i\sinh z}$ se prolonge en une fonction holomorphe en $0$; quel est le
rayon de son développement en $0$ ?
\finenonce{006720}



\finexercice


\exercice{6721, queffelec, 2011/10/16}
\enonce{006721}{}
\begin{enumerate}
\item Déterminer les développements en série de Laurent de $f(z)={1\over
1-z^2}+{1\over 3-z}$ dans les domaines $D=D(0,1)$, $C_1=\{1<|z|<3\}$ puis $C_2=
\{|z|>3\}$.

\item  Déterminer les développements en série de Laurent de $f(z)={z\over
z-1}e^z $ dans les domaines  $C_1=\{|z|<1\}$ puis $C_2=
\{|z|>1\}$.
\end{enumerate}
\finenonce{006721}



\finexercice


\exercice{6722, queffelec, 2011/10/16}
\enonce{006722}{}
Soit $\alpha\in\Rr$. Montrer que le développement en série de Laurent en $0$ de
la fonction
$f(z)=\hbox{exp}{\alpha\over2}(z+{1\over z})$ est de la forme
$a_0+\sum_{n\geq1}a_n(z^n+{1\over z^n})$, où 
$$a_0={1\over 2\pi}\int_0^{2\pi}\hbox{exp}(\alpha\cos t)\ dt,\quad a_n={1\over
2\pi}\int_0^{2\pi}\hbox{exp}(\alpha\cos t)\cos(nt)\ dt \ \hbox{si}\ n\geq1.$$

(Calculer $f$ pour $|z|=1$ et conclure avec le prolongement analytique.)
\finenonce{006722}



\finexercice


\exercice{6723, queffelec, 2011/10/16}
\enonce{006723}{} 
Développer les fonctions suivantes en série de Laurent dans
chacun des ouverts donnés

\begin{enumerate}  
\item $\displaystyle f(z)={1\over (z-1)(z-2)}$ dans $\vert z\vert <1$ ; 
$1<\vert z\vert <2$ ; $2<\vert z\vert $ ;

\item $\displaystyle f(z)={1\over (z-a)^k}$ ($k\in {\Nn}^*$) dans $\vert
z\vert <\vert a\vert$ et dans $\vert z\vert >\vert a\vert$ ;

\item $\displaystyle f(z)={1\over z(z-a)}$ dans $0<\vert
z\vert <\vert a\vert$ et dans $\vert a\vert <\vert z\vert$ ;

\item $\displaystyle f(z)={1\over (z-a)(z-b)}$ ($0<\vert a\vert <\vert
b\vert$) dans $0<\vert z\vert <\vert a\vert$ ; $\vert a\vert<\vert z\vert
<\vert b\vert$ ; $\vert b\vert <\vert z\vert$ ;

\item une détermination holomorphe $f$ de $[(z-a)(z-b)]^{1\over 2}$
($0<\vert a\vert =\vert b\vert$) dans $0<\vert z\vert <\vert a\vert$ ;
$\vert b\vert <\vert z\vert$ ;

\item $f(z)=z^2\exp{(z^{-1})}$ dans $0<\vert z\vert$.

\item $f(z)=\exp{(z+z^{-1})}$ dans $0<\vert z\vert$.

\item $f(z)=\sin z\cdot\sin{\left(z^{-1}\right)}$ dans $0<\vert z\vert$.

\item $f(z)=\mathrm{cotan} z$ dans $k\pi<\vert z\vert<(k+1)\pi$ ($k\in {\Nn}$) on
pourra exprimer le résultat en fonction des nombres $B_n$ de Bernoulli,
définis par : $${z\over \exp{(z)}-1}=\sum_{n\ge 0}B_n{z^n\over n!}$$
($B_0=1$,
$B_1=-1/2$, et $B_{2n+1}=0$ pour $n\ge 1$).
\end{enumerate}
\finenonce{006723}



\finexercice       


\exercice{6724, queffelec, 2011/10/16}
\enonce{006724}{} 
Déterminer la couronne de convergence des séries de
Laurent $\sum_{n\in{\Zz}}a^{\vert n\vert}z^n$,
$\sum_{n\in{\Zz}}\vert n\vert !z^n$, $\sum_{n\in{\Zz}}R(n)z^n$ ($R$
fonction rationnelle sans pôles dans ${\Zz}$).
\finenonce{006724}



\finexercice       


\exercice{6725, queffelec, 2011/10/16}
\enonce{006725}{} 
Soit un ouvert $U$ de $\C$ et $f$ une fonction définie sur $U$
admettant en tout point $a$ de $U$ un développement de Laurent
$$f(z)=\sum_{n\ge r_a}c_n(z-a)^n$$
($r_a\in{\Zz}$) convergeant dans un disque pointé $0<\vert z-a\vert
<R_a$. Montrer que pour tout compact $K\subset U$, il existe une
fonction rationnelle $g_K$ nulle à l'infini et telle que la fonction
$f-g_K$ soit holomorphe sur un voisinage de $K$ ($f$ est dite
méromorphe sur $U$).
\finenonce{006725}



\finexercice       


\exercice{6726, queffelec, 2011/10/16}
\enonce{006726}{} 
Déterminer les points singuliers des fonctions suivantes, puis
donner la nature de ces points singuliers (singularité effaçable, pôle
d'ordre $n$, singularité essentielle isolée, accumulation de points
singuliers).
\begin{enumerate} 
\item $$ z\mapsto {1\over z(z^2+4)^2} $$
\item $$ z\mapsto {1\over \exp{(z)}-1}-{1\over z} $$
\item $$ z\mapsto \sin{1\over 1-z}$$
\item $$ z\mapsto \exp{z\over 1-z}$$
\item $$ z\mapsto \mathrm{cotan} z-{1\over z}$$
\item $$ z\mapsto \mathrm{cotan}{1\over z}$$
\item $$ z\mapsto {1\over \sin z-\sin a}$$
\item $$ z\mapsto \sin{\left({1\over \sin{1\over z}}\right) }$$
\item $$ z\mapsto \exp{\left(\tan{1\over z}\right)} $$
\end{enumerate} 
\finenonce{006726}



\finexercice       


\exercice{6727, queffelec, 2011/10/16}
\enonce{006727}{} 
Exhiber des fonctions n'ayant dans le plan complexe que les
sin\-gu\-la\-rités suivantes:

\begin{enumerate}  
\item un pôle triple en 0, un pôle simple en 1 et un point singulier
essentiel en $i$ et $-i$.

\item un point singulier essentiel en tout entier.
\end{enumerate} 
\finenonce{006727}



\finexercice       


\exercice{6728, queffelec, 2011/10/16}
\enonce{006728}{} 
Déterminer les singularités isolées $a$ des fonctions $f$ suivantes
 et calculer $\mathrm{Res} (f,a)$.

$$\begin{array}{lll}
1. \displaystyle z\mapsto {1\over z^3-z^5}\ &
2.\ \displaystyle z\mapsto {z^2\over (z^2+1)^2}\ &
3.\ \displaystyle z\mapsto \exp{(z+z^{-1})}\\
4.\ \displaystyle z\mapsto {\sin{(2z)}\over (z+1)^3}\ &
5.\ \displaystyle z\mapsto \cos{\left( {z^2+4z-1\over z-3}\right) }\ &
6.\ \displaystyle z\mapsto z^n\sin{({1\over z})}\\
\end{array}$$
\finenonce{006728}



\finexercice       


\exercice{6729, queffelec, 2011/10/16}
\enonce{006729}{} 
Soit $U$ un ouvert simplement connexe de $\C$, $f$ une fonction
holomorphe sur $U\setminus S$ où $S$ est une partie fermée discrète de
$U$. Montrer que $f$ a une primitive sur $U\setminus S$ si et seulement
si pour tout point $s$ de $S$, le résidu de $f$ au point $s$ est nul.
\finenonce{006729}



\finexercice       


\exercice{6730, queffelec, 2011/10/16}
\enonce{006730}{} 
Calculer les intégrales suivantes, où les chemins fermés simples
$\gamma $ sont parcourus dans le sens direct.

\begin{enumerate}  
\item $\displaystyle \int_\gamma {z\over (z-1)(z-2)^2}dz$ où $\gamma $
est le cercle $\vert z-2\vert ={1\over 2}$ ;

\item $\displaystyle \int_\gamma {\exp z\over z^2(z-9)^2}dz$ où $\gamma
$ est le cercle $\vert z\vert =1$ ;

\item $\displaystyle \int_\gamma \exp{\left( {1\over z}\right) }\ dz$ où
$\gamma $ est le cercle $\vert z\vert =1$ ;

\item $\displaystyle \int_\gamma \sin^2{\left( {1\over z}\right) }\ dz$ où
$\gamma $ est le cercle $\vert z\vert =r$ ;

\item $\displaystyle \int_\gamma (z^2+z+1)^{-1/2}\ dz$ où $\gamma $
est le cercle $\vert z\vert =r\ne 1$.
\end{enumerate}
\finenonce{006730}



\finexercice 


\exercice{7573, mourougane, 2021/08/10}
\enonce{007573}{Exemples de singularités isolées}
Décrire le type singularité (apparente, pôle ou essentielle) des applications suivantes 
 $$f(z)=\frac{1}{\sin (z)^2}\ \ , \ \ g(z)=\frac{z}{e^z-1}\ \ \ \ \quad \text{ et } \quad h(z)=e^{1/z}.$$
\finenonce{007573}
\finexercice


\exercice{7574, mourougane, 2021/08/10}
\enonce{007574}{Théorème d'identité}
Si $f$ est une fonction méromorphe sur un ouvert $D$ de $\Cc$, on notera $P(f)$ le lieu de ses pôles.
\begin{enumerate}
 \item Soit $D$ un ouvert de $\Cc$ et $f,g$ deux applications méromorphes sur $D$.
Montrer qu'on a équivalence entre
\begin{enumerate}
 \item $f=g$ 
\item $\{z\in D-P(f)-P(g), f(z)=g(z)\}$ a un point d'accumulation dans $D$.
\end{enumerate}
\item Peut-on construire une application méromorphe sur $\Cc^\star$ non nulle
mais nulle en tous les réels de la forme $\frac{1}{n}$ avec $n\in\Nn$ ?
\end{enumerate}
\finenonce{007574}
\finexercice


\exercice{7575, mourougane, 2021/08/10}
\enonce{007575}{Fonctions rationnelles}
Si $f$ est une fonction méromorphe sur un ouvert $D$ de $\Cc$, on notera $P(f)$ le lieu de ses pôles.
Soit $f$ une application méromorphe sur $\Cc$ et $n$ un entier naturel et $r$ un nombre réel strictement positif tels que 
$$\forall z\in\Cc-P(f)-\Delta_r, |f(z)|\leq |z|^n.$$
Montrer que $f$ est une application rationnelle.
\finenonce{007575}
\finexercice

\section{444.00 Théorème des résidus}


\exercice{2672, matexo1, 2002/02/01}

\enonce{002672}{}
Calculer par la
m{\'e}thode des r{\'e}sidus
$$ I = \int_0^\pi  {a\,d\varphi\over a^2 + \sin^2\varphi} \qquad (a>0)$$
\finenonce{002672}
\finexercice


\exercice{2673, matexo1, 2002/02/01}

\enonce{002673}{}
Calculer les int{\'e}grales
$$ I = \int_{-\infty }^{+\infty } {dx\over x^4+x^2+1}
\qquad J = \int_{-\infty }^{+\infty } {\sin^2 x\over x^4+x^2+1}\,dx.$$
\finenonce{002673}
\finexercice


\exercice{2674, matexo1, 2002/02/01}

\enonce{002674}{}
Calculer par la m{\'e}thode des r{\'e}sidus l'int{\'e}grale de Wallis
$$ W_n = \int_0^{\pi \over 2} \cos^{2n} \theta\,d\theta. $$
\finenonce{002674}
\finexercice


\exercice{2675, matexo1, 2002/02/01}

\enonce{002675}{}
D{\'e}velopper la fonction
$$f(x) = {1\over 2-\cos x} $$ 
en s{\'e}rie de Fourier, en calculant les coefficients par
la m{\'e}thode des r{\'e}sidus.
\finenonce{002675}
\finexercice


\exercice{2676, matexo1, 2002/02/01}

\enonce{002676}{}
R{\'e}soudre l'{\'e}quation $\cos z = a$, o{\`u} $a$ est un r{\'e}el $>1$. Donner le sinus des
solutions. En d{\'e}duire la valeur de
$$ I_a = \int_0^{+\infty } {dx\over (1+x^2) (a-\cos x)}.$$
\finenonce{002676}
\finexercice


\exercice{2677, matexo1, 2002/02/01}

\enonce{002677}{}
Soit $a\in[0,1[$ un r{\'e}el. En int{\'e}grant $e^{az}/\cosh z$ le long du rectangle de
sommets
$-R$, $+R$, $R+i\pi $, $-R+i\pi $, montrer que l'on a
$$ I= \int_{-\infty }^{+\infty } {e^{ax}\over \cosh x} \,dx = {\pi \over \cos(\pi a/2)}.$$
\finenonce{002677}
\finexercice


\exercice{2678, matexo1, 2002/02/01}

\enonce{002678}{}
En int{\'e}grant $e^{2iaz -z^2}$ le long du rectangle de sommets $0$, $R$, $R+ia$, $ia$, et
en faisant tendre $R$ vers $+\infty $, montrer que l'on a
$$ \int_0^{+\infty } e^{-x^2}\cos 2ax \,dx = {\sqrt \pi \over 2} e^{-a^2}.$$
(On admettra la formule $\int_0^{+\infty }e^{-x^2}\,dx = \sqrt \pi /2$.)
\finenonce{002678}
\finexercice


\exercice{2679, matexo1, 2002/02/01}

\enonce{002679}{}
Soit $R$ une fraction rationnelle, ou plus g{\'e}n{\'e}ralement une fonction
 m{\'e}romorphe sur $\C$, sans p{\^o}le r{\'e}el.
\begin{itemize}
\item On souhaite calculer
$$ I = \int_0^{+\infty } x R(x^4)\,dx $$
Montrer que cela peut se faire par la m{\'e}thode des r{\'e}sidus, en int{\'e}grant sur un
contour form{\'e} du bord du quart de cercle  $\{ 0< \arg z<{\pi \over 2}, |z|<a\}$.

\item Application: calculer
$$  \int_0^{+\infty } {x\over 1+x^8}\,dx .$$

\item Plus g{\'e}n{\'e}ralement, montrer que si $n$ et $p$ sont des entiers, et $p\geq 3$, on
peut calculer
$$ I(n,p) = \int_0^{+\infty } x^n R(x^p)\, dx $$
sous une condition sur $n$, $p$ que l'on pr{\'e}cisera.

\item Application: calculer
$$ I_{p} = \int_0^{+\infty } {x\over 1+x^{2p}} \,dx.$$
\end{itemize}
\finenonce{002679}
\finexercice


\exercice{2820, burnol, 2009/12/15}
\enonce{002820}{} 
\label{ex:burnol1.1}
Soit $\Omega =
\Cc\setminus\{]-\infty,0]\}$. Déterminer en tout
$z_0\in\Omega$ la série de Taylor de la fonction holomorphe
$z\mapsto\mathrm{Log}\, z$ ainsi que son rayon de convergence. Soit
$z_0$ avec $\Re(z_0)<0$. Soit $R_0$ le rayon de
convergence pour $z_0$ et soit $f(z)$ la somme de la série
dans $D(z_0,R_0)$. A-t-on $f(z) = \mathrm{Log}\, z$ dans $D(z_0,R_0)$?
\finenonce{002820}
\finexercice


\exercice{2821, burnol, 2009/12/15}
\enonce{002821}{} 
On considère la fonction analytique $f(z) =
\frac1{\sin(z)}$ sur l'ouvert $U$ complémentaire de
$\pi\Zz$. Vérifier que la fonction $\sin(z)$ ne s'annule
jamais sur $U$. Déterminer  en tout $z_0\in U$ donné le
rayon de convergence du développement en série de Taylor de
$f$. 
\emph{Remarque :} il est déconseillé de chercher à résoudre ce
problème en déterminant explicitement les coefficients des
séries de  Taylor.
\finenonce{002821}
\finexercice


\exercice{2822, burnol, 2009/12/15}
\enonce{002822}{}
  Soient $f$ et $g$ deux fonctions entières avec $\forall z\
f(z)g(z) = 0$. Montrer que l'une des deux est identiquement
nulle.
\finenonce{002822}
\finexercice


\exercice{2823, burnol, 2009/12/15}
\enonce{002823}{} 
\label{ex:burnol1.4}
Soit $f$ une fonction holomorphe sur un ouvert
\textbf{convexe} $U$. Soit $z_1\in U$, on suppose que le
rayon de convergence de la série de Taylor de $f$ en $z_1$
est $R_1$. De même, en $z_2\in U$, on suppose que le rayon de
convergence de la série de Taylor de $f$ est $R_2$. Soit
$g_1$ sur le disque ouvert $D(z_1,R_1)$ la somme de la série
de Taylor de $f$ en $z_1$ et de même $g_2$ sur
$D(z_2,R_2)$. Soit  $V = D(z_1,R_1)\cap D(z_2,R_2)$. Montrer
que si $V$ est non vide alors $g_1 = g_2$ sur $V$. On
commencera par montrer que $V\cap U$ est non vide
aussi. 
\emph{Attention}: en général, sans hypothèse
spéciale comme la convexité de $U$ cela est complètement
faux; donner un exemple, avec $U$ connexe, mais
pas convexe, tel que  $g_1\neq g_2$ sur $V$ (et on peut même
faire avec $V\cap U\neq\emptyset$). Il suffira d'utiliser
l'exercice \ref{ex:burnol1.1}.
\finenonce{002823}
\finexercice


\exercice{2824, burnol, 2009/12/15}
\enonce{002824}{}
\begin{enumerate}
  \item 
Soit $\Omega$ l'ouvert habituel sur lequel est défini $\mathrm{Log}\,
z$. Justifier pour tout $z\in\Omega$
\[ \mathrm{Log}(z) = \int_0^1 \frac{z-1}{1+t(z-1)}dt\;,\]
et donner une formule intégrale explicite pour le reste
$R_N(z)$ dans:
\[ \mathrm{Log}(z) = (z-1) - \frac{(z-1)^2}2 + \frac{(z-1)^3}3 -
\dots + (-1)^{N-1}\frac{(z-1)^{N}}N + R_N(z) \;.\]  
  \item  On suppose $\Re(z)\geq\delta$ pour un certain
$\delta\in]0,1[$. Prouver :
\[|R_N(z)|\leq \frac1\delta\frac{|z-1|^{N+1}}{N+1}\]
On minorera $|1+t(z-1)|$ par $\delta$. 
  \item En déduire que la série de Taylor de $\mathrm{Log}$ au point $1$ est
uniformément convergente sur le compact $\{|z-1|\leq1,
\delta\leq\Re(z)\}$.
  \item Pour $-\pi<\phi<+\pi$ on pose $z = 1 +
e^{i\phi}$. Déterminer les coordonnées polaires $|z|$ et
$\mathrm{Arg}(z)$ de $z$ en fonction de $\phi$. Déduire de ce qui
précède les identités suivantes, pour tout
$\phi\in]-\pi,+\pi[$:
\begin{align*}
\log(2\cos \frac\phi2) &= \sum_{k=1}^\infty (-1)^{k-1}
\frac{\cos k\phi}k \\
\frac\phi2 &= \sum_{k=1} (-1)^{k-1} \frac{\sin k\phi}k
\end{align*}
et le fait que ces séries sont uniformément convergentes sur
tout intervalle $[-\pi+\epsilon,+\pi-\epsilon]$ ($0<\epsilon<\pi$).  
\end{enumerate}
\finenonce{002824}
\finexercice


\exercice{2825, burnol, 2009/12/15}
\enonce{002825}{}
\begin{enumerate}
  \item Soit $f$ une fonction continue sur $\overline{D(0,1)}$,
holomorphe sur $D(0,1)$, nulle sur le cercle de rayon
$1$. Montrer que $f$ est identiquement nulle.
  \item Plus fort: on
ne suppose plus que $f(e^{i\theta})$ est nulle pour tout
$\theta$ mais seulement pour $0\leq\theta\leq\pi$. Montrer
que $f$ est identiquement nulle. 
\emph{Indication :} $f(z)f(-z)$.  
\end{enumerate}
\finenonce{002825}
\finexercice


\exercice{2826, burnol, 2009/12/15}
\enonce{002826}{}
  Soit $\phi(z) = \frac{4z + 3}{4 +
3z}$. Montrer: $\forall\theta\in\Rr\quad  |\phi(e^{i\theta})|
= 1$. En déduire $|z|<1 \implies |\phi(z)|<1$.
\finenonce{002826}
\finexercice


\exercice{2827, burnol, 2009/12/15}
\enonce{002827}{}
Soit $F$ une fonction entière telle que $|F(z)|\leq \frac1n$
pour $|z| = n$, $n\geq1$. Montrer que $F$ est identiquement
nulle.
\finenonce{002827}
\finexercice


\exercice{2828, burnol, 2009/12/15}
\enonce{002828}{} 
\begin{enumerate}
  \item Soit $f$ analytique sur un disque $|z - z_0|\leq R$ et
telle qu'il existe un certain $z_1$ avec $|z_1
- z_0|<R$ tel que $|f(z)|>|f(z_1)|$ pour $|z-z_0| = R$.
Montrer que $f$ s'annule au moins une fois dans le
disque ouvert $D(z_0,R)$. 
\emph{Indication :} considérer sinon ce
que dit le principe du maximum pour la fonction
$\frac1f$.
\item \emph{Théorème de Hurwitz.} Soit $f_n$ des
fonctions holomorphes sur un voisinage commun $U$ de
$\overline{D(0,1)}$ qui convergent uniformément sur
$U$. Soit $F$ la fonction limite. On suppose que $F$ n'a
aucun zéro sur le cercle $|z|=1$, et qu'elle a au moins un
zéro dans le disque ouvert $D(0,1)$. Montrer  en appliquant
la question précédente à $f_n$ que pour $n\gg1$ la fonction
$f_n$ a au moins un zéro dans $D(0,1)$.\footnote{On verra
plus tard en cours ou en exercice que pour $n\gg1$ chaque
$f_n$ a, comptés avec leurs multiplicités, exactement le même
nombre de zéros que $F$ dans $D(0,1)$.} 
Ce résultat est
souvent appliqué sous sa forme réciproque: \emph{si des
fonctions holomorphes $f_n$ sans zéro convergent
uniformément sur un ouvert connexe vers $F$ alors soit $F$
est identiquement nulle soit $F$ n'a aucun zéro.} Justifier
cette dernière reformulation.
\end{enumerate}
\finenonce{002828}
\finexercice


\exercice{2829, burnol, 2009/12/15}
\enonce{002829}{}
Montrer que si une fonction entière $f$ a sa partie réelle
bornée supérieurement  alors elle est constante (considérer
$\exp(f)$). 
\finenonce{002829}
\finexercice


\exercice{2830, burnol, 2009/12/15}
\enonce{002830}{}
  Soit $f$ une fonction entière telle que $|f(z)|\leq
M\,(1+ |z|)^n$ pour un certain $M$ et un certain
$n\in\Nn$. Donner plusieurs démonstrations que $f$ est un
polynôme de degré au plus $n$:
\begin{itemize}
\item en utilisant une formule intégrale de Cauchy pour
  $f^{(n+1)}(z)$, avec comme contour les cercles de rayon
  $R$ centrés en l'origine, ou en $z$ si l'on veut,
\item en utilisant les formules de Cauchy pour
  $f^{(m)}(0)$, avec $m\geq n+1$,
\item en appliquant le théorème de Liouville à $(f(z) -
  P(z))/z^{n+1}$ avec $P$ le polynôme de McLaurin-Taylor à
  l'origine à l'ordre $n$.
\end{itemize} 
\finenonce{002830}
\finexercice


\exercice{2831, burnol, 2009/12/15}
\enonce{002831}{}
  Soit $f$ une fonction entière vérifiant $\lim_{|z|\to\infty}
|f(z)| = +\infty$. Donner plusieurs démonstrations que $f$
est un polynôme:
\begin{itemize}
\item en montrant, par un théorème du cours, que $w=0$ est
  une singularité polaire de $g(w) = f(\frac1w)$, et en en
  déduisant qu'il existe un polynôme $P$ tel que $f(z) -
  P(z)$ tende vers $0$ pour $|z|\to\infty$, puis Liouville,
\item ou en montrant que $f$ n'a qu'un nombre fini de zéros
  $z_j$, $1\leq j\leq n$, et en appliquant à $(z-z_1)\dots
  (z- z_n)/f(z)$ le résultat de l'exercice précédent, plus
  quelques réflexions de conclusion pour achever la preuve.
\end{itemize}
  Montrer que la fonction entière $z+e^z$ tend vers l'infini
le long de tout rayon partant de l'origine. D'après
ce qui précède $z+e^z$ est donc un
polynôme. Commentaires?
\finenonce{002831}
\finexercice


\exercice{2832, burnol, 2009/12/15}
\enonce{002832}{} Déterminer les séries de Laurent et les
résidus à l'origine des fonctions suivantes:
\begin{enumerate}
\item $f(z) = \frac1z$
\item $f(z) = \frac1{z^2+1}$
\item $f(z) = \frac1{z(z^2+1)}$
\end{enumerate}
\finenonce{002832}
\finexercice


\exercice{2833, burnol, 2009/12/15}
\enonce{002833}{} Déterminer la série de Laurent à l'origine de
la fonction analytique $\exp(\frac 1z)$, et son résidu à
l'origine. En $z_0\neq0$ quel est le résidu de cette fonction?
\finenonce{002833}
\finexercice


\exercice{2834, burnol, 2009/12/15}
\enonce{002834}{}
  Déterminer la partie singulière, le résidu, et
  le terme constant des séries de Laurent à
l'origine pour les fonctions:
\begin{enumerate}
\item $f(z) = \frac1{\sin z}$
\item $f(z) = \frac1{\sin z - \sh z}$
\item $f(z) = \frac1{z\sin(z)\sh(z)}$
\end{enumerate}
\finenonce{002834}
\finexercice


\exercice{2835, burnol, 2009/12/15}
\enonce{002835}{}
  Déterminer les séries de Laurent de $f(z) = \frac1{(z-1)(z-2)}$ dans chacune
des trois couronnes ouvertes  $0<|z|<1$, $1<|z|<2$,
$2<|z|<\infty$, ainsi que les séries de Laurent de $f$ aux
points $0$, $1$, $2$, et $3$. Quels sont les résidus en $z=0$, $z=1$,
$z=2$ et $z=3$?
\finenonce{002835}
\finexercice


\exercice{2836, burnol, 2009/12/15}
\enonce{002836}{}
  Montrer que tout lacet est homotopiquement trivial dans $\Cc$.
\finenonce{002836}
\finexercice


\exercice{2837, burnol, 2009/12/15}
\enonce{002837}{} 
Justifier les affirmations du polycopié
relatives à l'invariance de l'indice d'un lacet par rapport
à un point, lorsque l'on déforme continûment soit le lacet,
soit le point. Montrer que lorsque $\gamma$ est un lacet il
existe $R$ tel que $|z|>R\implies \mathrm{Ind}(\gamma,z) = 0$.
\finenonce{002837}
\finexercice


\exercice{2838, burnol, 2009/12/15}
\enonce{002838}{} 
\begin{enumerate}
  \item Soit $\gamma:[0,1]\to\Cc\setminus\{0\}$ un
lacet et soit $N\in \Zz$ son indice par rapport à $0$.  En
utilisant  la notion de variation de
l'argument, montrer qu'il existe une fonction
continue $g:[0,1]\to\Cc$ telle que $\forall t\quad \gamma(t)
= e^{g(t)}$ et  $g(1) - g(0) = 2\pi i N$.  Montrer
 que toute autre fonction continue $G$ avec
$\forall t\quad \gamma(t) = e^{G(t)}$  est de la
forme $g+2\pi i k$ pour un certain $k\in\Zz$. On pose
$h(t,u) = (1-u)\;2\pi i N\,t + ug(t)$ puis $H(t,u) =
e^{h(t,u)}$. Montrer que pour chaque $u\in[0,1]$
l'application $t\mapsto H(t,u)$ est un lacet. En déduire que
le lacet $c_N(t) = e^{2\pi i\; Nt}$ et $\gamma$ sont
homotopes dans $\Cc\setminus\{0\}$.
  \item On considère le lacet obtenu en suivant d'abord $c_N$ puis
  $c_M$. Montrer que ce lacet est homotope dans
  $\Cc\setminus\{0\}$ au lacet $c_{N+M}$ (il suffit de
  calculer son indice!).
\end{enumerate}
\finenonce{002838}
\finexercice


\exercice{2839, burnol, 2009/12/15}
\enonce{002839}{} 
On considère un lacet $\gamma:[a,b]
\to\Cc\setminus\{0\}$ (donc ne passant pas par
l'origine). On suppose qu'il n'existe qu'un nombre fini de
$t\in [a,b]$ avec $\gamma(t) \in \Delta = ]-\infty,0[$. On
les note $t_0<t_1<\dots <t_N$. Pour simplifier on supposera
que $\gamma(a)$ est sur $\Delta$, donc $t_0 =a$ et
$t_N=b$. Montrer que pour $t=t_j-\epsilon$, $\epsilon>0$
suffisamment petit, le signe $\mu_j$ de
$\Im(\gamma(t_j-\epsilon))$ ne dépend pas de $\epsilon$,
et de même pour le signe $\mu_j'$ de
$\Im(\gamma(t_j+\epsilon))$ (préciser ce que l'on fait
pour $j=0$ et $j=N$). 

 Si $\mu_j=+$ et $\mu_j' = -$ on dit
que $\gamma$ traverse $\Delta$ en $t=t_j$ dans le sens
direct, si $\mu_j=-$ et $\mu_j'=+$ on dit que $\gamma$
traverse $\Delta$ en $t=t_j$ dans le sens rétrograde. Sinon
on dit que $\gamma$ touche mais ne traverse pas $\Delta$. En
utilisant la relation entre la fonction $\mathrm{Log}(\gamma(t))$ et
la variation de l'argument de $\gamma(t)$ sur chaque
intervalle $]t_j,t_{j+1}[$, prouver $\Delta_{\gamma_j}
\arg(z) = \pi(\mu_{j+1} - \mu_j')$ avec $\gamma_j=\gamma$
restreint à $[t_j,t_{j+1}]$.

 En déduire que $\mathrm{Ind}(\gamma,0)$
est égal au nombre de valeurs de $t$ ($a$ et $b$ ne comptent
que pour un seul) pour lesquelles $\gamma$ traverse
$\Delta$, comptées positivement si la traversée est directe,
négativement si la traversée est rétrograde. 

Dans la
pratique vous pourrez  utiliser n'importe quelle demi-droite
issue de l'origine à la place de $\Delta$ à partir du moment
où elle n'intersecte le lacet $\gamma$ qu'en un nombre fini
de points (si on n'impose pas au lacet d'être régulier,
c'est-à-dire d'avoir un vecteur vitesse partout non nul,
alors il peut rester figé en un même point un certain temps,
et donc il faut modifier un petit peu la discussion
ci-dessus qui suppose qu'il n'y a qu'un nombre fini de
valeurs de $t$ pour lesquels $\gamma(t)$ est sur la
demi-droite).
\finenonce{002839}
\finexercice


\exercice{2840, burnol, 2009/12/15}
\enonce{002840}{}
\label{ex:burnol6.1}
  Justifier les formules suivantes :
lorsque $f$ présente en $z_0$ un pôle simple on a:
$$\mathrm{Res}(f,z_0) = \lim_{z\to z_0} (z-z_0)f(z)$$
Lorsque $f$ présente en $z_0$ un pôle d'ordre au plus $N$ on
a:
$$\mathrm{Res}(f,z_0) = \lim_{z\to z_0} \frac1{(N-1)!}\left(\frac{d}{dz}\right)^{N-1}(z-z_0)^N f(z)$$
\finenonce{002840}
\finexercice


\exercice{2841, burnol, 2009/12/15}
\enonce{002841}{}
\begin{enumerate}
  \item Soit $g$ une fonction analytique ayant un zéro simple en
$z_0$, et $f$ une autre fonction analytique définie dans un
voisinage de $z_0$. Montrer 
\[ \mathrm{Res}(\frac fg, z_0) =
\frac{f(z_0)}{g'(z_0)}\;.\] 
  \item On suppose que $g$ a un zéro d'ordre
$n$: $g(z_0+h) = h^n (c_0 + c_1 h + \dots)$, $c_0\neq0$, et l'on
écrit $f(z_0+h) = a_0 + a_1 h +\dots$. Montrer:
\[ \mathrm{Res}(\frac fg, z_0) = e_{n-1}\]
avec $e_0$, $e_1$, \dots, obtenus par  la division suivant les puissances
croissantes (comme dans les calculs de développement
limités):
\[ \frac{a_0 + a_1 h + a_2 h^2 + \dots}{c_0 + c_1 h + c_2
h^2 + \dots} = e_0 +
e_1 h + e_2 h^2 +\dots \;.\]
\end{enumerate}
\finenonce{002841}
\finexercice


\exercice{2842, burnol, 2009/12/15}
\enonce{002842}{}
\label{ex:burnol6.4}
  Soit $0<a<b<c$ et soit $C$ le cercle de rayon $r$ centré en
l'origine, parcouru dans le sens direct. Calculer $\int_C
\frac1{(z-a)(z-b)(z-c)}dz$ selon la valeur de $r$. On
donnera deux preuves, soit en utilisant le théorème des
résidus, soit en décomposant 
en éléments simples.
\finenonce{002842}
\finexercice


\exercice{2843, burnol, 2009/12/15}
\enonce{002843}{}
  Soit $\mathcal{R} = \{x_0\leq x\leq x_1, y_0\leq y\leq y_1\}$ un
  rectangle. En utilisant le théorème des résidus justifier
  la formule intégrale de Cauchy pour $z$ dans l'intérieur
  du rectangle et $f$ holomorphe sur le rectangle fermé:
\[ f(z) = \frac1{2\pi i}\int_{\partial\mathcal{R}}
\frac{f(w)}{w-z}dw\]
Démontrer ce résultat de manière plus
simple, directement à partir du théorème de Cauchy-Goursat pour les
fonctions holomorphes sur les rectangles, en utilisant la
fonction $w\mapsto (f(w)-f(z))/(w-z)$ (et aussi la notion
d'indice d'un lacet).  Dans le cas où $z$
est à l'\emph{extérieur} du rectangle $\mathcal{R}$, que vaut $\int_{\partial\mathcal{R}}
\frac{f(w)}{w-z}dw$?
\finenonce{002843}
\finexercice


\exercice{2844, burnol, 2009/12/15}
\enonce{002844}{} Soit $\Omega$ un domaine, de
bord le cycle $\partial\Omega$ orienté dans le sens direct. 
Soit $f$ une fonction holomorphe sur
$\overline\Omega$, soient $z_1$ et $z_2$ deux points de
$\Omega$. Que vaut
\[ \int_{\partial\Omega}
\frac{f(z)\,dz}{(z-z_1)(z-z_2)}\;?\]
Qu'obtient-on pour $z_2\to z_1$, $z_1$ fixé ?
\finenonce{002844}
\finexercice


\exercice{2845, burnol, 2009/12/15}
\enonce{002845}{}
Que vaut, en fonction de $R>0$:
\[ \int_{|z| = R} \frac{dz}{2z^2 - 5z + 2}\;?\]
On précisera les valeurs exclues de $R$.
\finenonce{002845}
\finexercice


\exercice{2846, burnol, 2009/12/15}
\enonce{002846}{}
Déterminer, $C$ désignant tour à tour le cercle $|z-i| = 1$,
ou le cercle $|z+i|=1$, ou encore $|z|=2$, parcourus dans le
sens direct, les valeurs des intégrales:
\[ \int_C \frac1{z^2+1}\,dz\]
Même question pour:
\[ \int_C \frac1{z^3-1}\,dz\qquad\text{et}\qquad\int_C
\frac1{z^4-1}\,dz\qquad\text{et}\qquad\int_C
\frac1{z^5-1}\,dz \] 
\finenonce{002846}
\finexercice


\exercice{2847, burnol, 2009/12/15}
\enonce{002847}{}
Que vaut $\int_{|z| = N} \tan(\pi z) \,dz$, pour $N\in\Nn$, $N\geq1$?
\finenonce{002847}
\finexercice


\exercice{2848, burnol, 2009/12/15}
\enonce{002848}{}
Déterminer pour $A, B, C$ réels, avec $A^2 > B^2 + C^2$ la
valeur de :
\[ \frac1{2\pi}\int_0^{2\pi} \frac{d\theta}{A + B\sin\theta + C\cos\theta}\]
On aura intérêt, comme première étape, à poser $B =
R\cos\phi$, $C=R\sin\phi$, mais on peut aussi se frotter
plus directement au résidu (utiliser bien sûr
$z=e^{i\theta}$ ou dans ce genre).
\finenonce{002848}
\finexercice


\exercice{2849, burnol, 2009/12/15}
\enonce{002849}{}
On considère dans le plan complexe un chemin fermé paramétré
$\gamma$ qui parcourt la figure ci-dessus dans le sens indiqué. 
$$  \centerline{\includegraphics{../images/img002849-1}} $$

Pour $j = 0, 1, 2, 3, 4$ on note \[A_j = \frac1{2\pi i} \int_\gamma \frac{dz}{z -z_j}\quad\text{et}\quad B_j = \frac1{2\pi i} \int_\gamma \frac{dz}{(z - z_j)^2}\] Déterminer, en le justifiant, les valeurs de $A_0$, $A_1$, $A_2$, $A_3$, $A_4$, et de $B_0$, $B_1$, $B_2$, $B_3$, $B_4$. On précisera aussi quel est le nom que l'on donne aux quantités données par les intégrales $A_j$, $j= 0\dots4$.
\finenonce{002849}



\finexercice


\exercice{2850, burnol, 2009/12/15}
\enonce{002850}{}
Soit $\gamma$ le contour, parcouru dans le sens direct, dessiné ci-contre. 
$$\centerline{\includegraphics{../images/img002850-1}}$$

Déterminer (avec justification) en fonction de $w_1$, $w_2$, $w_3$ les intégrales suivantes:
\begin{align*}
A &= \int_\gamma \frac{dz}{(z-w_1)(z-w_2)(z-w_3)}\\
B &= \int_\gamma \sin(z) dz\\
C &= \int_\gamma \frac{dz}{(z-w_1)^2(z-w_3)}
\end{align*}
\finenonce{002850}

\finexercice



\exercice{2853, burnol, 2009/12/15}
\enonce{002853}{}
Prouver pour $a>1$:
\[ \frac1{2\pi}\int_0^{2\pi} \frac{\sin\theta}{a+\sin\theta}d\theta =
\frac{\sqrt{a^2-1} - a}{\sqrt{a^2 - 1}}\]
En utilisant l'un des exercices précédents montrer que la
formule a un sens et est valable pour $a\in\Cc\setminus[-1,+1]$.
\finenonce{002853}
\finexercice


\exercice{2854, burnol, 2009/12/15}
\enonce{002854}{}
Que vaut en fonction de $R>0$
\[ \int_{|z| = R} \frac{z^2+1}{z^3 - z^2 - 4z + 4}dz\;?\]
\finenonce{002854}
\finexercice


\exercice{2855, burnol, 2009/12/15}
\enonce{002855}{}
Soit $P(z) = A z^4 + \dots$ un polynôme de degré au plus
$4$. Montrer que
$ \int_{|z| = R}
\frac{P(z)}{z^5-1}\,dz$ est indépendant de $R$ pour
$R>1$. En faisant tendre $R$ vers l'infini en déduire que
cette valeur constante est $2\pi iA$. Prouver alors via le
théorème des résidus:
$A = \frac15 \sum_{w^5 = 1} wP(w)$.
\finenonce{002855}
\finexercice


\exercice{2856, burnol, 2009/12/15}
\enonce{002856}{Résidu à l'infini}
\label{exo:residuinfini}
  Soit $f$ une fonction analytique pour
$\{|z|>R\}$. On pose:
\[ \mathrm{Res}(f,\infty) = - \frac1{2\pi i}\int_{C_r} f(z)dz\] avec
$C_r$ le cercle $\{|z| = r\}$ parcouru dans le sens
direct. Montrer que le terme de droite est bien indépendant
de $r>R$. On notera le signe $-$. On dit que
$\mathrm{Res}(f,\infty)$ est le ``résidu à l'infini'' de
$f$. Soit $f$ une fonction holomorphe sur $\Cc$ à
l'exception d'un nombre fini de singularités
isolées. Montrer le théorème suivant: \emph{la somme de tous
les résidus (y compris celui à l'infini) de $f$ est nulle}.
\finenonce{002856}
\finexercice


\exercice{2857, burnol, 2009/12/15}
\enonce{002857}{} 
Soit $f$ une fonction holomorphe sur
$\overline\Omega\setminus\{z_1, \dots, z_N\}$, avec $\Omega$
le domaine intérieur à une courbe de Jordan $\gamma$. Soit
$g_n(z)$ la partie principale (partie singulière) de $f$ en
la singularité isolée $z_n$. Prouver \emph{la formule
intégrale générale de Cauchy}:
\[\forall z\in
\Omega\setminus\{z_1,\dots,z_N\}\qquad f(z) = \sum_{1\leq
n\leq N} g_n(z) + \frac1{2\pi i
}\int_{\partial\Omega}\frac{f(w)}{w-z}dw\]
Pour cela, remarquer d'abord $\mathrm{Res}(\frac{f(w)}{w-z},z_n)=
 \mathrm{Res}(\frac{g_n(w)}{w-z},z_n)$; puis montrer que le résidu à
 l'infini de la fonction $\frac{g_n(w)}{w-z}$  de $w \in
 \Cc\setminus\{z_n\}$, 
 est nul. On pourra utiliser l'exercice \ref{exo:residuinfini}.
\finenonce{002857}
\finexercice


\exercice{2858, burnol, 2009/12/15}
\enonce{002858}{Morceaux de Résidus}
\label{ex:burnol2.1}
Soit $f$ présentant en $z_0$ un \emph{pôle simple}. Soit
$C_r(\alpha,\beta)$ l'arc de cercle $w= z_0 + r
e^{i\theta}$, $\alpha\leq \theta\leq\beta$, parcouru dans le
sens direct des $\theta$ et avec $0<\beta - \alpha\leq
2\pi$. Prouver:
\[ \lim_{r\to0} \int_{C_r(\alpha,\beta)} f(z)\,dz = 2\pi
i\;\frac{\beta-\alpha}{2\pi}\;\mathrm{Res}(f,z_0)\] 
Que se passe-t-il si le pôle est d'ordre plus élevé? 
\finenonce{002858}
\finexercice


\exercice{2859, burnol, 2009/12/15}
\enonce{002859}{Lemme de Jordan}
\label{ex:burnol2.2}
 Soit $f$ une fonction définie
et continue sur le domaine $\{\Im(z)>0, |z|>R\}$, ou
seulement sur une suite de demi-cercles $\{\Im(z)>0,
|z|=R_n\}$ de rayons tendant vers l'infini. On
suppose $\lim_{\substack{|z|\to\infty \\ \Im(z)>0}} |f(z)| = 0$
(ou $\displaystyle\lim_{n\to\infty}
\sup_{\Im(z)>0,\,|z|=R_n} |f(z)| = 0$.)  Montrer (on
utilisera  $\sin(\theta)\geq \frac2\pi \theta$ pour
$0\leq\theta\leq\frac\pi2$):
\[ \lim_{R\to\infty} \int_{z = R e^{i\theta},\,
0<\theta<\pi} f(z)e^{iz}\,dz = 0\qquad (\text{ou l'analogue avec
les\ }R_n)\]
\finenonce{002859}
\finexercice


\exercice{2860, burnol, 2009/12/15}
\enonce{002860}{}
  En considérant
l'intégrale de $\frac{e^{iz}}z$ sur un contour allant de
$-R$ à $+R$ le long de l'axe réel en contournant $0$ par un
petit demi-cercle, puis qui revient de $+R$ à $-R$ par le
demi-cercle dans le demi-plan supérieur, démontrer
$\int_0^\infty \frac{\sin x}x\,dx = \frac\pi2$. 
\finenonce{002860}
\finexercice


\exercice{2861, burnol, 2009/12/15}
\enonce{002861}{}
Déterminer les intégrales (semi-convergentes) de Fresnel
$\int_0^\infty \cos(x^2)dx$ et $\int_0^\infty \sin(x^2)dx$
en considérant l'intégrale de $\exp(-z^2)$ sur le contour
$z=x$, $0\leq x\leq R$, $z = R\exp(i\theta)$, $0\leq
\theta\leq\frac\pi4$, $z = xe^{i\frac\pi4}$, $R\geq
x\geq0$. On rappelle l'identité $\int_\Rr \exp(-\pi u^2)du
=1$.  
\finenonce{002861}
\finexercice


\exercice{2862, burnol, 2009/12/15}
\enonce{002862}{}
Que vaut $\Gamma(\frac12) = \int_0^\infty
\frac{e^{-t}}{\sqrt t}dt$? (faire un changement de variable
$t=\pi u^2$ pour se ramener à la Gaussienne). En considérant
un contour passant par l'axe réel, puis un quart
de cercle, puis l'axe imaginaire, puis un petit quart de
cercle évitant l'origine prouver:
\[ \int_0^\infty
\frac{e^{-t}}{\sqrt t}dt =
\exp(i\frac\pi4)\int_0^\infty\frac{e^{-ix}}{\sqrt x}dx\] et
en déduire les valeurs des intégrales  $\int_0^\infty
\frac{\cos x}{\sqrt x}dx$ et $\int_0^\infty \frac{\sin
x}{\sqrt x}dx$ (qui ne sont que semi-convergentes).
Comparer aux intégrales de Fresnel.  
\finenonce{002862}
\finexercice


\exercice{2863, burnol, 2009/12/15}
\enonce{002863}{}
 Reprendre l'exercice précédent et déterminer
pour $0<a<1$ les valeurs des intégrales (semi-convergentes)
\[ \int_0^\infty \frac{\cos x}{x^a}dx\qquad\text{et}\qquad\int_0^\infty
\frac{\sin x}{x^a}dx\] en utilisant la fonction Gamma. \`A
propos prouver que ces intégrales ne sont que
semi-convergentes (\textit{i.e.} pas absolument
convergentes). 
\finenonce{002863}
\finexercice


\exercice{2864, burnol, 2009/12/15}
\enonce{002864}{}
Confirmer par le calcul des résidus la valeur connue
($\Arctan\dots$!):
\[\int_\Rr \frac{dx}{1 + x^2} = \pi\]
On appliquera le théorème des résidus au contour direct
comportant le segment $[-R,+R]$ et le semi-cercle de rayon
$R$ dans le demi-plan supérieur, pour $R\to+\infty$.
\finenonce{002864}
\finexercice


\exercice{2865, burnol, 2009/12/15}
\enonce{002865}{}
Justifier $\int_\Rr \frac{e^{i\xi x}}{1 + x^2}dx = \int_\Rr
\frac{\cos(\xi x)}{1 + x^2}dx$ pour $\xi\in\Rr$. Prouver
par un calcul de résidu
\[\int_\Rr \frac{e^{i\xi x}}{1 + x^2}dx = \pi e^{-|\xi|}\;.\]
Suivant le cas $\xi\geq0$ ou $\xi<0$ on complètera le
segment $[-R,+R]$ par un semi-cercle dans le demi-plan supérieur,
ou inférieur, afin que la contribution du
semi-cercle tende vers $0$ pour $R\to\infty$. On peut aussi
observer que l'intégrale est une fonction paire de $\xi$ et que l'on peut
donc se restreindre à $\xi\geq0$.
\finenonce{002865}
\finexercice


\exercice{2866, burnol, 2009/12/15}
\enonce{002866}{}
Prouver, pour tout $x\in \Rr$: 
\[\frac1{2\pi}\int_\Rr e^{-i\xi x}(\pi e^{-|\xi|})d\xi =
\frac1{1+x^2}\;.\] Il suffit d'évaluer séparément
$\int_{-\infty}^0$ et $\int_0^\infty$ en utilisant le fait
que $\exp$ est sa propre primitive (ce calcul n'utilise donc
pas la notion de fonction analytique et le théorème des
résidus). On remarquera que l'on retombe sur la fonction
$1/(1+x^2)$, ce qui n'est pas un hasard (formule d'inversion
pour les transformations intégrales de Fourier).
\finenonce{002866}
\finexercice


\exercice{2867, burnol, 2009/12/15}
\enonce{002867}{}
Déterminer \[\int_\Rr \frac{1}{1 + x^4}\,dx %
\qquad \int_\Rr \frac{1+x^2}{1 + x^4}\,dx
\qquad \int_\Rr \frac1{1+x^2+x^4}\,dx\]
\finenonce{002867}
\finexercice


\exercice{2868, burnol, 2009/12/15}
\enonce{002868}{}
Préciser pourquoi $\int_\Rr \frac{e^{i\xi x}}{1 + x^4}dx$
est une intégrale convergente pour $\xi\in\Rr$, est une
fonction réelle et paire de $\xi$, et utiliser un calcul de
résidus pour établir, pour $\xi\geq0$:
\[ \int_0^{+\infty} \frac{\cos(\xi x)}{1+x^4}dx = \frac\pi2
e^{-\xi/\sqrt2}\sin(\frac\xi{\sqrt2} + \frac\pi4)\]
Cette formule est-elle valable pour $\xi<0$?
\finenonce{002868}
\finexercice


\exercice{2869, burnol, 2009/12/15}
\enonce{002869}{}
\begin{enumerate}
  \item Déterminer \[\int_0^{+\infty} \frac{dx}{1 + x^3}\] Pour ce
calcul, on considérera le contour allant le long de l'axe
réel de $0$ à $R$ puis de $R$ à $j R$ le long d'un cercle
puis de $jR$ à $0$ par un segment ($j =
\exp(i\frac{2\pi}3)$). On écrira d'une part chacune des trois
contributions à l'intégrale de contour, en faisant attention
au sens de parcours, et l'on utilisera  d'autre part le
théorème des résidus.
  \item On note, pour $|w| = 1$ et certaines
valeurs spéciales de $w$ (que l'on précisera) étant exclues,
$J(w)$ l'intégrale $\int \frac{dz}{1+z^3}$ le long du segment
infini $w\Rr^+$. Déterminer $J(w)$ en fonction de $w$.
\end{enumerate}
\finenonce{002869}
\finexercice

\exercice{2879, burnol, 2009/12/15}
\enonce{002879}{}
\begin{enumerate}
  \item Prouver pour $n\in\Nn$, $n>1$:
\[ \int_0^\infty \frac{dx}{1+x^n} = \frac{\pi/n}{\sin
  (\pi/n)}\] en utilisant le secteur angulaire $0\leq \mathrm{Arg}
z\leq \frac{2\pi}n$, $0\leq |z|\leq R$, $R\to+\infty$, et en
montrant que la contribution de l'arc de cercle tend vers
zéro pour $R\to+\infty$. 
  \item   Montrer, en utilisant les contours
$\epsilon\leq x\leq R$, $z= Re^{i\theta}$ ($0\leq \theta
\leq \frac{2\pi}a$), $z=r e^{i\frac{2\pi}a}$ ($R\geq r\geq
  \epsilon)$, $z = \epsilon e^{i\theta}$ ($\frac{2\pi}a \geq \theta
\geq 0$):
\[ a\in\Rr,\; a>1\quad\implies\qquad\int_0^\infty \frac{dx}{1+x^a} = 
  \frac{\pi/a}{\sin (\pi/a)}\;.\]
 Pour définir $z^a $ comme
fonction holomorphe 
sur $\{ z = r e^{i\alpha}\,|\, 0<r<\infty,\;0\leq \alpha
\leq \frac{2\pi}a\}$, on pose $z^a = r^a
e^{a\,i\alpha} = \exp(a(\log r+i\alpha))$ (car $\log r + i\alpha = \mathrm{Log}(z
e^{-i\frac\pi a}) + i\frac\pi a$; no comments).
  \item Soit $J(a) = \int_0^\infty
\frac{dx}{1+x^a}$; justifier que l'intégrale définissant
$J(a)$ est convergente et analytique comme fonction de $a$
pour $\Re(a)>1$ et prouver $J(a)
= \frac{\pi/a}{\sin (\pi/a)}$.
  \item  On définit maintenant 
\[ K(p) = \int_{-\infty}^{+\infty}
\frac{e^{pt}}{1+e^t}dt\]
 pour $0<p<1$.  Justifier les identités (pour $0<p<1$):
\[ K(p) = \int_{-\infty}^{+\infty} \frac{e^{pt}}{1+e^t}dt =
\int_0^{+\infty} \frac{t^{p-1}}{1+t}dt = \frac1p \int_0^{+\infty}
\frac{dt}{1+t^{1/p}} = \frac1p J(\frac1p) = \frac\pi{\sin(\pi p)}
\] 
  \item Expliquer pourquoi l'intégrale $K(p) = \int_{-\infty}^{+\infty}
\frac{e^{pt}}{1+e^t}dt$ est convergente et analytique pour
$p$ complexe avec $0<\Re(p)<1$ et établir la formule
$K(p) = \frac\pi{\sin(\pi p)}$ pour  $0<\Re(p)<1$.
  \item Donner une preuve simple directe de la
formule  $K(p) = \frac\pi{\sin(\pi p)}$ pour tout $p$
complexe avec $0<\Re(p)<1$ en appliquant le théorème des
résidus avec des contours liés aux  droites $z=x$,
$x\in\Rr$ et $z = x+2\pi i$, $x\in\Rr$.
  \item Déduire de ce qui précède avec $p = \frac12 + i\xi$,
$\xi\in\Rr$:
\[ \int_{-\infty}^{+\infty} \frac{\cos(\xi t)}{\ch(t/2)}dt
= \frac{2\pi}{\ch(\pi\xi)}\;,\]
Montrer  que  la transformation de Fourier $\widehat f(\xi)=\int_\Rr
e^{2\pi i \xi x} f(x)\,dx$ appliquée à la fonction $f(x) = \frac1{\ch(\pi x)}$
donne simplement $\widehat f = f$ (remarque: c'est aussi le cas
avec $f(x) = e^{-\pi x^2}$).
  \item On revient à la formule générale $K(p)= \frac\pi{\sin(\pi p)}$. En 
séparant parties réelles et imaginaires dans $\int_{-\infty}^{+\infty}
\frac{e^{pt}}{1+e^t}dt$ déterminer (en simplifiant le plus
possible) les valeurs de :
\[ \int_{-\infty}^{+\infty} \frac{e^{u
t}\cos(vt)}{1+e^t}dt\qquad,\qquad
\int_{-\infty}^{+\infty} \frac{e^{u
t}\sin(vt)}{1+e^t}dt\;,\]
pour $0<u<1$, $v\in\Rr$.
\end{enumerate}
\finenonce{002879}
\finexercice

\exercice{2880, burnol, 2009/12/15}
\enonce{002880}{}
  Déterminer 
$\displaystyle \int_{-\infty}^{+\infty}
\frac1{(1+x^2)(2+e^{ix})}\,dx$.
\finenonce{002880}
\finexercice


\exercice{2881, burnol, 2009/12/15}
\enonce{002881}{}
  Déterminer
$\displaystyle \int_{-\infty}^\infty
e^{ix}\;\frac{x-i}{x+i}\;\frac1{x^2+1}\,dx$ et $\displaystyle
\int_{-\infty}^\infty
e^{-ix}\;\frac{x-i}{x+i}\;\frac1{x^2+1}\,dx$.
\finenonce{002881}
\finexercice


\exercice{2882, burnol, 2009/12/15}
\enonce{002882}{} Déterminer $ \displaystyle
\int_{-\infty}^\infty \frac{\sin(x)}{x(x^2+1)}\,dx$.
\finenonce{002882}
\finexercice

\exercice{6587, hueb, 2011/10/16}
\enonce{006587}{}
\begin{enumerate}
\item Démontrer que
  $$\int_{\gamma} \frac{dz}{(z-a)^n} = \begin{cases}
    2\pi i \quad &n=1,\\
    0 &n=2,3,4,\dots \end{cases}$$
  $\gamma$ étant une courbe
  fermée simple ayant $a$ dans son intérieur et orientée
  positivement.
\item Quelle est la valeur de l'intégrale si $n=0,-1,-2, \dots$?
\end{enumerate}
\finenonce{006587}



\finexercice


\exercice{6588, hueb, 2011/10/16}
\enonce{006588}{} Evaluer
$$\int_{\gamma} \frac{5z^2 -3z +2}{(z-1)^3} dz$$
où $\gamma$ est une
courbe fermée simple quelconque entourant $z=1$ et orientée
positivement.  \finenonce{006588}



\finexercice


\exercice{6589, hueb, 2011/10/16}
\enonce{006589}{} Evaluer
\begin{enumerate}
\item[(a)]
  $$\int_{\gamma} \frac{\cos z}{z-\pi} dz, $$
\item[(b)]
  $$\int_{\gamma} \frac{e^z}{z(z+1)} dz, $$
\end{enumerate}
où
\begin{enumerate}
\item[(i)] $\gamma$ est le cercle positif $\{z; |z-1| = 3\}$,
\item[(ii)] $\gamma$ est le cercle positif $\{z; |z-1|=2.1\}$.
\end{enumerate}
\finenonce{006589}



\finexercice


\exercice{6590, hueb, 2011/10/16}
\enonce{006590}{} Evaluer
\begin{enumerate}
\item
  $$
  \int_{\gamma}\frac{\sin \pi z^2 +\cos \pi z^2}{(z-1)(z-2)} dz, $$
\item $$\int_{\gamma} \frac{e^{2z}}{(z+1)^4} dz, $$
\end{enumerate}
où $\gamma$ est le cercle positif $\{z; |z| = 3\}$.  \finenonce{006590}



\finexercice


\exercice{6591, hueb, 2011/10/16}
\enonce{006591}{} Démontrer que
$$\int_0^{2\pi} \cos^{2n} \vartheta d\vartheta = 2 \pi \cdot \frac{1
  \cdot 3 \cdot 5 \cdot \dots \cdot (2n-1)} {2 \cdot 4 \cdot 6 \cdot
  \dots \cdot (2n)}, \quad n=1,2,3, \dots .$$
\finenonce{006591}



\finexercice


\exercice{6592, hueb, 2011/10/16}
\enonce{006592}{} Localiser les
singularités de chacune des fonctions suivantes et les
caractériser.
\begin{enumerate}
\item $\frac {z^2}{(z+1)^3}$,
\item $\frac {2z^3-z+1}{(z-4)^2 (z-i)(z-1+2i)}$,
\item $\frac {\sin mz}{z^2 +2z +2}$,
\item $\frac {1-\cos z}z$,
\item $e^{-\frac 1{(z-1)^2}}$,
\end{enumerate}
\finenonce{006592}



\finexercice


\exercice{6593, hueb, 2011/10/16}
\enonce{006593}{} Trouver les séries de
Laurent par rapport aux singularités indiquées pour chacune des
fonctions suivantes. Caractériser la singularité dans chaque cas
et donner le domaine de convergence de chaque série.
\begin{enumerate}
\item $\frac {e^z}{(z-1)^2}, \qquad z=1,$
\item $z \cos \frac 1z,\qquad z=0,$
\item $\frac {\sin z}{z -\pi},\qquad z=\pi,$
\item $\frac z{(z+1)(z+2)},\qquad z=-1.$
\end{enumerate}
\finenonce{006593}



\finexercice


\exercice{6594, hueb, 2011/10/16}
\enonce{006594}{} Déterminer les résidus
de chacune des fonctions suivantes, au p\^ole indiqué.
\begin{enumerate}
\item $\frac {z^2}{(z-2)(z^2+1)}, \qquad z=2,\ z=i,\ z=-i,$
\item $\frac 1{z(z+2)^3},\qquad z=0,\ z=-2,$
\item $\frac {z e^{zt}}{(z -3)^2},\qquad z=3,$
\item $\textrm{cotg} z,\qquad z=-5 \pi.$
\end{enumerate}
\finenonce{006594}



\finexercice


\exercice{6595, hueb, 2011/10/16}
\enonce{006595}{} Trouver les séries de
Laurent de $\frac {z^2}{(z-2)(z^2+1)}$ par rapport à ses poles.
\finenonce{006595}



\finexercice


\exercice{6596, hueb, 2011/10/16}
\enonce{006596}{} Evaluer
$$\int_{\gamma} \frac {e^z}{(z-1)(z+3)^2} dz$$
où $\gamma$ est le
cercle positif donné par
\begin{enumerate}
\item $|z|= \frac 32$,
\item $|z|= 10$.
\end{enumerate}
\finenonce{006596}



\finexercice


\exercice{6597, hueb, 2011/10/16}
\enonce{006597}{} Evaluer
$$\int_0^{2\pi} \frac {d\theta}{5+3 \sin \theta} .$$
\finenonce{006597}



\finexercice


\exercice{6598, hueb, 2011/10/16}
\enonce{006598}{} 
Etude de la fonction
holomorphe $f(z)= \cos z$.
\begin{enumerate}
\item Trouver l'image d'une droite $x=c$.
\item Trouver l'image d'une droite $y=c$.
\item Trouver un ouvert maximal $U$ tel que le restriction de $f$ à
  cet ouvert soit injective.
\item Trouver un domaine de définition maximal pour la fonction
  réciproque, $\textrm{arccos}$.
\item Vérifier que parmi les branches de $\textrm{arccos}$ il y en a
  deux qui s'écrivent sous la forme $\textrm{arccos} (w) = \pm i
  \textrm{Log}(w + \sqrt {w^2-1})$.
\item Trouver toutes les branches de la fonction $\textrm{arccos}$.
\end{enumerate}
\finenonce{006598}



\finexercice


\exercice{6599, hueb, 2011/10/16}
\enonce{006599}{} 
Vérifier que
$$\int_{-\infty}^{\infty} \frac{x^2 dx}{(x^2+1)^2(x^2+2x+2)} = \frac{7
  \pi}{50}.$$
\finenonce{006599}



\finexercice


\exercice{6600, hueb, 2011/10/16}
\enonce{006600}{} 
Déterminer les
singularités de la fonction $f(z)= \frac 1z - \frac 1{\sin z}$ et
les classer.  \finenonce{006600}



\finexercice


\exercice{6601, hueb, 2011/10/16}
\enonce{006601}{} 
Démontrer que
\begin{enumerate}
\item $\mathcal{L}(1) (s)= \frac 1 s$
\item $\mathcal{L}(t^n)(s) = \frac {n!} {s^{n+1}}$
\item $\mathcal{L}(e^{at})(s) = \frac 1 {s-a}$
\item $\mathcal{L}(\sin at)(s) = \frac a {s^2+a^2}$
\item $\mathcal{L}(\cos at)(s) = \frac s {s^2+a^2}$
\item $$\mathcal{L}(U_a)(s) = \frac {e^{-as}} s, \ 
\text{où}\  U_a(t) = \begin{cases} 1,&t\geq a,\\
    0,&t < a,\\
         \end{cases}
         a \in \Rr, a \geq 0.$$
\end{enumerate}
Préciser les domaines de définition de ces transformées de
Laplace.  \finenonce{006601}



\finexercice


\exercice{6742, queffelec, 2011/10/16}
\enonce{006742}{}
\begin{enumerate}
\item Calculer les résidus aux différents p\^oles de la fonction $f(z)={z^2+z+1\over
z(z^2+1)^2};\ f(z)=\displaystyle{1\over 1+z+\cdots+z^{n-1}};\ f(z)={1\over
e^z-1}-{1\over z}$.

\item Montrer que si $f(z)=(z-a)^{-n} g(z)$ où $g$ est holomorphe dans $\Omega$
ouvert contenant $a$, alors Res$(f,a)=\displaystyle{g^{(n-1)}(a)\over (n-1)!}$. 
Trouver les p\^oles et résidus des fonctions suivantes :
$\displaystyle{1-\cos z\over z^3},\ \displaystyle{e^{2z}\over (z-1)^3},\
\displaystyle{1\over (1+z^2)^n}\ $.
\end{enumerate}
\finenonce{006742}



\finexercice


\exercice{6743, queffelec, 2011/10/16}
\enonce{006743}{Intégrales de fonctions trigonométriques}
Calculer $\displaystyle\int_0^{2\pi} {\cos 2t\over 3-2\cos t}\ dt$ (poser
$z=e^{it}$).
\finenonce{006743}



\finexercice


\exercice{6744, queffelec, 2011/10/16}
\enonce{006744}{Intégrales de fractions rationnelles sans p\^ole sur $\Rr$}
\begin{enumerate}
\item Calculer $\displaystyle\int_{\Rr}{dx\over1+x^{2n}}$, $n\geq2$, en
intégrant
$\displaystyle{1\over1+z^{2n}}$ sur un demi-cercle.

\item Calculer $\displaystyle\int_{\Rr}{dx\over(1+x^2)^n}$.
\end{enumerate}
\finenonce{006744}



\finexercice


\exercice{6745, queffelec, 2011/10/16}
\enonce{006745}{Utilisation du logarithme}
\begin{enumerate}
\item Calculer  $\displaystyle\int_0^{+\infty}{dx\over1+x^{3}}$, en intégrant
$\displaystyle{\log z\over1+z^3}$ sur un cercle privé de $\Rr^+$, $\hbox{log}$
dési\-gnant ici la détermination du logarithme avec coupure sur $\Rr^+$.

\item En choisissant la même détermination du logarithme et le même contour,
calculer simutanément les intégrales $I=\displaystyle\int_{\Rr}
{dx\over1+x^4}$ et $J=\displaystyle\int_{\Rr}{\hbox{Ln} x\over1+x^4}\ dx$.
(Intégrer cette fois $\displaystyle{(\log z)^2\over1+z^4}$.)

\item Calculer $\displaystyle\int_0^{+\infty}{dx\over1+x^{n}}$ et
$\displaystyle\int_0^{+\infty}{\hbox{Ln} x\over1+x^{n}}\ dx$, $n\geq2$, en
intégrant
$\displaystyle{\log z\over1+z^{n}}$ sur un secteur épointé bien choisi.
\end{enumerate}
\finenonce{006745}



\finexercice


\exercice{6746, queffelec, 2011/10/16}
\enonce{006746}{septembre 1999}
Soit $a$ un réel tel que $0\leq a<1$
\begin{enumerate}
\item Démontrer que les intégrales $I(a)=\int_0^{+\infty}{\sinh(ax)\over\sinh(x)}\
dx$ et $J(a)=\int_0^{+\infty}{\cosh(ax)\over\cosh(x)}\ dx$ sont convergentes.

\item Soit $\epsilon$ et $R$ des réels tels que $0<\epsilon<{\pi\over2}<R$,
$K_{\epsilon,R}\subset\Cc$ le compact obtenu en \^otant du rectangle de
sommets $R,R+i{\pi\over2},R+i{\pi\over2},-R$, la demi-boule ouverte de centre
$0$ et de rayon $\epsilon$, et $f(z)=\displaystyle{e^{az}\over e^z-e^{-z}}$.
  \begin{enumerate}

  \item Montrer que $\lim _{R\to+\infty} \int_\gamma f(z)\ dz=0$ lorsque
  $\gamma$ est le segment $[R,R+i{\pi\over2}]$; puis le segment
  $[-R+i{\pi\over2},-R]$.

  \item Calculer $\int_{\partial K_{\epsilon,R}} f(z)\ dz$ et la limite
  lorsque $\epsilon\to 0$ et $R\to+\infty$; en déduire les expressions de $I(a)$
  et $J(a)$.
  \end{enumerate}
\end{enumerate}
\finenonce{006746}



\finexercice


\exercice{6747, queffelec, 2011/10/16}
\enonce{006747}{}
Soit $a$ un réel $>0$ et $
\hbox{ Log}z$ la détermination principale du log sur
$\Cc\backslash R^-$. En intégrant la fonction
$f(z)={1\over (z^2+a^2)\hbox{ Log}z}$ sur un contour $\Gamma_{\varepsilon, R}$
constitué du cercle de rayon $R$ évitant le demi-axe $\Rr^-$, montrer que 
$$\int_0^\infty {dx\over (x^2+a^2)(\hbox {Ln}^2x+\pi^2)}={\pi\over2a(\hbox
{Ln}^2a+\pi^2/4)}-{1\over1+a^2}.$$
\finenonce{006747}



\finexercice


\exercice{6748, queffelec, 2011/10/16}
\enonce{006748}{Calcul d'intégrales semi-convergentes}
\begin{enumerate}
\item Calculer $\displaystyle\int_{\Rr}{\sin
x\over x}\ dx$ en intégrant $f(z)=\displaystyle{e^{iz}\over z}$ sur un contour
bien choisi.  

\item Calculer de même $\displaystyle\int_{\Rr}{x^2-1\over x^2+1}{\sin x\over x}\
dx$.
\end{enumerate}
\finenonce{006748}



\finexercice


\exercice{6749, queffelec, 2011/10/16}
\enonce{006749}{Calcul de transformées de Fourier}
\begin{enumerate}
\item
 Calculer la transformée de Fourier de $\displaystyle{1\over1+x^4}$ en
intégrant
$\displaystyle e^{itz}\over1+z^4$ sur un demi-cercle dans un demi-plan bien
choisi.

\item Calculer $\displaystyle I(m,a)=\int_0^{+\infty}{\cos mx\over(1+x^2)(x^2+a^2)}\
dx$ en distinguant les cas 
$a=1, a\not=1$. Vérifier que
$I(m,1)=\lim_{a\to1}I(m,a)$. 

En déduire sans nouveaux calculs la valeur de
$\displaystyle\int_0^{+\infty}{x\sin mx\over(1+x^2)^2}\ dx$. 
\end{enumerate}
\finenonce{006749}



\finexercice


\exercice{6750, queffelec, 2011/10/16}
\enonce{006750}{}
Soit $P$ et $Q$ deux polyn\^omes tels que $\deg Q>\deg P$.
\begin{enumerate}
\item Exprimer $\sum\hbox{Res}({P\over Q})$ à l'aide des coefficients de $P$ et $Q$.

\item Soit $P(z)=z^n+\sum_0^{n-1}a_kz^k$ un polyn\^ome dont toutes les racines sont
dans $D(0,R)$. Montrer que $f(x)=\displaystyle{1\over
2i\pi}\int_{\Gamma_R}{e^{xz}\over P(z)}\ dz$ est la solution de l'équation
différentielle d'ordre $n$, $y^{(n)}+ a_{n-1}y^{(n-1)}+\cdots+a_0y=0$ de
condition initiale $y^{(j)}(0)=0$ si $j<n-1$ et $y^{(n-1)}(0)=1$.
\end{enumerate}
\finenonce{006750}



\finexercice


\exercice{6751, queffelec, 2011/10/16}
\enonce{006751}{} Calculer

\begin{enumerate}  
\item $\displaystyle \int_0^{2\pi}{dt\over a+\cos t}$ ($a>1$) ;

\item $\displaystyle \int_0^{2\pi}{dt\over (a+b\cos t)^2}$ ($a>b>0$) ;

\item $\displaystyle \int_0^{2\pi}{\cos^3{3t}\over 1-2a\cos t+a^2}\ dt$
($a\in\C\setminus \{\pm 1\}$) ;

\item $\displaystyle \int_0^{2\pi}\exp{(\cos t)}\cos{(nt-\sin t)}\ dt$
($n\in {\Zz}$).
\end{enumerate}
\finenonce{006751}



\finexercice       


\exercice{6752, queffelec, 2011/10/16}
\enonce{006752}{} 
Dans cet exercice, on justifiera soigneusement chaque passage à la
limite. Calculer les intégrales :

\begin{enumerate} 
\item $\displaystyle \int_{-\infty}^{+\infty}{1\over 1+x^6}\ dx$

\item $\displaystyle \int_0^{+\infty}{1\over x^\alpha (1+x)}\ dx\ ,\ 0<\alpha <1$

\item $\displaystyle \int_0^{+\infty}{\sin^2x\over x^2}\ dx$
(on pourra considérer la fonction $\displaystyle {1-e^{2ix}\over x^2}$)

\item $\displaystyle \int_{-\infty}^{+\infty}{\cos x\over e^x+e^{-x}}\ dx$
(on pourra utiliser le rectangle de sommets $-R$, $R$, $R+i\pi$,
$-R+i\pi$)

\item $\displaystyle \int_0^{+\infty}{\mathrm{Log}^2 x\over 1+x^2}\ dx$

\item $\displaystyle \int_0^{+\infty}{\mathrm{Log} x\over (1+x)^3}\ dx$
\end{enumerate}
\finenonce{006752}



\finexercice  


\exercice{6753, queffelec, 2011/10/16}
\enonce{006753}{}
Soit $f(z) = z^5 +5z^3+z-2$. Montrer que $f$ a trois de  ses
zéros dans le disque $D(0,1)$ et tous ses zéros dans le disque $D(0,3)$.
\finenonce{006753}



\finexercice


\exercice{6754, queffelec, 2011/10/16}
\enonce{006754}{}
Soit $c$ un nombre complexe vérifiant $|c|>e$. Montrer que
l'équation $e^z=cz$ a une solution et une seule dans $H = \{\,z\in 
\Cc\mid \Re\!e z<1\,\}$.
(On pourra considérer $D_r = H \cap D(1,r)$ avec $r\ge2$ et les fonctions
$e^z-cz$ et $-cz$.)
\finenonce{006754}



\finexercice


\exercice{6755, queffelec, 2011/10/16}
\enonce{006755}{}
En utilisant le théorème de Rouché, démontrer le
théorème de d'Alembert.
\finenonce{006755}



\finexercice


\exercice{6756, queffelec, 2011/10/16}
\enonce{006756}{}
Soit $P(z) = z^n+a_1z^{n-1} + \cdots +a_n$, $n\ge 1$, $a_j\in
\Cc$. Montrer qu'il existe un point $c$ de $\partial D(0,1)$ tel que l'on
ait $|P(c)| \ge 1$.
\finenonce{006756}



\finexercice


\exercice{6757, queffelec, 2011/10/16}
\enonce{006757}{}
Montrer que, si $f$ est holomorphe au voisinage de
$\overline{D(0,1)}$ et si $f(\partial D(0,1)) \subset D(0,1)$, alors il existe
un $z_0$ et un seul dans $D(0,1)$ tel que l'on ait $f(z_0)=z_0$.
\finenonce{006757}



\finexercice


\exercice{6758, queffelec, 2011/10/16}
\enonce{006758}{}
En utilisant le théorème de Rouché, on se propose de
donner une preuve du théorème d'inversion local ``holomorphe'' n'utilisant
pas le théorème d'inversion local dans $\Rr^2$. 

Soit donc $f$ une fonction holomorphe au voisinage d'un point $z_0$ de $\Cc$ telle que $f'(z_0)\ne0$. On suppose sans restreindre la généralité que
$z_0=0$, $f(z_0)=0$, et $f'(z_0)=1$.
\begin{enumerate}
\item Montrer qu'il existe un voisinage $V$ de $z_0$ et une constante $K>0$
tel que l'on ait, pour tout $z$ dans $V$, $|f(z) - z|\le K\,|z^2|$.
\item Montrer qu'il existe $r>0$ tel que, si l'on a $|\alpha|<r/2$,
l'équation $f(z) = \alpha$ a une solution unique dans le disque $D(0,r)$. 
\item Montrer enfin qu'il existe un ouvert $U$ de $\Cc$ tel que $f$
soit une bijection de $U$ sur le disque $D(0, r/2)$.
\item En déduire que $g = f^{-1}$ est continue sur le disque $D(0,r/2)$
et, de là, que $g$ est holomorphe dans ce disque.
\end{enumerate}
\finenonce{006758}



\finexercice


\exercice{6759, queffelec, 2011/10/16}
\enonce{006759}{}
\label{gijsexoav}
On considère la série $\displaystyle \sum_{n\in \Zz}
\frac{1}{(z-n)^2}$. 
\begin{enumerate}
\item Montrer que cette série converge normalement sur tout compact $K$
de $\Cc$.
En déduire que $\displaystyle f(z) = \sum_{n\in \Zz}
\frac{1}{(z-n)^2}$ est une fonction méromorphe sur $\Cc$. Vérifier que
l'on a $f(z+1) = f(z)$.
\item Déterminer le résidu de $f$ en chacun de ses p\^oles. Montrer
que, si l'on note $z=x+iy$, $f(z)$ tend vers 0, uniformément par rapport à
$x$, lorsque $|y|$ tend vers $\infty$.
\item Montrer que $\displaystyle f(z) = \left(\frac{\pi}{\sin(\pi z)}
\right)^2$. (On pourra utiliser le théorème de Liouville.) En déduire que
l'on a $\displaystyle \sum_{n=1}^\infty \frac1{n^2} = \frac{\pi^2}{6}$.
\end{enumerate}
\finenonce{006759}



\finexercice


\exercice{6760, queffelec, 2011/10/16}
\enonce{006760}{}
Soit $\gamma_n$ le chemin dont l'image $\gamma_n^*$ est le
rectangle de sommets $\pm(n+\frac12)\pm in$ parcouru une fois dans le sens
direct. Evaluer pour $a\notin \Zz$
$$\int_{\gamma_n} \frac{\pi \mathrm{cotan}(\pi z)}{(z+a)^2} \,dz\ .
$$
Montrer que l'on a
$$
\lim_{n\to \infty} \int_{\gamma_n} \frac{\pi \mathrm{cotan}(\pi z)}{(z+a)^2} \,dz = 0\ .
$$
(On pourra établir auparavant que, si $z$ appartient à $\gamma_n^*$ et si
$n$ est assez grand, on a $|\mathrm{cotan}(\pi z)|\le 2$.) En déduire (cf. exercice \ref{gijsexoav})
que $\displaystyle \sum_{n\in \Zz} \frac1{(a+n)^2} = \left(
\frac{\pi}{\sin(\pi a)} \right)^2$.
\finenonce{006760}



\finexercice


\exercice{6808, gijs, 2011/10/16}
\enonce{006808}{}
Soit $\alpha$ un réel tel que $-1< \alpha < 2$, et soit
$f: \Cc\setminus\{it\mid t\in \,]-\infty,0]\,\}$ définie
par $$f(z) =   \frac{e^{iz} e^{\alpha\log z}}{1+z^2}\ ,$$ 
où $\log z$ désigne la branche uniforme
du logarithme complexe qui est réelle pour $z$ réel
strictement positif, avec $-\pi/2 < \arg z < 3\pi/2$.

\begin{enumerate}
\item Montrer que, pour tout $\theta$ tel que $0\le
\theta \le \pi/2$, on a~: $0\le 2\theta/\pi \le \sin
\theta \le \theta$.

\item 
Soit $\gamma_\epsilon$ le demi-cercle de rayon
$\epsilon>0$, de centre 0, situé dans le demi-plan
$\Im\  z\ge 0$. Démontrer que
$ \lim_{\epsilon\downarrow 0} \int_{\gamma_\epsilon}
f(z)\,dz = 0$.

\item
Soit $\gamma_R$ le demi-cercle de rayon
$R>0$, de centre 0, situé dans le demi-plan $\Im 
z\ge 0$. Démontrer que $ \lim_{R\to \infty}
\int_{\gamma_R} f(z)\,dz = 0$.

\item
En intégrant $f$ sur le bord du domaine $\epsilon\le
 |z| \le R$, $0\le \arg(z) \le \pi$, déduire de ce
qui précède que l'intégrale $$I(\alpha) =
  \int_0^\infty \frac{x^\alpha
\cos(x-\frac{\alpha\pi}2)}{1+x^2}\,dx$$ est convergente et
en même temps calculer sa valeur.
\end{enumerate}
\finenonce{006808}



\finexercice


\exercice{6811, gijs, 2011/10/16}
\enonce{006811}{}
Soit $a$ un réel tel que $0 \le a < 1$.
\begin{enumerate}
\item
Démontrer que les intégrales $%  
I(a) =
\int_0^{+\infty} \frac{\sinh(ax)}{\sinh x} \,dx$ et
$%  
J(a) = \int_0^{+\infty} \frac{\cosh(ax)}{\cosh x}
\,dx$ sont convergentes ($\sinh$ et $\cosh$ désignent
les sinus et cosinus hyperboliques).

\item
Soit $\epsilon$ et $R$ des réels tels que $0<\epsilon <
\frac\pi2 < R$, soit $f(z)$ la fonction $f(z) =
\frac{e^{az}}{e^z - e^{-z}}$ et soit
$K_{\epsilon,R} \subset \Cc$ le compact obtenu en \^otant
la demi-boule ouverte de centre 0 et de rayon $\epsilon$ du
rectangle de sommets $R$, $R + i\frac\pi2$, $-R +
i\frac\pi2$, $-R$.
  \begin{enumerate}
  \item
Démontrer que ${ \lim_{R\to
+\infty}}\int_\gamma f(z)\,dz = 0$ lorsque $\gamma$ est
(i) le c\^oté du rectangle joignant $R$ à $R +
i\frac\pi2$ et (ii) le c\^oté du rectangle joignant
$-R + i\frac\pi2$ à $-R$.

  \item
Calculer le résidu de $f$ en 0.

  \item
Du calcul de l'intégrale de $f$ le long du bord 
orienté de $K_{\epsilon,R}$ et de sa limite quand
$\epsilon
\to 0$ et
$R
\to +\infty$, déduire $J(a)$ et $I(a)$.
  \end{enumerate}
\end{enumerate}
\finenonce{006811}



\finexercice


\exercice{7554, mourougane, 2021/08/10}
\enonce{007554}{Pour apprendre le cours}
\begin{enumerate}
\item Soit $I=[a,b]$ ($a<b$) un intervalle de $\Rr$ et $\phi : I\to \Cc$ une application continue.
Montrer que $$\left|\int_a^b \phi\right|\leq \int_a^b |\phi|.$$
On pourra considérer un nombre complexe $c$ tel que $\left|\int_a^b \phi\right|=c\int_a^b \phi$.

\item Vérifier que la relation définie sur les chemins paramétrés par $\gamma_1\equiv \gamma_2$
s'il existe une bijection $\alpha : J\to I$ dérivable à dérivée continue et partout strictement positive telle que $\gamma_2=\gamma_1\circ\alpha$
est une relation d'équivalence.
\item Donner l'exemple de $D$ un ouvert de $\Cc$, $f : D \to \Cc$ une application continue, $I$ un intervalle de $\Rr$
et $\gamma_1 : I\to D$ une application continue, $J$ un intervalle de $\Rr$
et $\alpha : J\to I$ une application continue bijective et $\gamma_2=\gamma_1\circ\alpha$
tels que 
$$\int_I f( \gamma_1(t)) dt\not =\int_J f(\gamma_2(\tau)) d\tau.$$
\item Donner l'exemple d'une fonction $f :D\to\Cc$ holomorphe et de deux chemins $\Gamma_1$ et $\Gamma_2$ de même origine et fin,
tels que $$\int_{\Gamma_1}f(z)dz\not=\int_{\Gamma_2}f(z)dz.$$
\item Soit $D$ un ouvert de $\Cc$, $f : D \to \Cc$ une application continue, $\Gamma$ un chemin dans $D$.
A-t-on $$re\left(\int_\Gamma f(z)dz\right)=\int_\Gamma re(f(z)) dz ?$$
\end{enumerate}
\finenonce{007554}
\finexercice


\exercice{7555, mourougane, 2021/08/10}
\enonce{007555}{Calcul d'intégrales}
On paramètre le cercle $C_r$ de centre $0$ et de rayon $r>0$ orienté dans le sens trigonométrique 
du plan complexe en définissant pour $t\in\Rr$, $\xi(t)$ 
comme le point d'intersection de la droite d'équation $y=t(x+r)$ avec le cercle $C_r$ différent de $-r$.
\begin{enumerate}
\item Montrer que $\xi(t)=r\frac{1+it}{1-it}$.
\item Vérifier que $\xi$ est dérivable sur $\Rr$ et calculer $\frac{\xi'(t)}{\xi(t)}$.
\item En déduire que $$\int_{C_r} \frac{dz}{z}=2i\int_{-\infty}^{+\infty}\frac{dt}{1+t^2}.$$
\item En déduire que $$\int_{-\infty}^{+\infty}\frac{dt}{1+t^2}=\pi.$$
\end{enumerate}
\finenonce{007555}
\finexercice


\exercice{7556, mourougane, 2021/08/10}
\enonce{007556}{Calculs}
\begin{enumerate}
 \item On considère le bord $\mathcal{C}$ du triangle de sommet : $ z=0 $, $ z=1$, et $ z=i$, orienté dans le sens direct.
Calculer les intégrales 
$ \int_{\mathcal{C}} xdz \quad \text{ et } \quad \int_{\mathcal{C}} ze^{z}dz.$

\item Soit le cercle unité $\mathcal{C}$ parcouru dans le sens direct. Pour tout entier relatif $n$, calculer l'intégrale 
$ \int_{\mathcal{C}} z^{n}dz.$
Expliquer le cas particulier où $n=-1$.
\item
Montrer $\left|\int_\gamma \frac{dz}{z^2+1}\right| \le \frac{\pi}{3}, \mbox{ lorsque }\gamma(t)=2e^{it}, \; 0\le t\le \frac{\pi}{2}.$
\item Montrer $\left|\int_\gamma \frac{e^z}{z} dz\right| \le \pi e, \mbox{ lorsque }\gamma(t)=e^{it}, \; 0\le t\le \pi.$
\end{enumerate}
\finenonce{007556}
\finexercice


\exercice{7557, mourougane, 2021/08/10}
\enonce{007557}{Intégrale elliptique}
Soient deux réels strictement positifs $a$ et $b$, on pose $ \gamma = a \cos (t) +ib \sin (t) $ pour $ t \in [ 0, 2 \pi ] $. Calculer l'intégrale
$
\int_{\gamma} \frac{dz}{z} 
$
et en déduire : 
$
\int_{0}^{2\pi} \frac{ dt }{a^{2} \cos ^{2}(t) + b^{2} \sin ^{2}(t)}.
$
\finenonce{007557}
\finexercice


\exercice{7558, mourougane, 2021/08/10}
\enonce{007558}{Existence de primitive}
Soit $D:=\Cc-[0,1]$ et $f : D\to \Cc$, $z\mapsto\frac{1}{z(z-1)}$.
\begin{enumerate}
 \item $D$ est-il un ouvert étoilé ?
 \item La fonction $F : z\mapsto\log (1-\frac{1}{z})$ est-elle bien définie et continue sur $D$ ?
 \item Montrer que pour tout chemin orienté $\Gamma$ de $D$,
 $\int_\Gamma f(\zeta)d\zeta=0.$
\end{enumerate}
\finenonce{007558}
\finexercice


\exercice{7559, mourougane, 2021/08/10}
\enonce{007559}{Lemme de Goursat}
Soit $D$ un ouvert de $\Cc$ et $f : D\to\Cc$ une application holomorphe.
Soit $T$ un triangle (plein) inclus dans le domaine $D$. On notera $p$ son périmètre.
Le but de l'exercice est de montrer que 
$$\int_{\partial T}f(z)dz=0.$$
\begin{enumerate}
\item En considérant les quatre triangles obtenus en traçant les segments entre deux milieux de côtés de $T$,
montrer que pour l'un de ces triangles noté $T_1$,
$$|\int_{\partial T}f(z)dz|\leq 4|\int_{\partial T_1}f(z)dz|.$$
\item En itérant cette construction, montrer que pour tout $n$, il existe un triangle $T_n\subset T_{n-1}$ de périmètre $\frac{p}{2^n}$
tel que $$|\int_{\partial T}f(z)dz|\leq 4^n|\int_{\partial T_1}f(z)dz|.$$
\item Montrer que l'intersection des $T_n$ est un point $c$ de $D$.
\item Montrer qu'un existe une fonction continue $h$ sur $D$ nulle en $c$ telle que pour tout $z\in D$,
$$f(z)=f(c)+(z-c)f'(c)+(z-c)h(z).$$
\item Calculer $\int_{\partial T_n}f(c)dz$ et $\int_{\partial T_n}(z-c)f'(c)dz$.
\item Montrer que pour tout $\epsilon >0$ il existe $N$ tel que pour $n\geq N$,$$|\int_{\partial T_n}(z-c)h(z)dz|\leq (\frac{p}{2^n})^2\varepsilon.$$
\item Conclure.
\item Montrer en choisissant un découpage de $T$ avec un petit triangle autour de $a$ que 
 si $f : D\to\Cc$ est une application continue et holomorphe sur $D-\{a\}$, alors pour tout triangle de sommet $a$ dans $D$,
 $\int_{\partial T}f(z)dz=0$.
\end{enumerate}
\finenonce{007559}
\finexercice


\exercice{7560, mourougane, 2021/08/10}
\enonce{007560}{Calcul d'intégrales par chemin fermé}
Soit $a$ un nombre réel tel que $|a|\leq 1$. 
On considère la fonction $f :\Cc\to\Cc$, $z\mapsto e^{-z^2}$ et pour tout réel strictement positif $r$,
les chemins suivants dans le plan complexe
$$
\setlength{\unitlength}{.4in}
\begin{picture}(7,3.5)(0,-.25)
\linethickness{1pt}
\put(0,0){\line(1,0){4}}
\put(4,0){\line(0,1){3}}
\put(0,0){\line(4,3){4}}
%\put(2,0){\makebox(0,0){$\to$}}
\put(2,-.25){\makebox(0,0){$\gamma_1\to$}}
\put(4.5,1.5){\makebox(0,0){$\uparrow\gamma_2$}}
\put(2,2){\makebox(0,0){$\gamma_3\nearrow$}}
\put(-.25,-.25){\makebox(0,0){$0$}}
\put(4.25,-.25){\makebox(0,0){$r$}}
\put(4.25,3.25){\makebox(0,0){$r(1+ai)$}}
\end{picture}
$$
\begin{enumerate}
 \item Calculer $\int_{\gamma_1}f(z)dz$.
 \item Montrer que $|\int_{\gamma_2}f(z)dz|\leq \frac{1}{r}.$
 \item En déduire $$\int_0^{+\infty}e^{-(1+ai)^2t^2}dt=\frac{\sqrt{\pi}}{2}\frac{1}{1+ia}.$$
 \item En déduire les valeurs des intégrales $\int_0^{+\infty}\cos(t^2)dt$ et $\int_0^{+\infty}\sin(t^2)dt$.
\end{enumerate}
\finenonce{007560}
\finexercice


\exercice{7561, mourougane, 2021/08/10}
\enonce{007561}{Calcul d'intégrales par la formule intégrale de Cauchy}
Calculer $$\int_{\partial \Delta_2} \frac{e^z}{(z-3)^2} dz\ ,\ \int_{\partial \Delta_2} \frac{e^z}{(z+1)} dz\quad \text{ et } \quad \int_{\partial \Delta_2} \frac{e^z}{(z+1)(z-3)^2} dz.$$
\finenonce{007561}
\finexercice


\exercice{7562, mourougane, 2021/08/10}
\enonce{007562}{Calcul d'intégrales et théorème de Liouville}
\begin{enumerate}
 \item Soit $r>0$ et $D$ un voisinage ouvert de $\overline{\Delta_r}$.
 Soit $f : D\to\Cc$ une fonction holomorphe. Soit $a$ et $b$ deux points distincts dans $\Delta_r$.
 Calculer $$\int_{\partial\Delta_r}\frac{f(\zeta)}{(\zeta-a)(\zeta-b)}d\zeta.$$
 \item En déduire le théorème de Liouville : Toute fonction holomorphe bornée sur $\Cc$ est constante.
\end{enumerate}
\finenonce{007562}
\finexercice


\exercice{7563, mourougane, 2021/08/10}
\enonce{007563}{}
\begin{enumerate}
  \item Déterminer toutes les fonctions holomorphes définies sur le plan complexe tout entier vérifiant : $|f(z)| \geq 1 $ . 
 \item On considère une fonction $f$ holomorphe dans le disque unité, vérifiant $|f(z)| \leq 1 $, que peut on dire de $|f'(0)| $ ? 
 \item Soit $n_0$ un nombre entier naturel et $f$ une fonction holomorphe définie sur le plan complexe tout entier vérifiant $|f(z)|\leq |z|^{n_{0}}$. 
Montrer que $f$ est un polynôme de degré au plus $n_{0}$.
 \item Soient $\Omega$ un ouvert connexe et $D$ une droite, $f$ une fonction continue sur $\Omega$, holomorphe sur la restriction $\Omega - \{D\} $. 
Montrer que $f$ se prolonge en une fonction holomorphe sur $\Omega$ tout entier. 
\end{enumerate}
\finenonce{007563}
\finexercice


\exercice{7577, mourougane, 2021/08/10}
\enonce{007577}{Propriété}
Soient deux fonctions $g$ et $h$ holomorphes sur un ouvert connexe du plan complexe, $a$ un pôle de $g/h$ tel que $h(a)=0$ et $h'(a)$ soit non nul. Montrer que :
$$Res(a, \frac{g}{h} ) = \frac{g(a)}{h'(a)}$$
\finenonce{007577}
\finexercice


\exercice{7578, mourougane, 2021/08/10}
\enonce{007578}{Calculs de résidus}
\begin{enumerate} 
 \item L'application $\cot(\pi z):=\frac{\cos (\pi z)}{\sin (\pi z)}$ est-elle méromorphe ?
Si oui, déterminer ses résidus.
\item Calculer les résidus en tous les pôles de 
$$f(z)=\frac{z^2}{z^4+1} \quad \text{ et } \quad g(z)=\frac{z^2+z+5}{z(z^2+1)^2}.$$
\end{enumerate}
\finenonce{007578}
\finexercice


\exercice{7579, mourougane, 2021/08/10}
\enonce{007579}{Calcul d'intégrales}
\begin{enumerate} 
\item Montrer que $$\int_{\partial\Delta_2}\frac{dz}{\sin^2z\cos z}=0.$$
\item Calculer les intégrales suivantes, où les chemins fermés simples $\gamma$ sont parcourues dans le sens direct.
\begin{enumerate}
\item $\displaystyle \int \limits_{\gamma} \frac{z}{(z - 1)(z - 2)} dz $ où $\gamma=\partial\Delta_{1/2}(2)$ est le cercle centré en 2 de rayon 1/2 ; 
\item $\displaystyle \int \limits_{\gamma} \frac{\exp(z)}{z^{2}(z - 9)^{2}} dz $ où $\gamma=\partial\Delta$ est le cercle centré en l'origine de rayon 1. 
\end{enumerate}
\end{enumerate}
\finenonce{007579}
\finexercice


\exercice{7580, mourougane, 2021/08/10}
\enonce{007580}{Somme de résidus}
Soit $f(z)$ une fonction rationnelle, quotient d'un polynôme $P(z)$ par un polynôme $Q(z)$.
On suppose que $\deg P +2 < \deg Q$. 
Montrer que 
$$\sum_{c\in\Cc}Res_c(f)=0.$$
\finenonce{007580}
\finexercice


\exercice{7581, mourougane, 2021/08/10}
\enonce{007581}{Calcul d'intégrales trigonométriques}
\begin{enumerate} 
 \item Soit $a$ un nombre complexe de module différent de $1$.
 Montrer que $$i\int_0^{2\pi}\frac{d\theta}{1-2a\cos\theta + a^2}=\int_{\partial\Delta}\frac{dz}{(z-a)(1-az)}.$$
 \item Déterminer les résidus des pôles de $\frac{1}{(z-a)(1-az)}$ contenus dans le disque unité ouvert.
 \item En déduire la valeur de $\int_0^{2\pi}\frac{d\theta}{1-2a\cos\theta + a^2}$.
\end{enumerate}
\finenonce{007581}
\finexercice


\exercice{7582, mourougane, 2021/08/10}
\enonce{007582}{Calcul d'intégrales rationnelles}
On considère la fonction méromorphe $f(z)=\frac{z^2}{1+z^4}$ sur $\Cc$.
\begin{enumerate} 
 \item Montrer que $zf(z)$ a une limite quand $|z|$ tend vers $+\infty$.
 \item Montrer que $\int_{-\infty}^{+\infty}\frac{x^2}{1+x^4}dx$ converge.
 \item Déterminer les pôles de $f$ contenus dans le demi-plan $\mathbb{H}:=\{z\in\Cc, Im(z)>0\}$ et leurs résidus.
 \item En intégrant sur le chemin $\Gamma$ défini par
 $$\begin{tikzpicture}
\draw[help lines, step=1, very thin] (-6,0) grid (6,4);
\draw[shift={(3.4,0)}, thick] (6:0) arc (0:180:3.4); \draw (0,3.4) node[above]{$\leftarrow$} ;
\draw[thick] (-3.4,0) -- (3.4,0) ;
\draw (0,0) node[below]{$0$} ;
\draw (3.5,0) node[below]{$r$} ;\draw (-3.5,0) node[below]{$-r$} ;\draw (-1.5,0) node[below]{$\to$} ;
\draw (1,0) node[below]{$1$} ;
\end{tikzpicture}$$
montrer que 
$$\int_{-r}^rf(x)dx+\int_{\partial(\Delta_r\cap\mathbb{H})} f(z)dz=2i\pi\sum_{c\in\mathbb{H}}Res_c(f).$$
\item Montrer que $$\int_{-\infty}^{+\infty}\frac{x^2}{1+x^4}dx=\frac{\pi}{\sqrt{2}}.$$
\end{enumerate}
\finenonce{007582}
\finexercice


\exercice{7583, mourougane, 2021/08/10}
\enonce{007583}{Contraintes}
\begin{enumerate}
 \item Montrer que le polynôme $f(z)=3+7z+2z^4$ a, comme le polynôme $3+7z$, exactement un zéro dans $\Delta$.
 \item Déterminer le nombre de zéros (comptés avec multiplicité) de $z^5+\frac{z^3}{3}+\frac{z^2}{4}+\frac{1}{3}$ dans $\Delta$.
\item Déterminer le nombre de zéros (comptés avec multiplicité) de $z^5+\frac{z^3}{3}+\frac{z^2}{4}+\frac{1}{3}$ dans $\Delta_{\frac{1}{2}}$.
\item Déterminer le nombre de zéros du polynôme $ z^{4} -5z-1 $ dans la couronne $ 1<\arrowvert z \arrowvert < 2 $. 
\end{enumerate}
\finenonce{007583}
\finexercice

\section{445.00 Tranformée de Laplace et de Fourier}


\exercice{6602, hueb, 2011/10/16}
\enonce{006602}{} 
Calculer les
transformées de Laplace des fonctions suivantes, où $a,b,A,T$ sont
des nombres réels positifs.
\begin{enumerate}
\item $f(t)= \begin{cases}
    \frac Aa t,&0\leq t \leq a\\
    A,&t \geq a \end{cases}$
\item $f(t)= \begin{cases} nA,&\text{pour}\ (n-1) T \leq t < nT,\ 
    \text{où}\ n\in \Nn, n\geq 1, \ \text{c.a.d}
    \\
    A, &0\leq t <T\\
    2A, &T\leq t <2T, \ \text{etc.}  \end{cases}$
\item $f(t)= \begin{cases}
    \frac {4A}T t,&0\leq t < \frac T4\\
    -\frac {4A}T t + 2A,&\frac T4 \leq t < \frac 34 T\\
    \frac {4A}T t - 4A,&\frac 34 T \leq t < T\\
    0,& t \geq T \end{cases}$
\item $f(t)= \begin{cases}
    0,&0\leq t <a\\
    A,&a\leq t < a+b\\
    -A,&a+b\leq t < a+2b\\
    0,& t \geq a+2b \end{cases}$
\end{enumerate}
\finenonce{006602}



\finexercice


\exercice{6603, hueb, 2011/10/16}
\enonce{006603}{} 
La transformée de
Laplace de la \lq\lq fonction\rq\rq\ impulsion de Dirac: Pour
$\varepsilon >0$, on pose
$$
F_\varepsilon (t) = \begin{cases}
  \frac 1 \varepsilon,\quad & 0 \leq t \leq \varepsilon\\
  0,\quad & t > \varepsilon \phantom{ \leq 0}
                    \end{cases}.
                    $$
\begin{enumerate}
\item Trouver $\mathcal{L}(F_\varepsilon)$.
\item Montrer que $\lim_{\varepsilon \to 0}\mathcal{L}(F_\varepsilon)
  = 1$.
\end{enumerate}
\finenonce{006603}



\finexercice


\exercice{6604, hueb, 2011/10/16}
\enonce{006604}{} 
Calculer la transformée
de Laplace de la fonction $F$ où $a,b, \omega, k \in \Rr$, $a,b >
0$:
\begin{enumerate}
\item $\quad F(t) = a \sin \omega t,$
\item $\quad F(t) = a (1-e^{- bt}),$
\item $\quad F(t) = a \cos (bt-k).$
  N.B. Ici la formule qui exprime la transformée de Laplace de la fonction \\
  $G(t) = \begin{cases}
    F(t-a),&t \geq a,\\
    0,\quad &t < a,
                \end{cases}$
                en fonction de celle de $F$ ne s'applique pas à (3).
                Pourquoi pas?
\end{enumerate}
\finenonce{006604}



\finexercice


\exercice{6605, hueb, 2011/10/16}
\enonce{006605}{} 
Vérifier les
propriétés suivantes de la transformation de Laplace $\mathcal{L}$
:
\begin{enumerate}
\item $\mathcal{L}(c_1 F_1 + c_2 F_2) =c_1 \mathcal{L}(F_1) + c_2
  \mathcal{L}(F_2)$.
\item $\mathcal{L}(e^{at}F(t))(s) =(\mathcal{L} F)(s-a)$.
\item Pour $G(t) = \begin{cases} F(t-a),\quad &t \geq a,\\
    0, \quad &t \leq a,
                \end{cases}$
                on a $(\mathcal{L}G)(s) =e^{-as}(\mathcal{L}F)(s)\ (a
                \geq 0)$.
              \item Pour $F_a(t) = F(at)$ on a $(\mathcal{L}F_a)(s) =
                \frac 1a (\mathcal{L}F)(\frac sa)$.
              \item Pour $G(t) = \int_0^t F(u) du$ on a
                $(\mathcal{L}G) (s) = \frac {(\mathcal{L} F)(s)} s$.
              \item $\mathcal{L}(t^n F(t)) = (-1)^n (\mathcal{L}
                F)^{(n)}$
              \item Si $\lim_{t\to 0} \frac {F(t)}t$ existe,
                $\mathcal{L}(\frac {F(t)}t)(s) = \int_s^{\infty}
                (\mathcal{L} F)(\zeta) d \zeta$.\ 
              \item Si $F$ est périodique, $F(t+T) = F(t)$, alors
                $$(\mathcal{L}F)(s) = \frac 1{1-e^{-sT}}\int_0^T F(t)
                e^{-st} dt$$
              \item $\lim_{s \to +\infty} (\mathcal{L} F)(s) = 0$.
              \item $\lim_{s \to +\infty} (\mathcal{L}F)(s) = 0$.
              \item $\lim_{t \to 0, t > 0} F(t) = \lim_{s \to \infty}
                s (\mathcal{L}F)(s)$ si les limites indiquées
                existent (Théorème de la valeur initiale).
              \item $\lim_{t \to \infty} F(t) = \lim_{s \to 0} s
                (\mathcal{L} F)(s)$ si les limites indiquées
                existent (Théorème de la valeur finale).
\end{enumerate}
\finenonce{006605}



\finexercice


\exercice{6606, hueb, 2011/10/16}
\enonce{006606}{} Trouver toutes les
solutions $Y$ de l'équation différentielle
$$tY'' + 2 Y' + tY =0,\quad Y(0) =1,$$
définies sur la droite
réelle entière (i.~e. $D_Y= \Rr$).  \finenonce{006606}



\finexercice


\exercice{6607, hueb, 2011/10/16}
\enonce{006607}{} Etudier la convergence des
séries
$$\sum_1^{\infty} \frac{\cos nx} {n^2}, \quad \sum_1^{\infty}
\frac{\sin nx} {n^2}, \quad \sum_1^{\infty} \frac{e^{inx}} {n!}. $$
\finenonce{006607}



\finexercice


\exercice{6608, hueb, 2011/10/16}
\enonce{006608}{} Calculer les coefficients
de Fourier de la fonction périodique $f$, définie sur $[-\pi,\pi]$
par
$$f(x) = \pi - |x|,\quad |x| \leq \pi.$$
Etudier la convergence de la
série de Fourier qui en résulte; est-elle absolue ou peut-être
uniforme?  En déduire la valeur de $\sum_1^{\infty} \frac 1{n^2}$.
\finenonce{006608}



\finexercice


\exercice{6609, hueb, 2011/10/16}
\enonce{006609}{}
\begin{enumerate}
\item Trouver les coefficients de Fourier en $\sin$ et $\cos$ de la
  fonction périodique $F$, donnée sur $]-5,5[\setminus \{0\}$ par
  $$F(x)= \begin{cases}
    0,\quad &-5<x<0\\
    3,\quad & \phantom{-}0<x<5
                    \end{cases}.$$
                  \item Vérifier que $F$ satisfait aux conditions de
                    Dirichlet.  Comment doit $F$ être définie en
                    $x=-5,0,5$ pour que sa série de Fourier converge
                    vers $F(x)$ pour tout $x \in [-5, 5]$?
\end{enumerate}
\finenonce{006609}



\finexercice


\exercice{6610, hueb, 2011/10/16}
\enonce{006610}{} Développer la fonction
périodique $F$, donnée sur $]-2,2[$ par $F(x) = x$ (fonction en
dents de scie), en série trigonométrique.  \finenonce{006610}



\finexercice


\exercice{6611, hueb, 2011/10/16}
\enonce{006611}{}
\label{exo43}
\begin{enumerate}
\item Calculer la transformée de Fourier de la fonction porte
  symétrique
  $$f(x)= \begin{cases}
    A,\quad & -\frac \tau 2<x<\frac \tau 2\\
    0,\quad & |x| > \frac \tau 2
\end{cases}.$$
\item Vérifier que $f$ satisfait aux conditions de Dirichlet.
  Comment doit $f$ être définie en $x=\pm \frac \tau 2$ pour que
  l'intégrale de Fourier converge vers $f(x)$ pour tout $x$?
\end{enumerate}
\finenonce{006611}



\finexercice


\exercice{6612, hueb, 2011/10/16}
\enonce{006612}{}
\begin{enumerate}
\item Utiliser les résultats de l'exercice \ref{exo43} pour
  évaluer
  $$\int_{-\infty}^{\infty} \frac{\sin a \lambda \cos b
    \lambda}{\lambda} d\lambda.$$
\item En déduire la valeur de $\int_{0}^{\infty} \frac{\sin u} u
  du$.
\end{enumerate}
\finenonce{006612}



\finexercice


\exercice{6613, hueb, 2011/10/16}
\enonce{006613}{}
\begin{enumerate}
\item Trouver la transformée de Fourier en cosinus de $f(x) =
  e^{-m|x|},\ m>0$.
\item Utiliser le résultat de (1) pour montrer que, pour $p>0,\ 
  \beta >0$,
  $$\int_{0}^{\infty} \frac{\cos pv} {v^2 + \beta^2} dv = \frac \pi {2
    \beta} e^{-p \beta}.$$
\end{enumerate}
\finenonce{006613}



\finexercice


\exercice{6614, hueb, 2011/10/16}
\enonce{006614}{} Trouver une fonction $f$
de sorte que l'équation intégrale suivante soit vérifiée:
$$\int_{0}^{\infty} f(x) \sin \alpha x dx =
\begin{cases} 1-\alpha,\quad &0 \leq \alpha \leq 1,\\
  0,\quad &\alpha > 1 . \end{cases} $$
\finenonce{006614}



\finexercice


\exercice{6615, hueb, 2011/10/16}
\enonce{006615}{} Montrer que, pour $x \geq
0$,
$$\int_{0}^{\infty} \frac{\cos \lambda x} {\lambda^2 + 1} d\lambda =
\frac \pi 2 e^{-x}.$$
\finenonce{006615}



\finexercice


\exercice{6616, hueb, 2011/10/16}
\enonce{006616}{} Evaluer
\begin{enumerate}
\item $\int_{0}^{\infty} \frac {dx} {(x^2 +1)^2},$
\item $\int_{0}^{\infty} \frac {x^2 dx} {(x^2 +1)^2},$ en appliquant
  l'identité de Parseval.
\end{enumerate}
\finenonce{006616}



\finexercice


\exercice{6617, hueb, 2011/10/16}
\enonce{006617}{} Calculer la transformée
inverse de Laplace $\mathcal{L}^{-1}(f)$ pour :
\begin{enumerate}
\item $f(s) = \frac s{(s+1)^3(s-1)^2}$
\item $f(s) = \frac 1{(s^2+1)^2}$
\item $f(s) = \frac {s+1}{(s+2)(s+3)}$
\item $f(s) = \frac {s+1}{(s+2)(s+3)^2(s^2+4s+5)}$
\item $f(s) = \frac {1}{(s+3)(s+4)}$
\item $f(s) = \frac {s+2}{(s+1)^2(s+3)}$
\item $f(s) = \frac {s+1}{s^2 +4s+16}$
\end{enumerate}
\finenonce{006617}



\finexercice


\exercice{6618, hueb, 2011/10/16}
\enonce{006618}{} Soient $a,x$ des nombres
réels, $0<x<a$, et posons $f(s) = \frac
{\textrm{sh}sx}{s^2\textrm{ch} sa}$. Déterminer
$\mathcal{L}^{-1}(f)$.  \finenonce{006618}



\finexercice


\exercice{6619, hueb, 2011/10/16}
\enonce{006619}{} Résoudre l'équation
différentielle partielle :
$$\frac{\partial ^2 Y}{\partial t^2} = c^2\frac{\partial ^2
  Y}{\partial x^2}, \quad 0 \leq x < \ell, \ t \geq 0,$$
avec les
conditions initiales suivantes :
$$Y(x,0) = 0,\ Y_t(x,0) = 0,\ Y(0,t) = 0, \ Y_x(\ell,t) = \frac
{F_0}{E}\ \text{(i.e. constant)}$$
\finenonce{006619}



\finexercice

\section{446.00 Autre}


\exercice{2870, burnol, 2009/12/15}
\enonce{002870}{}
Le but de ce problème est d'établir les deux formules importantes:
\[ \forall z\in\Cc\setminus\Zz\quad \frac\pi{\sin \pi z} = \frac1{z} +
\lim_{\substack{N\to+\infty \\ M\to+\infty}} 
\sum_{\substack{-M\leq n \leq N \\ n\neq0}}
\frac{(-1)^n}{z-n} = \frac1z + \sum_{n=1}^\infty (-1)^n\frac{2z}{z^2-n^2}\] 
\[\forall z\in\Cc\qquad \sin(\pi z) = \lim_{N\to\infty} \pi z\prod_{n=1}^N (1 -
\frac{z^2}{n^2})\]
\begin{enumerate}
  \item Montrer la convergence de la série
$\sum_{n=1}^\infty \frac{(-1)^n}{z-n}$ (regarder les sommes
partielles pour les indices pairs).
  \item On pose $f(w) =  \frac\pi{\sin \pi
w}$. Soit $z\notin\Zz$ fixé, soit $N>|z|-\frac12$ et ${\cal R}_N$ le carré
$\{|x|\leq N+\frac12$, $|y|\leq N+\frac12\}$, et $C_N =
\partial {\cal R}_N$ son
bord parcouru dans le sens direct. Exprimer $\frac1{2\pi i}
\int_{C_N} \frac{f(w)}{w-z}dw$ à l'aide du Théorème
des résidus.
  \item Montrer $\int_{C_N} \frac{f(w)}w dw = 0$ (on notera que $f$ est
impaire) et en déduire :
\[ \frac1{2\pi i} \int_{C_N} \frac{f(w)}{w-z}dw =
\frac1{2\pi i} \int_{C_N} \frac\pi{\sin \pi 
w}\frac{z}{w(w-z)}dw\]
  \item On rappelle l'identité $\sin(w) = \sin(x)\ch(y) + i
\cos(x)\sh(y)$ pour $w=x+iy$. Montrer $|\sin w|^2 = \sin^2
x + \sh^2 y$ ($x,y\in\Rr\dots$). En déduire $|\sin(\pi w)| = \ch(\pi
y)\geq 1$ sur les bords verticaux du carré et $|\sin(\pi w)|
\geq \sh(\pi(N+\frac12))\geq \sh(\pi\frac12) =
2.301\dots\geq1$ sur les bords horizontaux. 
Conclure la preuve de
\[ \frac\pi{\sin \pi z} = \frac1z + \sum_{n=1}^\infty
(-1)^n\frac{2z}{z^2-n^2}\]  
avec convergence uniforme pour $|z|$ borné.
  \item Reprendre la même technique et prouver:
\[ \forall z\in\Cc\setminus\Zz\quad\frac{\pi \cos(\pi
z)}{\sin(\pi z)} = \lim_{N\to\infty} 
\sum_{-N\leq n\leq N} \frac1{z-n} = \frac1z +
\sum_{n=1}^\infty \frac{2z}{z^2-n^2}\;,\]
avec convergence uniforme pour $|z|$ borné.
  \item On veut maintenant prouver: $\sin(\pi z) = \lim_{N\to\infty}
\pi z\prod_{n=1}^N (1 - 
\frac{z^2}{n^2})$
%
On fixe une fois pour toutes $R>0$, et on va
montrer la formule pour $|z|<R$. 
Soit $N$ avec $N>R$ et notons $f_N(z) =
\frac{\pi z}{\sin(\pi z)}\prod_{n=1}^N (1 - \frac{z^2}{n^2})$, prolongé
par continuité en les $n$, $|n|\leq N$. Montrer que $f_N$
est holomorphe et ne s'annule pas sur $D(0,R)$.
  \item Soit $\gamma:[0,1]\to\Cc^*$ le chemin $\gamma(t) =
f_N(tz)$. On a donc $\gamma(0) = 1$, $\gamma(1) = f_N(z)$,
et $\gamma(t)\neq0$ pour tout $t$. Par un théorème démontré
en cours (lequel?)
on a 
$\gamma(1) = \gamma(0)\exp\left(\int_\gamma \frac{dw}w\right) $.
En déduire $f_N(z) = \exp\left(\int_0^1 \frac{f_N'(tz)}{f_N(tz)}
zdt\right)$.
  \item Soit $\epsilon>0$. En utilisant la convergence uniforme pour
 $|z|$ borné
 du développement en fractions de $\pi\cotg(\pi z)$,
montrer que pour $N$ suffisamment grand, on a $|f_N'(w)|\leq
 \epsilon |f_N(w)|$ pour tout $w\in D(0,R)$, puis en déduire
\[ N\gg0\quad |z|<R \implies\quad |f_N(z)| \leq e^{\epsilon|z|} \leq
 e^{\epsilon R} \] 
  \item En déduire $\lim_{N\to\infty} f_N(z) = 1$,
uniformément sur $D(0,R)$. Conclure la
preuve du produit infini de Euler pour $\sin(z)$. 
\end{enumerate}
\finenonce{002870}
\finexercice


\exercice{2871, burnol, 2009/12/15}
\enonce{002871}{Produit absolument convergent}
Soit $u_n$, $n\geq1$ des nombres complexes. Montrer:
$1 + \sum_{n=1}^N |u_n| \leq \prod_{n=1}^N (1+|u_n|) \leq
e^{\sum_{n=1}^N |u_n| }$. En déduire que la suite
croissante $\prod_{n=1}^N (1+|u_n|)$ a une limite finie si
et seulement si $\sum_{n=1}^\infty |u_n|
< \infty$. On suppose maintenant être dans ce cas, et de
plus $\forall n\; 1+u_n\neq0$. Montrer
alors $\sum_{n=1}^\infty |\mathrm{Log}(1+u_n)| < \infty$, et en déduire que 
$\prod_{n=1}^\infty (1+ u_n)$ converge. On conviendra donc
de dire 
 que  $\prod_{n=1}^\infty (1+ u_n)$ est
``absolument convergent'' si $\sum_{n=1}^\infty |u_n|<\infty$,
et on vient donc de prouver qu'un produit 
absolument convergent est
convergent. C'est principalement, la seule chose que vous
ayez à savoir sur ce sujet.
\finenonce{002871}
\finexercice


\exercice{2872, burnol, 2009/12/15}
\enonce{002872}{}
  Pour quelles valeurs de $p$ (réel) $\prod_{k=1}^\infty (1 +
k^{-p})$ converge?
\finenonce{002872}
\finexercice


\exercice{2873, burnol, 2009/12/15}
\enonce{002873}{}
  \'Etant admis que 
$\sin(\pi z) = \pi z \prod_{k=1}^\infty \left(1 -
  \frac{z^2}{k^2}\right)$, 
prouver:
\[ \sin(\pi z) = \pi z \prod_{\substack{k=-\infty \\ k\neq0}}^{+\infty} \left(1 -
  \frac{z}{k}\right)e^{\frac zk} \]
et justifier la convergence absolue du produit.
\finenonce{002873}
\finexercice 


\exercice{2874, burnol, 2009/12/15}
\enonce{002874}{}
\'Etant admis  
$\sin(\pi z) = \pi z \prod_{k=1}^\infty \left(1 -
  \frac{z^2}{k^2}\right)$, prouver:
\[ \sin(\pi z) = \pi z \lim_{N\to\infty} \prod_{k=-N, k\neq0}^{+N}
  \frac{z - k}{-k}\;,\]
puis établir pour tout $\alpha\notin\Zz$:
\[ \sin(\pi (z-\alpha)) =
  -\sin(\pi\alpha)\lim_{N\to\infty}\prod_{k=-N}^{+N} \left(1
  -\frac{z}{\alpha+k}\right)\]
Montrer que le résultat reste valable si l'on remplace dans
le produit $-N$ par $-N\pm1$ ou $+N$ par $+N\pm1$. En
déduire:
\[ \cos(\pi z) = \prod_{k=0}^\infty \left(1 -\frac{z^2}{(\frac12 +
  k)^2}\right)\]
avec un produit absolument convergent.   
\finenonce{002874}
\finexercice


\exercice{2875, burnol, 2009/12/15}
\enonce{002875}{}
On rappelle la formule $\pi\cotg(\pi \alpha) =
\lim_{N\to\infty}
\sum_{k=-N}^{+N}\frac1{\alpha-k}$, pour $\alpha\in\Cc\setminus\Zz$. Montrer:
\[ \frac{\sin(\pi (\alpha - z))}{ \sin(\pi\alpha)}=
  e^{-\pi\cotg(\pi\alpha)z}\prod_{k=-\infty}^{+\infty}
  \left(1 -\frac{z}{\alpha+k}\right)e^{\frac z{\alpha+k}}\]
avec un produit absolument convergent.  
\finenonce{002875}
\finexercice


\exercice{2876, burnol, 2009/12/15}
\enonce{002876}{}
 \'Etablir la convergence et évaluer les produits infinis
suivants:
\begin{alignat*}{2}
&\prod_{n=1}^\infty \left( 1 + \frac1{n(n+2)}\right)&\qquad
&\prod_{n=2}^\infty \left( 1 - \frac2{n(n+1)}\right)\\
&\prod_{n=2}^\infty\frac{n^3 - 1}{n^3 + 1}&\qquad
&\prod_{n=1}^\infty \frac{n^2+1}{n^2}
\end{alignat*}
Les trois premiers s'obtiennent par des réarrangements
simples. Pour le dernier, utiliser le produit infini de
$\sin z$. 
\finenonce{002876}
\finexercice


\exercice{2877, burnol, 2009/12/15}
\enonce{002877}{}
On suppose $\sum_{n\geq1} |u_n|^2 < \infty$. Montrer que les
deux séries $\sum u_n$ et $\sum \mathrm{Log}(1+u_n)$ sont soit
toutes deux convergentes soit toutes deux divergentes (on
suppose $\forall n\ u_n\neq -1$). Donc si $\sum_{n\geq1} |u_n|^2 <
\infty$ le produit infini $\prod_{n=1}^\infty (1+u_n)$ est
convergent si et seulement si la série $\sum_{n=1}^\infty
u_n$ converge.
\finenonce{002877}
\finexercice


\exercice{2878, burnol, 2009/12/15}
\enonce{002878}{}
  Montrer que $\prod_{k=1}^\infty (1 + \frac ik)$ diverge
tandis que $\prod_{k=1}^\infty \left|1 + \frac ik\right|$
converge.
\finenonce{002878}
\finexercice

\exercice{2883, burnol, 2009/12/15}
\enonce{002883}{}
  Montrer que les racines du polynôme $P(z) = z^{111}+ 3
z^{50} + 1$ vérifiant $|z|<1$ sont simples et qu'il y en a
exactement $50$. 
\emph{Indication :} utiliser le théorème de Rouché en
écrivant $P(z) = 3 z^{50} + (z^{111}+1)$ et calculer $P'$ pour
s'assurer que les racines avec $|z|<1$ sont simples.
\finenonce{002883}
\finexercice


\exercice{2884, burnol, 2009/12/15}
\enonce{002884}{}
  Déterminer l'image par $z\mapsto \frac{3z+5}{z+2}$ du cercle
unité, du cercle de rayon $2$ centré en $1$, du cercle de
rayon $2$ centré en l'origine; de la droite imaginaire, de
la droite d'équation $x=y$, de la droite verticale passant
en $3$, de la droite verticale passant en $-2$.
\finenonce{002884}
\finexercice


\exercice{2885, burnol, 2009/12/15}
\enonce{002885}{}
\label{ex:burnol2.6}
  Question de cours: quels sont les automorphismes
  de $D(0,1)$ avec $0$ comme point fixe? 
\finenonce{002885}
\finexercice


\exercice{2886, burnol, 2009/12/15}
\enonce{002886}{} Soit $\alpha$ avec $|\alpha|<1$. On sait que
$z\mapsto \phi_\alpha(z) = \frac{\alpha - z}{1 -
\overline\alpha z}$ est un automorphisme du disque unité
$D(0,1)$. Trouver $z_1$ et $z_2$ avec $\phi_\alpha(z_1) =
z_2$, $\phi_\alpha(z_2) = z_1$. Deux points distincts
arbitraires $z_1$ et $z_2$ étant donnés dans $D(0,1)$,
montrer qu'il existe un automorphisme les échangeant et que
cet automorphisme est unique à une rotation près (on se ramènera au cas où l'un
des points est l'origine).
\finenonce{002886}
\finexercice


\exercice{2887, burnol, 2009/12/15}
\enonce{002887}{}
  Trouver l'unique automorphisme du
premier quadrant qui échange $1+i$ et $2+2i$. On remarquera
que $z\mapsto z^2$ est une bijection analytique du premier
quadrant sur le demi-plan supérieur, et que l'on peut donc
ramener le problème à une question dans le demi-plan supérieur.
\finenonce{002887}
\finexercice


\exercice{2888, burnol, 2009/12/15}
\enonce{002888}{} Soit $f$  holomorphe sur
$\overline{D(0,1)}$. On suppose  $|f(w)|\leq 8$ pour tout
$|w|\leq1 $ et $ f(\frac34) = 0$. Montrer $ |f(0)|\leq
6$. 
\emph{Indication :} trouver un automorphisme $\phi$ du disque
avec $\phi(0) = \frac34$ et utiliser le Lemme de Schwarz
pour la fonction $\frac18 f(\phi(z))$. Trouver le $z$ avec
$\phi(z) = 0$.
\finenonce{002888}
\finexercice

\exercice{6731, queffelec, 2011/10/16}
\enonce{006731}{} 
Calculer $\int_C(x+2y)\ dx+(y-2x)\ dy$ le long de l'ellipse $C$
définie par $x=4\cos\theta $, $y=3\sin \theta $, $0\le \theta \le 2\pi$.
\finenonce{006731}



\finexercice       


\exercice{6732, queffelec, 2011/10/16}
\enonce{006732}{} 
Calculer $\int_C(z^2+3z)\ dz$ le long des chemins suivants :
\begin{enumerate}
\item le cercle $\vert z\vert=2$ du point (2,0) au point (0,2)
\item le segment de droite joignant les points (2,0) et (0,2)
\item le contour polygonal formé par les segments de droite joignant
(2,0) à (2,2) et (2,2) à (0,2)
\end{enumerate}
\finenonce{006732}



\finexercice       


\exercice{6733, queffelec, 2011/10/16}
\enonce{006733}{} 
Calculer $\int_C{dz\over (z-a)^n}$ pour $n=1,2,3\dots$, où $C$ est
un cercle de centre $a$.
\finenonce{006733}



\finexercice       


\exercice{6734, queffelec, 2011/10/16}
\enonce{006734}{} 
Soit $P(x,y)$ et $Q(x,y)$ des fonctions continues à valeurs réelles
et à dérivées partielles continues sur un ouvert connexe $\Omega $ et
sur sa frontière $C$. La formule de Green établit que
$$\int_CPdx+Qdy= \int\!\int_\Omega \left({\partial Q\over \partial
x}-{\partial P\over \partial y}\right) dxdy$$
\begin{enumerate}
\item Montrer la formule de Green pour une courbe fermée simple $C$ ayant
la propriété d'être rencontrée par des parallèles aux axes de
coordonnées en deux points au plus.

\item Si $f(z,\overline{z})=u(x,y)+iv(x,y)$ est continue et possède des
dérivées partielles continues dans un ouvert connexe $\Omega $ et sur
sa frontière $C$, montrer que la formule de Green peut s'écrire sous la
forme complexe suivante
$$\int_Cf(z,\overline{z})dz=2i\int\!\int_\Omega {\partial f\over
\partial \overline{z}}dxdy$$

\item Si $C$ est une courbe fermée simple délimitant un ouvert d'aire $A$,
montrer que $A={1\over 2i}\int_C\overline{z}dz$.

\item Calculer $\int_C\overline{z}dz$ le long
\begin{enumerate}
\item du cercle $\vert z-2\vert =3$
\item du carré de sommets $z=0$, $z=2$, $z=2i$ et $z=2+2i$
\item de l'ellipse $\vert z-3\vert +\vert z+3\vert =10$
\end{enumerate}
\end{enumerate}
\finenonce{006734}



\finexercice       


\exercice{6735, queffelec, 2011/10/16}
\enonce{006735}{} 
Soit $f$ une fonction holomorphe sur un ouvert connexe
$\Omega $. Soit $c$ un point de $\Omega $ et $r_0>0$ tel que
$D(c,r_0)\subset \Omega $, où $D(c,r_0)$ est le disque ouvert de centre
$c$ et de rayon
$r_0$. On pose $\mu (r)={1\over 2\pi}\int_0^{2\pi}f(c+re^{i\theta
})d\theta $ pour $0<r<r_0$.
\begin{enumerate}
\item Montrer que 
$$\lim_{r\to 0^+}\mu (r)=f(c)$$

\item On suppose $f'(z)$ continue. Montrer que $\mu $ est constante (on
montrera que $\displaystyle{d\mu \over dr}=0$ en dérivant sous le signe
d'intégration).

Soit maintenant $M=\sup_{z\in \Omega }\vert f(z)\vert$ et on suppose
qu'il existe $c\in \Omega $ tel que $\vert f(c)\vert =M$. 

\item Montrer que $M=\vert f(c+re^{i\theta })$ où $r>0$ est tel que
$D(c,r)\subset \Omega $.

\item Soit $V=\{ z\in \Omega \vert \ \vert f(z)\vert =M\}$. Montrer que
$V$ est à la fois un ouvert et un fermé de $\Omega $. En déduire le
principe du maximum : si $f$ atteint son maximum en un point d'un
ouvert connexe $\Omega $, alors $f$ est constante.
\end{enumerate}
\finenonce{006735}



\finexercice       


\exercice{6736, queffelec, 2011/10/16}
\enonce{006736}{} 
\begin{enumerate}
\item Soit $f(z)={1\over 1+z^4}$. Montrer que  $\int_{\gamma _R}f(z)\
dz$ tend vers 0 quand $R$ tend vers $+\infty$, où $\gamma _R
(t)=Re^{it}$ ($t\in [0,\pi]$).

\item Déduire de 1. la valeur de
$$\int_{-\infty}^{+\infty}{dx\over 1+x^4}$$
\end{enumerate}
\finenonce{006736}



\finexercice       


\exercice{6737, queffelec, 2011/10/16}
\enonce{006737}{} 
Si $\gamma $ est l'arc de courbe  
$\gamma (t)=t+i(t^3-3t^2+4t-1)$ joignant 
les points (1,1) et (2,3), trouver la valeur de
$$\int_\gamma (12z^2-4iz)\ dz$$
\finenonce{006737}



\finexercice


\exercice{6738, queffelec, 2011/10/16}
\enonce{006738}{}
\begin{enumerate}
\item Calculer l'intégrale $I=\int_{\gamma} \bar z\ dz$ où $\gamma$ est le chemin
joignant le point $(1,1)$ au point $(2,4)$ suivant la parabole $y=x^2$; puis le
segment joignant ces points. Qu'obtient-on avec $\int_{\gamma}  z\ dz$ ?

\item Soit $\gamma=e^{2i\pi t}, t\in [0,1]$ et  $f$
continue sur
$\gamma^*$ le cercle unité dans $\Cc$. Comparer $\int_\gamma
\overline{f(z)}\ {dz\over z^2}$ et $\int_\gamma
f(z)\ dz$. 
\end{enumerate}
\finenonce{006738}



\finexercice


\exercice{6739, queffelec, 2011/10/16}
\enonce{006739}{}
Soit $f$ une fonction continue du quart de plan $\{x+iy;\ x,y\geq0\}$ dans $\Cc$ et $C$ le quart de cercle paramétré par $\{re^{it};\ 0\leq t\leq
{\pi\over2}\}$ avec
$r>0$. On pose $M(r)=\sup_{z\in C} |f(z)|$. Montrer que  si $\lim_{r\to\infty}
M(r)=0$, alors $\lim_{r\to\infty}\int_{C}f(z)e^{iz}\ dz=0$.
\finenonce{006739}



\finexercice


\exercice{6740, queffelec, 2011/10/16}
\enonce{006740}{}
Soit $w=|w|e^{i\theta}$ un nombre complexe.
\begin{enumerate}
\item Montrer que $e^w-1=\int_{[0,w]}e^z\ dz=\int_0^{|w|}
e^{te^{i\theta}}e^{i\theta}\ dt$; montrer ainsi l'inégalité
$|e^w-1|\leq|w|e^{|w|}$ (autre démonstration ?).

\item Application. On considère $K$ un compact du plan inclus dans $\Cc^*$ et la suite de fonctions $v_n$
définie sur $K$ par $v_n(z)={1\over z(1+{1\over n})^z }$. Montrer que $v_n(z)$
tend vers ${1\over z}$ uniformément sur $K$.
\end{enumerate}
\finenonce{006740}



\finexercice


\exercice{6741, queffelec, 2011/10/16}
\enonce{006741}{}
Soit $\varphi$ une fonction continue sur le bord orienté $\partial K$ d'un
compact
$K$; soit $\Omega ={\Cc}\backslash \partial K$ : pour
$z\in \Omega$, on définit
$$f(z)=\int_{\partial K}{\varphi(u)\over u-z}\ du.$$ On va établir que $f$ est
holomorphe dans
$\Omega$. Fixons $a\in\Omega$ et
posons $r=d(a, \partial K)>0$.
\begin{enumerate}
\item Soit $0<\rho<r$ et $z\in \bar B(a,\rho)$. Montrer, en développant ${1\over
u-z}$ en série entière de $z-a$, que $f$ est somme d'une série entière
au voisinage de $a$; en déduire que $f\in H(\Omega)$.

\item Montrer que $f$ est indéfiniment dérivable en tout point $a$ de $\Omega$ et
que 
$$f^{(n)}(a)=n!\int_I{\varphi(u)\over (u-a)^{n+1}}\ du.$$
\end{enumerate}
\finenonce{006741}



\finexercice


\exercice{6761, queffelec, 2011/10/16}
\enonce{006761}{} 
Montrer que 
$$\prod_{n\ge 1}\left( 1+{e^{-nz}\over n^2}\right)$$
définit une fonction holomorphe sur $\Omega =\{z\in\C\vert \ \Re\!
ez>0\}$.
\finenonce{006761}



\finexercice       


\exercice{6762, queffelec, 2011/10/16}
\enonce{006762}{} 
On se propose de démontrer que pour tout $z$ de $\C$,
$$\prod_{n\ge 1}\left( 1-{z^2\over n^2}\right)={\sin{(\pi z)}\over
\pi z}$$
\begin{enumerate}  
\item On pose
$$F(z) ={1\over z}+\sum_{n\in {\Zz}, n\ne 0}\left({1\over z-n}+{1\over
n}\right)$$
Montrer que $F$ est une fonction méromorphe sur $\C$.

En utilisant 
$$\sum_{n\in{\Zz}}{1\over (z-n)^2}=\left({\pi\over \sin{(\pi
z)}}\right)^2$$
 montrer que $\displaystyle F(z)-{\pi\over \tan{(\pi z)}}$ est constante
sur $\C\setminus {\Zz}$, puis calculer cette constante par un argument
de parité. En déduire que
$${\pi\over \tan{(\pi z)}}={1\over z}+\sum_{n\ge 1}{2z\over z^2-n^2}$$

\item Pour $n\ge 1$, soit $\displaystyle f_n(z)=1-{z^2\over n^2}$. Montrer que
$\prod_{n\ge 1}f_n(z)$ définit une fonction entière $f$.

Montrer que
$$\sum_{n\ge 1}{f_n'(z)\over f_n(z)}={g'(z)\over g(z)}$$
avec $\displaystyle g(z)={\sin{(\pi z)}\over \pi z}$. En déduire le résultat voulu.


\item  Déduire de la décomposition de $\displaystyle {\sin{(\pi z)}\over
\pi z}$ en produit que
\begin{eqnarray*}
{\sh z\over z}&=&\prod_{n\ge 1}\left( 1+{z^2\over n^2\pi^2}\right)\\
\ch z&=&\prod_{n\ge 1}\left( 1+{4z^2\over (2n-1)^2\pi^2}\right)\\
\cos z&=&\prod_{n\ge 1}\left( 1-{4z^2\over (2n-1)^2\pi^2}\right)\\
\end{eqnarray*}
\end{enumerate}
\finenonce{006762}



\finexercice       


\exercice{6763, queffelec, 2011/10/16}
\enonce{006763}{} 
On pose $\displaystyle F(z)=\int_0^{+\infty} e^{-zx^2}\
dx$.

\begin{enumerate}  
\item Montrer que $F$ est holomorphe sur $\Omega =\{z\in{\Cc}\vert \Re\!
ez>0\}$. Calculer $F'(z)$ en fonction de $F(z)$.

\item En déduire une autre expression de $F(z)$ pour $z\in \Omega$ (on
pourra utiliser $\displaystyle\int_0^{+\infty} e^{-x^2}\ dx =\sqrt\pi/2$).
Conclure enfin que la fonction $F$ se prolonge en une fonction
holomorphe sur $\C\setminus \mathopen]-\infty, 0]$.
\end{enumerate}
\finenonce{006763}



\finexercice       


\exercice{6764, queffelec, 2011/10/16}
\enonce{006764}{Fonction $\zeta $ de Riemann}

On introduit la fonction ``Zeta'' :
$$\zeta (s)=\sum_{n\ge 1}{1\over n^s}$$

\subparagraph{I  Produit d'Euler}

Montrer que $\zeta $ est holomorphe
dans l'ouvert $\Omega =\{ z\in\C\vert \ \Re z>1\}$.

Soient $p_1=2,p_2=3,\dots,p_n,\dots$ la suite des nombres premiers.
Montrer que dans $\Omega $, on a $$\zeta (s)=\prod_{n\ge 1}{1\over
1-p_n^{-s}}$$ 
(produit d'Euler).

\subparagraph{II Relation de $\zeta $ avec la répartition des
nombres premiers}

\begin{enumerate}  
\item Montrer que $${\zeta '(s)\over \zeta (s)}=-\sum_{n\ge1}
\lambda (n)n^{-s}$$ où $\lambda (n)=\ln p$ si $n$ est une puissance
d'un nombre $p$ premier et $\lambda (n)=0$ si $n$ a au moins deux
diviseurs premiers distincts.

\item On a le théorème suivant :\\
Théorème des nombres premiers (Hadamard-De la Vallée Poussin 1896) :
Lorsque $x$ tend vers $+\infty$, la somme des $\lambda (n)$ pour $n\le
x$ est équivalente à $x$.

Démontrer que cette assertion est équivalente à dire que le nombre de
nombres premiers plus petits que $x$ est équivalent à $x/\ln x$.
\end{enumerate}


\subparagraph{III Equation fonctionnelle de $\zeta $, démonstration
par la formule de Poisson}

Soit 
$$\theta (t)=\sum_{n\in {\Zz}}e^{-\pi n^2t}$$
$\theta $ vérifie l'équation fonctionnelle suivante :
$$\forall t>0,\ \theta \left( {1\over t}\right)=\sqrt t\theta (t)$$
On pose
$$\psi (t)=\sum_{n\ge 1}e^{-\pi n^2t}$$
On a évidemment $\theta (t)=1+2\psi (t)$.

\begin{enumerate}  
\item Soit $s$ tel que $\Re s>1/2$. Calculer $\int_0^{+\infty}e^{-\pi
n^2t}t^{s-1}\ dt$ en fonction de $n$ et $\Gamma (s)$. En déduire que
$$\int_0^{+\infty}\psi (t)t^{s-1}\ dt=\pi^{-s}\Gamma (s)\zeta (2s)$$

\item Montrer que $F$ définie par
$$F(z)=\int_1^{+\infty}\psi (t)t^z\ dt$$
est une fonction entière. En écrivant
$$\int_0^{+\infty}\psi (t)t^z\ dt=\int_0^1\psi (t)t^z\
dt+\int_1^{+\infty}\psi (t)t^z\ dt$$
et en effectuant dans la première partie de la somme le changement de
variable $u=1/t$, puis en utilisant l'équation fonctionnelle de $\theta $,
montrer que la fonction $L$ :
$$L(s)=s(s-1)\pi^{-s/2}\Gamma \left({s\over 2}\right)\zeta (s)$$
se prolonge en une fonction entière, invariante par la transformation
$s$ donne $1-s$. Montrer alors que $\zeta $ se prolonge en une fonction
méromorphe sur $\C$, ayant 1 pour seule singularité (pôle simple).
Montrer enfin que $\mathrm{Res} (\zeta ,1)=1$.
\end{enumerate}

\subparagraph{IV  Equation fonctionnelle de $\zeta $, démonstration
par le théorème des résidus}

\begin{enumerate}  
\item Soit $s$ tel que $\Re(s)>1$. Par la même méthode qu'au III.1.,
montrer que
$$\int_0^{+\infty}{t^{s-1}\over e^t-1}\ dt=\Gamma (s)\zeta (s)$$

\item Pour tout $s$ de $\C$, on définit $F_s$ sur $\C\setminus
[0,+\infty\mathclose[$ : 
$$F_s(z)={\exp\left[ (s-1)\mathrm{Log}{(-z)}\right]\over
\exp z-1}$$
où $\mathrm{Log}$ est la détermination principale du logarithme. On définit 
également le contour $A_{\varepsilon ,\varphi }$ :



Calculer 
$$\lim_{\varphi \to 0}\int_{A_{(\varepsilon ,\varphi )}}F_s(z)\ dz$$
Montrer que pour $\Re s>1$ :
$$\lim_{\varepsilon \to 0}\left(\lim_{\varphi \to
0}\int_{A_{(\varepsilon ,\varphi )}}F_s(z)\ dz\right)=2i\sin{(\pi
z)}\Gamma (s)\zeta (s)$$ Montrer que $\int_{A_{(\varepsilon ,\varphi
)}}F_s(z)\ dz$ est en fait indépendant de $\varepsilon \in
\mathopen]0,1\mathclose[$, puis que $$s\mapsto \lim_{\varphi \to
0}\int_{A_{(\varepsilon ,\varphi )}}F_s(z)\ dz$$ est une fonction entière.
En déduire que $\sin{(\pi s)}\Gamma (s)\zeta (s)$ se prolonge en une
fonction entière.

\item Soit $C_{n,\varepsilon ,\varphi }$ le chemin fermé :


En appliquant la formule des résidus à $F_s$ sur $C_{n,\varepsilon
,\varphi }$ pour $\Re s<0$, montrer que
$$\forall s\in\C\setminus\{-{\Nn}\},\ \zeta
(1-s)=2^{1-s}\pi^{-s}\Gamma (s)\cos{\left({\pi s\over 2}\right)}\zeta (s)
$$

\item Déduire de 3. que
$$\forall s\in \C\setminus {\Zz},\ \Gamma (s)\Gamma
(1-s)={\pi\over \sin{(\pi s)}}$$
En utilisant l'équation fonctionnelle de {\bf III}, montrer alors que
$$\forall s\in \C\setminus {\Zz},\ \Gamma (s)={2^{s-1}\over
\sqrt\pi}\Gamma \left({1+s\over 2}\right)\Gamma \left({s\over
2}\right)$$
\end{enumerate}
\finenonce{006764}



\finexercice       


\exercice{6765, queffelec, 2011/10/16}
\enonce{006765}{} 
Soit $f$ holomorphe sur $\C$, réelle sur l'axe réel, imaginaire pure
sur l'axe imaginaire. Montrer que $f$ est impaire.
\finenonce{006765}



\finexercice       


\exercice{6766, queffelec, 2011/10/16}
\enonce{006766}{} 
Montrer que toute fonction $f$ holomorphe dans un ouvert connexe
$\Omega $ symétrique par rapport à l'axe réel peut s'écrire
$f=f_1+if_2$, où $f_1$ et $f_2$ sont holomorphes dans $\Omega $ et
réelles sur l'axe réel.
\finenonce{006766}



\finexercice       


\exercice{6767, queffelec, 2011/10/16}
\enonce{006767}{} 
Soit $f$ une fonction holomorphe dans le disque unité $D$, continue
sur $\overline{D}$, telle que $\vert f(z)\vert =1$ si $\vert z\vert=1$.
Montrer que $f$ est rationnelle.
\finenonce{006767}



\finexercice    


\exercice{6768, queffelec, 2011/10/16}
\enonce{006768}{} 
Donner un biholomorphisme entre les ouverts
$$\Omega _1=\{ z\in\C\vert\ \vert z\vert >1,\ \Im\! mz>0\}\mbox{ et }
\Omega _2=\{ z\in\C\vert\ \vert z\vert <1,\ \Im\! mz>0\}$$
\finenonce{006768}



\finexercice       


\exercice{6769, queffelec, 2011/10/16}
\enonce{006769}{} Montrer que la transformation
$$w=\left({1+z^m\over 1-z^m}\right)^2,$$
où $m\in {\Nn}^*$, définit une représentation conforme de
$$\Omega =\left\{ z\in\C\vert\ z=re^{i\theta },\ 0<r<1,\ 0<\theta
<{\pi\over m}\right\}$$ sur le demi-plan supérieur.
\finenonce{006769}



\finexercice       


\exercice{6770, queffelec, 2011/10/16}
\enonce{006770}{} 
Soit $\Omega $ le demi-plan $\Im\! mz>0$ privé du disque fermé de
centre $i$ et de rayon 1. Trouver l'image de $\Omega $ par la
transformation $w=\coth{\left(\pi /z\right)}$.
\finenonce{006770}



\finexercice       


\exercice{6771, queffelec, 2011/10/16}
\enonce{006771}{} 
Soit $C$ une couronne circulaire excentrique. Montrer qu'il existe
une transformation homographique appliquant $C$ sur une couronne
concentrique.
\finenonce{006771}



\finexercice


\exercice{6812, gijs, 2011/10/16}
\enonce{006812}{}
Soit $f:\Cc \to \Cc$ une fonction entière
(c'est-à-dire holomorphe).
\begin{enumerate}
\item
Soit $\alpha$ un réel, $\alpha>0$. On suppose que pour
tout $z\in \Cc$ on a $\vert f(z) \vert \ne \alpha$.
Démontrer que, ou bien $\forall z\in \Cc : \vert f(z)
\vert < \alpha$, ou bien $\forall z\in \Cc : \vert f(z)
\vert > \alpha$. En déduire que $f$ est constante.

\item
On suppose ici que $f$ n'est pas constante.

  \begin{enumerate}
  \item Démontrer que $\inf_{z\in \Cc} \vert
f(z) \vert = 0$ et $\sup_{z\in \Cc} \vert f(z) \vert = +
\infty$.
  \item
Démontrer que $\{ \ \vert f(z) \vert \  ,\ z\in
\Cc\,\}$ est soit $]0,+\infty[$, soit $[0,+\infty[$.
Donner un exemple dans chacun des deux cas.
  \end{enumerate}
\item
On suppose toujours que $f$ n'est pas constante.
Démontrer que $\{f(z) \ ,\ z\in \Cc \}$ est partout
dense dans $\Cc$. 
\end{enumerate}
\finenonce{006812}



\finexercice


\exercice{6825, gijs, 2011/10/16}
\enonce{006825}{}
\begin{enumerate}
\item  Démontrer que l'intégrale
$ \int_0^\infty \dfrac{\ln(x)}{(1+x^2)^2}\,dx$
converge.

\item Soit $f(z) = \dfrac{\log(z)}{(1+z^2)^2}$ avec
$\log(z) = \mathrm{Log}(-iz) + i\pi/2$ le logarithme défini sur
$\Cc\setminus i\,]-\infty,0]$ et tel que $\log(1) = 0$.
Déterminer les points singuliers isolés de $f$ et
pour chaque point singulier isolé déterminer son
résidu.

\item Soit $D = \{\,z\in \Cc \mid \epsilon< |z|
< R\ \&\ \Im z >0\,\}$. \`A l'aide de $\int_{\partial D}
f(z)\,dz$, déterminer la valeur de $ \int_0^\infty
\dfrac{\ln(x)}{(1+x^2)^2}\,dx$. N'oubliez pas de
justifier les passages à la limite que vous effectuez.
\end{enumerate}
\finenonce{006825}



\finexercice


\exercice{6826, gijs, 2011/10/16}
\enonce{006826}{}
\begin{enumerate}
\item \'Enoncer le principe du maximum.

\item Soit $a_1, \dots, a_n \in \Cc$ des points sur le
cercle unité ($|a_j| = 1$). Démontrer qu'il
existe un point $z_o$ sur le cercle unité tel que le
produit des distances de $z_o$ à $a_j$ ($j=1,\dots,n$)
est supérieur ou égal à 1. 
\end{enumerate}
\finenonce{006826}



\finexercice


\exercice{6830, gijs, 2011/10/16}
\enonce{006830}{}
Soit $f$ une fonction holomorphe sur $B_r(0) \subset \Cc$
pour un certain rayon $r>1$. Démontrer les égalités
suivantes~: 
\begin{equation*}
\tfrac2\pi \int_0^{2\pi} f(e^{i\theta})
\cos^2(\tfrac\theta2) \, d\theta = 2f(0) + f'(0)
\\
\tfrac2\pi \int_0^{2\pi} f(e^{i\theta})
\sin^2(\tfrac\theta2) \, d\theta = 2f(0) - f'(0)
\end{equation*}
Indication~: contempler les deux intégrales
$  \int_{|z|=1} \Bigl(2\pm (z + \dfrac1z)\Bigr) \,
\dfrac{f(z)}z\, dz$.
\finenonce{006830}



\finexercice


\exercice{6831, gijs, 2011/10/16}
\enonce{006831}{} 
Soit $a$ un réel, $0<a<2$.
\begin{enumerate}
\item  Démontrer que l'intégrale
$ \int_0^\infty \dfrac{x^a}{x(1+x^2)}\,dx$
converge.

\item Soit $f(z) = \dfrac{e^{a\log(z)}}{z(1+z^2)}$
avec $\log(z) = \mathrm{Log}(-iz) + i\pi/2$, c'est-à-dire que
$\log$ est le logarithme défini sur $\Omega
= \Cc\setminus i\,]-\infty,0]$ et tel que $\log(1) = 0$.
Déterminer les points singuliers isolés de $f$ dans
$\Omega$ et pour chaque point singulier isolé
déterminer son résidu.

\item Soit $D = \{\,z\in \Cc \mid \epsilon< |z|
< R\ \&\ \Im z >0\,\}$. \`A l'aide de $\int_{\partial D}
f(z)\,dz$, déterminer la valeur de 
$ \int_0^\infty \dfrac{x^a}{x(1+x^2)}\,dx$.
N'oubliez pas de justifier les passages à la limite que
vous effectuez.
\end{enumerate}
\finenonce{006831}



\finexercice


\exercice{6836, gijs, 2011/10/16}
\enonce{006836}{}
Soit $P$ et $Q$ deux polyn\^omes à coefficients complexes
sans zéro commun et soit $z_1, \dots, z_k \in \Cc$
les zéros de $Q$ (ce qui implique que le degré de $Q$
est supérieur ou égal à $k$). On définit la
fonction $f:\Cc \setminus \{z_1, \dots, z_k\} \to \Cc$
par $f(z) = P(z)/Q(z)$.
\begin{enumerate}
\item
Démontrer qu'il existe
une fonction continue $g: \Cc_\infty \to \Cc_\infty$ telle
que $\forall z\in \Cc\setminus\{z_1, \dots, z_k\}$~: $g(z)
= f(z)$. Quelle est la valeur de $g$ en un point $z_i$~?
Quelle est la valeur de $g$ en $\infty$~? (N'oubliez pas
de démontrer que la fonction $g$ que vous définissez
est continue.)

\item
Le résultat reste-t-il vrai si $P$ et
$Q$ ont des zéros  communs? Si non, donner un contre
exemple~; si oui, esquisser votre raisonnement.
\end{enumerate}
\finenonce{006836}



\finexercice


\exercice{6838, gijs, 2011/10/16}
\enonce{006838}{}
\begin{enumerate}
\item
\ Quand dit-on qu'une partie $S$ de $\Rr^n$ est
discrète~?

\item
Soit $K\subset \Rr^n$ un compact et soit $S \subset
\Rr^n$ une partie discrète et fermée. 
Dé\-mon\-trer que
$K \cap S$ est fini.

\item
Soit $\Omega\subset \Cc$ un ouvert connexe,
$f:\Omega\to \Cc$ une fonction holomorphe
non-identiquement nulle et soit $K\subset \Omega$ un
compact. Démontrer que $f$ n'a qu'un nombre fini de
zéros dans $K$.
\end{enumerate}
\finenonce{006838}



\finexercice


\exercice{6842, gijs, 2011/10/16}
\enonce{006842}{}
Soit $f(z) = \dfrac{1-e^{2iz^2}}{z^3(1+z^4)}$,  soit
$0<\epsilon<1<R$ et soit $D = B_\epsilon(0) \cup \{z\in
\Cc \mid {\Re z>0,}\ {\Im z > 0,}\ |z|<R\,\}$.
\begin{enumerate}
\item Démontrer que l'intégrale
$ \int_0^\infty
\dfrac{\sin^2(x^2)}{x^3(1+x^4)}\,dx$ converge.

\item Dessiner $D$, déterminer les singularités
 de $f$
%, calculer le résidu en chacun des point
%singuliers isolés 
et calculer $  \int_{\partial D}
f(z) \, dz$.

\item Déduire de 2. la valeur de 
$ \int_0^\infty 
\dfrac{\sin^2(x^2)}{x^3(1+x^4)}\,dx$. 
N'oubliez pas de justifier les passages à la limite que
vous effectuez.
\end{enumerate}
\finenonce{006842}



\finexercice


\exercice{6844, gijs, 2011/10/16}
\enonce{006844}{}
Soit $f:\Cc \to \Cc$ une application holomorphe et soit
$a$ un réel strictement positif.
\begin{enumerate}
\item Démontrer que $\{ |f(z)| \ \mid z \in \Cc\,\}$
est connexe.

\item Démontrer que si pour tout $z\in \Cc$ on a
$|f(z)| \neq a$, alors on a ou bien $\forall z\in \Cc :
|f(z)| > a$, ou bien $\forall z\in \Cc : |f(z)| < a$. En
déduire que $f$ est constante.

\item En utilisant le résultat de 2.,
démontrer que si $f$ n'est pas constante, alors $ 
\inf_{z\in
\Cc} |f(z)| = 0$ et $  \sup_{z\in \Cc} |f(z)|=
+\infty$.
\end{enumerate}
\finenonce{006844}



\finexercice


\exercice{6847, gijs, 2011/10/16}
\enonce{006847}{}
Soit $n$, $p$ deux entiers tels que $n\ge p+2 \ge 2$,
soit $f(z) = \dfrac{z^p}{1+z^n}$,  soit
$R>0$ et soit $D =  \{z\in
\Cc \mid 0<\mathrm{Arg}(z)<2\pi/n\ ,\  0<|z|<R\,\}$.
\begin{enumerate}
\item Démontrer que l'intégrale
$ \int_0^\infty
\dfrac{x^p}{1+x^n}\,dx$ converge.

\item Dessiner $D$, déterminer les singularités
isolés de $f$
%, calculer le résidu en chacun des point
%singuliers isolés 
et calculer $  \int_{\partial D}
f(z) \, dz$.

\item Déduire de 2. la valeur de 
$ \int_0^\infty 
\dfrac{x^p}{1+x^n}\,dx$. 
N'oubliez pas de justifier les passages à la limite que
vous effectuez.
\end{enumerate}
\finenonce{006847}



\finexercice


\exercice{6848, gijs, 2011/10/16}
\enonce{006848}{}
Soit $R>1$ et $\gamma _R$ le chemin fermé de
classe $C^1$ : $\gamma _R:[0,2\pi]\to \C$ tel que $\gamma
_R(t)=Re^{it}$.

On note respectivement $\gamma ^+_R$ et $\gamma ^-_R$ la restriction
de $\gamma _R$ à $[0,\pi]$ et à $[\pi,2\pi]$. On rappelle que $[a,b]$
désigne le chemin $\zeta $ de $[0,1]$ dans $\C$ défini par $\zeta
(t)=bt+(1-t)a$. On pose $C^+_R=\gamma ^+_R+[-R,R]$ et $C^-_R=\gamma
^-_R+[R,-R]$.

Montrer que $\mathrm{Ind}_{C^+_R}(i)+\mathrm{Ind}_{C^-_R}(i)=\mathrm{Ind}_{\gamma _R}(i)$ et
en déduire la valeur de $\mathrm{Ind}_{C^+_R}(i)$. Calculer, pour $u>0$ et $R>1$,
$$\int_{C^+_R}{e^{iuz}\over 1+z^2}dz.$$
En déduire la valeur de
$$\int_{-\infty}^{+\infty}{e^{iux}\over 1+x^2}dx,$$
pour $u>0$, puis pour tout $u$ réel.
\finenonce{006848}



\finexercice       


\exercice{6849, gijs, 2011/10/16}
\enonce{006849}{}
Soit $\alpha $, $0<\alpha <{\pi\over 2}$ et
$\Delta =\{ z\in\C, z=re^{i\theta }\vert\ 0<r,\vert \theta \vert<\alpha
\}$. On note, pour $\eta >0$, $\Delta_\eta  =\{ z\in\C, z=re^{i\theta
}\vert\ \eta <r,\vert \theta \vert<\alpha \}$. Dessiner $\Delta $ et
$\Delta _\eta $.

Dans tout ce qui suit, $f$ désigne une fonction holomorphe sur $\Delta
$, bornée et continue sur $\overline{\Delta }$, l'adhérence de $\Delta $.
On note $\partial \Delta $, la frontière de $\Delta $.

On pose $\displaystyle\tilde M=\sup_{z\in \partial \Delta }\vert f(z)\vert$ et, pour
tout $r>0$, $\displaystyle M_r=\sup_{z\in \Delta,\vert z\vert =r }\vert
f(z)\vert$.

\begin{enumerate}
\item On suppose
$$\lim_{z\in \Delta ,\vert z\vert\to\infty}\vert f(z)\vert =0.\leqno{(*)}$$

Soit $R>0$, montrer que l'on a
$$\sup_{z\in\overline{\Delta },\vert z\vert \le R}\vert f(z)\vert\le
\max{(\tilde M,M_R)}$$
et donc
$$\sup_{z\in \overline{\Delta }}\vert f(z)\vert \le \sup_{z\in \partial
\Delta }\vert f(z)\vert.$$

\item On ne suppose plus (*). On pose, pour tout $n\in {\Nn}$, $n>0$,
et tout $z\ne 0$, $g_n(z)={f(z)^n\over z}$. En majorant, comme dans la
question 1. $\sup_{z\in\overline{\Delta _\eta },\vert z\vert\le
R}\vert g_n(z)\vert$, montrer que l'on a, pour tout $n>0$, tout $\eta >0$
et tout $z$ de $\Delta _\eta $,
$$\left\vert{f(z)^n\over z}\right\vert\le {1\over \eta }\max{\left(
(\tilde M)^n,(M_\eta )^n\right) }$$
et de là
$$\sup_{z\in \overline{\Delta _\eta }}\vert f(z)\vert \le
\max{\left(\tilde M,M_\eta \right) }.$$
En déduire que l'on a encore
$$\sup_{z\in \overline{\Delta }}\vert f(z)\vert \le \sup_{z\in \partial
\Delta }\vert f(z)\vert.$$

\item On suppose maintenant que l'on a $\lim_{r\to\infty}\vert
f(re^{i\alpha })\vert =0$ et $\lim_{r\to\infty}\vert
f(re^{-i\alpha })\vert =0$. On se propose de montrer qu'alors
$\lim_{z\in \Delta ,\vert z\vert \to\infty}\vert f(z)\vert =0$.

En majorant la fonction $z\mapsto \left\vert {z\over z+A}\right\vert$,
$A>0$, montrer que, pour tout $\varepsilon >0$, il existe $R>0$ tel que
l'on ait, pour tout $A>0$,
$$\sup_{z\in\partial \Delta ,\vert z\vert>R}\left\vert{z\over
z+A}f(z)\right\vert\le \varepsilon .$$
Montrer qu'alors il existe $A>0$ tel que l'on ait
$$\sup_{z\in\partial \Delta }\left\vert{z\over
z+A}f(z)\right\vert\le \varepsilon .$$
En utilisant la question 2., en déduire que l'on a
$\lim_{z\in \Delta ,\vert z\vert \to\infty}\vert f(z)\vert =0$.
\end{enumerate}
\finenonce{006849}



\finexercice       


\exercice{6850, gijs, 2011/10/16}
\enonce{006850}{}
Soit $t$ un réel, $\vert t\vert \le \pi$.
\begin{enumerate}
\item On considère la série de fonctions holomorphes
$$\sum_{k=-\infty}^{+\infty}{(-1)^ke^{ikt}\over k^2-z^2}.$$
Montrer que sa somme $S(z)$ est une fonction méromorphe dans $\C$.

\item Soit $\gamma _n$, $n\ge 0$, le chemin parcouru dans le sens
positif dont l'image $\gamma _n^*$ dans $\C$ est le carré de sommets
$\left(n+{1\over 2}\right)(1+i)$, $\left(n+{1\over 2}\right)(-1+i)$,
$\left(n+{1\over 2}\right)(-1-i)$, $\left(n+{1\over 2}\right)(1-i)$.

  \begin{enumerate}  
  \item Montrer que, quel que soit $z$ vérifiant $z=\left(n+{1\over
2}\right)+iy$, $-\left(n+{1\over 2}\right)\le y\le n+{1\over 2}$, on
a 
$$\left\vert{e^{itz}\over \sin{\pi z}}\right\vert\le 2$$
et que, quel que soit $z$ vérifiant $z=x+i\left(n+{1\over
2}\right)$, $-\left(n+{1\over 2}\right)\le x\le n+{1\over 2}$, on
a 
$$\left\vert{e^{itz}\over \sin{\pi z}}\right\vert\le {2\over 1-e^{-\pi}}.$$
En déduire que $\left\vert{e^{itz}\over \sin{\pi z}}\right\vert$ est
bornée sur $\gamma _n^*$.

  \item Soit $a$ appartenant à $\C\setminus {\Zz}$. On pose
$$f(z)={e^{itz}\over (z^2-a^2)\sin{\pi z}}.$$
Montrer que l'on a
$$\lim_{n\to +\infty}\int_{\gamma _n}f(z)dz=0.$$
  \end{enumerate}
\item Calculer le résidu de $f$ en chacun de ses pôles. En déduire la
valeur de
$$S(a)=\sum_{k=-\infty}^{+\infty}{(-1)^ke^{ikt}\over k^2-a^2}.$$

\item Calculer de même
$$\sum_{k=-\infty}^{+\infty}{(-1)^ke^{ikt}\over (k+a)^2}.$$
\end{enumerate}
\finenonce{006850}



\finexercice


\exercice{6851, gijs, 2011/10/16}
\enonce{006851}{}
Soit $a$ un réel, $a>1$. On considère la série
entière
$$\sum_{n=0}^{+\infty}{z^n\over a^{n^2}}.$$
\begin{enumerate}
\item Montrer que sa somme, notée $f(z)$, est holomorphe dans $\C$.
Montrer qu'il existe $a_0>1$ tel que l'on ait, pour tout $a\ge a_0$,
$$\sum_{k=1}^{+\infty}{1\over a^{k^2}}\le {1\over 100}.$$

\item \emph{Dans toute la suite, on supposera $a\ge a_0$.}\\
Soit $p$ entier, $p\ge 1$. Montrer que l'on a, pour tout $z$ vérifiant
$\vert z\vert =a^{2p}$,
$${1\over a^{p^2}}\left\vert f(z)-{z^p\over a^{p^2}}\right\vert\le
{2\over 100}.$$
En déduire que $f(z)$ a $p$ zéros, $z_1, \dots ,z_p$ dans le disque
ouvert $D(0,a^{2p})$. Montrer que, quel que soit $p\ge 1$, $z_p$ a les
propriétés suivantes :
\begin{enumerate}
\item $a^{2(p-1)}<\vert z\vert <a^{2p}$,
\item $z_p$ est un zéro simple,
\item $z_p$ est un réel négatif.
\end{enumerate}
Pour établir (c), on pourra raisonner par l'absurde.
\end{enumerate}
\finenonce{006851}



\finexercice


\exercice{6852, gijs, 2011/10/16}
\enonce{006852}{}
Soit $f$ une fonction holomorphe et bornée dans
le disque ouvert $D(0,1)$, vérifiant $f(0)=0$ et $f'(0)=1$. On pose
$M=\sup_{\vert z\vert <1}\vert f(z)\vert$.
\begin{enumerate}
\item Montrer que les modules des coefficients du développement de
$f$ en série entière au voisinage de 0 sont majorés par $M$.

\item Utiliser 1. pour montrer que, si l'on pose $g(z)=f(z)-z$,
on a
$$\vert g'(z)\vert\le{M\over (1-r)^2}-M,{\rm \ si\ }\vert z\vert \le
r<1.$$
En déduire l'existence d'un réel $\rho $, $0<\rho <1$, \emph{dépendant
seulement de M}, tel que l'on ait
$$\vert g'(z)\vert <1{\rm \ si\ }\vert z\vert <\rho .$$

\item Montrer alors que la restriction de $f$ au disque ouvert
$D(0,\rho )$ est injective (on pourra, pour $z_1$ et $z_2$ appartenant
au disque ouvert $D(0,\rho )$, exprimer $g(z_1)-g(z_2)$ sous forme
d'une intégrale).
\end{enumerate}
\finenonce{006852}



\finexercice       


\exercice{6853, gijs, 2011/10/16}
\enonce{006853}{Question de cours}
Montrer en utilisant le lemme de
Schwarz que l'on énoncera soigneusement, que tout automorphisme du
disque unité $D=\{z\in \C\vert\ \vert z\vert<1\}$ est de la forme
$$f(z)=e^{i\theta }{z-a\over 1-\overline{a}z},$$
où $\theta \in{\Rr}$, $a\in\C$, $\vert a\vert <1$. On admettra que
$g_a$ défini par $\displaystyle g_a(z)={z-a\over 1-\overline{a}z}$, pour
$a\in\C$, $\vert a\vert <1$, est un automorphisme de $D$.
\finenonce{006853}



\finexercice       


\exercice{6854, gijs, 2011/10/16}
\enonce{006854}{}
\begin{enumerate}
\item Montrer que le produit infini
$$P(z)=\prod_{n=1}^{+\infty}{1+{z\over 2n}\over 1+{z\over 2n-1}}$$
converge normalement sur tout compact de $U=\{z\in\C\vert\ \Re\!
ez>0\}$.
\item On pose $f(z)=zP(z)$. Calculer $f(1)$. On rappelle la formule
$$\sin{\pi z}=\pi z\prod_{n=1}^{+\infty}\left(1-{z^2\over
n^2}\right) .$$
\item On pose $g=f'/f$. Ecrire $g$ sous forme de série.
Montrer que $g(z)+g(z+1)=1/z$ (on travaillera sur les sommes
partielles). En déduire que l'on a $f(z)f(z+1)={2\over \pi}z$.
\end{enumerate}
\finenonce{006854}



\finexercice       


\exercice{6855, gijs, 2011/10/16}
\enonce{006855}{}
Soit $\alpha $ un nombre réel, $\alpha >0$.
\subparagraph{I}
Pour tout $z$ dans $\C$ et tout $t$ dans ${\Rr}$, on pose
$$f(z,t)={t^{\alpha -1}\over e^t-z}.$$
\begin{enumerate}
\item
  \begin{enumerate}
  \item On note $[1,+\infty\mathclose[=\{x\in{\Rr}\vert \ x\ge
1\}$ et $\Omega =\C\setminus [1,+\infty\mathclose[$. Soit $K$ un
compact inclus dans $\Omega $. Démontrer brièvement qu'il existe un
réel $c$, $c>0$, tel que l'on ait $\inf_{z\in K}\inf_{x\in
[1,+\infty\mathclose[} \vert x-z\vert \ge c$.
  \item  Montrer que la fonction $F$ définie par
$$F(z)=\int_0^{+\infty}f(z,t)\ dt$$
 est holomorphe dans $\Omega $.
  \end{enumerate}
\item 
  \begin{enumerate}
  \item Déterminer le rayon de convergence de la série entière
$\sum_{k=0}^{+\infty}{z^k\over (k+1)^\alpha }$.\\
  \item Montrer que l'on a, pour tout $z$, $\vert z\vert<1$,
$$F(z)=F(0)\sum_{k=0}^{+\infty}{z^k\over (k+1)^\alpha }.$$
  \item La fonction $F$ admet-elle un prolongement holomorphe au
voisinage de 1 ?
  \end{enumerate}
\end{enumerate}

\subparagraph{II} Dans toute cette partie, $t$ appartient à $\C$.
\begin{enumerate}
\item Soit $a$ un réel, $0<a<\pi$ et $R$ un réel, $R>0$. On note
$Q(a,R)=\{t\in \C\vert \ 0<\Re t<R,\vert \Im mt\vert <a\}$ et
$Q(a)=\{t\in\C\vert \ 0<\Re t,\vert \Im\! mt\vert <a\}$. Soit $\mathrm{Log} z$ la 
détermination du logarithme de $z$ holomorphe dans $\C\setminus
\mathopen]-\infty, 0]$. On a donc $\mathrm{Log} z=\ln{\vert z\vert}+i\mathrm{Arg} z$,
avec $-\pi<\mathrm{Arg} z<\pi$.
  \begin{enumerate}
  \item Déterminer $S(a,R)=\{z\in\C\vert \ \mathrm{Log} z\in Q(a,R)\}$ et $S(a)
=\{z\in\C\vert \ \mathrm{Log} z\in Q(a)\}$. Faire des dessins.
On note, pour $t$ dans $\C\setminus
\mathopen]-\infty, 0]$, $t^{\alpha -1}=e^{(\alpha -1)\mathrm{Log} t}$. Déterminer,
pour $z$ dans $S(a)$, le pôle de la fonction
$$t\mapsto f(z,t)={t^{\alpha -1}\over e^t-z}$$
situé dans la bande $Q(a)$. Quel est le résidu de cette fonction en ce
pôle ?
  \item Pour tout $z$ dans $S(a,R)$, calculer en fonction de $z$
$$\int_{\partial Q_\varepsilon (a,R)}f(z,t)\ dt$$
où $\partial Q_\varepsilon (a,R)$ désigne le chemin parcouru dans le
sens positif dont l'image dans $\C$ est la frontière du domaine
$Q(a,R)\setminus \{t\in\C\vert\ \vert t\vert <\varepsilon \}$.
Quelle est la limite lorsque $\varepsilon $ tend vers 0 de 
$\int_{\partial Q_\varepsilon (a,R)}f(z,t)\ dt$ ?
  \end{enumerate}
\item Soit $\Gamma _a$ le chemin dont l'image $\Gamma _a^*$ et le
sens de parcours sont représentés ci-dessous.


On pose
$$F_a(z)=\int_{\Gamma _a}f(z,t)\ dt.$$
  \begin{enumerate}
  \item En reprenant brièvement les idées utilisées dans I.1.,
montrer que $F_a$ est holomorphe dans $\C$ privé d'un chemin que l'on
précisera et que l'on dessinera.\\
  \item Montrer que, si $z$ appartient à $S(a)$ et $\Im\! m z<0$, on a
$F_a(z)=F(z)$ (on pourra intégrer $f(z,t)$ le long d'un contour bien
choisi).
  \item Soit $\Gamma _a$ le chemin dont l'image $\Gamma _{-a}^*$ et
le sens de parcours sont représentés ci-dessous.


On pose
$$F_{-a}(z)=\int_{\Gamma _{-a}}f(z,t)\ dt.$$
Montrer brièvement que, si $z$ appartient à $S(a)$ et $\Im\! mz>0$, on a
$F_{-a}=F(z)$.
  \item  Montrer, en utilisant II.1.(b), que, pour tout $z$ dans
$S(a)$, on a
$$F_{-a}(z)-F_a(z)=2i\pi{(\mathrm{Log} z)^{\alpha -1}\over z}.$$
En déduire que $F(z)$ n'a pas de limite lorsque $z$ tend vers un point de
la demi-droite réelle $\mathopen]1,+\infty\mathclose[$.
  \end{enumerate}
\end{enumerate}
\finenonce{006855}



\finexercice       


\exercice{6856, gijs, 2011/10/16}
\enonce{006856}{Question de cours}
En énonçant avec précision les
différentes formes du prin\-ci\-pe du maximum utilisées, démontrer le
lemme de Schwarz.

\textbf{Application : }Soit $F$ une fonction holomorphe dans
$D(0,1)=\{ z\in\C\vert \ \vert z\vert <1\}$, nulle à l'origine et vérifiant
$\Re F(z)\le 1$, pour tout $z$ dans $D(0,1)$. Montrer que la fonction
$$f={F\over 2-F}$$
est bornée par 1 dans $D(0,1)$. En déduire une majoration de $\vert
F(z)\vert$ en fonction de $\vert z\vert$. Déterminer $F$ en supposant
de plus que l'on a $F'(0)=2$.
\finenonce{006856}



\finexercice       


\exercice{6857, gijs, 2011/10/16}
\enonce{006857}{}
Calculer, en utilisant le contour (I) ci-contre, l'intégrale
$$I_1=\int_0^{+\infty}{x^3\over 1+x^7}dx$$
et, en utilisant le contour (II) ci-contre, l'intégrale
$$I_2=\int_{-\infty}^{+\infty}{\cos x\over e^x+e^{-x}}dx.$$
\finenonce{006857}



\finexercice       


\exercice{6858, gijs, 2011/10/16}
\enonce{006858}{}
Si $\rho $ est un réel strictement positif, on note $D_\rho $ le disque
ouvert de centre 0 et de rayon $\rho $ et $\gamma _\rho $ le chemin
$t\mapsto \rho e^{it}$, $0\le t\le 2\pi$.

On considère une fonction $f$ holomorphe sur $D_1$, telle que l'on ait
$f(0)=0$ et $f'(0)\ne 0$. Pour tout $\rho \in\mathopen]0,1\mathclose[$,
on pose $\displaystyle m(\rho )=\inf_{\vert z\vert =\rho }\vert f(z)\vert
$.
\begin{enumerate}
\item Montrer qu'il existe un nombre réel
$r\in\mathopen]0,1\mathclose[$ tel que, pour tout $\rho
\in\mathopen]0,r\mathclose[$, on ait $$m(\rho )>0.$$

Dans toute la suite, on suppose que $r$ et $\rho $ sont fixés et qu'ils
vérifient les conclusions de 1.

\item Montrer que, pour tout nombre complexe $w$ vérifiant $\vert
w\vert <m(\rho )$, la fonction
$$z\mapsto f(z)-w$$
a un seul zéro, noté $g(w)$, dans $D_\rho $.

\item Montrer que l'on a
$${1\over 2i\pi}\int_{\gamma _\rho }{zf'(z)\over f(z)-w}dz=g(w).$$

\item Montrer que, pour tout $w$ vérifiant $\vert w\vert <m(\rho )$,
on a $$g(w)=\sum_{n=0}^{+\infty}c_nw^n$$
et l'on exprimera les coefficients $c_n$ au moyen d'intégrales faisant
intervenir $f$ et $f'$.
\end{enumerate}
\finenonce{006858}



\finexercice


\exercice{6859, gijs, 2011/10/16}
\enonce{006859}{}
Soit $g$ une fonction méromorphe dans $\C$. On suppose que $g$ n'est
pas égale à la fonction constante $g(z)=i$ et que $g$ vérifie l'équation
différentielle
$$g'(z)=g^2(z)=-1.$$

On considère la fonction $f$ définie par 
$$f(z)={g(z)+i\over g(z)-i}.$$

On note $P$ l'ensemble des pôles de $f$ et $U=\C\setminus P$.

Quelle équation différentielle vérifie $f$ ?

Déterminer $f(z)$ sur $U$, puis sur $\C$. En déduire que, si l'on suppose
de plus que $g$ a un pôle en 0, l'on a
$$g(z)=\mathrm{cotan} z.$$
\finenonce{006859}



\finexercice       


\exercice{6860, gijs, 2011/10/16}
\enonce{006860}{}
On se propose de calculer, à l'aide du théorème
des résidus, la valeur de l'intégrale
 $$I=\int_{-\infty}^{+\infty}\exp{(-\pi t^2)}\ dt.$$
\begin{enumerate}
\item Soit $a\in {\Rr}^*$ ; en intégrant la fonction $g(z)=\exp{(-\pi z^2)}$
sur le rectangle de sommets $-R$, $R$, $R+ia$, $-R+ia$ ($R>0$), montrer
que
$$I=\int_{-\infty}^{+\infty}\exp{\left(-\pi(t+ia)^2\right)}\ dt.$$

\medskip

On pose : $f(z)=\exp{(i\pi z^2)}\tan{(\pi z)}$.

\item  Quels sont les pôles de $f$ ? Préciser leur ordre.

\medskip

Soit $R>0$ ; on considère le parallélogramme $\Gamma _R$ de sommets
$A=R+1+iR$, $B=R+iR$, $C=-R-iR$ et $D=-R+1-iR$, orienté dans le sens
direct.



\item Montrer que
$$\int_{\Gamma _R}f(z)\ dz=2e^{-i\pi/4}.$$

\item 
  \begin{enumerate}
  \item   Montrer que 
$$\forall t\in [0,1],\ \vert \tan{(\pi (R+t+iR))\vert \le \coth{(\pi R)}}.$$

  \item  En déduire que l'intégrale de $f$ sur le segment orienté $[A,B]$
tend vers 0 quand $R$ tend vers $+\infty$.

  \item  Montrer de même que l'intégrale de $f$ sur $[C,D]$ tend vers 0
quand $R$ tend vers $+\infty$.
  \end{enumerate}
\item  On note $J$ et $K$ les intégrales de $f$ sur les segments
orientés $[B,C]$ et $[D,A]$.
  \begin{enumerate}
  \item   Montrer que
$$J+K=(-1+i)\int_{-R}^R e^{-2\pi t^2}\left( e^{-2\pi t+2i\pi t}-1\right)\
dt.$$

  \item En déduire que, quand $R$ tend vers $+\infty$, $J+K$ tend vers
$$L=e^{3i\pi/4}\left(\int_{-\infty}^{+\infty} e^{-\pi (u^2+u\sqrt
2-iu\sqrt 2)}\ du-\int_{-\infty}^{+\infty} e^{-\pi u^2}\ du\right) .$$

  \item l'aide de la question 1. et d'un changement de variable,
vérifier que $L=2Ie^{-i\pi/4}$.
  \end{enumerate}
\end{enumerate}
\item Conclure.
\finenonce{006860}



\finexercice       


\exercice{6861, gijs, 2011/10/16}
\enonce{006861}{}
 Soit $c$ un point singulier essentiel d'une fonction
$f$ holomorphe dans un disque pointé $U =\{ z\in {\Cc}\vert\ 0<\vert
z-c\vert <\rho \}$. Le but de l'exercice est de démontrer que $f$ n'est
injective dans aucun voisinage pointé de $c$.

\begin{enumerate}  
\item Montrer que pour tout $\gamma \in{\Cc}$ et tout $\varepsilon >0$, il
existe $z'\in U$ et $\varepsilon '>0$ tels que
$$\overline{D (f(z'),\varepsilon ')}\subset f(U) \cap 
D(\gamma ,\varepsilon ),$$
où $D(a,r)$ désigne le disque ouvert de centre $a$ et de rayon $r$.
On pourra utiliser le théorème de Casorati-Weierstrass, puis remarquer que
$f(U)$ est ouvert (la fonction $f$ est holomorphe donc ouverte).

\item Pour $n\ge 1$, soit $U_n $ le disque pointé $\{ z\in {\Cc}\vert\ 0<\vert
z-c\vert <\rho /n \}$. Soient $\gamma _0\in{\Cc}$ et $\varepsilon
_0>0$. Construire par récurrence une suite strictement décroissante
$(\varepsilon _n)_{n\ge 1}$ de réels strictement positifs et une suite
$(z_n)_{n\ge 1}$, $z_n\in U_{n-1}$, vérifiant
\begin{eqnarray*}
\overline{D (f(z_1),\varepsilon_1 )}&\subset& f(U) \cap 
D(\gamma_0 ,\varepsilon _0),\\
\overline{D (f(z_{n+1}),\varepsilon_{n+1} )}&\subset& f(U_{n+1}) \cap 
D(f(z_n) ,\varepsilon _n)\ \ {\rm pour\ }n\ge 1.\\
\end{eqnarray*}
En déduire qu'il existe  $a\in D(\gamma _0,\varepsilon _0)$  et une suite
$(c_n)_{n\ge 0}$ de points de $U$ distincts deux à deux tels que
$$\lim_{n\to +\infty}c_n= c{\rm\ \ et\ \ }\forall n,\ \ f(c_n)=a.$$
Conclure.
\end{enumerate}
\finenonce{006861}



\finexercice       


\exercice{6862, gijs, 2011/10/16}
\enonce{006862}{}
Les parties I et II sont indépendantes.
Les parties III et IV utilisent les résultats établis dans II.

\subparagraph{I}
\begin{enumerate}
\item
 Déterminer l'ensemble des zéros de la fonction $\sin$ dans
$\C$. Quel est leur ordre de multiplicité ?

\item On note, pour tout $n\ge 0$, $\gamma_n^*$ le bord du carré
dans
$\C$ de sommets $(n+{1\over 2})\pi (1+i)$, $(n+{1\over 2})\pi (-1+i)$,
$(n+{1\over 2})\pi (-1-i)$, $(n+{1\over 2})\pi (1-i)$. Montrer que, si
$z\in \gamma_n^*$, on a $\vert \sin z\vert ^2\ge 1$.
\end{enumerate}


\subparagraph{II} 
Soit $f$ une fonction entière dans $\C$. On note
$f(z)=\sum_{n=0}^{\infty} a_nz^n$, pour tout $z$ dans $\C$.

\begin{enumerate}
\item Montrer que, pour tout $n\ge 1$ et tout $r>0$, on a
$$a_n={1\over \pi}\int_0^{2\pi}\Re\left( f(re^{i\theta}\right)
r^{-n}e^{-in\theta}d\theta .$$
On rappelle que $\Re f={1\over 2}(f+\overline f)$.

On note, pour $n\ge 1$, $\varphi_n$ un réel vérifiant $\vert
a_n\vert =a_ne^{i\varphi_n}$. En déduire que l'on a, pour tout $n\ge 1$,
$$\vert a_n\vert ={1\over \pi}\int_0^{2\pi}\Re\left(
f(re^{i(\theta +{1\over n}\varphi_n )}\right)
r^{-n}\left( 1+e^{-in\theta}\right) d\theta .$$

\item On suppose $f(0)=0$. Montrer que l'on a, pour tout $n\ge 1$ et
tout $r>0$,
$$\vert a_n\vert ={1\over \pi}\int_0^{2\pi}\Re\left(
f(re^{i(\theta +{1\over n}\varphi_n )}\right) r^{-n}e^{-in\theta}d\theta
$$
et donc
$$\vert a_n\vert r^n\le 2\sup_{\vert z\vert =r}\Re f(z).$$

\item On note toujours $f$ une fonction entière dans $\C$. On
suppose maintenant qu'il existe une suite $(r_j)_{j\ge 0}$ de réels
strictement positifs tendant vers $+\infty$ avec $j$ et qu'il existe des
constantes $A>0$ et $\beta >0$ telles que l'on ait, pour tout $j\ge 0$ et
tout $\theta$ réel,
$$\Re f(r_je^{i\theta})\le Ar_j^\beta.$$
En déduire que $f$ est un polynôme de degré au plus $\beta$ (on étudiera
d'abord le cas où $f(0)=0$).
\end{enumerate}

\subparagraph{III} 
Soit $g$ une fonction entière nulle seulement aux points
$k\pi$, $k\in {\Zz}$ et telle que chacun de ces zéros soit simple. On
suppose de plus qu'il existe $C>0$ tel que, pour tout $z$ dans $\C$, on
ait $\vert g(z)\vert \le \exp{(C\vert z\vert )}$.
\begin{enumerate}
\item Montrer que ${g\over \sin }$ est une fonction entière qui ne
s'annule pas dans $\C$. En déduire, en énonçant avec précision le
théorème du cours utilisé, qu'il existe une fonction $h$ entière telle
que, pour tout $z$ dans $\C$, on ait
$${g(z)\over \sin z}=\exp{h(z)}.$$

\item On note $C_n$ le chemin orienté positivement et défini, pour
tout $t\in [0,2\pi]$, par $C_n(t)=(n+{1\over 2})\pi e^{it}$ et $C_n^*$ son
image. Montrer que
$$\sup_{z\in C_n^*}\left\vert {g(z)\over \sin z}\right\vert \le
\sup_{z\in \gamma_n^*}\left\vert {g(z)\over \sin z}\right\vert $$
(où $\gamma_n^*$ a été défini en I.2. En déduire que l'on a, pour
tout $z$ de $C_n^*$,
$$\Re h(z)\le C\sqrt 2\left( n+{1\over 2}\right)\pi$$
et qu'il existe des nombres complexes $\lambda$ et $\mu$ tels que,
pour tout $z$ dans $\C$, on ait
$$g(z)=\lambda \sin z\exp{(\mu z)}.$$
\end{enumerate}

\subparagraph{IV} Soit $\alpha$ un nombre complexe non nul. On se propose de
montrer que l'équation
$$\alpha z-\exp z=0$$
a au moins une racine dans $\C$.
\begin{enumerate}
\item  On raisonne par l'absurde. Montrer qu'alors il existe une
fonction $\delta$ entière telle que l'on ait, pour tout $z$ dans $\C$,
$$\alpha z-\exp z=\exp {\delta (z)}.$$

\item  Etablir, pour tout $z$ dans $\C$,
$$\vert \alpha z-\exp z\vert \le (\vert \alpha\vert +1)\exp {\vert
z\vert }.$$
En déduire qu'il existe des nombres complexes $\rho$ et $\sigma$ tels
que, pour tout $z$ dans $\C$, on ait
$$\alpha z-\exp z=\rho\exp{(\sigma z)}.$$
Conclure.
\end{enumerate}
\finenonce{006862}



\finexercice       


\exercice{7576, mourougane, 2021/08/10}
\enonce{007576}{Série de Laurent}
Calculer la série de Laurent de $ 1 / (z^{2} +z^{3}) $ au voisinage de l'origine. 
\finenonce{007576}
\finexercice

\exercice{7584, mourougane, 2021/08/10}
\enonce{007584}{Questions de cours}
\begin{enumerate}
\item Rappeler la définition du rayon de convergence de la série entière $\sum a_n z^n$.
\item Soit $R$ le rayon de convergence de la série $\sum a_n z^n$.
Montrer que la série $\sum a_n z^n$ diverge sur $\Cc-\Delta_R(0)$.
\end{enumerate}
\finenonce{007584}
\finexercice


\exercice{7585, mourougane, 2021/08/10}
\enonce{007585}{Une équation}
Résoudre dans $\Cc$, l'équation $z^6+2z^3-3=0$.
\finenonce{007585}
\finexercice


\exercice{7586, mourougane, 2021/08/10}
\enonce{007586}{Applications holomorphes}
Déterminer toutes les applications holomorphes sur $\Cc$ dont la partie réelle est
$u(x,y)=x^2-y^2 +3x^2y-y^3.$
\finenonce{007586}
\finexercice


\exercice{7587, mourougane, 2021/08/10}
\enonce{007587}{Biholomorphismes}
Soit $\mathbb{H}:=\{z\in\Cc / Im(z)>0\}$.
 Soit $A =\begin{pmatrix} a&b\\c&d\end{pmatrix}$ une matrice de $SL(2,\Rr)$ (i.e. à coefficients réels et de déterminant $1$.)
Soit l'application linéaire fractionnaire
$\begin{array}{cccc}
 f_A :& z&\longmapsto& \frac{az+b}{cz+d}.
\end{array}$
 Montrer que $f_A$ est définie sur $\mathbb{H}$ et réalise un biholomorphisme de $\mathbb{H}$ dans lui-même.
\finenonce{007587}
\finexercice


\exercice{7588, mourougane, 2021/08/10}
\enonce{007588}{Séries entières}
\begin{enumerate}
\item Soit $\sum_{n\geq 1}a_nz^n$ une série entière convergente (avec un rayon de convergence strictement positif).
Déterminer le rayon de convergence de la série $\sum_{n\geq 1}\frac{a_n}{n!}z^n$.
\item Soit deux séries entières centrées en $0$ de rayon de convergence $R>0$ et de somme respective $f$ et $g$.
On suppose que pour tout $x\in]-R,R[, f(x)=g(x).$ Montrer que pour tout $z\in\Delta_R(0), f(z)=g(z)$. 
\end{enumerate}
\finenonce{007588}
\finexercice


\exercice{7589, mourougane, 2021/08/10}
\enonce{007589}{Questions de cours}
\begin{enumerate}
\item Rappeler la définition de la branche principale du logarithme.
\item Rappeler la définition d'une primitive d'une fonction continue $f$ sur $\Cc$.
\item Donner si possible l'exemple d'une fonction holomorphe sur $\Cc$ qui n'admet pas de primitive.
Sinon, appliquer un théorème pour montrer que toute fonction holomorphe sur $\Cc$ admet une primitive.
\end{enumerate}
\finenonce{007589}
\finexercice


\exercice{7590, mourougane, 2021/08/10}
\enonce{007590}{Une équation}
On fixe la branche principale du logarithme. Résoudre dans $\Cc^-$, l'équation $z^{i}=-1$.
\finenonce{007590}
\finexercice


\exercice{7591, mourougane, 2021/08/10}
\enonce{007591}{Applications holomorphes}
Développer $z\mapsto \frac{z^2+z-1}{z+1}$ en éléments simples et calculer $\int_{\partial \Delta_2(0)}\frac{z^2+z-1}{z+1}dz$.
\finenonce{007591}
\finexercice


\exercice{7592, mourougane, 2021/08/10}
\enonce{007592}{Questions de cours}
\begin{enumerate}
 \item L'application $\Cc-\{3\}\to\Cc$, $z\mapsto \frac{1}{z-3}$ admet-elle une primitive sur $\Cc-\{3\}$ ? Justifier.
 \item Donner l'exemple d'un ouvert non étoilé de $\Cc$. Justifier que l'exemple proposé n'est pas étoilé.
\end{enumerate}
\finenonce{007592}
\finexercice


\exercice{7593, mourougane, 2021/08/10}
\enonce{007593}{Calcul d'intégrales}
Calculer $$\int_{\partial \Delta_1} \frac{\sin(z)}{(z+2)^2} dz\ ,\ \int_{\partial \Delta_3} \frac{\sin(z)}{(z+2)^2} dz\quad \text{ et } \quad \int_{\partial \Delta_3} \frac{\sin(z)}{(z+1)^2(z+2)^2} dz.$$
\finenonce{007593}
\finexercice


\exercice{7594, mourougane, 2021/08/10}
\enonce{007594}{Singularités}
\begin{enumerate}
 \item Soit $f:\Cc^\star\to\Cc$ une application holomorphe. On suppose que $\lim_{z\to 0}|f(z)|=+\infty$.
 Quelle est la nature de la singularité isolée $0$ ? (apparente, polaire ou essentielle) Justifier.
 \item Donner l'exemple d'une application holomorphe avec une singularité essentielle. Justifier.
\end{enumerate}
\finenonce{007594}
\finexercice


\exercice{7595, mourougane, 2021/08/10}
\enonce{007595}{Applications holomorphes}
\begin{enumerate}
\item Rappeler la définition de l'indice $Ind_\Gamma (a)$ d'un point $a$ de $\Cc$ par rapport à un chemin fermé compact orienté $\Gamma$ de $\Cc$.
\item On suppose que $\Gamma$ est paramétré sous la forme $\gamma(t)=r(t)e^{i\theta(t)}$ où $r$ et $\theta$ sont deux fonctions de classe $\mathcal{C}^\infty$ par morceaux sur $[0,1]$,
$r$ à valeurs strictement positives et $r(0)=r(1)$, $\theta(0)\equiv\theta(1)\mod{2\pi}$
montrer que $$Ind_\Gamma (a)=\frac{1}{2\pi}\int_0^1 \theta'(t) dt$$ et correspond donc au nombre de tours, comptés positivement dans le sens direct, que fait $\Gamma$ autour de $a$.
\end{enumerate}
\finenonce{007595}
\finexercice


\exercice{7596, mourougane, 2021/08/10}
\enonce{007596}{Questions de cours}
\begin{enumerate}
\item Démontrer que la partie imaginaire d'une fonction holomorphe est harmonique.
\item Soit $f$ une fonction holomorphe sur $\Cc$ non identiquement nulle.
Montrer qu'il existe $\epsilon >0$ tel que $f$ ne s'annule pas sur $\Delta_\epsilon -\{0\}$.
\item Soient $\Omega$ un ouvert connexe et $f$ une fonction holomorphe sur $\Omega$. 
L'intégrale de $f$ le long de tout chemin fermé contenu dans $\Omega$ est-elle nulle ?
\end{enumerate}
\finenonce{007596}
\finexercice


\exercice{7597, mourougane, 2021/08/10}
\enonce{007597}{Deux équations}
\begin{enumerate}
\item Trouver toutes les solutions complexes de l'équation $\cosh(z)=4i$.

\item Trouver toutes les solutions complexes de l'équation $z^i=-1$.

\end{enumerate}
\finenonce{007597}
\finexercice


\exercice{7598, mourougane, 2021/08/10}
\enonce{007598}{Logarithmes}
\begin{enumerate}
\item Peut on définir une détermination de la fonction logarithme sur $ \mathbb{C} \backslash \{ \; x \in \Rr \; / \; x \geq 0 \;\} $ ? 
Dans le cas affirmatif, donner une définition de cette détermination. 


%\item Plus généralement peut on définir une détermination du logarithme sur le plan complexe privé d'une demi droite quelconque issue de l'origine ? 
\item M\^ eme question sur le plan complexe privé de l'ensemble $$ L=[ \; 0,1\; ] \cup \{ \; z \in \mathbb{C} \; / \; \mid z-2 \mid =1  \; {et} \; Im(z) \geq 0 \; \} \cup [\; 3, \infty \; [ .$$


$$\begin{tikzpicture}
\draw[help lines, step=1, very thin] (-6,0) grid (6,4);
\draw[shift={(3,0)}, thick] (2:0) arc (0:180:1cm);
\draw[thick] (0,0) -- (1,0) ;\draw[thick] (3,0) -- (6,0) ;
\draw (0,0) node[below]{$0$} ;
\draw (3,0) node[below]{$3$} ;
\draw (1,0) node[below]{$1$} ;
\end{tikzpicture}$$



\end{enumerate}
\finenonce{007598}
\finexercice


\exercice{7599, mourougane, 2021/08/10}
\enonce{007599}{Intégrales}
\begin{enumerate}
\item Paramétrer le cercle unité parcouru dans le sens trigonométrique.

\item Calculer à l'aide du paramétrage précédent $\int_{\partial\Delta} \frac{\cosh z}{z} dz$.

\item La fonction $\frac{\cosh(z)}{z}$ admet-elle une primitive sur $\Cc^\times$ ? Si oui, l'expliciter.

% \item Reprendre les questions précédentes avec la fonction $\frac{\sinh (z)}{z}$.
\end{enumerate}
\finenonce{007599}
\finexercice


\exercice{7600, mourougane, 2021/08/10}
\enonce{007600}{Intégrales}
\begin{enumerate}
 \item Enoncer la formule de Cauchy pour les disques, en précisant les hypothèses.
 
 \item Calculer $I_1:=\int_{\partial \Delta}\frac{\sin(z)}{z}dz.$
 
 \item Calculer $I_2:=\int_{\partial \Delta}\frac{\sin(z)}{z^2}dz.$ 
  
  \item Calculer $I_3:=\int_{\partial \Delta}\frac{\sin(z)}{z^3}dz.$ 
  
  \item Parmi les applications $\sin(z)$, $\frac{\sin(z)}{z}$, $\frac{\sin(z)}{z^2}$
  et $\frac{\sin(z)}{z^3}$ sur $\Cc^\times$, lesquelles ont une primitive holomorphe sur $\Cc^\times$ ?
\end{enumerate}
\finenonce{007600}
\finexercice


\exercice{7601, mourougane, 2021/08/10}
\enonce{007601}{Logarithmes}
\begin{enumerate}
\item Rappeler la définition de la branche principale du logarithme.

\item Peut on définir un logarithme sur $\Cc-C$ le plan complexe privé de l'ensemble
$$ C:= [ \; 0,1\; ] \cup \{ \; z \in \mathbb{C} \; / \; \mid z-2 \mid =1 \; \mathbb{H}ox{et} \; Im(z) \geq 0 \; \} \cup [\; 3, \infty \; [ ?$$
On appelle $U:=\{z\in\Cc, \mid z-2 \mid < 1, Im(z) \geq 0\}$ la partie grisée,
et $V :=\{z\in\Cc, z\not\in C, z\not\in U\}$ le reste de $\Cc-C$. 
$$\begin{tikzpicture}
\draw[white, fill=gray!10] (-6,-2) rectangle (6,4);
\draw[shift={(3,0)}, thick, fill=gray!100] (2:0) arc (0:180:1cm);
\draw (2,1) node[below]{$U$} ;
\draw[thick] (0,0) -- (1,0) ;
\draw[thick] (3,0) -- (6,0) ;
\draw (0,0) node[below]{$0$} ;
\draw (-2,2) node[below]{$V$} ;
\draw (3,0) node[below]{$3$} ;
\draw (5.5,0) node[above]{$C$} ;
\draw (1,0) node[below]{$1$} ;
\draw (-4,4) node[below]{$\Cc-C=U\cup V$} ;
\end{tikzpicture}$$
On précisera sur $U$ et sur $V$ l'argument choisi dans la définition du logarithme.


\end{enumerate}
\finenonce{007601}
\finexercice


\exercice{7602, mourougane, 2021/08/10}
\enonce{007602}{Application entière dans $\mathbb{H}$}
\begin{enumerate}
\item Soit $f :\Cc\to\Cc$ une application holomorphe non constante. Montrer que l'image du plan complexe par $f$ 
rencontre tous les disques ouverts non vides $\Delta_r(a)$ de $\mathbb{C}$.

\item En déduire que toute application holomorphe de $\Cc$ dans $\mathbb{H}$ est constante.

\item On considère $$\begin{array}{cccc}
 h :& \Cc-\{-i\}&\longrightarrow&\Cc \\&z&\longmapsto& \frac{z-i}{z+i}.
\end{array}$$
Montrer que 
$
 \forall z\in\Cc-\{-i\}, \ \ \ 1-|h(z)|^2=\frac{4 Im(z)}{|z+i|^2}.
$

\item En déduire que l'image du demi-plan de Poincaré $\mathbb{H}$ par $h$ est une partie bornée de $\Cc$.

\item En déduire par une nouvelle démonstration que toute application holomorphe de $\Cc$ dans $\mathbb{H}$ est constante.
\end{enumerate}
\finenonce{007602}
\finexercice


\exercice{7603, mourougane, 2021/08/10}
\enonce{007603}{Primitive et résidus}
\begin{enumerate}
 \item Soit $D$ un ouvert de $\Cc$, $c$ un point de $D$ et $f : D-\{c\}\to \Cc$ une application holomorphe.
Rappeler la définition du résidu de $f$ en $c$.

 \item Soit $D$ un ouvert de $\Cc$ et $f : D-\{c\}\to \Cc$ une application holomorphe.
 On suppose que $f$ admet une primitive sur $D-\{c\}$.
 Que peut-on dire du résidu de $f$ en $c$ ? 
 
  \item Soit $c$ un point de $\Delta$ et $f :\Delta-\{c\}\to\Cc$ 
  une application holomorphe qui admet un pôle d'ordre $3$ en $c$.
  Sous quelle condition $f$ admet-elle une primitive sur $\Delta-\{c\}$ ?  
  
\end{enumerate}
\finenonce{007603}
\finexercice


\exercice{7604, mourougane, 2021/08/10}
\enonce{007604}{Résidus et représentation}
\begin{enumerate}
\item \'Enoncer le théorème des résidus.

\item \`A l'aide du théorème des résidus, retrouver la formule de représentation de Cauchy pour les disques :
Soit $D$ un ouvert de $\Cc$ et $f : D\to \Cc$ une application holomorphe.
Alors, pour tout disque $\Delta_r(a)$ dont l'adhérence est incluse dans $D$,
$$\forall b\in\Delta_r(a), \ \ \ \frac{1}{2i\pi}\int_{\partial\Delta_r(a)}\frac{f(z)dz}{z-b}=f(b).$$

\item Montrer plus généralement : Soit $D$ un ouvert de $\Cc$ et $f : D\to \Cc$ une application holomorphe.
Alors, pour tout disque $\Delta_r(a)$ dont l'adhérence est incluse dans $D$,
$$\forall k\in\Nn,\ \ \ \forall b\in\Delta_r(a), \ \ \ \
\frac{1}{2i\pi}\int_{\partial\Delta_r(a)}\frac{f(z)dz}{(z-b)^{k+1}}=\frac{f^{(k)}(b)}{k!}.$$


\end{enumerate}
\finenonce{007604}
\finexercice


\exercice{7605, mourougane, 2021/08/10}
\enonce{007605}{Racine de polynômes}
\begin{enumerate}
\item \'Enoncer le théorème de Rouché.

\item Montrer que les zéros du polynôme $ p(z) = z^{4}-7z-1 $ 
sont tous inclus dans le disque $\Delta_2(0)$ centré en l'origine de rayon 2.
On vérifiera soigneusement toutes les hypothèses du théorème utilisé.

\end{enumerate}
\finenonce{007605}
\finexercice


\exercice{7606, mourougane, 2021/08/10}
\enonce{007606}{Questions de cours}
\begin{enumerate}
\item Donner l'exemple d'une fonction $f :\Cc\to\Cc$ indéfiniment dérivable au sens réel, mais non holomorphe. On justifiera ces deux propriétés.

 \item L'application $\Cc-\{3\}\to\Cc$, $z\mapsto \frac{1}{z-3}$ admet-elle une primitive sur $\Cc-\{3\}$ ? Justifier. 

 \item Donner l'exemple d'un ouvert étoilé de $\Cc$. On précisera le point par rapport auquel l'ouvert est étoilé.

 \item Donner l'exemple d'un ouvert non étoilé de $\Cc$. On justifiera que l'exemple proposé n'est pas étoilé.
 
 
\end{enumerate}
\finenonce{007606}
\finexercice


\exercice{7607, mourougane, 2021/08/10}
\enonce{007607}{Calcul d'intégrales}
On paramètre le cercle $C_r$ privé du point ${-r}$ de centre $0$ et de rayon $r>0$ orienté dans le sens trigonométrique 
du plan complexe en définissant pour $t\in\Rr$, $\xi(t)$ 
comme le point d'intersection de la droite d'équation $y=t(x+r)$ avec le cercle $C_r$ différent de $-r$.
\begin{enumerate}
\item Montrer que $\xi(t)=r\frac{1+it}{1-it}$.


\item Vérifier que $\xi$ est dérivable sur $\Rr$ et calculer $\frac{\xi'(t)}{\xi(t)}$.

\item En déduire que 
$$\int_{C_r} \frac{dz}{z}=2i\int_{-\infty}^{+\infty}\frac{dt}{1+t^2}.$$

\item En déduire la valeur de $\int_{-\infty}^{+\infty}\frac{dt}{1+t^2}$.

\end{enumerate}
\finenonce{007607}
\finexercice


\exercice{7608, mourougane, 2021/08/10}
\enonce{007608}{Biholomorphisme de $\mathbb{H}$}
On rappelle qu'à toute matrice $A=\begin{pmatrix} a&b\\c&d\end{pmatrix}$ de $SL(2,\Rr)$,
on associe l'application linéaire fractionnaire
$$\begin{array}{cccc}
 h_A :& \Cc-\{-\frac{d}{c}\}&\longrightarrow&\Cc-\{\frac{a}{c}\}\\&z&\longmapsto& \frac{az+b}{cz+d}
\end{array}$$
\begin{enumerate}

\item Montrer que $h_A$ envoie $\mathbb{H}$ sur $\mathbb{H}$.

\item Montrer que pour tout élément $z$ de $\mathbb{H}$, il existe $A\in SL(2,\Rr)$ 
tel que $h_A(i)=z$.
                                            
\end{enumerate}
\finenonce{007608}
\finexercice


\exercice{7609, mourougane, 2021/08/10}
\enonce{007609}{Représentation}
\begin{enumerate}
\item \'Enoncer la formule de représentation de Cauchy pour les disques.

% \item \'Enoncer le théorème des résidus.

\item \`A l'aide de la formule de représentation de Cauchy, montrer que si une fonction holomorphe sur $\Delta$
ne prend que des valeurs réelles sur le cercle $\partial\Delta_{\frac{1}{2}}$,
alors $f$ prend une valeur réelle en $0$.

\item \`A l'aide du théorème de représentation de Cauchy et sans le théorème des zéros isolés, montrer que si une fonction holomorphe sur $\Delta$
est constante sur le cercle $\partial\Delta_{\frac{1}{2}}$,
alors $f$ est constante sur $\overline{\Delta_{\frac{1}{2}}}$.

\end{enumerate}
\finenonce{007609}
\finexercice


\exercice{7610, mourougane, 2021/08/10}
\enonce{007610}{Croissance à l'infini}
On rappelle la formule de Gutzmer : soit $f$ la somme de la série entière $\sum_{n\geq 0} a_n z^n$ de rayon de convergence $R$.
Alors, pour tout $r<R$,
$$\sum_{n\geq 0} |a_n|^2r^{2n}=\frac{1}{2\pi}\int_0^{2\pi}|f(re^{i\theta})|^2d\theta.$$
\begin{enumerate}
\item Démontrer à l'aide de la formule de Gutzmer que toute application holomorphe $f:\Cc\to\Cc$ bornée est constante..

\item Soit $f:\Cc\to\Cc$ application holomorphe. On suppose que 
$$\forall r\in ]0,+\infty[, M(r):=\sup_{|z|<r} |f(z)|\leq r.$$
Montrer que $f$ est une application affine.
\end{enumerate}
\finenonce{007610}
\finexercice


\exercice{7611, mourougane, 2021/08/10}
\enonce{007611}{Séries entières}
\begin{enumerate}
\item Soit $\sum_{n\geq 1}a_nz^n$ une série entière convergente (avec un rayon de convergence strictement positif).
Déterminer le rayon de convergence de la série $\sum_{n\geq 1}\frac{a_n}{n!}z^n$.

\item Soit deux séries entières centrées en $0$ de rayon de convergence $R>0$ et de somme respective $f$ et $g$.
On suppose que pour tout $x\in]-R,R[, f(x)=g(x).$ Montrer que pour tout $z\in\Delta_r, f(z)=g(z)$. 

\item En déduire que pour tout $a\in \Rr$, et pour tout $z\in\Cc$, $\sin(a+z)=\sin(a)\cos (z)+\cos(a)\sin (z).$

\end{enumerate}
\finenonce{007611}
\finexercice


\exercice{7612, mourougane, 2021/08/10}
\enonce{007612}{\'Equations}
\begin{enumerate}
\item Soit $c\in \Cc$. Montrer que si $z\in\Cc$, alors $\sin z=c\iff (e^{iz})^2-2ic e^{iz}-1=0$.

\item Soit $c\in [-1,1]$. Montrer que toutes les solutions dans $\Cc$ de $\sin z=c$ sont réelles.

\item Soit $a$ et $b$ deux nombres complexes. Calculer $e^{i(a+b)}$ et $e^{-i(a+b)}$ en fonction de $\sin a$, $\sin b$, $\cos a$ et $\cos b$.

\item Soit $a$ et $b$ deux nombres complexes. Démontrer la formule pour $\sin (a+b)$ en fonction de $\sin a$, $\sin b$, $\cos a$ et $\cos b$.

\item Résoudre dans $\Cc$, l'équation $\cos z+\sin z=2$.
\end{enumerate}
\finenonce{007612}
\finexercice


\exercice{7613, mourougane, 2021/08/10}
\enonce{007613}{Biholomorphisme de $\mathbb{H}$}
\textit{On rappelle qu'à toute matrice $A=\begin{pmatrix} a&b\\c&d\end{pmatrix}$ de $SL(2,\Rr)$,
on associe l'application linéaire fractionnaire
$$\begin{array}{cccc}
 h_A :& \Cc-\{-\frac{d}{c}\}&\longrightarrow&\Cc-\{\frac{a}{c}\}\\&z&\longmapsto& \frac{az+b}{cz+d}
\end{array}$$}
\begin{enumerate}
\item Montrer que $h_A$ envoie $\mathbb{H}$ sur $\mathbb{H}$.

\item Montrer que pour tout élément $z$ de $\mathbb{H}$, il existe $A\in SL(2,\Rr)$ 
tel que $h_A(i)=z$.

\end{enumerate}
\finenonce{007613}
\finexercice


\exercice{7614, mourougane, 2021/08/10}
\enonce{007614}{Logarithmes}
\begin{enumerate}
\item Déterminer et représenter l'ensemble $D:=\{z\in\Cc, z^3\in \Cc^-\}$.

\item Soit $f : D\to \Cc$, $z\mapsto \log (z^3)$. 
En quels points de $\Cc-D$ peut-on prolonger $f$ par continuité ?

\item Comparer les fonctions $f$ et $3\log$ sur l'intersection $D\cap \Cc^-$ de leur domaine de définition.
\end{enumerate}
\finenonce{007614}
\finexercice


\exercice{7615, mourougane, 2021/08/10}
\enonce{007615}{À partir des définitions}
\begin{enumerate}
\item Rappeler la définition de l'holomorphie d'une fonction $f:\Cc\to\Cc$ en un point $a$ de $\Cc$.
\item Soit $f : \Cc \to \Cc$ une application holomorphe.
Déterminer si
l'application $g : \Cc \to \Cc, z\mapsto \overline{f(\overline{z})}$ est holomorphe sur $\Cc$ ou pas. Si oui, déterminer sa $\Cc$-dérivée en fonction de celle de $f$.
\end{enumerate}
\finenonce{007615}
\finexercice


\exercice{7616, mourougane, 2021/08/10}
\enonce{007616}{Séries entières}
\begin{enumerate}
\item Rappeler la définition du rayon de convergence de la série entière $\sum a_n z^n$.
\item Soit $R$ le rayon de convergence de la série entière $\sum a_n z^n$.
Déterminer le rayon de convergence de la série $\sum (1+\frac{1}{n+1})a_n z^n$.
\end{enumerate}
\finenonce{007616}
\finexercice


\exercice{7617, mourougane, 2021/08/10}
\enonce{007617}{Applications holomorphes}
On rappelle les équations de Cauchy-Riemann en coordonnées polaires $(r,\theta)$ pour une fonction $f=u+iv$ d'un ouvert $D$ de $\Cc$ dans $\Cc$~:
\begin{center}
 $\displaystyle\frac{\partial u}{\partial r}=\frac{1}{r}\frac{\partial v}{\partial \theta}
\quad \text{ et } \quad \frac{1}{r}\frac{\partial u }{\partial \theta}=-\frac{\partial v}{\partial r}$.
\end{center}
\begin{enumerate}
\item L'application $u: \Rr^2-\{(0,0)\}\to\Rr, (x,y)\mapsto \log_\Rr(x^2+y^2)$ est-elle harmonique sur son domaine de définition ?
\item En utilisant les équations de Cauchy-Riemann en coordonnées polaires, déterminer s'il existe ou pas une fonction holomorphe $f$ sur $\Cc-\{0\}$ dont la partie réelle est $u$.
\end{enumerate}
\finenonce{007617}
\finexercice


\exercice{7618, mourougane, 2021/08/10}
\enonce{007618}{Biholomorphismes}
Soit $a$ un nombre complexe de module strictement inférieur à $1$.
Montrer que l'application
$\begin{array}{cccc}
 f_a :& z&\longmapsto& \frac{z-a}{\overline{a}z-1}
\end{array}$
 réalise un biholomorphisme du disque $\Delta$ sur lui-même.
 \finenonce{007618}
\finexercice


\exercice{7619, mourougane, 2021/08/10}
\enonce{007619}{Séries entières}
\begin{enumerate}
\item Soit $\sum_{n\geq 1}a_nz^n$ une série entière convergente (avec un rayon de convergence strictement positif).
Déterminer le rayon de convergence de la série $\sum_{n\geq 1}\frac{a_n}{n!}z^n$.
\item Soit deux séries entières centrées en $0$ de rayon de convergence $R>0$ et de somme respective $f$ et $g$.
On suppose que pour tout $x\in]-R,R[, f(x)=g(x).$ Montrer que pour tout $z\in\Delta_R(0), f(z)=g(z)$. 
\end{enumerate}
\finenonce{007619}
\finexercice


\exercice{7620, mourougane, 2021/08/10}
\enonce{007620}{Questions de cours}
\begin{enumerate}
\item Rappeler la définition de la branche principale du logarithme.
\item Rappeler la définition d'une primitive d'une fonction continue $f$ sur $\Cc$.
\item Donner si possible l'exemple d'une fonction holomorphe sur $\Cc$ qui n'admet pas de primitive.
Sinon, appliquer un théorème pour montrer que toute fonction holomorphe sur $\Cc$ admet une primitive.
\end{enumerate}
\finenonce{007620}
\finexercice


\exercice{7621, mourougane, 2021/08/10}
\enonce{007621}{Une équation}
On fixe la branche principale du logarithme. Résoudre dans $\Cc^-$, l'équation $z^{i}=-1$.
\finenonce{007621}
\finexercice


\exercice{7622, mourougane, 2021/08/10}
\enonce{007622}{Applications holomorphes}
Développer $z\mapsto \frac{z^2+z-1}{z+1}$ en éléments simples et calculer $\int_{\partial \Delta_2(0)}\frac{z^2+z-1}{z+1}dz$.
\finenonce{007622}
\finexercice


\exercice{7623, mourougane, 2021/08/10}
\enonce{007623}{Questions de cours}
\begin{enumerate}
 \item L'application $\Cc-\{3\}\to\Cc$, $z\mapsto \frac{1}{z-3}$ admet-elle une primitive sur $\Cc-\{3\}$ ? Justifier.
 \item Donner l'exemple d'un ouvert non étoilé de $\Cc$. Justifier que l'exemple proposé n'est pas étoilé.
\end{enumerate}
\finenonce{007623}
\finexercice


\exercice{7624, mourougane, 2021/08/10}
\enonce{007624}{Calcul d'intégrales}
Calculer $$\int_{\partial \Delta_1} \frac{\sin(z)}{(z+2)^2} dz\ ,\ \int_{\partial \Delta_3} \frac{\sin(z)}{(z+2)^2} dz\quad \text{ et } \quad \int_{\partial \Delta_3} \frac{\sin(z)}{(z+1)^2(z+2)^2} dz.$$
\finenonce{007624}
\finexercice


\exercice{7625, mourougane, 2021/08/10}
\enonce{007625}{Singularités}
\begin{enumerate}
 \item Soit $f:\Cc^\star\to\Cc$ une application holomorphe. On suppose que $\lim_{z\to 0}|f(z)|=+\infty$.
 Quelle est la nature de la singularité isolée $0$ ? (apparente, polaire ou essentielle) Justifier.
 \item Donner l'exemple d'une application holomorphe avec une singularité essentielle. Justifier.
\end{enumerate}
\finenonce{007625}
\finexercice


\exercice{7626, mourougane, 2021/08/10}
\enonce{007626}{Applications holomorphes}
\begin{enumerate}
\item Rappeler la définition de l'indice $Ind_\Gamma (a)$ d'un point $a$ de $\Cc$ par rapport à un chemin fermé compact orienté $\Gamma$ de $\Cc$.
\item On suppose que $\Gamma$ est paramétré sous la forme $\gamma(t)=r(t)e^{i\theta(t)}$ où $r$ et $\theta$ sont deux fonctions de classe $\mathcal{C}^\infty$ par morceaux sur $[0,1]$,
$r$ à valeurs strictement positives et $r(0)=r(1)$, $\theta(0)\equiv\theta(1)\mod{2\pi}$
montrer que $$Ind_\Gamma (a)=\frac{1}{2\pi}\int_0^1 \theta'(t) dt$$ et correspond donc au nombre de tours, comptés positivement dans le sens direct, que fait $\Gamma$ autour de $a$.
\end{enumerate}
\finenonce{007626}
\finexercice


\exercice{7627, mourougane, 2021/08/10}
\enonce{007627}{Questions de cours}
\begin{enumerate}
\item Rappeler la définition de la branche principale du logarithme.

\item Rappeler la définition d'une primitive d'une fonction continue $f$ sur $\Cc$.

\item Donner si possible l'exemple d'une fonction continue non holomorphe $f$ sur $\Cc$ qui admet une primitive.

\item Donner si possible l'exemple d'une fonction holomorphe non constante de $\Delta$ dans $\Cc$. Dans le cas contraire, démontrer la non-existence.

\item Donner si possible l'exemple d'une fonction holomorphe non constante de $\Cc$ dans $\Delta$. Dans le cas contraire, démontrer la non-existence.

\item Donner si possible l'exemple d'une fonction holomorphe non constante de $\mathbb{H}$ dans $\Delta$. Dans le cas contraire, démontrer la non-existence.
\end{enumerate}
\finenonce{007627}
\finexercice


\exercice{7628, mourougane, 2021/08/10}
\enonce{007628}{Calcul d'intégrales}
 Calculer $$\int_{\partial \Delta_2(1)}\frac{z^2+z-1}{z^2(z^2-4)}dz.$$
      


\finenonce{007628}
\finexercice


\exercice{7629, mourougane, 2021/08/10}
\enonce{007629}{Applications entières}
Soit $f :\Cc\to\Cc$ une application holomorphe non constante.
\begin{enumerate}
 \item Montrer que le point $0$ est dans l'adhérence de l'image de $f$.

 \item Déterminer l'adhérence de l'image de $f$.
\end{enumerate}
\finenonce{007629}
\finexercice 


\exercice{7630, mourougane, 2021/08/10}
\enonce{007630}{Singularités}
 Soit $D=\Cc^\star$ et $f : D\to \Cc$ l'application définie par $f(z)=\exp(\frac{1}{z})-\frac{1}{z}$.
\begin{enumerate}
  \item Déterminer la nature de la singularité de $f$ en $0$ (apparente, polaire ou essentielle)

  \item L'application admet-elle une primitive sur $D$ ?

  \item Calculer $\int_{\partial\Delta}\exp(\frac{1}{z})dz$.
\end{enumerate}
\finenonce{007630}
\finexercice


\exercice{7631, mourougane, 2021/08/10}
\enonce{007631}{Applications sur le disque}
Soit $f$ une fonction holomorphe sur le disque unité $\Delta$. 
Le but de l'exercice est de montrer que $|f(z)|$ ne tend pas vers $+\infty$ quand $|z|$ tend vers $1$.
\begin{enumerate}
\item Vérifier l'énoncé pour l'application $\Delta\to\Cc$, $z\mapsto \frac{1}{z^6-1}$.
 
\item Supposons d'abord que $f$ n'a pas de zéro dans $\Delta$. 
Soit $r_n$ une suite de réels de $[0,1[$, tels que pour tout entier naturel non nul $n$ et tout $z$ de $\Delta$ de module supérieur à $r_n$, $f(z)$ est de module supérieur à $n$. 
Montrer alors que pour tout $n$, $|f(0)|\geq n$ et conclure.
 
\item Dans le cas général, supposons que $|f(z)|$ tende vers $+\infty$ quand $|z|$ tend vers $1$.
Montrer qu'on peut écrire $f$ comme produit sur $\Delta$ d'un polynôme et d'une application holomorphe $g$ qui n'a pas de zéro dans $\Delta$. Conclure.

\end{enumerate}
\finenonce{007631}
\finexercice


\exercice{7632, mourougane, 2021/08/10}
\enonce{007632}{Questions de cours}
\begin{enumerate}
\item Calculer les intégrales $\int_{\partial \Delta}\frac{dz}{z}$ et $\int_{\partial \Delta}\frac{dz}{z-2}$ à l'aide d'un paramétrage et d'un développement en séries entières, sans utiliser les théorèmes généraux.
\item \'Enoncer le lemme de Goursat.
\item Soit $D$ un ouvert de $\Cc$ et $f : D\to \Cc$ une application holomorphe. Donner la définition d'un logarithme de $f$.
\item Soit $D$ un ouvert de $\Cc$ et $f : D\to \Cc$ une application holomorphe qui admet un logarithme $g$ sur $D$. Donner l'exemple d'un autre logarithme de $f$. Montrer que $f$ n'a pas de zéro dans $D$ et calculer la dérivée complexe de $g$.
\item Soit $D$ un ouvert de $\Cc$ et $a$ un point de $D$. Soit $f : D-\{a\}\to \Cc$ une application holomorphe. Définir le résidu de $f$ en $a$.
\item Définir l'indice d'un point par rapport à un chemin fermé orienté dans $\Cc$.
\end{enumerate}
\finenonce{007632}
\finexercice


\exercice{7633, mourougane, 2021/08/10}
\enonce{007633}{Applications entières}
Existe-t-il une application $f :\Cc\to\Cc$ holomorphe non constante et bi-périodique de périodes $1$ et $i$ i.e. 
$$\forall z\in\Cc,\ \ f(z+1)=f(z+i)=f(z).$$
Si oui, donner un exemple. Si non, démontrer la non-existence d'une telle application.
\finenonce{007633}
\finexercice


\exercice{7634, mourougane, 2021/08/10}
\enonce{007634}{Applications proportionnelles}
Soit $D$ un ouvert connexe de $\Cc$ contenant le disque unité fermé.
Soient $f$ et $g$ deux applications holomorphes sur $D$ telles que pour tout $z\in\partial\Delta$, $|f(z)|=|g(z)|$. 
\begin{enumerate}
\item On suppose que $f$ et $g$ n'ont pas de zéro dans $D$. Montrer qu'il existe un nombre complexe $\lambda$ de module $1$ tel que $f=\lambda g$ sur $D$. 

\item La conclusion est-elle encore vraie si on ne suppose plus que $f$ et $g$ n'ont pas de zéro dans $D$ ?
\end{enumerate}
\finenonce{007634}
\finexercice


\exercice{7635, mourougane, 2021/08/10}
\enonce{007635}{Singularités} 
On considère le contour $\Gamma$ et les points $w_1$, $w_2$, $w_3$ comme sur la figure ci-dessous
 \begin{center} 
 \begin{tikzpicture}
 %\draw [step=1cm, gray, very thin](-2,-4) grid (8,4);
 %\draw (-4,0) -- (4,0); \draw (0,-4) -- (0,4);
 \draw (5,-1) node {$\bullet$}; \draw (5,-1) node[below left] {$w_1$};
 \draw (2,1) node {$\bullet$};\draw (2,1) node [below right] {$w_2$};
 \draw (-1,-1) node {$\bullet$};\draw (-1,-1) node [below left] {$w_3$};
%\draw (1,0) node {$\bullet$}; \draw (1,0) node [below right] {$C$} circle (2) ;
\draw (2.2,2) node[below right] {${\Gamma}$};
%\draw [>=latex] (0,0) -- (1,1) arc (180:0:1) -- (4,0);
\draw[thick, draw=black](0,0) -- plot [domain=0:2*pi] (\x, {2*sin(\x r)}) -- cycle;
\draw [thick] (0.9,0.1)--(1,0)--(0.9,-0.1);
\draw [thick] (4.9,0.1)--(5,0)--(4.9,-0.1);
\draw [thick] (0.7,1.1)--(0.5,1)--(0.5,1.2);
\draw [thick] (3.85,-1.1)--(3.65,-1)--(3.65,-1.2);
\end{tikzpicture}
\end{center}
 Exprimer la valeur des intégrales suivantes à l'aide des nombres complexes $w_1,w_2,w_3$ :

 \begin{enumerate}
 \item $A=\int_\Gamma \frac{dz}{(z-w_1)(z-w_2)(z-w_3)}$.

 \item $B=\int_\Gamma \sin(z)dz$.
 
 \item $C=\int_\Gamma \frac{\sin(z) }{(z-w_1)^2}dz$.
 
\end{enumerate}
\finenonce{007635}
\finexercice


\exercice{7636, mourougane, 2021/08/10}
\enonce{007636}{Localisations de racines}
On considère le polynôme $P(z)=z^4+6z+3$.
\begin{enumerate}
 \item Démontrer que $P$ a ses quatre racines dans le disque $\Delta_2$.
 
 \item Démontrer que $P$ n'admet qu'une racine dans $\Delta$.

 \item Démontrer que $P$ n'admet pas de racine dans $\Delta_{\frac{1}{3}}$.

 \item Soit $a$ la racine de $P$ dans le disque $\Delta$. Démontrer que
$$ a=\frac{1}{2i\pi}\int_{\partial\Delta} \frac{4z^3+6}{z^4+6z+3}zdz.$$

\end{enumerate}
\finenonce{007636}
\finexercice


\exercice{7637, mourougane, 2021/08/10}
\enonce{007637}{Questions de cours}
\begin{enumerate}
\item Une application holomorphe sur $\Cc$ de dérivée partout non nulle est-elle nécessairement injective ?
\item Rappeler la définition de la branche (ou détermination) principale du logarithme.
\item Ayant fixé la détermination principale $\mathbb{H}ox{Log} z$ du logarithme, définir les fonctions puissances $z\mapsto z^\alpha$ pour tout nombre complexe $\alpha$ en précisant leur domaine de définition.
\item Montrer à l'aide d'un théorème du cours qu'une application holomorphe sur un ouvert connexe $D$ de $\Cc$ dont toutes les valeurs sont de module $1$ est constante.
\item Définir une singularité essentielle pour une application holomorphe définie sur un ouvert de $\Cc$.
\item Donner un exemple d'automorphisme du demi-plan de Poincaré qui n'est pas une translation.
\end{enumerate}
\finenonce{007637}
\finexercice


\exercice{7638, mourougane, 2021/08/10}
\enonce{007638}{Logarithme}
Soit $\mathbb{H}ox{Log} z$ la détermination principale du logarithme dans 
$\Cc^-$.
\begin{enumerate}
\item On considère $z=e^{2i\pi\over3}$; les nombres complexes $\mathbb{H}ox{Log} (z^2)$ et $2\mathbb{H}ox{Log}
z$ sont-ils égaux ?

\item On considère $z=e^{3i\pi\over4}$; les nombres complexes $z^{2i}$, $(z^2)^i$ et $(z^i)^2$ sont-ils égaux ?
\end{enumerate}
 \finenonce{007638}
\finexercice


\exercice{7639, mourougane, 2021/08/10}
\enonce{007639}{Laplacien}
On rappelle que le Laplacien $\Delta\varphi$ d'une fonction $\varphi~: U\to\Rr$ définie sur un ouvert $U$ de $\Rr^2$ est la fonction sur $U$ donnée par $\Delta\varphi (x,y)=\frac{\partial^2}{\partial x^2}\varphi (x,y)+\frac{\partial^2}{\partial y^2}\varphi(x,y)$.
 \begin{enumerate}
 \item Soit $D$ un ouvert connexe de $\Cc$ et $f : D\to\Cc$ une application holomorphe.
 Exprimer $\Delta |f|^2$ à l'aide de $f'$.
 

 \item Soit $D$ un ouvert connexe et $(f_i : D\to\Cc)_{i=1}^N$ une famille finie d'applications holomorphes. On suppose que 
 $$\forall z\in D, \ \ \sum_{i=1}^N |f_i(z)|^2=1.$$
 Montrer que toutes les $f_i$ sont constantes sur $D$.
\end{enumerate}
\finenonce{007639}
\finexercice


\exercice{7640, mourougane, 2021/08/10}
\enonce{007640}{Application bornée}
Montrer que pour tout nombre complexe $z$ dans le disque unité $\Delta$, $$\left|\frac{4z+3}{4+3z}\right|\leq 1.$$
\finenonce{007640}
\finexercice


\exercice{7641, mourougane, 2021/08/10}
\enonce{007641}{Singularités}
 Que vaut, en fonction du nombre réel $r>0$, l'intégrale
 $$I_r:=\int_{\Delta_r} \frac{dz}{2z^2-5z+2} \ \ ?$$
 On précisera les valeurs de $r$ exclues.
\finenonce{007641}
\finexercice

\section{450.00 Interpolation polynomiale}

\section{451.00 Courbe de Bézier, spline}

\section{452.00 Intégration numérique}

\section{453.00 Méthode de Newton}

\section{454.00 Résolution d'équation différentielle}

\section{455.00 Résolution de systèmes linéaires : méthode directe}

\section{456.00 Résolution de systèmes linéaires : méthode itérative}

\section{457.00 Résolution de systèmes linéaires : méthode de gradient}

\section{458.00 Calcul de valeurs propres et de vecteurs propres}

\section{459.00 Autre}


\exercice{2210, matos, 2008/04/23}
\enonce{002210}{Matrices triangulaires \'el\'ementaires}
Soit $n \in \Nn$ et on d\'efinit les matrices suivantes dans $\Rr^{n\times n}$:
\begin{itemize}
\item $E_{ij}$ matrice avec un $1$ dans la position $(i,j)$ et $0$ partout ailleurs;
\item $V_{ij}(\lambda )=I +  \lambda E_{ij}$ , $\lambda \in\Rr , i>j$;
\item $L(l_i) = I+l_ie_i^T$, $l_i\in\Rr^n$ tel que ses premi\`eres $i$ composantes sont nulles.
\end{itemize}
\begin{enumerate}
\item Quels sont les r\'esultats des op\'erations suivantes sur la matrice $A$:
     $$B=V_{ij}(\lambda )A  , C=AV_{ij}(\lambda ) ?$$
\item Quelle est la forme de la matrice 
$$V_{ij}(\lambda )V_{kj}(\lambda ') ,  k>i ?$$
\item Repr\'esenter $L(l_i)$ et montrer que $L(l_i)^{-1}=L(-l_i)$.
\item D\'ecomposer $L(l_i)$ comme produit de matrices de la forme $V_{km}(\lambda )$.
\item Calculer $L=\prod_{i=1}^{n-1}L(l_i)$ et son inverse $L^{-1}$
\item On suppose les $l_i$ stock\'es dans un tableau bidimensionnel $Z$ et $b\in\Rr^n$ stock\'e dans un tableau unidimensionnel $B$. Donner un algorithme permettant de calculer dans $B$ la solution de $Lx=b$:
\begin{enumerate}
\item en utilisant l'expression de $L^{-1}$;
\item en r\'esolvant le syst\`eme triangulaire.
\end{enumerate}
Quelle est la conclusion?
\end{enumerate}
\finenonce{002210} 


\finexercice


\exercice{2211, matos, 2008/04/23}
\enonce{002211}{Quelques identit\'es pour le calcul d'inverses}
D\'emontrer l'identit\'e
$$ (A + UBV)^{-1} = A^{-1} -A^{-1} U( I+BVA^{-1}U)^{-1} BVA^{-1}$$
en pr\'ecisant:
\begin{itemize}
\item son domaine de validit\'e;
\item les types des matrices $A, U, B, V$.
\end{itemize}
Quelques cas particuliers:
\begin{enumerate}
\item Supposons $B=\beta$ scalaire, $U=u\in\Rr^n$, $V=v^T \in\Rr^n$. Retrouver la formule de Shermann--Morrison qui permet le calcul de l'inverse d'une matrice qui apparait comme perturbation de rang 1 d'une matrice dont on connait l'inverse.
\item Soient $A\in\Rr^{n\times n} $ r\'eguli\`ere et $u,v\in\Rr^n$ tels que $1+v^TA^{-1}u=0$. Montrer que
$$B=\left(\begin{array}{cc}
A+uv^T&u\\
v^T& 0\end{array}\right) \mbox{ est r\'eguli\`ere.}$$
Calculer $B^{-1}$ en remarquant que
$$B=\left[\begin{array}{cc}
A&0\\
0&-1\end{array}\right]
 +\left[\begin{array}{c}
u\\
1\end{array}\right] \left[\begin{array}{cc}
v^T&1
\end{array}\right]$$
\item Soit 
$$D=\left[\begin{array}{cc}
P&Q\\
R&S\end{array}\right] \mbox{ matrice inversible  avec } P\in\Rr^{p\times p}, Q\in\Rr^{p\times q}, S\in\Rr^{q\times q}$$
Calculer $D^{-1}$ en remarquant que
$$D=\left[\begin{array}{cc}
P&0\\
R&I\end{array}\right] +\left[\begin{array}{c}
Q\\S-I\end{array}\right]\left[\begin{array}{cc}
0&I\end{array}\right]$$
\item {\it Calcul r\'ecursif de l'inverse}: on pose
$$A_n=\left[\begin{array}{cc}A_{n-1}&v\\
u^T&s\end{array}\right] \mbox{ avec } A_{n-1}\in\Rr^{(n-1)\times (n-1)} u,v\in\Rr^{n-1}, s\in\Rr$$
Utiliser la formule pr\'ec\'edente pour calculer $A_n^{-1}$ en fonction de $A_{n-1}^{-1}$. En d\'eduire un algorithme r\'ecursif pour le calcul de l'inverse d'une matrice carr\'ee de taille $n$.
\end{enumerate}
\finenonce{002211} 


\finexercice


\exercice{2212, matos, 2008/04/23}
\enonce{002212}{Quelques propri\'et\'es des normes matricielles}
\begin{enumerate}
\item Soit $A$ une matrice d'ordre $(m,n)$. D\'emontrer les in\'egalit\'es suivantes pour les normes $p$, $p=1,2, \infty$ et la norme de Frobenius:
\begin{enumerate}
\item $\|A\|_2 \leq \|A|_F \leq \sqrt{n} \|A|_2$
\item $\max |a_{ij}| \leq \|A\|_2 \leq \sqrt{mn} \max |a_{ij}|$
\item $ \displaystyle{\frac{1}{\sqrt{n}}}\|A\|_\infty \leq \|A\|_2 \leq \sqrt{m} \|A\|_\infty$
\item $\displaystyle{\frac{1}{\sqrt{m}}}\|A\|_1 \leq \|A\|_2 \leq \sqrt{n} \|A\|_1$
\end{enumerate}
\item Soit $u\in\Rr^m$, $v\in\Rr^n$ et $E=uv^T$. Montrer que
$$\|E\|_F = \|E\|_2 = \|u\|_2\|v\|_2$$
$$\|E\|_\infty =\|u\|_\infty \|v\|_1$$
\end{enumerate}
\finenonce{002212} 


\finexercice


\exercice{2213, matos, 2008/04/23}
\enonce{002213}{}
Montrer que si $\rho (A) <1$ alors
\begin{itemize}
\item $I-A$ est r\'eguli\`ere;
\item $(I-A)^{-1}=\lim_{k \rightarrow \infty} C_k$ avec $C_k=I+A+\cdots +A^k$.
\end{itemize}
\finenonce{002213} 


\finexercice


\exercice{2214, matos, 2008/04/23}
\enonce{002214}{Estimation de l'erreur dans le calcul de l'inverse}
Soit $A$ une matrice carr\'ee d'ordre $n$ inversible et $B$ une approximation de $A^{-1}$ On pose $X=I-AB$ et on suppose que $\|X\|<1$. Montrer que
$$\|A^{-1}-B\| \leq \frac{\|BX\|}{1-\|X\|}.$$
\finenonce{002214} 


\finexercice


\exercice{2215, matos, 2008/04/23}
\enonce{002215}{Projection orthogonale sur un sous-espace vectoriel de $\Rr^n$}
Soient
\begin{itemize}
\item $H=$span$\{v_1, \cdots , v_r\}$ le sous--espace vectoriel de $\Rr^n$ engendr\'e par les vecteurs $\{v_i\}$ suppos\'es ind\'ependants;
\item $V=\left[ \begin{array}{cccc}
v_1&v_2 & \cdots & v_r\end{array}\right]$ la matrice de type $n\times r$ dont les colonnes sont les composantes des $v_i$ dans la base canonique $\epsilon = (e_1, \cdots , e_n)$
\end{itemize}
Pour tout $x\in \Rr^n$ on d\'esigne par $y$ sa projection orthogonale sur $H$ et par $X$ et $Y$ les matrices colonnnes des composantes de $x$ et $y$ dans la base $\epsilon$.
On pose
$$y=\sum_{i=1}^r \alpha_i v_i .$$
\begin{enumerate}
\item Montrer que la matrice $G=G(v_1, \cdots ,v_r)=V^TV$ est inversible.
 \item Montrer que les $\alpha_i$ v\'erifient le syst\`eme
$$G\left( \begin{array}{c}\alpha_1 \\ \vdots \\ \alpha_r\end{array}
\right) = V^TX$$
\item En d\'eduire que $Y=VG^{-1}V^TX= PX$ avec $P=VG^{-1}V^T$ ($P$ est donc la matrice de la projection orthogonale de $\Rr^n$ sur $H$
\item \emph{Application}: on consid\`ere $n=3, v_1=e_1, v_2=e_1+e_2+e_3$. D\'eterminer la projection orthogonale sur $H=$ span $\{v_1, v_2\}$ de $x=2e_1-e_2+e_3$.
\item Quelle est la matrice de la projection orthogonale sur $H=$span $\{v\}$?
\item Montrer que, pour $x\in\Rr^n$
$$d^2(x,H)= \frac{\mbox{det} G(x,v_1,\cdots ,v_r)}{\mbox{det} G(v_1, \cdots , v_r)}$$
\end{enumerate}
\finenonce{002215} 



\finexercice


\exercice{2216, matos, 2008/04/23}
\enonce{002216}{}
Soit $A\in\Rr^{m\times n}$ de rang $r\leq p=\min (m,n)$. On consid\`ere la d\'ecomposition en valeurs singuli\`eres de $A$
 $$U^TAV = \mbox{diag} (\sigma_1, \cdots , \sigma_p )$$
o\`u les $\sigma_i$ sont les valeurs singuli\`eres de $A$
\begin{enumerate}
\item Montrer que Im$(A)=$ span$\{u^1, u^2, \cdots , u^r\}$ et 
 Ker$(A)=$ span$\{v^{r+1}, \cdots , v^n\}$.
\item  Montrer que Im$(^TA)=$ span$\{v^1, u^2, \cdots , v^r\}$ et 
 Ker$(A^T)=$ span$\{u^{r+1}, \cdots , u^m\}$.
\item D\'eterminer les matrices des projections orthogonales sur Im$(A)$, Ker$(A)$, Im$(A^T)$, Ker$(A^T)$ \`a l'aide de $U$ et $V$.
\item \emph{Application:} calculer la d\'ecomposition en valeurs singuli\`eres de la matrice
$$A=\left(\begin{array}{cc}1&1\\
2&1\\
-1&1\end{array}\right)$$
et les matrices correspondantes aux projections orthogonales de l'exercice pr\'ec\'edent.
\end{enumerate}
\finenonce{002216} 


\finexercice


\exercice{2217, matos, 2008/04/23}
\enonce{002217}{Pseudo--inverse d'une matrice}
\emph{D\'efinition}: Soit $\Sigma$ une matrice diagonale de type $(m\times n)$:
$$\Sigma =\left( \begin{array}{cccccc}
\mu_1&&&&&\\
&\ddots&&&&\\
&&\mu_r&&&\\
&&&0&&\\
&&&&\ddots &\\
&&&&&0\\
&&\bigcirc&&&
\end{array}\right)$$
On appelle pseudo--inverse de $\Sigma$ la matrice $\Sigma^\dagger$ de type $(n\times m)$ d\'efinie par
$$\Sigma^\dagger=\left(\begin{array}{cccc}
\mu_1^{-1}&&0&\\
&\ddots &&\bigcirc\\
0&&\mu_r^{-1}&
\end{array}\right)$$
Soit $A$ une matrice de type $(m\times n)$ dont la d\'ecomposition en valeurs singuli\`eres est $A=U\Sigma V^*$.

On appelle {\it pseudo-inverse } de la matrice $A$ la matrice $A^\dagger$ de type $(n\times m)$ d\'efinie par
$$A^\dagger = V\Sigma^\dagger U^* .$$
\begin{enumerate}
\item Quelle application repr\'esente la restriction de $\Sigma^\dagger \Sigma$ au sous-espace span$\{e_1, \cdots , e_r\}$ ?
\item Montrer que si $A$ est carr\'ee r\'eguli\`ere alors $A^\dagger =A^{-1}$.
\item Montrer que
$$A^\dagger =\sum_{i=1}^r \frac{1}{\mu_i}v_iu_i^*.$$
\item Montrer que
\begin{itemize}
\item $AA^\dagger $ est la matrice de la projection orthogonale sur Im$(A)$;
\item $A^\dagger A $ est la matrice de la projection orthogonale sur Im$(A^*)$
\end{itemize}
\item Montrer que la restriction \`a Im$(A^*)=$Ker$(A)^\bot$ de $A^*A$ est une matrice inversible et
$$(A^*A)^{-1} = \sum_{i=1}^r \mu_i^{-2} v_iv_i^* .$$
\end{enumerate}
\finenonce{002217} 


\finexercice


\exercice{2218, matos, 2008/04/23}
\enonce{002218}{}
Montrer que, pour $A\in\Cc^{n\times m}$
\begin{enumerate}
\item $\|A\|_2=\sigma_1$, la plus grande valeur singuli\`ere de $A$
\item $\|A\|_F=\sqrt{\sigma_1^2 +\sigma_2^2 +\cdots +\sigma_r^2}$ o\`u les $\sigma_i$ sont les valeurs singuli\`eres de $A$.
\item Les valeurs singuli\`eres non nulles de $A$ sont les racines carr\'es des valeurs propres non nulles de $A^*A$ et $AA^*$.
\item pour $A\in\Cc^{m\times m}$ , $|\det (A)|=\prod_{i=1}^m \sigma_i$.
\item Si $A=A^*$ alors les valeurs singuli\`eres de $A$ sont les valeurs absolues des valeurs propres de $A$
\end{enumerate}
\finenonce{002218} 


\finexercice


\exercice{2219, matos, 2008/04/23}
\enonce{002219}{}
Montrer que 
\begin{enumerate} 
\item cond$_2 (A) = \mu_n(A)/\mu_1(A)$ avec $\mu_n(A)$ et $\mu_1(A)$ respectivement la plus grande et la plus petite valeur singuli\`ere de $A$;
\item si $A$ est normale alors
$$\mbox{cond}_2(A) =\frac{\max_i |\lambda_i(A)|}{\min_i |\lambda_i(A)|} ;$$
\item Si $A \in\Rr^{n\times n}$ est inversible, $Q\in\Rr^{n\times n}$ orthogonale alors
$$\mbox{cond}_2(A)=\mbox{cond}_2(AQ)=\mbox{cond}_2(QA)$$
\end{enumerate}
\finenonce{002219} 


\finexercice


\exercice{2220, matos, 2008/04/23}
\enonce{002220}{}
Soit $A=\left(\begin{array}{cc}
1&0\\
0&10^{-6}
\end{array}\right)$
\begin{enumerate}
\item Calculer $\mbox{cond}_2(A)$, $\mbox{cond}_1(A)$ et $\mbox{cond}_\infty(A)$;
\item R\'esoudre: 
\begin{itemize}
\item $Ax=b$ pour $b=\left(\begin{array}{c}1\\10^{-6}\end{array}\right)$
\item $Ay=b+\delta b$ pour $\delta b=\left(\begin{array}{c}10^{-6}\\0\end{array}\right)$ et $Az=b+\Delta b$ pour $\Delta b=\left(\begin{array}{c}0\\10^{-6}\end{array}\right)$
\end{itemize}
\item Pour chacune des trois normes consid\'er\'ees, trouver une majoration th\'eorique de
$$\frac{\|y-x\|}{\|x\|} \mbox{ et } \frac{\|z-x\|}{\|x\|}$$
et comparer avec les valeurs exactes. Quelle conclusion?
\end{enumerate}
\finenonce{002220} 



\finexercice


\exercice{2221, matos, 2008/04/23}
\enonce{002221}{Conditionnement du probl\`eme de l'inversion d'une matrice}

Soit $A$ une matrice inversible donn\'ee.
\begin{enumerate}
\item si $(A+\delta A)$ est une matrice inversible, d\'emontrer
$$\frac{\|(A+\delta A)^{-1}-A^{-1}\|}{\|(A+\delta A)^{-1}\|} \leq \mbox{cond} (A) \frac{\|\delta A\|}{\|A\|}$$
\item D\'emontrer que
$$\frac{\|(A+\delta A)^{-1}-A^{-1}\|}{\|A^{-1}\|} \leq \mbox{cond} (A) \frac{\|\delta A\|}{\|A\|} (1+{ \cal O} (\|A\|)) $$
\end{enumerate}
\finenonce{002221} 


\finexercice

\exercice{2222, matos, 2008/04/23}
\enonce{002222}{Taille des \'el\'ements dans l'\'elimination de Gauss}

 Notons $\tilde{A}_k$ la matrice carr\'ee d'ordre $(n-k+1)$ form\'ee des \'el\'ements $a_{ij}^k , k\leq i,j \leq n$ de la matrice $A_k=(a_{ij}^k)$ obtenue come r\'esultat de la $(k-1)$--\`eme \'etape de l'\'elimination de Gauss. On suppose $A=A_1$ sym\'etrique d\'efinie positive. 
\begin{enumerate} 
\item Notant $(. , .)$ le produit scalaire euclidien et $v'\in\Rr^{n-k}$ le vecteur form\'e par les $(n-k)$ derni\`eres composantes d'un vecteur $v=(v_i)_{i=k}^n \in \Rr^{n-k+1}$ quelconque, \'etablir l'identit\'e 
$$ (\tilde{A}_k v,v)=(\tilde{A}_{k+1} v', v') + \frac{1}{a_{kk}^k} \left|a_{kk}^k v_k +\sum_{i=k+1}^n a_{ik}^k v_i\right|^2 .$$ 
\item Montrer que chaque matrice $\tilde{A_k}$ est sym\'etrique d\'efinie positive. 
\item Etablir les in\'egalit\'es suivantes: 
$$ 0< a_{ii}^{k+1} \leq a_{ii}^k ,\ \ \ k+1\leq i \leq n$$ 
$$\max_{k+1\leq i \leq n} a_{ii}^{k+1} = \max_{k+1\leq i,j \leq n}\left|a_{ij}^{k+1}\right| \leq \max_{k \leq i,j \leq n}\left|a_{ij}^k\right| =\max_{k\leq i\leq n}a_{ii}^k$$ 
\end{enumerate} 
\finenonce{002222} 


\finexercice


\exercice{2223, matos, 2008/04/23}
\enonce{002223}{Strat\'egie de pivotage} 
\ \\
\begin{enumerate} 
\item Montrer que pour une matrice quelconque $A=(a_{ij})$ de type $(2\times 2 )$ on a  
$$\mbox{cond}_2(A) =\sigma + (\sigma^2-1)^{1/2} \mbox{ avec } \sigma =\frac{\sum_{i,j=1}^2 |a_{ij}|^2}{2|\det (A)|}$$ 
\item Calculer les conditionnements cond$_p(.)$ pour $p=1,2,\infty $ des matrices exactes obtenues \`a la premi\`ere \'etape de la proc\'edure d'\'elimination de Gauss pour r\'esoudre le syst\`eme lin\'eaire 
$$\left\{\begin{array}{l} 
10^{-4} u_1 + u_2 =1\\ 
u_1 + u_2 =2 
\end{array}\right.$$ 
selon que l'on commence, ou non, par \'echanger les deux \'equations. Conclusion? 
\end{enumerate} 
\finenonce{002223} 



\finexercice


\exercice{2224, matos, 2008/04/23}
\enonce{002224}{Factorisation LU d'une matrice bande} 
 
Montrer que la factorisation LU pr\'eserve la structure des matrices bande au sens suivant : 
$$a_{ij}=0 \mbox{ pour } |i-j| \geq p \Rightarrow \left\{ \begin{array}{ll} 
l_{ij}=0 & \mbox{ pour } i-j \geq p\\ 
u_{ij}=0 & \mbox{ pour } j-i \geq p 
\end{array}\right.$$ 
\finenonce{002224} 


\finexercice


\exercice{2225, matos, 2008/04/23}
\enonce{002225}{Factorisation d'une matrice sym\'etrique} 
Soit $A$ une matrice sym\'etrique inversible admettant une factorisation LU. Montrer que l'on peut \'ecrire $A$ sous la forme 
$$A=B\tilde{B}^T  \mbox{ o\`u }$$ 
\begin{itemize} 
\item $B$ est une matrice triangulaire inf\'erieure ;
\item $\tilde{B}$ est une matrice o\`u chaque colonne est soit \'egale \`a la colonne correspondante de $B$, soit \'egale \`a la colonne correspondante de $B$ chang\'ee de signe. 
\end{itemize} 
\emph{Application num\'erique} 
$$A=\left(\begin{array}{cccc} 
1&2&1&1\\ 
2&3&4&3\\ 
1&4&-4&0\\ 
1&3&0&0 
\end{array}\right).$$ 
\finenonce{002225} 


\finexercice


\exercice{2226, matos, 2008/04/23}
\enonce{002226}{Quelques factorisations LU} 
\begin{enumerate} 
\item Soit $A=LU$ la d\'ecomposition LU d'une matrice $A\in\Rr^{n\times n}$ avec $|l_{ij}|\leq 1$. Soient $a_i^T$ et $ u_i^T$ les lignes {\it i} de $A$ et $U$ respectivement. Montrer que 
$$u_i^T=a_i^T-\sum_{j=1}^{i-1} l_{ij}u_j^T$$ 
et que 
$$ \|U\|_\infty \leq 2^{n-1}\|A\|_\infty$$ 
\item Soit $A\in\Rr^{n\times n}$ d\'efinie par 
$$a_{ij}=\left\{\begin{array}{ll} 
1& \mbox{ si } i=j \mbox{ ou } j=n\\ 
-1 & \mbox{ si } i>j\\ 
0 & \mbox{ sinon } 
\end{array}\right.$$ 
Montrer que $A$ a une d\'ecomposition LU avec $|l_{ij}|\leq 1$ et $u_{nn}=2^{n-1}$. 
\end{enumerate} 
\finenonce{002226} 



\finexercice


\exercice{2227, matos, 2008/04/23}
\enonce{002227}{}
On suppose $A\in\Rr^{n\times n}$ inversible. Montrer que si $PA\Pi =LU$ est obtenue par la m\'ethode de Gauss avec pivotage total, alors 
$$\begin{array}{l} 
\forall i,j=1, \cdots ,n \ \ |l_{ij}|\leq 1 \\ 
\forall i=1, \cdots  , n , \forall j=i, \cdots , n ,\ \ \ |u_{ij}| \leq |u_{ii}| 
\end{array}$$ 
\finenonce{002227} 



\finexercice


\exercice{2228, matos, 2008/04/23}
\enonce{002228}{}
Soit $A\in\Rr^{n\times n}$ telle que $A^T$ soit \`a diagonale strictement dominante. Montrer que $A$ admet une d\'ecomposition LU avec $L^T$ \`a diagonale strictement dominante.
\finenonce{002228} 


\finexercice


\exercice{2229, matos, 2008/04/23}
\enonce{002229}{Matrices de Householder}

\begin{enumerate}
\item Soit $v$ un vecteur r\'eel v\'erifiant $v^Tv=1$. Montrer que la matrice de Householder
$$H(v)=I-2vv^T$$
repr\'esente une sym\'etrie par rapport au sous--espace vectoriel form\'e par les vecteurs orthogonaux aux vecgteurs $v$. En d\'eduire que $\det (H(v))=-1$.
\item D\'emontrer que toute matrice orthogonale est le produit de au plus $n$ matrices de Householder. En d\'eduire une interpr\'etation g\'eom\'etrique des matrices orthogonales.
\end{enumerate}
\finenonce{002229} 


\finexercice


\exercice{2230, matos, 2008/04/23}
\enonce{002230}{Algorithme de Gram--Schmidt et Gram--Schmidt modifi\'e}

Etant donn\'es $n$ vecteurs lin\'eairement ind\'ependants de $\Rr^m$, $\{a_1, \cdots , a_n\}$, on veut calculer une base orthonormale pour span$\{a_1, \cdots , a_n\}$.

 On pose $A=[ a_1, a_2, \cdots , a_n] \in \Rr^{m\times n}$ et on consid\`ere la factorisation QR de $A$,
$$ A=QR,\ \ \ Q=[q_1, \cdots , q_n] ,\ \ r_i^T, i=1, \cdots , n \mbox{ les lignes de }R$$
\begin{enumerate}
\item Montrer que 
$$\mbox{Im} A= \mbox{span} \{q_1, \cdots , q_n\} .$$
\item Montrer que
$$q_k=\frac{1}{r_{kk}} \left( a_k -\sum_{i=1}^{k-1}r_{ik}q_i\right) \ \ \ k=1, \cdots , n$$
\item En d\'eduire un algorithme pour le calcul r\'ecursif des $q_i$ (algorithme de Gram--Schmidt).
\item Algorithme de Gram--Schmidt modifi\'e

L'algorithme pr\'ec\'edent est instable num\'eriquement d\^u \`a la perte d'orthogonalit\'e dans le calcul des $q_i$. On va reformuler l'algorithme pour le rendre stable.

Pour $k=1, \cdots , n-1$, on d\'efinit $A^{(k)}\in\Rr^{m\times (n-k+1)}$ de la fa\c con suivante:
$$[0 , A^{(k)}] =A-\sum_{i=1}^{k-1} q_ir_i^T = \sum_{i=k}^n q_ir_i^T$$
et on va d\'ecrire l'\'etape $k$ de l'algorithme.
\begin{enumerate}
\item Montrer que si on pose
$$ A^{(k)} =[z,B] ,\ \ \ z\in\Rr^m , \ \ \ B\in\Rr^{m\times (n-k)}$$
alors
$$r_{kk}=\|z\|_2, \ \ \ q_k=z/r_{kk}.$$
\item Comment peut--on calculer la ligne $k$ de $R$  \`a partir de $A^{(k)}$?

\item Calculer $A^{(k+1)}$.
\item A partir des questions pr\'ec\'edentes, d\'ecrire l'algorithme qui permet le calcul de la factorisation $A=Q_1R_1$, $Q_1\in\Rr^{m \times n}$ orthonormale, $R_1\in\Rr^{n\times n}$ triangulaire sup\'erieure (Gram--Schmidt modifi\'e). Le calcul de $Q_1$ doit se faire sur place.
\item Quelle est la complexit\'e de l'algorithme pr\'ec\'edent?
\end{enumerate}
\end{enumerate}
\finenonce{002230} 


\finexercice


\exercice{2231, matos, 2008/04/23}
\enonce{002231}{Rotation de Givens}

Soient $p,q : $ $1\leq p < q\leq n$, $c,s\in\Rr :$ $c^2+s^2=1$.\\
On consid\`ere les matrices
$$G=G_{p,q}(c,s)=\left(\begin{array}{ccccrccc}
1&&&&&&&\\
&\ddots &&&&&&\\
&&1&&&&&\\
&&&c&\cdots &-s&&\\
&&&&\ddots  &&&\\
&&&s&\cdots& c&&\\
&&&&&&\ddots &\\
&&\cdots &&&&&1
\end{array}\right)$$
\begin{enumerate}
\item Ecrire $G$ comme perturb\'ee de $I$ par des matrices de rang 1.
\item Montrer que $G$ est inversible, calculer $G^{-1}$, montrer que $G$ est orthogonale.
\item Quelle est l'action de $G$ sur $A\in\Rr^{n\times n}$?
\item Soit $A\in\Rr^{n\times n}$ avec $a_{pj}=\alpha , a_{qj}=\beta$. Peut--on trouver $G$ telle que $A'=GA$  v\'erifie:
$$a'_{pj}=0 =\alpha ', \ \ \ a'_{qj}=0=\beta ' ?$$
Est--ce que la solution est unique?
\end{enumerate}
\finenonce{002231} 


\finexercice


\exercice{2232, matos, 2008/04/23}
\enonce{002232}{}
Soit $Z=\left(\begin{array}{cc} c&s\\ -s&c\end{array}\right)$ avec $c^2 + s^2=1$.
On d\'efinit $\rho$ par
$$\rho = \left\{\begin{array}{lrr}
1 &\mbox{si} & c=0\\
1/2\mbox{sign}(c) s  &\mbox{si} & |s|< |c|\\
2\mbox{sign}(s)/c  &\mbox{si} & |c| \leq |s|
\end{array}\right.$$
\begin{enumerate}
\item Comment reconstruire $\pm Z$ \`a partir de $\rho$?
\item Soit $Q$ une matrice orthogonale produit de $n$ rotations de Givens: $Q=J_1 \cdots J_n$. Comment peut--on stocker de la fa\c con la plus \'economique $Q$ sous forme factoris\'ee?
\item Modifier l'algorithme de Givens pour r\'eduire $A$ \`a la forme triangulaire sup\'erieure ($QA=R$, $Q$ matrice produit de rotations de Givens) en stockant sur place ( donc dans $A$) toute l'information n\'ecessaire \`a reconstruire $Q$.
\item Ecrire l'algorithme qui, \`a partir des r\'esultats de l'algorithme pr\'ec\'edent permet de reconstruire $Q$.
\end{enumerate}
\finenonce{002232} 



\finexercice


\exercice{2233, matos, 2008/04/23}
\enonce{002233}{}
Soient $x$ et $y$ deux vecteurs unitaires. Donner un algorithme qui utilise les transformations de Givens pour calculer une matrice $Q$ telle que 
$Qx=y .$
\finenonce{002233} 



\finexercice


\exercice{2234, matos, 2008/04/23}
\enonce{002234}{M\'ethode de Givens Rapide}

Soit $A\in\Rr^{m\times n}$. On veut construire une matrice $M\in \Rr^{m\times m}$ telle que
\begin{itemize}
\item $MA=S$ triangulaire sup\'erieure;
\item $MM^T=D=$diag$(d_1,\cdots , d_m)$ , $d_i>0$
\end{itemize}
et appliquer cette factorisation de $A$ dans la r\'esolution de syst\`emes au sens des moindres carr\'es.
\begin{enumerate}
\item Donner la factorisation $QR$ de $A$ en termes de $M, D$ et $S$.
\item On consid\`ere maintenant $m=2$. Soient $x=(x_1, x_2)^T$ et $D=$diag$(d_1,d_2)$ ($d_i>0$) donn\'es.
\begin{enumerate}
\item On d\'efinit
$$M_1=\left(\begin{array}{cc}
\beta_1 & 1\\ 1&\alpha_1\end{array}\right).$$
Supposons $x_2\neq 0$. Calculer $M_1x$ et $M_1DM_1^T$.

Comment choisir $\alpha_1$ et $\beta_1$ de fa\c con \`a ce que la deuxi\`eme composante de $M_1x$ soit nulle et que $M_1DM_1^T$ soit diagonale?

Pour le choix pr\'ec\'edent d\'eterminer $\gamma_1$ tel que
$$M_1x=\left(\begin{array}{c}
x_2(1+\gamma_1)\\0\end{array}\right) \mbox{ et } M_1DM_1^T=\left(\begin{array}{cc}
d_2(1+\gamma_1)&0\\0 & d_1(1+\gamma_1)\end{array}\right)$$
\item Supposons $x_1\neq 0$. On d\'efinit
$$M_2=\left(\begin{array}{cc}
1&\alpha_2\\ \beta_2 &1\end{array}\right).$$
Choisir $\alpha_2$ et $\beta_2$ de fa\c con \`a ce que
$$M_2x=\left(\begin{array}{c}
x_1(1+\gamma_2)\\0\end{array}\right) \mbox{ et } M_2DM_2^T=\left(\begin{array}{cc}
d_1(1+\gamma_2)&0\\0 & d_2(1+\gamma_2)\end{array}\right)$$
et d\'eterminer $\gamma_2$
\item Montrer que l'on peut toujours choisir $M_i$ ($i=1,2$) de fa\c con \`a ce que le ``facteur de croissance'' $(1+\gamma_i)$ soit inf\'erieur \`a 2.
\end{enumerate}
\item Soit maintenant $m\in\Nn$ quelconque. D\'efinir les matrices $M_1(p,q)$ et $M_2(p,q)$ telles que
$$\left(\begin{array}{cc}
m_{pp}& m_{pq}\\
m_{qp} & m_{qq}
\end{array}\right)=\left(\begin{array}{cc}
\beta_1& 1\\
1 & \alpha_1
\end{array}\right) \mbox{ ou } =\left(\begin{array}{cc}
1& \alpha_2\\
\beta_2 & 1
\end{array}\right)$$
\begin{itemize}
\item $e_q^T M_i(p,q)x = 0$;
\item $M_iDM_i^T$ matrice diagonale, avec $D=$diag$(d_1, \cdots ,d_n), d_i>0$
\end{itemize}
Ces matrices $M_i$ sont appel\'ees {\it matrice de Givens rapide}.
\item D\'ecrire l'algorithme qui utilise les transformations de Givens rapides pour r\'eduire $A\in\Rr^{m\times n}$ \`a la forme triangulaire sup\'erieure ({\it m\'ethode de Givens rapide}):
$$MA=R ,\ \ \ \ MM^T =\mbox{diag}(d_1, \cdots , d_m) .$$
Les calculs doivent \^etre faits sur place.

Quel est le co\^ut de cet algorithme? Comparer avec le co\^ut de la m\'ethode de Householder pour r\'eduire $A$ \`a la forme triangulaire sup\'erieure.
\item Application \`a la r\'esolution d'un syst\`eme lin\'eaire au sens des moindres carr\'es.
\begin{enumerate}
\item Comment profiter des r\'esultats fournis par l'algorithme pr\'ec\'edent pour r\'esoudre
$$\min_{x\in\Rr^n}\|Ax-b\|_2 \mbox{ avec } A\in\Rr^{m\times n} \ \ \ (m>n) ,\ \ \ b\in \Rr^m ?$$
\item Quelles modifications introduire dans l'algorithme de la m\'ethode de Givens rapide pour qu'il r\'esolve le probl\`eme de moindres carr\'es de la question pr\'ec\'edente?
\end{enumerate}
\item {\it Application num\'erique}: r\'esoudre au sens des moindres carr\'es par la m\'ethode de Givens rapide le syst\`eme
$$Ax=b, A=\left(\begin{array}{cc}1&4\\2&5\\3&6\end{array}\right) ,
 b=\left(\begin{array}{c}7\\8\\9\end{array}\right)$$
\item Consid\'erons maintenant le probl\`eme de moindres carr\'es
\begin{equation}\label{ls}
\min_{x\in\Rr^n}\|D(Ax-b)\|_2
\end{equation}
avec $A\in\Rr^{m\times n}$, $b\in\Rr^m$ , $D=$diag$(d_i)$ ($d_i>0$). Cela correspond \`a donner un poids diff\'erent \`a chaque \'equation du syst\`eme.

Soit $M$ une matrice produit de matrices de Givens rapide v\'erifiant
$$\left\{ \begin{array}{l}
MA=R \mbox{ triangulaire sup\'erieure}\\
MD^{-2}M^T =\tilde{D}=\mbox{diag}(\tilde{d}_i) , \  \ \ \tilde{d}_i>0
\end{array}\right.$$
Comment peut--on r\'esoudre le probl\`eme (\ref{ls})?

Quelles adaptations faire \`a l'algorithme pr\'ec\'edent?
\end{enumerate}
\finenonce{002234} 


\finexercice


\exercice{2235, matos, 2008/04/23}
\enonce{002235}{}
Soit $a\in\Rr$ et $A=\left( \begin{array}{ccc}
1&a&a\\ a&1&a\\ a&a&1\end{array}\right)$
\begin{enumerate} 
\item Pour qu'elles valeurs de $a$ $A$ est--elle d\'efinie positive?
\item Pour qu'elles valeurs de $a$ la m\'ethode de Gauss--Seidel est--elle convergente?
\item Ecrire la matrice $J$ de l'it\'eration de Jacobi.
\item Pour qu'elles valeurs de $a$ la m\'ethode de Jacobi converge--t--elle?
\item Ecrire la matrice ${\cal L}_1$ de l'it\'eration de Gauss--Seidel. Calculer $\rho ({\cal L}_1)$.
\item Pour quelles valeurs de $a$ la m\'ethode de Gauss--Seidel converge--t--elle plus vite que celle de Jacobi?
\end{enumerate}
\finenonce{002235} 



\finexercice


\exercice{2236, matos, 2008/04/23}
\enonce{002236}{}
Soit $A$ une matrice hermitienne inversible d\'ecompos\'ee en $A=M-N$ o\`u $M$ est inversible. Soit $B=I-M^{-1}A$ la matrice de l'it\'eration:
$$x_{n+1}= Bx_n +c .$$
Supposons que $M+M^* -A$ soit d\'efinie positive.
\begin{enumerate}
\item Soit $x$ un vecteur quelconque et on pose $y=Bx$. Montrer l'identit\'e:
$$(x,Ax) - (y,Ay) = ((x-y), (M+M^*-A)(x-y)).$$
\item Supposons que $A$ est d\'efinie positive. Soit $x\neq 0$ un vecteur propre de $B$ associ\'e \`a la valeur propre $\lambda$, $y=Bx=\lambda x$. Utiliser l'identit\'e pr\'ec\'edente pour montrer que $|\lambda|<1$. Que peut--on conclure sur la convergence de la m\'ethode?
\item Supposons maintenant que $\rho (B)<1$. montrer que $A$ est d\'efinie positive.
\item Supposons $A$ d\'ecompos\'ee par points ou par blocs sous la forme
$$A=D-E-F \mbox{ avec } D \mbox{ d\'efinie positive} .$$
Montrer que la m\'ethode de relaxation par points ou par blocs pour $0<w<2$ converge si et seulement si $A$ est d\'efinie positive.
\end{enumerate}
\finenonce{002236} 


\finexercice


\exercice{2237, matos, 2008/04/23}
\enonce{002237}{}
Soit $A=I-E-E^*$ une matrice carr\'ee d'ordre $N$ o\`u $E$ est une matrice strictement triangulaire inf\'erieure ($e_{ij}=0$ pour $i\leq j$). Pour r\'esoudre le syst\`eme $Ax=b
$, on propose la m\'ethode it\'erative d\'efinie par
$$\left\{ \begin{array}{ccc}
(I-E)x_{2k+1}&=& E^* x_{2k} + b\\
(I-E^*)x_{2k+2}&=& E x_{2k+1} + b
\end{array}\right.$$
\begin{enumerate}
\item D\'eterminer $B$ et $c$ pour que l'on ait:
$$x_{2k+2} =Bx_{2k} +c .$$
V\'erifier que $B=M^{-1}N $ et $A=M-N$ avec $M=(I-E)(I-E^*)$ , $N=EE^*$.
\item Montrer que $M^*+N$ est une matrice d\'efinie positive. En d\'eduire une condition n\'ecessaire et suffisante pour la convergence de la m\'ethode.
\end{enumerate}
\finenonce{002237} 


\finexercice


\exercice{2238, matos, 2008/04/23}
\enonce{002238}{}
Soient $A$ et $B$ deux matrices r\'eelles d'ordre $N$ et $a,b$ deux vecteurs de $\Rr^n$. On consid\`ere les deux it\'erations suivantes:
\begin{equation}\label{ite}
\left\{ \begin{array}{ccc}
x_{k+1}&=& By_k +a\\
y_{k+1}&=& Ax_k + b
\end{array}\right.  \ \ k=0, 1, \cdots
\end{equation}
avec $x_0, y_0 \in\Rr^n$ donn\'es.
\begin{enumerate}
\item D\'eterminer une condition n\'ecessaire et suffisante de convergence des deux suites de vecteurs.
\item Soit $z_k= ( x_k, y_k)^T \in\Rr^{2n}$. Montrer que (\ref{ite}) peut s'\'ecrire
$$z_{k+1}= Cz_k+c$$
o\`u  $C$ est une matrice d'ordre $2n$. Expliciter $C$ et $c$.
\item Montrer que $\rho^2(C)=\rho(AB)$.
\item On consid\`ere maintenant les deux it\'erations suivantes:
\begin{equation}\label{ite2}
\left\{ \begin{array}{ccc}
x_{k+1}&=& By_k +a\\
y_{k+1}&=& Ax_{k+1} + b
\end{array}\right.  \ \ k=0, 1, \cdots
\end{equation}
Donner une condition n\'ecessaire et suffisante de convergence.

Montrer que (\ref{ite2}) est \'equivalent \`a
$$z_{k+1}= Dz_k+d$$
o\`u  $D$ est une matrice d'ordre $2N$.

Montrer que $\rho (D)=\rho (AB)$.
\item {\bf Taux de convergence}

On appelle taux de convergence asymptotique de la matrice it\'erative $M$ le nombre
$$R(M) =-\ln (\rho (M))) .$$
On pose $e^k=x^k-x^*$ l'erreur de l'it\'er\'e d'ordre $k$.
\begin{enumerate}
\item Montrer que le nombre d'it\'erations $k$ pour r\'eduire l'erreur d'un facteur $\epsilon$ , i.e., $\frac{\|e^k\|}{\|e^0\|}\leq \epsilon $ v\'erifie
$$k\geq \frac{-\ln \epsilon}{R(M)} .$$
\item Comparer le taux de convergence des algorithmes (\ref{ite}) et (\ref{ite2}).
\end{enumerate}
\end{enumerate}
\finenonce{002238} 


\finexercice


\exercice{2239, matos, 2008/04/23}
\enonce{002239}{}
On consid\`ere le syst\`eme $Ax=b$ avec
\begin{equation}\label{sys}
A=\left(\begin{array}{ccccc}
3&1&0&0&0\\
1&2&1&0&0\\
0&2&3&1&0\\
0&0&1&4&3\\
0&0&0&1&1
\end{array}\right)
\end{equation}
\begin{enumerate}
\item D\'ecomposer $A$ sous la forme $LU$ et en d\'eduire que (\ref{sys}) admet une solution unique $x^*$.
\item Ecrire l'it\'eration de Gauss--Seidel pour ce syst\`eme, c'est--\`a--dire, le syst\`eme lin\'eaire donnant $X_{n+1}=(x_{n+1}, y_{n+1}, z_{n+1}, t_{n+1}, u_{n+1})$ en fonction de  $X_{n}=(x_{n}, y_{n}, z_{n}, t_{n}, u_{n})$.
\item Pour tout $n\in\Nn$ on pose $e_n=X_n-x^*$. Montrer qu'il existe $a\in [0,1[$ tel que:

$$\forall n\in\Nn \ \ \  \|e_{n+1}\|_\infty \leq a \|e_n\|_\infty .$$
En d\'eduire la convergence de la suite.
\item D\'eterminer la matrice de Gauss--Seidel ${\cal L}_1$ associ\'ee \`a $A$. Calculer $\|{\cal L}_1\|_\infty$. En d\'eduire la convergence de $(X_n)$ vers $x^*$.
\item Soit $A\in\Rr^{n\times n} $ v\'erifiant la propri\'et\'e suivante:
$$\begin{array}{ccl}
|a_{ij}| & \geq & \sum_{j\neq i}|a_{ij}| \ \ i=2, \cdots , n\\
|a_{11}|& > & \sum_{j\neq 1}|a_{1j}|
\end{array}$$
et sur chaque ligne de $A$ il existe il existe un terme non nul $a_{ij}$ pour $i\geq 2$ et $j<i$.

Montrer qu'alors la m\'ethode de Gauss--Seidel converge.

\end{enumerate}
\finenonce{002239} 


\finexercice

\section{470.00 Fonction convexe}


\exercice{1729, ridde, 1999/11/01}

\enonce{001729}{}
Soient $n \in \Nn^*$ et $x_1, \ldots, x_n \in ]0,  + \infty[$.
\begin{enumerate}
\item En utilisant la concavit\'e du $\log$, montrer que
$ (x_1 \ldots x_n)^{\frac 1n} \leq \frac{x_1 + \ldots + x_n}n$.
\item Montrer que $ (x_1 \ldots x_n)^{\frac 1n} \geq \frac 1{\frac 1{x_1} + \ldots +
\frac 1{x_n}}$.
\item En d\'eduire que $ n ! \leq (\frac{n + 1}2)^n$.
\end{enumerate}
\finenonce{001729}



\finexercice



\exercice{1730, ridde, 1999/11/01}

\enonce{001730}{}
Soit $f$ une fonction $C^2$ sur $\Rr$ convexe croissante et non constante.
Montrer que $\lim\limits_{ + \infty}f = $ tûûûûûût.
\finenonce{001730}



\finexercice



\exercice{1731, ridde, 1999/11/01}

\enonce{001731}{}
Soient $p$ et $q \in ]0,  + \infty[$ tels que $\frac 1p  + \frac 1q = 1$.
\begin{enumerate}

\item Montrer que $\forall x, y>0$ $xy \leq \frac{x^p}p + \frac{y^q}q$.

\item Soient $x_1, \ldots, x_n, y_1, \ldots, y_n>0$ tels que $\sum\limits_{i = 1}
^n x_i^p =\sum\limits_{i = 1}^n y_i^q  = 1$. Montrer que $\sum\limits_{i = 1}
^n x_i y_i \leq 1$.

\item Soient $x_1, \ldots, x_n, y_1, \ldots, y_n>0$. Montrer l'in\'egalit\'e de Hölder :
$$\sum\limits_{i = 1} ^n x_i y_i \leq (\sum\limits_{i = 1} ^n x_i^p )^{\frac 1p}
(\sum\limits_{i = 1} ^n y_i^q )^{\frac 1q}$$

\item Soit $p>1$. En \'ecrivant $ (x_i + y_i)^p = x_i (x_i + y_i)^{p-1} +
y_i (x_i + y_i)^{p-1}$, montrer l'in\'egalit\'e de Minkowski :
$$(\sum\limits_{i = 1} ^n { (x_i + y_i)}^p )^{\frac 1p} \leq
(\sum\limits_{i = 1} ^n x_i^p )^{\frac 1p}  + (\sum\limits_{i = 1} ^n y_i^p )^{\frac 1p} $$

\item Soit $ (a_n)$ une suite strictement positive, $u_n = \sum\limits_{k = 1}^n
a_k^2$ et $v_n= \sum\limits_{k = 1}^n \frac{a_k}k$. Montrer que si $ (u_n)$ converge
alors $ (v_n)$ aussi.

\end{enumerate}
\finenonce{001731}



\finexercice



\exercice{1732, ridde, 1999/11/01}

\enonce{001732}{}
Soit $f \in C^2 (\Rr)$ convexe.
\begin{enumerate}
\item Montrer que $f'$ admet une limite dans $\bar{\Rr}$ en $ + \infty$.
\item En d\'eduire que $\frac{f (x)}x$ admet une limite en $ + \infty$
 (on pourra utiliser des $\epsilon$ et une formule de Taylor \`a l'ordre 1).
\end{enumerate}
\finenonce{001732}



\finexercice



\exercice{1733, gourio, 2001/09/01}

\enonce{001733}{}
$I\subset {\Rr}^{+*} $ un intervalle de $ {\Rr},J=\left\{ x;\frac{1}{x}\in I\right\} .$

Montrer que $J$ est un intervalle de ${\Rr}^{+*},$ puis que si $(x,y)\in I^{2},$ alors :
$$\forall \lambda \in [0,1],\exists \mu \in [0,1],\frac{1}{\lambda
x+(1-\lambda )y}=\mu \frac{1}{x}+(1-\mu )\frac{1}{y}. $$
Soit $f$ continue sur $I,$ et $g$ d\'{e}finie sur $J$ par $g(x)=f(\frac{1}{x}),h$
d\'{e}finie sur $I $ par $h(x)=xf(x). $ Montrer que $g $ est convexe
$\Leftrightarrow h$ est convexe.
\finenonce{001733}



\finexercice



\exercice{1734, gourio, 2001/09/01}

\enonce{001734}{}
Soit $f :{\Rr}\rightarrow {\Rr}$ convexe major\'{e}e.\ Que dire de $f $ ? Et si
$f:{\Rr}^{+}\rightarrow {\Rr }$ ?
\finenonce{001734}



\finexercice



\exercice{1735, gourio, 2001/09/01}

\enonce{001735}{}
Soit $(a_{n})_{n\in \Nn}\in \left( {\Rr}^{+*}\right) ^{\Nn},
u_{n}=\sum\limits_{k=1}^{n}a_{k}^{2},v_{n}=\sum\limits_{k=1}^{n}\frac{a_{k}}{k}.$
Montrer que si $(u_{n})_{n}$ converge alors $(v_{n})_{n}$
aussi.
\finenonce{001735}



\finexercice



\exercice{1736, gourio, 2001/09/01}

\enonce{001736}{}
Montrer que :
$$\forall n\in \Nn^{*},\forall (x_{1},...,x_{n})\in \left( {\Rr}^{+*}\right)^{n},
1+\left( \prod\limits_{k=1}^{n}x_{k}\right) ^{\frac{1}{n}}\leq
\left( 1+\prod\limits_{k=1}^{n}x_{k}\right) ^{\frac{1}{n}}. $$
\finenonce{001736}



\finexercice



\exercice{1737, gourio, 2001/09/01}

\enonce{001737}{}
Soit $f:{\Rr}\rightarrow {\Rr }$ continue telle que :
$$\forall (x,y)\in {\Rr}^{2}, f\left( \frac{x+y}{2}\right) \leq \frac{%
f(x)+f(y)}{2}. $$
Montrer que $f$ est convexe.
\finenonce{001737}



\finexercice



\exercice{1738, gourio, 2001/09/01}

\enonce{001738}{}
Soit $f:I\rightarrow {\Rr}$ convexe ou $I$ est un intervalle ouvert de ${\Rr}$,
d\'{e}rivable en $x_{0}\in I$ et telle que $f^{\prime}(x_{0})=0.$
Montrer que $x_{0}$ minimise $f$ sur $I.$
\finenonce{001738}



\finexercice



\exercice{1739, gourio, 2001/09/01}

\enonce{001739}{}
Soit $g\in C({\Rr},{\Rr})$, montrer que $g$ est convexe si et seulement
si :
$$\forall h\in CM([0,1],{\Rr}), g\left( \int_{0}^{1}h\right) \leq
\int_{0}^{1}g(h). $$
\finenonce{001739}



\finexercice


\section{471.00 Multiplicateurs de Lagrange}

\section{472.00 Algorithme d'Uzawa}

\section{473.00 Algorithme du simplexe}

\section{474.00 Autre}

\section{480.00 Loi, indépendance, loi conditionnelle}


\exercice{6938, ruette, 2013/01/24}
\enonce{006938}{Indépendance}
Soient $X$ et $Y$ des variables aléatoires réelles indépendantes ayant 
des lois continues. Montrer que $P(X=Y)=0$.
\finenonce{006938}

\finexercice


\exercice{6939, ruette, 2013/01/24}
\enonce{006939}{Calculs de loi}
Soient $X$, $Y$ deux variables aléatoires indépendantes de même loi exponentielle
$\mathcal{E}(\lambda)$, $\lambda>0$. Déterminer la loi de $|X-Y|$.
\finenonce{006939}
\finexercice


\exercice{6940, ruette, 2013/01/24}
\enonce{006940}{Calcul de loi, fonction de répartition}
\begin{enumerate}
\item Soit $X$ une variable aléatoire de loi normale $\mathcal{N}(0,1)$
et $Z=X^2$. Calculer la fonction de répartition et la densité
de $Z$.

\textit{Remarque : la loi de $Z$ est appelée loi $\chi_2$ à 1 degré de liberté}.

\item 
Soit $Y$ une variable aléatoire de loi exponentielle $\mathcal{E}(\lambda)$. Déterminer 
la loi de $Y^3$.
\end{enumerate}
\finenonce{006940}
\finexercice


\exercice{6941, ruette, 2013/01/24}
\enonce{006941}{Fonction de répartition}

Soit $X$ une variable aléatoire réelle positive intégrable. Montrer que
$\displaystyle E(X)=\int_0^{+\infty}(1-F(t))\,dt$.
\finenonce{006941}
\finexercice


\exercice{6942, ruette, 2013/01/24}
\enonce{006942}{Indépendance d'événements}

Soient $A$ et $B$ deux événements indépendants. Montrer que
$A\bot B^c$. En déduire que $A^c\bot B$ et $A^c\bot B^c$.
\finenonce{006942}
\finexercice


\exercice{6943, ruette, 2013/01/24}
\enonce{006943}{La loi exponentielle est sans mémoire}
\begin{enumerate}
\item Soit $X$ une variable aléatoire de loi  $\mathcal{E}(\lambda)$. Montrer que
$P(X> t+s\mid X> t)=P(X> s)$ pour tous $s,t\geq 0$.
\item  Soit $X$ une variable aléatoire positive avec une densité continue sur
$\Rr_+$. Si
$P(X> t+s\mid X> t)=P(X> s)$  pour tous $s,t\geq 0$, montrer
que $X$ suit une loi exponentielle.

\medskip
\textit{
Ce résultat montre que la loi exponentielle est une loi sans mémoire,
et que c'est la seule sous l'hypothèse du 2.
En fait, cette hypothèse n'est pas nécessaire
mais le résultat est alors plus difficile à montrer.}
\end{enumerate}
\finenonce{006943}
\finexercice


\exercice{6944, ruette, 2013/01/24}
\enonce{006944}{Loi d'un temps d'arrêt avec un jeu de pile ou face}
Soit $(X_n)_{n\geq 1}$ une suite de variables aléatoires indépendantes de même loi de
Bernoulli de paramètre $p$. On définit la variable aléatoire $T_1$ à valeurs dans
$\Nn\cup\{+\infty\}$ par $T_1=\inf\{k>0 \mid X_k=1\}$, avec la
convention $\inf(\emptyset)=+\infty$.\\ \textit{(si on joue à pile $=0$ ou
 face $=1$,
$T_1$ est le temps nécessaire pour obtenir face une première fois)}
\begin{enumerate}
\item
Montrer que $T_1$ est fini presque sûrement.
\item
Déterminer la loi et l'espérance de $T_1$ \textit{(cette loi est appelée loi géométrique, $E(T_1)$ est le nombre
moyen de lancers qu'il faut effectuer pour obtenir face une première fois)}.
\item
Pour tout $n\geq 2$, on définit par récurrence
$T_n=\inf\{k>T_{n-1}\mid X_k=1\}$.\\
\textit{(si on joue à pile  ou face, $T_n$
est le temps nécessaire pour obtenir exactement $n$ fois face)}\\
Montrer que les variables aléatoires $T_1,\ (T_2-T_1), \ldots,\ (T_n-T_{n-1}),\ldots$
sont indépendantes et de même loi.
\item
Quelle est la loi de $T_n$ ?
\end{enumerate}
\smallskip
\textit{Définition générale d'un temps d'arrêt : une variable aléatoire $T$ à valeurs
dans $\Nn$ est un temps d'arrêt relativement à la suite de variable aléatoire 
$(X_n)_{n\geq 1}$ si pour tout $n$
l'événement $\{T\leq n\}$ appartient à la tribu 
engendrée par les variables aléatoires $X_1,\ldots, X_n$ (autrement dit, il suffit de 
connaître les valeurs de $X_1,\ldots, X_n$ pour savoir si $T\leq n$).}
\finenonce{006944}

\finexercice


\exercice{6945, ruette, 2013/01/24}
\enonce{006945}{Fonctions génératrices}
Soit $X$ une variable aléatoire de loi binomiale $\mathcal{B}(n,p)$.
\begin{enumerate}
\item
Calculer la fonction génératrice de $X$.
\item
Soit $Y$ une variable aléatoire de loi $\mathcal{B}(m,p)$, indépendante de $X$.
Quelle est la fonction génératrice de $X+Y$ ? En déduire que
$\mathcal{B}(n,p)*\mathcal{B}(m,p)=\mathcal{B}(n+m,p)$.
\end{enumerate}

{\it C'est une des façons de démontrer la stabilité des lois binomiales
par convolution. La méthode marche aussi pour les lois de Poisson.}
\finenonce{006945}

\finexercice


\exercice{6946, ruette, 2013/01/24}
\enonce{006946}{Fonctions génératrices}
Soit $N$ et $(X_n)_{n\geq 1}$ des variables aléatoires indépendantes intégrables 
à valeurs dans $\Nn$, les $X_n$ ayant tous la même loi et $P(N=0)=0$. On
pose $Y(\omega)=\sum_{n=1}^{N(\omega)}X_n(\omega)$.

Exprimer la fonction génératrice de $Y$ en fonction des fonctions génératrices
de $N$ et de $X_1$. Puis exprimer l'espérance de $Y$ en fonction de $E(N)$
et de $E(X_1)$.
\finenonce{006946}

\finexercice


\exercice{6947, ruette, 2013/01/24}
\enonce{006947}{Fonction génératrice}
Soit $X$ une variable aléatoire de loi de Poisson $\mathcal{P}(\lambda)$.
\begin{enumerate}
\item
Calculer la fonction génératrice de $X$.
\item
Calculer $E(X)$ et $\text{Var}(X)$.
\item
Soit $Y$ une variable aléatoire de loi de Poisson $\mathcal{P}(\lambda')$, indépendante
de $X$. Quelle est la loi de $X+Y$ ? En déduire que
$\mathcal{P}(\lambda) *\mathcal{P}(\lambda')=\mathcal{P}(\lambda+\lambda')$.
\end{enumerate}
\finenonce{006947}

\finexercice


\exercice{6948, ruette, 2013/01/24}
\enonce{006948}{Fonction caractéristique}
Soit $X$ une variable aléatoire de loi exponentielle $\mathcal{E}(\lambda)$. Calculer sa
fonction caractéristique.
\finenonce{006948}

\finexercice


\exercice{6949, ruette, 2013/01/24}
\enonce{006949}{Fonctions caractéristiques}
\begin{enumerate}
\item
Soit $X$ une variable aléatoire de loi de Poisson $\mathcal{P}(\lambda)$. Calculer sa
fonction caractéristique.
\item
Soit $Y$ une variable aléatoire indépendante de $X$ telle que $P_Y=\mathcal{P}(\lambda')$.
Quelle est la fonction caractéristique de $X+Y$ ? En déduire
que $\mathcal{P}(\lambda)*\mathcal{P}(\lambda')=\mathcal{P}(\lambda+\lambda')$.
\end{enumerate}
\finenonce{006949}

\finexercice

\section{481.00 Variance, covariance, fonction génératrice}


\exercice{6952, ruette, 2013/01/24}
\enonce{006952}{Calcul d'espérance et de variance}
\begin{enumerate}
\item Soit $X$ une variable aléatoire de loi exponentielle $\mathcal{E}(\lambda)$, 
$\lambda>0$. Calculer l'espérance de $X$.
\item
Soit $X$ une variable aléatoire de loi de Poisson $\mathcal{P}(\lambda)$, $\lambda>0$. Calculer
la variance de $X$.
\end{enumerate}
\finenonce{006952}
\finexercice


\exercice{6953, ruette, 2013/01/24}
\enonce{006953}{Variance}
Soit $X$ une variable aléatoire réelle dans $L^2$.
Montrer que $\displaystyle \text{Var}(X)=\min_{t\in\Rr}E((X-t)^2)$.
\finenonce{006953}
\finexercice

\section{482.00 Convergence de variables aléatoires}


\exercice{6950, ruette, 2013/01/24}
\enonce{006950}{Lemme de Borel-Cantelli **}
Soit $(X_n)_{n\geq 1}$ une suite de variables aléatoires indépendantes de même
loi exponentielle $\mathcal{E}(1)$.
%de densité $\I1_{\Rr_+}(x)e^{-x}$. 
\begin{enumerate}
\item
En utilisant le lemme de Borel-Cantelli, montrer que
$$
P(X_n>\alpha \ln n\mbox{ pour une infinité de }n)=\left\{
\begin{array}[c]{ll}0 \mbox{ si } \alpha>1\\ 1\mbox{ si }\alpha\leq 1
\end{array}\right.
$$
\item
En déduire que $\displaystyle\limsup_{n\to+\infty} \frac{X_n}{\ln n}=1$ 
presque sûrement.
\end{enumerate}
\finenonce{006950}

\finexercice


\exercice{6951, ruette, 2013/01/24}
\enonce{006951}{Convergence en probabilité **}
\begin{enumerate}
\item
Montrer que pour tout $x>0$,
$$
e^{-x^2/2}\left(\frac{1}{x}-\frac{1}{x^3}\right)\leq
\int_x^{+\infty}e^{-t^2/2}dt\leq e^{-x^2/2}\frac{1}{x}.
$$
{\it Indication : on pourra intégrer par parties $t^{-1}te^{-t^2/2}$.}
\item
Soit $(X_n)_{n\geq 1}$ une suite de variables aléatoires indépendantes de même loi
$\mathcal{N}(0,1)$. Montrer que $\displaystyle\frac{\max_{1\leq i\leq n} X_i}{\sqrt{2\ln n}}$
tend vers~$1$ en probabilité.
\end{enumerate}
\finenonce{006951}

\finexercice

\section{483.00 Lois des grands nombres, théorème central limite}

\section{484.00 Estimateur}

\section{485.00 Tests sur la moyenne, test du chi2}

\section{486.00 Chaînes de Markov}

\section{487.00 Autre}

\finfiche



 \finenonces 



 \finindications 

\noindication
\noindication
\indication{000106}
Attention : la n\'egation d'une in\'egalit\'e stricte est une in\'egalit\'e large
(et r\'eci\-pro\-quement).
\finindication
\noindication
\noindication
\indication{000109}
Faire un dessin de $F_1$ et de $F_2$.
Essayer de voir si la difficult\'e pour r\'ealiser les assertions vient de
$\epsilon$ ``petit'' (c'est-\`a-dire proche de $0$) ou de $\epsilon$ ``grand''
(quand il tend vers $+\infty$).
\finindication
\noindication
\noindication
\noindication
\noindication
\noindication
\noindication
\noindication
\noindication
\noindication
\indication{000119}
En fait, on a toujours : $\frac{2n+1}{n+2} \leq 2$.
Puis chercher une condition sur $n$ pour que
l'in\'egalit\'e 
$$2-\epsilon < \frac{2n + 1}{n + 2}$$
soit vraie.
\finindication
\noindication
\noindication
\noindication
\noindication
\noindication
\noindication
\indication{000123}
Il est plus facile de raisonner en prenant un \'el\'ement $x\in E$.
Par exemple, soit $F, G$ des sous-ensembles de $E$. Montrer que $F\subset G$
revient à montrer que pour tout $x\in F$ alors $x\in G$.
Et montrer $F=G$ est \'equivalent \`a $x\in F$ si et seulement si $x\in G$,
et ce pour tout $x$ de $E$.
Remarque : pour montrer $F=G$ on peut aussi montrer $F\subset G$ puis $G\subset F$.

Enfin, se rappeler que $x\in \complement F$ si et seulement si $x\notin F$.
\finindication
\noindication
\noindication
\noindication
\noindication
\noindication
\noindication
\noindication
\noindication
\noindication
\noindication
\noindication
\noindication
\noindication
\noindication
\noindication
\noindication
\noindication
\noindication
\noindication
\noindication
\noindication
\noindication
\noindication
\noindication
\noindication
\noindication
\noindication
\noindication
\noindication
\noindication
\noindication
\noindication
\indication{000150}
Par l'absurde, supposer qu'il existe $p\in \Nn$ tel que $f=f_p$.
Puis pour un tel $p$, \'evaluer $f$ et $f_p$ en une valeur bien choisie.
\finindication
\indication{000151}
Pour la premi\`ere question vous pouvez raisonner par contraposition ou par l'absurde.
\finindication
\noindication
\noindication
\noindication
\indication{000155}
\begin{enumerate}
    \item R\'ecurrence : calculer $x_{n+1}-3$.
    \item Calculer  $x_{n+1}-3 - \frac{3}{2}(x_n-3)$.
    \item R\'ecurrence.
\end{enumerate}
\finindication
\noindication
\indication{000157}
Pour les deux questions, travailler par r\'ecurrence.
\finindication
\noindication
\noindication
\noindication
\noindication
\noindication
\noindication
\noindication
\noindication
\noindication
\noindication
\noindication
\noindication
\noindication
\noindication
\noindication
\noindication
\noindication
\noindication
\noindication
\noindication
\noindication
\noindication
\noindication
\noindication
\noindication
\indication{007014}
Calculer les premiers termes de la suite.
\finindication
\indication{007015}
Calculer les premiers termes de la suite.
\finindication
\indication{007016}
Récurrence double.
\finindication
\noindication
\noindication
\noindication
\noindication
\noindication
\noindication
\noindication
\indication{007024}
Procéder par récurrence forte.
\finindication
\indication{007025}
Procéder par récurrence forte.
\finindication
\noindication
\indication{007036}
Montrer par récurrence sur $m$ que pour tout $N\geq 2$ et tout $m \in \llbracket 2, N\rrbracket$, 
\[ \sqrt{m\sqrt{(m+1)\sqrt{...\sqrt{(N-1)\sqrt{N}}}}} < m+1.\]
\finindication
\noindication
\noindication
\indication{007039}
Considérer le sous-ensemble (sous-monoïde) de $\N$ suivant :
\[ \{3k+4l\mid k,l \in \N\}.\]
Pour le cas général, considérer le pgcd de $a$ et $b$.
\finindication
\noindication
\noindication
\indication{007042}
Calculer les premiers termes de la suite.
\finindication
\noindication
\noindication
\noindication
\indication{007046}
Montrer que pour $n>1$, le réel $u_n$ s'écrit comme le quotient d'un entier impair par un entier pair. Pour $n$ pair, exprimer $u_{n+1}$ en fonction de $u_n$. Pour $n$ impair, utiliser le fait que
\[ u_{n+1} := \left(1+\frac13 + \frac 15 + ... + \frac1n \right) + 
\left(\frac12 + \frac 14 + ... + \frac{1}{n+1}\right).\]
\finindication
\noindication
\noindication
\indication{000209}
Un dessin permettra d'avoir une bonne id\'ee de ce qui se passe...
\finindication
\indication{000210}
Il faut trouver l'erreur dans ce raisonnement,
car bien s\^ur s'il y a trois axiomes pour la d\'efinition d'une relation d'\'equivalence, c'est que deux ne suffisent pas !
\finindication
\noindication
\indication{000212}
\begin{enumerate}
 \item  Pour la transitivit\'e on pourra calculer $xye^z$.
 \item Poser la fonction $t \mapsto \frac t {e^t}$, apr\`es une \'etude de fonction on calculera le nombre d'ant\'ec\'edents possibles.
 \end{enumerate}
\finindication
\noindication
\noindication
\noindication
\noindication
\noindication
\noindication
\noindication
\noindication
\noindication
\noindication
\noindication
\noindication
\noindication
\noindication
\noindication
\noindication
\noindication
\noindication
\noindication
\noindication
\noindication
\noindication
\noindication
\noindication
\noindication
\noindication
\noindication
\noindication
\noindication
\indication{007191}

Regarder par exemple la divisibilité sur $\N$.

\finindication
\indication{007192}

Considérer une certaine relation d'équivalence sur $E$ et écrire que $E$ est réunion disjointe des classes d'équivalence.

\finindication
\noindication
\noindication
\noindication
\noindication
\noindication
\noindication
\noindication
\noindication
\noindication
\noindication
\noindication
\noindication
\noindication
\noindication
\noindication
\noindication
\noindication
\noindication
\noindication
\noindication
\noindication
\indication{000185}
Prouver que l'\'egalit\'e est fausse.
\finindication
\noindication
\noindication
\noindication
\noindication
\noindication
\noindication
\noindication
\noindication
\noindication
\noindication
\noindication
\noindication
\noindication
\noindication
\noindication
\indication{000190}
\begin{enumerate}
\item $f$ est injective mais pas surjective.
\item $g$ est bijective. 
\item $h$ aussi.
\item $k$ est injective mais par surjective.
\end{enumerate}
\finindication
\indication{000191}
\begin{enumerate}
    \item $f$ n'est ni injective, ni surjective.
    \item Pour $y\in \Rr$, r\'esoudre l'\'equation $f(x)=y$.
    \item On pourra exhiber l'inverse.
\end{enumerate}
\finindication
\noindication
\indication{000193}
Pour la premi\`ere assertion le d\'ebut du raisonnement est : ``supposons
que $g\circ f$ est injective, soient $a,a'\in A$ tels que $f(a)=f(a')$'',... 
\`a vous de travailler, cela se termine par
``...donc $a=a'$, donc $f$ est injective.''
\finindication
\noindication
\noindication
\noindication
\noindication
\noindication
\noindication
\noindication
\indication{000199}
$id$ est l'application identité définie par $id(x)=x$ pour tout $x\in[0,1]$.
Donc $f \circ f= id$ signifie $f\circ f(c) = x$ pour tout $x\in[0,1]$.
\finindication
\indication{000200}
Montrer que la restriction de $f$ définie par : 
$[0,2\pi[ \longrightarrow \mathbb{U}$, $t\mapsto e^{it}$ est une bijection.
Ici $\mathbb{U}$ est le cercle unit\'e de $\Cc$, c'est-\`a-dire 
l'ensemble des nombres complexes de module \'egal \`a $1$. 
\finindication
\noindication
\indication{000202}
Montrer que $f$ est injective et surjective.
\finindication
\noindication
\noindication
\noindication
\noindication
\noindication
\noindication
\noindication
\noindication
\noindication
\indication{000220}
\'Evaluer $(1+x)^n$ en $x=1$, d'une part directement et ensuite avec la formule du bin\^ome de Newton.
Pour la deuxi\`eme \'egalit\'e commencer par d\'eriver $x\mapsto(1+x)^n$.
\finindication
\noindication
\indication{000222}
Commencer par $2^{n} = (3-1)^{n}$.
\finindication
\noindication
\noindication
\noindication
\noindication
\noindication
\noindication
\indication{000229}
$1+i= \frac{\sqrt2}2e^{\frac{2i\pi}{4}}$
\finindication
\noindication
\noindication
\noindication
\noindication
\noindication
\noindication
\noindication
\noindication
\noindication
\noindication
\noindication
\noindication
\noindication
\noindication
\noindication
\noindication
\noindication
\noindication
\noindication
\noindication
\noindication
\noindication
\noindication
\noindication
\noindication
\noindication
\noindication
\indication{000236}
Tout d'abord faire un dessin (avec des patates !).

Pour $A$ et $B$ deux ensembles finis quelconques, commencer par (re)démontrer la formule : 
$\text{Card\,} A\cup B = \text{Card\,} A + \text{Card\,} B - \text{Card\,} A\cap B$.
\finindication
\indication{000237}
Combien y-a-t'il de choix pour l'\'el\'ement de $A$ ?
Combien y-a-t'il de choix pour le sous-ensemble de $E\setminus A$ ?
\finindication
\noindication
\indication{000239}
Petits rappels : dans un jeu de $52$ cartes il y a $4$ ``couleurs'' (pique, c{\oe}ur, carreau, trèfle)
et $13$ ``valeurs'' ($1 =$ As, $2,3,\ldots,10$, Valet, Dame, Roi).
Une ``main'' c'est juste choisir $5$ cartes parmi les $52$, l'ordre du choix n'important pas.
\finindication
\noindication
\noindication
\noindication
\noindication
\noindication
\noindication
\noindication
\noindication
\noindication
\indication{002918} 
Si $\{x_1, \dots, x_p\}$ est une telle partie avec $x_1 <
x_2 < \dots < x_p$, consid{\'e}rer l'ensemble $\{ x_1-1, \dots, x_p-p\}$.
\finindication
\noindication
\noindication
\noindication
\noindication
\noindication
\noindication
\noindication
\indication{005284}
Coder un chemin par un mot : $D$ pour droite, $H$ pour haut.
\finindication
\noindication
\noindication
\noindication
\noindication
\noindication
\noindication
\noindication
\noindication
\noindication
\noindication
\noindication
\noindication
\noindication
\noindication
\noindication
\noindication
\noindication
\noindication
\noindication
\noindication
\noindication
\noindication
\noindication
\noindication
\indication{000249}
Il ne faut surtout pas chercher \`a calculer $15!=1\times2\times3\times4\times\cdots\times15$, mais profiter du fait
qu'il est d\'ej\`a ``presque'' factoris\'e.
\finindication
\indication{000250}
Il faut travailler modulo $13$, tout d'abord r\'eduire $100$ modulo $13$.
Se souvenir que si $a\equiv b \pmod{13}$ alors $a^k\equiv b^k \pmod{13}$.
Enfin calculer ce que cela donne pour les exposants $k=1,2,3,\ldots$
en essayant de trouver une r\`egle g\'en\'erale.
\finindication
\indication{000251}
Attention le reste d'une division euclidienne est plus petit que le quotient !
\finindication
\noindication
\noindication
\indication{000254}
Utiliser les modulos (ici modulo $8$),
un entier est divisible par $8$ si et seulement si
il est \'equivalent \`a $0$ modulo $8$.
Ici vous pouvez commencer par calculer $7^n \pmod{8}$.
\finindication
\noindication
\noindication
\noindication
\noindication
\noindication
\noindication
\noindication
\noindication
\noindication
\noindication
\noindication
\noindication
\noindication
\noindication
\noindication
\noindication
\noindication
\noindication
\noindication
\noindication
\noindication
\noindication
\noindication
\noindication
\noindication
\noindication
\noindication
\noindication
\noindication
\noindication
\indication{000285}
\begin{enumerate}
  \item \'Ecrire $n=2p+1$.
  \item  \'Ecrire $n=2p$ et discuter selon que $p$ est pair ou impair.
  \item Utiliser la premi\`ere question.
  \item Par l'absurde supposer que cela s'\'ecrive comme un carr\'e, par exemple $a^2+b^2+c^2=n^2$
puis discuter selon que $n$ est pair ou impair.
\end{enumerate}
\finindication
\noindication
\noindication
\noindication
\noindication
\noindication
\noindication
\noindication
\noindication
\noindication
\noindication
\noindication
\noindication
\noindication
\noindication
\noindication
\noindication
\noindication
\noindication
\noindication
\noindication
\noindication
\noindication
\noindication
\noindication
\noindication
\noindication
\noindication
\noindication
\noindication
\noindication
\noindication
\noindication
\noindication
\noindication
\noindication
\noindication
\noindication
\noindication
\noindication
\noindication
\noindication
\noindication
\indication{000305}
Commencer par simplifier l'équation !
Ensuite trouver une solution particulière $(x_0,y_0)$
à l'aide de l'algorithme d'Euclide par exemple. Ensuite trouver
un expression pour une solution générale.
\finindication
\noindication
\noindication
\noindication
\noindication
\noindication
\noindication
\noindication
\noindication
\noindication
\noindication
\noindication
\noindication
\noindication
\noindication
\noindication
\noindication
\noindication
\noindication
\noindication
\noindication
\noindication
\noindication
\noindication
\noindication
\noindication
\noindication
\noindication
\noindication
\noindication
\noindication
\noindication
\noindication
\noindication
\noindication
\noindication
\noindication
\noindication
\noindication
\noindication
\noindication
\noindication
\noindication
\noindication
\noindication
\noindication
\noindication
\noindication
\noindication
\noindication
\noindication
\indication{000336}
Pour 1. utiliser l'\'egalit\'e 
$$x^b-1 = (x-1)(x^{b-1}+\cdots+x+1).$$

Pour 2. raisonner par contraposition et utiliser la question 1.

La question 3. est difficile ! Supposer $a\ge b$.
Commencer par 
montrer que $\pgcd(2^a-1,2^b-1) = \pgcd(2^a-2^b,2^b-1) = \pgcd(2^{a-b}-1,2^b - 1)$.
Cela vour permettra de comparer l'agorithme d'Euclide pour le calcul de $\pgcd(a,b)$ avec
l'algorithme d'Euclide pour le calcul de  $\pgcd(2^a-1,2^b-1)$.
\finindication
\indication{000337}
Raisonner par l'absurde et utiliser le lemme de Gauss.
\finindication
\noindication
\indication{000339}
\begin{enumerate}
  \item \'Ecrire $$C_p^i = \frac{p(p-1)(p-2)\ldots(p-(i+1))}{i!}$$
et utiliser le lemme de Gauss ou le lemme d'Euclide.
  \item Raisonner avec les modulos, c'est-\`a-dire prouver $a^p \equiv a \pmod{p}$.

\end{enumerate}
\finindication
\noindication
\indication{000341}
\begin{enumerate}
  \item Il faut \^etre tr\`es soigneux : $n$ est fix\'e une fois pour toute, la r\'ecurrence se fait sur $k \ge 1$.
  \item Utiliser la question pr\'ec\'edente avec $m=n+k$.
  \item Par l'absurde, supposer qu'il y a seulement $N$ nombres premiers, consid\'erer
$N+1$ nombres du type $F_i$. Appliquer le ``principe du tiroir'' : \emph{si vous avez $N+1$ chaussettes rang\'ees dans $N$ tiroirs alors il existe (au moins) un tiroir contenant (plus de) deux chaussettes.}
\end{enumerate}
\finindication
\noindication
\noindication
\noindication
\noindication
\noindication
\noindication
\noindication
\indication{000349}
Raisonner par contraposition (ou par l'absurde) : supposer que $n$ n'est pas de la forme $2^k$,
alors $n$ admet un facteur irr\'eductible $p>2$.
Utiliser aussi $x^p+1 = (x+1)(1-x+x^2-x^3+\ldots+x^{p-1})$ avec $x$ bien choisi.
\finindication
\noindication
\noindication
\noindication
\noindication
\noindication
\noindication
\noindication
\noindication
\noindication
\noindication
\noindication
\noindication
\noindication
\noindication
\noindication
\noindication
\noindication
\noindication
\noindication
\noindication
\indication{000001}
Pour se ``débarrasser'' d'un dénominateur écrivez $\frac{z_1}{z_2} = \frac{z_1}{z_2}\cdot \frac{\bar z_2}{\bar z_2} =  \frac{z_1 \bar z_2}{ |z_2|^2}$.
\finindication
\noindication
\indication{000003}
Il faut bien connaître ses formules trigonométriques.
En particulier si l'on connait $\cos(2\theta)$ ou $\sin(2\theta)$ on sait
calculer $\cos \theta$ et $\sin \theta$.
\finindication
\noindication
\noindication
\noindication
\noindication
\noindication
\noindication
\noindication
\indication{000011}
Passez à la forme trigonométrique.
Souvenez-vous des formules sur les produits de puissances :
$$e^{ia}e^{ib} = e^{i(a+b)}\text{ et  } e^{ia} / e^{ib} = e^{i(a-b)}.$$
\finindication
\noindication
\indication{000013}
Pour calculer un somme du type
$e^{iu}+e^{iv}$ il est souvent utile de factoriser par
$e^{i\frac{u+v}{2}}$.
\finindication
\noindication
\noindication
\noindication
\noindication
\noindication
\noindication
\indication{000020}
Utiliser la formule d'Euler pour faire appara\^{\i}tre des cosinus.
\finindication
\noindication
\noindication
\noindication
\noindication
\noindication
\noindication
\noindication
\noindication
\noindication
\noindication
\noindication
\noindication
\noindication
\indication{007209}
Remarquer que la conjugaison est involutive. On peut aussi utiliser la forme algébrique.
\finindication
\indication{007210}
Remarquer que si $z\in\C$ est solution, alors $z^2$ est forcément un nombre réel.
\finindication
\noindication
\noindication
\noindication
\noindication
\indication{000027}
Pour $z= a+ib$ on cherche $\omega = \alpha + i \beta$ tel que  $(\alpha +i \beta)^2 = a+ib$.
Développez et indentifiez. Utilisez aussi que $|\omega|^2 = |z|$.
\finindication
\noindication
\indication{000029}
Il s'agit de calculer les racines carrées de $\frac{1+i}{\sqrt{2}} =  e^{i\frac{\pi}{4}}$ de deux
fa\c{c}ons différentes.
\finindication
\noindication
\indication{000031}
Pour les équation du type $az^4+bz^2+c=0$, poser $Z=z^2$.
\finindication
\noindication
\noindication
\noindication
\noindication
\noindication
\noindication
\noindication
\noindication
\noindication
\noindication
\noindication
\noindication
\noindication
\noindication
\noindication
\noindication
\noindication
\noindication
\noindication
\noindication
\noindication
\noindication
\noindication
\noindication
\noindication
\indication{000047}
Calculer $(1-z)S_n$.
\finindication
\noindication
\noindication
\noindication
\noindication
\noindication
\noindication
\noindication
\noindication
\noindication
\noindication
\noindication
\noindication
\noindication
\noindication
\noindication
\noindication
\noindication
\noindication
\noindication
\noindication
\noindication
\noindication
\noindication
\noindication
\indication{000060}
Le premier ensemble est une droite le second est un cercle.
\finindication
\noindication
\noindication
\noindication
\noindication
\noindication
\noindication
\noindication
\noindication
\indication{000069}
Pour l'interprétation géométrique cherchez le parallélogramme.
\finindication
\noindication
\noindication
\noindication
\noindication
\noindication
\noindication
\noindication
\noindication
\noindication
\noindication
\noindication
\noindication
\noindication
\noindication
\noindication
\noindication
\noindication
\noindication
\noindication
\noindication
\noindication
\noindication
\noindication
\noindication
\noindication
\noindication
\noindication
\noindication
\noindication
\noindication
\noindication
\noindication
\noindication
\noindication 
\noindication      
\indication{007147}   
Les deux similitudes se déduisent l'une de l'autre par une certaine réflexion.
\finindication      
\noindication   
\noindication   
\noindication   
\noindication   
\noindication 
\noindication
\noindication
\noindication
\noindication
\indication{000080}
Appliquer deux fois la formule de Moivre en remarquant
$e^{i5\theta}=(e^{i\theta})^5$.
\finindication
\noindication
\noindication
\noindication
\noindication
\noindication
\noindication
\noindication
\noindication
\noindication
\noindication
\noindication
\noindication
\noindication
\noindication
\noindication
\noindication
\noindication
\noindication
\noindication
\noindication
\noindication
\noindication
\noindication
\noindication
\noindication
\noindication
\noindication
\noindication
\noindication
\noindication
\noindication
\noindication
\noindication
\noindication
\noindication
\noindication
\noindication
\noindication
\noindication
\noindication
\noindication
\noindication
\noindication
\noindication
\noindication
\noindication
\noindication
\noindication
\noindication
\noindication
\noindication
\noindication
\noindication
\noindication
\noindication
\noindication
\noindication
\noindication
\noindication
\noindication
\noindication
\noindication
\noindication
\noindication
\noindication
\noindication
\noindication
\noindication
\noindication
\noindication
\noindication
\noindication
\noindication
\noindication
\noindication
\noindication
\noindication
\noindication
\noindication
\noindication
\indication{000378}
Si $P=P'Q$ avec $P\not=0$, regarder le degré de $Q$.
\finindication
\noindication
\noindication
\noindication
\noindication
\noindication
\noindication
\noindication
\noindication
\noindication
\noindication
\noindication
\noindication
\noindication
\noindication
\noindication
\noindication
\noindication
\noindication
\noindication
\noindication
\noindication
\noindication
\noindication
\noindication
\noindication
\noindication
\noindication
\noindication
\noindication
\noindication
\noindication
\noindication
\noindication
\noindication
\noindication
\noindication
\noindication
\noindication
\noindication
\indication{006957}
Le calcul du pgcd se fait par l'algorithme d'Euclide, et la "remontée" 
de l'algorithme permet d'obtenir $U$ et $V$.
\finindication
\indication{006958}
Calculer $\pgcd(P,P')$.
\finindication
\noindication
\noindication
\noindication
\noindication
\noindication
\noindication
\noindication
\noindication
\noindication
\noindication
\noindication
\noindication
\noindication
\noindication
\noindication
\noindication
\noindication
\noindication
\noindication
\noindication
\noindication
\noindication
\noindication
\noindication
\noindication
\noindication
\noindication
\noindication
\noindication
\noindication
\noindication
\noindication
\noindication
\noindication
\noindication
\noindication
\noindication
\noindication
\noindication
\noindication
\noindication
\noindication
\noindication
\noindication
\noindication
\noindication
\noindication
\noindication
\noindication
\noindication
\noindication
\noindication
\noindication
\noindication
\noindication
\noindication
\noindication
\noindication
\noindication
\noindication
\noindication
\noindication
\noindication
\noindication
\noindication
\noindication
\noindication
\indication{003253} 
Si $X^p-a = PQ$ avec $P,Q\in K[X]$ unitaires non
constants, factoriser $P$ dans~$\C$ et consid{\'e}rer $P(0)$.
\finindication
\noindication
\noindication
\noindication
\noindication
\noindication
\noindication
\noindication
\noindication
\noindication
\noindication
\noindication
\noindication
\indication{006960}
Montrer que si $P$ est un polynôme non constant vérifiant la relation, 
alors ses seules racines possibles sont $0$ et $1$.
\finindication
\indication{006961}
Pour l'existence, preuve par récurrence sur $n$. Pour les racines,
montrer que $P(x)=2\cos(n\Arccos(x/2))$.
\finindication
\noindication
\noindication
\noindication
\noindication
\indication{000445}
Attention il y a une partie enti\`ere, la fraction s'\'ecrit 
$$\Phi = x+1+{4x^2-6x+1\over2x^3-x^2}.$$
\finindication
\indication{000446}
Il y a une partie enti\`ere qui vaut $2$.
\finindication
\noindication
\noindication
\noindication
\noindication
\noindication
\noindication
\noindication
\noindication
\noindication
\noindication
\noindication
\noindication
\noindication
\noindication
\noindication
\noindication
\noindication
\noindication
\noindication
\noindication
\noindication
\noindication
\noindication
\noindication
\noindication
\noindication
\noindication
\noindication
\noindication
\noindication
\noindication
\noindication
\noindication
\noindication
\noindication
\noindication
\noindication
\noindication
\indication{006964}
\'Ecrire $F=\frac{P}{Q}$ sous forme irréductible.
\finindication
\indication{006965}
\'Ecrire $G=\frac{A}{B}$ sous forme irréductible (on pourra choisir par exemple 
$n=\mathrm{max}(\mathrm{deg}A,\mathrm{deg}B)$).
\finindication
\indication{006966}
Considérer $P'/P$ et sa dérivée, et enfin $P''/P$.
\finindication
\indication{006967}
Pour $G$ et $H$, commencer par faire une division euclidienne 
pour trouver la partie polynomiale.
\finindication
\indication{006968}
Les fractions $F, K$ ont une partie polynomiale, elles s'écrivent

$F=X^2+X+1+\frac{X^2+X+1}{X^3-X}$

$K=X+1+\frac{4X^2-6X+1}{2X^3-X^2}$
\finindication
\indication{006969}
Pour $F$, commencer par écrire $F=\frac{a}{X}+F_1$ où $F_1=\frac{N}{(X^2+1)^3}$ puis diviser $N$ par $X^2+1$.
Pour $K$, commencer par obtenir $K=1+\frac{1}{X}+K_1$, puis faire le changement d'indéterminée dans $K_1$.
\finindication
\noindication
\noindication
\indication{006972}
Pour 1. exprimer $\cos\big((n+2)\theta\big)$ et $\cos(n\theta)$ 
en fonction de $\cos\big((n+1)\theta\big)$.
Pour 3. chercher les racines de $T_n$ : 
$\omega_k=\cos\left(\frac{(2k+1)\pi}{2n}\right)$ pour $k=0,\ldots,n-1$.
\finindication
\noindication
\noindication
\noindication
\noindication
\noindication
\noindication
\noindication
\noindication
\noindication
\noindication
\noindication
\noindication
\noindication
\noindication
\noindication
\noindication
\noindication
\noindication
\noindication
\noindication
\noindication
\noindication
\noindication
\noindication
\noindication
\noindication
\noindication
\noindication
\noindication
\noindication
\noindication
\noindication
\noindication
\noindication
\noindication
\noindication
\noindication
\noindication
\noindication
\noindication
\noindication
\noindication
\noindication
\noindication
\noindication
\noindication
\noindication
\noindication
\noindication
\noindication
\noindication
\noindication
\noindication
\noindication
\noindication
\noindication
\noindication
\noindication
\noindication
\noindication
\noindication
\noindication
\noindication
\noindication
\noindication
\noindication
\noindication
\noindication
\noindication
\noindication
\noindication
\noindication
\noindication
\noindication
\noindication
\noindication
\noindication
\noindication
\noindication
\noindication
\noindication
\noindication
\noindication
\noindication
\indication{000886}
\begin{enumerate}
\item $E_1$ est un sous-espace vectoriel.
\item $E_2$ n'est pas un sous-espace vectoriel.
\item $E_3$ est un sous-espace vectoriel.
\item $E_4$ n'est pas un sous-espace vectoriel.
\end{enumerate}
\finindication
\noindication
\indication{000888}
\begin{enumerate}
\item $E_1$ est un sous-espace vectoriel de $\Rr^3$ si et seulement si
  $a=0$.
\item $E_2$ est un sous-espace vectoriel.
\item $E_3$ n'est pas un espace vectoriel.
\item $E_4$ n'est pas un espace vectoriel.
\end{enumerate}
\finindication
\noindication
\noindication
\noindication
\noindication
\indication{000893}
\begin{enumerate}
\item Pour le sens $\Rightarrow$ : raisonner par l'absurde et prendre
  un vecteur de $F\setminus G$ et un de $G\setminus F$. Regarder la
  somme de ces deux vecteurs.
\item Raisonner par double inclusion, revenir aux vecteurs.
\end{enumerate}
\finindication
\noindication
\noindication
\noindication
\noindication
\noindication
\noindication
\noindication
\noindication
\noindication
\noindication
\noindication
\noindication
\noindication
\noindication
\noindication
\noindication
\noindication
\noindication
\noindication
\noindication
\noindication
\noindication
\noindication
\noindication
\noindication
\noindication
\noindication
\indication{006868}
On v\'erifiera sur ces exemples la d\'efinition donn\'ee en cours.
\finindication
\indication{006869}
\begin{enumerate}
  \item Discuter suivant la dimension des sous-espaces.
  \item Penser aux droites vectorielles.
\end{enumerate}

\finindication
\noindication
\noindication
\noindication
\indication{000900}
On ne peut pas pour le premier, mais on peut pour le second.
\finindication
\indication{000901}
$E$ est un sous-espace vectoriel de $\Rr^4$. Une base comporte trois vecteurs.
\finindication
\noindication
\noindication
\noindication
\noindication
\noindication
\noindication
\indication{000908}
Montrer la double inclusion. 
Utiliser le fait que de mani\`ere g\'en\'erale pour  $E=\mathrm{Vect}(v_1,\ldots,v_n)$
alors : 
$$E\subset F \iff \forall i=1,\ldots,n \quad v_i\in F.$$
\finindication
\noindication
\noindication
\noindication
\noindication
\noindication
\noindication
\noindication
\indication{000916}
Supposer qu'il existe des r\'eels $\lambda_1,\ldots, \lambda_n$
et des indices  $\alpha_1,\ldots,\alpha_n$ (tout cela en nombre fini !)
tels que 
$$\lambda_1f_{\alpha_1}+\cdots+\lambda_nf_{\alpha_n} = 0.$$
Ici le $0$ est la fonction constante \'egale \`a $0$. \'Evaluer cette expression en
des valeurs bien choisies.
\finindication
\indication{000917}
Supposer qu'il existe des r\'eels $\lambda_1,\ldots, \lambda_n$
et des indices  $\alpha_1 > \alpha_2 > \cdots > \alpha_n$ (tout cela en nombre fini !)
tels que 
$$\lambda_1f_{\alpha_1}+\cdots+\lambda_nf_{\alpha_n} = 0.$$
Ici le $0$ est la fonction constante \'egale \`a $0$. 
Regarder quel terme est dominant et factoriser.
\finindication
\noindication
\noindication
\noindication
\noindication
\noindication
\noindication
\noindication
\noindication
\noindication
\noindication
\noindication
\noindication
\noindication
\noindication
\noindication
\noindication
\noindication
\noindication
\noindication
\noindication
\noindication
\noindication
\indication{006870}
\begin{enumerate}
  \item On pensera \`a poser un syst\`eme.
  \item Trouver un vecteur non-nul commun aux deux plans.
\end{enumerate}
\finindication
\noindication
\indication{000919}
\begin{enumerate}
\item Vrai.
\item Vrai.
\item Faux.
\item Faux.
\item Vrai.
\end{enumerate}
\finindication
\indication{000920}
\begin{enumerate}
\item Non.
\item Oui.
\item Non.
\item Non. 
\end{enumerate}
\finindication
\noindication
\noindication
\indication{000923}
  Soit
  $$G= \left\lbrace x \mapsto ax+b \mid (a,b) \in \Rr^2 \right\rbrace.$$
  Montrer que $G$ est un suppl\'ementaire de $F$ dans $E$.
\finindication
\noindication
\noindication
\indication{000926}
Pour une suite $(u_n)$ qui converge vers $\ell$ regarder la suite $(u_n-\ell)$.
\finindication
\noindication
\noindication
\noindication
\noindication
\noindication
\noindication
\noindication
\noindication
\noindication
\noindication
\noindication
\noindication
\noindication
\indication{006871}
\begin{enumerate}
  \item Jamais.
  \item Jamais.
  \item Considérer un vecteur directeur de la droite.
\end{enumerate}
\finindication
\noindication
\noindication
\noindication
\indication{000981}
\^Etre une base, c'est être libre et génératrice.
Chacune de ces conditions se vérifie par un système linéaire.
\finindication
\noindication
\noindication
\noindication
\indication{000985}
\begin{enumerate}
  \item Faux.
  \item Vrai.
\end{enumerate}
\finindication
\noindication
\noindication
\noindication
\noindication
\noindication
\noindication
\indication{000992}
Il suffit de montrer que la famille est libre (pourquoi ?). 
Prendre ensuite une combinaison linéaire nulle et regarder le terme de plus haut degré.
\finindication
\noindication
\noindication
\noindication
\indication{000996}
C'est une base pour $t\neq \pm 1$.
\finindication
\noindication
\noindication
\noindication
\noindication
\noindication
\noindication
\noindication
\noindication
\noindication
\indication{001006}
Il n'y a aucune difficulté. C'est comme dans $\Rr^3$ sauf qu'ici les coefficients sont des nombres complexes.
\finindication
\noindication
\noindication
\noindication
\noindication
\noindication
\noindication
\noindication
\noindication
\noindication
\noindication
\noindication
\noindication
\noindication
\noindication
\noindication
\noindication
\indication{001015}
Partir d'une base $(e_1,\ldots, e_k)$ de $F\cap G$ et la compl\'eter par des vecteurs $(f_1,\ldots,f_\ell)$ en une base de $F$.
Repartir de $(e_1,\ldots, e_k)$ pour la compléter par  des vecteurs $(g_1,\ldots,g_m)$ en une base de $G$.
Montrer que $(e_1,\ldots, e_k,f_1,\ldots,f_\ell,g_1,\ldots,g_m)$ est une base de $F+G$.
\finindication
\indication{001016}
On peut utiliser des familles libres.
\finindication
\noindication
\noindication
\indication{001019}
Calculer d'abord les dimensions de $F$ et $G$.
Pour celles de $F\cap G$ et $F+G$ servez-vous de la formule $\dim (F+G) = \dim F + \dim G - \dim (F\cap G)$.
\finindication
\noindication
\noindication
\noindication
\noindication
\noindication
\noindication
\noindication
\noindication
\noindication
\noindication
\noindication
\noindication
\noindication
\noindication
\noindication
\noindication
\noindication
\noindication
\noindication
\noindication
\noindication
\noindication
\noindication
\noindication
\noindication
\noindication
\noindication
\noindication
\noindication
\noindication
\noindication
\noindication
\indication{000929}
Une seule application n'est pas lin\'eaire.
\finindication
\indication{000930}
Prendre une combinaison lin\'eaire nulle et l'\'evaluer par $\phi^{n-1}$.
\finindication
\noindication
\noindication
\noindication
\noindication
\noindication
\noindication
\noindication
\noindication
\noindication
\noindication
\noindication
\noindication
\noindication
\indication{000934}
Faire un dessin de l'image et du noyau pour $f: \Rr\times \Rr \longrightarrow \Rr$.
Montrer que le noyau est isomorphe à $E_1 \cap E_2$.
\finindication
\noindication
\noindication
\noindication
\noindication
\noindication
\noindication
\noindication
\noindication
\indication{000943}
Pour chacune des implications utiliser la formule du rang.
\finindication
\noindication
\noindication
\noindication
\indication{000947}
Dire qu'un sous-espace $F$ est stable par $g$ signifie que $g(F) \subset F$.
\finindication
\noindication
\indication{000949}
Montrer la double inclusion.
\finindication
\noindication
\noindication
\noindication
\noindication
\indication{000954}
$t=0$ est un cas à part.
\finindication
\noindication
\noindication
\noindication
\noindication
\indication{000959}
Résultats utiles d'arithmétique des polynômes : la division euclidienne, le théorème de Bézout,
le lemme de Gauss.
\finindication
\noindication
\noindication
\noindication
\indication{000963}
Pour une base $\mathcal{B} =\{ e_1, \ldots , e_n \}$ de $E$ considérer la famille 
$\{ \phi (e_1), \ldots, \phi (e_n) \} $.
\finindication
\noindication
\noindication
\noindication
\noindication
\noindication
\noindication
\noindication
\noindication
\noindication
\noindication
\noindication
\noindication
\noindication
\noindication
\noindication
\noindication
\noindication
\noindication
\noindication
\noindication
\noindication
\noindication
\noindication
\noindication
\noindication
\noindication
\noindication
\noindication
\noindication
\noindication
\noindication
\noindication
\noindication
\noindication
\noindication
\noindication
\noindication
\noindication
\noindication
\noindication
\noindication
\noindication
\noindication
\noindication
\noindication
\noindication
\noindication
\noindication
\noindication
\indication{000974}
Pour une fonction $f$ on peut \'ecrire 
$$f(x)= \frac{f(x)+f(-x)}{2}+\frac{f(x)-f(-x)}{2}.$$

Le projecteur sur $P$ de direction $I$ est l'application $\pi : E \longrightarrow E$
qui vérifie $\pi(f)\in P$, $\pi \circ \pi = \pi$ et $\ker \pi = I$.
\finindication
\noindication
\indication{000976}
$P'$ désigne la dérivée de $P$.
Pour trouver le noyau, résoudre une équation différentielle.
Pour l'image calculer les $f(X^k)$.
\finindication
\noindication
\noindication
\noindication
\noindication
\noindication
\noindication
\noindication
\noindication
\noindication
\noindication
\noindication
\noindication
\noindication
\noindication
\noindication
\noindication
\noindication
\noindication
\noindication
\noindication
\noindication
\noindication
\noindication
\noindication
\noindication
\noindication
\noindication
\noindication
\noindication
\noindication
\noindication
\noindication
\noindication
\noindication
\noindication
\noindication
\noindication
\noindication
\noindication
\noindication
\noindication
\noindication
\noindication
\noindication
\noindication
\noindication
\noindication
\noindication
\noindication
\noindication
\noindication
\noindication
\noindication
\noindication
\noindication
\noindication
\noindication
\noindication
\indication{001052}
Une fois que l'on a calculé $A^2$ et $A^3$ on peut en déduire $A^{-1}$ sans calculs.
\finindication
\noindication
\noindication
\noindication
\noindication
\noindication
\noindication
\noindication
\indication{001061}
Il faut connaître les formules de $\cos(\theta+\theta')$ et $\sin(\theta+\theta')$.
\finindication
\noindication
\indication{001063}
Essayer avec $X$ la matrice élémentaire $E_{ij}$ (des zéros partout sauf le coefficient $1$ à la
 $i$-ème ligne et la $j$-ème colonne).
\finindication
\indication{001064}
Appliquer la formule du produit pour calculer les coefficients diagonaux de $A\ {}^{t}\!{A}$
\finindication
\noindication
\noindication
\noindication
\noindication
\indication{001069}
Prendre un vecteur $X=\left(
\begin{array}{c}
x_1\\
\vdots\\
x_n
\end{array}
\right)$ tel que $AX=0$, considérer le rang $i_0$ tel $|x_{i_0}|=\max\big\{|x_i| \mid i=1,...,n\big\}$.
\finindication
\noindication
\noindication
\noindication
\noindication
\noindication
\noindication
\noindication
\noindication
\noindication
\noindication
\noindication
\noindication
\noindication
\indication{002444}
$A$ est \emph{idempotente} s'il existe un $n$ tel que $A^n=I$ (la matrice identité).

$A$ est \emph{nilpotente} s'il existe un $n$ tel que $A^n=(0)$ (la matrice nulle).
\finindication
\noindication
\noindication
\noindication
\noindication
\noindication
\noindication
\noindication
\noindication
\noindication
\noindication
\noindication
\noindication
\noindication
\noindication
\noindication
\noindication
\noindication
\noindication
\noindication
\noindication
\noindication
\noindication
\noindication
\noindication
\noindication
\noindication
\noindication
\noindication
\noindication
\noindication
\noindication
\noindication
\noindication
\noindication
\noindication
\noindication
\noindication
\noindication
\noindication
\noindication
\indication{001087}
$f$ est l'application qui à $\begin{pmatrix}x\\y\end{pmatrix}$ associe
 $\begin{pmatrix}x-y\\0\end{pmatrix}$.
\finindication
\noindication
\noindication
\noindication
\noindication
\noindication
\noindication
\noindication
\noindication
\noindication
\noindication
\noindication
\noindication
\noindication
\noindication
\noindication
\noindication
\noindication
\noindication
\noindication
\noindication
\noindication
\noindication
\noindication
\noindication
\noindication
\noindication
\noindication
\noindication
\noindication
\noindication
\noindication
\noindication
\indication{003380}
$M$ antisymétrique signifie ${}^{t}\!{M}=-M$.
\begin{enumerate}
  \item Si $Y$ est un vecteur alors ${}^t\!Y Y = \| Y\|^2$ est un réel positif ou nul.
  \item $I-M$ et $(I+M)^{-1}$ commutent.
\end{enumerate}
\finindication
\noindication
\noindication
\noindication
\noindication
\noindication
\noindication
\noindication
\noindication
\noindication
\noindication
\noindication
\noindication
\noindication
\noindication
\noindication
\noindication
\noindication
\noindication
\noindication
\noindication
\noindication
\noindication
\noindication
\noindication
\noindication
\noindication
\noindication
\noindication
\noindication
\noindication
\noindication
\indication{002475}
Il faut trouver les propriétés de l'application linéaire $f$ associée à chacune de ces matrices.
Les résultats s'expriment en explicitant une (ou plusieurs) matrice $M'$ qui est la matrice de $f$ dans une base bien choisie
et ensuite en montrant que toutes les autres matrices sont de la forme
$M=P^{-1}M'P$.

Plus en détails pour chacun des cas :
\begin{enumerate}
  \item $\Im f \subset \Ker f$ et discuter suivant la dimension du noyau.
  \item Utiliser l'exercice \ref{exo1093} : $\Ker f \oplus \Im f$ et il existe une base telle que 
$f(e_i)=0$ ou $f(e_i)=e_i$.
  \item Poser $N= \frac{I+M}{2}$ (et donc $M=\cdots$) chercher à quelle condition $M^2=I$.
\end{enumerate}
\finindication
\noindication
\noindication
\noindication
\noindication
\noindication
\noindication
\noindication
\noindication
\noindication
\noindication
\noindication
\noindication
\noindication
\noindication
\noindication
\noindication
\noindication
\noindication
\noindication
\noindication
\noindication
\noindication
\noindication
\noindication
\noindication
\noindication
\noindication
\noindication
\noindication
\noindication
\noindication
\noindication
\noindication
\noindication
\noindication
\noindication
\noindication
\noindication
\noindication
\noindication
\noindication
\noindication
\noindication
\indication{000451}
\begin{enumerate}
  \item Raisonner par l'absurde.
  \item Raisonner par l'absurde en \'ecrivant $\sqrt2=\frac pq$ avec $p$ et $q$ premiers entre eux. Ensuite plusieurs méthodes sont possibles par exemple essayer de montrer que $p$ et $q$ sont tous les deux pairs.
  \item Considérer $r + \frac{\sqrt 2}{2}(r'-r)$ (faites un dessin !) pour deux rationnels $r,r'$. Puis utiliser les deux questions pr\'ec\'edentes.

\end{enumerate}
\finindication
\noindication
\noindication
\noindication
\noindication
\noindication
\indication{000457}
\begin{enumerate}
  \item Calculer $\beta^n p(\frac \alpha \beta)$ et utiliser le lemme de Gauss.
  \item Utiliser la premi\`ere question avec $p(x)=(x^2-5)^2-24$.

\end{enumerate}
\finindication
\noindication
\indication{000459}
\begin{enumerate}
  \item Mutiplier $N_n$ par une puissance de $10$ suffisament grande pour obtenir un nombre entier.
  \item Mutiplier $M$ par une puissance de $10$ suffisament grande (pas trop grande) puis soustraire $M$ pour obtenir un nombre entier.
\end{enumerate}
\finindication
\noindication
\indication{000461}
Raisonner par l'absurde !
\finindication
\noindication
\noindication
\indication{000506}
\'Ecrire la d\'efinition de la convergence d'une suite $(u_n)$ avec les ``$\epsilon$''. Comme on a une proposition qui est vraie pour tout $\epsilon >0$, c'est en particulier vrai pour $\epsilon =1$. Cela nous donne un ``$N$''. Ensuite s\'eparez la suite en deux : regardez les $n<N$ (il n'y a qu'un nombre fini de termes) et les $n\geqslant N$ (pour lequel on utilise notre $\epsilon=1$).
\finindication
\indication{000507}
On prendra garde \`a ne pas parler de limite d'une suite sans savoir au pr\'ealable qu'elle converge !

Vous pouvez utiliser le r\'esultat du cours suivant :
Soit $(u_n)$ une suite convergeant  vers la limite $\ell$ alors toute sous-suite $(v_n)$ de $(u_n)$ a pour limite $\ell$.
\finindication
\noindication
\noindication
\noindication
\noindication
\noindication
\noindication
\noindication
\noindication
\noindication
\noindication
\noindication
\noindication
\noindication
\noindication
\noindication
\noindication
\indication{000465}
$\inf A =0$, $A$ n'a pas de borne supérieure.
\finindication
\noindication
\noindication
\noindication
\noindication
\noindication
\noindication
\noindication
\noindication
\noindication
\noindication
\indication{000476}
Il faut revenir à la définition de la borne supérieure d'un ensemble borné :
c'est le plus petit des majorants. En particulier la borne supérieure est un majorant.
\finindication
\indication{000477}
Deux propositions sont fausses...
\finindication
\noindication
\noindication
\noindication
\noindication
\noindication
\noindication
\noindication
\noindication
\noindication
\noindication
\noindication
\noindication
\noindication
\noindication
\noindication
\noindication
\noindication
\indication{000491}
\'Elever l'in\'egalit\'e au carr\'e.
\finindication
\noindication
\noindication
\noindication
\noindication
\noindication
\indication{000497}
\begin{enumerate}
  \item $f(2)= f(1+1)= \cdots$, faire une r\'ecurrence.
  \item $f((-n)+n)=\cdots$.
  \item Si $q = \frac ab$, calculer $f(\frac ab + \frac ab + \cdots +\frac ab)$
avec $b$ termes dans cette somme.
  \item Utiliser la densit\'e
de $\Qq$ dans $\Rr$ : pour $x\in\Rr$ fix\'e, prendre une suite de rationnels qui croit vers $x$, et une autre qui d\'ecroit vers $x$.
\end{enumerate}
\finindication
\noindication
\noindication
\noindication
\noindication
\noindication
\noindication
\noindication
\noindication
\noindication
\noindication
\noindication
\noindication
\noindication
\noindication
\noindication
\noindication
\noindication
\noindication
\noindication
\noindication
\noindication
\noindication
\noindication
\noindication
\indication{005982}
\begin{enumerate}
 \item Rappelez-vous que la partie entière de $x$ est le plus grand entier, inférieur ou égal à $x$. Mais il est ici préférable de donner la définition de $E(x)$ en disant que $E(x) \in \Zz$ et que $x$ vérifie un certain encadrement...

 \item Encadrer $E(kx)$, pour $k=1,\ldots,n$.

 \item Rappelez-vous d'abord de la formule $1+2+\cdots+n$ puis utilisez le fameux théorème des gendarmes.

 \item Les $u_n$ ne seraient-ils pas des rationnels ? 
\end{enumerate}
\finindication
\indication{007169} 
Utiliser l'inégalité arithmético-géométrique.
\finindication
\indication{007170} 
Utiliser l'inégalité arithmético-géométrique.
\finindication
\indication{007171} 
Utiliser l'inégalité arithmético-géométrique.
\finindication
\indication{007172} 
Utiliser l'inégalité arithmético-géométrique.
\finindication
\indication{007173} 
Utiliser l'inégalité arithmético-géométrique.
\finindication
\indication{007174} 
Utiliser l'inégalité arithmético-géométrique.
\finindication
\indication{007175} 
Utiliser l'inégalité arithmético-géométrique.
\finindication
\indication{007176} 
Utiliser l'inégalité arithmético-géométrique.
\finindication
\indication{007177} 
Utiliser l'inégalité arithmético-géométrique.
\finindication
\indication{007178} 
Utiliser l'inégalité arithmético-géométrique.
\finindication
\indication{007179} 
Utiliser l'inégalité arithmético-géométrique.
\finindication
\indication{007180} 
Utiliser l'inégalité arithmético-géométrique.
\finindication
\indication{007181} 
$2=1+1$...
\finindication
\indication{007182} 
Utiliser l'inégalité arithmético-géométrique.
\finindication
\indication{007183} 
Utiliser l'inégalité arithmético-géométrique.
\finindication
\indication{007184} 
Utiliser l'inégalité arithmético-géométrique sur des groupes de termes.
\finindication
\indication{007185} 
Utiliser l'inégalité arithmético-géométrique.
\finindication
\noindication
\indication{000505}
Dans l'ordre c'est vrai, faux et vrai. Lorsque c'est faux chercher un contre-exemple, lorsque c'est vrai il faut le prouver.
\finindication
\noindication
\noindication
\noindication
\noindication
\noindication
\noindication
\noindication
\indication{000519}
\'Ecrire la convergence de la suite et fixer $\epsilon = \frac 12$.
Une suite est \emph{stationnaire} si, à partir d'un certain rang, elle est constante.
\finindication
\indication{000520}
\begin{enumerate}
\item En se rappelant que l'int\'egrale calcule une aire montrer : 
$$\frac{1}{n+1} \leqslant \int_n^{n+1} \frac{dt}{t} \leqslant \frac 1n.$$
\item Pour chacune des majorations, il s'agit de faire la somme de l'in\'egalit\'e pr\'ec\'edente et de s'apercevoir que d'un cot\'e on calcule $H_n$ et de l'autre les termes s'\'eliminent presque tous deux \`a deux.
\item La limite est $+\infty$.
\item Calculer $u_{n+1}-u_n$.
\item Que fait une suite décroissante et minorée ?
\end{enumerate}
\finindication
\noindication
\noindication
\noindication
\indication{000524}
Pour la deuxi\`eme question, raisonner par l'absurde et trouver deux sous-suites ayant des limites distinctes.
\finindication
\noindication
\noindication
\noindication
\noindication
\noindication
\noindication
\noindication
\noindication
\noindication
\noindication
\noindication
\noindication
\noindication
\noindication
\noindication
\noindication
\noindication
\noindication
\noindication
\noindication
\noindication
\noindication
\noindication
\noindication
\noindication
\noindication
\noindication
\noindication
\noindication
\noindication
\noindication
\noindication
\noindication
\noindication
\noindication
\noindication
\noindication
\noindication
\noindication
\noindication
\noindication
\noindication
\noindication
\noindication
\noindication
\noindication
\noindication
\noindication
\noindication
\noindication
\noindication
\noindication
\noindication
\noindication
\noindication
\noindication
\noindication
\noindication
\noindication
\noindication
\noindication
\noindication
\noindication
\noindication
\noindication
\noindication
\noindication
\noindication
\noindication
\noindication
\noindication
\noindication
\noindication
\noindication
\noindication
\noindication
\noindication
\noindication
\noindication
\indication{000539}
Pour la premi\`ere question : attention on ne demande pas de calculer $\alpha$ !
L'existence vient du th\'eor\`eme des valeurs interm\'ediaires. L'unicit\'e vient du fait que la fonction est strictement croissante.
  

Pour la derni\`ere question : il faut d'une part montrer que $(x_n)$ converge et on note $\ell$ sa limite
et d'autre part il faut montrer que $\ell = \alpha$.
\finindication
\noindication
\noindication
\noindication
\noindication
\noindication
\noindication
\noindication
\noindication
\noindication
\noindication
\noindication
\noindication
\noindication
\noindication
\noindication
\noindication
\noindication
\noindication
\noindication
\noindication
\noindication
\noindication
\noindication
\noindication
\noindication
\noindication
\noindication
\noindication
\indication{004708}
\begin{enumerate}
\item
\item $\sqrt a - \sqrt b = \frac {a-b}{\sqrt
      a + \sqrt b}$
\end{enumerate}
\finindication
\noindication
\noindication
\noindication
\noindication
\noindication
\noindication
\noindication
\noindication
\noindication
\noindication
\noindication
\noindication
\noindication
\noindication
\noindication
\noindication
\noindication
\indication{000563}
Remarquer que $1-\frac{1}{k^2} = \frac{(k-1)(k+1)}{k.k}$.
Puis simplifier l'\'ecriture de $u_n$.
\finindication
\noindication
\noindication
\noindication
\noindication
\noindication
\indication{000569}
\begin{enumerate}
\item C'est un calcul de r\'eduction au m\^eme d\'enominateur.
\item Pour montrer la d\'ecroisance, montrer $\frac{u_{n+1}}{ u_n} \leqslant 1$.
\item Montrer d'abord que la suite converge, montrer ensuite que la limite est $\sqrt a$.
\item Penser \`a \'ecrire $u_{n+1}^2-a = (u_{n+1}-\sqrt a)(u_{n+1}+\sqrt a)$.
\item Raisonner par r\'ecurrence.
\item Pour $u_0= 3$ on a $u_1=3,166\ldots$, donc $3\leqslant \sqrt{10} \leqslant u_1$
et on peut prendre $k=0.17$ par exemple et $n=4$ suffit pour la pr\'ecision demand\'ee.
\end{enumerate}
\finindication
\indication{000570}
\begin{enumerate}
  \item Montrer que $(u_n)$ est croissante et $(v_n)$ d\'ecroissante.
  \item Montrer que $(u_n)$ est major\'ee et $(v_n)$ minor\'ee. Montrer que ces suites ont la m\^eme limite.
  \item Raisonner par l'absurde : si la limite $\ell = \frac pq$
alors multiplier l'in\'egalit\'e $u_q \leq \frac pq \leq v_q$  par $q!$ et raisonner avec des entiers.
\end{enumerate}
\finindication
\indication{000571}
Pour la premi\`ere question et la monotonie il faut raisonner par r\'ecurrence.
Pour la troisi\`eme question, remarquer que si $f$ est d\'ecroissante alors $f\circ f$ est croissante
et appliquer la premi\`ere question.
\finindication
\indication{000572}
\begin{enumerate}
  \item Regarder ce que donne l'inégalité en élevant au carré de chaque coté.
  \item Petites manipulations des inégalités.
  \item 
  \begin{enumerate}
     \item Utiliser 1.
     \item Utiliser 2.
     \item Une suite croissante et majorée converge ; une suite décroissante et minorée aussi.
  \end{enumerate}
\end{enumerate}
\finindication
\noindication
\indication{000574}
On notera  $f_n : [0,1] \longrightarrow \Rr$ la fonction
d\'efinie par $ f_n(x) = \sum_{k=1}^n x^k \ \ - 1.$
\begin{enumerate}
    \item C'est une \'etude de la fonction $f_n$.
    \item On sait que $f_n(a_n)=0$. Montrer par un calcul que $f_n(a_{n-1}) > 0$,
en d\'eduire la d\'ecroissance de $(a_n)$. En calculant $f_n(\frac 12)$ montrer que 
la suite  $(a_n)$ est minor\'ee par $\frac 12$.
    \item Une fois \'etablie la convergence de  $(a_n)$ vers une limite $\ell$,
composer l'in\'egalit\'e $ \frac 12 \leqslant \ell < a_n$ par $f_n$. Conclure.
\end{enumerate}
\finindication
\noindication
\noindication
\noindication
\noindication
\noindication
\noindication
\noindication
\noindication
\noindication
\noindication
\noindication
\noindication
\noindication
\noindication
\noindication
\noindication
\noindication
\noindication
\noindication
\noindication
\noindication
\noindication
\noindication
\noindication
\noindication
\noindication
\noindication
\noindication
\noindication
\noindication
\noindication
\noindication
\noindication
\noindication
\noindication
\noindication
\noindication
\noindication
\noindication
\noindication
\noindication
\noindication
\noindication
\noindication
\noindication
\noindication
\noindication
\noindication
\noindication
\noindication
\noindication
\noindication
\noindication
\noindication
\noindication
\noindication
\noindication
\noindication
\noindication
\noindication
\noindication
\noindication
\noindication
\noindication
\noindication
\noindication
\noindication
\noindication
\noindication
\noindication
\noindication
\noindication
\indication{001938}
\begin{enumerate}
\item Calculer $(1-3^{-1})\sum_{n \geq 0} (n+1)3^{-n}$.
\end{enumerate}
\finindication
\noindication
\noindication
\noindication
\noindication
\noindication
\noindication
\noindication
\noindication
\noindication
\noindication
\noindication
\noindication
\noindication
\noindication
\noindication
\noindication
\noindication
\noindication
\noindication
\noindication
\noindication
\noindication
\noindication
\noindication
\noindication
\noindication
\noindication
\noindication
\noindication
\noindication
\noindication
\noindication
\noindication
\noindication
\noindication
\noindication
\noindication
\noindication
\noindication
\noindication
\noindication
\noindication
\noindication
\noindication
\noindication
\noindication
\noindication
\noindication
\noindication
\noindication
\noindication
\noindication
\noindication
\noindication
\noindication
\noindication
\noindication
\noindication
\noindication
\noindication
\noindication
\noindication
\noindication
\noindication
\noindication
\noindication
\noindication
\noindication
\noindication
\noindication
\noindication
\noindication
\noindication
\noindication
\noindication
\noindication
\noindication
\noindication
\noindication
\noindication
\noindication
\noindication
\noindication
\noindication
\noindication
\noindication
\noindication
\noindication
\noindication
\noindication
\noindication
\noindication
\noindication
\noindication
\noindication
\noindication
\noindication
\noindication
\noindication
\noindication
\noindication
\noindication
\noindication
\noindication
\noindication
\noindication
\noindication
\noindication
\noindication
\noindication
\noindication
\noindication
\noindication
\noindication
\noindication
\noindication
\noindication
\noindication
\noindication
\noindication
\noindication
\noindication
\noindication
\noindication
\noindication
\noindication
\noindication
\noindication
\noindication
\noindication
\noindication
\noindication
\noindication
\noindication
\noindication
\noindication
\noindication
\noindication
\noindication
\noindication
\noindication
\noindication
\noindication
\noindication
\noindication
\noindication
\noindication
\noindication
\noindication
\noindication
\noindication
\noindication
\noindication
\noindication
\noindication
\noindication
\noindication
\noindication
\noindication
\noindication
\noindication
\indication{000639}
\begin{enumerate}
    \item On pourra utiliser la variante de l'in\'egalit\'e triangulaire $|x-y|\geq \big| |x|-|y| \big|$.
    \item Utiliser la premi\`ere question pour montrer que $|f-g|$ est continue.
\end{enumerate}
\finindication
\noindication
\noindication
\noindication
\noindication
\noindication
\indication{000645}
Ce n'est pas tr\`es dur mais il y a quand m\^eme quelque chose \`a d\'emontrer :
ce n'est pas parce que $f(x)$ vaut $+1$ ou $-1$ que la fonction est constante.
Raisonner par l'absurde et utiliser le th\'eor\`eme des valeurs interm\'ediaires.
\finindication
\indication{000646}
Il faut raisonner en deux temps : d'abord \'ecrire la d\'efinition de la limite en $+\infty$, en fixant par exemple $\epsilon =1$, cela donne une borne sur $[A,+\infty]$. Puis travailler sur $[0,A]$.
\finindication
\noindication
\noindication
\noindication
\noindication
\noindication
\noindication
\indication{000653}
Un \emph{point fixe} est une valeur $c \in [0,1]$ telle que $f(c)=c$.
Montrer que $c = \sup E$ est un point fixe. 
Pour cela montrer que $f(c) \leqslant c$ puis $f(c) \geqslant c$.
\finindication
\noindication
\noindication
\noindication
\noindication
\noindication
\noindication
\noindication
\noindication
\indication{000662}
Non, trouver un contre-exemple.
\finindication
\noindication
\noindication
\noindication
\noindication
\noindication
\noindication
\indication{000669}
\begin{enumerate}
 \item On pourra montrer que $\sup_{a\leq x \leq b} f(x)$ est
un majorant de $f$ sur $]a,b[$.

 \item Dans le cas $x_0=a$, par exemple, on pourra consid\'erer la suite de r\'eels
$a_n=a+1/n$ et \'etudier la suite $(f(a_n))$.
\end{enumerate}

\finindication
\noindication
\noindication
\noindication
\noindication
\noindication
\noindication
\noindication
\noindication
\noindication
\noindication
\noindication
\noindication
\noindication
\noindication
\noindication
\noindication
\noindication
\noindication
\noindication
\noindication
\noindication
\noindication
\noindication
\noindication
\noindication
\noindication
\noindication
\noindication
\noindication
\noindication
\noindication
\noindication
\noindication
\noindication
\noindication
\noindication
\noindication
\noindication
\noindication
\noindication
\noindication
\noindication
\noindication
\indication{000670}
Le ``$\epsilon$'' vous est donné, il ne faut pas y toucher.
Par contre c'est à vous de trouver le ``$\delta$''.
\finindication
\indication{000671}
Distinguer trois intervalles pour la formule définissant $f^{-1}$.
\finindication
\indication{000672}
Le seul probl\`eme est en $x=0$. Montrer que la fonction est bien continue en ce point.
\finindication
\noindication
\noindication
\noindication
\noindication
\indication{000677}
Oui pour le deux premi\`eres en posant $f(0)=0$, $g(0)=0$, non pour la troisi\`eme.
\finindication
\noindication
\noindication
\indication{000680}
Pour $x$ fix\'e, \'etudier la suite $f(\frac 1{2^n} x)$.
\finindication
\noindication
\noindication
\noindication
\noindication
\noindication
\noindication
\noindication
\noindication
\noindication
\indication{000609}
Utiliser l'expression conjugu\'ee.
\finindication
\noindication
\noindication
\indication{000612}
\begin{enumerate}
    \item Raisonner par l'absurde.
    \item Montrer que la limite est la borne sup\'erieure de l'ensemble des valeurs atteintes $f(\Rr)$.
\end{enumerate}
\finindication
\noindication
\noindication
\noindication
\indication{000616}
Réponses :
\begin{enumerate}
  \item La limite \`a droite vaut $+2$, la limite \`a gauche $-2$ donc il n'y a pas de limite.
  \item $-\infty$
  \item $4$
  \item $2$
  \item $\frac 12$
  \item $0$
  \item $\frac 13$ en utilisant par exemple que $a^3-1 = (a-1)(1+a+a^2)$ pour $a = \sqrt[3]{1+x^2}$.
  \item $\frac 1n$
\end{enumerate}
\finindication
\noindication
\noindication
\noindication
\noindication
\noindication
\noindication
\noindication
\noindication
\noindication
\noindication
\noindication
\indication{000628}
\begin{enumerate}
    \item Calculer d'abord la limite de $f(x) = \frac{x^k-\alpha^k}{x-\alpha}$.
    \item Utiliser $\cos 2x =  2\cos^2 x - 1$ et faire un changement de variable $u = \cos x$.
    \item Utiliser l'expression conjugu\'ee.
    \item Diviser num\'erateur et d\'enominateur par $\sqrt{x-\alpha}$ puis utiliser l'expression conjugu\'ee.
    \item On a toujours $y-1 \leq E(y) \leq y$, poser $y=1/x$.
    \item Diviser num\'erateur et d\'enominateur par $x-2$.
    \item  Pour $\alpha \geq 4$ il n'y a pas de  limite, pour $\alpha <4$ la limite est $+\infty$.
\end{enumerate}
\finindication
\noindication
\noindication
\noindication
\noindication
\noindication
\noindication
\indication{000635}
R\'{e}ponses : $0,\frac{1}{e},e.$
\begin{enumerate}
 \item Borner $\sin\frac 1x$.
 \item Utiliser que $\ln(1+t) = t \cdot \mu(t)$, pour une certaine fonction $\mu$ qui vérifie $\mu(t) \to 1$ lorsque $t\to 0$.
 \item Utiliser que $e^t-1 = t \cdot \mu(t)$, pour une certaine fonction $\mu$ qui vérifie $\mu(t) \to 1$ lorsque $t\to 0$.
\end{enumerate}

\finindication
\noindication
\indication{000637}
R\'{e}ponse: $\max(a,b)$.
\finindication
\indication{000638}
R\'{e}ponse: $\sqrt{ab}$.
\finindication
\noindication
\noindication
\noindication
\noindication
\noindication
\noindication
\noindication
\noindication
\noindication
\noindication
\noindication
\noindication
\noindication
\noindication
\noindication
\noindication
\noindication
\noindication
\noindication
\noindication
\noindication
\noindication
\noindication
\noindication
\noindication
\noindication
\noindication
\noindication
\noindication
\noindication
\noindication
\noindication
\noindication
\noindication
\noindication
\noindication
\noindication
\noindication
\noindication
\noindication
\noindication
\noindication
\noindication
\noindication
\noindication
\noindication
\noindication
\noindication
\noindication
\noindication
\noindication
\noindication
\noindication
\noindication
\noindication
\noindication
\noindication
\noindication
\noindication
\noindication
\noindication
\noindication
\noindication
\noindication
\noindication
\noindication
\noindication
\noindication
\noindication
\noindication
\noindication
\noindication
\noindication
\noindication
\noindication
\noindication
\noindication
\noindication
\noindication
\noindication
\indication{000698}
Les probl\`emes sont seulement en $0$ ou $1$. 
$f_1$ est d\'erivable en $0$ mais pas $f_2$.  $f_3$ n'est dérivable ni en $0$, ni en $1$.
\finindication
\indication{000699}
Vous avez deux conditions : il faut que la fonction soit continue (car on veut qu'elle soit d\'erivable donc elle doit \^etre continue) et ensuite la condition de d\'erivabilit\'e proprement dite.
\finindication
\indication{000700}
$f$ est continue en $0$ en la prolongeant par $f(0)=0$.
$f$ est alors d\'erivable en $0$ et $f'(0)=0$.
\finindication
\indication{000701}
On ne cherchera pas \`a utiliser la formule de Leibniz mais \`a lin\'eariser les expressions trigonom\'etriques.
\finindication
\noindication
\noindication
\noindication
\noindication
\noindication
\noindication
\noindication
\noindication
\noindication
\noindication
\noindication
\noindication
\noindication
\noindication
\noindication
\noindication
\noindication
\noindication
\noindication
\noindication
\noindication
\noindication
\indication{000715}
Il faut appliquer le th\'eor\`eme de Rolle une fois au polyn\^ome $(1-t^2)^n$,
puis deux fois \`a sa d\'eriv\'ee premi\`ere, puis trois fois \`a sa d\'eriv\'ee seconde,...
\finindication
\noindication
\indication{000717}
On peut appliquer le th\'eor\`eme de Rolle plusieurs fois.
\finindication
\indication{000718}
C'est encore Rolle de nombreuses fois
\finindication
\noindication
\noindication
\noindication
\noindication
\noindication
\indication{000724}
\begin{enumerate}
    \item Utiliser le th\'eor\`eme des accroissements finis avec la fonction $t \mapsto \ln t$
    \item Montrer d'abord que $f''$ est n\'egative. Se servir du th\'eor\`eme des valeurs interm\'ediaires pour $f'$.
\end{enumerate}
\finindication
\indication{000725}
Une fois le th\'eor\`eme des accroissements finis utilis\'e vous obtenez une somme t\'el\'escopique.
\finindication
\noindication
\indication{000727}
Le th\'eor\`eme des accroissements finis donne un r\'esultat proche de celui souhait\'e \`a un facteur pr\`es.
Pour obtenir la majoration demand\'ee on peut utiliser le th\'eor\`eme de Rolle avec une fonction bien choisie.
\finindication
\noindication
\noindication
\noindication
\noindication
\noindication
\noindication
\noindication
\noindication
\noindication
\noindication
\noindication
\noindication
\noindication
\noindication
\noindication
\noindication
\noindication
\noindication
\noindication
\noindication
\noindication
\noindication
\noindication
\noindication
\noindication
\noindication
\noindication
\noindication
\noindication
\noindication
\noindication
\noindication
\noindication
\noindication
\noindication
\noindication
\noindication
\noindication
\noindication
\noindication
\noindication
\noindication
\noindication
\noindication
\noindication
\indication{000734}
On distinguera les cas $\lambda \geq 0$ et $\lambda < 0$.
Pour le cas $\lambda < 0$ on considerera des sous-cas.
\finindication
\noindication
\noindication
\noindication
\indication{000738}
\begin{enumerate}
 \item Raisonner par l'absurde et appliquer le th\'eor\`eme de Rolle.
 \item Calculer $h(a)$ et $h(b)$.
 \item Appliquer la question 2. sur l'intervalle $[x,b]$.
 \item Calculer $f'$ et $g'$.
\end{enumerate}
\finindication
\noindication
\noindication
\noindication
\noindication
\noindication
\noindication
\noindication
\noindication
\noindication
\noindication
\noindication
\noindication
\noindication
\noindication
\noindication
\noindication
\noindication
\noindication
\noindication
\noindication
\noindication
\noindication
\noindication
\noindication
\noindication
\noindication
\noindication
\noindication
\noindication
\noindication
\noindication
\noindication
\noindication
\noindication
\noindication
\noindication
\noindication
\noindication
\noindication
\noindication
\noindication
\noindication
\noindication
\noindication
\noindication
\noindication
\noindication
\noindication
\noindication
\noindication
\noindication
\noindication
\noindication
\noindication
\noindication
\noindication
\noindication
\noindication
\noindication
\noindication
\noindication
\indication{001267}
Calculer d'abord le dl puis utiliser une formule de Taylor.
\finindication
\indication{001268}
\begin{enumerate}
  \item La formule à appliquer est celle de Taylor-Lagrange à l'ordre $2$.
  \item \'Etudier la fonction $\phi(h) = \frac{h}{2}M_2+\frac{2}{h}M_0$ et trouver $\inf_{h>0} \phi(h)$.
  \item Il faut choisir un $a>0$ tel que $g(x)$ soit assez petit sur $]a,+\infty[$ ; puis appliquer
les questions précédentes à $g$ sur cet intervalle.
\end{enumerate}
\finindication
\noindication
\noindication
\noindication
\noindication
\noindication
\noindication
\noindication
\noindication
\noindication
\noindication
\noindication
\noindication
\noindication
\noindication
\noindication
\noindication
\noindication
\noindication
\noindication
\noindication
\noindication
\noindication
\noindication
\noindication
\noindication
\noindication
\noindication
\indication{001243}
Pour la première question vous pouvez appliquer la formule de Taylor ou bien 
poser $h=x-1$ et considérer un dl au voisinage de $h=0$.
\finindication
\indication{001244}
En $x=0$ c'est le quotient de deux dl.
En $x=+\infty$, on pose $h=\frac1x$ et on calcule un dl en $h=0$.
\finindication 
\noindication
\noindication
\indication{002657}
Faites un développement faisant intervenir des $x$ et des $\ln x$.
Trouvez $\ell=1$.
\finindication
\noindication
\noindication
\noindication
\noindication
\noindication
\noindication
\noindication
\noindication
\noindication
\noindication
\noindication
\noindication
\noindication
\noindication
\noindication
\noindication
\noindication
\noindication
\indication{006888}
\begin{enumerate}
  \item $\cos x \cdot \exp x = 1 + x - \frac13 x^3 + o(x^3)$

  \item $\left( \ln (1+x) \right)^2= x^2-x^3+\frac{11}{12}x^4+ o(x^4)$

  \item $\frac{\sh x-x}{x^3} = \frac{1}{3!}+\frac{1}{5!}x^2+\frac{1}{7!}x^4+\frac{1}{9!}x^6 + o(x^6)$

  \item $\exp\big(\sin(x)\big)=1+x + \frac12 x^2 - \frac18 x^4+ o(x^4)$

  \item $\sin^6 (x)  =  x^6-x^8 + o(x^9)$

  \item $\ln(\cos x) =  - \frac{1}{2} x^2 -\frac{1}{12}x^4 -\frac{1}{45}x^6  + o(x^6)$

  \item $\frac{1}{\cos x} = 1+\frac{1}{2}x^2+\frac{5}{24}x^4 + o(x^4)$

  \item $\tan x = x + \frac{x^3}{3} + \frac{2x^5}{15} + \frac{17x^7}{315} + o(x^7)$

  \item $(1+x)^{\frac{1}{1+x}} =  \exp\left( \frac{1}{1+x} \ln(1+x) \right) = 1+x-x^2 + \frac{x^3}{2} + o(x^3)$

  \item  $\arcsin \left ( \ln(1+x^2) \right ) = x^2-\frac{x^4}{2}+\frac{x^6}{2}+o(x^6)$
\end{enumerate}
\finindication
\indication{001247}
Il s'agit bien sûr de calculer d'abord des dl afin d'obtenir la limite.
On trouve : 
\begin{enumerate}
  \item $\lim_{x\rightarrow 0}\frac{e^{x^2}-\cos x}{x^2} = \frac32$
  \item $\lim_{x\rightarrow 0}\frac{\ln (1+x)-\sin x}{x} = 0$
  \item $\lim_{x\rightarrow 0}\frac{\cos x-\sqrt{1-x^2}}{x^4}=\frac16$
\end{enumerate}
\finindication
\noindication
\indication{001249}
Faire un dl en $x=0$ à l'ordre $2$ cela donne $f(0)$, $f'(0)$ et la position par rapport à la tangente
donc tout ce qu'il faut pour répondre aux questions. Idem en $x=1$.
\finindication
\noindication
\noindication
\noindication
\noindication
\noindication
\noindication
\noindication
\noindication
\noindication
\noindication
\noindication
\noindication
\noindication
\noindication
\noindication
\indication{001265}
Il s'agit de faire un dl afin de trouver la limite.
\begin{enumerate}
  \item 
  \begin{enumerate}
    \item $\displaystyle \lim_{x \rightarrow +\infty} \sqrt {x^2+3x+2} +x = + \infty$
    \item $\displaystyle \lim_{x \rightarrow -\infty} \sqrt {x^2+3x+2} +x = -\frac32$
  \end{enumerate}
  \item $\displaystyle \lim_{x \rightarrow 0^+}(\Arctan x )^{\frac{1}{x^2}}=0$
  \item $\displaystyle \lim_{x \rightarrow 0} \frac{(1+3x)^{\frac{1}{3}}-1-\sin x}{1-\cos x}=-2$
\end{enumerate}

\finindication
\noindication
\noindication
\noindication
\noindication
\noindication
\noindication
\noindication
\noindication
\noindication
\noindication
\noindication
\noindication
\noindication
\noindication
\noindication
\noindication
\noindication
\noindication
\noindication
\indication{004045}
Identifier les dl de $\cos x$ et $\frac{1+ax^2}{1+bx^2}$ en $x=0$.
\finindication
\noindication
\noindication
\noindication
\noindication
\noindication
\noindication
\noindication
\noindication
\noindication
\noindication
\noindication
\noindication
\noindication
\noindication
\noindication
\noindication
\noindication
\noindication
\noindication
\indication{000745}
Faire un dessin. Calculer l'angle d'observation $\alpha$ en fonction de la distance $x$
et étudier cette fonction. Pour simplifier l'expression de $\alpha_0$, 
calculer $\tan\alpha_0$ à l'aide de la formule donnant $\tan(a-b)$.
\finindication
\indication{000746}
On pourra \'etudier les fonctions d\'efinies par la diff\'erence des deux termes de chaque in\'egalit\'e.
\finindication
\indication{000747}
Il faut utiliser les identités trigonométriques classiques.
\finindication
\noindication
\indication{000749}
On compose les équations par la bonne fonction (sur le bon domaine de définition), 
par exemple cosinus pour la premi\`ere. Pour la dernière, commencer par 
étudier la fonction pour montrer qu'il existe une unique solution.
\finindication
\noindication
\noindication
\indication{000752}
Faire une étude de fonction.
La fonction $\text{sgn}(x)$ est la \emph{fonction signe} : elle vaut $+1$ si $x> 0$,
$-1$ si $x < 0$ (et $0$ si $x=0$).
\finindication
\noindication
\noindication
\noindication
\noindication
\noindication
\noindication
\noindication
\noindication
\noindication
\noindication
\noindication
\noindication
\noindication
\noindication
\noindication
\noindication
\noindication
\noindication
\noindication
\noindication
\noindication
\noindication
\noindication
\noindication
\noindication
\noindication
\noindication
\noindication
\noindication
\noindication
\noindication
\noindication
\noindication
\noindication
\indication{006973}
Dériver la différence des deux expressions.
\finindication
\noindication
\noindication
\indication{000758}
\begin{enumerate}
  \item Regarder ce qui se passe en deux valeurs oppos\'ees $x$ et $-x$.
  \item Poser $X = e^x$.
 \end{enumerate}
\finindication
\indication{000759}
R\'eponses :
\begin{enumerate}
    \item $\frac 34$ ;
    \item $\ln 2$.
\end{enumerate}
\finindication
\noindication
\noindication
\noindication
\indication{000764}
Pour la première question calculer $\frac{1}{\cos^2t}$.
Pour la seconde question, vérifier que 
$y=\ln \left(\tan\left(\frac{t}{2}+\frac{\pi}{4}\right)\right)$ 
est bien défini et calculer $\sh y$.
\finindication
\noindication
\noindication
\noindication
\noindication
\noindication
\noindication
\noindication
\noindication
\noindication
\noindication
\noindication
\indication{000776}
Montrer que l'\'equation $x^y=y^x$ est \'equivalente \`a $\frac{\ln x}{x}= \frac{\ln y}{y}$,
puis \'etudier la fonction $x \mapsto \frac{\ln x}{x}$.
\finindication
\noindication
\noindication
\noindication
\noindication
\noindication
\indication{006975}
On trouve $-\frac{1+e^{-2x}}{\ln(1+e^{-2x})}$.
\finindication
\indication{006976}
Commencer par calculer $C_n+S_n$ et $C_n-S_n$ à l'aide des fonctions $\ch$ et $\sh$.
\finindication
\indication{006977}
Poser $X=e^x$ et $Y=e^y$ et se ramener à un système d'équations du type somme-produit.
\finindication
\noindication
\indication{006979}
On trouve $f(x)=|\ln x|$ pour tout $x>0$.
\finindication
\indication{006980}
Faire le tableau de variations de $f:x\mapsto\Argsh x+\Argch x$.
\finindication
\noindication
\noindication
\noindication
\noindication
\noindication
\noindication
\noindication
\noindication
\noindication
\noindication
\noindication
\noindication
\noindication
\noindication
\noindication
\noindication
\noindication
\noindication
\noindication
\noindication
\noindication
\noindication
\noindication
\noindication
\noindication
\noindication
\noindication
\indication{002082}
Les fonctions continues ne seraient-elles pas intégrables ?

\medskip

Il faut se souvenir de ce que vaut la somme des $n$ premiers entiers, la somme des carrés des $n$ premiers entiers 
et la somme d'une suite géométrique.
La formule générale pour les sommes de Riemann est que $\int_a^bf(x)d x$
est la limite (quand $n \to +\infty$) de 
$$S_n = \frac {b-a}n \sum_{k=0}^{n-1} f\left(a+k\frac {b-a}n\right).$$
\finindication
\indication{002083}
  \begin{enumerate}
  \item On pourra penser que le cosinus et le sinus sont les parties r\'eelles et imaginaires de la fonction $t \mapsto e^{it}$. On chercha donc d'abord \`a calculer $\int_0^{\frac \pi 2}  e^{it}\, dt$.
  \item On choisira $q$ tel que $q^n = \frac ba$.
  \end{enumerate}
\finindication
\noindication
\indication{002085}
  \begin{enumerate}
  \item Revenir à la définition de la continuité en $x_0$ en prenant $\epsilon = \frac {f(x_0)}{2}$ par exemple.
  \item Soit $f$ est tout le temps de même signe (et alors utiliser la première question), soit ce n'est pas le cas (et alors utiliser un théorème classique...).
  \item On remarquera que $\int_0^1 f(x) \, dx - \frac 12 = \int_0^1 (f(x) - x) dx$.
  \end{enumerate}
\finindication
\indication{002086}
Essayez d'encadrer $\int_a^b \frac{f(t)^n}{m^n}dt.$
\finindication
\indication{002087}
Il s'agit de montrer que la limite vaut $0$. Pour un $\alpha >0$ fix\'e on s\'eparera l'int\'egrale en deux partie selon que $f$ est plus petit ou plus grand que $1-\alpha$. 
\finindication
\noindication
\indication{002092}
Se ramener \`a une composition de fonctions ou revenir \`a la d\'efinition de la d\'eriv\'ee avec le taux d'accroissement.
\finindication
\indication{002093}
  \begin{enumerate}
  \item Soit faire comme l'exercice \ref{ex:compint}, soit s\'eparer l'int\'egrale en deux, et pour l'une faire un changement de variable $u=x^2$.
  \item $H(x)$ se calcule explicitement et montrer qu'en fait $H$ est une fonction constante, ensuite il faut comparer $H(x)$ et $F(x)$.
  \end{enumerate}
\finindication
\noindication
\noindication
\noindication
\noindication
\noindication
\noindication
\noindication
\noindication
\noindication
\noindication
\noindication
\noindication
\noindication
\noindication
\noindication
\noindication
\noindication
\noindication
\noindication
\noindication
\noindication
\noindication
\noindication
\noindication
\noindication
\noindication
\noindication
\noindication
\noindication
\noindication
\noindication
\noindication
\noindication
\indication{002100}
On pourra essayer de reconnaître des sommes de Riemann, puis calculer des intégrales. 
Pour le produit composer par la fonction $\ln$, afin de transformer le produit en une somme.
\finindication
\noindication
\noindication
\noindication
\noindication
\noindication
\noindication
\noindication
\noindication
\noindication
\noindication
\noindication
\noindication
\noindication
\noindication
\indication{002099}
Un dessin ne fait pas de mal !
Il faut ensuite résoudre l'équation $\frac{x^2}2=\frac 1{x^2+1}$
puis calculer deux intégrales.
\finindication
\noindication
\noindication
\noindication
\noindication
\indication{006863}
Il faut se ramener au calcul de $\displaystyle \int_0^a b\sqrt{1-\frac{x^2}{a^2}} dx$.
\finindication
\noindication
\noindication
\noindication
\noindication
\noindication
\noindication
\noindication
\noindication
\noindication
\noindication
\noindication
\noindication
\indication{006865}
\begin{enumerate}
  \item $\int \cos^{1234} x   \sin x \, d x = -\frac 1{1235}\cos^{1235}x+c$ (changement de variable $u = \cos x$)

  \item $\int \frac 1{x\ln x} \, dx=\ln \left| \ln x\right| +c$  (changement de variable $u=\ln x$)

  \item $\int \frac 1{3+\exp \left( -x\right) }dx=\frac 13\ln \left( 3\exp
x+1\right) +c$ (changement de variable  $u=\exp x$)

  \item $\int \frac{1}{\sqrt{4x-x^2}}dx=\arcsin \left( \frac 12x-1\right) +c$  (changement de variable $u=\frac 12x-1$)
\end{enumerate}
\finindication
\noindication
\noindication
\noindication
\noindication
\noindication
\noindication
\noindication
\noindication
\indication{006864}
\begin{enumerate}
  \item Pour $\int x^2 \ln x \, dx$ poser $v'=x^2$, $u=\ln x$.

  \item Pour $\int x \arctan x \, dx$ poser $v'=x$ et $u= \arctan x$.

  \item Pour les deux il faut faire une intégration par parties avec $v'=1$.

  \item Pour $\int \cos x\exp x \,dx$ il faut faire deux intégrations par parties.
\end{enumerate}
\finindication
\noindication
\noindication
\noindication
\indication{002096}
  \begin{enumerate}
  \item Faire une intégration par parties afin d'exprimer $I_{n+2}$ en fonction de $I_n$. 
Pour le calcul explicite on distinguera le cas des $n$ pairs et impairs.
  \item Rappel : $u_n\sim v_n$ est équivalent à $\frac{u_n}{v_n} \to 1$. 
Utiliser la décroissance de $I_n$ pour encadrer $\frac{I_{n+1}}{I_n}$.
  \end{enumerate}
\finindication
\noindication
\noindication
\noindication
\noindication
\noindication
\noindication
\indication{002097}
  \begin{enumerate}
  \item Majorer par $x^n$.
  \item
  \item On pourra calculer $(I_0+I_1)-(I_1+I_2)+(I_2+I_3)- \cdots$
  \end{enumerate}
\finindication
\noindication
\noindication
\noindication
\indication{002089}
Calculer la somme et la diff\'erence de ces deux int\'egrales.
\finindication
\indication{002095}
$\int_0^{\frac \pi 2}\frac 1{1+\sin x}dx=1$ (changement de variables $t=\tan \frac x2$).

$\int_0^{\frac \pi 2}\frac{\sin x}{1+\sin x}dx=\frac \pi 2-1$ (utiliser la précédente).
\finindication
\noindication
\noindication
\noindication
\noindication
\noindication
\noindication
\noindication
\noindication
\noindication
\noindication
\noindication
\noindication
\noindication
\noindication
\noindication
\noindication
\noindication
\noindication
\noindication
\noindication
\noindication
\noindication
\noindication
\noindication
\noindication
\noindication
\noindication
\noindication
\indication{006866}
\begin{enumerate}
  \item $\int \frac{x+2}{x^2-3x-4}\,dx=-\frac 15\ln \left| x+1\right| +\frac 65\ln
\left| x-4\right| +c$ (décomposition en éléments simples)

  \item $\int \frac{x-1}{x^2+x+1}\,dx= \frac 12 \ln|x^2+x+1| - 
\sqrt3 \arctan \left( \frac{2}{\sqrt3}\left(x+\frac 12\right) \right) + c$

  \item $\int \sin ^8x\cos ^3x \, dx=\frac 19\sin ^9x-\frac 1{11}\sin ^{11}x+c$

  \item $\int \frac 1{\sin x} \, dx=\frac 12\ln \left| \frac{1-\cos x}{1+\cos x}%
\right| +c=\ln \left| \tan \frac x2\right| +c$ (changement de variable $u=\cos x$ ou $u=\tan \frac
x2$)

  \item $\int \frac{3-\sin x}{2\cos x+3\tan x}\,dx=-\frac 15\ln |2-\sin x| + \frac 75 \ln |1+2\sin x| + c$  
(changement de variable $u=\sin x$)
\end{enumerate}
\finindication
\indication{006867}
\begin{enumerate}
  \item $\int_0^{\frac \pi 2}x\sin x \, dx=1$ (intégration par parties $v'=\sin x$, $u=x$)

  \item $\int_0^1 \frac{e^x}{\sqrt{e^x+1}} \,  dx=2\sqrt{e+1} -2\sqrt 2$ (à l'aide du changement de variable $u=e^x$)

  \item $\int_0^1\frac 1{\left( 1+x^2\right) ^2} \, dx=\frac \pi 8+\frac 14$
(changement de variable $x=\tan t$, $dx = (1+\tan^2 t) dt$ et $1+\tan^2 t = \frac{1}{\cos^2 t}$)

  \item $\int_0^1\frac{3x+1}{\left( x+1\right) ^2} \, dx=3\ln 2-1$ (décomposition en
éléments simples de la forme
$\frac{3x+1}{\left( x+1\right) ^2} = \frac{\alpha}{x+1}+\frac{\beta}{(x+1)^2}$)

  \item $\int_{\frac 12}^2\left( 1+\frac 1{x^2}\right) \arctan x \, dx=\frac{3\pi }4$
(changement de variables $u=\frac 1x$ et $\arctan x+\arctan \frac 1x= \pm \frac \pi 2$)

\end{enumerate}
\finindication
\noindication
\noindication
\noindication
\noindication
\noindication
\noindication
\noindication
\noindication
\noindication
\noindication
\noindication
\noindication
\noindication
\noindication
\noindication
\noindication
\noindication
\noindication
\noindication
\noindication
\noindication
\noindication
\noindication
\noindication
\noindication
\noindication
\noindication
\noindication
\noindication
\noindication
\noindication
\noindication
\noindication
\noindication
\noindication
\noindication
\noindication
\noindication
\noindication
\noindication
\noindication
\noindication
\noindication
\noindication
\noindication
\noindication
\noindication
\noindication
\noindication
\noindication
\noindication
\noindication
\noindication
\noindication
\noindication
\noindication
\noindication
\noindication
\noindication
\noindication
\noindication
\noindication
\noindication
\noindication
\noindication
\noindication
\noindication
\noindication
\noindication
\noindication
\noindication
\noindication
\noindication
\noindication
\noindication
\noindication
\noindication
\noindication
\noindication
\noindication
\noindication
\noindication
\noindication
\noindication
\noindication
\noindication
\noindication
\noindication
\noindication
\noindication
\noindication
\noindication
\noindication
\noindication
\noindication
\noindication
\noindication
\noindication
\noindication
\noindication
\noindication
\noindication
\noindication
\noindication
\noindication
\noindication
\noindication
\noindication
\noindication
\noindication
\noindication
\noindication
\noindication
\noindication
\noindication
\noindication
\noindication
\noindication
\noindication
\noindication
\noindication
\noindication
\noindication
\noindication
\noindication
\noindication
\noindication
\noindication
\noindication
\noindication
\noindication
\noindication
\noindication
\noindication
\noindication
\noindication
\noindication
\noindication
\noindication
\noindication
\noindication
\noindication
\noindication
\noindication
\noindication
\noindication
\noindication
\noindication
\noindication
\noindication
\noindication
\noindication
\noindication
\noindication
\noindication
\noindication
\noindication
\noindication
\noindication
\noindication
\noindication
\noindication
\noindication
\noindication
\noindication
\noindication
\noindication
\noindication
\noindication
\indication{001143}
Développer par rapport à la dernière colonne.
\finindication
\noindication
\indication{001145}
Développer par rapport à la première colonne
pour obtenir $\Delta_{n-1}$ et un autre déterminant facile à calculer
en développant par rapport à sa première ligne.
\finindication
\noindication
\noindication
\noindication
\noindication
\noindication
\noindication
\noindication
\indication{002453}
Faire les opérations suivantes sur les colonnes
$C_n \leftarrow C_n-t_n C_{n-1}$,
puis $C_{n-1} \leftarrow C_{n-1}-t_n C_{n-2}$,...,
$C_2 \leftarrow  C_2-t_nC_1$.
Développer par rapport a la bonne ligne et reconnaître
que l'on obtient le déterminant recherché mais au rang $n-1$.
\finindication
\noindication
\noindication
\noindication
\noindication
\noindication
\noindication
\noindication
\noindication
\noindication
\noindication
\noindication
\noindication
\noindication
\noindication
\noindication
 \noindication
 \noindication
\noindication
\noindication
\noindication
\noindication
\noindication
\noindication
\noindication
\noindication
\noindication
\noindication
\noindication
\noindication
\noindication
\noindication
\noindication
\noindication
\noindication
\noindication
\noindication
\noindication
\noindication
\noindication
\noindication
\noindication
\noindication
\noindication
\noindication
\noindication
\noindication
\noindication
\noindication
\noindication
\indication{006886}
\begin{enumerate}
  \item Règle de Sarrus.
  \item Développer par rapport à la deuxième ligne.
  \item Faire apparaître des $0$ sur la première colonne.
  \item Utiliser la linéarité par rapports à chaque ligne et chaque colonne
pour simplifier les coefficients.
  \item Faire apparaître des $0$...
  \item Faire apparaître des $0$...
  \item Permuter les lignes et les colonnes pour faire apparaître une matrice triangulaire
par blocs.
\end{enumerate}
\finindication
\noindication
\noindication
\noindication
\noindication
\noindication
\noindication
\noindication
\noindication
\noindication
\noindication
\noindication
\noindication
\noindication
\noindication
\noindication
\noindication
\noindication
\noindication
\noindication
\noindication
\noindication
\noindication
\noindication
\noindication
\noindication
\noindication
\noindication
\noindication
\noindication
\noindication
\noindication
\noindication
\noindication
\noindication
\noindication
\noindication
\noindication
\noindication
\noindication
\noindication
\noindication
\noindication
\noindication
\noindication
\noindication
\noindication
\noindication
\noindication
\noindication
\noindication
\noindication
\noindication
\noindication
\noindication
\noindication
\noindication
\noindication
\noindication
\noindication
\noindication
\noindication
\noindication
\noindication
\indication{003424}
\'Ecrire les polynômes sous la forme 
$P(x)=ax^3+bx^2+cx+d$.
Calculer $  \int_{2}^4 P(x)\; dx$ d'une part 
et $\alpha P(2) + \beta P(3) + \gamma P(4)$ d'autre part.
L'identification conduit à un système linéaire à quatre équations,
 d'inconnues $\alpha,\beta,\gamma$.
\finindication
\noindication
\noindication
\noindication
\noindication
\noindication
\noindication
\noindication
\noindication
\noindication
\noindication
\noindication
\noindication
\noindication
\noindication
\noindication
\noindication
\noindication
\noindication
\noindication
\noindication
\noindication
\noindication
\noindication
\noindication
\noindication
\noindication
\noindication
\noindication
\noindication
\noindication
\noindication
\noindication
\indication{007026}
Se ramener à des systèmes linéaires soit en éliminant les quotients, soit en effectuant un changement de variable.
\finindication
\noindication
\indication{007028}
Penser à $(a+b)^2$ et $(a-b)^2$.
\finindication
\noindication
\indication{007030}
Prendre le logarithme des équations. 
\finindication
\noindication
\noindication
\noindication
\noindication
\noindication
\noindication
\noindication
\noindication
\noindication
\noindication
\noindication
\noindication
\noindication
\noindication
\noindication
\noindication
\noindication
\noindication
\noindication
\noindication
\noindication
\noindication
\noindication
\noindication
\noindication
\noindication
\noindication
\noindication
\noindication
\noindication
\noindication
\noindication
\noindication
\noindication
\noindication
\noindication
\noindication
\noindication
\noindication
\noindication
\noindication
\noindication
\noindication
\noindication
\noindication
\noindication
\noindication
\noindication
\noindication
\noindication
\noindication
\noindication
\noindication
\noindication
\noindication
\noindication
\noindication
\noindication
\noindication
\noindication
\noindication
\noindication
\noindication
\noindication
\noindication
\noindication
\noindication
\noindication
\noindication
\noindication
\noindication
\noindication
\noindication
\noindication
\noindication
\noindication
\noindication
\noindication
\noindication
\noindication
\noindication
\noindication
\noindication
\noindication
\noindication
\noindication
\noindication
\noindication
\noindication
\noindication
\noindication
\noindication
\noindication
\noindication
\noindication
\noindication
\noindication
\noindication
\noindication
\noindication
\noindication
\noindication
\noindication
\noindication
\noindication
\noindication
\noindication
\noindication
\noindication
\noindication
\noindication
\noindication
\noindication
\noindication
\noindication
\noindication
\noindication
\noindication
\noindication
\noindication
\noindication
\noindication
\noindication
\noindication
\noindication
\noindication
\noindication
\noindication
\noindication
\noindication
\noindication
\noindication
\noindication
\noindication
\noindication
\noindication
\noindication
\noindication
\noindication
\noindication
\noindication
\noindication
\noindication
\noindication
\noindication
\noindication
\noindication
\noindication
\noindication
\noindication
\noindication
\noindication
\noindication
\noindication
\noindication
\noindication
\noindication
\noindication
\noindication
\noindication
\noindication
\noindication
\noindication
\noindication
\noindication
\noindication
\noindication
\noindication
\noindication
\noindication
\noindication
\noindication
\noindication
\noindication
\noindication
\noindication
\noindication
\noindication
\noindication
\noindication
\noindication
\noindication
\noindication
\noindication
\noindication
\noindication
\noindication
\noindication
\noindication
\noindication
\noindication
\noindication
\noindication
\noindication
\noindication
\noindication
\noindication
\noindication
\noindication
\noindication
 \noindication
 \noindication
 \noindication
\noindication
\noindication
\noindication
\noindication
\noindication
\noindication
\noindication
\noindication
\noindication
\noindication
\noindication
\noindication
\noindication
\noindication
\noindication
\noindication
\noindication
\noindication
\noindication
\noindication
\noindication
\noindication
\noindication
\noindication
\noindication
\noindication
\noindication
\noindication
\noindication
\noindication
\noindication
\noindication
\noindication
\noindication
\noindication
\noindication
\noindication
\noindication
\noindication
\noindication
\noindication
\noindication
\noindication
\noindication
\noindication
\noindication
\noindication
\noindication
\noindication
\noindication
\noindication
\noindication
\noindication
\noindication
\noindication
\noindication
\noindication
\noindication
\noindication
\noindication
\noindication
\noindication
\noindication
\noindication
\noindication
\noindication
\noindication
\noindication
\noindication
\noindication
\noindication
\noindication
\noindication
\noindication
\noindication
\noindication
\noindication
\noindication
\noindication
\noindication
\noindication
\noindication
\noindication
\noindication
\noindication
\noindication
\noindication
\noindication
\noindication
\noindication
\noindication
\noindication
\noindication
\noindication
\noindication
\noindication
\noindication
\noindication
\noindication
\noindication
\noindication
\noindication
\noindication
\noindication
\noindication
\noindication
\noindication
\noindication
\noindication
\noindication
\noindication
\noindication
\noindication
\noindication
\noindication
\noindication
\noindication
\noindication
\noindication
\noindication
\noindication
\noindication
\noindication
\noindication
\noindication
\noindication
\noindication
\noindication
\noindication
\noindication
\noindication
\noindication
\noindication
\noindication
\noindication
\noindication
\noindication
\noindication
\noindication
\noindication
\noindication
\noindication
\noindication
\noindication
\noindication
\noindication
\noindication
\noindication
\noindication
\noindication
\noindication
\noindication
\noindication
\noindication
\noindication
\noindication
\noindication
\noindication
\noindication
\noindication
\noindication
\noindication
\noindication
\noindication
\noindication
\noindication
\noindication
\noindication
\noindication
\noindication
\noindication
\noindication
\noindication
\noindication
\noindication
\noindication
\noindication
\noindication
\noindication
\noindication
\noindication
\noindication
\noindication
\noindication
\noindication
\noindication
\noindication
\noindication
\noindication
\noindication
\noindication
\noindication
\noindication
\noindication
\noindication
\noindication
\noindication
\noindication
\noindication
\noindication
\noindication
\noindication
\noindication
\noindication
\noindication
\noindication
\noindication
\noindication
\noindication
\noindication
\noindication
\noindication
\noindication
\noindication
\noindication
\noindication
\noindication
\noindication
\noindication
\noindication
\noindication
\noindication
\noindication
\noindication
\noindication
\noindication
\noindication
\noindication
\noindication
\noindication
\noindication
\noindication
\noindication
\noindication
\noindication
\noindication
\noindication
\noindication
\noindication
\noindication
\noindication
\noindication
\noindication
\noindication
\noindication
\noindication
\noindication
\noindication
\noindication
\noindication
\noindication
\noindication
\noindication
\noindication
\noindication
\noindication
\noindication
\noindication
\noindication
\noindication
\noindication
\noindication
\noindication
\noindication
\noindication
\noindication
\noindication
\noindication
\noindication
\noindication
\noindication
\noindication
\noindication
\noindication
\noindication
\noindication
\noindication
\noindication
\noindication
\noindication
\noindication
\noindication
\noindication
\noindication
\noindication
\noindication
\noindication
\noindication
\noindication
\noindication
\noindication
\noindication
\noindication
\noindication
\noindication
\noindication
\noindication
\noindication
\noindication
\noindication
\noindication
\noindication
\noindication
\noindication
\noindication
\noindication
\noindication
\noindication
\noindication
\noindication
\noindication
\noindication
\noindication
\noindication
\noindication
\noindication
\noindication
\noindication
\noindication
\noindication
\noindication
\noindication
\noindication
\noindication
\noindication
\noindication
\noindication
\noindication
\noindication
\noindication
\noindication
\noindication
\noindication
\noindication
\noindication
\noindication
\noindication
\noindication
\noindication
\noindication
\noindication
\noindication
\noindication
\noindication
\noindication
\noindication
\noindication
\noindication
\noindication
\noindication
\noindication
\noindication
\noindication
\noindication
\noindication
\noindication
\noindication
\noindication
\noindication
\noindication
\noindication
\noindication
\noindication
\noindication
\noindication
\noindication
\noindication
\noindication
\noindication
\noindication
\noindication
\noindication
\noindication
\noindication
\noindication
\noindication
\noindication
\noindication
\noindication
\noindication
\noindication
\noindication
\noindication
\noindication
\noindication
\noindication
\noindication
\noindication
\noindication
\noindication
\noindication
\noindication
\noindication
\noindication
\noindication
\noindication
\noindication
\noindication
\noindication
\noindication
\noindication
\noindication
\noindication
\noindication
\noindication
\noindication
\noindication
\noindication
\noindication
\noindication
\noindication
\noindication
\noindication
\noindication
\noindication
\noindication
\noindication
\noindication
\noindication
\noindication
\noindication
\noindication
\noindication
\indication{002985}
Utiliser la formule de B{\'e}zout.
\finindication
\noindication
\noindication
\noindication
\noindication
\noindication
\noindication
\noindication
\noindication
\noindication
\noindication
\noindication
\noindication
\noindication
\noindication
\noindication
\noindication
\noindication
\noindication
\noindication
\noindication
\noindication
\noindication
\noindication
\noindication
\noindication
\noindication
\noindication
\noindication
\noindication
\noindication
\noindication
\noindication
\noindication
\noindication
\noindication
\noindication
\noindication
\noindication
\noindication
\noindication
\noindication
\noindication
\noindication
\noindication
\noindication
\noindication
\noindication
\noindication
\noindication
\noindication
\noindication
\noindication
\noindication
\noindication
\noindication
\noindication
\noindication
\noindication
\noindication
\noindication
\noindication
\noindication
\noindication
\noindication
\noindication
\noindication
\noindication
\noindication
\noindication
\noindication
\noindication
\noindication
\noindication
\noindication
\noindication
\noindication
\noindication
\noindication
\noindication
\noindication
\noindication
\noindication
\noindication
\noindication
\noindication
\noindication
\noindication
\noindication
\noindication
\noindication
\noindication
\noindication
\noindication
\noindication
\noindication
\noindication
\noindication
\noindication
\noindication
\noindication
\noindication
\noindication
\noindication
\noindication
\noindication
\noindication
\noindication
\noindication
\noindication
\noindication
\noindication
\noindication
\noindication
\noindication
\noindication
\noindication
\noindication
\noindication
\noindication
\noindication
\noindication
\noindication
\noindication
\noindication
\noindication
\noindication
\noindication
\noindication
\noindication
\noindication
\noindication
\noindication
\noindication
\noindication
\noindication
\noindication
\noindication
\noindication
\noindication
\noindication
\noindication
\noindication
\noindication
\noindication
\noindication
\noindication
\noindication
\noindication
\noindication
\noindication
\noindication
\noindication
\noindication
\noindication
\noindication
\noindication
\noindication
\noindication
\noindication
\noindication
\noindication
\noindication
\noindication
\noindication
\noindication
\noindication
\noindication
\noindication
\noindication
\noindication
\noindication
\noindication
\noindication
\noindication
\noindication
\noindication
\indication{002988}
Montrer que l'ensemble des $x$ tq $x^2 \ne e$ est de cardinal pair.
\finindication
\noindication
\noindication
\noindication
\noindication
\noindication
\noindication
\noindication
\noindication
\noindication
\noindication
\noindication
\noindication
\noindication
\noindication
\noindication
\noindication
\noindication
\noindication
\noindication
\noindication
\noindication
\noindication
\noindication
\noindication
\noindication
\noindication
\noindication
\noindication
\noindication
\noindication
\noindication
\noindication
\noindication
\noindication
\noindication
\noindication
\noindication
\noindication
\noindication
\noindication
\noindication
\noindication
\noindication
\noindication
\noindication
\noindication
\noindication
\noindication
\noindication
\noindication
\noindication
\noindication
\noindication
\noindication
\noindication
\noindication
\noindication
\noindication
\noindication
\noindication
\noindication
\noindication
\noindication
\noindication
\noindication
\noindication
\noindication
\noindication
\noindication
\noindication
\noindication
\noindication
\noindication
\noindication
\noindication
\noindication
\noindication
\noindication
\noindication
\noindication
\noindication
\noindication
\noindication
\noindication
\noindication
\noindication
\noindication
\noindication
\noindication
\noindication
\noindication
\noindication
\noindication
\noindication
\noindication
\noindication
\noindication
\noindication
\noindication
\noindication
\noindication
\noindication
\noindication
\noindication
\noindication
\noindication
\noindication
\noindication
\noindication
\noindication
\noindication
\noindication
\noindication
\noindication
\noindication
\noindication
\noindication
\noindication
\noindication
\noindication
\noindication
\noindication
\noindication
\noindication
\noindication
\noindication
\noindication
\noindication
\noindication
\noindication
\noindication
\noindication
\noindication
\noindication
\noindication
\noindication
\noindication
\noindication
\noindication
\noindication
\noindication
\noindication
\noindication
\noindication
\noindication
\noindication
\noindication
\noindication
\noindication
\noindication
\noindication
\noindication
\noindication
\noindication
\noindication
\noindication
\noindication
\noindication
\noindication
\noindication
\noindication
\noindication
\noindication
\noindication
\noindication
\noindication
\noindication
\noindication
\noindication
\noindication
\noindication
\noindication
\noindication
\noindication
\noindication
\noindication
\noindication
\noindication
\noindication
\noindication
\noindication
\noindication
\noindication
\noindication
\noindication
\noindication
\noindication
\noindication
\noindication
\noindication
\noindication
\noindication
\noindication
\noindication
\noindication
\noindication
\noindication
\noindication
\noindication
\noindication
\noindication
\noindication
\noindication
\noindication
\noindication
\noindication
\noindication
\noindication
\noindication
\noindication
\noindication
\noindication
\noindication
\noindication
\noindication
\noindication
\noindication
\noindication
\noindication
\noindication
\noindication
\noindication
\noindication
\noindication
\noindication
\noindication
\noindication
\noindication
\noindication
\noindication
\noindication
\noindication
\noindication
\noindication
\noindication
\noindication
\noindication
\noindication
\noindication
\noindication
\noindication
\noindication
\noindication
\noindication
\noindication
\noindication
\noindication
\noindication
\noindication
\noindication
\noindication
\noindication
\noindication
\noindication
\noindication
\noindication
\noindication
\noindication
\noindication
\noindication
\noindication
\noindication
\noindication
\noindication
\noindication
\noindication
\noindication
\noindication
\noindication
\noindication
\noindication
\noindication
\noindication
\noindication
\noindication
\noindication
\noindication
\noindication
\noindication
\noindication
\noindication
\noindication
\noindication
\noindication
\noindication
\noindication
\noindication
\noindication
\noindication
\noindication
\noindication
\noindication
\noindication
\noindication
\noindication
\noindication
\noindication
\noindication
\noindication
\noindication
\noindication
\noindication
\noindication
\noindication
\noindication
\noindication
\noindication
\noindication
\noindication
\noindication
\noindication
\noindication
\noindication
\noindication
\noindication
\noindication
\noindication
\noindication
\noindication
\noindication
\noindication
\noindication
\noindication
\noindication
\noindication
\noindication
\noindication
\noindication
\noindication
\noindication
\noindication
\noindication
\noindication
\noindication
\noindication
\noindication
\noindication
\noindication
\noindication
\noindication
\noindication
\noindication
\noindication
\noindication
\noindication
\noindication
\noindication
\noindication
\noindication
\noindication
\noindication
\noindication
\noindication
\noindication
\noindication
\noindication
\noindication
\noindication
\noindication
\noindication
\noindication
\noindication
\noindication
\noindication
\noindication
\noindication
\noindication
\noindication
\noindication
\noindication
\noindication
\noindication
\noindication
\noindication
\noindication
\noindication
\noindication
\noindication
\noindication
\noindication
\noindication
\noindication
\noindication
\noindication
\noindication
\noindication
\noindication
\noindication
\noindication
\noindication
\noindication
\noindication
\noindication
\noindication
\noindication
\noindication
\noindication
\noindication
\noindication
\noindication
\noindication
\noindication
\noindication
\noindication
\noindication
\noindication
\noindication
\noindication
\noindication
\indication{003151}
\begin{enumerate}
\item
\item $\dot6^2 = -\dot1$.
\end{enumerate}
\finindication
\noindication
\noindication
\noindication
\noindication
\noindication
\noindication
\noindication
\noindication
\noindication
\noindication
\noindication
\noindication
\noindication
\noindication
\noindication
\noindication
\noindication
\noindication
\noindication
\noindication
\noindication
\noindication
\noindication
\noindication
\noindication
\noindication
\noindication
\noindication
\noindication
\noindication
\noindication
\noindication
\noindication
\noindication
\noindication
\noindication
\noindication
\noindication
\noindication
\noindication
\noindication
\noindication
\noindication
\noindication
\noindication
\noindication
\noindication
\noindication
\noindication
\noindication
\noindication
\noindication
\noindication
\noindication
\noindication
\noindication
\noindication
\noindication
\noindication
\noindication
\noindication
\noindication
\noindication
\noindication
\noindication
\noindication
\noindication
\noindication
\noindication
\noindication
\noindication
\noindication
\noindication
\noindication
\noindication
\noindication
\noindication
\noindication
\noindication
\noindication
\noindication
\noindication
\noindication
\noindication
\noindication
\noindication
\noindication
\noindication
\noindication
\noindication
\noindication
\noindication
\noindication
\noindication
\noindication
\noindication
\noindication
\noindication
\noindication
\noindication
\noindication
\noindication
\noindication
\noindication
\noindication
\noindication
\noindication
\noindication
\noindication
\noindication
\noindication
\noindication
\noindication
\noindication
\noindication
\noindication
\noindication
\noindication
\noindication
\noindication
\noindication
\noindication
\noindication
\noindication
\noindication
\noindication
\noindication
\noindication
\noindication
\noindication
\noindication
\noindication
\noindication
\noindication
\noindication
\noindication
\noindication
\noindication
\noindication
\noindication
\noindication
\noindication
\noindication
\noindication
\noindication
\noindication
\noindication
\noindication
\noindication
\noindication
\noindication
\noindication
\noindication
\noindication
\noindication
\noindication
\noindication
\noindication
\noindication
\noindication
\noindication
\noindication
\noindication
\noindication
\noindication
\noindication
\noindication
\noindication
\noindication
\noindication
\noindication
\noindication
\noindication
\noindication
\noindication
\noindication
\noindication
\noindication
\noindication
\noindication
\noindication
\noindication
\noindication
\noindication
\noindication
\noindication
\noindication
\noindication
\noindication
\noindication
\noindication
\noindication
\noindication
\noindication
\noindication
\noindication
\noindication
\noindication
\noindication
\noindication
\noindication
\noindication
\noindication
\noindication
\noindication
\noindication
\noindication
\noindication
\noindication
\noindication
\noindication
\noindication
\noindication
\noindication
\noindication
\noindication
\noindication
\noindication
\noindication
\noindication
\noindication
\noindication
\noindication
\noindication
\noindication
\noindication
\noindication
\noindication
\noindication
\noindication
\noindication
\noindication
\noindication
\noindication
\noindication
\noindication
\noindication
\noindication
\noindication
\noindication
\noindication
\noindication
\noindication
\noindication
\noindication
\noindication
\noindication
\noindication
\noindication
\indication{001784}

\begin{enumerate}
\item Raisonner \`a l'aide d'une fonction $f$ de la variable $x$
telle que $x+y=f(x)$ et $\lim_{x \to 0} f(x)=0$.

\item Trouver deux courbes  dans
\[
\R^3 \setminus \{(x,y,z);2x^3+yz^2 =0\}
\]
qui tendent vers l'origine
telle que les limites, calcul\'ees le long de ces courbes,
existent mais ont des valeurs distinctes.

\item Utiliser le fait que le num\'erateur et le d\'enominateur sont toujours positifs
et que l'ordre du d\'enominateur est strictement plus grand que celui du num\'erateur.

\item
Raisonner \`a l'aide d'une fonction $h$ de la variable $y$
telle que $x^2-y^2=h(y)$ et $\lim_{y \to 0} h(y)=0$.

\item Chercher deux courbes dans
le domaine de d\'efinition
qui tendent vers l'origine
telle que les limites, calcul\'ees le long de ces courbes,
existent mais ont des valeures distinctes.

\end{enumerate}
\finindication
\indication{001785} 
Diviser le num\'erateur et le d\'enominateur
par $x^2$ resp. $y^2$ pour d\'eterminer $\lim_{y\to 0}f(x,y)$
resp. $\lim_{x\to 0}f(x,y)$. Montrer que, 
calcul\'ee le long d'une autre courbe
convenable, $\lim_{(x,y)\to (0,0)} f(x,y)$ existe et ne vaut pas z\'ero.
\finindication
\noindication
\indication{001787} 
\begin{enumerate}
\item 
R\'efuter l'existence de la limite \`a l'aide
de l'\'etude des limites le long de deux courbes adapt\'ees.

\item Utiliser les coordonn\'ees polaires dans le plan.

\item Si $\lim_{(x,y) \to (x_0,y_0)}h(x,y)$ existe et est non nul
alors 
\[
\lim_{(x,y) \to (x_0,y_0)}\frac{f(x,y)}{h(x,y)}=
\frac{\lim_{(x,y) \to (x_0,y_0)}f(x,y)}{\lim_{(x,y) \to (x_0,y_0)}h(x,y)} .
\]
\item
Chercher deux courbes dans
le domaine de d\'efinition
qui tendent vers l'origine
telles que les limites, calcul\'ees le long de ces courbes,
existent mais ont des valeures distinctes.
\end{enumerate}
\finindication
\indication{001788}
\begin{enumerate}
\item Raisonner \`a l'aide d'une fonction $h$ des variables $x$ et $y$
telle que $x+y+z=h(x,y)$ et $\lim_{(x,y) \to (0,0)} h(x,y)=0$.

\item Montrer que, d\'ej\`a sous la contrainte suppl\'ementaire $z=0$,
la limite ne peut pas exister.
\end{enumerate}
\finindication
\noindication
\noindication
\noindication
\noindication
\indication{002621} 
Pour \'etablir ou r\'efuter
l'existences d'une limite  particuli\`ere dans le plan et pour ensuite
d\'eterminer une limite pourvu qu'elle existe, utiliser le fait que
pour que $\lim_{n \to \infty}(x_n,y_n)$ existe dans le plan $\R^2$ 
il faut et il suffit que
chacune des limites $\lim_{n \to \infty}x_n$ et $\lim_{n \to \infty}y_n$ existe
en tant que limite finie.
\finindication
\noindication
\noindication
\noindication
\noindication
\noindication
\noindication
\noindication
\noindication
\noindication
\noindication
\noindication
\noindication
\noindication
\noindication
\indication{001800}
Il est \'evident que, en tout point $(x,y)$ distinct de l'origine,
la fonction $f$ est continue et que les d\'eriv\'ees partielles y
existent et sont continues. Il suffit de montrer que
$f$ est continue en $(0,0)$ et que les d\'eriv\'ees partielles y
existent et y sont continues.
\finindication
\indication{001801} 
Pour calculer les  d\'eriv\'ees partielles par rapport 
\`a une variable, interpr\'eter les autres variables comme param\`etres
et utiliser les r\`egles de calcul de la d\'eriv\'ee ordinaires.
\finindication
\noindication
\indication{001803} 
Distinguer tout de suite la partie triviale et la partie non triviale
de l'exercice.
\finindication
\noindication
\noindication
\noindication
\noindication
\noindication
\noindication
\noindication
\noindication
\noindication
\noindication
\noindication
\noindication
\noindication
\noindication
\noindication
\noindication
\noindication
\noindication
\noindication
\indication{002628} 
Le plan tangent \`a la surface d'\'equation 
$f(x,y,z)=0$ au point $(x_0,y_0,z_0)$ 
est donn\'e par l'\'equation
\begin{equation}
\frac{\partial f}{\partial x}(x_0,y_0,z_0) (x-x_0) 
+\frac{\partial f}{\partial y}(x_0,y_0,z_0)
(y-y_0)+\frac{\partial f}{\partial z}(x_0,y_0,z_0)
(z-z_0) =0.
\label{tang1}
\end{equation}
Dans le cas (1.), les calculs deviennt plus simples avec
l'\'equation 
\[
z^2=19-x^2-y^2.
\]
\finindication
\indication{002629} 
Ne pas confondre les variables pour l'\'equation de la surface,
les variables pour l'\'equation de la tangente en un point,
et les coordonn\'ees du point de contact.
\finindication
\indication{002630}
Le plan tangent \`a la surface d'\'equation 
$z=f(x,y)$ au point $(x_0,y_0,z_0)$ 
est donn\'e par l'\'equation
\begin{equation}
z-z_0=\frac{\partial f}{\partial x}(x_0,y_0) (x-x_0) 
+\frac{\partial f}{\partial y}(x_0,y_0)
(y-y_0) .
\label{tang3}
\end{equation}

\finindication
\indication{002631}   Le vecteur normal de la surface d'\'equation 
$f(x,y,z)=0$ au point $(x_0,y_0,z_0)$ 
est le vecteur
\begin{equation}
\left (\frac{\partial f}{\partial x}(x_0,y_0,z_0), 
\frac{\partial f}{\partial y}(x_0,y_0,z_0),
\frac{\partial f}{\partial z}(x_0,y_0,z_0)
\right) .
\label{normal}
\end{equation}
%Utiliser plut\^ot l'\'equation $f(x,y,z)=x^2+y^2-z^2=0$ du c\^one $C$.
\finindication
\indication{002632} 
Utiliser  la version
\eqref{tang3} de l'\'equation d'un plan tangent \`a une surface en un point.
\finindication
\indication{002633}
Pour les majorations, utiliser les coordonn\'ees
polaires $(r,\varphi)$ dans le plan.
Distinguer tout de suite les parties triviales des parties non triviales
de l'exercice.
\finindication
\indication{002634} Prendre
\begin{align*}
f(x,y)&=\exp[\sin(\pi+x)\cos y]=\exp[-\sin x \cos y ],
\\
h(x,y)&=\arctan[\sqrt{4+x}-2\exp(y)] .
\end{align*}
\finindication
\noindication
\noindication
\noindication
\noindication
\noindication
\noindication
\noindication
\noindication
\indication{002622} 
Pour calculer les  d\'eriv\'ees partielles par rapport 
\`a une variable, interp\'eter les autres variables comme param\`etres
et utiliser les r\`egles de calcul de la d\'eriv\'ee ordinaires.
\finindication
\indication{002623} 
Interpr\'eter
la d\'eriv\'ee directionnelle \`a l'aide de l'intersection du graphe 
de la fonction avec un plan convenable.
\finindication
\indication{002624} 
\begin{enumerate} 
\item Utiliser les coordonn\'ees polaires $(r,\varphi)$ dans le plan
et le fait que $\lim_{\begin{smallmatrix} r \to 0\\ r >0
\end{smallmatrix}} r \log r =0$.
\end{enumerate}
\finindication
\indication{002625}
 \begin{enumerate}
\item Pour r\'efuter la diff\'erentiabilit\'e de $f$ en $(0,0)$, il suffit de 
trouver
une d\'eriv\'ee directionnelle qui n'est pas combinaison lin\'eaire
des d\'eriv\'ees partielles (par rapport aux deux variables).

\item Le plan tangent au point $(x_0,y_0,f(x_0,y_0))$
du graphe $z=f(x,y)$ de $F$ est donn\'ee par l'\'equation
\begin{equation}
z-f(x_0,y_0) =\frac{\partial f}{\partial x}(x_0,y_0)(x-x_0)+
\frac{\partial f}{\partial y}(x_0,y_0)(y-y_0) .
\label{eqt}
\end{equation}
\end{enumerate}
\finindication
\indication{002626}
Calculer m\`ene \`a la v\'erit\'e.
\finindication
\indication{002627} 
\'Ecrire
$f(x,y)=(\sin (xy), y\cos x, xy\sin(xy)\exp(y^2))=(u,v,w)$.
\finindication
\noindication
\noindication
\noindication
\noindication
\noindication
\noindication
\noindication
\noindication
\noindication
\noindication
\noindication
\noindication
\noindication
\noindication
\noindication
\noindication
\noindication
\noindication
\noindication
\noindication
\noindication
\noindication
\noindication
\noindication
\noindication
\noindication
\noindication
\noindication
\noindication
\noindication
\noindication
\noindication
\noindication
\noindication
\noindication
\noindication
\noindication
\noindication
\noindication
\noindication
\noindication
\noindication
\noindication
\noindication
\noindication
\noindication
\noindication
\noindication
\noindication
\noindication
\noindication
\noindication
\noindication
\noindication
\noindication
\noindication
\noindication
\noindication
\indication{002635} 
Utiliser les r\`egles
\begin{align*}
d(f+g)&= df+dg , \\ 
 d(fg)&=fdg+gdf ,\\
d(f\circ h)&= (f'\circ h) dh .
\end{align*}
\finindication
\indication{002636} 
Soient $h$, $u$, $v$ des fonctions des deux 
variables $x$ et $y$. Rappeler
que
\begin{align*}
\mathrm dh&=\frac{\partial h}{\partial x}\mathrm dx +
\frac{\partial h}{\partial y} \mathrm dy ,
\\
\mathrm d(u\mathrm dx + v \mathrm dy) 
&=
\left(\frac{\partial v}{\partial x} -\frac{\partial u}{\partial y}\right)
\mathrm dx \mathrm dy ,
\\
\mathrm dx \mathrm dy &=-\mathrm dy\mathrm dx .
\end{align*}
\finindication
\indication{002637} 
On va d\'eterminer une primitive d'une forme diff\'erentielle
de degr\'e 1 par un changement de variables tel que,
dans les nouvelles variables, la primitive soit presque \'evidente.
\finindication
\indication{002638} 
Rappeler que la matrice hessienne est la matrice
constitu\'ee des d\'eriv\'ees partielles secondes.
\finindication
\indication{002639}
\begin{enumerate}
 \item  Montrer que
$$\frac{\partial^2F}{\partial r^2}+\frac{1}{r}\frac{\partial F}{\partial r}=
\frac1r\frac{\partial}{\partial r} r\frac{\partial}{\partial r}F.$$
 \item Montrer que
$$r\frac{\partial F}{\partial r}=x\frac{\partial f}{\partial x}+y\frac{\partial f}{\partial y}.$$
 \item Montrer que
$$\frac{\partial F}{\partial \theta}=x\frac{\partial f}{\partial y}-y\frac{\partial f}{\partial x}$$
 \item Utiliser ces r\'esultats, puis calculer encore un peu pour obtenir le r\'esultat souhait\'e. 
\end{enumerate}
\finindication
\indication{002640}
\begin{enumerate}
\item Grace au changement de variables
\[
\R^2 \longrightarrow \R^2,
\ (u,v) \longmapsto (x,y)=\left(\frac{u-v}{2}, \frac{u+v}{2}\right),
\]
la fonction $f$ s'\'ecrit 
$F(u,v)=f(\frac{u-v}{2}, \frac{u+v}{2})$.
Montrer que pour que $f$ soit solution de (\ref{eq:ondes}) 
il faut et il suffit que
\begin{equation}\label{eq:ondesbis}
\frac{\partial^2F}{\partial u\partial v}=0 .
\end{equation}
 \item  Montrer que, si $F$ satisfait \`a (\ref{eq:ondesbis}), il existe deux fonctions $g_1,g_2\colon\R\to\R$ telles que
\[
F(u,v)=g_1(u)+g_2(v).
\]
 \item  \'Ecrire la solution g\'en\'erale de (\ref{eq:ondes}) et expliquer la phrase: ``En une dimension d'espace, toute solution de l'\'equation des ondes s'\'ecrit comme somme d'une onde qui se d\'eplace vers la droite et une qui se d\'eplace vers la gauche.''
\end{enumerate}
\finindication
\noindication
\noindication
\noindication
\indication{002641} 
Rappel: Pour qu'un point critique non d\'eg\'en\'er\'e
pr\'esente un maximum relatif (resp. minimum relatif) il faut et il suffit que
la forme hessienne en ce point soit n\'egative (resp. positive) ;
pour qu'un point critique non d\'eg\'en\'er\'e
pr\'esente un point selle il faut et il suffit que
la forme hessienne en ce point soit (non d\'eg\'en\'er\'ee et)
ind\'efinie.
\finindication
\indication{002642} 
Voir l'exercice pr\'ec\'edent.
\finindication
\indication{002643}  
Voir les exercices pr\'ec\'edents.
\finindication
\indication{002644}
Le plan tangent \`a la surface d'\'equation $f(x,y,z)=0$ au point $(x_0,y_0,z_0)$ 
est donn\'e par l'\'equation
\begin{equation}
\frac{\partial f}{\partial x}(x_0,y_0,z_0) (x-x_0) +\frac{\partial f}{\partial y}(x_0,y_0,z_0)
(y-y_0)+\frac{\partial f}{\partial z}(x_0,y_0,z_0)(z-z_0) =0.
\label{tang4}
\end{equation}
\finindication
\indication{002645} 
Rappel du th\'eor\`eme des fonctions implicites
pour une fonction $f$ de classe $C^1$
de deux variables d\'efinie dans un ouvert du plan: {\em Soit $(x_0,y_0)$ un point tel que
$\frac{\partial f}{\partial y}(x_0,y_0)~\ne~0$.
Au voisinage de $x_0$,
il existe une fonction $h$ de classe $C^1$ de la variable $x$
d\'efinie dans un intervalle ouvert 
appropri\'e telle que $h(x_0)=y_0$ et telle que,
pour qu'au voisinage de $(x_0,y_0)$
les coordonn\'ees $x$ et $y$ du point $(x,y)$
satisfassent \`a l'\'equation
$f(x,y)=0$
il faut et il suffit que $y=h(x)$ et, s'il en est ainsi,}
\[
h'(x_0)=-\frac{\frac{\partial f}{\partial x}(x_0,y_0)}
{\frac {\partial f}{\partial y}(x_0,y_0)}.
\]
{\em D\`es que l'intervalle de d\'efinition de la fonction $h$ est fix\'e la fonction
$h$ est unique.\/}


\finindication
\indication{002646} Voir l'exercice pr\'ec\'edent.
\finindication
\noindication
\noindication
\noindication
\noindication
\noindication
\noindication
\noindication
\noindication
\noindication
\noindication
\noindication
\noindication
\noindication
\noindication
\noindication
\noindication
\noindication
\noindication
\noindication
\noindication
\noindication
\noindication
\noindication
\noindication
\noindication
\noindication
\noindication
\noindication
\noindication
\noindication
\noindication
\noindication
\noindication
\noindication
\noindication
\noindication
\noindication
\noindication
\noindication
\noindication
\noindication
\noindication
\noindication
\noindication
\noindication
\noindication
\noindication
\noindication
\noindication
\noindication
\noindication
\noindication
\noindication
\noindication
\noindication
\noindication
\noindication
\noindication
\noindication
\noindication
\noindication
\noindication
\noindication
\noindication
\noindication
\noindication
\noindication
\noindication
\noindication
\noindication
\noindication
\noindication
\noindication
\noindication
\noindication
\noindication
\noindication
\noindication
\noindication
\noindication
\noindication
\noindication
\noindication
\noindication
\noindication
\noindication
\noindication
\noindication
\noindication
\noindication
\noindication
\noindication
\noindication
\noindication
\noindication
\noindication
\noindication
\noindication
\noindication
\noindication
\noindication
\noindication
\noindication
\noindication
\noindication
\noindication
\noindication
\noindication
\noindication
\noindication
\noindication
\noindication
\noindication
\noindication
\noindication
\noindication
\noindication
\noindication
\noindication
\noindication
\noindication
\noindication
\noindication
\noindication
\noindication
\noindication
\noindication
\noindication
\noindication
\noindication
\noindication
\noindication
\noindication
\noindication
\noindication
\noindication
\noindication
\noindication
\noindication
\noindication
\noindication
\noindication
\noindication
\noindication
\noindication
\noindication
\noindication
\noindication
\noindication
\noindication
\noindication
\noindication
\noindication
\noindication
\noindication
\noindication
\noindication
\noindication
\noindication
\noindication
\noindication
\noindication
\noindication
\noindication
\noindication
\noindication
\noindication
\indication{000853}
Chercher deux réels $a,b$ tels que, si on pose $u=x+ay$ et $v=x+by$, alors
$\dot u(t)=A u(t)$ et $\dot v(t)=B v(t)$, où $A,B$ sont des constantes.
\finindication
\noindication
\noindication
\noindication
\noindication
\noindication
\noindication
\noindication
\noindication
\noindication
\noindication
\indication{006992}
Une telle fonction $f$ est solution d'une équation différentielle $y'+y=c$.
\finindication
\indication{006993}\ 
\begin{enumerate}
  \item $x$ est solution particulière
  \item $\cos$ est solution particulière
\end{enumerate}
\finindication
\indication{006994}
Solution particulière :
\begin{enumerate}
  \item $-\frac1{2x}$
  
  \item $\frac{x^{k+1}}{k+1} \exp(x)$
  
  \item $\frac{\ln x}{1+\ln^2(x)}$
\end{enumerate}

\finindication
\indication{006995}
1. C'est une équation à variables séparées.
\finindication
\indication{006996}\ 
\begin{enumerate}
  \item une infinité de solutions
  \item une solution
\end{enumerate}
\finindication
\indication{007001}
\begin{enumerate}
  \item
  \begin{enumerate}
    \item Se ramener à $\frac{1}{1-n}z'+a(x)z+b(x)=0$.
    \item $y = \pm\frac{1}{\sqrt{\lambda x^2 + 2x}}$ ou $y=0$.
  \end{enumerate}
    \item 
  \begin{enumerate}
    \item Remplacer $y$ par $u+y_0$.
    \item $y=\frac1x+\frac1{x\ln|x|+\lambda x}$ ou $y = \frac1x$.
  \end{enumerate}
\end{enumerate}    
\finindication
\noindication
\noindication
\noindication
\noindication
\noindication
\noindication
\noindication
\noindication
\noindication
\noindication
\noindication
\noindication
\noindication
\noindication
\noindication
\noindication
\noindication
\noindication
\noindication
\noindication
\noindication
\noindication
\noindication
\noindication
\noindication
\noindication
\noindication
\noindication
\noindication
\noindication
\noindication
\indication{006998}
Pour la fin: principe de superposition.
\finindication
\indication{006999}
Utiliser la méthode de variation de la constante.
\finindication
\noindication
\noindication
\noindication
\noindication
\noindication
\noindication
\noindication
\noindication
\noindication
\noindication
\noindication
\noindication
\noindication
\noindication
\noindication
\noindication
\noindication
\noindication
\noindication
\noindication
\noindication
\noindication
\noindication
\noindication
\noindication
\noindication
\noindication
\noindication
\noindication
\noindication
\noindication
\noindication
\noindication
\noindication
\noindication
\noindication
\noindication
\noindication
\noindication
\noindication
\noindication
\noindication
\noindication
\noindication
\noindication
\noindication
\noindication
\noindication
\noindication
\noindication
\noindication
\noindication
\noindication
\noindication
\noindication
\noindication
\noindication
\noindication
\noindication
\noindication
\noindication
\noindication
\noindication
\noindication
\noindication
\noindication
\noindication
\noindication
\noindication
\noindication
\noindication
\noindication
\noindication
\noindication
\noindication
\noindication
\noindication
\noindication
\noindication
\noindication
\noindication
\noindication
\noindication
\noindication
\noindication
\noindication
\noindication
\noindication
\noindication
\noindication
\noindication
\noindication
\noindication
\noindication
\noindication
\noindication
\noindication
\noindication
\noindication
\noindication
\noindication
\noindication
\noindication
\noindication
\noindication
\noindication
\noindication
\noindication
\noindication
\noindication
\noindication
\noindication
\noindication
\noindication
\noindication
\noindication
\noindication
\noindication
\noindication
\noindication
\noindication
\noindication
\noindication
\noindication
\noindication
\noindication
\noindication
\noindication
\noindication
\noindication
\noindication
\noindication
\noindication
\noindication
\noindication
\noindication
\noindication
\noindication
\noindication
\noindication
\noindication
\noindication
\noindication
\noindication
\indication{002616}
Utiliser le langage de la g\'eom\'etrie \'el\'ementaire, y compris
les notions de surface de  r\'evolution, d'axe de
 r\'evolution, de sommet d'un parabolo\" \i de,
de sommet d'un c\^one, de concavit\'e  vers le haut ou vers le bas,
d'h\'elice, de spirale, etc.
\finindication
\indication{002617}
Exploiter les propri\'et\'es g\'eom\'etriques
des parties du plan qui d\'efinissent $A_1$ et $A_2$. Par exemple,
une courbe qui est d\'efinie comme \'etant l'image r\'eciproque d'un point
relativement \`a une fonction continue est une partie ferm\'ee du plan.
\finindication
\indication{002618} 
Raisonner \`a partir de la d\'efinition d'un ouvert dans le plan.
\finindication
\indication{002619} 
Exploiter le fait que le compl\'ementaire d'un ouvert est ferm\'e
et que le compl\'ementaire d'un ferm\'e est ouvert.
\finindication
\indication{002620}
Distinguer la partie triviale de l'exercice de la partie 
non triviale. 
Dans cet exercice,  le seul point d\'elicat est pour le param\`etre $t$ proche de $0$.
\finindication
\noindication
\noindication
\noindication
\noindication
\noindication
\noindication
\noindication
\noindication
\noindication
\noindication
\noindication
\noindication
\noindication
\noindication
\noindication
\noindication
\noindication
\noindication
\noindication
\noindication
\noindication
\noindication
\noindication
\noindication
\noindication
\noindication
\noindication
\noindication
\noindication
\noindication
\noindication
\noindication
\noindication
\noindication
\noindication
\noindication
\noindication
\noindication
\noindication
\noindication
\noindication
\noindication
\noindication
\noindication
\noindication
\noindication
\noindication
\noindication
\noindication
\noindication
\noindication
\noindication
\noindication
\noindication
\noindication
\noindication
\noindication
\noindication
\noindication
\noindication
\noindication
\noindication
\noindication
\noindication
\noindication
\noindication
\noindication
\noindication
\noindication
\noindication
\noindication
\noindication
\noindication
\noindication
\noindication
\noindication
\noindication
\noindication
\noindication
\noindication
\noindication
\noindication
\noindication
\noindication
\noindication
\noindication
\noindication
\noindication
\noindication
\noindication
\noindication
\noindication
\noindication
\noindication
\noindication
\noindication
\noindication
\noindication
\noindication
\noindication
\noindication
\noindication
\noindication
\noindication
\noindication
\noindication
\noindication
\noindication
\noindication
\noindication
\noindication
\noindication
\noindication
\noindication
\noindication
\noindication
\noindication
\noindication
\noindication
\noindication
\noindication
\noindication
\noindication
\noindication
\noindication
\noindication
\noindication
\indication{004811}
D{\'e}composer en {\'e}l{\'e}ments simples la fraction $\frac1{1-X^n}$.
\finindication
\noindication
\noindication
\noindication
\noindication
\noindication
\noindication
\noindication
\noindication
\noindication
\noindication
\noindication
\noindication
\noindication
\noindication
\noindication
\noindication
\noindication
\noindication
\noindication
\noindication
\noindication
\noindication
\noindication
\noindication
\noindication
\noindication
\noindication
\noindication
\noindication
\noindication
\noindication
\indication{004852} 
Consid{\'e}rer les disques ferm{\'e}s associ{\'e}e {\`a} un recouvrement
\og circulaire \fg\ du plan et mettre en {\'e}vidence une suite de
disques emboit{\'e}s dont les rayons tendent vers z{\'e}ro.
\finindication
\noindication
\noindication
\noindication
\noindication
\noindication
\noindication
\noindication
\noindication
\noindication
\noindication
\noindication
\noindication
\noindication
\noindication
\noindication
\noindication
\noindication
\noindication
\noindication
\noindication
\noindication
\noindication
\noindication
\noindication
\noindication
\noindication
\noindication
\noindication
\noindication
\noindication
\indication{001960}
Les médianes sont les droites $(AA')$, $(BB')$, $(CC')$.
\finindication
\noindication
\noindication
\noindication
\noindication
\noindication
\noindication
\noindication
\noindication
\noindication
\noindication
\noindication
\noindication
\noindication
\noindication
\noindication
\noindication
\noindication
\noindication
\noindication
\noindication
\noindication
\noindication
\noindication
\noindication
\noindication
\noindication
\noindication
\noindication
\noindication
\noindication
\noindication
\noindication
\noindication
\noindication
\noindication
\noindication
\noindication
\noindication
\noindication
\indication{002009}
Utiliser par exemple les nombres complexes.
\finindication
\noindication
\noindication
\noindication
\noindication
\noindication
\noindication
\noindication
\noindication
\noindication
\noindication
\noindication
\noindication
\noindication
\noindication
\noindication
\noindication
\noindication
\noindication
\noindication
\noindication
\noindication
\noindication
\noindication
\noindication
\noindication
\noindication
\noindication
\noindication
\noindication
\noindication
\noindication
\noindication
\noindication
\noindication
\noindication
\noindication
\noindication
\noindication
\noindication
\noindication
\noindication
\noindication
\noindication
\noindication
\noindication
\noindication
\noindication
\noindication
\noindication
\noindication
\noindication
\noindication
\noindication
\noindication
\noindication
\noindication
\noindication
\noindication
\noindication
\noindication
\noindication
\noindication
\noindication
\noindication
\noindication
\noindication
\noindication
\noindication
\noindication
\noindication
\noindication
\noindication
\noindication
\noindication
\noindication
\noindication
\noindication
\noindication
\noindication
\noindication
\noindication
\noindication
\noindication
\noindication
\noindication
\noindication
\noindication
\noindication
\noindication
\noindication
\noindication
\noindication
\noindication
\noindication
\noindication
\noindication
\noindication
\noindication
\noindication
\noindication
\noindication
\noindication
\noindication
\noindication
\noindication
\noindication
\noindication
\noindication
\noindication
\noindication
\noindication
\noindication
\noindication
\noindication
\noindication
\noindication
\noindication
\noindication
\noindication
\noindication
\noindication
\noindication
\noindication
\noindication
\noindication
\noindication
\noindication
\noindication
\noindication
\noindication
\noindication
\noindication
\noindication
\noindication
\noindication
\noindication
\noindication
\noindication
\noindication
\noindication
\noindication
\noindication
\noindication
\noindication
\noindication
\noindication
\noindication
\noindication
\noindication
\noindication
\noindication
\noindication
\noindication
\noindication
\noindication
\noindication
\noindication
\noindication
\noindication
\noindication
\noindication
\indication{007059}   
Le triangle rectangle $ADA’$ est inscrit dans un demi-cercle. 
% donc il y a un autre angle droit
\finindication      
\indication{007060}   
Déterminer les angles $\widehat{AMB}$ et $\widehat{ANB}$.
\finindication      
\indication{007061}   
Montrer que $(AB) \bot (RC)$.
\finindication      
\indication{007062}
\begin{enumerate}
\item Penser à un trapèze.
\item On peut obtenir une telle droite comme hauteur d'un triangle $ABC$ adéquat.
\end{enumerate}
\finindication      
\indication{007063}   
\begin{enumerate}
\item Penser au triangle de l'écolier.
\end{enumerate}
\finindication      
\indication{007064}    Pour l'aire, considérer l'aire du complémentaire de $IJKL$ par exemple, ou bien utiliser les diagonales de $ABCD$.
\finindication      
\indication{007065}   
Décomposer l'aire comme la somme des aires de deux triangles.
\finindication 
\indication{007066}   
\'Ecrire chacune des distances à l'aide d'aires de triangles.

\finindication      
\indication{007067}   
Utiliser des triangles isocèles.
\finindication      
\indication{007068}   
Où se trouve le point $M$ sur le segment $[AC]$ ?
\finindication      
\indication{007069}   
Le segment $[EC]$ est parallèle à $(AC)$, et sa longueur vaut la moitié du périmètre de $OAC$.
\finindication  
\indication{007070}
Pythagore.
\finindication  
\indication{007071}   
Utiliser la caractérisation des triangles isocèles à l'aide d'angles.
% Utiliser des angles opposés par le sommet. 
\finindication      
\indication{007072}   
\'Ecrire les distances aux sommets en fonction des angles du triangle.
\finindication      
\indication{007073}   
Décomposer les longueurs suivant les points de tangence du cercle inscrit.
\finindication
\indication{007074}
Triangles isocèles.
\finindication
\indication{007075}   
Partitionner le triangle en plusieurs triangles pour calculer l'aire.
\finindication   
\noindication
\noindication
\noindication
\noindication
\noindication
\noindication
\noindication
\noindication
\noindication
\noindication
\noindication
\noindication
\noindication
\noindication
\noindication
\noindication
\noindication
\noindication
\noindication
\noindication
\noindication
\noindication
\noindication
\noindication
\noindication
\noindication
\indication{004952}
La distance d'un point $M_0(x_0,y_0)$ à une droite $D$ d'équation $ax+by+c=0$ est donnée par la formule
$d(M_0,D) = \frac{|ax_0+by_0+c_0|}{\sqrt{a^2+b^2}}$.
\finindication
\noindication
\noindication
\noindication
\noindication
\noindication
\noindication
\noindication
\noindication
\indication{007047}
Pour la parallèle, considérer deux points $A$ et $B$ sur la droite et construire un parallélogramme $ABPQ$.

Pour la perpendiculaire, construire un cerf-volant $APBQ$.
\finindication
\indication{007048}
Utiliser une ou plusieurs cordes.
\finindication
\indication{007049}   Tracer le milieu $I$ de $AB$ puis procéder par analyse-synthèse.
\finindication      
\noindication
\indication{007051}   
2. Penser à l'intersection des médianes.
\finindication
\indication{007052}   
Le barycentre est associatif. Donc construire le milieu de deux côtés opposés, puis le milieu de ces milieux, puisque :
\[\frac{A+B+C+D}{4} = \frac{\frac{A+B}{2} + \frac{C+D}{2}}{2}.\]

Pour les autres questions, utiliser également l'associativité.
\finindication      
\indication{007053}   
\begin{enumerate}
\item Utiliser des triangles particuliers.
\item Utiliser le théorème de Thalès.
\end{enumerate}
\finindication
\indication{007054}   
\begin{enumerate}
\item À quelle distance de $O'$ se situe ce point d'intersection ?
\item Utiliser  le cercle de centre $O'$ et de rayon $r'-r$ pour une tangente \og extérieure\fg{} ou bien $r'+r$ pour une  tangente \og intérieure\fg, et un autre cercle.
\end{enumerate}
\finindication      
\indication{007055}   
1. Commencer par trouver le point de la droite qui va appartenir au cercle.\\

2. Idem.
\finindication  
\indication{007056}   
\begin{enumerate}
\item Test de méthodologie : quelles droites peut-on tracer à partir de ce qui est donné ?
\item Construire la droite équidistante (à distance $r$) des deux parallèles, puis les deux droites parallèles à la troisième et à distance $r$.
\end{enumerate}
\finindication      
\indication{007057} 
Soient $O_1$, $O_2$ et $O_3$ les centres des trois cercles. Considérer le centre le centre du cercle circonscrit à $O_1O_2O_3$. \\
\finindication
\noindication
\noindication
\noindication
\noindication
\indication{007080}   
Considérer les symétries centrales $\sigma_i$ en les points $P_i$, puis la composée $\sigma_1\circ \sigma_2\circ ... \circ \sigma_n$. Pour la deuxième partie, penser au théorème de Varignon.
%on doit avoir un parallélogramme, et alors tout point de départ est solution.
\finindication      
\indication{007081}
De la même façon qu'il est souvent utile de compléter un triangle rectangle en un rectangle, il est souvent utile de compléter un trapèze rectangle en un rectangle.

\finindication
\noindication
\indication{007083}   
Si $\phi$ est une homothétie envoyant $\mathcal C$ sur $\mathcal C'$, alors $O'=\phi(O)$. Il suffit d'avoir un deuxième couple $(M, \phi(M))$ pour pouvoir tracer le centre de l'homothétie $\phi$.
\finindication 
\indication{007084}   
Si les cercles ne sont pas confondus, il peut y avoir entre $0$ et $4$ tangentes. Faire des figures pour tous les cas possibles. Puis, déterminer les homothéties (ou les translations) envoyant un cercle sur l'autre et tracer leur centre.
\finindication
\indication{007085}   
En faisant  une figure avec le carré déjà construit, on voit alors deux segments parallèles, ce qui invite à utiliser  une homothétie.
\finindication      
\indication{007086}   
\begin{enumerate}
\item Considérer une homothétie de centre $G$.
\item Considérer une homothétie de centre $H$. On rappelle que $\overrightarrow{G\Omega}=-\frac12\overrightarrow{GH}$.
\end{enumerate}
\finindication      
\noindication
\indication{007088}   
Considérer le quadrilatère $ABB'A'$ et ses diagonales : elles se croisent sur la droite.
\finindication 
\noindication
\indication{007090}    Utiliser une homothétie et une symétrie centrale.
\finindication
\noindication
\indication{007092}   
Utiliser des homothéties.
\finindication
\indication{007093}   
Considérer la translation de vecteur $\overrightarrow{CB}$ et l'image de $I$ par cette translation.
\finindication
\indication{007094}    Des homothéties de même centre commutent, de même que des translations.
\finindication
\indication{007095}   
Homothéties ou translations.
\finindication
\indication{007096}   
Quelle est la partie linéaire de $f\circ g^{-1}$ ? 
\finindication      
\indication{007097}    Considérer la translation $\tau$ de distance $a$ suivant la direction de la droite.
\finindication      
\indication{007098}
Commencer par construire un cercle tangent aux deux droites.
\finindication
\indication{007099}   
Sans la condition sur $A$, l'exercice est facile. Tracer n'importe quel cercle tangent aux droites. Ensuite, appliquer la méthodologie classique. % translation dans un cas, homothétie dans l'autre.
\finindication
\indication{007100}   
Les images de trois points non alignés déterminent de façon unique une transformation affine. En déduire une condition sur le quatrième point.
\finindication 
\indication{007101}   
Considérer $I$ et $J$ les milieux de $MA$ et $AM'$ ainsi qu'une projection affine. Une telle projection réduit les distances.
\finindication      
\indication{007102}   
Faire tourner le carré circonscrit par rapport au carré inscrit.
\finindication      
\indication{007103}   
Procéder par analyse-synthèse et considérer des rotations.
\finindication
\indication{007104}   
Il suffit de construire une des droites d'appui du carré. Pour cela, il suffit de construire un deuxième point sur cette droite.
\finindication
\noindication
\indication{007106}   
L'aire de l'intersection vaut le quart de l'aire du carré $ABCD$.
\finindication      
\indication{007107}   
Considérer la rotation de centre $O$ et d'angle $\pi/2$.
\finindication
\indication{007108}   
Considérer la rotation de centre $O$ et d'angle $\pi/2$.
\finindication
\indication{007109}   
Considérer la rotation de centre $O$ et d'angle $\pi/3$.
\finindication
\indication{007110}   
Considérer  une rotation de centre $A$.
\finindication
\indication{007111}   
Si $ABC$ est un tel triangle, considérer les rotations d'angles $\pm \pi/3$ et centrées sur les sommets. Déterminer les images des différents points et droites par ces rotations.
\finindication
\indication{007112}   
Similitude.
\finindication
\indication{007113}   
Considérer les cercles de diamètre $[AB]$ et $[AC]$, le second point d'intersection, et une similitude de centre $A$.
\finindication
\indication{007114}   
Utiliser une similitude.
\finindication      
\indication{007115}   
Considérer des similitudes centrées sur $A$, $B$ et $C$ ainsi que leur compositions. (Ou alors, utiliser les nombres complexes.)
\finindication
\indication{007116}   
Considérer la similitude de centre $A$, d'angle $\pi/4$ et de rapport $\sqrt 2$, et celle de centre $C$, d'angle $-\pi/4$, et de rapport $1/\sqrt 2$. Ou alors, utiliser les nombres complexes.
\finindication      
\noindication
\indication{007118}
Triangles isocèles et rectangles
\finindication
\noindication
\indication{007120}   
Il y a deux tels octogones. En notant $O$ le centre d'un tel octogone, on doit avoir $\widehat{AOB}=\pm \pi/4$.
\finindication      
\noindication
\noindication
\noindication
\indication{007124}   
Rédiger avec des angles de droites et ne pas faire de distinction entre bissectrice extérieure et intérieure.
\finindication      
\indication{007125}
Introduire la tangente commune $\mathcal T$ aux deux cercles. 
%Utiliser le cas limite du théorème de l'angle au centre.
\finindication
\indication{007126}   
Soient $\mathcal C$ et $\mathcal C'$ les cercles circonscrits à $ARQ$ et $BPR$. 
Ils se coupent en $R$ et en un deuxième point $T$, ou alors ils sont tangents en $R$. Montrer que $T$ (ou $R$ dans le second cas) est sur le cercle circonscrit à $CQP$.
\finindication      
\indication{007127}
Les points $A$ et $E$ sont les extrémités d'un diamètre.  Ensuite, décomposer en triangles.
\finindication
\noindication
\noindication
\noindication
\indication{007131}   
Utiliser les angles droits pour montrer que des points sont cocycliques, puis utiliser le théorème de l'angle inscrit.
\finindication      
\indication{007132}   
Pour la conclusion, utiliser $AC = AK + KC$.
\finindication      
\noindication
\indication{007134}   
Décomposer $(\overrightarrow{MA},\overrightarrow{MC})$ en $(\overrightarrow{MA},\overrightarrow{AD}) + (\overrightarrow{AD},\overrightarrow{MC})$.% puis angles inscrits / angles au centre.
\finindication  
\noindication 
\indication{007136}
La somme des angles d'un quadrilatère convexe vaut $2\pi$.
\finindication
\noindication
\noindication
\indication{007139}
Utiliser les différentes caractérisations des triangles isocèles.
\finindication
\indication{007140}
Où se trouve le centre du triangle circonscrit d'un triangle rectangle ?
\finindication
\noindication
\indication{007142}
Considérer les puissances par rapport aux deux cercles.
\finindication
\indication{007143}
Utiliser des triangles semblables, ou bien utiliser la puissance d'un point par rapport à un cercle.
\finindication
\indication{007144}
Utiliser le second point d'intersection de la droite avec le cercle, s'il existe.
\finindication
\noindication
\indication{007153}
Considérer la similitude directe qui envoie le couple $(B,D)$ sur le couple $(C,E)$.
\finindication
\noindication
\indication{007155}
Si on n'utilise pas les nombres complexes, on pourra considérer $C'$ le milieu de $[AB]$, $s_1$ la similitude directe de centre $A$, de rapport $\sqrt 2$ et d'angle $\pi/4$ et $s_2$ la similitude de centre $B$, de rapport $1/\sqrt 2$ et d'angle $\pi/4$.

Dans le contexte de cet exercice, les similitudes directes seront utilisées comme dans l'exemple suivant : comme AQC est rectangle isocèle en $Q$, la similitude directe de centre $A$, d'angle $\pi/4$ et de rapport $\sqrt 2$ envoie $Q$ sur $C$.
\finindication
\indication{007156}
Reformuler une égalité de produits en une égalité de quotients.

Utiliser des similitudes de centre $M$.
\finindication
\indication{007157} 
Pour les translations ou symétries glissées, commencer par montrer que le vecteur doit être horizontal.
\finindication
\noindication
\noindication
\noindication 
\noindication 
\noindication
\indication{007165} 
Dans la dernière question, faire apparaître le point $P$ dans la condition d'alignement à l'aide de la relation de Chasles, par exemple.
\finindication
\noindication
\noindication
\noindication
\noindication
\noindication
\noindication
\noindication
\noindication
\noindication
\noindication
\noindication
\noindication
\noindication
\noindication
\noindication
\noindication
\noindication
\noindication
\noindication
\noindication
\noindication
\noindication
\noindication
\noindication
\noindication
\noindication
\noindication
\noindication
\noindication
\noindication
\noindication
\noindication
\noindication
\noindication
\noindication
\noindication
\noindication
\noindication
\noindication
\noindication
\noindication
\noindication
\noindication
\noindication
\noindication
\noindication
\noindication
\noindication
\noindication
\noindication
\noindication
\noindication
\noindication
\noindication
\noindication
\noindication
\noindication
\noindication
\noindication
\noindication
\noindication
\noindication
\noindication
\noindication
\noindication
\noindication
\noindication
\noindication
\noindication
\noindication
\noindication
\noindication
\noindication
\noindication
\noindication
\noindication
\noindication
\noindication
\noindication
\noindication
\noindication
\noindication
\noindication
\noindication
\noindication
\noindication
\noindication
\noindication
\noindication
\noindication
\noindication
\noindication
\noindication
\noindication
\noindication
\noindication
\noindication
\noindication
\noindication
\noindication
\noindication
\noindication
\noindication
\noindication
\noindication
\noindication
\noindication
\noindication
\noindication
\noindication
\noindication
\noindication
\noindication
\noindication
\noindication
\noindication
\noindication
\noindication
\noindication
\noindication
\noindication
\noindication
\noindication
\noindication
\noindication
\noindication
\noindication
\noindication
\noindication
\noindication
\noindication
\noindication
\noindication
\noindication
\noindication
\noindication
\noindication
\noindication
\noindication
\noindication
\noindication
\noindication
\noindication
\noindication
\noindication
\noindication
\noindication
\noindication
\noindication
\noindication
\noindication
\noindication
\noindication
\noindication
\noindication
\noindication
\noindication
\noindication
\noindication
\noindication
\noindication
\noindication
\noindication
\noindication
\noindication
\noindication
\noindication
\noindication
\noindication
\noindication
\noindication
\noindication
\noindication
\noindication
\noindication
\noindication
\noindication
\noindication
\noindication
\noindication
\noindication
\noindication
\noindication
\noindication
\noindication
\noindication
\noindication
\noindication
\noindication
\noindication
\noindication
\noindication
\noindication
\noindication
\noindication
\noindication
\noindication
\noindication
\noindication
\noindication
\noindication
\noindication
\noindication
\noindication
\noindication
\noindication
\noindication
\noindication
\noindication
\noindication
\noindication
\noindication
\noindication
\noindication
\noindication
\noindication
\noindication
\noindication
\noindication
\noindication
\noindication
\noindication
\noindication
\noindication
\noindication
\noindication
\noindication
\noindication
\noindication
\noindication
\noindication
\noindication
\noindication
\noindication
\noindication
\noindication
\noindication
\noindication
\noindication
\noindication
\noindication
\noindication
\noindication
\indication{006986}
Un point $M(t)$ est singulier si $x'(t)=0$ et $y'(t)=0$.
\finindication
\noindication
\noindication
\noindication
\noindication
\noindication
\noindication
\noindication
\noindication
\noindication
\noindication
\noindication
\noindication
\noindication
\noindication
\noindication
\noindication
\noindication
\noindication
\noindication
\noindication
\noindication
\indication{006990}
Utiliser le repère de Frenet $(\vec{u}_\theta,\vec{v}_\theta)$.
\finindication
\noindication
\noindication
\noindication
\noindication
\noindication
\noindication
\noindication
\noindication
\noindication
\noindication
\noindication
\noindication
\noindication
\noindication
\noindication
\noindication
\noindication
\noindication
\noindication
\noindication
\noindication
\noindication
\noindication
\noindication
\noindication
\noindication
\noindication
\noindication
\noindication
\noindication
\noindication
\noindication
\noindication
\noindication
\noindication
\noindication
\noindication
\noindication
\noindication
\noindication
\noindication
\noindication
\noindication
\noindication
\noindication
\noindication
\noindication
\noindication
\noindication
\noindication
\noindication
\noindication
\noindication
\noindication
\noindication
\noindication
\noindication
\noindication
\noindication
\noindication
\noindication
\noindication
\noindication
\noindication
\noindication
\noindication
\noindication
\noindication
\noindication
\noindication
\noindication
\noindication
\noindication
\noindication
\noindication
\noindication
\noindication
\noindication
\noindication
\noindication
\noindication
\noindication
\noindication
\noindication
\noindication
\noindication
\noindication
\noindication
\noindication
\noindication
\noindication
\noindication
\noindication
\noindication
\noindication
\noindication
\noindication
\noindication
\noindication
\noindication
\noindication
\noindication
\noindication
\noindication
\noindication
\noindication
\noindication
\noindication
\noindication
\noindication
\noindication
\indication{006881}
On rappelle la formule de Green-Riemann qui permet de faire le lien entre intégrale
double et intégrale curviligne :

\textbf{Théorème.} 
Soit $\mathcal{D}$ un domaine de $\Rr^2$ limité par une courbe
fermée $\mathcal{C}$ que l'on suppose coupée
par toute parallèle aux axes en deux points au plus. On considère une forme
différentielle
$\omega=Pdx+Qdy$  définie sur $\mathcal{D}$. Si
les fonctions $P$ et $Q$ sont de classe $C^{1}$, on a :
$$\int_{\mathcal{ C}^+} Pdx+Qdy=\iint_{\mathcal{D}} ( \frac{\partial Q}{\partial x}-\frac{\partial P}{\partial y} ) dx dy$$
o\`u l'on a noté $\mathcal{C}^{+}$ la courbe $\mathcal{C}$ que l'on a
orientée dans le sens direct.
\finindication
\noindication
\noindication
\noindication
\noindication
\noindication
\noindication
\noindication
\noindication
\noindication
\noindication
\noindication
\noindication
\noindication
\noindication
\noindication
\noindication
\noindication
\noindication
\noindication
\noindication
\noindication
\noindication
\noindication
\noindication 
\noindication
\noindication
\noindication
\noindication
\noindication
\noindication
\noindication
\noindication
\noindication
\noindication
\noindication
\noindication
\noindication
\noindication
\noindication
\noindication
\noindication
\noindication
\noindication
\noindication
\noindication
\noindication
\noindication
\noindication
\noindication
\noindication
\noindication
\noindication
\noindication
\noindication
\noindication
\noindication
\noindication
\noindication
\noindication
\noindication
\noindication
\noindication
\noindication
\noindication
\noindication
\noindication
\noindication
\noindication
\noindication
\noindication
\noindication
\noindication
\noindication
\noindication
\noindication
\noindication
\noindication
\noindication
\noindication
\noindication
\noindication
\noindication
\noindication
\noindication
\noindication
\noindication
\noindication
\noindication
\noindication
\noindication
\noindication
\noindication
\noindication
\noindication
\noindication
\noindication
\noindication
\noindication
\noindication
\noindication
\noindication
\noindication
\noindication
\noindication
\noindication
\noindication
\noindication
\noindication
\noindication
\noindication
 \noindication
\noindication
\noindication
\noindication
\noindication
\noindication
\noindication
\noindication
\noindication
\noindication
\noindication
\noindication
\noindication
\noindication
\noindication
\noindication
\noindication
\noindication
\noindication
\noindication
\noindication
\noindication
\noindication
\noindication
\noindication
\noindication
\noindication
\noindication
\noindication
\noindication
\noindication
\noindication
\indication{002101}
Consid\'erer la couleur des cases exclues.
\finindication
\indication{002102}
Pour la question (II) (b) on consid\`erera la partie $A_0$ minimale associ\'ee \`a $\varphi$
et l'on montrera que $A_0$ et $h(Y-g(A_0))$ forment une partition de $X$. La bijection sera
d\'efinie par $g$ sur $A_0$ et par $h^{-1} $ sur $h(Y-g(A_0))$.
\finindication
\noindication
\indication{002104}
Ne voir dans le mot ``rang\'ee'' qu'une condition d'alignement.
\finindication
\indication{002105}
Compter, dans un ensemble $E$ \`a $n$ \'el\'ements, le nombre de parties \`a $p$
\'el\'ements obtenues en r\'eunissant une partie $X$ \`a $k$ \'el\'ements \`a une 
partie \`a $p-k$ \'el\'ements du compl\'ementaire de $X$ dans $E$, $k$ d\'ecrivant
$\{0,\ldots, p\}$.
\finindication
\indication{002106}
$n^2-1= (n-1)(n+1)$ et $24= 2^3\cdot 3$.
\finindication
\indication{002107}
Les premi\`eres questions ne pr\'esentent aucune difficult\'e.

Pour la derni\`ere, le plus difficile (et le plus int\'eressant) est de deviner
la formule. Pour cela, calculer la puissance $n$-i\`eme pour $n=1,2,3,4,5\ldots$.
(La formule est donn\'ee dans la page ``Corrections''). 
\finindication
\noindication
\indication{002109}
On pourra montrer les points suivants:

(a) $x\star y=e \Rightarrow y\star x =e$

(b) L'\'el\'ement neutre \`a gauche est unique.

(c) L'\'el\'ement neutre \`a gauche est un \'el\'ement neutre \`a droite aussi.

(d) Tout \'el\'ement est inversible.
\finindication
\indication{002110}
Oui.
\finindication
\indication{002111}
Aucune difficult\'e.
\finindication
\indication{002112}
Pour l'existence d'un inverse pour toute matrice $n\times n$ de d\'eterminant non nul,
noter que $\hbox{\rm det}(A)\not=0$ entra\^\i ne que la matrice $A$ est inversible (comme
matrice) et que la matrice $A^{-1}$, qui est de d\'eterminant $1/\hbox{\rm det}(A)\not=0$ est
alors l'inverse de $A$ pour le groupe en question.
\finindication
\indication{002113}
Aucune difficult\'e.
\finindication
\noindication
\indication{002115}
Consid\'erer la partition de $G$ en sous-ensembles  du type $\{ x, x^{-1} \}$.
\finindication
\indication{002116}
On commence par montrer que $f$ est surjective, en notant que si $|G|=2m+1$, alors pour
tout $y\in G$ on a $y=(y^{m+1})^2$. 
\finindication
\indication{002117}
$x^m=a \Leftrightarrow x=a^u$ o\`u $um+v|G|=1$.
\finindication
\noindication
\indication{002119}
Standard.
\finindication
\noindication
\noindication
\indication{002122}
Pour le (c), introduire le morphisme $\Z \rightarrow < x >$ qui associe $nx$ \`a tout
entier $n\in \Z$. Ce morphisme est surjectif et de noyau $d\Z$ o\`u $d$ est l'ordre de $x$.
\finindication
\indication{002123}
Aucune difficult\'e.
\finindication
\noindication
\indication{002125}
Cons\'equence de l'exercice \ref{ex:deb24}.
\finindication
\indication{002126}
$\{1\}, \mu_2 \times \{1\}, \{1\}\times \mu_2, \{(1,1), (i,i)\}, \mu_2 \times \mu_2$. 
\finindication
\noindication
\indication{002128}
Standard.
\finindication
\indication{002129}
Pour la seconde question, noter que si $x$ est d'ordre $2$ dans $G$, alors $yxy^{-1}$ l'est aussi, pour tout $y\in G$.
\finindication
\noindication
\noindication
\indication{002132}
Commencer par analyser l'ordre possible des \'el\'ements de $G$.
\finindication
\noindication 
\indication{002134} 
Trouver l'ordre de $2$ dans $(\Zz/p\Zz)^\times$.
\finindication
\indication{002135} 
Trouver l'ordre de $2$ modulo $2^n-1$.
\finindication
\noindication
\noindication
\noindication
\noindication
\noindication
\noindication
\noindication
\noindication
\noindication
\noindication
\noindication
\noindication
\noindication
\noindication
\noindication
\noindication
\noindication
\noindication
\noindication
\noindication
\indication{002136}
$(xy)^{-1} = x^{-1} y^{-1} \Rightarrow xy=yx$.
\finindication
\noindication
\noindication
\noindication 
\noindication 
\noindication 
\noindication
\indication{002143}
(a) est standard. En utilisant (a), on obtient $(\Zz/15\Zz)^\times \simeq \Zz/2\Zz \times \Zz/4\Zz$, lequel n'est pas cyclique puisque tous les \'el\'ements sont d'ordre $1$, $2$ ou $4$. Le reste ne pose pas de grandes difficult\'es.
\finindication
\indication{002144}
(a) B\'ezout. (b) $\phi$ est injectif et ensembles de d\'epart et d'arriv\'ee ont m\^eme
cardinal.
\finindication
\noindication
\indication{002146}
$e^{2ik\pi/d} = \left(e^{2ik\pi/n}\right)^{n/d}\ (k\in \Z)$.
\finindication
\indication{002147}
$f(G^\prime)$ est un sous-groupe de $H$ isomorphe \`a $G^\prime /(\hbox{\rm
ker}(f) \cap G^\prime)$.
\finindication
\indication{002148}
R\'esulte de l'exercice \ref{ex:le12}.
\finindication
\noindication
\noindication
\indication{002151}
Les morphismes du groupe $(\Q,+)$ dans lui-m\^eme sont de la forme $x\rightarrow ax$
avec $a\in \Q$. Les morphismes du groupe $(\Q,+)$ dans $(\Z,+)$ sont, parmi les
pr\'ec\'edents, ceux dont l'image est dans $\Z$; il n'y a que le morphisme nul. Les 
morphismes du groupe $(\Z/m\Z,+)$ dans $(\Z,+)$ sont d\'etermin\'es par l'entier    
$f(1)$ qui doit v\'erifier $m f(1) = 0$; il n'y a que le morphisme nul, si $m\not=0$.
\finindication
\indication{002152}
L'ensemble $\textrm{Hom}(\Zz/m\Zz,\Zz/n\Zz)$ des morphismes de groupe de $\Zz/m\Zz$ dans $\Zz/n\Zz$ est un groupe ab\'elien pour l'addition naturelle des morphismes. On note $\delta$ le pgcd de $m$ et $n$ et $m^\prime$ et $n^\prime$ les entiers $m/\delta$ et $n/\delta$. Si $p:\Zz\rightarrow \Zz/m\Zz$ d\'esigne la surjection canonique, la correspondance associant \`a tout $f\in \textrm{Hom}(\Zz/m\Zz,\Zz/n\Zz)$ l'\'el\'ement $f\circ p(1)$ induit un isomorphisme de groupe entre $\textrm{Hom}(\Zz/m\Zz,\Zz/n\Zz)$ et le sous-groupe $n^\prime \Zz/n\Zz$ du groupe additif $ \Zz/n\Zz$, lequel est isomorphe \`a $\Zz/\delta \Zz$. 

L'ensemble $\textrm{Aut}(\Zz/n\Zz)$ des automorphismes de $\Zz/n\Zz$ est un groupe pour la composition. La correspondance pr\'ec\'edente induit un isomorphisme entre $\textrm{Aut}(\Zz/n\Zz)$ et le groupe $(\Zz/n\Zz)^\times$ des inversibles de $\Zz/n\Zz$.
\finindication
\noindication
\noindication 
\indication{002155} 
Le morphisme ``d\'eterminant'' de $\textrm{GL}_n(\Rr)$ dans $\Rr^\times$ est surjectif et de noyau $\textrm{SL}_n(\Rr)$.
\finindication
\noindication
\noindication
\indication{002158} 
Si $\zeta$ est un \'el\'ement de $G$ dont la classe modulo $H$ engendre
$G/H$, alors tout \'el\'ement de $G$ peut s'\'ecrire $h \zeta^m$ avec $h\in H$ et
$m\in \Z$.
\finindication
\indication{002159}
Appliquer l'exercice \ref{ex:le23} avec $H=Z(G)$.
\finindication
\noindication
\indication{002161} Exercice classique d'alg\`ebre lin\'eaire: $Z(\hbox{\rm GL}_n(\mathbb{F}_p)) =
\mathbb{F}_p^\times \cdot \hbox{\rm Id}_n$ (o\`u $\hbox{\rm Id}_n$ d\'esigne la matrice identit\'e   
d'ordre $n$).
\finindication
\noindication
\indication{002163}
Les questions (a) et (b) ne pr\'esentent aucune difficult\'e.
\smallskip

Pour la question (c), noter que, pour tout $x\in G$, on a $(\tau_x)^{|G|}=1$, et que la
restriction de $\tau_x$ \`a $H$ appartient \`a $\hbox{\rm Aut}(H)\simeq \hbox{\rm
Aut}(\Z/p\Z)$ (et utiliser l'exercice \ref{ex:le25}).
\finindication
\indication{002164}
Aucune difficult\'e. Observer que tout conjugu\'e d'un commutateur est un
commutateur et qu'un quotient $G/H$ est ab\'elien si et seulement si pour tous
$u,v\in G$, on a $uvu^{-1}v^{-1} \in H$.
\finindication
\noindication
\noindication
\noindication
\noindication
\noindication
\noindication
\noindication
\noindication
\noindication
\noindication
\noindication
\noindication
\noindication
\noindication
\noindication
\noindication
\noindication
\noindication
\noindication
\noindication
\noindication
\noindication
\noindication
\noindication
\noindication
\noindication
\noindication
\noindication
\noindication
\noindication
\indication{002166} 
Aucune difficult\'e.
\finindication
\noindication
\indication{002168}
(a) est une simple v\'erification.
\smallskip

(b) Les trois permutations s'\'ecrivent respectivement $(1\hskip 2pt 3\hskip 2pt 7\hskip 2pt
5)\hskip 2pt (2\hskip 2pt 6\hskip 2pt 4)$, $(1\hskip 2pt 7)\hskip 2pt (2\hskip 2pt 4\hskip
2pt 3)$ et $(2\hskip 2pt 3\hskip 2pt 7\hskip 2pt 6\hskip 2pt 5\hskip 2pt 4)$.
\smallskip

(c) est une simple v\'erification.
\smallskip

(d) {\bf Rappel:} De fa\c con g\'en\'erale, on dit qu'une permutation $\omega \in S_n$ est de
type $1^{r_1}$-$2^{r_2}$-$\cdots$-$d^{r_d}$ o\`u $d, r_1,\ldots,r_d$ sont des entiers $\geq 0$
tels que $r_1+\cdots+r_d=n$, si dans la d\'ecomposition de $\omega$ en cycles \`a support
disjoints, figurent $r_1$ $1$-cycles (ou points fixes), $r_2$ $2$-cycles, ... et $r_d$
$d$-cycles. En utilisant la question (c), il n'est pas difficile de montrer que deux
permutations sont conjugu\'ees dans $S_n$ si et seulement si elles sont de m\^eme type.    
Les classes de conjugaison de $S_n$ correspondent donc exactement \`a tous les types
possibles.

\smallskip
On obtient ainsi facilement les classes de conjugaison de $S_5$. Soit maintenant $H$ un
sous-groupe distingu\'e non trivial de $S_5$. D\`es que $H$ contient un \'el\'ement de $S_5$,
il contient sa classe de conjugaison; $H$ est donc une r\'eunion de classes de conjugaison.
En consid\'erant toutes les classes possibles que peut contenir $H$, on montre que $H=A_5$ ou
$H=S_5$. Par exemple, si $H$ contient la classe 1-2-2, alors $H$ contient $(1\hskip 2pt
2)\hskip 2pt (3\hskip 2pt 4) \times (1\hskip 2pt 3)\hskip 2pt (2\hskip 2pt 5) = (1\hskip 2pt
4\hskip 2pt 3\hskip 2pt 2\hskip 2pt 5)$ et donc la classe des $5$-cycles. D'apr\`es
l'exercice \ref{ex:deb67}, $H$ contient alors $A_5$. Le groupe $H$ est donc $A_5$ ou $S_5$. Les autres cas
sont similaires.
\finindication
\noindication
\noindication
\noindication
\noindication
\indication{002173}
Une puissance impaire d'une permutation impaire ne peut pas \^etre \'egale \`a $1$.
\finindication
\noindication
\noindication
\noindication
\indication{002177} 
(a) Aucune difficult\'e.

(b) Le nombre cherch\'e est l'orbite de $H$ sous l'action de $G$ par conjugaison sur ses sous-groupes et $\textrm{Nor}_G(H)$ est le fixateur de $H$ pour cette action.
\finindication
\indication{002178} 
Etudier l'action du groupe par translation sur l'ensemble quotient des classes modulo le sous-groupe.
\finindication
\indication{002179} 
Le seul point non imm\'ediat est que $H^\prime$ est d'indice fini dans $G$. Pour cela
consid\'erer le morphisme de $G$ \`a valeurs dans le groupe des permutations des
classes \`a gauche de $G$ modulo $H$, qui \`a $g\in G$ associe la permutation
$aH\rightarrow gaH$ et montrer que le noyau de ce morphisme est le groupe $H^\prime$.
\finindication
\noindication
\noindication
\noindication
\noindication
\indication{002184}
Question (d): Si $K$ le fixateur d'un \'el\'ement $x\in X$, alors $K$ est un sous-groupe
propre maximal de $G$ et $X$ est isomorphe \`a $G/\cdot K$ en tant que $G$-ensemble. D\'eduire
du fait que $H$ n'est pas contenu dans $K$ que $HK=G$ et que
$H/\cdot H\cap K\simeq G/\cdot K$.
\finindication
\indication{002185}
Soit $H$ un tel sous-groupe. On peut supposer sans perte de
g\'en\'eralit\'e que $H$ contient la transposition $(12)$. On 
pourra ensuite proc\'eder comme suit.

- montrer que $H$ est engendr\'e par le fixateur $H_1$ de $1$ et
par $(12)$.

- montrer que l'orbite de $2$ sous $H$ est l'union de l'orbite de
$2$ sous $H_1$ et de $1$.

- en d\'eduire que $H_1$ agit transitivement sur l'ensemble $\{ 2,
\dots , n \} $ et que $H$
agit $2$-transitivement sur $\{ 1, \dots , n \} $.

- d\'eduire du point pr\'ec\'edent que $H$ contient toutes
les transpositions.
\finindication
\indication{002186}
(a) est trivial.
\smallskip

(b): Noter d'abord que la condition sur le fixateur de $x$ est ind\'ependante de $x\in X$: en effet si $g$ est un \'el\'ement de $G$ envoyant $x$ sur un autre \'el\'ement $x^\prime\in X$ (qui existe par transitivit\'e de $G$), alors $G(x^\prime) = g G(x) g^{-1}$  et la correspondance $h\rightarrow g h g^{-1}$ permet d'identifier les actions de $G(x^\prime)$ sur $X\setminus \{x^\prime\}$ et celle de $G(x)$ sur $X\setminus \{x\}$. Supposons maintenant v\'erifi\'ee la condition sur le fixateur de $x$. Si $(x,y)$ et $(x^\prime,y^\prime)$ sont deux couples d'\'el\'ements distincts de $X$, il existe $\sigma\in G$ tel que $\sigma (x)=x^\prime$ (transitivit\'e de $G$) et il existe $\tau \in G$ tel que $\tau (x^\prime)=x^\prime$ et $\tau(\sigma(y)) = y^\prime$ (transitivit\'e de $G(x^\prime)$ sur $X\setminus \{x^\prime\}$ (noter que $\sigma
(y) \not=x^\prime$ car $\sigma (x)=x^\prime$)). La permutation $\tau \sigma$ v\'erifie
$\tau \sigma (x)=x^\prime$ et $\tau \sigma (y)=y^\prime$. Cela montre que $X$ est $2$-transitif. La r\'eciproque est triviale.

\smallskip

(c) Si l'action de $G$ sur $X$ est imprimitive et $X=\bigcup_{i=1}^r X_i$ est une partition de
$X$ comme dans la d\'efinition, alors il n'existe pas d'\'el\'ement $g\in G$ envoyant un
premier \'el\'ement $x_1\in X_1$ dans $X_1$ et un second \'el\'ement $x_1^\prime \in X_1$
dans $X_2$.
\smallskip

(d) L'action par translation d'un groupe cyclique $C$ sur lui-m\^eme est transitive, elle est
primitive si $|C|$ est premier (toute partition de $C$ en sous-ensembles de m\^eme cardinal
est forc\'ement triviale) mais elle n'est pas $2$-transitive (le fixateur de
tout \'el\'ement est trivial, ce qui contredit le (c) de l'exercice \ref{ex:deb84}).
\smallskip

(e) et (f) ne pr\'esentent aucune difficult\'e.
\finindication
\indication{002187}
On se ram\`ene \`a la situation o\`u le polygone est inscrit dans le plan
complexe et a pour sommets les racines de l'unit\'e $e^{2ik\pi/n}, k=0,1,\ldots,n-1$. Une
isom\'etrie laissant invariant le polygone fixe n\'ecessairement l'origine. Elle est donc de
la forme $z\rightarrow az$ ou $z\rightarrow a \overline z$ avec $|a|=1$. On voit ensuite que
$a$ est n\'ecessairement une racine $n$-i\`eme de $1$. Notons $\sigma$ l'isom\'etrie
$z\rightarrow e^{2i\pi/n} z$ et $\tau$ la conjugaison complexe. On a $D_n=\{\sigma^k
\tau^\varepsilon \hskip 2pt | \hskip 2pt k=0,\ldots,n-1, \varepsilon = \pm 1\}$. On v\'erifie
que $\sigma$ et $\tau$ engendrent le groupe $D_n$ et satisfont les relations $\sigma^n=1$,
$\tau^2=1$ et $\tau \sigma \tau^{-1}= \sigma^{-1}$. Autrement dit, $D_n$ est isomorphe au
groupe di\'edral d'ordre $2n$. Si $n$ est impair, son centre est trivial et si $n=2m$ est
pair, son centre est $\{1,\sigma^m\}$. Le groupe $D_n$ se plonge naturellement dans $S_n$;
comme $|D_3|=|S_3|=6$, ce plongement est un isomorphisme pour $n=3$. 
\finindication
\noindication
\noindication
\noindication
\noindication
\noindication
\noindication
\noindication
\noindication
\noindication
\noindication
\noindication
\noindication
\noindication
\noindication
\noindication
\noindication
\noindication
\noindication
\noindication
\noindication
\noindication
\noindication
\noindication
\noindication
\noindication
\noindication
\noindication
\noindication
\noindication
\indication{002190}
$|G| = |G/H| \hskip 2pt |H|$.
\finindication
\noindication
\indication{002192}
Pour les trois \'enonc\'es (a), (b) et (c), raisonner par r\'ecurrence sur $r$    
en utilisant le fait que le centre d'un $p$-groupe n'est pas trivial.
\finindication
\noindication
\noindication
\indication{002195}
On a 
$$D_6 = \Z/6\Z \times \hskip -6pt {\raise 1.4pt\hbox{${\scriptscriptstyle |}$}}
\Z/2\Z \simeq (\Z/2\Z \times \Z/3\Z) \times \hskip -6pt {\raise
1.4pt\hbox{${\scriptscriptstyle |}$}}
\Z/2\Z \simeq \Z/2\Z \times (\Z/3\Z \times \hskip -6pt {\raise
1.4pt\hbox{${\scriptscriptstyle |}$}} \Z/2\Z) \simeq \mu_2 \times S_3$$

Le premier isomorphisme est une application standard du lemme chinois. Pour le
deuxi\`eme, noter que le premier $\Z/2\Z$ est dans le centre du groupe et
donc que l'action sur lui par conjugaison du second $\Z/2\Z$ est triviale.
L'isomorphisme $\Z/3\Z \times \hskip -6pt {\raise
1.4pt\hbox{${\scriptscriptstyle |}$}} \Z/2\Z \simeq S_3$ est classique.
\finindication
\noindication
\indication{002197} 
Pour tout $g\in G$, $gSg^{-1}$ est un $p$-Sylow de $gHg^{-1}=H$
et donc $gSg^{-1}=S$.
\finindication
\noindication
\noindication
\noindication
\indication{002201}
Pour le (c), pour $H\not=\{1\}$ sous-groupe
distingu\'e de $A_5$, raisonner sur les \'el\'ements d'ordre $2$,   
$3$ et $5$ contenus dans $H$.
\finindication
\indication{002202}
L'identification de chacun des $p$-Sylow ne pose pas de difficult\'e. Observer ensuite que les sous-groupes de Sylow sont deux \`a deux d'intersection r\'eduite \`a $\{1\}$ et determiner leur nombre en comptant les \'el\'ements d'ordre $2$, $3$ et $5$.
\finindication
\noindication
\indication{002204} Les th\'eor\`emes de Sylow montrent qu'il n'y a qu'un seul
$q$-Sylow, n\'ecessairement distingu\'e. La suite est standard. Pour le dernier
point, utiliser que $\hbox{\rm Aut}(\Z/q\Z) \simeq (\Z/q\Z)^\times$ (exercice \ref{ex:le17}) et donc que $\Z/p\Z$ ne peut agir non trivialement sur
$\Z/q\Z$ que si
$p$ divise $q-1$.
\finindication
\indication{002205} D'apr\`es l'exercice \ref{ex:deb104}, un groupe d'ordre $35$ est isomorphe
\`a $\Z/5\Z \times \Z/7\Z$, lequel est isomorphe au groupe cyclique $\Z/35\Z$ par
le lemme chinois.
\finindication
\noindication
\indication{002207} 
Soit $G$ un groupe d'ordre $p^2q$ qu'on suppose simple. On distinguera deux cas:
$p>q$ et $p<q$. Dans le premier, montrer que $G$ admet $q$ $p$-Sylow d'ordre $p^2$
et que l'action par conjugaison de $G$ sur les $p$-Sylow d\'efinit un morphisme
injectif $G \hookrightarrow S_q$ et aboutir \`a une contradiction. Dans le second,
raisonner sur le nombre de $q$-Sylow pour aboutir \`a une contradiction (on sera
notamment amen\'e \`a \'eliminer le cas $p=2$ et $q=3$).
\finindication
\noindication
\indication{002209}
(a) Si $K$ est un sous-groupe d'ordre $20$, $K$ a un seul $5$-Sylow $L$ et donc
$K\subset \hbox{\rm Nor}_G(L)$ ce qui entraine que l'ordre de $\hbox{\rm Nor}_G
(L)$ est $20$ ou $60$. Mais alors il y aurait $1$ ou $3$ $5$-Sylow dans $G$. Or $1$
est impossible car $G$ est simple et $3$ contredit les pr\'edictions du
th\'eor\`eme de Sylow.
\smallskip

(b) Si $K$ a un unique $3$-Sylow $L$, $K\subset \hbox{\rm Nor}_G(L)$, et donc
l'ordre de $\hbox{\rm Nor}_G (L)$ serait $12$ ou $60$. Il y aurait alors $5$ ou
$1$ $3$-Sylow dans $G$. Comme ci-dessus, c'est impossible.
\smallskip

(c) Supposons que $H\cap K =<a>$ soit d'ordre $2$. Le centralisateur $\hbox{\rm
Cen}_G(a)$ de $a$ contient $H$ et $K$, donc $H\cup K$. Son ordre est au moins  
$6$ et est divisible par $4$. Les seules possibilit\'es sont $12$, $20$, $60$:

- $60$ est impossible, car $<a>$ serait distingu\'e dans $G$

- $20$ est impossible, d'apr\`es la question (a)

- $12$ est impossible, car $\hbox{\rm Cen}_G(a)$ aurait $4$ $3$-Sylow d'apr\`es la
question (b). Il ne resterait de la place que pour un seul $2$-Sylow ce qui
contredit $H\cup K \subset \hbox{\rm Cen}_G(a)$.
\smallskip

(d) Si $H=\hbox{\rm Nor}_G(H)$, il y a $15$ $2$-Sylow, et donc $46$
\'el\'ements d'ordre une puissance de $2$. Or il y a $6$ $5$-Sylow d'intersections
deux \`a deux triviales, et donc $24$ \'el\'ements d'ordre $5$. L'in\'egalit\'e
$46+24 > 60$ fournit une contradiction.
\smallskip

(e) Si $H$ est un $2$-Sylow, l'ordre de $\hbox{\rm Nor}_G(H)$ est $12$, $20$ ou
$60$. Mais $20$ est exclu (question (a)) de m\^eme que $60$ ($G$ est simple). La
seule possibilit\'e est $12$; il y a donc $5$ $2$-Sylow.
\smallskip

(f) L'action de $G$ par conjugaison sur les $5$-Sylow fournit un morphisme
$c:G\rightarrow S_5$ qui est injectif (car $G$ est simple). Le groupe $G$ est donc
isomorphe \`a son image $c(G)$ qui est un sous-groupe d'ordre $60$, donc d'indice
$2$ dans $S_5$. C'est donc $A_5$.
\finindication
\noindication
\noindication
\noindication
\noindication
\noindication
\noindication
\noindication
\noindication
\noindication
\noindication
\noindication
\noindication
\noindication
\noindication
\noindication
\noindication
\noindication
\noindication
\noindication
\noindication
\noindication
\noindication
\noindication
\noindication
\noindication
\noindication
\noindication
\noindication
\noindication
\noindication
\noindication
\noindication
\noindication
\noindication
\noindication
\noindication
\noindication
\noindication
\noindication
\noindication
\noindication
\noindication
\noindication
\noindication
\noindication
\noindication
\noindication
\noindication
\noindication
\noindication
\noindication
\noindication
\noindication
\noindication
\noindication
\noindication
\noindication
\noindication
\noindication
\noindication
\noindication
\noindication
\noindication
\noindication
\noindication
\noindication
\noindication
\noindication
\noindication
\noindication
\noindication
\noindication
\noindication
\noindication
\noindication
\noindication
\noindication
\noindication
\noindication
\noindication
\noindication
\noindication
\noindication
\noindication
\noindication
\noindication
\noindication
\noindication
\noindication
\noindication
\noindication
\noindication
\noindication
\noindication
\noindication
\noindication
\noindication
\noindication
\noindication
\noindication
\noindication
\noindication
\noindication
\noindication
\noindication
\noindication
\noindication
\noindication
\noindication
\noindication
\noindication
\noindication
\noindication
\noindication
\noindication
\noindication
\noindication
\noindication
\noindication
\noindication
\noindication
\noindication
\noindication
\noindication
\noindication
\noindication
\noindication
\noindication
\noindication
\noindication
\noindication
\noindication
\noindication
\noindication
\noindication
\noindication
\noindication
\noindication
\noindication
\noindication
\noindication
\noindication
\noindication
\noindication
\noindication
\noindication
\noindication
\noindication
\noindication
\noindication
\noindication
\noindication
\noindication
\noindication
\noindication
\noindication
\noindication
\noindication
\noindication
\noindication
\noindication
\noindication
\noindication
\noindication
\noindication
\noindication
\noindication
\noindication
\noindication
\noindication
\noindication
\noindication
\noindication
\noindication
\noindication
\noindication
\noindication
\noindication
\noindication
\noindication
\noindication
\noindication
\noindication
\noindication
\noindication
\noindication
\noindication
\noindication
\noindication
\noindication
\noindication
\noindication
\noindication
\noindication
\noindication
\noindication
\noindication
\noindication
\noindication
\noindication
\noindication
\noindication
\noindication
\noindication
\noindication
\noindication
\noindication
\noindication
\noindication
\indication{002253}
Voir la solution de l'exercice \ref{ptfermat}, deuxième question.
\finindication
\noindication
\noindication
\noindication
\noindication
\noindication
\noindication
\noindication
\noindication
\noindication
\noindication
\noindication
\noindication
\noindication
\noindication
\noindication
\noindication
\noindication
\noindication
\noindication
\noindication
\noindication
\noindication
\noindication
\noindication
\noindication
\noindication
\noindication
\noindication
\noindication
\noindication
\noindication
\noindication
\noindication
\noindication
\noindication
\noindication
\noindication
\noindication
\noindication
\noindication
\noindication
\noindication
\noindication
\noindication
\noindication
\noindication
\noindication
\noindication
\noindication
\noindication
\noindication
\noindication
\noindication
\noindication
\noindication
\noindication
\noindication
\noindication
\noindication
\noindication
\noindication
\noindication
\noindication
\noindication
\noindication
\noindication
\noindication
\noindication
\noindication
\noindication
\noindication
\noindication
\noindication
\noindication
\noindication
\noindication
\noindication
\noindication
\noindication
\noindication
\noindication
\noindication
\noindication
\noindication
\noindication
\noindication
\noindication
\noindication
\noindication
\noindication
\noindication
\noindication
\noindication
\noindication
\noindication
\noindication
\noindication
\noindication
\noindication
\noindication
\noindication
\noindication
\noindication
\noindication
\noindication
\noindication
\noindication
\noindication
\noindication
\noindication
\noindication
\noindication
\noindication
\noindication
\noindication
\noindication
\noindication
\noindication
\noindication
\noindication
\noindication
\noindication
\noindication
\noindication
\noindication
\noindication
\noindication
\noindication
\noindication
\noindication
\noindication
\noindication
\noindication
\noindication
\noindication
\noindication
\noindication
\noindication
\noindication
\noindication
\noindication
\noindication
\noindication
\noindication
\noindication
\noindication
\noindication
\noindication
\noindication
\noindication
\noindication
\noindication
\noindication
\noindication
\noindication
\noindication
\noindication
\noindication
\noindication
\noindication
\noindication
\noindication
\noindication
\noindication
\noindication
\noindication
\noindication
\noindication
\noindication
\noindication
\noindication
\noindication
\noindication
\noindication
\noindication
\noindication
\noindication
\noindication
\noindication
\noindication
\noindication
\noindication
\noindication
\noindication
\noindication
\noindication
\noindication
\noindication
\noindication
\noindication
\noindication
\noindication
\noindication
\noindication
\noindication
\noindication
\noindication
\noindication
\noindication
\noindication
\noindication
\noindication
\noindication
\noindication
\noindication
\noindication
\noindication
\noindication
\noindication
\noindication
\noindication
\noindication
\noindication
\noindication
\noindication
\noindication
\noindication
\noindication
\noindication
\noindication
\noindication
\noindication
\noindication
\noindication
\noindication
\noindication
\noindication
\noindication
\noindication
\noindication
\noindication
\noindication
\noindication
\noindication
\noindication
\noindication
\noindication
\noindication
\noindication
\noindication
\noindication
\noindication
\noindication
\noindication
\noindication
\noindication
\noindication
\noindication
\noindication
\noindication
\noindication
\noindication
\noindication
\noindication
\noindication
\noindication
\noindication
\noindication
\noindication
\noindication
\noindication
\noindication
\noindication
\noindication
\noindication
\noindication
\noindication
\noindication
\noindication
\noindication
\noindication
\noindication
\noindication
\noindication
\noindication
\noindication
\noindication
\noindication
\noindication
\noindication
\noindication
\noindication
\noindication
\noindication
\noindication
\noindication
\noindication
\noindication
\noindication
\noindication
\noindication
\noindication
\noindication
\noindication
\noindication
\noindication
\noindication
\noindication
\noindication
\noindication
\noindication
\noindication
\noindication
\noindication
\noindication
\noindication
\noindication
\noindication
\noindication
\noindication
\noindication
\noindication
\noindication
\noindication
\noindication
\noindication
\noindication
\noindication
\noindication
\noindication
\noindication
\noindication
\noindication
\noindication
\noindication
\noindication
\noindication
\noindication
\noindication
\noindication
\noindication
\noindication
\noindication
\noindication
\noindication
\noindication
\noindication
\noindication
\noindication
\noindication
\noindication
\noindication
\noindication
\noindication
\noindication
\noindication
\noindication
\noindication
\noindication
\noindication
\noindication
\noindication
\noindication
\noindication
\noindication
\noindication
\noindication
\noindication
\noindication
\noindication
\noindication
\noindication
\noindication
\noindication
\noindication
\noindication
\noindication
\noindication
\noindication
\noindication
\noindication
\noindication
\noindication
\noindication
\noindication
\noindication
\noindication
\noindication
\noindication
\noindication
\noindication
\noindication
\noindication
\noindication
\noindication
\noindication
\noindication
\noindication
\noindication
\noindication
\noindication
\noindication
\noindication
\noindication
\noindication
\noindication
\noindication
\noindication
\noindication
\noindication
\noindication
\noindication
\noindication
\noindication
\noindication
\noindication
\noindication
\noindication
\noindication
\noindication
\noindication
\noindication
\noindication
\noindication
\noindication
\noindication
\noindication
\noindication
\noindication
\noindication
\noindication
\noindication
\noindication
\noindication
\noindication
\noindication
\noindication
\noindication
\noindication
\noindication
\noindication
\noindication
\noindication
\noindication
\noindication
\noindication
\noindication
\noindication
\noindication
\noindication
\noindication
\noindication
\noindication
\noindication
\noindication
\noindication
\noindication
\noindication
\noindication
\noindication
\noindication
\noindication
\noindication
\noindication
\noindication
\noindication
\noindication
\noindication
\noindication
\noindication
\noindication
\noindication
\noindication
\noindication
\noindication
\noindication
\noindication
\noindication
\noindication
\noindication
\noindication
\noindication
\noindication
\noindication
\noindication
\noindication
\noindication
\noindication
\noindication
\noindication
\noindication
\noindication
\noindication
\noindication
\noindication
\noindication
\noindication
\noindication
\noindication
\noindication
\noindication
\noindication
\noindication
\noindication
\noindication
\noindication
\noindication
\noindication
\noindication
\noindication
\noindication
\noindication
\noindication
\noindication
\noindication
\noindication
\noindication
\noindication
\noindication
\noindication
\noindication
\noindication
\noindication
\noindication
\noindication
\noindication
\noindication
\noindication
\noindication
\noindication
\noindication
\noindication
\noindication
\noindication
\noindication
\noindication
\noindication
\noindication
\noindication
\noindication
\noindication
\noindication
\noindication
\noindication
\noindication
\noindication
\noindication
\noindication
\noindication
\noindication
\noindication
\noindication
\noindication
\noindication
\noindication
\noindication
\noindication
\noindication
\noindication
\noindication
\noindication
\noindication
\noindication
\noindication
\noindication
\noindication
\noindication
\noindication
\noindication
\noindication
\noindication
\noindication
\noindication
\noindication
\noindication
\noindication
\noindication
\noindication
\noindication
\noindication
\noindication
\noindication
\noindication
\noindication
\noindication
\noindication
\noindication
\noindication
\noindication
\noindication
\noindication
\noindication
\noindication
\noindication
\noindication
\noindication
\noindication
\noindication
\noindication
\noindication
\noindication
\noindication
\noindication
\noindication
\indication{002340}
Vérifier que :
  \begin{enumerate}
  \item $\sup(A+B)=\sup A+\sup B$ ;
  \item $\sup(A \cup B)=\max(\sup A,\sup B)$ ;
  \item $\max(\inf A, \inf B) \leq \sup(A\cap B) \leq \min(\sup A, \sup B)$ si $A\cap B\neq \varnothing$ ;
  \item $\inf(A \cup B)=\min(\inf A, \inf B)$ ;
  \item $\max(\inf A, \inf B) \leq \inf(A\cap B) \leq \min(\sup A, \sup B)$ si $A\cap B\neq \varnothing$ ;
  \end{enumerate}
\finindication
\indication{002341}
Montrer que $J_x$ est un intervalle ouvert ; que $J_x=J_y$ ou $J_x\cap J_y=\varnothing$.
Et penser que $\Qq$ est dénombrable.
\finindication
\indication{002342}
Pour trouver $m$, que prendriez-vous si on voulait seulement $m\in \Rr$ ?
\finindication
\indication{002343}
Revenir à la définition de ce qu'est un ``ensemble fermé'' et de ce
qu'est une ``boule fermée''.
\finindication
\noindication
\noindication
\indication{002346}
Une suite de $l^\infty$ est notée $(x^p)_{p\in \Nn}$, pour chaque $p \ge 0$, $x^p$ est elle m\^eme une suite $x^p = (x^p(0), x^p(1), x^p(2), \ldots)$. 
\finindication
\indication{002347}
Montrer
\begin{itemize}
  \item $\| f \| \le N(f)$ ;
  \item $\|f'\|_\infty \le \| f \|_\infty + \|f\|$ ;
  \item $\|f\|_\infty \le \|f\|$.
\end{itemize}
\finindication
\noindication
\indication{002349}
\begin{itemize}
  \item Montrer $\| f \|_1 \le \|f\|_\infty$.
  \item Par un contre-exemple, montrer qu'il n'existe 
aucune constante $C>0$ tel que $\|f\|_\infty \le C \| f \|_1$ pour tout $f$.
\end{itemize}
\finindication
\indication{002350}
Les seules relations sont :
$$N_1 \le N_2 \le 2N_1 \le 2N_4 \le 2N_3.$$
\finindication
\noindication
\noindication
\noindication
\noindication
\noindication
\noindication
\noindication
\noindication
\noindication
\noindication
\noindication
\noindication
\noindication
\noindication
\noindication
\noindication
\noindication
\noindication
\noindication
\noindication
\noindication
\noindication
\noindication
\noindication
\noindication
\noindication
\noindication
\noindication
\noindication
\noindication
\noindication
\noindication
\noindication
\noindication
\noindication
\noindication
\noindication
\noindication
\noindication
\noindication
\noindication
\noindication
\noindication
\noindication
\noindication
\noindication
\noindication
\noindication
\noindication
\noindication
\noindication
\noindication
\noindication
\noindication
\noindication
\noindication
\noindication
\noindication
\noindication
\noindication
\noindication
\noindication
\noindication
\noindication
\noindication
\noindication
\noindication
\noindication
\noindication
\noindication
\noindication
\noindication
\noindication
\noindication
\noindication
\indication{002370}
\begin{enumerate}
  \item Remarquer si $U_a$ est un voisinage de $a$, alors $A \subset \bigcup_{a\in A} U_a$.
  \item Raisonner par l'absurde et construire une suite $(x_n)$ dont aucun élément  n'est dans $U$ et
une suite $(y_n)$ de $K$. Quitte à extraire une sous-suite se débrouiller pour 
qu'elle converge vers la m\^eme limite.
\end{enumerate}
\finindication
\indication{002371}
Utiliser  qu'un ensemble $K$
est compact si et seulement si de toute suite d'éléments de $K$ on peut extraire une sous-suite convergente vers un élément de $K$.
\finindication
\indication{002372}
Extraire des sous-suites...
\finindication
\noindication
\noindication
\noindication
\indication{002376}
On pourra utiliser la caractérisation de la fermeture par des suites.
\finindication
\indication{002377}
\begin{enumerate}
  \item Utiliser la caractérisation de la fermeture par des suites.
  \item Remarquer que  ``$ \|f(x)\|\to \infty$ quand $\|x\|\to \infty$'' est équivalent à
$$ `` \forall M >0\quad \exists m >0 \forall x \in \Rr^n \quad \quad (x \notin B(0,m) \Rightarrow  f(x)\notin B(0,M)). ''$$
\end{enumerate}
\finindication
\noindication
\noindication
\indication{002380}
\begin{enumerate}
  \item ...
  \item Utiliser l'exercice \ref{exo 2}.
  \item Montrer $f(Y) \subset Y$ puis $Y \subset f(Y)$.
  \item Diamètre zéro implique ensemble réduit à un singleton.
\end{enumerate}
\finindication
\noindication
\noindication
\noindication
\noindication
\noindication
\noindication
\noindication
\noindication
\noindication
\noindication
\noindication
\noindication
\noindication
\noindication
\noindication
\noindication
\noindication
\noindication
\noindication
\noindication
\noindication
\noindication
\noindication
\noindication
\noindication
\noindication
\noindication
\noindication
\noindication
\noindication
\noindication
\noindication
\noindication
\noindication
\noindication
\noindication
\noindication
\noindication
\noindication
\noindication
\noindication
\noindication
\noindication
\noindication
\noindication
\noindication
\noindication
\noindication
\noindication
\indication{002353}
\begin{enumerate}
  \item Utiliser le fait que tout ouvert de $\Rr$ est l'union dénombrable d'intervalles ouverts.
  \item \'Ecrire un intervalle fermé comme union dénombrable d'intervalles
ouverts, puis utiliser la même remarque que ci-dessus.
   \end{enumerate}
\finindication
\indication{002354}
\begin{enumerate}
  \item $\ldots$.
  \item Pour montrer que $c_0$ est fermé, l'écrire comme image réciproque de quelque chose.
   \end{enumerate}
\finindication
\indication{002355}
Montrer que le complémentaire est un ouvert. Si vous le souhaitez, placez-vous dans des espaces métriques.
\finindication
\indication{002356}
\begin{enumerate}
  \item Pour un polynôme $P$, la limite de $P(x)$ ne vaut $\pm \infty$
que lorsque $x$ tend vers $\pm \infty$. 
\end{enumerate}
\finindication
\indication{002357}
\begin{enumerate}
\item Pour le sens direct utiliser la caractérisation de l'adhérence par les suites. Pour le sens réciproque, montrer que l'image réciproque d'un fermé est un fermé. 
   \end{enumerate}
\finindication
\noindication
\noindication
\indication{002360}
\begin{enumerate}
  \item Par l'absurde, considérer $I(x) = \int_0^x f$. Trouver une suite $(p_n)$
telle que $(I(p_n))$ ne soit pas une suite de Cauchy.
  \item Pour montrer que cette intégrale converge utiliser le changement de variable $u=t^2$ puis faire une intégration par partie.
   \end{enumerate}
\finindication
\noindication
\noindication
\noindication
\noindication
\noindication
\noindication
\noindication
\noindication
\noindication
\noindication
\noindication
\noindication
\noindication
\noindication
\noindication
\noindication
\noindication
\noindication
\noindication
\noindication
\noindication
\noindication
\noindication
\noindication
\noindication
\noindication
\noindication
\noindication
\noindication
\noindication
\noindication
\noindication
\noindication
\noindication
\noindication
\noindication
\noindication
\noindication
\noindication
\noindication
\noindication
\noindication
\noindication
\noindication
\noindication
\noindication
\noindication
\indication{002361}
Si la relation est vérifiée montrer que $B$ est continue en $x$ en calculant $B(x+y)-B(x)$.
Si $B$ est continue alors en particulier $B$ est continue en $(0,0)$, fixer 
le $\epsilon$ de cette continuité,...
\finindication
\indication{002362}
La continuité de $L$ sur $E$ équivaut  la continuité en $0$.
Par l'absurde supposer que $L$ n'est pas continue en $0$ et construire une suite
$(x_n)$  qui tend vers $0$ mais avec $(L(x_n))$ non bornée.
\finindication
\indication{002363}
Il faut montrer $f(\lambda x)=\lambda f(x)$ pour $\lambda\in\Rr$.
Le faire pour $\lambda \in \Nn$, puis $\lambda \in \Zz$, puis 
$\lambda \in \Qq$ et enfin $\lambda \in \Rr$.
\finindication
\indication{002364}
\begin{enumerate}
  \item $\| S \| = 1$ ;
  \item $\| T \| = \|g \|_\infty$ ;
  \item $\|u\| = \int_0^1|g|$, on distinguera les cas o\`u $g$ reste de signe constant et $g$ change de signe ;
  \item $\| u\| = \|a_n \|_2$ ;
  \item $\| u \| = \|a \|_\infty$ ;
  \item $\| u \| = 1$.
\end{enumerate}
\finindication
\indication{002365}
$U$ est continue et $\| U \|=1$, $V$ n'est pas continue.
\finindication
\noindication
\indication{002367}
\begin{enumerate}
  \item Montrer d'abord que $X$ se décompose sous la forme $H+\Rr.a$.
  \item ...
  \item Non ! Chercher un contre-exemple dans les exercices précédents.
\end{enumerate}
\finindication
\noindication
\indication{002369}
Montrer que $\|L\|=\pi$.
\finindication
\noindication
\noindication
\noindication
\noindication
\noindication
\noindication
\noindication
\noindication
\noindication
\noindication
\noindication
\noindication
\noindication
\noindication
\noindication
\noindication
\noindication
\noindication
\noindication
\noindication
\noindication
\noindication
\noindication
\noindication
\noindication
\noindication
\noindication
\noindication
\noindication
\noindication
\noindication
\noindication
\noindication
\noindication
\noindication
\noindication
\noindication
\noindication
\noindication
\noindication
\noindication
\noindication
\noindication
\noindication
\noindication
\indication{002395}
\begin{enumerate}
  \item C'est une suite de Cauchy. Essayer de se ramener à une suite de Cauchy de $(\Rr,|.|)$.

  \item Regarder la suite définie par $u_n=-n$.

  \item Comme la première question.
\end{enumerate}
\finindication
\indication{002396}
$f$ est injective uniquement afin que $d$ soit bien une distance. Raisonner par double implication.
Utiliser la caractérisation d'un fermé par les suites.
\finindication
\indication{002397}
\begin{enumerate}
  \item $(X,d_\omega)$ est complet. La démonstration est presque la m\^eme que pour montrer que
$({\cal C}([a,b]), \|.\|_\infty)$ est complet.

  \item Prendre par exemple, la fonction $f_n$ définie sur $[0,1]$ par $f_n(t) = 1$ pour $t\in[0,\frac 12]$, $f_n(t)=(1-n(t-\frac12))$ pour $t\in[\frac 12,\frac 12+\frac 1n]$
et $f(t)=0$ si $t\ge \frac 12+\frac 1n$.
\end{enumerate}
\finindication
\indication{002398}
\begin{enumerate}
  \item Intégrer l'exemple de l'exercice \ref{exocomp}.

  \item Oui cet espace est complet, montrer-le !
\end{enumerate}
\finindication
\indication{002399}
\begin{enumerate}
  \item Prendre la suite $(x^p)$ définie par $x^p= (1,1,\ldots,1,1,0,0,0,\ldots).$
($(x^p)_{p\in\Nn}$ est donc une suite de suite).

  \item Prendre $Y$ l'espace de toutes les suites.

  \item Considérer $x^p= (1,\frac 12,\ldots,\frac1p,0,0,\ldots)$.
\end{enumerate}
\finindication
\indication{002400}
\begin{enumerate}
  \item \'Ecrire ce que donne la définition de ``$(x_n)$ est une suite de Cauchy''
pour $\epsilon=1$, puis $\epsilon = \frac 12$, ..., puis $\epsilon = \frac 1{2^k}$.
Faire la somme. Remarquer que si $T_N = \sum _{k=0}^N u_k$ alors $T_N =  x_{n_{N+1}}-x_{n_0}$.

  \item ...
\end{enumerate}
\finindication
\noindication
\noindication
\noindication
\noindication
\noindication
\noindication
\noindication
\noindication
\noindication
\noindication
\noindication
\noindication
\noindication
\noindication
\noindication
\noindication
\noindication
\noindication
\noindication
\noindication
\noindication
\noindication
\noindication
\noindication
\noindication
\noindication
\noindication
\noindication
\noindication
\indication{002404}
C'est à peu prés la même démonstration que pour le théorème du point fixe d'une fonction contractante.
\finindication
\indication{002405}
Montrer que l'unique point fixe $x$ de $f^n$, est un point fixe de $f$.
Pour cela écrire l'égalité $f^n(x)=x$ et composée habilement cette égalité.
Pour conclure utiliser l'unicité du point fixe de $f^n$.
\finindication
\indication{002406}
Faire soigneusement le calcul : $(T\circ T f)(x)=1+x+\int_0^x\int_0^{t-t^2}f(u-u^2)dudt$.
Se souvenir que $X$ est complet et utiliser l'exercice \ref{exoiter}.
\finindication
\noindication
\noindication
\noindication
\noindication
\noindication
\noindication
\noindication
\noindication
\noindication
\noindication
\noindication
\noindication
\noindication
\indication{002392}
Raisonner par l'absurde et montrer que $\omega_x$ est un ouvert dense.
\finindication
\indication{002393}
\begin{enumerate}
  \item Une application $f :  X \to \Rr$ est \emph{semi-continue inf\'erieurement}
si 
$$\forall \lambda \in \Rr \qquad  \{ x\in X \mid f(x) > \lambda \} \qquad \text{est un ouvert.}$$
De façon équivalente  $f$ est \emph{semi-continue inf\'erieurement} si pour tout $x \in X$
$$\forall \epsilon >0 \quad \exists \delta >0 \quad \forall y \in X \quad (d(x,y)< \delta \Rightarrow f(x)-f(y) < \epsilon).$$
Attention il n'y a pas de valeur absolue autour de $f(x)-f(y)$.

  \item Pour la première question considérer $O_n =  \{ x\in X \mid f(x) > n \}$ et utiliser le théorème de Baire.

  \item Pour l'application utiliser la première question avec la fonction 
$$\phi : B \to \Rr, \text{ définie par } \phi(x) = \sup_{n\in\Nn} |f_n(x)|.$$
\end{enumerate}
\finindication
\noindication
\noindication
\noindication
\noindication
\noindication
\noindication
\noindication
\noindication
\noindication
\noindication
\noindication
\indication{002383}
Utiliser la première question pour les deux suivantes.
\finindication
\noindication
\indication{002385}
Utiliser la partition $X = \mathring A \cup \mathrm{Fr}\, A  \cup (X \setminus \bar A)$
o\`u  $\mathrm{Fr}\, A = \bar A \setminus \mathring A$ est la frontière de $A$.
\finindication
\indication{002386}
Faites un dessin de $T$.
Pour la dernière question, raisonner par l'absurde. O\`u peuvent
s'envoyer les points de la deuxième question ?
\finindication
\indication{002387}
\begin{enumerate}
  \item Pour la surjection, pensez à l'exponentielle ou aux sinus et cosinus...
Pour l'injection, raisonner par l'absurde et utiliser la connexité du cercle privé d'un point.

  \item   Raisonner par l'absurde et utiliser la connexité de $\Rr^2$ privé d'un point.
\end{enumerate}
\finindication
\indication{002388}
 Définir $g : \Rr \longrightarrow \{0,1\}$ tel que $g(x)$ prend  la valeur
qu'a $f$ sur $B_x$.
Montrer pour chaque points de $\Rr\setminus \Qq$, $g$ est constante dans un voisinage 
de ce point,
puis faire la m\^eme chose pour un point de $\Qq$. Conclure.
\finindication
\indication{002389}
\begin{enumerate}
  \item Faire un dessin !

  \item Utiliser le théorème des accroissements finis d'une part. La définition de la dérivée d'autre part.

  \item Utiliser l'exercice \ref{exocon} ou refaire la demonstration.
\end{enumerate}
\finindication
\noindication
\indication{002391}
\begin{enumerate}
  \item Faire un dessin !!

  \item Voir l'exercice \ref{exocon}.

  \item Raisonner par l'absurde. Prendre un chemin qui relie le point $(0,0)$ 
au point $(\frac{1}{2\pi},0)$ (par exemple). Ce chemin va quitter à un instant $t_0$ le segment 
$\{0\}\times [-1,1]$. Chercher une contradiction à ce moment l\`a.

\end{enumerate}
\finindication
\noindication
\noindication
\noindication
\noindication
\noindication
\noindication
\noindication
\noindication
\noindication
\noindication
\noindication
\noindication
\noindication
\noindication
\noindication
\noindication
\noindication
\noindication
\noindication
\noindication
\noindication
\noindication
\noindication
\noindication
\noindication
\noindication
\noindication
\noindication
\indication{002408}
Approcher $f$ par une suite de polynômes, et se rappeler que si l'intégrale d'une fonction positive et continue est nulle alors...
\finindication
\indication{002409}
Raisonner par l'absurde.
\finindication
\indication{002410}
Considérer l'application $\Phi : E \to \Rr^n$ définie par $\Phi = (f_1,\ldots, f_n)$.
\finindication
\indication{002411}
Appliquer le théorème de Stone-Weierstrass.
\finindication
\indication{002412}
Pour la deuxième question :
\begin{enumerate}
  \item Montrer que $\{ f_n \mid n\in \Nn \}$ est équicontinue.
  \item Montrer que $ \{ f_n(x) \mid n\in \Nn \}$ est borné.
  \item Applique le théorème d'Acoli sur le compact $\bar B(0,R)$.
  \item Utiliser le procédé diagonal de Cantor ($R=1,2,3,\ldots$).
\end{enumerate}
\finindication
\indication{002413}
Démarrer avec l'inégalité : 
$$|f_n(x_n)-b| \le |f_n(x_n)-f_n(a)|+|f_n(a)-b|.$$

Si $(f_n)$ n'est pas équicontinue le résultat peut \^etre faux. Prendre $f_n(x) =  (1+x)^n$
et $x_n= \frac 1 n$.
\finindication
\noindication
\indication{002415}
\begin{enumerate}
  \item Pour ouvert et fermé, écrire l'\'equicontinuité pour $\epsilon=1$ en un point $x$ (à fixer).

  \item Ascoli...
\end{enumerate}
\finindication
\indication{002416}
\begin{enumerate}
  \item Pour l'équicontinuité utiliser le théorème des  accroissement finis.
Pour la convergence simple montrer que pour $t$ fixé :
$f_n(t)= \sin (\frac{t}{4n\pi})+o(\frac1n)$.

  \item Montrer que $(f_n)$ ne converge par vers la fonction nulle pour la norme $\|.\|_\infty$
(c'est-\`a-dire il y a convergence simple mais pas convergence uniforme).
Le théorème d'Acoli serait-il faux ? 
\end{enumerate}
\finindication
\noindication
\noindication
\noindication
\noindication
\noindication
\noindication
\noindication
\noindication
\noindication
\noindication
\noindication
\noindication
\noindication
\noindication
\noindication
\noindication
\noindication
\noindication
\noindication
\noindication
\noindication
\noindication
\noindication
\noindication
\noindication
\noindication
\noindication
\noindication
\noindication
\noindication
\noindication
\noindication
\noindication
\noindication
\noindication
\noindication
\noindication
\noindication
\noindication
\noindication
\noindication
\noindication
\noindication
\noindication
\noindication
\noindication
\noindication
\noindication
\noindication
\noindication
\noindication
\noindication
\noindication
\noindication
\noindication
\noindication
\noindication
\noindication
\noindication
\noindication
\noindication
\noindication
\noindication
\noindication
\noindication
\noindication
\noindication
\noindication
\noindication
\noindication
\noindication
\noindication
\noindication
\noindication
\noindication
\noindication
\noindication
\noindication
\noindication
\noindication
\indication{007224}
Sans l'hypothèse de connexité, l'anneau n'est plus intègre.
\finindication
\noindication
\noindication
\noindication
\noindication
\noindication
\noindication
\noindication
\noindication
\noindication
\noindication
\indication{007532}
Pour la dernière question, on pourra montrer que $u\frac{\partial u}{\partial x}+v\frac{\partial v}{\partial x}=0$
puis que $u(u\frac{\partial u}{\partial x}+v\frac{\partial v}{\partial x})=\frac{\partial u}{\partial x}$.
\finindication
\noindication
\noindication
\noindication
\noindication
\noindication
\noindication
\noindication
\noindication
\noindication
\noindication
\noindication
\noindication
\noindication
\noindication
\noindication
\noindication
\noindication
\noindication
\noindication
\noindication
\noindication
\noindication
\noindication
\noindication
\noindication
\noindication
\noindication
\noindication
\noindication
\noindication
\noindication
\noindication
\noindication
\noindication
\noindication
\noindication
\noindication
\noindication
\noindication
\noindication
\noindication
\noindication
\noindication
\noindication
\noindication
\noindication
\noindication
\noindication
\noindication
\noindication
\noindication
\noindication
\noindication
\noindication
\noindication
\noindication
\noindication
\noindication
\noindication
\noindication
\noindication
\noindication
\noindication
\noindication
\noindication
\noindication
\noindication
\noindication
\noindication
\noindication
\noindication
\noindication
\noindication
\noindication
\noindication
\noindication
\noindication
\noindication
\noindication
\noindication
\noindication
\noindication
\noindication
\noindication
\noindication
\noindication
\noindication
\noindication
\noindication
\noindication
\noindication
\noindication
\noindication
\noindication
\noindication
\noindication
\noindication
\noindication
\noindication
\noindication
\noindication
\noindication
\noindication
\noindication
\noindication
\noindication
\noindication
\noindication
\noindication
\noindication
\noindication
\noindication
\noindication
\noindication
\noindication
\noindication
\noindication
\indication{007240}
Par l'absurde, supposer qu'il existe \(a\in \C\setminus \overline{f(\C)}\) et étudier la fonction \(z\mapsto \frac{1}{f(z)-a}\).)
\finindication
\noindication
\noindication
\noindication
\noindication
\indication{007245}
Utiliser le théorème de Morera et faire une disjonction de cas selon la position du triangle par rapport à \(L\). 
\finindication
\noindication
\noindication
\noindication
\noindication
\noindication
\noindication
\noindication
\noindication
\noindication
\noindication
\noindication
\noindication
\noindication
\noindication
\noindication
\noindication
\noindication
\noindication
\noindication
\noindication
\noindication
\noindication
\noindication
\noindication
\noindication
\noindication
\noindication
\noindication
\noindication
\noindication
\noindication
\noindication
\noindication
\noindication
\noindication
\noindication
\noindication
\noindication
\noindication
\noindication
\noindication
\noindication
\noindication
\noindication
\noindication
\noindication
\noindication
\noindication
\noindication
\noindication
\noindication
\noindication
\noindication
\noindication
\noindication
\noindication
\noindication
\noindication
\noindication
\noindication
\noindication
\noindication
\noindication
\noindication
\noindication
\noindication
\noindication
\noindication
\noindication
\noindication
\noindication
\noindication
\noindication
\noindication
\noindication
\noindication
\noindication
\noindication
\noindication
\noindication
\noindication
\noindication
\noindication
\noindication
\noindication
\noindication
\noindication
\noindication
\noindication
\noindication
\noindication
\noindication
\noindication
\noindication
\noindication
\noindication
\noindication
\noindication
\noindication
\noindication
\noindication
\noindication
\noindication
\noindication
\noindication
\noindication
\noindication
\noindication
\noindication
\noindication
\noindication
\noindication
\noindication
\noindication
\noindication
\noindication
\noindication
\noindication
\noindication
\noindication
\noindication
\noindication
\noindication
\noindication
\noindication
\noindication
\noindication
\noindication
\noindication
\noindication
\noindication
\noindication
\noindication
\noindication
\noindication
\noindication
\noindication
\noindication
\noindication
\noindication
\noindication
\noindication
\noindication
\noindication
\noindication
\noindication
\noindication
\noindication
\noindication
\noindication
\noindication
\noindication
\noindication
\noindication
\noindication
\noindication
\noindication
\noindication
\noindication
\noindication
\noindication
\noindication
\noindication
\noindication
\noindication
\noindication
\noindication
\noindication
\noindication
\noindication
\noindication
\noindication
\noindication
\noindication
\noindication
\noindication
\noindication
\noindication
\noindication
\noindication
\noindication
\noindication
\noindication
\noindication
\noindication
\noindication
\noindication
\noindication
\noindication
\noindication
\noindication
\noindication
\noindication
\noindication
\noindication
\noindication
\noindication
\noindication
\noindication
\noindication
\noindication
\noindication
\noindication
\noindication
\noindication
\noindication
\noindication
\noindication
\noindication
\noindication
\noindication
\noindication
\noindication
\noindication
\noindication
\noindication
\noindication
\noindication
\noindication
\noindication
\noindication
\noindication
\noindication
\noindication
\noindication
\noindication
\noindication
\noindication
\noindication
\noindication
\noindication
\noindication
\noindication
\noindication
\noindication
\noindication
\noindication
\noindication
\noindication
\noindication
\noindication
\noindication
\noindication
\noindication
\noindication
\noindication
\noindication
\noindication
\noindication
\noindication
\noindication
\noindication
\noindication
\noindication
\noindication
\noindication
\noindication
\noindication
\noindication
\noindication
\noindication
\noindication
\noindication
\noindication
\noindication
\noindication
\noindication
\noindication
\noindication
\noindication
\noindication
\noindication
\noindication
\noindication
\noindication
\noindication
\noindication
\noindication
\noindication
\noindication
\noindication
\noindication
\noindication
\noindication
\noindication
\noindication
\noindication
\noindication
\noindication
\noindication
\noindication
\noindication
\noindication
\noindication
\noindication
\noindication
\noindication
\noindication
\noindication
\noindication
\noindication
\noindication
\noindication
\noindication
\noindication
\noindication
\noindication
\noindication
\noindication
\noindication
\noindication
\noindication
\noindication
\noindication
\noindication
\noindication
\noindication
\noindication
\noindication
\noindication
\noindication
\noindication
\noindication
\noindication
\noindication
\noindication
\noindication


\newpage

\nocorrection
\correction{000105}
Il ne faut pas se laisser impressionner par l'allure de cette
assertion. En effet $A \Rightarrow B$ est une \'ecriture pour $B
\text{ ou } (\text{non} A)$ ;
 ici $A$ (la proposition $(1=2)$) est fausse, donc $(\text{non} A)$ est vraie
et $B \text{ ou } (\text{non} A)$ l'est \'egalement. Donc
l'assertion $A \Rightarrow B$ est vraie, quand $A$ est fausse et
quelque soit la proposition $B$.
\fincorrection
\correction{000106}
\begin{enumerate}
\item (a) est fausse. Car sa n\'egation qui est 
 $\forall x\in \Rr \ \exists y\in \Rr \quad x+y \leq 0$
est vraie. \'Etant donn\'e $x\in \Rr$ il existe toujours un $y\in\Rr$ tel que
$x+y \leq 0$, par exemple on peut prendre $y=-(x+1)$ et alors $x+y=x-x-1=-1 \leq 0$.

\item (b) est vraie, pour un $x$ donn\'e, on peut prendre (par exemple) $y=-x+1$
et alors $x+y=1>0$.
La n\'egation de (b) est 
 $\exists x\in \Rr \ \forall y\in \Rr \quad x+y \leq 0$.

\item (c) : $\forall x\in \Rr \ \forall y\in \Rr \quad x+y > 0$
est fausse, par exemple $x=-1$, $y=0$. La n\'egation est 
$\exists x\in \Rr\ \exists y\in \Rr\ x+y \leq 0$.

\item (d) est vraie, on peut prendre $x=-1$. La n\'egation est:
$\forall x\in \Rr \ \exists y\in \Rr \quad y^2 \leq x$. 
\end{enumerate}
\fincorrection
\correction{000107}
Dans ce corrig\'e, nous donnons une justification, ce qui n'\'etait
pas demand\'e.
\begin{enumerate}
\item Cette assertion se d\'ecompose de la mani\`ere suivante : ( Pour tout $x\in
  {\Rr}$) ($f(x)\leq 1$). La n\'egation de ``( Pour tout $x\in
  {\Rr}$)" est ``Il existe $x\in {\Rr}$" et la n\'egation de ``($f(x)\leq 1$)" est
$ f(x)>1$. Donc la n\'egation de l'assertion compl\`ete est : ``Il
existe $x\in {\Rr},  f(x)>1$".

\item Rappelons comment se traduit l'assertion ``L'application $f$ est
  croissante" : ``pour tout couple de r\'eels $(x_1,x_2)$, si $x_1\leq x_2$
  alors $f(x_1) \leq f(x_2)$". Cela se d\'ecompose en : ``(pour tout couple de
  r\'eels $x_1$ et $x_2$) ($x_1\leq x_2$ implique $f(x_1) \leq f(x_2)$)".
La n\'egation de la premi\`ere partie est : ``(il existe un couple de
r\'eels $(x_1,x_2)$)" et la n\'egation de la deuxi\`eme partie est :
``($x_1\leq x_2$ et $f(x_1) > f(x_2)$)". Donc la n\'egation de
l'assertion compl\`ete est : ``Il existe $x_1\in \Rr$ et $x_2\in
\Rr$ tels que $x_1 \leq x_2$ et $f(x_1) > f(x_2)$".

\item La n\'egation est : ``l'application $f$ n'est pas croissante ou n'est pas
  positive". On a d\'ej\`a traduit ``l'application $f$ n'est pas croissante",
  traduisons ``l'application $f$ n'est pas positive" : ``il existe $x \in \Rr,
  f(x)<0$". Donc la n\'egation de l'assertion compl\`ete est : ``
Il existe $x_1\in \Rr$ et $x_2\in \Rr$ tels que $x_1<x_2$ et
$f(x_1) \geq f(x_2)$, ou il existe $x \in \Rr,
  f(x)<0$".

\item Cette assertion se d\'ecompose de la mani\`ere suivante : ``(Il existe $x\in
  {\Rr}^+$) ($f(x)\leq 0$)". La n\'egation de la premi\`ere partie est : ``(pour
  tout $x\in  {\Rr}^+$)", et celle de la seconde est :``($f(x)> 0$)".
  Donc la n\'egation de l'assertion compl\`ete est : ``Pour tout $x\in {\Rr}^+$,  $f(x)> 0$".

\item Cette assertion se d\'ecompose de la mani\`ere suivante : 
``($\exists x \in {\Rr}$)($\forall y \in {\Rr}$)($ x<y \Rightarrow f(x)>f(y)$)".  
La n\'egation de la premi\`ere partie est ``($\forall x \in  {\Rr}$)", celle de la
seconde est ``($\exists y \in {\Rr}$)", et celle de la troisi\`eme est
``($ x<y \mathrm{\ et\ } f(x)\leq f(y)$)". Donc la n\'egation de
l'assertion compl\`ete est : `` $\forall x \in {\Rr}, \exists y \in
{\Rr} \text{ tels que } x<y \text{ et } f(x)\leq f(y)$".
\end{enumerate}
\fincorrection
\correction{000108}
\begin{enumerate}
  \item $\Leftarrow$
  \item $\Leftrightarrow$
  \item $\Rightarrow$
\end{enumerate}
\fincorrection
\correction{000109}
\begin{enumerate}
\item
 Cette proposition est vraie. En effet soit $\epsilon > 0$,  d\'efinissons $M_1 = (\frac{2}{\epsilon},0)  \in F_1$ et
 $M_2 = (\frac{2}{\epsilon},\frac{\epsilon}{2}) \in F_2$, alors $M_1M_2=\frac{\epsilon}{2} < \epsilon$. Ceci \'etant vrai quelque soit
$\epsilon >0$ la proposition est donc d\'emontr\'ee.


\item
Soit deux points fix\'es $M_1$, $M_2$ v\'erifiant cette
proposition, la distance $d= M_1M_2$ est
 aussi petite que l'on veut donc elle est nulle, donc $M_1 = M_2$ ; or les ensembles $F_1$ et $F_2$
sont disjoints. Donc la proposition est fausse. La n\'egation de
cette proposition est :
$$ \forall M_1 \in F_1 \ \  \forall M_2 \in F_2 \quad \exists \epsilon \in ]0,+\infty[ \quad  \quad M_1M_2 \geqslant \epsilon $$ et cela exprime le fait que les ensembles $F_1$ et $F_2$ sont disjoints.

\item
Celle ci est \'egalement fausse, en effet supposons qu'elle soit
vraie, soit alors $\epsilon$ correspondant \`a cette proposition.
Soit $M_1=(\epsilon +2,0)$ et $M_2 = (1,1)$, on a $M_1M_2 >
\epsilon+1$ ce qui est absurde. La n\'egation est :
$$\forall \epsilon \in ]0,+\infty[ \quad  \exists M_1 \in F_1 \ \  \exists M_2 \in F_2 \quad   \quad M_1M_2 \geqslant \epsilon $$
C'est-\`a-dire que l'on peut trouver deux points aussi
\'eloign\'es l'un de l'autre  que l'on veut.

\item
Cette proposition est vraie, il suffit de choisir
$\epsilon=M_1M_2+1$. Elle signifie que la distance entre deux
points donn\'es est un nombre fini !

\end{enumerate}
\fincorrection
\correction{000110}
``Il existe un habitant de la rue du Havre qui a les yeux bleus, qui
ne gagnera pas au loto ou qui prendra sa retraite apr\`es 50 ans.''
\fincorrection
\correction{000111}
\begin{enumerate}
  \item $P$ et non $Q$;
  \item ``non P ou $Q$" ce qui la même chose que ``$P \Rightarrow Q$";
  \item (non $P$) ou ((non $Q$) ou (non $R$)) (on peut supprimer les parenthèses);
  \item non $P$ et (non $Q$ ou non $R$) (ici les parenthèses sont importantes);
  \item $P$ et $Q$ et $R$ et non $S$;
\end{enumerate}
\fincorrection
\correction{000112}
\begin{enumerate}
    \item ``Il existe un triangle rectangle qui n'a pas d'angle droit." Bien sûr cette dernière phrase est fausse !
    \item ``Il existe une \'ecurie dans laquelle il y a (au moins) un cheval
dont la couleur n'est pas noire."
    \item Sachant que la proposition en langage math\'ematique s'\'ecrit
$$\forall x\in\Zz \ \  \exists y\in\Zz\ \  \forall z\in\Zz \quad (z<x \Rightarrow z<x+1),$$
la n\'egation est
$$\exists x\in\Zz\ \  \forall y\in\Zz\ \  \exists z\in\Zz \quad (z<x \text{ et } z\geq x+1).$$
    \item $\exists \epsilon>0\ \  \forall \alpha>0 \quad (|x-7/5|<\alpha \text{ et } |5x-7|\geq\epsilon).$
\end{enumerate}
\fincorrection
\nocorrection
\nocorrection
\nocorrection
\nocorrection
\nocorrection
\nocorrection
\correction{000119}
Remarquons d'abord que pour $n \in \Nn$, $\frac{2n+1}{n+2} \leq 2$
car $2n+1 \leq 2(n+2)$.
\'Etant donn\'e $\epsilon > 0$, nous avons donc 
$$\forall n \in \Nn \quad \frac{2n+1}{n+2} < 2 + \epsilon$$
Maintenant nous cherchons une condition sur $n$ pour que
l'in\'egalit\'e 
$$2-\epsilon < \frac{2n + 1}{n + 2}$$
soit vraie.
\begin{align*}
2-\epsilon < \frac{2n + 1}{n + 2} 
    &\Leftrightarrow (2-\epsilon)(n+2) < 2n+1 \\
    &\Leftrightarrow 3  < \epsilon (n+2) \\
   &\Leftrightarrow n >  \frac{3}{\epsilon}-2 \\
\end{align*}

Ici $\epsilon$ nous est donn\'e, nous prenons un $N\in \Nn$ tel
que $N > \frac{3}{\epsilon}-2$, alors pour tout $n \geq N$ nous avons
$n \geq N >  \frac{3}{\epsilon}-2$ et par
cons\'equent: $2-\epsilon < \frac{2n + 1}{n + 2}$.
Conclusion: \'etant donn\'e $\epsilon > 0$, nous avons trouv\'e un 
$N\in \Nn$ tel que pour tout $n \geq N$ on ait
$2-\epsilon < \frac{2n + 1}{n + 2}$ et $\frac{2n+1}{n+2} < 2 + \epsilon$.

En fait nous venons de prouver que la suite 
de terme $(2n+1)/(n+2)$ tend vers $2$ quand $n$ tend vers $+\infty$.
\fincorrection
\correction{000120}
\begin{enumerate}
  \item $\exists M \in \Rr \quad \forall x \in \Rr \qquad f(x) \leq M$;
  \item $\exists M \in \Rr\quad \exists m \in \Rr \quad \forall x \in \Rr \qquad m \leq f(x) \leq M$;
  \item $\forall x \in \Rr \qquad f(x) = f(-x)$;
  \item $\forall x \in \Rr \qquad f(-x) = -f(x)$;
  \item $\forall x \in \Rr \qquad f(x) \not= 0$;
  \item $\exists a \in \Rr^* \quad \forall x \in \Rr \qquad f(x+a) = f(x)$;
  \item $\forall (x,y) \in \Rr^2 \qquad (x\leq y \Rightarrow f(x) \leq f(y))$;
  \item $\forall (x,y) \in \Rr^2 \qquad (x < y \Rightarrow f(x) > f(y))$;
  \item $\exists x \in \Rr \qquad  f(x) \not= 0$;
  \item $\forall (x,y) \in \Rr^2 \qquad (x\not= y \Rightarrow f(x) \not= f(y))$;
  \item $\forall n\in \Nn \quad \exists x \in \Rr \qquad f(x)=n$;
  \item $\forall x \in \Rr \qquad f(x) \leq g(x)$;
  \item $\exists x \in \Rr \qquad f(x) > g(x)$.
\end{enumerate}
\fincorrection
\correction{005103}
\begin{enumerate}
 \item 

\begin{enumerate}
 \item  $(f=Id_{\mathcal{P}}\Leftrightarrow\forall M\in\mathcal{P},\;f(M)=M)$ et $(f\neq Id_{\mathcal{P}}\Leftrightarrow\exists
M\in\mathcal{P}/\;f(M)\neq M)$.
  \item $(f\;\mbox{a au moins un point fixe}\Leftrightarrow \exists M\in\mathcal{P}/\;f(M)=M)$ et $(f\;\mbox{n'a pas de point
fixe}\Leftrightarrow \forall M\in\mathcal{P},\;f(M)\neq M)$.

\end{enumerate}
Constatez que les phrases $f(M)=M$ ou $f(M)\neq M$ \textbf{n'ont aucun sens} si elles ne sont pas accompagnées de
quantificateurs.
 \item 
  \begin{enumerate}
  \item $(f=0\Leftrightarrow\forall x\in\Rr,\;f(x)=0)$ et $(f\neq0\Leftrightarrow\exists x\in\Rr/\;f(x)\neq0)$.
  \item (L'équation $f(x)=0$ a (au moins) une solution si et seulement si $\exists x\in\Rr/\;f(x)=0)$ et (l'équation
$f(x)=0$ n'a pas de solution si et seulement si $\forall x\in\Rr/\;f(x)\neq0)$.
  \item (L'équation $f(x)=0$ a exactement une solution si et seulement si $\exists!x\in\Rr/\;f(x)=0)$ et (l'équation
$f(x)=0$ n'a pas exactement une solution si et seulement si $\forall x\in\Rr/\;f(x)\neq0$ ou
$\exists(x,x')\in\Rr^2/\;(x\neq x'\;\mbox{et}\;f(x)=f(x')=0)$.
\end{enumerate}
 \item  
  \begin{enumerate}
  \item  $((u_n)_{n\in\Nn}\;\mbox{bornée}\Leftrightarrow\exists M\in\Rr/\;\forall n\in\Nn,\;|u_n|\leq M)$ et
$((u_n)_{n\in\Nn}\;\mbox{non bornée}\Leftrightarrow\forall M\in\Rr/\;\exists n\in\Nn,\;|u_n|>M)$.
   \item $((u_n)_{n\in\Nn}\;\mbox{croissante}\Leftrightarrow\forall n\in\Nn/\;u_{n+1}\geq u_n)$ et
$((u_n)_{n\in\Nn}\;\mbox{non croissante}\Leftrightarrow\exists n\in\Nn/\;u_{n+1}<u_n)$.

  \item  $((u_n)_{n\in\Nn}\;\mbox{monotone}\Leftrightarrow(\forall n\in\Nn/\;u_{n+1}\geq u_n)\;\mbox{ou}\;(\forall
n\in\Nn/\;u_{n+1}\leq u_n))$ et $((u_n)_{n\in\Nn}\;\mbox{non monotone}\Leftrightarrow((\exists
n\in\Nn/\;u_{n+1}<u_n)\;\mbox{et}\;(\exists
n\in\Nn/\;u_{n+1}>u_n))$.
  \end{enumerate}

\end{enumerate}

\fincorrection
\correction{005104}
Le contraire de $x\geq3$ est $x<3$. Le contraire de $0<x\leq 2$ est $((x\leq0)\;\mbox{ou}\;x>2)$.
\fincorrection
\correction{005105}
\begin{enumerate}
 \item  Oui. Dans les deux cas, chaque fois que l'on se donne un réel $x_0$, $f(x_0)$ et $g(x_0)$ sont tous deux nuls.
 \item  Non. La deuxième affirmation implique la première mais la première n'implique pas la deuxième. La première
phrase est la traduction avec des quantificateurs de l'égalité $fg=0$. La deuxième phrase est la traduction avec
quantificateurs de $(f=0\;\mbox{ou}\;g=0)$.
Voici un exemple de fonctions $f$ et $g$ toutes deux non nulles dont le produit est nul.
Soient $\begin{array}[t]{cccc}
f~:&\Rr&\rightarrow&\Rr\\
 &x&\mapsto&\left\{
\begin{array}{l}
0\;\mbox{si}\;x<0\\
x\;\mbox{si}\;x\geq0
\end{array}
\right.
\end{array}$ et $\begin{array}[t]{cccc}
g~:&\Rr&\rightarrow&\Rr\\
 &x&\mapsto&\left\{
\begin{array}{l}
0\;\mbox{si}\;x>0\\
x\;\mbox{si}\;x\leq0
\end{array}
\right.
\end{array}$. Pour chaque valeur de $x$, on a soit $f(x)=0$ (quand $x\leq0$), soit $g(x)=0$ (quand $x\geq0$). On a
donc~:~$\forall x\in\Rr,\;(f(x)=0\;\mbox{ou}\;g(x)=0)$ ou encore $\forall x\in\Rr,\;f(x)g(x)=0$ ou enfin, $fg=0$.
Cependant, $f(1)=1\neq0$ et donc $f\neq0$, et $g(-1)=-1\neq0$ et donc $g\neq0$. Ainsi, on n'a pas $(f=0\;\mbox{ou}\;g=0)$ ou
encore, on n'a pas $((\forall x\in\Rr,\;f(x)=0)\;\mbox{ou}\;(\forall x\in\Rr,\;g(x)=0))$.
\end{enumerate}
\fincorrection
\nocorrection
\correction{000122}
Nous allons d\'emontrer l'assertion $1.$ de deux mani\`eres diff\'erentes.
\begin{enumerate}
  \item Tout d'abord de fa\c{c}on ``directe". Nous supposons
que $A$ et $B$ sont tels que $A\cap B = A \cup B$. Nous devons montrer 
que $A=B$. 

Pour cela \'etant donn\'e $x \in A$ montrons qu'il est aussi dans $B$.
Comme $x\in A$ alors $x \in A\cup B$ donc $x \in A \cap  B$ (car $A\cup B= A \cap B$). Ainsi $x \in B$. 

Maintenant nous prenons $x\in B$ et le m\^eme raisonnement implique $x \in A$.
Donc tout \'el\'ement de $A$ est dans $B$ et tout \'el\'ement de $B$ est dans $A$.
Cela veut dire $A=B$.

  \item Ensuite, comme demand\'e, nous le montrons par contraposition.
Nous supposons que $A\not= B$ et nous devons montrer que 
$A\cap B \not= A\cup B$.

Si $A\not= B$ cela veut dire qu'il existe un \'el\'ement $x \in A\setminus B$
ou alors un \'el\'ement $x\in B \setminus A$. Quitte \`a \'echanger $A$ et $B$, nous supposons qu'il existe $x\in A \setminus B$. Alors $x \in A \cup B$ mais
$x \notin A\cap B$. Donc $A\cap B \not= A\cup B$.
\end{enumerate}
\fincorrection
\correction{000123}
\begin{align*}
x\in\complement (A\cup B) & \Leftrightarrow x\notin A\cup B \\
&\Leftrightarrow x\notin A \text{ et } x\notin B\\
&\Leftrightarrow x\in\complement A \text{ et } x\in\complement B\\
&\Leftrightarrow x\in \complement A \cap \complement B.\\
\end{align*}
\begin{align*}
x\in\complement (A\cap B) & \Leftrightarrow x\notin A\cap B\\
&\Leftrightarrow x\notin A \text{ ou } x\notin B\\
&\Leftrightarrow x\in\complement A \text{ ou } x\in\complement\\
&\Leftrightarrow x\in \complement A \cup \complement B.\\
\end{align*}
\fincorrection
\correction{000124}
Montrons quelques assertions.

$f(A\cap B) \subset f(A)\cap f(B)$.\\
Si $y\in f(A\cap B)$, il existe $x\in A\cap B$ tel que $y=f(x)$,
or $x\in A$ donc $y=f(x) \in f(A)$ et de m\^eme $x\in B$ donc
$y\in f(B)$. D'o\`u $y\in f(A)\cap f(B)$. Tout \'el\'ement de
$f(A\cap B)$ est un \'el\'ement de $f(A)\cap f(B)$ donc $f(A\cap
B) \subset f(A)\cap f(B)$.

Remarque : l'inclusion r\'eciproque est fausse. Exercice : trouver
un contre-exemple.

\bigskip

$f^{-1}(F\setminus A) = E\setminus f^{-1}(A)$.\\
\begin{align*}
x\in f^{-1}(F\setminus A) &\Leftrightarrow f(x) \in F\setminus A\\
&\Leftrightarrow f(x) \notin A\\
&\Leftrightarrow x \notin f^{-1}(A) \quad \text{ car } f^{-1}(A) = \{ x\in E \ / \ f(x) \in A \}\\
&\Leftrightarrow x\in E\setminus f^{-1}(A)\\
\end{align*}
\fincorrection
\nocorrection
\nocorrection
\nocorrection
\nocorrection
\nocorrection
\nocorrection
\nocorrection
\nocorrection
\nocorrection
\nocorrection
\nocorrection
\correction{000136}
$I_1=3\;\mbox{ et
}\; I_2=\left[-2,5\right].$
\fincorrection
\correction{000137}
$I=\left[0,2\right]\;\mbox{
et }\;J=\left]1,+\infty\right[.$
\fincorrection
\nocorrection
\nocorrection
\nocorrection
\nocorrection
\nocorrection
\nocorrection
\correction{000144}
\begin{enumerate}
  \item $B\setminus A \subset X \subset B$.
  \item $B \subset X \subset B \cup \complement{A}$.
\end{enumerate}
\fincorrection
\nocorrection
\nocorrection
\nocorrection
\correction{005112}
\begin{enumerate}
 \item  Si $A=B=\varnothing$ alors $A\Delta B=\varnothing=A\cap B$.
Si $A\Delta B=A\cap B$, supposons par exemple $A\neq\varnothing$.
Soit $x\in A$. Si $x\in B$, $x\in A\cap B=A\Delta B$ ce qui est absurde et si $x\notin B$, $x\in A\Delta B=A\cap B$ ce
qui est absurde. Donc $A=B=\varnothing$. Finalement, $A\Delta B=A\cap B\Leftrightarrow A=B=\varnothing$.

 \item  Par distributivité de $\cap$ sur $\cup$,

\begin{align*}
(A\cup B)\cap(B\cup C)\cap(C\cup A)&=((A\cap B)\cup(A\cap C)\cup(B\cap B)\cup(B\cap C))\cap(C\cup A)\\
 &=((A\cap C)\cup B)\cap(C\cup A)\;(\mbox{car}\;B\cap B=B\;\mbox{et}\;A\cap B\subset B\;\mbox{et}\;B\cap C\subset B)\\
 &=\left((A\cap C)\cap C\right)\cup\left((A\cap C)\cap A\right)\cup\left(B\cap C\right)\cup\left(B\cap A\right)\\
 &=(A\cap C)\cup(A\cap C)\cup(B\cap C)\cup(B\cap A)\\
 &=(A\cap B)\cup(B\cap C)\cup(C\cap A)
\end{align*}

 \item  $A\Delta B=(A\setminus B)\cup(B\setminus A)=(B\setminus A)\cup(A\setminus B)=B\Delta A$.

 \item 
\begin{align*}
x\in(A\Delta B)\Delta C&\Leftrightarrow x\;\mbox{est dans}\;A\Delta B\;\mbox{ou dans}\;C\;\mbox{mais pas dans les deux}\\
 &\Leftrightarrow((x\in A\;\mbox{et}\;x\notin B\;\mbox{et}\;x\notin C)\;\mbox{ou}\;(x\in B\;\mbox{et}\;x\notin
A\;\mbox{et}\;x\notin C)\;\mbox{ou}\;(x\in C\;\mbox{et}\;x\notin A\Delta B)\\
 &\Leftrightarrow x\;\mbox{est dans une et une seule des trois parties ou dans les trois}.
\end{align*}
Par symétrie des rôles de $A$, $B$ et $C$, $A\Delta(B\Delta C)$ est également l'ensemble des éléments qui sont dans une
et une seule des trois parties $A$, $B$ ou $C$ ou dans les trois. Donc $(A\Delta B)\Delta C=A\Delta(B\Delta C)$. Ces
deux ensembles peuvent donc se noter une bonne fois pour toutes $A\Delta B\Delta C$.

 \item  $A=B\Rightarrow A\setminus B=\varnothing$ et $B\setminus A=\varnothing\Rightarrow A\Delta B =\varnothing$.

$A\neq B\Rightarrow\exists x\in E/\;((x\in A\;\mbox{et}\;x\notin B)\;\mbox{ou}\;(x\notin A\;\mbox{et}\;
x\in B))\Rightarrow\exists x\in E/\;x\in(A\setminus B)\cup(B\setminus A)=A\Delta B\Rightarrow A\Delta B\neq\varnothing$.

 \item 
\begin{itemize}
\item[$\Leftarrow$] Immédiat.
\item[$\Rightarrow$] Soit $x$ un élément de $A$.

Si $x\notin C$ alors $x\in A\Delta C=B\Delta C$ et donc $x\in B$ car $x\notin C$.

Si $x\in C$ alors $x\notin A\Delta C=B\Delta C$. Puis $x\notin B\Delta C$ et $x\in C$ et donc $x\in B$. Dans tous les
cas, $x$ est dans $B$. Tout élément de $A$ est dans $B$ et donc $A\subset B$.
En échangeant les rôles de $A$ et $B$, on a aussi $B\subset A$ et finalement $A=B$.
\end{itemize}
\end{enumerate}
\fincorrection
\correction{005113}
\begin{enumerate}
 \item  Soit $x\in E$.

\begin{align*}
x\in f\left(\bigcup_{i\in I}A_i\right)&\Leftrightarrow\exists y\in\bigcup_{i\in I}A_i/\;x=f(y)\Leftrightarrow\exists i\in I,\;\exists y\in
A_i/\;x=f(y)\\
 &\Leftrightarrow\exists i\in I/\;x\in f(A_i)\Leftrightarrow x\in\bigcup_{i\in I}f(A_i)
\end{align*}
Donc

\begin{center}
\shadowbox{
$f\left(\displaystyle\bigcup_{i\in I}A_i\right)=\displaystyle\bigcup_{i\in I}f(A_i)$.
}
\end{center}

 \item  Soit $x\in E$.

\begin{align*}
x\in f\left(\bigcap_{i\in I}A_i\right)&\Leftrightarrow\exists y\in\bigcap_{i\in I}A_i/\;x=f(y)\Leftrightarrow\exists y\in E/\;\forall i\in I,\;
y\in A_i\;\mbox{et}\;x=f(y)\\
 &\Rightarrow\forall i\in I/\;\exists y\in A_i/\;x=f(y)\Leftrightarrow\forall i\in I/\;x\in f(A_i)\\
 &\Leftrightarrow x\in\bigcap_{i\in I}f(A_i)
\end{align*}
Donc

\begin{center}
\shadowbox{
$f\left(\displaystyle\bigcap_{i\in I}A_i\right)\subset\displaystyle\bigcap_{i\in I}f(A_i)$.
}
\end{center}
L'inclusion contraire n'est pas toujours vraie. Par exemple, pour $x$ réel on pose $f(x)=x^2$ puis $A=\{-1\}$ et
$B=\{1\}$. $A\cap B=\varnothing$ et donc $f(A\cap B)=\varnothing$ puis $f(A)=f(B)=\{1\}$ et donc $f(A)\cap f(B)=\{1\}$.

 \item  Il n'y a aucune inclusion vraie entre $f(E\setminus A)$ et $F\setminus f(A)$. Par exemple, soit
$\begin{array}[t]{cccc}
f~:&\Rr&\rightarrow&\Rr\\
 &x&\mapsto&x^2
\end{array}$ et $A=[-1,2]$.
$f(A)=[0,4]$ et donc $C_{\Rr}(f(A))=]-\infty,0[\cup]4,+\infty[$ mais $f(C_{\Rr}A)=f(]-\infty,-1[\cup]2,+\infty[)
=]1,+\infty[$ et aucune inclusion entre les deux parties n'est vraie.
 \item  Soit $x\in E$.

$$x\in f^{-1}\left(\bigcap_{i\in I}B_i\right)\Leftrightarrow f(x)\in\bigcap_{i\in I}B_i\Leftrightarrow\forall i\in I,\;f(x)\in B_i\Leftrightarrow\forall i\in
I,\;x\in f^{-1}(B_i)\Leftrightarrow x\in\bigcap_{i\in I}f^{-1}(B_i).$$
Donc,

\begin{center}
\shadowbox{
$f^{-1}(\displaystyle\bigcap_{i\in I}B_i)=\displaystyle\bigcap_{i\in I}f^{-1}(B_i)$.
}
\end{center}

 \item  Soit $x\in E$.

$$x\in f^{-1}(\bigcup_{i\in I}B_i)\Leftrightarrow f(x)\in\bigcup_{i\in I}B_i\Leftrightarrow\exists i\in I,\;f(x)\in B_i\Leftrightarrow\exists i\in
I,\;x\in f^{-1}(B_i)\Leftrightarrow x\in\bigcup_{i\in I}f^{-1}(B_i).$$
Donc,

\begin{center}
\shadowbox{
$f^{-1}(\displaystyle\bigcup_{i\in I}B_i)=\displaystyle\bigcup_{i\in I}f^{-1}(B_i)$.
}
\end{center}
 \item  Soit $x\in E$.

$$x\in f^{-1}(F\setminus B_i)\Leftrightarrow f(x)\in F\setminus B_i\Leftrightarrow f(x)\notin B_i\Leftrightarrow x\notin f^{-1}(B_i)\Leftrightarrow x\in
E\setminus f^{-1}(B_i).$$
Donc,

\begin{center}
\shadowbox{
$f^{-1}(F\setminus B_i)=E\setminus f^{-1}(B_i)$.
}
\end{center}
\end{enumerate}
\fincorrection
\correction{005117}
\begin{enumerate}
 \item  Il y a l'injection triviale $\begin{array}[t]{cccc}
f~:&E&\rightarrow&\mathcal{P}(E)\\
 &x&\mapsto&\{x\}
\end{array}$.
 \item  Soit $f$ une application quelconque de $E$ dans $\mathcal{P}(E)$. Montrons que $f$ ne peut être
surjective.
Soit $A=\{x\in E/\;x\notin f(x)\}$. Montrons que $A$ n'a pas d'antécédent par $f$. Supposons par
l'absurde que $A$ a un antécédent $a$. Dans ce cas, où est $a$~?~

$$a\in A\Rightarrow a\notin f(a)=A,$$
ce qui est absurde et

$$a\notin A\Rightarrow a\in f(a)=A,$$
ce qui est absurde. Finalement, $A$ n'a pas d'antécédent et $f$ n'est pas surjective. On a montré le théorème de
\textsc{Cantor}~:~pour tout ensemble $E$ (vide, fini ou infini), il n'existe pas de bijection de $E$ sur
$\mathcal{P}(E)$.
\end{enumerate}
\fincorrection
\nocorrection
\nocorrection
\nocorrection
\nocorrection
\nocorrection
\correction{000150}
Par l'absurde, supposons qu'il existe $p\in \Nn$ tel que $f=f_p$.
Deux applications sont \'egales si et seulement si elles prennent
les m\^emes valeurs.
$$\forall n\in \Nn\ \ f(n) = f_p(n).$$
En particulier pour $n=p$, $f(p)=f_p(p)$. D'autre part la
d\'efinition de $f$ nous donne $f(p) = f_p(p)+1$. Nous obtenons
une contradiction car $f(p)$ ne peut prendre deux valeurs
distinctes. En conclusion, quelque soit $p\in \Nn$, $f\not= f_p$.
\fincorrection
\correction{000151}
\begin{enumerate}
\item
Montrons en fait la contrapos\'ee.

S'il existe $i$ tel que $p_i$ divise $N=p_1p_2 \ldots p_r +1$ ($i$
est fix\'e) alors il existe $k \in \Zz$ tel que $N = kp_i$ donc
$$p_i(k-p_1p_2\ldots p_{i-1}p_{i+1}\ldots p_r) = 1$$
 soit $p_iq = 1$ (avec
$q = k-p_1p_2\ldots p_{i-1}p_{i+1}\ldots p_r $ un nombre entier).
Donc $p_i \in \Zz$ et $1/p_i = q \in \Zz$,  alors $p_i$ vaut $1$
ou $-1$. Et donc $p_i$ n'est pas un nombre premier.

Conclusion : par contraposition il est vrai que $N$ n'est
divisible par aucun des $p_i$

\item Raisonnons par l'absurde : s'il n'existe qu'un nombre fini $r$ de nombres premiers $p_1,\ldots,p_r$
alors $N=p_1p_2 \ldots p_r +1$ est un nombre premier car divisible
par aucun nombre premier autre que lui m\^eme (c'est le 1.).

 Mais $N$ est strictement sup\'erieur \`a tous les $p_i$. Conclusion on a construit un
nombre premier $N$ diff\'erent des $p_i$, il y a donc au moins
$r+1$ nombres premiers, ce qui est absurde.
\end{enumerate}
\fincorrection
\nocorrection
\correction{000153}
R\'edigeons la deuxi\`eme \'egalit\'e.
Soit $\mathcal{A}_n$, $n\in \Nn^*$ l'assertion suivante:
$$(\mathcal{A}_n) \ \ \ \sum_{k=1}^n k^2= \frac{n(n+1)(2n+1)}{6}.$$
\begin{itemize}
  \item $\mathcal{A}_0$ est vraie ($1=1$).
  \item \'Etant donn\'e $n\in\Nn^*$ supposons que $\mathcal{A}_n$
soit vraie. Alors
\begin{align*}
\sum_{k=1}^{n+1} k^2
  &= \sum_{k=1}^n k^2 \ \ + (n+1)^2 \\
  &= \frac{n(n+1)(2n+1)}{6} + (n+1)^2\\
  &= \frac{n(n+1)(2n+1)+6(n+1)^2}{6} \\
  &= \frac{(n+1)(n(2n+1)+6(n+1))}{6} \\
  &= \frac{(n+1)(n+2)(2(n+1)+1)}{6} \\
\end{align*}
Ce qui prouve $\mathcal{A}_{n+1}$.
  \item Par le principe de r\'ecurrence nous venons de montrer 
que $\mathcal{A}_n$ est vraie pour tout $n\in\Nn^*$.
\end{itemize}
\fincorrection
\nocorrection
\correction{000155}
\begin{enumerate}
    \item Montrons par r\'ecurrence $\forall n \in \Nn\  x_n > 3$.
Soit l'hypoth\`ese de r\'ecurrence :
$$(\mathcal{H}_n) : \quad x_n >3.$$

\begin{itemize}
    \item[$\bullet$] La proposition $\mathcal{H}_0$ est vraie car $x_0 = 4 > 3$.

    \item[$\bullet$] Soit $n\geq 0$, supposons $\mathcal{H}_n$ vraie et montrons
que $\mathcal{H}_{n+1}$ est alors vraie.

$$x_{n+1}-3 = \frac{2{x_n}^2-3}{x_n+2}-3 = \frac{2{x_n}^2-3x_n-9}{x_n+2}.$$
Par hypoth\`ese de r\'ecurrence $x_n > 3$, donc $x_n+2 > 0$ et
$2{x_n}^2-3x_n-9>0$ (ceci par \'etude de la fonction $x \mapsto
2{x}^2-3x-9$ pour $x>3$). Donc $x_{n+1}-3 $ et $\mathcal{H}_{n+1}$
est vraie.

    \item[$\bullet$] Nous avons montré
$$\forall n \in \Nn \quad \mathcal{H}_{n} \Rightarrow \mathcal{H}_{n+1}$$
et comme $\mathcal{H}_{0}$ est vraie alors $\mathcal{H}_{n}$ est
vraie quelque soit $n$. Ce qui termine la d\'emonstration.
\end{itemize}

    \item Montrons que  $x_{n+1}-3 - \frac{3}{2}(x_n-3)$ est positif.
$$x_{n+1}-3 - \frac{3}{2}(x_n-3) =
\frac{2{x_n}^2-3}{x_n+2}-\frac{3}{2}(x_n-3) =
\frac{1}{2}\frac{{x_n}^2-3x_n}{x_n+2}$$ Ce dernier terme est
positif car $x_n >3$.

    \item
 Montrons par r\'ecurrence $\forall n \in \Nn\  x_n > \left(\frac{3}{2}\right)^n+3$.
Soit notre nouvelle l'hypoth\`ese de r\'ecurrence :
$$(\mathcal{H}_n) \quad x_n >\left( \frac{3}{2}\right)^n+3.$$

\begin{itemize}
    \item[$\bullet$] La proposition $\mathcal{H}_0$ est vraie.

    \item[$\bullet$] Soit $n\geq 0$, supposons que $\mathcal{H}_n$ vraie et montrons
que $\mathcal{H}_{n+1}$ est v\'erifi\'ee.

D'apr\`es la question pr\'ec\'edente $ x_{n+1}-3 >
\frac{3}{2}(x_n-3)$ et par hypoth\`ese de r\'ecurrence $x_n
>\left( \frac{3}{2}\right)^n+3$ ; en r\'eunissant ces deux
in\'egalit\'es nous avons $ x_{n+1}-3 >
\frac{3}{2}(\left(\frac{3}{2}\right)^n) =
\left(\frac{3}{2}\right)^{n+1}$.

    \item[$\bullet$] Nous concluons en r\'esumant la situation :\\
$\mathcal{H}_{0}$ est vraie, et $\mathcal{H}_{n} \Rightarrow
\mathcal{H}_{n+1}$ quelque soit $n$. Donc $\mathcal{H}_{n}$ est
toujours vraie.
\end{itemize}
    \item La suite $(x_n)$ tend vers $+\infty$ et n'est donc pas convergente.
\end{enumerate}
\fincorrection
\correction{000156}
Montrons par r\'ecurrence sur $n \geqslant 1$ la proposition suivante :
$$\mathcal{H}_n :  \quad n \text{\  droites en position g\'en\'erale d\'ecoupent le plan en \ } R_n = \frac{n(n+1)}{2}+1
\text{\  r\'egions.}$$

\begin{itemize}
\item[$\bullet$] pour $n=1$ alors une droite divise le plan en deux r\'egions. $\mathcal{H}_1$ est vraie.

\item[$\bullet$] Soit $n\geqslant 2$ et supposons que $\mathcal{H}_{n-1}$ soit vraie, et montrons $\mathcal{H}_n$.
Soient $\Delta_1,\ldots,\Delta_n$ $n$ droites en position
g\'en\'erale, la droite $\Delta_n$ rencontre les droites
$\Delta_1,\ldots,\Delta_{n-1}$ en $n-1$ points, donc $\Delta_n$
traverse (et d\'ecoupe en deux) $n$ r\'egions du d\'ecoupage
$\Delta_1,\ldots,\Delta_{n-1}$. Le d\'ecoupage par $\Delta_n$
donne donc la relation $R_n=R_{n-1}+n$.

Or par hypoth\`ese de r\'ecurrence $\mathcal{H}_{n-1}$ : $R_{n-1}
= \frac{(n-1)n}{2}+1$ donc
$$  R_n = R_{n-1}+n =  \frac{(n-1)n}{2}+1+n=\frac{n(n+1)}{2}+1 $$
Et $\mathcal{H}_n$ est vraie.\\
Ainsi $\forall n\in\Nn^* \quad \mathcal{H}_{n-1}\Rightarrow
\mathcal{H}_{n}$.

\item[$\bullet$] Conclusion :  par r\'ecurrence on a montr\'e que $\mathcal{H}_n$ est vraie quelque soit $n \geqslant 1$.

\end{itemize}
\fincorrection
\correction{000157}
\begin{enumerate}
\item Montrons la proposition demand\'ee par r\'ecurrence:
soit $\mathcal{A}_{n}$ l'assertion $f^{n + 1} = f \circ f^n$.
Cette assertion est vraie pour $n=0$. Pour $n\in \Nn$ supposons
$\mathcal{A}_{n}$ vraie. Alors
$$f^{n + 2} = f^{n + 1} \circ f = (f \circ f^n) \circ f = f \circ (f^n \circ f) = f \circ f^{n + 1}.$$
Nous avons utiliser la definition de $f^{n + 2}$, puis la proposition $\mathcal{A}_{n}$,
puis l'associativit\'e de la composition, puis la d\'efinition de $f^{n + 1}$.
Donc $\mathcal{A}_{n+1}$ est vraie. Par le principe de r\'ecurrence
$$\forall \in \Nn \ \ f^n\circ f = f\circ f^n.$$

\item On proc\`ede de m\^eme par r\'ecurrence:
soit $\mathcal{A}_{n}$ l'assertion $ (f^{-1})^n  = (f^n)^{-1}$.
Cette assertion est vraie pour $n=0$. Pour $n\in \Nn$ supposons
$\mathcal{A}_{n}$ vraie. Alors
$$(f^{-1})^{n+1} = (f^{-1})^{n} \circ f^{-1} = (f^n)^{-1} \circ f^{-1} = (f\circ f^n)^{-1} =
 ( f^n \circ f)^{-1} = ( f^{n+1})^{-1} .$$
Donc $\mathcal{A}_{n+1}$ est vraie. Par le principe de r\'ecurrence
$$\forall \in \Nn \ \ (f^{-1})^n
 = (f^n)^{-1}.$$
\end{enumerate}
\fincorrection
\nocorrection
\nocorrection
\nocorrection
\nocorrection
\nocorrection
\nocorrection
\nocorrection
\nocorrection
\nocorrection
\nocorrection
\nocorrection
\nocorrection
\nocorrection
\nocorrection
\nocorrection
\nocorrection
\nocorrection
\nocorrection
\nocorrection
\nocorrection
\nocorrection
\nocorrection
\nocorrection
\nocorrection
\nocorrection
\nocorrection
\nocorrection
\nocorrection
\nocorrection
\nocorrection
\nocorrection
\nocorrection
\nocorrection
\nocorrection
\nocorrection
\nocorrection
\nocorrection
\nocorrection
\nocorrection
\nocorrection
\nocorrection
\nocorrection
\nocorrection
\nocorrection
\nocorrection
\nocorrection
\nocorrection
\nocorrection
\nocorrection
\nocorrection
\nocorrection
\correction{000209}
\begin{enumerate}
    \item Soient $z,z',z''$ des complexes quelconques.
\begin{itemize}
    \item[$\bullet$] Reflexivit\'e : $ z\mathcal{R}z$ car $|z|=|z|$.
    \item[$\bullet$] Sym\'etrie : $z\mathcal{R}z' \Rightarrow z'\mathcal{R}z$ car
$|z|=|z'|$ et donc $|z'|=|z|$.
    \item[$\bullet$] Transitivit\'e : $z\mathcal{R}z'$ et $z'\mathcal{R}z''$ alors
$|z|=|z'|=|z''|$ donc $z\mathcal{R}z''$.
\end{itemize}
En fait, nous avons juste retranscrit que l'\'egalit\'e ``$=$'' est
une relation d'\'equivalence.

    \item La classe d'\'equivalence d'un point $z\in\Cc$
est l'ensemble des complexes qui sont en relation avec $z$,
\emph{i.e.} l'ensemble des complexes dont le module est \'egal \`a
$|z|$. G\'eom\'etriquement la classe d'\'equivalence de $z$ est le
cerlce $\mathcal{C}$ de centre $0$ et de rayon $|z|$ :
$$\mathcal{C} = \left\{ |z|e^{i\theta} \ / \ \theta \in \Rr \right\}.$$
\end{enumerate}
\fincorrection
\correction{000210}
Le raisonnement est faux.

L'erreur est due au manque de quantification. En effet, rien ne
prouve que pout tout $x$ un tel $y$ existe. Il peut exister un
\'el\'ement $x$ qui n'est en relation avec personne (m\^eme pas
avec lui).
\fincorrection
\nocorrection
\correction{000212}
 \begin{enumerate}
 \item 
   \begin{itemize}
   \item Reflexivit\'e : Pour tout $x\in \Rr$, $xe^x=xe^x$ donc $x\mathcal{R}x$.
   \item Sym\'etrie : Pour $x,y \in \Rr$, si $x\mathcal{R} y$ alors $xe^y=ye^x$ donc 
$ye^x=xe^y$ donc $y\mathcal{R} x$.
   \item Transitivit\'e : Soient $x,y,z \in \Rr$ tels que $x\mathcal{R} y$ et $y\mathcal{R} z$,
alors $xe^y=ye^x$ et $ye^z=ze^y$.
Calculons $xye^z$ : 
$$xye^z= x(ye^z) = x(ze^y)= z(xe^y)=z(ye^x)=yze^x.$$
Donc $xye^z=yze^x$. Si $y \neq 0$ alors en divisant par $y$ on vient de montrer que
$xe^z=ze^x$ donc $x\mathcal{R} z$ et c'est fini.
Pour le cas $y=0$ alors $x=0$ et $z=0$ donc $x\mathcal{R} z$ \'egalement.
   \end{itemize}

 \item Soit $x\in \Rr$ fix\'e. On note $\mathcal{C}(x)$ la classe d'\'equivalence de $x$ modulo $\mathcal{R}$ :
$$\mathcal{C}(x) := \left\lbrace y \in \Rr \mid y\mathcal{R} x \right\rbrace.$$
Donc 
$$\mathcal{C}(x) = \left\lbrace y \in \Rr \mid xe^y = ye^x \right\rbrace.$$
Soit la fonction $f : \Rr \to \Rr$ d\'efinie par 
$$f(t) = \frac t {e^t}.$$
Alors 
$$\mathcal{C}(x) = \left\lbrace y \in \Rr \mid f(x)=f(y) \right\rbrace.$$
Autrement dit $\mathcal{C}(x)$ est l'ensemble des $y\in \Rr$ qui par $f$ prennent la m\^eme valeur que $f(x)$ ; en raccourci :
$$\mathcal{C}(x) = f^{-1}\left( f(x) \right).$$

\'Etudions maintenant la fonction $f$ afin de d\'eterminer le nombre d'ant\'ec\'edents:
par un calcul de $f'$ on montrer que $f$ est strictement croissante sur $]-\infty,1]$ puis
strictement d\'ecroissante sur $[1,+\infty[$. De plus en $-\infty$ la limite de $f$ est $-\infty$,
$f(1) = \frac 1e$, et la limite en $+\infty$ est $0$.

C'est le moment de dessiner le graphe de $f$ !!

   \begin{itemize}
      \item Pour $x\leqslant 0$ alors $f(x) \in ]-\infty,0]$ et alors $f(x)$ a un seul ant\'ec\'edent.

      \item Pour $x>0$ avec $x\neq 1$ alors $f(x) \in ]0,\frac 1e[$ et alors $f(x)$ a deux ant\'ec\'edents.
      
      \item pour $x=1$, alors $f(x) = 1/e$ n'a qu'un seul antécédent.
   \end{itemize}
   
   Bilan : si $x \in {}]0,1[ \cup ]1,+\infty[$ alors $\text{Card\,} \mathcal{C}(x) = \text{Card\,} f^{-1}\left( f(x) \right) = 2$,
   si $x\leqslant 0$ ou $x=1$ alors $\text{Card\,} \mathcal{C}(x) = \text{Card\,} f^{-1}\left( f(x) \right) = 1$.

 \end{enumerate}
\fincorrection
\nocorrection
\nocorrection
\nocorrection
\nocorrection
\correction{000217}
 \begin{itemize}
 \item Reflexivit\'e : pour tout $X\in\mathcal{P} (E)$ on a $X \prec X$ car $X=X$.
 \item Anti-sym\'etrie : pour $X,Y\in\mathcal{P} (E)$ tels que $X \prec Y$ et $Y \prec X$, alors par d\'efinition de $\prec$ on a 
$$\forall x \in X \quad \forall y \in Y \quad x\leqslant y \text{ et } y \leqslant x.$$
Comme la relation $\le$ est une relation d'ordre alors $x\leqslant y$ et $y \leqslant x$ implique $x=y$.
Donc 
$$\forall x \in X \quad \forall y \in Y \quad x = y,$$
ce qui implique que $X=Y$ (dans ce cas en fait $X$ est vide ou un singleton).

 \item Transitivit\'e : soit $X,Y, Z \in\mathcal{P} (E)$ tels que  $X \prec Y$ et $Y \prec Z$.
Si $X=Y$ ou $Y=Z$ alors il est clair que $X \prec Z$. Supposons que $X\neq Y$ et $Y\neq Z$
alors 
$$\forall x \in X \quad \forall y \in Y \quad x\leqslant y \qquad \text{ et } \qquad \forall y \in Y \quad \forall z \in Z \quad y\leqslant z.$$

Donc on a 
$$\forall x \in X \quad \forall y \in Y \quad \forall z \in Z \quad x \leqslant y \text{ et }  y \leqslant z,$$ 
alors par transitivit\'e de  la relation $\le$ on obtient :
$$\forall x \in X \quad \forall z \in Z \quad x \leqslant z.$$ 
Donc $X \prec Z$.
 \end{itemize}
\fincorrection
\nocorrection
\nocorrection
\nocorrection
\nocorrection
\nocorrection
\nocorrection
\nocorrection
\nocorrection
\nocorrection
\nocorrection
\nocorrection
\nocorrection
\nocorrection
\correction{003042}
Soit $A$ non vide et minor{\'e}e, et $B = \{$minorants de $A \}$.\par
$B$ n'est pas vide et est major{\'e}e par $A$ donc $\beta = \sup(B)$ existe.\par
Soit $a \in A$ : $\forall\ b \in B,\ b \le a$ donc $\beta \le a$.\par
Par cons{\'e}quent, $\beta$ minore $A$, donc $\beta = \max(B)$.
\fincorrection
\correction{003043}
\begin{enumerate}
  \item 
  \item
  \item 
  \item Si $(a,b)$ majore $A$, alors $(a,b) \gg (\pm\frac 1{\sqrt2}, \frac 1{\sqrt2})$
     donc $(a,b) \gg (0,\sqrt2)$.
     \par
     R{\'e}ciproque : si $x^2+y^2 \le 1$, alors $(x+y)^2 + (x-y)^2 \le 2$, donc
     $y \pm x \le \sqrt2$, et $(x,y) \ll (0,\sqrt2)$.
     \par
     Finalement, $\sup(A) = (0,\sqrt2)$.
\end{enumerate}
\fincorrection
\nocorrection
\nocorrection
\nocorrection
\nocorrection
\nocorrection
\nocorrection
\nocorrection
\nocorrection
\nocorrection
\correction{007191}

Le contraire de $2|n$ n'est pas que $n$ divise strictement $2$.

\fincorrection
\nocorrection
\nocorrection
\nocorrection
\nocorrection
\nocorrection
\correction{007197}

L'application $f$ est constante sur les classes d'équivalence de $\mathcal R$, donc par définition, elle descend au quotient en une application $[f] : \R/2\pi\Z \to \mathbb U$, qui vérifie  $f= [f]\circ p$. Comme $f$ est surjective, $[f]$ aussi. Montrons l'injectivité.

Soient $\alpha$ et $\beta$ dans $\R/2\pi\Z$ tels que $[f](\alpha) = [f](\beta)$. Si $x$ et $y$ sont des représentants de $\alpha$ et $\beta$, on a donc $f(x)=f(y)$, c'est-à-dire $e^{ix}=e^{iy}$, d'où par le cours sur l'exponentielle complexe, $x\equiv y \pmod{2\pi}$, d'où $[x]=[y]$, ou encore $\alpha=\beta$.

\fincorrection
\nocorrection
\nocorrection
\nocorrection
\nocorrection
\nocorrection
\nocorrection
\nocorrection
\nocorrection
\nocorrection
\nocorrection
\nocorrection
\nocorrection
\nocorrection
\nocorrection
\nocorrection
\nocorrection
\correction{000185}
Si $f\circ g=g\circ f$ alors 
$$\forall x \in \Rr \ \ f\circ g (x) = g\circ f(x).$$
Nous allons montrer que c'est faux, en exhibant un contre-exemple.
Prenons $x=0$. Alors $f\circ g (0) = f(-1) = -2$, et
$g\circ f(0) = g(1) = 0$ donc $f\circ g (0) \not= g\circ f(0)$.
Ainsi $f\circ g \not= g\circ f$.
\fincorrection
\nocorrection
\nocorrection
\nocorrection
\nocorrection
\nocorrection
\nocorrection
\nocorrection
\nocorrection
\nocorrection
\nocorrection
\nocorrection
\nocorrection
\nocorrection
\nocorrection
\nocorrection
\correction{000190}
\begin{enumerate}
\item $f$ n'est pas surjective car $0$ n'a pas d'antécédent : en effet il n'existe pas de $n\in\Nn$ tel que $f(n)=0$ (si ce $n$ existait ce serait $n=-1$
qui n'est pas un élément de $\Nn$). Par contre $f$ est injective : soient $n,n' \in \Nn$ tels que $f(n)=f(n')$ alors
$n+1=n'+1$ donc $n=n'$. Bilan $f$ est injective, non surjective et donc non bijective.

\item Pour montrer que $g$ est bijective deux méthodes sont possibles. Première méthode : montrer que $g$ est à la fois injective et surjective.
En effet soient $n,n'\in \Zz$ tels que $g(n)=g(n')$ alors $n+1=n'+1$ donc $n=n'$, alors $g$ est injective. Et $g$ est surjective car chaque $m\in \Zz$
admet un antécédent par $g$ : en posant $n=m-1 \in \Zz$ on trouve bien $g(n)=m$.
Deuxième méthode : expliciter directement la bijection réciproque. Soit la fonction $g' : \Zz \to \Zz$ définie par $g'(m)=m-1$
alors $g' \circ g(n) = n$ (pour tout $n\in \Zz$) et $g \circ g'(m) = m$ (pour tout $m\in \Zz$). Alors $g'$ est la bijection réciproque de $g$
et donc $g$ est bijective.

\item Montrons que $h$ est injective. Soient $(x,y), (x',y') \in \Rr^2$ tels que $h(x,y)=h(x',y')$.
Alors $(x + y, x-y)=(x' + y', x'-y')$ donc 
$$\begin{cases}
    x+y &= x'+y'\\
    x-y &= x'-y'\\
  \end{cases}$$
En faisant la somme des lignes de ce système on trouve $2x=2x'$ donc $x=x'$ et avec la différence on obtient $y=y'$.
Donc les couples $(x,y)$ et $(x',y')$ sont égaux. Donc $h$ est injective.

Montrons que $h$ est surjective. Soit $(X,Y) \in \Rr^2$, cherchons lui un antécédent $(x,y)$ par $h$.
Un tel antécédent vérifie $h(x,y)=(X,Y)$, donc $(x+y,x-y)=(X,Y)$ ou encore :
$$\begin{cases}
    x+y &= X\\
    x-y &= Y\\
  \end{cases}$$
Encore une fois on faisant la somme des lignes on obtient $x=\frac{X+Y}{2}$ et avec la différence $y = \frac{X-Y}{2}$,
donc $(x,y) = (\frac{X+Y}{2},\frac{X-Y}{2})$. La partie ``analyse'' de notre raisonnement en finie passons à la ``synthèse'' :
il suffit de juste de vérifier que le couple $(x,y)$ que l'on a obtenu est bien solution (on a tout fait pour !).
Bilan pour $(X,Y)$ donné, son antécédent par $h$ existe et est $(\frac{X+Y}{2},\frac{X-Y}{2})$. Donc $h$ est surjective.

En fait on pourrait montrer directement que $h$ est bijective en exhibant sa bijection réciproque $(X,Y) \mapsto (\frac{X+Y}{2},\frac{X-Y}{2})$.
Mais vous devriez vous convaincre qu'il s'agit là d'une différence de rédaction, mais pas vraiment d'un raisonnement différent.


        
\item Montrons d'abord que $k$ est injective : soient $x,x' \in \Rr\setminus \{1\}$ tels que $k(x)=k(x')$ alors
$\frac{x + 1}{x - 1}=\frac{x' + 1}{x' - 1}$ donc $(x+1)(x'-1)=(x-1)(x'+1)$. En développant nous obtenons
$xx'+x'-x=xx'-x'+x$, soit $2x=2x'$ donc $x=x'$.

Au brouillon essayons de montrer que $k$ est surjective : soit $y\in \Rr$ et cherchons $x\in \Rr\setminus \{1\}$
tel que $f(x)=y$. Si un tel $x$ existe alors il vérifie $\frac{x + 1}{x - 1}=y$ donc
$x+1=y(x-1)$, autrement dit $x(y-1)=y+1$. Si l'on veut exprimer $x$ en fonction de $y$ cela se fait par la formule
$x = \frac{y+1}{y-1}$. Mais attention, il y a un piège ! Pour $y=1$ on ne peut pas trouver d'antécédent $x$ 
(cela revient à diviser par $0$ dans la fraction précédente).
Donc $k$ n'est pas surjective car $y=1$ n'a pas d'antécédent.

Par contre on vient de montrer que s'il l'on considérait la restriction $k_| :  {\Rr \setminus \left\{ 1\right\}} \to {\Rr \setminus \left\{ 1\right\}}$
qui est définie aussi par $k_|(x) = {\frac{x + 1}{x - 1}}$ (seul l'espace d'arrivée change par rapport à $k$) alors
cette fonction $k_|$ est injective et surjective, donc bijective (en fait sa bijection réciproque est elle même).

\end{enumerate}
\fincorrection
\correction{000191}
\begin{enumerate}
  \item $f$ n'est pas injective car $f(2) = \frac 45 = f(\frac12)$.
  $f$ n'est pas surjective car $y=2$ n'a pas d'ant\'ec\'edent: en effet 
  l'\'equation $f(x)=2$ devient $2x=2(1+x^2)$ soit $x^2-x+1=0$ qui n'a pas de solutions r\'eelles.
  \item $f(x)=y$ est \'equivalent \`a l'\'equation 
  $yx^2-2x+y=0$. Cette \'equation a des solutions $x$ si et seulement si
  $\Delta = 4-4y^2 \geq 0$ donc il y a des solutions si et seulement si $y\in[-1,1]$. Nous venons de montrer que $f(\Rr)$ est exactement $[-1,1]$.
  \item Soit $y\in[-1,1]\setminus\{0\}$ alors les solutions $x$ possibles de l'\'equation $g(x)=y$ sont $x=\frac{1-\sqrt{1-y^2}}{y}$ ou
  $x=\frac{1+\sqrt{1-y^2}}{y}$. La seule solution $x\in[-1,1]$ est
  $x=\frac{1-\sqrt{1-y^2}}{y}$ en effet
  $x=\frac{1-\sqrt{1-y^2}}{y}=\frac{y}{1+\sqrt{1-y^2}} \in[-1,1]$.
  Pour $y=0$, la seule solution de l'\'equation $g(x)=0$ est $x=0$.
  Donc pour $g : [-1,1] \longrightarrow [-1,1]$ nous avons trouv\'e un inverse $h : [-1,1] \longrightarrow [-1,1]$ d\'efini par
  $h(y) = \frac{1-\sqrt{1-y^2}}{y}$ si $y\neq0$ et $h(0)=0$. Donc $g$ est une bijection.
  \item $f'(x) = \frac{2-2x^2}{1+x^2}$, donc $f'$ est strictement positive sur $]-1,1[$ donc $f$ est strictement croissante sur $[-1,1]$ avec $f(-1)=-1$ et $f(1)=1$. Donc la restriction
  de $f$, appelée $g : [-1,1] \longrightarrow [-1,1]$, est une bijection.
\end{enumerate}
\fincorrection
\nocorrection
\correction{000193}
\begin{enumerate}
\item Supposons $g\circ f$ injective, et montrons que $f$ est injective :
soient $a,\ a' \in A$ avec $f(a)=f(a')$ donc $g\circ f(a)=g\circ
f(a')$ or $g\circ f$ est injective donc $a= a'$. Conclusion on a
montr\'e : $$\forall a,a'  \in A \quad  f(a)=f(a') \Rightarrow
a=a'$$ c'est la d\'efinition de $f$ injective.

\item Supposons $g\circ f$ surjective, et montrons que $g$ est surjective :
soit $c \in C$ comme $g\circ f$ est surjective il existe $a \in A$
tel que  $g\circ f(a)=c$ ; posons $b = f(a)$, alors $g(b)=c$, ce
raisonnement est valide quelque soit $c \in C$ donc $g$ est
surjective.

\item  Un sens est simple $(\Leftarrow)$ si $f$ et $g$ sont bijectives alors $g\circ f$ l'est \'egalement. De m\^eme avec $h\circ g$.

\par

Pour l'implication directe $(\Rightarrow)$ : si $g\circ f$ est
bijective alors en particulier elle est surjective et donc
d'apr\`es la question 2. $g$ est surjective.

 Si $h\circ g$ est bijective, elle est en particulier  injective, donc $g$ est injective (c'est le 1.). Par cons\'equent $g$ est \`a la fois injective et
surjective donc bijective.

 Pour finir $f=g^{-1} \circ (g\circ f)$ est bijective comme compos\'ee d'applications bijectives, de m\^eme pour $h$.

\end{enumerate}
\fincorrection
\nocorrection
\nocorrection
\nocorrection
\correction{000197}
\begin{enumerate}
     \item Pour $z=x+iy$, le  module de $e^z=e^{x+iy}=e^xe^{iy}$ est $e^x$ et
son argument est $y$.
    \item Les r\'esultats : $e^{z+z'} = e^ze^{z'}$, $e^{\overline{z}} =
\overline{e^z}$, $e^{-z}= \left( e^{z} \right)^{-1}$,
$(e^z)^n=e^{nz}$.
    \item La fonction $\exp$ n'est pas surjective car $|e^z| = e^x >0$ et
donc $e^z$ ne vaut jamais $0$. La fonction $\exp$ n'est pas non
plus injective car pour $z\in\Cc$, $e^z=e^{z+2i\pi}$.
\end{enumerate}
\fincorrection
\correction{005110}
\begin{enumerate}
 \item  Soit $(x_1,x_2)\in E^2$.

\begin{align*}
f(x_1)=f(x_2)&\Rightarrow g(f(x_1))=g(f(x_2))\;(\mbox{car}\;g\;\mbox{est une application})\\
 &\Rightarrow x_1=x_2\;(\mbox{car}\;g\circ f\;\mbox{est injective}).
\end{align*}
On a montré que $\forall(x_1,x_2)\in E^2,\;f(x_1)=f(x_2)\Rightarrow x_1=x_2$, et donc $f$ est injective.

 \item  Soit $y\in H$. Puisque $g\circ f$ est surjective, il existe un élément $x$ dans $E$ tel que $g(f(x))=y$. En
posant $z=f(x)\in G$, on a trouvé $z$ dans $G$ tel que $g(z)=y$. On a montré~:~$\forall y\in H,\;\exists z\in G/\;g(z)=
y$, et donc $g$ est surjective.

\end{enumerate}
\fincorrection
\correction{005114}
\textbf{1) $\Rightarrow$ 2)} Soit $X\in\mathcal{P}(E)$. On a toujours $X\subset f^{-1}(f(X))$. (En effet, pou $x\in E$, $x\in
X\Rightarrow f(x)\in f(X)\Rightarrow x\in f^{-1}(f(X))$).
Réciproquement, soit $x\in E$.

\begin{align*}
x\in f^{-1}(f(X))&\Rightarrow f(x)\in f(X)\Rightarrow\exists x'\in X/\;f(x)=f(x')\Rightarrow\exists x'\in
X/\;x=x'\;(\mbox{puisque}\;f\;\mbox{est injective})\\
 &\Rightarrow x\in X.
\end{align*}
Finalement, $f^{-1}(f(X))\subset X$ et donc $f^{-1}(f(X))=X$.
\textbf{2) $\Rightarrow$ 1)} Soit $x\in X$. Par hypothése, $f^{-1}\{f(x)\}=f^{-1}(f(\{x\}))=\{x\}$ ce qui signifie que
$f(x)$ a un et un seul antécédent à savoir $x$. Par suite, tout élément de l'ensemble d'arrivée a au plus un antécédent
par $f$ et $f$ est injective.

\textbf{1) $\Rightarrow$ 3)} Soit $(X,Y)\in(\mathcal{P}(E))^2$. On a toujours $f(X\cap Y)\subset f(X)\cap
f(Y)$ ($X\cap Y\subset X\Rightarrow f(X\cap Y)\subset f(X)$ et de même, $f(X\cap Y)\subset 
f(Y)$ et finalement, $f(X\cap Y)\subset f(X)\cap f(Y)$).
Réciproquement, soit $y\in F$. $y\in f(X)\cap f(Y)\Rightarrow\exists(x,x')\in X\times Y/\;y=f(x)=f(x')$. Mais alors,
puisque $f$ est injective, $x=x'\in X\cap Y$ puis $y=f(x)\in f(X\cap Y)$. Finalement, $f(X\cap Y)=f(X)\cap f(Y)$.

\textbf{3) $\Rightarrow$ 4)} Soit $(X,Y)\in(\mathcal{P}(E))^2$. $X\cap Y=\varnothing\Rightarrow f(X)\cap f(Y)=f(X\cap
Y)=f(\varnothing)=\varnothing$.

\textbf{4) $\Rightarrow$ 5)} Soit $(X,Y)\in(\mathcal{P}(E))^2$ tel que $Y\subset X$.
Puisque $X\setminus Y\subset X$, on a $f(X\setminus Y)\subset f(X)$. Mais, puisque $Y\cap(X\setminus Y)=\varnothing$, par
hypothèse $f(X\setminus Y)\cap f(Y)=\varnothing$. Finalement,$f(X\setminus Y)\subset f(X)\setminus f(Y)$.
Inversement, si $f(X)\setminus f(Y)=\varnothing$, l'inclusion contraire est immédiate et si $f(X)\setminus
f(Y)\neq\varnothing$, un élémént de $f(X)\setminus f(Y)$ est l'image d'un certain élément de $X$ qui ne peut être dans $Y$
et donc est l'image d'un élément de $X\setminus Y$ ce qui montre que $f(X)\setminus f(Y)\subset f(X\setminus Y)$ et
finalement que $f(X)\setminus f(Y)=f(X\setminus Y)$.

\textbf{5) $\Rightarrow$ 1)} Soit $(x_1,x_2)\in E^2$ tel que $x_1\neq x_2$. Posons $X=\{x_1,x_2\}$ et $Y=\{x_2\}$.
On a
donc $Y\subset X$. Par hypothèse $f(X\setminus Y)=f(X)\setminus f(Y)$ ce qui fournit
$f(\{x_1\})=f(\{x_1,x_2\})\setminus f(\{x_2\})$ ou encore, $\{f(x_1)\}=\{f(x_1),f(x_2)\}\setminus\{f(x_2)\}$.
Maintenant, si $f(x_1)=f(x_2)$ alors $\{f(x_1),f(x_2)\}\setminus\{f(x_2)\}=\varnothing$ (et pas $\{f(x_1)\}$). Donc
$f(x_1)\neq f(x_2)$.
On a montré que $f$ est injective.
\fincorrection
\correction{000198}
L'inverse de $f_{a,b}$ est $g_{a,b}$ avec $g_{a,b}(y) = \frac1ay-\frac ba$.
Autrement dit $f_{a,b}^{-1} = g_{a,b} = f_{\frac 1a, -\frac ba}$.
\fincorrection
\correction{000199}
Soit $x\in [0,1]\cap\Qq$ alors $f(x) = x$ donc $f \circ f (x) = f(x) = x$.
Soit $x\notin [0,1]\cap\Qq$ alors $f(x) = 1-x$ donc $f \circ f (x) = f(1-x)$,
mais $1-x \notin [0,1]\cap\Qq$ (vérifiez-le !) donc  $f \circ f (x) = f(1-x) = 1 - (1-x) = x$.
Donc pour tout $x \in [0,1]$ on a  $f \circ f (x) = x$. Et donc $f \circ f= id$.
\fincorrection
\correction{000200}
Considérons la restriction suivante de $f$ : $f_|:[0,2\pi[ \longrightarrow \mathbb{U}$, 
$t\mapsto e^{it}$. Montrons que cette nouvelle application $f_|$ est bijective. Ici $\mathbb{U}$
est le cercle unit\'e de $\Cc$ donn\'e par l'\'equation $(|z|=1)$.
\begin{itemize}
    \item[$\bullet$] $f_|$ est surjective car tout nombre complexe de $\mathbb{U}$ s'\'ecrit
sous la forme polaire $e^{i\theta}$, et l'on peut choisir $\theta
\in [0,2\pi[$.

    \item[$\bullet$] $f_|$ est injective :
\begin{align*}
f_|(t) = f_|(t') &\Leftrightarrow e^{it}=e^{it'}\\
&\Leftrightarrow t=t' +2k\pi \text{ avec } k\in \Zz\\
&\Leftrightarrow t=t' \text{ car } t,t'\in[0,2\pi[ \text{ et donc }k=0.\\
\end{align*}
\end{itemize}
En conclusion $f_|$ est injective et surjective donc bijective.
\fincorrection
\nocorrection
\correction{000202}
\begin{itemize}
    \item[$\bullet$] $f$ est injective : soient $x,y \in [1,+\infty[$ tels que $f(x)=f(y)$ :
\begin{align*}
f(x)=f(y) &\Rightarrow x^2-1=y^2-1\\
&\Rightarrow x = \pm y \text{ or $x,y\in [1,+\infty[$ donc $x,y$ sont de m\^eme signe}\\
&\Rightarrow x =y.\\
\end{align*}

    \item[$\bullet$] $f$ est surjective : soit $y\in [0,+\infty[$.
Nous cherchons un \'el\'ement $x\in [1,+\infty[$ tel que $y = f(x)
= x^2-1$ . Le r\'eel $x= \sqrt{y+1}$ convient !
\end{itemize}
\fincorrection
\nocorrection
\nocorrection
\nocorrection
\nocorrection
\correction{005106}
\begin{enumerate}
 \item  $f$ est dérivable sur $I=]-\infty,2]$, et pour $x\in]-\infty,2[$, $f'(x)=2x-4<0$. $f$ est donc continue
et strictement décroissante sur $]-\infty,2]$.
Par suite, $f$ réalise une bijection de $]-\infty,2]$ sur
$f(]-\infty,2])=[f(2),\underset{-\infty}{\mbox{lim}}\;f[=[-1,+\infty[=J$. On note $g$ l'application de $I$ dans $J$ qui, à
$x$ associe $x^2-4x+3(=f(x))$. $g$ est bijective et admet donc une réciproque. Déterminons $g^{-1}$.
Soit 
$y\in[-1,+\infty[$ et $x\in]-\infty,2]$.

$$y=g(x)\Leftrightarrow y=x^2-4x+3\Leftrightarrow x^2-4x+3-y=0.$$
Or, $\Delta'=4-(3-y)=y+1\geq0$. Donc, $x=2+\sqrt{y+1}$ ou $x=2-\sqrt{y+1}$. Enfin, $x\in]-\infty,2]$ et donc,
$x=2-\sqrt{y+1}$. En résumé,

$$\forall x\in]-\infty,2],\;\forall y\in[-1,+\infty[,\;y=g(x)\Leftrightarrow x=2-\sqrt{y+1}.$$
On vient de trouver $g^{-1}$~:

\begin{center}
\shadowbox{$\forall x\in[-1,+\infty[,\;g^{-1}(x)=2-\sqrt{x+1}$.}
\end{center}
 \item  On vérifie facilement que $f$ réalise une bijection de $]-2,+\infty[$ sur $]-\infty,2[$, notée $g$.
Soient
alors $x\in]-2,+\infty[$ et $y\in]-\infty,2[$.

$$y=g(x)\Leftrightarrow y=\frac{2x-1}{x+2}\Leftrightarrow x(-y+2)=2y+1\Leftrightarrow x=\frac{2y+1}{-y+2}.$$
(on a ainsi trouvé au plus une valeur pour $x$ à savoir $x=\frac{2y+1}{-y+2}$, mais il n'est pas nécessaire 
de vérifier que cette expression est bien définie et élément de $]-2,+\infty[$ car on sait à l'avance que $y$ admet au
moins un antécédent dans $]-2,+\infty[$, et c'est donc nécessairement le bon). En résumé,

$$\forall x\in]-2,+\infty[,\;\forall y\in]-\infty,2[,\;y=g(x)\Leftrightarrow x=\frac{2y+1}{-y+2}.$$
On vient de trouver $g^{-1}$~:~

\begin{center}
\shadowbox{$\forall x\in]-\infty,2[,\;g^{-1}(x)=\frac{2x+1}{-x+2}$}.
\end{center}
 \item  $f$ est continue et strictement croissante sur $\left[-\frac{3}{2},+\infty\right[$.
$f$ est donc bijective de
$\left[-\frac{3}{2},+\infty\right[$ sur
$f\left(\left[-\frac{3}{2},+\infty\right[\right)=\left[f\left(-\frac{3}{2}\right),\underset{+\infty}{\mbox{lim}}f\right[=[-1,+\infty[$. Notons encore $f$
l'application de $\left[-\frac{3}{2},+\infty\right[$ dans $[-1,+\infty[$ qui à $x$ associe $\sqrt{2x+3}-1$. Soient alors
$x\in[-\frac{3}{2},+\infty[$ et $y\in[-1,+\infty[$.

$$f(x)=y\Leftrightarrow\sqrt{2x+3}-1=y\Leftrightarrow x=\frac{1}{2}(-3+(y+1)^2)\Leftrightarrow x=\frac{y^2}{2}+y-1.$$
En résumé, $\forall x\in\left[-\frac{3}{2},+\infty\right[,\;\forall y\in[-1,+\infty[,\;y=g(x)\Leftrightarrow x=\frac{y^2}{2}+y-1$. On vient
de trouver $g^{-1}$~:~

\begin{center}
\shadowbox{
$\forall x\in[-1,+\infty[,\;g^{-1}(x)=\frac{x^2}{2}+x-1$.
}
\end{center}
 \item  $f$ est définie sur $\Rr$, impaire.
Pour $x\in[0,+\infty[$, $0\leq f(x)=\frac{x}{1+x}<\frac{1+x}{1+x}=1$. Donc, $f([0,+\infty[)\subset[0,1[$.
Par parité, $f(]-\infty,0])\subset]-1,0]$ et même $f(]-\infty,0[)\subset]-1,0[$ car l'image par $f$ d'un réel
strictement négatif est un réel strictement négatif.
Finalement, $f(\Rr)\subset]-1,1[$.
Vérifions alors que $f$ réalise une bijection de $\Rr$ sur $]-1,1[$.
Soit $y\in[0,1[$ et $x\in\Rr$. L'égalité $f(x)=y$ impose à $x$ d'être dans $[0,+\infty[$. Mais alors

$$f(x)=y\Leftrightarrow\frac{x}{1+x}=y\Leftrightarrow x=\frac{y}{1-y}.$$
Le réel $x$ obtenu est bien défini, car $y\neq1$, et positif, car $y\in[0,1[$. On a montré que~:

$$\forall y\in[0,1[,\;\exists!x\in\Rr/\;y=f(x)\;(\mbox{à savoir}\;x=\frac{y}{1-y}).$$
Soit $y\in]-1,0[$ et $x\in\Rr$. L'égalité $f(x)=y$ impose à $x$ d'être dans $]-\infty,0[$. Mais alors

$$f(x)=y\Leftrightarrow\frac{x}{1-x}=y\Leftrightarrow x=\frac{y}{1+y}.$$
Le réel $x$ obtenu est bien défini, car $y\neq-1$, et strictement négatif, car $y\in]-1,0[$. On a montré que~:

$$\forall y\in]-1,0[,\;\exists!x\in\Rr/\;y=f(x)\;(\mbox{à savoir}\;x=\frac{y}{1+y}).$$
Finalement,

$$\forall y\in]-1,1[,\;\exists!x\in\Rr/\;y=f(x),$$
ce qui montre que $f$ réalise une bijection de $\Rr$ sur $]-1,1[$. De plus, pour $y\in]-1,1[$ donné,
$f^{-1}(y)=\frac{y}{1-y}$ si $y\geq0$ et $f^{-1}(y)=\frac{y}{1+y}$ si $y<0$. Dans tous les cas, on a
$f^{-1}(y)=\frac{y}{1-|y|}$.

En notant encore $f$ l'application de $\Rr$ dans $]-1,1[$ qui à $x$ associe $\frac{x}{1+|x|}$, on a donc
\begin{center}
\shadowbox{
$\forall x\in]-1,1[,\;f^{-1}(x)=\frac{x}{1-|x|}.$
}
\end{center}
\end{enumerate}
\fincorrection
\correction{005107}
\begin{enumerate}
 \item  Montrons que la restriction de $f$ à $D$, notée $g$, est bien une application de $D$ dans $P$.
Soit $z\in D$. On a $|z|<1$ et en particulier $z\neq i$. Donc, $f(z)$ existe. De plus,

$$\Re(f(z))=\frac{1}{2}(f(z)+\overline{f(z)})=\frac{1}{2}\left(\frac{z+i}{z-i}+\frac{{\bar z}-i}{{\bar z}+i}\right)=
\frac{1}{2}\frac{2z{\bar z}-2}{(z-i)(\overline{z-i})}=\frac{|z|^2-1}{|z-i|^2}<0.$$
Donc, $f(z)$ est élément de $P$. $g$ est donc une application de $D$ dans $P$.
 \item  Montrons que $g$ est injective.
Soit $(z,z')\in D^2$.

$$g(z)=g(z')\Rightarrow\frac{z+i}{z-i}=\frac{z'+i}{z'-i}\Rightarrow iz'-iz=iz-iz'\Rightarrow 2i(z'-z)=0\Rightarrow
z=z'.$$
 \item  Montrons que $g$ est surjective.
Soient $z\in D$ et $Z\in P$.

$$g(z)=Z\Leftrightarrow\frac{z+i}{z-i}=Z\Leftrightarrow z=\frac{i(Z+1)}{Z-1}\;(\text{car}\;Z\neq1,$$
(ce qui montre que $Z$ admet au plus un antécédent dans $D$, à savoir $z=\frac{i(Z+1)}{Z-1}$ (mais on le sait déjà car
$g$ est injective). Il reste cependant à vérifier que $\frac{i(Z+1)}{Z-1}$ est effectivement dans $D$).
Réciproquement, puisque
$\Re(Z)<0$, 

\begin{center}
$\left|\frac{i(Z+1)}{Z-1}\right|=\frac{|Z+1|}{|Z-1|}<1$
\end{center}
($Z$ étant strictement plus proche de $-1$ que
de $1$) et $z\in D$. Finalement $g$ est une bijection de $D$ sur $P$, et~:~

\begin{center}
\shadowbox{
$\forall z\in
P,\;g^{-1}(z)=\frac{i(z+1)}{z-1}$.
}
\end{center}
\end{enumerate}
\fincorrection
\correction{005111}
On peut supposer sans perte de généralité que $f\circ g\circ h$ et $g\circ h\circ f$ sont injectives et que $h\circ
f\circ g$ est surjective. D'après l'exercice \ref{exo:suprou8}, puisque $f\circ g\circ h=(f\circ g)\circ h$ est injective, $h$ est injective
et puisque $h\circ f\circ g=h\circ(f\circ g)$ est surjective, $h$ est surjective. Déjà $h$ est bijective. Mais alors,
$h^{-1}$ est surjective et donc $f\circ g=h^{-1}\circ(h\circ f\circ g)$ est surjective en tant que composée de
surjections. Puis $h^{-1}$ est injective et donc $f\circ g=(f\circ g\circ h)\circ h^{-1}$ est injective. $f\circ g$ est
donc bijective. $f\circ g$ est surjective donc $f$ est surjective. $g\circ h\circ f$ est injective donc $f$ est
injective. Donc $f$ est bijective. Enfin $g=f^{-1}\circ(f\circ g)$ est bijective en tant que composée de bijections.
\fincorrection
\correction{005118}
$f$ est bien une application de $\Nn^2$ dans $\Nn$ car, pour tout couple $(x,y)$ d'entiers
naturels, l'un des deux entiers $x+y$ ou $x+y+1$ est pair et donc, $\frac{(x+y)(x+y+1)}{2}$ est bien un entier naturel
(on peut aussi constater que $\frac{(x+y)(x+y+1)}{2}=1+2+...+(x+y)$ est entier pour $x+y\geq1$).

\textbf{Remarque.} La numérotation de $\Nn^2$ a été effectuée de la façon suivante~:


$$\includegraphics{../images/img005118-1}$$



Sur une parallèle à la droite d'équation $y = -x$, la somme $x+y$ est constante. Il en est de même de l'expression
$\frac{(x+y)(x+y+1)}{2}$ et quand on descend de $1$ en $y$, on avance de $1$ dans la numérotation.

\textbf{Lemme}. $\forall n\in\Nn,\;\exists!p\in\Nn/\;\frac{p(p+1)}{2}\leq n<\frac{(p+1)(p+2)}{2}$.

\textbf{Démonstration}. Pour démontrer ce lemme, on pourrait se contenter de constater que la suite des
nombres triangulaires $\left(\frac{p(p+1)}{2}\right)_{p\geq0}$ est strictement croissante. Néanmoins, on va fournir explicitement
$p$ en fonction de $n$.
Soient $n$ et $p$ deux entiers naturels.
\begin{align*}
\frac{p(p+1)}{2}\leq n<\frac{(p+1)(p+2)}{2}&\Leftrightarrow p^2+p-2n\leq0\;\mbox{et}\;p^2+3p+2-2n>0\\
 &\Leftrightarrow p\leq\frac{-1+\sqrt{8n+1}}{2}\;\mbox{et}\;p>\frac{-3+\sqrt{8n+1}}{2}=-1+\frac{-1+\sqrt{8n+1}}{2}\\
 &\Leftrightarrow p\leq\frac{-1+\sqrt{8n+1}}{2}<p+1\Leftrightarrow p= E\left(\frac{-1+\sqrt{8n+1}}{2}\right).
\end{align*}
Le lemme est démontré.

Montrons que $f$ est surjective (et au passage, déterminons l'antécédent d'un entier $n$ donné).
Soient $n$ un entier naturel et $p=E\left(\frac{-1+\sqrt{8n+1}}{2}\right)$ ($p$ est un entier naturel). On pose $\left\{
\begin{array}{l}
x+y=p\\
y=n-\frac{p(p+1)}{2}
\end{array}
\right.$ ou encore $\left\{
\begin{array}{l}
y=n-\frac{p(p+1)}{2}\\
x=p-y=\frac{p(p+3)}{2}-n
\end{array}
\right.$. Tout d'abord, $y+\frac{(x+y)(x+y+1)}{2}=n-\frac{p(p+1)}{2}+\frac{p(p+1)}{2}=n$. Mais il reste encore à
vérifier que $x$ et $y$ ainsi définis (qui sont à l'évidence des entiers relatifs) sont bien des entiers
naturels. Puisque $\frac{p(p+1)}{2}$ est un entier naturel et que $n\geq\frac{p(p+1)}{2}$, $y$ est bien un entier
naturel. Ensuite, $\frac{p(p+3)}{2}=\frac{p(p+1)}{2}+p$ est aussi un entier naturel et de plus,

$$\frac{p(p+3)}{2}-n\geq\frac{p(p+3)}{2}-\left(\frac{(p+1)(p+2)}{2}-1\right)=0,$$
et $x$ est bien un entier naturel. Ainsi, pour $n$ naturel donné, en posant $p= E\left(\frac{-1+\sqrt{8n+1}}{2}\right)$ puis
$x=\frac{p(p+3)}{2}-n$ et $y=n-\frac{p(p+1)}{2}$, $x$ et $y$ sont des entiers naturels tels que $f((x,y))=n$. $f$ est
donc surjective.
Montrons que $f$ est injective.
Pour cela, on montre que si $x$ et $y$ sont des entiers naturels vérifiant
$y+\frac{(x+y)(x+y+1)}{2}=n$, alors nécessairement, $x+y=p$ (et $y=n-\frac{p(p+1)}{2})$.
Soient donc $x$ et $y$ deux entiers naturels. On a~:

$$\frac{(x+y)(x+y+1)}{2}\leq\frac{(x+y)(x+y+1)}{2}+y=n<\frac{(x+y)(x+y+1)}{2}+(x+y+1)=\frac{(x+y+1)(x+y+2)}{2},$$
et le lemme montre que $x+y=p$. L'unicité du couple $(x,y)$ est donc démontrée. $f$ est une application injective et
surjective et donc $f$ est bijective. Sa réciproque est $\begin{array}[t]{cccc}
f^{-1}~:&\Nn&\rightarrow&\Nn^2\\
 &n&\mapsto&(\frac{p(p+3)}{2}-n,n-\frac{p(p+1)}{2})
\end{array}$ où $p=E\left(\frac{-1+\sqrt{8n+1}}{2}\right)$.
\fincorrection
\nocorrection
\correction{000220}
Soit $f : \Rr \longrightarrow \Rr$ la fonction $f(x) = (1+x)^n$.
Par la formule du bin\^ome de Newton nous savons que
$$f(x) = (1+x)^n = \sum_{k=0}^{n} C_n^k x^k.$$
\begin{enumerate}
  \item En calculant $f(1)$ nous avons $2^n = \sum_{k=0}^{n} C_n^k$.
  \item En calculant $f(-1)$ nous avons $0 = \sum_{k=0}^{n} (-1)^k C_n^k$.
  \item Maintenant calculons  $f'(x) = n(1+x)^{n-1} = \sum_{k=1}^{n} kC_n^k x^{k-1}$. \'Evaluons $f'(1) = n2^{n-1} = \sum_{k=1}^{n} kC_n^k$.
  \item Il s'agit ici de calculer la primitive $F$ de $f$ qui correspond \`a la somme :
$F(x) = \frac{1}{n+1}(1+x)^{n+1}-\frac{1}{n+1} = \sum_{k=0}^{n} \frac{1}{k+1}C_n^k x^{k+1}$. En $F(1) = \frac{1}{n+1}(2^{n+1}-1) = \sum_{k=0}^{n} \frac{1}{k+1}C_n^k $.
\end{enumerate}
\fincorrection
\nocorrection
\correction{000222}
L'astuce consiste \`a \'ecrire $2=3-1$ (!)
$$2^{n} = (3-1)^{n} = 3\times p+ (-1)^{n}$$
O\`u $3\times p$ ($p\in\Zz$) repr\'esente les $n$ premiers termes
de $\sum_{k=0}^{n}C^k_{n}3^k(-1)^{n-k}$ et $(-1)^{n}$ est le
dernier terme. 
Donc $2^n -(-1)^n=3p$. 
Si $n$ est impair l'\'egalit\'e s'\'ecrit
$2^n+1=3p$ et donc $2^n+1$ est divisible par 3.
Si $n$ est pair $2^n-1=3p$ donc $2^n+1=3p+2$ qui n'est pas
divisible par $3$.


Pour l'autre assertion regarder $3=7-4$.
\fincorrection
\nocorrection
\nocorrection
\nocorrection
\nocorrection
\nocorrection
\correction{000228}
Il s'agit de comparer les deux \'ecritures de la fonction
$$f(x) = (1+x)^n = \sum_{k=0}^{n}C_n^k x^k.$$

Pour $x=1$ et $x=-1$ nous obtenons respectivement les assertions
$(a)$ et $(b)$. En d\'erivant la fonction $f$ et en calculant
$f'(1)$, nous obtenons $(b)$. Pour $(d)$ il faut d\'eriver une
nouvelle fois.
\fincorrection
\correction{000229}
 $A=(1+i)^n$ a pour module $2^{n/2}$ et pour argument $n{\pi \over 4}$
(et $B$ est son conjugu\'e).
On en tire gr\^ace \`a la formule du bin\^ome, et en s\'eparant partie r\'eelle 
et partie imaginaire :
$S_1 = 2^{n/2}\cos n{\pi \over4}$ et 
et $S_2 = 2^{n/2}\sin n{\pi \over4}$. 
On a aussi 
$S_1={A+B\over 2}$ et $S_2={B-A\over 2}i$.
\fincorrection
\correction{000230}
L'application $\Phi$ est une bijection : son inverse est $\Phi$
elle-m\^eme.

Supposons que $E$ soit un ensemble fini. Notre bijection $\Phi$
envoie un ensemble $\mathcal{Q} \subset \mathcal{P}(E)$ sur un
ensemble de m\^eme cardinal.

Choisissons $E$ un ensemble \`a $n$ \'el\'ements, et soit $p\le
n$. Soit $\mathcal{Q} \subset \mathcal{P}(E)$ :
$$\mathcal{Q} = \left\lbrace F \subset E,\ \ \mathrm{Card} F = p \right\rbrace.$$

Nous savons que $\mathrm{Card} \mathcal{Q} =C_n^p$ (c'est la d\'efinition
de $C^p_n$).
 De plus
\begin{align*}
 \Phi(\mathcal{Q})
  &= \left\lbrace \Phi(F),\ \ F \subset E,\ \ \mathrm{Card} F = p \right\rbrace \\
  &= \left\lbrace \complement F,\ \ F \subset E,\ \ \mathrm{Card} F = p \right\rbrace \\
  &= \left\lbrace G\subset E,\ \ \mathrm{Card} G = n-p \right\rbrace.
\end{align*}

Donc $\mathrm{Card} \Phi(\mathcal{Q}) = C_n^{n-p}$. Et comme $\Phi$ est
une bijection, $\mathrm{Card} \Phi(\mathcal{Q}) = \mathrm{Card} (\mathcal{Q})$,
donc $C_n^{n-p}=C_n^p$.
\fincorrection
\nocorrection
\nocorrection
\nocorrection
\nocorrection
\correction{002900}
 $\sum_{k=0}^n kC_n^k = n2^{n-1}$, \qquad $\sum_{k=0}^n
\frac{C_n^k}{k+1} = \frac {2^{n+1}-1}{n+1}$.  
\fincorrection
\correction{002901}
\begin{enumerate}
  \item 
  \item $0$ si $p < n$, $(-1)^n$ si $p = n$.
  \item 
\end{enumerate}
\fincorrection
\correction{002902}
$(-1)^pC_{n-1}^p$.
\fincorrection
\correction{002903}
$n2^{n-1}$, $n4^{n-1}$, $3n4^{n-1}$.
\fincorrection
\correction{002904}
$\frac {n(n^2-1)}6$, $\frac {n(n^2-1)(3n^2-12)}{360}$.
\fincorrection
\correction{002905}
\begin{enumerate}
 \item $\Gamma_n^0 = 1$, $\Gamma_n^1 = n$, $\Gamma_n^2 = \frac {n(n+1)}2$,
    $\Gamma_2^n = n+1$.
 \item 
  \item 
\end{enumerate}
\fincorrection
\nocorrection
\correction{002907}
$p = \left[\frac {n+1}2 \right]$.
\fincorrection
\nocorrection
\nocorrection
\nocorrection
\nocorrection
\correction{005137}
\begin{enumerate}
\item  D'après la formule du binôme de \textsc{Newton},
\begin{center}
\shadowbox{
$\forall n\in\Nn,\;\sum_{k=0}^{n}\binom{n}{k}=(1+1)^n=2^n.$
}
\end{center}

\item  Soit $n$ un entier naturel non nul. Posons $S_1=\sum_{k=0}^{E(n/2)}\binom{n}{2k}$ et
$S_2=\sum_{k=0}^{E((n-1)/2)}\binom{n}{2k+1}$. Alors

$$S_1-S_2=\sum_{k=0}^{n}(-1)^k\binom{n}{k}=(1-1)^n=0\;(\mbox{car}\;n\geq1),$$
et donc $S_1=S_2$. Puis $S_1+S_2=\sum_{k=0}^{n}\binom{n}{k}=2^n$, et donc $S_1=S_2=2^{n-1}$.
\begin{center}
\shadowbox{
$\forall n\in\Nn^*,\;\binom{n}{0}+\binom{n}{2}+\binom{n}{4}+\ldots=\binom{n}{1}+\binom{n}{3}+\binom{n}{5}+\ldots=2^{n-1}.$
}
\end{center}

\item  En posant $j=e^{2i\pi/3}$, on a~:

$$\sum_{k=0}^{n}\binom{n}{k}=(1+1)^n=2^n,\;\sum_{k=0}^{n}\binom{n}{k}j^k=(1+j)^n\;\mbox{et}\;\sum_{k=0}^{n}\binom{n}{k}j^{2k}
=(1+j^2)^n.$$

En additionnant ces trois égalités, on obtient

$$\sum_{k=0}^{n}\binom{n}{k}(1+j^k+j^{2k})=2^n+(1+j)^n+(1+j^2)^n.$$

Maintenant,
\begin{itemize}
\item[-] si $k\in3\Nn$, il existe $p\in\Nn$ tel que $k=3p$ et $1+j^k+j^{2k}=1+(j^3)^p+(j^3)^{2p}=3$ car $j^3=1$.
\item[-] si $k\in3\Nn+1$, il existe $p\in\Nn$ tel que $k=3p+1$ et $1+j^k+j^{2k}=1+j(j^3)^p+j^2(j^3)^{2p}=1+j+j^2=0$
\item[-] si $k\in3\Nn+2$, il existe $p\in\Nn$ tel que $k=3p+2$ et
$1+j^k+j^{2k}=1+j^2(j^3)^p+j^4(j^3)^{2p}=1+j^2+j=0$.
\end{itemize}

Finalement, $\sum_{k=0}^{n}\binom{n}{k}(1+j^k+j^{2k})=3\sum_{k=0}^{E(n/3)}\binom{n}{3k}$. Par suite,

\begin{align*}
\sum_{k=0}^{E(n/3)}\binom{n}{3k}&=\frac{1}{3}(2^n+(1+j)^n+(1+j^2)^n)=\frac{1}{3}(2^n+2\Re((1+j)^n))\\
 &=\frac{1}{3}(2^n+2\Re((-j^2)^n))=\frac{1}{3}(2^n+2\cos\frac{n\pi}{3})
\end{align*}

\item  Pour $1\leq k\leq n$, on a

$$k\binom{n}{k}=k\frac{n!}{k!(n-k)!}=n\frac{(n-1)!}{(k-1)!((n-1)-(k-1))!}=k\binom{n-1}{k-1}.$$

\item  $\binom{2n}{n}$ est le coefficient de $x^n$ dans le développement de $(1+x)^{2n}$. Mais d'autre part ,

$$(1+x)^{2n}=(1+x)^n(1+x)^n=(\sum_{k=0}^{n}\binom{n}{k}x^k)(\sum_{k=0}^{n}\binom{n}{k}x^k).$$

Dans le développement de cette dernière expression, le coefficient de $x^n$ vaut $\sum_{k=0}^{n}\binom{n}{k}\binom{n}{n-k}$ ou
encore $\sum_{k=0}^{n}\binom{n}{k}^2$. Deux polynômes sont égaux si et seulement si ils ont mêmes
coefficients et donc

$$\binom{2n}{n}=\sum_{k=0}^{n}\binom{n}{k}^2.$$

\item 
\begin{itemize}
\item[\textbf{1ère solution.}] Pour $x$ réel, posons $P(x)=\sum_{k=1}^{n}k\binom{n}{k}x^{k-1}$.

Pour $x$ réel, $$P(x)=(\sum_{k=0}^{n}\binom{n}{k}x^{k})'=((1+x)^n)'=n(1+x)^{n-1}.$$

En particulier, pour $x=1$, on obtient~:

$$\sum_{k=1}^{n}k\binom{n}{k}=n(1+1)^{n-1}=n2^{n-1}.$$

\item[\textbf{2ème solution.}] D'après 4),

$$\sum_{k=1}^{n}k\binom{n}{k}=\sum_{k=1}^{n}n\binom{n-1}{k-1}=n\sum_{k=0}^{n-1}\binom{n-1}{k}=n(1+1)^{n-1}=n2^{n-1}.$$
\end{itemize}

\begin{itemize}
\item[\textbf{1ère solution.}] Pour $x$ réel, posons $P(x)=\sum_{k=0}^{n}\binom{n}{k}\frac{x^{k+1}}{k+1}$. On a

$$P'(x)=\sum_{k=0}^{n}\binom{n}{k}x^k=(1+x)^n,$$

et donc, pour $x$ réel,

$$P(x)=P(0)+\int_{0}^{x}P'(t)\;dt=\int_{0}^{1}(1+t)^n\;dt=\frac{1}{n+1}((1+x)^{n+1}-1).$$

En particulier, pour $x=1$, on obtient

$$\sum_{k=0}^{n}\frac{\binom{n}{k}}{k+1}=\frac{2^{n+1}-1}{n+1}.$$

\item[\textbf{2ème solution.}] D'après 4), $(n+1)\binom{n}{k}=(k+1)\binom{n+1}{k+1}$ et donc

$$\sum_{k=0}^{n}\frac{\binom{n}{k}}{k+1}=\sum_{k=0}^{n}\frac{\binom{n+1}{k+1}}{n+1}=\frac{1}{n+1}\sum_{k=1}^{n+1}\binom{n+1}{k}
=\frac{1}{n+1}((1+1)^{n+1}-1)=\frac{2^{n+1}-1}{n+1}.$$
\end{itemize}

\item  Pour $1\leq k\leq n-p$, $\binom{p+k}{p}=\binom{p+k+1}{p+1}-\binom{p+k}{p+1}$ (ce qui reste vrai pour $k=p$ en tenant
compte de $\binom{p}{p+1}=0$). Par suite,

\begin{align*}
\sum_{k=0}^{n-p}\binom{p+k}{p}&=1+\sum_{k=1}^{n-p}\binom{p+k+1}{p+1}-\binom{p+k}{p+1}=1+\sum_{k=2}^{n-p+1}\binom{p+k}{p+1}
-\sum_{k=1}^{n-p}\binom{p+k}{p+1}\\
 &=1+\binom{n+1}{p+1}-1=\binom{n+1}{p+1}.
\end{align*}

Interprétation dans le triangle de \textsc{Pascal}. Quand on descend dans le triangle de \textsc{Pascal}, le long de la
colonne $p$, du coefficient $\binom{p}{p}$ (ligne $p$) au coefficient $\binom{p}{n}$ (ligne $n$), et que l'on additionne ces
coefficients, on trouve$\binom{n+1}{p+1}$ qui se trouve une ligne plus bas et une colonne plus loin.

\item 
\begin{enumerate}
\item Pour $n$ naturel donné, posons $I_n=\int_{0}^{1}(1-x^2)^n\;dx$. Une intégration par parties fournit:

\begin{align*}
I_n-I_{n+1}&=\int_{0}^{1}((1-x^2)^n-(1-x^2)^{n+1})\;dx=\int_{0}^{1}x^2(1-x^2)^n\;dx=\int_{0}^{1}x.x(1-x^2)^{n+1}\;dx\\
 &=\left[-x\frac{(1-x^2)^{n+1}}{2(n+1)}\right]_{0}^{1}+\frac{1}{2(n+1)}\int_{0}^{1}(1-x^2)^{n+1}\;dx
=\frac{1}{2(n+1)}I_{n+1}
\end{align*}

et donc $2(n+1)(I_n-I_{n+1})=I_{n+1}$ ou encore~:

$$\forall n\in\Nn,\;(2n+3)I_{n+1}=2(n+1)I_n.$$

On a déjà $I_0=1$. Puis, pour $n\geq1$,

\begin{align*}
I_n=\frac{2n}{2n+1}I_{n-1}=\frac{2n}{2n+1}\frac{2n-2}{2n-1}...\frac{2}{3}I_0=\frac{(2n)(2n-2)...2}{(2n+1)(2n-1)...3
.1}.
\end{align*}

\item Pour $n$
naturel non nul donné~:

\begin{align*}
1-\frac{\binom{n}{1}}{3}+\frac{\binom{n}{2}}{5}+...+(-1)^n\frac{\binom{n}{n}}{2n+1}&=\int_{0}^{1}(1-\binom{n}{1}x^2+\binom{n}{2}x^4+...
+(-1)^n\binom{n}{n}x^{2n})\;dx\\
 &=\int_{0}^{1}(1-x^2)^n\;dx=I_n
=\frac{(2n)(2n-2)...2}{(2n+1)(2n-1)...3.1}.
\end{align*}

\end{enumerate}
\end{enumerate}
\fincorrection
\correction{005138}
La formule du binôme de \textsc{Newton} fournit

$$(a-b+2c)^9=\sum_{k=0}^{9}\binom{9}{k}(a-b)^k(2c)^{9-k}=(a-b)^9+...+\binom{9}{6}(a-b)^6(2c)^3+...+(2c)^9.$$

Ensuite,

$$(a-b)^6=\sum_{k=0}^{6}\binom{6}{k}a^k(-b)^{6-k}=a^6-...+\binom{6}{4}a^4b^2-..+b^6.$$

Le coefficient cherché est donc

$$\binom{9}{6}\binom{6}{4}2^3=\frac{9.8.7}{3.2}\frac{6.5}{2}.2^3=3.4.7.3.5.8=10080.$$
\fincorrection
\correction{005139}
 $$(a+b+c+d)^2=a^2+b^2+c^2+d^2+2(ab+ac+ad+bc+bd+cd)$$ et

$$(a+b+c)^3=a^3+b^3+c^3+3(a^2b+ab^2+a^2c+ac^2+b^2c+bc^2)+6abc.$$
\fincorrection
\correction{005140}
Soit $n$ un entier naturel non nul. Le terme général du développement de $(a+b)^n$ est $u_k=\binom{n}{k}a^kb^{n-k}$, $0\leq
k\leq n$. Pour $0\leq k\leq n-1$, on a~:

$$\frac{u_{k+1}}{u_k}=\frac{\binom{n}{k+1}a^{k+1}b^{n-k-1}}{\binom{n}{k}a^kb^{n-k}}=\frac{n-k}{k+1}\frac{a}{b}.$$

Par suite,

$$\frac{u_{k+1}}{u_k}>1\Leftrightarrow\frac{n-k}{k+1}\frac{a}{b}>1\Leftrightarrow(n-k)a>(k+1)b\Leftrightarrow k<\frac{na-b}{a+b}.$$

\begin{itemize}
\item[1er cas.] Si $\frac{na-b}{a+b}>n-1$ (ce qui équivaut à $n<\frac{a}{b}$), alors la suite $(u_k)_{0\leq k\leq n}$
est strictement croissante et le plus grand terme est le dernier~:~$a^n$.

\item[2ème cas.] Si $\frac{na-b}{a+b}\leq0$ (ce qui équivaut à $n\leq\frac{b}{a}$), alors la suite $(u_k)_{0\leq k\leq
n}$ est strictement décroissante et le plus grand terme est le premier~:~$b^n$.

\item[3ème cas.] Si $0<\frac{na-b}{a+b}\leq n-1$. Dans ce cas, la suite est strictement croissante puis éventuellement
momentanément constante, suivant que $\frac{na-b}{a+b}$ soit un entier ou non, puis strictement décroissante (on dit
que la suite u est unimodale).

Si $\frac{na-b}{a+b}\notin\Nn$, on pose $k=E(\frac{na-b}{a+b})+1$, la suite $u$ croit strictement jusqu'à ce rang
puis redécroit strictement. Le plus grand des termes est celui d'indice $k$, atteint une et une seule fois.

Si $\frac{na-b}{a+b}\in\Nn$, le plus grand des termes est atteint deux fois à l'indice $k$ et à l'indice $k+1$.
\end{itemize}
\fincorrection
\correction{005141}
Pour $n\geq3$,
\begin{align*}
\binom{n}{1}+\binom{n}{2}+\binom{n}{2}=5n&\Leftrightarrow n+\frac{n(n-1)}{2}+\frac{n(n-1)(n-2)}{6}=5n\\
 &\Leftrightarrow n(-24+3(n-1)+(n-1)(n-2))=0\Leftrightarrow n^2-25=0\\
 &\Leftrightarrow n=5.
\end{align*}
\fincorrection
\correction{005147}
$(1+a)^n=(1+a)\ldots(1+a)=1+na+...\geq1+na$.
\fincorrection
\correction{005158}
Soit $n\in\Nn^*$.
\begin{enumerate}
\item  La formule du binôme de \textsc{Newton} permet d'écrire

\begin{align*}
(2+\sqrt{3})^n&=\sum_{k=0}^{n}\binom{n}{k}\sqrt{3}^k2^{n-k}=\sum_{k=0}^{E(n/2)}\binom{n}{2k}\sqrt{3}^{2k}2^{n-2k}+
\sum_{k=0}^{E((n-1)/2)}\binom{n}{2k+1}\sqrt{3}^{2k+1}2^{n-2k-1}\\
 &=\sum_{k=0}^{E(n/2)}\binom{n}{2k}3^{k}2^{n-2k}+
\sqrt{3}\sum_{k=0}^{E((n-1)/2)}\binom{n}{2k+1}3^k2^{n-2k-1}.
\end{align*}

Ainsi, en posant $a_n=\sum_{k=0}^{E(n/2)}\binom{n}{2k}3^{k}2^{n-2k}$ et $b_n=\sum_{k=0}^{E((n-1)/2)}\binom{n}{2k+1}3^k2^{n-2k-1}$,
$a_n$ et $b_n$ sont des \textbf{entiers} tels que $(2+\sqrt{3})^n=a_n+b_n\sqrt{3}$. En remplaçant $\sqrt{3}$ par
$-\sqrt{3}$, on a aussi $(2-\sqrt{3})^n=a_n-b_n\sqrt{3}$. Mais alors,

$$a_n^2-3b_n^2=(a_n+b_n\sqrt{3})(a_n-b_n\sqrt{3})=(2+\sqrt{3})^n(2-\sqrt{3})^n=(4-3)^n=1.$$
\item  On note que $(2+\sqrt{3})^n+(2-\sqrt{3})^n=(a_n+b_n\sqrt{3})+(a_n-b_n\sqrt{3})=2a_n$. Mais,

$$0<(2-\sqrt{3})^n<1.$$

Par suite,

$$(2+\sqrt{3})^n<(2+\sqrt{3})^n+(2-\sqrt{3})^n=2a_n<(2+\sqrt{3})^n+1,$$

ou encore

$$2a_n-1<(2+\sqrt{3})^n<2a_n.$$

On en déduit que $E((2+\sqrt{3})^n)=2a_n-1$ et donc que$E((2+\sqrt{3})^n)$ est un entier impair.
\end{enumerate}
\fincorrection
\correction{005278}
\begin{enumerate}
\item  Soit $E$ un ensemble à $n$ éléments, $n\geq 1$, et $a$ un élément fixé de $E$. Soit $\begin{array}[t]{cccc}
f~:&\mathcal{P}(E)&\rightarrow&\mathcal{P}(E)\\
 &A&\mapsto&\left\{
 \begin{array}{l}
 A\setminus\{a\}\;\mbox{si}\;a\in A\\
 A\cup\{a\}\;\mbox{si}\;a\notin A
 \end{array}
 \right.
\end{array}$.

Montrons que $f$ est involutive (et donc bijective). Soit $A$ un élément de $\mathcal{P}(E)$.

Si $a\notin A$, $f(A)=A\cup\{a\}$ et donc, puisque $a\in A\cup\{a\}$, $f(f(A))=(A\cup\{a\})\setminus\{a\}=A$.

Si $a\in A$, $f(A)=A\setminus\{a\}$ et $f(f(A))=(A\setminus\{a\})\cup\{a\}=A$.

Ainsi, $\forall A\in\mathcal{P}(E),\;f\circ f(A)=A$ ou encore, $f\circ f=Id_{\mathcal{P}(E)}$.

Maintenant clairement, en notant $\mathcal{P}_p(E)$ (resp. $\mathcal{P}_i(E)$) l'ensemble des parties de $E$ de cardinal pair (resp. impair), $f(\mathcal{P}_p(E))\subset\mathcal{P}_i(E)$ et $f(\mathcal{P}_i(E))\subset\mathcal{P}_p(E)$. Donc, puisque $f$ est bijective 
$$\mbox{card}(\mathcal{P}_p(E))=\mbox{card}(f(P_p(E))\leq\mbox{card}\mathcal{P}_i(E)$$

et de même $\mbox{card}(\mathcal{P}_i(E))\leq\mbox{card}\mathcal{P}_p(E)$. Finalement, $\mbox{card}(\mathcal{P}_i(E))=\mbox{card}\mathcal{P}_p(E)$.

\item  
Soient $E=\{a_1,...,a_n\}$ un ensemble à $n$ éléments et $a$ un élément fixé de $E$. Soit $k\in\{1,...,n-1\}$.

Il y a  $C_{n-1}^{k-1}$ parties à $k$ éléments qui contiennent $a$. Donc, $nC_{n-1}^{k-1}(=C_{n-1}^{k-1}+...+C_{n-1}^{k-1})$ est donc la somme du nombre de parties à $k$ éléments qui contiennent $a_1$ et du nombre de parties à $k$ éléments qui contiennent $a_2$ ... et du nombre de parties à $k$ éléments qui contiennent $a_n$.

Dans cette dernière somme, chaque partie à $k$ éléments de $E$ a été comptée plusieurs fois et toutes les parties à $k$ éléments (en nombre égal à $C_n^k$) ont été comptés un même nombre de fois. Combien de fois a été comptée $\{a_1,a_2...a_k\}$~?~Cette partie a été comptée une fois en tant que partie contenant $a_1$, une fois en tant que partie contenant $a_2$... et une fois comme partie contenant $a_k$ et donc a été comptée $k$ fois.

Conclusion~:~$kC_n^k=nC_{n-1}^{k-1}$. 

\item  Soit $E=\{a_1,...,a_n,b_1,...,b_n\}$ un ensemble à $2n$ éléments. Il y a $C_{2n}^n$ parties à $n$ éléments de $E$. Une telle partie a $k$ éléments dans $\{a_1,...,a_n\}$ et $n-k$ dans $\{b_1,...,b_n\}$ pour un certain $k$ de $\{0,...,n\}$. Il y a $C_n^k$ choix possibles de $k$ éléments dans $\{a_1,...,a_n\}$ et $C_n^{n-k}$ choix possibles de $n-k$ éléments dans $\{b_1,...,b_n\}$ pour $k$ donné dans $\{0,...,n\}$ et quand $k$ varie de $0$ à $n$, on obtient~:

$$C_{2n}^n=\sum_{k=0}^{n}C_n^kC_n^{n-k}=\sum_{k=0}^{n}(C_n^k)^2.$$

\end{enumerate}
\fincorrection
\correction{005280}
Clairement, $\forall n\in\Nn^*,\;a_{n,0}=1$ (unique solution~:~$0+0+...+0=0$) et $\forall k\in\Nn,\;a_{1,k}=1$ (unique solution~:~$k=k$).

Soient $n\geq 1$ et $k\geq 0$ fixés. $a_{n+1,k}$ est le nombre de solutions en nombre entiers positifs $x_i$ de l'équation $x_1+...+x_n+x_{n+1}=k$.
Il y a $a_{n,k}$ solutions telles que $x_{n+1}=0$ puis $a_{n,k-1}$ solutions telles que $x_{n+1}=1$ ... puis $a_{n,0}$ solutions telles que $x_{n+1}=k$.

Donc, $\forall n\in\Nn^*,\;\forall k\in\Nn,\;a_{n+1,k}=a_{n,k}+a_{n,k-1}+...+a_{n,0}$ (et on rappelle $a_{n,0}=a_{1,k}=1$).

Montrons alors par récurrence sur $n$, entier naturel non nul, que~:~$\forall n\in\Nn^*,\;\forall k\in\Nn,\;a_{n,k}=C_{n+k-1}^k$.
 
Pour $n=1$, on a pour tout naturel $k$, $a_{1,k}=1=C_{1+k-1}^k$.

Soit $n\geq1$, supposons que $\forall k\in\Nn,\;a_{n,k}=C_{n+k-1}^k$. Soit $k\geq1$.

$$a_{n+1,k}=\sum_{i=0}^{k}a_{n,i}=\sum_{i=0}^{k}C_{n+i-1}^i=1+\sum_{i=1}^{k}(C_{n+i}^{i+1}-C_{n+i}^i)=1+C_{n+k}^{k+1}-1=
C_{n+k}^{k+1},$$

ce qui reste clair pour $k=0$.
 
On a montré par récurrence que $\forall n\in\Nn^*,\;\forall k\in\Nn,\;a_{n,k}=C_{n+k-1}^k$.
\fincorrection
\nocorrection
\correction{000236}
Tout d'abord si deux ensembles finis $A$ et $B$ sont disjoints
alors $\text{Card\,} A\cup B = \text{Card\,} A + \text{Card\,} B$.

Si maintenant $A$ et $B$ sont deux ensembles finis quelconques : nous décomposons $A\cup B$ en trois ensembles :
$$A \cup B = (A \setminus (A\cap B)) \cup (B \setminus (A\cap B))  \cup (A\cap B).$$
Ces trois ensembles sont disjoints deux à deux donc :
$\text{Card\,} A\cup B = \text{Card\,} A\setminus (A\cap B) + \text{Card\,} B \setminus (A\cap B) + \text{Card\,} A\cap B$.

Mais pour $R \subset S$ nous avons
$\text{Card\,} S\setminus R = \text{Card\,} S - \text{Card\,} R$.

Donc 
$\text{Card\,} A\cup B = \text{Card\,} A - \text{Card\,} A\cap B  + \text{Card\,}B - \text{Card\,} A\cap B  + \text{Card\,} A\cap B$.

Donc $\text{Card\,} A\cup B = \text{Card\,} A + \text{Card\,} B - \text{Card\,} A\cap B$.


Appliquons ceci à $A \Delta B  = (A \cup B) \setminus (A \cap B)$ :
$$ \text{Card\,} A \Delta B = \text{Card\,} A\cup B - \text{Card\,} A\cap B =
\text{Card\,} A + \text{Card\,} B - 2\text{Card\,} A\cap B.$$
\fincorrection
\correction{000237}
Fixons un \'el\'ement de $A$ ; dans $E\setminus A$ (de cardinal
$n-p$), nous pouvons choisir $C_{n-p}^k$ ensembles \`a $k$
\'el\'ements ($k = 0,1,\ldots,n$). Le nombre d'ensembles dans le
compl\'ementaire de $A$ est donc
$$ \sum_{k=0}^{n-p} C_{n-p}^{k} = 2^{n-p}.$$

Pour le choix d'un \'el\'ement de $A$ nous avons $p$ choix, donc
le nombre total d'ensembles qui v\'erifie la condition est :
$$p2^{n-p}.$$
\fincorrection
\nocorrection
\correction{000239}
\begin{enumerate}
  \item  Il s'agit donc de choisir $5$ cartes parmi $52$ : il y a donc $C_{52}^5$ mains différentes.
Ceci peut être calculé : $C_{52}^5 = \frac{52\cdot51\cdot50\cdot49\cdot48}{5!} = 2598960$.

  \item Il y a $4$ choix pour l'as (l'as de pique ou l'as de c{\oe}ur ou ...), puis il faut choisir les $4$ cartes restantes
parmi $48$ cartes (on ne peut pas rechoisir un as). Bilan $4 \times C_{48}^4$ mains comprenant exactement un as.

  \item Il est beaucoup plus facile de compter d'abord les mains qui ne contiennent aucun valet :
il faut choisir $5$ cartes parmi $48$ (on exclut les valets) ; il y a donc $C_{48}^5$ mains ne contenant aucun valet.
Les autres mains sont les mains qui contiennent au moins un valet : il y en a donc $C_{52}^5 - C_{48}^5$.

  \item Nous allons d'abord compter le nombre de mains que ne contiennent pas de roi ou pas de dame.
Le nombre de mains qui ne contiennent pas de roi est $C_{48}^5$ (comme la question 3.). Le nombre de mains qui ne contiennent
pas de dame est aussi $C_{48}^5$. 
Le nombre de mains ne contenant pas de roi ou pas de dame \emph{n'est pas} $C_{48}^5 + C_{48}^5$,
car on aurait compté deux fois les mains ne contenant ni roi, ni dame (il y a $C_{44}^5$ telles mains).
Le nombre de mains ne contenant pas de roi ou pas de dame est donc : $2C_{48}^5-C_{44}^5$ (on retire une fois les mains comptées deux fois !).
Ce que nous cherchons ce sont toutes les autres mains : celles qui contiennent au moins un roi et au moins une dame.
Leur nombre est donc : $C_{52}^5 - 2C_{48}^5 + C_{44}^5$.

\end{enumerate}

\fincorrection
\nocorrection
\nocorrection
\nocorrection
\correction{002912}
\begin{enumerate}
  \item $(6!)^2$
  \item $4!\times8!$
  \item $2!2!4!4!$
  \item $6^6\times12^6$, $4^4\times12^8$, $2^2\times4^2\times6^4\times12^4$.
\end{enumerate}
\fincorrection
\correction{002913}
\begin{enumerate}
  \item $(2n)!$.
  \item $2(n!)^2$.
  \item $2^{n+1}\times n!$.
  \item $4\times n!$.
\end{enumerate}
\fincorrection
\correction{002914}
\begin{enumerate}
  \item $n^{n^2}$.
  \item $n^{n(n+1)/2}$.
  \item $n\times n^{(n-1)^2}$.
  \item $n\times n^{n(n-1)/2}$.
\end{enumerate}
\fincorrection
\correction{002915}
\begin{enumerate}
  \item 
  \item
  \item 
  \begin{enumerate}
    \item $\sum_{k=1}^n (-1)^kC_n^k(n-k)^p$.
    \item $\sum_{k=1}^n \frac{(-1)^kn!}{k!}$.
  \end{enumerate}
\end{enumerate}
\fincorrection
\correction{002916}
\begin{enumerate}
  \item 
  \item R{\'e}currence. {\'E}galit{\'e} pour $n \le 2$ ou les $A_i$ 3 {\`a} 3 disjoints.
\end{enumerate}
\fincorrection
\correction{002917}
$3^n$.
\fincorrection
\correction{002918}
\begin{enumerate}
\item Comme $\{ x_1-1, \dots, x_p-p\}$ est une partie quelconque de $\{0,
  \dots, n-p\}$, on a $N = C_{n-p+1}^p$.
  \item  \begin{enumerate}
         \item 
         \item  32951280099.
         \end{enumerate}	 
\end{enumerate}
\fincorrection
\correction{002919}
\begin{enumerate}
  \item $R_n = \sum_{k=0}^{n-1}C_{n-1}^kR_k$ avec $R_0 = 1$.
  \item 1,1,2,5,15,52,203.
\end{enumerate}
\fincorrection
\nocorrection
\nocorrection
\nocorrection
\correction{002923}
\begin{enumerate}
  \item 1,2,5.
  \item $t_n = \sum_{k=1}^{n-1}\,t_kt_{n-k}$.
\end{enumerate}
\fincorrection
\correction{005279}
Il y a $C_{pq}^q$ choix possibles d'une première classe. Cette première classe étant choisie, il y a $C_{pq-q}^{q}=C_{(p-1)q}^q$ choix possibles de la deuxième classe... et $C_{q}^q$ choix possibles de la $p$-ième classe. Au total, il y a $C_{pq}^qC_{(p-1)q}^q...C_{q}^q$ choix possibles d'une première classe, puis d'une deuxième ...puis d'une $p$-ième.

Maintenant dans le nombre $C_{pq}^qC_{(p-1)q}^q...C_{q}^q$, on a compté plusieurs fois chaque partition, chacune ayant été compté un nombre égal de fois.

On a compté chaque partition autant de fois qu'il y a de permutations des $p$ classes à savoir $p!$. Le nombre cherché est donc~:  

$$\frac{1}{p!}C_{pq}^qC_{(p-1)q}^q...C_{q}^q=\frac{1}{p!}\frac{(pq)!}{q!((p-1)q)!}\frac{((p-1)q)!}{q!((p-2)q)!}...
\frac{(2q)!}{q!q!}\frac{q!}{q!0!}=\frac{(pq)!}{p!(q!)^p}.$$
\fincorrection
\correction{005281}
On place le $0$ soit au chiffre des unités, soit au chiffre des dizaines, soit au chiffre des centaines, soit au chiffre des milliers (mais pas au chiffre des dizaines de milliers) et le $0$ étant placé, on n'y a plus droit.

Réponse~:~$4.9.9.9.9=4.9^4=4.(80+1)^2=4.6561=26244$.

\fincorrection
\correction{005284}

On pose $H=$ "vers le haut"  et $D=$ "vers la droite". 
Un exemple de chemin de $(0,0)$ à $(p,q)$ est le mot $DD...DHH...H$ 
où $D$ est écrit $p$ fois et $H$ est écrit $q$ fois. 
Le nombre de chemins cherché est clairement le nombre d'anagrammes 
du mot précédent. 


Le nombre de choix de l'emplacement du $H$ est $C_{p+q}^q$. Une fois que les lettres $H$
sont placées il n'y a plus de choix pour les lettres $D$.
Il y a donc $C_{p+q}^q$ chemins possibles.

Remarque : si on place d'abord les lettres $D$ alors on a $C_{p+q}^p$ choix possibles.
Mais on trouve bien sûr le même nombre de chemins car $C_{p+q}^p = C_{p+q}^{(p+q) - p} = C_{p+q}^q$.

\fincorrection
\correction{005285}
On note respectivement $x$, $y$ et $z$ le nombre de pièces de $10$, $20$ et $50$ centimes. Il s'agit de résoudre dans $\Nn^3$ l'équation $10x+20y+50z=10000$ ou encore $x+2y+5z=1000$.

Soit $k\in\Nn$. $x+2y=k\Leftrightarrow x=k-2y$ et le nombre de solutions de cette équation est~:

$$\sum_{k=0}^{E(k/2)}1=E(\frac{k}{2})+1.$$

Pour $0\leq z\leq 200$ donné, le nombre de solutions de l'équation $x+2y=1000-5z$ est donc $E(\frac{1000-5z}{2})+1$. Le nombre de solutions en nombres entiers de l'équation $x+2y+5z=1000$ est donc 

$$\sum_{z=0}^{200}(E(\frac{1000-5z}{2})+1)=\sum_{z=0}^{200}(E(\frac{-5z}{2})+ 501)=201.501+\sum_{z=0}^{200}E(\frac{-5z}{2})=100701+\sum_{z=0}^{200}E(\frac{-5z}{2}).$$
Maintenant  

$$\sum_{z=0}^{200}E(\frac{-5z}{2})=\sum_{k=1}^{100}(E(\frac{-5(2k-1)}{2})+E(\frac{-5(2k)}{2}))=\sum_{k=1}^{100}(E(-5k+\frac{5}{2})-5k)=\sum_{k=1}^{100}(-10k+2)=200-10\frac{100.101}{2}.$$

Le nombre de solutions cherchés est donc $100701-50300=50401$. Il y a $50401$ façons de payer $100$ euros avec des pièces de $10$, $20$ et $50$ centimes.
\fincorrection
\correction{005286}
\begin{enumerate}
\item  
\begin{align*}\ensuremath
\chi_{A_1\cup...\cup A_n}&=1-\chi_{\overline{A_1\cup...\cup A_n}}
=1-\chi_{\overline{A_1}\cap...\cap\overline{A_n}}=1-\chi_{\overline{A_1}}\times...\times\chi_{\overline{A_n}}\\
 &=1-(1-\chi_{A_1})...(1-\chi_{A_n})=1-(\sum_{k=1}^{n}(-1)^{k-1}(\sum_{1\leq i_1<...<i_k\leq n}^{}\chi_{A_{i_1}}...\chi_{A_{i_k}}))
\end{align*}

et en sommant sur l'ensemble des $x$ de $E$, on obtient le résultat.

\item  Pour $1\leq k\leq n$, posons $A_k=\{\sigma\in S_n/\;\sigma(k)=k\}$. L'ensemble des permutations ayant au moins un point fixe est $A_1\cup A_2\cup...\cup A_n$. L'ensemble des permutations sans points fixes est le complémentaire dans $S_n$ de $A_1\cup A_2\cup...\cup A_n$.

D'après 1), leur nombre est donc~:

\begin{align*}\ensuremath
\mbox{card}(S_n)-\mbox{card}(A_1\cup A_2...\cup A_n)
&=\mbox{card}(S_n)-\sum_{i=1}^{n}\mbox{card}(A_i)+\sum_{i<j}^{}\mbox{card}(A_i\cap A_j)\\
 &-...+(-1)^{k}\sum_{i_1<i_2<...<i_k}^{}\mbox{card}(A_{i_1}\cap...A_{i_k})+...\\
 &+(-1)^{n}\mbox{card}(A_1\cap...\cap A_n).
\end{align*} 

$A_i$ est l'ensemble des permutations qui fixent $i$. Il y en a $(n-1)!$ ( nombre de permutations de $\{1,...,n\}\setminus\{i\}$). $A_i\cap A_j$ est l'ensemble des permutations qui fixent $i$ et $j$. Il y en a $(n-2)!$. Plus généralement, $\mbox{card}(A_{i_1}\cap...\cap A_{i_k})=(n-k)!$.

D'autre part, il y a $n=C_n^1$ entiers $i$ dans $\{1,n\}$ puis $C_n^2$ couples $(i,j)$ tels que $i<j$ et plus généralement, il y a $C_n^k$ $k$-uplets $(i_1,...,i_k)$ tels que $i_1<i_2< ...<i_k$. Le nombre de dérangements est  

$$n!+\sum_{k=1}^{n}(-1)^k\frac{n!}{k!(n-k)!}(n-k)!=n!\sum_{k=0}^{n}\frac{(-1)^k}{k!}.$$
 
Ainsi le \og~problème des chapeaux~\fg admet pour réponse

$$p_n=\sum_{k=0}^{n}\frac{(-1)^k}{k!}.$$ 

Montrons que cette suite tend très rapidement vers $\frac{1}{e}=0,36...$ quand $n$ tend vers l'infini.

(On adapte un  calcul déjà mené pour le nombre $e$.)

Montrons que $\forall n\in\Nn,\;e^{-1}=\sum_{k=0}^{n}\frac{(-1)^k}{k!}+(-1)^{n+1}\int_{0}^{1}\frac{(1-t)^n}{n!}e^{-t}\;dt$.

Pour $n=0$, $(-1)^{n+1}\int_{O}^{1}\frac{(1-t)^n}{n!}e^{-t}\;dt=-\int_{0}^{1}e^{-t}\;dt=-1+e^{-1}$ et donc, on a bien $e^{-1}=1-\int_{0}^{1}e^{-t}\;dt$.

Soit $n\geq0$. Supposons que $e^{-1}=\sum_{k=0}^{n}\frac{(-1)^k}{k!}+(-1)^{n+1}\int_{0}^{1}\frac{(1-t)^n}{n!}e^{-t}\;dt$.

Une intégration par parties fournit

$$\int_{0}^{1}\frac{(1-t)^n}{n!}e^{-t}\;dt=\left[-\frac{(1-t)^{n+1}}{(n+1)!}e^{-t}\right]_0^1-\int_{0}^{1}\frac{(1-t)^{n+1}}{(n+1)!}e^{-t}\;dt=\frac{1}{(n+1)!}-\int_{0}^{1}\frac{(1-t)^{n+1}}{(n+1)!}e^{-t}\;dt.$$

Mais alors,

$$e^{-1}=\sum_{k=0}^{n+1}\frac{(-1)^k}{k!}+(-1)^{n+2}\int_{0}^{1}\frac{(1-t)^{n+1}}{(n+1)!}e^{-t}\;dt.$$

Le résultat est démontré par récurrence.

On en déduit que

$$|p_n-\frac{1}{e}|=\left|(-1)^{n+1}\int_{0}^{1}\frac{(1-t)^n}{n!}e^{-t}\;dt\right|
=\int_{0}^{1}\frac{(1-t)^n}{n!}e^{-t}\;dt\leq\int_{0}^{1}\frac{(1-t)^n}{n!}\;dt=\frac{1}{(n+1)!}.$$

Ceci montre que $p_n$ tend très rapidement vers $\frac{1}{e}$.

\end{enumerate}
\fincorrection
\correction{005287}
Soit $n$ un naturel non nul. Dire que $f$ est une surjection de $\{1,...,n+1\}$ sur $\{1,...,n\}$ équivaut à dire que deux des entiers de $\{1,...,n+1\}$ ont même image $k$ par $f$ et que les autres ont des images deux à deux distinctes et distinctes de $k$. On choisit ces deux entiers : $C_{n+1}^2$ choix et leur image commune : $n$ images possibles ce qui fournit $nC_{n+1}^2$ choix d'une paire de $\{1,...,n+1\}$ et de leur image commune. Puis il y a $(n-1)!$ choix des images des $n-1$ éléments restants. Au total, il y a $n!\frac{n(n+1)}{2}=\frac{n.(n+1)!}{2}$ surjections de $\{1,...,n+1\}$ sur $\{1,...,n\}$.
\fincorrection
\correction{005288}
Soit $n\geq5$. De chaque sommet part $n-1$ droites (vers les $n-1$ autres sommets) dont $2$ sont des cotés et $n-3$ des diagonales. Comme chaque diagonale passe par 2 sommets , il y a $\frac{n(n-3)}{2}$ diagonales.

Ces diagonales se recoupent en $C_{n(n-3)/2}^2$ points distincts ou confondus. Dans ce décompte, chaque sommet a été compté autant de fois que l'on a choisi une paire de deux diagonales passant par ce sommet à savoir $C_{n-3}^2$. Maintenant, il y a $n$ sommets.

Réponse : 
\begin{align*}\ensuremath 
C_{n(n-3)/2}^2-nC_{n-3}^2&=\frac{1}{2}\frac{n(n-3)}{2}(\frac{n(n-3)}{2}-1)-n\frac{(n-3)(n-4)}{2}=\frac{n(n-3)}{8}(n(n-3)-2-4(n-4))\\
 &=\frac{n(n-3)}{8}(n^2-7n+14)
\end{align*}

Les diagonales se recoupent en $\frac{n(n-3)(n^2-7n+14)}{8}$ points distincts ou confondus et distincts des sommets (ou encore en $\frac{n(n-3)(n^2-7n+14)}{8}$ points au maximum).
\fincorrection
\correction{005289}
\begin{enumerate}
\item  On a bien sûr $P(1)=2$. Soit $n\geq1$. On trace $n$ droites vérifiant les conditions de l'énoncé. Elles partagent le plan en $P(n)$ régions. On trace ensuite $D_{n+1}$, une $(n+1)$ème droite. Par hypothèse, elle coupe chacune des $n$ premières droites en $n$ points deux à deux distincts. Ces $n$ points définissent $(n+1)$ intervalles sur la droite $D_{n+1}$. Chacun de ces $(n+1)$ intervalles partage une des $P(n)$ régions déjà existantes en deux régions et rajoute donc une nouvelle région. Ainsi, $P(n+1)=P(n)+(n+1)$.

Soit $n\geq2$.

\begin{align*}\ensuremath
P(n)&=P(1)+\sum_{k=1}^{n-1}(P(k+1)-P(k))=2+\sum_{k=1}^{n-1}(k+1)=1+\sum_{k=1}^{n}k=1+\frac{n(n+1)}{2}\\
 &=\frac{n^2+n+2}{2}
\end{align*}

ce qui reste vrai pour $n=1$.

\item  On a bien sûr $Q(1)=2$.Soit $n\geq1$. On trace $n$ plans vérifiant les conditions de l'énoncé. Ils partagent l'espace en $Q(n)$ régions. On trace ensuite $P_{n+1}$, un $(n+1)$ème plan. Par hypothèse, il recoupe chacun des $n$ premiers plans en $n$ droites vérifiant les conditions du 1). Ces $n$ droites délimitent $P(n)=1+\frac{n(n+1)}{2}$ régions sur le plan $P_{n+1}$. Chacune de ces régions partage une des $Q(n)$ régions déjà existantes en deux régions et rajoute donc une nouvelle région. Ainsi, $Q(n+1)=Q(n)+P(n)=Q(n)+\frac{n^2+n+2}{2}$.

Soit $n\geq2$.

\begin{align*}\ensuremath
Q(n)&=P(1)+\sum_{k=1}^{n-1}(Q(k+1)-Q(k))=2+\sum_{k=1}^{n-1}\frac{k^2+k+2}{2})=2+(n-1)+\frac{1}{2}\sum_{k=1}^{n-1}k^2+\frac{1}{2}\sum_{k=1}^{n-1}k\\
 &=(n+1)+\frac{(n-1)n(2n-1)}{12}+\frac{n(n-1)}{4}=\frac{n^3+5n+6}{6}
\end{align*}
\end{enumerate}
\fincorrection
\correction{005290}
Soient $n$ et $k$ des entiers naturels tels que $2\leq k\leq n-1$.

Soit $E$ un ensemble à $n$ éléments et $a$ un élément fixé de $E$.

Il y a $P_n^k$ partitions de $E$ en $k$ classes. Parmi ces partitions, il y a celles dans lesquelles $a$ est dans un singleton. Elles s'identifient aux partitions en $k-1$ classes de $E\setminus\{a\}$ et sont au nombre de $P_{n-1}^{k-1}$. Il y a ensuite les partitions dans lesquelles $a$ est élément d'une partie de cardinal au moins $2$. Une telle partition est obtenue en partitionnant $E\setminus\{a\}$ en $k$ classes puis en adjoignant à l'une de ces $k$ classes au choix l'élément $a$. Il y a $kP_{n-1}{k}$ telles partitions. Au total, $P_n^k=P_{n-1}^{k-1}+kP_{n-1}^k$.

Valeurs de $P_n^k$ pour $1\leq k,n\leq 5$.

$$\begin{array}{|c|c|c|c|c|c|}
\hline
n\;/\;k&1&2&3&4&5\\
\hline
1&1&0&0&0&0\\
\hline
2&1&1&0&0&0\\
\hline
3&1&3&1&0&0\\
\hline
4&1&7&6&1&0\\
\hline
5&1&15&25&10&1\\
\hline
\end{array}$$
Exprimons maintenant en fonction des $P_n^k$, le nombre de surjections d'un ensemble à $n$ éléments dans un ensemble à $p$ éléments.

Si $p>n$, il n'y a pas de surjections de $E_n$ dans $E_p$ (où $E_n$ et $E_p$ désignent des ensembles à $n$ et $p$ éléments respectivement).

On suppose dorénavant $p\leq n$. La donnée d'une surjection $f$ de $E_n$ sur $E_p$ équivaut à la donnée d'une partition de l'ensemble $E_n$ en $p$ classes (chaque élément d'une même classe ayant même image par $f$) puis d'une bijection de l'ensemble des parties de la partition vers $E_p$.

Au total, il y a donc $p!P_n^k$ surjections d'un ensemble à $n$ éléments dans un ensemble à $p$ éléments pour $1\leq p\leq n$.
\fincorrection
\nocorrection
\nocorrection
\nocorrection
\nocorrection
\nocorrection
\nocorrection
\nocorrection
\nocorrection
\nocorrection
\nocorrection
\nocorrection
\nocorrection
\correction{003057}
\begin{enumerate}
 \item  $a=33$, $b=-200$.
 \item
\end{enumerate}
\fincorrection
\nocorrection
\nocorrection
\correction{003060}
\begin{enumerate}
 \item R{\'e}currence sur $n$.
 \item Si $a > n$, alors $a = b+1$ avec $b \ge n$, donc $f(b) \ge n$,
    donc $f(f(b)) \ge f(a)$. Contradiction.
 \item  
\end{enumerate}
\fincorrection
\correction{005282}
Si $n\geq366$, on a clairement $p_n=1$ (Principe des tiroirs~:~si $366$ personnes sont à associer à $365$ dates d'anniversaire, alors $2$ personnes au moins sont à associer à la même date d'anniversaire).

Si $2\leq n\leq 365$, on a $p_n=1-q_n$ où $q_n$ est la probabilité que les dates d'anniversaire soient deux à deux distinctes. Il y a $(365)^n$ répartitions possibles des dates d'anniversaires (cas possibles) et parmi ces répartitions, il y en a $365.364.363....(365-n+1)$ telles que les dates d'anniversaire soient deux à deux distinctes. Finalement 

$$p_n=1-\frac{1}{(365)^n}365.364.363....(365-n+1)=1-\prod_{k=1}^{n-1}\frac{365-k}{365}=1-\prod_{k=1}^{n-1}(1-\frac{k}{365}).$$

Ensuite,

$$p_n\geq\frac{1}{2}\Leftrightarrow\prod_{k=1}^{n-1}(1-\frac{k}{365})\leq\frac{1}{2}\Leftrightarrow\sum_{k=1}^{n-1}\ln(1-\frac{k}{365}) \leq\ln\frac{1}{2}\Leftrightarrow\sum_{k=1}^{n-1}-\ln(1-\frac{k}{365})\geq\ln2.$$

Maintenant, soit $x\in[0,1[$. On a

$$-\ln(1-x)=\int_{0}^{x}\frac{1}{1-t}dt\geq\int_{0}^{x}\frac{1}{1-0}dt=x.$$

Pour $k$ élément de $\{1,...,n-1\}(\subset\{1,...,364\})$, $\frac{k}{365}$ est un réel élément de $[0,1[$.

En appliquant l'inégalité précédente, on obtient 

$$\sum_{k=1}^{n-1}-\ln(1-\frac{k}{365})\geq\sum_{k=1}^{n-1}\frac{k}{365}=\frac{n(n-1)}{730}.$$

Ainsi,

$$p_n\geq\frac{1}{2}\Leftarrow\frac{n(n-1)}{730}\geq\ln 2\Leftrightarrow n^2-n-730\ln2\geq0\Leftrightarrow n\geq\frac{1+\sqrt{1+2920\ln2}}{2}=22,99...\Leftrightarrow n\geq23.$$

Finalement, dans un groupe d'au moins $23$ personnes, il y a plus d'une chance sur deux que deux personnes au moins aient la même date d'anniversaire.

\fincorrection
\correction{005283}

\begin{enumerate}
\item  Notre calendrier est $400$ ans périodique (et presque $4.7=28$ ans périodique).
En effet,
\begin{enumerate}
\item la répartition des années bissextiles est $400$ ans périodique ($1600$ et $2000$ sont bissextiles mais $1700$, $1800$ et $1900$ ne le sont pas (entre autre pour regagner $3$ jours tous les $400$ ans et coller le plus possible au rythme du soleil))
\item il y a un nombre entier de semaines dans une période de $400$ ans. En effet, sur $400$ ans, le quart des années, soit $100$ ans, moins $3$ années sont bissextiles et donc sur toute période de $400$ ans il y a $97$ années bissextiles et $303$ années non bissextiles.

Une année non bissextile de $365$ jours est constituée de $52.7+1$ jours ou encore d'un nombre entier de semaines plus un jour et une année bissextile est constituée d'un nombre entier de semaine plus deux jours.

Une période de $400$ ans est donc constituée d'un nombre entier de semaines plus~:~$97.2+303.1=194+303=497=7.71$ jours qui fournit encore un nombre entier de semaines.
\end{enumerate}

\item  Deux périodes consécutives de $28$ ans ne contenant pas d'exception (siècles non bissextiles) reproduisent le même calendrier. En effet, les $7$ années bissextiles fournissent un nombre entiers de semaines plus $2.7$ jours $=2$ semaines et les $21$ années non bissextiles fournissent un nombre entier de semaines plus $21.1$ jours $=3$ semaines.

\item  D'après ce qui précède, il suffit de compter les 1ers de l'an qui tombe un dimanche ou un samedi sur une période de $400$ ans donnée, par exemple de $1900$ à $2299$ (inclus).

On décompose cette période comme suit~:

$$\begin{array}{c}
1900,\;1901\rightarrow1928,\;1929\rightarrow1956,\;1957\rightarrow1984,\;1985\rightarrow2012,\;2013\rightarrow2040,\;
2041\rightarrow2068,\;2069\rightarrow2096,\\
2097\rightarrow2100,\;2101\rightarrow2128,\;2129\rightarrow2156,\;2157\rightarrow2184,\;2185\rightarrow2200,\;2201\rightarrow2228\;2229\rightarrow2256,\\
2257\rightarrow2284,\;2285\rightarrow2299.
\end{array}$$

\item  On montre ensuite que sur toute période de $28$ ans sans siècle non bissextile, le premier de l'an tombe un même nombre de fois chaque jour de la semaine (Lundi, mardi,..). (La connaissance des congruences modulo $4$ et $7$ seraient bien utile). Quand on passe d'une année non bissextile à l'année suivante, comme une telle année contient un nombre entier de semaines plus un jour, le 1er de l'an tombe un jour plus tard l'année qui suit et deux jours plus tard si l'année est bissextile. Par exemple,

1er janvier 1998~:~jeudi  1999~:~vendredi  2000~:~samedi  2001~:~Lundi  2002~:~Mardi 2003~:~Mercredi 2004~:~Jeudi 2005~:~samedi...

Notons A,B,C,D,E,F,G les jours de la semaine. Sur une période de 28 ans sans siècle non bissextile finissant par exemple une année bissextile, on trouve la séquence suivante~:

ABCD FGAB DEFG BCDE GABC EFGA CDEF (puis çà redémarre ABCD...) soit 4A, 4B, 4C, 4D, 4E, 4F, et 4G.

\item  Il reste à étudier les périodes à exception (soulignées dans le 3)).

Détermination du 1er janvier 1900.
Le 1er janvier 1998 était un jeudi . Il en est donc de même du 1er janvier 1998-28 = 1970 et des premiers janvier 1942 et 1914 puis on remonte~:

1914 Jeudi 1913 Mercredi 1912 Lundi 1911 Dimanche 1910 Samedi 1909 Vendredi 1908 Mercredi 1907 Mardi
1906 Lundi 1905 Dimanche 1904 Vendredi 1903 Jeudi 1902 Mercredi 1901 Mardi 1900 Lundi (1900 n'est pas bissextile)

Les premiers de l'an 2000, 2028 , 2056 et 2084 sont des samedis, 2088 un jeudi, 2092 un mardi, 2096 un dimanche et donc 2097 mardi 2098 mercredi 2099 jeudi 2100 vendredi.

2101 est un samedi de même que 2129, 2157, 2185 ce qui donne de 2185 à 2200 inclus la séquence :

S D L Ma J V S D Ma Me J V D L Ma Me

2201 est un jeudi de même que 2285 ce qui donne de 2285 à 2299 inclus la séquence :

J V S D Ma Me J V D L Ma Me V S D
\end{enumerate}

Le décompte des Lundis , mardis ... soulignés est : 6D 4L 6Ma 5Me 5J 6V 4S. Dans toute période de 400 ans, le 1er de l'an tombe $2$ fois de plus le dimanche que le samedi et donc plus souvent le dimanche que le samedi.
\fincorrection
\correction{000249}
\'Ecrivons la d\'ecomposition de  $15 !=1.2.3.4\ldots15$ en facteurs premiers. $15 !  = 2^{11}.3^6.5^3.7^2 .11.13$.
Un diviseur de $15 !$ s'\'ecrit $d = 2^{\alpha}.3^\beta.5^\gamma.7^\delta .11^\epsilon.13^\eta$
avec $0 \leq \alpha \leq 11$, $0 \leq \beta \leq 6$, $0 \leq \gamma \leq 3$, $0 \leq \delta \leq 2$,
$0 \leq \epsilon \leq 1$, $0 \leq \eta \leq 1$. De plus tout nombre $d$ de cette forme est un diviseur de $15 !$.
Le nombre de diviseurs est donc $(11+1)(6+1)(3+1)(2+1)(1+1)(1+1) = 4032$.
\fincorrection
\correction{000250}
Il sagit de calculer $100^{1000}$ modulo $13$.
Tout d'abord  $100 \equiv 9 \pmod{13}$ donc  $100^{1000} \equiv 9^{1000} \pmod{13}$.
Or  $9^{2} \equiv 81 \equiv 3 \pmod{13}$, $9^{3} \equiv 9^2.9 \equiv 3.9 \equiv 1 \pmod{13}$, Or 
 $9^{4} \equiv 9^3.9 \equiv 9 \pmod{13}$, $9^{5} \equiv 9^4.9 \equiv 9.9 \equiv 3 \pmod{13}$. 
Donc 
$100^{1000} \equiv 9^{1000} \equiv 9^{3.333+1} \equiv (9^3)^{333}.9  \equiv 1^{333}.9 \equiv 9 \pmod{13}$.
\fincorrection
\correction{000251}
La seule chose \`a voir est que pour une division euclidienne le reste doit \^etre plus petit que le quotient.
Donc les divisions euclidiennes s'\'ecrivent :
$96842 = 256 \times 378 + 74$ et $96842 = 258 \times 375 + 92$.
\fincorrection
\nocorrection
\nocorrection
\correction{000254}
Raisonnons modulo $8$ :
$$7 \equiv -1 \pmod{8}.$$
Donc
$$7^n +1 \equiv (-1)^n + 1 \pmod{8}.$$

Le reste de la division euclidienne de $7^n+1$ par $8$ est donc
$(-1)^n+1$ donc Si $n$ est impair alors $7^n+1$ est divisible par
$8$. Et si $n$ est pair $7^n+1$ n'est pas divisible par $8$.
\fincorrection
\nocorrection
\nocorrection
\correction{000257}
Il suffit de constater que pour $4$ nombres cons\'ecutifs il y a
n\'ecessairement : un multiple de $2$, un multiple de $3$, un
multiple de $4$ (distinct du mutliple de $2$). Donc le produit de $4$ nombres
cons\'ecutifs est divisible par $2\times 3\times 4 = 24$.
\fincorrection
\nocorrection
\nocorrection
\nocorrection
\nocorrection
\nocorrection
\nocorrection
\nocorrection
\nocorrection
\nocorrection
\correction{000267}
Ecrire $n=p^2+q^2$ et \'etudier le reste de la division euclidienne de
$n$ par $4$ en distinguant les diff\'erents cas de parit\'e de $p$ et $q$. 
\fincorrection
\nocorrection
\nocorrection
\correction{000270}
Pour \ref{div}. Si $p$ divise $b-a$ alors $p$ divise aussi $b^n-a^n$
d'apr\`es la formule $(*)$.


Pour \ref{divv}. On utilise le r\'esultat de la question pr\'ec\'edente
avec $n=p-k-1$ pour \'ecrire $b^{p-k-1}$ en fonction de $a^{p-k-1}$
modulo $p$ dans $$\sum_{k=0}^{p-1}a^{k}b^{p-k-1}.$$ 
On peut alors conclure.
\fincorrection
\nocorrection
\nocorrection
\nocorrection
\nocorrection
\nocorrection
\nocorrection
\nocorrection
\nocorrection
\nocorrection
\nocorrection
\nocorrection
\nocorrection
\nocorrection
\nocorrection
\correction{000285}
\begin{enumerate}
  \item Soit $n$ un nombre impair, alors il s'\'ecrit $n=2p+1$ avec $p\in \Nn$.
Maintenant $n^2 = (2p+1)^2 = 4p^2+4p+1 = 4p(p+1) + 1$. Donc $n^2 \equiv 1 \pmod{8}$.
  \item Si $n$ est pair alors il existe $p\in \Nn$ tel que $n=2p$. Et $n^2 = 4p^2$.
Si $p$ est pair alors $p^2$ est pair et donc $n^2 = 4p^2$ est divisible par $8$, donc
$n^2 \equiv 0 \pmod{8}$. Si $p$ est impair alors $p^2$ est impair et donc $n^2 = 4p^2$ est divisible par $4$ mais pas par $8$, donc
$n^2 \equiv 4 \pmod{8}$. 
  \item Comme $a$ est impair alors d'apr\`es la premi\`ere question $a^2 \equiv 1 \pmod{8}$, et de m\^eme
$c^2 \equiv 1 \pmod{8}$, $c^2 \equiv 1 \pmod{8}$. Donc $a^2+b^2+c^2 \equiv 1+1+1 \equiv 3 \pmod{8}$. 
Pour l'autre reste, \'ecrivons $a = 2p+1$ et $b=2q+1$, $c=2r+1$, alors $2ab = 2(2p+1)(2q+1) = 8pq + 4(p+q)+2$.
Alors $2(ab+bc+ca)= 8pq+8qr+8pr + 8(p+q+r)+6$, donc $2(ab+bc+ca) \equiv 6 \pmod{8}$.
  \item Montrons par l'absurde que le nombre $a^2+b^2+c^2$ n'est pas le carr\'e d'un nombre entier.
Supposons qu'il existe $n\in \Nn$ tel que $a^2+b^2+c^2=n^2$. Nous savons que 
$a^2+b^2+c^2 \equiv 3 \pmod{8}$. Si $n$ est impair alors $n^2 \equiv 1 \pmod{8}$ et si $n$ est pair alors
$n^2 \equiv 0 \pmod{8}$ ou $n^2 \equiv 4 \pmod{8}$. Dans tous les cas $n^2$ n'est pas congru \`a $3$ modulo $8$.
Donc il y a une contradiction. La conclusion est que l'hypoth\`ese de d\'epart est fausse donc 
$a^2+b^2+c^2$ n'est pas un carr\'e.
Le m\^eme type de raisonnement est valide pour $2(ab+bc+ca)$.

Pour $ab+bc+ca$ l'argument est similaire : 
d'une part $2(ab+bc+ca)\equiv 6 \pmod{8}$
et d'autre part si, par l'absurde, on suppose $ab+bc+ca=n^2$ alors
selon la parit\'e de $n$ nous avons $2(ab+bc+ca)\equiv 2n^2 \equiv 2 \pmod{8}$
ou \`a $0 \pmod{8}$. Dans les deux cas cela aboutit \`a une contradiction. Nous avons montrer que
$ab+bc+ca$ n'est pas un carr\'e.
\end{enumerate}
\fincorrection
\nocorrection
\nocorrection
\nocorrection
\nocorrection
\nocorrection
\correction{003095}
$7^4\equiv 1(\mathrm{mod}\,{10})$, $7^2\equiv 1(\mathrm{mod}\, 4)$ et $7^{7^{7^{7^7}}}$
est impair donc $7^{7^{7^{7^{7^7}}}}\equiv 7\equiv 3(\mathrm{mod}\, 4)$ et
$7^{7^{7^{7^{7^{7^7}}}}}\equiv 7^3 \equiv 3(\mathrm{mod}\,{10})$.
\fincorrection
\correction{003096}
\begin{enumerate}
  \item 
  \item $x=1,y=1$ ou $x=2,y=3$.
\end{enumerate}
\fincorrection
\nocorrection
\correction{003098}
reste$=2$.
\fincorrection
\nocorrection
\nocorrection
\nocorrection
\correction{003102}
$n^2 + 3n + 5 \equiv (n-59)^2 - 88 (\mathrm{mod}\, {121})$.

Si $121$ divise $n^2 + 3n + 5$, alors $11\mid n-59  \Rightarrow $
contradiction.
\fincorrection
\nocorrection
\nocorrection
\correction{003105}
$x \equiv y (\mathrm{mod}\, 4)$.
\fincorrection
\correction{003106}
\begin{enumerate}
  \item 
  \item R{\'e}currence : $a^{4\times 10^{k+1}} - 1 = (a^{4\times10^k} - 1)
                           (a^{4\times10^k\times 9} + \dots + a^{4\times10^k\times 0})$.
  \item $x = 123456789^{800000001/3}$.
\end{enumerate}
\fincorrection
\correction{003107}
$56786730 = 2 \times 3 \times 5 \times 7 \times 11 \times 13 \times 31 \times 61$.
         Pour tous ces facteurs premiers, on a $\varphi(p) \mid 60$.
\fincorrection
\correction{003108}
L'ordre de $\dot 2$ dans $(\Z/q\Z)^*$ divise $p$ donc est
         {\'e}gal {\`a}~$p$ et cet ordre divise $\varphi(q) = q-1$.
\fincorrection
\nocorrection
\nocorrection
\correction{005116}
Montrons par récurrence que, pour $n\geq2$, $H_n$ peut s'écrire sous la forme $\frac{p_n}{q_n}$ où $q_n$ est un
entier pair et $p_n$ est un entier impair (la fraction précédente n'étant pas nécessairement irréductible mais à coup
sûr pas un entier).
Pour $n=2$, $H_2=\frac{3}{2}$ et $H_2$ est bien du type annoncé.
Soit $n\geq2$. Supposons que pour tout entier $k$ tel que $2\leq k\leq n$, on ait $H_k=\frac{p_k}{q_k}$ où $p_k$ est un
entier impair et $q_k$ est un entier pair et montrons que $H_{n+1}=\frac{p_{n+1}}{q_{n+1}}$ où $p_{n+1}$ est un 
entier impair et $q_{n+1}$ est un entier pair.
(Recherche. L'idée
$H_{n+1}=\frac{p_n}{q_n}+\frac{1}{n+1}=\frac{(n+1)p_n+q_n}{(n+1)q_n}$ ne marche à coup sur que si $(n+1)p_n+q_n$ est
impair ce qui est assuré si $n+1$ est impair et donc $n$ pair)
\begin{itemize}
\item[\textbf{1er cas.}] Si $n$ est pair, on peut poser $n=2k$ où $k\in\Nn^*$. Dans ce cas,
$H_{n+1}=\frac{(2k+1)p_n+q_n}{(2k+1)q_n}$ et $H_{n+1}$ est bien le quotient d'un entier impair par un
entier pair.

\item[\textbf{2ème cas.}] Si $n$ est impair, on pose $n=2k-1$ où $k\geq2$ (de sorte que $2k-1\geq3$).

\begin{align*}
H_{n+1}&=\sum_{i=1}^{2k}\frac{1}{i}=\sum_{i=1}^{k}\frac{1}{2i}+\sum_{i=0}^{k-1}\frac{1}{2i+1}\;\\
 &(\mbox{en séparant les
fractions de dénominateurs pairs des fractions de dénominateurs impairs})\\
 &=\frac{1}{2}\sum_{i=1}^{k}\frac{1}{i}+\sum_{i=0}^{k-1}\frac{1}{2i+1}=\frac{1}{2}H_k+\sum_{i=0}^{k-1}\frac{1}{2i+1}.
\end{align*}

Maintenant, en réduisant au même dénominateur et puisque un produit de nombres impairs est impair, on voit 
que $\sum_{i=0}^{k-1}\frac{1}{2i+1}$ est du type $\frac{K}{2K'+1}$ où $K$ et $K'$ sont des entiers. Ensuite, puisque
$2\leq k\leq2k-1=n$, par hypothèse de récurrence, $H_k=\frac{p_k}{q_k}$ où $p_k$ est un entier impair et $q_k$ un
entier pair. Après réduction au même dénominateur, on obtient

$$H_{n+1}=\frac{p_k}{2q_k}+\frac{K}{2K'+1}=\frac{(2K'+1)p_k+2Kq_k}{2q_k(2K'+1)}.$$

$2Kq_k$ est un entier pair et $(2K'+1)p_k$ est un entier impair en tant que produit de deux nombres impairs. Donc le
numérateur est bien un entier impair et puisque $2qk(2K'+1)$ est un entier pair, $H_{n+1}$ est bien dans tous les cas de
la forme désirée.
\end{itemize}
On a montré par récurrence que pour tout entier naturel $n\geq2$, $H_n$ est le quotient d'un entier impair par un entier pair
et donc n'est pas un entier.
\fincorrection
\correction{005157}
Soit $n\in\Nn$. La division euclidienne de $n$ par $25$ fournit un quotient entier $q$ et et un
reste $r$ élément de $\{0,1,...,24\}$ tels que $n=25q+r$.

On a alors

$$E(\frac{1}{3}(n+2-E(\frac{n}{25})))=E(\frac{25q+r+2-q}{3})=E(8q+\frac{r+2}{3})=8q+E(\frac{r+2}{3}),$$

et

$$E(\frac{8n+24}{25})=E(\frac{8(25q+r)+24}{25})=8q+E(\frac{8r+24}{25}).$$

Pour montrer l'égalité de l'énoncé, il reste donc à vérifier les $25$ égalités $E(\frac{r+2}{3})=E(\frac{8r+24}{25})$,
$0\leq r\leq 24$, $(*)$, ce qui peut déjà se vérifier \og~à la main~\fg.

Diminuons encore le nombre de vérifications. La division euclidienne de $r$ par $3$ s'écrit $r=3k+l$ avec $0\leq
l\leq2$. Mais alors,

$$E(\frac{r+2}{3})=k+E(\frac{l+2}{3})\;\mbox{et}\;E(\frac{8r+24}{25})=E(\frac{25k-k+8l+24}{25})=k+E(\frac{-k+8l+
24}{25}).$$

Si $l=0$, $k$ varie de $0$ à $8$ et dans ce cas, $0\leq\frac{-k+24}{25}=\frac{-k+8l+24}{25}\leq\frac{24}{25}<1$. Par
suite,

$$E(\frac{-k+8l+24}{25})=0=E(\frac{2}{3})=E(\frac{l+2}{3}).$$

On a ainsi vérifié $(*)$ quand $r\in\{0,3,6,9,12,15,18,21,24\}$.

Si $l=1$ ou $l=2$, $E(\frac{l+2}{3})=1$ et d'autre part, $k$ varie de $0$ à $7$. Dans ce cas,

$$1=\frac{-7+8+24}{25}\leq\frac{-k+8l+24}{25}\leq\frac{16+24}{25}<2$$

et donc

$$E(\frac{-k+8l+24}{25})=1=E(\frac{l+2}{3}).$$

On a ainsi vérifié $(*)$ pour les autres valeurs de $r$. Finalement, on a montré que
\begin{center}
\shadowbox{
$\forall n\in\Nn,\;E(\frac{1}{3}(n+2-E(\frac{n}{25})))=E(\frac{8n+24}{25}).$
}
\end{center}
\fincorrection
\nocorrection
\nocorrection
\nocorrection
\nocorrection
\correction{000290}
Il s'agit ici d'utiliser la d\'ecomposition des nombres en facteurs premiers.
\begin{enumerate}
  \item $126 = 2.3^2.7$ et $230 = 2.5.23$ donc le pgcd de $126$ et $230$ est $2$.
  \item $390 = 2.3.5.13$, $720 = 2^4.3^2.5$, $450 = 2.3^2.5^2$ et donc le pgcd de ces trois nombres est
$2.3.5=30$.
  \item $\pgcd(180,606,750) = 6$.
\end{enumerate}
\fincorrection
\nocorrection
\correction{000292}
Soient $a,b$ deux entiers de pgcd $18$ et de somme $360$. Soit
$a',b'$ tel que $a = 18a'$ et $b=18b'$. Alors $a'$ et $b'$ sont
premiers entre eux, et leur somme est $360/18 = 20$.

Nous pouvons facilement \'enum\'erer tous les couples d'entiers
naturels $(a',b')$ ($a'\leqslant b'$) qui v\'erifient cette condition,
ce sont les couples  :
$$(1,19),(3,17),(7,13),(9,11).$$

Pour obtenir les couples $(a,b)$ recherch\'es ($a\leqslant b$), il
suffit de multiplier les couples pr\'ec\'edents par $18$ :
$$(18,342),(54,306),(126,234),(162,198).$$
\fincorrection
\nocorrection
\nocorrection
\nocorrection
\correction{000296}
\begin{enumerate}
  \item $\pgcd(18480,9828) = 84$;
  \item $25 \times 18480 + (-47) \times 9828 = 84$.
\end{enumerate}
\fincorrection
\nocorrection
\correction{000298}
Comme le pgcd de $955$ et $183$ est $1$, donc d'après le théorème de Bézout
cette équation a des solutions. Par exemple une solution particulière est 
$(m_0,n_0) = (-32,167)$.
Les solutions sont exactement les couples $(m,n) = (m_0 - 83k,n_0+37k)$, pour $k\in \Zz$.
\fincorrection
\nocorrection
\nocorrection
\nocorrection
\nocorrection
\correction{000303}
\begin{enumerate}
  \item $a=9b+10$.
  \item Calculons le pgcd par l'algorithme d'Euclide. 
$a = 9b+10$, $b = 12345678 \times 10 + 9$, $10 = 1 \times 9 +1$.
Donc le pgcd vaut $1$; 
  \item Nous reprenons les \'equations pr\'ec\'edentes en partant de la fin:
$1 = 10 - 9$, puis nous rempla\c{c}ons $9$ gr\^{a}ce \`a la deuxi\`eme \'equation 
de l'algorithme d'Euclide:
$1 = 10 - (b - 12345678 \times 10) = -b + 1234679 \times 10$. Maintenant nous 
rempla\c{c}ons $10$ gr\^{a}ce \`a la premi\`ere \'equation:
$1 = -b+12345679 (a-9b) = 12345679a - 111111112b$.
\end{enumerate}
\fincorrection
\nocorrection
\correction{000305}
En divisant par $45$ (qui est le pgcd de $1665, 1035, 45$) nous obtenons l'\'equation \'equivalente :
$$37x+23y=1 \qquad (E)$$
Comme le pgcd de $37$ et $23$ est $1$, alors d'apr\`es le th\'eor\`eme de B\'ezout
cette \'equation $(E)$ a des solutions.

L'algorithme d'Euclide pour le calcul du pgcd de $37$ et $23$ fourni
les coefficients de Bézout: $37\times 5 + 23 \times (-8) = 1$.
Une solution particuli\`ere de $(E)$ est donc 
$(x_0,y_0) = (5,-8)$.

Nous allons maintenant trouver l'expression générale pour les solutions de l'équation $(E)$.
Soient $(x,y)$ une solution de l'équation $37x+23y=1$.
Comme $(x_0,y_0)$ est aussi solution, nous avons $37x_0+23y_0=1$.
Faisons la différence de ces deux égalités pour obtenir $37(x-x_0)+23(y-y_0)=0$.
Autrement dit 
$$37(x-x_0)=-23(y-y_0) \quad (*)$$
On en déduit que $37 | 23 (y-y_0)$, or $\pgcd(23,37)=1$ donc par le lemme de Gauss,
$37 | (y -y_0)$. (C'est ici qu'il est important d'avoir divisé par $45$ dès le début !)
Cela nous permet d'écrire $y-y_0 = 37 k$ pour un $k \in \Zz$.

Repartant de l'égalité $(*)$ : nous obtenons $37(x-x_0)=-23 \times 37 \times k$.
Ce qui donne $x-x_0 = -23 k$.
Donc si $(x,y)$ est solution de $(E)$ alors elle est de la forme :
$(x,y)= (x_0 - 23k,y_0+37k)$, avec $k \in \Zz$.

Réciproquement pour chaque  $k\in\Zz$, si $(x,y)$ est de cette forme alors c'est une solution de $(E)$
(vérifiez-le !).

Conclusion : les solutions sont 
 $$\big\lbrace (5 - 23k,-8+37k) \mid k\in \Zz\big\rbrace.$$
\fincorrection
\nocorrection
\nocorrection
\nocorrection
\nocorrection
\nocorrection
\nocorrection
\nocorrection
\nocorrection
\nocorrection
\nocorrection
\nocorrection
\nocorrection
\nocorrection
\nocorrection
\nocorrection
\nocorrection
\nocorrection
\nocorrection
\nocorrection
\nocorrection
\nocorrection
\nocorrection
\nocorrection
\nocorrection
\nocorrection
\nocorrection
\nocorrection
\nocorrection
\nocorrection
\nocorrection
\nocorrection
\correction{003111}
$(a+b) \wedge m=d$.
\fincorrection
\correction{003112}
$= |a-b|(a \wedge b)^2$ ou $3|a-b|(a \wedge b)^2$.
\fincorrection
\correction{003113}
$8(n^3 + n) = (2n+1)(4n^2-2n+5) - 5  \Rightarrow  d = (2n+1) \wedge 5
  \Rightarrow  \begin{cases} \text{si } n \equiv     2 (\mathrm{mod}\, 5), &$d = 5$ \cr
             \text{si } n \not\equiv 2 (\mathrm{mod}\, 5), &$d = 1$.\end{cases}$
\fincorrection
\correction{003114}
$(15n^2+8n+6) \wedge (30n^2+21n+13)=1$.
\fincorrection
\correction{003115}
$\{a,b\} \in \bigl\{\{50,600\},\{150,200\}\bigr\}$.
\fincorrection
\correction{003116}
$\{a,b\} \in \bigl\{\{1,192\},\{3,32\},\{7,126\},\{14,63\}\bigr\}$.
\fincorrection
\correction{003117}
$\{x,y\} = \{147,252\}$.
\fincorrection
\correction{003118}
$x = 14k,\quad y = 15k$.
\fincorrection
\correction{003119}
$x$ impair, $y = 2-x$.
\fincorrection
\correction{003120}
$x=1$ ou $y=1$.
\fincorrection
\correction{003121}
$(300,150)$, $(150,100)$, $(100,75)$, $(75,60)$, $(60,50)$.
\fincorrection
\correction{003122}
\begin{enumerate}
  \item $a^m - 1\mid (a^{qm} - 1)a^r = a^n - a^r$.
  \item $A\wedge(AQ+R) = A\wedge R$.
      Algorithme d'Euclide sur les exposants de $a$.
  \item ssi $m\mid n$.
\end{enumerate}
\fincorrection
\nocorrection
\correction{003124}
\begin{enumerate}
  \item 
  \item
  \item $x=1+11k,\quad y=-1+7k$.
\end{enumerate}
\fincorrection
\correction{003125}
\begin{enumerate}
  \item $x = 67-71k,\quad y = -89 + 95k$.
  \item $x = 24 + 53k,\quad y = 9 + 20k$.
  \item $x = -49 + 20k - 5m,\quad y = 49 - 20k + 4m,\quad z = -7 + 3k$.

\end{enumerate}
\fincorrection
\nocorrection
\correction{003127}
\begin{enumerate}
  \item $x\equiv 7422 (\mathrm{mod}\, {13860})$.
  \item $x\equiv 7 (\mathrm{mod}\, {60})$.
\end{enumerate}
\fincorrection
\correction{003128}
785.
\fincorrection
\nocorrection
\correction{000336}

\begin{enumerate}
 \item Nous savons que
$$x^b-1 = (x-1)(x^{b-1}+\cdots+x+1),$$
pour $x=2^a$ nous obtenons :
$$2^{ab}-1 = {(2^{a})}^b -1
     = (2^a-1)\left( 2^{a(b-1)}+\cdots+2^{a}+1 \right).$$
Donc $(2^a-1) | (2^{ab}-1)$.

 \item Montrons la contraposée.
Supposons que $p$ ne soit pas premier. Donc $p=ab$ avec $1<p,q<a$.
Par la question précédente $2^a-1$ divise $2^p-1$
(et $1 < 2^a-1 < 2^p-1)$. Donc $2^p - 1$ n'est pas un nombre premier.

 
 \item  Nous supposons $a \ge b$.
Nous allons montrer que faire l'algorithme d'Euclide pour le couple $(2^a-1,2^b-1)$
revient à faire l'algorithme d'Euclide pour $(a,b)$.
Tout d'abord rappellons la formule qui est à la base de l'algorithme d'Euclide :
$\pgcd(a,b) = \pgcd(a-b,b)$. 
Appliqué à $2^a-1$ et $2^b-1$ cela donne directement $\pgcd(2^a-1,2^b-1) = \pgcd(2^a-2^b,2^b-1)$.
Mais $2^a-2^b = 2^b(2^{a-b}-1)$ d'où
 $\pgcd(2^a-1,2^b-1) = \pgcd(2^b(2^{a-b}-1),2^b-1) = \pgcd(2^{a-b}-1,2^b-1)$.
La dernière égalité vient du fait $2^b$ et $2^b-1$ sont premiers entre eux (deux entiers consécutifs
sont toujours premiers entre eux).

Nous avons montrer : $\pgcd(2^a-1,2^b-1) =\pgcd(2^{a-b}-1,2^b-1)$.
Cette formule est à mettre en parallèle de $\pgcd(a,b) = \pgcd(a-b,b)$.
En itérant cette formule nous obtenons que si $a=bq+r$ alors :
 $\pgcd(2^a-1,2^b-1) = \pgcd(2^{a-bq}-1,2^b-1) = \pgcd(2^r-1,2^b-1)$
à comparer avec 
$\pgcd(a,b)=\pgcd(a-bq,b)=\pgcd(r,b)$.
Nous avons notre première étape de l'algorithme d'Euclide.
En itérant l'algorithme d'Euclide pour $(a,b)$,
nous nous arêtons au dernier reste non nul:
$\pgcd(a,b) = \pgcd(b,r) = \cdots = \pgcd(r_n,0)=r_n$.
Ce qui va donner pour nous
$\pgcd(2^a-1,2^b-1) = \pgcd(2^b-1,2^r-1) = \cdots = \pgcd(2^{r_n}-1,2^0-1) = 2^{r_n}-1$.

Bilan : $\pgcd(2^a-1,2^b-1) =  2^{\pgcd(a,b)}-1$.

\end{enumerate}

\fincorrection
\correction{000337}
Soit $a$ et $b$ des entiers premiers entre eux. Raisonnons par
l'absurde et supposons que $ab$ et $a+b$ ne sont pas premiers
entre eux. Il existe alors $p$ un nombre premier divisant
$ab$ et $a+b$. Par le lemme d'Euclide comme $p|ab$ alors
$p|a$ ou $p|b$. Par exemple supposons que $p|a$.
Comme $p|a+b$ alors $p$ divise aussi $(a+b)-a$, donc $p|b$.
$\delta$ ne divise pas $b$ cela implique que $\delta$ et $b$ sont
premiers entre eux.

D'apr\`es le lemme de Gauss, comme $\delta$ divise $ab$ et
 $\delta$ premier avec $b$ alors $\delta$ divise $a$.
Donc $p$ est un facteur premier de $a$ et de
$b$ ce qui est absurde.
\fincorrection
\nocorrection
\correction{000339}
\begin{enumerate}
  \item \'Etant donn\'e $0< i < p$, nous avons 
$$C_p^i = \frac{p!}{i!(p-i)!} = \frac{p(p-1)(p-2)\ldots(p-(i+1))}{i!}$$
Comme $C_p^i$ est un entier alors $i!$ divise $ p(p-1)\ldots(p-(i+1))$.
Mais $i!$ et $p$ sont premiers entre eux (en utilisant l'hypoth\`ese $0 < i < p$).
Donc d'apr\`es le th\'eor\`eme de Gauss: $i!$ divise $(p-1)\ldots(p-(i+1))$, autrement dit
il existe $k\in\Zz$ tel que $k i! = (p-1)\ldots(p-(i+1))$. Maintenant nous avons
$C_p^i = p k$ donc $p$ divise $C_p^i$.
  \item Il s'agit de montrer le petit th\'eor\`eme de Fermat: pour $p$ premier et $a\in\Nn^*$, alors
$a^p \equiv a \pmod{p}$. Fixons $p$. Soit l'assertion
$$(\mathcal{H}_a) \ \ \ a^p \equiv a \pmod{p}.$$
Pour $a=1$ cette assertion est vraie !
\'Etant donn\'e $a \geq 1$ supposons que $\mathcal{H}_a$ soit vraie.
Alors 
$$(a+1)^p = \sum_{i=0}^p {C_p^i}a^i.$$
Mais d'apr\`es la question pr\'ec\'edente pour $0 < i < p$, $p$ divise $C_p^i$.
En termes de modulo nous obtenons:
$$ (a+1)^p \equiv C_p^0 a^0 + C_p^pa^p \equiv 1+a^p \pmod{p}.$$
Par l'hypoth\`ese de r\'ecurrence nous savons que $a^p \equiv a \pmod{p}$, donc
$$(a+1)^p \equiv a+1 \pmod{p}.$$ Nous venons de prouver que $\mathcal{H}_{a+1}$ est vraie.
Par le principe de r\'ecurrence alors quelque soit $a\in \Nn^*$ nous avons:
$$a^p \equiv a \pmod{p}.$$
\end{enumerate}
\fincorrection
\nocorrection
\correction{000341}
\begin{enumerate}
  \item Fixons $n$ et montrons la r\'ecurrence sur $k \ge 1$.
La formule est vraie pour $k=1$.
Supposons la formule vraie au rang $k$.
Alors
\begin{align*}
(2^{2^n}-1) \times \prod_{i=0}^{k}{(2^{2^{n+i}}+1)} 
&= (2^{2^n}-1) \times \prod_{i=0}^{k-1}{(2^{2^{n+i}}+1)} \times (2^{2^{n+k}}+1) \\
&= (2^{2^{n+k}}-1)\times (2^{2^{n+k}}+1) = (2^{2^{n+k}})^2-1
= 2^{2^{n+k+1}}-1.\\
\end{align*}
Nous avons utiliser l'hypoth\`ese de r\'ecurrence dans ces \'egalit\'es.
Nous avons ainsi montrer la formule au rang $k+1$. Et donc par
le principe de r\'ecurrence elle est vraie.
  \item \'Ecrivons $m=n+k$, alors l'\'egalit\'e pr\'ec\'edente devient:
$$F_m+2 = (2^{2^n}-1) \times \prod_{i=n}^{m-1} {F_i}.$$
Soit encore :
$$F_n \times (2^{2^n}-1) \times \prod_{i=n+1}^{m-1} {F_i} \ \ \  - \ \ F_m = 2.$$
Si $d$ est un diviseur de $F_n$ et $F_m$ alors $d$ divise $2$
(ou alors on peut utiliser le th\'eor\`eme de B\'ezout). En cons\'equent $d=1$ ou $d=2$. Mais $F_n$ est impair donc $d=1$. Nous avons montrer
que tous diviseurs de $F_n$ et $F_m$ est $1$, cela signifie que
$F_n$ et $F_m$ sont premiers entre eux.
  \item Supposons qu'il y a un nombre fini de nombres premiers.
Nous les notons alors $\{p_1,\ldots,p_N\}$. Prenons alors $N+1$
nombres de la famille $F_i$, par exemple $\{F_1,\ldots,F_{N+1}\}$.
Chaque $F_i$, $i=1,\ldots,N+1$ est divisible par (au moins) un facteur 
premier $p_j$, $j=1,\ldots,N$. Nous avons $N+1$ nombres $F_i$ et seulement $N$ facteurs premiers $p_j$. Donc par le principe des tiroirs
il existe deux nombres distincts $F_k$ et $F_{k'}$ 
(avec $1 \leq k,k' \leq N+1$) qui ont un facteur premier en commun.
En cons\'equent $F_k$ et $F_{k'}$ ne sont pas premiers entre eux. Ce qui contredit la question pr\'ec\'edente. Il existe donc une infinit\'e de nombres premiers.
\end{enumerate}
\fincorrection
\nocorrection
\nocorrection
\nocorrection
\nocorrection
\nocorrection
\nocorrection
\correction{000348}
\begin{enumerate}
  \item $X$ est non vide car, par exemple pour $k=2$, $4k+3=11$ est premier.
  \item $(4k+1)(4\ell+1) = 16k\ell + 4(k+\ell)+1 = 4(4k\ell+k+\ell)+1$.
Si l'on note l'entier $k' = 4k\ell+k+\ell$ alors $(4k+1)(4\ell+1) = 4k'+1$, 
ce qui est bien de la forme voulue.
  \item Remarquons que $2$ est le seul nombre premier pair, les autres sont de la forme
$4k+1$ ou $4k+3$. Ici $a$ n'est pas divisible par $2$, supposons --par l'absurde-- 
que $a$ n'a pas de diviseur de la forme $4k+3$, alors 
tous les diviseurs de $a$ sont de la forme $4k+1$. C'est-\`a-dire que $a$ s'\'ecrit comme produit
de nombre de la forme $4k+1$, et par la question pr\'ec\'edente $a$ peut s'\'ecrire $a=4k'+1$.
Donc $a \equiv 1 \pmod{4}$. Mais comme $a = 4p_1p_2\ldots p_n -1$, $a \equiv -1 \equiv 3 \pmod{4}$.
Nous obtenons une contradiction. Donc $a$ admet une diviseur premier $p$ de la forme $p=4\ell+3$. 
  \item Dans l'ensemble $X = \{p_1,\ldots,p_n\}$ il y a tous les nombres premiers de la formes $4k+3$.
Le nombre $p$ est premier et s'\'ecrit $p = 4\ell+3$ donc $p$ est un \'el\'ement de $X$, donc
il existe $i\in \{1,\ldots,n\}$ tel que $p=p_i$.
Raisonnons modulo $p=p_i$: $a \equiv 0 \pmod{p}$ car $p$ divise $a$.
D'autre part $a=4p_1\ldots p_n - 1$ donc $a\equiv -1 \pmod{p}$. (car $p_i$ divise $p_1\ldots p_n$).
Nous obtenons une contradiction, donc $X$ est infini: il existe une infinit\'e de nombre 
premier de la forme $4k+3$.
Petite remarque, tous les nombres de la forme $4k+3$ ne sont pas des nombres premiers, par
exemple pour $k=3$, $4k+3 = 15$ n'est pas premier.
\end{enumerate}
\fincorrection
\correction{000349}
\begin{enumerate}
  \item Supposons que $a^n + 1$ est premier. Nous allons montrer la contrapos\'ee. Supposons
que $n$ n'est pas de la forme $2^k$, c'est-\`a-dire que $n=p\times q$ avec
$p$ un nombre premier $>2$ et $q\in\Nn$.
Nous utilisons la formule
$$x^p+1 = (x+1)(1-x+x^2-x^3+\ldots+x^{p-1})$$
avec $x = a^q$ :
$$a^n+1 = a^{pq}+1 = (a^q)^p+1 = (a^q+1)(1-a^q+(a^q)^2+\cdots+(a^q)^{p-1}).$$
Donc $a^q+1$ divise $a^n+1$ et comme $1 < a^q+1 < a^{n}+1$ alors $a^n+1$ n'est pas premier.
Par contraposition si $a^n+1$ est premier alor $n = 2^k$.
  \item Cette conjecture est fausse, mais pas facile \`a v\'erifier
sans une bonne calculette ! En effet pour $n=5$ nous obtenons :
$$2^{2^5} + 1 = 4294967297 = 641 \times 6700417.$$
\end{enumerate}
\fincorrection
\nocorrection
\nocorrection
\correction{003130}
$a,b,c$ 2 {\`a} 2 premiers entre eux.
\fincorrection
\correction{003131}
D{\'e}composer en facteurs premiers.
\fincorrection
\nocorrection
\nocorrection
\correction{003134}
R{\'e}currence.
\fincorrection
\correction{003135}
On suppose $a,r$ entiers sup{\'e}rieurs ou {\'e}gaux {\`a}~$2$.

$a-1\mid a^r-1$ donc $a=2$. Si $r=pq$ alors $2^p-1\mid 2^r-1$ donc $r$ est premier.

La r{\'e}ciproque est fausse, $2^{11}-1 = 23\times 89$.
\fincorrection
\correction{003136}
\begin{enumerate}
  \item 
  \item $M_{11} = 23\times 89$.
  \item 
  \item 
\end{enumerate}
\fincorrection
\nocorrection
\nocorrection
\nocorrection
\correction{003140}
\begin{enumerate}
  \item $(-1)^{(p-1)/2}\equiv 1(\mathrm{mod}\, p)$.
  \item
  \item 
  \item 
\end{enumerate}
\fincorrection
\nocorrection
\correction{003142}
498.
\fincorrection
\correction{003143}
$H_n  \Rightarrow  H_{2n}  \Rightarrow  H_{2n+1}$.
\fincorrection
\nocorrection
\nocorrection
\nocorrection
\nocorrection
\correction{000001}
Remarquons d'abord que pour $z \in \Cc$, $z \overline{z} = |z|^2$
est un nombre r\'eel, ce qui fait qu'en multipliant le dénominateur par son conjugué nous obtenons un nombre réel.
 
$$ \frac{3+6i}{3-4i}
= \frac{(3+6i)(3+4i)}{(3-4i)(3+4i)} = \frac{9-24+12i+18i}{9+16} =
\frac{-15+30i}{25} = -\frac{3}{5}+\frac{6}{5}i.$$

\bigskip

Calculons
$$\frac{1+i}{2-i}
= \frac{(1+i)(2+i)}{5} = \frac{1+3i}{5},
$$
et
$$ \left( \frac{1+i}{2-i}  \right)^2
= \left( \frac{1+3i}{5}  \right)^2 = \frac{-8+6i}{25}=
-\frac{8}{25}+\frac{6}{25}i.
$$

Donc
$$ \left( \frac{1+i}{2-i}  \right)^2 + \frac{3+6i}{3-4i}
= -\frac{8}{25}+\frac{6}{25}i  -\frac{3}{5}+\frac{6}{5}i
=-\frac{23}{25}+\frac{36}{25}i.$$


\bigskip
Soit $z = \frac{2+5i}{1-i}$. Calculons $z + \overline{z}$, nous
savons d\'ej\`a que c'est un nombre r\'eel, plus pr\'ecis\'ement :
$z = -\frac{3}{2}+\frac{7}{2}i$ et donc $ z + \overline{z} = -3$.
\fincorrection
\nocorrection
\correction{000003}
\begin{enumerate}
\item $z_1 = 2 e^{i\frac \pi 3} = 2(\cos \frac \pi 3 + i \sin \frac \pi 3) = 2 (\frac 12+ i\frac{\sqrt3}{2}) = 1+i\sqrt 3$.
\item $z_2 = 3e^{-i\frac \pi 8} = 3\cos {\pi \over 8}-3i\sin{\pi \over8}={3\sqrt{2+\sqrt2}\over 2}-{3i\sqrt{2-\sqrt2}\over 2}$.


Il nous reste à expliquer comment nous avons calculé $\cos \frac \pi 8$ et $\sin \frac\pi 8$: 
posons $\theta=\frac{\pi}{8}$,
alors $2\theta = \frac \pi 4$ et donc $\cos(2\theta)= \frac {\sqrt 2}{2} = \sin(2\theta)$.
Mais $\cos(2\theta) = 2\cos^2 \theta - 1$. Donc $\cos^2\theta = \frac{\cos(2\theta) + 1}{2} = \frac 14 (2 + \sqrt 2)$.
Et ensuite  $\sin^2 \theta = 1- \cos^2 \theta = \frac 14 (2 - \sqrt 2)$.
Comme $0 \le \theta = \frac \pi 8 \le \frac \pi 2$, $\cos \theta$ et $\sin \theta$ sont des nombres positifs. Donc
$$\cos \frac \pi 8 = \frac 12 \sqrt{2 + \sqrt 2} \quad, \quad \sin \frac \pi 8 = \frac 12 \sqrt{2 - \sqrt 2}.$$ 

\end{enumerate}
\fincorrection
\nocorrection
\nocorrection
\nocorrection
\correction{000007}
$9-7i$; \qquad
 $-6i$; \qquad
 $-0{,}3+1{,}1i$; \qquad
 $-\frac{\sqrt3}{3}-\frac{i}{3}$.
\fincorrection
\correction{000008}
$\rho =\sqrt{4+2\sqrt2}$, $\theta ={3\pi \over8}$; \qquad
$\rho =4$, $\theta =-{\pi \over 10}$; \qquad
$\rho =1$, $\theta =2\varphi +\pi $.

\fincorrection
\nocorrection
\correction{000010}
Il s'agit juste d'appliquer la formule de Moivre :
$$e^{i\theta} = \cos \theta +i \sin \theta ;$$
ainsi que les formules sur les produits de puissances :
$$e^{ia}e^{ib} = e^{i(a+b)}\text{ et  } e^{ia} / e^{ib} = e^{i(a-b)}.$$
\fincorrection
\correction{000011}
Nous avons
$$ u = \frac{\sqrt{6}-\sqrt{2}i}{2}
= \sqrt{2}\left( \frac{\sqrt{3}}{2}-\frac{i}{2} \right) =
\sqrt{2}\left( \cos\frac{\pi}{6} -i\sin\frac{\pi}{6} \right)
=\sqrt{2} e^{-i\frac{\pi}{6}}.$$ puis
$$v = 1-i = \sqrt{2}e^{-i\frac{\pi}{4}}.$$
Il ne reste plus qu'\`a calculer le quotient :
$$ \frac{u}{v} = \frac{\sqrt{2}e^{-i\frac{\pi}{6}}}{\sqrt{2}e^{-i\frac{\pi}{4}}}
= e^{-i\frac{\pi}{6}+i\frac{\pi}{4}} = e^{i\frac{\pi}{12}}.$$
\fincorrection
\nocorrection
\correction{000013}
D'apr\`es la formule de Moivre pour $e^{i\alpha}$ nous avons :
$$e^{e^{i\alpha}} = e^{\cos \alpha + i\sin \alpha}
= e^{\cos \alpha}e^{i\sin \alpha}.$$ Or $e^{\cos \alpha} > 0$ donc
l'\'ecriture pr\'ec\'edente est bien de la forme
``module-argument''.

\bigskip

De fa\c{c}on g\'en\'erale pour calculer une somme du type
$e^{iu}+e^{iv}$ il est souvent utile de factoriser par
$e^{i\frac{u+v}{2}}$. En effet
\begin{align*}
e^{iu}+e^{iv} &= e^{i\frac{u+v}{2}}\left( e^{i
\frac{u-v}{2}}+ e^{-i \frac{u-v}{2}}\right) \\
&= e^{i\frac{u+v}{2}} 2 \cos  \frac{u-v}{2} \\
&=  2 \cos  \frac{u-v}{2} e^{i\frac{u+v}{2}}.\\
\end{align*}
Ce qui est proche de l'\'ecriture en coordon\'ees polaires.

Pour le cas qui nous concerne :
$$z = e^{i\theta} + e^{2i\theta}
= e^{\frac{3i\theta}{2}} \left[ e^{-\frac{i\theta}{2}} +
e^{\frac{i\theta}{2}} \right] = 2\cos \frac{\theta}{2}
e^{\frac{3i\theta}{2}}.$$ Attention le module dans une
d\'ecomposion en forme polaire doit \^etre positif ! Donc si $\cos
\frac{\theta}{2}  \geq 0$  alors $2\cos\frac{\theta}{2}$ est le module de $z$
et $3\theta/2$ est son argument ; par contre si $\cos \frac{\theta}{2}  <
0$ le module est $2|\cos\frac{\theta}{2}|$ et l'argument $3\theta/2+\pi$ (le
$+\pi$ compense le changement de signe car $e^{i\pi} = -1$).
\fincorrection
\correction{000014}
$$\frac{1+i}{1-i} = \frac{\sqrt{2}e^{i\pi/4}}{\sqrt{2}e^{-i\pi/4}}
= e^{i\pi/2} = i.$$ On remarque $1 = i^0 = i^4 = i^8 = \cdots =
i^{32}$.
\fincorrection
\nocorrection
\nocorrection
\nocorrection
\nocorrection
\nocorrection
\correction{000020}
\'Ecrivons $z = \rho e^{i\theta}$, alors $\overline{z} = \rho
e^{-i\theta}$. Donc
\begin{align*}
P &= \prod_{k=1}^n \left(z^k+{\overline{z}}^k \right)\\
&= \prod_{k=1}^n \rho^k \left(  (e^{i\theta})^k + (e^{-i\theta})^k \right)\\
&= \prod_{k=1}^n \rho^k \left(  e^{ik\theta} + e^{-ik\theta}) \right)\\
&= \prod_{k=1}^n 2 \rho^k \cos {k\theta}\\
&= 2^n.\rho.\rho^2.\ldots.\rho^n \prod_{k=1}^n \cos {k\theta}\\
&= 2^n\rho^{\frac{n(n+1)}{2}} \prod_{k=1}^n \cos {k\theta}.\\
\end{align*}
\fincorrection
\correction{000021}
Soit $(\alpha,\beta)\in\R^2$ et $z$ le nombre complexe
$z=e^{i\alpha}+e^{i\beta}$. Soit $u=\frac{\alpha+\beta}{2}$ et
$v=\frac{\alpha-\beta}{2}$. Alors, $\alpha=u+v$ et $\beta=u-v$ et
:
\begin{align*}
   z &= e^{i\alpha}+e^{i\beta}\\
     &= e^{iu+iv} + e^{iu-iv} \\
     &= e^{iu} (e^{iv}+e^{-iv}) \\
     &= 2 \cos(v)e^{iu} \\
     &= 2 \cos(\frac{\alpha-\beta}{2})e^{i\frac{\alpha+\beta}{2}}
\end{align*}
On en d\'eduit la forme trigonom\'etrique de $z$ :
$$
   |z|=2|\cos(\frac{\alpha-\beta}{2})|
 \quad\text{ et, lorsque $\cos(\frac{\alpha-\beta}{2})\neq0$ :}\quad
 $$
$$ Arg(z)=
   \begin{cases}
      \frac{\alpha+\beta}{2} [2\pi] & \text{ si $\cos\frac{\alpha-\beta}{2}> 0$}\\
      \pi+\frac{\alpha+\beta}{2} [2\pi] & \text{ si $\cos\frac{\alpha-\beta}{2}< 0$}\\
   \end{cases}
$$
(Attention, si $\cos\frac{\alpha-\beta}{2}<0$, $z=2\cos ve^{iu}$
n'est pas la forme trigonom\'etrique de $z$ !).

Soit $n\in\N$. Calculons $z^n$ de deux façons diff\'erentes :
d'une part
$$z^n=(e^{i\alpha}+e^{i\beta})^n = \sum_{p=0}^n{C_{n}^p e^{ip\alpha}e^{i(n-p)\beta}},$$ et d'autre
part, en utilisant la forme obtenue plus haut: $z^n=2^n
\cos^n{v}\,e^{inu}$. En comparant les parties r\'eelles des
expressions obtenues on obtient :
$$
\sum_{p=0}^n C_{n}^p \cos[p\alpha+(n-p)\beta] =
   2^n\cos^n\frac{\alpha-\beta}{2}\cos(n\frac{\alpha+\beta}{2}).
$$
\fincorrection
\nocorrection
\correction{000023}
$$1+e^{i\theta}= e^{\frac{i\theta}{2}}(e^{-\frac{i\theta}{2}}+e^{\frac{i\theta}{2}})
=2\cos \frac{\theta}{2} e^{\frac{i\theta}{2}}.$$
Comme $\theta \in ]-\pi,+\pi[$ alors
  le module est 
$2\cos \frac{\theta}{2} \geq 0$ et l'argument
est $\frac{\theta}{2}$.
Géométriquement,  on trace le cercle de centre $1$ et de rayon
$1$. L'angle en $0$ du triangle $(0,1,1+e^{i\theta})$ est
$\frac{\theta}{2}$ et donc est le double de l'angle en $0$
du triangle $(0,2,1+e^{i\theta})$ qui vaut $\theta$.

C'est le résulat géométrique (théorème de l'angle au centre)
qui affirme que pour un cercle l'angle au centre est le double de l'angle inscrit.
\fincorrection
\nocorrection
\nocorrection
\nocorrection
\correction{002924}
\begin{enumerate}
  \item $|u+v| + |u-v| \ge 2|u|$ et
             $|u+v| + |u-v| \ge 2|v|$. Il y a {\'e}galit{\'e} ssi $u = \pm v$.
  \item $|z_1| + |z_2| + |z_3| + |z_4| \le
              |z_1+z_2| + |z_1-z_2| + |z_3+z_4| + |z_3-z_4|$,
              \\
             $|z_1-z_2| + |z_3-z_4| \le
              |z_1-z_2 + z_3-z_4| + |z_1-z_2 - z_3+z_4| \le
              |z_1+z_3| + |z_2+z_4| + |z_1+z_4| + |z_2+z_3|$.
\end{enumerate}
\fincorrection
\correction{002927}
$\overline{u} = -u$.
\fincorrection
\correction{005119}
D'abord on a $1+i=\sqrt{2}e^{i\pi/4}$. Les racines carrées de $1+i$ dans $\Cc$ sont donc
$\sqrt[4]{2}e^{i\pi/8}$ et $-\sqrt[4]{2}e^{i\pi/8}$.
On a aussi, pour $(x,y)\in\Rr^2$,
\begin{align*}
(x+iy)^2=1+i&\Leftrightarrow\left\{
\begin{array}{l}
x^2-y^2=1\\
x^2+y^2=\sqrt{2}\\
xy>0
\end{array}
\right.
\Leftrightarrow
\left\{
\begin{array}{l}
\rule[-4mm]{0mm}{0mm}x^2=\frac{1}{2}(\sqrt{2}+1)\\
y^2=\frac{1}{2}(\sqrt{2}-1)\\
xy>0
\end{array}
\right.\Leftrightarrow
(x,y)\in\left\{\pm\left(\sqrt{\frac{\sqrt{2}+1}{2}},\sqrt{\frac{\sqrt{2}-1}{2}}\right)\right\}.
\end{align*}
Les racines carrées de $1+i$ sont donc aussi
$\pm\left(\sqrt{\frac{\sqrt{2}+1}{2}}+i\sqrt{\frac{\sqrt{2}-1}{2}}\right)$. Puisque
$\Re(e^{i\pi/8})=\cos\frac{\pi}{8}>0$, on
obtient $\sqrt[4]{2}e^{i\pi/8}=\sqrt{\frac{\sqrt{2}+1}{2}}+i\sqrt{\frac{\sqrt{2}-1}{2}}$, ou encore

$$e^{i\pi/8}=\sqrt{\frac{\sqrt{2}+1}{2\sqrt{2}}}+i\sqrt{\frac{\sqrt{2}-1}{2\sqrt{2}}}=\frac{1}{2}\left(\sqrt{2+\sqrt{2}}+i
\sqrt{2-\sqrt{2}}\right)$$
et donc, par identification des parties réelles et imaginaires,
\begin{center}
\shadowbox{
$\cos\frac{\pi}{8}=\frac{1}{2}\sqrt{2+\sqrt{2}}\;\mbox{et}\;\sin\frac{\pi}{8}
=\frac{1}{2}\sqrt{2-\sqrt{2}}.$
}
\end{center}

\fincorrection
\correction{005127}
Soient $z$ un complexe non nul, $M$ le point d'affixe $z$ et $A$ le point d'affixe $1$.

$$|z|=\left|\frac{1}{z}\right|\Leftrightarrow|z|=\frac{1}{|z|}\Leftrightarrow|z|^2=1\Leftrightarrow|z|=1,$$
et

$$|z|=|z-1|\Leftrightarrow OM=AM\Leftrightarrow M\in\mbox{med}[OA]\Leftrightarrow x_M=\frac{1}{2}\Leftrightarrow\Re(z)=\frac{1}{2}.$$
Donc,

$$|z|=\left|\frac{1}{z}\right|=|z-1|\Leftrightarrow|z|=1\;\mbox{et}\;\Re(z)=\frac{1}{2}\Leftrightarrow
z=\frac{1}{2}\pm i\frac{\sqrt{3}}{2}\Leftrightarrow z=-j\;\mbox{ou}\;z=-j^2.$$
\fincorrection
\correction{005128}
Soient $x\in\Rr$ et $z=\frac{1+ix}{1-ix}$. Puisque $1-ix\neq0$, $z$ est bien défini et
$|z|=\frac{|1+ix|}{|1-ix|}=\frac{|1+ix|}{|\overline{1+ix}|}=1$. Enfin,
$z=\frac{-1+ix+2}{1-ix}=-1+\frac{2}{1-ix}\neq-1$. On a montré que~:
$$\forall x\in\Rr,\;\frac{1+ix}{1-ix}\in U\setminus\{-1\}.$$
Réciproquement, soit $z\in U\setminus\{-1\}$. Il existe un réel $\theta\notin\pi+2\pi\Zz$ tel que $z=e^{i\theta}$. Mais
alors,

$$z=e^{i\theta}=\frac{e^{i\theta/2}}{e^{-i\theta/2}}=\frac{\cos\frac{\theta}{2}+i\sin\frac{\theta}{2}}
{\cos\frac{\theta}{2}-i\sin\frac{\theta}{2}}=\frac{\cos\frac{\theta}{2}(1+i\tan\frac{\theta}{2})}
{\cos\frac{\theta}{2}(1-i\tan\frac{\theta}{2})}=\frac{1+i\tan\frac{\theta}{2}}
{1-i\tan\frac{\theta}{2}}\;(\cos\frac{\theta}{2}\neq0\;\mbox{car}\;\frac{\theta}{2}\notin\frac{\pi}{2}+\pi\Zz),
$$
et $z$ est bien sous la forme voulue avec $x=\tan\frac{\theta}{2}$.
\fincorrection
\correction{005129}
\begin{enumerate}
 \item Soit $\theta\in\Rr$.

$$1-\cos\theta+i\sin\theta=0\Leftrightarrow\cos\theta=1\;\mbox{et}\;\sin\theta=0\Leftrightarrow\theta\in0+2\pi\Zz.$$
Donc, $\frac{1+\cos\theta-i\sin\theta}{1-\cos\theta+i\sin\theta}$ existe pour $\theta\notin 2\pi\Zz$. Pour un tel
$\theta$,

$$\frac{1+\cos\theta-i\sin\theta}{1-\cos\theta+i\sin\theta}=\frac{2\cos^2\frac{\theta}{2}-2i\sin\frac{\theta}{2}\cos
\frac{\theta}{2}}{2\sin^2\frac{\theta}{2}+2i\sin\frac{\theta}{2}\cos\frac{\theta}{2}}
=\frac{\cos\frac{\theta}{2}}{\sin\frac{\theta}{2}}\frac{\cos(\theta/2)-i\sin(\theta/2)}{\sin(\theta/2)+i\cos(\theta/2)}
=\frac{\cos\frac{\theta}{2}}{\sin\frac{\theta}{2}}\frac{e^{-i\theta/2}}{e^{i(\pi-\theta)/2}}=-i\cotan\left(\frac{\theta}{2}\right).$$
\textbf{- 1er cas.}
$\cotan\frac{\theta}{2}>0\Leftrightarrow\frac{\theta}{2}\in\displaystyle\bigcup_{k\in\Zz}]k\pi,\frac{\pi}{2}+k\pi[\Leftrightarrow\theta\in
\displaystyle\bigcup_{k\in\Zz}]2k\pi,\pi+2k\pi[$.
Dans ce cas, la forme trigonométrique de $\frac{1+\cos\theta-i\sin\theta}{1-\cos\theta+i\sin\theta}$ est 
$\cotan(\frac{\theta}{2})e^{-i\pi/2}$ (module$=\cotan(\frac{\theta}{2})$ et argument$=-\frac{\pi}{2}\;(2\pi)$).

$$\frac{1+\cos\theta-i\sin\theta}{1-\cos\theta+i\sin\theta}=\left[\cotan\left(\frac{\theta}{2}\right),-\frac{\pi}{2}\right].$$
\textbf{- 2ème cas.} $\cotan\frac{\theta}{2}<0\Leftrightarrow\theta\in
\displaystyle\bigcup_{k\in\Zz}]\pi+2k\pi,2(k+1)\pi[$.
Dans ce cas,

$$\frac{1+\cos\theta-i\sin\theta}{1-\cos\theta+i\sin\theta}=-\cotan(\frac{\theta}{2}).e^{i\pi/2}
=|\cotan(\frac{\theta}{2})|e^{i\pi/2},$$
et donc,

$$\frac{1+\cos\theta-i\sin\theta}{1-\cos\theta+i\sin\theta}=\left[-\cotan\left(\frac{\theta}{2}\right),\frac{\pi}{2}\right].$$
\textbf{- 3ème cas.} $\cotan\frac{\theta}{2}=0\Leftrightarrow\theta\in\pi+2\pi\Zz$. Dans ce cas, on
a $\frac{1+\cos\theta-i\sin\theta}{1-\cos\theta+i\sin\theta}=0$.
\item Pour $\theta\notin2\pi\Zz$, on a

$$\frac{1+e^{i\theta}}{1-e^{i\theta}}=\frac{e^{i\theta/2}(e^{-i\theta/2}+e^{i\theta/2})}{e^{i\theta/2}(e^{-i\theta/2}-
e^{i\theta/2})}=\frac{2\cos\frac{\theta}{2}}{-2i\sin\frac{\theta}{2}}=i\cotan\frac{\theta}{2}.$$
Si $\theta\in\displaystyle\bigcup_{k\in\Zz}]2k\pi,\pi+2k\pi[$,
$\frac{1+e^{i\theta}}{1-e^{i\theta}}=[\cotan\frac{\theta}{2},\frac{\pi}{2}]$.
Si $\theta\in\displaystyle\bigcup_{k\in\Zz}]\pi+2k\pi,2(k+1)\pi[$,
$\frac{1+e^{i\theta}}{1-e^{i\theta}}=[-\cotan\frac{\theta}{2},-\frac{\pi}{2}]$.
Si $\theta\in\pi+2\pi\Zz$, $\frac{1+e^{i\theta}}{1-e^{i\theta}}=0$.
\end{enumerate}

\fincorrection
\correction{005130}
$(1+i\sqrt{3})^9=(2e^{i\pi/3})^9=2^9e^{3i\pi}=-512$.

\begin{center}
\shadowbox{
\begin{tabular}{l}
La forme algébrique d'un complexe est particulièrement bien adaptée à l'addition.\\
La forme trigonométrique d'un complexe est particulièrement bien adaptée à la multiplication.
\end{tabular}
}
\end{center}
\fincorrection
\correction{007209}
Soit $z\in\C$. Si $\overline{z}=1-z^2$, alors en conjuguant une deuxième fois on obtient : 
\[ z=1-\overline{z}^2=1-(1-z^2)^2=2z^2-z^4.\]
Donc $z$ est une racine du polynôme $X^4-2X^2+X$. Ce polynôme a deux racines évidentes : $0$ et $1$, et deux autres racines $\frac{-1\pm\sqrt 5}{2}$. Les deux dernières sont les seules à vérifier l'équation d'origine.

Autre approche : on écrit l'inconnue $z\in\C$ sous forme cartésienne $z=a+ib$ avec $a$ et $b$ réels. On obtient le système $\begin{cases}a^2-b^2+a&=1\\2ab-b&=0\end{cases}$.

La seconde équation équivaut à $b=0$ ou $a=\frac{1}{2}$. Si $b=0$, la première équation devient $a^2+a-1=0$, dont les solutions (réelles) sont $\frac{-1\pm\sqrt 5}{2}$. Si $a=1/2$, la première équation devient $b^2=-\frac{1}{4}$, qui n'a pas de solutions réelles.
\fincorrection
\correction{007210}
Comme $z^2$ doit être réel, on en déduit que $z$ est soit réel, soit imaginaire pur. Dans chacun de ces deux cas, on obtient une équation d'inconnue réelle qui fait intervenir la valeur absolue. En séparant suivant le signe de l'inconnue, on obtient dinalement quatre cas qui correspondent à quatre trinômes réels.
\fincorrection
\nocorrection
\nocorrection
\nocorrection
\nocorrection
\correction{000027}
\textbf{Racines carr\'ees.} Soit $z= a+ib$ un nombre complexe avec
$a,b \in \Rr$ ; nous cherchons les complexes $\omega \in \Cc$ tels
que $\omega^2 = z$. \'Ecrivons $\omega = \alpha + i \beta$. Nous
raisonnons par \'equivalence :
\begin{align*}
\omega^2 = z    & \Leftrightarrow (\alpha +i \beta)^2 = a+ib \\
        & \Leftrightarrow \alpha^2-\beta^2+2i\alpha\beta = a +i b \\
\intertext{Soit en identifiant les parties r\'eelles entre elles ainsi que les parties imaginaires :}\\
        & \Leftrightarrow
              \begin{cases}
                 \alpha^2-\beta^2 = a \\
                 2\alpha\beta = b
              \end{cases} \\
\intertext{Sans changer l'\'equivalence nous rajoutons la
condition $|\omega|^2 = |z|$.} & \Leftrightarrow
  \begin{cases}
     \alpha^2 + \beta^2 = \sqrt{a^2+b^2} \\
     \alpha^2-\beta^2 = a \\
     2\alpha\beta = b
  \end{cases} \\
\intertext{Par somme et différence des deux premières lignes :}
& \Leftrightarrow
  \begin{cases}
     \alpha^2  =\frac{a+ \sqrt{a^2+b^2}}{2} \\
     \beta^2 =  \frac{-a+ \sqrt{a^2+b^2}}{2} \\
     2\alpha\beta = b
  \end{cases} \\
& \Leftrightarrow
  \begin{cases}
     \alpha  =\pm \sqrt{\frac{a+ \sqrt{a^2+b^2}}{2}} \\
     \beta = \pm \sqrt{ \frac{-a+ \sqrt{a^2+b^2}}{2}} \\
     \alpha\beta \text{ est du m\^eme signe que } b
  \end{cases} \\
\end{align*}
Cela donne deux couples $(\alpha,\beta)$ de solutions et donc deux
racines carr\'ees (opposées) $\omega = \alpha + i\beta$ de $z$.

En pratique on r\'ep\`ete facilement ce raisonnement, par exemple
pour  $z = 8-6i$,
\begin{align*}
\omega^2 = z
&\Leftrightarrow (\alpha +i \beta)^2 = 8-6i\\
&\Leftrightarrow \alpha^2-\beta^2+2i\alpha\beta = 8-6i\\
&\Leftrightarrow
  \begin{cases}
     \alpha^2-\beta^2 = 8\\
     2\alpha\beta = -6
  \end{cases} \\
&\Leftrightarrow
  \begin{cases}
     \alpha^2 + \beta^2 = \sqrt{8^2+(-6)^2} = 10 \text{ le module de $z$} \\
     \alpha^2-\beta^2 = 8 \\
     2\alpha\beta = -6
  \end{cases}\\
&\Leftrightarrow
  \begin{cases}
     2\alpha^2  =18 \\
     \beta^2 = 1 \\
     2\alpha\beta = -6
  \end{cases}\\
&\Leftrightarrow
  \begin{cases}
     \alpha  =\pm \sqrt{9} = \pm 3 \\
     \beta = \pm 1 \\
     \alpha \text{ et } \beta \text{ de signes oppos\'es}
  \end{cases} \\
&\Leftrightarrow
  \begin{cases}
     \ \ & \alpha  = 3 \text{ et } \beta = - 1\\
     \text{ou}&\\
    \ \ &\alpha  = -3 \text{ et } \beta = +1 \\
\end{cases} \\
\end{align*}

Les racines de $ z = 8-6i$ sont donc $\omega_1 = 3-i$ et $\omega_2 = -\omega_1 =
-3+i$.

Pour les autres : 
\begin{itemize}
 \item Les racines carrées de $1$ sont : $+1$ et $-1$.
 \item  Les racines carrées de $i$ sont : $\frac{\sqrt 2}{2}(1+i)$ et $-\frac{\sqrt 2}{2}(1+i)$.
 \item  Les racines carrées de $3+4i$ sont : $2+i$ et $-2-i$. 
 \item  Les racines carrées de $7+24i$ sont : $4+3i$ et $-4-3i$. 
\end{itemize}

\fincorrection
\correction{000028}
$2-i$ et $-2+i$; \qquad
$5-i$ et $-5+i$.
\fincorrection
\correction{000029}
Par la m\'ethode usuelle nous calculons les racines carr\'ees
$\omega, -\omega$ de $z = \frac{1+i}{\sqrt{2}}$, nous obtenons
$$\omega = \sqrt{\frac{\sqrt{2}+1}{2\sqrt{2}}}+i\sqrt{\frac{\sqrt{2}-1}{2\sqrt{2}}},$$
qui peut aussi s'écrire :
$$\omega = \frac 12 \sqrt{2+\sqrt 2} +i \frac 12 \sqrt{2-\sqrt 2}.$$

Mais nous remarquons que $z$ s'\'ecrit \'egalement
$$z = e^{i\frac{\pi}{4}}$$
et  $e^{i\frac{\pi}{8}}$ v\'erifie
$$\left (e^{i\frac{\pi}{8}}\right)^2 = e^{\frac{2i\pi}{8}}
= e^{i\frac{\pi}{4}}.$$
 Cela signifie que $e^{i\frac{\pi}{8}}$ est
une racine carr\'ee de $z$, donc $ e^{i\frac{\pi}{8}} = \cos
\frac{\pi}{8}+i\sin\frac{\pi}{8}$
  est \'egal \`a  $\omega$ ou $-\omega$. Comme $ \cos \frac{\pi}{8} > 0$
alors $ e^{i\frac{\pi}{8}} =  \omega$ et donc par identification
des parties r\'eelles et imaginaires :
$$\cos \frac{\pi}{8} = \frac 12 \sqrt{2+\sqrt 2}
\quad \text{ et }\quad \sin \frac{\pi}{8} =
\frac 12 \sqrt{2-\sqrt 2}.$$
\fincorrection
\correction{000030}
Soit $P(z) = az^2+bz+c$, et $\Delta = b^2-4ac$, si $\Delta \geq 0$
alors les racines sont r\'eelles, seul le cas o\`u $\Delta < 0$
nous int\'eresse. Premi\`ere m\'ethode : il suffit de regarder les
deux solutions et de v\'erifier qu'elles sont conjugu\'ees...

Seconde m\'ethode : si $z$ est une racine de $P$ \emph{i.e.} $P(z)
= 0$, alors
$$ P(\overline{z}) = a{\overline{z}}^2+b\overline{z}+c =
\overline{az^2+b^z+c} = \overline{P(z)} = 0.$$ Donc $\overline{z}$
est aussi une racine de $P$. Or $z$ n'est pas un nombre r\'eel
(car $\Delta < 0$ ) donc $\overline{z} \not= z$. Sachant que le
polyn\^ome $P$ de degr\'e $2$ a exactement $2$ racines, ce sont
$z$ et $\overline{z}$ et elles sont conjugu\'ees.
\fincorrection
\correction{000031}
\textbf{\'Equations du second degr\'e.} La m\'ethode g\'enerale
pour r\'esoudre les \'equations du second degr\'e $az^2+bz+c= 0$
(avec $a,b,c \in \Cc$ et $a\not=0$) est la suivante : soit $\Delta
= b^2-4ac$ le discriminant complexe et $\delta$ une racine
carr\'ee de $\Delta$ ($\delta^2 = \Delta$) alors les solutions
sont :
$$z_1 = \frac{-b+\delta}{2a} \quad \text{ et } \quad z_2 = \frac{-b-\delta}{2a}.$$
Dans le cas o\`u les coefficients sont r\'eels, on retrouve la
m\'ethode bien connue. Le seul travail dans le cas complexe est de
calculer une racine $\delta$ de $\Delta$.

Exemple : pour $z^2-\sqrt{3}z-i =0$, $\Delta = 3+4i$, dont une
racine carr\'ee est $\delta = 2+i$,  les solutions sont donc :
$$z_1 = \frac{\sqrt{3}+2+i}{2}\quad  \text{ et }\quad  z_2 = \frac{\sqrt{3}-2-i}{2}.$$

Les solutions des autres équations sont :
\begin{itemize}
 \item L'équation $z^2+z+1=0$ a pour solutions : $\frac12 (-1+i\sqrt{3})$, $\frac12 (-1-i\sqrt{3})$.
 \item L'équation $z^2-(1+2i)z+i-1=0$ a pour solutions : $1+i$, $i$.
 \item L'équation $z^2-\sqrt{3}z-i=0$ a pour solutions : $\frac12(2-\sqrt{3}+i)$,  $\frac12(-2-\sqrt{3}-i)$
 \item L'équation $z^2-(5-14i)z-2(5i+12)=0$ a pour solutions : $5-12i$, $-2i$.
 \item L'équation $z^2-(3+4i)z-1+5i =0$ a pour solutions : $2+3i$, $1+i$.
 \item L'équation $4z^2-2z+1=0$ a pour solutions : $\frac 14(1+i\sqrt{3})$, $\frac 14(1-i\sqrt{3})$.
 \item L'équation $z^4+10z^2 +169=0$ a pour solutions : $2+3i$, $-2-3i$, $2-3i$, $-2+3i$.
 \item L'équation $z^4+2z^2 +4=0$ a pour solutions : $\frac{\sqrt{2}}{2}(1+i\sqrt{3})$, $\frac{\sqrt{2}}{2}(1-i\sqrt{3})$,  $\frac{\sqrt{2}}{2}(-1+i\sqrt{3})$, $\frac{\sqrt{2}}{2}(-1-i\sqrt{3})$.
\end{itemize}

\fincorrection
\nocorrection
\nocorrection
\nocorrection
\nocorrection
\correction{000036}
\begin{enumerate}
\item $\Delta =-2i$ dont les racines carr\'ees sont $1-i$ et $-1+i$, d'o\`u les
racines $z_1=5-2i$ et $z_2=6-3i$.
\item Une racine ``\'evidente'' $z_1=i$, d'o\`u la r\'esolution compl\`ete en effectuant
la division par $z-i$. On trouve $z_2=i$
et $z_3=-2i$.
\end{enumerate}
\fincorrection
\nocorrection
\nocorrection
\correction{002945}
Cercle circonscrit $ \Rightarrow $ ssi $|z| = 1$.
\fincorrection
\correction{002946}
$z_1 = -z_2 = 3-2i$, $z_3 = -z_4 = 1-i$.
\fincorrection
\correction{002947}
$z=1\pm2i$, $z=-4\pm2i$.
\fincorrection
\correction{002948}
$m=2i$.
\fincorrection
\correction{002949}
\begin{enumerate}
  \item 
  \item $\alpha = 0$ ou $\beta = t\alpha^2,\ t \ge \frac 14$.
\end{enumerate}
\fincorrection
\correction{002950}
\begin{enumerate}
  \item 
  \item {\'E}l{\'e}ver au carr{\'e} :
             $|\alpha|^2 + |\beta|^2 + 2\underbrace{|\alpha\beta|}_{|\mu|^2} =
             \underbrace{|m-\mu|^2 + |m+\mu|^2}_{2|m|^2 + 2|\mu|^2} +
             2\underbrace{|m^2-\mu^2|}_{|\alpha-\beta|^2/4}$.
\end{enumerate}
\fincorrection
\correction{005120}
\begin{enumerate}
 \item  $z^2+z+1=0\Leftrightarrow z=-\frac{1}{2}+i\frac{\sqrt{3}}{2}=j\;\mbox{ou}\;z=-\frac{1}{2}-i\frac{\sqrt{3}}{2}=j^2$.
 \item  $\Delta'=1^2-2=-1=i^2$. L'équation a donc deux solutions non réelles et conjuguées, à savoir
$z_1=\frac{1}{2}(-1+i)$ et $z_2=\frac{1}{2}(-1-i)$.\rule[-5mm]{0mm}{0mm}
 \item  Soit $\theta\in\Rr$. Pour tout complexe $z$, on a
\begin{align*}
z^2-2z\cos\theta+1&=(z-\cos\theta)^2+1-\cos^2\theta=(z-\cos\theta)^2+\sin^2\theta=(z-\cos\theta)^2-(i\sin\theta)^2\\
 &=(z-\cos\theta-i\sin\theta)(z-\cos\theta+i\sin\theta)=(z-e^{i\theta})(z-e^{-i\theta})
\end{align*}
L'équation proposée a donc deux solutions (pas nécessairement distinctes) $z_1=e^{i\theta}$ et $z_2=e^{-i\theta}$. De
plus,
$\Delta'=\cos^2\theta-1=-\sin^2\theta$ et ces solutions sont distinctes si et seulement si
$\theta\notin\pi\Zz$.
 \item  Soit $(E)$ l'équation $z^2-(6+i)z+(11+3i)=0$.
Son discriminant est
$\Delta=(6+i)^2-4(11+13i)=-9-40i$. Comme $40=2\times20=2\times(4\times5)$ et que $4^2-5^2=16-25=-9$, on est en droit de deviner que
$\Delta=(4-5i)^2$. L'équation $(E)$ a deux solutions distinctes dans $\Cc$ à savoir $z_1=\frac{6+i+4-5i}{2}=5-2i$ et
$z_2=\frac{6+i-4+5i}{2}=1+3i$.\rule[-5mm]{0mm}{0mm}
 \item  Soit (E) l'équation $2z^2-(7+3i)z+(2+4i)=0$.
Son discriminant est $\Delta=(7+3i)^2-8(2+4i)=24+10i$. Comme $10
=2\times5=2\times(5\times1)$ et que $5^2-1^2=24$, on est en droit de deviner que $\Delta=(5+i)^2$. L'équation proposée a deux solutions
distinctes dans $\Cc$ à savoir $z_1=\frac{7+3i+5+i}{4}=3+i$ et $z_2=\frac{7+3i-5-i}{4}=\frac{1}{2}(1+i)$.
\end{enumerate}

\fincorrection
\correction{005125}
 Le discriminant de l'équation $Z^2-(5-14i)Z-2(5i+12)=0$ vaut

\begin{center}
$\Delta=(5-14i)^2+8(5i+12)=-75-100i=25(-3-4i)=(5(1-2i))^2$.
\end{center}
Cette équation admet donc les deux solutions
$Z_1=\frac{5-14i+5-10i}{2}=5-12i$ et $Z_2=\frac{5-14i-5+10i}{2}=-2i$.
Ensuite,

\begin{align*}
z\;\mbox{est solution de l'équation proposée}&\Leftrightarrow z^2=5-12i=(3-2i)^2\;\mbox{ou}\;z^2=-2i=(1-i)^2\\
 &\Leftrightarrow z = 3-2i\;\mbox{ou}\;z=-3+2i\;\mbox{ou}\;z=1-i\;\mbox{ou}\;z=-1+i.
\end{align*}
\fincorrection
\nocorrection
\nocorrection
\nocorrection
\nocorrection
\nocorrection
\correction{000042}
$\frac 14(-1+i) = \frac{1}{(\sqrt 2)^3}e^{\frac{3i\pi}{4}}= (\frac{1}{\sqrt 2}e^{\frac{i\pi}{4}})^3$.
Les solutions sont les complexes $z_k=\frac{1}{\sqrt 2}e^{\frac{i\pi}{4}+ \frac{2i k\pi}{3}}$
pour $0\leq k\leq 2$. Et seul $z_0 = \frac{1}{2}(1+i)$ a une puissance quatri\`eme r\'eelle.
\fincorrection
\correction{000043}
\begin{enumerate}
\item Les trois racines cubiques ont m\^eme module $\sqrt2$, et leurs arguments
sont $-\pi /12$, $7\pi /12$ et $5\pi/4$. Des valeurs approch\'ees sont
$1{,}36603-0{,}36603 i$, $-0{,}36603+1{,}36603 i$ et $-1-i$.
\item $-1-2i$, $(-1-2i)j$ et $(-1-2i)j^2$ o\`u $j={-1+i\sqrt3\over2}$ (racine
cubique de 1).
\end{enumerate}
\fincorrection
\correction{000044}
$\cos{\pi \over 12}={1+\sqrt3\over 2\sqrt2}$ ;
$\sin{\pi \over 12}={-1+\sqrt3\over 2\sqrt2}$ ;
$\tan{\pi \over 12}=2-\sqrt3$ ;
$\tan{5\pi \over12}=2+\sqrt3$.



Les racines de $z^{24}=1$ sont donn\'ees par $z_k=e^{2ki\pi /24}$ pour
$k=0,1,\ldots,23$. Ce sont donc $1$, $\cos{\pi \over 12}+i\sin{\pi
\over 12}$, etc.
\fincorrection
\correction{000045}
\begin{enumerate}
\item $3$, $3i$, $-3$ et $-3i$.
\item ${3\sqrt2\over 2}(1+i)$, ${3\sqrt2\over 2}(-1+i)$, ${3\sqrt2\over 2}(-1-i)$ 
et ${3\sqrt2\over 2}(1-i)$.
\end{enumerate}
\fincorrection
\correction{000046}
Pour 2. Utiliser la formule d'Euler pour $\sin\left(x/2\right)$.

Pour 3.
Pour $x\neq
2k\pi$, $k\in\mathbb{Z}$,
$$Z_n=\frac{\sin\left(\frac{nx}{2}\right)}{\sin\left(\frac{x}{2}\right)}\exp\left(i\left(n-1\right)\frac{x}{2}\right),$$
et pour $x=2k\pi$, $k\in\mathbb{Z}$, $Z_n=n$.

Remarquer que $Z_n=X_n+iY_n$ pour en déduire que
$$X_n=\frac{\cos\left(\frac{(n-1)x}{2}\right)\sin\left(\frac{nx}{2}\right)}{\sin\left(\frac{x}{2}\right)}\;\mbox{
et }\;Y_n=\frac{\sin\left(\frac{(n-1)x}{2}\right)\sin\left(\frac{nx}{2}\right)}{\sin\left(\frac{x}{2}\right)}.$$
\fincorrection
\correction{000047}
$$S_n = 1+z+z^2+\cdots+z^n = \sum_{k=0}^{n}z^k.$$
Nous devons retrouver le r\'esultat sur la somme $S_n =
\frac{1-z^{n+1}}{1-z}$d'une suite g\'eom\'etrique dans le cas o\`u
$z\not=1$ est un r\'eel. Soit maintenant $z \not= 1$ un nombre
complexe. Calculons $S_n(1-z)$.
\begin{align*}
S_n(1-z) & =(1+z+z^2+\cdots+z^n)(1-z) \text{ d\'eveloppons }\\
         &= 1+z+z^2+\cdots+z^n - z-z^2-\cdots-z^{n+1} \text{ les termes interm\'ediaires s'annulent }\\
         &= 1-z^{n+1}.
\end{align*}
Donc $$S_n = \frac{1-z^{n+1}}{1-z}, \text{ pour } z\not= 1.$$
\fincorrection
\correction{000048}
\textbf{Calcul de racine $n$-i\`eme.} Soit $z\in\Cc$ tel que
$z^n=1$, d\'ej\`a $|z|^n=1$ et donc $|z|=1$. \'Ecrivons $z =
e^{i\theta}$. L'\'equation devient
$$e^{in\theta} = e^{0} =1
\Leftrightarrow n\theta = 0 + 2k\pi, \ k\in \Zz \Leftrightarrow \theta =
\frac{2k\pi}{n}, \ k\in \Zz.$$ Les solution sont donc
$$\mathcal{S} = \left\lbrace e^{\frac{2ik\pi}{n}}, \ k\in \Zz\right\rbrace.$$
Comme le polyn\^ome $z^n-1$ est de degr\'e $n$ il a au plus $n$
racines. Nous choisissons pour repr\'esentants :
$$\mathcal{S} = \left\lbrace e^{\frac{2ik\pi}{n}}, \ k= 0,\ldots,n-1\right\rbrace.$$
De plus si $\epsilon = e^{\frac{2i\pi}{n}}$ alors $\mathcal{S} =
\left\lbrace \epsilon^k, \ k= 0,\ldots,n-1\right\rbrace.$ Ces
racines sont les sommets d'un polygone r\'egulier \`a $n$
c\^ot\'es inscrit dans le cercle unit\'e.



Soit $P(z) = \sum_{k=0}^{n-1}z^k=\frac{1-z^{n}}{1-z}$ pour
$z\not=1$. Donc quelque soit $z\in\mathcal{S}\setminus\{ 1 \}$
$P(z) = 0$, nous avons ainsi trouver $n-1$ racines pour $P$ de
degr\'e $n-1$, donc l'ensemble des racines de $P$ est exactement
$\mathcal{S}\setminus\{ 1 \}$.




Pour conclure soit  $Q_p(z) = \sum_{k=0}^{n-1}\epsilon^{kp}$.\\
Si $p = 0 +\ell n$, $\ell \in \Zz$ alors $\epsilon^{kp}=\epsilon^{k\ell n}=(\epsilon^n)^{k \ell}
=1^{k \ell}=1$. Donc $Q_p(z)  = \sum_{k=0}^{n-1} 1 = n$. \\
Sinon $Q_p(z)$ est la somme d'une suite g\'eom\'etrique de raison
$\epsilon^p$ :
$$ Q_p(z) = \frac{1-\left( \epsilon^p \right)^n}{1-\epsilon^p}
= \frac{1-\left( \epsilon^n \right)^p}{1-\epsilon^p} =
\frac{1-1}{1-\epsilon^p}=0.$$
\fincorrection
\nocorrection
\nocorrection
\nocorrection
\nocorrection
\nocorrection
\nocorrection
\nocorrection
\correction{000056}
Soient $z_{1},z_{2},z_{3}$ trois nombres complexes
\emph{distincts} ayant le m\^eme cube.
\begin{enumerate}
\item
$z_{1}\neq0$ car sinon on aurait $z_{1}=z_{2}=z_{3}=0$. Ainsi
$(\frac{z_{2}}{z_{1}})^3=(\frac{z_{3}}{z_{1}})^3=1$. Comme les
trois nombres $1,(\frac{z_{2}}{z_{1}})$ et $(\frac{z_{3}}{z_{1}})$
sont distincts on en d\'eduit que ce sont les trois racines
cubiques de 1. Ces racines sont $1, j=e^{\frac{2i\pi}{3}}$ et $
j^2=e^{-\frac{2i\pi}{3}}$. A une permutation pr\`es des indices 2
et 3 on a donc :
$$
  z_{2}=jz_{1} \qquad\text{ et }\qquad z_{3}=j^2z_{1}.
$$
\item Soit $z\in\C$. On a les \'equivalences suivantes :
\begin{align*}
z^6+(7-i)z^3-8-8i=0 &\Leftrightarrow z^3\text{ est solution de }Z^2+(7-i)Z-8-8i=0\\
\intertext{ Etudions l'\'equation $Z^2+(7-i)Z-8-8i=0$.
$\Delta=(7-i)^2+4(8+8i)=80+18i=(9+i)^2$. Les solutions sont donc
$-8$ et $1+i$. Nous pouvons reprendre notre suite d'\'equivalences
: }
z^6+(7-i)z^3-8-8i=0 &\Leftrightarrow z^3\in\{-8,1+i\} \\
  &\Leftrightarrow z^3=(-2)^3 \quad\text{ ou }\quad z^3=(\sqrt[6]{2}e^{i\frac{\pi}{12}})^3\\
  &\Leftrightarrow z\in\{-2,
              -2e^{\frac{2i\pi}{3}},
              -2e^{-\frac{2i\pi}{3}}\}
         \text{ ou }
        z\in\{\sqrt[6]{2}e^{i\frac{\pi}{12}},
              \sqrt[6]{2}e^{i\frac{9\pi}{12}},
              \sqrt[6]{2}e^{i\frac{17\pi}{12}}\}\\
  &\Leftrightarrow z\in\{-2,
               2e^{\frac{i\pi}{3}},
               2e^{-\frac{i\pi}{3}},
              \sqrt[6]{2}e^{i\frac{\pi}{12}},
              \sqrt[6]{2}e^{i\frac{3\pi}{4}},
              \sqrt[6]{2}e^{i\frac{17\pi}{12}}\}.
\end{align*}
L'ensemble des solutions est donc :
           $$\{-2,
               2e^{\frac{i\pi}{3}},
               2e^{-\frac{i\pi}{3}},
              \sqrt[6]{2}e^{i\frac{\pi}{12}},
              \sqrt[6]{2}e^{i\frac{3\pi}{4}},
              \sqrt[6]{2}e^{i\frac{17\pi}{12}}\}.$$
\end{enumerate}
\fincorrection
\nocorrection
\nocorrection
\nocorrection
\correction{002939}
\begin{enumerate}
  \item $z=-i\mathrm{cotan}\frac {k\pi}n$.
  \item $6\mid n \Rightarrow  z = j$ ou $j^2$. Sinon, pas de solution.
  \item $z = \exp\frac{(2k+1)i\pi}5$, $k = 0,1,3,4$.
  \item $z = -1$ ou $z = \exp\frac{2ik\pi}n$, $1\le k< n$.
  \item $x = \tan\left(\frac {a+2k\pi}n\right)$.
  \item
  \item $z = \pm i,\ \pm i(2\pm\sqrt3)$.
\end{enumerate}
\fincorrection
\correction{002940}
\begin{enumerate}
  \item D{\'e}velopper. $S=2n$.
  \item $\frac {1 - (1+\omega)^n}{1 - \omega - \omega^2} =
              \frac {1 + (2\cos(\pi/n))^n}{1 - \omega - \omega^2}$.
\end{enumerate}
\fincorrection
\correction{002941}
\begin{enumerate}
  \item $\sum=n$ si $p\not\equiv 0 (\mathrm{mod}\, n)$, 0 sinon.
  \item $a_k = \sum_{x\in\mathbb{U}_n} \frac {P(x)}{nx^k}$
\end{enumerate}
\fincorrection
\correction{002942}
$n$ impair $ \Rightarrow  |Z|^2 = n$,\par
         $n$ pair $ \Rightarrow  |Z|^2 = n(1+(-1)^{n/2})$.
\fincorrection
\correction{002943}
\begin{enumerate}
  \item $u+v = -1$, $u^2 = u+2v = -2-u$.
  \item $\Sigma = \Im(u) = \frac {\sqrt7}2$.
\end{enumerate}
\fincorrection
\correction{002944}
$x = \frac {2(n^2+n+1)}{3n(n+1)}$.
\fincorrection
\correction{005122}
Soit $\alpha\in\left]\frac{\pi}{2},\frac{\pi}{2}\right[$.
$\frac{1+i\tan\alpha}{1-i\tan\alpha}=\frac{\cos\alpha+i\sin\alpha}{\cos\alpha-i\sin\alpha}=e^{2i\alpha}$. Donc,

$$\left(\frac{1+iz}{1-iz}\right)^3=\frac{1+i\tan\alpha}{1-i\tan\alpha}\Leftrightarrow\exists
k\in\{-1,0,1\}/\;\frac{1+iz}{1-iz}=e^{i(\frac{2\alpha}{3}+\frac{2k\pi}{3})}=\omega_k\Leftrightarrow\exists
k\in\{-1,0,1\}/\;i(\omega_k+1)z=\omega_k-1.$$
Maintenant, pour $k\in\{-1,0,1\}$,

$$\omega_k=-1\Leftrightarrow\frac{2\alpha}{3}+\frac{2k\pi}{3}\in\pi+2\pi\Zz\Leftrightarrow\alpha\in-k\pi+\frac{3\pi}{2}+3\pi\Zz,$$
ce qui est exclu pour $\alpha\in]-\frac{\pi}{2},\frac{\pi}{2}[$. Donc,

\begin{align*}
\left(\frac{1+iz}{1-iz}\right)^3=\frac{1+i\tan\alpha}{1-i\tan\alpha}&\Leftrightarrow\exists
k\in\{-1,0,1\}/\;z=\frac{\omega_k-1}{i(\omega_k+1)}\Leftrightarrow\exists
k\in\{-1,0,1\}/\;z=\frac{e^{i(\frac{\alpha}{3}+\frac{k\pi}{3})}}{e^{i(\frac{\alpha}{3}+\frac{k\pi}{3})}}
\frac{e^{i(\frac{\alpha}{3}+\frac{k\pi}{3})}-e^{-i(\frac{\alpha}{3}+\frac{k\pi}{3})}}{i(e^{i(\frac{\alpha}{3}
+\frac{k\pi}{3})}+e^{-i(\frac{\alpha}{3}+\frac{k\pi}{3})})}\\
 &\Leftrightarrow\exists
k\in\{-1,0,1\}/\;z=\frac{2i\sin(\frac{\alpha}{3}+\frac{k\pi}{3})}{i(2\cos(\frac{\alpha}{3}+\frac{k\pi}{3}))}\Leftrightarrow
\exists k\in\{-1,0,1\}/\;z=\tan(\frac{\alpha}{3}+\frac{k\pi}{3})
\end{align*}
\fincorrection
\correction{005126}
Posons, pour $n$ naturel non nul, $P=(X^2+1)^n-(X-1)^{2n}$.

\begin{align*}
P&=X^{2n}+\left(\mbox{termes de degré}\leq2n-2\right)-X^{2n}+2nX^{2n-1}+\left(\mbox{termes de degré}\;\leq2n-2\right)\\
 &=2nX^{2n-1}+\left(\mbox{termes de degré}\leq2n-2\right).
\end{align*}
Donc $\mbox{deg}(P)=2n-1$ et $P$ admet dans $\Cc$, $2n-1$ racines, distinctes ou confondues.

\begin{align*}
(z^2+1)^n=(z-1)^{2n}&\Leftrightarrow\exists k\in\{0,...,n-1\}/\;z^2+1=\omega_k(z-1)^2\;\mbox{où}\;\omega_k=e^{2ik\pi/n}\\
 &\Leftrightarrow\exists k\in\{0,...,n-1\}/\;(1-\omega_k)z^2+2\omega_kz+(1-\omega_k)=0
\end{align*}
Si $k=0$, l'équation précédente s'écrit $2z=0$ ou encore $z=0$.
Si $k$ est élément de $\llbracket1,n-1\rrbracket$, $\Delta_k'=\omega_k^2-(1-\omega_k)^2=2\omega_k-1=2e^{2ik\pi/n}-1$.
Soit $d_k$ une racine carrée dans $\Cc$ de $\Delta_k'$ (difficile à expliciter semble-t-il). On a
$S=\{0\}\cup\left\{\frac{-e^{2ik\pi/n}\pm d_k}{1-e^{2ik\pi/n}},k\in\llbracket1,n-1\rrbracket\right\}$.

\fincorrection
\correction{005131}
$i=e^{i\pi/2}$ et les racines quatrièmes de $i$ sont donc les $e^{i(\frac{\pi}{8}+\frac{k\pi}{2})}$,
$k\in\{0,1,2,3\}$.
Ensuite, $\frac{-4}{1+i\sqrt{3}}=\frac{-2}{e^{i\pi/3}}=-2e^{-i\pi/3}=2e^{2i\pi/3}$. Les racines
sixièmes de $\frac{-4}{1+i\sqrt{3}}$ sont donc les $\sqrt[6]{2}e^{i(\frac{\pi}{9}+\frac{k\pi}{3})}$,
$k\in\{0,1,2,3,4,5\}$.
\fincorrection
\correction{005135}
\begin{enumerate}
 \item  Soit $z\in\Cc$. Soient $M$, $A$ et $B$ les points d'affixes resectives $z$, $1$ et $-1$.

\begin{align*}
z\;\mbox{solution
de}\;(E)&\Rightarrow(z-1)^n=(z+1)^n\Rightarrow|(z-1)^n|=|(z+1)^n|\Rightarrow|z-1|^n=|z+1|^n
\Rightarrow|z-1|=|z+1|\\
 &\Rightarrow AM=BM\Rightarrow M\in\mbox{med}[AB]\Rightarrow M\in(Oy)\Rightarrow z\in i\Rr.
\end{align*}
 \item  Soit $z\in\Cc$.

$$(-z-1)^n-(-z+1)^n=(-1)^n((z+1)^n-(z-1)^n)=-(-1)^n((z-1)^n-(z+1)^n).$$
Par suite,

$$z\;\mbox{solution de}\;(E)\Leftrightarrow(z-1)^n-(z+1)^n=0\Leftrightarrow(-z-1)^n-(-z+1)^n=0\Leftrightarrow-z\;\mbox{solution de}\;(E).$$

\item  Soit $z\in\Cc$.
\begin{align*}
z\;\mbox{solution de}\;(E)&\Leftrightarrow(z-1)^n=(z+1)^n\Leftrightarrow\exists k\in\llbracket0,n-1\rrbracket/\;z-1=e^{2ik\pi/n}(z+1)\\
&\Leftrightarrow\exists k\in\llbracket1,n-1\rrbracket/\;z=-\frac{e^{2ik\pi/n}+1}{e^{2ik\pi/n}-1}
\Leftrightarrow\exists k\in\llbracket1,n-1\rrbracket/\;z=-\frac{e^{ik\pi/n}+e^{-ik\pi/n}}{e^{ik\pi/n}-e^{-ik\pi/n}}\\
&\Leftrightarrow\exists k\in\llbracket1,n-1\rrbracket/\;z=-\frac{2\cos\frac{k\pi}{n}}{2i\sin\frac{k\pi}{n}}
\Leftrightarrow\exists k\in\llbracket1,n-1\rrbracket/\;z=i\cotan\frac{k\pi}{n}
\end{align*}
\end{enumerate}
\fincorrection
\correction{005313}
\begin{enumerate}
\item  Soit $n\geq2$. On a 

$$a_n=\prod_{k=1}^{n-1}\frac{1}{2i}(e^{ik\pi/n}-e^{-ik\pi/n})=\frac{1}{(2i)^{n-1}}\prod_{k=1}^{n-1}e^{ik\pi/n}\prod_{k=1}^{n-1}(1-e^{-2ik\pi/n}).$$

Maintenant, 

$$\prod_{k=1}^{n-1}e^{ik\pi/n}=e^{\frac{i\pi}{n}(1+2+...+(n-1))}=e^{i\pi(n-1)/2}(e^{i\pi/2})^{n-1}=i^{n-1},$$

et donc $\frac{1}{(2i)^{n-1}}\prod_{k=1}^{n-1}e^{ik\pi/n}=\frac{1}{2^{n-1}}$.

Il reste à calculer $\prod_{k=1}^{n-1}(1-e^{-2ik\pi/n})$.

\begin{itemize}
\item[\textbf{1ère solution.}]
Les $e^{-2ik\pi/n}$, $1\leq k\leq n-1$, sont les $n-1$ racines $n$-ièmes de $1$ distinctes de $1$ et puisque 
$X^n-1=(X-1)(1+X+...+X^{n-1})$, ce sont donc les $n-1$ racines deux deux distinctes du polynôme $1+X+...+X^{n-1}$. Par suite, $1+X+...+X^{n-1}=\prod_{k=1}^{n-1}(X-e^{-2ik\pi/n})$, et en particulier $\prod_{k=1}^{n-1}(1-e^{-2ik\pi/n})=1+1...+1=n$.

\item[\textbf{2ème solution.}]
Pour $1\leq k\leq n-1$, posons $z_k=1-e^{-2ik\pi/n}$. Les $z_k$ sont deux à deux distincts et racines du polynôme 
$P=(1-X)^n-1=-X+...+(-1)^nX^n=X(-n+X-...+(-1)^nX^{n-1})$. Maintenant, $z_k=0\Leftrightarrow e^{-2ik\pi/n}=1\leq k\in n\Zz$ (ce qui n'est pas pour $1\leq k\leq n-1$). Donc, les $z_k$, $1\leq k\leq n-1$, sont $n-1$ racines deux à deux distinctes du polynôme de degré $n-1$~:~$-n+X-...+(-1)^nX^{n-1}$. Ce sont ainsi toutes les racines de ce polynôme ou encore

$$-n+X-...+(-1)^nX^{n-1}=(-1)^n\prod_{k=1}^{n-1}(X-z_k).$$

En particulier, en égalant les coefficients constants,

$$(-1)^n.(-1)^{n-1}\prod_{k=1}^{n-1}z_k=-n,$$

et donc encore une fois $\prod_{k=1}^{n-1}(1-e^{-2ik\pi/n})=n$.
\end{itemize}

Finalement,

$$\forall n\geq2,\;\prod_{k=1}^{n-1}\sin\frac{k\pi}{n}=\frac{n}{2^{n-1}}.$$

\item  Soit $n$ un entier naturel non nul.

$$b_n=\prod_{k=1}^{n}\frac{1}{2}(e^{i(a+\frac{k\pi}{n})}+e^{-i(a+\frac{k\pi}{n})})=\frac{1}{2^n}\prod_{k=1}^{n}e^{-i(a+\frac{k\pi}{n})}\prod_{k=1}^{n}(e^{2i(a+\frac{k\pi}{n})}+1).$$

Ensuite, 

$$\prod_{k=1}^{n}e^{-i(a+\frac{k\pi}{n})}=e^{-ina}e^{-\frac{i\pi}{n}(1+2+...+n)}=e^{-ina}e^{-i(n+1)\pi/2}.$$

D'autre part, soit $P=\prod_{k=1}^{n}(X+e^{2i(a+\frac{k\pi}{n})})=\prod_{k=1}^{n}(X-(-e^{2i(a+\frac{k\pi}{n})}))$.
Pour tout $k$, on a $(-e^{2i(a+\frac{k\pi}{n}})^n=(-1)^ne^{2ina}$. Par suite, les $n$ nombres deux à deux distincts $-e^{2i(a+\frac{k\pi}{n}}$, $1\leq k\leq n$ sont racines du polynôme $X^n-(-1)^ne^{2ina}$, de degré $n$. On en déduit que, $P=X^n-(-1)^ne^{2ina}$.

Par suite, $\prod_{k=1}^{n}(e^{2i(a+\frac{k\pi}{n})}+1)=P(1)=1-(-1)^ne^{2ina}=1-e^{2ina+n\pi}$, puis 

\begin{align*}\ensuremath
b_n&=\frac{1}{2^n}e^{-ina}e^{-i(n+1)\pi/2}(1-e^{2ina+n\pi})=\frac{1}{2^n}(e^{-i(na+(n+1)\frac{\pi}{2})}-e^{i(na+(n-1)\frac{\pi}{2})})\\
 &=\frac{1}{2^n}(e^{-i(na+(n+1)\frac{\pi}{2})}+e^{i(na+(n+1)\frac{\pi}{2})})=\frac{\cos(na+(n+1)\frac{\pi}{2})}{2^{n-1}}.
\end{align*}

\item  
\begin{align*}\ensuremath
c_n\;\mbox{est défini}\Leftrightarrow\forall k\in\{1,...,n\},\;a+\frac{k\pi}{n}\notin\frac{\pi}{2}+\pi\Zz\Leftrightarrow\forall k\in\Nn,\;a-\frac{k\pi}{n}+\frac{\pi}{2}+\pi\Zz\Leftrightarrow a\notin\frac{\pi}{2}+\frac{\pi}{n}\Zz
\end{align*}

Pour les $a$ tels que $c_n$ est défini, on a $c_n=\prod_{k=1}^{n}\frac{1}{i}\frac{e^{2i(a+k\pi/n)}-1}{e^{2i(a+k\pi/n)}+1}$.

Pour $1\leq k\leq n$, posons $\omega_k=e^{2i(a+k\pi/n)}$ puis $z_k=\frac{\omega_k-1}{\omega_k+1}$. On a donc $c_n=\frac{1}{i^n}\prod_{k=1}^{n}z_k$.

Puisque $z_k=\frac{\omega_k-1}{\omega_k+1}$, on a $\omega_k(1-z_k)=1+z_k$ et donc, pour $1\leq k\leq n$, $\omega_k^n(1-z_k)^n=(1+z_k)^n$ ou encore, les $z_k$ sont racines du polynôme $P=(1+X)^n-e^{2ina}(1-X)^n$. Maintenant, les $a+\frac{k\pi}{n}$ sont dans $[a,a+\pi[$ et donc deux à deux distincts puisque la fonction tangente est injective sur tout intervalle de cette forme.

\begin{itemize}
\item[1er cas.] Si $e^{2ina}\neq(-1)^n$ alors $P$ est de degré $n$ et $P=(1-(-1)^ne^{2ina})\prod_{k=1}^{n}(X-z_k)$. En évaluant en $0$, on obtient

$$(1-(-1)^ne^{2ina})\prod_{k=1}^{n}(-z_k)=1-e^{2ina}.$$

D'où,

$$\prod_{k=1}^{n}z_k=\frac{1-e^{2ina}}{(-1)^n-e^{2ina}}=\frac{1-e^{2ina}}{e^{in\pi}-e^{2ina}}=\frac{e^{ina}}{e^{in\pi/2}e^{ina}}\frac{-2i\sin(na)}{-2i\sin n(a-\frac{\pi}{2})}=\frac{1}{i^n} \frac{\sin(na)}{\sin n(a-\frac{\pi}{2})}.$$

Finalement, $c_n=(-1)^n\frac{\sin(na)}{\sin(n(a-\frac{\pi}{2}))}$.

Si $n$ est pair, posons $n=2p$, $p\in\Nn^*$. $c_n=c_{2p}=\frac{\sin(2pa)}{\sin (2pa-p\pi)}=(-1)^p$.

Si $n$ est impair, posons $n=2p+1$. $c_n=c_{2p+1}=(-1)^p\tan((2p+1)a)$.

\item[2ème cas.] Si $e^{2ina}=(-1)^n$, alors $2na\in n\pi+2\pi\Zz$ ou encore $a\in\frac{\pi}{2}+\pi\Zz$. Dans ce cas, $c_n$ n'est pas défini.
\end{itemize}
\end{enumerate}

\fincorrection
\nocorrection
\correction{000060}
Nous identifions $\Cc$ au plan affine et $z=x+iy$ \`a  $(x,y) \in
\Rr\times \Rr$.

Remarquons que pour les deux ensembles $z=5$ n'est pas solution,
donc
$$\left| \frac{z-3}{z-5} \right| = 1
\Leftrightarrow |z-3| = |z-5|.$$ Ce qui signifie pr\'eci\'sement que les
points d'affixe $z$ sont situ\'es \`a \'egale distance des points
$A,B$ d'affixes respectives $3 = (3,0)$ et $5=(5,0)$. L'ensemble
solution est la m\'ediatrice du segment $[A,B]$.

\bigskip

Ensuite pour
\begin{align*}
\left| \frac{z-3}{z-5} \right| = \frac{\sqrt{2}}{2}
&\Leftrightarrow |z-3|^2 = \frac{1}{2}|z-5|^2 \\
&\Leftrightarrow (z-3)\overline{(z-3)} = \frac{1}{2}(z-5)\overline{(z-5)}\\
&\Leftrightarrow z\overline{z}-(z+\overline{z})=7\\
&\Leftrightarrow |z-1|^2=8\\
&\Leftrightarrow |z-1|=2\sqrt{2}\\
\end{align*}
L'ensemble solution est donc le cercle de centre le point d'affixe
$1 = (1,0)$ et de rayon $2\sqrt{2}$.
\fincorrection
\nocorrection
\nocorrection
\nocorrection
\nocorrection
\correction{000065}
 En exprimant qu'un nombre complexe de module 1 peut s'\'ecrire $e^{i\theta
}$, on trouve $z={a-be^{i\theta }\over 1-e^{i\theta }}$. On peut encore \'ecrire
$z=A+B\cot{\theta \over 2}$, o\`u $A$ et $B$ sont ind\'ependants de $\theta $, ce qui
montre que le point d'affixe $z$ d\'ecrit une droite. G\'eom\'etriquement, cette droite
est bien entendu la m\'ediatrice du segment qui joint les points d'affixes $a$
et $b$.
\fincorrection
\correction{000066}
 M\'ethode analogue \`a celle de l'exercice \ref{exo:compl}. On trouve $z=
{a-bke^{i\theta }\over 1-ke^{i\theta }}$.
On peut v\'erifier que le point d'affixe $z$ d\'ecrit le cercle dont un diam\`etre
joint les points correspondant \`a $\theta =0$ et \`a $\theta =\pi $
(v\'erifier en cherchant le milieu $z_0$ de ce segment et en \'etudiant $\vert
z-z_0\vert$).
\fincorrection
\correction{000067}
\begin{enumerate}
\item R\'eciproque : $a+jb+j^2c=0$ ou $a+j^2b+jc=0$ (cela d\'epend de
l'orientation du triangle).
\item $ADOE$ est un parall\'elogramme. Les trois triangles $OBC$, $DBA$ et $EAC$
sont directement isom\'etriques, ce qui d'ailleurs se v\'erifie imm\'ediatement \`a
l'aide de rotations. 
\end{enumerate}
\fincorrection
\nocorrection
\correction{000069}
$$|u+v|^2+|u-v|^2=(u+v)(\bar u +\bar v) + (u-v)(\bar u -\bar v)
= 2u\bar u+2v\bar v = 2|u|^2+2|v|^2.$$
G\'eom\'etriquement il s'agit de l'identit\'e du parall\'elogramme.
Les points d'affixes $0,u,v,u+v$ forment un parall\'elogramme.
$|u|$ et $|v|$ sont les longueurs des cot\'es, et
$|u+v|, |u-v|$ sont les longueurs des diagonales.
Il n'est pas \'evident de montrer ceci sans les nombres complexes !!
\fincorrection
\nocorrection
\nocorrection
\nocorrection
\nocorrection
\nocorrection
\nocorrection
\nocorrection
\correction{000077}
\begin{enumerate}
\item
Comme $(A_{0},\ldots,A_{4})$ est un pentagone r\'egulier, on a
$OA_{0}=OA_{1}=OA_{2}=OA_{3}=OA_{4}=1$ et $
  (\overrightarrow{OA_{0}},\overrightarrow{OA_{1}})=\frac{2\pi}{5}[2\pi],
  (\overrightarrow{OA_{0}},\overrightarrow{OA_{2}})=\frac{4\pi}{5}[2\pi],
  (\overrightarrow{OA_{0}},\overrightarrow{OA_{3}})=-\frac{4\pi}{5}[2\pi],
  (\overrightarrow{OA_{0}},\overrightarrow{OA_{4}})=-\frac{2\pi}{5}[2\pi],
 $.
On en d\'eduit:
 $
  \omega_{0}=1,
  \omega_{1}=e^{\frac{2i\pi}{5}},
  \omega_{2}=e^{\frac{4i\pi}{5}},
  \omega_{3}=e^{-\frac{4i\pi}{5}}=e^{\frac{6i\pi}{5}},
  \omega_{4}=e^{-\frac{2i\pi}{5}}=e^{\frac{8i\pi}{5}},
 $.
On a bien $\omega_{i}=\omega_{1}^i$. Enfin, comme
$\omega_{1}\neq0$, $1+\omega_{1}+\ldots+\omega_{1}^4=
\frac{1-\omega_{1}^5}{1-\omega_{1}}=\frac{1-1}{1-\omega_{1}}=0$.

\item $\mathop{\mathrm{Re}}\nolimits(1+\omega_{1}+\ldots+\omega_{1}^4)=
1+2\cos(\frac{2\pi}{5})+2\cos(\frac{4\pi}{5})$. Comme
$\cos(\frac{4\pi}{5})=2\cos^2(\frac{2\pi}{5})-1$ on en d\'eduit:
$4\cos^2(\frac{2\pi}{5})+2\cos(\frac{2\pi}{5})-1=0$.
$\cos(\frac{2\pi}{5})$ est donc bien une solution de l'\'equation
$4z^2+2z-1=0$. Etudions cette \'equation: $\Delta=20=2^2.5$. Les
solutions sont donc $\frac{-1-\sqrt{5}}{4}$ et
$\frac{-1+\sqrt{5}}{4}$. Comme $\cos(\frac{2\pi}{5})>0$, on en
d\'eduit que $\cos(\frac{2\pi}{5})=\frac{\sqrt{5}-1}{4}$.

\item
 $
  BA_{2}^2=|\omega_{2}+1|^2
          =|\cos(\frac{4\pi}{5})+i\sin(\frac{4\pi}{5})+1|^2
          =1+2\cos(\frac{4\pi}{5})+\cos^2(\frac{4\pi}{5})+\sin^2(\frac{4\pi}{5})
          =4\cos^2(\frac{2\pi}{5})
  $. Donc $BA_{2}=\frac{\sqrt{5}-1}{2}$.

\item
$BI=|i/2+1|=\frac{\sqrt{5}}{2}$. $BJ=BI-1/2=\frac{\sqrt{5}-1}{2}$.

\item
Pour tracer un pentagone r\'egulier, on commence par tracer un
cercle $C_{1}$ et deux diam\`etres orthogonaux, qui jouent le r\^ole
du cercle passant par les sommets et des axes de coordonn\'ees. On
trace ensuite le milieu d'un des rayons: on obtient le point I de
la question 4. On trace le  cercle de centre $I$ passant par le
centre de $C_{1}$: c'est le cercle $\mathcal{C}$. On trace le
segment $BI$ pour obtenir son point $J$ d'intersection avec
$\mathcal{C}$. On trace enfin le cercle de centre $B$ passant par
$J$: il coupe $C_{1}$ en $A_{2}$ et $A_{3}$, deux sommets du
pentagone. Il suffit pour obtenir tous les sommets de reporter la
distance $A_{2}A_{3}$ sur $C_{1}$, une fois depuis $A_{2}$, une
fois depuis $A_{3}$. (en fait le cercle de centre $B$ et passant
par $J'$, le point de $\mathcal{C}$ diam\'etralement oppos\'e \`a
$J$, coupe $C_{1}$ en $A_{1}$ et $A_{4}$, mais nous ne l'avons pas
justifi\'e par le calcul : c'est un exercice !)
\end{enumerate}
% $$
% \includegraphics[65mm,50mm]{penta.eps}
% $$
\fincorrection
\correction{002925}
\begin{enumerate}
  \item 
  \item si $|a| \ne |b|$ : une solution unique, \\
         si $|a| = |b|$ : une droite ou $\varnothing$.
\end{enumerate}
\fincorrection
\correction{002926}
\begin{enumerate}
  \item 
  \item $\mathbb{U}\setminus\{1\}$, $i\R$, $\R\setminus\{1\}$.
\end{enumerate}
\fincorrection
\nocorrection
\correction{002929}
Les diagonales se coupent en leurs milieux, ont m{\^e}me longueur,
         et sont perpendiculaires $ \Rightarrow $ carr{\'e}.
\fincorrection
\correction{002930}
\begin{enumerate}
  \item $z \in \R$ ou $z \in -\frac 12 + i\R$.
  \item $z \in -1 + i\R$ ou $z \in i\R$
             ou $\left|z+\frac 12\right| = \frac 12$.
  \item $z\in i\R$ ou $|z-i| = \sqrt2$.
\end{enumerate}

\fincorrection
\correction{002931}
$(0,a,a+b,a+b+c = 1)$ forme un losange donc l'un des nombres vaut $1$
         et les deux autres sont oppos{\'e}s $ \Rightarrow  \{a,b,c\} = \{1,i,-i\}$.
\fincorrection
\nocorrection
\correction{002933}
$z=x+iy  \Rightarrow $ cercles $(\pm i, \sqrt2)$ (laborieux).
\fincorrection
\correction{002934}
On peut exprimer que $(AB)$ est la médiatrice du segment $[MM']$, mais on choisit ici d'utiliser le cours sur les similitudes. Cherchons donc à écrire la réflexion d'axe $(AB)$ en coordonnée complexe.
Cette réflexion s'écrit $z\mapsto \alpha \bar z + \beta$, avec $|\alpha|=1$ et $arg(\alpha)=2 arg(b-a)$. On obtient donc $\alpha= \frac{b-a}{\overline b - \overline a}$. 

Pour déterminer $\beta$, on peut exprimer le fait que $a$ est fixe:
\[
a = \frac{b-a}{\overline b - \overline a} \bar a + \beta 
\Leftrightarrow 
\beta = \frac{a\overline b - b\overline a }{\overline b - \overline a}
\]
Finalement, la réflexion s'écrit:
\[
z\mapsto \frac{b-a}{\overline b - \overline a} \cdot z + \frac{a\overline b - b\overline a }{\overline b - \overline a}.
\]

\fincorrection
\correction{002935}
$d = $ orthocentre de $abc$.
\fincorrection
\nocorrection
\correction{002937}
$\omega =
          \frac {a(c\bar c-b\bar b) + b(a\bar a-c\bar c) + c(b\bar b-a\bar a)}
                 {a(\bar c-\bar b) + b(\bar a-\bar c) + c(\bar b-\bar a)}$.
\fincorrection
\correction{002938}
\begin{enumerate}
  \item $z\in\R \Leftrightarrow \exists\ \alpha\in\R$ tq $u=\alpha v$.\par
    $x,y\in\R \Leftrightarrow \alpha = \frac1{|v|^2} \Leftrightarrow u = \frac 1{\bar v}$.
  \item 
  \item Il manque seulement les deux p{\^o}les.
\end{enumerate}

\fincorrection
\correction{005121}

\begin{enumerate}
 \item  On a $a=2\cos\left(\frac{2\pi}{5}\right)$ et $b=2\cos\left(\frac{4\pi}{5}\right)$. $1$, $z$, $z^2$, $z^3$ et $z^4$ sont les cinq
racines cinquièmes de $1$ dans $\Cc$. Par suite, $1+z+z^2+z^3+z^4=0$. Mais alors

$$a+b=z+z^2+z^3+z^4=-1$$
et

$$ab=(z+z^4)(z^2+z^3)=z^3+z^4+z^6+z^7=z+z^2+z^3+z^4=-1\;(\mbox{car}\;z^5=1).$$
$a$ et $b$ sont donc les solutions de l'équation $X^2+X-1=0$ dont les racines sont $\frac{-1+\sqrt{5}}{2}$ et
$\frac{-1-\sqrt{5}}{2}$. Enfin, puisque $\frac{2\pi}{5}\in\left]0,\frac{\pi}{2}\right[$, on a $a>0$. Par
suite, $\cos\left(\frac{2\pi}{5}\right)=\frac{-1+\sqrt{5}}{4}$ et $\cos\left(\frac{4\pi}{5}\right)=\frac{-1-\sqrt{5}}{4}$. D'autre part,
$\sin\left(\frac{2\pi}{5}\right)>0$ et donc,

$$\sin\left(\frac{2\pi}{5}\right)=+\sqrt{1-\cos^2(\frac{2\pi}{5})}=\sqrt{1-\left(\frac{-1+\sqrt{5}}{4}\right)^2}=\frac{1}{4}
\sqrt{10+2\sqrt{5}}.$$

\begin{center}
\shadowbox{
$\cos\left(\frac{2\pi}{5}\right)=\frac{\sqrt{5}-1}{4}$ et $\sin\left(\frac{2\pi}{5}\right)=\frac{1}{4}
\sqrt{10+2\sqrt{5}}$.
}
\end{center}
De même, en remplaçant $\sqrt{5}$ par $-\sqrt{5}$, $\cos\left(\frac{4\pi}{5}\right)=\frac{-1-\sqrt{5}}{4}$ et
$\sin\left(\frac{4\pi}{5}\right)=\frac{1}{4}\sqrt{10-2\sqrt{5}}$. Enfin,
$\cos\left(\frac{\pi}{5}\right)=-\cos\left(\pi-\frac{\pi}{5}\right)=-\cos\left(\frac{4\pi}{5}\right)=\frac{1+\sqrt{5}}{4}$
et $\sin\left(\frac{\pi}{5}\right)=\sin\left(\pi-\frac{\pi}{5}\right)=\sin\left(\frac{4\pi}{5}\right)=\frac{1}{4}\sqrt{10-2\sqrt{5}}$.\rule[-5mm]{0mm}{0mm}

 \item  Le rayon du grand cercle vaut, d'après le théorème de \textsc{Pythagore}~:

$$R=\sqrt{\Omega O^2+OM^2}=\frac{\sqrt{5}}{2}.$$
Donc $x_I=x_\Omega+R=\frac{-1+\sqrt{5}}{2}$ et $x_J=x_\Omega-R=\frac{-1-\sqrt{5}}{2}$. Par suite,
$x_I=2\cos\left(\frac{2\pi}{5}\right)$ et $x_J=2\cos\left(\frac{4\pi}{5}\right)$. Ceci montre que les médiatrices des segments $[O,I]$ et
$[O,J]$ coupent le cercle de centre $O$ et de rayon $1$ en quatre des cinq sommets du pentagone.

$$\includegraphics{../images/img005121-1}$$


 \item  Posons $x=\frac{AF}{AC}$. D'après le théorème de \textsc{Thales} (je vous laisse
vérifier les parallélismes),
$$x=\frac{AF}{AC}=\frac{HK}{HC}=\frac{FG}{FC}=\frac{AC-2AF}{AC-AF}=\frac{1-2x}{1-x}.$$
Donc $x^2-3x+1 = 0$ et puisque $x<1$, $x=\frac{3-\sqrt{5}}{2}$. Puis

$\frac{AG}{AC}=\frac{AC-AF}{AC}=1-x=\frac{-1+\sqrt{5}}{2}$ et
$\frac{FG}{AF}=\frac{AC-2AF}{AF}=\frac{1}{x}-2=\frac{2}{3-\sqrt{5}}-2=\frac{3+\sqrt{5}}{2}-2=\frac{-1+\sqrt{5}}{2}
$.

$$\includegraphics{../images/img005121-2}$$

Définition du \emph{nombre d'or}.

$$\includegraphics{../images/img005121-3}$$

On veut que $C$ partage le segment $[A,B]$ de telle sorte que $\frac{BC}{AC}=\frac{AC}{AB}$
(\og~$\frac{\mbox{petit}}{\mbox{moyen}}=\frac{\mbox{moyen}}{\mbox{grand}}$~\fg) c'est-à-dire, en posant $a=AB$ et
$x=AC$, $\frac{x}{a}=\frac{a-x}{x}$ ou encore $\left(\frac{x}{a}\right)^2+\frac{x}{a}-1=0$ et donc, puisque $\frac{x}{a}>0$,
$\frac{x}{a}=\frac{-1+\sqrt{5}}{2}$. 

\begin{center}
\shadowbox{
Le nombre d'or (ou proportion dorée) est le 
nombre $\frac{-1+\sqrt{5}}{2}=0,618...$
}
\end{center}
On peut aussi prendre pour le nombre d'or le
rapport $\frac{a}{x}=\frac{1+\sqrt{5}}{2}=1,618...$
\end{enumerate}
\fincorrection
\correction{005123}
\begin{enumerate}
 \item  On note $I_1$ le point d'intersection de la bissectrice $(\Delta_1)$ de l'angle
$\widehat{BAC}$ et de la droite $(BC)$. La parallèle à $(AC)$ passant par $B$ coupe $\Delta_1$ (puisque $(AC)$ n'est
pas parallèle à $(\Delta_1)$) en un point $A_1$. Les angles alternes-internes $\widehat{CAA_1}$ et $\widehat{AA_1B}$
sont alors égaux. Puisque d'autre part, $\widehat{CAA_1}=\widehat{A_1AB}$, on en déduit que
$\widehat{AA_1B}=\widehat{A_1AB}$ et donc que le triangle $(ABA_1)$ est isocèle en $B$. D'après le théorème de
\textsc{Thalès}, on a alors

$$\frac{I_1B}{I_1C}=\frac{A_1B}{AC}=\frac{AB}{AC}=\frac{c}{b},$$
et donc puisque $I_1$ est entre $B$ et $C$, $b\overrightarrow{I_1B}+c\overrightarrow{I_1C}=\vec{0}$, ou enfin
$I_1=\mbox{bar}\{B(b),C(c)\}$.

$$\includegraphics{../images/img005123-1}$$


On a aussi bien sûr les deux autres égalités $I_2=\mbox{bar}\{A(a),C(c)\}$ 
et $I_3=\mbox{bar}\{A(a),B(b)\}$ où $I_2$ et $I_3$ sont les points d'intersection des deux
autres bissectrices avec $(AC)$ et $(AB)$ respectivement.
Soit alors $I'=\mbox{bar}\{A(a),B(b),C(c)\}$. D'après le théorème du barycentre partiel, on a

$$I'=\mbox{bar}\{A(a),I_1(b+c)\}=\mbox{bar}\{B(b),I_2(a+c)\}=\mbox{bar}\{C(c),I_3(a+b)\},$$
ce qui montre que $I'$ est sur $(AI_1)$, $(BI_2)$ et $(CI_3)$, c'est-à-dire sur les trois bissectrices. Par suite,
$I'=I$.
 \item  Soit $z\in\Cc$.
$$z,\;z^2\;\mbox{et}\;z^3\;\mbox{ne sont pas deux à deux distincts}\Leftrightarrow
z^2=z\;\mbox{ou}\;z^3=z\;\mbox{ou}\;z^3=z^2\Leftrightarrow z\in\{-1,0,1\}.$$
Ensuite, pour $z\notin\{-1,0,1\}$,

$$z,\;z^2\;\mbox{et}\;z^3\;\mbox{sont
alignés}\Leftrightarrow\exists\lambda\in\Rr/\;z^3-z=\lambda(z^2-z)\Leftrightarrow\frac{z^3-z}{z^2-z}\in\Rr\Leftrightarrow z+1\in\Rr\Leftrightarrow z\in\Rr.$$
Finalement, $(z,z^2,z^3)$ est un \og vrai \fg~triangle si et seulement si $z$ n'est pas réel.
Soit alors $z$ un complexe non réel.

\begin{align*}
O\;\mbox{centre du cercle inscrit au triangle}\;(PQR)&\Leftrightarrow O=\mbox{bar}\{P(QR),Q(PR),R(PQ)\}\\
 &\Leftrightarrow z|z^2-z^3|+z^2|z-z^3|+z^3|z-z^2|=0
\Leftrightarrow z.|z|.|1-z|(|z|+z|1+z|+z^2)=0\\
 &\Leftrightarrow|z|+z|1+z|+z^2=0\;(E)\;(\mbox{car}\;z\notin\Rr)
\end{align*}
Ensuite,

\begin{align*}
|z|+z|1+z|+z^2=0&\Leftrightarrow(z+\frac{|z|}{z})+|1+z|=0\Rightarrow z+\frac{|z|}{z}\in\Rr\Leftrightarrow
z+\frac{|z|}{z}={\bar z}+\frac{|z|}{{\bar z}}\\
 &\Leftrightarrow z-{\bar z}-|z|\frac{z-{\bar z}}{z{\bar z}}=0\Leftrightarrow(z-{\bar z})(1-\frac{1}{|z|})=0
\Leftrightarrow1-\frac{1}{|z|}=0\;(\mbox{car}\;z\neq{\bar z})\\
 &\Leftrightarrow|z|= 1
\end{align*}
Posons donc $z=e^{i\theta}$ où $\theta\notin\pi\Zz$. En reportant dans $(E)$, on obtient

\begin{align*}
z\;\mbox{solution de}\;(E)&\Leftrightarrow e^{i\theta}+e^{-i\theta}+|1+e^{i\theta}|=
0\Leftrightarrow2\cos\theta+|e^{i\theta/2}|.|2\cos\frac{\theta}{2}|=0\\
 &\Leftrightarrow\cos\theta+|\cos\frac{\theta}{2}|=0\Leftrightarrow2|\cos\frac{\theta}{2}|^2+|\cos\frac{\theta}{2}|-1=0
\Leftrightarrow|\cos\frac{\theta}{2}|\;\text{est solution de l'équation}\;2X^2+X-1=0\\
 &\Leftrightarrow\left|\cos\frac{\theta}{2}\right|\in\left\{\frac{1}{2},-1\right\}\Leftrightarrow\left|\cos\frac{\theta}{2}\right|=\frac{1}{2}
\Leftrightarrow2\cos^2\frac{\theta}{2}-1=-\frac{1}{2}
\Leftrightarrow\cos\theta=-\frac{1}{2}\Leftrightarrow\theta\in\pm\frac{2\pi}{3}+2\pi\Zz\\
 &\Leftrightarrow
z\in\{j,j^2\}
\end{align*}
Les nombres complexes solutions sont donc $j$ et $j^2$.
\end{enumerate}
\fincorrection
\correction{005124}
\begin{align*}
(A,B,C)\;\mbox{équilatéral}&\Leftrightarrow C=r_{A,\pi/3}(B)\;\mbox{ou}\;C=r_{A,-\pi/3}(B)
\Leftrightarrow c-a=(-j^2)(b-a)\;\mbox{ou}\;c-a=(-j)(b-a)\\
 &\Leftrightarrow(-1-j^2)a+j^2b+c=0\;\mbox{ou}\;(-1-j)a+jb+c=0
\Leftrightarrow ja+j^2b+c=0\;\mbox{ou}\;j^2a+jb+c=0\\
 &\Leftrightarrow(j^2)^2a+j^2b+c=0\;\mbox{ou}\;j^2a+jb+c=0
\Leftrightarrow j\;\mbox{ou}\;j^2\;\mbox{sont solutions de l'équation}\;az^2+bz+c=0.
\end{align*}
Ensuite

\begin{align*}
(A,B,C)\;\mbox{équilatéral}&\Leftrightarrow ja+j^2b+c=0\;\mbox{ou}\;j^2a+jb+c=0\\
 &\Leftrightarrow(ja+j^2b+c)(j^2a+jb+c)=0\Leftrightarrow a^2+b^2+c^2+(j+j^2)(ab+ac+bc)=0\\
 &\Leftrightarrow a^2+b^2+c^2=ab+ac+bc,
\end{align*}
puis

\begin{align*}
(A,B,C)&\;\mbox{équilatéral}\Leftrightarrow a^2+b^2+c^2-ab-ac-bc=0\\
 &\Leftrightarrow-a^2+ab+ac-bc-b^2+bc+ba-ac-c^2+ca+cb-ab=0\\
 &\Leftrightarrow(c-a)(a-b)+(a-b)(b-c)+(b-c)(c-a)=0
\Leftrightarrow\frac{(c-a)(a-b)+(a-b)(b-c)+(b-c)(c-a)}{(b-c)(c-a)(a-b)}=0\\
 &\Leftrightarrow\frac{1}{b-c}+\frac{1}{c-a}+\frac{1}{a-b}=0.
\end{align*}
\fincorrection
\correction{005133}
\textbf{A- Solutions algébriques.}] Pour $z\in\Cc$, posons $z=x+iy$ où $(x,y)\in\Rr^2$.
\begin{enumerate}
 \item  $$|Z|=1\Leftrightarrow\frac{|1+z|^2}{|1-z|^2}=1\Leftrightarrow(1+x)^2+y^2=(1-x)^2+y^2\;\text{et}\;(x,y)\neq(1,0)\Leftrightarrow4x=0\Leftrightarrow x=0.$$
L'ensemble cherché est la droite $(Oy)$.

 \item 

\begin{align*}
|Z|=2&\Leftrightarrow(1+x)^2+y^2=4((1-x)^2+y^2)\Leftrightarrow3x^2+3y^2-10x+3=0\;\text{et}\;(x,y)\neq(1,0)\\
 &\Leftrightarrow x^2+y^2-\frac{10}{3}x+1=0\;\text{et}\;(x,y)\neq(1,0)\\
 &\Leftrightarrow\left(x-\frac{5}{3}\right)^2+y^2=\frac{16}{9}\;\text{et}\;(x,y)\neq(1,0)\\
 &\Leftrightarrow\left(x-\frac{5}{3}\right)^2+y^2=\frac{16}{9}.
\end{align*}
L'ensemble cherché est le cercle de centre $\Omega\left(\frac{5}{3},0\right)$ et de rayon $\frac{4}{3}$.

 \item 

\begin{align*}
Z\in\Rr&\Leftrightarrow Z=\overline{Z}\Leftrightarrow\frac{1+z}{1-z}=\frac{1+{\bar z}}{1-{\bar z}}\\
 &\Leftrightarrow(1+z)(1-{\bar z})=(1-z)(1+{\bar z})\;\text{et}\;z\neq1\Leftrightarrow z-{\bar z}={\bar z}-z\;\text{et}\;z\neq1\\
  &\Leftrightarrow z={\bar z}\;\text{et}\;z\neq1\Leftrightarrow z\in\Rr\;\text{et}\;z\neq1.
\end{align*}
L'ensemble cherché est la droite $(Ox)$ privé du point $(1,0)$.
 \item 

\begin{align*}
Z\in i\Rr&\Leftrightarrow Z=-\overline{Z}\Leftrightarrow\frac{1+z}{1-z}=-\frac{1+{\bar z}}{1-{\bar z}}\Leftrightarrow(1+z)(1-{\bar z})=-(1-z)(1+{\bar z})\;\text{et}\;z\neq1\\
 &\Leftrightarrow 1-z{\bar z}=-1+z{\bar z}\;\text{et}\;z\neq1\Leftrightarrow|z|^2=1\;\text{et}\;z\neq1\Leftrightarrow|z|=1\;\text{et}\;z\neq1.
\end{align*}
L'ensemble cherché est donc le cercle de centre $O$ et de rayon $1$ privé du point $(1,0)$.
\end{enumerate}
\textbf{B- Solutions géométriques.} Soient $A$ et $B$ les points d'affixes respectives $-1$ et $1$ et
$\mathcal{E}$ l'ensemble cherché. Soit $M$ un point du plan distinct de $B$ d'affixe $z$.
\begin{enumerate}
 \item 

$$M\in\mathcal{E}\Leftrightarrow|z+1|=|z-1|\Leftrightarrow AM=BM\Leftrightarrow M\in\mbox{med}[AB]=(Oy).$$
 \item  Soit $\Omega=\mbox{bar}(A(1),B(-4))$. On a $x_\Omega=\frac{-1}{5}(x_A-4x_B)=\frac{5}{3}$ et
$y_\Omega=\frac{-1}{5}(y_A-4y_B)=0$.

\begin{align*}
M\in\mathcal{E}&\Leftrightarrow|z+1|^2=4|z-1|^2\Leftrightarrow AM^2=4BM^2\\
 &\Leftrightarrow\overrightarrow{AM}^2-4\overrightarrow{BM}^2=0\Leftrightarrow(\overrightarrow{A\Omega}+\overrightarrow{\Omega
M})^2-4(\overrightarrow{B\Omega}+\overrightarrow{\Omega M})^2=0\\
 &\Leftrightarrow-3\overrightarrow{\Omega M}^2+2(\overrightarrow{A\Omega}-4\overrightarrow{B\Omega}).\overrightarrow{\Omega
M}+\overrightarrow{A\Omega}^2-4\overrightarrow{B\Omega}^2=0\\
 &\Leftrightarrow\Omega M^2=\frac{1}{3}(\Omega A^2-4\Omega B^2)
\end{align*}
Or, $\Omega A^2=(\frac{5}{3}+1)^2=\frac{64}{9}$ et $\Omega B^2=(\frac{5}{3}-1)^2=\frac{4}{9}$. Par suite,

$$\frac{1}{3}(\Omega A^2-4\Omega B^2)=\frac{1}{3}(\frac{64}{9}-\frac{16}{9})=\frac{16}{9}.$$
Ainsi,

$$M\in\mathcal{E}\Leftrightarrow\Omega M^2=\frac{16}{9}\Leftrightarrow\Omega M=\frac{4}{3},$$
et on retrouve le cercle de centre $\Omega(\frac{5}{3},0)$ et de rayon $\frac{4}{3}$.
 \item 

\begin{align*}
M\in\mathcal{E}&\Leftrightarrow z=-1\;\mbox{ou}\;\mbox{arg}\left(\frac{1+z}{1-z}\right)=0\;(\pi)
\Leftrightarrow M=A\;\mbox{ou}\;(\overrightarrow{BM},\overrightarrow{AM})=0\;(\pi)\\
 &\Leftrightarrow M\in(AB)\setminus\{B\}.
\end{align*}
et on retrouve la droite $(Ox)$ privée du point $(1,0)$.
 \item 

\begin{align*}
M\in\mathcal{E}&\Leftrightarrow z=-1\;\mbox{ou}\;\mbox{arg}\left(\frac{1+z}{1-z}\right)=\frac{\pi}{2}\;(\pi)
\Leftrightarrow M=A\;\mbox{ou}\;\left(\overrightarrow{BM},\overrightarrow{AM}\right)=\frac{\pi}{2}\;(\pi)\\
 &\Leftrightarrow M\;\mbox{est sur le cercle de diamètre}\;[AB]\;\mbox{privé de}\;B.
\end{align*}
et on retrouve le cercle de centre $O$ et de rayon $1$ privé du point $(1,0)$.

\end{enumerate}

\fincorrection
\correction{005134}
Soit $f$ la transformation considérée.
\begin{enumerate}
 \item  $f$ est la translation de vecteur $\overrightarrow{u}(3,-1)$.
 \item  $\omega=2\omega+3\Leftrightarrow\omega=-3$. $f$ est l'homothétie de rapport $2$ et de centre $\Omega(-3,0)$.
 \item  $\omega=i\omega+1\Leftrightarrow\omega=\frac{1}{2}(1+i)$. Comme $i=e^{i\pi/2}$, $f$ est la rotation d'angle
$\frac{\pi}{2}$ et de centre $\Omega(\frac{1}{2},\frac{1}{2})$.
 \item  $\omega=(1-i)\omega+2+i\Leftrightarrow\omega=1-2i$. Comme $1-i=\sqrt{2}e^{-i\pi/4}$, $f$ est la similitude de centre
$\Omega(1,-2)$, de rapport $\sqrt{2}$ et d'angle $-\frac{\pi}{4}$.
\end{enumerate}
\fincorrection
\correction{007004}
\begin{enumerate}
\item L'égalité $p=\frac{a-ib}{1-i}$ est équivalente à deux autres assertions ayant un sens géométrique plus clair. D'une part :
\[p=\frac{a-ib}{1-i} \Leftrightarrow p=\frac{a+b}{2}+i\frac{(a-b)}{2}\]
(Attention, $a$ et $b$ sont complexes : ceci n'est pas une forme algébrique.)  D'autre part, 
\[p=\frac{a-ib}{1-i} \Leftrightarrow (1-i)p=a-ib.\]
Ceci fournit deux façons de prouver le résultat demandé.

\begin{enumerate}
\item Dans le premier cas, on reconnaît que $(a+b)/2$ est l'affixe du milieu $M$ de $A$ et $B$, que $(a-b)/2$ est l'affixe du vecteur $\frac12 \overrightarrow{BA}$, et que $i(a-b)/2$ correspond à la rotation de $\pi/2$ de ce vecteur. On voit donc que le point d'affixe $\frac{a+b}{2}+i\frac{(a-b)}{2}$ est le point obtenu en translatant $M$ du vecteur  $\frac12\overrightarrow{AA'}$ (avec la notation $A'$ introduite plus haut). Ce point est bien le milieu $P$ du carré.

\item Dans le second cas, on réécrit l'égalité sous la forme 
\[a-p=i(b-p).\]
Ceci est vrai car cela signifie que $A$ est l'image de $B$ par la rotation de centre $P$ est d'angle $\pi/2$, ce qui est vrai dans un carré.
\end{enumerate}


\item Il s'agit de montrer que $\frac{s-q}{r-p}=i$. On calcule donc $\frac{s-q}{r-p}$ en utilisant la question  précédente:
\[
\frac{s-q}{r-p}
= \frac{d-ia -(b-ic)}{(c-id)-(a-ib)} 
=  \frac{d-b+i(c-a)}{c-a+i(b-d)}
\]
Là, on voit directement que le numérateur est égal à $i$ fois le dénominateur et c'est terminé (on le voit d'autant plus facilement que l'on sait que c'est ce  que l'on doit montrer : sinon, il faut une étape de simplification supplémentaire).

\item Considérer par exemple des rotations de centres $P, Q, R$ ou $S$ et d'angle $\pi/2$.
\end{enumerate}

Remarque : il existe une version un peu plus générale du théorème pour un quadrilatère $ABCD$ quelconque.
\fincorrection
\correction{007005}

\begin{enumerate}

\item Dans chaque carré, les sommets sont obtenus les uns des autres par rotations de $\pi/2$ par rapport aux centres. On a donc $a-p = i(b-p)$, $b-q = i(c-q)$ et $c-r = i(a-r)$. En développant ces expressions on obtient $p=\frac{a-ib}{1-i}$, $q=\frac{b-ic}{1-i}$ et $r=\frac{c-ia}{1-i}$.



\item Le centre de gravité de $ABC$ en est l'isobarycentre, donc son affixe est $\frac13(a+b+c)$. Celui de $PQR$ a pour affixe $\frac13(p+q+r)$. Il s'agit donc simplement de montrer que $a+b+c=p+q+r$. Or d'après la question précédente, $p+q+r =  \frac{a-ib+b-ic+c-ia}{1-i} = a+b+c$, ce qu'il fallait démontrer.



\item Il s'agit de montrer que l'argument de $\frac{q-a}{r-p}$ est $\pm \pi/2$. Pour cela, il suffit de calculer:
\[
\frac{q-a}{r-p}
=\frac{b-a+i(a-c)}{c-a+i(b-a)}
=-i
\]

On remarque que non seulement les segments sont perpendiculaires, mais ils sont de même longueur (ce qui n'était pas demandé).

La question précédente montre que la droite $(AQ)$ est la hauteur  de $PQR$ issue de $Q$. Le même raisonnement appliqué aux deux autres carrés permet de montrer que $(AQ)$, $(BR)$ et $(CP)$ sont les hauteurs de $PQR$. Elles sont donc concourantes.

\end{enumerate}

\fincorrection
\correction{007006}
%Cet exercice est semblable à celui sur le théorème de Van Aubel. 

Par construction, $A$ est l'image de $B$ par la rotation de centre $W$ et d'angle $2\pi/3$. En posant  $j=e^{2i\pi/3}$, on a donc 
\[ a-w = j(b-w), \]
autrement dit 

\[ w = \frac{a-jb}{1-j}.\]
On obtient de même $u=\frac{b-jc}{1-j}$ et $v=\frac{c-ja}{1-j}$.





On en déduit tout d'abord que 

\[ u+v+w = \frac{1}{1-j}\left(a-jb+b-jc+c-ja\right) = a+b+c. \]

Donc $\frac{1}{3} \left(u+v+w\right) = \frac{1}{3} \left(a+b+c\right)$, c'est-à-dire que les triangles $ABC$ et $UVW$ ont même centre de gravité. Ensuite, par la caractérisation des triangles équilatéraux, $UVW$ est équilatéral direct ssi $u+jv+j^2w=0$. Or:

\[ u+jv+j^2w = \frac{1}{1-j}\left(b-jc+j(c-ja)+j^2(c-ja)\right)=0.\]

Donc $UVW$ est bien équilatéral direct.

\fincorrection
\correction{007007}
\begin{enumerate}
\item Développons puis refactorisons $y(z-x)+z(x-y) = x(z-y)$. On en déduit le résultat, par inégalité triangulaire.

\item Si deux des points sont égaux, alors le membre de droite n'a qu'un seul terme et l'inégalité est une égalité, et trois points non alignés sont toujours cocycliques.

\item On a $(b-a)(d-c) + (d-a)(c-b) =ad-cd-ab+cb=(a-c)(d-b)$. Par inégalité triangulaire, on a donc $|(a-c)(d-b)| \leq |(b-a)(d-c)| + |(d-a)(c-b)|$  c'est-à-dire $AC\cdot BD \leq AB\cdot CD + AD\cdot BC$.
\item Rappelons le cas d'égalité de l'inégalité triangulaire : si $z$ et $z'$ sont des complexes non nuls, alors $|z+z'| \leq |z|+|z'|$, avec égalité ssi $\exists \lambda \in \R_+,\: z' = \lambda z$. Ici,  on a donc égalité ssi $\exists \lambda \in \R_+$ tel que $(b-a)(d-c)=\lambda (d-a)(c-b)$, autrement dit ssi $[a,b,c,d] \in \R_+$, ce qu'il fallait démontrer.
\end{enumerate}

\fincorrection
\nocorrection
\nocorrection
\correction{007010}

Si $\theta \equiv 0 [\pi]$, c'est la droite $(AB)$. Sinon, c'est un cercle dont $[AB]$ est une corde, par le théorème de l'angle inscrit.

Si $\theta \equiv \pi/2 [\pi]$, c'est le cercle de diamètre $[AB]$.

Sinon, par le théorème de l'angle au centre, le centre de ce cercle a pour affixe $\frac{a+b}{2} \pm i \frac{b-a}{2 \tan \theta}$, selon si $\theta$ est aigu ou obtus.

\fincorrection
\correction{007145} 
\begin{enumerate}
\item C'est une similitude directe de rapport $|1+i|=\sqrt 2$, d'angle $-\pi/4$. Son centre $\Omega$ est son (unique) point fixe, son affixe $\omega$ vérifie donc $\omega = (1-i)\omega + i $ c'est-à-dire $\omega=\frac{i}{i}=1$.
\item C'est une similitude indirecte de rapport $|i|=1$, donc un antidéplacement. C'est donc une réflexion ou une réflexion glissée. Cherchons d'éventuels points fixes.

Un point $z=x+iy$ est fixe ssi $x+iy = i(x-iy)+1-i = y+1+i(x-1)$ autrement dit ssi $x-y=1$.



(Rédaction alternative de la recherche de points fixes, sans prendre la partie réelle et imaginaire : un point d'affixe $z$ est fixe ssi:
\begin{align*}
z=i\bar z +1-i
&\Leftrightarrow z - i\bar z -1+i=0\\
&\Leftrightarrow \overline{(1-i)}z +(1-i)\overline z -2=0
\end{align*}
On reconnaît l'équation complexe de la même droite.)

L'antidéplacement est donc la réflexion d'axe d'équation $x-y-1=0$.
\item C'est une similitude indirecte (notons-la $s$) de rapport $2$. Elle admet donc un point fixe d'affixe $\omega$ vérifiant $ \omega = 2i\bar \omega +3$, ce qui donne après calcul $\omega=-1-2i$.

Soit $h$ l'homothétie de centre $\Omega$ et de rapport $2$. Alors la similitude $s$ s'écrit $h\circ \sigma$, avec $\sigma$ une réflexion que l'on peut obtenir comme $h^{-1}\circ s$. En coordonnée complexe, $h^{-1}$ s'écrit $z\mapsto \frac12 (z+1+2i)-1-2i = \frac12 z -\frac{1+2i}{2}$. En composant, on trouve que $\sigma$ est représentée par 
\[z\mapsto \frac12 (2i\bar z+3) -\frac{1+2i}{2}
=i\bar z +1-i.
\]
D'après la question précédente, c'est la réflexion d'axe $y=x-1$.
\item C'est la réflexion glissée composée de la translation d'affixe $1$ et de la réflexion suivant l'axe des abscisses.
\end{enumerate}
\fincorrection  
\correction{007146}  
\begin{enumerate}
\item 
Par le cours, la rotation d'angle $\theta$ et de centre d'affixe $\omega$ est représentée par l'application 
\[\C\to \C, \: z\mapsto e^{i\theta}(z-\omega)+\omega.\]

La rotation d'angle $\pi/4$ et de centre d'affixe $2+3i$ est donc représentée par l'application
\[ \C\to \C,\: z\mapsto e^{i\pi/4}(z-(2+3i)) + 2+3i = e^{i\pi/4}z +\frac{4 + \sqrt 2 + i(6-5\sqrt 2)}{2}.\]

\item La réflexion d'axe d'équation  $y=2x+1$ est représentée par une application de la forme
\[ \C\to \C, z\mapsto a\overline z+b,\]
où $a \in \C^*$ et $b\in \C$ sont des paramètres à déterminer. On sait que $a$ est le nombre complexe de module un dont l'argument est le double de l'angle entre l'axe des abscisses et l'axe de la réflexion. Donc dans ce cas, $a=e^{2i \arctan(2)}$. On peut aussi calculer la forme algébrique de $a$ en écrivant $a=\frac{(1+2i)^2}{|(1+2i)^2|} = \frac{-3+4i}{5}$.

Calculons $b$. On peut par exemple injecter un affixe particulier $z$ et son image $z'$ dans l'équation $z'=\frac{-3+4i}{5} \overline z+b$ et obtenir $b$. Le plus simple est de prendre l'affixe d'un point sur l'axe de la réflexion, qui est donc un point fixe de la réflexion. Prenons par exemple $z=i=z'$. On obtient l'égalité $i =\frac{3-4i}{5} i+b$, d'où $b=i\left(1-\frac{3-4i}{5}\right) = i\frac{2+4i}{5} = \frac{-4+2i}{5}$.

Finalement, la réflexion d'axe d'équation  $y=2x+1$ est représentée par l'application:
\[ \C\to \C, z\mapsto \frac{(-3+4i)\overline z-4+2i}{5}.\]
\end{enumerate}
\fincorrection
\correction{007147} 
Manifestement le rapport est $2$ car le segment donné est envoyé sur un segment deux fois plus long.
\begin{enumerate}
\item Si la similitude est directe, son angle est $\pi/2$, donc la similitude s'écrit $z\mapsto 2i z+ b$ en coordonnée complexe. Comme $s(2)=i$, on en déduit que $4i+b=i$ c'est-à-dire $b=-3i$.

Le point fixe a alors pour affixe $\dfrac{b}{1-2i} = \dfrac{-3i}{1-2i} = \dfrac{6-3i}{5}$.

\begin{center}
\includegraphics{../images/img007147-1}
\end{center}

\item La similitude indirecte est obtenue en composant la similitude directe précédente par la symétrie d'axe passant par $i$ et $3i$, c'est-à-dire l'axe des ordonnées. Cette symétrie s'écrit $z\mapsto -\bar z$. On en déduit que la similitude indirecte envoyant les points d'affixes $2$ et $3$ sur ceux d'affixes $i$ et $3i$ s'écrit
\[z\mapsto -\overline{(2iz-3i)} = 2i\bar z-3i.
\]
(Alternativement, et sans utiliser la question précédente, on sait d'après l'énoncé que la similitude indirecte doit s'écrire $z\mapsto 2i \overline z + b$, et comme $s(2)=i$, on en déduit bien $i = 4i+b$ d'où $b=-3i$.)

Cette similitude possède un unique point fixe d'affixe $\omega$, solution de l'équation $\omega = 2i\bar\omega -3i$.

On remarque que $\omega$ est donc aussi l'unique point fixe de $s\circ s$, ce qui donne l'équation $\omega = 2i\overline{(2i\bar \omega-3i)} -3i = 4\omega-6-3i$ d'où $\omega = 2+i$.


Finalement, comme on sait que l'axe est dirigé par le vecteur d'affixe $1+i$ et que cet axe passe par le centre de la similitude, cet axe est la droite d'équation $y=x-1$

\begin{center}
\includegraphics{../images/img007147-2}
\end{center}

\end{enumerate} 
\fincorrection  
\correction{007148} 
La similitude est:
\begin{itemize}
\item une translation ssi $a^2=1$ c'est-à-dire ssi $a \in \{-1,1\}$.
\item une homothétie de rapport $-4$ ssi $a^2=-4$ c'est-à-dire ssi $a\in \{2i,-2i\}$.
\item une rotation d'angle $\pi/2$ ssi $a^2=e^{i\pi/2}$ c'est-à-dire ssi $a\in \{e^{i\pi/4},e^{-3i\pi/4}\}$.
\end{itemize}
\fincorrection  
\correction{007149} 
\begin{enumerate}
\item Il y a exactement deux similitudes envoyant le couple $(A,B)$ sur le couple $(B,C)$ : une directe et une indirecte, et l'une se déduit de l'autre en composant (à gauche) par la réflexion d'axe $(BC)$.

D'après l'énoncé, le triangle $ABC$ est rectangle isocèle en $A$. On a donc $BC = \sqrt 2 AB$ et donc $s$ est de rapport $\sqrt 2$.
De plus, une similitude conservant les angles non orientés, on a $(AB)\bot (AC) \Rightarrow (s(A)s(B)) \bot (s(A)s(C))$, c'est-à-dire que $s(C)$ appartient à la perpendiculaire à $(s(A)s(B)) = (BC)$ passant par $s(A)=B$. Comme d'autre part on sait que $Cs(C)=\sqrt 2 BC$, cela donne deux possibilités, suivant que la similitude est directe ou indirecte.
\item Si la similitude est directe, son angle est $3\pi/4$ puisque $(\overrightarrow{AB},\overrightarrow{BC}) = 3\pi/4$.\\

On fixe maintenant un repère orthonormé direct du plan dont l'origine est $A$. On note $a$, $b$ et $c$ les affixes des trois points et on a donc $a=0$ et $c=ib$ d'après l'énoncé.\\
Soit $(\alpha,\beta) \in \C^*\times \C$ tel que $s$ s'écrive $z\mapsto \alpha z+\beta$ en coordonnée complexe. Par ce qui précède, on a $\alpha = \sqrt 2 e^{3i\pi/4} = -1+i$. D'autre part, comme $A$ est envoyé sur $B$, on a $b = \beta$.\\
La similitude s'écrit donc $z\mapsto (i-1)z + b$. Son unique point fixe $\Omega$ a pour affixe 
\[
\omega = \frac{b}{1-(i-1)} 
= \frac{b}{2-i} 
= \frac{2b+ib}{5} 
= \frac{2b+c}{5}. 
\]
On peut reformuler ceci sous la forme $\omega = \frac{2}{5} a + \frac{2}{5} b + \frac{1}{5} c $, ce qui montre que $\Omega$ est le barycentre  de $(A,2/5)$, $(B,2/5)$ et $(C,1/5)$.

\item On fixe maintenant un repère orthonormé direct du plan dont l'origine est $A$, et tel que $B$ soit sur l'axe des abscisses (donc son affixe est réel). Soit $(\alpha,\beta) \in \C^*\times \C$ tel que $s$ s'écrive $z\mapsto \alpha \bar z+\beta$ en coordonnée complexe. On a $s(0)=b$ donc $\beta=b$, et $s(b)=c=ib$ donc $\alpha\bar b + b = ib$ ce qui donne $\alpha = i-1$. La transformation s'écrit donc dans ce repère 
\[ z\mapsto (i-1)\bar z + b,\]
 et $c$ est envoyé sur $(i+2)b$.

Cherchons un point fixe $\omega$ pour $s$ : si $s(\omega)=\omega$, alors en particulier $s(s(\omega))=\omega$, ce qui s'écrit $(i-1)\overline{(i-1)\overline{z}+b}+b=\omega $, autrement dit $2\omega+ib = \omega$ et finalement $\omega=-ib$. Réciproquement, on vérifie que le point d'affixe $-ib$ est bien fixe sous $S$. Comme on peut écrire $\omega = -ib = -c = 2\cdot a - c$, on en déduit que $\Omega$ est le barycentre de $(A,2)$ et $(C,-1)$.

Cherchons maintenant l'axe de la similitude indirecte. C'est une droite passant par le centre de la similitude et d'après l'écriture de la similitude, elle est dirigée par un vecteur d'affixe $e^{3i\pi/8}$. Un paramétrage en coordonnée complexe de cette droite est donc :
\[\Delta = \{ -c+t.e^{3i\pi/8}\:|\: t\in \R\}
\]

\end{enumerate}
\fincorrection  
\correction{007150} 
Soit $z\mapsto \alpha z + \beta$ la représentation complexe d'une similitude directe. En coordonnées cartésiennes, elle s'écrit $\begin{pmatrix}x\\y\end{pmatrix} \mapsto  
\begin{pmatrix}
\Re(\alpha) &-\Im(\alpha) \\ \Im(\alpha) & \Re(\alpha)\end{pmatrix} \cdot \begin{pmatrix}x\\ y\end{pmatrix} + \begin{pmatrix}\Re(\beta)\\ \Im(\beta)\end{pmatrix}$.

On en déduit que $\begin{pmatrix}x\\y\end{pmatrix} \mapsto  \begin{pmatrix}ax+cy+e\\ bx+dy+f\end{pmatrix}$ représente une similitude directe ssi :
\[ c=-b \text{ et } a=d\]
Réciproquement, toute transformation
$\begin{pmatrix}x\\y\end{pmatrix} \mapsto  \begin{pmatrix}ax-by+e\\ bx+ay+f\end{pmatrix}$ peut s'écrire $z\mapsto (a+ib)z+e+if$ et correspond donc à une similitude directe.


Pour les similitudes indirectes, on trouve la condition \[ c=b \text{ et } a=-d.\]

L'application de $\C$ dans $\C$ qui représente $\phi$ est
$z=x+iy \mapsto (-2x-y-1)+i(x-2y+1) = -2x-y + i(x-2y)-1+i = (-2+i)z-1+i$.  C'est une similitude directe de rapport $\sqrt 5$, d'angle $arg(-2+i)$ et de centre d'affixe $\frac{-1+i}{3-i}$.

La deuxième est la composée d'une réflexion suivant une droite vectorielle et d'une homothétie de rapport $2$ et de centre $O$.

\fincorrection  
\correction{007151} 
\begin{enumerate}
\item L'application $s$ est une similitude directe, de rapport $2/3$, d'angle $arg(i)=\pi/2$, et de centre d'affixe $\frac{\frac{1-5i}{3}}{1-\frac{2}{3i}} = \frac{13-13i}{13}=1-i=z_A$.
\item 
\begin{enumerate}
\item On a $AB_{n+1}=s(A)s(B_n) = \frac{2}{3}AB_n$.
\item D'après la question précédente, la suite réelle $AB_n$ est géométrique de raison $2/3$ et de terme initial $AB$, donc :
\[ \forall n \in \N, AB_{n} = \left(\frac{2}{3}\right)^n AB.\]

Soit maintenant $n\in \N$. Le point $B_n$ appartient au disque de centre $A$ et de rayon $10^{-2}$ ssi $AB_n \leq 10^{-2}$.
Or, on a 
\begin{align*}
AB_n \leq 10^{-2} & \Leftrightarrow \left(\frac{2}{3}\right)^n AB \leq 10^{-2} \\
& \Leftrightarrow n \ln\left(\frac{2}{3}\right)+\ln(AB) \leq -\ln(100)\\
& \Leftrightarrow n\ln\left(\frac{2}{3}\right) \leq -\ln(AB)-\ln(100)\\
& \Leftrightarrow n \geq \frac{\ln(AB)+\ln(100)}{\ln(3)-\ln(2)}
\end{align*}

On en déduit que le plus petit entier $N$ tel que $\forall n\geq n, AB_n \leq 10^{-2}$ est $\left \lceil \frac{\ln(AB)+\ln(100)}{\ln(3)-\ln(2)} \right \rceil$ (partie entière supérieure du réel).

Comme $AB = 15/2$, on peut calculer explicitement $n$ à l'aide d'une calculatrice et on trouve $N= 17$, mais ce n'était pas demandé.
\item Pour tout $n$, les points $A$, $B$ et $B_n$ sont distincts. Ils sont alignés ssi
\[ 2\cdot arg\left(\frac{b_n-a}{b-a}\right)=0 \:(\in \R/2\pi\Z).\]
Or on a :
\begin{align*}
b_{n+1}-a
&= s(b_n)-s(a)\\
&= \frac23i b_{n}+\frac13 - \frac53i - \left( \frac23i a+\frac13 - \frac53i\right)\\
&= \frac23 i (b_n-a).
\end{align*}
La suite complexe $(b_n-a)_{n\in\N}$ est donc une suite géométrique de raison $\frac23i$, et donc pour tout $n\in \N$ on a 

\[ b_n-a = \left(\frac23i\right)^n \cdot (b-a).\]

On en déduit que  dans $\R/2\pi\Z$, on a :
\[ 2\cdot arg\left[\frac{b_n-a}{b-a}\right]
= 2\cdot arg \left[\left(\frac23i\right)^n \right]
= 2n\cdot arg \left(\frac23i\right)
= 2n\cdot \frac{\pi}{2} 
= n\pi,
\] et donc que $A$, $B$ et $B_n$ sont alignés ssi $n\pi\equiv 0 [2\pi] \Leftrightarrow n\equiv 0 [2]$, autrement dit ssi $n$ est pair.
\end{enumerate}
\end{enumerate} 
\fincorrection  
\correction{007163} 
\begin{enumerate}
\item Le triangle $ABC$ est équilatéral direct ssi $C$ est l'image de $B$ par la rotation de centre $A$ et d'angle $\pi/3$. Ceci est équivalent à $c-a = e^{i\pi/3}(b-a) = -j^2(b-a)$.
Or, on a :
\begin{align*}
c-a  = -j^2(b-a) 
&\Leftrightarrow c-a(1+j^2)+bj^2=0\\
&\Leftrightarrow aj+bj^2+c=0\\
&\Leftrightarrow  a+bj+cj^2=0.
\end{align*}

\item 
Fixons un repère orthonormé direct de centre $O$. 
\begin{center}
\includegraphics{../images/img007163-1}
\end{center}
D'après l'énoncé, on a $b=-j^2a$, $d=-j^2c$ et $f=-j^2e$. On en déduit:
\begin{align*}
m+jn+j^2p 
&= \frac{b+c+j(d+e)+j^2(f+a)}{2}\\
&=\frac{-j^2a+c+j(-j^2c+e)+j^2(-j^2e+a)}{2}\\
&=\frac{a(j^2-j^2)+c(1-j^3)+e(j-j^4)}{2}\\
&=0.\\
\end{align*}
D'après la première question, $MNP$ est équilatéral direct.
\end{enumerate}
\fincorrection  
\correction{007164} 

Soit $O$ le centre de gravité de $ABC$. On choisit un repère de centre $O$ tel que l'affixe $a$ de $A$ soit réel. On a alors $b = ja$ et $c=j^2a$.
\begin{enumerate}
\item On a l'égalité d'angles de droites $(OA,BA)=-\pi/6$ (angle inscrit ou calcul en coordonnées). On en déduit que la réflexion d'axe $(AB)$ s'écrit $z\mapsto \alpha\bar z+\beta$, avec $\alpha=e^{-\pi/3}=-j$. Ensuite, l'image de l'origine $O$ par la réflexion $\sigma_{AB}$ est le point d'affixe $-j^2a$. Finalement, vis-à-vis du repère choisi, la réflexion $\sigma_{AB}$ est représentée en coordonnée complexe par:
\[ z\mapsto -j\bar z -j^2.\]
\item On calcule de même que les réflexions suivant $(BC)$ et $(CA)$ s'écrivent $z\mapsto -\bar z -1$ et $z\mapsto -j^2 \bar z -j$.  Soit $m$ l'affixe de $M$. D'après ce qui précède, ses images par les trois réflexions sont:
\[ 
a'= -\bar m-1,\quad
b'=-j^2\bar m-j,\quad
c'=-j\bar m-j^2.
\]

La relation $1+j+j^2=0$ donne alors $\frac{1}{3}(a'+b'+c')=0$, ce qu'il fallait démontrer.
\end{enumerate}
\fincorrection  
\correction{007166}
\begin{enumerate}
\item Le point $E$ est l'image de $A$ par la composée des quatre symétries centrales de centre $M_1$,  ... $M_4$. Montrons que cette composée est l'identité. Par le cours, le composée de ces quatre rotations est une translation. Il suffit alors de tester un point particulier pour vérifier que c'est l'identité. Si $A=M_1$ par exemple, on voit que $A=E$ en utilisant le théorème de Thalès.
\item On a $2z = (z+z')+(z-z')$, d'où par inégalité triangulaire $2|z| \leq |z+z'|+|z-z'|$.
\item On a $a=m_1-m_2+t$, $b=m_1+m_2-t$, $c=-m_1+m_2+t$ et $d=-m_1-m_2-t$.

Le périmètre de $ABCD$ est  
\begin{align*}
p&=|b-a|+|c-b|+|d-c|+|a-d|\\
&=|2m_2-2t| + |-2m_1+2t|+|-2m_2-2t|+|2m_1+2t|\\
&=2\left(|m_1+t|+|m_1-t|+|m_2+t|+|m_2-t|\right).
\end{align*}

\item Le quadrilatère $AM_1OM_4$ est un parallélogramme direct ssi $a=m_1-m_2$ c'est-à-dire ssi $t=0$. Il s'agit de montrer que le périmètre calculé plus haut est minimal lorsque $t=0$.
D'après la question 3, on a bien 
\[
|m_1-t|+|m_1+t| \geq 2|m_1| \text{ et }
|m_2-t|+|m_2+t| \geq 2|m_2|,
\]
ce qu'il fallait démontrer.
\end{enumerate}
\fincorrection
\nocorrection
\nocorrection
\correction{000080}
Nous avons par la formule de Moivre
$$\cos5\theta+i\sin5\theta=e^{i5\theta}=(e^{i\theta})^5=
(\cos\theta+i\sin\theta)^5.$$
On d\'eveloppe ce dernier produit, puis on identifie parties r\'eelles
et parties imaginaires. On obtient :
\begin{eqnarray*}
\cos 5\theta & = & \cos^5\theta -10\cos^3\theta \sin^2\theta
                    +5\cos\theta \sin^4\theta\\
\sin 5\theta & = & 5\cos^4\theta \sin \theta -10\cos^2\theta \sin^3\theta
                   +\sin^5\theta 
\end{eqnarray*}

Remarque : Gr\^ace \`a la formule $\cos^2\theta +\sin^2\theta =1$, on pourrait
 continuer les
calculs et exprimer $\cos 5\theta $ en fonction de $\cos\theta $, et
$\sin 5\theta $ en fonction de $\sin\theta $.
\fincorrection
\nocorrection
\nocorrection
\nocorrection
\nocorrection
\nocorrection
\nocorrection
\nocorrection
\nocorrection
\correction{000089}
\begin{enumerate}
\item  $\sin\left(5x\right)=\sin\left(\frac{2\pi}{3}+x\right)$
ssi $x=\pi/6+k\pi/2$ ou $x=\pi/18+k\pi/3$, avec $k\in\mathbb{Z}$.
\item  $
  \sin\left(2x-\frac{\pi}{3}\right)=\cos\left(\frac{x}{3}\right)$ ssi
$x=5\pi/14+6k\pi/7$ ou $x=\pi/2+6k\pi/5$, avec $k\in\mathbb{Z}$.
 \item $\cos\left(3x\right)=\sin\left(x\right)$ ssi $x=\pi/8+k\pi/2$ ou
  $x=-\pi/4+k\pi$, avec $k\in\mathbb{Z}$.
\end{enumerate}
\fincorrection
\correction{000090}
L'équation
$\sqrt{3}\cos(x)-\sin(x)=m$ a des solutions ssi $m\in\left[-2,2\right]$ et pour
  $m=\sqrt{2}$, les solutions sont $x=\pi/12+2k\pi$ ou
  $x=-5\pi/12+2k\pi$, $k\in\mathbb{Z}$.
\fincorrection
\correction{000091}
$\cos(5x) + \cos(3x) \leq \cos x$ ssi
$2\cos(4x)\cos(x)\leq \cos x$ et $2\cos^2(x) -9\cos(x)
+4 >0$ ssi $\cos x>1/2$ ssi
$x\in\left]-\pi/6+2k\pi,\pi/6+2k\pi\right[$, $k\in\mathbb{Z}$.
\fincorrection
\correction{000092}
\begin{enumerate}
\item  $\cos^2(x)-\sin^2(x)=\sin(3x)$ ssi $x=\pi/2+2k\pi$ ou
 $x=-\pi/10+2k\pi/5$, avec $k\in\mathbb{Z}$.
\item  $\cos^4(x)-\sin^4(x)=1$ ssi $x=k\pi$, avec $k\in\mathbb{Z}$.
\end{enumerate}
\fincorrection
\correction{002951}
\begin{enumerate}
  \item $\frac {2^n+2\cos(n\pi/3)}3$.
  \item $2^n\cos(n\pi/3)$.
  \item $\left(2\cos\frac\theta2\right)^n\cos\frac{n\theta}2$.
  \item $\left(2\cos\frac\theta2\right)^n\sin\frac{(n+2)\theta}2$.
  \item $\left(2\cos\frac b2\right)^n\cos\left(a+\frac {nb}2\right)$.
\end{enumerate}
\fincorrection
\correction{002952}
\begin{enumerate}
  \item $\frac {n\sin\left(\frac{(2n+1)\theta}2\right)\sin\frac\theta2
             -\sin^2\frac{n\theta}2}
            {2\sin^2\frac\theta2}$
     si $\theta\not\equiv 0 (\mathrm{mod}\, {2\pi})$.
  \item $\frac {3\sin(n\theta/2)\sin\bigl((n+1)\theta/2\bigr)}{4\sin(\theta/2)}
    -\frac {\sin(3n\theta/2)\sin\bigl(3(n+1)\theta/2\bigr)}{4\sin(3\theta/2)}$.
\end{enumerate}
\fincorrection
\correction{002953}
$x \equiv -y \equiv \pm \frac {2\pi}3 (\mathrm{mod}\, {2\pi})$.
\fincorrection
\correction{002954}
\begin{enumerate}
  \item $= 3/2$.
  \item $32\cos^6(\theta) = \cos6\theta + 6\cos4\theta + 15\cos2\theta
    + 10 \Rightarrow \Sigma = \frac {15}8$.
  \item $\Sigma_p = \frac {pC_{2p}^p}{2^{2p-1}}$.
\end{enumerate}
\fincorrection
\correction{002955}
$x \equiv 0 (\mathrm{mod}\, {\frac\pi n})$, $x \not\equiv \frac \pi2 (\mathrm{mod}\, \pi)$.
\fincorrection
\correction{002956}
$\frac 12\Bigl( (x+e^{i\alpha})^n + (x+e^{-i\alpha})^n \Bigr) = 0
 \Leftrightarrow x = \mathrm{cotan}\left(\frac{(2k+1)\pi}{2n}\right)\sin\alpha - \cos\alpha$.
\fincorrection
\correction{002957}
$S = \frac {u-u^{-1}}{u^2-u^{-2}} +
     \frac {u-u^{-1}}{u^4-u^{-4}} + \dots +
     \frac {u-u^{-1}}{u^{2^n}-u^{-2^n}}$.
\par
$\frac {u-u^{-1}}{u^2-u^{-2}} +
 \frac {u-u^{-1}}{u^4-u^{-4}} =
 \frac {u^3-u^{-3}}{u^4-u^{-4}}$.
\par
$\frac {u^3-u^{-3}}{u^4-u^{-4}} +
 \frac {u-u^{-1}}{u^8-u^{-8}} =
 \frac {u^7-u^{-7}}{u^8-u^{-8}}$
$\ldots$
\par
$ \Rightarrow  S = \frac {u^{2^n-1}-u^{-2^n+1}}{u^{2^n}-u^{-2^n}}
       = \frac {\sin\bigl((2^n-1)\theta\bigr)}{\sin(2^n\theta)}$.
\fincorrection
\correction{002958}
$\tan(nx) = \frac {C_n^1\tan x - C_n^3\tan^3 x + \cdots}
                            {C_n^0       - C_n^2\tan^2 x + \cdots}$.
\fincorrection
\correction{002959}
$z = e^{i\theta}  \Rightarrow  a = \tan\frac\theta2$
         pour $\theta \not \equiv \pi (\mathrm{mod}\,{2\pi})$.
\fincorrection
\correction{005063}

\begin{enumerate}
 \item  $\sin x=0\Leftrightarrow x\in\pi\Zz$. De plus, $\mathcal{S}_{[0,2\pi]}=\{0,\pi,2\pi\}$.
 \item  $\sin x=1\Leftrightarrow x\in\frac{\pi}{2}+2\pi\Zz$. De plus, $\mathcal{S}_{[0,2\pi]}=\left\{\frac{\pi}{2}\right\}$.
 \item  $\sin x=-1\Leftrightarrow x\in-\frac{\pi}{2}+2\pi\Zz$. De plus, $\mathcal{S}_{[0,2\pi]}=\left\{\frac{3\pi}{2}\right\}$.
 \item  $\cos x=1\Leftrightarrow x\in2\pi\Zz$. De plus, $\mathcal{S}_{[0,2\pi]}=\{0,2\pi\}$.

 \item  $\cos x=-1\Leftrightarrow x\in\pi+2\pi\Zz$. De plus, $\mathcal{S}_{[0,2\pi]}=\{\pi\}$.
 \item  $\cos x=0\Leftrightarrow x\in\frac{\pi}{2}+\pi\Zz$. De plus,
$\mathcal{S}_{[0,2\pi]}=\left\{\frac{\pi}{2},\frac{3\pi}{2}\right\}$.
 \item  $\tan x=0\Leftrightarrow x\in\pi\Zz$. De plus, $\mathcal{S}_{[0,2\pi]}=\left\{0,\pi,2\pi\right\}$.
 \item  $\tan x=1\Leftrightarrow x\in\frac{\pi}{4}+\pi\Zz$. De plus,
$\mathcal{S}_{[0,2\pi]}=\left\{\frac{\pi}{4},\frac{5\pi}{4}\right\}$.
\end{enumerate}
\fincorrection
\correction{005064}
\begin{enumerate}
 \item $\sin x=\frac{1}{2} \Leftrightarrow x\in\left(\frac{\pi}{6}+2\pi\Zz\right)\cup\left(\frac{5\pi}{6}+2\pi\Zz\right)$. De plus,
$\mathcal{S}_{[0,2\pi]}=\left\{\frac{\pi}{6},\frac{5\pi}{6}\right\}$.

 \item $\sin x=-\frac{1}{\sqrt{2}}\Leftrightarrow x\in\left(-\frac{\pi}{4}+2\pi\Zz\right)\cup\left(-\frac{3\pi}{4}+2\pi\Zz\right)$. De plus,
$\mathcal{S}_{[0,2\pi]}=\left\{-\frac{\pi}{4},-\frac{3\pi}{4}\right\}$.

 \item $\tan x=-1\Leftrightarrow x\in-\frac{\pi}{4}+\pi\Zz$. De plus, $\mathcal{S}_{[0,\pi]}=\left\{\frac{3\pi}{4}\right\}$.
 
 \item $\tan x=\frac{1}{\sqrt{3}}\Leftrightarrow x\in\frac{\pi}{6}+\pi\Zz$. De plus,
$\mathcal{S}_{[0,\pi]}=\left\{\frac{\pi}{6}\right\}$.

 \item $\cos x=\frac{\sqrt{3}}{2}\Leftrightarrow x\in\left(-\frac{\pi}{6}+\pi\Zz\right)\cup\left(\frac{\pi}{6}+\pi\Zz\right)$. De
plus, $\mathcal{S}_{[0,2\pi]}=\left\{\frac{\pi}{6},\frac{11\pi}{6}\right\}$.

 \item $\cos x=-\frac{1}{\sqrt{2}}\Leftrightarrow x\in\left(-\frac{3\pi}{4}+\pi\Zz\right)\cup\left(\frac{3\pi}{4}+\pi\Zz\right)$. De
plus, $\mathcal{S}_{[0,2\pi]}=\left\{\frac{3\pi}{4},\frac{5\pi}{4}\right\}$.

\end{enumerate}
\fincorrection
\correction{005065}

\begin{enumerate}
 \item  $\sin(2x)=\frac{1}{2}\Leftrightarrow2x\in\left(\frac{\pi}{6}+2\pi\Zz\right)\cup\left(\frac{5\pi}{6}+2\pi\Zz\right)\Leftrightarrow
x\in\left(\frac{\pi}{12}+\pi\Zz\right)\cup\left(\frac{5\pi}{12}+\pi\Zz\right)$. De plus, $\mathcal{S}_{[0,2\pi]}=\left\{\frac{\pi}{12},\frac{5\pi}{12},\frac{13\pi}{12},\frac{17\pi}{12}\right\}$.
 \item  $\sin\frac{x}{2}=-\frac{1}{\sqrt{2}}\Leftrightarrow\frac{x}{2}\in\left(\frac{5\pi}{4}+2\pi\Zz\right)\cup\left(\frac{7\pi}{4}
+2\pi\Zz\right)\Leftrightarrow x\in\left(\frac{5\pi}{2}+4\pi\Zz)\cup(\frac{7\pi}{2}+4\pi\Zz\right)$. De
plus, $\mathcal{S}_{[0,4\pi]}=\left\{\frac{5\pi}{2},\frac{7\pi}{2}\right\}$.
 \item  $\tan(5x)=1\Leftrightarrow5x\in\frac{\pi}{4}+\pi\Zz\Leftrightarrow x\in\frac{\pi}{20}+\frac{\pi}{5}\Zz$. De plus,
$\mathcal{S}_{[0,\pi]}=\left\{\frac{\pi}{20},\frac{\pi}{4},\frac{9\pi}{20},\frac{13\pi}{20},\frac{17\pi}{20}\right\}$.
 \item  $\cos(2x)=\cos^2x\Leftrightarrow \cos(2x)=\frac{1}{2}(1+\cos(2x))\Leftrightarrow\cos(2x)=1\Leftrightarrow2x\in2\pi\Zz\Leftrightarrow x\in\pi\Zz$. De
plus, $\mathcal{S}_{[0,2\pi]}=\{0,\pi,2\pi\}$.
 \item  $2\cos^2x-3\cos x+1=0\Leftrightarrow(2\cos x-1)(\cos x-1)=0\Leftrightarrow\cos x=\frac{1}{2}\;\mbox{ou}\;\cos x=1\Leftrightarrow
x\in\left(-\frac{\pi}{3}+2\pi\Zz\right)\cup\left(\frac{\pi}{3}+2\pi\Zz\right)\cup2\pi\Zz$. De plus,
$\mathcal{S}_{[0,2\pi]}=\left\{0,\frac{\pi}{3},\frac{5\pi}{3},2\pi\right\}$.
 \item  $\cos(nx)=0\Leftrightarrow nx\in\frac{\pi}{2}+\pi\Zz\Leftrightarrow x\in\frac{\pi}{2n}+\frac{\pi}{n}\Zz$.
 \item  $|\cos(nx)|=1\Leftrightarrow nx\in\pi\Zz\Leftrightarrow x\in\frac{\pi}{n}\Zz$.
 \item  $\sin(nx)=0\Leftrightarrow nx\in\pi\Zz\Leftrightarrow x\in\frac{\pi}{n}\Zz$.
 \item  $|\sin(nx)|=1\Leftrightarrow nx\in\frac{\pi}{2}+\pi\Zz\Leftrightarrow x\in\frac{\pi}{2n}+\frac{\pi}{n}\Zz$.
 \item  $\sin x=\tan x\Leftrightarrow\sin x-\frac{\sin x}{\cos x}=0\Leftrightarrow\sin x\frac{\cos x-1}{\cos x}=0\Leftrightarrow\sin x=0\;\mbox{ou}\;\cos x=1\Leftrightarrow x\in\pi\Zz$. De
plus, $\mathcal{S}_{[0,2\pi]}=\{0,\pi,2\pi\}$.
 \item
\begin{align*}
\sin(2x)+\sin x=0&\Leftrightarrow\sin(2x)=\sin(x+\pi)\Leftrightarrow(\exists k\in\Zz/\;2x=x+\pi+2k\pi)\;\mbox{ou}\;(\exists
k\in\Zz/\;2x=-x+2k\pi)\\
 &\Leftrightarrow (\exists k\in\Zz/\;x=\pi+2k\pi)\;\mbox{ou}\;(\exists k\in\Zz/\;x=\frac{2k\pi}{3})
\end{align*}
De plus, $\mathcal{S}_{[0,2\pi]}=\{0,\frac{2\pi}{3},\pi,\frac{4\pi}{3},2\pi\}$.
 \item

\begin{align*}
12\cos^2x-8\sin^2x=2&\Leftrightarrow6\cos^2x-4(1-\cos^2x)=1\Leftrightarrow\cos^2x=\frac{1}{2}\Leftrightarrow\cos
x=\frac{1}{\sqrt{2}}\;\mbox{ou}\;\cos=-\frac{1}{\sqrt{2}}\\
 &\Leftrightarrow
x\in\left(-\frac{\pi}{4}+\pi\Zz\right)\cup\left(\frac{\pi}{4}+\pi\Zz\right)\Leftrightarrow x\in\frac{\pi}{4}+\frac{\pi}{2}\Zz.
\end{align*}
\end{enumerate}
\fincorrection
\correction{005066}

\begin{enumerate}
 \item   Pour $x\in[-\pi,\pi]$, $\cos x\leq\frac{1}{2}\Leftrightarrow x\in\left[-\pi,-\frac{\pi}{3}\right]\cup\left[\frac{\pi}{3},\pi\right]$.
 \item   Pour $x\in\Rr$, $\sin
x\geq-\frac{1}{\sqrt{2}}\Leftrightarrow x\in\displaystyle\bigcup_{k\in\Zz}\left[-\frac{\pi}{4}+2k\pi,\frac{5\pi}{4}+2k\pi\right]$.
 \item   Pour $x\in[0,2\pi]$,

\begin{align*}
\cos x>\cos\frac{x}{2}&\Leftrightarrow2\cos^2\frac{x}{2}-\cos\frac{x}{2}-1>0\Leftrightarrow(2\cos\frac{x}{2}+1)(\cos\frac{x}{2}-1)>0\Leftrightarrow
2\cos\frac{x}{2}+1<0\;\mbox{et}\;\cos\frac{x}{2}\neq1\\
 &\Leftrightarrow\cos\frac{x}{2}<-\frac{1}{2}\;\mbox{et}\;\frac{x}{2}\notin2\pi\Zz
\Leftrightarrow\frac{x}{2}\in\displaystyle\bigcup_{k\in\Zz}\left]\frac{2\pi}{3}+2k\pi,\frac{4\pi}{3}+2k\pi\right[\;\mbox{et}\;
 x\notin4\pi\Zz\\
 &\Leftrightarrow x\in\displaystyle\bigcup_{k\in\Zz}\left]\frac{4\pi}{3}+4k\pi,\frac{8\pi}{3}+4k\pi\right[\;\mbox{et}\;
 x\notin4\pi\Zz
\Leftrightarrow x\in]\frac{4\pi}{3},2\pi]
\end{align*}
 \item   Pour $x\in[-\pi,\pi]$, $\cos^2x\geq\cos(2x)\Leftrightarrow\frac{1}{2}(1+\cos(2x))\geq\cos(2x)\Leftrightarrow\cos(2x)\leq1\Leftrightarrow
x\in[-\pi,\pi]$.
 \item   Pour $x\in[0,2\pi]$, $\cos^2x\leq\frac{1}{2}\Leftrightarrow-\frac{1}{\sqrt{2}}\leq\cos x\leq\frac{1}{\sqrt{2}}\Leftrightarrow
x\in\left[\frac{\pi}{4},\frac{3\pi}{4}\right]\cup\left[\frac{5\pi}{4},\frac{7\pi}{4}\right]$.
 \item   Pour $x\in[0,2\pi]$,

\begin{align*}
\cos\frac{x}{3}\leq\sin\frac{x}{3}&\Leftrightarrow\frac{1}{\sqrt{2}}\sin\frac{x}{3}-\frac{1}{\sqrt{2}}\cos\frac{x}{3}\geq0
\Leftrightarrow\sin\left(\frac{x}{3}-\frac{\pi}{4}\right)\geq0\Leftrightarrow\exists k\in\Zz/\;2k\pi\leq\frac{x}{3}-\frac{\pi}{4}\leq\pi+2k\pi\\
 &\Leftrightarrow\exists k\in\Zz/\;\frac{3\pi}{4}+6k\pi\leq x\leq3\pi+\frac{3\pi}{4}+6k\pi\Leftrightarrow\frac{3\pi}{4}\leq x\leq2\pi
\end{align*}
\end{enumerate}
\fincorrection
\correction{005067}
$\cos^2\frac{\pi}{8}=\frac{1}{2}\left(1+\cos(2\times\frac{\pi}{8})\right)=\frac{1}{2}\left(1+\frac{\sqrt{2}}{2}\right)=\frac{2+\sqrt{2}}{4}$,
et puisque $\cos\frac{\pi}{8}>0$,

\begin{center}
\shadowbox{$\cos\frac{\pi}{8}=\frac{1}{2}\sqrt{2+\sqrt{2}}$.}
\end{center}
De même, 
puisque $\sin\frac{\pi}{8}>0$, $\sin\frac{\pi}{8}=\sqrt{\frac{1}{2}\left(1-\cos(2\times\frac{\pi}{8})\right)}$ 
et
\begin{center}
\shadowbox{
$\sin\frac{\pi}{8}=\frac{1}{2}\sqrt{2-\sqrt{2}}$.
}
\end{center}

\fincorrection
\correction{005068}
$$\cos\frac{\pi}{12}=\cos\left(\frac{\pi}{3}-\frac{\pi}{4}\right)=\cos\frac{\pi}{3}\cos\frac{\pi}{4}+\sin\frac{\pi}{3}
\sin\frac{\pi}{4}=\frac{\sqrt{6}+\sqrt{2}}{4}.$$
De même,
$$\sin\frac{\pi}{12}=\sin\left(\frac{\pi}{3}-\frac{\pi}{4}\right)=\sin\frac{\pi}{3}\cos\frac{\pi}{4}-\sin\frac{\pi}{3}
\sin\frac{\pi}{4}=\frac{\sqrt{6}-\sqrt{2}}{4}.$$

\begin{center}
\shadowbox{
$\cos\frac{\pi}{12}=\frac{\sqrt{6}+\sqrt{2}}{4}$ et $\sin\frac{\pi}{12}=\frac{\sqrt{6}-\sqrt{2}}{4}.
$}\end{center}
\fincorrection
\correction{005069}
Pour $n$ naturel non nul, on pose $S_n=\sum_{}^{}e^{i(\pm a_1\pm...\pm a_n)}$.
\textbullet~$S_1=e^{ia_1}+e^{-ia_1}=2\cos a_1$
\textbullet~Soit $n\geq1$. Supposons que $S_n=2^n\cos a_1...\cos a_n$ alors

\begin{align*}
S_{n+1}&=\sum_{}^{}e^{i(\pm a_1\pm...\pm a_{n+1})}=e^{ia_{n+1}}\sum_{}^{}e^{i(\pm a_1\pm...\pm
a_n)}+e^{-ia_{n+1}}\sum_{}^{}e^{i(\pm a_1\pm...\pm a_n)}\\
 &=2\cos(a_{n+1})S_n=2^{n+1}\cos a_1...\cos a_{n+1}.
\end{align*}
On a montré par récurrence que : $\forall n\geq1,\;S_n=2^n\cos a_1...\cos a_n$.
Ensuite, pour $n\geq1$, $\sum_{}^{}\cos(\pm a_1\pm...\pm a_n)=\Re(S_n)=2^n\cos a_1...\cos a_n$ (et on obtient aussi
$\sum_{}^{}\sin(\pm a_1\pm...\pm a_n)=\Im(S_n)=0$).

\begin{center}
\shadowbox{
$\forall n\in\Nn^*$, $\sum_{}^{}\cos(\pm a_1\pm...\pm a_n)=2^n\cos a_1...\cos a_n$.
}
\end{center}
\fincorrection
\correction{005070}
\begin{enumerate}
 \item  Soit $n\in\Nn^*$. Puisque $a$ est dans $]0,\pi[$ alors, pour tout entier naturel non nul $k$, $\frac{a}{2^k}$ est dans $]0,\pi[$ et donc
$\sin\frac{a}{2^k}\neq0$. De plus, puisque $\sin\left(\frac{a}{2^{k-1}}\right)=\sin\left(2\times\frac{a}{2^k}\right)=2\sin\left(\frac{a}{2^k}\right)\cos\left(\frac{a}{2^k}\right)$, on a~:

$$\prod_{k=1}^{n}\cos\left(\frac{a}{2^k}\right)=\prod_{k=1}^{n}\frac{\sin\left(\frac{a}{2^{k-1}}\right)}{2\sin\left(\frac{a}{2^k}\right)}=\frac{1}{2^{
n}}
\frac{\sin(a)\sin\left(\frac{a}{2}\right)\ldots\sin\left(\frac{a}{2^{n-1}}\right)}
{\sin\left(\frac{a}{2}\right)\ldots\sin\left(\frac{a}{2^{n-1}}\right)\sin\left(\frac{a}{2^{n}}\right)}=\frac{\sin
a}{2^{n}\sin\frac{a}{2^n}}.$$

\begin{center}
\shadowbox{
$\forall a\in]0,\pi[$, $\forall n\in\Nn^*$, $\prod_{k=1}^{n}\cos\left(\frac{a}{2^k}\right)=\frac{\sin
a}{2^{n}\sin\frac{a}{2^n}}$.
}
\end{center}

 \item  $\forall k\in\Nn^*,\;\cos\left(\frac{a}{2^k}\right)>0$ car $\frac{a}{2^k}$ est dans $]0,\frac{\pi}{2}[$. Puis

\begin{align*}
\sum_{k=1}^{n}\ln\left(\cos(\frac{a}{2^k})\right)&=\ln\left(\prod_{k=1}^{n}\cos(\frac{a}{2^k})\right)=\ln\left(\frac{\sin
a}{2^{n}\sin\frac{a}{2^n}}\right)=\ln\left(\frac{\sin a}{a}\right)-\ln\left(\frac{\sin\frac{a}{2^n}}{\frac{a}{2^n}}\right).
\end{align*}
Maintenant, $\lim_{n\rightarrow +\infty}\frac{\sin\frac{a}{2^n}}{\frac{a}{2^n}}=\lim_{x\rightarrow 0}\frac{\sin x}{x}=1$ et donc,

$$\lim_{n\rightarrow +\infty}\sum_{k=1}^{n}\ln\left(\cos(\frac{a}{2^k})\right)=\lim_{n\rightarrow +\infty}\left(\ln(\frac{\sin
a}{a})-\ln(\frac{\sin\frac{a}{2^n}}{\frac{a}{2^n}})\right)=\ln\left(\frac{\sin a}{a}\right).$$

\begin{center}
\shadowbox{
$\forall a\in]0,\pi[$, $\lim_{n\rightarrow +\infty}\sum_{k=1}^{n}\ln\left(\cos(\frac{a}{2^k})\right)=\ln\left(\frac{\sin a}{a}\right)$.
}
\end{center}
\end{enumerate}
\fincorrection
\correction{005071}
Soit $x\in\Rr$.

\begin{align*}
2^{4\cos^2x+1}+16.2^{4\sin^2x-3}=20&\Leftrightarrow2^{4\cos^2x+1}+16.2^{1-4\cos^2x}=20\Leftrightarrow2^{4\cos^2x}-10+16\times2^{-4\cos^2x}=0\\
 &\Leftrightarrow2^{4\cos^2x}-10+\frac{16}{2^{4\cos^2x}}=0\Leftrightarrow(2^{4\cos^2x})^2-10\times2^{4\cos^2x}+16=0\\
 &\Leftrightarrow2^{4\cos^2x}=2\;\mbox{ou}\;2^{4\cos^2x}=8\Leftrightarrow4\cos^2x=1\;\mbox{ou}\;4\cos^2x=3\\
 &\Leftrightarrow\cos x=\frac{1}{2}\;\mbox{ou}\;\cos x=-\frac{1}{2}\;\mbox{ou}\;\cos x=\frac{\sqrt{3}}{2}\;\mbox{ou}\;\cos
x=-\frac{\sqrt{3}}{2}\\
 &\Leftrightarrow x\in\left(\frac{\pi}{6}+\frac{\pi}{2}\Zz\right)\cup\left(\frac{\pi}{3}+\frac{\pi}{2}\Zz\right).
\end{align*}
\fincorrection
\correction{005072}
\begin{enumerate}
 \item  Tout d'abord, d'après la formule de \textsc{Moivre},

$$\cos(3\theta)+i\sin(3\theta)=(\cos\theta+i\sin\theta)^3=(\cos^3\theta-3\cos\theta\sin^2\theta)
+i(3\cos^2\theta\sin\theta-\sin^3\theta),$$

et par identification des parties réelles et imaginaires,

\begin{center}
\shadowbox{
$\forall\theta\in\Rr,\;\cos(3\theta)=\cos^3\theta-3\cos\theta\sin^2\theta\;\mbox{et}\;\sin(3\theta)
=3\cos^2\theta\sin\theta-\sin^3\theta.$
}
\end{center}
Ensuite, $\tan(3\theta)$ et $\tan\theta$ existent $\Leftrightarrow3\theta\notin\frac{\pi}{2}+\pi\Zz$ et
$\theta\notin\frac{\pi}{2}+\pi\Zz\Leftrightarrow3\theta\notin\frac{\pi}{2}+\pi\Zz\Leftrightarrow\theta\notin\frac{\pi}{6}+
\frac{\pi}{3}\Zz$.
Soit donc $\theta\notin\frac{\pi}{6}+\frac{\pi}{3}\Zz$.

$$\tan(3\theta)=\frac{\sin(3\theta)}{\cos(3\theta)}
=\frac{3\cos^2\theta\sin\theta-\sin^3\theta}{\cos^3\theta-3\cos\theta\sin^2\theta}
=\frac{3\tan\theta-\tan^3\theta}{1-3\tan^2\theta},$$
après division du numérateur et du dénominateur par le réel non nul $\cos^3\theta$.

\begin{center}
\shadowbox{
$\forall\theta\in\Rr\setminus\left(\frac{\pi}{6}+\frac{\pi}{3}\Zz\right),\;\tan(3\theta)
=\frac{3\tan\theta-\tan^3\theta}{1-3\tan^2\theta}.$
}\end{center}
 \item  Soit $a\neq\pm\frac{1}{\sqrt{3}}$.
\textbf{1ère méthode.} $a$ est bien sûr racine de l'équation proposée, ce qui permet d'écrire~:

\begin{align*}
\frac{3x-x^3}{1-3x^2}=\frac{3a-a^3}{1-3a^2}&\Leftrightarrow(3x-x^3)(1-3a^2)=(1-3x^2)(3a-a^3)\;(\mbox{car}\;\pm\frac{1}{\sqrt{3}}
\;\mbox{ne sont pas solution de l'équation})\\
 &\Leftrightarrow(x-a)((3a^2-1)x^2+8ax-a^2+3)=0.
\end{align*}
Le discriminant réduit du trinôme $(3a^2-1)x^2+8ax-a^2+3$ vaut~:

$$\Delta'=16a^2-(3a^2-1)(-a^2+3)=3a^4+6a^2+3=(\sqrt{3}(a^2+1))^2>0.$$
L'équation proposée a donc trois racines réelles~:

\begin{center}
\shadowbox{
$\mathcal{S}=\left\{a,\frac{4a-\sqrt{3}(a^2+1)}{1-3a^2},\frac{4a+\sqrt{3}(a^2+1)}{1-3a^2}\right\}.$
}
\end{center}
\textbf{2ème méthode.} Il existe un unique réel
$\alpha\in\left]-\frac{\pi}{2},\frac{\pi}{2}\right[\setminus\left\{-\frac{\pi}{6},\frac{\pi}{6}\right\}$ tel que $a=\tan\alpha$. De même,
si $x$ est un réel distinct de $\pm\frac{1}{\sqrt{3}}$, il existe un unique réel
$\theta\in\left]-\frac{\pi}{2},\frac{\pi}{2}\right[\setminus\left\{-\frac{\pi}{6},\frac{\pi}{6}\right\}$ tel que $x=\tan\theta$ (à savoir
$\alpha=\Arctan a$ et $\theta=\arctan x$). Comme $\pm\frac{1}{\sqrt{3}}$ ne sont pas solution de l'équation proposée,
on a~:

\begin{align*}
\frac{3x-x^3}{1-3x^2}=\frac{3a-a^3}{1-3a^2}&\Leftrightarrow\frac{3\tan\theta-\tan^3\theta}{1-3\tan^2\theta}
=\frac{3\tan\alpha-\tan^3\alpha}{1-3\tan^2\alpha}\Leftrightarrow\tan(3\theta)=\tan(3\alpha)\\
 &\Leftrightarrow3\theta\in3\alpha+\pi\Zz\Leftrightarrow\theta\in\alpha+\frac{\pi}{3}\Zz.
\end{align*}
Ceci refournit les solutions $x=\tan\alpha=a$, puis

$$x=\tan\left(\alpha+\frac{\pi}{3}\right)=\frac{\tan\alpha+\tan\frac{\pi}{3}}{1-\tan\alpha\tan\frac{\pi}{3}}=\frac{a+\sqrt{3}}{1
-\sqrt{3}a}=\frac{(a+\sqrt{3})(1+\sqrt{3}a)}{1-3a^2}=\frac{4a+\sqrt{3}(a^2+1)}{1-3a^2},$$
et $x=\tan\left(\alpha-\frac{\pi}{3}\right)=\frac{4a-\sqrt{3}(a^2+1)}{1-3a^2}$.

\end{enumerate}
\fincorrection
\correction{005073}
\begin{enumerate}
 \item  Pour $x\notin\frac{\pi}{10}+\frac{\pi}{5}\Zz$,

$$\tan(5x)=\frac{\Im((e^{ix})^5)}{\Re((e^{ix})^5)}=\frac{5\cos^4x\sin
x-10\cos^2x\sin^3x+\sin^5x}{\cos^5x-10\cos^3x\sin^2x+5\cos x\sin^4x}=\frac{5\tan
x-10\tan^3x+\tan^5x}{1-10\tan^2x+5\tan^4x},$$
après division du numérateur et du dénominateur par le réel non nul $\cos^5x$.

\begin{center}
\shadowbox{
$\forall x\in\Rr\setminus\left(\frac{\pi}{10}+\frac{\pi}{5}\Zz\right),\;\tan(5x)=\frac{5\tan
x-10\tan^3x+\tan^5x}{1-10\tan^2x+5\tan^4x}.$
}
\end{center}

 \item  $9^\circ$, $-27^\circ$, $-63^\circ$ et $81^\circ$ vérifient
$\tan(5\times9^\circ)=\tan(5\times(-27^\circ))=\tan(5\times(-63^\circ))=\tan(5\times81^\circ)=1$. On résoud donc l'équation~:

$$\tan(5x)=1\Leftrightarrow5x\in\left(\frac{\pi}{4}+\pi\Zz\right)\Leftrightarrow x\in\left(\frac{\pi}{20}+\frac{\pi}{5}\Zz\right).$$
Les solutions, exprimées en degrés et éléments de $]-90^\circ,90^\circ[$, sont $-63^\circ$, $-27^\circ$, $9^\circ$,
$45^\circ$ et $81^\circ$. Ainsi, les cinq nombres $\tan(-63^\circ)$, $\tan(-27^\circ)$, $\tan(9^\circ)$,
$\tan(45^\circ)$ et $\tan(81^\circ)$ sont deux à deux distincts et solutions de l'équation
$\frac{5X-10X^3+X^5}{1-10X^2+5X^4}=1$ qui s'écrit encore~:

$$X^5-5X^4-10X^3+10X^2+5X-1=0.$$

Le polynôme $X^5-5x^4-10X^3+10X^2+5X-1$ admet déjà $\tan(45^\circ)=1$ pour racine et on a

$$X^5-5X^4-10X^3+10X^2+5X-1=(X-1)(X^4-4X^3-14X^2-4X+1).$$
Les quatre nombres $\tan(-63^\circ)$, $\tan(-27^\circ)$, $\tan(9^\circ)$ et $\tan(81^\circ)$ sont ainsi les racines du
polynôme $X^4-4X^3-14X^2-4X+1$. Ce dernier peut donc encore s'écrire
$(X-\tan(9^\circ))(X+\tan(27^\circ))(X+\tan(63^\circ))(X-\tan(81^\circ))$. L'opposé du coefficient de $X^3$ à savoir $4$
vaut donc également $\tan(9^\circ)-\tan(27^\circ)-\tan(63^\circ)+\tan(81^\circ)$ et on a montré que~:

\begin{center}
\shadowbox{
$\tan(9^\circ)-\tan(27^\circ)-\tan(63^\circ)+\tan(81^\circ)=4.$
}
\end{center}
\end{enumerate}
\fincorrection
\correction{005074}
Pour $x\in[0,\pi]$, posons $f(x)=\tan x+\tan(2x)+\tan(3x)+\tan(4x)$.

\begin{align*}
f(x)\;\mbox{existe}&\Leftrightarrow\tan x,\;\tan(2x),\;\tan(3x)\;\mbox{et}\;\tan(4x)\;\mbox{existent}\\
 &\Leftrightarrow(x\notin\frac{\pi}{2}+\pi\Zz),\;(2x\notin\frac{\pi}{2}+\pi\Zz),\;(3x\notin\frac{\pi}{2}+\pi\Zz)\;\mbox{et}\;
 (4x\notin\frac{\pi}{2}+\pi\Zz)\\
 &\Leftrightarrow(x\notin\frac{\pi}{2}+\pi\Zz),\;(x\notin\frac{\pi}{4}+\frac{\pi}{2}\Zz),\;(x\notin\frac{\pi}{6}+
 \frac{\pi}{3}\Zz)\;\mbox{et}\;(x\notin\frac{\pi}{8}+\frac{\pi}{4}\Zz)\\
 &\Leftrightarrow x\notin\left\{\frac{\pi}{8},\frac{\pi}{6},\frac{\pi}{4},\frac{3\pi}{8},
\frac{\pi}{2},\frac{5\pi}{8},\frac{3\pi}{4},\frac{5\pi}{6},\frac{7\pi}{8}\right\}.
\end{align*}
$f$ est définie et continue sur

$$\left[0,\frac{\pi}{8}\right[\cup\left]\frac{\pi}{8},\frac{\pi}{6}\right[\cup\left]\frac{\pi}{6},\frac{\pi}{4}\right[\cup\left]\frac{\pi}{4},
\frac{3\pi}{8}\right[\cup\left]\frac{3\pi}{8},\frac{\pi}{2}\right[\cup\left]\frac{\pi}{2},\frac{5\pi}{8}\right[\cup\left]\frac{5\pi}{8},
\frac{3\pi}{4}\right[\cup\left]\frac{3\pi}{4},\frac{5\pi}{6}\right[\cup\left]\frac{5\pi}{6},\frac{7\pi}{8}\right[\cup\left]\frac{7\pi}{8},\pi\right].$$
Sur chacun des dix intervalles précédents, $f$ est définie, continue et strictement croissante en tant que somme de
fonctions strictement croissantes. La restriction de $f$ à chacun de ces dix intervalles est donc bijective de
l'intervalle considéré sur l'intervalle image, ce qui montre déjà que l'équation proposée, que l'on note dorénavant
$(E)$, a au plus une solution par intervalle et donc au plus dix solutions dans $[0,\pi]$.
Sur $I=\left[0,\frac{\pi}{8}\right[$ ou $I=\left]\frac{7\pi}{8},\pi\right]$, puisque $f(0)=f(\pi)=0$, $(E)$ a exactement une
solution dans $I$. Ensuite, dans l'expression de somme $f$, une et une seule des quatre fonctions est un infiniment
grand en chacun des nombres considérés ci-dessus, à l'exception de $\frac{\pi}{2}$. En chacun de ces nombres, $f$ est un
infiniment grand. L'image par $f$ de chacun des six intervalles ouverts n'ayant pas $\frac{\pi}{2}$ pour borne est donc
$]-\infty,+\infty[$ et $(E)$ admet exactement une solution dans chacun de ces intervalles d'après le théorème des valeurs intermédiaires. Ceci porte le total à $6+2=8$
solutions.
En $\frac{\pi}{2}^-$, $\tan x$ et $\tan(3x)$ tendent vers $+\infty$ tandis que $\tan(2x)$ et $\tan(4x)$ tendent vers
$0$. $f$ tend donc vers $+\infty$ en $\frac{\pi}{2}^-$, et de même $f$ tend vers $-\infty$ en $\frac{\pi}{2}^+$.
L'image par $f$ de chacun des deux derniers intervalles est donc encore une fois $]-\infty,+\infty[$. Finalement,

\begin{center}
\shadowbox{
$(E)\;\text{admet exactement dix solutions dans}\;[0,\pi].$
}
\end{center}
\fincorrection
\correction{005075}
\begin{enumerate}
 \item  D'après les formules d'\textsc{Euler},

$$z+z^4=e^{2i\pi/5}+e^{8i\pi/5}=e^{2i\pi/5}+e^{-2i\pi/5}=2\cos\frac{2\pi}{5}=a.$$

De même,

$$z^2+z^3=e^{4i\pi/5}+e^{6i\pi/5}=e^{4i\pi/5}+e^{-4i\pi/5}=2\cos\frac{4\pi}{5}=b.$$
 \item  Puisque $z\neq1$ et $z^5=e^{2i\pi}=1$,

$$1+z+z^2+z^3+z^4=\frac{1-z^5}{1-z}=\frac{1-1}{1-z}=0.$$

 \item  $a+b=z+z^2+z^3+z^4=-1$ et $ab=(z+z^4)(z^2+z^3)=z^3+z^4+z^6+z^7=z+z^2+z^3+z^4=-1$. Donc,

$$a+b=-1\;\mbox{et}\;ab=-1.$$
Ainsi, $a$ et $b$ sont les solutions de l'équation $X^2+X-1=0$ à savoir les nombres $\frac{-1\pm\sqrt{5}}{2}$. Puisque
$\frac{2\pi}{5}\in\left]0,\frac{\pi}{2}\right[$ et $\frac{4\pi}{5}\in\left]\frac{\pi}{2},\pi\right[$, on a $a>0$ et $b>0$. Finalement,

\begin{center}
\shadowbox{
$\cos\frac{2\pi}{5}=\frac{-1+\sqrt{5}}{4}\;\mbox{et}\;\cos\frac{4\pi}{5}=\frac{-1-\sqrt{5}}{4}.$
}
\end{center}
\end{enumerate}
\fincorrection
\correction{005076}
\begin{enumerate}
 \item  $\cos^2x=\frac{1}{2}(1+\cos(2x))$ et une primitive de $x\mapsto\cos^2x$ est
$x\mapsto\frac{1}{2}(x+\frac{1}{2}\sin(2x))$.

 \item  D'après les formules d'\textsc{Euler},

\begin{align*}
\cos^4x&=\left(\frac{1}{2}(e^{ix}+e^{-ix})\right)^4=\frac{1}{16}(e^{4ix}+4e^{2ix}+6+4e^{-2ix}+e^{-4ix})\\
 &=\frac{1}{16}(2\cos(4x)+8\cos(2x)+6)=\frac{1}{8}(\cos(4x)+4\cos(2x)+3)
\end{align*}
Donc, une primitive de la fonction $x\mapsto\cos^4x$ est $x\mapsto\frac{1}{8}(\frac{1}{4}\sin(4x)+2\sin(2x)+3x)$.

 \item  

\begin{align*}
\sin^4x&=\left(\frac{1}{2i}(e^{ix}-e^{-ix})\right)^4=\frac{1}{16}(e^{4ix}-4e^{2ix}+6-4e^{-2ix}+e^{-4ix})\\
 &=\frac{1}{16}(2\cos(4x)-8\cos(2x)+6)=\frac{1}{8}(\cos(4x)-4\cos(2x)+3)
\end{align*}
Donc, une primitive de la fonction $x\mapsto\sin^4x$ est $x\mapsto\frac{1}{8}(\frac{1}{4}\sin(4x)-2\sin(2x)+3x)$.

 \item   $\cos^2x\sin^2x=\frac{1}{4}\sin^2(2x)=\frac{1}{8}(1-\cos(4x))$ et une primitive de la fonction $x\mapsto\cos^2x\sin^2x$
est $x\mapsto\frac{1}{8}(x-\frac{1}{4}\sin(4x))$.

 \item  

\begin{align*}
\sin^6x&=\left(\frac{1}{2i}(e^{ix}-e^{-ix})\right)^6=-\frac{1}{64}(e^{6ix}-6e^{4ix}+15e^{2ix}-20+15e^{-2ix}-6e^{-4ix}+
e^{-6ix})\\
 &=-\frac{1}{64}(2\cos(6x)-12\cos(4x)+30\cos(2x)-20)=\frac{1}{32}(-\cos(6x)+6\cos(4x)-15\cos(2x)+10)
\end{align*}
Donc, une primitive de la fonction $x\mapsto\sin^6x$ est
$x\mapsto\frac{1}{32}(-\frac{1}{6}\sin(6x)+\frac{3}{2}\sin(4x)-\frac{15}{2}\sin(2x)+10x)$.

 \item   $\cos x\sin^6x=\sin'x\sin^6x$ et une primitive de $x\mapsto\cos x\sin^6x$ est $x\mapsto\frac{1}{7}\sin^7x$.

 \item   $\cos^5x\sin^2x=\cos x(1-\sin^2x)^2\sin^2x=\sin'x\sin^2x-2\sin'x\sin^4x+\sin'x\sin^6x$ et une primitive de
$x\mapsto\cos^5x\sin^2x$ est $x\mapsto\frac{1}{3}\sin^3x-\frac{2}{5}\sin^5x+\frac{1}{7}\sin^7x$.

 \item   $\cos^3x=\sin'x-\sin'x\sin^2x$ et une primitive de $x\mapsto\cos^3x$ est $x\mapsto\sin x-\frac{1}{3}\sin^3x$.
\end{enumerate}
\fincorrection
\correction{005077}
\begin{enumerate}
 \item   Pour $x$ réel , on a~:
\begin{align*}
\cos^4x\sin^6x&=\left(\frac{1}{2}(e^{ix}+e^{-ix})\right)^4\left(\frac{1}{2i}(e^{ix}-e^{-ix})\right)^6\\
 &=-\frac{1}{2^{10}}(e^{4ix}+4e^{2ix}+6+4e^{-2ix}+e^{-4ix})
(e^{6ix}-6e^{4ix}+15e^{2ix}-20+15e^{-2ix}-6e^{-4ix}+e^{-6ix})\\
 &= -\frac{1}{2^{10}}(e^{10ix}-2e^{8ix}-3e^{6ix}+8e^{4ix}+2e^{2ix}-12
+2e^{-2ix}+8e^{-4ix}-3e^{-6ix}-2e^{-8ix}+e^{-10ix})\\
 &=-\frac{1}{2^9}(\cos10x-2\cos8x-3\cos6x+8\cos4x+2\cos2x-6)\\
 &=-\frac{1}{512}(\cos10x-2\cos8x-3\cos6x+8\cos4x+2\cos2x-6)
\end{align*}
(Remarque. La fonction proposée était paire et l'absence de sinus était donc prévisible. Cette remarque guidait aussi
les calculs intermédiaires~:~les coefficients de $e^{-2ix}$, $e^{-4ix}$,... étaient les mêmes que ceux
de $e^{2ix}$, $e^{4ix}$,...) Par suite,

\begin{align*}
I&=-\frac{1}{512}\left(\left[\frac{\sin10x}{10}-\frac{\sin8x}{4}-\frac{\sin6x}{2}+2\sin4x+\sin2x\right]_{\pi/6}^{\pi/3}
-6\left(\frac{\pi}{3}-\frac{\pi}{6}\right)\right)\\
 &=-\frac{1}{512}\left(\frac{1}{10}(-\frac{\sqrt{3}}{2}+\frac{\sqrt{3}}{2})-\frac{1}{4}(\frac{\sqrt{3}}{2}+\frac{\sqrt{3}}{2})-\frac{1}{2}(0-0)+2(-\frac{\sqrt{3}}{2}-\frac{\sqrt{3}}{2})+(\frac{\sqrt{3}}{2}-\frac{\sqrt{3}}{2})-\pi\right)\\
 &=-\frac{1}{512}\left(-\frac{\sqrt{3}}{4}-2\sqrt{3}-\pi\right)=\frac{9\sqrt{3}+4\pi}{2048}.
\end{align*}

 \item   Pour $x$ réel, on a

\begin{align*}
\cos^4x\sin^7x&=\cos^4x\sin^6x\sin x=\cos^4x(1-\cos^2x)^3\sin x\\
 &=\cos^4x\sin x-3\cos^6x\sin x+3\cos^8x\sin x-\cos^{10}x\sin x.
\end{align*}

Par suite,

\begin{align*}
J&=\left[-\frac{\cos^5x}{5}+\frac{3\cos^7x}{7}-\frac{\cos^9x}{3}+\frac{\cos^{11}x}{11}\right]_{\pi/6}^{\pi/3}\\
 &=-\frac{1}{5}\times\frac{1-9\sqrt{3}}{32}+\frac{3}{7}\times\frac{1-27\sqrt{3}}{128}-\frac{1}{3}\times\frac{1-81\sqrt
{3}}{512}+\frac{1}{11}\times\frac{1-243\sqrt{3}}{2048}\\
 &=\frac{1}{2^{11}\times3\times5\times7\times11}(-14784(1-9\sqrt{3})+7920(1-27\sqrt{3})-154
0(1-81\sqrt{3})+105(1-243\sqrt{3}))\\
 &=\frac{1}{{2365440}}(-{8299}+{18441}\sqrt{3}).
\end{align*}
\end{enumerate}
\fincorrection
\correction{005078}
\begin{enumerate}
 \item   $\tan\frac{x}{2}$ existe si et seulement si $x\notin\pi+2\pi\Zz$ et $\frac{1-\cos x}{\sin x}$ existe si et
seulement si $x\notin\pi\Zz$. Pour $x\notin\pi\Zz$,

$$\frac{1-\cos x}{\sin x}=\frac{2\sin^2\frac{x}{2}}{2\sin\frac{x}{2}\cos\frac{x}{2}}=\tan\frac{x}{2}.$$
 \item  \textbf{1 ère solution.} Pour tout réel $x$,

$$\sin(x-\frac{2\pi}{3})+\sin x+\sin(x+\frac{2\pi}{3})=-\frac{1}{2}\sin x-\frac{\sqrt{3}}{2}\cos x+\sin
x-\frac{1}{2}\sin x+\frac{\sqrt{3}}{2}\cos x=0,$$
\textbf{2 ème solution.}

$$\sin\left(x-\frac{2\pi}{3}\right)+\sin x+\sin(x+\frac{2\pi}{3})=\Im(e^{i(x-\frac{2\pi}{3})}+e^{ix}+e^{i(x+\frac{2\pi}{3})})
=\Im(e^{ix}(j^2+1+j))=0.$$
 \item  $\tan\left(\frac{\pi}{4}-x\right)$, $\tan\left(\frac{\pi}{4}+x\right)$ et $\frac{2}{\cos(2x)}$ existent si et seulement si
$\frac{\pi}{4}-x$, $\frac{\pi}{4}+x$ et $2x$ ne sont pas dans $\frac{\pi}{2}+\pi\Zz$, ce qui équivaut à
$x\notin\frac{\pi}{4}+\frac{\pi}{2}\Zz$. Donc, pour $x\notin\frac{\pi}{4}+\frac{\pi}{2}\Zz$,

\begin{align*}
\tan\left(\frac{\pi}{4}-x\right)+\tan\left(\frac{\pi}{4}+x\right)&=\frac{1-\tan x}{1+\tan x}+\frac{1+\tan x}{1-\tan x}\;(\text{pour}\;x\;\text{vérifiant de plus}\;x\notin\frac{\pi}{2}+\pi\Zz)\\
 &=\frac{\cos x-\sin x}{\cos x+\sin x}+\frac{\cos x+\sin x}{\cos x-\sin x}=\frac{(\cos x-\sin x)^2+(\cos x+\sin x)^2}{\cos^2
x-\sin^2x}=\frac{2(\cos^2x+\sin^2x)}{\cos(2x)}\\
 &=\frac{2}{\cos(2x)}\;(\text{ce qui reste vrai pour}\;x\in\frac{\pi}{2}+\pi\Zz).
\end{align*}
 \item  Pour $x\notin\frac{\pi}{4}\Zz$,

$$\frac{1}{\tan x}-\tan x=\frac{\cos x}{\sin x}-\frac{\sin x}{\cos x}=\frac{\cos^2 x-\sin^2x}{\sin x\cos
x}=\frac{2\cos(2x)}{\sin(2x)}=\frac{2}{\tan(2x)}.$$
\end{enumerate}
\fincorrection
\correction{005079}
\begin{enumerate}
 \item  \textbullet~Pour tout réel $x$, $1-2k\cos x+k^2=(k-\cos x)^2+\sin^2x\geq0$. De plus,

$$1-2k\cos x+k^2=0\Rightarrow k-\cos x=\sin x=0\Rightarrow x\in\pi\Zz\;\mbox{et}\;k=\cos x\Rightarrow k\in\{-1,1\},$$
ce qui est exclu. Donc,

$$\forall k\in\Rr\setminus\{-1,1\},\;\forall x\in\Rr,\;1-2k\cos x+k^2>0.$$
\textbullet$f_k$ est donc définie sur $\Rr$, dérivable sur $\Rr$ en vertu de théorèmes généraux, impaire et $2\pi$-périodique. On
l'étudie dorénavant sur $[0,\pi]$. Pour $x\in[0,\pi]$, on a~:

\begin{align*}
f'_k(x)&=\cos x(1-2k\cos x+k^2)^{-1/2}-\frac{1}{2}\sin x(2k\sin x)(1-2k\cos x+k^2)^{-3/2}\\
 &=(1-2k\cos x+k^2)^{-3/2}(\cos x(1-2k\cos x+k^2)-k\sin^2x)\\
 &=(1-2k\cos x+k^2)^{-3/2}(-k\cos^2x+(1+k^2)\cos x-k)\\
 &=(1-2k\cos x+k^2)^{-3/2}(k\cos x-1)(k-\cos x)
\end{align*}

\begin{center}
\shadowbox{
$\forall x\in\Rr,\;f_k'(x)=\frac{(k\cos x-1)(k-\cos x)}{(1-2k\cos x+k^2)^{3/2}}.$
}
\end{center}

\begin{itemize}
\item[\textbf{1er cas}~:~$|k|<1$ et $k\neq0$.] (si $k=0$, $f_k(x)=\sin x$) Pour tout réel $x$, $(1-2k\cos x+k^2)^{-3/2}(k\cos
x-1)<0$ et $f_k'(x)$ est du signe de $\cos x-k$.

$$\includegraphics{../images/img005079-1}$$

(car $f_k(\Arccos k)=\frac{\sqrt{1-k^2}}{\sqrt{1-2k^2+k^2}}=1$).

\item[\textbf{2ème cas}~:~$k>1$.] Pour tout réel $x$, $(1-2k\cos x+k^2)^{-3/2}(k-\cos x)>0$ et $f_k'(x)$ est du signe de $k\cos
x-1$.

$$\includegraphics{../images/img005079-2}$$

(car $f_k(\Arccos\frac{1}{k})=\frac{\sqrt{1-\frac{1}{k^2}}}{\sqrt{1-2+k^2}}=\frac{1}{k}$).
\item[\textbf{3ème cas}~:~$k<-1$.] Pour tout réel $x$, $(1-2k\cos x+k^2)^{-3/2}(k-\cos x)<0$ et $f_k'(x)$ est du signe de
$1-k\cos x$.

$$\includegraphics{../images/img005079-3}$$



(car $f_k(\Arccos\frac{1}{k})=\frac{\sqrt{1-\frac{1}{k^2}}}{\sqrt{1-2+k^2}}=-\frac{1}{k}$).

\end{itemize}

 \item  Pour $k\in\Rr\setminus\{-1,1\}$, posons $I_k=\int_{0}^{\pi}f_k(x)\;dx$.

Si $k=0$, $I_k=\int_{0}^{\pi}\sin x\;dx=2$. Sinon,

\begin{align*}
I_k&=\frac{1}{k}\int_{0}^{\pi}\frac{2k\sin x}{2\sqrt{1-2k\cos x+k^2}}\;dx=\frac{1}{k}
\left[\sqrt{1-2k\cos x+k^2}\right]_0^\pi\\
 &=\frac{1}{k}(\sqrt{1+2k+k^2}-\sqrt{1-2k+k^2})=\frac{1}{k}(|k+1|-|k-1|).
\end{align*}
Plus précisément, si $k\in]-1,1[\setminus\{0\}$, $I_k=\frac{1}{k}((1+k)-(1-k))=2$, ce qui reste vrai pour $k=0$. Si
$k>1$, $I_k=\frac{1}{k}((1+k)-(k-1))=\frac{2}{k}$, et enfin, si $k<-1$, $I_k=\frac{-2}{k}$. En résumé,

\begin{center}
\shadowbox{
$\mbox{Si}\;k\in]-1,1[,\;I_k=2\;\mbox{et si}\;k\in]-\infty,-1[\cup]1,+\infty[,\;I_k=\frac{2}{|k|}.$
}
\end{center}
\end{enumerate}
\fincorrection
\correction{005080}
\begin{enumerate}
 \item  Soient $n\in\Nn$ et $x\in\Rr$. Posons $S_n=\sum_{k=0}^{n}\cos(kx)$ et $S_n'=\sum_{k=0}^{n}\sin(kx)$.
\begin{itemize}
\item[\textbf{1ère solution.}]

$$S_n+iS_n'=\sum_{k=0}^{n}(\cos(kx)+i\sin(kx))=\sum_{k=0}^{n}e^{ikx}=\sum_{k=0}^{n}(e^{ix})^k.$$
Maintenant, $e^{ix}=1\Leftrightarrow x\in2\pi\Zz$. Donc,

\begin{itemize}
\item[\textbf{1er cas.}] Si $x\in2\pi\Zz$, on a immédiatement $S_n=n+1$ et $S_n'=0$.
\item[\textbf{2ème cas.}] Si x$\notin2\pi\Zz$,

\begin{align*}
S_n+iS_n'&=\frac{1-e^{i(n+1)x}}{1-e^{ix}}=\frac{e^{i(n+1)x/2}}{e^{ix/2}}\frac{e^{-i(n+1)x/2}-e^{i(n+1)x/2}}
{e^{-i(n+1)x/2}+e^{i(n+1)x/2}}=e^{inx/2}\frac{-2i\sin\frac{(n+1)x}{2}}{-2i\sin\frac{x}{2}}\\
 &=e^{inx/2}\frac{\sin\frac{(n+1)x}{2}}{\sin\frac{x}{2}}
\end{align*}
Par identification des parties réelles et imaginaires, on obtient

\begin{center}
\shadowbox{
$
\sum_{k=0}^{n}\cos(kx)=
\left\{
\begin{array}{l}
\frac{\cos\frac{nx}{2}\sin\frac{(n+1)x}{2}}{\sin\frac{x}{2}}\;\mbox{si}\;x\notin2\pi\Zz\\
n+1\;\mbox{si}\;x\in2\pi\Zz
\end{array}
\right.
\;\mbox{et}\;\sum_{k=0}^{n}\sin(kx)=
\left\{
\begin{array}{l}
\frac{\sin\frac{nx}{2}\sin\frac{(n+1)x}{2}}{\sin\frac{x}{2}}\;\mbox{si}\;x\notin2\pi\Zz\\
0\;\mbox{si}\;x\in2\pi\Zz
\end{array}
\right.
$
}
\end{center}

\end{itemize}

\item[\textbf{2ème solution.}]
\begin{align*}
2\sin\frac{x}{2}\sum_{k=0}^{n}\cos(kx)&=\sum_{k=0}^{n}2\sin\frac{x}{2}\cos(kx)=\sum_{k=0}^{n}(\sin(k+\frac{1}{2})x-\sin(
k-\frac{1}{2})x)\\
 &=\left(\sin\frac{x}{2}-\sin\frac{-x}{2}\right)+\left(\sin\frac{3x}{2}-\sin\frac{x}{2}\right)+\ldots+\left(\sin\frac{(2n-1)x}{2}-\sin\frac{(2n-3)x}{2}\right)\\
  &+\left(\sin\frac{(2n+1)x}{2}-\sin\frac{(2n-1)x}{2}\right)\\
 &=\sin\frac{(2n+1)x}{2}+\sin\frac{x}{2}=2\sin\frac{(n+1)x}{2}\cos\frac{nx}{2}
\end{align*}
et donc, si $x\notin2\pi\Zz$,...

\end{itemize}

 \item  Soient $n\in\Nn$ et $x\in\Rr$. Posons $S_n=\sum_{k=0}^{n}\cos^2(kx)$ et $S_n'=\sum_{k=0}^{n}\sin^2(kx)$. On
a~:

$$S_n+S_n'=\sum_{k=0}^{n}(\cos^2(kx)+\sin^2(kx))=\sum_{k=0}^{n}1=n+1,$$

et

$$S_n-S_n'=\sum_{k=0}^{n}(\cos^2(kx)-\sin^2(kx))=\sum_{k=0}^{n}\cos(2kx).$$

D'après 1), si $x\in\pi\Zz$, on trouve immédiatement,

$$\sum_{k=0}^{n}\cos^2(kx)=n+1\;\mbox{et}\;\sum_{k=0}^{n}\sin^2(kx)=0,$$

et si $x\notin\pi\Zz$,

$$S_n+S_n'=n+1\;\mbox{et}\;S_n-S_n'=\frac{\cos(nx)\sin(n+1)x}{\sin x},$$

de sorte que

$$S_n=\frac{1}{2}\left(n+1+\frac{\cos(nx)\sin(n+1)x}{\sin
x}\right)\;\mbox{et}\;S_n'=\frac{1}{2}\left(n+1-\frac{\cos(nx)\sin(n+1)x}{\sin x}\right).$$

 \item  Soient $n\in\Nn$ et $x\in\Rr$.

\begin{align*}
\left(\sum_{k=0}^{n}C_n^k\cos(kx)\right)+i\left(\sum_{k=0}^{n}C_n^k\sin(kx)\right)&=\sum_{k=0}^{n}C_n^ke^{ikx}
=\sum_{k=0}^{n}C_n^k(e^{ix})^k1^{n-k}\\
 &=(1+e^{ix})^n=(e^{ix/2}+e^{-ix/2})^ne^{inx/2}=2^n\cos^n\left(\frac{x}{2}\right)\left(\cos\frac{nx}{2}+i\sin\frac{nx}{2}\right).
\end{align*}
Par identification des parties réelles et imaginaires, on obtient alors

\begin{center}
\shadowbox{
$\sum_{k=0}^{n}C_n^k\cos(kx)=2^n\cos^n\left(\frac{x}{2}\right)\cos\left(\frac{nx}{2}\right)\;\mbox{et}\;\sum_{k=0}^{n}C_n^k\sin(kx)=
2^n\cos^n\left(\frac{x}{2}\right)\sin\left(\frac{nx}{2}\right).$
}
\end{center}
\end{enumerate}
\fincorrection
\correction{005081}
\begin{align*}
\left\{
\begin{array}{l}
\cos a+\cos b+\cos c=0\\
\sin a+\sin b+\sin c=0
\end{array}
\right.&\Leftrightarrow(\cos a+\cos b+\cos c)+i(\sin a+\sin b+\sin c)=0\Leftrightarrow e^{ia}+e^{ib}+e^{ic}=0\\
 &\Rightarrow|e^{ia}+e^{ib}|=|-e^{ic}|=1\Leftrightarrow|e^{ia/2}e^{ib/2}(e^{i(a-b)/2}+e^{-i(a-b)/2})|=1\\
 &\Leftrightarrow|\cos\frac{a-b}{2}|=\frac{1}{2}\\
 &\Leftrightarrow\frac{a-b}{2}\in\left(\frac{\pi}{3}+\pi\Zz\right)\cup\left(-\frac{\pi}{3}+\pi\Zz\right)\Leftrightarrow
a-b\in\left(\frac{2\pi}{3}+2\pi\Zz\right)\cup\left(-\frac{2\pi}{3}+2\pi\Zz\right)\\
 &\Leftrightarrow\exists k\in\Zz,\;\exists\varepsilon\in\{-1,1\}/\;b=a+\varepsilon\frac{2\pi}{3}+2k\pi.
\end{align*}

Par suite, nécessairement, $e^{ib}=je^{ia}$ ou $e^{ib}=j^2e^{ia}$. Réciproquement, si $e^{ib}=je^{ia}$ ou
encore $b=a+\frac{2\pi}{3}+2k\pi$,

$$e^{ia}+e^{ib}+e^{ic}=0\Leftrightarrow e^{ic}=-(e^{ia}+e^{ib})=-(1+j)e^{ia}=j^2e^{ia}\Leftrightarrow\exists
k'\in\Zz/\;c=a-\frac{2\pi}{3}+2k'\pi,$$

et si $e^{ib}=j^2e^{ia}$ ou encore
$b=a-\frac{2\pi}{3}+2k\pi$,

$$e^{ia}+e^{ib}+e^{ic}=0\Leftrightarrow e^{ic}=-(e^{ia}+e^{ib})=-(1+j^2)e^{ia}=je^{ia}\Leftrightarrow\exists
k'\in\Zz/\;c=a+\frac{2\pi}{3}+2k'\pi.$$

\begin{center}
\shadowbox{
$\mathcal{S}=\{(a,a+\varepsilon\frac{2\pi}{3}+2k\pi,a-\varepsilon\frac{2\pi}{3}+2k'\pi),\;a\in\Rr,\;
\varepsilon\in\{-1,1\},\;(k,k')\in\Zz^2\}.$
}
\end{center}
\fincorrection
\correction{005082}
\begin{align*}
\cos^4\frac{\pi}{8}+\cos^4\frac{3\pi}{8}+\cos^4\frac{5\pi}{8}+\cos^4\frac{7\pi}{8}&
=2(\cos^4\frac{\pi}{8}+\cos^4\frac{3\pi}{8})=2(\cos^4\frac{\pi}{8}+\sin^4\frac{\pi}{8})\\
 &=2\left((\cos^2\frac{\pi}{8}+\sin^2\frac{\pi}{8})^2-2\cos^2\frac{\pi}{8}\sin^2\frac{\pi}{8}\right)
=2\left(1-\frac{1}{2}\sin^2\frac{\pi}{4}\right)\\
 &=2(1-\frac{1}{4})=\frac{3}{2}
\end{align*}
\fincorrection
\correction{005083}
\begin{enumerate}
 \item 

\begin{align*}
\cos(3x)=\sin(2x)&\Leftrightarrow\cos(3x)=\cos(\frac{\pi}{2}-2x)\Leftrightarrow(\exists k\in\Zz/\;3x=\frac{\pi}{2}-2x+2k\pi)\;\mbox{ou}\;
(\exists k\in\Zz/\;3x=-\frac{\pi}{2}+2x+2k\pi)\\
 &\Leftrightarrow(\exists k\in\Zz/\;x=\frac{\pi}{10}+\frac{2k\pi}{5})\;\mbox{ou}\;(\exists k\in\Zz/\;x=-\frac{\pi}{2}+2k\pi)
\end{align*}

\begin{center}
\shadowbox{
$\mathcal{S}_{[0,2\pi]}=\left\{\frac{\pi}{10},\frac{\pi}{2},\frac{9\pi}{10},\frac{13\pi}{10},\frac{3\pi}{2},
\frac{17\pi}{10}\right\}.$
}
\end{center}

 \item  $\cos(3x)=\Re(e^{3ix})=\Re((\cos x+i\sin x)^3)=\cos^3x-3\cos x\sin^2x=\cos^3x-3\cos
x(1-\cos^2x)=4\cos^3x-3\cos x$.

\begin{center}
\shadowbox{
$\forall x\in\Rr,\;\cos(3x)=4\cos^3x-3\cos x.$
}
\end{center}
Par suite,

\begin{align*}
\cos(3x)=\sin(2x)&\Leftrightarrow4\cos^3x-3\cos x=2\sin x\cos x\Leftrightarrow\cos x(4\cos^2x-3-2\sin x)=0\\
 &\Leftrightarrow\cos x(-4\sin^2x-2\sin x+1)=0\Leftrightarrow(\cos x=0)\;\mbox{ou}\;(4\sin^2x+2\sin x-1=0).
\end{align*}
D'après 1), l'équation $4\sin^2x+2\sin x-1=0$ admet entre autre pour solutions $\frac{\pi}{10}$ et $\frac{13\pi}{10}$
(car, dans chacun des deux cas, $\cos x\neq0$), ou encore, l'équation $4X^2+2X-1=0$ admet pour solutions les deux
nombres \textbf{distincts} $X_1=\sin\frac{\pi}{10}$ et $X_2=\sin\frac{13\pi}{10}$, qui sont donc les deux solutions de
cette équation. Puisque $X_1>0$ et que $X_2<0$, on obtient

$$X_1=\frac{-1+\sqrt{5}}{4}\;\mbox{et}\;X_2=\frac{-1-\sqrt{5}}{4}.$$
Donc, (puisque $\sin\frac{13\pi}{10}=-\sin\frac{3\pi}{10}$),

\begin{center}
\shadowbox{
$\sin\frac{\pi}{10}=\frac{-1+\sqrt{5}}{4}\;\mbox{et}\;\sin\frac{3\pi}{10}=\frac{1+\sqrt{5}}{4}.$
}
\end{center}
Ensuite, $\sin\frac{3\pi}{10}=\cos\left(\frac{\pi}{2}-\frac{3\pi}{10}\right)=\cos\frac{\pi}{5}$, et donc

\begin{center}
\shadowbox{
$\cos\frac{\pi}{5}=\frac{1+\sqrt{5}}{4}.$
}
\end{center}
Puis

\begin{center}
\shadowbox{$\cos\frac{\pi}{10}=\sqrt{1-\sin^2\frac{\pi}{10}}=\frac{1}{4}\sqrt{10+2\sqrt{5}}$}
\end{center}
et de même 

\begin{center}
\shadowbox{$\sin\frac{\pi}{5}=\frac{1}{4}\sqrt{10-2\sqrt{5}}=\cos\frac{3\pi}{10}$}.
\end{center}
\end{enumerate}

\fincorrection
\correction{005162}
Soit $n\in\Nn^*$. Puisque, pour tout entier $k$, $|\cos k\;|\in[0,1]$, on a alors

\begin{align*}
\sum_{k=1}^{n}|\cos k\;|&\geq\sum_{k=1}^{n}\cos^2k=\sum_{k=1}^{n}\frac{1}{2}(1+\cos(2k))=\frac{n}{2}+\frac{1}{2}
\Re(\sum_{k=1}^{n}e^{2ik})
=\frac{n}{2}+\frac{1}{2}\Re(e^{2i}\frac{1-e^{2in}}{1-e^{2i}})\\
 &=\frac{n}{2}+\frac{1}{2}\Re(e^{i(n-1+2)}
 \frac{\sin n}{\sin 1})
=\frac{n}{2}+\frac{\cos(n+1)\sin n}{2\sin 1}\geq\frac{n}{2}-\frac{1}{2\sin1}.
\end{align*}

Maintenant, $\frac{1}{2\sin1}=0,594...$. Par suite, pour $n\geq3$,
$\frac{1}{2\sin1}\leq0,75=\frac{3}{4}\leq\frac{n}{4}$, et donc

$$\frac{n}{2}-\frac{1}{2\sin1}\geq\frac{n}{2}-\frac{n}{4}=\frac{n}{4}.$$

Enfin, si $n=1$, $|\cos1|=0.5...\geq0.25=\frac{1}{4}$ et si $n=2$, $|\cos1|+|\cos2|=0.9...\geq0.5=\frac{2}{4}$.
Finalement,
\begin{center}
\shadowbox{
$\forall n\in\Nn^*,\;\sum_{k=1}^{n}|\cos k|\geq\frac{n}{4}.$
}
\end{center}
\fincorrection
\nocorrection
\nocorrection
\nocorrection
\nocorrection
\nocorrection
\nocorrection
\correction{000096}
\begin{enumerate}
\item Soit $\alpha,\beta\in\Z[i]$. Notons $\alpha=a+ib$ et $\beta=c+id$ avec $a,b,c,d\in\Z$. Alors
$\alpha+\beta=(a+c)+i(b+d)$ et $a+c\in\Z$, $b+d\in\Z$ donc
$\alpha+\beta\in\Z[i]$. De m\^eme, $\alpha\beta=(ac-bd)+i(ad+bc)$ et
$ac-bd\in\Z$, $ad+bc\in\Z$ donc $\alpha\beta\in\Z[i]$.

\item
Soit $\alpha\in\Z[i]$ inversible. Il existe donc $\beta\in\Z[i]$
tel que $\alpha\beta=1$. Ainsi, $\alpha\neq0$ et
$\frac{1}{\alpha}\in\Z[i]$. Remarquons que tout \'el\'ement non
nul de $\Z[i]$ est de module sup\'erieur ou \'egal \`a 1: en effet
$\forall z\in\C, |z|\geq \sup(|\mathop{\mathrm{Re}}\nolimits(z)|,|\mathop{\mathrm{Im}}\nolimits(z)|)$ et si
$z\in\Z[i]\setminus\{0\}$, $\sup(|\mathop{\mathrm{Re}}\nolimits(z)|,|\mathop{\mathrm{Im}}\nolimits(z)|)\geq 1$. Si
$|\alpha|\neq 1$ alors $|\alpha|>1$ et $|1/\alpha|<1$. On en
d\'eduit $1/\alpha=0$ ce qui est impossible. Ainsi $|\alpha|=1$,
ce qui implique $\alpha\in\{1,-1,i,-i\}$.

R\'eciproquement,
$1^{-1}=1\in\Z[i],(-1)^{-1}=-1\in\Z[i],i^{-1}=-i\in\Z[i],(-i)^{-1}=i\in\Z[i]$.
Les \'el\'ements inversibles de $\Z[i]$ sont donc $1,-1,i$ et
$-i$.

\item
Soit $\omega\in\C$. Notons $\omega=x+iy$ avec $x,y\in\R$. soit
$E(x)$ la partie enti\`ere de $x$, i.e. le plus grand entier
inf\'erieur ou \'egal \`a $x$: $E(x)\leq x<E(x)+1$. Si $x\leq E(x)+1/2$, 
notons $n_{x}=E(x)$, et si $x> E(x)+1/2$, notons
$n_{x}=E(x)+1$. $n_{x}$ est le, ou l'un des s'il y en a deux,
nombre entier le plus proche de $x$: $|x-n_{x}|\leq1/2$. Notons
$n_{y}$ l'entier associ\'e de la m\^eme mani\`ere \`a $y$. Soit
alors $\alpha=n_{x}+i \cdot n_{y}$. $z\in\Z[i]$ et
$|\omega-\alpha|^2=(x-n_{x})^2+(y-n_{y})^2\leq 1/4+1/4=1/2$. Donc
$|\omega-\alpha|<1$.

% $$
% \includegraphics[64mm,44mm]{gauss.wmf}
% $$
\item
Soit $\alpha,\beta\in\Z[i]$, avec $\beta\neq0$. Soit alors
$q\in\Z[i]$ tel que $|\frac{\alpha}{\beta}-q|<1$. Soit
$r=\alpha-\beta q$. Comme $\alpha\in\Z[i]$ et $\beta q\in\Z[i]$,
$r\in\Z[i]$. De plus
$|\frac{r}{\beta}|=|\frac{\alpha}{\beta}-q|<1$ donc $|r|<|\beta|$.
\end{enumerate}
\fincorrection
\nocorrection
\nocorrection
\nocorrection
\nocorrection
\nocorrection
\nocorrection
\correction{000103}
\begin{enumerate}
\item \`A permutation pr\`es, $x=-2$, $y=-2j$ et $z=-2j^2$ ($j$ d\'esigne la
racine cubique de l'unit\'e ${-1+i\sqrt3\over 2}$).
\item \`A permutation pr\`es, $x=1$, $y={3+i\sqrt3\over 2}$ et
$z={3-i\sqrt3\over 2}$.
\end{enumerate}
\fincorrection
\correction{005132}
Soit $z\in\Cc$ tel que $|z|>1$.

$$|1+z+...+z^{n-1}|\leq1+|z|+|z|^2+...+|z|^{n-1}<|z|^n+|z|^n+...+|z|^n=n|z|^n=|nz^n|,$$
et en particulier, $1+z+...+z^{n-1}\neq nz^n$. Donc, si $1+z+...+z^{n-1}-nz^n=0$, alors $|z|\leq1$.
\fincorrection
\correction{005136}

\begin{enumerate}
 \item  Soit $z\in\Cc$. $\sh z$ et $\ch z$ sont définis et donc, $\tanh z$ existe si et seulement si $\ch z\neq0$. Or,

$$\ch z=0\Leftrightarrow e^z+e^{-z}=0\Leftrightarrow e^{2z}=-1\Leftrightarrow e^{2z}=e^{i\pi}\Leftrightarrow2z\in i\pi+2i\pi\Zz\Leftrightarrow z\in i\left(\frac{\pi}{2}+\pi\Zz\right).$$
$\tanh z$ existe si et seulement si $z\notin i\left(\frac{\pi}{2}+\pi\Zz\right)$.
 \item  Soit $z\notin i\left(\frac{\pi}{2}+\pi\Zz\right)$.

$$\tanh z=0\Leftrightarrow\sh z=0\Leftrightarrow e^z=e^{-z}\Leftrightarrow e^{2z}=1\Leftrightarrow 2z\in2i\pi\Zz\Leftrightarrow z\in i\pi\Zz.$$
Comme $i\left(\frac{\pi}{2}+\pi\Zz\right)\cap i\pi\Zz=\varnothing$, $\tanh z=0$ si et seulement si $z\in i\pi\Zz$.
 \item  Soit $z\notin i\left(\frac{\pi}{2}+\pi\Zz\right)$. Posons $z=x+iy$ où $(x,y)\in\Rr^2$.

\begin{align*}
|\tanh z|<1&\Leftrightarrow|e^z-e^{-z}|^2<|e^z+e^{-z}|^2\Leftrightarrow(e^z-e^{-z})(e^{{\bar z}}-e^{-{\bar z}})<(e^z+e^{-z})(e^{{\bar z}}+e^{-{\bar z}})\\
 &\Leftrightarrow-e^{z-{\bar z}}-e^{-(z-{\bar z})}<e^{z-{\bar z}}+e^{-(z-{\bar z})}\Leftrightarrow2(e^{2iy}+e^{-2iy})>0\\
 &\Leftrightarrow\cos(2y)>0
\end{align*}
Par suite,

$$\left\{
\begin{array}{l}
|\Im z|<\frac{\pi}{2}\\
|\tanh z|<1
\end{array}
\right.
\Leftrightarrow
\left\{
\begin{array}{l}
|y|<\frac{\pi}{2}\\
\cos(2y)>0
\end{array}
\right.
\Leftrightarrow
|y|<\frac{\pi}{4}\Leftrightarrow z\in\Delta.$$
 \item  Soit $z\in\Delta$.
D'après 1), $\tanh z$ existe et d'après 3), $|\tanh z|<1$. Donc $z\in\Delta\Rightarrow\tanh z\in U$. Ainsi, $\tanh$ est une
application de $\Delta$ dans $U$.
Soit alors $Z\in U$ et $z\in\Delta$.

$$\tanh z=Z\Leftrightarrow\frac{e^{2z}-1}{e^{2z}+1}=Z\Leftrightarrow e^{2z}=\frac{1+Z}{1-Z}.$$
Puisque $Z\neq-1$, $\frac{1+Z}{1-Z}\neq0$ et on peut poser $\frac{1+Z}{1-Z}=re^{i\theta}$ où $r\in\Rr_+^*$ et
$\theta\in]-\pi,\pi]$.

Par suite,

\begin{align*}
e^{2z}=\frac{1+Z}{1-Z}&\Leftrightarrow e^{2z}=re^{i\theta}\Leftrightarrow e^{2x}=r\;\mbox{et}\;2y\in\theta+2\pi\Zz\\
 &\Leftrightarrow x=\frac{1}{2}\ln r\;\mbox{et}\;y\in\frac{\theta}{2}+\pi\Zz
\end{align*}
Maintenant, on ne peut avoir $\theta=\pi$. Dans le cas contraire, on aurait $\frac{1+Z}{1-Z}=-r\in\Rr_-^*$ puis
$Z=\frac{r+1}{r-1}\in\Rr$. Par suite, puisque $|Z|<1$, on aurait $Z\in]-1,1[$ et donc $\frac{1+Z}{1-Z}\in\Rr_+^*$ ce
qui est une contradiction. Donc, $\theta\in]-\pi,\pi[$ puis $\frac{\theta}{2}\in]-\frac{\pi}{2},\frac{\pi}{2}[$.
Mais alors,

$$\left\{
\begin{array}{l}
\tanh z=Z\\
z\in\Delta
\end{array}
\right.
\Leftrightarrow\left\{
\begin{array}{l}
x=\frac{1}{2}\ln r\\
y=\frac{\theta}{2}\rule{0mm}{6mm}
\end{array}
\right..$$
Ainsi, tout élément $Z$ de $U$ a un et un seul antécédent $z$ dans $\Delta$ (à savoir
$z=\frac{1}{2}\ln\left|\frac{1+Z}{1-Z}\right|+\frac{i}{2}\mbox{Arg}\left(\frac{1+Z}{1-Z}\right)$ où 
$\mbox{Arg}\left(\frac{1+Z}{1-Z}\right)$ désigne l'argument de $\frac{1+Z}{1-Z}$ qui est dans $]-\pi,\pi[$).
Finalement, $\tanh$ réalise donc une bijection de $\Delta$ sur $U$.
\end{enumerate}

\fincorrection
\nocorrection
\nocorrection
\nocorrection
\nocorrection
\nocorrection
\nocorrection
\nocorrection
\nocorrection
\correction{000364}
\begin{enumerate}
  \item $A =  3X^5+4X^2+1$, $B = X^2+2X+3$, le quotient de $A$ par $B$ est
$3{X}^{3}-6{X}^{2}+3X+16$ et le reste $-47-41X$.
  \item $A = 3X^5+2X^4-X^2+1$, $B = X^3+X+2$ le quotient de $A$ par $B$ est
$ 3{X}^{2}+2X-3 $ et le reste est $ 7-9{X}^{2}-X$.
 \item $ A = X^4-X^3-X-2$, $B = X^2-2X+4$, le quotient de $A$ par $B$ est
$ {X}^{2}+X-2$ de reste $ 6-9X$.
\end{enumerate}
\fincorrection
\nocorrection
\correction{000366}
$X^{4}+X^{3}-2X+1 = (X^{2}+X+1)(2X^{2}-3X+1)+X^{3}(2-X)$.
\fincorrection
\nocorrection
\nocorrection
\nocorrection
\correction{000370}
\
\begin{enumerate}
\item On remarque que si $P$ est solution, alors $P+1=(X-1)^4A$ et par ailleurs $P-1=(X+1)^4B$, ce qui donne $1=\frac{A}{2}(X-1)^4+\frac{-B}{2}(X+1)^4$. Cherchons des polynômes $A$ et $B$ qui conviennent: pour cela, on écrit la relation de Bézout entre $(X-1)^4$ et $(X+1)^4$ qui sont premiers entre eux, et on obtient 
$$\frac{A}{2}=\frac{5}{32}X^3+\frac{5}{8}X^2+\frac{29}{32}X+\frac{1}{2}$$
$$\frac{-B}{2}=-\frac{5}{32}X^3+\frac{5}{8}X^2-\frac{29}{32}X+\frac{1}{2}$$
On a alors par construction
$$(X-1)^4A-1=2\bigg(1+(X+1)^4\frac{-B}{2}\bigg)=1+(X+1)^4B$$
et $P_0=(X-1)^4A-1$ convient. En remplaçant, on obtient après calculs :
$$P_0 = \frac{5}{16}X^7-\frac{21}{16}X^5+\frac{35}{16}X^3-\frac{35}{16}X$$


\item Si $(X-1)^4$ divise $P+1$, alors $1$ est racine de multiplicité au moins $4$
de $P+1$, et donc racine de multiplicité au moins $3$ de $P'$ : 
alors $(X-1)^3$ divise $P'$. De même $(X+1)^3$ divise $P'$. 
Comme  $(X-1)^3$ et  $(X+1)^3$ sont premiers entre eux, nécessairement $(X-1)^3(X+1)^3$ divise $P'$. 
Cherchons un polynôme de degré minimal : on remarque que les primitives de 
$$\lambda(X-1)^3(X+1)^3=\lambda(X^2-1)^3=\lambda(X^6-3X^4+3X^2-1)$$
sont de la forme $P(X)=\lambda(\frac{1}{7}X^7-\frac{3}{5}X^5+X^3-X+a)$. 
Si $P$ convient, nécessairement $1$ est racine de $P+1$ et $-1$ est racine de $P-1$, ce qui donne
$\lambda(\frac{-16}{35}+a)=-1$ et $\lambda(\frac{16}{35}+a)=1$. 
D'où $\lambda a=0$ et comme on cherche $P$ non nul, il faut $a=0$ et $\lambda=\frac{35}{16}$. 
On vérifie que 
$$P_0(X)=\frac{35}{16}(\frac{1}{7}X^7-\frac{3}{5}X^5+X^3-X)=\frac{5}{16}X^7-\frac{21}{16}X^5
+\frac{35}{16}X^3-\frac{35}{16}X$$
 est bien solution du problème: le polynôme $A=P_0+1$ admet $1$ comme racine, i.e. $A(1)=0$, 
 et sa dérivée admet $1$ comme racine triple donc $A'(1)=A''(1)=A'''(1)=0$, 
 ainsi $1$ est racine de multiplicité au moins $4$ de $A$ et donc $(X-1)^4$ divise $A=P+1$. 
 De même, $(X+1)^4$ divise $P-1$.
\end{enumerate}


Supposons que $P$ soit une solution du problème. On note toujours $P_0$ la solution particulière obtenue ci-dessus.
Alors $P+1$ et $P_0+1$ sont divisibles par $(X-1)^4$, et $P-1$ et $P_0-1$ sont divisibles par $(X+1)^4$. 
Ainsi $P-P_0=(P+1)-(P_0+1)=(P-1)-(P_0-1)$ est divisible par $(X-1)^4$ et par $(X+1)^4$. 
Comme $(X-1)^4$ et $(X+1)^4$ sont premiers entre eux, nécessairement $P-P_0$ est divisible 
par $(X-1)^4(X+1)^4$. Réciproquement, si $P=P_0+(X-1)^4(X+1)^4A$, alors $P+1$ est 
bien divisible par $(X-1)^4$ et $P-1$ est divisible par $(X+1)^4$. 

Ainsi les solutions sont exactement les polynômes de la forme
$$P_0(X)+(X-1)^4(X+1)^4A(X)$$
où $P_0$ est la solution particulière trouvée précédemment, et $A$ un polynôme quelconque.
\fincorrection
\correction{000371}
\begin{enumerate}
    \item Quotient $Q=X^3-X^2-X+1$, reste $R=X$.
    \item Quotient $Q=1-X^2-X^4$, reste $R=X^5(1+2X+X^2)$.
\end{enumerate}
\fincorrection
\nocorrection
\nocorrection
\correction{000375}
Soient $ A = X^5-7 X^4-X^2-9 X+9$, $B = X^2-5 X+4$, le quotient de
$A$ par $B$ est $ {X}^{3}-2\,{X}^{2}-14\,X-63$, le reste \'etant $
261-268\,X$.
\fincorrection
\nocorrection
\nocorrection
\correction{000378}
Le polynôme nul convient. Dans la suite on suppose que $P$ n'est pas le polynôme nul.

Notons $n = \deg P$ son degré. Comme $P'$ divise $P$, alors $P$ est non constant, 
donc $n\ge 1$. Soit $Q\in\Cc[X]$ tel que $P=P'Q$.
Puisque $\deg(P')=\deg(P)-1\ge 0$, alors $Q$ est de degré $1$.
Ainsi $Q(X)=aX+b$ avec $a\neq0$, et donc  
$$P(X)=P'(X)(aX+b)=aP'(X)(X+\frac{b}{a})$$
Donc si $z\not=\frac{-b}{a}$ et si $z$ est racine de $P$ de multiplicité $k \ge 1$, 
alors $z$ est aussi racine de $P'$ avec la même multiplicité, ce qui est impossible. 
Ainsi la seule racine possible de $P$ est $\frac{-b}{a}$.

Réciproquement, soit $P$ un polynôme avec une seule racine $z_0 \in \Cc$ : 
il existe $\lambda\not=0$, $n\ge 1$ tels que $P=\lambda(X-z_0)^n$, 
qui est bien divisible par son polynôme dérivé.
\fincorrection
\nocorrection
\nocorrection
\correction{003198}
$\Im\varphi = \{ P \in E$ tel que $X-1 \mid P \}$ (B{\'e}zout g{\'e}n{\'e}ralis{\'e}). \par
$\mathrm{Ker}\varphi = \text{vect}(X^3+X^2+X)$.
\fincorrection
\correction{003199}
\begin{enumerate}
  \item $\frac {\alpha(b-X) + \beta(X-a)}{b-a}$.
  \item $\cos n\theta + X\sin n\theta$.
\end{enumerate}
\fincorrection
\correction{003200}
$P = \lambda((X+2)(X+3)(X+4)-6)$.
\fincorrection
\correction{003201}
\begin{enumerate}
  \item $X+1$
  \item $1$
  \item $X^2-iX+1$
\end{enumerate}
\fincorrection
\correction{003202}
\begin{enumerate}
  \item $7U= X+3,\quad 7V= -X^3-3X^2+X+4$
  \item $3U= 2X^2-X+1,\quad 3V = -2X^2-X+2$
\end{enumerate}
\fincorrection
\correction{003203}
Substituer $j$ {\`a} $X$ $ \Rightarrow $
         $R = \begin{cases}(-1)^n - 2           &\text{ si }n \equiv 0 (\mathrm{mod}\, 3) \hfill\cr
                     ((-1)^{n+1}-1)(X+1)  &\text{ si }n \equiv 1 (\mathrm{mod}\, 3) \hfill\cr
                     ((-1)^n+1)X          &\text{ si }n \equiv 2 (\mathrm{mod}\, 3).\end{cases}$
\fincorrection
\correction{003204}
$n\equiv 0(\mathrm{mod}\, 6)$.
\fincorrection
\correction{003205}
\begin{enumerate}
  \item 
  \item
  \item Faire le produit.
\end{enumerate}
\fincorrection
\correction{003206}
\begin{enumerate}
  \item $(2^{50}-1)X + 2-2^{50}$.
  \item $2^{16}\bigl(X-\sqrt3\bigr)$.
  \item $192\bigl(X-\sqrt2\bigr)^2$.
  \item 
\end{enumerate}
\fincorrection
\correction{003207}
$\lambda = \mu = -1$.
\fincorrection
\correction{003208}
$-3X^3 + X^2 - X - 1$.
\fincorrection
\correction{003209}
$n = qm+r  \Rightarrow  X^n-1 \equiv X^r-1 (\mathrm{mod}\,{X^m-1})$.
On applique la m{\'e}thode des divisions euclidiennes entre $n$ et $m$
$ \Rightarrow $ \hbox{pgcd$ {}= X^{n \wedge m} - 1$.}
\fincorrection
\correction{003210}
\begin{enumerate}
  \item 
  \item R{\'e}currence.
\end{enumerate}
\fincorrection
\nocorrection
\nocorrection
\nocorrection
\correction{005323}
On prend $n\geq2$ (sinon tout est clair).

$Q=(X-e^{ia})(X-e^{-ia})$ est à racines simples si et seulement si $e^{ia}\neq e^{-ia}$ ou encore $e^{2ia}\neq 1$ ou enfin, $a\notin\pi\Zz$.

1er cas. Si $a\in\pi\Zz$ alors, $P=0=0.Q$.

2ème cas. Si $a\notin\pi\Zz$, alors 

\begin{align*}\ensuremath
P(e^{ia})&=\sin a(\cos(na)+i\sin(na))-\sin(na)(\cos a+i\sin a)+\sin((n-1)a)\\
 &=\sin((n-1)a)-(\sin(na)\cos a-\cos(na)\sin a)=0.
\end{align*}

Donc, $e^{ia}$ est racine de $P$ et de même, puisque $P$ est dans $\Rr[X]$, $e^{-ia}$ est racine de $P$. $P$ est donc divisible par $Q$.

\begin{align*}\ensuremath
P&=P-P(e^{ia})=\sin a(X^n-e^{ina})-\sin(na)(X-e^{ia})=(X-e^{ia})(\sin a\sum_{k=0}^{n-1}X^{n-1-k}e^{ika}-\sin(na))\\
 &=(X-e^{ia})S.
\end{align*}
 
Puis,

\begin{align*}\ensuremath
S&=S-S(e^{-ia})=\sin a\sum_{k=0}^{n-1}e^{ika}(X^{n-1-k}-e^{-i(n-1-k)a})=\sin a(X-e^{-ia})\sum_{k=0}^{n-2}e^{ika}(\sum_{j=0}^{n-2-k}X^{n-2-k-j}e^{-ija})\\
 &=\sin a(X-e^{-ia})\sum_{k=0}^{n-2}(\sum_{j=0}^{n-2-k}X^{n-2-k-j}e^{i(k-j)a})
=\sin a(X-e^{-ia})\sum_{l=0}^{n-2}(\sum_{k+j=l}^{}e^{i(k-j)a})X^{n-2-l}\\
 &=\sin a(X-e^{-ia})\sum_{l=0}^{n-2}(\sum_{k=0}^{l}e^{i(2k-l)a})X^{n-2-l}
\end{align*}
 
Maintenant,

$$\sum_{k=0}^{l}e^{i(2k-l)}a=e^{-ila}\frac{1-e^{2i(l+1)a}}{1-e^{2ia}}=\frac{\sin((l+1)a)}{\sin a}.$$

Donc

$$S=\sin a(X-e^{-ia})\sum_{l=0}^{n-2}\frac{\sin((l+1)a)}{\sin a}X^{n-2-l}=(X-e^{-ia})\sum_{l=0}^{n-2}\sin((l+1)a)X^{n-2-l},$$

et finalement 

$$P=(X-e^{ia})(X-e^{-ia})\sum_{k=0}^{n-2}\sin((k+1)a)X^{n-2-k}=(X^2-2X\cos a+1)\sum_{k=0}^{n-2}\sin((k+1)a).$$
\fincorrection
\correction{006955}
\
\begin{enumerate}
\item \begin{enumerate}
\item $3X^5+4X^2+1=(X^2+2X+3)(3X^3-6X^2+3X+16)-41X-47$
\item $3X^5+2X^4-X^2+1=(X^3+X+2)(3X^2+2X-3)-9X^2-X+7$
\item $X^4-X^3+X-2=(X^2-2X+4)(X^2+X-2)-7X+6$
\item $X^5-7X^4-X^2-9X+9$\\
\ \ \ $=(X^2-5X+4)(X^3-2X^2-14X-63)-268X+261$
\end{enumerate}
\item \begin{enumerate}
\item $1-2X+X^3+X^4=(1+2X+X^2)(1-4X+7X^2)+X^3(-9-6X)$
\item $1+X^3-2X^4+X^6=(1+X^2+X^3)(1-X^2-X^4)+X^5(1+2X+X^2)$
\end{enumerate}
\end{enumerate}
\fincorrection
\correction{006956}
La division euclidienne de $A=X^4+aX^2+bX+c$ par $B=X^2+X+1$ donne
$$X^4+aX^2+bX+c=(X^2+X+1)(X^2-X+a)+(b-a+1)X+c-a$$
Or $A$ est divisible par $B$ si et seulement si le reste 
$R=(b-a+1)X+c-a$ est le polynôme nul, 
c'est-à-dire si et seulement si $b-a+1=0$ et $c-a=0$.
\fincorrection
\correction{000379}
\begin{enumerate}
  \item $\pgcd(X^3-X^2-X-2,X^5-2 X^4+X^2-X-2) = X-2$.
  \item $\pgcd(X^4+X^3-2 X+1,X^3+X+1) = 1$.
\end{enumerate}
\fincorrection
\correction{000380}
\begin{enumerate}
  \item $\pgcd(X^5+3 X^4+X^3+X^2+3 X+1,X^4+2 X^3+X+2)=X^3+1$.
  \item $\pgcd(X^4+X^3-3 X^2-4 X-1,X^3+X^2-X-1)= X+1$
 \item  $\pgcd(X^5+5 X^4+9 X^3+7 X^2+5 X+1,X^4+2 X^3+2 X^2+X+1)= 1$.
\end{enumerate}
\fincorrection
\nocorrection
\nocorrection
\nocorrection
\nocorrection
\nocorrection
\nocorrection
\correction{000387}
\begin{enumerate}
 
    \item $D = X^2+3X+2 =
A ({1\over18} X- {1\over6})+B (-{1\over18} X^2+{1\over9} X+{5\over18})$.
    \item $D = 1 = A(-X^3)+B(X^5+X^3+X+1)$.
\end{enumerate}
\fincorrection
\nocorrection
\nocorrection
\nocorrection
\nocorrection
\nocorrection
\nocorrection
\nocorrection
\nocorrection
\correction{005317}
$X^6-7X^4+8X^3-7X+7=(X^6+8X^3+7)-(7X^4+7X)=(X^3+1)(X^3+7)-7X(X^3+1)=(X^3+1)(X^3-7X+7)$ et $3X^5-7X^3+3X^2-7=3X^2(X^3+1)-7(X^3+1)=(X^3+1)(3X^2-7)$. Donc,

$$(X^6-7X^4+8X^3-7X+7)\wedge(3X^5-7X^3+3X^2-7)=(X^3+1)((X^3-7X+7)\wedge(3X^2-7)).$$

Maintenant, pour $\varepsilon\in\{-1,1\}$, $(\varepsilon\sqrt{\frac{7}{3}})^3-7(\varepsilon\sqrt{\frac{7}{3}})+7=-(\varepsilon\frac{14}{3}\sqrt{\frac{7}{3}})+7\neq0$.

Les polynômes $(X^3-7X+7)$ et $(3X^2-7)$ n'ont pas de racines communes dans $\Cc$ et sont donc premiers entre eux. Donc, $(X^6-7X^4+8X^3-7X+7)\wedge(3X^5-7X^3+3X^2-7)=X^3+1$.
\fincorrection
\correction{006957}
\
\begin{enumerate}
\item L'algorithme d'Euclide permet de calculer le pgcd par 
une suite de divisions euclidiennes.

\begin{enumerate}
\item $X^5-2X^4+X^2-X-2=(X^3-X^2-X-2)(X^2-X)+2X^2-3X-2$ 

puis $X^3-X^2-X-2=(2X^2-3X-2)(\frac{1}{2}X+\frac{1}{4})+\frac{3}{4}X-\frac{3}{2}$ 

puis $2X^2-3X-2=(\frac{3}{4}X-\frac{3}{2})(\frac{8}{3}X+\frac{4}{3})$


Le pgcd est le dernier reste non nul, 
divisé par son coefficient dominant: 
$$\pgcd(X^3-X^2-X-2,X^5-2X^4+X^2-X-2)=X-2$$


\item $X^4+X^3-2X+1=(X^3+X+1)(X+1)-X^2-4X$

puis $X^3+X+1=(-X^2-4X)(-X+4)+17X+1$

$$\begin{array}{rl}\text{donc}\ \pgcd&(X^4+X^3-2X+1,X^3+X+1)\\
&=\pgcd(-X^2-4X,17X+1)=1\end{array}$$

car $-X^2-4X$ et $17X+1$ n'ont pas de racine (même complexe) commune.


\item $X^5+3X^4+X^3+X^2+3X+1=(X^4+2X^3+X+2)(X+1)-X^3-1$

puis $X^4+2X^3+X+2=(-X^3-1)(-X-2)+2X^3+2$

$$\pgcd(X^5+3X^4+X^3+X^2+3X+1,X^4+2X^3+X+2)=X^3+1$$


\item $nX^{n+1}-(n+1)X^n+1$

\ \ \ \ \ \ \ \ \ \ \ $=(X^n-nX+n-1)(nX-(n+1))+n^2(X-1)^2$ 

Si $n=1$ alors $X^n-nX+n-1=0$ et le $\pgcd$ vaut $(X-1)^2$.
On constate que $1$ est racine de $X^n-nX+n-1$, 
et on trouve $X^n-nX+n-1=(X-1)(X^{n-1}+X^{n-2}+\cdots+X^2+X-(n-1))$.

Si $n\ge 2$: $1$ est racine de $X^{n-1}+X^{n-2}+\cdots+X^2+X-(n-1)$ et on trouve 

$X^{n-1}+X^{n-2}+\cdots+X^2+X-(n-1)$

\ \ \ \ \ \ \ $=(X-1)(X^{n-2}+2X^{n-3}+\cdots+(n-1)X^2+nX+(n+1))$, donc finalement 
$(X-1)^2$ divise $X^n-nX+n-1$ (on pourrait aussi remarquer 
que $1$ est racine de multiplicité au moins deux de 
$X^n-nX+n-1$, puisqu'il est racine de ce polynôme et de sa dérivée). 
Ainsi 
$$\text{si}\ n\ge 2,\ \pgcd(nX^{n+1}-(n+1)X^n+1,X^n-nX+n-1)=(X-1)^2$$


\end{enumerate}


\item \begin{enumerate}
\item $A=X^5+3X^4+2X^3-X^2-3X-2$ et $B=X^4+2X^3+2X^2+7X+6$

donc $A=BQ_1+R_1$ avec $Q_1=X+1$, $R_1=-2X^3-10X^2-16X-8$

puis $B=R_1Q_2+R_2$ avec $Q_2=-\frac{1}{2}X+\frac{3}{2}$ et $R_2=9X^2+27X+18$

et enfin $R_1=R_2Q_3$ avec $Q_3=-\frac{2}{9}X-\frac{4}{9}$

Donc $D=X^2+3X+2$, et on obtient
$$9D=B-R_1Q_2=B-(A-BQ_1)Q_2=-AQ_2+B(1+Q_1Q_2)$$
soit 
$$\left\{\begin{array}{l}U=\frac{1}{9}(-Q_2)=\frac{1}{18}X-\frac{1}{6}\\
V=\frac{1}{9}(1+Q_1Q2)=-\frac{1}{18}X^2+\frac{1}{9}X+\frac{5}{18}
\end{array}\right.$$

\item On a $A=BQ_1+R_1$ avec $Q_1=X^2+1$, $R_1=X^2-X-1$

puis $B=R_1Q_2+R_2$ avec $Q_2=X^2-X+1$ et $R_2=-X+2$

et enfin $R_1=R_2Q_3+R_3$ avec $Q_3=-X-1$ et $R_3=1$

Donc $D=1$, et on obtient
\begin{eqnarray*}
1&=&R_1-R_2Q_3=R_1-(B-R_1Q_2)Q_3=R_1(1+Q_2Q_3)-BQ_3\\
 &=&(A-BQ_1)(1+Q_2Q_3)-BQ_3\\
 &=&A(1+Q_2Q_3)-B(Q_1(1+Q_2Q_3)+Q_3)
\end{eqnarray*}
soit 
$$\left\{\begin{array}{l}
U=1+Q_2Q_3=-X^3\\
V=-Q_1(1+Q_2Q_3)-Q_3=1+X+X^3+X^5\end{array}\right.$$
\end{enumerate}
\end{enumerate}
\fincorrection
\correction{006958}
\
\begin{enumerate}
\item Lorsqu'on effectue la division euclidienne $A=BQ+R$, 
les coefficients de $Q$ sont obtenus par des opérations 
élémentaires (multiplication, division, addition) à partir 
des coefficients de $A$ et $B$ : ils restent donc dans $\Qq$. 
De plus, $R=A-BQ$ est alors encore à coefficients rationnels. 

Alors $\pgcd(A,B)=\pgcd(B,R)$ et pour l'obtenir, on fait 
la division euclidienne de $B$ par $R$ (dont le quotient 
et le reste sont encore à coefficients dans $\Qq$), puis on 
recommence... Le pgcd est le dernier reste non nul, c'est 
donc encore un polynôme à coefficients rationnels.

\item Notons $P_1=\pgcd(P,P')$: comme $P$ est à coefficients 
rationnels, $P'$ aussi et donc $P_1$ aussi. 
Or $P_1(X)=(X-a)^{p-1}(X-b)^{q-1}(X-c)^{r-1}$. En itérant 
le processus, on obtient que $P_{r-1}(X)=(X-c)$ est à 
coefficients rationnels, donc $c\in\Qq$.

On remonte alors les étapes: $P_{q-1}(X)=(X-b)(X-c)^{r-q+1}$ 
est à coefficients rationnels, et $X-b$ aussi en tant que 
quotient de $P_{q-1}$ par le polynôme à coefficients rationnels 
$(X-c)^{r-q+1}$, donc $b\in\Qq$. De même, en considérant 
$P_{p-1}$, on obtient $a\in\Qq$.
\end{enumerate}
\fincorrection\nocorrection
\nocorrection
\nocorrection
\nocorrection
\nocorrection
\correction{000401}
$$\left ({x}^{2}+\sqrt {2}x+1\right )\left ({x}^{2}-\sqrt {2}x+1\right )$$
\fincorrection
\nocorrection
\nocorrection
\nocorrection
\nocorrection
\nocorrection
\nocorrection
\nocorrection
\correction{000409}
 L'ordre de multiplicit\'e est~2.
\fincorrection
\correction{000410}
Soit $x\in\R$ ; $x$ est une racine multiple de $P$ si et seulement si $P(x)=0$ et $P'(x)=0$:
$$
\begin{array}{rcl}
P(x)=P'(x) 0 
&\iff& \left\{\begin{array}{l}(x+1)^7-x^7-a=0\\7(x+1)^6-7x^6=0\end{array}\right.\\
&\iff& \left\{\begin{array}{l}(x+1)x^6-x^7-a=0\qquad \text{ en utilisant la deuxième équation}\\(x+1)^6=x^6\end{array}\right.\\
&\iff& \left\{\begin{array}{l}x^6=a\\(x+1)^3=\pm x^3 \qquad \text{ en prenant la racine carrée} \end{array}\right.\\
&\iff& \left\{\begin{array}{l}x^6=a\\x+1=\pm x \qquad \text{ en prenant la racine cubique} \end{array}\right.\\
\end{array}$$
qui admet une solution ($x=-\frac{1}{2}$) si et seulement si $a=\frac{1}{64}$.
\fincorrection
\nocorrection
\correction{000412}
\begin{enumerate}
    \item $\left \lbrace\begin{array}{rcl}
  X^3-3 & = &(X-\sqrt[3]{3})(X^2+\sqrt[3]{3}\, X+\sqrt[3]{9})\\
  & = & (X-\sqrt[3]{3}) 
         (X+{\sqrt[3]{3}\over2}- i{\sqrt3\sqrt[3]{3}\over2})
         (X+{\sqrt[3]{3}\over2}+ i{\sqrt3\sqrt[3]{3}\over2}).
\end{array}\right.$

    \item $\left\lbrace\begin{array}{rcl}
  X^{12}-1 & = & \textstyle(X-1) (X+1) (X^2+1) (X^2-X+1) (X^2+X+1)\times{}\\
           &   & \qquad\textstyle(X^2-\sqrt3\,X+1) (X^2+\sqrt3\,X+1)\\
           & = & \textstyle(X-1) (X+1) (X- i) (X+ i)\times{}\\
           &   & \qquad\textstyle\bigl(X-{1+ i\sqrt3\over2}\bigr)
                 \bigl(X-{1- i\sqrt3\over2}\bigr)
                 \bigl(X-{-1+ i\sqrt3\over2}\bigr) 
                 \bigl(X-{-1- i\sqrt3\over2}\bigr)\times{}\\
           &   & \qquad\textstyle\bigl(X-{\sqrt3+ i\over2}\bigr) 
                 \bigl(X-{\sqrt3- i\over2}\bigr)
                 \bigl(X-{-\sqrt3+ i\over2}\bigr) 
                 \bigl(X-{-\sqrt3- i\over2}\bigr).
\end{array}\right.$
\end{enumerate}
\fincorrection
\nocorrection
\nocorrection
\nocorrection
\nocorrection
\nocorrection
\nocorrection
\nocorrection
\nocorrection
\nocorrection
\nocorrection
\correction{000423}
\begin{enumerate}
  \item   $X^6+1 = -\left( {X}^{2}+1 \right) \left( {X}^{2}+X\sqrt {3}+1 \right) \left(-{X}^{2}+X\sqrt {3}-1 \right)$.
  \item $X^9+X^6+X^3+1= -\left( {X}^{2}+1 \right) \left({X}^{2}-X+1 \right) \left( {X}^{2}+X \sqrt {3}+1 \right) \left(-{X}^{2}+X\sqrt {3}-1 \right)
\left( X+1 \right)$.
\end{enumerate}
\fincorrection
\correction{003214}
\begin{enumerate}
  \item $\frac n{2^{n-1}}$.
  \item $\frac {\sin(n\theta)}{2^{n-1}}$.
  \item $-(-n)^n$.
\end{enumerate}
\fincorrection
\correction{003215}
$\omega ^{2k}-2\omega ^k \cos \theta +1 = (\omega^k-e^{i\theta})(\omega^k-e^{-i\theta})$
et $\prod_{k=0}^{n-1}(\omega^k-x) = (-1)^n(x^n-1)$.
\fincorrection
\nocorrection
\correction{003217}
$X^{2n}-2X^n\cos n\theta+1 = (X^n-e^{in\theta})(X^n-e^{-in\theta})$.

$Q$
$=\Bigl(\sum_{k=0}^{n-1} X^ke^{i(n-1-k)\theta}\Bigr)
  \Bigl(\sum_{\ell=0}^{n-1} X^le^{-i(n-1-\ell)\theta}\Bigr)$ \par
$= \sum_{k=0}^{n-1} X^k\Bigl({\sum_{p=0}^k e^{i(k-2p)\theta}}\Bigr)
 + \sum_{k=n}^{2n-2} X^k\Bigl({\sum_{p=k-n+1}^{n-1} e^{i(k-2p)\theta}}\Bigr)$ \par
$= \sum_{k=0}^{n-1} X^k \frac {\sin(k+1)\theta}{\sin\theta}
 + \sum_{k=n}^{2n-2} X^k \frac {\sin(2n-k-1)\theta}{\sin\theta}$. \par

\fincorrection
\correction{003218}
Division de proche en proche : $Q = \sum_{k=0}^{n-1} X^k\cos k\theta$.
\fincorrection
\correction{003219}
$\iff X^2 + X + 1\ |\ X^{2n} + pX^n + q \iff j^{2n} + pj^n + q = 0$.
\fincorrection
\correction{003220}
$(X+1)(3X-1)(X^2+3X+5)$.
\fincorrection
\correction{003221}
On calcule pgcd$({P(X),Q(X)}) = X^2+5$.

$ \Rightarrow  x_1 = i\sqrt5$ et $x_2 = -i\sqrt5$.

On obtient alors : $P(X) = (X^2+5)(X^2-3X+1)$.

Les deux derni{\`e}res racines sont $x_3 = \frac {3+\sqrt5}2$ et
$x_4 = \frac {3-\sqrt5}2$.
\fincorrection
\correction{003222}
$P = (X-2)^2(X-3)^3$.
\fincorrection
\correction{003223}
$a = 10i$.
Racines : $i,i,i,\frac {-3i+\sqrt{15}}2, \frac {-3i-\sqrt{15}}2$.
\fincorrection
\correction{003224}
$\lambda = 0$, $x=1$.
\fincorrection
\correction{003225}
$P$ doit {\^e}tre divisible par $X^2-X+r$,
         $ \Rightarrow  r^2-3r+p+1 = 2r^2 - r + q = 0$.\par
         On calcule le pgcd de ces expressions $ \Rightarrow $ CNS :
         $4p^2 - 4pq + q^2 + 3p + 11q - 1 = 0$.
\fincorrection
\correction{003226}
$ = (-1)^{n+1}\frac {(X-1)\cdots(X-n)\bigl(X-(n+1)\bigr)}{(n+1)!}$.
\fincorrection
\nocorrection
\correction{003228}
Pour $x \in \R$, on a $P(x) = \Im\bigl((1+xe^{i\theta})^n\bigr)$.
\\
Donc $P(x) = 0 \iff \exists\ k \in \{0,\dots,n-1\}$ et $\lambda \in \R$
tels que : $1+xe^{i\theta} = \lambda e^{ik\pi/n}$.
\\
On obtient $x_k = \frac{\sin(k\pi/n)}{\sin(\theta-k\pi/n)}$,\quad
$0\le k \le n-1$.
\fincorrection
\nocorrection
\correction{003230}
$P = a(X-b)^\alpha$.
\fincorrection
\correction{003231}
\begin{enumerate}
  \item si $P(x) = 0$, alors $P\bigl((x-1)^2\bigr) = P\bigl((x+1)^2\bigr) = 0$.
\\
On a toujours $|x| < \max\{ |x-1|,|x+1| \}$ donc, s'il y a une racine
de module $> 1$, il n'y a pas de racine de module maximal $ \Rightarrow  P = 0$.
\\
Or $\max\{ |x-1|,|x+1| \} \ge 1$ avec {\'e}galit{\'e} ssi $x = 0$.
Donc $P = 0$ ou $P = 1$.
  \item Si $x$ est racine, alors $x^2$ et $(x+1)^2$ le sont aussi.
\\
$ \Rightarrow  |x| = 0$ ou $1  \Rightarrow  |x+1| = 0$ ou $1  \Rightarrow  x \in \{0,-1,j,j^2 \}$.
\\
$x = 0$ ou $x = -1  \Rightarrow  P(1) = 0$ : exclus.
\\
Donc $P = a(X-j)^\alpha (X-j^2)^\beta$. On remplace
$ \Rightarrow  P = (X^2 + X + 1)^\alpha$.
  \item Seule racine possible : 1 $ \Rightarrow  P=-(X-1)^k$.
\end{enumerate}
\fincorrection
\nocorrection
\correction{003233}
M{\^e}mes racines avec les m{\^e}mes multiplicit{\'e}s.
\fincorrection
\correction{003234}
\begin{enumerate}
  \item $P = \frac \Phi{X-z_k}  \Rightarrow 
     \frac {\Phi(z_0)}{z_0-z_k}  = \frac {\Phi'(z_k)}n$.
  \item Les deux membres sont {\'e}gaux en $z_0,\dots,z_n$.
  \item D{\'e}composer $\Phi$ sur la base $\bigl((X-z_0)^k\bigr)$.
  \item $\sum_k e^{2ikp/n} = 0$ pour $p < n  \Rightarrow $ OK.
\end{enumerate}
\fincorrection
\nocorrection
\nocorrection
\nocorrection
\correction{003238}
$x = z + \frac 1z$ avec $z^6 + z^5 + \dots + 1 = 0$.
\\
Autres racines : $2\cos\frac{4\pi}7$ et $2\cos\frac{6\pi}7$.
\fincorrection
\correction{003239}
\begin{enumerate}
  \item Soit $P(x) = a_n\prod_{k=1}^n(x-x_k)$.
      On a : $\sum_{k=0}^n \frac 1{(x-x_k)^2}
               = -\frac d{dx}\Bigl(\frac {P'}P\Bigr)(x) = \frac {P'^2 - PP''}{P^2}(x)$.
  \item Pour $k=1,x=0$, on a : $a_0a_2 \le \frac 12a_1^2$.\\
      Pour $k$ quelconque : on applique le cas pr{\'e}c{\'e}dent {\`a} $P^{(k-1)}$ dont les
      racines sont encore r{\'e}elles simples : \\
      $(k-1)!a_{k-1} \times \frac {(k+1)!}2a_{k+1} \le \frac 12(k!a_k)^2
        \Rightarrow 
       a_{k-1}a_{k+1} \le \frac {k}{k+1}a_k^2$.
\end{enumerate}
\fincorrection
\correction{003240}
$Q = X^2 -3X +1$,
$P = \left(X-\frac {5+\sqrt{33}}2\right) \left(X-\frac {5-\sqrt{33}}2\right)
     (X^2 - X + 4)$.
\fincorrection
\nocorrection
\correction{003242}
\begin{enumerate}
  \item $p$ est premier car $ K$ est int{\`e}gre.

    On a $1^p = 1$,
    $(xy)^p = x^py^p$ (un corps est commutatif) et
    $(x+y)^p = x^p + y^p + \sum_{k=1}^{p-1}C_p^kx^ky^{p-k} = x^p + y^p$
    car $p$ divise $C_p^k$ si $1\le k \le p-1$.
  \item Remarquer que $P'=0 \Leftrightarrow P\in  K[X^p]$.
    
    On suppose $\sigma$ surjectif. Soit $P(X) = Q(X^p) = a_0 + \dots + a_kX^{kp}$
    un polyn{\^o}me non constant {\`a} d{\'e}riv{\'e}e nulle. Il existe $b_0,\dots,b_k$
    tels que $b_i^p = a_i$. Alors $P(X) = Q(X)^p$ est r{\'e}ductible.
    
    On suppose que tout polyn{\^o}me irr{\'e}ductible a une d{\'e}riv{\'e}e non nulle.
    Soit $a\in  K$ et $P(X) = X^p - a$. $P'=0$ donc $P$ est r{\'e}ductible.
    Soit $Q$ un facteur unitaire irr{\'e}ductible de $X^p-a$. Alors $Q^p$ et $X^p-a$
    ont $Q$ en facteur commun donc leur pgcd, $D$, est non constant.
    Mais $Q^p$ et $X^p-a$ appartiennent {\`a} $ K[X^p]$ donc $D$, obtenu par l'algorithme
    d'Euclide aussi, d'o{\`u} $D = X^p-a$ et $X^p-a$ divise $Q^p$. Par unicit{\'e} de
    la d{\'e}composition de $Q^p$ en facteurs irr{\'e}ductibles, il existe $r\in\N$
    tel que $X^p-a = Q^r$. Par examen du degr{\'e} on a $r=p$ donc $\deg Q = 1$,
    $Q = X-b$ et finalement $b^p=a$.
\end{enumerate}
\fincorrection
\correction{003243}
$V(\alpha)$ est pair si et seulement si $P(\alpha)$ et $P^{(n)}(\alpha)$
ont m{\^e}me signe, de m{\^e}me pour $V(\beta)$. Comme $P^{(n)}(\alpha) = P^{(n)}(\beta)$
on en d{\'e}duit que $V(\alpha)-V(\beta)$ est pair si et seulement si $P(\alpha)$
et $P(\beta)$ ont m{\^e}me signe, donc si et seulement si $P$ a un nombre pair
de racines dans $[\alpha,\beta]$.

D{\'e}croissance de~$V$~: $V$ est constant sur tout intervalle ne contenant
aucune racine de $P,P',\dots,P^{(n-1)}$. Consid{\'e}rons $x_0\in{[\alpha,\beta[}$
tel que $P^{(k)}(x_0) \ne 0$, $P^{(k+1)}(x_0) = \dots = P^{(\ell-1)}(x_0) = 0$
et $P^{(\ell)}(x_0) \ne 0$. Alors pour $x$ proche de $x_0$ avec $x>x_0$, $P^{(k)}(x)$
a m{\^e}me signe que $P^{(k)}(x_0)$ et $P^{(k+1)}(x),\dots,P^{(\ell)}(x)$ ont
m{\^e}me signe que $P^{(\ell)}(x_0)$ donc les nombres de changements de signe dans
les sous-suites $(P^{(k)}(x),\dots,P^{(\ell)}(x))$ et $(P^{(k)}(x_0),\dots,P^{(\ell)}(x_0))$
sont {\'e}gaux. De m{\^e}me si $P(x_0) = \dots = P^{(\ell-1)}(x_0) = 0$ et
$P^{(\ell)}(x_0) \ne 0$. Ceci prouve que $V(x_0^+) = V(x_0)$ pour tout
$x_0\in{[\alpha,\beta[}$.

On consid{\`e}re {\`a} pr{\'e}sent $x_0\in{]\alpha,\beta]}$ tel que
$P^{(k)}(x_0) \ne 0$, $P^{(k+1)}(x_0) = \dots = P^{(\ell-1)}(x_0) = 0$
et $P^{(\ell)}(x_0) \ne 0$. Alors pour $x$ proche de $x_0$ avec $x<x_0$
la sous-suite $(P^{(k)}(x),\dots,P^{(\ell)}(x))$ a $\ell-k-1$ changements de
signe si $P^{(k)}(x_0)$ et  $P^{(\ell)}(x_0)$ ont m{\^e}me signe, $\ell-k$
changements de signe sinon tandis que la sous-suite $(P^{(k)}(x_0),\dots,P^{(\ell)}(x_0))$
en a un ou z{\'e}ro. De m{\^e}me, si $P(x_0) = \dots = P^{(\ell-1)}(x_0) = 0$ et
$P^{(\ell)}(x_0) \ne 0$ on trouve $\ell$ changements de signe pour
$(P(x),\dots,P^{(\ell)}(x))$ et z{\'e}ro pour $(P(x_0),\dots,P^{(\ell)}(x_0))$
donc dans tous les cas $V(x_0^-) \ge V(x_0)$. Ceci ach{\`e}ve la d{\'e}monstration.
\fincorrection
\correction{003244}
Pour $z\in\mathbb{U}$, on a $Q(z) = 0 \Leftrightarrow {P(z)}/{z^d\overline P(\overline
z)} = -\omega$. Comme $\overline P(\overline z) = \overline{P(z)}$, les
deux membres ont m{\^e}me module pour tout~$z\in\mathbb{U}$, il faut et il suffit donc
que les arguments soient {\'e}gaux modulo $2\pi$.
Pour $a\in\C$ avec $|a| < 1$, une d{\'e}termination continue de
$\mathrm{Arg}(e^{i\theta}-a)$ augmente de~$2\pi$ lorsque $\theta$ varie de~$0$ {\`a}~$2\pi$
donc, vu l'hypoth{\`e}se sur les racines de~$P$, une d{\'e}termination continue
de $\mathrm{Arg}(P(z)/z^d\overline{P(z)}\,)$ augmente de $2\pi d$ lorsque $\theta$ varie
de~$0$ {\`a}~$2\pi$. Une telle d{\'e}termination prend donc au moins $d$ fois
une valeur congrue {\`a} $\mathrm{Arg}(-\omega)$ modulo~$2\pi$, ce qui prouve que $Q$
admet au moins $d$ racines distinctes dans~$\mathbb{U}$.
\fincorrection
\correction{003245}
$f(2k\pi/n) > 0 > f((2k+1)\pi/n)$ pour $k\in\Z$ donc $f$ admet
$2n$ racines dans $[0,2\pi[$. En posant $z=e^{ix}$, $z^nf(x)$ est un polyn{\^o}me
en~$z$ de degr{\'e} $2n$ ayant $2n$ racines sur le cercle unit{\'e}~; il n'en n'a pas ailleurs.
\fincorrection
\correction{003246}
$(X^2-X+1) (X^2+X+1) (X^2-X\sqrt3+1) (X^2+X\sqrt3+1)$.
\fincorrection
\correction{003247}
Racines :
$\alpha = 2 + \sqrt{\frac {\sqrt2+1}2} + i\sqrt{\frac {\sqrt2-1}2},\
 \overline{\alpha},\
 \beta =  2 - \sqrt{\frac {\sqrt2+1}2} - i\sqrt{\frac {\sqrt2-1}2},\
 \overline{\beta}$.

Factorisation de $P$ sur $\R$ :
$P = (X^2 - 2\Re(\alpha)X + |\alpha|^2) (X^2 - 2\Re(\beta)X + |\beta|^2)$
et les facteurs sont irrationnels.
\fincorrection
\correction{003248}
\begin{enumerate}
  \item $P = |Q+iR|^2$.
  \item Factoriser $P$.
  \item Avec Maple~: $P=\frac1{65\strut}Q\overline Q$ avec
    $Q=65X^2+(49i-67)X+(42+11i)$ et $Q$ est irr{\'e}ductible sur $Q[i]$.
    \let\l\lambda\def\m{\overline\lambda}%
    Donc si $P=A^2+B^2 = (A+iB)(A-iB)$ avec $A,B$ polyn{\^o}mes {\`a} coefficients entiers
    alors, quitte {\`a} changer $B$ en $-B$, il existe $\l\in\Q[i]$ tel que~:
    $A+iB = \l Q$ et $A-iB=\m\overline Q$ d'o{\`u}~:
    $$\begin{aligned}
      2A  &= 65(\l+\m)X^2 + ((49i-67)\l-(49i+67)\m)X + ((42+11i)\l+(42-11i)\m)\cr
      2iB &= 65(\l-\m)X^2 + ((49i-67)\l+(49i+67)\m)X + ((42+11i)\l-(42-11i)\m)\cr
      \l\m &=65.\cr
    \end{aligned}$$
    En particulier $65\l\in \Z[i]$, {\'e}crivons $\l = \frac{u+iv}{\strut 65}$
    avec $u,v\in\Z$~:
    $$\begin{aligned}
      A &= uX^2 -\frac{67u+49v}{\strut 65}X + \frac{42u-11v}{\strut 65}\cr
      B &= vX^2 +\frac{49u-67v}{\strut 65}X + \frac{11u+42v}{\strut 65}\cr
      u^2+v^2 &=65.\cr
    \end{aligned}$$
    $67u+49v$ est divisible par $65$ si et seulement si $u\equiv 8v(\mathrm{mod}\,{65})$
    et dans ce cas les autres num{\'e}rateurs sont aussi multiples de~$65$.
    La condition $u^2+v^2 = 65$ donne alors $v=\pm1, u=\pm8$ d'o{\`u}~:
    $$A = \pm(8X^2 -9X + 5),   \qquad B = \pm(X^2 +5X + 2).$$
\end{enumerate}
\fincorrection
\nocorrection
\nocorrection
\correction{003251}
\begin{enumerate}
  \item Si $P = QR$ alors $Q(a_i)R(a_i) = -1  \Rightarrow  Q(a_i) = -R(a_i) = \pm1$,
    donc $Q+R$ a $n$ racines, donc est nul, et $P = -Q^2$~:
    contradiction pour $x \to \infty$.
  \item M{\^e}me raisonnement : $P = Q^2$, donc
    $Q^2-1 = (Q-1)(Q+1) = (X-a_1) \dots (X-a_n)$.
    \par
    On r{\'e}partit les facteurs entre $Q-1$ et $Q+1$ : $n = 2p$, contradiction.
\end{enumerate}
\fincorrection
\correction{003252}
Soit $P = QR$ avec
$Q=X^{n_1}+b_{n_1-1}X^{n_1-1}+\dots+b_0X^0$ et
$R=X^{n_2}+c_{n_2-1}X^{n_2-1}+\dots+c_0X^0$.

Par hypoth{\`e}se sur $a_0 = b_0c_0$, $p$ divise un et un seul des entiers
$b_0$, $c_0$.
Supposons que $p$ divise $b_0,b_1,\dots,b_{k-1}$~:
alors $a_k \equiv b_kc_0 (\mathrm{mod}\, p)$ donc $p$ divise $b_k$.
On aboutit {\`a} \og$p$ divise le coefficient dominant
de $Q$\fg, ce qui est absurde.
\fincorrection
\correction{003253}
On suppose $a\ne 0$ et $X^p-a = PQ$ avec $P,Q\in K[X]$ unitaires
non constants. Soit $n=\deg(P)\in{[[1,p-1]]}$ et $b=(-1)^nP(0)\in K$.
$b$ est le produit de cetraines $p$-{\`e}mes de~$a$, donc $b^p = a^n$.
De plus $n\wedge p = 1$~; soit $nu+pv=1$ une relation de B{\'e}zout.
On a alors $b^{pu} = a^{nu} = a^{1-pv}$ d'o{\`u} $a = (b^u/a^v)^p$ 
donc $b^u/a^v\in K$ est racine de $X^p-a$.
\fincorrection
\correction{005318}
Soit $n\in\Nn$.

\begin{align*}\ensuremath
(X+1)^n-X^n-1\;\mbox{est divisible par}\;X^2+X+1&\Leftrightarrow j\;\mbox{et}\;j^2\;\mbox{sont racines de}\;(X+1)^n-X^n-1\\
 &\Leftrightarrow j\;\mbox{est racine de}\;(X+1)^n-X^n-1\\
 &(\mbox{car}\;(X+1)^n-X^{n-1}\;\mbox{est dans}\;\Rr[X])\\
 &\Leftrightarrow(j+1)^n-j^n-1=0\Leftrightarrow(-j^2)^n-j^n-1=0.
\end{align*}

Si $n\in6\Zz$, $(-j^2)^n-j^n-1=-3\neq0$.
 
Si $n\in1+6\Zz$, $(-j^2)^n-j^n-1=-j^2-j-1=0$.
 
Si $n\in2+6\Zz$, $(-j^2)^n-j^n-1=j-j^2-1=2j\neq0$.
 
Si $n\in3+6\Zz$, $(-j^2)^n-j^n-1=-3\neq0$.
 
Si $n\in4+6\Zz$, $(-j^2)^n-j^n-1=j^2-j-1=2j^2\neq0$.
 
Si $n\in5+6\Zz$, $(-j^2)^n-j^n-1=-j-j^2-1=0$.

En résumé, $(X+1)^n-X^n-1$ est divisible par $X^2+X+1$ si et seulement si $n$ est dans $(1+6\Zz)\cup(5+6\Zz)$.

\fincorrection
\correction{005319}
Soit $P$ un polynôme non nul à coefficients réels.

Pour tout réel $x$, on peut écrire 

$$P(x)=\lambda\prod_{i=1}^{k}(x-a_i)^{\alpha_i}\prod_{j=1}^{l}((x-z_j)(x-\overline{z_j}))^{\beta_j},$$

où $\lambda$ est un réel non nul, $k$ et $l$ sont des entiers naturels, les $a_i$ sont des réels deux à deux distincts, les $\alpha_i$ et les $\beta_i$ des entiers naturels et les $(x-z_j)(x-\overline{z_j})$ des polynômes deux à deux premiers entre eux à racines non réelles.

Tout d'abord, pour tout réel $x$, $\prod_{j=1}^{l}((x-z_j)(x-\overline{z_j}))^{\beta_j}>0$ (tous les trinomes du second degré considérés étant unitaires sans racines réelles.)

Donc, $(\forall x\in\Rr,\;P(x)\geq0)\Leftrightarrow(\forall x\in\Rr,\;\lambda\prod_{i=1}^{k}(x-a_i)^{\alpha_i}\geq0)$.

Ensuite, si $\forall x\in\Rr,\;P(x)\geq0$, alors $\lim_{x\rightarrow +\infty}P(x)\geq0$ ce qui impose $\lambda>0$. Puis, si un exposant $\alpha_i$ est impair, $P$ change de signe en $a_i$,  ce qui contredit l'hypothèse faite sur $P$. Donc, $\lambda>0$ et tous les $\alpha_i$ sont pairs. Réciproquement, si $\lambda>0$ et si tous les $\alpha_i$ sont pairs, alors bien sûr, $\forall x\in\Rr,\;P(x)\geq0$.

Posons $A=\sqrt{\lambda}\prod_{i=1}^{k}(x-a_i)^{\alpha_i/2}$. $A$ est un élément de $\Rr[X]$ car $\lambda>0$ et car les $\alpha_i$ sont des entiers pairs. Posons ensuite $Q_1=\prod_{j=1}^{l}(x-z_j)^{\beta_j}$ et $Q_2=\prod_{j=1}^{l}(x-\overline{z_j})^{\beta_j}$. $Q_1$ admet après développement une écriture de la forme $Q_1=B+iC$ où $B$ et $C$ sont des polynômes à coefficients réels. Mais alors, $Q_2=B-iC$. Ainsi, $$P=A^2Q_1Q_2=A^2(B+iC)(B-iC)=A^2(B^2+C^2)=(AB)^2+(AC)^2=R^2+S^2,$$

où $R$ et $S$ sont des polynômes à coefficients réels.
\fincorrection
\correction{005324}
Soit $P$ un polynôme de degré $n$ supèrieur ou égal à $2$.

Posons $P=\lambda(X-z_1)(X-z_2)...(X-z_n)$ où $\lambda$ est un complexe non nuls et les $z_k$ des complexes pas nécessairement deux à deux distincts.

$$P'=\lambda\sum_{i=1}^{n}(\prod_{j\neq i}^{}(X-z_j))=\sum_{i=1}^{n}\frac{P}{X-z_i},$$

et donc 
 
$$\frac{P'}{P}=\sum_{i=1}^{n}\frac{1}{X-z_i}.$$

Soit alors $z$ une racine de $P'$ dans $\Cc$. Si z est racine de $P$ (et donc racine de $P$ d'ordre au moins $2$) le résultat est clair. Sinon,

$$0=\frac{P'(z)}{P(z)}=\sum_{i=1}^{n}\frac{1}{z-z_i}=\sum_{i=1}^{n}\frac{\overline{z-z_i}}{|z-z_i|^2}.$$

En posant $\lambda_i=\frac{1}{|z-z_i|^2}$, ($\lambda_i$ est un réel strictement positif) et en conjugant, on obtient
$\sum_{i=1}^{n}\lambda_i(z-z_i)=0$ et donc 
$$z=\frac{\sum_{i=1}^{n}\lambda_iz_i}{\sum_{i=1}^{n}\lambda_i}=\mbox{bar}(z_1(\lambda_1),...,z_n(\lambda_n)).$$
\fincorrection
\correction{005325}
On suppose que $n=\mbox{deg}P\geq1$.

On pose $P=\lambda(X-z_1)(X-z_2)...(X-z_n)$ où $\lambda$ est un complexe non nul et les $z_k$ sont des complexes pas nécessairement deux à deux distincts.

D'après l'exercice précédent, $\frac{P'}{P}=\sum_{k=1}^{n}\frac{1}{X-z_k}$.

Si $P$ est divisible par $P'$, $\exists(a,b)\in\Cc^2\setminus\{(0,0)\}/\;P=(aX+b)P'$ et donc $\exists(a,b)\in\Cc^2\setminus\{(0,0)\}/\;\frac{P'}{P}=\frac{1}{aX+b}$ ce qui montre que la fraction rationelle $\frac{P'}{P}$ a exactement un et un seul pôle complexe et donc que les $z_k$ sont confondus.

En résumé, si $P'$ divise $P$, $\exists(a,\lambda)\in\Cc^2/\;P=\lambda(X-a)^n$ et $\lambda\neq0$.

Réciproquement, si $P=\lambda(X-a)^n$ avec $\lambda\neq0$, alors $P'=n\lambda(X-a)^{n-1}$ divise $P$.

Les polynômes divisibles par leur dérivée sont les polynômes de la forme $\lambda(X-a)^n$, $\lambda\in\Cc\setminus\{0\}$, $n\in\Nn^*$, $a\in\Cc$.
\fincorrection
\correction{005328}
$a$ est solution du problème si et seulement si $X^5-209X+a$ est divisible par un polynôme de la forme $X^2+\alpha X+1$. Mais 

$$X^5-209X+a=(X^2+\alpha X+1)(X^3-\alpha X^2+(\alpha^2-1)X-(\alpha^3-2\alpha))+(\alpha^4-3\alpha^2-208)X+a+(\alpha^3-2\alpha).$$

Donc a est solution $\Leftrightarrow\exists\alpha\in\Cc/\;\left\{
\begin{array}{l}
\alpha^4-3\alpha^2-208=0\\
a=-\alpha^3+2\alpha
\end{array}
\right.$. Mais, $\alpha^4-3\alpha^2-208=0\Leftrightarrow\alpha^2\in\{-13,16\}\Leftrightarrow\alpha\in\{-4,4,i\sqrt{13},-i\sqrt{13}\}$ et la deuxième équation fournit $a\in\{56,-56,15i\sqrt{13},-15i\sqrt{13}\}$.

\fincorrection
\correction{005329}
On note que $P(1)=1\neq0$ et donc que l'expression proposée a bien un sens.

$$\sum_{k=1}^{5}\frac{a_k+2}{a_k-1}=\sum_{k=1}^{5}(1+\frac{3}{a_k-1})=5-3\sum_{k=1}^{5}\frac{1}{1-a_k}=5-3\frac{P'(1)}{P(1)}=5-3\frac{12}{1}=-31.$$

\fincorrection
\correction{005342}
\begin{align*}\ensuremath
P&=X^6-2X^3\cos a+1=(X^3-e^{ia})(X^3-e^{-ia})\\
 &=(X-e^{ia/3})(X-je^{ia/3})(X-j^2e^{ia/3})(X-e^{-ia/3})(X-je^{-ia/3})(X-j^2e^{-ia/3})\\
 &=(X^2-2X\cos\frac{a}{3}+1)(X^2-2X\cos(\frac{a}{3}+\frac{2\pi}{3})+1)(X^2-2X\cos(\frac{a}{3}-\frac{2\pi}{3})+1)
\end{align*}
 
Il reste à se demander 1) si les facteurs précédents sont irréductibles sur $\Rr$ et 2) si ces facteurs sont deux à deux distincts.

Les trois facteurs de degré $2$ ont un discriminant réduit du type $\Delta'=\cos^2\alpha-1=-\sin^2\alpha$ et $\Delta'$ est nul si et seulement si $\alpha$ est dans $\pi\Zz$.

Les cas particuliers sont donc ($\frac{a}{3}$ est dans $\pi\Zz$ et donc $a=0$) et ($\frac{a+2\pi}{3}$ est dans $\pi\Zz$ et donc $a=\pi$) et ($\frac{a-2\pi}{3}$ est dans $\pi\Zz$ ce qui n'a pas de solution dans $[0,\pi]$).

\begin{itemize}
\item[1er cas.] Si $a=0$.

$$P=(X^2-2X+1)(X^2+X+1)(X^2+X+1)=(X-1)^2(X^2+X+1)^2.$$

\item[2ème cas.] Si $a=\pi$, en remplaçant $X$ par $-X$ on obtient :

$$P=(X+1)^2(X^2-X+1)^2.$$

\item[3ème cas.] Si $a$ est dans $]0,\pi[$, les trois facteurs de degré $2$ sont irréductibles sur $\Rr$ et clairement deux à deux distincts. Donc 

$$P=(X^2-2X\cos\frac{a}{3}+1)(X^2-2X\cos\frac{a+2\pi}{3}+1)(X^2-2X\cos\frac{a-2\pi}{3}+1).$$
\end{itemize}
\fincorrection
\correction{005345}
Pour $k$ élément de $\{-3,-2,-1,1,2,3\}$, posons $x_k=\sin\frac{k\pi}{7}$ (les $x_k$ sont deux à deux opposés). Il faut calculer les coefficients du polynôme

\begin{align*}\ensuremath
P&=(X-\sin\frac{\pi}{7})(X-\sin\frac{2\pi}{7})(X-\sin\frac{3\pi}{7})(X+\sin\frac{\pi}{7})(X+\sin\frac{2\pi}{7})(X+\sin\frac{3\pi}{7})\\
 &=(X^2-\sin^2\frac{\pi}{7})(X^2-\sin^2\frac{2\pi}{7})(X^2-\sin^2\frac{3\pi}{7})\\
 &=(X^2-\frac{1}{2}(1-cos\frac{2\pi}{7}))(X^2-\frac{1}{2}(1-cos\frac{4\pi}{7}))
 (X^2-\frac{1}{2}(1-cos\frac{6\pi}{7}))\\
  &=\frac{1}{8}Q(-2X^2+1)
\end{align*}

où $Q(Y)=(\cos\frac{2\pi}{7}-Y)(\cos\frac{4\pi}{7}-Y)(\cos\frac{8\pi}{7}-Y)$.

Posons $\omega=e^{2i\pi/7}$.

\begin{align*}\ensuremath
\cos\frac{2\pi}{7}\cos\frac{4\pi}{7}\cos\frac{6\pi}{7}&=\frac{1}{8}(\omega+\omega^6)(\omega^2+\omega^5)(\omega^3+\omega^4)
=\frac{1}{8}(6+\omega^7+\omega^9+\omega^{10}+\omega^{11}+\omega^{12}+\omega^{14}+\omega^{15})\\
 &=\frac{1}{8}(\omega^6+\omega^7+\omega^2+\omega^3+\omega^4+\omega^5+1+\omega)=\frac{1}{8}.
\end{align*}

Puis,

\begin{align*}\ensuremath
\cos\frac{2\pi}{7}\cos\frac{4\pi}{7}+\cos\frac{2\pi}{7}\cos\frac{6\pi}{7}+
\cos\frac{6\pi}{7}\cos\frac{4\pi}{7}&=\frac{1}{4}((\omega+\omega^6)(\omega^2+\omega^5)+(\omega+\omega^6)(\omega^3+\omega^4)+(\omega^3+\omega^4)(\omega^2+\omega^5))\\
 &=\frac{1}{4}(2\omega+2\omega^2+2\omega^3+2\omega^4+2\omega^5+2\omega^6)=\frac{-2}{4}=-\frac{1}{2}.
\end{align*}

Enfin,

\begin{align*}\ensuremath
\cos\frac{2\pi}{7}+\cos\frac{4\pi}{7}+\cos\frac{6\pi}{7}&=\frac{1}{2}(\omega+\omega^2+\omega^3+\omega^4+\omega^5+\omega^6)=-\frac{1}{2}
\end{align*}

Donc, $Q=\frac{1}{8}-(-\frac{1}{2})Y+(-\frac{1}{2})Y^2-Y^3=\frac{1}{8}(-8Y^3-4Y^2+4Y+1)$ puis,
 
$$P=\frac{1}{64}(-8(-2X^2+1)^3-4(-2X^2+1)^2+4(-2X^2+1)+1)=\frac{1}{64}(64X^6-112X^4+54X^2-7).$$
Une équation du $6$ème degré dont les solutions sont les sin est $64x^6-112x^4+54x^2-7=0$.

Maintenant, si $r=$ ($p$ entier relatif non nul, $q$ entier naturel  non nul, $p$ et $q$ premiers entre eux) est une racine rationnelle de cette équation, alors, d'après l'exercice \ref{exo:suprou9}, $p$ divise $-7$ et $q$ divise $64$ et donc 
$p$ est élément de $\{1,-1,7,-7\}$ et $q$ est élément de $\{1,2,4,8,16,32,64\}$. On vérifie aisémént qu'aucun des rationnels $r$ obtenu n'est racine de $P$ et donc les racines de $P$ sont irrationnelles.
\fincorrection
\correction{005349}
Posons $P=X^4-4X^3-36X^2+\lambda X+\mu$.
 
\begin{align*}\ensuremath
(\lambda,\mu)\;\mbox{solution}&\Leftrightarrow\exists(z,r)\in\Cc^2/\;\mbox{les racines de}\;P\;\mbox{soient}\;z,z+r,z+2r,z+3r\\
  &\Leftrightarrow\exists(z,r)\in\Cc^2/\;\left\{
 \begin{array}{l}
 \sigma_1=4\\
 \sigma_2=-36\\
 \sigma_3=-\lambda\\
 \sigma_4=\mu
 \end{array}
 \right.
 \\
 &\Leftrightarrow\exists(z,r)\in\Cc^2/\;\left\{
 \begin{array}{l}
 4z+6r=4\\
 z(3z+6r)+(z+r)(2z+5r)+(z+2r)(z+3r)=-36\\
 \sigma_3=-\lambda\\
 \sigma_4=\mu
 \end{array}
 \right.
 \\
 &\Leftrightarrow\exists(z,r)\in\Cc^2/\;\left\{
 \begin{array}{l}
 2z+3r=2\\
 6z^2+18rz+11r^2=-36\\
 \sigma_3=-\lambda\\
 \sigma_4=\mu
 \end{array}
 \right.
 \\
 &\Leftrightarrow\exists(z,r)\in\Cc^2/\;\left\{
 \begin{array}{l}
 z=1-\frac{3}{2}r\\
 6(1-\frac{3}{2}r)^2+18(1-\frac{3}{2}r)r+11r^2=-36\\
 \sigma_3=-\lambda\\
 \sigma_4=\mu
 \end{array}
 \right.
 \\
 &\Leftrightarrow\exists(z,r)\in\Cc^2/\;\left\{
 \begin{array}{l}
 -\frac{5}{2}r^2+42=0\\
 z=1-\frac{3}{2}r\\
 \sigma_3=-\lambda\\
 \sigma_4=\mu
 \end{array}
 \right.
\end{align*}

D'où la solution (les deux valeurs opposées de $r$ fournissent évidemment la même progression arithmétique)
$r=2\sqrt{\frac{21}{5}}$ puis $z=1-3\sqrt{\frac{21}{5}}$ puis les racines $z_1=1-3\sqrt{\frac{21}{5}}$, $z_2=1-\sqrt{\frac{21}{5}}$, $z_3=1+\sqrt{\frac{21}{5}}$ et $z_4=1+3\sqrt{\frac{21}{5}}$, obtenues pour 

$$\lambda=z_1z_2z_3z_4=(1-3\sqrt{\frac{21}{5}})(1-\sqrt{\frac{21}{5}})(1+\sqrt{\frac{21}{5}})
(1+3\sqrt{\frac{21}{5}})= (1-9\frac{21}{5})(1-\frac{21}{5})=\frac{2994}{25},$$
 
et 

\begin{align*}\ensuremath
\mu&=(1-3\sqrt{\frac{21}{5}})(1-\frac{21}{5})+(1-9\frac{21}{5})(1-\sqrt{\frac{21}{5}})+(1-9\frac{21}{5})(1+\sqrt{\frac{21}{5}})+(1-\frac{21}{5})(1+3\sqrt{\frac{21}{5}})\\
 &=2(1-\frac{21}{5})+2(1-9\frac{21}{5})=2(2-10\frac{21}{5})=-80
\end{align*}
\fincorrection
\correction{005350}
Pour $i\in\{1,2,3\}$, on a $x_i^3+2x_i-1=0$ et donc $x_i^4+2x_i^2-x_i=0$. En additionnant ces trois égalités, on obtient $S_4+2S_2-S_1=0$ et donc

$$S_4=-2((\sigma_1^2-2\sigma_2)+\sigma_1=(-2)(-2.2)=8.$$
\fincorrection
\correction{005351}
Pour chacun des $8$ numérateurs possibles, il y a $C_7^2=21$ dénominateurs et donc au total, $8\times21=168$ termes.

$$\sum_{}^{}\frac{x_1}{x_2x_3}=\sum_{}^{}\frac{x_1^2x_4x_5x_6x_7x_8}{x_1x_2...x_8}=\frac{1}{\sigma_8}\sum_{}^{}x_1^2x_2x_3x_4x_5x_6=\frac{1}{3}\sum_{}^{}x_1^2x_2x_3x_4x_5x_6.$$

Ensuite,

$$\sigma_1\sigma_6=(\sum_{}^{}x_1)(\sum_{}^{}x_1x_2x_3x_4x_5x_6)=\sum_{}^{}x_1^2x_2x_3x_4x_5x_6+\sum_{}^{}x_1x_2x_3x_4x_5x_6x_7,$$
 
et donc, 

$$\sum_{}^{}x_1^2x_2x_3x_4x_5x_6=\sigma_1\sigma_6-\sigma_7=(-1)(0)-1=-1.$$

Donc, $\sum_{}^{}\frac{x_1}{x_2x_3}=-\frac{1}{3}$.
\fincorrection
\correction{006959}
\
\begin{enumerate}
\item \begin{enumerate}
\item $X^3-3=(X-3^{1/3})(X^2+3^{1/3}X+3^{2/3})$ où $X^2+3^{1/3}X+3^{2/3}$ 
est irréductible sur $\R$. On cherche ses racines complexes 
pour obtenir la factorisation sur $\Cc$ :
$$X^3-3=(X-3^{1/3})(X+\frac{1}{2}3^{1/3}-\frac{i}{2}3^{5/6})(X+\frac{1}{2}3^{1/3}+\frac{i}{2}3^{5/6})$$

\item Passons à $X^{12}-1$. $z=re^{i\theta}$ vérifie $z^{12}=1$ si et seulement 
si $r=1$ et $12\theta\equiv 0 [2\pi]$, on obtient donc comme racines 
complexes les $e^{ik\pi/6}$ ($k=0,\ldots,11$), 
parmi lesquelles il y en a deux réelles ($-1$ et $1$) et cinq couples de 
racines complexes conjuguées 
($e^{i\pi/6}$ et $e^{11i\pi/6}$, $e^{2i\pi/6}$ et $e^{10i\pi/6}$, 
$e^{3i\pi/6}$ et $e^{9i\pi/6}$, $e^{4i\pi/6}$ et $e^{8i\pi/6}$, 
$e^{5i\pi/6}$ et $e^{7i\pi/6}$), d'où la factorisation sur $\Cc[X]$:
$$\begin{array}{rl}
X^{12}-1=&(X-1)(X+1)(X-e^{i\pi/6})(X-e^{11i\pi/6})(X-e^{2i\pi/6})\\
         & \ \ (X-e^{10i\pi/6})(X-e^{3i\pi/6})(X-e^{9i\pi/6})(X-e^{4i\pi/6})\\
 & \ \ \ \ \ \ \ \ \ \ \ \ \ (X-e^{8i\pi/6})(X-e^{5i\pi/6})(X-e^{7i\pi/6})
\end{array}$$

Comme $(X-e^{i\theta})(X-e^{-i\theta})=(X^2-2\cos(\theta)X+1)$, on en déduit la factorisation dans $\Rr[X]$ :
$$\begin{array}{rl}
X^{12}-1&=(X-1)(X+1)(X^2-2\cos(\pi/6)X+1)\\ 
 &\ \ \ (X^2-2\cos(2\pi/6)X+1)(X^2-2\cos(3\pi/6)X+1)\\
 &\ \ \ \ \ \ \ \ (X^2-2\cos(4\pi/6)X+1)(X^2-2\cos(5\pi/6)X+1)\\
 &=(X-1)(X+1)(X^2-\sqrt{3}X+1)\\
 &\ \ \ \ \ \ (X^2-X+1)(X^2+1)(X^2+X+1)(X^2+\sqrt{3}X+1)
\end{array}$$

\item Pour $X^6+1$, $z=re^{i\theta}$ vérifie $z^{6}=-1$ si et seulement 
si $r=1$ et $6\theta\equiv \pi [2\pi]$, on obtient donc comme racines 
complexes les $e^{i(\pi+2k\pi)/6}$ ($k=0,\ldots,5$). D'où la factorisation dans $\Cc[X]$ :
$$\begin{array}{rl}
X^6+1 &=(X-e^{i\pi/6})(X-e^{3i\pi/6})(X-e^{5i\pi/6})(X-e^{7i\pi/6})\\
  &\ \ \ \ \ (X-e^{9i\pi/6})(X-e^{11i\pi/6})
\end{array}$$

Pour obtenir la factorisation dans $\Rr[X]$, on regroupe les paires de racines complexes conjuguées :
$$X^6+1=(X^2+1)(X^2-\sqrt{3}X+1)(X^2+\sqrt{3}X+1)$$

\item $X^9+X^6+X^3+1=P(X^3)$ où $P(X)=X^3+X^2+X+1=\frac{X^4-1}{X-1}$ : 
les racines de $P$ sont donc les trois racines quatrièmes de l'unité 
différentes de $1$ ($i$, $-i$, $-1$) et 
$$\begin{array}{rcl}
X^9+X^6+X^3+1&=&P(X^3)\\
 &=&(X^3+1)(X^3-i)(X^3+i)\\
 &=&(X^3+1)(X^6+1)
\end{array}$$
On sait déjà factoriser $X^6+1$, il reste donc à factoriser le polynôme 
$X^3+1=(X+1)(X^2-X+1)$, où $X^2-X+1$ n'a pas de racine réelle. Donc
$$\begin{array}{rl}
X^9+X^6+X^3+1&=(X+1)(X^2-X+1)(X^2+1)\\
 & \ \ \ \ \ (X^2-\sqrt{3}X+1)(X^2+\sqrt{3}X+1)
\end{array}$$

Pour la factorisation sur $\Cc$ : les racines de $X^2-X+1$ sont $e^{i\pi/3}$ et $e^{5i\pi/3}$, ce qui donne
$$\begin{array}{rl}
X^9+X^6+X^3+1&=(X+1)(X-e^{i\pi/3})(X-e^{5i\pi/3})\\
 &\ \ \ (X-e^{i\pi/6})(X-e^{3i\pi/6})(X-e^{5i\pi/6})\\
 &\ \ \ (X-e^{7i\pi/6})(X-e^{9i\pi/6})(X-e^{11i\pi/6})
\end{array}$$
\end{enumerate}

\item 
\begin{enumerate}
\item Pour $X^2+(3i-1)X-2-i$, on calcule le discriminant 
$$\Delta=(3i-1)^2-4(-2-i)=-2i$$ 
et on cherche les racines carrées (complexes!) de $\Delta$: $w=a+ib$ 
vérifie $w^2=\Delta$ si et seulement si $w=1-i$ ou $w=-1+i$. 
Les racines du polynômes sont donc $\frac{1}{2}(-(3i-1)\pm(1-i))$ 
et $P(X)=(X+i)(X-1+2i)$.

\item Pour $X^3+(4+i)X^2+(5-2i)X+2-3i$: $-1$ est racine évidente, et $P(X)=(X+1)(X^2+(3+i)X+2-3i)$. 
Le discriminant du polynôme $X^2+(3+i)X+2-3i$ vaut $\Delta=18i$, 
ses deux racines carrées complexes sont $\pm(3+3i)$ et finalement 
on obtient $P(X)=(X+1)(X-i)(X+3+2i)$.
\end{enumerate}
\end{enumerate}
\fincorrection
\correction{006960}
Si $P$ est constant égal à $c$, il convient si et seulement si $c=c^2$,
et alors $c\in\{0;1\}$. 

Dans la suite on suppose $P$ non constant.
Notons $Z$ l'ensemble des racines de $P$. On sait que $Z$ est un ensemble non vide, fini.


\emph{Analyse}

Si $z\in Z$, alors $P(z)=0$ et la relation $P(X^2)=P(X)P(X+1)$ implique $P(z^2)=0$, donc $z^2\in Z$.
En itérant, on obtient $z^{2k}\in Z$ (pour tout $k\in\N^*$). 
Si $|z|>1$, la suite $(|z^{2k}|)_k$ est strictement croissante 
donc $Z$ contient une infinité d'éléments, ce qui est impossible. 
De même si $0<|z|<1$, la suite $(|z^{2k}|)_k$ est strictement décroissante, 
ce qui est impossible pour la même raison. 
Donc les éléments de $Z$ sont soit $0$, soit des nombres complexes de module $1$.

De plus, si $P(z)=0$, alors toujours par la relation $P(X^2)=P(X)P(X+1)$,
on a que $P((z-1)^2)=0$ donc $(z-1)^2\in Z$. Par le même raisonnement que précédemment,
alors ou bien $z-1=0$ ou bien $|z-1|=1$. 

En écrivant $z=a+ib$, on vérifie que $|z|=|z-1|=1$ équivaut à 
$z=e^{\pm i\pi/3}$. Finalement, $Z\subset\big\{0,1,e^{i\pi/3},e^{-i\pi/3}\big\}$. 
Or si $e^{\pm i\pi/3}$ était racine de $P$, alors $(e^{\pm i\pi/3})^2$ devrait aussi être dans $Z$, 
mais ce n'est aucun des quatre nombres complexes listés ci-dessus. 
Donc ni $e^{i\pi/3}$, ni $e^{-i\pi/3}$ ne sont dans $Z$. 
Les deux seules racines (complexes) possibles sont donc $0$ et $1$. 
Conclusion : le polynôme $P$ est nécessairement de la forme $\lambda X^k(X-1)^\ell$. 

\bigskip

\emph{Synthèse}

La condition $P(X^2)=P(X)P(X+1)$ devient
$$\lambda X^{2k}(X^2-1)^\ell=\lambda^2X^k(X-1)^\ell(X+1)^kX^\ell$$
qui équivaut à $\left\{\begin{array}{l}\lambda^2=\lambda\\2k=k+\ell\\k=\ell\end{array}\right.$.

Autrement dit $k=\ell$ et $\lambda=1$ (puisqu'on a supposé $P$ non constant). 

{\it Conclusion}
Finalement, les solutions sont le polynôme nul et les polynômes $(X^2-X)^k$, 
$k\in\Nn$ ($k=0$ donne le polynôme $1$).
\fincorrection
\correction{006961}
\begin{enumerate}
  \item Commençons par remarquer que si $P$ et $Q$ sont deux polynômes 
  qui conviennent, alors pour tout $z\in\Cc^*$, 
$P\left(z+\frac{1}{z}\right)-Q\left(z+\frac{1}{z}\right)=0$. 
En appliquant cette égalité à $z=e^{i\theta}$, on obtient 
$(P-Q)(2\cos\theta)=0$. Le polynôme $P-Q$ a une infinité 
de racines, donc il est nul, ce qui montre $P=Q$.
  
  \item Montrons l'existence de $P$ par récurrence forte sur $n$:
  \begin{itemize}
    \item Pour $n=0$, $P=2$ convient et pour $n=1$, $P=X$ convient.
    
    \item Passage des rangs $k\le n$ au rang $n+1$. 
    Si on note $P_k$ le polynôme construit pour $k\le n$, on a 
$$z^{n+1}+\frac{1}{z^{n+1}}=(z+\frac{1}{z})(z^n+\frac{1}{z^n})-(z^{n-1}+\frac{1}{z^{n-1}})
=(z+\frac{1}{z})P_n(z+\frac{1}{z})-P_{n-1}(z+\frac{1}{z})$$
donc $P_{n+1}(X)=XP_n(X)-P_{n-1}(X)$ convient.

    \item On a ainsi construit $P_n$ pour tout $n$ (avec $\deg P_n =n$). 
  \end{itemize}

  \item Fixons $n$ et notons $P$ le polynôme obtenu.
  Pour tout $\theta\in\R$, $P(e^{i\theta}+e^{-i\theta})=e^{in\theta}+e^{-in\theta}$ 
  donc $P(2\cos(\theta))=2\cos(n\theta)$. 
  
  En posant $x=2\cos(\theta)$ et donc $\theta = \Arccos(\frac x2)$ on obtient la relation
  Ainsi,
$$P(x)=2\cos(n\Arccos(\frac x2)) \qquad \forall x\in[-2,2]$$
Le polynôme dérivée est $P'(x)=\frac{n}{\sqrt{1-(\frac{x}{2})^2}}\sin(n\Arccos(\frac x2))$, 
il s'annule en changeant de signe en chaque 
$\alpha_k = 2\cos(\frac{k\pi}{n})$, ainsi $P'(\alpha_k)=0$ pour $k = 0,\ldots,n$.

On calcule aussi que $P(\alpha_k) = \pm 2$.
Le tableau de signe montre que $P$ est alternativement croissante 
(de $-2$ à $+2$) puis décroissante (de $+2$ à $-2$) 
sur chaque intervalle
$[\alpha_{k+1}, \alpha_k]$, qui forment une partition de $[-2,2]$.
D'après le théorème des valeurs intermédiaires, $P$ possède $n$ racines simples 
(une dans chaque intervalle $[\alpha_{k+1}, \alpha_k]$) dans $[-2,2]$. 
Puisque $P$ est de degré $n$, on a ainsi obtenu toutes ses racines.

\end{enumerate}
\fincorrection
\correction{006962}
\begin{enumerate}
\item Si $k\in\Zz$ est racine de $P$, alors $k^n+a_{n-1}k^{n-1}+\cdots+a_1k=-a_0$ 
ce qui donne $k(k^{n-1}+\cdots+a_1)=-a_0$, donc $k$ divise $a_0$.

\item Si $X^3-X^2-109X-11$ a une racine $k\in\Zz$, nécessairement $k$ divise $11$, 
donc $k$ vaut $-1$, $1$, $-11$ ou $11$. En testant ces quatre valeurs, on trouve 
que seul $11$ est racine. 

De même, si $X^{10}+X^5+1$ admettait une racine entière $k$, celle-ci diviserait $1$ donc vaut 
$k = \pm 1$, or on vérifie que ni $+1$, ni $-1$ ne sont racines.
Ainsi $X^{10}+X^5+1$ n'a pas de racine entière.
\end{enumerate}
\fincorrection
\correction{006963}
On a 
$$L_i(a_i)=\prod_{\substack{1\le j\le n \\ j\not= i}}\frac{a_i-a_j}{a_i-a_j}=1
\qquad \text{ et } \quad L_i(a_j)=0  \text{ si } j\not=i$$
puisque le produit contient un facteur qui est nul: $(a_j-a_j)$. 
Puisque les $L_i$ sont tous de degré $n$, le polynôme $P$ est de degré inférieur ou égal à $n$, et 
$P(a_j)=\sum_{i=0}^nb_iL_i(a_j)=b_i$. 

Il reste à montrer qu'un tel polynôme est unique. Supposons que $Q$ convienne aussi, 
alors $P-Q$ est de degré inférieur ou égal à 
$n$ et s'annule en $n+1$ points (les $a_i$), donc il est identiquement nul, i.e. $P=Q$.

\bigskip

Pour l'application on utilise utilise les polynômes interpolateurs de Lagrange avec
$a_0=0$, $b_0=1$ ; $a_1=1$, $b_1=0$ ; $a_2=-1$, $b_2=-2$ ; $a_3=2$, $b_3=4$. 
On sait qu'un tel polynôme $P(X)$ est unique et s'écrit 
$$P(X)=1\cdot L_0(X)+0\cdot L_1(X)-2\cdot L_2(X)+4L_3(X)$$
où

$$L_0(X)=\frac{(X-1)(X+1)(X-2)}{(0-1)(0+1)(0-2)}=\frac{1}{2}(X^3-2X^2-X+2)$$

$$L_1(X)=\frac{(X-0)(X+1)(X-2)}{(1-0)(1+1)(1-2)}=\frac{-1}{2}(X^3-X^2-2X)$$

$$L_2(X)=\frac{(X-0)(X-1)(X-2)}{(-1-0)(-1-1)(-1-2)}=\frac{-1}{6}(X^3-3X^2+2X)$$

$$L_3(X)=\frac{(X-0)(X-1)(X+1)}{(2-0)(2-1)(2+1)}=\frac{1}{6}(X^3-X)$$

Ainsi :
$$P(X)=\frac{3}{2}X^3-2X^2-\frac{1}{2}X+1.$$

\fincorrection
\nocorrection
\correction{000444}
\begin{enumerate}
  \item
$\frac{X^3-3X^2+X-4}{X-1} = 
X^2-2X-1-\frac{5}{X-1}$.
  \item
$\frac{2X^3+X^2-X+1}{X^2-3X+2} = 
2X+7-\frac{3}{X-1}+\frac{19}{X-2}$.
  \item
$\frac{2X^3+X^2-X+1}{X^2-2X+1} =
2X+5+\frac{3}{(X-1)^2}+\frac{7}{X-1}$.
  \item
$\frac{X^4+2X^2+1}{X^2-1} = 
X^2+3+\frac{2}{X-1}-\frac{2}{X+1}$.
  \item
$\frac{X}{X^2-4} = 
\frac{1/2}{X+2} + \frac{1/2}{X-2}$.
  \item
$\frac{X^5+X^4+1}{X^3-X} = 
X^2+X+1-\frac{1}{X}+\frac{1/2}{X+1}+\frac{3/2}{X-1}$.
  \item
$\frac{X^5+X^4+1}{X(X-1)^4} = 
1 + \frac{1}{X} + \frac{3}{(X-1)^4} + \frac{6}{(X-1)^3} + 
\frac{10}{(X-1)^2} + \frac{4}{X-1}$.
  \item
$\frac{X^5+X^4+1}{(X-1)^3(X+1)^2} = 
1 + \frac{3/4}{(X-1)^3} + \frac{3/2}{(X-1)^2} + \frac{37/16}{X-1} - 
\frac{1/8}{(X+1)^2} -\frac{5/16}{X+1}$.
  \item
$\frac{X^7+3}{(X^2+X+2)^3} = 
X-3 + \frac{7X+13}{(X^2+X+2)^3} - \frac{7X+21}{(X^2+X+2)^2} + 
\frac{14}{X^2+X+2}$.
  \item
$\frac{(3-2 i )X-5+3 i }{X^2+ i  X+2} = 
\frac{2+ i }{X- i }+\frac{1-3 i }{X+2 i }$.
  \item
$\frac{X+ i }{X^2+ i } = 
\frac{\frac{-\sqrt2+2}{4}+\frac{\sqrt2}{4} i }{X-\frac{\sqrt2-\sqrt2 i }{2}} +
\frac{\frac{\sqrt2+2}{4}-\frac{\sqrt2}{4} i }{X-\frac{-\sqrt2+\sqrt2 i }{2}}$.
  \item
$\frac{X}{(X+ i )^2} = 
\frac{1}{X+ i }-\frac{ i }{(X+ i )^2}$.
  \item
$\frac{X^2+1}{X^4+1} = 
\frac{1/2}{X^2+\sqrt{2}X+1} + \frac{1/2}{X^2-\sqrt{2}X+1} =  
\frac{-\frac{\sqrt{2}}{4} i }{X-\frac{\sqrt{2}}{2}-\frac{\sqrt{2}}{2} i } + 
\frac{ \frac{\sqrt{2}}{4} i }{X-\frac{\sqrt{2}}{2}+\frac{\sqrt{2}}{2} i } + 
\frac{ \frac{\sqrt{2}}{4} i }{X+\frac{\sqrt{2}}{2}+\frac{\sqrt{2}}{2} i } + 
\frac{-\frac{\sqrt{2}}{4} i }{X+\frac{\sqrt{2}}{2}-\frac{\sqrt{2}}{2} i }$.
  \item
$\frac{X}{X^4+1} = 
- \frac{\sqrt{2}/4}{X^2+\sqrt{2}X+1} + \frac{\sqrt{2}/4}{X^2-\sqrt{2}X+1} = 
\frac{-\frac{1}{4} i }{X-\frac{\sqrt{2}}{2}-\frac{\sqrt{2}}{2} i } + 
\frac{ \frac{1}{4} i }{X-\frac{\sqrt{2}}{2}+\frac{\sqrt{2}}{2} i } + 
\frac{-\frac{1}{4} i }{X+\frac{\sqrt{2}}{2}+\frac{\sqrt{2}}{2} i } + 
\frac{ \frac{1}{4} i }{X+\frac{\sqrt{2}}{2}-\frac{\sqrt{2}}{2} i }$.
  \item
$\frac{X^2+X+1}{X^4+1} = 
\frac{(2-\sqrt{2})/4}{X^2+\sqrt{2}X+1} + 
\frac{(2+\sqrt{2})/4}{X^2-\sqrt{2}X+1} = 
\frac{-\frac{1+\sqrt{2}}{4} i }{X-\frac{\sqrt{2}}{2}-\frac{\sqrt{2}}{2} i } + 
\frac{ \frac{1+\sqrt{2}}{4} i }{X-\frac{\sqrt{2}}{2}+\frac{\sqrt{2}}{2} i } + 
\frac{-\frac{1-\sqrt{2}}{4} i }{X+\frac{\sqrt{2}}{2}+\frac{\sqrt{2}}{2} i } + 
\frac{ \frac{1-\sqrt{2}}{4} i }{X+\frac{\sqrt{2}}{2}-\frac{\sqrt{2}}{2} i }$.
  \item
$\frac{X^5+X+1}{X^4-1} = 
X + \frac{3/4}{X-1} + \frac{1/4}{X+1} - \frac{X+\frac{1}{2}}{X^2+1} = 
X + \frac{3/4}{X-1} + \frac{1/4}{X+1} + 
\frac{-\frac{1}{2}+\frac{1}{4} i }{X- i } + 
\frac{-\frac{1}{2}-\frac{1}{4} i }{X+ i }$.
  \item
$\frac{X^5+X+1}{X^6-1} = 
\frac{1/2}{X-1} + \frac{1/6}{X+1} + \frac{\frac{1}{3}X-\frac{2}{3}}{X^2-X+1} = 
\frac{1/2}{X-1} + \frac{1/6}{X+1} - 
\frac{\frac{1}{3} j }{X+ j } - \frac{\frac{1}{3} j ^2}{X+ j ^2}$, o\`u on 
a pos\'e de fa\c con standard 
$ j =-\frac{1}{2}+\frac{\sqrt{3}}{2} i $.
  \item
$\frac{X^3-2}{X^4(X^2+X+1)^2} = 
-\frac{2}{X^4}+\frac{4}{X^3}-\frac{2}{X^2}-\frac{3}{X} + 
\frac{X+1}{(X^2+X+1)^2} + \frac{3X+5}{X^2+X+1} = \newline
-\frac{2}{X^4}+\frac{4}{X^3}-\frac{2}{X^2}-\frac{3}{X} + 
\frac{\frac{1}{3} j ^2}{(X- j )^2} + \frac{\frac{1}{3} j }{(X- j ^2)^2} + 
\frac{\frac{3}{2}-\frac{23\sqrt{3}}{18} i }{X- j } + 
\frac{\frac{3}{2}+\frac{23\sqrt{3}}{18} i }{X- j ^2}$, o\`u on 
a pos\'e de fa\c con standard 
$ j =-\frac{1}{2}+\frac{\sqrt{3}}{2} i $.
  \item
$\frac{X}{(X^2+1)(X^2+4)} = 
\frac{\frac{1}{3}X}{X^2+1} - \frac{\frac{1}{3}X}{X^2+4} = 
\frac{1/6}{X- i } + \frac{1/6}{X+ i } - \frac{1/6}{X-2 i } - 
\frac{1/6}{X+2 i }$.
  \item
$\frac{X^2-3}{(X^2+1)(X^2+4)} = 
-\frac{4/3}{X^2+1} + \frac{7/3}{X^2+4} = 
\frac{\frac{2}{3} i }{X- i } + \frac{-\frac{2}{3} i }{X+ i } + 
\frac{-\frac{7}{12} i }{X-2 i } + \frac{\frac{7}{12} i }{X+2 i }$.
\end{enumerate}
\fincorrection
\correction{000445}
Commencer bien
s\^ur par la division suivant les puissances d\'ecroissantes (la faire faire
par les \'etudiants)~: $\Phi=x+1+\Phi_1$ avec $\Phi_1={4x^2-6x+1\over2x^3-x^2}.$\\
Puis factoriser le d\'enominateur et faire donner le type de d\'ecomposition de
$\Phi_1$~:
\begin{equation}
\label{eq11}
\Phi_1={A\over x^2}+{B\over x}+{C\over x-{1\over2}}.
\end{equation}
Expliquer qu'on
obtient alors $A$ en multipliant les deux membres de~(\ref{eq11}) par
$x^2$ et en passant \`a la limite quand $x$ tend vers 0 ($A=-1$). On obtient de
m\^eme $C$ par multiplication par $x-{1\over2}$ et calcul de la limite quand $x$
tend vers ${1\over2}$ ($C=-2$). Enfin on trouve $B$ en identifiant pour une
valeur particuli\`ere non encore utilis\'ee, par exemple $x=1$, ou mieux en
multipliant les deux membres de~(\ref{eq11}) 
par $x$ et en passant \`a la limite pour
$x\to\infty$ ($B=4$). Faire remarquer que pour un cas aussi simple, les calculs
peuvent se faire \emph{de t\^ete} en \'ecrivant simplement les coefficients $A$,
$B$, $C$ au fur et \`a mesure qu'on les obtient.
$${2x^4+x^3+3x^2-6x+1\over2x^3-x^2}=x+1-{1\over x^2}+{4\over x}-{2\over
x-{1\over2}}.$$
\fincorrection
\correction{000446}
La
division suivant les puissances d\'ecroissantes\linebreak %%%%
donne~: $\Phi=2+\Phi_1$ avec
$$
\Phi_1={4x^4-10x^3+8x^2-4x+1\over x^3(x-1)^2}=
{A\over x^3}+{B\over x^2}+{C\over x}+{D\over(x-1)^2}+{E\over x-1}.
$$
Faire remarquer que la m\'ethode de l'exercice pr\'ec\'edent permettrait
d'obtenir facilement $A$ et~$D$ par multiplication par $x^3$ et par $(x-1)^2$,
mais qu'il resterait encore 3 coefficients \`a d\'eterminer. \\Il y a ici une
m\'ethode plus efficace~: effectuer la division suivant les puissances
croissantes, \`a l'ordre~3 
(qui est l'exposant du facteur $x$) du num\'erateur
$1-4x+8x ^2-10x^3+4x^4$ par $(x-1)^2$, ou plut\^ot par $1-2x+x^2$~:
\begin{equation}
\label{eq22}
1-4x+8x ^2-10x^3+4x^4=(1-2x+x^2)\times(1-2x+3x^2)+(-2x^3+x^4).
\end{equation}
En divisant les deux membres de~(\ref{eq22}) par $x^3(x-1)^2$, 
on obtient $A$, $B$
et~$C$ d'un seul coup~:
$$
\Phi_1={1\over x^3}-{2\over x^2}+{3\over x}+{x-2\over(x-1)^2}.
$$
Le calcul de $D$ et $E$ est alors imm\'ediat par d\'ecomposition de
${x-2\over(x-1)^2}$~: m\'ethode de l'exercice pr\'ec\'edent, ou division
suivant les puissances d\'ecroissantes de $x-2$ par $x-1$~: $x-2=(x-1)-1$.
$${2x^5-8x^3+8x^2-4 x+1\over x^3(x-1)^2}=2+
{1\over x^3}-{2\over x^2}+{3\over x}-{1\over(x-1)^2}+{1\over x-1}.$$

\smallskip
Remarque~: cette m\'ethode est efficace pour un exposant assez grand (en gros
\`a partir de~3). Elle peut \^etre utilis\'ee pour une fraction du type 
${P(x)\over(x-a)^nQ(x)}$, mais il faut commencer par le changement de variable
$u=x-a$ avant de faire la division, puis bien entendu revenir ensuite \`a la
variable $x$.
\fincorrection
\correction{000447}
Pas de division pr\'eliminaire dans ce cas\dots{} Forme de la d\'ecomposition~:
\begin{equation}
\label{eq31}
\Phi={A\over x}+{Bx+C\over(x^2+1)^3}+{Dx+E\over(x^2+1)^2}+{Fx+G\over x^2+1}.
\end{equation}
La m\'ethode du premier exercice permet d'obtenir $A$, puis $B$ et~$C$
(pour ces derniers~: multiplication des deux membres de~(\ref{eq31}) par $x^2+1$, puis
limite quand $x$ tend vers $i$ ou vers $-i$, avec s\'eparation des parties
r\'eelle et imaginaire), mais c'est bien insuffisant pour conclure~: il faut
encore soustraire ${Bx+C\over(x^2+1)^3}$, simplifier par $x^2+1$, calculer $D$
et $E$\dots{} (le faire faire quand m\^eme \`a titre d'entra\^inement).

On va ici se contenter de trouver $A$ ($A=3$), puis faire la soustraction
$\Phi_1=\Phi-{A\over x}$. Faire faire le calcul aux \'etudiants~; leur faire
remarquer que, sauf erreur de calcul, la fraction $\Phi_1$ \emph{doit} se
simplifier par $x$. On trouve~:
$$
\Phi={3\over x}+{x^5-2x^4+2x^3-x^2+2x+2\over(x^2+1)^3}.
$$
La fin de la d\'ecomposition se fait par divisions successives suivant les
puissances d\'ecroissantes~: division du num\'erateur $x^5-2x^4+2x^3-x^2+2x+2$
par $x^2+1$, puis du quotient obtenu par $x^2+1$.
$${4x^6-2x^5+11x^4-x^3+11x^2+2x+3\over x(x^2+1)^3}=
{3\over x}+{x+1\over(x^2+1)^3}+{3\over(x^2+1)^2}+{x-2\over x^2+1}.$$

\smallskip
Remarque~: cette m\'ethode des divisions successives est tr\`es pratique quand
la fraction \`a d\'ecomposer a un d\'enominateur \emph{simple}, c'est \`a dire
comportant un d\'enominateur du type $Q^n$ o\`u $Q$ est du premier degr\'e, ou
du second degr\'e sans racine r\'eelle. Faire remarquer aussi comment on
peut simplifier petit \`a petit en \'eliminant du d\'enominateur un
d\'enominateur \emph{simple} (m\'ethode utilis\'ee dans l'exercice~3 par le
calcul de $\Phi-{A\over x}$).
\fincorrection
\nocorrection
\nocorrection
\nocorrection
\correction{003270}
\begin{enumerate}
  \item 
  \item
  \item ssi $\exists\ G \in { K(X)}$ tel que $G\circ F = X  \Rightarrow  P\circ F = XQ\circ F$.
             \par
             $F = \frac AB,\ A \wedge B = 1  \Rightarrow 
              \begin{cases}A\mid(p_0 - Xq_0)\cr B\mid(p_n - Xq_n)\end{cases}  \Rightarrow 
              F$ est homographique.
  \item $F = \phi(X)$.
\end{enumerate}
\fincorrection
\nocorrection
\correction{003272}
$F = \frac PQ$.
Si $P = \lambda Q$ : $\Im F = \{\lambda\}$.\par
Si $P = \lambda Q + \mu$ : $\Im F = \C \setminus \{\lambda\}$.\par
Sinon, $\Im F = \C$.
\fincorrection
\correction{003273}
1) $G = $cste.\par
2) $F$ a un seul p{\^o}le $a$
   $ \Rightarrow  F = \frac P{(X-a)^k}$ et $G = a + \frac 1Q$ avec $\deg P \le k$.\par
3) $F \in {\C[X]}  \Rightarrow  G \in {\C[X]}$.
\fincorrection
\correction{003274}
\begin{enumerate}
  \item 
  \item$n\frac {X^n+1}{X^n-1}$.
\end{enumerate}
\fincorrection
\correction{003275}
\begin{enumerate}
\item $ \Rightarrow \frac{P(\scriptstyle X)}{Q(\scriptstyle X)} =
  \frac{P(\frac 1X)}{Q(\frac 1X)} = \frac{P(\scriptstyle X)+P(\frac
    1X)}{Q(\scriptstyle X)+Q(\frac 1X)}$.
  \item
  \item 
\end{enumerate}
\fincorrection 
\nocorrection
\nocorrection
\correction{003278}
$I_k = \{ F$ tels que $\deg F \le -k \}$.
\fincorrection
\nocorrection
\correction{003280}
$$\begin{aligned}\frac1{(X^2+2X+1)(X^3-1)}
&= \frac{-1/2}{(X+1)^2} + \frac{-3/4}{X+1} + \frac{1/12}{X-1} + \frac{1/3}{X-j} + \frac{1/3}{X-j^2}\cr
&= \frac{-1/2}{(X+1)^2} + \frac{-3/4}{X+1} + \frac{1/12}{X-1} + \frac13\frac{2X+1}{X^2+X+1}.\cr\end{aligned}$$
\fincorrection
\correction{003281}
\begin{enumerate}
  \item $\sum_{k=0}^n \frac {(-1)^{k+p}p!}{k!\,(n-k)!\,(X+k)^{p+1}}$.
  \item $\frac{(-1)^pp!}{2i\sin\alpha}\Bigl(\frac 1{(X-e^{i\alpha})^{p+1}}
             - \frac 1{(X-e^{-i\alpha})^{p+1}}\Bigr)
             = \frac{\sum_{k=0}^p C_{p+1}^kp!\,{(-1)^k}\frac{\sin(p+1-k)\alpha}{\sin\alpha} X^k}
               {(X^2-2X\cos\alpha+1)^{p+1}}$.
  \item $\frac{\sum_{k\text{ pair}} C_{p+1}^kp!\,(-1)^{p+1}
             \frac {\sh k\alpha}{\ch\alpha} X^{p+1-k}
             +\sum_{k\text{ impair}} C_{p+1}^kp!\,(-1)^p
             \frac {\ch k\alpha}{\ch\alpha} X^{p+1-k}}
             {(X^2-2X\sh\alpha-1)^{p+1}}.$
\end{enumerate}
\fincorrection
\correction{003282}
\begin{enumerate}
  \item $1$.
  \item $1/4$.
  \item $1/2$.
\end{enumerate}
\fincorrection
\correction{003283}
$\frac 1{Q(a)(X-a)^2}    - \frac {Q'(a)}{Q^2(a)(X-a)} =
\frac 2{R''(a)(X-a)^2} - \frac {2R'''(a)}{3{R''}^2(a)(X-a)}$.
\fincorrection
\correction{003284}
\begin{enumerate}
  \item $\sum_{i=1}^n \Bigl( \frac{(1+a_i^2)^n}{P'^2(a_i)(X-a_i)^2}
             + \frac{2na_i - P''(a_i)(1+a_i^2)/P'(a_i)}{P'^2(a_i)(X-a_i)} \Bigr)$.
  \item
\end{enumerate}
\fincorrection
\correction{003285}
\begin{enumerate}
  \item D{\'e}composer $1/P$ en {\'e}l{\'e}ments simples, et prendre $x \to \infty$.
  \item Idem avec  $X^k/P  \Rightarrow  \Sigma = \begin{cases}0 &\text{ si } 0 \le k < n-1 \cr
                                                   1 &\text{ si } k = n-1.\end{cases}$

\end{enumerate}
\fincorrection
\correction{003286}
\begin{enumerate}
  \item $P' = \sum_{i=1}^n \frac {m_iP}{X-x_i}  \Rightarrow 
             \frac {P'}P = \sum_{i=1}^n \frac {m_i}{X-x_i}$.
  \item $P'(z) = 0 \Leftrightarrow \sum_{i=1}^n m_i\frac {\overline{z-x_i}}{|z-x_i|^2} = 0
             \Leftrightarrow z = \text{Bar}\biggl( x_i, \frac {m_i}{|z-x_i|^2} \biggr)$.
  \item 
  \item 
\end{enumerate}
\fincorrection
\nocorrection
\correction{003288}
$F(X+1) - F(X) = \frac 2{X-1} - \frac 3X + \frac 1{X+1}
          \Rightarrow 
         F(X) = \frac 1X -\frac 2{X-1} + \text{cste}$.
\fincorrection
\correction{003289}
$F = \sum_{j=1}^n \frac {x_j}{X-b_j} - \frac 1{X-c}
 = \lambda \frac {\prod(X-a_i)}{(X-c)\prod(X-bj)}$ o{\`u}
$\lambda = -\prod \frac {c-b_i}{c-a_i}$.
\fincorrection
\correction{003290}
$Q = (XP + P')(XP' + P) = XP^2\Bigl( X + \frac {P'}P \Bigr) \Bigl(\frac 1X + \frac {P'}P \Bigr)$.

$\frac {P'}P = \sum \frac 1{X-a_i}$, donc les expressions :
$x + \frac {P'(x)}{P(x)}$ et $\frac 1x + \frac {P'(x)}{P(x)}$ changent de signe entre
$a_i$ et $a_{i+1}$.

Cela fait au moins $2n-3$ racines distinctes ($2n-2$ si 1 n'est pas racine),
plus encore une racine pour $\frac 1x + \frac {P'(x)}{P(x)}$ entre
$0$ et $a_1$.
\fincorrection
\correction{003291}
$\frac PQ = \sum_{k=0}^n\frac{P(k)}{(X-k)\prod_{i\ne k}(k-i)}$
donc $\sum_{k=0}^n\frac{P(k)}{\prod_{i\ne k}(k-i)} = \lim_{x\to\infty}\frac{xP(x)}{Q(x)}=1$.

Si l'on suppose $|P(k)|< \frac{n!}{2^n}$ pour tout $k\in{[[0,n]]}$ alors
$\left|\sum_{k=0}^n\frac{P(k)}{\prod_{i\ne k}(k-i)}\right| <
\frac{1}{2^n}\sum_{k=0}^n\frac{n!}{k!\,(n-k)\!} = 1$, contradiction.
\fincorrection
\correction{003292}
\begin{enumerate}
\item On suppose $P\ne 0$.  Soient $\alpha_1,\dots,\alpha_p$ les
  racines de $P$ de multiplicit{\'e}s $m_1,\dots,m_p$ et
  $n=m_1+\dots+m_p=\deg(P)$.  On a $\frac{P'}P =
  \sum_i\frac{m_i}{X-\alpha_i}$ et $\sum_i\frac{-m_i}{(X-\alpha_i)^2}
  = \Bigl(\frac{P'}P\Bigr)' = \frac{P''}P - \Bigl(\frac{P'}P\Bigr)^2 =
  \frac{n(n-1)}{(X-\alpha)(X-\beta)} -
  \Bigl(\sum_i\frac{m_i}{X-\alpha_i}\Bigr)^2$ o{\`u} $\alpha,\beta$ sont
  les deux racines de~$P$ manquant dans~$P''$.  Si
  $\alpha_i\notin\{\alpha,\beta\}$ alors en comparant les termes en
  $1/(X-\alpha_i)^2$ des membres extr{\^e}mes on trouve $m_i=1$. De m{\^e}me
  si $\alpha_i=\alpha\ne\beta$ ou l'inverse. Reste le cas
  $\alpha_i=\alpha=\beta$ qui donne $-m_i = n(n-1) - m_i^2$ donc $m_i
  = n$ ce qui contredit l'hypoth{\`e}se ``$P$ a deux racines distinctes''.
\item Soient $\alpha_i<\alpha_j$ les deux plus petites racines r{\'e}elles
  de~$P$. Si $\alpha_i$ est aussi racine de $P''$ alors $P$ et $P''$
  changent de signe en $\alpha_i$ et, en rempla\c cant au besoin $P$
  par $-P$, $P$ est convexe positif sur $]-\infty,\alpha_i[$ et
  concave n{\'e}gatif sur $]\alpha_i,\alpha_j[$ ce qui est absurde.  Donc
  $\alpha_i\in\{\alpha,\beta\}$. De m{\^e}me pour la plus grande racine
  r{\'e}elle de~$P$, ce qui prouve que $\alpha$ et $\beta$ sont r{\'e}els. En
  identifiant les {\'e}l{\'e}ments de premi{\`e}re esp{\`e}ce dans les deux
  d{\'e}compositions de $P'/P$ on obtient~:
  $$\forall\ i\in{[[1,n]]},\ \sum_{j\ne i}\frac2{\alpha_i-\alpha_j} =
  \begin{cases}n(n-1)/(\alpha-\beta) & \text{ si } \alpha_i=\alpha,\cr
    n(n-1)/(\beta-\alpha) &\text{ si } \alpha_i=\beta,\cr 0 &\text{
      sinon.}\cr\end{cases}$$ En particulier pour
  $\alpha_i\notin\{\alpha,\beta\}$ on a~: $\sum_{j\ne
    i}\frac{\overline{\alpha_i}-\overline{\alpha_j}}{|\alpha_i-\alpha_j|^2}
  = 0$ ce qui signifie que $\overline{\alpha_i}$ est barycentre des
  $\overline{\alpha_j}$ avec des coefficients positifs, donc
  appartient {\`a} l'enveloppe convexe des $\overline{\alpha_j}$, $j\ne
  i$. Il en va de m{\^e}me sans les barres, et donc l'ensemble des racines
  de $P$ n'a pas d'autres points extr{\'e}maux que $\alpha$ et $\beta$~;
  il est inclus dans $[\alpha,\beta]$ donc dans~$\R$.
\end{enumerate}
\fincorrection
\correction{003293}
\begin{enumerate}
  \item $X^3-1 = (X^2+1)(X^3+X^2-1) - X^4(X+1)$.
  \item $F(x) = \ln\Bigl(\frac x{\sqrt{x^2+1}}\Bigr) - \arctan x
              + \frac 1{3x^3} - \frac 1x$.
\end{enumerate}
\fincorrection
\correction{003294}
\begin{enumerate}
  \item $1 = (1-X)^2 (1 + 2X + 3X^2 + \dots + nX^{n-1}) + (n+1)X^n - nX^{n+1}$.
  \item $=\frac {-n\cos n\theta + (n+1)\cos (n-1)\theta -\cos\theta}
                      {4\sin^2{\textstyle{\frac\theta2}}}.$
\end{enumerate}
\fincorrection
\correction{003295}
\begin{enumerate}
  \item $1-X^2 = (1-2X\cos\theta+X^2)(1 + 2X\cos\theta + \dots + 2X^n\cos n\theta)
                    + 2X^{n+1}\cos(n+1)\theta - 2X^{n+2}\cos n\theta$.
  \item $=\frac{\cos n\theta - \cos(n+1)\theta}{1-\cos\theta}$.
\end{enumerate}
\fincorrection
\correction{003296}
\begin{enumerate}
  \item
  \item Division de 1 par $P  \Rightarrow  U = 1-2X+X^2+X^3-X^4,\quad
             V = -1+X^2+X^3+X^4$.
\end{enumerate}
\fincorrection
\correction{005335}
\begin{enumerate}
\item  Soit $F=\frac{X^2+3X+5}{X^2-3X+2}=\frac{X^2+3X+5}{(X-1)(X-2)}$.

$1$ et $2$ ne sont pas racines du polynôme $X^2+3X+5$ et donc, $F$ est bien sous forme irréductible. La partie entière de $F$ étant clairement $1$, $F$ s'écrit sous la forme~: 

$$F=1+\frac{a}{X-1}+\frac{b}{X-2},$$

où $a$ et $b$ sont deux réels.

$a=\lim_{x\rightarrow 1}(x-1)F(x)=\frac{1+3+5}{1-2}=-9$ et $b=\lim_{x\rightarrow 2}(x-2)F(x)=\frac{4+6+5}{2-1}=15$. Donc,
 
$$F=1-\frac{9}{X-1}+\frac{15}{X-2}.$$

\item  Soit $F=\frac{X^2+1}{(X-1)(X-2)(X-3)}$. La décomposition en éléments simples de $F$ s'écrit sous la forme~:

$$F=\frac{a}{X-1}+\frac{b}{X-2}+\frac{c}{X-3},$$

où $a$, $b$ et $c$ sont trois réels.

$a=\lim_{x\rightarrow 1}(x-1)F(x)=\frac{1+1}{(1-2)(1-3)}=1$, puis  $b=\lim_{x\rightarrow 2}(x-2)F(x)=\frac{4+1}{(2-1)(2-3)}=-5$ et

$c=\lim_{x\rightarrow 3}(x-3)F(x)=\frac{9+1}{(3-1)(3-2)}=5$. Donc,

$$F=\frac{1}{X-1}-\frac{5}{X-2}+\frac{5}{X-3}.$$

\item  Soit $F=\frac{1}{X(X-1)^2}$.

$$F=\frac{a}{X}+\frac{b}{X-1}+\frac{c}{(X-1)^2},$$ avec 

$a=\lim_{x\rightarrow 0}xF(x)=1$ et $c=\lim_{x\rightarrow 1}(x-1)^2F(x)=1$. Enfin, $x=-1$ fournit $-1-\frac{b}{2}+\frac{1}{4}=-\frac{1}{4}$ et donc $b=-1$.

Pour trouver $b$, on peut aussi écrire (le meilleur) $0=\lim_{x\rightarrow +\infty}xF(x)=a+b$ et donc que $b=-a=-1$.

On peut encore écrire (le moins bon ici)

$$\frac{1}{X(X-1)^2}-\frac{1}{X}-\frac{1}{(X-1)^2}=\frac{1-(X-1)^2-X}{X(X-1)^2}=\frac{-X^2+X}{X(X-1)^2}
=-\frac{1}{X-1}.$$

Donc,

$$F=\frac{1}{X}-\frac{1}{X-1}+\frac{1}{(X-1)^2}.$$

Autre démarche.

\begin{align*}\ensuremath
\frac{1}{X(X-1)^2}&=\frac{X-1-X}{X(X-1)^2}=\frac{1}{X(X-1)}-\frac{1}{(X-1)^2}
=\frac{X-1-X}{X(X-1)}-\frac{1}{(X-1)^2}\\
 &=\frac{1}{X}-\frac{1}{X-1}+\frac{1}{(X-1)^2}.
\end{align*}

\item  Soit $F=\frac{X^2+1}{(X-1)^2(X+1)^2}$. Puisque $F$ est paire, la décomposition en éléments simples de $F$ est de la forme~:

$$F=\frac{a}{X-1}+\frac{b}{(X-1)^2}-\frac{a}{X+1}+\frac{b}{(X+1)^2}.$$

$b=\lim_{x\rightarrow 1}(x-1)^2F(x)=\frac{1}{2}$ puis, $x=0$ fournit $-2a+2b=1$ et donc $a=0$.
 
$$F=\frac{1}{2}(\frac{1}{(X-1)^2}+\frac{1}{(X+1)^2}).$$

\item  Soit $F=\frac{1}{(X-2)^3(X+2)^3}$. Puisque $F$ est paire, la décomposition en éléments simples de $F$ est de la forme~:

$$F=\frac{a}{X-2}+\frac{b}{(X-2)^2}+\frac{c}{(X-2)^3}-\frac{a}{X+2}+\frac{b}{(X+2)^2}-\frac{c}{(X+2)^3}.$$

$c=\lim_{x\rightarrow 2}(x-2)^3F(x)=\frac{1}{64}$ puis,

\begin{align*}\ensuremath
F-\frac{1}{64}(\frac{1}{(X-2)^3}-\frac{1}{(X+2)^3})&=\frac{64-(X+2)^3+(X-2)^3}{64(X-2)^3(X+2)^3}=
\frac{-12X^2+48}{64(X-2)^3(X+2)^3}\\
 &=-\frac{3}{16}\frac{X^2-4}{(X-2)^3(X+2)^3}=-\frac{3}{16}\frac{1}{(X-2)^2(X+2)^2}
\end{align*}

Puis, $b=\lim_{x\rightarrow 2}(x-2)^2(F(x)-\frac{1}{64}(\frac{1}{(x-2)^3}-\frac{1}{(x+2)^3})=-\frac{3}{16}\frac{1}{(2+2)^2}
=-\frac{3}{256}$. Enfin, $x=0$ fournit $-\frac{1}{64}=-a-\frac{3}{512}-\frac{1}{256}$ et $a=\frac{1}{64}-\frac{5}{512}=\frac{3}{512}$. Donc,

$$F=\frac{1}{512}(\frac{3}{X-2}-\frac{6}{(X-2)^2}+\frac{8}{(X-2)^3}-\frac{3}{X+2}-\frac{6}{(X+2)^2}-\frac{8}{(X+2)^3}).$$

\item  Soit $F=\frac{X^6}{(X^3-1)^2}$. On a déjà $(X^3-1)^2=(X-1)^2(X-j)^2(X-j^2)^2$. Puisque $F$ est réelle, la décomposition en éléments simples de $F$ s'écrit

$$F=1+\frac{a}{X-1}+\frac{b}{(X-1)^2}+\frac{c}{X-j}+\frac{d}{(X-j)^2}+\frac{\overline{c}}{X-j^2}+\frac{\overline{d}}{(X-j^2)^2}.$$

$b=\lim_{z\rightarrow 1}(z-1)^2F(z)=\frac{1}{9}$ et 

\begin{align*}\ensuremath
d&=\lim_{z\rightarrow j}(z-j)^2F(z)=\frac{j^6}{(j-1)^2(j-j^2)^2}=\frac{1}{j^2(j-1)^4}=\frac{1}{j^2(j^2-2j+1)^2}\\
 &= \frac{1}{j^2(-3j)^2}=\frac{j^2}{9}
\end{align*}

Puis,

\begin{align*}\ensuremath
\frac{1}{(X-1)^2}+\frac{j^2}{(X-j)^2}+\frac{j}{(X-j^2)^2}&=\frac{1}{(X-1)^2}+\frac{(j+j^2)X^2-2(j+j^2)X+2}{(X-j)^2(X-j^2)}\\
 &=\frac{1}{(X-1)^2}+\frac{-X^2+2X+2}{(X-j)^2(X-j^2)}=\frac{(X^2+X+1)^2+(X-1)^2(-X^2+2X+2)}{(X^3-1)^2}\\
 &=\frac{6X^3+3}{(X^3-1)^2}
\end{align*}

Par suite,

\begin{align*}\ensuremath
F-1-\frac{1}{9}( \frac{1}{(X-1)^2}+\frac{j^2}{(X-j)^2}+\frac{j}{(X-j^2)^2})&=\frac{X^6}{(X^3-1)^2}-1-\frac{2X^3+1}{3(X^3-1)^2}\\
 &=\frac{3X^6-3(X^3-1)^2-2X^3-1}{3(X^3-1)^2}=\frac{4X^3-4}{3(X^3-1)^2}\\
 &=\frac{4}{3}\frac{1}{X^3-1}.
\end{align*}

Mais alors, $a=\lim_{z\rightarrow 1}(z-1)(F(z)-1-\frac{1}{9}(\frac{1}{(z-1)^2}+\frac{j^2}{(z-j)^2}+\frac{j}{(z-j^2)^2})
=\frac{4}{3}\frac{1}{1+1+1}=\frac{4}{9}$. De même,

$$c=\lim_{z\rightarrow j}(z-j)(F(z)-1-\frac{1}{9}(\frac{1}{(z-1)^2}+\frac{j^2}{(z-j)^2}+\frac{j}{(z-j^2)^2})
=\frac{4}{3}\frac{1}{(j-1)(j-j^2)}=\frac{4j^2}{9}.$$
  
Donc,

$$F=1+\frac{1}{9}(\frac{4}{X-1}+\frac{1}{(X-1)^2}+\frac{4j^2}{X-j}+\frac{j^2}{(X-j)^2}+\frac{4j}{X-j^2}+\frac{j}{(X-j^2)^2}).$$

Si on veut la décomposition sur $\Rr$, on peut regrouper les conjugués~:

\begin{align*}\ensuremath
F&=1+\frac{1}{9}(\frac{4}{X-1}+\frac{1}{(X-1)^2}+\frac{4j^2(X-j^2)+4j(X-j)}
{X^2+X+1}+\frac{j^2(X-j^2)^2+j(X-j)^2}{(X^2+X+1)^2})\\
 &=1+\frac{1}{9}(\frac{4}{X-1}+\frac{1}{(X-1)^2}+\frac{-4X+4}{X^2+X+1}+\frac{-X^2+2X+2}{(X^2+X+1)^2})\\
 &=1+\frac{1}{9}(\frac{4}{X-1}+\frac{1}{(X-1)^2}+\frac{-4X+4}{X^2+X+1}+\frac{-X^2-X-1+3X+3}{(X^2+X+1)^2})\\
 &=1+\frac{1}{9}(\frac{4}{X-1}+\frac{1}{(X-1)^2}+\frac{-4X+3}{X^2+X+1}+\frac{3X+3}{(X^2+X+1)^2})
\end{align*}

\item  Soit $F=\frac{1}{X^6+1}$. 

$$F=\sum_{k=0}^{5}\frac{\lambda_k}{X-\omega_k},$$ 

où $\omega_k=e^{i(\frac{\pi}{6}+\frac{2k\pi}{6})}$. Mais, 

$$\lambda_k=\frac{1}{6\omega_k^5}=\frac{\omega_k}{6\omega_k^6}=-\frac{\omega_k}{6}.$$

Donc,

$$\frac{1}{X^6+1}=\frac{1}{6}(-\frac{i}{X-i}+\frac{i}{X+i}-\frac{e^{i\pi/6}}{X-e^{i\pi/6}}-\frac{e^{-i\pi/6}}{X-e^{-i\pi/6}}+
\frac{e^{i\pi/6}}{X+e^{i\pi/6}}+\frac{e^{-i\pi/6}}{X+e^{-i\pi/6}}).$$
\item  Soit $F=\frac{X^2+3}{X^5-3X^4+5X^3-7X^2+6X-2}$.

\begin{align*}\ensuremath
X^5-3X^4+5X^3-7X^2+6X-2&=(X-1)(X^4-2X^3+3X^2-4X+2)=(X-1)^2(X^3-X^2+2X-2)\\
 &= (X-1)^2(X^2(X-1)+2(X-1))=(X-1)^3(X^2+2).
\end{align*}

La décomposition en éléments simples de $F$ est donc de la forme

$$F=\frac{a}{X-1}+\frac{b}{(X-1)^2}+\frac{c}{(X-1)^3}+\frac{d}{X-i\sqrt{2}}+\frac{\overline{d}}{X+i\sqrt{2}}.$$

Puis,

\begin{align*}\ensuremath
d&=\lim_{z \rightarrow i\sqrt{2}}(z-i\sqrt{2})F(z)=\frac{(i\sqrt{2})^2+3}{(i\sqrt{2}-1)^3(i\sqrt{2}+i\sqrt{2})}=
\frac{1}{(2i\sqrt{2})(-2i\sqrt{2}+6+3i\sqrt{2}-1)}=\frac{1}{-4+10i\sqrt{2}}\\
 &=-\frac{2+5i\sqrt{2}}{108}.
\end{align*}

Ensuite,

$$\frac{d}{X-i\sqrt{2}}+\frac{\overline{d}}{X+i\sqrt{2}}=-\frac{1}{108}\frac{(2+5i\sqrt{2})(X+i\sqrt{2})+
(2-5i\sqrt{2})(X-i\sqrt{2})}{X^2+2}=-\frac{1}{108}\frac{4X-20}{X^2+2}=\frac{-X+5}{27(X^2+2)}.$$

Mais alors,

\begin{align*}\ensuremath
\frac{a}{X-1}+\frac{b}{(X-1)^2}+\frac{c}{(X-1)^3}&=\frac{X^2+3}{(X-1)^3(X^2+2)}-\frac{-X+5}{27(X^2+2)}\\
 &=\frac{27(X^2+3)-(-X+5)(X-1)^3}{27(X-1)^3(X^2+2)}=\frac{X^4-8X^3+45X^2-16X+86}{27(X-1)^3(X^2+2)}\\
 &=\frac{(X^2+2)(X^2-8X+43)}{27(X-1)^3(X^2+2)}=\frac{X^2-8X+43}{27(X-1)^3}\\
 &=\frac{X^2-2X+1-6X+6+36}{27(X-1)^3}\\
 &=\frac{1}{27}(\frac{1}{X-1}-6\frac{1}{(X-1)^2}+\frac{36}{(X-1)^3}).
\end{align*}

Finalement,

$$F=\frac{1}{27}(\frac{1}{X-1}-6\frac{1}{(X-1)^2}+\frac{36}{(X-1)^3})-\frac{1}{108}(\frac{2+5i\sqrt{2}}{X-i\sqrt{2}}+\frac{2-5i\sqrt{2}}{X+i\sqrt{2}}).$$

\item  Soit $F=\frac{X}{(X^2+1)^3(X^2-1)}$. Puisque $F$ est réelle et impaire, la décomposition en éléments simples de $F$ est de la forme

$$F=\frac{a}{X-1}+\frac{a}{X+1}+\frac{b}{X-i}+\frac{c}{(X-i)^2}+\frac{d}{(X-i)^3}+\frac{\overline{b}}{X+i}+\frac{\overline{c}}{(X+i)^2}+\frac{\overline{d}}{(X+i)^3}.$$

$a=\lim_{x\rightarrow 1}(x-1)F(x)=\frac{1}{(1+1)^3(1+1)}=\frac{1}{16}$. Puis,

\begin{align*}\ensuremath
F-\frac{1}{16}(\frac{1}{X-1}+\frac{1}{X+1})&=\frac{8X-X(X^2+1)^3}{8(X^2+1)^3(X^2-1)}
=\frac{-X^7-3X^5-3X^3+7X}{8(X^2+1)^3(X^2-1)}\\
 &=\frac{X(X^2-1)(-X^4-4X^2-7)}{8(X^2+1)^3(X^2-1)}=-\frac{1}{8}\frac{X^4+4X^2+7}{(X^2+1)^3}
\end{align*}

Mais alors, 

\begin{align*}\ensuremath
d&=\lim_{x\rightarrow i}(x-i)^3F(x)=\lim_{x\rightarrow i}(x-i)^3(F(x)-\frac{1}{16}(\frac{1}{x-1}+\frac{1}{x+1}))\\
 &=-\frac{1}{8}\frac{i^4+4i^2+7}{(i+i)^3}=-\frac{i}{16}
\end{align*}

Puis,

\begin{align*}\ensuremath
-\frac{1}{8}\frac{X^4+4X^2+7}{(X^2+1)^3}+\frac{i}{16}\frac{1}{(X-i)^3}-\frac{i}{16}\frac{1}{(X+i)^3}=
-\frac{1}{8}\frac{X^4+4X^2+7}{(X^2+1)^3}+\frac{1}{8}\frac{3X^2-1}{(X^2+1)^3}=\frac{X^2+6}{8(X^2+1)^2}.
\end{align*}

Ensuite, $c=\frac{i^2+6}{8(i+i)^2}=-\frac{5}{32}$. Puis,

$$\frac{X^2+6}{8(X^2+1)^2}+\frac{5}{32}(\frac{1}{(X-i)^2}+\frac{1}{(X+i)^2})
=\frac{2(X^2+6)+5(X^2-1)}{16(X^2+1)^2}=\frac{7}{16(X^2+1)}.$$

Enfin, $b=\frac{7}{16(i+i)}=-\frac{7i}{32}$. Finalement,

$$F=\frac{1}{16}(\frac{1}{X-1}+\frac{1}{X+1})-\frac{7i}{32}(\frac{1}{X-i}-\frac{1}{X+i})-\frac{5}{32}(\frac{1}{(X-i)^2}+\frac{1}{(X+i)^2})-\frac{i}{16}(\frac{1}{(X-i)^3}-\frac{1}{(X+i)^3}).$$

\item  Soit $F=\frac{X^6+1}{X^5-X^4+X^3-X^2+X-1}$.

\begin{align*}\ensuremath
X^5-X^4+X^3-X^2+X-1&=X^4(X-1)+X^2(X-1)+(X-1)=(X-1)((X^4+2X^2+1)-X^2)\\
 &=(X-1)(X^2+X+1)(X^2-X+1)\\
 &=(X-1)(X-j)(X-j^2)(X+j)(X+j^2).
\end{align*}

Püisque $F$ est réelle, la décomposition en éléments simples de $F$ est de la forme

$$F=aX+b+\frac{c}{X-1}+\frac{d}{X-j}+\frac{\overline{d}}{X-j^2}+\frac{e}{X+j}+\frac{\overline{e}}{X-j^2}.$$

$a=\lim_{x\rightarrow +\infty}\frac{F(x)}{x}=1$, puis $b=\lim_{x\rightarrow +\infty}(F(x)-x)=\lim_{x\rightarrow +\infty}\frac{x^5...}{x^5...}=1$. Puis, $c=\frac{1^6+1}{5-4+3-2+1}=\frac{2}{3}$,
$d=\frac{j^6+1}{5j^4-4j^3+3j^2-2j+1}=\frac{2}{3j^2+3j-3}=\frac{2}{-6}=-\frac{1}{3}$ et $e=\frac{(-j)^6+1}{5j^4+4j^3+3j^2+2j+1}=\frac{2}{3j^2+7j+5}=\frac{1}{2j+1}$. Donc,

$$F=X+1+\frac{2}{3}\frac{1}{X-1}-\frac{1}{3}\frac{1}{X-j}-\frac{1}{3}\frac{1}{X-j^2}+\frac{1}{2j+1}\frac{1}{X+j}+\frac{1}{2j^2+1}\frac{1}{X-j^2}.$$

\item  Soit $F=\frac{X^7+1}{(X^2+X+1)^3}$.

La décomposition sur $\Rr$ (hors programme) s'obtiendrait de la façon suivante

\begin{align*}\ensuremath
X^7+1&=(X^2+X+1)(X^5-X^4+X^2-X)+X+1\\
 &=(X^2+X+1)[(X^2+X+1)(X^3-2X^2+X+2)-4X-2]+X+1\\
 &=(X^2+X+1)^2(X^3-2X^2+X+2)-(4X+2)(X^2+X+1)+X+1\\
 &=(X^2+X+1)^2[(X^2+X+1)(X-3)+3X+5]-(4X+2)(X^2+X+1)+X+1\\
 &=X+1-(4X+2)(X^2+X+1)+(3X+5)(X^2+X+1)^2+(X-3)(X^2+X+1)^3
\end{align*}

Donc,

$$F=X-3+\frac{3X+5}{X^2+X+1}-\frac{4X+2}{(X^2+X+1)^2}+\frac{X+1}{(X^2+X+1)^3}.$$

\item  Soit $F=\frac{X^2+1}{X(X-1)^4(X^2-2)^2}$. La décomposition de $F$ en éléments simples est de la forme

$$F=\frac{a}{X}+\frac{b_1}{X-1}+\frac{b_2}{(X-1)^2}+\frac{b_3}{(X-1)^3}+\frac{b_4}{(X-1)^4}+\frac{c_1}{X-\sqrt{2}}+\frac{c_2}{(X-\sqrt{2})^2}+\frac{d_1}{X+\sqrt{2}}+\frac{d_2}{(X+\sqrt{2})^2}.$$

$a=\lim_{x\rightarrow 0}xF(x)=\frac{1}{4}$. Puis, 

\begin{align*}\ensuremath
c_2&=\lim_{x \rightarrow \sqrt{2}}(x-\sqrt{2})^2F(x)=\frac{2+1}{\sqrt{2}(\sqrt{2}-1)^4(\sqrt{2}+\sqrt{2})^2}
=\frac{3}{8\sqrt{2}(4-8\sqrt{2}+12-4\sqrt{2}+1)}\\
 &=\frac{3}{8\sqrt{2}(17-12\sqrt{2})}=\frac{3}{8(-24+17\sqrt{2})}=\frac{3}{16}(24+17\sqrt{2}).
\end{align*}

Un calcul conjugué fournit $d_2=\frac{3}{16}(24-17\sqrt{2})$. On a ensuite

\begin{align*}\ensuremath
\frac{3}{16}(\frac{24+17\sqrt{2}}{(X-\sqrt{2})^2}+\frac{24-17\sqrt{2}}{(X+\sqrt{2})^2})&=\frac{3}{2}\frac{6X^2+17X+12}{(X^2-2)^2}.
\end{align*}

Puis,

\begin{align*}\ensuremath
F-\frac{3}{2}\frac{6X^2+17X+12}{(X^2-2)^2}&=\frac{2(X^2+1)-3(6X^2+17X+12)X(X-1)^4}{2X(X-1)^4(X^2-2)^2}\\
 &=\frac{-18X^7+21X^6+60X^5-90X^4-30X^3+95X^2-36X+2}{2X(X-1)^4(X^2-2)^2}\\
 &=\frac{-18X^5+21X^4+24X^3-48X^2+18X-1}{2X(X-1)^4(X^2-2)}
\end{align*}

Mais alors,

$$c_1=\frac{-18.4\sqrt{2}+21.4+24.2\sqrt{2}-48.2+18\sqrt{2}-1}{2\sqrt{2}(\sqrt{2}-1)^4(\sqrt{2}+\sqrt{2})}
=\frac{-13-6\sqrt{2}}{8(17-12\sqrt{2})}=-\frac{1}{8}(365+258\sqrt{2}),$$

et par un calcul conjugué, $d_1=-\frac{1}{8}(365-258\sqrt{2})$. Ensuite,

$$-\frac{1}{8}(\frac{365+258\sqrt{2}}{X-\sqrt{2}}+\frac{365-258\sqrt{2}}{X+\sqrt{2}})=
-\frac{1}{4}\frac{365X+516}{X^2-2}.$$

Puis,

\begin{align*}\ensuremath
F-\frac{3}{2}\frac{6X^2+17X+12}{(X^2-2)^2}&+\frac{1}{4}\frac{365X+516}{X^2-2}=
\frac{-18X^5+21X^4+24X^3-48X^2+18X-1}{2X(X-1)^4(X^2-2)}+\frac{365X+516}{4(X^2-2)}\\
 &=\frac{2(-18X^5+21X^4+24X^3-48X^2+18X-1)+(365X+516)X(X-1)^4}{4X(X-1)^4(X^2-2)}\\
 &=\frac{365X^6-980X^5+168X^4+1684X^3-1795X^2+552X-2}{4X(X-1)^4(X^2-2)}\\
 &=\frac{365X^4-980X^3+898X^2-276X+1}{4X(X-1)^4}
\end{align*}

Ensuite,

\begin{align*}\ensuremath
F&-\frac{3}{2}\frac{6X^2+17X+12}{(X^2-2)^2}+\frac{1}{4}\frac{365X+516}{X^2-2}-\frac{1}{4X}=
\frac{365X^4-980X^3+898X^2-276X+1}{4X(X-1)^4}-\frac{1}{4X}\\
 &=\frac{(365X^4-980X^3+898X^2-276X+1)-(X-1)^4}{4X(X-1)^4}=\frac{364X^4-976X^3+892X^2-272X}{4X(X-1)^4}\\
 &=\frac{182X^3-488X^2+446X-136}{2(X-1)^4}
\end{align*}

Enfin, $b_4=\lim_{x\rightarrow 1}(x-1)^4\frac{182x^3-488x^2+446x-136}{2(x-1)^4}=2$, puis

$$\frac{182X^3-488X^2+446X-136}{2(X-1)^4}-\frac{2}{(X-1)^4}=\frac{91X^3-244X^2+223X-70}{(X-1)^4}
=\frac{91X^2-153X+70}{(X-1)^3}.$$

Puis, $b_3=\lim_{x\rightarrow 1}(x-1)^3\frac{91x^2-153x+70}{(x-1)^3}=8$, puis

\begin{align*}\ensuremath
\frac{91X^2-153X+70}{(X-1)^3}-\frac{8}{(X-1)^3}&=\frac{91X^2-153X+62}{(X-1)^3}=\frac{91X-62}{(X-1)^2}\\
 &=\frac{91X-91+29}{(X-1)^2}=\frac{91}{X-1}+\frac{29}{(X-1)^2}
\end{align*}

ce qui fournit $b_2=29$ et $b_1=91$.

\begin{align*}\ensuremath
F&=\frac{1}{4X}+\frac{91}{X-1}+\frac{29}{(X-1)^2}+\frac{8}{(X-1)^3}+\frac{2}{(X-1)^4}
-\frac{365+258\sqrt{2}}{8}\frac{1}{X-\sqrt{2}}+\frac{3(24+17\sqrt{2})}{16}\frac{1}{(X-\sqrt{2})^2}\\
 &-\frac{365-258\sqrt{2}}{8}\frac{1}{X+\sqrt{2}}+\frac{3(24-17\sqrt{2})}{16}\frac{1}{(X+\sqrt{2})^2}.
\end{align*}

\item  Soit $F=\frac{1}{(X+1)^7-X^7-1}$.

\begin{align*}\ensuremath
(X+1)^7-X^7-1&=7X^6+21X^5+35X^4+35X^3+21X^2+7X=7X(X^5+3X^4+5X^3+5X^2+3X+1)\\
 &=7X(X+1)(X^4+2X^3+3X^2+2X+1)=7X(X+1)(X^2+X+1)^2
\end{align*}

Si on n'a pas deviné que $X^4+2X^3+3X^2+2X+1=(X^2+X+1)^2$ (par exemple, en repérant que $j$ est racine ou encore en manipulant l'identité $(a+b+c)^2=a^2+b^2+c^2+2(ab+ac+bc)$), on peut pratiquer comme suit

\begin{align*}\ensuremath
X^4+2X^3+3X^2+2X+1&=X^2(X^2+\frac{1}{X^2}+2(X+\frac{1}{X})+3)=X^2((X+\frac{1}{X})^2+2(X+\frac{1}{X})+1)\\
 &=X^2(X+\frac{1}{X}+1)^2=(X^2+X+1)^2
\end{align*}
La déccomposition en éléments simples de $7F$ est donc de la forme

$$7F=\frac{a}{X}+\frac{b}{X+1}+\frac{c}{X-j}+\frac{d}{(X-j)^2}+\frac{\overline{c}}{X-j^2}+\frac{\overline{d}}{(X-j^2)^2}.$$

$a=\frac{1}{(0+1)(0^2+0+1)^2}=1$ et $b=\frac{1}{(-1)(1-1+1)^2}=-1$. Puis,

$$d=\frac{1}{j(j+1)(j-j^2)^2}=\frac{1}{j(-j^2)j^2(1-2j+j^2)}=\frac{1}{-j^2(-3j)}=\frac{1}{3}.$$

Ensuite, 

$$\frac{1}{3}(\frac{1}{(X-j)^2}+\frac{1}{(X-j^2)^2})=\frac{(X-j^2)^2+(X-j)^2}{3(X^2+X+1)^2}
=\frac{2X^2+2X-1}{3(X^2+X+1)^2}.$$

Puis,

\begin{align*}\ensuremath
7F-\frac{2X^2+2X-1}{3(X^2+X+1)^2}&=\frac{3-X(X+1)(2X^2+2X-1)}{3X(X+1)(X^2+X+1)^2}
=\frac{-2X^4-4X^3-4X^2+X+3}{3X(X+1)(X^2+X+1)^2}\\
 &=\frac{-2X^2-2X+3}{3X(X+1)(X^2+X+1)}.
\end{align*}

Mais alors, 

$$c=\frac{-2j^2-2j+3}{3j(j+1)(j-j^2)}=\frac{5}{-3(j-j^2)}=\frac{5i}{3\sqrt{3}}.$$

Finalement,

$$F=\frac{1}{7}(\frac{1}{X}-\frac{1}{X+1}+\frac{5i}{3\sqrt{3}(X-j)}+\frac{3}{3(X-j)^2}-\frac{5i}{3\sqrt{3}(X-j^2)}+\frac{1}{3(X-j^2)^2}).$$

\end{enumerate}
\fincorrection
\correction{005336}
\begin{enumerate}
\item  Soit $P=X^n-1$ et $F=\frac{1}{P}$. La partie entière de $F$ est nulle et les pôles de F sont simples (car $P=X^n-1$ et $P'=nX^{n-1}$ n'ont pas de racines communes dans $\Cc$). De plus, $P=\prod_{k=0}^{n-1}(X-\omega_k)$ où $\omega_k=e^{2ik\pi/n}$. Donc, $F=\sum_{k=0}^{n-1}\frac{\lambda_k}{X-\omega_k}$ où 

$$\lambda_k=\frac{1}{P'(\omega_k)}=\frac{1}{n\omega_k^{n-1}}=\frac{\omega_k}{n\omega_k^n}=\frac{\omega_k}{n}.$$

Ainsi,

$$\frac{1}{X^n-1}=\frac{1}{n}\sum_{k=0}^{n-1}\frac{e^{2ik\pi/n}}{X-e^{2ik\pi/n}}.$$

\item   Soit $P=(X-1)(X^n-1)=(X-1)^2\prod_{k=1}^{n-1}\omega_k$ où $\omega_k=e^{2ik\pi/n}$. Soit $F=\frac{1}{P}$. La partie entière de $F$ est nulle. D'autre part, F admet un pôle double, à savoir $1$ et $n-1$ pôles simples à savoir les $\omega_k=e^{2ik\pi/n}$, $\leq k\leq n-1$. Donc, 

$$F=\frac{a}{X-1}+\frac{b}{(X-1)^2}+\sum_{k=1}^{n-1}\frac{\lambda_k}{X-\omega_k}.$$ 

$\lambda_k=\frac{1}{(n+1)\omega_k^n-n\omega_k^{n-1}-1}=\frac{1}{n(1-\omega_k^{n-1})}=\frac{\omega_k}{n(\omega_k-1)}$. Ensuite, 

$$b=\lim_{x\rightarrow 1}(x-1)^2F(x)=\frac{1}{1^{n-1}+...+1^1+1}=\frac{1}{n}.$$

Il reste à calculer $a$.

\begin{align*}\ensuremath
F-\frac{1}{n(X-1)^2}&=\frac{n-(X^{n-1}+X^{n-2}+...+X+1)}{n(X-1)^2(X^{n-1}+...+X+1)}
=\frac{-X^{n-2}-2X^{n-3}-...-(n-2)X-(n-1))}{n(X-1)(X^{n-1}+...+X+1)}.
\end{align*}

Donc, $a=\lim_{x\rightarrow 1}(x-1)(F(x)-\frac{1}{n(x-1)^2})=\frac{-[(n-1)+(n-2)+...+2+1]}{n(1+1...+1)}=-\frac{n-1}{2n}$.

Finalement,

$$F=\frac{1}{n}(-\frac{n-1}{2n(X-1)}+\frac{1}{(X-1)^2}+\sum_{k=1}^{n-1}\frac{\omega_k}{\omega_k-1}\frac{1}{X-\omega_k}).$$

\item  $\frac{n!}{(X-1)...(X-n)}=\sum_{k=1}^{n}\frac{\lambda_k}{X-k}$ avec

$$\lambda_k=\lim_{x\rightarrow k}(x-k)F(x)=\frac{n!}{\prod_{j\neq k}^{}(j-k)}=\frac{n!}{(-1)^{n-k}(k-1)!(n-k)!}=n(-1)^{n-k}C_{n-1}^{k-1}.$$

Donc,

$$\frac{n!}{(X-1)...(X-n)}=\sum_{k=1}^{n}\frac{(-1)^{n-k}nC_{n-1}^{k-1}}{X-k}.$$

\item  Posons $P=X^4-2X^2\cos(2a)+1$.

\begin{align*}\ensuremath
X^4-2X^2\cos(2a)+1&=(X^2-e^{2ia})(X^2-e^{-2ia})=(X-e^{ia})(X-e^{-ia})(X+e^{ia})(X+e^{-ia})\\
 &(=(X^2-2X\cos a+1)(X^2+2X\cos a+1)).
\end{align*}

$P$ est à racines simples si et seulement si $e^{ia}\neq\pm e^{-ia}$ ce qui équivaut à $a\notin\frac{\pi}{2}\Zz$.

\begin{itemize}
\item[1er cas.] Si $a\in\pi\Zz$,

$$F=\frac{X^2}{(X^2-1)^2}=\frac{a}{X-1}+\frac{b}{(X-1)^2}-\frac{a}{X+1}+\frac{b}{(X+1)^2}.$$

$b=\lim_{x\rightarrow 1}(x-1)^2F(x)=\frac{1}{4}$ puis $x=0$ fournit $0=-2a+2b$ et donc $a=b=\frac{1}{4}$.

$$F=\frac{X^2}{(X^2-1)^2}=\frac{1}{4}(\frac{1}{X-1}+\frac{1}{(X-1)^2}-\frac{1}{X+1}+\frac{1}{(X+1)^2}).$$

\item[2ème cas.] Si $a\in\frac{\pi}{2}+\pi\Zz$,

$$F=\frac{X^2}{(X^2+1)^2}=\frac{a}{X-i}+\frac{b}{(X-i)^2}-\frac{a}{X+i}+\frac{b}{(X+i)^2}.$$

$b=\lim_{x\rightarrow i}(x-i)^2F(x)=\frac{i^2}{(i+i)^2}=\frac{1}{4}$ puis $x=0$ fournit $0=2ia-2b$ et donc $a=-ib=-\frac{i}{4}$.

$$F=\frac{X^2}{(X^2+1)^2}=\frac{1}{4}(-\frac{i}{X-i}+\frac{1}{(X-i)^2}+\frac{i}{X+i}+\frac{1}{(X+i)^2}).$$

\item[3ème cas.] Si $a\notin\frac{\pi}{2}\Zz$, puisque $F$ est réelle et paire,

$$F=\frac{A}{X-e^{ia}}+\frac{\overline{A}}{X-e^{-ia}}-\frac{A}{X+e^{ia}}-\frac{\overline{A}}{X+e^{-ia}},$$

avec 

$$A=\frac{e^{2ia}}{(e^{ia}-e^{-ia})(e^{ia}+e^{ia})(e^{ia}+e^{-ia})}=\frac{e^{2ia}}{8i\sin a\cos ae^{ia}}=\frac{-ie^{ia}}{4\sin(2a)}.$$

Donc,

$$F=\frac{1}{4\sin(2a)}(-\frac{ie^{ia}}{X-e^{ia}}+\frac{ie^{-ia}}{X-e^{-ia}}+\frac{ie^{ia}}{X+e^{ia}}+\frac{ie^{-ia}}{X+e^{-ia}}).$$

\end{itemize}

\item  Le polynôme $X^{2n}+1=\prod_{k=0}^{2n-1}(X-e^{i(\frac{\pi}{2n}+\frac{2k\pi}{2n})})$ est à racines simples car n'a pas de racine commune avec sa dérivée. En posant $\omega_k=e^{i(\frac{\pi}{2n}+\frac{2k\pi}{2n})}$, on a

$$\frac{1}{X^{2n}+1}=\sum_{k=0}^{2n-1}\frac{\lambda_k}{X-\omega_k},$$

où 

$$\lambda_k=\frac{1}{2n\omega_k^{2n-1}}=\frac{\omega_k}{2n\omega_k^{2n}}=-\frac{\omega_k}{2n}.$$

Finalement,

$$\frac{1}{X^{2n}+1}=-\frac{1}{2n}\sum_{k=0}^{2n-1}\frac{\omega_k}{X-\omega_k}.$$
\end{enumerate}

\fincorrection
\correction{005337}
Pour $k$ élément de $\{0,...,n-1\}$, posons $\omega_k=e^{2ik\pi/n}$. Décomposons $F$ en éléments simples (sur $\Cc$).

$$\frac{\omega X+1}{\omega^2X+\omega X+1}=\frac{\omega X+1}{(\omega X)^2+\omega X+1}=\frac{\omega X+1}{(\omega X-j)(\omega X-j^2)}=\frac{a}{\omega X-j}+\frac{b}{\omega X-j^2},$$
 
avec $a=\frac{\omega\frac{j}{\omega}}{\omega\frac{j}{\omega}-j^2}=\frac{j+1}{j-j^2}=-\frac{-j^2}{j-j^2}=\frac{j}{j-1}$ et de même $b=\frac{j^2+1}{j^2-j}=-\frac{1}{j-1}$.
Donc,

\begin{align*}\ensuremath
F&=\frac{1}{j-1}\sum_{k=0}^{n-1}(\frac{j}{\omega_k X-j}-\frac{1}{\omega_k X-j^2})=\frac{1}{j-1}\sum_{k=0}^{n-1}(\frac{j\omega_{-k}}{X-j\omega_{-k}}-\frac{\omega_{-k}}{X-j^2\omega_{-k}})
 &=\frac{1}{j-1}\sum_{k=0}^{n-1}(\frac{j\omega_{k}}{X-j\omega_{k}}-\frac{\omega_{k}}{X-j^2\omega_{k}})
\end{align*}
 
Maintenant les $n$ nombres $j\omega_k$ sont deux à deux distincts et vérifient $(j\omega_k)^n=j^n$ et donc,

$$\prod_{k=0}^{n-1}(X-j\omega_k)=X^n-j^n.$$

$\sum_{k=0}^{n-1}\frac{j\omega_k}{X-j\omega_k}$ est donc la décomposition en éléments simples d'une fraction du type $\frac{P}{X^n-j^n}$ avec $\mbox{deg}P\leq n-1$. De plus, on sait que $j\omega_k=\frac{P(j\omega_k)}{n(j\omega_k)^{n-1}}$ et donc, $\forall k\in\{0,...,n-1\},\;P(j\omega_k)=nj^n$.
Le polynôme $P-nj^n$ est de degré inférieur ou égal à $n-1$, admet les $n$ racines deux à deux distinctes $j\omega_k$ et est donc le polynôme nul. Par suite 

$$\sum_{k=0}^{n-1}\frac{j\omega_{k}}{X-j\omega_{k}}=\frac{nj^n}{X^n-j^n}.$$

De même, $\sum_{k=0}^{n-1}\frac{\omega_{k}}{X-j^2\omega_{k}}=\frac{nj^{2n-2}}{X^n-j^{2n}}$, puis 

$$F=\frac{n}{j-1}(\frac{j^n}{X^n-j^n}-\frac{j^{2n-2}}{X^n-j^{2n}}).$$

Si $n\in3\Zz$, posons $n=3p$, $p\in\Zz$. Dans ce cas,

$$F=\frac{3p}{j-1}(\frac{1}{X^{3p}-1}-\frac{j}{X^{3p}-1})=\frac{3p}{1-X^{3p}}.$$
 
Si $n\in3\Zz+1$, posons $n=3p+1$, $p\in\Zz$. Dans ce cas,

$$F=\frac{3p+1}{j-1}(\frac{j}{X^{3p+1}-j}-\frac{1}{X^{3p+1}-j^2})=\frac{(3p+1)(X^{3p+1}+1)}{X^{6p+2}+X^{3p+1}+1}.$$
 
Si $n\in3\Zz+2$, posons $n=3p+2$, $p\in\Zz$. Dans ce cas,

$$F=\frac{3p+2}{j-1}(\frac{j^2}{X^{3p+2}-j^2}-\frac{j^2}{X^{3p+2}-j})=\frac{3p+2}{X^{6p+4}+X^{3p+2}+1}.$$

\fincorrection
\correction{005338}
Soient $P$ et $Q$ deux polynômes non nuls et premiers entre eux, puis soit $F=\frac{P}{Q}$. Si $F$ est paire, alors $\frac{P(-X)}{Q(-X)}=\frac{P(X)}{Q(X)}$, ou encore $P(-X)Q(X)=P(X)Q(-X)$ $(*)$.

Par suite, $P(X)$ divise $P(X)Q(-X)=Q(X)P(-X)$ et $P(X)$ est premier à $Q(X)$. D'après le théorème de \textsc{Gauss}, $P(X)$ divise $P(-X)$. Donc, il existe $\lambda\in\Cc^*$ tel que $P(-X)=\lambda P(X)$ (car $\mbox{deg}(P(-X))=\mbox{deg}(P)$). L'analyse des coefficients dominants des deux membres fournit $\lambda=(-1)^n$ où $n=\mbox{deg}P$. Ceci s'écrit $P(-X)=(-1)^nP(X)$. En reportant dans $(*)$, on obtient encore $Q(-X)=(-1)^n=Q(X)$. Ainsi, si $F$ est paire, alors $P$ et $Q$ sont ou bien tous deux pairs, ou bien tous deux impairs. Ce dernier cas est exclu, car alors $P$ et $Q$ admettraient tous deux $0$ pour racine contredisant le fait qu'ils sont premiers entre eux. Finalement, si $F$ est paire, alors $P$ et $Q$ sont pairs. La réciproque est claire.

$$F\;\mbox{paire}\Leftrightarrow(P\;\mbox{et}\;Q\;\mbox{sont pairs.})$$

Je vous laisse établir que

$$F\;\mbox{impaire}\Leftrightarrow(P\;\mbox{est impair et}\;Q\;\mbox{est pair})\;\mbox{ou}(P\;\mbox{est pair et}\;Q\;\mbox{est impair.})$$

\fincorrection
\correction{005339}
C'est du cours (unicité de la décomposition en éléments simples).
\fincorrection
\correction{005340}
Soit $n\in\Nn^*$. $\frac{1}{X^2+1}=\frac{1}{2i}(\frac{1}{X-i}-\frac{1}{X+i})$. Donc,

\begin{align*}\ensuremath
\left(\frac{1}{X^2+1}\right)^{(n)}&=\frac{1}{2i}(\left(\frac{1}{X-i}\right)^{(n)}-\left(\frac{1}{X+i}\right)^{(n)})=\frac{1}{2i}(\frac{(-1)(-2)...(-n)}{(X-i)^{n+1}}-\frac{(-1)(-2)...(-n)}{(X+i)^{n+1}})\\
 &=(-1)^n.n!\mbox{Im}(\frac{1}{(X-i)^{n+1}})=(-1)^n.n!\mbox{Im}(\frac{(X+i)^{n+1}}{(X^2+1)^{n+1}})
=\frac{(-1)^n.n!\sum_{}^{}C_{n+1}^{2k+1}(-1)^kX^{2n-2k}}{(X^2+1)^{n+1}}
\end{align*}
\fincorrection
\correction{005341}
$P'=a\sum_{k=1}^{n}\prod_{j\neq k}(X-x_j)=\sum_{k=1}^{n}\frac{P}{X-x_k}$, et donc

$$\frac{P'}{P}=\sum_{k=1}^{n}\frac{1}{X-x_k}.$$

Regroupons maintenant les pôles identiques, ou encore posons $P=a(X-z_1)^{\alpha_1}...(X-z_k)^{\alpha_k}$ où cette fois-ci les $z_j$ sont deux à deux distincts. La formule ci-dessus s'écrit alors

$$\frac{P'}{P}=\sum_{j=1}^{k}\frac{\alpha_j}{X-z_j}\quad(*).$$

Déterminons les polynômes divisibles par leur dérivée. Soit $P$ un tel polynôme. Nécessairement $\mbox{deg}P\ge1$ puis, il existe deux complexes $a$ et $b$, $a\neq0$ tel que $P=(aX+b)P'$ ou encore $\frac{P'}{P}=\frac{1}{aX+b}$. $(*)$ montre que $P$ a une et une seule racine. Par suite, $P$ est de la forme $\lambda(X-a)^n$, $\lambda\neq0$, $n\geq1$ et $a$ quelconque.

Réciproquement, on a dans ce cas $P=\frac{1}{n}(X-a)n(X-a)^{n-1}=(\frac{1}{n}X-\frac{a}{n})P'$ et $P'$ divise effectivement $P$.

Les polynômes divisibles par leur dérivée sont les polynômes de la forme $\lambda(X-a)^n$, $\lambda\in\Cc^*$, $n\in\Nn^*$ et $a\in\Cc$.
\fincorrection
\correction{006964}
\'Ecrivons $F(X)=\frac{P(X)}{Q(X)}$ avec $P$ et $Q$ deux polynômes premiers entre eux, avec $Q$ unitaire.
La condition $\big(F(X)\big)^2=(X^2+1)^3$ devient $P^2=(X^2+1)^3Q^2$. 
Ainsi $Q^2$ divise $P^2$. D'où $Q^2=1$, puisque $P^2$ et $Q^2$ sont premiers entre eux. 
Donc $Q=1$ (ou $-1$). Ainsi $F=P$ est un polynôme et $P^2=(X^2+1)^3$. 

En particulier $P^2$ est de degré $6$, donc $P$ doit être de degré 3. 
\'Ecrivons $P=aX^3+bX^2+cX+d$,
on développe l'identité $P^2=(X^2+1)^3$ :

{\small
$$\begin{array}{c}
X^6+3X^4+3X^2+1 = \\
a^2X^6 + 2abX^5 + (2ac+b^2)X^4 + (2ad+2bc)X^3 + (2bd+c^2)X^2 + 2cdX+d^2
 \end{array}
$$
}

On identifie les coefficients :
pour le coefficient de $X^6$, on a $a=\pm1$,
puis pour le coefficient de $X^5$, on a $b =0$ ;
pour le coefficient de $1$, on a $d=\pm 1$, 
puis pour le coefficient de $X$, on a $c=0$.
Mais alors le coefficient de $X^3$ doit vérifier $2ad+2bc=0$, ce qui est faux.

Ainsi aucun polynôme ne vérifie l'équation $P^2=(X^2+1)^3$, et 
par le raisonnement du début, aucune fraction non plus.
\fincorrection
\correction{006965}\ 
\begin{enumerate}
\item Posons $G=\frac{A}{B}$ et $F = \frac{P}{Q}$ (avec $A,B,P,Q$ des polynômes).
On réécrit l'identité $G(F(X))=X$ sous la forme $A(F(X))=XB(F(X))$. 
Posons $n=\mathrm{max}(\mathrm{deg}A,\mathrm{deg}B)$.
Alors $n\ge 1$ car sinon, 
$A$ et $B$ seraient constants et $G(\frac{P}{Q})=X$ aussi.

On a donc $A=\sum_{k=0}^na_kX^k$ et $B=\sum_{k=0}^nb_kX^k$, où $(a_n,b_n)\neq(0,0)$, et l'identité devient
$$\sum_{k=0}^na_k\left(\frac{P}{Q}\right)^k=X\sum_{k=0}^nb_k\left(\frac{P}{Q}\right)^k$$
En multipliant par $Q^n$, cela donne
$$\sum_{k=0}^n a_kP^kQ^{n-k}= \sum_{k=0}^n b_k X P^kQ^{n-k}.$$
Donc 
$$(a_0-b_0X)Q^n \quad + \quad (\cdots + (a_k-b_kX) P^kQ^{n-k} + \cdots)\quad + \quad (a_n-b_nX)P^n = 0$$
où les termes dans la parenthèse centrale sont tous divisibles par $P$ et par $Q$. 
Comme $Q$ divise aussi le premier terme, alors $Q$ divise $(a_n-b_nX)P^n$.
D'après le lemme de Gauss, puisque $P$ et $Q$ sont premiers entre eux, alors $Q$ divise $(a_n-b_nX)$. 
De même, $P$ divise tous les termes de la parenthèse centrale et le dernier, donc $P$ divise aussi $(a_0-b_0X)Q^n$, 
donc $P$ divise $(a_0-b_0X)$.

\item Supposons de plus qu'on a écrit $G=\frac{A}{B}$ sous forme irréductible, 
c'est-à-dire avec $\pgcd(A,B)=1$. 
Vu que $a_n$ et $b_n$ ne sont pas tous les deux nuls, alors $a_n-b_nX$ n'est pas le polynôme nul.
Comme $Q$ divise $a_n-b_nX$ alors nécessairement $Q$ est de degré au plus $1$ ; on écrit $Q(X)=cX+d$. 
Par ailleurs, $a_0-b_0X$ n'est pas non plus le polynôme nul, car sinon on aurait $a_0=b_0=0$ et 
donc $A$ et $B$ seraient tous les deux sans terme constant, donc divisibles par $X$ 
(ce qui est impossible puisqu'ils sont premiers entre eux).  
Donc $P$ est aussi de degré au plus $1$ et on écrit $P(X)=aX+b$. 
Conclusion : $F(X)=\frac{aX+b}{cX+d}$.
Notez que $a$ et $b$ ne sont pas tous les deux nuls en même temps (de même pour $b$ et $d$).

\item Si $Y = \frac{aX+b}{cX+d}$ avec $(a,b) \neq (0,0)$,
alors $X = -\frac{dY-b}{cY-a}$. 
Autrement dit si on note
$\phi(X)= \frac{aX+b}{cX+d}$, alors sa bijection réciproque est 
$\phi^{-1}(Y) = -\frac{dY-b}{cY-a}$.

Nous avons prouvé que $G\left( \frac{aX+b}{cX+d}\right) =X$.
Cette identité s'écrit $G\big( \phi(X) \big)=X$.
Appliquée en $X = \phi^{-1}(Y)$ elle devient
$G\big( \phi( \phi^{-1}(Y) ) \big)=\phi^{-1}(Y)$, c'est-à-dire
$G(Y) = \phi^{-1}(Y)$.
Ainsi $G(Y) = -\frac{dY-b}{cY-a}$.
\end{enumerate}
\fincorrection
\correction{006966}
\begin{enumerate}
\item 
  \begin{enumerate}
    \item Puisque $P(X)=c(X-a_1)\cdots(X-a_n)$ : 
\begin{eqnarray*}
P'(X)=&c(X-a_2)\cdots(X-a_n)+c(X-a_1)(X-a_3)\cdots(X-a_n)\\
 &+\cdots+c(X-a_1)\cdots(X-a_{k-1})(X-a_{k+1})\cdots(X-a_n)\\
 &\ \ \ \ \ \ +\cdots+c(X-a_1)\cdots(X-a_{n-1})
\end{eqnarray*}
La dérivée est donc la somme des termes de la forme : $\frac{c(X-a_1)\cdots(X-a_n)}{X-a_k} = \frac{P(X)}{X-a_k}$.

Ainsi 
$$P'(X) = \frac{P(X)}{X-a_1}+ \cdots + \frac{P(X)}{X-a_k}+ \cdots + \frac{P(X)}{X-a_n}.$$
Donc :
$$\frac{P'}{P}=\sum_{k=1}^n\frac{1}{X-a_k}$$

  \item Puisque $\sum_{k=1}^n\frac{1}{(X-a_k)^2}$ est la dérivée de $-\sum_{k=1}^n\frac{1}{X-a_k}$, on obtient
  par dérivation de $-\frac{P'}{P}$ :
$$\frac{P'^2-PP''}{P^2}=\sum_{k=1}^n\frac{1}{(X-a_k)^2}$$

  \item On a remarqué que la dérivée de $P'$ est la somme de facteurs $c(X-a_1)\cdots(X-a_n)$
avec un des facteurs en moins, donc de la forme $\frac{c(X-a_1)\cdots(X-a_n)}{X-a_k} = \frac{P}{X-a_k}$. 
De même $P''$ est la somme de facteurs $c(X-a_1)\cdots(X-a_n)$
avec deux facteurs en moins, c'est-à-dire de la forme $\frac{c(X-a_1)\cdots(X-a_n)}{(X-a_k)(X-a_\ell)} = \frac{P}{(X-a_k)(X-a_\ell)}$ :
$$P'' = \sum_{\substack{1\le k,\ell\le n \\ k\not= \ell}}\frac{P}{(X-a_k)(X-a_\ell)} \quad \text{ donc } \quad
\frac{P''}{P} = \sum_{\substack{1\le k,\ell\le n \\ k\not= \ell}}\frac{1}{(X-a_k)(X-a_\ell)}$$

% \begin{eqnarray*}
% \sum_{k\not=l}\frac{1}{(X-a_k)(X-a_l)}&=&\sum_{k=1}^n\frac{1}{X-a_k}\left(\sum_{l\not= k}\frac{1}{X-a_l}\right)\\
%  &=&\sum_{k=1}^n\frac{1}{X-a_k}\left(\sum_{l=1}^n\frac{1}{X-a_l}-\frac{1}{X-a_k}\right)\\
%  &=&\sum_{k=1}^n\frac{1}{X-a_k}\frac{P'}{P}-\sum_{k=1}^n\frac{1}{(X-a_k)^2}\\
%  &=&\frac{P'^2}{P^2}-\frac{P'^2-PP''}{P^2}=\frac{P''}{P}
% \end{eqnarray*}
  \end{enumerate}

\item On applique l'identité  $\frac{P'(X)}{P(X)}=\sum_{k=1}^n\frac{1}{X-a_k}$ en $z$ avec les hypothèses 
$P(z)\not=0$ et $P'(z)=0$.
On en déduit $\displaystyle \sum_{k=1}^n\frac{1}{z-a_k} = 0$.
L'expression conjuguée est aussi nulle :
$$\sum_{k=1}^n\frac{1}{\overline{z-a_k}}=\sum_{k=1}^n\frac{z-a_k}{|z-a_k|^2} = 0$$
Posons $\mu_k = \frac{1}{|z-a_k|^2}$.
Alors 
$$\sum_{k=1}^n \mu_k(z-a_k) = 0 \quad \text{ donc } \left(\sum_{k=1}^n \mu_k \right) z = \sum_{k=1}^n \mu_k a_k$$
Posons $\lambda_k = \mu_k / \left(\sum_{k=1}^n \mu_k \right)$,
alors : 
  \begin{itemize}
    \item Les $\lambda_k$ sont des réels positifs.
    \item $\sum_{k=1}^n\lambda_k=1$
    \item Et $z=\sum_{k=1}^n\lambda_ka_k$.
  \end{itemize}

En particulier si les $a_k$ sont tous des nombres réels, alors $z$ est aussi un nombre réel.
On vient de prouver que si un polynôme $P$ a toutes ses racines réelles, 
alors $P'$ a aussi toutes ses racines réelles. On a même plus : si on ordonne les racines réelles 
de $P$ en $a_1 \le a_2 \le \cdots \le a_n$ alors une racine $z$ de $P'$ est réelle et vérifie $a_1 \le z \le a_n$.

Plus généralement, l'interprétation géométrique de ce que l'on vient de prouver s'appelle le théorème de Gauss-Lucas :
<<Les racines de $P'$ sont dans l'enveloppe convexe des racines (réelles ou complexes) de $P$.>>
\end{enumerate}
\fincorrection
\correction{006967}\
\begin{enumerate}
\item $F=\frac{X}{X^2-4}$. 

Commençons par factoriser le dénominateur : $X^2-4=(X-2)(X+2)$, 
d'où une décomposition en éléments simples du type
$F=\frac{a}{X-2}+\frac{b}{X+2}$. 
En réduisant au même dénominateur, il vient
$\frac{X}{X^2-4}=\frac{(a+b)X+2(a-b)}{X^2-4}$ et en identifiant les coefficients, on obtient le système
$\left\{\begin{array}{l}a+b=1\\2(a-b)=0\end{array}\right.$. Ainsi $a=b=\frac{1}{2}$ et 
$$\frac{X}{X^2-4} = \frac{\frac12}{X-2} + \frac{\frac12}{X+2}$$

\item $G=\frac{X^3-3X^2+X-4}{X-1}$. 

Lorsque le degré du numérateur (ici $3$) est supérieur 
ou égal au degré du dénominateur (ici $1$), il
faut effectuer la division euclidienne du numérateur par le dénominateur pour 
faire apparaître la partie polynomiale (ou partie entière).
Ici la division euclidienne s'écrit $X^3-3X^2+X-4=(X-1)(X^2-2X-1)-5$. 
Ainsi en divisant les deux membres par $X-1$ on obtient 
$$\frac{X^3-3X^2+X-4}{X-1} = X^2-2X-1-\frac{5}{X-1}$$
La fraction est alors déjà décomposée en éléments simples.

\item $H=\frac{2X^3+X^2-X+1}{X^2-2X+1}$. 

Commençons par faire la division euclidienne du numérateur par le dénominateur :
$2X^3+X^2-X+1=(X^2-2X+1)(2X+5)+7X-4$, ce qui donne
$H=2X+5+\frac{7X-4}{X^2-2X+1}$. 
Il reste à décomposer en éléments simples la fraction rationnelle $H_1=\frac{7X-4}{X^2-2X+1}$. 
Puisque le dénominateur se factorise en $(X-1)^2$, elle sera de la forme 
$H_1=\frac{a}{(X-1)^2}+\frac{b}{X-1}$.
En réduisant au même dénominateur, il vient
$\frac{7X-4}{X^2-2X+1}=\frac{bX+a-b}{X^2-2X+1}$ et en identifiant les coefficients, on obtient 
$b=7$ et $a=3$. Finalement,
$$\frac{2X^3+X^2-X+1}{X^2-2X+1} =
2X+5+\frac{3}{(X-1)^2}+\frac{7}{X-1}$$

\item $K=\frac{X+1}{X^4+1}$. 

Ici, il n'y a pas de partie polynomiale puisque le degré du numérateur 
est strictement inférieur au degré du dénominateur. 
Le dénominateur admet quatre racines complexes 
$e^{\frac{i\pi}{4}}$, $e^{\frac{3i\pi}{4}}$, $e^{\frac{5i\pi}{4}}=e^{-\frac{3i\pi}{4}}$ et 
$e^{\frac{7i\pi}{4}}=e^{-\frac{i\pi}{4}}$. 
En regroupant les racines complexes conjuguées, on obtient sa factorisation sur $\Rr$:
\begin{eqnarray*}
X^4+1 &=& \big( (X-e^{\frac{i\pi}{4}})(X-e^{-\frac{i\pi}{4}}) \big)\big( (X-e^{\frac{3i\pi}{4}})(X-e^{-\frac{3i\pi}{4}}) \big) \\
      &=&\big(X^2-2\cos\tfrac{\pi}{4}+1\big)\big(X^2-2\cos\tfrac{3\pi}{4}+1\big)\\
      &=&(X^2-\sqrt{2}X+1)(X^2+\sqrt{2}X+1)
\end{eqnarray*}
Puisque les deux facteurs $(X^2-\sqrt{2}X+1)$ et $(X^2+\sqrt{2}X+1)$ sont irréductibles sur $\R$, la décomposition en éléments simples de $K$ est de la forme
$K=\frac{aX+b}{X^2-\sqrt{2}X+1}+\frac{cX+d}{X^2+\sqrt{2}X+1}$.

En réduisant au même dénominateur et en identifiant les coefficients avec ceux de $K=\frac{X+1}{X^4+1}$, 
on obtient le système
$$\left\{\begin{array}{l}
a+c=0\\
\sqrt{2}a+b-\sqrt{2}c+d=0\\
a+\sqrt{2}b+c-\sqrt{2}d=1\\
b+d=1
\end{array}\right.$$
Système que l'on résout en $a=\frac{-\sqrt{2}}{4}$, $c=\frac{\sqrt{2}}{4}$, $b=\frac{2+\sqrt{2}}{4}$ et $d=\frac{2-\sqrt{2}}{4}$. Ainsi
$$\frac{X+1}{X^4+1} = \frac{-\frac{\sqrt{2}}{4}X+\frac{2+\sqrt{2}}{4}}{X^2-\sqrt{2}X+1} + \frac{\frac{\sqrt{2}}{4}X+\frac{2-\sqrt{2}}{4}}{X^2+\sqrt{2}X+1}$$ 
\end{enumerate}
\fincorrection
\correction{006968}\ 
\begin{enumerate}

\item $F=\frac{X^5+X^4+1}{X^3-X}$.

Pour obtenir la partie polynomiale, on fait une division euclidienne : 
$X^5+X^4+1=(X^3-X)(X^2+X+1)+ X^2+X+1$. Ce qui donne $F=X^2+X+1+F_1$, 
où $F_1=\frac{X^2+X+1}{X^3-X}$. Puisque $X^3-X=X(X-1)(X+1)$, 
la décomposition en éléments simples est de la forme 
$$F_1 = \frac{X^2+X+1}{X(X-1)(X+1)}=\frac{a}{X}+\frac{b}{X-1}+\frac{c}{X+1}$$

Pour obtenir $a$ :
\begin{itemize}
  \item on multiplie l'égalité par $X$ : $\frac{X(X^2+X+1)}{X(X-1)(X+1)}=X \left(\frac{a}{X}+\frac{b}{X-1}+\frac{c}{X+1}\right)$,
  \item on simplifie $\frac{X^2+X+1}{(X-1)(X+1)}= a+\frac{bX}{X-1}+\frac{cX}{X+1}$,
  \item on remplace $X$ par $0$ et on obtient $-1= a+0+0$, donc $a=-1$.
\end{itemize}

 De même, en multipliant par $X-1$ et en remplaçant $X$ par $1$, il vient $b=\frac{3}{2}$.
 Puis en multipliant par $X+1$ et en remplaçant $X$ par $-1$, on trouve $c=\frac{1}{2}$.

D'où
$$\frac{X^5+X^4+1}{X^3-X} = X^2+X+1-\frac{1}{X}+\frac{\tfrac12}{X+1}+\frac{\tfrac32}{X-1}$$

\item $G=\frac{X^3+X+1}{(X-1)^3(X+1)}$. 

La partie polynomiale est nulle. La décomposition en éléments simples est de la forme
$G=\frac{a}{(X-1)^3}+\frac{b}{(X-1)^2}+\frac{c}{X-1}+\frac{d}{X+1}$.

\begin{itemize}
  \item En multipliant les deux membres de l'égalité par $(X-1)^3$, 
en simplifiant puis en remplaçant $X$ par $1$, 
on obtient $a=\frac32$. 

  \item De même, en multipliant par $X+1$, 
en simplifiant puis en remplaçant $X$ par $-1$, 
on obtient $d=\frac18$.

  \item En multipliant par $X$ et en regardant la limite 
  quand $X\to +\infty$, on obtient $1=c+d$. Donc $c=\frac78$.

  \item En remplaçant $X$ par $0$, il vient $-1=-a+b-c+d$.
  Donc $b = \frac54$.
\end{itemize}

Ainsi :
$$G = \frac{X^3+X+1}{(X-1)^3(X+1)} = \frac{\tfrac32}{(X-1)^3} + \frac{\tfrac54}{(X-1)^2} 
+ \frac{\tfrac78}{X-1}  + \frac{\tfrac18}{X+1}$$

\item $H=\frac{X}{(X^2+1)(X^2+4)}$.

Puisque $X^2+1$ et $X^2+4$ sont irréductibles sur $\Rr$, la décomposition en éléments simples sera de la forme 
$$\frac{X}{(X^2+1)(X^2+4)} = \frac{aX+b}{X^2+1} + \frac{cX+d}{X^2+4}$$

  \begin{itemize}
    \item En remplaçant $X$ par $0$, on obtient $0=b+\frac{1}{4}d$.
    \item En multipliant les deux membres par $X$, on obtient 
    $\frac{X^2}{(X^2+1)(X^2+4)} = \frac{aX^2+bX}{X^2+1} + \frac{cX^2+dX}{X^2+4}$.
    En calculant la limite quand $X\to +\infty$, on a $0=a+c$.
    \item Enfin, en évaluant les fractions en $X=1$ et $X=-1$, 
    on obtient $\frac{1}{10}=\frac{a+b}{2}+\frac{c+d}{5}$ et 
    $\frac{-1}{10}=\frac{-a+b}{2}+\frac{-c+d}{5}$.
  \end{itemize}

La résolution du système donne $b=d=0$, $a=\frac{1}{3}$, $c=-\frac{1}{3}$ et donc 
$$\frac{X}{(X^2+1)(X^2+4)} = \frac{\frac{1}{3}X}{X^2+1} - \frac{\frac{1}{3}X}{X^2+4}$$

\item $K=\frac{2X^4+X^3+3X^2-6X+1}{2X^3-X^2}$.

Pour la partie polynomiale, on fait la division euclidienne:
$$2X^4+X^3+3X^2-6X+1=(2X^3-X^2)(X+1)+(4X^2-6X+1)$$ 
ce qui donne 
$K=X+1+K_1$ où $K_1=\frac{4X^2-6X+1}{2X^3-X^2}$.
Pour trouver la décomposition en éléments simples de $K_1$, on factorise son numérateur: 
$2X^3-X^2=2X^2(X-\frac{1}{2})$, ce qui donne une décomposition de la forme
$K_1=\frac{a}{X^2}+\frac{b}{X}+\frac{c}{X-\frac{1}{2}}$.

On obtient alors $a$ en multipliant les deux membres de l'égalité par
$X^2$ puis en remplaçant $X$ par 0: $a=-1$. On obtient de
m\^eme $c$ en multipliant par $X-\frac{1}{2}$ et en remplaçant $X$
par $\frac{1}{2}$: $c=-2$. Enfin on trouve $b$ en identifiant pour une
valeur particuli\`ere non encore utilisée, par exemple $X=1$, ou mieux en
multipliant les deux membres par $X$ et en passant \`a la limite pour
$X\to+\infty$: $b=4$. Finalement:
$$\frac{2X^4+X^3+3X^2-6X+1}{2X^3-X^2}=X+1-\frac{1}{X^2}+\frac{4}{X}-\frac{2}{X-\frac{1}{2}}$$
\end{enumerate}

\fincorrection
\correction{006969}\ 
\begin{enumerate}
\item $F=\frac{4X^6-2X^5+11X^4-X^3+11X^2+2X+3}{X(X^2+1)^3}$.

  \begin{enumerate}
  \item 
La décomposition en éléments simples de $F$ est de la forme
$F=\frac{a}{X}+\frac{bX+c}{(X^2+1)^3}+\frac{dX+e}{(X^2+1)^2}+\frac{fX+g}{X^2+1}$.
Il est difficile d'obtenir les coefficients par substitution.
% $a,b,c$, (pour ces derniers : multiplication des deux membres par $X^2+1$, simplification, puis remplacement de $X$ par $i$ ou $-i$, avec séparation des parties
% réelle et imaginaire), mais c'est insuffisant pour conclure: il faut
% encore soustraire $\frac{bX+c}{(X^2+1)^3}$, simplifier par $X^2+1$, calculer $d$
% et $e$,\ldots

  \item
On va ici se contenter de trouver $a$ :
on multiplie $F$ par $X$, puis on remplace $X$ par $0$, on obtient $a=3$.

  \item On fait la soustraction $F_1=F-\frac{a}{X}$. 
  On sait que la fraction $F_1$ \emph{doit} se simplifier par $X$. 
  On trouve $F_1=\frac{X^5-2X^4+2X^3-X^2+2X+2}{(X^2+1)^3}$.
  
  \item La fin de la décomposition se fait par divisions euclidiennes successives.
  Tout d'abord la division du numérateur $X^5-2X^4+2X^3-X^2+2X+2$
par $X^2+1$:
$$X^5-2X^4+2X^3-X^2+2X+2=(X^2+1)(X^3-2X^2+X+1)+X+1$$
puis on recommence en divisant le quotient obtenu par $X^2+1$, pour obtenir 
$$X^5-2X^4+2X^3-X^2+2X+2=(X^2+1)\big((X^2+1)(X-2)+3\big)+X+1$$
On divise cette identité par $(X^2+1)^3$ :

$F_1 = \frac{(X^2+1)\big((X^2+1)(X-2)+3\big)+X+1}{(X^2+1)^3} 
= \frac{X+1}{(X^2+1)^3}+\frac{3}{(X^2+1)^2}+\frac{X-2}{X^2+1}$

Ainsi 
$$F=
\frac{3}{X}+\frac{X+1}{(X^2+1)^3}+\frac{3}{(X^2+1)^2}+\frac{X-2}{X^2+1}$$
  \end{enumerate}
  
\smallskip
Remarque~: cette méthode des divisions successives est tr\`es pratique quand
la fraction \`a décomposer a un dénominateur \emph{simple}, c'est \`a dire
comportant un dénominateur du type $Q^n$ o\`u $Q$ est du premier degré, ou
du second degré sans racine réelle. 

\item $G=\frac{4X^4-10X^3+8X^2-4X+1}{X^3(X-1)^2}$.

La décomposition en éléments simples de $G$ est de la forme
$\frac{a}{X^3}+\frac{b}{X^2}+\frac{c}{X}+\frac{d}{(X-1)^2}+\frac{e}{X-1}$.
% On pourrait obtenir facilement $a$ et $d$ par multiplication par $X^3$ et par $(X-1)^2$,
% mais qu'il resterait encore trois coefficients à déterminer. Il y a ici une
La méthode la plus efficace pour déterminer les coefficients est d'effectuer une division suivant les puissances
croissantes, ici à l'ordre $2$ (de sorte que le reste soit divisible par $X^3$ comme le dénominateur).
On calcule la division suivant les puissances
croissantes, à l'ordre $2$ du numérateur $1-4X+8X ^2-10X^3+4X^4$ par $(X-1)^2$, ou plutôt par $1-2X+X^2$ :
$$1-4X+8X^2-10X^3+4X^4=(1-2X+X^2)(1-2X+3X^2)+(-2X^3+X^4)$$
Remarquer que le reste $-2X^3+X^4$ est divisible par $X^3$.

En divisant les deux membres de cette identité par $X^3(X-1)^2$, 
on obtient $a$, $b$ et $c$ d'un seul coup:
$$
\begin{array}{rcl}
G & = & \frac{4X^4-10X^3+8X^2-4X+1}{X^3(X-1)^2}   \\
  & = & \frac{(X-1)^2(1-2X+3X^2)+(-2X^3+X^4)}{X^3(X-1)^2} \\
  & = & \frac{1}{X^3}-\frac{2}{X^2}+\frac{3}{X}+\frac{X-2}{(X-1)^2}
\end{array}
$$

Il reste à trouver $d$ et $e$: par exemple en faisant la division euclidienne de $X-2$ par $X-1$: $X-2=(X-1)-1$.
$$G=\frac{1}{X^3}-\frac{2}{X^2}+\frac{3}{X}-\frac{1}{(X-1)^2}+\frac{1}{X-1}$$

\item $H=\frac{X^4+2X^2+1}{X^5-X^3} = \frac{X^4+2X^2+1}{X^3(X-1)(X+1)}$.

La décomposition sera de la forme 
$H=\frac{a}{X^3}+\frac{b}{X^2}+\frac{c}{X}+\frac{d}{X-1}+\frac{e}{X+1}$. 
Pour obtenir $a, b, c$, on fait la division du numérateur par $(X-1)(X+1) = X^2-1$ 
selon les puissances croissantes, à l'ordre $2$ (de sorte que le reste soit divisible par $X^3$ qui est la puissance de $X$ au dénominateur de $H$,
en fait on l'obtient à l'ordre $3$) :
$$1+2X^2+X^4=(-1+X^2)(-1-3X^2)+4X^4$$
ce qui donne directement
$$
\begin{array}{rcl}
H 
&=& \frac{X^4+2X^2+1}{X^3(X-1)(X+1)} \\
&=& \frac{(X^2-1)(-1-3X^2)+4X^4}{X^3(X^2-1)} \\
&=& -\frac{1}{X^3}-\frac{3}{X}+\frac{4X}{X^2-1}  
\end{array}
$$
Il reste à décomposer $\frac{4X}{X^2-1}=\frac{d}{X-1}+\frac{e}{X+1}$. On trouve $d=e=2$, d'où
$$H = \frac{X^4+2X^2+1}{X^5-X^3} =-\frac{1}{X^3}-\frac{3}{X}+\frac{2}{X-1}+\frac{2}{X+1}$$ 

\medskip
Remarque : la méthode de division selon les puissances croissantes est 
efficace pour un exposant assez grand (en gros
à partir de $3$) dans une fraction du type $\frac{P(X)}{X^n Q(X)}$. 
Elle peut être utilisée pour une fraction du type 
$\frac{P(X)}{(X-a)^nQ(X)}$, mais il faut commencer par le changement de variable
$Y=X-a$ avant de faire la division, puis bien entendu revenir à la
variable $X$.


\item $K=\frac{X^5+X^4+1}{X(X-1)^4}$.

Puisque le degré du numérateur $N$ est supérieur ou égal à celui du dénominateur $D$, 
il y a une partie polynomiale.
$N$ et $D$ étant de même degré, avec le même coefficient dominant, 
la partie polynomiale vaut $1$ et $K$ se décompose sous la forme 
$K = 1+\frac{a}{X} + \frac{b}{(X-1)^4} + \frac{c}{(X-1)^3} + \frac{d}{(X-1)^2} + \frac{e}{X-1}$. 
Le coefficient $a$ s'obtient facilement en multipliant $K$ par $X$ puis en remplaçant $X$ par $0$ : $a=1$. 

Soit $K_1=K-1-\frac{1}{X}=\frac{4X^3-2X^2-2X+3}{(X-1)^4}$. Le changement d'indéterminée $X=Y+1$ donne 
$K_1=\frac{4(Y+1)^3-2(Y+1)^2-2(Y+1)+3}{Y^4} $. En développant, on obtient directement
$$K_1=\frac{4Y^3+10Y^2+6Y+3}{Y^4}= \frac{3}{Y^4} + \frac{6}{Y^3} + 
\frac{10}{Y^2} + \frac{4}{Y}$$
et donc (avec $Y=X-1$) :
$K_1=\frac{4X^3-2X^2-2X+3}{(X-1)^4}=\frac{3}{(X-1)^4} + \frac{6}{(X-1)^3} + 
\frac{10}{(X-1)^2} + \frac{4}{X-1}$ .
Finalement,
$$K = \frac{X^5+X^4+1}{X(X-1)^4} = 1 + \frac{1}{X} + \frac{3}{(X-1)^4} + \frac{6}{(X-1)^3} + 
\frac{10}{(X-1)^2} + \frac{4}{X-1}$$
\end{enumerate}
\fincorrection
\correction{006970}\
\begin{enumerate}
\item 
$$\begin{array}{rcl}
\frac{(3-2 i )X-5+3 i}{X^2+ i  X+2} &=&\frac{2+ i}{X- i}+\frac{1-3 i}{X+2 i}\\
\frac{X+ i }{X^2+ i} &=&\frac{\frac{2-\sqrt2}{4}+\frac{\sqrt2}{4} i}{X-\frac{\sqrt2-\sqrt2 i }{2}} +
\frac{\frac{2+\sqrt2}{4}-\frac{\sqrt2}{4} i }{X+\frac{\sqrt2-\sqrt2 i}{2}}\\
\frac{X}{(X+ i)^2} &=&\frac{1}{X+ i}+\frac{-i}{(X+ i)^2}
\end{array}$$

\item 
$$\begin{array}{rcl}
\frac{X^5+X+1}{X^4-1} 
& =& X + \frac{\frac34}{X-1} + \frac{\frac14}{X+1} - \frac{X+\frac{1}{2}}{X^2+1} \\
 &=& X + \frac{\frac34}{X-1} + \frac{\frac14}{X+1} + 
\frac{-\frac{1}{2}+\frac{1}{4} i }{X- i } + 
\frac{-\frac{1}{2}-\frac{1}{4} i }{X+ i }\\
\frac{X^2-3}{(X^2+1)(X^2+4)} 
&=& -\frac{\frac43}{X^2+1} + \frac{\frac73}{X^2+4} \\
 & =& \frac{\frac{2}{3} i }{X- i } + \frac{-\frac{2}{3} i }{X+ i } + 
\frac{-\frac{7}{12} i }{X-2 i } + \frac{\frac{7}{12} i }{X+2 i }\\
\frac{X^2+1}{X^4+1} 
&=&\frac{\frac12}{X^2-\sqrt{2}X+1}+\frac{\frac12}{X^2+\sqrt{2}X+1}\\
& =&\frac{-\frac{\sqrt{2}}{4} i }{X-\frac{\sqrt{2}}{2}-\frac{\sqrt{2}}{2} i } + 
\frac{ \frac{\sqrt{2}}{4} i }{X-\frac{\sqrt{2}}{2}+\frac{\sqrt{2}}{2} i } \\
 & &\ \ \ \ \ \  \ \ \ + 
\frac{ \frac{\sqrt{2}}{4} i }{X+\frac{\sqrt{2}}{2}+\frac{\sqrt{2}}{2} i } + 
\frac{-\frac{\sqrt{2}}{4} i }{X+\frac{\sqrt{2}}{2}-\frac{\sqrt{2}}{2} i }
\end{array}$$

\end{enumerate}
\fincorrection
\correction{006971}
La décomposition en élément simple s'écrit :
$$\frac{1}{X(X-1)(X-2)^2}=\frac{-\frac14}{X}+\frac{1}{X-1}+\frac{\frac12}{(X-2)^2}+\frac{-\frac34}{X-2}.$$ 
En multipliant cette identité par le dénominateur $X(X-1)(X-2)^2$, il vient :
$$1 =-\tfrac{1}{4}Q_0+Q_1+\left(\tfrac{1}{2}-\tfrac{3}{4}(X-2)\right)Q_2.$$
Ainsi $A_0=-\frac{1}{4}$, $A_1=1$ et $A_1=(2-\frac{3}{4}X)$ conviennent. 
On a obtenu une relation de Bézout entre $Q_1$, $Q_2$ et $Q_3$ qui prouve que 
ces trois polynômes sont premiers dans leur ensemble : $\pgcd(Q_1,Q_2,Q_3)=1$.
\fincorrection
\correction{006972}

\begin{enumerate}
  \item 
  \begin{enumerate}
    \item Si on pose $x=\cos \theta$ alors l'égalité $T_n(x)=\cos\big(n \arccos(x)\big)$
    devient $T_n(\cos\theta)=\cos(n\theta)$, car $\arccos (\cos \theta) = \theta$ pour $\theta \in [0,\pi]$.
    
    \item $T_0(x) = 1$, $T_1(x)=x$.
    
    \item En écrivant $(n+2)\theta = (n+1)\theta + \theta$ et $n\theta = (n+1)\theta - \theta$ on obtient :
    $$\begin{array}{c}
    \cos\big( (n+2)\theta \big) = \cos\big((n+1)\theta\big)\cos \theta - \sin\big((n+1)\theta\big)\sin \theta     \\
    \cos\big( n\theta \big) = \cos\big((n+1)\theta\big)\cos \theta + \sin\big((n+1)\theta\big)\sin \theta     \\    
      \end{array}$$
    Lorsque l'on fait la somme de ces deux égalités on obtient :
    $$\cos\big( (n+2)\theta \big) + \cos\big( n\theta \big) = \cos\big((n+1)\theta\big)\cos \theta$$
    
    Avec $x = \cos \theta$ cela donne :
    $$T_{n+2}(x) + T_n(x) = 2x T_{n+1}(x)$$
    
    \item $T_0$ et $T_1$ étant des polynômes alors, par récurrence,  $T_n(x)$ est un polynôme.
    De plus, toujours par la formule de récurrence, il est facile de voir que le degré de $T_n$ est $n$.
  \end{enumerate}

% Puisque $\arccos (\cos\theta)=\theta$, on a bien $T_n(\cos\theta)=\cos(n\theta)$. D'après la formule de Moivre,
% $$\cos(n\theta)=\mathrm{Re}(e^{in\theta)})=\mathrm{Re}((\cos\theta+i\sin\theta)^n)$$
% $$=\mathrm{Re}\left(\sum_{k=0}^nC_n^ki^k(\sin\theta)^k(\cos\theta)^{n-k}\right)$$
% Or $i^k$ est réel si et seulement si $k$ est pair (et dans ce cas, $i^k=(-1)^{k/2}$), donc en faisant le changement d'indice $k=2p$:
% \begin{eqnarray*}
%  \cos(n\theta)&=&\sum_{k\ pair,\ 1\le k\le n}C_n^ki^k(\sin\theta)^k(\cos\theta)^{n-k}\\
%  &=&\sum_{p=0}^{E(n/2)}C_n^{2p}(-1)^p(\sin\theta)^{2p}(\cos\theta)^{n-2p}\\
%  &=&\sum_{p=0}^{E(n/2)}C_n^{2p}(-1)^p(1-(\cos\theta)^2)^{p}(\cos\theta)^{n-2p}
% \end{eqnarray*}
% qui est polynomiale en $\cos\theta$.
% Comme tout $x\in[-1;1]$ peut s'écrire sous la forme $x=\cos\theta$, on a 
% $$T_n(x)=\sum_{p=0}^{E(n/2)}C_n^{2p}(-1)^p(1-x^2)^{p}x^{n-2p}$$
% Cette fonction polynomiale (sur $[-1;1]$) est de degré au plus $n$, et le coefficient du terme de degré $n$ vaut
% $a_n=\sum_{p=0}^{E(n/2)}C_n^{2p}$ qui est non nul (somme de termes strictement positifs).

\item Puisque les racines de $P=\lambda(X-a_1)\cdots(X-a_n)$ sont deux à deux distinctes, 
la décomposition en éléments simples de $\frac{1}{P}$ est de la forme 
$\frac{c_1}{X-a_1}+\cdots+\frac{c_n}{X-a_n}$. 

Expliquons comment calculer le coefficient $c_1$. On multiplie la fraction 
$\frac{1}{P}$ par $(X-a_1)$ ce qui donne  

$
\frac{X-a_1}{P} = c_1 + c_2 \frac{X-a_1}{X-a_2}+\cdots+ c_n \frac{X-a_1}{X-a_n}
\text{ et }
\frac{X-a_1}{P} = \frac{1}{\lambda(X-a_2)\cdots (X-a_n)}
$

On évalue ces égalités en $X=a_1$ ce qui donne
$$c_1 = \frac{1}{\lambda(a_1-a_2)\cdots (a_1-a_n)} = \frac{1}{\lambda\prod_{j\not= 1}(a_1-a_j)}$$

On obtiendrait de même le coefficient $c_k$ en multipliant $\frac{1}{P}$ par $(X-a_k)$,
puis en remplaçant $X$ par $a_k$, ce qui donne: $c_k=\frac{1}{\lambda\prod_{j\not= k}(a_k-a_j)}$.

Or la dérivée de $P$ est 
\begin{eqnarray*}
P'(X)=&\lambda(X-a_2)\cdots(X-a_n)+\lambda(X-a_1)(X-a_3)\cdots(X-a_n)\\
 &+\cdots+\lambda(X-a_1)\cdots(X-a_{k-1})(X-a_{k+1})\cdots(X-a_n)\\
 &\ \ \ \ \ +\cdots+\lambda(X-a_1)\cdots(X-a_{n-1})
\end{eqnarray*}
et donc $P'(a_1)=\lambda\prod_{j\not= 1}(a_1-a_j)$
et plus généralement $P'(a_k)=\lambda\prod_{j\not= k}(a_k-a_j)$. 
On a bien prouvé $c_k = \frac{1}{P'(a_k)}$ et ainsi la décomposition en éléments simple 
de $1/P$ est :
$$\frac{1}{P(X)}=\sum_{k=1}^n\frac{\frac{1}{P'(a_k)}}{X-a_k}$$

\item 
  \begin{enumerate}
    \item Cherchons d'abord les racines de $T_n(x)$.
Soit $x\in[-1,1]$ :
\begin{eqnarray*}
T_n(x)=0
  &\Longleftrightarrow& \cos\big(n \arccos(x)\big) = 0 \\
  &\Longleftrightarrow& n\arccos(x) \equiv \frac{\pi}{2} \pmod{\pi}\\
  &\Longleftrightarrow& \exists k\in\Zz\ \ \arccos(x)= \frac{\pi}{2n}+\frac{k\pi}{n}
\end{eqnarray*}
Comme par définition $\arccos(x) \in[0,\pi]$, les entiers $k$ possibles sont $k=0,\ldots,n-1$. Ainsi
$$\arccos x=\frac{\pi}{2n}+\frac{k\pi}{n} \quad \Longleftrightarrow \quad x=\cos\left(\frac{\pi}{2n}+\frac{k\pi}{n}\right)$$ 
Posons donc $\omega_k=\cos\left(\frac{(2k+1)\pi}{2n}\right)$ pour $k=0,\ldots,n-1$. 
Les $\omega_k$ sont les racines de $T_n$. Finalement $T_n(x)= \lambda \prod_{k=0}^{n-1}(x-\omega_k)$.

  \item Ainsi $T_n(x)= \lambda \prod_{k=0}^{n-1}(x-\omega_k)$ et les $\omega_k$ sont deux à deux distincts.
  On sait par la question précédente que la décomposition en éléments simple de $1/T_n$ s'écrit 
$$\frac{1}{T_n(X)}=\sum_{k=0}^{n-1} \frac{1/T_n'(\omega_k)}{X-\omega_k}$$
Repartant de $T_n(x)=\cos\big(n \arccos(x)\big)$, on calcule 
$T_n'(x)=\frac{n}{\sqrt{1-x^2}}\sin\big(n\arccos(x)\big)$.
En utilisant que 
$$\sin\big(n \arccos(\omega_k)\big) = \sin \left( n\left(\frac{(2k+1)\pi}{2n}\right)\right) 
= \sin\left(\frac{\pi}{2}+k\pi \right) = (-1)^k$$
et 
$$\sqrt{1-\cos^2 \theta} = \sqrt{\sin^2 \theta} = \sin \theta \qquad \text{ pour } \theta \in [0,\pi]$$
on trouve que 
$T_n'(\omega_k)=\frac{n(-1)^k}{\sin\left(\frac{(2k+1)\pi}{2n}\right)}$.

Ainsi 
$$\frac{1}{T_n(X)}=\sum_{k=0}^{n-1}\frac{\frac{(-1)^k}{n}\sin\left(\frac{(2k+1)\pi}{2n}\right)}{X-\cos\left(\frac{(2k+1)\pi}{2n}\right)}$$
  \end{enumerate}
\end{enumerate}
\fincorrection
\nocorrection
\nocorrection
\correction{000426}
 Utiliser la formule d'interpolation de Lagrange~!
$P={1\over3}(X^2-4X-3)$.
\fincorrection
\correction{000427}
On cherche $P$ sous la forme $P(X)=aX^3+bX^2+cX+d$, ce qui donne le système linéaire suivant à résoudre:
$$\left\{\begin{array}{rcrcrcrcr} &&&&&&d&=&1\\a&+&b&+&c&+&d&=&0\\-a&+&b&-&c&+&d&=&-2\\8a&+&4b&+&2c&+&d&=&4\end{array}\right.$$
Après calculs, on trouve une unique solution :
 $a=\frac{3}{2}$, $b=-2$, $c=-\frac{1}{2}$, $d=1$ c'est-à-dire 
 $$P(X)=\frac{3}{2}X^3-2X^2-\frac{1}{2}X+1.$$
\fincorrection
\nocorrection
\nocorrection
\nocorrection
\nocorrection
\nocorrection
\nocorrection
\nocorrection
\nocorrection
\nocorrection
\nocorrection
\nocorrection
\nocorrection
\nocorrection
\nocorrection
\nocorrection
\nocorrection
\nocorrection
\nocorrection
\correction{003167}
$P(X) = -1 + Q(X)\times(X-1)^n \Leftrightarrow (X+1)^n \mid Q(X)(X-1)^n-2
\Leftrightarrow X^n \mid Q(X-1)(X-2)^n-2$.
Soit $2 = A(X)(X-2)^n + X^nB(X)$ la division suivant les puissances croissantes
de $2$ par $(X-2)^n$ {\`a} l'ordre~$n$.
On obtient $X^n\mid Q(X-1)-A(X)$ soit $Q(X) = A(X+1) + X^nR(X)$ et
$\deg(P)< 2n \Leftrightarrow R=0$.
Calcul de $A(X)$ par d{\'e}veloppement limit{\'e}~:  
$\frac{1}{(1+x)^n} = \sum_{k=0}^{n-1}\binom{-n}{k}x^k + O(x^n)$ donc~:
$$
A(X) = \frac{(-1)^n}{2^{n-1}}\sum_{k=0}^{n-1}\binom{-n}{k}\frac{(-1)^kX^k}{2^k}
     = \sum_{k=0}^{n-1}C_{n+k-1}^k(-1)^n\frac{X^k}{2^{n+k-1}}
$$
\fincorrection
\nocorrection
\correction{003169}
vect$(P(X), P(X+1), \dots, P(X+n))$ contient
$P$, $\Delta P$, $\Delta^2P$, ..., $\Delta^nP$ donc ${ K_n[X]}$
d'apr{\`e}s le thm des degr{\'e}s {\'e}tag{\'e}s.
\fincorrection
\correction{003170}
D{\'e}j{\`a} il est n{\'e}cessaire que $k=n$. Supposant ceci r{\'e}alis{\'e}, la matrice de~$(P_0,\dots,P_k)$
dans la base canonique de~$\C_n[X]$ est {\'e}quivalente {\`a} la matrice de Vandermonde de $z_0,\dots,z_k$.
Donc une CNS est~: $k=n$ et $z_0,\dots,z_k$ sont distincts.
\fincorrection
\correction{003171}
\begin{enumerate}
  \item $P\circ P - X = (P\circ P - P) + (P - X)$.
  \item $P(z) = z^2+3z+1  \Rightarrow  z=-1,-1,-2\pm i$.
\end{enumerate}
\fincorrection
\nocorrection
\correction{003173}
$\mathrm{Ker}\Phi =  K_0[X]$, $\Im\Phi = (X-a) K_{n-1}[X]$.
\fincorrection
\nocorrection
\correction{003175}
Formule de Taylor : $\frac {P^{(k)}}{k!} \in {\Z[X]}$.
\fincorrection
\correction{003176}
\begin{enumerate}
  \item 
  \item Appliquer le { 1)} {\`a} $P(X+k)$.
\end{enumerate}
\fincorrection
\nocorrection
\nocorrection
\correction{003179}
$\begin{cases}P = a(X^2+1) + bX + c \cr Q = a'(X^2+1) + b'X + c'\end{cases}
  \Rightarrow 
 \begin{cases}P = \cos\theta (X^2-1) + 2X\sin\theta \cr
        Q = \sin\theta (X^2-1) - 2X\cos\theta. \end{cases}$\par
$P \wedge Q = 1$ car $\pm i$ ne sont pas racines de $P$ et $Q$.
\fincorrection
\correction{003180}
$\deg P < 2  \Rightarrow  P \in \{1,X,X+1\}$.
\fincorrection
\correction{003181}
\begin{enumerate}
  \item Isomorphisme $P \longmapsto P(X) + P(X+1)$.
  \item $P_n' = nP_{n-1}$.
  \item $P_n(X+1) = \sum_{k=0}^n C_n^k P_k$\quad (Taylor).
  \item $Q_n(X) = P_n(1-X)  \Rightarrow  Q_n(X) + Q_n(X+1) = 2(-1)^nX^n$.
\end{enumerate}
\fincorrection
\correction{003182}
\begin{enumerate}
  \item Bezout g{\'e}n{\'e}ralis{\'e}.
  \item $\bigl((1-X)P' - nP\bigr) (1-X)^{n-1} + \bigl(nQ + XQ'\bigr)X^{n-1} = 0$.
  \item $P^{(k+1)}(0) = (n+k)P^{(k)}(0)  \Rightarrow 
               P = \sum_{k=0}^{n-1} C_{n+k-1}^k X^k$.
  \item 
\end{enumerate}
\fincorrection
\nocorrection
\correction{003184} $Q = P + P' + P'' + \dots$ : $Q(x)\to+\infty$ lorsque
$x\to+\infty$, donc il existe $\alpha \in \R$ tel que $Q(\alpha)$ soit
minimal.\par
Alors $0 = Q'(\alpha) = Q(\alpha) - P(\alpha) \Rightarrow \min Q \ge
0$.  \fincorrection \finexercice
\correction{003185}
oui ssi $P$ est pair.
\fincorrection
\correction{003186}
\begin{enumerate}
\item $P_0(u) = 2$, $P_1(u) = u$, $P_{n+1}(u) = uP_n(u) - P_{n-1}(u)$.
  \item $u_k = 2\cos\Bigl(\frac {(2k+1)\pi}{2n}\Bigr),\ k = 0, \dots, n-1$.
  \item 
\end{enumerate}
\fincorrection
\correction{003187}
\begin{enumerate}
  \item 
  \item Trivialement vrai ou trivialement faux selon le choix qu'on
    a fait en {\bf 1}.
  \item 
  \item Soit $Q\in \C[X]$ et $F_Q = \{\overline{RQ},\ R\in\C_n[X]\}$.
    On a $F_Q = \{\overline{RQ},\ R\in\C[X]\}$ de mani{\`e}re {\'e}vidente, donc
    $F_Q$ est stable par la multiplication modulaire par~$X$.
    
    Soit r{\'e}ciproquement $F$ un sev de $\C_n[X]$ stable par la multiplication
    modulaire par $X$. Si $(P_1,\dots,P_k)$ est une famille g{\'e}n{\'e}ratrice de~$F$
    alors $Q = \mathrm{pgcd}(P_1,\dots,P_k)\in F$ d'apr{\`e}s la relation de
    B{\'e}zout et la stabilit{\'e} de~$F$ donc $F_Q\subset F$ et $P_i\in F_Q$ puisque
    $Q$ divise $P_i$ d'o{\`u} $F\subset F_Q$ et $F=F_Q$.
\end{enumerate}
\fincorrection
\correction{003188}
Tout polyn{\^o}me {\`a} coefficients complexes non constant est surjectif sur~$\C$
donc $P(\C)\subset\R \Leftrightarrow P = a$ (constante r{\'e}elle).

On a par interpolation de Lagrange~: $P(\Q)\subset\Q \Leftrightarrow P\in\Q[X]$.

Montrons que $P(\Q)=\Q \Leftrightarrow P=aX+b$ avec $a\in\Q^*$, $b\in\Q$~: la condition
est clairement suffisante. Pour prouver qu'elle est n{\'e}cessaire, consid{\'e}rons
un polyn{\^o}me {\'e}ventuel $P$ de degr{\'e} $n\ge 2$ tel que $P(\Q)=\Q$.
On sait d{\'e}j{\`a} que $P$ est {\`a} coefficients rationnels, donc on peut
l'{\'e}crire sous la forme~: $P = \frac{a_0+a_1X+\dots+a_nX^n}d$
avec $a_i\in\Z$, $a_n\ne 0$ et $d\in\N^*$. Soit $\pi$ un nombre premier ne divisant
ni $a_n$ ni $d$, et $x=p/q$ (forme irr{\'e}ductible) un rationnel tel que
$P(x) = 1/\pi$. On a donc~:
$\pi(a_0q^n+\dots+a_np^n) = dq^n$ ce qui implique que $\pi$ divise $q$.
Il vient alors~: $a_np^n = dq^n/\pi - a_0q^n - \dots -a_{n-1}qp^{n-1}$ ce qui
est impossible puisque $\pi$ est facteur du second membre ($n\ge 2$) mais
pas du premier ($p\wedge q = 1$).
\fincorrection
\correction{003189}
Clairement $E=\varnothing$ si les $y_i$ ne sont pas distincts.
Si $y_1,\dots,y_n$ sont distincts, soit $P\in E$, $n=\deg(P)$ et $\lambda$
le coefficient dominant de~$P$ ($P\ne0$ car les $y_i$ ne sont pas tous nuls).
Alors $P(X) - y_i$ a pour seule racine $x_i$ donc $P(X) - y_i = \lambda(X-x_i)^n$.
Pour $n=1$ on obtient $P(X) = y_1 + \lambda(X-x_1)$ avec $\lambda\in\C^*$.
Pour $n\ge2$ on obtient $y_2-y_1 = \lambda(X-x_1)^n - \lambda(X-x_2)^n = n\lambda X^{n-1}(x_2-x_1) + \dots$
ce qui est impossible donc $E=\varnothing$.
\fincorrection
\correction{003190}
\begin{enumerate}
  \item R{\'e}currence sur $\mathrm{Card}\,(S)$ en mettant le terme de plus bas degr{\'e} en
facteur et en d{\'e}rivant le quotient.
  \item Appliquer la question pr{\'e}c{\'e}dente aux suites $(\Re(a_s))$ et $(\Im(a_s))$.
\end{enumerate}
\fincorrection
\correction{003191}
Soit~$f(x) = \sum_{k=1}^{100}\frac k{x-k}$.
$f$ est strictement d{\'e}croissante de~$0$ {\`a}~$-\infty$ sur~$]-\infty,0[$,
de~$+\infty$ {\`a} $-\infty$ sur chaque
intervalle~$]k,k+1[$,~$1\le k\le 100$ et de~$+\infty$ {\`a}~$0$ sur~$]100,+\infty[$.
Donc il existe $1<\alpha_1<2<\alpha_2<\dots<\alpha_{99}<100<\alpha_{100}$
tels que $E = \{x\in\R\text{ tq }f(x)\ge 1\} = \bigcup_{k=1}^{100}]k,\alpha_k]$.

La somme des longueurs est $L = \sum_{k=1}^{100}\alpha_k-\sum_{k=1}^{100}k$
et $\alpha_1,\dots,\alpha_{100}$ sont les racines du polyn{\^o}me :
$$P(X) = \prod_{k=1}^{100}(X-k) - \sum_{k=1}^{100}k\prod_{i\ne k}(X-i)
= X^{100}-2X^{99}\sum_{k=1}^{100}k + \dots$$
D'o{\`u} $\sum_{k=1}^{100}\alpha_k = 2\sum_{k=1}^{100}k$ et $L=\sum_{k=1}^{100}k=5050$.
\fincorrection
\correction{003192}
Le sens $\Leftarrow$ est trivial. Pour le sens $ \Rightarrow $, il suffit de v{\'e}rifier
la propri{\'e}t{\'e} lorsque $P$ est irr{\'e}ductible, strictement positif sur~$\R^+$, et
le seul cas non trivial est celui o{\`u} $P$ est de la forme~:
$P = (X-a)^2 + b^2$ avec $a>0$, $b>0$. Dans ce cas, le coefficient de~$X^k$
dans $(X+1)^\ell P(X)$ est~:
$C_\ell^k(a^2+b^2) - 2aC_\ell^{k-1} + C_\ell^{k-2}$, en convenant que $C_x^y$
vaut~$0$ si l'on n'a pas $0\le y\le x$. En mettant ce qui peut l'{\^e}tre en facteur
et en ordonnant le reste suivant les puissances de~$k$, on
est rammen{\'e} {\`a} montrer que la quantit{\'e}~:
$$k^2(a^2+b^2+2a+1) - k((a^2+b^2)(2\ell+3) + 2a(\ell+2) + 1) + \ell^2(a^2+b^2)$$
est strictement positive pour tout $k\in[[0,\ell+2]]$ si~$\ell$ est choisi convenablement.
Or le discriminant par rapport {\`a}~$k$ est {\'e}quivalent {\`a} $-4\ell^2(2a+1)$ lorsque
$\ell$ tend vers $+\infty$ donc un tel choix de~$\ell$ est possible.
\fincorrection
\correction{003193}
On suppose $E$ fini et on montre que~$P$ est constant~: il existe~$a\in\Z$
tel que~$P(a)\ne 0$. Soit $N = \prod_{p\in E}p^{1+v_p(P(a))}$.
Alors pour tout~$k\in\Z$, $P(a+kN)\equiv P(a) (\mathrm{mod}\, N)$ (formule de Taylor),
donc $v_p(P(a+kN)) = v_p(P(a))$ pour tous~$k\in\Z$ et~$p\in E$.
Comme $P(a+kN)$ est produit d'{\'e}l{\'e}ments de~$E$, on en d{\'e}duit que $P(a+kN) = \pm P(a)$
pour tout~$k$, donc $P$ prend une infinit{\'e} de fois la m{\^e}me valeur.
\fincorrection
\correction{003194}
Prendre pour~$P_n$ la partie r{\'e}guli{\`e}re du d{\'e}veloppement limit{\'e} {\`a} l'ordre~$n$
de~$\sqrt{1+x}$.
\fincorrection
\correction{003195}
Analyse~: on pose $P_j = a_0 + a_1X + \dots + a_nX^n$ et on consid{\`e}re la fraction rationnelle
$$F(X) = \frac{a_0}{X} + \frac{a_1}{X+1} + \dots + \frac{a_n}{X+n}
= \frac{P(X)}{X(X+1)\dots(X+n)}.$$
Alors $ \int_{t=0}^1 t^iP_j(t)\,dt = F(i+1) = \frac{i!\,P(i+1)}{(i+n+1)!}$
donc $P(j+1) = \frac{(j+n+1)!}{j!}$ et $P(k) = 0$ pour $k\in{[[1,n+1]]}\setminus\{j+1\}$,
soit $$
\begin{aligned}
  P(X) &= \frac{(j+n+1)!}{j!}\prod_{k\ne j+1}\frac{X-k}{j+1-k} \\&=
  (-1)^{n-j}\frac{(j+n+1)!}{(j!)^2(n-j)!}\prod_{k\ne j+1}(X-k) =
  Q_j(X).
\end{aligned}
$$

Synth{\`e}se : soit $Q_j$ le polyn{\^o}me ci-dessus et $a_0,\dots,a_n$ les
coefficients de la d{\'e}composition en {\'e}l{\'e}ments simples de $\frac{Q_j(X)}{X(X+1)\dots(X+n)}$.
On doit juste v{\'e}rifier que les $a_i$ sont entiers. Calcul~:
$$
\begin{aligned}
  a_i &= \frac{Q_j(-i)}{(-1)^ii!\,(n-i)!}  =
  (-1)^{i+j}\frac{(i+j)!\,(i+n+1)!\,(j+n+1)!}{(i+j+1)!\,(i!)^2\,(j!)^2\,(n-i)!\,(n-j)!}
  \\ &= (-1)^{i+j}C_{i+j}^iC_{i+n+1}^{i+j+1}C_{j+n+1}^{j}C_n^i(n+1) \in\Z.
\end{aligned}
$$
\fincorrection
\correction{003254}
$3\Bigl(\frac ab + \frac ac + \frac ba + \frac bc + \frac ca
          + \frac cb\big) = -9$.
\fincorrection
\correction{003255}
$-\frac 53$.
\fincorrection
\correction{003256}
$\frac 13$.
\fincorrection
\correction{003257}
$x^7 = -2x^2 + 2x - 1  \Rightarrow  a^7 + b^7 + c^7 = -7$.
\fincorrection
\correction{003258}
$\{ a,b,c \} = \left\{ 1, -\frac {1+i}2, -\frac {1-i}2 \right\}$.
\fincorrection
\correction{003259}
$\{x,y,z\}=\{-1,1,2\}$.
\fincorrection
\correction{003260}
$d = \frac 23$, $\{ a,b,c \} = \left\{ 0, \pm\frac 1{\sqrt2} \right\}$.
\fincorrection
\nocorrection
\correction{003262}
\begin{enumerate}
  \item $50p^3 = 27q^2$.
  \item $a^2+b^2+c^2 = 2c^2 + 1 = -2p  \Rightarrow $ l'une des racines de l'{\'e}quation
      aux carr{\'e}s ($-Y^3 -2pY^2 -p^2Y + q^2$) doit {\^e}tre $-p - \frac 12$.
      CNS $\Leftrightarrow 2p+1 + 8q^2 = 0$.
\end{enumerate}
\fincorrection
\correction{003263}
$20p^3 + 27q^2 = 0$.
\fincorrection
\correction{003264}
$b=0,\ c = \frac 9{100}a^2  \Rightarrow $ racines : $-3x,-x,x,3x$ avec
$x = \sqrt{\frac {-a}{10}}$.
\fincorrection
\correction{003265}
$-P(-2-X) = X^3 + 4X^2 + 7X + 2$.
\fincorrection
\correction{003266}
$X^3 + (2b-a^2)X^2 + (b^2-2ac)X -c^2$.
\fincorrection
\correction{003267}
 racine 1 : $\lambda = -6$.\par
 racine -1 : $\lambda = -4$.\par
 racine $\alpha \in \mathbb{U}\setminus\{\pm1\}$ : les autres sont $\frac 1\alpha$ et
 $-\lambda  \Rightarrow  \lambda = 6$, $\alpha = \frac{1+i\sqrt{15}}4$.
\fincorrection
\correction{003268}
\begin{enumerate}
  \item Les coefficients de $P$ sont born{\'e}s.
  \item $\widetilde P(X^2) = (-1)^nP(X)P(-X)  \Rightarrow  \widetilde P \in {\Z[X]}$.
  \item La suite $(\widetilde {\dot{\dot{\widetilde P}}})$ prend un nombre
      fini de valeurs.
\end{enumerate}
\fincorrection
\correction{003269}
Soit $y = \gamma x + \delta$ l'{\'e}quation de la droite en question.
On veut que l'{\'e}quation $x^4 + ax^3 + bx^2 + (c-\gamma)x + (d-\delta) = 0$
ait quatre racines distinctes en progression arithm{\'e}tique. Si $r$ est la raison
de cette progression alors les racines sont
$-\frac a4 -\frac32r$, $-\frac a4 -\frac12r$, $-\frac a4 +\frac12r$ et $-\frac a4 +\frac32r$.
On doit donc chercher {\`a} quelle condition sur $a,b,c,d$ il existe $\gamma,\delta,r$ r{\'e}els
tels que ~:
$$\begin{aligned}
&\Bigl(-\frac a4 -\frac32r\Bigr)\Bigl(-\frac a4 -\frac12r\Bigr) +
\dots + \Bigl(-\frac a4 +\frac12r\Bigr)\Bigl(-\frac a4 +\frac32r\Bigr)
= -b\\ 
&
\Bigl(-\frac a4 -\frac32r\Bigr)\Bigl(-\frac a4
-\frac12r\Bigr)\Bigl(-\frac a4 + \frac12r\Bigr) + \dots \\
&\qquad\qquad\qquad+ \bigl(-\frac
a4 - \frac12r\Bigr)\Bigl(-\frac a4 +\frac12r\Bigr)\Bigl(-\frac a4
+\frac32r\Bigr) = c-\gamma\\ 
&\Bigl(-\frac a4 -\frac32r\Bigr)\Bigl(-\frac a4
-\frac12r\Bigr)\Bigl(-\frac a4 + \frac12r\Bigr)\Bigl(-\frac a4 +
\frac32r\Bigr) = \delta - d 
\end{aligned}$$
Les deux derni{\`e}res {\'e}quations sont satisfaites {\`a} $r$ donn{\'e} en choisissant convenablement $\gamma$ et $\delta$.
La premi{\`e}re s'{\'e}crit apr{\`e}s simplifications~: $\frac52r^2 = \frac38a^2 + b$ et la condition
demand{\'e}e est $3a^2 + 8b> 0$.
\fincorrection
\correction{005314}
Tout d'abord
$$Q=(1+X+...+X^n)'=(\frac{X^{n+1}-1}{X-1})'=\frac{(n+1)X^n(X-1)-X^{n+1}}{(X-1)^2}=\frac{nX^{n+1}-(n+1)X^n+1}{(X-1)^2}.$$

Ensuite, $\omega_0=1$ et donc, $Q(\omega_0)=1+2+...+n=\frac{n(n+1)}{2}$. Puis, pour $1\leq k\leq n-1$, $\omega_k\neq1$ et donc, puisque $\omega_k^n=1$,

$$Q(\omega_k)=\frac{n\omega_k^{n+1}-(n+1)\omega_k^n+1}{(\omega_k-1)^2}=\frac{n\omega_k-(n+1)+1}{(\omega_k-1)^2}
=\frac{n}{\omega_k-1}.$$

Par suite, 

$$\prod_{k=0}^{n-1}Q(\omega_k)=\frac{n(n+1)}{2}\prod_{k=1}^{n-1}\frac{n}{\omega_k-1}=\frac{n^n(n+1)}{2\prod_{k=1}^{n-1}(\omega_k-1)}.$$

Mais, $X^n-1=(X-1)(1+X+...+X^{n-1})$ et d'autre part $X^n-1=\prod_{k=0}^{n-1}(X-e^{2ik\pi/n})=(X-1)\prod_{k=1}^{n-1}(X-\omega_k)$. Par intégrité de $\Rr[X]$, $\prod_{k=1}^{n-1}(X-e^{2ik\pi/n})=1+X+...+X^{n-1}$ (Une autre rédaction possible est~:~$\forall z\in\Cc,\;(z-1)\prod_{k=1}^{n-1}(z-\omega_k)=(z-1)(1+z+...+z^{n-1})$ et donc $\forall z\in\Cc\setminus\{1\}$, $\prod_{k=1}^{n-1}(z-\omega_k)=1+z+...+z^{n-1}$ et finalement $\forall z\in\Cc,\;\prod_{k=1}^{n-1}(z-\omega_k)=1+z+...+z^{n-1}$ car les deux polynômes ci-contre coincident en une infinité de valeurs de $z$.)

En particulier, $\prod_{k=1}^{n-1}(1-\omega_k)=1+1^2+...+1^{n-1}=n$ ou encore $\prod_{k=1}^{n-1}(\omega_k-1)=(-1)^{n-1}n$.
Donc,

$$\prod_{k=0}^{n-1}Q(\omega_k)=\frac{n^n(n+1)}{2}\frac{1}{(-1)^{n-1}n}=\frac{(-1)^{n-1}n^{n-1}(n+1)}{2}.$$

\fincorrection
\correction{005320}
Si $P$ est de degré inférieur ou égal à $0$, c'est clair.

Sinon, posons $P=\sum_{k=0}^{n}a_kX^k$ avec $n\in\Nn^*$.

\begin{align*}\ensuremath
P(P(X))-X&=P(P(X))-P(X)+P(X)-X=\sum_{k=0}^{n}a_k((P(X))^k-X^k)+(P(X)-X)\\
 &=\sum_{k=1}^{n}a_k((P(X))^k-X^k)+(P(X)-X).
\end{align*}

Mais, pour $1\leq k\leq n$, $(P(X))^k-X^k)=(P(X)-X)((P(X))^{k-1}+X(P(X))^{k-2}+...+X^{k-1})$ est divisible par $P(X)-X$ et il en est donc de même de $P(P(X))-X$.
\fincorrection
\correction{005321}
\begin{enumerate}
\item  Posons $P=\sum_{i=0}^{l}a_iX_i$ où $l\geq1$ et où les $a_i$ sont des entiers relatifs avec $a_l\neq0$.

$$P(n+km)=\sum_{i=0}^{l}a_i(n+km)^i=\sum_{i=0}^{l}a_i(n^i+K_im)=\sum_{i=0}^{l}a_in^i+Km=m+Km=m(K+1),$$

où $K$ est un entier relatif. $P(n+km)$ est donc un entier relatif multiple de $m=P(n)$.

\item  Soit $P\in\Zz[X]$ tel que $\forall n\in\Nn,\;P(n)$ est premier.

Soit $n$ un entier naturel donné et $m=P(n)$ (donc, $m\geq2$ et en particulier $m\neq0$). Pour tout entier relatif $k$, $P(n+km)$ est divisible par $m$ mais $P(n+km)$ est un nombre premier ce qui impose $P(n+km)=m$. Par suite, le polynôme $Q=P-m$ admet une infinité de racines deux à deux distinctes (puisque $m\neq0$) et est donc le polynôme nul ou encore $P$ est constant.
\end{enumerate}
\fincorrection
\correction{005322}
\begin{enumerate}
\item  Déjà, $P_0$ est dans $E$.

Soit $n$ un naturel non nul. $P_n=\frac{1}{n!}(X+1)...(X+n)$ et donc, si $k$ est élément de $\{-1,...,-n\}$, $P_n(k)=0\in\Zz$.

Si $k$ est un entier positif, $P_n(k)=\frac{1}{n!}(k+1)...(k+n)=C_{n+k}^n\in\Zz$.

Enfin, si $k$ est un entier strictement plus petit que $-n$, 

$$P_n(k)=\frac{1}{n!}(k+1)...(k+n)=(-1)^n\frac{1}{n!}(-k-1)...(-k-n)=(-1)^nC_{-k-1}^n\in\Zz.$$

Ainsi, $\forall k\in\Zz,\;P_(k)\in\Zz$, ou encore $P_(\Zz)\subset\Zz$.

\item  Evident

\item  Soit $P\in\Cc[X]\setminus\{0\}$ tel que $\forall k\in\Zz,\;P(k)\in\Zz$ (si $P$ est nul, $P$ est combinaison linéaire à coefficients entiers des $P_k$).

Puisque $\forall k\in\Nn,\;\mbox{deg}(P_k)=k$, on sait que pour tout entier naturel $n$, $(P_k)_{0\leq k\leq n}$ est une base de $\Cc_n[X]$ et donc, $(P_k)_{k\in\Nn}$ est une base de $\Cc[X]$ (tout polynôme non nul ayant un degré $n$, s'écrit donc de manière unique comme combinaison linéaire des $P_k$).

Soit $n=\mbox{deg}P$.

Il existe $n+1$ nombres complexes $a_0$,..., $a_n$ tels que $P=a_0P_0+...+a_nP_n$. Il reste à montrer que les $a_i$ sont des entiers relatifs.

L'égalité $P(-1)$ est dans $\Zz$, fournit~:~$a_0$ est dans $\Zz$.

L'égalité $P(-2)$ est dans $\Zz$, fournit~:~$a_0-a_1$ est dans $\Zz$ et donc $a_1$ est dans $\Zz$.

L'égalité $P(-3)$ est dans $\Zz$, fournit~:~$a_0-2a_1+a_2$ est dans $\Zz$ et donc $a_2$ est dans $\Zz$...

L'égalité $P(-(k+1))$ est dans $\Zz$, fournit~:~$a_0-a_1+...+(-1)^ka_k$ est dans $\Zz$ et si par hypothèse de récurrence, $a_0$,..., $a_{k-1}$ sont des entiers relatifs alors $a_k$ l'est encore.

Tous les coefficients $a_k$ sont des entiers relatifs et $E$ est donc constitué des combinaisons linéaires à coefficients entiers relatifs des $P_k$.
\end{enumerate}
\fincorrection
\correction{005326}
Soit $P$ un tel polynôme. $-2$ est racine de $P+10$ d'ordre au moins trois et donc racine de $(P+10)'= P'$ d'ordre au moins deux.

De même, $2$ est racine de $P'$ d'ordre au moins deux et puisque $P'$ est de degré $4$, il existe un complexe $\lambda$ tel que $P'=\lambda(X-2)^2(X+2)^2=\lambda(X^2-4)^2=\lambda(X^4-8X^2+16)$ et enfin, nécessairement,

$$\exists(\lambda,\mu)\in\Cc^2/\;P=\lambda(\frac{1}{5}X^5-\frac{8}{3}X^3+16X)+\mu\;\mbox{avec}\;\lambda\neq0.$$

Réciproquement, soit $P=\lambda(\frac{1}{5}X^5-\frac{8}{3}X^3+16X)+\mu$ avec $\lambda\neq0$.

\begin{align*}\ensuremath
P\;\mbox{solution}&\Leftrightarrow P+10\;\mbox{divisible par}\;(X+2)^3\;\mbox{et}\;P-10\;\mbox{est divisible par}\;(X-2)^3\\
 &\Leftrightarrow P(-2)+10=0=P'(-2)=P''(-2)\;\mbox{et}\;P(2)+10=0=P'(2)=P''(2)\Leftrightarrow P(-2)=-10\;\mbox{et}\;P(2)=10\\
 &\left\{
 \begin{array}{l}
 \lambda(-\frac{32}{5}+\frac{64}{3}-32)+\mu=-10\\
  \lambda(\frac{32}{5}-\frac{64}{3}+32)+\mu=10
  \end{array}
  \right.
  \Leftrightarrow\mu=0\;\mbox{et}\;\lambda(\frac{32}{5}-\frac{64}{3}+32)+\mu=10\\
  &\Leftrightarrow\mu=0\;\mbox{et}\;\lambda=\frac{75}{128}
\end{align*}

On trouve un et un seul polynôme solution à savoir $P=\frac{75}{128}(\frac{1}{5}X^5-\frac{8}{3}X^3+16X)=\frac{15}{128}X^5-\frac{25}{16}X^3+\frac{75}{8}X$.
\fincorrection
\correction{005327}
Les polynômes de degré inférieur ou égal à $0$ solutions sont clairement $0$ et $1$.

Soit $P$ un polynôme de degré supérieur ou égal à $1$ tel que $P(X^2)=P(X)P(X+1)$.

Soit $a$ une racine de $P$ dans $\Cc$. Alors, $a^2$, $a^4$, $a^8$..., sont encore racines de $P$. Mais, $P$ étant non nul, $P$ ne doit admettre qu'un nombre fini de racines. La suite $(a^{2^n})_{n\in\Nn}$ ne doit donc prendre qu'un nombre fini de valeurs ce qui impose $a=0$ ou $|a|=1$ car si $|a|\in]0,1[\cap]1,+\infty[$, la suite $(|a^{2^n}|)$ est strictement monotone et en particulier les $a^{2^n}$ sont deux à deux distincts.

De même, si $a$ est racine de $P$ alors $(a-1)^2$ l'est encore mais aussi $(a-1)^4$, $(a-1)^8$..., ce qui impose $a=1$ ou $|a-1|=1$.

En résumé,

$$(a\;\mbox{racine de}\;P\;\mbox{dans}\;\Cc)\Rightarrow((a=0\;\mbox{ou}\;|a|=1)\;\mbox{et}\;(a=1\;\mbox{ou}\;|a-1|= 1))\Rightarrow(a=0\;\mbox{ou}\;a=1\;\mbox{ou}\;|a|=|a-1|=1).$$

Maintenant, $|a|=|a-1|=1\Leftrightarrow|a|=1\;\mbox{et}\;|a|=|a-1|\Leftrightarrow a\in\mathcal{C}((0,0),1)\cap\mbox{med}[(0,0),(1,0)]=\{-j,-j^2\}$.

Donc, si $P\in\Rr[X]$ est solution, il existe $K$, $\alpha$, $\beta$, $\gamma$, $K$ complexe non nul et $\alpha$, $\beta$ et $\gamma$ entiers naturels tels que $P=KX^\alpha(X-1)^\beta(X+j)^\gamma(X+j^2)^\gamma$ ($-j$ et $-j^2$ devant avoir même ordre de multiplicité).

Réciproquement, si $P=KX^\alpha(X-1)^\beta(X+j)^\gamma(X+j^2)^\gamma=KX^\alpha(X-1)^\beta(X^2-X+1)^\gamma$.
$$P(X^2)=KX^{2\alpha}(X^2-1)^\beta(X4-X^2+1)^\gamma=KX^{2\alpha}(X-1)^\beta(X+1)^\beta(X^2-\sqrt{3}X+1)^\gamma(X^2+
\sqrt{3}X+1)^\gamma,$$

et 

\begin{align*}\ensuremath
P(X)P(X+1)&=KX^\alpha(X-1)^\beta(X^2-X+1)\gamma K(X+1)^\alpha X^\beta(X^2+X+1)^\gamma\\
 &=K^2X^{\alpha+\beta}(X-1)^\beta(X+1)^\alpha(X^2-X+1)^\gamma(X^2+X+1)^\gamma.
\end{align*}

Par unicité de la décompôsition en produit de facteurs irréductibles d'un polynôme non nul, $P$ est solution si et seulement si $P=0$ ou $K=1$ et $\alpha=\beta$ et $\gamma=0$.

Les polynômes solutions sont $0$ et les $(X^2-X)^\alpha$ où $\alpha$ est un entier naturel quelconque.
\fincorrection
\correction{005330}
\begin{enumerate}
\item  
\begin{align*}\ensuremath
S&\Leftrightarrow\left\{
\begin{array}{l}
x+y+z=1\\
\frac{xy+xz+yz}{xyz}=1\\
xyz=-4
\end{array}\right.\Leftrightarrow\sigma_1=1,\;\sigma_2=\sigma_3=-4\\
 &\Leftrightarrow x,\;y\;\mbox{et}\;z\;\mbox{sont les trois solutions de l'équation}\;X^3-X^2-4X+4=0\\
 &\Leftrightarrow x,\;y\;\mbox{et}\;z\;\mbox{sont les trois solutions de l'équation}\;(X-1)(X-2)(X+2)=0\\
 &\Leftrightarrow(x,y,z)\in\{(1,2,-2),(1,-2,2),(2,1,-2),(2,-2,1),(-2,1,2),(-2,2,1)\}
\end{align*}

\item  Pour $1\leq k\leq 4$, posons $S_k=x^k+y^k+z^k+t^k$. On a $S_2=\sigma_1^2-2\sigma_2$. Calculons $S_3$ en fonction des $\sigma_k$. On a $\sigma_1^3=S_3+3\sum_{}^{}x^2y+6\sum_{}^{}xyz=S_3+3\sum_{}^{}x^2y+6\sigma_3$ $(*)$. Mais on a aussi $S_1S_2=S_3+\sum_{}^{}x^2y$. Donc, $\sum_{}^{}x^2y=\sigma_1(\sigma_1^2-2\sigma_2)-S_3$. En reportant dans $(*)$, on obtient $\sigma_1^3=S_3+3(\sigma_1^3-2\sigma_1\sigma_2-S_3)+6\sigma_3$ et donc,

$$S_3=\frac{1}{2}(-\sigma_1^3+3(\sigma_1^3-2\sigma_1\sigma_2-S_3)+6\sigma_3)=\sigma_1^3-3\sigma_1\sigma_2+3\sigma_3.$$

Calculons $S_3$ en fonction des $\sigma_k$. Soit $P=(X-x)(X-y)(X-z)(X-t)=X^4-\sigma_1X^3+\sigma_2X^2-\sigma_3X+\sigma_4$.

\begin{align*}\ensuremath
P(x)+P(y)+P(z)+P(t)=0&\Leftrightarrow S_4-\sigma_1S_3+\sigma_2S_2-\sigma_3S_1+4\sigma_4=0\\
 &\Leftrightarrow S_4=\sigma_1(\sigma_1^3-3\sigma_1\sigma_2+3\sigma_3)-\sigma_2(\sigma_1^2-2\sigma_2)+\sigma_3\sigma_1-4\sigma_4\\
 &\Leftrightarrow S_4=\sigma_1^4-4\sigma_1^2\sigma_2+4\sigma_1\sigma_3+2\sigma_2^2-4\sigma_4.
\end{align*}

Par suite,

\begin{align*}\ensuremath
S&\Leftrightarrow
\left\{
\begin{array}{l}
\sigma_1=0\\
-2\sigma_2=10\\
3\sigma_3=0\\
2\sigma_2^2-4\sigma_4=26
\end{array}
\right.\Leftrightarrow\left\{
\begin{array}{l}
\sigma_1=0\\
\sigma_2=-5\\
\sigma_3=0\\
\sigma_4=6
\end{array}
\right.\\
 &\Leftrightarrow x,\;y,\;z,\;\mbox{et}\;t\;\mbox{sont les 4 solutions de l'équation}\;X^4-5X^2+6=0\\
 &\Leftrightarrow(x,y,z,t)\;\mbox{est l'une des 24 permutations du quadruplet}\;(\sqrt{2},-\sqrt{2},\sqrt{3},-\sqrt{3})
\end{align*}

\end{enumerate}
\fincorrection
\correction{005331}
Le polynôme nul est solution. Soit $P$ un polynôme non nul de degré $n$ solution alors $n=n-1+n-2$ et donc $n=3$. Posons donc $P=aX^3+bX^2+cX+d$ avec $a\neq0$.

\begin{align*}\ensuremath
P(2X)=P'(X)P''(X)&\Leftrightarrow8aX^3+4bX^2+2cX+d=(3aX^2+2bX+c)(6aX+2b)\\
 &\Leftrightarrow(18a^2-8a)X^3+(18ab-4b)X^2+(4b^2+6ac-2c)X+2bc-d=0\\
 &\Leftrightarrow18a^2-8a=18ab-4b=4b^2+6ac-2c=2bc-d=0\\
 &\Leftrightarrow a=\frac{4}{9}\;\mbox{et}\;b=c=d=0.
\end{align*}

Les polynômes solutions sont $0$ et $\frac{4}{9}X^3$.

\fincorrection
\correction{005332}
$0$ n'est pas racine de $P$.

On rappelle que si $r=\frac{p}{q}$, ($p\in\Zz^*$, $q\in\Nn^*$, $p\wedge q=1$) est racine de $P$, alors $p$ divise le coefficient constant de $P$ et $q$ divise son coefficient dominant. Ici, $p$ divise $4$ et $q$ divise $12$ et donc, $p$ est élément de $\{\pm1,\pm2,\pm4\}$ et $q$ est élément 
de $\{1,2,3,4,6,12\}$ ou encore $r$ est élément de $\{\pm1,\pm2,\pm4,\pm\frac{1}{2},\pm\frac{1}{3},\pm\frac{2}{3},\pm\frac{4}{3},\pm\frac{1}{4},\pm\frac{1}{6},\pm\frac{1}{12}\}$.

Réciproquement, on trouve $P(\frac{2}{3})=P(\frac{1}{4})=0$. $P$ est donc divisible par

$$12(X-\frac{2}{3})(X-\frac{1}{4})=(3X-2)(4X-1)=12X^2-11X+2.$$

Plus précisément, $P=(12X^2-11X+2)(X^2+X+2)=(3X-2)(4X-1)(X-\frac{-1+i\sqrt{7}}{2})(X-\frac{-1-i\sqrt{7}}{2})$.
\fincorrection
\correction{005333}
Pour $n\geq0$, posons $P_n=(X-1)^{2n}-X^{2n}+2X-1$. $P_n(0)=P_n(1)=P_n(\frac{1}{2})=0$. $P_n$ admet $0$, $1$ et $\frac{1}{2}$ pour racines et est donc divisible par $X(X-1)(2X-1)=2X^3-3X^2+X$.

Si $n=0$ ou $n=1$, le quotient est nul. Si $n=2$, le quotient vaut $-2$.

Soit $n\geq 3$. On met succesivement $2X-1$ puis $X-1$ puis $X$ en facteur~:
\begin{align*}\ensuremath
P_n&=((X-1)^2)^n-(X^2)^n+(2X-1)=((X-1)2-X2)\sum_{k=0}^{n-1}(X-1)^{2k}X^{2(n-1-k)}+(2X-1)\\
 &=(2X-1)(-\sum_{k=0}^{n-1}(X-1)^{2k}X^{2(n-1-k)}+1)=(2X-1)(-\sum_{k=1}^{n-1}(X-1)^{2k}X^{2(n-1-k)}+1-X^{2n-2})\\
 &=(2X-1)(-(X-1)\sum_{k=1}^{n-1}(X-1)^{2k-1}X^{2(n-1-k)}-(X-1)\sum_{k=0}^{2n-1}X^k)\\
 &=(2X-1)(X-1)(-\sum_{k=1}^{n-1}(X-1)^{2k-1}X^{2(n-1-k)}-\sum_{k=0}^{2n-1}X^k)\\
 &=(2X-1)(X-1)(-\sum_{k=1}^{n-2}(X-1)^{2k-1}X^{2(n-1-k)}-\sum_{k=1}^{2n-1}X^k-1-(X-1)^{2n-3})\\
 &=(2X-1)(X-1)(-\sum_{k=1}^{n-2}(X-1)^{2k-1}X^{2(n-1-k)}-\sum_{k=1}^{2n-3}X^k-\sum_{k=1}^{2n-3}(-1)^{2n-3-k}C_{2n-3}^kX^k)\\
 &= X(2X-1)(X-1)(-\sum_{k=1}^{n-2}(X-1)^{2k-1}X^{2n-2k-3}-
 \sum_{k=1}^{2n-3}X^{k-1}-\sum_{k=1}^{2n-3}(-1)^{2n-3-k}C_{2n-3}^kX^{k-1})
\end{align*}

\fincorrection
\correction{005334}
\begin{align*}\ensuremath
1&=(X+(1-X))^{n+m-1}=\sum_{k=0}^{n+m-1}C_{n+m-1}^{k}X^k(1-X)^{n+m-1-k}\\
 &=\sum_{k=0}^{n-1}C_{n+m-1}^{k}X^k(1-X)^{n+m-1-k}
+\sum_{k=n}^{n+m-1}C_{n+m-1}^{k}X^k(1-X)^{n+m-1-k}\\ 
 &=(1-X)^m\sum_{k=0}^{n-1}C_{n+m-1}^{k}X^k(1-X)^{n-1-k}+X^n\sum_{k=n}^{n+m-1}C_{n+m-1}^{k}X^{k-n}(1-X)^{n+m-1-k}
\end{align*}

Soient $U=\sum_{k=n}^{n+m-1}C_{n+m-1}^{k}X^{k-n}(1-X)^{n+m-1-k}$ et $V=\sum_{k=0}^{n-1}C_{n+m-1}^{k}X^k(1-X)^{n-1-k}$. $U$ et $V$ sont des polynômes tels que $UX^n+V(1-X)^m=1$. De plus, pour $n\leq k\leq n+m-1$, $\mbox{deg}(X^{k-n}(1-X)^{n+m-1-k})=k-n+n+m-1-k=m-1<m$ et donc $\mbox{deg}(U)<m$ et de même pour $0\leq k\leq n-1$, $\mbox{deg}(X^k(1-X)^{n-1-k})=k+n-1-k=n-1<n$ et $\mbox{deg}(V)<n$.
\fincorrection
\correction{005343}
On suppose $a_0\neq0$ de sorte que $0$ n'est pas racine de $P$. Soient $p$ un relatif non nul et $q$ un entier naturel non nul tels que $p$ et $q$ soient premiers entre eux.

Si $r=\frac{p}{q}$ est racine de $P$ alors $a_n(\frac{p}{q})^n+...+a_0=0$ et donc
 
$$a_np^n=-q(a_0q^{n-1}+...+a_{n-1}p^{n-1})\;\mbox{et}\;a_0q^n=-p(a^np^{n-1}+...+a_1q^{n-1}).$$

Donc $p$ divise $a_0q^n$, mais $p$ est premier à $q^n$ et d'après le théorème de \textsc{Gauss}, $p$ divise $a_0$. De même $q$ divise $a_n$.

Application. Soit $P=9X^4-3X^3+16X^2-6X-4$ et soit $p$ un entier relatif non nul et $q$ un entier naturel non nul tels que $p\wedge q=1$. Si $\frac{p}{q}$ est racine de $P$, $p$ divise $-4$ et $q$ divise $9$ de sorte que $p$ est élément de $\{1,-1,2,-2,4,-4\}$ et $q$ est élément de $\{1,3,9\}$ puis $\frac{p}{q}$ est élément de $\{\pm1,\pm2,\pm4,\pm\frac{1}{3},\pm\frac{2}{3},\pm\frac{4}{3},\pm\frac{1}{9},\pm\frac{2}{9},\pm\frac{4}{9}\}$.
On trouve $P(\frac{2}{3})=P(-\frac{1}{3})=0$ et P est divisible par $(3X-2)(3X+1)=9X^2-3X-2$.
Plus précisément $P=9X^4-3X^3+16X^2-6X-4=(9X^2-3X-2)(X^2+2)$ et les racines de $P$ dans $\Cc$ sont $\frac{2}{3}$, $-\frac{1}{3}$, $i\sqrt{2}$ et $-i\sqrt{2}$.
\fincorrection
\correction{005344}
\begin{enumerate}
\item  
\begin{align*}\ensuremath
P&=X^4+2X^3+3X^2+2X+1=X^2(X^2+\frac{1}{X^2}+2(X+\frac{1}{X})+3)=X^2((X+\frac{1}{X})^2+2(X+\frac{1}{X})+1)\\
 &=X^2(X+\frac{1}{X}+1)^2=(X^2+X+1)^2=(X-j)^2(X-j^2)^2.
\end{align*}

\item  $1$ et $-1$ sont racines de $P$. On écrit donc $P=(X^2-1)(X^4-5X^3+6X^2-5X+1)$ puis

\begin{align*}\ensuremath
X^4-5X^3+6X^2-5X+1&=X^2((X^2+\frac{1}{X^2})-5(X+\frac{1}{X})+6)=X^2((X+\frac{1}{X})^2-5(X+\frac{1}{X})+4)\\
 &=X^2(X+\frac{1}{X}-1)(X+\frac{1}{X}-4)=(X^2-X+1)(X^2-4X+1)
\end{align*}

et donc, $P=(X-1)(X+1)(X+j)(X+j^2)(X-2+\sqrt{3})(X-2-\sqrt{3})$.

\item  

\begin{align*}\ensuremath
P&=X^7-X^6-7X^5+7X^4+7X^3-7X^2-X+1=(X^2-1)(X^5-X^4-6X^3+6X^2+X-1)\\
 &=(X-1)^2(X+1)(X^4-6X^2+1)\\
 &=(X-1)^2(X+1)(X^2(3+2\sqrt{2}))(X^2-(3-2\sqrt{2}))
\end{align*}

Les racines de $P$ dans $\Cc$ sont $1$, $-1$, $\sqrt{3+2\sqrt{2}}$,$-\sqrt{3+2\sqrt{2}}$, $\sqrt{3-2\sqrt{2}}$ et $-\sqrt{3-2\sqrt{2}}$.
\end{enumerate}
\fincorrection
\correction{005346}
Soit $P$ un polynôme à coefficients complexes de degré $4$. On suppose $P$ unitaire sans perte de généralité. On note $z_1$, $z_2$, $z_3$ et $z_4$ les racines de $P$ dans $\Cc$.

Si $z_1$, $z_2$, $z_3$ et $z_4$ forment un parallélogramme, notons $a$ le centre de ce parallélogramme. Les racines de $P$ s'écrivent alors $z_1$, $z_2$, $2_a-z_1$, $2a-z_2$ et si $Q=P(X+a)$ alors 
$Q(-a+z_1)=Q(a-z_1)=Q(-a+z_2)=Q(a-z_2)=0$. Les racines du polynôme $Q$ sont deux à deux opposées, ce qui équivaut à dire que le polynôme $Q$ est bicarré ou encore de la forme $X^4+\alpha X^2+\beta$ ou enfin que

$$P=(X-a)^4+\alpha(X-a)^2+\beta.$$

Mais alors $a$ est racine de $P'=4(X-a)^3+2\alpha(X-a)$ et de $P^{(3)}=24(X-a)$.

Réciproquement, si $P'$ et $P^{(3)}$ ont une racine commune $a$. $P^{(3)}$ est de degré $1$ et de coefficient dominant $24$ et donc $P^{(3)}=24(X-a)$ puis en intégrant $P''=12(X-a)^2+\lambda$ puis $P'=4(X-a)^3+\lambda(X-a)+\mu$.
La condition $a$ est racine de $P'$ fournit $\mu=0$ et donc $P=(X-a)^4+\alpha(X-a)^2+\beta$. Donc, le polynôme $Q=P(X+a)$ est bicarré et ses racines sont deux à deux opposées et donc de la forme $Z_1=a-z_1$, $Z_2=z1-a$, $Z_3=a-z_2$, $Z_4=z_2-a$ et on a bien $Z_1-Z_3=Z_4-Z_2$.

\fincorrection
\correction{005347}
Si $(x,y,z)$ est solution du système proposé noté $(S)$, alors $x$, $y$ et $z$ sont deux à deux distincts. En effet, si par exemple $x=y$ alors $7=y^2+yz+z^2=x^2+xz+z^2=13$ ce qui est impossible. Donc,

$$(S)\Leftrightarrow\left\{
\begin{array}{l}
y^3-z^3=7(y-z)\\
z^3-x^3=13(z-x)\\
x^3-y^3=3(x-y)
\end{array}
\right..$$
  
En additionnant les trois équations, on obtient $-10x+4y+6z=0$ ou encore $-5x+2y+3z=0$. Donc,

\begin{align*}\ensuremath
(S)&\Leftrightarrow\left\{
\begin{array}{l}
y=\frac{1}{2}(5x-3z)\\
(\frac{1}{2}(5x-3z))^2+\frac{1}{2}(5x-3z)z+z^2=7\\
z^2+zx+x^2=13\\
x^2+\frac{1}{2}(5x-3z)x+(\frac{1}{2}(5x-3z))^2=3
\end{array}
\right.\Leftrightarrow\left\{
\begin{array}{l}
y=\frac{1}{2}(5x-3z)\\
25x^2-20xz+7z^2=28\\
z^2+zx+x^2=13\\
39x^2-36xz+9z^2=12
\end{array}
\right.\\
 &\Leftrightarrow\left\{
\begin{array}{l}
y=\frac{1}{2}(5x-3z)\\
xz=13-x^2-z^2\\
25x^2-20(13-x^2-z^2)+7z^2=28\\
39x^2-36(13-x^2-z^2)+9z^2=12
\end{array}
\right.
\Leftrightarrow\left\{
\begin{array}{l}
y=\frac{1}{2}(5x-3z)\\
xz=13-x^2-z^2\\
5x^2+3z^2=32
\end{array}
\right.
\end{align*}

Soit $(S')$ le système formé des deux dernières équations. On note que $x=0$ ne fournit pas de solution et donc

\begin{align*}\ensuremath
(S')&\Leftrightarrow\left\{
\begin{array}{l}
z^2=\frac{1}{3}(32-5x^2)\\
xz=13-x^2-\frac{1}{3}(32-5x^2)
\end{array}
\right.\Leftrightarrow\left\{
\begin{array}{l}
z=\frac{2x^2+7}{3x}\\
\frac{(2x^2+7)^2}{9x^2}=\frac{1}{3}(32-5x^2)
\end{array}
\right.
\end{align*}

La deuxième équation s'écrit $(2x^2+7)^2=3x^2(32-5x^2)$ puis $19x^4-68x^2+49=0$ puis $x^2=\frac{34\pm15}{19}$
D'où les solutions $x=1$ ou $x=-1$ ou $x=\sqrt{\frac{49}{19}}$ ou $x=-\sqrt{\frac{49}{19}}$. Puis, les quatre triplets solutions du système~:~$(1,-2,3)$, $(-1,2,-3)$, $(\frac{7}{\sqrt{19}},\frac{1}{\sqrt{19}},\frac{11}{\sqrt{19}})$ et $(-\frac{7}{\sqrt{19}},-\frac{1}{\sqrt{19}},-\frac{11}{\sqrt{19}})$.
\fincorrection
\correction{005348}
Soit $P=\prod_{k=0}^{n-1}(X-\omega_k)=X^n-1$ (où $\omega_k=e^{2ik\pi/n}$)
\begin{enumerate}
\item  $\prod_{k=0}^{n-1}(1+\frac{2}{2-\omega_k})=\frac{\prod_{k=0}^{n-1}(4-\omega_k)}{\prod_{k=0}^{n-1}(2-\omega_k)}
=\frac{P(4)}{P(2)}=\frac{4^n-1}{2^n-1}=2^n+1$.
\item  

\begin{align*}\ensuremath
\prod_{k=0}^{n-1}(\omega_k^2-2\omega_k\cos a+1)&=\prod_{k=0}^{n-1}(e^{ia}-\omega_k)(e^{-ia}-\omega_k)=P(e^{ia})P(e^{-ia})=(e^{ina}-1)(e^{-ina}-1)\\
 &=2-e^{ina}-e^{-ina}=2(1-\cos na).
\end{align*}
\end{enumerate}
\fincorrection
\correction{005352}
L'équation proposée admet deux solutions inverses l'une de l'autre si et seulement si il existe deux complexes $a$ et $b$ tels que

$$X^4-21X+8=(X^2+aX+1)(X^2+bX+8)=X^4+(a+b)X^3+(9+ab)X^2+(8a+b)X+8\;(*)$$

$(*)\Leftrightarrow b=-a$ et $ab=-9$ et $8a+b=-21\Leftrightarrow a=3$ et $b=-3$. Ainsi,

$$X^4-21X+8=(X^2+3X+1)(X^2-3X+8)=(X-\frac{-3+\sqrt{5}}{2})(X-\frac{-3-\sqrt{5}}{2})(X-\frac{3+i\sqrt{15}}{2})(X-\frac{3-i\sqrt{15}}{2}).$$
\fincorrection
\correction{000886}
\begin{enumerate}
  \item 

    \begin{enumerate}
    \item $(0,0,0) \in E_1$.
    \item Soient $(x,y,z)$ et $(x',y',z')$ deux \' el\' ements
de $E_1$. On a donc $3x-7y=z$ et $3x'-7y'=z'$. Donc $3(x+x')-7(y+y')=(z+z')$, d'où $(x+x',y+y',z+z')$  
appartient \`a  $E_1$.

\item Soit $\lambda \in {\R}$ et $(x,y,z)\in E_1$.  Alors la relation $3x-7y=z$ implique
  que $3 (\lambda x) -7(\lambda y)=\lambda z$
  donc que $\lambda(x,y,z)=(\lambda x,\lambda  y,\lambda  z)$ appartient \`a  $E_1$.
    \end{enumerate}


  
  \item $E_2 =\{ (x,y,z)\in {\R}^3 \mid x^2-z^2=0 \} $ c'est-\`a-dire $E_2
    =\{ (x,y,z)\in {\R}^3 \mid x=z \hbox{ ou } x=-z \} $. Donc
    $(1, 0,-1)$ et $(1, 0, 1)$ appartiennent \`a $ E_2$ mais
    $(1, 0,-1)+(1, 0, 1)=(2, 0, 0)$ n'appartient pas \`a $ E_2$ qui n'est en cons\' equence pas un
sous-espace vectoriel de ${\R}^3$.



\item $E_3$ est un sous-espace vectoriel de ${\R}^3$. En effet :
    \begin{enumerate}
    \item $(0, 0, 0) \in E_3$.
    \item Soient $(x, y, z)$ et
      $(x', y', z')$ deux \' el\'ements
      de $E_3$. On a donc $x+y-z=x+y+z=0$ et $x'+y'-z'=x'+y'+z'=0$.
      Donc $(x+x')+(y+y')-(z+z')= (x+x')+(y+y')+(z+z')=0$ et
      $(x, y, z)+(x', y', z')=(x+x',y+y',z+z')$  appartient \`a  $E_3$.
\item Soit $\lambda \in {\R}$ et $(x,y,z)\in E_3$.  Alors la relation $x+y-z=x+y+z=0$ implique
  que $\lambda x+\lambda y-\lambda z=\lambda x+\lambda y+\lambda z=0$
  donc que $\lambda (x, y, z)=(\lambda x,\lambda  y,\lambda  z)$ appartient \`a  $E_3$.
    \end{enumerate}


\item Les vecteurs $(1, 0, 0)$ et
  $(0, 0, 1)$ appartiennent \`a $ E_4$
  mais leur somme $(1, 0, 0)+(0, 0, 1)=(1, 0, 1)$  ne lui
appartient pas donc $E_4$ n'est pas un sous-espace vectoriel de
${\R}^3$.
\end{enumerate}
\fincorrection
\nocorrection
\correction{000888}
\begin{enumerate}
\item $E_1$ : non si $a \neq 0$ car alors $0 \notin E_1$ ; oui, si $a
  = 0$ car alors $E_1$ est l'intersection des sous-espaces vectoriels
  $\{(x,y,z)\in \Rr^3 \mid x+y=0 \}$ et $\{(x,y,z)\in \Rr^3 \mid x=0 \}$.
\item $E_2$ est un sous-espace vectoriel de $\mathcal{F}(\Rr,\Rr)$.
\item $E_3$ : non, car la fonction nulle n'appartient pas \`a $E_3$.
\item $E_4$ : non, en fait $E_4$ n'est m\^eme pas un sous-groupe de
  $(\Rr^2,+)$ car $(2,0)\in E_4$ mais $-(2,0)=(-2,0) \notin E_4$.
\end{enumerate}
\fincorrection
\nocorrection
\nocorrection
\nocorrection
\nocorrection
\correction{000893}
\begin{enumerate}
\item Sens $\Leftarrow$. Si $F\subset G$ alors $F\cup G=G$ donc $F\cup
  G$ est un sous-espace vectoriel.  De m\^eme si $G\subset F$.
  
  Sens $\Rightarrow$. On suppose que $F\cup G$ est un sous-espace
  vectoriel.  Par l'absurde supposons que $F$ n'est pas inclus dans
  $G$ et que $G$ n'est pas inclus dans $F$. Alors il existe $x\in
  F\setminus G$ et $y\in G\setminus F$. Mais alors $x\in F\cup G$,
  $y\in F\cup G$ donc $x+y\in F\cup G$ (car $F\cup G$ est un
  sous-espace vectoriel). Comme $x+y\in F\cup G$ alors
$x+y\in F$ ou $x+y\in G$.
\begin{itemize}
 \item Si $x+y\in F$ alors, comme $x\in F$, $(x+y)+(-x) \in F$ donc $y\in F$, ce qui est absurde.
 \item  Si $x+y\in G$ alors, comme $y\in G$, $(x+y)+(-y) \in G$ donc $x\in G$, ce qui est absurde.
\end{itemize}  
Dans les deux cas nous obtenons une contradiction. Donc $F$ est inclus
dans $G$ ou $G$ est inclus dans $F$.

\item Supposons $G\subset F$.
\begin{itemize}
 \item Inclusion $\supset$. Soit $x\in G+(F\cap H)$. Alors il existe $a\in G$, $b\in F\cap H$ tels que $x=a+b$. Comme $G\subset F$ alors $a\in F$, de plus $b\in F$ donc $x=a+b\in F$.
D'autre part $a\in G$, $b\in H$, donc $x=a+b\in G+H$. Donc $x\in F\cap(G+H)$.
 \item Inclusion $\subset$. Soit $x\in F\cap(G+H)$. $x\in G+H$ alors il existe $a\in G$,
$b\in H$ tel que $x=a+b$. Maintenant $b=x-a$ avec $x\in F$ et $a\in G\subset F$, donc $b\in F$,
donc $b\in F\cap H$. Donc $x=a+b\in G+(F\cap H)$.
\end{itemize}  
  
\end{enumerate}
\fincorrection
\nocorrection
\nocorrection
\nocorrection
\nocorrection
\nocorrection
\nocorrection
\nocorrection
\nocorrection
\nocorrection
\nocorrection
\correction{003298}
Il y a égalité.
\fincorrection
\correction{003299}
L'intersection contient $F$.

Soit $\vec u \in \bigl( F+(G\cap F')\bigr) \cap \bigl( F+(G\cap G')\bigr)$ :
$\vec u = \vec a + \vec b = \vec a{\,}' + \vec b{}'$
avec $\vec a,\vec a{\,}' \in F$, $\vec b\in G\cap F'$ et $\vec b{}'\in G\cap G'$.

Alors $\vec b-\vec b{}' = \vec a{\,}'-\vec a \in F\cap G = F'\cap G'$,
donc $\vec b \in G'$, donc $\vec b \in F'\cap G' \subset F$.
\fincorrection
\nocorrection
\nocorrection
\correction{005164}
\begin{enumerate}
\item  La fonction nulle est dans $F$ et en particulier, $F\neq\varnothing$.
Soient alors $(f,g)\in F^2$ et $(\lambda,\mu)\in\Rr^2$.

$$(\lambda f+\mu g)(0)+(\lambda f+\mu g)(1)=\lambda(f(0)+f(1))+\mu(g(0)+g(1))=0.$$

Par suite, $\lambda f+\mu g$ est dans $F$. On a montré que~:

$$F\neq\varnothing\;\mbox{et}\;\forall(f,g)\in F^2,\;\forall(\lambda,\mu)\in\Rr^2,\;\lambda f+\mu g\in F.$$

$F$ est donc un sous-espace vectoriel de $E$.

\item  Même démarche et même conclusion .

\item  $F$ ne contient pas la fonction nulle et n'est donc pas un sous-espace vectoriel de $E$.

\item  La fonction nulle est dans $F$ et en particulier, $F\neq\varnothing$.
Soient alors $(f,g)\in F^2$ et $(\lambda,\mu)\in\Rr^2$.

Pour $x$ élément de $[0,1]$,
$$(\lambda f+\mu g)(x)+(\lambda f+\mu g)(1-x)= \lambda(f(x)+f(1-x))+\mu(g(x)+g(1-x))=0$$ et $\lambda f+\mu g$ est dans
$F$. $F$ est un sous-espace vectoriel de $E$.

\textbf{Remarque.} Les graphes des fonctions considérés sont symétriques par rapport au point $(\frac{1}{2},0)$.

\item  $F$ contient la fonction constante $1$ mais pas son opposé la fonction constante $-1$ et n'est donc pas un
sous-espace vectoriel de $E$.

\item  $F$ ne contient pas la fonction nulle et n'est donc pas un sous-espace vectoriel de $E$.
\end{enumerate}
\fincorrection
\correction{005165}
Dans les cas où $F$ est un sous-espace, on a à chaque fois trois démarches possibles pour le
vérifier~:

\begin{itemize}
\item[-] Utiliser la caractérisation d'un sous-espace vectoriel.
\item[-] Obtenir $F$ comme noyau d'une forme linéaire ou plus généralement, comme noyau d'une application linéaire.
\item[-] Obtenir $F$ comme sous-espace engendré par une famille de vecteurs.
\end{itemize}

Je vous détaille une seule fois les trois démarches.

\begin{enumerate}
\item 
\begin{itemize}
\item[\textbf{1ère démarche.}] $F$ contient le vecteur nul $(0,...,0)$ et donc $F\neq\varnothing$.
Soient alors $((x_1,...,x_n),(x_1' ,...,x_n'))\in F^2$ et $(\lambda,\mu)\in\Rr^2$. On a

$$\lambda(x_1,...,x_n)+\mu(x_1',...,x_n')=(\lambda x_1+\mu x_1',...,\lambda x_n+\mu x_n')$$

avec $\lambda x_1+\mu x_1'=0$. Donc, $\lambda(x_1,...,x_n)+\mu(x_1',...,x_n')\in F$. $F$ est un sous-espace
vectoriel de $\Rr^n$.
\item[\textbf{2ème démarche.}] L'application $(x_1,...,x_n)\mapsto x_1$ est une forme linéaire sur $\Rr^n$ et $F$ en
est lenoyau. $F$ est donc un sous-espace vectoriel de $\Rr^n$.
\item[\textbf{3ème démarche.}]

\begin{align*}
F&=\{(0,x_2,...,x_n),\;(x_2,...,x_n)\in\Rr^{n-1}\}
=\{x_2(0,1,0,...,0)+...+x_n(0,...,0,1),\;(x_2,...,x_n)\in\Rr^{n-1}\}\\
 &=\mbox{Vect}((0,1,0,...,0),...,(0,...,0,1)).
\end{align*}

$F$ est donc un sous-espace vectoriel de $\Rr^n$.

\end{itemize}

\item  $F$ ne contient pas le vecteur nul et n'est donc pas un sous-espace vectoriel de $\Rr^n$.

\item  (Ici, $n\geq2$). L'application $(x_1,...,x_n)\mapsto x_1-x_2$ est une forme linéaire sur $\Rr^n$ et $F$ en
est le noyau. $F$ est donc un sous-espace vectoriel de $\Rr^n$.

\item  L'application $(x_1,...,x_n)\mapsto x_1+...+x_n$ est une forme linéaire sur $\Rr^n$ et $F$ en est le
noyau. $F$ est donc un sous-espace vectoriel de $\Rr^n$.

\item  (Ici, $n\geq2$). Les vecteurs $e_1=(1,0,...,0)$ et $e_2=(0,1,0...,0)$ sont dans $F$ mais $e_1+e_2=
(1,1,0...0)$ n'y est pas. $F$ n'est donc pas un sous espace vectoriel de $E$.

\textbf{Remarque.} $F$ est la réunion des sous-espaces $\{(x_1,...,x_n)\in\Rr^n/\;x_1=0\}$ et
$\{(x_1,...,x_n)\in\Rr^n/\;x_2=0\}$.
\end{enumerate}
\fincorrection
\correction{005166}
Il suffit de montrer que $C\subset B$.

Soit $x$ un élément de $C$. Alors $x\in A+C=A+B$ et il existe $(y,z)\in A\times B$ tel que $x=y+z$. Mais $z\in
B\subset C$ et donc, puisque $C$ est un sous-espace vectoriel de $E$, $y=x-z$ est dans $C$. Donc, $y\in A\cap C=A\cap B$
et en particulier $y$ est dans $B$. Finalement, $x=y+z$ est dans $B$. On a montré que tout élément de $C$ est dans $B$
et donc que, $C\subset B$. Puisque d'autre part $B\subset C$, on a $B=C$.
\fincorrection
\correction{005168}
Soit $(\lambda,\mu)\in\Rr^2$.

\begin{align*}
(\lambda,\mu,-37,-3)\in F&\Leftrightarrow\exists(a,b)\in\Rr^2/\;(\lambda,\mu,-37,-3)=au+bv
\Leftrightarrow\exists(a,b)\in\Rr^2/\;
\left\{
\begin{array}{l}
a+2b=\lambda\\
2a-b=\mu\\
-5a+4b=-37\\
3a+7b=-3
\end{array}
\right.\\
 &\Leftrightarrow\exists(a,b)\in\Rr^2/\;
\left\{
\begin{array}{l}
a+2b=\lambda\\
2a-b=\mu\\
a=\frac{247}{47}\\
b=-\frac{126}{47}
\end{array}
\right.
\Leftrightarrow
\left\{
\begin{array}{l}
\lambda=\frac{247}{47}+2(-\frac{126}{47})\\
\mu=2\frac{247}{47}+\frac{126}{47}
\end{array}
\right.\Leftrightarrow\left\{
\begin{array}{l}
\lambda=-\frac{5}{47}\\
\mu=\frac{620}{47}
\end{array}
\right..
\end{align*}
\fincorrection
\correction{005169}
Posons $F=\mbox{Vect}(a,b)$ et $G=\mbox{Vect}(c,d)$.

Montrons que $c$ et $d$ sont dans $F$.

\begin{align*}
c\in F\Leftrightarrow\exists(\lambda,\mu)\in\Rr^2/\;c=\lambda a+\mu b
\Leftrightarrow\exists(\lambda,\mu)\in\Rr^2/\;
\left\{
\begin{array}{l}
\lambda+2\mu=1\\
2\lambda-\mu=0\\
3\lambda+\mu=1
\end{array}
\right.
\Leftrightarrow\exists(\lambda,\mu)\in\Rr^2/\;
\left\{
\begin{array}{l}
\lambda=\frac{1}{5}\\
\mu=\frac{2}{5}\\
3\lambda+\mu=1
\end{array}
\right..
\end{align*}

Puisque $3.\frac{1}{5}+\frac{2}{5}=1$, le système précédent admet bien un couple $(\lambda,\mu)$ solution
et $c$ est dans $F$. Plus précisément, $c=\frac{1}{5}a+\frac{2}{5}b$.

\begin{align*}
d\in F\Leftrightarrow\exists(\lambda,\mu)\in\Rr^2/\;d=\lambda a+\mu b
\Leftrightarrow\exists(\lambda,\mu)\in\Rr^2/\;
\left\{
\begin{array}{l}
\lambda+2\mu=0\\
2\lambda-\mu=1\\
3\lambda+\mu=1
\end{array}
\right.
\Leftrightarrow\exists(\lambda,\mu)\in\Rr^2/\;
\left\{
\begin{array}{l}
\lambda=\frac{2}{5}\\
\mu=-\frac{1}{5}\\
3\lambda+\mu=1
\end{array}
\right..
\end{align*}

Puisque $3.\frac{2}{5}-\frac{1}{5}=1$, le système précédent admet bien un couple $(\lambda,\mu)$
solution et $d$ est dans $F$. Plus précisément, $d=\frac{2}{5}a-\frac{1}{5}b$. En résumé,
$\{c,d\}\subset F$ et donc $G=\mbox{Vect}(c,d)\subset F$.

Montrons que $a$ et $b$ sont dans $G$ mais les égalités $c=\frac{1}{5}a+\frac{2}{5}b$ et
$d=\frac{2}{5}a-\frac{1}{5}b$ fournissent $a=c+2d$ et $b=2c-d$. Par suite, $\{a,b\}\subset G$ et donc
$F=\mbox{Vect}(a,b)\subset G$. Finalement $F=G$.
\fincorrection
\correction{005172}
\begin{enumerate}
\item  Soit $x\in E$.

$x\in(A\cap B)+(A\cap C)\Rightarrow\exists y\in A\cap B,\;\exists z\in A\cap C/\;x=y+z$.

$y$ et $z$ sont dans $A$ et donc $x=y+z$ est dans $A$ car $A$ est un sous-espace vectoriel de $E$.

Puis $y$ est dans $B$ et $z$ est dans $C$ et donc $x=y+z$ est dans $B+C$.
Finalement,

$$\forall x\in E,\;[x\in(A\cap B)+(A\cap C)\Rightarrow x\in A\cap(B+C)].$$

Autre démarche.

$(A\cap B\subset B$ et $A\cap C\subset C)\Rightarrow(A\cap B)+(A\cap C)\subset B+C$ puis
$(A\cap B\subset A$ et $A\cap C\subset A\Rightarrow (A\cap B)+(A\cap C)\subset A+A=A$, et finalement
$(A\cap B)+(A\cap C)\subset A\cap(B+C)$.

\item  Si on essaie de démontrer l'inclusion contraire, le raisonnement coince car la somme $y+z$ peut être dans
$A$ sans que ni $y$, ni $z$ ne soient dans $A$.

Contre-exemple. Dans $\Rr^2$, on considère $A=\Rr.(1,0)=\{(x,0),\;x\in\Rr\}$, $B=\Rr.(0,1)$ et $C=\Rr.(1,1)$.

$B+C=\Rr^2$ et $A\cap(B+C)=A$ mais $A\cap B=\{0\}$ et $A\cap C=\{0\}$ et donc $(A\cap B)+(A\cap C)=\{0\}\neq
A\cap(B+C)$.

\item  $A\cap B\subset B\Rightarrow (A\cap B)+(A\cap C)\subset B+(A\cap C)$ mais aussi $(A\cap B)+(A\cap C)\subset
A+A=A$. Donc, $(A\cap B)+(A\cap C)\subset A\cap(B+(A\cap C))$.

Inversement, soit $x\in A\cap(B+(A\cap C))$ alors $x=y+z$ où $y$ est dans $B$ et $z$ est dans $A\cap C$. Mais alors,
$x$ et $z$ sont dans $A$ et donc $y=x-z$ est dans $A$ et même plus précisément dans $A\cap B$. Donc,
$x\in(A\cap B)+(A\cap C)$. Donc,  $A\cap(B+(A\cap C))\subset(A\cap B)+(A\cap C)$ et finalement, $A\cap(B+(A\cap
C))=(A\cap B)+(A\cap C)$.
\end{enumerate}
\fincorrection
\correction{005173}
\begin{enumerate}
\item  Pour $(x,y,z,t)\in\Rr^4$, on pose $f((x,y,z,t))=x-2y$, $g((x,y,z,t))=y-2z$ et $h((x,y,z,t))=x-y+z-t$. $f$,
$g$ et $h$ sont des formes linéaires sur $\Rr^4$. Donc, $V=\mbox{Ker}f\cap\mbox{Ker}g$ est un sous-espace vectoriel de
$\Rr^4$ en tant qu'intersection de sous-espaces vectoriels de $\Rr^4$ et $W=\mbox{Ker}h$ est un sous-espace vectoriel
de $\Rr^4$.

\item  Soit $(x,y,z,t)\in\Rr^4$.

\begin{align*}
(x,y,z,t)\in V\Leftrightarrow
\left\{
\begin{array}{l}
x=2y\\
y=2z
\end{array}
\right.\Leftrightarrow
\Leftrightarrow
\left\{
\begin{array}{l}
x=4z\\
y=2z
\end{array}
\right.
.
\end{align*}

Donc, $V=\{(4z,2z,z,t),\;(z,t)\in\Rr^2\}=\mbox{Vect}(e_1,e_2)$ où $e_1=(4,2,1,0)$ et $e_2=(0,0,0,1)$. Montrons alors
que $(e_1,e_2)$ est libre. Soit $(z,t)\in\Rr^2$.

$$ze_1+te_2=0\Rightarrow(4z,2z,z,t)=(0,0,0,0)\Rightarrow z=t=0.$$

Donc, $(e_1,e_2)$ est une base de $V$.

Pour $(x,y,z,t)\in\Rr^4$, $(x,y,z,t)\in W\Leftrightarrow t=x-y+z$. Donc,
$W=\{(x,y,z,x-y+z),\;(x,y,z)\in\Rr^3\}=\mbox{Vect}(e_1',e_2',e_3')$ où $e_1'=(1,0,0,1)$, $e_2'=(0,1,0,-1)$ et
$e_3'=(0,0,1,1)$.

Montrons alors que $(e_1',e_2',e_3')$ est libre. Soit $(x,y,z)\in\Rr^3$.

$$xe_1'+ye_2'+ze_3'=0\Rightarrow(x,y,z,x-y+z)=(0,0,0,0)\Rightarrow x=y=z=0.$$

Donc, $(e_1',e_2',e_3')$ est une base de $W$.
Soit $(x,y,z,t)\in\Rr^4$.

$$(x,y,z,t)\in V\cap W\Leftrightarrow
\left\{
\begin{array}{l}
x=2y\\
y=2z\\
x-y+z-t=0
\end{array}
\right.\Leftrightarrow
\left\{
\begin{array}{l}
x=4z\\
y=2z\\
t=3z
\end{array}
\right..$$

Donc, $V\cap W=\{(4z,2z,z,3z),\;z\in\Rr\}=\mbox{Vect}(e)$ où $e=(4,2,1,3)$. De plus, $e$ étant non nul, la famille
$(e)$ est libre et est donc une base de $V\cap W$.

\item  Soit $u=(x,y,z,t)$ un vecteur de $\Rr^4$.

On cherche $v=(4\alpha,2\alpha,\alpha,\beta)\in V$ et $w=(a,b,c,a-b+c)\in W$ tels que $u=v+w$.

$$u=v+w\Leftrightarrow\left\{
\begin{array}{l}
4\alpha+a=x\\
2\alpha+b=y\\
\alpha+c=z\\
\beta+a-b+c=t
\end{array}
\right.\Leftrightarrow
\left\{
\begin{array}{l}
a=x-4\alpha\\
b=y-2\alpha\\
c=z-\alpha\\
\beta=-x+y-z+t-3\alpha
\end{array}
\right..$$

et $\alpha=0$, $\beta=-x+y-z+t$, $a=x$, $b=y$ et $c=z$ conviennent. Donc, $\forall u\in\Rr^4,\;\exists(v,w)\in
V\times W/\;u=v+w$. On a montré que $\Rr^4=V+W$.
\end{enumerate}
\fincorrection
\correction{005174}
\begin{enumerate}
\item  Pour tout $(y,y')$ élément de $[0,2\pi[^2$, $f((0,y))=f((0,y'))$ et f n'est pas injective.

Montrons que $f$ est surjective.

Soit $(X,Y)\in\Rr^2$.
\begin{itemize}
\item[-] Si $X=Y=0$, $f((0,0))=(0,0)$.
\item[-] Si $X=0$ et $Y>0$, $f((Y,\frac{\pi}{2}))=(0,Y)$ avec $(Y,\frac{\pi}{2})$ élément de
$[0,+\infty[\times[0,2\pi[$.
\item[-] Si $X=0$ et $Y<0$, $f((-Y,\frac{3\pi}{2}))=(0,Y)$ avec $(-Y,\frac{3\pi}{2})$ élément
de $[0,+\infty[\times[0,2\pi[$.
\item[-] Si $X>0$ et $Y\geq0$, $f((\sqrt{X^2+Y^2},\Arctan\frac{Y}{X}))=(X,Y)$ avec
$(\sqrt{X^2+Y^2},\Arctan\frac{Y}{X})$ élément de

$[0,+\infty[\times[0,2\pi[$.
\item[-] Si $X<0$ et $Y\geq0$, $f((\sqrt{X^2+Y^2},\pi+\Arctan\frac{Y}{X}))=(X,Y)$ avec
$(\sqrt{X^2+Y^2},\pi+\Arctan\frac{Y}{X})$ élément de $[0,+\infty[\times[0,2\pi[$.
\item[-] Si $X>0$ et $Y<0$, $f((\sqrt{X^2+Y^2},2\pi+\Arctan\frac{Y}{X}))=(X,Y)$ avec $(\sqrt{X^2+Y^2},
2\pi+\Arctan\frac{Y}{X})$ élément de $[0,+\infty[\times[0,2\pi[$.
\item[-] Si $X<0$ et $Y<0$, $f((\sqrt{X^2+Y^2},\pi+\Arctan\frac{Y}{X}))=(X,Y)$ avec
$(\sqrt{X^2+Y^2},\pi+\Arctan\frac{Y}{X})$ élément de $[0,+\infty[\times[0,2\pi[$.
\end{itemize}

\item  Pour tout réel $x$, on a $a\cos(x-\alpha)+b\cos(x-\beta)=(a\cos\alpha+b\cos\beta)\cos
x+(a\sin\alpha+b\sin\beta)\sin x$.

D'après 1), $f$ est surjective et il existe $(c,\gamma)$ élément de $[0,+\infty[\times[0,2\pi[$ tel que
$a\cos\alpha+b\cos\beta=c\cos\gamma$ et $a\sin\alpha+b\sin\beta=c\sin\gamma$. Donc,

$$\exists(c,\gamma)\in[0,+\infty[\times[0,2\pi[/\;\forall
x\in\Rr,\;a\cos(x-\alpha)+b\cos(x-\beta)=c(\cos xcos\gamma+\sin x\sin\gamma)=c\cos(x-\gamma).$$
\item  $F$ est non vide car contient l'application nulle et est contenu dans $E$. De plus, pour $x$ réel,

\begin{align*}
a\cos(x-\alpha)+b\cos(2x-\beta)&+a'\cos(x-\alpha')+b'\cos(2x-\beta')\\
 &=a\cos(x-\alpha)+a'\cos(x-\alpha')+b\cos(2x-\beta)
+b'\cos(2x-\beta')\\
 &=a''cos(x-\alpha'')+b''\cos(2x-\beta''),
\end{align*}

pour un certain $(a',b'',\alpha'',\beta'')$ (d'après 2)). $F$ est un sous-espace vectoriel de $E$.

\item 

Pour tout réel $x$, $\cos x=1.\cos(x-0)+0.\cos(2x-0)$ et $x\mapsto\cos x$ est élément de $F$.

Pour tout réel $x$, $\sin x=1.\cos(x-\frac{\pi}{2})+0.\cos(2x-0)$ et $x\mapsto\sin x$ est élément de $F$.

Pour tout réel $x$, $\cos(2x)=0.\cos(x-0)+1.\cos(2x-0)$ et $x\mapsto\cos(2x)$ est élément de $F$.

Pour tout réel $x$, $\sin(2x)=0.\cos(x-0)+1.\cos(2x-\frac{\pi}{2})$ et $x\mapsto\sin(2x)$ est élément de $F$.

D'autre part, pour tout réel $x$, $\cos(2x)=2\cos^2x-1=1-2\sin^2x$ et donc,

$$x\mapsto1\in F\Leftrightarrow x\mapsto\cos^2x\in F\Leftrightarrow x\mapsto\sin^2x\in F.$$

Montrons alors que $1\notin F$.

On suppose qu'il existe $(a,b,\alpha,\beta)\in\Rr^4$ tel que

$$\forall x\in\Rr,\;a\cos(x-\alpha)+b\cos(2x-\beta)=1.$$

En dérivant deux fois, on obtient~:

$$\forall x\in\Rr,\;-a\cos(x-\alpha)-4b\cos(2x-\beta)=0,$$

et donc en additionnant

$$\forall x\in\Rr,\;-3b\cos(2x-\beta)=1,$$

ce qui est impossible (pour $x=\frac{\pi}{4}+\frac{\beta}{2}$, on trouve $0$). Donc, aucune des trois
dernières fonctions n'est dans $F$.

\item  On a vu que $(x\mapsto\cos x,\;x\mapsto\sin x,\;x\mapsto\cos(2x),\;x\mapsto\sin(2x))$ est une famille
d'éléments de $F$. Montrons que cette famille est libre.

Soit $(a,b,c,d)\in\Rr^4$.

Supposons que $\forall x\in\Rr,\;a\cos x+b\sin x+c\cos(2x)+d\sin(2x)=0$. En dérivant deux fois, on obtient $\forall
x\in\Rr,\;-a\cos x-b\sin x-4c\cos(2x)-4d\sin(2x)=0$ et en additionnant~:~$\forall
x\in\Rr,\;-3c\cos(2x)-3d\sin(2x)=0$. Donc,

$$\forall x\in\Rr,\;\left\{
\begin{array}{l}
a\cos x+b\sin x=0\\
c\cos(2x)+d\sin(2x)=0
\end{array}
\right..$$

$x=0$ fournit $a=c=0$ puis $x=\frac{\pi}{4}$ fournit $b=d=0$. Donc, $(x\mapsto\cos x,\;x\mapsto\sin
x,\;x\mapsto\cos(2x),\;x\mapsto\sin(2x))$ est une famille libre d'éléments de $F$.
\end{enumerate}
\fincorrection
\correction{005175}
\begin{enumerate}
\item  $C$ contient l'identité de $\Rr$, mais ne contient pas son opposé. Donc, $C$ n'est pas un espace vectoriel.

\item  Montrons que $V$ est un sous-espace vectoriel de l'espace vectoriel des applications de $\Rr$ dans $\Rr$.
$V$ est déjà non vide car contient la fonction nulle $(0=0-0)$.

Soit $(f_1,f_2)\in V^2$. Il existe $(g_1,g_2,h_1,h_2)\in C^4$ tel que $f_1=g_1-h_1$ et $f_2=g_2-h_2$. Mais alors,
$f_1+f_2=(g_1+g_2)-(h_1+h_2)$. Or, une somme de fonctions croissantes sur $\Rr$ est croissante sur $\Rr$, et 
donc, $g_1+g_2$ et $h_1+h_2$ sont des éléments de $C$ ou encore $f_1+f_2$ est dans $V$.

Soit $f\in V$ et $\lambda\in\Rr$. Il existe $(g,h)\in V^2$ tel que $f=g-h$ et donc $\lambda f=\lambda g-\lambda h$.

Si $\lambda\geq 0$, $\lambda g$ et $\lambda h$ sont croissantes sur $\Rr$ et $\lambda f$ est dans $V$.

Si $\lambda<0$, on écrit $\lambda f=(-\lambda h)-(-\lambda g)$, et puisque $-\lambda g$ et $-\lambda h$ sont 
croissantes sur $\Rr$, $\lambda f$ est encore dans $V$. $V$ est donc un sous-espace vectoriel de l'espace vectoriel
des applications de $\Rr$ dans $\Rr$.
\end{enumerate}
\fincorrection
\correction{005176}
Soit $(x,y)\in E^2$.

$(1+1).(x+y)=1.(x+y)+1.(x+y)=(x+y)+(x+y)=x+y+x+y$ mais aussi
$(1+1).(x+y)=(1+1).x+(1+1).y=x+x+y+y$.

Enfin, $(E,+)$ étant un groupe, tout élément est régulier et en particulier $x$ est régulier à gauche et $y$ est
régulier à droite. On a montré que pour tout couple $(x,y)$ élément de $E^2$, $x+y=y+x$.
\fincorrection
\correction{005177}
Soit $F=(A\cap B)+(A\cap C)+(B\cap C)$.

$F\subset A+A+B=A+B$ puis $F\subset A+C+C=A+C$ puis $F\subset B+C+C=B+C$ et finalement
$F\subset(A+B)\cap(A+C)\cap(B+C)$.
\fincorrection
\correction{005563}
$\Leftarrow)$ Si $F\subset G$ ou $G\subset F$ alors $F\cup G=G$ ou $F\cup G=F$. Dans tous les cas, $F\cup G$ est un sous-espace vectoriel.

$\Rightarrow)$ Supposons que $F\not\subset G$ et que $F\cup G$ est un sous-espace vectoriel de $E$ et montrons que $G\subset F$.

$F$ n'est pas inclus dans $G$ et donc il existe $x$ élément de $E$ qui est dans $F$ et pas dans $G$.

Soit $y$ un élément de $G$. $x+y$ est dans $F\cup G$ car $x$ et $y$ y sont et car $F\cup G$ est un sous-espace vectoriel de $E$. Si $x+y$ est élément de $G$ alors $x=(x+y)-y$ l'est aussi ce qui est exclu. Donc $x+y$ est élément de $F$ et par suite $y=(x+y)-x$ est encore dans $F$. Ainsi, tout élément de $G$ est dans $F$ et donc $G\subset F$.
\fincorrection
\correction{005564}
$\Leftarrow)$ Immédiat .

$\Rightarrow)$ On raisonne par récurrence sur $n$.

Pour $n=2$, c'est l'exercice \ref{ex:rou1}.

Soit $n\geqslant 2$. Supposons que toute réunion de $n$ sous-espaces de $E$ est un sous-espace de $E$ si et seulement si l'un de ces sous-espaces contient tous les autres.

Soient $F_1$,..., $F_n$, $F_{n+1}$ $n+1$ sous-espaces vectoriels de $E$ tels que $F_1\cup...\cup F_{n+1}$ soit un sous-espace vectoriel de $E$. Posons $F=F_1\cup...\cup F_n$.

\textbullet~Si $F_{n+1}$ contient $F$, c'est fini.

\textbullet~Si $F_{n+1}\subset F$, alors $=F_1\cup...\cup F_n=F_1\cup...\cup F_n\cup F_{n+1}$ est un sous-espace vectoriel de $E$. Par hypothèse de récurrence, $F$ est l'un des $F_i$ pour un certain $i$ élément de $\llbracket1,n\rrbracket$. $F_i=F$ contient également $F_{n+1}$ et contient donc tous les $F_j$ pour $j$ élément de $\llbracket1,n+1\rrbracket$.

\textbullet~Supposons dorénavant que $F\not\subset F_{n+1}$ et que $F_{n+1}\not\subset F$ et montrons que cette situation est impossible.

Il existe un vecteur $x$ qui est dans $F_{n+1}$ et pas dans $F$ et un vecteur $y$ qui est dans $F$ et pas dans $F_{n+1}$.

Soit $\lambda$ un élément de $\Kk$. $y-\lambda x$ est un élément de $F\cup F_{n+1}$ (puisque $F\cup F_{n+1}$ est un sous-espace) mais $y-\lambda x$ n'est pas dans $F_{n+1}$ car alors $y=(y -\lambda x)+\lambda x$ y serait ce qui n'est pas.

Donc $\forall\lambda\in\Kk$, $y-\lambda x\in F$. On en déduit que pour tout scalaire $\lambda$, il existe un indice $i(\lambda)$ élément de $\llbracket1,n\rrbracket$ tel que $y-\lambda x\in F_{i(\lambda)}$. Remarquons enfin que si $\lambda\neq\mu$ alors $i(\lambda)\neq i(\mu)$. En effet, si pour $\lambda$ et $\mu$ deux scalaires distincts donnés, il existe un indice $i$ élément de $\llbracket1,n\rrbracket$ tel que $y-\lambda x$ et $y-\mu x$ soient dans $F_i$, alors $x=\frac{(y-\mu x)-(y-\lambda x)}{\mu-\lambda}$  est encore dans $F_i$ et donc dans $F$, ce qui n'est pas. 

Comme l'ensemble des scalaires est infini et que l'ensemble des indices ne l'est pas, on vient de montrer que cette dernière situation n'est pas possible, ce qui achève la démonstration.
\fincorrection
\correction{006868}
Pour qu'un ensemble $E$, muni d'une addition $x+y \in E$ (pour tout $x,y \in E$)
et d'une multiplication par un scalaire $\lambda \cdot x \in E$ (pour tout $\lambda \in K$, $x\in E$),
soit un $K$-espace vectoriel il faut qu'il vérifie les huit points suivants.

\begin{enumerate}
  \item $x+(y+z)=(x+y)+z$ (pour tout $x,y,z \in E$)
  \item il existe un vecteur nul $0 \in E$ tel que $x+0=x$ (pour tout $x\in E$)
  \item il existe un opposé $-x$ tel que $x+(-x)=0$ (pour tout $x\in E$)
  \item $x+y=y+x$ (pour tout $x,y \in E$) \\
Ces quatre premières propriétés font de $(E,+)$ un groupe abélien.

  \item $1\cdot x = x$ (pour tout $x\in E$)
  \item $\lambda \cdot (x+y) = \lambda\cdot x + \lambda \cdot y$ (pour tout $\lambda \in K=$, pour tout $x,y \in E$)
  \item $(\lambda+\mu) \cdot x = \lambda\cdot x+ \mu \cdot x$ (pour tout $\lambda, \mu \in K$, pour tout $x \in E$)
  \item $(\lambda\times\mu) \cdot x = \lambda\cdot (\mu\cdot x)$ (pour tout $\lambda, \mu \in K$, pour tout $x \in E$)
\end{enumerate}


Il faut donc vérifier ces huit points pour chacun des ensembles (ici $K=\Rr$).

Commençons par $E_1$.
\begin{enumerate}
  \item $f+(g+h)=(f+g)+h$ ; en effet on bien pour tout $t\in[0,1]$ : $f(t)+\big(g(t)+h(t)\big)=\big(f(t)+g(t)\big)+h(t)$
d'où l'égalité des fonctions $f+(g+h)$ et $(f+g)+h$. Ceci est vrai pour tout $f,g,h \in E_1$.
  \item le vecteur nul est ici la fonction constante égale à $0$, que l'on note encore $0$, on a bien $f+0=f$ 
(c'est-à-dire pour tout $x\in[0,1]$, $(f+0)(t)=f(t)$, ceci pour toute fonction $f$).
  \item il existe un opposé $-f$ définie par $-f(t) = - \big(f(t)\big)$ tel que $f+(-f)=0$ 
  \item $f+g=g+f$ (car $f(t)+g(t)=g(t)+f(t)$ pour tout $t\in[0,1]$).
  \item $1\cdot f = f$ ; en effet pour tout $t\in[0,1]$, $(1\cdot f)(t) = 1\times f(t) = f(t)$.
Et une fois que l'on compris que $\lambda\cdot f$ vérifie par définition 
$(\lambda\cdot f)(t) = \lambda\times f(t)$ les autres points se vérifient sans peine.
  \item $\lambda \cdot (f+g) = \lambda\cdot f + \lambda \cdot g$
  \item $(\lambda+\mu) \cdot f = \lambda\cdot f+ \mu \cdot f$
  \item $(\lambda\times\mu) \cdot f = \lambda\cdot (\mu\cdot f)$ ; en effet pour tout $t\in [0,1]$,
 $(\lambda\times\mu)  f (t) = \lambda (\mu f(t))$
\end{enumerate}

\bigskip

Voici les huit points à vérifier pour $E_2$ en notant $u$ la suite $(u_n)_{n\in\Nn}$

\begin{enumerate}
  \item $u+(v+w)=(u+v)+w$ 
  \item le vecteur nul est la suite dont tous les termes sont nuls.
  \item La suite $-u$ est définie par $(-u_n)_{n\in\Nn}$
  \item $u+v=v+u$

  \item $1\cdot u = u$
  \item $\lambda \cdot (u+v) = \lambda\cdot u + \lambda \cdot v$ : montrons celui-ci en détails
par définition $u+v$ est la suite $(u_n+v_n)_{n\in\Nn}$ et par définition de la multiplication par un scalaire
$\lambda \cdot (u+v)$ est la suite $\big(\lambda\times (u_n+v_n)\big)_{n\in\Nn}$ qui est bien la suite
$\big(\lambda u_n+ \lambda v_n)\big)_{n\in\Nn}$ qui est exactement la suite $\lambda\cdot u + \lambda\cdot v$.
  \item $(\lambda+\mu) \cdot u = \lambda\cdot u+ \mu \cdot v$ 
  \item $(\lambda\times\mu) \cdot u = \lambda\cdot (\mu\cdot u)$ 
\end{enumerate}

\bigskip

Voici ce qu'il faut vérifier pour $E_3$, après avoir remarqué que la somme de deux polynômes de degré 
$\le n$ est encore un polynôme de degré $\le n$
(même chose pour $\lambda\cdot P$), on vérifie :

\begin{enumerate}
  \item $P+(Q+R)=(P+Q)+R$ 
  \item il existe un vecteur nul $0 \in E_3$ : c'est le polynôme nul
  \item il existe un opposé $-P$ tel que $P+(-P)=0$ 
  \item $P+Q=Q+P$ 
  \item $1\cdot P = P$ 
  \item $\lambda \cdot (P+Q) = \lambda\cdot P + \lambda \cdot Q$ 
  \item $(\lambda+\mu) \cdot P = \lambda\cdot P+ \mu \cdot P$ 
  \item $(\lambda\times\mu) \cdot P = \lambda\cdot (\mu\cdot P)$ 
\end{enumerate}


\fincorrection
\correction{006869}
\begin{enumerate}
  \item L'espace vectoriel $\mathbb{R}$ a deux sous-espaces : celui formé du vecteur nul $\{0\}$  et $\Rr$ lui-m\^eme.
\\
L'espace vectoriel $\mathbb{R}^2$ a trois types de sous-espaces: $\{0\}$, 
une infinit\'e de sous-espaces de dimension $1$ (ce sont les droites vectorielles) et $\Rr^2$ lui-m\^eme.
\\
Enfin, l'espace $\mathbb{R}^3$ a quatre types de sous-espaces: le vecteur nul, les droites vectorielles, 
les plans vectoriels et lui-m\^eme.

  \item On consid\`ere deux droites vectorielles de $\mathbb{R}^3$ dont des vecteurs directeurs 
$u$ et $v$ ne sont pas colin\'eaires alors le vecteur $u+v$ n'appartient 
\`a aucune de ces deux droites, l'union de celles-ci n'est pas un espace vectoriel.
\end{enumerate}
\fincorrection
\nocorrection
\nocorrection
\nocorrection
\correction{000900}
\begin{enumerate}
\item 
\begin{align*}
& (x,1,y,1) \in \text{Vect}\{v_1,v_2\} \\ 
\iff& \exists \lambda,\mu\in \Rr \qquad (x,1,y,1) = \lambda(1,2,3,4)+\mu(1,-2,3,-4) \\
\iff& \exists \lambda,\mu\in \Rr \qquad (x,1,y,1) = (\lambda,2\lambda,3\lambda,4\lambda)+(\mu,-2\mu,3\mu,-4\mu) \\
\iff& \exists \lambda,\mu\in \Rr \qquad (x,1,y,1) = (\lambda+\mu,2\lambda-2\mu,3\lambda+3\mu,4\lambda-4\mu) \\
\implies& \exists \lambda,\mu\in \Rr \qquad 1 = 2(\lambda-\mu) \text{ et } 1=4(\lambda-\mu) \\
\implies& \exists \lambda,\mu\in \Rr \qquad \lambda-\mu= \frac 12 \text{ et } \lambda-\mu=\frac14 \\
\end{align*}
Ce qui est impossible (quelque soient $x,y$). Donc on ne peut pas trouver de tels $x,y$.

\item On fait le m\^eme raisonnement : 
\begin{align*}
& (x,1,1,y) \in \text{Vect}\{v_1,v_2\} \\ 
iff&\ \exists \lambda,\mu\in \Rr \qquad (x,1,1,y) = (\lambda+\mu,2\lambda-2\mu,3\lambda+3\mu,4\lambda-4\mu) \\
\iff& \exists \lambda,\mu\in \Rr \qquad 
\begin{cases}
  x &= \lambda + \mu \\
  1  &= 2\lambda -2\mu \\
  1 &= 3\lambda +3\mu \\
  y &= 4\lambda -4 \mu \\
 \end{cases} \\
\iff& \exists \lambda,\mu\in \Rr \qquad 
\begin{cases}
  \lambda &= \frac 5{12} \\
  \mu &=  -\frac1{12} \\
  x &= \frac 13 \\
  y &=  2 \\  
\end{cases}.\\
\end{align*}  
Donc le seul vecteur $(x,1,1,y)$ qui convienne est $(\frac13,1,1,2)$.
\end{enumerate}
\fincorrection
\correction{000901}
\begin{enumerate}
\item On v\'erifie les propri\'et\'es qui font de $E$ un sous-espace
  vectoriel de $\Rr^4$ : 
  \begin{enumerate}
    \item l'origine $(0,0,0,0)$ est dans $E$,
    \item si $v=(x_1,x_2,x_3,x_4) \in E$ et $v'=(x_1',x_2',x_3',x_4')\in E$ alors $v+v'=(x_1+x_1',x_2+x_2',x_3+x_3',x_4+x_4')$
a des coordonnées qui vérifient l'équation et donc $v+v' \in E$.
    \item si $v=(x_1,x_2,x_3,x_4) \in E$ et $\lambda \in \Rr$ alors les coordonnées de 
$\lambda\cdot v = (\lambda x_1,\lambda x_2,\lambda x_3,\lambda x_4)$ vérifient 
l'équation et donc $\lambda \cdot v \in E$.
  \end{enumerate}

\item Il faut trouver une famille libre de vecteurs qui engendrent
  $E$. Comme $E$ est dans $\Rr^4$, il y aura moins de $4$ vecteurs
  dans cette famille. On prend un vecteur de $E$ (au hasard), par
  exemple $v_1 = (1,-1,0,0)$. Il est bien clair que $v_1$ n'engendre
  pas tout $E$, on cherche donc un vecteur $v_2$ lin\'eairement
  ind\'ependant de $v_1$, prenons $v_2 = (1,0,-1,0)$. Alors $\{v_1,v_2\}$
  n'engendrent pas tout $E$ ; par exemple $v_3 = (1,0,0,-1)$ est dans $E$
  mais n'est pas engendr\'e par $v_1$ et $v_2$. Montrons que
  $(v_1,v_2,v_3)$ est une base de $E$.
  \begin{enumerate}
  \item $(v_1,v_2,v_3)$ est une famille libre. En effet soient
    $\alpha,\beta,\gamma \in \Rr$ tels que $\alpha v_1+\beta
    v_2+\gamma v_3=0$. Nous obtenons donc :
\begin{align*}
& \alpha v_1+\beta v_2+\gamma v_3=0 \\
&\Rightarrow \alpha \begin{pmatrix}1 \cr -1 \cr 0 \cr 0 \cr
\end{pmatrix} + \beta
 \begin{pmatrix}1 \cr 0 \cr -1 \cr 0 \cr \end{pmatrix} + \gamma 
\begin{pmatrix}1 \cr 0 \cr 0 \cr -1 \cr \end{pmatrix}= 
\begin{pmatrix}0 \cr 0 \cr 0 \cr 0 \cr \end{pmatrix} \\
 &\Rightarrow \begin{cases}
  \alpha+\beta +\gamma &= 0 \\
  -\alpha &= 0 \\
  -\beta &=  0 \\
  -\gamma &=  0 \\
\end{cases}\\
 &\Rightarrow \alpha=0, \beta=0, \gamma = 0 \\
\end{align*}
Donc la famille est libre.

\item Montrons que la famille est g\'en\'eratrice : soit $v =
  (x_1,x_2,x_3,x_4) \in E$. Il faut \'ecrire $v$ comme combinaison
  lin\'eaire de $v_1,v_2,v_3$.  On peut r\'esoudre un syst\`eme comme
  ci-dessus (mais avec second membre) en cherchant
  $\alpha,\beta,\gamma$ tels que $\alpha v_1+\beta v_2+\gamma v_3=v$.
  On obtient que $v = -x_2v_1-x_3v_2-x_4v_4$ (on utilise
  $x_1+x_2+x_3+x_4=0$).
  \end{enumerate}
  
  Bien s\^ur vous pouvez choisir d'autres vecteurs de base (la seule
  chose qui reste ind\'ependante des choix est le nombre de vecteurs
  dans une base : ici $3$).
\end{enumerate}
\fincorrection
\nocorrection
\nocorrection
\nocorrection
\nocorrection
\nocorrection
\nocorrection
\correction{000908}

Montrons d'abord que $E \subset F$.
On va d'abord montrer que $v_1 \in F$ et $v_2 \in F$.

Tout d'abord 
$v_1\in F \iff v_1 \in \text{Vect}\{w_1,w_2\} \iff \exists \lambda, \mu \quad v_1 = \lambda w_1+\mu w_2 $.

Il s'agit donc de trouver ces $\lambda,\mu$. Cela se fait en résolvant un système
(ici on peut même le faire de tête) on trouve la relation 
$7 (2, 3, -1) = 3(3, 7, 0)-(5, 0, -7)$
ce qui donne la relation $v_1 = \frac37 w_1 -\frac 17 w_2$ et donc $v_1\in F$.

De même $7v_2=-w_1+2w_2$ donc $v_2\in F$.

Maintenant $v_1$ et $v_2$ sont dans l'espace vectoriel $F$, donc toute combinaison linéaire
de $v_1$ et $v_2$ aussi, c'est-à-dire : pour tout $\lambda,\mu$, on a $\lambda v_1+\mu v_2 \in F$.
Ce qui implique $E \subset F$.


\bigskip

Il reste à montrer $F\subset E$. Il s'agit donc d'écrire $w_1$ (puis $w_2$) en fonction de $v_1$ et $v_2$.
On trouve $w_1= 2v_1-v_2$ et $w_2=v_1+3v_2$. Encore une fois cela entraîne $w_1 \in E$ et $w_2 \in E$ donc
$\text{Vect}\{w_1,w_2\} \subset E$ d'où $F\subset E$.

Par double inclusion on a montré $E=F$. 

\fincorrection
\nocorrection
\nocorrection
\nocorrection
\nocorrection
\nocorrection
\correction{000914}
  $v\in \text{Vect}(e_1,e_2)$ est \'equivalent \`a l'existence de deux
  r\'eels $\lambda, \mu$ tels que $v = \lambda e_1+\mu e_2$.

  Alors $(-2,x,y,3)= \lambda (1,-1,1,2) + \mu (-1,2,3,1) $ est
  \'equivalent \`a
  $$
\begin{cases}
  -2 &= \lambda -\mu \\
  x  &= -\lambda +2\mu \\
  y &= \lambda +3\mu \\
  3 &= 2\lambda + \mu \\
 \end{cases}
 \qquad \Leftrightarrow \qquad
\begin{cases}
  \lambda &= 1/3 \\
  \mu &=  7/3 \\
  x &= 13/3\\
  y &=  22/3 \\
 \end{cases}.
 $$
 Le couple qui convient est donc $(x,y) = (13/3,22/3)$.
\fincorrection
\nocorrection
\correction{000916}
  \`A partir de la famille $(f_\alpha)_{\alpha\in \Rr}$ nous
  consid\'erons une combinaison lin\'eaire (qui ne correspond qu'\`a un
  nombre \emph{fini} de termes).

  
  Soient $\alpha_1,\ldots,\alpha_n$ des r\'eels distincts, consid\'erons la
  famille (finie) : $(f_{\alpha_i})_{i=1,\ldots,n}$. Supposons qu'il
  existe des r\'eels $\lambda_1,\ldots, \lambda_n$ tels que
  $\sum_{i=1}^n \lambda_i f_{\alpha_i}=0$. Cela signifie que, quelque
  soit $x \in \Rr$, alors $\sum_{i=1}^n \lambda_i f_{\alpha_i}(x) = 0$
  ; en particulier pour $x = \alpha_j$ l'\'egalit\'e devient $\lambda_j =
  0$ car $f_{\alpha_i}(\alpha_j)$ vaut $0$ si $i\neq j$ et $1$ si
  $i=j$. En appliquant le raisonnement ci-dessus pour $j=1$ jusqu'\`a
  $j=n$ on obtient : $\lambda_j = 0$, $j=1,\ldots,n$. Donc la famille
  $(f_\alpha)_{\alpha}$ est une famille libre.
\fincorrection
\correction{000917}
  \`A partir de la famille $(f_\alpha)_{\alpha\in \Rr}$ nous
  consid\'erons une combinaison lin\'eaire (qui ne correspond qu'\`a un
  nombre \emph{fini} de termes).

  
  Soient $\alpha_1 > \alpha_2 > \ldots > \alpha_n$ des r\'eels distincts que nous avons ordonnés, consid\'erons la
  famille (finie) : $(f_{\alpha_i})_{i=1,\ldots,n}$. Supposons qu'il
  existe des r\'eels $\lambda_1,\ldots, \lambda_n$ tels que
  $\sum_{i=1}^n \lambda_i f_{\alpha_i}=0$. Cela signifie que, quelque
  soit $x \in \Rr$, alors $\sum_{i=1}^n \lambda_i f_{\alpha_i}(x) = 0$,
autrement dit pour tout $x\in \Rr$ :
$$\lambda_1 e^{\alpha_1 x} + \lambda_2 e^{\alpha_2 x} + \cdots + \lambda_n e^{\alpha_n x}=0.$$
Le terme qui domine est $e^{\alpha_1 x}$ (car $\alpha_1>\alpha_2>\cdots$).
Factorisons par $e^{\alpha_1 x}$ :
$$e^{\alpha_1 x} \Big( \lambda_1  + \lambda_2 e^{(\alpha_2-\alpha_1) x} + \cdots 
+ \lambda_n e^{(\alpha_n-\alpha_1) x} \Big) =0.$$

Mais $e^{\alpha_1 x}\neq 0$ donc :
$$\lambda_1  + \lambda_2 e^{(\alpha_2-\alpha_1) x} + \cdots 
+ \lambda_n e^{(\alpha_n-\alpha_1) x} =0.$$
Lorsque $x\to +\infty$ alors $e^{(\alpha_i-\alpha_1) x} \to 0$ 
(pour tout $i\ge 2$, car $\alpha_i-\alpha_1<0$).
Donc pour $i\ge 2$, $\lambda_i e^{(\alpha_i-\alpha_1) x} \to 0$ et
en passant à la limite dans l'égalité ci-dessus on trouve :
$$\lambda_1=0.$$

Le premier coefficients est donc nul. On repart de la combinaison linéaire qui est maintenant
$\lambda_2 f_2+\cdots + \lambda_n f_n=0$ et en appliquant le raisonnement ci-dessus 
on prouve par récurrence $\lambda_1=\lambda_2=\cdots=\lambda_n=0$.
 Donc la famille $(f_\alpha)_{\alpha\in\Rr}$ est une famille libre.
\fincorrection
\nocorrection
\nocorrection
\nocorrection
\nocorrection
\nocorrection
\nocorrection
\nocorrection
\nocorrection
\correction{003297}
$F = G$.
\fincorrection
\nocorrection
\nocorrection
\correction{003303}
CNS $\Leftrightarrow \sum_{i=1}^n \alpha_i \ne -1$.
\fincorrection
\nocorrection
\correction{005167}
Soit $u'=(\sin(n\theta))_{n\in\Nn}$. On a~:~$u=1.u+0.u'$, puis $v=\cos a.u-\sin a.u'$, puis $w=\cos b.u-\sin b.u'$.
Les trois vecteurs $u$, $v$ et $w$ sont donc combinaisons linéaires des deux vecteurs $u$ et $u'$ et constituent par
suite une famille liée ($p+1$ combinaisons linéaires de $p$ vecteurs constituent une famille liée).
\fincorrection
\correction{005180}
\begin{enumerate}
\item  Notons respectivement $g$ et $h$, les fonctions sinus et cosinus.

$f_a=\cos a.g+\sin a.h$, $f_b=\cos b.g+\sin b.h$ et $f_c=\cos c.g+\sin c.h$. Donc, $f_a$, $f_b$ et $f_c$ sont trois
combinaisons linéaires des deux fonctions $g$ et $h$ et constituent donc une famille liée ($p+1$ combinaisons linéaires
de $p$ vecteurs donnés constituent une famille liée).

\item  $f_0$, $f_1$ et $f_2$ sont trois combinaisons linéaires des deux fonctions $x\mapsto1$ et $x\mapsto x$. Donc,
la famille $(f_0,f_1,f_2)$ est une famille liée puis la famille $(f_n)_{n\in\Zz}$ est liée en tant que sur-famille
d'une famille liée.

\item  Pour $\alpha$ réel donné et $x>0$, posons $f_\alpha(x)=x^\alpha$.

Soient $n$ un entier naturel supérieur ou égal à $2$, puis $(\alpha_1,...,\alpha_n)\in\Rr^n$ tel que
$\alpha_1<...<\alpha_n$. Soit encore $(\lambda_1,...,\lambda_n)\in\Rr^n$.

$$\sum_{k=1}^{n}\lambda_kf_{\alpha_k}=0\Rightarrow\forall
x\in]0;+\infty[,\;\sum_{k=1}^{n}\lambda_kx^{\alpha_k}=0\Rightarrow\forall
x\in]0;+\infty[,\;\sum_{k=1}^{n}\lambda_kx^{\alpha_k-\alpha_n}=0,$$

(en divisant les deux membres par $x^{\alpha_n}$). Dans cette dernière égalité, on fait tendre $x$ vers
$+\infty$ et on obtient $\lambda_n=0$. Puis, par récurrence descendante, $\lambda_{n-1}=...=\lambda_1=0$. On a montré
que toute sous-famille finie de la famille $(f_\alpha)_{\alpha\in\Rr}$ est libre et donc, la famille
$(f_\alpha)_{\alpha\in\Rr}$ est libre.

\item  Pour $a$ réel donné et $x$ réel, posons $f_a(x)=|x-a|$. Soient $n$ un
entier naturel supérieur ou égal à $2$, puis $a_1$,...,$a_n$, $n$ réels deux à deux distincts. Soit
$(\lambda_1,...,\lambda_n)\in\Rr^n$ tel que $\sum_{k=1}^{n}\lambda_kf_{a_k}=0$.

S'il existe $i\in\{1,...,n\}$ tel que $\lambda_i\neq 0$ alors,

$$f_{a_i}=-\frac{1}{\lambda_i}\sum_{k\neq i}^{}\lambda_kf_{a_k}.$$

Mais cette dernière égalité est impossible car $f_{a_i}$ n'est pas dérivable en $a_i$ alors que
$-\frac{1}{\lambda_i}\sum_{k\neq i}^{}\lambda_kf_{a_k}$ l'est. Donc, tous les $\lambda_i$ sont nuls.

\end{enumerate}
\fincorrection
\correction{005566}
\begin{enumerate}
 \item  La matrice de la famille $(e_1,e_2,e_3)$ dans la base canonique de $\Rr^4$ est $\left(
\begin{array}{ccc}
3&1&7\\
0&5&5\\
1&0&2\\
-2&-1&1
\end{array}
\right)$. Les trois dernières équations du système $\lambda e_1+\mu e_2+\nu e_3=0$ d'inconnues $\lambda$, $\mu$ et $\nu$ forment un sous-système de matrice $A=\left(
\begin{array}{ccc}
0&5&5\\
1&0&2\\
-2&-1&1
\end{array}
\right)$

En développant le déterminant de cette matrice suivant sa première colonne,  on obtient $\text{det}(A)=-10-2\times10=-30\neq0$. Ce sous-système est de \textsc{Cramer} et admet donc l'unique solution $(\lambda,\mu,\nu)=(0,0,0)$. Par suite, la famille $(e_1,e_2,e_3)$ est libre.

\item  

\begin{align*}\ensuremath
\left|
\begin{array}{cccc}
1&1&1&1\\
1&1&1&-1\\
1&1&-1&1\\
1&-1&1&1
\end{array}
\right|
&=\left|
\begin{array}{cccc}
1&1&1&1\\
0&0&0&-2\\
0&0&-2&0\\
0&-2&0&0
\end{array}
\right|\;(\text{pour}\;2\leqslant i\leqslant4,\;L_i\leftarrow L_i-L_1)\\
 &=\left|
\begin{array}{ccc}
0&0&-2\\
0&-2&0\\
-2&0&0
\end{array}
\right|=-8\neq0.
\end{align*}

Donc la famille $(e_1,e_2,e_3,e_4)$ est une famille libre (et donc une base de $E$).

\item  Notons $(u_1,u_2,u_3,u_4)$ la base canonique de $\Rr^4$.

La famille $(e_1,e_2,e_3,e_4)=(u_3,u_4,u_1,u_2)$ a même rang que la famille $(u_1,u_2,u_3,u_4)$ c'est-à-dire $4$. La famille $(e_1,e_2,e_3,e_4)$ est donc une base de $\Rr^4$.

\item  La matrice de la famille $(e_2,e_1,e_3,e_4)$ dans la base canonique de $\Rr^4$ est 
$\left(
\begin{array}{cccc}
1&2&4&1\\
1&-1&1&-2\\
1&3&5&2\\
1&1&3&0
\end{array}
\right)$. Cette matrice a même rang que les matrices suivantes :

$\left(
\begin{array}{cccc}
1&0&0&0\\
1&-3&-3&-3\\
1&1&1&1\\
1&-1&-1&-1
\end{array}
\right)$  ($e_5=e_1-2e_2$, $e_6=e_3-4e_2$ et $e_7=e_4-e_2$)

$\left(
\begin{array}{cccc}
1&0&0&0\\
1&-3&0&0\\
1&1&10&0\\
1&-1&0&0
\end{array}
\right)$ ($e_8=e_6-e_5$ et $e_9=e_7-e_5$).

La matrice ci-dessus est de rang $2$. Il en est de même de la famille $(e_1,e_2,e_3,e_4)$ qui est en particulier liée. La nullité de la troisième colonne fournit $0=e_8=e_6-e_5=(e_3-4e_2)-(e_1-2e_2)=-e_1-2e_2+e_3$ et donc $e_3=e_1+2e_2$. La nullité de la quatrième colonne fournit $0=e_9=e_7-e_5=(e_4-e_2)-(e_1-2e_2)=e_4+e_2-e_1$ et donc $e_4=e_1-e_2$.
\end{enumerate}
\fincorrection
\correction{005567}
Soit $(a,b,c)\in\Qq^3$.

\begin{center}
$a+b\sqrt{2}+\sqrt{3}= 0\Rightarrow(a+b\sqrt{2})^2=(-c\sqrt{3})^2\Rightarrow a^2+2b^2+2ab\sqrt{2}=3c^2\Rightarrow2ab\sqrt{2}\in\Qq.$
\end{center}

Mais $\sqrt{2}$ est irrationnel donc $ab=0$.

Si $b=0$, puisque $a+c=0$ et que $\sqrt{3}$ est irrationnel, on en déduit que $c = 0$ (sinon $\sqrt{3}$  serait rationnel) puis $a=0$ et finalement $a=b=c=0$.

Si $a=0$, il reste $2b^2=3c^2$. Mais $\sqrt{\frac{3}{2}}$ est irrationnel (dans le cas contraire, il existe deux entiers $p$ et $q$ non nuls tels que $3q^2 = 2p^2$ et par exemple l'exposant du nombre premier $2$ n'a pas la même parité dans les deux membres de l'égalité ce qui est impossible) et donc $b=c=0$ puis encore une fois $a=b=c=0$.  

On a montré que $\forall(a,b,c)\in\Qq^3$, $(a+b\sqrt{2}+\sqrt{3}= 0\Rightarrow a=b=c=0)$.
Donc la famille $(1,\sqrt{2},\sqrt{3})$ est une famille de réels $\Qq$-libre.
\fincorrection
\correction{005568}
Les fonctions $f_1$, $f_2$ et $f_3$ sont bien définies sur $\Rr^+$.

Soient $a$, $b$ et $c$ trois réels tels que $af_1+bf_2 +cf_3=0$.

\textbf{Première solution.} Si $a$ est non nul, la fonction $af_1+bf_2+cf_3$ est équivalente au voisinage de $+\infty$ à $a\ln x$ et ne peut donc être égale à la fonction nulle. Donc $a=0$. Puis si $b$ est non nul, la fonction $af_1+bf_2+cf_3=bf_2+cf_3$ est équivalente à $b\ln(\ln x)$ et ne peut être égale à la fonction nulle. Donc $b =0$. Puis $c=0$.

\textbf{Deuxième solution.} On effectue un développement limité à un ordre suffisant de la fonction $af_1+bf_2+cf_3$ quand $x$ tend vers $0$ :

$f_1(x)=\ln(1+x)\underset{x\rightarrow0}{=}x-\frac{x^2}{2}+\frac{x^3}{3}+o(x^3)$ puis

\begin{align*}\ensuremath
f_2(x)&=\ln(1+f_1(x))\underset{x\rightarrow0}{=}\ln\left(1+x-\frac{x^2}{2}+\frac{x^3}{3}+o(x^3)\right) =\left(x-\frac{x^2}{2}+\frac{x^3}{3}\right)-\frac{1}{2}\left(x-\frac{x^2}{2}\right)^2+\frac{1}{3}x^3+o(x^3)\\
 &=x-x^2+\frac{7}{6}x^3+o(x^3)
\end{align*}

puis

\begin{align*}\ensuremath
f_3(x)&=\ln(1+f_2(x))\underset{x\rightarrow0}{=}\ln\left(1+x-x^2+\frac{7}{6}x^3+o(x^3)\right)=\left(x-x^2+\frac{7}{6}x^3\right)-\frac{1}{2}\left(x-x^2\right)^2+\frac{1}{3}x^3+o(x^3)\\
 &=x-\frac{3}{2}x^2+\frac{5}{2}x^3+o(x^3).
\end{align*}

Par suite, $af_1(x)+bf_2(x)+cf_3(x)\underset{x\rightarrow0}{=}(a+b+c)x+\left(-\frac{a}{2}-b-\frac{3c}{2}\right)x^2+\left(\frac{a}{3}+\frac{7b}{6}+\frac{5c}{2}\right)x^3+o(x^3)$. L'égalité $af_1+bf_2+cf_3=0$ fournit, par identification des parties régulières des développements limités à l'ordre trois en zéro :

\begin{center}
$\left\{
\begin{array}{l}
a+b+c=0\\
-\frac{a}{2}-b-\frac{3c}{2}=0\\
\rule{0mm}{6mm}\frac{a}{3}+\frac{7b}{6}+\frac{5c}{2}
\end{array}
\right.$ ou encore $\left\{
\begin{array}{l}
a+b+c=0\\
a+2b+3c=0\\
2a+7b+15c=0
\end{array}
\right.$.
 \end{center}
 
 
Comme $\left|
\begin{array}{ccc}
1&1&1\\
1&2&3\\
2&7&15
\end{array}
\right|=\left|
\begin{array}{ccc}
1&0&0\\
1&1&1\\
2&5&8
\end{array}
\right|=3\neq0$, on a donc $a=b=c=0$.
\fincorrection
\correction{005569}
Soient $n$ un entier naturel non nul puis $a_1$,..., $a_n$ $n$ réels deux à deux distincts et $\lambda_1$,..., $\lambda_n$ $n$ réels.

Supposons $\lambda_1f_{a_1}+...+\lambda_nf_{a_n}=0$. Soit $i$ un élément de $\llbracket1,n\rrbracket$. On a $\lambda_if_{a_i}=-\sum_{j\neq i}^{}\lambda_jf_{a_j}$ et on ne peut avoir $\lambda_i\neq0$ car alors le membre de gauche est une fonction non dérivable en $a_i$ tandis que le membre de droite l'est. Par suite, tous les $\lambda_i$ sont nuls et donc la famille $(f_{a_i})_{1\leqslant i\leqslant n}$ est libre.

On a montré que toute sous-famille finie de la famille $(f_a)_{a\in\Rr}$ est libre et donc la famille $(f_a)_{a\in\Rr}$ est libre.
\fincorrection
\correction{005570}
Soient $a_1< ... <a_n$ $n$ réels deux à deux distincts et $\lambda_1$,..., $\lambda_n$ n réels tels que $\sum_{i=1}^{n}\lambda_if_{a_i}=0$ $(*)$.

\textbf{Première solution.} Après multiplication des deux
 membres de $(*)$ par $e^{-a_nx}$ puis passage à la limite quand $x$ tend vers $+\infty$, on obtient $\lambda_n=0$. En réitérant, on obtient donc $\lambda_n=\lambda_{n-1}=...=\lambda_1=0$.
 

\textbf{Deuxième solution.} On note $f$ la fonction apparaissant au premier membre de $(*)$.

\begin{align*}\ensuremath
f=0&\Rightarrow \forall k\in\llbracket0,,n-1\rrbracket,\;f^{(k)}(0)=0\\
 &\Rightarrow\forall k\in\llbracket0,n-1\rrbracket,\;\lambda_1a_1^k+...+\lambda_na_n^k=0.
\end{align*}

Le système prédédent d' inconnues $\lambda_i$, $1\leqslant n$, est un système linéaire homogène à $n$ équations et $n$ inconnues. Son déterminant est le déterminant de \text{Vandermonde} des $a_i$ et est non nul puisque les $a_i$ sont deux à deux distincts. Le système est donc de \textsc{Cramer} et admet l'unique solution $(0,...,0)$.

\textbf{Troisième solution.} ( dans le cas où on se restreint à démontrer la liberté de la famille $(x\mapsto e^{nx})_{n\in\Nn}$).

Soient $n_1< ... < n_p$ $p$ entiers naturels deux à deux distincts.
Supposons que pour tout réel $x$ on ait $\sum_{i=1}^{n}\lambda_ie^{n_ix}=0$. On en déduit que pour tout réel strictement positif $t$, on a $\sum_{i=1}^{n}\lambda_it^{n_i}=0$ et donc le polynôme  $\sum_{i=1}^{n}\lambda_iX^{n_i}$ est nul (car a une infinité de racines) ou encore les coefficients du polynôme $\sum_{i=1}^{n}\lambda_iX^{n_i}$ à savoir les $\lambda_i$ sont tous nuls.

\textbf{Quatrième solution.} (pour les redoublants)
L'application $\varphi$ qui à $f$ de classe $C^\infty$ fait correspondre sa dérivée est un endomorphisme de l'espace des fonctions de classe $C^\infty$ sur $\Rr$ à valeurs dans $\Rr$.
Pour $a$ réel donné, $\varphi(f_a)=af_a$ et la famille $(f_a)_{a\in\Rr}$ est constituée de vecteurs propres de $\varphi$ (les $f_a$ sont non nulles) associés à des valeurs propres deux à deux distinctes. On sait qu'une telle famille est libre.

\fincorrection
\correction{005571}
Soient $n$ un entier naturel non nul puis $P_1$,..., $P_n$ $n$ polynômes non nuls de degrés respectifs $d_1<...<d_n$.

Soit $(\lambda_1,\ldots,\lambda_n)\in\Kk^n$ tel que $\lambda_1P_1 + ... +\lambda_nP_n=0$. Supposons par l'absurde que les $\lambda_i$ ne soient pas tous nuls et posons $k=\text{Max}\left\{i\in\llbracket1,n\rrbracket/\;\lambda_i\neq0\right\}$. On ne peut avoir $k=1$ car $P_1\neq0$ puis

\begin{center}
$\lambda_1P_1 + ... +\lambda_nP_n=0\Rightarrow\lambda_1P_1 + ... +\lambda_kP_k=0\Rightarrow\lambda_kP_k=-\sum_{i<k}^{}\lambda_iP_i$.
\end{center}

Cette dernière égalité est impossible car $\lambda_kP_k$ est un polynôme de degré $d_k$ (car $\lambda_k\neq0$) et $-\sum_{i<k}^{}\lambda_iP_i$ est un polynôme de degré au plus $d_{k-1}<d_k$. Donc tous les $\lambda_k$ sont nuls.

La même démarche tient en remplaçant degré par valuation et en s'intéressant à la plus petite valuation au lieu du plus grand degré.

\fincorrection
\correction{005574}
\begin{enumerate}
 \item  Pour $p$ et $q$ entiers relatifs, posons $I(p,q)=\int_{0}^{2\pi}e^{i(p-q)x}\;dx$.

Si $p\neq q$, $I(p,q)=\frac{1}{i(p-q)}\left[e^{i(p-q)x}\right]_0^{2\pi}=0$. Soient alors $p$ et $q$ deux entiers naturels.

Donc si $p\neq q$, $J(p,q)\frac{1}{2}\text{Re}(I(p,q) + I(p,-q))=0$ puis $K(p,q)=\frac{1}{2}\text{Im}(I(p,-q)-I(p,q))=0$ puis $L(p,q)=\frac{1}{2}\text{Re}(I(p,-q)-I(p,q))=0$.

Si $p=q$, $J(p,p)=2\pi$ si $p=0$ et $\pi$ si $p\neq 0$ puis $K(p,p)=0$ puis $L(p,p)=\pi$ si $p\neq0$ et $0$ si $p=0$.

\item  Sur l'espace $E$ des fonctions continues sur $\Rr$ à valeurs dans $\Rr$ et $2\pi$-périodiques, l'application qui à $(f,g)$ élément de $E^2$ associe $\int_{0}^{2\pi}f(t)g(t)\;dt$ est classiquement un produit scalaire. La famille de fonctions proposée est une famille orthogonale pour ce produit scalaire et ne contient pas le vecteur nul de $E$. Cette famille est donc est libre.
\end{enumerate}
\fincorrection
\correction{006870}
\begin{enumerate}
  \item 
\begin{align*}
     & \alpha v_1 + \beta v_2 + \gamma v_3 = 0 \\ 
\iff & \alpha (2,1,4) + \beta (1,-1,2) + \gamma (3,3,6) = (0,0,0) \\
\iff &  \Big(2\alpha+\beta+3\gamma,\alpha-\beta+3\gamma,4\alpha+2\beta+6\gamma\Big) = (0,0,0) \\
\iff &
\begin{cases}
  2\alpha+\beta+3\gamma &= 0 \\
  \alpha-\beta+3\gamma  &= 0 \\
  4\alpha+2\beta+6\gamma &= 0 \\
 \end{cases} \\
\iff & \cdots  \qquad  \text{(on résout le système)} \\
\iff & \alpha=-2t, \beta = t, \gamma = t \quad t \in \Rr \\
\end{align*}  

Si l'on prend $t=1$ par exemple alors $\alpha=-2$, $\beta = 1$, $\gamma = 1$
donne bien $-2v_1+v_2+v_3=0$.

Cette solution n'est pas unique, les autres coefficients qui conviennent sont les 
$(\alpha=-2t, \beta = t, \gamma = t)$ pour tout $t \in \Rr$.

  \item 

Il s'agit donc de trouver un vecteur $v=(x,y,z)$ dans $P_1$ et $P_2$ et donc qui doit vérifier 
$x-y+z=0$ et $x-y=0$ :

\begin{align*}
     & v=(x,y,z) \in P_1 \cap P_2 \\ 
\iff & x-y+z=0 \text{ et } x-y=0 \\
\iff &
\begin{cases}
  x-y-z = 0 \\
  x-y = 0 \\
 \end{cases} \\
\iff & \cdots \qquad  \text{(on résout le système)} \\
\iff & (x=t, y = t, z = 0) \quad t \in \Rr \\
\end{align*}  

Donc, si l'on fixe par exemple $t=1$, alors $v=(1,1,0)$ 
est un vecteur directeur de la droite vectorielle $D$,
une équation paramétrique étant $D=\{(t,t,0) \mid t\in \Rr \}$.
\end{enumerate}
\fincorrection
\correction{007410}
\begin{enumerate}
 \item On observe que $v_3 = v_2 - v_1$ et que $v_4=v_1+v_2$ donc $V=Vect(v_1,v_2,v_3,v_4)=Vect(v_1,v_2)$. Comme de plus $v_1$ et $v_2$ sont linéairement indépendants, $(v_1,v_2)$ est une base de $V$. Donc $\dim(V)=2$. Le système n'est pas libre car par exemple $v_3=v_2-v_1$ est une relation de dépendance linéaire non triviale. Le système n'est pas générateur de $\Rr^4$ car $\dim(V)=2<4$ donc $Vect(v_1,v_2,v_3,v_4)=V\neq \Rr^4$.

 \item Comme nous l'avons vu dans la question 1, $(v_1,v_2)$ est une base de $V$. On échelonne la matrice dont les lignes sont les vecteurs $v_1$ et $v_2$.

\[
\left( {\begin{array}{cccc}
 1 & 0 & -1& 1 \\
 2 & 1 & 0 &1 \\
 \end{array} } \right)
\to
\left( {\begin{array}{cccc}
 1 & 0 & -1& 1 \\
 0 & 1 & 2 & -1 \\
 \end{array} } \right)
\]

Maintenant on voit que l'on peut compléter la base $(v_1,v_2)$ en une base de $\Rr^4$ en ajoutant les vecteurs $(0,0,1,0)$ et $(0,0,0,1)$ car la matrice
\[
\left( {\begin{array}{cccc}
 1 & 0 & -1& 1 \\
 0 & 1 & 2 &-1 \\
 0& 0& 1 & 0 \\
 0 & 0 & 0 & 1\\
 \end{array} } \right)
\]
est de rang 4.

 \item $V=Vect(v_1,v_2)$. Maintenant soit $(x,y,z,t)\in \Rr^4$. $(x,y,z,t)\in V$ si et seulement si $\exists \lambda,\mu \in \Rr $ tels que $(x,y,z,t)=\lambda v_1 + \mu v_2$.
Il nous suffit donc de voir pour quelles valeurs de $(x,y,z,t)$ ce système d'équations a une solution. Le système s'écrit :
\[
\left( {\begin{array}{cc|c}
 1 & 2 & x \\
 0 & 1 & y \\
 -1& 0& z \\
 1 & 1 & t\\
 \end{array} } \right)
 \to
 \left( {\begin{array}{cc|c}
 1 & 2 & x \\
 0 & 1 &y \\
 0& 2& z+x \\
 0 & -1 & t-x\\
 \end{array} } \right)
 \to
 \left( {\begin{array}{cc|c}
 1 & 2 & x \\
 0 & 1 & y \\
 0& 0& z+x-2y \\
 0 & 0 & t-x+y\\
 \end{array} } \right)
\]
Le système a donc une solution si et seulement si $z+x-2y=0$ et $t-x+y=0$,ce sont des équations cartesiennes pour $V$, c'est à dire que
$$V=\{(x,y,z,t)\in\Rr^4 \ / \ \ x-2y+z=0 \ \ et \ \ -x+y+t=0 \}.$$

 \item $H$ est un sous-espace vectoriel de $\Rr^4$ car c'est l'ensemble des solutions d'une équation linéaire et homogène.

 \item \begin{eqnarray}
H &=& \{(x,y,z,t)\in\Rr^4 \ / \ \ -3x+y+2z-t=0 \} \nonumber \\
 &=& \{(x,y,z,t)\in\Rr^4 \ / \ \ y=3x-2z+t \} \nonumber \\
 &=& \{(x,3x-2z+t,z,t) \ / \ \ x,z,t\in \Rr \} \nonumber \\
 &=& \{x(1,3,0,0)+z(0,-2,1,0)+t(0,1,0,1) \ / \ \ x,z,t\in \Rr \} \nonumber
\end{eqnarray}
On a donc $H=Vect((1,3,0,0),(0,-2,1,0),(0,1,0,1))$. De plus le système $((1,3,0,0),(0,-2,1,0),(0,1,0,1))$ est une base de $H$ car 
 \[
rang \left( {\begin{array}{cccc}
 1 & 3 & 0& 0 \\
 0 & 1 & 0 &1 \\
 0 & -2 & 1 &0 \\
 \end{array} } \right)
=
rang \left( {\begin{array}{cccc}
 1 & 3 & 0& 0 \\
 0 & 1 & 0 &1 \\
 0 & 0 & 1 &2 \\
 \end{array} } \right)=3.
\]
Et donc $\dim(H)=3$.

 \item $v_3\in H$ car $-3\times 1 +1\times 1 + 2\times 1 -1\times 0 = 0$. Mais $v_1\notin H$ car $-3\times 1 +1\times 0 + 2\times -1 -1\times 1 = -6\neq 0$.

Comme $v_3\in V$ et que l'on vient de voir que $v_3\in H$, on en déduit que $v_3\in V\cap H$ et donc que $Vect(v_3)\subseteq V\cap H$.
 On a donc obtenu que $\dim (V\cap H)\geq 1$. Comme de plus $V\cap H \subseteq V$ on sait que $\dim (V\cap H)\leq 2$. Maintenant si on avait $\dim (V\cap H)=2$ cela impliquerait que $V\cap H =V$ et donc que $V\subseteq H$, ce qui est faux car $v_1\in V$ mais $v_1\notin H$. 
 On obtient donc $\dim (V \cap H)=1$.
 
En utilisant la formule de Grassmann on obtient alors que $$\dim (V+H)=\dim (V)+\dim (H)-\dim (V\cap H)=2+3-1=4.$$

 \item On a vu que $\dim(V\cap H)=1$ et que $Vect(v_3)\in V\cap H$ donc on obtient $Vect(v_3)= V\cap H$ et donc $(v_3)$ est une base de $V\cap H$.
 \end{enumerate}
\fincorrection
\correction{000919}
Faisons d'abord une remarque qui va simplifier les calculs : 
$$v_3 = 2v_1+3v_2.$$
 Donc en fait nous avons 
$\text{Vect} \{v_1,v_2,v_3\}= \text{Vect}\{v_1,v_2\}$ et c'est un espace de dimension $2$, c'est-à-dire un plan vectoriel.
Par la m\^eme relation on trouve que $\text{Vect} \{v_1,v_2,v_3\}= \text{Vect}\{v_2,v_3\}$.

\begin{enumerate}
\item Vrai. $\text{Vect}\{(1,1,0,0),(-1,1,-4,2)\}$ est inclus dans $\text{Vect} \{v_1,v_2,v_3\}$, car
$(1,1,0,0) = v_1+v_2$ et $(-1,1,-4,2)=v_1-v_2$. Comme ils sont de m\^eme dimension ils sont \'egaux 
(autrement dit : comme un plan est inclus dans un autre alors ils sont égaux).
\item Vrai. On a $(1,1,0,0) = v_1+v_2$ donc  $(1,1,0,0)\in \text{Vect}\{v_1,v_2\}$, or 
$\text{Vect}\{v_1,v_2\}=\text{Vect}\{v_2,v_3\}\subset \text{Vect}\{v_2,v_3,v_4\}$. 
Donc $(1,1,0,0) \in \text{Vect}\{v_1, v_2\} \cap    \text{Vect}\{v_2, v_3, v_4\}$.
\item Faux. Toujours la m\^eme relation nous donne que 
$\text{Vect} \{v_1, v_2\} \cap \text{Vect} \{v_2, v_3, v_4\} = \text{Vect}\{v_1, v_2\}$ donc est de dimension $2$.
C'est donc un plan vectoriel et pas une droite.
 \item Faux. Encore une fois la relation donne que 
$\text{Vect}\{v_1, v_2\}+\text{Vect}\{v_2, v_3, v_4\} = \text{Vect}\{v_1, v_2, v_4\}$, or $3$ vecteurs ne
   peuvent engendrer $\Rr^4$ qui est de dimension $4$.
\item Vrai.  Faire le calcul : l'intersection est $\{0\}$ et la somme est $\Rr^4$.
\end{enumerate}
\fincorrection
\correction{000920}
\begin{enumerate}
\item Non. 
Tout d'abord par définition
$\text{Vect}\{v_1,v_2\}+\text{Vect}\{v_3\} = \text{Vect}\{v_1,v_2,v_3\}$,
Nous allons trouver un vecteur de $\Rr^4$ qui n'est pas dans 
$\text{Vect}\{v_1,v_2\}+\text{Vect}\{v_3\}$.
Il faut tâtonner un peu pour le choix, par exemple faisons le calcul avec $u=(0,0,0,1)$.

$u \in \text{Vect}\{v_1,v_2,v_3\}$ si et seulement si 
il existe des réels $\alpha,\beta,\gamma$ tels que $u=\alpha v_1 + \beta v_2 + \gamma v_3$.
Si l'on écrit les vecteurs verticalement, on cherche donc $\alpha,\beta,\gamma$ tels que :
$$\begin{pmatrix} 0 \\ 0 \\ 0 \\1\end{pmatrix}= \alpha \begin{pmatrix} 1 \\ 0 \\ 0 \\1\end{pmatrix}
+ \beta \begin{pmatrix} 0 \\ 0 \\ 1 \\ 0 \end{pmatrix}+\gamma\begin{pmatrix} 0 \\ 1 \\ 0 \\0\end{pmatrix}$$
Ce qui est équivalent à trouver $\alpha,\beta,\gamma$ vérifiant le système linéaire :
$$\begin{cases}
0 = \alpha \cdot 1 + \beta \cdot 0 + \gamma \cdot 0   \\
0 = \alpha \cdot 0 + \beta \cdot 0 + \gamma \cdot 1   \\
0 = \alpha \cdot 0 + \beta \cdot 1 + \gamma \cdot 0   \\
1 = \alpha \cdot 1 + \beta \cdot 0 + \gamma \cdot 0   \\
\end{cases}
\text{ qui équivaut à } 
\begin{cases}
0 = \alpha \\
0 = \gamma \\
0 = \beta  \\
1 = \alpha \\
\end{cases}
$$
Il n'y a clairement aucune solution à ce système (les trois premières lignes impliquent $\alpha=\beta=\gamma=0$
et cela rentre alors en contradiction avec la quatrième).

\bigskip


Un autre type de raisonnement, beaucoup plus rapide, 
est de dire que ces deux espaces ne peuvent engendr\'es tout $\Rr^4$ 
car il n'y pas assez de vecteurs en effet $3$ vecteurs ne peuvent engendrer 
l'espace $\Rr^4$ de dimension $4$.

\item Oui. Notons $F=\mbox{Vect}\{v_1,v_2\}$ et $G=\mbox{Vect}\{v_4, v_5\}$.
Pour montrer $F \oplus G = \Rr^4$ il faut montrer $F\cap G =\{ (0,0,0,0) \}$ et
$F+G=\Rr^4$.

  \begin{enumerate}
     \item Montrons $F\cap G =\{ (0,0,0,0 \}$. Soit $u \in F\cap G$,
d'une part $u \in F = \mbox{Vect}\{v_1,v_2\}$ donc il existe $\alpha,\beta \in \Rr$ tels que
$u = \alpha v_1+\beta v_2$. D'autre part $u \in G = \mbox{Vect}\{v_4, v_5\}$ donc il existe
$\gamma,\delta \in \Rr$ tels que $u=\gamma v_4 + \delta v_5$. On a écrit $u$ de deux façons donc on a l'égalité
$\alpha v_1+\beta v_2 = \gamma v_4 + \delta v_5$. En écrivant les vecteurs comme des vecteurs colonnes cela donne
$$\alpha \begin{pmatrix} 1 \\ 0 \\ 0 \\1\end{pmatrix}
+ \beta \begin{pmatrix} 0 \\ 0 \\ 1 \\ 0 \end{pmatrix} 
= \gamma \begin{pmatrix} 0 \\ 0 \\ 0 \\1\end{pmatrix} + \delta \begin{pmatrix} 0 \\ 1 \\ 0 \\1\end{pmatrix}$$
Donc $(\alpha,\beta,\gamma,\delta)$ est solution du système linéaire suivant :
$$\begin{cases}
\alpha = 0   \\
0 = \delta \\
\beta = 0 \\
\alpha =  \gamma + \delta \\
\end{cases}
$$
Cela implique $\alpha=\beta=\gamma=\delta=0$ et donc $u = (0,0,0,0)$.
Ainsi le seul vecteur de $F\cap G$ est le vecteur nul.


     \item Montrons $F+G=\Rr^4$.
$F+G = \mbox{Vect}\{v_1,v_2\}+\mbox{Vect}\{v_4, v_5\}= \mbox{Vect}\{v_1,v_2,v_4, v_5\}$.
Il faut donc montrer que n'importe quel vecteur $u = (x_0,y_0,z_0,t_0)$ de $\Rr^4$ s'écrit comme une combinaison linéaire
de $v_1,v_2,v_4,v_5$. Fixons $u$ et cherchons $\alpha,\beta,\gamma,\delta \in \Rr$ tels que 
$\alpha v_1+\beta v_2 + \gamma v_4 + \delta v_5=u$.
Après avoir considéré les vecteurs comme des vecteurs colonnes cela revient à résoudre le système linéaire :
$$\begin{cases}
\alpha = x_0 \\
\delta = y_0 \\
\beta = z_0 \\
\alpha + \gamma + \delta = t_0 \\    
  \end{cases}$$
Nous étant donné un vecteur $u = (x_0,y_0,z_0,t_0)$ on a calculé qu'en choisissant 
$\alpha = x_0$, $\beta=z_0$, $\gamma= t_0 -x_0-y_0$, $\delta=y_0$ on obtient bien
$\alpha v_1+\beta v_2 + \gamma v_4 + \delta v_5=u$.
Ainsi tout vecteur est engendré par $F+G$.
  \end{enumerate}
Ainsi $F\cap G =\{ (0,0,0,0) \}$ et $F+G=\Rr^4$ donc $F \oplus G = \Rr^4$.


\item Non. Ces deux espaces ne sont pas suppl\'ementaires car il y a trop de vecteurs !
Il engendrent tout, mais l'intersection n'est pas triviale. En effet on remarque assez vite que 
$v_5=v_3+v_4$ est dans l'intersection. On peut aussi obtenir ce r\'esultat en résolvant un syst\`eme.

\item Non. Il y a bien quatre vecteurs mais il existe des relations entre eux.

On peut montrer $\mbox{Vect}\{v_1,v_4\}$ et $\mbox{Vect}\{v_3, v_5\}$ ne sont pas supplémentaires
de deux façons. Première méthode : leur intersection est non nulle, par exemple
$v_4=v_5-v_3$ est dans l'intersection. Deuxième méthode : les deux espaces n'engendrent pas tout,
en effet il est facile de voir que $(0,0,1,0) \notin \mbox{Vect}\{v_1,v_4\} + \mbox{Vect}\{v_3, v_5\} 
= \mbox{Vect}\{v_1,v_4,v_3, v_5\}$.


\end{enumerate}
\fincorrection
\nocorrection
\nocorrection
\correction{000923}
  Analysons d'abord les fonctions de $E$ qui ne sont pas dans $F$ : ce sont 
  les fonctions $h$ qui v\'erifient $h(0) \neq 0$ \textbf{ou} $h'(0) \neq 0$. Par exemple
  les fonctions constantes $x \mapsto b$, ($b \in \Rr^*$) ou les
  homoth\'eties $x \mapsto a x$, ($a \in \Rr^*$) n'appartiennent pas \`a
  $F$.
  
  Cela nous donne l'idée de poser
  $$G= \left\lbrace x \mapsto ax+b \mid (a,b) \in \Rr^2 \right\rbrace.$$
  Montrons que $G$ est un suppl\'ementaire de $F$ dans $E$.
  
  Soit $f \in F \cap G$, alors $f(x) = ax+b$ (car $f\in G$) et $f(0) =
  b$ et $f'(0)=a$ ; mais $f\in F$ donc $f(0) = 0$ donc $b=0$ et
  $f'(0)=0$ donc $a=0$. Maintenant $f$ est la fonction nulle : $F\cap
  G = \{ 0 \}$.
  
  Soit $h \in E$, alors remarquons que pour $f(x) = h(x) - h(0)
  -h'(0)x$ la fonction $f$ v\'erifie $f(0) = 0$ et $f'(0)=0$ donc $f \in
  F$. Si nous \'ecrivons l'\'egalit\'e diff\'eremment nous obtenons
  $$
  h(x) = f(x) +h(0)+h'(0)x.$$
  Posons $g(x) = h(0)+h'(0)x$, alors la
  fonction $g \in G$ et
  $$h =f +g,$$
  ce qui prouve que toute fonction de $E$ s'\'ecrit comme
  somme d'une fonction de $F$ et d'une fonction de $G$ : $E = F + G$.
  
  En conclusion nous avons montré que $E = F \oplus G$.
\fincorrection
\nocorrection
\nocorrection
\correction{000926}
 On note $F$ l'espace vectoriel des suites constantes et $G$ l'espace
  vectoriel des suites convergeant vers $0$.
\begin{enumerate}
\item $F\cap G = \{0\}$. En effet une suite constante qui converge
  vers $0$ est la suite nulle.
  \item $F+G=E$. Soit $(u_n)$ un \'el\'ement de $E$. Notons $\ell$ la
    limite de $(u_n)$.  Soit $(v_n)$ la suite d\'efinie par
    $v_n=u_n-\ell$, alors $(v_n)$ converge vers $0$. Donc $(v_n)\in
    G$.  Notons $(w_n)$ la suite constante \'egale \`a $\ell$. Alors nous
    avons $u_n=\ell+u_n-\ell$, ou encore $u_n=w_n+v_n$, ceci pour tout
    $n\in\Nn$. En terme de suite cela donne $(u_n)=(w_n)+(v_n)$. Ce
    qui donne la d\'ecomposition cherch\'ee.
\end{enumerate}
Bilan : $F$ et $G$ sont en somme directe dans $E$ : $E=F\oplus G$.
\fincorrection
\correction{002433}
Notons l'ancienne base $\mathcal{B}=(e_1,e_2,e_3)$
et ce qui sera la nouvelle base $\mathcal{B}'=(e'_1,e'_2,e'_3)$.
Soit $P$ la matrice de passage qui contient -en colonnes- les coordonnées des vecteurs
de la nouvelle base $\mathcal{B}'$ exprimés dans l'ancienne base $\mathcal{B}$
$$P=\begin{pmatrix}
2 & 3 & 1 \\
3 & 4 & 2 \\
1 & 1 & 2 \\      
\end{pmatrix}$$

On vérifie que $P$ est inversible (on va même calculer son inverse) donc
$\mathcal{B}'$ est bien une base.
De plus 
$$P^{-1} = \begin{pmatrix}
-6 & 5 & -2 \\
4 & -3 & 1 \\
1 & -1 & 1 \\             
           \end{pmatrix}
\text{ et on calcule  } B=P^{-1} A P = 
\begin{pmatrix}
1 & 0 & 0 \\
0 & 2 & 0 \\
0 & 0 & 3 \\  
\end{pmatrix}$$

$B$ est la matrice de $f$ dans la base $\mathcal{B}'$.
\fincorrection
\correction{003305}
\begin{enumerate}
  \item
  \begin{enumerate}
    \item Soit $P\in\R[X]$ que l'on décompose en $P=P_1(X^2) + XP_2(X^2)$ .
    Alors $P = (P_1+P_2)(X^2) - (1-X)P_2(X^2) = (1-X)P_1(X^2) + X(P_1+P_2)(X^2)$,
    ce qui prouve que les deux sommes sont égales à $\R[X]$. Ces sommes
    sont facilement directes.
    \item Cela ne change pas~$A$~: les éléments de~$A$ sont ceux dont les parties
    paire et impaire sont opposées (au facteur $X$ près), in\-dé\-pen\-da\-ment du fait (vrai) que ces parties sont
    des polynômes.
  \end{enumerate}
  \item Soit $f$ un isomorphisme de $E_1$ sur $E_2$ et $F = \{x - f(x)$ tq $x\in E_1 \}$.
Alors $E = E_1 \oplus F = E_2\oplus F$.
\end{enumerate}
 \fincorrection
\nocorrection
\nocorrection
\nocorrection
\nocorrection
\nocorrection
\nocorrection
\nocorrection
\correction{003659}
Les $f_i$ sont des projecteurs commutant deux à deux, ils sont simmultanément
diagonalisables. Soit $e_1$ tel que $f_1(e_1) = e_1$~: $f_i(e_1) = f_i\circ f_1(e_1) = 0$
si $i\ge 2$ donc les supports des restrictions des $f_i$ à une base propre
commnue sont deux à deux disjoints non vides, ce sont des singletons.
\fincorrection
\correction{005178}
$F=\mbox{Vect}(u)$ est un sous espace vectoriel de $\Rr^n$ et $G$ est un sous espace vectoriel de $\Rr^n$, car est le
noyau de la forme linéaire $(x_1,...,x_n)\mapsto x_1+...+x_n$.

Soit $x=(x_1,...,x_n)\in\Rr^n$ et soit $\lambda\in\Rr$.

$$x-\lambda u\in G\Leftrightarrow(x_1-\lambda,...,x_n-\lambda)\in
G\Leftrightarrow\sum_{k=1}^{n}(x_k-\lambda)=0\Leftrightarrow\lambda=\frac{1}{n}\sum_{k=1}^{n}x_k.$$

Donc,

$$\forall x\in\Rr^n,\;\exists!\lambda\in\Rr/\;x-\lambda u\in G,$$

et donc,

$$\Rr^n=F\oplus G.$$

Le projeté sur $F$ parallèlement à $G$ d'un vecteur $x=(x_1,...,x_n)$ est

$$\frac{1}{n}\sum_{k=1}^{n}x_k.u=(\frac{1}{n}\sum_{k=1}^{n}x_k,...,\frac{1}{n}\sum_{k=1}^{n}x_k)$$

et le projeté du même vecteur sur $G$ parallèlement à $F$ est

$$x-(\frac{1}{n}\sum_{k=1}^{n}x_k).u =(x_1-\frac{1}{n}\sum_{k=1}^{n}x_k,...,x_n-\frac{1}{n}\sum_{k=1}^{n}x_k).$$
\fincorrection
\correction{005185}
On a
$$n=\mbox{dim }E=\mbox{dim }(\mbox{Ker }f+\mbox{Ker }g)=\mbox{dim }(\mbox{Ker }f)+\mbox{dim }(\mbox{Ker }g)-\mbox{dim }(\mbox{Ker }f\cap\mbox{Ker }g),$$
mais aussi,

$$n=\mbox{dim }(\mbox{Im }f)+\mbox{dim }(\mbox{Im }g)-\mbox{dim }(\mbox{Im }f\cap\mbox{Im }g)=2n-\mbox{dim }\mbox{Ker }f-
\mbox{dim }(\mbox{Ker }g)-\mbox{dim }(\mbox{Im }f\cap\mbox{Im }g).$$
Par suite,

$$n+\mbox{dim }(\mbox{Ker }f\cap\mbox{Ker }g)=\mbox{dim }(\mbox{Ker }f)+\mbox{dim }\mbox{Ker }g=n-\mbox{dim }(\mbox{Im }f\cap
\mbox{Im }g)$$
puis $n+\mbox{dim }(\mbox{Ker }f\cap\mbox{Ker }g)=n-\mbox{dim }(\mbox{Im }f\cap
\mbox{Im }g)\Rightarrow\mbox{dim }(\mbox{Ker }f\cap\mbox{Ker }g)+\mbox{dim }(\mbox{Im }f\cap\mbox{Im }g)=0$ ou
encore $\mbox{dim }(\mbox{Ker }f\cap\mbox{Ker }g)=\mbox{dim }(\mbox{Im }f\cap\mbox{Im }g)=0$, et finalement,
$\mbox{Ker }f\cap\mbox{Ker }g=\mbox{Im }f\cap\mbox{Im }g=\{0\}$. Ceci montre que les sommes proposées sont directes.

\fincorrection
\correction{005565}
\textbf{1ère solution.} $F$ est le noyau d'une forme linéaire non nulle sur $E$ et est donc un hyperplan de $E$.

Soit $x=(x_1,...,x_n)$ un élément de $F\cap G$. Il existe $\lambda\in\Kk$ tel que $x=(\lambda,...,\lambda)$ et $n\lambda=0$ et donc $\lambda=0$ puis $x=0$. Donc $F\cap G=\{0\}$. De plus $\text{dim}(F)+\text{dim}(G)=n-1+1=n=\text{dim}(E)<+\infty$ et donc $F\oplus G=E$.

Soit $x=(x_1,...,x_n)$ un vecteur de $E$. Soit $\lambda\in \Kk$. $x -(\lambda,...,\lambda)\in F\Leftrightarrow (x_1-\lambda)+...+(x_n-\lambda)=0\Leftrightarrow\lambda=\frac{1}{n}\sum_{i=1}^{n}x_i$. Le projeté de $x$ sur $G$ parallèlement à $F$ est donc $\frac{1}{n}\sum_{i=1}^{n}x_i(1,...,1)$ et
le projeté de $x$ sur $F$ parallèlement à $G$ est $x-\frac{1}{n}\sum_{i=1}^{n}x_i(1,...,1)$.

\textbf{2ème solution} (dans le cas où $\Kk=\Rr$). On munit $\Rr^n$ de sa structure euclidienne canonique. Posons $\overrightarrow{u}=(1,\ldots,1)$.

On a $F=\overrightarrow{u}^\bot=G^\bot$. Par suite, $F$ est le supplémentaire orthogonal de $F$.

Soit $x\in E$. Le projeté orthogonal de $x$ sur $G$ est $\frac{x.u}{\|u\|^2}u=\frac{x_1+\ldots+x_n}{n}(1,\ldots,1)$.
\fincorrection
\correction{006871}
\begin{enumerate}
  \item Si les deux droites vectorielles sont distinctes alors elles engendrent un plan vectoriel 
et donc pas $\Rr^3$ tout entier. Si elles sont confondues c'est pire : elles n'engendrent qu'une droite.
Dans tout les cas elles n'engendrent pas $\Rr^3$ et ne sont donc pas supplémentaires.
  \item Si $P$ et $P'$ sont deux plans vectoriels alors $P\cap P'$ est une droite vectorielle si $P \neq P'$
ou le plan $P$ tout entier si $P=P'$. Attention, tous les plans vectoriels ont une équation du type 
$ax+by+cz=0$ et doivent passer par l'origine, il n'existe donc pas deux plans parallèles par exemple.
Donc l'intersection $P\cap P'$ n'est jamais réduite au vecteur nul. Ainsi $P$ et $P'$ ne sont pas supplémentaires.

  \item Soit $D$ une droite et $P$ un plan, $u$ un vecteur directeur de $D$. 
Si le vecteur $u$ appartient au plan $P$ alors $D\subset P$ et les espaces ne sont pas suppl\'ementaires 
(ils n'engendrent pas tout $\Rr^3$). 
Si $u \notin P$ alors d'une part $D\cap P$ est juste le vecteur nul
d'autre part $D$ et $P$ engendrent tout $\Rr^3$ ; $D$ et $P$ sont supplémentaires.

Détaillons un exemple : si $P$ est le plan d'équation $z=0$ alors il est engendré par les deux vecteurs $v=(1,0,0)$ et $w=(0,1,0)$.
Soit $D$ une droite de vecteur directeur $u=(a,b,c)$.

Alors  $u \notin P \iff u \notin \text{Vect}\{v,w\} \iff c \neq 0$.
Dans ce cas on bien que d'une part que $D = \text{Vect}\{ u\}$ intersecté avec $P$ est réduit au vecteur nul.
Ainsi $D\cap P = \{(0,0,0)\}$.
Et d'autre part tout vecteur $(x,y,z)\in \Rr^3$ appartient à $D+P = \text{Vect}\{u,v,w\}$.
Il suffit de remarquer que $(x,y,z) - \frac zc (a,b,c) = (x-\frac{za}{c},y-\frac{zb}{c},0) = (x-\frac{za}{c}) (1,0,0) + 
(y-\frac{zb}{c})(0,1,0)$. Et ainsi $(x,y,z)= \frac zc u + (x-\frac{za}{c}) v + 
(y-\frac{zb}{c}) w$. Donc $D+P = \Rr^3$.

Bilan on a bien $D\oplus P = \Rr^3$ : $D$ et $P$ sont en somme directe.

\end{enumerate}
\fincorrection
\correction{007414}
\begin{enumerate}
	\item On v\'{e}rifie ais\'{e}ment que $F$ et $G$ sont des sous-ensembles non vides de $E$ stables par combinaison lin\'{e}aire.

	\item 
	\begin{itemize}
\item L'intersection des sous-espaces $F$ et $G$ est triviale:
si $f\in F\cap G$, $f\in F$ implique que $f$ est une fonction polynomiale de degr\'{e} 1, soit $x\mapsto ax+b$, pour certains r\'{e}els $a$ et $b$,
mais alors on a $f(0)=b$ et $f'(0)=a$ donc $f\in G$ implique $a=b=0$ et finalement $f=0$.\\
\item Il reste \`{a} v\'{e}rifier que $E$ est la somme de $F$ et $G$. Comme $F$ et $G$ sont des sous-espaces vectoriels de $E$ on a \'{e}videmment $F+G\subset E$.
Pour l'inclusion inverse, soit $f\in E$, soit $g$ la fonction d\'{e}finie sur $\Rr$ par $x\mapsto f'(0)x+f(0)$.
Soit $h=f-g$. On a $g\in F$, et $h\in G$ puisque la d\'{e}rivation est lin\'{e}aire; de plus, $f=g+h$.
On a donc montr\'{e} $E=F+G$. 
\end{itemize}
Des deux points pr\'{e}c\'{e}dents, il r\'{e}sulte que $E=F\oplus G$.
\end{enumerate}
\fincorrection
\correction{000979}
$\hbox{det }\begin{pmatrix}  1 & -1 & 1 \cr
 1 & 1  & 0 \cr
 1 & 0 & -1 \cr \end{pmatrix}= 3\not = 0$  donc
la famille  $\mathcal{B} =\{ \begin{pmatrix}1 \cr 1 \cr 1 \cr
\end{pmatrix} ,
 \begin{pmatrix}-1 \cr 1 \cr 0\cr \end{pmatrix}  ,  \begin{pmatrix}1 \cr 0 \cr -1 \cr \end{pmatrix}\}$  est une base
de   ${\R}^3$.

$ \begin{pmatrix}1 \cr 0 \cr 0 \cr
\end{pmatrix}=\frac{1}{3}\begin{pmatrix}1 \cr 1 \cr 1 \cr
\end{pmatrix} -\frac{1}{3}
 \begin{pmatrix}-1 \cr 1 \cr 0\cr \end{pmatrix} +\frac{1}{3} \begin{pmatrix}1 \cr 0 \cr -1 \cr \end{pmatrix}$. Ses coordonn\' ees
 dans  $\mathcal{B} $  sont donc  $(1/3 ,  -1/3  ,  1/3)$.

$\begin{pmatrix}0 \cr 0 \cr 1 \cr \end{pmatrix}
=\frac{1}{3}\begin{pmatrix}1 \cr 1 \cr 1 \cr
\end{pmatrix}-\frac{1}{3}
 \begin{pmatrix}-1 \cr 1 \cr 0\cr \end{pmatrix} -\frac{2}{3} \begin{pmatrix}1 \cr 0 \cr -1 \cr \end{pmatrix}$.
 Ses coordonn\' ees
 dans  $\mathcal{B} $  sont donc  $(1/3 ,  -1/3  ,  -2/3)$.


$ \begin{pmatrix}1 \cr 0 \cr 1 \cr \end{pmatrix}= \begin{pmatrix}1
\cr 0 \cr 0 \cr \end{pmatrix}+ \begin{pmatrix}0 \cr 0 \cr 1 \cr
\end{pmatrix}$. Donc ses coordonn\' ees dans  $\mathcal{B} $  sont
$(2/3 ,  -2/3  ,  -1/3)$.
\fincorrection
\nocorrection
\correction{000981}
\begin{enumerate}
  \item Pour montrer que la famille $\{ v_1, v_2, v_3\}$ est une base nous allons 
montrer que cette famille est libre et génératrice.


  \begin{enumerate}
    \item Montrons que la famille $\{ v_1, v_2, v_3\}$ est libre.
Soit une combinaison linéaire nulle $a v_1+b v_2 + c v_3 = 0$, nous devons montrer qu'alors
les coefficients $a,b,c$ sont nuls. Ici le vecteur nul est $0=(0,0,0)$

\begin{align*}
  &  a v_1+b v_2 + c v_3 = (0,0,0) \\
\iff & a (0,1,1) + b(1,0,1) + c(1,1,0) = (0,0,0) \\
\iff & (b+c,a+c,a+b)=(0,0,0) \\
\iff & \begin{cases}
       b+c = 0 \\
       a+c = 0 \\
       a+b = 0  \\
       \end{cases} 
\iff \begin{cases}
       a = 0 \\
       b = 0 \\
       c = 0  \\
       \end{cases} \\
\end{align*}
Ainsi les coefficients vérifient $a=b=c=0$, cela prouve que la famille est libre.

    \item Montrons que la famille $\{ v_1, v_2, v_3\}$ est génératrice. Pour n'importe quel
vecteur $v=(x,y,z)$ de $\Rr^3$ on doit trouver $a,b,c\in\Rr$ tels que $a v_1+b v_2 + c v_3 = v$.
\begin{align*}
  &  a v_1+b v_2 + c v_3 = v \\
\iff & a (0,1,1) + b(1,0,1) + c(1,1,0) = (x,y,z) \\
\iff & (b+c,a+c,a+b)=(x,y,z) \\
\iff & \begin{cases}
       b+c = x \\
       a+c = y \quad (L_2) \\
       a+b = z \quad (L_3) \\
       \end{cases} 
\iff  \begin{cases}
       b+c = x \quad (L_1')\\
       a+c = y \\
       b-c = z-y  \quad(L_3')=(L_3-L_2)\\
       \end{cases} \\
\iff & \begin{cases}
       2b = x+z-y \quad (L_1'+L_3') \\
       a+c = y \\
       2c = x-(z-y) \quad  (L_1'-L_3')  \\
       \end{cases} 
\iff  \begin{cases}
       a = \frac12 (-x+y+z) \\
       b = \frac12 (x-y+z) \\
       c = \frac12 (x+y-z) \\
       \end{cases} \\
\end{align*}


Pour $a= \frac12 (-x+y+z)$, $b= \frac12 (x-y+z)$, $c = \frac12 (x+y-z)$ nous avons donc
la relation $av_1+bv_2+cv_3=(x,y,z)=v$. Donc la famille $\{ v_1, v_2, v_3\}$ est génératrice.


    \item La famille est libre et génératrice  donc c'est une base.

    \item Pour écrire  $w=(1,1,1)$ dans la base $(v_1,v_2,v_3)$  on peut résoudre le système correspondant
à la relation $a v_1+b v_2 + c v_3 = w$.
Mais en fait nous l'avons déjà résolu pour tout vecteur $(x,y,z)$, en particulier pour le vecteur
$(1,1,1)$ la solution est $a=\frac 12$, $b=\frac 12$, $c=\frac 12$.
Autrement dit $\frac 12 v_1+\frac 12 v_2 + \frac 12 v_3 = w$. Les coordonnées de $w$ dans la base $(v_1,v_2,v_3)$
sont donc $(\frac12,\frac12,\frac12)$.

  \end{enumerate} 


  \item Pour montrer que la famille est libre et génératrice les calculs 
sont similaires à ceux de la question précédente. Notons $\mathcal{B}$ la base $(v_1,v_2,v_3)$.

Exprimons ensuite $e_1$ dans cette base, les calculs donnent : $e_1 = \frac 13 v_1-\frac 13 v_2 +\frac 13 v_3$. 
Ses coordonnées dans la base $\mathcal{B}$ sont $(\frac13,-\frac13,\frac13)$.


$e_2 = \frac 13 v_1+\frac 23 v_2 +\frac 13 v_3$. 
Ses coordonnées dans $\mathcal{B}$ sont $(\frac13,\frac23,\frac13)$.

$e_3 = \frac 13 v_1-\frac 13 v_2 -\frac 23 v_3$. 
Ses coordonnées dans  $\mathcal{B}$ sont $(\frac13,-\frac13,-\frac23)$.


Les calculs sont ensuite terminés, on remarque que $w=(1,2,-3)$ vaut en fait
$w=e_1+2e_2-3e_3$ donc par nos calculs précédents $w=\frac 13 v_1-\frac 13 v_2 +\frac 13 v_3
+2(\frac 13 v_1+\frac 23 v_2 +\frac 13 v_3)-3(\frac 13 v_1-\frac 13 v_2 -\frac 23 v_3)
=  2v_2 +3 v_3$. Les coordonnées de $w$ dans  $\mathcal{B}$ sont $(0,2, 3)$.

  \item Par exemple la famille $\{(1,0,0),(0,1,0)\}$ est libre dans $\Rr^3$ mais pas g\'en\'eratrice.

  \item La famille  $\{(1,0,0),(0,1,0),(0,0,1),(1,1,1)\}$ est g\'en\'eratrice dans $\Rr^3$ mais pas libre.
\end{enumerate}
\fincorrection
\nocorrection
\nocorrection
\nocorrection
\correction{000985}
\begin{enumerate}
  \item Faux. Par exemple dans $\Rr^3$, $x=(1,0,0)$, $y=(0,1,0)$, $z=(1,1,0)$.
  \item Vrai. Soit une combinaison lin\'eaire  nulle   
$\lambda_1x_1+\cdots \lambda_p x_p =0.$
Supposons qu'un des coefficient est non nul: par exemple $\lambda_1 \neq 0$.
Alors on \'ecrit $x_1 = -\frac{\lambda_2}{\lambda_1} x_2-\cdots - \frac{\lambda_p}{\lambda_1}x_p.$
Donc $x_1$ est une combinaison lin\'eaire de $\{x_2,\ldots,x_p\}$. Ce qui contredit l'hypoth\`ese de l'\'enonc\'e, donc  tous les coefficients sont nuls. Donc $\{x_1,\ldots,x_p\}$ est une famille libre. 
\end{enumerate}
\fincorrection
\nocorrection
\correction{000987}
\begin{enumerate}
  \item C'est une base.
  \item Ce n'est pas une base : $v_3=4v_1-v_2$. Donc l'espace $\mathrm{Vect} (v_1,v_2,v_3)=\mathrm{Vect}(v_1,v_2)$.
  \item Ce n'est pas une base : $v_3=5v_1-4v_2$. Donc l'espace $\mathrm{Vect} (v_1,v_2,v_3)=\mathrm{Vect}(v_1,v_2)$.
\end{enumerate}
\fincorrection
\nocorrection
\nocorrection
\nocorrection
\nocorrection
\correction{000992}
\begin{enumerate}
  \item Tout d'abord la famille $\{P_0,P_1,P_2,\ldots,P_n\}$ contient $n+1$ vecteurs 
dans l'espace $E=\Rr_n[X]$ de dimension $n+1$. Ici un vecteur est un polynôme :
$P_0$ est un polynôme constant non nul, $P_1$ est un polynôme de degré exactement $1$,...
Rappelons que lorsque le nombre de vecteurs égal la dimension de l'espace  nous avons les équivalences, entre 
\emph{être une famille libre} et \emph{être une famille génératrice} et donc aussi \emph{être une base}.

\medskip

Nous allons donc montrer que $\{P_0,P_1,\ldots,P_n\}$ est une famille libre.
Soit une combinaison linéaire nulle :
$$\lambda_0 P_0+\lambda_1 P_1 + \cdots + \lambda_n P_n = 0.$$

Introduisons l'hypothèse concernant les degrés :
$\deg P_0 = 0$, $\deg P_1=1$, \ldots, $\deg P_n=n$.
Définissons le polynôme $P(X) = \lambda_0 P_0+\lambda_1 P_1 + \cdots + \lambda_n P_n$.

\medskip

Nous allons montrer successivement $\lambda_n=0$ puis $\lambda_{n-1}=0$,\ldots, $\lambda_0=0$.

Par l'absurde supposons $\lambda_n \neq 0$ et écrivons 
$P_n(X) = a_n X^n + a_{n-1}X^{n-1}+\cdots+a_1X+a_0$, comme
$\deg P_n(X)=n$ alors $a_n\neq0$.
Maintenant 
$P(X)$ est aussi un polynôme de degré exactement $n$ qui s'écrit
$$P(X) = \lambda_n \cdot a_n \cdot X^n + \text{termes de plus bas degré}$$
La combinaison linéaire nulle implique que $P(X)=0$ (le polynôme nul).
Donc en identifiant les coefficients devant $X^n$ on obtient $\lambda_n\cdot a_n=0$
On obtient  $a_n=0$ ou $\lambda_n=0$. Ce qui est une contradiction.
Conclusion $\lambda_n=0$.

\medskip

Maintenant la combinaison linéaire nulle s'écrit $\lambda_0 P_0+\lambda_1 P_1 + \cdots + \lambda_{n-1} P_{n-1}=0$.
Par récurrence descendante on trouve $\lambda_{n-1}=0$, \ldots, jusqu'à $\lambda_0=0$.

Bilan : $\lambda_0=0$, \ldots, $\lambda_n=0$ donc la famille $\{P_0,P_1,\ldots,P_n\}$ est libre,
elle donc aussi génératrice ; ainsi $\{P_0,P_1,\ldots,P_n\}$ est une base de $E=\Rr_n[X]$.

\medskip

Un point que nous avons utilisé et qu'il est peut-être utile de détailler est le suivant : 
si un polynôme égal le polynôme nul alors tous ces coefficients sont nul.

Voici une justification : écrivons $a_nX^n+ a_{n-1}X^{n-1}+\cdots+a_1X+a_0=0$
et divisons par $X^n$ :
$$a_n + \frac {a_{n-1}}{X} + \frac{a_{n-2}}{X^2} + \cdots + \frac{a_1}{X^{n-1}} + \frac{a_0}{X^n} = 0$$
Lorsque l'on fait tendre $X$ vers $+\infty$ alors le terme de gauche tend vers $a_n$ et celui de droite vaut $0$ donc 
par unicité de la limite $a_n=0$.
On fait ensuite une récurrence descendante pour prouver $a_{n-1}=0$,\ldots, $a_0=0$.

Une conséquence est que si deux polynômes sont égaux alors leurs coefficients sont égaux.
Et une autre formulation est de dire que $\{1,X,X^2,\ldots,X^n\}$ est une base de $\Rr_n[X]$.

  \item On trouve $a=10, b= -10, c = -7, d= -8$.
Puis $\alpha=-3,\beta=4,\gamma=-9,\delta=8$.
\end{enumerate}
\fincorrection
\nocorrection
\nocorrection
\nocorrection
\correction{000996}

Quand le nombre de vecteurs égal la dimension de l'espace  nous avons les équivalences, entre 
\emph{être une famille libre} et \emph{être une famille génératrice} et donc aussi \emph{être une base}.

Trois vecteurs dans $\Rr^3$ forment donc une base si et seulement s'ils forment une famille libre.
Vérifions quand c'est le cas.


\begin{align*}
      & a (1, 0, t) + b (1, 1, t) + c (t,0,1) = (0,0,0) \\
\iff  & (a+b+tc,b,at+bt+c)=(0,0,0) \\
\iff  & \begin{cases}
         a+b+tc = 0 \\
         b = 0 \\
         at+bt+c = 0 \\
        \end{cases} 
\iff \begin{cases}
         b = 0 \\
         a+tc = 0 \\
         at+c = 0 \\
        \end{cases} \\
\iff & \begin{cases}
         b = 0 \\
         a=-tc \\
         (-tc)t+c = 0 \\
        \end{cases} 
\iff \begin{cases}
         b = 0 \\
         a=-tc \\
         (t^2-1)c = 0 \\
        \end{cases} \\
\end{align*}

Premier cas : si $t\neq \pm 1$. Alors $t^2-1 \neq 0$ et donc
la seule solution du système est $(a=0,b=0,c=0)$.
Dans ce cas la famille est libre et est donc aussi une base.


Deuxième cas : si $t=\pm1$. Alors la dernière ligne du système disparaît
et il existe des solutions non triviales (par exemple si $t=1$, $(a=1,b=0,c=-1)$ est une solution).
La famille n'est pas libre et n'est donc pas une base. 
\fincorrection
\nocorrection
\nocorrection
\nocorrection
\nocorrection
\nocorrection
\nocorrection
\nocorrection
\nocorrection
\nocorrection
\correction{001006}
\begin{enumerate}
  \item C'est bien une base. Comme nous avons trois vecteurs et nous souhaitons 
montrer qu'ils forment un base d'un espace vectoriel de dimension $3$, 
il suffit de montrer que soit la famille est libre, soit elle est génératrice
(ces conditions sont équivalentes pour $n$ vecteurs dans un espace vectoriel de dimension $n$).

Il est plus simple de montrer que la famille est libre. Soit une combinaison linéaire nulle
$a v_1+b v_2+c v_3=0$  il faut montrer que $a=b=c=0$. Mais attention ici le corps de base est $K=\Cc$
donc $a,b,c$ sont des nombres complexes.

\begin{align*}
     &  a v_1+b v_2+c v_3=0 \\
\iff & a (1,-1,i) + b (-1,i,1) + c (i,1,-1) = (0,0,0) \\
\iff & (a-b+ic,-a+ib+c,ia+b-c)=(0,0,0) \\
\iff & \begin{cases}
        a-b+ic  = 0 \\
        -a+ib+c = 0  \\
        ia+b-c = 0 \\
       \end{cases} \\
\iff & \cdots \text{ on résout le système } \\
\iff & a=0, b=0, c=0 \\
\end{align*}

La famille $(v_1,v_2,v_3)$ est libre, donc aussi génératrice ; c'est donc une base de $\Cc^3$.




  \item On cherche $a,b,c \in \Cc$ tels que $a v_1+b v_2+c v_3=v$.
Il s'agit donc de r\'esoudre le syst\`eme :
$$ \begin{cases} 
a  -b +ic = 1+i \\ 
-a + ib  +c = 1-i \\ 
ia+b-c = i \\
\end{cases}$$
On trouve $a=0$, $b=\frac12(1-i)$, $c=\frac12(1-3i)$.
Nous avons donc $v = \frac12(1-i) v_2 + \frac12(1-3i) v_3$ et ainsi
 les coordonn\'ees de $v$ dans la base $(v_1,v_2,v_3)$ sont
$(0,\frac12(1-i),\frac12(1-3i))$.
\end{enumerate}
\fincorrection
\nocorrection
\nocorrection
\nocorrection
\nocorrection
\nocorrection
\nocorrection
\nocorrection
\correction{003317}
$$\left\{\begin{array}{lllllll} 
x'  &= &     & & 2y &+& z  \cr
                 3y' &= & -x  & &    &+& z  \cr
                 3z' &= & -x  &+& 3y &+& z. \cr
       \end{array}\right.$$
\fincorrection
\correction{003318}
$r = 3$,\quad $2\vec a - 3\vec b + 5\vec c = \vec 0$,
                 \quad $\vec b - 2\vec d - \vec e = \vec 0$.
\fincorrection
\nocorrection
\correction{003320}
3. $\pi_H$ : $$\left\{
  \begin{array}{lllllll}
  4x' &=& 3x &-&  y &-& z  \cr
                     4y' &=& -x &+& 3y &-& z  \cr
                     4z' &=&-2x &-& 2y &+&2z, \cr
                     \end{array}\right.$$  \qquad
    $s_H$   : $$\left\{
    \begin{array}{lllllll}
    		2x' &=& x  &-&  y &-& z  \cr
                     2y' &=&-x  &+&  y &-& z  \cr
                     2z' &=&-2x &-& 2y.       \cr
                     \end{array}\right.$$
 \fincorrection
\nocorrection
\correction{005179}
\begin{enumerate}
\item  Soit $n$ un entier naturel supèrieur ou égal à $2$.

Si $\sqrt{n}\in\Qq$, il existe $(a,b)\in(\Nn^*)^2$ tel que $\sqrt{n}=\frac{a}{b}$ ou encore tel que $n.b^2=a^2$. Mais
alors, par unicité de la décomposition d'un entier naturel supèrieur ou égal à 2 en facteurs premiers, tous les
facteurs premiers de $n$ ont un exposant pair ce qui signifie exactement que $n$ est un carré parfait.

Si $n=0$ ou $n=1$, $\sqrt{n}\in\Qq$ et $n$ est d'autre part un carré parfait. On a montré que~:

$$\forall n\in\Nn,\;(\sqrt{n}\in\Qq\Rightarrow n\;\mbox{est un carré parfait})$$

ou encore par contraposition

$$\forall n\in\Nn,\;(n\;\mbox{n'est pas un carré parfait}\Rightarrow\sqrt{n}\notin\Qq).$$

\item  D'après 1), $\sqrt{2}$, $\sqrt{3}$ et $\sqrt{6}$ sont irrationnels.

$E=\mbox{Vect}_\Qq(1,\sqrt{2},\sqrt{3},\sqrt{6})$ et donc, $E$ est un $\Qq$-espace
vectoriel et $(1,\sqrt{2},\sqrt{3},\sqrt{6})$ en est une famille génératrice.

Montrons que cette famille est $\Qq$-libre.

Soit $(a,b,c,d)\in\Qq^4$.

\begin{align*}
a+b\sqrt{2}+c\sqrt{3}+d\sqrt{6}=0&\Rightarrow(a+d\sqrt{6})^2=(-b\sqrt{2}-c\sqrt{3})^2
\Rightarrow a^2+2ad\sqrt{6}+6d^2=2b^2+2bc\sqrt{6}+3c^2\\
 &\Rightarrow a^2-2b^2-3c^2+6d^2=2(-ad+bc)\sqrt{6}
\end{align*}

Puisque $\sqrt{6}\notin\Qq$, on obtient $a^2-2b^2-3c^2+6d^2=2(-ad+bc)=0$ (car si $bc-ad\neq0$,
$\sqrt{6}=\frac{a^2-2b^2-3c^2+6d^2}{2(-ad+bc)}\in\Qq$) ou encore,

$$\left\{
\begin{array}{l}
a^2-3c^2=2b^2-6d^2\quad(1)\\
ad=bc\quad(2)
\end{array}
\right..$$

De même,

\begin{align*}
a+b\sqrt{2}+c\sqrt{3}+d\sqrt{6}=0&\Rightarrow(a+c\sqrt{3})^2=(-b\sqrt{2}-d\sqrt{6})^2
\Rightarrow(a^2+2ac\sqrt{3}+3c^2=2b^2+4bd\sqrt{3}+6d^2\\
 &\Rightarrow\left\{
\begin{array}{l}
a^2+3c^2=2b^2+6d^2\quad(3)\\
ac=2bd\quad(4)
\end{array}
\right..
\end{align*}

(puisque $\sqrt{3}$ est irrationnel). En additionnant et en retranchant (1) et (3), on obtient $a^2=2b^2$ et
$c^2=2d^2$.Puisque $\sqrt{2}$ est irrationnel, on ne peut avoir $b\neq0$ (car alors $\sqrt{2}=\pm\frac{a}{b}\in\Qq$) 
ou $d\neq0$. Donc, $b=d=0$ puis  $a=c=0$. Finalement, la famille $(1,\sqrt{2},\sqrt{3},\sqrt{6})$ est $\Qq$-libre et
est donc une base de E.
\end{enumerate}
\fincorrection
\correction{005572}
\textbf{Première solution.} Chaque $P_k$, $0\leqslant k\leqslant n$, est de degré $k+n-k=n$ et est donc dans $\Rr_n[X]$.

Les polynômes $P_k$, $0\leqslant k\leqslant n$ ont des valuations deux à deux distinctes et donc constituent une famille libre. Comme de plus $\text{card}(P_k)_{0\leqslant k\leqslant n}=n+1=\text{dim}(E)<+\infty$, la famille $(P_k)_{0\leqslant k\leqslant n}$ est une base de $E$.

\textbf{Deuxième solution.} La matrice carrée $M$ de la famille $(P_k)_{0\leqslant k\leqslant n}$ dans la base canonique de $\Rr_n[X]$ est triangulaire inférieure. Ses coefficients diagonaux sont tous non nuls car égaux à $1$. $M$ est donc inversible et $(P_k)_{0\leqslant k\leqslant n}$ est une base de $E$.
\fincorrection
\correction{005573}
\textbf{Unicité.} Soit $i\in\llbracket0,n\rrbracket$. $L_i$ doit admettre les $n$ racines deux à deux distinctes $a_j$ où $j$ est différent de $i$ et donc $L_i$ est divisible par le polynôme $\prod_{j\neq i}^{}(X-a_j)$.
$L_i$ doit être de degré $n$ et donc il existe un réel non nul $\lambda$ tel que $L_i=\lambda
\prod_{j\neq i}(X-a_j)$. Enfin $L_i(a_i)=1$ fournit $\lambda=\frac{1}{\prod_{j\neq i}^{}(a_i-a_j)}$. Ainsi nécessairement $L_i=\prod_{j\neq i}^{}\frac{X-a_j}{a_i-a_j}$.

\textbf{Existence.} Les $L_i$ ainsi définis conviennent.

\begin{center}
\shadowbox{
$\forall i\in\llbracket0,n\rrbracket,\;L_i=\prod_{j\neq i}^{}\frac{X-a_j}{a_i-a_j}$.
}
\end{center}

Montrons que la famille $(L_i)_{0\leqslant i\leqslant n}$ est libre.

Soient $\lambda_0$,..., $\lambda_n$ $n+1$ nombres complexes tels que $\lambda_0L_0 + ... +\lambda_nL_n =0$. En particulier, pour un indice $i$ de $\llbracket0,n\rrbracket$ donné, $\sum_{j=0}^{n}\lambda_jL_j(a_i)=0$ et donc $\lambda_i=0$ au vu des égalités définissant les $L_j$. La famille $(L_i)_{0\leqslant i\leqslant n}$ est libre.

De plus les $L_i$ sont tous dans $\Cc_n[X]$ et vérifient $\text{card}(L_i)_{0\leqslant i\leqslant n}= n+1=\text{dim}\Cc_n[X]<+\infty$. Donc la famille $(L_i)_{0\leqslant i\leqslant n}$ est une base de $\Cc_n[X]$.

Soit $P$ un polynôme quelconque de degré inférieur ou égal à $n$.

On écrit $P$ dans la base $(L_j)_{0\leqslant j\leqslant n}$ : $P=\sum_{j=0}^{n}\lambda_jL_j$. 
En prenant la valeur en $a_i$, $i$ donné dans $\llbracket0,n\rrbracket$, on obtient $\lambda_i= P(a_i)$. D'où l'écriture générale d'un polynôme de degré inférieur ou égal à $n$ dans la base $(L_i)_{0\leqslant i\leqslant n}$ : 

\begin{center}
\shadowbox{
$\forall P\in\Cc_n[X]$, $P=P(a_0)L_0 + ... + P(a_n)L_n$.
}
\end{center}

Mais alors : $(\forall i\in\llbracket0,n\rrbracket,\;P(a_i)=b_i)\Rightarrow P=b_0L_0 + ... + b_nL_n$.

 
Réciproquement le polynôme $P= b_0L_0+ ... +b_nL_n$ vérifie bien sûr les égalités demandées et est de degré inférieur ou égal à $n$.

Ainsi, il existe un et un seul polynôme de degré inférieur ou égal à $n$ vérifiant $\forall i\in\llbracket0,n\rrbracket$, $P(a_i)=b_i$ à savoir $P_0=\sum_{i=0}^{n}b_iL_i$.

Soient $P\in\Cc[X]$ et $R=(X-a_0) ... (X-a_n)$ ($\text{deg}(R)=n+1$).

\begin{align*}\ensuremath
(\forall i\in\llbracket0,n\rrbracket\;P(a_i) = b_i)&\Leftrightarrow(\forall i\in\llbracket0,n\rrbracket\;P(a_i)=P_0(a_i))\\
 &\Leftrightarrow P-P_0\;\text{admet les}\;n+1\;\text{racines deux à deux distinctes}\;a_0,...,a_n\\
  &\Leftrightarrow P-P_0\;\text{est divisible par}\;R\Leftrightarrow\exists Q\in\Cc[X]/\;P= P_0+QR.
\end{align*}

Les polynômes cherchés sont les $P_0+QR$ où $Q$ décrit $\Cc[X]$.
\fincorrection
\nocorrection
\correction{001015}
\begin{enumerate}
  \item
$F\cap G$ est un sous-espace vectoriel de $E$ donc est de dimension finie.
Soit $(e_1,\ldots, e_k)$ une base de $F\cap G$ avec $k=\dim F\cap G$.

$(e_1,\ldots, e_k)$ est une famille libre dans $F$ donc on peut la compl\'eter en une base de $F$ 
par le th\'eor\`eme de la base incompl\`ete.
Soient donc $(f_1,\ldots,f_\ell)$ des vecteurs de $F$ tels que 
$(e_1,\ldots, e_k,f_1,\ldots,f_\ell)$ soit une base de $F$.
Nous savons que $k+\ell = \dim F$.
Remarquons que les vecteurs $f_i$ sont dans $F\setminus G$ (car ils sont dans $F$ mais pas dans $F\cap G$).

Nous repartons de la famille $(e_1,\ldots, e_k)$ mais cette fois nous la compl\'etons en une base de $G$ : soit donc $(g_1,\ldots,g_m)$ des vecteurs de $G$ tels que $(e_1,\ldots, e_k,g_1,\ldots,g_m)$ soit une base de $G$.
Nous savons que $k+m = \dim G$.
Remarquons que cette fois les vecteurs $g_i$ sont dans $G\setminus F$.

\item Montrons que $\mathcal{B}=(e_1,\ldots, e_k,f_1,\ldots,f_\ell,g_1,\ldots,g_m)$ est une base de $F+G$.

C'est une famille g\'en\'eratrice car $F=\mathrm{Vect}(e_1,\ldots, e_k,f_1,\ldots,f_\ell) \subset \mathrm{Vect}(\mathcal{B})$ et $G=\mathrm{Vect}(e_1,\ldots, e_k,g_1,\ldots,g_m) \subset \mathrm{Vect}(\mathcal{B})$. Donc $F+G \subset  \mathrm{Vect}(\mathcal{B})$.

C'est une famille libre : en effet soit une combinaison lin\'eaire nulle
$$a_1e_1+\cdots+ a_ke_k\quad + \quad b_1f_1+\cdots +b_\ell f_\ell \quad + \quad c_1 g_1+\cdots +c_m g_m=0.$$
Notons $e=a_1e_1+\ldots +a_ke_k$, $f=b_1f_1+\cdots +b_\ell f_\ell$, $g=c_1 g_1+\cdots +c_m g_m$.
Donc la combinaison lin\'eaire devient :
$$e+f+g=0.$$
Donc $g=-e-f$, or $e$ et $f$ sont dans $F$ donc $g$ appartient \`a $F$.
Or les vecteurs $g_i$ ne sont pas dans $F$. Donc 
$g=c_1 g_1+\cdots+ c_m g_m$ est n\'ecessairement le vecteur nul.
Nous obtenons $c_1 g_1+\cdots+ c_m g_m=0$ c'est donc une combinaison lin\'eaire nulle pour la famille libre 
$(g_1,\ldots,g_m)$. Donc tous les coefficients $c_1,\ldots,c_m$ sont nuls.

Le reste de l'\'equation devient $a_1e_1+\cdots +a_ke_k+b_1f_1+\cdots +b_\ell f_\ell=0$,
or $(e_1,\ldots, e_k,f_1,\ldots,f_\ell)$ est une base de $F$ donc tous les coefficients 
$a_1,\ldots,a_k,b_1,\ldots,b_\ell$ sont nuls.

Bilan : tous les coefficients sont nuls donc la famille est libre. Comme elle \'etait g\'en\'eratrice, c'est une base.


  \item Puisque $\mathcal{B}$ est une base de $F+G$ alors la dimension de $F+G$ est le nombre de vecteurs de la base $\mathcal{B}$:
$$\dim (F+G) = k + \ell + m.$$
Or $k=\dim F\cap G$, $\ell = \dim F - k$, $m=\dim G - k$, donc 
$$\dim (F+G) = \dim F + \dim G - \dim (F\cap G).$$

\end{enumerate}
\fincorrection
\correction{001016}
Soit $E$ un espace vectoriel de dimension $n$ et $F$ un sous-espace vectoriel de $E$.
Par l'absurde supposons que $F$ ne soit pas de dimension finie, 
alors il existe $v_1,\ldots,v_{n+1}$, $n+1$ vecteurs de $F$ lin\'eairement ind\'ependants dans $F$.
Mais ils sont aussi lin\'eairement ind\'ependants dans $E$.
Donc la dimension de $E$ est au moins $n+1$. Contradiction.

\bigskip

Deux remarques :
\begin{itemize}
  \item En fait on a m\^eme  montré que la dimension de $F$ est plus petite que la dimension de $E$.
  \item On a utilisé le r\'esultat suivant : si $E$ admet une famille libre \`a $k$ \'el\'ements alors la dimension de $E$ 
est plus grande que $k$ (ou est infini). Ce r\'esultat est une cons\'equence imm\'ediate du th\'eor\`eme de la base incompl\`ete.
\end{itemize}
\fincorrection
\nocorrection
\nocorrection
\correction{001019}
\begin{enumerate}
  \item $G$  est engendr\'e par deux vecteurs donc $\dim G \leq 2$. Clairement  $v_4$ et  $v_5$  ne sont pas li\'es donc
$\dim G \geq 2$  c'est-\`a-dire  $\dim G=2$. 

  \item $F$  est engendr\'e par trois vecteurs donc  $\dim F \leq 3$.
Un calcul montre que la famille $\{ v_1,  v_2, v_3 \} $  est libre, 
d'où $\dim F\geq 3$  et donc $\dim F= 3$.

  \item Essayons d'abord d'estimer la dimension de $F\cap G$. 
D'une part $F\cap G \subset G$  donc  $\dim(F\cap G)\leq 2$.
Utilisons d'autre part la formule $\dim(F+G) =\dim F + \dim G - \dim(F\cap G)$. Comme  $F+G\subset  {\R}^4$, on a
$\dim(F+G)\leq 4$  d'o\`u on tire l'in\'egalit\' e  
$\dim(F\cap G) \ge 1$. Donc soit  $\dim(F\cap G)=1$
ou bien $\dim(F\cap G)=2$.

Supposons que  $\dim(F\cap G)$  soit \'egale \`a $2$. 
Comme  $F\cap G \subset G$  on aurait dans ce cas  $F\cap G =
G$ et donc $G\subset F$.  En particulier il existerait  $\alpha  , \beta , \gamma \in
{\R}$  tels que  $v_4=\alpha v_1 + \beta v_2 + \gamma v_3$. On v\'erifie ais\'ement que ce n'est pas le cas, 
ainsi $\dim (F\cap G)$  n'est pas \'egale \`a  $2$.
On peut donc conclure $\dim(F\cap G)=1$


  \item Par la formule $\dim (F+G) = \dim F + \dim G - \dim (F\cap G)$,
on obtient $\dim (F+G) = 2+3-1=4$. Cela entraîne $F+G=\Rr^4$.
\end{enumerate}
\fincorrection
\nocorrection
\nocorrection
\nocorrection
\nocorrection
\nocorrection
\nocorrection
\nocorrection
\correction{001027}
\begin{enumerate}
  \item Par la formule $\dim(G+H) = \dim(G)+\dim(H)-
\dim(G\cap H)$, on sait que $\dim(G+H) \leqslant \dim(G)+\dim(H)$.
Pour $G=\Im u$ et $H=\Im v$ on obtient :
$\dim (\Im u+\Im v) \leqslant \dim \Im u +\dim \Im v$.
Or $\Im (u+v) \subset \Im u+\Im v$.
Donc $\text{rg} (u + v) \leq \text{rg} (u) + \text{rg} (v)$.
  \item On applique la formule pr\'ec\'edente \`a $u+v$ et $-v$ :
$\text{rg}  ((u+v)+(-v)) \leqslant \text{rg}  (u+v)+\text{rg} (-v)$, or $\text{rg} (-v)=\text{rg} (v)$
donc $\text{rg} (u) \leqslant \text{rg} (u+v)+\text{rg} (v)$.
Donc $\text{rg} (u)-\text{rg} (v)\leqslant \text{rg} (u+v)$.
On recommence en \'echangeant $u$ et $v$ pour obtenir :
$\left|\text{rg} (u) - \text{rg} (v)\right| \leq \text{rg} (u + v)$.
\end{enumerate}
\fincorrection
\nocorrection
\nocorrection
\nocorrection
\nocorrection
\nocorrection
\nocorrection
\nocorrection
\nocorrection
\nocorrection
\nocorrection
\nocorrection
\nocorrection
\correction{003322}
codim$\,H = 0 :$ supplémentaire = $\{\vec 0\}$.\par
codim$\,H = p :$ Soit $\vec u \in E \setminus (H \cup K)$ :
Alors $H \oplus  K\vec u$ et $K \oplus  K\vec u$
ont un supplémentaire commun, $L$, donc $H$ et $K$ ont un supplémentaire
commun : $L \oplus  K\vec u$.
\fincorrection
\correction{005183}
\textbullet~$e_4$ et $e_5$ ne sont clairement pas colinéaires. Donc $(e_4,e_5)$ est une famille libre et
$\mbox{dim }G=\mbox{rg }(e_4,e_5)=2$.
Ensuite, puisque $e_1$ et $e_2$ ne sont pas colinéaires, on a $2\leq\mbox{dim }F\leq3$.
Soit alors $(\lambda,\mu,\nu)\in\Rr^3$.

$$\lambda e_1+\mu e_2+\nu e_3=0\Rightarrow
\left\{
\begin{array}{l}
\lambda+\mu+2\nu=0\quad(1)\\
2\lambda+\mu+\nu=0\quad(2)\\
3\lambda+\mu+\nu=0\quad(3)\\
4\lambda+3\mu+\nu=0\quad(4)
\end{array}
\right.
\Rightarrow
\left\{
\begin{array}{l}
\lambda=0\;((3)-(2))\\
\nu-\lambda=0\;((1)-(2))\\
\lambda+\mu+2\nu=0\;(1)
\end{array}
\right.
\Rightarrow\lambda=\mu=\nu=0
.$$
On a montré que~:~
$\forall(\lambda,\mu,\nu)\in\Rr^3,\;(\lambda e_1+\mu e_2+\nu e_3=0\Rightarrow\lambda=\mu=\nu=0)$.
$(e_1,e_2,e_3)$ est donc libre et $\mbox{dim }F=\mbox{rg }(e_1,e_2,e_3)=3$.
\textbullet~Comme $F\subset F+G$, $\mbox{dim }(F+G)\geq3$ ou encore $\mbox{dim }(F+G)=3$ ou $4$. De plus~:

$$\mbox{dim }(F+G)=3\Leftrightarrow F=F+G\Leftrightarrow G\subset F\Leftrightarrow\{e_4,e_5\}\subset F.$$
On cherche alors $(\lambda,\mu,\nu)\in\Rr^3$ tel que $e_4=\lambda e_1+\mu e_2+\nu e_3$ ce qui fournit le
système~:

$$
\left\{
\begin{array}{l}
\lambda+\mu+2\nu=-1\quad(1)\\
2\lambda+\mu+\nu=0\quad(2)\\
3\lambda+\mu+\nu=-1\quad(3)\\
4\lambda+3\mu+\nu=2\quad(4)
\end{array}
\right..$$
$(3)-(2)$ fournit $\lambda=-1$ puis $(1)-(2)$ fournit $\nu=-2$ puis $(2)$ fournit $\mu=4$.
Maintenant, $(4)$ n'est pas vérifiée car $4\times(-1)+3\times4-2=6\neq 2$. Le système proposé n'admet pas de solution et donc
$e_4\notin\mbox{Vect}(e_1,e_2,e_3)=F$. Par suite, $\mbox{dim }(F+G)=4$.
Enfin,

$$\mbox{dim }(F\cap G)=\mbox{dim }F+\mbox{dim }G-\mbox{dim }(F+G)=3+2-4=1.$$

\begin{center}
\shadowbox{
$\text{dim}(F)=3$, $\text{dim}(G)=2$, $\text{dim}(F+G)=4$ et $\text{dim}(F\cap G)=1$.
}
\end{center}

\fincorrection
\correction{005184}
On a $H_1\subset H_1+H_2$ et donc $\mbox{dim }(H_1+H_2)\geq n-1$ ou encore $\mbox{dim }(H_1+H_2)\in\{n-1,n\}$. Donc

$$\mbox{dim }(H_1\cap H_2)=\mbox{dim }H_1+\mbox{dim }H_2-\mbox{dim }(H_1+H_2)=
\left\{
\begin{array}{l}
(n-1)+(n-1)-(n-1)=n-1\\
\quad\mbox{ou}\\
(n-1)+(n-1)-n=n-2
\end{array}
\right..$$
Maintenant, si $\mbox{dim }(H_1+H_2)=n-1=\mbox{dim }H_1=\mbox{dim }H_2$, alors $H_1=H_1+H_2=H_2$ et donc en particulier,
$H_1=H_2$. Réciproquement, si $H_1=H_2$ alors $H_1+H_2=H_1$ et $\mbox{dim }(H_1+H_2)=n-1$.
En résumé, si $H_1$ et $H_2$ sont deux hyperplans distincts, $\mbox{dim }(H_1\cap H_2)=n-2$ et bien sûr, si 
$H_1=H_2$, alors $\mbox{dim }(H_1\cap H_2)=n-1$.
Si $n=2$, les hyperplans sont des droites vectorielles et l'intersection de deux droites vectorielles distinctes du
plan vectoriel est de dimension $0$, c'est-à-dire réduite au vecteur nul.
Si $n=3$, les hyperplans sont des plans vectoriels et l'intersection de deux plans vectoriels distincts de l'espace de
dimension $3$ est une droite vectorielle.

\fincorrection
\correction{005575}
Soit $f$ l'application de $F\times G$ dans $E$ qui à un élément $(x,y)$ de $F\times G$ associe $x+y$.

$f$ est clairement linéaire et d'après le thèorème du rang 

\begin{center}
$\text{dim}(F\times G)=\text{dim}(\text{Ker}f)+\text{dim}(\text{Im}f)$ avec $\text{dim}(FxG)=\text{dim}F +\text{dim}G$ et $\text{dim}(\text{Im}f)=\text{dim}(F+G)$.
\end{center}

Il reste à analyser $\text{Ker}f$.

Soit $(x,y)\in E^2$. $(x,y)$ est élément de $\text{Ker}f$ si et seulement si $x$ est dans $F$, $y$ est dans $G$ et $x+y=0$ ou encore si et seulement si $x$ et $y$ sont dans $F\cap G$ et $y=-x$. Donc $\text{Ker}f=\{(x,-x),\;x\in F\cap G\}$.

Montrons enfin que $\text{Ker}f$ est isomorphe à $F\cap G$. Soit $\varphi$ l'application de $F\cap G$ dans $\text{Ker}f$ qui à l'élément $x$ de $F\cap G$ associe $(x,-x)$ dans $\text{Ker}f$. $\varphi$ est clairement une application linéaire, clairement injective et clairement surjective. Donc $\varphi$ est un isomorphisme de $F\cap G$ sur $\text{Ker}f$ et en particulier $\text{dim}(\text{Ker}f)=\text{dim}(F\cap G)$. Finalement

\begin{center}
\shadowbox{
$\text{dim}(F+G)=\text{dim}F+\text{dim}G-\text{dim}(F\cap G)$.
}
\end{center}
\fincorrection
\correction{005576}
\begin{align*}\ensuremath
\text{dim}(F+G+H)&=\text{dim}((F+G)+H)=\text{dim}(F+G)+\text{dim}H -\text{dim}((F+G)\cap H)\\
 &=\text{dim}F+\text{dim}G+\text{dim}H -\text{dim}(F\cap G)-\text{dim}((F+G)\cap H).
\end{align*}

Maintenant , $F\cap H+G\cap H\subset (F+G)\cap H$ (car si $x$ est dans $F\cap H+G\cap H$ il existe $y$ dans $F$ et dans $H$ et $z$ dans $G$ et dans $H$ tel que $x = y+z$ et $x$ est bien dans $F+G$ et aussi dans $H$). Donc

\begin{align*}\ensuremath
\text{dim}((F+G)\cap H)&\geqslant\text{dim}(F\cap H + G\cap H)=\text{dim}(F\cap H)+\text{dim}(G\cap H)-\text{dim}((F\cap H)\cap (G\cap H))\\
 &=\text{dim}(F\cap H)+\text{dim}(G\cap H)-\text{dim}(F\cap G\cap H)
\end{align*}

et finalement

\begin{center}
\shadowbox{
$\text{dim}(F+G+H)\leqslant\text{dim}F+\text{dim}G+\text{dim}H-\text{dim}(F\cap G)-\text{dim}(F\cap H)-\text{dim}(G\cap H)+\text{dim}(F\cap G\cap H)$.
}
\end{center}

Le cas de trois droites vectorielles de $\Rr^2$ deux à deux distinctes fournit un cas d'inégalité stricte.
\fincorrection
\correction{005577}
Montrons par récurrence que $\forall n\geqslant2$, $\text{dim}(F_1+\ldots+F_n)\leqslant\text{dim}(F_1)+\ldots+\text{dim}(F_n)$.

\textbullet~Pour $n=2$, $\text{dim}(F_1+F_2)=\text{dim}(F_1)+\text{dim}(F_2)-\text{dim}(F_1\cap F_2)\leqslant\text{dim}(F_1)+\text{dim}(F_2)$.

\textbullet~Soit $n\geqslant2$. Supposons que si $F_1$,\ldots, $F_n$ sont $n$ sous-espaces de $E$, $\text{dim}(F_1+\ldots+F_n)\leqslant\text{dim}(F_1)+\text{dim}(F_n)$.

Soient $F_1$,\ldots, $F_{n+1}$ $n+1$ sous-espaces de $E$.

\begin{align*}\ensuremath
\text{dim}(F_1+F_2+...+F_{n+1})&\leqslant\text{dim}(F_1+...+F_n) +\text{dim}(F_{n+1})\;(\text{d'après le cas}\;n=2)\\
 &\leqslant \text{dim}(F_1)+ ... +\text{dim}(F_{n+1})\;(\text{par hypothèse de récurrence}).
\end{align*} 

Le résultat est démontré par récurrence.

On sait que si la somme $F_1+\ldots+F_n$ est directe, on a $\text{dim}(F_1+\ldots+F_n)=\text{dim}(F_1)+\ldots+\text{dim}(F_n)$.

Montrons par récurrence que $\forall n\geqslant,\;2\;[\text{dim}(F_1+...+F_n)=\text{dim}F_1 + ... + \text{dim}F_n]\Rightarrow\text{la somme}\;F_1+\ldots+F_n\;\text{est directe}]$.

\textbullet~Pour n=2 , d'après le \ref{ex:rou14}, $\text{dim}(F_1+F_2)=\text{dim}(F_1)+\text{dim}(F_2)\Rightarrow\text{dim}(F_1\cap F_2)=0\Rightarrow F_1\cap F_2=\{0\}$.

\textbullet~Soit $n\geqslant2$. Soient $F_1$,...,$F_{n+1}$ $n+1$ sous-espaces de $E$ tels que $\text{dim}(F_1+...+F_{n+1})=\text{dim}(F_1)+...+\text{dim}(F_{n+1})$.

On sait que

\begin{align*}\ensuremath
\text{dim}(F_1)+...+\text{dim}(F_{n+1})&=\text{dim}(F_1+...+F_{n+1})\\ 
 &=\text{dim}(F_1+...+F_n) +\text{dim}(F_{n+1})-\text{dim}((F_1+...+F_n)\cap F_{n+1})\\
 &\leqslant\text{dim}(F_1)+...+\text{dim}(F_{n+1})-\text{dim}((F_1+...+F_n)\cap F_{n+1}),
\end{align*}

 
et donc $\text{dim}((F_1+...+F_n)\cap F_{n+1})\leqslant0$ puis $\text{dim}((F_1+...+F_n)\cap F_{n+1})= 0$. Par suite $(F_1+...+F_n)\cap F_{n+1}=\{0\}$ et aussi $\text{dim}(F_1)+...+\text{dim}(F_{n+1})=\text{dim}(F_1+...+F_n) +\text{dim}(F_{n+1})$ et donc $\text{dim}(F_1+...+F_n)=\text{dim}(F_1)+ ... +\text{dim}(F_n)$.

Mais alors, par hypothèse de récurrence, la somme $F_1+...+F_n$ est directe et si l'on rappelle que 
$(F_1+...+F_n)\cap F_{n+1}=\{0\}$, on a montré que la somme $F_1+...+F_{n+1}$ est directe.

Le résultat est démontré par récurrence.
\fincorrection
\correction{005578}
Soit $n\geqslant3$. Montrons par récurrence que $\forall k\in\llbracket2,n-1\rrbracket$, si $H_1$,..., $H_k$ sont $k$ hyperplans de $E$, alors $\text{dim}(H_1\cap...\cap H_k)\geqslant n-k$.

\textbullet~Pour $k=2$. Soient $H_1$ et $H_2$ deux hyperplans de $E$.

$\text{dim}(H_1\cap H_2)=\text{dim}(H_1)+\text{dim}(H_2)-\text{dim}(H_1+H_2)\geqslant(n-1)+(n-1)-n=n-2$.

\textbullet~Soit $k\in\llbracket2,n-3\rrbracket$. Supposons que la dimension d'une intersection de $k$ hyperplans de $E$ soit supérieure ou égale à $n-k$.

Soient $H_1$,..., $H_k$, $H_{k+1}$ $k+1$ hyperplans de $E$.

\begin{center}
$\text{dim}(H_1\cap...\cap H_k\cap H_{k+1})=\text{dim}(H_1\cap...\cap H_k)+\text{dim}(H_{k+1})-\text{dim}((H_1\cap...\cap H_k)+H_{k+1})\geqslant(n-k)+(n-1)-n=n-(k+1)$,
\end{center}

ce qui démontre le résultat par récurrence.

Pour $k=n-1$, on obtient en particulier $\text{dim}(H_1\cap...\cap H_{n-1})\geqslant n-(n-1)=1>0$ et donc $H_1\cap...\cap H_{n-1}\neq\{0\}$.
\fincorrection
\correction{005579}
Si $m=n$, c'est immédiat.

Supposons $m<n$.

\begin{align*}\ensuremath
r&=\text{dim}(\text{Vect}(x_1,...,x_n))=\text{dim}\left(\text{Vect}(x_1,...,x_m)+\text{Vect}(x_{m+1},...x_n)\right)\\
 &\leqslant\text{dim}(\text{Vect}(x_1,...,x_m))+\text{dim}(\text{Vect}(x_{m+1},...,x_n))\\
&\leqslant s + (n-m)
\end{align*}

et donc $s\geqslant r+m-n$. On a l'égalité si et seulement si chaque inégalité est une égalité, c'est à dire si et seulement si
$\text{Vect}(x_1,...,x_m)\cap\text{Vect}(x_{m+1},...x_n)=\{0\}$ (pour la première) et la famille $(x_{m+1},...,x_n)$ est libre (pour la deuxième).
\fincorrection
\correction{005580}
$\text{Im}(f+g)=\{f(x)+g(x),\;x\in E\}\subset\{f(x)+g(x'),\;(x,x')\in E^2\}=\text{Im}f+\text{Im}g$. Donc 

\begin{center}
$\text{rg}(f+g)\leqslant\text{dim}(\text{Im}f+\text{Im}g)\leqslant\text{rg}f+\text{rg}g$
\end{center}

puis $\text{rg}f=\text{rg}((f+g)+(-g))\leqslant\text{rg}(f+g)+\text{rg}(-g)=\text{rg}(f+g)+\text{rg}g$ (car $\text{Im}(-g)=\{-g(x),\;x\in E\}=\{g(-x),\;x\in E\}=\{g(x'),\;x'\in E\}=\text{Im}g$) et donc $\text{rg}(f+g)\geqslant\text{rg}f-\text{rg}g$. De même , en échangeant les rôles de $f$ et $g$, $\text{rg}(f+g)\geqslant\text{rg}g-\text{rg}f$ et finalement $\text{rg}(f+g)\geqslant|\text{rg}f -\text{rg}g|$.

\begin{center}
\shadowbox{
$\forall(f,g)\in\mathcal{L}(E,F)^2,\;|\text{rg}f -\text{rg}g|\leqslant\text{rg}(f+g)\leqslant\text{rg}f+\text{rg}g$.
}
\end{center}
\fincorrection
\correction{005581}
$\text{Im}(g\circ f)=g(f(E))\subset g(F)$ fournit $\text{rg}(gof)\leqslant\text{rg}g$.

Soit $g'=g_{/f(E)}$. D'après le théorème du rang, on a 

\begin{center}
$\text{rg}f=\text{dim}(f(E))=\text{dim}\text{Ker}g'+\text{dim}\text{Im}g'\geqslant\text{dim}\text{Im}g')=\text{rg}(g\circ f)$
\end{center}

et donc $\text{rg}(g\circ f)\leqslant\text{Min}\{\text{rg}f,\text{rg}g\}$.

A partir du théorème du rang, on voit que l'inégalité $\text{rg}f+\text{rg}g-\text{dim}F\leqslant\text{rg}(g\circ f)$ est équivalente à l'inégalité $\text{dim}(\text{Ker}(g\circ f))\leqslant\text{dim}\text{Ker}f+ \text{dim}\text{Ker}g$.

Soit $f'=f_{/\text{Ker}(g\circ f)}$. D'après le théorème du rang, $\text{dim}(\text{Ker}(g\circ f))=\text{dim}\text{Ker}f'+\text{dim}\text{Im}f'$. Mais $\text{Ker}f'\subset\text{Ker}f$ puis $\text{Im}f'=\{f(x)/\;x\in E\;\text{et}\;g(f(x))= 0\}\subset\{y\in F/\;g(y) = 0\}=\text{Ker}g$ et finalement $\text{dim}\text{Ker}(g\circ f)\leqslant\text{dim}\text{Ker}f+\text{dim}\text{Ker}g$.

\begin{center}
\shadowbox{
$\forall(f,g)\in\mathcal{L}(E,F)\times\mathcal{L}(F,G),\;\text{rg}f+\text{rg}g-\text{dim}F\leqslant\text{rg}(g\circ f)\leqslant\text{Min}\{\text{rg}f,\text{rg}g\}$.
}
\end{center}
\fincorrection
\nocorrection
\nocorrection
\correction{000929}
\begin{enumerate}
  \item $f_1$ est linéaire. Pour $(x,y) \in \Rr^2$ et $(x',y')\in \Rr^2$ :
\begin{align*}
f_1\big( (x,y)+(x',y') \big) 
  & = f_1\big( x+x',y+y' \big) \\
  & = \big(2(x+x')+ (y+y'), (x+x')-(y+y') \big) \\
  & = \big(2x+y +2x'+y', x-y+x'-y' \big) \\
  & = \big(2x+y,x-y \big) + \big(2x'+y',x'-y' \big) \\
  & = f_1(x,y)+f_1(x',y')
\end{align*}

Pour $(x,y)\in \Rr^2$ et $\lambda\in \Rr$ :
$$f_1\big( \lambda \cdot (x,y)\big) = f_1\big(\lambda x,\lambda y\big)
 = \big( 2\lambda x+\lambda y,\lambda x - \lambda y \big) = \lambda \cdot\big(2x+y,x-y \big) 
= \lambda \cdot f_1(x,y).$$

  \item $f_2$ n'est pas lin\'eaire, en effet par exemple $f_2(1,1,0)+f_2(1,1,0)$ n'est pas \'egal \`a $f_2(2,2,0)$.

  \item $f_3$ est linéaire : il faut vérifier d'abord que pour tout $(x,y,z)$ et $(x',y',z')$ alors
$f_3\big( (x,y,z) + (x',y',z') \big) = f_3(x,y,z)+f_3(x',y',z')$. Et ensuite que pour tout $(x,y,z)$ et $\lambda$
on a $f_3\big(\lambda\cdot(x,y,z) \big) = \lambda \cdot f_3(x,y,z)$.


  \item $f_4$ est linéaire : il faut vérifier d'abord que pour tout $(x,y)$ et $(x',y')$ alors
$f_4\big( (x,y) + (x',y') \big) = f_4(x,y)+f_4(x',y')$. Et ensuite que pour tout $(x,y)$ et $\lambda$
on a $f_4\big(\lambda\cdot(x,y) \big) = \lambda \cdot f_4(x,y)$.

  \item $f_5$ est lin\'eaire : soient $P,P' \in \Rr_3[X]$ alors 

\begin{align*}
f_5\big(P+P'\big) 
  & = \big( (P+P')(-1), (P+P')(0), (P+P')(1) \big) \\
  & = \big( P(-1)+P'(-1), P(0)+P'(0), P(1)+P'(1) \big) \\
  & = \big( P(-1), P(0), P(1) \big)  + \big( P'(-1), P'(0), P'(1) \big) \\
  & = f_5(P)+f_5(P')  \\
\end{align*}

Et si $P\in \Rr_3[X]$ et $\lambda \in \Rr$ :
\begin{align*}
f_5\big( \lambda \cdot P\big) 
  & =  \big( (\lambda P)(-1), (\lambda P)(0), (\lambda P)(1) \big) \\
  & = \big( \lambda \times P(-1), \lambda \times P(0), \lambda \times P(1) \big) \\
  & = \lambda \cdot \big( P(-1), P(0), P(1) \big) \\
  & = \lambda \cdot f_5(P) \\
\end{align*}
\end{enumerate}
\fincorrection
\correction{000930}
Montrons que la famille  $\{ x,\phi(x),\phi^2(x), \ldots , \phi ^{n-1}(x) \}$ est
libre. Soient  $\lambda _0, \ldots , \lambda _{n-1} \in {\R}$ tels
que  $\lambda _0 x + \lambda_1\phi(x)+\cdots + \lambda _{n-1} \phi ^{n-1}(x)=0$.
Alors : $\phi ^{n-1} \big(\lambda _0 x+ \lambda_1\phi(x) + \cdots + \lambda _{n-1}
\phi ^{n-1}(x)\big)=0$. Mais comme de plus $\phi ^n=0 $, on a
l'\'egalit\'e $\phi ^{n-1} \big(\lambda _0 x+ \lambda_1\phi(x) + \cdots + \lambda
_{n-1} \phi ^{n-1}(x)\big)=\phi ^{n-1} (\lambda _0 x ) + \phi
^n \big(\lambda _1 x+ \cdots + \lambda _{n-1} \phi
^{n-2}(x)\big)= \phi ^{n-1}(\lambda _0x)=\lambda _0 \phi ^{n-1}(x)$. Comme  $\phi
^{n-1}(x) \not= 0$  on obtient  $\lambda _0=0$.

\noindent En calculant ensuite  $\phi ^{n-2}\big(\lambda _1 \phi
(x )+ \cdots + \lambda _{n-1} \phi ^{n-1}(x)\big)$  on obtient
$\lambda _1=0$  puis, de proche en proche,  $\lambda_2=2$,\ldots, $\lambda _{n-1}=0$. 
La famille  $\{ x, \phi(x), \ldots , \phi
^{n-1}(x) \}$  est donc libre. En plus elle compte  $n$  vecteurs, comme
$\dim E=n$  elle est libre et maximale et forme donc une base
de $E$.
\fincorrection
\nocorrection
\nocorrection
\nocorrection
\nocorrection
\nocorrection
\correction{003312}
3. $|a| \ne |b|$.
 \fincorrection
\correction{005170}
\begin{enumerate}
\item  Si $f$ existe alors nécessairement, pour tout $(x,y,z)\in\Rr^3$~:

$$f((x,y,z))=xf((1,0,0))+yf((0,1,0))+zf((0,0,1))=x(1,1)+y(0,1)+z(-1,1)=(x-z,x+y+z).$$

On en déduit l'unicité de $f$.

Réciproquement, $f$ ainsi définie vérifie bien les trois égalités de l'énoncé. Il reste donc à se convaincre que $f$ est
linéaire.

Soient $((x,y,z),(x',y',z'))\in(\Rr^3)^2$ et $(\lambda,\mu)\in\Rr^2$.

\begin{align*}
f(\lambda(x,y,z)+\mu(x',y',z'))&=f((\lambda x+\mu x',\lambda y+\mu y',\lambda z+\mu z'))\\
 &=((\lambda x+\mu x')-(\lambda z+\mu z'),(\lambda x+\mu x')+(\lambda y+\mu y')+(\lambda z+\mu
z'))\\
 &=\lambda(x-z,x+y+z)+\mu(x'-z',x'+y'+z')\\
 &=\lambda f((x,y,z))+\mu f((x',y',z')).
\end{align*}

$f$ est donc linéaire et convient. On en déduit l'existence de $f$. On a alors $f((3,-1,4))=(3-4,3-1+4)=(-1,6)$.

\textbf{Remarque.} La démonstration de la linéarité de $f$ ci-dessus est en fait superflue car le cours donne
l'expression générale d'une application linéaire de $\Rr^n$ dans $\Rr^p$.

\item  Détermination de $\mbox{Ker}f$.

Soit $(x,y,z)\in\Rr^3$.

\begin{align*}
(x,y,z)\in\Rr^3\Leftrightarrow f((x,y,z))=(0,0)\Leftrightarrow(x-z,x+y+z)=(0,0)\Leftrightarrow
\left\{
\begin{array}{l}
x-z=0\\
x+y+z=0
\end{array}
\right.\Leftrightarrow
\left\{
\begin{array}{l}
z=x\\
y=-2x
\end{array}
\right.
\end{align*}

Donc, $\mbox{Ker}f=\{(x,-2x,x),\;x\in\Rr\}=\{x(1,-2,1),\;x\in\Rr\}=\mbox{Vect}((1,-2,1))$.
La famille $((1,-2,1))$ engendre $\mbox{Ker}f$ et est libre. Donc, la famille $((1,-2,1))$ est une base de
$\mbox{Ker}f$.
Détermination de $\mbox{Im}f$.

Soit $(x',y')\in\Rr^2$.

\begin{align*}
(x',y')\in\mbox{Im}f&\Leftrightarrow\exists(x,y,z)\in\Rr^3/\;f((x,y,z))=(x',y')\\
 &\Leftrightarrow\exists(x,y,z)\in\Rr^3/\;
\left\{
\begin{array}{l}
x-z=x'\\
x+y+z=y'
\end{array}
\right.
\Leftrightarrow\exists(x,y,z)\in\Rr^3/\;
\left\{
\begin{array}{l}
z=x-x'\\
y=-2x+x'+y'
\end{array}
\right.\\
 &\Leftrightarrow\mbox{le système d'inconnue}\;(x,y,z)~:~\left\{
\begin{array}{l}
z=x-x'\\
y=-2x+x'+y'
\end{array}
\right.
\;\mbox{a au moins une solution.}
\end{align*}

Or, le triplet $(0,x'+y',-x')$ est solution et le système proposé admet une solution. Par suite, tout $(x',y')$ de
$\Rr^2$ est dans $\mbox{Im}f$ et finalement, $\mbox{Im}f=\Rr^2$.

Détermination d'un supplémentaire de $\mbox{Ker}f$.

Posons $e_1=(1,-2,1)$, $e_2=(1,0,0)$ et $e_3=(0,1,0)$ puis $F=\mbox{Vect}(e_2,e_3)$ et montrons que
$\Rr^3=\mbox{Ker}f\oplus F$.

Tout d'abord, $\mbox{Ker}f\cap F=\{0\}$. En effet~:

\begin{align*}
(x,y,z)\in\mbox{Ker}f\cap F&\Leftrightarrow\exists(a,b,c)\in\Rr^3/\;(x,y,z)=ae_1=be_2+ce_3\\
 &\Leftrightarrow\exists(a,b,c)\in\Rr^3/\;
\left\{
\begin{array}{l}
x=a=b\\
y=-2a=c\\
z=a=0
\end{array}
\right.
\Leftrightarrow x=y=z=0
\end{align*}

Vérifions ensuite que $\mbox{Ker}f+F=\Rr^3$.

\begin{align*}
(x,y,z)\in\mbox{Ker}f+F&\Leftrightarrow\exists(a,b,c)\in\Rr^3/\;(x,y,z)=ae_1+be_2+ce_3\\
 &\Leftrightarrow\exists(a,b,c)\in\Rr^3/\;
\left\{
\begin{array}{l}
a+b=x\\
-2a+c=y\\
a=z
\end{array}
\right.
\Leftrightarrow
\exists(a,b,c)\in\Rr^3/\;\left\{
\begin{array}{l}
a=z\\
b=x-z\\
c=y+2z
\end{array}
\right.
\end{align*}

Le système précédent (d'inconnue $(a,b,c)$) admet donc toujours une solution et on a montré que $\Rr^3=\mbox{Ker}f+F$.
Finalement, $\Rr^3=\mbox{Ker}f\oplus F$ et $F$ est un supplémentaire de $\mbox{Ker}f$ dans $\Rr^3$.

Vérifions enfin que $F$ est isomorphe à $\mbox{Im}f$. Mais, $F=\{(x,y,0),\;(x,y)\in\Rr^2\}$ et
$\begin{array}[t]{cccc}
\varphi~:&F&\rightarrow&\Rr^2\\
 &(x,y,0)&\mapsto&(x,y)
\end{array}$ est clairement un isomorphisme de $F$ sur $\mbox{Im}f(=\Rr^2)$.
\end{enumerate}
\fincorrection
\correction{005186}
 \begin{enumerate}
 \item  Si $P$ est un polynôme de degré inférieur ou égal à $n$, alors $P(X+1)-P(X)$ est encore un polynôme de degré
inférieur ou égal à $n$. Par suite, $\varphi$ est bien une application de $E$ dans lui-même.
Soient alors $(P,Q)\in E^2$ et $(\lambda,\mu)\in\Rr^2$.

\begin{align*}
\varphi(\lambda P+\mu Q)&=(\lambda P+\mu Q)(X+1)-(\lambda P+\mu
Q)(X)=\lambda(P(X+1)-P(X))+\mu(Q(X+1)-Q(X))\\
 &=\lambda\varphi(P)+\mu\varphi(Q).
\end{align*}
$\varphi$ est linéaire de $E$ vers lui-même et donc un endomorphisme de $E$.
 \item  Soit $P\in E$.
$P\in\mbox{Ker }\varphi\Leftrightarrow\forall x\in\Rr,\;P(x+1)=P(x)$. Montrons alors que $P$ est constant.
Soit $Q=P-P(0)$. $Q$ est un polynôme de degré inférieur ou égal à $n$ s'annulant en les entiers naturels 0, 1, 2,...
(car $P(0)=P(1)=P(2)=...$) et a ainsi une infinité de racines deux à deux distinctes. $Q$ est donc le polynôme nul ou
encore $\forall x\in\Rr,\;P(x)=P(0)$. Par suite, $P$ est un polynôme constant.
Réciproquement, les polynômes constants sont clairement dans $\mbox{Ker }\varphi$ et donc

$$\mbox{Ker }\varphi=\{\mbox{polynômes constants}\}=\Rr_0[X].$$
Pour déterminer $\mbox{Im }\varphi$, on note tout d'abord que si $P$ est un polynôme de degré inférieur ou égal à $n$,
alors $\varphi(P)=P(X+1)-P(X)$ est un polynôme de degré inférieur ou égal à $n-1$. En effet, si
$P=a_nX^n+\sum_{k=0}^{n-1}a_kX^k$ (avec $a_n$ quelconque, éventuellement nul) alors

\begin{align*}
\varphi(P)&=a_n((X+1)^n-X^n)+\mbox{termes de degré inférieur on égal à}\;n-1\\
 &=a_n(X^n-X^n)+\mbox{termes de degré
inférieur on égal à}\;n-1\\
 &=\mbox{termes de degré inférieur on égal à}\;n-1.
\end{align*}
Donc, $\mbox{Im }(\varphi)\subset\Rr_{n-1}[X]$. Mais d'après le théorème du rang,

$$\mbox{dim }\mbox{Im }(\varphi)=\mbox{dim }\Rr_n[X]-\mbox{dim }\mbox{Ker }(\varphi)=(n+1)-1=n=\mbox{dim }\Rr_{n-1}[X]
<+\infty,$$
et donc $\mbox{Im }\varphi=\Rr_{n-1}[X]$. (On peut noter que le problème difficile~\og~soit $Q\in\Rr_{n-1}[X]$.
Existe-t-il $P\in\Rr_n[X]$ tel que $P(X+1)-P(X)=Q$~?~\fg~a été résolu simplement par le théorème du rang.)
\end{enumerate}
\fincorrection
\correction{005188}
Soient $(z,z')\in\Cc^2$ et $(\lambda,\mu)\in\Rr^2$.

$$f(\lambda z+\mu z')=(\lambda z+\mu z')+a(\overline{\lambda z+\mu
z'})=\lambda(z+a{\bar z})+\mu(z'+a\overline{z'})=\lambda f(z)+\mu f(z').$$
$f$ est donc $\Rr$-linéaire. On note que $f(ia)=i(a-|a|^2)$ et que $if(a)=i(a+|a|^2)$. Comme $a\neq0$, on a
$f(ia)\neq if(a)$. $f$ n'est pas $\Cc$-linéaire.
Soit $z\in\Cc\setminus\{0\}$. Posons $z=re^{i\theta}$ où $r\in\Rr^*_+$ et $\theta\in\Rr$.

$$z\in\mbox{Ker }f\Leftrightarrow z+a{\bar z}=0\Leftrightarrow e^{i\theta}+ae^{-i\theta}=0\Leftrightarrow e^{2i\theta}=-a.$$
\textbf{1er cas.} Si $|a|\neq1$, alors, pour tout réel $\theta$, $e^{2i\theta}\neq-a$. Dans ce cas, $\mbox{Ker }f=\{0\}$ et
d'après le théorème du rang, $\mbox{Im }f=\Cc$.
\textbf{2ème cas.} Si $|a|=1$, posons $a=e^{i\alpha}$.

$$e^{2i\theta}=-a\Leftrightarrow
e^{2i\theta}=e^{i(\alpha+\pi)}\Leftrightarrow2\theta\in\alpha+\pi+2\pi\Zz\Leftrightarrow\theta\in\frac{\alpha+\pi}{2}+\pi\Zz.$$
Dans ce cas, $\mbox{Ker }f=\mbox{Vect}(e^{i(\alpha+\pi)/2})$. D'après le théorème du rang, $\mbox{Im }f$ est une droite
vectorielle et pour déterminer $\mbox{Im }f$, il suffit d'en fournir un vecteur non nul, comme par exemple $f(1)=1+a$.
Donc, si $a\neq-1$, $\mbox{Im }f=\mbox{Vect}(1+a)$. Si $a=-1$, $\forall z\in\Cc,\;f(z)=z-{\bar z}=2i\mbox{Im }(z)$ et
$\mbox{Im }f=i\Rr$.
\fincorrection
\correction{005189}
\begin{enumerate}
 \item  Pour $(x,y)\in\Rr^2$, posons $f((x,y))=(x',y')$.

\begin{center}
$f\in\mathcal{L}(\Rr^2)\Leftrightarrow\exists(\alpha,\beta,\gamma,\delta)\in\Rr^4/\;\forall(x,y)\in\Rr^2,\;\left\{
\begin{array}{l}
x'=\alpha x+\gamma y\\
y'=\beta x+\delta y
\end{array}
\right.$.
\end{center}
 \item  Avec les notations précédentes,

\begin{align*}
z'&=x'+iy'=(\alpha x+\gamma y)+i(\beta x+\delta y)
=(\alpha\frac{z+{\bar z}}{2}+\gamma\frac{z-{\bar z}}{2i})+i(\beta\frac{z+{\bar z}}{2}+\delta\frac{z-{\bar z}}{2i})\\
 &=\left(\frac{\alpha+\delta}{2}+i\frac{\beta-\gamma}{2}\right)z+\left(\frac{\alpha-\delta}{2}+i\frac{\beta+\gamma}{2}\right){\bar z}
=az+b{\bar z}
\end{align*}
où $a=\frac{\alpha+\delta}{2}+i\frac{\beta-\gamma}{2}$ et $b=\frac{\alpha-\delta}{2}+i\frac{\beta+\gamma}{2}$.
 \item  Réciproquement, si $z'=az+b{\bar z}$, en posant $a=a_1+ia_2$ et $b=b_1+ib_2$ où $(a_1,a_2,b_1,b_2)\in\Rr^4$, on
obtient~:

$$x'+iy'=(a_1+ia_2)(x+iy)+(b_1+ib_2)(x-iy)=(a_1+b_1)x+(-a_2+b_2)y+i\left((a_2+b_2)x+(a_1-b_1)y\right)$$
et donc,

$$\left\{
\begin{array}{l}
x'=(a_1+b_1)x+(b_2-a_2)y\\
y'=(a_2+b_2)x+(a_1-b_1)y
\end{array}
\right..$$
\end{enumerate}
\fincorrection
\nocorrection
\nocorrection
\nocorrection
\correction{000934}
\begin{enumerate}
  \item Aucun problème...
  \item Par d\'efinition de $f$ et de ce qu'est la somme de deux sous-espaces vectoriels, l'image est
$$\Im f =  \{ f(x_1,x_2) \mid x_1 \in E_1, x_2\in E_2 \} = \{ x_1+x_2 \mid x_1 \in E_1, x_2\in E_2 \} = E_1 + E_2.$$

Pour le noyau :
$$\ker f = \{ (x_1,x_2) \mid f(x_1,x_2)=0 \} = \{ (x_1,x_2) \mid x_1+x_2=0 \}$$

Mais on peut aller un peu plus loin. En effet un \'el\'ement $(x_1,x_2) \in \ker f$,
v\'erifie $x_1\in E_1$, $x_2\in E_2$ et $x_1=-x_2$. Donc $x_1 \in E_2$. Donc $x_1\in E_1 \cap E_2$.
R\'eciproquement si $x\in E_1 \cap E_2$, alors $(x,-x)\in \ker f$.
Donc 
$$\ker f = \{ (x,-x) \mid x \in  E_1 \cap E_2 \}. $$
De plus l'application $x \mapsto (x,-x)$ montre que $\ker f$ est isomorphe \`a $E_1 \cap E_2$.

\item Le th\'eor\`eme du rang s'\'ecrit :
$$\dim \ker f+ \dim \Im f = \dim (E_1\times E_2).$$
Compte tenu de  l'isomorphisme entre $\ker f$ et $E_1 \cap E_2$ on obtient :
$$\dim (E_1 \cap E_2) + \dim (E_1+E_2) = \dim (E_1\times E_2).$$
Mais $\dim (E_1\times E_2) = \dim E_1 + \dim E_2$, donc on retrouve ce que l'on appelle le th\'eor\`eme des quatre dimensions :
$$\dim (E_1+E_2) = \dim E_1+\dim E_2-\dim (E_1 \cap E_2).$$

\end{enumerate}
\fincorrection
\nocorrection
\nocorrection
\nocorrection
\nocorrection
\nocorrection
\nocorrection
\correction{000941}
Montrons ceci par r\'ecurence : Pour $n=1$, l'assertion est triviale
: $x\notin \ker \varphi \Rightarrow \varphi(x) \neq 0$. Supposons
que si $x \notin \ker \varphi$ alors  $\varphi^{n-1}(x) \neq 0$,
 ($n \geqslant 2$).
Fixons $x \notin \ker \varphi$,
 Alors par hypoth\`eses de r\'ecurrence
$\varphi^{n-1}(x) \neq 0$, mais $\varphi^{n-1}(x) = \varphi(
\varphi^{n-2}(x)) \in \mathop{\mathrm{Im}}\nolimits \varphi$ donc $\varphi^{n-1}(x) \notin
\ker \varphi$ gr\^ace \`a l'hypoth\`ese sur $\varphi$. Ainsi
$\varphi(\varphi^{n-1}(x)) \neq 0$, soit $\varphi^{n}(x)\neq 0$.
Ce qui termine la r\'ecurrence.
\fincorrection
\nocorrection
\correction{000943}
\begin{itemize}
  \item[(i) $\Rightarrow$ (ii)] Supposons
$\ker f = \Im f$. Soit $x\in E$, alors $f(x) \in \Im f$ donc $f(x)
\in \ker f$, cela entraîne $f(f(x)) = 0$ ; donc $f^2=0$. De plus
d'apr\`es la formule du rang $\dim \ker f + \text{rg}  (f) = n$, mais $\dim
\ker f = \dim \Im f = \text{rg}  f$, ainsi $2\text{rg}  (f)=n$.
  \item[(ii) $\Rightarrow$ (i)] Si $f^2 = 0$ alors
$\Im f \subset \ker f$ car pour $y\in \Im f$ il existe $x$ tel que
$y=f(x)$ et $f(y)=f^2(x)=0$. De plus si $2 \text{rg}  (f) = n$ alors la
formule du rang donne $\dim \ker f = \text{rg}  (f)$ c'est-\`a-dire $\dim \ker f =
\dim \Im f$. Nous savons donc que $\Im f$ est inclus dans $\ker f$
mais ces espaces sont de m\^eme dimension donc sont \'egaux :
$\ker f = \Im f$.
 \end{itemize}
\fincorrection
\nocorrection
\nocorrection
\nocorrection
\correction{000947}
On va montrer $g(\ker f) \subset \ker f$.
Soit $y \in g(\ker f)$. Il existe $x\in \ker f$ tel que $y=g(x)$.
Montrons $y\in \ker f$:
$$f(y)=f(g(x))=f\circ g (x) = g\circ f(x) = g(0)=0.$$

\bigskip

On fait un raisonnement similaire pour montrer $g(\Im f) \subset \Im f$.
Soit $z\in g(\Im f)$, il existe $y \in \Im f$ tel que $z=g(y)$. Il existe alors $x\in E$ tel que $y=f(x)$.
Donc 
$$z=g(y)=g(f(x))=g\circ f(x)= f\circ g (x) = f(g(x)) \in \Im f.$$

\fincorrection
\nocorrection
\correction{000949}
Pour montrer l'\'egalit\'e $\ker f \cap \mathop{\mathrm{Im}}\nolimits f = f(\ker f^2)$, nous
montrons la double inclusion.

Soit $y\in \ker f \cap \mathop{\mathrm{Im}}\nolimits f$, alors $f(y) = 0$ et il existe $x$
tel que $y=f(x)$. De plus $f^2(x) = f(f(x))=f(y) =0$ donc $x\in
\ker f^2$. Comme $y = f(x)$ alors $y \in f(\ker f^2)$. Donc $\ker
f \cap \mathop{\mathrm{Im}}\nolimits f \subset f(\ker f^2)$.


Pour l'autre inclusion, nous avons d\'ej\`a que $ f(\ker f^2) \subset
f(E) = \mathop{\mathrm{Im}}\nolimits f$. De plus $ f(\ker f^2) \subset \ker f$, car si $y\in
f(\ker f^2)$ il existe $x\in \ker f^2$ tel que $y=f(x)$, et
$f^2(x) = 0$ implique $f(y)=0$ donc $y\in \ker f$. Par cons\'equent
$ f(\ker f^2) \subset \ker f \cap \mathop{\mathrm{Im}}\nolimits f $.
\fincorrection
\nocorrection
\correction{000951}
\begin{enumerate}
  \item Par exemple $f(x,y)=(0,x)$ alors $\mathop{\mathrm{Ker}}\nolimits f=\mathop{\mathrm{Im}}\nolimits f = \{0\}\times \Rr = \{ (0,y) \mid y\in \Rr\}$.
  \item Par exemple l'identit\'e : $f(x,y)=(x,y)$. En fait un petit exercice est de montrer que les seules applications possibles sont les applications bijectives (c'est tr\`es particulier aux applications de $\Rr^2$ dans $\Rr^2$).
  \item L'application nulle : $f(x,y)=(0,0)$. Exercice : c'est la seule possible !
\end{enumerate}
\fincorrection
\nocorrection
\nocorrection
\correction{000954}
\begin{enumerate}
  \item Comment est d\'efinie $\phi$ \`a partir de la d\'efinition sur les \'el\'ements de la base ?
Pour  $x\in E$ alors $x$ s'\'ecrit dans la base  $\{e_1,e_2,e_3\}$, $x=\alpha_1e_1+\alpha_2e_2+\alpha_3e_3$. 
Et $\phi$ est d\'efinie sur $E$ par la formule
$$\phi(x)=\alpha_1 \phi(e_1) + \alpha_2 \phi(e_2) + \alpha_3 \phi(e_3).$$
Soit ici :
$$\phi(x) = (\alpha_1+\alpha_2+\alpha_3) e_1 + (\alpha_1-\alpha_2)e_2 + t\alpha_3e_3.$$

Cette d\'efinition rend automatiquement $\phi$ lin\'eaire (v\'erifiez-le si vous n'\^etes pas convaincu !).
  \item On cherche \`a savoir si $\phi$ est injective.
Soit $x\in E$ tel que $\phi(x)=0$ donc
$(\alpha_1+\alpha_2+\alpha_3) e_1 + (\alpha_1-\alpha_2)e_2 + t\alpha_3e_3=0$. Comme $\{e_1,e_2,e_3\}$ est une base alors tous les coefficients sont nuls :
$$\alpha_1+\alpha_2+\alpha_3=0, \quad \alpha_1-\alpha_2=0, \quad  t\alpha_3 = 0.$$
Si $t \neq 0$ alors en résolvant le syst\`eme on obtient $\alpha_1=0$, $\alpha_2=0$,
$\alpha_3=0$. Donc $x=0$ et $\phi$ est injective.

Si $t =0$, alors $\phi$ n'est pas injective, en résolvant le m\^eme syst\`eme on obtient
des solutions non triviales, par exemple $\alpha_1=1$, $\alpha_2=1$, $\alpha_3=-2$.
Donc pour $x= e_1+e_2-2e_3$ on obtient $\phi(x)=0$.

  \item Pour la surjectivité on peut soit faire des calculs, soit appliquer la formule du rang. 
Examinons cette deuxi\`eme m\'ethode. $\phi$ est surjective si et seulement si 
la dimension de $\Im \phi$ est \'egale
\`a la dimension de l'espace d'arriv\'ee (ici $E$ de dimension $3$).
Or on a une formule pour $\dim \Im \phi$ :
$$\dim \ker \phi + \dim \Im \phi = \dim E.$$
Si $t \neq 0$, $\phi$ est injective donc $\ker \phi = \{0\}$ est de dimension $0$.
Donc $\dim \Im \phi =3$ et $\phi$ est surjective.

Si $t = 0$ alors $\phi$ n'est pas injective donc $\ker \phi$ est de dimension au moins $1$
(en fait $1$ exactement), donc  $\dim \Im \phi \leqslant 2$. Donc $\phi$ n'est pas surjective.


On remarque que $\phi$ est injective si et seulement si elle est surjective.
Ce qui est un r\'esultat du cours pour les applications ayant l'espace 
de d\'epart et d'arriv\'ee de m\^eme dimension (finie).
\end{enumerate}
\fincorrection
\nocorrection
\correction{000956}
Calculer le noyau revient à résoudre un système linéaire,
et calculer l'image aussi. On peut donc tout faire ``à la main''.

Mais on peut aussi appliquer un peu de théorie ! Noyau et image sont liés par la formule du rang : 
$\dim \ker f + \dim \Im f = \dim E$
pour $f : E \to F$. Donc si on a trouvé le noyau alors on connaît la dimension de l'image.
Et il suffit alors de trouver autant de vecteur de l'image.


\begin{enumerate}
  \item $f_1$ est injective, surjective (et donc bijective).

  \begin{enumerate}
  \item Faisons tout à la main. Calculons le noyau :
\begin{align*}
(x,y) \in \ker f_1
  & \iff f_1(x,y) = (0,0) 
  \iff (2x+y,x-y) = (0,0) \\
  &\iff \begin{cases}
         2x+y=0 \\
        x-y=0 \\   
        \end{cases}
  \iff (x,y)=(0,0) \\
\end{align*}
Ainsi $\ker f_1 = \{ (0,0) \}$ et donc $f_1$ est injective.

   \item Calculons l'image. Quels éléments $(X,Y)$ peuvent s'écrire $f_1(x,y)$ ?
\begin{align*}
f_1(x,y) = (X,Y) 
  & \iff (2x+y,x-y) = (X,Y) \\
  &\iff \begin{cases}
         2x+y=X \\
         x-y=Y \\   
        \end{cases} 
  \iff \begin{cases}
         x=\frac{X+Y}{3} \\
         y=\frac{X-2Y}{3} \\   
        \end{cases} \\
  & \iff (x,y)=\left(\frac{X+Y}{3},\frac{X-2Y}{3}\right) \\
\end{align*}
Donc pour n'importe quel $(X,Y)\in\Rr^2$ on trouve un antécédent $(x,y)=(\frac{X+Y}{3},\frac{X-2Y}{3})$
qui vérifie donc $f_1(x,y)=(X,Y)$. Donc $\Im f_1 = \Rr^2$. Ainsi $f_1$ est surjective.

  \item Conclusion : $f_1$ est injective et surjective donc bijective.
  \end{enumerate}


  \item 
\begin{enumerate}
  \item Calculons d'abord le noyau :
\begin{align*}
(x,y,z) \in \ker f_2
  & \iff f_2(x,y,z) = (0,0,0) \\
  & \iff (2x+y+z,y-z,x+y) = (0,0,0) \\
  &\iff \begin{cases}
         2x+y+z = 0 \\
         y-z = 0 \\   
         x+y = 0 \\ 
        \end{cases} \\
  & \iff  \begin{cases}
         x = -z \\
         y = z \\   
        \end{cases} \\
  & \iff 
  \begin{pmatrix}x\\y\\z\end{pmatrix} = \begin{pmatrix}-z\\z\\z\end{pmatrix}\\ 
  &\iff 
  \begin{pmatrix}x\\y\\z\end{pmatrix} \in 
  \text{Vect} \begin{pmatrix}-1\\1\\1\end{pmatrix} 
  = \left\lbrace \lambda\begin{pmatrix}-1\\1\\1\end{pmatrix} \mid \lambda \in \Rr \right\rbrace\\ 
\end{align*}
Ainsi $\ker f_2 =  \text{Vect}(-1,1,1) $ et donc $f_2$ n'est pas injective.

  \item Maintenant nous allons utiliser que $\ker f_2 = \text{Vect}(-1,1,1)$ , autrement dit
$\dim  \ker f_2 = 1$.
La formule du rang, appliquée à $f_2 : \Rr^3 \to \Rr^3$ s'écrit 
$\dim \ker f_2 + \dim \Im f_2 = \dim \Rr^3$. Donc $\dim \Im f_2 = 2$.
Nous allons trouver une base de $\Im f_2$. Il suffit donc de trouver deux vecteurs linéairement indépendants.
Prenons par exemple
$v_1 = f_2(1,0,0) =  (2,0,1) \in\Im f_2$
et $v_2 = f_2(0,1,0) = (1,1,1) \in \Im f_2$. Par construction ces vecteurs sont dans l'image de $f_2$ et 
il est clair qu'ils sont linéairement indépendants. Donc $\{v_1,v_2\}$ est une base de $\Im f_2$.

  \item $f_2$ n'est ni injective, ni surjective (donc pas bijective).
  \end{enumerate}


  \item Sans aucun calcul on sait $f_3: \Rr^2 \to \Rr^4$ ne peut être surjective 
car l'espace d'arrivée est de dimension strictement
supérieur à l'espace de départ. 
  \begin{enumerate}
  \item Calculons le noyau :
\begin{align*}
(x,y) \in \ker f_3
  & \iff f_3(x,y) = (0,0,0,0) \\
  & \iff  (y,0,x-7y,x+y) = (0,0,0,0) \\
  & \iff \begin{cases}
         y = 0 \\
         0  = 0 \\   
         x-7y = 0 \\ 
         x+y = 0 \\ 
        \end{cases} \\
  & \iff \cdots \\ 
  & \iff (x,y)=(0,0) \\
\end{align*}
Ainsi $\ker f_3 = \{ (0,0) \}$ et donc $f_3$ est injective.

  \item La formule du rang, appliquée à $f_3 : \Rr^2 \to \Rr^4$ s'écrit 
$\dim \ker f_3 + \dim \Im f_3 = \dim \Rr^2$. Donc $\dim \Im f_3 = 2$.
Ainsi $\Im f_3$ est un espace vectoriel de dimension $2$ inclus dans $\Rr^3$,
 $f_3$ n'est pas surjective.

Par décrire $\Im f_3$ nous allons trouver deux vecteurs indépendants de $\Im f_3$.
Il y a un nombre infini de choix : prenons par exemple 
$v_1 = f(1,0) = (0,0,1,1)$. Pour $v_2$ on cherche (un peu à tâtons) un vecteur linéairement indépendant de $v_1$.
Essayons $v_2 = f(0,1)=(1,0,-7,1)$. Par construction $v_1,v_2 \in \Im f$ ; ils sont clairement linéairement indépendants
et comme $\dim \Im f_3=2$ alors $\{v_1,v_2\}$ est une base de $\Im f_3$.

Ainsi $\Im f_3= \text{Vect}\{v_1,v_2\} =\big\{ \lambda(0,0,1,1) + \mu (1,0,-7,1) \mid \lambda,\mu \in \Rr \big\}$.
  \end{enumerate}

  
  
  \item $f_4 : \Rr_3[X] \to \Rr^3$ va d'un espace de dimension $4$ vers un espace de dimension strictement plus petit
et donc $f_4$ ne peut être injective.

\begin{enumerate}
  \item Calculons le noyau. \'Ecrivons un polynôme $P$ de degré $\le 3$ sous la forme
$P(X)= aX^3+bX^2+cX+d$. Alors $P(0) = d$, $P(1)=a+b+c+d$, $P(-1)=-a+b-c+d$.
\begin{align*}
P(X) \in \ker f_4 
  & \iff \big( P(-1), P(0), P(1) \big)= (0,0,0)  \\
  & \iff (-a+b-c+d,d,a+b+c+d) = (0,0,0) \\
  &\iff \begin{cases}
          -a+b-c+d = 0 \\
          d = 0 \\   
          a+b+c+d = 0 \\ 
        \end{cases}\\
 & \iff \cdots \\
 & \iff \begin{cases}
          a = -c \\
          b = 0 \\  
          d = 0 \\ 
        \end{cases} \\
  & \iff (a,b,c,d)=(t,0,-t,0) \quad t\in\Rr \\
  \end{align*}

Ainsi le noyau $\ker f_4 = \big\{ tX^3 - tX \mid t\in \Rr\big\} = \text{Vect}\{ X^3-X \}$.
$f_4$ n'est pas injective son noyau étant de dimension $1$.

  \item La formule du rang pour $f_4 : \Rr_3[X] \to \Rr^3$ s'écrit 
$\dim \ker f_4 + \dim \Im f_4 = \dim \Rr_3[4]$. Autrement dit
$1+ \dim \Im f_4 = 4$. Donc $\dim \Im f_4 = 3$. 
Ainsi $\Im f_4$ est un espace de dimension $3$ dans $\Rr^3$ donc 
$\Im f_4=\Rr^3$. Conclusion $f_4$ est surjective.
  \end{enumerate}

\end{enumerate}
\fincorrection
\nocorrection
\nocorrection
\correction{000959}
\begin{enumerate}
  \item Soit $P \in E$ et $\lambda \in \Rr$, alors la division euclidienne de $AP$ par $B$ s'\'ecrit $AP= Q\cdot B + R$, donc en multipliant par $\lambda$ on obtient :
$A\cdot (\lambda P)= (\lambda Q)B+\lambda R$.
ce qui est la division euclidienne de $A\cdot (\lambda P)$ par $B$, 
donc si $f(P)= R$ alors $f(\lambda P)=\lambda R$. Donc $f(\lambda P)=\lambda f(P)$.

Soient $P, P' \in E$. On \'ecrit les divisions euclidiennes :
$$AP= Q\cdot B + R,\quad  AP'=Q'\cdot B+R'.$$
En additionnant :
$$A(P+P')=(Q+Q')B+(R+R')$$
qui est la division euclidienne de $A(P+P')$ par $B$.
Donc si $f(P)=R$, $f(P')=R'$ alors $f(P+P')=R+R'=f(P)+f(P')$.

Donc $f$ est lin\'eaire.
  \item Sens $\Rightarrow$. Supposons $f$ est bijective, donc en particulier $f$ est surjective,
    en particulier il existe $P\in E$ tel que $f(P)=1$ ($1$ est le polyn\^ome constant \'egale \`a $1$). La division euclidienne est donc $AP=BQ+1$, autrement dit 
$AP-BQ=1$. Par le th\'eor\`eme de B\'ezout, $A$ et $B$ sont premiers entre eux.

  \item Sens $\Leftarrow$. Supposons $A, B$ premiers entre eux. Montrons que $f$ est injective.
Soit $P\in E$ tel que $f(P)=0$. Donc la division euclidienne s'\'ecrit : $AP=BQ+0$.
Donc $B$ divise $AP$. Comme $A$ et $B$ sont premiers entre eux, par le lemme de Gauss,
alors $B$ divise $P$. Or $B$ est de degr\'e $n+1$ et $P$ de degr\'e moins que $n$, 
donc la seule solution est $P=0$. Donc $f$ est injective. 
Comme $f : E\longrightarrow E$ est injective et $E$ est de dimension finie, alors $f$ est bijective.
\end{enumerate}
\fincorrection
\nocorrection
\nocorrection
\nocorrection
\correction{000963}
\begin{enumerate}
    \item Montrons que si  $\phi $  est un isomorphisme,
l'image de toute base de  $E$  est une base de  $F$ : soit
$\mathcal{B} =\{ e_1, \ldots , e_n \} $  une base de  $E$  et
nommons  $\mathcal{B} ' $  la famille  $\{ \phi (e_1), \ldots ,
\phi (e_n) \} $.
    \begin{enumerate}
        \item $\mathcal{B} '$  est libre. Soient en effet  $\lambda _1 , \ldots , \lambda _n\in {\R}$  tels
que  $\lambda _1\phi (e_1)+ \cdots + \lambda _n \phi (e_n)
=0 $. Alors  $ \phi (\lambda _1e_1+ \cdots + \lambda _ne_n) =0
$  donc, comme  $\phi$  est injective,  $\lambda _1e_1+ \cdots
+ \lambda _ne_n=0$  puis, comme  $\mathcal{B} $  est libre,
$\lambda _1=\cdots =\lambda _n=0$.
        \item $\mathcal{B} '$  est g\' en\' eratrice. Soit  $y\in F$. Comme  $\phi $  est surjective, il
existe  $x\in E$  tel que  $y=\phi (x)$. Comme  $\mathcal{B}$
est g\' en\' eratrice, on peut choisir   $\lambda _1 , \cdots ,
\lambda _n\in {\R}$  tels que  $x=\lambda _1 e_1 +\cdots + \lambda
_n e_n $. Alors  $y=\lambda _1\phi (e_1)+ \cdots + \lambda _n
\phi (e_n)  $.
    \end{enumerate}
    \item Supposons que l'image par  $\phi $  de toute base de  $E$  soit une base  $F$. Soient  $\mathcal{B}
=\{ e_1,\ldots , e_n\} $  une base de  $E$  et  $\mathcal{B} ' $
la base  $\{ \phi (e_1), \ldots , \phi (e_n) \} $.
    \begin{enumerate}
        \item $\Im \phi$  contient  $\mathcal{B} '$  qui est une partie g\' en\' eratrice de $F$. Donc  $\phi $  est surjective.
        \item Soit maintenant  $x\in E$  tel que  $\phi (x)=0$.
Comme  $\mathcal{B} $  est une base, il existe  $\lambda _1 ,
\ldots , \lambda _n\in {\R}$  tels que  $x= \lambda _1e_1+ \cdots
+ \lambda _ne_n$. Alors  $\phi (x)=0=\lambda _1\phi (e_1)+
\cdots + \lambda _n \phi (e_n) $ donc puisque  $\mathcal{B} '$
est libre :  $\lambda _1=\cdots =\lambda _n=0$. En cons\' equence
si  $\phi (x)=0$  alors  $x=0$ : $\phi$  est injective.
    \end{enumerate}
\end{enumerate}

En fait on montrerait de la même façon que ``$\phi $  est un isomorphisme si et seulement si l'image par 
$\phi $  \textbf{d'une} base de  $E$  est une base de  $F$''.
\fincorrection
\nocorrection
\nocorrection
\nocorrection
\nocorrection
\nocorrection
\nocorrection
\nocorrection
\nocorrection
\nocorrection
\correction{003310}
\begin{enumerate}
  \item $\vec u = (\vec u - g\circ f(\vec u)) + g\circ f(\vec u)$.
  \item $f(\Im g) \subset \Im f$.\par
             $f = (f\circ g) \circ f  \Rightarrow  \Im f \subset \Im(f\circ g)= f(\Im g)$.
\end{enumerate}
 \fincorrection
\nocorrection
\nocorrection
\correction{003316}
\begin{enumerate}
  \item $\varphi(v\circ w) = \varphi(v)\circ w + v\circ \varphi(w)$.
  \item Par récurrence $\varphi^n(v\circ w) = \sum_{k=0}^n C_n^k\varphi^k(v)\circ\varphi^{n-k}(w)$
donc si $v\in c_p$ et $w\in c_q$ alors $v\circ w\in c_{p+q-1}$.
\end{enumerate}
 \fincorrection
\nocorrection
\nocorrection
\nocorrection
\nocorrection
\nocorrection
\correction{003332}
  2. $\mathrm{rg}(f+g) = \mathrm{rg}(f) +\mathrm{rg}(g)$
           $\iff \begin{cases}\Im f \cap \Im g = \{\vec 0_F \} \cr
                        \Im( f+g ) = \Im f + \Im g       \cr\end{cases}$ \par
           $\iff \begin{cases}\Im f \cap \Im g = \{\vec 0_F \} \cr
                        \forall\ \vec x,\vec y,\ \exists\ \vec z\text{ tq } f(\vec x)+g(\vec y)=(f+g)(\vec z).\end{cases}$

    

    $ \Rightarrow $ : Donc $f(\vec x - \vec z) = g(\vec z - \vec y) = \vec 0$.
            Pour $\vec y = \vec 0 : \vec x = (\vec x - \vec z) + \vec z \in \mathrm{Ker} f + \mathrm{Ker} g$.

    $\Leftarrow$ : Soient $\vec x = \vec x_f + \vec x_g$ et $\vec y = \vec y_f + \vec y_g$ :
            Alors $f(\vec x) + g(\vec y) = f(\vec x_g) + g(\vec y_f) = (f+g)(\vec x_g + \vec y_f)$.
 \fincorrection
\nocorrection
\nocorrection
\correction{003335}
$\Im f \subset \mathrm{Ker} g  \Rightarrow  \mathrm{rg} f + \mathrm{rg} g \le \dim E$.\par
         $f+g$ est surjective $ \Rightarrow  \Im f + \Im g = E
                                \Rightarrow  \mathrm{rg} f + \mathrm{rg} g \ge \dim E$.
\fincorrection
\correction{003336}
\begin{enumerate}
  \item $\dim H + \dim K = \dim E$.
  \item Si $H \oplus K \ne E$ alors $\cal E$ n'est pas stable pour
              $\circ$.
\end{enumerate}
 \fincorrection
\nocorrection
\nocorrection
\correction{003350}
3. $\dim {\cal K} =(\dim E)(\dim \mathrm{Ker} f) = \dim {\cal I}$, \quad
         $\dim({\cal K} \cap {\cal I})= (\mathrm{rg} f)^2$.
 \fincorrection
\correction{003351}
$(\mathrm{rg} u)(\mathrm{rg} v)$.
\fincorrection
\nocorrection
\correction{005171}
\begin{enumerate}
\item  On a toujours $\mbox{Ker}f\subset\mbox{Ker}f^2$.

En effet, si $x$ est un vecteur de $\mbox{Ker}f$, alors $f^2(x)=f(f(x))=f(0)=0$ (car $f$ est linéaire) et $x$ est 
dans $\mbox{Ker}f^2$.

Montrons alors que~:~$[\mbox{Ker}f=\mbox{Ker}f^2\Leftrightarrow\mbox{Ker}f\cap\mbox{Im}f=\{0\}]$.
Supposons que $\mbox{Ker}f=\mbox{Ker}f^2$ et montrons que $\mbox{Ker}f\cap\mbox{Im}f=\{0\}$.

Soit $x\in\mbox{Ker}f\cap\mbox{Im}f$. Alors, d'une part f(x) = 0 et d'autre part, il existe $y$ élément de $E$ tel que
$x=f(y)$. Mais alors, $f^2(y)=f(x)=0$ et $y\in\mbox{Ker}f^2=\mbox{Ker}f$. Donc, $x=f(y)=0$. On a montré
que $\mbox{Ker}f=\mbox{Ker}f^2\Rightarrow\mbox{Ker}f\cap\mbox{Im}f=\{0\}$.

Supposons que $\mbox{Ker}f\cap\mbox{Im}f=\{0\}$ et montrons que $\mbox{Ker}f=\mbox{Ker}f^2$.

Soit $x\in\mbox{Ker}f^2$. Alors $f(f(x))=0$ et donc $f(x)\in\mbox{Ker}f\cap\mbox{Im}f=\{0\}$. Donc, $f(x)=0$ et $x$ est
dans $\mbox{Ker}f$. On a ainsi montré que $\mbox{Ker}f^2\subset\mbox{Ker}f$ et, puisque l'on a toujours
$\mbox{Ker}f\subset\mbox{Ker}f^2$, on a finalement $\mbox{Ker}f=\mbox{Ker}f^2$. On a montré
que $\mbox{Ker}f\cap\mbox{Im}f=\{0\}\Rightarrow\mbox{Ker}f=\mbox{Ker}f^2$.

On a toujours $\mbox{Im}f^2\subset\mbox{Im}f$. En effet~:~$y\in\mbox{Im}f^2\Rightarrow\exists x\in E/\;y=f^2(x)=f(f(x))
\Rightarrow y\in\mbox{Im}f$.

Supposons que $\mbox{Im}f=\mbox{Im}f^2$ et montrons que $\mbox{Ker}f+\mbox{Im}f=E$.
Soit $x\in E$. Puisque $f(x)\in\mbox{Im}f=\mbox{Im}f^2$, il existe $t\in E$ tel que $f(x)=f^2(t)$. Soit
alors $z=f(t)$ et $y=x-f(t)$. On a bien $x=y+z$ et $z\in\mbox{Im}f$. De plus, $f(y)=f(x)-f(f(t))=0$ et $y$ est bien
élément de $\mbox{Ker}f$. On a donc montré que $E=\mbox{Ker}f+\mbox{Im}f$.

Supposons que $\mbox{Ker}f+\mbox{Im}f=E$ et montrons que $\mbox{Im}f=\mbox{Im}f^2$.

Soit $x\in E$. Il existe $(y,z)\in\mbox{Ker}f\times\mbox{Im}f$ tel que $x=y+z$. Mais alors $f(x)=f(z)\in\mbox{Im}f^2$
car $z$ est dans $\mbox{Im}f$. Ainsi, pour tout $x$ de $E$, $f(x)$ est dans $\mbox{Im}f^2$ ce qui montre que
$\mbox{Im}f\subset\mbox{Im}f^2$ et comme on a toujours $\mbox{Im}f^2\subset\mbox{Im}f$, on a montré que
$\mbox{Im}f=\mbox{Im}f^2$.

\item  $Id-p\;\mbox{projecteur}\Leftrightarrow(Id-p)^2=Id-p\Leftrightarrow Id-2p+p^2=Id-p\Leftrightarrow p^2=p\Leftrightarrow p\;\mbox{projecteur}$.

Soit $x$ un élément de $E$.
$x\in\mbox{Im}p\Rightarrow\exists y\in E/\;x=p(y)$. Mais alors $p(x)=p^2(y)=p(y)=x$. Donc,
$\forall x\in E,\;(x\in\mbox{Im}p\Rightarrow p(x)=x)$.

Réciproquement, si $p(x)=x$ alors bien sûr, $x$ est dans $\mbox{Im}p$.

Finalement, pour tout vecteur $x$ de $E$, $x\in\mbox{Im}p\Leftrightarrow p(x)=x\Leftrightarrow(Id-p)(x)=0\Leftrightarrow x\in\mbox{Ker}(Id-p)$. On a
montré que $\mbox{Im}p=\mbox{Ker}(Id-p)$.

En appliquant ce qui précède à $Id -p$ qui est également un projecteur, on obtient
$\mbox{Im}(Id-p)=\mbox{Ker}(Id-(Id-p))=\mbox{Ker}p$.

Enfin, puisque $p^2=p$ et donc en particulier que $\mbox{Ker}p=\mbox{Ker}p^2$ et $\mbox{Im}p=\mbox{Im}p^2$, le 1) montre
que $E=\mbox{Ker}p\oplus\mbox{Im}p$.

\item 
\begin{align*}
p=p\circ q\;\mbox{et}\;q=q\circ p&\Leftrightarrow
p\circ(Id-q)=0\;\mbox{et}\;q\circ(Id-p)=0\Leftrightarrow\mbox{Im}(Id-q)\subset\mbox{Ker}p\;\mbox{et}\;\mbox{Im}(Id-p)\subset
\mbox{Ker}q\\
 &\Leftrightarrow\mbox{Ker}q\subset\mbox{Ker}p\;\mbox{et}\;\mbox{Ker}p\subset\mbox{Ker}q\;(\mbox{d'après 2)})\\
 &\Leftrightarrow\mbox{Ker}p=\mbox{Ker}q.
\end{align*}

\item  $p\circ q+q\circ p=0\Rightarrow p\circ q=(p\circ p)\circ q=p\circ(p\circ q)=-p\circ (q\circ p)$ et de même,
$q\circ p=q\circ p\circ p=-p\circ q\circ p$. En particulier, $p\circ q=q\circ p$ et donc
$0=p\circ q+q\circ p=2p\circ q=2q\circ p$ puis $p\circ q=q\circ p=0$.

La réciproque est immédiate.

$p+q\;\mbox{projecteur}\;\Leftrightarrow(p+q)^2=p+q\Leftrightarrow p^2+pq+qp+q^2=p+q\Leftrightarrow pq+qp= 0\Leftrightarrow pq=qp=0$ (d'après ci-dessus).
Ensuite, $\mbox{Im}(p+q)=\{p(x)+q(x),\;x\in E\}\subset\{p(x)+q(y),\;(x,y)\in E^2\}=\mbox{Im}p+\mbox{Im}q$.

Réciproquement, soit $z$ un élément de $\mbox{Im}p+\mbox{Im}q$. Il existe deux vecteurs $x$ et $y$ de $E$ tels que
$z=p(x)+q(y)$. Mais alors, $p(z)=p^2(x)+pq(y)=p(x)$ et $q(z)=qp(x)+q^2(y)=q(y)$ et donc

$$z=p(x)+p(y)=p(z)+q(z)=(p+q)(z)\in\mbox{Im}(p+q).$$

Donc, $\mbox{Im}p+\mbox{Im}q\subset\mbox{Im}(p+q)$ et finalement, $\mbox{Im}(p+q)=\mbox{Im}p+\mbox{Im}q$.

$\mbox{Ker}p\cap\mbox{Ker}q=\{x\in E/\;p(x)=q(x)=0\}\subset\{x\in E/\;p(x)+q(x)=0\}=\mbox{Ker}(p+q)$.

Réciproquement, si $x$ est élément de $\mbox{Ker}(p+q)$ alors $p(x)+q(x)=0$. Par suite,
$p(x)=p^2(x)+pq(x)=p(p(x)+q(x))=p(0)=0$ et $q(x)=qp(x)+q^2(x)=q(0)=0$. Donc, $p(x)=q(x)=0$ et
$x\in\mbox{Ker}p\cap\mbox{Ker}q$. Finalement, $\mbox{Ker}(p+q)=\mbox{Ker}p\cap\mbox{Ker}q$.
\end{enumerate}
\fincorrection
\correction{005181}
\begin{enumerate}
\item 
\begin{itemize}
\item[$\Leftarrow$] Soit $(u,v)((\mathcal{L}(E))^2$. On suppose qu'il existe $w\in\mathcal{L}(E)$ tel que $u=w\circ
v$. Soit $x$ un élément de $\mbox{Ker}v$. Alors $v(x)=0$ et donc $u(x)=w(v(x))=w(0)=0$. Mais alors, $x$ est dans
$\mbox{Ker}u$. Donc $\mbox{Ker}v\subset{Ker}u$.

\item[$\Rightarrow$] Supposons que $\mbox{Ker}v\subset\mbox{Ker}u$. On cherche à définir $w$, élément de
$\mathcal{L}(E)$ tel que $w\circ v=u$. Il faut définir précisément $w$ sur $\mbox{Im}v$ car sur $E\setminus\mbox{Im}v$,
on a aucune autre contrainte que la linéarité.

Soit $y$ un élément de $\mbox{Im}v$. (Il existe $x$ élément de $E$ tel que $y=v(x)$. On
a alors envie de poser $w(y)=u(x)$ mais le problème est que $y$, élément de $\mbox{Im}v$ donné peut avoir plusieurs
antécédents $x$, $x'$... et on peut avoir $u(x)\neq u(x')$ de sorte que l'on n'aurait même pas défini une application
$w$.)
\end{itemize}

Soient $x$ et $x'$ deux éléments de $E$ tels que $v(x)=v(x')=y$ alors $v(x-x')=0$ et donc
$x-x'\in\mbox{Ker}v\subset\mbox{Ker}u$. Par suite, $u(x-x')=0$ ou encore $u(x)=u(x')$. En résumé, pour $y$ élément donné
de $\mbox{Im}v$, il existe $x$ élément de $E$ tel que $v(x)=y$. On pose alors $w(y)=u(x)$ en notant que $w(y)$ est bien
uniquement défini, car ne dépend pas du choix de l'antécédent $x$ de $y$ par $v$. $w$ n'est pas encore défini sur $E$
tout entier. Notons $F$ un supplémentaire quelconque de $\mbox{Im}v$ dans $E$ (l'existence de $F$ est admise).

Soit $X$ un élément de $E$. Il existe deux vecteurs $y$ et $z$, de $\mbox{Im}v$ et $F$ respectivement, tels que $X=y+z$.
On pose alors $w(X)=u(x)$ où $x$ est un antécédent quelconque de $y$ par $v$ (on a pris pour restriction de $w$ à $F$
l'application nulle). $w$ ainsi définie est une application de $E$ dans $E$ car, pour $X$ donné $y$ est uniquement
défini puis $u(x)$ est uniquement défini (mais pas nécessairement $x$).

Soit $x$ un élément de $E$ et $y=v(x)$. $w(v(x))=w(y)=w(y+0)=u(x)$ (car 1)$y$ est dans $\mbox{Im}v$ 2)$0$ est dans $F$
3) $x$ est un antécédent de $y$ par $v$) et donc $w\circ v=u$.

Montrons que $w$ est linéaire. Soient, avec les notations précédentes, $X_1=y_1+z_1$ et $X_2=y_2+z_2$ ...

\begin{align*}
w(X_1+X_2)&=w((y_1+y_2)+(z_1+z_2))=u(x_1+x_2)\quad(\mbox{car}\;y_1+y_2=v(x_1)+v(x_2)=v(x_1+x2))\\
 &=u(x_1)+u(x_2)=w(X_1)+w(X_2)
\end{align*}

et

$$w(\lambda X)=w(\lambda y+\lambda z)=u(\lambda x)=\lambda u(x)=\lambda w(X).$$

\item  On applique 1) à $u=Id$.

$$v\;\mbox{injective}\Leftrightarrow\mbox{Ker}v=\{0\}\Leftrightarrow\mbox{Ker}v=\mbox{Ker}Id\Leftrightarrow\exists w\in\mathcal{L}(E)/\;w\circ v=Id.$$

\end{enumerate}
\fincorrection
\correction{005182}
\begin{enumerate}
\item  $\forall P\in E$, $f(P)=P'$ est un polynôme et donc $f$ est une application de $E$ vers $E$.

$\forall(P,Q)\in E^2,\;\forall(\lambda,\mu)\in\Rr^2,\;f(\lambda P+\mu Q)=(\lambda P+\mu Q)'=\lambda P'+\mu
Q'=\lambda f(P)+\mu f(Q)$ et $f$ est un endomorphisme de $E$.

Soit $P\in E$. $P\in\mbox{Ker}f\Leftrightarrow P'=0\Leftrightarrow P$ est constant. $\mbox{Ker}f$ n'est pas nul et $f$ n'est pas injective.

Soient $Q\in E$ puis $P$ le polynôme défini par~:~$\forall x\in\Rr,\;P(x)=\int_{0}^{x}Q(t)\;dt$. $P$ est bien un
polynôme tel que $f(P)=Q$. $f$ est surjective.

Soit $F=\{P\in E/\;P(0) = 0\}$. $F$ est un sous espace de $E$ en tant que noyau de la forme linéaire $P\mapsto
P(0)$. $\mbox{Ker}f\cap F=\{0\}$ car si un polynôme est constant et s'annule en $0$, ce polynôme est nul. Enfin, si $P$
est un polynôme quelconque, $P=P(0)+(P-P(0))$ et $P$ s'écrit bien comme la somme d'un polynôme constant et d'un
polynôme s'annulant en $0$. Finalement $E=\mbox{Ker}f\oplus F$.

\item  On montre facilement que $g$ est un endomorphisme de $E$.

$P\in\mbox{Ker}g\Rightarrow\forall x\in\Rr,\;\int_{0}^{x}P(t)\;dt=0\Rightarrow\forall x\in\Rr,\;P(x)=0$
(en dérivant). Donc, $\mbox{Ker}g=\{0\}$ et donc $g$ est injective.

Si P est dans $\mbox{Im}g$ alors $P(0)=0$ ce qui montre que $g$ n'est pas surjective.
De plus, si $P(0)=0$ alors $\int_{0}^{x}P'(t)\;dt=P(x)-P(0)=P(x)$ ce qui montre que $P=g(P')$ est dans $\mbox{Im}g$
et donc que $\mbox{Im}g=\{P\in E/\;P(0)=0\}$.
\end{enumerate}
\fincorrection
\correction{005187}
Soit $u=(x,y,z,t)=xe_1+ye_2+ze_3+te_4\in\Rr^4$. Alors,

\begin{align*}
f(u)&=xf(e_1)+yf(e_2)+zf(e_3)+tf(e_4)=x(2e_1+e_3)+y(-e_2+e_4)+z(e_1+2e_3)+t(e_2-e_4)\\
 &=(2x+z)e_1+(-y+t)e_2+(x+2z)e_3+(y-t)e_4.
\end{align*}
Par suite,

$$u\in\mbox{Ker }f\Leftrightarrow
\left\{
\begin{array}{l}
2x+z=0\\
-y+t=0\\
x+2z=0\\
y-t=0
\end{array}
\right.\Leftrightarrow
\left\{
\begin{array}{l}
x=z=0\\
y=t
\end{array}
\right..$$
Donc, $\mbox{Ker }f=\{(0,y,0,y),\;y\in\Rr\}=\mbox{Vect}((0,1,0,1))$.

\begin{center}
\shadowbox{
$\mbox{Ker }f=\mbox{Vect}((0,1,0,1))$.
}
\end{center}

Soit $u'=(x',y',z',t')\in\Rr^4$.

\begin{align*}
u'= (x',y',z',t')\in\mbox{Im }f&\Leftrightarrow\exists(x,y,z,t)\in\Rr^4/\;\left\{
\begin{array}{l}
2x+z=x'\\
-y+t=y'\\
x+2z=z'\\
y-t=t'
\end{array}
\right.\Leftrightarrow\exists(x,y,z,t)\in\Rr^4/\;\left\{
\begin{array}{l}
x=\frac{1}{3}(2x'-z')\\
z=\frac{1}{3}(-x'+2z')\\
t=y+y'\\
y'+t'=0
\end{array}
\right.\\
 &\Leftrightarrow y'=-t'
\end{align*}
(si $y'\neq-t'$, le système ci-dessus, d'inconnues $x$, $y$, $z$ et $t$, n'a pas de solution et si $y'=-t'$, le système
ci-dessus admet au moins une solution comme par exemple
$(x,y,z,t)=\left(\frac{1}{3}(2x'-z'),0,\frac{1}{3}(-x'+2z'),y')\right)$.
Donc,
$\mbox{Im }f=\{(x,y,z,t)\in\Rr^4/\;y+t=0\}=\{(x,y,z,-y)/(x,y,z)\in\Rr^3\}=\{xe_1+y(e_2-e_4)+ze_3,\;(x,y,z)\in\Rr^4\}=
\mbox{Vect}(e_1,e_2-e_4,e_3)$.

\begin{center}
\shadowbox{
$\mbox{Im }f==\{(x,y,z,-y)/(x,y,z)\in\Rr^3\}=\mbox{Vect}(e_1,e_2-e_4,e_3)$.
}
\end{center}

\textbf{Autre solution} pour la détermination de $\mbox{Im }f$.
$\mbox{Im }f=\mbox{Vect}(f(e_1),f(e_2),f(e_3),f(e_4))=\mbox{Vect}(2e_1+e_3,-e_2+e_4,e_1+2e_3,e_2-e_4)
=\mbox{Vect}(2e_1+e_3,e_1+2e_3,e_2-e_4)$. Mais d'autre part, d'après le théorème du rang,
$\mbox{dim }(\mbox{Im }f)=4-1=3$. Donc, $(2e_1+e_3,e_1+2e_3,e_2-e_4)$ est une base de $\mbox{Im }f$.
\fincorrection
\correction{005190}
Par définition, $\mbox{rg }(u+v)=\mbox{dim }\left(\mbox{Im }(u+v)\right)$.

\begin{center}
$\mbox{Im }(u+v)=\{u(x)+v(x),\;x\in E\}\subset\{u(x)+v(y),\;(x,y)\in E^2\}=\mbox{Im }u+\mbox{Im }v$.
\end{center}
Donc,

\begin{align*}
rg(u+v)&=\mbox{dim }\left(\mbox{Im }(u+v)\right)\\
 &\leq\mbox{dim }(\mbox{Im }u+\mbox{Im }v)=\mbox{dim }(\mbox{Im }u)+\mbox{dim }(\mbox{Im }v)-\mbox{dim }(\mbox{Im }u\cap\mbox{
Im}v)\\
 &\leq\mbox{dim }(\mbox{Im }u)+\mbox{dim }(\mbox{Im }v)=\mbox{rg }u+\mbox{rg }v.
\end{align*}
On a montré que~:

\begin{center}
\shadowbox{
$\forall(u,v)\in(\mathcal{L}(E,F))^2,\;\mbox{rg }(u+v)\leq\mbox{rg }u+\mbox{rg }v.$
}
\end{center}
Ensuite, $$\mbox{rg }u=\mbox{rg }(u+v-v)\leq\mbox{rg }(u+v)+\mbox{rg }(-v)=\mbox{rg }(u+v)+\mbox{rg }v,$$
(il est clair que $\mbox{Im }(-v)=\mbox{Im }v)$ et donc $\mbox{rg }u-\mbox{rg }v\leq\mbox{rg }(u+v)$. En échangeant les rôles de
$u$ et $v$, on a aussi $\mbox{rg }v-\mbox{rg }u=\mbox{rg }(u+v)$ et finalement

\begin{center}
\shadowbox{
$\forall(u,v)\in(\mathcal{L}(E,F))^2,\;|\mbox{rg }u-\mbox{rg }v|\leq\mbox{rg }(u+v).$
}
\end{center}

\fincorrection
\correction{005191}
 \begin{enumerate}
 \item  
\textbullet~\textbf{(1)$\Rightarrow$(2).} Si $\mbox{Ker }f=\mbox{Im }f$, alors pour tout élément $x$ de $E$, $f(x)$ est dans
$\mbox{Im }f=\mbox{Ker }f$ et donc $f(f(x))=0$. Par suite, $f^2=0$. De plus, d'après le théorème du rang,
$n=\mbox{dim }(\mbox{Ker }f)+\mbox{rg }f=2\;\mbox{rg }f$ ce qui montre que $n$ est nécessairement pair et que
$\mbox{rg }f=\frac{n}{2}$.
\textbullet~\textbf{(2)$\Rightarrow$(3).} Si $f^2=0$ et $n=2\;\mbox{rg }f\;(\in2\Nn)$, cherchons un endomorphisme $g$ de $E$ tel que
$f\circ g+g\circ f=Id_E$.
Posons $r=\mbox{rg }f$ et donc $n=2r$, puis $F=\mbox{Ker }f=\mbox{Im }f$ ($\mbox{dim }F=r$).

Soit $G$ un supplémentaire de $F$ dans $E$ ($\mbox{dim }G=r$).
Soit $(e_1',...,e_r')$ une base de $G$. Pour $i\in\llbracket1,r\rrbracket$, on pose $e_i=f(e_i')$.
Montrons que la famille $(e_1,...,e_r)$ est libre. Soit $(\lambda_1,...,\lambda_r)\in\Rr^r$.

$$\sum_{i=1}^{r}\lambda_ie_i=0\Rightarrow f\left(\sum_{i=1}^{r}\lambda
e_i'\right)=0\Rightarrow\sum_{i=1}^{r}\lambda_ie_i'\in\mbox{Ker }f\cap G=\{0\}\Rightarrow\forall
i\in\{1,...,r\},\;\lambda_i=0,$$
car la famille $(e_i')_{1\leq i\leq r}$ est libre. $(e_1,...,e_r)$ est une famille libre de $F=\mbox{Im }f$ de cardinal
$r$ et donc une base de $F=\mbox{Ker }f=\mbox{Im }f$. Au passage, puisque $E=F\oplus G$, $(e_1,...,e_r,e_1',...,e_r')$ est
une base de $E$.
Soit alors $g$ l'endomorphisme de $E$ défini par les égalités~:~$\forall i\in\llbracket1,r\rrbracket,\;g(e_i)=e_i'$ et
$g(e_i')=e_i$ ($g$ est entièrement déterminé par les images des vecteurs d'une base de $E$). Pour $i$ élément de
$\llbracket1,r\rrbracket$, on a alors~:~

$$(f\circ g+g\circ f)(e_i)=f(e_i')+g(0)=e_i+0= e_i,$$

et

$$(f\circ g+g\circ f)(e_i')=f(e_i)+g(e_i)=0+e_i'=e_i'.$$
Ainsi, les endomorphismes $f\circ g+g\circ f$ et $Id_E$ coïncident sur une base de $E$ et donc $f\circ g+g\circ f=Id_E$.
\textbullet~\textbf{(3)$\Rightarrow$(1).} Supposons que $f^2=0$ et qu'il existe $g\in\mathcal{L}(E)$ tel que $f\circ g+g\circ
f=Id_E$.
Comme $f^2=0$, on a déjà $\mbox{Im }f\subset\mbox{Ker }f$. D'autre part, si $x$ est un élément de $\mbox{Ker }f$, alors
$x=f(g(x))+g(f(x))=f(g(x))\in\mbox{Im }f$ et on a aussi $\mbox{Ker }f\subset\mbox{Im }f$. Finalement,
$\mbox{Ker }f=\mbox{Im }f$.

 \item  L'existence d'une base $(e_1,...,e_p,e_1',...,e_p')$ de $E$ vérifiant les conditions de l'énoncé a été établie
au passage (avec $p=r=\mbox{rg }f$).
\end{enumerate}
\fincorrection
\correction{005192}
\begin{enumerate}
 \item  Soient $k$ un entier naturel et $x$ un élément de $E$.

$$x\in N_k\Rightarrow f^k(x)=0\Rightarrow f(f^k(x))=f(0)=0\Rightarrow f^{k+1}(x)=0\Rightarrow x\in N_{k+1}.$$
On a montré que~:~$\forall k\in\Nn,\;N_k\subset N_{k+1}$. Ensuite,

$$x\in I_{k+1}\Rightarrow\exists y\in E/\;x=f^{k+1}(y)\Rightarrow\exists z(=f(y))\in E/\;x=f^k(z)\Rightarrow x\in I_k.$$
On a montré que~:~$\forall k\in\Nn,\;I_{k+1}\subset I_k$.
 \item  
  \begin{enumerate}
  \item Soit $k$ un entier naturel. Supposons que $N_k=N_{k+1}$.

On a déjà $N_{k+1}\subset N_{k+2}$. Montrons que $N_{k+2}\subset N_{k+1}$.

Soit $x$ un élément de $E$.

\begin{align*}
x\in N_{k+2}&\Rightarrow f^{k+2}(x)=0\Rightarrow f^{k+1}(f(x))=0\Rightarrow f(x)\in N_{k+1}=N_k
\Rightarrow f^k(f(x))=0\\
 &\Rightarrow f^{k+1}(x)=0\Rightarrow x\in N_{k+1}.
\end{align*}
  \item  On a $\{0\}=N_0\subset N_1\subset N_2...$ Supposons que chacune de ces inclusions soient strictes. Alors,
$0=\mbox{dim }N_0<\mbox{dim }N_1<\mbox{dim }N_2$... Donc $\mbox{dim }N_1\geq1$, $\mbox{dim }N_2\geq2$ et par une récurrence
facile, $\forall k\in\Nn,\;\mbox{dim }N_k\geq k$. En particulier, $\mbox{dim }N_{n+1}\geq n+1>n=\mbox{dim }E$, ce qui est
impossible. Donc, il existe $k$ entier naturel tel que $N_k=N_{k+1}$.

Soit $p$ le plus petit de ces entiers $k$ (l'existence de $p$ est démontrée proprement de la façon suivante~:~si
$K=\in\{k\in\Nn/\;N_k=N_{k+1}\}$, $K$ est une partie non vide de $\Nn$ et admet donc un plus petit élément). On note
que puisque $f$ est non injectif, $\{0\}=N_0\underset{\neq}{\subset}N_1$ et donc $p\in\Nn^*$. Par définition de $p$,
pour $k<p$, $N_k\underset{\neq}{\subset}N_{k+1}$ et, d'après le a) et puisque $N_p=N_{p+1}$, on montre par récurrence
que pour $k=p$, on a $N_k=N_p$.
  \item  $0<\mbox{dim }N_1<...<\mbox{dim }N_p$ montre que pour $k\leq p$, on a $\mbox{dim }N_k=k$ et en particulier $p
\leq\mbox{dim }N_p=n$.
  \end{enumerate}
 \item  Puisque $N_k\subset N_{k+1}$, $I_{k+1}\subset I_k$ et que $\mbox{dim }E<+\infty$, on a~:

$$N_k=N_{k+1}\Leftrightarrow\mbox{dim }N_k=\mbox{dim }N_{k+1}\Leftrightarrow
n-\mbox{rg }(f^k)=n-\mbox{rg }(f^{k+1})\Leftrightarrow\mbox{dim }(I_k)=\mbox{dim }(I_{k+1})\Leftrightarrow I_k=I_{k+1}.$$

Donc, pour $k<p$, $I_k\underset{neq}{\subset}I_{k+1}$ et pour $k=p$, $I_k=I_{k+1}$.
 \item  Soit $x\in I_p\cap N_p$. Alors, $f^p(x)=0$ et $\exists y\in E/\;x=f^p(y)$. D'où, $f^{2p}(y)=0$ et $y\in
N_{2p}=N_p$ (puisque $2p\geq p$) et donc $x=f^p(y)=0$. On a montré que $I_p\cap N_p=\{0\}$. Maintenant, le théorème du
rang montre que $E=\mbox{dim }(I_p)+\mbox{dim }(N_p)$ et donc $E=I_p\oplus N_p$.

Posons $f_{/I_p}=f'$. $f'$ est déjà un endomorphisme de $I_p$ car $f'(I_p)=f(I_p)=I_{p+1}=I_p$.

Soit alors $x\in I_p$. $\exists y\in E/\;x=f^p(y)$.

$$x\in\mbox{Ker }f'\Rightarrow f'(x)=0\Rightarrow f(f^p(y))=0\Rightarrow y\in N_{p+1}=N_p x=f_p(y)=0.$$

Donc $\mbox{Ker }f'=\{0\}$ et donc, puisque $\mbox{dim }I_p<+\infty$, $f'\in\mathcal{GL}(I_p)$.
 \item  Soient $k$ un entier naturel et $g_k$ la restriction de $f$ à $I_k$.

D'après le théorème du rang, $d_k=\mbox{dim }(I_k)=\mbox{dim }(\mbox{Ker }g_k)+\mbox{dim }(\mbox{Im }g_k)$.
Maintenant, $\mbox{Im }g_k=g_k(I_k)=f(I_k)=I_{k+1}$ et donc $\mbox{dim }(\mbox{Im }g_k)=d_{k+1}$. D'autre part, $\mbox{Ker }g_k=\mbox{Ker }f_{/I_k}=\mbox{Ker }f\cap I_k$.
 
Ainsi, pour tout entier naturel $k$, $d_k-d_{k+1}=\mbox{dim }(\mbox{Ker }f\cap I_k)$. Puisque la suite $(I_k)_{k\in\Nn}$ est décroissante pour l'inclusion, la suite d'entiers naturels $(\mbox{dim }(\mbox{Ker }f\cap I_k))_{k\in\Nn}=(d_k-d_{k+1})_{k\in\Nn}$ est décroissante.
\end{enumerate}
\fincorrection
\correction{005194}
Soit $x\in E$.

$$x\in\mbox{Ker }(f-2Id)\cap\mbox{Ker }(f-3Id)\Rightarrow f(x)=2x\;\mbox{et}\;f(x)=3x\Rightarrow 3x-2x=f(x)-f(x)= 0\Rightarrow x=0.$$

Donc, $\mbox{Ker }(f-2Id)\cap\mbox{Ker }(f-3Id)=\{0\}$ (même si $f^2-5f+6Id\neq0$).

Soit $x\in E$. On cherche $y$ et $z$ tels que $y\in\mbox{Ker }(f-2Id)$, $z\in\mbox{Ker }(f-3Id)$ et $x=y+z$.

Si $y$ et $z$ existent, $y$ et $z$ sont solution du système $\left\{
\begin{array}{l}
y+z=x\\
2y+3z=f(x)
\end{array}
\right.$ et donc $\left\{
\begin{array}{l}
y=3x-f(x)\\
z=f(x)-2x
\end{array}
\right.$.

Réciproquement . Soient $x\in E$ puis $y=3x-f(x)$ et $z=f(x)-2x$.

On a bien $y+z=x$ puis

\begin{align*}
f(y)&=3f(x)-f^2(x)=3f(x)-(5f(x)-6x)\quad(\mbox{car}\;f^2=5f-6Id)\\
 &= 6x-2f(x)=2(3x-f(x))=2y
\end{align*}  
  
et $y\in\mbox{Ker }(f-2Id)$. De même,

$$f(z)=f^2(x)-2f(x)=(5f(x)-6x)-2f(x)=3(f(x)-2x)=3z,$$
 
et $z\in\mbox{Ker }(f-3Id)$. On a montré que $E=\mbox{Ker }(f-2Id)+\mbox{Ker }(f-3d)$, et finalement que 

$$E=\mbox{Ker }(f-2Id)\oplus\mbox{Ker }(f-3d).$$

\fincorrection
\correction{005582}
Une condition nécessaire est bien sur $\text{dim}F +\text{dim}G =\text{dim}E$ (et non pas $F\oplus G = E$).

Montrons que cette condition est suffisante. Soient $F$ et $G$ deux sous-espaces de $E$ tels que $\text{dim}F+\text{dim}G =\text{dim}E$.

Soit $F'$ un supplémentaire de $F$ dans $E$ ($F'$ existe car $E$ est de dimension finie).

Si $G=\{0\}$ (et donc $F = E$), $f=0$ convient.

Si $G\neq\{0\}$, il existe un isomorphisme $\varphi$ de $F'$ sur $G$ (car $F'$ et $G$ ont même dimension finie) puis il existe un unique endomorphisme de $E$ vérifiant : $f_{/F}=0_{/F}$ et $f_{/F'}=\varphi$.

Mais alors $\text{Im}f=f(F\oplus F') = f(F)+f(F') =\{0\}+G=G$ puis $F\subset\text{Ker}f$ et pour des raisons de dimension, $F=\text{Ker}f$.
\fincorrection
\correction{005583}
\begin{enumerate}
\item 

 $\Leftarrow$/ Si $f=0$, $f$ n'est pas injective (car $E\neq\{0\}$).

Si $f\neq 0$ et s'il existe un endomorphisme non nul $g$ de $E$ tel que $f\circ g= 0$ alors il existe un vecteur $x$ de E tel que $g(x)\neq 0$ et $f(g(x)) = 0$. Par suite $\text{Ker}f\neq\{0\}$ et $f$ n'est pas injective.

$\Rightarrow$/ Supposons $f$ non injective et non nulle. Soient $F=\text{Ker}f$ et $G$ un supplémentaire quelconque de $F$ dans $E$. Soit $p$ la projection sur $F$ parallèlement à $G$.

Puisque $F=\text{Ker}f$, on a $f\circ p=0$ et puisque $f$ n'est pas nul, $F$ est distinct de $E$ et donc $G$ n'est pas nul ($E$ étant de dimension finie) ou encore $p$ n'est pas nul. $f$ est donc diviseur de zéro à gauche.

\item  $\Leftarrow$/ Si $f=0$, $f$ n'est pas surjective.

Si $f$ n'est pas nul et s'il existe un endomorphisme non nul $g$ de $E$ tel que $g\circ f=0$ alors $f$ ne peut être surjective car sinon $g(E)=g(f(E))=\{0\}$ contredisant $g\neq0$.

$\Rightarrow$/ Supposons $f$ non surjective et non nulle.

Soient $G=\text{Im}f$ et $F$ un supplémentaire quelconque de $G$ dans $E$ puis $p$ la projection sur $F$ parallèlement à $G$. $F$ et $G$ sont non nuls et distincts de $E$ et donc $p$ n'est pas nulle et vérifie $p\circ f= 0$. $f$ est donc diviseur de zéro à droite.
\end{enumerate}
\fincorrection
\correction{005586}
\begin{enumerate}
 \item 
  \begin{enumerate}
  \item Soient $k\in\Nn$ et $x\in E$. $x\in N_k\Rightarrow f^k(x)=0\Rightarrow f(f^k(x)) = 0\Rightarrow x\in N_{k+1}$. 

\begin{center}
\shadowbox{
$\forall k\in\Nn$, $N_k\subset N_{k+1}$.
}
\end{center}

Soient $k\in\Nn$ et $y\in I_{k+1}\Rightarrow\exists x\in E/\;y = f^{k+1}(x)\Rightarrow\exists x\in E/\;y = f^k(f(x))\Rightarrow y\in I_k$.

\begin{center}
\shadowbox{
$\forall k\in\Nn$, $I_{k+1}\subset I_{k}$.
}
\end{center}

  \item Soit $x\in N$. Il existe un entier $k$ tel que $x$ est dans $N_k$ ou encore tel que $f^k(x)=0$. Mais alors $f^k(f(x)) = f(f^k(x))=0$ et $f(x)$ est dans $N_k$ et donc dans $N$. Ainsi, $N$ est stable par $f$.

Soit $y\in I$. Alors, pour tout naturel $k$, il existe $x_k\in E$ tel que $y=f^k(x_k)$. Mais alors, pour tout entier $k$, $f(y)=f(f^k(x_k))=f^k(f(x))$ est dans $I_k$, et donc $f(y)$ est dans $I$. $I$ est stable par $f$.

  \item Si $N_k=N_{k+1}$, on a déjà $N_{k+1}\subset N_{k+2}$. Montrons que $N_{k+2}\subset N_{k+1}$.

Soit $x\in N_{k+2}$. Alors $f^{k+1}(f(x))=0$ et donc $f(x)\in N_{k+1}= N_k$. Donc, $f^k(f(x))=0$ ou encore $x$ est dans $N_{k+1}$. On a montré que

\begin{center}
\shadowbox{
$\forall k\in\Nn$, $[(N_k = N_{k+1})\Rightarrow (N_{k+1}= N_{k+2})]$.
}
\end{center}
  \end{enumerate}

\item   
  \begin{enumerate}
  \item  Notons tout d'abord que, pour tout entier naturel $k$, $N_k\subset N_{k+1}$ et $I_{k+1}\subset I_k$. Si de plus, on est en dimension finie, alors d'après le théorème du rang,

\begin{center}
$N_k =N_{k+1}\Leftrightarrow I_{k+1}=I_k\Leftrightarrow\text{dim}N_k=\text{dim}N_{k+1}$.
\end{center}

Donc $A=B$ (éventuellement $=\varnothing$).

La suite des noyaux itérés ne peut être strictement croissante pour l'inclusion car alors la suite des dimensions de ces noyaux serait une suite strictement croissante d'entiers naturels, vérifiant par une récurrence facile $\text{dim}N_k\geqslant k$ pour tout naturel $k$, et en particulier $\text{dim}N_{n+1} >\text{dim}E$ ce qui est exclu.

Donc il existe un entier $k$ tel que $N_k=N_{k+1}$. Soit $p$ le plus petit de ces entiers $k$.

Par définition de $p$, $N_k$ est strictement inclus dans $N_{k+1}$ pour $k<p$, puis $N_p=N_{p+1}$ et d'après 1)c) pour tout entier naturel $k$ supérieur ou égal à $p$ on a $N_k=N_p$ (par récurrence sur $k\geqslant p$). Donc $A=\{p,p+1,p+2,...\}$.

Enfin, $\text{dim}(N_0)<\text{dim}(N_1)<...<\text{dim}(N_p)$ et donc $\text{dim}(N_p)\geqslant p$ ce qui impose $p\leqslant n$.

  \item On a déjà $\text{dim}N_p +\text{dim}I_p=\text{dim}E$. Il reste à vérifier que $I_p\cap N_p= \{0\}$.

Soit $x$ un élément de $I_p\cap N_p$. Donc $f^p(x) = 0$ et il existe $y\in E$ tel que $x = f^p(y)$.
Mais alors $f^{2p}(y) = 0$ et $y$ est dans $N_{2p}= N_p$ (car $2p\geqslant p$) ou encore $x = f^p(y) = 0$.

\begin{center}
\shadowbox{
$E=I_p\oplus N_p$.
}
\end{center}

  \item Ici $N=N_p=\text{Ker}f^p$ et $I = I_p =\text{Im}f^p$.

Soit $f'=f_{/N}$. D'après 1)b), $f'$ est un endomorphisme de $N$ puis immédiatement $f'^p = 0$. Donc $f_{/N}$ est nilpotent.

Soit $f''=f_{/I}$ . $f''$ est d'après 1)b) un endomorphisme de $I$. Pour montrer que $f''$ est un automorphisme de $I$, il suffit de vérifier que $\text{Ker}f''=\{0\}$. Mais $\text{Ker}f''\subset\text{Ker}f\subset N$ et aussi $\text{Ker}f''\subset I$. Donc $\text{Ker}f''\subset N\cap I =\{0\}$. Donc $f_{/I}\in\mathcal{GL}(I)$.
  \end{enumerate}

\item  Il faut bien sûr chercher les exemples en dimension infinie.

  \begin{enumerate}
  \item Soit $f$ de $\Rr[X]$ dans lui-même qui à un polynôme $P$ associe sa dérivée $P'$.
On vérifie aisément que $\forall k \in\Nn$, $N_k =\Rr_k[X]$ et donc la suite des noyaux itérés est strictement croissante. La suite des $I_k$ est par contre constante : $\forall k\in\Nn$, $I_k=\Rr[X]$. Dans ce cas, $A$ est vide et $B=\Nn$.

  \item A un polynôme $P$, on associe le polynôme $XP$. Les $N_k$ sont tous nuls et pour $k\in\Nn$ donné, $I_k$ est constitué des polynômes de valuation supérieure ou égale à $k$ ou encore $I_k=X^k\Rr[X]$. Dans ce cas, $A=\Nn$ et $B=\varnothing$.

  \item Soit $f$ l'endomorphisme de $\Rr[X]$ qui à Xn associe Xn+1 si $n$ n'est pas une puissance de $2$ et $0$ si $n$ est une puissance de $2$ ($f(1) = X$, $f(X) = 0$, $f(X^2) = 0$, $f(X^3) = X^4$, $f(X^4) = 0$, ...)

Soit $k$ un entier naturel.

\begin{center}
$f^{2^k-1}(X^{2^k+1})=X^{2^k+1+2^k-1}=X^{2^{k+1}}\neq0$ et $f^{2^k}(X^{2^k+1})=f(X^{2^{k+1}})=0$.
\end{center}

Donc, pour tout entier naturel $k$, $N_{2^k-1}$ est strictement inclus dans $N_k$. $A$ est vide.

Ensuite, $X^{2^{k+1}}\in I_{2^k-1}$ mais $X^{2^{k+1}}\notin I_{2^k}$. En effet, si $l\geqslant 2^{k+1} +1$, $f^{2^k}(X^l)$ est ou bien nul ou bien de degré supérieur ou égal à $2^k+2^{k+1}+1>2^{k+1}$ et si $l\leqslant 2^{k+1}$, $f^{2^k}(X^l)=0$ car entre $l$ et $2^{k}+l-1$, il y a une puissance de $2$ (il y a $2^k$ nombres entre $l$ et $2^k+l-1$, ensuite $2^k+l-1 < 2^k+2^{k+1}= 3\times2^k <2^{k+2}$ et enfin l'écart entre deux puissances de $2$ inférieures à $2^{k+1}$ vaut au maximum $2^{k+1}-2^k = 2^k$) . Donc, $I_{2^k}$ contient le polynôme nul ou des polynômes de degré strictement supérieur à $2^{k+1}$ et ne contient donc pas $X^{2^{k+1}}$. Finalement, pour tout entier naturel $k$, $I_{2^k}$ est strictement inclus dans $I_{2^k-1}$ et $B$ est vide.

  \end{enumerate}

\item  Pour $k$ entier naturel donné, on note $f_k$ la restriction de $f$ à $I_k$. D'après le théorème du rang, on a

\begin{center}
$\text{dim}I_k=\text{dim}\text{Ker}f_k+\text{dim}\text{Im}f_k$ avec $\text{Im}f_k=f(I_k)=I_{k+1}$.
\end{center}

Donc, pour tout entier naturel $k$, $d_k -d_{k+1}=\text{dim}\text{Ker}f_k$.

Or, pour tout entier naturel $k$, $\text{Ker}f_{k+1}=\text{Ker}f\cap I_{k+1}\subset\text{Ker}f\cap I_{k}=\text{Ker}f_k$ et donc $d_{k+1}-d_{k+2}=\text{dim}\text{Ker}f_{k+1}\leqslant\text{dim}\text{Ker}f_k=d_k-d_{k+1}$.

Finalement, pour tout entier naturel $k$, $d_{k+1}-d_{k+2}\leqslant d_k-d_{k+1}$ et la suite des images itérées décroît de moins en moins vite.
\end{enumerate}
\fincorrection
\correction{005601}
\begin{enumerate}
 \item  Si $N=\text{Ker}f\neq\{0\}$, considérons $g$ non nul tel que $\text{Im}g\neq\{0\}$ et $\text{Im}g\subset\text{Ker}f$.

Pour un tel $g$, $f\circ g = 0$ puis $f\circ g\circ f = 0$ et donc $g = 0$ par hypothèse, contredisant $g$ non nulle. Donc $\text{Ker}f =\{0\}$.

Si $\text{Im}f\neq F$, on choisit $g$ nulle sur $\text{Im}f$ et non nulle  sur un supplémentaire de $\text{Im}f$ (dont l'existence est admise en dimension infinie). Alors, $g\circ f = 0$ puis $f\circ g\circ f = 0$ et donc $g = 0$ contredisant $g$ non nulle. Donc $\text{Im}f = F$.

Finalement, $f$ est bien un isomorphisme de $E$ sur $F$.

\item  Soit $A=\{g\in\mathcal{L}(F,E)/\;f\circ g\circ f = 0\}$. Tout d'abord $A$ est bien un sous-espace vectoriel de $\mathcal{L}(F,E)$ car contient l'application nulle et est stable par combinaison linéaire (ou bien $A$ est le noyau de l'application linéaire de $\mathcal{L}(F,E)$ dans $\mathcal{L}(E,F)$ qui à $g$ associe $f\circ g\circ f$).

Soit $J$ un supplémentaire de $I=\text{Im}f$ dans $F$. Un élément $g$ de $\mathcal{L}(F,E)$ est entièrement déterminé par ses restrictions à $I$ et $J$.

\begin{center}
$f\circ g\circ f = 0\Leftrightarrow(f\circ g)_{/I}= 0\;\text{et}\;g_{/J}\;\text{est quelconque}\Leftrightarrow g(I)\subset N$.
\end{center}

Pour être le plus méticuleux possible, on peut alors considérer l'application $G$ de $\mathcal{L}(I,N)\times\mathcal{L}(J,E)$ dans $\mathcal{L}(F,E)$ qui à un couple $(g_1,g_2)$ associe l'unique application linéaire $g$ de $F$ dans $E$ telle que $g_{/I}= g_1$ et $g_{/J}=g_2$. $G$ est linéaire et injective d'image $A$. Donc 

\begin{center}
$\text{dim}A =\text{dim}\mathcal{L}(I,N)\times\text{dim}\mathcal{L}(J,E)=\text{dim}\mathcal{L}(I,N) + \text{dim}\mathcal{L}(J,E) = r(p - r) + (n - r)p = pn - r^2$.
\end{center}
\end{enumerate}
\fincorrection
\correction{005602}
\begin{enumerate}
 \item  $u$ est dans L(E) car $u$ est linéaire et si $P$ est un polynôme de degré au plus $n$ alors $u(P)$ est un polynôme de degré au plus $n$.

\textbullet~Les polynômes constants sont dans $\text{Ker}u$. Réciproquement , soit $P$un élément de $\text{Ker}u$ puis $Q = P - P(0)$.

Par hypothèse, $P(0) = P(1) = P(2) = ...$ et donc $0$, $1$, $2$, ... sont des racines de $Q$. Puisque le polynôme $Q$ admet une infinité de racines, $Q$ est nul et donc $P = P(0)$ et $P\in\Kk_0[X]$. Ainsi, $\text{Ker}u =\Kk_0[X]$. 

\textbullet~Mais alors, d'après le théorème du rang, $\text{rg}u =(n+1)-1 = n$.
D'autre part, si $P$ est dans $\Kk_n[X]$, $P(X+1)-P$ est dans $\Kk_{n-1}[X]$ (si on pose $P=a_nX^n+\ldots$, le coefficient de $X^n$ dans $u(P)$ est $a_n-a_n=0$).

En résumé, $\text{Im}u\subset\Kk_{n-1}[X]$ et $\text{dim}\text{Im}u=\text{dim}\Kk_{n-1}[X]<+\infty$ et donc $\text{Im}u=\Rr_{n-1}[X]$.

\begin{center}
\shadowbox{
$\text{Ker}u=\Kk_0[X]$ et $\text{Im}u=\Kk_{n-1}[X]$.
}
\end{center}

\item  On part de $P_0 = 1$ et aussi de $P_1 = X$ qui vérifient bien $u(P_0) = 0$ et $u(P_1) = P_0$.

Trouvons $P_2=aX^2 + bX$ tel que $u(P_2) = P_1$ (il est clair que si $\text{deg}(P)\geqslant1$, $\text{deg}(u(P))=\text{deg}(P)-1$ et d'autre part, les constantes sont inutiles car $\text{Ker}u=\Kk_0[X]$).

\begin{center}
$u(P_2) =P_1\Leftrightarrow a(X+1)^2 +b(X+1) - aX^2 - bX =X\Leftrightarrow (2a-1)X +a+b =0\Leftrightarrow a =\frac{1}{2}\;\text{et}\;b = -a$.
\end{center}

On prend $P_2=\frac{1}{2}(X^2-X)=\frac{1}{2}X(X-1)$.

Trouvons $P_3=aX^3 + bX^2+cX$ tel que $u(P_3) = P_2$.

\begin{align*}\ensuremath
u(P_3) =P_2&\Leftrightarrow a(X+1)^3 +b(X+1)^2+c(X+1) - aX^3 - bX^2-cX =\frac{1}{2}X^2-\frac{1}{2}X\\
 &\Leftrightarrow \left(3a-\frac{1}{2}\right)X^2+\left(3a+2b-\frac{1}{2}\right)X+a+b+c=0\\
 &\Leftrightarrow a =\frac{1}{6}\;\text{et}\;b =-\frac{1}{2}\;\text{et}\;c=\frac{1}{3}.
\end{align*}

On prend $P_3=\frac{1}{6}(X^3-3X^2+2X)=\frac{1}{6}X(X-1)(X-2)$.

Essayons, pour $1\leqslant k\leqslant n$, $P_k=\frac{1}{k!}\prod_{i=0}^{k-1}(X-i)$. Pour $1\leqslant k\leqslant n-1$,

\begin{align*}\ensuremath
u(P_{k+1})&=\frac{1}{(k+1)!}\prod_{i=0}^{k}(X+1-i)-\frac{1}{(k+1)!}\prod_{i=0}^{k}(X-i)=\frac{1}{(k+1)!}((X+1)-(X-k))\prod_{i=0}^{k-1}(X-i)\\
 &=\frac{1}{k!}\prod_{i=0}^{k-1}(X-i)=P_k.
\end{align*}

Enfin, les $P_k$, $0\leqslant k\leqslant n$, constituent une famille de $n+1=\text{dim}\Kk_n[X]$ polynômes de degrés échelonnés de $\Kk_n[X]$ et donc la famille $(P_k)_{0\leqslant k\leqslant n}$ est une base de $\Kk_n[X]$. Dans cette base, la matrice de $u$ a la forme désirée.
\end{enumerate}
\fincorrection
\nocorrection
\nocorrection
\nocorrection
\nocorrection
\nocorrection
\nocorrection
\nocorrection
\nocorrection
\correction{000974}
\begin{enumerate}
  \item La seule fonction qui est \`a la fois paire et impaire est la fonction nulle : $P\cap I = \{0\}$. Montrons qu'une fonction $f:\Rr \longrightarrow \Rr$ se d\'ecompose en une fonction paire et une fonction impaire. 
En effet : 
 $$f(x)= \frac{f(x)+f(-x)}{2}+\frac{f(x)-f(-x)}{2}.$$
La fonction $x \mapsto \frac{f(x)+f(-x)}{2}$ est paire (le v\'erifier !),
la fonction $x \mapsto \frac{f(x)-f(-x)}{2}$ est impaire.
Donc $P+I=E$.
Bilan : $E=P\oplus I.$
  \item Le projecteur sur $P$ de direction $I$ est l'application $\pi : E \longrightarrow E$
    qui \`a $f$ associe la fonction $x \mapsto \frac{f(x)+f(-x)}{2}$, c'est-à-dire à $f$ on associe la partie paire
 de $f$.
Nous avons bien 
\begin{itemize}
  \item $\pi(f)\in P$. Par définition de $\pi$, $\pi(f)$ est bien une fonction paire.
  \item $\pi \circ \pi = \pi$. Si $g$ est une fonction paire alors $\pi(g)=g$. 
Appliquons ceci avec $g=\pi(f)$ (qui est bien est une fonction paire) donc $\pi(\pi(f))=\pi(f)$. 
  \item $\ker \pi = I$. Si $\pi(f)=0$ alors cela signifie exactement que la fonction $x \mapsto \frac{f(x)+f(-x)}{2}$
est la fonction nulle. Donc pour tout $x$ : $\frac{f(x)+f(-x)}{2}=0$ donc $f(x)=-f(-x)$ ; 
cela implique que $f$ est une fonction impaire. Réciproquement si $f\in I$ est une fonction impaire, 
sa partie paire est nulle donc $f\in \ker f$.
\end{itemize}

\end{enumerate}
\fincorrection
\nocorrection
\correction{000976}
\begin{enumerate}
  \item $f$ est bien lin\'eaire...
  \item Soit $P$ tel que $f(P)=0$. Alors $P$ v\'erifie l'\'equation diff\'erentielle
$$P+(1-X)P'=0.$$ Dont la solution est $P = \lambda(X-1)$, $\lambda \in \Rr$.
Donc $\ker f$ est de dimension $1$ et une base est donn\'ee par un seul vecteur : $X-1$.

\item Par le th\'eor\`eme du rang la dimension de l'image est :
$$\dim \Im f = \dim \Rr_n[X]-\dim \ker f = (n+1) - 1 = n.$$
Il faut donc trouver $n$ vecteurs lin\'eairement ind\'ependants dans $\Im f$.
\'Evaluons $f(X^k)$, alors 
$$f(X^k) = (1-k)X^k+kX^{k-1}.$$
Cela donne $f(1)=1, f(X)=1, f(X^2)=-X^2+2X,...$
on remarque que pour $k= 2,\ldots n$, $f(X^k)$ est de degr\'e $k$ sans terme constant.
Donc l'ensemble 
$$\big\{ f(X), f(X^2), \ldots, f(X^n)\big\}$$
est une famille de $n$ vecteurs, appartenant \`a $\Im f$, et libre (car les degr\'es sont distincts).
Donc ils forment une base de $\Im f$.
\end{enumerate}
\fincorrection
\nocorrection
\nocorrection
\nocorrection
\nocorrection
\nocorrection
\nocorrection
\nocorrection
\nocorrection
\nocorrection
\nocorrection
\nocorrection
\correction{003342}
2. (c) $g(\vec x) = \vec z$.

 \fincorrection
\correction{003343}

    3. (a) $(f+g)(\vec x) = \vec0  \Rightarrow  f^k(\vec x) + g\circ f^{k-1}(\vec x) = \vec0$. \par
             Pour $k = p$   : $f^{k-1}(\vec x) = \vec0$, puis
             pour $k = p-1$ : $f^{k-2}(\vec x) = \vec 0$, etc, jusqu'à $\vec x = \vec 0$.\\
    3. (b) Même principe sur l'équation : $(f+g)(\vec x) = \vec y$.\par
             On obtient : $(f+g)^{-1} = g^{-1} \circ \bigl( \mathrm{id} - g^{-1}\circ f
             + g^{-2}\circ f^2 - \dots + (-1)^{p-1}g^{1-p}\circ f^{p-1} \bigr)$.

 \fincorrection
\nocorrection
\nocorrection
\nocorrection
\nocorrection
\correction{003348}
$f(\vec x) = \alpha\vec x + \beta(\vec x)\vec u$, $\beta \in E^*$.
\fincorrection
\correction{003354}
\begin{enumerate}
  \item $\psi_{ij} \circ \psi_{k\ell} = \delta_{jk}\psi_{i\ell}$.
  \item $\psi_{11}$ est un projecteur non trivial.
  \item Si $\sum \lambda_k\vec u_k = \vec 0$, alors en appliquant
              $\psi_{1j}$ : $\lambda_j\vec u_1 = \vec 0  \Rightarrow  \lambda_j = 0$.
  \item Décomposer $g$ sur la base $(\varphi_{ij})$.
\end{enumerate}
 \fincorrection
\correction{003355}
$g =$ une projection sur $\Im f$ et $h=f$.
\fincorrection
\nocorrection
\nocorrection
\nocorrection
\nocorrection
\correction{003489}
$f^2 = f \circ g\circ f = f\circ g = f$ donc $f$ est une projection. $g$ idem.

$f\circ g = f  \Rightarrow  \mathrm{Ker} g \subset \mathrm{Ker} f$ et donc, par symétrie, $\mathrm{Ker} f = \mathrm{Ker} g$.

Réciproquement, si $f,g$ sont deux projections de même direction, alors
$f\circ g$ et $f$ coïncident sur la base et la direction de $g$, donc sont
égales. De même, $g\circ f = g$.
\fincorrection
\correction{003490}
Direction = $\mathrm{Ker} f$ et Base = $\mathrm{Im} g$.
\fincorrection
\correction{003491}
Direction = $\mathrm{Ker} p \cap \mathrm{Ker} q$ et Base = $\Im p \oplus \Im q$.
\fincorrection
\correction{003492}
Si $\lambda \neq 1$, $(\mathrm{id}-f)^{-1} = \mathrm{id} + \frac1{1-\lambda}f$.
\fincorrection
\nocorrection
\correction{003494}
\begin{enumerate}
  \item 
  \item $\begin{cases} 2x' = x - 2y - z  \cr
                  2y' = -x -z       \cr
                  2z' = -x -2y + z. \cr \end{cases}$
\end{enumerate}
 \fincorrection
\correction{003495}
Si $A=\varnothing$ c'est évident.\par
Sinon, $A$ est un sous-groupe de $GL(E)$ donc $\frac{u}{\mathrm{Card}\, A}$ est
un projecteur et $\mathrm{tr}(u) = \mathrm{Card}\,(A)\mathrm{rg}(u)$.
\fincorrection
\correction{005193}
\begin{enumerate}
\item  Soit $p(\in\Nn^*)$ l'indice de nilpotence de $u$.

Par définition, $u^{p-1}\neq0$ et plus généralement, pour $1\leq k\leq p-1$, $u^k\neq0$ car si $u^k= 0$ alors $u^{p-1}=u^k\circ u^{p-1-k}=0$ ce qui n'est pas.

Puisque $u^{p-1}\neq0$, il existe au moins un vecteur $x$ non nul tel que $u^{p-1}(x)\neq0$.

Montrons que la famille $(u^k(x))_{0\leq k\leq p-1}$ est libre.

Soit $(\lambda_k)_{0\leq k\leq p-1}\in\Kk^p$ tel que $\sum_{k=0}^{p-1}\lambda_ku^k(x)=0$. Supposons qu'au moins un des coefficients $\lambda_k$ ne soit pas nul. Soit $i=\mbox{Min }\{k\in\{0,...,p-1\}/\;\lambda_k\neq0\}$.

\begin{align*}
\sum_{k=0}^{p-1}\lambda_ku^k(x)=0&\Rightarrow\sum_{k=i}^{p-1}\lambda_ku^k(x)=0
\Rightarrow u^{p-1-i}(\sum_{k=i}^{p-1}\lambda_ku^k(x))=0\Rightarrow\sum_{k=i}^{p-1}\lambda_ku^{p-1-i+k}(x)=0\\
 &\Rightarrow\lambda_iu^{p-1}(x)=0\quad(\mbox{car pour}\;k\geq i+1,\;p-1-i+k\geq p\;\mbox{et donc}\;u^{p-1-i+k}=0)\\
 &\Rightarrow\lambda_i=0\quad(\mbox{car}\;u^{p-1}(x)\neq0)
\end{align*} 

ce qui contredit la définition de $i$.

Donc tous les coefficients $\lambda_k$ sont nuls et on a montré que la famille $(u^k(x))_{0\leq k\leq p-1}$ est libre.

\item  Le cardinal d'une famille libre est inférieur ou égal à la dimension de l'espace et donc $p\leq n$. 
Par suite, $u^n=u^p\circ u^{n-p}=0$.

\item  On applique l'exerice \ref{exo:suprou10}.

Puisque $u^{n-1}\neq0$, on a $N_{n-1}\underset{\neq}{\subset}N_n$.
Par suite (d'après l'exercice \ref{exo:suprou12}, 2), c)), les inclusions $N_0\subset N_1\subset...\subset N_n=E$ sont toutes strictes et donc 

$$0<\mbox{dim }N_1<\mbox{dim }N_2 ...<\mbox{dim }N_n=n.$$

Pour $k\in\{0,...,n\}$, notons $d_k$ est la dimension de $N_k$. Pour $k\in\{0,...,n-1\}$, on a $d_{k+1}\geq d_k$ et une récurrence facile montre que, pour $k\in\{0,...,n\}$, on a $d_k\geq k$.

Mais si de plus, pour un certain indice $i$ élément de $\{1,...,n-1\}$, on a $d_i=\mbox{dim }N_i>i$, alors, par une récurrence facile, pour $i\leq k\leq n$, on a $d_k>k$ et en particulier $d_n>n$ ce qui n'est pas. Donc,

$$\forall k\in\{0,...,n\},\;\mbox{dim }(N_k)=k,$$

ou encore, d'après le théorème du rang,

$\forall k\in\{0,...,n\},\;\mbox{rg }(u^k)=n-k$, et en particulier $\mbox{rg }(u)=n-1$.
\end{enumerate}
\fincorrection
\correction{005584}
\textbf{1ère solution.} Si $f=0$, c'est immédiat . Sinon, soit $p$ l'indice de nilpotence de $f$ ($p\geqslant2$).

Par définition de $p$, il existe un vecteur $x_0$ tel que $f^{p-1}(x_0)\neq 0$ (et $f^p(x_0)=0$).

Montrons que la famille $(f^k(x_0))_{0\leqslant k\leqslant p-1}$ est libre. Dans le cas contraire, il existe $a_0$,..., $a_{p-1}$ $p$ scalaires non tous nuls tels que $a_0x_0+...+a_{p-1}f^{p-1}(x_0)=0$.

Soit $k=\text{Min}\{i\in\llbracket0,p-1\rrbracket/\;a_i\neq0\}$.

\begin{align*}\ensuremath
\sum_{i=0}^{p-1}a_if^i(x_0)=0&\Rightarrow\sum_{i=k}^{p-1}a_if^i(x_0)=0\Rightarrow f^{p-1}\left(\sum_{i=k}^{p-1}a_if^i(x_0)\right)=0\\
 &\Rightarrow a_kf^{p-1}(x_0)=0\;(\text{car pour}\;i\geqslant p,\;f^i=0\\
 &\Rightarrow a_k=0\;(\text{car}\;f^{p-1}(x_0)\neq0).
\end{align*}

Ceci contredit la définition de $k$ et donc la famille $(f^k(x_0))_{0\leqslant k\leqslant p-1}$ est libre. Puisque le cardinal d'une famille libre est inférieur à la dimension de
l'espace, on a montré que $p\leqslant n$ ou, ce qui revient au même, $f^n=0$.

\textbf{2ème solution.} (pour les redoublants)

Soit $p\in\Nn^*$ l'indice de nilpotence de $f$. Le polynôme $X^p$ est annulateur de $f$. Son polynôme minimal est un diviseur de $X^p$ et donc égal à $X^k$ pour un certain $k\in\llbracket1,p\rrbracket$. Par définition de l'indice de nilpotence, $k=p$ puis $\mu_f=X^p$. D'après le théorème de \textsc{Cayley}-\textsc{Hamilton}, $\mu_f$ divise $\chi_f$ qui est de degré $n$ et en particulier $p\leqslant n$.
\fincorrection
\correction{005587}
On transforme légèrement l'énoncé.

Si $x$ est un vecteur non nul tel que $(x,f(x))$ est liée alors il existe un scalaire $\lambda_x$ tel que $f(x) =\lambda_x x$. Si $x = 0$, $f(x)=0=0x$ et encore une fois il existe un scalaire $\lambda_x$ tel que $f(x)=\lambda_x x$.

Inversement, si pour tout $x$ de $E$, il existe $\lambda_x\in\Kk$ tel que $f(x)=\lambda_xx$, alors  la famille $(x,f(x))$ est liée. Donc 

\begin{center}
$[(\forall x\in E,\;(x,f(x))\;\text{liée})\Leftrightarrow(\forall x\in E,\exists\lambda_x\in\Kk/\;f(x) =\lambda_xx)]$.
\end{center}

Notons de plus que dans le cas où $x\neq0$, la famille $(x)$ est une base de la droite vectorielle $\text{Vect}(x)$ et en particulier, le nombre $\lambda_x$ est uniquement défini.

Montrons maintenant que $f$ est une homothétie c'est à dire montrons que : $\exists\lambda\in\Kk/\;\forall x\in E,\;f(x) =\lambda x$.

Soient $x_0$ un vecteur non nul et fixé de $E$ puis $x$ un vecteur quelconque de $E$.

\textbf{1er cas.} Supposons la famille $(x_0,x)$ libre. On a $f(x+x_0)=\lambda_{x+x_0}(x+x_0)$ mais aussi $f(x+x_0)=f(x)+f(x_0) =\lambda_xx +\lambda_{x_0}x_0$ et donc

\begin{center}
$(\lambda_{x+x_0}-\lambda_x)x+ (\lambda_{x+x_0}-\lambda_{x_0})x_0=0$.
\end{center}

Puisque la famille $(x_0,x)$ est libre, on obtient  $\lambda_{x+x_0}-\lambda_x=\lambda_{x+x_0}-\lambda_{x_0}= 0$ et donc $\lambda_x=\lambda_{x+x_0}=\lambda_{x_0}$. Ainsi, pour tout vecteur $x$ tel que $(x,x_0)$ libre, on a $f(x)=\lambda_{x_0}x$.

 
\textbf{2ème cas.} Supposons la famille $(x_0,x)$ liée. Puisque $x_0$ est non nul, il existe un scalaire $\mu$ tel que $x =\mu x_0$. Mais alors 

\begin{center}
$f(x) =\mu f(x_0)=\mu \lambda_{x_0}x_0 =\lambda_{x_0}x$.
\end{center}

Finalement, il existe un scalaire $k=\lambda_{x_0}$ tel que pour tout vecteur $x$, $f(x) =kx$ et $f$ est une homothétie. La réciproque étant claire, on a montré que

\begin{center}
\shadowbox{
$\forall f\in\mathcal{L}(E)$, $[(f\;\text{homothétie})\Leftrightarrow(\forall x\in E,\;(x,f(x))\;\text{liée})]$.
}
\end{center}

\fincorrection
\correction{005588}
Remarques. 1) Soit $(G,*)$ un groupe. Le centre de $G$ est l'ensemble des éléments de $G$ qui commutent avec tous les éléments de $G$. Ce centre, souvent noté $Z$, est un sous-groupe de $(G,*)$.

2) $(\mathcal{L}(E),\circ)$ est un magma associatif et unitaire mais non commutatif (pour $\text{dim}E> 1)$ mais $(\mathcal{L}(E),\circ)$ n'est pas un groupe. Par contre $(\mathcal{GL}(E),\circ)$ est un groupe (groupe des inversibles de $(\mathcal{L}(E),\circ)$).

Soit $f$ un endomorphisme (resp. automorphisme) de E commutant avec tous les endomorphismes (resp. les automorphismes) de $E$. $f$ commute en particulier avec toutes les symétries.

Soit $x$ un vecteur non nul de $E$ et $s$ la symétrie par rapport à $\text{Vect}(x)$ parallèlement à un supplémentaire donné de $\text{Vect}(x))$.

\begin{center}
$s(f(x)) =f(s(x))=f(x)$.
\end{center}

Par suite, $f(x)$ est invariant par $s$ et appartient donc à $Vect(x)$. Ainsi, si $f$ commute avec tout endomorphisme (resp. automorphisme) de $E$, $f$ vérifie nécessairement $\forall x\in E$, $(x,f(x))$ liée et d'après le \ref{ex:rou25}, $f$ est nécessairement une homothétie.
Réciproquement , les homothéties de $E$ commutent effectivement avec tout endomorphisme de $E$. 

\begin{center}
\shadowbox{
Les endomorphismes de $E$ qui commutent avec tout endomorphismes de $E$ sont les homothéties.
}
\end{center}

Pour le centre de $\mathcal{GL}(E)$, il faut enlever l'application nulle qui est une homothétie mais qui n'est pas inversible.
\fincorrection
\correction{005589}
$\Rightarrow$/ Si $p+q$ est un projecteur alors l'égalité $(p+q)^2=p+q$ founit $pq+qp=0$. En composant par $p$ à droite ou à gauche , on obtient  $pqp+qp=0=pq+pqp$ et donc  $pq = qp$.

 
Cette égalité jointe à l'égalité $pq + qp = 0$ fournit $pq = qp = 0$.

$\Leftarrow$/ Si $ pq = qp = 0$, alors $(p+q)^2=p^2 + pq + qp + q^2 = p + q$ et $p + q$ est un projecteur.

\begin{center}
\shadowbox{
Pour tous projecteurs $p$ et $q$, ($p+q$ projecteur$\Leftrightarrow p\circ q=q\circ p=0\Leftrightarrow\text{Im}q\subset\text{Ker}p\;\text{et}\;\text{Im}p\subset\text{Ker}q$).
}
\end{center}

Dorénavant, $p+q$ est un projecteur ou ce qui revient au même $pq=qp=0$.

On a $\text{Ker}p\cap\text{Ker}q\subset\text{Ker}(p+q)$. Inversement, pour $x\in E$,

\begin{center}
$x\in\text{Ker}(p+q)\Rightarrow(p+q)(x)=0\Rightarrow p(p(x)+q(x))=0\Rightarrow p(x)=0$,
\end{center}

et de même $q(x)=0$. Ainsi, $\text{Ker}(p+q)\subset\text{Ker}p\cap\text{Ker}q$ et donc $\text{Ker}(p+q)=\text{Ker}p\cap\text{Ker}q$.

On a $\text{Im}(p+q)\subset\text{Im}p+\text{Im}q$. Inversement, pour $x\in E$,

\begin{center}
$x\in\text{Im}p+\text{Im}q\Rightarrow\exists(x_1,x_2)\in E^2/\;x=p(x_1)+q(x_2)$.
\end{center}

Mais alors, $(p+q)(x)=p^2(x_1)+pq(x_1)+qp(x_2)+q^2(x_2)=p(x_1)+q(x_2)=x$ et donc $x\in\text{Im}(p+q)$. Ainsi, $\text{Im}p+\text{Im}q\subset\text{Im}(p+q)$ et donc $\text{Im}(p+q)=\text{Im}p+\text{Im}q$. En résumé, si $p$ et $q$ sont deux projecteurs tels que $p+q$ soit un projecteur, alors

\begin{center}
\shadowbox{
$\text{Ker}(p+q)=\text{Ker}p\cap\text{Ker}q$ et $\text{Im}(p+q)=\text{Im}p+\text{Im}q$.
}
\end{center}
\fincorrection
\correction{005590}
Soit $p$ un projecteur de $E$. Si $p=0$, $\text{Tr}(p)=\text{rg}(p)=0$ et si $p=Id_E$, $\text{Tr}(p)=\text{rg}(p)=n$.

Dorénavant, $p$ est un projecteur de rang $r\in\llbracket1,n-1\rrbracket$. On choisit une base de $E$ $\mathcal{B}$ adaptée à la décomposition $E=\text{Im}(p)\oplus\text{Ker}(p)$. Dans cette base, la matrice de $p$ s'écrit $\left(
\begin{array}{cccccc}
1&0&\ldots& &\ldots&0\\
0&\ddots&\ddots& & &\vdots\\
\vdots&\ddots&1& & & \\
 & & &$0$&\ddots&\vdots\\
\vdots& & &\ddots&\ddots&0\\
0&\ldots& &\ldots&0&0
\end{array}
\right)$ où le nombre de $1$ est $\text{dim}(\text{Im}(p))=r$. Mais alors $\text{Tr}(p)=r$.

\begin{center}
\shadowbox{
En dimension finie, la trace d'un projecteur est son rang.
}
\end{center}
\fincorrection
\correction{005591}
$\Leftarrow$/ Si $\forall i\neq j$, $p_i\circ p_j= 0$ alors 

\begin{center}
$(p_1+...+p_n)^2=p_1^2+...+p_n^2+\sum_{i\neq j}^{}p_i\circ p_j = p_1+...+p_n$,
\end{center}

et $p_1+...+p_n$ est un projecteur.

$\Rightarrow$/ Supposons que $p=p_1 + ... + p_n$ soit un projecteur. Posons $F_i=\text{Im}p_i$, $1\leqslant i\leqslant n$, puis $F=F_1 + ... + F_n$ et $G =\text{Im}p$.
On sait que la trace d'un projecteur est son rang. Par linéarité de la trace, on obtient

\begin{center}
$\text{rg}p=\text{Tr}p=\text{Tr}(p_1) + ...+\text{Tr}(p_n)=\text{rg}(p_1)+... +\text{rg}(p_n)$,
\end{center}

et donc $\text{dim}G=\text{dim}F_1+...+\text{dim}F_n\geqslant\text{dim}F$. D'autre part, $G=\text{Im}(p_1 + ... + p_n)\subset\text{Im}p_1 + ... +\text{Im}p_n=F_1 + ... +F_n =F$.

On obtient donc $G = F$ et aussi $\text{dim}(F_1+...+F_n)=\text{dim}F_1+...+\text{dim}F_n$. D'après l'exercice \ref{ex:rou15}, $F=F_1\oplus...\oplus F_n$ c'est-à-dire 

\begin{center}
$\text{Im}p =\text{Im}(p_1)\oplus ...\oplus\text{Im}(p_n)$.
\end{center}

Il reste à vérifier que pour $i\neq j$ et $x$ dans $E$, $p_i(p_j(x))=0$ ou ce qui revient au même que pour $i\neq j$ et $y$ dans $\text{Im}(p_j)$, $p_i(y)=0$.

Soit $y$ dans $\text{Im}(p_j)$ (et donc dans $\text{Im}p$). Les égalités $y=p_j(y)=p(y)$ fournissent  $\sum_{i\neq j}^{}p_i(y) = 0$. La somme $\sum_{i}^{}\text{Im}(p_i)$ étant directe, on a donc $p_i(y) = 0$ pour chaque $i\neq j$ ce qu'il fallait démontrer.

\begin{center}
\shadowbox{
$p_1+\ldots+p_n$ projecteur$\Leftrightarrow\forall i\neq j,\;p_i\circ p_j=0$.
}
\end{center}
\fincorrection
\correction{005592}
\begin{enumerate}
 \item  D'après le \ref{ex:rou29}, $\text{Im}(p_1 + ... + p_n) =\text{Im}(p_1)+ ... +\text{Im}(p_n)=\text{Im}(p_1)\oplus ... \oplus\text{Im}(p_n)$.

Chaque $p_i$ est de rang au moins $1$, mais si l'un des $p_i$ est de rang supérieur ou égal à $2$ alors $n =\text{dim}E\geqslant\text{rg}(p_1+...+p_n) =\text{rg}(p_1)+ ... +\text{rg}(p_n)> n$ ce qui est impossible. Donc chaque $p_i$ est de rang $1$.

\item  Les images des $p_i$ (resp. $q_i$) sont des droites vectorielles. Pour chaque $i$, notons $e_i$ (resp. $e_i'$) un vecteur non nul de $\text{Im}(p_i)$ (resp. $\text{Im}(q_i)$). D'après 1), $E=
\text{Vect}(e_1)\oplus  ... \oplus\text{Vect}(e_n)$ ou encore $(e_i)_{1\leqslant i\leqslant n}$ (resp. $(e_i')_{1\leqslant i\leqslant n}$) est une base de $E$.

Soit $f$ l'automorphisme de $E$ défini par $f(e_i) =e_i'$ ($f$ est un automorphisme car l'image par $f$ d'une base de $E$ est une base de $E$).

Soit $(i,j)\in\llbracket1,n\rrbracket^2$. $f\circ p_i\circ f^{-1}(e_j') = f(p_i(e_j)) =f(\delta_{i,j}e_i)=\delta_{i,j}e_i'=q_i(e_j)$. Ainsi, les endomorphismes $q_i$ et $f\circ p_i\circ f^{-1}$ coïncident sur une base de $E$ et sont donc égaux.
\end{enumerate}
\fincorrection
\correction{005593}
Soit $q=\frac{1}{n}\sum_{g\in G}^{}g\circ p\circ g^{-1}$.

\begin{center}
$q^2=\frac{1}{n}\sum_{(g,h)\in G^2}^{}h\circ p\circ h^{-1}\circ g\circ p\circ g^{-1}$.
\end{center}

Mais si $g$ et $h$ sont deux éléments de $G$ et $x$ est un vecteur quelconque de $E$, $p(g^{-1}(x))$ est dans $F$ et donc par hypothèse $h^{-1}\circ g\circ p\circ g^{-1}(x)$ est encore dans $F$ ($h^{-1}$ est dans $G$ puisque $G$ est un groupe). On en déduit que

\begin{center}
$h\circ p\circ h^{-1}\circ g\circ p\circ g^{-1}=h\circ h^{-1}\circ g\circ p\circ g^{-1}=g\circ p\circ g^{-1}$.
\end{center}

Mais alors

\begin{center}
$q^2=\frac{1}{n^2}\sum_{(g,h)\in G^2}^{}g\circ p\circ g^{-1}=\frac{1}{n^2}\times n\sum_{g\in G}^{}g\circ p\circ g^{-1}=\frac{1}{n}\sum_{g\in G}^{}g\circ p\circ g^{-1}= q$
\end{center}

et $q$ est un projecteur.

Montrons que $F\subset\text{Im}q$. Soit $x$ un élément de $F$. Pour chaque $g\in G$, $g^{-1}(x)$ est encore dans $F$ et donc $p(g^{-1}(x))=g^{-1}(x)$ puis $g(p(g^{-1}(x)))=x$. Mais alors

\begin{center}
$q(x)=\frac{1}{n}\sum_{g\in G}^{}x=x$,
\end{center}

ou encore $x$ est dans $\text{Im}q$. On a montré que $F\subset\text{Im}q$.

Montrons que $\text{Im}q\subset F$. Soit $x$ un élément de $\text{Im}q$. 

\begin{align*}\ensuremath
p(x)&=p(q(x))=\frac{1}{n}\sum_{g\in G}^{}p\circ g\circ p\circ g^{-1}(x)\\
 &=\frac{1}{n}\sum_{g\in G}^{}g\circ p\circ g^{-1}(x)\;(\text{car}\;p\circ g^{-1}(x)\in F\;\text{et donc}\;g\circ p\circ g^{-1}(x)\in F)\\
 &=q(x) = x,
\end{align*}

et $x$ est dans $F$. On a montré que $\text{Im}q\subset F$ et finalement que $\text{Im}q=F$.

\begin{center}
\shadowbox{
$q$ est un projecteur d'image $F$.
}
\end{center}
\fincorrection
\correction{005598}
Deux cas particuliers se traitent immédiatement.

Si $f=0$, on prend $p=0$ et $g=Id_E$ et si $f\in\mathcal{GL}(E)$, on prend $p=Id_E$ et $g=f$. 

On se place dorénavant dans le cas où $\text{Ker}f$ et $\text{Im}f$ ne sont pas réduit à ${0}$.

Soit $F$ un supplémentaire de $\text{Ker}f$ dans $E$ et $G$ un  supplémentaire de $\text{Im}f$ dans $E$.

On sait que la restriction $f'$ de $f$ à $F$ réalise un isomorphisme de $F$ sur $\text{Im}f$. D'autre part $\text{dim}\text{Ker}f=\text{dim}G<+\infty$ et donc $\text{Ker}f$ et $G$ sont isomorphes. Soit $\varphi$ un isomorphisme de $\text{Ker}f$ sur $G$.

On définit une unique application linéaire $g$ en posant $g_{/\text{Ker}f}=\varphi$ et $g_{/F}=f'$.

$g$ est un automorphisme de $E$. En effet,

\begin{center}
$g(E)=g(\text{Ker}f +F)=g(\text{Ker}f) +g(F)=\varphi(\text{Ker}f) + f'(F)=G +\text{Im}f=E$,
\end{center}

(puisque $\varphi$ et $f'$ sont des isomorphismes) et donc $g$ est surjective. Par suite $g$ est bijective de $E$ sur lui-même puisque $\text{dim}E<+\infty$.

Soit $p$ la projection sur $F$ parallèlement à $\text{Ker}f$. On a

\begin{center}
$(g\circ p)_{/\text{Ker}f}= g\circ0_{/\text{Ker}f}= 0_{/\text{Ker}f}=f_{/\text{Ker}f}$ et $(g\circ p)_{/F}=g\circ Id_{/F}=f'=f_{/F}$.
\end{center}

Ainsi les endomorphismes $g\circ p$ et $f$ coïncident sur deux sous espaces supplémentaires de $E$ et donc $g\circ p=f$. Finalement, si on note $P(E)$ l'ensemble des projecteurs de $E$,

\begin{center}
\shadowbox{
$\forall f\in\mathcal{L}(E)$, $\exists g\in\mathcal{GL}(E)$, $\exists p\in P(E)/\;f=g\circ p$.
}
\end{center}
\fincorrection
\correction{005599}
(Ne pas confondre : ($\forall x\in E,\; \exists p\in\Nn^*/\;f^p(x)=0$) et ($\exists p\in\Nn^*/\;\forall x\in E,\;f^p(x) = 0)$. Dans le deuxième cas, $p$ est indépendant de $x$ alors que dans le premier cas, $p$ peut varier quand $x$ varie).

Soit $\mathcal{B}=(e_i)_{1\leqslant i\leqslant n}$ une base de $E$. Pour chaque $i\in\llbracket1,n\rrbracket$, il existe un entier non nul $p_i$ tel que $f^{p_i}(e_i)=0$. Soit $p=\text{Max}\{p_1,...,p_n\}$. $p$ est un entier naturel non nul et pour $i$ dans $\llbracket1,n\rrbracket$, on a

\begin{center}
$f^{p}(e_i)=f^{p-p_i}(f^{p_i}(e_i))=f^{p-p_i}(0)=0$.
\end{center}

Ainsi l'endomorphisme $f^p$ s'annule sur une base de $E$ et on sait que $f^p=0$.

On a donc trouvé un entier non nul $p$ tel que $f^p=0$ et par suite $f$ est nilpotent.
\fincorrection
\correction{005600}
\begin{enumerate}
 \item  A partir de $fg-gf=af+bg$ $(1)$, on obtient après composition à droite par $g$, $fg-gfg = afg + bg$ ou encore $fg=g\circ\frac{1}{1-a}(fg + bId)$ (puisque $1 - a\neq0$). On en déduit

\begin{center} 
$\text{Im}(fg)\subset\text{Im}g$.
\end{center}

Mais alors en écrivant $(1)$ sous la forme $f=\frac{1}{a}(fg- gf -bg)$ (puisque $a$ n'est pas nul), on obtient

\begin{center}
$\text{Im}f\subset\text{Im}g$.
\end{center}

L'égalité $\text{Im}f\subset\text{Im}g$ montre que tout vecteur de $\text{Im}f$ est invariant par $g$ et  fournit donc l'égalité $gf = f$. On compose alors (1) à droite par $f$ et en tenant compte de $gf = f$ et de $f^2 = f$, on obtient $f - f = af + bf$ et donc $(a+b)f = 0$ puis $b=-a$ puisque $f$ n'est pas nul.

(1) s'écrit alors $fg - f = a(f-g)$. En composant à droite par $g$, on obtient : $a(fg - g) = 0$ et donc $fg = f$ puisque $a$ n'est pas nul. (1) s'écrit maintenant $g-f = a(f-g)$ ou encore $(a+1)(g - f) = 0$ et donc, puisque $f$ et $g$ sont distincts, $a = -1$.

\item  (D'après 1), si a est distinct de $0$ et de $1$, nécessairement $a = -1$ et (1) s'écrit $fg - gf =-f +bg$).

Soit $x$ un élément de $\text{Ker}g$. $(1)$ fournit $-g(f(x)) = af(x)$ $(*)$ puis en prenant l'image par $g$, $(a+1)g(f(x)) = 0$. Puisque $a$ est distinct de $-1$, on obtient $g(f(x)) = 0$ et $(*)$ fournit $af(x) = 0$ puis $f(x) = 0$. Donc $x$ est élément de $\text{Ker}f$. On a montré que $\text{Ker}g\subset\text{Ker}f$.

On en déduit $\text{Im}(g-Id)\subset\text{Ker}f$ et donc $f(g-Id) =0$ ou encore $fg = f$. $(1)$ s'écrit $f - gf = af + bg$ et en composant à gauche par $f$, on obtient $f - fgf = af + bfg$. En tenant compte de $fg = f$, on obtient  $(a+b)f = 0$ et donc $b = -a$.

$(1)$ s'écrit alors $f - gf = a(f - g)$ et en composant à gauche par $g$, on obtient $0 = a(gf - g)$ et donc $gf = g$. (1) s'écrit enfin $f-g = a(f-g)$ et donc $a = 1$.

\item  Si $a = 0$, $(1)$ s'écrit $fg - gf = bg$. En composant à gauche ou à droite par $g$, on obtient $gfg - gf = bg$ et $fg - gfg = bg$. En additionnant ces deux égalités, on obtient $fg - gf = 2bg$. D'où, en tenant compte de $(1)$, $bg = 2bg$ et puisque $g$ n'est pas nul, $b = 0$. Par suite $fg - gf = 0$ ce qui est exclu par l'énoncé. Donc, on ne peut avoir $a = 0$. D'après 1) et 2), $(a,b)\in\{(-1,1),(1,-1)\}$.

\textbf{1er cas.} $(a,b) = (-1,1)$. C'est le 1) : $fg - gf = -f +g$. On a vu successivement que $gf = f$  puis que $fg = g$ fournissant $(g-Id)f = 0$ et $(f-Id)g = 0$ ou encore 
$\text{Im}f \subset\text{Ker}(g-Id) =\text{Im}g$ et $\text{Im}g \subset\text{Im}f$ et donc $\text{Im}f =\text{Im}g$. Réciproquement, si $f$ et $g$ sont deux projecteurs de même image alors $gf = f$, $fg = g$ et donc $fg - gf = -f + g$. Le premier cas est donc le cas de deux projecteurs de même image .

\textbf{2ème cas.} $(a,b) = (1,-1)$. C'est le cas de deux projecteurs de même noyau.
\end{enumerate}
\fincorrection
\correction{003315}
Non, ils n'ont pas même dimension si $E\ne\{\vec 0\}$ ou
$F\ne\{\vec 0\}$.
\fincorrection
\correction{003356}
On veut $\Im g \subset\mathrm{Ker} f$ et $\mathrm{Ker} g\supset\Im f$
donc $g$ est entièrement définie par sa restriction à un supplémentaire
de~$\Im f$, application linéaire à valeurs dans~$\mathrm{Ker} f$.
On en déduit $\dim F = (\mathrm{codim}\Im f)(\dim\mathrm{Ker} f) = (\dim\mathrm{Ker} f)^2$.
\fincorrection
\correction{003357}
Soit $p=\frac1{\mathrm{Card}\, G}\sum_{g\in G}g$. Alors $g\circ p = p$,
pour tout $g\in G$ donc $p^2=p$,
$F\subset \Im p$ et si $x\in\Im p$, on a $p(x)=x$ d'où $g(x) = x$ pour tout $g\in G$
c'est-à-dire $x\in F$. Donc $F = \Im p$ et $\dim F = \mathrm{rg}(p) = \mathrm{tr}(p)$
(trace d'un projecteur).
\fincorrection
\correction{001040}
Si $C = A \times B$ alors on obtient le coefficient $c_{ij}$ (situé à la $i$-ème ligne et la $j$-ème colonne
de $C$) en effectuant le produit scalaire du $i$-ème vecteur-ligne de $A$ avec le $j$-éme vecteur colonne de $B$.

On trouve  
$$\left( \begin{array}{cc} 2 & 1  \\ 3& 2 \end{array} \right)\times
 \left( \begin{array}{cc}1 & -1  \\ 1& 2 \end{array} \right)
= \begin{pmatrix}
  3 & 0 \\
  5 & 1 \\  
  \end{pmatrix}$$

$$\left( \begin{array}
 {ccc} 1 & 2 & 0  \\ 3 & 1 & 4 \end{array} \right) 
\times
\left( \begin{array}{ccc} -1 &-1& 0 \\ 1 & 4 & -1 \\ 2 & 1 & 2\end{array} \right)
= \begin{pmatrix}
  1 & 7 & -2 \\
  6 & 5 & 7 \\ 
  \end{pmatrix}$$

$$\left( \begin{array}{ccc} a &b& c \\ c & b & a \\ 1 & 1 & 1\end{array} \right)
\times
 \left( \begin{array}{ccc} 1 &a& c \\ 1 & b & b \\ 1 & c & a\end{array} \right) 
 = \begin{pmatrix}
    a + b + c & a^2+b^2+c^2 & 2ac + b^2 \\
    a + b + c & 2ac + b^2   & a^2+b^2+c^2 \\
    3         &  a+b+c      & a+b+c \\
\end{pmatrix}$$
\fincorrection
\nocorrection
\nocorrection
\nocorrection
\nocorrection
\nocorrection
\nocorrection
\nocorrection
\nocorrection
\nocorrection
\nocorrection
\nocorrection
\correction{001052}

On trouve 
$$A^2 =
\begin{pmatrix} 
3 & -4 & 2 \cr
1 & -1 & -1 \cr
1 & 2 & 0 \cr
\end{pmatrix}
\qquad \text{ et } \qquad 
A^3=
\begin{pmatrix} 
5 & 0 & 2 \cr
0 & 3 & 1 \cr
1 & -2 & 4 \cr
\end{pmatrix}
.$$

Un calcul donne $A^3-A = 4 I$.
En factorisant par $A$ on obtient $A\times (A^2-I) = 4I$.
Donc $A \times \frac 1 4 (A^2-I) = I$, ainsi $A$ est inversible et
$$A^{-1} = \frac 1 4 (A^2-I) = \frac 1 4
\begin{pmatrix}
2&-4&2\\
1&-2&-1\\
1&2&-1\\
\end{pmatrix}.
$$
\fincorrection
\correction{001053}
$$ A =
\begin{pmatrix}
0&1\\
0&0\\
\end{pmatrix}, \qquad
B =
\begin{pmatrix}
1&0\\
0&0\\
\end{pmatrix}.
$$
\fincorrection
\nocorrection
\nocorrection
\correction{001057}
$F$  est un sous espace vectoriel de  $M_2 ({\R })$  donc
$\hbox{dim }(F) \in \{ 0,\ldots , 4\} $. Comme  $F\not = M_2 ({\R
})$  on a aussi  $\hbox{dim }(F) \not = 4$. D'autre part les
matrices  $M_1=\begin{pmatrix} 0 & 1 \cr
                                      0 & 0 \end{pmatrix} ,  M_2=\begin{pmatrix} 0 & 0 \cr
                                      1 & 0 \end{pmatrix} ,
M_3=\begin{pmatrix} 1 & 0 \cr
                                      0 & -1 \end{pmatrix} $  appartiennent \`a  $F$  et sont lin\' eairement
ind\' ependantes. En effet, si  $\alpha M_1 + \beta M_2 + \gamma
M_3 =0 $  alors  $\begin{pmatrix}  \gamma & \alpha \cr
                      \beta & -\gamma \cr \end{pmatrix}= \begin{pmatrix} 0 & 0 \cr 0 & 0 \end{pmatrix}$  c'est \`a dire  $\alpha = \beta =\gamma =0$. Donc
$\hbox{dim }(F)\geq 3 $  c'est \`a dire $\hbox{dim }(F)= 3 $.
Enfin $\{ M_1
 ,
 M_2  ,  M_3 \} $  est une famille libre de trois vecteurs dans  $F$  qui est un espace de
dimension $3$. C'est donc une base de  $F$.
\fincorrection
\nocorrection
\nocorrection
\nocorrection
\correction{001061}
\begin{align*}
A(\theta)\times A(\theta')  
  & =  \begin{pmatrix} \cos \theta & -\sin \theta \\ \sin \theta
& \cos \theta \end{pmatrix} \times \begin{pmatrix} \cos \theta' & -\sin \theta' \\ \sin \theta'
& \cos \theta' \end{pmatrix} \\
  & = \begin{pmatrix} 
\cos \theta\cos \theta' -\sin \theta \sin\theta' & - \cos \theta \sin \theta' - \sin \theta \cos \theta' \\
 \sin \theta \cos \theta'+\cos \theta \sin \theta'   & -\sin \theta \sin\theta'+\cos \theta\cos \theta'  \\
 \end{pmatrix} \\
  & =
\begin{pmatrix} \cos (\theta+\theta') & -\sin (\theta+\theta')  \\ \sin  (\theta+\theta')
& \cos  (\theta+\theta') \end{pmatrix} \\
  & = A(\theta+\theta') \\
\end{align*}

Bilan : $A(\theta)\times A(\theta') = A(\theta+\theta')$.

\bigskip

Nous allons montrer par récurrence sur $n\ge 1$ que $\big(A(\theta)\big)^n = A(n\theta)$.

\begin{itemize}
  \item C'est bien sûr vrai pour $n=1$.
  \item Fixons $n\ge 1$ et supposons que $\big(A(\theta)\big)^n = A(n\theta)$ alors
$$\big(A(\theta)\big)^{n+1} = \big(A(\theta)\big)^n \times A(\theta) = A(n\theta) \times A(\theta) = A(n\theta+\theta) = A((n+1)\theta)$$
  \item C'est donc vrai pour tout $n\ge 1$.
\end{itemize}


Remarques :
\begin{itemize}
  \item On aurait aussi la formule $A(\theta')\times A(\theta) = A(\theta+\theta') = A(\theta)\times A(\theta')$.
Les matrices $A(\theta)$ et $A(\theta')$ commutent.

  \item En fait il n'est pas plus difficile de montrer que $\big(A(\theta)\big)^{-1}=A(-\theta)$.
On sait aussi que par définition $\big(A(\theta)\big)^{0}=I$. Et on en déduit que pour $n\in \Zz$ on a 
$\big(A(\theta)\big)^n = A(n\theta)$.

  \item En terme géométrique $A(\theta)$ est la matrice de la rotation d'angle $\theta$ (centrée à l'origine).
On vient de montrer que si l'on compose un rotation d'angle $\theta$ avec un rotation d'angle $\theta'$
alors on obtient une rotation d'angle $\theta+\theta'$.
\end{itemize}

\fincorrection
\nocorrection
\correction{001063}
Notons $E_{ij}$ la matrice élémentaire (des zéros partout sauf le coefficient $1$ à la
 $i$-ème ligne et la $j$-ème colonne).

Soit $A = (a_{ij})  \in \mathcal{M}_n(\Rr)$.
Alors 
$$A \times E_{ij} = 
\begin{pmatrix}  
0 & 0 & \cdots & 0 & a_{1i} & 0 & \cdots \\
0 & 0 & \cdots & 0 & a_{2i} & 0 & \cdots \\
\vdots&  & \cdots & & \vdots &  & \cdots \\
0 & 0 & \cdots & 0 & a_{ji} & 0 & \cdots \\
\vdots&  & \cdots & & \vdots &  & \cdots \\
0 & 0 & \cdots & 0 & a_{ni} & 0 & \cdots \\
\end{pmatrix}$$
La seule colonne non nulle est la $j$-ème colonne.

La trace est la somme des éléments sur la diagonale. Ici le seul élément non nul de la diagonale est
$a_{ji}$, on en déduit donc 
$$\text{tr} (A \times E_{ij})=a_{ji}$$
(attention à l'inversion des indices).

Maintenant prenons deux matrices $A, B$ telles que $\text{tr} (AX) = \text{tr} (BX)$
pour toute matrice $X$. Alors pour $X=E_{ij}$ on en déduit $a_{ji}=b_{ji}$.
On fait ceci pour toutes les matrices élémentaires $E_{ij}$ avec $1 \le i,j \le n$
ce qui implique $A=B$.
\fincorrection
\correction{001064}
Notons $A=(a_{ij})$, notons $B = {}^{t}\!{A}$ si les coefficients sont $B=(b_{ij})$ 
alors par définition de la transposée on a $b_{ij}= a_{ji}$.

Ensuite notons $C = A \times B$ alors par définition du produit de matrices 
le coefficients $c_{ij}$ de $C$ s'obtient par la formule :
$$c_{ij} = \sum_{k=1}^n a_{ik}b_{kj}.$$

Appliquons ceci avec $B = {}^{t}\!{A}$
$$c_{ij} = \sum_{k=1}^n a_{ik}b_{kj} = \sum_{k=1}^n a_{ik}a_{jk}.$$
Et pour un coefficient de la diagonale on a $i=j$ donc 
$$c_{ii} =  \sum_{k=1}^n a_{ik}^2.$$

La trace étant la somme des coefficients sur la diagonale on a :
$$\text{tr} (A \ {}^{t}\!{A}) = \text{tr} (C) 
= \sum_{i=1}^n c_{ii} =  \sum_{i=1}^n \sum_{k=1}^n a_{ik}^2 = \sum_{1\le i,k \le n} a_{ik}^2.$$

Si on change l'indice $k$ en $j$ on obtient 
$$\text{tr} (A \ {}^{t}\!{A})  = \sum_{1\le i,j \le n} a_{ij}^2.$$

Donc cette trace vaut la somme des carrés de tous les coefficients.

Conséquence : si $\text{tr} (A \ {}^{t}\!{A}) = 0$ alors la somme des carrés 
$\sum_{1\le i,j \le n} a_{ij}^2$ est nulle 
donc chaque carré $a_{ij}^2$ est nul. Ainsi $a_{ij}=0$ (pour tout $i,j$) autrement dit $A$ est la matrice nulle.
\fincorrection
\nocorrection
\nocorrection
\nocorrection
\nocorrection
\correction{001069}
Soit $X=\left(
\begin{array}{c}
x_1\\
\vdots\\
x_n
\end{array}
\right)$  un vecteur tel que $AX=0$. 
Nous allons montrer qu'alors $X$ est le vecteur nul ce qui entraîne que $A$ est inversible.

Par l'absurde supposons $X\neq0$. 
Alors, si $i_0$ est un indice tel que $|x_{i_0}|=\max\big\{|x_i| \mid i=1,\ldots,n\big\}$, on a  $|x_{i_0}|>0$.

Mais alors comme $AX=0$ on a pour tout $i=1,\ldots,n$ :
$$\sum_{j=1}^{n}a_{i,j}x_j=0$$
donc 
$$|a_{i_0,i_0}x_{i_0}|=\left|-\sum_{j\neq i_0}^{}a_{i_0,j}x_j\right|\leq\sum_{j\neq i_0}^{}|a_{i_0,j}|.|x_j|
 \leq|x_{i_0}|\sum_{j\neq i_0}^{}|a_{i_0,j}|
$$
et, puisque $|x_{i_0}|> 0$, on obtient $|a_{i_0,i_0}|\leq\sum_{j\neq i_0}^{}|a_{i_0,j}|$ contredisant les hypothèses de l'énoncé. 
Ainsi $X=0$. On a donc prouvé <<$AX=0 \Rightarrow X=0$>> ce qui équivaut à $A$ inversible.
\fincorrection
\nocorrection
\nocorrection
\nocorrection
\nocorrection
\nocorrection
\nocorrection
\nocorrection
\nocorrection
\nocorrection
\nocorrection
\correction{001080}
\begin{enumerate}
\item $A=
  \begin{pmatrix}
    -9&-18\\6&12
  \end{pmatrix}$
\item $U_n=A^nU_0$
\item C'est la droite engendrée par $
  \begin{pmatrix}
    -2\\1
  \end{pmatrix}$. Le rang est 1.
\item C'est la droite engendrée par $
  \begin{pmatrix}
    -3\\2
  \end{pmatrix}$.
\item Ce sont deux vecteurs non colinéaires. On a
$$P^{-1}AP=D=
\begin{pmatrix}
  3&0\\0&0
\end{pmatrix}$$
\item On a $A=PDP^{-1}$ donc $A^n=PD^nP^{-1}=
  \begin{pmatrix}
   -3^{n+1} & -2\cdot 3^{n+1} \\
2\cdot 3^n & 4\cdot 3^n \\
  \end{pmatrix}
$
\item Donc
$$\Big\{
\begin{array}{rcc}
x_n &=& -137\cdot 3^{n+1}-36\cdot 3^{n+1}\\
y_n &=& 274(3^n)+72\cdot 3^n
\end{array}$$
\end{enumerate}
\fincorrection
\correction{002442}
\begin{enumerate}
  \item Notons $C=AB$ et $D=BA$.
Alors par la définition du produit de matrice :
$$c_{ij}=\sum_{1\le k \le n} a_{ik}b_{kj} \quad \text{ donc } c_{ii}=\sum_{1\le k \le n} a_{ik}b_{ki}$$
Ainsi 
$$\textrm{tr}(AB) = \textrm{tr}\, C = \sum_{1\le i \le n} c_{ii} = \sum_{1\le i \le n} \sum_{1\le k \le n} a_{ik}b_{ki}$$

De même $$\textrm{tr}(BA) = \textrm{tr}\, D = \sum_{1\le i \le n} \sum_{1\le k \le n} b_{ik}a_{ki}$$
Si dans cette dernière formule on renomme l'indice $i$ en $k$ et l'indice $k$ en $i$ (ce sont des variables muettes 
donc on leur donne le nom qu'on veut) alors on obtient :
$$\textrm{tr}(BA) =\sum_{1\le k \le n} \sum_{1\le i \le n} b_{ki}a_{ik} =  \sum_{1\le i \le n}\sum_{1\le k \le n} a_{ik}b_{ki} =\textrm{tr}(AB)$$

  \item $M$ et $M'$ sont semblables donc il existe une matrice de passage $P$ telle que $M'=P^{-1}MP$ donc
$$\textrm{tr}\, M' = \textrm{tr}\big( P^{-1}(MP) \big) = \textrm{tr}\big( (MP)P^{-1} \big) = \textrm{tr} ( M I ) = \textrm{tr}\, M$$

  \item La trace a aussi la propriété évidente que 
$$\textrm{tr}(A+B)=\textrm{tr}\,A+\textrm{tr}\,B.$$

Fixons une base de $E$. Notons $A$ la matrice de $f$ dans cette base et $B$ la matrice de $g$ dans cette même base.
Alors $AB$ est la matrice de $f\circ g$ et $BA$ est la matrice de $g\circ f$.
Ainsi la matrice de $f\circ g - g\circ f$ est $AB-BA$
Donc
$$\textrm{tr}(f\circ g - g\circ f)= \textrm{tr}(AB-BA) = \textrm{tr}(AB)-\textrm{tr}(BA)=0.$$
\end{enumerate}


\fincorrection
\nocorrection
\correction{002444}
Soit $A,B$ tel que $B = P^{-1} A P$.
\begin{enumerate}
  \item Supposons $A$ inversible, alors il existe $A'$ tel que $A\times A'=I$ et $A'\times A=I$.
Notons alors $B'= P^{-1} A'P$. On a
$$B \times B' = \big(P^{-1}A P \big)\times \big(P^{-1} A' P\big)=P^{-1}A \big(P P^{-1}\big)A' P 
= P^{-1}A A' P=P^{-1} I P=I$$
De même $B' \times B=I$. Donc $B$ est inversible d'inverse $B'$.

  \item Supposons que $A^n=I$. Alors 
$$\begin{array}{rcl}
B^n 
&=&\big(P^{-1} A P\big)^n= \big(P^{-1} A P \big)\big(P^{-1} A P \big)\cdots \big(P^{-1} A P \big) \\
&=&P^{-1} A (P P^{-1}) A (P P^{-1}) \cdots  A P \\
&=& P^{-1} A^n P  \\
&=& P^{-1} I P = I \\
\end{array}$$
Donc $B$ est idempotente.

  \item Si $A^n=(0)$ alors le même calcul qu'au-dessus conduit à $B^n=(0)$.

  \item Si $A = \lambda I$ alors $B=P^{-1} (\lambda I) P = \lambda I \times P^{-1}P= \lambda I$
(car la matrice $\lambda I$ commute avec toutes les matrices).
\end{enumerate}

\fincorrection
\nocorrection
\nocorrection
\nocorrection
\nocorrection
\nocorrection
\nocorrection
\nocorrection
\nocorrection
\nocorrection
\nocorrection
\correction{002774}
Avant toute, un coup d'\oe il sur la matrice nous informe de deux choses :
(a) $A$  n'est pas la matrice nulle donc $\textrm{rg}(A)\ge 1$ ;
(b) il y a $3$ lignes donc $\textrm{rg}(A)\le 3$ (le rang est plus petit que le nombre 
de colonnes et que le nombre de lignes).

\begin{enumerate}
  \item Montrons de différentes façons que $\textrm{rg}(A)\ge 2$.
  \begin{itemize}
     \item \textbf{Première méthode : sous-déterminant non nul.}
On trouve une sous-matrice $2\times 2$ dont le déterminant est non nul.
Par exemple la sous-matrice extraite du coin en bas à gauche vérifie
$\begin{vmatrix}3 & 0\\ 5 & 4\end{vmatrix}= 12 \neq 0$ donc $\textrm{rg}(A)\ge 2$.

     \item \textbf{Deuxième méthode : espace vectoriel engendré par les colonnes.}
On sait que l'image de l'application linéaire associée à la matrice $A$
est engendrée par les vecteurs colonnes. Et le rang est la dimension de cette image.
On trouve facilement deux colonnes linéairement indépendantes : 
la deuxième $\begin{pmatrix}2\\0\\4\end{pmatrix}$
et la troisième $\begin{pmatrix}-1\\1\\-1\end{pmatrix}$ colonne.
Donc $\textrm{rg}(A)\ge 2$.

     \item \textbf{Troisième méthode : espaces vectoriel engendré par les lignes.}
Il se trouve que la dimension de l'espace vectoriel engendré par les lignes
égal la dimension de l'espace vectoriel engendré par les colonnes (car $\textrm{rg}(A)=\textrm{rg}({}^tA)$).
Comme les deuxième et troisième lignes sont linéairement indépendantes alors
$\textrm{rg}(A)\ge 2$. 

Attention : les dimensions des espaces vectoriels engendrés sont égales mais les espaces sont différents !
  \end{itemize}

  \item En utilisant la dernière méthode : le rang est exactement $2$ si la première ligne est dans le sous-espace engendré
par les deux autres.
Donc 
\begin{align*}
\textrm{rg}(A) = 2
 & \iff (a,2,-1,b) \in \textrm{Vect} \big\{ (3,0,1,-4), (5,4,-1,2) \big\} \\
 & \iff \exists \lambda,\mu \in \Rr \quad (a,2,-1,b) = \lambda (3,0,1,-4) + \mu(5,4,-1,2) \\
 & \iff \exists \lambda,\mu \in \Rr \quad 
\left\{
\begin{array}{rcl}
3\lambda+5\mu &=& a \\
4\mu &=& 2 \\
\lambda-\mu &=& -1 \\
-4\lambda+2\mu &=& b \\
\end{array}
\right. 
 \iff  \left\{
\begin{array}{rcl}
\lambda &=& -\frac12 \\
\mu &=& \frac12 \\
a &=& 1 \\
b &=& 3 \\
\end{array} 
\right.\\  
\end{align*}
Conclusion la rang de $A$ est $2$ si $(a,b)=(1,3)$. Sinon le rang de $A$ est $3$.
\end{enumerate}

\fincorrection
\nocorrection
\nocorrection
\correction{003359}
2. Si $A,B \in {\cal D}$ et $AB = I$, alors pour $i \ne j$,
    $\forall\ k,\ a_{ik}b_{kj} = 0$.

    Soit $a_{i1} \ne 0$ : alors $b_{1j} = 0$ pour tout $j \ne i$, donc
    $a_{i1} = b_{1i} = 1$.

    Donc chaque colonne de $A$ contient $n-1$ fois $0$ et une fois $1$.
    $A$ est inversible $ \Rightarrow  A$ est une matrice de permutation.
 \fincorrection
\nocorrection
\correction{003393}
\begin{enumerate}
  \item 
  \item $X = \begin{pmatrix} \alpha & 1 + \beta \cr -2\alpha &1 - 2\beta \cr
                            1 + \alpha & \beta \cr\end{pmatrix}$.
\end{enumerate}
 \fincorrection
\correction{003394}
$B = \begin{pmatrix} a&2a-1&a \cr b+2&2b+3&b \cr c+2&2c+1&c \cr \end{pmatrix}$.
\fincorrection
\correction{003395}
$A^n = \begin{pmatrix} 1 & 0 & 0 \cr C_n^2 & 1 & n \cr n & 0 & 1 \cr\end{pmatrix}$.
\fincorrection
\nocorrection
\correction{003397}
\begin{enumerate}
  \item $\begin{pmatrix} a & &b    \cr
                          &\ddots \cr
                        b & &a    \cr\end{pmatrix}$ avec
            $\begin{cases} a = \frac1n\bigl( (2n-1)^k + (n-1)(-1)^k \bigr) \cr
                     b = \frac1n\bigl( (2n-1)^k -      (-1)^k \bigr).\cr \end{cases}$
  \item $\begin{pmatrix} 1 &2k &2k^2+k &\frac 13(4k^3 + 6k^2 + 2k) \cr
                       0 &1  &2k     &2k^2+k                     \cr
                       0 &0  &1      &2k                         \cr
                       0 &0  &0      &1                          \cr \end{pmatrix}$.
  \item $\left[{ \begin{pmatrix} x\cr y\cr z \cr\end{pmatrix} \begin{pmatrix} x &y &z \cr\end{pmatrix} }\right]^k            = (x^2+y^2+z^2)^{k-1}A$.
\end{enumerate}
 \fincorrection
\correction{005258}
Soient $x$ et $y$ deux réels.

\begin{align*}\ensuremath
A(x)A(y)&=\left(
\begin{array}{cc}
\ch x&\sh x\\
\sh x&\ch x
\end{array}
\right)\left(
\begin{array}{cc}
\ch y&\sh y\\
\sh y&\ch y
\end{array}
\right)=\left(
\begin{array}{cc}
\ch x\ch y+\sh x\sh y&\sh x\ch y+\ch x\sh y\\
\sh x\ch y+\ch x\sh y&\ch x\ch y+\sh x\sh y
\end{array}
\right)\\
 &=\left(
\begin{array}{cc}
\ch(x+y)&\sh(x+y)\\
\sh(x+y)&\ch(x+y)
\end{array}
\right).
\end{align*}

En particulier,

$$A(x)A(-x)=A(-x)A(x)=A(0)=\left(
\begin{array}{cc}
1&0\\
0&1
\end{array}
\right)=I_2,$$

et $A(x)$ est inversible d'inverse $A(-x)$.

On a aussi, pour $n$ entier naturel non nul donné~:

$$(A(x))^n=A(x)A(x)...A(x)=A(x+x...+x)=A(nx),$$

ce qui reste clair pour $n=0$ car $A(x)^0=I_2=A(0)$. Enfin, $(A(x))^{-n}=(A(x)^{-1})^n=A(-x)^n=A(-nx)$. Finalement, 

$$\forall n\in\Zz,\;(A(x))^n=A(nx)=\left(
\begin{array}{cc}
\ch(nx)&\sh(nx)\\
\sh(nx)&\ch(nx)
\end{array}
\right).$$
\fincorrection
\correction{005262}
Soit $f$ l'endomorphisme de $\Rr^p$ de matrice $A$ dans la base canonique $\mathcal{B}$ de $\Rr^p$. Pour $1\leq k\leq p$, on a $f(e_k)=e_{p+1-k}$ et donc $f^2(e_k)=e_k$. Ainsi, $A^2=I_p$. Mais alors, il est immédiat que, pour $n$ entier naturel donné, $A^n=I_p$ si $n$ est pair et $A^n=A$ si $n$ est impair.
\fincorrection
\correction{005263}
Pour $x\in]-1,1[$, posons $M(x)=\frac{1}{\sqrt{1-x^2}}\left(
\begin{array}{cc}
1&x\\
x&1
\end{array}
\right)$. Posons ensuite $G=\{M(x),\;x\in]-1,1[\}$.

Soit alors $x\in]-1,1[$. Posons $a=\Argth x$ de sorte que $x=\mbox{th}a$. On a 

$$M(x)=\frac{1}{\sqrt{1-x^2}}\left(
\begin{array}{cc}
1&x\\
x&1
\end{array}
\right)=\ch a\left(
\begin{array}{cc}
1&\tanh a\\
\tanh a&1
\end{array}
\right)=\left(
\begin{array}{cc}
\ch a&\sh a\\
\sh a&\ch a
\end{array}
\right).$$

Posons, pour $a\in\Rr$, $N(a)=\left(
\begin{array}{cc}
\ch a&\sh a\\
\sh a&\ch a
\end{array}
\right)$. On a ainsi $\forall x\in]-1,1[,\;M(x)=N(\Argth x)$ ou aussi, $\forall a\in\Rr,\;N(a)=M(\mbox{th}a)$. Par suite, $G=\{N(a),\;a\in\Rr\}$.

Soit alors $(a,b)\in\Rr^2$.

\begin{align*}\ensuremath
N(a)N(b)&=\left(
\begin{array}{cc}
\ch a&\sh a\\
\sh a&\ch a
\end{array}
\right)\left(
\begin{array}{cc}
\ch b&\sh b\\
\sh b&\ch b
\end{array}
\right)=\left(
\begin{array}{cc}
\ch a\ch b+\sh a\sh b&\sh a\ch b+\sh b\ch a\\
\sh a\ch b+\sh b\ch a&\ch a\ch b+\sh a\sh b
\end{array}
\right)
\\
 &=\left(
\begin{array}{cc}
\ch(a+b)&\sh(a+b)\\
\sh(a+b)&\ch(a+b)
\end{array}
\right)
=N(a+b).
\end{align*}

Montrons alors que $G$ est un sous-groupe de $(\mathcal{GL}_2(\Rr),\times)$.

$N(0)=I_2\in G$ et donc $G$ est non vide.

$\forall a\in\Rr,\;\mbox{det}(N(a))=\ch^2a-\sh^2a=1\neq0$ et donc $G\subset\mathcal{GL}_2(\Rr)$.

$\forall(a,b)\in\Rr^2,\;N(a)N(b)=N(a+b)\in G$.

$\forall a\in\Rr,\;(N(a))^{-1}=\left(
\begin{array}{cc}
\ch a&-\sh a\\
-\sh a&\ch a
\end{array}
\right)=N(-a)\in G$.

On a montré que $G$ est un sous-groupe de $(\mathcal{GL}_2(\Rr),\times)$.
\fincorrection
\correction{005611}
Par hypothèse, $a_{i,j}= 0$ pour $j\leqslant i + r - 1$ et $bi,j = 0$ pour $j\leqslant i + s - 1$.

Soient $i$ et $j$ deux indices tels que $j\leqslant i + r + s - 1$. Le coefficient ligne $i$, colonne $j$, de $AB$ vaut $\sum_{k=1}^{n}a_{i,k}b_{k,j}$.

 
Dans cette somme, si $k\leqslant i + r -1$, $a_{i,k}= 0$. Sinon $k\geqslant i + r$ et donc $j\leqslant i + r + s - 1\leqslant k + s - 1$ et dans ce cas $b_{k,j}= 0$.

Finalement, le coefficient ligne $i$, colonne $j$, de $AB$ est bien nul si $j\leqslant i + r+s-1$.
\fincorrection
\nocorrection
\nocorrection
\correction{005257}
\begin{enumerate}
\item  Soit $X=\left(
\begin{array}{c}
2\\
-3\\
5
\end{array}
\right)$. $MX=\left(
\begin{array}{ccc}
2&1&0\\
-3&-1&1\\
1&0&-1
\end{array}
\right)\left(
\begin{array}{c}
2\\
-3\\
5
\end{array}
\right)=\left(
\begin{array}{c}
1\\
2\\
-3
\end{array}
\right)$  et $u(2i-3j+5k)=i+2j-3k$.
\item  Soit $X=\left(
\begin{array}{c}
x\\
y\\
z
\end{array}
\right)\in\mathcal{M}_{3,1}(\Rr)$ . 

$$MX=0\Leftrightarrow\left(
\begin{array}{ccc}
2&1&0\\
-3&-1&1\\
1&0&-1
\end{array}
\right)\left(
\begin{array}{c}
x\\
y\\
z
\end{array}
\right)=\left(
\begin{array}{c}
0\\
0\\
0
\end{array}
\right)\Leftrightarrow
\left\{
\begin{array}{l}
2x+y=0\\
-3x-y+z=0\\
x-z=0
\end{array}
\right.\Leftrightarrow\left\{
\begin{array}{l}
y=-2x\\
z=x
\end{array}
\right..$$

Donc, $\mbox{Ker}u=\mbox{Vect}(i-2j+k)$. En particulier, $\mbox{dim}(\mbox{Ker}u)=1$ et, d'après le théorème du rang, $\mbox{rg}u=2$. Or, $u(j)=i-j$ et $u(k)=j+k$ sont deux vecteurs non colinéaires de $\mbox{Im}u$ qui est un plan vectoriel et donc $\mbox{Im}u=\mbox{Vect}(i-j,j-k)$.

\item  $$M^2=\left(
\begin{array}{ccc}
2&1&0\\
-3&-1&1\\
1&0&-1
\end{array}
\right)\left(
\begin{array}{ccc}
2&1&0\\
-3&-1&1\\
1&0&-1
\end{array}
\right)=\left(
\begin{array}{ccc}
1&1&1\\
-2&-2&-2\\
1&1&1
\end{array}
\right)$$

et 

$$M^3=M^2.M=\left(
\begin{array}{ccc}
1&1&1\\
-2&-2&-2\\
1&1&1
\end{array}
\right)\left(
\begin{array}{ccc}
2&1&0\\
-3&-1&1\\
1&0&-1
\end{array}
\right)=0.$$

\item  $\mbox{Ker}u^2$ est à l'évidence le plan d'équation $x+y+z=0$. Une base de $\mbox{Ker}u^2$ est $(i-j,j-k)$ et donc $\mbox{Ker}u^2=\mbox{Im}u=\mbox{Vect}(i-j,j-k)$.

D'après le théorème du rang, $\mbox{Im}u^2$ est une droite vectorielle. Mais $u^3=0$ s'écrit encore $u\circ u^2=0$, et donc $\mbox{Im}u^2$ est contenue dans $\mbox{Ker}u$ qui est une droite vectorielle. Donc, $\mbox{Im}u^2=\mbox{Ker}u=\mbox{Vect}(i-2j+k)$.

\item  $(I-M)(I+M+M^2)=I-M^3=I$. Par suite, $I-M$ est inversible à droite et donc inversible et 

$$(I-M)^{-1}=I+M+M^2=\left(
\begin{array}{ccc}
1&0&0\\
0&1&0\\
0&0&1
\end{array}
\right)+\left(
\begin{array}{ccc}
2&1&0\\
-3&-1&1\\
1&0&-1
\end{array}
\right)+\left(
\begin{array}{ccc}
1&1&1\\
-2&-2&-2\\
1&1&1
\end{array}
\right)=\left(
\begin{array}{ccc}
4&2&1\\
-5&-2&-1\\
2&1&1
\end{array}
\right).$$

\end{enumerate}
\fincorrection
\correction{005260}
\begin{enumerate}
\item  Pour $P$ élément de $\Rr_n[X]$, 

$$f(P)=e^{X^2}(Pe^{-X^2})'=e^{X^2}(P'e^{-X^2}-2XPe^{-X^2})=P'-2XP.$$

Ainsi, si $P$ est un polynôme de degré infèrieur ou égal à $n$, $f(P)=P'-2XP$ est un polynôme de degré inférieur ou égal à $n+1$, et $f$ est bien une application de $\Rr_n[X]$ dans $\Rr_{n+1}[X]$.

De plus, pour $(\lambda,\mu)\in\Rr^2$ et $(P,Q)\in\Rr_n[X]$, on a~:

$$f(\lambda P+\mu Q)=(\lambda P+\mu Q)'-2X(\lambda P+\mu Q)=\lambda(P'-2XP)+\mu(Q'-2XQ)=\lambda f(P)+\mu f(Q).$$

$f$ est élément de $\mathcal{L}(\Rr_n[X],\Rr_{n+1}[X])$.

\item  La matrice $A$ cherchée est élément de $\mathcal{M}_{n+1,n}(\Rr)$.

Pour $k=0$, $f(X^k)=f(1)=-2X$ et pour $1\leq k\leq n$, $f(X^k)=kX^{k-1}-2X^{k+1}$. On a donc~:

$$A=
\left(
\begin{array}{cccccc}
0&1&0&\ldots&\ldots&0\\
-2&0&2&0& &\vdots\\
0&-2&0&\ddots&\ddots&\vdots\\
\vdots&\ddots&\ddots&\ddots&\ddots&0\\
 & & &\ddots&\ddots&n\\
\vdots& & &\ddots&-2&0\\
0&\ldots& &\ldots&0&-2
\end{array}
\right).$$

\item  Soit $P\in\Rr_n[X]$ tel que $f(P)=0$.

Si $P$ n'est pas nul, $-2XP$ a un degré strictement plus grand que $P'$ et donc $f(P)$ n'est pas nul. Par suite, $\mbox{Ker}f=\{0\}$ ($f$ est donc injective) et d'après le théorème du rang, $\mbox{rg}f=\mbox{dim}(\Rr_n[X])-0=n+1$, ce qui montre que $\mbox{Im}f$ n'est pas $\Rr_{n+1}[X]$ ($f$ n'est pas surjective).
\end{enumerate}
\fincorrection
\correction{005269}
\begin{enumerate}
\item 
 
\begin{align*}\ensuremath
\mbox{rg}\left(
\begin{array}{ccc}
1&1/2&1/3\\
1/2&1/3&1/4\\
1/3&1/4&m
\end{array}
\right)&=\mbox{rg}\left(
\begin{array}{ccc}
1&0&0\\
1/2&1/12&1/12\\
1/3&1/12&m-\frac{1}{9}
\end{array}
\right)\quad(\mbox{rg}(C_1,C_2,C_3)=\mbox{rg}(C_1,C_2-\frac{1}{2}C_1,C_3-\frac{1}{3}C_1))
\\
 &=\mbox{rg}\left(
\begin{array}{ccc}
1&0&0\\
1/2&1/12&0\\
1/3&1/12&m-\frac{7}{36}
\end{array}
\right)\quad(\mbox{rg}(C_1,C_2,C_3)=\mbox{rg}(C_1,C_2,C_3-C_2))
\end{align*}

Si $m=\frac{7}{36}$, $\mbox{rg}A=2$ (on note alors que $C_1=6(C_2-C_3)$) et si $m\neq\frac{7}{36}$, $\mbox{rg}A=3$ et $A$ est inversible.

\item  
\begin{align*}\ensuremath
\mbox{rg}\left(
\begin{array}{ccc}
1&1&1\\
b+c&c+a&a+b\\
bc&ca&ab
\end{array}
\right)&=\mbox{rg}\left(
\begin{array}{ccc}
1&0&0\\
b+c&a-b&a-c\\
bc&c(a-b)&b(a-c)
\end{array}
\right)\quad(\mbox{rg}(C_1,C_2,C_3)=\mbox{rg}(C_1,C_2-C_1,C_3-C_1))
\end{align*}

\begin{itemize}
\item[1er cas.] si $a$, $b$ et $c$ sont deux à deux distincts. 

$$\mbox{rg}\left(
\begin{array}{ccc}
1&0&0\\
b+c&1&1\\
bc&c&b
\end{array}
\right)=\mbox{rg}\left(
\begin{array}{ccc}
1&0&0\\
b+c&1&0\\
bc&c&b-c
\end{array}
\right)=\mbox{rg}\left(
\begin{array}{ccc}
1&0&0\\
b+c&1&0\\
bc&c&1
\end{array}
\right).$$

Donc, si $a$, $b$ et $c$ sont deux à deux distincts alors $\mbox{rg}A=3$.

\item[2ème cas.] Si $b=c\neq a$ (ou $a=c\neq b$ ou $a=b\neq c$).
$A$ a même rang que $\left(
\begin{array}{ccc}
1&0&0\\
b+c&1&1\\
bc&c&b
\end{array}
\right)$ puis que $\left(
\begin{array}{ccc}
1&0&0\\
b+c&1&0\\
bc&c&0
\end{array}
\right)$. Donc, si $b=c\neq a$ ou $a=c\neq b$ ou $a=b\neq c$, $\mbox{rg}A=2$.

\item[3ème cas.] Si $a=b=c$, il est clair dès le départ que $A$ est de rang $1$.
\end{itemize}

\item  Puisque $\mbox{rg}(C_1,C_2,C_3,C_4)=\mbox{rg}(C_1,C_2-aC_1,C_3-C_1,C_4-bC_1)$,

\begin{align*}\ensuremath
\mbox{rg}\left(
\begin{array}{cccc}
1&a&1&b\\
a&1&b&1\\
1&b&1&a\\
b&1&a&1
\end{array}
\right)&=\mbox{rg}\left(
\begin{array}{cccc}
1&0&0&0\\
a&1-a^2&b-a&1-ab\\
1&b-a&0&a-b\\
b&1-ab&a-b&1-b^2
\end{array}
\right)=1+\mbox{rg}\left(
\begin{array}{ccc}
1-a^2&b-a&1-ab\\
b-a&0&a-b\\
1-ab&a-b&1-b^2
\end{array}
\right)
\end{align*}

\begin{itemize}
\item[1er cas.] Si $a\neq b$.

\begin{align*}\ensuremath
\mbox{rg}A&=1+\mbox{rg}\left(
\begin{array}{ccc}
1-a^2&b-a&1-ab\\
b-a&0&a-b\\
1-ab&a-b&1-b^2
\end{array}
\right)=1+\mbox{rg}\left(
\begin{array}{ccc}
1-a^2&1&1-ab\\
1&0&-1\\
1-ab&-1&1-b^2
\end{array}
\right)
\\
 &=1+\mbox{rg}\left(
\begin{array}{ccc}
1&0&-1\\
1-a^2&1&1-ab\\
1-ab&-1&1-b^2
\end{array}
\right)\quad(\mbox{rg}(L_1,L_2,L_3)=\mbox{rg}(L_2,L_1,L_3)).
\\
 &=1+\mbox{rg}\left(
\begin{array}{ccc}
1&0&0\\
1-a^2&1&2-a^2-ab\\
1-ab&-1&2-b^2-ab
\end{array}
\right)=1+\mbox{rg}\left(
\begin{array}{ccc}
1&0&0\\
1-a^2&1&0\\
1-ab&-1&(2-b^2-ab)-(2-a^2-ab)
\end{array}
\right)
\\
 &=1+\mbox{rg}\left(
\begin{array}{ccc}
1&0&0\\
1-a^2&1&0\\
1-ab&-1&(a-b)(a+b)
\end{array}
\right)
\end{align*}
Si $|a|\neq|b|$, $\mbox{rg}A=4$ et si $a=-b\neq0$, $\mbox{rg}A=3$.

\item[2ème cas.] Si $a=b$.

\begin{align*}\ensuremath
\mbox{rg}A&=1+\mbox{rg}\left(
\begin{array}{ccc}
1-a^2&0&1-a^2\\
0&0&0\\
1-a^2&0&1-a^2
\end{array}
\right)=1+\mbox{rg}\left(
\begin{array}{cc}
1-a^2&1-a^2\\
0&0\\
1-a^2&1-a^2
\end{array}
\right)=1+\mbox{rg}\left(
\begin{array}{cc}
1-a^2&1-a^2\\
1-a^2&1-a^2
\end{array}
\right)
\end{align*}

Si $a=b=\pm1$, $\mbox{rg}A=1$ et si $a=b\neq\pm1$, $\mbox{rg}A=2$.
\end{itemize}

\item  Pour $n\geq2$ et $j\in\{1,...,n\}$, notons $C_j$ la $j$-ème colonne de la matrice proposée.

$$C_j=(i+j+ij)_{1\leq i\leq n}=(i)_{1\leq i\leq n}+j(i+1)_{1\leq i\leq n}=jU+V,$$

avec $U=\left(\begin{array}{c}
2\\
3\\
\vdots\\
i+1\\
\vdots\\
n+1
\end{array}
\right)$ 
et $V=\left(\begin{array}{c}
1\\
2\\
\vdots\\
i\\
\vdots\\
n
\end{array}
\right)$.

Ainsi, $\forall j\in\{1,...,n\},\;C_j\in\mbox{Vect}(U,V)$ ce qui montre que $\mbox{rg}A\leq2$.
De plus, la matrice extraite $\left(
\begin{array}{cc}
3&5\\
5&8
\end{array}
\right)$ (lignes et colonnes 1 et 2) est inversible et finalement $\mbox{rg}A=2$.

\item  On suppose $n\geq2$. La $j$-ème colonne de la matrice s'écrit

$$C_j=(\sin i\cos j+\sin j\cos i)_{1\leq i\leq n}=\sin jC+\cos jS\;\mbox{avec}\;C=(\cos i)_{1\leq i\leq n}\;\mbox{et}\;S=(\sin i)_{1\leq i\leq n}.$$

Par suite, $\forall j\in\{1,...,n\}$, $C_j\in\mbox{Vect}(C,S)$ ce qui montre que $\mbox{rg}A\leq2$. De plus, la matrice extraite formée des termes lignes et colonnes 1 et 2 est inversible car son déterminant vaut
$\sin2\sin4-\sin^23=-0,7...\neq0$ et finalement $\mbox{rg}A=2$.

\item  Déterminons $\mbox{Ker}A$. Soit $(x_i)_{1\leq i\leq n}\in\mathcal{M}_{n,1}(\Cc)$.

$$(x_i)_{1\leq i\leq n}\in\mbox{Ker}A\Leftrightarrow\forall i\in\{1,...,n-1\},\;ax_i+bx_{i+1}=0\;\mbox{et}bx_1+ax_n=0\;(S).$$

\begin{itemize}
\item[1er cas.] Si $a=b=0$, alors clairement $\mbox{rg}A=0$.

\item[2ème cas.] Si $a=0$ et $b\neq0$, alors $(S)\Leftrightarrow\forall i\in\{1,...,n\}\;x_i=0$. Dans ce cas, $\mbox{Ker}A=\{0\}$ et donc $\mbox{rg}A=n$.

\item[3ème cas.] Si $a\neq0$. Posons $\alpha=-\frac{b}{a}$.

\begin{align*}\ensuremath
(S)&\Leftrightarrow\forall k\in\{1,...,n-1\},\;x_k=\alpha x_{k+1}\;\mbox{et}\;x_n=\alpha x_1\\
 &\Leftrightarrow\forall k\in\{1,...,n\},\;x_k=\alpha^{-(k-1)}x_1\;\mbox{et}\;x_n=\alpha x_1\\
 &\Leftrightarrow\forall k\in\{1,...,n\},\;x_k=\alpha^{k-1}x_1\;\mbox{et}\;\alpha^n x_1=x_1
\end{align*}

Mais alors, si $\alpha^n\neq1$, le système $(S)$ admet l'unique solution $(0,...,0)$ et $\mbox{rg}A=n$, et si $\alpha^n=1$, $\mbox{Ker}A=\mbox{Vect}((1,\alpha^{n-1},...,\alpha^2,\alpha))$ est de dimension $1$ et $\mbox{rg}A=n-1$.
\end{itemize}

En résumé, si $a=b=0$, $\mbox{rg}A=0$ et si $a=0$ et $b\neq0$, $\mbox{rg}A=n$. Si $a\neq0$ et $-\frac{b}{a}\in U_n$, $\mbox{rg}A=n-1$ et si $a\neq0$ et $-\frac{b}{a}\notin U_n$, $\mbox{rg}A=n$.
\end{enumerate}

\fincorrection
\correction{005603}
(C'est en fait un exercice sur les polynômes de \textsc{Tchebychev} de 1ère espèce et vous pouvez généraliser cet exercice en passant au format $n$ au lieu du format $4$.)

Si on note $C_j$, $j\in\{1,2,3,4\}$, la $j$-ème colonne de $A$ alors $C_j= (\cos(i+j-2)a)_{1\leqslant i\leqslant4}$ puis pour $j$ élément de $\{1,2\}$, 

\begin{center}
$C_{j+2}+ C_j =(2\cos(i+j-1)a\cos a)_{1\leqslant i\leqslant4}= 2\cos aC_{j+1}$
\end{center}

et donc $C_3 =2\cos aC_2 - C_1\in\text{Vect}(C_1,C_2)$ et $C_4= 2\cos a C_3 -C_2\in\text{Vect}(C_2,C_3)\subset\text{Vect}(C_1,C_2)$.

Donc $\text{Vect}(C_1,C_2,C_3,C_4)=\text{Vect}(C_1,C_2)$ et $\text{rg}A=\text{rg}(C_1,C_2)\leqslant 2$.

Enfin $\left|
\begin{array}{cc}
1&\cos(a)\\
\cos(a)&\cos(2a)
\end{array}
\right|=\cos(2a) -\cos^2a=\cos^2a - 1= -\sin^2a$.

\textbullet~Si $a$ n'est pas dans $\pi\Zz$, ce déterminant n'est pas nul et donc les deux premières colonnes ne sont pas colinéaires. Dans ce cas, $\text{rg}A = 2$. 

\textbullet~Si $a$ est dans $\pi\Zz$, la première colonne n'est pas nulle et les autres colonnes lui sont colinéaires. Dans ce cas, $\text{rg}A = 1$.

\begin{center}
\shadowbox{
$\text{rg}(A)=2$ si $a\notin\pi\Zz$ et $\text{rg}(A)=1$ si $a\in\pi\Zz$.
}
\end{center}
\fincorrection
\correction{005607}
Pour $j\in\llbracket1,n\rrbracket$, notons $C_j$ la $j$-ème colonne de la matrice $A$. Posons encore $U=\left(
\begin{array}{c}
1\\
2\\
\vdots\\
n
\end{array}
\right)$  et $V=\left(
\begin{array}{c}
2\\
3\\
\vdots\\
n+1
\end{array}
\right)$. Pour $j\in\llbracket1,n\rrbracket$, on a

\begin{center} 
$Cj = (i+j(i+1))_{1\leqslant i\leqslant n}= (i)_{1\leqslant i\leqslant n}+j(i+1)_{1\leqslant i\leqslant n}=U+jV$.
\end{center}

Donc $\text{Vect}(C_1,...,C_n)\subset\text{Vect}(U,V)$ et en particulier, $\text{rg}A\leqslant 2$.
Maintenant, si $n\geqslant 2$, les deux premières colonnes de $A$ ne sont pas colinéaires car $\left|
\begin{array}{cc}
3&5\\
5&8
\end{array}
\right|=-1\neq0$. Donc, si $n\geqslant2$, $\text{rg}A=2$ et si $n=1$, $\text{rg}A =1$.

\begin{center}
\shadowbox{
Si $n\geqslant 2$, $\text{rg}(i+j+ij)_{1\leqslant i,j\leqslant n}=2$ et si $n=1$, $\text{rg}(i+j+ij)_{1\leqslant i,j\leqslant n}=1$.
}
\end{center}
\fincorrection
\correction{005622}
On note $r$ le rang de $A$. Si $r=0$, $A$ est nulle et donc $B$ est nulle.

Sinon, il existe deux matrices carrées inversibles $P$ et $Q$ de format $n$ telles que $A=PJ_rQ$ où $J_r=\left(
\begin{array}{cc}
I_r&0\\
0&0
\end{array}
\right)$. Soient $P'=\left(
\begin{array}{cccc}
P&0&\ldots&0\\
0&\ddots&\ddots&\vdots\\
\vdots&\ddots&\ddots&0\\
0&\ldots&0&P
\end{array}
\right)\in\mathcal{M}_{np}(\Cc)$ et $Q'=\left(
\begin{array}{cccc}
P&0&\ldots&0\\
0&\ddots&\ddots&\vdots\\
\vdots&\ddots&\ddots&0\\
0&\ldots&0&Q
\end{array}
\right)\in\mathcal{M}_{np}(\Cc)$. Puisque $\text{det}(P')=(\text{det}(P))^p\neq0$ et $\text{det}(Q')=(\text{det}(Q))^p\neq0$, les matrices $P'$ et $Q'$ sont inversibles. De plus, un calcul par blocs montre que

\begin{center}
$B=\left(
\begin{array}{cccc}
PJ_rQ&0&\ldots&0\\
0&\ddots&\ddots&\vdots\\
\vdots&\ddots&\ddots&0\\
0&\ldots&0&PJ_rQ
\end{array}
\right)=P'J_r'Q'$ où $J_r'=\left(
\begin{array}{cccc}
J_r&0&\ldots&0\\
0&\ddots&\ddots&\vdots\\
\vdots&\ddots&\ddots&0\\
0&\ldots&0&J_r
\end{array}
\right)$.
\end{center}

La matrice $B$ est équivalente a la matrice $J_r'$ et a donc même rang que $J_r'$. Enfin, en supprimant les lignes nulles et les colonnes nulles, on voit que la matrice $J_r'$ a même rang que la matrice $I_{pr}$ à savoir $pr$. Dans tous les cas, on a montré que  

\begin{center}
\shadowbox{
$\text{rg}B =p\text{rg}A$.
}
\end{center}
\fincorrection
\correction{005623}
Soit $r$ le rang de $H$. Il existe deux matrices carrées inversibles $P$ et $Q$ de format $n$ telles que $H =PJ_rQ$ où $J_r=\left(
\begin{array}{cc}
I_r&0\\
0&0
\end{array}
\right)$. L'égalité $HAH =\lambda_AH$ s'écrit après simplifications $J_rQAPJ_r=\lambda_AJ_r$.
Maintenant , quand $A$ décrit $\mathcal{M}_n(\Kk)$, la matrice $B=QAP$ décrit également $\mathcal{M}_n(\Kk)$ (par exemple, l'application qui à $A$ associe $QAP$ est une permutation de $\mathcal{M}_n(\Kk)$ de réciproque l'application qui à $A$ associe $Q^{-1}AP^{-1}$).

L'énoncé s'écrit maintenant de manière plus simple : montrons que $(\forall B\in\mathcal{M}_n(\Kk),\;\exists\lambda_B\in\Kk/\;J_rBJ_r =\lambda_BJ_r)\Rightarrow r\leqslant 1$.

Un calcul par blocs fournit en posant $B=\left(
\begin{array}{cc}
B_1&B_3\\
B_2&B_4
\end{array}
\right)$

\begin{center}
$J_rBJ_r=\left(
\begin{array}{cc}
I_r&0\\
0&0
\end{array}
\right)\left(
\begin{array}{cc}
B_1&B_3\\
B_2&B_4
\end{array}
\right)\left(
\begin{array}{cc}
I_r&0\\
0&0
\end{array}
\right)=\left(
\begin{array}{cc}
B_1&0\\
B_2&0
\end{array}
\right)\left(
\begin{array}{cc}
I_r&0\\
0&0
\end{array}
\right)=\left(
\begin{array}{cc}
B_1&0\\
0&0
\end{array}
\right)$
\end{center} 

Mais si $r\geqslant 2$, il existe des matrices carrées $B_1$ de format $r$ qui ne sont pas des matrices scalaires et donc telles que $B_1$ n'est pas colinéaire à $I_r$. Donc $r\leqslant 1$.
\fincorrection
\correction{005624}
(1) $\Rightarrow$ (2).

$M^2 = 0\Rightarrow\text{Im}M\subset\text{Ker}M\Rightarrow\text{rg}M\leqslant\text{dim}(\text{Ker}M)=3 -\text{rg}M$ et donc $\text{rg}M\leqslant1$.

Si $\text{rg}M=0$ alors $\text{Tr}M = 0$. On suppose maintenant que $\text{rg}M=1$ et donc $\text{dim}(\text{Ker}M)= 2$.

Soit $e_1$ un vecteur non nul de $\text{Im}M$ alors il existe un vecteur $e_3$ (non nul) tel que $Me_3 = e_1$.

On complète la famille libre $(e_1)$ de $\text{Im}M\subset\text{Ker}M$ en $(e_1,e_2)$ base de $\text{Ker}M$. La famille $(e_1,e_2,e_3)$ est une base de $\mathcal{M}_{3,1}(\Rr)$ car

\begin{center}
$ae_1+be_2+ce_3 = 0\Rightarrow M(ae_1+be_2+ce_3) = 0\Rightarrow ce_1 = 0\Rightarrow c = 0$,
\end{center}

puis $a = b = 0$ car la famille $(e_1,e_2)$ est libre.

$M$ est donc semblable à la matrice $\left(
\begin{array}{ccc}
0&0&1\\
0&0&0\\
0&0&0
\end{array}
\right)$ et en particulier $\text{Tr}M = 0$.

(2) $\Rightarrow$ (1).

 
Si $\text{rg}M = 0$, $M^2 = 0$.

Si $\text{rg}M = 1$, on peut se rappeler de l'écriture générale d'une matrice de rang 1 : il existe trois réels $u_1$, $u_2$ et $u_3$ non tous nuls et trois réels $v_1$, $v_2$ et $v_3$ non tous nuls tels que $M=\left(
\begin{array}{ccc}
u_1v_1&u_1v_2&u_1v_3\\
u_2v_1&u_2v_2&u_2v_3\\
u_3v_1&u_3v_2&u_3v_3
\end{array}
\right)$ ou encore il existe deux vecteurs colonnes, tous deux non nuls $U$ et $V$ tels que $M=U{^t}V$. L'égalité $\text{Tr}M =0$ fournit $u_1v_1+u_2v_2+u_3v_3= 0$ ou encore ${^t}UV=0$. Mais alors

\begin{center}
$M^2=U{^t}VU{^t}V=U{^t}({^t}UV){^t}V=0$
\end{center}

Cet exercice admet des solutions bien plus brèves avec des connaissances sur la réduction .
\fincorrection
\nocorrection
\nocorrection
\nocorrection
\nocorrection
\nocorrection
\nocorrection
\correction{001087}
L'expression de $f$ dans la base $\mathcal{B}$ est la suivante $f(x,y)=(x-y,0)$.
Autrement dit à un vecteur $\begin{pmatrix}x\\y\end{pmatrix}$ on associe
le vecteur  $\begin{pmatrix}x-y\\0\end{pmatrix}$.
On note que $f$ est bien une application linéaire.
Cette expression nous permet de calculer les matrices demandées.


Remarque : comme $\mathcal{B}$ est la base canonique on note
$\begin{pmatrix}x\\y\end{pmatrix}$ pour $\begin{pmatrix}x\\y\end{pmatrix}_\mathcal{B}$
qui est le vecteur $x \vec{i}+y\vec{j}$.



\begin{enumerate}
  \item Calcul de $\textrm{Mat}(f,\mathcal{B},\mathcal{B})$.
Comme $\mathcal{B}=(\vec{i}, \vec{j})$, la matrice s'obtient en calculant $f(\vec{i})$ et $f(\vec{j})$ :
$$f(\vec{i})=f\begin{pmatrix}1\\0\end{pmatrix} = \begin{pmatrix}1\\0\end{pmatrix} = \vec{i}
\quad 
f(\vec{j})=f\begin{pmatrix}0\\1\end{pmatrix} = \begin{pmatrix}-1\\0\end{pmatrix} = -\vec{i}$$
donc
$$\textrm{Mat}(f,\mathcal{B},\mathcal{B}) = \begin{pmatrix} 1 & -1 \\ 0 & 0 \end{pmatrix}$$

  \item On garde la même application linéaire mais la base de départ change (la base d'arrivée reste $\mathcal{B}$).
Si on note $\vec{u} = \vec{i}-\vec{j}$ et $\vec{v} = -2\vec{i}+3\vec{j}$, on a 
$\mathcal{B'}=(\vec{i} - \vec{j}, -2\vec{i}+3\vec{j}) = (\vec{u},\vec{v})$. On exprime 
$f(\vec{u})$ et $f(\vec{v})$ dans la base d'arrivée $\mathcal{B}$.
$$f(\vec{u})=f(\vec{i}- \vec{j})=f\begin{pmatrix}1\\-1\end{pmatrix} = \begin{pmatrix}2\\0\end{pmatrix}
\quad 
f(\vec{v})=f(-2\vec{i}+3\vec{j})=f\begin{pmatrix}-2\\3\end{pmatrix} = \begin{pmatrix}-5\\0\end{pmatrix}$$
donc
$$\textrm{Mat}(f,\mathcal{B}',\mathcal{B}) = \begin{pmatrix} 2 & -5 \\ 0 & 0 \end{pmatrix}$$
  
  \item Toujours avec le même $f$ on prend $\mathcal{B}'$ comme base de départ et d'arrivée,
il s'agit donc d'exprimer $f(\vec{u})$ et $f(\vec{v})$ dans la base $\mathcal{B}'=(\vec{u},\vec{v})$.
Nous venons de calculer que 
$$f(\vec{u})=f(\vec{i}- \vec{j})=f\begin{pmatrix}1\\-1\end{pmatrix} = \begin{pmatrix}2\\0\end{pmatrix}=2\vec{i}
\quad 
f(\vec{v})=f(2\vec{i}+3\vec{j})=f\begin{pmatrix}-2\\3\end{pmatrix} = \begin{pmatrix}-5\\0\end{pmatrix}=-5\vec{i}$$
Mais il  nous faut obtenir une expression en fonction de la base $\mathcal{B}'$.
Remarquons que 
$$\left\{\begin{array}{lcr}
\vec{u} &=& \vec{i}-\vec{j} \\
\vec{v} &=& -2\vec{i}+3\vec{j} \\           
         \end{array}\right.
\implies
\left\{\begin{array}{lcr}
\vec{i} &=& 3\vec{u}+\vec{v} \\
\vec{j} &=& 2\vec{u}+\vec{v} \\           
         \end{array}\right.$$
Donc 
$$f(\vec{u})=f(\vec{i}- \vec{j})=2\vec{i}=6\vec{u}+2\vec{v} = \begin{pmatrix}6\\2\end{pmatrix}_{\mathcal{B}'}
\quad
f(\vec{v})=f(-2\vec{i}+3\vec{j})=-5\vec{i}=-15\vec{u}-5\vec{v} = \begin{pmatrix}-15\\-5\end{pmatrix}_{\mathcal{B}'}$$
Donc
$$\textrm{Mat}(f,\mathcal{B}',\mathcal{B}') = \begin{pmatrix} 6 & -15 \\ 2 & -5 \end{pmatrix}$$

Remarque :
$\begin{pmatrix}x\\y\end{pmatrix}_{\mathcal{B}'}$
désigne le vecteur $x \vec{u}+y\vec{v}$.
\end{enumerate}

\fincorrection
\nocorrection
\nocorrection
\nocorrection
\nocorrection
\nocorrection
\correction{001093}
\begin{enumerate}
  \item Nous devons montrer $\Ker f \cap \Im f = \{0\}$ et 
$\Ker f + \Im f = E$.
  \begin{enumerate}
     \item Si $x \in \Ker f \cap \Im f$ alors d'une part $f(x)=0$ et d'autre part il existe $x'\in E$ tel que 
$x=f(x')$. Donc $0=f(x)=f\big(f(x')\big)= f(x')=x$ donc $x=0$ (on a utilisé $f\circ f=f$). 
Donc $\Ker f \cap \Im f = \{0\}$.

     \item Pour $x\in E$ on le réécrit $x=x-f(x) + f(x)$. Alors $x-f(x) \in \Ker f$ 
(car $f\big(x-f(x) \big)=f(x)-f\circ f (x)=0$)
et $f(x)\in \Im f$. Donc $x \in \Ker f + \Im f$. Donc $\Ker f + \Im f = E$.

     \item Conclusion  : $E= \Ker f \oplus \Im f$.
  \end{enumerate}

  \item Notons $r$ le rang de $f$ : $r=\dim \Im f$.
Soit $\{e_1,\ldots,e_r\}$ une base de $\Im f$ et 
soit $\{e_{r+1},\ldots, e_n\}$ une base de $\Ker f$.
Comme 
$E= \Ker f \oplus \Im f$ alors
$(e_1,\ldots ,e_n)$ est une base de $E$.
Pour $i > r$ alors $e_i \in \Ker f$ donc $f(e_i)=0$.

Comme $f\circ f=f$ alors pour n'importe quel $x\in \Im f$ on a $f(x)=x$ :
en effet comme $x\in \Im f$, il existe $x'\in E$ tel que $x=f(x')$ 
ainsi $f(x)=f\big(f(x')\big)=f(x')=x$. En particulier si $i\le r$ alors
$f(e_i)=e_i$.

  \item La matrice de $f$ dans la base $(e_1,\ldots,e_n)$ est donc :
$$\begin{pmatrix}
  I & (0) \\
  (0) & (0) \\  
  \end{pmatrix}$$
où $I$ désigne la matrice identité de taille $r\times r$
et les $(0)$ désignent des matrices nulles.

\end{enumerate}

\fincorrection
\correction{001094}
\begin{enumerate}
  \item Il est facile de voir que $f(\lambda P + \mu Q) = \lambda f(P)+\mu f(Q)$ donc $f$ est linéaire,
de plus, $P$ étant un polynôme de degré $\le n$ alors $f(P)$ aussi.

  \item Pour $n=3$ on calcule l'image de chacun des éléments de la base :
$$f(1)=1+1-2=0,\quad f(X)=(X+1)+(X-1)-2X=0,$$
$$f(X^2)=(X+1)^2+(X-1)^2-2X^2=2,
\quad f(X^3)=(X+1)^3+(X-1)^3-2X^3=6X.$$
Donc la matrice de $f$ dans la base $(1, X, X^2, X^3)$ est
$$\begin{pmatrix}
0 & 0 & 2 & 0 \\
0 & 0 & 0 & 6 \\
0 & 0 & 0 & 0 \\
0 & 0 & 0 & 0 \\     
  \end{pmatrix}$$

Pour le cas général on calcule 
\begin{align*}
f(X^p)
 &=(X+1)^p+(X-1)^p-2X^p \\
 &= \sum_{k=1}^p \binom{p}{k}X^k + \sum_{k=1}^p \binom{p}{k}X^k(-1)^{p-k} -2X^p\\
 &= \sum_{p-k \text{ pair et } k<p} 2\binom{p}{k}X^k  
\end{align*}

Donc la matrice est 
$$\begin{pmatrix} 
0 & 0 & 2\binom{2}{0} & 0             & \cdots & 2\binom{p}{0} & 0               &        \\
  & 0 & 0             & 2\binom{3}{1} &        & 0             & 2\binom{p+1}{1} &        \\
  &   & 0             & 0             & \cdots & 2\binom{p}{2} & 0               &        \\
  &   &               & 0             &        & 0             & 2\binom{p+1}{3} & \vdots \\  
  &   &               &               & \ddots & \vdots        & 0               &        \\
  &   &               &               &        & 0             & \vdots          &        \\
  &   &               &               &        &               & 0               &        \\
  &   &               &               &        &               &                 & 0      \\
  \end{pmatrix}$$
Dans cet exemple de matrice, $p$ est pair.
Chaque colonne commence en alternant une valeur nulle/une valeur non-nulle 
jusqu'à l'élément diagonal (qui est nul).


  \item Nous savons que $f(1)=0$ et $f(X)=0$ donc $1$ et $X$ sont dans le noyau $\Ker f$.
Il est aussi clair que les colonnes de la matrices $f(X^2),\cdots, f(X^n)$ sont linéairement indépendantes
(car la matrice est échelonnée). Donc $\Im f = \textrm{Vect}\{f(X^2),f(X^3),\ldots,f(X^n)\}$ 
et $\dim \Im f = n-1$.

Par la formule du rang $\dim \Ker f + \dim \Im f = \dim \Rr_n[X]$ donc
$\dim \Ker f = 2$. Comme nous avons déjà deux vecteurs du noyau alors 
$\Ker f =  \textrm{Vect}\{1,X\}$.

   \item 
   \begin{enumerate}
      \item Soit $Q \in \Im f$. Il existe donc $R\in\Rr_n[X]$ tel que $f(R)=Q$.
On pose ensuite $P(X)=R(X)-R(0)-R'(0)X$.
On a tout fait pour que $P(0)=0$ et $P'(0)=0$.
De plus par la linéarité de $f$ et son noyau alors
$$f(P)= f\big( R(X)-R(0)-R'(0)X\big) = f\big( R(X)\big)-R(0)f(1) -R'(0)f(X)=f(R)=Q.$$
Donc notre polynôme $P$ convient.

      \item Montrons l'unicité. Soient $P$ et $\tilde P$ tels que $f(P)=f(\tilde P)=Q$
avec $P(0)=P'(0)=0 = \tilde P(0)=\tilde P'(0)$.
Alors  $f(P-\tilde P) = Q-Q=0$ donc $P-\tilde P \in \Ker f = \textrm{Vect}\{1,X\}$.
Ainsi $P-\tilde P$ s'écrit $P-\tilde P = aX+b$. Mais comme $(P-\tilde P)(0)=0$ alors 
$b=0$, et comme  $(P-\tilde P)'(0)=0$ alors $a=0$. Ce qui prouve $P = \tilde P$.

    \end{enumerate}
\end{enumerate}
\fincorrection
\nocorrection
\nocorrection
\correction{001097}
\begin{enumerate}
  \item 
  On note la base $\mathcal{B}=(e_1,e_2,e_3)$
  et $X=\begin{pmatrix}x\\y\\z\end{pmatrix}_{\mathcal{B}}= x e_1+y e_2+z e_3$.
  La matrice $A=\textrm{Mat}_{\mathcal{B}}(f)$ est composée des vecteurs colonnes $\phi(e_i)$,
on sait 
$$\phi(e_1)=e_3 = \begin{pmatrix}0\\0\\1\end{pmatrix}_{\mathcal{B}} \quad
\phi(e_2)=-e_1+e_2+e_3 = \begin{pmatrix}-1\\1\\1\end{pmatrix}_{\mathcal{B}} \quad 
\phi(e_3)=e_3 = \begin{pmatrix}0\\0\\1\end{pmatrix}_{\mathcal{B}} \quad
$$

$$\text{donc } \quad A=\begin{pmatrix}
0 & -1 & 0 \\
0 & 1  & 0 \\
1 & 1  & 1 \\    
\end{pmatrix}$$

Le noyau de $\phi$ (ou celui de $A$) est l'ensemble de $X=\begin{pmatrix}x\\y\\z\end{pmatrix}$
tel que $AX=0$.

$$AX=0 \iff \begin{pmatrix}
0 & -1 & 0 \\
0 & 1  & 0 \\
1 & 1  & 1 \\  
\end{pmatrix} \times \begin{pmatrix}x\\y\\z\end{pmatrix}
=\begin{pmatrix}0\\0\\0\end{pmatrix}
\iff \left\{
\begin{array}{rcl}
-y&=&0\\
y&=&0\\
x+y+z&=&0\\
\end{array}\right.
$$
Donc $\Ker \phi = \big\{ \begin{pmatrix}x \\ 0 \\-x\end{pmatrix}_{\mathcal{B}}  \in \Rr^3 \mid x\in \Rr \big\}= 
\textrm{Vect} \begin{pmatrix}1\\0\\-1\end{pmatrix}_{\mathcal{B}} = \textrm{Vect} (e_1-e_3)$.
Le noyau est donc de dimension $1$.


  \item On applique le pivot de Gauss comme si c'était un système linéaire :
$$\left\{
\begin{array}{cccccclr}
e_1  & &     &-& e_3  &=& f_1 &_{L_1}\\
e_1  &-& e_2 & &      &=& f_2 &_{L_2}\\
-e_1 &+& e_2 &+& e_3  &=& f_3 &_{L_3}\\
\end{array}\right.
\iff  \left\{
\begin{array}{cccccclr}
e_1  & &     &-& e_3  &=& f_1 &\\
     &-& e_2 &+& e_3  &=& f_2-f_1 &_{L_2-L_1}\\
     & & e_2 & &      &=& f_3+f_1 &_{L_3+L_1}\\
\end{array}\right.
$$
On en déduit
$$\left\{
\begin{array}{rcl}
 e_1 &=& f_1+f_2+f_3 \\
 e_2 &=& f_1+f_3\\
 e_3 &=& f_2+f_3 \\
\end{array}\right.
$$

Donc tous les vecteurs de la base $\mathcal{B}=(e_1,e_2,e_3)$ s'expriment en fonction
de $(f_1,f_2,f_3)$, ainsi la famille $(f_1,f_2,f_3)$ est génératrice.
Comme elle a exactement $3$ éléments dans l'espace vectoriel $\Rr^3$ de dimension $3$ alors
$\mathcal{B}'=(f_1,f_2,f_3)$ est une base.


  \item
$$\phi(f_1)=\phi(e_1-e_3)=\phi(e_1)-\phi(e_3)=e_3-e_3=0$$

$$\phi(f_2)=\phi(e_1-e_2)= \phi(e_1)-\phi(e_2)=e_3 - (-e_1+e_2+e_3) = e_1-e_2 = f_2$$

$$\phi(f_3)=\phi(-e_1+e_2+e_3)=-\phi(e_1)+\phi(e_2)+\phi(e_3)=-e_1+e_2+e_3=f_3$$

Donc, dans la base $\mathcal{B}'=(f_1,f_2,f_3)$, nous avons
$$\phi(f_1)=0=\begin{pmatrix}0\\0\\0\end{pmatrix}_{\mathcal{B}'}\quad
\phi(f_2)=f_2=\begin{pmatrix}0\\1\\0\end{pmatrix}_{\mathcal{B}'}
\phi(f_3)=f_3=\begin{pmatrix}0\\0\\1\end{pmatrix}_{\mathcal{B}'}$$

Donc la matrice de $\phi$ dans la base $\mathcal{B}'$ est
$$B=\begin{pmatrix}
0 & 0 & 0 \\
0 & 1 & 0 \\
0 & 0 & 1 \\    
\end{pmatrix}$$

$\phi$ est la projection sur  $\textrm{Vect} (f_2,f_3)$ parallèlement à $\textrm{Vect} (f_1)$ (autrement dit
c'est la projection sur le plan d'équation $(x'=0)$, parallèlement à l'axe des $x'$, ceci dans la base $\mathcal{B}'$).

  \item $P$ est la matrice de passage de $\mathcal{B}$ vers $\mathcal{B}'$.
En effet la matrice de passage contient -en colonnes- les coordonnées des vecteurs
de la nouvelle base $\mathcal{B}'$ exprimés dans l'ancienne base $\mathcal{B}$.


Si un vecteur a pour coordonnées $X$ dans la base $\mathcal{B}$ et $X'$ dans la base $\mathcal{B}'$
alors $PX'=X$ (attention à l'ordre).
Et si $A$ est la matrice de $\phi$ dans la base $\mathcal{B}$ et $B$ est la matrice de $\phi$ dans la base
$\mathcal{B}'$ alors
$$B=P^{-1}AP$$
(Une matrice de passage entre deux bases est inversible.)

Ici on calcule l'inverse de $P$ :
$$P^{-1} = \begin{pmatrix}
1 & 1 & 0 \\
1 & 0 & 1 \\
1 & 1 & 1 \\    
\end{pmatrix}
\quad \text{ donc } \quad 
B=P^{-1}AP=\begin{pmatrix}
0 & 0 & 0 \\
0 & 1 & 0 \\
0 & 0 & 1 \\   
\end{pmatrix}
$$

On retrouve donc bien les mêmes résultats que précédemment.

\end{enumerate}

\fincorrection
\correction{001098}
Posons : $ e_1=\begin{pmatrix}1 \cr 2 \cr -1 \cr \end{pmatrix}  ,
  e_2=\begin{pmatrix}3 \cr -1 \cr 1 \cr \end{pmatrix}  ,
  e_{3,\alpha }=\begin{pmatrix}\alpha \cr 2 \cr 2 \cr \end{pmatrix} ,
  e_{4, \beta }=\begin{pmatrix}\beta \cr 1 \cr 0 \cr \end{pmatrix}$. Notons  $\varphi _{\alpha , \beta }$
l'application lin\' eaire associ\' ee \`a   $M _{\alpha , \beta }$
et  $F=\hbox{Vect }\{ e_1 ,  e_2 \}$. Par d\' efinition de la
matrice associ\' ee \`a une application lin\' eaire, $\hbox{Im
}(\varphi _{\alpha , \beta })=\hbox{Vect }\{ e_1 ,  e_2  ,  e_{3,
\alpha }  ,  e_{4, \beta }\}$. En particulier,
 $F\subset \hbox{Im }(\varphi _{\alpha , \beta })$.
Comme  $e_1$  et  $e_2$  sont lin\' eairement ind\' ependants,
$\hbox{rg}(\varphi _{\alpha , \beta })\geq 2$. Ainsi  $\varphi
_{\alpha , \beta }$  est surjective si et seulement si l'un des
deux vecteurs  $e_{3, \alpha }$  ou  $e_{4, \beta }$  n'appartient
pas \`a  $F$. En ce cas en effet, $\hbox{rg}(\varphi _{\alpha ,
\beta })=3=\hbox{dim }{\R}^3$. Or  $e_{3, \alpha }$  et  $e_{4,
\beta }$  appartiennent \`a  $F$
 si et seulement si il existe  $\lambda  , \lambda '  ,  \mu  ,  \mu ' \in
{\R}$  tels que : $e_{3, \alpha }=\lambda e_1 + \mu e_2$  et
$e_{4, \beta }=\lambda 'e_1 + \mu 'e_2$. Un petit calcul montre
donc que  $\varphi _{\alpha , \beta }$  n'est pas surjective si et
seulement si  $\alpha =22$  et  $\beta=4$. Donc  $\varphi _{\alpha
, \beta }$  est surjective si et seulement si  $\alpha \not =22$
ou  $\beta\not =4$.
\fincorrection
\correction{001099}
\begin{enumerate}
  \item 
  \begin{enumerate}
     \item Commençons par des remarques élémentaires : la matrice est non nulle donc $\textrm{rg}(A) \ge 1$
et comme il y a $p=4$ lignes et $n=3$ colonnes alors $\textrm{rg}(A) \le \min(n,p)=3$.

     \item Ensuite on va montrer $\textrm{rg}(A) \ge 2$ en effet le sous-déterminant $2\times 2$ 
(extrait du coin en haut à gauche) :
$\begin{vmatrix} 
1 & 2 \cr
3 & 4 \cr
\end{vmatrix}= -2$ est non nul.

     \item Montrons que $\textrm{rg}(A)=2$. Avec les déterminants il faudrait vérifier que pour toutes
les sous-matrices $3\times 3$ les déterminants sont nuls. Pour éviter de nombreux calculs on remarque ici
que les colonnes sont liées par la relation $v_2=v_1+v_3$. Donc $\textrm{rg}(A)=2$.

     \item L'application linéaire associée à la matrice $A$ est l'application 
$f_A : \Rr^3 \to \Rr^4$. Et le théorème du rang 
$\dim \Ker f_A+ \dim \Im f_A = \dim \Rr^3$ donne ici
$\dim \Ker f_A = 3 - \textrm{rg}(A)=1$.

Mais la relation $v_2=v_1+v_3$ donne immédiatement un élément du noyau :
en écrivant $v_1-v_2+v_3=0$ alors $A\begin{pmatrix}1\\-1\\1\end{pmatrix}=\begin{pmatrix}0\\0\\0\end{pmatrix}$
Donc $\begin{pmatrix}1\\-1\\1\end{pmatrix} \in \Ker f_A$. Et comme le noyau est de dimension $1$ alors
$$\Ker f_A = \textrm{Vect} \begin{pmatrix}1\\-1\\1\end{pmatrix}$$

     \item Pour un base de l'image, qui est de dimension $2$, 
     il suffit par exemple de prendre les deux premiers vecteurs colonnes de la matrice $A$ (ils sont clairement non colinéaires) :
$$\Im f_A = \textrm{Vect} \left\{ v_1, v_2 \right\} = 
\textrm{Vect} \left\{  \begin{pmatrix}1\\3\\5\\7\end{pmatrix},  \begin{pmatrix}1\\1\\1\\1\end{pmatrix} \right\}
$$

  \end{enumerate}


  \item On fait le même travail avec $B$ et $f_B$.
  \begin{enumerate}
     \item Matrice non nulle avec $4$ lignes et $4$ colonnes donc $1 \le \textrm{rg}(B) \le 4$.

     \item Comme le sous-déterminant (du coin supérieur gauche)
$\begin{vmatrix} 
2 & 2 \cr
4 & 3 \cr
\end{vmatrix}= -2$ est non nul alors $\textrm{rg}(B) \ge 2$.

     \item Et pareil avec le sous-déterminant $3\times 3$ :
$$\begin{vmatrix} 
2 & 2 & -1 \cr
4 & 3 & -1 \cr
0 & -1 & 2 \cr
\end{vmatrix} = -2$$
qui est non nul donc $\textrm{rg}(B) \ge 3$.

     \item Maintenant on calcule le déterminant de la matrice $B$ et 
on trouve $\det B = 0$, donc $\textrm{rg}(B) < 4$. Conclusion $\textrm{rg}(B) = 3$.
Par le théorème du rang alors $\dim \Ker f_B=1$.



     \item Cela signifie que les colonnes (et aussi les lignes) sont liées, comme il n'est pas clair
de trouver la relation à la main on résout le système $B X = 0$ pour trouver cette relation ; autrement dit :
$$\begin{pmatrix} 
2 & 2 & -1 & 7  \cr
4 & 3 & -1 & 11 \cr
0 & -1 & 2 & -4 \cr
3 & 3 & -2 & 11 \cr 
\end{pmatrix}
\cdot 
\begin{pmatrix} x\\y\\z\\t \end{pmatrix} = 
\begin{pmatrix} 0\\0\\0\\0 \end{pmatrix}
\text{ ou encore }
\left\{ 
\begin{array}{rcl}
2x + 2y -z + 7t  &=& 0  \cr
4x + 3y -z + 11t &=& 0 \cr
     -y +2z -4t  &=& 0  \cr
3x + 3y -2z+ 11t &=& 0  \cr   
\end{array}
\right.$$
Après résolution de ce système on trouve que 
les solutions s'écrivent $(x,y,z,t)= (-\lambda,-2\lambda,\lambda,\lambda)$.
Et ainsi 
$$\Ker f_B = \textrm{Vect} \begin{pmatrix}-1\\-2\\1\\1\end{pmatrix}$$
Et pour une base de l'image il suffit, par exemple, de prendre les $3$ premiers vecteurs colonnes $v_1,v_2,v_3$ 
de la matrice $B$, car ils sont linéairement indépendants :
$$\Im f_B = \textrm{Vect} \left\{ v_1, v_2, v_3 \right\} = 
\textrm{Vect} \left\{  
\begin{pmatrix}2\\4\\0\\3\end{pmatrix},  
\begin{pmatrix}2\\3\\-1\\3\end{pmatrix},
\begin{pmatrix}-1\\-1\\2\\-2\end{pmatrix} 
\right\}
$$

  \end{enumerate}
\end{enumerate}


\fincorrection
\correction{001100}
$\mathcal{L} (E)$  est isomorphe \`a  $M_n ({\R})$  donc est de
dimension finie  $n^2$. La famille  $\{ id_E ,  \varphi  , \ldots
,  \varphi ^{n^2}\} $  compte  $n^2+1$  vecteurs donc est li\' ee
c'est \`a dire : il existe  $\lambda _0 , \ldots , \lambda _{n^2}
$  dans  ${\R}$, non tous nuls et tels que  $\lambda _0id_E +
\lambda _1 \varphi  + \cdots +\lambda _{n^2} \varphi ^{n^2}=0$.
 Le polyn\^ome  $P(X)=\lambda _0 + \lambda _1 X  + \cdots
+\lambda _{n^2} X^{n^2}$  r\' epond donc \`a la question.
\fincorrection
\nocorrection
\nocorrection
\correction{001104}
\begin{enumerate}
  \item Notons $P$ la matrice de passage de la base canonique 
$\mathcal{B}=\big((1,0),(0,1)\big)$ vers (ce qui va être) la base
$\mathcal{B}' = (e_1, e_2)$. C'est la matrice composée des vecteurs colonnes
$e_1$ et $e_2$ :
$$P = \begin{pmatrix} -2 & -2 \\ 3 & 5 \\  \end{pmatrix}$$
$\det P=-4 \neq 0$ donc $P$ est inversible et ainsi $\mathcal{B}'$ est bien une base.

Alors la matrice de $f$ dans la base $\mathcal{B}'$ est :
$$B= P^{-1}AP = -\frac14\begin{pmatrix} 5 & 2 \\ -3 & -2 \\  \end{pmatrix} \begin{pmatrix} 2&\frac 23\\
-\frac 52&-\frac 23 \end{pmatrix}\begin{pmatrix} -2 & -2 \\ 3 & 5 \\  \end{pmatrix}=
\begin{pmatrix} 1 & 0 \\ 0 & \frac13 \\  \end{pmatrix}$$

  \item Il est très facile de calculer la puissance d'une matrice diagonale :
$$B^n=\begin{pmatrix} 1 & 0 \\ 0 & \big(\frac13\big)^n \\  \end{pmatrix}$$

Comme $A=PBP^{-1}$ on va en déduire $A^n$ : 

$$A^n = \big( PBP^{-1} \big)^n = P B^n P^{-1} = 
\frac14\begin{pmatrix} 
10- \frac{6} {3^n} & 4- \frac{4} {3^n}\\
-15 + \frac{15} {3^n}& -6 + \frac{10} {3^n}\\                    
       \end{pmatrix}$$

  \item Si l'on note $X_n = \begin{pmatrix}x_n \\ y_n \end{pmatrix}$
alors les équations que vérifient les suites s'écrivent en terme matriciel :
$$X_{n+1}=AX_n.$$

Si l'on note les conditions initiales $X_0 = \begin{pmatrix}x_0 \\ y_0 \end{pmatrix} \in \Rr^2$ alors
$X_n = A^n X_0$.
On en déduit 
$$
\left\{\begin{array}{rcl}
x_n &=& \frac14 \Big(( 10- \frac{6} {3^n}) x_0 + (4- \frac{4} {3^n})  y_0 \Big)          \\
y_n &=& \frac14 \Big( (-15 + \frac{15} {3^n})x_0 + (-6 + \frac{10} {3^n}) y_0 \Big)   
\end{array}\right.
$$
\end{enumerate}


\fincorrection
\nocorrection
\nocorrection
\correction{002565}
Soit $A=(a_{ij})_{1\leq i,j\leq n}$ une matrice carr\'ee $n\times n$. On veut d\'emontrer le
 r\'esultat suivant d\^u \`a Hadamard : Supposons que pour tout $i\in\{1,\cdots,n\}$, on ait
$$|a_{ii}|>\sum_{j=1,j\neq i}^{n}|a_{ij}| $$
alors $A$ est inversible.



\begin{enumerate}  
\item {\it Montrons le r\'esultat pour $n=2$.}

Dans ce cas, la matrice $A$ s'\'ecrit 
$$A=\begin{pmatrix}a_{11}&a_{12} \\  a_{21}&a_{22}\end{pmatrix}$$ et les hypoth\`eses deviennent
$$|a_{11}|>|a_{12}|\ {\hbox{et}}\ |a_{22}|>|a_{21}|.$$
La matrice $A$ est inversible si et seulement si son d\'eterminant est non nul, or 
$$\det A=a_{11}a_{22}-a_{12}a_{21},$$
 et, compte tenu des hypoth\`eses,
$$|a_{11}a_{22}|=|a_{11}||a_{22}|>|a_{12}||a_{21}|=|a_{12}a_{21}|,$$
ainsi$|a_{11}a_{22}|>|a_{12}a_{21}|$ donc $a_{12}a_{21}\neq a_{12}a_{21}$ et le d\'eterminant est non nul.

  \item Soit $B$, la matrice obtenue en rempla\c cant, pour $j\geq 2$, chaque colonne $c_j$ de $A$ par la colonne
$$ c_j-{\frac{a_{1j}}{a_{11}}}c_1 ,$$ {\it Calculons les $b_{ij}$ en fonction des $a_{ij}$. Montrons que si les
coefficients de $A$ satisfont les in\'egalit\'es ci-dessus, alors pour $i\geq 2$, on a
$$|b_{ii}|>\sum_{j=2,j\neq i}^{n}|b_{ij}| .$$}

On a $$b_{ij}=a_{ij}-\frac{{a_{1j}}{a_{11}}}a_{i1}\ \ {\hbox{si}}\ \ j\geq2\ \ {\hbox{et}}\ \ b_{i1}=a_{i1}.$$
 par l'in\'egalit\'e triangulaire, on a
 \begin{align*}\sum_{j=2,j\neq i}|b_{ij}|&=\sum_{j=2,j\neq i}|a_{ij}-{\frac{a_{1j}}{a_{11}}}a_{i1}| \\ 
&\leq \sum_{j=2,j\neq i}|a_{ij}|+\frac{|{a_{1j}|}{|a_{11}|}}|a_{i1}| \\ 
&= \sum_{j=2,j\neq i}|a_{ij}|+\frac{|{a_{i1}|}{|a_{11}|}} \sum_{j=2,j\neq i}|a_{1j}|. 
\end{align*}
Mais, par hypoth\`ese, pour $i=1$, on a
$$\sum_{j=2}^{n}|a_{1j}|<|a_{11}|,$$
donc
$$\sum_{j=2,j\neq i}^{n}|a_{1j}|<|a_{11}|-|a_{1i}|.$$
D'o\`u, en rempla\c cant dans l'in\'egalit\'e pr\'ec\'edente
\begin{align*}\sum_{j=2,j\neq i}|b_{ij}|
&< \sum_{j=2,j\neq i}|a_{ij}|+ |a_{i1}|-\frac{{|a_{i1}|}{|a_{11}|}}|a_{1i}| \\ 
&= \sum_{j=1,j\neq i}|a_{ij}|-\frac{{|a_{i1}|}{|a_{11}|}}|a_{1i}| \\ 
&<|a_{ii}|-\frac{{|a_{i1}|}{|a_{11}|}}|a_{1i}| \\ 
&\leq \left|a_{ii}-\frac{{a_{i1}}{a_{11}}}a_{1i}\right|=|b_{ii}|.
\end{align*}

  \item{\it  D\'emontrons le r\'esultat de Hadamard pour $n$ quelconque.}

Soit $A=(a_{ij})_{1\leq i,j\leq n}$ une matrice carr\'ee $n\times n$, v\'erifiant pour tout $i\in\{1,\cdots,n\}$, 
$$|a_{ii}|>\sum_{j=1,j\neq i}^{n}|a_{ij}| $$
On veut d\'emontrer que $A$ est inversible.

Le r\'esultat est vrai pour $n=2$, d'apr\`es la question $1)$. Soit $n$ arbitrairement fix\'e, supposons le r\'esultat vrai pour
$n-1$ et d\'emontrons le pour $n$.

On a $\det A=\det B$ o\`u $B$ est la matrice construite dans la question $2)$
$$B=\begin{pmatrix}a_{11}&0&\cdots&0 \\ 
\vdots& &({b_{ij}}_{(2\leq i,j\leq n)}) \\ 
a_{n1}\end{pmatrix}$$
Or, la matrice $({b_{ij}}_{(2\leq i,j\leq n)})$ est une matrice carr\'ee d'ordre $n-1$ qui v\'erifie les hypoth\`eses de
Hadamard, d'apr\`es la question $2)$. Elle est donc inversible par hypoth\`ese de r\'ecurrence. Et, par cons\'equent, la 
matrice $A$ est inversible car $a_{11}\neq 0$.
\end{enumerate}
\fincorrection
\correction{002585}
Soient $A$ et $B$ des matrices non nulles de $M_n(\R)$. On suppose que $A.B=0$.
\begin{enumerate}
 \item D\'emontrons que $\Im B\subset \ker A$.

Soit $y\in \Im B$, il existe $x\in\R^n$ tel que $y=Bx$, d'o\`u $Ay=ABx=0$, ainsi $y\in\ker A$ ce qui prouve l'inclusion.
 \item On suppose que le rang de $A$ est \'egal \`a $n-1$, d\'eterminons le rang de $B$.

On a $\mathrm{rg} B=\dim \Im B$ et on sait que $\dim\Im A+\dim\ker A=n$ par cons\'equent, si $\mathrm{rg} A=n-1$ on a $\dim\ker A=1$ et l'inclusion $\Im B\subset \ker A$ implique $\dim\Im B\leq 1$ or, $B$ est suppos\'ee non nulle d'o\`u $\dim\Im B=1=\mathrm{rg} B$.
\end{enumerate}
\fincorrection
\nocorrection
\nocorrection
\correction{003407}
\begin{enumerate}
  \item $\phi(P) = (-X-1)^{n-1}P\left(-\frac 1{X+1}\right)$.
  \item $I$.
\end{enumerate}
 \fincorrection
\correction{005261}
$f$ n'est pas nul et donc $\mbox{dim}(\mbox{Ker}f)\leq 2$. Puisque $f^2=0$, $\mbox{Im}f\subset\mbox{Ker}f$. En particulier, $\mbox{dim}(\mbox{Ker}f)\geq\mbox{rg}f=3-\mbox{dim}(\mbox{Ker}f)$ et $\mbox{dim}(\mbox{Ker}f)\geq\frac{3}{2}$.

Finalement, $\mbox{dim}(\mbox{Ker}f)=2$. $\mbox{Ker}f$ est un plan vectoriel et $\mbox{Im}f$ est une droite vectorielle contenue dans $\mbox{Ker}f$.

$f$ n'est pas nul et donc il existe $e_1$ tel que $f(e_1)\neq0$ (et en particulier $e_1\neq0$). Posons $e_2=f(e_1)$. 
Puisque $f^2=0$, $f(e_2)=f^2(e_1)=0$ et $e_2$ est un vecteur non nul de $\mbox{Ker}f$. D'après le théorème de la base incomplète, il existe un vecteur $e_3$ de $\mbox{Ker}f$ tel que $(e_2,e_3)$ soit une base de $\mbox{Ker}f$.

Montrons que $(e_1,e_2,e_3)$ est une base de $\Rr^3$.

Soit $(\alpha,\beta,\gamma)\in\Rr^3$.

$$\alpha e_1+\beta e_2+\gamma e_3=0\Rightarrow f(\alpha e_1+\beta e_2+\gamma e_3)=0\Rightarrow\alpha e_2=0\Rightarrow\alpha=0\;(\mbox{car}\;e_2\neq0).$$

Puis, comme $\beta e_2+\gamma e_3=0$, on obtient $\beta=\gamma=0$ (car la famille $(e_2,e_3)$ est libre).

Finalement, $\alpha=\beta=\gamma=0$ et on a montré que $(e_1,e_2,e_3)$ est libre. Puisque cette famille est de cardinal $3$, c'est une base de $R^3$. Dans cette base, la matrice $A$ de $f$ s'écrit~:~$A= \left(
\begin{array}{ccc}
0&0&0\\
1&0&0\\
0&0&0
\end{array}
\right)$.
\fincorrection
\correction{007411}
\begin{enumerate}
 \item $$A = [f]_{\mathcal{B}}^{\mathcal{B}} = \left( {\begin{array}{ccc}
 4 & 1 & 1 \\
 4 & 7 & 2 \\
 -6 & -6 & -1 \\
 \end{array} } \right)$$

 \item Pour montrer que les vecteurs $v_1$, $v_2$ et $v_3$ forment un système libre de $\Rr^3$, on considère $\lambda_1,~\lambda_2,~\lambda_3 \in \Rr$ tels que 
	$$\lambda_1 v_1 + \lambda_2 v_2 + \lambda_3 v_3 = 0$$
et on veut montrer qu'alors nécessairement $\lambda_1 = \lambda_2 = \lambda_3 = 0$. Cela revient à montrer que le système
	$$\left( {\begin{array}{ccc}
 1 & 1 & 0 \\
 0 & -1 & -1 \\
 -2 & 0 & 1 \\
 \end{array} } \right) 
\left( {\begin{array}{c}
\lambda_1 \\
\lambda_2 \\
\lambda_3
 \end{array} } \right) = 
 \left( {\begin{array}{c}
0 \\
0 \\
0
 \end{array} } \right)$$ 
possède pour unique solution $\lambda_1 = \lambda_2 = \lambda_3 = 0$.
On applique alors l'algorithme de Gauss à la matrice 
$$\left( {\begin{array}{ccc}
 1 & 1 & 0 \\
 0 & -1 & -1 \\
 -2 & 0 & 1 \\
 \end{array} } \right) $$
 Après avoir échelonné complètement cette matrice, on obtient le système équivalent suivant
 $$\left( {\begin{array}{ccc}
 1 & 0 & 0 \\
 0 & 1 & 0 \\
 0 & 0 & 1 \\
 \end{array} } \right)
\left( {\begin{array}{c}
\lambda_1 \\
\lambda_2 \\
\lambda_3
 \end{array} } \right) = 
 \left( {\begin{array}{c}
0 \\
0 \\
0
 \end{array} } \right)$$ 
 qui a pour unique solution $\lambda_1 = \lambda_2 = \lambda_3 = 0$. Les vecteurs $v_1$, $v_2$ et $v_3$ forment donc un système libre de $3$ éléments dans l'espace vectoriel $\Rr^3$ qui est de dimension $3$, donc ils forment un système également générateur de $\Rr^3$ et c'est une base de $\Rr^3$.

 \item La matrice de passage $P := [Id_{R^3}]_{\mathcal{B}'}^{\mathcal{B}}$ de $\mathcal{B}$ à $\mathcal{B}'$ est 
	$$P = 
\left( {\begin{array}{ccc}
 1 & 1 & 0 \\
 0 & -1 & -1 \\
 -2 & 0 & 1 \\
 \end{array} } \right)$$
 La matrice de passage $[Id_{\Rr^3}]_{\mathcal{B}}^{\mathcal{B}'}$ de $\mathcal{B}'$ à $\mathcal{B}$ est l'inverse de la matrice de passage $[Id_{\Rr^3}]_{\mathcal{B}'}^{\mathcal{B}}$ de $\mathcal{B}$ à $\mathcal{B}'$ i.e. $[Id_{R^3}]_{\mathcal{B}}^{\mathcal{B}'} = P^{-1}$. Pour calculer l'inverse de $P$, on applique l'algorithme de Gauss à la matrice 
$$\left( {\begin{array}{ccc | ccc}
 1 & 1 & 0 & 1 & 0 & 0 \\
 0 & -1 & -1 & 0 & 1 & 0 \\
 -2 & 0 & 1 & 0 & 0 & 1 \\
 \end{array} } \right)$$
 Pour ce faire, on reprend les opérations effectuées sur la matrice $P$ lors de la question précédente pour obtenir 
$$\left( {\begin{array}{ccc | ccc}
 1 & 0 & 0 & -1 & -1 & -1 \\
 0 & 1 & 0 & 2 & 1 & 1 \\
 0 & 0 & 1 & -2 & -2 & -1 \\
 \end{array} } \right)$$ 
et ainsi 
$$[Id_{R^3}]_{\mathcal{B}}^{\mathcal{B}'} = P^{-1} = 
\left( {\begin{array}{ccc}
  -1 & -1 & -1 \\
 2 & 1 & 1 \\
 -2 & -2 & -1 \\
 \end{array} } \right)$$

 \item On utilise la formule 
 $$[f]_{\mathcal{B}'}^{\mathcal{B}'} = [Id_{\Rr^3}]_{\mathcal{B}}^{\mathcal{B}'} [f]_{\mathcal{B}}^{\mathcal{B}} [Id_{\Rr^3}]_{\mathcal{B}'}^{\mathcal{B}}$$
 et on a alors 
\begin{equation*}
M = P^{-1}A P = 
 \left( {\begin{array}{ccc}
 2 & 0 & 0 \\
 0 & 3 & 0 \\
 0 & 0 & 5 \\
 \end{array} } \right)\end{equation*}

 \item On a
 $$A = P M P^{-1}$$
 donc pour tout $n \in \mathbb{N}^*$
 $$A^n = P M^n P^{-1} = 
 \left( {\begin{array}{ccc}
 1 & 1 & 0 \\
 0 & -1 & -1 \\
 -2 & 0 & 1 \\
 \end{array} } \right) 
 \left( {\begin{array}{ccc}
  2^n & 0 & 0 \\
 0 & 3^n & 0 \\
 0 & 0 & 5^n \\
 \end{array} } \right) 
 \left( {\begin{array}{ccc}
  -1 & -1 & -1 \\
 2 & 1 & 1 \\
 -2 & -2 & -1 \\
 \end{array} } \right) $$
 $$= 
 \left( {\begin{array}{ccc}
  -2^n + 2.3^n & -2^n + 3^n & -2^n + 3^n \\
 -2.3^n + 2.5^n & -3^n + 2.5^n & -3^n + 5^n \\
 2.2^n - 2.5^n & 2.2^n - 2.5^n & 2.2^n - 5^n \\
 \end{array} } \right)$$

\end{enumerate}
\fincorrection
\correction{007412}
\begin{enumerate}
	\item Soit $\lambda,\,\mu\in\R$ et $P,\, Q\in \R[X]$. Comme la dérivation est une application linéaire on a bien :
\begin{equation*}
\Phi(\lambda\,P+\mu\,Q)=\lambda\Phi(P)+\mu\Phi(Q).
\end{equation*}
$\Phi$ est une application linéaire. Examinons l'image des vecteurs de base $(1,X,X^2)$ par $\Phi$ :
\begin{equation*}
\Phi(1)(X)=2X+1,\qquad \Phi(X)(X)=X^2+X+1,\qquad \Phi(X^2)(X)=X^2+2X.
\end{equation*}
donc $\Phi(X^{j})\in\R_{2}[X]$ pour $j\in\{0,1,2\}$. C'est bien un endomorphisme.

	\item En appliquant la question précédente, on a bien :
\begin{equation*}
[\Phi]_{\mathcal{B}}^{\mathcal{B}}=\left[[\Phi(1)]_{\mathcal{B}}\,|\,[\Phi(X)]_{\mathcal{B}}\,|\,[\Phi(X^2)]_{\mathcal{B}}\right]=A
\end{equation*}

	\item $rg(A)=rg(\Phi)$. On applique une méthode de Gauss à la matrice $A$:
\begin{equation*}
rg(A)=rg\begin{pmatrix}
1&1&0\\
2&1&2\\
0&1&1
\end{pmatrix}=rg\begin{pmatrix}
1&1&0\\
0&-1&2\\
0&1&1
\end{pmatrix}=rg\begin{pmatrix}
1&1&0\\
0&-1&2\\
0&0&3
\end{pmatrix}=3
\end{equation*}

	\item Comme $rg\,A=3=dim\R_{2}[X]$, l'application $\Phi$ est surjective. La matrice $A$ est carrée donc en appliquant le théorème du rang, elle est aussi injective donc bijective.

	\item Le noyau est donnée par :

\begin{align*}
Ker\, (A-I_{3})&=\left\{(x,y,z)\in\R^3,\ \text{tels que}\, (A-I_{3})^t(x,y,z)=0\right\}.
\end{align*}
\begin{equation*}
(A-I_{3})^t(x,y,z)=0\Leftrightarrow y=0,\quad 2x=-2z.
\end{equation*}
Ainsi $Ker\, (A-I_{3})=\left\{(x,0,-x),\, x\in\R\right\}=vect\left((1,0,-1)\right)$. $(-1,0,1)$ est une base de $Ker\, (A-I_{3})$. Une base de $Ker\,(\Phi-I_{d})$ est donc : $X^2-1$.

	\item Par le théorème du rang : $dim\ Im(\Phi-Id)=dim \R_{2}[X]-dim\ Ker\,(\Phi-Id)=3-1=2$. Il reste à déterminer une base de l'image. 
\begin{align*}
Im(A-I_{3})&=\left\{(A-I_{3})^t(x,y,z),\quad (x,y,z)\in\R^3\right\}=\left\{(y,2(x+z),y),\quad (x,y,z)\in\R^3\right\}\\
&=\left\{(y,x,y),\quad (x,y)\in\R^2\right\}=\left\{y(1,0,1)+x(0,1,0),\quad (x,y)\in\R^2\right\}\\
&=vect\left((1,0,1),\, (0,1,0)\right).
\end{align*}
La famille $((1,0,1),\, (0,1,0))$ est une famille génératrice de l'image qui est de dimension $2$. C'est donc une base. On en déduit qu'une base de $Ker\,(\Phi-Id)$ est $((X^2+1),X)$.

	\item On a l'égalité ensembliste :
\begin{align*}
\left\{P\in\R_{2}[X],\, \text{tels que}\ 2X\,P=(X^2-1)\,P'\right\}&=\left\{P\in\R_{2}[X],\ \text{tels que}\, (\Phi-Id)P=0\right\}\\
&=Ker\,(\Phi-Id).
\end{align*}
Or une base de $Ker\, (\Phi-Id)$ est $X^2-1$. Ainsi on obtient :
\begin{equation*}
\left\{P\in\R_{2}[X],\, \text{tels que}\ 2X\,P=(X^2-1)\,P'\right\}=\R(X^2-1).
\end{equation*}
\end{enumerate}
\fincorrection
\nocorrection
\nocorrection
\correction{003362}
\begin{enumerate}
  \item 
  \item $M^{-1} = \frac {-a}{b(na+b)}U + \frac 1b I$.
  \item 
  \item $$M^n = \frac {(na+b)^n-b^n}nU + b^nI  \Rightarrow 
             \left\{
             \begin{array} {ll}
             n \text{ pair }  &: a = 0, \text{ ou }
                                       -\frac{\smash{2b}}n,\ b = \pm1 \cr
                    n \text{ impair} &: a = 0,\ b = 1. \cr
                    \end{array}\right.$$
\end{enumerate}
 \fincorrection
\correction{003365}
2. $M^{-1} = \begin{pmatrix} A^{-1} & -A^{-1}BC^{-1} \cr 0 & C^{-1} \cr\end{pmatrix}$.
 \fincorrection
\nocorrection
\correction{003376}
\begin{enumerate}
  \item 
  \item 
    Si $A$ est diagonale :
    $M = \begin{pmatrix}0  &\dots  &0       &1      \cr
                  1  &\ddots &        &0      \cr
                     &\ddots &\ddots  &\vdots \cr
                  0  &       &1       &0      \cr \end{pmatrix}$.
    Si $a_{k\ell} \ne 0$ : $M = I - \frac {\mathrm{tr} A}{a_{k\ell}}E_{\ell k}$.
\end{enumerate}
 \fincorrection
\correction{003380}

Avant de commencer la résolution nous allons faire une remarque importante :
pour $X = \begin{pmatrix}x_1 \\ x_2 \\ \vdots \\ x_n \end{pmatrix}$ un vecteur 
(considéré comme une matrice à une seule colonne)
alors nous allons calculer ${}^t\!X X$ : 
$${}^t\!X X = (x_1,x_2,\cdots,x_n) \begin{pmatrix}x_1 \\ x_2 \\ \vdots \\ x_n \end{pmatrix} = x_1^2+x_2^2+\cdots+x_n^2.$$
On note $ \| X \|^2 = {}^t\!X X$ : $\| X \|$ est la \emph{norme} ou la \emph{longueur} du vecteur $X$.
De ce calcul on déduit d'une part que ${}^t\!X X \ge 0$. Et aussi que  ${}^t\!X X \ge 0$ si et seulement si
$X$ est le vecteur nul.


\begin{enumerate}
  \item Nous allons montrer que $I+M$ est inversible en montrant que si un vecteur $X$ vérifie
$(I+M)X = 0$ alors $X=0$.

Nous allons estimer ${}^t\!(MX)(MX)$ de deux façons.
D'une part c'est un produit de la forme ${}^t\!Y Y = \| Y\|^2$ et donc ${}^t\!(MX)(MX) \ge 0$.

D'autre part :
\begin{align*}
{}^t\!(MX)(MX) 
 & = {}^t\!(MX) (-X) \quad \text{ car } (I+M)X = 0 \text{ donc } MX = -X \\
 & = {}^t\!X {}^t\!M (-X)    \quad \text{ car } {}^t\!(AB)={}^t\! B {}^t\!A \\
 & = {}^t\!X (-M) (-X)    \quad \text{ car } {}^t\!M=-M \\
 & = {}^t\!X MX \\
 & = {}^t\!X (-X) \\
 & = -{}^t\!XX \\
 & = - \|X\|^2 \\
\end{align*}
Qui est donc négatif.

Seule possibilité $\|X\|^2=0$ donc $X=0$ (= le vecteur nul) et donc $I+M$ inversible.

\item 
  \begin{enumerate}
    \item Calculons $A^{-1}$.
$$A^{-1} =\big( (I-M)\times(I+M)^{-1} \big)^{-1} = \big( (I+M)^{-1} \big)^{-1} \times(I-M)^{-1} = (I+M)\times(I-M)^{-1}$$
(n'oubliez pas que $(AB)^{-1}=B^{-1}A^{-1}$).


     \item Calculons ${}^t\!A$.
\begin{align*}
{}^t\!A 
  & = {}^t\!\big( (I-M) \times (I+M)^{-1} \big) \\
  & = {}^t\!\big( (I+M)^{-1} \big) \times {}^t\!(I-M) \qquad \text{ car } {}^t\!(AB)={}^t\!B{}^t\!A \\
  & = \big( {}^t\!(I+M) \big)^{-1}\times{}^t\!(I-M) \qquad \text{ car } {}^t\!(A^{-1})=\big({}^tA\big)^{-1} \\
  & = \big( I+{}^t\!M) \big)^{-1}\times(I-{}^t\!M) \qquad \text{ car } {}^t\!(A+B)={}^t\!A+{}^t\!B \\
  & = (I-M) ^{-1}\times (I+M) \qquad \text{ car ici } {}^{t}\!{M} = -M \\  
\end{align*}

      \item Montrons que $I+M$ et $(I-M) ^{-1}$ commutent.

Tout d'abord $I+M$ et $I-M$ commutent car $(I+M)(I-M) = I-M^2 = (I-M)(I+M)$.
Maintenant nous avons le petit résultat suivant :

\textbf{Lemme.} Si $AB=BA$ alors $AB^{-1}=B^{-1}A$.

Pour la preuve on écrit :
$$AB=BA \Rightarrow B^{-1}(AB)B^{-1}=B^{-1}(BA)B^{-1} \Rightarrow B^{-1}A=AB^{-1}.$$

En appliquant ceci à $I+M$ et $I-M$ on trouve $(I+M)\times (I-M) ^{-1}= (I-M) ^{-1}\times(I+M)$
et donc $A^{-1}={}^t\!A$.

    \end{enumerate}
\end{enumerate}

\fincorrection
\correction{003381}
3. $X = -A$ ou $X = \frac 12A$ ou $X = A-I$ ou $X = -\frac 12A - I$.
 \fincorrection
\correction{003382}
\begin{enumerate}
  \item $J^2 = J$.
  \item $JM = MJ$.
  \item $k$ = $\mathrm{rg} J$.
\end{enumerate}
 \fincorrection
\correction{003398}
\begin{enumerate}
  \item $\frac {A + (2-n)I}{n-1}$.
  \item $\frac 1{(a-b)(a+(n-1)b)}
              \begin{pmatrix} a+(n-2)b &       &(-b)     \cr
                                 &\ddots           \cr
                       (-b)      &       &a+(n-2)b \cr\end{pmatrix}$.
  \item $\begin{pmatrix} 1 &-1     &1      &\dots   &\pm1   \cr
                          &\ddots &\ddots &\ddots  &\vdots \cr
                          &       &\ddots &\ddots  &1      \cr
                          &       &       &\ddots  &-1     \cr
                       (0)&       &       &        &1      \cr \end{pmatrix}$.
  \item $\frac 1{1-\alpha\bar\alpha}
            \begin{pmatrix} 1       &-\bar\alpha        &0           \cr
                      -\alpha &1+\alpha\bar\alpha &-\bar\alpha \cr
                      0       &-\alpha            &1           \cr\end{pmatrix}$.
  \item $\begin{pmatrix} (0) &&1/a_n \cr &\cdots\cr 1/a_1 &&(0) \cr\end{pmatrix}$.
  \item $\text{diag}(\lambda_i) -            \frac 1{1 + \lambda_1 + \dots + \lambda_n}(\lambda_i\lambda_j)$.
\end{enumerate}
 \fincorrection
\correction{003399}
$A^{-1} = \begin{pmatrix} 9   &-36  &30   \cr
                              -36 &192  &-180 \cr
                              30  &-180 &180  \cr \end{pmatrix}$,
         $B^{-1} \approx \begin{pmatrix} 55.6   &-277.8  &255.6   \cr
                                    -277.8 &1446.0  &-1349.2 \cr
                                    255.6  &-1349.2 &1269.8  \cr \end{pmatrix}$.
\fincorrection
\correction{005267}
Soit $\mathcal{B}=(e_i)_{1\leq i\leq n}$ la base canonique de $\Cc^n$ et $(e_i')_{1\leq i\leq n}$ la famille d'éléments de $\Cc^n$ de matrice $A$ dans la base $\mathcal{B}$.

Par définition, on a

$$\forall i\in\{1,...,n-1\},\;e_i'=ie_i+\sum_{j=i+1}^{n}e_j\;\mbox{et}\;e_n'=ne_n.$$

En retranchant membre à membre ces égalités, on obtient

$$\forall i\in\{1,...,n-1\},\;e_i'-e_{i+1}'=i(e_i-e_{i+1})\;\mbox{et}\;e_n'=ne_n,$$

ou encore

$$\forall i\in\{1,...,n-1\},\;e_i-e_{i+1}=\frac{1}{i}(e_i'-e_{i+1}')\;\mbox{et}\;e_n=\frac{1}{n}e_n'.$$

Mais alors, pour $i\in\{1,...,n-1\}$, on a

\begin{align*}\ensuremath
e_i&=\sum_{j=i}^{n-1}(e_j-e_{j+1})+e_n=\sum_{j=i}^{n-1}\frac{1}{j}(e_j'-e_{j+1}')+\frac{1}{n}e_n'
=\sum_{j=i}^{n-1}\frac{1}{j}e_j'-\sum_{j=i+1}^{n}\frac{1}{j-1}e_j'+\frac{1}{n}e_n'\\
 &=\frac{1}{i}e_i'+\sum_{j=i+1}^{n}\frac{1}{j}e_j'-\sum_{j=i+1}^{n}\frac{1}{j-1}e_j'\\
 &=\frac{1}{i}e_i'-\sum_{j=i+1}^{n}\frac{1}{j(j-1)}e_j'
\end{align*}

Mais alors, $\Cc^n=\mbox{Vect}(e_1,...,e_n)\subset\mbox{Vect}(e_1',...,e_n')$, ce qui montre que la famille $\mathcal{B}'=(e_1',...,e_n')$ est génératrice de $\Cc^n$ et donc une base de $\Cc^n$. Par suite, $A$ est inversible et 

$$A^{-1}=\mbox{Mat}_{\mathcal{B}'}\mathcal{B}=(a_{i,j}')_{1\leq i,j\leq n}\;\mbox{où}\;
a_{i,j}'=
\left\{
\begin{array}{l}
\frac{1}{i}\;\mbox{si}\;i=j\\
-\frac{1}{i(i-1)}\;\mbox{si}\;i>j\\
0\;\mbox{si}\;i<j
\end{array}
\right..
$$
\fincorrection
\correction{005272}
Soit $X=\left(
\begin{array}{c}
x_1\\
\vdots\\
x_n
\end{array}
\right)$  un vecteur du noyau de $A$. Supposons $X\neq0$. Alors, si $i_0$ est un indice tel que $|x_{i_0}|=\mbox{Max}\{|x_i|,\;i\in\{1,...,n\}\}$, on a  $|x_{i_0}|>0$.

Mais alors,

\begin{align*}\ensuremath
AX=0&\Rightarrow\forall i\in\{1,...,n\},\;\sum_{j=1}^{n}a_{i,j}x_j=0
\\
 &\Rightarrow|a_{i_0,i_0}x_{i_0}|=|-\sum_{j\neq i_0}^{}a_{i_0,j}x_j|\leq\sum_{j\neq i_0}^{}|a_{i_0,j}|.|x_j|
 \leq|x_{i_0}|\sum_{j\neq i_0}^{}|a_{i_0,j}|
\end{align*}

et, puisque $|x_{i_0}|> 0$, on obtient $|a_{i_0,i_0}\leq\sum_{j\neq i_0}^{}|a_{i_0,j}|$ contredisant les hypothèses de l'énoncé. Donc, il est absurde de supposer que $\mbox{Ker}A$ contient un vecteur non nul et $A$ est bien inversible.
\fincorrection
\correction{005274}
Soient $k$ et $l$ deux entiers tels que $1\leq k\leq n$ et $1\leq l\leq n$. Le coefficient ligne $k$, colonne $l$ de $A\overline{A}$ vaut~:

$$\sum_{j=1}^{n}\omega^{(k-1)(j-1)}\omega^{-(j-1)(l-1)}=\sum_{j=1}^{n}(\omega^{k-l})^{j-1}.$$

\begin{itemize}
\item[1er cas.]  Si $k=l$, $\omega^{k-l}=1$, et le coefficient vaut $\sum_{j=1}^{n}1=n$.

\item[2ème cas.] Si $k\neq l$. On a $-(n-1)\leq k-l\leq n-1$ avec $k-l\neq0$ et donc, $k-l$ n'est pas multiple de $n$. Par suite, $\omega^{k-l}\neq1$ et 

$$\sum_{j=1}^{n}(\omega^{k-l})^{j-1}=\frac{1-(\omega^{k-l})^n}{1-\omega}=\frac{1-1^{k-l}}{1-\omega}=0.$$
\end{itemize}

En résumé, $A\overline{A}=nI_n$. Donc $A$ est inversible à gauche et donc inversible et $A^{-1}=\frac{1}{n}\overline{A}$.
\fincorrection
\correction{005276}
Soit $f$ l'endomorphisme de $\Rr_n[X]$ qui, à un polynôme $P$ de degré inférieur ou égal à $n$, associe le polynôme $P(X+1)$.

Par la formule du binôme de \textsc{Newton}, on voit que $A$ est la matrice de $f$ dans la base canonique $(1,X,...,X^n)$ de $\Rr_n[X]$. $f$ est clairement un automorphisme de $\Rr_n[X]$, sa réciproque étant l'application qui, à un polynôme $P$ associe le polynôme $P(X-1)$.

$A$ est donc inversible et $A^{-1}=(b_{i,j})_{0\leq i,j\leq n}$ où $b_{i,j}=0$ si $i>j$ et $b_{i,j}=(-1)^{i+j}C_{j}^{i}$ si $i\leq j$.
\fincorrection
\correction{005597}
Soit $H$ un hyperplan de $\mathcal{M}_n(\Rr)$. $H$ est le noyau d'une forme linéaire non nulle $f$.

Pour $M=(m_{i,j})_{1\leqslant i,j\leqslant n}$, posons $f(M)=\sum_{1\leqslant i,j\leqslant n}^{}a_{i,j}m_{i,j}$ où les $a_{i,j}$ sont $n^2$ scalaires indépendants de $M$ et non tous nuls.

\textbf{1er cas.} Supposons qu'il existe deux indices distincts $k$ et $l$ tels que $a_{k,l}\neq0$.
Soit $M=I_n-\frac{\sum_{i=1}^{n}a_{i,i}}{a_{k,l}}E_{k,l}$. $M$ est inversible car triangulaire à coefficients diagonaux tous non nuls et $M$ est dans $H$ car $f(M)=\sum_{i=1}^{n}a_{i,i}-a_{k,l}\frac{\sum_{i=1}^{n}a_{i,i}}{a_{k,l}}=0$.

\textbf{2ème cas.} Si tous les $a_{k,l}$, $k\neq l$, sont nuls, $H$ contient la matrice inversible $\left(
\begin{array}{ccccc}
0&1&0&\ldots&0\\
\vdots&\ddots&\ddots&\ddots&\vdots\\
\vdots& & &\ddots&0\\
0& & &\ddots&1\\
1&0&\ldots&\ldots&0
\end{array}
\right)$.

\fincorrection
\correction{005604}
$A= \left(
\begin{array}{cccccc}
1&2&0&\ldots&\ldots&0\\
0&1&3&\ddots& &\vdots\\
\vdots&\ddots&\ddots&\ddots&\ddots&\vdots\\
 & & & &\ddots&0\\
\vdots& & &\ddots&\ddots&n-1\\
0&\ldots& &\ldots&0&1
\end{array}
\right)= I + N$ où $N=\left(
\begin{array}{cccccc}
0&2&0&\ldots&\ldots&0\\
0&0&3&\ddots& &\vdots\\
\vdots&\ddots&\ddots&\ddots&\ddots&\vdots\\
 & & & &\ddots&0\\
\vdots& & &\ddots&\ddots&n-1\\
0&\ldots& &\ldots&0&0
\end{array}
\right)$.

$N$ est nilpotente et donc $N^n = 0$. Par suite, 

\begin{center}
$I=I-(-N)^n=(I+N)(I -N + ... +(-N)^{n-1})$.
\end{center}

Ainsi $A$ est inversible à gauche et donc inversible, d'inverse $I -N + ... +(-N)^{n-1}$.

Calcul de $N^p$ pour $1\leqslant p\leqslant n$.

\begin{center}
$N^2 =\left(\sum_{j=2}^{n}jE_{j-1,j}\right)^2=\sum_{2\leqslant j,k\leqslant n}^{}jkE_{j-1,j}E_{k-1,k}=\sum_{j=2}^{n-1}j(j+1)E_{j-1,j}E_{j,j+1}=\sum_{j=3}^{n}j(j-1)jE_{j-2,j}$.
\end{center}

c'est-à-dire $N^2=\left(
\begin{array}{cccccc}
0&0&2\times3&0&\ldots&0\\
\vdots& &\ddots&3\times4&\ddots&\vdots\\
 & & &\ddots&\ddots&0\\
 & & & &\ddots&(n-1)n\\
\vdots& & & & &0\\
0&\ldots& &\ldots&0&0
\end{array}
\right)$.

Ensuite, $N^3=\left(\sum_{j=3}^{n}(j-1)jE_{j-2,j}\right)\left(\sum_{k=2}^{n}kE_{k-1,k}\right)) =\sum_{j=4}^{n}j(j-1)(j-2)E_{j-3,j}$.

Supposons que pour $p$ donné dans $\llbracket1,n-1\rrbracket$, $N^p=\sum_{j=p+1}^{n}j(j-1)...(j-p+1)E_{j-p,j}$.

Alors $N^{p+1}=\left(\sum_{j=p+1}^{n}j(j-1)...(j-p+1)E_{j-p,j}\right)\left(\sum_{k=2}^{n}kE_{k-1,k}\right)=\sum_{j=p+2}^{n}j(j-1)...(j-p)E_{j-p-1,j}$. Ainsi 

\begin{center}
\shadowbox{
$A^{-1}= (a_{i,j})_{1\leqslant i,j\leqslant n}$ où $a_{i,j}= 0$ si $i > j$, $1$ si $i = j$ et $(-1)^{i+j-2}\prod_{k=0}^{j-i-1}(j-k)$ sinon.
}
\end{center}
\fincorrection
\correction{005610}
On inverse $A$ en l'interprétant comme une matrice de passage.

Soit $\mathcal{B}=(e_1,...,e_n)$ la base canonique de $\Rr^n$ et $(e_1',...,e_n)$ la famille de vecteurs de $\Rr^n$ de matrice $A$ dans la base $\mathcal{B}$.

\begin{center}
$A\;\text{inversible}\Leftrightarrow(e_1',...,e_n')\;\text{base de}\;E\Leftrightarrow\text{Vect}(e_1,...,e_n)\subset\text{Vect}(e_1',...,e_n')\Leftrightarrow\forall i\in\llbracket1,n\rrbracket,\;e_i\in\text{Vect}(e_1',...,e_n')$.
\end{center}

Dans ce cas, $A^{-1}$ est la matrice de passage de $\mathcal{B}'$ à $\mathcal{B}$.

Soit $u = e_1 + ... + e_n$. Pour tout $i\in\llbracket1,n\rrbracket$, $e_i'=a_ie_i + u$ ce qui fournit $e_i =\frac{1}{a_i}(e_i'-u)$.

En additionnant membre à membre ces $n$ égalités, on obtient $u=\sum_{i=1}^{n}\frac{1}{a_i}e_i'-\left(\sum_{i=1}^{n}\frac{1}{a_i}\right)u$ et donc $\lambda u=\sum_{i=1}^{n}\frac{1}{a_i}e_i'$ où
$\lambda=1+\sum_{i=1}^{n}\frac{1}{a_i}$.

\textbf{1er cas.} Si $\lambda\neq0$, on peut exprimer $u$ en fonction des $e_i'$, $1\leqslant i\leqslant n$, et donc les $e_i$ fonction des $e_i'$. Dans ce cas $A$ est inversible. Plus précisément, $ u=\frac{1}{\lambda}\sum_{i=1}^{n}\frac{1}{a_i}e_i'$ puis, $\forall i\in\llbracket1,n\rrbracket$, $e_i =\frac{1}{a_i}\left(e_i'-\frac{1}{\lambda}\sum_{j=1}^{n}\frac{1}{a_j}e_j'\right)$ et enfin

\begin{center}
$A^{-1}=\left(
\begin{array}{ccccc}
\frac{1}{a_1}-\frac{1}{\lambda a_1^2}&-\frac{1}{\lambda a_2a_1}&\ldots&\ldots&-\frac{1}{\lambda a_na_1}\\
-\frac{1}{\lambda a_1a_2}&\frac{1}{a_2}-\frac{1}{\lambda a_2^2}& & &\vdots\\
\vdots&-\frac{1}{\lambda a_2a_3}&\ddots& &\vdots\\
\vdots&\vdots& &\frac{1}{a_{n-1}}-\frac{1}{\lambda a_{n-1}^2}&-\frac{1}{\lambda a_na_{n-1}}\\
-\frac{1}{\lambda a_1a_n}&-\frac{1}{\lambda a_2a_n}&\ldots&-\frac{1}{\lambda a_na_{n-1}}&\frac{1}{a_n}-\frac{1}{\lambda a_n^2}
\end{array}
\right)$  où $\lambda=1+\sum_{i=1}^{n}\frac{1}{a_i}$.
\end{center}

\textbf{2ème cas.} Si $\lambda=0$, on a $\sum_{i=1}^{n}\frac{1}{a_i}e_i'=0$ ce qui montre que la famille $(e_i')_{1\leqslant i\leqslant n}$ est liée et donc que $A$ n'est pas inversible.
\fincorrection
\correction{005612}
Notons $A$ la matrice de l'énoncé. Soit $f$ l'endomorphisme de $\Rr_n[X]$ de matrice $A$ dans la base canonique $\mathcal{B}$ de $\Rr_n[X]$. D'après la formule du binôme de \textsc{Newton}, $\forall k\in\llbracket0,n\rrbracket$, $f(X^k)=(X+1)^k$. $f$ coïncide donc sur la base $\mathcal{B}$ avec l'endomorphisme de $\Rr_n[X]$ qui à un polynôme $P$ associe $P(X+1)$ et $f$ est donc cet endomorphisme.

f est un automorphisme de $\Rr_n[X]$ de réciproque  l'application qui à un polynôme $P$ associe $P(X-1)$. Par suite, $A$ est inversible d'inverse la matrice de $f^{-1}$ dans la base $\mathcal{B}$.

Le coefficient ligne $i$, colonne $j$, de $A^{-1}$ vaut donc $0$ si $i > j$ et $(-1)^{i+j}\dbinom{j}{i}$ si $i\leqslant j$.
\fincorrection
\correction{005613}
Calculons $A\overline{A}$. Soit $(j,k)\in\llbracket1,n\rrbracket^2$. Le coefficient ligne $j$, colonne $k$ de $A\overline{A}$ vaut

\begin{center}
$\sum_{u=1}^{n}\omega^{(j-1)(u-1)}\omega^{-(u-1)(k-1)}=\sum_{u=1}^{n}\left(\omega^{j-k}\right)^{u-1}$.
\end{center}

\textbullet~Si $j=k$, ce coefficient vaut $n$.

\textbullet~Si $j\neq k$, puisque $j-k$ est strictement compris entre $-n$ et $n$ et que $j-k$ n'est pas nul, $\omega^{j-k}$ est différent de $1$. Le coefficient ligne $j$, colonne $k$, de $A\overline{A}$ est donc égal à $\frac{1-\left(\omega^{j-k}\right)^{n}}{1-\omega^{j-k}}=\frac{1-1}{1-\omega^{j-k}}= 0$.

Finalement, $A\overline{A}=nI_n$. Ainsi, $A$ est inversible à gauche et donc inversible, d'inverse $A^{-1}=\frac{1}{n}\overline{A}$.
\fincorrection
\correction{005617}
Montrons que $\text{Ker}A$ est réduit à $\{0\}$. Dans le cas contraire, on dispose d'un vecteur colonne non nul $X_0$ tel que $AX_0 = 0$. Posons $X_0=(x_i)_{1\leqslant i\leqslant n}$. Pour tout $i\in\llbracket1,n\rrbracket$,

\begin{center}
$\sum_{j=1}^{n}a_{i,j}x_j= 0\Rightarrow a_{i,i}x_i =-\sum_{j\neq i}^{}a_{i,j}x_j\Rightarrow |a_{i,i}||x_i|\leqslant\sum_{j\neq i}^{}|a_{i,j}||x_j|$.
\end{center}

On prend alors pour $i$ un indice $i_0$ tel que $|x_{i_0}| =\text{Max}\{|x_1|,...,|x_n|\}$. Puisque $X\neq0$, on a $|x_{i_0}|>0$. De plus,

\begin{center}
$|a_{i_0,i_0}||x_{i_0}|\leqslant\sum_{j\neq i_0}^{}|a_{i_0,j}||x_j|\leqslant\left(\sum_{j\neq i_0}^{}|a_{i_0,j}|\right)|x_{i_0}|$,
\end{center}

et puisque $|x_{i_0}|>0$, on obtient après simplification $|a_{i_0,i_0}|\leqslant\sum_{j\neq i_0}^{}|a_{i_0,j}|$ ce qui contredit les hypothèses.

Donc $\text{Ker}A=\{0\}$ et $A$ est inversible.

\fincorrection
\correction{006872}

\begin{enumerate}
  \item  si le déterminant $ad-bc$ est non nul l'inverse est 
$\frac{1}{ad-bc}\begin{pmatrix} d & -b \\ -c & a \\ \end{pmatrix}$
  \item $\frac14 \begin{pmatrix} -4 & 0 & -4 \\ 3 & 1 & 2 \\2 & -2 & 0 \\  \end{pmatrix}$

  \item si $|\alpha|\neq 1$ alors l'inverse est $\frac 1{1-\alpha\bar\alpha}
            \begin{pmatrix} 1       &-\bar\alpha        &0           \cr
                      -\alpha &1+\alpha\bar\alpha &-\bar\alpha \cr
                      0       &-\alpha            &1           \cr\end{pmatrix}$

   \item $\frac13 \begin{pmatrix}
-2 & 1 & 1 & 1 \\ 1 & -2 & 1 & 1 \\ 1 & 1 & -2 & 1 \\ 1 & 1 & 1 & -2 \\
                                      \end{pmatrix}$

  \item $\begin{pmatrix}
1 & -1 & 0 & \cdots & 0 \\
0 & 1 & -1 & 0 & \cdots\\
 & \ddots & \ddots & \ddots &  \\
\vdots & & 0 & 1 & -1  \\
0 & \cdots & \cdots & 0 & 1
\end{pmatrix}$


  \item $\begin{pmatrix}
1 & -2 & 1 & 0 &\cdots & 0 \\
 & 1 & -2 & 1 & 0 & \cdots\\
 &  & 1 & -2 & 1 & 0 \\
 & &  & \ddots & \ddots & \ddots \\
 &  & (0) &  & 1 & -2 \\
 &  &  & &  & 1
\end{pmatrix}$
\end{enumerate}
\fincorrection
\correction{007413}
\begin{enumerate}
	\item $det(D_2)=x^2-1=(x+1)(x-1)$, $det(D_3)=x^3-3x+2=(x+2)(x-1)^2$

	\item Sommer toutes les lignes à la première ligne, on a :

\[
 det(D_n)= \left | {\begin{array}{ccccc}
 x+n-1 & x+n-1 & \ldots & \ldots & x+n-1 \\
 1 & x & 1& \ldots & 1 \\
 \vdots & \ddots & \ddots &\ddots& \vdots \\
 \vdots &\ldots & \ddots & \ddots & 1 \\
 1 & \ldots & \ldots & 1 & x
 \end{array} } \right |,
\]
et comme tous les éléments sur la première ligne sont $x+n-1$, donc 
\[
 det(D_n)= (x+n-1)\left | {\begin{array}{cccc}
 1 & 1 & \ldots & 1 \\
 1 & x & \ddots & \vdots \\
 \vdots & \ddots & \ddots & 1 \\
 1 & \ldots & 1 & x
 \end{array} } \right |
\]

	\item Par la méthode du pivot de Gauss, $L_2-L_1,L_3-L_1 \ldots, L_n-L_1,$, on obtient un détermiant de la forme triangulaire 
\[
 det(D_n)= (x+n-1)\left | {\begin{array}{cccc}
 1 & 1 & \ldots & 1 \\
 0 & x-1 & \ddots & \vdots \\
 \vdots & \ddots & \ddots & 0 \\
 0 & \ldots & 0 & x-1
 \end{array} } \right |.
\]
Donc $det(D_n)= (x+n-1)(x-1)^{n-1}$.

  \item $D_n$ est inversible si et seulement si $det(D_n) \neq 0$, donc $x \neq 1-n$ et $x \neq 1$.
\end{enumerate}
\fincorrection
\nocorrection
\correction{003368}
Il n'y a pas de solution.
\fincorrection
\correction{003400}
$M = \begin{pmatrix} 1 &0 &0 &0 \cr
                        0 &1 &0 &0 \cr
		                        0 &0 &0 &0 \cr \end{pmatrix}$.
\fincorrection
\correction{003401}
$P = \begin{pmatrix} a &b  &c     &d        \cr
               0 &2a &3a+2b &4a+3b+2c \cr
               0 &0  &4a    &12a+4b   \cr
               0 &0  &0     &8a       \cr \end{pmatrix}$est inversible pour $a\ne 0$.
\fincorrection
\correction{003402}
$N = P^{-1}AP$ avec $P =
         \begin{pmatrix} 1 &1  &1  \cr
                    1 &-1 &1  \cr
                    1 &1  &-1 \cr \end{pmatrix}$.
\fincorrection
\correction{003403}
$B-I$ est inversible.
\fincorrection
\nocorrection
\correction{003405}
oui, $P = \begin{pmatrix}0&1&0&0\cr 0&0&1&0\cr 0&0&0&1\cr 1&0&0&0\cr\end{pmatrix}$.
\fincorrection
\nocorrection
\nocorrection
\correction{005259}
\begin{enumerate}
\item  $\mbox{rg}u=\mbox{rg}(u(i),u(j),u(k))=rg(u(j),u(k),u(i))$. La matrice de cette dernière famille dans la base $(i,j,k)$ est $\left(
\begin{array}{ccc}
1&0&0\\
0&1&0\\
-3&3&1
\end{array}
\right)$. Cette dernière famille est de rang $3$. Donc, $\mbox{rg}u=3$ et $u$ est bien un automorphisme de $\Rr^3$. Posons $e_1=u(i)$, $e_2=u(j)$ et $e_3=u(k)$.

$$\left\{
\begin{array}{l}
e_1=k\\
e_2=i-3k\\
e_3=j+3k
\end{array}\right.
\Leftrightarrow\left\{
\begin{array}{l}
k=e_1\\
i=3e_1+e_2\\
j=-3e_1+e_3
\end{array}\right.
\Leftrightarrow
\left\{
\begin{array}{l}
u^{-1}(k)=i\\
u^{-1}(i)=3i+j\\
u^{-1}(j)=-3i+k
\end{array}\right.
$$

et 

$$A^{-1}=\mbox{Mat}_{\mathcal{B}}(u^{-1})=\left(
\begin{array}{ccc}
3&-3&1\\
1&0&0\\
0&1&0
\end{array}
\right).$$ 

\item (Questions 2) et 3)). Posons $e_1=xi+yj+zk$ ($e_1$, $e_2$ et $e_3$ désignent d'autres vecteurs que ceux du 1)).

$$u(e_1)=e_1\Leftrightarrow(u-Id)(e_1)=0\Leftrightarrow\left(
\begin{array}{ccc}
-1&1&0\\
0&-1&1\\
1&-3&2
\end{array}
\right)\left(
\begin{array}{c}
x\\
y\\
z
\end{array}
\right)=\left(
\begin{array}{c}
0\\
0\\
0
\end{array}
\right)\Leftrightarrow
\left\{
\begin{array}{l}
-x+y=0\\
-y+z=0\\
x-3y+2z=0
\end{array}
\right.\Leftrightarrow x=y=z.$$

On prend $e_1=i+j+k$.

Posons $e_2=xi+yj+zk$.

$$u(e_2)=e_1+e_2\Leftrightarrow(u-Id)(e_2)=e_1\Leftrightarrow
\left\{
\begin{array}{l}
-x+y=1\\
-y+z=1\\
x-3y+2z=1
\end{array}
\right.\Leftrightarrow y=x+1\;\mbox{et}\;z=x+2.$$

On prend $e_2=j+2k$.

Posons $e_3=xi+yj+zk$.

$$u(e_3)=e_2+e_3\Leftrightarrow(u-Id)(e_3)=e_2\Leftrightarrow
\left\{
\begin{array}{l}
-x+y=0\\
-y+z=1\\
x-3y+2z=2
\end{array}
\right.\Leftrightarrow y=x\;\mbox{et}\;z=x+1.$$

On prend $e_3=k$.

La matrice de la famille $(e_1,e_2,e_3)$ dans la base $(i,j,k)$ est $P=
\left(
\begin{array}{ccc}
1&0&0\\
1&1&0\\
1&2&1
\end{array}
\right)$. Cette matrice est de rang $3$ et est donc inversible. Par suite $(e_1,e_2,e_3)$ est une base de $\Rr^3$. Enfin,

$$\left\{
\begin{array}{l}
e_1=i+j+k\\
e_2=j+2k\\
e_3=k
\end{array}
\right.
\Leftrightarrow
\left\{
\begin{array}{l}
k=e_3\\
j=e_2-2e_3\\
i=e_1-e_2+e_3
\end{array}
\right.
,$$

et 

$$P^{-1}=\left(
\begin{array}{ccc}
1&0&0\\
-1&1&0\\
1&-2&1
\end{array}
\right).$$

\item Voir question précédente.

\item  Soit $T$ est la matrice de $u$ dans la base $(e_1,e_2,e_3)$. $T=
\left(
\begin{array}{ccc}
1&1&0\\
0&1&1\\
0&0&1
\end{array}
\right)$. Les formules de changement de bases s'écrivent $T=P^{-1}AP$ ou encore $A=PTP^{-1}$. Par suite, pour tout relatif $n$, $A^n=PT^nP^{-1}$.

Posons $N=\left(
\begin{array}{ccc}
0&1&0\\
0&0&1\\
0&0&0
\end{array}
\right)
$. On a $N^2=\left(
\begin{array}{ccc}
0&0&1\\
0&0&0\\
0&0&0
\end{array}
\right)$  puis $N^3=0$.

Donc, pour $n$  entier naturel supérieur ou égal à $2$ donné, puisque $I$ et $N$ commutent, la formule du binôme de \textsc{Newton} fournit 

$$T^n=(I+N)^n=I+nN+\frac{n(n-1)}{2}N^2=\left(
\begin{array}{ccc}
1&n&n(n-1)/2\\
0&1&n\\
0&0&1
\end{array}
\right).$$

Cette formule reste claire pour $n=0$ et $n=1$.
Pour $n=-1$, $(I+N)(I-N+N^2)=I+N^3=I$ et donc

$$T^{-1}=(I+N)^{-1}=I-N+N^2=\left(
\begin{array}{ccc}
1&-1&1\\
0&1&-1\\
0&0&1
\end{array}
\right)=\left(
\begin{array}{ccc}
1&-1&\frac{(-1)(-1-1)}{2}\\
0&1&-1\\
0&0&1
\end{array}
\right),$$

et la formule reste vraie pour $n=-1$. Enfin, pour $n$ entier naturel non nul donné, $T^{-n}=(I+nN+\frac{n(n-1)}{2}N^2)^{-1}$ mais $(I+nN+\frac{n(n-1)}{2}N^2)(I-nN+\frac{-n(-n-1)}{2}N^2)=I$ et donc
$T^{-n}=I-nN+\frac{-n(-n-1)}{2}N^2$. Finalement, 

$$\forall n\in\Zz,\;T^n=I+nN+\frac{n(n-1)}{2}N^2=\left(
\begin{array}{ccc}
1&n&n(n-1)/2\\
0&1&n\\
0&0&1
\end{array}
\right).$$

Puis 

\begin{align*}\ensuremath
A^n&=PT^nP^{-1}=\left(
\begin{array}{ccc}
1&0&0\\
1&1&0\\
1&2&1
\end{array}
\right)\left(
\begin{array}{ccc}
1&n&n(n-1)/2\\
0&1&n\\
0&0&1
\end{array}
\right)\left(
\begin{array}{ccc}
1&0&0\\
-1&1&0\\
1&-2&1
\end{array}
\right)\\
 &=\left(
\begin{array}{ccc}
1&n&n(n-1)/2\\
1&n+1&n(n+1)/2\\
1&n+2&(n+1)(n+2)/2
\end{array}
\right)\left(
\begin{array}{ccc}
1&0&0\\
-1&1&0\\
1&-2&1
\end{array}
\right)
\\
 &=\left(
\begin{array}{ccc}
(n-1)(n-2)/2&-n(n-2)&n(n-1)/2\\
n(n-1)/2&-(n-1)(n+1)&n(n+1)/2\\
n(n+1)/2&-n(n+2)&(n+1)(n+2)/2
\end{array}
\right)
\end{align*}

ce qui fournit $u^n(i)$, $u^n(j)$ et $u^n(k)$.
\end{enumerate}
\fincorrection
\correction{005626}
Si $M(a)$ et $N(a)$ sont semblables alors nécessairement $\text{Tr}(M(a))=\text{Tr}(N(a))$. Or, pour tout scalaire $a$, $\text{Tr}(M(a))=4-3a =\text{Tr}(N(a))$. La trace ne fournit aucun renseignement.

On doit aussi avoir $\text{det}(M(a))=\text{det}(N(a))$. Or, $\text{det}(N(a))=(1-a)^2(2-a)$ et 

\begin{align*}\ensuremath
\text{det}(M(a))&= (4-a)(a^2-1-2)+6(1-a+1)+2(2-1-a) = (4-a)(a^2-3)+14-8a=-a^3+4a^2-5a+2\\
 &= (a-1)^2(2-a)=\text{det}(N(a)).
\end{align*}

Le déterminant ne fournit aucun renseignement.

Soit $f$ l'endomorphisme de $\Kk^3$ de matrice $M(a)$ dans la base canonique $\mathcal{B}_0=(i,j,k)$ de $\Kk^3$.

Le problème posé équivaut à l'existence d'une base $\mathcal{B}=(e_1,e_2,e_3)$ de $\Kk^3$ telle que 
$f(e_1)=(1-a)e_1$, $f(e_2) =(1-a)e_2 +e_1$ et $f(e_3)= (2-a)e_3$. Soit $(x,y,z)$ un élément de $\Kk^3$.

\textbullet~$f((x,y,z)) =(1-a)(x,y,z)\Leftrightarrow\left\{
\begin{array}{l}
3x+y-z=0\\
-6x-2y+2z=0\\
2x+y=0
\end{array}
\right.\Leftrightarrow\left\{
\begin{array}{l}
y=-2x\\
z=x
\end{array}
\right.$. On peut prendre $e_1=(1,-2,1)$.

\textbullet~$f((x,y,z)) =(1-a)(x,y,z)+(1,-2,1)\Leftrightarrow\left\{
\begin{array}{l}
3x+y-z=1\\
-6x-2y+2z=-2\\
2x+y=1
\end{array}
\right.\Leftrightarrow\left\{
\begin{array}{l}
y=-2x-1\\
z=x-2
\end{array}
\right.$. On peut prendre $e_2=(0,-1,-2)$.

\textbullet~$f((x,y,z)) =(2-a)(x,y,z)\Leftrightarrow\left\{
\begin{array}{l}
2x+y-z=0\\
-6x-3y+2z=0\\
2x+y-z=0
\end{array}
\right.\Leftrightarrow\left\{
\begin{array}{l}
y=-2x\\
z=0
\end{array}
\right.$. On peut prendre $e_3=(1,-2,0)$.

La matrice de la famille $\mathcal{B}=(e_1,e_2,e_3)$ dans la base $\mathcal{B}_0$ est $P=\left(
\begin{array}{ccc}
1&0&1\\
-2&-1&-2\\
1&-2&0
\end{array}
\right)$. $\text{det}P=-4+4+1 =1\neq 0$ et donc la famille $\mathcal{B}=(e_1,e_2,e_3)$ est une base de $\Kk^3$. Puisque $\text{Mat}_{\mathcal{B}_0}f=
M(a)$ et $\text{Mat}_{\mathcal{B}}f=N(a)$, les matrices $M(a)$ et $N(a)$ sont semblables.
\fincorrection
\correction{005627}
Soient $A$ et $B$ deux matrices carrées réelles de format $n$ semblables dans $\mathcal{M}_n(\Cc)$.

Il existe $P$ élément de $\mathcal{GL}_n(\Cc)$ telle que $PB =AP$ (bien plus manipulable que $B=P^{-1}AP$).

Posons $P=Q+iR$ où $Q$ et $R$ sont des matrices réelles. 
Par identification des parties réelles et imaginaires, on a $QB=AQ$ et $RB=AR$ mais cet exercice n'en est pas pour autant achevé car $Q$ ou $R$ n'ont aucune raison d'être inversibles.

On a $QB=AQ$ et $RB= AR$ et donc plus généralement pour tout réel $x$, $(Q+xR)B=A(Q+xR)$.

Maintenant, $\text{det}(Q+xR)$ est un polynôme à coefficients réels en $x$ mais n'est pas le polynôme nul car sa valeur en $i$ (tel que $i^2 = -1$) est $\text{det}P$ qui est non nul. Donc il n'existe qu'un nombre fini de réels $x$, éventuellement nul, tels que $\text{det}(Q+xR) = 0$. En particulier, il existe au moins un réel $x_0$ tel que la matrice $P_0=Q+x_0R$ soit inversible. $P_0$ est une matrice réelle inversible telle que $P_0A=BP_0$ et $A$ et $B$ sont bien semblables dans $\mathcal{M}_n(\Rr)$.

\fincorrection
\nocorrection
\nocorrection
\correction{002475}
\begin{enumerate}
  \item Soit $M$ une matrice telle que $M^2=0$ et soit $f$ l'application linéaire associée à $M$.
Comme $M^2=0$ alors $f\circ f = 0$. Cela entraîne $\Im f \subset \Ker f$. Discutons suivant la dimension
du noyau :
   \begin{enumerate}
      \item Si $\dim \Ker f=3$ alors $f=0$ donc $M=0$ (la matrice nulle).

      \item Si $\dim \Ker f=2$ alors prenons une base de $\Rr^3$ formée de deux vecteurs du noyau et d'un troisième vecteur.
Dans cette base la matrice de $f$ est
$M'=\begin{pmatrix}0&0&a\\0& 0 & b\\0&0&c\end{pmatrix}$ mais comme $f\circ f=0$ alors $M'^2=0$ ;
un petit calcul implique $c=0$. Donc $M$ et $M'$ sont les matrices de la même application linéaire $f$ mais exprimées dans des bases différentes,
donc $M$ et $M'$ sont semblables.

      \item Si $\dim \Ker f=1$ alors comme $\Im f \subset \Ker f$ on a $\dim \Im f \le 1$ mais alors cela contredit le théorème
du rang : $\dim \Ker f + \dim \Im f = \dim \Rr^3$. Ce cas n'est pas possible.

      \item Conclusion : $M$ est une matrice qui vérifie $M^2=0$ si et seulement si 
il existe une matrice inversible $P$ et des réels $a,b$ tels que 
$$M=P^{-1} \begin{pmatrix}0&0&a\\0& 0 & b\\0&0&0\end{pmatrix}P$$

   \end{enumerate}

  \item On va s'aider de l'exercice \ref{exo1093}. Si $M^2=M$ et $f$ est l'application linéaire associée 
alors $f\circ f= f$. On a vu dans l'exercice \ref{exo1093} qu'alors $\Ker f \oplus \Im f$
et que l'on peut choisir une base $(e_1,e_2,e_3)$ telle que $f(e_i)=e_i$ 
puis $f(e_i)=0$. Suivant la dimension du noyau cela donne que la matrice $M'$ de $f$ dans cette base est
$$A_0=\begin{pmatrix}0&0&0\\0&0&0\\0&0&0\end{pmatrix}\quad
A_1=\begin{pmatrix}1&0&0\\0&0&0\\0&0&0\end{pmatrix} \quad 
A_2=\begin{pmatrix}1&0&0\\0&1&0\\0&0&0\end{pmatrix} \quad 
A_3=\begin{pmatrix}1&0&0\\0&1&0\\0&0&1\end{pmatrix}.$$
Maintenant $M$ est semblable à l'une de ces matrices : il existe $P$ inversible telle que $M=P^{-1}M'P$
où $M'$ est l'une des quatre matrices $A_i$ ci-dessus.

Géométriquement notre application est une projection (projection sur une droite pour la seconde matrice 
et sur un plan pour la troisième).

  \item Posons $N= \frac{I+M}{2}$ et donc $M=2N-I$.
  Alors $M^2=I \iff (2N-I)^2=I \iff 4N^2-4N-I=I \iff N^2=N$.
  Donc par la deuxième question $N$ est semblable à l'une des matrice $A_i$ :
  $N= P^{-1}A_iP$.
  Donc $M=2P^{-1}A_iP-I = P^{-1}(2A_i-I)P$.
  Ainsi $M$ est semblable à l'une des matrices $2A_i-I$ suivantes :
$$\begin{pmatrix}-1&0&0\\0&-1&0\\0&0&-1\end{pmatrix}\quad
\begin{pmatrix}1&0&0\\0&-1&0\\0&0&-1\end{pmatrix} \quad 
\begin{pmatrix}1&0&0\\0&1&0\\0&0&-1\end{pmatrix} \quad 
\begin{pmatrix}1&0&0\\0&1&0\\0&0&1\end{pmatrix}.$$  
Ce sont des matrices de symétrie (par rapport à l'origine pour la première matrice, par rapport à une droite pour la seconde matrice 
et par rapport à un plan pour la troisième).

L'idée de poser $N= \frac{I+M}{2}$ est la suivante : si $M^2=I$ alors géométriquement l'application
linéaire $s$ associée à $M$ est une \emph{symétrie}, alors que si $N^2=N$ alors l'application
linéaire $p$ associée est une \emph{projection}. Et projection et symétrie sont liées par $p(x) =\frac{x+s(x)}{2}$ 
(faites un dessin !) c'est-à-dire $p =\frac{\text{id}+s}{2}$ ou encore $N= \frac{I+M}{2}$.


\end{enumerate}

\fincorrection
\nocorrection
\nocorrection
\nocorrection
\nocorrection
\correction{003371}
\begin{enumerate}
  \item 
  \item Pour $i<j$, on doit avoir $M(I+E_{ij}) = (I+E_{ij})M  \Rightarrow 
    \begin{cases} a_{ki} = 0 &\text{si } k \ne i \cr
            a_{jk} = 0 &\text{si }k \ne j \cr\end{cases} 
	    \Rightarrow 
    M = \begin{pmatrix} 1 &0      &\dots  &0      &*      \cr
                      &\ddots &\ddots &       &0      \cr
                      &       &\ddots &\ddots &\vdots \cr
                      &0      &       &\ddots &0      \cr
                      &       &       &       &1      \cr \end{pmatrix}$.
\end{enumerate}
 \fincorrection
\correction{003372}
$(\alpha + \mathrm{tr} A)\mathrm{tr} X = \mathrm{tr} B$.

\leavevmode\vbox{
\halign{Si $#$ : \hfil &#\hfil\cr
\alpha(\alpha + \mathrm{tr} A) \ne 0 &solution unique :
  $X = \frac 1\alpha\left(B - \frac {\mathrm{tr} B}{\alpha + \mathrm{tr} A}A\right)$. \cr
\alpha = 0                   &solutions ssi $A$ et $B$ sont proportionnelles.                               \cr
\alpha + \mathrm{tr} A = 0           &solutions ssi $\mathrm{tr} B = 0$ :
  $X = \frac 1\alpha B + \lambda A$.               \cr}}
\fincorrection
\nocorrection
\nocorrection
\nocorrection
\nocorrection
\nocorrection
\correction{003383}
2. $u_{|\Im v} = \mathrm{id}  \Rightarrow  \mathrm{tr}\big(u_{|\Im v}\big) = \mathrm{rg} v
                                 \Rightarrow  \mathrm{tr}\big(v_{|\Im v}\big) = k\mathrm{rg} v$.
 \fincorrection
\nocorrection
\correction{003385}
$M_k = A^kM_0 + S_kB$ avec
$S_k = I + A + \dots + A^{k-1} = (I-A^k)(I-A)^{-1}$ si $I-A$ est inversible.
\fincorrection
\correction{003386}
\begin{enumerate}
  \item $A^3 - (\lambda+\mu)A^2 + \lambda\mu A = 0$.
  \item $U=\frac{\mu A-A^2}{\lambda(\mu-\lambda)}$, $V=\frac{\lambda A-A^2}{\mu(\lambda-\mu)}$
    et la valeur propre est $0$, $\lambda$ ou $\mu$.
\end{enumerate}
 \fincorrection
\nocorrection
\nocorrection
\correction{003389}
\begin{enumerate}
  \item Compacité.
  \item    Si $x_1 = 0$, on pose $Y = \begin{pmatrix}\alpha \cr x_2 \cr \vdots \cr x_n \cr\end{pmatrix}$ :

    $R(Y) \ge \text{min}\Big( a_{11} + \frac {a_{12}x_2 + \dots + a_{1n}x_n}{\alpha},
    \frac {\alpha a_{21}}{x_2} + R(X_0), \dots, \frac {\alpha a_{n1}}{x_n} + R(X_0) \Big)
    > R(X_0)$ pour $\alpha > 0$ assez petit.
  \item Si $y_1 > 0$, on pose $X = X_0 + \begin{pmatrix}\alpha\cr0\cr\vdots\cr0\cr\end{pmatrix}$ :

    $AX - RX = Y + \alpha\begin{pmatrix}a_{11}-R\cr a_{21}\cr\vdots\cr a_{n1}\cr\end{pmatrix}$,
    donc pour $\alpha > 0$ assez petit, $R(X) > R$.
  \item Inégalité triangulaire.
\end{enumerate}
 \fincorrection
\correction{003390}
\begin{enumerate}
  \item 
  \item La base canonique de $E$ est $(F_{ij} = E_{ij}-E_{ji})_{1\le i<j \le n}$ où $(E_{ij})$
est la base canonique de $\mathcal{M}_n(\R)$ :
Si $M \in E$, la coordonnée de $M$ suivant $F_{ij}$ est le coefficient d'indices
$i,j$ de $M$. En particulier, en notant $A = (a_{ij})$, la coordonnée de
$f(F_{ij})$ suivant $F_{ij}$ est $a_{ii}+a_{jj}$, donc :
$$\mathrm{tr} f = \sum_{i,j} (a_{ii} + a_{jj}) = (n-1)\mathrm{tr} A.$$
\end{enumerate}
 \fincorrection
\nocorrection
\correction{003392}
Soit $\varphi$ un tel morphisme. Alors pour toute matrice $M\in GL_n(\R)$
on a $\dot 0 = p\varphi(M) = \varphi(M^p)$, donc $\varphi$ s'annule
sur toute matrice qui est une puissance $p$-ème.
Notons $P(i,j,\alpha)$ la matrice de l'opération élémentaire $L_i \leftarrow L_i + \alpha L_j$,
qui est aussi la matrice de l'opération élémentaire $C_j \leftarrow C_j + \alpha C_i$.
Toute matrice $M\in GL_n(\R)$ peut être transformée, à l'aide de ces seules
opérations élémentaires, en une matrice $M' = \mathrm{diag}(1,\dots,1,\det(M))$
par une adaptation de l'algorithme de Gauss. Comme $P(i,j,\alpha) = P(i,j,\alpha/p)^p$
et $\det(M) = \pm(|\det(M)|^{1/p})^p$, on obtient~:
$\varphi(M) = \dot 0$ si $\det(M) > 0$ et $\varphi(M) = \varphi(\mathrm{diag}(1,\dots,1,-1)) = x$
si $\det(M) < 0$. Réciproquement, la fonction $\varphi$ ainsi définie est
effectivement un morphisme de groupe si et seulement si $2x = \dot 0$, soit
$x=\dot 0$ pour $p$ impair, et $x\in\{\dot 0,\dot q\}$ pour $p=2q$.
\fincorrection
\correction{005264}
\begin{enumerate}
\item  La démonstration la plus simple apparaîtra dans le chapitre suivant~:~le déterminant d'une matrice triangulaire est le produit de ses coefficients diagonaux. Cette matrice est inversible si et seulement si son déterminant est non nul ou encore si et seulement si aucun des coefficients diagonaux n'est nul.

Pour l'instant, le plus simple est d'utiliser le rang d'une matrice. Si aucun des coefficients diagonaux n'est nul, on sait que le rang de la matrice est son format et donc que cette matrice est inversible.

Réciproquement, notons $(e_1,...,e_n)$ la base canonique de $\mathcal{M}_{n,1}(\Kk)$. Supposons que $A$ soit une matrice triangulaire inférieure dont le coefficient ligne $i$, colonne $i$, est nul. Si $i=n$, la dernière colonne de $A$ est nulle et $A$ n'est pas de rang $n$ et donc n'est pas inversible. Si $i<n$, alors les $n-i+1$ dernières colonnes sont dans $\mbox{Vect}(e_{i+1},...,e_n)$ qui est de dimension au plus $n-i(<n-i+1)$, et encore une fois, la famille des colonnes de $A$ est liée.

\item  Soit $A=(a_{i,j})_{1\leq i,j\leq n}$ une matrice triangulaire supérieure et $f$ l'endomorphisme de $\Kk^n$ de matrice $A$ dans la base canonique $\mathcal{B}=(e_1,...,e_n)$ de $\Kk^n$. Soit $\mathcal{B'}=(e_n,...,e_1)$. $\mathcal{B'}$ est encore une base de $\Kk^n$. Soit alors $P$ la matrice de passage de $\mathcal{B}$ à $\mathcal{B}'$ puis $A'$ la matrice de $f$ dans la base $\mathcal{B}'$. Les formules de changement de bases permettent d'affirmer que $A'=P^{-1}AP$ et donc que $A$ et $A'$ sont semblables.

Vérifions alors que $A'$ est une matrice triangulaire inférieure. Pour $i\in\{1,...,n\}$, posons $e_i'=e_{n+1-i}$. $A$ est triangulaire supérieure. Donc, pour tout $i$, $f(e_i)\in\mbox{Vect}(e_1,...,e_i)$. Mais alors, pour tout $i\in\{1,...,n\}$, $f(e_{n+1-i}')\in\mbox{Vect}(e_n',...,e_{n+1-i}')$ ou encore, pour tout $i\in\{1,...,n\}$, $f(e_{i}')\in\mbox{Vect}(e_n',...,e_{i}')$. Ceci montre que $A'$ est une matrice triangulaire inférieure.
\end{enumerate}
\fincorrection
\correction{005265}
\begin{enumerate}
\item  $E=\mbox{Vect}(I,J)$. Donc, $E$ est un sous-espace vectoriel de $\mathcal{M}_2(\Rr)$. La famille $(I,J)$ est clairement libre et donc est une base de $E$. Par suite, $\mbox{dim}E=2$.
\item  $J^2=\left(
\begin{array}{cc}
1&1\\
0&1
\end{array}
\right)
\left(
\begin{array}{cc}
1&1\\
0&1
\end{array}
\right)=\left(
\begin{array}{cc}
1&2\\
0&1
\end{array}
\right)=2J-I$. Plus généralement, pour $(x,y,x',y')\in\Rr^4$,

$$M(x,y)M(x',y')=(xI+yJ)(x'I+y'J)=xx'I+(xy'+yx')J+yy'J^2=(xx'-yy')I+(xy'+yx'+2yy')J\;(*).$$

Montrons alors que $(E,+,\times)$ est un sous-anneau de $(\mathcal{M}_2(\Rr),+,\times)$.

$E$ contient $I=1.I+0.J$. $(E,+)$ est un sous-groupe de $(\mathcal{M}_2(\Rr),+)$ et, d'après $(*)$, $E$ est stable pour $\times$. Donc, $(E,+,\times)$ est un sous-anneau de $(\mathcal{M}_2(\Rr),+,\times)$.

\item  Soit $((x,y),(x',y'))\in(\Rr^2)^2$.

$$M(x,y)M(x',y')=I\Leftrightarrow(xx'-yy')I+(xy'+yx'+2yy')J=I\Leftrightarrow
\left\{
\begin{array}{l}
xx'-yy'=1\\
yx'+(x+2y)y'=0
\end{array}
\right..$$

Le déterminant de ce dernier système d'inconnues $x'$ et $y'$ vaut $x(x+2y)+y^2=x^2+2xy+y^2=(x+y)^2$. Si $y\neq-x$, ce système admet un et seule couple solution. Par suite, si $y\neq -x$, il existe $(x',y')\in\Rr^2$ tel que $M(x,y)M(x',y')=I$. Dans ce cas, la matrice $M(x,y)$ est inversible dans $E$.

Si $y=-x$, le système s'écrit $\left\{
\begin{array}{l}
x(x'+y')=1\\
-x(x'+y')=0
\end{array}
\right.$ et n'a clairement pas de solution.
\item 
\begin{enumerate}
\item Soit $(x,y)\in\Rr^2$.

$$M(x,y)^2=I\Leftrightarrow\left\{
\begin{array}{l}
x^2-y^2=1\\
2y(x+y)=0
\end{array}
\right.\Leftrightarrow\left\{
\begin{array}{l}
y=0\\
x^2=1
\end{array}
\right.
\;\mbox{ou}
\left\{
\begin{array}{l}
x^2-y^2=1\\
x+y=0
\end{array}
\right.
\Leftrightarrow\left\{
\begin{array}{l}
y=0\\
x=1
\end{array}
\right.
\;\mbox{ou}\;\left\{
\begin{array}{l}
y=0\\
x=-1
\end{array}
\right..$$

Dans $E$, l'équation $X^2=I$ admet exactement deux solutions à savoir $I$ et $-I$.

\item Soit $(x,y)\in\Rr^2$.

$$M(x,y)^2=0\Leftrightarrow\left\{
\begin{array}{l}
x^2-y^2=0\\
2y(x+y)=0
\end{array}
\right.\Leftrightarrow\left\{
\begin{array}{l}
y=0\\
x^2=0
\end{array}
\right.
\;\mbox{ou}
\left\{
\begin{array}{l}
y=-x\\
0=0
\end{array}
\right.
\Leftrightarrow y=-x.$$

Dans $E$, l'équation $X^2=0$ admet pour solutions les matrices de la forme $\lambda(J-I)=\left(
\begin{array}{cc}
0&\lambda\\
0&0
\end{array}
\right)$, $\lambda\in\Rr$.

\item Soit $(x,y)\in\Rr^2$.

\begin{align*}\ensuremath
M(x,y)^2=M(x,y)&\Leftrightarrow\left\{
\begin{array}{l}
x^2-y^2=x\\
2y(x+y)=y
\end{array}
\right.
\Leftrightarrow\left\{
\begin{array}{l}
x^2-y^2=x\\
y(2x+2y-1)=0
\end{array}
\right.
\\
 &
\Leftrightarrow\left\{
\begin{array}{l}
y=0\\
x^2=x
\end{array}
\right.
\;\mbox{ou}\;
\left\{
\begin{array}{l}
y=-x+\frac{1}{2}\\
x^2-(-x+\frac{1}{2})^2=x
\end{array}
\right.
\\
 &\Leftrightarrow\left\{
\begin{array}{l}
y=0\\
x=0
\end{array}
\right.
\;\mbox{ou}\;\left\{
\begin{array}{l}
y=0\\
x=1
\end{array}
\right.
\;\mbox{ou}\;\left\{
\begin{array}{l}
\frac{1}{4}=0\\
y=-x+\frac{1}{2}
\end{array}
\right.\Leftrightarrow\left\{
\begin{array}{l}
y=0\\
x=0
\end{array}
\right.
\;\mbox{ou}\;\left\{
\begin{array}{l}
y=0\\
x=1
\end{array}
\right.
.
\end{align*}

Dans $E$, l'équation $X^2=X$ admet exactement deux solutions à savoir $0$ et $I$.
\end{enumerate}

\end{enumerate}
\fincorrection
\correction{005266}
Soit $(i,j)$ la base canonique de $\Rr^2$ et $(e_1,e_2,e_3)$ la base canonique de $\Rr^3$. On cherche $f\in\mathcal{L}(\Rr^2,\Rr^3)$ et $g\in\mathcal{L}(\Rr^3,\Rr^2)$ tels que

$$f\circ g(e_1)=-e_2+e_3,\;f\circ g(e_2)=-e_1+e_3\;\mbox{et}\;f\circ g(e_3)=-e_1-e_2+2e_3(=f\circ g(e_1+e_2)).$$

On pose $g(e_1)=i$, $g(e_2)=j$ et $g(e_3)=i+j$, puis $f(i)=-e_2+e_3$ et $f(j)=-e_1+e_3$. Les applications linéaires $f$ et $g$ conviennent, ou encore si on pose

$$A=\left(
\begin{array}{cc}
0&-1\\
-1&0\\
1&1
\end{array}
\right)\;\mbox{et}\;B=\left(
\begin{array}{ccc}
1&0&1\\
0&1&1
\end{array}
\right),$$

alors $AB=\left(
\begin{array}{cc}
0&-1\\
-1&0\\
1&1
\end{array}
\right)\left(
\begin{array}{ccc}
1&0&1\\
0&1&1
\end{array}
\right)=\left(
\begin{array}{ccc}
0&-1&-1\\
-1&0&-1\\
1&1&2
\end{array}
\right)$.

$A$ et $B$ désignent maintenant deux matrices quelconques, éléments de $\mathcal{M}_{3,2}(\Rr)$ et $\mathcal{M}_{2,3}(\Rr)$ respectivement, telles que $AB=\left(
\begin{array}{ccc}
0&-1&-1\\
-1&0&-1\\
1&1&2
\end{array}
\right)$. Calculons $(AB)^2$. On obtient

$$(AB)^2=\left(
\begin{array}{ccc}
0&-1&-1\\
-1&0&-1\\
1&1&2
\end{array}
\right)\left(
\begin{array}{ccc}
0&-1&-1\\
-1&0&-1\\
1&1&2
\end{array}
\right)=\left(
\begin{array}{ccc}
0&-1&-1\\
-1&0&-1\\
1&1&2
\end{array}
\right)=AB.$$

Mais alors, en multipliant les deux membres de cette égalité par $B$ à gauche et $A$ à droite, on obtient

$$(BA)^3=(BA)^2\;(*).$$

Notons alors que 

$$\mbox{rg}(BA)\geq\mbox{rg}(ABAB)=\mbox{rg}((AB)^2)=\mbox{rg}(AB)=2,$$

et donc, $BA$ étant une matrice carrée de format $2$, $\mbox{rg}(BA)=2$. $BA$ est donc une matrice inversible. Par suite, on peut simplifier les deux membres de l'égalité $(*)$ par $(BA)^2$ et on obtient $BA=I_2$.
\fincorrection
\correction{005268}
Soit $A=(a_{k,l})_{1\leq k,l\leq n}\in\mathcal{M}_n(\Kk)$.

Si $A$ commute avec toute matrice, en particulier~:~$\forall(i,j)\in\{1,...,n\}^2,\;AE_{i,j}=E_{i,j}A$. Maintenant,

$$AE_{i,j}=\sum_{k,l}^{}a_{k,l}E_{k,l}E_{i,j}=\sum_{k=1}^{n}a_{k,i}E_{k,j}\;\mbox{et}\;E_{i,j}A=\sum_{k,l}^{}a_{k,l}E_{i,j}E_{k,l}=\sum_{l=1}^{n}a_{j,l}E_{i,l}.$$

On note que si $k\neq i$ ou $l\neq j$, $E_{k,j}\neq E_{i,l}$. Puisque la famille $(E_{i,j})$ est libre, on peut identifier les coefficients et on obtient~:~si $k\neq i$, $a_{k,i}=0$.
D'autre part, le coefficient de $E_{i,j}$ est $a_{i,i}$ dans la première somme et $a_{j,j}$ dans la deuxième. Ces coefficients doivent être égaux.

Finalement, si $A$ commute avec toute matrice, ses coefficients non diagonaux sont nuls et ses coefficients diagonaux sont égaux. Par suite, il existe un scalaire $\lambda\in\Kk$ tel que $A=\lambda I_n$. Réciproquement, si $A$ est une matrice scalaire, $A$ commute avec toute matrice.

\fincorrection
\correction{005270}
Soit $H$ un hyperplan de $\mathcal{M}_n(\Kk)$ et $f$ une forme linéaire non nulle sur $\mathcal{M}_n(\Kk)$ telle que $H=\mbox{Ker}f$.

Pour $A=(a_{i,j})_{1\leq i,j\leq n}$, posons $f(A)=\sum_{1\leq i,j\leq n}^{}\alpha_{i,j}a_{i,j}$.

\begin{itemize}
\item[1er cas.] Supposons $\exists(i,j)\in\{1,...,n\}^2/\;i\neq j\;\mbox{et}\;\alpha_{i,j}\neq 0$. On pose alors $S=\sum_{k=1}^{n}\alpha_{k,k}$ et on considère $A=\sum_{k=1}^{n}E_{k,k}-\frac{S}{\alpha_{i,j}}E_{i,j}$. $A$ est triangulaire à coefficients diagonaux tous non nuls et est donc inversible.

De plus, $f(A)=\sum_{k=1}^{n}\alpha_{k,k}-\frac{S}{\alpha_{i,j}}\alpha_{i,j}=S-S=0$ et $A$ est élément de $H$.

\item[2ème cas.] Supposons $\forall(i,j)\in\{1,...,n\}^2,\;(i\neq j\Rightarrow \alpha_{i,j}=0)$. Alors, $\forall A\in\mathcal{M}_n(\Kk),\;f(A)=\sum_{i=1}^{n}\alpha_{i,i}a_{i,i}$. Soit $A=E_{n,1}+E_{2,1}+E_{3,2}+...+E_{n-1,n}$. $A$ est inversible car par exemple égale à la matrice de passage de la base canonique $(e_1,e_2,...,e_n)$ de $\Kk^n$ à la base $(e_n,e_1,...,e_{n-1})$. De plus, $f(A)=0$.
\end{itemize}
\fincorrection
\nocorrection
\correction{005273}
\begin{enumerate}
\item  Soit $(i,j)\in\{1,...,p\}\times\{1,...,r\}$. Le coefficient ligne $i$, colonne $j$, de la matrice $M+N$ est la somme du coefficient ligne $i$, colonne $j$, de la matrice $M$ et du coefficient ligne $i$, colonne $j$, de la matrice $N$ ou encore la somme du coefficient ligne $i$, colonne $j$, de la matrice $A$ et du coefficient ligne $i$, colonne $j$, de la matrice $A'$. On a des résultats analogues pour les autres valeurs du couple $(i,j)$ et donc

$$M+N=\left(
\begin{array}{cc}
A+A'&B+B'\\
C+C'&D+D'
\end{array}
\right).$$

\item  Posons $M=\left(
\begin{array}{cc}
A&B\\
C&D
\end{array}
\right)$ et $N=\left(
\begin{array}{cc}
A'&B'\\
C'&D'
\end{array}
\right)$ où $A\in\mathcal{M}_{p,r}(\Kk)$, $B\in\mathcal{M}_{q,r}(\Kk)$, $C\in\mathcal{M}_{p,s}(\Kk)$, $D\in\mathcal{M}_{q,s}(\Kk)$, puis $A'\in\mathcal{M}_{t,p}(\Kk)$, $B'\in\mathcal{M}_{u,p}(\Kk)$, $C'\in\mathcal{M}_{t,q}(\Kk)$, $D'\in\mathcal{M}_{u,q}(\Kk)$ (le découpage de $M$ en colonne est le même que le découpage de $N$ en lignes).

Soit alors $(i,j)\in\{1,...,r\}\times\{1,...,t\}$. Le coefficient ligne $i$, colonne $j$ de la matrice $MN$ vaut 

$$\sum_{k=1}^{p+q}m_{i,k}n_{k,j}=\sum_{k=1}^{p}m_{i,k}n_{k,j}+\sum_{k=p+1}^{p+q}m_{i,k}n_{k,j}.$$

Mais, $\sum_{k=1}^{p}m_{i,k}n_{k,j}$ est le coefficient ligne $i$, colonne $j$ du produit $AA'$ et $\sum_{k=p+1}^{p+q}m_{i,k}n_{k,j}$ est le coefficient ligne $i$, colonne $j$ du produit $BC'$. Finalement, $\sum_{k=1}^{p+q}m_{i,k}n_{k,j}$ est le coefficient ligne $i$, colonne $j$ du produit $AA'+BC'$. On a des résultats analogues pour les autres valeurs du couple $(i,j)$ et donc

$$MN=\left(
\begin{array}{cc}
AA'+BC'&AB'+BD'\\
CA'+DC'&CB'+DD'
\end{array}
\right).$$
\end{enumerate}
\fincorrection
\correction{005275}
\begin{enumerate}
\item  Un vecteur non nul $x$ est colinéaire à son image si et seulement si il existe $\lambda\in\Cc$ tel que $u(x)=\lambda x$. Les nombres $\lambda$ correspondants sont les complexes tels qu'il existe un vecteur $x\neq0$ dans $\mbox{Ker}(u-\lambda Id)$ ou encore tels que $A-\lambda I_4\notin\mathcal{GL}_4(\Cc)$.

Le déterminant de $A-\lambda I_4$ vaut :

\begin{align*}\ensuremath
\left|
\begin{array}{cccc}
7-\lambda&4&0&0\\
-12&-7-\lambda&0&0\\
20&11&-6-\lambda&-12\\
-12&-6&6&11-\lambda
\end{array}
\right|&=(7-\lambda)\left|
\begin{array}{ccc}
-7-\lambda&0&0\\
11&-6-\lambda&-12\\
-6&6&11-\lambda
\end{array}
\right|-4\left|
\begin{array}{ccc}
-12&0&0\\
20&-6-\lambda&-12\\
-12&6&11-\lambda
\end{array}
\right|\\
 &=(7-\lambda)(-7-\lambda)\left|
\begin{array}{cc}
-6-\lambda&-12\\
6&11-\lambda
\end{array}
\right|-4(-12)\left|
\begin{array}{cc}
-6-\lambda&-12\\
6&11-\lambda
\end{array}
\right|\\
 &=(\lambda-7)(\lambda+7)(\lambda^2-5\lambda+6)+48(\lambda^2-5\lambda+6)\\
 &=(\lambda^2-5\lambda+6)(\lambda^2-49+48)=(\lambda-2)(\lambda-3)(\lambda-1)(\lambda+1)
\end{align*}

Ainsi, $A-\lambda I_4\notin\mathcal{GL}_4(\Cc)\Leftrightarrow\lambda\in\{-1,1,2,3\}$.

\begin{itemize}
\item[- Cas $\lambda=-1$.] Soit $(x,y,z,t)\in\Cc^4$.

\begin{align*}\ensuremath
(x,y,z,t)\in\mbox{Ker}(u+Id)&\Leftrightarrow\left\{
\begin{array}{l}
8x+4y=0\\
-12x-6y=0\\
20x+11y-5z-12t=0\\
-12x-6y+6z+12t=0
\end{array}
\right.\Leftrightarrow\left\{
\begin{array}{l}
y=-2x\\
-2x-5z-12t=0\\
z+2t=0\end{array}
\right.
\\
 &\Leftrightarrow\left\{
\begin{array}{l}
y=-2x\\
z=-2t\\
-2x-2t=0
\end{array}
\right.
\Leftrightarrow
\left\{
\begin{array}{l}
y=-2x\\
t=-x\\
z=2x
\end{array}
\right..
\end{align*}

Donc, $\mbox{Ker}(u+Id)=\mbox{Vect}(e_1)$ où $e_1=(1,-2,2,-1)$.

\item[- Cas $\lambda=1$.]
Soit $(x,y,z,t)\in\Cc^4$.

\begin{align*}\ensuremath
(x,y,z,t)\in\mbox{Ker}(u-Id)&\Leftrightarrow\left\{
\begin{array}{l}
6x+4y=0\\
-12x-8y=0\\
20x+11y-7z-12t=0\\
-12x-6y+6z+10t=0
\end{array}
\right.\Leftrightarrow\left\{
\begin{array}{l}
3x+2y=0\\
20x+11y-7z-12t=0\\
-6x-3y+3z+5t=0
\end{array}
\right.
\\
 &\Leftrightarrow\left\{
\begin{array}{l}
y=-\frac{3}{2}x\\
14z+24t=7x\\
6z+10t=3x
\end{array}
\right.
\Leftrightarrow
\left\{
\begin{array}{l}
y=-\frac{3}{2}x\\
z=\frac{1}{2}x\\
t=0
\end{array}
\right..
\end{align*}

Donc, $\mbox{Ker}(u-Id)=\mbox{Vect}(e_2)$ où $e_2=(2,-3,1,0)$.

\item[- Cas $\lambda=2$.]

Soit $(x,y,z,t)\in\Cc^4$.

\begin{align*}\ensuremath
(x,y,z,t)\in\mbox{Ker}(u-Id)&\Leftrightarrow\left\{
\begin{array}{l}
5x+4y=0\\
-12x-9y=0\\
20x+11y-8z-12t=0\\
-12x-6y+6z+9t=0
\end{array}
\right.\Leftrightarrow\left\{
\begin{array}{l}
x=0\\
y=0
2z+3t=0
\end{array}
\right.
\\
 &\Leftrightarrow\left\{
\begin{array}{l}
x=y=0\\
z=-\frac{3}{2}t
\end{array}
\right.
.
\end{align*}

Donc, $\mbox{Ker}(u-2Id)=\mbox{Vect}(e_3)$ où $e_3=(0,0,3,-2)$.

\item[-Cas $\lambda=3$.]

Soit $(x,y,z,t)\in\Cc^4$.

\begin{align*}\ensuremath
(x,y,z,t)\in\mbox{Ker}(u-Id)&\Leftrightarrow\left\{
\begin{array}{l}
4x+4y=0\\
-12x-10y=0\\
20x+11y-9z-12t=0\\
-12x-6y+6z+8t=0
\end{array}
\right.\Leftrightarrow\left\{
\begin{array}{l}
x=0\\
y=0
3z+4t=0
\end{array}
\right.
\\
 &\Leftrightarrow\left\{
\begin{array}{l}
x=y=0\\
z=-\frac{4}{3}t
\end{array}
\right.
.
\end{align*}

Donc, $\mbox{Ker}(u-3Id)=\mbox{Vect}(e_4)$ où $e_4=(0,0,4,-3)$.

\end{itemize}

Soit $P$ la matrice de la famille $(e_1,e_2,e_3,e_4)$ dans la base canonique $(i,j,k,l)$. On a $P=\left(
\begin{array}{cccc}
1&2&0&0\\
-2&-3&0&0\\
2&1&3&4\\
-1&0&-2&-3
\end{array}
\right)$.

Montrons que $P$ est inversible et déterminons son inverse.

\begin{align*}\ensuremath
\left\{
\begin{array}{l}
e_1=i-2j+2k-l\\
e_2=2i-3j+k\\
e_3=3k-2l\\
e_4=4k-3l
\end{array}
\right.
&\Leftrightarrow
\left\{
\begin{array}{l}
k=3e_3-2e_4\\
l=4e_3-3e_4\\
e_1=i-2j+2(3e_3-2e_4)-(4e_3-3e_4)\\
e_2=2i-3j+(3e_3-2e_4)
\end{array}
\right.
\\
 &\Leftrightarrow
\left\{
\begin{array}{l}
k=3e_3-2e_4\\
l=4e_3-3e_4\\
i-2j=e_1-2e_3+e_4\\
2i-3j=e_2-3e_3+2e_4
\end{array}
\right.\Leftrightarrow
\left\{
\begin{array}{l}
k=3e_3-2e_4\\
l=4e_3-3e_4\\
i=-3e_1+2e_2+e_4\\
j=-2e_1+e_2+e_3
\end{array}
\right.
\end{align*}
Ainsi, $\Cc^4=\mbox{Vect}(i,j,k,l)\subset\mbox{Vect}(e_1,e_2,e_3,e_4)$. Donc, la famille $(e_1,e_2,e_3,e_4)$ est génératrice de $\Cc^4$ et donc une base de $\Cc^4$. Ainsi, $P$ est inversible et 

$$P^{-1}=
\left(
\begin{array}{cccc}
-3&-2&0&0\\
2&1&0&0\\
0&1&3&4\\
1&0&-2&-3
\end{array}
\right)
.$$

\item  Les formules de changement de bases s'écrivent $A=PDP^{-1}$ avec $D=\mbox{diag}(-1,1,2,3)$.

\item  Soit $n\in\Nn^*$. Calculons $A^n$.

\begin{align*}\ensuremath
A^n&=PD^nP^{-1}=\left(
\begin{array}{cccc}
1&2&0&0\\
-2&-3&0&0\\
2&1&3&4\\
-1&0&-2&-3
\end{array}
\right)\left(
\begin{array}{cccc}
(-1)^n&0&0&0\\
0&1&0&0\\
0&0&2^n&0\\
0&0&0&3^n
\end{array}
\right)\left(
\begin{array}{cccc}
-3&-2&0&0\\
2&1&0&0\\
0&1&3&4\\
1&0&-2&-3
\end{array}
\right)
\\
 &=\left(
\begin{array}{cccc}
1&2&0&0\\
-2&-3&0&0\\
2&1&3&4\\
-1&0&-2&-3
\end{array}
\right)\left(
\begin{array}{cccc}
-3(-1)^n&-2(-1)^n&0&0\\
2&1&0&0\\
0&2^n&3.2^n&4.2^n\\
3^n&0&-2.3^n&-3.3^n
\end{array}
\right)\\
 &=\left(
\begin{array}{cccc}
-3(-1)^n+4&-2(-1)^n+2&0&0\\
6(-1)^n-6&4(-1)^n-3&0&0\\
-6(-1)^n+2+4.3^n&-4(-1)^n+1+3.2^n&9.2^n-8.3^n&12(2^n-3^n)\\
3((-1)^n-3^n)&2((-1)^n-2^n)&6(3^n-2^n)&-8.2^n+9.3^n
\end{array}
\right)
\end{align*}
\end{enumerate}
\fincorrection
\correction{005585}
Cherchons une matrice $A$ de format $(3,2)$ et une matrice $B$ de format $(2,3)$ telles que $AB=\left(
\begin{array}{ccc}
8&2&-2\\
2&5&4\\
-2&4&5
\end{array}
\right)$.

Posons $E=\Rr^2$ et notons $(i,j)$ la base canonique de $E$.

Posons $F=\Rr^3$ et notons $(e_1,e_2,e_3)$ la base canonique de $F$.

Le problème posé matriciellement peut aussi s'énoncer en termes d'applications linéaires :
trouvons $f\in\mathcal{L}(E,F)$ et $g\in\mathcal{L}(F,E)$ telles que $f\circ g(e_1)=8e_1+2e_2-2e_3$, $f\circ g(e_2)= 2e_1+5e_2+4e_3$ et
$f\circ g(e_3) = -2e_1+4e_2+5e_3$.

Remarquons tout d'abord que le problème posé n'a pas nécessairement de solution car par exemple $\text{rg}(f\circ g)\leqslant\text{Min}\{f,g\}\leqslant\text{dim}E=2$ et si la matrice proposée est de rang $3$ (c'est à dire inversible), le problème posé n'a pas de solution.

Ici, $\left|
\begin{array}{ccc}
8&2&-2\\
2&5&4\\
-2&4&5
\end{array}
\right|=8\times9-2\times18-2\times18=0$ et la matrice proposée est de rang au plus $2$ puis de rang $2$ car ses deux premières colonnes ne sont pas colinéaires.

Une relation de dépendance des colonnes est $C_1=2C_2-2C_3$.

Un couple $(f,g)$ solution devra vérifier $f\circ g(e_1)=2f\circ g(e_2)-2f\circ g(e_3)$.

Prenons n'importe quoi ou presque pour $g(e_2)$ et $g(e_3)$ mais ensuite prenons $g(e_1)= 2g(e_2) - 2g(e_3)$.

Par exemple, posons $g(e_2) = i$, $g(e_3) =j$ et $g(e_1) = 2i-2j$ puis $f(i) = 2e_1+5e_2+4e_3$ et $f(j)=-2e_1+4e_2+5e_3$ ou encore soient $A=\left(
\begin{array}{cc}
2&-2\\
5&4\\
4&5
\end{array}
\right)$ et $B=\left(
\begin{array}{ccc}
2&1&0\\
-2&0&1
\end{array}
\right)$. On a $AB=\left(
\begin{array}{ccc}
8&2&-2\\
2&5&4\\
-2&4&5
\end{array}
\right)$.

Soient $A$ et $B$ deux matrices de formats respectifs $(3,2)$ et $(2,3)$ telles que $AB=\left(
\begin{array}{ccc}
8&2&-2\\
2&5&4\\
-2&4&5
\end{array}
\right)$. Calculons $BA$ (il n'y a bien sûr pas unicité de $A$ et $B$, mais l'énoncé suggère que le produit $BA$ doit être indépendant de $A$ et $B$).

Tout d'abord 

\begin{center}
$(AB)^2=\left(
\begin{array}{ccc}
8&2&-2\\
2&5&4\\
-2&4&5
\end{array}
\right)\left(
\begin{array}{ccc}
8&2&-2\\
2&5&4\\
-2&4&5
\end{array}
\right)=\left(
\begin{array}{ccc}
72&18&-18\\
18&45&36\\
-18&36&45
\end{array}
\right)=9AB$.
\end{center}

De plus, $\text{rg}(BA)\geqslant\text{rg}(A(BA)B)=\text{rg}((AB)^2)=\text{rg}(9AB) =\text{rg}(AB)=2$ et donc $\text{rg}(BA)=2$ puis $BA\in\mathcal{GL}_2(\Rr)$.

De l'égalité $(AB)^2=9AB$, on tire après multiplication à gauche par $B$ et à droite par $A$, $(BA)^3=9(BA)^2$ et, puisque $BA$ est une matrice carrée inversible et donc simplifiable pour la multiplication des matrices, $BA=9I_2$.

\begin{center}
\shadowbox{
$BA=9I_2$.
}
\end{center}
\fincorrection
\correction{005594}
Soit $A=\sum_{M\in G}^{}M$. Alors $A^2=\sum_{(M,N)\in G^2}^{}MN$.

Soit $M\in G$ fixée. Considérons l'application $\varphi$ de $G$ dans $G$ qui à un élément $N$ de $G$ associe $MN$. Puisque $G$ est stable pour le produit, $\varphi$ est bien une application. Plus précisément, $\varphi$ est une permutation de $G$ car l'application $\psi$ de $G$ dans lui-même  qui à un élément $N$ de $G$ associe $M^{-1}N$ vérifie $\psi\circ\varphi=\varphi\circ\psi=Id_G$. On en déduit que

\begin{center}
$A^2=\sum_{M\in G}^{}\left(\sum_{N\in G}^{}MN\right)=\sum_{M\in G}^{}A=pA$ où $p=\text{card}(G)$.
\end{center}

Finalement, la matrice $P=\frac{1}{p}A$ est idempotente car $\left(\frac{1}{p}A\right)^2=\frac{1}{p^2}pA=\frac{1}{p}A$. Comme $A$ est une matrice de projection, on sait que $\text{rg}P=\text{Tr}P =\sum_{M\in G}^{}\text{Tr}M=0$ et donc $P=0$ ou encore $\sum_{M\in G}^{}M= 0$.
\fincorrection
\correction{005595}
Par la même méthode qu'au \ref{ex:rou32}, on voit que $f=\frac{1}{p}\sum_{g\in G}^{}g$ est un projecteur et donc $\frac{1}{p}\sum_{g\in G}^{}\text{Tr}g=\text{rg}f$. Maintenant, si $x$ est un élément de $F$ alors pour tout $g$ dans $G$, $g(x)=x$ et donc $f(x) = x$. Ainsi, un élément $x$ de $F$ est dans $\text{Im}f$. 

Inversement, soit $x$ un élément de $\text{Im}f$. Pour $g\in G$, 

\begin{center}
$g(x)=g(f(x))=\frac{1}{p}\sum_{h\in G}^{}g\circ h(x)=\frac{1}{p}\sum_{h\in G}^{}h(x)=f(x)= x$.
\end{center}

(Comme au \ref{ex:rou32}, l'application qui, pour $g\in G$ 
fixé, associe à un élément $h$ de $G$ l'élément $g\circ h$, est une permutation de $G$).

Ainsi, l'élément $x$ de $\text{Im}f$ est dans $F$. On a montré que $F=\text{Im}f$. Puisque $f$ est un projecteur, on en déduit que

\begin{center}
$\text{dim}F=\text{rg}f=\text{Tr}f=\frac{1}{p}\sum_{g\in G}^{}\text{Tr}g$. 
\end{center}
\fincorrection
\correction{005596}
Comme à l'exercice \ref{ex:rou32}, la matrice $A =\frac{1}{p}\sum_{k=1}^{p}A_k$ est idempotente et donc $\text{Tr}A=\text{rg}A$ d'après le \ref{ex:rou28}. Par suite, $\text{Tr}(A_1)+...+\text{Tr}A_p =p\text{rg}A$ est un entier divisible par $p$.
\fincorrection
\correction{005605}
On note $\mathcal{B}=(E_{i,j})_{1\leqslant i,j\leqslant n}$ la base canonique de $\mathcal{M}_n(\Kk)$.

$\text{Tr}f=\sum_{1\leqslant i,j\leqslant n}^{}\alpha_{i,j}$ où $\alpha_{i,j}$ désigne la $(i,j)$-ème coordonnée de $f(E_{i,j})=AE_{i,j}+E_{i,j}A$ dans la base $\mathcal{B}$.

Mais pour $(i,j)\in\llbracket1,n\rrbracket^2$ donné, 

\begin{center}
$AE_{i,j}=\sum_{1\leqslant k,l\leqslant n}^{}a_{k,l}E_{k,l}E_{i,j}=\sum_{k=1}^{n}a_{k,i}E_{k,j}$
\end{center}

et de même,

\begin{center}
$E_{i,j}A=\sum_{1\leqslant k,l\leqslant n}^{}a_{k,l}E_{i,j}E_{k,l}=\sum_{l=1}^{n}a_{j,l}E_{i,l}$.
\end{center}

Donc $\forall(i,j)\in\llbracket1,n\rrbracket^2$, $\alpha_{i,j}=a_{i,i}+a_{j,j}$ puis

\begin{center}
$\text{Tr}f=\sum_{1\leqslant i,j\leqslant n}^{}(a_{i,i}+a_{j,j})=2\sum_{1\leqslant i,j\leqslant n}^{}a_{i,i}=2\sum_{j=1}^{n}\left(\sum_{i=1}^{n}a_{i,i}\right)=2\sum_{j=1}^{n}\text{Tr}A=2n\text{Tr}A$.
\end{center}

\begin{center}
\shadowbox{
$\text{Tr}f=2n\text{Tr}A$.
}
\end{center}
\fincorrection
\correction{005606}
Si $M$ est solution, nécessairement $a\text{Tr}M +(\text{Tr}M)(\text{Tr}A)=\text{Tr}B$ ou encore $(\text{Tr}M)(a+\text{Tr}A) =\text{Tr}B$.

\textbf{1er cas.} Si $\text{Tr}A\neq-a$ alors nécessairement $\text{Tr}M=\frac{\text{Tr}B}{a+\text{Tr}A}$ puis $M=\frac{1}{a}\left(B-\frac{\text{Tr}B}{a+\text{Tr}A}A\right)$.

Réciproquement, si $M=\frac{1}{a}\left(B-\frac{\text{Tr}B}{a+\text{Tr}A}A\right)$ alors

\begin{center}
$aM +(TrM)A=B-\frac{\text{Tr}B}{a+\text{Tr}A}A +\frac{1}{a}\left(\text{Tr}B-\frac{\text{Tr}B}{a+\text{Tr}A}\text{Tr}A\right)A = B$.
\end{center}

\begin{center}
\shadowbox{
Si $\text{Tr}A\neq-a$, $\mathcal{S}=\left\{\frac{1}{a}\left(B-\frac{\text{Tr}B}{a+\text{Tr}A}A\right)\right\}$.
}
\end{center}

\textbf{2ème cas.} Si $\text{Tr}A=-a$ et $\text{Tr}B\neq 0$, il n'y a pas de solution .

\textbf{3ème cas.} Si $\text{Tr}A=-a$ et $\text{Tr}B=0$, $M$ est nécessairement de la forme $\frac{1}{a}B +\lambda A$ où $\lambda$ est un réel quelconque.

Réciproquement, soient $\lambda\in\Rr$ puis $M=\frac{1}{a}B+\lambda A$. Alors 

\begin{center}
$aM+(\text{Tr}M)A=B+a\lambda A+\left(\frac{1}{a}\text{Tr}B+\lambda\text{Tr}A\right)A=B+a\lambda A-a\lambda A= B$,
\end{center}

et toute matrice de la forme $B +\lambda A$, $\lambda\in\Rr$, est solution.

\begin{center}
\shadowbox{
Si $\text{Tr}A=-a$, $\mathcal{S}=\varnothing$ si $\text{Tr}B\neq0$ et $\mathcal{S}=\left\{B+\lambda A,\;\lambda\in\Rr\right\}$ si $\text{Tr}B=0$.
}
\end{center}
\fincorrection
\correction{005608}
\begin{enumerate}
 \item  $E=\text{Vect}(I,J)$ est un sous-espace vectoriel de $\mathcal{M}_2(\Rr)$ de dimension inférieure ou égale à 2. De plus, la famille $(I,J)$ est libre car la matrice $J$ n'est pas une matrice scalaire et donc $\text{dim}E= 2$.

\item  Puisque $(E,+,.)$ est un espace vectoriel, $(E,+)$ est un groupe commutatif.

Ensuite, $I^2=I\in E$, $IJ=JI=J\in E$ et $J^2=(I + E_{1,2})^2=I +2E_{1,2}= I+2(J-I) - I =2J -I\in E$. Par bilinéarité du produit matriciel, la multiplication est interne dans $E$ et commutative. De plus, $I\in E$ et finalement $E$ est un sous-anneau commutatif de $\mathcal{M}_2(\Rr)$.

 
\textbf{Remarque.} $M(x,y)M(x',y')= xx'I +(xy'+yx')J +yy'(2J-I)= (xx'-yy')I + (xy'+yx'+2yy')J$.

\item  Soit $(x,y)\in\Rr^2$. 

\begin{align*}\ensuremath
M(x,y)\;\text{est inversible dans}\;E&\Leftrightarrow\exists(x',y')\in\Rr^2/;(xx'-yy')I + (xy'+yx'+2yy')J=I\\
 &\Leftrightarrow\exists(x',y')\in\Rr^2/;\left\{
 \begin{array}{l}
 xx'-yy'=1\\
 yx'+(x+2y)y'=0
 \end{array}
 \right.\;(\text{car la famille}\;(I,J)\;\text{est libre})\quad(*).
\end{align*}

Le déterminant de ce système d'inconnue $(x',y')$ est $x(x+2y)+y^2=(x+y)^2$.

 

\textbullet~Si $x+y\neq0$, le système $(*)$ admet une et une seule solution. Dans ce cas, $M(x,y)$ est inversible dans $E$.

\textbullet~Si $x+y=0$, le système $(*)$ s'écrit $\left\{
 \begin{array}{l}
 x(x'+y')=1\\
 -x(x'+y')=0
 \end{array}
 \right.$ et n'a pas de solution. Dans ce cas, $M(x,y)$ n'est pas inversible dans $E$.
 
 \begin{center}
 \shadowbox{
 $M(x,y)$ est inversible dans $E\Leftrightarrow x+y\neq0$.
 }
 \end{center}
 

\textbf{Remarque.} Puisque $I\in E$, $M(x,y)$ est inversible dans $E$ si et seulement si $M(x,y)$ est inversible dans $\mathcal{M}_2(\Rr)$.

\item  Posons $X= xI+yJ$, $(x,y)\in\Rr^2$. 

  \begin{enumerate}
  \item D'après 1), $X^2=(x^2-y^2)I+(2xy+2y^2)J$. Donc

\begin{align*}\ensuremath
X^2= I&\Leftrightarrow x^2-y^2 = 1\;\text{et}\;2xy+2y^2= 0\;(\text{car la famille}\;(I,J)\;\text{est libre})\\
 &\Leftrightarrow (y = 0\;\text{et}\;x^2 = 1)\;\text{ou}\;(y = -x et 0 = 1)\Leftrightarrow(y = 0\;\text{et}\;x = 1)\;\text{ou}\;(y = 0\;\text{et}\;x = -1)\\
  &\Leftrightarrow X = I\;\text{ou}\;X = -I.
\end{align*}
 
\begin{center}
\shadowbox{
$\mathcal{S}=\{I,-I\}$.
}
\end{center}

  \item
\begin{align*}\ensuremath
X^2=0&\Leftrightarrow x^2-y^2 =0\;\text{et}\;2xy+2y^2= 0\Leftrightarrow (y = 0\;\text{et}\;x^2 = 0)\;\text{ou}\;(y = -x et 0 = 0)\Leftrightarrow y=-x.
\end{align*}

\begin{center}
\shadowbox{
$\mathcal{S}=\{x(I-J),\;x\in\Rr\}$.
}
\end{center}

\textbf{Remarque.} L'équation $X^2=0$, de degré 2, admet une infinité de solutions dans $E$ ce qui montre une nouvelle fois que $(E,+,\times)$ n'est pas un corps.

  \item 

\begin{align*}\ensuremath
X^2=X&\Leftrightarrow x^2-y^2 =x\;\text{et}\;2xy+2y^2= y\Leftrightarrow y(2x+2y-1)= 0\;\text{et}\;x^2-y^2=x\\
 &\Leftrightarrow(y = 0\;\text{et}\;x^2= x)\;\text{ou}\;(2(x+y)=1\;\text{et}\;(x+y)(x-y)=x)\Leftrightarrow(X = 0\;\text{ou}\;X = I)\;\text{ou}\;(2(x+y)=1\;\text{et}\;x-y=2x)\\
  &\Leftrightarrow X = 0\;\text{ou}\;X =I.
\end{align*}

\begin{center}
\shadowbox{
$\mathcal{S}=\{0,I\}$.
}
\end{center}
  \end{enumerate}

\item  Soit $n\in\Nn^*$. On pose $N=J-I=\left(
\begin{array}{cc}
0&1\\
0&0
\end{array}
\right)$. Alors $M(x,y)=xI+y(I+N)=(x+y)I + yN$.

Puisque $I$ et $N$ commutent, la formule du binôme de \textsc{Newton} fournit

\begin{align*}\ensuremath
(M(x,y))^n&=((x+y)I+yN)=(x+y)^nI +ny(x+y)^{n-1}N\;(\text{car}\;N^k=0\;\text{pour}\;k\geqslant 2)\\
 &=\left(
\begin{array}{cc}
(x+y)^n&ny(x+y)^{n-1}\\
0&(x+y)^n
\end{array}
\right).
\end{align*}

\begin{center}
\shadowbox{
$\forall n\in\Nn^*,\;(M(x,y))^n=\left(
\begin{array}{cc}
(x+y)^n&ny(x+y)^{n-1}\\
0&(x+y)^n
\end{array}
\right)$.
}
\end{center}
\end{enumerate}
\fincorrection
\correction{005609}
$\{0\}$ est un idéal bilatère de l'anneau $\mathcal{M}_n(\Kk),+,\times)$.

Soit $I$ un idéal non nul de de l'anneau $\mathcal{M}_n(\Kk),+,\times)$. Montrons que $I=\mathcal{M}_n(\Kk)$.

Il existe une matrice $A$ non nulle dans $I$. Pour tout quadruplet d'indices $(i,j,k,l)$, $I$ contient le produit

\begin{center}
$E_{i,j}AE_{k,l}= \sum_{1\leqslant u,v\leqslant n}^{}a_{u,v}E_{i,j}E_{u,v}E_{k,l}=a_{j,k}E_{i,l}$.
\end{center}

$A$ est non nulle et on peut choisir $j$ et $k$ tels que $a_{j,k}$ soit non nul. $I$ contient alors $a_{j,k}E_{i,l}\frac{1}{a_{j,k}}I_n= E_{i,l}$. Finalement $I$ contient toutes les matrices élémentaires et donc encore toutes les sommes du type $\sum_{1\leqslant i,j\leqslant n}^{}m_{i,j}I_nE_{i,j}=(m_{i,j})_{1\leqslant i,j\leqslant n}$, c'est-à-dire $\mathcal{M}_n(\Kk)$ tout entier.

\begin{center}
\shadowbox{
Les idéaux bilatères de l'anneau $\mathcal{M}_n(\Kk),+,\times)$ sont $\{0\}$ et $\mathcal{M}_n(\Kk)$.
}
\end{center}
\fincorrection
\correction{005618}
Non, car $\text{Tr}(AB-BA)=\text{Tr}(AB) -\text{Tr}(BA) = 0\neq n =\text{Tr}(I_n)$.
\fincorrection
\correction{005619}
Soit $f$ une forme linéaire sur $\mathcal{M}_n(\Cc)$. Pour $A=(a_{i,j})_{1\leqslant i,j\leqslant n}$, posons $f(A)=\sum_{1\leqslant i,j\leqslant n}^{}\alpha_{i,j}a_{i,j}$ où les $\alpha_{i,j}$ sont indépendants de $A$ (les $\alpha_{i,j}$ sont les $f(E_{i,j})$).

Soient $i$ et $j$ deux entiers distincts pris dans $\llbracket1,n\rrbracket$.

\begin{center}
$\alpha_{i,i}=f\left(E_{i,i}\right) = f\left(E_{i,j}E_{j,i}\right) = f\left(E_{j,i}E_{i,j}\right) =f\left(E_{j,j}\right)=\alpha_{j,j}$,
\end{center}

et

\begin{center}
$\alpha_{i,j}=f(E_{i,j}) = f(E_{i,i}E_{i,j}) = f(E_{i,j}E_{i,i}) =f(0) =0$.
\end{center}

Finalement en notant $\alpha$ la valeur commune des $\alpha_{i,i}$, $1\leqslant i\leqslant n$, pour toute matrice $A$ on a $f(A)=\alpha\sum_{i=1}^{n}a_{i,i}=\alpha\text{Tr}A$ où $\alpha$ est indépendant de $A$. (Réciproquement, les $f=\alpha\text{Tr}$, $\alpha\in\Cc$, sont des formes linéaires vérifiant $\forall(A,B)\in\mathcal{M}_n(\Rr)^2$, $f(AB)=f(BA)$.)
\fincorrection
\correction{005620}
 Puisque $\left(\frac{1}{\sqrt{1+\frac{a^2}{n^2}}}\right)^2+\left(\frac{\frac{a}{n}}{\sqrt{1+\frac{a^2}{n^2}}}\right)^2=1$, il existe un unique réel $\theta_n\in[-\pi,\pi[$ tel que

\begin{center}
$\cos\theta_n=\frac{1}{\sqrt{1+\frac{a^2}{n^2}}}$ et $\sin\theta_n=\frac{\frac{a}{n}}{\sqrt{1+\frac{a^2}{n^2}}}$.
\end{center}

La matrice $A_n$ s'écrit alors $A_n=\sqrt{1+\frac{a^2}{n^2}}\left(
\begin{array}{cc}
\cos\theta_n&-\sin\theta_n\\
\sin\theta_n&\cos\theta_n
\end{array}
\right)$ et donc

\begin{center}
$(A_n)^n=\left(1+\frac{a^2}{n^2}\right)^{n/2}\left(
\begin{array}{cc}
\cos(n\theta_n)&-\sin(n\theta_n)\\
\sin(\theta_n)&\cos(n\theta_n)
\end{array}
\right)$.
\end{center}

Maintenant, 

\begin{center}
$\left(1+\frac{a^2}{n^2}\right)^{n/2}=\text{exp}\left(\frac{n}{2}\ln\left(1+\frac{a^2}{n^2}\right)\right)\underset{n\rightarrow+\infty}{=}\text{exp}\left(\frac{n}{2}\times o\left(\frac{1}{n}\right)\right)\underset{n\rightarrow+\infty}{=}\text{exp}(o(1))\underset{n\rightarrow+\infty}{\rightarrow}1$.
\end{center}

Ensuite, en notant $\varepsilon$ le signe de $a$, $\theta_n=\varepsilon\Arccos\left(\frac{1}{\sqrt{1+\frac{a^2}{n^2}}}\right)\underset{n\rightarrow+\infty}{\rightarrow}0$ et on en déduit que

\begin{center}
$n\theta_n \underset{n\rightarrow+\infty}{\sim}n\sin(\theta_n)=n\frac{\frac{a}{n}}{\sqrt{1+\frac{a^2}{n^2}}}\underset{n\rightarrow+\infty}{\rightarrow}a$.
\end{center}

Finalement

\begin{center}
\shadowbox{
$\lim_{n \rightarrow +\infty}(A_n)^n=\left(
\begin{array}{cc}
\cos(a)&-\sin(a)\\
\sin(a)&\cos(a)
\end{array}
\right)$.
}
\end{center}
\fincorrection
\correction{005621}
Soient $i$ et $j$ deux indices pris dans $\llbracket1,n\rrbracket$.

\begin{center}
$f(E_{i,j})=E_{i,j}\sum_{1\leqslant k,l\leqslant n}^{}a_{k,l}E_{k,l}=\sum_{l=1}^{n}a_{j,l}E_{i,l}$,
\end{center}

et en remplissant coefficient à coefficient, on trouve la matrice 
définie par blocs $\left(
\begin{array}{cccc}
{^t}A&0&\ldots&0\\
0&\ddots&\ddots&\vdots\\
\vdots&\ddots&\ddots&0\\
0&\ldots&0&{^t}A
\end{array}
\right)$.
\fincorrection
\correction{005625}
Soit $p$ un entier supérieur ou égal à $2$.

\begin{align*}
A^pB -BA^p&=A^pB-A^{p-1}BA+A^{p-1}BA-A^{p-2}BA^2+A^{p-2}BA^2-...+ABA^{p-1}-BA^p\\
 &=\sum_{k=0}^{}(A^{p-k}BA^k-A^{p-k-1}BA^{k+1})=\sum_{k=0}^{}A^{p-k-1}(AB-BA)A^k=\sum_{k=0}^{}A^{p-k-1}AA^k\sum_{k=0}^{}A^p\\
 &=pA^p.
\end{align*}

Donc $2010\times\text{Tr}(A^{2010})=\text{Tr}(2010\;A^{2010})=\text{Tr}(A^{2010}B)-\text{Tr}(BA^{2010})=0$ et $\text{Tr}(A^{2010})=0$.
\fincorrection
\correction{005628}
\begin{enumerate}
 \item  Soient $p$ l'indice de nilpotence de $A$ et $q$ l'indice de nilpotence de $B$. Puisque $A$ et $B$ commutent, la formule du binôme de \textsc{Newton} fournit

\begin{center}
$(A+B)^{p+q-1}=\sum_{k=0}^{p+q-1}\dbinom{p+q-1}{k}A^kB^{p+q-1-k}$ 
\end{center}

Dans cette somme, 

\textbullet~si $k\geqslant p$, $A^k=0$ et donc $A^kB^{p+q-1-k}= 0$

\textbullet~si $k\leqslant p-1$ alors $p+q-1-k\geqslant q$ et encore une fois $B^{p+q-1-k}= 0$.

Finalement, $(A+B)^{p+q-1}=\sum_{k=0}^{p+q-1}\dbinom{p+q-1}{k}A^kB^{p+q-1-k}=0$ et $A+B$ est nilpotente d'indice inférieur ou égal à $p+q-1$. 

Les sommes définissant $\text{exp}A$, $\text{exp}B$ et $\text{exp}(A+B)$ sont finies car $A$, $B$ et $A+B$ sont nilpotentes et

\begin{align*}\ensuremath
\text{exp}(A+B)&=\sum_{k=0}^{+\infty}\frac{1}{k!}(A+B)^k=\sum_{k=0}^{+\infty}\sum_{i+j=k}^{}\frac{1}{i!j!}A^iB^j\\
 &\left(\sum_{i=0}^{+\infty}\frac{1}{i!}A^i\right)\left(\sum_{j=0}^{+\infty}\frac{1}{j!}B^j\right)\;(\text{toutes les sommes sont finies})\\
 &=\text{exp}A\times\text{exp}B.
\end{align*}

\item  Si $A$ est nilpotente, $-A$ l'est aussi et commute avec $A$. Donc $\text{exp}A\times\text{exp}(-A)=\text{exp}(A-A)=\text{exp}(0)=I_n$.

$\text{exp}A$ est inversible à gauche et donc inversible et $(\text{exp}A)^{-1}=\text{exp}(-A)$.

\item  Les puissances de $A$ sont bien connues et on trouve immédiatement

\begin{center}
$\text{exp}A=\left(
\begin{array}{ccccc}
1&\frac{1}{1!}&\frac{1}{2!}&\ldots&\frac{1}{(n-1)!}\\
0&\ddots&\ddots&\ddots&\vdots\\
\vdots&\ddots& & &\frac{1}{2!}\\
\vdots& &\ddots&\ddots&\rule[-4mm]{0mm}{11mm}\frac{1}{1!}\\
0&\ldots&\ldots&0&1
\end{array}
\right)$.
\end{center}
\end{enumerate}
\fincorrection
\correction{000451}
\begin{enumerate}
\item Soit $r=\frac pq\in \Qq$ et $x\notin\Qq$.
Par l'absurde supposons que $r+x\in \Qq$ alors il existe deux
entiers $p', q'$ tels que $r+x =\frac {p'}{q'}$. Donc $x = \frac
{p'}{q'}-\frac pq = \frac{qp'-pq'}{qq'}\in\Qq$ ce qui est absurde
car $x\notin \Qq$.

De la m\^eme fa\c{c}on  si $r \cdot x \in \Qq$ alors $r \cdot x = \frac{p'}{q'}$
Et donc $x = \frac {p'}{q'}\frac {q}{p}$. Ce qui est absurde.

\item \emph{Méthode ``classique''.} Supposons, par l'absurde, que $\sqrt2 \in \Qq$ alors il existe deux entiers $p,q$ tels que $\sqrt2=\frac pq$. De plus nous pouvons supposer que la fraction est irr\'eductible ($p$ et $q$ sont premiers entre eux). En \'elevant l'\'egalit\'e au carr\'e nous obtenons $q^2\times 2=p^2$. Donc $p^2$ est un nombre pair, cela implique que $p$ est un nombre pair (si vous n'\^etes pas convaincu \'ecrivez la contrapos\'ee ``$p$ impair $\Rightarrow$ $p^2$ impair''). Donc $p = 2\times p'$ avec $p'\in \Nn$, d'o\`u $p^2= 4\times {p'}^2$. Nous obtenons $q^2=2\times {p'}^2$. Nous en d\'eduisons maintenant que $q^2$ est pair et comme ci-dessus que $q$ est pair.
Nous obtenons ainsi une contradiction car  $p$ et $q$ \'etant tous
les deux pairs la fraction $\frac pq$ n'est pas irr\'eductible et
aurait pu \^etre simplifiée. Donc $\sqrt 2\notin\Qq$.

\emph{Autre méthode.} Supposons par l'absurde que $\sqrt 2 \in \Qq$. Alors $\sqrt 2 = \frac pq$ pour deux entiers $p,q \in \Nn^*$.
Alors nous avons $q \cdot \sqrt 2 \in \Nn$. Considérons l'ensemble suivant :
$$\mathcal{N} = \left\lbrace n \in \Nn^* \mid n\cdot \sqrt 2 \in \Nn \right\rbrace.$$
Cet ensemble $\mathcal{N}$ est une partie de $\Nn^*$ qui est non vide car $q\in\mathcal{N}$.
On peut alors prendre le plus petit élément de $\mathcal{N}$ : $n_0 = \min \mathcal{N}$.
En particulier $n_0 \cdot \sqrt 2 \in \Nn$.
Définissons maintenant $n_1$ de la façon suivante : $n_1 = n_0 \cdot \sqrt 2 - n_0$.
Il se trouve que $n_1$ appartient aussi à $\mathcal{N}$ car d'une part
 $n_1 \in \Nn$ (car $n_0$ et $n_0 \cdot \sqrt{2}$ sont des entiers) et d'autre part
$n_1 \cdot \sqrt 2 = n_0 \cdot 2 - n_0 \cdot \sqrt 2 \in \Nn$.
Montrons maintenant que $n_1$ est plus petit que $n_0$.
Comme $0 < \sqrt 2 -1 < 1$ alors $n_1 = n_0 (\sqrt 2 -1) < n_0$ et est non nul.

Bilan : nous avons trouvé $n_1 \in \mathcal{N}$ strictement plus petit que $n_0 = \min \mathcal{N}$.
Ceci fournit une contradiction. Conclusion : $\sqrt{2}$ n'est pas un nombre rationnel.

\item Soient $r,r'$ deux rationnels avec $r<r'$. Notons $x=r + \frac{\sqrt2}{2}(r'-r)$.
D'une part $x\in]r,r'[$ (car $0 < \frac{\sqrt2}{2} < 1$) et d'apr\`es les deux premi\`eres questions
$\sqrt2\left(\frac{r'-r}{2}\right) \notin \Qq$ donc $x\notin \Qq$. Et donc $x$ est un
nombre irrationnel compris entre $r$ et $r'$.
\end{enumerate}
\fincorrection
\nocorrection
\nocorrection
\nocorrection
\nocorrection
\nocorrection
\correction{000457}
\begin{enumerate}
\item
  Soit $\frac \alpha \beta \in \Qq$ avec $\pgcd(\alpha,\beta) = 1$.
Pour $p(\frac \alpha \beta) = 0$, alors $\sum_{i=0}^n {a_i
\left(\frac \alpha \beta \right)^i} = 0$. Apr\`es multiplication par
$\beta^n$ nous obtenons l'\'egalit\'e suivante :
$$
a_n\alpha^n+a_{n-1}\alpha^{n-1}\beta + \cdots +
a_1\alpha\beta^{n-1}+a_0\beta^n = 0.$$ 
En factorisant tous les termes de cette somme sauf le premier par $\beta$, nous \'ecrivons
$a_n\alpha^n+\beta q=0$. Ceci entra\^{\i}ne que $\beta$ divise
$a_n\alpha^n$, mais comme $\beta$ et $\alpha^n$ sont premier entre
eux alors par le lemme de Gauss
$\beta$ divise $a_n$. De m\^eme en factorisant par $\alpha$ tous les termes
de la somme ci-dessus, sauf le dernier,  nous obtenons $\alpha q'
+a_0\beta^n = 0$ et par un raisonnement similaire $\alpha$ divise
$a_0$.
  \item  Notons $\gamma = \sqrt 2+\sqrt 3$.
Alors $\gamma^2 = 5 +2\sqrt 2 \sqrt 3$ Et donc
$\left(\gamma^2-5\right)^2= 4\times 2 \times 3$, Nous choisissons
$p(x) = (x^2-5)^2-24$, qui s'\'ecrit aussi $p(x)=x^4-10x^2+1$. Vu
notre choix de $p$, nous avons $p(\gamma)=0$. Si nous supposons
que $\gamma$ est rationnel, alors $\gamma = \frac \alpha \beta$ et
d'apr\`es la premi\`ere question $\alpha$ divise le terme constant de
$p$, c'est-\`a-dire $1$. Donc $\alpha=\pm 1$. De m\^eme $\beta$
divise le coefficient du terme de plus haut degr\'e de $p$, donc
$\beta$ divise $1$, soit $\beta = 1$. Ainsi $\gamma = \pm 1$, ce
qui est \'evidemment absurde !
\end{enumerate}
\fincorrection
\nocorrection
\correction{000459}
\begin{enumerate}
\item
Soit $p = 1997\,1997\,\ldots 1997$ et $q = 1\, 0000\, 0000\, \ldots
0000 = 10^{4n}$. Alors $N_n = \frac pq$.
\item Remarquons que $10\,000 \times M = 1997,1997\,1997\,\ldots$ Alors
$10\,000 \times M -M=1997$ ; donc $9999\times M = 1997$ d'o\`u $M
= \frac{1997}{9999}$.
\item $0,111\ldots = \frac19$, $0,222\ldots = \frac 29$, etc.
D'o\`u $P = \frac 19 + \frac 29 +\cdots + \frac 99 =
\frac{1+2+\cdots+9}{9}= \frac {45}{9}= 5$.
\end{enumerate}
\fincorrection
\nocorrection
\correction{000461}
Par l'absurde supposons que $\frac{\ln 3}{\ln 2}$ soit un rationnel.
Il s'\'ecrit alors $\frac pq$ avec $p\geqslant 0,q>0$ des entiers.
On obtient $q\ln 3=p\ln 2$. En prenant l'exponentielle nous obtenons :
$\exp (q\ln 3) = \exp(p\ln 2)$ soit
$3^q=2^p$. Si $p \ge 1$ alors $2$ divise $3^q$ donc $2$ divise $3$, ce qui est absurde.
Donc $p=0$. Ceci nous conduit à l'égalité $3^q=1$, donc $q=0$. La seule solution possible est
$p=0$, $q=0$.  Ce qui contredit $q\neq 0$.
Donc $\frac{\ln 3}{\ln 2}$ est irrationnel.
\fincorrection
\nocorrection
\nocorrection
\correction{000506}
Soit $(u_n)$ une suite convergeant vers $\ell \in \Rr$. Par
d\'efinition
$$\forall \epsilon > 0 \quad \exists N \in \Nn \quad  \forall n\geqslant N \qquad |u_n-\ell| < \epsilon.$$
Choisissons $\epsilon = 1$, nous obtenons le  $N$ correspondant.
Alors pour $n\geqslant N$, nous avons $|u_n-\ell| < 1$ ; 
autrement dit $\ell -1 <
u_n < \ell + 1$. Notons $M = \max_{n=0,\ldots,N-1}  \{u_n\}$  et
puis $ M' = \max (M,\ell+1)$. Alors  pour tout $n \in \Nn$ $u_n
\leq M'$. De m\^eme en posant $m = \min_{n=0,\ldots,N-1} \{u_n\}$ et
$m' = \min(m,\ell -1)$ nous obtenons pour tout $n\in \Nn$, $u_n
\geq m'$.
\fincorrection
\correction{000507}
Il est facile de se convaincre que $(u_n)$ n'a pas de
 limite, mais plus délicat d'en donner une d\'emonstration
 formelle. En effet, d\`es lors qu'on ne sait pas qu'une suite $(u_n)$
converge, on ne peut pas \'ecrire $\lim u_n$, c'est un nombre qui
n'est pas d\'efini. Par exemple l'\'egalit\'e $$\lim_{n \rightarrow
\infty}\: (-1)^n+1/n=\lim_{n\rightarrow \infty} (-1)^n$$ n'a pas de
sens. Par contre voil\`a ce qu'on peut dire : \emph{Comme la
suite $1/n$ tend vers $0$ quand $n \rightarrow \infty$, la suite $u_n$ est
convergente si et seulement si la suite $(-1)^n$ l'est. De plus,
dans le cas o\`u elles sont toutes les deux convergentes, elles
ont m\^eme limite.} Cette affirmation provient tout simplement du
th\'eor\`eme suivant

\textbf{Th\'eor\`eme} : Soient $(u_n)$ et $(v_n)$ deux suites
convergeant vers deux limites $\ell$ et $\ell'$. Alors la suite $(w_n)$ définie par
$w_n=u_n+v_n$ est convergente (on peut donc parler de sa limite)
et $\lim w_n=\ell+\ell'$.

De plus, il n'est pas vrai que toute suite convergente
doit forc\'ement \^etre croissante et major\'ee ou d\'ecroissante
et minor\'ee. Par exemple, $(-1)^n/n$ est une suite qui converge
vers $0$ mais qui n'est ni croissante, ni d\'ecroissante. 

\bigskip

Voici maintenant un exemple de r\'edaction de l'exercice.  On veut
montrer que la suite $(u_n)$ n'est pas convergente. Supposons donc
par l'absurde qu'elle soit convergente et notons
$\ell=\lim_{n\rightarrow\infty} u_n$. (Cette expression a un sens puisqu'on
suppose que $u_n$ converge).

{\bf  Rappel.} Une {\em sous-suite} de $(u_n)$ (on dit aussi  {\em
suite extraite} de $(u_n)$) est une suite $(v_n)$ de la forme
$v_n=u_{\phi(n)}$ o\`u $\phi$ est une application strictement
croissante de $\N$ dans $\N$. Cette fonction $\phi$ correspond
``au choix des indices qu'on veut garder'' dans notre sous-suite.
Par exemple, si on ne veut garder dans la suite $(u_n)$ que les
termes pour lesquels $n$ est un multiple de trois, on pourra poser
$\phi(n)=3n$, c'est \`a dire $v_n=u_{3n}$.

\vspace{0.3cm}

Consid\'erons maintenant les sous-suites $v_n=u_{2n}$ et
$w_n=u_{2n+1}$ de $(u_n)$. On a que $v_n=1+1/2n\rightarrow1$ et que
$w_n=-1+1/(2n+1)\rightarrow -1$. Or on a le th\'eor\`eme suivant sur les
sous-suites d'une suite convergente:

\textbf{Th\'eor\`eme} : Soit $(u_n)$ une suite convergeant  vers la
limite $\ell$ (le th\'eor\`eme est encore vrai si $\ell=+\infty$ ou
$\ell=-\infty$). Alors, toute sous-suite $(v_n)$ de $(u_n)$ a pour limite
$\ell$.


Par cons\'equent, ici, on a que $\lim v_n=\ell$ et $\lim w_n=\ell$  donc
$\ell=1$ et $\ell=-1$ ce qui est une contradiction. L'hypoth\`ese disant
que $(u_n)$ \'etait convergente est donc fausse. Donc $(u_n)$ ne
converge pas.
\fincorrection
\nocorrection
\nocorrection
\nocorrection
\nocorrection
\correction{003066}
Supposons d'abord $\bigl|r-\sqrt a\bigr| \le \frac 1{q^2}$.
Cela implique $|r| \le \sqrt a + 1$.

Majorons $|r^2-a|$ :
$$|r^2-a|=|r-\sqrt a| \times |r+\sqrt a| \le |r-\sqrt a| \times \big( |r|+\sqrt a \big) \le |r-\sqrt a| \times \big( 2\sqrt a+1 \big)$$

Minorons $|r^2-a|$, en posant $r=\frac pq$, $a=\frac mn$.
$$|r^2-a|=\left|\left(\frac pq\right)^2-\frac mn \right| =\left| \frac{np^2-mq^2}{nq^2}\right| \ge \frac{1}{nq^2}$$
La dernière inégalité provient que le numérateur $np^2-mq^2$ n'est pas nul (sinon $\sqrt a$ serait rationnel).

On déduit de ces deux majorations : 
$$\frac{1}{nq^2} \le |r^2-a| \le |r-\sqrt a| \times \big( 2\sqrt a+1 \big)$$
Et donc :
$$ |r-\sqrt a| \ge  \frac{1}{n( 2\sqrt a+1)}\frac{1}{q^2}.$$


Cette inégalité est aussi clairement vérifier si $\bigl|r-\sqrt a\bigr| > \frac 1{q^2}$.

La constante $C=\frac{1}{n( 2\sqrt a+1)}$ convient.

\fincorrection
\correction{003067}
$\sqrt x + \sqrt y = \sqrt{a+\sqrt b} \Leftrightarrow
          x+y+2\sqrt{xy} = a+\sqrt b \Leftrightarrow b+4xy-4\sqrt{bxy} = (x+y-a)^2$.\par
          $ \Rightarrow $ : $bxy = r^2  \Rightarrow  \sqrt b\left(1-\frac{2r}b\right) = x+y-a
                   \Rightarrow  r=\frac b2$ et $x+y = a  \Rightarrow  (x-y)^2 = a^2-b$.\par
          $\Leftarrow$ : $a^2 - b = u^2$. On prend $x = \frac{a+u}2$ et
                  $y = \frac{a-u}2$ $ \Rightarrow  x+y+2\sqrt{xy} = a+\sqrt b$.
\fincorrection
\correction{003144}
$=\frac{q-1}2$.
\fincorrection
\correction{003145}
Si $p \in P$ : $\exists\ n \in\Z$ tel que $\frac np \in A$ avec $n\wedge p=1$.

Alors pour tous $x,y \in \Z$, on a $\frac{nx+py}p \in A$,
donc $\frac 1p\Z \subset A$.
\fincorrection
\nocorrection
\correction{003147}
\begin{enumerate}
  \item Si $x_n = \frac pq \ne 0$ : $\frac 1{k_n}\le \frac pq < \frac 1{k_n-1},
     \Rightarrow  x_{n+1} = \frac {k_np-q}{k_nq}$ et $0\le k_np-q < p$.
    Donc la suite des num{\'e}rateurs est strictement d{\'e}croissante.
  \item Car $x_{n+1} = \frac pq-\frac1{k_n} < \frac1{k_n-1}-\frac1{k_n}$.
  \item $n_p > n_{p-1}(n_{p-1}-1)  \Rightarrow  \frac 1{n_{p-1}} + \frac 1{n_p} < \frac 1{n_{p-1}-1}$.\par
    $n_{p-1}-1 \ge n_{p-2}(n_{p-2}-1)  \Rightarrow  \frac 1{n_{p-2}} + \frac 1{n_{p-1}-1} \le \frac 1{n_{p-2}-1}$, etc.\par
    Finalement, $x < \frac 1{n_0-1}  \Rightarrow  n_0 = k_0$.\par
    (cf. INA opt. 1977)
\end{enumerate}
\fincorrection
\correction{003148}
\begin{enumerate}
  \item Pour $x = \frac pq$, on peut prendre :
             $m = qc-pd$, et $n = pb-qa$.
  \item $(m,n)$ est unique {\`a} un facteur pr{\`e}s.
  \item $\frac{ma+nc}{mb+nd} - \frac ab = \frac{n(bc-ad)}{b(mb+nd)}$,
          et $\frac cd - \frac{ma+nc}{mb+nd} = \frac{m(bc-ad)}{d(mb+nd)}$.
\end{enumerate}
\fincorrection
\correction{003149}
\begin{enumerate}
  \item $x=-\frac 12$.
  \item $x=\frac 23$.
  \item Pas de solution.
\end{enumerate}
\fincorrection
\correction{003150}
\begin{enumerate}
  \item $x^y = y^x \Leftrightarrow \frac{p^{p'q}}{q^{p'q}} = \frac{p'^{pq'}}{q'^{pq'}}$
    (formes irr{\'e}ductibles) \\
    $ \Rightarrow  \begin{cases} p^{p'q} = p'^{pq'} \cr q^{p'q} = q'^{pq'}\end{cases}
      \Rightarrow  \begin{cases} p^b = p'^a \cr q^b = q'^a.\end{cases}$\par
    Comme $a\wedge b = 1$, on d{\'e}compose $p,p',q,q'$ en facteurs premiers
    $ \Rightarrow $ le r{\'e}sultat.
  \item $\Bigl( pq' = m^a n^b,\ p'q = m^b n^a,\ m\wedge n = 1,\ a<b \Bigr)
              \Rightarrow  d=m^a n^a$, $a=n^{b-a}$ et $b = m^{b-a}$.
  \item $m\ge n+1$ donc si $b-a \ge 2$, on a :
             $m^{b-a}-n^{b-a} = (m-n)(m^{b-a-1} + \dots + n^{b-a-1}) > b-a$.\par
             Donc $b-a=1=m-n$, $a=n$, $x=\Bigl(1+\frac1n\Bigr)^n$ et
             $y=\Bigl(1+\frac1n\Bigr)^{n+1}$.\par
             Ces valeurs conviennent.

\end{enumerate}
\fincorrection
\correction{005209}

\begin{enumerate}
 \item  Soient $m$ et $n$ deux entiers naturels supérieurs à $2$.

$$\sqrt[n]{m}\in\Qq\Leftrightarrow\exists(a,b)\in(\Nn^*)^2/\;\sqrt[n]{m}=\frac{a}{b}\Leftrightarrow\exists(a,b)\in(\Nn^*)^2/\;a^n=m\times b^n.$$
Tout d'abord, si $b=1$, $m=a^n$ et $m$ est une puissance $n$-ième parfaite. Ensuite, $a=1$ est impossible car $m\times b^n\geq2$.
Supposons alors que $a$ et $b$ soient des entiers supérieurs à $2$ (et que $a^n=m\times b^n$). L'exposant de tout facteur premier de $a^n$ ou de $b^n$ est un multiple de $n$ et par unicité de la décomposition en facteurs premiers, il en est de même de tout facteur premier de $m$. Ceci montre que, si $\sqrt[n]{m}$ est rationnel, $m$ est une puissance $n$-ième parfaite.
Réciproquement, si $m$ est une puissance $n$-ième parfaite, $\sqrt[n]{m}$ est un entier et en particulier un rationnel. En résumé~:~

\begin{center}
\shadowbox{
$\forall(m,n)\in(\Nn\setminus\{0,1\})^2,\;\sqrt[n]{m}\in\Qq\Leftrightarrow\sqrt[n]{m}\in\Nn\Leftrightarrow m\;\mbox{est une puissance}\;n\;\mbox{-ième parfaite}.$
}
\end{center}
Par suite, si $m$ n'est pas une puissance $n$-ième parfaite, $\sqrt[n]{m}$ est irrationnel.
 \item 

\begin{align*}
\log 2\in\Qq&\Rightarrow\exists(a,b)\in(\Nn^*)^2/\;\log2=\frac{a}{b}\Rightarrow\exists(a,b)\in(\Nn^*)^2/\;10^{a/b}=2
\Rightarrow\exists(a,b)\in(\Nn^*)^2/\;10^a=2^b\\
 &\Rightarrow\exists(a,b)\in(\Nn^*)^2/\;5^a=2^{b-a}.
\end{align*} 
Puisque $5^a>1$, ceci impose $b-a\in\Nn^*$. Mais alors, l'égalité ci-dessus est impossible pour $a\neq0$ et $b\neq0$ par unicité de la décomposition en facteurs premiers d'un entier naturel supérieur ou égal à $2$.
On a montré par l'absurde que

\begin{center}
\shadowbox{
$\log2$ est irrationnel.
}
\end{center}
 \item  Supposons par l'absurde que $\pi$ soit rationnel. Il existe alors deux entiers naturels non nuls $p$ et $q$ tels que $\pi=\frac{p}{q}$.
Pour $n$ entier naturel non nul donné, posons 

$$I_n=\frac{1}{n!}\int_{0}^{\pi}x^n(p-qx)^n\sin x\;dx=\frac{1}{n!}\int_{0}^{p/q}x^n(p-qx)^n\sin x\;dx.$$
\textbullet~Tout d'abord, pour $0\leq x\leq\frac{p}{q}$, on a $0\leq x(p-qx)=\frac{p}{2q}\left(p-\frac{p}{2q}\times q\right)=\frac{p^2}{4q}$, 
et donc (puisque $0\leq\sin x\leq1$ pour $x\in[0,\pi]$), 

$$0\leq I_n\leq\frac{1}{n!}\int_{0}^{p/q}\left(\frac{p^2}{4q}\right)^n\;dx=\frac{\pi}{n!}\left(\frac{p^2}{4q}\right)^n.$$
D'après le résultat admis par l'énoncé, $\frac{\pi}{n!}\left(\frac{p^2}{4q}\right)^n$ tend vers $0$ quand $n$ tend vers $+\infty$, et donc d'après le théorème de la limite par encadrement, la suite $(I_n)$ converge et $\lim_{n\rightarrow +\infty}I_n=0$.
\textbullet~Ensuite, puisque pour $x$ élément de $[0,\pi]$, on a $x^n(p-qx)^n\sin x\geq0$, pour $n$ entier naturel non nul donné, on a

\begin{align*}
I_n&=\frac{1}{n!}\int_{0}^{\pi}x^n(p-qx)^n\sin x\;dx\geq\frac{1}{n!}\int_{\pi/4}^{3\pi/4}x^n(p-qx)^n\sin x\;dx\geq
\frac{1}{n!}\left(\frac{3\pi}{4}-\frac{\pi}{4}\right)\left(\frac{p}{4q}\left(p-\frac{p}{4q}\times q\right)\right)^n\frac{1}{\sqrt{2}}\\
 &=\frac{\pi}{2\sqrt{2}n!}\left(\frac{3p^2}{16q}\right)^n>0.
\end{align*}
Donc, $\forall n\in\Nn,\;I_n>0$.
\textbullet~Vérifions enfin que, pour tout entier naturel non nul $n$, $I_n$ est un entier (relatif).
Soit $P_n=\frac{1}{n!}x^n(p-qx)^n$. $P_n$ est un polynôme de degré $2n$ et $0$ et $\frac{p}{q}$ sont racines d'ordre $n$ de $P_n$ et donc, pour $0\leq k\leq n$, racines d'ordre $n-k$ de $P_n^{(k)}$. En particulier, $P_n^{(k)}(0)$ et $P_n^{(k)}\left(\frac{p}{q}\right)$ sont, pour $0\leq k<n$, des entiers relatifs. De même, puisque $\mbox{deg }P_n=2n$, pour $k\geq2n+1$, $P_n^{(k)}\geq0$ et en particulier, $P_n^{(k)}(0)$ et $Pn^{(k)}\left(\frac{p}{q}\right)$ sont, pour $k\geq2n+1$, des entiers relatifs.
Soit $k$ un entier tel que $n\leq k\leq2n$.

$$\frac{1}{n!}x^n(p-qx)^n=\frac{1}{n!}x^n\sum_{i=0}^{n}C_n^ip^{n-i}(-1)^{i}q^ix^i
=\sum_{i=0}^{n}\frac{C_n^i}{n!}p^{n-i}(-1)^{i}q^ix^{n+i}=
\sum_{k=n}^{2n}\frac{C_n^{k-n}}{n!}p^{2n-k}(-1)^{k-n}q^{k-n}x^{k}.$$
On sait alors que 

$$P_n^{(k)}(0)=k!\times(\mbox{coefficient de}\;x^k)=(-1)^{k-n}\frac{k!}{n!}C_n^{k-n}p^{2n-k}q^{k-n}.$$
ce qui montre que $P_n^{(k)}(0)$ est entier relatif (puisque $n\leq k\leq2n$). Puis, comme $P_n\left(\frac{p}{q}-x\right)=P_n(x)$, on a encore 
$P_n^{(k)}\left(\frac{p}{q}-x\right)=(-1)^kP_n^{(k)}(x)$  et en particulier $P_n^{(k)}\left(\frac{p}{q}\right)=(-1)^kP_n^{(k)}(0)\in\Zz$.
On a montré que pour tout entier naturel $k$, $P_n^{(k)}(0)$ et $P_n^{(k)}\left(\frac{p}{q}\right)$ sont des entiers relatifs.
Montrons alors que $I_n$ est un entier relatif.
Une première intégration par parties fournit~:~$I_n=\left[-P_n(x)\cos x\right]_{0}^{p/q}+\int_{0}^{p/q}P_n'(x)\cos x\;dx$.
$\cos$ prend des valeurs entières en $0$ et  $\frac{p}{q}=\pi$ de même que $P_n$. Par suite,
 
$$I_n\in\Zz\Leftrightarrow\int_{0}^{p/q}P_n'(x)\cos x\;dx\in\Zz.$$
Une deuxième intégration par parties fournit~:~$\int_{0}^{p/q}P_n'(x)\cos x\;dx=\left[P_n'(x)\sin x\right]_{0}^{p/q}-\int_{0}^{p/q}P_n''(x)\sin x\;dx$.
sin prend des valeurs entières en $0$ et $\frac{p}{q}=\pi$, de même que $P_n'$ et 

$$I_n\in\Zz\Leftrightarrow\int_{0}^{p/q}P_n''(x)\sin x\;dx\in\Zz.$$
En renouvelant les intégrations par parties et puisque sin et cos prennent des valeurs entières en $0$ et $\pi$ 
de même que les dérivées succesives de $P_n$, on en déduit que~:

$$I_n\in\Zz\Leftrightarrow\int_{0}^{p/q}P_n^{(2n)}(x)\sin x\;dx\in\Zz.$$
Mais, 

$$\int_{0}^{p/q}P_n^{(2n)}(x)\sin x\;dx=\int_{0}^{p/q}\frac{1}{n!}(-q)^n(2n)!\sin x\;dx=2(-q)^n(2n)(2n-1)...(n+1)\in\Zz.$$
Donc pour tout naturel $n$, $I_n$ est un entier relatif, strictement positif d'après plus haut. On en déduit que pour tout naturel $n$, $I_n\geq1$. Cette dernière constatation contredit le fait que la suite $(I_n)$ converge vers $0$.
L'hypothèse $\pi$ est rationnel est donc absurde et par suite,

\begin{center}
\shadowbox{
$\pi$ est irrationnel.
}
\end{center}
 \item  Montrons par récurrence que~:~$\forall n\in\Nn,\;e=\sum_{k=0}^{n}\frac{1}{k!}+\int_{0}^{1}\frac{(1-t)^n}{n!}e^t\;dt$.
\textbullet~Pour $n=0$, $\int_{0}^{1}\frac{(1-t)^n}{n!}e^t\;dt=\int_{0}^{1}e^t\;dt=e-1$ et donc, $e=1+\int_{0}^{1}e^t\;dt=\sum_{k=0}^{0}\frac{1}{k!}+\int_{0}^{1}\frac{(1-t)^0}{0!}e^t\;dt$.
\textbullet~Soit $n\geq0$. Supposons que $e=\sum_{k=0}^{n}\frac{1}{k!}+\int_{0}^{1}\frac{(1-t)^n}{n!}e^t\;dt$.
Une intégrations par parties fournit~:

$$\int_{0}^{1}\frac{(1-t)^n}{n!}e^t\;dt=\left[-\frac{(1-t)^{n+1}}{(n+1)\times n!}e^t\right]_0^1+\int_{0}^{1}\frac{(1-t)^{n+1}}{(n+1)!}e^t\;dt=\frac{1}{(n+1)!}+\int_{0}^{1}\frac{(1-t)^{n+1}}{(n+1)!}e^t\;dt,$$
et donc,

$$e=\sum_{k=0}^{n}\frac{1}{k!}+\frac{1}{(n+1)!}+\int_{0}^{1}\frac{(1-t)^{n+1}}{(n+1)!}e^t\;dt=
\sum_{k=0}^{n+1}\frac{1}{k!}+\int_{0}^{1}\frac{(1-t)^{n+1}}{(n+1)!}e^t\;dt.$$
Le résultat est ainsi démontré par récurrence.
Soit $n$ un entier naturel non nul. D'après ce qui précède,

$$0<e-\sum_{k=0}^{n}\frac{1}{k!}=\int_{0}^{1}\frac{(1-t)^n}{n!}e^t\;dt<e\int_{0}^{1}\frac{(1-t)^n}{n!}\;dt=\frac{e}{(n+1)!}<\frac{3}{(n+1)!}.$$
Supposons alors par l'absurde que $e$ soit rationnel. Alors, il existe $(a,b)\in(\Nn^*)^2/\;e=\frac{a}{b}$.
Soit $n$ un entier naturel non nul quelconque. D'après ce qui précède, on a 
$0<\frac{a}{b}-\sum_{k=0}^{n}\frac{1}{k!}<\frac{3}{(n+1)!}$, ce qui s'écrit encore après multiplication des trois membres par $bn!$

\begin{center}
$0<a\times n!-b\sum_{k=0}^{n}\frac{n!}{k!}<\frac{3b}{n+1}$.
\end{center}
En particulier, pour $n=3b$, on a $0<a\times(3b)!-b\sum_{k=0}^{3b}\frac{(3b)!}{k!}<\frac{3b}{3b+1}<1$. Mais ceci est impossible car $a\times n!-b\sum_{k=0}^{3b}\frac{(3b)!}{k!}$ est un entier relatif. Il etait donc absurde de supposer que $e$ est rationnel et finalement,

\begin{center}
\shadowbox{
$e$ est irrationnel.
}
\end{center}
 \item  Une équation du troisième degré dont les solutions sont $\cos\frac{2\pi}{7}$, $\cos\frac{4\pi}{7}$ et $\cos\frac{6\pi}{7}$ est 

\begin{center}
$(X-\cos\frac{2\pi}{7})(X-\cos\frac{4\pi}{7})(X-\cos\frac{6\pi}{7})=0$,
\end{center}
ou encore

$$X^3-\left(\cos\frac{2\pi}{7}+\cos\frac{4\pi}{7}+\cos\frac{6\pi}{7}\right)X^2+\left(\cos\frac{2\pi}{7}\cos\frac{4\pi}{7}+
\cos\frac{2\pi}{7}\cos\frac{6\pi}{7}+\cos\frac{4\pi}{7}\cos\frac{6\pi}{7}\right)X-\cos\frac{2\pi}{7}\cos\frac{4\pi}{7}\cos\frac{6\pi}{7}=0.$$
Calculons alors ces trois coefficients.
Soit $\omega=e^{2i\pi/7}$. Puisque $\omega^7=1$ et que $\omega+\omega^2+\omega^3+\omega^4+\omega^5+\omega^6=-1$, on a d'après les formules d'\textsc{Euler}

\begin{align*}
\cos\frac{2\pi}{7}+\cos\frac{4\pi}{7}+\cos\frac{6\pi}{7}&=\frac{1}{2}(\omega+\omega^6+\omega^2+\omega^5+\omega^3+\omega^4)=-\frac{1}{2},
\end{align*}
puis,

\begin{align*}
\cos\frac{2\pi}{7}\cos\frac{4\pi}{7}+&
\cos\frac{2\pi}{7}\cos\frac{6\pi}{7}+\cos\frac{4\pi}{7}\cos\frac{6\pi}{7}=\frac{1}{4}
((\omega+\omega^6)(\omega^2+\omega^5)+(\omega+\omega^6)(\omega^3+\omega^4)+(\omega^2+\omega^5)(\omega^3+\omega^4))\\
 &=\frac{1}{4}((\omega^3+\omega^6+\omega+\omega^4)+(\omega^4+\omega^5+\omega^2+\omega^3)+(\omega^5+\omega^6+\omega+\omega^2))\\
 &=\frac{2(-1)}{4}=-\frac{1}{2},
\end{align*}
et enfin,

\begin{align*}
\cos\frac{2\pi}{7}\cos\frac{4\pi}{7}\cos\frac{6\pi}{7}&=\frac{1}{8}(\omega+\omega^6)(\omega^2+\omega^5)(\omega^3+\omega^4)\\
 &=\frac{1}{8}(\omega^3+\omega^6+\omega+\omega^4)(\omega^3+\omega^4)=\frac{1}{8}(\omega^6+1+\omega^2+\omega^3+\omega^4+\omega^5+1+\omega)=\frac{1}{8}
\end{align*}
Les trois nombres $\cos\frac{2\pi}{7}$, $\cos\frac{4\pi}{7}$ et $\cos\frac{6\pi}{7}$ sont donc solution de l'équation $X^3+\frac{1}{2}X^2-\frac{1}{2}X-\frac{1}{8}=0$ ou encore de l'équation

$$8X^3+4X^2-4X-1=0.$$
Montrons que cette équation n'admet pas de racine rationnelle. Dans le cas contraire, si, pour $p$ entier relatif non nul et $q$ entier naturel non nul tels que $p$ et $q$ sont premiers entre eux, le nombre $r=\frac{p}{q}$ est racine de cette équation, alors $8p^3+4p^2q-4pq^2-q^3=0$. Ceci peut encore s'écrire $8p^3=q(-4p^2+4pq+q^2)$ ce qui montre que $q$ divise $8p^3$. Comme $q$ est premier avec $p$ et donc avec $p^3$, on en déduit d'après le théorème de \textsc{Gauss} que $q$ divise $8$. De même, l'égalité $q^3=p(8p^2+4pq-4q^2)$ montre que $p$ divise $q^3$ et donc que $p$ divise $1$.
Ainsi, $p\in\{-1,1\}$ et $q\in\{1,2,4,8\}$ ou encore $r\in\left\{1,-1,\frac{1}{2},-\frac{1}{2},\frac{1}{4},-\frac{1}{4},\frac{1}{8},-\frac{1}{8}\right\}$.
On vérifie alors aisément qu'aucun de ces nombres n'est racine de l'équation considérée et donc cette équation n'a pas de racine rationnelle. En particulier,

\begin{center}
\shadowbox{
$\cos\frac{2\pi}{7}$ est irrationnel.
}
\end{center}
 \item  On sait que $\sqrt{2}$, $\sqrt{3}$ et $\sqrt{5}$ sont irrationnels mais ceci n'impose rien à la somme $\sqrt{2}+\sqrt{3}+\sqrt{5}$.
Soit $\alpha=\sqrt{2}+\sqrt{3}+\sqrt{5}$.

\begin{align*}
\alpha=\sqrt{2}+\sqrt{3}+\sqrt{5}&\Rightarrow(\alpha-\sqrt{2})^2=(\sqrt{3}+\sqrt{5})^2 \Rightarrow\alpha^2-2\sqrt{2}\alpha+2=8+2\sqrt{15}\\
 &\Rightarrow(\alpha^2-2\sqrt{2}\alpha-6)^2=60\Rightarrow\alpha^4+8\alpha^2-24=4\sqrt{2}\alpha(\alpha^2-6)
\end{align*}
 

Si maintenant, on suppose que $\alpha$ est rationnel, puisque $\sqrt{2}$ est irrationnel, on a nécessairement $\alpha(\alpha^2-6)=0$ (dans le cas contraire, $\sqrt{2}=\frac{\alpha^4+8\alpha^2-24}{4\alpha(\alpha^2-6)}\in\Qq$). Mais $\alpha$ n'est ni $0$, ni $-\sqrt{6}$, ni $\sqrt{6}$ (car $\alpha^2>2+3+5=10>6$). Donc

\begin{center}
\shadowbox{
$\sqrt{2}+\sqrt{3}+\sqrt{5}$ est irrationnel.
}
\end{center}
\end{enumerate}
\fincorrection
\correction{005214}
Soient $k$ un entier naturel non nul et $n$ un entier naturel supérieur ou égal à $k$.

\begin{align*}
\binom{n+10\times k!}{k}&=\frac{(n+10\times k!)(n+10\times k!-1)...(n+10\times k!-k+1)}{k!}\\
 &=\frac{n(n-1)...(n-k+1)+10\times k!\times K}{k!}\quad(\mbox{pour un certain entier}\;K)\\
 &=\frac{n(n-1)...(n-k+1)}{k!}+10K=\binom{n}{k}+10K.
\end{align*}
La différence $\dbinom{n+10\times k!}{k}-\dbinom{n}{k}$ est donc divisible par $10$. Par suite, $\dbinom{n+10\times k!}{k}$ et $\dbinom{n}{k}$ ont même chiffre des unités en base 10. Ainsi, $\forall n\geq k,\;u_{n+10\times k!}=u_n$ et donc la suite $u$ est donc $10k!$-périodique. On sait alors que 

\begin{center}
\shadowbox{
$0,u_ku_{k+1}u_{k+2}...$ est rationnel.
}
\end{center}
\fincorrection
\correction{005243}
Soit $x$ un irrationnel et $(\frac{p_n}{q_n})_{n\in\Nn}$ une suite de rationnels tendant vers $x$ ($p_n$ entier relatif et $q_n$ entier naturel non nul, la fraction $\frac{p_n}{q_n}$ n'étant pas nécessairement irréductible). Supposons que la suite $(q_n)_{n\in\Nn}$ ne tende pas vers $+\infty$. Donc~:

$$\exists A>0/\;(\forall n_0\in\Nn)(\exists n\geq n_0/\;q_n\geq A)$$ 

ou encore, il existe une suite extraite $(q_{\varphi}(n))_{n\in\Nn}$ de la suite $(q_n)_{n\in\Nn}$ qui est bornée.

La suite $(q_{\varphi}(n))_{n\in\Nn}$ est une suite d'entiers naturels qui est bornée, et donc cette suite ne prend qu'un nombre fini de valeurs. Mais alors, on peut extraire de la suite $(q_{\varphi}(n))_{n\in\Nn}$ et donc de la suite $(q_n)_{n\in\Nn}$ une suite $(q_{\psi(n)})_{n\in\Nn}$ qui est constante et en particulier convergente.

La suite $(p_{\psi(n)})_{n\in\Nn}=(\frac{p_{\psi(n)}}{q_{\psi(n)}})_{n\in\Nn}(q_{\psi(n)})_{n\in\Nn}$ est aussi une suite d'entiers relatifs convergente et est donc constante à partir d'un certain rang.

Ainsi, on peut extraire de la suite $(p_{\psi(n)})_{n\in\Nn}$ et donc de la suite $(p_n)_{n\in\Nn}$ une suite $(p_{\sigma(n)})_{n\in\Nn}$ constante.
La suite $((q_{\sigma(n)})_{n\in\Nn}$ est également constante car extraite de la suite constante $(q_{\psi(n)})_{n\in\Nn}$ et finalement, on a extrait de la suite $(\frac{p_n}{q_n})_{n\in\Nn}$  une sous suite $(\frac{p_{\sigma(n)}}{q_{\sigma(n)}})_{n\in\Nn}$ constante.

Mais la suite $(\frac{p_n}{q_n})_{n\in\Nn}$ tend vers $x$ et donc la suite extraite $(\frac{p_{\sigma(n)}}{q_{\sigma(n)}})_{n\in\Nn}$ tend vers $x$. Puisque $(\frac{p_{\sigma(n)}}{q_{\sigma(n)}})_{n\in\Nn}$ est constante, on a $\forall n\in\Nn,\;\frac{p_{\sigma(n)}}{q_{\sigma(n)}}=x$ et donc $x$ est rationnel. Contradiction .

Donc la suite $(q_n)_{n\in\Nn}$ tend vers $+\infty$. Enfin si $(|p_n|)_{n\in\Nn}$ ne tend pas vers $+\infty$, on peut extraire de $(p_n)_{n\in\Nn}$ une sous-suite bornée $(p_{\varphi}(n))_{n\in\Nn}$. Mais alors, la suite $(\frac{p_{\varphi(n)}}{q_{\varphi(n)}})_{n\in\Nn}$ tend vers $x=0$ contredisant l'irrationnalité de $x$. Donc, la suite $(|p_n|)_{n\in\Nn}$ tend vers $+_infty$.
\fincorrection
\correction{000465}
$(u_{2k})_k$ tend vers $+\infty$ et donc 
$A$ ne possède pas de majorant, ainsi
$A$ n'a pas de borne supérieure (cependant certains écrivent alors $\sup A = +\infty$). 
D'autre part toutes les valeurs de $(u_n)$ sont
positives et $(u_{2k+1})_k$ tend vers $0$, donc $\inf A =0$.
\fincorrection
\correction{000466}
\begin{enumerate}
\item $[0,1]\cap \Qq$. Les majorants   : $[1,+\infty[$. Les minorants : $]-\infty,0]$. La borne sup\'erieure :
$1$. La borne inf\'erieure : $0$. Le plus grand \'el\'ement : $1$. Le
plus petit \'el\'ement $0$.
\item $]0,1[\cap \Qq$. Les majorants   : $[1,+\infty[$. Les minorants : $]-\infty,0]$. La borne sup\'erieure :
$1$. La borne inf\'erieure : $0$. Il nexiste pas de plus grand
\'el\'ement ni de plus petit \'el\'ement.
\item $\Nn$. Pas de majorants, pas de borne sup\'erieure, ni de plus grand \'el\'ement. Les minorants : $]-\infty,0]$.  La borne inf\'erieure : $0$. Le plus petit \'el\'ement : $0$.
\item $\Big\lbrace (-1)^n+\frac{1}{n^2} \mid n \in \Nn^* \Big\rbrace$. Les majorants   : $[\frac54,+\infty[$. Les minorants : $]-\infty,-1]$. La borne sup\'erieure :
$\frac54$. La borne inf\'erieure : $-1$. Le plus grand \'el\'ement :
$\frac54$. Pas de  plus petit \'el\'ement.
\end{enumerate}
\fincorrection
\nocorrection
\nocorrection
\nocorrection
\nocorrection
\nocorrection
\nocorrection
\nocorrection
\nocorrection
\nocorrection
\correction{000476}
\begin{enumerate}

    \item  Soient $A$ et $B$ deux parties born\'ees de $\R$.
On sait que $\sup A$ est un majorant de $A$, c'est-\`a-dire,
pour tout $a\in A$, $a\leqslant \sup A$. De m\^eme, pour tout $b\in B$, $b\le
\sup B$. On veut montrer que $\sup A+\sup B$ est un majorant de
$A+B$. Soit donc $x\in A+B$. Cela signifie que $x$ est de la forme
$a+b$ pour un $a\in A$ et un $b\in B$. Or $a\leqslant \sup A$, et $b \le
\sup B$, donc $x=a+b\leqslant \sup A+\sup B$. Comme ce raisonnement est
valide pour tout $x\in A+B$ cela signifie que  $\sup A+\sup B$ est
un majorant de $A+B$.

    \item On veut montrer que, quel que soit
$\epsilon>0$, $\sup A +\sup B-\epsilon$ n'est pas un majorant de $A+B$. On
prend donc un $\epsilon >0$ quelconque, et on veut montrer que $\sup A
+\sup B-\epsilon$ ne majore pas $A+B$. On s'interdit donc dans la
suite de modifier $\epsilon$. Comme $\sup A$ est le plus petit des
majorants de $A$, $\sup A-\epsilon/2$ n'est pas un majorant de $A$.
Cela signifie qu'il existe un \'el\'ement $a$ de $A$ tel que
$a>\sup A-\epsilon/2$. {\em Attention: $\sup A-\epsilon/2$ n'est pas
forc\'ement dans $A$ ; $\sup A$ non plus.} De la m\^eme mani\`ere, il existe $b\in B$ tel que
$b>\sup B-\epsilon/2$. Or l'\'el\'ement $x$ d\'efini par $x=a+b$ est
un \'el\'ement de $A+B$, et il v\'erifie $x>(\sup A-\epsilon/2)+(\sup
B-\epsilon/2)=\sup A +\sup B-\epsilon.$ Ceci implique que $\sup A +\sup
B-\epsilon$ n'est pas un majorant de $A+B$.

    \item  $\sup A+\sup B$
est un majorant de $A+B$ d'apr\`es la partie 1. Mais, d'apr\`es la
partie 2., d\`es qu'on prend un $\epsilon>0$, $\sup A+\sup B -\epsilon$
n'est pas un majorant de $A+B$. Donc $\sup A+\sup B$ est bien le
plus petit des majorants de $A+B$, c'est donc la borne supérieure de $A+B$. Autrement dit
 $\sup (A+B)= \sup A +\sup B$.
\end{enumerate}
\fincorrection
\correction{000477}
\begin{enumerate}
  \item Vrai.
  \item Faux. C'est vrai avec l'hypothèse $B \subset A$ et non $A \subset B$.
  \item Vrai.
  \item Faux. Il y a égalité.
  \item Vrai.
  \item Vrai.
\end{enumerate}
\fincorrection
\nocorrection
\nocorrection
\nocorrection
\nocorrection
\nocorrection
\nocorrection
\nocorrection
\nocorrection
\nocorrection
\correction{005210}
$A$ et $B$ sont deux parties non vides et majorées de $\Rr$ et admettent donc des bornes supérieures notées respectivement $\alpha$ et $\beta$.
Pour tout $(a,b)\in A\times B$, on a $a+b\leq\alpha+\beta$. Ceci montre que $A+B$ est une partie non vide et majorée de $\Rr$, et donc que $\mbox{sup}(A+B)$ existe dans $\Rr$. (De plus, puisque $\alpha+\beta$ est un majorant de $A+B$, on a déjà $\mbox{sup}(A+B)\leq\alpha+\beta$).
Soit alors $\varepsilon>0$.
Il existe $a\in A$ et $b\in B$ tels que $\alpha-\frac{\varepsilon}{2}<a\leq\alpha$ et $\beta-\frac{\varepsilon}{2}<b\leq\beta$, et donc tels que $\alpha+\beta-\varepsilon<a+b\leq\alpha+\beta$.

En résumé,

\begin{center}
\textbf{(1)} $\forall(a,b)\in A\times B,\;a+b\leq\alpha+\beta$ et \textbf{(2)} $\forall\varepsilon>0,\;\exists(a,b)\in A\times B/\;a+b>\alpha+\beta-\varepsilon$.
\end{center}

On en déduit que

\begin{center}
\shadowbox{
$\mbox{sup}(A+B)=\alpha+\beta=\mbox{sup }A+\mbox{sup }B$.
}
\end{center}
Pour les bornes inférieures, on peut refaire le travail précédent en l'adaptant ou appliquer le résultat précédent aux ensembles $-A$ et $-B$ car $\text{Inf}A=-\text{sup}(-A)$.
\fincorrection
\correction{005211}

Posons pour $n$ entier naturel non nul $u_n=\frac{1}{n}+(-1)^n$ de sorte que $A=\{u_n,\;n\in\Nn^*\}=\left\{0,\frac{1}{2}+1,\frac{1}{3}-1,\frac{1}{4}+1,\frac{1}{5}-1,...\right\}$.
Pour $n\geq1$, $u_{2n}=1+\frac{1}{2n}$. Donc $\forall n\in\Nn^*$, $1<u_{2n}\leq\frac{3}{2}$.
Pour $n\geq1$, $u_{2n-1}=-1+\frac{1}{2n-1}$. Donc $\forall n\in\Nn^*$, $-1<u_{2n-1}\leq0$.
Par suite, $\forall n\in\Nn^*,\;-1<u_n\leq\frac{3}{2}$. Donc, $\mbox{sup }A$ et $\mbox{inf }A$ existent dans $\Rr$ et de plus $-1\leq\mbox{inf }A\leq\mbox{sup }A\leq\frac{3}{2}$.
Ensuite, $\frac{3}{2}=u_2\in A$. Donc,

\begin{center}
\shadowbox{
$\mbox{sup }A=\mbox{max }A=\frac{3}{2}$.
}
\end{center}
Enfin, pour chaque entier naturel non nul $n$, ona $-1\leq\mbox{inf }A\leq u_{2n-1}=-1+\frac{1}{2n-1}$. On fait tendre  $n$ tend vers l'infini dans cet encadrement, on obtient

\begin{center}
\shadowbox{
$\mbox{inf }A=-1$
}
\end{center}
(cette borne inférieure n'est pas un minimum).
\fincorrection
\correction{005212}
 Posons $B=\{|y-x|,\;(x,y)\in A^2\}$.
$A$ est une partie non vide et bornée de $\Rr$, et donc $m=\mbox{inf }A$ et $M=\mbox{sup }A$ existent dans $\Rr$.
Pour $(x,y)\in A^2$, on a $m\leq x\leq M$ et $m\leq y M$, et donc $y-x\leq M-m$ et $x-y\leq M-m$ ou encore $|y-x|\leq M-m$.
Par suite, $B$ est une partie non vide et majorée de $\Rr$. $B$ admet donc une borne supérieure.
Soit $\varepsilon>0$. Il existe $(x_0,y_0)\in A^2$ tel que $x_0<\mbox{inf }A+\frac{\varepsilon}{2}$ et $y_0>\mbox{sup }A-\frac{\varepsilon}{2}$.

Ces deux éléments $x_0$ et $y_0$ vérifient, 

$$|y_0-x_0|\geq y_0-x_0>\left(\mbox{sup }A-\frac{\varepsilon}{2}\right)-\left(\mbox{inf }A+\frac{\varepsilon}{2}\right)=\mbox{sup }A-\mbox{inf }A-\varepsilon.$$
En résumé, 
\begin{enumerate}
 \item  $\forall(x,y)\in A^2,\;|y-x|\leq\mbox{sup }A-\mbox{inf }A$ et  
 \item  $\forall\varepsilon>0,\;\exists(x,y)\in A^2/\;|y-x|>\mbox{sup }A-\mbox{inf }A-\varepsilon$.
\end{enumerate}
Donc, $\mbox{sup }B=\mbox{sup }A-\mbox{inf }A$.

\begin{center}
\shadowbox{
$\mbox{sup }\{|y-x|,\;(x,y)\in A^2\}=\mbox{sup }A-\mbox{inf }A$.
}
\end{center}
\fincorrection
\correction{005213}

\begin{enumerate}
 \item  $A\cap B$ peut être vide et on n'a rien à dire. Supposons donc $A\cap B$ non vide.
Pour $x\in A\cap B$, on a $x\leq\mbox{sup }A$ et $x\leq\mbox{sup }B$ et donc $x\leq\mbox{min}\{\mbox{sup }A,\mbox{sup }B\}$.
Dans ce cas, $\mbox{sup}(A\cap B)$ existe et $\mbox{sup}(A\cap B)\leq\mbox{min}\{\mbox{sup }A,\mbox{sup }B\}$.
On ne peut pas améliorer. Par exemple, soit $A=[0,1]\cap\Qq$ et $B=([0,1]\cap(\Rr\setminus\Qq))\cup\{0\}$. On a $\mbox{sup }A=1$, $\mbox{sup }B=1$, $A\cap B=\{0\}$ et donc $\mbox{sup}(A\cap B)=0<1=\mbox{min}\{\mbox{sup }A,\mbox{sup }B\}$.
 \item  Pour $x\in A\cup B$, on a $x\leq\mbox{max}\{\mbox{sup }A,\mbox{sup }B\}$. Donc $\mbox{sup}(A\cup B)$ existe dans $\Rr$ et $\mbox{sup}(A\cup B)\leq\mbox{max}\{\mbox{sup }A,\mbox{sup }B\}$.
Inversement, supposons par exemple $\mbox{sup }A\geq\mbox{sup }B$ de sorte que $\mbox{max}\{\mbox{sup }A,\mbox{sup }B\}=\mbox{sup }A$.
Soit alors $\varepsilon>0$. Il existe $a\in A$ tel que $\mbox{sup }A-\varepsilon<a\leq\mbox{sup }A$. $a$ est dans $A$ et donc dans $A\cup B$.
En résumé, $\forall x\in(A\cup B),\;x\leq\mbox{max}\{\mbox{sup }A,\mbox{sup }B\}$ et $\forall\varepsilon>0,\;\exists x\in(A\cup B)/\;\mbox{max}\{\mbox{sup }A,\mbox{sup }B\}-\varepsilon<x$ et donc

\begin{center}
\shadowbox{
$\mbox{sup}(A\cup B)=\mbox{max}\{\mbox{sup }A,\mbox{sup }B\}$.
}
\end{center}
 \item  D'après l'exercice \ref{exo:suprou2bis}, $\mbox{sup}(A+B)=\mbox{sup }A+\mbox{sup }B$.
 \item  Pour $\mbox{sup}(AB)$, tout est possible. Par exemple, si $A=B=]-\infty,0]$ alors $\mbox{sup }A=\mbox{sup }B=0$, mais $AB=[0,+\infty[$ et $\mbox{sup}(AB)$ n'existe pas dans $\Rr$.
\end{enumerate}
\fincorrection
\nocorrection
\nocorrection
\nocorrection
\nocorrection
\correction{000491}
$$ \sqrt a +\sqrt b \leqslant 2\sqrt{a+b} \Leftrightarrow (\sqrt a +\sqrt b)^2 \leqslant 2(a+b)$$
 car les termes sont positifs, et la fonction $x \mapsto x^2$ est croissante sur $\Rr_+$.
\'Evaluons la différence $2(a+b) - (\sqrt a +\sqrt b)^2$ :
$$2(a+b) - (\sqrt a +\sqrt b)^2 = a+b-2\sqrt a \sqrt b = (\sqrt a - \sqrt b)^2 \geqslant 0.$$
Donc par l'\'equivalence, nous obtenons l'in\'egalit\'e recherch\'ee.
\fincorrection
\nocorrection
\nocorrection
\nocorrection
\nocorrection
\nocorrection
\correction{000497}
\begin{enumerate}
\item  Calculons d'abord $f(0)$. Nous savons $f(1) = f(1+0) = f(1) +f(0)$, donc $f(0) = 0$.
Montrons le r\'esultat demand\'e par r\'ecurrence : pour $n=1$, nous
avons bien $f(1)=1\times f(1)$. Si $f(n) = n f(1)$ alors $f(n+1) =
f(n) + f(1) = nf(1) + f(1) = (n+1)f(1)$.
\item $0 = f(0) = f(-1 + 1) = f(-1) + f(1)$. Donc $f(-1) = - f(1)$. Puis comme ci-dessus $f(-n) = n f(-1)= -n f(1)$.
\item Soit $q = \frac ab$. Alors $f(a) = f(\frac ab + \frac ab + \cdots +\frac ab) = f(\frac ab ) + \cdots + f(\frac ab)$
($b$ termes dans ces sommes). Donc $f(a) = b f(\frac ab)$. Soit
$a f(1) = b f(\frac ab)$. Ce qui s'\'ecrit aussi $f(\frac ab) =
\frac ab f(1)$.
\item Fixons $x \in \Rr$. Soit $(\alpha_i)$ une suite croissante de rationnels qui tend vers $x$. Soit
$(\beta_i)$ une suite d\'ecroissante de rationnels qui tend vers $x$
:
$$\alpha_1\leq \alpha_2 \leq \alpha_3 \leq \ldots \leq x \leq \cdots \leq \beta_2 \leq \beta_1.$$
Alors comme $\alpha_i \leq x \leq \beta_i$ et que $f$ est
croissante nous avons $f(\alpha_i)\leq f(x) \leq f(\beta_i)$.
D'apr\`es la question pr\'ec\'edent cette in\'equation devient : $\alpha_i
f(1)\leq f(x)\leq \beta_i f(1)$. Comme $(\alpha_i)$ et $(\beta_i)$
tendent vers $x$. Par le ``th\'eor\`eme des gendarmes'' nous obtenons
en passant \`a la limite : $x f(1) \leq f(x) \leq xf(1)$.  Soit
$f(x) = xf(1)$.
\end{enumerate}
\fincorrection
\nocorrection
\nocorrection
\nocorrection
\nocorrection
\nocorrection
\nocorrection
\nocorrection
\nocorrection
\nocorrection
\nocorrection
\correction{003065}
\begin{enumerate}
\item $a=bq+r \Rightarrow \sum = \underbrace{q + q + \dots + q}_{b-r}
  + \underbrace{(q+1) + \dots + (q+1)}_r = bq+r = a$.
\item
\end{enumerate}
\fincorrection
\correction{005146}
Soient $x$ et $y$ deux réels tels que $0<x\leq y$.
\begin{enumerate}
\item  On a déjà $x=\frac{x+x}{2}\leq\frac{x+y}{2}=m\leq\frac{y+y}{2}=y$ et donc \shadowbox{$x\leq m\leq y$}.

(on peut aussi écrire~:~$m-x=\frac{x+y}{2}-x=\frac{y-x}{2}\geq0$).
\item  On a ensuite $x=\sqrt{x.x}\leq\sqrt{xy}=g\leq\sqrt{y.y}=y$ et donc \shadowbox{$x\leq g\leq y$}.
\item 
$m-g=\frac{x+y}{2}-\sqrt{xy}=\frac{1}{2}((\sqrt{x})^2-2\sqrt{xy}+(\sqrt{y})^2)=\frac{1}{2}(\sqrt{y}-\sqrt{x})^2
\geq0$ et donc, \shadowbox{$x\leq g\leq m\leq y$}.
\item  D'après 1), la moyenne arithmétique de $\frac{1}{x}$ et $\frac{1}{y}$ est comprise entre $\frac{1}{x}$ et
$\frac{1}{y}$, ce qui fournit $\frac{1}{y}\leq\frac{1}{h}\leq\frac{1}{x}$, ou encore \shadowbox{$x\leq h\leq y$}.
\item  D'après 3), la moyenne géométrique des deux réels $\frac{1}{x}$ et $\frac{1}{y}$ est inférieure ou égale à leur moyenne arithmétique. Ceci
fournit $\sqrt{\frac{1}{x}.\frac{1}{y}}\leq\frac{1}{2}(\frac{1}{x}+\frac{1}{y})$ ou encore
$\frac{1}{g}\leq\frac{1}{h}$ et finalement 
\begin{center}
\shadowbox{$x\leq h\leq g\leq m\leq y$ où $\frac{1}{h}=\frac{1}{2}\left(\frac{1}{x}+\frac{1}{y}\right)$, $g=\sqrt{xy}$ et $m=\frac{x+y}{2}$.}
\end{center}
\end{enumerate}

\textbf{Remarque 1.} On a $h=\frac{2xy}{x+y}$, mais cette expression ne permet pas de comprendre que $\frac{1}{h}$ est
la moyenne arithmétique de $\frac{1}{x}$ et $\frac{1}{y}$.

\textbf{Remarque 2.} On peut visualiser l'inégalité entre moyenne arithmétique et géométrique.

Si $(ABC)$ est un triangle rectangle en $A$ et $A'$ est le pied de la hauteur issue de $A$, on sait que $AA'^2=A'B.A'C$. On
se sert de cette remarque pour construire $g$ et la comparer graphiquement à $m$.

On accolle deux segments de longueurs respectives $x$ et $y$. On construit alors un triangle rectangle d'hypothénuse
ce segment (de longueur $x+y$) noté [BC], tel que le troisième sommet $A$ ait une projection orthogonale $A'$ sur $(BC)$
vérifiant $BA'=x$ et $CA'=y$.


$$\includegraphics{../images/img005146-1}$$


La moyenne arithmétique de $x$ et $y$ est $m=\frac{x+y}{2}$, le rayon du cercle, et la moyenne géométrique de $x$
et $y$ est $g=\sqrt{xy}=\sqrt{A'B.A'C}=AA'$, la hauteur issue de $A$ du triangle $(ABC)$.
\fincorrection
\correction{005151}
Si l'un des réels $a$, $b$ ou $c$ est strictement plus grand que $1$, alors l'un au moins des
trois réels $a(1-b)$, $b(1-c)$, $c(1-a)$ est négatif (puisque $a$, $b$ et $c$ sont positifs) et donc inférieur ou égal à
$\frac{1}{4}$.

Sinon, les trois réels $a$, $b$ et $c$ sont dans $[0,1]$. Le produit des trois réels $a(1-b)$, $b(1-c)$ et $c(1-a)$ vaut

$$a(1-a)b(1-b)c(1-c).$$

Mais, pour $x\in[0,1]$, $x(1-x)$ est positif et d'autre part, $x(1-x)=-(x-\frac{1}{2})^2+\frac{1}{4}\leq\frac{1}{4}$.
Par suite,

$$a(1-a)b(1-b)c(1-c)\leq\frac{1}{4^3}.$$

Il est alors impossible que les trois réels $a(1-b)$, $b(1-c)$ et $c(1-a)$ soient strictement plus grand que
$\frac{1}{4}$, leur produit étant dans ce cas strictement plus grand que $\frac{1}{4^3}$.

On a montré dans tous les cas que l'un au moins des trois réels $a(1-b)$, $b(1-c)$ et $c(1-a)$ est inférieur ou égal à
$\frac{1}{4}$.
\fincorrection
\correction{005152}
\begin{enumerate}
\item  Soit $x\in\Rr$. Alors, $E(x)\leq x<E(x)+1$ puis $E(x)+1\leq x+1<(E(x)+1)+1$. Comme $E(x)+1\in\Zz$, on a
bien $E(x+1)=E(x)+1$.

\item  Soient $(x,y)\in\Rr^2$. On a $E(x)+E(y)\leq x+y$. Ainsi, $E(x)+E(y)$ est un entier relatif inférieur ou égal
à $x+y$. Comme $E(x+y)$ est le plus grand entier relatif inférieur ou égal à $x+y$, on a donc $E(x)+E(y)\leq E(x+y)$.

Améliorons. $E(x)\leq x<E(x)+1$ et $E(y)\leq y<E(y)+1$ fournit $E(x)+E(y)\leq x+y<E(x)+E(y)+2$ et donc $E(x+y)$ vaut,
suivant le cas, $E(x)+E(y)$ ou $E(x)+E(y)+1$ (et est dans tous les cas supérieur ou égal à $E(x)+E(y)$).

\item  Soit $(x,y)\in\Rr^2$. Posons $k=E(x)$ et $l=E(y)$.
\begin{itemize}
\item[\textbf{1er cas.}] Si $x\in[k,k+\frac{1}{2}[$ et $y\in[l,l+\frac{1}{2}[$, alors $x+y\in[k+l,k+l+1[$ et donc $E(x+y)=k+l$,
puis $E(x)+E(y)+E(x+y)=k+l+k+l=2k+2l$. D'autre part, $2x\in[2k,2k+1[$ et $2y\in[2l,2l+1[$. Par suite,
$E(2x)+E(2y)=2k+2l$. Dans ce cas, $E(x)+E(y)+E(x+y)=E(2x)+E(2y)$.
\item[\textbf{2ème cas.}] Si $x\in[k+\frac{1}{2},k+1[$ et $y\in[l,l+\frac{1}{2}[$, alors
$x+y\in[k+l+\frac{1}{2},k+l+\frac{3}{2}[$ et donc $E(x+y)=k+l$ ou $k+l+1$,puis $E(x)+E(y)+E(x+y)=2k+2l$ ou $2k+2l+1$.
D'autre part, $2x\in[2k+1,2k+2[$ et $2y\in[2l,2l+1[$. Par suite, $E(2x)+E(2y)=2k+2l+1$. Dans ce cas,
$E(x)+E(y)+E(x+y)\leq E(2x)+E(2y)$.
\item[\textbf{3ème cas.}] Si $x\in[k,k+\frac{1}{2}[$ et $y\in[l+\frac{1}{2},l+1[$, on a de
même $E(x)+E(y)+E(x+y)\leq E(2x)+E(2y)$.
\item[\textbf{4ème cas.}] Si $x\in[k+\frac{1}{2},k+1[$ et $y\in[l+\frac{1}{2},l+1[$, on a
$E(x)+E(y)+E(x+y)=2k+2l+2=E(2x)+E(2y)$.
\end{itemize}
Finalement, on a dans tous les cas $E(x)+E(y)+E(x+y)\leq E(2x)+E(2y)$.
\end{enumerate}
\fincorrection
\correction{005153}
 $p$ est déterminé par l'encadrement~:~$10^p\leq n<10^{p+1}$ qui s'écrit encore $p\leq\frac{\ln
n}{\ln 10}<p+1$. Par suite,
\begin{center}
\shadowbox{
$p=E(\mbox{log}_{10}(n)).$
}
\end{center}

\begin{center}
\shadowbox{
Le nombre de chiffres d'un entier $n$ en base $10$ est donc $E(\mbox{log}_{10}(n))+1$.
}
\end{center}
\fincorrection
\correction{005155}
\begin{enumerate}
\item  Par définition d'un entier, il y a $n$ entiers entre $1$ et $n$. Ensuite, pour tout entier naturel $k$, on a

$$1\leq k\leq x\Leftrightarrow 1\leq k\leq E(x).$$

Il y a donc $E(x)$ entiers entre $1$ et $x$.

\item  Il y a $n+1$ entiers entre $0$ et $n$ et $E(x)+1$ entiers entre $0$ et $x$.

\item  Les entiers naturels pairs sont les entiers de la forme $2k$, $k\in\Nn$. Or,

$$0\leq2k\leq x\Leftrightarrow 0\leq k\leq \frac{x}{2}.$$

Le nombre des entiers pairs compris entre $0$ et $x$ est encore le nombre des entiers $k$ compris au sens large entre
$0$ et $\frac{x}{2}$. D'après 2), il y a $E(\frac{x}{2})+1$ entiers pairs entre $0$ et $x$. De même, il y a
$E(\frac{x}{3})+1$ multiples de $3$ entre $0$ et $x$.

De même,

$$0\leq2k+1\leq x\Leftrightarrow-\frac{1}{2}\leq k\leq\frac{x-1}{2}\Leftrightarrow0\leq k\leq E(\frac{x-1}{2}).$$

Il y a donc $E(\frac{x-1}{2})+1=E(\frac{x+1}{2})$ entiers impairs entre $0$ et $x$.

\item  Il y a $E(\frac{x}{3})+1$ multiples de $3$ entre $0$ et $x$.

\item  Soient $n\in\Nn$ et $(x,y)\in\Nn^2$. On a

$$x+2y=n\Leftrightarrow x=n-2y.$$

Donc, $(x,y)$ est solution si et seulement si $y\in\Nn$ et $n-2y\in\Nn$ ou encore si et seulement si $0\leq2y\leq n$.
Il y a donc $E(\frac{n}{2})+1$ couples solutions.

\item  Si $x$ et $y$ sont respectivement le nombre de pièces de $10$ centimes d'euros et le nombre de pièces de $20$
centimes d'euros, le nombre cherché est le nombre de couples d'entiers naturels solutions de l'équation $10x+20y=1000$
qui s'écrit encore $x+2y=100$. D'après 5), il y a $E(\frac{100}{2})+1=51$ façons de payer $10$ euros avec des pièces de
$10$ et $20$ centimes d'euros.

\item  Soient $n\in\Nn$ et $(x,y)\in\Nn^2$. On a

$$2x+3y=n\Leftrightarrow x=\frac{n-3y}{2}.$$

Donc,

$$(x,y)\;\mbox{solution}\Leftrightarrow x=\frac{n-3y}{2}\;\mbox{et}\;y\in\Nn\;\mbox{et}\;n-3y\in2\Nn.$$

Maintenant, comme $n-3y=(n-y)-2y$ et que $2y$ est un entier pair, $n-3y$ est pair si et seulement si $n-y$ est pair ce
qui revient à dire que $y$ a la parité de $n$. Ainsi,

$$(x,y)\;\mbox{solution}\Leftrightarrow x=\frac{n-3y}{2}\;\mbox{et}\;y\in\Nn\;\mbox{et}\;0\leq y\leq
\frac{n}{3}\;\mbox{et}\;y\;\mbox{a la parité de}\;n.$$

\begin{itemize}
\item[\textbf{1er cas.}] Si $n$ est pair, le nombre de couples solutions est encore le nombre d'entiers pairs $y$
compris au sens large entre $0$ et $\frac{n}{3}$. Il y a $E(\frac{n}{6}))+1=E(\frac{n+6}{6})$ tels entiers.
\item[\textbf{2ème cas.}] Si $n$ est impair, le nombre de couples solutions est encore le nombre d'entiers impairs $y$
compris au sens large entre $0$ et $\frac{n}{3}$. Il y a  $E(\frac{\frac{n}{3}-1}{2}))+1=E(\frac{n+3}{6})$ tels
entiers.
\end{itemize}
Finalement, le nombre cherché est $E(\frac{n+6}{6})$ si $n$ est pair et $E(\frac{n+3}{6})$ si $n$ est impair.

\end{enumerate}

\fincorrection
\correction{005156}
\begin{enumerate}
\item  Par définition d'un entier, il y a $n$ entiers entre $1$ et $n$. Ensuite, pour tout entier naturel $k$, on a

$$1\leq k\leq x\Leftrightarrow 1\leq k\leq E(x).$$

Il y a donc $E(x)$ entiers entre $1$ et $x$.

\item  Il y a $n+1$ entiers entre $0$ et $n$ et $E(x)+1$ entiers entre $0$ et $x$.

\item  Les entiers naturels pairs sont les entiers de la forme $2k$, $k\in\Nn$. Or,

$$0\leq2k\leq x\Leftrightarrow 0\leq k\leq \frac{x}{2}.$$

Le nombre des entiers pairs compris entre $0$ et $x$ est encore le nombre des entiers $k$ compris au sens large entre
$0$ et $\frac{x}{2}$. D'après 2), il y a $E(\frac{x}{2})+1$ entiers pairs entre $0$ et $x$. De même, il y a
$E(\frac{x}{3})+1$ multiples de $3$ entre $0$ et $x$.

De même,

$$0\leq2k+1\leq x\Leftrightarrow-\frac{1}{2}\leq k\leq\frac{x-1}{2}\Leftrightarrow0\leq k\leq E(\frac{x-1}{2}).$$

Il y a donc $E(\frac{x-1}{2})+1=E(\frac{x+1}{2})$ entiers impairs entre $0$ et $x$.

\item  Il y a $E(\frac{x}{3})+1$ multiples de $3$ entre $0$ et $x$.

\item  Soient $n\in\Nn$ et $(x,y)\in\Nn^2$. On a

$$x+2y=n\Leftrightarrow x=n-2y.$$

Donc, $(x,y)$ est solution si et seulement si $y\in\Nn$ et $n-2y\in\Nn$ ou encore si et seulement si $0\leq2y\leq n$.
Il y a donc $E(\frac{n}{2})+1$ couples solutions.

\item  Si $x$ et $y$ sont respectivement le nombre de pièces de $10$ centimes d'euros et le nombre de pièces de $20$
centimes d'euros, le nombre cherché est le nombre de couples d'entiers naturels solutions de l'équation $10x+20y=1000$
qui s'écrit encore $x+2y=100$. D'après 5), il y a $E(\frac{100}{2})+1=51$ façons de payer $10$ euros avec des pièces de
$10$ et $20$ centimes d'euros.

\item  Soient $n\in\Nn$ et $(x,y)\in\Nn^2$. On a

$$2x+3y=n\Leftrightarrow x=\frac{n-3y}{2}.$$

Donc,

$$(x,y)\;\mbox{solution}\Leftrightarrow x=\frac{n-3y}{2}\;\mbox{et}\;y\in\Nn\;\mbox{et}\;n-3y\in2\Nn.$$

Maintenant, comme $n-3y=(n-y)-2y$ et que $2y$ est un entier pair, $n-3y$ est pair si et seulement si $n-y$ est pair ce
qui revient à dire que $y$ a la parité de $n$. Ainsi,

$$(x,y)\;\mbox{solution}\Leftrightarrow x=\frac{n-3y}{2}\;\mbox{et}\;y\in\Nn\;\mbox{et}\;0\leq y\leq
\frac{n}{3}\;\mbox{et}\;y\;\mbox{a la parité de}\;n.$$

\begin{itemize}
\item[\textbf{1er cas.}] Si $n$ est pair, le nombre de couples solutions est encore le nombre d'entiers pairs $y$
compris au sens large entre $0$ et $\frac{n}{3}$. Il y a $E(\frac{n}{6}))+1=E(\frac{n+6}{6})$ tels entiers.
\item[\textbf{2ème cas.}] Si $n$ est impair, le nombre de couples solutions est encore le nombre d'entiers impairs $y$
compris au sens large entre $0$ et $\frac{n}{3}$. Il y a  $E(\frac{\frac{n}{3}-1}{2}))+1=E(\frac{n+3}{6})$ tels
entiers.
\end{itemize}
Finalement, le nombre cherché est $E(\frac{n+6}{6})$ si $n$ est pair et $E(\frac{n+3}{6})$ si $n$ est impair.

\end{enumerate}

\fincorrection
\correction{005159}
Soient $n\in\Nn^*$ et $x\in\Rr$.

\begin{align*}
E(x)\leq x<E(x)+1&\Rightarrow nE(x)\leq nx<nE(x)+n\Rightarrow nE(x)\leq E(nx)<nE(x)+n
\Rightarrow E(x)\leq\frac{E(nx)}{n}<E(x)+1\\
 &\Rightarrow E(\frac{E(nx)}{n})=E(x).
\end{align*}

\fincorrection
\correction{005160}
Soient $n\in\Nn^*$ et $(x_1,x_2,...,x_n )\in[-1,1]^n$ tels que $x_1+x_2+...+x_n=0$.

On écrit

$$(x_1+2x_2+...+nx_n)=(x_1+x_2+...+x_n)+(x_2+x_3+....+x_n)+(x_3+...+x_n)+...+(x_{n-1}+x_n)+x_n,$$

avec $x_1+...+x_n=0$ et donc $x_2+...+x_n=-x_1$ ...

\begin{itemize}
\item[\textbf{1er cas.}] Si $n=2p$ est pair, alors $\frac{n^2}{4}=p^2$ et
donc, $E(\frac{n^2}{4})=p^2=\frac{n^2}{4}$. Dans ce cas, on peut écrire

\begin{align*}
|x_1+2x_2+...+nx_n|&\leq|x_1+x_2+...+x_{2p}|+|x_2+....+x_{2p}|+...+|x_p+...+x_{2p}|\\
 &\;+|x_{p+1}+...+x_{2p}|...+|x_{2p-1}+
x_{2p}|+|x_{2p}|\\
 &=0+|-x_1|+|-x_1-x_2|+...+|-x_{1}+...-x_{p-1}|\\
 &+|x_{p+1}+...+x_{2p}|...+|x_{2p-1}+x_{2p}|+|x_{2p}|\\
 &\leq0+1+2+...+(p-1)+p+(p-1)+...+1=2\frac{p(p-1)}{2}+p=p^2=E(\frac{n^2}{4})
\end{align*}

\item[\textbf{2ème cas.}] Si $n=2p+1$ est impair, alors $\frac{n^2}{4}=p^2+p+\frac{1}{4}$ et
donc, $E(\frac{n^2}{4})=p^2+p=\frac{n^2-1}{4}$. Dans ce cas, on peut écrire

\begin{align*}
|x_1+2x_2+...+nx_n|&\leq|x_1+x_2+...+x_{2p+1}|+...+|x_{p+1}+...+x_{2p+1}|\\
 &+|x_{p+2}+...+x_{2p+1}|...+|x_{2p+1}|\\
 &=0+|-x_1|+|-x_1-x_2|+...+|-x_{1}+...-x_{p}|\\
 &+|x_{p+2}+...+x_{2p+1}|...+|x_{2p+1}|\\
 &\leq0+1+2+...+(p-1)+p+p+(p-1)+...+1=2\frac{p(p+1)}{2}=p^2+p=E(\frac{n^2}{4})
\end{align*}

\end{itemize}

Dans tous les cas, on a montré que
\begin{center}
\shadowbox{
$\forall n\in\Nn^*,\;|x_1+2x_2+...+nx_n|\leq E(\frac{n^2}{4}).$
}
\end{center}
\fincorrection
\correction{005215}
Pour $n=1$, $\sum_{k=0}^{2n-1}\frac{(-1)^k}{k+1}=1-\frac{1}{2}=\frac{1}{2}$ et $\sum_{k=n+1}^{2n}\frac{1}{k}=\frac{1}{2}$. L'identité proposée est donc vraie pour $n=1$.

Soit $n\geq1$. Supposons que $\sum_{k=0}^{2n-1}\frac{(-1)^k}{k+1}=\sum_{k=n+1}^{2n}\frac{1}{k}$.

On a alors
  
\begin{align*}
\sum_{k=0}^{2(n+1)-1}\frac{(-1)^k}{k+1}&=\sum_{k=0}^{2n-1}\frac{(-1)^k}{k+1}+\frac{1}{2n+1}-\frac{1}{2n+2}
=\sum_{k=n+1}^{2n}\frac{1}{k}+\frac{1}{2n+1}-\frac{1}{2(n+1)}\\
 &=\frac{1}{n+1}+\sum_{k=n+2}^{2n+1}\frac{1}{k}+\frac{1}{2(n+1)}
=\sum_{k=n+2}^{2n+1}\frac{1}{k}+\frac{1}{2n+2}=\sum_{k=n+2}^{2(n+1)}\frac{1}{k}
\end{align*}
On a montré par récurrence que $\forall n\geq1,\;\sum_{k=0}^{2n-1}\frac{(-1)^k}{k+1}=\sum_{k=n+1}^{2n}\frac{1}{k}$ (identité de \textsc{Catalan}).

\fincorrection
\correction{005216}
\begin{enumerate}
 \item  Si les $b_k$ sont tous nuls, l'inégalité est claire.
Sinon, pour $x$ réel, posons

\begin{center}
$f(x)=\sum_{k=1}^{n}(a_k+xb_k)^2=\left(\sum_{k=1}^{n}b_k^2\right)x^2+2\left(\sum_{k=1}^{n}a_kb_k\right)x+\sum_{k=1}^{n}a_k^2$.
\end{center}
$f$ est un trinôme du second degré de signe constant sur $\Rr$. Son discriminant réduit est donc négatif ou nul ce qui fournit~:

$$0\geq\Delta'=\left(\sum_{k=1}^{n}a_kb_k\right)^2-\left(\sum_{k=1}^{n}a_k^2\right)\left(\sum_{k=1}^{n}b_k^2\right),$$ 
ou encore $\left|\sum_{k=1}^{n}a_kb_k\right|\leq\sqrt{\sum_{k=1}^{n}a_k^2}\sqrt{\sum_{k=1}^{n}b_k^2}$, qui est l'inégalité de \textsc{Cauchy}-\textsc{Schwarz}.
 \item 

\begin{align*}
\sum_{k=1}^{n}(a_k+b_k)^2&=\sum_{k=1}^{n}a_k^2+2\sum_{k=1}^{n}a_kb_k+\sum_{k=1}^{n}b_k^2
\leq\sum_{k=1}^{n}a_k^2+2\left|\sum_{k=1}^{n}a_kb_k\right|+\sum_{k=1}^{n}b_k^2\\
 &\leq\sum_{k=1}^{n}a_k^2+2\sqrt{\sum_{k=1}^{n}a_k^2}\sqrt{\sum_{k=1}^{n}b_k^2}+\sum_{k=1}^{n}b_k^2
 \quad(\mbox{\textsc{Cauchy}-\textsc{Schwarz}})\\
 &=\left(\sqrt{\sum_{k=1}^{n}a_k^2}+\sqrt{\sum_{k=1}^{n}b_k^2}\right)^2
\end{align*}
et donc, $\sqrt{\sum_{k=1}^{n}(a_k+b_k)^2}\leq\sqrt{\sum_{k=1}^{n}a_k^2}+\sqrt{\sum_{k=1}^{n}b_k^2}$, qui est l'inégalité de \textsc{Minkowski}.

\end{enumerate}
\fincorrection
\correction{005217}
Pour $x\geq1$, $x+2\sqrt{x-1}=x-1+2\sqrt{x-1}+1=(\sqrt{x-1}+1)^2\geq0$. De même, $x-2\sqrt{x-1}=(\sqrt{x-1}-1)^2\geq0$.
Donc, si on pose $f(x)=\sqrt{x+2\sqrt{x-1}}+\sqrt{x-2\sqrt{x-1}}$, $f(x)$ existe si et seulement $x\geq1$ et pour $x\geq1$, $f(x)=\sqrt{x-1}+1+|\sqrt{x-1}-1|$. Par suite, 

$$f(x)=1\Leftrightarrow\sqrt{x-1}+|\sqrt{x-1}-1|=0\Leftrightarrow\sqrt{x-1}=0\;\mbox{et}\;\sqrt{x-1}-1=0\Leftrightarrow\sqrt{x-1}=0\;\mbox{et}\;\sqrt{x-1}=1,$$
ce qui est impossible. L'équation proposée n'a pas de solution.
\fincorrection
\correction{005218}
\begin{enumerate}
 \item  Soit $G$ un sous groupe non nul de $(\Rr,+)$ ($\{0\}=0.\Zz$ est du type voulu).
Il existe dans $G$ un réel non nul $x_0$.
Puisque $G$ est un sous groupe de $(\Rr,+)$, le réel $-x_0$ est aussi dans $G$ et l'un des deux réels $x_0$ ou $-x_0$ est strictement positif. Soit alors $A=G\cap]0,+\infty[$.
D'après ce qui précède, $A$ est une partie non vide et minorée (par $0$) de $\Rr$. $A$ admet donc une borne inférieure que l'on note $a$.

\textbf{1er cas.} Si $a=0$, montrons dans ce cas que $G$ est dense dans $\Rr$ (c'est par exemple le cas de $(\Qq,+)$).
Soient $x$ un réel et $\varepsilon$ un réel strictement positif.
Puisque $\mbox{inf }A=\mbox{inf}(G\cap]0,+\infty[)=0$, il existe dans $G$ un élément $g$ tel que $0<g<\varepsilon$. 
Puis il existe un entier relatif $n$ tel que $ng\leq x-\varepsilon<(n+1)g$ à savoir $n=E\left(\frac{x-\varepsilon}{g}\right)$.
Soit $y=(n+1)g$. D'une part, $y$ est dans $G$ (si $n+1=0$, $(n+1)g=0\in G$, si $n+1>0$, $(n+1)g=g+g+...+g\in G$ et si $n+1<0$, $(n+1)g=-(-(n+1)g)\in G$) et d'autre part

$$x-\varepsilon<(n+1)g=ng+g<x-\varepsilon+\varepsilon=x.$$
On a montré que $\forall x\in\Rr,\;\forall\varepsilon>0,\;\exists y\in G/\;x-\varepsilon<y<x$ et donc 
\begin{center}
\shadowbox{
si $G\neq\{0\}$ et si $\text{inf}\left(G\cap]0,+\infty[\right)=0$, $G$ est dense dans $\Rr$.
}
\end{center}
\textbf{2ème cas.} Si $a>0$, montrons dans ce cas que $G=a\Zz$. Pour cela, montrons tout d'abord que $a$ est dans $G$.
Mais si $a$ n'est pas élément de $G$, par définition de $a$, il existe un réel $x$ dans $G\cap]a,2a[$ puis il existe un réel $y$ dans $G\cap]a,x[$. Le réel $x-y$ est alors dans $G\cap]0,a[$ ce qui est impossible. Donc $a$ est élément de $G$.
Montrons alors que $G=a\Zz$. Puisque $a$ est dans $G$, $G$ contient encore $a+a=2a$, puis $a+a+a=3a$ et plus généralement tous les $na$, $n\in\Nn^*$. Puisque $G$ contient aussi les opposés de ces nombres et également $0=0\times a$, $G$ contient finalement tous les $na$, $n\in\Zz$. On a ainsi montré que $a\Zz\subset G$.
Réciproquement, soit $x$ un élément de $G$ et $n=E\left(\frac{x}{a}\right)(\in\Zz)$. Alors, $n\leq\frac{x}{a}<n+1$ puis $0\leq x-na<a$. 
Or, $x$ est dans $G$ et $na$ est dans $G$. Donc, $x-na$ est dans $G\cap[0,a[=\{0\}$, puis $x=na\in a\Zz$. On a ainsi montré l'inclusion contraire et donc $G=a\Zz$.

\begin{center}
\shadowbox{
si $\text{inf}\left(G\cap]0,+\infty[\right)=a>0$, $G=a\Zz$.
}
\end{center}

 \item  Soit $G=\{a+b\sqrt{2},\;(a,b)\in\Zz^2\}$. On vérifie aisément que $G$ est un sous-groupe de $(\Rr,+)$. Maintenant, la formule du binôme de \textsc{Newton} montre que, pour chaque entier naturel $n$,

\begin{center}
$(\sqrt{2}-1)^n\in G\cap]0,+\infty[$.
\end{center}
Or, $0<\sqrt{2}-1<1$ et donc $\lim_{n\rightarrow +\infty}(\sqrt{2}-1)^n=0$. Ceci montre que $\mbox{inf}(G\cap]0;+\infty[)=0$ et donc que $G$ est dense dans $\Rr$.
 \item  
  \begin{enumerate}
  \item  Soit $f$ une application de $\Rr$ dans $\Rr$ et $G_f=\{T\in\Rr/\;\forall x\in\Rr,\;f(x+T)=f(x)\}$.
$0$ est élément de $G_f$ (et c'est même le seul élément de $G_f$ si $f$ n'est pas périodique) et donc $G\neq\varnothing$.
De plus, si $T$ et $T'$ sont deux éléments de $G$ alors, pour $x$ réel donné~:

$$f(x+(T-T'))=f((x-T')+T)=f(x-T')=f(x-T'+T')=f(x),$$
et $T-T'$ est encore un élément de $G$. On a montré que

\begin{center}
\shadowbox{
$G_f$ est un sous groupe de $(\Rr,+)$.
}
\end{center}
  \item Soit $f$ une application de $\Rr$ dans $\Rr$ admettant $1$ et $\sqrt{2}$ pour périodes. $G_f$ contient encore tous les nombres de la forme $a+b\sqrt{2},\;(a,b)\in\Zz^2$ et est donc dense dans $\Rr$.
Montrons que si de plus $f$ est continue sur $\Rr$, $f$ est constante.
Soit $x$ un réel quelconque. On va montrer que $f(x)=f(0)$.
Remarque préliminaire : soit $T$ une période strictement positive de $f$.
Il existe un entier relatif $p$ tel que $pT\leq x<(p+1)T$ à savoir $p=E\left(\frac{x}{T}\right)$. On a alors $f(x)=f(x-pT)$ avec $0\leq x-pT<T$.
Soit alors $n\in\Nn^*$. Puisque $G_f$ est dense dans $\Rr$, il existe dans $G_f$ un réel $T_n$ tel que $0<T_n<\frac{1}{n}$ (ce qui implique que $\lim_{n\rightarrow +\infty}T_n=0$).
Mais alors, puisque $0<x-E\left(\frac{x}{T_n}\right)T_n<T_n$, on a aussi $\lim_{n\rightarrow +\infty}x-E\left(\frac{x}{T_n}\right)T_n=0$.
Maintenant, la suite $\left(f\left(x-E\left(\frac{x}{T_n}\right)T_n\right)\right)_{n\in\Nn^*}$ est constante égale à $f(x)$) et donc convergente vers $f(x)$. On en déduit que

\begin{align*}
f(x)&=\lim_{n\rightarrow +\infty}f\left(x-E\left(\frac{x}{T_n}\right)T_n\right)= f\left(\lim_{n\rightarrow +\infty}\left(x-E\left(\frac{x}{T_n}\right)T_n\right)\right)\quad(\mbox{par continuité de}\;f\;\text{en}\;0)\\
 &=f(0)
\end{align*}
ce qu'il fallait démontrer.
  \end{enumerate}
\end{enumerate}
\fincorrection
\correction{005219}
Soient $x$ un réel et $\varepsilon$ un réel strictement positif. On a $\sqrt[3]{x}<\sqrt[3]{x+\varepsilon}$. Puisque $\Qq$ est dense dans $\Rr$, il existe un rationnel $r$ tel que $\sqrt[3]{x}<r<\sqrt[3]{x+\varepsilon}$ et donc tel que $x<r^3<x+\varepsilon$, par stricte croissance de la fonction $t\mapsto t^3$ sur $\Rr$. On a montré que

\begin{center}
\shadowbox{
$\{r^3,\;r\in\Qq\}$ est dense dans $\Rr$.
}
\end{center}
\fincorrection
\correction{005982}
\begin{enumerate}
 \item Par définition est l'unique nombre $E(x) \in \Zz$ tel que 
 $$E(x) \le x < E(x)+1.$$

 \item Pour le réel $kx$, ($k=1,\ldots,n$) l'encadrement précédent s'écrit $E(kx) \le kx < E(kx)+1$.
 Ces deux inégalités s'écrivent aussi $E(kx) \le kx$ et $E(kx) > kx - 1$, d'où l'encadrement
 $kx-1 < E(kx) \le kx$. On somme cet encadrement, $k$ variant de $1$ à $n$, pour obtenir :
 $$\sum_{k=1}^n (kx-1) < \sum_{k=1}^n E(kx) \le \sum_{k=1}^n kx.$$
Ce qui donne 
$$ x \cdot \sum_{k=1}^n k \quad - n < n^2 \cdot u_n \le x \cdot  \sum_{k=1}^n k.$$

 \item On se rappelle que $\sum_{k=1}^n k = \frac{n(n+1)}{2}$ donc 
 nous obtenons l'encadrement :
 $$ x\cdot  \frac{1}{n^2} \cdot \frac{n(n+1)}{2} -   \frac{1}{n} < u_n \le x \cdot  \frac{1}{n^2} \cdot  \frac{n(n+1)}{2}.$$
 $\frac{1}{n^2} \cdot \frac{n(n+1)}{2}$ tend vers $\frac 12$, donc par le théorème des gendarmes $(u_n)$ tend vers $\frac x2$.

 \item Chaque $u_n$ est un rationnel (le numérateur et le dénominateur sont des entiers).
Comme la suite $(u_n)$ tend vers $\frac x 2$, alors la suite de rationnels $(2u_n)$ tend vers $x$.
Chaque réel $x\in \Rr$ peut être approché d'aussi près que l'on veut par des rationnels, donc 
$\Qq$ est dense dans $\Rr$. 
\end{enumerate}
\fincorrection
\correction{007169} 
On applique l'inégalité arithmético-géométrique à $\frac{a^2}{4}$ et à $b^2$, ce qui donne
\[
\frac{a^2}{4}+ b^2
\geq
2\sqrt{\frac{a^2}{4}\cdot b^2} = |ab| \geq ab.
\]

On peut généraliser en remplaçant $4$ par un réel strictement positif. Une formulation est la suivante. Pour $a$ et $b$ réels et $\lambda>0$, on a 
\[
2|ab|  \leq \frac{a^2}{\lambda}+ \lambda b^2.
\]
\fincorrection  
\correction{007170} 
Notons $a$ et $b$ les nombres de l'énoncé. On a $ab=100$. L'inégalité arithmético-géométrique fournit :
\[ \frac{a+b}{2}\geq \sqrt{ab}=10\]
avec égalité ssi $a=b$, donc la somme est supérieure à $20$, avec égalité ssi $a=b=10$.
\fincorrection  
\correction{007171} 
Notons $a_1$, ..., $a_n$ les nombres de l'énoncé. On a 
\[
a_1 a_2 \dots a_n =  \prod_{i=1}^n a_i=1.\]
L'inégalité arithmético-géométrique fournit :
\[ \frac{\sum_{i=1}^n a_i}{n} = \frac{a_1+..+a_n}{n}\geq \sqrt[n]{\prod_{i=1}^n a_i} = \sqrt[n]{1}=1\]
avec égalité ssi tous les $a_i$ sont égaux. On en déduit que la somme est supérieure à $n$, avec égalité ssi tous les réels sont égaux (et donc égaux à  $1$).
\fincorrection  
\correction{007172} 
Il s'agit de savoir laquelle des deux quantités
\[ 
a^3+b^3+c^3
\text{ et }
3abc
\]
est la plus grande.

Or, en appliquant l'inégalité arithmético-géométrique à $a^3$, $b^3$ et $c^3$, on obtient directement:
\[
\frac13\left(  a^3+b^3+c^3\right)
\geq 
\sqrt[3]{a^3b^3c^3} = abc.
\]
Il est donc préférable d'acheter les trois cubes.
\fincorrection  
\correction{007173} 
On applique l'inégalité arithmético-géométrique à $\frac{a^2}{bc}$, $\frac{b^2}{ca}$ et $\frac{c^2}{ab}$ ce qui donne
\[ \frac{a^2}{bc}+\frac{b^2}{ca}+\frac{c^2}{ab}
\geq3\sqrt[3]{\frac{a^2}{bc}\frac{b^2}{ca}\frac{c^2}{ab}}
=3
\]
\fincorrection  
\correction{007174} 
On applique l'inégalité arithmético-géométrique aux réels $\frac{a_i}{b_i}$ ce qui donne
\[ \frac{a_1}{b_1} + ... + \frac{a_n}{b_n}
\geq n\sqrt[n]{\frac{a_1}{b_1} \cdot ... \cdot \frac{a_n}{b_n}}
=n \sqrt[n]{1}=n.
\]
\fincorrection  
\correction{007175} 
On applique l'inégalité arithmético-géométrique à chacun des deux facteurs ce qui donne:
\[
(1+a^2)(1+b^2) \geq 
(2\sqrt{a^2})(2\sqrt{b^2})
= 4|ab| \geq  4ab.
\]

Remarque : on aurait également pu développer le membre de gauche et minorer par une seule utilisation de l'inégalité arithmético-géométrique à quatre variables:
\[
(1+a^2)(1+b^2) 
=1+a^2+b^2+a^2b^2
\geq  
4\sqrt[4]{a^2b^2a^2b^2} =4\sqrt[4]{a^4b^4}
= 4|ab| \geq 4ab.
\]
\fincorrection  
\correction{007176} 
On applique l'inégalité arithmético-géométrique à chacun des deux facteurs ce qui donne:
\[
a^2b+b^2c+c^2a
\geq
3\sqrt[3]{a^2b \cdot b^2c \cdot c^2a} = 3abc
\]
et
\[
ab^2+bc^2+ca^2
 \geq 
 3\sqrt[3]{ab^2 \cdot bc^2 \cdot ca^2} = 3abc.
\]
En multipliant, on obtient le résultat.
\fincorrection  
\correction{007177} 
On peut essayer d'appliquer l'inégalité arithmético-géométrique à chaque facteur. Ceci donne
\[ 
\left(1+\frac1a\right)\left(1+\frac1b\right)
\geq 
\left(\frac{2}{\sqrt a}\right)\left(\frac{2}{\sqrt b}\right)
= \frac{4}{\sqrt{ab}},
\]
ce qui minore la quantité par $1$ après une deuxième utilisation de l'inégalité arithmético-géométrique sur le dénominateur et utilisation de $a+b=8$, mais cette dernière minoration est évidente vu la forme initiale de l'expression : les deux facteurs sont supérieurs à $1$.

(Remarque : lors de la première utilisation de l'inégalité arithmético-géométrique, il y avait égalité ssi $a=1$ et $b=1$ ce qui est impossible vu l'énoncé. L'inégalité est donc toujours stricte, ce qui indique que la minoration n'est sans doute pas très précise.)

Commençons donc plutôt par développer la quantité à minorer. On a 
\[ \left(1+\frac1a\right)\left(1+\frac1b\right)
=
\frac{1+a+b+ab}{ab} = \frac{9+ab}{ab} = 1+\frac{9}{ab}.\]
Il s'agit donc de majorer le produit $ab$. L'inégalité arithmético-géométrique donne $\sqrt{ab} \leq \frac{a+b}{2}=4$, donc $ab\leq 16$. On en déduit que $\frac{1}{ab} \geq \frac{1}{16}$ et donc que 
\[
\left(1+\frac1a\right)\left(1+\frac1b\right) \geq 1+\frac{9}{16} = \frac{25}{16},
\]
avec égalité ssi $a=b=4$.

Remarque : il est préférable d'utiliser les contraintes (ici $a+b=8$) le plus tôt possible dans les majorations ou minorations successives, pour gagner en précision.
\fincorrection  
\correction{007178} 
En développant, l'inégalité est équivalente à 
\[
a^2d^2+b^2c^2 \geq 2abcd.
\]
En appliquant l'inégalité arithmético-géométrique à $a^2d^2$ et $b^2c^2$, on obtient:
\[
a^2d^2+b^2c^2 \geq 
2\sqrt{a^2d^2b^2c^2}
=2|abcd|
\geq 2abcd.
\]

Deuxième solution : en fait, on a l'identité remarquable
\[ (a^2+b^2)(c^2+d^2)=  (ac+bd)^2 + (ad-bc)^2.\]

Remarque : majorer chacun des deux facteurs dans le membre de gauche donne juste 
\[ (a^2+b^2)(c^2+d^2) \geq 4|ab|\cdot|cd|,\]
ce qui ne permet pas de conclure puisque par ailleurs on a également 
\[ (ac+bd)^2 \geq 4|acbd|.\]
\fincorrection  
\correction{007179} 
L'équation est équivalente à 
\[ 2^x +\frac{1}{2^x} = 2-x^2.\]
Or, par inégalité arithmético-géométrique, on a 
\[ 2^x +\frac{1}{2^x} \geq 2\sqrt{2^x\cdot \frac{1}{2^x}} \geq 2,\]
avec égalité ssi $2^x = \frac{1}{2^x}$ c'est-à-dire ssi $x=0$.

D'autre part, $2-x^2\geq 2$ avec égalité ssi $x=0$ là aussi.

On en déduit que l'équation admet bien une solution, unique, égale à $0$.
\fincorrection  
\correction{007180} 
Si $(x,y,z)$ est une solution, alors en sommant les trois équation on obtient
\[ 4x + \frac{18}{y}
+2y + \frac{9}{z}
+9z + \frac{16}{x}
=14+15+17 = 46.
\]
D'autre part, par inégalité arithmético-géométrique, on a 

\[
4x+\frac{16}{x} \geq 2\sqrt{4x\cdot \frac{16}{x}} = 16,\quad
2y+\frac{18}{y} \geq 12,\quad
9z+\frac{9}{z} \geq 18,
\]
d'où 
\[
4x+\frac{16}{x} +
2y+\frac{18}{y} +
9z+\frac{9}{z} \geq 46,
\]
avec égalité ssi chacune des trois inégalité sont des égalités donc ssi $(x,y,z)=(2,3,1)$.

Le système d'équations admet donc une unique solution, $(2,3,1)$.
\fincorrection  
\correction{007181} 
Appliquer l'inégalité arithmético-géométrique à $2a^3$ et à $b^3$ ne semble pas donner le résultat.

On a $2a^3+b^3 = a^3+a^3+b^3$. Appliquons l'inégalité arithmético-géométrique à trois variables. On obtient:
\[
2a^3+b^3 =a^3+a^3+b^3 \geq 3\sqrt[3]{a^3a^3b^3} = 3a^2b.
\]

\fincorrection  
\correction{007182} 
Partir du membre de gauche et appliquer l'inégalité arithmético-géométrique à trois variables ne donne pas immédiatement le résultat. 
On peut par contre essayer de majorer séparément chacun des trois termes du membre de droite :
\[ a^2b =\sqrt[3]{a^3a^3b^3} \leq \frac{2a^3+b^3}{3},\]
\[ b^2c =\sqrt[3]{b^3b^3c^3} \leq \frac{2b^3+c^3}{3},\]
\[ c^2a =\sqrt[3]{c^3c^3a^3} \leq \frac{2c^3+a^3}{3},\]
ce qui donne le résultat en sommant les trois inégalités.

\fincorrection  
\correction{007183} 
Partir du membre de gauche et appliquer l'inégalité arithmético-géométrique à trois variables ne donne pas immédiatement le résultat. Si on développe le membre de droite, on obtient
\[ a^2bc+b^2ca+c^2ab,\]
que l'on peut essayer de minorer par trois utilisations indépendantes de l'inégalité arithmético-géométrique à quatre variables:
\[ a^2bc = \sqrt[4]{a^4a^4b^4c^4} \leq \frac14\left(a^4+a^4+b^4+c^4\right) = \frac14\left(2a^4+b^4+c^4\right). \]
De même, 
\[ b^2ac\leq \frac14\left(a^4+2b^4+c^4\right)
\text{ et }
c^2ab \leq \frac14\left(a^4+b^4+2c^4\right).\]
En sommant ces trois inégalités, on obtient le résultat.
\fincorrection  
\correction{007184} 
L'inégalité arithmético-géométrique appliquée à $\frac{a^2}{b^2}+\frac{b^2}{c^2}$ uniquement donne:
\[
\frac{a^2}{b^2}+\frac{b^2}{c^2}
\geq 
2\sqrt{\frac{a^2}{c^2}}=2\left|\frac{a}{c}\right| \geq 2\frac{a}{c}.
\]
On obtient de même
\[
\frac{b^2}{c^2}+\frac{c^2}{a^2}
\geq 
\frac{b}{a}
\text{ et }
\frac{c^2}{a^2}+\frac{a^2}{b^2}
\geq 
\frac{c}{b}.
\]
En sommant ces trois inégalités, on obtient le résultat.
\fincorrection  
\correction{007185} 
En appliquant l'inégalité arithmético-géométrique, on voit que les $a$ et $c$ se simplifient mais pas $b$ et $b+c$. L'idée est alors de modifier la forme de l'inégalité pour obtenir la simplification. Or, l'inégalité est équivalente à 
\[ \frac{c}{a} + \frac{a}{b+c} +\frac{b}{c}+1 \geq 3,
\]
c'est-à-dire à 
\[ \frac{c}{a} + \frac{a}{b+c} +\frac{b+c}{c} \geq 3.
\]
Sous cette forme, l'inégalité arithmético-géométrique donne le résultat:
\[ \frac{c}{a} + \frac{a}{b+c} +\frac{b+c}{c} \geq 
3\sqrt[3]{\frac{c}{a}  \frac{a}{b+c} \frac{b+c}{c}}
=3\sqrt[3]{1}=3.
\]
\fincorrection  
\nocorrection
\correction{000505}
\begin{enumerate}
  \item Vrai. Toute sous-suite d'une suite convergente est convergente
et admet la m\^eme limite (c'est un résultat du cours).
  \item Faux. Un contre-exemple est la suite $(u_n)_n$ d\'efinie
par $u_n = (-1)^n$. Alors $(u_{2n})_n$ est la suite constante
(donc convergente) de valeur $1$, et $(u_{2n+1})_n$ est constante
de valeur $-1$. Cependant la suite $(u_n)_n$ n'est pas
convergente.
  \item Vrai.
La convergence de la suite $(u_n)_n$ vers $\ell$, que nous
souhaitons d\'emontrer, s'\'ecrit :
$$ \forall \epsilon > 0\  \ \exists N \in \Nn \ \text{\  tel que\ \  }
(n \geqslant N \Rightarrow |u_n-\ell| < \epsilon).$$ 
Fixons $\epsilon >
0$. Comme, par hypoth\`ese, la suite $(u_{2p})_p$  converge vers
$\ell$ alors il existe $N_1$ tel
$$ 2p \geqslant N_1 \Rightarrow |u_{2p}-\ell| < \epsilon.$$
Et de m\^eme, pour la suite $(u_{2p+1})_p$ il existe $N_2$ tel que
$$ 2p+1 \geqslant N_2 \Rightarrow |u_{2p+1}-\ell| < \epsilon.$$
Soit $N = \max(N_1,N_2)$, alors
$$n \geqslant N\Rightarrow |u_{n}-\ell| < \epsilon.$$
Ce qui prouve la convergence de $(u_n)_n$ vers $\ell$.
\end{enumerate}
\fincorrection
\nocorrection
\nocorrection
\nocorrection
\nocorrection
\nocorrection
\nocorrection
\correction{000518}
\begin{enumerate}
\item Suite non convergente car non bornée.
\item Suite convergente vers $0$.
\item Suite non convergente car la sous-suite $u_{2p} = 1+\frac {1}{2p}$ est
toujours plus grande que $1$. Alors que la sous-suite $u_{2p+1} =
-1+\frac {1}{2p+1}$ est toujours plus petite que $0$.
\end{enumerate}
\fincorrection
\correction{000519}
Soit $(u_n)$ une suite d'entiers qui converge vers $\ell \in
\Rr$.
Dans l'intervalle $I = ] \ell - \frac12, \ell +\frac12[$ de
longueur $1$, il existe au plus un \'el\'ement de $\Nn$. Donc $I \cap
\Nn$ est soit vide soit un singleton $\{a \}$.

La convergence de  $(u_n)$ s'\'ecrit :
$$ \forall \epsilon > 0 \ \ \exists N \in \Nn \text{\  \ tel que\ \  }
(n \geqslant N \Rightarrow |u_n-\ell| < \epsilon).$$ Fixons $\epsilon =
\frac 12$, nous obtenons un $N$ correspondant. Et pour $n \geqslant N$,
$u_n \in I$. Mais de plus $u_n$ est un entier, donc
  $$ n \geqslant N \Rightarrow u_n \in I \cap \Nn.$$
En cons\'equent, $I\cap \Nn$ n'est pas vide (par exemple $u_N$ en
est un \'el\'ement) donc $I \cap \Nn = \{ a \}$. L'implication
pr\'ec\'edente s'\'ecrit maintenant :
$$ n \geqslant N \Rightarrow u_n = a.$$
Donc la suite $(u_n)$ est stationnaire (au moins) \`a partir de
$N$. En prime, elle est bien \'evidemment convergente vers $\ell = a
\in \Nn$.
\fincorrection
\correction{000520}
\begin{enumerate}
\item La fonction $t \mapsto \frac 1 t$ est d\'ecroissante
sur $[n,n+1]$ donc
$$\frac{1}{n+1} \leqslant \int_n^{n+1} \frac{dt}{t} \leqslant \frac 1n$$
(C'est un encadrement de l'aire de l'ensemble des points $(x,y)$
du plan tels que $x\in[n,n+1]$ et $0\leqslant y\leqslant 1/x$ par l'aire de
deux rectangles.) Par calcul de l'intégrale nous obtenons l'in\'egalit\'e :
$$\frac{1}{n+1} \leqslant \ln(n+1)-\ln(n)  \leqslant \frac 1n.$$
\item $H_n = \frac1n+\frac{1}{n-1}+\cdots +\frac12+1$, nous majorons chaque terme de cette somme en utilisant l'in\'egalit\'e $\frac1k \leqslant \ln(k)-\ln (k-1)$ obtenue pr\'ec\'edemment : nous obtenons
$H_n \leqslant \ln(n)-\ln (n-1) + \ln(n-1)-\ln (n-2)+\cdots-\ln(2) +\ln (2) - \ln
(1) + 1$. Cette somme est t\'elescopique (la plupart des termes
s'\'eliminent et en plus $\ln (1) =0$) et donne $H_n \leqslant \ln (n) + 1$.

L'autre in\'egalit\'e  s'obtient de la fa\c{c}on similaire en
utilisant l'in\'egalit\'e $ \ln(k+1)-\ln(k) \leqslant \frac{1}{k}$ .
\item Comme $H_n \geqslant \ln (n+1)$ et que $\ln(n+1) \rightarrow +\infty$ quand $n\rightarrow +\infty$ alors $H_n \rightarrow +\infty$ quand $n\rightarrow +\infty$.
\item $u_{n+1}-u_n = H_{n+1}-H_n - \ln(n+1)+\ln(n) = \frac{1}{n+1}-(\ln (n+1)-\ln (n))\leqslant 0$ d'apr\`es la premi\`ere question. Donc $u_{n+1}-u_n  \leqslant 0$. Ainsi $u_{n+1} \leqslant u_n$ et la suite
$(u_n)$ est d\'ecroissante.

Enfin comme $H_n \geqslant \ln(n+1)$ alors $H_n \geqslant \ln (n)$ et donc
$u_n\geqslant 0$.
\item La suite $(u_n)$ est d\'ecroissante et minor\'ee (par $0$) donc elle converge
vers un r\'eel $\gamma$. Ce r\'eel $\gamma$ s'appelle \emph{la constante d'Euler}
(d'après Leonhard Euler, 1707-1783, math\'ematicien d'origine suisse). Cette
constante vaut environ $0,5772156649\ldots$ mais on ne sait pas si
$\gamma$ est rationnel ou irrationnel.
\end{enumerate}
\fincorrection
\nocorrection
\nocorrection
\nocorrection
\correction{000524}
\begin{enumerate}
  \item $u_{n+q} = \cos \left( \frac{2(n+q)\pi}{q} \right) = \cos \left(\frac{2n\pi}{q}+2\pi\right) = \cos \left(\frac{2n\pi}{q}\right) = u_n$.
  \item $u_{nq} = \cos \left(\frac{2nq\pi}{q}\right) = \cos \left({2n\pi}\right)= 1 = u_0$ et $u_{nq+1} = \cos \left(\frac{2(nq+1)\pi}{q}\right) = \cos \left(\frac{2\pi}{q}\right) = u_1$.
Supposons, par l'absurde que $(u_n)$ converge vers $\ell$.
Alors la sous-suite $(u_{nq})_n$ converge vers $\ell$ comme 
$u_{nq}= u_0 = 1$ pour tout $n$ alors $\ell = 1$. D'autre part
la sous-suite $(u_{nq+1})_n$ converge aussi vers $\ell$, mais 
$u_{nq+1}= u_1 = \cos \frac{2\pi}{q}$, donc $\ell = \cos \frac{2\pi}{q}$. Nous obtenons une contradiction car pour $q\geq 2$, nous avons $\cos \frac{2\pi}{q} \not= 1$. Donc la suite $(u_n)$ ne converge pas.
\end{enumerate}
\fincorrection
\nocorrection
\nocorrection
\nocorrection
\nocorrection
\nocorrection
\nocorrection
\nocorrection
\nocorrection
\nocorrection
\nocorrection
\nocorrection
\nocorrection
\nocorrection
\nocorrection
\nocorrection
\nocorrection
\nocorrection
\nocorrection
\nocorrection
\nocorrection
\nocorrection
\nocorrection
\nocorrection
\nocorrection
\correction{004672}
\begin{enumerate}
  \item 
  \item $\ell = \pi$.
\end{enumerate}
 \fincorrection
\nocorrection
\nocorrection
\nocorrection
\correction{004676}
$\frac x2$.
\fincorrection
\nocorrection
\nocorrection
\nocorrection
\correction{004680}
\begin{enumerate}
  \item $u_n = \frac {(2n)!}{4^n(n!)^2}$.
  \item 
  \item 
  \item 
\end{enumerate}
 \fincorrection
\correction{004681}
Si $|a| < 1$, {\`a} partir d'un certain rang, $\left|\frac a{1+a^n} \right|<\alpha < 1$
$ \Rightarrow  \lim = 0$.

Si $|a| >1$, $\frac{u_{n+1}}{u_n} \to 0  \Rightarrow  u_n \to 0$.
\fincorrection
\nocorrection
\nocorrection
\correction{004684}
$\lim S_n = \lim u_n$.\par
Si $u_n = (3i/2)^n$, alors $(u_n)$ diverge, mais $S_n \xrightarrow[n\to\infty]{} 0$.
\fincorrection
\nocorrection
\nocorrection
\nocorrection
\nocorrection
\correction{004689}
\begin{enumerate}
  \item $1 \le S(n+1) \le S(n) + 1  \Rightarrow  0 \le \frac {S(n+1)}{S(n)} \le 2$.
  \item $\inf = 0$ $(99\dots99)$, $\sup = 2$ $(100\dots00)$.
  \item 
\end{enumerate}
 \fincorrection
\nocorrection
\nocorrection
\nocorrection
\correction{004693}
\begin{enumerate}
  \item Sinon, on construit une sous-suite strictement croissante.
  \item La suite $(\min(x_0,\dots,x_n))$ converge vers $0$, et prend
             une infinit{\'e} de valeurs diff{\'e}rentes.
\end{enumerate}
 \fincorrection
\nocorrection
\nocorrection
\correction{004696}
$u_{2n}-u_{n} \ge \frac {1 + \dots + n}{4n^2} \ge \frac 18$.
\fincorrection
\correction{004697}
\begin{enumerate}
  \item Par r{\'e}currence, $\sqrt n \le u_n \le \sqrt{2n}$.
  \item $\sqrt n \le u_n \le \sqrt{n+\sqrt{2(n-1)}}  \Rightarrow  \lim = 1$.
  \item $\sqrt{n+\sqrt{n-1}} \le u_n \le \sqrt{n+\sqrt{n-1+\sqrt{2(n-2)}}}
               \Rightarrow  \lim = \frac 12$.
\end{enumerate}
 \fincorrection
\correction{004698}
Soit $\ell =\inf \{b_n\,|\, n\in\N^*\}$, $\varepsilon>0$
et $p\in\N^*$ tel que $b_p\le \ell+\varepsilon$. Pour $n\in\N^*$ on
effectue la division euclidienne de $n$ par~$p$~: $n=pq+r$ d'o{\`u}
$a_n\le a_p^qa_r$ et $b_n\le b_p+\frac{\ln a_r}{\strut n}\le \ell+2\varepsilon$
pour $n$ assez grand.
\fincorrection
\correction{004699}
Si $e^{i\alpha}$ n'est pas valeur d'adh{\'e}rence alors il existe
$\delta>0$ tel que $|e^{ih(n)}-e^{i\alpha}|> \delta$ pour tout~$n$
assez grand donc l'ensemble $\bigcup_{k\in\N}[\alpha-\delta+2k\pi,\alpha-\delta+2k\pi]$
ne contient aucun terme de la suite $(h(n))$ pour $n$ assez grand ce qui
contredit les hypoth{\`e}ses $h(n)\xrightarrow[n\to\infty]{}+\infty$ et
$h(n+1)-h(n)\xrightarrow[n\to\infty]{}0$.
\fincorrection
\correction{004700}
Soit $E$ l'ensemble des valeurs d'adh{\'e}rence de~$(u_n)$. Si $u_{n_k}\xrightarrow[k\to\infty]{}\lambda$
alors $u_{n_k+1}\xrightarrow[k\to\infty]{}\lambda+\lambda^2$ donc $E$ est stable
par l'application $f$ : $x \mapsto x+x^2$. En fait $E$ est invariant par cette application
car la suite $(u_{n_k-1})$ admet une valeur d'adh{\'e}rence $\mu \in E$ et on a
$\mu^2+\mu = \lambda$. En particulier l'intervalle $[\inf(E),\sup(E)]$ est invariant par~$f$ ce qui implique $\sup(E) = \inf(E) = 0$.
\fincorrection
\correction{004701}
Soit $\ell$ une valeur d'adh{\'e}rence de~$(x_n)$.
Si l'on suppose que~$(x_n)$ ne converge pas vers~$\ell$ alors il existe un voisinage~$[a,b]$
de~$\ell$ tel qu'il y a une infinit{\'e} de termes dans~$[a,b]$ et une infinit{\'e}
hors de~$[a,b]$. Ceci implique que $[c,d] = f([a,b])$ n'est pas inclus dans~$[a,b]$
et que $[c,d]\setminus[a,b]$ contient une infinit{\'e} de termes, donc $(x_n)$
a une deuxi{\`e}me valeur d'adh{\'e}rence dans $[c,d]\setminus]a,b[$.
\fincorrection
\correction{004702}
Les suites $(\ln(u_n)/2^n)$ et $(\ln(1+u_n)/2^n)$ sont adjacentes.
\fincorrection
\correction{005221}
Supposons sans perte de généralité $u$ croissante (quite à remplacer $u$ par $-u$).
Dans ce cas, ou bien $u$ converge, ou bien $u$ tend vers $+\infty$.
Supposons que $u$ tende vers $+\infty$, et montrons qu'il en est de même pour la suite $v$.
Soit $A\in\Rr$. Il existe un rang $n_0$ tel que pour n naturel supérieur ou égal à $n_0$, $u_n\geq2A$.
Pour $n\geq n_0+1$, on a alors,

\begin{align*}
v_n&=\frac{1}{n+1}\left(\sum_{k=0}^{n_0}u_k+\sum_{k=n_0+1}^{n}u_k\right)\geq \frac{1}{n+1}\sum_{k=0}^{n_0}u_k+\frac{(n-n_0)2A}{n+1}
\end{align*}
Maintenant, quand $n$ tend vers $+\infty$, $\frac{1}{n+1}\sum_{k=0}^{n_0}u_k+\frac{(n-n_0)2A}{n+1}$ tend vers $2A$ et donc, il existe un rang $n_1$ à partir duquel $v_n\geq\frac{1}{n+1}\sum_{k=0}^{n_0}u_k+\frac{(n-n_0)2A}{n+1}>A$.
On a montré que~:~$\forall n\in\Nn,\;\exists n_1\in\Nn/\;(\forall n\in\Nn),\;(n\geq n_1\Rightarrow v_n>A)$. Par suite, $\lim_{n\rightarrow +\infty}v_n=+\infty$. Par contraposition, si $v$ ne tend pas vers $+\infty$, la suite $u$ ne tend pas vers $+\infty$ et donc converge, d'après la remarque initiale.
\fincorrection
\correction{005222}
\begin{enumerate}
 \item  La fonction $x\mapsto\frac{1}{x}$ est continue et décroissante sur $]0,+\infty[$ et donc, pour $k\in\Nn^*$, on a~:
 
$$\frac{1}{k+1}=(k+1-k)\frac{1}{k+1}\leq\int_{k}^{k+1}\frac{1}{x}\;dx\leq(k+1-k)\frac{1}{k}=\frac{1}{k}.$$
Donc, pour $k\geq1$,  $\frac{1}{k}\geq\int_{k}^{k+1}\frac{1}{x}\;dx$ et, pour $k\geq2$, $\frac{1}{k}\leq\int_{k-1}^{k}\frac{1}{x}\;dx$.
En sommant ces inégalités, on obtient pour $n\geq1$,

$$H_n=\sum_{k=1}^{n}\frac{1}{k}\geq\sum_{k=1}^{n}\int_{k}^{k+1}\frac{1}{x}\;dx=\int_{1}^{n+1}\frac{1}{x}\;dx=\ln(n+1),$$ 
et pour $n\geq2$,

$$H_n=1+\sum_{k=2}^{n}\frac{1}{k}\leq1+\sum_{k=2}^{n}\int_{k-1}^{k}\frac{1}{x}\;dx=1+\int_{1}^{n}\frac{1}{x}\;dx=1+\ln n,$$
cette dernière inégalité restant vraie quand $n=1$. Donc,

\begin{center}
\shadowbox{
$\forall n\in\Nn^*,\;\ln(n+1)\leq H_n\leq1+\ln n.$
}
\end{center}
 \item  Soit $n$ un entier naturel non nul.
$$u_{n+1}-u_n=\frac{1}{n+1}-\ln(n+1)+\ln n=\frac{1}{n+1}-\int_{n}^{n+1}\frac{1}{x}\;dx=\int_{n}^{n+1}\left(\frac{1}{n+1}-\frac{1}{x}\right)\;dx\leq0$$
car la fonction $x\mapsto\frac{1}{x}$ décroit sur $[n,n+1]$. De même,

$$v_{n+1}-v_n=\frac{1}{n+1}-\ln(n+2)+\ln(n+1)=\frac{1}{n+1}-\int_{n+1}^{n+2}\frac{1}{x}\;dx=\int_{n+1}^{n+2}\left(\frac{1}{n+1}-\frac{1}{x}\right)\;dx\geq0$$
car la fonction $x\mapsto\frac{1}{x}$ décroit sur $[n+1,n+2]$. Enfin,

$$u_n-v_n=\ln(n+1)-\ln n=\ln\left(1+\frac{1}{n}\right)$$ 
et donc la suite $u-v$ tend vers 0 quand $n$ tend vers $+\infty$.
Finalement, la suite $u$ décroit, la suite $v$ croit et la suite $u-v$ tend vers $0$. On en déduit que les suites $u$ et $v$ sont adjacentes, et en particulier convergentes et de même limite. Notons $\gamma$ cette limite.
Pour tout entier naturel non nul $n$, on a $v_n\leq\gamma\leq u_n$, et en particulier, $v_3\leq\gamma\leq u_1$ avec $v_3=0,5...$ et $u_1=1$. Donc, $\gamma\in\left[\frac{1}{2},1\right]$.
Plus précisément, pour $n$ entier naturel non nul donné, on a

$$0\leq u_n-v_n\leq\frac{10^{-2}}{2}\Leftrightarrow\ln\left(1+\frac{1}{n}\right)\leq0,005\Leftrightarrow\frac{1}{n}\leq e^{0,005}-1\Leftrightarrow n\geq\frac{1}{e^{0,005}-1}=199,5...\Leftrightarrow n\geq200.$$
Donc $0\leq\gamma-v_{100}\leq\frac{10^{-2}}{2}$ et une valeur approchée de $v_{200}$ à $\frac{10^{-2}}{2}$ près (c'est-à-dire arrondie à la 3 ème décimale la plus proche) est une valeur approchée de $\gamma$ à $10^{-2}$ près. On trouve $\gamma=0,57$ à $10^{-2}$ près par défaut. Plus précisémént,

\begin{center}
\shadowbox{
$\gamma=0,5772156649...$ ($\gamma$ est la constante d'\textsc{Euler}).
}
\end{center}
\end{enumerate}
\fincorrection
\correction{005225}
Posons $\alpha=\Arccos\frac{a}{b}$. $\alpha$ existe car $0<\frac{a}{b}<1$ et est élément de $\left]0,\frac{\pi}{2}\right[$. De plus, $a=b\cos\alpha$. Enfin, pour tout entier naturel $n$, $\frac{\alpha}{2^n}\in\left]0,\frac{\pi}{2}\right[$ et donc, $\cos\frac{\alpha}{2^n}>0$.
On a $u_0=b\cos\alpha$ et $v_0=b$ puis $u_1=\frac{1}{2}(u_0+v_0)=\frac{b}{2}(1+\cos\alpha)=b\cos^2\frac{\alpha}{2}$ et $v_1=\sqrt{u_1v_0}=\sqrt{b\cos^2\frac{\alpha}{2}\times b}=b\cos\frac{\alpha}{2}$ puis $u_2=\frac{b}{2}\cos\frac{\alpha}{2}(1+\cos\frac{\alpha}{2})=b\cos\frac{\alpha}{2}\cos^2\frac{\alpha}{2^2}$ et 
$v_2=\sqrt{b\cos\frac{\alpha}{2}\cos^2\frac{\alpha}{2^2}\times b\cos\frac{\alpha}{2}}=b\cos\frac{\alpha}{2}\cos\frac{\alpha}{2^2}$...
Montrons par récurrence que pour tout entier naturel non nul $n$, $v_n=b\prod_{k=1}^{n}\cos\frac{\alpha}{2^k}$ et $u_n=v_n\cos\frac{\alpha}{2^n}$.
C'est vrai pour $n=1$ et si pour $n\geq1$ donné, on a $v_n=b\prod_{k=1}^{n}\cos\frac{\alpha}{2^k}$ et $u_n=v_n\cos\frac{\alpha}{2^n}$ alors, 

$$u_{n+1}=\frac{1}{2}(v_n\cos\frac{\alpha}{2^n}+v_n)=v_n\cos^2\frac{\alpha}{2^{n+1}}$$ puis 

$$v_{n+1}=\sqrt{u_{n+1}v_n}=v_n\cos\frac{\alpha}{2^{n+1}}\;(\mbox{car}\;\cos\frac{\alpha}{2^{n+1}}>0),$$
et donc, $v_{n+1}=b\prod_{k=1}^{n+1}\cos\frac{\alpha}{2^k}$ puis $u_{n+1}=v_{n+1}\cos\frac{\alpha}{2^{n+1}}$.
On a montré par récurrence que

\begin{center}
\shadowbox{
$\forall n\in\Nn^*,\;v_n=b\prod_{k=1}^{n}\cos\frac{\alpha}{2^k}\;\mbox{et}\;u_n=v_n\cos\frac{\alpha}{2^n}.$
}
\end{center}
Pour tout entier naturel non nul $n$, on a $v_n>0$ et $\frac{v_{n+1}}{v_n}=\cos\frac{\alpha}{2^{n+1}}<1$. La suite $v$ est donc strictement décroissante. Ensuite, pour tout entier naturel non nul $n$, on a $u_n>0$ et
 
$$\frac{u_{n+1}}{u_n}=\frac{v_{n+1}}{v_n}\frac{\cos\frac{\alpha}{2^{n+1}}}{\cos\frac{\alpha}{2^n}}
=\frac{\cos^2\frac{\alpha}{2^{n+1}}}{\cos\frac{\alpha}{2^n}}
=\frac{1}{2}\left(1+\frac{1}{\cos\frac{\alpha}{2^n}}\right)>\frac{1}{2}(1+1)=1.
$$
La suite $u$ est strictement croissante. Maintenant, pour $n\in\Nn^*$,

\begin{align*}
v_n&=b\prod_{k=1}^{n}\cos\frac{\alpha}{2^k}=b\prod_{k=1}^{n}\frac{\sin\frac{\alpha}{2^{k-1}}}
{2\sin\frac{\alpha}{2^k}}\\
 &=\frac{\sin\alpha}{2^n\sin\frac{\alpha}{2^n}}
\end{align*}
Donc, quand $n$ tend vers $+\infty$, $v_n\sim\frac{\sin\alpha}{2^n\frac{\alpha}{2^n}}=\frac{\sin\alpha}{\alpha}$, puis $u_n=v_n\cos\frac{\alpha}{2^n}\sim v_n\sim\frac{\sin\alpha}{\alpha}$.
Ainsi, les suites $u$ et $v$ sont adjacentes de limite commune $b\frac{\sin\alpha}{\alpha}=\frac{\sqrt{b^2-a^2}}{\Arccos\left(\frac{a}{b}\right)}$.
\fincorrection
\correction{005226}
\begin{enumerate}
 \item  Pour $n\in\Nn^*$, $\left|\frac{\sin n}{n}\right|\leq\frac{1}{n}$. Comme $\frac{1}{n}\underset{n\rightarrow+\infty}{\longrightarrow}0$, $\frac{\sin n}{n}\underset{n\rightarrow+\infty}{\longrightarrow}0$.

\begin{center}
\shadowbox{
$\lim_{n\rightarrow +\infty}\frac{\sin n}{n}=0$.
}
\end{center}
 \item  Quand $n$ tend vers $+\infty$, $\ln\left(\left(1+\frac{1}{n}\right)^n\right)=n\ln\left(1+\frac{1}{n}\right)\sim n\times \frac{1}{n}=1$. Donc, $\ln\left(\left(1+\frac{1}{n}\right)^n\right)$ tend vers $1$ puis, $\left(1+\frac{1}{n}\right)^n=e^{n\ln(1+1/n)}$ tend vers $e^1=e$.

\begin{center}
\shadowbox{
$\lim_{n\rightarrow +\infty}\left(1+\frac{1}{n}\right)^n=e$.
}
\end{center}
 \item  Pour $n\in\Nn^*$, posons $u_n=\frac{n!}{n^n}$. Pour $n$ entier naturel non nul, on a

$$\frac{u_{n+1}}{u_n}=\frac{(n+1)!}{n!}\times\frac{n^n}{(n+1)^{n+1}}=\left(\frac{n}{n+1}\right)^n=\left(1+\frac{1}{n}\right)^{-n}.$$
Donc, quand $n$ tend vers $+\infty$, $\frac{u_{n+1}}{u_n}=e^{-n\ln(1+1/n)}=e^{-n(1/n+o(1/n))}=e^{-1+o(1)}$. Ainsi, $\frac{u_{n+1}}{u_n}$ tend vers $\frac{1}{e}=0.36...<1$. On sait alors que $\lim_{n\rightarrow +\infty}u_n=0$.

\begin{center}
\shadowbox{
$\lim_{n\rightarrow +\infty}\frac{n!}{n^n}=0$.
}
\end{center}

 \item  Pour $n\geq1$, $\frac{(n+\frac{1}{2})^2-1}{(n-\frac{1}{2})^2}\leq u_n\leq\frac{(n+\frac{1}{2})^2}{(n-\frac{1}{2})^2-1}$. Or, $\frac{(n+\frac{1}{2})^2-1}{(n-\frac{1}{2})^2}$ et $\frac{(n+\frac{1}{2})^2}{(n-\frac{1}{2})^2-1}$ tendent vers $1$ quand $n$ tend vers $+\infty$ et donc, d'après le théorème de la limite par encadrement, la suite $u$ converge et a pour limite $1$.

\begin{center}
\shadowbox{
$\lim_{n\rightarrow +\infty}\frac{E\left(\left(n+\frac{1}{2}\right)^2\right)}{E\left(\left(n-\frac{1}{2}\right)^2\right)}=1$.
}
\end{center}

 \item  Quand $n$ tend vers $+\infty$, $\sqrt[n]{n^2}=e^{\frac{1}{n}\ln(n^2)}=e^{2\ln n/n}=e^{o(1)}$, et donc $\sqrt[n]{n^2}$ tend vers $1$.

\begin{center}
\shadowbox{
$\lim_{n\rightarrow +\infty}\sqrt[n]{n^2}=1$.
}
\end{center}
 \item  $\sqrt{n+1}-\sqrt{n}=\frac{1}{\sqrt{n+1}+\sqrt{n}}\rightarrow0$.
 \item  $\frac{1}{n^3}\sum_{k=1}^{n}k^2=\frac{n(n+1)(2n+1)}{6n^3}\sim\frac{2n^3}{6n^3}=\frac{1}{6}$.
 \item  $\prod_{k=1}^{n}2^{k/2^k}=2^{\frac{1}{2}\sum_{k=1}^{n}\frac{k}{2^{k-1}}}$. Pour $x$ réel, posons $f(x)=\sum_{k=1}^{n}kx^{k-1}$. $f$ est dérivable sur $\Rr$ en tant que polynôme et pour tout réel $x$, 

$$f(x)=\left(\sum_{k=1}^{n}x^k\right)'(x)=\left(\sum_{k=0}^{n}x^k\right)'(x).$$
Pour $x\neq1$, on a donc 

$$f(x)=\left(\frac{x^{n+1}-1}{x-1}\right)'(x)=\frac{(n+1)x^n(x-1)-(x^{n+1}-1)}{(x-1)^2}=\frac{nx^{n+1}-(n+1)x^n+1}{(x-1)^2}.$$
En particulier, $\sum_{k=1}^{n}\frac{k}{2^{k-1}}=f\left(\frac{1}{2}\right)=\frac{\frac{n}{2^{n+1}}-\frac{n+1}{2^n}+1}{(\frac{1}{2}-1)^2}
\rightarrow4$ (d'après un théorème de croissances comparées). Finalement,

$$\prod_{k=1}^{n}2^{k/2^k}\rightarrow2^{4/2}=4.$$
\end{enumerate}
\fincorrection
\correction{005227}
Soit $n\in\Nn$.

\begin{align*}
\frac{1}{2\sqrt{n+u_n}}=\sqrt{n+1}-\sqrt{n}&\Leftrightarrow 2\sqrt{n+u_n}=\frac{1}{\sqrt{n+1}-\sqrt{n}}\Leftrightarrow2\sqrt{n+u_n}=\sqrt{n+1}+\sqrt{n}\\
 &4(n+u_n)=(\sqrt{n+1}+\sqrt{n})^2\Leftrightarrow u_n=-n+\frac{1}{4}(2n+1+2\sqrt{n(n+1)})\\
 &\Leftrightarrow u_n=\frac{1}{4}(-2n+1+2\sqrt{n(n+1)})
\end{align*}
Par suite, quand $n$ tend vers $+\infty$,

\begin{align*}
u_n&=-\frac{n}{2}+\frac{1}{4}+\frac{1}{2}\sqrt{n^2+n}=
\frac{1}{4}+\frac{n}{2}\left(\sqrt{1+\frac{1}{n}}-1\right)=\frac{1}{4}+\frac{n}{2}\frac{1/n}{\sqrt{1+\frac{1}{n}}+1}\\
 &=
\frac{1}{4}+\frac{1}{2}\frac{1}{\sqrt{1+\frac{1}{n}}+1}=\frac{1}{4}+\frac{1}{4}+o(1)=\frac{1}{2}+o(1).
\end{align*}
La suite $(u_n)$ converge et a pour limite $\frac{1}{2}$.

\fincorrection
\correction{005232}
Supposons que la suite $(\sqrt[n]{v_n})$ tende vers le réel positif $\ell$.

\begin{itemize}
\item[\textbullet] Supposons que $0\leq\ell<1$. Soit $\varepsilon=\frac{1-\ell}{2}$.

$\varepsilon$ est un réel strictement positif et donc, $\exists n_0\in\Nn/\;\forall n\in\Nn,(n\geq n_0\Rightarrow\sqrt[n]{v_n}<\ell+\frac{1-\ell}{2}=\frac{1+\ell}{2})$.

Pour $n\geq n_0$, par croissance de la fonction $t\mapsto t^n$ sur $\Rr^+$, on obtient $|u_n|<\left(\frac{1+\ell}{2}\right)^n$. Or, $0<\frac{1+\ell}{2}<\frac{1+1}{2}=1$ et donc 
$\left(\frac{1+\ell}{2}\right)^n$ tend vers $0$ quand $n$ tend vers $+\infty$. Il en résulte que $u_n$ tend vers  $0$ quand $n$ tend vers $+\infty$.

\item[\textbullet] Supposons que $\ell>1$. $\exists n_0\in\Nn/\;\forall n\in\Nn,\;(n\geq n_0\Rightarrow\sqrt[n]{v_n}>\ell-\frac{\ell-1}{2}=\frac{1+\ell}{2})$. Mais alors, pour $n\geq n_0$, $|u_n|>\left(\frac{1+\ell}{2}\right)^n$. Or, $\frac{1+\ell}{2}>\frac{1+1}{2}=1$, et donc $\left(\frac{1+\ell}{2}\right)^n$ tend vers $+\infty$ quand $n$ tend vers $+\infty$. Il en résulte que $|u_n|$ tend vers $+\infty$ quand $n$ tend vers $+\infty$.
\end{itemize}
Soit, pour $\alpha$ réel et $n$ entier naturel non nul, $u_n=n^\alpha$. $\sqrt[n]{u_n}=e^{\alpha\frac{\ln n}{n}}$ tend vers $1$ quand $n$ tend vers $+\infty$, et ceci pour toute valeur de $\alpha$. Mais, si $\alpha<0$, $u_n$ tend vers $0$, si $\alpha=0$, $u_n$ tend vers $1$ et si $\alpha>0$, $u_n$ tend vers $+\infty$. Donc, si $\ell=1$, on ne peut rien conclure.
\fincorrection
\correction{005233}
\begin{enumerate}
 \item  Supposons $\ell>0$. Soit $\varepsilon$ un réel strictement positif, élément de $]0,\ell[$.
$\exists n_0\in\Nn/\;\forall n\in\Nn,\;(n\geq n_0\Rightarrow0< \ell-\frac{\varepsilon}{2}<\frac{u_{n+1}}{u_n}<\ell+\frac{\varepsilon}{2})$.
Pour $n>n_0$, puisque $u_n=\frac{u_n}{u_{n-1}}\frac{u_{n-1}}{u_{n-2}}\frac{u_{n-2}}{u_{n-3}}...\frac{u_{n_0+1}}{u_{n_0}}u_{n_0}$, on a 
$u_{n_0}\left(\ell-\frac{\varepsilon}{2}\right)^{n-n_0}\leq  u_n\leq u_{n_0}\left(\ell+\frac{\varepsilon}{2}\right)^{n-n_0}$, et donc 

$$(u_{n_0})^{1/n}\left(\ell-\frac{\varepsilon}{2}\right)^{-n_0/n}\left(\ell-\frac{\varepsilon}{2}\right)\leq \sqrt[n]{u_n}\leq(u_{n_0})^{1/n}\left(\ell+\frac{\varepsilon}{2}\right)^{-n_0/n}\left(\ell+\frac{\varepsilon}{2}\right).$$
Maintenant, le membre de gauche de cet encadrement tend vers $\ell-\frac{\varepsilon}{2}$, et le membre de droite rend vers $\ell+\frac{\varepsilon}{2}$. Par suite, on peut trouver un entier naturel $n_1\geq n_0$ tel que, pour $n\geq n_1$, $(u_{n_0})^{1/n}\left(\ell-\frac{\varepsilon}{2}\right)^{-n_0/n}\left(\ell-\frac{\varepsilon}{2}\right)>\ell-\varepsilon$, et $(u_{n_0})^{1/n}\left(\ell+\frac{\varepsilon}{2}\right)^{-n_0/n}\left(\ell+\frac{\varepsilon}{2}\right)<\ell+\varepsilon$. Pour $n\geq n_1$, on a alors $\ell-\varepsilon<\sqrt[n]{u_n}<\ell+\varepsilon$.
On a montré que $\forall\varepsilon>0,\;\exists n_1\in\Nn/\;(\forall n\in\Nn),\;(n\geq n_1\Rightarrow\ell-\varepsilon<\sqrt[n]{u_n}<\ell+\varepsilon)$. Donc, $\sqrt[n]{u_n}$ tend vers $\ell$.
On traite de façon analogue le cas $\ell=0$.
 \item  Soient $a$ et $b$ deux réels tels que $0<a<b$. Soit $u$ la suite définie par 

$$\forall p\in\Nn,\;u_{2p}=a^pb^p\;\mbox{et}\;u_{2p+1}=a^{p+1}b^p.$$
(on part de $1$ puis on multiplie alternativement par $a$ ou $b$).
Alors, $\sqrt[2p]{u_{2p}}=\sqrt{ab}$ et $\sqrt[2p+1]{u_{2p+1}}=a^{\frac{p+1}{2p+1}}b^{\frac{p}{2p+1}}\rightarrow\sqrt{ab}$. Donc, $\sqrt[n]{u_n}$ tend vers $\sqrt{ab}$ (et en particulier converge).
On a bien sûr $\frac{u_{2p+1}}{u_{2p}}=a$ et $\frac{u_{2p+2}}{u_{2p+1}}=b$. La suite $\left(\frac{u_{n+1}}{u_n}\right)$ admet donc deux suites extraites convergentes de limites distinctes et est ainsi divergente. La réciproque du 1) est donc fausse.
 \item  
  \begin{enumerate}
  \item  Pour $n$ entier naturel donné, posons $u_n=\dbinom{2n}{n}$.
 
$$\frac{u_{n+1}}{u_n}=\frac{(2n+2)!}{(2n)!}\frac{n!^2}{(n+1)!^2}=\frac{(2n+2)(2n+1)}{(n+1)^2}=\frac{4n+2}{n+1}.$$
Ainsi, $\frac{u_{n+1}}{u_n}$ tend vers $4$ quand $n$ tend vers $+\infty$, et donc $\sqrt[n]{\dbinom{2n}{n}}$ tend vers $4$ quand $n$ tend vers $+\infty$.
  \item Pour $n$ entier naturel donné, posons $u_n=\frac{n^n}{n!}$.
 
$$\frac{u_{n+1}}{u_n}=\frac{(n+1)^{n+1}}{n^n}\frac{n!}{(n+1)!}=\left(1+\frac{1}{n}\right)^n.$$
Ainsi, $\frac{u_{n+1}}{u_n}$ tend vers $e$ quand $n$ tend vers $+\infty$, et donc $\sqrt[n]{u_n}=\frac{n}{\sqrt[n]{n!}}$ tend vers $e$ quand $n$ tend vers $+\infty$.
  \item Pour $n$ entier naturel donné, posons $u_n=\frac{(3n)!}{n^{2n}n!}$.

\begin{align*}
\frac{u_{n+1}}{u_n}&=\frac{(3n+3)!}{(3n)!}\frac{n^{2n}}{(n+1)^{2n+2}}\frac{n!}{(n+1)!}
=\frac{(3n+3)(3n+2)(3n+1)}{(n+1)^2(n+1)}\left(\frac{n}{n+1}\right)^{2n}\\
 &=
\frac{3(3n+2)(3n+1)}{(n+1)^2}\left(1+\frac{1}{n}\right)^{-2n}.
\end{align*}
Maintenant, $\left(1+\frac{1}{n}\right)^{-2n}=e^{-2n\ln(1+1/n)}=e^{-2n(\frac{1}{n}+o(\frac{1}{n}))}=e^{-2+o(1)}$, et donc $\frac{u_{n+1}}{u_n}$ tend vers $27e^{-2}$. Par suite, $\frac{1}{n^2}\sqrt[n]{\frac{(3n)!}{n!}}$ tend vers $\frac{27}{e^2}$.
  \end{enumerate}
\end{enumerate}
\fincorrection
\correction{005234}
D'après le théorème de la limite par encadrement~:

$$0\leq u_nv_n\leq u_n\leq1\Rightarrow u\;\mbox{converge et tend vers}\;1.$$
Il en est de même pour $v$ en échangeant les rôles de $u$ et $v$.
\fincorrection
\correction{005235}
Si $u_n^2\rightarrow0$, alors $|u_n|=\sqrt{|u_n^2|}\rightarrow0$ et donc $u_n\rightarrow0$.
Si $u_n^2\rightarrow\ell\neq0$, alors $(u_n)=(\frac{u_n^3}{u_n^2})$ converge.
(L'exercice n'a d'intérêt que si la suite $u$ est une suite complexe, car si $u$ est une suite réelle, on écrit immédiatement $u_n=\sqrt[3]{u_n^3}$ 
(et non pas $u_n=\sqrt{u_n^2}$)).
\fincorrection
\correction{005236}
Les suites $u$ et $v$ sont définies à partir du rang $1$ et strictement positives.
Pour tout naturel non nul $n$, on a~:

$$\frac{u_{n+1}}{u_n}=\left(\frac{n+2}{n+1}\right)^{n+1}\left(\frac{n}{n+1}\right)^n=e^{(n+1)\ln(n+2)+n\ln n-(2n+1)\ln(n+1)}.$$
Pour $x$ réel strictement positif, posons alors $f(x)=(x+1)\ln(x+2)+x\ln x-(2x+1)\ln(x+1)$.
$f$ est dérivable sur $]0,+\infty[$ et pour $x>0$,

\begin{align*}
f'(x)&=\frac{x+1}{x+2}+\ln(x+2)+1+\ln x-\frac{2x+1}{x+1}-2\ln(x+1)\\
 &=\frac{x+2-1}{x+2}+\ln(x+2)+1+\ln x-\frac{2x+2-1}{x+1}-2\ln(x+1)\\
 &=-\frac{1}{x+2}+\frac{1}{x+1}+\ln x+\ln(x+2)-2\ln(x+1).
\end{align*}
De même, $f'$ est dérivable sur $]0,+\infty[$ et pour $x>0$,

\begin{align*}
f''(x)&=\frac{1}{(x+2)^2}-\frac{1}{(x+1)^2}+\frac{1}{x}+\frac{1}{x+2}-\frac{2}{x+1}\\ 
 &=\frac{x(x+1)^2-x(x+2)^2+(x+1)^2(x+2)^2+x(x+1)^2(x+2)-2x(x+1)(x+2)^2}{x(x+1)^2(x+2)^2}\\
 &=\frac{-2x^2-3x+(x^2+2x+1)(x^2+4x+4)+(x^2+2x)(x^2+2x+1)-2(x^2+x)(x^2+4x+4)}{x(x+1)^2(x+2)^2}\\
 &=\frac{3x+4}{x(x+1)^2(x+2)^2}>0.
\end{align*}
$f'$ est strictement croissante sur $]0,+\infty[$ et donc, pour $x>0$, 

$$f'(x)<\lim_{t\rightarrow +\infty}f'(t)=\lim_{t\rightarrow +\infty}\left(-\frac{1}{t+2}+\frac{1}{t+1}+\ln\frac{t(t+2)}{(t+1)^2}\right)=0.$$
Donc, $f$ est strictement décroissante sur $]0,+\infty[$. Or, pour $x>0$,

\begin{align*}
f(x)&=(x+1)\ln(x+2)+x\ln x-(2x+1)\ln(x+1)\\
 &=(x+(x+1)-(2x+1))\ln x+(x+1)\ln\left(1+\frac{2}{x}\right)-(2x+1)\ln\left(1+\frac{1}{x}\right)\\
 &=\ln\left(1+\frac{2}{x}\right)-\ln\left(1+\frac{1}{x}\right)+2\frac{\ln\left(1+\frac{2}{x}\right)}{\frac{2}{x}}-2\frac{\ln\left(1+\frac{1}{x}\right)}{\frac{1}{x}}.
\end{align*}
On sait que $\lim_{u\rightarrow 0}\frac{\ln(1+u)}{u}=1$, et donc, quand $x$ tend vers $+\infty$, $f(x)$ tend vers $0+0+2-2=0$. Comme $f$ est strictement décroissante sur $]0,+\infty[$, pour tout réel $x>0$, on a $f(x)>\lim_{t\rightarrow +\infty}f(t)=0$.
f est donc strictement positive sur $]0,+\infty[$. Ainsi, $\forall n\in\Nn^*,\;f(n)>0$ et donc $\frac{u_{n+1}}{u_n}=e^{f(n)}>1$. La suite $u$ est strictement croissante.
(Remarque. On pouvait aussi étudier directement la fonction $x\mapsto\left(1+\frac{1}{x}\right)^x$ sur $]0,+\infty[$.)
On montre de manière analogue que la suite $v$ est strictement décroissante. Enfin, puisque $u_n$ tend vers $e$, et que $v_n=\left(1+\frac{1}{n}\right)u_n$ tend vers $e$, les suites $u$ et $v$ sont adjacentes.
(Remarque. En conséquence, pour tout entier naturel non nul $n$, $\left(1+\frac{1}{n}\right)^n<e<\left(1+\frac{1}{n}\right)^{n+1}$. Par exemple, pour $n=10$, on obtient $\left(\frac{11}{10}\right)^{10}<e<\left(\frac{11}{10}\right)^{11}$ et donc, $2,59...<e<2,85...$ et pour $n=100$, on obtient $1,01^{100}<e<1,01^{101}$ et donc $2,70...<e<2,73...$ Ces deux suites convergent vers $e$ lentement).
\fincorrection
\correction{005237}
Il est immédiat que $u$ croit strictement et que $v-u$ est strictement positive et tend vers $0$.
De plus, pour $n$ entier naturel donné, 

$$v_{n+1}-v_n=\frac{1}{(n+1)!}+\frac{1}{(n+1)\times(n+1)!}-\frac{1}{n\times n!}=\frac{n(n+1)+n-(n+1)^2}{n(n+1)\times(n+1)!}=\frac{-1}{n(n+1)\times(n+1)!}< 0,$$
et la suite $v$ est strictement décroissante. Les suites $u$ et $v$ sont donc adjacentes et convergent vers une limite commune (à savoir $e$).
 

(Remarque. Dans ce cas, la convergence est très rapide. On a pour tout entier naturel non nul $n$,  $\sum_{k=0}^{n}\frac{1}{k!}<e<\sum_{k=0}^{n}\frac{1}{k!}+\frac{1}{n\times n!}$ et $n=5$ fournit par exemple $2,716...<e<2,718...$).
\fincorrection
\correction{005238}
Pour $n$ entier naturel non nul donné, on a
$$u_{n+1}-u_n=\frac{1}{\sqrt{n+1}}-2\sqrt{n+2}+2\sqrt{n+1}=\frac{1}{\sqrt{n+1}}-\frac{2}{\sqrt{n+1}+\sqrt{n+2}}>\frac{1}{\sqrt{n+1}}-\frac{2}{\sqrt{n+1}+\sqrt{n+1}}=0.$$
De même,

$$v_{n+1}-v_n=\frac{1}{\sqrt{n+1}}-2\sqrt{n+1}+2\sqrt{n}=\frac{1}{\sqrt{n+1}}-\frac{2}{\sqrt{n+1}+\sqrt{n}}<\frac{1}{\sqrt{n+1}}-\frac{2}{\sqrt{n+1}+\sqrt{n+1}}=0.$$
La suite $u$ est strictement croissante et la suite $v$ est strictement décroissante. Enfin, 

$$v_n-u_n=2\sqrt{n+1}-2\sqrt{n}=\frac{2}{\sqrt{n}+\sqrt{n+1}},$$
et la suite $v-u$ converge vers $0$. Les suites $u$ et $v$ sont ainsi adjacentes et donc convergentes, de même limite.
\fincorrection
\correction{005241}
L'égalité proposée est vraie pour $n=2$ car $\cos\frac{\pi}{2^2}=\cos\frac{\pi}{4}=\frac{\sqrt{2}}{2}$.

Soit $n\geq2$. Supposons que $\cos(\frac{\pi}{2^{n}})=\frac{1}{2}\sqrt{2+\sqrt{2+...\sqrt{2}}}$ ($n-1$ radicaux).

Alors, puisque $\cos(\frac{\pi}{2^{n+1}})>0$ (car $\frac{\pi}{2^{n+1}}$ est dans $]0,\frac{\pi}{2}[$), 

$$\cos(\frac{\pi}{2^{n+1}})=\sqrt{\frac{1+\cos(\frac{\pi}{2^{n}})}{2}}=\sqrt{\frac{1}{2}(1+\frac{1}{2}\sqrt{2+\sqrt{2+...\sqrt{2}}})}=\frac{1}{2}\sqrt{2+\sqrt{2+...\sqrt{2}}},\;(n\;\mbox{radicaux}).$$

On a montré par récurrence que, pour $n\geq2$, $\cos(\frac{\pi}{2^{n}})=\frac{1}{2}\sqrt{2+\sqrt{2+...\sqrt{2}}}$ ($n-1$ radicaux).
 
Ensuite, pour $n\geq2$, 

$$\sin(\frac{\pi}{2^{n}})=\sqrt{\frac{1}{2}(1-\cos(\frac{\pi}{2^{n-1}})}=\frac{1}{2}\sqrt{2-\sqrt{2+...\sqrt{2}}}\;(n-1\;\mbox{radicaux})$$

Enfin, 

$$2^n\sqrt{2-\sqrt{2+...\sqrt{2}}}=2^n.2\sin\frac{\pi}{2^{n+1}}\sim2^{n+1}\frac{\pi}{2^{n+1}}=\pi.$$

Donc, $\lim_{n\rightarrow +\infty}2^n\sqrt{2-\sqrt{2+...\sqrt{2}}}=\pi$.
\fincorrection
\correction{005242}
\begin{enumerate}
\item  Pour $x$ réel positif, posons $f(x)=x-\ln(1+x)$ et $g(x)=(x+1)\ln(x+1)-x$.
$f$ et $g$ sont dérivables sur $[0,+\infty[$ et pour $x>0$, on a

$$f'(x)=1-\frac{1}{x+1}=\frac{x}{x+1}>0,$$

et

$$g'(x)=\ln(x+1)+1-1=\ln(x+1)>0.$$

$f$ et $g$ sont donc strictement croissantes sur $[0,+\infty[$ et en particulier, pour $x>0$, $f(x)>f(0)=0$ et de même, $g(x)>g(0)=0$. Finalement, $f$ et $g$ sont strictement positives sur $]0,+\infty[$ ou encore,

$$\forall x>0,\;\ln(1+x)<x<(1+x)\ln(1+x).$$

\item  Soit $k$ un entier naturel non nul.

D'après 1), $\ln(1+\frac{1}{k})<\frac{1}{k}<(1+\frac{1}{k})\ln(1+\frac{1}{k})$, ce qui fournit $k\ln(1+\frac{1}{k})<1<(k+1)Ln(1+\frac{1}{k})$, puis, par stricte croissance de la fonction exponentielle sur $\Rr$, 

$$\forall k\in\Nn^*,\;0<(1+\frac{1}{k})^k<e<(1+\frac{1}{k})^{k+1}.$$

En multipliant membre à membre ces encadrements, on obtient pour tout naturel non nul $n$~:

$$\prod_{k=1}^{n}(1+\frac{1}{k})^k<e^n<\prod_{k=1}^{n}(1+\frac{1}{k})^{k+1}.$$

Maintenant, 

$$\prod_{k=1}^{n}(1+\frac{1}{k})^k=\prod_{k=1}^{n}\left(\frac{k+1}{k}\right)^k=\frac{\prod_{k=2}^{n+1}k^{k-1}}{\prod_{k=1}^{n}k^k}=\frac{(n+1)^n}{n!}.$$

De même,
$$\prod_{k=1}^{n}(1+\frac{1}{k})^{k+1}=\frac{\prod_{k=2}^{n+1}k^{k}}{\prod_{k=1}^{n}k^{k+1}}=\frac{(n+1)^{n+1}}{n!}.$$

On a montré que $\forall n\in\Nn^*,\;\frac{(n+1)^n}{n!}<e^n<\frac{(n+1)^{n+1}}{n!}$ et donc 
 
$$\forall n\in\Nn^*,\;\frac{1}{e}\frac{n+1}{n}<\frac{\sqrt[n]{n!}}{n}<\frac{1}{e}\frac{n+1}{n}(n+1)^{1/n}.$$  

D'après le théorème de la limite par encadrements, comme $\frac{n+1}{n}$ tend vers 1 quand $n$ tend vers l'infini de même que $(n+1)^{1/n}=e^{\ln(n+1)/n}$, on a montré que $\frac{\sqrt[n]{n!}}{n}$ tend vers $\frac{1}{e}$ quand $n$ tend vers $+\infty$. 
\end{enumerate}
\fincorrection
\correction{005244}
On pose $u_0=0$, $u_1=0$, $u_2=1$, $u_3=1$, $u_4=0$, $u_5=1$,... c'est-à-dire 

$$\forall n\in\Nn,\;u_n=\left\{
\begin{array}{l}
0\;\mbox{si}\;n\;\mbox{n'est pas premier}\\
1\;\mbox{si}\;n\;\mbox{est premier}
\end{array}
\right..$$
 
Soit $k$ un entier naturel supérieur ou égal à $2$. Pour $n\geq2$, l'entier $kn$ est composé et donc, pour $n\geq 2$, $u_{kn}=0$. En particulier, la suite $(u_{kn})_{n\in\Nn}$ converge et a pour limite $0$. Maintenant, l'ensemble des nombres premiers est infini et si $p_n$ est le $n$-ième nombre premier, la suite $(p_n)_{n\in\Nn}$ est strictement croissante. La suite $(u_{p_n})_{n\in\Nn}$ est extraite de $(u_n)_{n\in\Nn}$ et est constante égale à $1$. En particulier, la suite $(u_{p_n})_{n\in\Nn}$ tend vers $1$. Ainsi la suite $(u_n)_{n\in\Nn}$ admet au moins deux suites extraites convergentes de limites distinctes et donc la suite $(u_n)_{n\in\Nn}$ diverge bien que toutes les suites $(u_{kn})_{n\in\Nn}$ convergent vers $0$ pour $k\geq2$.
\fincorrection
\correction{005245}
Soit $f$ une application de $\Nn$ dans lui-même, injective. Montrons que $\lim_{n\rightarrow +\infty}f(n)=+\infty$.

Soient $A$ un réel puis $m=\mbox{Max}(0,1+E(A))$.

Puisque $f$ est injective, on a $\mbox{card}(f^{-1}(\{0,1,...,m\})\geq m+1$. En particulier, $f^{-1}(\{0,1,...,m\})$ est fini (éventuellement vide).

Posons $n_0=1+\left\{
\begin{array}{l}
0\;\mbox{si}\;f^{-1}(\{0,1,...,m\})=\emptyset\\
\mbox{Max}f^{-1}(\{0,1,...,m\})\;\mbox{sinon}
\end{array}
\right.$.

Par définition de $n_0$, si $n\geq n_0$, $n$ n'est pas élément de $f^{-1}(\{0,1,...,m\})$ et donc $f(n)>m>A$.

On a montré que $\forall A\in\Rr,\;\exists n_0\in\Nn/\;(\forall n\in\Nn),\;(n\geq n_0\Rightarrow f(n)>A)$ ou encore 
$\lim_{n\rightarrow +\infty}f(n)=+\infty$.

\fincorrection
\correction{005247}
\begin{enumerate}
\item  Posons $a=\frac{2p\pi}{q}$ où $p\in\Zz$, $q\in\Nn^*$ et $\mbox{PGCD}(p,q)=1$. Pour tout entier naturel $n$, on a

$$u_{n+q}=\cos\left((n+q)\frac{2p\pi}{q}\right)=\cos\left(n\frac{2p\pi}{q}+2p\pi\right)=\cos(na)=u_n.$$

La suite $u$ est donc $q$-périodique et de même la suite $v$ est $q$-périodique. Maintenant, une suite périodique converge si et seulement si elle est constante (en effet, soient $T$ une période strictement positive de $u$ et $\ell$ la limite de $u$. Soit $k\in\{0,...,T-1\}$. $|u_k-u_0|=|u_{k+nT}-u_{nT}|\rightarrow|\ell-\ell|=0$ quand $n$ tend vers l'infini).

Or, si $a=\frac{2p\pi}{q}$ où $p\in\Zz$, $q\in\Nn^*$, $\mbox{PGCD}(p,q)=1$ et $\frac{p}{q}\in\Zz$, alors $u_1\neq u_0$ et la suite $u$ n'est pas constante et donc diverge, et si $a\in2\pi\Zz$, la suite $u$ est constante et donc converge.

\item  (a) et b)) Pour tout entier naturel $n$, 

$$v_{n+1}=\sin((n+1)a)=\sin(na)\cos a+\cos(na)\sin a=u_n\sin a+v_n\cos a.$$

Puisque $\frac{a}{2\pi}\notin\Zz$, $\sin a\neq0$ et donc $u_n=\frac{v_{n+1}-v_n\cos a}{\sin a}$. Par suite, si $v$ converge alors $u$ converge. De même, à partir de $\cos((n+1)a)=\cos(na)\cos a-\sin(na)\sin a$, on voit que si $u$ converge alors $v$ converge. Les suites $u$ et $v$ sont donc simultanément convergentes ou divergentes.

Supposons que la suite $u$ converge, alors la suite $v$ converge. Soient $\ell$ et $\ell'$ les limites respectives de $u$ et $v$. D'après ce qui précède, $\ell$ et $\ell'$ sont solutions du système~:
 
$$\left\{
\begin{array}{l}
\ell\sin a+\ell'\cos a=\ell'\\
\ell\cos a-\ell'\sin a=\ell.
\end{array}
\right.\Leftrightarrow\left\{
\begin{array}{l}
\ell\sin a+\ell'(\cos a-1)=0\\
\ell(\cos a-1)-\ell'\sin a=0.
\end{array}
\right..$$

Le déterminant de ce système vaut $-\sin^2a-(\cos a-1)^2<0$ car $a\notin2\pi\Zz$. Ce système admet donc l'unique solution $\ell=\ell'=0$ ce qui contredit l'égalité $\ell^2+{\ell'}^2=1$. Donc, les suites $u$ et $v$ divergent.

\item 
\begin{enumerate}
\item Soit $E'=\{na+2k\pi,\;n\in\Nn,\;k\in\Zz\}$. Supposons que $E'$ est dense dans $\Rr$ et montrons que $\{u_n,\;n\in\Nn\}$ et $\{v_n,\;n\in\Nn\}$ sont dense dans $[-1,1]$.

Soient $x$ un réel de $[-1,1]$ et $b=\Arccos x$, de sorte que $b\in[0,\pi]$ et que $x=\cos b$.

Soit $\varepsilon>0$. Pour $n$ entier naturel et $k$ entier relatif donnés, on a~:

\begin{align*}\ensuremath
|u_n-x|&=|\cos(na)-\cos b|=|\cos(na+2k\pi)-\cos b|=2|\sin(\frac{na+2k\pi-b}{2})\sin(\frac{na+2k\pi+b}{2})|\\
 &\leq2\left|\frac{na+2k\pi-b}{2}\right|\;(\mbox{l'inégalité}\;|\sin x|\leq|x|\;\mbox{valable pour tout réel}\;x\;\mbox{est classique})\\
 &=|na+2k\pi-b|
\end{align*}

En résumé, $\forall k\in\Zz,\;\forall n\in\Nn,\;|u_n-x|\leq|na+2k\pi-b|$. Maintenant, si $E'$ est dense dans $\Rr$, on peut trouver $n\in\Nn$ et $k\in\Zz|$ tels que $|na+2k\pi-b|<\varepsilon$ et donc $|u_n-x|<\varepsilon$.

Finalement, $\{u_n,\;n\in\Nn\}$ est dense dans $[-1,1]$. De même, on montre que $\{v_n,\;n\in\Nn\}$ est dense dans $[-1,1]$.

Il reste donc à démontrer que $E'$ est dense dans $\Rr$.

\item Soit $E=\{na+2k\pi,\;n\in\Zz,\;k\in\Zz\}$. $E$ est un sous groupe non nul de $(\Rr,+)$ et donc est soit de la forme $\alpha\Zz$ avec $\alpha=\mbox{inf}(E\cap]0,+\infty[)>0$, soit dense dans $\Rr$ si $\mbox{inf}(E\cap]0,+\infty[)=0$.

Supposons par l'absurde que $\mbox{inf}(E\cap]0,+\infty[)>0$. Puisque $E=\alpha\Zz$ et que $2\pi$ est dans $E$, il existe un entier naturel non nul $q$ tel que $2\pi=q\alpha$, et donc tel que $\alpha=\frac{2\pi}{q}$.

Mais alors, $a$ étant aussi dans $E$, il existe un entier relatif $p$ tel que $a=p\alpha=\frac{2p\pi}{q}\in2\pi\Qq$. Ceci est exclu et donc, $E$ est dense dans $\Rr$.

\item Soit $x$ dans $[-1,1]$. D'après ce qui précède, pour $\varepsilon>0$ donné, il existe $n\in\Zz$ tel que $|\cos(na)-x|<\varepsilon$ et donc $|u_{|n|}-x|<\varepsilon$, ce qui montre que $\{u_n,\;n\in\Nn\}$ est dense dans $[-1,1]$. De même, $\{v_n,\;n\in\Nn\}$ est dense dans $[-1,1]$.
\end{enumerate}
\end{enumerate}
\fincorrection
\correction{005250}
La suite $u$ n'est pas majorée. Donc, $\forall M\in\Rr,\;\exists n\in\Nn/\;u_n> M$. En particulier, $\exists n_0\in\Nn/\;u_{n_0}\geq0$.

Soit $k=0$. Supposons avoir construit des entiers $n_0$, $n_1$,..., $n_k$ tels que $n_0<n_1<...<n_k$ et $\forall i\in\{0,...,k\},\;u_{n_i}\geq i$.

On ne peut avoir~:~$\forall n>n_k,\;u_n<k+1$ car sinon la suite $u$ est majorée par le nombre 
$\mbox{Max}\{u_0,u_1,...,u_{n_k},k+1\})$. Par suite, $\exists n_{k+1}>n_k/\;u_{n_{k+1}}\geq k+1$.

On vient de construire par récurrence une suite $(u_{n_k})_{k\in\Nn}$ extraite de la suite $u$ telle que $\forall k\in\Nn,\;u_{n_k}\geq k$ et en particulier telle que $\lim_{k\rightarrow +\infty}u_{n_k}=+\infty$.

\fincorrection
\correction{005251}
Si $u$ converge vers un réel $\ell$, alors $\ell\in[0,1]$ puis, par passage à la limite quand $n$ tend vers $+\infty$,  $\ell(1-\ell)\geq\frac{1}{4}$, et donc $(\ell-\frac{1}{2})^2\leq0$ et finalement $\ell=\frac{1}{2}$. Par suite, si $u$ converge, $\lim_{n\rightarrow +\infty}u_n=\frac{1}{2}$.

De plus, puisque la suite $u$ est à valeurs dans $]0,1[$, pour $n$ naturel donné, on a~:

$$u_n(1-u_n)=\frac{1}{4}-(\frac{1}{2}-u_n)^2\leq\frac{1}{4}<u_{n+1}(1-u_n),$$

et puisque $1-u_n>0$, on a donc $\forall n\in\Nn,\;u_n<u_{n+1}$.

$u$ est croissante et majorée. Donc $u$ converge et $\lim_{n\rightarrow +\infty}u_n=\frac{1}{2}$ (amusant).
\fincorrection
\correction{005316}
\begin{enumerate}
\item  Pour tout réel $a$,

$$e^{i(2p+1)a}=(\cos a+i\sin a)^{2p+1}=\sum_{j=0}^{2p+1}C_{2p+1}^{j}\cos^{2p+1-j}a(i\sin a)^j$$

puis 

$$\sin((2p+1)a)=\mbox{Im}(e^{i(2p+1)a})=\sum_{j=0}^{p}C_{2p+1}^{2j+1}\cos^{2(p-j)}a(-1)^j\sin^{2j+1}a.$$

Pour $1\leq k\leq p$, en posant $a=\frac{k\pi}{2p+1}$, on obtient~: 

$$\forall k\in\{1,...,p\},\;\sum_{j=0}^{p}C_{2p+1}^{2j+1}\cos^{2(p-j)}\frac{k\pi}{2p+1}(-1)^j\sin^{2j+1}
\frac{k\pi}{2p+1}=0.$$
 
Ensuite, pour $1\leq k\leq p$, $0<\frac{k\pi}{2p+1}<\frac{\pi}{2}$ et donc $\sin^{2p+1}
\frac{k\pi}{2p+1}\neq0$. En divisant les deux membres de $(*)$ par $\sin^{2p+1}
\frac{k\pi}{2p+1}$, on obtient~:

$$\forall k\in\{1,...,p\},\;\sum_{j=0}^{p}(-1)^jC_{2p+1}^{2j+1}\cotan^{2(p-j)}\frac{k\pi}{2p+1}=0.$$

Maintenant, les $p$ nombres $\cotan^2\frac{k\pi}{2p+1}$ sont deux à deux distincts. En effet, pour $1\leq k\leq p$, $0<\frac{k\pi}{2p+1}<\frac{\pi}{2}$. Or, sur $]0,\frac{\pi}{2}[$, la fonction $x\mapsto\cotan x$ est strictement décroissante et strictement positive, de sorte que la fonction $x\mapsto\cotan^2x$ est strictement décroissante et en particulier injective.

Ces $p$ nombres deux à deux distintcs sont racines du polynôme $P=\sum_{j=0}^{p}(-1)^jC_{2p+1}^{2j+1}X^{p-j}$, qui est de degré $p$. Ce sont donc toutes les racines de $P$ (ces racines sont par suite simples et réelles). D'après les relations entre les coefficients et les racines d'un polynôme scindé, on a~:

\begin{align*}\ensuremath
\sum_{k=1}^{p}\cotan^2\frac{k\pi}{2p+1}=-\frac{-C_{2p+1}^3}{C_{2p+1}^1}=\frac{p(2p-1)}{3}.
\end{align*}

puis,

$$\sum_{k=1}^{p}\frac{1}{\sin^2\frac{k\pi}{2p+1}}=\sum_{k=1}^{p}(1+\cotan^2\frac{k\pi}{2p+1})=p+\frac{p(2p-1)}{3}=\frac{2p(p+1)}{3}.$$

\item  Pour $n$ entier naturel non nul donné, on a

$$u_{n+1}-u_n=\sum_{k=1}^{n+1}\frac{1}{k^2}-\sum_{k=1}^{n}\frac{1}{k^2}=\frac{1}{(n+1)^2}>0,$$

et la suite $(un)$ est strictement croissante. De plus, pour $n\geq2$,
 
$$u_n=\sum_{k=1}^{n}\frac{1}{k^2}=1+\sum_{k=2}^{n}\frac{1}{k^2}<1+\sum_{k=2}^{n}\frac{1}{k(k-1)}=1+\sum_{k=2}^{n}(\frac{1}{k-1}-\frac{1}{k})=1+1-\frac{1}{n}<2.$$

La suite $(u_n)$ est croissante et est majorée par $2$. Par suite, la suite $(u_n)$ converge vers un réel inférieur ou égal à $2$.

\item  Pour $x$ élément de $[0,\frac{\pi}{2}]$, posons $f(x)=x-\sin x$ et $g(x)=\tan x-x$.
$f$ et $g$ sont dérivables sur $[0,\frac{\pi}{2}]$ et pour $x$ élément de $[0,\frac{\pi}{2}]$, $f'(x)=1-\cos x$ et $g'(x)=\tan^2x$. $f'$ et $g'$ sont strictement positives sur $]0,\frac{\pi}{2}]$ et donc strictement croissantes sur $[0,\frac{\pi}{2}]$. Comme $f(0)=g(0)=0$, on en déduit que $f$ et $g$ sont strictement positives sur $]0,\frac{\pi}{2}[$.

Donc, $\forall x\in]0,\frac{\pi}{2}[,\;0<\sin x<x<\tan x$ et par passage à l'inverse $\forall x\in]0,\frac{\pi}{2}[,\;0<\cotan x<\frac{1}{x}<\frac{1}{\sin x}$.

\item  Pour $1\leq k\leq p$, $0<\frac{k\pi}{2p+1}<\frac{\pi}{2}$  et donc $0<\cotan\frac{k\pi}{2p+1}<\frac{2p+1}{k\pi}<\frac{1}{\sin\frac{k\pi}{2p+1}}$. Puis, $\cotan^2\frac{k\pi}{2p+1}<(\frac{(2p+1)^2}{\pi^2})\frac{1}{k^2}<\frac{1}{\sin\frac{k\pi}{2p+1}}$. En sommant ces inégalités, on obtient

$$\frac{\pi^2p(2p-1)}{3(2p+1)^2}=\frac{\pi^2}{(2p+1)^2}\sum_{k=1}^{p}\cotan^2\frac{k\pi}{2p+1}<u_p=\sum_{k=1}^{p}
\frac{1}{k^2}<\frac{\pi^2}{(2p+1)^2}\sum_{k=1}^{p}\frac{1}{\sin^2\frac{k\pi}{2p+1}}=\frac{2p(p+1)\pi^2}{3(2p+1)^2}.$$

Les membres de gauche et de droite tendent vers $\frac{\pi^2}{6}$ quand $p$ tend vers l'infini et donc la suite $(u_p)$ tend vers $\frac{\pi^2}{6}$.
\end{enumerate}
\fincorrection
\correction{005459}
\begin{enumerate}
\item  Soit $n\in\Nn$. Pour $x\in[0,\frac{\pi}{2}]$, $0\leq\Arcsin^x\leq(\frac{\pi}{2})^n$ et donc, par croissance de l'intégrale,

$$0\leq u_n\leq\frac{1}{n!}\int_{0}^{1}(\frac{\pi}{2})^ndx=\frac{1}{n!}(\frac{\pi}{2})^n.$$

D'après un théorème de croissances comparées, $\frac{1}{n!}(\frac{\pi}{2})^n$ tend vers $0$ quand $n$ tend vers $+\infty$. D'après le théorème des gendarmes, $u_n$ tend vers $0$ quand $n$ tend vers $+\infty$.

\item  $0\leq\int_{0}^{1}\frac{x^n}{1+x}\;dx\leq\int_{0}^{1}\frac{x^n}{1+0}\;dx=\frac{1}{n+1}$. Comme $\frac{1}{n+1}$ tend vers $0$ quand $n$ tend vers $+\infty$, $\int_{0}^{1}\frac{x^n}{1+x}\;dx$ tend vers $0$ quand $n$ tend vers $+\infty$.

\item  Soit $n\in\Nn^*$.

\begin{align*}\ensuremath
\left|\int_{0}^{\pi}\frac{n\sin x}{x+n}\;dx-\int_{0}^{\pi}\sin x\;dx\right|=\left|
\int_{0}^{\pi}\frac{-x\sin x}{x+n}\;dx\right|\leq\int_{0}^{\pi}\left|\frac{-x\sin x}{x+n}\;dx\right|
\leq\int_{0}^{\pi}\frac{\pi}{0+n}\;dx=\frac{\pi^2}{n}.
\end{align*}

Or, $\frac{\pi^2}{n}$ tend vers $0$ quand $n$ tend vers $+\infty$, et donc $\int_{0}^{\pi}\frac{n\sin x}{x+n}\;dx$ tend vers $\int_{0}^{\pi}\sin x\;dx=2$ quand $n$ tend vers $+\infty$.

\end{enumerate}
\fincorrection
\nocorrection
\nocorrection
\nocorrection
\correction{000539}
\begin{enumerate}
\item
 La fonction polynomiale $P (x) := x^{3} - 3 x + 1$ est continue et
d\'erivable sur $\R$ et sa d\'eriv\'ee est $P' (x) = 3 x^{2} - 3,$
qui est strictement n\'egative sur $]- 1 , + 1[.$ Par cons\'equent
$P$ est strictement d\'ecroissante sur $]- 1 , + 1[.$ Comme $P (0)
= 1 > 0$ et $P (1\slash 2) = - 3 \slash 8 < 0$ il en r\'esulte
gr\^ace au th\'eor\`eme des valeurs interm\'ediaires qu'il existe
un r\'eel unique $\alpha \in ]0,1\slash 2[$ tel que $P (\alpha) =
0.$
\item Comme $ f (x) - x = (x^{3} - 3 x + 1) \slash 9$ il en
r\'esulte que $\alpha$ est l'unique solution de l'\'equation $f
(x) = x$ dans $]0,1\slash 2[.$
\item
Comme $f' (x) = (x^{2} +2)\slash 3 > 0$ pour tout $x \in \R,$ on
en d\'eduit que $f$ est strictement croissante sur $\R.$ Comme $f
(0) = 1 \slash 9$ et $\lim_{x \to + \infty} f (x) = + \infty,$ on
en d\'eduit que $ f (\R^{+}) = [1\slash 9 , + \infty[.$ Comme
$x_{1} = f (x_{0}) = 1 \slash 9  >  0$ alors $x_1 > x_0=0$ ;   $f$ étant
strictement croissante sur $\R^{+},$ on en d\'eduit par
r\'ecurrence que $x_{n + 1} > x_{n}$ pour tout $n \in \N$ ce qui
prouve que la suite $(x_{n})$ est croissante.
\item  Un calcul simple montre que  $f (1 \slash 2) < 1
\slash 2.$ Comme $0 = x_{0} < 1 \slash 2$ et que $f$ est
croissante on en d\'eduit par r\'ecurrence que $ x_{n} < 1 \slash
2$ pour tout $n \in \N$ (en effet si $x_n < 1/2$ alors 
$x_{n+1} = f(x_n) < f(1/2) < 1/2$).
\item D'apr\`es les questions
pr\'ec\'edentes, la suite $(x_{n})$  est croissante et major\'ee,
elle converge donc vers un nombre r\'eel $\ell \in ]0, 1 \slash 2].$
De plus comme $x_{n + 1} = f (x_{n})$ pour tout $n \in \N,$ on en
d\'eduit par continuit\'e de $f$ que $\ell  = f (\ell).$ Comme $f
(1 \slash 2) < 1\slash 2,$ On en d\'eduit que $\ell \in ]0, 1
\slash 2[$ et v\'erifie l'\'equation $f (\ell) = \ell.$ D'apr\`es
la question 2, on en d\'eduit que $\ell = \alpha$ et donc $(x_{n})
$ converge vers $\alpha.$
\end{enumerate}
\fincorrection
\nocorrection
\nocorrection
\nocorrection
\nocorrection
\nocorrection
\nocorrection
\nocorrection
\nocorrection
\nocorrection
\nocorrection
\nocorrection
\nocorrection
\nocorrection
\nocorrection
\nocorrection
\nocorrection
\nocorrection
\nocorrection
\nocorrection
\nocorrection
\nocorrection
\nocorrection
\nocorrection
\correction{004703}
\begin{enumerate}
  \item $u_n\searrow\sqrt a$, et $u_n-\sqrt a < \frac{a-\sqrt a}{(2\sqrt a)^{2^n-1}}$.
  \item $u_{2n}\to0$, $u_{2n+1}\to1$.
  \item Si $0\le u_0 \le 1$ : $u_n\searrow0$, sinon $u_n\searrow-\infty$.
  \item \begin{itemize}	
	  \item $\frac{1}{4} < \alpha$ : $u_n\to\infty$;
	  \item $- \frac{3}{4} < \alpha \le \frac{1}{4}$: $u_n\to\frac{1-\sqrt{1-4a}}2$,
	  \item $-1 < \alpha \le - \frac{3}{4}$: 
	  $1$ point fixe et deux points réciproques. $(u_n)$ ne converge pas.
	  \end{itemize}
  \item Si $u_0 > {\scriptstyle-}\frac12$, $u_n\to\infty$;
             si $u_0 < {\scriptstyle-}\frac12$, $u_n\to-1$.
  \item $u_n\to\frac12$.
  \item Thm du point fixe sur $]-\infty, \frac74]  \Rightarrow  u_n\to 1$.
  \item Si $u_0\ne 1$, $\exists\ n \text{ tq } 4-3u_n < 0  \Rightarrow $ suite finie.
  \item $u_n\to\alpha \approx 0.39754$.
  \item 1 est point fixe, il y a deux points réciproques. $(u_n)$ ne converge pas.

  \item \begin{itemize}
       \item $1 < \alpha$: $u_n\to0$ si $u_0<1$, $u_n\to\infty$ si $u_0>1$
       \item $-1 < \alpha <1$: $u_n\to 1$
       \item $\alpha \le -1$: si $u_0 \ne 1$, $(u_n)$ diverge.
    	\end{itemize}
  \item \begin{itemize}
       $e^{1/e}  < \alpha$: $u_n\to\infty$.
       $1 < \alpha < e^{1/e}$: 2 pts fixes, $\beta<\gamma$. $u_n\to\beta$ si $u_0<\gamma$, 
       et $u_n\to\infty$ si $u_0>\gamma$. 
       $e^{-e} \le \alpha <  1$: 1 pt fixe, $\beta$, et $u_n\to\beta$.
       $\alpha <  e^{-e}$: 1 point fixe et deux points réciproques. $(u_n)$ ne converge pas.
    	\end{itemize} 
\end{enumerate}
 \fincorrection
\correction{004704}
CV (vers 0) ssi $|ka_0| < 1$.
\fincorrection
\nocorrection
\correction{004706}
$\ell = \frac{1+\sqrt5}2$.
\fincorrection
\correction{004707}
$u_{n+1} = \sqrt{ u_n^2 + u_n } \ge u_n$. Il n'y a pas de point fixe.
\fincorrection
\nocorrection
\nocorrection
\correction{004710}
$y_n-x_n = \text{cste}  \Rightarrow  x_n \to 0,\ y_n \to y_0-x_0$.
\fincorrection
\correction{004711}
$y_n-x_n = \frac {y_0-x_0}{3^n}$ et $y_n+x_n = y_0 + x_0
 \Rightarrow  x_n,y_n \longrightarrow \frac {y_0 + x_0}2$.
\fincorrection
\correction{004712}
$x_ny_n = \text{c}^{\text{te}}  \Rightarrow  x_n,y_n \to \sqrt{x_0y_0}$.
\fincorrection
\correction{004713}
\begin{enumerate}
  \item 
  \item $\ell = b\frac {\sin\varphi}{\varphi}$.
\end{enumerate}
 \fincorrection
\correction{004714}
\begin{enumerate}
  \item $x^3+y^3+z^3 - 3xyz = \frac 12 (x+y+z)\bigl((x-y)^2 + (y-z)^2 + (z-x)^2\bigr)$.
  \item $3c_{n+1} - 3b_{n+1} = a_n + b_n + c_n - 3\sqrt[3]{a_nb_nc_n} \ge 0  \Rightarrow  b_{n+1} \le c_{n+1}$.\par
    $\frac 3{a_{n+1}} - \frac 3{b_{n+1}} =
    \frac 1{a_n} + \frac 1{b_n} + \frac 1{c_n} - \frac 3{\sqrt[3]{a_nb_nc_n}}
    \ge 0  \Rightarrow  a_{n+1} \le b_{n+1}$.
    \par
    Donc $(a_n)$ cro{\^\i}t et $(c_n)$ d{\'e}cro{\^\i}t :
    $a_n \longrightarrow a$,
    $b_n \longrightarrow b$,
    $c_n \longrightarrow c$
    avec $\begin{cases} \frac 2a = \frac 1b + \frac 1c \cr b^2 = ac \cr 2c = a+b\cr\end{cases}
          \Rightarrow  a=b=c$.
\end{enumerate}
 \fincorrection
\correction{004715}
Pour $u_0>0$ on a $u_n\searrow0$ et pour $u_0<0$ on a $u_n\nearrow0$.
$f'(0)=\frac12$ donc $u_{n+1}\sim\frac12u_n$ et la s{\'e}rie $\sum u_n$ converge absolument
(d'Alembert).
\fincorrection
\correction{004716}
L'ensemble des valeurs d'adh{\'e}rence de la suite est un intervalle
dont tous les {\'e}l{\'e}ments sont points fixes par $f$.
S'il y a plusieurs valeurs d'adh{\'e}rence il faut passer de l'une {\`a} l'autre
avec une longueur de saut qui tend vers z{\'e}ro, on doit tomber sur point fixe
entre les deux, contradiction.
\fincorrection
\correction{005228}
\begin{enumerate}
 \item  Calcul formel de $u_n$.
Soit $x\in\Rr$. $\frac{x}{3-2x}=x\Leftrightarrow2x^2-2x=0\Leftrightarrow x=0\;\mbox{ou}\;x=1$. Pour $n$ entier naturel donné, on a alors

$$\frac{u_{n+1}-1}{u_{n+1}}=\frac{\frac{u_n}{3-2u_n}-1}{\frac{u_n}{3-2u_n}}=\frac{3u_n-3}{u_n}
=3\frac{u_n-1}{u_n}.$$
Par suite, $\frac{u_n-1}{u_n}=3^n\frac{u_0-1}{u_0}$, puis $u_n=\frac{u_0}{u_0-3^n(u_0-1)}$.
 \item   Calcul formel de $u_n$.
Soit $x\in\Rr$. $\frac{4(x-1)}{x}=x\Leftrightarrow x^2-4x+4=0\Leftrightarrow x=2$. Pour $n$ entier naturel donné, on a alors

$$\frac{1}{u_{n+1}-2}=\frac{1}{\frac{4(u_n-1)}{u_n}-2}=\frac{u_n}{2(u_n-2)}=\frac{u_n-2+2}{2(u_n-2)}=\frac{1}{2}+\frac{1}{u_n-2}.$$
Par suite, $\frac{1}{u_n-2}=\frac{n}{2}+\frac{1}{u_0-2}$, puis $u_n=2+\frac{2(u_0-2)}{(u_0-2)n+2}$.
\end{enumerate}
\fincorrection
\correction{005229}

Pour tout entier naturel $n$, on a $\left\{
\begin{array}{l}
u_{n+1}-u_n=\frac{1}{3}(v_n-u_n)\\
v_{n+1}-v_n=-\frac{1}{3}(v_n-u_n)\\
v_{n+1}-u_{n+1}=\frac{1}{3}(v_n-u_n)
\end{array}
\right.$.
La dernière relation montre que la suite $v-u$ garde un signe constant puis les deux premières relations montrent que pour tout entier naturel $n$, $\mbox{sgn}(u_{n+1}-u_n)=\mbox{sgn}(v_n-u_n)$ et 
$\mbox{sgn}(v_{n+1}-v_n)=-\mbox{sgn}(v_n-u_n)$. Les suites $u$ et $v$ sont donc monotones de sens de variation opposés.
Si par exemple $u_0\leq v_0$, alors, pour tout naturel $n$, on a~:

$$u_0\leq u_n\leq u_{n+1}\leq v_{n+1}\leq v_n\leq v_0.$$
Dans ce cas, la suite $u$ est croissante et majorée par $v_0$ et donc converge vers un certain réel $\ell$. De même, la suite $v$ est décroissante et minorée par $u_0$ et donc converge vers un certain réel $\ell'$. Enfin, puisque pour tout entier naturel $n$, on a $u_{n+1}=\frac{2u_n+v_n}{3}$, on obtient par passage à la limite quand $n$ tend vers l'infini,  $\ell=\frac{2\ell+\ell'}{3}$ et donc $\ell=\ell'$. Les suites $u$ et $v$ sont donc adjacentes. Si $u_0>v_0$, il suffit d'échanger les rôles de $u$ et $v$.
\textbf{Calcul des suites $u$ et $v$.}
Pour $n$ entier naturel donné, on a $v_{n+1}-u_{n+1}=\frac{1}{3}(v_n-u_n)$. La suite $v-u$ est géométrique de raison $\frac{1}{3}$. Pour tout naturel $n$, on a donc $v_n-u_n=\frac{1}{3^n}(v_0-u_0)$.
D'autre part, pour $n$ entier naturel donné, $v_{n+1}+u_{n+1}=v_n+u_n$. La suite $v+u$ est constante et donc, pour tout entier naturel $n$, on a $v_n+u_n=v_0+u_0$.
En additionnant et en retranchant les deux égalités précédentes, on obtient pour tout entier naturel $n$~:

$$u_n=\frac{1}{2}\left(v_0+u_0+\frac{1}{3^n}(v_0-u_0)\right)\;\mbox{et}\;v_n=\frac{1}{2}\left(v_0+u_0-\frac{1}{3^n}(v_0-u_0)\right).$$
En particulier, $\ell=\ell'=\frac{u_0+v_0}{2}$.
\fincorrection
\correction{005230}
Pour tout entier naturel $n$, on a $u_{n+1}-v_{n+1}=-\frac{1}{2}(u_n-v_n)$  et donc, $u_n-v_n=\left(-\frac{1}{2}\right)^n(u_0-v_0)$.
De même, en échangeant les rôles de $u$, $v$ et $w$, $v_n-w_n=\left(-\frac{1}{2}\right)^n(v_0-w_0)$ et $w_n-u_n=\left(-\frac{1}{2}\right)^n(w_0-v_0)$ (attention, cette dernière égalité n'est autre que la somme des deux premières et il manque encore une équation).
On a aussi, $u_{n+1}+v_{n+1}+w_{n+1}=u_n+v_n+w_n$ et donc, pour tout naturel $n$, $u_n+v_n+w_n=u_0+v_0+w_0$.
Ainsi, $u_n$, $v_n$ et $w_n$ sont solutions du système

$$\left\{
\begin{array}{l}
v_n-u_n=\left(-\frac{1}{2}\right)^n(v_0-u_0)\\
\rule{0mm}{7mm}w_n-u_n=\left(-\frac{1}{2}\right)^n(w_0-u_0)\\
\rule{0mm}{4mm}u_n+v_n+w_n=u_0+v_0+w_0
\end{array}
\right..$$
Par suite, pour tout entier naturel $n$, on a

$$\left\{
\begin{array}{l}
u_n=\frac{1}{3}\left((u_0+v_0+w_0)+\left(-\frac{1}{2}\right)^n(2u_0-v_0-w_0)\right)\\
v_n=\frac{1}{3}\left((u_0+v_0+w_0)+\left(-\frac{1}{2}\right)^n(-u_0+2v_0-w_0)\right)\\
w_n=\frac{1}{3}\left((u_0+v_0+w_0)+\left(-\frac{1}{2}\right)^n(-u_0-v_0+2w_0)\right)
\end{array}
\right..$$
Les suites $u$, $v$ et $w$ convergent vers $\frac{u_0+v_0+w_0}{3}$.
\fincorrection
\correction{005231}
Montrons tout d'abord que~:

$$\forall(x,y,z)\in]0,+\infty[^3,\;(x\leq y\leq z\Rightarrow\frac{3}{\frac{1}{x}+\frac{1}{y}+\frac{1}{z}}\leq\sqrt[3]{xyz}\leq\frac{x+y+z}{3}).$$
Posons $m=\frac{x+y+z}{3}$, $g=\sqrt[3]{xyz}$ et $h=\frac{3}{\frac{1}{x}+\frac{1}{y}+\frac{1}{z}}$.
Soient $y$ et $z$ deux réels strictement positifs tels que $y\leq z$. Pour $x\in]0,y]$, posons

\begin{center}
$u(x)=\ln m-\ln g=\ln\left(\frac{x+y+z}{3}\right)-\frac{1}{3}\left(\ln x+\ln y+\ln z\right)$.
\end{center}
$u$ est dérivable sur $]0,y]$ et pour $x\in]0,y]$, 

$$u'(x)=\frac{1}{x+y+z}-\frac{1}{3x}\leq\frac{1}{x+x+x}-\frac{1}{3x}=0.$$
$u$ est donc décroissante sur $]0,y]$ et pour $x$ dans $]0,y]$, $u(x)\geq u(y)=\ln\left(\frac{2y+z}{3}\right)-\frac{1}{3}(2\ln y+\ln z)$.
Soit $z$ un réel strictement positif fixé. Pour $y\in]0,z]$, posons $v(y)=\ln\left(\frac{2y+z}{3}\right)-\frac{1}{3}(2\ln y+\ln z)$. $v$ est dérivable sur $]0,z]$ et pour $y\in]0,z]$, 

$$v'(y)=\frac{2}{2y+z}-\frac{2}{3z}\leq\frac{2}{3z}-\frac{2}{3z}=0.$$
$v$ est donc décroissante sur $]0,z]$ et pour $y$ dans $]0,z]$, on a $v(y)\geq v(z)=0$. On vient de montrer que $g\leq m$.
En appliquant ce résultat à $\frac{1}{x}$, $\frac{1}{y}$ et $\frac{1}{z}$, on obtient $\frac{1}{g}\leq\frac{1}{h}$ et donc $h\leq g$.
Enfin, $m\leq\frac{z+z+z}{3}=z$ et $h\geq\frac{3}{\frac{1}{x}+\frac{1}{x}+\frac{1}{x}}=x$. Finalement,

\begin{center}
\shadowbox{
$x\leq h\leq g\leq m\leq z.$
}
\end{center}
Ce résultat préliminaire étant établi, puisque $0<u_0<v_0<w_0$, par récurrence, les suites $u$, $v$ et $w$ sont définies puis, pour tout naturel $n$, on a $u_n\leq v_n\leq w_n$, et de plus $u_0\leq u_n\leq u_{n+1}\leq w_{n+1}\leq w_n\leq w_0$.
La suite $u$ est croissante et majorée par $w_0$ et donc converge. La suite $w$ est décroissante et minorée par $u_0$ et donc converge. Enfin, puisque pour tout entier naturel $n$, $v_n=3w_{n+1}-u_n-w_n$, la suite $v$ converge.
Soient alors $a$, $b$ et $c$ les limites respectives des suites $u$, $v$ et $w$.
Puisque pour tout entier naturel $n$, on a $0<u_0\leq u_n\leq v_n\leq w_n$, on a déjà par passage à la limite $0<u_0\leq a\leq b\leq c$.
Toujours par passage à la limite quand $n$ tend vers $+\infty$~:

$$
\left\{
\begin{array}{l}
\frac{3}{a}=\frac{1}{a}+\frac{1}{b}+\frac{1}{c}\\
\rule{0mm}{5mm}b=\sqrt[3]{abc}\\
c=\frac{a+b+c}{3}
\end{array}
\right.
\Leftrightarrow
\left\{
\begin{array}{l}
2bc=ab+ac\\
b^2=ac\\
a+b=2c
\end{array}
\right.
\Leftrightarrow
\left\{
\begin{array}{l}
b=2c-a\\
a^2-5ac+4c^2=0\\
\end{array}
\right.
\Leftrightarrow(a=c\;\mbox{et}\;b=c)\;\mbox{ou}\;(a=4c\;\mbox{et}\;b=-2c).$$
$b=-2c$ est impossible car $b$ et $c$ sont strictement positifs et donc, $a=b=c$.
Les suites $u$, $v$ et $w$ convergent vers une limite commune.
\fincorrection
\correction{005240}
Tout d'abord , on montre facilement par récurrence que, pour tout entier naturel non nul $n$, $u_n$ existe et $u_n\geq1$.
Mais alors, pour tout entier naturel non nul $n$, $1\leq u_{n+1}=1+\frac{n}{u_n}\leq1+n$. Par suite, pour $n\geq2$, $1\leq u_n\leq n$, ce qui reste vrai pour $n=1$.

$$\forall n\in\Nn^*,\;1\leq u_n\leq n.$$
Supposons momentanément que la suite $(u_n-\sqrt{n})_{n\geq1}$ converge vers un réel $\ell$. Dans ce cas :

$$1+\frac{n}{u_n}=1+\frac{n}{\sqrt{n}+\ell+o(1)}=1+\sqrt{n}\frac{1}{1+\frac{\ell}{\sqrt{n}}+o\left(\frac{1}{\sqrt{n}}\right)}= 1+\sqrt{n}\left(1-\frac{\ell}{\sqrt{n}}+o\left(\frac{1}{\sqrt{n}}\right)\right)=\sqrt{n}+1-\ell+o(1).$$
D'autre part,

$$u_{n+1}=\sqrt{n+1}+\ell+o(1)=\sqrt{n}\left(1+\frac{1}{n}\right)^{1/2}+\ell+o(1)=\sqrt{n}+\ell+o(1),$$
et donc $\ell-(1-\ell)=o(1)$ ou encore $2\ell-1=0$. Donc, si la suite $(u_n-\sqrt{n})_{n\geq1}$ converge vers un réel $\ell$, alors $\ell=\frac{1}{2}$.
Il reste à démontrer que la suite $(u_n-\sqrt{n})_{n\geq1}$ converge.
On note que pour tout entier naturel non nul,

$$u_{n+1}-u_n=\frac{1}{u_n}(-u_n^2+u_n+n)=\frac{1}{u_n}\left(\frac{1}{2}(1+\sqrt{4n+1})-u_n\right)\left(u_n-\frac{1}{2}
(1-\sqrt{4n+1})\right).$$
Montrons par récurrence que pour $n\geq1$, $\frac{1}{2}(1+\sqrt{4n-3})\leq u_n\leq\frac{1}{2}(1+\sqrt{4n+1})$. Posons $v_n=\frac{1}{2}(1+\sqrt{4n-3})$ et $w_n=\frac{1}{2}(1+\sqrt{4n+1})$.

Si $n=1$, $v_1=1\leq u_1=1\leq\frac{1}{2}(1+\sqrt{5})=w_1$.

Soit $n\geq1$. Supposons que $v_n\leq u_n\leq w_n$. Alors, 

$$1+\frac{2n}{\sqrt{4n+1}+1}\leq u_{n+1}=1+\frac{n}{u_n}\leq1+\frac{2n}{\sqrt{4n-3}+1}.$$

Mais, pour $n\geq1$,

\begin{align*}
\mbox{sgn}(\frac{1}{2}(1+\sqrt{4n+5})-&(1+\frac{2n}{\sqrt{4n-3}+1}))=\mbox{sgn}((1+\sqrt{4n+5})(1+\sqrt{4n-3})-2(2n+1+\sqrt{4n-3}))\\
 &=\mbox{sgn}(\sqrt{4n+5}(1+\sqrt{4n-3})-(4n+1+\sqrt{4n-3}))\\
 &=\mbox{sgn}((4n+5)(1+\sqrt{4n-3})^2-(4n+1+\sqrt{4n-3})^2)\;(\mbox{par croissance de}\;x\mapsto x^2\;\mbox{sur}\;[0,+\infty[)\\
 &=\mbox{sgn}((4n+5)(4n-2+2\sqrt{4n-3})-((4n+1)^2+2(4n+1)\sqrt{4n-3}+4n-3))\\
 &=\mbox{sgn}(-8+8\sqrt{4n-3}) =\mbox{sgn}(\sqrt{4n-3}-1) =\mbox{sgn}((4n-3)-1)=\mbox{sgn}(n-1)=+
\end{align*}

Donc, $u_{n+1}\leq 1+1+\frac{2n}{\sqrt{4n-3}+1}\leq w_{n+1}$.

D'autre part,

$$1+\frac{2n}{\sqrt{4n+1}+1}=\frac{2n+1+\sqrt{4n+1}}{\sqrt{4n+1}+1}=\frac{(\sqrt{4n+1}+1)^2}{2(\sqrt{4n+1}+1)}
=\frac{1}{2}(1+\sqrt{4n+1})=v_{n+1},$$ 

et donc $v_{n+1}\leq u_{n+1}\leq w_{n+1}$.

On a montré par récurrence que 

$$\forall n\in\Nn^*,\;\frac{1}{2}(1+\sqrt{4n-3})\leq u_n\leq\frac{1}{2}(1+\sqrt{4n+1}),$$

(ce qui montre au passage que $u$ est croissante).

Donc, pour $n\geq1$,

$$\frac{1}{2}+\sqrt{n-\frac{3}{4}}-\sqrt{n}\leq u_n-\sqrt{n}\leq\frac{1}{2}+\sqrt{n+\frac{1}{4}}-\sqrt{n},$$

ou encore, pour tout $n\geq1$,
 
$$\frac{1}{2}-\frac{3}{4}\frac{1}{\sqrt{n-\frac{3}{4}}+\sqrt{n}}\leq u_n-\sqrt{n}\leq\frac{1}{2}+\frac{1}{4}\frac{1}{\sqrt{n+\frac{1}{4}}+\sqrt{n}}.$$

Maintenant, comme les deux suites 
$(\frac{1}{2}-\frac{3}{4}\frac{1}{\sqrt{n-\frac{3}{4}}+\sqrt{n}})$ et $(\frac{1}{2}+\frac{1}{4}\frac{1}{\sqrt{n+\frac{1}{4}}+\sqrt{n}})$ convergent toutes deux vers $\frac{1}{2}$, d'après le théorème de la limite par encadrements, la suite $(u_n-\sqrt{n})_{n\geq1}$ converge vers $\frac{1}{2}$.
\fincorrection
\correction{005277}
\begin{enumerate}
\item 
Posons $J=\left(
\begin{array}{cc}
1&1\\
1&1
\end{array}
\right)
$ de sorte que $A=I+J$. On a $J^2=2j$ et donc, plus généralement~:~$\forall k\geq1,\;J^k=2^{k-1}J$. Mais alors, puisque $I$ et $J$ commutent, la formule du binôme de \textsc{Newton} fournit pour $n$ entier naturel non nul donné~:

\begin{align*}\ensuremath
A^n&=(I+J)^n=I+\sum_{k=1}^{n}C_n^kJ^k=I+(\sum_{k=1}^{n}C_n^k2^{k-1})J=I+\frac{1}{2}(\sum_{k=0}^{n}C_n^k2^k-1)J\\
 &=I+\frac{1}{2}((1+2)^n-1)J=I+\frac{1}{2}(3^n-1)J=\frac{1}{2}\left(
\begin{array}{cc}
3^n+1&3^n-1\\
3^n-1&3^n+1
\end{array}
\right)
\end{align*}

ce qui reste vrai pour $n=0$. Donc,

$$\forall n\in\Nn,\;A^n=\frac{1}{2}\left(
\begin{array}{cc}
3^n+1&3^n-1\\
3^n-1&3^n+1
\end{array}
\right).$$

Poour $n$ entier naturel donné, posons $X_n=\left(
\begin{array}{c}
u_n\\
v_n
\end{array}
\right)$. Pour tout entier naturel $n$, on a alors $X_{n+1}=A.X_n$ et donc,

\begin{align*}\ensuremath
X_n=A^n.X_0=\frac{1}{2}\left(
\begin{array}{cc}
3^n+1&3^n-1\\
3^n-1&3^n+1
\end{array}
\right)\left(
\begin{array}{c}
1\\
0
\end{array}
\right)=\left(
\begin{array}{c}
\frac{3^n+1}{2}\\
\frac{3^n-1}{2}
\end{array}
\right).
\end{align*}

Donc,

$$\forall n\in\Nn,\;u_n=\frac{3^n+1}{2}\;\mbox{et}\;v_n=\frac{3^n-1}{2}.$$

\item  Soit $n\in\Nn$. $u_{n+1}+v_{n+1}=3(u_n+v_n)$. Donc, la suite $u+v$ est une suite géométrique de raison $3$ et de premier terme $u_0+v_0=1$. On en déduit que 

$$\forall n\in\Nn,\;u_n+v_n=3^n\;(I).$$

De même, pour tout entier naturel $n$ $u_{n+1}-v_{n+1}=u_n-v_n$. Donc, la suite $u+v$ est une suite constante. Puisque $u_0-v_0=1$, on en déduit que

$$\forall n\in\Nn,\;u_n-v_n=1\;(II).$$

En additionnant et en retranchant $(I)$ et $(II)$, on obtient

$$\forall n\in\Nn,\;u_n=\frac{3^n+1}{2}\;\mbox{et}\;v_n=\frac{3^n-1}{2}.$$

\end{enumerate}

\fincorrection
\correction{005425}
\begin{enumerate}
\item  Pour $x\geq-1$, posons $f(x)=\sqrt{1+x}$ et $g(x)=f(x)-x$.

Soit $u_0\in I=[-1,+\infty[$. $f$ est définie sur $I$ et de plus $f(I)=[0,+\infty[\subset[-1,+\infty[$. On en déduit, par une démonstration par récurrence, que la suite $u$ est définie.

Si la suite $u$ converge, puisque $\forall n\in\Nn,\;u_n\geq-1$, sa limite $\ell$ vérifie $\ell\geq-1$. Puisque $f$ est continue sur $[-1,+\infty[$ et donc en $\ell$,

$$\ell=\lim_{n\rightarrow +\infty}u_{n+1}=\lim_{n\rightarrow +\infty}f(u_n)=f(\lim_{n\rightarrow +\infty}u_n)=f(ell).$$

et $\ell$ est un point fixe de $f$. Or, pour $x\geq-1$,

\begin{align*}\ensuremath
\sqrt{1+x}=x&\Leftrightarrow1+x=x^2\;\mbox{et}\;x\geq0\Leftrightarrow(x=\frac{1-\sqrt{5}}{2}\;\mbox{ou}\;x=\frac{1+\sqrt{5}}{2})\;
\mbox{et}\;x\geq0\\
 &\Leftrightarrow x=\frac{\sqrt{5}+1}{2}.
\end{align*}

Ainsi, si la suite $(u_n)$ converge, c'est vers le nombre $\alpha=\frac{\sqrt{5}+1}{2}$.

Pour $x\geq-1$,

\begin{align*}\ensuremath
\mbox{sgn}(f(x)-\alpha)&=\mbox{sgn}(\sqrt{1+x}-\sqrt{1+\alpha})=\mbox{sgn}((1+x)-(1+\alpha))\quad(\mbox{par croissance de}\;x\mapsto x^2\;\mbox{sur}\;[0,+\infty[)\\
 &=\mbox{sgn}(x-\alpha).
\end{align*}

Ainsi, les intervalles $[-1,\alpha[$ et $]\alpha,+\infty[$ sont stables par $f$. Donc, si $-1\leq u_0<\alpha$, alors par récurrence $\forall n\in\Nn,\;-1\leq u_n<\alpha$ et si $u_0>\alpha$, alors par récurrence, $\forall n\in\Nn,\;u_n>\alpha$.

Soit $x\geq-1$. Si $x\in[-1,0]$, $\sqrt{1+x}-x\geq0$ et si $x\geq0$,

\begin{align*}\ensuremath
\mbox{sgn}(g(x))&=\mbox{sgn}(\sqrt{1+x}-x)\\
 &=\mbox{sgn}((1+x)-x^2)\quad(\mbox{par croissance de}\;x\mapsto x^2\;\mbox{sur}\;[0,+\infty[)\\
 &=\mbox{sgn}(x+\frac{\sqrt{5}-1}{2})(-x+\frac{1+\sqrt{5}}{2}-x)=\mbox{sgn}(\alpha-x)\;(\mbox{car ici}\;x\geq0).
\end{align*}

On en déduit que, si $x\in[-1,\alpha[$, $f(x)>x$, et si $x\in]\alpha,+\infty[$, $f(x)<x$. Mais alors, 
si $-1\leq u_0<\alpha$, puisque $\forall n\in\Nn,\;-1\leq u_n<\alpha$, pour $n$ entier naturel donné, on a

$$u_{n+1}=f(u_n)>u_n.$$

La suite $u$ est donc strictement croissante, majorée par $\alpha$ et donc convergente. On sait de plus que sa limite est nécessairement $\alpha$.

Si $u_0>\alpha$, puisque $\forall n\in\Nn,\;u_n>\alpha$, pour $n$ entier naturel donné, on a

$$u_{n+1}=f(u_n)<u_n.$$

La suite $u$ est donc strictement décroissante, minorée par $\alpha$ et donc convergente. On sait de plus que sa limite est nécessairement $\alpha$. Enfin, si $u_0=\alpha$, la suite $u$ est constante.

En résumé,

si $u_0\in[-1,\frac{\sqrt{5}+1}{2}[$, la suite $u$ est strictement croissante, convergente de limite $\frac{\sqrt{5}+1}{2}[$,

si $u_0\in]\frac{\sqrt{5}+1}{2},+\infty[$, la suite $u$ est strictement décroissante, convergente de limite $\frac{\sqrt{5}+1}{2}[$,

si $u_0=\frac{\sqrt{5}+1}{2}[$, la suite $u$ est constante et en particulier convergente de limite $\frac{\sqrt{5}+1}{2}$.

Ainsi, dans tous les cas, la suite $u$ est convergente et $\lim_{n\rightarrow +\infty}u_n=\frac{1+\sqrt{5}}{2}$.

$$\includegraphics{../images/img005425-1}$$


\item  Si $u_0>0$, alors puisque $f$ est définie sur l'intervalle $I=]0,+\infty[$ et que $I$ est stable par $f$ ($\forall x>0,\;\ln(1+x)>\ln1=0$), la suite $u$ est définie et est strictement positive. Si la suite $u$ converge, sa limite $\ell$ est un réel positif \textbf{ou nul}. Par continuité de $f$ sur $[0,+\infty[$ et donc en $\ell$,

$$\ell=\lim_{n\rightarrow +\infty}u_{n+1}=\lim_{n\rightarrow +\infty}f(u_n)=f(\lim_{n\rightarrow +\infty}u_n)=f(\ell).$$

Pour $x>-1$, posons $g(x)=\ln(1+x)-x$. $g$ est définie et dérivable sur $]-1,+\infty[$ et pour $x>-1$,

$$g'(x)=\frac{1}{1+x}-1=-\frac{x}{1+x}.$$

$g'$ est strictement positive sur $]-1,0[$ et strictement négative sur $]0,+\infty[$. $g$ est donc strictement croissante sur $]-1,0]$ et strictement décroissante sur $[0,+\infty[$. Par suite, si $x\in]-1,0[\cup]0,+\infty[$, $g(x)<0$. En particulier, pour $x\in]-1,0[\cup]0,+\infty[$, $f(x)\neq x$. Puisque $f(0)=0$, $f$ admet dans $]-1,+\infty[$ un et un seul point fixe à savoir $0$.

En résumé, si $u_0>0$, la suite $u$ est définie, strictement positive, et de plus, si la suite $u$ converge, alors $\lim_{n\rightarrow +\infty}u_n=0$.

Mais, pour $n$ entier naturel donné,

$$u_{n+1}-u_n=\ln(1+u_n)-u_n<0.$$

Par suite, la suite $u$ est strictement décroissante, minorée par $0$ et donc, d'après ce qui précède, converge vers $0$.

Si $u_0=0$, la suite $u$ est constante. Il reste donc à étudier le cas où $u_0\in]-1,0[$. Montrons par l'absurde qu'il existe un rang $n_0$ tel que $u_{n_0}\leq-1$. Dans le cas contraire, $\forall n\in\Nn,\;u_n>-1$. Comme précédemment, par récurrence, la suite $u$ est à valeurs dans $]-1,0[$ et strictement décroissante. Etant minorée par $-1$, la suite $u$ converge vers un certain réel $\ell$.

Puisque $\forall n\in\Nn,\;-1<u_n\leq u_0<0$, on a $-1\leq\ell\leq u_0<0$. Donc, ou bien $\ell=-1$, ou bien $f$ est continue en $\ell$ et $\ell$ est un point fixe de $f$ élément de $]-1,0[$.

On a vu que $f$ n'admet pas de point fixe dans $]-1,0[$ et donc ce dernier cas est exclu. Ensuite, si $\ell=-1$, il existe un rang $N$ tel que $u_N\leq -0.9$. Mais alors, $u_{N+1}=\leq\ln(-0,9+1)=-2,3...<-1$ ce qui constitue de nouveau une contradiction.

Donc, il existe un rang $n_0$ tel que $u_{n_0}\leq-1$ et la suite $u$ n'est pas définie à partir d'un certain rang.

En résumé,

si $u_0\in]0,+\infty[$, la suite $u$ est strictement décroissante, convergente et $\lim_{n\rightarrow +\infty}u_n=0$,

si $u_0=0$, la suite $u$ est constante,

et si $u_0\in]-1,0[$, la suite $u$ n'est pas définie à partir d'un certain rang.

$$\includegraphics{../images/img005425-2}$$


\item  Pour tout choix de $u_0$, $u_1\in[-1,1]$. On supposera dorénavant que $u_0\in[-1,1]$. Si $u_0=0$, la suite $u$ est constante. Si $u_0\in[-1,0[$, considérons la suite $u'$ définie par $u_0'=-u_0$ et $\forall n\in\Nn,\;u_{n+1}'=\sin(u_n')$. La fonction $x\mapsto\sin x$, il est clair par récurrence que $\forall n\in\Nn,\;u_n'=-u_n$. On supposera dorénavant que $u_0\in]0,1]$.

Puisque $]0,1]\subset]0,\frac{\pi}{2}]$, on a $\sin]0,1]\subset]0,1]$ et l'intervalle $I=]0,1]$ est stable par $f$. Ainsi, si $u_0\in]0,1]$, alors, $\forall n\in\Nn,\;u_n\in]0,1]$.

Pour $x\in[0,1]$, posons $g(x)=\sin x-x$. $g$ est dérivable sur $[0,1]$ et pour $x\in[0,1]$, $g'(x)=\cos x-1$. $g'$ est strictement négative sur $]0,1]$ et donc strictement décroissante sur $[0,1]$. On en déduit que pour $x\in]0,1]$, $g(x)<g(0)=0$.

Mais alors, pour $n$ entier naturel donné, $u_{n+1}=\sin(u_n)<u_n$. La suite $u$ est ainsi strictement décroissante, minorée par $0$ et donc converge vers $\ell\in[0,1]$. La fonction $x\mapsto\sin x$ est continue sur $[0,1]$ et donc, $\ell$ est un point fixe de $f$. L'étude de $g$ montre que $f$ a un et un seul point fixe dans $[0,1]$ à savoir $0$. La suite $u$ est donc convergente et $\lim_{n\rightarrow +\infty}u_n=0$.

L'étude préliminaire montre la suite $u$ converge vers $0$ pour tout choix de $u_0$.

$$\includegraphics{../images/img005425-3}$$


\item  Si $u_0$ est un réel quelconque, $u_1\in[-1,1]\subset[-\frac{\pi}{2},\frac{\pi}{2}]$ puis $u_2\in[0,1]$. On supposera dorénavant que $u_0\in[0,1]$.

On a $\cos([0,1])=[\cos 1,\cos0]=[0,504...,1]\subset[0,1]$. Donc, la fonction $x\mapsto\cos x$ laisse stable l'intervalle $I=[0,1]$. On en déduit que $\forall n\in\Nn,\;u_n\in[0,1]$.

Pour $x\in[0,1]$, on pose $g(x)=\cos x-x$. $g$ est somme de deux focntions strictement décroissantes sur $[0,1]$ et est donc strictement décroissante sur $[0,1]$. De plus, $g$ est continue sur $[0,1]$ et vérifie $g(0)=\cos0>0$ et $g(1)=\cos1-1<0$. $g$ s'annule donc une et une seule fois sur $[0,1]$ en un certain réel $\alpha$. Ainsi, $f$ admet sur $[0,1]$ un unique point fixe, à savoir $\alpha$. Puisque $f$ est continue sur le segment $[0,1]$, on sait que si la suite $u$ converge, c'est vers $\alpha$.

La fonction $f~:~x\mapsto\cos x$ est dérivable sur $[0,1]$ et pour $x\in[0,1]$,

$$|f'(x)|=|-\sin x|\leq\sin1<1.$$

L'inégalité des accroissements finis montre alors que $\forall(x,y)\in[0,1]^2,\;|\cos x-\cos y|\leq\sin1|x-y|$. Pour $n$ entier naturel donné, on a alors

$$|u_{n+1}-\alpha|=|f(u_n)-f(\alpha)|\leq\sin1|u_n-\alpha|,$$

et donc, pour tout entier naturel $n$,

$$|u_n-\alpha|\leq(\sin1)^n|u_0-\alpha|\leq(\sin1)^n.$$

Comme $0\leq\sin1<1$, la suite $(\sin1)^n$ converge vers $0$, et donc la suite $(u_n)_{n\in\Nn}$ converge vers $\alpha$. On peut noter que puisque la fonction $x\mapsto\cos x$ est strictement décroissante sur $[0,1]$, les deux suites $(u_{2n})_{n\in\Nn}$ et $(u_{2n+1})_{n\in\Nn}$ sont strictement monotones, de sens de variations contraires (dans le cas où $u_0\in[0,1]$. On peut noter également que si $n>\frac{\ln(10^{-2})}{\ln(\sin1)}=26,6...$, alors $(\sin1)^n<10^{-2}$. Par suite, $u_{27}$ est une valeur approchée de $\alpha$ à $10^{-2}$ près. La machine fournit $\alpha=0,73...$ (et même $\alpha=0,739087042.....$).

$$\includegraphics{../images/img005425-4}$$

\item  Si $u_0$ est un réel quelconque, alors $\forall n\in\Nn^*,\;u_n\in[-1,1]$. On supposera sans perte de généralité que $u_0\in[-1,1]$. Si $u_0=0$, la suite $u$ est constante et d'autre part, l'étude du cas $u_0\in[-1,0[$ se ramème, comme en 3), à l'étude du cas $u_0\in]0,1]$. On supposera dorénavant que $u_0\in]0,1]$.

Si $x\in]0,1]$, alors $2x\in]0,2]\subset]0,\pi[$ et donc $\sin(2x)\in]0,1]$. L'intervalle $I=]0,1]$ est donc stable par la fonction $f~:~x\mapsto\sin(2x)$. On en déduit que $\forall n\in\Nn,\;u_n\in]0,1]$.

Pour $x\in[0,1]$, posons $g(x)=\sin(2x)-x$. $g$ est dérivable sur $[0,1]$ et pour $x\in[0,1]$, $g'(x)=2\cos(2x)-1$. $g$ est donc strictement croissante sur $[0,\frac{\pi}{4}]$ et strictement décroissante sur $[\frac{\pi}{4},1]$. On en déduit que si $x\in]0,\frac{\pi}{4}]$, $g(x)>g(0)=0$. D'autre part, $g$ est continue et strictement décroissante sur $[\frac{\pi}{4},1]$ et vérifie $g(\frac{\pi}{4})=1-\frac{\pi}{4}>0$ et $g(1)=\sin2-1<0$. $g$ s'annule donc une et une seule fois en un certain réel $\alpha\in]\frac{\pi}{4},1[$.

En résumé, $g$ s'annule une et une seule fois sur $]0,1]$ en un certain réel $\alpha\in]\frac{\pi}{4},1[$, $g$ est strictement positive sur $]0,\alpha[$ et strictement négative sur $]\alpha,1]$.

Supposons que $u_0\in]0,\frac{\pi}{4}[$ et montrons par l'absurde que $\exists n_0\in\Nn/\;u_{n_0}\in[\frac{\pi}{4},1]$. Dans le cas contraire, tous les $u_n$ sont dans $]0,\frac{\pi}{4}[$. Mais alors, pour tout entier naturel $n$,

$$u_{n+1}-u_n=f(u_n)-u_n=g(u_n)>0.$$

La suite $u$ est donc strictement croissante. Etant majorée par $\frac{\pi}{4}$, la suite $u$ converge. Comme $g$ est continue sur $[u_0,\frac{\pi}{4}]$ et que $\forall n\in\Nn,\;u_n\in[u_0,\frac{\pi}{4}]$, on sait que la limite de $u$ est un point fixe de $f$ élément de $[u_0,\frac{\pi}{4}]$. Mais l'étude de $g$ a montré que $f$ n'admet pas de point fixe dans cet intervalle ($u_0$ étant strictement positif). On aboutit à une contradiction.

Donc, ou bien $u_0\in[\frac{\pi}{4},1]$, ou bien $u_0\in]0,\frac{\pi}{4}[$ et dans ce cas, $\exists n_0\in\Nn/\;u_{n_0}\in[\frac{\pi}{4},1]$. Dans tous les cas, $\exists n_0\in\Nn/\;u_{n_0}\in[\frac{\pi}{4},1]$. Mais alors, puisque $f([\frac{\pi}{4},1])=[\sin2,\sin\frac{\pi}{2}]\subset[\frac{\pi}{4},1]$ (car $\sin2=0,909...>0,785...=\frac{\pi}{4}$), pour tout entier $n\geq n_0$, $u_n\in[\frac{\pi}{4},1]$.

Pour $x\in[\frac{\pi}{4},1]$, $|g'(x)|=|2\cos(2x)|\leq|2\cos2|$. L'inégalité des accroissements finis montre alors que $\forall n\geq n_0,\;|u_{n+1}-\alpha|\leq|2\cos2|.|u_n-\alpha|$, puis que 

$$\forall n\geq n_0,\;|u_n-\alpha|\leq|2\cos2|^{n-n_0}|u_{n_0}-\alpha|.$$

Comme $|2\cos2|=0,83...<1$, on en déduit que la suite $u$ converge vers $\alpha$. La machine donne par ailleurs $\alpha=0,947...$.

$$\includegraphics{../images/img005425-5}$$


\item  Pour $x\in\Rr$,

$$x^2-2x+2=x\Leftrightarrow x^2-3x+2=0\Leftrightarrow(x-1)(x-2)=0\Leftrightarrow x=1\;\mbox{ou}\;x=2.$$

Donc, si la suite $u$ converge, ce ne peut être que vers $1$ ou $2$.

Pour $n\in\Nn$,

$$\begin{array}{l}
u_{n+1}-u_n=(u_n^2-2u_n+2)-u_n=(u_n-1)(u_n-2)\quad(I)\\
u_{n+1}-1=u_n^2-2u_n+1=(u_n-1)^2\quad(II)\\
u_{n+1}-2=u_n^2-2u_n=u_n(u_n-2)\quad(III).
\end{array}
$$

\begin{itemize}
\item[\textbf{1er cas.}] Si $u_0=1$ ou $u_0=2$, la suite $u$ est constante.
\item[\textbf{2ème cas.}] Si $u_0\in]1,2[$, $(II)$ et $(III)$ permettent de montrer par récurrence que $\forall n\in\Nn,\;u_n\in]1,2[$. $(I)$ montre alors que la suite $u$ est strictement décroissante. Etant minorée par $1$, elle converge vers un réel $\ell\in[1,u_0]\subset[1,2[$. Dans ce cas, la suite $(u_n)$ converge vers $1$.
\item[\textbf{3ème cas.}] Si $u_0\in]2,+\infty[$, $(III)$ permet de montrer par récurrence que $\forall n\in\Nn,\;u_n>2$. Mais alors, $(I)$ montre que la suite $u$ est strictement croissante. Si $u$ converge, c'est vers un réel $\ell\in[u_0,+\infty[\subset]2,+\infty[$. $f$ n'ayant pas de point fixe dans cet intervalle, la suite $u$ diverge et, $u$ étant strictement croissante, on a $\lim_{n\rightarrow +\infty}u_n=+\infty$.
\item[\textbf{4ème cas.}] Si $u_0\in]0,1[$, alors $u_1=(u_0-1)^2+1\in]1,2[$ ce qui ramène au deuxième cas. La suite $u$ converge vers $1$.
\item[\textbf{5ème cas.}] Si $u_0=0$, alors $u_1=2$ et la suite $u$ est constante à partir du rang $1$. Dans ce cas, la suite $u$ converge vers $2$.
\item[\textbf{6ème cas.}] Si $u_0<0$, alors $u_1=u_n^2-2u_n+2>2$, ce qui ramène au troisième cas. La suite $u$ tend vers $+\infty$.
\end{itemize}

En résumé, si $u_0\in]0,2[$, la suite $u$ converge vers $1$, si $u_0\in\{0,2\}$, la suite $u$ converge vers $2$ et si $u_0\in]-\infty,0[\cup]2,+\infty[$, la suite $u$ tend vers $+\infty$.

$$\includegraphics{../images/img005425-6}$$

\end{enumerate}
\fincorrection
\correction{005435}
\begin{enumerate}
 \item  
Pour $x\in\left[0,\frac{\pi}{2}\right]$, posons $f(x)=\sin x$. On a $f\left(\left]0,\frac{\pi}{2}\right]\right)=]0,1]\subset\left]0,\frac{\pi}{2}\right]$. Donc, puisque $u_0\in\left]0,\frac{\pi}{2}\right]$, on en déduit que $\forall n\in\Nn,\;u_n\in\left]0,\frac{\pi}{2}\right]$.\rule[-5mm]{0mm}{10mm}
Il est connu que $\forall x\in\left]0,\frac{\pi}{2}\right]$, $\sin x<x$ et de plus, pour $x\in\left[0,\frac{\pi}{2}\right]$, $\sin x=x\Leftrightarrow x=0$.
La suite $u$ est à valeurs dans $\left]0,\frac{\pi}{2}\right]$ et donc $\forall n\in\Nn,\;u_{n+1}=\sin(u_n)<u_n$. La suite $u$ est donc strictement décroissante et, étant minorée par $0$, converge vers un réel $\ell$ de $\left[0,\frac{\pi}{2}\right]$ qui vérifie ($f$ étant continue sur le segment $\left[0,\frac{\pi}{2}\right]$) $f(\ell)=\ell$ ou encore $\ell=0$.
En résumé,

\begin{center}
\shadowbox{
la suite $u$ est strictement positive, strictement décroissante et converge vers $0$.
}
\end{center}
 \item  Soit $\alpha$ un réel quelconque. Puisque la suite $u$ tend vers 0 , on a

\begin{align*}\ensuremath
u_{n+1}^{\alpha}-u_n^{\alpha}=(\sin u_n)^{\alpha}-u_n^{\alpha}&\underset{n\rightarrow+\infty}{=}\left(u_n-\frac{u_n^3}{6}+o(u_n^3)\right)^{\alpha}-u_n^{\alpha}\\
 &=u_n^{\alpha}\left(\left(1-\frac{u_n^2}{6}+o(u_n^2)\right)^{\alpha}-1\right)=u_n^{\alpha}\left(-\alpha\frac{u_n^2}{6}+o(u_n^2)\right)\\
 &=-\alpha\frac{u_n^{\alpha+2}}{6}+o(u_n^{\alpha+2})
\end{align*}
Pour $\alpha=-2$ on a donc 

$$\frac{1}{u_{n+1}^2}-\frac{1}{u_n^2}=\frac{1}{3}+o(1).$$
D'après le lemme de \textsc{Cesaro}, $\frac{1}{n}\sum_{k=0}^{n-1}\left(\frac{1}{u_{k+1}^2}-\frac{1}{u_k^2}\right)=\frac{1}{3}+o(1)$ ou encore $\frac{1}{n}\left(\frac{1}{u_n^2}-\frac{1}{u_0^2}\right)=\frac{1}{3}+o(1)$ ou enfin, 

\begin{center}
$\frac{1}{u_n^2}\underset{n\rightarrow+\infty}{=}\frac{n}{3}+\frac{1}{u_0^2}+o(n)\underset{n\rightarrow+\infty}{=}\frac{n}{3}+o(n)\underset{n\rightarrow+\infty}{\sim}\frac{n}{3}$.
\end{center}
Par suite, puisque la suite $u$ est strictement positive, 

\begin{center}
\shadowbox{
$u_n\underset{n\rightarrow+\infty}{\sim}\sqrt{\frac{3}{n}}.$
}
\end{center}
\end{enumerate}
\fincorrection
\correction{005436}
Il est immédiat par récurrence que $\forall n\in\Nn,\;u_n>0$. Donc, $\forall n\in\Nn^,\;\frac{u_{n+1}}{u_n}=e^{-u_n}<1$ et donc, puisque la suite $u$ est stritement positive, $u_{n+1}<u_n$. La suite $u$ est strictement décroissante, minorée par $0$ et donc converge vers un réel $\ell$ vérifiant $\ell=\ell e^{-\ell}$ ou encore $\ell(1-e^{-\ell})=0$ ou encore $\ell=0$.

\begin{center}
\shadowbox{
$u$ est strictement positive, strictement décroissante et converge vers $0$.
}
\end{center}
Soit $\alpha$ un réel quelconque. Puisque la suite $u$ tend vers $0$,

$$u_{n+1}^{\alpha}-u_n^{\alpha}=u_n^{\alpha}(e^{-\alpha u_n}-1)=u_n^{\alpha}(-\alpha u_n+o(u_n))=-\alpha u_n^{\alpha+1}+o(u_n^{\alpha+1}).$$
Pour $\alpha=-1$, on obtient en particulier $\frac{1}{u_{n+1}}-\frac{1}{u_n}=1+o(1)$. Puis, comme au numéro précédent, $\frac{1}{u_n}=n+\frac{1}{u_0}+o(n)\underset{n\rightarrow+\infty}{\sim}n$ et donc 

\begin{center}
\shadowbox{
$u_n\underset{n\rightarrow+\infty}{\sim}\frac{1}{n}$.
}
\end{center}
\fincorrection
\correction{000563}
Remarquons d'abord que $1-\frac{1}{k^2} = \frac{1-k^2}{k^2} = \frac{(k-1)(k+1)}{k.k}$.
En \'ecrivant les fractions de $u_n$ sous la cette forme, l'\'ecriture va se simplifier radicalement:
$$u_n = \frac{(2-1)(2+1)}{2.2}\frac{(3-1)(3+1)}{3.3}\cdots \frac{(k-1)(k+1)}{k.k}\frac{(k)(k+2)}{(k+1).(k+1)}\cdots \frac{(n-1)(n+1)}{n.n}$$
Tous les termes des num\'erateurs se retrouvent au d\'enominateur (et vice-versa), sauf aux extr\'emit\'es. D'o\`u:
$$u_n = \frac12\frac{n+1}{n}.$$
Donc $(u_n)$ tends vers $\frac12$ lorsque $n$ tend vers $+\infty$.
\fincorrection
\nocorrection
\nocorrection
\nocorrection
\nocorrection
\correction{000568}
\begin{enumerate}
  \item $0$.
  \item $1$.
  \item $7/30$.
  \item $1/2$.
  \item $1$.
  \item $-3/2$.
  \item $1$.
  \item $3$.
  \item $1$~; $2$.
  \item $3/4$.
  \item $0$.
  \item $0$.
  \item $1/3$.
\end{enumerate}
\fincorrection
\correction{000569}
\begin{enumerate}
\item \begin{align*}
       u_{n+1}^2-a &= \frac14\left(\frac{u_n^2+a}{u_n}\right)^2-a\\
                   &= \frac1{4u_n^2}(u_n^4-2au_n^2+a^2)\\
                   &= \frac14 \frac{(u_n^2-a)^2}{u_n^2}\\
  \end{align*}
\item Il est clair que pour $n\geqslant 0$ on a $u_n > 0$.
D'apr\`es l'\'egalit\'e pr\'ec\'edente pour $n\geqslant 0$, $u_{n+1}^2-a \geqslant 0$ et
comme $ u_{n+1}$ est positif alors $u_{n+1}\geqslant \sqrt a$.

Soit $n\geqslant 1$. Calculons le quotient de $u_{n+1}$ par $u_n$ :
$\frac{ u_{n+1}}{ u_n} = \frac12\left(1+\frac{a}{u_n^2}\right)$ or
$\frac{a}{u_n^2}\leqslant 1$ car $u_n \geqslant \sqrt a$. Donc $\frac{
u_{n+1}}{ u_n} \leqslant 1$ et donc $u_{n+1} \leqslant  u_n $. La suite
$(u_n)_{n\geqslant 1}$ est donc d\'ecroissante.
\item La suite $(u_n)_{n\geqslant 1}$ est d\'ecroissante et minor\'ee par $\sqrt a$ donc elle converge vers une limite $\ell>0$.
D'apr\`es la relation
$$u_{n+1} = \frac12\left(u_n+\frac{a}{u_n}\right)$$
quand $n\rightarrow + \infty$ alors $u_n \rightarrow \ell$ et
$u_{n+1} \rightarrow \ell$. \`A la limite nous obtenons la
relation
$$\ell = \frac12\left(\ell+\frac{a}{\ell}\right).$$
La seule solution positive est $\ell = \sqrt a$. Conclusion
$(u_n)$ converge vers $\sqrt a$.
\item La relation
$$ u_{n+1}^2-a =  \frac{(u_n^2-a)^2}{4u_n^2}$$
s'\'ecrit aussi
$$ (u_{n+1}-\sqrt a)(u_{n+1}+\sqrt a) = \frac{(u_n-\sqrt a)^2(u_n+\sqrt a)^2}{4u_n^2}.$$
Donc
\begin{align*}
       u_{n+1}-\sqrt a &= (u_n-\sqrt a)^2 \frac{1}{4(u_{n+1}+\sqrt a)}\left(\frac{u_n+\sqrt a}{u_n}\right)^2\\
                   &\leqslant (u_n-\sqrt a)^2 \frac{1}{4(2\sqrt a)}\left(1+\frac{\sqrt a}{u_n}\right)^2\\
                   &\leqslant (u_n-\sqrt a)^2  \frac{1}{2\sqrt a}\\
  \end{align*}
\item Par r\'ecurrence pour $n=1$, $u_1-\sqrt a \leqslant k$.
Si la proposition est vraie rang $n$, alors
\begin{align*}
       u_{n+1}-\sqrt a &\leqslant \frac{1}{2\sqrt a} (u_n-\sqrt a)^2  \\
       &\leqslant \frac{1}{2\sqrt a} (2\sqrt a)^2\left(\left( \frac{k}{2\sqrt a} \right)^{2^{n-1}} \right)^2\\
       &\leqslant 2\sqrt a \left( \frac{k}{2\sqrt a} \right)^{2^n}\\
\end{align*}
\item Soit $u_0=3$, alors $u_1 = \frac12(3+\frac{10}{3}) = 3,166\ldots$.
Comme $3\leqslant \sqrt{10} \leqslant u_1$ donc $u_1-\sqrt{10} \le
0.166\ldots$. Nous pouvons choisir $k=0,17$. Pour que l'erreur
$u_n-\sqrt a$ soit inf\'erieure \`a $10^{-8}$ il suffit de calculer le
terme $u_4$ car alors l'erreur (calcul\'ee par la formule de la
question pr\'ec\'edente) est inf\'erieure \`a $1,53\times 10^{-10}$. Nous
obtenons $u_4 = 3,16227766\ldots$
Bilan $\sqrt{10} =  3,16227766\ldots$ avec une précision de $8$ chiffres après la virgule. 
Le nombre de chiffres exacts double à chaque itération, avec $u_5$ nous aurions (au moins) $16$ chiffres exacts,
et avec $u_6$ au moins $32$\ldots
\end{enumerate}
\fincorrection
\correction{000570}
\begin{enumerate}
  \item La suite $(u_n)$ est strictement croissante, en effet $u_{n+1}-u_n = \frac{1}{(n+1)!} > 0$. La suite $(v_n)$ est strictement d\'ecroissante :
$$v_{n+1}-v_n = u_{n+1}-u_n + \frac{1}{(n+1)!}-\frac{1}{n!}= \frac{1}{(n+1)!}+ \frac{1}{(n+1)!}-\frac{1}{n!}= \frac{1}{n!}(\frac 2n-1).$$
Donc \`a partir de $n\geq 2$, la suite $(v_n)$ est strictement d\'ecroissante.
  \item Comme $u_n \leq v_n \leq v_2$, alors $(u_n)$ est une suite croissante et major\'ee. Donc elle converge vers $\ell \in \Rr$.
De m\^eme $v_n \geq u_n \geq u_0$, donc  $(v_n)$ est une suite d\'ecroissante et minor\'ee. Donc elle converge vers $\ell' \in \Rr$.
De plus $v_n -u_n = \frac1{n!}$. Et donc $(v_n-u_n)$ tend vers $0$
ce qui prouve que $\ell=\ell'$.
  \item Supposons que $\ell \in \Qq$, nous \'ecrivons alors $\ell = \frac pq$ avec $p,q \in \Nn$. Nous obtenons pour $n\geq 2$:
$$u_n \leq \frac pq \leq v_n.$$
Ecrivons cette \'egalit\'e pour $n=q$: 
$u_q \leq \frac pq \leq v_q$ et multiplions par $q!$:
$q! u_q \leq q!\frac pq \leq q! v_q$. Dans cette double in\'egalit\'e toutes les termes sont des entiers ! De plus $v_q = u_q +\frac 1{q!}$ donc:
$$q! u_q \leq q! \frac pq \leq q! u_q + 1.$$
Donc l'entier $q! \frac pq$ est \'egal \`a l'entier $q! u_q$
ou \`a $q! u_q + 1 = q! v_q$. Nous obtenons que $\ell = \frac pq$
est \'egal \`a $u_q$ ou \`a $v_q$. Supposons par exemple que $\ell = u_q$,
comme la suite $(u_n)$ est strictement croissante alors $u_q  < u_{q+1} < \cdots < \ell$, ce qui aboutit \`a une contradiction. Le m\^eme raisonnement s'applique en supposant $\ell = v_q$ car la suite $(v_n)$ est strictement d\'ecroissante. Pour conclure nous avons montré que $\ell$ n'est pas un nombre rationnel.
\end{enumerate}

En fait $\ell$ est le nombre $e = \exp(1)$.
\fincorrection
\correction{000571}
\begin{enumerate}
    \item Si $u_0 \leqslant u_1$ alors comme $f$ est croissante $f(u_0)\leqslant f(u_1)$ donc $u_1 \leqslant u_2$, ensuite $f(u_1)\leqslant f(u_2)$ soit $u_2 \leqslant u_3$,... Par r\'ecurrence on montre que $(u_n)$ est d\'ecroissante. Comme elle est minor\'ee par $a$ alors elle converge. Si $u_0 \leqslant u_1$ alors la suite $(u_n)$ est croissante et major\'ee par $b$ donc converge.

Notons $\ell$ la limite de $(u_n)_n$. Comme $f$ est continue alors
$(f(u_n))$ tend vers $f(\ell)$. De plus la limite de $(u_{n+1})_n$
est aussi $\ell$. En passant \`a la limite dans l'expression
$u_{n+1}=f(u_n)$ nous obtenons l'\'egalit\'e $\ell = f(\ell)$.

  \item La fonction $f$ définie par $f(x) = \frac{4x+5}{x+3}$ est continue et d\'erivable sur l'intervalle $[0,4]$ et $f([0,4])\subset [0,4]$.
La fonction $f$ est croissante (calculez sa d\'eriv\'ee). Comme $u_0 =
4$ et $u_1= 3$ alors $(u_n)$ est d\'ecroissante. Calculons la valeur
de sa limite $\ell$. $\ell$ est solution de l'\'equation $f(x)=x$
soit $4x+5=x(x+3)$. Comme $u_n \geqslant 0$ pour tout $n$ alors $\ell
\geqslant 0$. La seule solution positive de l'équation du second degré $4x+5=x(x+3)$ est $\ell =
\frac{1+\sqrt{21}}{2}=2,7912\ldots$

  \item Si $f$ est d\'ecroissante alors $f\circ f$ est croissante (car $x\leqslant y \Rightarrow
f(x)\geqslant f(y) \Rightarrow f\circ f(x)\leqslant  f\circ f(y)$). Nous
appliquons la premi\`ere question avec la fonction $f\circ f$. La
suite $(u_0, u_2 = f\circ f(u_0),u_4 = f\circ f(u_2),\ldots)$ est
monotone et convergente. De m\^eme pour la suite $(u_1, u_3 =
f\circ f(u_1),u_5 = f\circ f(u_3),\ldots)$.

  \item La fonction $f$ définie par $f(x) = (1-x)^2$ est continue et d\'erivable de $[0,1]$ dans $[0,1]$.
Elle est d\'ecroissante sur cet intervalle. Nous avons $u_0 =
\frac12$, $u_1=\frac14$, $u_2=\frac{9}{16}$, $u_3 =
0,19\ldots$,... Donc la suite $(u_{2n})$ est croissante, nous
savons qu'elle converge et notons $\ell$ sa limite. La suite
$(u_{2n+1})$ et d\'ecroissante,  notons $\ell'$ sa limite. Les
limites $\ell$ et $\ell'$ sont des solutions de l'\'equation
$f\circ f(x)=x$. Cette \'equation s'\'ecrit $(1-f(x))^2=x$, ou encore
$(1-(1-x)^2)^2=x$ soit $x^2(2-x)^2=x$. Il y a deux solutions
\'evidentes $0$ et $1$. Nous factorisons le polyn\^ome
$x^2(2-x)^2-x$ en $x(x-1)(x-\lambda)(x-\mu)$ avec $\lambda$ et
$\mu$ les solutions de l'\'equation $x^2-3x+1$ : $\lambda =
\frac{3-\sqrt{5}}{2} = 0,3819\ldots$ et $\mu =
\frac{3+\sqrt{5}}{2} > 1$. Les solutions de l'\'equation $f\circ
f(x)=x$ sont donc $\{ 0,1,\lambda, \mu\}$. Comme $(u_{2n})$ est
croissante et que $u_0 = \frac12$ alors  $(u_{2n})$ converge vers
$\ell=1$ qui est le seul point fixe de $[0,1]$ sup\'erieur \`a
$\frac12$. Comme $(u_{2n+1})$ est d\'ecroissante et que $u_1 =
\frac14$ alors  $(u_{2n+1})$ converge vers $\ell'=0$ qui est le
seul point fixe de $[0,1]$ inf\'erieur \`a $\frac14$.
\end{enumerate}
\fincorrection
\correction{000572}
\begin{enumerate}
    \item  Soient $a,b >0$. On veut d\'emontrer
    que $\sqrt{ab}\leqslant \frac{a+b}{2}$. Comme les deux membres de cette
in\'egalit\'e sont positifs, cette in\'egalit\'e est \'equivalente
\`a $ ab\leqslant (\frac{a+b}{2})^2$. De plus,
$$ ab\leqslant \left( \frac{a+b}{2}\right)^2  \Leftrightarrow 4ab\leqslant a^2+2ab+b$$
$$ \Leftrightarrow 0\leqslant a^2-2ab+b^2$$ ce qui
est toujours vrai car $a^2-2ab+b^2=(a-b)^2$ est un carr\'e parfait. On a
donc bien l'in\'egalit\'e voulue.

  \item  Quitte \`a \'echanger
$a$ et $b$ (ce qui ne change pas les moyennes arithm\'etique et
g\'eom\'etrique, et qui pr\'eserve le fait d'\^etre compris entre
$a$ et $b$), on peut supposer que $a\leqslant b$. Alors en ajoutant les
deux in\'egalit\'es $$a/2 \leqslant a/2 \leqslant b/2$$ $$a/2 \leqslant b/2 \leqslant
b/2,$$ on obtient $$a\leqslant \frac{a+b}{2}\leqslant b.$$

De m\^eme, comme tout est positif, en multipliant les deux
in\'egalit\'es
$$\sqrt{a} \leqslant \sqrt{a} \leqslant \sqrt{b}$$ $$\sqrt{a} \leqslant \sqrt{b} \leqslant
\sqrt{b}$$ on obtient $$a\leqslant \sqrt{ab} \leqslant b.$$

    \item  Il faut avant tout remarquer que pour tout $n$,
    $u_n$ et $v_n$ sont strictement positifs, ce qui
permet de dire que les deux suites sont bien d\'efinies. On le
d\'emontre par r\'ecurrence: c'est clair pour $u_0$ et $v_0$, et
si $u_n$ et $v_n$ sont strictement positifs alors leurs moyennes
g\'eom\'etrique (qui est $u_{n+1}$) et arithm\'etique (qui est $v_{n+1}$) sont
strictement positives.
    \begin{enumerate}
\item  On veut montrer que pour chaque $n$, $u_n\leqslant v_n$. L'in\'egalit\'e est claire pour $n=0$
     gr\^ace aux hypoth\`eses faites sur $u_0$ et $v_0$.
     Si maintenant $n$ est plus grand que 1, $u_{n}$ est la
     moyenne g\'eom\'etrique de $u_{n-1}$ et $v_{n-1}$ et $v_{n}$
     est la moyenne arithm\'etique de $u_{n-1}$ et $v_{n-1}$,
     donc, par 1., $u_n\leqslant v_n$.

 \item  On sait d'apr\`es 2. que $u_n\leqslant u_{n+1}\leqslant v_n$.
 En particulier, $u_n\leqslant u_{n+1}$ i.e. $(u_n)$ est croissante.
 De m\^eme, d'apr\`es 2., $u_n\leqslant v_{n+1}\leqslant v_n$. En particulier,
 $v_{n+1}\leqslant v_n$ i.e. $(v_n)$ est d\'ecroissante.

        \item  Pour tout $n$, on a $u_0\leqslant u_n\leqslant v_n\leqslant v_0$.
        $(u_n)$ est donc croissante et major\'ee, donc converge
        vers une limite $\ell$. Et $(v_n)$ est d\'ecroissante et
        minor\'ee et donc converge vers une limite $\ell'$. 
Nous savons maintenant que 
$u_{n} \rightarrow \ell$, donc aussi $u_{n+1} \rightarrow \ell$, et $v_{n} \rightarrow \ell'$ ;
la relation $u_{n+1}=\sqrt{u_n v_n}$ s'écrit à la limite :
$$\ell=\sqrt{\ell\ell'}.$$
De même la relation $v_{n+1}=\frac{u_n+v_n}{2}$ donnerait à la limite :
$$\ell'=\frac{\ell+\ell'}{2}.$$
Un petit calcul avec l'une ou l'autre de ces égalités implique $\ell=\ell'$.
    \end{enumerate}
\end{enumerate}

Il y a une autre m\'ethode un peu plus longue mais toute aussi
valable.

{\bf D\'efinition} Deux suites $(u_n)$ et $(v_n)$ sont dites {\em
adjacentes} si
\begin{enumerate}
\item $u_n\leqslant v_n$,
\item $(u_n)$ est croissante et $(v_n)$ est d\'ecroissante,
\item $\lim (u_n-v_n)=0$.
\end{enumerate}

Alors, on a le th\'eor\`eme suivant:\\
 \textbf{Th\'eor\`eme} : Si
$(u_n)$ et $(v_n)$ sont deux suites adjacentes, elles sont toutes les
deux convergentes et ont la m\^eme limite.

Pour appliquer ce th\'eor\`eme, vu qu'on sait d\'ej\`a que $(u_n)$
et $(v_n)$ v\'erifient les points 1 et 2 de la d\'efinition, il
suffit de d\'emontrer que $\lim (u_n-v_n)=0$. On a d'abord que
$v_n-u_n\geqslant 0$. Or, d'apr\`es (a) $$v_{n+1}-u_{n+1}{\le}
v_{n+1}-u_n=\frac{v_n-u_n}{2}. $$

Donc, si on note $w_n= v_n-u_n$, on a que $0\leqslant w_{n+1}\leqslant w_n/2$.
Donc, on peut d\'emontrer (par r\'ecurrence) que $0\leqslant w_n\le
\frac{w_0}{2^n}$, ce qui implique que $\lim_{n\rightarrow\infty}w_n=0$.
Donc $v_n-u_n$ tend vers 0, et ceci termine de d\'emontrer que les
deux suites $(u_n)$ et $(v_n)$ sont convergentes et ont m\^eme limite
en utilisant le th\'eor\`eme sur les suites adjacentes.
\fincorrection
\nocorrection
\correction{000574}
Notons $f_n : [0,1] \longrightarrow \Rr$ la fonction d\'efinie par :
$$ f_n(x) = \sum_{k=1}^n x^k \ \ - 1.$$
\begin{enumerate}
  \item La fonction $f_n$ est continue
sur $[0,1]$. De plus $f_n(0) = -1 < 0$ et $f_n(1) = n-1\geqslant 0$.
D'apr\`es le th\'eor\`eme des valeurs interm\'ediaires, $f_n$, admet un
z\'ero dans l'intervalle $[0,1]$. De plus elle strictement
croissante (calculez sa d\'eriv\'ee) sur $[0,1]$ donc ce z\'ero est
unique.
  \item Calculons $f_n(a_{n-1})$.
\begin{align*}
f_n(a_{n-1}) &= \sum_{k=1}^{n} a_{n-1}^k  - 1 \\
     &= a_{n-1}^n + \sum_{k=1}^{n-1} a_{n-1}^k  - 1 \\
     &= a_{n-1}^n + f_{n-1}(a_{n-1})  \\
     &= a_{n-1}^n \text{\ \  (car $f_{n-1}(a_{n-1})=0$ par d\'efinition de $a_{n-1}$).}
\end{align*}

Nous obtenons l'in\'egalit\'e
$$ 0 = f_n(a_n) < f_n(a_{n-1}) = a_{n-1}^n.$$
Or $f_n$ est strictement croissante, l'in\'egalit\'e ci-dessus
implique donc $ a_n < a_{n-1}$.
Nous venons de d\'emontrer que la suite $(a_n)_n$ est d\'ecroissante.


Remarquons avant d'aller plus loin que $f_n(x)$ est la somme d'une
suite g\'eom\'etrique :
$$f_n(x) = \frac{1-x^{n+1}}{1-x}-2.$$

\'Evaluons maintenant $f_n(\frac12)$, \`a l'aide de l'expression
pr\'ec\'edente
$$f_n(\frac12) = \frac{1-(\frac12)^{n+1}}{1-{\frac12}}-2 = -\frac 1 {2^n} < 0.$$
Donc $ f_n(\frac12) < f_n(a_n)=0$ entra\^{\i}ne $\frac12 < a_n$.

Pour r\'esumer, nous avons montré que la suite $(a_n)_n$ est
strictement d\'ecroissante et minor\'ee par $\frac12$.

\item Comme $(a_n)_n$ est
d\'ecroissante et minor\'ee par $\frac12$ alors elle converge, nous
notons $\ell$ sa limite :
$$ \frac 12 \leqslant \ell < a_n.$$
Appliquons $f_n$ (qui est strictement croissante) \`a cette
in\'egalit\'e :
 $$ f_n\left(\frac 12\right) \leqslant f_n(\ell) < f_n(a_n),$$
qui s'\'ecrit aussi :
$$ -\frac 1 {2^n} \leqslant f_n(\ell) < 0,$$
et ceci quelque soit $n\geqslant 1$. La suite $(f_n(\ell))_n$ converge
donc vers $0$ (th\'eor\`eme des  ``gendarmes''). Mais nous savons
aussi que
$$f_n(\ell) = \frac{1-\ell^{n+1}}{1-\ell}-2 ;$$
donc $(f_n(\ell))_n$ converge vers $\frac{1}{1-\ell}-2$ car
$(\ell^n)_n$ converge vers $0$. Donc
$$\frac{1}{1-\ell}-2 = 0, \text{ d'o\`u } \ell = \frac12.$$
\end{enumerate}
\fincorrection
\nocorrection
\nocorrection
\nocorrection
\nocorrection
\nocorrection
\nocorrection
\nocorrection
\nocorrection
\nocorrection
\nocorrection
\nocorrection
\nocorrection
\nocorrection
\nocorrection
\nocorrection
\nocorrection
\nocorrection
\nocorrection
\correction{005252}
\begin{enumerate}
\item  Tout d'abord, pour $n\geq1$, $\frac{n-1}{n}$ existe et est élément de $[-1,1]$. Donc, $\Arccos\frac{n-1}{n}$ existe pour tout entier naturel non nul $n$.

Quand $n$ tend vers $+\infty$, $\frac{n-1}{n}$ tend vers $1$ et donc $\Arccos\frac{n-1}{n}$ tend vers $0$. Mais alors,

$$\Arccos\frac{n-1}{n}\sim\sin(\Arccos\frac{n-1}{n})=\sqrt{1-(\frac{n-1}{n})^2}=\frac{\sqrt{2n-1}}{n}\sim\frac{\sqrt{2}}{\sqrt{n}}.$$

\item  $\Arccos\frac{1}{n}$ tend vers $1$ et donc $\Arccos\frac{1}{n}\sim1$.

\item  $\ch(\sqrt{n})=\frac{1}{2}(e^{\sqrt{n}}+e^{-\sqrt{n}})\sim\frac{1}{2}e^{\sqrt{n}}$.

\item  $n\ln(1+\frac{1}{n})\sim n.\frac{1}{n}=1$ et donc, $\left(1+\frac{1}{n}\right)^n=e^{n\ln(1+1/n)}$ tend vers $e$. Par suite, $\left(1+\frac{1}{n}\right)^n\sim e$.

\item  $\Argch n$ existe pour $n\geq1$ et comme, pour $n\geq1$, $n^4+n^2-1\geq n^4>0$, $\frac{\Argch n}{\sqrt{n^4+n^2-1}}$ existe pour $n\geq1$.

$$\Argch n=\ln(n+\sqrt{n^2-1})\sim\ln(n+n)=\ln(2n)=\ln n+\ln2\sim\ln n.$$

Donc, $\frac{\Argch n}{\sqrt{n^4+n^2-1}}\sim\frac{\ln n}{\sqrt{n^4}}=\frac{\ln n}{n^2}$.

\item   $-\sqrt{n}\ln(\sqrt{n}+1)=-\sqrt{n}\ln(\sqrt{n})-\sqrt{n}\ln(1+\frac{1}{\sqrt{n}})=-\sqrt{n}\ln(\sqrt{n})-\sqrt{n}(\frac{1}{\sqrt{n}}+o(\frac{1}{\sqrt{n}}))=-\sqrt{n}\ln(\sqrt{n})-1+o(1)$, et donc

$$(1+\sqrt{n})^{-\sqrt{n}}=e^{-\sqrt{n}\ln(\sqrt{n})-1+o(1)}\sim e^{-\sqrt{n}\ln(\sqrt{n})-1}=\frac{1}{e}\frac{1}{\sqrt{n}^{\sqrt{n}}}.$$
\item 

$$\ln(\cos\frac{1}{n})(\ln\sin\frac{1}{n})\sim(\cos\frac{1}{n}-1)\ln(\frac{1}{n})\sim(-\frac{1}{2n^2})(-\ln n)=\frac{\ln n}{2n^2}.$$

\item  $(\Arctan n)^{3/5}=(\frac{\pi}{2}-\Arctan\frac{1}{n})^{3/5}=(\frac{\pi}{2})^{3/5}(1-\frac{2}{\pi}(\frac{1}{n}+o(\frac{1}{n})))^{3/5}=(\frac{\pi}{2})^{3/5}(1-\frac{6}{5n\pi}+o(\frac{1}{n}))$, et donc

$$(\frac{\pi}{2})^{3/5}-(\Arctan n)^{3/5}=(\frac{\pi}{2})^{3/5}(1-1+\frac{6}{5n\pi}+o(\frac{1}{n}))
\sim(\frac{\pi}{2})^{3/5}\frac{6}{5n\pi}$$

\item  Tout d'abord, pour $n\geq1$, $\left|\frac{(-1)^n}{\sqrt{n}}\right|=\frac{1}{\sqrt{n}}\leq 1$, et donc $1+\frac{(-1)^n}{\sqrt{n}}\geq0$, puis $\sqrt{1+\frac{(-1)^n}{\sqrt{n}}}-1$ existe. Ensuite, quand $n$ tend vers $+\infty$,

$$\sqrt{1+\frac{(-1)^n}{\sqrt{n}}}-1\sim\frac{(-1)^n}{2\sqrt{n}}.$$
\end{enumerate}
\fincorrection
\correction{005253}
Pour $n\geq2$, on a

$$\frac{1}{n!}\sum_{k=0}^{n}k!=1+\frac{1}{n}+\sum_{k=0}^{n-2}\frac{k!}{n!}.$$

Mais, pour $0\leq k\leq n-2$, $\frac{k!}{n!}=\frac{1}{n(n-1)...(k+1)}\leq\frac{1}{n(n-1)}$ (le produit contenant au moins les deux premiers facteurs. Par suite,

$$0\leq\sum_{k=0}^{n-2}\frac{k!}{n!}\leq\frac{n-2}{n(n-1)}.$$

On en déduit que $\sum_{k=0}^{n-2}\frac{k!}{n!}$ tend vers $0$ quand $n$ tend vers $+\infty$. Comme $\frac{1}{n}$ tend aussi vers $0$ quand $n$ tend vers $+\infty$, on en déduit que $\frac{1}{n!}\sum_{k=0}^{n}k!$ tend vers $1$ et donc que

$$\sum_{k=0}^{n}k!\sim n!.$$
\fincorrection
\correction{005254}
\begin{enumerate}
\item  Soit $\varepsilon>0$.

Les suites $u$ et $v$ sont équivalentes et la suite $v$ est strictement positive. Donc, il existe un rang $n_0$ tel que, pour $n\geq n_0$, $|u_n-v_n|<\frac{\varepsilon}{2}v_n$. Soit $n>n_0$.

\begin{align*}\ensuremath
\left|\frac{U_n}{V_n}-1\right|&=\frac{|U_n-V_n|}{V_n}\leq\frac{1}{V_n}\sum_{k=0}^{n}|u_k-v_k|\\
 &\leq \frac{1}{V_n}(\sum_{k=0}^{n_0}|u_k-v_k|+\frac{\varepsilon}{2}\sum_{k=n_0+1}^{n}v_k)\\
 &\leq\frac{1}{V_n}(\sum_{k=0}^{n_0}|u_k-v_k|+\frac{\varepsilon}{2}V_n)=\frac{1}{V_n}\sum_{k=0}^{n_0}|u_k-v_k|+\frac{\varepsilon}{2}
\end{align*}

Maintenant, l'expression $\sum_{k=0}^{n_0}|u_k-v_k|$ est constante quand $n$ varie, et d'autre part, $V_n$ tend vers $+\infty$ quand $n$ tend vers $+\infty$. On en déduit que $\frac{1}{V_n}\sum_{k=0}^{n_0}|u_k-v_k|$ tend vers $0$ quand $n$ tend vers $+\infty$. Par suite, il existe un rang $n_1>n_0$ tel que, pour $n\geq n_1$, $\frac{1}{V_n}\sum_{k=0}^{n_0}|u_k-v_k|<\frac{\varepsilon}{2}$.

Pour $n\geq n_1$, on a alors $\left|\frac{U_n}{V_n}-1\right|<\frac{\varepsilon}{2}+\frac{\varepsilon}{2}=\varepsilon$.

On a montré que

$$\forall\varepsilon>0,\;\exists n_1\in\Nn/\;\forall n\in\Nn,\;(n\geq n_1\Rightarrow\left|\frac{U_n}{V_n}-1\right|<\varepsilon.$$

Ainsi, la suite $\frac{U_n}{V_n}$ tend vers $1$ quand $n$ tend vers $+\infty$ et donc $U_n\sim V_n$.

\item  $$2(\sqrt{n+1}-\sqrt{n})=\frac{2}{\sqrt{n+1}+\sqrt{n}}\sim\frac{2}{2\sqrt{n}}=\frac{1}{\sqrt{n}}.$$

De plus, 

$$\sum_{k=1}^{n}2(\sqrt{k+1}-\sqrt{k})=2\sqrt{n+1}-2\sqrt{1}.$$

Cette dernière expression tend vers $+\infty$ quand $n$ tend vers $+\infty$.

En résumé, pour $n\geq1$, $\frac{1}{\sqrt{n}}>0$, $2(\sqrt{n+1}-\sqrt{n})>0$, de plus quand $n$ tend vers $+\infty$, 
$\frac{1}{\sqrt{n}}\sim2(\sqrt{n+1}-\sqrt{n})$ et enfin, $\sum_{k=1}^{n}2(\sqrt{k+1}-\sqrt{k})$ tend vers $+\infty$ quand $n$ tend vers $+\infty$. D'après 1),

$$\sum_{k=1}^{n}\frac{1}{\sqrt{k}}\sim\sum_{k=1}^{n}2(\sqrt{k+1}-\sqrt{k})=2\sqrt{n+1}-2\sqrt{1}\sim2\sqrt{n}.$$

$$(n+1)\ln(n+1)-n\ln n=(n+1-n)\ln n+(n+1)\ln(1+\frac{1}{n})=\ln n+1+o(1)\sim\ln n.$$

Comme $\sum_{k=1}^{n}((k+1)\ln(k+1)-k\ln k)=(n+1)\ln(n+1)$ tend vers $+\infty$ et que les suites considéres sont positives, on en déduit que 

$$\ln(n!)=\sum_{k=1}^{n}\ln k\sim\sum_{k=1}^{n}((k+1)\ln(k+1)-k\ln k)=(n+1)\ln(n+1)\sim n\ln n.$$
\end{enumerate}
\fincorrection
\correction{005255}
Pour $n\geq1$, posons $u_n=\frac{(-1)^n}{\ln n}+\frac{1}{n}$. On a alors 

\begin{align*}
n(u_n+u_{n+1}-\frac{2}{n})&=1+\frac{n}{n+1}-2+n(-1)^n(\frac{1}{\ln n}-\frac{1}{\ln(n+1)})=\frac{(-1)^nn(\ln(n+1)-\ln n)}{\ln n\ln(n+1)}+o(1)\\
 &=\frac{(-1)^nn\ln(1+1/n)}{\ln n\ln(n+1)}+o(1)=\frac{(-1)^n(1+o(1))}{\ln n\ln(n+1)}+o(1)=o(1).
\end{align*}

Donc, $n(u_n+u_{n+1}-\frac{2}{n})=o(1)$, ou encore $u_n+u_{n+1}=\frac{2}{n}+o(\frac{1}{n})$, ou enfin, $u_n+u_{n+1}\sim\frac{2}{n}$. Pourtant, $u_n$ est équivalent à $\frac{(-1)^n}{\ln n}$ et pas du tout à $\frac{1}{n}$ ($|nu_n|=\frac{n}{\ln n}\rightarrow+\infty$).

Supposons maintenant que $u_n+u_{2n}\sim\frac{3}{2n}$ et montrons que $u_n\sim\frac{1}{n}$.

On pose $v_n=u_n-\frac{1}{n}$. Il s'agit maintenant de montrer que $v_n=o(\frac{1}{n})$ sous l'hypothèse $v_n+v_{2n}=o(\frac{1}{n})$.

Soit $\varepsilon>0$. Il existe $n_0\in\Nn$ tel  que, pour $n\geq n_0$, $n|v_n+v_{2n}|<\frac{\varepsilon}{4}$.

Soient $n\geq n_0$ et $p\in\Nn$.

\begin{align*}\ensuremath
|v_n|&=|v_n+v_{2n}-v_{2n}-v_{4n}+...+(-1)^p(v_{2^pn}+v_{2^{p+1}n})+(-1)^{p+1}v_{2^{p+1}n}|
\leq\sum_{k=0}^{p}|v_{2^kn}+v_{2^{k+1}n}|+|v_{2^{p+1}n}|\\
 &\frac{\varepsilon}{4}\sum_{k=0}^{p}\frac{1}{2^kn}+|v_{2^{p+1}n}|=\frac{\varepsilon}{4n}\frac{1-\frac{1}{2^{p+1}}}{1-\frac{1}{2}}+|v_{2^{p+1}n}|\\
 &\leq\frac{\varepsilon}{2n}+|v_{2^{p+1}n}|
\end{align*}

Maintenant, la suite $u$ tend vers $0$, et il en est de même de la suite $v$. Par suite, pour chaque $n\geq n_0$, il est possible de choisir $p$ tel que $|v_{2^{p+1}n}|<\frac{\varepsilon}{2n}$.

En résumé, si $n$ est un entier donné supérieur ou égal à $n_0$, $n|v_n|<\frac{\varepsilon}{2}+\frac{\varepsilon}{2}=\varepsilon$. On a montré que

$$\forall\varepsilon>0,\;\exists n_0\in\Nn/\;\forall n\in\Nn,\;(n\geq n_0\Rightarrow|nv_n|<\varepsilon.$$

Par suite, $v_n=o(\frac{1}{n})$ et donc $u_n=\frac{1}{n}+o(\frac{1}{n})$, ou encore $u_n\sim\frac{1}{n}$.
\fincorrection
\correction{005256}
\begin{enumerate}
\item  Il est immédiat que la suite $u$ est définie et à valeurs dans $[-1,\frac{\pi}{2}]$.

Plus précisément, $u_0\in]0,\frac{\pi}{2}]$, et si pour $n\geq0$, $u_n\in]0,\frac{\pi}{2}]$, alors $u_{n+1}\in]0,1]\subset]0,\frac{\pi}{2}]$. On a montré par récurrence que, $\forall n\in\Nn,\;u_n\in]0,\frac{\pi}{2}]$.

Montrons que pour tout réel $x\in]0,\frac{\pi}{2}]$, on a $\sin x>x$. Pour $x\in[0,\frac{\pi}{2}]$, posons $f(x)=x-\sin x$. $f$ est dérivable sur $[0,\frac{\pi}{2}]$ et pour $x\in[0,\frac{\pi}{2}]$, $f'(x)=1-\cos x$. Par suite, $f'$ est strictement positive sur $]0,\frac{\pi}{2}]$ et donc strictement croissante sur $[0,\frac{\pi}{2}]$. Mais alors, pour $x\in]0,\frac{\pi}{2}]$, on a $f(x)>f(0)=0$.

Soit $n\in\Nn$. Puisque $u_n\in]0,\frac{\pi}{2}]$, on a $u_{n+1}=\sin(u_n)<u_n$. La suite $u$ est donc strictement décroissante. Puisque la suite $u$ est d'autre part minorée par $0$, la suite $u$ converge vers un réel noté $\ell$. Puisque pour tout $n\in\Nn$, $0<u_n\leq\frac{\pi}{2}$, on a $0\leq\ell\leq\frac{\pi}{2}$. Mais alors, par continuité de la fonction $x\mapsto\sin x$ sur $[0,\frac{\pi}{2}]$ et donc en $\ell$, on a

$$\ell=\lim_{n\rightarrow +\infty}u_{n+1}=\lim_{n\rightarrow +\infty}\sin(u_n)=\sin(\lim_{n\rightarrow +\infty}u_n)=\sin(\ell).$$

Or, si $x\in]0,\frac{\pi}{2}]$, $\sin x<x$ et en particulier $\sin x\neq x$. Donc, $\ell=0$.

La suite $u$ est strictement positive, strictement décroissante, de limite nulle.

\item  Soit $\alpha\in\Rr$. Puisque $u_{n}$ tend vers $0$ quand $n$ tend vers $+\infty$,

$$u_{n+1}^\alpha=(\sin(u_n))^\alpha=(u_n-\frac{u_n^3}{6}+o(u_n^3))^\alpha=u_n^\alpha(1-\frac{u_n^2}{6}+o(u_n^2))^\alpha
=u_n^\alpha(1-\frac{\alpha u_n^2}{6}+o(u_n^2))=u_n^\alpha-\frac{\alpha u_n^{2+\alpha}}{6}+o(u_n^{2+\alpha}).$$

et donc, $u_{n+1}^\alpha-u_n^\alpha=-\frac{\alpha u_n^{2+\alpha}}{6}+o(u_n^{2+\alpha})$. En prenant $\alpha=-2$, on obtient alors

$$v_n=\frac{1}{u_{n+1}^2}-\frac{1}{u_n^2}=\frac{1}{3}+o(1).$$

D'après le lemme de \textsc{Césaro}, $\frac{1}{n}\sum_{k=0}^{n-1}v_k$ tend également vers $\frac{1}{3}$. Mais,

$$\frac{1}{n}\sum_{k=0}^{n-1}v_k=\frac{1}{n}\sum_{k=0}^{n-1}(\frac{1}{u_{k+1}^2}-\frac{1}{u_k^2})=\frac{1}{n}
(\frac{1}{u_n^2}-\frac{1}{u_0^2}).$$

Ainsi, $\frac{1}{n}(\frac{1}{u_n^2}-\frac{1}{u_0^2})=\frac{1}{3}+o(1)$ puis, $\frac{1}{u_n^2}=\frac{n}{3}+\frac{1}{u_0^2}+o(n)=\frac{n}{3}+o(n)$. Donc, $\frac{1}{u_n^2}\sim\frac{n}{3}$, puis $u_n^2\sim\frac{3}{n}$ et enfin, puisque la suite $u$ est strictement positive,

$$u_n\sim\sqrt{\frac{3}{n}}.$$
\end{enumerate}
\fincorrection
\nocorrection
\nocorrection
\nocorrection
\nocorrection
\nocorrection
\nocorrection
\nocorrection
\nocorrection
\nocorrection
\nocorrection
\nocorrection
\nocorrection
\nocorrection
\correction{001202}
L'équation caractéristique est :
$$r^2-r-1=0$$
dont les solution sont $\lambda = \frac{1-\sqrt5}{2}$ et $\mu =
\frac{1+\sqrt5}{2}$. Donc $u_n$ est de la forme
$$u_n = \alpha \lambda^n + \beta\mu^n$$
pour $\alpha, \beta$ des réels que nous allons calculer gr\^ace à
$u_0$ et $u_1$. En effet $u_0 = 1 = \alpha \lambda^0 + \beta\mu^0$
donc $\alpha+\beta = 1$. Et comme $u_1 =1= \alpha \lambda^1 +
\beta\mu^1$ nous obtenons $\alpha \frac{1-\sqrt5}{2} + \beta
\frac{1+\sqrt5}{2}=1$. En résolvant ces deux équations nous
obtenons $\alpha = \frac{1}{\sqrt5}\frac{\sqrt5-1}{2}=
\frac{1}{\sqrt5}(-\lambda)$ et $\beta =
\frac{1}{\sqrt5}\frac{1+\sqrt5}{2}= \frac{1}{\sqrt5}(\mu)$. Nous
écrivons donc pour finir :
$$u_n = \frac{1}{\sqrt5}\big( \mu^{n+1}-\lambda^{n+1}\big).$$
\fincorrection
\nocorrection
\correction{001204}
L'équation caractéristique est :
$$r^2-3r+2=0$$
dont les solutions sont $\lambda = 2$ et $\mu = 1$. Donc $u_n$ est
de la forme
$$u_n = \alpha 2^n + \beta 1^n = \alpha 2^n+\beta$$
Or la suite $(2^n)_n$ tend vers $+\infty$. Donc si $(u_n)_n$ est
bornée alors $\alpha = 0$. Donc $(u_n)_n$ est la suite constante
égale à $\beta$. Réciproquement toute suite constante qui vérifie
$u_n = \beta$ pour $n\in \Nn$ vérifie bien la relation de
récurrence $u_{n+2} = 3u_{n+1}-2u_n$. Donc les suites cherchées
sont les suites constantes.
\fincorrection
\nocorrection
\nocorrection
\nocorrection
\correction{003068}
\begin{enumerate}
  \item $u_n = \frac 13((a+2b) + 2(a-b)(-\frac12)^n)$.
  \item
  \item  $v_n = \lambda\times\mu^{(-\frac12)^n}$.
  \item 
\end{enumerate}
\fincorrection
\correction{003069}
\begin{enumerate}
  \item $u_n = \frac {n^2}4 - n + \frac 38(1-(-1)^n)$.
  \item $u_n = \frac {n-1}3 + aj^n + bj^{2n}$.
\end{enumerate}
\fincorrection
\correction{003070}
$2u_n = u_0 + v_0 + \left(-\frac 15\right)^n (u_0 - v_0)$,\quad
$2v_n = u_0 + v_0 - \left(-\frac 15\right)^n (u_0 - v_0)$.
\fincorrection
\correction{003071}
\begin{enumerate}
 \item $u^{(k)}_n = \sum_{p=0}^k C_k^p(-1)^{k-p}u_{n+p}$.
  \item 
\end{enumerate}
\fincorrection
\correction{003072}
\begin{enumerate}
  \item 
  \item $6T_n = \bigl(3+\sqrt6\bigr)\bigl(5+2\sqrt6\bigr)^n
                     + \bigl(3-\sqrt6\bigr)\bigl(5-2\sqrt6\bigr)^n$.
\end{enumerate}
\fincorrection
\nocorrection
\correction{005239}
\begin{enumerate}
 \item  L'équation caractéristique est $4z^2-4z-3=0$. Ses solutions sont $-\frac{1}{2}$ et $\frac{3}{2}$. Les suites cherchées sont les suites de la forme $(u_n)=\left(\lambda\left(-\frac{1}{2}\right)^n+\mu\left(\frac{3}{2}\right)^n\right)$ où $\lambda$ et $\mu$ sont deux réels (ou deux complexes si on cherche toutes les suites complexes). Si $u_0$ et $u_1$ sont les deux premiers termes de la suite $u$, $\lambda$ et $\mu$ sont les solutions du système $\left\{
\begin{array}{l}
\lambda+\mu=u_0\\
-\frac{\lambda}{2}+\frac{3\mu}{2}=u_1
\end{array}
\right.$ et donc $\lambda=\frac{1}{4}(3u_0-2u_1)$ et $\mu=\frac{1}{4}(u_0+2u_1)$.

\begin{center}
\shadowbox{
$\forall n\in\Nn,\;u_n=\frac{1}{4}(3u_0-2u_1)\left(-\frac{1}{2}\right)^n+\frac{1}{4}(u_0+2u_1)\left(\frac{3}{2}\right)^n.$
}
\end{center}
 \item  Clairement $u_{2n}=\frac{1}{4^n}u_0$ et $u_{2n+1}=\frac{1}{4^n}u_1$ et donc $u_n=\frac{1}{2}\left(\frac{1}{2^n}(1+(-1)^n)u_0+2\times\frac{1}{2^n}(1-(-1)^n)u_1\right)$.

\begin{center}
\shadowbox{
$\forall n\in\Nn$, $u_n=\frac{1}{2^{n+1}}\left((1+(-1)^n)u_0+2(1-(-1)^n)u_1\right)$.
}
\end{center}
 \item  Les solutions de l'équation homogène associée sont les suites de la forme $\lambda\left(-\frac{1}{2}\right)^n+\mu\left(\frac{3}{2}\right)^n$.
Une solution particulière de l'équation proposée est une constante $a$ telle que $4a=4a+3a+12$ et donc $a=-4$.
Les solutions de l'équation proposée sont donc les suites de la forme $\left(-4+\lambda\left(-\frac{1}{2}\right)^n+\mu\left(\frac{3}{2}\right)^n\right)$ où $\lambda$ et $\mu$ sont les solutions du système $\left\{
\begin{array}{l}
\lambda+\mu=4+u_0\\
-\frac{\lambda}{2}+\frac{3\mu}{2}=4+u_1
\end{array}
\right.$ et donc $\lambda=\frac{1}{4}(4+3u_0-2u_1)$ et $\mu=\frac{1}{4}(12+u_0+2u_1)$.

\begin{center}
\shadowbox{
$\forall n\in\Nn,\;u_n=-4+\frac{1}{4}(4+3u_0-2u_1)\left(-\frac{1}{2}\right)^n+\frac{1}{4}(12+u_0+2u_1)\left(\frac{3}{2}\right)^n.$
}
\end{center}

 \item  La suite $v=\frac{1}{u}$ est solution de la récurrence $2v_{n+2}=v_{n+1}-v_n$ et donc,
$(v_n)$ est de la forme $\left(\lambda\left(\frac{1+i\sqrt{7}}{4}\right)^n+\mu\left(\frac{1-i\sqrt{7}}{4}\right)^n\right)$ et donc $u_n=\frac{1}{\lambda\left(\frac{1+i\sqrt{7}}{4}\right)^n+\mu\left(\frac{1-i\sqrt{7}}{4}\right)^n}$.
 \item  Les solutions de l'équation homogène associée sont les suites de la forme $(\lambda+\mu2^n)$.
$1$ est racine simple de l'équation caractéristique et donc il existe une solution particulière de l'équation proposée de la forme $u_n=an^4+bn^3+cn^2+dn$. Pour $n\geq2$, on a

\begin{align*}
u_n-3u_{n-1}+2u_{n-2}&=(an^4+bn^3+cn^2+dn)-3(a(n-1)^4+b(n-1)^3+c(n-1)^2+d(n-1))\\
 &\;+2(a(n-2)^4+b(n-2)^3+c(n-2)^2+d(n-2))\\
 &= a(n^4-3(n-1)^4+2(n-2)^4)+b(n^3-3(n-1)^3+2(n-2)^3)\\
 &\;+c(n^2-3(n-1)^2+2(n-2)^2)+d(n-3(n-1)+2(n-2))\\
  &= a(-4n^3+30n^2-52n+29)+b(-3n^2+15n-13)+c(-2n+5)+d(-1)\\
 &=n^3(-4a)+n^2(30a-3b)+n(-52a+15b-2c)+29a-13b+5c-d.
\end{align*}

\begin{align*}
u\;\mbox{est solution}&\Leftrightarrow-4a=1\;\mbox{et}\;30a-3b=0\;\mbox{et}\;-52a+15b-2c=0\;\mbox{et}\;29a-13b+5c-d=0\\
 &\Leftrightarrow a=-\frac{1}{4},\;b=-\frac{5}{2},\;c=-\frac{49}{4},\;d=-36.
\end{align*}
Les suites cherchées sont les suites de la forme $\left(-\frac{1}{4}(n^3+10n^2+49n+144)+\lambda+\mu2^n\right)$.
 \item  Pour tout complexe $z$, $z^3-6z^2+11z-6=(z-1)(z-2)(z-3)$ et les suites solutions sont les suites de la forme 
$(\alpha+\beta2^n+\gamma3^n)$.
 \item  Pour tout complexe $z$, $z^4-2z^3+2z^2-2z+1=(z^2+1)^2-2z(z^2+1)=(z-1)^2(z^2+1)$.
Les solutions de l'équation homogène associée sont les suites de la forme $\alpha+\beta n+\gamma i^n+\delta(-i)^n$.
$1$ est racine double de l'équation caractéristique et donc l'équation proposée admet une solution particulière de la forme $u_n=an^7+bn^6+cn^5+dn^4+en^3+fn^2$. Pour tout entier naturel $n$, on a

\begin{align*}
u_{n+4}-2u_{n+3}&+2u_{n+2}-2u_{n+1}+u_n=a((n+4)^7-2(n+3)^7+2(n+2)^7-2(n+1)^7+n^7)\\
 &+b((n+4)^6-2(n+3)^6+2(n+2)^6-2(n+1)^6+n^6)\\
 &+c((n+4)^5-2(n+3)^5+2(n+2)^5-2(n+1)^5+n^5)\\
 &+d((n+4)^4-2(n+3)^4+2(n+2)^4-2(n+1)^4+n^4)\\
 &+e((n+4)^3-2(n+3)^3+2(n+2)^3-2(n+1)^3+n^3)\\
 &+f((n+4)^2-2(n+3)^2+2(n+2)^2-2(n+1)^2+n^2)\\
 &=a(84n^5+840n^4+4340n^3+12600n^2+19348n+12264)\\
 &+b(60n^4+480n^3+1860n^2+3600n+2764)\\
 &+c(40n^3+240n^2+620n+600)+d(24n^2+96n+124)+e(12n+24)+4f\\
 &=n^5(84a)+n^4(840a+60b)+n^3(4340a+480b+40c)+n^2(12600a+1860b+240c+24d)\\
 &+n(19348a+3600b+620c+96d+12e)+(12264a+2764b+600c+124d+24e+4f)
\end{align*}
$u$ est solution si et seulement si $84a=1$ et donc $a=\frac{1}{84}$, puis $840a+60b=0$ et donc $b=-\frac{1}{6}$,
puis $4340a+480b+40c=0$ et donc $c=\frac{17}{24}$, puis $12600a+1860b+240c+24d=0$ et donc $d=-\frac{5}{12}$
puis $19348a+3600b+620c+96d+12e=0$ et donc $e=-\frac{59}{24}$ puis $12264a+2764b+600c+124d+24e+4f=0$ et donc $f=\frac{1}{12}$.
La solution générale de l'équation avec second membre est donc~:

$$\forall n\in\Nn,\;u_n=\frac{1}{168}(2n^7-28n^6+119n^5-70n^4-413n^3+14n^2)+\alpha+\beta n+\gamma i^n+\delta(-i)^n,\; (\alpha,\beta,\gamma,\delta)\in\Cc^4.$$
\end{enumerate}
\fincorrection
\nocorrection
\nocorrection
\nocorrection
\nocorrection
\nocorrection
\nocorrection
\nocorrection
\nocorrection
\nocorrection
\nocorrection
\correction{005145}
\begin{enumerate}
\item  Soit $n\in\Nn^*$.

\begin{align*}
\prod_{k=1}^{n}(1+\frac{1}{k})&=\prod_{k=1}^{n}\frac{k+1}{k}=\frac{\prod_{k=1}^{n}(k+1)}{\prod_{k=1}^{n}k}=\frac{(n+1)!}
{n!}=n+1
\end{align*}

\item  Soit $a\in]0,\pi[$ et $n\in\Nn^*$. Alors, pour tout naturel non nul $k$, on a
$0<\frac{a}{2^k}\leq\frac{a}{2}<\frac{\pi}{2}$ et donc $\sin\frac{a}{2^k}\neq0$.

On sait alors que pour tout réel $x$, $\sin(2x)=2\sin x\cos x$. Par suite, pour tout naturel $k$,

$$\sin(2.\frac{a}{2^k})=2\sin\frac{2^k}\cos\frac{a}{2^k}\quad\mbox{et donc}\quad
\cos\frac{a}{2^k}=\frac{\sin(a/2^{k-1})}{2\sin(a/2^k)}.$$

Mais alors,

\begin{align*}
\prod_{k=1}^{n}\cos\frac{a}{2^k}&=\prod_{k=1}^{n}\frac{\sin(a/2^{k-1})}{2\sin(a/2^k)}=\frac{1}{2^n}\frac{\prod_{k=1}^{n
}\sin(a/2^{k-1})}{\prod_{k=1}^{n}\sin(a/2^{k})}
 =\frac{1}{2^n}\frac{\prod_{k=0}^{n-1}
\sin(a/2^{k})}{\prod_{k=1}^{n}\sin(a/2^{k})}
=\frac{\sin a}{2^n\sin(a/2^n)}.
\end{align*}
\end{enumerate}
\fincorrection
\correction{005148}
Pour $n\in\Nn^*$, $(1+\frac{1}{n})^n=\sum_{k=0}^{n}\frac{C_n^k}{n^k}$. Pour $k\in\{0,...,n\}$,
posons $u_k=\frac{C_n^k}{n^k}$ puis $v_k=\frac{u_{k+1}}{u_k}$.
Pour $k\in\{1,...,n-1\}$, on a alors

\begin{align*}
v_k&=\frac{C_n^{k+1}.n^k}{C_n^k.n^{k+1}}=\frac{1}{n}.\frac{n!k!(n-k)!}{n!(k+1)!(n-k-1)!}=\frac{n-k}{n(k+1)}
=\frac{(n+1)-(k+1)}{n(k+1)}=-\frac{1}{n}+\frac{n+1}{n(k+1)}\\
 &\leq-\frac{1}{n}+\frac{n+1}{2n}\;(\mbox{car}\;k\geq1)\\
 &=\frac{1}{2}-\frac{1}{2n}<\frac{1}{2}
\end{align*}

Ainsi, pour $k\in\{1,...,n-1\}$, $u_{k+1}\leq\frac{1}{2}u_k$ et donc, immédiatement par récurrence,

$$u_k\leq\frac{1}{2^{k-1}}u_1=\frac{1}{2^{k-1}}\frac{n}{n}=\frac{1}{2^{k-1}}.$$

En tenant compte de $u_0=1$, on a alors pour $n\in\Nn^*$,

$$(1+\frac{1}{n})^n=\sum_{k=0}^{n}u_k\leq1+\sum_{k=1}^{n}\frac{1}{2^{k-1}}=1+\frac{1-\frac{1}{2^n}}{1-\frac{1}{2}}
=1+2(1-\frac{1}{2^n})=3-\frac{1}{2^{n-1}}<3.$$
\fincorrection
\correction{005154}
Soient $x\in\Rr$ et $n\in\Nn^*$. Pour $1\leq k\leq n$, on a

$$kx-1<E(kx)\leq kx.$$

En sommant ces inégalités, on obtient

$$\frac{E(x)+E(2x)+...+E(nx)}{n^2}\leq\frac{x+2x+...+nx}{n^2}=\frac{n(n+1)x}{2n^2}=\frac{(n+1)x}{2n},$$

et aussi,

$$\frac{E(x)+E(2x)+...+E(nx)}{n^2}>\frac{(x-1)+(2x-1)+...+(nx-1)}{n^2}=\frac{n(n+1)x/2-n}{n^2}
=\frac{(n+1)x}{2n}-\frac{1}{n}.$$

Finalement, pour tout naturel non nul,

$$\frac{(n+1)x}{2n}-\frac{1}{n}<\frac{E(x)+E(2x)+...+E(nx)}{n^2}\leq\frac{(n+1)x}{2n}.$$

Les deux membres extrêmes de cet encadrement tendent vers $\frac{x}{2}$ quand $n$ tend vers $+\infty$. D'après le
théorème des gendarmes, on peut affirmer que
\begin{center}
\shadowbox{
$\forall x\in\Rr,\;\lim_{n\rightarrow +\infty}\frac{E(x)+E(2x)+...+E(nx)}{n^2}=\frac{x}{2}.$
}
\end{center}
\fincorrection
\correction{005220}
\begin{enumerate}
 \item  Soit $\varepsilon>0$. Il existe un rang $n_0$ tel que, si $n\geq n_0$ alors $|u_n-\ell|<\frac{\varepsilon}{2}$. Soit $n$ un entier naturel strictement supérieur à $n_0$.

\begin{align*}
|v_n-\ell|&=\left|\frac{1}{n+1}\sum_{k=0}^{n}u_k-\ell\right|=\left|\frac{1}{n+1}\sum_{k=0}^{n}(u_k-\ell)\right|\\
 &\leq\frac{1}{n+1}\sum_{k=0}^{n}|u_k-\ell|=\frac{1}{n+1}\sum_{k=0}^{n_0}|u_k-\ell|+\frac{1}{n+1}\sum_{k=n_0+1}^{n}|u_k-\ell|\\
 &\leq\frac{1}{n+1}\sum_{k=0}^{n_0}|u_k-\ell|+\frac{1}{n+1}\sum_{k=n_0+1}^{n}\frac{\varepsilon}{2}\leq\frac{1}{n+1}\sum_{k=0}^{n_0}|u_k-\ell|+\frac{1}{n+1}\sum_{k=0}^{n}\frac{\varepsilon}{2}\\
 &=\frac{1}{n+1}\sum_{k=0}^{n_0}|u_k-\ell|+\frac{\varepsilon}{2}
\end{align*}
Maintenant, $\sum_{k=0}^{n_0}|u_k-\ell|$ est une expression constante quand $n$ varie et donc, $\lim_{n\rightarrow +\infty}\frac{1}{n+1}\sum_{k=0}^{n_0}|u_k-\ell|=0$. Par suite, il existe un entier $n_1\geq n_0$ tel que pour $n\geq n_1$, $\frac{1}{n+1}\sum_{k=0}^{n_0}|u_k-\ell|<\frac{\varepsilon}{2}$.
Pour $n\geq n_1$, on a alors $|v_n-\ell|<\frac{\varepsilon}{2}+\frac{\varepsilon}{2}=\varepsilon$.
On a montré que $\forall\varepsilon>0,\;\exists n_1\in\Nn/\;(\forall n\in\Nn)(n\geq n_1\Rightarrow|v_n-\ell|<\varepsilon)$. La suite $(v_n)$ est donc convergente et $\lim_{n\rightarrow +\infty}v_n=\ell$.

\begin{center}
\shadowbox{
Si la suite $u$ converge vers $\ell$ alors la suite $v$ converge vers $\ell$.
}
\end{center}
La réciproque est fausse. Pour $n$ dans $\Nn$, posons $u_n=(-1)^n$. La suite $(u_n)$ est divergente. D'autre part, pour $n$ dans $\Nn$, $\sum_{k=0}^{n}(-1)^k$ vaut $0$ ou $1$ suivant la parité de $n$ et donc, dans tous les cas, $|v_n|\leq\frac{1}{n+1}$. Par suite, la suite $(v_n)$ converge et $\lim_{n\rightarrow +\infty}v_n=0$.
 \item  Si $u$ est bornée, il existe un réel $M$ tel que, pour tout naturel $n$, $|u_n|\leq M$.
Pour $n$ entier naturel donné, on a alors

$$|v_n|\leq\frac{1}{n+1}\sum_{k=0}^{n}|u_k|\leq\frac{1}{n+1}\sum_{k=0}^{n}M=\frac{1}{n+1}(n+1)M=M.$$
La suite $v$ est donc bornée.

\begin{center}
\shadowbox{
Si la suite $u$ est bornée alors la suite $v$ est bornée.
}
\end{center}
La réciproque est fausse. Soit $u$ la suite définie par~:~$\forall n\in\Nn,\;u_n=(-1)^nE\left(\frac{n}{2}\right)=
\left\{
\begin{array}{l}
p\;\mbox{si}\;n=2p,\;p\in\Nn\\
-p\;\mbox{si}\;n=2p+1,\;p\in\Nn
\end{array}
\right.$.
$u$ n'est pas bornée car la suite extraite $(u_{2p})$ tend vers $+\infty$ quand $p$ tend vers $+\infty$. 
Mais, si $n$ est impair, $v_n=0$, et si $n$ est pair, $v_n=\frac{1}{n+1}\times u_n=\frac{n}{2(n+1)}$, et dans tous les cas $|v_n|\leq\frac{1}{n+1}\frac{n}{2}\leq\frac{1}{n+1}\frac{n+1}{2}=\frac{1}{2}$ et la suite $v$ est bornée.
 \item  Si $u$ est croissante, pour $n$ entier naturel donné on a~:

\begin{align*}
v_{n+1}-v_n&=\frac{1}{n+2}\sum_{k=0}^{n+1}u_k-\frac{1}{n+1}\sum_{k=0}^{n}u_k=\frac{1}{(n+1)(n+2)}\left((n+1)\sum_{k=0}^{n+1}u_k-(n+2)\sum_{k=0}^{n}u_k\right)\\
 &=\frac{1}{(n+1)(n+2)}\left((n+1)u_{n+1}-\sum_{k=0}^{n}u_k\right)=\frac{1}{(n+1)(n+2)}\sum_{k=0}^{n}(u_{n+1}-u_k)\geq0.
\end{align*}
La suite $v$ est donc croissante.

\begin{center}
\shadowbox{
Si la suite $u$ est croissante alors la suite $v$ est croissante.
}
\end{center}
\end{enumerate}
\fincorrection
\correction{005246}
Pour $n$ naturel non nul et $x$ réel positif, posons $f_n(x)=x^n+x-1$.

Pour $x\geq0$, $f_1(x)=0\Leftrightarrow x=\frac{1}{2}$ et donc $u_1=\frac{1}{2}$.

Pour $n\geq2$, $f_n$ est dérivable sur $\Rr^+$ et pour $x\geq 0$, $f_n'(x)=nx^{n-1}+1>0$.

$f_n$ est ainsi continue et strictemnt croissante sur $\Rr^+$ et donc bijective de $\Rr^+$ sur $f_n(\Rr^+)=[f(0),\lim_{x\rightarrow +\infty}f_n(x)[=[-1,+\infty[$, et en particulier,

$$\exists!x\in[0,+\infty[/\;f_n(x)=0.$$

Soit $u_n$ ce nombre. Puisque $f_n(0)=-1<0$ et que $f_n(1)=1>0$, par stricte croissance de $f_n$ sur $[0,+\infty[$, on a~:

$$\forall n\in\Nn,\;0<u_n<1.$$
 
La suite $u$ est donc bornée.

Ensuite, pour $n$ entier naturel donné et puisque $0<u_n<1$~:

$$f_{n+1}(u_n)=u_n^{n+1}+u_n-1<u_n^n+u_n-1=f_n(u_n)=0=f_{n+1}(u_{n+1}),$$

et donc $f_{n+1}(u_n)<f_{n+1}(u_{n+1})$ puis, par stricte croissance de $f_{n+1}$ sur $\Rr^+$, on obtient~:

$$\forall n\in\Nn,\;u_n<u_{n+1}.$$

La suite $u$ est bornée et strictement croissante. Donc, la suite $u$ converge vers un réel $\ell$, élément de $[0,1]$.

Si $0\leq\ell<1$, il existe un rang $n_0$ tel que pour $n\geq n_0$, on a~:~$u_n\leq\ell+\frac{1-\ell}{2}=\frac{1+\ell}{2}$. Mais alors, pour $n\geq n_0$, on a $1-u_n=u_n^n\leq(\frac{1+\ell}{2})^n$ et quand $n$ tend vers vers $+\infty$, on obtient $1-\ell\leq0$ ce qui est en contradiction avec $0\leq\ell<1$. Donc, $\ell=1$.

\fincorrection
\correction{005248}
Soit $x$ dans $[-1,1]$ et $\varepsilon>0$.

Soit $\theta=\Arcsin x$ (donc $\theta$ est élément de $[-\frac{\pi}{2},\frac{\pi}{2}]$ et $x=\sin\theta$). Pour $k$ entier naturel non nul donné, il existe un entier $n_k$ tel que $\ln(n_k)\leq \theta+2k\pi<\ln(n_k+1)$ à savoir $n_k=E(e^{\theta+2k\pi})$.

Mais,

$$0<\ln(n_k+1)-\ln(n_k)=\ln(1+\frac{1}{n_k})<\frac{1}{n_k}$$

(d'après l'inégalité classique $\ln(1+x)<x$ pour $x>0$, obtenue par exemple par l'étude de la fonction $f~:~x\mapsto\ln(1+x)-x$).
Donc,

$$0\leq\theta+2k\pi-\ln(n_k)<\ln(n_k+1)-\ln(n_k)<\frac{1}{n_k},$$

puis 

\begin{align*}\ensuremath
|\sin(\theta)-\sin(\ln(n_k))|&=2|\sin(\frac{\theta+2k\pi-\ln(n_k)}{2})\cos(\frac{\theta+2k\pi+\ln(n_k)}{2})|\\
 &\leq2\left|\frac{\theta+2k\pi-\ln(n_k)}{2}\right|=|\theta+2k\pi-\ln(n_k)|<\frac{1}{n_k}.
\end{align*}

Soit alors $\varepsilon$ un réel strictement positif.

Puisque $n_k=E(e^{\theta+2k\pi})$ tend vers $+\infty$ quand $k$ tend vers $+\infty$, on peut trouver un entier $k$ tel que  $\frac{1}{n_k}<\varepsilon$ et pour cet entier $k$, on a $|\sin\theta-\sin(\ln(n_k))|<\varepsilon$.

On a montré que $\forall x\in[-1,1],\;\forall\varepsilon>0,\;\exists n\in\Nn^*/\;|x-\sin(\ln n)|<\varepsilon$, et donc $\{\sin(\ln n),\;n\in\Nn^*\}$ est dense dans $[-1,1]$.
\fincorrection
\correction{005249}
Pour $\alpha\in]0,\pi[$, posons $f(\alpha)=\mbox{sup}_{n\in\Nn}(|\sin(n\alpha)|)$. $\{(\sin(n\alpha),\;n\in\Nn\}$ est une partie non vide et majorée (par $1$) de $\Rr$. Donc, pour tout réel $\alpha$ de $]0,\pi[$, $f(\alpha)$ existe dans $\Rr$.

Si $\alpha$ est dans $[\frac{\pi}{3},\frac{2\pi}{3}]$,

$$f(\alpha)=\mbox{sup}_{n\in\Nn}(|\sin(n\alpha)|)\geq\sin\alpha\geq\frac{\sqrt{3}}{2}=f(\frac{\pi}{3}).$$

Si $\alpha$ est dans $]0,\frac{\pi}{3}]$. Soit $n_0$ l'entier naturel tel que $(n_0-1)\alpha<\frac{\pi}{3}\leq n_0\alpha$ ($n_0$ existe car la suite $(n\alpha)_{n\in\Nn}$ est strictement croissante). Alors, 
 
$$\frac{\pi}{3}\leq n_0\alpha=(n_0-1)\alpha+\alpha<\frac{\pi}{3}+\alpha\leq\frac{\pi}{3}+\frac{\pi}{3}=\frac{2\pi}{3}.$$

Mais alors,

$$f(\alpha)=\mbox{sup}_{n\in\Nn}(|\sin(n\alpha)|)\geq|\sin(n_0\alpha)|\geq\frac{\sqrt{3}}{2}=f(\frac{\pi}{3}).$$

Si $\alpha$ est dans $[\frac{2\pi}{3},\pi[$, on note que 
$$f(\alpha)=\mbox{sup}_{n\in\Nn}(|\sin(n\alpha)|)=\mbox{sup}_{n\in\Nn}(|\sin(n(\pi-\alpha)|)=f(\pi-\alpha)\geq f(\frac{\pi}{3}),$$

car $\pi-\alpha$ est dans $]0,\frac{\pi}{3}]$.

On a montré que $\forall\alpha\in]0,\pi[,\;f(\alpha)\geq f(\frac{\pi}{3})=\frac{\sqrt{3}}{2}$. Donc, $\mbox{inf}_{\alpha\in]0,\pi[}(\mbox{sup}_{n\in\Nn}(|\sin(n\alpha)|))$ existe dans $\Rr$ et 

$$\mbox{inf}_{\alpha\in]0,\pi[}(\mbox{sup}_{n\in\Nn}(|\sin(n\alpha)|))= \mbox{Min}_{\alpha\in]0,\pi[}(\mbox{sup}_{n\in\Nn}(|\sin(n\alpha)|))=f(\frac{\pi}{3})=\frac{\sqrt{3}}{2}.$$
\fincorrection
\nocorrection
\nocorrection
\nocorrection
\nocorrection
\correction{001936} 

\smallskip
\begin{enumerate}

\item Pour $n \ge 4$
\[
\frac{n!}{n^n} = \frac{2}{n} \times \frac{3}{n} \times 
\underbrace{\frac{4}{n} \times \cdots\times\frac{n}{n}}_{\le 1} 
\le \frac{6}{n^2} .
\]
Or $\sum \frac{6}{n^2}$ est convergente.	
Donc $\sum \frac{n!}{n^n}$ est aussi convergente par comparaison.

\item
Montrons que $(\ch \sqrt{\ln n })^{-2} \ge \left(\sqrt n+\frac{1}{\sqrt{n}}\right)^{-2}$ 
pour $n$ assez grand.	On a :
\[ \begin{array}{rcl}
4\ln n  & \le & \ln^{2}{n} \quad \text{pour $n$ assez grand}
\\
\ln n & \le & \left(\frac{1}{2}\ln n \right)^{2}
\\
\sqrt{\ln n} & \le & \frac{1}{2}\ln n = \ln{\sqrt{n}} 
\\
\mathrm{ch}(\sqrt{\ln n }) & \le & \mathrm{ch}(\ln{\sqrt{n}}) 
= \frac{1}{2} \left( \sqrt{n} + \frac{1}{\sqrt{n}} \right)
\quad \text{car $ x \mapsto \mathrm{ch} x$ est croissante}
\\
\mathrm{ch}(\sqrt{\ln n })^{-2} & \ge & 4 \left(\sqrt{n} + \frac{1}{\sqrt{n}} \right)^{-2}
\end{array} \]
Or $\sqrt{n} \sim \sqrt{n} + \frac{1}{\sqrt{n} }$, et
$\left(\sqrt{n} + \frac{1}{\sqrt{n}} \right)^{-2} \sim \frac{1}{n}$, 
donc $\sum \left(\sqrt{n} + \frac{1}{\sqrt{n}} \right)^{2}$ est divergente.
Par comparaison, la série de terme  général $(\mathrm{ch}\sqrt{\ln n })^{2}$ est divergente.

\item
Montrons que $n^{-\left(1+\frac{1}{n}\right) } \sim n^{-1}$. On a : 
$\displaystyle 
\frac{ n^{-\left(1+\frac{1}{n}\right) } }{ n^{-1} } = n^{-\frac{1}{n}} = e^{-\frac{\ln n }{n}}$.
Or $\lim_{n \to + \infty} \frac{\ln n }{n} = 0$, d'où $\lim_{n \to + \infty} e^{-\frac{\ln n }{n}} = 1$.
Par équivalence, la série de terme  général $n^{-\left(1+\frac{1}{n}\right) } $ est donc divergente 
car la série harmonique est divergente.

\item
Montrons que $\frac{1}{\sqrt{n}}\ln{\left(1+\frac{1}{\sqrt{n}}\right)} \sim \frac{1}{n}$.
En utilisant le développement limité de $\ln (1 +x)$ en $0$, on a :
$\ln \left(1+\frac{1}{\sqrt{n}}\right) = \frac{1}{\sqrt{n}} + o\left(\frac{1}{\sqrt{n}}\right)$.
De là on tire que $\frac{1}{\sqrt{n}}\ln{\left(1+\frac{1}{\sqrt{n}}\right)} 
\sim \frac{1}{\sqrt{n}} \times \frac{1}{\sqrt{n}} = \frac{1}{n}$.
Par équivalence, la série de terme général $n^{-\left(1+\frac{1}{n}\right) }$ est donc divergente.

\item 
On sait que :
$\ln{(e^n-1)} \le \ln{e^n} = n$. De plus, $ \ln n \ge 1$ pour $n$  assez grand, par conséquent
$\frac{\ln n}{\ln (e^n-1)} \ge \frac{1}{n}$. On conclut par comparaison 
que la série $\sum \frac{\ln n }{\ln{(e^n-1)}}$ est divergente.

\item
Montrons que $n^{\ln n }e^{-\sqrt{n}} \le n^{-2}$.
On remarque que $ n^{\ln n }e^{-\sqrt{n}} = e^{(\ln n)^2}e^{-\sqrt{n}}$. 
Or pour $u$ assez grand $4u^2 + 4u \le e^u $, soit $4u^2 - e^u \le - 4u$.
En posant $u = \ln \sqrt{n} = \frac{1}{2} \ln n$, il vient $\ln^2{n} - \sqrt{n} \le -2\ln n$.
D'où 
\[
\underbrace{e^{\ln^2{n}-\sqrt{n}}}_{n^{\ln n }e^{-\sqrt{n}}} \le \underbrace{e^{-2\ln n }}_{\frac{1}{n^2}}
\]

Par comparaison, la série de terme général $n^{\ln n }e^{-\sqrt{n} }$ est donc convergente 
car la série de terme général $\frac{1}{n^2}$ est convergente.
\end{enumerate}
\medskip

(\emph{Corrigé de Lévi Operman})

\fincorrection
\correction{001937} 

\begin{itemize}
\item Pour $p = 0$:
\[ 
u_n = \frac{ 1! + 2! + \cdots + n! }{ n! } = 1 + \frac{ 1! + 2! + \cdots + (n-1)! }{ n! } > 1 
\]
$u_n$ ne tend pas vers $ 0 $ donc, $ \sum u_n $ diverge grossièrement pour $ p=0 $. 

\item Pour $p = 1$:
\[ 
u_n = \frac{1}{ (n+1)! } + \frac{2!}{ (n+1)! } + \cdots + \frac{ (n-1)! }{ (n+1)! } + \frac{n!}{ (n+1)! } 
\]
\[ 
u_n \geq \frac{n!}{ (n+1)! } = \frac{1}{ n + 1 }  \cdotp
\]
Or $\sum \frac{1}{n + 1}$ diverge, donc $\sum u_n$ diverge pour $p = 1$ .

\item Pour $p = 2$:
\[ u_n = \frac{1}{ (n+2)! } + \frac{2!}{ (n+2)! } + \cdots + \frac{ (n-1)! }{ (n+2)! } + \frac{n!}{ (n+2)! } 
\]
On serait tenté de dire que l'on a une somme de séries convergentes, donc $\sum u_n$ converge. 
Pas de chance, le nombre de terme croît en fonction de $n$, donc à l'infini, on en a une infinité 
et on ne peut rien conclure.
\[ 
u_n = \sum_{k=1}^{n} \frac{k!}{ (n + 2)!} = \sum_{k=1}^{n-1} \frac{k!}{(n + 2)!} + \frac{n!}{(n + 2)!}  
   \leq \frac{n (n - 1)!}{(n + 2)!} + \frac{n!}{(n + 2)!} \]
\[ u_n \leq 2\frac{n!}{(n + 2)!} = \frac{2}{(n + 1)(n + 2)} \sim \frac{2}{n^2} \]
Or $\sum \frac{1}{n^2}$ converge, donc $\sum u_n$ converge pour $p = 2$. 

\item Pour $p \geq 3$:
\[ 
u_n = \frac{ 1! + 2! + \cdots + n! }{( n + p)! } \leq \frac{ n \, n! }{ (n + p)! } = \frac{ n \, n! }{n! (n+1) \cdots (n+p) } 
\]
En simplifiant par $ n! $ et en posant $ u_n \leq \frac{ n }{ (n+1) \cdots (n+p)}$ et
\[ 
\frac{ n }{ (n+1) \cdots (n+p)} \thicksim  \frac{n}{n^p} 
= \frac{1}{n^{p - 1}} \text{ avec } p \geq 3
\] 
Or $\sum \frac{1}{n^{p - 1}}$ est une série de Riemann convergente car $p-1 \geq 2$, 
donc $\sum u_n $ converge pour $p \geq 3$. 

Note: on peut aussi remarquer que $u_n$ (quand $p\geq3$) est majoré par $u_n$ (quand $p=2$), or ce dernier est convergent.
\end{itemize}
\medskip

(\emph{Corrigé de Eugène Ndiaye})

\fincorrection
\correction{001938}
\begin{enumerate}

\item Posons $S_n = \sum_{k=0}^{n} (k+1)3^{-k}$.
L'idée est de calculer la somme de $(1-3^{-1})S_n$.
On a ainsi:
\begin{eqnarray*}
(1-3^{-1})S_n & = & (1-3^{-1})\sum_{k=0}^{n} (k+1)3^{-k} \\
 & = & \sum_{k=0}^{n} (k+1)3^{-k}-\sum_{k=0}^{n} (k+1)3^{-(k+1)} \\
 & = & \sum_{k=0}^{n} k 3^{-k} + \sum_{k=0}^{n} 3^{-k} - \sum_{k=0}^{n} (k+1)3^{-(k+1)}
{}\end{eqnarray*}
En réindexant les sommes, on obtient: 
\begin{eqnarray*}
(1-3^{-1})S_n & = & \sum_{k=1}^{n} k3^{-k} + \sum_{k=0}^{n} \left(\frac{1}{3}\right)^{k} - \sum_{k=1}^{n} k3^{-k} - (n+1)3^{-(n+1)}\\
 & = & \sum_{k=0}^{n} \frac{1}{3^k} - \frac{n+1}{3^{n+1}} 
 = \cfrac{1 - \left(\frac{1}{3}\right)^{n+1}}{\frac{2}{3}} 
 - \underbrace{\frac{n+1}{3^{n+1}}}_{\to 0}
\end{eqnarray*}
somme des termes d'une suite géométrique de raison $\frac{1}{3} \in ]-1,1[$. Et donc, 
\[ \lim_{n \to +\infty} \left(1-\frac{1}{3}\right)S_n = \cfrac{3}{2} \]
d'où 
\[ 
\sum_{k=0}^{+\infty}(k+1)3^{-k} = \frac{9}{4} \, \cdotp
\]
Remarque: On reconnaît la série géométrique dérivée première de raison $\frac{1}{3}$.

\item Posons  $\displaystyle u_n = \frac{n}{n^4+n^2+1} $ et cherchons à la décomposer en éléments simples.
\begin{eqnarray*}
n^4 + n^2 + 1 &=& (n^4 + 2n^2 + 1) - n^2 = (n^2 + 1)^2 - n^2 
\\
&=& (n^2 + n + 1)(n^2 - n + 1)
{}\end{eqnarray*}
d'où $ u_n = \frac{n}{ (n^2+n+1)(n^2-n+1) } $~. 
Trouvons maintenant $A$ et $B \in \mathbb{R}$ tel que $u_n = \frac{A}{n^2+n+1} + \frac{B}{n^2-n+1}$,
soit tels que $A(n^2-n+1) + B(n^2+n+1) = n$, ce qui équivaut à $(A+B)n^2 + (B-A)n + (A+B) = n$.
Par identification, on a: 
\[
\left\{
\begin{array}{r c l}
A+B = 0\\
B-A = 1\\
A+B = 0
\end{array}
\right.
\iff \left\{
\begin{array}{r c l}
A & = & -\frac{1}{2}\\
& & \\
B & = & \frac{1}{2}
\end{array}
\right.
\]
D'où:  
\[ 
u_n = \frac{n}{n^4+n^2+1} = \frac{1}{2} \left( \frac{1}{n^2-n+1} - \frac{1}{n^2+n+1} \right  ) 
\]
\[  
\sum_{n=0}^{N} u_n = \sum_{n = 0}^{N} \frac{1}{2} \left( \frac{1}{n^2-n+1} - \frac{1}{n^2+n+1} \right) = \frac{1}{2} \left(\sum_{n = 0}^{N} \frac{1}{n^2-n+1} - \sum_{n = 0}^{N} \frac{1}{n^2+n+1} \right) 
\] 
Or $ n^2+n+1 = (n+1)^2 - (n+1) + 1 $. En réindexant la deuxième somme:
\begin{eqnarray*}
\sum_{n=0}^{N} u_n & = & \frac{1}{2} \left(\sum_{n = 0}^{N} \frac{1}{n^2-n+1} - \sum_{n = 1}^{N+1} \frac{1}{n^2-n+1} \right) \\
& = & \frac{1}{2} \left( 1 - \frac{1}{(N+1)^2 - (N+1) + 1} \right) \quad \text{par télescopage.}
{}\end{eqnarray*}
La série est donc convergente et la somme vaut
\[ 
\sum_{n = 0}^{+\infty} \frac{n}{n^4+n^2+1} = \frac{1}{2}  \,.
\]

\item Décomposons $\displaystyle v_n=\frac{2n-1}{n^3-4n}$ en éléments simples.
Comme 
\[
n^3-4n = n(n^2-4) = n (n-2) (n + 2),
\] 
cherchons $\alpha$, $\beta$ et $\gamma \in \mathbb{R}$ tels que:
\[ 
v_n = \frac{\alpha}{n} + \frac{\beta}{n+2} + \frac{\gamma}{n-2} .
\]
Soit
\begin{eqnarray*}
 2n-1 & = & \alpha(n-2)(n+2) + \beta n(n-2) + \gamma n(n+2)\\
 & = & (\alpha+\beta+\gamma) n^2 + (2\gamma-2\beta)n - 4\alpha
{}\end{eqnarray*}
Par identification:
\[
\left\{
\begin{array}{r c l}
\alpha + \beta + \gamma = 0\\
2(\gamma - \beta) = 2\\
-4 \alpha = - 1
\end{array}
\right.
\Longleftrightarrow \left\{
\begin{array}{r c l}
\alpha & = & \frac{1}{4}\\
\gamma & = & 1 + \beta \\
\alpha + 2\beta + 1 & = & 0
\end{array}
\right.
\Longleftrightarrow \left\{
\begin{array}{r c l}
\alpha & = & \frac{1}{4}\\
\beta & = & - \frac{5}{8}\\
\gamma & = & \frac{3}{8}
\end{array}
\right.
\]
D'où 
\[
\frac{2n-1}{n^3-4n} = \frac{1}{4n} - \frac{5}{8(n+2)} + \frac{3}{8(n-2)} 
\]
et
\begin{eqnarray*}
\sum_{n = 3}^{N} v_n & = & \sum_{n = 3}^{N} \left(\frac{1}{4n} - \frac{5}{8(n+2)} + \frac{3}{8(n-2)}\right) \\
& = & \sum_{n = 3}^{N}\frac{1}{4n} - \sum_{n = 3}^{N} \frac{5}{8(n+2)} + \sum_{n = 3}^{N} \frac{3}{8(n-2)} \,.
{}\end{eqnarray*}
Soit en réindexant les 2 dernières sommes:
\[ 
\sum_{n = 3}^{N} v_n = \frac{1}{8} \left[2 \sum_{n=3}^{N} \frac{1}{n} - 5 \sum_{n=5}^{N+2} \frac{1}{n} + 3 \sum_{n=1}^{N-2} \frac{1}{n} \right] 
\]
Puis,  par télescopage,
\[ 
\sum_{n = 3}^{N} v_n = \frac{1}{8} \left(\frac{89}{12}  - \frac{3}{N-1} - \frac{3}{N} - \frac{5}{N+1} - \frac{5}{N+2}\right) \xrightarrow[n \to +\infty]{} \frac{89}{12}
\, . 
\]
Donc, la série converge et $\displaystyle \sum_{n = 3}^{+\infty} \frac{2n-1}{n^3-4n} = \frac{89}{96}$.

\end{enumerate}

\medskip

(\emph{Corrigé d'Antoine Poulain})

\fincorrection
\nocorrection
\nocorrection
\nocorrection
\nocorrection
\nocorrection
\correction{001949}
Convergence de $W_n = \ln (u_n n^{b-a})$.

On remarque que $W_n$ est la somme partielle de la suite de terme général 
\begin{eqnarray*}
w_n &=& W_{n+1} - W_n = \ln \left[ \frac{ u_{n+1} }{ u_n } \left(\frac{n+1}{n}\right)^{b-a} \right]
\\
&=& \ln \left[ \frac{ n + a }{ n + b  } \left(\frac{n+1}{n}\right)^{b-a} \right] 
= \ln \left[ \frac{ n( 1+\frac{a}{n}) }{ n( 1+\frac{b}{n}) } \left(1+\frac{1}{n}\right)^{b-a} \right] 
\\
&=& (b-a) \ln \left( 1 + \frac{1}{n} \right) 
+ \ln \left( 1 + \frac{a}{n} \right) - \ln \left( 1 + \frac{b}{n} \right)
{}\end{eqnarray*}
Il suffit donc de montrer que cette série converge pour montrer que $(W_n)$ converge. 
On utilise le développement limité de $\ln{( 1 + x )} $ en 0 , ce qui donne 
\[
w_n = (b-a)\left( \frac{1}{n} + O\left(\frac{1}{n^2}\right) \right) + \left(\frac{a}{n} + O\left(\frac{1}{n^2}\right)\right)  - \left(\frac{b}{n} + O\left(\frac{1}{n^2}\right)\right)
= O\left(\frac{1}{n^2}\right) .
\]
donc $\sum{w_n}$ est une série convergente et $(W_n)$ converge. Soit $\ell$ sa limite.

Condition sur $a,b$ pour que $\sum u_n$ converge.

On sait que $\lim \ln{u_n n^{b-a}} = \ell$ ; par composée des limites,  $\lim u_n n^{b-a}  = e^\ell$, 
donc $u_n \sim \frac{e^\ell}{n^{b-a}}$. Or $\sum \frac{1}{n^{b-a}} $ est une série de Riemann, 
qui converge si et seulement si $b-a  > 1$.
Ainsi, par équivalence, $\sum^{\infty}u_n $ converge si et seulement si $b-a  > 1$.

Calcul somme partielle de $s_n$.

Par hypothèse $\frac{u_{n+1}}{u_n} = \frac{n+a}{n+b}$, d'où $[u_{n+1}( n+b )] = [u_n( n+a )]$ et
\[
\sum_{j=0}^n[u_{j+1}( ( j+1 )+( b-1 ) )] = \sum_{j=0}^n[u_j( j+a ) ] \, .
\]
En effectuant un changement d'indice on a :
\[
\sum_{j=1}^{n+1}[u_{j}( j+b-1 )] = \sum_{j=0}^n[u_j( j+a ) ]
\] 
\[
\sum_{j=1}^{n}[u_{j}( j+b-1 )] + u_{n+1}( n+b ) = \sum_{j=1}^n[u_j( j+a ) ] + au_0
\]
\[
u_{n+1}(n+b)-au_0=\sum_{j=1}^n[u_j(a-b+1)]
\]
Si $b-a \ne 1$, on obtient donc que $s_n = \frac{u_{n+1}(n+b)-a}{a-b+1}$.


Valeur de la somme.

 
On se place dans le cas où la série converge, i.e. $b-a>1$. Alors $\lim u_{n+1}(n+b) = 0$.
On sait que $u_n\sim\frac{e^\ell}{n^{b-a}}$, de plus, $n+b\sim n$.
Donc $u_{n+1}(n+b)\sim e^ln^{1+a-b}$. Or $1 + a - b < 0$, donc $\lim e^\ell n^{1+a-b} = 0$.
Finalement
\[
s_n = \frac{u_{n+1}(n+b)-a}{a-b+1} \xrightarrow[n \to +\infty]{}
\frac{-a}{a-b+1}
\]
et on conclut que $\sum_{k=0}^\infty{u_n} = \frac{a}{b-a-1}$.

\medskip

(\emph{Corrigé de Lévi Operman})


\fincorrection
\nocorrection
\nocorrection
\nocorrection
\nocorrection
\correction{002722}

Soit $\alpha, \beta \in \mathbb{R}$ et 
$I_{(\alpha,\beta)} = \int_{0}^{\pi} (\alpha t + \beta t^2) \cos(nt) \, dt$.
Une intégration par partie nous donne 
\[ 
I_{(\alpha,\beta)} = \underbrace{  \left[ (\alpha t + \beta t^2)\frac{\sin(nt)}{n} \right]_{0}^{\pi} }_{ = 0 }
                       - \frac{1}{n} \int_{0}^{\pi} (\alpha + 2\beta t)\sin(nt)\, dt 
\]
En faisant une intégration par partie sur la deuxième intégrale, on a: 
\begin{eqnarray*}
I_{(\alpha,\beta)} &=& 
\left[ \frac{\alpha + 2\beta t}{n^2}\cos(nt) \right]_{0}^{\pi}  + \int_{0}^{\pi} \frac{2\beta \cos(nt)}{n^2}\, dt 
\\
&=& \left( \frac{\alpha + 2\beta \pi}{n^2}\cos(n\pi) - \frac{\alpha}{n^2} \right) + \underbrace{ \frac{2\beta}{n^2} \left[ \frac{\sin(nt)}{n} \right]_{0}^{\pi} }_{ = 0}  \; \cdotp
\end{eqnarray*}
On obtient $I_{(\alpha,\beta)} = \frac{\alpha + 2 \beta \pi}{n^2}\cos(n \pi) - \frac{\alpha}{n^2}$.
\begin{eqnarray*} 
I_{(\alpha,\beta)} = \frac{1}{n^2} & \Longleftrightarrow &
(\alpha + 2 \beta \pi) \underbrace{ \cos(n\pi) }_{ = (-1)^n} - \alpha = 1 
\\
& \Longleftrightarrow & (\alpha + 2 \beta \pi)(-1)^n - (1 +\alpha) = 0 
{}\end{eqnarray*}
Donc pour tout $n\in \mathbb{N}$
\[
\left\{
\begin{array}{rcl}
(\alpha + 2 \beta \pi)(-1)^n &=& 0 \\
(1 +\alpha) &=& 0
\end{array}
\right.
\]
Ainsi en prenant $\alpha = -1$ et $\beta = \frac{1}{2\pi}$, on obtient: 
\[ 
I_{ (-1,\frac{1}{2\pi}) } = \int_{0}^{\pi} \left(-t + \frac{1}{2\pi} t^2 \right) \cos(nt)\, dt = \frac{1}{n^2} 
\; \cdotp
\]
D'où,
\[ 
\sum_{k = 1}^{n} \frac{1}{k^2} = \int_{0}^{\pi} \left(-t + \frac{1}{2\pi} t^2 \right) \sum_{k = 1}^{n}\cos(kt)\, dt  \qquad \qquad (1) 
\]
Or 
\begin{eqnarray*}
\sum_{k = 1}^{n} \cos(kt) &=& 
\Re \left(\sum_{k = 0}^{n}e^{ikt}\right) - 1 = \cfrac{ \sin(\frac{n+1}{2}t)} { \sin(\frac{t}{2}) }\text{Re}(e^{i \frac{nt}{2} })  - 1 
\\
&=& \cos \left( \frac{nt}{2} \right) \times \cfrac{ \sin( \frac{n}{2}t )\cos( \frac{t}{2} ) + \cos( \frac{n}{2}t )\sin( \frac{t}{2} ) } { \sin(\frac{t}{2}) } - 1 
\\
&=& \cos \left( \frac{nt}{2} \right) \sin \left( \frac{nt}{2} \right) \cot \left( \frac{t}{2} \right) + \underbrace{ \cos^2 \left( \frac{nt}{2} \right) - 1 }_{ = -\sin^2(\frac{nt}{2}) }  
\end{eqnarray*}
Donc $\sum_{k = 1}^{n} \cos(kt) 
= \frac{1}{2}\sin(nt)\cot \left( \frac{t}{2} \right) - \sin^2 \left( \frac{nt}{2} \right)$.

En appliquant ce résultat à (1) on obtient:
\[  
\sum_{k = 1}^{n} \frac{1}{k^2} 
=  \int_{0}^{\pi} \frac{1}{2} \left(-t + \frac{t^2}{2\pi} \right) \sin(nt)\cot \left( \frac{t}{2} \right) \, dt 
- \int_{0}^{\pi} \left(-t + \frac{t^2}{2\pi} \right) \sin^2 \left(\frac{nt}{2} \right) \, dt  
\] 
En posant $ \phi (t) = (-t + \frac{1}{2\pi} t^2)\cot(\frac{t}{2}) $, on a:
\[ \sum_{k = 1}^{+\infty} \frac{1}{k^2} 
= \lim_{n \to +\infty} \int_{0}^{\pi} \phi (t)\sin(nt)\, dt 
- \lim_{n \to +\infty}\int_{0}^{\pi} \left(-t + \frac{1}{2\pi} t^2 \right) \sin^2 \left(\frac{nt}{2} \right) \, dt 
\]
Comme
\[ 
(-t + \frac{1}{2\pi} t^2)\cot \frac{t}{2}
= (-t + \frac{1}{2\pi} t^2) \cfrac{\cos \frac{t}{2}}{\sin \frac{t}{2}}  
\thicksim_{t \to 0} - t \cfrac{\cos \frac{t}{2}} {\frac{t}{2}} = - 2 \cos \frac{t}{2}
\thicksim - 2
\] 
l'application $\phi$ se prolonge par continuité en $0$.
Utilisons le résultat classique suivant: si $h$ est une fonction continue sur $[0, \pi]$, alors
\[ 
\lim_{n \to +\infty} \int_{0}^{\pi} h(t)\sin(nt)\, dt = 0 \quad \text{ et } \quad
\lim_{n \to +\infty} \int_{0}^{\pi} h(t)\cos(nt)\, dt = 0 \; .
\]
appliqué à $\phi$: $\lim_{n \to +\infty} \int_{0}^{\pi} \phi (t)\sin(nt)\, dt = 0$.
De plus, \[ \int_{0}^{\pi} \left(-t + \frac{t^2}{2\pi} \right) \sin^2 \left(\frac{nt}{2} \right) \, dt 
          = \int_{0}^{\pi} \left(-t + \frac{t^2}{2\pi} \right) \left(\frac{1 - \cos(nt)}{2} \right)\, dt 
\]
par conséquent 
\[
\lim_{n \to +\infty} \int_{0}^{\pi} \left(-t + \frac{t^2}{2\pi} \right) \sin^2 \left( \frac{nt}{2} \right) \, dt 
= \int_{0}^{\pi} \frac{1}{2} \left( -t + \frac{t^2}{2\pi} \right) \, dt \; .
\]
Finalement
\[ 
\sum_{k = 1}^{+\infty} \frac{1}{k^2} 
= - \int_{0}^{\pi} \frac{1}{2} \left(-t + \frac{t^2}{2\pi} \right) \, dt 
= -\frac{1}{2} \left( \left[ -\frac{t^2}{2} \right]_{0}^{\pi} + \frac{1}{2\pi} 
\left[ \frac{1}{3}t^3 \right]_{0}^{\pi} \right) 
= \frac{\pi^2}{4} - \frac{\pi^2}{12} = \frac{\pi^2}{6}
\]

\medskip

(\emph{Corrigé de Eugène Ndiaye})

\fincorrection
\nocorrection
\nocorrection
\correction{005108}
Montrons par récurrence que $\forall
n\geq1,\;\sum_{k=1}^{n}\frac{1}{k(k+1)(k+2)}=\frac{n(n+3)}{4(n+1)(n+2)}$.
\textbullet~Pour $n=1$,
$\sum_{k=1}^{1}\frac{1}{k(k+1)(k+2)}=\frac{1}{6}=\frac{1\times(1+3)}{4\times(1+1)(1+2)}$ et la formule proposée est vraie
pour $n=1$.
Soit $n\geq1$. Supposons que $\sum_{k=1}^{n}\frac{1}{k(k+1)(k+2)}=\frac{n(n+3)}{4(n+1)(n+2)}$ et montrons
que $\sum_{k=1}^{n+1}\frac{1}{k(k+1)(k+2)}=\frac{(n+1)(n+4)}{4(n+2)(n+3)}$.

\begin{align*}
\sum_{k=1}^{n+1}\frac{1}{k(k+1)(k+2)}&=\sum_{k=1}^{n}\frac{1}{k(k+1)(k+2)}+\frac{1}{(n+1)(n+2)(n+3)}\\
 &=\frac{n(n+3)}{4(n+1)(n+2)}+\frac{1}{(n+1)(n+2)(n+3)}\;(\mbox{par hypothèse de récurrence})\\
 &=\frac{n(n+3)^2+4}{4(n+1)(n+2)(n+3)}=\frac{n^3+6n^2+9n+4}{4(n+1)(n+2)(n+3)}\\
 &=\frac{(n+1)(n^2+5n+4)}{4(n+1)(n+2)(n+3)}=\frac{(n+1)(n+4)}{4(n+2)(n+3)}
\end{align*}
On a montré par récurrence que :
\begin{center}
\shadowbox{
$\forall n\geq1,\;\sum_{k=1}^{n}\frac{1}{k(k+1)(k+2)}=\frac{n(n+3)}{4(n+1)(n+2)}.$
}
\end{center}
Démonstration directe. Pour $k\geq1$,

$$\frac{1}{k(k+1)(k+2)}=\frac{1}{2}\frac{(k+2)-k}{k(k+1)(k+2)}=\frac{1}{2}\left(\frac{1}{k(k+1)}-\frac{1}{(k+1)(k+2)}\right),
$$
et donc,

\begin{align*}
\sum_{k=1}^{n}\frac{1}{k(k+1)(k+2)}&=\frac{1}{2}(\sum_{k=1}^{n}\frac{1}{k(k+1)}-\sum_{k=1}^{n}\frac{1}{(k+1)(k+2)})
=\frac{1}{2}(\sum_{k=1}^{n}\frac{1}{k(k+1)}-\sum_{k=2}^{n+1}\frac{1}{k(k+1)})\\
 &=\frac{1}{2}\left(\frac{1}{2}-\frac{1}{(n+1)(n+2)}\right)=\frac{n^2+3n}{4(n+1)(n+2)}
=\frac{n(n+3)}{4(n+1)(n+2)}
\end{align*}
\fincorrection
\correction{005109}
\begin{enumerate}
 \item  Montrons par récurrence que~:~$\forall n\geq1,\;\sum_{k=1}^{n}k=\frac{n(n+1)}{2}$.
Pour $n=1$, $\sum_{k=1}^{1}k=1=\frac{1\times(1+1)}{2}$.
Soit $n\geq1$. Supposons que $\sum_{k=1}^{n}k=\frac{n(n+1)}{2}$ et montrons que
$\sum_{k=1}^{n+1}k=\frac{(n+1)(n+2)}{2}$.

\begin{align*}
\sum_{k=1}^{n+1}k&=\sum_{k=1}^{n}k+(n+1)
=\frac{n(n+1)}{2}+(n+1)\;(\mbox{par hypothèse de récurrence})\\
 &=(n+1)(\frac{n}{2}+1)=\frac{(n+1)(n+2)}{2}
\end{align*}
On a montré par récurrence que~:
\begin{center}
\shadowbox{
$\forall n\geq1,\;\sum_{k=1}^{n}k=\frac{n(n+1)}{2}$.
}
\end{center}
On peut donner plusieurs démonstrations directes.

\begin{itemize}
\item[\textbf{1ère demonstration.}] Pour $k\geq1$, $(k+1)^2-k^2=2k+1$ et donc
$\sum_{k=1}^{n}((k+1)^2-k^2)=2\sum_{k=1}^{n}k+\sum_{k=1}^{n}1$ ce qui s'écrit $(n+1)^2-1=2\sum_{k=1}^{n}k+n$ ou encore
$2\sum_{k=1}^{n}k=n^2+n$ ou enfin $\sum_{k=1}^{n}k=\frac{n(n+1)}{2}$.

\item[\textbf{2ème demonstration.}]   On écrit

$$\begin{array}{ccccccccccccc}
1&+&2&+&3&+&\ldots&+&(n-1)&+&n&=&S\\
n&+&(n-1)&+&(n-2)&+&\ldots&+&2&+&1&=&S
\end{array}
$$
et en additionnant (verticalement), on obtient $2S=(n+1)+(n+1)+\ldots+(n+1)=n(n+1)$ d'où le résultat. La même
démonstration s'écrit avec le symbole
sigma~:~

$$2S=\sum_{k=1}^{n}k+\sum_{k=1}^{n}(n+1-k)=\sum_{k=1}^{n}(k+n+1-k)=\sum_{k=1}^{n}(n+1)=n(n+1).$$

\item[\textbf{3ème demonstration.}] On compte le nombre de points d'un rectangle ayant $n$ points de large et $n+1$
points de long. Il y en a $n(n+1)$. Ce rectangle se décompose en deux triangles isocèles contenant chacun $1+2+...+n$ 
points. D'où le résultat.

$$\begin{array}{ccccccccc}
{*}& &{*}&{*}&\ldots& & &\ldots&{*}\\
{*}&{*}& &\ddots& & & & &\vdots\\
{*}&{*}&{*}& & & & & & \\
\vdots& & &\ddots& &\ddots& & &\vdots\\
 & & & & & &{*}&{*}&{*}\\
\vdots& & & & &\ddots& &{*}&{*}\\
*&\ldots& & &\ldots&{*}&{*}& &{*}
\end{array}
$$

\item[\textbf{4ème démonstration.}] Dans le triangle de \textsc{Pascal}, on sait que pour $n$ et
$p$ entiers naturels donnés,

\begin{center}
$C_n^p+C_n^{p+1}=C_{n+1}^{p+1}$.
\end{center}
Donc, pour $n\geq2$ (le résultat est clair pour $n=1$),

$$1+2+...+n=1+\sum_{k=2}^{n}C_k^1=1+\sum_{k=2}^{n}\left(C_{k+1}^2-C_k^2\right)=1+(C_{n+1}^2-1)=\frac{n(n+1)}{2}.$$

\end{itemize}

 \item  Pour $k\geq1$, $(k+1)^3-k^3=3k^2+3k+1$. Donc, pour $n\geq1$~:

$$3\sum_{k=1}^{n}k^2+3\sum_{k=1}^{n}k+\sum_{k=1}^{n}1=\sum_{k=1}^{n}((k+1)^3-k^3)=(n+1)^3-1.$$
D'où,

$$\sum_{k=1}^{n}k^2=\frac{1}{3}\left((n+1)^3-1-3\frac{n(n+1)}{2}-n\right)=\frac{1}{6}(2(n+1)^3-3n(n+1)-2(n+1))=
\frac{1}{6}(n+1)(2n^2+n),$$
et donc
\begin{center}
\shadowbox{
$\forall n\geq1,\;\sum_{k=1}^{n}k^2=\frac{n(n+1)(2n+1)}{6}.$
}
\end{center}
Pour $k\geq1$, $(k+1)^4-k^4=4k^3+6k^2+4k+1$. Donc, pour $n\geq1$, on a

$$4\sum_{k=1}^{n}k^3+6\sum_{k=1}^{n}k^2+4\sum_{k=1}^{n}k+\sum_{k=1}^{n}1=\sum_{k=1}^{n}((k+1)^4-k^4)=(n+1)^4-1.$$
D'où~:

\begin{align*}
\sum_{k=1}^{n}k^3&=\frac{1}{4}((n+1)^4-1-n(n+1)(2n+1)-2n(n+1)-n)=\frac{1}{4}((n+1)^4-(n+1)(n(2n+1)+2n+1)\\
 &=\frac{1}{4}((n+1)^4-(n+1)^2(2n+1))=\frac{(n+1)^2((n+1)^2-(2n+1))}{4}=\frac{n^2(n+1)^2}{4}
\end{align*}
\begin{center}
\shadowbox{
$\forall n\geq1,\;\sum_{k=1}^{n}k^3=\frac{n^2(n+1)^2}{4}=\left(
\sum_{k=1}^{n}k\right)^2.$
}
\end{center}
Pour $k\geq1$, $(k+1)^5-k^5=5k^4+10k^3+10k^2+5k+1$. Donc, pour $n\geq1$,

$$5\sum_{k=1}^{n}k^4+10\sum_{k=1}^{n}k^3+10\sum_{k=1}^{n}k^2+5\sum_{k=1}^{n}k+\sum_{k=1}^{n}1=\sum_{k=1}^{n}((k+1)^5-k^5)
=(n+1)^5-1.$$
D'où~:

\begin{align*}
\sum_{k=1}^{n}k^4&=\frac{1}{5}((n+1)^5-1-\frac{5}{2}n^2(n+1)^2-\frac{5}{3}n(n+1)(2n+1)-\frac{5}{2}n(n+1)-n)\\
 &=\frac{1}{30}(6(n+1)^5-15n^2(n+1)^2-10n(n+1)(2n+1)-15n(n+1)-6(n+1))\\
 &=\frac{1}{30}(n+1)(6n^4+9n^3+n^2-n)=\frac{n(n+1)(6n^3+9n^2+n-1)}{30}
\end{align*}

Finalement,
\begin{center}
\shadowbox{
$
\begin{array}{c}
\forall n\in\Nn^*,\;\sum_{k=1}^{n}k=\frac{n(n+1)}{2}\\
\forall n\in\Nn^*,\;\sum_{k=1}^{n}k^2=\frac{n(n+1)(2n+1)}{6}\\
\forall n\in\Nn^*,\;\sum_{k=1}^{n}k^3=\frac{n^2(n+1)^2}{4}=\left(\sum_{k=1}^{n}k\right)^2\\
\forall n\in\Nn^*,\;\sum_{k=1}^{n}k^4=\frac{n(n+1)(6n^3+9n^2+n-1)}{30}.
\end{array}
$
}
\end{center}
 \item  Soit $p$ un entier naturel. Pour $k\geq1$,

$$(k+1)^{p+1}-k^{p+1}=\sum_{j=0}^{p}C_{p+1}^jk^j.$$
Donc, pour $n\geq1$~:

$$\sum_{j=0}^{p}C_{p+1}^j(\sum_{k=1}^{n}k^j)=\sum_{k=1}^{n}(\sum_{j=0}^{p}C_{p+1}^jk^j)=\sum_{k=1}^{n}((k+1)^{p+1}-k^{p+1})=(
n+1)^{p+1}-1.$$
D'où la formule de récurrence~:

\begin{center}
\shadowbox{$\forall p\in\Nn,\;\forall n\in\Nn^*,\sum_{j=0}^{p}C_{p+1}^{j}S_j=(n+1)^{p+1}-1.$}
\end{center}
\end{enumerate}
\fincorrection
\correction{005143}
\begin{enumerate}
\item  Pour tout naturel non nul $k$, on a $\frac{1}{k(k+1)}=\frac{(k+1)-k}{k(k+1)}=\frac{1}{k}-\frac{1}{k+1}$,
et donc

$$\sum_{k=1}^{n}\frac{1}{k(k+1)}=\sum_{k=1}^{n}(\frac{1}{k}-\frac{1}{k+1})=1-\frac{1}{n+1}=\frac{n}{n+1}.$$

Pour tout naturel non nul $k$, on a
$\frac{1}{k(k+1)(k+2)}=\frac{1}{2}\frac{(k+2)-k}{k(k+1)(k+2)}=\frac{1}{2}(\frac{1}{k(k+1)}-\frac{1}{(k+1)(k+2)})$,
et donc

$$\sum_{k=1}^{n}\frac{1}{k(k+1)(k+2)}=\frac{1}{2}\sum_{k=1}^{n}(\frac{1}{k(k+1)}-\frac{1}{(k+1)(k+2)})=\frac{1}{2}(
\frac{1}{2}-\frac{1}{(n+1)(n+2)})=
\frac{n(n+3)}{4(n+1)(n+2)}.$$

\item  Soit $n\in\Nn^*$.
\begin{itemize}
\item[\textbf{- Calcul de} $\bf{S_1}$.] Posons $P_1=aX^2+bX+c$. On a

$$P_1(X+1)-P_1(X)=a((X+1)^2-X^2)+b((X+1)-X)=2aX+(a+b).$$

Par suite,

\begin{align*}
P_1(X+1)-P_1(X)=X&\Leftrightarrow 2a=1\;\mbox{et}\;a+b=0\Leftrightarrow a=\frac{1}{2}\;\mbox{et}\;b=-\frac{1}{2}\\
 &\Leftarrow
P_1=\frac{X^2}{2}-\frac{X}{2}=\frac{X(X-1)}{2}.
\end{align*}

Mais alors,

$$\sum_{k=1}^{n}k=\sum_{k=1}^{n}(P_1(k+1)-P_1(k))=P_1(n+1)-P_1(1)=\frac{n(n+1)}{2}.$$

\item[\textbf{- Calcul de} $\bf{S_2}$.] Posons $P_2=aX^3+bX^2+cX+d$. On a

$$P_2(X+1)-P_2(X)=a((X+1)^3-X^3)+b((X+1)^2-X^2)+c((X+1)-X)=3aX^2+(3a+2b)X+a+b+c.$$

Par suite,

\begin{align*}
P_2(X+1)-P_2(X)=X^2&\Leftrightarrow 3a=1\;\mbox{et}\;3a+2b=0\;\mbox{et}\;a+b+c=0\Leftrightarrow
a=\frac{1}{3}\;\mbox{et}\;b=-\frac{1}{2}\;\mbox{et}\;c=\frac{1}{6}\\
 &\Leftarrow P_2=\frac{X^3}{3}-\frac{X^2}{2}+\frac{X}{6}=\frac{X(X-1)(2X-1)}{6}.
\end{align*}

Mais alors,

$$\sum_{k=1}^{n}k^2=\sum_{k=1}^{n}(P_2(k+1)-P_2(k))=P_2(n+1)-P_2(1)=\frac{n(n+1)(2n+1)}{6}.$$

\item[\textbf{- Calcul de} $\bf{S_3}$.] Posons $P_3=aX^4+bX^3+cX^2+dX+e$. On a

\begin{align*}
P_3(X+1)-P_3(X)&=a((X+1)^4-X^4)+b((X+1)^3-X^3)+c((X+1)^2-X^2)+d((X+1)-X)\\
 &=4aX^3+(6a+3b)X^2+(4a+3b+2c)X+a+b+c+d.
\end{align*}

Par suite,

\begin{align*}
P_3(X+1)-P_3(X)=X^3&\Leftrightarrow4a=1,\;6a+3b=0,\;4a+3b+2c=0\;\mbox{et}\;a+b+c+d=0\\
 &\Leftrightarrow
a=\frac{1}{4},\;b=-\frac{1}{2},\;c=\frac{1}{4}\;\mbox{et}\;d=0\\
 &\Leftarrow P_3=\frac{X^4}{4}-\frac{X^3}{2}+\frac{X^2}{4}=\frac{X^2(X-1)^2}{4}.
\end{align*}

Mais alors,

$$\sum_{k=1}^{n}k^3=\sum_{k=1}^{n}(P_3(k+1)-P_3(k))=P_3(n+1)-P_3(1)=\frac{n^2(n+1)^2}{4}.$$

\item[\textbf{- Calcul de} $\bf{S_4}$.] Posons $P_4=aX^5+bX^4+cX^3+dX^2+eX+f$. On a

\begin{align*}
P_4(X+1)-P_4(X)&=a((X+1)^5-X^5)+b((X+1)^4-X^4)+c((X+1)^3-X^3)+d((X+1)^2-X^2)\\
 &\;+e((X+1)-X)\\
 &=5aX^4+(10a+4b)X^3+(10a+6b+3c)X^2+(5a+4b+3c+2d)X+a+b+c+d+e.
\end{align*}

Par suite,

\begin{align*}
P_4(X+1)-P_4(X)=X^4&\Leftrightarrow5a=1,\;10a+4b=0,\;10a+6b+3c=0,\;5a+4b+3c+2d=0\\
 &\;\;\mbox{et}\;a+b+c+d+e=0\\
 &\Leftrightarrow
a=\frac{1}{5},\;b=-\frac{1}{2},\;c=\frac{1}{3},\;d=0\;\mbox{et}\;e=-\frac{1}{30}\\
 &\Leftarrow P_4=\frac{X^5}{5}-\frac{X^4}{2}+\frac{X^3}{3}-\frac{X}{30}=\frac{X(X-1)(6X^3-9X^2+X+1)}{30}.
\end{align*}

Mais alors,

$$\sum_{k=1}^{n}k^4=\sum_{k=1}^{n}(P_4(k+1)-P_4(k))=P_4(n+1)-P_4(1)=\frac{n(n+1)(6n^3+9n^2+n-1)}{30}.$$
\end{itemize}
\begin{center}
\shadowbox{
\begin{tabular}{c}
$\forall n\in\Nn^*$,\\
$\sum_{k=1}^{n}k=\frac{n(n+1)}{2},\;\sum_{k=1}^{n}k^2=\frac{n(n+1)(2n+1)}{6},\;
\sum_{k=1}^{n}k^3=\frac{n^2(n+1)^2}{4}=\left(\sum_{k=1}^{n}k\right)^2$\\
$\text{et}\;\sum_{k=1}^{n}k^4=\frac{n(n+1)(6n^3+9n^2+n-1)}{30}$.
\end{tabular}
}
\end{center}

\item  Soit $n\in\Nn^*$.

On rappelle que 
\begin{center}
\shadowbox{
$\forall(a,b)\in]0,+\infty[^2,\;\Arctan a-\Arctan b=\Arctan\frac{a-b}{1+ab}.$
}
\end{center}

Soit alors $k$ un entier naturel non nul. On a

$$\Arctan\frac{1}{k^2+k+1}=\Arctan\frac{(k+1)-k}{1+k(k+1)}=\Arctan(k+1)-\Arctan k.$$

Par suite,

$$\sum_{k=1}^{n}\Arctan\frac{1}{k^2+k+1}=\sum_{k=1}^{n}(\Arctan(k+1)-\Arctan
k)=\Arctan(n+1)-\Arctan1=\Arctan(n+1)-\frac{\pi}{4}.$$

\item  Soit $n\in\Nn^*$.

Pour $k$ entier naturel non nul donné, on a

$$\Arctan\frac{2}{k^2}=\Arctan\frac{(k+1)-(k-1)}{1+(k-1)(k+1)}=\Arctan(k+1)-\Arctan(k-1).$$

Par suite,

\begin{align*}
\sum_{k=1}^{n}\Arctan\frac{2}{k^2}&=\sum_{k=1}^{n}(\Arctan(k+1)-\Arctan
(k-1))=\sum_{k=1}^{n}\Arctan(k+1)-\sum_{k=1}^{n}\Arctan(k-1)\\
 &=\sum_{k=2}^{n+1}\Arctan k-\sum_{k=0}^{n-1}\Arctan k=\Arctan(n+1)+\Arctan n-\Arctan1-\Arctan0\\
 &=\Arctan(n+1)+\Arctan n-\frac{\pi}{4}.
\end{align*}

\end{enumerate}

\fincorrection
\correction{005144}
\begin{enumerate}
\item  Soit $n$ un entier supérieur ou égal à $2$. Parmi les $n^2$ couples $(i,j)$ tels que $1\leq i,j\leq n$, il y
en a $n$ tels que $i=j$ et donc $n^2-n=n(n-1)$ tels que $1\leq i,j\leq n$ et $i\neq j$. Comme il y a autant de couples
$(i,j)$ tels que $i>j$ que de couples $(i,j)$ tels que $i<j$, il y a $\frac{n(n-1)}{2}$ couples $(i,j)$ tels que $1\leq
i<j\leq n$. Finalement,
$$\sum_{1\leq i<j\leq n}^{}1=\frac{n(n-1)}{2}.$$

\item  Soit $n\in\Nn^*$.

$$\sum_{1\leq i,j\leq n}^{}j=\sum_{j=1}^{n}\left(\sum_{i=1}^{n}j\right)=\sum_{j=1}^{n}nj=n\sum_{j=1}^{n}j=n.\frac{n(n+1)}{2}
=\frac{n^2(n+1)}{2}.$$

Soit $n$ un entier supérieur ou égal à $2$.

\begin{align*}
\sum_{1\leq i<j\leq
n}^{}j&=\sum_{j=2}^{n}\left(\sum_{i=1}^{j-1}j\right)=\sum_{j=2}^{n}(j-1)j=\sum_{j=2}^{n}j^2-\sum_{j=2}^{n}j\\
 &=(\frac{n(n+1)(2n+1)}{6}-1)-(\frac{n(n+1)}{2}-1)=\frac{n(n+1)}{2}(\frac{2n+1}{3}-1)\\
 &=\frac{n(n+1)^2}{6}.
\end{align*}

\item  Soit $n\in\Nn^*$.

$$\sum_{1\leq i,j\leq n}^{}ij=(\sum_{1\leq i\leq n}^{}i)(\sum_{1\leq j\leq n}^{}j)=\frac{n^2(n+1)^2}{4}.$$

\item  Soit $n\in\Nn^*$.

$$\sum_{1\leq h,k\leq
n}^{}h^2k^2=\sum_{h=1}^{n}(h^2\sum_{k=1}^{n}k^2)=(\sum_{k=1}^{n}k^2)(\sum_{h=1}^{n}h^2)=\left(\frac{n(n+1)(2n+1)}{6}\right)^
2.$$

Comme d'autre part, $\sum_{h=1}^{n}h^4=\sum_{k=1}^{n}k^4=\frac{n(n+1)(6n^3+9n^2+n-1)}{30}$, on a

$$\sum_{1\leq h,k\leq
n}^{}h^4=\sum_{h=1}^{n}(\sum_{k=1}^{n}h^4)=\sum_{h=1}^{n}nh^4=n\sum_{h=1}^{n}h^4=\frac{n^2(n+1)(6n^3+9n^2+n-1)}{30},$$

et bien sûr $\sum_{1\leq h,k\leq
n}^{}k^4=\frac{n^2(n+1)(6n^3+9n^2+n-1)}{30}$. Par suite,

\begin{align*}
u_n&=\frac{1}{n^5}\left(2.5\frac{n^2(n+1)(6n^3+9n^2+n+14)}{30}-18\frac{n^2(n+1)^2(2n+1)^2}{36}\right)\\
 &=\frac{1}{n^5}(2n^6-2n^6+n^5(\frac{15}{3}-\frac{12}{2})+\mbox{termes de degré au plus}\;4)\\
 &=-1+\mbox{termes tendant vers}\;0
\end{align*}

Par suite,

$$\lim_{n\rightarrow +\infty}u_n=-1.$$

\end{enumerate}
\fincorrection
\correction{005149}
Soient $n\in\Nn^*$ et $a_1$, $a_2$,..., $a_n$, $n$ réels strictement positifs.

\begin{align*}
\left(\sum_{i=1}^{n}a_i\right)\left(\sum_{j=1}^{n}\frac{1}{a_j}\right)&=\sum_{1\leq i,j\leq n}^{}\frac{a_i}{a_j}
=\sum_{i=1}^{n}\frac{a_i}{a_i}+\sum_{1\leq i<j\leq n}^{}(\frac{a_i}{a_j}+\frac{a_j}{a_i})
=n+\sum_{1\leq i<j\leq n}^{}(\frac{a_i}{a_j}+\frac{a_j}{a_i})
\end{align*}

Pour $x>0$, posons alors $f(x)=x+\frac{1}{x}$. $f$ est dérivable sur $]0,+\infty[$ et pour $x>0$,
$f'(x)=1-\frac{1}{x^2}=\frac{(x-1)(x+1)}{x^2}$. $f$ est donc strictement décroissante sur $]0,1]$ et strictement
croissante sur $[1,+\infty[$. $f$ admet ainsi un minimum en $1$. Par suite,

$$\forall x>0,\;f(x)\geq f(1)=1+\frac{1}{1}=2.$$

(\textbf{Remarque.} L'inégalité entre moyenne géométrique et arithmétique permet aussi d'obtenir le résultat~:~
$$\frac{1}{2}(x+\frac{1}{x})\geq\sqrt{x.\frac{1}{x}}=1.)$$

On en déduit alors que

$$\sum_{i=1}^{n}a_i\sum_{j=1}^{n}\frac{1}{a_j}\geq n+\sum_{1\leq i<j\leq n}^{}2=n+2\frac{n^2-n}{2}=n^2.$$
\fincorrection
\correction{005223}
Soit $r$ la raison de la suite $u$.
Pour tout entier naturel $k$, on a 

\begin{center}
$\frac{r}{u_ku_{k+1}}=\frac{u_{k+1}-u_k}{u_ku_{k+1}}=\frac{1}{u_k}-\frac{1}{u_{k+1}}$.
\end{center}
En sommant ces égalités, on obtient~:

$$r\sum_{k=0}^{n}\frac{1}{u_ku_{k+1}}=\sum_{k=0}^{n}\left(\frac{1}{u_k}-\frac{1}{u_{k+1}}\right)=\frac{1}{u_0}-\frac{1}{u_{n+1}}=\frac{u_{n+1}-u_0}{u_0u_{n+1}}=\frac{(n+1)r}{u_0u_{n+1}}.$$
Si $r\neq0$, on obtient $\sum_{k=0}^{n}\frac{1}{u_ku_{k+1}}=\frac{(n+1)}{u_0u_{n+1}}$, et si $r=0$ (et $u_0\neq0$), $u$ est constante et le résultat est immédiat.
\fincorrection
\correction{005224}
Soit $k$ un entier naturel non nul. On sait que $\sum_{i=1}^{k}i^2=\frac{k(k+1)(2k+1)}{6}$. Déterminons alors trois réels $a$, $b$ et $c$ tels que, pour entier naturel non nul $k$, 

$$\frac{6}{k(k+1)(2k+1)}=\frac{a}{k}+\frac{b}{k+1}+\frac{c}{2k+1}\;(*).$$
Pour $k$ entier naturel non nul donné,

$$\frac{a}{k}+\frac{b}{k+1}+\frac{c}{2k+1}=\frac{a(k+1)(2k+1)+bk(2k+1)+ck(k+1)}{k(k+1)(2k+1)}=
\frac{(2a+2b+c)k^2+(3a+b+c)k+a}{k(k+1)(2k+1)}.$$
Par suite,

$$(*)\Leftarrow\left\{
\begin{array}{l}
2a+2b+c=0\\
3a+b+c=0\\
a=6
\end{array}
\right.\Leftrightarrow\left\{
\begin{array}{l}
a=6\\
b=6\\
c=-24
\end{array}
\right.,$$
et donc,

$$\forall n\in\Nn^*,\;\sum_{k=1}^{n}\frac{6}{k(k+1)(2k+1)}=6\left(\sum_{k=1}^{n}\frac{1}{k}+\sum_{k=1}^{n}\frac{1}{k+1}-4\sum_{k=1}^{n}\frac{1}{2k+1}\right).$$
Ensuite, d'après l'exercice \ref{exo:suprou3bis}, quand $n$ tend vers $+\infty$, $\sum_{k=1}^{n}\frac{1}{k}=\ln n+\gamma+o(1)$ puis 

$$\sum_{k=1}^{n}\frac{1}{k+1}=\sum_{k=2}^{n+1}\frac{1}{k}=H_{n+1}-1=-1+\ln(n+1)+\gamma+o(1)=\ln n+\ln\left(1+\frac{1}{n}\right)+\gamma-1+o(1)=\ln n+\gamma-1+o(1).$$
Enfin,
  
\begin{align*}
\sum_{k=1}^{n}\frac{1}{2k+1}&=-1+\sum_{k=1}^{2n+1}\frac{1}{k}-\sum_{k=1}^{n}\frac{1}{2k}=-1+H_{2n+1}-\frac{1}{2}H_n\\
 &=\ln(2n+1)+\gamma-\frac{1}{2}(\ln n+\gamma)-1+o(1)=\ln2+\ln n+\ln\left(1+\frac{1}{2n}\right)+\gamma-\frac{1}{2}\ln n-\frac{1}{2}\gamma-1+o(1)\\
 &=\frac{1}{2}\ln n+\ln2+\frac{1}{2}\gamma-1+o(1)
\end{align*}
Finalement, quand $n$ tend vers $+\infty$, on a

$$\sum_{k=1}^{n}\frac{1}{1^2+2^2+...+k^2}=6\left(\ln n+\gamma+\ln n+\gamma-1-4\left(\frac{1}{2}\ln n+\ln2+\frac{1}{2}\gamma-1\right)\right)=6(3-4\ln2)+o(1).$$
Donc,

\begin{center}
\shadowbox{
$\lim_{n\rightarrow +\infty}\sum_{k=1}^{n}\frac{1}{1^2+2^2+...+k^2}=6(3-4\ln2)$.
}
\end{center}
\fincorrection
\correction{005697}
Pour $n\in\Nn^*$, on a $\left(\sqrt{u_n}-\frac{1}{n}\right)^2$ et donc $0\leqslant\frac{\sqrt{u_n}}{n}\leqslant\frac{1}{2}\left(u_n+\frac{1}{n^2}\right)$. Comme la série terme général $\frac{1}{2}\left(u_n+\frac{1}{n^2}\right)$ converge, la série de terme général $\frac{\sqrt{u_n}}{n}$  converge.
\fincorrection
\correction{005698}
 Pour $n\geqslant2$, $v_n=\frac{u_n+1-1}{(1+u_1)\ldots(1+u_n)}=\frac{1}{(1+u_1)\ldots(1+u_{n-1})}-\frac{1}{(1+u_1)\ldots(1+u_n)}$  et d'autre part $v_1=1-\frac{1}{1+u_1}$. Donc, pour $n\geqslant2$

\begin{center}
$\sum_{k=1}^{n}v_k =1-\frac{1}{(1+u_1)\ldots(1+u_n)}$ (somme télescopique).
\end{center}

Si la série de terme général $u_n$ converge alors $\lim_{n \rightarrow +\infty}u_n=0$ et donc $0<u_n\underset{n\rightarrow+\infty}{\sim}\ln(1+u_n)$. Donc la série de terme général $\ln(1+u_n)$ converge ou encore la suite  $\left(\ln\left(\prod_{k=1}^{n}(1+u_k)\right)\right)_{n\geqslant1}$
converge vers un certain réel $\ell$. Mais alors la suite $\left(\prod_{k=1}^{n}(1+u_k)\right)_{n\geqslant1}$ converge vers le réel strictement positif $P=e^{\ell}$.
Dans ce cas, la suite $\left(\sum_{k=1}^{n}v_k\right)_{n\geqslant1}$ converge vers $1-\frac{1}{P}$.

Si la série de terme général $u_n$ diverge alors la série de terme général $\ln(1+u_n)$ diverge vers $+\infty$ et il en est de même que la suite $\left(\prod_{k=1}^{n}(1+u_k)\right)_{n\geqslant1}$. Dans ce cas, la suite $\left(\sum_{k=1}^{n}v_k\right)_{n\geqslant1}$ converge vers $1$.
\fincorrection
\correction{005699}
Etudions tout d'abord la convergence de la série de terme général $\frac{u_n}{S_n}$.

Si  $\frac{u_n}{S_n}$ tend vers $0$ alors

\begin{center}
$0<\frac{u_n}{S_n}\underset{n\rightarrow+\infty}{\sim}-\ln\left(1-\frac{u_n}{S_n}\right)=\ln\left(\frac{S_{n-1}}{S_n}\right)=\ln(S_n) -\ln(S_{n-1})$.
\end{center}

Par hypothèse, $\lim_{n \rightarrow +\infty}S_n=+\infty$. On en déduit que la série de terme général $\ln(S_n) - \ln(S_{n-1})$ est divergente car  $\sum_{k=1}^{n}\ln(S_k) - \ln(S_{k-1}) =\ln(S_n)-\ln(S_0)\underset{n\rightarrow+\infty}{\rightarrow}+\infty$. Dans ce cas, la série de terme général $\frac{u_n}{S_n}$  diverge ce qui est aussi le cas si $\frac{u_n}{S_n}$ ne tend pas vers $0$.

Donc, dans tous les cas, la série de terme général $\frac{u_n}{S_n}$  diverge.

Si $\alpha\leqslant1$, puisque $S_n$ tend vers $+\infty$, à partir d'un certain rang on a $S_n^\alpha\leqslant S_n$ et donc  $\frac{u_n}{S_n^\alpha}\geqslant\frac{u_n}{S_n}$. Donc, si $\alpha\leqslant1$, la série de terme général  $\frac{u_n}{S_n^\alpha}$ diverge.

Si $\alpha> 1$, puisque la suite $(S_n)$ est croissante,

\begin{center}
$0<\frac{u_n}{S_n^\alpha}=\frac{S_n-S_{n-1}}{S_n^\alpha}=\int_{S_{n-1}}^{S_n}\frac{dx}{S_n^\alpha}\leqslant\frac{dx}{x^\alpha}=\frac{1}{\alpha-1}\left(\frac{1}{S_{n-1}^{\alpha-1}}-\frac{1}{S_n^{\alpha-1}}\right)$,
\end{center}

qui est le terme général d'une série télescopique convergente puisque $\frac{1}{S_n^{\alpha-1}}$ tend vers $0$ quand $n$ tend vers l'infini. Dans ce cas, la série de terme général $\frac{u_n}{S_n^\alpha}$  converge.

\begin{center}
\shadowbox{
La série de terme général $\frac{u_n}{S_n^\alpha}$ converge si et seulement si $\alpha>1$.
}
\end{center}
\fincorrection
\correction{005704}
Pour tout entier naturel non nul $n$, $0<\frac{1}{2^pn^{p-1}}=\sum_{k=1}^{n}\frac{1}{(2n)^p}\leqslant\sum_{k=1}^{n}\frac{1}{(n+k)^p}\leqslant\sum_{k=1}^{n}\frac{1}{n^p}=\frac{1}{n^{p-1}}$ et la série de terme général $u_n$ converge si et seulement si $p > 2$.
\fincorrection
\correction{005705}
(On applique la règle de \textsc{Raabe}-\textsc{Duhamel} qui n'est pas un résultat de cours.) Pour $n\in\Nn$, posons $u_n=\frac{n!}{(a+1)(a+2)\ldots(a+n)}$.

\begin{center}
$\frac{u_{n+1}}{u_n}=\frac{n+1}{a+n+1}=\left(1+\frac{1}{n}\right)\left(1+\frac{a+1}{n}\right)^{-1}\underset{n\rightarrow+\infty}{=}\left(1+\frac{1}{n}\right)\left(1-\frac{a+1}{n}+O\left(\frac{1}{n^2}\right)\right)\underset{n\rightarrow+\infty}{=}1-\frac{a}{n}+O\left(\frac{1}{n^2}\right)$,
\end{center}

et \og on sait \fg~qu'il existe un réel strictement positif $K$ tel que $u_n\underset{n\rightarrow+\infty}{\sim}\frac{K}{n^a}$.  

\fincorrection
\nocorrection
\nocorrection
\nocorrection
\nocorrection
\nocorrection
\nocorrection
\nocorrection
\nocorrection
\correction{005710}
La suite $\left((-1)^n\frac{1}{3n+1}\right)_{n\in\Nn}$ est alternée en signe et sa valeur absolue tend vers $0$ en décroissant. Donc la série de terme général $(-1)^n\frac{1}{3n+1}$, $n\geqslant 1$, converge en vertu du critère spécial aux séries alternées.

Soit $n\in\Nn$.

\begin{center}
$\sum_{k=0}^{n}\frac{(-1)^k}{3k+1}=\sum_{k=0}^{n}(-1)^k\int_{0}^{1}t^{3k}\;dt=\int_{0}^{1}\frac{1-(-t^3)^{n+1}}{1-(-t^3)}\;dt=\int_{0}^{1}\frac{1}{1+t^3}\;dt+(-1)^n\int_{0}^{1}\frac{t^{3n+3}}{1+t^3}\;dt$.
\end{center}

Mais $\left|(-1)^n\int_{0}^{1}\frac{t^{3n+3}}{1+t^3}\;dt\right|=\int_{0}^{1}\frac{t^{3n+3}}{1+t^3}\;dt\leqslant\int_{0}^{1}t^{3n+3}\;dt=\frac{1}{3n+4}$. On en déduit que $(-1)^n\int_{0}^{1}\frac{t^{3n+3}}{1+t^3}\;dt$ tend vers $0$ quand $n$ tend vers $+\infty$ et donc que

\begin{center}
$\sum_{n=0}^{+\infty}\frac{(-1)^n}{3n+1}=\int_{0}^{1}\frac{1}{1+t^3}\;dt$.
\end{center}

Calculons cette dernière intégrale.

\begin{align*}\ensuremath
\frac{1}{X^3+1}&=\frac{1}{(X+1)(X+j)(X+j^2)}=\frac{1}{3}\left(\frac{1}{X+1}+\frac{j}{X+j}+\frac{j^2}{X+j^2}\right)=\frac{1}{3}\left(\frac{1}{X+1}+\frac{-X+2}{X^2-X+1}\right)\\
 &\frac{1}{3}\left(\frac{1}{X+1}-\frac{1}{2}\frac{2X-1}{X^2-X+1}+
 \frac{3}{2}\frac{1}{\left(X-\frac{1}{2}\right)^2+\left(\frac{\sqrt{3}}{2}\right)^2}
 \right).
\end{align*}

 
Donc, 

\begin{center}
$\sum_{n=0}^{+\infty}\frac{(-1)^n}{3n+1}=\frac{1}{3}\left[\ln(t+1)-\frac{1}{2}\ln(t^2-t+1)+\sqrt{3}\Arctan\left(\frac{2t-1}{\sqrt{3}}\right)\right]_0^1=\frac{1}{3}\left(\ln2+\sqrt{3}\left(\frac{\pi}{6}-\left(-\frac{\pi}{6}\right)\right)\right)=\frac{3\ln2+\pi\sqrt{3}}{9}$.
\end{center}

\begin{center}
\shadowbox{
$\sum_{n=0}^{+\infty}\frac{(-1)^n}{3n+1}=\frac{3\ln2+\pi\sqrt{3}}{9}$.
}
\end{center}
\fincorrection
\nocorrection
\correction{004488}
\begin{enumerate}
  \item 
  \item 
  \item Soit $(r_n)$ une énumération de~$Q$.
On pose $f(x) = \sum_{r_n<x} \frac 1{(n+1)^2}$.
$f$ est strictement croissante car pour $x< y$ il existe $n\in\N$
tel que $x< r_n< y$ donc $f(y)-f(x)\ge \frac1{(n+1)^2}$.
Si $x\in\Q$, $x=r_k$ alors $f(x^+)-f(x^-)\ge \frac 1{(k+1)^2}$
d'où $f$ est discontinue en~$x$.
Si $x\in\R\setminus\Q$ et $n\in\N$ alors il existe un voisinage de~$x$
ne contenant aucun $r_i$, $i\le n$ d'où

$f(x^+)-f(x^-)\le \sum_{i> n}\frac1{(i+1)^2}$
et $f$ est continue en~$x$.
\end{enumerate}
 \fincorrection
\correction{004489}
Soit $[a,b]$ de longueur supérieure ou égale à $2\zeta(2)$ et
$F_n = [a,b]\setminus(I_1\cup\dots\cup I_n)$. Alors $(F_n)$ vérifie le
théorème des fermés emboités dans un compact.
\fincorrection
\correction{004490}
\begin{enumerate}
  \item Regroupement à $i+j$ constant $ \Rightarrow $ CV ssi $\alpha > 2$.
  \item Pour $\alpha\ge 1$ on a par convexité~:
             $2^{1-\alpha}(i+j)^\alpha \le i^\alpha+j^\alpha \le (i+j)^\alpha$
             donc il y a convergence ssi $\alpha > 2$.
  \item Il y a une infinité de termes supérieurs à~$1/4$.
  \item $\frac1{a^p+b^q} \le \frac1{2\sqrt a^p\sqrt b^q}  \Rightarrow $ sommable.
\end{enumerate}
 \fincorrection
\correction{004491}
$\sum_{n=0}^{+\infty}\sum_{k=n}^{+\infty}\frac1{k!}=
\sum_{k=0}^{+\infty}\frac{k+1}{k!}=2e$.
\fincorrection
\correction{004492}
$-\frac78\zeta(3)$.
\fincorrection
\correction{004493}
\begin{enumerate}
  \item $\sum_{n=1}^\infty a_{n,n-1}$ diverge.
  \item $\sum_{p=0}^\infty a_{n,p} = \frac1{4n^2}$ si $n\ne 0$, $-\frac{\pi^2}6$ si $n=0$.
             $\sum_{n=0}^\infty\sum_{p=0}^\infty a_{n,p} = -\frac{\pi^2}{8}
             =-\sum_{p=0}^\infty\sum_{n=0}^\infty a_{n,p}$.
\end{enumerate}
 \fincorrection
\correction{004494}
$\sum_{n=0}^{+\infty}\frac{x^{2n+1}}{1-x^{2n+1}}=
\sum_{(n,p)\in\N}^{+\infty}x^{(p+1)(2n+1)}=\sum_{p=0}^{+\infty}\frac{x^{p+1}}{1-x^{2p+2}}$.
\fincorrection
\nocorrection
\correction{004496}
\begin{enumerate}
  \item $|t|<1$.
  \item $S(t) = \sum_{n=1}^\infty \frac{t^n}{1+t^n}
                   = \sum_{n=1}^\infty \sum_{k=1}^\infty (-1)^{k-1}t^{kn}$
    et on peut échanger les deux sommes car il y a convergence absolue.
  \item On suppose $t\in{]0,1[}$.
    $\frac{d}{d x}\Bigl(\frac{t^x}{1-t^x}\Bigr) = \frac{t^x\ln t}{(1-t^x)^2} < 0$
    donc le critère des séries alternées s'applique,
    le reste est majoré en valeur absolue par le premier terme du reste.
    $0\le\frac{t^k(1-t)}{1-t^k} = \frac t{1+\frac1t+\dots+\frac1{t^{k-1}}}\le\frac1k$
    donc le terme général converge uniformément vers~$0$.
    
    Par interversion de limite (puisqu'il y a convergence uniforme) on
    obtient $\lim_{t\to1^-}(1-t)S(t) = \sum_{k=1}^\infty\frac{(-1)^{k-1}}k=\ln 2$.
  \item $S(t) = \sum_{n=1}^\infty \sum_{k=1}^\infty (-1)^{k-1}t^{kn}
                   = \sum_{p=1}^\infty \sum_{kn=p} (-1)^{k-1}t^p
                   = \sum_{p=1}^\infty \sigma(p)t^p$

    avec $\sigma(p) = (\text{nombre de diviseurs impairs de }p)
                    -(\text{nombre de diviseurs pairs de }p) = \sigma_i(p)-\sigma_p(p)$.
   Si $p=2^\alpha q$ avec $q$ impair alors $\sigma_p(p) = \alpha\sigma_i(p)= \alpha\sigma_i(q)$
   donc $\sigma(p) > 0$ ssi $p$ est impair {\it (très joli exercice)}.
\end{enumerate}
 \fincorrection
\correction{004497}
\begin{enumerate}
  \item Il y a convergence si $|z|<1$.
On a alors $f(z) = \sum_{(n,p)\in\N^2}z^{anp + bn+cp}$.
Il y a aussi convergence pour $|z|>1$ lorsque $a>b$ et on a dans ce cas~:
$f(z) = \sum_{(n,p)\in\N\times\N^*}z^{-anp + bn - cp}$ (non symétrique en $b,c$).
  \item $f(z) = \sum_{n=2}^\infty d_nz^n$ avec $d_n =$ nombre de diviseurs de~$n$
dans~$[[1,n-1]]$.
\end{enumerate}
 \fincorrection
\correction{004498}
$A=\zeta(2)^2$, $B=\zeta(2)\zeta(4)$, $C=A/\zeta(4) = 5/2$.
\fincorrection
\correction{004499}
$\ln2 + \frac12\ln(p/q)$.
\fincorrection
\nocorrection
\nocorrection
\correction{004502}
La série converge pour tout $x\notin\{-2,-3,\dots\}$ car le critère
des séries alternées s'applique à partir d'un certain rang (fonction de~$x$).
Il en va de même pour toutes les séries obtenues par dérivations successives
terme à terme, et ces séries convergent localement uniformément (le reste
d'une série vérifant le CSA est majoré en valeur absolue par la valeur
absolue du premier terme figurant dans le reste) donc $f$ est $\mathcal{C}^\infty$.

Pour $|x|<2$ on a
$$f(x) = \sum_{k=2}^\infty\sum_{n=0}^\infty\frac{(-1)^k}k\Bigl(\frac{-x}k\Bigr)^n
      = \sum_{k=2}^\infty\sum_{n=0}^\infty\frac{(-1)^{k+n}}{k^{n+1}}x^n
      = (1-\ln2) + \sum_{n=1}^\infty(-1)^n(1-(1-2^{-n})\zeta(n+1))x^n.$$
\fincorrection
\correction{004403}
$\cos(x+iy) = \cos x \ch y - i\sin x \sh y  \Rightarrow  \cos z \in {[-1,1]}$
          si et seulement si $z \in \R$.
\fincorrection
\correction{004404}
Mettre $1+\frac zn$ sous forme trigonométrique.
\fincorrection
\correction{004405}
Développement en série.
\fincorrection
\correction{004406}
$\Bigl|\frac{e^z-1}z\Bigr|^2 = \frac{e^{2x}+1-2e^x\cos y}{x^2+y^2}$.
Après simplifications, on est ramené à prouver que
$x^2(1-\cos y)\le y^2(\ch x - 1)$, ce qui est vrai car on peut caser
$\frac12x^2y^2$ entre les deux. Il y a égalité si et seulement si $y=0$.
\fincorrection
\nocorrection
\correction{004408}
$e^{x+iy} = x+iy \Leftrightarrow \begin{cases}x = y/\tan y \cr e^{-y/\tan y}
         = \sin y/y.\cr\end{cases}$
         Au voisinage de $2k\pi^+$, $e^{-y/\tan y} < \sin y/y$ (point plat)
         et au voisinage de $(2k+1)\pi^-$, $e^{-y/\tan y} > \sin y/y$
         (limite infinie).
\fincorrection
\correction{004409}
\begin{enumerate}
  \item $z \equiv \pm i\ln(2+\sqrt3\,) (\mathrm{mod}\,{2\pi})$.
  \item $z \equiv i\pi(\mathrm{mod}\,{2i\pi})$.
  \item $z \equiv 0 (\mathrm{mod}\, {2\pi})$ ou
             $z \equiv 0 (\mathrm{mod}\, {2j\pi})$ ou
             $z \equiv 0 (\mathrm{mod}\, {2j^2\pi})$.
  \item $\Leftrightarrow 6e^{2iz} -(7+5i)e^{iz} + 2 = 0 \Leftrightarrow
               \begin{cases} e^{iz} = 1+i\hfill:
                       & z \equiv \pi/4 - i\ln\sqrt2      (\mathrm{mod}\,{2\pi})\cr
                       e^{iz} = (1-i)/6  :
                       & z \equiv -\pi/4 - i\ln(\sqrt2/6) (\mathrm{mod}\,{2\pi}).\cr\end{cases}$
\end{enumerate}
 \fincorrection
\correction{004410}
$|\cos(x+iy)|^2 = \cos^2x + \sh^2y = \ch^2y - \sin^2x  \Rightarrow  \sup = \ch1$.
         \par
         $|\sin(x+iy)|^2 = \sin^2x + \sh^2y = \ch^2y - \cos^2x$.
         \`A $x$ fixé, le module augmente avec $|y|$, donc le maximum est atteint
         au bord du disque.\par
         $\varphi(\theta) = \sin^2\cos\theta + \sh^2\sin\theta  \Rightarrow 
          \varphi'(\theta) = \sin2\theta\left(
             \frac{\sh(2\sin\theta)}{2\sin\theta}
           - \frac{\sin(2\cos\theta)}{2\cos\theta}\right) \Rightarrow  \sup = \sh1$.
\fincorrection
\nocorrection
\correction{004412}
Si $x$ est vecteur propre de $M$ il l'est aussi de $\exp(M)$ donc
$x = ke_1$ et la valeur propre associée est $\alpha\in\C$ tel que $e^\alpha = 2i$
($\alpha = \ln2 + i(\frac\pi2+2k\pi)$, $k\in\Z$).
On a donc $M = \begin{pmatrix}\alpha&\beta\cr0&\alpha\cr\end{pmatrix}$,
$\exp(M) = \begin{pmatrix}e^\alpha&e^\alpha\beta\cr0&e^\alpha\cr\end{pmatrix}$ d'où
$\beta = \frac{1-i}2$.
\fincorrection
\nocorrection
\nocorrection
\nocorrection
\nocorrection
\correction{004413}
\begin{enumerate}
  \item $\sim -\frac e{2n}  \Rightarrow $ DV.
  \item $\sim \frac \alpha{2^{\alpha-1}}e^{n(\alpha-2)}  \Rightarrow $
             CV ssi $\alpha < 2$.
  \item $\sim -\frac 3{n^2}  \Rightarrow $ CV.
  \item $\sim \frac 1{n^2}  \Rightarrow $ CV.
  \item $\sim \sqrt{\frac 2{n^3}}  \Rightarrow $ CV.
  \item cv ssi $|a| \ne 1$.
  \item Série alternée $ \Rightarrow $ CV.
  \item Série alternée $ \Rightarrow $ CV.
  \item Harmonique + alternée $ \Rightarrow $ DV.
  \item d'Alembert $ \Rightarrow $ CV.
  \item $\le \frac{(n-1)(n-1)!+n!}{(n+2)!} \le \frac 2{(n+1)(n+2)}  \Rightarrow $ CV.
  \item $= \frac{(-1)^{n-1}}{n+1} +  O\Bigl(\frac1{n^2}\Bigr)  \Rightarrow $ CV.
  \item Décomposition en 3 séries alternées $ \Rightarrow $ CV.
  \item $=\frac{(-1)^n}{2\sqrt n} - \frac1{8n} +  O(n^{-3/2})  \Rightarrow $ DV.
  \item Regroupement de termes $ \Rightarrow $ DV.
  \item Regroupement par paquets + CSI $ \Rightarrow $ CV.
  \item Terme général ne tend pas vers zéro, DV.
  \item $= \frac1{n^{\ln\ln n}}  \Rightarrow $ CV.
\end{enumerate}
 \fincorrection
\correction{004414}
$P(n) = n^3+\frac34n+C$.
\fincorrection
\correction{004415}
$=\frac{(-1)^n}{\sqrt{n(n+1)}}+O\Bigl(\frac1{n^2}\Bigr)  \Rightarrow $ converge.
\fincorrection
\correction{004416}
$u_n = \frac{(-1)^n}{n^{\alpha/2}}-\frac1{2n^{3\alpha/2}}
+  o\Bigl(\frac1{n^{3\alpha/2}}\Bigr)$,
il y a convergence ssi $\alpha>\frac23$.
\fincorrection
\correction{004417}
Effectuer un développement asymptotique pour les deux premières. Elles convergent si et seulement si $\alpha>\frac12$.
La troisième diverge par comparaison série-intégrale.
\fincorrection
\correction{004418}
$\frac{u_{2n+1}}{u_{2n-1}}\to ab$ et
$\frac{u_{2n}}{u_{2n-2}}\to ab$ (lorsque $n\to\infty$) donc il y a convergence
si $|ab| < 1$.
\fincorrection
\nocorrection
\correction{004420}
$n=21$, $S \approx 0.65314389$.
\fincorrection
\correction{004421}
$1 \Rightarrow 2$ par comparaison série-intégrale. Contre-exemple pour $(2)\not\Rightarrow(1)$~:
$u_n = e^{(n+1)^2}-e^{n^2}$, $S_n = e^{(n+1)^2}-1$, $f(t)=\frac1{(t+2)\ln(t+2)}$.
\fincorrection
\correction{004422}
$\frac{n^2}{(n^2+1)^2}-\frac{1}{n^2-1}
         = -\frac{3n^2+1}{(n^2+1)^2(n^2-1)} \ge -\frac{4}{n^4}$ pour
         $n\ge 3$.

Donc $S = \sum_{n=1}^N \frac{n^2}{(n^2+1)^2}
        + \sum_{n=N+1}^\infty \frac{1}{n^2-1} + R_N$ avec
     $-\frac{4}{3N^3}\le R_N\le 0$ et
     $\sum_{n=N+1}^\infty \frac{1}{n^2-1} = \frac{N+\frac12}{N(N+1)}$.

Pour $N=25$ on obtient~: $0.76981 < S < 0.76990$.
\fincorrection
\nocorrection
\correction{004424}
Si $\sum u_n$ et $\sum v_n$ convergent alors $n^2u_n\to \infty$ (lorsque $n\to\infty$)
donc $u_nv_n \sim 1/n^2$. Alors les suites $(\sqrt{u_n})$ et $(\sqrt{v_n})$
sont de carrés sommables tandis que la suite $(\sqrt{u_nv_n})$ n'est pas sommable,
c'est absurde.

Si $\sum u_n$ diverge on ne peut rien dire~: avec $u_n=1$ on a $\sum v_n$
convergente tandis qu'avec $u_n=\frac1{n}$ on a $\sum v_n$ divergente.
\fincorrection
\correction{004425}
\begin{enumerate}
  \item $u_1 + \dots + u_n = 1 - \frac1{(1+a_1)\dots(1+a_n)} \le 1$.
  \item $\ln\bigl((1+a_1)\dots(1+a_n)\bigr)
              = \sum_{k=1}^n \ln\left(1+\frac1{\sqrt k}\right)
              \to +\infty  \Rightarrow  \sum u_n = 1$ lorsque $n\to\infty$.
\end{enumerate}
 \fincorrection
\correction{004426}
Regroupement de termes par valeur constante de $p_k$
         $ \Rightarrow  \sum_{k=1}^\infty \frac1{a^{p_k}}
         = \sum_{p=1}^\infty \frac{10^p-10^{p-1}}{a^p}
         = \frac{9}{a-10}$.
\fincorrection
\nocorrection
\correction{004428}
\begin{enumerate}
  \item 
  \item $|v_n| =  O(n^{-3/2})  \Rightarrow $ CV.
\end{enumerate}
 \fincorrection
\correction{004429}
Série alternée.
\fincorrection
\correction{004430}
\begin{enumerate}
  \item $\frac 34$.
  \item $\frac 14$.
  \item $S_p - (p+1)S_{p+1} = S_p - \frac 1{(p+1)!} \Rightarrow 
             S_p=\frac 1{pp!}$.
  \item $\frac{23}{144}$.
  \item $\ln 3$.
  \item $-\ln 2$.
  \item $\ln\left(\frac{\sin2\alpha}{2\alpha}\right)$.
  \item $\frac1\alpha-2\mathrm{cotan}(2\alpha)$.
  \item $109 - 40e$.
  \item $\frac{x^p}{(1-x)^{p+1}}$ pour $|x|<1$ par récurrence.
  \item $\frac x{(1-x)^2}$ si $|x| <1$,
             $\frac 1{(1-x)^2}$ si $|x| > 1$.
  \item $S_n = \sum_{q=0}^\infty\sum_{r=1}^{n-1}
                       \frac r{(qn+r)(qn+r+1)}
                  = \sum_{q=0}^\infty\sum_{r=1}^{n-1}
                       \frac r{qn+r} - \frac r{qn+r+1}$.
                  \par
             $S_n = \sum_{q=0}^\infty \left(\frac 1{qn+1} + \frac 1{qn+2}
                       + \dots + \frac 1{qn+n} - \frac 1{q+1}\right)
                  = \lim_{N\to\infty}\left(\sum_{k=1}^{(N+1)n} \frac 1k
                       - \sum_{k=1}^{N+1} \frac 1k \right)
                  = \ln n$.
\end{enumerate}
 \fincorrection
\correction{004431}
Si $n+1$ n'est pas un carré alors $u_n=0$ donc
$\sum_{n=1}^\infty u_n = \sum_{k=2}^\infty u_{k^2-1}
= \sum_{k=1}^\infty \frac1{k^2-1} = \frac34$.
\fincorrection
\correction{004432}
\begin{enumerate}
  \item $y = e^{-x}(a\cos x + b\sin x)$,
             $y = e^{-x}\sin x + e^{-2x}(cx+d)$.
  \item $u_n = \frac{(-1)^ne^{-n\pi}(e^\pi+1)}2$,
             $\sum_{n=0}^\infty u_n = \frac12$.
\end{enumerate}
 \fincorrection
\correction{004433}
$\frac{\pi^2}3 - 3$.
\fincorrection
\correction{004434}
$\frac 1{1^2 + 2^2 + \dots + k^2} = \frac 6k + \frac 6{k+1}
         - \frac {24}{2k+1}
          \Rightarrow  s_n = 18 - 24\sum_{k=1}^{2n+1} \frac{(-1)^{k+1}}k
         + \frac 6{n+1} \to 18 - 24\ln 2$ lorsque $n\to\infty$.
\fincorrection
\correction{004435}
$\begin{cases} a+b = -1 \cr a+2b = 0\cr \end{cases} \Leftrightarrow a=-2$, $b=1$, $S = -\ln 2$.
\fincorrection
\correction{004436}
$\tan s_n = n+1$ par récurrence et
         $s_n \le \sum_{k=0}^\infty \frac 1{k^2+k+1} \le 1
         + \sum_{k=0}^\infty \frac1{n(n+1)} = 2$.
\fincorrection
\correction{004437}
\begin{enumerate}
  \item $\sim \frac a{n^2}$.
  \item $S(a) \ge \sum_{k=0}^n \Arctan(k+a)-\Arctan k \to
              \frac\pi2 + \Arctan 1 + \Arctan\frac12 + \dots
              + \Arctan\frac1n
               \Rightarrow  S(a) \to +\infty$ lorsque $a\to+\infty$ .
\end{enumerate}
 \fincorrection
\correction{004438}
Le déport maximal entre la première pièce et la dernière pour une
pile de $n$ pièces est $\frac 12 + \frac 14 + \dots + \frac 1{2(n-1)}$
(en diamètre d'une pièce).
Il dépasse 1 pour $n > 4$.
\fincorrection
\correction{004439}
\begin{enumerate}
  \item $2(\sqrt 2-1)\sqrt n$.
  \item $\ln(\ln n)$.
\end{enumerate}
 \fincorrection
\correction{004440}
$u_n \sim n\ln^2 n  \Rightarrow $ CV.
\fincorrection
\correction{004441}
2997.
\fincorrection
\correction{004442}
$\sqrt{\frac n2}$.
\fincorrection
\nocorrection
\correction{004444}
$T_{n+1}-T_n = \frac1{n+1}-\ln\Bigl(\frac{n+1}n\Bigr) = \frac1{n+1} -  \int_{t=n}^{n+1}\frac{d t}t < 0$

$S_{n+1}-S_n = \frac2n-\frac1{n+1}-\ln\Bigl(\frac{n+1}n\Bigr)
=\frac1n -  \int_{t=n}^{n+1}\Bigl(\frac1t-\frac1{t^2}\Bigr)\,d t > 0$.
\fincorrection
\correction{004445}
$\frac{u_{n,k}}k = \frac nk - \bigl[\frac nk\bigr]$, donc
$v_n = \frac1n\sum_{k=1}^n\frac{u_{n,k}}k$ est une somme de Riemann pour
l'intégrale $I =  \int_{t=0}^1 \Bigl(\frac1t - \bigl[\frac1t\bigr]\Bigr)\,d t$.
La fonction $\varphi$ : $t \mapsto\frac1t - \bigl[\frac1t\bigr]$ est Riemann-intégrable
sur $[0,1]$, donc $v_n\to I$ lorsque $n\to\infty$.

Calcul de~$I$~: $I_n =  \int_{t=1/n}^1 \Bigl(\frac1t - \bigl[\frac1t\bigr]\Bigr)\,d t
                     = \ln n - \sum_{k=1}^n \int_{t=\frac1{k+1}}^{\frac1k}k\,d t
                     = \ln n - \sum_{k=1}^n\frac1{k+1} \to 1-\gamma = I$ lorsque $n\to\infty$.
\fincorrection
\correction{004446}
\begin{enumerate}
  \item Comparaison série-intégrale~: $u_n\sim\frac{\ln^2 n}2$.
  \item Comparaison série-intégrale encore ($v_n$ est la somme des aires entre
    les rectangles aux points entiers et la courbe de~$t\to\ln(t)/t$).
  \item $v_n-\ell = -\sum_{k=n}^\infty\Bigl( \int_{t=k}^{k+1}\frac{\ln t}t\,d t - \frac{\ln(k+1)}{k+1}\Bigr) = -\sum_{k=n}^\infty w_k$
    avec $w_k\sim\frac{\ln k}{2k^2}$ donc $v_n-\ell\sim- \int_{t=n}^{+\infty}\frac{\ln t}{2t^2}\,d t\sim-\frac{\ln n}{2n}$.
\end{enumerate}
 \fincorrection
\correction{004447}
$\sum_{k=0}^{n-1}\frac1{n^2-k^2}
= \frac1{2n}\sum_{k=0}^{n-1}\Bigl(\frac1{n-k} + \frac1{n+k}\Bigr)
= \frac1{2n}\Bigl(\sum_{k=1}^{2n-1}\frac1k +\frac1n\Bigr)
\sim \frac{\ln n}{2n}$.
\fincorrection
\correction{004448}
\begin{enumerate}
  \item 
  \item 
  \item 
  \item 
  \item $S_n + \frac1{\ln(n+1)} \le S \le S_n + \frac1{\ln n}$.
             Pour $n = 60$ : $2.06857 < S < 2.06956$.
\end{enumerate}
 \fincorrection
\correction{004449}

La fonction $t \mapsto \frac{1}{t^x}$ étant décroissante sur l'intervalle $[n, n+1]$,
$\displaystyle 0 < \frac{1}{t^x} \leq \frac{1}{n^x}$. Donc,
\begin{eqnarray*}
\int_{n}^{n+1} \frac{dt}{t^x}  & \leq & \int_{n}^{n+1} \frac{dt}{n^x} = \frac{1}{n^x} \quad \text{par positivité de l'intégrale}
\end{eqnarray*}
puis
\[
\int_{1}^{N+1} \frac{dt}{t^x} = \sum_{n=1}^{N} \int_{n}^{n+1} \frac{dt}{t^x}
\le \sum_{n=1}^{N} \frac{1}{n^x} \quad \text{ par la relation de Chasles.}
\]
De même sur $[n-1,n]$, $\displaystyle \frac{1}{t^x} \geq \frac{1}{n^x} > 0$ et
\[
\int_{n-1}^{n} \frac{dt}{t^x} \ge \int_{n-1}^{n} \frac{dt}{n^x} = \frac{1}{n^x}
\]
\[
\int_{1}^{N} \frac{dt}{t^x} = \sum_{n=2}^{N} \int_{n-1}^{n} \frac{dt}{t^x}
\ge \sum_{n=2}^{N} \frac{1}{n^x} \quad \text{par la relation de Chasles.}
\]
Finalement
\[
1 + \int_{1}^{N} \frac{dt}{t^x} \ge \sum_{n=1}^{N} \frac{1}{n^x} \, \cdotp
\]
\begin{equation}
\text{Donc, on a:} \qquad \int_{1}^{N+1} \frac{dt}{t^x} \leq \sum_{n=1}^{N} \frac{1}{n^x} \leq 1 + \int_{1}^{N} \frac{dt}{t^x} \, \cdotp \label{I}
\end{equation}
Calculons dans un premier temps, $ \int_{1}^{N} \frac{dt}{t^x}$:
\[ 
\int_{1}^{N} \frac{dt}{t^x} = \left[\frac{t^{1-x}}{1-x}\right]_{1}^{N} = \frac {N^{1-x}-1}{1-x} 
\xrightarrow[N \to +\infty]{} \frac{1}{x-1} 
\]
avec la même limite pour $\int_{1}^{N+1} \frac{dt}{t^x}$.
Ainsi on a montré que $\sum_{n \geq 0} \frac{1}{n^x}$, série de Riemann avec $x > 1$, est convergente.
On déduit alors de \eqref{I}, en faisant tendre $N$ vers $+\infty$:
\[ \begin{array}{lrcccl}
& \frac{1}{x - 1} & \leq & \sum_{n = 1}^{+\infty} \dfrac{1}{n^x} & \leq & 1 + \frac{1}{x - 1}
\\
\text{donc} \quad & \frac{x- 1}{x - 1} & \leq & (x - 1) \zeta (x) & \leq & (x - 1)\left(1 + \frac{1}{x - 1}\right) \\
\text{puis} \qquad & 1 & \leq & (x - 1) \zeta (x) & \leq & x \, .
\end{array} \]
D'où, par le théorème des gendarmes:
\[ 
\lim_{x \rightarrow 1^{+}} (x - 1) \zeta (x) = 1 \, . 
\]

\medskip

(\emph{Corrigé d'Antoine Poulain})

\fincorrection
\correction{004450}
Si $u_n \to 0$, alors $v_n \sim u_n$;
         sinon, $v_n$ \hbox to 0pt{\hskip 3mm$/$\hss} $\longrightarrow 0$.
\fincorrection
\correction{004451}
\begin{enumerate}
  \item 
  \item $\frac r{(1-r)^2}$.
\end{enumerate}
 \fincorrection
\correction{004452}
On remarque déjà que $\sum u_i$ diverge car $u_n\sim\frac{U_n}{n\alpha}\ge \frac{U_1}{n\alpha}$.
On calcule $\sum_{k=0}^n ku_k$ par parties~:
$$\sum_{k=0}^n ku_k = \sum_{k=1}^n k(U_k - U_{k-1}) = nU_n -\sum_{k=0}^n U_k$$
Comme $U_n\sim\alpha nu_n$, terme général strictement positif d'une série divergente,
on a $\sum_{k=0}^n U_k \sim \alpha\sum_{k=0}^n ku_k$ d'où~:
$(1+\alpha)\sum_{k=0}^n ku_k\sim nU_n$ et lorsque $n\to\infty$ :
$$\frac1{n^2u_n}\sum_{k=0}^n ku_k\sim \frac{nU_n}{(1+\alpha)n^2u_n} \to \frac\alpha{1+\alpha}.$$
\fincorrection
\correction{004453}
$S_n = \sum_{k=0}^n ku_k  \Rightarrow 
\sum_{k=0}^n u_k = \sum_{k=1}^{n-1}\frac{S_k}{k(k+1)} - S_0 + \frac{S_n}n$.
\fincorrection
\correction{004454}
\begin{enumerate}
  \item 
  \item 
  \item     $kr^k = k(u_k - u_{k+1})$ avec $u_k = \frac{r^k}{1-r}$
    donc $\sum_{k=1}^\infty kr^k = \sum_{k=1}^\infty\frac{r^k}{1-r} = \frac r{(1-r)^2}$.
    \par
    De même, $S_n = \sum_{k=n}^\infty kr^k = \frac{(n-1)r^n}{1-r} +
     \sum_{k=n}^\infty\frac{r^k}{1-r} = \frac{nr^n}{1-r} + \frac{r^{n+1}}{(1-r)^2}$.
    
    $k^2r^k = k(S_k - S_{k+1})$ et $(S_k)$ décroît d'où
    $\sum_{k=1}^\infty k^2r^k = \sum_{k=1}^\infty S(k)
    = \sum_{k=1}^\infty \Bigl(\frac{kr^k}{1-r}+\frac{r^{k+1}}{(1-r)^2}\Bigr)
    = \frac{r+r^2}{(1-r)^3}$.
\end{enumerate}
 \fincorrection
\correction{004455}
\begin{enumerate}
  \item 
  \item $p_n = \frac{u_0}{S_n} \to 0$ donc la série de terme général
              $\ln\left(1-\frac{u_n}{S_n}\right)$ diverge.
\end{enumerate}
 \fincorrection
\correction{004456}
Méthode des rectangles~:
$\sum_{k=0}^n\frac{a_k-a_{k+1}}{a_{k+1}}\ge
 \int_{t=a_{n+1}}^{a_0}\frac{d t}t \to +\infty$ lorsque $k\to\infty$.

Si $a_k\sim a_{k+1}$ la série donnée diverge donc. Sinon, elle diverge
aussi car son terme général ne tend pas vers~$0$.
\fincorrection
\correction{004457}
$\sum_{n=1}^N v_n = \sum_{k=1}^{2N-1} u_k\sum_{k/2<n\le k}\frac 1n
           \Rightarrow  \frac12\sum_{k=1}^{2N-1} u_k \le \sum_{n=1}^N v_n \le
              2\sum_{k=1}^{2N-1} u_k$.
\fincorrection
\correction{004458}
$\sum_{k=1}^n v_k + nv_n = \sum_{k=1}^n u_k$.\par
         Si $\sum u_n$ converge, $\sum v_n$ converge aussi (SP majorées) et
         $nv_n \to \ell  \Rightarrow  \ell = 0$.\par
         Si $\sum u_n$ diverge et $\sum v_n$ converge, alors $nv_n \to +\infty$,
         contradiction.
\fincorrection
\correction{004459}
$\sum_{k=1}^n v_k = \sum_{k=1}^n ku_k\sum_{p=k}^n \frac 1{p^2} \le
         \sum_{k=1}^n \frac{ku_k}{k-1}  \Rightarrow $ CV.
\fincorrection
\correction{004460}
$\frac12\sum_{k=1}^{n+1} v_k \le \sum_{k=1}^{2^{n+1}} u_k \le
         \sum_{k=0}^n v_k$.
\fincorrection
\correction{004461}
Pour $n > 2$, $u_{n+1} < \frac1n$
     donc $u_{n+2} > \frac1{(n+1)e^{1/n}} \sim \frac1n$ donc la série
     diverge.
\fincorrection
\nocorrection
\nocorrection
\correction{004464}
\begin{enumerate}
  \item 
  \item $\ln(v_{n+1})-\ln(v_n) = \ln\left(1+\frac{a-b+1}{n+b-1}\right)  \Rightarrow 
             \begin{cases} \text{si } a-b+1 > 0, v_n\to+\infty \cr
                     \text{si } a-b+1 = 0, v_n = \text{cste}\cr
                     \text{si } a-b+1 < 0, v_n\to0.\cr \end{cases}$
  \item $(n+b)u_{n+1} - (n+a)u_n = 0  \Rightarrow 
              (n+b)u_{n+1} + (b-a-1)\sum_{k=1}^n u_k - au_0 = 0  \Rightarrow 
              \sum_{k=0}^\infty u_k = \frac{(b-1)u_0}{b-a-1}$.
\end{enumerate}
 \fincorrection
\correction{004465}
La suite $(u_n)$ est croissante donc tend vers $\ell\in{]0,+\infty]}$.
On a $\ell$ fini si et seulement si la série télescopique
$\sum(u_{n+1}-u_n) = \sum\frac1{n^au_n}$ est convergente, soit si et
seulement si $a>1$.

Pour $a< 1$ on a $u_{n+1}^2 = u_n^2 + \frac2{n^a} +  o\Bigl(\frac2{n^a}\Bigr)$
donc $u_{n+1}^2-u_n^2\sim \frac2{n^a}$ et $u_n\sim\sqrt{\frac{2n^{1-a}}{1-a}}$
(sommation des relations de comparaison).

Pour $a=1$ on a de même $u_n\sim \sqrt{2\ln n}$.
\fincorrection
\correction{004466}
$\alpha > 1  \Rightarrow  \sum u_n$ cv et vaut $\zeta(\alpha)^2$.\par
         $\alpha < 1  \Rightarrow  \sum_{n=1}^{2N} u_n \ge
                         \sum_{k=1}^N \frac 1{k^\alpha}
                      \Rightarrow  \sum u_n$ dv.
\fincorrection
\nocorrection
\correction{004468}
$\frac{a}{(1-a)^2}$ et $\frac{a+a^2}{(1-a)^3}$.
\fincorrection
\nocorrection
\correction{004470}
\begin{enumerate}
  \item Césaro.
  \item $v_0 + v_1 + \dots + v_n = 2(u_0 + u_1 + \dots + u_n) - v_n$.
\end{enumerate}
 \fincorrection
\correction{004471}
$|a_n| \le M  \Rightarrow  \left|\sum_{n=2}^\infty \frac {a_n}{n^p} \right|
        \le M\sum_{n=2}^\infty \frac 1{n^p}
        \le M \int_{t=1}^\infty \frac{d t}{t^p}
        = \frac M{p-1}  \Rightarrow  a_1 = 0$.
\fincorrection
\correction{004472}
Démonstration pour $x_1$ :
$\sum x_n = 0$, $\sum x_{2n} = 0  \Rightarrow  \sum_{n\text{ impair}} x_n = 0$.
On retire les multiples impairs de 3 ($\sum x_{3n} - \sum x_{6n} = 0$)
$ \Rightarrow  \sum_{n\not\equiv 0[2] ; n\not\equiv 0[3]} x_n = 0$.
On retire les multiples restants de $5,7,\dots$
On obtient ainsi une suite $(s_p)_{p \text{ premier}}$ nulle qui converge vers
$x_1$, donc $x_1 = 0$.\par
Peut-on se passer de la convergence absolue ?
\fincorrection
\correction{004473}
\begin{enumerate}
  \item Récurrence sur $p$.
  \item Transformation d'Abel et interversion de sommations :
         $\sum_{n=0}^p v_n
         = \sum_{k=0}^p \frac{C_{p+1}^{k+1}}{2^p}\sum_{n=0}^k u_n$.
         \par Théorème de Césaro $ \Rightarrow  \sum v_n = 2\sum u_n$ .
\end{enumerate}
 \fincorrection
\correction{004474}
\begin{enumerate}
  \item $nu_{2n} \le \sum_{k=n+1}^{2n} u_k$,
             $nu_{2n+1} \le \sum_{k=n+2}^{2n+1} u_k$.
  \item $\varepsilon > 0$ : Pour $k$ suffisament grand,
             $u_k \le \frac \varepsilon k$,
             donc $u_k \ge \frac 1n  \Rightarrow  k\le n\varepsilon$.
             Alors $\sum_{u_k \ge 1/n} \frac 1{u_k} \le n^2\varepsilon + Kn$.
\end{enumerate}
 \fincorrection
\correction{004475}
\begin{enumerate}
  \item TAF : $\exists\ x_n \in {[R_{n+1},R_n]} \text{ tel que }
             R_n^{1-p} - R_{n+1}^{1-p} = (1-p) \frac {R_n-R_{n+1}}{x_n^p}
             \ge (1-p) \frac {a_n}{R_n^p}$.
             Donc, $\sum_{n=0}^{\infty} \frac {a_n}{R_n^p} \le
             \frac {A^{1-p}}{1-p}$.
  \item C'est $\frac 1{1-p}$ : Pour $a_n = k^n$,
             $A^{p-1}\sum_{n=0}^{\infty} \frac {a_n}{R_n^p}
             = \frac {1-k}{1-k^{1-p}} \to \frac 1{1-p}$ lorsque $k\to1^-$.
\end{enumerate}
 \fincorrection
\correction{004476}
$(u_n)$ est croissante. Si la suite $(u_n)$ converge alors
         $a_n = u_n(u_{n+1} - u_n) \le M(u_{n+1} - u_n)$ donc les sommes
         partielles de $\sum a_n$ sont bornées.\par
         Si $\sum a_n$ converge, alors $u_{n+1} - u_n = \frac{a_n}{u_n}
         \le \frac{a_n}{u_0}$ donc $\sum(u_{n+1}-u_n)$ converge.
\fincorrection
\nocorrection
\nocorrection
\nocorrection
\nocorrection
\correction{004481}
Transformaton d'Abel.
\fincorrection
\correction{004482}
Transformation d'Abel $+$ découpage, $v_n\to0$ lorsque $n\to\infty$.
\fincorrection
\correction{004483}
$|u_n|+|v_n|\le(|u_1|+|v_1|)\prod_{k=1}^{n-1}\Bigl(1+\frac1{k(k+1)}\Bigr)$
et le produit infini est trivialement convergent.
\fincorrection
\correction{004484}
\begin{enumerate}
  \item 
  \item   
  \begin{enumerate}
    \item $1+S_N\le P_N$ n'est plus triviale mais reste vraie par récurrence
    (la différence est une fonction décroissante de~$a_1$).
    \item
  \end{enumerate}
  \item La suite $(P_Ne^{-S_N})$ est positive décroissante donc converge,
    ce qui entraîne la convergence de $(P_N)$.
    On a $P_N \to 0$ ssi $P_Ne^{-S_N} \to 0$ (lorsque $N\to\infty$)
    soit ssi la série de terme général $\ln(1+a_n)-a_n\sim-\frac{a_n^2}{\strut 2}$
    diverge.
  \item 
   \begin{enumerate}
     \item Démontrer l'inégalité en développant les deux membres.
    Sachant que la suite $(P_N)$ est bornée on en déduit qu'elle est de Cauchy
    donc converge.
     \item
   \end{enumerate}
\end{enumerate}
 \fincorrection
\correction{004485}
On a $f(x) = \sum_{n=2}^\infty \frac{f(x^n)}{2^{n-1}}$.
Soit $a\in{[0,1[}$ et $M_a$, $m_a$ le maximum et le minimum de $f$ sur $[0,a]$.
D'après la relation précédente, $m_a \ge m_{a^2}$ et $M_a \le M_{a^2}$ donc en fait
$m_a = m_{a^2}$ et $M_a = M_{a^2}$.

On en déduit $f([0,a]) = f([0,a^2]) = \dots = f([0,a^{2^k}])
= \dots = \{f(0)\}$. Donc $f$ est constante et réciproquement les fonctions
constantes conviennent.
\fincorrection
\correction{004486}
Soit $(p_0,p_1,\dots)$ la suite croissante des nombres premiers
et $S_k = \sum_{P(n)\le k}\frac1n$.
On a $S_k = S_{k-1}\sum_{i=0}^\infty \frac1{p_k^i} = \frac{p_k}{p_k-1}S_{k-1}$,
ce qui prouve que $S_k$ est fini. La série demandée est
$\frac{S_0}{p_0}+ \sum_{k=1}^\infty \frac{S_k-S_{k-1}}{p_k} =
\frac{S_0}{p_0}+ \sum_{k=1}^\infty\frac{S_k}{p_k^2}$.

Montrons que $S_k\le2\sqrt{p_k}$, ceci prouvera la convergence. C'est vrai pour $k=0$ et $k=1$, et si
c'est vrai pour $k-1$ avec $k\ge 2$ alors on obtient $S_k\le 2\sqrt{p_k}\sqrt{\frac{p_kp_{k-1}}{(p_k-1)^2}}
\le 2\sqrt{p_k}\sqrt{\frac{p_k(p_k-2)}{(p_k-1)^2}} \le 2\sqrt{p_k}$.

Remarque~: on a en réalité $S_k\sim e^\gamma\ln(p_k)$ où $\gamma$ est la constante
d'Euler (formule de Mertens).
\fincorrection
\correction{005142}
\begin{enumerate}
\item  Soit $n\geq3$.

$$\sum_{i=3}^{n}i=\frac{(3+n)(n-2)}{2}=\frac{(n-2)(n+3)}{2}.$$

Soit $n\in\Nn^*$.

$$\sum_{i=1}^{n}(2i-1)=\frac{(1+(2n-1))n}{2}=n^2$$

et

$$\sum_{k=4}^{n+1}(3k+7)=\frac{(19+3n+10)(n-2)}{2}=\frac{1}{2}(3n+29)(n-2)=\frac{1}{2}(3n^2+23n-58).$$

\item  Soit $n\in\Nn^*$. Posons $u_n=1,\underbrace{11...1}_n$. On a

$$u_n=1+\sum_{k=1}^{n}\frac{1}{10^k}=1+\frac{1}{10}\frac{1-\frac{1}{10^n}}{1-\frac{1}{10}}
=1+\frac{1}{9}(1-\frac{1}{10^n})=\frac{10}{9}-\frac{1}{9.10^n}.$$

Quand $n$ tend vers $+\infty$, $\frac{1}{9.10^n}$ tend vers $0$, et donc, $u_n$ tend vers $\frac{10}{9}$.
\begin{center}
\shadowbox{
$1,11111....=\frac{10}{9}.$
}
\end{center}

Soit $n\in\Nn^*$. Posons $u_n=0,\underbrace{99...9}_n$. On a

$$u_n=\sum_{k=1}^{n}\frac{9}{10^k}=\frac{9}{10}\frac{1-\frac{1}{10^n}}{1-\frac{1}{10}}
=1-\frac{1}{10^n}.$$

Quand $n$ tend vers $+\infty$, $\frac{1}{10^n}$ tend vers $0$, et donc, $u_n$ tend vers $1$.
\begin{center}
\shadowbox{
$0,9999....=1.$
}
\end{center}

\item  Soit $n\in\Nn^*$. Posons $u_n=\underbrace{1-1+1-...+(-1)^{n-1}}_n$.  On a

$$u_n=\sum_{k=0}^{n-1}(-1)^k=\frac{1-(-1)^n}{1-(-1)}=\frac{1}{2}(1-(-1)^n)=\left\{
\begin{array}{l}
0\;\mbox{si}\;n\;\mbox{est pair}\\
1\;\mbox{si}\;n\;\mbox{est impair}
\end{array}
\right..$$

\item  Soit $n\in\Nn^*$.
$\sum_{k=1}^{n}\frac{1}{2^k}=\frac{1}{2}\frac{1-\frac{1}{2^n}}{1-\frac{1}{2}}=1-\frac{1}{2^n}$. Quand $n$ tend vers
$+\infty$, on obtient
\begin{center}
\shadowbox{
$\frac{1}{2}+\frac{1}{4}+\frac{1}{8}+...=1.$
}
\end{center}
\item  Soit $n\in\Nn$.

\begin{align*}
\sum_{k=0}^{n}\cos\frac{k\pi}{2}&=\Re(\sum_{k=0}^{n}e^{ik\pi/2})(=\Re(\sum_{k=0}^{n}i^k))\\
 &=\Re(\frac{1-e^{(n+1)i\pi/2}}{1-e^{i\pi/2}})=\Re(\frac{e^{i(n+1)\pi/4}}{e^{i\pi/4}}\frac{-2i\sin\frac{(n+1)\pi}{4}}
 {-2i\sin\frac{\pi}{4}})=\sqrt{2}\sin\frac{(n+1)\pi}{4}\cos\frac{n\pi}{4}\\
 &=\frac{1}{\sqrt{2}}\sin\frac{(2n+1)\pi}{4}+\frac{1}{2}
=\left\{
\begin{array}{l}
1\;\mbox{si}\;n\in4\Nn\cup(4\Nn+1)\\
0\;\mbox{si}\;n\in(4\Nn+2)\cup(4\Nn+3)
\end{array}
\right.
\end{align*}

En fait, on peut constater beaucoup plus simplement que $\cos0+\cos\frac{\pi}{2}+\cos\pi+\cos\frac{3\pi}{2}=1+0-1+0=0$,
on a immédiatement $S_{4n}=1$, $S_{4n+1}=S_{4n}+0=1$, $S_{4n+2}=S_{4n+1}-1=0$ et $S_{4n+3}=S_{4n+2}+0=0$.
\item  Soient $n\in\Nn$ et $\theta\in\Rr$. Posons $C_n=\sum_{k=0}^{n}\cos(k\theta)$ et
$S_n=\sum_{k=0}^{n}\sin(k\theta)$. Alors, d'après la formule de \textsc{Moivre},

$$C_n+iS_n=\sum_{k=0}^{n}(\cos(k\theta)+i\sin(k\theta))=\sum_{k=0}^{n}e^{ik\theta}=\sum_{k=0}^{n}(e^{i\theta})^k.$$

\begin{itemize}
\item[\textbf{- 1er cas.}] Si $\theta\notin2\pi\Zz$, alors $e^{i\theta}\neq1$. Par suite,

$$C_n+iS_n=\frac{1-e^{i(n+1)\theta}}{1-e^{i\theta}}=e^{i\theta(n+1-1)/2}\frac{-2i\sin\frac{(n+1)\theta}{2}}{-2i\sin
\frac{\theta}{2}}=e^{in\theta/2}\frac{\sin\frac{(n+1)\theta}{2}}{\sin
\frac{\theta}{2}}.$$

Par suite,

$$C_n=\Re(C_n+iS_n)=\frac{\cos\frac{n\theta}{2}\sin\frac{(n+1)\theta}{2}}{\sin\frac{\theta}{2}}\;\mbox{et}\;S_n=\Im(C_
n+iS_n)=\frac{\sin\frac{n\theta}{2}\sin\frac{(n+1)\theta}{2}}{\sin\frac{\theta}{2}}.$$

\item[\textbf{- 2ème cas.}] Si $\theta\in2\pi\Zz$, on a immédiatement $C_n=n+1$ et $S_n=0$.
\end{itemize}

Finalement,
\begin{center}
\shadowbox{
$\forall n\in\Nn,\;\sum_{k=0}^{n}\cos(k\theta)=
\left\{
\begin{array}{l}
\frac{\cos\frac{n\theta}{2}\sin\frac{(n+1)\theta}{2}}{\sin\frac{\theta}{2}}\;\mbox{si}\;\theta\notin2\pi\Zz\\
n+1\;\mbox{si}\;\theta\in2\pi\Zz
\end{array}
\right.
\;\text{et}\;\sum_{k=0}^{n}\sin(k\theta)=
\left\{
\begin{array}{l}
\frac{\sin\frac{n\theta}{2}\sin\frac{(n+1)\theta}{2}}{\sin\frac{\theta}{2}}\;\mbox{si}\;\theta\notin2\pi\Zz\\
0\;\mbox{si}\;\theta\in2\pi\Zz
\end{array}
\right.
.$
}
\end{center}

\item  Soient $x\in[0,1]$ et $n\in\Nn^*$. Puisque $-x\neq1$, on a

\begin{align*}
S_n'(x)&=\sum_{k=1}^{n}(-1)^{k-1}x^{k-1}=\sum_{k=0}^{n-1}(-x)^{k}=\frac{1-(-x)^{n}}{1-(-x)}=\frac{1}{1+x}(1-(-x)^{n}).
\end{align*}

Par suite,

$$S_n(x)=S_n(0)+\int_{0}^{x}S_n'(t)\;dt=\int_{0}^{x}\frac{1-(-t)^{n}}{1+t}\;dt=\int_{0}^{x}\frac{1}{1+t}\;dt-
\int_{0}^{x}\frac{(-t)^{n}}{1+t}\;dt=\ln(1+x)-\int_{0}^{x}\frac{(-t)^{n}}{1+t}\;dt.$$

Mais alors,

$$|S_n(x)-\ln(1+x)|=\left|\int_{0}^{x}\frac{(-t)^{n}}{1+t}\;dt\right|\leq\int_{0}^{x}
\left|\frac{(-t)^{n}}{1+t}\right|\;dt=\int_{0}^{x}\frac{t^n}{1+t}\;dt
\leq\int_{0}^{x}t^{n}dt=\frac{x^{n+1}}{n+1}\leq\frac{1}{n+1}.$$

Comme $\frac{1}{n+1}$ tend vers $0$ quand $n$ tend vers $+\infty$, on en déduit que
\begin{center}
\shadowbox{
$\forall x\in[0,1],\;\lim_{n\rightarrow +\infty}\sum_{k=1}^{n}(-1)^{k-1}x^{k-1}=\ln(1+x).$
}
\end{center}
En particulier, 
\begin{center}
\shadowbox{$\ln2=\lim_{n\rightarrow +\infty}(1-\frac{1}{2}+\frac{1}{3}-...+\frac{(-1)^{n-1}}{n})$}.
\end{center}

\item 

\begin{enumerate}
\item Soit $n\in\Nn$. $u_{n+1}-3=2u_n-6=2(u_n-3)$. La suite $(u_n-3)_{n\in\Nn}$ est donc une suite géométrique, de
raion $q=2$ et de premier terme $u_0-3=-2$. On en déduit que, pour $n$ enteir naturel donné, $u_n-3=-2.2^n$. Donc,

$$\forall n\in\Nn,\;u_n=3-2^{n+1}.$$

\item Soit $n\in\Nn$.

$$\sum_{k=0}^{n}u_k=\sum_{k=0}^{n}3-2\sum_{k=0}^{n}2^k=3(n+1)-2\frac{2^{n+1}-1}{2-1}=-2^{n+2}+3n+5.$$
\end{enumerate}
\end{enumerate}

\fincorrection
\correction{005150}
Pour $x$ réel, posons $f(x)=\sum_{k=1}^{n}(a_k+b_kx)^2$. On remarque que pour tout réel $x$,
$f(x)\geq0$. En développant les $n$ carrés, on obtient,
$$f(x)=\sum_{k=1}^{n}(b_k^2x^2+2a_kb_kx+a_k^2)=(\sum_{k=1}^{n}b_k^2)x^2+2(\sum_{k=1}^{n}a_kb_k)x+(\sum_{k=1}^{n}a_k^2).$$

\begin{itemize}
\item[\textbf{1er cas.}] Si $\sum_{k=1}^{n}b_k^2\neq0$, $f$ est un trinôme du second degré de signe constant sur $\Rr$.
Son discriminant réduit est alors négatif ou nul. Ceci fournit

$$0\geq\Delta'=(\sum_{k=1}^{n}a_kb_k)^2-(\sum_{k=1}^{n}b_k^2)(\sum_{k=1}^{n}a_k^2),$$

et donc

$$\left|\sum_{k=1}^{n}a_kb_k\right|\leq\sqrt{\sum_{k=1}^{n}a_k^2}\sqrt{\sum_{k=1}^{n}b_k^2}.$$

\item[\textbf{2ème cas.}] Si $\sum_{k=1}^{n}b_k^2=0$, alors tous les $b_k$ sont nuls et l'inégalité est immédiate.\\
Finalement, dans tous les cas,
\begin{center}
\shadowbox{
$\left|\sum_{k=1}^{n}a_kb_k\right|\leq\sqrt{\sum_{k=1}^{n}a_k^2}\sqrt{\sum_{k=1}^{n}b_k^2}.$
}
\end{center}
\end{itemize}

Cette inégalité est encore valable en remplaçant les $a_k$ et les $b_k$ par leurs valeurs absolues, ce qui fournit
les inégalités intermédiaires.

Retrouvons alors l'inégalité de l'exercice \ref{exo:suprou4bis}. Puisque les $a_k$ sont strictement positifs, on peut écrire~:

$$\left(\sum_{i=1}^{n}a_i\right)\left(\sum_{i=1}^{n}\frac{1}{a_i}\right)=\left(\sum_{i=1}^{n}\sqrt{a_i}^2\right)\left(
\sum_{i=1}^{n}\sqrt{\frac{1}{a_i}}^2\right)\geq\left(\sum_{i=1}^{n}\sqrt{a_i}\sqrt{\frac{1}{a_i}}\right)^2=n^2.$$
\fincorrection
\correction{005458}
Soit $n\in\Nn^*$.

\begin{align*}\ensuremath
\sum_{k=1}^{n}\sin\frac{k}{n^2}&=\mbox{Im}(\sum_{k=1}^{n}e^{ik/n^2})=\mbox{Im}\left(e^{i/n^2}\frac{1-e^{ni/n^2}}{1-e^{i/n^2}}\right)
=\mbox{Im}\left(e^{i(1+\frac{n}{2}-\frac{1}{2})/n^2}\frac{\sin\frac{1}{2n}}{\sin\frac{1}{2n^2}}\right)
=\frac{\sin\frac{n+1}{2n^2}\sin\frac{1}{2n}}{\sin\frac{1}{2n^2}}\\
 &\underset{n\rightarrow+\infty}{=}(\frac{1}{2n}+\frac{1}{2n^2}+o(\frac{1}{n^2}))(\frac{1}{2n}+o(\frac{1}{n^2}))(\frac{1}{2n^2}+o(\frac{1}{n^3}))^{-1}\\
 &=(1+\frac{1}{n}+o(\frac{1}{n}))(\frac{1}{2}+o(\frac{1}{n}))(1+o(\frac{1}{n}))^{-1}=\frac{1}{2}+\frac{1}{2n}
 +o(\frac{1}{n}),
\end{align*}

(on peut aussi partir de l'encadrement $\frac{k}{n^2}-\frac{k^3}{6n^6}\leq\sin\frac{k}{n^2}\leq\frac{k}{n^2}$).
\fincorrection
\correction{005688}
\begin{enumerate}
 \item  Pour $n\geqslant1$, on pose $u_n =\ln\left(\frac{n^2+n+1}{n^2+n-1}\right)$. $\forall n\geqslant 1$, $u_n$ existe

\begin{center} 
$u_n=\ln\left(1+\frac{1}{n}+\frac{1}{n^2}\right)-\ln\left(1+\frac{1}{n}-\frac{1}{n^2}\right)\underset{n\rightarrow+\infty}{=}\left(\frac{1}{n}+O\left(\frac{1}{n^2}\right)\right)-\left(\frac{1}{n}+O\left(\frac{1}{n^2}\right)\right)=O\left(\frac{1}{n^2}\right)$.
\end{center}

Comme la série de terme général $\frac{1}{n^2}$, $n\geqslant1$, converge (série de \textsc{Riemann} d'exposant $\alpha>1$), la série de terme général $u_n$ converge.

\item  Pour $n\geqslant2$, on pose $u_n =\frac{1}{n+(-1)^n\sqrt{n}}$. $\forall n\geqslant 2$, $u_n$ existe et de plus $u_n\underset{n\rightarrow+\infty}{\sim}\frac{1}{n}$. Comme la série de terme général $\frac{1}{n}$, $n\geqslant 2$, diverge et est positive, la série de terme général $u_n$ diverge.

\item  Pour $n\geqslant1$, on pose $u_n =\left(\frac{n+3}{2n+1}\right)^{\ln n}$. Pour $n\geqslant1$, $u_n > 0$ et 

\begin{align*}\ensuremath
\ln(u_n)&=\ln(n)\ln\left(\frac{n+3}{2n+1}\right) =\ln(n)\left(\ln\left(\frac{1}{2}\right)+\ln\left(1+\frac{3}{n}\right) -\ln\left(1+\frac{1}{2n}\right)\right)\\
 &\underset{n\rightarrow+\infty}{=}\ln(n)\left(-\ln2+O\left(\frac{1}{n}\right)\right)\underset{n\rightarrow+\infty}{=}-\ln2\ln(n)+o(1).
\end{align*}

Donc $u_n=e^{\ln(u_n)}\underset{n\rightarrow+\infty}{\sim}e^{-\ln2\ln n}=\frac{1}{n^{\ln2}}$.  Comme la série de terme général $\frac{1}{n^{\ln 2}}$, $n\geqslant1$, diverge (série de \textsc{Riemann} d'exposant $\alpha\leqslant1$) et est positive, la série de terme général $u_n$ diverge.

\item  Pour $n\geqslant2$, on pose $u_n=\frac{1}{\ln(n)\ln(\ch n)}$. $u_n$ existe pour $n\geqslant2$. $\ln(\ch n)\underset{n\rightarrow+\infty}{\sim}\ln\left(\frac{e^n}{2}\right)=n -\ln2\underset{n\rightarrow+\infty}{\sim}n$ et $un\underset{n\rightarrow+\infty}{\sim}\frac{1}{n\ln(n)}>0$.

Vérifions alors que la série de terme général $\frac{1}{n\ln n}$, $n\geqslant2$, diverge. La fonction $x\rightarrow x\ln x$ est continue, croissante et strictement positive sur $]1,+\infty[$ (produit de deux fonctions strictement positives et croissantes sur $]1,+\infty[$). Par suite, la fonction $x\rightarrow\frac{1}{x\ln x}$ est continue et décroissante sur $]1,+\infty[$ et pour tout entier $k$ supérieur ou égal à $2$, 

\begin{center}
$\frac{1}{k\ln k}\geqslant\int_{k}^{k+1}\frac{1}{x\ln x}\;dx$
\end{center}

Par suite, pour $n\geqslant2$, 

\begin{center}
$\sum_{k=2}^{n}\frac{k\ln k}\geqslant\sum_{k=2}^{n}\int_{k}^{k+1}\frac{1}{x\ln x}\;dx=\int_{2}^{n+1}\frac{1}{x\ln x}\;dx=\ln(\ln(n+1)) -\ln(\ln(2)\underset{n\rightarrow+\infty}{\rightarrow}+\infty.$
\end{center}

Donc $u_n$ est positif et équivalent au terme général d'une série divergente. La série de terme général $u_n$ diverge.

\item  Pour $n\geqslant1$, on pose $u_n=\Arccos\sqrt[3]{1-\frac{1}{n^2}}$. $u_n$ existe pour $n\geqslant 1$. De plus $u_n\underset{n\rightarrow+\infty}{\rightarrow}0$. On en déduit que 

\begin{align*}\ensuremath
u_n&\underset{n\rightarrow+\infty}{\sim}\sin(u_n)=\sin\left(\Arccos\sqrt[3]{1-\frac{1}{n^2}}\right) =\sqrt{1-\left(1-\frac{1}{n^2}\right)^{2/3}}\underset{n\rightarrow+\infty}{=}\sqrt{1-1+\frac{2}{3n^2}+o\left(\frac{1}{n^2}\right)}\\
 &\underset{n\rightarrow+\infty}{\sim}\sqrt{\frac{2}{3}}\frac{1}{n}>0
\end{align*}

terme général d'une série de \textsc{Riemann} divergente. La série de terme général un diverge.

\item  Pour $n\geqslant1$, on pose $u_n=\frac{n^2}{(n-1)!}$. $u_n$ existe  et $u_n \neq0$ pour $n\geqslant1$. De plus,

\begin{center}
$\left|\frac{u_{n+1}}{u_n}\right|=\frac{(n+1)^2}{n^2}\times\frac{(n-1)!}{n!}=\frac{(n+1)^2}{n^3}  \underset{n\rightarrow+\infty}{\sim}\frac{1}{n}\underset{n\rightarrow+\infty}{\rightarrow}0< 1$.
\end{center}

D'après la règle de d'\textsc{Alembert}, la série de terme général $u_n$ converge.

\item  Pour $n\geqslant1$, on pose $u_n=\left(\cos\frac{1}{\sqrt{n}}\right)^n-\frac{1}{\sqrt{e}}$. $u_n$ est défini pour $n\geqslant1$ car pour $n\geqslant1$, $\frac{1}{\sqrt{n}}\in\left]0,\frac{\pi}{2}\right[$ et donc $\cos\frac{1}{\sqrt{n}}>0$. Ensuite

\begin{align*}\ensuremath
\ln\left(\cos\frac{1}{\sqrt{n}}\right)&\underset{n\rightarrow+\infty}{=}\ln\left(1-\frac{1}{2n}+\frac{1}{24n^2}+o\left(\frac{1}{n^2}\right)\right)\underset{n\rightarrow+\infty}{=}-\frac{1}{2n}+\frac{1}{24n^2}-\frac{1}{8n^2}+o\left(\frac{1}{n^2}\right)\\
 &\underset{n\rightarrow+\infty}{=}-\frac{1}{2n}-\frac{1}{12n^2}+o\left(\frac{1}{n^2}\right).
\end{align*}

Puis $n\ln\left(\cos\frac{1}{\sqrt{n}}\right)\underset{n\rightarrow+\infty}{=}-\frac{1}{2}-\frac{1}{12n}+o\left(\frac{1}{n}\right)$ et donc

\begin{center}
$u_n=e^{n\ln(\cos(1/\sqrt{n})}-\frac{1}{\sqrt{e}}\underset{n\rightarrow+\infty}{=}\frac{1}{\sqrt{e}}\left(e^{-\frac{1}{12n}+o\left(\frac{1}{n}\right)}-1\right)\underset{n\rightarrow+\infty}{\sim}-\frac{1}{12n\sqrt{e}}<0$.
\end{center}

La série de terme général $-\frac{1}{12n\sqrt{e}}$ est  divergente et donc la série de terme général $u_n$ diverge.

\item 

\begin{align*}\ensuremath
\ln\left(\frac{2}{\pi}\Arctan\left(\frac{n^2+1}{n}\right)\right)&=\ln\left(1-\frac{2}{\pi}\Arctan\left(\frac{n}{n^2+1}\right)\right)\\
 &\underset{n\rightarrow+\infty}{\sim}-\frac{2}{\pi}\Arctan\left(\frac{n}{n^2+1}\right)\underset{n\rightarrow+\infty}{\sim}-\frac{2}{\pi}\frac{n}{n^2+1}\underset{n\rightarrow+\infty}{\sim}-\frac{2}{n\pi}<0.
\end{align*}

Donc, la série de terme général $u_n$ diverge.

\item  Pour $n\geqslant1$, on pose $u_n=\int_{0}^{\pi/2}\frac{\cos^2x}{n^2+\cos^2x}\;dx$.

Pour $n\geqslant1$, la fonction $x\mapsto\frac{\cos^2x}{n^2+\cos^2x}\;dx$ est continue sur $\left[0,\frac{\pi}{2}\right]$ et positive et donc, $u_n$ existe et est positif. De plus, pour $n\geqslant1$,

\begin{center}
$0\leqslant u_n\leqslant\int_{0}^{\pi/2}\frac{1}{n^2+0}\;dx=\frac{\pi}{2n^2}$.
\end{center}

La série de terme général $\frac{\pi}{2n^2}$ converge et donc la série de terme général $u_n$ converge.

\item  $-\sqrt{2}\sin\left(\frac{\pi}{4}+\frac{1}{n}\right) =-\sin\left(\frac{1}{n}\right)-\cos\left(\frac{1}{n}\right)\underset{n\rightarrow+\infty}{=}-1+O\left(\frac{1}{n}\right)$ puis

\begin{center}
$-\sqrt{2}\sin\left(\frac{\pi}{4}+\frac{1}{n}\right)\ln n\underset{n\rightarrow+\infty}{=}-\ln(n)+O\left(\frac{\ln n}{n}\right)\underset{n\rightarrow+\infty}{=}-\ln(n)+o(1)$.
\end{center}

Par suite,

\begin{center}
$0< u_n=e^{-\sqrt{2}\sin\left(\frac{\pi}{4}+\frac{1}{n}\right)\ln n}\underset{n\rightarrow+\infty}{\sim}e^{-\ln n}=\frac{1}{n}$.
\end{center}

La série de terme général $\frac{1}{n}$ diverge et la série de terme général $u_n$ diverge.

\item  $n\ln\left(1+\frac{1}{n}\right)\underset{n\rightarrow+\infty}{=}1-\frac{1}{2n}+o\left(\frac{1}{n}\right)$ et donc

\begin{center}
$u_n\underset{n\rightarrow+\infty}{=}e-e^{1-\frac{1}{2n}+o\left(\frac{1}{n}\right)}\underset{n\rightarrow+\infty}{=}e\left(1-1+\frac{1}{2n}+o\left(\frac{1}{n}\right)\right)\underset{n\rightarrow+\infty}{\sim}\frac{e}{2n}>0$.
\end{center}

La série de terme général $\frac{e}{2n}$ diverge et la série de terme général $u_n$ diverge.
\end{enumerate}
\fincorrection
\correction{005689}
\begin{enumerate}
 \item  Si $P$ n'est pas unitaire de degré $3$, $u_n$ ne tend pas vers $0$ et la série de terme général $u_n$ diverge grossièrement.

Soit $P$ un polynôme unitaire de degré $3$. Posons $P =X^3+aX^2+bX+c$.

\begin{align*}\ensuremath
u_n&=n\left(\left(1+\frac{2}{n^2}\right)^{1/4}-\left(1+\frac{a}{n}+\frac{b}{n^2}+\frac{c}{n^3}\right)^{1/3}\right)\\
 &\underset{n\rightarrow+\infty}{=}n\left(\left(1+\frac{1}{2n^2}+O\left(\frac{1}{n^3}\right)\right)-\left(1+\frac{a}{3n}+\frac{b}{3n^2}-\frac{a^2}{9n^2}+O\left(\frac{1}{n^3}\right)\right)\right)\\
 &\underset{n\rightarrow+\infty}{=}-\frac{a}{3}+\left(\frac{1}{2}-\frac{b}{3}+\frac{a^2}{9}\right)\frac{1}{n}+O\left(\frac{1}{n^2}\right).
\end{align*}

\textbullet~Si $a\neq0$, $u_n$ ne tend pas vers $0$ et la série de terme général $u_n$ diverge grossièrement.

\textbullet~Si $a=0$ et $\frac{1}{2}-\frac{b}{3}\neq0$, $u_n\underset{n\rightarrow+\infty}{\sim}\left(\frac{1}{2}-\frac{b}{3}\right)\frac{1}{n}$. $u_n$ est donc de signe constant pour $n$ grand et est équivalent au terme général d'une série divergente. Donc la série de terme général $u_n$ diverge.

\textbullet~Si $a = 0$ et $\frac{1}{2}-\frac{b}{3}= 0$, $u_n\underset{n\rightarrow+\infty}{=}O\left(\frac{1}{n^2}\right)$. Dans ce cas, la série de terme général $u_n$ converge (absolument).

En résumé, la série de terme général $u_n$ converge si et seulement si $a = 0$ et  $b=\frac{3}{2}$ ou encore la série de terme général $u_n$ converge si et seulement si $P$ est de la forme $X^3+\frac{3}{2}X+c$, $c\in\Rr$.

\item  Pour $n\geqslant2$, posons $u_n=\frac{1}{n^\alpha}S(n)$. Pour $n\geqslant2$,

\begin{center}
$0<S(n+1)=\sum_{p=2}^{+\infty}\frac{1}{p}\times\frac{1}{p^n}\leqslant\frac{1}{2}\sum_{p=2}^{+\infty}\frac{1}{p^n}=\frac{1}{2}S(n)$
\end{center}

et donc $\forall n\geqslant2$, $S(n)\leqslant\frac{S(2)}{2^{n-2}}$. Par suite,

\begin{center}
$u_n\leqslant\frac{1}{n^\alpha}\frac{S(2)}{2^{n-2}}\underset{n\rightarrow+\infty}{=}o\left(\frac{1}{n^2}\right)$.
\end{center}

Pour tout réel $\alpha$, la série de terme général $u_n$ converge.

\item  $\forall u_0\in\Rr$, $\forall n\in\Nn^*$, $u_n > 0$. Par suite, $\forall n\geqslant 2$, $0< u_n<\frac{1}{n}$.

On en déduit que $\lim_{n \rightarrow +\infty}u_n=0$ et par suite $u_n\underset{n\rightarrow+\infty}{\sim}\frac{1}{n}>0$. La série de terme général $u_n$ diverge.

\item  On sait qu'il existe une infinité de nombres premiers.

Notons $(p_n)_{n\in\Nn^*}$ la suite croissante des nombres premiers. La suite $(p_n)_{n\in\Nn^*}$ est une suite strictement croissante d'entiers et donc $\lim_{n \rightarrow +\infty}p_n= +\infty$ ou encore $\lim_{n \rightarrow +\infty}\frac{1}{p_n}=0$.

Par suite, $0<\frac{1}{p_n}\underset{n\rightarrow+\infty}{\sim}\ln\left(\left(1-\frac{1}{p_n}\right)^{-1}\right)$ et les séries de termes généraux $\frac{1}{p_n}$ et $\ln\left(\left(1-\frac{1}{p_n}\right)^{-1}\right)$ sont de même nature.

Il reste donc à étudier la nature de la série de terme général $\ln\left(\left(1-\frac{1}{p_n}\right)^{-1}\right)$.

Montrons que $\forall N\in\Nn^*$,  $\sum_{n=1}^{+\infty}\ln\left(\left(1-\frac{1}{p_n}\right)^{-1}\right)\geqslant\ln\left(\sum_{k=1}^{N}\frac{1}{k}\right)$.

Soit $n\geqslant$. Alors $\frac{1}{p_n}<1$ et la série de terme général $\frac{1}{p_n^k}$, $k\in\Nn$, est une série géométrique convergente de somme : $\sum_{k=0}^{+\infty}\frac{1}{p_n^k}= \left(1-\frac{1}{p_n}\right)^{-1}$.

Soit alors $N$ un entier naturel supérieur ou égal à $2$ et $p_1 < p_2... < p_n$ la liste des nombres premiers inférieurs ou égaux à $N$.

Tout entier entre $1$ et $N$ s'écrit de manière unique $p_1^{\beta_1}\ldots p_k^{\beta_k}$ où $\forall i\in\llbracket1,n\rrbracket$, $0\leqslant\beta_i\leqslant\alpha_i=E\left(\frac{\ln(N)}{\ln(p_i)}\right)$ et deux entiers distincts ont des décompositions distinctes. Donc

\begin{align*}\ensuremath
\sum_{k=1}^{+\infty}\ln\left(\left(1-\frac{1}{p_k}\right)^{-1}\right)&\geqslant\sum_{k=1}^{n}\ln\left(\left(1-\frac{1}{p_k}\right)^{-1}\right)\;(\text{car}\;\forall k\in\Nn^*,\;\left(1-\frac{1}{p_k}\right)^{-1}> 1)\\
 &=\sum_{k=1}^{n}\ln\left(\sum_{i=0}^{+\infty}\frac{1}{p_k^i}\right)\geqslant\sum_{k=1}^{n}\ln\left(\sum_{i=0}^{\alpha_k}\frac{1}{p_k^i}\right)\\
  &=\ln\left(\prod_{k=1}^{n}\left(\sum_{i=0}^{\alpha_k}\frac{1}{p_k^i}\right)\right)=\ln\left(\sum_{0\leqslant\beta_1\leqslant\alpha_1,\ldots,\ldots0\leqslant\beta_n\leqslant\alpha_n}^{}\frac{1}{p_1^{\beta_1}\ldots,\;p_n^{\beta_n}}\right)\\
   &\geqslant\ln\left(\sum_{k=1}^{N}\frac{1}{k}\right).
\end{align*}

Or $\lim_{N \rightarrow +\infty}\ln\left(\sum_{k=1}^{N}\frac{1}{k}\right)=+\infty$ et donc $\sum_{k=1}^{+\infty}\ln\left(\left(1-\frac{1}{p_k}\right)^{-1}\right)=+\infty$.

La série de terme général $\ln\left(1-\frac{1}{p_k}\right)^{-1}$ diverge et il en est de même de la série de terme général $\frac{1}{p_n}$.

(Ceci montre qu'il y a beaucoup de nombres premiers et en tout cas beaucoup plus de nombres premiers que de carrés parfaits par exemple).

\item  Soit $n\in\Nn^*$. Posons $n =a_p\times10^p+\ldots+a_1\times10+a_0$ où $\forall i\in\llbracket0,p\rrbracket$, $a_i\in\{0,1;...,9\}$ et $a_p\neq 0$. Alors $c(n) = p+1$.

Déterminons $p$ est  en fonction de $n$. On a $10^p\leqslant n <10^{p+1}$ et donc $p=E\left(\log(n)\right)$. Donc

\begin{center}
$\forall n\in\Nn^*$, $u_n=\frac{1}{n(E(\log n)+1)^\alpha}$.
\end{center}

Par suite, $u_n\underset{n\rightarrow+\infty}{\sim}\frac{\ln^\alpha(10)}{n\ln^\alpha(n)}$ et la série de terme général $u_n$ converge si et seulement si $\alpha> 1$ (séries de \textsc{Bertrand}). Redémontrons ce résultat qui n'est pas un résultat de cours.

La série de terme général $\frac{1}{n\ln n}$ est divergente (voir l'exercice \ref{ex:rou1ter}, 4)). Par suite, si $\alpha\leqslant1$, la série de terme général $\frac{1}{n\ln^\alpha(n)}$ est divergente car $\forall n\geqslant2$, $\frac{1}{n\ln^\alpha(n)}\geqslant\frac{1}{n\ln n}$.

Soit $\alpha>1$. Puisque la fonction $x\mapsto\frac{1}{x\ln^\alpha x}$ est continue et strictement décroissante sur $]1,+\infty[$, pour $k\geqslant3$,

\begin{center}
$\frac{1}{k\ln^\alpha k}\leqslant\int_{k-1}^{k}\frac{1}{x\ln^\alpha x}\;dx$
\end{center}

puis, pour $n\geqslant3$, en sommant pour $k\in\llbracket3,n\rrbracket$

\begin{center}
$\sum_{k=3}^{n}\frac{1}{k\ln^\alpha k}\leqslant\sum_{k=3}^{n}\int_{k-1}^{k}\frac{1}{x\ln^\alpha x}\;dx=\int_{2}^{n}\frac{1}{x\ln^\alpha x}\;dx=\frac{1}{\alpha-1}\left(\frac{1}{\ln^{\alpha-1}(2)}-\frac{1}{\ln^{\alpha-1}(n)}\right)\leqslant\frac{1}{\alpha-1}\frac{1}{\ln^{\alpha-1}(2)}$.
\end{center}

Ainsi, la suite des sommes partielles de la série à termes positifs, de terme général $\frac{1}{k\ln^\alpha k}$, est majorée et donc la série de terme général $\frac{1}{k\ln^\alpha k}$ converge.

\textbf{6} Soit $n\geqslant2$.

\begin{center}
$\left|\frac{u_{n+1}}{u_n}\right|=\frac{\ln^a(n+1)}{(n+1)^b}\underset{n\rightarrow+\infty}{\rightarrow}0 < 1$
\end{center}

et d'après la règle de d'\textsc{Alembert}, la série de terme général $u_n$ converge.

\item  $\lim_{n \rightarrow +\infty}u_n=\frac{\pi}{4}-\frac{\pi}{4}= 0$. Donc

\begin{align*}\ensuremath
u_n&\underset{n\rightarrow+\infty}{\sim}\tan(u_n)\\
 &=\frac{\left(1+\frac{1}{n}\right)^a-\left(1-\frac{1}{n}\right)^a}{1+\left(1-\frac{1}{n^2}\right)^a}\underset{n\rightarrow+\infty}{=}\frac{\frac{2a}{n}+O\left(\frac{1}{n^2}\right)}{2+O\left(\frac{1}{n^2}\right)}\underset{n\rightarrow+\infty}{=}\frac{a}{n}+O\left(\frac{1}{n^2}\right).
\end{align*}

Par suite, la série de terme général $u_n$ converge si et seulement si $a = 0$.

\item  La fonction $x\mapsto x^{3/2}$ est continue et croissante sur $\Rr^+$. Donc pour $k\geqslant 1$, $\int_{k-1}^{k}x^{3/2}\;dx\leqslant k^{3/2}\leqslant\int_{k}^{k+1}x^{3/2}\;dx$ puis pour $n\in\Nn^*$ :

\begin{center}
$\int_{0}^{n}x^{3/2}\;dx\sum_{k=1}^{n}\int_{k-1}^{k}x^{3/2}\;dx\leqslant\sum_{k=1}^{n}k^{3/2}\leqslant\sum_{k=1}^{n}\int_{k}^{k+1}x^{3/2}\;dx=\int_{1}^{n+1}x^{3/2}\;dx$
\end{center}

ce qui fournit

\begin{center}
$\frac{2}{5}n^{5/2}\leqslant\sum_{k=1}^{n}k^{3/2}\leqslant\frac{2}{5}((n+1)^{5/2}-1)$ et donc $\sum_{k=1}^{n}k^{3/2}\underset{n\rightarrow+\infty}{\sim}\frac{2n^{5/2}}{5}$.
\end{center}

Donc $u_n\underset{n\rightarrow+\infty}{\sim}\frac{2n^{\frac{5}{2}-\alpha}}{5}>0$. La série de terme général $u_n$ converge si et seulement si $\alpha>\frac{7}{2}$.

\item  Pour $n\geqslant1$,

\begin{center}
$u_n =\left(1+\frac{1}{n^\alpha}\right)\left(1+\frac{2}{n^\alpha}\right)\ldots\left(1+\frac{n}{n^\alpha}\right)-1\geqslant\frac{1}{n^\alpha}+\frac{2}{n^\alpha}+\ldots+\frac{n}{n^\alpha}=\frac{n(n+1)}{2n^\alpha}>0$.
\end{center}

Comme $\frac{n(n+1)}{2n^\alpha}\underset{n\rightarrow+\infty}{\sim}\frac{1}{2n^{\alpha-2}}$, si $\alpha\leqslant3$, on a $\alpha-2\leqslant1$ et la série de terme général $u_n$ diverge.

Si $\alpha> 3$,

\begin{align*}\ensuremath
0 < u_n&\leqslant\left(1+\frac{n}{n^\alpha}\right)^n -1= e^{n\ln\left(1+\frac{1}{n^{\alpha-1}}\right)}-1\\
 &\underset{n\rightarrow+\infty}{\sim}n\ln\left(1+\frac{1}{n^{\alpha-1}}\right)\\
 &\underset{n\rightarrow+\infty}{\sim}\frac{1}{n^{\alpha-2}}\;\text{terme général d'une série de \textsc{Riemann} convergente},
\end{align*}
			   
			   
et, puisque $\alpha-2>1$,  la série de terme général $u_n$ converge. Finalement, la série de terme général $u_n$ converge si et seulement si $\alpha > 3$.
\end{enumerate}
\fincorrection
\correction{005690}
\begin{enumerate}
 \item  Pour $n\in\Nn$,

\begin{center}
$u_n =\sin\left(\frac{\pi n^2}{n+1}\right)=\sin\left(\frac{\pi(n^2-1+1)}{n+1}\right)=\sin\left(\frac{\pi}{n+1}+(n-1)\pi\right)=(-1)^{n-1}\sin\left(\frac{\pi}{n+1}\right)$.
\end{center}

La suite $\left((-1)^{n-1}\sin\left(\frac{\pi}{n+1}\right)\right)_{n\in\Nn}$ est alternée en signe et sa valeur absolue tend vers 0 en décroissant. La série de terme général $u_n$ converge donc en vertu du critère spécial aux séries alternées.

\item  (la suite $\left(\frac{1}{n+(-1)^{n-1}}\right)_{n\in\Nn}$ n'est pas décroisante à partir d'un certain rang).

\begin{center}
$u_n=\frac{(-1)^n}{n}\frac{1}{1+\frac{(-1)^{n-1}}{n}}\underset{n\rightarrow+\infty}{=}\frac{(-1)^n}{n}\left(1+O\left(\frac{1}{n}\right)\right)\underset{n\rightarrow+\infty}{=}\frac{(-1)^n}{n}+O\left(\frac{1}{n^2}\right)$.
\end{center}

La série de terme général $\frac{(-1)^n}{n}$ converge en vertu du critère spécial aux séries alternées et la série de terme général $O\left(\frac{1}{n^2}\right)$ est absolument convergente. On en déduit que la série de terme général $u_n$ converge.

\item  $u_n=\ln\left(1+\frac{(-1)^n}{\sqrt{n}}\right)\underset{n\rightarrow+\infty}{=}\frac{(-1)^n}{\sqrt{n}}-\frac{1}{2n}+ O\left(\frac{1}{n^{3/2}}\right)$. Les séries de termes généraux respectifs $\frac{(-1)^n}{\sqrt{n}}$  et $O\left(\frac{1}{n^{3/2}}\right)$ sont convergentes et la série de terme général $-\frac{1}{2n}$  est divergente. Si la série de terme général $u_n$ convergeait alors la série de terme général $-\frac{1}{2n}=u_n-\frac{(-1)^n}{\sqrt{n}}-O\left(\frac{1}{n^{3/2}}\right)$ convergerait ce qui n'est pas. Donc la série de terme général $u_n$ diverge.

\textbf{Remarque.} La série de terme général $u_n$ diverge bien que $u_n$ soit équivalent au terme général d'une série convergente.

\item  Si $\alpha\in2\pi\Zz$, alors les deux premières séries divergent et la dernière converge.

Soit $\alpha\notin2\pi\Zz$. Pour $n\in\Nn^*$, posons $v_n=e^{in\alpha}$ et $\varepsilon_n=\frac{1}{n}$ de sorte que $u_n=\varepsilon_nv_n$. Pour $n\in\Nn^*$, posons encore $V_n =\sum_{k=1}^{n}v_k$.

Pour $(n ,p)\in(\Nn^*)^2$, posons enfin $R_n^p=\sum_{k=1}^{n+p}u_k-\sum_{k=1}^{n}u_k=\sum_{k=n+1}^{n+p}u_k$. (On effectue alors une transformation d'\textsc{Abel}).

\begin{align*}\ensuremath
R_n^p&=\sum_{k=n+1}^{n+p}\varepsilon_kv_k =\sum_{k=n+1}^{n+p}\varepsilon_k(V_k-V_{k-1}) =\sum_{k=n+1}^{n+p}\varepsilon_kV_k-\sum_{k=n+1}^{n+p}\varepsilon_kV_{k-1} =\sum_{k=n+1}^{n+p}\varepsilon_kV_k-\sum_{k=n}^{n+p-1}\varepsilon_{k+1}V_{k}\\
 &=\varepsilon_{n+p}V_{n+p}-\varepsilon_{n+1}V_n+ \sum_{k=n+1}^{n+p-1}(\varepsilon_k-\varepsilon_{k+1})V_k.
\end{align*}

Maintenant, pour $n\in\Nn^*$, $V_n=e^{i\alpha}\frac{e^{in\alpha}-1}{e^{i\alpha}-1}=e^{i\alpha}\frac{\sin(n\alpha/2)}{\sin(\alpha/2)}$  et donc $\forall n\in\Nn^*$, $|V_n|\leqslant\frac{1}{|\sin(\alpha/2)|}$. Par suite, pour $(n,p)\in(\Nn^*)^2$

\begin{align*}\ensuremath
|R_n^p|&=\left|\frac{1}{n+p}V_{n+p}-\frac{1}{n+1}V_n+ \sum_{k=n+1}^{n+p-1}\left(\frac{1}{k}-\frac{1}{k+1}\right)V_k\right|&\\
 &\leqslant\frac{1}{|\sin(\alpha/2)|}\left(\frac{1}{n+p}+\frac{1}{n+1}+ \sum_{k=n+1}^{n+p-1}\left(\frac{1}{k}-\frac{1}{k+1}\right)\right)\\
 &=\frac{1}{|\sin(\alpha/2)|}\left(\frac{1}{n+p}+\frac{1}{n+1}+\frac{1}{n+1}-\frac{1}{n+p}\right)=\frac{2}{|\sin(\alpha/2)|(n+1)}\\
  &\leqslant\frac{2}{n|\sin(\alpha/2)|}.
\end{align*}

Soit alors $\varepsilon$ un réel strictement positif. Pour $n\geqslant E\left(\frac{2}{\varepsilon|\sin(\alpha/2)|}\right)+ 1$ et p entier naturel non nul quelconque, on a $|R_n^p|<\varepsilon$.

On a montré que $\forall\varepsilon>0$, $\exists n_0\in\Nn^*/$ $\forall(n,p)\in\Nn^*$, $(n\geqslant n_0\Rightarrow\left|\sum_{k=1}^{n+p}u_k-\sum_{k=1}^{n}u_k\right|<\varepsilon$.

Ainsi, la série de terme général $u_n$ vérifie le critère de \textsc{Cauchy} et est donc convergente.
Il en est de même des séries de termes généraux respectifs  $\frac{\cos(n\alpha)}{n}=\text{Re}\left(\frac{e^{in\alpha}}{n}\right)$ et $\frac{\sin(n\alpha)}{n}=\text{Im}\left(\frac{e^{in\alpha}}{n}\right)$.

\item  Pour $x\in]0,+\infty[$, posons $f(x)=\frac{\ln x}{x}$. $f$ est dérivable sur $]0,+\infty[$ et $\forall x>e$, $f'(x)=\frac{1-\ln x}{x}< 0$.

Donc, la fonction $f$ est décroissante sur $[e,+\infty[$. On en déduit que la suite $\left(\frac{\ln n}{n}\right)_{n\geqslant3}$ est une suite décroissante. Mais alors la série de terme général $(-1)^n\frac{\ln n}{n}$ converge en vertu du critère spécial aux séries alternées.

\item  \textbullet~Si $\text{deg}P\geqslant\text{deg}Q$, $u_n$ ne tend pas vers $0$ et la série de terme général $u_n$ est grossièrement divergente.

\textbullet~Si $\text{deg}P\leqslant\text{deg}Q - 2$, $u_n=O\left(\frac{1}{n^2}\right)$ et la série de terme général $u_n$ est absolument convergente.

\textbullet~Si $\text{deg}P =\text{deg}Q - 1$, $u_n\underset{n\rightarrow+\infty}{=}(-1)^n\frac{\text{dom}P}{n\;\text{dom}Q}+O\left(\frac{1}{n^2}\right)$. $u_n$ est alors somme de deux termes généraux de séries convergentes et la série de terme général $u_n$ converge.

En résumé, la série de terme général $u_n$ converge si et seulement si $\text{deg}P <\text{deg}Q$.

\item  $e=\sum_{k=0}^{+\infty}\frac{1}{k!}$  puis pour $n\geqslant2$, $n!e=1+ n+\sum_{k=0}^{n-2}\frac{n!}{k!}+\sum_{k=n+1}^{+\infty}\frac{n!}{k!}$.

Pour $0\leqslant k\leqslant n-2$, $\frac{n!}{k!}$ est un entier divisible par $n(n-1)$ et est donc un entier pair que l'on note $2K_n$. Pour $n\geqslant2$, on obtient

\begin{center}
$\sin(n!\pi e)=\sin\left(2K_n\pi+(n+1)\pi+\pi\sum_{k=n+1}^{+\infty}\frac{n!}{k!}\right)=(-1)^{n+1}\sin\left(\pi\sum_{k=n+1}^{+\infty}\frac{n!}{k!}\right)$.
\end{center}

Déterminons un développement limité à l'ordre $2$ de $\sum_{k=n+1}^{+\infty}\frac{n!}{k!}$ quand $n$ tend vers $+\infty$.

\begin{center}
$\sum_{k=n+1}^{+\infty}\frac{n!}{k!}=\frac{1}{n+1}+\frac{1}{(n+1)(n+2)}+\sum_{k=n+3}^{+\infty}\frac{n!}{k!}$.
\end{center}

Maintenant, pour $k\geqslant n+3$, $\frac{n!}{k!}=\frac{1}{k(k-1)\ldots(n+1)}\leqslant\frac{1}{(n+1)^{k-n}}$ et donc

\begin{center} 
$\sum_{k=n+3}^{+\infty}\frac{n!}{k!}\leqslant\sum_{k=n+3}^{+\infty}\frac{1}{(n+1)^{k-n}}=\frac{1}{(n+1)^3}\times\frac{1}{1-\frac{1}{n+1}}=\frac{1}{n(n+1)^2}\leqslant\frac{1}{n^3}$.
\end{center}

On en déduit que $\sum_{k=n+3}^{+\infty}\frac{n!}{k!}\underset{n\rightarrow+\infty}{=}o\left(\frac{1}{n^2}\right)$. Il reste

\begin{center}
$\sum_{k=n+1}^{+\infty}\frac{n!}{k!}\underset{n\rightarrow+\infty}{=}
\frac{1}{n+1}+\frac{1}{(n+1)(n+2)}+o\left(\frac{1}{n^2}\right)\underset{n\rightarrow+\infty}{=}\frac{1}{n}\left(1+\frac{1}{n}\right)^{-1}+\frac{1}{n^2}+o\left(\frac{1}{n^2}\right)
\underset{n\rightarrow+\infty}{=}\frac{1}{n}+o\left(\frac{1}{n^2}\right)
$.
\end{center}

Finalement , $\sin(n!\pi e)\underset{n\rightarrow+\infty}{=}(-1)^{n+1}\sin\left(\frac{\pi}{n}+o\left(\frac{1}{n^2}\right)\right)=\frac{(-1)^{n+1}\pi}{n}+\left(\frac{1}{n^2}\right)$.

$\sin(n!\pi e)$ est somme de deux termes généraux de séries convergentes et la série de terme général $\sin(n!\pi e)$ converge.

Si $p\geqslant2$, $|\sin^p(n!\pi e)|\underset{n\rightarrow+\infty}{\sim}\frac{\pi^p}{n^p}$ et la série de terme général $\sin^p(n!\pi e)$ converge absolument.
\end{enumerate}
\fincorrection
\correction{005691}
\begin{enumerate}
 \item  $\frac{n+1}{3^n}\underset{n\rightarrow+\infty}{=}o\left(\frac{1}{n^2}\right)$. Par suite, la série de terme général $\frac{n+1}{3^n}$ converge.

\textbf{1er calcul.} Soit $S=\sum_{n=0}^{+\infty}\frac{n+1}{3^n}$. Alors

\begin{align*}\ensuremath
\frac{1}{3}S&=\sum_{n=0}^{+\infty}\frac{n+1}{3^{n+1}}=\sum_{n=1}^{+\infty}\frac{n}{3^{n}}=\sum_{n=1}^{+\infty}\frac{n+1}{3^{n}}-\sum_{n=1}^{+\infty}\frac{1}{3^{n}}\\
 &=(S-1)-\frac{1}{3}\frac{1}{1-\frac{1}{3}}=S-\frac{3}{2}.
\end{align*}

On en déduit que $S=\frac{9}{4}$.

\begin{center}
\shadowbox{
$\sum_{n=0}^{+\infty}\frac{n+1}{3^n}=\frac{9}{4}$.
}
\end{center}

\textbf{2ème calcul.} Pour $x\in\Rr$ et $n\in\Nn$, on pose $f_n(x)=\sum_{k=0}^{n}x^k$.

Soit $n\in\Nn^*$. $f_n$ est dérivable sur $\Rr$ et pour $x\in\Rr$,

\begin{center}
$f_n'(x)=\sum_{k=1}^{n}kx^{k-1}=\sum_{k=0}^{n-1}(k+1)x^k$.
\end{center}

Par suite, pour $n\in\Nn^*$ et $x\in\Rr\setminus\{1\}$

\begin{center}
$\sum_{k=0}^{n-1}(k+1)x^k=f_n'(x)=\left(\frac{x^n-1}{x-1}\right)'(x)=\frac{nx^{n-1}(x-1)-(x^n-1)}{(x-1)^2}=\frac{(n-1)x^n-nx^{n-1}+1}{(x-1)^2}$.
\end{center}

Pour $x=\frac{1}{3}$, on obtient  $\sum_{k=0}^{n-1}\frac{k+1}{3^k}=\frac{\frac{n-1}{3^n}-\frac{n}{3^{n-1}}+1}{\left(\frac{1}{3}-1\right)^2}$ et quand $n$ tend vers l'infini, on obtient de nouveau $S=\frac{9}{4}$.

\item  Pour $k\geqslant3$, $\frac{2k-1}{k^3-4k}=\frac{3}{8(k-2)}+\frac{1}{4k}-\frac{5}{8(k+2)}$. Puis

\begin{align*}\ensuremath
\sum_{k=3}^{n}\frac{2k-1}{k^3-4k}&=\frac{3}{8}\sum_{k=3}^{n}\frac{1}{k-2}+\frac{1}{4}\sum_{k=3}^{n}\frac{1}{k}-\frac{5}{8}\sum_{k=3}^{n}\frac{1}{k+2}=\frac{3}{8}\sum_{k=1}^{n-2}\frac{1}{k}+\frac{1}{4}\sum_{k=3}^{n}\frac{1}{k}-\frac{5}{8}\sum_{k=5}^{n+2}\frac{1}{k}\\
 &\underset{n\rightarrow+\infty}{=}\frac{3}{8}\left(1+\frac{1}{2}+\sum_{k=3}^{n}\frac{1}{k}\right)+\frac{1}{4}\sum_{k=3}^{n}\frac{1}{k}-\frac{5}{8}\left(-\frac{1}{3}-\frac{1}{4}\sum_{k=3}^{n}\frac{1}{k}\right)+o(1)\\
 &\underset{n\rightarrow+\infty}{=}\frac{3}{8}\times\frac{3}{2}+\frac{5}{8}\times\frac{7}{12}+o(1)\underset{n\rightarrow+\infty}{=}\frac{89}{96}+ o(1) .
\end{align*}

La série proposée est donc convergente de somme $\frac{89}{96}$.

\begin{center}
\shadowbox{
$\sum_{n=3}^{+\infty}\frac{2n-1}{n^3-4n}=\frac{89}{96}$.
}
\end{center}

\item  Pour $k\in\Nn$, on a $1^{3k}+j^{3k}+(j^2)^{3k}=3$ puis $1^{3k+1}+j^{3k+1}+(j^2)^{3k+1} =1+j+j^2= 0$ et $1^{3k+2}+j^{3k+2}+(j^2)^{3k+2}=1+j^2 + j^4 = 0$. Par suite,

\begin{center}
$e+e^j+e^{j^2}=\sum_{n=0}^{+\infty}\frac{1^n+j^n+(j^2)^n}{n!}=3\sum_{n=0}^{+\infty}\frac{1}{(3n)!}$,
\end{center}

et donc

\begin{align*}\ensuremath
\sum_{n=0}^{+\infty}\frac{1}{(3n)!}&=\frac{1}{3}(e+e^j+e^{j^2})=\frac{1}{3}\left(e+e^{-\frac{1}{2}+i\frac{\sqrt{3}}{2}}+e^{-\frac{1}{2}-i\frac{\sqrt{3}}{2}}\right) =\frac{1}{3}\left(e + 2e^{-1/2}\text{Re}(e^{-i\sqrt{3}/2})\right)\\
 &=\frac{1}{3}\left(e +2e^{-1/2}\cos\left(\frac{\sqrt{3}}{2}\right)\right).
\end{align*}

\begin{center}
\shadowbox{
$\sum_{n=0}^{+\infty}\frac{1}{(3n)!}=\frac{1}{3}\left(e +\frac{2}{\sqrt{e}}\cos\left(\frac{\sqrt{3}}{2}\right)\right)$.
}
\end{center}

\item 

\begin{align*}\ensuremath
\sum_{k=2}^{n}\left(\frac{1}{\sqrt{k-1}}+\frac{1}{\sqrt{k+1}}-\frac{2}{\sqrt{k}}\right)&=\sum_{k=2}^{n}\left(\left(\frac{1}{\sqrt{k-1}}-\frac{1}{\sqrt{k}}\right)-\left(\frac{1}{\sqrt{k}}-\frac{1}{\sqrt{k+1}}\right)\right)\\
 &=\left(1-\frac{1}{\sqrt{2}}\right)-\left(\frac{1}{\sqrt{n}}-\frac{1}{\sqrt{n+1}}\right)\;(\text{somme télescopique})\\
 &\underset{n\rightarrow+\infty}{=}1-\frac{1}{\sqrt{2}}+o(1) 
\end{align*}

\begin{center}
\shadowbox{
$\sum_{n=2}^{+\infty}\left(\frac{1}{\sqrt{n-1}}+\frac{1}{\sqrt{n+1}}-\frac{2}{\sqrt{n}}\right)=1-\frac{1}{\sqrt{2}}$.
}
\end{center}

\item  $\ln\left(1+\frac{(-1)^n}{n}\right)\underset{n\rightarrow+\infty}{=}\frac{(-1)^n}{n}+O\left(\frac{1}{n^2}\right)$. Donc la série de terme général $\ln\left(1+\frac{(-1)^n}{n}\right)$ converge.

Posons $S=\sum_{k=2}^{+\infty}\ln\left(1+\frac{(-1)^k}{k}\right)$ puis pour $n\geqslant2$, $S_n=\sum_{k=2}^{n}\ln\left(1+\frac{(-1)^k}{k}\right)$. Puisque la série converge $S=\lim_{n \rightarrow +\infty}S_n=\lim_{p \rightarrow +\infty}S_{2p+1}$ avec

\begin{align*}\ensuremath
S_{2p+1}&=\sum_{k=2}^{2p+1}\ln\left(1+\frac{(-1)^k}{k}\right)=\sum_{k=1}^{p}\left(\ln\left(1-\frac{1}{2k+1}\right)+\ln\left(1+\frac{1}{2k}\right)\right)\\
 &=\sum_{k=1}^{p}(\ln(2k)-\ln(2k+1)+\ln(2k+1)-\ln(2k))=0
\end{align*}

et quand $p$ tend vers $+\infty$, on obtient $S = 0$.

\begin{center}
\shadowbox{
$\sum_{n=2}^{+\infty}\ln\left(1+\frac{(-1)^n}{n}\right)=0$.
}
\end{center}

\item  Si $a\in\left]0,\frac{\pi}{2}\right[$ alors, pour tout entier naturel $n$, $\frac{a}{2^n}\in\left]0,\frac{\pi}{2}\right[$ et donc $\cos\left(\frac{a}{2^n}\right)>0$.

Ensuite, $\ln\left(\cos\left(\frac{a}{2^n}\right)\right)\underset{n\rightarrow+\infty}{=}\ln\left(1+O\left(\frac{1}{2^{2n}}\right)\right)\underset{n\rightarrow+\infty}{=}O\left(\frac{1}{2^{2n}}\right)$ et la série converge. Ensuite,

\begin{align*}\ensuremath
\sum_{k=0}^{n}\ln\left(\cos\left(\frac{a}{2^k}\right)\right)&=\ln\left(\prod_{k=0}^{n}\cos\left(\frac{a}{2^k}\right)\right)=\ln\left(\prod_{k=0}^{n}\frac{\sin\left(2\times\frac{a}{2^k}\right)}{2\sin\left(\frac{a}{2^k}\right)}\right)=\ln\left(\frac{1}{2^{n+1}}\prod_{k=0}^{n}\frac{\sin\left(\frac{a}{2^{k-1}}\right)}{\sin\left(\frac{a}{2^k}\right)}\right)\\
 &=\ln\left(\frac{\sin(2a)}{2^{n+1}\sin\left(\frac{a}{2^n}\right)}\right)\;(\text{produit télescopique})\\
 &\underset{n\rightarrow+\infty}{\sim}\ln\left(\frac{\sin(2a)}{2^{n+1}\times\frac{a}{2^n}}\right)=\ln\left(\frac{\sin(2a)}{2a}\right).
\end{align*}

\begin{center}
\shadowbox{
$\forall a\in\left]0,\frac{\pi}{2}\right[$, $\sum_{n=0}^{+\infty}\ln\left(\cos\left(\frac{a}{2^n}\right)\right)=\ln\left(\frac{\sin(2a)}{2a}\right)$.
}
\end{center}

\item  Vérifions que pour tout réel $x$ on a $\tanh(2x)=\frac{2\tanh x}{1+\tanh^2x}$. Soit $x\in\Rr$.

$\ch^2x+\sh^2x=\frac{1}{4}((e^x+e^{-x})^2+(e^x-e^{-x})^2)=\frac{1}{2}(e^{2x}+e^{-2x})=\ch(2x)$ et $2\sh x\ch x=\frac{1}{2}(e^{x}-e^{-x})(e^x+e^{-x})=\frac{1}{2}(e^{2x}-e^{-2x})=\sh(2x)$ puis

\begin{center}
$\frac{2\tanh x}{1+\tanh^2x}=\frac{2\sh x\ch x}{\ch^2x+\sh^2x}=\frac{\sh(2x)}{\ch(2x)}=\tanh(2x)$.
\end{center}

Par suite, pour $x\in\Rr^*$, $\tanh x=\frac{2}{\tanh(2x)}-\frac{1}{\tanh x}$. Mais alors, pour $a\in\Rr^*$ et $n\in\Nn$

\begin{align*}\ensuremath
\sum_{k=0}^{n}\frac{1}{2^k}\tanh\left(\frac{a}{2^k}\right)&=\sum_{k=0}^{n}\frac{1}{2^k}\left(\frac{2}{\tanh\frac{a}{2^{k-1}}}-\frac{1}{\tanh\frac{a}{2^{k}}}\right)=\sum_{k=0}^{n}\left(\frac{1}{2^{k-1}\tanh\frac{a}{2^{k-1}}}-\frac{1}{2^k\tanh\frac{a}{2^{k}}}\right)\\
 &=\frac{2}{\tanh(2a)}-\frac{1}{2^n\tanh\frac{a}{2^{n}}}\;(\text{somme télescopique})\\
 &\underset{n\rightarrow+\infty}{\rightarrow}\frac{2}{\tanh(2a)}-\frac{1}{a},
\end{align*}

ce qui reste vrai quand $a=0$.

\begin{center}
\shadowbox{
$\forall a\in\Rr$, $\sum_{n=0}^{+\infty}\frac{1}{2^n}\tanh\left(\frac{a}{2^n}\right)=\frac{2}{\tanh(2a)}-\frac{1}{a}$.
}
\end{center}
\end{enumerate}	
\fincorrection
\correction{005692}
Il faut vérifier que $nu_n\underset{n\rightarrow+\infty}{\rightarrow}0$. Pour $n\in\Nn$, posons $S_n =\sum_{k=0}^{n}u_k$. Pour $n\in\Nn$, on a

\begin{align*}\ensuremath
0<(2n)u_{2n}&=2(\underbrace{u_{2n}+\ldots+u_{2n}}_{n})\leqslant2\sum_{k=n+1}^{2n}u_k\;(\text{car la suite}\;u\;\text{est décroissante})\\
 &= 2(S_{2n} - S_n).
\end{align*}

Puisque la série de terme général $u_n$ converge, $\lim_{n \rightarrow +\infty}2(S_{2n} - S_n)=0$ et donc $\lim_{n \rightarrow +\infty}(2n)u_{2n}=0$.

Ensuite, $0 < (2n+1)u_{2n+1}\leqslant(2n+1)u_{2n}=(2n)u_{2n}+u_{2n}\underset{n\rightarrow+\infty}{\rightarrow}0$. Donc les suites des termes de rangs pairs et impairs extraites de la suite $(nu_n)_{n\in\Nn}$ convergent et ont même limite à savoir $0$. On en déduit que $\lim_{n \rightarrow +\infty}nu_n=0$ ou encore que $u_n\underset{n\rightarrow+\infty}{=}o\left(\frac{1}{n}\right)$.

Contre exemple avec $u$ non monotone. Pour $n\in\Nn$, on pose $u_n=\left\{
\begin{array}{l}
0\;\text{si}\;n=0\\
\rule[-4mm]{0mm}{10mm}\frac{1}{n}\;\text{si}\;n\;\text{est un carré parfait non nul}\\
0\;\text{sinon}
\end{array}
\right.$. La suite $u$ est positive et $\sum_{n=0}^{+\infty}u_n=\sum_{p=1}^{+\infty}\frac{1}{p^2}<+\infty$. Pourtant, $p^2u_{p^2}=1\underset{p\rightarrow+\infty}{\rightarrow}1$ et la suite $(nu_n)$ admet une suite extraite convergeant vers $1$. On a donc pas $\lim_{n \rightarrow +\infty}nu_n=0$.
\fincorrection
\correction{005693}
Soit $\sigma$ une permutation de $\llbracket1,n\rrbracket$. Montrons que la suite $S_n=\sum_{k=1}^{n}\frac{\sigma(k)}{k^2}$, $n\geqslant1$, ne vérifie pas le critère de \textsc{Cauchy}. Soit $n\in\Nn^*$.

\begin{align*}\ensuremath
S_{2n}-S_n&=\sum_{k=n+1}^{2n}\frac{\sigma(k)}{k^2}\geqslant\frac{1}{(2n)^2}\sum_{k=n+1}^{2n}\sigma(k)\\
 &\geqslant\frac{1}{4n^2}(1+2+...+n)\;(\text{car les}\;n\;\text{entiers}\;\sigma(k),\;1\leqslant k\leqslant n,\;\text{sont strictement positifs et deux à deux distincts})\\
 &=\frac{n(n+1)}{8n^2}\geqslant\frac{n^2}{8n^2}=\frac{1}{8}.
\end{align*}

Si la suite $(S_n)$ converge, on doit avoir $\lim_{n \rightarrow +\infty}(S_{2n}-S_n)=0$ ce qui contredit l'inégalité précédente. Donc la série de terme général $\frac{\sigma(n)}{n^2}$, $n\geqslant1$, diverge.
\fincorrection
\correction{005694}
Pour $n\in\Nn$, posons $v_n=\ln(1+u_n)$, $w_n=\frac{u_n}{1+u_n}$ et $t_n=\int_{0}^{u_n}\frac{dx}{1+x^e}$.

\textbullet~Si $u_n\underset{n\rightarrow+\infty}{\rightarrow}0$, alors $0\leqslant u_n\underset{n\rightarrow+\infty}{\sim}v_n\underset{n\rightarrow+\infty}{\sim}w_n$. Dans ce cas, les séries de termes généraux $u_n$, $v_n$ et $w_n$ sont de même nature.

D'autre part, pour $n\in\Nn$, $\frac{u_n}{1+u_n^e}\leqslant t_n\leqslant u_n$ puis $\frac{1}{1+u_n^e}\leqslant\frac{t_n}{u_n}\leqslant1$et donc $t_n\underset{n\rightarrow+\infty}{\sim}u_n$. Les séries de termes généraux $u_n$ et $t_n$ sont aussi de même nature.

\textbullet~Si $u_n$ ne tend pas vers $0$, la série de terme général $u_n$ est grossièrement divergente. Puisque $u_n =e^{v_n}-1$, $v_n$ ne tend pas vers $0$ et la série de terme général $v_n$ est grossièrement divergente. Dans ce cas aussi, les séries de termes généraux sont de même nature.

De même, puisque $w_n=\frac{u_n}{1+u_n}<1$, on a $u_n =\frac{w_n}{1-w_n}$ et $w_n$ ne peut tendre vers $0$.

Enfin, puisque $u_n$ ne tend pas vers $0$, il existe $\varepsilon> 0$ tel que pour tout entier naturel $N$, il existe $n = n(N)\geqslant N$ tel que $u_n\geqslant\varepsilon$. Pour cet $\epsilon$ et ces $n$, on a $t_n\geqslant\int_{0}^{\varepsilon}\frac{dx}{1+x^e}>0$ (fonction continue, positive et non nulle) et la suite $t_n$ ne tend pas vers $0$. Dans le cas où $u_n$ ne tend pas vers $0$, les quatre séries sont grossièrement divergentes.

\fincorrection
\correction{005695}
Pour $n\in\Nn$, posons $u_n= (n+1)!\left(e-\sum_{k=0}^{n}\frac{1}{k!}\right)$. Soit $n\in\Nn$.

\begin{align*}\ensuremath
u_n&=\sum_{k=n+1}^{+\infty}\frac{(n+1)!}{k!}\\
 &=1+\frac{1}{n+2}+\frac{1}{(n+2)(n+3)}+\frac{1}{(n+2)(n+3)(n+4)}+\frac{1}{(n+2)(n+3)(n+4)(n+5)}+\sum_{k=n+6}^{+\infty}\frac{1}{(n+2)(n+3)\ldots k}
\end{align*}
	

On a $0 <\sum_{k=n+6}^{+\infty}\frac{1}{(n+2)(n+3)\ldots k}=\sum_{k=n+1}^{+\infty}\frac{1}{(n+2)^{k-(n+1)}}=\frac{1}{(n+2)^5}\frac{1}{1-\frac{1}{n+2}}=\frac{1}{(n+2)^4(n+1)}\leqslant\frac{1}{n^5}$. On en déduit que $\sum_{k=n+6}^{+\infty}\frac{1}{(n+2)(n+3)\ldots k}\underset{n\rightarrow+\infty}{=}o\left(\frac{1}{n^4}\right)$. Donc

\begin{align*}\ensuremath
u_n&\underset{n\rightarrow+\infty}{=}1+\frac{1}{n+2}+\frac{1}{(n+2)(n+3)}+\frac{1}{(n+2)(n+3)(n+4)}+\frac{1}{(n+2)(n+3)(n+4)(n+5)}+o\left(\frac{1}{n^4}\right)\\
 &\underset{n\rightarrow+\infty}{=}1+\frac{1}{n}\left(1+\frac{2}{n}\right)^{-1}+\frac{1}{n^2}\left(1+\frac{2}{n}\right)^{-1}\left(1+\frac{3}{n}\right)^{-1}+\frac{1}{n^3}\left(1+\frac{2}{n}\right)^{-1}\left(1+\frac{3}{n}\right)^{-1}\left(1+\frac{4}{n}\right)^{-1}+\frac{1}{n^4}+o\left(\frac{1}{n^4}\right)\\
  &\underset{n\rightarrow+\infty}{=}1+\frac{1}{n}\left(1-\frac{2}{n}+\frac{4}{n^2}-\frac{8}{n^3}\right)+\frac{1}{n^2}\left(1-\frac{2}{n}+\frac{4}{n^2}\right)\left(1-\frac{3}{n}+\frac{9}{n^2}\right)+\frac{1}{n^3}\left(1-\frac{2}{n}\right)\left(1-\frac{3}{n}\right)\left(1-\frac{4}{n}\right)\\
  &\;+\frac{1}{n^4}+o\left(\frac{1}{n^4}\right)\\
  &\underset{n\rightarrow+\infty}{=}1+\frac{1}{n}\left(1-\frac{2}{n}+\frac{4}{n^2}-\frac{8}{n^3}\right)+\frac{1}{n^2}\left(1-\frac{5}{n}+\frac{19}{n^2}\right)+\frac{1}{n^3}\left(1-\frac{9}{n}\right)+\frac{1}{n^4}+o\left(\frac{1}{n^4}\right)\\
  &\underset{n\rightarrow+\infty}{=}1+\frac{1}{n}-\frac{1}{n^2}+\frac{3}{n^4}+o\left(\frac{1}{n^4}\right).
\end{align*}

Finalement

\begin{center}
\shadowbox{
$(n+1)!\left(e-\sum_{k=0}^{n}\frac{1}{k!}\right)\underset{n\rightarrow+\infty}{=}1+\frac{1}{n}-\frac{1}{n^2}+\frac{3}{n^4}+o\left(\frac{1}{n^4}\right)$.
}
\end{center}
\fincorrection
\correction{005696}
 Pour $n\in\Nn$, posons $u_n=\sin\left(\pi(2+\sqrt{3})^n\right)$. D'après la formule du binôme de \textsc{Newton}, $(2+\sqrt{3})^n=A_n +B_n\sqrt{3}$ où $A_n$ et $B_n$ sont des entiers naturels. Un calcul conjugué fournit aussi $(2-\sqrt{3})^n=A_n-B_n\sqrt{3}$. Par suite, $(2+\sqrt{3})^n+(2-\sqrt{3})^n= 2A_n$ est un entier pair. Par suite, pour $n\in\Nn$,

\begin{center}
$u_n =\sin\left(2A_n\pi-\pi(2-\sqrt{3})^n\right)=-\sin\left(\pi(2-\sqrt{3})^n\right)$.
\end{center}

Mais $0< 2-\sqrt{3}< 1$ et donc $(2-\sqrt{3})^n\underset{n\rightarrow+\infty}{\rightarrow}0$. On en déduit que $|u_n|\underset{n\rightarrow+\infty}{\sim}\pi(2-\sqrt{3})^n$ terme général d'une série géométrique convergente. Donc la série de terme général $u_n$ converge.
\fincorrection
\correction{005700}
Si $\alpha<0$, $u_n\underset{n\rightarrow+\infty}{\sim}n^{-2\alpha}$  et si $\alpha=0$, $u_n=1+(-1)^n$. Donc si $\alpha\leqslant0$, $u_n$ ne tend pas vers $0$. La série de terme général $u_n$ diverge grossièrement dans ce cas.

On suppose dorénavant que $\alpha> 0$. Pour tout entier naturel non nul $n$, $|u_n|\underset{n\rightarrow+\infty}{\sim}\frac{1}{n^\alpha}$ et donc la série de terme général $u_n$ converge absolument si et seulement si $\alpha> 1$.

Il reste à étudier le cas où $0 <\alpha\leqslant1$. On a $u_n=\frac{(-1)^n}{n^\alpha}+\frac{1}{n^{2\alpha}}$. La suite $\left(\frac{1}{n^\alpha}\right)_{n\geqslant1}$ tend vers $0$ en décroissant et donc la série de terme général $\frac{(-1)^n}{n^\alpha}$ converge en vertu du critère spécial aux séries alternées. On en déduit que la série de terme général $u_n$ converge si et seulement si la série de terme général $\frac{1}{n^{2\alpha}}$ converge ou encore si et seulement si $\alpha>\frac{1}{2}$.

En résumé

\begin{center}
\shadowbox{
\begin{tabular}{l}
Si $\alpha\leqslant0$, la série de terme général $\frac{1+(-1)^nn^\alpha}{n^{2\alpha}}$ diverge grossièrement,\\
\rule[-4mm]{0mm}{10mm}si $0 <\alpha\leqslant\frac{1}{2}$, la série de terme général $\frac{1+(-1)^nn^\alpha}{n^{2\alpha}}$ diverge,\\
\rule[-4mm]{0mm}{10mm}si  $\frac{1}{2}<\alpha\leqslant1$, la série de terme général $\frac{1+(-1)^nn^\alpha}{n^{2\alpha}}$ est semi convergente,\\
si $\alpha> 1$, la série de terme général $\frac{1+(-1)^nn^\alpha}{n^{2\alpha}}$ converge absolument.
\end{tabular}
}
\end{center}
\fincorrection
\correction{005701}
Pour $n\in\Nn^*$, on note $S_n$ la somme des $n$ premiers termes de la série considérée et on pose $H_n=\sum_{k=1}^{n}\frac{1}{k}$. Il est connu que $H_n\underset{n\rightarrow+\infty}{=}\ln n+\gamma+o(1)$.

Soit $m\in\Nn^*$.

\begin{align*}\ensuremath
S_{m(p+q)}&=\left(1+\frac{1}{3}+\ldots+\frac{1}{2p-1}\right) -\left(\frac{1}{2}+\frac{1}{4}+\ldots+\frac{1}{2q}\right)+\left(\frac{1}{2p+1}+...+\frac{1}{4p-1}\right) -\left(\frac{1}{2q+2}+...+\frac{1}{4q}\right) +...\\
 &\;+\left(\frac{1}{2(m-1)p+1}+...+\frac{1}{2mp-1}\right) -\left(\frac{1}{2(m-1)q+2}+...+\frac{1}{2mq}\right)\\
 &=\sum_{k=1}^{mp}\frac{1}{2k-1}-\sum_{k=1}^{mq}\frac{1}{2k}=\sum_{k=1}^{2mp}\frac{1}{k}-\sum_{k=1}^{mp}\frac{1}{2k}-\sum_{k=1}^{mq}\frac{1}{2k}= H_{2mp} -\frac{1}{2}(H_{mp}+H_{mq})\\
  &\underset{m\rightarrow+\infty}{=}(\ln(2mp)+\gamma) -\frac{1}{2}(\ln(mp) +\gamma+\ln(mq) +\gamma)+ o(1)=\ln2 +\frac{1}{2}\ln\left(\frac{p}{q}\right)+o(1).
\end{align*}

Ainsi, la suite extraite $(S_{m(p+q)})_{m\in\Nn^*}$ converge vers $\ln2 +\frac{1}{2}\ln\left(\frac{p}{q}\right)$.

Montrons alors que la suite $(S_n)_{n\in\Nn^*}$ converge. Soit $n\in\Nn^*$. Il existe un unique entier naturel non nul $m_n$ tel que $m_n(p+q)\leqslant n <(m_n+1)(p+q)$ à savoir $m_n=E\left(\frac{n}{p+q}\right)$.

\begin{align*}\ensuremath
|S_n-S_{m_n(p+q)}|&\leqslant\frac{1}{2m_np+1}+\ldots+\frac{1}{2(m_n+1)p-1}+\frac{1}{2m_nq+2}+\frac{1}{2(m_n+1)q}\\
 &\leqslant\frac{p}{2m_np+1}+\frac{q}{2m_nq+2}\leqslant\frac{1}{2m_n}+\frac{1}{2m_n}=\frac{1}{m_n}.
\end{align*}

Soit alors $\varepsilon>0$.

Puisque $\lim_{n \rightarrow +\infty}m_n=+\infty$, il existe $n_0\in\Nn^*$ tel que pour $n\geqslant n_0$, $\frac{1}{m_n}<\frac{\varepsilon}{2}$ et aussi $\left|S_{m_n(p+q)}-\ln2-\frac{1}{2}\ln\left(\frac{p}{q}\right)\right|<\frac{\varepsilon}{2}$. Pour $n\geqslant n_0$, on a alors

\begin{align*}\ensuremath
\left|S_{n}-\ln2 -\frac{1}{2}\ln\left(\frac{p}{q}\right)\right|&\leqslant|S_n-S_{m_n(p+q)}|+\left|S_{m_n(p+q)}-\ln2-\frac{1}{2}\ln\left(\frac{p}{q}\right)\right|\leqslant\frac{1}{m_n}+\left|S_{m_n(p+q)}-\ln2-\frac{1}{2}\ln\left(\frac{p}{q}\right)\right|\\
 &<\frac{\varepsilon}{2}+\frac{\varepsilon}{2}=\varepsilon.
\end{align*}

On a montré que $\forall \varepsilon>0,\;\exists n_0\in\Nn^*/\;\forall n\in\Nn,\;(n\geqslant n_0\Rightarrow\left|S_{n}-\left(\ln2 +\frac{1}{2}\ln\left(\frac{p}{q}\right)\right)\right|<\varepsilon)$ et donc, la série proposée converge et a pour somme $\ln2 +\frac{1}{2}\ln\left(\frac{p}{q}\right)$.
\fincorrection
\correction{005702}
La série proposée est le produit de \textsc{Cauchy} de la série de terme général   $\frac{1}{n^\alpha}$, $n\geqslant1$, par elle même.

\textbullet~Si $\alpha>1$, on sait que la série de terme général $\frac{1}{n^\alpha}$ converge absolument et donc que la série proposée converge.

\textbullet~Si $0\leqslant\alpha\leqslant1$, pour $0 < k < n$ on a $0<k(n-k)\leqslant\frac{n}{2}\left(n-\frac{n}{2}\right)=\frac{n^2}{4}$. Donc $u_n\geqslant\frac{n-1}{\left(\frac{n^2}{4}\right)^\alpha}$ avec $\frac{n-1}{\left(\frac{n^2}{4}\right)^\alpha}\underset{n\rightarrow+\infty}{\sim}\frac{4^\alpha}{n^{2\alpha-1}}$. Comme $2\alpha-1\leqslant1$, la série proposée diverge.

\textbullet~Si $\alpha< 0$, $u_n\geqslant\frac{1}{(n-1)^\alpha}$  et donc $u_n$ ne tend pas vers $0$. Dans ce cas, la série proposée diverge grossièrement.

\fincorrection
\correction{005703}
\begin{enumerate}
 \item  Soit $n\in\Nn$.

\begin{align*}\ensuremath
2n^3-3n^2+1&=2(n+3)(n+2)(n+1)-15n^2 - 22n -11= 2(n+3)(n+2)(n+1)-15(n+3)(n+2) +53n+79\\
 &= 2(n+3)(n+2)(n+1)-15(n+3)(n+2)+53(n+3)-80
\end{align*}

Donc

\begin{align*}\ensuremath
\sum_{n=0}^{+\infty}\frac{2n^3-3n^2+1}{(n+3)!}&=\sum_{n=0}^{+\infty}\left(\frac{2}{n!}-\frac{15}{(n+1)!}+\frac{53}{(n+2)!}-\frac{80}{(n+3)!}\right)=2e-15(e-1)+53(e-2)-80\left(e-\frac{5}{2}\right)\\
 &=-40e + 111.
\end{align*}

\begin{center}
\shadowbox{
$\sum_{n=0}^{+\infty}\frac{2n^3-3n^2+1}{(n+3)!}=-40e + 111$.
}
\end{center}

\item  Pour $n\in\Nn$, on a $u_{n+1}=\frac{n+1}{a+n+1}u_n$. Par suite $(n+a+1)u_{n+1}= (n+1)u_n = (n+a)u_n + (1-a)u_n$ puis

\begin{center}
$(1-a)\sum_{k=1}^{n}u_k=\sum_{k=1}^{n}(k+a+1)u_{k+1}-\sum_{k=1}^{n}(k+a)u_k =(n+a+1)u_{n+1}-(a+1)u_1 =(n+a+1)u_{n+1}-1.$
\end{center}

Si $a = 1$, $\forall n\in\Nn^*$, $u_n=\frac{1}{n+1}$. Dans ce cas, la série diverge.

Si $a\neq 1$, $\forall n\in\Nn^*$, $\sum_{k=1}^{n}u_k=\frac{1}{1-a}((n+a+1)u_{n+1} - 1)=\frac{1}{a-1}- \frac{1}{a-1}(a+n+1)u_{n+1}$.

Si $a > 1$, la suite $u$ est strictement positive et la suite des sommes partielles $(S_n)$ est majorée par $\frac{1}{a-1}$. Donc la série de terme général $u_n$ converge. Il en est de même de la suite $((a+n+1)u_{n+1})$. Soit $\ell=\lim_{n \rightarrow +\infty}(a+n+1)u_{n+1}$.

Si $\ell\neq0$, $u_{n+1}\underset{n\rightarrow+\infty}{\sim}\frac{\ell}{n+a+1}$ contredisant la convergence de la série de terme général $u_n$. Donc $\ell= 0$ et

\begin{center}
si $a > 1$,  $\sum_{n=1}^{+\infty}u_n=\frac{1}{a-1}$.
\end{center}

Si $0<a<1$, pour tout $n\in\Nn^*$, $u_n\geqslant\frac{1\times2\times\ldots\times n}{2\times3\ldots\times(n+1)}=\frac{1}{n+1}$. Dans ce cas, la série diverge.
\end{enumerate}
\fincorrection
\correction{005706}
Pour tout entier naturel non nul $n$, $0<\frac{1}{2^pn^{p-1}}=\sum_{k=1}^{n}\frac{1}{(2n)^p}\leqslant\sum_{k=1}^{n}\frac{1}{(n+k)^p}\leqslant\sum_{k=1}^{n}\frac{1}{n^p}=\frac{1}{n^{p-1}}$ et la série de terme général $u_n$ converge si et seulement si $p > 2$.
\fincorrection
\correction{005707}
Pour $n\in\Nn^*$, posons $R_n=\sum_{k=n+1}^{+\infty}\frac{1}{k^2}$. Puisque la série de terme général $\frac{1}{k^2}$, $k\geqslant 1$, converge, la suite $(R_n)$ est définie et tend vers $0$ quand $n$ tend vers $+\infty$.

$0<\frac{1}{k^2}\underset{k\rightarrow+\infty}{\sim}\frac{1}{k(k-1)}=\frac{1}{k-1}-\frac{1}{k}$  et puisque la série de terme général $\frac{1}{k^2}$ converge, la règle de l'équivalence des restes de séries à termes positifs convergentes permet d'affirmer que

\begin{align*}\ensuremath
R_n&= \sum_{k=n+1}^{+\infty}\frac{1}{k^2}\underset{n\rightarrow+\infty}{\sim} \sum_{k=n+1}^{+\infty}\left(\frac{1}{k-1}-\frac{1}{k}\right)\\
 &=\lim_{N \rightarrow +\infty}\sum_{k=n+1}^{N}\left(\frac{1}{k-1}-\frac{1}{k}\right)\;(\text{surtout ne pas décomposer en deux sommes})\\
 &=\lim_{N \rightarrow +\infty}\left(\frac{1}{n}-\frac{1}{N}\right)\;(\text{somme télescopique})\\
 &=\frac{1}{n} 
\end{align*}

ou encore $R_n\underset{n\rightarrow+\infty}{=}\frac{1}{n}+o\left(\frac{1}{n}\right)$.

Plus précisément, pour $n\in\Nn^*$, $R_n-\frac{1}{n}=\sum_{k=n+1}^{+\infty}\frac{1}{k^2}-\sum_{k=n+1}^{+\infty}\frac{1}{k(k-1)}=-\sum_{k=n+1}^{+\infty}\frac{1}{k^2(k-1)}$.

Or $-\frac{1}{k^2(k-1)}+\frac{1}{k(k-1)(k-2)}=\frac{2}{k^2(k-1)(k-2)}$ puis

$\frac{2}{k^2(k-1)(k-2)}-\frac{2}{k(k-1)(k-2)(k-3)}=-\frac{6}{k^2(k-1)(k-2)(k-3)}$ et donc

\begin{align*}\ensuremath
R_n&=\frac{1}{n}-\sum_{k=n+1}^{+\infty}\frac{1}{k^2(k-1)}=\frac{1}{n}-\sum_{k=n+1}^{+\infty}\frac{1}{k(k-1)(k-2)}+\sum_{k=n+1}^{+\infty}\frac{2}{k^2(k-1)(k-2)}\\
 &=\frac{1}{n}-\sum_{k=n+1}^{+\infty}\frac{1}{k(k-1)(k-2)}+\sum_{k=n+1}^{+\infty}\frac{2}{k(k-1)(k-2)(k-3)}-\sum_{k=n+1}^{+\infty}\frac{6}{k^2(k-1)(k-2)(k-3)}
\end{align*}

Ensuite  $\sum_{k=n+1}^{+\infty}\frac{1}{k^2(k-1)(k-2)(k-3)}\underset{n\rightarrow+\infty}{\sim}\sum_{k=n+1}^{+\infty}\frac{1}{k^5}\underset{n\rightarrow+\infty}{\sim}\frac{1}{4n^4}$ ou encore $-\sum_{k=n+1}^{+\infty}\frac{6}{k^2(k-1)(k-2)(k-3)}\underset{n\rightarrow+\infty}{=}-\frac{3}{2n^4}+o\left(\frac{1}{n^4}\right)$. Puis

\begin{align*}\ensuremath
\sum_{k=n+1}^{+\infty}\frac{1}{k(k-1)(k-2)}&=\lim_{N \rightarrow +\infty}\frac{1}{2}\sum_{k=n+1}^{N}\left(\frac{1}{(k-1)(k-2)}-\frac{1}{k(k-1)}\right)=\lim_{N \rightarrow +\infty}\frac{1}{2}\left(\frac{1}{n(n-1)}-\frac{1}{N(N-1)}\right)=\frac{1}{2n(n-1)}\\
 &=\frac{1}{2n^2}\left(1-\frac{1}{n}\right)^{-1}\underset{n\rightarrow+\infty}{=}\frac{1}{2n^2}+\frac{1}{2n^3}+\frac{1}{2n^4}+o\left(\frac{1}{n^4}\right)
\end{align*}

et

\begin{align*}\ensuremath
\sum_{k=n+1}^{+\infty}\frac{2}{k(k-1)(k-2)(k-3)}&=\lim_{N \rightarrow +\infty}\frac{2}{3}\sum_{k=n+1}^{N}\left(\frac{1}{(k-1)(k-2)(k-3)}-\frac{1}{k(k-1)(k-2)}\right)\\
 &=\lim_{N \rightarrow +\infty}\frac{2}{3}\left(\frac{1}{n(n-1)(n-2)}-\frac{1}{N(N-1)(N-2)}\right)=\frac{2}{3n(n-1)(n-2)}\\
 &=\frac{2}{3n^3}\left(1-\frac{1}{n}\right)^{-1}\left(1-\frac{2}{n}\right)^{-1}\underset{n\rightarrow+\infty}{=}\frac{2}{3n^3}\left(1+\frac{1}{n}+o\left(\frac{1}{n}\right)\right)\left(1+\frac{2}{n}+o\left(\frac{1}{n}\right)\right)\\
  &\underset{n\rightarrow+\infty}{=}\frac{2}{3n^3}+\frac{2}{n^4}+o\left(\frac{1}{n^4}\right)
\end{align*}

et finalement

\begin{center}
$R_n\underset{n\rightarrow+\infty}{=}\frac{1}{n}-\left(\frac{1}{2n^2}+\frac{1}{2n^3}+\frac{1}{2n^4}\right)+\left(\frac{2}{3n^3}+\frac{2}{n^4}\right)-\frac{3}{2n^4}+o\left(\frac{1}{n^4}\right)\underset{n\rightarrow+\infty}{=}\frac{1}{n}-\frac{1}{2n^2}+\frac{1}{6n^3}+o\left(\frac{1}{n^4}\right)$.
\end{center}

\begin{center}
\shadowbox{
$\sum_{k=n+1}^{+\infty}\frac{1}{k^2}\underset{n\rightarrow+\infty}{=}\frac{1}{n}-\frac{1}{2n^2}+\frac{1}{6n^3}+o\left(\frac{1}{n^4}\right)$.
}
\end{center}
\fincorrection
\correction{005708}
\begin{enumerate}
 \item  La suite $\left(\frac{\ln n}{n}\right)_{n\in\Nn^*}$ tend vers $0$, en décroissant à partir du rang $3$ (fourni par l'étude de la fonction $x\mapsto\frac{\ln x}{x}$ sur $[e,+\infty[$) et donc la série de terme général $(-1)^n\frac{\ln n}{n}$, $n\geqslant1$, converge en vertu du critère spécial aux séries alternées. Pour $n\in\Nn^*$, on pose $R_n=\sum_{p=n+1}^{+\infty}(-1)^p\frac{\ln p}{p}$.

$(-1)^k\frac{\ln k}{k}$  n'est pas de signe constant à partir d'un certain rang et on ne peut donc lui appliquer la règle de l'équivalence des restes.

Par contre, puisque la série de terme général $(-1)^k\frac{\ln k}{k}$ converge, on sait que l'on peut associer les termes à volonté et pour $k\in\Nn^*$, on a

\begin{center}
$R_{2k-1}=\sum_{p=2k}^{+\infty}(-1)^p\frac{\ln p}{p}=\sum_{p=k}^{+\infty}\left(\frac{\ln(2p)}{2p}-\frac{\ln(2p+1)}{2p+1}\right)$.
\end{center}

Puisque la fonction  $x\mapsto\frac{\ln x}{x}$ est décroissante sur $[e,+\infty[$ et donc sur $[3,+\infty[$, pour $p\geqslant 2$, $\frac{\ln(2p)}{2p}-\frac{\ln(2p+1)}{2p+1}\geqslant0$ et on peut utiliser la règle de l'équivalence des restes de séries à termes positifs convergentes.

Cherchons déjà un équivalent plus simple de $\frac{\ln(2p)}{2p}-\frac{\ln(2p+1)}{2p+1}$ quand $p$ tend vers $+\infty$.

\begin{align*}\ensuremath
\frac{\ln(2p)}{2p}-\frac{\ln(2p+1)}{2p+1}&=\frac{\ln(2p)}{2p}-\frac{1}{2p}\left(\ln(2p)+\ln\left(1+\frac{1}{2p}\right)\right)\left(1+\frac{1}{2p}\right)^{-1}\\
 &\underset{p\rightarrow+\infty}{=}\frac{\ln(2p)}{2p}-\frac{1}{2p}\left(\ln(2p)+\frac{1}{2p}+o\left(\frac{1}{p}\right)\right)\left(1-\frac{1}{2p}+o\left(\frac{1}{p}\right)\right)\\
 &\underset{p\rightarrow+\infty}{=}\frac{\ln(2p)}{4p^2}+o\left(\frac{\ln p}{p^2}\right)\underset{p\rightarrow+\infty}{=}\frac{\ln p+\ln2}{4p^2}+o\left(\frac{\ln p}{p^2}\right)\\
 &\underset{p\rightarrow+\infty}{\sim}\frac{\ln p}{4p^2}.
\end{align*}   

et donc $R_{2k-1}\underset{k\rightarrow+\infty}{\sim}\frac{1}{4}\sum_{p=k}^{+\infty}\frac{\ln p}{p^2}$.

Cherchons maintenant un équivalent simple de $\frac{\ln p}{p^2}$  de la forme $v_p - v_{p+1}$. 

Soit $v_p=\frac{\ln p}{p}-\frac{\ln(p+1)}{p+1}$ (suggéré par $\left(\frac{\ln x}{x}\right)'=\frac{1-\ln x}{x^2}\underset{x\rightarrow+\infty}{\sim}-\frac{\ln x}{x^2}$). Alors 

\begin{align*}\ensuremath
v_p- v_{p+1}&=\frac{\ln p}{p}-\frac{1}{p}\left(\ln p+\ln\left(1+\frac{1}{p}\right)\right)\left(1+\frac{1}{p}\right)^{-1}\underset{p\rightarrow+\infty}{=}\frac{\ln p}{p}-\frac{1}{p}\left(\ln p+\frac{1}{p}+o\left(\frac{1}{p}\right)\right)\left(1-\frac{1}{p}+o\left(\frac{1}{p}\right)\right)\\
 &\underset{p\rightarrow+\infty}{\sim}\frac{\ln p}{p^2}.
\end{align*}

D'après la règle de l'équivalence des restes de séries à termes positifs convergentes, $R_{2k-1}\underset{k\rightarrow+\infty}{\sim}\frac{1}{4}\sum_{p=k}^{+\infty}\left(\frac{\ln p}{p}-\frac{\ln(p+1)}{p+1}\right)=\frac{\ln k}{4k}$ (série télescopique). 

Puis, $R_{2k}=R_{2k-1}-\frac{\ln(2k)}{2k}\underset{k\rightarrow+\infty}{\sim}\frac{\ln k}{4k}-\frac{\ln(2k)}{2k}+o\left(\frac{\ln k}{k}\right)\underset{k\rightarrow+\infty}{\sim}\frac{\ln k}{4k}-\frac{\ln k}{2k}+o\left(\frac{\ln k}{k}\right)\underset{k\rightarrow+\infty}{\sim}-\frac{\ln k}{4k}+o\left(\frac{\ln k}{k}\right)$.

En résumé, $R_{2k-1}\underset{k\rightarrow+\infty}{\sim}\frac{\ln k}{4k}$ et $R_{2k}\underset{k\rightarrow+\infty}{\sim}-\frac{\ln k}{4k}$.

On peut unifier : $R_{2k-1}\underset{k\rightarrow+\infty}{\sim}\frac{\ln k}{4k}\underset{k\rightarrow+\infty}{\sim}\frac{\ln(2k-1)}{2(2k-1)}$   et $R_{2k}\underset{k\rightarrow+\infty}{\sim}-\frac{\ln k}{4k}\underset{k\rightarrow+\infty}{\sim}-\frac{\ln(2k)}{2(2k)}$. Finalement,

\begin{center}
\shadowbox{
$\sum_{p=n+1}^{+\infty}(-1)^p\frac{\ln p}{p}\underset{n\rightarrow+\infty}{\sim}(-1)^{n-1}\frac{\ln n}{2n}$.
}
\end{center}

\item  $\sum_{}^{}n^n$ est une série à termes positifs grossièrement divergente.

\textbf{1 ère solution.}

$0< n^n\underset{n\rightarrow+\infty}{\sim} n^n - (n-1)^{n-1}$ car $\frac{n^n - (n-1)^{n-1}}{n^n}=  1-\frac{1}{n-1}\left(1-\frac{1}{n}\right)^{n}\underset{n\rightarrow+\infty}{=}1-\frac{1}{ne}+o\left(\frac{1}{n}\right)\underset{n\rightarrow+\infty}{\rightarrow}1$.

D'après la règle de l'équivalence des sommes partielles de séries à termes positifs divergentes,

\begin{center}
$\sum_{p=1}^{n}p^p\underset{n\rightarrow+\infty}{\sim}\sum_{p=2}^{n}p^p\underset{n\rightarrow+\infty}{\sim}\sum_{p=2}^{n}(p^p-(p-1)^{p-1}) = n^n-1\underset{n\rightarrow+\infty}{\sim}n^n$.
\end{center}

(La somme est équivalente à son dermier terme.)

\textbf{2 ème solution.} Pour $n\geqslant3$, $0\leqslant\frac{1}{n^n}\sum_{p=1}^{n-2}p^p\leqslant\frac{1}{n^n}\times(n-2)(n-2)^{n-2}\leqslant \frac{n^{n-1}}{n^n}=\frac{1}{n}$. Donc $\frac{1}{n^n}\sum_{p=1}^{n-2}p^p$. On en déduit que $\frac{1}{n^n}\sum_{p=1}^{n}p^p=1+\frac{(n-1)^{n-1}}{n^n}+\frac{1}{n^n}\sum_{p=1}^{n-2}p^p\underset{n\rightarrow+\infty}{=}1+o(1)+o(1)=1+o(1)$.

\begin{center}
\shadowbox{
$\sum_{p=1}^{n}p^p\underset{n\rightarrow+\infty}{\sim}n^n$.
}
\end{center}
\end{enumerate}
\fincorrection
\correction{005709}
Soit $p\in\Nn^*$. Pour $n\in\Nn^*\setminus\{p\}$, $\frac{1}{n^2-p^2}=\frac{1}{2p}\left(\frac{1}{n-p}-\frac{1}{n+p}\right)$. Donc pour $N > p$,

\begin{align*}
\sum_{1\leqslant n\leqslant N,\;n\neq p}^{}\frac{1}{n^2-p^2}&=\frac{1}{2p}\sum_{1\leqslant n\leqslant N,\;n\neq p}^{}\left(\frac{1}{n-p}-\frac{1}{n+p}\right)=\frac{1}{2p}\left(\sum_{1-p\leqslant k\leqslant N-p,\;k\neq 0}^{}\frac{1}{k}-\sum_{p+1\leqslant k\leqslant N+p,\;k\neq 2p}^{}\frac{1}{k}\right)\\
 &=\frac{1}{2p}\left(-\sum_{k=1}^{p-1}\frac{1}{k}+\sum_{k=1}^{N-p}\frac{1}{k}-\sum_{k=1}^{N+p}\frac{1}{k}+\frac{1}{2p}+\sum_{k=1}^{p}\frac{1}{k}\right)=\frac{1}{2p}\left(\frac{3}{2p}-\sum_{k=N-p+1}^{N+p}\frac{1}{k}\right)
\end{align*}

Maintenant,  $\sum_{k=N-p+1}^{N+p}\frac{1}{k}=\frac{1}{N-p+1}+\ldots+\frac{1}{N+p}$ est une somme de $2p-1$ termes tendant vers $0$ quand $N$ tend vers $+\infty$. Puisque $2p-1$ est constant quand $N$ varie, $\lim_{N \rightarrow +\infty}\sum_{k=N-p+1}^{N+p}\frac{1}{k}=0$ et donc

\begin{center}
$\sum_{n\in\Nn^*,\;n\neq p}^{}\frac{1}{n^2-p^2}=\frac{1}{2p}\times\frac{3}{2p}=\frac{3}{4p^2}$ puis $\sum_{p\in\Nn^*}^{}\left(\sum_{n\in\Nn^*,\;n\neq p}^{}\frac{1}{n^2-p^2}\right)=\sum_{p=1}^{+\infty}\frac{3}{4p^2}=\frac{\pi^2}{8}$.
\end{center}

Pour $n\in\Nn^*$ donné, on a aussi $\sum_{p\in\Nn^*,\;p\neq n}^{}\frac{1}{n^2-p^2}=-\sum_{p\in\Nn^*,\;p\neq n}^{}\frac{1}{p^2-n^2}=-\frac{3}{4n^2}$ et donc

\begin{center}
$\sum_{n\in\Nn^*}^{}\left(\sum_{p\in\Nn^*,\;p\neq n}^{}\frac{1}{n^2-p^2}\right)=-\frac{\pi^2}{8}$.
\end{center}

On en déduit que la suite double $\left(\frac{1}{n^2-p^2}\right)_{(n,p)\in(\Nn^*)^2,\;n\neq p}$ n'est pas sommable.
\fincorrection
\correction{005710}
La suite $\left((-1)^n\frac{1}{3n+1}\right)_{n\in\Nn}$ est alternée en signe et sa valeur absolue tend vers $0$ en décroissant. Donc la série de terme général $(-1)^n\frac{1}{3n+1}$, $n\geqslant 1$, converge en vertu du critère spécial aux séries alternées.

Soit $n\in\Nn$.

\begin{center}
$\sum_{k=0}^{n}\frac{(-1)^k}{3k+1}=\sum_{k=0}^{n}(-1)^k\int_{0}^{1}t^{3k}\;dt=\int_{0}^{1}\frac{1-(-t^3)^{n+1}}{1-(-t^3)}\;dt=\int_{0}^{1}\frac{1}{1+t^3}\;dt+(-1)^n\int_{0}^{1}\frac{t^{3n+3}}{1+t^3}\;dt$.
\end{center}

Mais $\left|(-1)^n\int_{0}^{1}\frac{t^{3n+3}}{1+t^3}\;dt\right|=\int_{0}^{1}\frac{t^{3n+3}}{1+t^3}\;dt\leqslant\int_{0}^{1}t^{3n+3}\;dt=\frac{1}{3n+4}$. On en déduit que $(-1)^n\int_{0}^{1}\frac{t^{3n+3}}{1+t^3}\;dt$ tend vers $0$ quand $n$ tend vers $+\infty$ et donc que

\begin{center}
$\sum_{n=0}^{+\infty}\frac{(-1)^n}{3n+1}=\int_{0}^{1}\frac{1}{1+t^3}\;dt$.
\end{center}

Calculons cette dernière intégrale.

\begin{align*}\ensuremath
\frac{1}{X^3+1}&=\frac{1}{(X+1)(X+j)(X+j^2)}=\frac{1}{3}\left(\frac{1}{X+1}+\frac{j}{X+j}+\frac{j^2}{X+j^2}\right)=\frac{1}{3}\left(\frac{1}{X+1}+\frac{-X+2}{X^2-X+1}\right)\\
 &\frac{1}{3}\left(\frac{1}{X+1}-\frac{1}{2}\frac{2X-1}{X^2-X+1}+
 \frac{3}{2}\frac{1}{\left(X-\frac{1}{2}\right)^2+\left(\frac{\sqrt{3}}{2}\right)^2}
 \right).
\end{align*}

 
Donc, 

\begin{center}
$\sum_{n=0}^{+\infty}\frac{(-1)^n}{3n+1}=\frac{1}{3}\left[\ln(t+1)-\frac{1}{2}\ln(t^2-t+1)+\sqrt{3}\Arctan\left(\frac{2t-1}{\sqrt{3}}\right)\right]_0^1=\frac{1}{3}\left(\ln2+\sqrt{3}\left(\frac{\pi}{6}-\left(-\frac{\pi}{6}\right)\right)\right)=\frac{3\ln2+\pi\sqrt{3}}{9}$.
\end{center}

\begin{center}
\shadowbox{
$\sum_{n=0}^{+\infty}\frac{(-1)^n}{3n+1}=\frac{3\ln2+\pi\sqrt{3}}{9}$.
}
\end{center}
\fincorrection
\correction{005711}
Pour tout entier $n\geqslant 2$, on a $nv_n-(n-1)v_{n-1}=u_n$ ce qui reste vrai pour $n = 1$ si on pose de plus $v_0 = 0$. Par suite, pour $n\in\Nn^*$

\begin{align*}\ensuremath
v_n^2 -2u_nv_n&=v_n^2 - 2(nv_n - (n-1)v_{n-1})v_n = -(2n-1) v_n^2 + 2(n-1)v_{n-1}v_n\\
 &\leqslant -(2n-1) v_n^2 +(n-1)(v_{n-1}2+v_n^2) =(n-1)v_{n-1}^2 - n v_n^2.
\end{align*}
	
        
Mais alors, pour $N\in\Nn^*$,

\begin{center}
$\sum_{n=1}^{N}(v_n^2 - 2u_nv_n)\leqslant\sum_{n=1}^{N}((n-1)v_{n-1}^2 - n v_n^2)= - nv_n^2\leqslant0$.
\end{center}

Par suite, 

\begin{center}
$\sum_{n=1}^{N}v_n^2\leqslant\sum_{n=1}^{N}2u_nv_n\leqslant2\left(\sum_{n=1}^{N}u_n^2\right)^{1/2}\left(\sum_{n=1}^{N}v_n^2\right)^{1/2}\;$ (inégalité de \textsc{Cauchy}-\textsc{Schwarz}).
\end{center}

Si $\left(\sum_{n=1}^{N}v_n^2\right)^{1/2}>0$, on obtient après simplification par $\left(\sum_{n=1}^{N}v_n^2\right)^{1/2}$ puis élévation au carré

\begin{center}
$\sum_{n=1}^{N}v_n^2\leqslant4\sum_{n=1}^{N}u_n^2$,
\end{center}

cette inégalité restant claire si $\left(\sum_{n=1}^{N}v_n^2\right)^{1/2}=0$. Finalement,

\begin{center}
$\sum_{n=1}^{N}v_n^2\leqslant4\sum_{n=1}^{N}u_n^2\leqslant4\sum_{n=1}^{+\infty}u_n^2$.
\end{center}

La suite des sommes partielles de la série de terme général $v_n^2(\geqslant0)$ est majorée. Donc la série de terme général $v_n^2$ converge et de plus, quand $N$ tend vers l'infini, on obtient

\begin{center}
$\sum_{n=1}^{+\infty}v_n^2\leqslant 4\sum_{n=1}^{+\infty}u_n^2$.
\end{center} 
\fincorrection
\correction{005712}
Soit $n\in\Nn$.

\begin{align*}\ensuremath
u_n&=\frac{\pi}{4}-\sum_{k=0}^{n}\frac{(-1)^k}{2k+1}=\int_{0}^{1}\frac{1}{1+t^2}\;dt-\sum_{k=0}^{n}(-1)^k\int_{0}^{1}t^{2k}\;dt=\int_{0}^{1}\frac{1}{1+t^2}\;dt-\int_{0}^{1}\frac{1-(-t^2)^{n+1}}{1+t^2}\;dt\\
 &= (-1)^{n+1}\int_{0}^{1}\frac{t^{2n+2}}{1+t^2}\;dt.
\end{align*}

Par suite, pour $N\in\Nn$,

\begin{align*}\ensuremath
\sum_{n=0}^{N}u_n=\int_{0}^{1}\sum_{n=0}^{N}\frac{(-t^2)^{n+1}}{1+t^2}\;dt  =\int_{0}^{1}(-t^2)\frac{1-(-t^2)^{N+1}}{(1+t^2)^2}\;dt=-\int_{0}^{1}\frac{t^2}{(1+t^2)^2}\;dt+(-1)^{N+1}\int_{0}^{1}\frac{t^{2N+2}}{(1+t^2)^2}\;dt.
\end{align*}

Or $\left|(-1)^{N+1}\int_{0}^{1}\frac{t^{2N+2}}{(1+t^2)^2}\;dt\right|=\int_{0}^{1}\frac{t^{2N+2}}{(1+t^2)^2}\;dt\leqslant\int_{0}^{1}t^{2N+2}\;dt=\frac{1}{2N+3}$. Comme $\frac{1}{2N+3}$ tend vers $0$ quand $N$ tend vers $+\infty$, il en est de même de $(-1)^{N+1}\int_{0}^{1}\frac{t^{2N+2}}{(1+t^2)^2}\;dt$. On en déduit que la série de terme général $u_n$, $n\in\Nn$, converge et de plus

\begin{align*}\ensuremath
\sum_{n=0}^{+\infty}u_n&=-\int_{0}^{1}\frac{t^2}{(1+t^2)^2}\;dt=\int_{0}^{1}\frac{t}{2}\times\frac{-2t}{(1+t^2)^2}\;dt\\
 &=\left[\frac{t}{2}\times\frac{1}{1+t^2}\right]_0^1-\int_{0}^{1}\frac{1}{2}\times\frac{1}{1+t^2}\;dt=\frac{1}{4}-\frac{\pi}{8}.
\end{align*}

\begin{center}
\shadowbox{
$\sum_{n=0}^{+\infty}\left(\frac{\pi}{4}-\sum_{k=0}^{n}\frac{(-1)^k}{2k+1}\right)=\frac{1}{4}-\frac{\pi}{8}$.
}
\end{center}
\fincorrection
\correction{000639}
\begin{enumerate}
    \item On a pour tout $x,y\in\R$ $|x-y|\geq \big| |x|-|y|\big|$
(c'est la deuxi\`eme formulation de l'in\'egalit\'e triangulaire).
Donc pour tout $x\in I$ :$ \big| |f(x)|-|f(a)| \big| \leq |f(x)-f(a)| $.
L'implication annonc\'ee r\'esulte alors imm\'ediatement de la
d\'efinition de l'assertion $\lim_{x\to a} f(x)=f(a). $
    \item  Si $f,g$ sont continues
alors $\alpha f+\beta g$ est continue sur $I$, pour tout
$\alpha,\beta\in\R$. Donc les fonctions $f+g$ et $f-g$ sont
continues sur $I$. L'implication de $1.$ prouve alors que $|f-g|$
est continue sur $I$, et finalement on peut
conclure :\\
La fonction $\sup (f,g) = \frac{1}{2}(f+g+|f-g|)$ est continue sur
$I$.
\end{enumerate}
\fincorrection
\nocorrection
\nocorrection
\correction{000642}
\begin{enumerate}
  \item $g(a) = f(\frac{a+b}{2})-f(a)$ et
$g(\frac{a+b}{2}) = f(b) - f(\frac{a+b}{2})$.
Comme $f(a) = f(b)$ alors nous obtenons que $g(a) = -g(\frac{a+b}{2})$.
Donc ou bien $g(a) \leq 0$ et $g(\frac{a+b}{2}) \geq 0$
ou bien $g(a) \geq 0$ et $g(\frac{a+b}{2}) \leq 0$.
D'après le théorème des valeurs intermédiaires, $g$
s'annule en $c$ pour un $c$ entre $a$ et $\frac{a+b}{2}$.
  \item Notons $t$ le temps (en heure) et $d(t)$ la distance parcourue (en km) entre les instants $0$ et $t$.
Nous supposons que la fonction $t \mapsto d(t)$ est continue.
Soit $f(t) = d(t) - 4t$. Alors $f(0) = 0$ et par hypothèse $f(1) = 0$.
Appliquons la question précédente avec $a=0$, $b=1$. 
Il existe $c\in [0,\frac 12]$ tel que $g(c) = 0$, c'est-à-dire
$f(c+\frac12)= f(c)$. Donc $d(c+\frac 12)-d(c) = 4(c+\frac12)-4c = 2$.
Donc entre $c$ et $c+\frac 12$, (soit 1/2 heure), la personne parcourt
exactement $2$ km.
\end{enumerate}
\fincorrection
\correction{000643}
Il existe $x < 0$ tel que $f(x) <0$ et $y>0$ tel que $f(y) > 0$,
d'après le théorème des valeurs intermédiaires,
il existe $z \in ]x,y[$ tel que $f(z) = 0$.
Donc $f$ s'annule. 
Les polynômes de degré impair vérifient les propriétés des limites, donc
s'annulent. Ceci est faux, en général,  pour les polynômes de degré pair, par exemple regardez $f(x) = x^2+1$.
\fincorrection
\nocorrection
\correction{000645}
Comme $f(x)^2 = 1$ alors $f(x) = \pm 1$. 
Attention ! Cela ne veut pas dire que 
la fonction est constante \'egale \`a $1$ ou $-1$.
Supposons, par exemple, qu'il existe $x$ tel que $f(x)=+1$.
Montrons que $f$ est constante \'egale \`a $+1$.
S'il existe $y \not= x$ tel que $f(y) = -1$
alors $f$ est positive en $x$, n\'egative en $y$
et continue sur $I$. Donc, par le th\'eor\`eme des valeurs interm\'ediaires,
il existe $z$ entre $x$ et $y$ tel que $f(z) = 0$, ce qui contredit
$f(z)^2 = 1$. Donc $f$ est constante \'egale \`a $+1$.
\fincorrection
\correction{000646}
Notons $\ell$ la limite de $f$ en $+\infty$:
$$\forall \epsilon > 0 \quad \exists A \in \Rr
\quad x>A \Rightarrow \ell - \epsilon \leq f(x) \leq \ell + \epsilon.$$
Fixons $\epsilon = +1$, nous obtenons un $A$ correspondant tel que
pour $x>A$, $\ell - 1 \leq f(x) \leq \ell +1$. Nous venons de montrer
que $f$ est born\'ee ``\`a l'infini''.
La fonction $f$ est continue sur l'intervalle ferm\'e born\'e $[0,A]$,
donc $f$ est born\'ee sur cet intervalle: il existe $m,M$ tels que
pour tout $x\in [0,A]$, $m \leq f(x) \leq M$.
En prenant $M' = \max (M,\ell+1)$, et $m' = \min(m,\ell-1)$ nous avons que pour tout $x\in \Rr$,
$m' \leq f(x) \leq M'$. Donc $f$ est born\'ee sur $\Rr$.

La fonction n'atteint pas n\'ecessairement ses bornes: regardez
$f(x) = \frac{1}{1+x}$.
\fincorrection
\nocorrection
\nocorrection
\nocorrection
\nocorrection
\nocorrection
\nocorrection
\correction{000653}
\begin{enumerate}
  \item  Si $f(0) = 0$ et c'est fini, on a trouver le point fixe !
Sinon $f(0)$ n'est pas nul.  Donc $f(0) > 0$ et $0 \in E$. Donc $E$ n'est pas vide. 
  \item Maintenant $E$ est un partie de $[0,1]$ non vide
donc $\sup E$ existe et est fini. Notons $c = \sup E \in [0,1]$.
Nous allons montrer que $c$ est un point fixe.
  \item Nous approchons ici $c = \sup E$ par des éléments de $E$ :
 Soit $(x_n)$ une suite de $E$ telle que $x_n \rightarrow c$
et $x_n \leq c$. Une telle suite existe d'apr\`es les propri\'et\'es de
$c= \sup E$. Comme $x_n \in E$ alors
$x_n < f(x_n)$. Et comme $f$ est croissante $f(x_n) \leq f(c)$.
Donc pour tout $n$, $x_n < f(c)$ ; comme $x_n \rightarrow c$ alors \`a la limite nous avons $c \leq f(c)$.
  \item Si $c=1$ alors $f(1)=1$ et nous avons notre point fixe. Sinon, nous utilisons maintenant le fait que 
les élements supérieurs à $\sup E$ ne sont pas dans $E$ :
Soit $(t_n)$ une suite telle que 
$t_n \rightarrow c$, $t_n \geq c$
et telle que $f(t_n) \leq t_n$. Une telle suite existe
car sinon $c$ ne serait pas \'egal \`a $\sup E$.
Nous avons $f(c) \leq f(t_n) \leq t_n$ et donc \`a la limite
$f(c) \leq c$.

Nous concluons donc que $c \leq f(c) \leq c$, donc $f(c) = c$
et $c$ est  un point fixe de $f$.
\end{enumerate}
\fincorrection
\nocorrection
\nocorrection
\nocorrection
\nocorrection
\nocorrection
\nocorrection
\correction{000660}
\begin{enumerate}
  \item Soit $x\in[0,1]$ et $y=f(x) \in [0,1]$.
Alors $f(y)=y$ car $f(f(x)) = f(x)$. Donc $E_f \not= \varnothing$.
Nous venons de montrer que $I = f([0,1])$ est inclus dans $E_f$.
  \item Montrons r\'eciproquement $E_f$ est inclus dans $I$.
Soit $x\in[0,1]$ tel que $f(x)=x$ alors $x \in I=f([0,1])$
(car $x = f(x)$ !). 
Ainsi $E_f = I = f([0,1])$. Mais l'image de l'intervalle $[0,1]$ par la fonction continue $f$ est un intervalle donc 
$E_f$ est un intervalle.
  \item La réciproque est vraie : une fonction continue pour laquelle $E_f = f([0,1])$  v\'erifie aussi $f\circ f = f$.
En effet pour $x \in [0,1]$ et $y = f(x)$ alors $y \in f([0,1])$ donc $y \in E_f$. Donc $f(y)=y$, autrement dit $f(f(x))=f(x)$.

Les fonctions continues qui vérifient $f \circ = f$ sont donc exactement les fonctions continues telles que $E_f = f([0,1])$.
Pour une telle fonction si l'on note $[a,b] = E_f$ alors 
$f$ est définie sur $[0,a]$ par n'importe qu'elle fonction continue prenant ses valeurs entre $a$ et $b$,
et valant $a$ en $a$ : $f([0,a]) \subset [a,b]$ et $f(a)=a$.
Elle est ensuite définie par l'identité sur $[a,b]$ : pour tout $x\in [a,b]$, $f(x)=x$.
Et enfin sur $[b,1]$ elle est définie par n'importe quelle fonction continue prenant ses valeurs entre $a$ et $b$,
et valant $b$ en $b$ : $f([b,1]) \subset [a,b]$ et $f(b)=b$.
\end{enumerate}
\fincorrection
\nocorrection
\correction{000662}
Non, par exemple $f : \Rr \longrightarrow \Rr$.
Avec $f(x) = \sin\frac1x$ pour $x \not=0$ et
$f(0) = 0$. $f$ n'est pas continue (en $0$), mais pour tout
$a,b$ et pour tout $y \in [f(a),f(b)]$, il existe $x \in [a,b]$ tel que $y=f(x)$.
\fincorrection
\nocorrection
\nocorrection
\nocorrection
\nocorrection
\nocorrection
\nocorrection
\correction{000669}
\begin{enumerate}
\item Pour tout $x \in ]a,b[$, on a $x \in [a,b]$ donc $f(x) \leq \sup_{a\leq
    t \leq b} f(t)$. Par cons\'equent $\sup_{a\leq t \leq b} f(t)$ est un
  majorant de $f$ sur l'intervalle $]a,b[$, donc il est plus grand que le plus
  petit des majorants : $ \sup_{a<x<b} f(x) \leq  \sup_{a\leq t \leq b} f(t)$.
\item $f$ est continue sur un intervalle ferm\'e et born\'e, donc elle est born\'ee
  et elle atteint ses bornes. Soit $x_0$ le r\'eel o\`u le maximum est atteint : 
$f(x_0)=\sup_{a\leq x \leq b} f(x)$.
  \begin{itemize}
  \item si $x_0=a$, consid\'erons la suite $a_n=a+1/n$. Pour $n\geq
    \frac{1}{b-a}$ on a $a_n \in [a,b]$, donc on peut consid\'erer la suite
    $(f(a_n))_{n\geq \frac{1}{b-a}}$. Or $a_n$ tend vers $a$ quand $n$ tend
    vers $+\infty$, et comme $f$ est continue, ceci implique que $f(a_n)$ tend
    vers $f(a)$ quand $n$ tend    vers $+\infty$. Donc $\forall \epsilon>0,
    \exists n \in \N, f(x_0)-\epsilon \leq f(a_n) \leq f(x_0)$, ce qui
    implique que $f(x_0)=\sup_{a< x< b}    f(x)$.
  \item si $x_0=b$ on obtient le r\'esultat de mani\`ere identique en consid\'erant
    la suite $b_n=b-1/n$.
  \item si $a<x_0<b$ : $f(x_0)$ est major\'e par le sup de $f$ sur $]a,b[$, donc
$$f(x_0) \leq \sup_{a<x<b} f(x) \leq  \sup_{a\leq x \leq b} f(x)=f(x_0)$$
donc $f(x_0) = \sup_{a<x<b} f(x)$.
  \end{itemize}
\item Avec la fonction $g$, on a $\sup_{0<x<1} g(x)=0$ car pour chaque $x \in
  ]0,1[$, $g(x)=0$, et $\sup_{0\leq x \leq 1} g(x)=1$ car $g(0)=0$ et
  $g(1)=1$. La propri\'et\'e d\'emontr\'ee pr\'ec\'edemment n'est pas vraie dans notre
  cas, car la fonction $g$ ne remplit pas la condition essentielle d'\^etre continue.
\end{enumerate}
\fincorrection
\nocorrection
\nocorrection
\nocorrection
\nocorrection
\nocorrection
\nocorrection
\nocorrection
\nocorrection
\nocorrection
\nocorrection
\correction{003856}
\begin{enumerate}
  \item $f(x)-x$ change de signe entre $0$ et $1$.
  \item Sinon $f-g$ est de signe constant, par exemple positif.
Si $a$ est le plus grand point fixe de~$f$ alors $g(a)>a$ et $g(a)$ est
aussi point fixe de~$f$, absurde.
\end{enumerate}
 \fincorrection
\nocorrection
\correction{003858}
En posant $b=f(a)$ on a $(f(a)-a) + (f(b)-b) = 0$
donc $x \mapsto f(x)-x$ s'annule entre $a$ et $b$.
De même, s'il existe $k\in\N^*$ et $a\in\R$ tels que $f^k(a)=a$
alors $(f(a)-a) + (f^2(a)-f(a)) + \dots + (f^{k}(a)-f^{k-1}(a)) = 0$
donc $f(x)-x$ s'annule entre $\min(a,f(a),\dots,f^{k-1}(a))$ et
$\max(a,f(a),\dots,f^{k-1}(a))$.
\fincorrection
\nocorrection
\nocorrection
\correction{003861}
\begin{enumerate}
  \item Soit $a\in\R$ tel que $f$ est discontinue en $a$. Il existe une suite
         $a_n$ telle que $a_n \to a$ et $|f(a_n)-f(a)| \ge \varepsilon$.
         Alors $f(a)\pm\varepsilon$ a une infinité d'antécédents.
  \item 
\end{enumerate}
 \fincorrection
\nocorrection
\nocorrection
\nocorrection
\nocorrection
\nocorrection
\nocorrection
\nocorrection
\correction{003869}
Si $f$ n'est pas identiquement nulle, alors $f(0) = \pm1$ et $f$ est paire,
de signe constant.

Par récurrence, $\forall\ p \in \N,\ f(px) = \pm f^{p^2}(x)$
$ \Rightarrow $ par densité, $f(x) = \pm\lambda^{x^2}$.
\fincorrection
\correction{003870}
$\omega(\delta) \to 0 \Leftrightarrow f$ (lorsque $\delta\to0^+$) est uniformément continue.

Continuité en $\delta>0$~: on remarque que $\omega$ est croissante
donc $\omega(\delta^-)$ et $\omega(\delta^+)$ existent et
encadrent $\omega(\delta)$.
Si $\delta_n \to \delta^+$ (lorsque $n\to\infty$), soient $x_n,y_n$ tels que
$\omega(\delta_n) = |f(x_n)-f(y_n)|$ et $|x_n-y_n|\le\delta_n$.
On extrait de $(x_n,y_n)$ une suite convergente vers~$(x,y)$ avec
$|x-y|\le\delta$ et $|f(x)-f(y)| = \omega(\delta^+)$ d'où
$\omega(\delta^+) \le \omega(\delta)$ puis
$\omega(\delta^+) = \omega(\delta)$.

On a aussi $\omega(\delta)= \sup\{|f(x)-f(y)|\text{ tq } |x-y|<\delta\}
\le \omega(\delta^-)$ d'où $\omega(\delta^-) = \omega(\delta)$.
\fincorrection
\correction{003871}
$f$ admet des points fixes car l'application $x \mapsto f(x)-x$ change
de signe entre $0$ et~$1$. Si $E$ est l'ensemble des points fixes de~$f$ alors $E$ est
stable par~$g$ donc $f-g$ a des signes opposés en $\min(E)$ et $\max(E)$.
\fincorrection
\correction{003872}
\begin{enumerate}
  \item Non. Si $\varphi$ est lipschitzienne alors $\varphi(x) = O(\|x\|)$
    lorsque $\|x\|$ tend vers l'infini, donc toute fonction $f$ à décroissance
    suffisament rapide vers $-\infty$ n'est pas minorable par une fonction
    lipschitzienne. Contre-exemple explicite~: $f(x) = -\|x\|^2$.
  \item CNS~: $x \mapsto f(x)+k\|x\|$ est minorée.
  \item On pose $\varphi(x) = \sup\{g(x),\ g\ k$-lipschitzienne minorant $f\}$.
    Il suffit de vérifier que~$\varphi$ est $k$-lipschitzienne, ce sera alors
    la plus grande fonction $k$-lipschitzienne minorant~$f$. Soient~$x,y\in\R^n$,
    $\varepsilon>0$ et~$g_x,g_y$ des fonctions $k$-lipschitziennes minorant~$f$ telles
    que $g_y(x)\le\varphi(x)\le g_x(x)+\varepsilon$ et $g_x(y)\le\varphi(y)\le g_y(y)+\varepsilon$.
    On a~: 
\begin{align*}
    \varphi(y)&\le g_y(y)+\varepsilon\le g_y(x)+k\|x-y\|+\varepsilon\le\varphi(x)+k\|x-y\|+\varepsilon,\\
    \varphi(y)&\ge g_x(y)\ge g_x(x)-k\|x-y\|\ge\varphi(x)-k\|x-y\|-\varepsilon.\\
\end{align*}
    Donc $|\varphi(x)-\varphi(y)|\le k\|x-y\|+\varepsilon$ et on fait tendre $\varepsilon$
    vers~$0^+$.
\end{enumerate}
 \fincorrection
\correction{003873}
Soit pour $x>0$, $\ell(x) = \lim_{n\to+\infty} f(nx)$. On a $\ell(kx) = \ell(x)$
pour tout~$k\in\N^*$ d'où aussi pour tout~$k\in\Q^{+*}$. Montrons alors que
 $f(x) \to \ell(1)$ lorsque $x\to+\infty$~:
soit $\varepsilon>0$ et $\delta$ associé dans la définition de l'uniforme
 continuité de~$f$. On choisit un rationnel $\alpha\in{]0,\delta[}$ et un
 entier $N$ tel que $|f(n\alpha)-\ell(1)|=|f(n\alpha)-\ell(\alpha)|\le\varepsilon$ pour tout~$n\ge N$.
Alors $|f(x)-\ell(1)|\le2\varepsilon$ pour tout~$x\ge N\alpha$.
\fincorrection
\correction{005382}
Il existe $a>0$ tel que $f$ est définie et continue sur $[a,+\infty[$.

1er cas. Supposons que $\ell$ est réel. Soit $\varepsilon>0$. 

$$\exists A_1\geq a/\;\forall X\in[a,+\infty[,\;(X\geq A_1\Rightarrow \ell-\frac{\varepsilon}{2}<f(X+1)-f(X)<\ell+\frac{\varepsilon}{2}).$$

Soit $X\geq A_1$ et $n\in\Nn^*$. On a~:
$$\sum_{k=0}^{n-1}(\ell-\frac{\varepsilon}{2})<\sum_{k=0}^{n-1}(f(X+k+1)-f(X+k))=f(X+n)-f(X)<\sum_{k=0}^{n-1}(\ell+\frac{\varepsilon}{2}),$$ et on a donc montré que 

$$\forall\varepsilon>0,\;\exists A_1\geq a/\;\forall X\geq A_1,\;\forall n\in\Nn^N*,\;n(\ell-\frac{\varepsilon}{2})<f(X+n)-f(X)<n(\ell+\frac{\varepsilon}{2}).$$ 

Soit de nouveau $\varepsilon>0$. Soit ensuite $x\geq A_1+1$ puis $n=E(x-A_1)\in\Nn^*$ puis $X=x-n$.

On a $X=x-E(x-A_1)\geq x-(x-A_1)=A_1$ et donc $n(\ell-\frac{\varepsilon}{2})<f(x)-f(x-n)<n(l+\frac{\varepsilon}{2})$ ou encore

$$\frac{f(x-n)}{x}+\frac{n}{x}(\ell-\frac{\varepsilon}{2})<\frac{f(x)}{x}<\frac{f(x-n)}{x}+\frac{n}{x}(\ell+\frac{\varepsilon}{2}).$$

Ensuite,

$$1-\frac{A_1+1}{x}=\frac{x-A_1-1}{x}\leq\frac{n}{x}=\frac{E(x-A_1)}{x}\leq\frac{x-A_1}{x}=1-\frac{A_1}{x},$$ et comme $1-\frac{A_1+1}{x}$ et $1-\frac{A_1}{x}$ tendent vers $1$ quand $x$ tend vers $+\infty$, on en déduit que $\frac{n}{x}$ tend vers $1$ quand $x$ tend vers $+\infty$.

Puis, puisque $f$ est continue sur le segment $[A_1,A_1+1]$, $f$ est bornée sur ce segment. Or $n\leq x-A_1<n+1$ s'écrit encore $A_1\leq x-n <A_1+1$ et donc, en posant $M=\mbox{sup}\{|f(t)|,\;t\in[A_1,A_1+1]\}$, on a $\left|\frac{x-n)}{x}\right|\leq\frac{M}{x}$ qui tend vers $0$ quand $x$ tend vers $+\infty$. En résumé, $\frac{f(x-n)}{x}+\frac{n}{x}(\ell-\frac{\varepsilon}{2})$ et $\frac{f(x-n)}{x}+\frac{n}{x}(\ell+\frac{\varepsilon}{2})$ tendent respectivement vers $\ell-\frac{\varepsilon}{2}$ et $\ell+\frac{\varepsilon}{2}$ quand $x$ tend vers $+\infty$. On peut donc trouver un réel $A_2\geq a$ tel que $x\geq A_2\Rightarrow\frac{f(x-n)}{x}+\frac{n}{x}(\ell+\frac{\varepsilon}{2})>(\ell-\frac{\varepsilon}{2})-\frac{\varepsilon}{2}=\ell-\varepsilon$ et un réel $A_3\geq a$ tel que $x\geq A_2\Rightarrow\frac{f(x-n)}{x}+\frac{n}{x}(\ell+\frac{\varepsilon}{2})<\ell+\varepsilon$.

Soit $A=\mbox{Max}(A_1,A_2,A_3)$ et $x\geq A$. On a $\ell-\varepsilon<\frac{f(x)}{x}<\ell+\varepsilon$. On a montré que $\forall\varepsilon>0,\;(\exists A\geq a/\;\forall x\geq A,\;\ell-\varepsilon<\frac{f(x)}{x}<\ell+\varepsilon$ et donc $\lim_{x\rightarrow +\infty}=\ell$.

2ème cas. Supposons $\ell=+\infty$ (si $\ell=-\infty$, remplacer $f$ par $-f$).

Soit $B>0$. $\exists A_1\geq a/\;\forall X\geq A_1,\;f(X+1)-f(X)\geq 2B$.

Pour $X\geq A_1$ et $n\in\Nn^*$, on a~:~$f(X+n)-f(X)=\sum_{k=0}^{n-1}(f(X+k+1)-f(X+k))\geq2nB$.

Soient $x\geq1+A_1$, $n=E(x-A_1)$ et $X=x-n$. On a $f(x)-f(x-n)\geq2nB$ et donc,

$$\frac{f(x)}{x}\geq\frac{f(x-n)}{x}+\frac{2nB}{x},$$

qui tend vers $2B$ quand $x$ tend vers $+\infty$ (démarche identique au 1er cas).

Donc $\exists A\geq A_1>a$ tel que $x\geq A\Rightarrow\frac{f(x-n)}{x}+\frac{2nB}{x}>B$.

Finalement~:~($\forall B>0,\;\exists A>a/\;(\forall x\geq A,\;\frac{f(x)}{x}>B$ et donc, $\lim_{x\rightarrow +\infty}\frac{f(x)}{x}=+\infty$.

\fincorrection
\correction{005383}
Pour $x\neq 0$, posons $g(x)=\frac{f(2x)-f(x)}{x}$. $f$ est définie sur un voisinage de $0$ et donc il existe $a>0$ tel que $]-a,a[\subset D_f$. Mais alors, $]-\frac{a}{2},\frac{a}{2}[\setminus\{0\}\subset D_g$.

Soit $x\in]-\frac{a}{2},\frac{a}{2}[\setminus\{0\}$ et $n\in\Nn^*$.

$$f(x)=\sum_{k=0}^{n-1}(f(\frac{x}{2^k})-f(\frac{x}{2^{k+1}}))+f(\frac{x}{2^n})=\sum_{k=0}^{n-1}\frac{x}{2^{k+1}}g(\frac{x}{2^{k+1}})+f(\frac{x}{2^n}).$$

Par suite, pour $x\in]-\frac{a}{2},\frac{a}{2}[\setminus\{0\}$ et $n\in\Nn^*$, on a~:

$$\left|\frac{f(x)}{x}\right|\leq\sum_{k=0}^{n-1}\frac{1}{2^{k+1}}\left|g(\frac{x}{2^{k+1}})\right|+\left|\frac{f(x/2^n)}{x}\right|.$$

Soit $\varepsilon>0$. Puisque par hypothèse, $g$ tend vers $0$ quand $x$ tend vers $0$,

$$\exists\alpha\in]0,\frac{a}{2}[/\;\forall X\in]-\alpha,\alpha[,\;|g(X)|<\frac{\varepsilon}{2}.$$

Or, pour $x\in]-\alpha,\alpha[\setminus\{0\}$ et pour $k$ dans $N*$, $\frac{x}{2^k}$ est dans $]-\alpha,\alpha[\setminus\{0\}$ et par suite,

$$\sum_{k=0}^{n-1}\frac{1}{2^{k+1}}\left|g(\frac{x}{2^{k+1}})\right|\leq\frac{\varepsilon}{2}\sum_{k=0}^{n-1}\frac{1}{2^{k+1}}=\frac{\varepsilon}{2}\frac{1}{2}\frac{1-\frac{1}{2^n}}{1-\frac{1}{2}}=\frac{\varepsilon}{2}(1-\frac{1}{2^n})<\frac{\varepsilon}{2},$$

et donc, $\left|\frac{f(x)}{x}\right|\leq\frac{\varepsilon}{2}+\left|\frac{f(x/2^n)}{x}\right|$. On a ainsi montré que 

$$\forall x\in]-\alpha,\alpha[\setminus\{0\},\;\forall n\in\Nn^*,\;\left|\frac{f(x)}{x}\right|\leq\frac{\varepsilon}{2}+\left|\frac{f(x/2^n)}{x}\right|.$$

Mais, à $x$ fixé, $\frac{f(x/2^n)}{x}$ tend vers $0$ quand $n$ tend vers $+\infty$. Donc, on peut choisir $n$ tel que  $\frac{f(x/2^n)}{x}<\frac{\varepsilon}{2}$ et on a alors  $\left|\frac{f(x)}{x}\right|<\frac{\varepsilon}{2}+\frac{\varepsilon}{2}=\varepsilon$. On a montré que 

$$\forall\varepsilon>0,\;\exists\alpha>0/\;(\forall x\in D_f,\;0<|x|<\alpha\Rightarrow\left|\frac{f(x)}{x}\right|<\varepsilon,$$ ce qui montre que ($f$ est dérivable en $0$ et que) $\lim_{x\rightarrow 0}\frac{f(x)}{x}=0$.

\fincorrection
\correction{005384}
$\mbox{Min}(f,g)=\frac{1}{2}(f+g-|f-g|)$ et $\mbox{Max}(f,g)=\frac{1}{2}(f-g+|f-g|)$ sont continues en $x_0$ en vertu de théorèmes généraux.
\fincorrection
\correction{005385}
Soit $(x,y)\in\Rr^2$ et $z\in A$. $|x-z|\leq|x-y|+|y-z|$. Or, $forall z\in A,\;|x-z|\geq d(x,A)$ et donc $d(x,A)-|x-y|$ est un minorant de $\{|y-z|,\;z\in A\}$. Par suite, $d(x,A)-|x-y|\leq d(y,A)$. On a montré que 

$$\forall(x,y)\in\Rr^2,\;d(x,A)-d(y,A)\leq|y-x|.$$

En échangeant les roles de $x$ et $y$, on a aussi montré que $\forall(x,y)\in\Rr^2,\;d(y,A)-d(x,A)\leq|y-x|$.

Finalement, $\forall(x,y)\in\Rr^2,\|f(y)-f(x)|\leq|y-x|$. Ainsi, $f$ est donc $1$-Lipschitzienne et en particulier continue sur $\Rr$.
\fincorrection
\correction{005386}
Soit $x_0\in\Rr\setminus\{5\}$. Pour $x\neq5$, 

$$|f(x)-f(x_0)|=\left|\frac{3x-1}{x-5}-\frac{3x_0-1}{x_0-5}\right|=\frac{14|x-x_0|}{|x-5|.|x_0-5|}.$$

Puis, pour $x\in]x_0-\frac{|x_0-5|}{2},x_0+\frac{|x_0-5|}{2}[$, on a $|x-5|>\frac{|x_0-5|}{2}[(>0)$, et donc,

$$\forall x\in]x_0-\frac{|x_0-5|}{2},x_0+\frac{|x_0-5|}{2}[,\;|f(x)-f(x_0)|=\frac{28}{(x_0-5)^2}|x-x_0|.$$

Soient $\varepsilon>0$ puis $\alpha=\mbox{Min}\{\frac{|x_0-5|}{2},\frac{(x_0-5)^2\varepsilon}{28}\}(>0)$.

$$|x-x_0|<\alpha\Rightarrow|f(x)-f(x_0)|\leq\frac{28}{(x_0-5)^2}|x-x_0|<\frac{28}{(x_0-5)^2}\frac{(x_0-5)^2\varepsilon}{28}=\varepsilon.$$

On a monté que $\forall\varepsilon>0,\;\exists\alpha>0/\;(\forall x\in\Rr\setminus\{5\},\;|x-x_0|<\alpha\Rightarrow|f(x)-f(x_0)|<\varepsilon)$. $f$ est donc continue sur $\Rr\setminus\{5\}$.
\fincorrection
\correction{005387}
Soit $\chi$ la fonction caractéristique de $\Qq$.
Soit $x_0$ un réel. On note que

$$x_0\in\Qq\Leftrightarrow\forall n\in\Nn^*,\;x_0+\frac{1}{n}\in\Qq Q\Leftrightarrow\forall n\in\Nn^*,\;x_0+\frac{\pi}{n}\notin\Qq.$$

Donc, $\lim_{n\rightarrow +\infty}\chi(x_0+\frac{1}{n})$ existe, $\lim_{n\rightarrow +\infty}\chi(x_0+\frac{\pi}{n})$ existe et$\lim_{n\rightarrow +\infty}\chi(x_0+\frac{1}{n})\neq\lim_{n\rightarrow +\infty}\chi(x_0+\frac{\pi}{n})$ (bien que $\lim_{n\rightarrow +\infty}x_0+\frac{1}{n}=\lim_{n\rightarrow +\infty}x_0+\frac{\pi}{n}=x_0$. Ainsi, pour tout réel $x_0\in\Rr$, la fonction caractéristique de $\Qq$ n'a pas de limite en $x_0$ et est donc discontinue en $x_0$.
\fincorrection
\correction{005388}
Soit $a$ un réel strcitement positif. On peut déjà noter que $\lim_{x\rightarrow a,\;x\in\Rr\setminus\Qq}f(x)=0$. Donc, si $f$ a une limite quand $x$ tend vers $a$, ce ne peut être que $0$ et $f$ est donc discontinue en tout rationnel strictement positif.

$a$ désigne toujours un réel strictement positif fixé. Soit $\varepsilon>0$.

Soit x un réel strictement positif tel que $f(x)\geq\varepsilon$.

$x$ est nécessairement rationnel, de la forme $\frac{p}{q}$ où $p$ et $q$ sont des entiers naturels non nuls premiers entre eux vérifiant $\frac{1}{p+q}\geq\varepsilon$ et donc 

$$2\leq p+q\leq\frac{1}{\varepsilon}.$$

Mais il n'y a qu'un nombre fini de couples d'entiers naturels non nuls $(p,q)$ vérifiant ces inégalités et donc, il n'y a qu'un nombre fini de réels strictement positifs $x$ tels que $f(x)\geq\varepsilon$.

Par suite, $\exists\alpha>0$ tel que aucun des réels $x$ de $]x_0-\alpha,x_0+\alpha[$ ne vérifie $f(x)\geq\varepsilon$. Donc, 

$$\forall a>0,\;\forall\varepsilon> 0,\;\exists\alpha>0/\;\forall x>0,\;(0<|x-a|<\alpha\Rightarrow|f(x)|<\varepsilon),$$

ou encore 

$$\forall a>0,\;\lim_{x\rightarrow a,\;x\neq a}f(x)=0.$$

Ainsi, $f$ est continue en tout irrationnel et discontinue en tout rationnel.
\fincorrection
\correction{005389}
Donnons tout d'abord une expression plus explicite de $f(x)$ pour chaque réel $x$.

Si $x>1$, alors $\frac{1}{x}\in]0,1[$ et donc, $f(x)=0$.

Si $\exists p\in\Nn^*/\;x\in]\frac{1}{p+1},\frac{1}{p}],\;f(x)=px$.

$f(0)=1$ (et plus généralement, $\forall p\in\Zz^*,\;f(\frac{1}{p})=1)$.

Si $x\leq-1$, alors $\frac{1}{x}\in[-1,0[$ et donc, $f(x)=-x$. 

Enfin, si $\exists p\in\Zz^\setminus\{-1\}$ tel que $x\in]\frac{1}{p+1},\frac{1}{p}]$, alors $\frac{1}{p+1}<x\leq\frac{1}{p}(<0)$ fournit, par décroissance de la fonction $x\mapsto\frac{1}{x}$ sur $]-\infty,0[$, $p\leq\frac{1}{x}<p+1(<0)$ et donc $f(x)=px$.

Etude en $0$. $\forall x\in\Rr^*,\;\frac{1}{x}-1<E(\frac{1}{x})\leq\frac{1}{x}$ et donc $1-x<f(x)\leq1$ si $x>0$ et $1\leq f(x)<1-x$ si $x<0$. Par suite, 

$$\forall x\in\Rr,\;|f(x)-1\leq|x|,$$

et $\lim_{x\rightarrow 0}f(x)=1$. $f$ est donc continue en $0$.

$f$ est affine sur chaque intervalle de la forme $]\frac{1}{p+1},\frac{1}{p}]$ pour $p$ élément de $\Zz\setminus\{-1,0\}$ et donc est continue sur ces intervalles et en particulier continue à gauche en chaque $\frac{1}{p}$. $f$ est affine sur $]-\infty,-1]$ et aussi sur $]1,+\infty[$ et est donc continue sur ces intervalles. Il reste donc à analyser la continuité à droite en $\frac{1}{p}$, pour $p$ entier relatif non nul donné. Mais,

$$f(\frac{1}{p}^+)=\lim_{x \rightarrow\frac{1}{p},\;x>\frac{1}{p}}(x(p-1))=1-\frac{1}{p}\neq 1=f(\frac{1}{p}).$$

$f$ est donc discontinue à droite en tout $\frac{1}{p}$ où $p$ est un entier relatif non nul donné.

Graphe de $f$~:

$$\includegraphics{../images/img005389-1}$$

\fincorrection
\correction{005390}
Soit $f(x)=\left\{
\begin{array}{l}
x\;\mbox{si}\;x\in(\Qq\cap[0,1])\setminus\{0,\frac{1}{2}\}\\
1-x\;\mbox{si}\;x\in(\Rr\setminus\Qq)\cap[0,1]\\
0\;\mbox{si}\;x=\frac{1}{2}\;\mbox{et}\;\frac{1}{2}\;\mbox{si}\;x=0
\end{array}
\right.$. $f$ est bien une application définie sur $[0,1]$ à valeurs dans $[0,1]$. De plus, si $x\in(\Qq\cap[0,1])\setminus\{0,\frac{1}{2}\}$, alors $f(f(x))=f(x)=x$.

Si $x\in(\Rr\setminus\Qq)\cap[0,1]$, alors $1-x\in(\Rr\setminus\Qq)\cap[0,1]$ et donc $f(f(x))=f(1-x)=1-(1-x)=x$.

Enfin, $f(f(0))=f(\frac{1}{2})=0$ et $f(f(\frac{1}{2}))=f(0)=\frac{1}{2}$.

Finalement, $f\circ f=Id_{[0,1]}$ et $f$, étant une involution de $[0,1]$, est une permutation de $[0,1]$.

Soit $a$ un réel de $[0,1]$. On note que $\lim_{x\rightarrow a,\;x\in(\Rr\setminus\Qq)}f(x)=1-a$ et $\lim_{x\rightarrow a,\;x\in\Qq}f(x)=a$. Donc, si $f$ a une limite en $a$, nécessairement $1-a=a$ et donc $a=\frac{1}{2}$. Mais, si $a=\frac{1}{2}$, $\lim_{x\rightarrow a,\;x\in\Qq,\;x\neq a}f(x)=a=\frac{1}{2}\neq0=f(\frac{1}{2})$ et donc $f$ est discontinue en tout point de $[0,1]$.
\fincorrection
\correction{005391}
Soit $T$  une période strictement positive de $f$. On note $\ell$ la limite de $f$ en $+\infty$.

Soit $x$ un réel. $\forall n\in\Nn,\;f(x)=f(x+nT)$ et quand $n$ tend vers $+\infty$, on obtient~:
 
$$f(x)=\lim_{n\rightarrow +\infty}f(x+nT)=\ell.$$ 

Ainsi, $\forall x\in\Rr,\;f(x)=\ell$ et donc, $f$ est constante sur $\Rr$.
\fincorrection
\correction{005392}
Voir la correction de l'exercice \ref{exo:roudist}.
\fincorrection
\correction{005395}
Soit $E=\{x\in[a,b]/\;f(x)\geq x\}$. $E$ est une partie non vide de $\Rr$ (car $a$ est dans $E$) et majorée (par $b$). Donc, $E$ admet une borne supèrieure $c$ vérifiant $a\leq c\leq b$.

Montrons que $f(c)=c$.

Si $c=b$, alors $\forall n\in\Nn^*,\;\exists x_n\in E/\;b-\frac{1}{n}<x_n\leq b$. Puisque $f$ est à valeurs dans $[a,b]$ et que les $x_n$ sont dans $E$, pour tout entier naturel non nul $n$, on a

$$x_n\leq f(x_n)\leq b\;(*).$$

Quand $n$ tend vers $+\infty$, la suite $(x_n)$ tend vers $b$ (théorème des gendarmes) et donc, $f$ étant croissante sur $[a,b]$, la suite $(f(x_n))$ tend vers $f(b^-)\leq f(b)$. Par passage à la limite quand $n$ tend vers $+\infty$ dans $(*)$, on obtient alors $b\leq f(b^-)\leq f(b)\leq b$ et donc $f(b)=b$. Finalement, dans ce cas, $b$ est un point fixe de $f$.

Si $c\in[a,b[$, par définition de $c$, pour $x$ dans $]c,b]$, $f(x)<x$ (car $x$ n'est pas dans $E$) et par passage à la limite quand $x$ tend vers $c$ par valeurs supérieures et d'après les propriétés usuelles des fonctions croissantes, on obtient~:~$f(c)(\leq f(c+))\leq c$.
 
D'autre part, $\forall n\in\Nn^*,\;\exists x_n\in E/\;c-\frac{1}{n}<x_n\leq c$. $x_n$ étant dans $E$, on a $f(x_n)\geq x_n$. Quand $n$ tend vers $+\infty$, on obtient~:~$f(c)\geq f(c^-)\geq c$. Finalement, $f(c)=c$ et dans tous les cas, $f$ admet au moins un point fixe.
\fincorrection
\correction{005396}
Puisque $f$ est croissante sur $[a,b]$, on sait que $f$ admet en tout point $x_0$ de $]a,b[$ une limite à gauche et une limite à droite réelles vérifiant $f(x_0^-)\leq f(x_0)\leq f(x_0^+)$ puis une limite à droite en $a$ élément de $[f(a),+\infty[$ et une limite à gauche en $b$ élément de $]-\infty,f(b)]$.

Si $f$ est discontinue en un $x_0$ de $]a,b[$, alors on a $f(x_0^-)<f(x_0)$ ou $f(x_0)<f(x_0^+)$. Mais, si par exemple $f(x_0^-)<f(x_0)$ alors, $\forall x\in[a,x_0[\;(\neq\emptyset),\;f(x)\leq f(x_0-)$ et $\forall x\in[x_0,b],\;f(x)\geq f(x_0)$.

Donc $]f(x_0^-),f(x_0)[\cap f([a,b])=\emptyset$ ce qui est exclu puisque d'autre part $]f(x_0^-),f(x_0)[\neq\emptyset$ et $]f(x_0^-),f(x_0)[\subset[f(a),f(b)]$ (la démarche est identique si $f(x_0^+)>f(x_0)$). Donc, $f$ est continue sur $]a,b[$. Par une démarche analogue, $f$ est aussi continue en $a$ ou $b$ et donc sur $[a,b]$.
\fincorrection
\correction{005398}
Posons $\ell=\lim_{x\rightarrow +\infty}f(x)$.

Soit $\varepsilon>0$. $\exists A>0/\;\forall x\in\Rr^+,\;(x\geq A\Rightarrow|f(x)-\ell|<\frac{\varepsilon}{3}$.
 
Soit $(x,y)\in[A,+\infty[^2$. Alors, $|f(x)-f(y)|\leq|f(x)-\ell|+|\ell-f(y)|<\frac{2\varepsilon}{3}(<\varepsilon)$. D'autre part, $f$ est continue sur le segment $[0,A]$ et donc est uniformément continue sur ce segment d'après le théorème de \textsc{Heine}. Donc, $\exists\alpha>0/\;\forall(x,y)\in[0,A]^2,\;|x-y|<\alpha\Rightarrow|f(x)-f(y)|<\frac{\varepsilon}{3}(<\varepsilon)$.

Résumons. $\alpha>0$ étant ainsi fourni, soient $x$ et $y$ deux réels de $[0,+\infty[$ vérifiant $|x-y|<\alpha$.

Si $(x,y)\in[0,A]^2$, on a $|f(x)-f(y)|<\frac{\varepsilon}{3}<\varepsilon$.

Si $(x,y)\in[A,+\infty[^2$, on a $|f(x)-f(y)|<\frac{2\varepsilon}{3}<\varepsilon$.

Si enfin on a $x\leq A\leq y$, alors, puisque $|A-x|\leq|x-y|<\alpha$, on a $|f(x)-f(A)|<\frac{\varepsilon}{3}$ et puisque $A$ et $y$ sont dans $[A,+\infty[$, on a $|f(y)-f(A)|<\frac{2\varepsilon}{3}$. Mais alors,

$$|f(x)-f(y)|\leq|f(x)-f(A)|+|f(y)-f(A)|<\frac{\varepsilon}{3}+\frac{2\varepsilon}{3}=\varepsilon.$$

On a montré que $\forall\varepsilon>0,\;\exists\alpha>0/\;\forall(x,y)\in[0,+\infty[^2,\;(|x-y|<\alpha\Rightarrow|f(x)-f(y)|<\varepsilon)$. $f$ est donc uniformément continue sur $[0,+\infty[$.
\fincorrection
\correction{005401}
Soit $T$ une pèriode strictement positive de $f$. $f$ est continue sur le segment $[0,T]$ et donc est bornée sur ce segment. $f$ est par suite bornée sur $\Rr$ par $T$-périodicité.

Soit $\varepsilon>0$.

$f$ est continue sur le segment $[0,T]$ et donc, d'après le théorème de \textsc{Heine}, $f$ est uniformément continue sur ce segment. Donc,

$$\exists\alpha\in]0,T[/\;\forall(x,y)\in[0,T],\;(|x-y|<\alpha\Rightarrow|f(x)-f(y)|<\frac{\varepsilon}{2}.$$

Soient $x$ et $y$ deux réels tels que $|x-y|<\alpha$. Si il existe un entier naturel $k$ tel que $(x,y)\in[kT,(k+1)T]$, alors $x-kT\in[0,T]$, $y-kT\in[0,T]$, puis $|(x-kT)-(y-kT)|=|y-x|<\alpha$ et donc $|f(x)-f(y)|<\frac{\varepsilon}{2}<\varepsilon$.

Sinon, en supposant par exemple que $x\leq y$, puisque l'on a choisi $\alpha<T$, 

$$\exists k\in\Zz/\;(k-1)T\leq x\leq kT\leq y\leq(k+1)T.$$

Mais alors, $|x-kT|=|y-x|<\alpha$ et $|y-kT|\leq|y-x|<\alpha$. Par suite,

$$|f(x)-f(y)|\leq|f(x)-f(kT)|+|f(y)-f(kT)|<\frac{\varepsilon}{2}+\frac{\varepsilon}{2}=\varepsilon.$$

Dans tous les cas, si $|x-y|<\alpha$, alors $|f(x)-f(y)|<\varepsilon$. On a montré que 

$$\forall\varepsilon>0,\;\exists\alpha>0/\;\forall(x,y)\in\Rr^2,\;(|x-y|<\alpha\Rightarrow|f(x)-f(y)|<\varepsilon).$$

$f$ est donc uniformément continue sur $\Rr$.
\fincorrection
\correction{000670}
Commençons par la fin, trouver un tel $\delta$ montrera que 
$$\forall \epsilon > 0 \quad \exists \delta > 0 \quad |x-x_0| < \delta \Rightarrow |f(x)-(-3)| < \epsilon$$
autrement dit la limite de $f$ en $x_0=0$ est $-3$.
Comme $f(0)=-3$ alors cela montre aussi que $f$ est continue en $x_0=0$.

\bigskip

On nous donne un $\epsilon>0$, à nous de trouver ce fameux $\delta$.
Tout d'abord 
$$\left| f(x)+3 \right| = \left| \frac{2x+3}{3x-1} + 3 \right| = \frac{11|x|}{|3x-1|}.$$
Donc notre condition devient :
$$ \left| f(x)+3 \right| < \epsilon \quad
 \Leftrightarrow \quad \frac{11|x|}{|3x-1|} < \epsilon 
\quad  \Leftrightarrow \quad |x| < \epsilon\frac{|3x-1|}{11}.$$

Comme nous voulons éviter les problèmes en $x = \frac 13$ pour lequel la fonction $f$ n'est pas définie, nous
allons nous placer ``loin'' de $\frac 13$.
Considérons seulement les $x \in \Rr$ tel que $|x| < \frac 16$.
Nous avons :
$$|x| < \frac 16 \Rightarrow -\frac 16 < x < + \frac 16 \quad  \Rightarrow \quad  -\frac 32 < 3x-1 < -\frac 12 \quad \Rightarrow \quad \frac 12 < |3x-1|.$$
Et maintenant explicitons $\delta$ :
prenons $\delta < \epsilon \cdot \frac{1}{2}\cdot \frac{1}{11}$.
Alors pour $|x| < \delta$ nous avons 
$$|x| < \delta = \epsilon \cdot  \frac{1}{2} \cdot \frac{1}{11} < \epsilon \cdot|3x-1|\cdot \frac{1}{11}$$
ce qui implique par les équivalences précédentes que 
$\left| f(x)+3 \right| < \epsilon$.

Il y a juste une petite correction à apporter à notre $\delta$ : au cours de nos calculs
nous avons supposé que $|x| < \frac 16$, mais rien ne garantie que $\delta \le \frac 16$
(car $\delta$ dépend de $\epsilon$ qui pourrait bien être très grand, même
si habituellement ce sont les $\epsilon$ petits qui nous intéressent).
Au final le $\delta$ qui convient est donc :
$$\delta = \min (\frac 16, \frac{\epsilon}{22}).$$

\bigskip

Remarque finale :
bien sûr on savait dès le début que $f$ est continue en $x_0=0$. En effet
$f$ est le quotient de deux fonctions continues, le dénominateur ne s'annulant
pas en $x_0$. Donc nous savons dès le départ qu'un tel $\delta$ existe,
mais ici nous avons fait plus, nous avons trouvé une formule explicite pour ce $\delta$.

\fincorrection
\correction{000671}
\begin{enumerate}
 \item Le graphe est composé d'une portion de droite au dessus des $x \in ]-\infty,1[$ ;
d'une portion de parabole pour les $x \in [1,4]$, d'une portion d'une autre parabole pour les 
$x \in ]4,+\infty$. (Cette dernière branche est bien une parabole, mais elle n'est pas dans le sens ``habituel'',
en effet si $y=8\sqrt x$ alors $y^2 = 64 x$ et c'est bien l'équation d'une parabole.)

On ``voit'' immédiatemment sur le graphe que la fonction est continue (les portions se recollent !).
On ``voit'' aussi que la fonction est bijective.

 \item La fonction est continue sur $]-\infty,1[$, $]1,4[$ et $]4,+\infty[$ car sur chacun des ces intervalles elle
y est définie par une fonction continue. Il faut examiner ce qui se passe en $x=1$ et $x=4$.
Pour $x<1$, $f(x)=x$, donc  la limite à gauche (c'est-à-dire $x\to 1$ avec $x<1$) est donc $+1$.
Pour $x\ge 1$, $f(x) = x^2$ donc la limite à droite vaut aussi $+1$. 
Comme on a $f(1) = +1$ alors les limites à gauche, à droite et la valeur en $1$ coïncident donc $f$
est continue en $x=1$.

Même travail en $x= 4$. Pour $x \in [1,4]$, $f(x) = x^2$ donc la limite à gauche en $x=4$ est $+16$.
On a aussi $f(4)=+16$. Enfin pour $x>4$, $f(x) = 8 \sqrt x$, donc la limite à droite en $x=4$ est aussi $+16$.
Ainsi $f$ est continue en $x=4$.

Conclusion : $f$ est continue en tout point $x \in \Rr$ donc $f$ est continue sur $\Rr$.
 
 \item Le graphe devrait vous aider : tout d'abord il vous aide à se convaincre que $f$ est bien bijective
et que la formule pour la bijection réciproque dépend d'intervalles. Petit rappel : le graphe de la bijection réciproque $f^{-1}$ s'obtient
comme symétrique du graphe de $f$ par rapport à la bissectrice d'équation $(y=x)$ (dans un repère orthonormal).



Ici on se contente de donner directement la formule de $f^{-1}$. 
Pour $x \in  ]-\infty,1[$, $f(x)=x$. Donc la bijection réciproque est définie par $f^{-1}(y)=y$ pour tout $y \in  ]-\infty,1[$.
Pour $x \in [1,4]$, $f(x)=x^2$. L'image de l'intervalle $[1,4]$ est l'intervalle $[1,16]$. Donc pour chaque $y \in [1,16]$,
la bijection réciproque est définie par $f^{-1}(y) = \sqrt y$. 
Enfin pour $x\in]4,+\infty[$, $f(x) = 8\sqrt x$. L'image de l'intervalle $]4,+\infty[$ est donc $]16,+\infty[$
et $f^{-1}$ est définie par $f^{-1}(y) = \frac{1}{64} y^2$ pour chaque $y \in ]16,+\infty[$.

Nous avons définie $f^{-1} : \Rr \to \Rr$ de telle sorte que $f^{-1}$ soit la bijection réciproque de $f$.

\bigskip

C'est un bon exercice de montrer que $f$ est bijective sans calculer $f^{-1}$ :
vous pouvez par exemple montrer que $f$ est injective et surjective.
Un autre argument est d'utiliser un résultat du cours : $f$ est continue, strictement croissante avec une limite $-\infty$ en$-\infty$
et $+\infty$ en $+\infty$ donc elle est bijective de $\Rr$ dans $\Rr$ (et on sait même que la bijection réciproque est continue).
\end{enumerate}
\fincorrection
\correction{000672}
Soit $x_0 \not= 0$, alors la fonction $f$ est continue en $x_0$, car elle s'exprime sous la forme d'un quotient de fonctions continues o\`u le d\'enominateur ne s'annule pas en $x_0$.
Reste \`a \'etudier la continuit\'e en $0$.
Mais 
$$ \lim_{x\rightarrow 0} \frac{\sin x}{x} = 1 = f(0)$$
donc $f$ est continue en $0$.
\fincorrection
\nocorrection
\nocorrection
\nocorrection
\nocorrection
\correction{000677}
\begin{enumerate}
  \item La fonction est d\'efinie sur $\Rr^*$ t elle est continue sur $\Rr^*$. Il faut d\'eterminer un \'eventuel prolongement par continuit\'e en $x=0$, c'est-\`a-dire savoir si $f$ a une limite en $0$.
$$|f(x)| = |\sin x| |\sin 1/x| \leq |\sin x|.$$
Donc $f$ a une limite en $0$ qui vaut $0$. 
Donc en posant $f(0) = 0$, nous obtenons une fonction $f : \Rr \longrightarrow \Rr$ qui est continue.
  \item La fonction $g$ est d\'efinie et continue sur $\Rr^*$. 
Etudions la situation en $0$. Il faut remarquer que $g$ est la taux d'accroissement en
$0$ de la fonction $k(x) = \ln \frac{e^x+e^{-x}}{2}$ : en effet $g(x)=\frac{k(x)-k(0)}{x-0}$.
Donc si $k$ est dérivable en $0$ alors la limite de $g$ en $0$ est \'egale \`a
la valeur de $k'$ en $0$. 

Or la fonction $k$ est dérivable sur $\Rr$ et $k'(x) = \frac{e^x-e^{-x}}{e^x+e^{-x}}$
donc $k'(0)=0$. Bilan : en posant $g(0)=0$ nous obtenons une fonction $g$ d\'efinie et continue sur $\Rr$.

  \item $h$ est d\'efinie et continue sur $\Rr \setminus \{ -1,1 \}$.
$$h(x) = \frac{1}{1-x} - \frac{2}{1-x^2} = \frac{1+x-2}{(1-x)(1+x)}
= \frac{-1+x}{(1-x)(1+x)} = \frac{-1}{(1+x)}.$$
Donc $h$ a pour limite $-\frac 12$ quand $x$ tend vers $1$.
Et donc en posant $h(1) = -\frac 12$, nous d\'efinissons une fonction
continue sur $\Rr \setminus \{ -1 \}$.
En $-1$ la fonction $h$ ne peut \^etre prolong\'ee continuement,
car en $-1$, $h$ n'admet de limite finie.
\end{enumerate}
\fincorrection
\nocorrection
\nocorrection
\correction{000680}
Fixons $x\in \Rr$ et soit $y = x/2$, comme $f(y) = f(2y)$ nous obtenons $f(\frac 12 x) = f(x)$. Puis en prenant $y = \frac 14 x$, nous obtenons
$f(\frac 14 x) = f(\frac 12 x)= f(x)$. Par une r\'ecurrence facile nous avons
$$\forall n \in \Nn \ \ \ \ f(\frac 1{2^n} x) = f(x).$$
Notons $(u_n)$ la suite d\'efinie par $u_n = \frac 1{2^n} x$ alors
$u_n \rightarrow 0$ quand $n \rightarrow +\infty$.
Par la continuit\'e de $f$ en $0$ nous savons alors que:
$ f(u_n) \rightarrow f(0)$ quand $n \rightarrow +\infty$.
Mais $f(u_n) = f(\frac 1{2^n} x) = f(x)$, donc $(f(u_n))_n$ est une suite constante \'egale \`a $f(x)$, et donc la limite de cette suite est $f(x)$ ! Donc $f(x) = f(0)$. Comme ce raisonnement est valable pour tout $x \in \Rr$ nous venons de montrer que $f$ est une fonction constante.
\fincorrection
\nocorrection
\nocorrection
\nocorrection
\nocorrection
\nocorrection
\nocorrection
\nocorrection
\nocorrection
\nocorrection
\correction{000609}
 G\'en\'eralement pour calculer des limites faisant intervenir des sommes de racines carr\'ees,  il est utile de faire intervenir ``l'expression conjugu\'ee": 
$$\sqrt a - \sqrt b = \frac{(\sqrt a - \sqrt b)(\sqrt a + \sqrt b)}{\sqrt a + \sqrt b} = \frac{a-b}{\sqrt a + \sqrt b}.$$
Les racines au num\'erateur ont ``disparu" en utilisant l'identit\'e
$(x-y)(x+y) = x^2-y^2$.

Appliquons ceci sur un exemple :
\begin{align*}
 f(x) &= 
  \frac{\sqrt{1+x^m}-\sqrt{1-x^m}}{x^n} \\
     &=  \frac{(\sqrt{1+x^m}-\sqrt{1-x^m})(\sqrt{1+x^m}+\sqrt{1-x^m})}{x^n(\sqrt{1+x^m}+\sqrt{1-x^m})} \\
   &= \frac{1+x^m-(1-x^m)}{x^n(\sqrt{1+x^m}+\sqrt{1-x^m})}  \\
   &= \frac{2x^m}{x^n(\sqrt{1+x^m}+\sqrt{1-x^m})}  \\
   &= \frac{2x^{m-n}}{\sqrt{1+x^m}+\sqrt{1-x^m}}  \\
\end{align*}
Et nous avons 
$$ \lim_{x \rightarrow 0} \frac 2 {\sqrt{1+x^m}+\sqrt{1-x^m}} = 1.$$
Donc l'\'etude de la limite de $f$ en $0$ est la m\^eme que celle de la fonction $x \mapsto x^{m-n}$.

Distinguons plusieurs cas pour la limite de $f$ en $0$.
\begin{itemize}
  \item Si $m > n$ alors $x^{m-n}$, et donc $f(x)$, tendent vers $0$.
  \item Si $m=n$ alors $x^{m-n}$ et $f(x)$ tendent vers $1$.
  \item Si $m < n$ alors $x^{m-n} = \frac 1 {x^{n-m}}  = \frac 1 {x^k}$ avec $k = n-m$ un exposant positif. Si $k$ est pair alors les limites 
\`a droite et \`a gauche de $\frac 1{x^k}$ sont $+\infty$. 
Pour $k$ impair la limite \`a droite vaut $+\infty$ et la limite \`a gauche vaut $-\infty$. Conclusion pour $k=n-m>0$ pair, la limite de $f$ en $0$ vaut $+\infty$ et pour $k = n-m>0$ impair $f$ \emph{n'a pas de limite en $0$} car les limites \`a droite et \`a gauche ne sont pas \'egales.
\end{itemize}
\fincorrection
\nocorrection
\nocorrection
\correction{000612}
\begin{enumerate}
  \item Soit $p>0$ la p\'eriode: pour tout $x\in \Rr$,
$f(x+p) = f(x)$. Par une r\'ecurrence facile on montre :
$$ \forall n \in \Nn \qquad \forall x\in \Rr \qquad f(x+np)=f(x).$$
Comme $f$ n'est pas constante il existe $a,b \in \Rr$ tels que $f(a)\not= f(b)$. Notons $x_n = a +np$ et $y_n = b+np$.
Supposons, par l'absurde, que $f$ a une limite $\ell$ en $+\infty$.
Comme $x_n \rightarrow +\infty$ alors $f(x_n) \rightarrow \ell$.
Mais $f(x_n) = f(a +np) = f(a)$, donc $\ell = f(a)$. 
De m\^eme avec la suite $(y_n)$: $y_n \rightarrow +\infty$ donc $f(y_n) \rightarrow \ell$ et $f(y_n) = f(b +np) = f(b)$, donc $\ell = f(b)$. 
Comme $f(a) \not= f(b)$ nous obtenons une contradiction.
  \item Soit $f: \Rr \longrightarrow \Rr$ une fonction croissante et major\'ee par $M\in \Rr$. 
Notons 
$$ F = f(\Rr) = \{ f(x) \ | \ x \in \Rr \}.$$
$F$ est un ensemble (non vide) de $\Rr$, notons $\ell = \sup F$.
Comme $M\in \Rr$ est un majorant de $F$, alors
$\ell < +\infty$.
Soit $\epsilon > 0$, par les propri\'et\'es du $\sup$
il existe $y_0 \in F$ tel que $\ell - \epsilon \leq y_0 \leq \ell$.
Comme $y_0\in F$, il existe $x_0 \in \Rr$ tel que $f(x_0) = y_0$.
Comme $f$ est croissante alors:
$$ \forall x\geq x_0 \qquad f(x) \geq f(x_0) = y_0 \geq \ell - \epsilon.$$
De plus par la d\'efinition de $\ell$:
$$\forall x\in \Rr \ \ f(x) \leq \ell.$$
Les deux propri\'et\'es pr\'ec\'edentes s'\'ecrivent:
$$ \forall x\geq x_0 \qquad \ell - \epsilon \leq f(x) \leq \ell.$$
Ce qui exprime bien que la limite de $f$ en $+\infty$ est $\ell$.
\end{enumerate}
\fincorrection
\nocorrection
\nocorrection
\nocorrection
\correction{000616}
\begin{enumerate}
  \item $\frac{x^2+2|x|}{x} = x + 2 \frac{|x|}{x}$.
Si $x > 0$ cette expression vaut $x+2$ donc la limite à droite en $x=0$ est $+2$.
Si $x<0$ l'expression vaut $-2$ donc la limite à gauche en $x=0$ est $-2$.
Les limites \`a droite et à gauche sont différentes donc il n'y a pas de limite en $x=0$.

  \item $\frac{x^2+2|x|}{x} = x + 2 \frac{|x|}{x} = x -2$ pour $x<0$.
Donc la limite quand $x \to -\infty$ est $-\infty$.

  \item $\frac{x^2-4}{x^2-3\,x+2}=\frac{(x-2)(x+2)}{(x-2)(x-1)} = \frac{x+2}{x-1}$,
lorsque $x\to 2$ cette expression tend vers $4$.

  \item $\frac{\sin^2 x}{1+\cos x} = \frac{1-\cos^2 x}{1+\cos x} = \frac{(1-\cos x)(1+\cos x)}{1+\cos x} = 1-\cos x$.
Lorsque $x\to \pi$ la limite est donc $2$.

  \item $\frac{\sqrt{1+x}-\sqrt{1+x^2}}{x} = \frac{\sqrt{1+x}-\sqrt{1+x^2}}{x}\times\frac{\sqrt{1+x}+\sqrt{1+x^2}}{\sqrt{1+x}+\sqrt{1+x^2}}
= \frac{1+x - (1+x^2)}{x(\sqrt{1+x}+\sqrt{1+x^2})} = \frac{x - x^2}{x(\sqrt{1+x}+\sqrt{1+x^2})} = \frac{1-x}{\sqrt{1+x}+\sqrt{1+x^2}}$.
Lorsque $x\to 0$ la limite vaut $\frac 12$.

  \item $\sqrt{x+5}-\sqrt{x-3} = \left( \sqrt{x+5}-\sqrt{x-3}\right) \times \frac{\sqrt{x+5}+\sqrt{x-3}}{\sqrt{x+5}+\sqrt{x-3}}
= \frac{x+5-(x-3)}{\sqrt{x+5}+\sqrt{x-3}} = \frac{8}{\sqrt{x+5}+\sqrt{x-3}}$. Lorsque $x \to +\infty$, la limite vaut $0$.

  \item Nous avons l'égalité $a^3-1 = (a-1)(1+a+a^2)$. Pour $a = \sqrt[3]{1+x^2}$
cela donne :
$$\frac{a-1}{x^2} = \frac{a^3-1}{x^2(1+a+a^2)} = \frac{1+x^2-1}{x^2(1+a+a^2)} = \frac{1}{1+a+a^2}.$$
Lors que $x\to 0$, alors $a \to 1$ et la limite cherchée est $\frac 13$.

Autre méthode : si l'on sait que la limite d'un taux d'accroissement correspond à la dérivée nous avons une méthode moins
astucieuse. Rappel (ou anticipation sur un prochain chapitre) : pour une fonction $f$ dérivable en $a$ alors
$$\lim_{x\to a} \frac{f(x)-f(a)}{x-a} = f'(a).$$
Pour la fonction $f(x) = \sqrt[3]{1+x} =(1+x)^{\frac 13}$ ayant $f'(x) = \frac 13 (1+x)^{-\frac 23}$ cela donne en $a=0$ :
$$\lim_{x\to 0} \frac{\sqrt[3]{1+x^2} -1}{x^2} = \lim_{x\to 0} \frac{\sqrt[3]{1+x} -1}{x} = \lim_{x\to 0} \frac{f(x)-f(0)}{x-0} = f'(0) = \frac 13.$$

  \item $\frac{x^n-1}{x-1} = 1+x+x^2+\cdots + x^n$. Donc si $x\to 1$ la limite de  $\frac{x^n-1}{x-1}$ est  $n$.
Donc la limite de $\frac{x-1}{x^n-1}$ en $1$ est $\frac 1n$.

La méthode avec le taux d'accroissement fonctionne aussi très bien ici. Soit $f(x) = x^n$, $f'(x)=nx^{n-1}$ et $a=1$.
Alors $\frac{x^n-1}{x-1} = \frac{f(x)-f(1)}{x-1}$ tend vers $f'(1)=n$.
\end{enumerate}
\fincorrection
\nocorrection
\nocorrection
\nocorrection
\nocorrection
\nocorrection
\nocorrection
\correction{000623}
\begin{enumerate}
\item $- \infty$
\item $0$
\item $+ \infty$
\item $+ \infty$ 
\item ${3 \over2}$
\item $- \infty$
\item $0$
\item $0$
\item $0$
\item $0$
\item $-2$
\item $- \infty$
\item $1$
\item $e^4$
\item $1$
\item $e$
\item $e$
\item $0$
\item $0$
\item $0$
\end{enumerate}
\fincorrection
\nocorrection
\nocorrection
\nocorrection
\nocorrection
\correction{000628}
\begin{enumerate}
  \item Montrons d'abord que la limite de $$f(x) = \frac{x^k-\alpha^k}{x-\alpha}$$ en $\alpha$ est 
$k\alpha^{k-1}$, $k$ étant un entier fixé. Un calcul montre que 
$f(x) = x^{k-1} + \alpha x^{k-2} + \alpha^2x^{k-3} + \cdots + \alpha^{k-1}$ ;
en effet $(x^{k-1} + \alpha x^{k-2} + \alpha^2x^{k-3} + \cdots + \alpha^{k-1})(x-\alpha) = x^k-\alpha^k$.
Donc la limite en $x=\alpha$ est $k\alpha^{k-1}$.
Une autre m\'ethode consiste \`a dire que $f(x)$ est la taux d'accroissement de la fonction $x^k$, et donc la limite de
$f$ en $\alpha$ est exactement la valeur de la d\'eriv\'ee de $x^k$ en $\alpha$, soit $k\alpha^{k-1}$.
Ayant fait ceci revenons \`a la limite de l'exercice : comme 
$$
\frac{x^{n+1}-\alpha^{n+1}}{x^n-\alpha^n} = \frac{x^{n+1}-\alpha^{n+1}}{x-\alpha} \times \frac{x-\alpha}{x^n-\alpha^n}.$$
Le premier terme du produit tend vers $(n+1)\alpha^n$ et le second
terme, \'etant l'inverse d'un taux d'accroissement, tend vers $1/(n\alpha^{n-1})$.
Donc la limite cherch\'ee est 
$$\frac {(n+1)\alpha^n}{n\alpha^{n-1}}= \frac {n+1}{n} \alpha.$$
  \item La fonction $f(x)=\frac{\tan x - \sin x}{\sin x(\cos 2x - \cos x)}$ s'\'ecrit aussi $f(x) = \frac{1-\cos x}{\cos x (\cos 2x- \cos x)}$. Or $\cos 2x =  2\cos^2 x - 1$. Posons $u = \cos x$, alors 
$$f(x) = \frac {1-u}{u(2u^2-u-1)} = \frac {1-u}{u(1-u)(-1-2u)}= \frac {1}{u(-1-2u)}$$
Lorsque $x$ tend vers $0$, $u = \cos x$ tend vers $1$, et donc $f(x)$ tend vers $-\frac 13$.

  \item 
\begin{align*}
\sqrt{x+\sqrt{x+\sqrt x}}-\sqrt x 
  &= \frac{\left(\sqrt{x+\sqrt{x+\sqrt x}}-\sqrt x\right)\left(\sqrt{x+\sqrt{x+\sqrt x}}+\sqrt x\right)}{\sqrt{x+\sqrt{x+\sqrt x}}+\sqrt x} \\
  &= \frac{\sqrt{x+\sqrt x}}{\sqrt{x+\sqrt{x+\sqrt x}}+\sqrt x} \\
 &= \frac{\sqrt{1+\frac 1 {\sqrt x}}}{\sqrt{1+\frac{\sqrt{x+\sqrt x}}{x}}+1} \\
\end{align*}
Quand $x \rightarrow +\infty$ alors $\frac 1 {\sqrt x} \rightarrow 0$
et $\frac{\sqrt{x+\sqrt x}}{x}=\sqrt{\frac 1x + \frac{1}{x\sqrt{x}}}\rightarrow 0$, donc la limite recherch\'ee 
est $\frac 12$.
  \item La fonction s'\'ecrit 
$$f(x) = \frac{\sqrt{x}-\sqrt{\alpha}-\sqrt{x-\alpha}}{\sqrt{x^2-\alpha^2}} = \frac{\sqrt x - \sqrt \alpha - \sqrt{x-\alpha}}{\sqrt{x-\alpha}\sqrt{x+\alpha}} = 
\frac{\frac{\sqrt x - \sqrt \alpha}{\sqrt{x-\alpha}}-1}{\sqrt{x+\alpha}}.$$
Notons $g(x) = \frac{\sqrt x - \sqrt \alpha}{\sqrt{x-\alpha}}$ alors 
\`a l'aide de l'expression conjugu\'ee $$g(x) = \frac{x -  \alpha}{(\sqrt{x-\alpha})(\sqrt x +\sqrt \alpha)} = \frac{\sqrt{x -  \alpha}}{\sqrt x +\sqrt \alpha}.$$
Donc $g(x)$ tend vers $0$ quand $x \rightarrow \alpha^+$. Et maintenant 
$f(x) = \frac{g(x)-1}{\sqrt{x+\alpha}}$ tend vers $-\frac {1}{\sqrt{2 \alpha}}$.
  \item Pour tout r\'eel $y$ nous avons la double in\'egalit\'e $y-1 < E(y) \leq y$. Donc pour $y>0$,
$\frac{y-1}{y} < \frac{E(y)}{y} \leq 1$.
On en d\'eduit que lorsque $y$ tend vers $+\infty$  alors
$\frac{E(y)}{y}$ tend $1$. On obtient le même résultat quand $y$ tend vers $-\infty$. 
En posant $y = 1/x$, et en faisant tendre
$x$ vers $0$, alors $xE(\frac 1x) = \frac{E(y)}{y}$ tend vers $1$.
  \item 
$$\frac {e^x-e^2}{x^2+x-6} = \frac {e^x-e^2}{x-2} \times \frac{x-2}{x^2+x-6} =\frac {e^x-e^2}{x-2} \times \frac{x-2}{(x-2)(x+3)} = \frac {e^x-e^2}{x-2} \times \frac{1}{x+3}.$$
La limite de $\frac {e^x-e^2}{x-2}$ en $2$ vaut $e^2$ ($\frac {e^x-e^2}{x-2}$ est la taux d'accroissement de la fonction $x \mapsto e^x$ en la valeur $x=2$), la limite voulue est $\frac {e^2}{5}$.

  \item Soit $f(x) = \frac{x^4}{1+x^\alpha\sin^2x}$. Supposons $\alpha \geq 4$, alors on prouve que $f$ n'a pas de limite en $+\infty$.
En effet pour pour $u_k = 2k\pi$, $f(2k\pi) = (2k\pi)^4$ tend vers $+\infty$ lorsque $k$ (et donc $u_k$) tend vers $+\infty$.
Cependant pour $v_k = 2k\pi + \frac\pi 2$, $f(v_k) = \frac{v_k^4}{1+v_k^\alpha}$ tend vers $0$ (ou vers $1$ si $\alpha = 4$) lorsque $k$ (et donc $v_k$) tend vers $+\infty$. Ceci prouve que $f(x)$ n'a pas de limite lorsque $x$ tend vers $+\infty$.

Reste le cas $\alpha < 4$. Il existe $\beta$ tel que
$\alpha < \beta < 4$. 
$$f(x)= \frac {x^4}{1+x^\alpha \sin^2 x}= \frac {x^{4-\beta}}{\frac 1 {x^\beta} +\frac{x^\alpha}{x^\beta} \sin^2 x}.$$
Le num\'erateur tend $+\infty$ car $4-\beta >0$.
$\frac 1 {x^\beta}$ tend vers $0$ ainsi que $\frac{x^{\alpha}}{x^\beta} \sin^2 x$ (car $\beta > \alpha$ et $\sin^2 x$ est born\'ee par $1$). Donc le d\'enominateur tend vers $0$ (par valeurs positives). La limite est donc de type $+\infty / 0^+$ (qui n'est pas ind\'etermin\'ee !) et vaut donc $+\infty$.
\end{enumerate}
\fincorrection
\nocorrection
\nocorrection
\nocorrection
\nocorrection
\nocorrection
\correction{000634}
R\'{e}ponse: $\frac{2}{3}$
\fincorrection
\correction{000635}
\begin{enumerate}
 \item Comme $-1 \le \sin \frac 1x \le +1$ alors $1 \le 2 + \sin \frac 1x \le +3$.
Donc pour $x>0$, nous obtenons $\frac x3 \le \frac{x}{2+\sin \frac{1}{x}} \le x$.
On obtient une inégalité similaire pour $x<0$.
Cela implique $\lim\limits_{x\rightarrow 0}\frac{x}{2+\sin \frac{1}{x}} = 0$.

 \item Sachant que $\frac{\ln (1+t)}{t} \to 1$ lorsque $t \to 0$, on peut le reformuler
ainsi $\ln(1+t) = t \cdot \mu(t)$, pour une certaine fonction $\mu$ qui vérifie $\mu(t) \to 1$ lorsque $t\to 0$.
Donc $\ln (1+e^{-x}) = e^{-x} \mu(e^{-x})$.
Maintenant 

\begin{align*}
(\ln (1+e^{-x}))^{\frac{1}{x}} 
  &= \exp \left(\frac{1}{x} \ln\left( \ln (1+e^{-x}) \right)  \right) \\
  &= \exp \left(\frac{1}{x} \ln \left( e^{-x} \mu(e^{-x}) \right)  \right) \\
  &= \exp \left(\frac{1}{x} \left( -x + \ln \mu(e^{-x}) \right)  \right) \\
  &= \exp \left( -1 + \frac{\ln \mu(e^{-x})}{x}  \right) \\
\end{align*}
$\mu(e^{-x}) \to 1$ donc $\ln\mu(e^{-x}) \to 0$, donc $\frac{\ln \mu(e^{-x})}{x} \to 0$ lorsque $x \to +\infty$.

Bilan : la limite est $\exp(-1)=\frac 1 e$.
 \item


 \item Sachant $\frac{e^x - 1}{x} \to 1$ lorsque $x\to 0$, on reformule ceci
en $e^x-1 = x \cdot \mu(x)$, pour une certaine fonction $\mu$ qui vérifie $\mu(x) \to 1$ lorsque $x\to 0$.
Cela donne $\ln (e^x-1) = \ln (x \cdot \mu(x)) = \ln x + \ln \mu(x).$
\begin{align*}
 x^{\frac{1}{\ln (e^{x}-1)}} 
  &= \exp\left( \frac{1}{\ln (e^{x}-1)} \ln x  \right) \\
  &= \exp\left( \frac{1}{\ln x + \ln \mu(x)} \ln x  \right) \\
  &= \exp\left( \frac{1}{1 + \frac{\ln \mu(x)}{\ln x}} \right) \\
\end{align*}
Maintenant $\mu(x) \to 1$ donc $\ln \mu(x) \to 0$, et $\ln x \to - \infty$ lorsque $x \to 0$.
Donc $\frac{\ln \mu(x)}{\ln x} \to 0$.
Cela donne 
$$\lim\limits_{x\rightarrow 0^{+}}x^{\frac{1}{\ln (e^{x}-1)}} = \lim\limits_{x\rightarrow 0^{+}} \exp\left( \frac{1}{1 + \frac{\ln \mu(x)}{\ln x}} \right) = \exp\left(1\right) = e.$$

\fincorrection
\correction{000636}
R\'{e}ponse: 1.
\fincorrection
\correction{000637}
Supposons $a \ge b$.
Alors 
$$\left(\frac{a^{x}+b^{x}}{2}\right)^{\frac{1}{x}} = \left( a^x \times \frac{1+ (\frac{b}{a})^x}{2}\right)^{\frac{1}{x}}
= a \left( \frac{1+ (\frac{b}{a})^x}{2}\right)^{\frac{1}{x}}.$$
Or $0 \le \frac b a \le 1$, donc $0 \le (\frac ba)^x \le 1$ pour tout $x \ge 1$.
Donc $(\frac 12)^{\frac{1}{x}} \le \left( \frac{1+ (\frac{b}{a})^x}{2}\right)^{\frac{1}{x}} \le 1^{\frac{1}{x}}$.
Les deux termes extrêmes tendent vers $1$ lorsque $x$ tend vers $+\infty$ donc le terme du milieu tend aussi vers $1$.
Conclusion : si $a \ge b$ alors  $\lim\limits_{x\rightarrow +\infty}\left(\frac{a^{x}+b^{x}}{2}\right)^{\frac{1}{x}} = a.$
Si $b \ge a$ alors cette limite vaudrait $b$.
Cela se résume dans le cas général o\`u $a,b$ sont quelconques par 
 $\lim\limits_{x\rightarrow +\infty}\left(\frac{a^{x}+b^{x}}{2}\right)^{\frac{1}{x}} = \max(a,b).$
\fincorrection
\correction{000638}

Soit 
$$
 f(x) = \left(\frac{a^{x}+b^{x}}{2}\right)^{\frac{1}{x}}
     = \exp\left( \frac 1x \ln \left(\frac{a^{x}+b^{x}}{2}\right) \right) 
$$

$a^x \to 1$, $b^x \to 1$ donc $\frac{a^{x}+b^{x}}{2} \to 1$ lorsque $x \to 0$
et nous sommes face à une forme indéterminée.
Nous savons que $\lim_{t \to 0} \frac{\ln(1+t)}{t} = 1$.
Autrement dit il existe un fonction $\mu$ telle que $\ln(1+t) = t \cdot \mu(t)$ avec
$\mu(t) \to 1$ lorsque $t\to 0$.

Appliquons cela à $g(x) = \ln \left(\frac{a^{x}+b^{x}}{2}\right)$.
Alors 
$$g(x) = \ln \left(1+ \left(\frac{a^{x}+b^{x}}{2}-1\right)\right) =  \left(\frac{a^{x}+b^{x}}{2}-1\right) \cdot \mu(x)$$
o\` u $\mu(x) \to 1$ lorsque $x\to 0$. (Nous écrivons pour simplifier $\mu(x)$ au lieu
de $\mu(\frac{a^{x}+b^{x}}{2}-1)$.)



\bigskip

Nous savons aussi que $\lim_{t \to 0} \frac{e^t - 1}{t} = 1$.
Autrement dit il existe un fonction $\nu$ telle que $e^t - 1 = t \cdot \nu(t)$ avec
$\nu(t) \to 1$ lorsque $t\to 0$.

Appliquons ceci :

\begin{align*}
 \frac{a^{x}+b^{x}}{2}-1 
     &= \frac 12 (e^{x \ln a} + e^{x\ln b})-1 \\
     &= \frac12 (e^{x \ln a}-1 + e^{x\ln b}-1) \\
     &= \frac12( x \ln a \cdot \nu(x \ln a) + x \ln b \cdot \nu(x \ln b)) \\
     &= \frac12 x\left(\ln a \cdot \nu(x \ln a) +  \ln b \cdot \nu(x \ln b) \right) \\
\end{align*}

\bigskip

Reste à rassembler tous les éléments du puzzle :
\begin{align*}
 f(x) &= \left(\frac{a^{x}+b^{x}}{2}\right)^{\frac{1}{x}} \\
      &= \exp\left( \frac 1x \ln \left(\frac{a^{x}+b^{x}}{2}\right) \right) \\
      &= \exp\left( \frac 1x g(x) \right) \\
      &= \exp\left( \frac 1x  \left(\frac{a^{x}+b^{x}}{2}-1\right) \cdot \mu(x) \right) \\
      &= \exp\left( \frac 1x  \cdot \frac12  \cdot x\left(\ln a \cdot \nu(x \ln a) +  \ln b \cdot \nu(x \ln b) \right) \cdot \mu(x) \right) \\
      &= \exp\left( \frac12 \left(\ln a \cdot \nu(x \ln a) +  \ln b \cdot \nu(x \ln b) \right) \cdot \mu(x) \right) \\
\end{align*}

Or $\mu(x) \to 1$, $\nu(x \ln a) \to 1$, $\nu(x \ln b) \to 1$ lorsque $x \to 0$.
Donc 
$$\lim_{x\to 0} f(x) = \exp\left( \frac12 \left(\ln a + \ln b \right) \right) = \exp\left( \frac12 \ln (ab) \right) = \sqrt{ab}.$$
\fincorrection
\nocorrection
\correction{005101}
Pour $x>0$, $(x^x)^x=e^{x\ln(x^x)}=e^{x^2\ln x}$ et $x^{(x^x)}=e^{x^x\ln x}$. Par suite,

$$\forall x>0,\;\frac{(x^x)^x}{x^{(x^x)}}=\text{exp}(\ln x(x^2-x^x)).$$
Or, $x^2-x^x=-x^x(1-x^{2-x})=-e^{x\ln x}(1-e^{(2-x)\ln x})$. Quand $x$ tend vers $+\infty$, $(2-x)\ln x$ tend vers
$-\infty$. Donc, $1-e^{(2-x)\ln x}$ tend vers $1$ puis $x^2-x^x$ tend vers $-\infty$. Mais alors, $\ln x(x^2-x^x)$ tend
vers $-\infty$, puis $\frac{(x^x)^x}{x^{(x^x)}}=\mbox{exp}(\ln x(x^2-x^x))$ tend vers $0$.

\begin{center}
\shadowbox{
$\lim_{x\rightarrow +\infty}\frac{(x^x)^x}{x^{(x^x)}}=0.$
}
\end{center}
\fincorrection
\correction{000686}
\begin{enumerate}
  \item Il faut que le d\'enominateur ne s'annule pas donc $x \not= \frac 52$. En plus il faut
que le terme sous la racine soit positif ou nul, c'est-\`a-dire $(2+3x)\times(5-2x) \geq 0$, soit
$x \in [-\frac23, \frac52]$. L'ensemble de d\'efinition est donc  $[-\frac23, \frac52[$.

  \item Il faut $x^2-2\,x-5\geq 0$, soit $x \in ]-\infty, 1-\sqrt6] \cup [1+\sqrt6,+\infty[$.
  \item Il faut $4x+3>0$ soit $x > -\frac34$, l'ensemble de d\'efinition \'etant $]-\frac34,+\infty[$.
\end{enumerate}
\fincorrection
\nocorrection
\nocorrection
\nocorrection
\correction{000690}
 Pour tout $x\in\R$ on a :
          $$ 0\leq |f(x)| = \frac{|\cos x|}{1+x^2} \leq \frac{1}{1+x^2}\leq 1. $$
Par cons\'equent, pour tout $x\in\R$, $f(x)\in [-1,1] $ donc $f$
est minor\'ee ($-1$ est un minorant), major\'ee ($1$ est un
majorant) et $\sup_{x\in\R}f(x)\leq 1$. Comme $f(0)=1$ on a
n\'ecessairement $\sup_{x\in\R}f(x) \geq 1$. Conclusion :
        $$   \sup_{x\in\R}f(x)=1.$$
\fincorrection
\nocorrection
\nocorrection
\nocorrection
\nocorrection
\nocorrection
\nocorrection
\nocorrection
\nocorrection
\nocorrection
\nocorrection
\nocorrection
\correction{003883}
$\alpha = \inf(A)  \Rightarrow  f(\alpha^+) \le \alpha \text{ et } f(\alpha^-) \ge \alpha$.
\fincorrection
\correction{003884}
Supposons qu'il existe $a,b\in\R$ avec $a<b$ et $f(b)<f(a)$.
On note $E=\{x\in{[a,b]}\text{ tq }f(x) < f(a)\}$ et $c=\inf(E)$.
On a $c\in E$ et $c>a$ par hypothèse et donc $f(c)=\lim_{x\to c^-}f(x)\ge f(a)$, absurde.
\fincorrection
\correction{003885}
Injectivité évidente.

Monotonie~: pour $a<b<c$ on a $|a-b| < |a-c|$ et $|c-b| < |c-a|$ d'où les
mêmes inégalités pour $f(a),f(b),f(c)$ ce qui prouve que $f(b)$ est strictement
compris entre $f(a)$ et $f(c)$.

Continuité~: soit $a\in\R$, $\delta>0$, $x=a-\delta$, $y=a+\delta$ et $z=a-4\delta$.
On a $2\delta=|x-y|<|x-z|=3\delta$ donc $|f(x)-f(y)|<|f(x)-f(z)|$ et
en faisant tendre $\delta$ vers~$0^+$~: $|f(a^-)-f(a^+)| \le |f(a^-)-f(a^-)| = 0$.

Affine~: soient $x\in\R$, $h>0$, $z=x+h$ et $x-h<y<x$. On a $|f(x)-f(y)|<|f(x)-f(x+h)|$
d'où en faisant tendre $y$ vers $(x-h)^+$~: $|f(x)-f(x-h)|\le|f(x)-f(x+h)|$.
On obtient l'inégalité inverse en permutant $y$ et $z$, donc $f(x-h)$ et $f(x+h)$
sont équidistants de $f(x)$ et, par injectivité de~$f$~:
$f(x) = \frac{f(x-h)+f(x+h)}2$ ce qui permet de conclure avec la continuité de~$f$.
\fincorrection
\nocorrection
\nocorrection
\correction{003888}
\begin{enumerate}
  \item \'Etudier les logs.
  \item Idem.
\end{enumerate}
 \fincorrection
\correction{003889}
$\frac{f'(x)}{f(x)} = \frac a{(1+ax)\ln(1+ax)} - \frac b{(1+bx)\ln(1+bx)}$.

Pour $x \ge 0$ fixé, la fonction $t\longmapsto \frac t{(1+tx)\ln(1+tx)}$ est
décroissante.
\fincorrection
\nocorrection
\nocorrection
\correction{003892}
\begin{enumerate}
  \item Oui ssi $|a| > |b|$.
  \item Oui ssi $|a| < |b|$.
\end{enumerate}
 \fincorrection
\correction{003893}
$=\frac{ \ch(nx/2)\sh((n+1)x/2) }{ \sh(x/2) }$.
\fincorrection
\correction{003894}
$x = -\frac23a$.
\fincorrection
\correction{003895}
$2\coth 2x - \frac 1x$.
\fincorrection
\correction{003896}
$\coth \frac x2 - 1$.
\fincorrection
\nocorrection
\correction{003898}
Poser $X = e^x$, $Y = e^y$:
$ \Rightarrow \begin{cases} X+Y = a+b\cr  XY = \frac{a+b}{a-b}.\end{cases}$
Il y a des solutions si et seulement si $a \ge \sqrt{b^2+4}$.
\fincorrection
\nocorrection
\correction{003900}
$= x + \ln\sqrt2$.
\fincorrection
\correction{003901}
$x=\frac54$.
\fincorrection
\correction{003902}
\begin{enumerate}
  \item $F(x) = \Argsh\left(\frac{2x+1}{\sqrt3}\right)$.
  \item $F(x) = \frac 1{\sqrt5}\ln\left|\frac{\sqrt5 -2x-1}{\sqrt5 +2x+1}\right|$.
\end{enumerate}
 \fincorrection
\correction{003903}
\begin{enumerate}
  \item \'Etude de $x \mapsto \left(\frac{a_1}a\right)^x + \dots + \left(\frac{a_p}a\right)^x$.
  \item $x_a > x_b$.
  \item $x_a\to\ell$.
               Si $\ell > 0$, $a^{a_a} \to +\infty$, mais
               $a_1^{x_a} + \dots + a_p^{x_a} \to a_1^\ell + \dots + a_p^\ell$.

               Donc $\ell = 0$, et $x_a\ln a \to \ln p$.
\end{enumerate}
 \fincorrection
\correction{003904}
\begin{enumerate}
  \item Pour $x=1$ on a $f\circ f(y) = yf(1)$ donc $f$ est injective
    et pour $y=1$~: $f(xf(1))=f(x)$ d'où $f(1)=1$.
  \item $f(xy) = f(xf(f(y))) = f(y)f(x)$.\par
    Pour $0<x<1$ on a $f(x^n) = f(x)^n \to +\infty$ (lorsque $n\to\infty$) donc
    $f(x)>1$ ce qui entraîne par morphisme la décroissance de~$f$.
    Enfin $f$ est monotone et $f(]0,+\infty[) = {]0,+\infty[}$ donc $f$
    n'a pas de saut et est continue.
  \item En tant que morphisme continu, $f$ est de la forme $x \mapsto x^\alpha$
    avec $\alpha\in\R$ et l'involutivité et la décroissance donnent $\alpha=-1$.
\end{enumerate}
 \fincorrection
\correction{003905}
\begin{enumerate}
  \item $4\cos\left(\theta-\frac\pi{12}\right) = 2 \iff
               \theta \equiv \pm\frac\pi3 + \frac\pi{12} (\mathrm{mod}\, {2\pi})$.
  \item     $\sin\theta + \dots + \sin4\theta
                  = 2\sin\theta\cos\theta(4\cos^2\theta + 2\cos\theta - 1)
                  = 4\sin( 5\theta/2 )\cos\theta\cos( \theta/2 )$

               $4\cos^2\theta + 2\cos\theta - 1 = 0 \iff
                  \cos\theta = \frac {\sqrt5 -1}4 = \cos( 2\pi/5 )
                  \text{ ou }
                  \cos\theta = -\frac {\sqrt5 +1}4 = \sin( 2\pi/5 )$

               $ \Rightarrow $ modulo $2\pi$,
               $\theta \in \{ 0, \pi, \pi/2, 3\pi/2, 2\pi/5, 4\pi/5,
                                   6\pi/5, 8\pi/5 \}$. 
  \item  $\cos\theta \in \left\{ -\frac 12, \pm\frac 1{\sqrt2} \right\} \iff
               \theta \equiv \pm\frac {2\pi}3 (\mathrm{mod}\, {2\pi})
               \text{ ou } \theta \equiv \frac \pi4 (\mathrm{mod}\, {\frac\pi2})$. 
  \item $2\sin(3\theta/2)\sin(\theta/2) = 2\sin(3\theta/2)\cos(3\theta/2)
             \iff
             \theta \equiv 0 (\mathrm{mod}\, {\frac {2\pi}3}) \text{ ou }
             \theta \equiv \frac \pi4 (\mathrm{mod}\, \pi) \text{ ou }
             \theta \equiv \frac {3\pi}2 (\mathrm{mod}\, {2\pi})$. 
  \item $2\cos4\theta\cos3\theta = \cos4\theta \iff
             \theta \equiv \frac \pi8 (\mathrm{mod}\, {\frac \pi4}) \text{ ou }
             \theta \equiv \pm\frac \pi9 (\mathrm{mod}\, {\frac{2\pi}3})$. 
  \item $2\cos7\theta\cos5\theta = \sqrt3\cos5\theta \iff
             \theta \equiv \frac \pi{10} (\mathrm{mod}\, {\frac \pi5}) \text{ ou }
             \theta \equiv \pm\frac \pi{42} (\mathrm{mod}\, {\frac{2\pi}7})$. 
  \item $\theta\equiv 0\mathrm{mod}\, {\frac\pi3} \text{ ou } \theta\equiv\pm\frac \pi{12} (\mathrm{mod}\, {\frac \pi2})$.
  \item     $\cos^3\theta\sin3\theta + \cos3\theta\sin^3\theta
                = \frac 34\sin4\theta
                 \Rightarrow  \theta \equiv \frac \pi8 (\mathrm{mod}\, {\frac \pi2})$. 
  \item $\theta\equiv 0 (\mathrm{mod}\, {\frac \pi8})$.
  \item $\sin\theta = \frac 12 \Leftrightarrow
           \theta\equiv \frac\pi6 (\mathrm{mod}\, {2\pi})$ ou
          $\theta\equiv \frac{5\pi}6 (\mathrm{mod}\, {2\pi})$.
  \item  $\theta \equiv 0,\pm\arctan\sqrt5 (\mathrm{mod}\,\pi)$. 
  \item  $\cos^2\theta - \sin^2\theta
           = \cos\theta\sin\theta(\cos\theta+\sin\theta)$.

          $\cos\theta + \sin\theta = 0 \iff
           \theta\equiv -\frac\pi4 (\mathrm{mod}\, \pi)$.

          $\cos\theta-\sin\theta = \cos\theta\sin\theta  \Rightarrow 
           (\cos\theta\sin\theta)^2 + 2\cos\theta\sin\theta = 1  \Rightarrow 
           \begin{cases}\cos\theta = \frac {\sqrt{2\sqrt2-1} + \sqrt2 - 1}2 \cr
        
                  \sin\theta = \frac {\sqrt{2\sqrt2-1} - \sqrt2 + 1}2.\end{cases}$

          Les valeurs trouvées conviennent.
  \item $\tan x = \tan y = \frac 12$.
\end{enumerate}
 \fincorrection
\correction{003906}
\begin{enumerate}
  \item $-\frac {2\pi}3 < \theta(\mathrm{mod}\,{2\pi}) < \frac \pi3$. 
  \item $2\alpha < \theta(\mathrm{mod}\,{2\pi}) < 2\pi$ avec
               $\begin{cases}\cos\alpha = 2/\sqrt5 \cr \sin\alpha = 1/\sqrt5.\end{cases}$
\end{enumerate}
 \fincorrection
\nocorrection
\correction{003908}
\begin{enumerate}
  \item $1 - \cos\alpha = 2\sin\frac\alpha2 \cos\frac {\beta+\gamma}2$,

$\cos\beta + \cos\gamma = 2\sin\frac\alpha2 \cos\frac {\beta-\gamma}2$.
  \item $=1$.
\end{enumerate}
 \fincorrection
\nocorrection
\correction{003910}
$\tan p - \tan q = \frac {\sin(p-q)}{\cos p\cos q}$.

$S_n = \begin{cases}\frac {\tan((n+1)\theta) - \tan\theta}{\sin\theta}
                     &\text{ si } \sin\theta \ne 0 \cr
              n      &\text{ si } \theta \equiv 0 (\mathrm{mod}\, {2\pi}) \cr
              -n     &\text{ si } \theta \equiv \pi (\mathrm{mod}\, {2\pi}).\end{cases}$
\fincorrection
\correction{003911}
linéariser :
         $\sum = \frac 14\left(3^n\sin\frac\alpha{3^n} - \sin\alpha\right)$.
\fincorrection
\correction{003912}
$\mathrm{cotan} x - 2\mathrm{cotan}2x = \tan x$,
         $\sum = \frac 1{2^n}\mathrm{cotan}\frac\alpha{2^n} - 2\mathrm{cotan}2\alpha$.
\fincorrection
\correction{003913}
 $\theta = \frac \pi7$ :
 $\frac 1{\sin\theta} - \frac 1{\sin3\theta} = \frac {2\cos2\theta}{\sin3\theta}
 = \frac 1{\sin2\theta}$.
\fincorrection
\correction{005097}

\begin{enumerate}
 \item  Soit $f$ une fonction dérivable sur $\Rr$ à valeurs dans $\Rr$. Si $f$ est paire, alors, pour tout réel $x$,
$f(-x)=f(x)$. En dérivant cette égalité, on obtient

$$\forall x\in\Rr,\;-f'(-x)=f'(x),$$
et donc $f'$ est impaire. De même, si $f$ est impaire, pour tout réel $x$, on a $f(-x)=-f(x)$, et par dérivation on
obtient pour tout réel $x$, $f'(-x)=f'(x)$. $f'$ est donc paire.

\begin{center}
\shadowbox{
$(f\;\mbox{paire}\Rightarrow f'\;\mbox{impaire) et}\;(f\;\mbox{impaire}\Rightarrow f'\;\mbox{paire.})$
}
\end{center}
 \item  Soient $n\in\Nn^*$ et $f$ une fonction $n$ fois dérivable sur $\Rr$ à valeurs dans $\Rr$. Supposons $f$
paire. Par suite, pour tout réel $x$, $f(-x)=f(x)$. Immédiatement par récurrence, on a

$$\forall x\in\Rr,\;f^{(n)}(-x)=(-1)^nf(x).$$
Ceci montre que $f^{(n)}$ a la parité de $n$, c'est-à-dire que $f^{(n)}$ est une fonction paire quand $n$ est un
entier pair et est une fonction impaire quand $n$ est un entier impair.
De même, si $f$ est impaire et $n$ fois dérivable sur $\Rr$, $f^{(n)}$ a la parité contraire de celle de $n$.
 \item  Soit $f$ une fonction continue sur $\Rr$ et impaire et $F$ une primitive de $f$. Montrons que $F$ est paire.
Pour $x$ réel, posons $g(x)=F(x)-F(-x)$. $g$ est dérivable sur $\Rr$ et pour tout réel $x$,

$$g'(x)=F'(x)+F'(-x)=f(x)+f(-x)=0.$$
$g$ est donc constante sur $\Rr$ et par suite, pour tout réel $x$, $g(x)=g(0)=F(0)-F(0)=0$. Ainsi, $g$ est la fonction nulle et donc, pour
tout réel $x$, $F(x)=F(-x)$. On a montré que $F$ est paire.
Par contre, si $f$ est paire, $F$ n'est pas nécessairement impaire. Par exemple, la fonction $f~:~x\mapsto1$ est paire,
mais $F~:~x\mapsto x+1$ est une primitive de $f$ qui n'est pas impaire.
 \item  On montre aisément en dérivant une ou plusieurs fois l'égalité~:~$\forall x\in\Rr,\;f(x+T)=f(x)$, que les
dérivées successives d'une fonction $T$-périodique sont $T$-périodiques. Par contre, il n'en est pas de même des
primitives. Par exemple, si pour tout réel $x$, $f(x)=\cos^2x=\frac{1}{2}(1+\cos(2x))$, $f$ est $\pi$-périodique, mais
la fonction $F~:~x\mapsto\frac{x}{2}+\frac{\sin(2x)}{4}$, qui est une primitive de $f$ sur $\Rr$, n'est pas
$\pi$-périodique ni même périodique tout court.
\end{enumerate}
\fincorrection
\correction{005098}
Pour $n\in\Nn^*$, posons $u_n=\sqrt[n]{n}$ puis, pour $x$ réel strictement positif, $f(x)=x^{1/x}$
de sorte que pour tout naturel non nul $n$, on a $u_n=f(n)$.
$f$ est définie sur $]0,+\infty[$ et pour $x>0$, $f(x)=e^{\ln x/x}$. $f$ est dérivable sur $]0,+\infty[$ et pour
$x>0$,

$$f'(x)=\frac{1-\ln x}{x^2}e^{\ln x/x}.$$
Pour $x>0$, $f'(x)$ est du signe de $1-\ln x$ et donc $f'$ est strictement positive sur $]0,e[$ et strictement négative
sur $]e,+\infty[$. $f$ est donc strictement croissante sur $]0,e]$ et strictement décroissante sur $[e,+\infty[$. En
particulier, pour $n\geq3$,

$$u_n=f(n)\leq f(3)=u_3=\sqrt[3]{3}.$$
Comme $u_2=\sqrt{2}>1=u_1$, on a donc $\mbox{Max}\{u_n,\;n\in\Nn^*\}=\mbox{Max}\{\sqrt{2},\sqrt[3]{3}\}$. Enfin,
$\sqrt{2}=1,41...<1,44..=\sqrt[3]{3}$ (on peut aussi constater que $(\sqrt{2})^6=8<9=(\sqrt[3]{3})^6$). Finalement,

\begin{center}
\shadowbox{
$\text{Max}\left\{\sqrt[n]{n},\;n\in\Nn^*\right\}=\sqrt[3]{3}=1,44...$
}
\end{center}
\fincorrection
\correction{005099}
\begin{enumerate}
 \item  Pour $x>0$, posons $f(x)=\frac{\ln x}{x}$. $f$ est définie et dérivable sur $]0,+\infty[$ et, pour $x>0$,
$f'(x)=\frac{1-\ln x}{x^2}$. $f$ est donc strictement croissante sur $]0,e]$ et strictement décroissante sur
$[e,+\infty[$. Le graphe de $f$ s'en déduit facilement~:

$$\includegraphics{../images/img005099-1}$$


 \item  Soient $a$ et $b$ deux entiers naturels non nuls tels que $a<b$. On a alors

$$a^b=b^a\Leftrightarrow\ln(a^b)=\ln(b^a)\Leftrightarrow b\ln a=a\ln b\Leftrightarrow\frac{\ln a}{a}=\frac{\ln b}{b}\Leftrightarrow f(a)=f(b).$$
Si $a\geq3$, puisque $f$ est strictement décroissante sur $[e,+\infty[$, on a alors $f(a)>f(b)$ et en particulier,
$f(a)\neq f(b)$. $a$ n'est donc pas solution.
$a=1$ n'est évidemment pas solution. Par exemple, $a^b=b^a\Rightarrow1^b=b^1\Rightarrow b=1=a$ ce qui est exclu.
Donc, nécessairement $a=2$ et $b$ est un entier supérieur ou égal à $3$, et donc à $e$, vérifiant $f(b)=f(2)$. Comme $f$
est strictement décroissante sur $[e,+\infty[$, l'équation $f(b)=f(2)$ a au plus une solution dans $[e,+\infty[$. Enfin,
comme $2^4=16=4^2$, on a montré que~:~il existe un et un seul couple $(a,b)$ d'entiers naturels non nuls tel que $a<b$
et $a^b=b^a$, à savoir $(2,4)$.
\end{enumerate}
\fincorrection
\correction{005100}
\begin{enumerate}
 \item  Soit $x\in\Rr$,

\begin{align*}
\ln|x+1|-\ln|2x+1|\leq\ln2&\Leftrightarrow\ln\left|\frac{x+1}{2x+1}\right|\leq\ln2\Leftrightarrow\left|\frac{x+1}{2x+1}\right|\leq2\;
\mbox{et}\;x+1\neq0\\
 &\Leftrightarrow-2\leq\frac{x+1}{2x+1}\leq2\;\mbox{et}\;x\neq-1\Leftrightarrow\frac{x+1}{2x+1}+2\geq0\;\mbox{et}\;\frac{x+1}{2x+1}-2
 \leq0\;\mbox{et}\;x\neq-1\\
 &\Leftrightarrow\frac{5x+3}{2x+1}\geq0\;\mbox{et}\;\frac{-3x-1}{2x+1}\leq0\;\mbox{et}\;x\neq-1\\
 &\Leftrightarrow\left(x\in\left]-\infty,-\frac{3}{5}\right]\cup\left]-\frac{1}{2},+\infty\right[\right)\;\mbox{et}\;\left(\left]-\infty,-\frac{1}{2}\right[\cup\left[-
\frac{1}{3},+\infty\right[\right)\;\mbox{et}\;x\neq-1\\
 &\Leftrightarrow x\in]-\infty,-1[\cup\left]-1,-\frac{3}{5}\right]\cup\left[-\frac{1}{3},+\infty\right[
\end{align*}

 \item  Pour $x>0$

\begin{align*}
x^{\sqrt{x}}=\sqrt{x}^x&\Leftrightarrow\sqrt{x}\ln x=x\ln\sqrt{x}\Leftrightarrow\ln x(\sqrt{x}-\frac{x}{2})=0\\
 &\Leftrightarrow\ln x\times\sqrt{x}(2-\sqrt{x})=0\Leftrightarrow x=1\;\mbox{ou}\;x=4.
\end{align*}

 \item  $\Argch3=\ln(3+\sqrt{3^2-1})=\ln(3+\sqrt{8})$ et
$\Argth\frac{7}{9}=\frac{1}{2}\ln\left(\frac{1+\frac{7}{9}}{1-\frac{7}{9}}\right)=\ln\sqrt{8}$. Donc,
$\Argch3-\Argth\frac{7}{9}=\ln\left(1+\frac{3}{\sqrt{8}}\right)$.
Par suite,

\begin{align*}
2\Argsh x=\Argch3-\Argth\frac{7}{9}&\Leftrightarrow x=\sh\left(\frac{1}{2}\ln\left(1+\frac{3}{\sqrt{8}}\right)\right)\\
 &\Leftrightarrow
x=\frac{1}{2}\left(\sqrt{1+\frac{3}{\sqrt{8}}}-\frac{1}{\sqrt{1+\frac{3}{\sqrt{8}}}}\right)=\frac{3}{2\sqrt{8}}
\frac{1}{\sqrt{1+\frac{3}{\sqrt{8}}}}=\frac{3}{2\sqrt[4]{8}}
\frac{1}{\sqrt{3+2\sqrt{2}}}\\
 &\Leftrightarrow x=\frac{3\sqrt[4]{2}}{4}\frac{1}{\sqrt{(1+\sqrt{2})^2}}=\frac{3\sqrt[4]{2}(\sqrt{2}-1)}{4}.
\end{align*}

 \item  Pour $x\in]0,+\infty[\setminus\left\{\frac{1}{100},\frac{1}{10},1\right\}$,

\begin{align*}
\mbox{ln}_x(10)+2\mbox{ln}_{10x}(10)&+3\mbox{ln}_{100x}(10)=0\Leftrightarrow\frac{\ln(10)}{\ln
x}+2\frac{\ln(10)}{\ln(10x)}+3\frac{\ln(10)}{\ln(100x)}=0\\
 &\Leftrightarrow\frac{(\ln x+\ln(10))(\ln x+2\ln(10))+2\ln x(\ln x+2\ln(10))+3\ln x(\ln x+\ln(10))}{\ln x(\ln x+\ln(10))(\ln
x+2\ln(10))}=0\\
 &\Leftrightarrow6\ln^2x+10\ln(10)\times\ln x+2\ln^2(10)=0\\
 &\Leftrightarrow\ln x\in\left\{\frac{-5\ln(10)+\sqrt{13\ln^2(10)}}{6},\frac{-5\ln(10)-\sqrt{13\ln^2(10)}}{6}\right\}\\
 &\Leftrightarrow x\in\left\{10^{(-5-\sqrt{13})/6},10^{(-5+\sqrt{13})/6}\right\}.
\end{align*}
Comme aucun de ces deux nombres n'est dans $\left\{\frac{1}{100},\frac{1}{10},1\right\}$, $\mathcal{S}=\left\{10^{(-5-\sqrt{13})/6},10^{(-5+\sqrt{13})/6}\right\}$.

 \item  Soit $x\in\Rr$.

\begin{align*}
2^{2x}-3^{x-\frac{1}{2}}=3^{x+\frac{1}{2}}-2^{2x-1}&\Leftrightarrow2^{2x}+2^{2x-1}=3^{x+\frac{1}{2}}+3^{x-\frac{1}{2}}\\
 &\Leftrightarrow2^{2x-1}(2+1)=3^{x-\frac{1}{2}}(3+1)\Leftrightarrow3\times2^{2x-1}=4\times3^{x-\frac{1}{2}}\\
 &\Leftrightarrow2^{2x-3}=3^{x-\frac{3}{2}}\Leftrightarrow(2x-3)\ln2=\left(x-\frac{3}{2}\right)\ln3\\
 &\Leftrightarrow x=\frac{3\ln2-\frac{3}{2}\ln3}{2\ln2-\ln3}\Leftrightarrow x=\frac{3}{2}.
\end{align*}
\end{enumerate}
\fincorrection
\correction{005102}
On notera $\mathcal{C}_i$ le graphe de $f_i$.
\begin{enumerate}
 \item  $f_1$ est définie et continue sur $\Rr$, dérivable sur $\Rr\setminus\left\{-2,\frac{1}{2}\right\}$. On précise dans
un tableau l'expression de $f_1(x)$ suivant les valeurs de $x$.

$$
\begin{array}{|c|lcc|ccc|ccr|}
\hline
x&-\infty& &\multicolumn{2}{c}{-2}& &\multicolumn{2}{c}{1/2}& &+\infty\\
\hline
|2x-1|& &-2x+1& & &-2x+1& & &2x-1& \\
\hline
|x+2|& &-x-2& & &x+2& & &x+2& \\
\hline
f_1(x)& &4& & &-2x& & &6x-4& \\
\hline
\end{array}
$$
On en déduit $\mathcal{C}_1$.


$$\includegraphics{../images/img005102-1}$$

 \item  Soit $x\in\Rr$. $\ch x\geq1$ et donc $f_2(x)$ existe et $f_2(x)\geq0$. $f_2$ est donc définie sur $\Rr$. De
plus, $f_2$ est continue et dérivable sur $\Rr$, paire.
Puisque la fonction $x\mapsto\ch x$ est strictement croissante sur $\Rr^+$ à valeurs dans $]0,+\infty[$ et que la fonction $x\mapsto\ln x$ est
strictement croissante sur $]0,+\infty[$, $f_2$ est strictement croissante sur $\Rr^+$ et, par parité, strictement
décroissante sur $\Rr^-$.
$f_2$ est paire et donc $f_2'$ est impaire. Par suite, $f_2'(0)=0$ et $\mathcal{C}_2$ admet l'axe des abscisses pour
tangente en $(0,f_2(0))=(0,0)$.
\textbf{Etude en} $+{\infty}$. Pour $x\geq0$,

$$f_2(x)=\ln\left(\frac{1}{2}(e^x+e^{-x}))=\ln(e^x+e^{-x}\right)-\ln2=\ln(e^x(1+e^{-2x}))-\ln2=x-\ln2+\ln(1+e^{-2x}).$$
Quand $x$ tend vers $+\infty$, $e^{-2x}$ tend vers $0$ et donc, $\ln(1+e^{-2x})$ tend vers $0$. On en déduit que
$\lim_{x\rightarrow +\infty}f_2(x)=+\infty$. De plus, $\lim_{x\rightarrow +\infty}(f_2(x)-(x-\ln2))=0$ et la droite $(D)$ d'équation
$y=x-\ln2$ est asymptote à $\mathcal{C}_2$ en $+\infty$. Par symétrie par rapport à la droite $(Oy)$, la droite $(D')$ d'équation $y=-x-\ln2$ est asymptote à $\mathcal{C}_2$ en $-\infty$. Enfin, pour tout réel $x$,

$$f_2(x)-(x-\ln2)=\ln(1+e^{-2x})>\ln1=0,$$
et $\mathcal{C}_2$ est strictement au-dessus de $(D)$ sur $\Rr$. De même, $\mathcal{C}_2$ est strictement au-dessus de
$(D')$ sur $\Rr$.
On en déduit $\mathcal{C}_2$.


$$\includegraphics{../images/img005102-2}$$

 \item  $f_3$ est définie et continue sur $\Rr$, dérivable sur $\Rr\setminus\{-1,1\}$.
\textbf{Etude en} $-\infty$. Soit $x\leq-1$.

$$f_3(x)=x+\sqrt{x^2-1}=\frac{(x+\sqrt{x^2-1})(x-\sqrt{x^2-1})}{x-\sqrt{x^2-1}}=\frac{1}{x-\sqrt{x^2-1}}.$$
Or, quand $x$ tend vers $-\infty$, $x-\sqrt{x^2-1}$ tend vers $-\infty$ et donc $\lim_{x\rightarrow -\infty}f_3(x)=0$.
\textbf{Etude en} $+\infty$. Immédiatement, $\lim_{x\rightarrow +\infty}f_3(x)=+\infty$. Ensuite, pour $x\geq1$,

$$\frac{f_3(x)}{x}=\frac{x+\sqrt{x^2-1}}{x}=1+\sqrt{1-\frac{1}{x^2}},$$
qui tend vers $2$ quand $x$ tend vers $+\infty$. Mais alors,

$$f_3(x)-2x=-x+\sqrt{x^2-1}=\frac{(-x+\sqrt{x^2-1})(-x-\sqrt{x^2-1})}{-x-\sqrt{x^2-1}}=-\frac{1}{x+\sqrt{x^2-1}}.$$
On en déduit que $\lim_{x\rightarrow +\infty}(f_3(x)-2x)=0$ et donc que la droite $(D)$ d'équation $y=2x$ est asymptote à $\mathcal{C}_3$ en $+\infty$.
\textbf{Etude en} $1$. Pour $x>1$,

$$\frac{f_3(x)-f_3(1)}{x-1}=\frac{(x-1)+\sqrt{(x-1)(x+1)}}{x-1}=1+\sqrt{\frac{x+1}{x-1}},$$
et pour $x\in]-1,1[$,

$$\frac{f_3(x)-f_3(1)}{x-1}=\frac{(x-1)+\sqrt{(-x+1)(x+1)}}{-(-x+1)}=1-\sqrt{\frac{x+1}{-x+1}}.$$
Par suite, $\lim_{x\rightarrow 1,\;x>1}\frac{f_3(x)-f_3(1)}{x-1}=+\infty$ et
$\lim_{x\rightarrow 1,\;x<1}\frac{f_3(x)-f_3(1)}{x-1}=-\infty$. On en déduit que $f_3$ n'est pas dérivable en $1$, mais que
$\mathcal{C}_3$ admet deux demi-tangentes parallèles à $(Oy)$ au point de $\mathcal{C}_3$ d'abscisse $1$. Les résultats
sont analogues en $-1$.
\textbf{Etude des variations de} $\bf{f_3}$. Pour $x\in]-\infty,-1[\cup]1,+\infty[$, $f_3(x)=x+\sqrt{x^2-1}$ et donc

$$f_3'(x)=1+\frac{x}{\sqrt{x^2-1}}=\frac{x+\sqrt{x^2-1}}{\sqrt{x^2-1}}.$$

Si $x>1$, on a $x+\sqrt{x^2-1}>0$ et donc, $f_3'(x)>0$. Si $x<-1$, on a

$$\sqrt{x^2-1}<\sqrt{x^2}=|x|=-x,$$
et donc, $x+\sqrt{x^2-1}<0$ puis $f_3'(x)<0$. Ainsi, $f_3$ est strictement décroissante sur $]-\infty,-1[$ et
strictement croissante sur $]1,+\infty[$.
Pour $x\in]-1,1[$, $f_3(x)=x+\sqrt{-x^2+1}$ et donc

$$f_3'(x)=1-\frac{x}{\sqrt{-x^2+1}}=\frac{\sqrt{-x^2+1}-x}{\sqrt{-x^2+1}}.$$
Si $x\in]-1,0]$, on a clairement $f_3'(x)>0$. Si x$\in[0,1[$, par stricte croissance de la fonction $x\mapsto x^2$ sur $\Rr^+$, on a

$$\mbox{sgn}(f_3'(x))=\mbox{sgn}(\sqrt{-x^2+1}-x)=\mbox{sgn}((-x^2+1)-x^2)=\mbox{sgn}(1-2x^2)=\mbox{sgn}((1-x\sqrt{2})
(1+x\sqrt{2}))=\mbox{sgn}[\frac{1}{\sqrt{2}}-x).$$
Donc, $f_3'$ est strictement positive sur $\left[0,\frac{1}{\sqrt{2}}\right[$, strictement négative sur $\left]\frac{1}{\sqrt{2}},1\right[$
et s'annule en $\frac{1}{\sqrt{2}}$.
En résumé, $f_3'$ est strictement négative sur $]-\infty,-1[$ et sur $\left]\frac{1}{\sqrt{2}},1\right[$ et strictement positive
sur $\left]-1,\frac{1}{\sqrt{2}}\right[$ et sur $]1,+\infty[$. $f_3$ est donc strictement croissante sur $]-\infty,-1]$ et sur
$\left[\frac{1}{\sqrt{2}},1\right[$ et strictement décroissante sur $\left[-1,\frac{1}{\sqrt{2}}\right]$ et sur $[1,+\infty[$.
On en déduit $\mathcal{C}_3$.

$$\includegraphics{../images/img005102-3}$$

 \item  $f_4$ est définie sur $\Rr\setminus\left(\frac{\pi}{2}+\pi\Zz\right)$, $2\pi$-périodique et paire. On étudie
donc $f_4$ sur $\left[0,\frac{\pi}{2}\right[\cup\left]\frac{\pi}{2},\pi\right]$.
\textbf{Etude des variations de} $\bf{f_4}$. Pour $x\in\left[0,\frac{\pi}{2}\right[$, $f_4(x)=\tan x+\cos x$ et donc,

$$f_4'(x)=\frac{1}{\cos^2x}-\sin x\geq 1-1=0,$$
avec égalité si et seulement si $\sin x=\cos^2x=1$ ce qui est impossible. Donc, $f_4'$ est strictement positive sur
$\left[0,\frac{\pi}{2}\right[$ et $f_4$ est strictement croissante sur $\left[0,\frac{\pi}{2}\right[$.
Pour $x\in\left]\frac{\pi}{2},\pi\right]$, $f_4(x)=-\tan x+\cos x$ et $f_4$ est strictement décroissante sur
$\left]\frac{\pi}{2},\pi\right]$ en tant que somme de deux fonctions strictement décroissantes sur $\left]\frac{\pi}{2},\pi\right]$.
On a immédiatement
$\displaystyle\lim_{\substack{x\rightarrow\frac{\pi}{2}\\
x<\frac{\pi}{2}}}f_4(x)=\displaystyle\lim_{\substack{x\rightarrow\frac{\pi}{2}\\
x>\frac{\pi}{2}}}f_4(x)=+\infty$.
On en déduit $\mathcal{C}_4$.

$$\includegraphics{../images/img005102-4}$$

 \item  Soit $x>0$. $x$ n'est pas nul donc $\frac{1}{x}$ existe puis $1+\frac{1}{x}>0$ et $f_6(x)$ existe.
\textbf{Etude en} $\bf{0.}$ Pour $x>0$, $x\ln(1+\frac{1}{x})=-x\ln x+x\ln(1+x)$. Par suite, $x\ln(1+\frac{1}{x})$ tend
vers $0$ quand $x$ tend vers $0$ par valeurs supérieures et donc $f_5(x)=\mbox{exp}(x\ln(1+\frac{1}{x}))$ tend vers
$1$.
Posons encore $f_5(0)=1$ et étudions la dérivabilité de $f_5$ en $0$. Pour $x>0$,

$$\frac{f_5(x)-f_5(0)}{x-0}=\frac{1}{x}\left(\text{exp}(x\ln(1+\frac{1}{x}))-1\right)
=\frac{\text{exp}\left(x\ln(1+\frac{1}{x})\right)-1}{x\ln\left(1+\frac{1}{x}\right)}\ln\left(1+\frac{1}{x}\right).$$
Or, $x\ln\left(1+\frac{1}{x}\right)$ tend vers $0$ quand $x$ tend vers $0$, et donc

$$\lim_{\substack{x\rightarrow0\\
x>0}}\frac{\mbox{exp}(x\ln\left(1+\frac{1}{x})\right)-1}{x\ln\left(1+\frac{1}{x}\right)}=\lim_{y\rightarrow 0}\frac{e^y-1}{y}=1.$$
D'autre part, $\ln\left(1+\frac{1}{x}\right)$ tend vers $+\infty$ quand $x$ tend vers $0$ par valeurs supérieures. Finalement,

$$\lim_{\substack{x\rightarrow0\\
x>0}}\frac{f_5(x)-f_5(0)}{x-0}=+\infty.$$
Ainsi, $f_5$ n'est pas dérivable en $0$ mais $\mathcal{C}_5$ admet l'axe des ordonnées pour tangente en
$(0,f_5(0))=(0,1)$.
\textbf{Etude en} $+\infty.$ Pour $x>0$,
$x\ln\left(1+\frac{1}{x}\right)=\frac{\ln\left(1+\frac{1}{x}\right)}{\frac{1}{x}}$ et 
donc $\lim_{x\rightarrow +\infty}x\ln\left(1+\frac{1}{x}\right)=\lim_{y\rightarrow 0}\frac{\ln(1+y)}{y}=1$.
Par suite,

$$\lim_{x\rightarrow +\infty}f_5(x)=e.$$
\textbf{Etude des variations de} $\bf{f_5.}$ Pour $x>0$, $f_5(x)>0$ puis $\ln(f_5(x))=x\ln(1+\frac{1}{x})$. Par suite,
pour $x>0$,

$$f_5'(x)=f_5(x)\ln(f_5)'(x)=f_5(x)\left(\ln\left(1+\frac{1}{x}\right)+\frac{x(-\frac{1}{x^2})}{1+\frac{1}{x}}\right)=f_5(x)g(x),$$
où $g(x)=\ln\left(1+\frac{1}{x}\right)-\frac{1}{1+x}$. Sur $]0,+\infty[$, $f_5'$ est du signe de $g$.
Pour déterminer le signe de $g$, étudions d'abord les variations de $g$ sur $]0,+\infty[$. $g$ est dérivable sur
$]0,+\infty[$ et pour $x>0$,

$$g'(x)=\frac{-\frac{1}{x^2}}{1+\frac{1}{x}}+\frac{1}{(x+1)^2}=-\frac{1}{x(x+1)}+\frac{1}{(x+1)^2}
=\frac{-1}{x(x+1)^2}<0.$$
$g$ est donc strictement décroissante sur $]0,+\infty[$, et puisque $\lim_{x\rightarrow +\infty}g(x)=0$, $g$ est strictement
positive sur $]0,+\infty[$. Il en est de même de $f_5'$. $f_5$ est strictement croissante sur $]0,+\infty($.
On en déduit $\mathcal{C}_5$.

$$\includegraphics{../images/img005102-5}$$


 \item  \textbf{Domaine de définition de} $\bf{f_6.}$ Soit $x\in\Rr$.

\begin{align*}
f_6(x)\;\mbox{existe}&\Leftrightarrow x^2-5x+6>0\;\mbox{et}\;1-\mbox{log}_{\frac{1}{2}}(x^2-5x+6)>0
\Leftrightarrow x^2-5x+6>0\;\mbox{et}\;\frac{\ln(x^2-5x+6)}{\ln\frac{1}{2}}<1\\
 &\Leftrightarrow x^2-5x+6>0\;\mbox{et}\;\ln(x^2-5x+6)>\ln\frac{1}{2}\Leftrightarrow x^2-5x+6>\frac{1}{2}\\
 &\Leftrightarrow x^2-5x+\frac{11}{2}>0\Leftrightarrow x\in]-\infty,\frac{5-\sqrt{3}}{2}[\cup]\frac{5+\sqrt{3}}{2},+\infty[=\mathcal{D}_f.
\end{align*}
\textbf{Variations de} $\bf{f_6.}$ La fonction $x\mapsto x^2-5x+6$ est strictement décroissante sur
$\left]-\infty,\frac{5}{2}\right]$ et strictement croissante sur $\left[\frac{5}{2},+\infty\right[$. Comme
$\frac{5+\sqrt{3}}{2}>\frac{5}{2}$ et que $\frac{5-\sqrt{3}}{2}<\frac{5}{2}$, la fonction $x\mapsto x^2-5x+6$ est
strictement décroissante sur$\left]-\infty,\frac{5-\sqrt{3}}{2}\right]$ et strictement croissante sur
$\left[\frac{5+\sqrt{3}}{2},+\infty\right[$, à valeurs dans $]0,+\infty[$, intervalle sur lequel la fonction logarithme néperien
est strictement croissante. La fonction $x\mapsto1+\frac{\ln(x^2-5x+6)}{\ln2}$ a le même sens de variations et
finalement $f_6$ est strictement décroissante sur $\left]-\infty,\frac{5-\sqrt{3}}{2}\right]$ et strictement croissante
sur $\left[\frac{5+\sqrt{3}}{2},+\infty\right[$.
\textbf{Axe de symétrie} Soit $x\in\Rr$. $x\in\mathcal{D}_f\Leftrightarrow\frac{5}{2}-x\in\mathcal{D}_f$ et de 
plus, $\left(\frac{5}{2}-x\right)^2-5\left(\frac{5}{2}-x\right)+6=x^2-5x+6$. Par suite,

$$\forall x\in D,\;f_6(\frac{5}{2}-x)=f_6(x).$$

$\mathcal{C}_6$ admet donc la droite d'équation $x=\frac{5}{2}$ pour axe de symétrie.

Le calcul des limites étant immédiat, on en déduit $\mathcal{C}_6$.

$$\includegraphics{../images/img005102-6}$$

\end{enumerate}
\fincorrection
\correction{005404}
On a $0\leq f(0)\leq 1$ et $0\leq f(1)\leq 1$. Donc $|f(1)-f(0)|\leq1$. Mais, par hypothèse, $|f(1)-f(0)|\geq1$. Par suite, $|f(1)-f(0)|= 1$ et nécessairement, $(f(0),f(1))\in\{(0,1),(1,0)\}$.

Supposons que $f(0)=0$ et $f(1)=1$ et montrons que $\forall x\in[0,1],\;f(x)=x$.

Soit $x\in[0,1]$. On a $|f(x) -f(0)|\geq|x-0|$ ce qui fournit $f(x)\geq x$. On a aussi $|f(x)-f(1)|=|x-1|$ ce qui fournit $1-f(x)\geq 1-x$ et donc $f(x)\leq x$. Finalement, $\forall x\in[0,1],\;f(x)=x$ et $f=Id$.

Si $f(0)=1$ et $f(1)=0$, posons pour $x\in[0,1]$, $g(x)=1-f(x)$. Alors, $g(0)=0$, $g(1)=1$ puis, pour $x\in[0,1]$, $g(x)\in[0,1]$. Enfin,
 
$$\forall(x,y)\in[0,1]^2,\;|g(y)-g(x)|=|f(y)-f(x)|\geq|y-x|.$$
 
D'après l'étude du premier cas, $g=Id$ et donc $f=1-Id$. Réciproquement, $Id$ et $1-Id$ sont bien bien solutions du problème.
\fincorrection
\correction{005443}

\begin{enumerate}

\item  $f_1$ est définie et de classe $C^\infty$ sur $\Rr^*$ en vertu de théorèmes généraux. De plus, $f_1$ est paire. On étudiera $f_1$ sur $[0,+\infty[$ (se méfier alors pour la dérivabilité en $0$).

\textbf{Etude en 0} (à gauche et à droite).
\begin{align*}\ensuremath
f_1(x)&\underset{x\rightarrow0}{=}\frac{1}{x^3}(1+x^2)[x-\frac{x^3}{3}+\frac{x^5}{5}+o(x^5)-x(1-x^2+x^4+o(x^4))]\\
 &=(1+x^2)(\frac{2x^3}{3}-\frac{4x^5}{5}+o(x^5))=(1+x^2)(\frac{2}{3}-\frac{4x^2}{5}+o(x^2))\\
 &=\frac{2}{3}-\frac{2x^2}{15}+o(x^2).
\end{align*}

Par suite, $f_1$ se prolonge par continuité en $0$ en posant $f_1(0)=\frac{2}{3}$. Puisque $f_1$ admet en $0$ un développement limité d'ordre $1$, le prolongement encore noté $f_1$ est dérivable en $0$ et $f_1'(0)=0$. $C_1$ admet au point d'abscisse $0$ une tangente parallèle à $(0x)$ d'équation $y=\frac{2}{3}$. Enfin, puisque  $f(x)-\frac{2}{3}$ est, au voisinage de $0$, du signe de $-\frac{2x^2}{15}$, la courbe est localement en dessous de sa tangente.

\textbf{Etude en} $\bf{+\infty}$ (et $-\infty$). $f_1(x)\underset{x\rightarrow+\infty}{\sim}\frac{\pi}{2x}\underset{x\rightarrow+\infty}{\rightarrow}0$, et de même $f_1(x)\underset{x\rightarrow-\infty}{\rightarrow}0$.

\textbf{Dérivée, variations.}

Pour $x>0$, 

\begin{align*}\ensuremath
f'(x)&=(-\frac{3}{x^4}-\frac{1}{x^2})
(\Arctan x-\frac{x}{1+x^2})+\frac{1+x^2}{x^3}(\frac{1}{1+x^2}-\frac{1-x^2}{(1+x^2)^2})\\
 &=-\frac{3+x^2}{x^4}(\Arctan x-\frac{x}{1+x^2})+\frac{1+x^2}{x^3}\frac{2x^2}{(1+x^2)^2}\\
 &=\frac{3+x^2}{x^4}(-\Arctan x+\frac{x}{1+x^2}+\frac{x^4}{3+x^2}\frac{2}{x(1+x^2)})\\
 &=\frac{3+x^2}{x^4}(-\Arctan x+\frac{x(3+x^2)+2x^3}{(1+x^2)(3+x^2))}=\frac{3+x^2}{x^4}g(x)
\end{align*}

où, pour tout réel $x$, $g(x)=-\Arctan x+\frac{3x}{3+x^2}$.

$g$ est dérivable sur $\Rr$ et pour $x$ réel,

$$g'(x)=3\frac{(3+x^2)-2x^2}{(3+x^2)^2}-\frac{1+x^2}=\frac{3(3-x^2)(1+x^2)-(3+x^2)^2}{(1+x^2)(3+x^2)^2}=\frac{-4x^4}{(3+x^2)^2(1+x^2)}.$$

$g'$ est donc strictement négative sur $]0,+\infty[$ et par suite, $g$ est donc strictement décroisante sur $[0,+\infty[$. Puisque $g(0)=0$, pour $x>0$, $g(x)<0$. Finalement, $f_1'$ est strictement négative sur $]0,+\infty[$ et $f_1$ est strictement décroissante sur $[0,+\infty[$.

Le tableau de variations de $f_1$ n'apporte rien de plus.

\textbf{Graphe}

$$\includegraphics{../images/img005443-4}$$

%%%%%%%%%%%%%%%%%%%%%%%%%%%%%%%%%%%%%%%%%%%%%%%%%%%%%%%%%%%%%%%%%%%%%%%%%%%%%%%%%%%%%%%%%%%%%%%%%%%%%%%%%%%%%%%%%%%%%

\item  $f_2$ est définie sur $D=\Rr\setminus(\frac{\pi}{2}+\pi\Zz)$, paire et $2\pi$-périodique. $f_2$ est continue sur $D$ en vertu de théorèmes généraux. On étudie $f_2$ sur $[0,\frac{\pi}{2}[\cup]\frac{\pi}{2},\pi]$.

\textbf{Etude en} $\bf{\frac{{\pi}}{2}}$.

$f(x)\underset{x\rightarrow\pi/2}{\sim}|\tan x|$ et donc, $\lim_{x\rightarrow \pi/2}f(x)=+\infty$. $C_2$ admet la droite d'équation $x=\frac{\pi}{2}$ pour droite asymptote.

\textbf{Dérivabilité et dérivée.}

$f_2$ est dérivable sur $\Rr\setminus\frac{\pi}{2}\Zz$ en vertu de théorèmes généraux et pour $x\notin\frac{\pi}{2}\Zz$, $f_2'(x)=\varepsilon\frac{1}{\cos^2x}-\sin x$ où $\varepsilon$ est le signe de $\tan x$.

$f_2$ est aussi dérivable à droite en $0$ et $(f_2)_d'(0)=1$. Par symétrie, $f_2$ est dérivable à gauche en $0$ et $(f_2)_g'(0)=-1$. $f_2$ n'est pas dérivable en $0$.

De même, $f_2$ est dérivable à gauche et à droite en $\pi$ avec $(f_2)_g'(\pi)=-1$ et $(f_2)_d'(\pi)=1$, et n'est donc pas dérivable en $\pi$.

\textbf{Variations}.

$f_2$ est strictement décroissante sur $]\frac{\pi}{2},\pi]$ en tant que somme de deux fonctions strictement décroissantes sur $]\frac{\pi}{2},\pi]$. Puis, pour $x$ élément de $]0,\frac{\pi}{2}[$, $f_2'(x)=\frac{1}{\cos^2x}-\sin x>1-1=0$. $f_2'$ est strictement positive sur $]0,\frac{\pi}{2}[$ et donc $f_2$ est strictement croissante sur $[0,\frac{\pi}{2}[$.

\textbf{Graphe.}

$$\includegraphics{../images/img005443-5}$$

%%%%%%%%%%%%%%%%%%%%%%%%%%%%%%%%%%%%%%%%%%%%%%%%%%%%%%%%%%%%%%%%%%%%%%%%%%%%%%%%%%%%%%%%%%%%%%%%%%%%%%%%%%%%%%%%%%%%

\item  Pour $x$ réel, posons $P(x)=x^3+12x^2+60x+120$. Pour tout réel $x$, on a $P'(x)=3(x^2+8x+20)=3((x+4)^2+4)>0$. $P$ est une fonction polynôme de degré $3$ strictement croissante sur $\Rr$ et s'annule donc une et une seule fois en un certain réel noté $\alpha$. De plus, $P(-5)P(-4)<0$ et $\alpha\in]-5,-4[$.

Enfin, $P$ est strictement négatif sur $]-\infty,\alpha[$ et strictement positif sur $]\alpha,+\infty[$.

$f_3$ est définie sur $\Rr\setminus\{-\alpha,\alpha\}$, et pour $x\in\Rr\setminus\{-\alpha,\alpha\}$, 

$$f_3(x)=x-\ln\left|\frac{P(x)}{P(-x)}\right|=x-\ln|P(x)|+\ln|P(-x)|.$$
 
Notons que $f_3$ est impaire.
 
\textbf{Dérivabilité et dérivée.}

$f_3$ est de classe $C^\infty$ sur $\Rr\setminus\{-\alpha,\alpha\}$ en vertu de théorèmes généraux et pour $x\in\Rr\setminus\{-\alpha,\alpha\}$,

$$f_3'(x)=1-\frac{P'(x)}{P(x)}-\frac{P'(-x)}{P(-x)}=\frac{P(x)P(-x)-P'(x)P(-x)-P'(-x)P(x)}{P(-x)P(x)}.$$

Puis,

\begin{align*}\ensuremath 
P(x)P(-x)&-P'(x)P(-x)-P'(-x)P(x)\\
 &=((12x^2+120)+(x^3+60x))((12x^2+120)-(x^3+60x))\\
 &\;-3((x^2+20)+8x)((12x^2+120)-(x^3+60x)))-
3((x^2+20)-8x))((12x^2+120)+(x^3+60x))\\
 &=144(x^2+10)^2-x^2(x^2+60)^2-6((x^2+20)(12x^2+120)-(8x)(x^3+60x))\\
 &=(-x^6+24x^4-720x^2+14400)-6(4x^4-120x^2+2400)=-x^6,
\end{align*}

et donc $f_3'(x)=\frac{-x^6}{P(x)P(-x)}$.

\textbf{Etude en} $\bf{+{\infty}}$.

$$f_3(x)-x\underset{x\rightarrow+\infty}{=}-\ln(1+\frac{12}{x}+o(\frac{1}{x}))+\ln(1-\frac{12}{x}+o(\frac{1}{x}))=-\frac{24}{x}+o(\frac{1}{x}).$$

On en déduit tout d'abord que $\lim_{x\rightarrow +\infty}f_3(x)=+\infty$ (resp.$\lim_{x\rightarrow -\infty}f_3(x)=-\infty$, puis que $C_3$ admet en $+\infty$ (resp.$-\infty$) la droite d'équation $y=x$ pour droite asymptote et que $C_3$ est au-dessous (resp. au-dessus) de cette droite au voisinage de $+\infty$ (resp.$-\infty$).

\textbf{Variations.}

D'une part, $f_3'(0)=0$. D'autre part, pour $x>0$, $P(x)>0$. $f_3'$ est donc du signe de $-P(-x)$ sur $]0,+\infty[\setminus\{\alpha\}$. Ainsi, $f_3'$ est strictement négative sur $]0,\alpha[$ et strictement positive sur $]\alpha,+\infty[$.

On en déduit le tableau de variations de $f_3$.


$$\includegraphics{../images/img005443-1}$$


\textbf{Graphe}.

$$\includegraphics{../images/img005443-6}$$


%%%%%%%%%%%%%%%%%%%%%%%%%%%%%%%%%%%%%%%%%%%%%%%%%%%%%%%%%%%%%%%%%%%%%%%%%%%%%%%%%%%%%%%%%%%%%%%%%%%%%%%%%%%%%%%%%%%%%%

\item  $f_4$ est définie sur $\Rr\setminus\{-1,1\}$. De plus, pour $x\neq 0$, 

$$f_4(\frac{1}{x})=\frac{1}{x}e^{\frac{2/x}{1/x^2-1}}=\frac{1}{x}e^{-\frac{2x}{x^2-1}}=\frac{1}{f_4(x)}.$$

Ce genre de constatation peut servir à calculer $\lim_{x\rightarrow +\infty}f_4(x)$ si l'on connait $\lim_{x\rightarrow 0,\;x>0}f_4(x)$, obtenir les variations de $f_4$ sur $]0,1[$ si on les connait sur $]1,+\infty[$ ...

On peut aussi noter que $\forall x\in\Rr\setminus\{-1,1\}$, $f_4(-x)f_4(x)=-x^2$ et donc, pour $x\neq 0$, $f_4(-x)=\frac{-x^2}{f_4(x)}$. Cette constatation pourra être utile pour déduire l'étude de $f_4$ en $-1$ de l'étude en $1$.

\textbf{Etude en} $\bf{+{\infty}}$ \textbf{et} $\bf{-{\infty}}$.

Puisque $\frac{2x}{x^2-1}\underset{x\rightarrow\pm\infty}{\rightarrow}0$, on a  $f_4(x)\underset{x\rightarrow\pm\infty}{\sim}x$ ce qui montre déjà que $\lim_{x\rightarrow +\infty}f_4(x)=+\infty$, $\lim_{x\rightarrow -\infty}f_4(x)=-\infty$ et que $C_4$ admet en $+\infty$ et $-\infty$, une direction asymptotique d'équation $y=x$. Plus précisément,

$$\frac{2x}{x^2-1}\underset{x\rightarrow\pm\infty}{=}\frac{2}{x}(1-\frac{1}{x^2})^{-1}=\frac{2}{x}+o(\frac{1}{x^2}),$$

puis,

$$e^{\frac{2x}{x^2-1}}\underset{x\rightarrow\pm\infty}{=}1+(\frac{2}{x})+(\frac{2}{x})^2+ o(\frac{1}{x^2})=1+\frac{2}{x}+\frac{2}{x^2}+ o(\frac{1}{x^2}).$$

On en déduit que

$$f_4(x)\underset{x\rightarrow\pm\infty}{=}x+2+\frac{2}{x}+o(\frac{1}{x}).$$

Par suite, $C_4$ admet la droite d'équation $y=x+2$ pour droite asymptote en $+\infty$ et $-\infty$. De plus, le signe de $f_4(x)-(x+2)$ étant localement le signe de $\frac{2}{x}$, $C_4$ est au-dessus de son asymptote au voisinage de $+\infty$ et au-dessous au voisinage de $-\infty$.

\textbf{Etude en 1 (et -1).}

Clairement, $\lim_{x\rightarrow 1,\;x>1}f_4(x)=+\infty$ et $\lim_{x\rightarrow -1,\;x>-1}f_4(x)=-\infty$. Ensuite, $\lim_{x\rightarrow 1,\;x<1}f_4(x)=0$ et $\lim_{x\rightarrow -1,\;x<-1}f_4(x)=0$.

On prolonge $f_4$ par continuité à gauche en $1$ en posant $f_4(1)=0$, et de même en $-1$ et on étudie la dérivabilité du prolongement encore noté $f_4$.

$f_4$ est continue sur $]-1,1]$, de classe $C^1$ sur $]-1,1[$ et pour $x\in]-1,1[$ (voir dérivée-variations), 

$$f_4'(x)=\frac{x^4-2x^3-2x^2-2x+1}{(x^2-1)^2}e^{\frac{2x}{x^2-1}}\underset{x\rightarrow1,\;x<1}{\rightarrow}0.$$

D'après un théorème classique d'analyse, $f_4$ est de classe $C^1$ sur $]-1,1]$ et en particulier dérivable à gauche en $1$ et $f_g'(1)=0$.

De même, $f_4$ est dérivable à gauche en $-1$ et $f_g'(-1)=0$. $C_4$ admet en ces points des demi-tangentes parallèles à l'axe $(Ox)$.

\textbf{Dérivée. Variations.}

$f_4$ est de classe $C^1$ sur $\Rr\setminus\{-1,1\}$ en vertu de théorèmes généraux et pour $x\neq 0$,

\begin{align*}\ensuremath
\frac{f_4'(x)}{f_4(x)}&=(\ln|f_4|)'(x)=(\ln|x|+\frac{2x}{x^2-1})'(x)=\frac{1}{x}+2\frac{(x^2-1)-x(2x)}{(x^2-1)^2}\\
 &=\frac{(x^2-1)^2-2x(x^2+1)}{x(x^2-1)^2}=\frac{x^4-2x^3-2x^2-2x+1}{x(x^2-1)^2},
\end{align*}
 
et donc 

$$\forall x\neq 0,\;f4'(x)=\frac{x^4-2x^3-2x^2-2x+1}{(x^2-1)^2}e^{\frac{2x}{x^2-1}},$$

ce qui reste vrai pour $x=0$ par continuité de $f_4'$ en $0$.

$f_4'$ est donc du signe de $P(x)=x^4-2x^3-2x^2-2x+1$. Or, pour $x\neq0$,

\begin{align*}\ensuremath 
P(x)&=x^2((x^2+\frac{1}{x^2})-2(x+\frac{1}{x})-2)=x^2((x+\frac{1}{x})^2-2(x+\frac{1}{x})-4)=\\
 &=x^2(x+\frac{1}{x}-(1-\sqrt{5})(x+\frac{1}{x}-(1+\sqrt{5}))=(x^2-(1-\sqrt{5})x+1)(x^2-(1+\sqrt{5})x+1),
\end{align*}

ce qui reste vrai pour $x=0$.

Le premier trinôme a un discriminant égal à $(\sqrt{5}-1)^2-4=2-2\sqrt{5}<0$ et donc $\forall x\in\Rr,\;x^2-(1-\sqrt{5})x+1>0$.

Le deuxième trinôme a un discriminant égal à $(\sqrt{5}+1)^2-4=2+2\sqrt{5}>0$ et admet donc deux racines réelles $\alpha=\frac{1}{2}(1+\sqrt{5}+\sqrt{2+\sqrt{5}})2,89...>1$ et $\beta=\frac{1}{2}(1+\sqrt{5}-\sqrt{2+2\sqrt{5}})=\frac{1}{\alpha}0,34...\in]0,1[$.
On en déduit le tableau de variation de $f_4$.

$$\includegraphics{../images/img005443-2}$$

\textbf{Graphe.}

$$\includegraphics{../images/img005443-7}$$


%%%%%%%%%%%%%%%%%%%%%%%%%%%%%%%%%%%%%%%%%%%%%%%%%%%%%%%%%%%%%%%%%%%%%%%%%%%%%%%%%%%%%%%%%%%%%%%%%%%%%%%%%%%%%%%%
\item  Si $x>0$, $e^x-1>0$ et si $x<0$, $e^x-1<0$. Donc, pour $x\neq 0$, $ > 0$ et $f_5$ est définie sur $\Rr^*$.
Pour $x\neq 0$, 

$$f_5(-x)=-\frac{1}{x}\ln\frac{e^{-x}-1}{-x}=-\frac{1}{x}\ln(e^{-x})-\frac{1}{x}\ln\frac{e^x-1}{x}=1-f(x).$$

Donc, pour tout réel non nul $x$, $f(x)+f(-x)=1$. Le point de coordonnées $(0,\frac{1}{2})$ est centre de 
symétrie de $C_5$.

\textbf{Etude en 0.}

$$f_5(x)\underset{x\rightarrow0}{=}\frac{1}{x}\ln(1+\frac{x}{2}+\frac{x^2}{6}+o(x^2))=\frac{1}{x}((\frac{x}{2}+\frac{x^2}{6})-\frac{1}{2}(\frac{x}{2})^2+o(x^2))=\frac{1}{2}+\frac{1}{24}x+o(x).$$

Ainsi, $f_5$ se prolonge par continuité en $0$ en posant $f_5(0)=\frac{1}{2}$. Le prolongement, encore noté $f_5$, admet en $0$ un développement limité d'ordre $1$ et est donc dérivable en $0$ avec $f_5'(0)=\frac{1}{24}$. Une équation de la tangente à $C_5$ en le point d'abscisse $0$ est $y=\frac{1}{24}x+\frac{1}{2}$. Par symétrie, ce point est un point d'inflexion.

\textbf{Etude en} $\bf{+{\infty}}$.

$$f_5(x)\underset{x\rightarrow+\infty}{=}\frac{1}{x}(\ln(e^x)+\ln(1-e^{-x})-\ln x)=1-\frac{\ln x}{x}+\frac{\ln(1-e^{-x}}{x}=1+o(1).$$

Donc, $\lim_{x\rightarrow +\infty}f_5(x)=1$. Par symétrie, $\lim_{x\rightarrow -\infty}f_5(x)=\lim_{x\rightarrow -\infty}(1-f_5(-x))=1-1=0$.

\textbf{Dérivée. Variations.}

$f_5$ est dérivable sur $\Rr^*$ en vertu de théorèmes généraux (et donc sur $\Rr$) et pour $x\neq 0$, (puisque $\ln\frac{e^x-1}{x}=\ln\left|\frac{e^x-1}{x}\right|=\ln|e^x-1|-\ln|x|$),

$$f_5'(x)=-\frac{1}{x^2}\ln\frac{e^x-1}{x}+\frac{1}{x}(\frac{e^x}{e^x-1}-\frac{1}{x})
=\frac{1}{x^2}(-\ln\frac{e^x-1}{x}+\frac{xe^x}{e^x-1}-1).$$

$f_5'$ est, sur $\Rr^*$, du signe de $g(x)=-\ln\frac{e^x-1}{x}+\frac{xe^x}{e^x-1}+-1$. $g$ est dérivable sur $\Rr^*$ et pour $x$ réel non nul,

\begin{align*}\ensuremath
g'(x)&=-\frac{e^x}{e^x-1}+\frac{1}{x}+\frac{(e^x+xe^x)(e^x-1)-xe^x.e^x}{(e^x-1)^2}
=\frac{-xe^x(e^x-1)+(e^x-1)^2+xe^x(e^x-x-1)}{x(e^x-1)^2}\\
 &=\frac{(e^x-1)^2-x^2e^x}{x(e^x-1)^2}=\frac{(e^{x/2}-e^{-x/2})^2-x^2}{x(e^{x/2}-e^{-x/2})^2}\\
 &=\frac{(2\sh\frac{x}{2})^2-x^2}{x(2\sh\frac{x}{2})^2}=\frac{\sh^2\frac{x}{2}-(\frac{x}{2})^2}{x\sh^2\frac{x}{2}}.
\end{align*} 
 
L'inégalité $\sh x>x$, valable pour $x>0$, est classique (par exemple, la formule de \textsc{Taylor}-\textsc{Laplace} à l'ordre 1 fournit pour $x>0$, $\sh x=x+\int_{0}^{x}(x-t)\sh t\;dt>x$.) Par suite, $g'$ est strictement positive sur $]0,+\infty[$, et donc $g$ est strictement croissante sur $]0,+\infty[$. En tenant compte de $g(0^+)= 0$, $g$ est donc strictement positive sur $]0,+\infty[$. Il en est de même de $f_5'$ et $f_5$ est strictement croissante sur $]0,+\infty[$.

Par symétrie et continuité en $0$, $f_5$ est strictement croissante sur $\Rr$.

\textbf{Graphe.}

$$\includegraphics{../images/img005443-8}$$


%%%%%%%%%%%%%%%%%%%%%%%%%%%%%%%%%%%%%%%%%%%%%%%%%%%%%%%%%%%%%%%%%%%%%%%%%%%%%%%%%%%%%%%%%%%%%%%%%%%%%%%%%%%%%%%%%%%%%

\item  $f_6$ est définie et continue sur $\Rr$, dérivable sur $\Rr\{-1,1\}$ en vertu de théorèmes généraux.

\textbf{Etude en 1.}

$f_6(x)-f_6(1)=x-1+\sqrt{|x^2-1|}\underset{x\rightarrow1}{\sim}\sqrt{2}\sqrt{|x-1|}$, ce qui montre que $f_6$ n'est pas dérivable en $1$ mais que $C_6$ admet au point d'abscisse $1$ deux demi-tangentes parallèles à $(Oy)$.

\textbf{Etude en -1.}

$f_6(x)-f_6(-1)=x+1+\sqrt{|x^2-1|}\underset{x\rightarrow-1}{\sim}\sqrt{2}\sqrt{|x+1|}$,  ce qui montre que $f_6$ n'est pas dérivable en $-1$ mais que $C_6$ admet au point d'abscisse $-1$ deux demi-tangentes parallèles à $(Oy)$.

\textbf{Etude en} $\bf{+{\infty}}$.

Au voisinage de $+\infty$, on a

$$f_6(x)=x+x(1-\frac{1}{x^2})^{1/2}=x+x(1-\frac{1}{2x^2}+o(\frac{1}{x^2}))=2x-\frac{1}{2x}+o(\frac{1}{x}),$$

ce qui montre tout à la fois que $\lim_{x\rightarrow +\infty}f_6(x)=+\infty$, puis que la droite d'équation $y=2x$ est asymptote à $C_6$ en $+\infty$ et que $C_6$ est au-dessous de cette droite au voisinage de $+\infty$.

\textbf{Etude en} $\bf{-{\infty}}$.

Au voisinage de $-\infty$, on a, $f_6(x)=x-x(1+o(\frac{1}{x}))=o(1)$, et $\lim_{x\rightarrow -\infty}f_6(x)=0$.

\textbf{Dérivée. Variations.}

Soit $\varepsilon$ le signe de $x^2-1$. Pour $x\neq\pm1$, 

$$f_6'(x)=1+\frac{2\varepsilon x}{2\sqrt{\varepsilon(x^2-1)}}=\frac{\sqrt{\varepsilon(x^2-1)}+\varepsilon x}{\sqrt{\varepsilon(x^2-1)}}.$$

Si $-1<x\leq0$, (de sorte que $\varepsilon x>0$) ou $x>1$, $f_6'(x)>0$.

Si $x<-1$, $\mbox{sgn}(f_6'(x))=\mbox{sgn}(x+\sqrt{x^2-1})=\mbox{sgn}(\frac{1}{x-\sqrt{x^2-1}}=-$ et $f_6'(x)<0$.

Si $0\leq x<1$. $\mbox{sgn}(f_6'(x))=\mbox{sgn}(-x+\sqrt{x^2-1})=\mbox{sgn}(-x^2-(x^2-1))=\mbox{sgn}(\frac{1}{\sqrt{2}}-x)$.

D'où le tableau de variations de $f_6$~:

$$\includegraphics{../images/img005443-3}$$



\textbf{Graphe.}

$$\includegraphics{../images/img005443-9}$$



%%%%%%%%%%%%%%%%%%%%%%%%%%%%%%%%%%%%%%%%%%%%%%%%%%%%%%%%%%%%%%%%%%%%%%%%%%%%%%%%%%%%%%%%%%%%%%%%%%%%%%%%%%%%%%%%%%%%%

\item 

$$\includegraphics{../images/img005443-10}$$


%%%%%%%%%%%%%%%%%%%%%%%%%%%%%%%%%%%%%%%%%%%%%%%%%%%%%%%%%%%%%%%%%%%%%%%%%%%%%%%%%%%%%%%%%%%%%%%%%%%%%%%%%%%%%%%%%%%%%

\item 

%%%%%%%%%%%%%%%%%%%%%%%%%%%%%%%%%%%%%%%%%%%%%%%%%%%%%%%%%%%%%%%%%%%%%%%%%%%%%%%%%%%%%%%%%%%%%%%%%%%%%%%%%%%%%%%%%%%%%

\item 

%%%%%%%%%%%%%%%%%%%%%%%%%%%%%%%%%%%%%%%%%%%%%%%%%%%%%%%%%%%%%%%%%%%%%%%%%%%%%%%%%%%%%%%%%%%%%%%%%%%%%%%%%%%%%%%%%%%%%

\item  

$$\includegraphics{../images/img005443-11}$$


%%%%%%%%%%%%%%%%%%%%%%%%%%%%%%%%%%%%%%%%%%%%%%%%%%%%%%%%%%%%%%%%%%%%%%%%%%%%%%%%%%%%%%%%%%%%%%%%%%%%%%%%%%%%%%%%%%%%%

\item  

$$\includegraphics{../images/img005443-12}$$


%%%%%%%%%%%%%%%%%%%%%%%%%%%%%%%%%%%%%%%%%%%%%%%%%%%%%%%%%%%%%%%%%%%%%%%%%%%%%%%%%%%%%%%%%%%%%%%%%%%%%%%%%%%%%%%%%%%%%

\item 

$$\includegraphics{../images/img005443-13}$$



%%%%%%%%%%%%%%%%%%%%%%%%%%%%%%%%%%%%%%%%%%%%%%%%%%%%%%%%%%%%%%%%%%%%%%%%%%%%%%%%%%%%%%%%%%%%%%%%%%%%%%%%%%%%%%%%%%%%%

\item 

%%%%%%%%%%%%%%%%%%%%%%%%%%%%%%%%%%%%%%%%%%%%%%%%%%%%%%%%%%%%%%%%%%%%%%%%%%%%%%%%%%%%%%%%%%%%%%%%%%%%%%%%%%%%%%%%%%%%%

\item 

%%%%%%%%%%%%%%%%%%%%%%%%%%%%%%%%%%%%%%%%%%%%%%%%%%%%%%%%%%%%%%%%%%%%%%%%%%%%%%%%%%%%%%%%%%%%%%%%%%%%%%%%%%%%%%%%%%%%%

\item 

%%%%%%%%%%%%%%%%%%%%%%%%%%%%%%%%%%%%%%%%%%%%%%%%%%%%%%%%%%%%%%%%%%%%%%%%%%%%%%%%%%%%%%%%%%%%%%%%%%%%%%%%%%%%%%%%%%%%%

\item

%%%%%%%%%%%%%%%%%%%%%%%%%%%%%%%%%%%%%%%%%%%%%%%%%%%%%%%%%%%%%%%%%%%%%%%%%%%%%%%%%%%%%%%%%%%%%%%%%%%%%%%%%%%%%%%%%%%%%

\item

%%%%%%%%%%%%%%%%%%%%%%%%%%%%%%%%%%%%%%%%%%%%%%%%%%%%%%%%%%%%%%%%%%%%%%%%%%%%%%%%%%%%%%%%%%%%%%%%%%%%%%%%%%%%%%%%%%%%%

\item

%%%%%%%%%%%%%%%%%%%%%%%%%%%%%%%%%%%%%%%%%%%%%%%%%%%%%%%%%%%%%%%%%%%%%%%%%%%%%%%%%%%%%%%%%%%%%%%%%%%%%%%%%%%%%%%%%%%%%

\item

%%%%%%%%%%%%%%%%%%%%%%%%%%%%%%%%%%%%%%%%%%%%%%%%%%%%%%%%%%%%%%%%%%%%%%%%%%%%%%%%%%%%%%%%%%%%%%%%%%%%%%%%%%%%%%%%%%%%%

\item

%%%%%%%%%%%%%%%%%%%%%%%%%%%%%%%%%%%%%%%%%%%%%%%%%%%%%%%%%%%%%%%%%%%%%%%%%%%%%%%%%%%%%%%%%%%%%%%%%%%%%%%%%%%%%%%%%%%%%

\item

%%%%%%%%%%%%%%%%%%%%%%%%%%%%%%%%%%%%%%%%%%%%%%%%%%%%%%%%%%%%%%%%%%%%%%%%%%%%%%%%%%%%%%%%%%%%%%%%%%%%%%%%%%%%%%%%%%%%%

\item

%%%%%%%%%%%%%%%%%%%%%%%%%%%%%%%%%%%%%%%%%%%%%%%%%%%%%%%%%%%%%%%%%%%%%%%%%%%%%%%%%%%%%%%%%%%%%%%%%%%%%%%%%%%%%%%%%%%%%

\item

$$\includegraphics{../images/img005443-14}$$


\end{enumerate}

\fincorrection
\nocorrection
\nocorrection
\nocorrection
\nocorrection
\nocorrection
\nocorrection
\nocorrection
\nocorrection
\correction{001216}
$$\lim _{x\rightarrow \infty }\sqrt [3]{{x}^{3}+1}-\sqrt {{x}^{2}+x+1}=-
1/2.$$
\fincorrection
\correction{001217}
\begin{enumerate}
  \item $\displaystyle{\lim _{x\rightarrow 0}{\frac
{\sin(x)\ln (1+{x}^{2})}{x\tan(x)}}=0 }$.
  \item $\displaystyle{\lim _{x\rightarrow 0}{\frac {\ln (1+\sin(x))}{\tan(6\,x)}}=1/6 }$.
  \item $\displaystyle{\lim _{x\rightarrow 0}\left (\ln ({e}+x)\right )^{{x}^{-1}}={e^{{e^{-1}}}} }$.
  \item $\displaystyle{\lim _{x\rightarrow \infty }\left (\ln (1+{e^{-x}})\right )^{{x}^{-1}}={e^{-1}} }$.
\end{enumerate}
\fincorrection
\nocorrection
\correction{001219}
\begin{enumerate}
  \item $\frac 23$
  \item ${\frac {\sqrt {2}}{8\,{x}^{3}}}$
  \item $\frac{a^3}{b^3}$
  \item $-1$
  \item $-{\frac {\sqrt {2}}{4}}$
  \item $\frac {1}{2}{x}^{2}$
  \item $-{\frac {3}{2}}\left (-{\frac {\pi }{4}}+x\right)$
  \item $\sqrt{e}$
  \item $\frac 1{\pi}$
  \item $1$
  \item $x$
\end{enumerate}
\fincorrection
\nocorrection
\correction{005161}
Soit $n\in\Nn\setminus\{0,1,2\}$.

$$n!^2=\prod_{k=1}^{n}(n+1-k)\prod_{k=1}^{n}k=\prod_{k=1}^{n}k(n+1-k).$$

Maintenant, la fonction $x\mapsto x(n+1-x)$ est strictement croissante sur $[0,\frac{n+1}{2}]$ et
strictement décroissante sur $[\frac{n+1}{2},n+1]$. Puisque $f(1)=f(n)=n$, on en déduit que pour $x\in[2,n-1]$,
$f(x)>n$. Puisque $n\geq3$, on a $n-1\geq2$ et on peut écrire

$$n!^2=n^2\prod_{k=2}^{n-1}k(n+1-k)>n^2\prod_{k=2}^{n-2}n=n^n,$$

et donc,

$$\sqrt[n]{n!}=(n!^2)^{1/(2n)}>(n^n)^{1/2n}=\sqrt{n}.$$
\fincorrection
\correction{005163}
Soit $x\in\Rr$. Montrons par récurrence que~:~$\forall n\in\Nn,\;|\sin(nx)|\leq
n|\sin x|$.
\begin{itemize}
\item[-] C'est clair pour $n=0$.
\item[-] Soit $n\geq0$. Supposons que $|\sin(nx)|\leq n|\sin x|$. Alors,

\begin{align*}
|\sin(n+1)x|&=|\sin nx\cos x+\cos nx\sin x|\leq|\sin nx|.|\cos x|+|\cos nx|.|\sin x|\\
 &\leq|\sin nx|+|\sin x|\\
 &\leq n|\sin x|+|\sin x|\quad(\mbox{par hypothèse de récurrence})\\
 &=(n+1)|\sin x|
\end{align*}

\end{itemize}

On a montré par récurrence que \shadowbox{$\forall n\in\Nn,\forall x\in\Rr,\;|\sin(nx)|\leq n|\sin x|$}.

\fincorrection
\correction{005393}
Pour $x\in[a,b]$, posons $g(x)=f(x)-x$. $g$ est continue sur $[a,b]$ puisque $f$ l'est. De plus, $g(a)=f(a)-a\geq0$ et $g(b)=f(b)-b\leq0$. D'après le théorème des valeurs intermédiaires, $g$ s'annule au moins une fois sur $[a,b]$ ou encore, l'équation $f(x)=x$ admet au moins une solution dans $[a,b]$.
\fincorrection
\correction{005394}
Puisque $\frac{f(x)}{x}$ tend vers $\ell\in[0,1[$, il existe $A>0$ tel que pour $x\geq A$, $\frac{f(x)}{x}\leq\frac{\ell+1}{2}<1$. Mais alors, $f(A)<A$ (et $f(0)\geq0$) ce qui ramène à la situation de l'exercice \ref{exo:rouconti}~:~pour $x\in[0,A]$, soit $g(x)=f(x)-x$...
\fincorrection
\correction{005397}
Soit $x>0$. Pour tout naturel $n$, $f(x)=f(\sqrt{x})=f(x^{1/4})=...=f(x^{1/2^n})$. Or, à $x$ fixé, $\lim_{n\rightarrow +\infty}x^{1/2^n}=1$ et, $f$ étant continue en $1$, on a~:

$$\forall x>0,\;f(x)=\lim_{n\rightarrow +\infty}f(x^{1/2^n})=f(1).$$

$f$ est donc constante sur $]0,+\infty[$, puis sur $[0,+\infty[$ par continuité de $f$ en $0$.

Pour $x\geq0$, posons $f(x)=0$ si $x\neq1$ et $f(x)=1$ si $x=1$. Pour $x\geq0$, on a $x^2=1\Leftrightarrow x=1$. $f$ vérifie donc~:~$\forall x\geq0,\;f(x^2)=f(x)$, mais $f$ n'est pas constante sur $\Rr^+$.
\fincorrection
\correction{005399}
Soit $f$ un morphisme de $(\Rr,+)$, c'est-à-dire que $f$ est une application de $\Rr$ dans $\Rr$ vérifiant 
$$\forall(x,y)\in\Rr^2,\;f(x+y)=f(x)+f(y).$$

On sait déjà $f(0)=f(0+0)=f(0)+f(0)$ et donc $f(0)=0$. Puis, pour $x$ réel donné, $f(-x)+f(x)=f(-x+x)=f(0)=0$ et donc, pour tout réel $x$, $f(-x)=-f(x)$ ($f$ est donc impaire). On a aussi $n\in\Nn^*$ et $x\in\Rr$, $f(nx)=f(x)+...+f(x)=nf(x)$. De ce qui précède, on déduit~:

$$\forall x\in\Rr,\;\forall n\in\Zz,\;f(nx)=nf(x).$$

Soit $a=f(1)$. D'après ce qui précède, $\forall n\in\Zz,\;f(n)=f(n.1)=nf(1)=an$.

Puis, pour $n\in\Nn^*$, $nf(\frac{1}{n})=f(n\frac{1}{n})=f(1)=a$ et donc $\forall n\in\Nn^*,\;f(\frac{1}{n})=a\frac{1}{n}$.

Puis, pour $p\in\Zz$ et $q\in\Nn^*$, $f(\frac{p}{q})=pf(\frac{1}{q})=pa\frac{1}{q}=a\frac{p}{q}$.

Finalement, 

$$\forall r\in\Qq,\;f(r)=ar.$$

Maintenant, si l'on n'a pas l'hypothèse de continuité, on ne peut aller plus loin. Supposons de plus que $f$ soit continue sur $\Rr$.

Soit $x$ un réel. Puisque $\Qq$ est dense dans $\Rr$, il existe une suite $(r_n)_{n\in\Nn}$ de rationnels, convergente de limite $x$.

$f$ étant continue en $x$, on a~:~

$$f(x)=f(\lim_{n\rightarrow +\infty}r_n)=\lim_{n\rightarrow +\infty}f(r_n)=\lim_{n\rightarrow +\infty}ar_n=ax.$$

Donc, si $f$ est un morphisme continu de $(\Rr,+)$, $f$ est une application linéaire de $\Rr$ dans $\Rr$. Réciproquement, les applications linéaires conviennent.
\fincorrection
\correction{005400}
Soient $a$ et $b$ deux réels fixés tels que $0<a<b$. Trouvons les entiers naturels non nuls $k$ tels que 
$]ka,kb[\cap](k+1)a,(k+1)b[\neq\emptyset$. Pour $k$ dans $\Nn^*$, posons $I_k=]ka,kb[$.

$$I_k\cap I_{k+1}\neq\emptyset\Leftrightarrow ka<(k+1)a<kb<(k+1)b\Leftrightarrow k>\frac{a}{b-a}\Leftrightarrow k\geq E(\frac{a}{b-a})+1.$$

Posons $k_0=E(\frac{a}{b-a})+1$. Pour $k\geq k_0$, on a donc $I_k\cap I_{k+1}\neq\emptyset$ et donc $\displaystyle\bigcup_{k\geq k_0}]ka,kb[=]k_0a,+\infty[$.

Maintenant, si $k_0=1$, $\displaystyle\bigcup_{k\geq k_0}]ka,kb[=]a,+\infty[$ et si $k_0>1$, $\displaystyle\bigcup_{k\geq k_0}]ka,kb[=(\displaystyle\bigcup_{k=1}^{k_0-1}]ka,kb[)\cup]k_0a,+\infty[$. Mais, si $x$ est dans $\displaystyle\bigcup_{k=1}^{k_0-1}]ka,kb[$, alors $x<(k_0-1)b<k_0a$ et donc, $(\displaystyle\bigcup_{k=1}^{k_0-1}]ka,kb[)\cap]k_0a,+\infty[=\emptyset$. La plus petite valeur de $A$ est donc $(E(\frac{a}{b-a})+1)a$.
\fincorrection
\correction{005402}
Si $f$ est strictement monotone sur $I$, on sait que $f$ est injective.

Réciproquement, supposons $f$ injective et continue sur $I$ et montrons que $f$ est strictement monotone.

Supposons par l'absurde que $f$ n'est pas strictement monotone. On peut alors trouver trois réels $a$, $b$ et $c$ dans l'intervalle $I$ tels que
 
$$a<b<c\;\mbox{et}\;((f(b)\geq f(a)\;\mbox{et}\;f(b)\geq f(c))\;\mbox{ou}\;(f(b)\leq f(a)\;\mbox{et}\;f(b)\leq f(c))).$$

Quitte à remplacer $f$ par $-f$, on supposera que $a<b<c$ et $f(b)\geq f(a)$ et $f(b)\geq f(c)$.

Puisque $f$ est injective, on a même $a<b<c$ et $f(b)>f(a)$ et $f(b)>f(c)$. Soit $M=\mbox{Max}\{f(a),f(c)\}$. On a $M<f(b)$. $M$ est élément de $[f(a),f(b)]$ et, puisque $f$ est continue sur $[a,b]$, le théorème des valeurs intermédiaires permet d'affirmer qu'il existe $\alpha\in[a,b]$ tel que $f(\alpha)=M$. De plus, on ne peut avoir $\alpha=b$ car $f(\alpha)=M\neq f(b)$ (et $f$ injective). Donc, 

$$\exists\alpha\in[a,b[/\;f(\alpha)=M.$$

De même, puisque $M$ est élément de $[f(c),f(b)]$, $\exists\beta\in]b,c]/\;f(\beta)=M$. Ainsi, on a trouvé dans $I$ deux réels $\alpha$ et $\beta$ vérifiant $\alpha\neq\beta$ et $f(\alpha)=f(\beta)$ ce qui contredit l'injectivité de $f$.

Donc, $f$ est strictement monotone sur $I$.
\fincorrection
\correction{005403}
Soit $f$ la fonction caractéristique de $Q$. Le groupe des périodes de $f$ est $\Qq$. En effet,

$$\forall x\in\Rr,\;\forall r\in\Qq,\;x+r\in\Qq\Leftrightarrow x\in\Qq,$$ et donc 

$\forall x\in\Rr,\;\forall r\in\Qq,\;f(x+r)=f(x)$. Mais on a aussi $$\forall x\in\Rr,\;\forall r\in(\Rr\setminus\Qq),\;x+r\in\Qq,\Leftrightarrow x\notin\Qq,$$ et donc $\forall x\in\Rr,\;\forall r\in(\Rr\setminus\Qq),\;f(x+r)\neq f(x)$.
\fincorrection
\correction{005405}
$Id$ est solution.

Réciproquement, soit $f$ une bijection de $[0,1]$ sur lui-même vérifiant $\forall x\in[0,1],\;f(2x-f(x))=x$. Nécessairement, $\forall x\in[0,1],\;0\leq 2x-f(x)\leq 1$ et donc $\forall x\in[0,1],\;2x-1\leq f(x)\leq 2x$.

Soit $f^{-1}$ la réciproque de $f$.

\begin{align*}\ensuremath
\forall x\in[0,1],\;f(2x-f(x))=x&\Leftrightarrow\forall x\in[0,1],\;2x-f(x)=f^{-1}(x)\\
 &\Leftrightarrow\forall y\in[0,1],\;f(f(y))-2f(y)+y=0\;(\mbox{car}\;\forall x\in[0,1],\;\exists!y\;[0,1]/\;x=f(y))
\end{align*}

Soit $y\in[0,1]$ et $u_0=y$. En posant $\forall n\in\Nn,\;u_{n+1}=f(u_n)$, on définit une suite de réels de $[0,1]$ (car $[0,1]$ est stable par $f$). La condition $\forall y\in[0,1],\;f(f(y))-2f(y)+ y=0$ fournit $\forall n\in\Nn, \;u_{n+2}-2u_{n+1}+u_n=0$, ou encore $\forall n\in\Nn,\;u_{n+2}-u_{n+1}=u_{n+1}-u_n$. La suite $(u_{n+1}-u_n)_{n\in\Nn}$ est constante ou encore $u$ est arithmétique. Mais, $u$ est également bornée et donc $u$ est constante.

En particulier, $u_1=u_0$ ce qui fournit $f(y)=y$. On a montré que $\forall y\in[0,1],\;f(y)=y$ et donc $f=Id$.
\fincorrection
\correction{005406}
\begin{enumerate}
\item  Soit $n$ un entier naturel non nul donné. Pour $x$ élément de $[0,1-\frac{1}{n}]$, posons $g(x)=f(x+\frac{1}{n})-f(x)$.

$g$ est définie et continue sur $[0,1-\frac{1}{n}]$. De plus,

$$\sum_{k=0}^{n-1}g(\frac{k}{n})=\sum_{k=0}^{n-1}(f(\frac{k+1}{n})-f(\frac{k}{n}))=f(1)-f(0)=0.$$ 

Maintenant, s'il existe un entier $k$ élément de $\{0,...,n-1\}$ tel que $g(\frac{k}{n})=0$, on a trouvé un réel $x$ de $[0,1]$ tel que $f(x+\frac{1}{n})=f(x)$ (à savoir $x=\frac{k}{n}$).

Sinon, tous les $g(\frac{k}{n})$ sont non nuls et, étant de somme nulle, il existe deux valeurs de la variable en lesquels $g$ prend des valeurs de signes contraires. Puisque $g$ est continue sur $[0,1-\frac{1}{n}]$, le théorème des valeurs intermédiares permet d'affirmer que $g$ s'annule au moins une fois dans cet intervalle ce qui fournit de nouveau une solution à l'équation $f(x+\frac{1}{n})=f(x)$.

\item  Soit $a\in]0,1[$ tel que $\frac{1}{a}\notin\Nn^*$. Soit, pour $x\in[0,1]$, 
$f(x)=|\sin\frac{\pi x}{a}|-x|\sin\frac{\pi}{a}|$. $f$ est continue sur $[0,1]$, $f(0)=f(1)=0$ mais,
 
$$\forall x\in\Rr,\;f(x+a)-f(x)=(|\sin\frac{\pi(x+a)}{a}|-|\sin\frac{\pi x}{a}|)-((x+a)-x)|\sin\frac{\pi}{a}|=-a|\sin\frac{\pi}{a}|\neq0.$$

\item (a) et b)) Soit $g(t)$ la distance, exprimée en kilomètres, parcourue par le cycliste à l'instant $t$ exprimé en heures, $0\leq t\leq1$, puis, pour $t\in[0,1]$, $f(t)=g(t)-20t$. $f$ est continue sur $[0,1]$ (si le cycliste reste un tant soit peu cohérent) et vérife $f(0)=f(1)=0$.

D'après 1), $\exists t_1\in[0,\frac{1}{2}]$, $\exists t_2\in[0,\frac{19}{20}]$ tels que $f(t_1+\frac{1}{2})=f(t_1)$ et $f(t_2+\frac{1}{20})=f(t_2)$ ce qui s'écrit encore
$g(t_1+\frac{1}{2})-g(t_1)=10$ et $g(t_2+\frac{1}{20})-g(t_2)=1$.

c) Posons pour $0\leq t\leq 1$, $f(t)=|\sin\frac{4\pi t}{3}|-\frac{t\sqrt{3}}{2}$ et donc, $g(t)=|\sin\frac{4\pi t}{3}|+(20-\frac{\sqrt{3}}{2})t$. $\forall t\in[0,\frac{1}{4}]$, $f(t+\frac{3}{4})-f(t)\neq 0$ ou encore $g(t+\frac{3}{4})-g(t)\neq 15$.
\end{enumerate}

\fincorrection
\correction{000698}
\begin{enumerate}
  \item La fonction $f_1$ est d\'erivable en dehors de $x=0$.
En effet $x \mapsto \frac 1x$ est dérivable sur $\Rr^*$ et $x \mapsto \cos x$ est dérivable sur $\Rr$,
donc par composition $x \mapsto  \cos \frac 1x$ est dérivable sur $\Rr^*$. Puis par multiplication par la fonction dérivable 
$x \mapsto x^2$, la fonction $f_1$ est dérivable sur $\Rr^*$. 
Par la suite on omet souvent ce genre de discussion ou on l'abrège sous la forme ``$f$ est dérivable sur $I$ comme somme, produit,
composition de fonctions dérivables sur $I$''.

Pour savoir si $f_1$ est d\'erivable en $0$ regardons le taux d'accroissement:
$$ \frac{f_1(x)-f_1(0)}{x-0}= x\cos \frac 1 x.$$
Mais $x \cos (1/x)$ tend vers $0$ (si $x\rightarrow 0$) car
$|\cos (1/x)| \leq 1$.
Donc le taux d'accroissement tend vers $0$. Donc $f_1$ est d\'erivable en $0$ et $f_1'(0)=0$.
  \item Encore une fois $f_2$ est d\'erivable en dehors de $0$.
Le taux d'accroissement en $x=0$ est :
$$ \frac{f_2(x)-f_2(0)}{x-0}= \frac{\sin x}{x} \sin \frac 1 x$$
Nous savons que $\frac{\sin x}{x} \rightarrow 1$ et que
$\sin 1/x$ n'a pas de limite quand $x\rightarrow 0$. Donc le taux d'accroissement n'a pas de limite, donc $f_2$ n'est pas d\'erivable en $0$.
  \item 
La fonction $f_3$ s'\'ecrit :
$$f_3(x) = \frac{|x||x-1|}{x-1}.$$
\begin{itemize}
  \item
  Donc pour $x \geq 1$ on a $f_3(x) = x$ ;
pour $0 \leq x < 1$ on a $f_3(x) = -x$ ;
pour $x <0$ on a $f_3(x) = x$.

\item La fonction $f_3$ est d\'efinie, continue et d\'erivable sur 
$\Rr \setminus \{0,1\}$. Attention ! La fonction $x \mapsto |x|$ n'est pas dérivable en $0$.

\item 
La fonction $f_3$ n'est pas continue en $1$, en effet
$\lim_{x \rightarrow 1+} f_3(x) = +1$ et  $\lim_{x \rightarrow 1-} f_3(x) = -1$. Donc la fonction n'est pas d\'erivable en $1$.

\item
La fonction $f_3$ est continue en $0$. 
Le taux d'accroissement pour $x>0$ est 
$$\frac{f_3(x)-f_3(0)}{x-0}= \frac{-x}{x} = -1$$
et pour $x <0$, 
$$\frac{f_3(x)-f_3(0)}{x-0}= \frac{x}{x} = +1.$$
Donc le taux d'accroissement n'a pas de limite en $0$ et donc $f_3$ n'est pas d\'erivable en $0$.
\end{itemize}
\end{enumerate}
\fincorrection
\correction{000699}
La fonction $f$ est continue et dérivable sur $]0,1[$ et sur
$]1,+\infty[$. Le seul problème est en $x=1$.

Il faut d'abord que la fonction soit continue en $x=1$.
La limite \`a gauche est
$\lim_{x\rightarrow 1^-} \sqrt x = +1$
et \`a droite 
$\lim_{x\rightarrow 1^+} ax^2+bx+1 = a+b+1$.
Donc  $a+b+1=1$. Autrement dit $b = -a$.

Il faut maintenant que les d\'eriv\'ees \`a droite et \`a gauche soient
\'egales.
Comme la fonction $f$ restreinte à $]0,1]$ est définie par $x \mapsto \sqrt{x}$ alors elle est dérivable
à gauche et la dérivée à gauche s'obtient en évaluant la fonction dérivée $x \mapsto \frac{1}{2\sqrt x}$
en $x=1$. Donc $f'_g(1)=\frac 12$.

Pour la dérivée à droite il s'agit de calculer la limite du taux d'accroissement
$\frac{f(x)-f(1)}{x-1}$, lorsque $x \to 1$ avec $x>1$.
Or 
$$
\frac{f(x)-f(1)}{x-1} = \frac{ax^2+bx+1 - 1}{x-1} = \frac{ax^2-ax}{x-1} = \frac{ax(x-1)}{x-1} = ax.
$$
Donc $f$ est dérivable à droite et $f'_d(1) = a$.
Afin que $f$ soit dérivable, il faut et il suffit que les dérivées à droite et à gauche
existent et soient égales, donc ici la condition est $a=\frac 12$.


Le seul couple $(a,b)$ que rend $f$ dérivable sur $]0,+\infty[$ 
est $(a=\frac 12, b= -\frac 12)$.
\fincorrection
\correction{000700}
$f$ est $C^\infty$ sur $\Rr^*$.
\begin{enumerate}
  \item Comme $|\sin (1/x)| \leq 1$ alors
$f$ tend vers $0$ quand $x\rightarrow 0$. Donc en prolongeant $f$ par $f(0) =0$,
la fonction  $f$ prolongée est continue sur $\Rr$.
  \item Le taux d'accroissement est
$$\frac{f(x)-f(0)}{x-0}= x \sin\frac{1}{x}.$$
Comme ci-dessus il y a une limite (qui vaut $0$) en $x=0$.
Donc $f$ est d\'erivable en $0$ et $f'(0)=0$.
  \item Sur $\Rr^*$, $f'(x) = 2x\sin (1/x) -\cos(1/x)$,
Donc $f'(x)$ n'a pas de limite quand $x\rightarrow 0$.
Donc $f'$ n'est pas continue en $0$.
\end{enumerate}
\fincorrection
\correction{000701}
\begin{enumerate}
  \item Selon que $n \equiv 0 \pmod 4, 1 \pmod 4, 2 \pmod 4 , 3\pmod 4$ alors
$f^{(n)}(x)$ vaut respectivement $\sin x$, $\cos x$, $-\sin x$,
$-\cos x$.
  \item La d\'eriv\'ee de $\sin^2 x$ est $2\sin x \cos x = \sin 2x$.
Et donc les d\'eriv\'ees suivantes seront :
$2\cos 2x, -4\sin 2x, -8\cos 2x, 16\sin 2x$,...
Et selon que $n \equiv 1 \pmod 4, 2 \pmod 4 , 3\pmod 4, 0 \pmod 4,$ alors
$g^{(n)}(x)$ vaut respectivement $2^{n-1}\sin 2x$, $2^{n-1}\cos 2x$, $-2^{n-1}\sin 2x$, $-2^{n-1}\cos 2x$.
  \item $\sin(x)^3+\cos(x)^3 = -\frac14\sin(3x)+\frac34\sin(x)+\frac14\cos(3x)+\frac34\cos(x)$ et on d\'erive...
\end{enumerate}
\fincorrection
\nocorrection
\nocorrection
\nocorrection
\nocorrection
\nocorrection
\nocorrection
\nocorrection
\correction{000709}
La limite de $f$ en $0$ est $0$, donc $f$ est continue en $0$.
De même le taux d'accroissement de $f$ en $0$ est
$f(x)/x$ qui tend vers $0$. Donc $f$ est dérivable en $0$ et 
$f'(0)=0$. En dehors de $0$, on a $f'(x) = 2{e^{-{x}^{-2}}}{x}^{-3}$
donc $f'$ est continue en $0$.

On continue de la même fa\c{c}on en remarquant que si
$f^{(n)}(x) = P(1/x) \exp(-1/x^2)$  
où $P$ est un polynôme et $f^{(n)}(0) = 0$.
Donc $f^{(n)}(x)$ tend vers $0$ si $x$ tend vers $0$. Donc $f^{(n)}$
est continue. De plus $f^{(n)}$ est dérivable en $0$ car son taux
d'accroissement vaut $1/x P(1/x) \exp(-1/x^2)$ qui tend vers $0$, donc 
$f^{(n+1)}(0) = 0$.
En dehors de $0$ on $f^{(n+1)}(x) = Q(1/x) \exp(-1/x^2)$ où $Q$ est un polynôme. 
Et on recommence...
\fincorrection
\nocorrection
\nocorrection
\nocorrection
\nocorrection
\nocorrection
\nocorrection
\nocorrection
\nocorrection
\nocorrection
\nocorrection
\nocorrection
\nocorrection
\nocorrection
\correction{005413}
\begin{enumerate}
\item  Pour $n\geq1$, on a d'après la formule de \textsc{Leibniz}~:

\begin{align*}\ensuremath
(x^{n-1}\ln(1+x))^{(n)}&=\sum_{k=0}^{n}\dbinom{n}{k}(x^{n-1})^{(k)}(\ln(1+x))^{(n-k)}\\
 &=\sum_{k=0}^{n-1}\dbinom{n}{k}(x^{n-1})^{(k)}(\ln(1+x))^{(n-k)}\;(\mbox{car}\;(x^{n-1})^{(n)})=0)\\
 &=\sum_{k=0}^{n-1}\dbinom{n}{k}\frac{(n-1)!}{(n-1-k)!}x^{n-1-k}(-1)^{n-1-k}\frac{(n-1-k)!}{(x+1)^{n-k}}\\
 &(\mbox{car}\;(\ln(1+x))^{(n-k)}=(\frac{1}{1+x})^{(n-k-1)}).
\end{align*}

Puis, pour $x=0$, $(x^{n-1}\ln(1+x))^{(n)}(0)=n.(n-1)!=n!$, et pour $x\neq0$, 

\begin{align*}\ensuremath
(x^{n-1}\ln(1+x))^{(n)}(x)&=-\frac{(n-1)!}{x}\sum_{k=0}^{n-1}\dbinom{n}{k}(-\frac{x}{x+1})^{n-k}
=-\frac{(n-1)!}{x}((1-\frac{x}{x+1})^n-1)\\
 &=\frac{(n-1)!}{x}\frac{(x+1)^n-1}{(x+1)^n}.
\end{align*}

\item  On sait dériver facilement des sommes ou plus généralement des combinaisons linéaires. Donc, on linéarise~:

\begin{align*}\ensuremath
\cos^3x\sin(2x)&=\frac{1}{8}(e^{ix}+e^{-ix})^3(-\frac{1}{4})(e^{2ix}-e^{-2ix})=-\frac{1}{32}
(e^{3ix}+3e^{ix}+3e^{-ix}+e^{-3ix})(e^{2ix}-2+e^{-2ix})\\
 &=-\frac{1}{32}(e^{5ix}+e^{3ix}-2e^{ix}-2e^{-ix}+e^{-3ix}+e^{-5ix})=-\frac{1}{16}(\cos(5x)+\cos(3x)-2\cos(x))
\end{align*}

Puis, pour $n$ naturel donné~:

$$(\cos^3x\sin2x)^{(n)}=-\frac{1}{16}(5^n\cos(5x+n\frac{\pi}{2})+3^n\cos(3x+n\frac{\pi}{2})-2cos(x+n\frac{\pi}{2})),$$

expression que l'on peut détailler suivant la congruence de $n$ modulo $4$.

\item  On sait dériver des objets simples et donc on décompose en éléments simples~:

$$\frac{X^2+1}{(X-1)^3}=\frac{X^2-2X+1+2X-2+2}{(X-1)^3}=\frac{1}{X-1}+\frac{2}{(X-1)^2}+\frac{2}{(X-1)^3}.$$

Puis, pour $n$ entier naturel donné,

\begin{align*}\ensuremath
\left(\frac{X^2+1}{(X-1)^3}\right)^{(n)}&=\frac{(-1)^nn!}{(X-1)^{n+1}}+2\frac{(-1)^n(n+1)!}{(X-1)^{n+2}}+\frac{(-1)^n(n+2)!}{(X-1)^{n+3}}\\
 &=\frac{(-1)^nn!}{(X-1)^{n+3}}((X-1)^2+2(n+1)(X-1)+(n+2)(n+1))\\
 &=\frac{(-1)^nn!(X^2+2nX+n^2+n+1)}{(X-1)^{n+3}}.
\end{align*}

\item  La fonction proposée est de classe $C^\infty$ sur $\Rr$ en vertu de théorèmes généraux. La formule de \textsc{Leibniz} fournit pour $n\geq3$~:

\begin{align*}\ensuremath
((x^3+2x-7)e^x)^{(n)}&=\sum_{k=0}^{n}\dbinom{n}{k}(x^3+2x-7)^{(k)}(e^x)^{(n-k)}=\sum_{k=0}^{3}\dbinom{n}{k}(x^3+2x-7)^{(k)}(e^x)^{(n-k)}\\
 &=((x^3+2x-7)+n(3x^2+2)+\frac{n(n-1)}{2}(6x)+\frac{n(n-1)(n-2)}{6}.6)e^x\\
 &=(x^3+3nx^2+(3n^2-3n+2)x+n^3-3n^2+4n-7)e^x.
\end{align*}
\end{enumerate}
\fincorrection
\correction{000715}
$Q_n(t) = (1-t^2)^n$ est un polyn\^ome de degr\'e $2n$, on le d\'erive $n$ fois, on obtient un polyn\^ome de degr\'e $n$.
Les valeurs $-1$ et $+1$ sont des racines d'ordre $n$ de $Q_n$, donc
$Q_n(1)=Q_n'(1)=\ldots = Q_n^{(n-1)}(1)=0$. M\^eme chose en $-1$.
Enfin $Q(-1)=0=Q(+1)$ donc d'apr\`es le th\'eor\`eme de Rolle il existe $c \in ]-1,1[$ telle que $Q_n'(c)=0$. 

Donc $Q_n'(-1)=0$, $Q_n'(c)=0$, $Q_n'(-1)=0$. En appliquant le th\'eor\`eme de Rolle deux fois 
(sur $[-1,c]$ et sur $[c,+1]$), on obtient l'existence
de racines $d_1,d_2$ pour $Q_n''$, qui s'ajoutent aux racines $-1$ et $+1$.


On continue ainsi par r\'ecurrence. On obtient pour $Q_n^{(n-1)}$,
$n+1$ racines: $-1, e_1,\ldots, e_{n-1}, +1$. Nous appliquons le th\'eor\`eme de Rolle $n$ fois. Nous obtenons $n$ racines pour 
$P_n = Q_n^{(n)}$.  Comme un polyn\^ome de degr\'e $n$ a au plus
$n$ racines, nous avons obtenu toutes les racines. Par constructions ces racines sont r\'eelles distinctes, donc simples.
\fincorrection
\nocorrection
\correction{000717}
\begin{enumerate}
\item Par l'absurde on suppose qu'il y a (au moins) quatre racines distinctes  pour 
$  P_n(X) =  X^n+aX+b  $. Notons les $x_1 <x_2<x_3<x_4$.
Par le th\'eor\`eme de Rolle appliqu\'e trois fois (entre $x_1$ et $x_2$, entre $x_2$ et $x_3$,...) il existe $x'_1<x'_2<x'_3$ des racines de $P'_n$. On applique deux fois le th\'eor\`eme Rolle entre $x'_1$ et $x'_2$ et entre $x'_2$ et $x'_3$. On obtient deux racines distinctes pour $P''_n$. Or 
$P''_n = n(n-1)X^{n-2}$ ne peut avoir que $0$ comme racines. Donc nous avons 
obtenu une contradiction. 
\item \emph{Autre m\'ethode :} 
Le r\'esultat est \'evident si $  n\leq 3  .$ On suppose donc $
n\geq 3  .$ Soit $  P_n  $ l'application $  X\mapsto X^n+aX+b  $
de $  \R   $ dans lui-m\^eme. Alors $  P'_n(X)=nX^{n-1}+a  $
s'annule en au plus deux valeurs. Donc $P_n$ est successivement
croissante-d\'ecroissante-croissante ou bien
d\'ecroissante-croissante-d\'ecroissante. Et donc $P_n$ s'annule au plus 
trois fois.
\end{enumerate}
\fincorrection
\correction{000718}
 Comme $  f'  $ est d\'erivable, elle est continue. Comme
$  f  $ s'annule $  n+1  $ fois, $  f'  $ change de signe (au
moins) $  n+1  $ fois donc s'annule (au moins) $  n   $ fois. On
peut bien s\^ur recommencer, le r\'esultat en d\'ecoule.
\fincorrection
\nocorrection
\nocorrection
\correction{000721}
La fonction $f$ est continue et dérivable sur $\Rr$ donc en particulier sur $[a,b]$.
Le théorème des accroissement finis assure l'existence d'un nombre $c \in ]a,b[$ tel que
$f(b)-f(a) = f'(c) (b-a)$.

Mais pour la fonction particulière de cet exercice nous pouvons expliciter ce $c$.
En effet $f(b)-f(a) = f'(c) (b-a)$ implique $\alpha(b^2-a^2)+\beta(b-a)= (2\alpha c+\beta)(b-a)$.
Donc $c = \frac{a+b}{2}$. 

G\'eom\'etriquement, le graphe $\mathcal{P}$ de $f$ est une parabole.
Si l'on prend deux points $A = (a,f(a))$ et $B = (b,f(b))$ appartenant 
à cette parabole, alors la droite $(AB)$
est parall\`ele \`a la tangente en $\mathcal{P}$ qui passe en $M=(\frac{a+b}{2}, f(\frac{a+b}{2}))$.
L'abscisse de $M$ étant le milieu des abscisses de $A$ et $B$.
\fincorrection
\nocorrection
\nocorrection
\correction{000724}
\begin{enumerate}
  \item Soit $g(t) = \ln t$. Appliquons le th\'eor\`eme des accroissements finis sur $[x,y]$. Il existe $c \in ]x,y[$, 
$g(y)-g(x) = g'(c)(y-x)$. Soit $\ln y - \ln x = \frac 1c (y-x)$.
Donc $\frac{\ln y - \ln x}{y-x} = \frac 1c$.
Or $x <c<y$ donc $\frac 1y < \frac 1c < \frac 1x$.
Ce qui donne les in\'egalit\'es recherch\'ees. 
  \item $f'(\alpha)= \frac{x-y}{\alpha x + (1-\alpha)y} - \ln x + \ln y$. Et $f''(\alpha) = -\frac{(x-y)^2}{(\alpha x + (1-\alpha)y)^2}$.
Comme $f''$ est n\'egative alors $f'$ est d\'ecroissante sur $[0,1]$.
Or $f'(0) = \frac{x-y - y(\ln x - \ln y)}{y} >0$ d'apr\`es la premi\`ere question et de m\^eme $f'(1) < 0$. Par le th\'eor\`eme des valeurs interm\'ediaires, il existe $c \in [x,y]$ tel que $f'(c) = 0$.
Maintenant $f'$ est positive sur $[0,c]$ et n\'egative sur $[c,1]$.
Donc $f$ est croissante sur $[0,c]$ et d\'ecroissante sur $[c,1]$.
Or $f(0)=0$ et $f(1)=0$ donc pour tout $x\in[0,1]$, $f(x) \geq 0$.
Cela prouve l'in\'egalit\'e demand\'ee.
  \item G\'eom\'etriquement nous avons prouv\'e que la fonction $\ln$
est concave, c'est-\`a-dire que la corde
(le segment qui va de $(x,f(x))$ \`a $(y,f(y)$) est sous la courbe d'\'equation $y=f(x)$.
\end{enumerate}
\fincorrection
\correction{000725}
Le th\'eor\`eme des accroissements finis donne : 
$\ln(n+1)-\ln(n) = \frac1{c_n} (n+1-n) = \frac1{c_n}$, avec $c_n \in [n,n+1]$. Or $c_n \geq n$ donc $\frac 1n \geq \frac 1{c_n}$.
Donc :
$$S_n=\sum_{k=1}^n \frac 1k \geq \sum_{k=1}^n \frac 1{c_k}
= \sum_{k=1}^n \ln(k+1)-\ln(k) = \ln(n+1).$$
La derni\`ere \'egalit\'e s'obtient car la somme est t\'el\'escopique et $\ln 1 = 0$.
Donc $S_n \geq \ln(n+1)$, donc $S_n \rightarrow +\infty$.
\fincorrection
\nocorrection
\correction{000727}
Pour simplifier nous supposons $x>0$.
\begin{enumerate}
  \item Appliquer le th\'eor\`eme des accroissements finis  ne va pas \^etre suffisant. En effet, soit $f(x) = e^x-1-x$. Alors
il existe $c \in ]0,x[$ tel que $f(x)-f(0)=f'(c)(x-0)$.
Soit $f(x)=(e^c-1)x$.
Soit maintenant $g(x) = e^x-1$ alors, par le th\'eor\`eme des accroissements finis sur $[0,c]$ il existe $d \in ]0,c[$
tel que $g(c)-g(0)=g'(d)(c-0)$, soit $e^c-1=e^dc$.
Donc $e^x-1-x = f(x) = (e^c-1)x = e^dcx$. 
Comme $d \leq c \leq x$, alors
$e^x-1-x \leq  e^xx^2$. 

Cela donne une in\'egalit\'e, mais il manque un facteur $1/2$.

  \item Nous allons obtenir l'in\'egalit\'e par application du th\'eor\`eme de Rolle. Soit maintenant $f(t) = e^t-1-t-k\frac{t^2}{2}$. Nous avons $f(0)=0$,
$x>0$ \'etant fix\'e, nous choisisons $k$ tel que $f(x)=0$,
(un tel $k$ existe car $e^x-1-x >0$ et $x^2>0$).
Comme $f(0)=0=f(x)$ alors par Rolle il existe $c \in ]0,x[$
tel que $f'(c)=0$.
Mais $f'(t) = e^t-t-kt$, donc $f'(0)=0$.
Maintenant $f'(0)=0=f'(c)$ donc il existe (par Rolle toujours !)
$d \in ]0,c[$ tel que $f''(d)=0$.
Or $f''(t) = e^t-k$, donc $f''(d)=0$ donne $k = e^d$.
Ainsi $f(x)=0$ devient $e^x-1-x = e^d \frac{x^2}{2}$.
Comme $d\leq x$ alors $e^x-1-x \leq e^x \frac{x^2}{2}$.
\end{enumerate}
\fincorrection
\nocorrection
\correction{005408}
On a déjà $g(b)=f(b)-f(b)=0$. Puisque $a\neq b$, on peut choisir $A$ tel que $g(a)=0$ (à savoir $A=\frac{(n+1)!}{(b-a)^{n+1}}(f(b)-\sum_{k=0}^{n}\frac{f^{(k)}(a)}{k!}(b-a)^k$).

Avec les hypothèses faites sur $f$, $g$ est d'autre part continue sur $[a,b]$ et dérivable sur $]a,b[$. Le théorème de \textsc{Rolle} permet alors d'affirmer qu'il existe $c\in]a,b[$ tel que $g'(c)=0$.

Pour $x\in]a,b[$, on a

\begin{align*}\ensuremath
g'(x)&=-\sum_{k=0}^{n}\frac{f^{(k+1)}(x)}{k!}(b-x)^k+\sum_{k=1}^{n}\frac{f^{(k)}(x)}{(k-1)!}(b-x)^{k-1}
+A\frac{(b-x)^n}{n!}\\
 &=-\sum_{k=0}^{n}\frac{f^{(k+1)}(x)}{k!}(b-x)^k+\sum_{k=0}^{n-1}\frac{f^{(k+1)}(x)}{k!}(b-x)^{k}
+A\frac{(b-x)^n}{n!}=-\sum_{k=0}^{n}\frac{f^{(n+1)}(x)}{n!}(b-x)^n+A\frac{(b-x)^n}{n!}\\
 &=\frac{(b-x)^n}{n!}(A-f^{(n+1)}(x)).
\end{align*}

Ainsi, il existe $c\in]a,b[$ tel que $\frac{(b-c)^n}{n!}(A-f^{(n+1)}(c))=0$, et donc, puisque $c\neq b$, tel que $A=f^{(n+1)}(c)$.

L'égalité $g(a)=0$ s'éxrit alors

$$f(b)=\sum_{k=0}^{n}\frac{f^{(k)}(a)}{k!}(b-a)^k+\frac{(b-a)^{n+1}f^{(n+1)}(c)}{(n+1)!},$$

pour un certain réel $c$ de $]a,b[$.
\fincorrection
\correction{005409}
Pour $x\in[a,b]$, posons $g(x)=f(x)-f(a)-\frac{x-a}{2}(f'(x)+f'(a))-A(x-a)^3$ où $A$ est choisi de sorte que 
$g(b)=g(a)=0$ (c'est-à-dire $A=\frac{1}{(b-a)^3}(f(b)-f(a)-\frac{b-a}{2}(f'(b)+f'(a)))$.

$f\in C^2([a,b],\Rr)\cap D^3(]a,b[,\Rr)$ et donc $g\in C^1([a,b],\Rr)\cap D^2(]a,b[,\Rr)$. Pour $x\in[a,b]$, on a~: 
$$g'(x)=f'(x)-\frac{1}{2}(f'(x)+f'(a))-\frac{x-a}{2}f''(x)-3A(x-a)^2,$$ 

puis 

$$g''(x)=\frac{1}{2}f''(x)-\frac{1}{2}f''(x)-\frac{x-a}{2}f^{(3)}(x)-6A(x-a)=\frac{x-a}{2}(-12A-f^{(3)}(x)).$$

$g$ est continue sur $[a,b]$, dérivable sur $]a,b[$ et vérifie de plus $g(a)=g(b)$. Donc, d'après le théorème de \textsc{Rolle}, il existe $d\in]a,b[$ tel que $g'(d)=0$. De même, $g'$ est continue sur $[a,d]\subset[a,b]$, dérivable sur $]a,d[(\neq\emptyset)$ et vérifie de plus $g'(a)=g'(d)(=0)$. D'après le théorème de \textsc{Rolle}, il existe $c\in]a,d[\subset]a,b[$ tel que $g''(c)=0$ ou encore tel que $A=-\frac{1}{12}f^{(3)}(c)$ (puisque $c\neq a$).

En écrivant explicitement l'égalité $g(b)=0$, on a montré que~:

$$\exists c\in]a,b[/\;f(b)=f(a)+\frac{b-a}{2}(f'(b)+f'(a))-\frac{1}{12}f^{(3)}(c)(b-a)^3.$$

Si $f\in C^1([a,b],\Rr)\cap D^2(]a,b[,\Rr)$ et si $F$ est une primitive de $f$ sur $[a,b]$, la formule précédente s'écrit~:

$$\int_{a}^{b}f(t)\;dt=F(b)-F(a)=\frac{b-a}{2}(F'(b)+F'(a))-\frac{1}{12}F^{(3)}(c)(b-a)^3=\frac{b-a}{2}(f(b)+f(a))-\frac{1}{12}f''(c)(b-a)^3.$$

Donc, si $f\in C^1([a,b],\Rr)\cap D^2(]a,b[,\Rr)$, 

$$\exists c\in]a,b[/\;\int_{a}^{b}f(t)\;dt=\frac{b-a}{2}(f(b)+f(a))-\frac{1}{12}f''(c)(b-a)^3.$$

Interprétation géométrique.

Si $f$ est positive, $A_1=\int_{a}^{b}f(t)\;dt$ est l'aire du domaine $D=\{M(x,y)\in\Rr^2/\;a\leq x\leq b\;\mbox{et}\;0\leq y\leq f(x)\}$ et $A_2=\frac{b-a}{2}(f(b)+f(a))$ est l'aire du trapèze $\left(
\begin{array}{c}
a\\
0
\end{array}
\right)\left(
\begin{array}{c}
b\\
0
\end{array}
\right)\left(
\begin{array}{c}
b\\
f(b)
\end{array}
\right)\left(
\begin{array}{c}
a\\
f(a)
\end{array}
\right)$. Si $M_2=\mbox{sup}\{|f''(x)|,\;x\in[a,b]\}$ existe dans $\Rr$, on a~:

$$|A_1-A_2|\leq M_2\frac{(b-a)^3}{12}.$$

$$\includegraphics{../images/img005409-1}$$

\fincorrection
\correction{005412}
\begin{enumerate}
\item  $(X^2-1)^n$ est de degré $2n$ et donc, $L_n$ est de degré $2n-n=n$. Puis, $\mbox{dom}(L_n)=\mbox{dom}((X^{2n})^{(n)})=\frac{(2n)!}{n!}$.

\item  $1$ et $-1$ sont racines d'ordre $n$ de $A_n$ et donc racines d'ordre $n-k$ de $A_n^{(k)}$, pour tout $k$ élément de $\{0,...,n\}$.

Montrons par récurrence sur $k$ que $\forall k\in\{0,...,n\}$, $A_n^{(k)}$ s'annule en au moins $k$ valeurs deux à deux distinctes de l'intervalle $]-1,1[$.

Pour $k=1$, $A_n$ est continu sur $[-1,1]$ et dérivable sur $]-1,1[$. De plus, $A_n(0)=A_n(1)=0$ et d'après le théorème de \textsc{Rolle}, $A_n'$ s'annule au moins une fois dans l'intervalle $]-1,1[$.

Soit $k$ élément de $\{1,...,n-1\}$. Supposons que $A_n^{(k)}$ s'annule en au moins $k$ valeurs de $]-1,1[$. $A_n^{(k)}$ s'annule de plus en $1$ et $-1$ car $k\leq n-1$ et donc s'annule en $k+2$ valeurs au moins de l'intervalle $[-1,1]$. D'après le théorème de \textsc{Rolle}, $A_n^{(k+1)}$ s'annule en au moins $k+1$ points de $]-1,1[$ (au moins une fois par intervalle ouvert).

On a montré que $\forall k\in\{0,...,n\}$, $A_n^{(k)}$ s'annule en au moins $k$ valeurs de $]-1,1[$. En particulier, $A_n^{(n)}=L_n$ s'annule en au moins $n$ réels deux à deux distincts de $]-1,1[$. Puisque $L_n$ est de degré $n$, on a trouvé toutes les racines de $L_n$, toutes réelles, simples et dans $]-1,1[$.
\end{enumerate}
\fincorrection
\correction{005415}
Montrons que ($\forall x>0,\;\left(1+\frac{1}{x}\right)^x<e<\left(1+\frac{1}{x}\right)^{x+1}$. Soit $x>0$.

\begin{align*}\ensuremath
\left(1+\frac{1}{x}\right)^x<e<\left(1+\frac{1}{x}\right)^{x+1}
&\Leftrightarrow x\ln(1+\frac{1}{x})<1<(x+1)\ln(1+\frac{1}{x})\\
 &\Leftrightarrow x(\ln(x+1)-\ln x)<1<(x+1)(\ln(x+1)-\ln x)\Leftrightarrow\frac{1}{x+1}<\ln(x+1)-\ln x<\frac{1}{x}.
\end{align*} 

Soit $x$ un réel strictement positif fixé. Pour $t\in[x,x+1]$, posons $f(t)=\ln t$. $f$ est continue sur $[x,x+1]$ et dérivable sur $]x,x+1[$. Donc, d'après le théorème des accroissements finis, il existe un réel $c$ dans $]x,x+1[$ tel que $f(x+1)-f(x)=(x+1-x)f'(c)$ ou encore

$$\exists c\in]x,x+1[/\;\ln(x+1)-\ln x=\frac{1}{c},$$

ce qui montre que $\forall x>0,\;\frac{1}{x+1}<\ln(x+1)-\ln x<\frac{1}{x}$, et donc que 

$$\forall x>0,\;\left(1+\frac{1}{x}\right)^x<e<\left(1+\frac{1}{x}\right)^{x+1}.$$
\fincorrection
\correction{005416}

$$\includegraphics{../images/img005416-1}$$


Soit $x_0$ un réel non nul. Une équation de la tangente $(T_{x_0})$ à la courbe représentative de $f$ au point d'abscisse $x_0$ est $y=f'(x_0)(x-x_0)+f(x_0)$. $(T_{x_0})$ passe par l'origine si et seulement si 

$$x_0f'(x_0)-f(x_0)=0.$$

Pour $x$ réel, on pose $g(x)=\left\{
\begin{array}{l}
\frac{f(x)}{x}\;\mbox{si}\;x\neq0\\
0\;\mbox{si}\;x=0
\end{array}
\right.$ ($g$ est \og~la fonction pente à l'origine~\fg).

Puisque $f$ est continue et dérivable sur $\Rr$, $g$ est déjà continue et dérivable sur $\Rr^*$.

Puisque $f$ est dérivable en $0$ et que $f(0)=f'(0)=0$, $g$ est de plus continue en $0$.

Finalement, $g$ est continue sur $[0,a]$, dérivable sur $]0,a[$ et vérifie $g(0)=g(a)(= 0)$. D'après le théorème de \textsc{Rolle}, il existe un réel $x_0$ dans $]0,a[$ tel que $g'(x_0)=0$. Puisque $x_0$ n'est pas nul, on a $g'(x_0)=\frac{x_0f'(x_0)-f(x_0)}{x_0^2}$. L'égalité $g'(x_0)=0$ s'écrit $x_0f'(x_0)-f(x_0)=0$ et, d'après le début de l'exercice, la tangente à la courbe représentative de $f$ au point d'abscisse $x_0$ passe par l'origine.
\fincorrection
\correction{005421}
En pensant à l'expression développée de $\Delta$, on voit que $\Delta$ est continue sur $[a,b]$, dérivable sur $]a,b[$ et vérifie $\Delta(a)=\Delta(b)(=0)$ (un déterminant ayant deux colonnes identiques est nul).

Donc, d'après le théorème de \textsc{Rolle}, $\exists c\in]a,b[/\;\Delta'(c)=0$.

Mais, pour $x\in]a,b[$, $\Delta'(x)=f'(x)(g(a)-g(b))-g'(x)(f(a)-f(b))$ (dérivée d'un déterminant). L'égalité $\Delta'(c)=0$ s'écrit~: $f'(c)(g(b)-g(a))=g'(c)(f(b)-f(a))$ ce qu'il fallait démontrer.

(Remarque. Ce résultat généralise le théorème des accroissements finis ($g=Id$ \og~est~\fg~le théorème des accroissements finis.))
\fincorrection
\nocorrection
\correction{001229}
\begin{enumerate}
    \item Pour tout $  n \geq 2  $ on a : $  P_n(0)=-1  $ et $  P_n(1)=3  .$ Comme l'application
$  X\mapsto P_n(X)  $ est continue, elle s'annulle en (au moins)
un point de l'intervalle $ \rbrack 0, 1 \lbrack .$  Comme par
ailleurs, pour tout $  X $ positif, $
P'_n(X)=nX^{n-1}+(n-1)X^{n-2}+2X+1    $ est strictement positif,
l'application $  X\mapsto P_n(X)  $ est strictement croissante sur
$  {\R}_+  $ et s'annule en au plus un point de $  {\R}_+  .$ En
cons\' equence $  P_n   $ a une unique racine positive $  \lambda
_n  $ qui de plus satisfait \`a l'in\' egalit\' e $  0<\lambda
_n<1   .$
    \item Pour tout $  X\in   \rbrack 0, 1
\lbrack   ,   P_n(X)-P_{n-1}(X)=X^n-X^{n-2}<0  .$ En particulier $
P_n(\lambda _{n-1})<0 $ donc $ \lambda _n>\lambda _{n-1}  .$ La
suite $  (\lambda _n)_{n\geq 2}  $ est donc croissante et major\'
ee (cf 1.) : elle est convergente.
    \item Pour tout $  n \geq 2  $ on a  :
$  \lambda _n^n+\lambda _n^{n-1}=-\lambda _n^2-\lambda _n+1  .$ Or
$ \displaystyle{ P_n\Bigl( \frac{3}{ 4}\Bigr) >\Bigl( \frac{3}{
4}\Bigr) ^2+\frac{3}{ 4}-1>0} $ donc la suite $  (\lambda
_n^n+\lambda _n^{n-1})_{n\in {\N}}  $ satisfait aux in\' egalit\'
es $  0<\lambda _n^n+\lambda _n^{n-1}<\displaystyle{ \Bigl(
\frac{3}{ 4}\Bigr) ^n+\Bigl(\frac {3}{ 4} \Bigr) ^{n+1}}  $ et
converge vers $  0  .$ Il en va de m\^eme de la suite $  (-\lambda
_n^2-\lambda _n+1)_{n\geq 2}  .$ En passant \`a la limite, on
obtient l'\' egalit\' e~: $  \ell ^2+\ell -1=0  .$ La seule
solution positive de cette \' equation \' etant $
\displaystyle{\frac { -1+\sqrt{5} }{ 2}}  ,$ on a l'\' egalit\'
e~: $  \ell =\displaystyle{ \frac{ -1+\sqrt{5}}{ 2}}  .$
\end{enumerate}

\bigskip
\noindent \textbf{Remarques}
\begin{enumerate}
    \item  L'in\'egalit\'e $ 0<\lambda _n <1 $ (pour tout $ n\geq 2 $) n'implique pas que $ (\lambda _n^n )_{n\geq 2} $ converge vers $ 0 .$ Par exemple
la suite $ (v_n)_{n \geq 1} $ d\'efinie par $ v_n=\displaystyle
{(1 -\frac{1}{ n})^n} $ converge vers $ \displaystyle{\frac{ 1}{ e
}} .$ (Pour le v\'erifier appliquer le 1. du probl\`eme \`a $
\hbox{log }(v_n) .$)
    \item La propri\'et\'e
$ \displaystyle{\lim _{n\rightarrow \infty} \frac{1}{
n+k}\epsilon \Bigl( \frac{1}{ n+k}  \Bigr) }
 =0  $ n'implique pas que
$ \displaystyle{\lim _{n\rightarrow \infty} \sum _{\scriptstyle
k=0}^{\scriptstyle  n}\frac{1}{ n+k}\epsilon \Bigl( \frac{1}{
n+k}  \Bigr) } =0 .  $ Par exemple... $ \displaystyle{\lim
_{n\rightarrow \infty} \sum _{\scriptstyle k=0}^{\scriptstyle
n}\frac{1}{ n+k}}=\hbox{log }(2) .$
\end{enumerate}

\fincorrection
\nocorrection
\nocorrection
\nocorrection
\correction{001233}
\begin{enumerate}
    \item $\epsilon(x)=\displaystyle{\frac{ f (x)-f (c)}{ x-c}-f '(c)}.$
Comme $  f   $ est continue, $  \epsilon   $ est continue sur $
]a,b[-\{ c\}   $ et la continuit\'e en $  c  $ de $  \epsilon
$ \'equivaut \`a la d\'erivabilit\'e de $  f  $ en $  c  .$
L'unicit\'e est \'evidente.
    \item Pour tout $  n \geq 1  ,   S_{n+1}-S_n=\displaystyle{\frac{1}{ 2n+2}}+
\displaystyle{\frac{1}{ 2n+1}}-\displaystyle{\frac{1}{ n}}<0  $
(par exemple parce que $  \displaystyle{\frac{1}{
2n+2}}<\displaystyle{\frac{1}{ 2n}}  $ et $
\displaystyle{\frac{1}{ 2n+1}}<\displaystyle{\frac{1}{ 2n}}  $
donc $  \displaystyle{\frac{1}{ 2n+2}}+ \displaystyle{\frac{1}{
2n+1}}<2\times \displaystyle{\frac{1}{2n}}  $) donc la suite $
(S_n)_{n\geq 1}  $ est d\' ecroissante. Elle est minor\' ee (par $
0  $) donc elle converge.
    \item Pour tout $  0\leq k\leq n  ,  \displaystyle{\frac{1}{ 2n}}\leq \displaystyle{\frac{1}{ n+k}}\leq
\displaystyle{\frac{1}{ n}}    $  donc $  (n+1)\times
\displaystyle{\frac{1}{ 2n}} \leq S_n \leq (n+1)\times
\displaystyle{\frac{1}{ n}}  $ d'o\`u, en passant \`a la limite,
l'in\'egalit\' e $
 \displaystyle{\frac{1}{ 2}}\leq S \leq 1  .$
    \item Soit $  \epsilon   $ l'application de $  \rbrack -1, 1
\lbrack    $ \`a valeurs dans $  {\R}  $ telle que $
f(x)=f'(0)x+\epsilon (x)  .$ Pour tous $  n, k \in {\N}  , n>0
,$ on a l'\' egalit\' e :
$$\displaystyle{ f\Bigl(\frac {1}{ n+k} \Bigr) =\frac{1}{ n+k}f'(0)+\frac{1}{ n+k}\epsilon \Bigl(
\frac{1}{ n+k} \Bigr)}
$$
donc $  \sigma _n(f)-f'(0)S_n= \displaystyle{\sum _{\scriptstyle
k=0}^{\scriptstyle  n}\frac{1}{ n+k}\epsilon \Bigl( \frac{1}{
n+k} \Bigr) }
  .  $ Comme, pour tout $  k\geq 0  ,$ on a $  \displaystyle{
\frac{1}{ n+k} \leq \frac{1}{ n}}  ,$ on en d\' eduit les in\'
egalit\' es~:

$$ \vert \sigma _n(f)-f'(0)S_n \vert \leq \displaystyle { \frac{1}{ n}
\sum _{\scriptstyle k=0}^{\scriptstyle  n}\Bigl \vert \epsilon
\Bigl( \frac{1}{ n+k} \Bigr) \Bigr \vert \leq  \displaystyle {
\frac{n+1}{ n}}\max  _{\scriptstyle 0\leq k \leq n }\Bigl \vert
\epsilon \Bigl( \frac{1}{ n+k}  \Bigr)\Bigr \vert }.$$ Comme $
\displaystyle{ \max  _{\scriptstyle 0\leq k \leq n }\Bigl \vert
\epsilon \Bigl( \frac{1}{ n+k}  \Bigr)\Bigr \vert \leq \sup
_{\scriptstyle x\in [ 0, \frac{1}{ n}] }\vert \epsilon ( x  )
\vert }  ,$ cette quantit\' e converge vers $  0   $ lorsque $  n
$ tend vers l'infini (puisque $  \epsilon   $ est continue et
s'annulle en $  0  $).
    \item Des \' egalit\' es $  \displaystyle { \hbox{log }\Bigl( 1+\frac{1}{n+k}
\Bigr) =\hbox{log }\Bigl( \frac{n+k+1}{ n+k} \Bigr) =\hbox{log
}(n+k+1)-\hbox{log }(n+k)}   $ on d\'eduit que~: $$ \displaystyle{
\sigma _n(f)= \hbox{log }(2n+1)-\hbox{log }(n)=\hbox{log
}\Bigl(\frac {2n+1}{ n} \Bigr) =\hbox{log }\Bigl( {2+\frac{1}{ n}}
\Bigr)}  . $$ Comme la fonction logarithme est continue, $ (
\sigma _n(f))_{n\geq 1}  $ converge vers $ \hbox{log }(2)  $
lorsque $  n  $ tend vers l'infini. Ainsi $  S  =\hbox{log }(2)
.$
    \item Par les deux questions qui pr\' ec\' edent il est imm\' ediat que $  \displaystyle{ \lim_{n\rightarrow \infty }\sigma _n =\hbox{log }(2)}  .$
    \item Soit $f  :  \rbrack -1, 1\lbrack \rightarrow   {\R}  $ une application continue, d\'erivable en $
 0  $ et telle que $  f(0)=0  .$ Soit $  \epsilon   $ l'application de $  \rbrack -1, 1
\lbrack    $ \`a valeurs dans $  {\R}  $ telle que $
f(x)=f'(0)x+\epsilon (x)  .$

 On pose, pour tous $  n, k\in {\N}  ,   n>0  $~:
$$\sigma _n(p,f) =\displaystyle{\sum _{\scriptstyle k=0}^{\scriptstyle  pn}f\bigl( \frac{1}{ n+k}
\bigr) }  \hbox{ et }S_{n,p} =\displaystyle{\sum _{\scriptstyle
k=0}^{\scriptstyle  pn} \frac{1}{ n+k}  }  .$$

Pour tous $  n, k \in {\N}  ,   n>0  $ on a l'\' egalit\' e : $
\displaystyle{ f\Bigl( \frac{1}{ n+k} \Bigr) =\frac{1}{
n+k}f'(0)+\frac{1}{ n+k}\epsilon \Bigl( \frac{1}{ n+k} \Bigr)}
  $
d'o\`u $$  \vert \sigma _{n}(p,f)-f'(0)S_{n,p}\vert \leq
\displaystyle { \frac{1}{ n} \sum _{\scriptstyle
k=0}^{\scriptstyle  pn}\Bigl\vert \epsilon \Bigl( \frac{1}{
n+k} \Bigr) \Bigr\vert \leq  \displaystyle { \frac{pn+1}{ n}}\sup
_{\scriptstyle x\in [ 0, \frac{1}{ n}] } \Bigl \vert \epsilon
\Bigl( \frac{1}{ n+k}  \Bigr)\Bigr \vert }$$ donc cette
diff\'erence converge vers $  0   $ lorsque $  n   $ tend vers
l'infini.


Lorsque $  f  $ est la fonction $  x\mapsto \hbox{log} (1+x)  $,
on obtient (comme pr\'ec\'edemment) que~:  $$ \displaystyle{
\sigma _n(p,f)= \hbox{log }( (p+1)n+1)-\hbox{log }(n)=\hbox{log
}\Bigl( {1+p +\frac{1}{ n}} \Bigr)}  $$ puis que $  \displaystyle{
\lim _{n\rightarrow \infty }\sigma _n(p,f)=\hbox{log }(p+1)}  $
c'est \`a dire $  S_p=\hbox{log }(p+1)  .$
\end{enumerate}
\fincorrection
\nocorrection
\nocorrection
\nocorrection
\nocorrection
\correction{003982}
$y_i = \frac{x_i}{x_{i+1}}  \Rightarrow 
          \frac n {y_1\cdots y_n} \le \frac{y_1 + \dots + y_n}n$.
\fincorrection
\nocorrection
\nocorrection
\nocorrection
\nocorrection
\nocorrection
\correction{003988}
\begin{enumerate}
  \item 
  \item $f(x) \to \ell \in\R$,\quad
             $f'(x) \to 0$.\par
             TAF entre $x$ et $x/2$ $ \Rightarrow  2\left(f(x) - f\left(\frac x2\right)\right) \le xf'(x) \le 0
              \Rightarrow  xf'(x) \to 0$.
\end{enumerate}
 \fincorrection
\nocorrection
\nocorrection
\nocorrection
\nocorrection
\nocorrection
\correction{003994}
\begin{enumerate}
  \item Fonction décroissante sur $\R^+$.
  \item $f(x) - xf'(x) = -x^2\frac {d}{d x}\left(\frac {f(x)}x\right)$.
    Donc, $x \mapsto \frac {f(x) - p}x \searrow$ et
          $x \mapsto \frac {f(x) - f(0)}x \nearrow$.
  \item $p\le f(x) - mx \le f(x) - xf'(x)$.
\end{enumerate}
 \fincorrection
\correction{003995}
\begin{enumerate}
  \item Soient $x < y$ : $\frac {f(x)-f(0)}{x-0} \ge \frac {f(y)-f(0)}{y-0}
     \Rightarrow  \frac {f(x)}x \ge \frac {f(y)}y + f(0)\left(\frac 1x - \frac 1y\right)$.
  \item    Pour $x < y$ : $f(x+y) \le tf(x) + (1-t)f(y)$ avec $t = \frac x{x-y} < 0$,

     donc $f(x+y) - f(x) - f(y) \le \frac {xy}{x-y}\left( \frac {f(x)}x -
     \frac {f(y)}y \right) \le 0$.
\end{enumerate}
 \fincorrection
\nocorrection
\nocorrection
\nocorrection
\correction{003999}
\begin{enumerate}
  \item 
  \item Prendre $x$ tel que $f(x)$ soit maximal.
  \item 
  \item 
\end{enumerate}
 \fincorrection
\correction{004000}
Pour $a_0$ : $\bigl| f(a_0+h) - f(a_0) - |h| \bigr| \le \frac {|h|}2$.
\fincorrection
\nocorrection
\nocorrection
\correction{004003}
Soit $F(x) = x^2 + xG(x)$.
On a pour $h>0$~: $f(x) \le \frac{F(x+xh)-F(x)}{xh} = 2x + xh + \frac{G(x+xh)-G(x)}h + G(x+xh)$.
Soit $\varepsilon>0$ et $A$ tel que $y\ge A \Rightarrow |G(y)|\le\varepsilon^2$. On prend
$h=\varepsilon/\sqrt{x}$ et on obtient
$f(x)- 2x -\varepsilon\sqrt x \le \varepsilon\sqrt x + \varepsilon^2$
d'où $f(x)\le 2x+ o(\sqrt x)$. L'inégalité inverse se montre de même.
\fincorrection
\correction{000728}
$f'(x) = 2(1-k)^3x+3(1+k)x^2$, $f''(x) = 2(1-k)^2+6(1+k)x$.
Nous avons $f'(0)=0$ et $f''(0)=2(1-k)^3$.
Donc si $k\not=1$ alors, la dérivée seconde étant non nulle en $x=0$, $0$ est un extremum (maximum ou minimum) local.
Si $k=1$ alors $f(x) = 2x^3$ et bien s\^ur $0$ n'est pas un extremum local.
Dans tous les cas $0$ n'est ni un minimum global, ni un maximum global (regardez les limites en $+\infty$ et $-\infty$).
\fincorrection
\nocorrection
\nocorrection
\nocorrection
\nocorrection
\correction{000733}
$f'(x) = 4x^3-3x^2 = x^2(4x-3)$ donc les extremums
appartiennent à  $\{0,\frac 34\}$. Comme $f''(x) = 12x^2-6x
= 6x(2x-1)$. Alors $f''$ ne s'annule pas en $\frac 34$, donc
$\frac 34$ donne un extremum local (qui est même un minimum global).
Par contre $f''(0) = 0$ et $f'''(0)\not=0$ donc $0$ est
un point d'inflexion qui n'est pas un extremum (m\^eme pas local,
pensez \`a un fonction du type  $x \mapsto x^3$). 
\fincorrection
\correction{000734}
\begin{enumerate}
  \item 
$f_\lambda'(x) = \lambda e^x+2x$, $f_\lambda''(x) = \lambda e^x+2$.
Les points d'inflexion sont les racines de $f_\lambda''$,
donc si $\lambda \geq 0$ il n'y a pas de point d'inflexion,
si $\lambda < 0$ alors il y a un point d'inflexion 
en $x_\lambda = \ln(-2/\lambda)$.
  \item Si $\lambda \geq 0$ alors $f_\lambda''$ est toujours strictement positive, donc $f_\lambda'$ est strictement croissante, en $-\infty$
la limite de $f_\lambda'$ est $-\infty$, en $+\infty$ la limite de $f_\lambda'$ est $+\infty$, donc il existe
un unique r\'eel $y_\lambda$ tel que $f'_\lambda(y_\lambda)=0$. 
$f_\lambda$ est d\'ecroissante sur $]-\infty,y_\lambda]$ et
croissante sur $[y_\lambda, +\infty[$. Et en $y_\lambda$
nous avons un minimum absolu.
  \item Nous supposons $\lambda < 0$. Alors $f_\lambda''$ s'annule 
seulement en $x_\lambda$. $f_\lambda'$ est croissante sur $]-\infty,x_\lambda]$ et d\'ecroissante sur $[x_\lambda, +\infty[$.
Donc $f_\lambda'$ a des racines si et seulement si
$f'(x_\lambda) \geq 0$. Or $f'(x_\lambda) = -2+2x_\lambda$.
  \begin{enumerate}
  \item Si $\lambda = -2/e$ alors $f_\lambda'(x_\lambda) = 0$.
  Comme $f_\lambda''(x_\lambda)=0$ et $f_\lambda'''$ ne s'annule pas alors $x_\lambda$ est un point d'inflexion qui 
n'est pas un extremum local.
  \item Si $\lambda>-2/e$ alors $f_\lambda'(x_\lambda) <0$
  donc $f_\lambda'$ est n\'egative donc $f$ est strictement d\'ecroissante.
  Il n'y a pas d'extremum local.
  \item Si $-2/e<\lambda<0$ alors $f_\lambda'(x_\lambda) > 0$.
  Donc $f_\lambda'$ s'annule en deux points,
  une fois sur $]-\infty,x_\lambda[$ et une sur $]x_\lambda,+\infty[$.
  Ce sont des extremums locaux (minimum et maximum respectivement).
  \end{enumerate}
\end{enumerate}
\fincorrection
\nocorrection
\nocorrection
\nocorrection
\correction{000738}

 Le th\'eor\`eme de Rolle dit que si $h :[a,b]
\longrightarrow \R$ est une fonction continue sur l'intervalle
ferm\'e $[a,b]$ et d\'erivable sur l'ouvert $]a,b[$ alors il
existe $c \in ]a,b[$ tel que $h' (c) = 0.$

  \begin{enumerate}
  \item Supposons par l'absurde, qu'il existe $x_{0}\in ]a,b]$  tel
que $g (x_{0}) = g (a).$ Alors en appliquant le th\'eor\`eme de
Rolle \`a la restriction de $g$ \`a l'intervalle $[a,x_{0}]$ (les
hypoth\`eses \'etant clairement v\'erifi\'ees), on en d\'eduit
qu'il existe $c \in ]a,x_{0}[$ tel que $g' (c) = 0,$ ce qui
contredit les hypoth\`eses faites sur $g.$ Par cons\'equent on a
d\'emontr\'e que $g (x) \neq g (a)$ pour tout $x \in ]a,b].$

  \item D'apr\`es la question pr\'ec\'edente, on a en particulier  $g (b)
\neq g (a)$ et donc $p$ est un nombre r\'eel bien d\'efini et $h =
f - p \cdot g$ est alors une fonction continue sur $[a,b]$ et
d\'erivable sur $]a,b[.$ Un calcul simple montre que $h (a) = h
(b).$ D'apr\`es le th\'eor\`eme de Rolle il en r\'esulte  qu'il
existe $c \in ]a,b[$ tel que $h' (c) = 0.$ Ce qui implique la
relation requise.

  \item  Pour chaque $x \in ]a,b[,$ on peut appliquer
la question 2. aux
 restrictions de $f$ et $g$ \`a l'intervalle $[x,b],$ on en d\'eduit qu'il
 existe un point $c (x) \in ]x,b[,$ d\'ependant de $x$ tel que
 $$\frac{f (x) - f (b)}{g (x) - g (b)} = \frac{f' (c (x))}{g' (c (x))} .
 \leqno (*)$$
 Alors, comme $\lim_{x \to b^{-}} \frac{f' (t)}{g' (t)} = \ell$ et
 $\lim_{x \to b^{-}} c (x) = b,$ (car $c(x) \in ]x,b[$) on en d\'eduit en passant \`a la
 limite dans $(*)$ que
 $$ \lim_{x \to b^{-}} \frac{f (x) - f (b)}{g (x) - g (b)} = \ell.$$
 Ce r\'esultat est connu sous le nom de ``r\`egle de
 l'H\^opital''.

 \item Consid\'erons les deux fonctions $f
 (x) = \Arccos x$ et $g (x) = \sqrt{1-x^2}$ pour $x \in [0,1].$ Ces fonctions sont continues sur  $[0,1]$ et
 d\'erivables sur $]0,1[$ et $f' (x) = - 1 \slash \sqrt{1-x^2}$,
 $g' (x) = - x \slash \sqrt{1-x^2} \neq 0$ pour tout $ x
 \in ]0,1[.$ En appliquant les r\'esultats de la question 3., on en
 d\'eduit que
 $$ \lim_{x \to 1^{-}} \frac{\Arccos x}{\sqrt{1-x^2}} = \lim_{x \to 1^{-}} \frac{\frac{-1}{\sqrt{1-x^2}}} {\frac{-x}{\sqrt{1-x^2}}}= \lim_{x \to 1^{-}} \frac{1}{x} = 1.$$
\end{enumerate}
\fincorrection
\correction{000739}
\begin{enumerate}
\item
  \begin{enumerate}
  \item
 Il est clair que la fonction $f$ est d\'erivable sur $\R^{+}$
puisque c'est une fonction rationnelle sans p\^ole dans cet
intervalle. De plus d'apr\`es la formule de la d\'eriv\'ee d'un
quotient, on obtient pour $x \geq 0$ :
$$ f' (x) = \frac{n (x^{n-1} - 1)}{(1 + x)^{n + 1}}.$$
  \item  Par l'expression pr\'ec\'edente $f' (x)$ est du signe de $x^{n - 1} - 1$ sur $\R^{+}.$ Par
cons\'equent on obtient: $ f' (x) \leq 0$ pour $0 \leq x \leq 1$
et $ f' (x) \geq 0$ pour $x \geq 1.$ Il en r\'esulte que $f$ est
d\'ecroissante sur $[0 , 1]$ et croissante sur $[1, + \infty[$ et
par suite $f$ atteint son minimum sur $\R^{+}$ au point $1$ et ce
minimum vaut $f (1) = 2^{1 - n}.$
  \end{enumerate}
\item
  \begin{enumerate}
  \item
 Il   r\'esulte de la
question 1.b que $f (x) \geq f (1)$ pour tout $x \in \R^{+}$ et
donc
$$ (1 + x)^n \leq 2^{n - 1} (1 + x^n).$$
 \item En appliquant l'in\'egalit\'e pr\'ec\'edente avec $ x = b \slash
a,$ on en d\'eduit imm\'ediatement l'in\'egalit\'e requise (le cas du couple $(0,0)$ étant trivial).
  \end{enumerate}
\end{enumerate}
\fincorrection
\correction{000740}
\begin{enumerate}
\item $f$ est d\'erivable sur $\R_-^*$ en tant que compos\'ee de fonctions
  d\'erivables, et sur $\R_+^*$ car elle est nulle sur cet intervalle ; \'etudions donc la
  d\'erivabilit\'e en $0$.

On a
$$\frac{f(t)-f(0)}{t}=\begin{cases}
  e^{1/t}/t & \mathrm{\ si\ } t<0\\
0  & \mathrm{\ si\ } t > 0
\end{cases}$$
or $e^{1/t}/t$ tend vers $0$ quand $t$ tend vers $0$ par valeurs
n\'egatives. Donc $f$ est d\'erivable \`a gauche et \`a droite en 0 et ces
d\'eriv\'ees sont identiques, donc $f$ est d\'erivable et $f'(0)=0$.
\item On a
$$f'(t)=\begin{cases}
  -e^{1/t}/t^2 & \mathrm{\ si\ } t<0\\
0  & \mathrm{\ si\ } t \geq 0
\end{cases}$$
donc le taux d'accroissement de $f'$ au voisinage de 0 est
$$\frac{f'(t)-f'(0)}{t}=\begin{cases}
  -e^{1/t}/t^3& \mathrm{\ si\ } t<0\\
0  & \mathrm{\ si\ } t > 0
\end{cases}$$
et il tend vers $0$ quand $t$ tend vers $0$ par valeurs sup\'erieures
comme inf\'erieures. Donc $f$ admet une d\'eriv\'ee seconde en $0$, et
$f''(0)=0$.

\item
  \begin{enumerate}
  \item On a d\'ej\`a trouv\'e que $f'(t)=-e^{1/t}/t^2$, donc  $f'(t)=P_1(t)/t^2
    e^{1/t}$ si on pose $P_1(t)=-1$.

Par ailleurs, $f''(t)=e^{1/t}/t^4+e^{1/t}
(2/t^3)=\frac{1+2t}{t^4}e^{1/t}$ donc la formule est vraie pour
$n=2$ en posant $P_2(t)=1+2t$.
\item Supposons que la formule est vraie au rang $n$.
Alors $f^{(n)}(t)=\frac{P_n(t)}{t^{2n}}e^{1/t}$ d'o\`u
\begin{equation*}
\begin{split}
  f^{(n+1)}(t)&=\frac{P'_n(t)t^{2n}-P_n(t)(2n)t^{2n-1}}{t^{4n}}e^{1/t}+\frac{P_n(t)}{t^{2n}}e^{1/t}(-1/t^2)\\
&= \frac{P'_n(t)t^{2}-(2n t+1) P_n(t)}{t^{2(n+1)}}e^{1/t}
\end{split}
\end{equation*}
donc la formule est vraie au rang $n+1$ avec
$$P_{n+1}(t)=P'_n(t)t^{2}-(2n t+1) P_n(t).$$
  \end{enumerate}
\item Sur $\R_-^*$ et sur $\R_+^*$, $f$ est ind\'efiniment d\'erivable, donc il
  suffit d'\'etudier ce qui se passe en $0$.

Montrons  par r\'ecurrence que $f$ est ind\'efiniment d\'erivable en $0$,
et que pour tout   $n \in \N,   f^{(n)}(0)=0$. On sait que
  c'est vrai au rang 1. Supposons que $f$ est $n$-fois d\'erivable, et que
  $f^{(n)}(0)=0$. Alors le taux d'accroissement de $f^{(n)}$ en $0$ est :
$$\frac{f^{(n)}(t)-f^{(n)}(0)}{t}=\begin{cases}
  P_n(t)e^{1/t}/t^{2n+1} & \mathrm{\ si\ } t<0\\
0  & \mathrm{\ si\ } t > 0
\end{cases}$$
et sa limite est $0$ quand  $t$ tend vers $0$ par valeurs sup\'erieures
comme inf\'erieures. Donc $f^{(n)}$ est d\'erivable  en $0$, et
$f^{(n+1)}(0)=0$. Donc l'hypoth\`ese de r\'ecurrence est v\'erifi\'ee  au
rang $n+1$.

Par cons\'equent, $f$ est de classe $C^\infty$.
\end{enumerate}
\fincorrection
\nocorrection
\nocorrection
\correction{003934}
Contre-exemple~: $f(t) = t^2$ si $t\ge 0$ et $f(t) = -t^2$ si $t<0$.
\fincorrection
\nocorrection
\nocorrection
\nocorrection
\nocorrection
\correction{003939}
Poser $g(x)=  \int_{t=0}^x f^2(t)\,d t$.
On obtient $(g^3)'(x)\to 3\ell^2$ lorsque $x\to+\infty$,
ce qui implique (classiquement) que $g^3(x)\sim 3\ell^2x$,
puis $f(x)\sim \sqrt[3]{\frac\ell{3x}}$.
\fincorrection
\nocorrection
\nocorrection
\nocorrection
\correction{003943}
TAF $ \Rightarrow  \ell = \sqrt2$.
\fincorrection
\nocorrection
\nocorrection
\nocorrection
\nocorrection
\nocorrection
\correction{003949}
Dériver par rapport à $a$ puis par rapport à $b$.
\fincorrection
\nocorrection
\nocorrection
\nocorrection
\correction{003953}
\begin{enumerate}
  \item 
  \item $f^{(n)}(1) = 2^nn!$, $f^{(n)}(-1) = (-2)^nn!$.
  \item 
\end{enumerate}
 \fincorrection
\nocorrection
\nocorrection
\nocorrection
\correction{003957}
$\frac {f'(0)}2$.
\fincorrection
\nocorrection
\correction{003959}
$f = \mathrm{id}$.
\fincorrection
\nocorrection
\nocorrection
\correction{003962}
\begin{enumerate}
  \item  AF $ \Rightarrow  \exists\ w(x)$ compris entre $u(x)$ et $v(x)$ tel que
    $\frac {u^v - v^v}{u - v} = vw^{v-1} \to a^a$.
  \item     $u^v - v^u = (u^v - v^v) + (v^v - v^u) = (u-v)(vw_1^{v-1} - (\ln v)v^{w_2})$

    $u^u - v^v = (u-v)w_3^{w_3}(1 + \ln w_2)$

    $ \Rightarrow  \lim = \frac {1-\ln a}{1 + \ln a}$.
\end{enumerate}
 \fincorrection
\correction{003963}
\begin{enumerate}
  \item 
  \item   Pour $k \ge 0$, la suite $(u_n)$ est croissante et $\ln u_n \le \frac k{1-k}$.

    Pour $k < 0$, $(u_{2n})$ décro\^\i t et converge, et $u_{2n+1} \sim u_{2n}$.
\end{enumerate}
 \fincorrection
\nocorrection
\correction{003965}
$(1+x^2)f^{(n+1)} + 2nxf^{(n)} + n(n-1)f^{(n-1)} = 0$ pour $n \ge 1$.\par
         $g^{(n+1)} = 3x^2g^{(n)} + 6nxg^{(n-1)} + 3n(n-1)g^{(n-2)}$ pour $n \ge 0$.
\fincorrection
\correction{003966}
$(-1)^ne^{-x}\bigl( x^3 + (2-3n)x^2 + (3n^2-7n)x + (-n^3+5n^2-4n-5)\bigr)$.
\fincorrection
\correction{003967}
$f^{(n)}(x) = 2^ne^{x\sqrt3}\sin(x+n\frac\pi6)$.
\fincorrection
\correction{003968}
$\frac{d^n}{d x^n}\bigl(x^n(1-x)^n\bigr) =
         \sum_{k=0}^n n!\,\bigl(C_n^k\bigr)^2 (-1)^{n-k} x^{n-k}(1-x)^k$.
         \par
         coefficient de $x^n = (-1)^n n!\, \sum_{k=0}^n \bigl(C_n^k\bigr)^2
                             = (-1)^n A_{2n}^n$.
\fincorrection
\correction{003969}
$\frac {(n-1)!}t$,
  \qquad
    $\frac {(-1)^n}{t^{n+1}}\exp(1/t)$.
 \fincorrection
\correction{003970}
\begin{enumerate}
  \item $a_{n+1,k} = a_{n,k-1} + 2(2k-n)a_{n,k}$.
  \item $a_{n,k} = \frac {n!}{(n-k)!(2k-n)!}$.
\end{enumerate}
 \fincorrection
\nocorrection
\nocorrection
\nocorrection
\nocorrection
\correction{003975}
\begin{enumerate}
  \item  Au point d'abscisse $\alpha \text{ tq } \ln\alpha = \frac {a\ln a}{1-a} + 1$ pour
    $\cal C$, et $\alpha' = a\alpha$ pour ${\cal C}'$.
  \item Centre = $\left(0,\frac {a\ln a}{1-a}\right)$, rapport = $a$.
\end{enumerate}
 \fincorrection
\correction{003976}
Si $f$ change de signe, soit par exemple $f(a) > 0$,
$f(b) < 0$, $a<b$ et $c = \sup\{x\text{ tq }f_{|[a,x]}\text{ est croissante}\}$.
Alors $f$ est croissante sur $[a,c]$ et $f(c) = 0$, contradiction.
\fincorrection
\correction{003977}
Si l'on pose $F(x)= \int_{t=0}^x e^{t^2}\, d t$, on constate que
$a(x)=F^{-1}(1+F(x))$ ce qui prouve l'existence, l'unicité et
le caractère $\mathcal{C}^\infty$ de~$a$. Pour la symétrie, il faut montrer que
$a(-a(x)) = -x$ soit $ \int_{t=-a(x)}^{-x} e^{t^2}\, d t=1$ ce qui est immédiat.
\fincorrection
\correction{003978}
Toute fonction linéaire $\varphi$ : $x \mapsto ax$ convient. Réciproquement, si~$\varphi$
est solution alors $\varphi(0)=0$. On note~$a=\varphi'(0)$ et $\psi(x) = \varphi(x)-ax$~:
$\psi$ est également solution et $\psi'(0)=0$. Si $x\in\R$ et~$n\in\N$ alors $\smash{\psi(x) = 2^n\psi(x/2^n) \to 0}$ lorsque $n\to\infty$,
d'où $\psi=0$ et $\varphi(x)=ax$.
\fincorrection
\correction{003979}
\begin{enumerate}
  \item $m=\frac{1+\sqrt5}2$, $c=m^{-1/m}$.
  \item 
  \item 
  \item $f$ et $f'$ ont des limites nulles en $0^+$ et infinies en~$+\infty$
    donc $f(x) = o_{x\to0^+}(x)$ et $x = o_{x\to+\infty}(f(x))$,
    ce qui implique que $f(x)-x$ s'annule sur~$]0,+\infty[$.
    S'il y a deux points fixes, $a<b$, alors par le thm. des accroissements
    finis l'équation $f'(x) = 1$ admet une solution dans~$]0,a[$ et une dans~$]a,b[$,
    en contradiction avec la bijectivité de~$f' = f^{-1}$.
  \item On note~$a$ le point fixe de~$f$, $b$ celui de~$g$ et on suppose
    $a\ne b$, par exemple $a<b$. On a $g(x) < x$ pour $x\in{]0,b[}$ donc $g(a) < a=f(a)$.
    Par conséquent $g(x) < x \le f(x)$ si $x\in{[a,b[}$~; soit $]c,b[$ le plus grand
    intervalle sur lequel $g(x) < f(x)$.
    On a $0\le c < a$, $g(c^+) = f(c^+) \le c$ et $f$, $g$ sont strictement
    croissantes, donc $g^{-1}(x) > f^{-1}(x)$ pour $x\in{]c,b[}$. Ainsi
    $g-f$ est strictement croissante sur~$]c,b[$ a une limite nulle en~$c^+$
    et est négative en~$b$, c'est absurde.
    
    Remarque~: le point fixe est le nombre d'or~$m$. De plus, si $f$ et $g$
    sont deux éléments de~$E$ distincts alors $f-g$ n'est de signe constant
    sur aucun voisinage de~$m^-$ (même démonstration).
\end{enumerate}
 \fincorrection
\correction{003980}
On a $f'\le 0$ donc $f$ est décroissante sur~$\R$,
et en particulier elle admet des limites finies, $a$ et $b$ en $-\infty$ et $+\infty$
avec $-1\le b\le a\le 1$.

Supposons $a>0$~: soit $\alpha\in{]0,a[}$. Il existe $x_0\in\R$
tel que $\forall\ x\le x_0,\ f(x)\ge a-\alpha > 0$, d'où
$f'(x)\le -1+\sqrt{a-\alpha}< 0$. Ceci est incompatible avec le caractère
borné de~$f$, donc on a en fait $a\le 0$.
On montre de même que $b\ge 0$ et comme $b\le a$, on a finalement
$a=b=0$.
\fincorrection
\correction{005407}
$f'$ est continue sur le segment $[a,b]$ et donc est bornée sur ce segment. Soit $M=\mbox{sup}\{f'(x),\;x\in[a,b]\}$, et soit $g$ la fonction affine qui prend les mêmes valeurs que $f$ en $a$ et $b$ (c'est-à-dire $\forall x\in[a,b],\;g(x)=\frac{f(b)-f(a)}{b-a}(x-a)+f(a)=$) puis $h=f-g$. On va montrer que $h=0$ sous l'hypothèse $M=\frac{f(b)-f(a)}{b-a}$.

$h$ est dérivable sur $[a,b]$ et, pour $x\in[a,b]$, $h'(x)=f'(x)-\frac{f(b)-f(a)}{b-a}=f'(x)-M\leq0$. $h$ est donc décroissante sur $[a,b]$. Par suite, $\forall x\in[a,b],\;0=h(b)\leq h(x)\leq h(a)=0$. Ainsi, $\forall x\in[a,b],\;h(x)=0$, ou encore $f=g$. $f$ est donc affine sur $[a,b]$.
\fincorrection
\correction{005410}
Supposons que $f$ est convexe sur un intervalle ouvert $I=]a,b[$ ($a$ et $b$ réels ou infinis).

Soit $x_0\in I$. On sait que la fonction pente en $x_0$ est croissante.

Pour $x\neq x_0$, posons $g(x)=\frac{f(x)-f(x_0)}{x-x_0}$. Soit $x'$ un élément de $]x_0,b[$. $\forall x\in]a,x_0[$, on a $g(x)<g(x')$, ce qui montre que $g$ est majorée au voisinage de $x_0$ à gauche. Etant croissante, $g$ admet une limite réelle quand $x$ tend vers $x_0$ par valeurs inférieures ou encore, $\lim_{x\rightarrow x_0,\;x<x_0}\frac{f(x)-f(x_0)}{x-x_0}$ existe dans $\Rr$. $f$ est donc dérivable à gauche en $x_0$. On montre de même que $f$ est dérivable à droite en $x_0$.

Finalement, $f$ est dérivable à droite et à gauche en tout point de $]a,b[$. En particulier, $f$ est continue à droite et à gauche en tout point de $]a,b[$ et donc continue sur $]a,b[$.
\fincorrection
\correction{005411}
\begin{enumerate}
\item  La fonction $f~:~x\mapsto\ln x$ est deux fois dérivable sur $]0,+\infty[$ et, pour $x>0$, $f''(x)=-\frac{1}{x^2}<0$. Par suite, $f$ est concave sur $]0,+\infty[$. On en déduit que~:

$$\forall n\in\Nn,\;\forall(x_1,...,x_n)\in(]0,+\infty[)^n,\;\forall(\alpha_1,...,\alpha_n)\in(]0,1[)^n,\;
(\sum_{k=1}^{n}\alpha_k=1\Rightarrow\ln(\sum_{k=1}^{n}\alpha_kx_k)\geq\sum_{k=1}^{n}\alpha_k\ln(x_k),$$

et donc par croissance de $f$ sur $]0,+\infty[$,

$$\prod_{k=1}^{n}x_k^{\alpha_k}\leq\sum_{k=1}^{n}\alpha_kx_k.$$

Si l'un des $x_k$ est nul, l'inégalité précédente est immédiate.

En choisissant en particulier $\alpha_1=...=\alpha_n=\frac{1}{n}$, de sorte que $(\alpha_1,...,\alpha_n)\in(]0,1[)^n$ et que $\sum_{k=1}^{n}\alpha_k=1$, on obtient
 
$$\forall n\in\Nn^*,\;\forall(x_1,...,x_n)\in([0,+\infty[)^n,\;\sqrt[n]{x_1...x_n}\leq\frac{1}{n}(x_1+...+x_n).$$

\item 
\begin{enumerate}
\item Soient $p$ et $q$ deux réels strictement positifs vérifiant $\frac{1}{p}+\frac{1}{q}=1$ (de sorte que l'on a même $\frac{1}{p}<\frac{1}{p}+\frac{1}{q}=1$ et donc $p>1$ et aussi $q>1$).

Si $a=0$ ou $b=0$, l'inégalité proposée est immédiate.

Soit alors $a$ un réel strictement positif puis, pour $x\geq0$, $f(x)=\frac{a^p}{p}+\frac{x^q}{q}-ax$.

$f$ est dérivable sur $[0,+\infty[$ (car $q>1$) et pour $x\geq0$, $f'(x)=x^{q-1}-a$.
$q$ étant un réel strictement plus grand que $1$, $q-1$ est strictement positif et donc, la fonction $x\mapsto x^{q-1}$ est strictement croissante sur $[0,+\infty[$. Par suite,

$$f'(x)>0\Leftrightarrow x^{q-1}>a\Leftrightarrow x>a^{1/(q-1)}.$$

$f$ est donc strictement décroissante sur $[0,a^{1/(q-1)}]$ et strictement croissante sur $[a^{1/(q-1)},+\infty[$. Ainsi,

$$\forall x\geq0,\;f(x)\geq f(a^{1/(q-1)})=\frac{1}{p}a^{p}+\frac{1}{q}a^{q/(q-1)}-a.a^{1/(q-1)}.$$

Maintenant, $\frac{1}{p}+\frac{1}{q}=1$ fournit $q=\frac{p}{p-1}$ puis $q-1=\frac{1}{p-1}$. Par suite, $\frac{q}{q-1}=p$. Il en résulte que

$$\frac{1}{p}a^{p}+\frac{1}{q}a^{q/(q-1)}-a.a^{1/(q-1)}=(\frac{1}{p}+\frac{1}{q}-1)a^p=0.$$

$f$ est donc positive sur $[0,+\infty[$, ce qui fournit $f(b)\geq0$. De plus, 

$$f(b)=0\Leftrightarrow b=a^{1/(q-1)}\Leftrightarrow b^q=a^{q/(q-1)}\Leftrightarrow b^q=a^p.$$

\item Soient $A=\sum_{k=1}^{n}|a_k|^p$ et $B=\sum_{k=1}^{n}|b_k|^q$.

Si $A=0$, alors $\forall k\in\{1,...,n\},\;a_k=0$ et l'inégalité est immédiate. De même, si $B=0$.

Si $A>0$ et $B>0$, montrons que $\sum_{k=1}^{n}\frac{|a_k|}{A^{1/p}}\frac{|b_k|}{B^{1/q}}\leq1$.

D'après a),

$$\sum_{k=1}^{n}\frac{|a_k|}{A^{1/p}}\frac{|b_k|}{B^{1/q}}\leq
\sum_{k=1}^{n}(\frac{1}{p}\frac{|a_k|^p}{A}+\frac{1}{q}\frac{|b_k|^q}{B})=\frac{1}{pA}.A+\frac{1}{qB}.B=1,$$

ce qu'il fallait démontrer.

\item Pour $p>1$, la fonction $x\mapsto x^p$ est deux fois dérivable sur $]0,+\infty[$ et $(x^p)''=p(p-1)x^{p-2}>0$. Donc, la fonction $x\mapsto x^p$ est strictement convexe sur $]0,+\infty[$ et donc sur $[0,+\infty[$ par continuité en $0$. Donc,

$$\forall(x_1,...,x_n)\in(]0,+\infty[)^n,\;\forall(\lambda_1,...,\lambda_n)\in([0,+\infty[)^n\setminus\{(0,...,0)\},
\;\left(\frac{\sum_{k=1}^{n}\lambda_kx_k}{\sum_{k=1}^{n}\lambda_k}\right)^p\leq
\frac{\sum_{k=1}^{n}\lambda_kx_k^p}{\sum_{k=1}^{n}\lambda_k},$$

et donc 
$$\sum_{k=1}^{n}\lambda_kx_k\leq(\sum_{k=1}^{n}\lambda_k)^{1-\frac{1}{p}}(\sum_{k=1}^{n}\lambda_kx_k^p)^{\frac{1}{p}}.$$

On applique alors ce qui précède à $\lambda_k=|b_k|^q$ puis $x_k=\lambda_k^{-1/p}|a_k|$ (de sorte que $\lambda_kx_k=|a_kb_k|$) et on obtient l'inégalité désirée.

\item Pour $p=q=2$, c'est l'inégalité de \textsc{Cauchy}-\textsc{Schwarz} démontrée dans une planche précédente.
\end{enumerate}
\end{enumerate}
\fincorrection
\correction{005414}
$f$ est de classe $ ^\infty$ sur $\Rr^*$ en vertu de théorèmes généraux.

Montrons par récurrence que $\forall n\in\Nn,\;\exists P_n\in\Rr[X]/\;\forall x\in\Rr^*,\;f^{(n)}(x)=\frac{P_n(x)}{x^{3n}}e^{-1/x^2}$.

C'est vrai pour $n=0$ avec $P_0=1$.

Soit $n\geq0$. Supposons que $\exists P_n\in\Rr[X]/\;\forall x\in\Rr^*,\;f^{(n)}(x)=\frac{P_n(x)}{x^{3n}}e^{-1/x^2}$. Alors, pour $x\in\Rr^*$,

$$f^{(n+1)}(x)=(\frac{2}{x^3}\frac{P_n(x)}{x^{3n}}+(P_n'(x)\frac{1}{x^{3n}}-3nP_n(x)\frac{1}{x^{3n+1}})e^{-1/x^2}
=\frac{P_{n+1}(x)}{3^{3(n+1)}}e^{-1/x^2},$$

où $P_n+1=2P_n+X^3P_n'-3nX^2P_n$ est un polynôme. On a montré que 

$$\forall n\in\Nn,\;\exists P_n\in\Rr[X]/\;\forall x\in\Rr^*,\;f^{(n)}(x)=\frac{P_n(x)}{x^{3n}}e^{-1/x^2}.$$

Montrons alors par récurrence que pour tout entier naturel $n$, $f$ est de classe $C^n$ sur $\Rr$ et que $f^{(n)}(0)=0$.

Pour $n=0$, $f$ est continue sur $\Rr^*$ et de plus, $\lim_{x\rightarrow 0,\;x\neq0}f(x)=0=f(0)$. Donc, $f$ est continue sur $\Rr$.

Soit $n\geq0$. Supposons que $f$ est de classe $C^n$ sur $\Rr$ et que $f^{(n)}(0)=0$. Alors, d'une part $f$ est de classe $C^n$ sur $\Rr$ et $C^{n+1}$ sur $\Rr^*$ et de plus, d'après les théorèmes de croissances comparées, $f^{(n+1)}(x)=\frac{P_{n+1}(x)}{x^{3n+3}}e^{-1/x^2}$ tend vers $0$ quand $x$ tend vers $0$, $x\neq 0$. D'après un théorème classique d'analyse, $f$ est de classe $C^{n+1}$ sur $\Rr$ et en particulier, $f^{(n+1)}(0)=\lim_{x\rightarrow 0,\;x\neq0}f^{(n+1)}(x)=0$.

On a montré par récurrence que $\forall n\in\Nn$, $f$ est de classe $C^n$ sur $\Rr$ et que $f^{(n)}(0)=0$. $f$ est donc de classe $C^\infty$ sur $\Rr$.
\fincorrection
\correction{005417}
\begin{enumerate}
\item  Soit $m$ un élément de $]f'(a),f'(b)[$. Puisque $\lim_{h\rightarrow 0}\frac{f(a+h)-f(a)}{h}=f'(a)$ et que 
$\lim_{h\rightarrow 0}\frac{f(b+h)-f(b)}{h}=f'(b)$, on a (en prenant par exemple $\varepsilon=\mbox{Min}\{m-f'(a),f'(b)-m\}>0$) 

$$\begin{array}{l}
\exists h_1>0/\;\forall h\in]0,h_1[,\;(a+h\in I\Rightarrow\frac{f(a+h)-f(a)}{h}<m\;\mbox{et}\\
\exists h_2>0/\;\forall h\in]0,h_2[\;(b+h\in I\Rightarrow\frac{f(b+h)-f(b)}{h}> m.
\end{array}$$

L'ensemble $E=\{h\in]0,\mbox{Min}\{h_1,h_2\}[/\;a+h\;\mbox{et}\;b+h\;\mbox{sont dans}\;I\}$ n'est pas vide (car $I$ est ouvert) et pour tous les $h$ de $E$, on a~:$\frac{f(a+h)-f(a)}{h}<m<\frac{f(b+h)-f(b)}{h}$.

$h>0$ est ainsi dorénavant fixé.

\item  La fonction $f$ est continue sur $I$ et donc, la fonction $g~:~x\mapsto\frac{f(x+h)-f(x)}{h}$ est continue sur $[a,b]$. D'après le théorème des valeurs intermédiaires, comme $g(a)<m<g(b)$, $\exists y\in[a,b]/\;g(y)=m$ ou encore $\exists y\in[a,b]/\;\frac{f(y+h)-f(y)}{h}=m$.

Maintenant, d'après le théorème des accroissements finis, $\exists x\in]y,y+h[\subset I/\;m=\frac{f(y+h)-f(y)}{h}=f'(x)$.

Donc une fonction dérivée n'est pas nécessairement continue mais vérifie tout de même le théorème des valeurs intermédiaires (Théorème de \textsc{Darboux}).
\end{enumerate}
\fincorrection
\correction{005419}
Quand $x$ tend vers $0$ par valeurs supérieures,

$$\frac{\cos(\sqrt{x})-1}{x}=\frac{1}{2}\frac{\cos(\sqrt{x})-1}{(\sqrt{x})^2/2}\rightarrow-\frac{1}{2}.$$

$f$ est donc dérivable à droite en $0$ et $f_d'(0)=-\frac{1}{2}$.

Autre solution. $f$ est continue sur $\Rr$ et de classe $C^1$ sur $\Rr^*$ en vertu de théorèmes généraux. Pour $x\neq0$, $f'(x)=-\frac{1}{2\sqrt{x}}\sin(\sqrt{x})$. Quand $x$ tend vers $0$, $f'$ tend vers $-\frac{1}{2}$. En résumé, $f$ est continue sur $\Rr$, de classe $C^1$ sur $\Rr^*$ et $f'$ a une limite réelle quand $x$ tend vers $0$ à savoir $0$. On en déduit que $f$ est de classe $C^1$ sur $\Rr$ et en particulier, $f$ est dérivable en $0$ et $f'(0)=-\frac{1}{2}$.

\fincorrection
\correction{005420}
Soit $n\geq2$ le degré de $P$.
\begin{enumerate}
\item  Si $P$ admet $n$ racines réelles simples, le théorème de \textsc{Rolle} fournit au moins $n-1$ racines réelles deux à deux distinctes pour $P'$. Mais, puisque $P'$ est de degré $n-1$, ce sont toutes les racines de $P'$, nécessairement toutes réelles et simples.

(Le résultat tombe en défaut si les racines de $P$ ne sont pas toutes réelles. Par exemple, $P=X^3-1$ est à racines simples dans $\Cc$ mais $P'=3X^2$ admet une racine double)

\item  Séparons les racines simples et les racines multiples de $P$. Posons $P=(X-a_1)...(X-a_k)(X-b_1)^{\alpha_1}...(X-b_l)^{\alpha_l}$ où les $a_i$ et les $b_j$ sont $k+l$ nombres réels deux à deux distincts et les $\alpha_j$ des entiers supérieurs ou égaux à $2$ (éventuellement $k=0$ ou $l=0$ et dans ce cas le produit vide vaut conventionnellement $1$).

$P$ s'annule déjà en $k+l$ nombres réels deux à deux distincts et le théorème de \textsc{Rolle} fournit $k+l-1$ racines réelles deux à deux distinctes et distinctes des $a_i$ et des $b_j$. D'autre part, les $b_j$ sont racines d'ordre $\alpha_j$ de $P$ et donc d'ordre $\alpha_j-1$ de $P'$. On a donc trouvé un nombre de racines (comptées en nombre de fois égal à leur ordre de multiplicité) égal à $k+l-1+\sum_{j=1}^{l}(\alpha_j-1)=k+\sum_{j=1}^{l}\alpha_j-1=n-1$ racines réelles et c'est fini.
\end{enumerate}
\fincorrection
\correction{005422}
Puisque $\lim_{x\rightarrow +\infty}xf'(x)=1$, $\exists A>0/\;\forall x>0,\;(x\geq A\Rightarrow xf'(x)\geq\frac{1}{2})$.

Soit $x$ un réel fixé supérieur ou égal à $A$. $\forall t\in[A,x],\;f'(t)\geq\frac{1}{2x}$ et donc, par croissance de l'intégrale, $\int_{A}^{x}f'(t)\;dt\geq\int_{A}^{x}\frac{1}{2t}\;dt$ ce qui fournit~:

$$\forall x\geq A,\;f(x)\geq f(A)+\frac{1}{2}(\ln x-\ln A),$$

et montre que $\lim_{x\rightarrow +\infty}f(x)=+\infty$.
\fincorrection
\correction{005423}
$$\forall x\in\Rr f(\frac{x}{2}+3)=f(f\circ f(x))=f\circ f(f(x))=\frac{f(x)}{2}+3.$$

Puisque $f$ est dérivable sur $\Rr$, on obtient en dérivant $\forall x\in\Rr,\;\frac{1}{2}f'(\frac{x}{2}+3)=\frac{1}{2}f'(x)$,

et donc

$$\forall x\in\Rr,\;f'(\frac{x}{2}+3)=f'(x).$$

Soit alors $x$ un réel donné et $u$ la suite définie par $u_0=x$ et $\forall n\in\Nn,\;u_{n+1}=\frac{1}{2}u_n+3$.

D'après ce qui précède, $\forall n\in\Nn,\;f'(x)=f'(u_n)$.
Maintenant, $u$ est une suite arithmético-géométrique et on sait que 

$$\forall n\in\Nn,\;u_n-6=\frac{1}{2^n}(u_0-6)$$ 

ce qui montre que la suite $u$ converge vers $6$. La suite $(f'(u_n))_{n\geq0}$ est constante, de valeur $f'(x)$. $f'$ étant continue sur $\Rr$, on en déduit que

$$\forall x\in\Rr,\;f'(x)=\lim_{n\rightarrow +\infty}f'(u_n)=f'(\lim_{n\rightarrow +\infty}u_n)=f'(6),$$

ce qui montre que la fonction $f'$ est constante sur $\Rr$ et donc que $f$ est affine.

Réciproquement, pour $x$ réel, posons $f(x)=ax+b$.
 
\begin{align*}\ensuremath
f\;\mbox{solution}&\Leftrightarrow\forall x\in\Rr,\;a(ax+b)+b=\frac{x}{2}+3\Leftrightarrow\forall x\in\Rr,\;(a^2-\frac{1}{2})x+ab+b-3=0\\
 &\Leftrightarrow a^2=\frac{1}{2}\;\mbox{et}\;(a+1)b=3\Leftrightarrow(a=\frac{1}{\sqrt{2}}\;\mbox{et}\;b=3(2-\sqrt{2}))\;\mbox{ou}\;(a=-\frac{1}{\sqrt{2}}\;\mbox{et}\;b=3(2+\sqrt{2})).
\end{align*}

On trouve deux fonctions solutions, les fonctions $f_1$ et $f_2$ définies par~:

$$\forall x\in\Rr,\;f_1(x)=\frac{1}{\sqrt{2}}x+3(2-\sqrt{2})\;\mbox{et}\;f_2(x)=-\frac{1}{\sqrt{2}}x+3(2+\sqrt{2}).$$

\fincorrection
\correction{005424}
Montrons que $\lim_{x\rightarrow +\infty}f(x)=0$.

Pour $x$ réel, posons $g(x)=e^xf(x)$. $g$ est dérivable sur $R$ et $\forall x\in\Rr,\;g'(x)=e^x(f(x)+f'(x))$. Il s'agit donc maintenant de montrer que si $\lim_{x\rightarrow +\infty}e^{-x}g'(x)=0$ alors $\lim_{x\rightarrow +\infty}e^{-x}g(x)=0$.

Soit $\varepsilon$ un réel strictement positif.

$$\exists A>0/\;\forall x\in\Rr,\;(x\geq A\Rightarrow-\frac{\varepsilon}{2}<e^{-x}g'(x)<\frac{\varepsilon}{2}\Rightarrow-\frac{\varepsilon}{2}e^x\leq g'(x)\leq \frac{\varepsilon}{2}e^x).$$
 
Pour $x$ réel donné supérieur ou égal à $A$, on obtient en intégrant sur $[A,x]$~:

$$-\frac{\varepsilon}{2}(e^x-e^A)=\int_{A}^{x}-\frac{\varepsilon}{2}e^t\;dt\leq\int_{A}^{x}g'(t)\;dt=g(x)-g(A)\leq
\frac{\varepsilon}{2}(e^x-e^A),$$

et donc 

$$\forall x\geq A,\;g(A)e^{-x}-\frac{\varepsilon}{2}(1-e^{A-x})\leq e^{-x}g(x)\leq g(A)e^{-x}+\frac{\varepsilon}{2}(1-e^{A-x}).$$

Maintenant, $g(A)e^{-x}-\frac{\varepsilon}{2}(1-e^{A-x})$ et $g(A)e^{-x}+\frac{\varepsilon}{2}(1-e^{A-x})$ tendent respectivement vers $-\frac{\varepsilon}{2}$ et $\frac{\varepsilon}{2}$ quand $x$ tend vers $+\infty$. Donc,

$$\exists B\geq A/\;\forall x\in\Rr,\;(x\geq B\Rightarrow g(A)e^{-x}-\frac{\varepsilon}{2}(1-e^{A-x})>-\varepsilon\;\mbox{et}\;<g(A)e^{-x}-\frac{\varepsilon}{2}(1-e^{A-x})
<\varepsilon.$$

Pour $x\geq B$, on a donc $-\varepsilon<e^{-x}g(x)<\varepsilon$.

On a montré que $\forall\varepsilon>0,\;\exists B>0/\;\forall x\in\Rr,\;(x\geq B\Rightarrow|e-xg(x)|<\varepsilon)$ et donc $\lim_{x\rightarrow +\infty}e^{-x}g(x)=0$ ce qu'il fallait démontrer.
\fincorrection
\correction{001267}
Habituellement on trouve le développement limité d'une fonction 
à partir des dérivées successives. Ici on va faire l'inverse.

Calcul du dl (à un certain ordre) :
\begin{align*}
f(x) & = \frac{x^3}{1+x^6} = x^3 \frac{1}{1+x^6}\\
     & =  x^3 \left( 1-x^6 + x^{12} - \cdots \pm x^{6\ell} \cdots \right) \\
     &= x^3 - x^9 + x^{15} -\cdots \pm x^{3+6\ell} \cdots \\
     &= \sum_{\ell \ge 0} (-1)^\ell x^{3+6\ell} \\
\end{align*}

Il s'agit d'identifier ce développement avec la formule de Taylor :
$$f(x)=f(0)+f'(0)x+\frac{f''(0)}{2!}x^2 + \cdots + \frac{f^{(n)}(0)}{n!} x^n + \cdots$$


Par unicité des DL, en identifiant les coefficients devant $x^n$ on trouve :
$$\begin{cases}
\frac{f^{(n)}(0)}{n!} = (-1)^\ell & \text{ si } n = 3+6\ell \\
\frac{f^{(n)}(0)}{n!} = 0         & \text{ sinon.} \\    
\end{cases}$$

Si $n=3+6\ell$ (avec $\ell\in \Nn$) alors on peut écrire $\ell = \frac{n-3}{6}$
et donc on peut conclure :
$$\begin{cases}
f^{(n)}(0) = (-1)^{\frac{n-3}{6}} \cdot n! & \text{ si } n \equiv 3 \pmod{6} \\
f^{(n)}(0) = 0         & \text{ sinon.} \\    
\end{cases}$$
\fincorrection
\correction{001268}
\begin{enumerate}
  \item La formule de Taylor-Lagrange à l'ordre $2$ entre $x$ et $x+h$ (avec $h>0$) donne :
$$f(x+h) = f(x) + f'(x) h  + f''(c_{x,h}) \frac{h^2}{2!} $$
où $c_{x,h} \in ]x,x+h[$.

Cela donne :
$$f'(x) h = f(x+h) - f(x)  - f''(c_{x,h}) \frac{h^2}{2!}.$$

On peut maintenant majorer $f'(x)$ :
\begin{align*}
h|f'(x)| 
   & \le \left|f(x+h) \right| + \left| f(x) \right| + \frac{h^2}{2}\left| f''(c_{x,h})\right|  \\
   & \le 2M_0 + \frac {h^2}{2} M_2 \\
\end{align*}
Donc
$$|f'(x)|\le \frac{2}{h}{M_0} + \frac{h}{2}M_2.$$

  \item Soit $\phi : ]0,+\infty[ \rightarrow \Rr$ la fonction définie par $\phi(h) = \frac {h}{2}M_2+\frac{2}{h}M_0$.
C'est une fonction continue et dérivable. La limite en $0$ et $+\infty$ est $+\infty$.
La dérivée $\phi'(h)=\frac12 M_2-\frac{2M_0}{h^2}$ s'annule en $h_0 = 2\sqrt{\frac{M_0}{M_2}}$ et en ce point 
$\phi$ atteint son minimum
$\phi(h_0) = 2\sqrt{M_0M_2}$.

Fixons $x>a$. Comme pour tout $h>0$ on a $|f'(x)| \le \frac {h}{2}M_2+\frac{2}{h}M_0 =\phi(h)$ alors en particulier pour 
$h=h_0$ on obtient $|f'(x)| \le \phi(h_0)= 2\sqrt{M_0M_2}$. Et donc $f'$ est bornée.

  \item Fixons $\epsilon >0$. $g''$ est bornée, notons $M_2 = \sup_{x> 0}\vert g''(x)\vert$. Comme $g(x)\to 0$ alors il 
existe $a>0$ tel que sur l'intervalle $]a,+\infty[$, $g$ soit aussi petit que l'on veut. Plus précisément nous choisissons $a$ de sorte que 
$$M_0 = \sup_{x>a}\vert g(x)\vert \le \frac{\epsilon^2}{4M_2}.$$


La première question appliquée à $g$ sur l'intervalle $]a,+\infty[$ implique 
que pour tout $h>0$ : 
$$|g'(x)| \le  \frac{2}{h}M_0 + \frac {h}{2} M_2 $$

En particulier pour $h = \frac{\epsilon}{M_2}$ et en utilisant $M_0  \le \frac{\epsilon^2}{4M_2}$
on obtient :
$$|g'(x)| \le  \frac{2}{\frac{\epsilon}{M_2}}\frac{\epsilon^2}{4M_2} + \frac{\frac{\epsilon}{M_2}}{2} M_2 \le \epsilon.$$

Ainsi pour chaque $\epsilon$ on a trouvé $a>0$ tel que si $x>a$ alors $|g'(x)|\le \epsilon$.
C'est exactement dire que $\lim_{x\to+\infty} g'(x)=0$.

\end{enumerate}
\fincorrection
\nocorrection
\nocorrection
\nocorrection
\nocorrection
\nocorrection
\correction{004004}
$f'(a)f'''(a) - f''(a)^2$.
\fincorrection
\nocorrection
\nocorrection
\correction{004007}
\begin{enumerate}
  \item Formule de Taylor Lagrange entre $\frac 12$ et $0$.
  \item Sinon, la fonction $g$ : $x \mapsto f(x) - (1-2x)^n$ est monotone
             sur $[0,\frac12]$ et nulle en $0$ et $\frac 12$, donc
             identiquement nulle. Impossible car $g^{(n)}(\frac12) \ne 0$.
\end{enumerate}
 \fincorrection
\correction{004008}
\begin{enumerate}
  \item Formule de Taylor pour $f$ et $f'$ $ \Rightarrow  \lambda = 1/180$.
  \item 
\end{enumerate}
 \fincorrection
\correction{004009}
\begin{enumerate}
  \item Formule de Taylor pour calculer $f(a)$ et $f(b)$ à partir de $f(x)$.
  \item \'Etudier $f(x) - f(a) - (x-a)f'(a) \pm M(x-a)^2/2$.
\end{enumerate}
 \fincorrection
\correction{004010}
Appliquer la formule de Taylor à l'ordre 2 de~$x$ à~$a$ et
de~$x$ à~$-a$.
\fincorrection
\correction{004011}
$f^{(n)}(a + h\theta_h) = f^{(n)}(a) + h\theta_hf^{(n+1)}(a+\theta'h)
                        = f^{(n)}(a) + \frac h{n+1}f^{(n+1)}(a+\theta''h)$.
\fincorrection
\nocorrection
\correction{004013}
\begin{enumerate}
  \item     $f(x+h) = f(x) + hf'(x) + \frac {h^2}2f''(x+\theta h)  \Rightarrow 
     f'(x) = \frac {f(x+h) - f(x)}h -\frac h2f''(x+\theta h)$.
  \item $h = 2\sqrt{\alpha/\beta}  \Rightarrow  |f'| \le 2\sqrt{\alpha\beta}$.
\end{enumerate}
 \fincorrection
\correction{004014}
\begin{enumerate}
  \item $= f(x+y) + \frac {y^2}2(M-f''(z))$.
  \item $\Delta \le 0$.
  \item $\sqrt f$ est affine.
\end{enumerate}
 \fincorrection
\nocorrection
\correction{004016}
Soit $\varepsilon > 0$ :
         $f(x+\varepsilon x) = f(x) + \varepsilon xf'(x) + \frac {\varepsilon^2 x^2}2 f''(x+\varepsilon\theta x)
          \Rightarrow 
         xf'(x) = \frac{ f(x+\varepsilon x) - f(x) }\varepsilon - \frac{\varepsilon x^2}2f''(x+\varepsilon\theta x)$.
\fincorrection
\correction{004017}
Si $Q = P\circ f$ alors $Q' = f'\times (P'\circ f)$ a autant de racines que $P'$
d'où $p = q$, $f(y_i) = x_i$ et $Q(y_i) = P(x_i)$.
De plus, au voisinage de~$y_i$~:
$$\lambda_i(y-y_i)^{n_i} \sim Q'(y) = f'(y)\times P'(f(y))
  \sim f'(y_i)\times \mu_i(f(y)-x_i)^{m_i} \sim \mu_if'(y_i)^{1+m_i}(y-y_i)^{m_i}$$
d'où $m_i = n_i$.

Réciproquement, si $p=q$, $P(x_i) = Q(y_i)$ et $m_i = n_i$ alors en posant
$x_0 = y_0 = -\infty$ et $x_{p+1} = y_{p+1} = +\infty$, $P$ induit un
$\mathcal{C}^1$-difféomorphisme de $]x_{i},x_{i+1}[$ sur $P(]x_{i},x_{i+1}[) = Q(]x_{i},x_{i+1}[)$
(les limites de $P$ et $Q$ en $+\infty$ sont égales à $+\infty$ vu les coefficients
dominants de $P$ et $Q$~; celles en $-\infty$ s'en déduisent en comptant les
changements de signe pour $P'$ ou pour $Q'$ et on trouve le même compte puisque $m_i = n_i$).
On note $f_i$ la fonction réciproque de $P_{|]x_{i},x_{i+1}[}$ et
$f$ définie par $f(y) = f_i(Q(y))$ si $y_i < y < y_{i+1}$ et $f(y_i) = x_i$.
$f$ ainsi définie est strictement croissante, de classe $\mathcal{C}^1$ à dérivée
non nulle sauf peut-être aux $y_i$, et $P\circ f = Q$. Reste à étudier le
caractère $\mathcal{C}^1$ en $y_i$ et à vérifier que $f'(y_i) \ne 0$.

Au voisinage de~$y_i$, par intégration des DL de $P$ et $Q$ on a~:
$$\frac{\lambda_i}{1+m_i}(y-y_i)^{1+m_i}\sim Q(y) - Q(y_i) = P(f(y)) - P(f(y_i)) \sim \frac{\mu_i}{1+m_i}(f(y) - f(y_i))^{1+m_i}$$
d'où $\frac{f(y)-f(y_i)}{y-y_i} \to \Bigl(\frac{\lambda_i}{\mu_i}\Bigr)^{1/(1+m_i)}$ lorsque $y\to y_i$, car les taux
d'accroissement de~$f$ sont positifs. Ceci prouve que $f$ est dérivable en~$y_i$
et $f'(y_i) = \Bigl(\frac{\lambda_i}{\mu_i}\Bigr)^{1/(1+m_i)} \ne 0$. Enfin on a, lorsque $y\to y_i$ :
$$f'(y) = \frac{Q'(y)}{P'(f(y))} \sim \frac{\lambda_i(y-y_i)^{m_i}}{\mu_i(f(y)-f(y_i))^{m_i}}\to  \frac{\lambda_i}{\mu_i}\Bigl(\frac{\lambda_i}{\mu_i}\Bigr)^{-m_i/(1+m_i)} = \Bigl(\frac{\lambda_i}{\mu_i}\Bigr)^{1/(1+m_i)}$$
et donc $f$ est $\mathcal{C}^1$ en $y_i$.
\fincorrection
\correction{005418}
Soit $(x,y)\in\Rr\times\Rr$. Puisque $f$ est de classe $C^3$ sur $\Rr$, la formule de \textsc{Taylor}-\textsc{Laplace} à l'ordre $2$ permet d'écrire

$$
\begin{array}{l}
f(x+y)=f(x)+yf'(x)+\frac{y^2}{2}f''(x)+\int_{x}^{x+y}\frac{(x+y-t)^2}{2}f^{(3)}(t)\;dt\;\mbox{et}\\
f(x-y)=f(x)-yf'(x)+\frac{y^2}{2}f''(x)\int_{x}^{x-y}\frac{(x-y-t)^2}{2}f^{(3)}(t)\;dt.
\end{array}$$

Donc,

\begin{align*}\ensuremath
(f(x)^2)&\geq f(x+y)f(x-y)\\
 &=(f(x)+yf'(x)+\frac{y^2}{2}f''(x)+\int_{x}^{x+y}\frac{(x+y-t)^2}{2}f^{(3)}(t)\;dt)\times\\
 &\quad(f(x)-yf'(x)+\frac{y^2}{2}f''(x)+\int_{x}^{x-y}\frac{(x-y-t)^2}{2}f^{(3)}(t)\;dt)\\
 &=(f(x))^2+y^2(f(x)f''(x)-(f'(x))^2)\\
 &\quad+(f(x)-yf'(x)+\frac{y^2}{2}f''(x))\int_{x}^{x+y}\frac{(x+y-t)^2}{2}f^{(3)}(t)\;dt\\
 &\quad+(f(x)+yf'(x)+\frac{y^2}{2}f''(x))\int_{x}^{x-y}\frac{(x-y-t)^2}{2}f^{(3)}(t)\;dt\;(*)
\end{align*}

Maintenant, pour $y\in[-1,1]$, ($f^{(3)}$ étant continue sur $\Rr$ et donc continue sur le segment $[-1,1]$),

$$\left|\int_{x}^{x+y}\frac{(x+y-t)^2}{2}f^{(3)}(t)\;dt\right|\leq|y|.\frac{y^2}{2}\mbox{Max}\{|f^{(3)}(t)|,\;t\in[x-1,x+1]\},$$

et donc,

$$\frac{1}{y^2}\left|\int_{x}^{x+y}\frac{(x+y-t)^2}{2}f^{(3)}(t)\;dt\right|\leq|y|\mbox{Max}\{|f^{(3)}(t)|,\;t\in[x-1,x+1]\}.$$

Cette dernière expression tend vers $0$ quand $y$ tend vers $0$. On en déduit que $\frac{1}{y^2}\left|\int_{x}^{x+y}\frac{(x+y-t)^2}{2}f^{(3)}(t)\;dt\right|$ tend vers $0$ quand $y$ tend vers $0$. De même, $\frac{1}{y^2}\left|\int_{x}^{x-y}\frac{(x-y-t)^2}{2}f^{(3)}(t)\;dt\right|$ tend vers $0$ quand $y$ tend vers $0$.

On simplifie alors $(f(x)^2$ dans les deux membres de $(*)$. On divise les deux nouveaux membres par $y^2$ pour $y\neq 0$ puis on fait tendre $y$ vers $0$ à $x$ fixé. On obtient $0\geq f(x)f''(x)-(f'(x))^2$, qui est l'inégalité demandée.
\fincorrection
\correction{005457}
Soit $x\in[0,1]\subset[0,\frac{\pi}{2}]$. 

D'après la formule de \textsc{Taylor}-\textsc{Laplace} à l'ordre 1 fournit
 
$$\sin x=x-\int_{0}^{x}(x-t)\sin t\;dt\leq x,$$

car pour $t\in[0,x]$, $(x-t)\geq0$ et pour $t\in[0,x]\subset[0,\frac{\pi}{2}]$, $\sin t\geq0$. 

De même, la formule de \textsc{Taylor}-\textsc{Laplace} à l'ordre 3 fournit
 
$$\sin x=x-\frac{x^3}{6}+\int_{0}^{x}\frac{(x-t)^3}{6}\sin t\;dt\geq x-\frac{x^3}{6}.$$

Donc, $\forall x\in[0,1],\;x-\frac{x^3}{6}\leq\sin x\leq x$.

Soient alors $n\geq1$ et $k\in\{1,...,n\}$. On a $0\leq\frac{1}{(n+k)^2}\leq1$ et donc 

$$\frac{1}{(n+k)^2}-\frac{1}{6(n+k)^6}\leq\sin\frac{1}{(n+k)^2}\leq\frac{1}{(n+k)^2},$$

puis en sommant 

$$\sum_{k=1}^{n}\frac{1}{(n+k)^2}-\sum_{k=1}^{n}\frac{1}{6(n+k)^6}\leq\sum_{k=1}^{n}\sin\frac{1}{(n+k)^2}\leq\sum_{k=1}^{n}\frac{1}{(n+k)^2}.$$

Maintenant, quand $n$ tend vers $+\infty$,

$$\sum_{k=1}^{n}\frac{1}{(n+k)^2}=\frac{1}{n}.\frac{1}{n}\sum_{k=1}^{n}\frac{1}{(1+\frac{k}{n})^2}
=\frac{1}{n}\left(\int_{0}^{1}\frac{1}{(1+x)^2}\;dx+o(1)\right)=\frac{1}{2n}+o(\frac{1}{n}).$$

D'autre part,

$$0\leq\sum_{k=1}^{n}\frac{1}{6(n+k)^6}\leq n.\frac{1}{6n^6}=\frac{1}{6n^5},$$

et donc, $\sum_{k=1}^{n}\frac{1}{6(n+k)^6}=o(\frac{1}{n})$.

On en déduit que $2n(\frac{1}{(n+k)^2}-\frac{1}{6(n+k)^6})=2n(\frac{1}{2n}+o(\frac{1}{n}))=1+o(1)$ et que $2n\frac{1}{(n+k)^2}=1+o(1)$. Mais alors, d'après le théorème des gendarmes, $2n\sum_{k=1}^{n}\sin\frac{1}{(n+k)^2}$ tend vers $1$ quand $n$ tend vers $+\infty$, ou encore

$$\sum_{k=1}^{n}\sin\frac{1}{(n+k)^2}\underset{n\rightarrow+\infty}{\sim}\frac{1}{2n}.$$
\fincorrection
\correction{001237}
\begin{enumerate}
  \item $\displaystyle{\ln(\cos x) = -{\frac {1}{2}}{x}^{2}-{\frac {1}{12}}{x}^{4}-{\frac {1}{45}}{x}^{6}+
o\left ({x}^{6}\right ) }$.
  \item $\displaystyle{\tan x = x+{\frac {1}{3}}{x}^{3}+{\frac {2}{15}}{x}^{5}+{\frac {17}{315}}{x}^{
7}+o\left ({x}^{7}\right )}$.
  \item $\displaystyle{\sin(\tan x) = x+{\frac {1}{6}}{x}^{3}-{\frac {1}{40}}{x}^{5}-{\frac {55}{1008}}{x}^
{7}+o\left ({x}^{7}\right ) }$.
  \item $\displaystyle{\left( \ln (1+x) \right)^2={x}^{2}-{x}^{3}+{\frac {11}{12}}{x}^{4}+o\left ({x}^{4}\right ) }$.
  \item $\displaystyle{\exp(\sin x) = 1+x+{\frac {1}{2}}{x}^{2}+o\left ({x}^{3}\right ) }$.
  \item $\displaystyle{\sin^6 x = {x}^{6}+o\left ({x}^{6}\right ) }$.
\end{enumerate}
\fincorrection
\nocorrection
\correction{001239}
$$\lim _{x\rightarrow 0}{\frac {\arctan(x)-\sin(x)}{\tan(x)-\arcsin(x)}}
=-1.$$
\fincorrection
\correction{001240}
\
\begin{enumerate}
  \item $\displaystyle{\ln \cos x = -{\frac {1}{2}}{x}^{2}-{\frac {1}{12}}{x}^{4}-{\frac {1}{45}}{x}^{6}+
o\left ({x}^{7}\right ) }$.
\item $\displaystyle{{\frac {\arctan(x)-x}{\sin(x)-x}} =
2-{\frac {11}{10}}{x}^{2}+o\left ({x}^{3}\right ) }$.
\item $\displaystyle{\ln (\tan(1/2\,x+1/4\,\pi )) = x+{\frac {1}{6}}{x}^{3}+o\left ({x}^{4}\right ) }$.
\item $\displaystyle{\ln \sin x = \ln (1/2\,\sqrt {2})+x-\frac \pi 4 -\left (x-\frac \pi 4 \right )^{2}+{\frac {2}{3}}\left (x- \frac \pi 4\right )^{3}+o\left (\left (x-\frac \pi 4 \right )^{3}\right ) }$.
\item $\displaystyle{\sqrt [3]{{x}^{3}+x}-\sqrt [3]{{x}^{3}-x}=
2/3\,\frac 1{x}+o(\frac 1 {x^4}) }$.
\item $\displaystyle{(1+x)^{\frac 1 x}={e^{{\frac {\ln (1+x)}{x}}}}=
{e}}-1/2\,{e}x+{\frac {11}{24}}\,{e}{x}^{2}-{\frac {7}{16}
}\,{e}{x}^{3}+o\left ({x}^{3}\right ) $.
\item $\displaystyle{x\left (\sqrt {{x}^{2}+\sqrt {{x}^{4}+1}}-
x\sqrt {2}\right ) = 1/8\,{\frac {\sqrt {2}}{{x}^{2}}}+o({x}^{-5})
}$.
\end{enumerate}
\fincorrection
\nocorrection
\nocorrection
\correction{001243}
\begin{enumerate}
  \item 
Première méthode.
On applique la formule de Taylor (autour du point $x=1$)
$$f(x)=f(1)+f'(1)(x-1)+\frac{f''(1)}{2!}(x-1)^2 + \frac{f'''(1)}{3!}(x-1)^3 + o((x-1)^3)$$
Comme $f(x) = \sqrt x= x^{\frac12}$ alors $f'(x) = \frac12 x^{-\frac12}$ et donc $f'(1)=\frac12$.
Ensuite on calcule $f''(x)$ (puis $f''(1)$), $f'''(x)$ (et enfin $f'''(1)$).

On trouve le dl de $f(x)=\sqrt x$ au voisinage de $x=1$ :
$$\sqrt x = 1 + \frac12 (x-1) - \frac18 (x-1)^2 + \frac{1}{16} (x-1)^3 + o((x-1)^3)$$


\bigskip

Deuxième méthode.
Posons $h=x-1$ (et donc $x=h+1$). On applique la formule du dl de $\sqrt{1+h}$ autour de $h=0$.
\begin{align*}
f(x)=\sqrt x  
  & = \sqrt{1+h} \\
  & = 1 + \frac 12 h - \frac18 h^2 + \frac{1}{16} h^3 + o(h^3)   \quad \text{ c'est la formule du dl de } \sqrt{1+h} \\
  & = 1 + \frac12 (x-1) - \frac18 (x-1)^2 + \frac{1}{16} (x-1)^3 + o((x-1)^3) \\
\end{align*}


  \item La première méthode consiste à calculer $g'(x)=\frac{1}{2\sqrt x}\exp{\sqrt x}$, $g''(x)$, $g'''(x)$
 puis $g(1)$, $g'(1)$, $g''(1)$, $g'''(1)$ pour pouvoir appliquer la formule de Taylor conduisant à :
$$\exp(\sqrt x)= e + \frac{e}{2} (x-1) + \frac{e}{48} (x-1)^3 + o((x-1)^3)$$
(avec $e=\exp(1)$).

\bigskip

Autre méthode. Commencer par calculer le dl de $k(x)=\exp x$ en $x=1$ ce qui est très facile car pour tout $n$, $k^{(n)}(x)=\exp x$
et donc $k^{(n)}(1)=e$ :
$$\exp x= e + e (x-1) + \frac{e}{2!} (x-1)^2 + \frac{e}{3!} (x-1)^3 + o((x-1)^3).$$

Pour obtenir le dl $g(x)= h(\sqrt x)$ en $x=1$ on écrit d'abord :
$$\exp (\sqrt x)= e + e (\sqrt{x}-1) + \frac{e}{2!} (\sqrt{x}-1)^2 + \frac{e}{3!} (\sqrt{x}-1)^3 + o((\sqrt{x}-1)^3).$$
Il reste alors à substituer $\sqrt{x}$ par son dl obtenu dans la première question.

  \item 

Posons $u=x-\frac\pi3$ (et donc $x=\frac\pi3+u$).
Alors 
$$\sin(x)=\sin(\frac\pi3+u) = \sin(\frac\pi3)\cos(u)+\sin(u)\cos(\frac\pi3) = \frac{\sqrt3}{2}\cos u +\frac12\sin u$$
On connaît les dl de $\sin u$ et $\cos u$ autour de $u=0$ (car on cherche un dl autour de $x=\frac\pi3$) donc

\begin{align*}
\sin x 
  & = \frac{\sqrt3}{2}\cos u +\frac12\sin u \\
  & = \frac{\sqrt3}{2} \bigg(1-\frac{1}{2!}u^2 + o(u^3) \bigg) + \frac{1}{2} \bigg( u -\frac{1}{3!}u^3 + o(u^3) \bigg) \\
  & = \frac{\sqrt3}{2} + \frac{1}{2} u -  \frac{\sqrt3}{4} u^2 -\frac{1}{12}u^3 + o(u^3) \\
  & = \frac{\sqrt3}{2} + \frac{1}{2} (x-\frac\pi3) -  \frac{\sqrt3}{4}(x-\frac\pi3)^2 -\frac{1}{12}(x-\frac\pi3)^3 + o((x-\frac\pi3)^3) \\
\end{align*}


\bigskip

Maintenant pour le dl de la forme $\ln(a+v)$ en $v=0$ on se ramène au dl de $\ln(1+v)$ ainsi :
$$\ln(a+v)=\ln\big(a(1 + \frac v a)\big) = \ln a + \ln(1 + \frac v a)
= \ln a + \frac v a - \frac12 \frac{v^2}{a^2}  + \frac13 \frac{v^3}{a^3} + o(v^3)$$

On applique ceci à $h(x)= \ln(\sin x)$ en posant toujours $u=x-\frac\pi3$ :
\begin{align*}
h(x)= \ln(\sin x)
  & = \ln\left(\frac{\sqrt3}{2} + \frac{1}{2} u -  \frac{\sqrt3}{4} u^2 -\frac{1}{12}u^3 + o(u^3)\right) \\
  & = \ln\left(\frac{\sqrt3}{2}\right) + \ln\left(1 + \frac{2}{\sqrt3}\left(\frac{1}{2} u -  \frac{\sqrt3}{4} u^2 -\frac{1}{12}u^3 + o(u^3)\right) \right) \\
  & = \ \cdots \qquad \text{ on effectue le dl du } \ln \text{ et on regroupe les termes} \\
  & = \ln\left(\frac{\sqrt3}{2}\right) + \frac{1}{\sqrt 3} u -\frac23 u^2 + \frac{4}{9\sqrt3}u^3 + o(u^3) \\
  & = \ln\left(\frac{\sqrt3}{2}\right) + \frac{1}{\sqrt 3}(x-\frac\pi3)  -\frac23 (x-\frac\pi3)^2 + \frac{4}{9\sqrt3}(x-\frac\pi3)^3 + o((x-\frac\pi3)^3) \\
\end{align*}

On trouve donc :
$$\ln(\sin x) = \ln\left(\frac{\sqrt3}{2}\right) + \frac{1}{\sqrt 3}(x-\frac\pi3) 
 -\frac23 (x-\frac\pi3)^2 + \frac{4}{9\sqrt3}(x-\frac\pi3)^3 + o((x-\frac\pi3)^3)$$


\bigskip Bien sûr une autre méthode consiste à calculer $h(1)$, $h'(1)$, $h''(1)$ et $h'''(1)$.

\end{enumerate}
\fincorrection
\correction{001244}
\begin{enumerate}
  \item Dl de $f(x)$ à l'ordre $2$ en $0$.
\begin{align*}
f(x) 
  & =   \frac{\sqrt{1+x^2}}{1+x+\sqrt{1+x^2}} \\
  & = \frac{1+\frac{x^2}{2} + o(x^2)}{1 + x + 1+\frac{x^2}{2} + o(x^2)} 
\quad \text{ car } \sqrt{1+x^2} = 1 + \frac12 x^2+o(x^2) \\  
  & = \big(1+\frac{x^2}{2} + o(x^2) \big) \times \frac 12 \frac{1}{1+\frac{x}{2}+\frac{x^2}{4} + o(x^4)} 
\quad \text{ on pose } u=\frac{x}{2}+\frac{x^2}{4} + o(x^4)  \\
  & = \frac 12 \big(1+\frac{x^2}{2} + o(x^2) \big) \times \frac{1}{1+u} \\
  & = \frac 12 \big(1+\frac{x^2}{2} + o(x^2) \big) \times \big(1-u+u^2+o(u^2)\big)  \\
  & = \frac 12 \big(1+\frac{x^2}{2} + o(x^2) \big) \times \big(1- \big(\frac{x}{2}+\frac{x^2}{4}\big) + \big(\frac{x}{2}+\frac{x^2}{4}\big)^2+o(x^2)\big)  \\
  & = \frac 12 \big(1+\frac{x^2}{2} + o(x^2) \big) \times \big(1- \frac{x}{2}+o(x^2)\big)  \\
  & = \frac 12 \big(1- \frac{x}{2}+\frac{x^2}{2} + o(x^2) \big)  \\
  & = \frac 12 - \frac{x}{4}+\frac{x^2}{4} + o(x^2)   \\ 
\end{align*}

  \item En $+\infty$ on va poser $h=\frac1x$ et se ramener à un dl en $h=0$.

$$f(x)= \frac{\sqrt{1+x^2}}{1+x+\sqrt{1+x^2}} =  \frac{x\sqrt{\frac{1}{x^2}+1}}{x\big(\frac1x+1+\sqrt{\frac{1}{x^2}+1}\big)}
=  \frac{\sqrt{1+h^2}}{1+h+\sqrt{1+h^2}} = f(h).$$

Ici -miraculeusement- on retrouve exactement l'expression de $f$ dont on a déjà calculé le dl en $h=0$ :
$f(h) =  \frac 12 - \frac{h}{4}+\frac{h^2}{4} + o(h^2)$. Ainsi 
$$f(x) = f(h) = \frac 12 - \frac{1}{4x}+\frac{1}{4x^2} + o(\frac1{x^2})$$

  \item Attention cela ne fonctionne plus du tout en $-\infty$.
Dans le calcul de la deuxième question on était on voisinage de $+\infty$ et nous avons considéré que $x$ était positif.
En $-\infty$ il faut faire attention au signe, par exemple $\sqrt{1+x^2}= |x|\sqrt{\frac{1}{x^2}+1} = -x \sqrt{\frac{1}{x^2}+1}$.

Ainsi toujours en posant $h=\frac1x$.
\begin{align*}
f(x) 
  & = \frac{\sqrt{1+x^2}}{x+1+\sqrt{1+x^2}} \\
  & = \frac{-x\sqrt{\frac{1}{x^2}+1}}{x\big(1+\frac1x-\sqrt{\frac{1}{x^2}+1}\big)} \\
  & = -\frac{\sqrt{1+h^2}}{1+h-\sqrt{1+h^2}} \\
  & = -\frac{1+\frac12 h^2 + o(h^2)}{1+h - \big(1+\frac12 h^2 + o(h^2) \big)} \\
  & = - \frac{1+\frac12 h^2 + o(h^2)}{h  - \frac12 h^2 + o(h^2)} \\
  & = -\frac1h \frac{1+\frac12 h^2 + o(h^2)}{1 - \frac12 h + o(h)} \\
  & = -\frac1h \big(1+\frac12 h^2 + o(h^2)\big)\times\big(1+\frac12h +\frac14 h^2 + o(h^2) \big) \\
  & = -\frac1h\big(1 +\frac12h +\frac34 h^2 + o(h^2) \big) \\
  & =  -\frac1h -\frac12 -\frac34 h + o(h) \\
  & = -x -\frac12 -\frac34 \frac1x + o(\frac1x) \\
\end{align*}  

Ainsi un développement (asymptotique) de $f$ en $-\infty$ est 
$$f(x) = -x -\frac12 -\frac34 \frac1x + o(\frac1x)$$
On en déduit par exemple que $f(x)$ se comporte essentiellement comme la fonction $-x$
en $-\infty$ et en particulier $\lim_{x\to -\infty} f = +\infty$.
\end{enumerate}

\fincorrection
\nocorrection
\nocorrection
\correction{002657}

$$\ln(x+1) = \ln \Big(x \times (1+\frac1x)\Big) = \ln x+\ln\big(1+{\frac1x}\big) = \ln x + \frac 1x + o(\frac1x)$$
Donc 
$$\frac{\ln(x+1)}{\ln x} = 1 + \frac1{x\ln x} + o(\frac1{x\ln x}).$$
Ainsi
\begin{align*}
\left(\frac{\ln(x+1)}{\ln x}\right)^x 
  & = \exp\left( x \ln \left( \frac{\ln(x+1)}{\ln x}\right) \right)  \\
  & = \exp\left( x \ln \left(1 + \frac1{x\ln x} + o(\frac1{x\ln x})\right)  \right)  \\
  & = \exp\left( x \left(\frac1{x\ln x} + o(\frac1{x\ln x})\right) \right)  \\
  & = \exp\left( \frac1{\ln x} + o(\frac1{\ln x})\right)  \\
  & = 1 + \frac1{\ln x} + o(\frac1{\ln x}) \\
\end{align*}
On en déduit immédiatement que 
$$\lim_{x\to+\infty}\left(\frac{\ln(x+1)}{\ln x}\right)^x = 1$$
et que lorsque $x\to + \infty$
$$\left(\frac{\ln(x+1)}{\ln x}\right)^x - 1 \sim \frac1{\ln x}.$$
\fincorrection
\correction{002683}
\vfill
{\leftskip=-1cm A\qquad  1.\par} 
$$\begin{aligned}
\qquad &{d\over dt}\sum_{m=1}^n (-1)^m (z-a)^m
                \phi^{(n-m)}(t) f^{(m)}\left(a+t(z-a)\right) \cr
        &=\sum_{m=1}^n (-1)^m (z-a)^m \Big[
                \phi^{(n-m+1)}(t)f^{(m)}\left(a+t(z-a)\right)  \cr
        &  \mkern260mu   +\phi^{(n-m)}(t) (z-a) f^{(m+1)}\left(a+t(z-a)\right)  \Big]  \cr
        &=\sum_{m=1}^n (-1)^m 
                (z-a)^m \phi^{(n-m+1)}(t)f^{(m)}\left(a+t(z-a)\right)   \cr
        &  \mkern220mu  + \sum_{m=1}^n (-1)^m(z-a)^{m+1} \phi^{(n-m)}(t)f^{(m+1)}\left(a+t(z-a)\right)   \cr
\end{aligned}
$$
Effectuons le changement d'indice de sommation $m=m'+1$ dans la deuxième somme; tous les
termes s'éliminent deux à deux, à l'exception du premier terme de la première somme
et du dernier terme de la deuxième somme, d'où le résultat demandé. 

\medskip
2.a. Plus généralement, on a le résultat suivant: si une fonction $f$ est nulle
en zéro et de classe $C^{n+1}$ sur un intervalle $I$ contenant zéro, $g(t)=f(t)/t$ est
prolongeable par continuité en zéro et son prolongement est de classe $C^n$.
Gardons nous de déduire fallacieusement ce résultat de l'existence d'un développement
limité d'ordre $n$ de $g$.
On montre à l'aide d'un développement limité d'ordre 1 de $f$ que $g$ est
prolongeable par continuité en $0$ en posant $\tilde g(0)=f'(0)$. Par ailleurs, $g$ est de
classe $C^{n+1}$ sur $I\backslash\{0\}$. 
Faisons l'hypothèse de
récurrence $g^{(k-1)}(t)={f^{(k)}(0) \over k} + {f^{(k+1)}(0) \over k+1} {t \over 1!}
+ \dots + {f^{(n+1)}(0) \over n+1} {t^{n-k+1} \over (n-k+1)!}+o(t^{n-k+1})$
$(1\leq k \leq n,\ t\neq0)$.
En dérivant $k$ fois l'identité $f(t)=tg(t)$, on obtient que
$\forall t\neq0 \ \ g^{(k)}(t)={ f^{(k)}(t) - k g^{(k-1)}(t) \over t }$.
Or $f^{(k)}(t)=f^{(k)}(0)+f^{(k+1)}(0){t\over1!}+\dots
+f^{(n+1)}(0){t^{n+1-k}\over(n-k)!}+o(t^{n+1-k})$.
En substituant ces développements limités dans l'identité précédente, on montre que
l'hypothèse de récurrence est vérifiée au rang suivant, donc pour tout entier
$k$ de 1 à $n+1$.
 Faisons l'hypothèse de récurrence que $g^{(k)}(0)$ existe
et est égale à $f^{(k+1)}(0)/(k+1)$ ($0\leq k\leq n$). Du développement limité
de $g^{(k)}$ (tronqué à l'ordre 1) et de l'hypothèse de
récurrence, il résulte que $g^{(k)}$ est continue et dérivable en zéro et que, si
$k<n$, $g^{(k+1)}(0)=f^{(k+2)}(0)/(k+2)$, ce qui prouve par récurrence que $g$ est
$n$ fois continument dérivable sur $I$.

Par conséquent, $t\mapsto{e^t-1 \over t}$ est prolongeable par continuité en une
fonction indéfiniment dérivable sur $\Rr$. Cette fonction ne s'annulant jamais,
son inverse est également définie et indéfiniment dérivable sur $\Rr$.



La fonction $t \mapsto {t \over e^t-1} + {t\over2}$ est paire car
$$\begin{aligned}
{-t \over e^{-t}-1} + {-t\over2}
&= {-t e^t \over 1-e^t} - {t\over2}  \cr
&= {t e^t - t + t \over e^t-1} - {t\over2}  \cr
&= {t \over e^t-1} + {t\over2}  \cr
\end{aligned}
$$
Donc  le développement limité de ${t \over e^t-1}$, dont l'existence est garantie par
le fait que la fonction est indéfiniment dérivable, est de la forme demandée par
l'énoncé.

Par conséquent
\begin{eqnarray*}
{t \over e^t-1} (e^{zt}-1) &=&
\Big( 1-{t\over2}+ {b_1 t^2 \over 2!} + {b_2 t^4 \over 4!} + \dots 
+  {b_N t^{2N} \over (2N)!} + o(t^n) \Big) 
\\
\times \Big( zt + {z^2 t^2 \over 2!} + \dots + {z^n t^n \over n!} + o(t^n)   \Big)   
 \end{eqnarray*}
et $\phi_n(z)/n!$ est le coefficient de $t^n$ dans ce développement:
$$\phi_n(z)/n!= {z^n\over n!} - {1\over2} {z^{n-1}\over(n-1)!}
+ {b_1\over2!}{z^{n-2}\over(n-2)!} + {b_2\over4!}{z^{n-4}\over(n-4)!}
+ \dots + {b_N\over(2N)!}{z^{n-2N}\over(n-2N)!},$$
d'où l'expression de $\phi_n$ demandée.

\medskip
2.b. 
$$\begin{aligned}
\phi_n(z+1)-\phi_n(z)&={d^n \over dt^n}_{| t=0}
                      \left(t{e^{zt}-1 \over e^t-1} - t{e^{(z+1)t}-1 \over e^t-1} \right) \cr
                     &={d^n \over dt^n}_{| t=0} \left(t e^{zt} \right) \cr
\end{aligned}
$$
Comme $t \mapsto t e^{zt}$ est de classe $C^\infty$ et que son développement limité
d'ordre $n$ en $t=0$ est $\sum_{k=0}^{n-1} {z^k\over k!} t^{k+1} + o(t^n)$, il vient
$\phi_n(z+1)-\phi_n(z)=nz^{n-1}$.


\medskip
3. $(i)$ est obtenue en dérivant autant de fois que nécessaire l'identité précédente
et en donnant à $z$ la valeur zéro.
   $(ii)$, $(iii)$, $(iv)$, $(v)$ et $(vi)$ sont des conséquences immédiates de 2.a.

4.a. On applique la question 1 au polynôme $\phi_{2n}$ de degré 2n et on
intègre entre $0$ et $1$. Il vient
\[
\sum_{m=1}^{2n} (-1)^m (z-a)^m \Big[ \phi_{2n}^{(2n-m)}(1) f^{(m)}(z)
- \phi_{2n}^{(2n-m)}(0) f^{(m)}(a)  \Big]
\]\[
 = -(2n)! \bigl( f(z)-f(a) \bigr) + (z-a)^{2n+1} \int_0^1 \phi_{2n}(t)f\Big( a+(z-a)t \Big)dt
\]
en tenant compte du fait que $\phi_{2n}^{(2n)}=(2n)!$

On obtient l'égalité demandée en substituant aux  dérivées itérées
de $\phi_{2n}$ les expressions déterminées dans la question 3.      

\medskip
4.b. Appliquons la question précédente en remplaçant $f$ par une
primitive de $F$ et $z$ par $\omega$. Il vient
\begin{eqnarray*}
0 &=& \int_a^{a+\omega} F(t)dt - {\omega\over2} (F(a+\omega)+F(a)) 
+ \sum_{m=1}^{n-1} b_m {(z-a)^{2m} \over (2m)!} \left[ F^{(2m-1)}(a+\omega) - F^{(2m-1)}(a) \right] 
\\
&& - {\omega^{2n+1}\over (2n)!} \int_0^1 \phi_{2n}(t)F^{(2n)}\Big( a+(z-a)t \Big)dt  
 \end{eqnarray*}
Lorsqu'on somme les égalités obtenues en remplaçant $a$ successivement
par lui-m{\^e}me, $a+\omega$, \dots, $a+(r-1)\omega$, on obtient le résultat demandé,
certains termes se simplifiant deux à deux.

\bigskip
{\leftskip=-1cm B\qquad 1. On a pour tout $x>0$ fixé \par}
$$u_k(x)=\ln \left( 1+{x\over k} \right)
                        + x \ln \left( 1-{k\over k+1} \right)
                        ={x\over k}-{x\over k+1}+O({1\over k^2})
                        =O({1\over k^2})$$
Donc la série $\sum_{k\geq1}u_k(x)$ converge.

\medskip
2. 
$$\begin{aligned}
\ln(x+1)+\sum_{k=1}^n &u_k(x+1)  \cr
        &=\ln(x+1) + \sum_{k=2}^{n+1} \ln(x+k)
                        + \sum_{k=1}^n \Big[ -\ln(k)
                                + (x+1) \big(\ln(k)-\ln(k+1)\big) \Big] \cr
        &=\sum_{k=1}^{n+1} \ln(x+k)                                
                        + \sum_{k=1}^n  x (\ln(k)-\ln(k+1)) 
                        - \sum_{k=1}^n \ln(k+1)        \cr
        &=\sum_{k=1}^n u_k(x) + \ln(x+n+1) - \ln(n+1)  \cr
\end{aligned}
$$
Lorsque $n$ tend vers $+\infty$,
$\ln(x+n+1) - \ln(n+1)=\ln \left( {x+n+1 \over n+1} \right)$ tend vers zéro;
on obtient donc par passage à la limite l'égalité souhaitée.

\medskip
3. Pour tout $k\geq1$ entier, $u_k(1)=0$, donc $G(1)=0$ et
on prouve aisément par récurrence à l'aide de la question précédente
que pour tout entier strictement positif $n$, $G(n+1)=\sum_{k=1}^n \ln(k)$,
égalité de laquelle on déduit immédiatement le résultat demandé.

\medskip
4. immédiat

\medskip
5. C'est une application directe de la question A.4.b.\par
$T_{p,n}(x,y)=-{1\over(2p)!}\int_0^1 \phi_{2p}(t)
                        \sum_{m=0}^{n-1} f^{(2p)}(m+t) dt
=-{1\over(2p)!} \int_0^n \phi_{(2p)} \big( t-E(t) \big) f^{(2p)}(t)dt$


\medskip
6. L'intégrande dans l'expression de $T_{p,n}(x,y)$ est majorée en valeur
absolue par le produit de la borne supérieure de la fonction continue
$\phi_{2p}$ sur le segment $[0,1]$ et de la valeur absolue de $f^{(2p)}$.

On prouve aisément par récurrence que
$f^{(m)}(t)=(-1)^{m-1}\left({1\over(y+t)^m}-{1\over(x+t)^m}\right)
= O \left( {1\over t^{m+1}} \right)$ (quand $t \to +\infty$).
Donc l'intégrale $\int_0^{+\infty} \phi_{(2p)} \big( t-E(t) \big) f^{(2p)}(t)dt$
est absolument convergente, ce qui prouve que $T_{p,n}(x,y)$ admet une limite
finie lorsque $n$ tend vers $+\infty$.



\medskip
7. D'après les questions 4 et 5,
$$\begin{aligned}
G(y)-G(x)&=\sum_{k=1}^{+\infty}\bigl[ \ln(y+k) - \ln(x+k)
                        + (y-x)\ln\left({k\over k+1}\right) \bigr] + \ln y - \ln x \cr
        &=\lim_{n\to +\infty}\Bigl\{\sum_{k=0}^n [\ln(y+k) - \ln(x+k)]
                 + (y-x)\ln{1\over n+1} \Bigr\}                          \cr 
        &\mkern-30mu =\lim_{n\to +\infty} \Biggl\{
                 (y+n)\ln(y+n) - (y+n) - y\ln y + y - (x+n)\ln(x+n)             \cr
        &\mkern40mu  + (x+n) + x\ln x - x
         + {1\over2} \bigl( \ln y - \ln x + \ln(y+n) - \ln(x+n) \bigr)          \cr
        &\mkern10mu +\sum_{h=1}^{p-1} {b_h\over(2h)!}
                \big( f^{(2h-1)}(n) - {1\over y^{2h-1}} + {1\over x^{2h-1}}\big) 
         + T_{p,n}(x,y) + (y-x)\ln{1\over n+1} \Biggr\}                  \cr
\end{aligned}
$$

Or $(y+n)\ln(y+n) - (x+n)\ln(x+n) = (y-x)\ln n + y - x + o(1)$
quand $n\to +\infty$. On obtient donc après simplifications
$G(y)-G(x)=g(x)-g(y)+R_n(x,y)$, ce qu'il fallait démontrer.

\medskip
8. D'après la question 5 et l'expression des dérivées successives de $f$ donnée dans
la question 6, $T_{p,n}(x,y)$ est majoré en valeur absolue par le produit d'une constante
et de l'intégrale
$\int_0^n \left| {1\over (y+t)^{2p}} - {1\over (x+t)^{2p}} \right| dt$.
$|R_p(x,y)|$ est majoré de la m{\^e}me façon en remplaçant la borne finale d'intégration
$n$ par $+\infty$. L'argument de la valeur absolue gardant un signe constant, l'intégrale
majorante est égale à
${1\over 2p-1} \left[ \ \left| {1\over (y+t)^{2p-1}}
                                - {1\over (x+t)^{2p-1}} \right| \ \right]_0^{+\infty}$
et on obtient ainsi l'estimée souhaitée.


\medskip
9. On a $g(m)=m\ln m - m - {1\over2}\ln m + o(1)$ et $G(m)=-\ln(m-1)!=-\ln m! + \ln m$ pour
$m$ entier, donc
le résultat demandé découle immédiatement de la formule de Stirling
$m! \sim \sqrt{2\pi} m^{m+{1\over2}} e^{-m}$.

\medskip
10. Le résultat demandé est obtenu à partir de l'égalité de la question 7 par passage à la
limite. On fait tendre $x$ vers $+\infty$ par valeurs entières et on tient compte de l'estimée
obtenue dans la question 8.

\medskip
11. En calculant les premiers termes du développement limité de la question A.2.a,
on trouve $b_1=1/6,\ b_2=-1/30,\ b_3=1/42$.
Des questions 3 et 10, il résulte que
$$\begin{aligned}
\ln(m!) &= -G(m) + \ln m        \cr
        &= m\ln m - m + {1\over2}\ln m + {1\over2}\ln2\pi
        + {1\over12} {1\over m} - {1\over360} {1\over m^3} + {1\over1260} {1\over m^5}
        + O({1\over m^7})  \cr
\end{aligned}
$$
\fincorrection
\nocorrection
\correction{004019}
$e^{1/3}\left(1+\frac7{90}x^2 +  o(x^3)\right)$.
\fincorrection
\correction{005426}
\begin{enumerate}
 \item  Si $x\in]0,\pi[$, $\sin x>0$, de sorte que la fonction proposée est bien définie sur un voisinage pointé de $\frac{\pi}{2}$ (c'est-à-dire un voisinage de $\frac{\pi}{2}$ auquel on a enlevé le point $\frac{\pi}{2}$) et de plus $(\sin x)^{1/(2x-\pi)}=e^{\ln(\sin x)/(2x-\pi)}$.
Quand $x$ tend vers $\frac{\pi}{2}$, $\sin x$ tend vers $1$ et donc

$$\ln(\sin x)\sim\sin x-1=-\left(1-\cos\left(\frac{\pi}{2}-x\right)\right)\sim-\frac{1}{2}\left(\frac{\pi}{2}-x\right)^2=-\frac{(2x-\pi)^2}{8}.$$
Donc, $\frac{\ln(\sin x)}{2x-\pi}\sim-\frac{2x-\pi}{8}\rightarrow0$ et enfin $(\sin x)^{1/(2x-\pi)}=e^{\ln(\sin x)/(2x-\pi)}\rightarrow e^0=1$.

\begin{center}
\shadowbox{
$\lim_{x \rightarrow \frac{\pi}{2}}(\sin x)^{1/(2x-\pi)}=1$.
}
\end{center}
 \item  Si $x\in]0,\pi[\setminus\left\{\frac{\pi}{2}\right\}$, $|\tan x|>0$, de sorte que la fonction proposée est bien définie sur un voisinage pointé de $\frac{\pi}{2}$ et de plus $|\tan x|^{\cos x}=e^{\cos x\ln(|\tan x|)}$. Quand $x$ tend vers $\frac{\pi}{2}$, 

$$\ln|\tan x|=\ln|\sin x|-\ln|\cos x|\sim-\ln|\cos x|,$$
puis $\cos x\ln|\tan x|\sim-\cos x\ln|\cos x|\rightarrow0$ (car, quand $u$ tend vers $0$, $u\ln u\rightarrow 0$).
Donc, $|\tan x|^{\cos x}=e^{\cos x\ln|\tan x|}\rightarrow e^0=1$.

\begin{center}
\shadowbox{
$\lim_{x \rightarrow \frac{\pi}{2}}|\tan x|^{\cos x}=1$.
}
\end{center}
 \item  Quand $n$ tend vers $+\infty$, $\cos\frac{n\pi}{3n+1}+\sin\frac{n\pi}{6n+1}\rightarrow\cos\frac{\pi}{3}+\sin\frac{\pi}{6}=1$ (et on est en présence d'une indétermination du type $1^{+\infty}$). Quand $n$ tend vers $+\infty$,

\begin{align*}\ensuremath
\cos\frac{n\pi}{3n+1}&=\cos\left(\frac{\pi}{3}\left(1+\frac{1}{3n}\right)^{-1}\right)
=\cos\left(\frac{\pi}{3}-\frac{\pi}{9n}+o\left(\frac{1}{n}\right)\right)\\
 &=\frac{1}{2}\cos\left(\frac{\pi}{9n}+o\left(\frac{1}{n}\right)\right)+\frac{\sqrt{3}}{2}\sin\left(\frac{\pi}{9n}+o\left(\frac{1}{n}\right)\right)=\frac{1}{2}\left(1+o\left(\frac{1}{n}\right)\right)+\frac{\sqrt{3}}{2}\left(\frac{\pi}{9n}+o\left(\frac{1}{n}\right)\right)\\
  &=\frac{1}{2}+\frac{\sqrt{3}\pi}{18n}+o\left(\frac{1}{n}\right)
\end{align*}
De même, 

\begin{align*}\ensuremath
\sin\frac{n\pi}{6n+1}&=\sin\left(\frac{\pi}{6}\left(1+\frac{1}{6n}\right)^{-1}\right)=\sin\left(\frac{\pi}{6}-\frac{\pi}{36n}
+o\left(\frac{1}{n}\right)\right)\\
 &=\frac{1}{2}\cos\left(\frac{\pi}{36n}+o\left(\frac{1}{n}\right)\right)-\frac{\sqrt{3}}{2}
 \sin\left(\frac{\pi}{36n}+o\left(\frac{1}{n}\right)\right)
 =\frac{1}{2}-\frac{\sqrt{3}\pi}{72n}+o\left(\frac{1}{n}\right).
\end{align*}
Puis, 

$$n\ln\left(\cos\frac{n\pi}{3n+1}+\sin\frac{n\pi}{6n+1}\right)=n\ln\left(1+\frac{\sqrt{3}\pi}{24n}+ o\left(\frac{1}{n}\right)\right)= n\left(\frac{\sqrt{3}\pi}{24n}+o\left(\frac{1}{n}\right)\right)=\frac{\sqrt{3}\pi}{24}+o(1),$$
et donc

\begin{center}
\shadowbox{
$\lim_{n\rightarrow +\infty}\left(\cos\frac{n\pi}{3n+1}+\sin\frac{n\pi}{6n+1}\right)^n=e^{\sqrt{3}\pi/24}$.
}
\end{center}
 \item  Quand $x$ tend vers $0$, $\ln(\cos x)\sim\cos x-1\sim-\frac{x^2}{2}$. Puis, $\ln|x|\ln(\cos x)\sim-\frac{x^2}{2}\ln|x|\rightarrow0$.
Donc, $(\cos x)^{\ln|x|}\rightarrow e^0=1$.

\begin{center}
\shadowbox{
$\lim_{x\rightarrow 0}(\cos x)^{\ln|x|}=1$.
}
\end{center}
 \item  Quand $x$ tend vers $\frac{\pi}{2}$, $\frac{1}{1-\sin x}$ tend vers $+\infty$. Posons $h=x-\frac{\pi}{2}$ puis $\varepsilon=\mbox{sgn}(h)$, de sorte que 

$$(\cos x)e^{1/(1-\sin x)}=-\varepsilon|\sin h|e^{1/(1-\cos h)}=-\varepsilon e^{\ln|\sin h|+\frac{1}{1-\cos h}}.$$
Or, quand $h$ tend vers $0$, 

$$\ln|\sin h|+\frac{1}{1-\cos h}=\frac{(1-\cos h)\ln|\sin h|+1}{1-\cos h}=\frac{(-\frac{h^2}{2}+o(h^2))(\ln|h|+o(\ln|h|))+1}{\frac{h^2}{2}+o(h^2)}=\frac{1+o(1)}{\frac{h^2}{2}+o(h^2)}
\sim\frac{2}{h^2},$$ et donc, quand $h$ tend vers $0$, $\ln|\sin h|+\frac{1}{1-\cos h}\sim\frac{2}{h^2}\rightarrow+\infty$. Par suite,

\begin{center}
\shadowbox{
$\lim_{x\rightarrow \pi/2,\;x<\pi/2}\cos(x)e^{1/(1-\sin x)}=+\infty$ et $\lim_{x\rightarrow \pi/2,\;x>\pi/2}\cos(x)e^{1/(1-\sin x)}=-\infty$.
}
\end{center}

 \item  Pour $x\in\Rr$, $2\cos^2x-3\cos x+1=(2\cos x-1)(\cos x-1)$ et donc 

$$\forall x\in\Rr,\;2\cos^2x-3\cos x+1=0\Leftrightarrow x\in\left(\pm\frac{\pi}{3}+2\pi\Zz\right)\cup2\pi\Zz.$$
Pour $x\notin\left(\pm\frac{\pi}{3}+2\pi\Zz\right)\cup2\pi\Zz$,
  
$$\frac{2\cos^2x+\cos x-1}{2\cos^2x-3\cos x+1}=\frac{(2\cos x-1)(\cos x+1)}{(2\cos x-1)(\cos x-1)}
=\frac{\cos x+1}{\cos x-1},$$
et donc, $\lim_{x\rightarrow \pi/3}\frac{2\cos^2x+\cos x-1}{2\cos^2x-3\cos x+1}=\frac{\frac{1}{2}+1}{\frac{1}{2}-1}=-3$.

\begin{center}
\shadowbox{
$\lim_{x\rightarrow \pi/3}\frac{2\cos^2x+\cos x-1}{2\cos^2x-3\cos x+1}=-3$.
}
\end{center}
 \item  Quand $x$ tend vers $0$, 
 
$$\frac{1+\tan x}{1+\tanh x}=\frac{1+x+o(x)}{1+x+o(x)}=(1+x+o(x)(1-x+o(x))=1+o(x).$$
Puis, quand $x$ tend vers $0$,

$$\frac{1}{\sin x}\ln\left(\frac{1+\tan x}{1+\tanh x}\right)=\frac{\ln(1+o(x))}{x+o(x)}=\frac{o(x)}{x+o(x)}=\frac{o(1)}{1+o(1)}\rightarrow0.$$
Donc, 

\begin{center}
\shadowbox{
$\lim_{x\rightarrow 0}\left(\frac{1+\tan x}{1+\tanh x}\right)^{1/\sin x}=1$.
}
\end{center}

 \item  Quand $x$ tend vers $e$ par valeurs inférieures, $\ln(x)$ tend vers $1$ et donc

$$\ln(\ln x)\sim\ln x-1=\ln\left(\frac{x}{e}\right)\sim\frac{x}{e}-1=-\frac{1}{e}(e-x),$$
puis,

$$\ln(e-x)\ln(\ln x)\sim-\frac{1}{e}(e-x)\ln(e-x)\rightarrow 0,$$
et donc $(\ln x)^{\ln(e-x)}=e^{\ln(e-x)\ln(\ln x)}\rightarrow 1$.

\begin{center}
\shadowbox{
$\displaystyle\lim_{\substack{x\rightarrow e\\ x<e}}(\ln x)^{\ln(e-x)}=1$.
}
\end{center}

 \item  Quand $x$ tend vers $1$ par valeurs supérieures, $x\ln x\rightarrow0$, et donc

$$x^x-1=e^{x\ln x}-1\sim x\ln x\sim1\times(x-1)=x-1.$$
Ensuite, $\sqrt{x^2-1}$ tend vers $0$ et donc

$$\ln(1-\sqrt{x^2-1})\sim-\sqrt{x^2-1}=-\sqrt{(x-1)(x+1)}\sim-\sqrt{2(x-1)}.$$
Finalement, quand $x$ tend vers $1$ par valeurs supérieures, 

$$\frac{x^x-1}{\ln(1-\sqrt{x^2-1})}\sim\frac{x-1}{-\sqrt{2(x-1)}}=-\frac{1}{\sqrt{2}}\sqrt{x-1}\rightarrow0.$$

\begin{center}
\shadowbox{
$\displaystyle\lim_{\substack{x\rightarrow 1\\ x>e}}\frac{x^x-1}{\ln(1-\sqrt{x^2-1})}=0$.
}
\end{center}
 \item  Quand $x$ tend vers $+\infty$, 

$$\ln(\ch x-1)\sim\ln(\ch x)\sim\ln\left(\frac{e^x}{2}\right)=x-\ln2\sim x,$$ et donc 

$$\frac{x\ln(\ch x-1)}{x^2+1}\sim\frac{x\times x}{x^2}=1.$$

\begin{center}
\shadowbox{
$\displaystyle\lim_{x\rightarrow+\infty}\frac{x\ln(\ch x-1)}{x^2+1}=1$.
}
\end{center}
\item  Quand $x$ tend vers $0$ par valeurs supérieures,

$$\ln(x-x^2)+x-\ln x=x+\ln(1-x)=-\frac{x^2}{2}+o(x^2)\sim-\frac{x^2}{2}.$$
Ensuite,

$$(\sin x)^x=e^{x\ln(\sin x)}=e^{x\ln(x-\frac{x^3}{6}+o(x^3))}=e^{x\ln x}e^{x\ln(1-\frac{x^2}{6}+o(x^2))}=x^xe^{-\frac{x^3}{6}+o(x^3)}=x^x\left(1-\frac{x^3}{6}+o(x^3)\right),$$
et,

$$x^{\sin x}=e^{(x-\frac{x^3}{6}+o(x^3))\ln x}=e^{x\ln x}e^{-\frac{x^3\ln x}{6}+o(x^3\ln x)}=x^x\left(1-\frac{x^3\ln x}{6}+o(x^3\ln x)\right).$$
Donc,

$$(\sin x)^x-x^{\sin x}=x^x\left(1-\frac{x^3}{6}+o(x^3)\right)-x^x\left(1-\frac{x^3\ln x}{6}+o(x^3\ln x)\right)=x^x\left(\frac{x^3\ln x}{6}+o(x^3\ln x)\right)\sim\frac{x^3\ln x}{6},$$
et enfin

$$\frac{(\sin x)^x-x^{\sin x}}{\ln(x-x2^)+x-\ln x}\sim\frac{x^3\ln x/6}{-x^2/2}=-\frac{x\ln x}{3}\rightarrow0.$$

\begin{center}
\shadowbox{
$\displaystyle\lim_{\substack{x\rightarrow 0\\ x>0}}\frac{(\sin x)^x-x^{\sin x}}{\ln(x-x2^)+x-\ln x}=0$.
}
\end{center}
\item Quand $x$ tend vers $+\infty$,

$$\ln(x+1)=\ln x+\ln\left(1+\frac{1}{x}\right)=\ln x+\frac{1}{x}+o\left(\frac{1}{x}\right),$$
puis
 
$$\frac{\ln(x+1)}{\ln x}=1+\frac{1}{x\ln x}+o\left(\frac{1}{x\ln x}\right).$$
 

Ensuite,

$$x\ln\left(\frac{\ln(x+1)}{\ln x}\right)=x\ln\left(1+\frac{1}{x\ln x}+o\left(\frac{1}{x\ln x}\right)\right)=\frac{1}{\ln x}+o\left(\frac{1}{\ln x}\right)\rightarrow0.$$
Donc, $\left(\frac{\ln(x+1)}{\ln x}\right)^x=\text{exp}\left(x\ln\left(\frac{\ln(x+1)}{\ln x}\right)\right)\rightarrow e^0=1$.

\begin{center}
\shadowbox{
$\displaystyle\lim_{x\rightarrow+\infty}\left(\frac{\ln(x+1)}{\ln x}\right)^x=1$.
}
\end{center}
 \item Quand $x$ tend vers $\frac{1}{\sqrt{2}}$,

\begin{align*}\ensuremath
\frac{(\Arcsin x)^2-\frac{\pi^2}{16}}{2x^2-1}&=\frac{1}{2}\times\frac{\Arcsin x+\frac{\pi}{4}}{x+\frac{1}{\sqrt{2}}}
\times\frac{\Arcsin x-\frac{\pi}{4}}{x-\frac{1}{\sqrt{2}}}\sim\frac{1}{2}\times\frac{\frac{\pi}{4}+\frac{\pi}{4}}{\frac{1}{\sqrt{2}}+\frac{1}{\sqrt{2}}}\times\frac{\Arcsin x-\frac{\pi}{4}}{x-\frac{1}{\sqrt{2}}}=\frac{\pi}{4\sqrt{2}}\frac{\Arcsin x-\frac{\pi}{4}}{x-\frac{1}{\sqrt{2}}}\\
 &\rightarrow\frac{\pi}{4\sqrt{2}}(\Arcsin)'(\frac{1}{\sqrt{2}})=\frac{\pi}{4\sqrt{2}}\frac{1}{\sqrt{1-\frac{1}{2}}}=\frac{\pi}{4}.
\end{align*}

\begin{center}
\shadowbox{
$\lim_{x \rightarrow 1/\sqrt{2}}\frac{(\Arcsin x)^2-\frac{\pi^2}{16}}{2x^2-1}=\frac{\pi}{4}
$.
}
\end{center}
 \item Quand $x$ tend vers $+\infty$,

\begin{align*}\ensuremath
x\ln\left(\frac{\cos\left(a+\frac{1}{x}\right)}{\cos a}\right)&=x\ln\left(\cos\frac{1}{x}-\tan a\sin\frac{1}{x}\right)=x\ln\left(1-\frac{\tan a}{x}+o\left(\frac{1}{x}\right)\right)=x\left(-\frac{\tan a}{x}+o\left(\frac{1}{x}\right)\right)\\
 &=-\tan a+o(1),
\end{align*}
et donc $\lim_{x\rightarrow +\infty}\left(\frac{\cos\left(a+\frac{1}{x}\right)}{\cos a}\right)^x=e^{-\tan a}$.

\begin{center}
\shadowbox{
$\lim_{x\rightarrow +\infty}\left(\frac{\cos\left(a+\frac{1}{x}\right)}{\cos a}\right)^x=e^{-\tan a}
$.
}
\end{center}
\end{enumerate}
\fincorrection
\correction{005427}
\begin{enumerate}
 \item 
\begin{align*}\ensuremath 
\frac{1}{1-x^2-x^3}\underset{x\rightarrow0}{=}1+(x^2+x^3)+(x^2+x^3)^2+(x^2+x^3)^3+o(x^7)= 1+x^2+x^3+x^4+2x^5+2x^6+3x^7+o(x^7).
\end{align*}

\begin{center}
\shadowbox{
$\frac{1}{1-x^2-x^3}\underset{x\rightarrow0}{=}1+x^2+x^3+x^4+2x^5+2x^6+3x^7+o(x^7)$.
}
\end{center}
 \item 

\begin{align*}\ensuremath
\frac{1}{\cos x}&\underset{x\rightarrow0}{=}\left(1-\frac{x^2}{2}+\frac{x^4}{24}-\frac{x^6}{720}+o(x^7)\right)^{-1}
=1+\left(\frac{x^2}{2}-\frac{x^4}{24}+\frac{x^6}{720}\right)+\left(\frac{x^2}{2}-\frac{x^4}{24}\right)^2+\left(\frac{x^2}{2}\right)^3+o(x^7)\\
 &=1+\frac{x^2}{2}+x^4\left(-\frac{1}{24}+\frac{1}{4}\right)+x^6\left(\frac{1}{720}-\frac{1}{24}+\frac{1}{8}\right)+o(x^7)
=1+\frac{1}{2}x^2+\frac{5}{24}x^4+\frac{61}{720}x^6+o(x^7).
\end{align*}

\begin{center}
\shadowbox{
$\frac{1}{\cos x}\underset{x\rightarrow0}{=}1+\frac{1}{2}x^2+\frac{5}{24}x^4+\frac{61}{720}x^6+o(x^7)$.
}
\end{center}
 \item  \textbf{Remarques.}

  \begin{enumerate}
  \item Pour $x\in\left]-\frac{\pi}{2},\frac{\pi}{2}\right[\setminus\{0\}$, on a $0<\frac{x}{\tan x}<1$ et donc la fonction $x\mapsto\Arccos\left(\frac{x}{\tan x}\right)$ est définie sur $\left]-\frac{\pi}{2},\frac{\pi}{2}\right[\setminus\{0\}$ (qui est un voisinage pointé de $0$).\rule[-2mm]{0mm}{6mm}
  \item Quand $x$ tend vers $0$, $\frac{x}{\tan x}\rightarrow1$ et donc $\Arccos\left(\frac{x}{\tan x}\right)=o(1)$ (développement limité à l'ordre $0$).
  \item La fonction $x\mapsto\Arccos x$ n'est pas dérivable en $1$ et n'admet donc pas en $1$ de développement limité d'ordre supérieur ou égal à $1$ (donc à priori, c'est mal parti).
  \item La fonction proposée est paire et, si elle admet en $0$ un développement limité d'ordre $3$, sa partie régulière ne contient que des exposants pairs.

  \end{enumerate}

\textbullet~Recherche d'un équivalent simple de $\Arccos x$ en $1$ à gauche.
Quand $x$ tend vers $1$ par valeurs inférieures, $\Arccos x\rightarrow 0$ et donc,

$$\Arccos x\sim\sin(\Arccos x)=\sqrt{1-x^2}=\sqrt{(1+x)(1-x)}\sim\sqrt{2}\sqrt{1-x}.$$ 
\textbullet~Déterminons un équivalent simple de $\Arccos\left(\sqrt{\frac{x}{\tan x}}\right)$ en $0$. D'après ce qui précède,

$$\Arccos\left(\sqrt{\frac{x}{\tan x}}\right)\sim\sin\left(\Arccos\left(\sqrt{\frac{x}{\tan x}}\right)\right)=\sqrt{1-\left(\sqrt{\frac{x}{\tan x}}\right)^2}=\sqrt{\frac{\tan x-x}{\tan x}}\sim\sqrt{\frac{x^3/3}{x}}=\frac{|x|}{\sqrt{3}}.$$
Ainsi, la fonction $x\mapsto\Arccos\left(\sqrt{\frac{x}{\tan x}}\right)$ n'est pas dérivable en $0$ (mais est dérivable à droite et à gauche) et n'admet donc pas de développement limité d'ordre supérieur ou égal à $1$ (mais admet éventuellement des développements limités à gauche et à droite pour lesquels la remarque initiale sur la parité des exposants ne tient plus).
\textbullet~Déterminons un équivalent simple de $f(x)=\Arccos\left(\sqrt{\frac{x}{\tan x}}\right)-\frac{x}{\sqrt{3}}$ quand $x$ tend vers $0$ par valeurs supérieures.

\begin{align*}\ensuremath
\Arccos\left(\sqrt{\frac{x}{\tan x}}\right)-\frac{x}{\sqrt{3}}&\sim\sin\left(\Arccos\left(\sqrt{\frac{x}{\tan x}}\right)-\frac{x}{\sqrt{3}}\right)\\
 &=\sin\left(\Arccos\left(\sqrt{\frac{x}{\tan x}}\right)\right)\cos\left(\frac{x}{\sqrt{3}}\right)-\sin\left(\frac{x}{\sqrt{3}}\right)\cos\left(\Arccos\left(\sqrt{\frac{x}{\tan x}}\right)\right)\\
 &=\sqrt{\frac{\tan x-x}{\tan x}}\cos\left(\frac{x}{\sqrt{3}}\right)-\sqrt{\frac{x}{\tan x}}\sin\left(\frac{x}{\sqrt{3}}\right)=g(x)
\end{align*}
Maintenant,

\begin{align*}\ensuremath
\sqrt{\frac{\tan x-x}{\tan x}}&=\left(\left(\frac{x^3}{3}+\frac{2x^5}{15}+o(x^5)\right)\left(x+\frac{x^3}{3}+o(x^3)\right)^{-1}\right)^{1/2}=
\frac{1}{\sqrt{x}}\left(\left(\frac{x^3}{3}+\frac{2x^5}{15}+o(x^5)\right)\left(1-\frac{x^2}{3}+o(x^2)\right)\right)^{1/2}\\
 &=\frac{1}{\sqrt{x}}\left(\frac{x^3}{3}+\frac{x^5}{45}+o(x^5)\right)^{1/2}=\frac{x}{\sqrt{3}}\left(1+\frac{x^2}{15}+o(x^2)\right)^{1/2}\\
 &=\frac{x}{\sqrt{3}}+\frac{x^3}{30\sqrt{3}}+o(x^3),
\end{align*}
et donc,

$$
\sqrt{\frac{\tan x-x}{\tan x}}\cos\left(\frac{x}{\sqrt{3}}\right)=\left(\frac{x}{\sqrt{3}}+\frac{x^3}{30\sqrt{3}}+o(x^3)\right)\left(1-\frac{x^2}{6}+o(x^2)\right)=\frac{x}{\sqrt{3}}-\frac{2x^3}{15\sqrt{3}}+o(x^3).
$$
Ensuite,

\begin{align*}\ensuremath
\sqrt{\frac{x}{\tan x}}\sin\left(\frac{x}{\sqrt{3}}\right)&=\left(1+\frac{x^2}{3}+o(x^2)\right)^{-1/2}\left(\frac{x}{\sqrt{3}}-\frac{x^3}{18\sqrt{3}}+o(x^3)\right)=\left(1-\frac{x^2}{6}+o(x^2)\right)\left(\frac{x}{\sqrt{3}}-\frac{x^3}{18\sqrt{3}}+o(x^3)\right)\\
 &=\frac{x}{\sqrt{3}}-\frac{2x^3}{9\sqrt{3}}+o(x^3),
\end{align*}
et finalement,

$$g(x)=\left(\frac{x}{\sqrt{3}}-\frac{2x^3}{15\sqrt{3}}+o(x^3)\right)-\left(\frac{x}{\sqrt{3}}-\frac{2x^3}{9\sqrt{3}}+o(x^3)\right)=\frac{4x^3}{45\sqrt{3}}+o(x^3)\sim\frac{4x^3}{45\sqrt{3}}.$$

Ainsi, quand $x$ tend vers $0$ par valeurs supérieures,

$$\Arccos\left(\sqrt{\frac{x}{\tan x}}\right)-\frac{x}{\sqrt{3}}=\frac{4x^3}{45\sqrt{3}}+o(x^3).$$
$f$ étant paire, on en déduit que

\begin{center}
\shadowbox{$\Arccos\left(\sqrt{\frac{x}{\tan x}}\right)\underset{x\rightarrow0}{=}\frac{|x|}{\sqrt{3}}+\frac{4|x|^3}{45\sqrt{3}}+o(x^3).$}
\end{center}
(Ce n'est pas un développement limité).
 \item  La fonction $x\mapsto\tan x$ est trois fois dérivable en $\frac{\pi}{4}$ et admet donc en $\frac{\pi}{4}$ un développement limité d'ordre $3$ à savoir son développement de \textsc{Taylor}-\textsc{Young}.
$\tan\frac{\pi}{4}=1$ puis $(\tan)'\left(\frac{\pi}{4}\right)=1+\tan^2\frac{\pi}{4}=2$. Ensuite, $(\tan)''(x)=2\tan x(1+\tan^2x)$ et $(\tan)''(\frac{\pi}{4})=4$. Enfin,

$$(\tan)^{(3)}(x)=2(1+\tan^2x)^2+4\tan^2x(1+\tan^2x),$$
et $(\tan)^{(3)}(\frac{\pi}{4})=16$. Finalement,

\begin{center}
\shadowbox{$\tan x\underset{x\rightarrow\pi/4}=1+2\left(x-\frac{\pi}{4}\right)+2\left(x-\frac{\pi}{4}\right)^2+\frac{8}{3}\left(x-\frac{\pi}{4}\right)^3+o\left(\left(x-\frac{\pi}{4}\right)^3\right).$
}
\end{center}
 \item 

$$\frac{1}{x^2}\ln(\ch x)\underset{x\rightarrow0}{=}\frac{1}{x^2}\ln\left(1+\frac{x^2}{2}+\frac{x^4}{24}+o(x^4)\right)=\frac{1}{x^2}\left(\frac{x^2}{2}+\frac{x^4}{24}-\frac{1}{2}\left(\frac{x^2}{2}\right)^2+o(x^4)\right)=\frac{1}{2}-\frac{x^2}{12}+o(x^2),$$
et donc 

$$(\ch x)^{1/x^2}=e^{\frac{1}{2}-\frac{x^2}{12}+o(x^2)}=e^{1/2}e^{-\frac{x^2}{12}+o(x^2)}=\sqrt{e}-\frac{\sqrt{e}}{12}x^2+o(x^2).$$

\begin{center}
\shadowbox{$(\ch x)^{1/x^2}=\underset{x\rightarrow0}=\sqrt{e}-\frac{\sqrt{e}}{12}x^2+o(x^2).$
}
\end{center}
 \item  $\tan^3x(\cos(x^2)-1)=\tan x\times\tan^2x(\cos(x^2)-1)$ et un équivalent de $\tan^2x(\cos(x^2)-1)$ en $0$ est $-\frac{x^6}{2}$. On écrit donc $\tan x$ à l'ordre $2$. De même, un équivalent de $\tan^3x$ est $x^3$ et on écrit donc $\cos(x^2)-1$ à l'ordre $5$.
$$\tan^3x(\cos(x^2)-1)\underset{x\rightarrow0}{=}(x+o(x^2))^3\left(-\frac{x^4}{2}+o(x^5)\right)=(x^3+o(x^4))\left(-\frac{x^4}{2}
+o(x^5)\right)=-\frac{x^7}{2}+o(x^8).$$

\begin{center}
\shadowbox{$\tan^3x(\cos(x^2)-1)\underset{x\rightarrow0}{=}-\frac{x^7}{2}+o(x^8).$
}
\end{center}

 \item  On pose $h=x-1$ ou encore $x=1+h$, de sorte que $x$ tend vers $1$ si et seulement si $h$ tend vers $0$.

\begin{align*}\ensuremath
\frac{\ln(1+x)}{x^2}&=\ln(2+h)(1+h)^{-2}\\
 &\underset{h\rightarrow0}{=}\left(\ln2+\ln\left(1+\frac{h}{2}\right)\right)\left(1-2h+\frac{(-2)(-3)}{2}h^2+\frac{(-2)(-3)(-4)}{6}h^3
 +o(h^3)\right)\\
  &=\left(\ln2+\frac{h}{2}-\frac{h^2}{8}+\frac{h^3}{24}+o(h^3)\right)(1-2h+3h^2-4h^3+o(h^3))\\
  &=\ln2+\left(\frac{1}{2}-2\ln2\right)h+\left(3\ln2-\frac{9}{8}\right)h^2+\left(-4\ln2+\frac{43}{24}\right)h^3 + o(h^3).
\end{align*}
Donc,

\begin{center}
\shadowbox{
$\frac{\ln(1+x)}{x^2}\underset{x\rightarrow1}{=}\ln2+\left(\frac{1}{2}-2\ln2\right)(x-1)+\left(3\ln2-\frac{9}{8}\right)(x-1)^2+\left(-4\ln2+\frac{43}{24}\right)(x-1)^3+o((x-1)^3)$.
}
\end{center}
 \item  Pour $x$ réel, posons $f(x)=\Arctan(\cos x)$.
$f$ est dérivable sur $\Rr$, et pour $x$ réel, $f'(x)=-\frac{\sin x}{1+\cos^2x}$. Puis,

\begin{align*}\ensuremath
f'(x)&\underset{x\rightarrow0}{=}-\left(x-\frac{x^3}{6}+o(x^4)\right)\left(1+\left(1-\frac{x^2}{2}+o(x^3)\right)^2\right)^{-1}\\
 &=-\left(x-\frac{x^3}{6}+o(x^4)\right)\left(2-x^2+o(x^3)\right)^{-1}=-\frac{1}{2}\left(x-\frac{x^3}{6}+o(x^4)\right)
\left(1-\frac{x^2}{2}+o(x^3)\right)^{-1}\\
 &=-\frac{1}{2}\left(x-\frac{x^3}{6}+o(x^4)\right)\left(1+\frac{x^2}{2}+o(x^3)\right)=-\frac{x}{2}-\frac{x^3}{6}+o(x^4).
\end{align*}
Donc, $f'$ admet un développement limité d'ordre $4$ en $0$ et on sait que $f$ admet en $0$ un développement limité d'ordre $5$ obtenu par intégration.

$$\Arctan(\cos x)\underset{x\rightarrow0}{=}\Arctan(\cos0)-\frac{x^2}{4}-\frac{x^4}{24}+o(x^5)=\frac{\pi}{4}-\frac{x^2}{4}-\frac{x^4}{24}+o(x^5).$$

\begin{center}
\shadowbox{
$\Arctan(\cos x)\underset{x\rightarrow0}{=}\frac{\pi}{4}-\frac{x^2}{4}-\frac{x^4}{24}+o(x^5).$
}
\end{center}
 \item  Pour $x>-1$, posons $f(x)=\Arctan\sqrt{\frac{x+1}{x+2}}$. $f$ est dérivable sur $]-1,+\infty[$ et pour $x>-1$,
 
\begin{align*}\ensuremath
f'(x)&=\frac{1}{(x+2)^2}\frac{1}{2\sqrt{\frac{x+1}{x+2}}}\frac{1}{1+\frac{x+1}{x+2}}=\frac{1}{2}\times\frac{1}{(2x+3)\sqrt{(1+x)(2+x)}}\\
 &=\frac{1}{2\times3\times\sqrt{2}}\left(1+\frac{2x}{3}\right)^{-1}(1+x)^{-1/2}\left(1+\frac{x}{2}\right)^{-1/2}=\frac{1}{6\sqrt{2}}\left(1-\frac{2x}{3}+o(x)\right)\left(1-\frac{x}{2}+o(x)\right)\left(1-\frac{x}{4}+o(x)\right)\\
 &=\frac{1}{6\sqrt{2}}\left(1-\left(\frac{2}{3}+\frac{1}{2}+\frac{1}{4}\right)x+o(x)\right)=\frac{1}{6\sqrt{2}}\left(1-\frac{17x}{12}+o(x)\right).
\end{align*}
Ainsi, $f'$ admet donc en $0$ un développement limité d'ordre $1$ et on sait alors que $f$ admet en $0$ un développement limité d'ordre $2$ obtenu par intégration.

\begin{center}
\shadowbox{
$\Arctan\sqrt{\frac{x+1}{x+2}}\underset{x\rightarrow0}{=}\Arctan\frac{1}{\sqrt{2}}+\frac{1}{6\sqrt{2}}x-\frac{17}{144\sqrt{2}}x^2+o(x^2).$
}
\end{center}
 \item  

\begin{align*}\ensuremath 
\frac{1}{\sqrt{1-x^2}}&=(1-x^2)^{-1/2}\underset{x\rightarrow0}{=}1+\left(-\frac{1}{2}\right)(-x^2)+\frac{\left(-\frac{1}{2}\right)\left(-\frac{3}{2}\right)}{2}(-x^2)^2+\frac{\left(-\frac{1}{2}\right)\left(-\frac{3}{2}\right)\left(-\frac{5}{2}\right)}{6}(-x^2)^3+o(x^7)\\
 &\underset{x\rightarrow0}{=}1+\frac{1}{2}x^2+\frac{3}{8}x^4+\frac{5}{16}x^6+o(x^7).
\end{align*}

Donc, $\Arcsin x\underset{x\rightarrow0}{=}x+\frac{x^3}{6}+\frac{3x^5}{40}+\frac{5x^7}{112}+o(x^8)$. Ensuite,

\begin{align*}\ensuremath
\frac{1}{\Arcsin^2x}&=(\Arcsin x)^{-2}\underset{x\rightarrow0}{=}\frac{1}{x^2}\left(1+\frac{x^2}{6}+\frac{3x^4}{40}+\frac{5x^6}{112}+o(x^7)\right)^{-2}\\
 &=\frac{1}{x^2}\left(1-2\left(\frac{x^2}{6}+\frac{3x^4}{40}+\frac{5x^6}{112}\right)+3\left(\frac{x^2}{6}+\frac{3x^4}{40}\right)^2-4\left(\frac{x^2}{6}\right)^3+o(x^7)\right)\\
 &=\frac{1}{x^2}-\frac{1}{3}+\left(-\frac{3}{20}+\frac{1}{12}\right)x^2+\left(-\frac{5}{56}+\frac{3}{40}-\frac{1}{54}\right)x^4+o(x^5)\\
 &=\frac{1}{x^2}-\frac{1}{3}-\frac{x^2}{15}-\frac{31x^4}{945}+o(x^5).
\end{align*}
Finalement,

\begin{center}
\shadowbox{
$\frac{1}{x^2}-\frac{1}{\Arcsin^2x}\underset{x\rightarrow0}{=}\frac{1}{3}+\frac{x^2}{15}+\frac{31x^4}{945}+o(x^5).$
}
\end{center}
\item  Pour $x$ réel, posons $f(x)=\frac{1}{\sqrt{1+x^4}}$. $f$ est continue sur $\Rr$ et admet donc des primitives sur $\Rr$. Soit $F$ la primitive de $f$ sur $\Rr$ qui s'annule en $0$ puis, pour $x$ réel, soit $g(x)=\int_{x}^{x^2}\frac{1}{\sqrt{1+t^4}}\;dt$.
$g$ est définie sur $\Rr$ et, pour $x$ réel $g(x)=F(x^2)-F(x)$. $g$ est dérivable sur $\Rr$ et, pour tout réel $x$,

$$g'(x)=2xF'(x^2)-F'(x)=2xf(x^2)-f(x)=\frac{2x}{\sqrt{1+x^8}}-\frac{1}{\sqrt{1+x^4}}.$$
Puis, 

$$g'(x)\underset{x\rightarrow0}{=}2x\left(1-\frac{1}{2}x^8+o(x^8)\right)-\left(1-\frac{1}{2}x^4+\frac{3}{8}x^8+o(x^9)\right)=-1+2x+\frac{1}{2}x^4-\frac{3}{8}x^8-x^9+o(x^9).$$
Ainsi $g'$ admet un développement limité d'ordre $9$ en $0$ et on sait que $g$ admet un développement limité d'ordre $10$ en $0$ obtenu par intégration. En tenant compte de $g(0)=0$, on obtient

\begin{center}
\shadowbox{
$g(x)\underset{x\rightarrow0}{=}-x+x^2+\frac{x^5}{10}-\frac{x^9}{24}-\frac{x^{10}}{10}+o(x^{10}).$
}\end{center}
\item 
\begin{align*}\ensuremath
\ln\left(\sum_{k=0}^{99}\frac{x^k}{k!}\right)&\underset{x\rightarrow0}{=}\ln\left(e^x-\frac{x^{100}}{100!}+o(x^{100})\right)=\ln(e^x)+\ln\left(1-e^{-x}\left(\frac{x^{100}}{100!}+o(x^{100})\right)\right)\\
 &=x+\ln\left(1-(1+o(1))\left(\frac{x^{100}}{(100)!}+o(x^{100})\right)\right)=x+\ln\left(1-\frac{x^{100}}{(100)!}+o(x^{100})\right)=x-\frac{x^{100}}{(100)!}+o(x^{100})
\end{align*}

\begin{center}
\shadowbox{
$\ln\left(\sum_{k=0}^{99}\frac{x^k}{k!}\right)\underset{x\rightarrow0}{=}x-\frac{x^{100}}{(100)!}+o(x^{100}).$
}\end{center}
\item Posons $h=x-\pi$ ou encore $x=\pi+h$ de sorte que $x$ tend vers $\pi$ si et seulement si $h$ tend vers $0$.

\begin{align*}\ensuremath
\sqrt[3]{4(\pi^3+x^3)}&=\sqrt[3]{4(\pi^3+(\pi+h)^3)}=\sqrt[3]{8\pi^3+12\pi^2h+12\pi h^2+4h^3}\\
 &\underset{h\rightarrow0}{=}2\pi\left(1+\frac{3h}{2\pi}+\frac{3h^2}{2\pi^2}+\frac{h^3}{2\pi^3}\right)^{1/3}\\
 &=2\pi\left(1+\frac{1}{3}\left(\frac{3h}{2\pi}+\frac{3h^2}{2\pi^2}+\frac{h^3}{2\pi^3}\right)-\frac{1}{9}\left(\frac{3h}{2\pi}+\frac{3h^2}{2\pi^2}\right)^2+\frac{5}{81}\left(\frac{3h}{2\pi}\right)^3+o(h^3)\right)\\
 &=2\pi+h+h^2\left(\frac{1}{\pi}-\frac{1}{2\pi}\right   )+h^3\left(\frac{1}{3\pi^2}-\frac{1}{\pi^2}+\frac{5}{12\pi^2}\right)+o(h^3)\\
 &=2\pi+h+\frac{h^2}{2\pi}-\frac{h^3}{4\pi^2}+o(h^3).
\end{align*}
Puis,

\begin{align*}\ensuremath
\tan(\sqrt[3]{4(\pi^3+x^3)})&=\tan\left(h+\frac{h^2}{2\pi}-\frac{h^3}{4\pi^2}+o(h^3)\right)\\
 &=\left(h+\frac{h^2}{2\pi}-\frac{h^3}{4\pi^2}\right)+\frac{1}{3}h^3+o(h^3)=h+\frac{h^2}{2\pi}+\left(\frac{1}{3}-\frac{1}{4\pi^2}\right)h^3+o(h^3).
\end{align*}
Finalement,

\begin{center}
\shadowbox{
$\tan(\sqrt[3]{4(\pi^3+x^3)})\underset{x\rightarrow\pi}{=}(x-1)+\frac{1}{2\pi}(x-1)^2+\left(\frac{1}{3}-\frac{1}{4\pi^2}\right)(x-1)^3+o((x-1)^3)$.
}
\end{center}
\end{enumerate}
\fincorrection
\correction{005428}
Puisque $a>0$, $b>0$ et que pour tout réel $x$, $\frac{a^x+b^x}{2}>0$, $f$ est définie sur $\Rr^*$, et pour

\begin{center}
$\forall x\in\Rr^*$, $f(x)=\text{exp}\left(\frac{1}{x}\ln\left(\frac{a^x+b^x}{2}\right)\right)$.
\end{center}
\textbf{Etude en 0}.

\begin{align*}\ensuremath
\ln\left(\frac{a^x+b^x}{2}\right)&=\ln\left(\frac{e^{x\ln a}+e^{x\ln b}}{2}\right)\underset{x\rightarrow0}{=}\ln\left(1+x\left(\frac{1}{2}\ln a+\frac{1}{2}\ln b\right)+x^2\left(\frac{1}{4}\ln^2a+\frac{1}{4}\ln^2b\right)+o(x^2)\right)\\
 &=\ln\left(1+x\ln\left(\sqrt{ab}\right)+x^2\frac{\ln^2a+\ln^2b}{4}+o(x^2)\right)=x\ln\left(\sqrt{ab}\right)+x^2\frac{\ln^2a+\ln^2b}{4}-\frac{1}{2}(x\ln\sqrt{ab})^2+o(x^2)\\
 &=x\ln\left(\sqrt{ab}\right)+\frac{1}{8}(\ln^2a-2\ln a\ln b+\ln^2b)x^2+o(x^2)=x\ln\left(\sqrt{ab}\right)+x^2\frac{1}{8}\ln^2\left(\frac{a}{b}\right)+o(x^2).
\end{align*}
Enfin,

$$f(x)=\left(\frac{a^x+b^x}{2}\right)^{1/x}=\mbox{exp}(\ln(\sqrt{ab})+\frac{1}{8}\ln^2\frac{a}{b}x+o(x))=\sqrt{ab}(1+\frac{1}{8}x\ln^2\frac{a}{b}+o(x)).$$
Ainsi, $f$ se prolonge par continuité en $0$ en posant $f(0)=\sqrt{ab}$. Le prolongement obtenu est dérivable en $0$ et
$f'(0)=\frac{\sqrt{ab}}{8}\ln^2\frac{a}{b}(>0)$.
\textbf{Etude en} $\bf{+\infty}$.
\begin{align*}\ensuremath
\frac{1}{x}\ln\left(\frac{1}{2}(a^x+b^x)\right)&=\frac{1}{x}\left(\ln(b^x)-\ln2+\ln\left(1+\left(\frac{a}{b}\right)^x\right)\right)
=\frac{1}{x}(x\ln b+o(x))\quad(\mbox{car}\;0<\frac{a}{b}<1)\\
 &=\ln b+o(1).
\end{align*}
et $\lim_{x\rightarrow +\infty}f(x)=b(=\mbox{Max}\{a,b\})$.
\textbf{Etude en} $\bf{-\infty}$. Pour tout réel $x$,

$$f(-x)=\left(\frac{a^{-x}+b^{-x}}{2}\right)^{-1/x}= \left(\frac{a^{x}+b^{x}}{2a^xb^x}\right)^{-1/x}=\frac{ab}{f(x)},$$
et donc,

$$\lim_{x\rightarrow -\infty}f(x)=\lim_{X\rightarrow +\infty}f(-X)=\lim_{X\rightarrow +\infty}\frac{ab}{f(X)}=\frac{ab}{b}=a\quad(=\mbox{Min}\{a,b\}).$$
\textbf{Dérivée et variations}.
$f$ est dérivable sur $]-\infty,0\cup]0,+\infty[$ en vertu de théorèmes généraux (et aussi en $0$ d'après l'étude faite plus haut), et pour $x\neq0$ (puisque $f>0$ sur $\Rr^*$),

$$\frac{f'(x)}{f(x)}=(\ln f)'(x)=\left(\frac{1}{x}\ln\left(\frac{a^x+b^x}{2}\right)\right)'(x)=-\frac{1}{x^2}\ln\left(\frac{a^x+b^x}{2}\right)+\frac{1}{x}
\frac{a^x\ln a+b^x\ln b}{a^x+b^x}.$$
$f'$ a le même signe que $(\ln f)'$ qui, elle-même, a le même signe que la fonction $g$ définie sur $\Rr$ par

$$\forall x\in\Rr,\;g(x)=-\ln\left(\frac{a^x+b^x}{2}\right)+x\frac{a^x\ln a+b^x\ln b}{a^x+b^x}.$$
$g$ est dérivable sur $\Rr$ et, pour $x$ réel,

\begin{align*}\ensuremath
g'(x)&=-\frac{a^x\ln a+b^x\ln b}{a^x+b^x}+\frac{a^x\ln a+b^x\ln b}{a^x+b^x}+x\frac{(a^x\ln^2a+b^x\ln^2b)(a^x+b^x)-(a^x\ln a+b^x\ln b)^2}{(a^x+b^x)^2}\\
 &=x\frac{(ab)^x(\ln a-\ln b)^2}{(a^x+b^x)^2}.
\end{align*}
$g'$ est donc strictement négative sur $]-\infty,0[$ et strictement positive sur $]0,+\infty[$. Par suite, $g$ est strictement décroissante sur $]-\infty,0]$ et strictement croissante sur $[0,+\infty[$. $g'$ admet donc un minimum global strict en $0$ et puisque $g(0)=0$, on en déduit que $g$ est strictement positive sur $\Rr^*$. De même, $f'$ est strictement positive sur $\Rr^*$. En tant compte de l'étude en $0$, on a montré que $f$ est dérivable sur $\Rr$ et que $f'$ est strictement positive sur $\Rr$. $f$ est donc strictement croissante sur $\Rr$.
Le \textbf{graphe de $f$} a l'allure suivante~:

$$\includegraphics{../images/img005428-1}$$

On peut noter que les inégalités $\lim_{x\rightarrow -\infty}f<f(-1)<f(0)<f(1)<\lim_{x\rightarrow +\infty}f$ fournissent~:

$$a<\frac{1}{\frac{1}{2}(\frac{1}{a}+\frac{1}{b})}<\sqrt{ab}<\frac{a+b}{2}<b.$$
\fincorrection
\correction{005429}
Quand $x$ tend vers $+\infty$,

$$\sqrt{x^2-3}=x\left(1-\frac{3}{x^2}\right)^{1/2}=x\left(1-\frac{3}{2x^2}+o(\frac{1}{x^2})\right)=x-\frac{3}{2x}+o\left(\frac{1}{x}\right),$$
et,

$$\sqrt[3]{8x^3+7x^2+1}=2x\left(1+\frac{7}{8x}+o\left(\frac{1}{x^2}\right)\right)^{1/3}=2x\left(1+\frac{7}{24x}-\frac{49}{576x^2}+o\left(\frac{1}{x^2}\right)\right)=2x+\frac{7}{12}-\frac{49}{288x}+o\left(\frac{1}{x}\right).$$
Donc,

$$f(x)\underset{x\rightarrow+\infty}{=}-x-\frac{7}{12}-\frac{383}{288x}+o\left(\frac{1}{x}\right).$$
La courbe représentative de $f$ admet donc en $+\infty$ une droite asymptote d'équation $y=-x-\frac{7}{12}$. De plus, le signe de $f(x)-\left(-x-\frac{7}{12}\right)$ est, au voisinage de $+\infty$, le signe de $-\frac{383}{288x}$. Donc la courbe représentative de $f$ est au-dessous de la droite d'équation $y=-x-\frac{7}{12}$ au voisinage de $+\infty$.
\fincorrection
\correction{005430}
$f$ est de classe $C^\infty$ sur son domaine $\Rr\setminus\{-1,1\}$ en tant que fraction rationelle et en particulier admet un développement limité à tout ordre en $0$. Pour tout entier naturel $n$, on a

$$f(x)\underset{x\rightarrow0}{=}x+x^3+...+x^{2n+1}+o(x^{2n+1}),$$
Par unicité des coefficients d'un développement limité et d'après la formule de \textsc{Taylor}-\textsc{Young}, on obtient

\begin{center}
\shadowbox{
$\forall n\in\Nn,\;f^{(2n)}(0)=0\;\mbox{et}\;f^{(2n+1)}(0)=(2n+1)!.$
}
\end{center}
Ensuite, pour $x\notin\{-1,1\}$, et $n$ entier naturel donné, 

$$f^{(n)}(x)=\frac{1}{2}\left(\frac{1}{1-x}-\frac{1}{1+x}\right)^{(n)}(x)=\frac{n!}{2}\left(\frac{1}{(1-x)^{n+1}}-
\frac{(-1)^{n}}{(1+x)^{n+1}}\right).$$
\fincorrection
\correction{005431}

\begin{enumerate}
 \item  $$\sqrt{x^2+3x+5}-x+1\underset{x\rightarrow-\infty}{\sim}-x-x=-2x,$$ 
et,

$$\sqrt{x^2+3x+5}-x+1=\frac{(x^2+3x+5)-(x-1)^2}{\sqrt{x^2+3x+5}+x-1}
\underset{x\rightarrow+\infty}{\sim}\frac{3x+2x}{x+x}=\frac{5}{2}.$$
 \item  $3x^2-6x\underset{x\rightarrow0}{\sim}-6x$ et $3x^2-6x\underset{x\rightarrow+\infty}{\sim}3x^2$. Ensuite, quand $x$ tend vers $1$, $3x^2-6x$ tend vers $-3\neq0$ et donc, $3x^2-6x\underset{x\rightarrow1}{\sim}-3$. Enfin, 
$3x^2-6x=3x(x-2)\underset{x\rightarrow2}{\sim}6(x-2)$.

\begin{center}
\shadowbox{
$3x^2-6x\underset{x\rightarrow0}{\sim}-6x$\qquad$3x^2-6x\underset{x\rightarrow+\infty}{\sim}3x^2$\qquad$3x^2-6x\underset{x\rightarrow0}{\sim}-3$\qquad$3x^2-6x\underset{x\rightarrow2}{\sim}6(x-2)$.
}
\end{center}
 \item  $$(x-x^2)\ln(\sin x)\underset{x\rightarrow0}{=}(x-x^2)\ln x+(x-x^2)\ln\left(1-\frac{x^2}{6}+o(x^2)\right)=x\ln x-x^2\ln x+o(x^2\ln x).$$ 
Ensuite,

$$\sin x\ln(x-x^2)\underset{x\rightarrow0}{=}\left(x-\frac{x^3}{6}+o(x^3)\right)(\ln x+\ln(1-x))=(x-\frac{x^3}{6}+o(x^3))(\ln x-x+o(x))=x\ln x+o(x^2\ln x).$$
Donc,

\begin{align*}\ensuremath
(\sin x)^{x-x^2}-(x-x^2)^{\sin x}&=e^{x\ln x}(e^{-x^2\ln x+o(x^2\ln x)}-e^{o(x^2\ln x)})=e^{x\ln x}(1-x^2\ln x-1+o(x^2\ln x))\\
 &=(1+o(1))(-x^2\ln x+o(x^2\ln x))\underset{x\rightarrow0}{\sim}-x^2\ln x.
\end{align*}

\begin{center}
\shadowbox{
$(\sin x)^{x-x^2}-(x-x^2)^{\sin x}\underset{x\rightarrow0}{\sim}-x^2\ln x$.
}
\end{center}
 \item  $\tanh x=\frac{1-e^{-2x}}{1+e^{-2x}}\underset{x\rightarrow+\infty}{=}=(1-e^{-2x})(1-e^{-2x}+o(e^{-2x}))=1-2e^{-2x}+o(e^{-2x})$, et donc $\tanh x\ln x=(1-2e^{-2x}+o(e^{-2x}))\ln x=\ln x+o(1)$. Par suite,

$$x^{\mbox{\scriptsize{th}}x}\underset{x\rightarrow+\infty}{\sim}e^{\ln x}=x.$$
 \item  \textbf{Tentative à l'ordre 3.}

$\tan(\sin x)\underset{x\rightarrow0}{=}\tan\left(x-\frac{x^3}{6}+o(x^3)\right)=\left(x-\frac{x^3}{6}\right)+\frac{1}{3}(x)^3+o(x^3)=x+\frac{x^3}{6}+o(x^3)$, et,

$\sin(\tan x)\underset{x\rightarrow0}{=}\sin\left(x+\frac{x^3}{3}+o(x^3)\right)=\left(x+\frac{x^3}{3}\right)-\frac{1}{6}(x)^3+o(x^3)=x+\frac{x^3}{6}+o(x^3)$.
Donc, $\tan(\sin x)-\sin(\tan x)\underset{x\rightarrow0}{=}o(x^3)$. L'ordre $3$ est insuffisant pour obtenir un équivalent.
\textbf{Tentative à l'ordre 5.}

\begin{align*}\ensuremath
\tan(\sin x)&\underset{x\rightarrow0}{=}\tan\left(x-\frac{x^3}{6}+\frac{x^5}{120}+o(x^5)\right)
=\left(x-\frac{x^3}{6}+\frac{x^5}{120}\right)+\frac{1}{3}\left(x-\frac{x^3}{6}\right)^3+\frac{2}{15}(x)^5+o(x^5)\\
 &=x+\frac{x^3}{6}+x^5\left(\frac{1}{120}-\frac{1}{6}+\frac{2}{15}\right)+o(x^5)=x+\frac{x^3}{6}-\frac{x^5}{40}+o(x^5),
\end{align*}
et,

\begin{align*}\ensuremath
\sin(\tan x)&\underset{x\rightarrow0}{=}\sin\left(x+\frac{x^3}{3}+\frac{2x^5}{15}+o(x^5)\right)
=\left(x+\frac{x^3}{3}+\frac{2x^5}{15}\right)-\frac{1}{6}\left(x+\frac{x^3}{3}\right)^3+\frac{1}{120}(x)^5+o(x^5)\\
 &=x+\frac{x^3}{6}+\left(\frac{2}{15}-\frac{1}{6}+\frac{1}{120}\right)x^5+o(x^5)=x+\frac{x^3}{6}-\frac{x^5}{40}+o(x^5).
\end{align*}
Donc, $\tan(\sin x)-\sin(\tan x)\underset{x\rightarrow0}{=}o(x^5)$. L'ordre $5$ est insuffisant pour obtenir un équivalent. Le contact entre les courbes représentatives des fonctions $x\mapsto\sin(\tan x)$ et $x\mapsto\tan(\sin x)$ est très fort.
\textbf{Tentative à l'ordre 7.}

\begin{align*}\ensuremath
\tan(\sin x)&\underset{x\rightarrow0}{=}\tan\left(x-\frac{x^3}{6}+\frac{x^5}{120}-\frac{x^7}{5040}+o(x^7)\right)\\
 &=\left(x-\frac{x^3}{6}+\frac{x^5}{120}-\frac{x^7}{5040}\right)+\frac{1}{3}\left(x-\frac{x^3}{6}+\frac{x^5}{120}\right)^3+\frac{2}{15}\left(x-\frac{x^3}{6}\right)^5+\frac{17}{315}x^7+o(x^7)\\
 &=x+\frac{x^3}{6}-\frac{x^5}{40}+\left(-\frac{1}{5040}+\frac{1}{3}\left(3\times\frac{1}{120}+3\times\frac{1}{36}\right)+\frac{2}{15}\left(-\frac{5}{6}\right)+\frac{17}{315}\right)x^7+o(x^7)\\
 &=x+\frac{x^3}{6}-\frac{x^5}{40}+\left(-\frac{1}{5040}+\frac{1}{120}+\frac{1}{36}-\frac{1}{9}+\frac{17}{315}\right)x^7+o(x^7),
\end{align*}
et,

\begin{align*}\ensuremath
\sin(\tan x)&\underset{x\rightarrow0}{=}\sin\left(x+\frac{x^3}{3}+\frac{2x^5}{15}+\frac{17}{x^7}{315}+o(x^7)\right)\\
 &=\left(x+\frac{x^3}{3}+\frac{2x^5}{15}+\frac{17x^7}{315}\right)-\frac{1}{6}\left(x+\frac{x^3}{3}+\frac{2x^5}{15}\right)^3+\frac{1}{120}\left(x+\frac{x^3}{3}\right)^5-\frac{1}{5040}(x)^7+o(x^7)\\
 &= x+\frac{x^3}{6}-\frac{x^5}{40}+\left(\frac{17}{315}
 -\frac{1}{6}\left(3\times\frac{2}{15}+3\times\frac{1}{9}\right)+\frac{1}{120}\times\frac{5}{3}-\frac{1}{5040}\right)x^7+o(x^7)\\
 &=x+\frac{x^3}{6}-\frac{x^5}{40}+\left(\frac{17}{315}-\frac{1}{15}-\frac{1}{18}+\frac{1}{72}-\frac{1}{5040}\right)x^7+o(x^7).
\end{align*}
Finalement,

\begin{align*}\ensuremath
\tan(\sin x)-\sin(\tan x)&\underset{x\rightarrow0}{=}\left(\frac{1}{120}+\frac{1}{36}
-\frac{1}{9}+\frac{1}{15}+\frac{1}{18}-\frac{1}{72}\right)x^7+o(x^7)
=\frac{x^7}{30}+o(x^7),
\end{align*}
et donc

\begin{center}
\shadowbox{
$\tan(\sin x)-\sin(\tan x)\underset{x\rightarrow0}{\sim}\frac{x^7}{30}$.
}
\end{center}
\end{enumerate}
\fincorrection
\correction{005432}
Pour $n\geq 5$, on a

$$u_n=1+\frac{1}{n}+\frac{1}{n(n-1)}+\frac{1}{n(n-1)(n-2)}+\frac{1}{n(n-1)(n-2)(n-3)}+\sum_{k=0}^{n-5}\frac{1}{n(n-1)...(k+1)}.$$
Ensuite, 

$$0\leq n^3\sum_{k=0}^{n-5}\frac{1}{n(n-1)...(k+1)}\leq n^3(n-4)\frac{1}{n(n-1)(n-2)(n-3)(n-4)}\underset{n\rightarrow+\infty}{\sim}\frac{1}{n}\underset{n\rightarrow+\infty}{\rightarrow}0.$$
Donc, $\sum_{k=0}^{n-5}\frac{1}{n(n-1)...(k+1)}\underset{n\rightarrow+\infty}{=}o\left(\frac{1}{n^3}\right)$ et de même $\frac{1}{n(n-1)(n-2)(n-3)}\underset{n\rightarrow+\infty}{=} o\left(\frac{1}{n^3}\right)$. Il reste

\begin{align*}\ensuremath
u_n&\underset{n\rightarrow+\infty}{=}1+\frac{1}{n}+\frac{1}{n^2}\left(1-\frac{1}{n}\right)^{-1}+\frac{1}{n^3}+o\left(\frac{1}{n^3}\right)=1+\frac{1}{n}
+\frac{1}{n^2}\left(1+\frac{1}{n}\right)+\frac{1}{n^3}+o\left(\frac{1}{n^3}\right)\\
 &=1+\frac{1}{n}+\frac{1}{n^2}+\frac{2}{n^3}+o\left(\frac{1}{n^3}\right).
\end{align*}

\begin{center}
\shadowbox{
$\frac{1}{n!}\sum_{k=0}^{n}k!\underset{n\rightarrow+\infty}{=}1+\frac{1}{n}+\frac{1}{n^2}+\frac{2}{n^3}+o\left(\frac{1}{n^3}\right)$.
}
\end{center}
\fincorrection
\correction{005433}
\begin{enumerate}
 \item 
\begin{align*}\ensuremath
\frac{1}{x(e^x-1)}-\frac{1}{x^2}&\underset{x\rightarrow0}{=}\frac{1}{x\left(x+\frac{x^2}{2}+\frac{x^3}{6}+\frac{x^4}{24}+\frac{x^5}{120}+o(x^5)\right)}-\frac{1}{x^2}=\frac{1}{x^2}\left(\left(1+\frac{x}{2}+\frac{x^2}{6}+\frac{x^3}{24}+\frac{x^4}{120}+o(x^4)\right)^{-1}-1\right)\\
 &=\frac{1}{x^2}\left(-\left(\frac{x}{2}+\frac{x^2}{6}+\frac{x^3}{24}+\frac{x^4}{120}\right)+\left(\frac{x}{2}+\frac{x^2}{6}+\frac{x^3}{24}\right)^2-\left(\frac{x}{2}+\frac{x^2}{6}\right)^3+\left(\frac{x}{2}\right)^4+o(x^4)\right)\\
 &=\frac{1}{x^2}\left(-\frac{x}{2}+x^2\left(-\frac{1}{6}+\frac{1}{4}\right)+x^3\left(-\frac{1}{24}+\frac{1}{6}-\frac{1}{8}\right)
 +x^4\left(-\frac{1}{120}+\left(\frac{1}{36}+\frac{1}{24}\right)-\frac{1}{8}+\frac{1}{16}\right)+o(x^4)\right)\\
 &=-\frac{1}{2x}+\frac{1}{12}-\frac{x^2}{720}+o(x^2).
\end{align*}

\begin{center}
\shadowbox{
$\frac{1}{x(e^x-1)}-\frac{1}{x^2}\underset{x\rightarrow0}{=}-\frac{1}{2x}+\frac{1}{12}-\frac{x^2}{720}+o(x^2)$.
}
\end{center}
 \item 
 
\begin{align*}\ensuremath
x\ln(x+1)-(x+1)\ln x&=x\left(\ln x+\ln\left(1+\frac{1}{x}\right)\right)-(x+1)\ln x\underset{x\rightarrow+\infty}{=}-\ln
 x+x\left(\frac{1}{x}-\frac{1}{2x^2}+\frac{1}{3x^3}-\frac{1}{4x^4}+o\left(\frac{1}{x^4}\right)\right)\\
 &=-\ln x+1-\frac{1}{2x}+\frac{1}{3x^2}-\frac{1}{4x^3}+o\left(\frac{1}{x^3}\right).
\end{align*}

\begin{center}
\shadowbox{
$x\ln(x+1)-(x+1)\ln x\underset{x\rightarrow+\infty}{=}-\ln x+1-\frac{1}{2x}+\frac{1}{3x^2}-\frac{1}{4x^3}+o\left(\frac{1}{x^3}\right)$.
}
\end{center}
\end{enumerate}
\fincorrection
\correction{005434}
\begin{enumerate}
 \item  

$$f_n(a)=\text{exp}\left(n\ln\left(1+\frac{a}{n}\right)\right)\underset{n\rightarrow+\infty}{=}\text{exp}\left(a-\frac{a^2}{2n}+o\left(\frac{1}{n}\right)\right)= e^a\left(1-\frac{a^2}{2n}+o\left(\frac{1}{n}\right)\right).$$
En remplaçant $a$ par $b$ ou $a+b$, on obtient

\begin{align*}\ensuremath
f_n(a+b)-f_n(a)f_n(b)&\underset{n\rightarrow+\infty}{=}e^{a+b}\left(1-\frac{(a+b)^2}{2n}\right)-e^a\left(1-\frac{a^2}{2n}\right)e^b\left(1-\frac{b^2}{2n}\right)+o\left(\frac{1}{n}\right)\\
 &=e^{a+b}\frac{-(a+b)^2+a^2+b^2}{2n}+o\left(\frac{1}{n}\right)=-\frac{ab\;e^{a+b}}{n}+o\left(\frac{1}{n}\right).
\end{align*}
Donc, si $ab\neq0$, $f_n(a+b)-f_n(a)f_n(b)\underset{n\rightarrow+\infty}{\sim}-\frac{ab\;e^{a+b}}{n}$. Si $ab=0$, il est clair que $f_n(a+b)-f_n(a)f_n(b)=0$.
 \item  $e^{-a}f_n(a)\underset{n\rightarrow+\infty}{=}\text{exp}\left(-a+\left(a-\frac{a^2}{2n}+\frac{a^3}{3n^2}\right)+o\left(\frac{1}{n^2}\right)\right)
=1+\left(-\frac{a^2}{2n}+\frac{a^3}{3n^2}\right)+\frac{1}{2}\left(-\frac{a^2}{2n}\right)^2+o\left(\frac{1}{n^2}\right)$, et donc

$$e^{-a}f_n(a)-1+\frac{a^2}{2n}\underset{n\rightarrow+\infty}{\sim}\left(\frac{a^3}{3}+\frac{a^4}{8}\right)\frac{1}{n^2}.$$
\end{enumerate}
\fincorrection
\correction{005437}
Pour $n$ entier naturel donné, posons $I_n=\left]-\frac{\pi}{2}+n\pi,\frac{\pi}{2}+n\pi\right[$.
\textbullet~Soit $n\in\Nn$. Pour $x\in I_n$, posons $f(x)=\tan x-x$. $f$ est dérivable sur $I_n$ et pour $x$ dans $I_n$, $f'(x)=\tan^2x$. Ainsi, $f$ est dérivable sur $I_n$ et $f'$ est strictement positive sur $I_n\setminus\{n\pi\}$. Donc $f$ est strictement croissante sur $I_n$.

\textbullet~Soit $n\in\Nn$. $f$ est continue et strictement croissante sur $I_n$ et réalise donc une bijection de $I_n$ sur $f(I_n)=\Rr$. En particulier, $\forall n\in\Nn,\;\exists!x_n\in I_n/\;f(x_n)=0$ (ou encore tel que $\tan x_n=x_n$.
\textbullet~On a $x_0=0$ puis pour $n\in\Nn^*$, $f(n\pi)=-n\pi<0$ et donc, $\forall n\in\Nn^*,\;x_n\in]n\pi,\frac{\pi}{2}+n\pi[$. En particulier,

\begin{center}
\shadowbox{
$x_n\underset{n\rightarrow+\infty}{=}n\pi+O(1)$.
}
\end{center}
\textbullet~Posons alors $y_n=x_n-n\pi$. $\forall n\in\Nn^*,\;y_n\in\left]0,\frac{\pi}{2}\right[$. De plus, $\tan(y_n)=\tan(x_n)=n\pi+y_n$ et donc, puisque $y_n\in\left]0,\frac{\pi}{2}\right[$,

$$\frac{\pi}{2}>y_n=\Arctan(y_n+n\pi)\geq\Arctan(n\pi).$$ 
Puisque $\Arctan(n\pi)$ tend vers $\frac{\pi}{2}$, on a $y_n=\frac{\pi}{2}+o(1)$ ou encore

\begin{center}
\shadowbox{
$x_n\underset{n\rightarrow+\infty}{=}n\pi+\frac{\pi}{2}+o(1)$.
}
\end{center}
\textbullet~Posons maintenant $z_n=y_n-\frac{\pi}{2}=x_n-n\pi-\frac{\pi}{2}$.
D'après ce qui précède, $\forall n\in\Nn^*,\;z_n\in\left]-\frac{\pi}{2},0\right[$ et d'autre part $z_n\underset{n\rightarrow+\infty}{=}o(1)$.
Ensuite, $\tan\left(z_n+\frac{\pi}{2}\right)=n\pi+\frac{\pi}{2}+z_n$ et donc $-\cotan(z_n)=n\pi+\frac{\pi}{2}+z_n\underset{n\rightarrow+\infty}{\sim}n\pi$. Puisque $z_n$ tend vers $0$, on en déduit que

\begin{center}
$-\frac{1}{z_n}\underset{n\rightarrow+\infty}{\sim}-\cotan(z_n)\underset{n\rightarrow+\infty}{\sim}n\pi$,
\end{center}
ou encore $z_n\underset{n\rightarrow+\infty}{=}-\frac{1}{n\pi}+o\left(\frac{1}{n}\right)$. Ainsi,

\begin{center}
\shadowbox{$x_n\underset{n\rightarrow+\infty}{=}n\pi+\frac{\pi}{2}-\frac{1}{n\pi}+o\left(\frac{1}{n}\right).$}
\end{center}
\textbullet~Posons enfin $t_n=z_n+\frac{1}{n\pi}=x_n-n\pi-\frac{\pi}{2}+\frac{1}{n\pi}$. On sait que $t_n=o\left(\frac{1}{n}\right)$ et que 

\begin{center}
$-\cotan\left(t_n-\frac{1}{n\pi}\right)=-\cotan(z_n)=n\pi+\frac{\pi}{2}+z_n=n\pi+\frac{\pi}{2}-\frac{1}{n\pi}+o(\frac{1}{n})$.
\end{center} Par suite,
 
$$-\tan\left(t_n-\frac{1}{n\pi}\right)=\frac{1}{n\pi}\left(1+\frac{1}{2n}+o(\frac{1}{n})\right)^{-1}=\frac{1}{n\pi}-\frac{1}{2n^2\pi}+o\left(\frac{1}{n^2}\right),$$
puis,

$$\frac{1}{n\pi}-t_n=\Arctan\left(\frac{1}{n\pi}-\frac{1}{2n^2\pi}+o(\frac{1}{n^2})\right)=\frac{1}{n\pi}-\frac{1}{2n^2\pi}+o\left(\frac{1}{n^2}\right),$$ 
et donc $t_n=\frac{1}{2n^2\pi}+o\left(\frac{1}{n^2}\right)$. Finalement, 

\begin{center}
\shadowbox{
$x_n\underset{n\rightarrow+\infty}{=}n\pi+\frac{\pi}{2}-\frac{1}{n\pi}+\frac{1}{2n^2\pi}+o\left(\frac{1}{n^2}\right).$
}
\end{center}

\fincorrection
\correction{005438}
\begin{enumerate}
 \item  Pour $x>0$, posons $f(x)=x+\ln x$. $f$ est continue sur $]0,+\infty[$, strictement croissante sur $]0,+\infty[$ en tant que somme de deux fonctions continues et strictement croissantes sur $]0,+\infty[$. $f$ réalise donc une bijection de $]0,+\infty[$ sur $f\left(]0,+\infty[\right)=\left]\lim_{x\rightarrow 0,\;x>0}f(x),\lim_{x\rightarrow +\infty}f(x)\right[=]-\infty,+\infty[=\Rr$. En particulier,

\begin{center}
\shadowbox{
$\forall k\in\Rr,\;\exists!x_k\in]0,+\infty[/\;f(x_k)=k.$
}
\end{center}
 \item  $f\left(\frac{k}{2}\right)=\frac{k}{2}+\ln\frac{k}{2}<k$ pour $k$ suffisament grand (car $k-(\frac{k}{2}+\ln\frac{k}{2})=\frac{k}{2}-\ln\frac{k}{2}\underset{k\rightarrow+\infty}{\rightarrow}+\infty$ d'après les théorèmes de croissances comparées). Donc, pour $k$ suffisament grand, $f\left(\frac{k}{2}\right)<f(x_k)$. Puisque $f$ est strictement croissante sur $]0,+\infty[$, on en déduit que $x_k>\frac{k}{2}$ pour $k$ suffisament grand et donc que $\lim_{k\rightarrow +\infty}x_k=+\infty$. Mais alors, $k=x_k+\ln x_k\sim x_k$ et donc, quand $k$ tend vers $+\infty$,

\begin{center}
\shadowbox{
$x_k\underset{k\rightarrow+\infty}{=}k+o(k).$
}
\end{center}
Posons $y_k=x_k-k$. On a $y_k=o(k)$ et de plus $y_k+\ln(y_k+k)=0$ ce qui s'écrit~:

$$y_k=-\ln(k+y_k)=-\ln(k+o(k))=-\ln k+\ln(1+o(1))=-\ln k+o(1).$$
Donc,

\begin{center}
\shadowbox{
$x_k\underset{k\rightarrow+\infty}{=}k-\ln k+o(1).$
}
\end{center}
Posons $z_k=y_k+\ln k=x_k-k+\ln k$. Alors, $z_k=o(1)$ et $-\ln k+z_k=-\ln(k-\ln k+z_k)$. Par suite,

$$z_k=\ln k-\ln(k-\ln k+o(1))=-\ln\left(1-\frac{\ln k}{k}+o\left(\frac{\ln k}{k}\right)\right)=\frac{\ln k}{k}+o\left(\frac{\ln k}{k}\right).$$

Finalement,

\begin{center}
\shadowbox{
$x_k\underset{k\rightarrow+\infty}{=}k-\ln k+\frac{\ln k}{k}+o\left(\frac{\ln k}{k}\right).$
}
\end{center}
\end{enumerate}
\fincorrection
\correction{005439}
\begin{enumerate}
 \item  $x^3\sin\frac{1}{x^2}\underset{x\rightarrow0}{=}O(x^3)$ et en particulier $x^3\sin\frac{1}{x}\underset{x\rightarrow0}{=}o(x^2)$. Donc, en tenant compte de $f(0)=1$,

\begin{center}
$f(x)\underset{x\rightarrow0}{=}1+x+x^2+o(x^2)$.
\end{center}
$f$ admet en $0$ un développement limité d'ordre $2$.
 \item  $f(x)\underset{x\rightarrow0}{=}1+x+o(x)$. Donc, $f$ admet en $0$ un développement limité d'ordre $1$. On en déduit que $f$ est continue et dérivable en $0$ avec $f(0)=f'(0)=1$. $f$ est d'autre part dérivable sur $\Rr^*$ en vertu de théorèmes généraux (et donc sur $\Rr$) et pour $x\neq 0$, $f'(x)=1+2x+3x^2\sin\frac{1}{x^2}-2\cos\frac{1}{x^2}$.
 \item  $f'$ est définie sur $\Rr$ mais n'a pas de limite en $0$. $f'$ n'admet donc même pas un développement limité d'ordre $0$ en $0$.
\end{enumerate}
\fincorrection
\correction{005440}
$\frac{1}{\sqrt{1-x^2}}\underset{x\rightarrow0}{=}1+\frac{x^2}{2}+\frac{3x^4}{8}+o(x^4)$, et donc 

$$\Arcsin x\underset{x\rightarrow0}{=}x+\frac{x^3}{6}+\frac{3x^5}{40}+o(x^5).$$
Puis,

$$\frac{1}{\Arcsin x}\underset{x\rightarrow0}{=}\frac{1}{x}\left(1+\frac{x^2}{6}+\frac{3x^4}{40}+o(x^4)\right)^{-1}=\frac{1}{x}\left(1-\frac{x^2}{6}-\frac{3x^4}{40}+\frac{x^4}{36}+o(x^4)\right)=\frac{1}{x}-\frac{x}{6}-\frac{17x^3}{360}+o(x^3),$$
et donc,

\begin{center}
\shadowbox{
$\frac{1}{x}-\frac{1}{\Arcsin x}\underset{x\rightarrow0}{=}\frac{x}{6}+\frac{17x^3}{360}+o(x^3).$
}
\end{center}
La fonction $f$ proposée se prolonge donc par continuité en $0$ en posant $f(0)=0$. Le prolongement est dérivable en $0$ et $f'(0)=\frac{1}{6}$. La courbe représentative de $f$ admet à l'origine une tangente d'équation $y=\frac{x}{6}$. Le signe de la différence $f(x)-\frac{x}{6}$ est, au voisinage de $0$, le signe de $\frac{17x^3}{360}$. La courbe représentative de $f$ admet donc à l'origine une tangente d'inflexion d'équation $y=\frac{x}{6}$.
\fincorrection
\correction{005441}
\begin{enumerate}
 \item  $\Arccos x\underset{x\rightarrow1^-}{=}o(1)$ (développement limité à l'ordre $0$). Mais la fonction $x\mapsto\Arccos x$ n'est pas dérivable en $1$ et n'admet donc pas en $1$ un développemement limité d'ordre $1$.

{2)} Puisque $\Arccos x\underset{x\rightarrow1^-}{=}o(1)$, 

$$\Arccos x\underset{x\rightarrow1^-}{\sim}\sin(\Arccos x)=\sqrt{1-x^2}=\sqrt{(1+x)(1-x)}\underset{x\rightarrow1^-}{\sim}\sqrt{2}\sqrt{1-x}.$$

\begin{center}
\shadowbox{
$\Arccos x\underset{x\rightarrow1^-}{\sim}\sqrt{2}\sqrt{1-x}.$
}
\end{center}
\end{enumerate}
\fincorrection
\correction{005442}
\begin{enumerate}
 \item  Quand $x$ tend vers $0$,
  
\begin{align*}\ensuremath
\frac{1}{(1-x)^2(1+x)}&=\frac{1}{4}\frac{1}{1-x}+\frac{1}{2}\frac{1}{(1-x)^2}+\frac{1}{4}\frac{1}{1+x}
\underset{x\rightarrow0}{=}\frac{1}{4}\left(\sum_{k=0}^{n}x^k+2\sum_{k=0}^{n}(k+1)x^k+\sum_{k=0}^{n}(-1)^kx^k\right)+o(x^n)\\
 &=\sum_{k=0}^{n}\frac{2k+3+(-1)^k}{4}x^k+o(x^n).
\end{align*}
 \item  On a aussi,

\begin{align*}\ensuremath
\frac{1}{(1-x)^2(1+x)}&=\frac{1}{(1-x)(1-x^2)}\underset{x\rightarrow0}{=}\left(\sum_{k=0}^{n}x^p\right)\left(\sum_{k=0}^{n}x^{2q}\right)+o(x^n)\\
 &=\sum_{k=0}^{n}\left(\sum_{p+2q=k}^{}1\right)x^k+o(x^n)=\sum_{k=0}^{n}a_kx^k+o(x^n).
\end{align*}
Par unicité des coefficients d'un développement limité, on a donc

\begin{center}
\shadowbox{
$\forall k\in\Nn$, $a_k=\frac{2k+3+(-1)^k}{4}$.
}
\end{center}
($a_k$ est le nombre de façons de payer $k$ euros en pièces de $1$ et $2$ euros).
\end{enumerate}
\fincorrection
\correction{006888}
\begin{enumerate}
  \item $\cos x \cdot \exp x$ (à l'ordre $3$).
  
Le dl de $\cos x$ à l'ordre $3$ est
$$\cos x = 1 - \frac{1}{2!} x^2 + \epsilon_1(x)x^3.$$

Le dl de $\exp x$ à l'ordre $3$ est
$$\exp x =1+x+\frac1{2!}x^2+\frac1{3!}x^3 + \epsilon_2(x)x^3.$$

Par convention toutes nos fonctions $\epsilon_i(x)$ vérifierons $\epsilon_i(x)\to 0$ lorsque $x\to0$.

\bigskip

On multiplie ces deux expressions  
\begin{align*}
\cos x \times \exp x 
  & =  \Big(1 - \frac{1}{2} x^2   + \epsilon_1(x)x^3\Big) \times \Big( 1+x+\frac1{2!}x^2+\frac1{3!}x^3 + \epsilon_2(x)x^3\Big) \\
  & =  1 \cdot \Big(1+x+\frac1{2!}x^2+\frac1{3!}x^3 + \epsilon_2(x)x^3 \Big)  \quad \text{on développe la ligne du dessus}\\  
  & \qquad  - \frac{1}{2} x^2 \cdot \Big( 1+x+\frac1{2!}x^2+\frac1{3!}x^3 + \epsilon_2(x)x^3 \Big) \\
  & \qquad  + \epsilon_1(x)x^3 \cdot \Big(1+x+\frac1{2!}x^2+\frac1{3!}x^3 + \epsilon_2(x)x^3 \Big) \\
\end{align*}

On va développer chacun de ces produits, par exemple pour le deuxième produit :
$$- \frac{1}{2!} x^2 \cdot \Big(  1+x+\frac1{2!}x^2+\frac1{3!}x^3 +  \epsilon_2(x)x^3\Big)
= - \frac{1}{2} x^2 - \frac{1}{2} x^3 - \frac14x^4  -\frac1{12}x^5 -\frac12x^2\cdot \epsilon_2(x)x^3.$$

Mais on cherche un dl à l'ordre $3$ donc tout terme en $x^4$, $x^5$ ou plus se met dans $\epsilon_3(x)x^3$,
y compris $x^2 \cdot \epsilon_2(x)x^3$ qui est un bien de la forme $\epsilon(x)x^3$.
Donc $$- \frac{1}{2} x^2 \cdot \Big(1+x+\frac1{2!}x^2+\frac1{3!}x^3 + \epsilon_2(x)x^3\Big)
= - \frac{1}{2} x^2 - \frac{1}{2} x^3  + \epsilon_3(x)x^3.$$

Pour le troisième produit on a
$$\epsilon_1(x)x^3 \cdot \Big(1+x+\frac1{2!}x^2+\frac1{3!}x^3 + \epsilon_2(x)x^3\Big) 
= \epsilon_1(x)x^3+x\epsilon_1(x)x^3+\cdots = \epsilon_4(x)x^3$$

On en arrive à :
\begin{align*}
\cos x \cdot \exp x 
  & =  \Big(1 - \frac{1}{2} x^2 + \epsilon_1(x)x^3 \Big) \times \Big( 1+x+\frac1{2!}x^2+\frac1{3!}x^3 + \epsilon_2(x)x^3\Big) \\
  & = 1+x+\frac1{2!}x^2+\frac1{3!}x^3 +  \epsilon_1(x)x^3\\
  & \qquad - \frac{1}{2} x^2- \frac{1}{2} x^3  + \epsilon_3(x)x^3 \\
  & \qquad  + \epsilon_4(x)x^3 \qquad \text{il ne reste plus qu'à regrouper les termes :}  \\    
  & =  1 + x + (\frac12-\frac12) x^2 + (\frac{1}{6}- \frac{1}{2})x^3 + \epsilon_5(x)x^3 \\
  & =  1 + x - \frac13 x^3 + \epsilon_5(x)x^3 \\
\end{align*}

Ainsi le dl de $\cos x \cdot \exp x$ en $0$ à l'ordre $3$ est :
$$\cos x \cdot \exp x = 1 + x - \frac13 x^3 + \epsilon_5(x)x^3.$$


\item $\left( \ln (1+x) \right)^2$ (à l'ordre $4$).

Il s'agit juste de multiplier le dl de $\ln(1+x)$ par lui-même.
En fait si l'on réfléchit un peu on s'aperçoit qu'un dl à l'ordre $3$ sera suffisant (car le terme constant est nul) :
$$\ln(1+x)=x-\frac12x^2+\frac13x^3+ \epsilon(x)x^3$$
 $\epsilon_5(x)\to 0$ lorsque $x\to0$.

\begin{align*}
\left( \ln (1+x) \right)^2 
  & = \ln (1+x)  \times \ln (1+x)  \\
  & = \left(x-\frac12x^2+\frac13x^3+ \epsilon(x)x^3\right) \times \left( x-\frac12x^2+\frac13x^3+ \epsilon(x)x^3\right) \\
  & = x \times \left( x-\frac12x^2+\frac13x^3+ \epsilon(x)x^3\right) \\
  & \qquad  -\frac12x^2\times \left( x-\frac12x^2+\frac13x^3+ \epsilon(x)x^3\right) \\
  & \qquad +\frac13x^3\times \left( x-\frac12x^2+\frac13x^3+ \epsilon(x)x^3\right) \\
  & \qquad + \epsilon(x)x^3\times \left( x-\frac12x^2+\frac13x^3+ \epsilon(x)x^3\right) \\
  & =  x^2-\frac12x^3+\frac13x^4+ \epsilon(x)x^4 \\
  & \qquad -\frac12x^3+\frac14x^4+ \epsilon_1(x)x^4 \\
  & \qquad +\frac13x^4 + \epsilon_2(x)x^4 \\
  & \qquad + \epsilon_3(x)x^4 \\  
  & =  x^2-x^3+\frac{11}{12}x^4+ \epsilon_4(x)x^4 \\
\end{align*}


\item $\displaystyle{\frac{\sh x-x}{x^3}}$ (à l'ordre $6$).

Pour le dl de $\displaystyle{\frac{\sh x-x}{x^3}}$ on commence par faire un dl du numérateur.
Tout d'abord :
$$\sh x = x+\frac{1}{3!}x^3+\frac{1}{5!}x^5+\frac{1}{7!}x^7+\frac{1}{9!}x^9 +\epsilon(x) x^9$$
donc 
$$\sh x - x = \frac{1}{3!}x^3+\frac{1}{5!}x^5+\frac{1}{7!}x^7+\frac{1}{9!}x^9 +\epsilon(x) x^9.$$

Il ne reste plus qu'à diviser par $x^3$ :
$$\frac{\sh x-x}{x^3} = \frac{\frac{1}{3!}x^3+\frac{1}{5!}x^5+\frac{1}{7!}x^7+\frac{1}{9!}x^9 +\epsilon(x) x^9 }{x^3} 
= \frac{1}{3!}+\frac{1}{5!}x^2+\frac{1}{7!}x^4+\frac{1}{9!}x^6 +\epsilon(x) x^6$$

Remarquez que nous avons commencé par calculer un dl du numérateur à l'ordre $9$,
pour obtenir après division un dl à l'ordre $6$.


  \item $\exp\big(\sin(x)\big)$ (à l'ordre $4$).

On sait $\sin x= x -\frac{1}{3!}x^3 + o(x^4)$
et $\exp(u)=1+u+\frac1{2!} u^2+\frac{1}{3!}u^3+\frac{1}{4!}u^4+o(u^4)$.


On note désormais toute fonction $\epsilon(x)x^n$ (où $\epsilon(x)\to 0$ lorsque $x\to0$) par $o(x^n)$.
Cela évite les multiples expressions $\epsilon_i(x)x^n$.


On substitue $u=\sin(x)$, il faut donc calculer $u, u^2, u^3$ et $u^4$ : 
$$u = \sin x= x -\frac{1}{3!}x^3 + o(x^4)$$
$$u^2 = \big( x -\frac{1}{3!}x^3 + o(x^4)\big)^2 = x^2-\frac13 x^4 + o(x^4)$$
$$u^3 = \big( x -\frac{1}{3!}x^3 + o(x^4)\big)^3 = x^3 + o(x^4)$$
$$u^3 = x^4 + o(x^4) \quad \text{ et } \quad o(u^4)=o(x^4)$$

Pour obtenir :
\begin{align*}
  \exp(\sin(x)) 
    & =  1+ x -\frac{1}{3!}x^3 + o(x^4)\\
    &  \qquad   + \frac1{2!}\big(x^2-\frac13 x^4 + o(x^4)\big) \\
    &  \qquad   + \frac1{3!}\big(x^3 + o(x^4)\big) \\
    &  \qquad   + \frac1{4!}\big(x^4 + o(x^4)\big) \\    
    & \qquad + o(x^4) \\
    & = 1+x + \frac12 x^2 - \frac18 x^4 + o(x^4).
\end{align*}



  \item  $\sin^6(x)$ (à l'ordre $9$).

On sait $\sin (x)= x -\frac{1}{3!}x^3 + o(x^4)$.



Si l'on voulait calculer un dl de $\sin^2(x)$ à l'ordre $5$ on écrirait :
$$\sin^2 (x)  =  \big(x -\frac{1}{3!}x^3 + o(x^4)\big)^2 =  
\big(x -\frac{1}{3!}x^3 + o(x^4)\big) \times  \big(x -\frac{1}{3!}x^3 + o(x^4)\big) 
= x^2 -2\frac{1}{3!}x^4 + o(x^5).$$
En effet tous les autres termes sont dans $o(x^5)$.


Le principe est le même pour $\sin^6(x)$:
$$\sin^6 (x)  =  \big(x -\frac{1}{3!}x^3 + o(x^4)\big)^6 =  
\big(x -\frac{1}{3!}x^3 + o(x^4) \big) \times  \big(x -\frac{1}{3!}x^3 + o(x^4) \big) 
\times  \big(x -\frac{1}{3!}x^3 + o(x^4) \big) \times \cdots$$
Lorsque l'on développe ce produit en commençant par les termes de plus petits degrés on obtient 
$$\sin^6 (x)  =  x^6 + 6 \cdot x^5 \cdot (-\frac1{3!} x^3) + o(x^9) = x^6-x^8 + o(x^9)$$



  \item $\ln \big(\cos(x)\big)$ (à l'ordre $6$).

Le dl de $\cos x$ à l'ordre $6$ est
$$\cos x = 1 - \frac{1}{2!} x^2 + \frac{1}{4!}x^4 - \frac{1}{6!}x^6 + o(x^6).$$
Le dl de $\ln(1+u)$ à l'ordre $6$ est
$\ln(1+u)=u-\frac12u^2+\frac13u^3-\frac14u^4+\frac15u^5-\frac16u^6+o(u^6)$.

On pose $u= - \frac{1}{2!} x^2 + \frac{1}{4!}x^4 - \frac{1}{6!}x^6 + o(x^6)$ de sorte que
$$\ln(\cos x) = \ln (1+u)=u-\frac12u^2+\frac13u^3-\frac14u^4+\frac15u^5-\frac16u^6+o(u^6).$$

Il ne reste qu'à développer les $u^k$, ce qui n'est pas si dur que cela si les calculs sont bien menés et 
les puissances trop grandes écartées.

Tout d'abord :
\begin{align*}
u^2
  & = \left(- \frac{1}{2!} x^2 + \frac{1}{4!}x^4 - \frac{1}{6!}x^6 + o(x^6)\right)^2 \\
  & = \left(- \frac{1}{2!} x^2 + \frac{1}{4!}x^4 \right)^2 + o(x^6) \\
  & = \left(- \frac{1}{2!} x^2\right)^2 + 2 \left(- \frac{1}{2!} x^2\right) \left(\frac{1}{4!}x^4 \right) + o(x^6) \\
  & = \frac14 x^4 - \frac1{24} x^6 + o(x^6) \\
\end{align*}

Ensuite :
\begin{align*}
u^3 
  & = \left(- \frac{1}{2!} x^2 + \frac{1}{4!}x^4 - \frac{1}{6!}x^6 + o(x^6)\right)^3 \\
  & = \left(- \frac{1}{2!} x^2 \right)^3 + o(x^6) \\
  & =  -\frac18 x^6 + o(x^6) \\
\end{align*}
En effet lorsque l'on développe $u^3$ le terme $(x^2)^6$ est le seul terme dont l'exposant est $\le 6$.

Enfin les autres termes $u^4$, $u^5$, $u^6$ sont tous des $o(x^6)$. Et en fait développer $\ln(1+u)$ à l'ordre $3$ est suffisant.

Il ne reste plus qu'à rassembler :
\begin{align*}
\ln(\cos x) 
  & = \ln (1+u) \\
  & = u-\frac12u^2+\frac13u^3+o(u^3) \\
  & = \left(- \frac{1}{2!} x^2 + \frac{1}{4!}x^4 - \frac{1}{6!}x^6 + o(x^6)\right)\\
  & \qquad   -\frac12 \left(\frac14 x^4 - \frac{1}{24} x^6 + o(x^6)\right) \\
  & \qquad   +\frac13 \left(-\frac18 x^6 + o(x^6)\right)\\
  & = - \frac{1}{2} x^2 -\frac{1}{12}x^4 -\frac{1}{45}x^6  + o(x^6)\\
\end{align*}

  \item $\displaystyle{\frac{1}{\cos x}}$ à l'ordre $4$.

Le dl de $\cos x$ à l'ordre $4$ est
$$\cos x = 1 - \frac{1}{2!} x^2 + \frac{1}{4!}x^4 + o(x^4).$$
Le dl de $\frac{1}{1+u}$ à l'ordre $2$ (qui sera suffisant ici) est
$\frac{1}{1+u}=1-u+u^2+o(u^2)$.

On pose $u=- \frac{1}{2!} x^2 + \frac{1}{4!}x^4 + o(x^4)$ et on a $u^2 = \frac14 x^4 + o(x^4)$.

\begin{align*}
\frac{1}{\cos x}
  & =  \frac{1}{1+u} \\
  & =  1-u+u^2+o(u^2) \\
  & = 1 -\big(- \frac{1}{2!} x^2 + \frac{1}{4!}x^4 + o(x^4)\big)+\big(- \frac{1}{2!} x^2 + \frac{1}{4!}x^4 + o(x^4)\big)^2 +o(x^4)  \\
  & = 1+\frac{1}{2}x^2+\frac{5}{24}x^4 + o(x^4) \\
\end{align*}


  \item $\tan x$ (à l'ordre $5$ (ou $7$ pour les plus courageux)).

Pour ceux qui souhaitent seulement un dl à l'ordre $5$ de $\tan x =\sin x \times \frac{1}{\cos x}$ alors
il faut multiplier le dl de $\sin x$ à l'ordre $5$ par le dl de $\frac{1}{\cos x}$ à l'ordre $4$ (voir question précédente).


Si l'on veut un dl de $\tan x$ à l'ordre $7$ il faut d'abord refaire le dl $\frac{1}{\cos x}$ mais cette fois à l'ordre $6$ :
$$\frac{1}{\cos x}=1+\frac{1}{2}x^2+\frac{5}{24}x^4 +\frac{61}{720}x^6 + o(x^6)$$

Le dl à l'ordre $7$ de $\sin x$ étant :
$$\sin x = x -\frac{1}{3!}x^3 +\frac{1}{5!}x^5 - \frac{1}{7!}x^7 +o(x^7)$$

Comme  $\tan x = \sin x \times \frac{1}{\cos x}$, il ne reste donc qu'à multiplier les deux dl
pour obtenir après calculs :
$$\tan x = x + \frac{x^3}{3} + \frac{2x^5}{15} + \frac{17x^7}{315} + o(x^7)$$



  \item $(1+x)^{\frac{1}{1+x}}$ (à l'ordre $3$).

Si l'on pense bien à écrire $(1+x)^{\frac{1}{1+x}}= \exp\left( \frac{1}{1+x} \ln(1+x) \right)$
alors c'est juste des calculs utilisant les dl à l'ordre $3$ de $\ln(1+x)$, $\frac{1}{1+x}$ et 
$\exp x$.

On trouve 
$$(1+x)^{\frac{1}{1+x}} = 1+x-x^2 + \frac{x^3}{2} + o(x^3).$$


  \item $\arcsin \left ( \ln(1+x^2) \right )$ (à l'ordre $6$).

Tout d'abord $\ln(1+x^2)=x^2-\frac{x^4}{2}+\frac{x^6}{3}+o(x^6)$.
Et $\arcsin u =  u + \frac{u^3}{6} + o(u^3)$.
Donc en posant $u=x^2-\frac{x^4}{2}+\frac{x^6}{3}+o(x^6)$ on a :
\begin{align*}
\arcsin \left ( \ln(1+x^2) \right ) 
  & =  \arcsin \left ( x^2-\frac{x^4}{2}+\frac{x^6}{3}+o(x^6) \right )  \\
  & = \arcsin u \\
  & = u + \frac{1}{6}u^3 + o(u^3) \\
  & = \left ( x^2-\frac{x^4}{2}+\frac{x^6}{3} \right ) + \frac{1}{6}\left (x^2-\frac{x^4}{2}+\frac{x^6}{3} \right )^3 + o(x^6) \\
  & = \left ( x^2-\frac{x^4}{2}+\frac{x^6}{3}\right ) + \frac{x^6}{6} + o(x^6) \\
  & = x^2-\frac{x^4}{2}+\frac{x^6}{2}+o(x^6) \\
\end{align*}


\end{enumerate}
\fincorrection
\correction{001247}
\begin{enumerate}
  \item 
On a 
$$e^{x^2} = 1+x^2+\frac{x^4}{2!} + o(x^4) \quad \text{ et } \quad  \cos x= 1-\frac{x^2}{2!}+\frac{x^4}{4!} + o(x^4)$$

On s'aperçoit qu'en fait un dl à l'ordre $2$ suffit :
$$e^{x^2}-\cos x 
= \big(1+x^2 + o(x^2) \big) - \big(1-\frac{x^2}{2}+ o(x^2) \big)
= \frac32 x^2 + o(x^2)$$
Ainsi $\frac{e^{x^2}-\cos x}{x^2} = \frac32 + o(1)$ (où $o(1)$ désigne une fonction qui tend vers $0$)
 et donc
$$\lim_{x\rightarrow 0}\frac{e^{x^2}-\cos x}{x^2} = \frac32$$


  \item 
On sait que 
$$\ln(1+x)=x-\frac{x^2}{2}+\frac{x^3}{3}+o(x^3) 
\quad \text{ et } \quad \sin x = x-\frac{x^3}{3!} + o(x^3).$$

Les dl sont distincts dès le terme de degré $2$ donc un dl à l'ordre $2$ suffit :
$$\ln (1+x)-\sin x = \big(x - \frac{x^2}{2} + o(x^2) \big) - \big(x + o(x^2) \big) = -\frac{x^2}{2} + o(x^2)$$
donc 
$$\frac{\ln (1+x)-\sin x}{x} = -\frac{x}{2} + o(x)$$
et ainsi 
$$\lim_{x\rightarrow 0}\frac{\ln (1+x)-\sin x}{x} = 0.$$


  \item 
Sachant
$$\cos x= 1-\frac{x^2}{2!}+\frac{x^4}{4!}+ o(x^4)$$
et 
$$\sqrt{1-x^2} = 1-\frac12x^2-\frac18x^4 + o(x^4)$$ alors

\begin{align*}
\frac{\cos x-\sqrt{1-x^2}}{x^4}
 & =\frac{\big(1-\frac{x^2}{2}+\frac{x^4}{24}+ o(x^4)\big)-\big(1-\frac12x^2-\frac18x^4 + o(x^4)\big)}{x^4} \\
 & = \frac{\frac 16 x^4 + o(x^4)}{x^4}  \\
 &= \frac16+o(1) \\
\end{align*}
Ainsi 
$$\lim_{x\rightarrow 0}\frac{\cos x-\sqrt{1-x^2}}{x^4}=\frac16$$
\end{enumerate}
\fincorrection
\correction{001248}
$$\lim _{x\rightarrow 0}{\frac {{e^{x}}-\cos(x)-x}{{x}^{2}}}=1 .$$

$$\lim _{x\rightarrow 0}{\frac {{x}^{3}\arctan(x)-{x}^{4}}{\cos({x}^{2})
-1}}=0 .$$
\fincorrection
\correction{001249}

Commençons en $x=0$, le dl de $f(x)=\ln(1+x+x^2)$ à l'ordre $2$
est 
$$\ln(1+x+x^2)
=(x+x^2)-\frac{(x+x^2)^2}{2} + o(x^2)
= x + \frac12 x^2 + o(x^2)$$
Par identification avec
$f(x)= f(0)+f'(0)x+f''(0)\frac{x^2}{2!}+o(x^2)$
cela entraîne donc $f(0)=0$, $f'(0)=1$ (et $f''(0)=1$).
L'équation de la tangente est donc
$y=f'(0)(x-0)+f(0)$ donc $y=x$.

La position par rapport à la tangente correspond à l'étude du signe de 
$f(x)-y(x)$ où $y(x)$ est l'équation de la tangente.
$$f(x)-y(x)=x + \frac12 x^2 + o(x^2) \  - \  x = \frac12 x^2 + o(x^2).$$

Ainsi pour $x$ suffisamment proche de $0$, $f(x)-y(x)$ est du signe de $\frac12 x^2$
et est donc positif. Ainsi dans un voisinage de $0$ la courbe de $f$ est au-dessus de la tangente
en $0$.

\bigskip

Même étude en $x=1$.

Il s'agit donc de faire le dl de $f(x)$ en $x=1$.
On pose $x=1+h$ (de sorte que $h=x-1$ est proche de $0$) :
\begin{align*}
f(x)=\ln(1+x+x^2) 
  & = \ln \big(1+ (1+h)+(1+h)^2\big)  \\
  & = \ln \big(3 + 3h + h^2 \big)  \\
  & = \ln \left(3 \big(1 + h + \frac{h^2}{3} \big)\right) \\
  & = \ln 3 + \ln \big(1 + h + \frac{h^2}{3} \big) \\
  & = \ln 3 +  \big( h + \frac{h^2}{3} \big) - \frac{\big( h + \frac{h^2}{3} \big)^2}{2} + o\big((h + \frac{h^2}{3})^2 \big) \\ 
  & = \ln 3 + h  + \frac{h^2}{3} -\frac{h^2}{2}  + o(h^2) \\
  & = \ln 3 + h -\frac16 h^2 + o(h^2) \\
  & = \ln 3 + (x-1) - \frac16 (x-1)^2 + o((x-1)^2) \\
\end{align*}

La tangente en $x=1$ est d'équation $y=f'(1)(x-1)+f(1)$ et est donc donnée par le dl à l'ordre $1$ : c'est  $y = (x-1) + \ln 3$.
Et la différence 
$f(x)-\big(\ln 3 + (x-1)\big) = - \frac16 (x-1)^2 + o((x-1)^2)$ est négative pour $x$ proche de $1$.
Donc, dans un voisinage de $1$, le graphe de $f$ est en-dessous de la tangente en $x=1$.
\fincorrection
\nocorrection
\nocorrection
\nocorrection
\nocorrection
\nocorrection
\nocorrection
\nocorrection
\nocorrection
\nocorrection
\nocorrection
\nocorrection
\nocorrection
\correction{001262}
\begin{enumerate}
\item La fonction
$g$ est d\'efinie en $x$ sauf si $\sin (x)=0$ ou $x=0$.
Son domaine de d\'efinition est donc $\Rr -\{ k\pi , k\in \Zz \} .$


\item
On peut prolonger $g$ en une fonction continue en $0$ si et seulement si elle y
admet une limite.
Elle est d\'erivable en ce point si et seulement si elle y admet un 
d\'eveloppement limit\'e \`a l'ordre
$1$. Toutefois, comme l'\'enonc\'e demande la position du graphe de 
$g$ par rapport
\`a sa tangente en $0$, nous allons calculer directement le 
d\'eveloppement limit\'e \`a l'ordre
$2$ de $g$ en $0$.



Le d\'eveloppement limit\'e en $0$ \`a l'ordre $5$ de $\arctan x = x 
-\dfrac{x^3}{3} +\dfrac{x^5}{5}+x^5\epsilon_1 (x).$

Or $\sin x =x-\dfrac{x^3}{3!} +\dfrac{x^5}{5!}+x^5\epsilon_2 (x).$ Donc
$\sin ^3 x =x^3-\dfrac{x^5}{2} + \dfrac{13x^7}{120}   +x^7\epsilon_3 (x)$
et $\dfrac{1}{\sin ^3 x}=\dfrac{1}{x^3} 
(1+\dfrac{x^2}{2}+\dfrac{9x^4}{40} +x^4\epsilon_4 (x)).$
On en d\'eduit que :

$ \dfrac{\arctan x}{(\sin x)^3}-\dfrac {1}{x^2}=\dfrac 
{1}{x^3}(x+\dfrac{x^3}{6}+\dfrac{31x^5}{120}+x^5\epsilon_5 (x))
-\dfrac {1}{x^2}=\dfrac{1}{6}+\dfrac{31x^2}{120}+x^2\epsilon_5 (x)
.$

Ainsi on peut prolonger $g$ en une fonction continue en $0$ en posant 
$g(0)=\dfrac{1}{6}$.
La fonction obtenue est d\'erivable en $0$ et sa d\'eriv\'ee est 
nulle. La tangente en $0$ \`a son graphe est la droite d'\'equation
$y=\dfrac{1}{6}$. Enfin le graphe de $g$ est au-dessus de cette 
droite au voisinage de $0$.
\end{enumerate}
\fincorrection
\nocorrection
\nocorrection
\correction{001265}
\begin{enumerate}
  \item 
  \begin{enumerate}
    \item La première limite n'est pas une forme indéterminée, en effet 
$$\lim_{x \rightarrow +\infty} \sqrt {x^2+3x+2}  = +\infty \quad \text{ et } \quad \lim_{x \rightarrow +\infty} x = +\infty$$
donc 
$$\lim_{x \rightarrow +\infty} \sqrt {x^2+3x+2} +x = + \infty$$
    \item Lorsque $x\to -\infty$ la situation est tout autre car
$$\lim_{x \rightarrow -\infty} \sqrt {x^2+3x+2} = +\infty \quad \text{ et } \quad \lim_{x \rightarrow -\infty} x = -\infty$$
donc $\sqrt {x^2+3x+2} +x$ est une forme indéterminée !

Calculons un développement limité à l'ordre $1$ en $-\infty$ en faisant très attention au signe (car par exemple $|x|=-x$):
\begin{align*}
\sqrt{x^2+3x+2} +x 
  & = |x| \left( \sqrt{1+\frac{3}{x}+\frac{2}{x^2}} -1 \right) \\
  & = |x| \left( 1+\frac12\left(\frac{3}{x}+\frac{2}{x^2}\right)+o(\frac1x) \ -1)  \right) \\
  & = |x| \left( \frac12\frac{3}{x}+o(\frac1x) \right) \\
  & = -\frac32 + o(1) \\
\end{align*}
Et donc  
$$\lim_{x \rightarrow -\infty} \sqrt {x^2+3x+2} +x = -\frac32$$

  \end{enumerate}
  \item 
Nous utiliserons que 

\begin{align*}
(\Arctan x)^{\frac{1}{x^2}}
  & = \exp\left(\frac{1}{x^2} \ln \left( \Arctan x \right) \right) \\
  & = \exp\left(\frac{1}{x^2} \ln \big( x + o(x)\big) \right) \quad \text{ car } \Arctan x = x + o(x) \\
\end{align*}
Mais lorsque $x\to 0^+$ on sait que $\ln (x+o(x)) \to -\infty$, $x^2 \to 0$ 
donc :
$$\lim_{x\to0^+} \frac{\ln (x+o(x))}{x^2} = -\infty$$

Composé avec l'exponentielle on trouve :
$$\lim_{x \rightarrow 0^+}(\Arctan x )^{\frac{1}{x^2}}=0$$



  \item 
Effectuons le dl à l'ordre $2$ :
comme 
$$(1+x)^\alpha = 1+\alpha x + \frac{\alpha(\alpha-1)}{2}x^2+o(x^2)$$
alors
$$(1+3x)^{\frac{1}{3}} = 1+x-x^2 + o(x^2).$$

$$\sin x = x + o(x^2) \quad \text{ et } \cos x = 1 - \frac{x^2}{2!} + o(x^2).$$
Ainsi

\begin{align*}
\frac{(1+3x)^{\frac{1}{3}}-1-\sin x}{1-\cos x}  
 & = \frac{-x^2+o(x^2)}{\frac12x^2 + o(x^2)} \\
 & = \frac{-1+o(1)}{\frac12 + o(1)} \quad \text{ après factorisation par } x^2 \\
 &= -2 + o(1) \\
\end{align*}
Donc 

$$\lim_{x \rightarrow 0} \frac{(1+3x)^{\frac{1}{3}}-1-\sin x}{1-\cos x}=-2$$

\end{enumerate}
\fincorrection
\nocorrection
\nocorrection
\nocorrection
\nocorrection
\nocorrection
\correction{004024}
$\sum_{k = 1}^n \frac {a_k}k = 0 \Rightarrow  L = \frac 13 \sum_{k = 1}^n a_k$.
\fincorrection
\correction{004025}
$v_p = \left(\frac{2^{1+1/p}-1}{1+1/p}\right)^p$,
         $w_n = \exp\left({\frac1n\sum_{k=0}^{n-1}\ln\Bigl(1+\frac kn\Bigr)}\right)$,
	   $v=w = \frac 4e$.
\fincorrection
\correction{004026}
$\sum_{k=1}^n f\Bigl(\frac k{n^2}\Bigr) \to \frac{f'(0)}2$ lorsque $n\to\infty$.
	     $\Bigl($Utiliser $|f(x)-xf'(0)| \le \frac12\sup\limits_{0\le t\le 1}|f''(t)|$
	     pour $0\le x \le 1\Bigr)$
\fincorrection
\correction{004027}
\begin{enumerate}
  \item $y = 1 + \frac x2 - \frac{x^3}8$.
  \item $y = -\frac x6 + \frac{7x^3}{360}$.
  \item $y = -\frac x6 - \frac{17x^3}{360}$.
  \item $y = e^3(1 - 4x + 16x^2)$.
  \item $y = -1 + \frac x3 - \frac{x^3}{45}$.
  \item $y = \sqrt2\left(1 + \frac x8 - \frac{5x^2}{128}\right)$.
\end{enumerate}
 \fincorrection
\correction{004028}
$h \le k \le g \le f$.
\fincorrection
\correction{004029}
\begin{enumerate}
  \item $y = x + \frac 12 - \frac 1{8x}$.
  \item $y = x + \frac12 + \frac3{8x}$.
  \item $y = 2x - \frac4{3x}$.
  \item $y = \frac{\pi x}4 + \frac\pi4+1 - \frac1{3x^2}$.
  \item $y = \frac {\pi x}2 + \frac \pi2-1 + \frac {\pi/4-1}x$.
  \item $y = \frac{\pi x}2 + \pi-1 + \frac{5\pi/4-2}x$.
  \item $y = x + \frac 12 - \frac 9{8x}$.
\end{enumerate}
 \fincorrection
\correction{004037}
\begin{enumerate}
  \item Notons $I_n$ l'intervalle  $\left]n\pi-\frac \pi2, n\pi+\frac \pi2\right[$.
Alors sur chaque $I_n$ la fonction définie par $f(x) = \tan x-x$ est un fonction continue et dérivable.
De plus $f'(x) = 1+\tan^2x -1 = \tan ^2 x$. La dérivée est strictement positive sauf en un point où elle est nulle et ainsi 
la fonction $f$ est strictement croissante sur $I_n$. La limite à gauche est $-\infty$ et
la limite à droite est $+\infty$. Par le théorème des valeurs intermédiaires il existe un unique $x_n \in I_n$
tel que $f(x_n)=0$ c'est-à-dire $\tan x_n=x_n$.

  \item $x \mapsto \arctan x$ est la bijection réciproque de la restriction de la tangente  
$\tan_| : ]-\frac\pi2,+\frac\pi2[ \to ]-\infty,+\infty[$. Sur ces intervalles on a bien
$\tan x = y \iff x = \arctan y$. Mais si $y \notin ]-\frac\pi2,+\frac\pi2[$ il faut d'abord se ramener dans l'intervalle
 $]-\frac\pi2,+\frac\pi2[$.

Ainsi $x_n \in I_n$ donc $x_n-n\pi \in ]-\frac\pi2,+\frac\pi2[$.
Maintenant $x_n = \tan (x_n)=\tan(x_n-n\pi)$. 

Donc $\arctan x_n = \arctan \big(\tan(x_n-n\pi) \big) = x_n - n\pi$.
Ainsi $$x_n = \arctan x_n + n\pi.$$


L'erreur classique est de penser que $\arctan (\tan x) = x$. Ce qui n'est vrai que pour 
$x \in ]-\frac\pi2,+\frac\pi2[$ !

  \item Comme $x_n \in I_n$ alors $x_n \to +\infty$ lorsque $n\to +\infty$.

On sait par ailleurs que pour $x>0$ on a $\arctan x + \arctan \frac 1x = \frac\pi 2$.
Ainsi $\arctan x_n = \frac \pi 2 - \arctan \frac{1}{x_n}$

Lorsque $n$ tend vers $+\infty$ alors $\frac{1}{x_n}\to 0$ donc 
$\arctan \frac{1}{x_n} \to 0$.

Ainsi 
$$x_n = n\pi + \arctan x_n = n\pi +\frac \pi 2 - \arctan \frac{1}{x_n} = n\pi +\frac \pi 2 + o(1).$$

  \item 

On va utiliser le dl obtenu précédemment pour obtenir un dl à un ordre plus grand :

\begin{align*}
x_n 
 & = n\pi + \arctan x_n \\
 & = n\pi +\frac \pi 2 - \arctan \frac{1}{x_n} \\
 & = n\pi +\frac \pi 2 - \arctan \frac{1}{n\pi +\frac \pi 2 + o(1)} \\
 & = n\pi +\frac \pi 2 - \frac{1}{n\pi +\frac \pi 2 + o(1)} +o(\frac1{n^2}) \qquad \text{ car } \arctan u = u + o(u^2) \text{ en } u=0\\
 & = n\pi +\frac \pi 2 - \frac{1}{n\pi} \frac{1}{1+\frac{1}{2n} + o(\frac1n)}+o(\frac1{n^2}) \\
 & = n\pi +\frac \pi 2 - \frac{1}{n\pi}\big(1-\frac{1}{2n} + o(\frac1n) \big) +o(\frac1{n^2})\\
 & = n\pi +\frac \pi 2 -\frac{1}{n\pi} + \frac{1}{2\pi n^2} + o(\frac1{n^2}) \\
\end{align*}

Ainsi en $+\infty$ on a le développement :
$$x_n  = n\pi +\frac \pi 2 -\frac{1}{n\pi} + \frac{1}{2\pi n^2} + o(\frac1{n^2}).$$
\end{enumerate}
 \fincorrection
\correction{004038}
\begin{enumerate}
  \item $f_n'(x) = 0 \iff \mathrm{cotan} x = nx$.
  \item $0$.
  \item $x_n\tan x_n = \frac 1n$.
  \item $\ln\left(\frac{y_n}{x_n}\right) \to -\frac 1e
                 \Rightarrow  y_n \sim \frac 1{\sqrt{ne}}$.
\end{enumerate}
 \fincorrection
\correction{004039}
\begin{enumerate}
  \item 
  \item
  \begin{enumerate} 
    \item 
    \item$a \sim e^{-b}  \Rightarrow  a\ln b \to 0  \Rightarrow  b^a \to 1$.
  \end{enumerate}    
\end{enumerate}
 \fincorrection
\correction{004040}
Existence et unicité de $x_n$ par étude de $f$ sur $[3,+\infty[$
(pour $x\le 3$ on ne peut pas avoir $0< f(x)< 1$). On a facilement
$x_n\to +\infty$ lorsque $n\to\infty$.

$\ln(x_n-2) = \Bigl(1-\frac1n\Bigr)\ln(x_n) \Rightarrow \ln\Bigl(1-\frac2{x_n}\Bigr) = -\frac{\ln(x_n)}n
 \Rightarrow x_n\ln(x_n)\sim 2n \Rightarrow x_n\sim\frac{2n}{\strut\ln n}$.
\fincorrection
\correction{004041}
$u_n = \frac1n - \frac1{n^6} +  o\Bigl(\frac1{n^6}\Bigr)$.
\fincorrection
\correction{004042}
Existence et unicité de~$x_n$ par étude de la fonction
$x \mapsto e^x+x$ sur $\R^+$.
On a clairement $x_n\to +\infty$ (lorsque $n\to\infty$) et $n=e^{x_n}+x_n$
d'où~: $$\ln n = \ln(e^{x_n}+x_n) = x_n + \ln(1+x_ne^{-x_n}) = 
x_n + x_ne^{-x_n} - \frac{x_n^2}2e^{-2x_n} +  o(x_n^2e^{-2x_n}).$$
On en déduit $x_n\sim \ln n$. \'Ecrivons $x_n = \ln n + y_n$~:
$$0 = y_n + x_ne^{-x_n} - \frac{x_n^2}2e^{-2x_n} +  o(x_n^2e^{-2x_n})$$
d'où $y_n \to 0$ (lorsque $n\to\infty$) et $y_n\sim -x_ne^{-x_n} \sim -\frac{\ln n}ne^{-y_n}
\sim-\frac{\ln n}n$. \'Ecrivons maintenant $y_n=-\frac{\ln n}n+z_n$~:
\begin{align*}0 
&= -\frac{\ln n}n + z_n + (\ln n + y_n)\frac{e^{-y_n}}n - \frac{x_n^2}2e^{-2x_n} +  o(x_n^2e^{-2x_n})\\
&=z_n + \frac{(\ln n)(-y_n+ o(y_n))}n + y_n\frac{e^{-y_n}}n - \frac{x_n^2}2e^{-2x_n} +  o(x_n^2e^{-2x_n})\\
&=z_n + \frac{(\ln n)(-y_n+ o(y_n))}n + y_n\frac{e^{-y_n}}n - \frac{x_n^2}{2n^2}e^{-2y_n} +  o\Bigl(\frac{x_n^2e^{-2y_n}}{n^2}\Bigr)\\
&=z_n + \frac{\ln^2n}{2n^2} +  o\Bigl(\frac{\ln^2n}{n^2}\Bigr)\cr
\end{align*}
d'où $z_n\sim-\frac{\ln^2n}{2n^2}$ et finalement,
$x_n = \ln n - \frac{\ln n}n - \frac{\ln^2n}{2n^2}+o\Bigl(\frac{\ln^2n}{n^2}\Bigr).$
\fincorrection
\correction{004043}
\begin{enumerate}
  \item 
  \item \begin{align*} nx_n^n(x_n - 1) = x_n^n+\frac 12 & \Rightarrow 
    f_{n+1}(x_n) = \frac{n+1}n\Bigl(x_n^{n+1} + \frac{x_n}2\Bigr) - x_n^{n+1} - \frac12
                 = \frac{x_n^{n+1}}n + \frac{(n+1)x_n - n}{\strut 2n} > 0\\ & \Rightarrow  x_{n+1} < x_n.\\ \end{align*}
    Donc la suite $(x_n)$ est décroissante et minorée par~$1$, elle converge vers~$\ell \ge 1$.
    
    $0 \le x_n - 1 = \frac 1n+ \frac1{2nx_n^n} \to 0$ (lorsque $n\to\infty$) donc $\ell = 1$.
    
    Soit $y_n = n(x_n-1) = 1 + \frac1{2x_n^n}$.
    On a $f(y_n) = \frac{\ln(2(y_n-1))}{y_n} = -\frac{\ln x_n}{x_n-1} = -g(x_n)$
    et $f,g$ sont strictement coissantes sur $[1,+\infty[$ donc les suites $(x_n)$
    et $(y_n)$ varient en sens contraire. On en déduit que la suite $(y_n)$
    décroît donc admet une limite $\lambda \ge 1$,
    soit $x_n = 1 + \frac\lambda n + o\Bigl(\frac1n\Bigr)$.
    Alors $x_n^n \to e^\lambda$ (lorsque $n\to\infty$) d'où $\lambda = 1 + \frac1{2e^\lambda}$.
\end{enumerate}
 \fincorrection
\correction{004044}

Il s'agit bien sûr de calculer un développement limité, le premier terme de ce développement donne l'équivalent cherché.

\begin{enumerate}
  \item 
Le dl à l'ordre $3$ en $0$ est 
$$2e^x - \sqrt{1+4x} - \sqrt{1+6x^2} = -\frac {11x^3}3 + o(x^3)$$
donc 
$$2e^x - \sqrt{1+4x} - \sqrt{1+6x^2} \sim -\frac {11x^3}3.$$

  \item De même 
$$(\cos x)^{\sin x} - (\cos x)^{\tan x} \sim \frac {x^5}4.$$

  \item On pose $h=x-\sqrt 3$ alors 
$$ \arctan x + \arctan \frac 3x -\frac {2\pi}3 = -\frac {h^2}{8\sqrt3} + o(h^2)$$
donc 
$$ \arctan x + \arctan \frac 3x -\frac {2\pi}3 \sim -\frac {(x-\sqrt 3)^2}{8\sqrt3}.$$

  \item En $+\infty$
$$\sqrt{x^2+1} -2\sqrt[3]{x^3+x} + \sqrt[4]{x^4+x^2}= \frac 1{12x} + o(\frac 1x)$$
donc $$\sqrt{x^2+1} -2\sqrt[3]{x^3+x} + \sqrt[4]{x^4+x^2} \sim \frac 1{12x}.$$

  \item Il faut distinguer les cas $x>0$ et $x<0$ pour trouver :
$$\Argch\left(\frac1{\cos x}\right) \sim |x|.$$
\end{enumerate}
 \fincorrection
\correction{004045}
Le dl de $\cos x$ en $0$ à l'ordre $6$ est :
$$\cos x = 1 - \frac{1}{2!} x^2 + \frac{1}{4!}x^4  - \frac{1}{6!} x^6 + o(x^6).$$

Calculons celui de $\frac{1+ax^2}{1+bx^2}$ :

\begin{align*}
\frac{1+ax^2}{1+bx^2} 
  & = (1+ax^2) \times \frac{1}{1+bx^2} \\
  & = (1+ax^2)\times \big(1-bx^2+b^2x^4-b^3x^6+o(x^6) \big) \quad \text{ car } \frac{1}{1+u} = 1-u+u^2 - u^3+o(u^3) \\
  &= \ \ \cdots  \qquad \text{ on développe } \\
  &= 1 + (a-b) x^2 - b(a-b) x^4 + b^2(a-b) x^6 + o(x^6) \\
\end{align*}

Notons $\Delta(x) = \cos x - \frac{1+ax^2}{1+bx^2}$ alors 
$$\Delta(x) = \big(-\frac12-(a-b)\big)x^2 + \big(\frac{1}{24} + b(a-b)\big) x^4 
+ \big(-\frac{1}{720}-b^2(a-b)\big) x^6 + o(x^6).$$

Pour que cette différence soit la plus petite possible (lorsque $x$ est proche de $0$)
il faut annuler le plus possible de coefficients de bas degré.
On souhaite donc avoir 
$$-\frac12-(a-b) = 0 \qquad \text{et} \qquad \frac{1}{24} + b(a-b)=0.$$
En substituant l'égalité de gauche dans celle de droite on trouve :
$$a=-\frac{5}{12}  \qquad \text{et} \qquad b=\frac{1}{12}.$$

On obtient alors 
$$\Delta(x) =  \big(-\frac{1}{720}-b^2(a-b)\big) x^6 + o(x^6) = \frac{1}{480} x^6 + o(x^6).$$

\bigskip

Avec notre choix de $a,b$ nous avons obtenu une très bonne approximation de $\cos x$.
Par exemple lorsque l'on évalue  $\frac{1+ax^2}{1+bx^2}$ (avec $a=-\frac{5}{12}$ et $b=\frac{1}{12}$)
en $x=0.1$ on trouve :
$$0.9950041631\ldots$$
Alors que 
$$\cos(0.1)=0.9950041652\ldots$$
En l'on trouve ici $\Delta(0.1) \simeq 2\times 10^{-9}$.
\fincorrection
\correction{004046}
$a=-7/60$, $b=1/20$, $\Delta \sim 11x^7/50400$.
\fincorrection
\nocorrection
\correction{004048}
\begin{enumerate}
  \item $(3a\alpha(\alpha-a) + 2(b\alpha-a\beta))X^7$.
  \item $-x^7/30 + o_{x\to0}(x^7)$.
\end{enumerate}
 \fincorrection
\nocorrection
\nocorrection
\correction{004051}
$\sim x$ pour $k$ impair, et
         $= 1 + x\ln(x) + o_{x\to0}(x\ln x) \sim 1$ si $k$ est pair $\ge 2$.
\fincorrection
\nocorrection
\nocorrection
\correction{004030}
\begin{enumerate}
  \item 
  \item $e\left(1 - \frac{\ln2}x + \frac{\ln^22}{2!\,x^2} - \dots
                       + (-1)^n\frac{\ln^n2}{n!\,x^n}\right)
                       + o (x^{-n})$.
\end{enumerate}
 \fincorrection
\nocorrection
\correction{004032}
$f^{-1}(y) = \frac{y}{a_1} - \frac{a_2y^2}{a_1^3} + o (y^2)$.
\fincorrection
\correction{004033}
$(1-e^x)^n = \sum_{k=0}^n(-1)^kC_n^ke^{kx}
 = \sum_{p=0}^{n+2} \Bigl(\sum_{k=0}^n(-1)^kC_n^kk^p \Bigr) \frac {x^p}{p!}
   + o (x^{n+2})$,

$(1-e^x)^n = (-x)^n\left({1 + \frac {nx}2 + \frac {n(3n+1)}{24}x^2 + o  (x^2)}\right)$.

\fincorrection
\nocorrection
\nocorrection
\nocorrection
\nocorrection
\nocorrection
\nocorrection
\nocorrection
\correction{000745}\ 
\begin{enumerate}
  \item 

On note $x$ la distance de l'observateur au pied de la statue.
On note $\alpha$ l'angle d'observation de la statue seule, et $\beta$
l'angle d'observation du piédestal seul.

\begin{center}
\begin{tikzpicture}[scale=1]

\filldraw[fill=gray!50, draw=gray] (5,0)--(6,0)--(6,4)--(5,4)--cycle;
\filldraw[fill=gray, draw=black] (5,4)--(5.5,4)--(5.5,6)--(5,6)--cycle;
\draw[thick] (0,0)--(5,0);
\draw[thick] (0,0)--(5,4);
\draw[thick] (0,0)--(5,6);
 
 \node at (5,5) [left] {$s$};    
 \node at (5,2) [left] {$p$};  
 \node at (2.5,0) [below] {$x$};  

\draw (2,0) arc(0:39:2) ;
\draw (39:2.5) arc(39:50:2.5) ;
\draw (39:2.6) arc(39:50:2.6) ;

\node at (45:2.9)  {$\alpha$};  
\node at (20:2.3)  {$\beta$};  
\end{tikzpicture}
\end{center}

Nous avons les relations trigonométriques dans les triangles rectangles :
$$\tan(\alpha+\beta) = \frac{p+s}{x} \qquad\text{ et }\qquad \tan\beta = \frac{p}{x}$$
On en déduit les deux identités :
$$\alpha+\beta = \Arctan\left(\frac{p+s}{x}\right) \qquad\text{ et }\qquad \beta = \Arctan\left(\frac{p}{x}\right)$$
à partir desquelles on obtient $\alpha=\alpha(x)=\Arctan\left(\frac{p+s}{x}\right)-\Arctan\left(\frac{p}{x}\right)$.

\'Etudions cette fonction sur $]0,+\infty[$ : elle est dérivable et 
$$\alpha'(x)=\frac{-\frac{s+p}{x^2}}{1+\left(\frac{s+p}{x}\right)^2}-\frac{-\frac{p}{x^2}}{1+\left(\frac{p}{x}\right)^2}=\frac{s}{(x^2+p^2)(x^2+(s+p)^2)}\big(p(p+s)-x^2\big)$$
Ainsi $\alpha'$ ne s'annule sur $]0,+\infty[$ qu'en $x_0 = \sqrt{p(p+s)}$.
Par des considérations physiques, à la limite en $0$ et en $+\infty$, 
l'angle $\alpha$ est nul,
alors en $x_0$ nous obtenons un angle $\alpha$ maximum.
Donc la distance optimale de vision est  $x_0 = \sqrt{p(p+s)}$.

\item
Pour calculer l'angle maximum $\alpha_0$ correspondant, on pourrait 
calculer $\alpha_0 = \alpha(x_0)$ à partir de la définition de la fonction
$\alpha(x)$. Pour obtenir une formule plus simple nous 
utilisons la formule trigonométrique suivante : si
$a$, $b$ et $a-b$ sont dans l'intervalle de définition de la fonction $\tan$, alors
$\tan(a-b)=\frac{\tan a-\tan b}{1+\tan a \tan b}$, ce qui donne ici
$$\tan\alpha_0 = \tan\big( (\alpha_0+\beta_0) - \beta_0 \big)
=\frac{\frac{p+s}{x_0}-\frac{p}{x_0}}{1+\frac{p+s}{x_0}\cdot\frac{p}{x_0}}
= \frac{s}{2x_0}=\frac{s}{2\sqrt{p(p+s)}}$$
Comme $\alpha_0\in]-\frac{\pi}{2},\frac{\pi}{2}[$, 
on en déduit $\alpha_0 = \Arctan \frac{s}{2x_0} = \Arctan \frac{s}{2\sqrt{p(p+s)}}$.

\item Pour la statue de la liberté, on a la hauteur de la statue $s=46$ mètres
et la hauteur du piédestal $p=47$ mètres.
On trouve donc 
$$x_0 = \sqrt{p(p+s)} \simeq 65,40 \text{mètres} \qquad 
\alpha_0 = \Arctan \frac{s}{2\sqrt{p(p+s)}} \simeq 19^\circ.$$

Voici les représentations de la statue et de la fonction $\alpha(x)$ pour ces valeurs de $s$
et $p$.
\end{enumerate}

\begin{center}
\begin{tikzpicture}[scale=0.5]

\filldraw[fill=gray!50, draw=gray] (6.54,0)--(7.54,0)--(7.54,4.7)--(6.54,4.7)--cycle;
\filldraw[fill=gray, draw=black] (6.54,4.7)--(7.04,4.7)--(7.04,9.3)--(6.54,9.3)--cycle;
\draw[thick] (0,0)--(6.54,0);
\draw[thick] (0,0)--(6.54,4.7);
\draw[thick] (0,0)--(6.54,9.3);
 
 \node at (6.5,6.5) [left] {$s$};    
 \node at (6.5,2) [left] {$p$};  
 \node at (3.5,0) [below] {$x_0$};  

\draw (2,0) arc(0:36:2) ;
\draw (36:2.5) arc(36:55:2.5) ;
\draw (36:2.6) arc(36:55:2.6) ;

\node at (45:3)  {$\alpha_0$};  
\node at (18:2.5)  {$\beta_0$};  
\end{tikzpicture} 
\qquad
\begin{tikzpicture}[scale=0.2]
      \draw[->,>=latex, gray] (-0.5,0)--(30,0) node[below,black] {$x$};
      \draw[->,>=latex, gray] (0,-0.5)--(0,23) node[left,black] {$\alpha(x)$};  
 
          \draw[ultra thick, color=red,domain=0.1:30,samples=200,smooth] plot (\x,{(atan(93/(\x*10))-atan(47/(\x*10)))}) node[above left] {$\alpha(x)$}; 

   \draw[dashed] (0,19.3)--(6.54,19.3);
   \draw[dashed] (6.54,0)--(6.54,19.3);

     \fill (0,0) circle (7pt);     
     \fill (6.54,0) circle (7pt);
     \fill (0,19.3) circle (7pt);
    
     \node at (0,19.3)[left] {$\alpha_0$};
     \node at (6.54,0)[below] {$x_0$};
     \node at (0,0)[below right] {$0$};
\end{tikzpicture}
\end{center}

\fincorrection
\correction{000746}
\begin{enumerate}
    \item Soit $f(a) = \Arcsin a - \frac{a}{\sqrt{1-a^2}}$ sur $]0,1[$.
Alors $f'(a) = \frac{1}{\sqrt{1-a^2}} - \frac{1}{\sqrt{1-a^2}(1-a^2)} = \frac{1}{\sqrt{1-a^2}}\cdot \frac{-a^2}{1-a^2}$
donc $f'(a) \leq 0$. Ainsi $f$ est strictement décroissante et $f(0)=0$
donc $f(a) < 0$ pout tout $a \in ]0,1[$.
  \item Si $g(a) = \Arctan a - \frac{a}{1+a^2}$
alors $g'(a) = \frac1{1+a^2}-\frac{1-a^2}{(1+a^2)^2}=\frac{2a^2}{(1+a^2)^2} > 0$.
Donc $g$ est strictement croissante et $g(0) = 0$ donc $g$ est strictement positive sur $]0,+\infty[$.
\end{enumerate}
\fincorrection
\correction{000747}\
\begin{enumerate}
    \item $\sin^2 y = 1-\cos^2 y$, donc
$\sin y = \pm\sqrt{1-\cos^2 y}$.
Avec $y=\Arccos x$, il vient $\sin(\Arccos x)= \pm\sqrt{1-x^2}$. Or $\Arccos x \in [0,\pi]$, 
donc $\sin(\Arccos x)$ est positif et finalement $\sin(\Arccos x)=  +\sqrt{1-x^2}$.
De la m\^eme mani\`ere on trouve $\cos(\Arcsin x) = \pm\sqrt{1-x^2}$. 
Or $\Arcsin x \in [-\frac{\pi}{2},\frac{\pi}{2}]$, donc $\cos(\Arcsin x)$ est positif 
et finalement $\cos(\Arcsin x)=+\sqrt{1-x^2}$.

Ces deux égalités sont à connaître ou à savoir retrouver très rapidement :
$$\sin(\Arccos x)  = \sqrt{1-x^2} = \cos(\Arcsin x).$$


Enfin, puisque $\cos(2y) = \cos^2 y - \sin^2 y$, on obtient avec $y=\Arcsin x$,
$$\cos(2\Arcsin x)=(\sqrt{1-x^2})^2-x^2=1-2x^2.$$

\item Commen\c{c}ons par calculer $\sin(\Arctan x)$, $\cos(\Arctan x)$.
On utilise l'identité $1+\tan^2y=\frac 1{\cos^2 y}$ avec $y = \Arctan x$, ce qui donne 
$\cos^2 y = \frac{1}{1+x^2}$ et
 $\sin^2 y = 1-\cos^2 y=\frac{x^2}{1+x^2}$.
 Il reste à déterminer les signes de $\cos(\Arctan x)= \pm\frac{1}{\sqrt{1+x^2}}$ et $\sin(\Arctan x) = \pm\frac{x}{\sqrt{1+x^2}}$
Or $y=\Arctan x$ donc $y\in]-\frac{\pi}{2},\frac{\pi}{2}[$ et $y$ a le même signe que $x$ : 
ainsi $\cos y>0$, et $\sin y$ a le même signe que $y$ et donc que $x$. Finalement, on a 
$\cos(\Arctan x)= \frac{1}{\sqrt{1+x^2}}$ et $\sin(\Arctan x) = \frac{x}{\sqrt{1+x^2}}$.

\medskip


Il ne reste plus qu'à linéariser $\sin(3y)$ :
\begin{eqnarray*}
\sin(3y) &=& \sin(2y+y)=\cos(2y)\sin(y)+\cos(y)\sin(2y)\\
 &=&(2\cos^2 y-1)\sin y +2\sin y\cos^2y\\
 &=&4 \sin y \cos^2 y - \sin y
\end{eqnarray*}



Maintenant
\begin{eqnarray*}
\sin(3\Arctan x) &=& \sin(3y)= 4 \sin y \cos^2 y - \sin y\\
 &=& 4{\frac {x}{\left (1+{x}^{2}\right )^{3/2}}}-{\frac {x}{\sqrt {1+{x}^{2}}}}={\frac {x(3-x^2)}{\left (1+{x}^{2}\right )^{3/2}}}
\end{eqnarray*}

\bigskip

\emph{Remarque :} la méthode générale pour obtenir la formule de linéarisation de $\sin(3y)$ est 
d'utiliser les nombres complexes
et la formule de Moivre. On développe 
$$\cos(3y) + i \sin(3y) = (\cos y + i \sin y)^3 = \cos^3 y + 3i\cos^2y \sin y + \cdots$$ 
puis on identifie les parties imaginaires pour avoir $\sin(3y)$, 
ou les parties réelles pour avoir $\cos(3y)$.

\end{enumerate}
\fincorrection
\nocorrection
\correction{000749}\
\begin{enumerate}
\item On vérifie d'abord que $2\Arccos \frac{3}{4}\in[0,\pi]$ (sinon, l'équation 
n'aurait aucune solution). En effet, par définition, la fonction $\Arccos$ est décroissante 
sur $[-1,1]$ à valeurs dans $[0,\pi]$, donc puisque $\frac{1}{2}\le\frac{3}{4}\le 1$ 
on a $\frac{\pi}{3}\ge\cos\left(\frac{3}{4}\right)\ge 0$.
Puisque par définition $\Arccos x\in[0,\pi]$, on obtient en prenant le cosinus:
$$\Arccos x = 2\Arccos\left(\frac{3}{4}\right)\Longleftrightarrow x 
= \cos\left(2\Arccos \frac{3}{4}\right)$$
En appliquant la formule $\cos 2u = 2\cos^2u-1$, on arrive  donc \`a une unique solution
$x = 2(\frac{3}{4})^2-1 = \frac{1}{8}$.
\item Vérifions d'abord que $-\frac{\pi}{2}\le\Arcsin \frac{2}{5} + \Arcsin \frac{3}{5}\le\frac{\pi}{2}$. 
En effet, la fonction $\Arcsin$ est strictement croissante et 
$0<\frac{2}{5}<\frac{1}{2}<\frac{3}{5}<\frac{\sqrt{2}}{2}$, ce qui donne
$0<\Arcsin\left(\frac{2}{5}\right)<\frac{\pi}{6}<\Arcsin\left(\frac{3}{5}\right)<\frac{\pi}{4}$, 
d'où l'encadrement
$0+\frac{\pi}{6}<\Arcsin \frac{2}{5} + \Arcsin \frac{3}{5}\le\frac{\pi}{6}+\frac{\pi}{4}$.

Puisque par définition on aussi $\Arcsin x\in [-\frac{\pi}{2},\frac{\pi}{2}]$, 
il vient en prenant le sinus:
\begin{eqnarray*}
\lefteqn{\Arcsin x=\Arcsin \frac{2}{5} + \Arcsin \frac{3}{5}}\\
 &\Longleftrightarrow& x=\sin\left(\Arcsin \frac{2}{5} + \Arcsin \frac{3}{5}\right)\\
 &\Longleftrightarrow& x = \frac{3}{5}\cos\left(\Arcsin \frac{2}{5}\right)+\frac{2}{5} \cos\left(\Arcsin \frac{3}{5}\right)
\end{eqnarray*}
La dernière équivalence vient de la formule de $\sin(a+b)= \cos a \sin b + \cos b \sin a$
et de l'identité $\sin\big( \Arcsin u\big) = u$.

En utilisant la formule $\cos\left(\Arcsin x\right) = \sqrt{1-x^2}$, 
on obtient une unique solution: $x = \frac35\sqrt{\frac{21}{25}}+\frac25\frac 45 
=\frac{3\sqrt{21}+8}{25}$.

\item Supposons d'abord que $x$ est solution. Remarquons d'abord que $x$ est 
nécessairement positif, puisque $\Arctan x$ a le même signe que $x$. 
Alors, en prenant la tangente des deux membres, on obtient $\tan\big(\Arctan(2x)+\Arctan(x)\big)=1$.

En utilisant la formule donnant la tangente d'une somme :
$\tan(a+b)=\frac{\tan a +\tan b}{1-\tan a \tan b}$, on obtient $\frac{2x+x}{1-2x\cdot x}=1$, 
et finalement $2x^2+3x-1=0$ qui admet une unique solution positive $x_0=\frac{-3+\sqrt{17}}{4}$. 
Ainsi, {\it si} l'équation de départ admet une solution, c'est nécessairement $x_0$. 

Or, en posant $f(x)=\Arctan(2x)+\Arctan(x)$, la fonction $f$ est continue sur $\Rr$. 
Comme $f(x)\xrightarrow[x\to -\infty]{}-\pi$ et $f(x)\xrightarrow[x\to +\infty]{}+\pi$, 
on sait d'après le théorème des valeurs intermédiaires que $f$ prend la valeur 
$\frac{\pi}{4}$ au moins une fois (et en fait une seule fois, puisque $f$ est 
strictement croissante comme somme de deux fonctions strictement croissantes). 
Ainsi l'équation de départ admet bien une solution, qui est $x_0$.
\end{enumerate}
\fincorrection
\nocorrection
\nocorrection
\correction{000752}\ 
\begin{enumerate}
    \item Soit $f$ la fonction définie sur $[-1,1]$ par 
$f(x) = \Arcsin x +\Arccos x$: $f$ est continue sur l'intervalle $[-1,1]$, 
et dérivable sur $]-1,1[$. Pour tout $x\in]-1,1[$, 
$f'(x)= \frac{1}{\sqrt{1-x^2}}+\frac{-1}{\sqrt{1-x^2}}  = 0$. 
Ainsi $f$ est constante sur $]-1,1[$, donc sur $[-1,1]$ (car continue aux extrémités). 
Or $f(0) = \Arcsin 0 +\Arccos 0 = \frac \pi 2$
donc pour tout $x\in[-1,1]$, $f(x) = \frac \pi 2$.

   \item Soit $g(x) = \Arctan x + \Arctan \frac 1x$.
Cette fonction est définie sur $]-\infty,0[$ et sur $]0,+\infty[$ (mais pas en $0$).
On a 
$$g'(x)= \frac{1}{1+x^2} + \frac{-1}{x^2} \cdot \frac{1}{1+\frac{1}{x^2}} = 0,$$
donc $g$ est constante sur chacun de ses intervalles de définition:
$g(x) = c_1$ sur $]-\infty,0[$ et $g(x) = c_2$ sur $]0,+\infty[$.
Sachant $\Arctan 1 = \frac\pi4$, on calcule $g(1)$ et $g(-1)$ on obtient $c_1 = -\frac \pi 2$
et $c_2 = +\frac \pi 2$.
\end{enumerate}
\fincorrection
\nocorrection
\nocorrection
\nocorrection
\nocorrection
\correction{003914}
$\arcsin x = \arctan\left( \frac x{\sqrt{1-x^2}} \right)$.
\fincorrection
\correction{003915}
$\arctan a + \arctan b \equiv \arctan\left(\frac {a+b}{1-ab}\right)(\mathrm{mod}\, \pi)$.
\fincorrection
\correction{003916}
$x=\frac{\sqrt{8}-\sqrt{15}}{12}$.
\fincorrection
\correction{003917}
$\cos 4x = -\sin x  \Rightarrow  x \equiv \frac{3\pi}{10} \pmod{\frac{2\pi}{5}}  \text{ ou } x \equiv \frac{\pi}{6} \pmod {\frac{2\pi}{3}}$.
Donc $x=\frac{3\pi}{10}$.
\fincorrection
\correction{003918}
\begin{enumerate}
  \item $x  > -1  \Rightarrow  = \frac \pi4   - \arctan x$,  \qquad
$x  < -1  \Rightarrow  =-\frac {3\pi}4- \arctan x$.
  \item $= \frac 12 \arccos x$.
  \item $-1 \le x < -\frac1{\sqrt2}  \Rightarrow  = \arcsin x + \frac{3\pi}4$,  \qquad
         $-\frac1{\sqrt2} < x \le 1   \Rightarrow  = \arcsin x - \frac{\pi}4$.
  \item $=\frac\pi4$.
  \item 
$f(x)=0$ pour $x\in ]-\infty,0[$ ; \\
$f(x)=\pi$ pour $x\in ]0,1[$ ; \\
$f(x)=0$ pour $x\in ]1,+\infty[$.
\end{enumerate}
 \fincorrection
\correction{003919}
$D = \left[-\frac 12, \frac{\sqrt3}2\right]$,\quad
   $f(x) = \frac12 - x^2- x\sqrt3\sqrt{1-x^2}$.
\fincorrection
\correction{003920}
$\cos(3\arctan x) = \frac {1-3x^2}{(1+x^2)^{3/2}}$.
\qquad
$\cos^2\left(\frac 12\arctan x\right) = \frac 12\left(1+\frac1{\sqrt{1+x^2}}\right)$.
\fincorrection
\correction{003921}
$f(x)=0$ pour $x\in ]0,\pi/2[$ ;
$f(x)=2x-\pi$ pour $x\in ]\pi/2,\pi[$ ;
$f(x)=3\pi-2x$ pour $x\in ]\pi,3\pi/2[$ ;
$f(x)=0$ pour $x\in ]3\pi/2,2[$.
\fincorrection
\correction{003922}
$f(x)=-8\arctan x - 2\pi$ pour $x\in ]-\infty,-1[$, solution $-\sqrt3$ ; \\
$f(x)=-4\arctan x $ pour $x\in ]-1,0[$, solution $-1/\sqrt3$ ; \\
$f(x)=4\arctan x$ pour $x\in ]0,1[$, solution $1/\sqrt3$ ; \\
$f(x)=2\pi$ pour $x\in ]1,+\infty[$.

\fincorrection
\correction{003923}
$f(x)=x+\pi/4$ pour $x\in ]-\pi,-\pi/2[$ ; \\
$f(x)=-\pi/4$ pour $x\in ]-\pi/2,\pi/2[$ ; \\
$f(x)=x-3\pi/4$ pour $x\in ]\pi/2,\pi[$.
\fincorrection
\correction{003924}
$= \sqrt{\frac {x^2+1}{x^2+2} }$.
\fincorrection
\correction{003925}
\begin{enumerate}
  \item $x = \frac 16$.
  \item $x = \pm1{\sqrt 2}$.
  \item $x \in {]-\infty,-1[} \cup ]0,+\infty[$.
  \item $x^3 - 3x^2 - 12x + 10 = 0  \Rightarrow  x = 5,\,-1\pm\sqrt3$.
Seule la solution $x=5$ convient.
\end{enumerate}
 \fincorrection
\nocorrection
\nocorrection
\correction{003928}
$\sin\bigl(2g(x)\bigr) = \sin(f(x))$.

$f(x)=-\pi-2g(x)$ pour $x\in ]-\infty,\sin a-\cos a[$ ; \\
$f(x)=2g(x)$ pour $x\in ]\sin a-\cos a,\sin a+\cos a[$ ; \\
$f(x)=\pi-2g(x)$ pour $x\in ]\sin a+\cos a,+\infty[$.
\fincorrection
\nocorrection
\nocorrection
\nocorrection
\correction{005084}
$\Arcsin x$ existe si et seulement si $x$ est dans $[-1,1]$. Donc, $\sin(\Arcsin x)$ existe si et seulement si
$x$ est dans $[-1,1]$ et pour $x$ dans $[-1,1]$, $\sin(\Arcsin x)=x$.
 \item  $\Arcsin(\sin x)$ existe pour tout réel $x$ mais ne vaut $x$ que si $x$ est 
dans $\left[-\frac{\pi}{2},\frac{\pi}{2}\right]$.
\textbullet~S'il existe un entier relatif $k$ tel que
$-\frac{\pi}{2}+2k\pi\leq x<\frac{\pi}{2}+2k\pi$, alors $-\frac{\pi}{2}\leq x-2k\pi<\frac{\pi}{2}$ et
donc

$$\Arcsin(\sin x)=\Arcsin(\sin(x-2k\pi))=x-2k\pi.$$
De plus, on a $k\leq\frac{x}{2\pi}+\frac{1}{4}<k+\frac{1}{2}$
et donc $k=E\left(\frac{x}{2\pi}+\frac{1}{4}\right)$.
\textbullet~S'il existe un entier relatif $k$ tel que $\frac{\pi}{2}+2k\pi\leq
x<\frac{3\pi}{2}+2k\pi$, alors $-\frac{\pi}{2}<\pi-x+2k\pi\leq\frac{\pi}{2}$ 
et donc

$$\Arcsin(\sin x)=\Arcsin(\sin(\pi-x+2k\pi))=\pi-x+2k\pi.$$
De plus, $k\leq\frac{x}{2\pi}-\frac{1}{4}<k+\frac{1}{2}$ et donc $k=E\left(\frac{x}{2\pi}-\frac{1}{4}\right)$.
 \item  $\Arccos x$ existe si et seulement si $x$ est dans $[-1,1]$. Donc, $\cos(\Arccos x)$ existe si et seulement si
$x$ est dans $[-1,1]$ et pour $x$ dans $[-1,1]$, $\cos(\Arccos x)=x$.
 \item  $\Arccos(\cos x)$ existe pour tout réel $x$ mais ne vaut $x$ que si $x$ est dans $[0,\pi]$.
\textbullet~S'il existe un
entier relatif $k$ tel que $2k\pi\leq x<\pi+2k\pi$, alors $\Arccos(\cos x)=x-2k\pi$ avec $k=E\left(\frac{x}{2\pi}\right)$.
\textbullet~S'il existe un entier relatif $k$ tel que $-\pi+2k\pi\leq x<2k\pi$ alors 
$\Arccos(\cos x)=\Arccos(\cos(2k\pi-x))=2k\pi-x$ avec $k=E\left(\frac{x+\pi}{2\pi}\right)$.
 \item  Pour tout réel $x$, $\tan(\Arctan x)=x$.
 \item  $\Arctan(\tan x)$ existe si et seulement si $x$ n'est pas
dans $\frac{\pi}{2}+\pi\Zz$ et pour ces $x$, il existe un entier relatif $k$ 
tel que $-\frac{\pi}{2}+k\pi<x<\frac{\pi}{2}+k\pi$. Dans ce cas, $\Arctan(\tan x)=\Arctan(\tan(x-k\pi))=x-k\pi$ avec
$k=E\left(\frac{x}{\pi}+\frac{1}{2}\right)$.
\fincorrection
\correction{005085}
\begin{enumerate}
 \item  \textbf{1ère solution}. Posons $f(x)=\Arccos x+\Arcsin x$ pour $x$ dans $[-1,1]$.
$f$ est définie et continue sur $[-1,1]$, dérivable sur $]-1,1[$. De plus, pour $x$ dans
$]-1,1[$,

$$f'(x)=\frac{1}{\sqrt{1-x^2}}-\frac{1}{\sqrt{1-x^2}}= 0.$$
Donc $f$ est constante sur $[-1,1]$ et pour $x$
dans $[-1,1]$, $f(x)=f(0)=\frac{\pi}{2}$.

\begin{center}
\shadowbox{
$\forall x\in[-1,1],\;\Arccos x+\Arcsin x=\frac{\pi}{2}.$
}
\end{center}
\textbf{2ème solution}. Il existe un unique réel $\theta$ dans $[0,\pi]$ tel que $x=\cos\theta$, à savoir
$\theta=\Arccos x$. Mais alors,

$$\Arccos x+\Arcsin x=\theta+\Arcsin\left(\sin(\frac{\pi}{2}-\theta)\right)
=\theta+\frac{\pi}{2}-\theta=\frac{\pi}{2}$$
(car $\frac{\pi}{2}-\theta$ est dans
$[-\frac{\pi}{2},\frac{\pi}{2}])$.
 \item  \textbf{1ère solution}. Pour $x$ réel non nul, posons $f(x)=\Arctan x+\Arctan\frac{1}{x}$. $f$ est
impaire.
$f$ est dérivable sur $\Rr^*$ et pour tout réel $x$ non
nul, $f'(x)=\frac{1}{1+x^2}-\frac{1}{x^2}\frac{1}{1+\frac{1}{x^2}}=0$. $f$ est donc constante sur $]-\infty,0[$ et
sur $]0,+\infty[$ (mais pas nécessairement sur $\Rr^*$). Donc, pour $x>0$, $f(x)=f(1)=2\Arctan1=\frac{\pi}{2}$, et 
puisque $f$ est impaire, pour $x< 0$, $f(x)=-f(-x)=-\frac{\pi}{2}$. Donc,

\begin{center}
\shadowbox{
$\forall
x\in\Rr^*,\;\Arctan 
x+\Arctan\frac{1}{x}=\left\{
\begin{array}{l}\frac{\pi}{2}\;\mbox{si}\;x>0\\-\frac{\pi}{2}\;\mbox{si}\;x<0
\end{array}\right.=\frac{\pi}{2}\text{sgn}(x).$
}
\end{center}
\textbf{2ème solution} Pour $x$ réel strictement positif donné, il existe un unique réel $\theta$ dans
$\left]0,\frac{\pi}{2}\right[$ tel que $x=\tan\theta$ à savoir $\theta=\Arctan x$.
Mais alors,

$$\Arctan x+\Arctan\frac{1}{x}=\theta+\Arctan\left(\frac{1}{\tan\theta}\right)
=\theta+\Arctan\left(\tan(\frac{\pi}{2}-\theta)\right)
=\theta+\frac{\pi}{2}-\theta=\frac{\pi}{2}$$
(car $\theta$ et $\frac{\pi}{2}-\theta$ sont éléments de
$\left]0,\frac{\pi}{2}\right[)$.
 \item  $\cos^2(\Arctan a)=\frac{1}{1+\tan^2(\Arctan a)}=\frac{1}{1+a^2}$. De plus , $\Arctan
a$ est dans $]-\frac{\pi}{2},\frac{\pi}{2}[$ et donc $\cos(\Arctan a)>0$.
On en déduit que pour tout réel $a$, $\cos(\Arctan a)=\frac{1}{\sqrt{1+a^2}}$ puis

$$\sin(\Arctan a)=\cos(\Arctan a)\tan(\Arctan a)=\frac{a}{\sqrt{1+a^2}}.$$

\begin{center}
\shadowbox{
$\forall a\in\Rr,\;\cos(\Arctan a)=\frac{1}{1+a^2}\;\mbox{et}\;\sin(\Arctan a)=\frac{a}{\sqrt{1+a^2}}.$
}
\end{center}
 \item  D'après 3),

$$\cos(\Arctan a+\Arctan b)=\cos(\Arctan a)\cos(\Arctan
b)-\sin(\Arctan a)\sin(\Arctan b)=\frac{1-ab}{\sqrt{1+a^2}\sqrt{1+b^2}},$$
ce qui montre déjà , puisque $ab\neq1$, que
$\cos(\Arctan a+\Arctan b)\neq0$ et donc que $\tan(\Arctan a+\Arctan b)$ existe. On a immédiatement,

$$\tan(\Arctan a+\Arctan b)=\frac{a+b}{1-ab}.$$
Maintenant, $\Arctan a+\Arctan b$ est dans
$\left]-\pi,-\frac{\pi}{2}\right[\cup\left]-\frac{\pi}{2},\frac{\pi}{2}\right[\cup\left]\frac{\pi}{2},\pi\right[$.

\begin{itemize}
\item[\textbf{1er cas.}] Si $ab<1$ alors $\cos(\Arctan a+\Arctan
b)>0$ et donc $\Arctan a+\Arctan b$ est dans $\left]-\frac{\pi}{2},\frac{\pi}{2}\right[$.
Dans ce cas, $\Arctan a+\Arctan b=\Arctan\left(\frac{a+b}{1-ab}\right)$.

\item[\textbf{2ème cas.}]
Si $ab>1$ alors $\cos(\Arctan a+\Arctan b)<0$ et donc $\Arctan a+\Arctan b$ est 
dans $\left]-\pi,-\frac{\pi}{2}\right[\cup\left]\frac{\pi}{2},\pi\right[$. Si de plus $a>0$, 
$\Arctan a+\Arctan b>-\frac{\pi}{2}$ et donc $\Arctan a+\Arctan b$ est dans 
$\left]\frac{\pi}{2},\pi\right[$. Dans ce cas, $\Arctan a+\Arctan b-\pi$ est dans
$\left]-\frac{\pi}{2},\frac{\pi}{2}\right[$ et a même tangente que 
$\Arctan\frac{a+b}{1-ab}$. Donc, $\Arctan a+\Arctan 
b=\Arctan\frac{a+b}{1-ab}+\pi$. Si $a<0$, on trouve de même $\Arctan a+\Arctan
b=\Arctan\frac{a+b}{1-ab}-\pi$.
\end{itemize}
En résumé,

\begin{center}
\shadowbox{
$\Arctan a+\Arctan b=
\left\{\begin{array}{l}\rule[-4mm]{0mm}{0mm}\Arctan\frac{a+b}{1-ab}\;\mbox{si}\;ab<1\\
\rule[-4mm]{0mm}{0mm}\Arctan\frac{a+b}{1-ab}
+\pi\;\mbox{si}\;ab>1\;\mbox{et}\;a>0\\
\Arctan\frac{a+b}{1-ab}-\pi\;\mbox{si}\;ab>1\;
\mbox{et}\;a<0\end{array}\right..$
}
\end{center}
\end{enumerate}
\fincorrection
\correction{005087}
Pour $x$ réel, on pose
$f(x)=\int_{0}^{\sin^2x}\Arcsin\sqrt{t}\;dt+\int_{0}^{\cos^2x}\Arccos\sqrt{t}\;dt$.\\
La fonction $t\mapsto\Arcsin\sqrt{t}$ est continue sur $[0,1]$. Donc, la fonction $y\mapsto\int_{0}^{y}\Arcsin\sqrt{t}\;dt$ est définie et dérivable sur $[0,1]$. De
plus, $x\mapsto\sin^2x$ est définie et dérivable sur $\Rr$ à valeurs dans 
$[0,1]$. Finalement, la fonction $x\mapsto\int_{0}^{\sin^2x}\Arcsin\sqrt{t}\;dt$ est définie
et dérivable sur $\Rr$.
De même, la fonction $t\mapsto\Arccos\sqrt{t}$ est continue sur $[0,1]$. Donc, la fonction $y\mapsto\int_{0}^{y}\Arccos\sqrt{t}\;dt$ est définie et dérivable sur $[0,1]$. De plus, la fonction $x\mapsto\cos^2x$ est définie
et dérivable sur $\Rr$, à valeurs dans $[0,1]$. Finalement, la fonction $x\mapsto\int_{0}^{\cos^2x}\Arccos\sqrt{t}\;dt$ est définie
et dérivable sur $\Rr$.
Donc, $f$ est définie et dérivable sur $\Rr$ et, pour tout réel $x$,

\begin{align*}
f'(x)&=2\sin x\cos x\Arcsin(\sqrt{\sin^2x})-2\sin x\cos x\Arccos(\sqrt{\cos^2x})\\
 &=2\sin x\cos x\left(\Arcsin(|\sin x|)-\Arccos(|\cos x|)\right).
\end{align*}
On note alors que $f$ est $\pi$-pérodique et paire. Pour $x$ élément de
$[0,\frac{\pi}{2}]$, $f'(x)=2\sin x\cos x(x-x)=0$. $f$ est donc constante sur $[0,\frac{\pi}{2}]$ et pour $x$ élément
de $[0,\frac{\pi}{2}]$,
$f(x)=f\left(\frac{\pi}{4}\right)=\int_{0}^{1/2}\Arcsin\sqrt{t}\;dt+\int_{0}^{1/2}\Arccos\sqrt{t}dt=\int_{0}^{1/2}\frac{\pi}{2}\;dt=
\frac{\pi}{4}$. Mais alors, par parité et $\pi$-périodicité,

\begin{center}
\shadowbox{
$\forall
x\in\Rr,\;\int_{0}^{\sin^2x}\Arcsin\sqrt{t}\;dt+\int_{0}^{\cos^2x}\Arccos\sqrt{t}\;dt=\frac{\pi}{4}.$
}
\end{center}
\fincorrection
\correction{005088}
\begin{enumerate}
 \item   \textbf{1ère solution.} Pour tout réel $x$, $\sqrt{x^2+1}>\sqrt{x^2}=|x|$ et donc $-1<\frac{x}{\sqrt{x^2+1}}<
1$. Ainsi $f_1$ est définie et dérivable sur $\Rr$, impaire, et pour tout réel $x$,

$$f_1'(x)=\left(\frac{1}{\sqrt{x^2+1}}-\frac{1}{2}x\frac{2x}{(x^2+1)\sqrt{x^2+1}}\right)\frac{1}{\sqrt{1-\frac{x^2}{1+x^2}}}
=\frac{1}{1+x^2}=\Arctan'(x).$$
Donc il existe une constante réelle $C$ telle que pour tout réel $x$, $f_1(x)=\Arctan x+C$.
$x=0$ fournit
$C= 0$ et donc,

\begin{center}
\shadowbox{
$\forall x\in\Rr,\;\Arcsin\left(\frac{x}{\sqrt{x^2+1}}\right)=\Arctan x.$
}
\end{center}
\textbf{2ème solution.} Pour $x$ réel donné, posons $\theta=\Arctan x$. $\theta$ est dans
$\left]-\frac{\pi}{2},\frac{\pi}{2}\right[$ et $x=\tan\theta$.

\begin{align*}
\frac{x}{\sqrt{x^2+1}}&=\frac{\tan\theta}{\sqrt{1+\tan^2\theta}}=\sqrt{\cos^2\theta}\tan\theta=\cos\theta\tan\theta\;
(\mbox{car}\;\cos\theta>0)\\
 &=\sin\theta
\end{align*}
et donc

\begin{align*}
f_1(x)&=\Arcsin(\sin\theta)=\theta\;(\mbox{car}\;\theta\;\mbox{est dans}\;\left]-\frac{\pi}{2},\frac{\pi}{2}\right[)\\
 &=\Arctan x.
\end{align*}
 \item  \textbf{1ère solution.} Pour tout réel $x$, $-1<-1+\frac{2}{1+x^2}=\frac{1-x^2}{1+x^2}\leq-1+2=1$ (avec égalité
si et seulement si $x=0$). $f_2$ est donc définie et continue sur $\Rr$, dérivable sur $\Rr^*$. Pour tout réel $x$ non nul,

$$f_2'(x)=\frac{-2x(1+x^2)-2x(1-x^2)}{(1+x^2)^2}\frac{-1}{\sqrt{1-\left(\frac{1-x^2}{1+x^2}\right)^2}}=\frac{4x}{1+x
^2}\frac{1}{\sqrt{4x^2}}=\frac{2\varepsilon}{1+x^2}$$
où $\varepsilon$ est le signe de $x$. Donc il existe une
constante réelle $C$ telle que pour tout réel positif $x$, $f_2(x)=2\Arctan x+C$ (y compris $x=0$ puisque $f$ est continue en $0$).\\
$x= 0$ fournit $C=0$ et donc, pour
tout réel positif $x$, $f_2(x)=2\Arctan x$. Par parité,

\begin{center}
\shadowbox{
$\forall x\in\Rr,\;\Arccos\left(\frac{1-x^2}{1+x^2}\right)=2\Arctan|x|.$
}
\end{center}
\textbf{2ème solution.} Soit $x\in\Rr$ puis $\theta=\Arctan x$. $\theta$ est dans
$\left]-\frac{\pi}{2},\frac{\pi}{2}\right[$ et $x=\tan\theta$.

$$\frac{1-x^2}{1+x^2}=\frac{1-\tan^2\theta}{1+\tan^2\theta}=\cos^2\theta(1-\tan^2\theta)=\cos^2\theta-\sin^2\theta=
\cos(2\theta).$$
Donc

$$f_2(x)=\Arccos(\cos(2\theta))=\left\{
\begin{array}{l}
2\theta\;\mbox{si}\;\theta\in\left[0,\frac{\pi}{2}\right[\\
-2\theta\;\mbox{si}\;\theta\in\left]-\frac{\pi}{2},0\right]
\end{array}
\right.
=\left\{
\begin{array}{l}
2\Arctan x\;\mbox{si}\;x\geq0\\
-2\Arctan x\;\mbox{si}\;x\leq0
\end{array}
\right.=2\Arctan|x|.$$

 \item  La fonction $x\mapsto\Arcsin\sqrt{1-x^2}$ est définie et continue sur $[-1,1]$, dérivable sur $[-1,1]\setminus\{0\}$ car
pour $x$ élément de $[-1,1]$, $1-x^2$ est élément de $[0,1]$ et vaut $1$ si et seulement si $x$ vaut $0$.
$\frac{1-x}{1+x}$ est défini et positif si et seulement si $x$ est dans $]-1,1]$, et nul si et seulement si $x=1$.
$f_3$ est donc définie et continue sur $]-1,1]$, dérivable sur $]-1,0[\cup]0,1[$. Pour $x$ dans $]-1,0[\cup]0,1[$, on
note $\varepsilon$ le signe de $x$ et on
a~:~

$$f_3'(x)=-\frac{x}{\sqrt{1-x^2}}\frac{1}{\sqrt{1-(1-x^2)}}-\frac{-(1+x)-(1-x)}{(1+x)^2}
\frac{1}{2\sqrt{\frac{1-x}{1+x}}}\frac{1}{1+\frac{1-x}{1+x}}=-\frac{\varepsilon}{\sqrt{1-x^2}}+\frac{1}{2}\frac{1
}{\sqrt{1-x^2}}.$$
Si $x$ est dans $]0,1[$, $f_3'(x)=-\frac{1}{2}\frac{1}{\sqrt{1-x^2}}=(-\frac{1}{2}\Arcsin)'(x)$.
Donc, il existe un réel $C$ tel que, pour tout $x$ de $[0,1]$ (par continuité) $f_3(x)=-\frac{1}{2}\Arcsin x+C$. $x=1$
fournit $C=\frac{\pi}{4}$. Donc,

\begin{center}
\shadowbox{
$\forall x\in[0,1],\;f_3(x)=\frac{\pi}{4}-\frac{1}{2}\Arcsin x=\frac{1}{2}\Arccos x.$
}
\end{center}
Si $x$ est dans $]-1,0[$, $f_3'(x)=\frac{3}{2}\frac{1}{\sqrt{1-x^2}}=(\frac{3}{2}\Arcsin)'(x)$. Donc il existe un
réel $C'$ tel que, pour tout $x$ de $]-1,0]$ (par continuité) $f_3(x)=\frac{3}{2}\Arcsin x+C'$.
$x=0$ fournit $\frac{\pi}{2}-\frac{\pi}{4}=C'$. Donc,

\begin{center}
\shadowbox{
$\forall x\in]-1,0],\;f_3(x)=\frac{3}{2}\Arcsin x+\frac{\pi}{4}.$
}
\end{center}
 \item  $f_4$ est dérivable sur $\mathcal{D}=\Rr\setminus\{-1,0\}$ et pour $x$ élément de $\mathcal{D}$, on a~:~

\begin{align*}
f_4'(x)&=-\frac{1}{x^3}\frac{1}{1+\frac{1}{4x^4}}-\frac{(x+1)-x}{(x+1)^2}\frac{1}{1+\frac{x^2}{(x+1)^2}}+\frac{x
-(x-1)}{x^2}\frac{1}{1+\frac{(x-1)^2}{x^2}}\\
 &=-\frac{4x}{4x^4+1}-\frac{1}{2x^2+1+2x}+\frac{1}{2x^2+1-2x}=-\frac{4x
}{4x^4+1}+\frac{4x}{(2x^2+1)^2-4x^2}=0.
\end{align*}
$f_4$ est donc constante sur chacun des trois intervalles
$]-\infty,-1[$, $]-1,0[$ et $]0,+\infty[$. Pour $x>0$, $f(x)=f(1)=0$. Pour $-1<x<0$,
$\displaystyle f(x)=\lim_{
\substack{
t\rightarrow-1\\
t>-1}
}f(t)=\Arctan\frac{1}{2}-(-\frac{\pi}{2})+\Arctan2=\frac{\pi}{2}+\frac{\pi}{2}=\pi$.
Pour $x<-1$, $f(x)=\lim_{t\rightarrow -\infty}f(t)=0$ et donc

\begin{center}
\shadowbox{
$
\forall x\in\Rr\setminus\{-1;0\},\;f_4(x)=\left\{
\begin{array}{l}
0\;\mbox{si}\;x\in]-\infty,-1[\cup]0,+\infty[\\
\pi\;\mbox{si}\;x\in]-1,0[
\end{array}\right.
.$
}
\end{center}
\end{enumerate}
\fincorrection
\correction{005089}
 $0\leq\Arctan\frac{1}{2}+\Arctan\frac{1}{5}<\Arctan1+\Arctan1=\frac{\pi}{2}$ et

$$\tan\left(\Arctan\frac{1}{2}+\Arctan\frac{1}{5}\right)=\frac{\frac{1}{2}+\frac{1}{5}}{1-\frac{1}{2}\times\frac{1}{5}}
=\frac{7}{9}.$$
Comme $\Arctan\frac{1}{2}+\Arctan\frac{1}{5}\in[0,\frac{\pi}{2}[$, on a donc
$\Arctan\frac{1}{2}+\Arctan\frac{1}{5}=\Arctan\frac{7}{9}$. De même,
$\Arctan\frac{7}{9}+\Arctan\frac{1}{8}\in[0,\frac{\pi}{2}[$ et

$$\tan\left(\Arctan\frac{7}{9}+\Arctan\frac{1}{8}\right)=\frac{\frac{7}{9}+\frac{1}{8}}{1-\frac{7}{9}\times\frac{1}{8}}=\frac{65}{65}
=1,$$
et donc $\Arctan\frac{7}{9}+\Arctan\frac{1}{8}=\Arctan1=\frac{\pi}{4}$. Finalement,

\begin{center}
\shadowbox{
$\Arctan\frac{1}{2}+\Arctan\frac{1}{5}+\Arctan\frac{1}{8}=\frac{\pi}{4}.$
}
\end{center}
\fincorrection
\correction{005090}
 (On va retrouver le résultat de l'exercice \ref{exo:suprou2} dans un cas particulier)
 
 
 Soient $a$ et $b$ deux réels positifs. Alors, $\Arctan a\in\left[0,\frac{\pi}{2}\right[$, $\Arctan b\in\left[0,\frac{\pi}{2}\right[$ et donc, $\Arctan a-\Arctan b\in\left]-\frac{\pi}{2},\frac{\pi}{2}\right[$. De plus, 

$$\tan(\Arctan a-\Arctan b)=\frac{\tan(\Arctan a)-\tan(\Arctan b)}{1+\tan(\Arctan a)\tan(\Arctan b)}=\frac{a-b}{1+ab},$$
et donc, puisque $\Arctan a-\Arctan b\in]-\frac{\pi}{2},\frac{\pi}{2}[$,

\begin{center}
\shadowbox{
$\forall a\geq0,\;\forall b\geq0,\;\Arctan a-\Arctan b
=\Arctan\left(\frac{a-b}{1+ab}\right).$
}
\end{center}
Soit alors $k$ un entier naturel non nul.
$\Arctan\frac{2}{k^2}=\Arctan\frac{(k+1)-(k-1)}{1+(k-1)(k+1)}=\Arctan(k+1)-\Arctan(k-1)$ (puisque $k-1$ et $k+1$ sont
positifs). Par suite, si $n$ est un entier naturel non nul donné,

\begin{align*}
u_n&=\sum_{k=1}^{n}\Arctan\frac{2}{k^2}=\sum_{k=1}^{n}(\Arctan(k+1)-\Arctan(k-1))=\sum_{k=2}^{n+1}\Arctan
k-\sum_{k=0}^{n-1}\Arctan k\\
 &=\Arctan(n+1)+\Arctan n-\frac{\pi}{4}.
\end{align*}
La limite de $u_n$ vaut donc
$\frac{\pi}{2}+\frac{\pi}{2}-\frac{\pi}{4}=\frac{3\pi}{4}$.

\begin{center}
\shadowbox{
$\lim_{n\rightarrow +\infty}u_n=\frac{3\pi}{4}$.
}
\end{center}
\fincorrection
\correction{005092}
\begin{enumerate}
 \item  $f$ est définie et dérivable sur $\mathcal{D}=\Rr\setminus\{\frac{1}{2}\}$.
 \item  Pour $x$ élément de $\mathcal{D}$,

$$f'(x)=2x\Arctan\frac{1}{2x-1}+(x^2-1)\frac{-2}{(2x-1)^2}\frac{1}{1+\frac{1}{(2x-1)^2}}=2x\Arctan\frac{1}{2x-1}-
\frac{x^2-1}{2x^2-2x+1}.$$
De plus, pour $x$ non nul~:~$f'(x)=2xg(x)$ où
$g(x)=\Arctan\frac{1}{2x-1}-\frac{1}{2x}\frac{x^2-1}{2x^2-2x+1}$.
 \item  Pour $x$ élément de $\mathcal{D}\setminus\{0\}$,

\begin{align*}
g'(x)&=-\frac{1}{2x^2-2x+1}-\frac{1}{2}\frac{2x(2x^3-2x^2+x)-(x^2-1)(6x^2-4x+1)}{x^2(2x^2-2x+1)^2}\\
 &=\frac{-2x^2(2x^
2-2x+1)+2x^4-7x^2+4x-1}{2x^2(2x^2-2x+1)^2}=-\frac{2x^4-4x^3+9x^2-4x+1}{2x^2(x^2-2x+1)^2}.
\end{align*}
Maintenant,

$$2x^4-4x^3+9x^2-4x+1=2x^2(x-1)^2+7x^2-4x+1=2x^2(x-1)^2+7\left(x-\frac{2}{7}\right)^2+\frac{3}{7}>0.$$
Donc, $g$ est
strictement décroissante sur $]-\infty,0[$, sur $\left]0,\frac{1}{2}\right[$ et sur $\left]\frac{1}{2},+\infty\right[$. En $+\infty$, $g(x)$
tend vers $0$. Donc $g$ est strictement positive sur $\left]\frac{1}{2},+\infty\right[$. Quand $x$ tend vers $\frac{1}{2}$ par
valeurs inférieures, $g$ tend vers $-\frac{\pi}{2}+\frac{3}{2}<0$ et quand $x$ tend vers $0$ par valeurs
supérieures, $g(x)$ tend vers $+\infty$. Donc $g$ s'annule une et une seule fois sur l'intervalle $]0,\frac{1}{2}[$ en
un certain réel $x_0$ de $\left]0,\frac{1}{2}\right[$. $g$ est de plus strictement négative sur $\left]x_0,\frac{1}{2}\right[$ 
et strictement positive sur $]0,x_0[$. Quand $x$ tend vers $-\infty$, $g(x)$ tend vers $0$. Donc $g$ est strictement
négative sur $]-\infty,0[$. Enfin, puisque $f'(x)=2xg(x)$ pour $x\neq0$, on a les résultats suivants~:
sur
$]-\infty,0[$, $f'> 0$, sur $]0,x_0[$, $f'> 0$, sur $\left]x_0,\frac{1}{2}\right[$, $f'< 0$, sur $\left]\frac{1}{2},+\infty\right[$, $f'>
0$. Comme $f'(0)=1>0$, on a donc~:~sur $]-\infty,x_0[$, $f'>0$,
sur $\left]x_0,\frac{1}{2}\right[$, $f'< 0$ et sur $\left]\frac{1}{2},+\infty\right[$, $f'> 0$. $f$ est strictement croissante sur
$]-\infty,x_0]$ et sur $\left]\frac{1}{2},+\infty\right[$ et est strictement décroissante sur $\left[x_0,\frac{1}{2}\right[$.
\end{enumerate}
\fincorrection
\correction{005095}
\begin{enumerate}
 \item  Pour tout réel $x$ de $[-1,1]$, $\sin(2\Arcsin x)=2\sin(\Arcsin x)\cos(\Arcsin x)=2x\sqrt{1-x^2}$.
 \item  Pour tout réel $x$ de $[-1,1]$, $\cos(2\Arccos x)=2\cos^2(\Arccos x)-1=2x^2-1$.
 \item  Pour tout réel $x$ de $[-1,1]$,
$\sin^2(\frac{1}{2}\Arccos x)=\frac{1}{2}(1-\cos(\Arccos x))=\frac{1-x}{2}$.
 \item  Soit $x\in\Rr$.

$$\sqrt{x^2+1}>\sqrt{x^2}=|x|=\mbox{Max}\{x,-x\}.$$
Donc, $\sqrt{x^2+1}+x>0$ et $\sqrt{x^2+1}-x>0$. L'expression proposée existe pour tout réel $x$. De plus,

$$\ln(\sqrt{x^2+1}+x)+\ln(\sqrt{x^2+1}-x)=\ln\left((\sqrt{x^2+1}+x)(\sqrt{x^2+1}-x)\right)=\ln(x^2+1-x^2)=\ln1=0.$$
 \item  L'expression proposée est définie sur $\Rr^*$ et impaire. Soit alors $x>0$.

\begin{align*}
\Argsh\left(\frac{x^2-1}{2x}\right)&=\ln\left(\frac{x^2-1}{2x}+\sqrt{\frac{(x^2-1)^2}{(2x)^2}+1}\right)
=\ln\left(\frac{1}{2x}(x^2-1+\sqrt{x^4-2x^2+1+4x^2})\right)\\
 &=\ln\left(\frac{1}{2x}(x^2-1+\sqrt{(x^2+1)^2})\right)=\ln\left(\frac{1}{2x}(x^2-1+x^2+1)\right)=\ln x
\end{align*}
Par imparité, si $x<0$, $\Argsh\left(\frac{x^2-1}{2x}\right)=-\ln(-x)$. En résumé, en notant $\varepsilon$ le signe de
$x$,

\begin{center}
\shadowbox{
$\forall x\in\Rr^*,\;\Argsh\left(\frac{x^2-1}{2x}\right)=\varepsilon\ln|x|.$
}
\end{center}
 \item  L'expression proposée existe si et seulement si $2x^2-1\in[1,+\infty[$ ou encore $x^2\in[1,+\infty[$ ou enfin
$x\in]-\infty,-1]\cup[1,+\infty[$. Cette expression est paire. Soit donc $x\in[1,+\infty[$.

\begin{align*}
\Argch(2x^2-1)&=\ln(2x^2-1+\sqrt{(2x^2-1)^2-1})=\ln(2x^2-1+2x\sqrt{x^2-1})=\ln\left(\left(x+\sqrt{x^2-1}\right)^2\right)\\
 &=2\ln\left(x+\sqrt{x^2-1}\right)=2\Argch x
\end{align*}
Par parité, on en déduit que

\begin{center}
\shadowbox{
$\forall x\in]-\infty,-1]\cup[1,+\infty[,\;\Argch(2x^2-1)=2\Argch|x|.$
}
\end{center}
 \item  Soit $x\in\Rr$.

\begin{align*}
\Argth\sqrt{\frac{\ch x-1}{\ch x+1}}\;\mbox{existe}&\Leftrightarrow\ch x+1\neq0\;\mbox{et}\;\frac{\ch x-1}{\ch
x+1}\geq0\;\mbox{et}\;\sqrt{\frac{\ch x-1}{\ch x+1}}\in]-1,1[\\
 &\Leftrightarrow\frac{\ch x-1}{\ch x+1}\in[0,1[
\end{align*}
Mais, d'une part, $\frac{\ch x-1}{\ch x+1}\geq0$ et d'autre part, $\frac{\ch x-1}{\ch x+1}=\frac{\ch x+1-2}
{\ch x+1}=1-\frac{2}{\ch x+1}<1$. L'expression proposée existe donc pour tout réel $x$ et est paire. Ensuite, pour $x$ réel positif, on a

\begin{align*}
\frac{1+\sqrt{\frac{\ch x-1}{\ch x+1}}}{1-\sqrt{\frac{\ch x-1}{\ch x+1}}}
&=\frac{\sqrt{\ch x+1}+\sqrt{\ch x-1}}{\sqrt{\ch x+1}-\sqrt{\ch x-1}}=
\frac{(\sqrt{\ch x+1}+\sqrt{\ch x-1})^2}{(\ch x+1)-(\ch x-1)}=\frac{2\ch x+2\sqrt{\ch^2x-1}}{2}\\
 &=\ch x+\sqrt{\sh^2x}=\ch x+|\sh x|=\ch x+\sh x=e^x
\end{align*}
Par suite, $x$ étant toujours positif,

$$\Argth\sqrt{\frac{\ch x-1}{\ch x+1}}=\frac{1}{2}\ln(e^x)=\frac{x}{2}.$$
Par parité, on a alors

\begin{center}
\shadowbox{
$\forall x\in\Rr,\;\Argth\left(\sqrt{\frac{\ch x-1}{\ch x+1}}\right)=\frac{|x|}{2}.$
}
\end{center}
(Remarque. Pour 5), 6) et 7), on peut aussi dériver chaque expression)
 \item  Pour $x>0$,

\begin{align*}
\frac{\ch(\ln x)+\sh(\ln x)}{x}&=\frac{1}{2x}\left(x+\frac{1}{x}+x-\frac{1}{x}\right)=1.
\end{align*}
\end{enumerate}
\fincorrection
\correction{005096}
\begin{enumerate}
 \item  $\ch x=2\Leftrightarrow x=\pm\Argch2=\pm\ln(2+\sqrt{2^2-1})=\pm\ln(2+\sqrt{3})$. Les solutions sont $\ln(2+\sqrt{3})$ et
$-\ln(2+\sqrt{3})$ (ou encore $\ln(2-\sqrt{3})$ car $(2+\sqrt{3})(2-\sqrt{3})=1$).
 \item  Une solution est nécessairement dans $\left[-\frac{1}{2},\frac{1}{2}\right]$. Soit donc $x\in\left[-\frac{1}{2},\frac{1}{2}\right]$.

\begin{align*}
\Arcsin(2x)=\Arcsin x+\Arcsin(x\sqrt{2})&\Rightarrow\sin(\Arcsin(2x))=\sin(\Arcsin x+\Arcsin(x\sqrt{2}))\\
 &\Leftrightarrow2x=x\sqrt{1-(x\sqrt{2})^2}+x\sqrt{2}\sqrt{1-x^2}\Leftrightarrow x=0\;\mbox{ou}\;\sqrt{1-2x^2}+\sqrt{2-2x^2}=2\\
 &\Leftrightarrow x=0\;\mbox{ou}\;1-2x^2+2-2x^2+2\sqrt{(1-2x^2)(2-2x^2)}=4\\
 &\Leftrightarrow x=0\;\mbox{ou}\;2\sqrt{(1-2x^2)(2-2x^2)}=1+4x^2\\
 &\Leftrightarrow x=0\;\mbox{ou}\;4(4x^4-6x^2+2)=(4x^2+1)^2\\
 &\Leftrightarrow x=0\;\mbox{ou}\;32x^2=7\Leftrightarrow x=0\;\mbox{ou}\;x=\sqrt{\frac{7}{32}}\;\mbox{ou}\;x=-\sqrt{\frac{7}{32}}\\
\end{align*}
Réciproquement, pour chacun des ces trois nombres $x$, la seule implication écrite est une équivalence si $x$ est dans
$[-\frac{1}{2},\frac{1}{2}]$ (ce qui est le cas puisque
$\left(\pm\sqrt{\frac{7}{32}}\right)^2=\frac{14}{64}\leq\frac{16}{64}=(\frac{1}{2})^2$) et $\Arcsin x+\Arcsin(x\sqrt{2})$ est
dans $[-\frac{\pi}{2},\frac{\pi}{2}]$. Mais,

$$0\leq\Arcsin\sqrt{\frac{7}{32}}+\Arcsin(\sqrt{\frac{7}{32}}\sqrt{2})=\Arcsin\sqrt{\frac{7}{32}}+\Arcsin\sqrt{\frac
{7}{16}}\leq2\Arcsin\sqrt{\frac{8}{16}}=2\Arcsin\frac{1}{\sqrt{2}}=\frac{\pi}{2}$$
et donc $\Arcsin\sqrt{\frac{7}{32}}+\Arcsin(\sqrt{\frac{7}{32}}\sqrt{2}\in[0,\frac{\pi}{2}]$. De même, par parité,
$\Arcsin(-\sqrt{\frac{7}{32}})+\Arcsin(-\sqrt{\frac{7}{32}}\sqrt{2})\in[-\frac{\pi}{2},0]$ ce qui achève la
résolution.

\begin{center}
\shadowbox{
$\mathcal{S}=\left\{0,\frac{\sqrt{14}}{8},-\frac{\sqrt{14}}{8}\right\}.$
}
\end{center}

 \item  Soit $x\in\Rr$. $\Arcsin x$ existe si et seulement si $x\in[-1,1]$. Ensuite,

\begin{align*}
\Arcsin(2x\sqrt{1-x^2})\;\mbox{existe}&\Leftrightarrow x\in[-1,1]\;\mbox{et}\;2x\sqrt{1-x^2}\in[-1,1]\\
 &\Leftrightarrow x\in[-1,1]\;\mbox{et}\;4x^2(1-x^2)\in[0,1]\Leftrightarrow x\in[-1,1]\;\mbox{et}\;4x^2(1-x^2)\leq1\\
 &\Leftrightarrow x\in[-1,1]\;\mbox{et}\;4x^4-4x^2+1\geq0\Leftrightarrow x\in[-1,1]\;\mbox{et}\;(2x^2-1)^2\geq0\\
 &\Leftrightarrow x\in[-1,1]
\end{align*}
Pour $x\in[-1,1]$, $\sin(2\Arcsin(x))=2\sin(\Arcsin x)\cos(\Arcsin x)=2x\sqrt{1-x^2}=\sin(\Arcsin(2x\sqrt{1-x^2}))$, et
de plus,\\
$\Arcsin(2x\sqrt{1-x^2})\in[-\frac{\pi}{2},\frac{\pi}{2}]$. Par suite,

\begin{align*}
x\;\mbox{solution}&\Leftrightarrow x\in[-1,1]\;\mbox{et}\;2\Arcsin(x)\in\left[-\frac{\pi}{2},\frac{\pi}{2}\right]\\
 &\Leftrightarrow x\in[-1,1]\;\mbox{et}\;\Arcsin(x)\in\left[-\frac{\pi}{4},\frac{\pi}{4}\right]\Leftrightarrow
x\in\left[-\frac{1}{\sqrt{2}},\frac{1}{\sqrt{2}}\right].
\end{align*}

\begin{center}
\shadowbox{
$\mathcal{S}=\left[-\frac{1}{\sqrt{2}},\frac{1}{\sqrt{2}}\right].$
}
\end{center}
\end{enumerate}
\fincorrection
\correction{005315}
Il faut prendre garde au fait que les nombres $x_k=\cotan^2(\frac{\pi}{2n}+\frac{k\pi}{n})$ ne sont pas nécessairement deux à deux distincts.

\begin{itemize}
\item[1er cas.] Si $n$ est pair, posons $n=2p$, $p\in\Nn^*$.

\begin{align*}\ensuremath
S_n&=\sum_{k=0}^{p-1}\cotan^2(\frac{\pi}{4p}+\frac{k\pi}{2p})+\sum_{k=p}^{2p-1}\cotan^2(\frac{\pi}{4p}+\frac{k\pi}{2p})\\
 &=\sum_{k=0}^{p-1}\cotan^2(\frac{\pi}{4p}+\frac{k\pi}{2p})+\sum_{k=0}^{p-1}\cotan^2(\frac{\pi}{4p}+\frac{(2p-1-k)\pi}{2p})
\end{align*}
 
Or, $\cotan^2(\frac{\pi}{4p}+\frac{(2p-1-k)\pi}{2p})=\cotan^2(\pi-\frac{\pi}{4p}-\frac{k\pi}{2p})=\cotan^2(\frac{\pi}{4p}+\frac{k\pi}{2p})$ et donc $S_n=2\sum_{k=0}^{p-1}\cotan^2(\frac{\pi}{4p}+\frac{k\pi}{2p})$.

Mais cette fois ci, 

$$0\leq k\leq p-1\Rightarrow 0<\frac{\pi}{4p}+\frac{k\pi}{2p}\leq\frac{\pi}{4p}+\frac{(p-1)\pi}{2p}=\frac{(2p-1)\pi}{4p}<\frac{2p\pi}{4p}=\frac{\pi}{2}.$$

et comme, la fonction $x\mapsto\cotan^2x$ est strictement décroissante sur $]0,\frac{\pi}{2}[$, les $x_k$, $0\leq k\leq p-1$, sont deux à deux distincts.

Pour $0\leq k\leq p-1$, posons $y_k=\cotan(\frac{\pi}{4p}+\frac{k\pi}{2p})$.

\begin{align*}\ensuremath
y_k&=i\frac{e^{(2k+1)i\pi/4p}+1}{e^{(2k+1)i\pi/4p}-1}\Rightarrow e^{(2k+1)i\pi/4p}(y-k-i)=y_k+i\\
 &\Rightarrow(y_k+i)^{2p}=e^{(2k+1)i\pi}(y_k-i)^{2p}=(-1)^{2k+1}(y_k-i)^{2p}=-(y_k-i)^{2p}\\
 &\Rightarrow(y_k+i)^{2p}+(y_k-i)^{2p}=0\Rightarrow2(y_k^{2p}-C_{2p}^2y_k^{2p-2}+...+(-1)^p)=0\\
 &\Rightarrow x_k^p-C_{2p}^2x_k^{p-1}+...+(-1)^p=0.
\end{align*}

Les $p$ nombres deux à deux distincts $x_k$ sont racines de l'équation de degré $p$~:~$z^p-C_{2p}^{2}z^{p-1}+...+(-1)^p=0$ qui est de degré $p$. On en déduit que 

$$S_n=2\sum_{k=0}^{p-1}x_k=2C_{2p}^{2}=n(n-1).$$

\item[2ème cas.] Si $n$ est impair, posons $n=2p+1$, $p\in\Nn$.

\begin{align*}\ensuremath
S_n&=\sum_{k=0}^{p-1}\cotan^2(\frac{\pi}{2(2p+1)}+\frac{k\pi}{2p+1})+\cotan^2\frac{\pi}{2}+\sum_{k=p+1}^{2p}\cotan^2(\frac{\pi}{2(2p+1)}+\frac{k\pi}{2p+1})\\
 &=2\sum_{k=0}^{p-1}\cotan^2(\frac{\pi}{2(2p+1)}+\frac{k\pi}{2p+1})
\end{align*}

La même démarche amène alors à $S_n=2C_{2p+1}^{2}=n(n-1)$.
\end{itemize}

Dans tous les cas, 

$$\sum_{k=0}^{n-1}\cotan^2(\frac{\pi}{2n}+\frac{k\pi}{n})=n(n-1).$$
\fincorrection
\correction{005852}
Posons $I_0=\left[0, \frac{\pi}{2}\right[$ puis pour $n\in\Nn^*$, $I_n=\left]- \frac{\pi}{2}+n\pi, \frac{\pi}{2}+n\pi\right[$ et enfin $D =\displaystyle\bigcup_{n\in\Nn}I_n$.

Pour $x\in D$, posons $f(x) =\tan x -x$. La fonction $f$ est dérivable sur $D$ et pour $x\in D$, $f'(x) =\tan^2x$. La fonction $f$ est ainsi strictement croissante sur chaque $I_n$ et s'annule donc au plus une fois dans chaque $I_n$.

$f(0) = 0$ et donc $f$ s'annule exactement une fois dans $I_0$ en $x_0=0$.

Pour $n\in\Nn^*$, $f$ est continue sur $I_n$ et de plus $f\left(\left(- \frac{\pi}{2}+n\pi\right)^+\right)\times f\left(\left( \frac{\pi}{2}+n\pi\right)^-\right)=-\infty\times+\infty<0$. D'après le théorème des valeurs intermédiaires, $f$ s'annule au moins une fois dans $I_n$ et donc exactement une fois dans $I_n$.

L'équation $\tan x = x$ admet donc dans chaque intervalle $I_n$, $n\in\Nn$, une et une seule solution notée $x_n$. De plus, $\forall n\geqslant1$, $f(n\pi) =-n\pi < 0$ et donc $x_n\in]n\pi, \frac{\pi}{2}+n\pi[$.

Pour $n\geqslant1$, $n\pi< x_n < \frac{\pi}{2}+n\pi$ et donc $\lim_{n \rightarrow +\infty}x_n=+\infty$ puis $x_n\underset{n\rightarrow+\infty}{\sim}n\pi$ et même

\begin{center}
$x_n\underset{n\rightarrow+\infty}=n\pi+ O(1)$.
\end{center}

Ensuite, puisque $x_n-n\pi\in\left]0, \frac{\pi}{2}\right[$ et que $x_n=\tan(x_n)=\tan(x_n-n\pi)$, $x_n-n\pi=\Arctan(x_n)\underset{n\rightarrow+\infty}{\rightarrow} \frac{\pi}{2}$. Donc

\begin{center}
$x_n\underset{n\rightarrow+\infty}{=} n\pi+ \frac{\pi}{2}+o(1)$.
\end{center}

Posons $y_n = x_n-n\pi- \frac{\pi}{2}$. Alors d'après ce qui précède, $y_n\in\left]- \frac{\pi}{2},0\right[$ et $y_n\underset{n\rightarrow+\infty}{=} o(1)$. De plus, l'égalité $\tan(x_n)=x_n$ fournit $\tan(n\pi+ \frac{\pi}{2}+y_n) =n\pi+ \frac{\pi}{2}+y_n$ ou encore 

\begin{center}
$n\pi+ \frac{\pi}{2}+y_n = -\mathrm{cotan}(y_n)$.
\end{center}

Puisque $y_n\underset{n\rightarrow+\infty}{=}o(1)$, on obtient $n\underset{n\rightarrow+\infty}{\sim}- \frac{1}{y_n}$  ou encore $y_n\underset{n\rightarrow+\infty}{=} - \frac{1}{n\pi}+o\left( \frac{1}{n}\right)$. Donc

\begin{center}
$x_n\underset{n\rightarrow+\infty}{=} n\pi+ \frac{\pi}{2}- \frac{1}{n\pi}+o\left( \frac{1}{n}\right)$.
\end{center}

Posons $z_n = y_n+ \frac{1}{n\pi}= x_n-n\pi- \frac{\pi}{2}+ \frac{1}{n\pi}$. D'après ce qui précède, $\tan\left(- \frac{1}{n\pi}+z_n\right)=- \frac{1}{n\pi+ \frac{\pi}{2}- \frac{1}{n\pi}+z_n}$ et aussi $z_n\underset{n\rightarrow+\infty}{=}o\left( \frac{1}{n}\right)$. On en déduit que

\begin{center}
$z_n= \frac{1}{n\pi}-\Arctan\left( \frac{1}{n\pi+ \frac{\pi}{2}- \frac{1}{n\pi}+z_n}\right)\underset{n\rightarrow+\infty}{=} \frac{1}{n\pi}-\Arctan\left( \frac{1}{n\pi}- \frac{1}{2\pi n^2}+o\left( \frac{1}{n^2}\right)\right)\underset{n\rightarrow+\infty}{=} \frac{1}{2\pi n^2}+o\left( \frac{1}{n^2}\right)$.
\end{center}

Finalement

\begin{center}
\shadowbox{
$x_n\underset{n\rightarrow+\infty}{=} n\pi+ \frac{\pi}{2}- \frac{1}{n\pi}+ \frac{1}{2\pi n^2}+o\left( \frac{1}{n^2}\right)$.
}
\end{center}

\fincorrection
\correction{005853}
\textbf{1ère solution.} Soit $z\in\Cc$. Posons $z = x+iy$ où $(x,y)\in\Rr^2$ et $1+ \frac{r}{n}=r_ne^{i\theta}$ où $r_n\geqslant0$ et $\theta_n\in]-\pi,\pi]$ de sorte que 

\begin{center}
$\left(1+ \frac{z}{n}\right)^n=r_n^n\;e^{in\theta_n}$.
\end{center}

Puisque $1+ \frac{z}{n}$ tend vers $1$ quand $n$ tend vers $+\infty$, pour $n$ assez grand on a $r_n > 0$ et $\theta_n\in\left]- \frac{\pi}{2}, \frac{\pi}{2}\right[$. Mais alors pour $n$ assez grand

\begin{center}
$r_n =\sqrt{\left(1+ \frac{x}{n}\right)^2+\left( \frac{y}{n}\right)^2}$ et $\theta_n=\Arctan\left( \frac{ \frac{y}{n}}{1+ \frac{x}{n}}\right)$.
\end{center}

Maintenant, $r_n^n=\exp\left( \frac{n}{2}\ln\left(\left(1+ \frac{x}{n}\right)^2+\left( \frac{y}{n}\right)^2 \right)\right)\underset{n\rightarrow+\infty}{=}\exp\left( \frac{n}{2}\ln\left(1+ \frac{2x}{n}+o\left( \frac{1}{n}\right)\right)\right) \underset{n\rightarrow+\infty}{=} \exp(x+o(1))$ et donc $r_n^n$ tend vers $e^x$ quand $n$ tend vers $+\infty$.

Ensuite $n\theta_n\underset{n\rightarrow+\infty}{=}n\Arctan\left( \frac{ \frac{y}{n}}{1+ \frac{x}{n}}\right)\underset{n\rightarrow+\infty}{=}n\Arctan\left( \frac{y}{n}+o\left( \frac{1}{n}\right)\right)\underset{n\rightarrow+\infty}{=}y+o(1)$ et donc $n\theta_n$ tend vers $y$ quand $n$ tend vers $+\infty$.

Finalement, $\left(1+ \frac{z}{n}\right)^n =r_n^n\;e^{in\theta_n}$ tend vers $e^x\times e^{iy}= e^z$.

\begin{center}
\shadowbox{
$\forall z\in\Cc$, $\lim_{n \rightarrow +\infty}\left(1+ \frac{z}{n}\right)^n=e^z$.
}
\end{center}

\textbf{2ème solution.} Le résultat est connu quand $z$ est réel. Soit $z\in\Cc$. Soit $n\in\Nn^*$.

\begin{center}
$\left|\sum_{k=0}^{n} \frac{z^k}{k!}-\left(1+ \frac{z}{n}\right)^n\right| =\left|\sum_{k=0}^{n}\left( \frac{1}{k!}- \frac{C_n^k}{n^k}\right)z^k\right|\leqslant\sum_{k=0}^{n}\left| \frac{1}{k!}- \frac{C_n^k}{n^k}\right||z|^k $.
\end{center}

Maintenant, $\forall k\in\llbracket0,n\rrbracket$, $ \frac{1}{k!}- \frac{C_n^k}{n^k}= \frac{1}{k!}\left(1-  \frac{\overbrace{n\times(n-1)\times\ldots\times(n-k+1)}^{k}}{\underbrace{n\times n\times\ldots\times n}_{k}}\right)\geqslant 0$. Donc,

\begin{center}
$\sum_{k=0}^{n}\left| \frac{1}{k!}- \frac{C_n^k}{n^k}\right||z|^k=\sum_{k=0}^{n} \frac{|z|^k}{k!}-\left(1+ \frac{|z|^n}{n}\right)^n\underset{n\rightarrow+\infty}{\rightarrow}e^{|z|}-e^{|z|}= 0$.
\end{center}

On en déduit que $\sum_{k=0}^{n} \frac{z^k}{k!}-\left(1+ \frac{z}{n}\right)^n$ tend vers $0$ quand $n$ tend vers $+\infty$ et puisque $\sum_{k=0}^{n} \frac{z^k}{k!}$ tend vers $e^z$ quand $n$ tend vers $+\infty$, il en est de même de $\left(1+ \frac{z}{n}\right)^n$.
\fincorrection
\correction{006973}
Posons $f(x)=\Arctan\left(\frac{1}{2x^2}\right)-\Arctan\left(\frac{x}{x+1}\right)+\Arctan\left(\frac{x-1}{x}\right)$ 
pour tout $x>0$. La fonction $f$ est dérivable, et
\begin{eqnarray*}
f'(x)
 &=& \frac{-\frac{2}{2x^3}}{1+\left(\frac{1}{2x^2}\right)^2}-
\frac{\frac{1}{(1+x)^2}}{1+\left( \frac{x}{x+1} \right)^2}
+\frac{\frac{1}{x^2}}{1+\left( \frac{x-1}{x} \right)^2}\\
 &=& \frac{-4x}{4x^4+1}-\frac{1}{(1+x)^2+x^2}+\frac{1}{x^2+(x-1)^2}\\
 &=& \frac{-4x}{4x^4+1}+\frac{-\big(x^2+(x-1)^2\big)+\big((1+x)^2+x^2\big)}{\big((1+x)^2+x^2\big)\big(x^2+(x-1)^2\big)}\\
 &=& 0
\end{eqnarray*}
Ainsi $f$ est une fonction constante. 
Or $f(x)\xrightarrow[x\to +\infty]{}\Arctan 0-\Arctan 1+\Arctan 1=0$. Donc la constante vaut $0$, 
d'où l'égalité cherchée. 

Alors :
\begin{eqnarray*}
S_n
 &=&\sum_{k=1}^n\Arctan\left(\frac{1}{2k^2}\right)\\
 &=&\sum_{k=1}^n\Arctan\left(\frac{k}{k+1}\right)-\sum_{k=1}^n\Arctan\left(\frac{k-1}{k}\right)
 \quad \text{(par l'identité prouvée)}\\
 &=&\sum_{k=1}^n\Arctan\left(\frac{k}{k+1}\right)-\sum_{k'=0}^{n-1}\Arctan\left(\frac{k'}{k'+1}\right)
 \quad \text{(en posant $k'=k-1$)}\\
 &=&\Arctan\left(\frac{n}{n+1}\right)-\Arctan\left(\frac{0}{0+1}\right)
 \quad \text{(les sommes se simplifient)}\\
 &=&\Arctan\left(1-\frac{1}{n+1}\right)
 \quad \text{(car $\tfrac{n}{n+1} = 1-\tfrac{1}{n+1}$)} \\
\end{eqnarray*} 
Ainsi $S_n\xrightarrow[n\to +\infty]{}\Arctan 1=\frac{\pi}{4}$.
\fincorrection
\correction{006974}\
\begin{enumerate}
\item Si $x>0$, alors $\frac{y}{x}$ est bien défini et $\Arctan\frac{y}{x}$ aussi. 
Comme $x=r\cos\theta$ et $y=r\sin\theta$, on a bien $\frac{y}{x}=\tan\theta$.
Puisque par hypothèse $\theta\in]-\pi,\pi]$ et que l'on a supposé $x>0$,
alors $\cos\theta>0$.
Cela implique $\theta\in]-\frac{\pi}{2},\frac{\pi}{2}[$. 
Donc $\theta=\Arctan(\tan\theta)=\Arctan\frac{y}{x}$. (Attention ! Il est important
d'avoir $\theta\in]-\frac{\pi}{2},\frac{\pi}{2}[$ pour considérer l'identité 
$\Arctan(\tan\theta) = \theta$.)


\item Si $\theta\in]-\pi,\pi[$ alors $\frac{\theta}{2}\in]-\frac{\pi}{2},\frac{\pi}{2}[$, 
donc $\frac{\theta}{2}=\Arctan\left(\tan\frac{\theta}{2}\right)$. Or 
$$\frac{\sin\theta}{1+\cos\theta}=
\frac{2\cos\frac{\theta}{2}\sin\frac{\theta}{2}}{1+\big(2\cos^2\left(\frac{\theta}{2}\right)-1\big)}
=\frac{\sin\frac{\theta}{2}}{\cos\frac{\theta}{2}}=\tan\frac{\theta}{2}$$
d'où $\frac\theta2=\Arctan\left(\tan\frac{\theta}{2}\right)=\Arctan\left(\frac{\sin\theta}{1+\cos\theta}\right)$.

\item Remarquons que $z=x+iy$, supposé non nul, est un nombre réel négatif si et seulement si 
($x=r\cos\theta< 0$ et $y=r\sin\theta=0$), c'est-à-dire $\theta=\pi$. 
Par conséquent, dire que $z$ n'est pas réel négatif ou nul signifie que 
$\theta\in]-\pi,\pi[$. On a alors $x+\sqrt{x^2+y^2}\not=0$ (sinon, on aurait $\sqrt{x^2+y^2}=-x$ 
et donc $y=0$ et $x\le 0$) et 
$$\frac{y}{x+\sqrt{x^2+y^2}}=\frac{r\sin\theta}{r\cos\theta+r}=\frac{\sin\theta}{1+\cos\theta}.$$
Par la question précédente :
$$\theta = 2\Arctan\left(\frac{\sin\theta}{1+\cos\theta}\right) = 2\Arctan\left(\frac{y}{x+\sqrt{x^2+y^2}}\right).$$

\end{enumerate}
\fincorrection
\nocorrection
\correction{000758}
\begin{enumerate}
    \item Si $f$ existe alors pour $x=1$ on a $f(\ch 1)= e$
et pour $x=-1$ on a $f(\ch -1)=f(\ch 1)=1/e$. Une fonction ne peut prendre deux valeurs diff\'erentes au m\^eme point (ici $t=\ch1$).
  \item Notons $X = e^x$, l'\'equation devient 
$$f(X)=\frac{e^x+e^{-x}}{2} = \frac12(X+\frac1X).$$
Comme la fonction exponentielle est une bijection de $\Rr$ sur $]0,+\infty[$, alors l'unique fa\c{c}on de d\'efinir $f$ sur $]0,+\infty[$ est par la formule $f(t)=\frac12(t+\frac1t)$.
  \item Comme $e^x$ est toujours non nul, alors $f$ peut prendre n'importe quelle valeur en $0$. $f(0)=c \in \Rr$ et
$f(t)=\frac12(t+\frac1t)$ pour $t>0$. Il y a une infinit\'e de solutions.
Mais aucune de ces solutions n'est continue car la limite de $f(t)$
quand $t>0$ et $t\rightarrow 0$ est $+\infty$.
\end{enumerate}
\fincorrection
\correction{000759}
\begin{enumerate}
  \item Par la formule du binôme de Newton nous avons $\ch^3 x = \left( \frac{e^x+e^{-x}}{2} \right)^3 = \frac18 \left( e^{3x} +3 e^x + 3e^{-x} + e^{-3x} \right)$. Et de même $\sh^3 x = \left( \frac{e^x-e^{-x}}{2} \right)^3 = \frac18 \left( e^{3x} - 3 e^x + 3e^{-x} - e^{-3x} \right)$.
Donc $e^{-x} (\ch ^{3}x-\sh ^{3}x) = \frac 18 e^{-x} (6 e^x + 2 e^{-3x}) = \frac 34 + \frac 14 e^{-4x}$
qui tend vers $\frac 34$ lorsque $x$ tend vers $+\infty$.



  \item $x- \ln(\ch x) = x - \ln(\frac{e^x + e^{-x}}{2}) = x - \ln(e^x + e^{-x}) + \ln 2 =
x - \ln(e^x(1+e^{-2x})) + \ln 2 = x - x + \ln(1+e^{-2x}) + \ln 2
= \ln(1+e^{-2x}) + \ln 2$.
Lorsque $x\rightarrow +\infty$, $\ln(1+e^{-2x}) \rightarrow 0$ donc $x- \ln(\ch x)\rightarrow \ln 2$.
\end{enumerate}
\fincorrection
\nocorrection
\nocorrection
\nocorrection
\correction{000764}\
\begin{enumerate}
  \item 
  \begin{enumerate}
    \item Remarquons d'abord que, par construction, $t\in]-\frac{\pi}{2},\frac{\pi}{2}[$,
    $t$ est donc dans le domaine de définition de la fonction $\tan$.
    En prenant la tangente de l'égalité $t=\Arctan(\sh x)$ on obtient directement $\tan t=
    \tan\big(\Arctan(\sh x)\big) =\sh x$. 

    \item Ensuite, $\frac{1}{\cos^2t}=1+\tan^2 t= 1 + \tan^2\big(\Arctan(\sh x)\big) 
    = 1+\sh^2x=\ch^2x$. Or la fonction $\ch$ ne prend que des valeurs positives, 
    et $t\in]-\frac{\pi}{2},\frac{\pi}{2}[$ donc $\cos t>0$. Ainsi $\frac{1}{\cos t}=\ch x$. 

    \item Enfin, $\sin t=\tan t \cdot \cos t= \sh x \cdot \frac{1}{\ch x}=\tanh x$.
  \end{enumerate}

\item Puisque $t\in]-\frac{\pi}{2},\frac{\pi}{2}[$, on a 
$0<\frac{t}{2}+\frac{\pi}{4}<\frac{\pi}{2}$, donc 
$\tan\left(\frac{t}{2}+\frac{\pi}{4}\right)$ est bien défini et 
strictement positif. Ainsi $y=\ln \left(\tan\left(\frac{t}{2}+\frac{\pi}{4}\right)\right)$ 
est bien défini.

Ensuite : 
\begin{eqnarray*}
\sh y
  &=& \frac{e^y-e^{-y}}{2}\\
  &=& \tfrac{1}{2}\tan\left(\tfrac{t}{2}+\tfrac{\pi}{4}\right)-
  \tfrac{1}{2}\frac{1}{\tan\left(\frac{t}{2}+\frac{\pi}{4}\right)}\\
  &=&\frac{\sin^2\left(\frac{t}{2}+\frac{\pi}{4}\right)-
 \cos^2\left(\frac{t}{2}+\frac{\pi}{4}\right)}{2\cos\left(\frac{t}{2}+\frac{\pi}{4}\right)\sin\left(\frac{t}{2}+\frac{\pi}{4}\right)}\\
 &=&\frac{-\cos\left(t+\frac{\pi}{2}\right)}{\sin\left(t+\frac{\pi}{2}\right)}
\end{eqnarray*}
car $\sin(2u)=2\cos u\sin u$ et $\cos(2u)=\cos^2u-\sin^2u$. 


Enfin, puisque $\cos\left(t+\frac{\pi}{2}\right)=-\sin t$ et 
$\sin\left(t+\frac{\pi}{2}\right)=\cos t$, on a $\sh y= \frac{\sin t}{\cos t} = \tan t=\sh x$. 
Puisque la fonction $\sh$ est bijective de $\R$ dans $\R$, on en déduit $y=x$.
Conclusion : $x = y = \ln \big(\tan\big(\frac{t}{2}+\frac{\pi}{4}\big)\big)$.
\end{enumerate}
\fincorrection
\nocorrection
\nocorrection
\nocorrection
\nocorrection
\nocorrection
\nocorrection
\nocorrection
\nocorrection
\correction{000773}
\begin{enumerate}
    \item Soit $f(x) = \ln(1+x)-x+x^2/2$ alors $f'(x)=\frac{1}{1+x}
-1+x=\frac{x^2}{1+x} >0$. Donc $f$ est strictement croissante sur $[0,+\infty[$ et comme $f(0)=0$ alors $f(x) > f(0)=0$ pour $x>0$. Ce qui donne l'inégalité recherchée.
  \item De même avec $g(x) = e^x-x-1$, $g'(x)=e^x-1$.
Sur $[0,+\infty[$ $g'(x) \geq 0$ et $g$ est croissante
sur $]-\infty,0]$, $g'(x)\leq 0$ et $g$ est décroissante. Comme
 $g(0)=0$ alors pour tout $x\in\Rr$ $g(x)\geq 0$.
\end{enumerate}
\fincorrection
\nocorrection
\nocorrection
\correction{000776}
$$x^y=y^x \Leftrightarrow e^{y\ln x}= e^{x\ln y}
\Leftrightarrow {y\ln x}= {x\ln y}\Leftrightarrow \frac{\ln x}{x}= \frac{\ln y}{y}$$
(la fonction exponentielle est bijective).
Etudions la fonction $f(x) = \frac{\ln x}{x}$ sur $[1,+\infty[$.
$$f'(x)= \frac{1-\ln x}{x^2},$$
donc $f$ est croissante sur $[1,e]$ et d\'ecroissante
sur $[e,+\infty[$. Donc pour $z \in ]0,f(e)[=]0,1/e[$, l'\'equation $f(x)=z$ a exactement deux solutions, une 
dans $]1,e[$ et une dans $]e,+\infty[$.

Revenons \`a l'\'equation $x^y=y^x$ \'equivalente \`a $f(x)=f(y)$.
Prenons $y$ un entier, nous allons distinguer trois cas :
$y=1$, $y=2$ et $y \geq 3$.
Si $y=1$ alors $f(y)=z=0$ on doit
donc r\'esoudre $f(x)=0$ et alors $x=1$.
Si $y=2$ alors il faut r\'esoudre l'\'equation $f(x) = \frac{\ln 2}{2} \in ]0,1/e[$.
Alors d'apr\`es l'\'etude pr\'ec\'edente, il existe deux solutions
une sur $]0,e[$ qui est $x=2$ (!) et une sur $]e,+\infty[$
qui est $4$, en effet $\frac {\ln 4}4= \frac {\ln 2}2$.
Nous avons pour l'instant les solutions correspondant à $2^2=2^2$ et $2^4=4^2$.

Si $y \geq 3$ alors $y> e$ donc il y a une solution $x$ 
de l'\'equation $f(x)=f(y)$ dans $]e,+\infty[$ qui est $x=y$,
et une solution $x$ dans l'intervalle $]1,e[$.
Mais comme $x$ est un entier alors $x=2$ (c'est le seul entier appartenant à $]1,e[$)
c'est un cas que nous avons d\'ej\`a \'etudi\'e conduisant à $4^2 = 2^4$.

Conclusion : les couples d'entiers qui v\'erifient l'\'equation $x^y=y^x$
sont les couples $(x,y=x)$ et les couples $(2,4)$ et $(4,2)$.
\fincorrection
\nocorrection
\correction{005086}
\ 
\begin{center}
\shadowbox{
$\begin{array}{ccc}
\ch(a+b)=\ch a\ch b+\sh a\sh b&\mbox{et}&\ch(a-b)=\ch a\ch b-\sh a\sh b,\\
\sh(a+b)=\sh a\ch b+\ch a\sh b&\mbox{et}&\sh(a-b)=\sh a\ch b-\sh b\ch a\\
\tanh(a+b)=\frac{\tanh a+\tanh
b}{1+\tanh a\tanh b}&\mbox{et}&\tanh(a-b)=\frac{\tanh a-\tanh b}{1-\tanh a\tanh b}.
\end{array}$
}
\end{center}
Deux
démonstrations~:

$$\ch a\ch b+\sh a\sh b=\frac{1}{4}((e^a+e^{-a})(e^b+e^{-b})+(e^a-e^{-a})(e^b-e^{-b})) =
\frac{1}{2}(e^{a+b}+e^{-a-b})=\ch(a+b).$$

$$\tanh(a+b)=\frac{\sh(a+b)}{\ch(a+b)}=\frac{\sh a\ch b+\sh b\ch
a}{\ch a\ch b+\sh a\sh b}=\frac{\tanh a+\tanh b}{1+\tanh a\tanh b}$$
après division
du numérateur et du dénominateur par le nombre non nul $\ch a\ch b$.
En appliquant à $a=b=x$, on
obtient~:~

\begin{center}
\shadowbox{
$\forall x\in\Rr,\;\ch(2x)=\ch^2x+sh^2x=2ch^2x-1=2sh^2x+1,\;\sh(2x)=2\sh x\ch x\;\mbox{et}\;
\tanh(2x)=\frac{2\tanh
x}{1+\tanh^2x}.$
}
\end{center}
En additionnant entre elles les formules d'addition, on obtient les formules de linéarisation~:

$$\ch a\ch
b=\frac{1}{2}(\ch(a+b)+\ch(a-b)),\;\sh a\sh
b=\frac{1}{2}(\ch(a+b)-\ch(a-b))\;\mbox{et}\;\sh a\ch b=\frac{1}{2}(\sh(a+b)+\sh(a-b)),$$
et en particulier

$$\ch^2x=\frac{\ch(2x)+1}{2}\;\mbox{et}\;\sh^2x=\frac{\ch(2x)-1}{2}.$$
\fincorrection
\correction{005091}
\textbullet~Pour tout réel $x$, $\ch x>0$. Donc $f$ est définie, continue et dérivable sur $\Rr$. Pour tout réel $x$,

\begin{center}
$f'(x)=\sh x\frac{1}{\ch x}-1=\tanh x-1<0.$
\end{center}
$f$ est donc strictement décroissante sur $\Rr$.
\textbullet~Etude en $-\infty$. $\lim_{x\rightarrow -\infty}\ch x=+\infty$ et donc $\lim_{x\rightarrow -\infty}f(x)=+\infty$.
Cherchons une éventuelle droite asymptote.

\begin{center}
$f(x)=\ln\frac{e^x+e^{-x}}{2}-x=\ln(e^x+e^{-x})-\ln2-x=\ln(e^{-x})-x-\ln2+\ln(1+e^{2x})=-2x-\ln2+\ln(1+e^{2x}).$
\end{center}
Donc, $f(x)-(-2x-\ln2)=\ln(1+e^{2x})$. Or, d'une part $\lim_{x\rightarrow -\infty}\ln(1+e^{2x})= 0$ et donc la droite
$(D)$ d'équation $y=-2x-\ln2$ est asymptote à la courbe représentative de $f$ en $-\infty$ et d'autre part, pour
tout réel $x$, $\ln(1+e^{2x})>0$ et la courbe représentative de f est strictement au dessus de $(D)$ sur $\Rr$.
\textbullet~Etude en $+\infty$.
\begin{center}
$f(x)=\ln\frac{e^x+e^{-x}}{2}-x=\ln(e^x+e^{-x})-\ln2-x=\ln(e^x)-x-\ln2+\ln(1+e^{-2x})=-\ln2+\ln(1+e^{-2x})$
\end{center}
et $f$ tend vers $-\ln2$ quand $x$ tend vers $+\infty$.
\textbullet~Graphe.

$$\includegraphics{../images/img005091-1}$$

\fincorrection
\correction{005093}
Soit $x$ un réel.

$$S=\sum_{k=1}^{100}\sh(2+kx)=\frac{1}{2}\left(e^2\sum_{k=1}^{100}e^{kx}-e^{-2}\sum_{k=1}^{100}e^{-kx}\right).$$
Si $x=0$
alors directement $S=100\sh2\neq0$. Si $x\neq0$ alors $e^x\neq1$ et $e^{-x}\neq1$. Dans ce
cas,

$$S=\frac{1}{2}\left(e^2e^x\frac{1-e^{100x}}{1-e^x}-e^{-2}e^{-x}\frac{1-e^{-100x}}{1-e^{-x}}\right)=\frac{1}{2}\left(e^2e^x
\frac{1-e^{100x}}{1-e^x}+e^{-2}\frac{1-e^{-100x}}{1-e^{x}}\right).$$
après multiplication du numérateur et du dénominateur
de la deuxième fraction par $e^x$. Pour $x\neq0$, on a donc~:
\begin{align*}
S=0&\Leftrightarrow e^{x+2}(1-e^{100x})+ e^{-2}(1-e^{-100x})=0\Leftrightarrow
e^{x+2}(1-e^{100x})+e^{-2-100x}(e^{100x}-1)=0\\
 &\Leftrightarrow(1-e^{100x})(e^{x+2}-e^{-100x-2})=0\Leftrightarrow e^{x+2}=e^{-100x-2}\;(\mbox{car}\;
x\neq0)\\
 &\Leftrightarrow x+2=-100x-2\Leftrightarrow x=-\frac{4}{101}.
\end{align*}

\begin{center}
\shadowbox{
$\mathcal{S}=\left\{-\frac{4}{101}\right\}.$
}
\end{center}
\fincorrection
\correction{005094}
On a vu au \ref{exo:suprou3} que pour tout réel $x$, $\tanh(2x)=\frac{2\tanh x}{1+\tanh^2x}$ ce qui s'écrit pour
$x$ non nul~:~$\frac{1+\tanh^2x}{\tanh x}=\frac{2}{\tanh(2x)}$ ou encore $\tanh x+\frac{1}{\tanh x}=\frac{2}{\tanh(2x)}$
ou finalement

\begin{center}
\shadowbox{
$\forall x\in\Rr^*,\;\tanh x=\frac{2}{\tanh(2x)}-\frac{1}{\tanh x}.$
}
\end{center}
Soient $n$ un entier naturel non
nul et $x$ un réel non nul. D'après ce qui précède,

$$u_n
=\sum_{k=0}^{n}2^k\tanh(2^kx)=\sum_{k=0}^{n}\left(\frac{2^{k+1}}{\tanh(2^{k+1}x)}-\frac{2^{k}}{\tanh(2^{k}x)}\right)
=\sum_{k=1}^{n+1}\frac{2^{k}}{\tanh(2^{k}x)}-\sum_{k=0}^{n}\frac{2^{k}}{\tanh(2^{k}x)}=\frac{2^{n+1}}{\tanh(2^{n+1}x)}
-\frac{1}{\tanh x}.$$
Ensuite, pour $x>0$, $\tanh(2^{n+1}x)$ tend vers $1$ quand $n$ tend vers l'infini. Donc $u_n$ tend vers $+\infty$ quand
$n$ tend vers $+\infty$ si $x>0$ et vers $-\infty$ quand $n$ tend vers $+\infty$ si $x<0$.
\fincorrection
\correction{006975}
Par définition des fonctions $\ch$ et $\sh$, on a 
\begin{eqnarray*}
2\ch^2(x)-\sh(2x)&=&2\left(\frac{e^x+e^{-x}}{2}\right)^2-\frac{e^{2x}-e^{-2x}}{2}\\
 &=&\frac{e^{2x}+2+e^{-2x}}{2}+\frac{e^{-2x}-e^{2x}}{2}\\
 &=&1+e^{-2x}
\end{eqnarray*}

Et en utilisant les deux relations $\ln(ab)=\ln a + \ln b$ et $\ln(e^x) = x$ on calcule :
\begin{eqnarray*}
x-\ln(\ch x)-\ln 2
  &=& x-\ln\left(\frac{e^x+e^{-x}}{2}\right)-\ln 2\\
  &=& x-\ln(e^x+e^{-x}) + \ln 2 - \ln 2\\
  &=& x-\ln\big(e^x(1+e^{-2x})\big)\\
  &=& x-\ln(e^x)-\ln(1+e^{-2x})\\
  &=& x-x-\ln(1+e^{-2x})\\
  &=& -\ln(1+e^{-2x})\\
\end{eqnarray*}
d'où 
$$\frac{2\ch^2(x)-\sh(2x)}{x-\ln(\ch x)-\ln 2}=-\frac{1+e^{-2x}}{\ln(1+e^{-2x})}$$
C'est une expression de la forme $-\frac{u}{\ln u}$ avec $u=1+e^{-2x}$ :
\begin{itemize}
  \item si $x\to +\infty$, alors $u\to 1^+$, $\frac{1}{\ln u} \to +\infty$ 
  donc $-\frac{u}{\ln u}\to -\infty$ ;
  \item si $x\to -\infty$, alors $u\to +\infty$ donc d'après les 
  relations de croissances comparées, $-\frac{u}{\ln u}\to -\infty$.
\end{itemize}
\fincorrection
\correction{006976}
Puisque $\ch x+\sh x=e^x$ et $\ch x-\sh x=e^{-x}$, 
les expressions $C_n+S_n=\sum_{k=1}^ne^{kx}$ et $C_n-S_n=\sum_{k=1}^ne^{-kx}$ 
sont des sommes de termes de suites géométriques, de raison respectivement $e^x$ et $e^{-x}$.

Si $x=0$, on a directement $C_n=\sum_{k=1}^n1=n$ et $S_n=\sum_{k=1}^n0=0$.

Supposons $x\not=0$, alors $e^{x}\not=1$. 
Donc
\begin{eqnarray*}
C_n+S_n
  &=& \sum_{k=1}^ne^{kx} = \frac{e^x-e^{(n+1)x}}{1-e^x}\\
  &=& e^x\,\frac{1-e^{nx}}{1-e^x}\\
  &=& e^x\ \frac{e^{\frac{nx}{2}}(e^{-\frac{nx}{2}}-e^{\frac{nx}{2}})}{e^{\frac x2}(e^{-\frac x2}-e^{\frac x2})}\\
  &=& e^{\frac{(n+1)x}{2}}\ \frac{e^{\frac{nx}{2}}-e^{-\frac{nx}{2}}}{e^{\frac x2}-e^{-\frac x2}}\\
  &=& e^{\frac{(n+1)x}{2}}\ \frac{\sh\frac{nx}{2}}{\sh\frac{x}{2}}
\end{eqnarray*}

De même $C_n-S_n = \sum_{k=1}^ne^{-kx}$ ; c'est donc la même formule que ci-dessus en remplaçant $x$ par $-x$.
Ainsi :
$$C_n-S_n = e^{-\frac{(n+1)x}{2}}\ \frac{\sh\frac{nx}{2}}{\sh\frac{x}{2}}$$

En utilisant $C_n=\frac{(C_n+S_n)+(C_n-S_n)}{2}$ et $S_n=\frac{(C_n+S_n)-(C_n-S_n)}{2}$, on récupère donc
$$C_n=\frac{e^{\frac{(n+1)x}{2}}+e^{-\frac{(n+1)x}{2}}}{2}\,\frac{\sh\frac{nx}{2}}{\sh\frac{x}{2}}=\ch\left(\tfrac{(n+1)x}{2}\right)\,\frac{\sh\frac{nx}{2}}{\sh\frac{x}{2}}$$
$$S_n=\frac{e^{\frac{(n+1)x}{2}}-e^{-(n+1)\frac x2}}{2}\,\frac{\sh\frac{nx}{2}}{\sh\frac{x}{2}}=\sh\left(\tfrac{(n+1)x}{2}\right)\,\frac{\sh\frac{nx}{2}}{\sh\frac{x}{2}}$$
\fincorrection
\correction{006977}
\begin{eqnarray*}
(S)\ \left\{\begin{array}{l}
\ch(x)+\ch(y)=2a\\
\sh(x)+\sh(y)=2b
\end{array}\right.
&\Longleftrightarrow&
\left\{\begin{array}{l}
e^x+e^{-x}+e^y+e^{-y}=4a\\
e^x-e^{-x}+e^y-e^{-y}=4b
\end{array}\right.\\
 &\Longleftrightarrow&
\left\{\begin{array}{l}
e^x+e^y=2a+2b\\
e^x-e^{-x}+e^y-e^{-y}=4b
\end{array}\right.\\
 &\Longleftrightarrow&
\left\{\begin{array}{l}
e^x+e^y=2a+2b\\
-e^{-x}-e^{-y}=2b-2a
\end{array}\right.\\
 &\Longleftrightarrow&
\left\{\begin{array}{l}
e^x+e^y=2(a+b)\\
\frac{1}{e^{x}}+\frac{1}{e^{y}}=2(a-b)
\end{array}\right.
\end{eqnarray*}
ce qui donne, en posant  $X=e^x$ et $Y=e^y$:
\begin{eqnarray*}
(S)&\Longleftrightarrow&
\left\{\begin{array}{l}
X+Y=2(a+b)\\
\frac{1}{X}+\frac{1}{Y}=2(a-b)
\end{array}\right.\\
&\Longleftrightarrow&
\left\{\begin{array}{l}
X+Y=2(a+b)\\
\frac{X+Y}{XY}=2(a-b)
\end{array}\right.\\
&\Longleftrightarrow&
\left\{\begin{array}{l}
X+Y=2(a+b)\\
\frac{2(a+b)}{XY}=2(a-b)
\end{array}\right.
\end{eqnarray*}
Or $a\not=b$ puisque par hypothèse, $a^2-b^2=1$. Ainsi,
\begin{eqnarray*}
(S)&\Longleftrightarrow&
\left\{\begin{array}{l}
X+Y=2(a+b)\\
XY=\frac{a+b}{a-b}
\end{array}\right.\\
&\Longleftrightarrow& X\ \text{et}\ Y\ \text{sont les solutions de}\ z^2-2(a+b)z+\frac{a+b}{a-b}=0
\end{eqnarray*}

\medskip

\emph{Remarque :} On rappelle
que si $X, Y$ vérifient le système 
$\left\{\begin{array}{l}
X+Y=S\\
XY=P
\end{array}\right.$, alors $X$ et $Y$ sont les solutions de l'équation $z^2-Sz+P=0$.
\medskip

Or le discriminant du trinôme $z^2-2(a+b)z+\frac{a+b}{a-b}=0$ vaut 
$$\Delta=4(a+b)^2-4\frac{a+b}{a-b}=4(a+b)\left(a+b-\frac{1}{a-b}\right)=\frac{4(a+b)(a^2-b^2-1)}{a-b}=0$$ 
Il y a donc une racine double qui vaut $\frac{2(a+b)}{2}$, ainsi $X=Y=a+b$ et donc :
$$(S)\Longleftrightarrow e^x=e^y=a+b$$
On vérifie que $a+b\ge 0$ (car $a\ge 0$ et $b\ge 0$) et $a+b\not=0$ (car $a^2-b^2=1$).
Conclusion : le système $(S)$ admet une unique solution, donnée par $\big(x=\ln(a+b),y=\ln(a+b)\big)$. 
\fincorrection
\correction{006978}\
\begin{enumerate}
  \item 
  \begin{enumerate}
    \item On sait que $\ch^2 u=1+\sh^2u$. Comme de plus la fonction $\ch$ est à valeurs positives, 
$\ch u=\sqrt{1+\sh^2u}$ et donc $\ch(\Argsh x)=\sqrt{1+\sh^2(\Argsh x)} = \sqrt{1+x^2}$. 

    \item Alors
$$\tanh(\Argsh x)=\frac{\sh(\Argsh x)}{\ch(\Argsh x)}=\frac{x}{\sqrt{1+x^2}}.$$

    \item Et $\sh(2\Argsh x)=2\ch(\Argsh x)\sh(\Argsh x)=2x\sqrt{1+x^2}$.
  \end{enumerate}
  
\item On pourrait, comme pour la question précédente, appliquer les formules 
trigonométriques hyperboliques. Pour changer, on va plutôt utiliser les expressions explicites 
des fonctions hyperboliques réciproques. Supposons $x\ge 1$, pour que $\Argch x$ 
soit bien défini, alors on a la formule (à connaître) :
$$\Argch x=\ln\big(x+\sqrt{x^2-1}\big).$$ 

Ainsi :
\begin{eqnarray*}
\sh(\Argch x)&=&\frac{e^{\Argch x}-e^{-\Argch x}}{2}\\
 &=& \frac{x+\sqrt{x^2-1}-\frac{1}{x+\sqrt{x^2-1}}}{2}\\
 &=& \frac{x+\sqrt{x^2-1}}{2}-\frac{x-\sqrt{x^2-1}}{2\big(x+\sqrt{x^2-1}\big)\big(x-\sqrt{x^2-1}\big)}\\
 &=& \frac{x+\sqrt{x^2-1}}{2}-\frac{x-\sqrt{x^2-1}}{2\big(x^2-(x^2-1)\big)}\\
 &=& \sqrt{x^2-1}
\end{eqnarray*}

\medskip

Donc $\displaystyle\tanh(\Argch x)=\frac{\sh(\Argch x)}{\ch(\Argch x)}=\frac{\sqrt{x^2-1}}{x}$. 

\medskip

Enfin, si $u=\Argch x$:
$\ch(3u)=\ch(2u+u)=\ch(2u)\ch u+\sh(2u)\sh u$, où 
$$\left\{\begin{array}{l}
\ch(2u)=\ch^2u+\sh^2u=x^2+(x^2-1)=2x^2-1\\
\sh(2u)=2\sh u\ch u=2x\sqrt{x^2-1}
\end{array}\right.$$ 
Donc 
$\ch(3\Argch x)=(2x^2-1)x+2x\sqrt{x^2-1}\sqrt{x^2-1}=x(4x^2-3)$. 
\end{enumerate}
\fincorrection
\correction{006979}
La fonction $\Argch$ est définie sur $[1,+\infty[$. Or
\begin{eqnarray*}
\frac{1}{2}\left(x+\frac{1}{x}\right)\ge 1
 &\Longleftrightarrow&\frac{x^2+1}{x}\ge 2\\
 &\Longleftrightarrow&\frac{x^2+1-2x}{x}\ge 0\\
 &\Longleftrightarrow&\frac{(x-1)^2}{x}\ge 0\\
 &\Longleftrightarrow& x > 0 
\end{eqnarray*}
donc $f$ est définie sur $]0,+\infty[$. 

\medskip

Soit $x>0$, alors $y=\frac{1}{2}\left(x+\frac{1}{x}\right)\ge 1$
et on sait que $\Argch y=\ln(y+\sqrt{y^2-1})$. Ainsi
$\sqrt{y^2-1}=\sqrt{\frac{1}{4}\left(x+\frac{1}{x}\right)^2-1}
=\sqrt{\frac{(x^2+1)^2}{4x^2}-1}=\sqrt{\frac{(x^2-1)^2}{4x^2}}
=\left|\frac{x^2-1}{2x}\right|$,
on obtient
$$f(x)=\Argch y=\ln(y+\sqrt{y^2-1})=\ln\left(\frac{x^2+1}{2x}+\left|\frac{x^2-1}{2x}\right|\right)$$


On a supposé $x>0$, il suffit donc de distinguer les cas $x\ge 1$ et $0<x\le 1$.
\begin{itemize}
  \item Si $x\ge 1$, $\displaystyle f(x)=\ln\left(\frac{x^2+1}{2x}+\frac{x^2-1}{2x}\right)=\ln x$.
  \item Si $0<x\le 1$, $\displaystyle f(x)=\ln\left(\frac{x^2+1}{2x}+\frac{1-x^2}{2x}\right)=\ln \frac{1}{x}=-\ln x$.
\end{itemize}

Puisque $\ln x$ est positif si $x\ge 1$ et négatif si $x\le 1$, 
on obtient dans les deux cas $f(x)=|\ln x|$.

\begin{center}
\begin{tikzpicture}[scale=1.5]
      \draw[->,>=latex, gray] (-1,0)--(5,0) node[below,black] {$x$};
      \draw[->,>=latex, gray] (0,-1)--(0,2.7) node[right,black] {$y$};  
 
     \draw[ultra thick, color=red,domain=0.1:5,samples=100,smooth] plot (\x,{abs(ln(\x))}) node[above left] {$f(x)$}; 
     
     \fill (0,0) circle (1.5pt);     
     \fill (1,0) circle (1.5pt);
     \fill (0,1) circle (1.5pt);     
     \node at (0,1)[above left] {$1$};
     \node at (1,0)[below] {$1$};
     \node at (0,0)[below right] {$0$};
\end{tikzpicture}  
\end{center}

\fincorrection
\correction{006980}
Soit $f(x)=\Argsh x+\Argch x$. La fonction $f$ est bien définie, continue, et strictement croissante, 
sur $[1,+\infty[$ (comme somme de deux fonctions continues strictement croissantes). 

\begin{center}
\begin{tikzpicture}[scale=2]
      \draw[->,>=latex, gray] (-0.5,0)--(2.5,0) node[below,black] {$x$};
      \draw[->,>=latex, gray] (0,-0.5)--(0,3) node[right,black] {$y$};  
 
      \draw[ultra thick, color=red,domain=1.0:2,samples=100,smooth] plot (\x,{ln(\x+sqrt(\x^2+1)) + ln(\x+sqrt(\x^2-1))}) node[above left] {$f(x)$}; 

   \draw (0,1.1)--(1.03,1.1)--(1.03,0);

   
   \draw (0,0.9)--(0.98,0.9)--(0.98,0);
   

     \fill (0,0) circle (1pt);
     \fill (0,0.9) circle (1pt); 
     \fill (0.96,0) circle (1pt);
     \fill (1.04,0) circle (1pt);
     
     \fill (0,1.1) circle (1pt);     
     \node at (0,1)[above left] {$1$};
     \node at (0,0.9)[below left] {$f(1)$};
     \node at (1,0)[below left] {$1$};
     \node at (1.03,0)[below right] {$a$};
     \node at (0,0)[below right] {$0$};
\end{tikzpicture}  
\end{center}

De plus, $f(x)\xrightarrow[x\to +\infty]{}+\infty$, donc $f$ atteint exactement une fois 
toute valeur de l'intervalle $[f(1),+\infty[$. Comme (par la formule logarithmique)
$f(1)=\Argsh 1=\ln(1+\sqrt{2})<\ln(e)=1$, 
on a $1\in[f(1),+\infty[$. Par le théorème des valeurs intermédiaires
l'équation $f(x)=1$ admet une unique solution, que l'on notera $a$. 

\medskip

Déterminons la solution :
\begin{eqnarray*}
\sh 1&=&\sh(\Argsh a+\Argch a)\\ 
 &=&\sh(\Argsh a)\ch(\Argch a)+\sh(\Argch a)\ch(\Argsh a)\\
 &=&a^2+\sqrt{a^2-1}\sqrt{a^2+1}=a^2+\sqrt{a^4-1}
\end{eqnarray*}
donc $\sqrt{a^4-1}=\sh 1-a^2$. En élevant au carré et en simplifiant, on obtient
$a^2=\frac{1+\sh^21}{2\sh 1}=\frac{\ch^2 1}{2\sh 1}$. Comme on cherche $a$ positif 
(et que $\ch 1>0$), on en déduit $a=\frac{\ch 1}{\sqrt{2\sh 1}}$. 
Cette valeur est la seule solution possible de l'équation $f(x)=1$, il faudrait normalement 
vérifier qu'elle convient bien, puisqu'on a seulement raisonné par implication (et pas par équivalence).
Or on sait déjà que l'équation admet une unique solution: c'est donc nécessairement 
$$a=\frac{\ch 1}{\sqrt{2\sh 1}}= \tfrac{1}{2}\frac{e+\frac1e}{\sqrt{e-\frac1e}} = 1,0065\ldots.$$
\fincorrection
\nocorrection
\nocorrection
\nocorrection
\nocorrection
\nocorrection
\nocorrection
\nocorrection
\nocorrection
\nocorrection
\nocorrection
\nocorrection
\nocorrection
\nocorrection
\nocorrection
\nocorrection
\nocorrection
\nocorrection
\nocorrection
\nocorrection
\nocorrection
\nocorrection
\nocorrection
\nocorrection
\nocorrection
\nocorrection
\nocorrection
\correction{002081}
  \begin{enumerate}
  \item On trouve $\int_0^4 f(t) \, dt = +7$. Il faut tout d'abord tracer le graphe de cette fonction. 
Ensuite la valeur d'une intégrale ne dépend pas de la valeur de la fonction en un point, 
c'est-à-dire ici les valeurs en $x=0$, $x=1$, $x=2$ n'ont aucune influence sur l'intégrale. 
Ensuite on revient à la définition de $\int_0^4f(t) \, dt$ : pour la subdivision de $[0,4]$
 définie par $\{x_0=0,x_1=1,x_2=2,x_3=3, x_4=4\}$, on trouve la valeur de l'intégrale (ici le sup et l'inf 
sont atteints et égaux pour cette subdivision et toute subdivision plus fine).
Une autre façon de faire est considérer que $f$ est une fonction en escalier (en <<oubliant>> les accidents en
$x=0$, $x=1$, $x=2$) dont on sait calculer l'intégrale.
  \item C'est la même chose pour $\int_0^x  f(t) \, dt$, mais au lieu d'aller jusqu'à $4$ on s'arrête à $x$, on trouve
\begin{equation*}
  F(x)=
  \begin{cases}
   x &\text{ si $0\leqslant x \leqslant 1$}\\
   3-2x &\text{ si $1<x\leqslant 2$}\\
   4x-9 &\text{ si $2 < x \leqslant 4$.}\\
  \end{cases}
\end{equation*}

  \item Les seuls points à discuter pour la continuité sont les points $x=1$ et $x=2$,
mais les limites à droite et à gauche de $F$ sont égales en ces points donc $F$ est continue. 
Par contre $F$ n'est pas dérivable en $x=1$ (les dérivées à droite et à gauche 
sont distinctes), $F$ n'est pas non plus dérivable en $x=2$.

  \end{enumerate}
\fincorrection
\correction{002082}
Les fonctions sont continues donc intégrables !

  \begin{enumerate}
  \item En utilisant les sommes de Riemann, on sait que $\int_0^1f(x)d x$
est la limite (quand $n \to +\infty$) de $\frac 1n\sum_{k=0}^{n-1}  f(\frac kn)$.
Notons $S_n = \frac 1n \sum_{k=0}^{n-1}  f(\frac kn)$.
Alors $S_n = \frac 1n\sum_{k=0}^{n-1}\frac k{n}=\frac 1{n^2} \sum_{k=0}^{n-1} k = \frac 1{n^2} \frac{n(n-1)}{2}$. On a utilisé que la somme des entiers de $0$ à $n-1$ vaut $\frac{n(n-1)}{2}$. Donc $S_n$ tend vers $\frac 12$. Donc $\int_0^1f(x)d x = \frac 12$.
  \item Même travail : $\int_1^2 g(x)
d x$ est la limite de $S'_n = \frac {2-1} n\sum_{k=0}^{n-1}  g(1+ k\frac {2-1}n) = \frac 1 n \sum_{k=0}^{n-1}  (1+ \frac kn)^2 =  \frac 1 n \sum_{k=0}^{n-1}  (1+ 2\frac k n + \frac {k^2}{n^2})$.
En séparant la somme en trois nous obtenons : $S'_n = \frac 1n (n+\frac 2 n\sum_{k=0}^{n-1} k + \frac 1 {n^2} \sum_{k=0}^{n-1} k^2)= 1+\frac 2 {n^2} \frac{n(n-1)}{2}+ \frac 1 {n^3} \frac{(n-1)n(2n-1)}{6}$. Donc à la limite on trouve
$S'_n \to 1+1 + \frac 13 = \frac 73$. Donc $\int_1^2 g(x)
d x = 7/3$. Remarque : on a utilisé que la somme des carrés des entiers de $0$ à $n-1$ est $\frac{(n-1)n(2n-1)}{6}$.

  \item Même chose pour $\int_0^x h(t) d t$ qui est la limite de 
$S''_n = \frac {x} n\sum_{k=0}^{n-1}  g(\frac {kx}n) = \frac {x} n\sum_{k=0}^{n-1}  e^{\frac {kx}n} = \frac {x} n\sum_{k=0}^{n-1}  (e^{\frac {x}n})^k$. Cette dernière somme est la somme d'une suite géométrique (si $x\neq 0$), donc 
$S''_n = \frac {x} n \frac{1-(e^\frac xn)^n}{1-e^\frac xn}=  \frac {x} n \frac{1-e^x}{1-e^\frac xn}=(1-e^x) \frac{\frac xn}{1-e^\frac xn}$ 
qui tend vers $e^x-1$. 
Pour obtenir cette dernière limite on remarque qu'en posant 
$u=\frac xn$ on a $\frac{\frac xn}{1-e^\frac xn}= -1/ \frac{e^u-1}{u}$ qui tend vers $-1$
lorsque $u\to 0$ (ce qui est équivalent à $n\to +\infty$).

  \end{enumerate}
\fincorrection
\correction{002083}
  \begin{enumerate}
  \item On calcul d'abord  $\int_0^{\frac \pi 2}  e^{it}\, dt$. Par le th\'eor\`eme de  Riemann-Darboux
c'est la limite de 
$$S_n = \sum_{k=0}^{n-1} (x_{k+1}-x_k) \cdot f(x_k).$$
Pour $x_k = \frac{k\pi }{2n}$ (on obtient en fait un somme de Riemann) :
$$S_n = \frac{\pi}{2n}\sum_{k=0}^{n-1}e^{\frac{ik\pi}{2n}}= \frac{\pi}{2n}\sum_{k=0}^{n-1}(e^{\frac{i\pi}{2n}})^k.$$
Ce qui est une somme g\'eom\'etrique de somme $S_n = (1-i) \frac{\frac {\pi}{2n}}{1-e^{i\frac {\pi}{2n}}}$. La limite de ce taux d'accroissement est $1+i$ (en posant $u=\frac {\pi}{2n}$ et en remarquant que $\frac{e^{iu}-1}{u} \to i$ quand $u\to 0$).
Donc $\int_0^{\frac \pi 2}  e^{it}\, dt=1+i$. Mais $e^{it} = \cos t+ i \sin t$ donc 
$\int_0^{\frac \pi 2} \cos t \, dt + \int_0^{\frac \pi 2}\sin t \, dt = 1+i$.
Par identification des parties r\'eelles et imaginaires on trouve : 
$\int_0^{\frac \pi 2} \cos t \, dt = 1$ et $\int_0^{\frac \pi 2}\sin t \, dt = 1$.

  \item On veut $x_k = aq^k$ ce qui donne bien $x_0=a$, mais il faut aussi $x_n = b$ donc
$aq^n=b$, donc $q^n = \frac ba$ soit $q = (\frac ba)^{\frac 1 n}$.
Nous cherchons la limite de $S'_n =  \sum_{k=0}^{n-1} (x_{k+1}-x_k) \cdot g(x_k).$
Il est n'est pas trop dur de montrer que $S'_n = n(q-1)$.
Pour trouver la limite quand $n\to +\infty$ c'est plus d\'elicat car $q$ d\'epend de $n$ :
$S'_n= n(q-1)= n((\frac ba)^{\frac 1 n}-1) = n(e^{\frac1n \ln\frac ba}-1)$.
En posant $u = \frac 1n$ et en remarquant que l'on obtient un taux d'accroissement on calcule :
$S'_n = \frac 1 u(e^{ u \ln\frac ba}-1) \to  \ln \frac ba = \ln b - \ln a$.
Donc $\int_a^b \frac {dt}{t} = \ln b - \ln a$.

  \item \`A l'aide des sommes g\'eom\'etrique est des taux d'accroissement on trouve 
$$\int_a^b \alpha^t \, dt = \frac{e^{\alpha b}- e^{\alpha a}}{\alpha}.$$

  \end{enumerate}
\fincorrection
\correction{002084}
  \begin{enumerate}
  \item Oui.
  \item Non.
  \item Non.
  \item Non.
  \end{enumerate}
\fincorrection
\correction{002085}
 \begin{enumerate}
  \item \'Ecrivons la continuité de $f$ en $x_0$ avec $\epsilon = \frac {f(x_0)}{2} > 0$ :
il existe $\delta >0$ tel que pour tout $x\in [x_0-\delta, x_0+\delta]$ on ait $|f(x)-f(x_0)| \leqslant \epsilon$.
Avec notre choix de $\epsilon$ cela donne pour $x\in [x_0-\delta, x_0+\delta]$ que $f(x) \geqslant \frac {f(x_0)}{2}$.
Pour évaluer $\int_a^b f(x) \, dx$ nous la coupons en trois morceaux par linéarité de l'intégrale :
$$\int_a^b f(x) \, dx = \int_a^{x_0-\delta} f(x) dx +  \int_{x_0-\delta}^{x_0+\delta} f(x) dx +  \int_{x_0+\delta}^b f(x) dx.$$
Comme $f$ est positive alors par positivité de l'intégrale $\int_a^{x_0-\delta} f(x) dx \geqslant 0$ et
$\int_{x_0+\delta}^b f(x) dx \geqslant 0$. Pour le terme du milieu on a $f(x) \geqslant \frac {f(x_0)}{2}$ donc
$\int_{x_0-\delta}^{x_0+\delta} f(x) dx \geqslant \int_{x_0-\delta}^{x_0+\delta} \frac {f(x_0)}{2} dx =
2\delta\frac {f(x_0)}{2}$ (pour la dernière équation on calcule juste l'intégrale d'une fonction constante !).
Le bilan de tout cela est que $\int_a^b f(x) \, dx \geqslant 2\delta\frac {f(x_0)}{2} >0$.

Donc pour une fonction continue et positive $f$, si elle est strictement positive en un point alors  $\int_a^b f(x) \, dx >0$.
Par contraposition pour une fonction continue et positive si $\int_a^b f(x) \, dx =0$ alors
$f$ est identiquement nulle.


  \item Soit $f$ est tout le temps positive, soit elle tout le temps négative, soit elle change 
(au moins un fois) de signe. Dans le premier cas $f$ est identiquement nulle par la première question, 
dans le second cas c'est pareil (en appliquant la première question à $-f$). Pour le troisième cas 
le théorème des valeurs intermédiaires affirme qu'il existe $c$ tel que $f(c)=0$.

  \item Posons $g(x) = f(x)-x$. Alors $\int_0^1 g(x) dx = \int_0^1 \big( f(x) - x \big)  dx= \int_0^1 f(x) dx - \frac 12 = 0$.
Donc par la question précédente, $g$ étant continue, il existe $d \in [0,1]$ tel que $g(d)=0$, ce qui est équivalent à $f(d)=d$.
  \end{enumerate}
\fincorrection
\correction{002086}
Notons $I = \int_a^b \frac{f(t)^n}{m^n}dt.$ Comme $f(t) \leqslant m$ pour tout $t\in [a,b]$ alors 
$I \leqslant 1$. Ceci implique que $\lim_{n\to+\infty} I^{\frac1n} \leqslant 1$.
Fixons $\alpha >0$ (aussi petit que l'on veut). Comme $f$ est continue et $m$ est sa borne sup\'erieure sur $[a,b]$ alors il existe
un intervalle $[x,y]$, ($x<y$), sur le quel $f(t) \geqslant m - \alpha$.
Comme $f$ est positive alors 
$$I \geqslant \int_x^y  \frac{f(t)^n}{m^n}dt \geqslant \int_x^y \frac{(m-\alpha)^n}{m^n} = (y-x)\left(\frac{m-\alpha}{m}\right)^n$$.

Donc $I^{\frac1n} \geqslant (y-x)^{\frac1n} \frac{m-\alpha}{m}$.
Quand $n\to+\infty$ on a $(y-x)^{\frac1n} \to 1$, donc \`a la limite
nous obtenons $\lim_{n\to+\infty} I^{\frac1n} \geqslant \frac{m-\alpha}{m}$.

Comme $\alpha$ est quelconque, nous pouvons le choisir aussi proche de $0$ de sorte que 
$\frac{m-\alpha}{m}$ est aussi proche de $1$ que d\'esir\'e.
Donc $\lim_{n\to+\infty} I^{\frac1n} \geqslant 1$.


En conclusion nous trouvons que $\lim_{n\to+\infty} I^{\frac1n} = 1$ ce qui \'etait l'\'egalit\'e recherch\'ee.

\fincorrection
\correction{002087}
Soit $\alpha>0$ fix\'e. Soit $0<x_0<1$ tel que pour tout $x \in [0,x_0]$, $f(x) \leqslant 1-\alpha$. 
Ce $x_0$ existe bien car $f$ est strictement croissante
et $f(0)=0$, $f(1)=1$.
S\'eparons l'int\'egrale en deux :
\begin{align*}
\int_{0}^{1}f^{n}(t) dt 
  &= \int_{0}^{x_0}f^{n}(t) dt+\int_{x_0}^{1}f^{n}(t) dt \\
  &\leqslant \int_{0}^{x_0} (1-\alpha)^n dt + \int_{x_0}^{1} 1^n dt \\
  &\leqslant x_0(1-\alpha)^n+(1-x_0) \\
  &\leqslant (1-\alpha)^n+(1-x_0) \quad \text{ car } x_0 \leqslant 1 \\
\end{align*}

Soit maintenant donn\'e un $\epsilon >0$, on choisit $\alpha>0$ tel que 
$1-x_0 \leqslant \frac \epsilon 2$ (en remarquant que si $\alpha \to 0$ alors $x_0(\alpha) \to 1$),
puis il existe $n$ assez grand tel que $(1-\alpha)^n \leqslant \frac \epsilon 2$.
Donc pour tout $\epsilon>0$ il existe $n$ assez grand tel que $\int_{0}^{1}f^{n}(t) dt \leqslant \epsilon$ .
Donc $\int_{0}^{1}f^{n}(t) dt \to 0$.
\fincorrection
\correction{002091}
  \begin{enumerate}
  \item Vrai.
  \item Vrai.
  \item Faux ! Attention aux valeurs négatives par exemple pour $f(x)=x$ alors $F$ est décroissante sur $]-\infty,0]$ et croissante sur $[0,+\infty[$.
  \item Faux. Attention aux valeurs négatives par exemple pour $f(x)=x^2$ alors $F$ est négative sur $]-\infty,0]$ et positive sur $[0,+\infty[$.
  \item Vrai.
  \item Faux. Faire le calcul avec la fonction $f(x) = 1+\sin(x)$ par exemple.
  \item Vrai.
  \end{enumerate}
\fincorrection
\correction{002092}
 \begin{enumerate}
  \item Commen\c{c}ons plus simplement avec la fonction 
$$H(x) =  \int_{a}^{v(x)}f(t)dt.$$
En fait $H$ est la composition de la fonction $x\mapsto v(x)$ avec la fonction
$G : x \mapsto \int_{a}^{x}f(t)dt$ : 
$$H = G \circ v.$$
La fonction $v$ est d\'erivable et la fonction $G$ aussi (c'est une primitive) donc
la compos\'ee $H=G\circ v$ est d\'erivable, de plus $H'(x) = v'(x)\cdot G'(v(x))$.
En pratique comme $G'(x) = f(x)$ cela donne $H'(x) = v'(x) f(v(x))$.

\emph{Remarque :} Il n'est pas n\'ecessaire de conna\^\i tre cette formule mais il est important de savoir refaire ce petit raisonnement.

On montrerait de m\^eme que la fonction $x\to\int_{u(x)}^{a}f(t)dt$ est d\'erivable de d\'eriv\'ee $-u'(x)f(u(x))$.
Revenons \`a notre fonction $F(x) = \int_{u(x)}^{v(x)}f(t)d t =  \int_{u(x)}^{a}f(t)dt+ \int_{a}^{v(x)}f(t)dt$, 
c'est la somme de deux fonctions d\'erivables donc est d\'erivable de d\'eriv\'ee :
$$F'(x) = v(x)f(v(x))-u'(x)f(u(x)).$$

  \item On applique ceci \`a $u(x) = x$ et $v(x) = 2x$ nous obtenons :
$$G'(x) = \frac{2}{1+(2x)^2+(2x)^4}-\frac{1}{1+x^2+x^4}.$$
  \end{enumerate}
\fincorrection
\correction{002093}
  \begin{enumerate}
  \item $F$ est d\'efinie sur $]0,1[\cup ]1,+\infty[$. $F$ est continue et d\'erivable sur $]0,1[$
et sur $]1,+\infty[$. Pour vois cela il suffit d'\'ecrire $F(x)= \int_x^a \frac{dt}{\ln t}+\int_a^{x^2} \frac{dt}{\ln t}$. La premi\`ere de ces fonctions est continue et d\'erivable (c'est une primitive), la seconde est la compos\'ee de $x\mapsto x^2$ avec $x\mapsto \int_a^{x} \frac{dt}{\ln t}$ et est donc aussi continue et d\'erivable.
On pourrait m\^eme calculer la d\'eriv\'ee.
  \item Notons $f(t) = \frac 1{\ln t}$ et $g(t) = \frac 1{t\ln t}$. On se place sur $]1,+\infty[$.
Bien \'evidemment $g(t) \leqslant f(t)$, mais nous avons aussi que pour $\epsilon >0$ fix\'e il existe $x>1$ tel que
pour tout $t\in [1,x^2]$ on ait $\frac 1t \leqslant 1+\epsilon$ donc sur $]1,x^2]$ nous avons $f(t) \leqslant (1+\epsilon)g(t)$.
Par int\'egration de l'in\'egalit\'e $g(t) \leqslant f(t) \leqslant (1+\epsilon)g(t)$ sur $[x,x^2]$ nous obtenons pour $x$ assez proche de $1$ :
$$H(x) \leqslant F(x)  \leqslant (1+\epsilon)H(x).$$

Il ne reste plus qu'a calculer $H(x)$. En fait $g(t) = \frac 1{t\ln t}$ est la d\'eriv\'ee de la fonction $h(t) = \ln (\ln t)$. Donc 
\begin{align*}
H(x) = \int_x^{x^2}\frac{dt}{t\ln t} 
 &= [ \ln (\ln t) ]_x^{x^2}  = \ln (\ln (x^2))-\ln (\ln x) \\
 &= \ln (2 \ln x)-\ln (\ln x) = \ln \frac {2 \ln x}{\ln x} \\
 &= \ln 2. \\
\end{align*}

Nous obtenons alors, pour $\epsilon >0$ fix\'e et $x>1$ assez proche de $1$, l'encadrement 
$$\ln 2 \leqslant F(x) \leqslant (1+\epsilon) \ln 2.$$

Donc la limite de $F(x)$ quand $x\to 1^+$ est $\ln 2$.
  \end{enumerate}
\fincorrection
\nocorrection
\nocorrection
\nocorrection
\nocorrection
\nocorrection
\nocorrection
\correction{004222}
\begin{enumerate}
  \item 
  \item 
  \item $\frac12f(0)$.
\end{enumerate}
 \fincorrection
\nocorrection
\nocorrection
\nocorrection
\correction{004226}
DL de $1-\cos u \Rightarrow \lim = \frac12\ln(b/a)$.
\fincorrection
\nocorrection
\correction{004232}
$a_{n+1} = \begin{cases} b_n &\text{ si } b_n < -1\cr
                   -(b_n-1)^2/4 &\text{ si } -1\le b_n\le1\cr
                   0 &\text{ si } b_n > 1,\cr\end{cases}$  

$b_{n+1} = \begin{cases} 0 &\text{ si } a_n < -1\cr
                   (a_n+1)^2/4 &\text{ si } -1\le a_n\le1\cr
                   a_n &\text{ si } a_n > 1.\cr\end{cases}$
\par
Donc $a_{n+1} = f(a_{n-1})$, $b_{n+1} = g(b_{n-1})$.
Point fixe : $a_n \to \sqrt8-3$, $b_n\to 3-\sqrt8$.
 \fincorrection
\nocorrection
\correction{004235}
$\frac {\pi^2}4$.
\fincorrection
\correction{004236}
$u = \pi-t  \Rightarrow  I = \frac \pi2 \int_{t=0}^\pi \ \frac 1{1+\sin t}\,d t
= \frac \pi2  \int_{t=-\pi/2}{\pi/2} \ \frac 1{1+\cos t}\,d t = \pi$.
\fincorrection
\nocorrection
\correction{005444}
$f$ est continue sur le segment $[a,b]$ et admet donc un maximum $M$ sur ce segment. Puisque $f$ est strictement positive sur $[a,b]$, ce maximum est strictement positif.

Pour $n\in\Nn^*$, posons $u_n=\left(\int_{a}^{b}(f(x))^n\;dx\right)^{1/n}$. Par croissance de l'intégrale, on a déjà 

$$u_n\leq\left(\int_{a}^{b}M^n\;dx\right)^{1/n}=M(b-a)^{1/n},$$
 
(car $\forall x\in[a,b],\;0\leq f(x)\leq M\Rightarrow \forall x\in[a,b],\;(f(x))^n\leq M^n$ par croissance de la fonction $t\mapsto t^n$ sur $[0,+\infty[$).

D'autre part, par continuité de $f$ en $x_0$ tel que $f(x_0)=M$, pour $\varepsilon\in]0,2M[$ donné, $\exists[\alpha,\beta]\subset[a,b]/\;\alpha<\beta\;\mbox{et}\;\forall x\in[\alpha,\beta],\;f(x)\geq M-\frac{\varepsilon}{2}$.

Pour $n$ élément de $\Nn^*$, on a alors 

$$u_n\geq\left(\int_{\alpha}^{\beta}(f(x))^n\;dx\right)^{1/n}\geq\left(\int_{\alpha}^{\beta}(M-\frac{\varepsilon}{2})^n\;dx\right)^{1/n}=(M-\frac{\varepsilon}{2})(\beta-\alpha)^{1/n}.$$

En résumé,

$$\forall\varepsilon\in]0,2M[,\;\exists(\alpha,\beta)\in[a,b]^2/\;\alpha<\beta\;\mbox{et}\;\forall n\in\Nn^*,\;
(M-\frac{\varepsilon}{2})(\beta-\alpha)^{1/n}\leq u_n\leq M(b-a)^{1/n}.$$

Mais, $\lim_{n\rightarrow +\infty}M(b-a)^{1/n}=M$ et $\lim_{n\rightarrow +\infty}(M-\frac{\varepsilon}{2})(\beta-\alpha)^{1/n}=(M-\frac{\varepsilon}{2})(\beta-\alpha)^{1/n}$.

Par suite, $\exists n_1\in\Nn^*/\;\forall n\geq n_1,\;M(b-a)^{1/n}<M+\varepsilon$ et $\exists n_2\in\Nn^*/\;\forall n\geq n_2,\;(M-\frac{\varepsilon}{2})(\beta-\alpha)^{1/n}>M-\varepsilon$.

Soit $n_0=\mbox{Max}\{n_1,n_2\}$. Pour $n\geq n_0$, on a $M-\varepsilon<u_n<M+\varepsilon$. On a montré que 

$$\forall\varepsilon>0,\;\exists n_0\in\Nn^*/\;\forall n\in\Nn,\;(n\geq n_0\Rightarrow|u_n-M|<\varepsilon),$$

et donc que $\lim_{n\rightarrow +\infty}u_n=M$.

Plus généralement, si $g$ continue sur $[a,b]$, $g$ admet un minimum $m_1$ et un maximum $M_1$ sur cet intervalle, tous deux strictement positifs puisque $g$ est strictement positive. Pour $n$ dans $\Nn^*$, on a
 
$$m_1^{1/n}\left(\int_{a}^{b}(f(x))^n\;dx\right)^{1/n}\leq\left(\int_{a}^{b}(f(x))^ng(x)\;dx\right)^{1/n}\leq M_1^{1/n}\left(\int_{a}^{b}(f(x))^n\;dx\right),$$

et comme d'après l'étude du cas $g=1$, on a $\lim_{n\rightarrow +\infty}m_1^{1/n}\left(\int_{a}^{b}(f(x))^n\;dx\right)^{1/n}=\lim_{n\rightarrow +\infty}M_1^{1/n}\left(\int_{a}^{b}(f(x))^n\;dx\right)^{1/n}=M$, le théorème de la limite par encadrements permet d'affirmer que $\lim_{n\rightarrow +\infty}\left(\int_{a}^{b}(f(x))^ng(x)\;dx\right)^{1/n}=M$. On a montré que 

$$\lim_{n\rightarrow +\infty}\left(\int_{a}^{b}(f(x))^ng(x)\;dx\right)^{1/n}=\mbox{Max}\{f(x),\;x\in[a,b]\}.$$

\fincorrection
\correction{005449}
\begin{enumerate}
\item  Soient $m$ un réel strictement positif et, pour $t\in\Rr$, $f_m(t)=e^{mt}$. $f_m$ est bien un élément de $E$ et de plus,

\begin{align*}\ensuremath
\varphi(f_m)&=\frac{1}{m^2}(e^{mb}-e^{ma})(e^{-ma}-e^{-mb})\\
 &=\frac{1}{m^2}e^{m(a+b)/2}(e^{m(b-a)/2}+e^{-m(b-a)/2})e^{-m(a+b)/2}(e^{m(b-a)/2}+e^{-m(b-a)/2})\\
 &=\frac{4\sh^2(m(b-a)/2)}{m^2}.
\end{align*} 

Cette expression tend vers $+\infty$ quand $m$ tend vers $+\infty$ et $\varphi(E)$ n'est pas majoré.

\item  Soit $f$ continue et strictement positive sur $[a,b]$. L'inégalité 
de \textsc{Cauchy}-\textsc{Schwarz} montre que~:

$$\varphi(f)=\int_{a}^{b}\left(\sqrt{f(t)}\right)^2\;dt\int_{a}^{b}\left(\frac{1}{\sqrt{f(t)}}\right)^2\;dt\geq\left(
\int_{a}^{b}\sqrt{f(t)}\frac{1}{\sqrt{f(t)}}\;dt\right)^2=(b-a)^2,$$

avec égalité si et seulement si la famille de fonctions $(\sqrt{f(t)},\frac{1}{\sqrt{f(t)}})$ est liée ou encore si et seulement si $\exists\lambda\in\Rr_+^*/\;\forall t\in[a,b],\;\sqrt{f(t)}=\lambda\frac{1}{\sqrt{f(t)}}$ ou enfin si et seulement si $\exists\lambda\in\Rr_+^*/\;\forall t\in[a,b],\;f(t)=\lambda$, c'est-à-dire que $f$ est une constante strictement positive.

Tout ceci montre que $\varphi(E)$ admet un minimum égal à $(b-a)^2$ et obtenu pour toute fonction $f$ qui est une  constante strictement positive.
\end{enumerate}
\fincorrection
\correction{005917}
\begin{enumerate}
\item Soit $\varepsilon>0$ donn\'e. Puisque $f$ est
Riemann-int\'egrable sur $[a,b]$, il existe une subdivision
$\sigma_{1}=\{a_{0}=a< a_{1}< \dots< a_{n}=b\}$ de $[a, b]$ telle
que
$\overline{S}_{f}^{\sigma_{1}}\leq\underline{S}_{f}^{\sigma_{1}} +
\frac{\varepsilon}{2}$. Puisque $g$ est Riemann-int\'egrable sur
$[a,b]$, il existe une subdivision $\sigma_{2}= \{b_{0}=a< b_{1}<
\dots< b_{p}=b\}$ de $[a, b]$ telle que
$\overline{S}_{g}^{\sigma_{2}}\leq\underline{S}_{g}^{\sigma_{2}} +
\frac{\varepsilon}{2}$. On note $\sigma_{1}\cup\sigma_2=\{c_{0}=a<
c_{1} < \dots < c_{q-1} < c_{q} = b\}$ la subdivision de $[a,b]$
obtenue en ordonnant l'ensemble $\{a_{0}, \dots, a_{n}, b_{0},
\dots, b_{n}\}$ par ordre croissant, puis en identifiant les
points qui apparaissent plusieurs fois (on obtient une subdivision
de $[a,b]$ en $q$ intervalles avec  $max\{n, p\}\leq q \leq n+p$).
Puisque $\sigma_{1}\cup\sigma_{2}$ est une subdivision \emph{plus
fine} que $\sigma_{1}$, on a~:
\begin{equation}\label{1}
\overline{S}_{f}^{\sigma_{1}\cup\sigma_{2}}\, \leq\,\,
\overline{S}_{f}^{\sigma_{1}}\quad\quad\text{et}\quad\quad
\underline{S}_{f}^{\sigma_{1}}\leq\,\,
\underline{S}_{f}^{\sigma_{1}\cup\sigma_{2}}.
\end{equation}
De m\^eme,
\begin{equation}\label{2}
\overline{S}_{g}^{\sigma_{1}\cup\sigma_{2}}\, \leq\,\,
\overline{S}_{g}^{\sigma_{2}}\quad\quad\text{et}\quad\quad
\underline{S}_{g}^{\sigma_{2}}\leq\,\,
\underline{S}_{g}^{\sigma_{1}\cup\sigma_{2}}.
\end{equation}
De plus, sur un intervalle $]c_{k-1}, c_{k}[$ donn\'e, on a~:
\begin{eqnarray*}
\sup\{f(x) + g(x), \,x \in \,]c_{k-1}, c_{k}[\} &\leq &
\sup\{f(x), \,x \in \,]c_{k-1}, c_{k}[\} \\& &+ \sup\{g(x), \,x
\in \,]c_{k-1}, c_{k}[\}.
\end{eqnarray*}
De m\^eme~:
\begin{eqnarray*}
\inf\{f(x) + g(x), \,x \in \,]c_{k-1}, c_{k}[\} & \geq &
\inf\{f(x),\,x \in \,]c_{k-1}, c_{k}[\}\\ & & + \inf\{g(x), \,x
\in \,]c_{k-1}, c_{k}[\}.
\end{eqnarray*}
On en d\'eduit que~:
\begin{equation}\label{3}
\overline{S}_{f+g}^{\sigma_{1}\cup\sigma_{2}} \,\leq\,
\overline{S}_{f}^{\sigma_{1}\cup\sigma_{2}}\,+\,\overline{S}_{g}^{\sigma_{1}\cup\sigma_{2}},
\end{equation}
et
\begin{equation}\label{4}
\underline{S}_{f}^{\sigma_{1}\cup\sigma_{2}}\,+\,
\underline{S}_{g}^{\sigma_{1}\cup\sigma_{2}} \,\leq\,
\underline{S}_{f+g}^{\sigma_{1}\cup\sigma_{2}}.
\end{equation}
En utilisant les in\'egalit\'es \eqref{1}, \eqref{2}, \eqref{3} et
\eqref{4}, il vient alors~:
\begin{equation*}
\overline{S}_{f+g}^{\sigma_{1}\cup\sigma_{2}}\,\leq\,
\overline{S}_{f}^{\sigma_{1}}+
\overline{S}_{g}^{\sigma_{2}}\,\leq\,
\underline{S}_{f}^{\sigma_{1}} + \underline{S}_{g}^{\sigma_{2}} +
\varepsilon \leq \underline{S}_{f+g}^{\sigma_{1}\cup\sigma_{2}} +
\varepsilon.
\end{equation*}
D'apr\`es le th\'eor\`eme rappel\'e en introduction, on en d\'eduit que $f+g$ est
Riemann-int\'egrable sur $[a, b]$. De plus, de l'in\'egalit\'e
\begin{equation*}
\underline{S}_{f}^{\sigma_{1}} + \underline{S}_{g}^{\sigma_{2}}
\leq \underline{S}_{f+g}^{\sigma_{1}\cup\sigma_{2}},
\end{equation*}
on d\'eduit que
\begin{equation*}
\sup_{\sigma_1, \sigma_{2}}\left(\underline{S}_{f}^{\sigma_{1}} +
\underline{S}_{g}^{\sigma_{2}} \right) \leq \sup_{\sigma_1,
\sigma_{2}}\underline{S}_{f+g}^{\sigma_{1}\cup\sigma_{2}}.
\end{equation*}
Or $$\sup_{\sigma_1,
\sigma_{2}}\left(\underline{S}_{f}^{\sigma_{1}} +
\underline{S}_{g}^{\sigma_{2}} \right) =
\sup_{\sigma_{1}}\underline{S}_{f}^{\sigma_{1}} +
\sup_{\sigma_{2}}\underline{S}_{g}^{\sigma_{2}} = \int_{a}^{b}
f(x)\,dx + \int_{a}^{b} g(x)\,dx $$ et
$$\sup_{\sigma_1,
\sigma_{2}}\underline{S}_{f+g}^{\sigma_{1}\cup\sigma_{2}} =
\sup_{\sigma}\underline{S}_{f+g}^{\sigma} = \int_{a}^{b}\left(f(x)
+ g(x)\right)\,dx.$$ Ainsi
\begin{equation*}
 \int_{a}^{b}
f(x)\,dx + \int_{a}^{b} g(x)\,dx \leq \int_{a}^{b}\left(f(x) +
g(x)\right)\,dx.
\end{equation*}
De m\^eme, l'in\'egalit\'e
\begin{equation*}
\overline{S}_{f+g}^{\sigma_{1}\cup\sigma_{2}}\,\leq\,
\overline{S}_{f}^{\sigma_{1}}+ \overline{S}_{g}^{\sigma_{2}}
\end{equation*}
implique $ \int_{a}^{b}\left(f(x) + g(x)\right)\,dx \leq
\int_{a}^{b} f(x)\,dx + \int_{a}^{b} g(x)\,dx$. En conclusion,
 $ \int_{a}^{b}\left(f(x) +
g(x)\right)\,dx = \int_{a}^{b} f(x)\,dx + \int_{a}^{b} g(x)\,dx$.

 \item $\cdot$ Pour $\lambda = 0$ il n'y a rien a d\'emontrer. \\
 $\cdot$ Si $f$
est Riemann-int\'egrable sur $[a,b]$ et $\lambda > 0$, alors pour
tout subdivision $\sigma=\{a_{0}=a<\dots<a_{n}=b\}$ de $[a,b]$, on
a:
\begin{eqnarray*}
\inf\{\lambda f(x),\,x \in\,]a_{k-1},a_{k}[\,\} = \lambda
\inf\{f(x),\,x \in\,]a_{k-1},a_{k}[\,\}\\
\sup\{\lambda f(x),\,x \in\,]a_{k-1},a_{k}[\,\} = \lambda
\sup\{f(x),\,x \in\,]a_{k-1},a_{k}[\,\}.
\end{eqnarray*}
Par cons\'equent, $\underline{S}_{\lambda f}^{\sigma} = \lambda
\underline{S}_{f}^{\sigma}$ et $\overline{S}_{\lambda f}^{\sigma}
= \lambda \overline{S}_{f}^{\sigma}$. On en d\'eduit que
\begin{equation*}
\sup_{\sigma}\underline{S}_{\lambda f}^{\sigma} = \lambda
\sup_{\sigma}\underline{S}_{f}^{\sigma} = \lambda \int_{a}^{b}
f(x)\,dx = \lambda \inf_{\sigma}\overline{S}_{f}^{\sigma} =
\inf_{\sigma}\overline{S}_{\lambda f}^{\sigma}.
\end{equation*}
En conclusion, $\lambda f$ est Riemann-int\'egrable et
$\int_{a}^{b} \lambda f(x)\,dx = \lambda \int_{a}^{b} f(x)\,dx$.\\
$\cdot$ Si $f$ est Riemann-int\'egrable sur $[a,b]$ et $\lambda
<0$, alors pour tout subdivision
$\sigma=\{a_{0}=a<\dots<a_{n}=b\}$ de $[a,b]$, on a:
\begin{eqnarray*}
\inf\{\lambda f(x),\,x \in\,]a_{k-1},a_{k}[\,\} = \lambda
\sup\{f(x),\,x \in\,]a_{k-1},a_{k}[\,\}\\
\sup\{\lambda f(x),\,x \in\,]a_{k-1},a_{k}[\,\} = \lambda
\inf\{f(x),\,x \in\,]a_{k-1},a_{k}[\,\}.
\end{eqnarray*}
Par cons\'equent, $\underline{S}_{\lambda f}^{\sigma} = \lambda
\overline{S}_{f}^{\sigma}$ et $\overline{S}_{\lambda f}^{\sigma} =
\lambda \underline{S}_{f}^{\sigma}$. On en d\'eduit que
\begin{equation*}
\sup_{\sigma}\underline{S}_{\lambda f}^{\sigma} = \lambda
\inf_{\sigma}\overline{S}_{f}^{\sigma} = \lambda \int_{a}^{b}
f(x)\,dx = \lambda \sup_{\sigma}\underline{S}_{f}^{\sigma} =
\inf_{\sigma}\overline{S}_{\lambda f}^{\sigma}.
\end{equation*}
En conclusion, $\lambda f$ est Riemann-int\'egrable et
$\int_{a}^{b} \lambda f(x)\,dx = \lambda \int_{a}^{b} f(x)\,dx$.

\item Soient $f$ et $g$  deux fonctions Riemann-int\'egrables sur
$[a,b]$ telles que, pour tout $t\in [a,b]$, $f(t)\leq g(t)$. Soit
$\sigma=\{a_{0}=a< \dots < a_{n}=b\}$ une subdivision de $[a,b]$.
Alors
\begin{equation*}
\inf\{f(x), \, x \in\, ]a_{k-1}, a_{k}[\,\} \leq \inf\{g(x), \, x
\in\, ]a_{k-1}, a_{k}[\,\}.
\end{equation*}
Il en d\'ecoule que
\begin{equation*}
\sup_{\sigma} \underline{S}_{f}^{\sigma} \leq \sup_{\sigma}
\underline{S}_{f}^{\sigma},
\end{equation*}
c'est-\`a-dire $\int_{a}^{b} f(x)\,dx \leq \int_{a}^{b} g(x)\,dx$.

\item Soit $\{f_{i}\}_{i\in\mathbb{N}}$ une suite de fonctions
Riemann-int\'egrables, qui converge uniform\'ement vers $f$ sur
$[a, b]$. Soit $\varepsilon>0$ donn\'e. Il existe $N>0$ tel que
$\forall i> N$, $\sup_{[a, b]}|f_{i}(t) - f(t)| < \varepsilon$. En
particulier, $f_{i}(t) - \varepsilon < f(t) < f_{i}(t) +
\varepsilon$.  Pour un tel $i$, on en d\'eduit que pour toute
subdivision $\sigma = \{a_{0}= a < \dots < a_{n} = b\}$, on a
\begin{eqnarray*}
\sup_{]a_{k-1}, a_{k}[}f \,\leq \,\sup_{]a_{k-1}, a_{k}[}f_{i} +
\varepsilon\quad \text{et}\quad \inf_{]a_{k-1}, a_{k}[}f \geq
\inf_{]a_{k-1}, a_{k}[}f_{i} - \varepsilon
\end{eqnarray*}
En particulier~:
\begin{eqnarray*}
\sup_{]a_{k-1}, a_{k}[}f - \inf_{]a_{k-1}, a_{k}[}f\,\leq\,
\sup_{]a_{k-1}, a_{k}[}f_{i} - \inf_{]a_{k-1}, a_{k}[}f_{i} +
2\varepsilon.
\end{eqnarray*}
Il en d\'ecoule que~:
\begin{equation*}
\overline{S}_{f}^{\sigma} - \underline{S}_{f}^{\sigma}\, \leq\,
\overline{S}_{f_{i}}^{\sigma} - \underline{S}_{f_{i}}^{\sigma} +
2\varepsilon (b- a).
\end{equation*}
Comme $f_{i}$ est Riemann-int\'egrable, d'apr\`es le th\'eor\`eme
de l'introduction, il existe une subdivision $\sigma$ de $[a, b]$ telle que
$\overline{S}_{f_{i}}^{\sigma} - \underline{S}_{f_{i}}^{\sigma}
\leq \varepsilon$. On en d\'eduit que
\begin{equation*}
\overline{S}_{f}^{\sigma} - \underline{S}_{f}^{\sigma} \leq
\varepsilon \left(1 + 2(b-a)\right),
\end{equation*}
ce qui implique que $f$ est Riemann-int\'egrable.
\end{enumerate}
\fincorrection
\correction{005918}
Soit $f$ une fonction croissante $[a,b]$. Pour montrer que $f$ est
Riemann-int\'egrable, il suffit de trouver, pour tout
$\varepsilon>0$ donn\'e, une subdivision de $[a, b]$ telle que
$\overline{S}_{f}^{\sigma} - \underline{S}_{f}^{\sigma} <
\varepsilon$. Soit $\sigma=\{a_{0}=a < \dots < a_{n} = b\}$ la
subdivision r\'eguli\`ere de $[a,b]$, de pas
$\left(\frac{b-a}{n}\right)$. On a
\begin{equation*}
\inf_{]a_{k-1}, a_{k}[} f = f(a_{k-1}) \quad \text{et} \quad
\sup_{]a_{k-1}, a_{k}[} f = f(a_{k}).
\end{equation*}
Ainsi~:
\begin{eqnarray*}
\overline{S}_{f}^{\sigma} - \underline{S}_{f}^{\sigma} &=&
\sum_{k=1}^{n} (a_{k} - a_{k-1})\left( f(a_{k}) - f(a_{k-1})
\right) \\ &= &\left(\frac{b-a}{n}\right) \sum_{k=1}^{n}\left(
f(a_{k}) - f(a_{k-1}) \right)\\& =&
\left(\frac{b-a}{n}\right)\left(f(b)- f(a)\right).
\end{eqnarray*}
Pour $n$ assez grand, la subdivision r\'eguli\`ere de $[a,b]$
satisfait $\overline{S}_{f}^{\sigma} - \underline{S}_{f}^{\sigma}
< \varepsilon$. D'autre part, si $g$ est d\'ecroissante, $f = -g$
est croissante, donc $g$ est Riemann-int\'egrable par l'exercice
pr\'ec\'edent (question 2.) avec $\lambda = -1$.
\fincorrection
\correction{005919}
Une fonction $f$ continue sur $[a, b]$ est uniform\'ement continue
sur $[a, b]$. En particulier, pour tout $\varepsilon>0$, il existe
$n>0$ tel que
\begin{equation*}
|x - y| < \left(\frac{b-a}{n}\right) \Rightarrow |f(x) - f(y)|<
\varepsilon.
\end{equation*}
Soit $\sigma=\{a_{0}=a < \dots < a_{n} = b\}$ la subdivision
r\'eguli\`ere de $[a,b]$, de pas $\left(\frac{b-a}{n}\right)$. On
a~:
\begin{equation*}
\sup_{]a_{k-1}, a_{k}[} f - \inf_{]a_{k-1}, a_{k}[} f \leq
2\varepsilon.
\end{equation*}
Il vient alors~:
\begin{equation*}
\overline{S}_{f}^{\sigma} - \underline{S}_{f}^{\sigma} \leq
\left(\frac{b-a}{n}\right)\sum_{k=1}^{n} 2 \varepsilon = (b-a)
2\varepsilon,
\end{equation*}
ce qui permet de conclure gr\^ace au th\'eor\`eme de l'introduction que $f$
est Riemann-int\'egrable sur $[a,b]$.
\fincorrection
\correction{005920}
\begin{enumerate}
\item Consid\'erons la fonction $f~:[0,1]\rightarrow \mathbb{R}$
d\'efinie par~:
$$
 f(x) = \left\{\begin{array}{l} 1 \quad\text{si}\quad  x\in\mathbb{Q}\\
0 \quad\text{si} \quad x\in\mathbb{R}\setminus\mathbb{Q}
\end{array}\right..
$$
Pour toute subdivision $\sigma$ de $[a,b]$, on a~:
\begin{equation*}
\overline{S}_{f}^{\sigma} = 1 \quad \text{et}\quad
\underline{S}_{f}^{\sigma} = 0.
\end{equation*}
On en d\'eduit que $1 = \sup_{\sigma}\overline{S}_{f}^{\sigma}
\neq \inf_{\sigma}\underline{S}_{f}^{\sigma} = 0$, ce qui implique
que $f$ n'est pas Riemann-int\'egrable sur $[0, 1]$.

\item Consid\'erons la fonction $g~:[0,1]\rightarrow \mathbb{R}$
d\'efinie par~:
$$
 g(x) = \left\{\begin{array}{l} \frac{1}{q} \quad\text{si}\quad
 \, x=\frac{p}{q} \quad\text{avec}\,\, p \,\,\text{et} \,\,q \,\,\text{premiers entre eux}\\
0 \quad\text{si} \quad x\in\mathbb{R}\setminus\mathbb{Q} \text{ ou } x=0
\end{array}\right..
$$
Pour toute subdivision $\sigma$ de $[a,b]$, on a~:
\begin{equation*}
\underline{S}_{g}^{\sigma} = 0.
\end{equation*}
Pour tout $\varepsilon>0$ donn\'e, la fonction $g$ prend des
valeurs sup\'erieures \`a $\frac{\varepsilon}{b-a}$ en un nombre fini de
points seulement (les points $\frac{k}{q}$, avec
$\frac{1}{q}>\frac{\varepsilon}{b-a}$ ce qui équivaut à $q < \frac{b-a}{\varepsilon}$). Notons $x_{i}$, $i= 1,\dots, p$ ces
points ordonn\'es par ordre (strictement) croissant. 

Sur $[0,1]\setminus \{x_1,\ldots,x_p\}$ la fonction $g$ prend des valeurs
$\le \varepsilon$ et $\ge 0$. Ainsi avec la subdivision
$\sigma = \{x_1,\ldots,x_p\}$ nous obtenons :
\begin{equation*}
 0 \leq \overline{S}_{g}^{\sigma} \leq \frac{\varepsilon}{b-a} (b-a) = \varepsilon
\end{equation*}

% On a $p \leq \frac{1}{\varepsilon^2}$.
%  Notons \'egalement $\delta$ le minimum
% de la distance entre deux de ces points. Consid\'erons la
% subdivision $\sigma=\{a_{0}=a<\dots<a_{2p}=b\}$ d\'efinie par
% $a_{2k-1}=x_{i}-\min(\frac{\delta}{3}, \varepsilon^3)$ pour $k=1,
% \dots, p$, $a_{2k}=x_{i}+\min(\frac{\delta}{3},\varepsilon^3)$
% pour $k = 1, \dots, p-1$, et $a_{2p} = b$. Alors~:
% \begin{equation*}
% \overline{S}_{g}^{\sigma} \leq \varepsilon +
% p\min\left(\frac{\delta}{3},\varepsilon^3\right) \leq \varepsilon
% + \frac{1}{\varepsilon^{2}}\varepsilon^3 \leq 2\varepsilon.
% \end{equation*}

Comme On en conclut que $g$ est Riemann-int\'egrable sur $[0,1]$.
\end{enumerate}
\fincorrection
\correction{005921}
cf Andr\'e Gramain, \emph{Int\'egration}, p.~7, Hermann (1998).
\fincorrection
\correction{005922}
Pour tout $n\in\mathbb{N}$, on d\'efinit
$f_{n}:]0,1]\rightarrow\mathbb{R}$ par~: $f_{n}(x)= n e^{-nx}$.
Pour tout $x\in]0,1]$, on a $\lim_{n\rightarrow+\infty}f_{n}(x) =
\lim_{n\rightarrow+\infty}ne^{-nx} = 0$. On en d\'eduit que la
suite de fonctions $f_{n}$ converge ponctuellement (ou
\emph{simplement}) vers la fonction identiquement nulle
$f\equiv0$. On a $\int_{0}^{1}f(x)\,dx = 0$ mais
\begin{equation*}
\int_{0}^{1} f_{n}(x)\,dx = 1- e^{-n},
\end{equation*}
et $\lim_{n\rightarrow+\infty}\int_{0}^{1}f_{n}(x)\,dx = 1$. La
suite $(f_{n})_{n\in\mathbb{R}}$ ne converge pas uniform\'ement
vers $f$ sur $]0,1]$, car pour tout $\varepsilon>0$, et pour tout
$n\in\mathbb{N}$, on a~:
\begin{equation*}
\sup_{]0,
-\frac{1}{n}\log\left(\frac{\varepsilon}{n}\right)[}|f_{n}(x)-f(x)|
> \varepsilon.
\end{equation*}
\fincorrection
\correction{005923}
Soit $f~:[a,b] \rightarrow \mathbb{R}$ une fonction int\'egrable
au sens de Riemann. 

Notons $x_k = a+k\frac{b-a}{n}$, $k=1,\ldots,n$ les points où

Soit $a_0=a$, $a_{n+1}=b$ et $a_k = a+\frac{2k+1}{2n}$ pour $k=1,\ldots,n$.

Consid\'erons la subdivision 
$\sigma=\{a_{0}=a< \cdots<a_k<\cdots < a_{n}=b\}$ de $[a,b]$. 
Cette subdivision est presque régulière, seul le premier intervalle et le dernier
ont des longueurs différentes. Pour $k=1,\ldots,n-1$, $x_k$
est le milieu de $]a_k,a_{k+1}[$.

Notons 
$m_{k} = \inf\{ f(x), x\in ]a_{k-1},a_k[\}$ et
$M_k = \sup\{ f(x), x\in ]a_{k-1},a_k[\}$. 

Donc pour $k=1,\ldots,n-1$ on a $m_k \le f(x_k) \le M_k$.
Mais il faut aussi tenir compte de $f(x_n)=f(b)$ et des premiers et derniers intervalles.
D'où pour la minoration:
\begin{equation*}
\underline{S}_{f}^{\sigma} = (m_0+m_n)\frac{b-a}{2n} + \frac{b-a}{n} \sum_{k=1}^{n} m_k
\le (m_0+m_n)\frac{b-a}{2n} + \frac{b-a}{n} \sum_{k=1}^{n-1} f(x_k).
\end{equation*}
Cela donne 
\begin{equation*}
 \underline{S}_{f}^{\sigma} - (m_0+m_n+2f(b)) \frac{b-a}{2n} \le  \frac{b-a}{n}\sum_{k=1}^{n} f(x_k).
\end{equation*}

Quand $n$ tend vers $+\infty$ on trouve que $\underline{S}_{f}^{\sigma} \to \int_a^b f$
et $(m_0+m_n+2f(b)) \frac{b-a}{2n}\to 0$ cela donne l'inégalité :
$$\int_a^b f \le \lim_{n\rightarrow+\infty} \frac{b-a}{n}\sum_{k=1}^{n} f(x_k).$$

La somme $\overline{S}_{f}^{\sigma}$ conduit de manière similaire à l'inégalité inverse,
d'o\`u :
\begin{equation*}
\int_{a}^{b} f(x)\,dx =
\lim_{n\rightarrow+\infty}\left(\frac{b-a}{n}\right)
\sum_{k=1}^{n}f\left(a+k\frac{b-a}{n}\right).
\end{equation*}
On a~:
\begin{eqnarray*}
a)\quad\lim_{n\rightarrow+\infty}
\frac{1}{n}\sum_{k=1}^{n}\tan\frac{k}{n} = -\log(\cos 1)
\quad\quad\quad &
b)\quad\lim_{n\rightarrow+\infty}\sum_{k=1}^{n}\frac{n}{n^2+k^2} =
\frac{\pi}{4} \quad\quad\quad\\
c)\quad\lim_{n\rightarrow+\infty}\sum_{k=1}^{n}\log\left(\frac{n}{n+k}\right)^{\frac{1}{n}}
= -2\ln 2 + 1. &
\end{eqnarray*}
\fincorrection
\correction{005924}
\begin{enumerate}
\item
\begin{multline*}
\overset{b}{\underset{a}{\int }}f(a+b-x)dx=-\overset{b}{\underset{a}{\int }}%
f(a+b-x)(a+b-x)^{\prime }dx\\
=-\overset{\varphi (b)}{\underset{\varphi (a)}{%
\int }}f(t)dt=-\overset{a}{\underset{b}{\int }}f(t)dt=\overset{b}{\underset{a%
}{\int }}f(t)dt
\end{multline*}

o\`{u} $\varphi :[a,b]\rightarrow \lbrack a,b]$, $\varphi
(x)=a+b-x$ est une fonction de classe $C^{1}$.

\item $a)$
\begin{multline*}
I:=\overset{\pi }{\underset{0}{\int }}\frac{x\sin x}{1+\cos ^{2}x}dx=\overset%
{\pi }{\underset{0}{\int }}\frac{(\pi -x)\sin (\pi -x)}{1+\cos ^{2}(\pi -x)}%
dx=\overset{\pi }{\underset{0}{\int }}\frac{(\pi -x)\sin x}{1+\cos ^{2}x}%
dx\\
=\pi \overset{\pi }{\underset{0}{\int }}\frac{\sin x}{1+\cos
^{2}x}dx-I
\end{multline*}

\begin{multline*}
I =\frac{\pi }{2}\overset{\pi }{\underset{0}{\int }}\frac{\sin
x}{1+\cos ^{2}x}dx=-\frac{\pi }{2}\overset{\pi }{\underset{0}{\int
}}\frac{(\cos
x)^{\prime }}{1+\cos ^{2}x}dx=-\frac{\pi }{2}\overset{\varphi (\pi )}{%
\underset{\varphi (0)}{\int }}\frac{1}{1+t^{2}}dt\\
=-\frac{\pi }{2}\overset{-1}%
{\underset{1}{\int }}\frac{1}{1+t^{2}}dt=\frac{\pi }{2}\overset{1}{\underset{-1}{\int }}\frac{1}{1+t^{2}}dt=\frac{%
\pi ^{2}}{4}.
\end{multline*}

o\`{u} $\varphi :[0,\pi ]\rightarrow \lbrack -1,1]$, $\varphi
(x)=\cos x$ est une fonction de classe $C^{1}$.

$b)$
\begin{multline*}
J :=\overset{\pi /4}{\underset{0}{\int }}\log (1+\tan x)dx=\overset{\pi /4}{%
\underset{0}{\int }}\log \left(1+\tan (\frac{\pi }{4}-x)\right)dx\\
=\overset{\pi /4}{%
\underset{0}{\int }}\log \left(1+\frac{1-\tan x}{1+\tan x}\right)dx
=\overset{\pi /4}{\underset%
{0}{\int }}\log \left(\frac{2}{1+\tan x}\right)dx =\frac{\pi }{4}\log
2-J
\end{multline*}

d'o\`{u} la valeur de l'int\'{e}grale est $J=\frac{\pi }{8}\log
2.$
\end{enumerate}
\fincorrection
\nocorrection
\nocorrection
\nocorrection
\nocorrection
\nocorrection
\nocorrection
\correction{002100}
  \begin{enumerate}
  \item Soit 
$$u_n =n \sum_{k=0}^{n-1}\frac 1{k^2+n^2} = \frac 1n  \sum_{k=0}^{n-1}\frac 1{1+\big(\frac k n \big)^2}.$$
En posant $f(x) = \frac 1 {1+x^2}$ nous venons d'écrire la somme de Riemann correspondant à 
$\int_0^1 f(x) dx$. Cette intégrale se calcule facilement : 
$$\int_0^1 f(t) dt = \int_0^1 \frac {dx} {1+x^2} = \big[\arctan x\big]_0^1 = \frac \pi 4.$$
La somme de Riemann $u_n$ convergeant vers $\int_0^1 f(x) dx$ nous venons de montrer que
$(u_n)$ converge vers $\frac \pi 4$.

  \item Soit $v_n=\prod\limits_{k=1}^n\left(1+\frac{k^2}{n^2}\right) ^{\frac 1n}$, notons 
$$w_n = \ln v_n = \sum_{k=1}^n \ln\left( \left(1+\frac{k^2}{n^2}\right)^{\frac 1n} \right) 
= \frac 1 n \sum_{k=1}^n \ln \left(1+\frac{k^2}{n^2}\right).$$
En posant $g(x) = \ln (1+x^2)$ nous reconnaissons la somme de Riemann correspondant à
$I = \int_0^1 g(x)dx$.

Calculons cette intégrale : 
\begin{align*}
 I &= \int_0^1 g(x)dx = \int_0^1 \ln(1+x^2) dx \\
   &= \big[x\ln(1+x^2)\big]_0^1 - \int_0^1 x \frac{2x}{1+x^2}dx \quad \text{par intégration par parties} \\
   &= \ln 2 -2 \int_0^1 1-\frac 1{1+x^2} dx \\
   &= \ln 2  - 2\big[x-\arctan x\big]_0^1 \\
   &= \ln 2 - 2 + \frac \pi 2. \\
\end{align*}


Nous venons de prouver que $w_n=\ln v_n$ converge vers $I=\ln 2 - 2 + \frac \pi 2$,
donc $v_n = \exp w_n$ converge vers $\exp(\ln 2 - 2 + \frac \pi 2) = 2e^{\frac \pi 2 -2}$.
Bilan $(v_n)$ a pour limite $2e^{\frac \pi 2 -2}$.

  \end{enumerate}
\fincorrection
\correction{004228}
\begin{enumerate}
  \item 
  \item $\ln k$.
  \item $\frac \pi8$.
  \item $\frac 4e$.
  \item $\frac13 \int_{t=0}^{3\pi}\frac{d t}{2+\cos t}
             =  \int_{t=0}^{\pi}\frac{d t}{2+\cos t} = \frac\pi{\sqrt3}$.
  \item $\frac 43n\sqrt n$.
  \item $\frac 4\pi$.
\end{enumerate}
 \fincorrection
\nocorrection
\nocorrection
\correction{004231}
\begin{enumerate}
  \item 
  \item $\exp\left(\frac\pi4\right)$.
\end{enumerate}
 \fincorrection
\nocorrection
\correction{005446}

\begin{enumerate}
\item  Pour $n\geq1$, 

$$u_n=\frac{1}{n^3}\sum_{k=1}^{n}k^2\sin\frac{k\pi}{n}=\frac{1}{n}\sum_{k=1}^{n}(\frac{k}{n})^2\sin\frac{k\pi}{n}
=\frac{1}{n}\sum_{k=1}^{n}f(\frac{k}{n}),$$

où $f(x)=x^2\sin(\pi x)$. $u_n$ est donc une somme de \textsc{Riemann} à pas constant associée à la fonction continue $f$ sur $[0,1]$. Quand $n$ tend vers $+\infty$, le pas $\frac{1}{n}$ tend vers $0$ et on sait que $u_n$ tend vers
 
\begin{align*}\ensuremath
\int_{0}^{1}x^2\sin(\pi x)\;dx&=\left[-\frac{1}{\pi}x^2\cos(\pi x)\right]_{0}^{1}+\frac{2}{\pi}\int_{0}^{1}x\cos(\pi x)\;dx
=\frac{1}{\pi}+\frac{2}{\pi}(\left[\frac{1}{\pi}x\sin(\pi x)\right]_{0}^{1}-\frac{1}{\pi}\int_{0}^{1}\sin(\pi x)\;dx)\\
  &=\frac{1}{\pi}-\frac{2}{\pi^2}\left[-\frac{1}{\pi}\cos(\pi x)\right]_{0}^{1}=\frac{1}{\pi}-\frac{2}{\pi^2}(\frac{1}{\pi}+\frac{1}{\pi})\\
  &=\frac{1}{\pi}-\frac{4}{\pi^3}.
\end{align*}

\item  On peut avoir envie d'écrire~:~

$$\ln(u_n)=\frac{1}{n}(\sum_{k=1}^{n}(\ln(a+k)-\ln k))=\frac{1}{n}\sum_{k=1}^{n}\ln(1+\frac{a}{k}).$$

La suite de nombres $a$, $\frac{a}{2}$,..., $\frac{a}{n}$ \og est une subdivision (à pas non constant) de $[0,a]$~\fg~mais malheureusement son pas $a-\frac{a}{2}=\frac{a}{2}$ ne tend pas vers $0$ quand $n$ tend vers $+\infty$. On n'est pas dans le même type de problèmes.

Rappel. (exo classique) Soit $v$ une suite strictement positive telle que la suite $(\frac{v_{n+1}}{v_n})$ tend vers un réel positif $\ell$, alors la suite $(\sqrt[n]{v_n})$ tend encore vers $\ell$.

Posons $v_n=\frac{1}{n!}\prod_{k=1}^{n}(a+k)$ puis $u_n=\sqrt[n]{v_n}$.

$$\frac{v_{n+1}}{v_n}=\frac{a+n+1}{n+1}\rightarrow1,$$

et donc $\lim_{n\rightarrow +\infty}u_n=1$.

\item  Encore une fois, ce n'est pas une somme de \textsc{Riemann}. On tente un encadrement assez large~:~pour $1\leq k\leq n$, 

$$\frac{n+k}{n^2+n}\leq\frac{n+k}{n^2+k}\leq\frac{n+k}{n^2}.$$

En sommant ces inégalités, il vient 

$$\frac{1}{n^2+n}\sum_{k=1}^{n}(n+k)\leq\sum_{k=1}^{n}\frac{n+k}{n^2+k}\leq\frac{1}{n^2}\sum_{k=1}^{n}(n+k),$$

et donc ((premier terme + dernier terme)$\times$nombre de termes/2),

$$\frac{1}{n^2+n}\frac{((n+1)+2n)n}{2}\leq u_n\leq\frac{1}{n^2}\frac{((n+1)+2n)n}{2},$$

et finalement, $\frac{3n+1}{2(n+1)}\leq u_n\leq\frac{3n+1}{2n}$. Or, $\frac{3n+1}{2(n+1)}$ et $\frac{3n+1}{2n}$ tendent tous deux vers $\frac{3}{2}$. Donc, $u_n$ tend vers $\frac{3}{2}$.

\item  Tout d'abord

$$u_n=\sum_{k=1}^{n}\frac{1}{\sqrt{n^2-k^2}}=\frac{1}{n}\sum_{k=1}^{n}\frac{1}{\sqrt{1-(\frac{k}{n})^2}}=
\frac{1}{n}\sum_{k=1}^{n}f(\frac{k}{n}),$$ 

où $f(x)=\frac{1}{\sqrt{1-x^2}}$ pour $x\in[0,1[$. $u_n$ est donc effectivement une somme de \textsc{Riemann} à pas constant associée à la fonction $f$ mais malheureusement, cette fonction n'est pas continue sur $[0,1]$, ou même prolongeable par continuité en $1$. On s'en sort néanmoins en profitant du fait que $f$ est croissante sur $[0,1[$.

Puisque $f$ est croissante sur $[0,1[$, pour $1\leq k\leq n-2$, on a $\frac{1}{n}\frac{1}{\sqrt{1-(\frac{k}{n})^2}}\leq\int_{k/n}^{(k+1)/n}\frac{1}{\sqrt{1-x^2}}\;dx$, et pour $1\leq k\leq n-1$, $\frac{1}{n}\frac{1}{\sqrt{1-(\frac{k}{n})^2}}\geq\int_{(k-1)/n}^{k/n}\frac{1}{\sqrt{1-x^2}}\;dx$. En sommant ces inégalités, on obtient

$$u_n=\frac{1}{n}\sum_{k=1}^{n-1}\frac{1}{\sqrt{1-(\frac{k}{n})^2}}\geq\sum_{k=1}^{n-1}\int_{(k-1)/n}^{k/n}\frac{1}{\sqrt{1-x^2}}\;dx=\int_{0}^{1-\frac{1}{n}}\frac{1}{\sqrt{1-x^2}}\;dx=\Arcsin(1-\frac{1}{n}),$$

et 

\begin{align*}\ensuremath
u_n&=\frac{1}{n}\sum_{k=1}^{n-2}\frac{1}{\sqrt{1-(\frac{k}{n})^2}}+\frac{1}{\sqrt{n^2-(n-1)^2}}\leq\int_{\frac{1}{n}}^{1-\frac{1}{n}}\frac{1}{\sqrt{1-x^2}}\;dx+\frac{1}{\sqrt{2n-1}}\\
 &=\Arcsin(1-\frac{1}{n})-\Arcsin\frac{1}{n}+\frac{1}{\sqrt{2n-1}}.
\end{align*}

Quand $n$ tend vers $+\infty$, les deux membres de cet encadrement tendent vers $\Arcsin1=\frac{\pi}{2}$, et donc $u_n$ tend vers $\frac{\pi}{2}$.

\item  Pour $1\leq k\leq n$, $\sqrt{k}-1\leq E(\sqrt{k})\leq \sqrt{k}$, et en sommant,

$$\frac{1}{n\sqrt{n}}\sum_{k=1}^{n}\sqrt{k}-\frac{1}{\sqrt{n}}\leq u_n\leq \frac{1}{n\sqrt{n}}\sum_{k=1}^{n}\sqrt{k}.$$

Quand $n$ tend vers $+\infty$, $\frac{1}{\sqrt{n}}$ tend vers $0$ et la somme de \textsc{Riemann} $\frac{1}{n\sqrt{n}}\sum_{k=1}^{n}\sqrt{k}=\frac{1}{n}\sum_{k=1}^{n}\sqrt{\frac{k}{n}}$ tend vers $\int_{0}^{1}\sqrt{x}\;dx=\frac{3}{2}$. Donc, $u_n$ tend vers $\frac{3}{2}$.

\item  $u_n=\frac{1}{n}\sum_{k=1}^{n}\frac{(k/n)^2}{1+8(k/n)^3}$ tend vers $\int_{0}^{1}\frac{x^2}{8x^3+1}\;dx=\left[\frac{1}{24}\ln|8x^3+1|\right]_{0}^{1}=\frac{\ln3}{12}$.

\item  $u_n=\sum_{k=0}^{n-1}\frac{1}{2(k+n)+1}=\frac{1}{2}\frac{2}{n}\sum_{k=0}^{n-1}\frac{1}{2+\frac{2k+1}{n}}$ tend vers $\frac{1}{2}\int_{0}^{2}\frac{1}{2+x}\;dx=\frac{1}{2}(\ln4-\ln2)=\ln2$.

\item  Soit $f(x)=\frac{1}{x^2}e^{-1/x}$ si $x>0$ et $0$ si $x=0$. $f$ est continue sur $[0,1]$ (théorèmes de croissances comparées). Donc, $u_n=\frac{1}{n}\sum_{k=1}^{n}f(\frac{k}{n})$ tend vers $\int_{0}^{1}f(x)\;dx$. Pour $x\in[0,1]$, posons $F(x)=\int_{x}^{1}f(t)\;dt$. Puisque $f$ est continue sur $[0,1]$, $F$ l'est et

$$\int_{0}^{1}f(x)\;dx=F(0)=\lim_{x\rightarrow 0,\;x>0}F(x)=\lim_{x\rightarrow 0,\;x>0}\left[e^{-1/t}\right]_{x}^{1}=\lim_{x\rightarrow 0,\;x>0}(e^{-1}-e^{-1/x})=\frac{1}{e}.$$

Donc, $u_n$ tend vers $\frac{1}{e}$ quand $n$ tend vers $+\infty$.

\end{enumerate}
\fincorrection
\correction{005447}
Supposons $f$ de classe $C^2$ sur $[0,1]$.
Soit $F$ une primitive de $f$ sur $[0,1]$. Soit $n$ un entier naturel non nul.

$$u_n=\int_{0}^{1}f(t)\;dt-\frac{1}{n}\sum_{k=0}^{n-1}f(\frac{k}{n})=\sum_{k=0}^{n-1}(\int_{k/n}^{(k+1)/n}f(t)\;dt-\frac{1}{n}f(\frac{k}{n}))=\sum_{k=0}^{n-1}(F(\frac{k+1}{n})-F(\frac{k}{n})-\frac{1}{n}F'(\frac{k}{n})).$$

$f$ est de classe $C^2$ sur le segment $[0,1]$. Par suite, $F^{(3)}=f''$ est définie et bornée sur ce segment. En notant $M_2$ la borne supérieure de $|f''|$ sur $[0,1]$, l'inégalité de \textsc{Taylor}-\textsc{Lagrange} à l'ordre 3 appliquée à $F$ sur le segment $[0,1]$ fournit

$$\left|F(\frac{k+1}{n})-F(\frac{k}{n})-\frac{1}{n}F'(\frac{k}{n})-\frac{1}{2n^2}F''(\frac{k}{n})\right|
\leq\frac{(1/n)^3M_2}{6},$$

et donc,
 
\begin{align*}\ensuremath
\left|\sum_{k=0}^{n-1}[F(\frac{k+1}{n})-F(\frac{k}{n})-\frac{1}{n}F'(\frac{k}{n})-\frac{1}{2n^2}F''(\frac{k}{n})]
\right|&\leq\sum_{k=0}^{n-1}|F(\frac{k+1}{n})-F(\frac{k}{n})-\frac{1}{n}F'(\frac{k}{n})-\frac{1}{2n^2}F''(\frac{k}{n})|\\
 &\leq\sum_{k=0}^{n-1}\frac{(1/n)^3M_2}{6}=\frac{M_2}{6n^2}.
\end{align*}

Ainsi, $\sum_{k=0}^{n-1}[F(\frac{k+1}{n})-F(\frac{k}{n})-\frac{1}{n}F'(\frac{k}{n})-\frac{1}{2n^2}F''(\frac{k}{n})]=O(\frac{1}{n^2})$, ou encore $\sum_{k=0}^{n-1}[F(\frac{k+1}{n})-F(\frac{k}{n})-\frac{1}{n}F'(\frac{k}{n})-\frac{1}{2n^2}F''(\frac{k}{n})]=o(\frac{1}{n})$, ou enfin,

$$u_n=\sum_{k=0}^{n-1}\frac{1}{2n^2}F''(\frac{k}{n})+o(\frac{1}{n}).$$

Maintenant, 

$$\sum_{k=0}^{n-1}\frac{1}{2n^2}F''(\frac{k}{n})=\frac{1}{2n}.\frac{1}{n}\sum_{k=0}^{n-1}f'(\frac{k}{n}).$$

Or, la fonction $f'$ est continue sur le segment $[0,1]$. Par suite, la somme de \textsc{Riemann} $\frac{1}{n}\sum_{k=0}^{n-1}f'(\frac{k}{n})$ tend vers $\int_{0}^{1}f'(t)\;dt=f(1)-f(0)$ et donc 

$$\frac{1}{2n}\frac{1}{n}\sum_{k=0}^{n-1}f'(\frac{k}{n})=\frac{1}{2n}(f(1)-f(0)+o(1))=\frac{f(1)-f(0)}{2n}+o(\frac{1}{n}).$$

Finalement,

$$\int_{0}^{1}f(t)\;dt-\frac{1}{n}\sum_{k=0}^{n-1}f(\frac{k}{n})=\frac{f(1)-f(0)}{2n}+o(\frac{1}{n}).$$

\fincorrection
\correction{000805}
 $\displaystyle S=\frac{\pi}{(1-\lambda^2)^{3/2}}$.
\fincorrection
\correction{000806}
 $\displaystyle L=\frac{3\pi a}{2},\quad
A_1=\frac{5\pi-9\sqrt3}{32}a^2,\quad
A_2=\frac{5\pi+18\sqrt3}{32}a^2$.
\fincorrection
\correction{000807}
 $A=4\pi^2Rr,\quad V=2\pi^2Rr^2$.
\fincorrection
\correction{000808}
 $\displaystyle L=8R,\quad A=3\pi R^2,\quad V_1=5\pi^2R^3,\quad
V_2=6\pi^3R^3,\quad A_1=\frac{64\pi R^2}{3},\quad
A_2=16\pi^2R^2$.
\fincorrection
\correction{000809}
 $\displaystyle L=8(n+1)r=8\frac{n+1}{n}R,\quad
A=\pi(n+1)(n+2)r^2=\pi\frac{(n+1)(n+2)}{n^2}R^2,\quad
S=\frac{128\pi R^2}{5},\quad
V=\frac{64\pi R^3}{3}$.
\fincorrection
\correction{000810}
 $L=4R\bigl(\sqrt2+\ln(1+\sqrt2)\bigr)$.
\fincorrection
\correction{002098}
$\int_{-R}^R\sqrt{R^2-x^2}dx=\frac \pi 2R^2$.
\fincorrection
\correction{002099}
La courbe d'équation $y=x^2/2$ est une parabole, la courbe 
d'équation $y=\frac 1{1+x^2}$ est une courbe en cloche. Dessinez les deux graphes.
Ces deux courbes délimitent une région dont nous allons calculer l'aire.
Tout d'abord ces deux courbes s'intersectent
aux points d'abscisses $x=+1$ et $x=-1$ : cela se devine sur le graphique puis
se vérifie en résolvant l'équation $\frac{x^2}2=\frac 1{x^2+1}$.

Nous allons calculer deux aires :
\begin{itemize}
  \item L'aire $\mathcal{A}_1$ de la région sous la parabole, au-dessus de l'axe des abscisses et 
entre les droites d'équation $(x=-1)$ et $(x=+1)$.
Alors 
$$\mathcal{A}_1 = \int_{-1}^{+1} \frac{x^2}2 \, dx = \left[ \frac{x^3}{6} \right]_{-1}^{+1} = \frac 13.$$

  \item L'aire $\mathcal{A}_2$ de la région sous la cloche, au-dessus de l'axe des abscisses et 
entre les droites d'équation $(x=-1)$ et $(x=+1)$.
Alors 
$$\mathcal{A}_2 = \int_{-1}^{+1} \frac 1{x^2+1} \, dx = \left[ \arctan x \right]_{-1}^{+1} = \frac{\pi}{2}.$$

  \item L'aire $\mathcal{A}$ sous la cloche et au-dessus de la parabole vaut maintenant
$$\mathcal{A}= \mathcal{A}_2 - \mathcal{A}_1 = \frac{\pi}{2} - \frac 13.$$
\end{itemize}
\fincorrection
\nocorrection
\nocorrection
\nocorrection
\nocorrection
\correction{006863}
 Calculons seulement un quart de l'aire : la partie du quadrant $x\ge 0, y\ge 0$.
Pour ce quadrant les points de l'ellipse ont une abscisse $x$ qui vérifie $0 \le x \le a$.
Et la relation $\frac{x^2}{a^2}+ \frac{y^2}{b^2} = 1$ donne $y = b\sqrt{1-\frac{x^2}{a^2}}$.

\medskip

Nous devons donc calculer l'aire sous la courbe d'équation $y = b\sqrt{1-\frac{x^2}{a^2}}$,
au-dessus de l'axe des abscisses et entre les droites d'équations $(x=0)$ et $(x=a)$ (faites un dessin !).

Cette aire vaut donc : $\displaystyle \int_0^a b\sqrt{1-\frac{x^2}{a^2}} dx$.
Nous allons calculer cette intégrale à l'aide du changement de variable 
$x=a \cos u$ qui donne $dx = -a\sin u \, du$. La variable $x$ variant de 
$x=0$ à $x=a$ alors la nouvelle variable $u$ varie du $u=\frac \pi 2$ (pour lequel on a bien
$a\cos \frac\pi 2 = 0$) à $u=\frac \pi 2$ (pour lequel on a bien
$a\cos 0 = a$). Autrement dit la fonction $u \mapsto a\cos u$ est une bijection
de $[\frac\pi2,0]$ vers $[0,a]$.

\begin{align*}
\int_0^a b\sqrt{1-\frac{x^2}{a^2}} dx
  & = \int_{\frac \pi 2}^0 b \sqrt{1-\cos^2 u} (-a \sin u \, du) \quad \text{ en posant } x=a \cos u \\
  & = \int_{\frac \pi 2}^0 b\sin u (-a \sin u \, du) \\
  & = ab \int_0^{\frac \pi 2} \sin^2 u \, du \\
  & = ab \int_0^{\frac \pi 2} \frac{1-\cos(2u)}{2} \, du \\ 
  &= ab \left[ \frac u2 - \frac{\sin(2u)}{4} \right]_0^{\frac \pi 2} \\
  &= \frac{\pi a b}{4} \\
\end{align*}
L'aire d'un quart d'ellipse est donc $\frac{\pi a b}{4}$. 

Conclusion : l'aire d'une ellipse est $\pi a b$, où $a$ et $b$ sont les longueurs des demi-axes.
Si $a=b=r$ on retrouve que l'aire d'un disque de rayon $r$ est $\pi r^2$.
\fincorrection
\nocorrection
\nocorrection
\nocorrection
\nocorrection
\nocorrection
\nocorrection
\nocorrection
\nocorrection
\nocorrection
\nocorrection
\nocorrection
\correction{005470}
\begin{enumerate}
\item  On pose $t=\frac{1}{x}$ et donc $x=\frac{1}{t}$ et $dx=-\frac{1}{t^2}\;dt$. On obtient

$$I=\int_{1/a}^{a}\frac{\ln x}{x^2+1}\;dx=-\int_{a}^{1/a}\frac{\ln(1/t)}{\frac{1}{t^2}+1}\frac{1}{t^2}\;dt=-\int_{1/a}^{a}\frac{\ln t}{t^2+1}\;dt=-I,$$

et donc, $I=0$.

\item  ($p$ et $q$ sont des entiers naturels)

$\cos(px)\cos(qx)=\frac{1}{2}(\cos(p+q)x+\cos(p-q)x)$ et donc,

Premier cas. Si $p\neq q$,

$$\int_{0}^{\pi}\cos(px)\cos(qx)\;dx=\frac{1}{2}\left[\frac{\sin(p+q)x}{p+q}+\frac{\sin(p-q)x}{p-q}\right]_{0}^{\pi}=0.$$

Deuxième cas. Si $p=q\neq0$,

$$\int_{0}^{\pi}\cos(px)\cos(qx)\;dx=\frac{1}{2}\int_{0}^{\pi}(1+\cos(2px))\;dx=\frac{1}{2}\int_{0}^{\pi}dx=\frac{\pi}{2}.$$

Troisième cas. Si $p=q=0$. $\int_{0}^{\pi}\cos(px)\cos(qx)\;dx=\int_{0}^{\pi}\;dx=\pi$.

La démarche est identique pour les deux autres et on trouve $\int_{0}^{\pi}\sin(px)\sin(qx)\;dx=0$ si $p\neq q$ et $\frac{\pi}{2}$ si $p=q\neq0$ puis $\int_{0}^{\pi}\sin(px)\cos(qx)\;dx=0$ pour tout choix de $p$ et $q$.

\item  La courbe d'équation $y=\sqrt{(x-a)(b-x)}$ ou encore $\left\{
\begin{array}{l}
x^2+y^2-(a+b)x+ab=0\\
y\geq0
\end{array}
\right.$ est le demi-cercle de diamètre $[\left(\begin{array}{c}
a\\
0
\end{array}
\right),\left(\begin{array}{c}
b\\
0
\end{array}
\right)]$. Par suite, si $a\leq b$, $I=\frac{\pi R^2}{2}=\frac{\pi(b-a)^2}{8}$ et si $a>b$, $I=-\frac{\pi(b-a)^2}{8}$.

\item  L'intégrale proposée est somme de quatre intégrales. Chacune d'elles est la somme des aires de deux triangles. Ainsi, $I=\frac{1}{2}((1^2+3^2)+(2^2+2^2)+(3^2+1^2)+4^2)=22$.
\item  On pose $u=\frac{1}{x}$. On obtient

\begin{align*}\ensuremath
I&=\int_{1/2}^{2}\left(1+\frac{1}{x^2}\right)\Arctan x\;dx=\int_{2}^{1/2}(1+u^2)\Arctan u\frac{-du}{u^2}=\int_{1/2}^{2}(1+\frac{1}{u^2})(\frac{\pi}{2}-\Arctan u)\;du\\
 &=\frac{\pi}{2}((2-\frac{1}{2})-(\frac{1}{2}-2))-I).
\end{align*}

Par suite, $I=\frac{3\pi}{2}-I$ et donc $I=\frac{3\pi}{4}$.

\item  $I=\int_{-1}^{1}\sqrt{1+|x(1-x)|}\;dx=\int_{-1}^{0}\sqrt{1+x(x-1)}\;dx+\int_{0}^{1}\sqrt{1+x(1-x)}\;dx=I_1+I_2$.

Pour $I_1$, $1+x(x-1)=x^2-x+1=(x-\frac{1}{2})^2+(\frac{\sqrt{3}}{2})^2$ et on pose $x-\frac{1}{2}=\frac{\sqrt{3}}{2}\sh t$ et donc $dx=\frac{\sqrt{3}}{2}\ch t\;dt$.

\begin{align*}\ensuremath
I_1&=\int_{\ln(2-\sqrt{3})}^{-\ln(\sqrt{3})}\frac{\sqrt{3}}{2}\sqrt{\sh^2t+1}\;\frac{\sqrt{3}}{2}\ch t\;dt=
\frac{3}{4}\int_{\ln(2-\sqrt{3})}^{-\ln(\sqrt{3})}\ch^2t\;dt=\frac{3}{16}\int_{\ln(2-\sqrt{3})}^{-\ln(\sqrt{3})}(e^{2t}+e^{-2t}+2)\;dt\\
 &=\frac{3}{16}(\frac{1}{2}(e^{-2\ln(\sqrt{3})}-e^{2\ln(2-\sqrt{3})})-\frac{1}{2}(e^{2\ln(\sqrt{3})}-e^{-2\ln(2-\sqrt{3})})+2(-\ln(\sqrt{3})-\ln(2-\sqrt{3})))\\
 &=\frac{3}{16}(\frac{1}{2}(\frac{1}{3}-(2-\sqrt{3})^2)-\frac{1}{2}(3-\frac{1}{(2-\sqrt{3})^2})-2\ln(2\sqrt{3}-3))\\
 &=\frac{3}{16}(-\frac{4}{3}+\frac{1}{2}(-(2-\sqrt{3})^2+(2+\sqrt{3})^2))-2\ln(2\sqrt{3}-3))\\
 &=-\frac{1}{4}+\frac{3\sqrt{3}}{4}-\frac{3}{8}\ln(2\sqrt{3}-3).
\end{align*}

Pour $I_2$, $1+x(1-x)=-x^2+x+1=-(x-\frac{1}{2})^2+(\frac{\sqrt{5}}{2})^2$ et on pose $x-\frac{1}{2}=\frac{\sqrt{3}}{2}\sin t$ et donc $dx=\frac{\sqrt{3}}{2}\cos t\;dt$.

\begin{align*}\ensuremath
I_2&=\int_{-\Arcsin\frac{1}{\sqrt{5}}}^{\Arcsin\frac{1}{\sqrt{5}}}\frac{\sqrt{3}}{2}\sqrt{1-\sin^2t}\;\frac{\sqrt{3}}{2}\cos t\;dt=
\frac{3}{4}\int_{-\Arcsin\frac{1}{\sqrt{5}}}^{\Arcsin\frac{1}{\sqrt{5}}}\cos^2t\;dt
=\frac{3}{8}\int_{-\Arcsin\frac{1}{\sqrt{5}}}^{\Arcsin\frac{1}{\sqrt{5}}}(1+\cos(2t))\;dt\\
 &=\frac{3}{8}(2\Arcsin\frac{1}{\sqrt{5}}+2\left[\sin t\cos t\right]_{0}^{\Arcsin\frac{1}{\sqrt{5}}}
 =\frac{3}{4}\Arcsin\frac{1}{\sqrt{5}}+\frac{3}{4}\frac{1}{\sqrt{5}}\sqrt{1-\frac{1}{5}}\\
 &=\frac{3}{4}\Arcsin\frac{1}{\sqrt{5}}+\frac{3}{10}...
\end{align*}

\item 

\begin{align*}\ensuremath
I&=\int_{0}^{\pi}\frac{x\sin x}{1+\cos^2x}\;dx=\int_{\pi}^{0}\frac{(\pi-u)\sin(\pi-u)}{1+\cos^2(\pi-u)}\;-du
=\pi\int_{0}^{\pi}\frac{\sin u}{1+\cos^2u}\;du-\int_{0}^{\pi}\frac{u\sin u}{1+\cos^2u}\;du\\
 &=-\pi\left[\Arctan(\cos u)\right]_{0}^{\pi}-I=\frac{\pi^2}{2}-I,
\end{align*}

et donc, $I=\frac{\pi^2}{4}$.

\item  Pour $n\in\Nn^*$, posons $I_n=\int_{1}^{x}\ln^nt\;dt$.

$$I_{n+1}=\left[t\ln^{n+1}t\right]_{1}^{x}-(n+1)\int_{1}^{x}t\ln^nt\frac{1}{t}\;dt=x\ln^{n+1}x-(n+1)I_n.$$

Donc, $\forall n\in\Nn^*,\;\frac{I_{n+1}}{(n+1)!}+\frac{I_n}{n!}=\frac{x(\ln x)^{n+1}}{(n+1)!}$, et de plus, $I_1=x\ln x-x+1$.

Soit $n\geq2$.

$$\sum_{k=1}^{n-1}(-1)^k(\frac{I_k}{k!}+\frac{I_{k+1}}{(k+1)!})=\sum_{k=1}^{n-1}(-1)^k\frac{I_k}{k!}+\sum_{k=2}^{n}(-1)^{k-1}\frac{I_k}{k!}=-I_1-(-1)^n\frac{I_n}{n!},$$

Par suite,

$$I_n=(-1)^nn!(\sum_{k=1}^{n-1}(-1)^k\frac{x(\ln x)^{k+1}}{(k+1)!}-x\ln x+x-1)=(-1)^nn!(1-\sum_{k=0}^{n}(-1)^{k}\frac{x(\ln x)^{k}}{k!}).$$
\end{enumerate}
\fincorrection
\correction{006865}
\begin{enumerate}
  \item $\int (\cos x) ^{1234} \sin x \, d x$

En posant le changement de variable $u = \cos x$ on a $x=\arccos u$ et $du = - \sin x \, dx$ et 
on obtient 
$$\int (\cos x) ^{1234} \sin x \, d x = \int u^{1234} (-du) = - \frac 1{1235} u^{1235} + c=  -\frac 1{1235}(\cos x)^{1235}+c$$
Cette primitive est définie sur $\Rr$.

  \item $\int \frac 1{x\ln x} \, dx$

En posant le changement de variable $u=\ln x$ on a $x=\exp u$ et $du = \frac {dx}{x}$ on écrit :
$$\int \frac 1{x\ln x} \, dx = \int \frac 1{\ln x} \frac{dx}{x} = \int \frac 1 u du= \ln |u| + c = \ln \left| \ln x\right| +c$$
Cette primitive est définie sur $\left] 0,1\right[$ ou sur $\left] 1,+\infty \right[$ (la constante peut être différente pour chacun des intervalles).


  \item $\int \frac 1{3+\exp \left( -x\right)}dx$

Soit le changement de variable $u=\exp x$. Alors $x=\ln u$ et $du = \exp x \, dx$
ce qui s'écrit aussi $dx = \frac{du}{u}$.
$$\int \frac 1{3+\exp \left( -x\right) }dx = \int \frac{1}{3+\frac{1}{u}} \frac{du}{u}
= \int \frac{1}{3u+1} du = \frac 13 \ln |3u+1| + c = \frac13\ln \left( 3\exp
x+1\right) +c$$
Cette primitive est définie sur $\Rr$.

  \item $\int \frac{1}{\sqrt{4x-x^2}}dx$

Le changement de variable a pour but de se ramener à quelque chose de connu.
Ici nous avons une fraction avec une racine carrée au dénominateur et sous la racine un polynôme de degré $2$.
Ce que l'on sait intégrer c'est
$$\int \frac{1}{\sqrt{1-u^2}}du = \arcsin u + c,$$
car on connaît la dérivée de la fonction $\arcsin(t)$ c'est $\arcsin'(t)=\frac{1}{\sqrt{1-t^2}}$.
On va donc essayer de s'y ramener.
Essayons d'écrire ce qu'il y a sous la racine, $4x-x^2$ sous la forme $1-t^2$ :
$4x-x^2 =  4 - (x-2)^2 = 4 \bigg( 1 - \big(\frac 12 x - 1\big)^2\bigg)$.
Donc il est naturel d'essayer le changement de variable $u= \frac 12 x - 1$
pour lequel $4x-x^2=4(1-u^2)$ et $dx = 2du$. 

$$\int \frac{1}{\sqrt{4x-x^2}}dx = \int \frac{1}{\sqrt{4(1-u^2)}} 2du = \int \frac{du}{\sqrt{1-u^2}}
= \arcsin u + c = \arcsin \left( \frac 12x-1\right) +c$$
La fonction $\arcsin u$ est définie et dérivable pour $u\in]-1,1[$ alors cette primitive est 
définie sur $x \in \left] 0,4\right[$.
\end{enumerate}
\fincorrection
\nocorrection
\nocorrection
\nocorrection
\correction{000818}
\begin{enumerate}

\item Par IPP, $I_{n+2} = \frac{n+1}{n+2} I_n$.

\item $I_0 = \pi/2$ et $I_1 = 1$ et
\[ 
I_{2p} = \frac{ (2p-1) \times (2p-3) \times \cdots \times 1}{2p \times (2p-2) \times \cdots \times 2} I_0
= \frac{(2p)!}{2^{2p} (p!)^2} \frac{\pi}{2} ,\quad
I_{2p+1} = \frac{2p \times (2p-2) \times \cdots \times 2}{ (2p+1) \times (2p-1) \times \cdots \times 1} I_1
= \frac{2^{2p} (p!)^2}{(2p+1)!}
\]

\item En regardant l'intégrand.

\item D'après la question précédente, $0 < I_{n+2} \leq I_{n+1} \leq I_n$ donc
\[
\frac{n+1}{n+2} = \frac{I_{n+2}}{I_n} \leq \frac{I_{n+1}}{I_n} \leq 1
\]
par conséquent $\frac{I_{n+1}}{I_n} \xrightarrow[n\to \infty]{} 1$.

\item 
\[
(2p-1) I_{2p-1} I_{2p} = \frac{2p-1}{2 p} \frac{\pi}{2}
, \quad
2p I_{2p} I_{2p+1} = \frac{2p}{2p+1} \frac{\pi}{2}
\]
soit $n I_n I_{n+1} = \frac{n}{n+1} \frac{\pi}{2}$, ce qui peut aussi se démontrer par récurrence.

\item Comme $\frac{\pi}{2(n+1)} I_n I_{n+1} \sim I_n^2$ on en déduit que $I_n \sim \sqrt{\frac{\pi}{2n}}$.

\end{enumerate}
\fincorrection
\nocorrection
\nocorrection
\nocorrection
\nocorrection
\correction{006864}
\begin{enumerate}
  \item $\int x^2 \ln x \, dx$

Considérons l'intégration par parties avec $u=\ln x$ et $v'=x^2$.
On a donc $u'=\frac 1x$ et $v = \frac{x^3}3$.
Donc
\begin{align*}
\int  \ln x \times x^2\, dx 
  &= \int uv' = \big[ uv \big] - \int u'v \\
  &= \left[ \ln x \times \frac{x^3}3 \right] - \int  \frac 1x\times\frac{x^3}3  \, dx \\
  &= \left[ \ln x  \times \frac{x^3}3\right] - \int \frac{x^2}3 \, dx \\
  &= \frac{x^3}3 \ln x - \frac{x^3}9 + c \\
\end{align*}


  \item $\int x \arctan x \, dx$

Considérons l'intégration par parties avec $u=\arctan x$ et $v'=x$.
On a donc $u'=\frac 1{1+x^2}$ et $v = \frac{x^2}2$.
Donc
\begin{align*}
\int  \arctan x \times x \, dx 
&= \int uv' = \big[ uv \big] - \int u'v \\
&= \left[ \arctan x \times \frac{x^2}2  \right] - \int \frac 1{1+x^2} \times\frac{x^2}2  \, dx  \\
&= \left[  \arctan x \times \frac{x^2}2\right] - \frac12 \int \left( 1 -  \frac 1{1+x^2} \right)\, dx  \\
&= \frac{x^2}2  \arctan x -\frac 12 x + \frac 12 \arctan x + c \\
&= \frac 12 (1+x^2) \arctan x -\frac 12 x+ c \\
\end{align*}

  \item $\int \ln x \, dx$ puis $\int (\ln x)^2 \, dx$

Pour la primitive $\int \ln x \, dx$, regardons l'intégration par parties avec $u=\ln x$ et $v'=1$.
Donc $u' = \frac 1x$ et $v=x$.
\begin{align*}
\int \ln x \, dx
&= \int uv' = \big[ uv \big] - \int u'v \\
&= \left[ \ln x \times x \right] - \int \frac 1x \times x \, dx \\
&= \left[  \ln x \times x \right] - \int 1 \, dx \\
&= x\ln x - x + c \\
\end{align*}

\bigskip

Par la primitive $\int (\ln x)^2 \, dx$ soit l'intégration par parties définie par $u=(\ln x)^2$ et  $v'=1$.
Donc $u' = 2 \frac 1x \ln x$ et $v=x$.
\begin{align*}
\int (\ln x)^2 \, dx
&= \int uv' = \big[ uv \big] - \int u'v \\
&= \left[ x (\ln x)^2 \right] - 2 \int \ln x \, dx \\
&= x(\ln x)^2 -2 (x\ln x - x) + c \\
\end{align*}
Pour obtenir la dernière ligne on a utilisé la primitive calculée précédemment.

  \item Notons $I=\int \cos x\exp x \, dx$.

Regardons l'intégration par parties avec $u=\exp x$ et $v'=\cos x$.
Alors  $u' = \exp x$ et $v=\sin x$.

Donc

$$I = \int \cos x\exp x \,dx= \big[ \sin x \exp x \big] - \int \sin x \exp x\,dx$$

Si l'on note $J = \int \sin x \exp x\,dx$, alors on a obtenu
\begin{equation}
\label{eq:intI}
I = \big[ \sin x \exp x \big] - J  
\end{equation}
 
Pour calculer $J$ on refait une deuxième intégration par parties
avec  $u=\exp x$ et $v'=\sin x$.
Ce qui donne
$$J = \int \sin x \exp x\,dx = \big[ -\cos x \exp x \big] - \int -\cos x \exp x\,dx
= \big[ -\cos x \exp x \big] + I$$
Nous avons ainsi une deuxième équation :
\begin{equation}
\label{eq:intJ}
J = \big[ -\cos x \exp x \big] + I
\end{equation}

Repartons de l'équation (\ref{eq:intI}) dans laquelle on remplace $J$ par la formule obtenue dans l'équation 
(\ref{eq:intJ}).


$$I=\big[ \sin x \exp x \big] - J  = \big[ \sin x \exp x \big] - \big[ -\cos x \exp x \big] - I$$
D'où
$$2I = \big[ \sin x \exp x \big] + \big[ \cos x \exp x \big]$$

Ce qui nous permet de calculer notre intégrale :
$$I= \frac12 (\sin x + \cos x) \exp x + c.$$
\end{enumerate}
\fincorrection
\nocorrection
\nocorrection
\correction{002090}
a-$\int \sin ^8x\cos ^3xdx=\frac 19\sin ^9x-\frac 1{11}\sin ^{11}x+c$ sur $%
\R$.

b-$\int \cos ^4xdx=\frac 1{32}\sin 4x+\frac 14\sin 2x+\frac 38x+c$ sur $\R$.

c-$\int \cos ^{2003}x\sin xdx=-\frac 1{2004}\cos ^{2004}x+c$ sur $\R$.

d-$\int \frac 1{\sin x}dx=\frac 12\ln \left| \frac{1-\cos x}{1+\cos x}%
\right| +c=\ln \left| \tan \frac x2\right| +c$ sur $\left] k\pi ,\left(
k+1\right) \pi \right[ $ (changement de variable $u=\cos x$ ou $u=\tan \frac
x2$).

e-$\int \frac 1{\cos x}dx=\frac 12\ln \left| \frac{1+\sin x}{1-\sin x}%
\right| +c=\ln \left| \tan \left( \frac x2+\frac \pi 4\right) \right| +c$
sur $\left] -\frac \pi 2+k\pi ,\frac \pi 2+k\pi \right[ $ (changement de
variable $u=\sin x$ ou $u=\tan \frac x2$).

f-$\int \frac{3-\sin x}{2\cos x+3\tan x}dx=-\frac 15\ln \left( 2-\sin
x\right) +\frac 7{10}\ln \left| 1+2\sin x\right| +c$ sur $\R\setminus
\left\{ \frac{2\pi }3\left[ 2\pi \right] ,-\frac{2\pi }3\left[ 2\pi \right]
\right\} $ (changement de variable $u=\sin x$).

g-$\int \frac 1{7+\tan x}dx=\frac 7{50}x+\frac 1{50}\ln \left| \tan
x+7\right| +\frac 1{50}\ln \left| \cos x\right| +c$ sur $\R\setminus\left\{ \arctan \left( -7\right) +k\pi\,,\,
\frac{\pi}{2}+k\pi\,, \;k\in
Z\right\} $ (changement de variable $u=\tan x$).

h-$\int \frac {1}{2+\sin
x+\cos x}dx=\sqrt{2}\Arctan\left(\frac{1+\tan(x/2)}{\sqrt{2}}\right)+c$ sur
$\R\setminus\{k\pi\,,\;k\in\Z\}$ (changement de variable $u=\tan
(x/2)$).
\fincorrection
\correction{002096}
  \begin{enumerate}

  \item 
  \begin{enumerate}
     \item $$ I_{n+2}  = \int_0^{\frac \pi 2} \sin^{n+1} x \cdot \sin x \,  dx.$$
En posant $u(x) = \sin^{n+1} x$ et $v'(x) = \sin x$ et en intégrant par parties nous obtenons
\begin{align*}
I_{n+2} &= \bigg[ -\cos x \sin^{n+1}x \bigg]_0^{\frac \pi 2} \ \  + \ \ (n+1)\int_0^{\frac \pi 2} \cos^2x \sin^nx \, dx \\
 &= 0 \ \  + \ \ (n+1)\int_0^{\frac \pi 2} (1-\sin^2x)\sin^nx \, dx \\
 &= (n+1)I_n-(n+1)I_{n+2}.  \\
\end{align*}

Donc $(n+2)I_{n+2}=(n+1)I_n$.
Conclusion 
$$I_{n+2} = \frac{n+1}{n+2} I_n.$$

     \item Nous avons donc une formule de récurrence pour $I_n$ qui s'exprime en fonction de $I_{n-2}$
qui a son tour s'exprime en fonction de $I_{n-4}$, etc. On se ramène ainsi à l'intégrale de $I_0$ (si $n$ est pair) 
ou bien de $I_1$ (si $n$ est impair). Un petit calcul donne $I_0=\frac \pi 2$ et $I_1=1$.
Par récurrence nous avons donc pour $n$ pair :
$$I_n = \frac{1\cdot3 \cdots (n-1) }{2 \cdot 4 \cdots n} \frac \pi 2,$$
et pour $n$ impair :
$$I_n = \frac{2 \cdot 4 \cdots (n-1)}{1 \cdot 3 \cdots n}.$$

     \item Pour calculer $\int_{-1}^1\left( 1-x^2\right) ^n d x$ nous allons nous ramener à une intégrale de Wallis.
Avec le changement de variable $x=\cos u$, on montre assez facilement que :
\begin{align*}
 \int_{-1}^1\left( 1-x^2\right) ^n d x 
   &=  2\int_0^1\left( 1-x^2\right) ^n d x\\
   &=  2\int_{\frac \pi 2}^{0}  \left( 1-\cos ^2 u\right) ^n (- \sin u \, du)  \quad \text{ avec } x=\cos u \\
   &=  2\int_0^{\frac \pi 2}  \sin^{2n+1} u \, du  \\
   &= 2I_{2n+1}\\
\end{align*}

  \end{enumerate}

  \item
  \begin{enumerate}
     \item Sur $[0,\frac \pi 2]$ la fonction sinus est positive donc $I_n$ est positive.
De plus, sur ce même intervalle $\sin x \leqslant 1$ donc  $(\sin x)^{n+1}  \le (\sin x)^n$.
Cela implique 
$$I_{n+1}=\int_0^{\frac \pi 2}(\sin x)^{n+1} d x \le \int_0^{\frac \pi 2}(\sin x)^n d x = I_n.$$

     \item Comme $(I_n)$ est décroissante alors $I_{n+2} \leqslant I_{n+1} \leqslant I_n$, en 
divisant le tout par $I_n>0$ nous obtenons $\frac{I_{n+2}}{I_n} \leqslant \frac{I_{n+1}}{I_n} \leqslant 1$.
Mais nous avons déjà calculé $\frac{I_{n+2}}{I_n} = \frac{n+1}{n+2}$ 
qui tend vers $1$ quand $n$ tend vers l'infini. Donc $ \frac{I_{n+1}}{I_n}$ tend vers $+1$ donc $I_{n} \sim I_{n+1}$.

  \end{enumerate}

  \item 
  \begin{enumerate}
     \item Nous allons calculer $I_n\cdot I_{n+1}$.
Supposons par exemple que $n$ est pair, alors par les formules obtenues précédemment :
$$I_n \times I_{n+1} = 
\frac{1\cdot3 \cdots (n-1) }{2 \cdot 4 \cdots n} \frac \pi 2 \times 
\frac{2 \cdot 4 \cdots n}{1 \cdot 3 \cdots (n+1)} = \frac \pi 2 \times \frac{1}{n+1}.$$

Si $n$ est impair nous obtenons la même fraction. 
On en déduit que pour tout $n$ : $I_n \cdot I_{n+1} = \frac{\pi}{2(n+1)}$.

     \item Maintenant 
$$I_n^2 = I_n \cdot I_n \sim I_n \cdot I_{n+1} = \frac \pi {2(n+1)} \sim  \frac \pi {2n},$$
donc $$I_n \sim \sqrt{\frac \pi {2n}}.$$

     \item $$\frac{1\cdot3 \cdots (2n+1) }{2\cdot 4 \cdots (2n) } 
= I_{2n} \cdot (2n+1) \cdot \frac 2 \pi  \sim \sqrt{\frac \pi {4n}} \cdot (2n+1) \cdot \frac 2 \pi 
 \sim 2\sqrt {\frac n \pi}.$$

  \end{enumerate}

\end{enumerate}
\fincorrection
\nocorrection
\correction{005475}
\begin{enumerate} 
\item  $I_0=\int_{0}^{\pi/4}dx=\frac{\pi}{4}$ et $I_{1}=\int_{0}^{\pi/4}\frac{\sin x}{\cos x}\;dx=\left[-\ln|\cos x|\right]_{0}^{\pi/4}=\frac{\ln2}{2}$.

Soit $n\in\Nn$.

$$I_n+I_{n+2}=\int_{0}^{\pi/4}(\tan^nx+\tan^{n+2}x)\;dx=\int_{0}^{\pi/4}\tan^nx(1+\tan^2x)\;dx=\left[\frac{\tan^{n+1}x}{n+1}\right]_{0}^{\pi/4}=\frac{1}{n+1}.$$

Soit $n\in\Nn^*$.

\begin{align*}\ensuremath
\sum_{k=1}^{n}\frac{(-1)^{k-1}}{2k-1}&=\sum_{k=1}^{n}(-1)^{k-1}(I_{2k-2}+I_{2k})=\sum_{k=1}^{n}(-1)^{k-1}I_{2k-2}
+\sum_{k=1}^{n}(-1)^{k-1}I_{2k}\\
 &=\sum_{k=0}^{n-1}(-1)^{k}I_{2k}-\sum_{k=1}^{n}(-1)^{k}I_{2k}=I_0-(-1)^nI_{2n}.
\end{align*}

Ainsi, $\forall n\in\Nn^*,\;I_{2n}=(-1)^n\left(\frac{\pi}{4}-\sum_{k=1}^{n}\frac{(-1)^{k-1}}{2k-1}\right)$.

De même, $\sum_{k=1}^{n}\frac{(-1)^{k-1}}{2k}=I_1-(-1)^nI_{2n+1}$ et donc, $\forall n\in\Nn^*,\;I_{2n+1}=\frac{(-1)^n}{2}\left(\ln2-\sum_{k=1}^{n}\frac{(-1)^{k-1}}{k}\right)$.

\item  Soient $\varepsilon\in]0,\frac{\pi}{2}[$ et $n\in\Nn^*$.

$$0\leq I_n=\int_{0}^{\pi/4-\varepsilon/2}\tan^nx\;dx+\int_{\pi/4-\varepsilon/2}^{\pi/4}\tan^nx\;dx\leq\frac{\pi}{4}\tan^n(\frac{\pi}{4}-\frac{\varepsilon}{2})+\frac{\varepsilon}{2}.$$

Maintenant, $0<\tan(\frac{\pi}{4}-\frac{\varepsilon}{2})<1$ et donc $\lim_{n\rightarrow +\infty}\tan^n(\frac{\pi}{4}-\frac{\varepsilon}{2})=0$. Par suite, il existe $n_0\in\Nn$ tel que, pour $n\geq n_0$, $0\leq\tan^n(\frac{\pi}{4}-\frac{\varepsilon}{2})<\frac{\varepsilon}{2}$. Pour $n\geq n_0$, on a alors $0\leq I_n<\varepsilon$.

Ainsi, $I_n$ tend vers $0$ quand $n$ tend vers $+\infty$. On en déduit immédiatement que $u_n$ tend vers $\ln2$ et $v_n$ tend vers $\frac{\pi}{4}$.

\end{enumerate}
\fincorrection
\correction{000824}
 Résultats valables sur chaque intervalle 
du domaine de définition.
\begin{enumerate}


\item  ${1 \over x^2+a^2}$ est un élément simple.
Primitives~: ${1 \over a} \arctan({x \over a})+k$.

\item  ${1 \over{(1+x^2)}^2}$ est un élément simple.
Primitives~: ${1 \over2} \arctan x + {x \over2(1+x^2)}+k$.

\item  ${x^3 \over x^2-4}=x+{2 \over x-2}+{2 \over x+2}$.
Primitives~: ${x^2 \over2}+ \ln(x^2-4)^2 +k$.

\item  ${4x \over{(x-2)}^2}={4 \over x-2}+{8 \over{(x-2)}^2}$.
Primitives~: $4 \ln\vert x-2\vert-{8 \over x-2}+k$.

\item  ${1 \over x^2+x+1}$ est un élément simple.
Primitives~: ${2 \over\sqrt3} \arctan{(2x+1) \over\sqrt3}+k$.

\item  $ {1 \over{(t^2+2t-1)}^2} = {1 \over8{(t+1+ \sqrt2)}^2 }+
{ \sqrt2 \over16(t+1+\sqrt2)} + {1 \over8{(t+1- \sqrt2)}^2}+
{ - \sqrt2 \over16(t+1- \sqrt2)}$.\newline
Primitives~: $-{t+1 \over4(t^2+2t-1)} +{\sqrt2 \over16} \ln
\left\vert{t+1+ \sqrt2 \over t+1- \sqrt2 } \right\vert+ k$.

\item  $ {3t+1 \over{(t^2-2t+10)}^2}$ est un élément simple.\newline
Primitives~: $-{3 \over2(t^2-2t+10)} +{2(t-1) \over9(t^2-2t+10)}
+{2 \over27} \arctan({t-1 \over3}) +k $.

\item  $ {3t+1 \over t^2-2t+10}$ est un élément simple.
Primitives~: ${3 \over2} \ln(t^2-2t+10) +{4 \over3} \arctan({t-1
\over3})+k$.

\item  ${1 \over t^3+1}={1 \over3(t+1)}-{t-2 \over3(t^2-t+1)}$.
Primitives~: $ {1 \over3}\ln\vert t+1 \vert-{1 \over6} \ln(
t^2-t+1) + {1 \over\sqrt3} \arctan({2t-1 \over\sqrt3}) +k$.

\item  $ {x^3+2 \over{(x+1)}^2}= x-2+{3 \over x+1}+{1 \over{(x+1)}^2}$.
Primitives~: ${x^2 \over2}-2x+3 \ln\vert x+1 \vert-{1 \over x+1} +k$.

\item  ${x+1 \over x{(x-2)}^2}= {1 \over4x}- {1 \over4(x-2)}
 +{3 \over2{(x-2)}^2}$.
Primitives~: $ {1 \over4} \ln\vert x \vert-{1 \over4} \ln\vert x-2 \vert
-{3 \over2(x-2)} +k$.

\item  ${(x^2-1)(x^3+3) \over2x+2x^2}= {1 \over2} (x^3-x^2+3) -{3 \over2x}
$.
Primitives~: ${x^4 \over8}-{x^3 \over6}+{3x \over2}-{3 \over2}\ln\vert x
\vert+k$.

\item  ${x^2 \over{(x^2+3)}^3(x+1)}={1 \over4^3(x+1)}+
{1-x \over4^3 (x^2+3)} +
{1-x \over4^2 {(x^2+3)}^2}
-{3(1-x) \over4{(x^2+3)}^3}$.\newline
Primitives~: $-{x+3 \over4^2{(x^2+3)}^2}-{2x-3 \over3.2^5(x^2+3)}
-{1 \over2^7}\ln(x^2+3) - {1 \over3 \sqrt3  \, 2^6} \arctan({x \over\sqrt3})
+{1 \over4^3}\ln\vert x+1 \vert+k$.

\item  ${x^7+x^3-4x-1 \over x{(x^2+1)}^2}= x^2-2-{1 \over x} +{x+4 \over x^2+1}+
{x-6 \over{(x^2+1)}^2}$.\newline
Primitives~: ${x^3 \over3}-2x -\ln\vert x \vert+{1 \over2}\ln
(1+x^2) +\arctan x -{6x+1 \over2(x^2+1)} +k$.

\item  ${3x^4-9x^3+12x^2-11x+7\over(x-1)^3(x^2+1)}={1\over
(x-1)^3}-{2\over(x-1)^2}+{3\over x-1}-{1\over x^2+1}$.\newline
Primitives~: $-{1/2\over(x-1)^2}+{2\over x-1}+3\ln\vert
x-1\vert-\arctan x +k$.
\end{enumerate}
\fincorrection
\correction{000825}
\begin{enumerate}
\item  $\frac{1}{x^2+2}$ est un élément simple.
$\int_0^1\frac{dx}{x^2+2} = \frac{1}{\sqrt{2}} \arctan \frac{1}{\sqrt{2}}$.

\item  Décomposition~: $\frac{1}{1-x^2} = \frac{1/2}{x+1} - \frac{1/2}{x-1}$. 
Intégrale~: $\int_{-1/2}^{1/2} \frac{dx}{1-x^2} = \ln 3$.

\item  Pas besoin de décomposer la fraction rationnelle, car $2x+1$ est 
la dérivée de $x^2+x-3$~! $\int_2^3 \frac{2x+1}{x^2+x-3}\,dx = \ln 3$.

\item  On peut évidemment décomposer la fraction rationnelle en 
éléments simples~: 
$\frac{x}{x^4+16} = \frac{\sqrt{2}/8}{x^2-2x\sqrt{2}+4} - 
\frac{\sqrt{2}/8}{x^2+2x\sqrt{2}+4}$, 
mais il est bien plus simple de 
faire le changement de variables $x^2=u$. 
Alors 
$\int_0^2 \frac{x\,dx}{x^4+16} = 
\frac{1}{2}\int_0^4 \frac{du}{u^2+16} = \frac{\pi}{32}$.

\item  La décomposition de $\frac{x^4+6x^3-5x^2+3x-7}{(x-4)^3}$ est
$x+18+\frac{163}{x-4}+\frac{507}{(x-4)^2}+\frac{565}{(x-4)^3}$~; les 
primitives sont 
$\frac{x^2}{2}+18x-\frac{1014x-3491}{2(x-4)^2}+163\ln\vert x-4\vert+C$.
Enfin,
$\int_0^3 \frac{x^4+6x^3-5x^2+3x-7}{(x-4)^3}\,dx =
\frac{5565}{32}-326\ln 2$.

\item  Décomposition~: 
$\frac{1}{x^3-7x+6} = \frac{1}{20(x+3)} -\frac{1}{4(x-1)} 
+\frac{1}{5(x-2)}$. Primitives~:
$\frac{1}{20}\ln\Bigl\vert\frac{(x-2)^4(x+3)}{(x-1)^5}\Bigr\vert+C$, d'où
 $\int_{-2}^0 \frac{dx}{x^3-7x+6} = 
\frac{1}{10}\ln(27/4)$.


\item  Décomposition~: 
$\frac{2x^4+3x^3+5x^2+17x+30}{x^3+8} = 
2x+3+\frac{2}{x+2}+\frac{3x-1}{x^2-2x+4}$. 
Les primitives sont~:
$x^2+3x+\ln(x+2)^2 + \frac{3}{2}\ln(x^2-2x+4) + 
\frac{2}{\sqrt{3}}\arctan\frac{x-1}{\sqrt{3}}+C$.
Intégrale~: 
$\int_{-1}^1 \frac{2x^4+3x^3+5x^2+17x+30}{x^3+8}\,dx = 
6+\frac{7\ln 3-3\ln 7}{2} + 
\frac{2}{\sqrt{3}}\arctan\frac{2}{\sqrt{3}}$.


\item  Décomposition~: 
$\frac{4x^2}{x^4-1} = 
\frac{2}{x^2+1} - \frac{1}{x+1} +\frac{1}{x-1}$.
Primitives~:
$\ln\Bigl\vert \frac{x-1}{x+1}\Bigr\vert +2\arctan x +C$, d'où 
$\int_2^3 \frac{4x^2}{x^4-1}\,dx =
\ln\frac{3}{2}+2\arctan \frac{1}{7}$.


\item  La décomposition est $\frac{x^3+2x+1}{x^3-3x+2} = 
1+\frac{4/3}{(x-1)^2} + \frac{11/9}{x-1} - \frac{11/9}{x+2}$. On 
trouve alors\linebreak %%%%
$\int_{-1}^0 \frac{x^3+2x+1}{x^3-3x+2}\,dx = 
\frac{5}{3}-\frac{22}{9}\ln 2$.

\item  La décomposition de 
$\frac{2x^8+5x^6-12x^5+30x^4+36x^2+24}{x^4(x^2+2)^3}$ est 
$\frac{3}{x^4} + \frac{2}{x^2+2} - \frac{6}{(x^2+2)^2} - 
\frac{12x-16}{(x^2+2)^3}$~; les primitives sont 
$-\frac{1}{x^3} + \frac{2x+3}{(x^2+2)^2} + \sqrt{2} 
\arctan\frac{x}{\sqrt{2}}+C$. 
Enfin 
$\int_{1}^2 \frac{2x^8+5x^6-12x^5+30x^4+36x^2+24}
    {x^4(x^2+2)^3}\,dx =
    \frac{37}{72} + 2\sqrt{2}\arctan\sqrt{2}-\frac{\pi}{\sqrt{2}}$.


\item  Décomposition de la fraction rationnelle~: 
$\frac{-2x^2+6x+7}{x^4+5x^2+4} = 
\frac{2x+3}{x^2+1}-\frac{2x+5}{x^2+4}$. Primitives~:
$\ln\Bigl\vert \frac{x^2+1}{x^2+4}\Bigr\vert+3\arctan 
x-\frac{5}{2}\arctan\frac{x}{2}+C$. 
Alors
$\int_0^a \frac{-2x^2+6x+7}{x^4+5x^2+4}\,dx = 
\ln\Bigl\vert \frac{a^2+1}{a^2+4}\Bigr\vert+3\arctan 
a-\frac{5}{2}\arctan\frac{a}{2} + 2\ln 2$. Enfin
$\lim_{a\to +\infty} \int_0^a \frac{-2x^2+6x+7}{x^4+5x^2+4}\,dx =
\frac{\pi}{4}+2\ln 2$.



\item  Pour factoriser le dénominateur, penser à 
faire $x^4+1=x^4+2x^2+1-2x^2$~; on trouve alors
$\frac{1}{x^4+1} = 
\frac{(x\sqrt{2}+2)/4}{x^2+x\sqrt {2}+1} - 
\frac{(x\sqrt{2}-2)/4}{x^2-x\sqrt{2}+1}$. Les primitives s'écrivent 
$$\textstyle \frac{1}{4\sqrt{2}} \ln \frac{x^2+x\sqrt{2}+1}{x^2-x\sqrt{2}+1} +
\frac{1}{2\sqrt{2}}\Bigl(\arctan(x\sqrt{2}+1)+\arctan(x\sqrt{2}-1)\Bigr)+C$$
ce qui donne 
$\int_0^2 \frac{dx}{x^4+1} = 
\frac{1}{4\sqrt{2}}\ln\frac{33+20\sqrt{2}}{17}
+\frac{1}{2\sqrt{2}}\Bigl(\pi-\arctan \frac{2\sqrt{2}}{3}\Bigr)$.
\end{enumerate}
\fincorrection
\nocorrection
\nocorrection
\correction{002097}
  \begin{enumerate}
  \item Pour $x>0$ on a $\frac{x^n}{1+x} \leqslant x^n$,
donc 
$$I_n  \leqslant \int_0^1 x^n dx = \left[ \frac{1}{n+1} x^{n+1} \right]_0^1=\frac{1}{n+1}.$$
Donc $I_n \to 0$ lorsque $n\to +\infty$.
  \item $I_n+I_{n+1}=\int_0^1 x^n \frac {1+x}{1+x} dx =  \int_0^1 x^n dx=\frac{1}{n+1}$.
  \item Soit $S_n = 1-\frac 12 + \frac13-\frac 14 +\cdots \pm \frac 1n = \sum_{k = 1}^n \frac{
 (-1)^{k + 1}}k$.
Par la question pr\'ec\'edente nous avons 
$S_n = (I_0+I_1)-(I_1+I_2)+(I_2+I_3)- \cdots \pm (I_{n-1}+I_n)$.
Mais d'autre part cette somme \'etant t\'elescopique cela conduit à $S_n = I_0 \pm I_n$.
 Alors la limite de $S_n$ et donc de $\sum_{k = 1}^n \frac{
 (-1)^{k + 1}}k$ (quand $n\to +\infty$) est $I_0$ car $I_n \to 0$.
Un petit calcul montre que $I_0 = \int_0^1 \frac {dx}{1+x} = \ln 2$.
Donc la somme altern\'ee des inverses des entiers converge vers $\ln 2$.
  \end{enumerate}
\fincorrection
\nocorrection
\nocorrection
\nocorrection
\correction{002089}
$\int \frac{\sin x}{\sin x+\cos x}dx=\frac 12\left( x-\ln \left| \cos x+\sin
x\right| \right) +c$ sur $\R$

$\int \frac{\cos x}{\sin x+\cos x}dx=\frac 12\left( x+\ln \left| \cos x+\sin
x\right| \right) +c$ sur $\R$ (en calculant la somme et la
diff\'erence).
\fincorrection
\correction{002095}
\begin{enumerate}
  \item Notons $I = \int_0^{\frac \pi 2}\frac 1{1+\sin x} dx$.
Le changement de variable $t = \tan \frac x2$ transforme toute fraction rationnelle 
de sinus et cosinus en une fraction rationnelle en $t$ (que l'on sait résoudre !).

En posant $t=\tan \frac{x}{2}$ on a $x=\arctan \frac t2$ ainsi que les formules suivantes : 
$$\cos x = \frac {1-t^2}{1+t^2}, \quad \sin x = \frac{2t}{1+t^2}, 
\quad \tan x = \frac{2t}{1-t^2}, \quad dx=\dfrac{2dt}{1+t^2}.$$

Ici, on a seulement à remplacer $\sin x$.
Comme $x$ varie de $x=0$ à $x=\frac\pi 2$ alors $t=\tan \frac{x}{2}$ varie de $t=0$ à $t=1$.

\begin{align*}  
  I &= \int_0^{\frac \pi 2}\frac 1{1+\sin x} dx 
     = \int_0^1  \frac 1 {1+ \frac{2t}{1+t^2}}  \dfrac{2dt}{1+t^2} \\ 
    &= \int_0^1 \frac{2}{1+t^2 + 2t} dt 
    = \int_0^1 \frac{2}{(1+t)^2} dt \\
    &= \left[ \frac{-2}{1+t} \right]_0^1 
    = 1
\end{align*}

  \item Notons $J = \int_0^{\frac \pi 2}\frac{\sin x}{1+\sin x} dx$.
Alors 
$$I+J =  \int_0^{\frac \pi 2}\frac 1{1+\sin x} dx + \int_0^{\frac \pi 2}\frac{\sin x}{1+\sin x} dx
= \int_0^{\frac \pi 2} \frac{1+\sin x}{1+\sin x} dx = \int_0^{\frac \pi 2} 1 \, dx = \big[ x \big]_0^{\frac \pi 2} = \frac \pi 2.$$
Donc $J = \frac \pi 2 - I= \frac \pi 2 - 1$.
\end{enumerate}

\fincorrection
\nocorrection
\nocorrection
\nocorrection
\nocorrection
\correction{000832}
\begin{enumerate}


\item   Changement de variable $u= \sin^2x$ (ou d'abord $u=\sin x$)~;
$e^{\sin^2x}+C$.

\item   Deux méthodes~: changement de variable $u=\sin t$ (ou $u=\sinh t$),
ou linéarisation.\newline
$ {1\over15}(15\sin t -10\sin^3 t+3\sin^5 t)+C$ ou
$ {1\over80}\sin5t+{5\over48}\sin3t+{5\over8}\sin t +C$~;\newline
$ \sinh t + {1\over3}\sinh^3 t +C$ ou ${1\over12}\sinh3t +{3\over4}\sinh t+C$~;\newline
${1\over32}(\sin4t+8\sin2t+12t) +C$~;
${1 \over32}(\sinh4t-8\sinh2t+12t)+C$.

\item   Intégrations par parties~: $(x^3-3x^2+6x-6)e^x+C$.

\item   Intégration par parties~: $x\ln x-x+C$~;
${x^2\over2}\ln x -{x^2\over4}+C$~; $x\arcsin x+\sqrt{1-x^2}+C$.

\item   Intégrations par parties~: ${1\over2}(\sinh t\sin t -\cosh t\cos t)+C$.

\item   Changement de variable $t=\tan{x\over2}$~;
$\ln\bigl\vert\tan{x\over2}\bigr\vert+C$ sur chaque intervalle\dots

\item   Changement de variable $x=a\sin u$~;
${a^2\over2}\arcsin{x\over a} + {x\over2}\sqrt{a^2-x^2} + C$.

\item   Changement de variable $u=e^x$~; ${2\over3}\sqrt{e^x+1}(e^x-2)+C$.

\item   Intégrations par parties~:
${1\over a^2+b^2}e^{ax}(a\cos bx + b\sin bx)+C$~;\newline
${1\over a^2+b^2}e^{ax}(-b\cos bx + a\sin bx)+C$.

\item   Changement de variable $t=\sqrt{x\over1-x}$~;
$2\sqrt{x\over1-x}-2\arctan\sqrt{x\over1-x}+C$.

\item   Changement de variable $t=\arcsin x$~;
${1\over2}(\arcsin x - x\sqrt{1-x^2}) + C$.

\item   Changements de variable $u=\tan{x\over2}$, $t=1+u$~;
$\arctan(\tan{x\over2}+1)+C$ sur chaque intervalle\dots{} Mais, au fait, ne
cherchait-on pas une primitive {\it sur} $\bf R$~?

\item   Changement de variable $x^3=u^2$~;
${2\over3}\arcsin\sqrt{x^3\over a^3}+C$.

\item  Multiplier et diviser par $\cosh x-\sinh x$, ou passer en $e^x$~;
${x\over2} +{\sinh2x\over4}-{\cosh2x\over4}+C$ ou
${x\over2}-{e^{-2x}\over4}+C$.
\end{enumerate}
\fincorrection
\nocorrection
\nocorrection
\nocorrection
\nocorrection
\nocorrection
\nocorrection
\nocorrection
\correction{002088}
a-$\int \arctan xdx=x\arctan x-\frac 12\ln \left( 1+x^2\right) +c$ sur $\Bbb{%
R}$ (int\'egration par parties)

b-$\int \tan ^2xdx=\tan x-x+c$ sur $\left] -\frac \pi 2+k\pi ,\frac \pi
2+k\pi \right[ $

c-$\int \frac 1{x\ln x}dx=\ln \left| \ln x\right| +c$ sur $\left] 0,1\right[
\cup \left] 1,+\infty \right[ $ (changement de variable : $u=\ln x$)

d-$\int \frac x{\sqrt{x+1}}dx=\frac 23\left( x-2\right) \left( x+1\right)
^{\frac 12}+c$ sur $\left] -1,+\infty \right[ $ (changement de variable : $u=\sqrt{x+1}$ ou 
int\'egration par parties)

e-$\int \arcsin xdx=x\arcsin x+\sqrt{1-x^2}+c$ sur $\left] -1,1\right[ $
(int\'egration par parties)

f-$\int \frac 1{3+\exp \left( -x\right) }dx=\frac 13\ln \left( 3\exp
x+1\right) +c$ sur $\Bbb{R}$ (changement de variable : $u=\exp x$)

g-$\int \frac{-1}{\sqrt{4x-x^2}}dx=\arccos \left( \frac 12x-1\right) +c$ sur 
$\left] 0,4\right[ $ (changement de variable : $u=\frac 12x-1$)

h-$\int \frac 1{x\sqrt{1-\ln ^2x}}dx=\arcsin \left( \ln x\right) +c$ sur $%
\left] \frac 1e,e\right[ $ (changement de variable : $u=lnx$)

i-$\int \frac 1{\sqrt{1+\exp x}}dx=x-2\ln \left( 1+\sqrt{\exp x+1}\right) +c$
sur $\Bbb{R}$ (changement de variable : $u=\sqrt{\exp x+1}$)

j-$\int \frac{x-1}{x^2+x+1}dx=\frac 12\ln \left( x^2+x+1\right) -\sqrt{3}%
\arctan \left( \frac 2{\sqrt{3}}\left( x+\frac 12\right) \right) +c$ sur $%
\Bbb{R}$

k-$\int \frac{x+2}{x^2-3x-4}dx=-\frac 15\ln \left| x+1\right| +\frac 65\ln
\left| x-4\right| +c$ sur $\R\setminus\left\{
-1,4\right\} $ (d\'ecomposition en \'el\'ements simples)

l-$\int \cos x\exp xdx=\frac 12\left( \cos x+\sin x\right) \exp x+c$ sur $%
\R$ (deux int\'egrations par parties)
\fincorrection
\correction{002094}
a-$\int_{0^{}}^1\frac{\arctan x}{1+x^2}dx=\frac{\pi ^2}{32}$ (changement de
variables ou int\'egration par parties).

b-$\int_{\frac 12}^2\left( 1+\frac 1{x^2}\right) \arctan xdx=\frac{3\pi }4$
(changement de variables $u=\frac 1x$ et $\arctan x+\arctan \frac 1x=\frac \pi 2$).

c-$\int_0^{\frac \pi 2}x\sin xdx=1$ (int\'egration par parties).

d-$\int_{-1}^1\left( \arccos x\right) ^2dx=\pi ^2+4$ (2 int\'egrations par
parties).

e-$\int_0^1\frac 1{\left( 1+x^2\right) ^2}dx=\frac \pi 8+\frac 14$
(changement de variables ou int\'egration par parties).

f-$\int_0^{\sqrt{3}}\frac{x^2}{\sqrt{4-x^2}}dx=\frac{2\pi }3-\frac{\sqrt{3}}%
2 $ (changement de variables $u=\arcsin \frac x2$).

g-$\int_1^2x^2\ln xdx=\frac 83\ln 2-\frac 79$ (int\'egration par parties).

h-$\int_{-1}^1\frac 1{x^2+4x+7}dx=\frac \pi {6\sqrt{3}}$ (changement de
variables $u=\frac{x+2}{\sqrt{3}}$).

i-$\int_0^1\frac{3x+1}{\left( x+1\right) ^2}dx=3\ln 2-1$ (d\'ecomposition en
\'el\'ements simples).
\fincorrection
\nocorrection
\nocorrection
\nocorrection
\nocorrection
\nocorrection
\nocorrection
\correction{005450}
Pour $t$ réel, posons $g(t)=\frac{t^2}{\sqrt{1+t^8}}$ puis, pour $x$ réel, $G(x)=\int_{1}^{x}g(t)\;dt$. Puisque $g$ est définie et continue sur $\Rr$, $G$ est définie sur $\Rr$ et de classe $C^1$ et $G'=g$ ($G$ est la primitive de $g$ sur $\Rr$ qui s'annule en $1$). Plus précisément, $g$ est de classe $C^\infty$ sur $\Rr$ et donc $G$ est de classe $C^\infty$ sur $\Rr$.

Finalement, $f$ est définie et de classe $C^\infty$ sur $]-\infty,1[\cup]1,+\infty[$.

\textbf{Etude en 1.}

Pour $x\neq1$, 

$$f(x)=\frac{G(x)}{x-1}=\frac{G(1)+G'(1)(x-1)+\frac{G''(1)}{2}(x-1)^2+o((x-1)^2)}{x-1}=g(1)+g'(1)(x-1)+o((x-1)).$$

Donc, $f$ admet en $1$ un développement limité d'ordre $1$. Par suite, $f$ se prolonge par continuité en $1$ en posant $f(1)=g(1)=\frac{1}{\sqrt{2}}$ puis le prolongement est dérivable en $1$ et $f'(1)=\frac{1}{2}g'(1)$. Or, pour tout réel $x$, $g'(x)=2x\frac{1}{\sqrt{1+x^8}}+x2.(-\frac{4x^7}{(1+x^8)\sqrt{1+x^8}})=2x\frac{1-x^8}{(1+x^8)\sqrt{1+x^8}}$ et $g'(1)=0$. Donc, $f'(1)=0$.

\textbf{Dérivée. Variations}

Pour $x\neq1$, $f'(x)=\frac{G'(x)(x-1)-G(x)}{(x-1)^2}$.

$f'(x)$ est du signe de $h(x)=G'(x)(x-1)-G(x)$ dont la dérivée est $h'(x)=G''(x)(x-1)+G'(x)-G'(x)=(x-1)g'(x)$.
$h'$ est du signe de $2x(1-x^8)(x-1)$ ou encore du signe de $-2x(1+x)$. $h$ est donc décroissante sur $]-\infty,-1]$ et sur $[0,+\infty[$ et croissante sur $[-1,0]$.

Maintenant, quand $x$ tend vers $+\infty$ (ou $-\infty$), $G'(x)(x-1)=g(x)(x-1)\sim x\frac{1}{x^2}=\frac{1}{x}$ et donc $G'(x)(x-1)$ tend vers $0$. Ensuite, pour $x\geq1$

$$0\leq G(x)\leq\int_{1}^{x}\frac{t^2}{\sqrt{t^8}}\;dt=1-\frac{1}{x}\leq1,$$

et $G$ est bornée au voisinage de $+\infty$ (ou de $-\infty$). Comme $G$ est croissante sur $\Rr$, $G$ a une limite réelle en $+\infty$ et en $-\infty$. Cette limite est strictement positive en $+\infty$ et strictement négative en $-\infty$. Par suite, $h$ a une limite strictement positive en $-\infty$ et une limite strictement négative en $+\infty$.
Sur $[0,+\infty[$, $h$ est décroissante et s'annule en $1$. Donc, $h$ est positive sur $[0,1]$ et négative sur $[1,+\infty[$.

Ensuite, 

$$h(-1)=\int_{-1}^{1}\frac{t^2}{\sqrt{1+t^8}}\;dt-\sqrt{2}=2\int_{0}^{1}\frac{t^2}{\sqrt{1+t^8}}\;dt-\sqrt{2}<2\int_{0}^{1}\frac{1}{\sqrt{2}}\;dt-\sqrt{2}=0,$$

et $h(-1)<0$. $h$ s'annule donc, une et une seule fois sur $]-\infty,-1[$ en un certain réel $\alpha$ et une et une seule fois sur $]-1,0[$ en un certain réel $\beta$. De plus, $h$ est strictement positive sur $]-\infty,\alpha[$, strictement négative sur $]\alpha,\beta[$, strictement positive sur $]\beta,1[$ et strictement négative sur $]1,+\infty[$.

$f$ est strictement croissante sur $]-\infty,\alpha]$, strictement décroissante sur $[\alpha,\beta]$, strictement croissante sur $[\beta,1]$ et strictement décroissante sur $[1,+\infty[$.

\textbf{Etude en l'infini.}

En $+\infty$ ou $-\infty$, $G$ a une limite réelle et donc $f$ tend vers $0$.
\fincorrection
\correction{005451}
\begin{enumerate}
\item  La fonction $t\mapsto e^{t^2}$ est de classe $C^\infty$ sur $\Rr$. Donc, la fonction $x\mapsto\int_{0}^{x}e^{t^2}\;dt$ est de classe $C^\infty$ sur $\Rr$ et il en est de même de $f$.

La fonction $t\mapsto e^{t^2}$ est paire et donc la fonction $x\mapsto\int_{0}^{x}e^{t^2}\;dt$ est impaire. Comme la fonction $x\mapsto e^{-x^2}$ est paire, $f$ est impaire.

\item  Pour $x$ réel, $f'(x)=-2xe^{-x^2}\int_{0}^{x}e^{t^2}\;dt+e^{-x^2}e^{x^2}=-2xf(x)+1$.
\item  Pour $x\geq1$, une intégration par parties fournit~:

$$\int_{1}^{x}e^{t^2}\;dt=\int_{1}^{x}\frac{1}{2t}.2te^{t^2}\;dt=\left[\frac{1}{2t}e^{t^2}\right]_{1}^{x}+\frac{1}{2}\int_{1}^{x}\frac{e^{t^2}}{t^2}\;dt=\frac{e^{x^2}}{2x}-\frac{e}{2}+\frac{1}{2}\int_{1}^{x}\frac{e^{t^2}}{t^2}\;dt,$$

et donc,
 
\begin{align*}\ensuremath
|1-2xf(x)|&=\left|1-2xe^{-x^2}\int_{1}^{x}e^{t^2}\;dt-2xe^{-x^2}\int_{0}^{1}e^{t^2}\;dt\right|\\
 &\leq xe^{-x^2}\int_{1}^{x}\frac{e^{t^2}}{t^2}\;dt+exe^{-x^2}+2xe^{-x^2}\int_{0}^{1}e^{t^2}\;dt.
\end{align*}

Les deux derniers termes tendent vers $0$ quand $x$ tend vers $+\infty$. Il reste le premier.

Pour $x\geq2$, 

\begin{align*}\ensuremath
0\leq xe^{-x^2}\int_{1}^{x}\frac{e^{t^2}}{t^2}\;dt&=xe^{-x^2}\int_{1}^{x-1}\frac{e^{t^2}}{t^2}\;dt+
xe^{-x^2}\int_{x-1}^{x}\frac{e^{t^2}}{t^2}\;dt\\
 &\leq x(x-1)e^{-x^2}\frac{e^{(x-1)^2}}{1^2}+xe^{-x^2}e^{x^2}\int_{x-1}^{x}\frac{1}{t^2}\;dt\\
 &=x(x-1)e^{-2x+1}+x\left(\frac{1}{x-1}-\frac{1}{x}\right)=x(x-1)e^{-2x+1}+\frac{1}{x-1}.
\end{align*}

Cette dernière expression tend vers $0$ quand $x$ tend vers $+\infty$. On en déduit que $xe^{-x^2}\int_{1}^{x}\frac{e^{t^2}}{t^2}\;dt$ tend vers $0$ quand $x$ tend vers $+\infty$. Finalement, $1-2xf(x)$ tend vers $0$ quand $x$ tend vers $+\infty$, ou encore, $f(x)\sim\frac{1}{2x}$.

\item  Pour $x>0$, $g(x)=\frac{e^{x^2}}{2x}(1-2xf(x))=\frac{e^{x^2}}{2x}-\int_{0}^{x}e^{t^2}\;dt$ puis,

$$g'(x)=e^{x^2}-\frac{e^{x^2}}{2x^2}-e^{x^2}=-\frac{e^{x^2}}{2x^2}<0.$$

$g$ est donc strictement décroissante sur $]0,+\infty[$ et donc, $g$ s'annule au plus une fois sur $]0,+\infty[$. Ensuite, $f'(1)=1-2f(1)=1-2e^{-1}\int_{0}^{1}e^{t^2}\;dt$. Or, la méthode des rectangles fournit $\int_{0}^{1}e^{t^2}\;dt =1,44... >1,35...=\frac{e}{2}$, et donc $f'(1)<0$ puis $g(1)<0$. Enfin, comme en $0^+$, $g(x)\sim\frac{1}{2x}f'(0)=\frac{1}{2x}$, $g(0^+)=+\infty$.

Donc, $g$ s'annule exactement une fois sur $]0,+\infty[$ en un certain réel $x_0$ de $]0,1[$.

\item  $g$ est strictement positive sur $]0,x_0[$ et strictement négative sur $]x_0,+\infty[$. Il en de même de $f'$. $f$ est ainsi strictement croissante sur $[0,x_0]$ et strictement décroissante sur $[x_0,+\infty[$.
\end{enumerate}

\fincorrection
\correction{005465}
\begin{enumerate}
\item  $f$ est continue sur $\Rr^*$ en tant que quotient de fonctions continues sur $\Rr^*$ dont le dénominateur ne s'annule pas sur $\Rr^*$. D'autre part, quand $t$ tend vers $0$, $f(t)\sim\frac{t^2}{t}=t$ et $\lim_{t\rightarrow 0,\;t\neq0}f(t)=0=f(0)$. Ainsi, $f$ est continue en $0$ et donc sur $\Rr$.

\item  $f$ est continue sur $\Rr$ et donc $F$ est définie et de classe $C^1$ sur $\Rr$. De plus, $F'=f$ est positive sur $[0,+\infty[$, de sorte que $F$ est croissante sur $[0,+\infty[$. On en déduit que $F$ admet en $+\infty$ une limite dans $]-\infty,+\infty]$.

Vérifions alors que $F$ est majorée sur $\Rr$. On contate que $t^2.\frac{t^2}{e^t-1}$ tend vers $0$ quand $t$ tend vers $+\infty$, d'après un théorème de croissances comparées. Par suite, il existe un réel $A$ tel que pour $t\geq A$, 
$0\leq t^2.\frac{t^2}{e^t-1}\leq1$ ou encore $0\leq f(t)\leq\frac{1}{t^2}$. Pour $x\geq A$, on a alors

\begin{align*}\ensuremath
F(x)&=\int_{0}^{A}f(t)\;dt+\int_{A}^{x}\frac{t^2}{e^t-1}\;dt\leq\int_{0}^{A}f(t)\;dt+\int_{A}^{x}\frac{1}{t^2}\;dt\\
 &=\int_{0}^{A}f(t)\;dt+\frac{1}{A}-\frac{1}{x}\leq\int_{0}^{A}f(t)\;dt+\frac{1}{A}.
\end{align*}

$F$ est croissante et majorée et donc a une limite réelle $\ell$ quand $n$ tend vers $+\infty$.

Soit $n\in\Nn^*$. Pour $t\in]0,+\infty[$,

\begin{align*}
f(t)&=t^2e^{-t}\frac{1}{1-e^{-t}}=t^2e^{-t}(\sum_{k=0}^{n-1}(e^{-t})^k+\frac{(e^{-t})^{n}}{1-e^{-t}})\\
 &=\sum_{k=0}^{n-1}t^2e^{-(k+1)t}+\frac{t^2e^{-t}}{1-e^{-t}}e^{-nt}=\sum_{k=1}^{n}t^2e^{-kt}+f_n(t)\;(*),
\end{align*}

où $f_n(t)=\frac{t^2e^{-t}}{1-e^{-t}}e^{-nt}$ pour $t>0$. En posant de plus $f_n(0)=0$, d'une part, $f_n$ est continue sur $[0,+\infty[$ et d'autre part, l'égalité $(*)$ reste vraie quand $t=0$. En intégrant, on obtient

$$\forall x\in[0,+\infty[,\;\forall n\in\Nn^*,\;F(x)=\sum_{k=1}^{n}\int_{0}^{x}t^2e^{-kt}\;dt+\int_{0}^{x}f_n(t)\;dt\;(**).$$

Soient alors $k\in\Nn^*$ et $x\in[0,+\infty[$. Deux intégrations par parties fournissent~:

\begin{align*}\ensuremath
\int_{0}^{x}t^2e^{-kt}\;dt&=\left[-\frac{1}{k}t^2e^{-kt}\right]_{0}^{x}+\frac{2}{k}\int_{0}^{x}te^{-kt}\;dt
=-\frac{1}{k}x^2e^{-kx}+\frac{2}{k}(\left[-\frac{1}{k}te^{-kt}\right]_{0}^{x}+\frac{1}{k}\int_{0}^{x}e^{-kt}\;dt)\\
 &=-\frac{1}{k}x^2e^{-kx}-\frac{2}{k^2}xe^{-kx}-\frac{2}{k^3}e^{-kx}+\frac{2}{k^3}.
\end{align*}

Puisque $k>0$, quand $x$ tend vers $+\infty$, on obtient $\lim_{x\rightarrow +\infty}\int_{0}^{x}t^2e^{-kt}\;dt=\frac{2}{k^3}$. On fait alors tendre $x$ vers $+\infty$ dans $(**)$ et on obtient

$$\forall n\in\Nn^*,\;\ell-2\sum_{k=1}^{n}\frac{1}{k^3}=\lim_{x\rightarrow +\infty}\int_{0}^{x}f_n(t)\;dt\;(***).$$

Vérifions enfin que $\lim_{n\rightarrow +\infty}\left(\lim_{x\rightarrow +\infty}\int_{0}^{x}f_n(t)\;dt\right)=0$. La fonction $t\mapsto\frac{t^2e^{-t}}{1-e^{-t}}$ est continue sur $]0,+\infty[$, se prolonge par continuité en $0$ et a une limite réelle en $+\infty$. On en déduit que cette fonction est bornée sur $]0,+\infty[$. Soit $M$ un majorant de cette fonction sur $]0,+\infty[$. Pour $x\in[0,+\infty[$ et $n\in\Nn^*$, on a alors

$$0\leq\int_{0}^{x}f_n(t)\;dt\leq M\int_{0}^{x}e^{-nt}\;dt=\frac{M}{n}(1-e^{-nx}).$$

A $n\in\Nn^*$ fixé, on passe à la limite quand $n$ tend vers $+\infty$ et on obtient

$$0\leq\lim_{x\rightarrow +\infty}\int_{0}^{x}f_n(t)\;dt\leq\frac{M}{n},$$

puis on passe à la limite quand $n$ tend vers $+\infty$ et on obtient

$$\lim_{n\rightarrow +\infty}\left(\lim_{x\rightarrow +\infty}\int_{0}^{x}f_n(t)\;dt\right)=0.$$

Par passage à la limite quand $x$ tend vers $+\infty$ puis quand $n$ tend vers $+\infty$ dans $(***)$, on obtient enfin

$$\lim_{x\rightarrow +\infty}\int_{0}^{x}\frac{t^2}{e^t-1}\;dt=2\lim_{n\rightarrow +\infty}\sum_{k=1}^{n}\frac{1}{k^3}.$$
 
\end{enumerate}
\fincorrection
\correction{005466}
\begin{enumerate}
\item  $I$ est l'un des deux intervalles $]-\infty,-1[$ ou $]-1,+\infty[$. $f$ est continue sur $I$ et admet donc des primitives sur $I$.

$$\frac{1}{X^3+1}=\frac{1}{(X+1)(X+j)(X+j^2)}=\frac{a}{X+1}+\frac{b}{X+j}+\frac{\overline{b}}{X+j^2},$$

où $a=\frac{1}{3(-1)^2}=\frac{1}{3}$ et $b=\frac{1}{3(-j)^2}=\frac{j}{3}$. Par suite,

\begin{align*}\ensuremath
\frac{1}{X^3+1}&=\frac{1}{3}(\frac{1}{X+1}+\frac{j}{X+j}+\frac{j^2}{X+j^2})
=\frac{1}{3}(\frac{1}{X+1}+\frac{-X+2}{X^2-X+1})
=\frac{1}{3}(\frac{1}{X+1}-\frac{1}{2}\frac{2X-1}{X^2-X+1}+\frac{3}{2}\frac{1}{X^2-X+1})\\
 &=\frac{1}{3}(\frac{1}{X+1}-\frac{1}{2}\frac{2X-1}{X^2-X+1}
 +\frac{3}{2}\frac{1}{(X-\frac{1}{2})^2+(\frac{\sqrt{3}}{2})^2}).
\end{align*}

Mais alors,

\begin{align*}\ensuremath
\int_{}^{}\frac{1}{x^3+1}\;dx&=\frac{1}{3}(\ln|x+1|-\frac{1}{2}\ln(x^2-x+1)+\frac{3}{2}\frac{2}{\sqrt{3}}\Arctan
\frac{x-\frac{1}{2}}{\frac{\sqrt{3}}{2}})
=\frac{1}{6}\ln\frac{(x-1)^2}{x^2-x+1}+\frac{1}{\sqrt{3}}\Arctan\frac{2x-1}{\sqrt{3}}+C.
\end{align*}

\item  $I$ est l'un des deux intervalles $]-\infty,-1[$ ou $]1,+\infty[$. Sur $I$, $\int_{}^{}\frac{x^2}{x^3+1}\;dx=\frac{1}{3}\ln(x^3+1)+C$.

\item  $X^3-X^2-X+1=X^2(X-1)-(X-1)=(X^2-1)(X-1)=(X-1)^2(X+1)$. Donc, la décomposition en éléments simples de $f=\frac{X^5}{X^3-X^2-X+1}$ est de la forme $aX^2+bX+c+\frac{d_1}{X-1}+\frac{d_2}{(X-1)^2}+\frac{e}{X+1}$.

Détermination de $a$, $b$ et $c$. La division euclidienne de $X^5$ par $X^3-X^2-X+1$ s'écrit $X^5=(X^2+X+2)(X^3-X^2-X+1)+2X^2+X-2$. On a donc $a=1$, $b=1$ et $c=2$.

$e=\lim_{x\rightarrow -1}(x+1)f(x)=\frac{(-1)^5}{(-1-1)^2}=-\frac{1}{4}$. Puis, $d_2=\lim_{x\rightarrow 1}(x-1)^2f(x)=\frac{1^5}{1+1}=\frac{1}{2}$. Enfin, $x=0$ fournit $0=c-d_1+d_2+e$ et donc, $d_1=-2-\frac{1}{2}+\frac{1}{4}=-\frac{9}{4}$. Finalement,

$$\frac{X^5}{X^3-X^2-X+1}=X^2+X+2-\frac{9}{4}\frac{1}{X-1}+\frac{1}{2}\frac{1}{(X-1)^2}-\frac{1}{4}\frac{1}{X+1},$$

et donc, $I$ désignant l'un des trois intervalles $]-\infty,-1[$, $]-1,1[$ ou $]1,+\infty[$, on a sur $I$

$$\int_{}^{}\frac{x^5}{x^3-x^2-x+1}\;dx=\frac{x^3}{3}+\frac{x^2}{2}+2x-\frac{1}{2(x-1)}-\frac{1}{4}\ln|x+1|+C.$$
\item  Sur $\Rr$,

\begin{align*}\ensuremath
\int_{}^{}\frac{1-x}{(x^2+x+1)^5}\;dx&=-\frac{1}{2}\int_{}^{}\frac{2x+1}{(x^2+x+1)^5}\;dx+\frac{3}{2}\int_{}^{}\frac{1}{(x^2+x+1)^5}\;dx
=\frac{1}{8(x^2+x+1)^4}+\frac{3}{2}\int_{}^{}\frac{1}{((x+\frac{1}{2})^2+\frac{3}{4})^5}\;dx\\
 &=\frac{1}{8(x^2+x+1)^4}+\frac{3}{2}\int_{}^{}\frac{1}{((\frac{\sqrt{3}}{2}u)^2+\frac{3}{4})^5}\;\frac{\sqrt{3}}{2}\;du\;(\mbox{en posant}\;x+\frac{1}{2}=\frac{u\sqrt{3}}{2})\\
 &=\frac{1}{8(x^2+x+1)^4}+\frac{2^8\sqrt{3}}{3^4}\int_{}^{}\frac{1}{(u^2+1)^5}\;du.
\end{align*}

Pour $n\in\Nn^*$, posons alors $I_n=\int_{}^{}\frac{du}{(u^2+1)^n}$. Une intégration par parties fournit

$$I_n=\frac{u}{(u^2+1)^n}+2n\int_{}^{}\frac{u^2+1-1}{(u^2+1)^{n+1}}\;du=\frac{u}{(u^2+1)^n}+2n(I_n-I_{n+1}),$$

et donc, $I_{n+1}=\frac{1}{2n}\left(\frac{u}{(u^2+1)^n}+(2n-1)I_n\right)$. Mais alors,

\begin{align*}\ensuremath
I_5&=\frac{1}{8}\frac{u}{(u^2+1)^4}+\frac{7}{8}I_4=\frac{1}{8}\frac{u}{(u^2+1)^4}+\frac{7}{8.6}\frac{u}{(u^2+1)^3}+\frac{7.5}{8.6}I_3\\
 &=\frac{1}{8}\frac{u}{(u^2+1)^4}+\frac{7}{8.6}\frac{u}{(u^2+1)^3}+\frac{7.5}{8.6.4}\frac{u}{(u^2+1)^2}+\frac{7.5.3}{8.6.4}I_2\\
 &=\frac{1}{8}\frac{u}{(u^2+1)^4}+\frac{7}{8.6}\frac{u}{(u^2+1)^3}+\frac{7.5}{8.6.4}\frac{u}{(u^2+1)^2}+\frac{7.5.3}{8.6.4.2}\frac{u}{u^2+1}+\frac{7.5.3.1}{8.6.4.2}I_1\\
 &=\frac{1}{8}\frac{u}{(u^2+1)^4}+\frac{7}{8.6}\frac{u}{(u^2+1)^3}+\frac{7.5}{8.6.4}\frac{u}{(u^2+1)^2}+\frac{7.5.3}{8.6.4.2}\frac{u}{u^2+1}+\frac{7.5.3.1}{8.6.4.2}\Arctan u+C.
\end{align*}

Maintenant,

$$u^2+1=(\frac{2}{\sqrt{3}}(x+\frac{1}{2}))^2+1=\frac{4}{3}x^2+\frac{4}{3}x+\frac{1}{3}+1=\frac{4}{3}(x^2+x+1).$$

Par suite,

\begin{align*}\ensuremath
\frac{2^8\sqrt{3}}{3^4}\int_{}^{}\frac{1}{(u^2+1)^5}\;du&=\frac{2^8\sqrt{3}}{3^4}
\left(
\frac{1}{8}\frac{3^4}{4^4}\frac{\frac{2}{\sqrt{3}}(x+\frac{1}{2})}{(x^2+x+1)^4}
+\frac{7}{8.6}\frac{3^3}{4^3}\frac{\frac{2}{\sqrt{3}}(x+\frac{1}{2})}{(x^2+x+1)^3}
+\frac{7.5}{8.6.4}\frac{3^2}{4^2}\frac{\frac{2}{\sqrt{3}}(x+\frac{1}{2})}{(x^2+x+1)^2}
\right.
\\
 &\;\left.+\frac{7.5.3}{8.6.4.2}\frac{3}{4}\frac{\frac{2}{\sqrt{3}}(x+\frac{1}{2})}{x^2+x+1}
 +\frac{7.5.3.1}{8.6.4.2}\Arctan\frac{2x+1}{\sqrt{3}}+C
 \right).\\
 &=\frac{1}{8}\frac{2x+1}{(x^2+x+1)^4}+\frac{7}{36}\frac{2x+1}{(x^2+x+1)^3}+\frac{35}{108}\frac{2x+1}{(x^2+x+1)^2}+\frac{35}{54}\frac{2x+1}{x^2+x+1}\\
 &\;+\frac{70\sqrt{3}}{81}\Arctan\frac{2x+1}{\sqrt{3}}+C,
\end{align*}

(il reste encore à réduire au même dénominateur).

\item  On pose $u=x^2$ et donc $du=2xdx$

\begin{align*}\ensuremath
\int_{}^{}\frac{1}{x(x^2+1)^2}\;dx&=\int_{}^{}\frac{x}{x^2(x^2+1)^2}\;dx=\frac{1}{2}\int_{}^{}\frac{du}{u(u+1)^2}
=\frac{1}{2}\int_{}^{}(\frac{1}{u}-\frac{1}{u+1}-\frac{1}{(u+1)^2})\;du\\
 &=\frac{1}{2}(\ln|u|-\ln|u+1|+\frac{1}{u+1})+C\\
 &=\frac{1}{2}(\ln\frac{x^2}{x^2+1}+\frac{1}{x^2+1})+C.
\end{align*}
\item  $\int_{}^{}\frac{x^2+x}{x^6+1}\;dx=\int_{}^{}\frac{x^2}{x^6+1}\;dx+\int_{}^{}\frac{x}{x^6+1}\;dx$.

Ensuite, en posant $u=x^3$ et donc $du=3x^2\;dx$,

$$\int_{}^{}\frac{x^2}{x^6+1}\;dx=\frac{1}{3}\int_{}^{}\frac{1}{u^2+1}\;du=\frac{1}{3}\Arctan u+C=\frac{1}{3}\Arctan(x^3)+C,$$

et en posant $u=x^2$ et donc $du=2x\;dx$,

\begin{align*}\ensuremath
\int_{}^{}\frac{x}{x^6+1}\;dx&=\frac{1}{2}\int_{}^{}\frac{1}{u^3+1}\;du
=\frac{1}{6}\ln\frac{(u-1)^2}{u^2-u+1}+\frac{1}{\sqrt{3}}\Arctan\frac{2u-1}{\sqrt{3}}+C\;(\mbox{voir}\;1))\\
 &=\frac{1}{6}\ln\frac{(x^2-1)^2}{x^4-x^2+1}+\frac{1}{\sqrt{3}}\Arctan\frac{2x^2-1}{\sqrt{3}}+C
\end{align*}

Finalement,

$$\int_{}^{}\frac{x^2+x}{x^6+1}\;dx=\frac{1}{3}\Arctan(x^3)+\frac{1}{6}\ln\frac{(x^2-1)^2}{x^4-x^2+1}+\frac{1}{\sqrt{3}}\Arctan\frac{2x^2-1}{\sqrt{3}}+C.$$

\item  $\frac{1}{X^4+1}=\sum_{k=0}^{3}\frac{\lambda_k}{X-z_k}$ où $z_k=e^{i(\frac{\pi}{4}+k\frac{\pi}{2})}$. De plus,
$\lambda_k=\frac{1}{4z_k^3}=\frac{z_k}{4z_k^4}=-\frac{z_k}{4}$. Ainsi,

\begin{align*}\ensuremath
\frac{1}{X^4+1}&=-\frac{1}{4}\left(\frac{e^{i\pi/4}}{X-e^{i\pi/4}}+\frac{e^{-i\pi/4}}{X-e^{-i\pi/4}}
+\frac{-e^{i\pi/4}}{X+e^{i\pi/4}}+\frac{-e^{-i\pi/4}}{X+e^{-i\pi/4}}\right)\\
 &=-\frac{1}{4}\left(\frac{\sqrt{2}X-2}{X^2-\sqrt{2}X+1}-\frac{\sqrt{2}X+2}{X^2+\sqrt{2}X+1}\right).
\end{align*}

Mais,

$$\frac{\sqrt{2}X-2}{X^2-\sqrt{2}X+1}=\frac{1}{\sqrt{2}}\frac{2X-\sqrt{2}}{X^2-\sqrt{2}X+1}-\frac{1}{(X-\frac{1}{\sqrt{2}})^2+(\frac{1}{\sqrt{2}})^2},$$

et donc,

$$\int_{}^{}\frac{\sqrt{2}x-2}{x^2-\sqrt{2}x+1}\;dx=\frac{1}{\sqrt{2}}\ln(x^2-\sqrt{2}x+1)-\sqrt{2}\Arctan(\sqrt{2}x-1)+C,$$

et de même,

$$\int_{}^{}\frac{\sqrt{2}x+2}{x^2+\sqrt{2}x+1}\;dx=\frac{1}{\sqrt{2}}\ln(x^2+\sqrt{2}x+1)+\sqrt{2}\Arctan(\sqrt{2}x+1)+C.$$

Finalement,

$$\int_{}^{}\frac{1}{x^4+1}\;dx=\frac{1}{\sqrt{2}}\ln\frac{x^2-\sqrt{2}x+1}{x^2+\sqrt{2}x+1}-\sqrt{2}(\Arctan(\sqrt{2}x-1)+\Arctan(\sqrt{2}x+1))+C.$$
\item  Une intégration par parties fournit

\begin{align*}\ensuremath
\int_{}^{}\frac{1}{x^4+1}\;dx&=\frac{x}{x^4+1}+\int_{}^{}\frac{4x^4}{(x^4+1)^2}\;dx=\frac{x}{x^4+1}+4\int_{}^{}\frac{x^4+1-1}{(x^4+1)^2}\;dx\\
 &=\frac{x}{x^4+1}+4\int_{}^{}\frac{1}{x^4+1}\;dx-4\int_{}^{}\frac{1}{(x^4+1)^2}\;dx\\
\end{align*}

Et donc,

$$\int_{}^{}\frac{1}{(x^4+1)^2}\;dx=\frac{1}{4}(\frac{x}{x^4+1}+3\int_{}^{}\frac{1}{x^4+1}\;dx)=...$$
\item  Posons $R=\frac{1}{X^8+X^4+1}$.

\begin{align*}\ensuremath
X^8+X^4+1&=\frac{X^{12}-1}{X^4-1}=\frac{\prod_{k=0}^{11}(X-e^{2ik\pi/12})}{(X-1)(X+1)(X-i)(X+i)}\\
 &=(X-e^{i\pi/6})(X-e^{-i\pi/6})(X+e^{i\pi/6})(X+e^{-i\pi/6})(X-j)(X-j^2)(X+j)(X+j^2).
\end{align*}

$R$ est réelle et paire. Donc,

$$R=\frac{a}{X-j}+\frac{\overline{a}}{X-j^2}-\frac{a}{X+j}-\frac{\overline{a}}{X+j^2}+\frac{b}{X-e^{i\pi/6}}
+\frac{\overline{b}}{X-e^{-i\pi/6}}-\frac{b}{X+e^{i\pi/6}}-\frac{\overline{b}}{X+e^{-i\pi/6}}.$$

$a=\frac{1}{8j^7+4j^3}=\frac{1}{4(2j+1)}=\frac{2j^2+1}{4(2j+1)(2j^2+1)}=\frac{-1-2j}{12}$ et donc,

$$\frac{a}{X-j}+\frac{\overline{a}}{X-j^2}=\frac{1}{12}(\frac{-1-2j}{X-j}+\frac{-1-2j^2}{X-j^2})
=\frac{1}{4}\frac{1}{X^2+X+1}=\frac{1}{4}\frac{1}{(X+\frac{1}{2})^2+(\frac{\sqrt{3}}{2})^2},$$

 et par parité,
 
$$\frac{a}{X-j}+\frac{\overline{a}}{X-j^2}-\frac{a}{X+j}-\frac{\overline{a}}{X+j^2}=
\frac{1}{4}(\frac{1}{(X+\frac{1}{2})^2+(\frac{\sqrt{3}}{2})^2}+\frac{1}{(X-\frac{1}{2})^2+(\frac{\sqrt{3}}{2})^2}).$$

Ensuite, $b=\frac{1}{8e^{7i\pi/6}+4e^{3i\pi/6}}=\frac{1}{4e^{i\pi/6}(-2-j^2)}=\frac{e^{-i\pi/6}}{4(-1+j)}
=\frac{e^{-i\pi/6}(-1+j^2)}{12}=\frac{e^{-i\pi/6}(-2-j)}{12}=\frac{-2e^{-i\pi/6}-i}{12}$, et donc,

\begin{align*}\ensuremath
\frac{b}{X-e^{i\pi/6}}
+\frac{\overline{b}}{X-e^{-i\pi/6}}&=\frac{1}{12}(\frac{-2e^{-i\pi/6}-i}{X-e^{i\pi/6}}+\frac{-2e^{i\pi/6}+i}{X-e^{-i\pi/6}})=\frac{1}{12}\frac{-2\sqrt{3}X+3}{X^2-\sqrt{3}X+1}=-\frac{1}{4\sqrt{3}}\frac{2X-\sqrt{3}}{X^2-\sqrt{3}X+1}.
\end{align*}

Par parité,

$$\frac{b}{X-e^{i\pi/6}}
+\frac{\overline{b}}{X-e^{-i\pi/6}}-\frac{b}{X+e^{i\pi/6}}-\frac{\overline{b}}{X+e^{-i\pi/6}}=
-\frac{1}{4\sqrt{3}}\frac{2X-\sqrt{3}}{X^2-\sqrt{3}X+1}+\frac{1}{4\sqrt{3}}\frac{2X+\sqrt{3}}{X^2+\sqrt{3}X+1}.$$

Finalement,

$$\int_{}^{}\frac{1}{x^8+x^4+1}=\frac{1}{2\sqrt{3}}(\Arctan\frac{2x-1}{\sqrt{3}}+\Arctan\frac{2x+1}{\sqrt{3}})+\frac{1}{4\sqrt{3}}\ln\frac{x^2+\sqrt{3}x+1}{x^2-\sqrt{3}x+1}+C.$$

\item  En posant $u=x^2$ et donc $du=2x\;dx$, on obtient $\int_{}^{}\frac{x}{(x^4+1)^3}\;dx=\frac{1}{2}\int_{}^{}\frac{1}{(u^2+1)^3}$.

Pour $n\geq1$, posons $I_n=\int_{}^{}\frac{1}{(u^2+1)^n}\;du$. Une intégration par parties fournit~:

\begin{align*}\ensuremath
I_n=\frac{u}{(u^2+1)^n}+\int_{}^{}\frac{u.(-n)(2u)}{(u^2+1)^{n+1}}\;du=\frac{u}{(u^2+1)^n}+2n\int_{}^{}\frac{u^2+1-1}{(u^2+1)^{n+1}}\;du\\
 &=\frac{u}{(u^2+1)^n}+2n(I_n-I_{n+1}),
\end{align*}

et donc, $\forall n\geq1,\;I_{n+1}=\frac{1}{2n}(\frac{u}{(u^2+1)^n}+(2n-1)I_n)$.

On en déduit que

$$I_3=\frac{1}{4}(\frac{u}{(u^2+1)^2}+3I_2)=\frac{u}{4(u^2+1)^2}+\frac{3}{8(u^2+1)}+\frac{3}{8}\Arctan u+C,$$

et finalement que

$$\int_{}^{}\frac{x}{(x^4+1)^3}\;dx=\frac{1}{16}(\frac{2x^2}{(x^4+1)^2}+\frac{3}{x^4+1}+3\Arctan(x^2))+C.$$

\item  
\begin{align*}\ensuremath
(X+1)^7-X^7-1&=7X^6+21X^5+35X^4+35X^3+21X^2+7X=7X(X^5+3X^4+5X^3+5X^2+3X+1)\\
 &=7X(X+1)(X^4+2X^3+3X^2+2X+1)=7X(X+1)(X^2+X+1)^2.
\end{align*}

Par suite,

$$\frac{7}{(X+1)^7-X^7-1}=\frac{1}{X(X+1)(X-j)^2(X-j^2)^2}=\frac{a}{X}+\frac{b}{X+1}+\frac{c_1}{X-j}
+\frac{c_2}{(X-j)^2}+\frac{\overline{c_1}}{X-j^2}+\frac{\overline{c_2}}{(X-j^2)^2}.$$

$a=\lim_{x\rightarrow 0}xR(x)=1$, $b=\lim_{x\rightarrow -1}(x+1)R(x)=-1$, et

$c_2=\lim_{x\rightarrow j}(x-j)^2R(x)=\frac{1}{j(j+1)(j-j^2)^2}=-\frac{1}{j^2(1-2j+j^2)}=\frac{1}{3}$. Puis, 

$$\frac{c_2}{(X-j)^2}+\frac{\overline{c_2}}{(X-j^2)^2}=\frac{1}{3}(\frac{(X-j^2)^2+(X-j)^2}{(X^2+X+1)^2}
=\frac{2X^2+2X-1}{3(X^2+X+1)^2},$$

et

\begin{align*}\ensuremath
R-(\frac{c_2}{(X-j)^2}+\frac{\overline{c_2}}{(X-j^2)^2})
&=\frac{1}{X(X+1)(X^2+X+1)^2}-\frac{2X^2+2X-1}{3(X^2+X+1)^2}
=\frac{3-X(X+1)(2X^2+2X-1)}{3X(X+1)(X^2+X+1)^2}\\
 &=\frac{-2X(X+1)(X^2+X+1)+3+3X(X+1)}{3X(X+1)(X^2+X+1)^2}
=\frac{-2X^2-2X+3}{3X(X+1)(X^2+X+1)}.
\end{align*}

Puis, $c_2=\frac{-2j^2-2j+3}{3j(j+1)(j-j^2)}=-\frac{5}{j-j^2}=\frac{5(j-j^2)}{(j-j^2)(j^2-j)}=\frac{5(j-j^2)}{3}$.

Ainsi,

\begin{align*}\ensuremath
\frac{1}{(X+1)^7-X^7-1}&=\frac{1}{7}
(\frac{1}{X}-\frac{1}{X+1}+\frac{1}{3}(\frac{5(j-j^2)}{X-j}+\frac{5(j^2-j)}{X-j^2}+\frac{1}{(X-j)^2}+\frac{1}{(X-j^2)^2}))\\
 &=\frac{1}{7}
(\frac{1}{X}-\frac{1}{X+1}-\frac{5}{X^2+X+1}+\frac{1}{3}(\frac{1}{(X-j)^2}+\frac{1}{(X-j^2)^2}))\\
 &=\frac{1}{7}
(\frac{1}{X}-\frac{1}{X+1}-\frac{5}{(X+\frac{1}{2})^2+(\frac{\sqrt{3}}{2})^2}+\frac{1}{3}(\frac{1}{(X-j)^2}+\frac{1}{(X-j^2)^2})).
\end{align*}

Finalement,

\begin{align*}\ensuremath
\int_{}^{}\frac{1}{(x+1)^7-x^7-1}\;dx&=\frac{1}{7}\left(\ln\left|\frac{x}{x+1}\right|-\frac{10}{\sqrt{3}}
\Arctan\frac{2x+1}{\sqrt{3}}-\frac{1}{3}(\frac{1}{x-j}+\frac{1}{x-j^2})\right)+C\\
 &=\frac{1}{7}\left(\ln\left|\frac{x}{x+1}\right|-\frac{10}{\sqrt{3}}
\Arctan\frac{2x+1}{\sqrt{3}}-\frac{2x+1}{3(x^2+x+1)}\right)+C.
\end{align*}

\end{enumerate}
\fincorrection
\correction{005467}
\begin{enumerate}

\item  On pose $t=\tan\frac{x}{2}$ et donc $dx=\frac{2dt}{1+t^2}$.

\begin{align*}\ensuremath
\int_{}^{}\frac{1}{\cos x}\;dx&=\int_{}^{}\frac{1+t^2}{1-t^2}\;\frac{2dt}{1+t^2}=\int_{}^{}2\frac{1}{1-t^2}\;dt
=\ln\left|\frac{1+t}{1-t}\right|+C=\ln\left|\frac{\tan\frac{\pi}{4}+\tan\frac{x}{2}}{1-\tan\frac{\pi}{4}\tan\frac{x}{2}}\right|+C\\
 &=\ln|\tan(\frac{x}{2}+\frac{\pi}{4})|+C.
\end{align*}

ou bien

$$\int_{}^{}\frac{1}{\cos x}\;dx=\int_{}^{}\frac{\cos x}{1-\sin^2x}\;dx=\ln\left|\frac{1+\sin x}{1-\sin x}\right|+C...$$

ou bien, en posant $u=x+\frac{\pi}{2}$, (voir 2))

$$\int_{}^{}\frac{1}{\cos x}\;dx=\int_{}^{}\frac{1}{\cos(u-\frac{\pi}{2})}\;du=\int_{}^{}\frac{1}{\sin u}\;du=
\ln|\tan\frac{u}{2}|+C=\ln|\tan(\frac{x}{2}+\frac{\pi}{4})|+C.$$

Ensuite, en posant $t=e^x$ et donc $dx=\frac{dt}{t}$,

$$\int_{}^{}\frac{1}{\ch x}\;dx
=\int_{}^{}\frac{2}{t+\frac{1}{t}}\frac{dt}{t}=2\int_{}^{}\frac{1}{1+t^2}\;dt=2\Arctan(e^x)+C,$$

ou bien

$$\int_{}^{}\frac{1}{\ch x}\;dx
=\int_{}^{}\frac{\ch x}{\sh^2x+1}\;dx=\Arctan(\sh x)+C.$$
\item  En posant $t=\tan\frac{x}{2}$,

$$\int_{}^{}\frac{1}{\sin x}\;dx=\int_{}^{}\frac{1+t^2}{2t}\frac{2dt}{1+t^2}=\int_{}^{}\frac{1}{t}\;dt=\ln|t|+C=\ln|\tan\frac{x}{2}|+C.$$

\item  $\int_{}^{}\frac{dx}{\tan x}=\int_{}^{}\frac{\cos x}{\sin x}\;dx=\ln|\sin x|+C$ et $\int_{}^{}\frac{1}{\tanh x}=\ln|\sh x|+C$.

\item  $\int_{}^{}\frac{\sin^2(x/2)}{x-\sin x}\;dx=\frac{1}{2}\int_{}^{}\frac{1-\cos x}{x-\sin x}\;dx=\frac{1}{2}\ln|x-\sin x|+C$.

\item  $\frac{1}{2+\sin^2x}\;dx=\frac{1}{\frac{2}{\cos^2x}+\tan^2x}\frac{dx}{\cos^2x}=\frac{1}{2+3\tan^2x}d(\tan x)$, et en posant $u=\tan x$,

$$\int_{}^{}\frac{1}{2+\sin^2x}\;dx=\int_{}^{}\frac{1}{2+3u^2}\;du=\frac{1}{3}\sqrt{\frac{3}{2}}\Arctan(\sqrt{\frac{3}{2}}u)+C=\frac{1}{\sqrt{6}}\Arctan(\sqrt{\frac{3}{2}}\tan x)+C.$$

\item  Posons $I=\int_{}^{}\frac{\cos x}{\cos x+\sin x}\;dx$ et $J=\int_{}^{}\frac{\sin x}{\cos x+\sin x}\;dx$. Alors, $I+J=\int_{}^{}dx=x+C$ et $I-J=\int_{}^{}\frac{-\sin x+\cos x}{\cos x+\sin x}\;dx=\ln|\cos x+\sin x|+C$. En additionnant ces deux égalités, on obtient~:

$$I=\int_{}^{}\frac{\cos x}{\cos x+\sin x}\;dx=\frac{1}{2}(x+\ln|\cos x+\sin x|)+C.$$

ou bien, en posant $u=x-\frac{\pi}{4}$,

\begin{align*}\ensuremath
I&=\int_{}^{}\frac{\cos x}{\cos x+\sin x}\;dx=\int_{}^{}\frac{\cos x}{\sqrt{2}\cos(x-\frac{\pi}{4})}\;dx
=\int_{}^{}\frac{\cos(u+\frac{\pi}{4})}{\sqrt{2}\cos u}\;du=\frac{1}{2}\int_{}^{}(1-\frac{\sin u}{\cos u})\;du
=\frac{1}{2}(u+\ln|\cos u|)+C\\
 &=\frac{1}{2}(x-\frac{\pi}{4}+\ln|\frac{1}{\sqrt{2}}(\cos x+\sin x)|)+C=\frac{1}{2}(x+\ln|\cos x+\sin x|)+C.
\end{align*}

\item  

$$\frac{\cos(3x)}{\sin x+\sin(3x)}\;dx=\frac{4\cos^3x-3\cos x}{4\sin x-4\sin^3x}=\frac{1}{4}\frac{4\cos^3x-3\cos x}{\sin x(1-\sin^2x)}=\frac{1}{4}(\frac{4\cos x}{\sin x}-\frac{3}{\sin x\cos x})=\frac{\cos x}{\sin x}-\frac{3}{2}\frac{1}{\sin(2x)}.
$$

Par suite,

$$\int_{}^{}\frac{\cos(3x)}{\sin x+\sin(3x)}\;dx=\ln|\sin x|-\frac{3}{4}\ln|\tan x|+C.$$

\item  $\cos^4x+\sin^4x=(\cos^2x+\sin^2x)^2-2\sin^2x\cos^2x=1-\frac{1}{2}\sin^2(2x)$, et donc

\begin{align*}\ensuremath
\int_{}^{}\frac{1}{\cos^4x+\sin^4x}\;dx&=\int_{}^{}\frac{1}{1-\frac{1}{2}\sin^2(2x)}\;dx
=\int_{}^{}\frac{1}{2-\sin^2u}\;du\;(\mbox{en posant}\;u=2x)\\
 &=\int_{}^{}\frac{1}{1+\cos^2u}\;du=\int_{}^{}\frac{1}{1+\frac{1}{1+v^2}}\;\frac{dv}{1+v^2}\;(\mbox{en posant}\;v=\tan u)\\
 &=\int_{}^{}\frac{dv}{v^2+2}=\frac{1}{\sqrt{2}}\Arctan\frac{v}{\sqrt{2}}+C
=\frac{1}{\sqrt{2}}\Arctan\frac{\tan(2x)}{\sqrt{2}}+C.
\end{align*}

\item \begin{align*}\ensuremath
\frac{\sin x\sin(2x)}{\sin^4x+\cos^4x+1}\;dx&=\frac{2\sin^2x}{1-2\sin^2x\cos^2x+1}\cos x\;dx=\frac{2\sin^2x}{2-2\sin^2x(1-\sin^2x)}\cos x\;dx\\
 &=\frac{u^2}{u^4-u^2+1}\;du\;(\mbox{en posant}\;u=\sin x).
\end{align*}

Maintenant, $u^4-u^2+1=\frac{u^6+1}{u^2+1}=(u-e^{i\pi/6})(u-e^{-i\pi/6})(u+e^{i\pi/6})(u+e^{-i\pi/6})$, et donc,

$$\frac{u^2}{u^4-u^2+1}=\frac{a}{u-e^{i\pi/6}}+\frac{\overline{a}}{u-e^{-i\pi/6}}-\frac{a}{u+e^{i\pi/6}}-\frac{\overline{a}}{u+e^{-i\pi/6}},$$

ou $a=\frac{(e^{i\pi/6})^2}{(e^{i\pi/6}-e^{-i\pi/6})(e^{i\pi/6}+e^{i\pi/6})(e^{i\pi/6}+e^{-i\pi/6})}=
\frac{(e^{i\pi/6})^2}{i.2e^{i\pi/6}.\sqrt{3}}=\frac{-ie^{i\pi/6}}{2\sqrt{3}}$, et donc

\begin{align*}\ensuremath
\frac{u^2}{u^4-u^2+1}&=\frac{1}{2\sqrt{3}}(\frac{-ie^{i\pi/6}}{u-e^{i\pi/6}}+\frac{ie^{-i\pi/6}}{u-e^{-i\pi/6}}+\frac{ie^{i\pi/6}}{u+e^{i\pi/6}}-\frac{ie^{-i\pi/6}}{u+e^{-i\pi/6}})\\
 &=\frac{1}{2\sqrt{3}}(\frac{u}{u^2-\sqrt{3}u+1}-\frac{u}{u^2+\sqrt{3}u+1})\\
 &=\frac{1}{2\sqrt{3}}(\frac{1}{2}\frac{2u-\sqrt{3}}{u^2-\sqrt{3}u+1}+\frac{\sqrt{3}}{2}\frac{1}
 {u^2-\sqrt{3}u+1}-\frac{1}{2}\frac{2u+\sqrt{3}}{u^2+\sqrt{3}u+1}+\frac{\sqrt{3}}{2}\frac{1}{u^2+\sqrt{3}u+1})\\
 &=\frac{1}{4\sqrt{3}}(\frac{2u-\sqrt{3}}{u^2-\sqrt{3}u+1}-\frac{2u+\sqrt{3}}{u^2+\sqrt{3}u+1})
 +\frac{1}{4}(\frac{1}{(u+\frac{\sqrt{3}}{2})^2+(\frac{1}{2})^2}+
 \frac{1}{(u-\frac{\sqrt{3}}{2})^2+(\frac{1}{2})^2})
\end{align*}

et donc,

$$\int_{}^{}\frac{\sin x\sin(2x)}{\sin^4x+\cos^4x+1}\;dx=\frac{1}{4\sqrt{3}}\ln\left|
\frac{\sin^2x-\sqrt{3}\sin x+1}{\sin^2x+\sqrt{3}\sin x+1}\right|+\frac{1}{2}(\Arctan(2\sin x-\sqrt{3})+\Arctan(2\sin x+\sqrt{3})+C.$$

\item  En posant $u=\sin x$, on obtient 

$$\frac{\tan x}{1+\sin(3x)}\;dx=\frac{\sin x}{1+3\sin x-4\sin^3x}\frac{1}{\cos^2x}\cos x\;dx=\frac{u}{(1+3u-4u^3)(1-u^2)}\;du$$

Or, $1+3u-4u^3=(u+1)(-4u^2-4u-1)=-(u-1)(2u+1)^2$ et donc, $(1+3u-4u^3)(1-u^2)=(u+1)(u-1)^2(2u+1)^2$ et donc,

$$\frac{u}{(1+3u-4u^3)(1-u^2)}=\frac{a}{u+1}+\frac{b_1}{u-1}+\frac{b_2}{(u-1)^2}+\frac{c_1}{2u+1}+\frac{c_2}{(2u+1)^2}.$$

$a=\lim_{u\rightarrow -1}(u+1)f(u)=\frac{-1}{(-1-1)^2(-2+1)^2}=-\frac{1}{4}$, $b_2=\frac{1}{(1+1)(2+1)^2}=\frac{1}{18}$

et $c_2=\frac{-1/2}{(-\frac{1}{2}+1)(-\frac{1}{2}-1)^2}=-\frac{4}{9}$.

Ensuite, $u=0$ fournit $0=a-b_1+b_2+c_1+c_2$ ou encore $c_1-b_1=\frac{1}{4}-\frac{1}{18}+\frac{4}{9}=\frac{23}{36}$. D'autre part, en multipliant par $u$, puis en faisant tendre $u$ vers $+\infty$, on obtient $0=a+b_1+c_1$ et donc $b_1+c_1=\frac{1}{4}$ et donc, $c_1=\frac{4}{9}$ et $b_1=-\frac{7}{36}$. Finalement,

$$\frac{u}{(u+1)(u-1)^2(2u+1)^2}=-\frac{1}{4(u+1)}-\frac{7}{36(u-1)}+\frac{1}{18(u-1)^2}+\frac{4}{9(2u+1)}-\frac{4}{9(2u+1)^2}.$$

Finalement,

$$\int_{}^{}\frac{\tan x}{1+\sin(3x)}\;dx=-\frac{1}{4}\ln(\sin x+1)-\frac{7}{36}\ln(1-\sin x)-\frac{1}{18(\sin x-1)}+\frac{2}{9}\ln|2\sin x+1|+\frac{2}{9}\frac{1}{2\sin x+1}+C$$

\item  (voir 6)) 

\begin{align*}\ensuremath
\int_{}^{}\frac{\cos x+2\sin x}{\sin x-\cos x}\;dx&=\int_{}^{}\frac{\frac{1}{2}((\sin x+\cos x)-(\sin x-\cos x))+((\sin x+\cos x)+(\sin x-\cos x)}{\sin x-\cos x}\;dx\\
 &=\frac{3}{2}\int_{}^{}\frac{\sin x+\cos x}{\sin x-\cos x}\;+\frac{1}{2}\int_{}^{}dx\\
 &=\frac{3}{2}\ln|\sin x-\cos x|+\frac{x}{2}+C.
\end{align*}

\item 

\begin{align*}\ensuremath
\int_{}^{}\frac{\sin x}{\cos(3x)}\;dx&=\int_{}^{}\frac{\sin x}{4\cos^3x-3\cos x}\;dx
=\int_{}^{}\frac{1}{3u-4u^3}\;du\;(\mbox{en posant}\;u=\cos x)\\
 &=\int_{}^{}(\frac{1}{3u}-\frac{1}{3(2u-\sqrt{3})}-\frac{1}{3(2u+\sqrt{3})})\;du\\
 &=\frac{1}{3}(\ln|\cos x|-\frac{1}{2}\ln|2\cos x-\sqrt{3}|-\frac{1}{2}\ln|2\cos x+\sqrt{3}|)+C.
\end{align*}
\item  Dans tous les cas, on pose $t=\tan x$ et donc $dx=\frac{dt}{1+t^2}$.

$$\int_{}^{}\frac{1}{\alpha\cos^2x+\beta\sin^2 x}\;dx=\int_{}^{}\frac{1}{\alpha+\beta\tan^2x}\frac{dx}{\cos^2x}=\int_{}^{}\frac{dt}{\alpha+\beta t^2}.$$

Si $\beta=0$ et $\alpha\neq0$, $\int_{}^{}\frac{1}{\alpha\cos^2x+\beta\sin^2x}\;dx=\frac{1}{\alpha}\tan x+C$.

Si $\beta\neq0$ et $\alpha\beta>0$, 

$$\int_{}^{}\frac{1}{\alpha\cos^2x+\beta\sin^2 x}\;dx=\frac{1}{\beta}\int_{}^{}\frac{1}{t^2+(\sqrt{\frac{\alpha}{\beta}})^2}\;dt=\frac{1}{\sqrt{\alpha\beta}}\Arctan(\sqrt{\frac{\beta}{\alpha}}\tan x)+C.$$

Si $\beta\neq0$ et $\alpha\beta<0$, 

$$\int_{}^{}\frac{1}{\alpha\cos^2x+\beta\sin^2 x}\;dx=\frac{1}{\beta}\int_{}^{}\frac{1}{t^2-(\sqrt{-\frac{\alpha}{\beta}})^2}\;dt=\frac{\mbox{sgn}(\beta)}
{2\sqrt{-\alpha\beta}}\ln\left|\frac{\tan x-\sqrt{-\frac{\alpha}{\beta}}}{\tan x+\sqrt{-\frac{\alpha}{\beta}}}\right|+C.$$

\item  \begin{align*}\ensuremath
\int_{}^{}\frac{\ch^3x}{1+\sh x}\;dx&=\int_{}^{}\frac{1+\sh^2x}{1+\sh x}\ch x\;dx\\
 &=\int_{}^{}\frac{u^2+1}{u+1}\;du\;(\mbox{en posant}\;u=\sh x)\\
 &=\int_{}^{}(u-1+\frac{2}{u+1})\;du=\frac{\sh^2x}{2}-\sh x+2\ln|1+\sh x|+C.
\end{align*}
\item  On peut poser $u=e^x$ mais il y a mieux.

\begin{align*}\ensuremath
\int_{}^{}\sqrt{\ch x-1}\;dx&=\int_{}^{}\frac{\sqrt{(\ch x-1)(\ch x+1)}}{\sqrt{\ch x+1}}\;dx=\mbox{sgn}(x)\int_{}^{}\frac{\sh x}{\sqrt{\ch x+1}}\;dx\\
 &=2\mbox{sgn}(x)\sqrt{\ch x+1}+C.
\end{align*}

\item
\begin{align*}\ensuremath
\int_{}^{}\frac{\tanh x}{\ch x+1}\;dx&=\int_{}^{}\frac{1}{\ch x(\ch x+1)}\sh x\;dx\\
 &=\int_{}^{}\frac{1}{u(u+1)}\;du\;(\mbox{en posant}\;u=\ch x)\\
 &=\int_{}^{}(\frac{1}{u}-\frac{1}{u+1})\;du=\ln\frac{\ch x}{\ch x+1}+C.
\end{align*}

\item $\int_{}^{}\frac{1}{\sh^5x}\;dx=\int_{}^{}\frac{\sh x}{\sh^6x}\;dx
=\int_{}^{}\frac{\sh x}{\sh^6x}\;dx=\int_{}^{}\frac{\sh x}{(\ch^2x-1)^3}\;dx=\int_{}^{}\frac{1}{(u^2-1)^3}\;du$ (en posant $u=\ch x$).

\item

\begin{align*}\ensuremath
\int_{}^{}\frac{1}{1-\ch x}\;dx&=\int_{}^{}\frac{1+\ch x}{1-\ch^2x}\;dx=-\int_{}^{}\frac{1}{\sh^2x}\;dx-\int_{}^{}\frac{\ch x}{\sh^2x}\;dx\\
 &=\mbox{coth}x+\frac{1}{\sh x}+C.
\end{align*}
\end{enumerate}
\fincorrection
\correction{005468}
\begin{enumerate}
\item  
\begin{align*}\ensuremath
\int_{}^{}\frac{1}{\sqrt{x^2+2x+5}}\;dx&=\int_{}^{}\frac{1}{\sqrt{(x+1)^2+2^2}}\;dx=\Argsh\frac{x+1}{2}+C\\
 &=\ln(\frac{x+1}{2}+\sqrt{(\frac{x+1}{2})^2+1})+C=\ln(x+1+\sqrt{x^2+2x+5})+C.
\end{align*}

Puis,

\begin{align*}\ensuremath
\int_{}^{}\sqrt{x^2+2x+5}\;dx&=(x+1)\sqrt{x^2+2x+5}-\int_{}^{}(x+1)\frac{2x+2}{2\sqrt{x^2+2x+5}}\;dx\\
 &=(x+1)\sqrt{x^2+2x+5}-\int_{}^{}\frac{x^2+2x+5-4}{\sqrt{x^2+2x+5}}\;dx\\
 &=(x+1)\sqrt{x^2+2x+5}-\int_{}^{}\sqrt{x^2+2x+5}\;dx+4\int_{}^{}\frac{1}{\sqrt{x^2+2x+5}}\;dx,
\end{align*}

et donc,

$$\int_{}^{}\sqrt{x^2+2x+5}\;dx=\frac{1}{2}(x+1)\sqrt{x^2+2x+5}+2\ln(x+1+\sqrt{x^2+2x+5})+C.$$

(On peut aussi poser $x+1=2\sh u$).

\item  $\int_{}^{}\frac{1}{\sqrt{2x-x^2}}\;dx=\int_{}^{}\frac{1}{\sqrt{1-(x-1)^2}}\;dx=\Arcsin(x-1)+C$.

\item  On pose $u=x^6$ puis $v=\sqrt{1+u}$ (ou directement $u=\sqrt{1+x^6}$) et on obtient~:

\begin{align*}\ensuremath
\int_{}^{}\frac{\sqrt{1+x^6}}{x}\;dx&=\int_{}^{}\frac{\sqrt{1+x^6}}{x^6}\;x^5dx=\frac{1}{6}\int_{}^{}\frac{\sqrt{1+u}}{u}\;du\\
 &=\frac{1}{6}\int_{}^{}\frac{v}{v^2-1}2v\;dv=\frac{1}{3}\int_{}^{}\frac{v^2}{v^2-1}\;dv=\frac{1}{3}(v+\int_{}^{}\frac{1}{v^2-1}\;dv)
=\frac{1}{3}(v+\frac{1}{2}\ln\left|\frac{v-1}{v+1}\right|)+C\\
 &=\frac{1}{3}(\sqrt{1+x^6}+\frac{1}{2}\ln\left|\frac{\sqrt{1+x^6}-1}{\sqrt{1+x^6}+1}\right|)+C\\
\end{align*}

\item  
\begin{align*}\ensuremath
\int_{}^{}\frac{1}{\sqrt{1+x}+\sqrt{1-x}}\;dx&=\int_{}^{}\frac{\sqrt{1+x}-\sqrt{1-x}}{(1+x)-(1-x)}\;dx
=\frac{1}{2}(\int_{}^{}\frac{\sqrt{1+x}}{x}\;dx-\int_{}^{}\frac{\sqrt{1-x}}{x}\;dx)\\
 &=\frac{1}{2}(\int_{}^{}\frac{u}{u^2-1}2u\;du+\int_{}^{}\frac{v}{1-v^2}2v\;dv)\;(\mbox{en posant}\;u=\sqrt{1+x}\;\mbox{et}\;v=\sqrt{1-x})\\
 &=\int_{}^{}(1+\frac{1}{u^2-1})\;du+\int_{}^{}(-1+\frac{1}{1-v^2}\;dv\\
 &=u-v+\frac{1}{2}(\ln\left|\frac{1-u}{1+u}\right|+\ln\left|\frac{1+v}{1-v}\right|)+C\\
 &=\sqrt{1+x}-\sqrt{1-x}+\frac{1}{2}(\ln\left|\frac{1-\sqrt{1+x}}{1+\sqrt{1+x}}\right|+\ln\left|\frac{1+\sqrt{1-x}}{1-\sqrt{1-x}}\right|)+C.
\end{align*}

\item  On pose $u=\sqrt{\frac{x+1}{x-1}}$ et donc $x=\frac{u^2+1}{u^2-1}$, puis $dx=\frac{2u(-2)}{(u^2-1)^2}\;du$. Sur $]1,+\infty[$, on obtient

\begin{align*}
\int_{}^{}\sqrt{\frac{x+1}{x-1}}\;dx&=-2\int_{}^{}u\frac{2u}{(u^2-1)^2}\;du\\
 &=2\frac{u}{u^2-1}-2\int_{}^{}\frac{u^2-1}\;du\\
 &=\frac{2u}{u^2-1}+2\ln\left||\frac{1+u}{1-u}\right|+C\\
 &=2\sqrt{x^2-1}+\ln\left|\frac{\sqrt{x+1}+1}{\sqrt{x+1}-1}\right|+C
\end{align*}

\item  On note $\varepsilon$ le signe de $x$.

$\sqrt{x^4-x^2+1}=\varepsilon x\sqrt{x^2+\frac{1}{x^2}-1}=\varepsilon x\sqrt{(x-\frac{1}{x})^2+1}$ puis, $\frac{x^2+1}{x}.\frac{1}{x}=1+\frac{1}{x^2}=(x-\frac{1}{x})'$. On pose donc $u=x-\frac{1}{x}$ et on obtient

\begin{align*}\ensuremath
\int_{}^{}\frac{x^2+1}{x\sqrt{x^4-x^2+1}}\;dx&=\varepsilon\int_{}^{}\frac{1}{\sqrt{(x-\frac{1}{x})^2+1}}.\frac{x^2+1}{x}\frac{1}{x}\;dx
=\varepsilon\int_{}^{}\frac{1}{\sqrt{u^2+1}}\;du=\varepsilon\Argsh(x-\frac{1}{x})+C\\
 &=\varepsilon\ln(\frac{x^2-1+\varepsilon\sqrt{x^4-x^2+1}}{x})+C.
\end{align*}

\item  Sur $]0,1]$, on pose déjà $u=\sqrt{x}$ et donc, $x=u^2$, $dx=2u\;du$.

$$\int_{}^{}\sqrt{\frac{1-\sqrt{x}}{\sqrt{x}}}\;dx=\int_{}^{}\sqrt{\frac{1-u}{u}}2u\;du=2\int_{}^{}\sqrt{u(1-u)}\;du
=2\int_{}^{}\sqrt{(\frac{1}{2})^2-(u-\frac{1}{2})^2}\;du.$$

Puis, on pose $u-\frac{1}{2}=\frac{1}{2}\sin v$ et donc $du=\frac{1}{2}\cos v\;dv$. 
On note que $x\in]0,1]\Rightarrow u\in]0,1]\Rightarrow v=\Arcsin(2u-1)\in]-\frac{\pi}{2},\frac{\pi}{2}]\Rightarrow\cos v\geq0.$

\begin{align*}\ensuremath
\int_{}^{}\sqrt{\frac{1-\sqrt{x}}{\sqrt{x}}}\;dx&=2\int_{}^{}\sqrt{\frac{1}{4}(1-\sin^2v)}\frac{1}{2}\cos v\;dv
=\frac{1}{2}\int_{}^{}\cos^2v\;dv=\frac{1}{4}\int_{}^{}(1+\cos(2v))\;dv\\
 &=\frac{1}{4}(v+\frac{1}{2}\sin(2v))+C=\frac{1}{4}(v+\sin v\cos v)+C\\
 &=\frac{1}{4}(\Arcsin(2\sqrt{x}-1)+(2\sqrt{x}-1)\sqrt{1-(2\sqrt{x}-1)^2})+C\\
 &\frac{1}{4}(\Arcsin(2\sqrt{x}-1)+2(2\sqrt{x}-1)\sqrt{\sqrt{x}-x})+C
\end{align*}

\item  On pose $x=\sh t$ puis $u=e^t$.

\begin{align*}\ensuremath
\int_{}^{}\frac{1}{1+\sqrt{1+x^2}}\;dx&=\int_{}^{}\frac{1}{1+\ch t}\ch t\;dt=\int_{}^{}\frac{\frac{1}{2}(u+\frac{1}{u})}{1+\frac{1}{2}(u+\frac{1}{u})}\frac{du}{u}
=\int_{}^{}\frac{u^2+1}{u(u^2+2u+1)}\;du=\int_{}^{}(\frac{1}{u}-\frac{2}{(u+1)^2})\;du\\
 &=\ln|u|+\frac{2}{u+1}+C.
\end{align*}

Maintenant, $t=\Argsh x=\ln(x+\sqrt{x^2+1})$ et donc, $u=x+\sqrt{x^2+1}$. Finalement,

$$\int_{}^{}\frac{1}{1+\sqrt{1+x^2}}\;dx=\ln(x+\sqrt{x^2+1})-\frac{2}{x+\sqrt{x^2+1}}+C.$$

\item  On pose $u=\frac{1}{x}$ puis $v=\sqrt[3]{u^3+1}=\frac{\sqrt[3]{x^3+1}}{x}$ et donc $v^3=u^3+1$ puis $v^2\;dv=u^2\;du$.

\begin{align*}\ensuremath
\int_{}^{}\frac{\sqrt[3]{x^3+1}}{x^2}\;dx&=\int_{}^{}\frac{\sqrt[3]{(\frac{1}{u})^3+1}}{\frac{1}{u^2}}\;\frac{-du}{u^2}=-\int_{}^{}\frac{\sqrt[3]{u^3+1}}{u}\;du=-\int_{}^{}\frac{\sqrt[3]{u^3+1}}{u^3}u^2\;du\\
 &=-\int_{}^{}\frac{v}{v^3-1}v^2\;dv=\int_{}^{}(-1-\frac{1}{(v-1)(v^2+v+1)})\;dv\\
 &=\int_{}^{}(-1-\frac{1}{3}\frac{1}{v-1}+\frac{1}{3}\frac{v+2}{v^2+v+1})\;dv\\
 &=-v-\frac{1}{3}\ln|v-1|+\frac{1}{6}\int_{}^{}\frac{2v+1}{v^2+v+1}\;dv+\frac{1}{2}\int_{}^{}\frac{1}{(v+\frac{1}{2})^2+(\frac{\sqrt{3}}{2})^2}\;dv\\
 &=-v-\frac{1}{3}\ln|v-1|+\frac{1}{6}\ln(v^2+v+1)+\sqrt{3}\Arctan(\frac{2v+1}{\sqrt{3}})+C...
\end{align*}
\end{enumerate}
\fincorrection
\correction{005469}
\begin{enumerate}
\item  $\int_{}^{}\frac{1}{x\ln x}\;dx=\ln|\ln x|+C$.
\item  $\int_{}^{}\Arcsin x\;dx=x\Arcsin x-\int_{}^{}\frac{x}{\sqrt{1-x^2}}\;dx=x\Arcsin x+\sqrt{1-x^2}+C$.
\item  $\int_{}^{}\Arctan x\;dx=x\Arctan x-\int_{}^{}\frac{x}{1+x^2}\;dx=x\Arctan x-\frac{1}{2}\ln(1+x^2)+C$.
\item  $\int_{}^{}\Arccos x\;dx=x\Arccos x+\int_{}^{}\frac{x}{\sqrt{1-x^2}}\;dx=x\Arccos x-\sqrt{1-x^2}+C$.
\item  $\int_{}^{}\Argsh x\;dx=x\Argsh x-\int_{}^{}\frac{x}{\sqrt{1+x^2}}\;dx=x\Argsh x-\sqrt{1+x^2}+C$.
\item  $\int_{}^{}\Argch x\;dx=x\Argch x-\int_{}^{}\frac{x}{\sqrt{x^2-1}}\;dx=x\Argch x-\sqrt{x^2-1}+C$.
\item  $\int_{}^{}\Argth x\;dx=x\Argth x-\int_{}^{}\frac{x}{1-x^2}\;dx=x\Argth x+\frac{1}{2}\ln(1-x^2)+C$ (on est sur $]-1,1[$).
\item  $\int_{}^{}\ln(1+x^2)\;dx=x\ln(1+x^2)-2\int_{}^{}\frac{x^2+1-1}{x^2+1}\;dx=x\ln(1+x^2)-2x+2\Arctan x+C$.
\item  
\begin{align*}\ensuremath
\int_{}^{}e^{\mbox{\small{Arccos}}\;x}\;dx&=xe^{\mbox{\small{Arccos}}\;x}+\int_{}^{}\frac{x}{\sqrt{1-x^2}}e^{\mbox{\small{Arccos}}\;x}\;dx\\
 &=xe^{\mbox{\small{Arccos}}\;x}-\sqrt{1-x^2}e^{\mbox{\small{Arccos}}\;x}+\int_{}^{}\sqrt{1-x^2}\frac{-1}{\sqrt{1-x^2}}
 e^{\mbox{\small{Arccos}}\;x}\;dx
\end{align*}

et donc, $\int_{}^{}e^{\mbox{\small{Arccos}}\;x}\;dx=\frac{1}{2}(xe^{\mbox{\small{Arccos}}\;x}-\sqrt{1-x^2}e^{\mbox{\small{Arccos}}\;x})+C$.

\item 
\begin{align*}\ensuremath
\int_{}^{}\cos x\ln(1+\cos x)\;dx&=\sin x\ln(1+\cos x)-\int_{}^{}\sin x\frac{-\sin x}{1+\cos x}\;dx=\sin x\ln(1+\cos x)-\int_{}^{}\frac{\cos^2x-1}{\cos x+1}\;dx\\
 &=\sin x\ln(1+\cos x)-\int_{}^{}(\cos x-1)\;dx=\sin x\ln(1+\cos x)-\sin x+x+C.
\end{align*}
\item  $\int_{}^{}\frac{\Arctan x}{\sqrt{x}}\;dx=2\sqrt{x}\Arctan x-2\int_{}^{}\frac{\sqrt{x}}{x^2+1}\;dx$.

Dans la dernière intégrale, on pose $u=\sqrt{x}$ et donc $x=u^2$ puis, $dx=2u\;du$. On obtient $\int_{}^{}\frac{\sqrt{x}}{x^2+1}\;dx=\int_{}^{}\frac{2u^2}{u^4+1}\;du$. Mais,
\begin{align*}\ensuremath
\frac{2u^2}{u^4+1}&=\frac{1}{\sqrt{2}}(\frac{u}{u^2-\sqrt{2}u+1}-\frac{u}{u^2+\sqrt{2}u+1})\\
 &=\frac{1}{2\sqrt{2}}(\frac{2u-\sqrt{2}}{u^2-\sqrt{2}u+1}-\frac{2u+\sqrt{2}}{u^2+\sqrt{2}u+1})+\frac{1}{2}(\frac{1}{(u-\frac{1}{\sqrt{2}})^2+(\frac{1}{\sqrt{2}})^2}+\frac{1}{(u+\frac{1}{\sqrt{2}})^2+(\frac{1}{\sqrt{2}})^2}).
\end{align*}

Par suite,

$$\int_{}^{}\frac{2u^2}{u^4+1}\;du=\frac{1}{2\sqrt{2}}\ln(\frac{u^2-\sqrt{2}u+1}{u^2+\sqrt{2}u+1})+\frac{1}{\sqrt{2}}
(\Arctan(\sqrt{2}u-1)+\Arctan(\sqrt{2}u+1))+C,$$

et donc,

$$\int_{}^{}\frac{\Arctan x}{\sqrt{x}}\;dx=2\sqrt{x}\Arctan x-\frac{1}{\sqrt{2}}\ln(\frac{x-\sqrt{2x}+1}{x+\sqrt{2x}+1})-\sqrt{2}(\Arctan(\sqrt{2x}-1)+\Arctan(\sqrt{2x}+1))+C.$$

\item  $\frac{x}{(x+1)^2}e^x=\frac{1}{x+1}e^x-\frac{1}{(x+1)^2}e^x=\left(\frac{1}{x+1}e^x\right)'$ et donc 
$\int_{}^{}\frac{xe^x}{(x+1)^2}\;dx=\frac{e^x}{x+1}+C$.

\item  $\int_{}^{}\left(\frac{x}{e}\right)^x\ln x\;dx=\int_{}^{}e^{x\ln x-x}\;d(x\ln x-x)=e^{x\ln x-x}+C=\left(\frac{x}{e}\right)^x\;dx$.

\item  $\int_{}^{}x^n\ln x\;dx=\frac{x^{n+1}}{n+1}\ln x-\frac{1}{n+1}\int_{}^{}x^n\;dx=\frac{x^{n+1}}{n+1}\ln x-\frac{x^{n+1}}{(n+1)^2}+C$.

\item  

\begin{align*}\ensuremath
\int_{}^{}e^{ax}\cos(\alpha x)\;dx&=\mbox{Re}\left(\int_{}^{}e^{(a+i\alpha)x}\;dx\right)=\mbox{Re}\left(\frac{e^{(a+i\alpha)x}}{a+i\alpha}\right)+C
=\frac{e^{ax}}{a^2+\alpha^2}\mbox{Re}((a-i\alpha)(\cos(\alpha x)+i\sin(\alpha x))+C\\
 &=\frac{e^{ax}(a\cos(\alpha x)+\alpha\sin(\alpha x))}{a^2+\alpha^2}+C
\end{align*}

\item $\int_{}^{}\sin(\ln x)\;dx=x\sin(\ln x)-\int_{}^{}\cos(\ln x)\;dx=x\sin(\ln x)-x\cos(\ln x)-\int_{}^{}\sin(\ln x)\;dx$ et donc 

$\int_{}^{}\sin(\ln x)\;dx=\frac{x}{2}(\sin(\ln x)-\cos(\ln x))+C$.

\item En posant $u=x^n$ et donc $du=nx^{n-1}dx$, on obtient

$$\int_{}^{}\frac{\sqrt{x^n+1}}{x}\;dx=\int_{}^{}\frac{\sqrt{x^n+1}}{x^n}x^{n-1}\;dx=\frac{1}{n}\int_{}^{}\frac{\sqrt{u+1}}{u}\;du,$$

puis en posant $v=\sqrt{u+1}$ et donc $u=v^2-1$ et $du=2vdv$, on obtient

$$\int_{}^{}\frac{\sqrt{u+1}}{u}\;du=\int_{}^{}\frac{v}{v^2-1}\;2vdv=2\int_{}^{}\frac{v^2-1+1}{v^2-1}\;dv
=2v+\ln\left|\frac{1-v}{1+v}\right|+C.$$

Finalement,

$$\int_{}^{}\frac{\sqrt{x^n+1}}{x}\;dx=\frac{1}{n}(2\sqrt{x^n+1}+\ln\left|
\frac{1-\sqrt{x^n+1}}{1+\sqrt{x^n+1}}\right|)+C.$$

\item $\int_{}^{}x^2e^x\sin x\;dx=\mbox{Im}(\int_{}^{}x^2e^{(1+i)x}\;dx)$. Or,

\begin{align*}
\int_{}^{}x^2e^{(1+i)x}\;dx&=x^2\frac{e^{(1+i)x}}{1+i}-\frac{2}{1+i}\int_{}^{}xe^{(1+i)x}\;dx
=x^2\frac{e^{(1+i)x}}{1+i}-\frac{2}{1+i}(x\frac{e^{(1+i)x}}{1+i}-\int_{}^{}e^{(1+i)x}\;dx)\\
 &=x^2\frac{(1-i)e^{(1+i)x}}{2}+ixe^{(1+i)x}-i\frac{e^{(1+i)x}}{1+i}+C\\
 &=e^x(\frac{1}{2}x^2(1-i)(\cos x+i\sin x)+ix(\cos x+i\sin x)-\frac{1}{2}(1+i)(\cos x+i\sin x)+C.
\end{align*}

Par suite,

$$\int_{}^{}x^2e^x\sin x\;dx=e^x(\frac{x^2}{2}(\cos x+\sin x)-x\sin x-\frac{1}{2}(\cos x-\sin x))+C.$$

\end{enumerate}
\fincorrection
\correction{005471}
Si $c\neq d$, les primitives considérées sont rationnelles si et seulement si il existe $A$ et $B$ tels que

$$\frac{(x-a)(x-b)}{(x-c)^2(x-d)^2}=\frac{A}{(x-c)^2}+\frac{B}{(x-d)^2}\;(*)$$

\begin{align*}\ensuremath
(*)&\Leftrightarrow\exists(A,B)\in\Rr^2/\;\left\{
\begin{array}{l}
A+B=1\\
-2(Ad+Bc)=-(a+b)\\
Ad^2+Bc^2=ab
\end{array}
\right.\Leftrightarrow\exists(A,B)\in\Rr^2/\;\left\{
\begin{array}{l}
B=1-A\\
A(d-c)+c=\frac{1}{2}(a+b)\\
Ad^2+Bc^2=ab
\end{array}
\right.
\\
 &\Leftrightarrow\exists(A,B)\in\Rr^2/\;\left\{
\begin{array}{l}
A=\frac{a+b-2c}{2(d-c)}\\
B=\frac{2d-a-b}{2(d-c)}\\
Ad^2+Bc^2=ab
\end{array}
\right.
\Leftrightarrow\frac{a+b-2c}{2(d-c)}d^2+\frac{2d-a-b}{2(d-c)}c^2=ab\\
 &\Leftrightarrow d^2(a+b-2c)+c^2(2d-a-b)=2ab(d-c)\Leftrightarrow(a+b)(d^2-c^2)-2cd(d-c)=2ab(d-c)\\
 &\Leftrightarrow2cd+(a+b)(c+d)=2ab\Leftrightarrow(a+b)(c+d)=2(ab-cd).
\end{align*}

Si $c=d$, il existe trois nombres $A$, $B$ et $C$ tels que $(x-a)(x-b)=A(x-c)^2+B(x-c)+C$ et donc tels que

$$\frac{(x-a)(x-b)}{(x-c)^4}=\frac{A}{(x-c)^2}+\frac{B}{(x-c)^3}+\frac{C}{(x-c)^4}.$$

Dans ce cas, les primitives sont rationnelles. Finalement, les primitives considérées sont rationnelles si et seulement si $c=d$ ou ($c\neq d$ et $(a+b)(c+d)=2(ab-cd)$).
\fincorrection
\correction{006866}
\begin{enumerate}
  \item $\int \frac{x+2}{x^2-3x-4}\,dx$

Pour calculer cette intégrale on décompose la fraction $\frac{x+2}{x^2-3x-4}$ en éléments simples, 
le dénominateur n'étant pas irréductible.
On sait que cette fraction rationnelle se décompose avec des dénominateurs de degré $1$ et des constantes aux numérateurs :
$$\frac{x+2}{x^2-3x-4} = \frac{x+2}{(x+1)(x-4)} = \frac{\alpha}{x+1} + \frac{\beta}{x-4}$$
Il ne reste plus qu'à calculer $\alpha$ et $\beta$ à l'aide de votre méthode favorite :
$$\frac{x+2}{x^2-3x-4}  = \frac{-\frac15}{x+1} + \frac{\frac65}{x-4}$$

Chacune de ces fractions est du type $\frac 1u$ qui s'intègre en $\ln |u|$, d'où :
$$\int \frac{x+2}{x^2-3x-4}\,dx = -\frac15\int \frac{1}{x+1}\,dx + \frac65\int \frac{1}{x-4}\,dx
= -\frac15 \ln |x+1| + \frac65 \ln |x-4| + c$$

Cette primitive est définie sur $\Rr\setminus\left\{-1,4\right\} $


  \item $\int \frac{x-1}{x^2+x+1}\,dx$

Le dénominateur $u=x^2+x+1$ est irréductible, la fraction est donc déjà décomposée en éléments simples.
On fait apparaître artificiellement une fraction du type $\frac {u'}{u}$
qui s'intégrera à l'aide du logarithme :

$$\frac{x-1}{x^2+x+1} =  \frac 12 \frac{2x+1}{x^2+x+1} - \frac 32 \frac{1}{x^2+x+1}$$
Chacune de ces fractions s’intègre, la première est du type $\frac {u'}{u}$ dont une primitive sera $\ln |u|$, la deuxième sera 
du type $\frac{1}{1+v^2}$ dont une primitive est $\arctan v$.

En détails cela donne :
\begin{align*}
\int \frac{x-1}{x^2+x+1}\,dx 
  &= \int  \frac 12 \frac{2x+1}{x^2+x+1} \, dx  - \frac 32\int \frac{1}{x^2+x+1} \, dx \\
  &= \frac 12 \big[ \ln |x^2+x+1| \big] -  \frac 32\int \frac{1}{\frac 34} \frac{1}{1+ \left( \frac{2}{\sqrt3}\left(x+\frac 12\right) \right)^2} \, dx \\
  &= \frac 12 \big[ \ln|x^2+x+1| \big] - 2 \int \frac{1}{1+v^2} \, \frac{\sqrt3}{2} dv \quad \text{en posant } v=\frac{2}{\sqrt3}\left(x+\frac 12\right) \\
  &= \frac 12 \big[ \ln|x^2+x+1| \big] - \sqrt3 \big[\arctan v\big] \\
  &= \frac 12 \ln|x^2+x+1| - \sqrt3 \arctan \left( \frac{2}{\sqrt3}\left(x+\frac 12\right) \right) + c \\
\end{align*}

Cette primitive est définie sur $\Rr$.


  \item $\int \sin ^8x\cos ^3x \, dx$

Lorsque l'on a une fonction qui s'exprime comme un polynôme (ou une fraction rationnelle),
on peut tester un des changements de variable $u=\cos x$, $u= \sin x$ ou $u= \tan x$.
Soit vous essayez les trois, soit vous appliquez les règles de Bioche. Ici,
si l'on change $x$ en $\pi-x$ alors $\sin ^8x\cos ^3x \, dx$
devient $\sin ^8(\pi-x)\cos ^3 (\pi-x) \, d(\pi -x)= \sin^8 x (-\cos^3 x) (-dx) = \sin ^8x\cos ^3x \, dx$.
Donc le changement de variable adéquat est $u=\sin x$.


Posons $u=\sin x$, $du = \cos x \, dx$.
\begin{align*}
\int \sin ^8x\cos ^3 x \, dx  
  &= \int \sin^8 x \cos^2 x (\cos x \, dx) = \int \sin^8 x (1-\sin^2 x) (\cos x \, dx) \\ 
  &= \int u^8 (1-u^2) du  = \int u^8 \, du - \int u^{10} \, du \\
  &= \big[\frac19 u^9 \big] - \big[ \frac1{11} u^{11}\big] = \frac19\sin^9 x-\frac1{11}\sin^{11} x + c \\
\end{align*}


Cette primitive est définie sur $\Rr$.


  \item $\int \frac 1{\sin x} \, dx$

Comme $\frac 1{\sin (-x)} \, (-dx) = \frac 1{\sin x} \, dx$ la règle de Bioche nous indique le changement de variable 
$u=\cos x$. Donc $du = -\sin x \, dx$.

Donc
\begin{align*}
\int \frac 1{\sin x} \, dx 
  &= \int \frac{-1}{\sin^2 x} (  -\sin x \, dx) \\
  &= \int \frac{-1}{1-\cos^2 x} (  -\sin x \, dx) \\
  &= -\int \frac{1}{1-u^2} \, du \\
\end{align*}

On décompose cette fraction en éléments simples : $\frac{1}{1-u^2}= \frac 12 \frac{1}{1+u} + \frac 12 \frac{1}{1-u}$.
Donc
\begin{align*}
\int \frac 1{\sin x} \, dx 
  &= - \frac 12 \int \frac{1}{1+u}\, du - \frac 12 \int \frac{1}{1-u}\, du \\
  &= - \frac 12 \big[ \ln |1+u| \big]- \frac 12 \big[ \ln |1-u| \big] \\
  &= - \frac 12 \ln |1+\cos x| - \frac 12 \ln |1-\cos x| + c  \\
\end{align*}

Cette primitive est définie sur tout intervalle du type $\left] k\pi ,\left(
k+1\right) \pi \right[ $, $k \in \Zz$.
Elle peut se réécrire sous différentes formes :
$$\int \frac 1{\sin x} \, dx = \frac 12 \ln \frac{1-\cos x}{1+\cos x} +c= \ln \left| \tan \frac x2\right| +c$$

Un autre changement de variable possible aurait été $t=\tan \frac x2$.

  \item $\int \frac{3-\sin x}{2\cos x+3\tan x}\,dx$

La règle de Bioche nous indique le changement de variable $u =\sin x$, $du = \cos x \, dx$.

\begin{align*}
\int \frac{3-\sin x}{2\cos x+3\tan x}\,dx  
  &= \int \frac{3-\sin x}{2\cos x+3\tan x}\frac{1}{\cos x}(\cos x \,dx) \\
  &= \int \frac{3-\sin x}{2\cos^2 x+3\sin x}(\cos x \,dx) \\
  &= \int \frac{3-\sin x}{2-2\sin^2 x+3\sin x}(\cos x \,dx) \\
  &= \int \frac{3-u}{2-2u^2 +3u} du \\ 
\end{align*}

Occupons nous de la fraction que l'on réduit en éléments simples :
$$\frac{3-u}{2-2u^2 +3u} = \frac{u-3}{(u-2)(2u+1)} = \frac{\alpha}{u-2}+\frac{\beta}{2u+1}$$
On trouve $\alpha = -\frac 15$ et $\beta = \frac 75$.

Ainsi 
\begin{align*}
\int \frac{3-\sin x}{2\cos x+3\tan x}\,dx  
  &=  \int \frac{\alpha \, du}{u-2}+\int\frac{\beta\, du}{2u+1}  \\ 
  &= \alpha \ln |u-2| + \beta \ln |2u+1| + c \\
  &= -\frac 15\ln |2-\sin x| + \frac 75 \ln |1+2\sin x| + c \\
\end{align*}

Cette primitive est définie pour les $x$ vérifiant $1+2\sin x >0$ donc sur tout intervalle du type $\left] -\frac\pi6 + 2k\pi , \frac{7\pi}{6} + 2k\pi\right[ $, $k \in \Zz$.
\end{enumerate}
\fincorrection
\correction{006867}
\begin{enumerate}
  \item $\int_0^{\frac \pi 2}x\sin x \, dx$

Par intégration par parties avec $u=x$, $v'=\sin x$ :
\begin{align*}
\int_0^{\frac \pi 2}x\sin x \, dx 
  &= \big[ uv \big]_0^{\frac \pi 2} - \int_0^{\frac \pi 2} u'v \\
  &= \big[ -x\cos x \big]_0^{\frac \pi 2}  + \int_0^{\frac \pi 2} \cos x \, dx \\
  &= \big[ -x\cos x \big]_0^{\frac \pi 2} +  \big[ \sin x \big]_0^{\frac \pi 2} \\
  &= 0-0 \quad + \quad 1-0 \\
  &= 1 \\
\end{align*}

  \item $\int_0^1 \frac{e^x}{\sqrt{e^x+1}} \,  dx$

Posons le changement de variable $u=e^x$ avec $x=\ln u$ et $du = e^x\, dx$.
La variable $x$ varie de $x=0$ à $x=1$, donc la variable $u=e^x$ varie de
$u=1$ à $u=e$.

\begin{align*}
\int_0^1 \frac{e^x \, dx}{\sqrt{e^x+1}} \,  dx 
  &= \int_1^e \frac{du}{\sqrt{u+1}} \\
  &= \big[ 2\sqrt{u+1} \big]_1^e \\
  &= 2\sqrt{e+1} -2\sqrt 2 \\
\end{align*}



  \item $\int_0^1\frac 1{\left( 1+x^2\right) ^2} \, dx$

Posons le changement de variable $x=\tan t$, 
alors on a $dx = (1+\tan^2 t) dt$, $t=\arctan x$ et  on sait aussi que $1+\tan^2 t = \frac{1}{\cos^2 t}$.
Comme $x$ varie de $x=0$ à $x=1$ alors $t$ doit varier de $t=\arctan 0=0$ à $t=\arctan 1 = \frac \pi4$.

\begin{align*}
\int_0^1\frac 1{\left( 1+x^2\right) ^2} \, dx 
  &= \int_0^{\frac\pi4}  \frac{1}{(1+\tan^2 t)^2} (1+\tan^2 t) dt \\
  &= \int_0^{\frac\pi4}  \frac{dt}{1+\tan^2 t}\\
  &=  \int_0^{\frac\pi4}  \cos^2 t \, dt \\
  & = \frac 12 \int_0^{\frac\pi4} (\cos(2t)+1) \, dt\\
  &= \frac12 \Big[ \frac12 \sin(2t) + t \Big]_0^{\frac\pi4} \\
  &= \frac 14 + \frac \pi8 \\
\end{align*}



  \item $\int_0^1\frac{3x+1}{\left( x+1\right) ^2} \,  dx$

Commençons par décomposer la fraction en éléments simples :
$$\frac{3x+1}{\left( x+1\right) ^2} = \frac{\alpha}{x+1}+\frac{\beta}{(x+1)^2}
= \frac{3}{x+1}-\frac{2}{(x+1)^2}$$
où l'on a trouvé $\alpha=3$ et $\beta=-2$.
La première est une intégrale du type $\int \frac 1u = [\ln |u|]$ et la seconde
$\int \frac 1{u^2} = [ -\frac 1u ]$.

\begin{align*}
\int_0^1\frac{3x+1}{\left( x+1\right) ^2} \,  dx
  &=   3 \int_0^1 \frac{1}{x+1} dx   - 2 \int_0^1 \frac{1}{(x+1)^2} \, dx \\
  &= 3 \Big[ \ln|x+1| \Big]_0^1  - 2 \Big[ - \frac{1}{x+1} \Big]_0^1 \\
  &= 3 \ln 2 - 0 \quad + \quad 1 - 2 \\
  &= 3 \ln 2 - 1 \\
\end{align*}


  \item Notons $I = \int_{\frac 12}^2\left( 1+\frac 1{x^2}\right) \arctan x \, dx$.

Posons le changement de variable $u=\frac 1x$ et 
on a $x=\frac 1u$, $dx = -\frac{du}{u^2}$.
Alors $x$ variant de $x=\frac 12$ à $x=2$,
$u$ varie lui de $u=2$ à $u=\frac 12$ (l'ordre est important !).

\begin{align*}
I &= \int_{\frac 12}^2\left( 1+\frac 1{x^2}\right) \arctan x \, dx  \\
  &= \int_2^{\frac 12}\left( 1+u^2 \right) \arctan \frac 1u  \, \left(-\frac{du}{u^2}\right)  \\
  &= \int_{\frac 12}^2 \left( \frac1{u^2} + 1\right) \arctan \frac 1u  \, du  \\
  &= \int_{\frac 12}^2 \left( \frac1{u^2} + 1\right) \left( \frac \pi 2 - \arctan u \right)   \, du  \quad \text{car} \quad \arctan u+\arctan \frac1u=\frac\pi2\\
  &= \frac \pi 2\int_{\frac 12}^2 \left( \frac1{u^2} + 1 \right)\, du  -  \int_{\frac 12}^2  \left( \frac1{u^2} + 1\right)\arctan u  \, du  \\  
  &= \frac \pi 2 \left[ -\frac1{u} + u\right]_{\frac 12}^2  -  I \\  
  &= \frac{3\pi}{2} - I\\
\end{align*}

Conclusion : $I = \frac{3\pi}{4}.$
\end{enumerate}
\fincorrection
\correction{001280}
\begin{enumerate}
 \item $\displaystyle{ \int _{0}^{\infty }\!
{\frac {{e^{-x}}}{\sqrt {x}}}{dx} }$ est convergente (en fait elle
vaut $\sqrt {\pi }$).
\item $\displaystyle{\int _{1}^{\infty }\!{x^x}{dx} }$ est divergente.
\item $\displaystyle{\int _{0}^{\infty }\!{\frac {\sqrt {x}\sin({x}^{-1})}{\ln (1+x)}}{dx} }$ est divergente.
\item $\displaystyle{\int _{1}^{2}\!{\frac {1}{\sqrt {{x}^{2}-1}}}{dx}=
\ln (2+\sqrt {3}) }$.
\item $\displaystyle{\int _{0}^{\infty }\!{\frac {{x}^{5}}{{x}^{12}+1}}{dx}=
1/12\,\pi  }$.
\item $\displaystyle{\int _{0}^{\infty }\!{e^{-\sqrt {x}}}{dx}=2 }$.
\item $\displaystyle{\int _{1}^{\infty }\! \frac 1 {\sinh(x)}{dx}
= - \ln \tanh(1/2) }$.
\end{enumerate}
\fincorrection
\nocorrection
\nocorrection
\nocorrection
\nocorrection
\nocorrection
\nocorrection
\nocorrection
\nocorrection
\nocorrection
\nocorrection
\correction{001291}
R\'{e}ponses : $\frac{\pi }{2}-\ln 2,\pi ,\frac{1}{\left( n-1\right) ^{2}}$.
\fincorrection
\nocorrection
\nocorrection
\nocorrection
\nocorrection
\nocorrection
\nocorrection
\nocorrection
\nocorrection
\nocorrection
\nocorrection
\nocorrection
\nocorrection
\nocorrection
\nocorrection
\nocorrection
\nocorrection
\nocorrection
\nocorrection
\nocorrection
\nocorrection
\nocorrection
\nocorrection
\nocorrection
\nocorrection
\nocorrection
\correction{004275}
$f(x) = \frac\pi{\sqrt{1-x^2}}$ pour $-1 < x < 1$.
\fincorrection
\correction{004276}
$\frac{\pi^2}4$.
\fincorrection
\correction{004277}
$2nI_n - (2n+2)I_{n+1}=0  \Rightarrow  I_n = -\frac\pi{4n}$.
\fincorrection
\correction{004278}
$I_n+I_{n+2} = \frac{2I_{n+1}}{\sin\alpha}  \Rightarrow 
          I_n = \frac\pi{\cos\alpha} \Bigl(\tan\frac\alpha2 \Bigr)^n$.
\fincorrection
\correction{004279}
$I_n = \pi\sqrt2 \prod_{k=1}^n \bigl( 1-\frac3{4k} \bigr) $.
\fincorrection
\correction{004280}
Décomposer en éléments simples et intégrer.
On obtient $I_n = \frac1{(n-1)!}\sum_{k=1}^n(-1)^kC_{n-1}^{k-1}\ln k$.
\fincorrection
\correction{004281}
\begin{enumerate}
  \item 
  \item 
  \item $A_n+A_{n+2} - 2A_{n+1} = 0  \Rightarrow  A_n = \frac{n\pi}2$.\par
             $A_n - B_n = \frac{\pi}4  \Rightarrow  B_n = \frac{n\pi}2 - \frac{\pi}4$
             pour $n \ge 1$.
  \item $J = \frac{\pi}2$.
\end{enumerate}
 \fincorrection
\nocorrection
\nocorrection
\nocorrection
\correction{004285}
$1-\gamma$.
\fincorrection
\correction{004286}
\begin{enumerate}
  \item Intégrations par parties successives,\par
\begin{align*}
 \int_{t=0}^{+\infty} e^{-t}\ln t\,d t
&=  \int_{t=0}^{+\infty} e^{-t}\frac{t^k}{k!}\left(\ln t - \frac 11 - \frac 12 - \dots - \frac 1k\right)\,d t\cr
&=  \int_{t=0}^{+\infty} e^{-t}\frac{t^k}{k!}\ln(t/k)\,d t + \left(\ln k - \frac 11 - \frac 12 - \dots - \frac 1k\right).\cr\end{align*}

Soit $I_k =  \int_{t=0}^{+\infty} e^{-t}\frac{t^k}{k!}\ln(t/k)\,d t$.
On pose $t=ku$~:
\begin{align*}I_k
          &= \frac{k^{k+1}}{k!} \int_{u=0}^{+\infty} (ue^{-u})^k\ln u\,d u\cr
          &= \frac{k^{k+1}}{k!} \int_{u=0}^{+\infty}  d\left(\frac{(ue^{-u})^{k-1}}{k-1}\right) \frac{u^2e^{-u}\ln u}{1-u}\cr
          &= -\frac{k^{k+1}}{k!(k-1)} \int_{u=0}^{+\infty} (ue^{-u})^{k-1}  d\left(\frac{u^2e^{-u}\ln u}{1-u}\right).\cr\end{align*}

Comme $0\le ue^{-u}\le e^{-1}$, il reste $\sqrt k$ au dénominateur multiplié par quelque chose de borné.
  \item 
\begin{align*}
 \int_{t=0}^1 \frac{1-e^{-t}-e^{-1/t}}t\,d t
&=  \int_{t=0}^1 \frac{1-e^{-t}}t\,d t -  \int_{t=1}^{+\infty}\frac{e^{-t}}t\,d t\cr
&= \lim_{x\to0^+}\left(-\ln x -  \int_{t=x}^{+\infty}\frac{e^{-t}}t\,d t\right)\cr
&= \lim_{x\to0^+}\left((e^{-x}-1)\ln x -  \int_{t=x}^{+\infty}e^{-t}\ln t\,d t\right)\cr
&= - \int_{t=0}^{+\infty} e^{-t}\ln t\,d t.\cr\end{align*}
  \item $ \int_{t=x}^{+\infty}\frac{e^{-t}}t\,d t = (t=-\ln(1-u)) =
 \int_{u=1-e^{-x}}^1 \frac{-d u}{\ln(1-u)}$.
\end{enumerate}
 \fincorrection
\nocorrection
\correction{004288}
$ \int_{t=0}^1 \frac {d t}{\sqrt{1-t^2}} = \frac \pi2$.
\fincorrection
\nocorrection
\correction{004290}
\begin{enumerate}
  \item En supposant $f$ positive décroissante, $ \int_0^{+\infty}f \le tg(t) \le tf(0) +  \int_0^{+\infty}f$.
  \item $ \int_{u=Pt}^{Qt}f(u)\,d u - \sum_{n=P}^{Q-1}tf(nt) =  \int_{u=Pt}^{Qt}(f(u)-f(t[u/t]))\,d u =  \int_{u=Pt}^{Qt}t(1-\{u/t\})f'(u)\,d u\to 0$ lorsque $P,Q\to\infty$.

Donc la série de terme général $tf(nt)$ est de Cauchy~; elle converge.

On a alors ${ \int_0^{+\infty}\!\!f} - tg(t) =  \int_{u=0}^{+\infty}t(1-\{u/t\})f'(u)\,d u\to 0$ lorsque $t\to0^+$.
  \item ${ \int_0^{+\infty}\!\!2f} - 2tg(t) =  \int_{u=0}^{+\infty}tf'(u)\,d u +  \int_{u=0}^{+\infty}t(1-2\{u/t\})f'(u)\,d u = tf(0) -  \int_{u=0}^{+\infty}t^2\{u/t\}(1-\{u/t\})f''(u)\,d u$.
\end{enumerate}
 \fincorrection
\nocorrection
\nocorrection
\correction{004293}
$0 \le xf(x) \le 2 \int_{t=x/2}^x f(t)\,d t \to 0$ (lorsque $x\to+\infty$).\par
         $ \int_{t=1}^x t(f(t)-f(t+1))\,d t =
           \int_{t=1}^2 tf(t)\,d t +  \int_{t=2}^x f(t)\,d t -  \int_{t=x}^{x+1} (t-1)f(t)\,d t
          \to
           \int_{t=1}^2 tf(t)\,d t +  \int_{t=2}^{+\infty} f(t)\,d t$ (lorsque $x\to+\infty$).
\fincorrection
\nocorrection
\correction{004295}
\begin{enumerate}
  \item $ \int_{t=x}^y \frac {f(at)}t \,d t =  \int_{t=ax}^{ay} \frac {f(t)}t \,d t  \Rightarrow 
  \int_{t=x}^y \frac {f(at)-f(t)}t \,d t =  \int_{t=ax}^x \frac {f(t)}t \,d t
 + \int_{t=y}^{ay} \frac {f(t)}t \,d t$.

On obtient $ \int_{t=0}^{+\infty} \frac {f(at) - f(t)}t \,d t = (L-\ell)\ln a$.
  \item $I =  \int_{t=0}^{+\infty} \frac{e^{-t}-e^{-2t}}t\,d t = \ln 2$.
\end{enumerate}
 \fincorrection
\correction{004296}
\begin{enumerate}
  \item 
  \item $\frac\pi4+\frac{\ln 2}2$.
\end{enumerate}
 \fincorrection
\correction{004297}
$ \int_{t=a}^b F(t)\,d t =  \int_{u=a-1}^{a+1} \frac{u-(a-1)}2f(u)\,d u
                        +  \int_{u=a+1}^{b-1} f(u)\,d u
                        +  \int_{u=b-1}^{b+1} \frac{b+1-u}2f(u)\,d u$.\par
$\phantom{ \int_{t=a}^b F(t)\,d t} =
\varphi(a+1) - \frac12 \int_{u=a-1}^{a+1}\varphi(u)\,d u
+  \int_{u=a+1}^{b-1} f(u)\,d u
+ \frac12 \int_{u=b-1}^{b+1}\varphi(u)\,d u - \varphi(b-1)$
où $\varphi$ est une primitive de $f$.
\fincorrection
\correction{004298}
Soit $F(x) =  \int_{t=0}^x f(t)\,d t$ :
         $\frac1x \int_{t=0}^x tf(t)\,d t = F(x) - \frac1x \int_{t=0}^x F(t)\,d t$.
\fincorrection
\nocorrection
\nocorrection
\correction{004301}
\begin{enumerate}
  \item 
  \item $e^{-x}/x$.
  \item $-\ln x$.
\end{enumerate}
 \fincorrection
\correction{004302}
\begin{enumerate}
  \item 
  \item $I_n = \sqrt nK_{2n+1}$,  $J_n = \sqrt nK_{2n-1}$.
  \item $\frac{\sqrt\pi}2$.
\end{enumerate}
 \fincorrection
\correction{004303}
$I_n = \frac{(2n)!\sqrt\pi}{2^{2n+1}n!}$.
\fincorrection
\correction{004304}
Intégrale trivialement convergente. Couper en $\int_0^1$ et $\int_1^{+\infty}$,
changer $x$ en $\frac 1x$ dans l'une des intégrales, regrouper et poser $u=x-\frac1x$.
On obtient $I= \int_{u=0}^{+\infty}e^{-u^2}\,d u = \frac{\sqrt\pi}2$.
\fincorrection
\correction{004305}
L'intégrale converge par parties.
\fincorrection
\correction{004306}
$ \int_{}^x \cos(P(t))\,d t = \left[\frac{\sin(P(t))}{P'(t)}\right]^x
         - \int_{}^x \sin(P(t))\frac{d}{d t}\left(\frac{1}{P'(t)}\right)d t$.
\fincorrection
\correction{004307}
$I=J$ par changement $u \mapsto1/u$.\par
$I+J = \frac\pi2  \Rightarrow  I = J = \frac\pi4$.
\fincorrection
\nocorrection
\nocorrection
\correction{004310}
\begin{align*}
 \int_{x=0}^{X} \int_{t=x}^{+\infty}\frac{\sin t}t\,d t\,d x
%
&=\Bigl[x \int_{t=x}^{+\infty}\frac{\sin t}t\,d t\Bigr]_{x=0}^{X}
+ \int_{x=0}^{X}\sin x\,d x\cr
%
&=X \int_{t=X}^{+\infty}\frac{\sin t}t\,d t + 1 - \cos X\cr
%
&=X\Bigl[\frac{-\cos t}t\Bigr]_{t=X}^{+\infty} - X \int_{t=X}^{+\infty}\frac{\cos t}{t^2}\,d t + 1 - \cos X\cr
%
&=-X\Bigl[\frac{\sin t}{t^2}\Bigr]_{t=X}^{+\infty} - X \int_{t=X}^{+\infty}\frac{2\sin t}{t^3}\,d t + 1\cr
%
&\to 1 \text{ lorsque } X\to+\infty\cr
\end{align*}
\fincorrection
\correction{004311}
$2|ff''| \le f^2 + f''^2$ donc $ff''$ est intégrable.
On en déduit que $f'^2$ admet une limite finie en $+\infty$, et cette limite
est nulle sans quoi $f^2$ ne serait pas intégrable
(si $f'(x)\to\ell$ lorsque $x\to+\infty$ alors $f(x)/x\to\ell$). Ainsi $f'$ est bornée sur~$\R^+$,
$f$ est lipschitzienne et donc uniformément continue.
De plus,
$$ \int_{t=0}^X f'^2(t)\,d t = f(X)f'(X) - f(0)f'(0) -  \int_{t=0}^X f(t)f''(t)\,d t$$
donc $f(X)f'(X)$ admet en $+\infty$ une limite finie ou $+\infty$, et le
cas $f(X)f'(X) = \frac12(f^2)'(X) \to +\infty$ lorsque $X\to+\infty$ contredit l'intégrabilité
de $f^2$ donc ce cas est impossible, ce qui prouve que $f'^2$ est intégrable
sur $\R^+$.
Enfin, $ff'$ est intégrable (produit de deux fonctions de carrés intégrables)
donc $f^2$ admet une limite finie en $+\infty$ et cette
limite vaut zéro par intégrabilité de~$f^2$.
\fincorrection
\correction{004312}
On pose $u_n= \int_{t=n\pi}^{(n+1)\pi} \Bigl|\frac{\sin t}{t}\Bigr|\,d t
= \int_{t=0}^{\pi} \frac{\sin t}{t+n\pi}\,d t$. Par encadrement du dénominateur
on a $u_n\sim \frac2{n\pi}$, d'où
$u_0+\dots+u_n\sim \frac{2\ln n}\pi$ et, par encadrement encore, l'intégrale
arrêtée en $x$ est équivalente à $\frac{2\ln x}\pi$.
\fincorrection
\correction{005713}
\begin{enumerate}
 \item  Pour $x\geqslant0$, $x^2 +4x +1\geqslant0$ et donc la fonction $f~:~x\mapsto x+2 -\sqrt{x^2+4x+1}$ est continue sur $[0,+\infty[$.
Quand $x$ tend vers  $+\infty$, $x+2--\sqrt{x^2+4x+1}=\frac{3}{x+2+\sqrt{x^2+4x+1}}\sim{3}{2x}$. Comme la fonction $x\mapsto\frac{3}{2x}$ est positive et non intégrable au voisinage de $+\infty$, $f$ n?est pas intégrable sur $[0;+\infty[$.

\item  Pour $x\geqslant1$, $1+\frac{1}{x}$ est défini et strictement positif. Donc la fonction $f~:~x\mapsto e-\left(1+\frac{1}{x}\right)^x$ est définie et continue sur $[1,+\infty[$.

Quand $x$ tend vers $+\infty$, $\left(1+\frac{1}{x}\right)^x=e^{x\ln(1+\frac{1}{x})}=e^{1-\frac{1}{2x}+o(\frac{1}{x})}=e-\frac{e}{2x}+ o\left(\frac{1}{x}\right)$ puis $f(x)\underset{x\rightarrow+\infty}{\sim}\frac{e}{2x}$. Puisque la fonction $x\mapsto\frac{e}{2x}$ est positive et non intégrable au voisinage de $+\infty$, $f$ n'est pas intégrable sur $[1,+\infty[$.

\item  La fonction $f~:~x\mapsto\frac{\ln x}{x+e^x}$ est continue et positive sur $]0,+\infty[$.

\textbullet~En $0$, $\frac{\ln x}{x+e^x}\sim\ln x$ et donc $f(x)\underset{x\rightarrow0}{=}o\left(\frac{1}{\sqrt{x}}\right)$. Comme $\frac{1}{2}<1$, la fonction $x\mapsto\frac{1}{\sqrt{x}}$ est intégrable sur un voisinage de $0$ et il en est de même de la fonction $f$.

\textbullet~En $+\infty$,  $f(x)\sim\frac{\ln x}{e^x}=o\left(\frac{1}{x^2}\right)$. Comme $2>1$, la fonction $x\mapsto\frac{1}{x^2}$ est intégrable sur un voisinage de $+\infty$ et il en est de même de la fonction $f$. 

Finalement, $f$ est intégrable sur $]0,+\infty[$.

\item  La fonction $x\mapsto\sqrt[3]{x+1}-\sqrt[3]{x}$ est continue et strictement positive sur $[0,+\infty[$. Donc la fonction $f~:~x\mapsto\left(\sqrt[3]{x+1}-\sqrt[3]{x}\right)$ est continue sur $[0,+\infty[$.

En $+\infty$, $\ln\left(\sqrt[3]{x+1}-\sqrt[3]{x}\right)=\frac{1}{3}\ln x+\ln\left(\left(1+\frac{1}{x}\right)^{1/3}-1\right)=\frac{1}{3}\ln x+\ln\left(\frac{1}{3x}+O\left(\frac{1}{x^2}\right)\right)= -\frac{2}{3}\ln x -\ln3 + O\left(\frac{1}{x}\right)$. Par suite, $\sqrt{x}\ln\left(\sqrt[3]{x+1}-\sqrt[3]{x}\right)=-\frac{2}{3}\sqrt{x}\ln x -\ln3\sqrt{x}+o(1)$.

Mais alors $x^2f(x)\underset{x\rightarrow+\infty}{=}\text{exp}\left(-\frac{2}{3}\sqrt{x}\ln x -\ln3\sqrt{x}+2\ln x+o(1)\right)$ et donc $\lim_{x \rightarrow +\infty}x^2f(x)=0$. Finalement $f(x)$ est négligeable devant $\frac{1}{x^2}$ en $+\infty$ et $f$ est intégrable sur $[0,+\infty[$.

\item  La fonction $f~:~x\mapsto e^{-\sqrt{x^2-x}}$ est continue sur $[1,+\infty[$.

Quand $x$ tend vers $+\infty$, $x^2f(x)=\text{exp}\left(-\sqrt{x^2-x}+2\ln x\right)=\text{exp}(-x+o(x))$ et donc $x^2f(x)\underset{x\rightarrow+\infty}{\rightarrow}0$. $f(x)$ est ainsi négligeable devant $\frac{1}{x^2}$ au voisinage de $+\infty$ et donc $f$ est intégrable sur $[1,+\infty[$.

\item  La fonction $f~:~x\mapsto x^{-\ln x}$ est continue sur $]0,+\infty[$.

\textbullet~Quand $x$ tend vers $0$, $x^{-\ln x}=e^{-\ln^2x}\rightarrow0$. La fonction $f$ se prolonge par continuité en $0$ et est en particulier intégrable sur un voisinage de $0$.

\textbullet~Quand $x$ tend vers $+\infty$, $x^2f(x)=\text{exp}\left(-\ln^2x+2\ln x\right)\rightarrow0$. Donc $f$ est négligeable devant  $\frac{1}{x^2}$ quand $x$ tend vers $+\infty$ et $f$ est intégrable sur un voisinage de $+\infty$.

Finalement, $f$ est intégrable sur $]0,+\infty[$.

\item  La fonction $f~:~x\mapsto\frac{\sin(5x)-\sin(3x)}{x^{5/3}}$  est continue sur $]0,+\infty[$.

\textbullet~Quand $x$ tend vers $0$,  $f(x)\sim\frac{5x-3x}{x^5/3}=\frac{2}{x^{2/3}}>0$. Puisque $\frac{2}{3}<1$, la fonction $x\mapsto\frac{2}{x^{2/3}}$ est positive et intégrable sur un voisinage de $0$ et il en est de même de la fonction $f$.

\textbullet~En $+\infty$, $|f(x)|\leqslant\frac{2}{x^{5/3}}$ et puisque $\frac{5}{3}>1$, la fonction $f$ est intégrable sur  un voisinage de $+\infty$.

Finalement, $f$ est intégrable sur $]0,+\infty[$.

\item  La fonction $f~:~x\mapsto\frac{\ln x}{x^2-1}$ est continue sur $]0,1[\cup]1,+\infty[$.

\textbullet~En $0$,  $f(x)\sim-\ln x=o\left(\frac{1}{\sqrt{x}}\right)$. Donc $f$ est intégrable sur un voisinage de $0$ à droite.

\textbullet~En $1$, $f(x)\sim\frac{\ln x}{2(x-1)}\sim\frac{1}{2}$. La fonction $f$ se prolonge par continuité en $1$ et est en particulier intégrable sur un voisinage de $1$ à gauche ou à droite.

\textbullet~En $+\infty$,  $x^{3/2}f(x)\sim\frac{\ln x}{\sqrt{x}}=o(1)$. Donc $f(x)$ est négligeable devant $\frac{1}{x^{3/2}}$ quand $x$ tend vers $+\infty$ et donc intégrable sur un voisinage de $+\infty$.

Finalement, $f$ est intégrable sur $]0,1[\cup]1,+\infty[$.

\item  La fonction$f~:~x\mapsto\frac{e^x}{\sqrt{|x|}}$ est continue sur $]-\infty,0[\cup]0,+\infty[$ et paire. Il suffit donc d'étudier l'intégrabilité de $f$ sur $]0,+\infty[$.

$f$ est positive et équivalente en $0$ à droite à $\frac{1}{\sqrt{x}}$ et négligeable devant $\frac{1}{x^2}$  en $+\infty$ d'après un théorème de croissances comparées.

$f$ est donc intégrable sur $]0,+\infty[$ puis par parité sur $]-\infty,0[\cup]0,+\infty[$. On en déduit que $\int_{-\infty}^{+\infty}\frac{e^x}{\sqrt{|x|}}\;dx$ existe dans $\Rr$ et vaut par parité $2\int_{0}^{+\infty}\frac{e^x}{\sqrt{|x|}}\;dx$.

\item   La fonction$f~:~x\mapsto\frac{1}{(1+x^2)\sqrt{1-x^2}}$ est continue et positive sur $]-1,1[$, paire et équivalente au voisinage de $1$ à droite à $\frac{1}{2\sqrt{2}\sqrt{1-x}}$. $f$ est donc intégrable sur $]-1,1[$.

\item  La fonction$f~:~x\mapsto\frac{1}{\sqrt[3]{x^2-x^3}}$ est continue et positive sur $]0,1[$, équivalente au voisinage de $0$ à droite à $\frac{1}{x^{2/3}}$ et au voisinage de 1 à gauche à $\frac{1}{(1-x)^{1/3}}$. $f$ est donc intégrable sur $]0,1[$.

\item  La fonction$f~:~x\mapsto\frac{1}{\Arccos(1-x)}$ est continue et positive sur $]0,1]$.

En $0$, $\Arccos(1-x) = o(1)$. Donc $\Arccos(1-x)\sim\sin\left(\Arccos(1-x)\right)=\sqrt{1-(1-x)^2}=\sqrt{2x-x^2}\sim\sqrt{2}\sqrt{x}$.

Donc $f(x)\underset{x\rightarrow0}{\sim}\frac{1}{\sqrt{2}\sqrt{x}}$ et $f$ est intégrable sur $]0,1[$.
\end{enumerate}
\fincorrection
\correction{005714}
\begin{enumerate}
 \item  Pour tout couple de réels $(a,b)$, la fonction $f~:~x\mapsto\frac{1}{x^a\ln^bx}$ est continue et positive sur $[2,+\infty[$. Etudions l'intégrabilité de $f$ au voisinage de $+\infty$.

\textbf{1er cas.} Si $a > 1$,  $x^{(a+1)/2}f(x)=\frac{1}{x^{(a-1)/2}\ln^bx}\underset{x\rightarrow+\infty}{\rightarrow}0$ car  $\frac{a-1}{2}>0$ et d'après un théorème de croissances comparées. Donc $f(x)\underset{x\rightarrow+\infty}{=}o\left(\frac{1}{x^{(a+1)/2}}\right)$. Comme $\frac{a+1}{2}>1$, la fonction $x\mapsto\frac{1}{x^{(a+1)/2}}$ est intégrable sur un voisinage de $+\infty$ et il en est de même de $f$. Dans ce cas, $f$ est intégrable sur $[2,+\infty[$.

\textbf{2ème cas.} Si $a <1$,  $x^{(a+1)/2}f(x)=\frac{x^{(1-a)/2}}{\ln^bx}\underset{x\rightarrow+\infty}{\rightarrow}+\infty$ car  $\frac{1-a}{2}>0$ et d'après un théorème de croissances comparées. Donc $f(x)$ est prépondérant devant $\frac{1}{x^{(a+1)/2}}$ en $+\infty$. Comme $\frac{a+1}{2}<1$, la fonction $x\mapsto\frac{1}{x^{(a+1)/2}}$ n'est pas intégrable sur un voisinage de $+\infty$ et il en est de même de $f$. Dans ce cas, $f$ n'est pas intégrable sur $[2,+\infty[$.

\textbf{3ème cas.} Si $a = 1$. Pour $X > 2$ fixé , en posant $t =\ln x$ et donc $dt=\frac{dx}{x}$ on obtient

\begin{center}
$\int_{2}^{X}\frac{1}{x\ln^b}\;dx=\int_{\ln2}^{\ln X}\frac{dt}{t^b}$.
\end{center}

Puisque $\ln X$ tend vers $+\infty$ quand $X$ tend vers $+\infty$ et que les fonctions considérées sont positives, f est intégrable sur $[2, +\infty[$ si et seulement si $b > 1$.

En résumé , 

\begin{center}
\shadowbox{
la fonction $x\mapsto\frac{1}{x^a\ln^bx}$ est intégrable sur $[2,+\infty[$ si et seulement si $a > 1$ ou ($a=1$ et $b > 1$).
}
\end{center}

(En particulier, la fonction $x\mapsto\frac{1}{x\ln x}$ n'est pas intégrable sur voisinage de $+\infty$ bien que négligeable devant $\frac{1}{x}$ en $+\infty$).

\item  Pour tout réel $a$,  la fonction $f~:~x\mapsto(\tan x)^a$ est continue et strictement positive sur $\left]0,\frac{\pi}{2}\right[$. De plus, pour tout réel $x$ de$\left]0,\frac{\pi}{2}\right[$, on a $f\left(\frac{\pi}{2}-x\right) =\frac{1}{f(x)}$.

\textbullet~\textbf{Etude en $0$ à droite.} $f(x)\underset{x\rightarrow0}{\sim}x^a$. Donc $f$ est intégrable sur un voisinage de $0$ à droite si et seulement si $a > -1$.

\textbullet~\textbf{Etude en $\frac{\pi}{2}$ à gauche.}$f(x)=\frac{1}{f\left(\frac{\pi}{2}-x\right)}\underset{x\rightarrow\frac{\pi}{2}}{\sim}\left(\frac{\pi}{2}-x\right)^{-a}$. Donc $f$ est intégrable sur un voisinage de $\frac{\pi}{2}$ à gauche si et seulement si $a < 1$. 

En résumé, $f$ est intégrable sur $\left]0,\frac{\pi}{2}\right[$ si et seulement si $-1 < a < 1$.

\item  Pour $x\geqslant1$, $1+\frac{1}{x}$ est défini et strictement positif. Donc pour tout couple $(a,b)$ de réels, la fonction $f~:~ x\mapsto\left(1+\frac{1}{x}\right)^{1+\frac{1}{x}}-a -\frac{b}{x}$  est continue sur $[1,+\infty[$.

En $+\infty$, $\left(1+\frac{1}{x}\right)\ln\left(1+\frac{1}{x}\right) =\left(1+\frac{1}{x}\right)\left(\frac{1}{x}+O\left(\frac{1}{x^2}\right)\right)=\frac{1}{x}+O\left(\frac{1}{x^2}\right)$. Donc 

\begin{center}
$f(x)\underset{x\rightarrow+\infty}{=}(1-a)+\frac{1-b}{x}+O\left(\frac{1}{x^2}\right)$.
\end{center}

\textbullet~Si $a\neq 1$, $f$ a une limite réelle non nulle en $+\infty$ et n'est donc pas intégrable sur $[1,+\infty[$.

\textbullet~Si $a = 1$ et $b\neq 1$, $f(x)\underset{x\rightarrow+\infty}{\sim}\frac{1-b}{x}$. En particulier, $f$ est de signe constant sur un voisinage de $+\infty$ et n'est pas intégrable sur $[1,+\infty[$.

\textbullet~Si $a = b = 1$, $f(x)\underset{x\rightarrow+\infty}{=}O\left(\frac{1}{x^2}\right)$ et dans ce cas, $f$ est intégrable sur $[1,+\infty[$.

En résumé, $f$ est intégrable sur $[1,+\infty[$ si et seulement si $a = b = 1$.

\item  Pour tout couple $(a,b)$ de réels, la fonction $f~:~ x\mapsto\frac{1}{x^a(1+x^b)}$ est continue et positive sur $]0,+\infty[$.

\textbullet~\textbf{Etude en $0$.} 

-Si $b>0$, $f(x)\underset{x\rightarrow0}{\sim}\frac{1}{x^a}$, et donc $f$ est intégrable sur un voisinage de $0$ si et seulement si $a<1$,

-si $b=0$, $f(x)\underset{x\rightarrow0}{\sim}\frac{1}{2x^a}$, et donc $f$ est intégrable sur un voisinage de $0$ si et seulement si $a<1$,

-si $b<0$, $f(x)\underset{x\rightarrow0}{\sim}\frac{1}{x^{a+b}}$, et donc $f$ est intégrable sur un voisinage de $0$ si et seulement si $a+b<1$.

\textbullet~\textbf{Etude en $+\infty$.} 

-Si $b>0$, $f(x)\underset{x\rightarrow0}{\sim}\frac{1}{x^{a+b}}$, et donc $f$ est intégrable sur un voisinage de $+\infty$ si et seulement si $a+b>1$,

-si $b=0$, $f(x)\underset{x\rightarrow0}{\sim}\frac{1}{2x^a}$, et donc $f$ est intégrable sur un voisinage de $+\infty$ si et seulement si $a>1$,

-si $b<0$, $f(x)\underset{x\rightarrow0}{\sim}\frac{1}{x^{a}}$, et donc $f$ est intégrable sur un voisinage de $+\infty$ si et seulement si $a>1$.

En résumé, $f$ est intégrable sur $]0,+\infty[$ si et seulement si (($b\geqslant0$ et $a<1$) ou ($b<0$ et $a+b<1$)) et (($b>0$ et $a+b>1$) ou ($b\leqslant0$ et $a>1$)) ce qui équivaut à ($b>0$ et $a+b>1$ et $a<1$) ou ($b<0$ et $a>1$ et $a+b<1$).

Représentons graphiquement l'ensemble des solutions. La zone solution est la zone colorée.

$$\includegraphics{../images/img005714-1}$$

\end{enumerate}
\fincorrection
\correction{005715}

\begin{enumerate}

 \item  Soient $\varepsilon$ et $X$ deux réels tels que $0<\varepsilon<X$. Les deux fonction $x\mapsto1-\cos x$ et $x\mapsto\frac{1}{x}$ sont de classe $C^1$ sur le segment $[\varepsilon,X]$. On peut donc effectuer une intégration par parties et on obtient



\begin{center}

$\int_{\varepsilon}^{X}\frac{\sin x}{x}\;dx=\left[\frac{1-\cos x}{x}\right]_1^X+\int_{\varepsilon}^{X}\frac{1-\cos x}{x^2}\;dx=\frac{1-\cos X}{X}-\frac{1-\cos \varepsilon}{\varepsilon}+\int_{\varepsilon}^{X}\frac{1-\cos x}{x^2}\;dx$.

\end{center}



\textbullet~La fonction $x\mapsto\frac{1-\cos x}{x^2}$ est continue sur $]0,+\infty[$, est prolongeable par continuité en $0$ car $\lim_{x \rightarrow 0}\frac{1-\cos x}{x^2}=\frac{1}{2}$ et donc intégrable sur un voisinage de $0$, est dominée par $\frac{1}{x^2}$ en $+\infty$ et donc intégrable sur un voisinage de $+\infty$. La fonction $x\mapsto\frac{1-\cos x}{x^2}$ est donc intégrable sur $]0,+\infty[$ et $\int_{\varepsilon}^{X}\frac{1-\cos x}{x^2}\;dx$ a une limite réelle quand $\varepsilon$ tend vers $0$ et $X$ tend vers $+\infty$.



\textbullet~$\left|\frac{1-\cos X}{X}\right|\leqslant\frac{1}{X}$ et donc $\lim_{X \rightarrow +\infty}\frac{1-\cos X}{X}=0$.



\textbullet~$\frac{1-\cos\varepsilon}{\varepsilon}\underset{\varepsilon\rightarrow0}{\sim}\frac{\varepsilon}{2}$ et donc $\lim_{\varepsilon \rightarrow \varepsilon}\frac{1-\cos \varepsilon}{\varepsilon}=0$.



On en déduit que $\int_{0}^{+\infty}\frac{\sin x}{x}\;dx$ est une intégrale convergente et de plus



\begin{center}

$\int_{0}^{+\infty}\frac{\sin x}{x}\;dx=\int_{0}^{+\infty}\frac{1-\cos x}{x^2}\;dx=\int_{0}^{+\infty}\frac{2\sin^2(x/2)}{x^2}\;dx=\int_{0}^{+\infty}\frac{2\sin^2(u)}{4u^2}\;2du=\int_{0}^{+\infty}\frac{\sin^2(u)}{u^2}\;du$.

\end{center}

  

\begin{center}

\shadowbox{

L'intégrale $\int_{0}^{+\infty}\frac{\sin x}{x}\;dx$ converge et de plus $\int_{0}^{+\infty}\frac{\sin x}{x}\;dx=\int_{0}^{+\infty}\frac{1-\cos x}{x^2}\;dx=\int_{0}^{+\infty}\frac{\sin^2x}{x^2}\;dx$.

}

\end{center}



\item  La fonction $f~:~ x\mapsto\frac{\sin x}{x^a}$ est continue sur $]0,+\infty[$.



\textbullet~Sur $]0,1[$, la fonction $f$ est de signe constant et l'existence de $\lim_{\varepsilon \rightarrow 0}\int_{\varepsilon}^{1}f(x)\;dx$ équivaut à l'intégrabilité de la fonction $f$ sur $]0,1]$. Puisque $f$ est équivalente en $0$ à $\frac{1}{x^{a-1}}$, l'intégrale impropre $\int_{0}^{1}f(x)\;dx$ converge en $0$ si et seulement si $a > 0$. On suppose dorénavant $a>0$.



\textbullet~Soit $X>1$. Les deux fonction $x\mapsto -\cos x$ et $x\mapsto\frac{1}{x^a}$ sont de classe $C^1$ sur le segment $[1,X]$. On peut donc effectuer une intégration par parties et on obtient



\begin{center}

$\int_{1}^{X}\frac{\sin x}{x^a}\;dx=\left[\frac{-\cos x}{x^a}\right]_1^X-a\int_{1}^{X}\frac{\cos x}{x^{a+1}}\;dx=-\frac{\cos X}{X^a}+\cos1-a\int_{1}^{X}\frac{\cos x}{x^{a+1}}\;dx$.

\end{center}



Maintenant, $\left|\frac{\cos x}{x^{a+1}}\right|\leqslant\frac{1}{x^{a+1}}$ et puisque $a+1>1$, la fonction $x\mapsto\frac{\cos x}{x^{a+1}}$ est intégrable sur un voisinage de $+\infty$. On en déduit que la fonction $X\mapsto\int_{1}^{X}\frac{\cos x}{x^{a+1}}\;dx$ a une limite réelle quand $X$ tend vers $+\infty$. Comme d'autre part, la fonction $X\mapsto-\frac{\cos X}{X^a}+\cos1$ a une limite réelle quand $X$ tend vers $+\infty$, on a montré que l'intégrale impropre $\int_{1}^{+\infty}f(x)\;dx$ converge en $+\infty$.



Finalement

  

\begin{center}

\shadowbox{

l'intégrale $\int_{0}^{+\infty}\frac{\sin x}{x^a}\;dx$ converge si et seulement si $a > 0$.

}

\end{center}



\item  Soit $X$ un réel strictement positif. Le changement de variables $t = x^2$ suivi d'une intégration par parties fournit :



\begin{center}

$\int_{1}^{X}e^{ix^2}\;dx=\int_{1}^{X^2}\frac{e^{it}}{2\sqrt{t}}\;dt =\frac{i}{2}\left(-\frac{e^{iX}}{\sqrt{X}}+e^i-\frac{1}{2}\int_{1}^{X}\frac{e^{it}}{t^{3/2}}\;dt\right)$

\end{center}



Maintenant, $\lim_{X \rightarrow +\infty}\frac{e^{iX}}{\sqrt{X}}=0$ car $\left|\frac{e^{iX}}{\sqrt{X}}\right|=\frac{1}{\sqrt{X}}$. D'autre part, la fonction $t\mapsto\frac{e^{it}}{t^{3/2}}$ est intégrable sur $[1,+\infty[$ car $\left|\frac{e^{it}}{t^{3/2}}\right|=\frac{1}{t^{3/2}}$.  Ainsi, $\int_{1}^{+\infty}e^{ix^2}\;dx$ est une intégrale convergente et puisque d'autre part la fonction $x\mapsto e^{ix^2}$ est continue sur $[0,+\infty[$, on a montré que



\begin{center}

\shadowbox{

l'intégrale $\int_{0}^{+\infty}e^{ix^2}\;dx$ converge.

}

\end{center}



On en déduit encore que les intégrales $\int_{0}^{+\infty}\cos(x^2)\;dx$ et $\int_{0}^{+\infty}\sin(x^2)\;dx$ sont des intégrales convergentes (intégrales de \textsc{Fresnel}).



\item  La fonction $f~:~x\mapsto x^3\sin(x^8)$ est continue sur $[0,+\infty[$.

Soit $X>0$. Le changement de variables $t = x^4$ fournit



\begin{center}

$\int_{0}^{X}x^3\sin(x^8)dx =\frac{1}{4}\int_{0}^{X^4}\sin(t^2)\;dt=\frac{1}{4}\text{Im}\left(\int_{0}^{X^4}e^{it^2}\;dt\right)$.

\end{center}



D'après 3), $\int_{0}^{+\infty}e^{it^2}\;dt$ est une intégrale convergente et donc $\int_{0}^{+\infty}x^3\sin(x^8)\;dx$ converge.



\item  La fonction $f~:~x\mapsto\cos(e^x)$ est continue sur $[0,+\infty[$.

Soit $X>0$. Le changement de variables $t = e^x$ fournit



\begin{center}

$\int_{0}^{X}\cos(e^x)\;dx=\int_{1}^{e^X}\frac{\cos t}{t}\;dt$.

 \end{center}

 



On montre la convergence en $+\infty$ de cette intégrale par une intégration par parties analogue à celle de la question 1). L'intégrale impropre  $\int_{0}^{+\infty}\cos(e^x)\;dx$ converge .



\item  Pour tout réel $x\geqslant0$, $1+x^3\sin^2x\geqslant 1>0$ et donc la fonction $f~:~x\mapsto\frac{1}{1+x^3\sin^2x}$ est continue sur $[0,+\infty[$.\rule[-6mm]{0mm}{0mm}



La fonction $f$ étant positive, la convergence de l'intégrale proposée équivaut à l'intégrabilité de la fonction $f$ sur $[0,+\infty[$, intégrabilité elle-même équivalente à la convergence de la série numérique de terme général $u_n =\int_{n\pi}^{(n+1)\pi}\frac{1}{1+x^3\sin^2x}\;dx$.



Soit $n\in\Nn^*$. On a $u_n\geqslant0$ et d'autre part



\begin{align*}\ensuremath

u_n&=\int_{n\pi}^{(n+1)\pi}\frac{1}{1+x^3\sin^2x}\;dx=\int_{0}^{\pi}\frac{1}{1+(u+n\pi)^3\sin^2u}\;du\\

 &\leqslant\int_{0}^{\pi}\frac{1}{1+n^3\pi^3\sin^2}\;du = 2\int_{0}^{\pi/2}\frac{1}{1+n^3\pi^3\sin^2u}\;du\\

 &\leqslant2\int_{0}^{\pi/2}\frac{1}{1+n^3\pi^3\left(\frac{2u}{\pi}\right)^2}\;du\;(\text{par concavité de la fonction sinus sur}\;[0,\pi])\\

 &= 2\times\frac{1}{2\sqrt{\pi}n^{3/2}}\int_{0}^{(n\pi)^{3/2}}\frac{1}{1+v^2}\;dv\leqslant\frac{1}{\sqrt{\pi}n^{3/2}}\int_{0}^{+\infty}\frac{1}{1+v^2}\;dv =\frac{\sqrt{\pi}}{2n^{3/2}}.

\end{align*}



	

Donc, pour $n\in\Nn^*$, $0\leqslant u_n\leqslant\frac{\sqrt{\pi}}{2n^{3/2}}$ et la série de terme général $u_n$ converge. On en déduit que la fonction $f~:~x\mapsto\frac{1}{1+x^3\sin^2x}$ est intégrable sur $[0,+\infty[$.

\end{enumerate}

\fincorrection
\correction{005716}
Existence et calcul de :

\begin{center}
\begin{tabular}{ll}
\textbf{1) (** I)} $I_n=\int_{0}^{+\infty}\frac{1}{(x^2+1)^n}\;dx$&\textbf{2) (très long)} $\int_{2}^{+\infty}\frac{1}{(x-1)^3(x^4+1)}\;dx$\\
\textbf{3) (** I)} $\int_{0}^{+\infty}\frac{1}{x^3+1}\;dx$&
\textbf{4) (***)} $\int_{0}^{+\infty}\frac{1}{(x+1)(x+2)\ldots(x+n)}\;dx$\\
\rule[-6mm]{0mm}{14mm}\textbf{5)(***)} $\int_{0}^{1}\frac{1}{\sqrt{(1-x)(1+ax)}}\;dx$&\textbf{6) (**)} $\int_{0}^{+\infty}\frac{1}{(e^x+1)(e^{-x}+1)}\;dx$\\
\rule[-6mm]{0mm}{0mm}\textbf{7) (**)} $\int_{0}^{+\infty}\frac{1}{5\ch x+3\sh x+4}\;dx$&\textbf{8) (***)} $\int_{0}^{+\infty}\left(2+(t+3)\ln\left(\frac{t+2}{t+4}\right)\right)dt$\\
\textbf{9) (** I)} $\int_{0}^{+\infty}\frac{x\Arctan x}{(1+x^2)^2}\;dx$&
\rule[-6mm]{0mm}{0mm}\textbf{10) (I très long)} $\int_{0}^{+\infty}\frac{x\ln x}{(x^2+1)^a}\;dx$ (calcul pour $a\in\left\{\frac{3}{2},2,3\right\}$)\\
\rule[-7mm]{0mm}{0mm}\textbf{11) (***)} $\int_{0}^{\pi/2}\sqrt{\tan x}\;dx$&\textbf{12) (*** I)} $\int_{0}^{+\infty}\frac{e^{-at}-e^{-bt}}{t}\;dt$ ($0 < a < b$)
\end{tabular}
\end{center}
\fincorrection
\correction{005717}
La fonction $f~:~x\mapsto\ln(\sin x)$ est continue sur $\left]0,\frac{\pi}{2}\right]$. De plus, quand $x$ tend vers $0$, $\ln(\sin x)\sim\ln x=o\left(\frac{1}{\sqrt{x}}\right)$. Par suite, $f$ est intégrable sur$\left]0,\frac{\pi}{2}\right]$.

\begin{enumerate}
 \item  Soient $I=\int_{0}^{\pi/2}\ln(\sin x)\;dx$ et $J =\int_{0}^{\pi/2}\ln(\cos x)\;dx$. Le changement de variables $x =\frac{\pi}{2}-t$ fournit $J$ existe et $J=I$. Par suite,

\begin{align*}\ensuremath
2I&= I+J =\int_{0}^{\pi/2}\ln(\sin x\cos x)\;dx=-\frac{\pi\ln2}{2}+\int_{0}^{\pi/2}\ln(\sin(2x))\;dx =-\frac{\pi\ln2}{2}+\frac{1}{2}\int_{0}^{\pi}\ln(\sin u)\;du\\
 &= -\frac{\pi\ln2}{2}+\frac{1}{2}\left(I+\int_{\pi/2}^{\pi}\ln(\sin u)\;du\right)
= -\frac{\pi\ln2}{2}+\frac{1}{2}\left(I+\int_{\pi/2}^{0}\ln(\sin (\pi-t))\;(-dt)\right)=-\frac{\pi\ln2}{2}+ I.
\end{align*}

Par suite, $I =-\frac{\pi\ln2}{2}$.

\begin{center}
\shadowbox{
$\int_{0}^{\pi/2}\ln(\sin x)\;dx=\int_{0}^{\pi/2}\ln(\cos x)\;dx=-\frac{\pi\ln2}{2}$.
}
\end{center}

\item  Pour $n\geqslant2$, posons $P_n =\prod_{k=1}^{n-1}\sin\left(\frac{k\pi}{2n}\right)$. Pour $1\leqslant k\leqslant n-1$, on a $0<\frac{k\pi}{2n}<\frac{\pi}{2}$ et donc $P_n > 0$. D'autre part, 
$\sin\left(\frac{(2n-k)\pi}{2n}\right)=\sin\left(\frac{k\pi}{2n}\right)$ et $\sin\frac{n\pi}{2n}=1$. On en déduit que

\begin{center}
$P_n^2=\prod_{k=1}^{2n-1}\sin\left(\frac{k\pi}{2n}\right)$,
\end{center}

puis

\begin{align*}\ensuremath
P_n^2&=\prod_{k=1}^{2n-1}\frac{e^{ik\pi/(2n)}-e^{-ik\pi/(2n)}}{2i}=\frac{1}{(2i)^{2n-1}}\prod_{k=1}^{2n-1}\left(-e^{-ik\pi/(2n)}\right)\prod_{k=1}^{2n-1}\left(1-e^{2ik\pi/(2n)}\right)\\
 &=\frac{1}{(2i)^{2n-1}}(-1)^{2n-1}(e^{-i\pi/2})^{2n-1}\prod_{k=1}^{2n-1}\left(1-e^{2ik\pi/(2n)}\right)=   \frac{1}{2^{2n-1}}\prod_{k=1}^{2n-1}\left(1-e^{2ik\pi/(2n)}\right)
\end{align*}

Maintenant, le polynôme $Q$ unitaire de degré $2n-1$ dont les racines sont les $2n-1$ racines $2n$-èmes de l'unité distinctes de $1$ est

\begin{center}
$\frac{X^{2n}-1}{X-1}= 1+X+X^2+...+X^{2n-1}$
\end{center}

et  donc $\prod_{k=1}^{2n-1}\left(1-e^{2ik\pi/(2n)}\right)=Q(1)=2n$. Finalement,

\begin{center}
$\prod_{k=1}^{n-1}\sin\left(\frac{k\pi}{2n}\right)=P_n=\sqrt{\frac{2n}{2^{2n-1}}}=\frac{\sqrt{n}}{2^{n-1}}$.
\end{center}

Pour $0\leqslant k\leqslant n$, posons alors $x_k=\frac{k\pi}{2n}$ de sorte que $0 = x_0 < x_1< ...< x_n =\frac{\pi}{2}$ est une subdivision de $\left[0,\frac{\pi}{2}\right]$
à pas constant égal à $\frac{\pi}{2n}$.

Puisque la fonction $x\mapsto\ln(\sin x)$ est continue et croissante sur $\left]0,\frac{\pi}{2}\right]$, pour $1\leqslant k\leqslant n-1$, on a  $\frac{\pi}{2n}\ln(\sin(x_k))\leqslant\int_{x_k}^{x_{k+1}}\ln(\sin x)\;dx$ puis en sommant  ces inégalités, on obtient 

\begin{center}
$\frac{\pi}{2n}\ln(P_n)\leqslant\int_{\pi/(2n)}^{\pi/2}\ln(\sin x)\;dx$
\end{center}

De même, pour $0\leqslant k\leqslant n-1$,  $\int_{x_k}^{x_{k+1}}\ln(\sin x)\;dx\leqslant\frac{\pi}{2n}\ln(\sin(x_{k+1}))$ et en sommant  

\begin{center}
$\int_{0}^{\pi/2}\ln(\sin x)\;dx\leqslant\frac{\pi}{2n}\ln(P_n)$.
\end{center}

Finalement,  $\forall n\geqslant2$, $\frac{\pi}{2n}\ln(P_n)+\int_{0}^{\pi/(2n)}\ln(\sin x)\;dx\leqslant I\leqslant\frac{\pi}{2n}\ln(P_n)$. Mais $\ln(P_n) =\ln n -(n-1)\ln2$ et donc  $\frac{\pi}{2n}\ln(P_n)$ tend vers $-\frac{\pi\ln2}{2}$ quand $n$ tend vers $+\infty$ et comme d'autre part, $\int_{0}^{\pi/(2n)}\ln(\sin x)\;dx$ tend vers $0$ quand $n$ tend vers $+\infty$ (puisque la fonction $x~:~\mapsto\ln(\sin x)$ est intégrable sur $\left]0,\frac{\pi}{2}\right]$), on a redémontré que $I = -\frac{\pi\ln2}{2}$.
\end{enumerate}
\fincorrection
\correction{005718}
La fonction $f~:~t\mapsto\frac{\ln t}{t-1}$   est continue et positive sur $]0,1[$, négligeable devant $\frac{1}{\sqrt{t}}$ quand $t$ tend vers  $0$ et prolongeable par continuité en $1$. La fonction $f$ est donc intégrable sur $]0;1[$.

\textbf{1ère solution.} (à la main, sans utilisation d'un théorème d'intégration terme à terme) Pour $t\in]0,1[$ et $n\in\Nn$, 

\begin{center}
$\frac{\ln t}{t-1}=\frac{-\ln t}{1-t}=-\sum_{k=0}^{n}t^k\ln t+\frac{t^{n+1}\ln t}{t-1}$
\end{center}

Pour $t\in]0,1]$ et $n\in\Nn$, posons $f_n(t)=-t^n\ln t$. 

Soit $n\in\Nn$. Chaque fonction $f_k$, $0\leqslant k\leqslant n$, est continue sur $]0,1]$ et négligeable en $0$ devant $\frac{1}{\sqrt{t}}$. Donc chaque fonction $f_k$ est intégrable sur $]0,1]$ et donc sur $]0,1[$. Mais alors, il en est de même de la fonction $t\mapsto\frac{t^{n+1}\ln t}{1-t}=\frac{\ln t}{t-1}+\sum_{k=0}^{n}t^k\ln t$ et

\begin{center}
$\int_{0}^{1}\frac{\ln t}{t-1}\;dt=-\sum_{k=0}^{n}\int_{0}^{1}t^k\ln t\;dt+\int_{0}^{1}\frac{t^{n+1}\ln t}{t-1}\;dt$
\end{center}

\textbullet~La fonction $g~:~t\mapsto\frac{t\ln t}{t-1}$ est continue sur $]0,1[$ et prolongeable par continuité en $0$ et en $1$. Cette fonction est en particulier bornée sur $]0,1[$. Soit $M$ un majorant de la fonction $|g|$ sur $]0,1[$. Pour $n\in\Nn$,

\begin{center}
$\left|\int_{0}^{1}\frac{t^{n+1}\ln t}{t-1}\;dt\right|\leqslant\int_{0}^{1}t^n|g(t)|\;dt\leqslant M\int_{0}^{1}t^n\;dt=\frac{M}{n+1}$.
\end{center}

Par suite, $\lim_{n \rightarrow +\infty}\int_{0}^{1}\frac{t^{n+1}\ln t}{t-1}\;dt=0$. On en déduit que la série de terme général $-\int_{0}^{1}t^k\ln t\;dt$ converge et que

\begin{center}
$\int_{0}^{1}\frac{\ln t}{t-1}\;dt=\sum_{k=0}^{+\infty}\int_{0}^{1}(-t^k\ln t)\;dt$.
\end{center}

\textbullet~Soit $\varepsilon\in]0,1[$. Pour $k\in\Nn$, une intégration par parties fournit

\begin{center}
$\int_{\varepsilon}^{1}(-t^k\ln t)\;dt=\left[-\frac{t^{k+1}\ln t}{k+1}\right]_\varepsilon^{1}+\frac{1}{k+1}\int_{\varepsilon}^{1}t^k\;dt=\frac{\varepsilon^{k+1}\ln \varepsilon}{k+1}+\frac{1-\varepsilon^{k+1}}{(k+1)^2}$.
\end{center}

Quand $\varepsilon$ tend vers $0$, on obtient $\int_{\varepsilon}^{1}(-t^k\ln t)\;dt=\frac{1}{(k+1)^2}$. Finalement,

\begin{center}
$\int_{0}^{1}\frac{\ln t}{t-1}\;dt=\sum_{k=0}^{+\infty}\frac{1}{(k+1)^2}=\sum_{n=1}^{+\infty}\frac{1}{n^2}=\frac{\pi^2}{6}$.
\end{center}

\begin{center}
\shadowbox{
$\int_{0}^{1}\frac{\ln t}{t-1}\;dt=\sum_{n=1}^{+\infty}\frac{1}{n^2}=\frac{\pi^2}{6}$.
}
\end{center}

\textbf{2ème solution.} (utilisation d'un théorème d'intégration terme à terme) Chaque fonction $f_n$ est continue et intégrable sur $]0,1[$ et la série de fonctions de terme général $f_n$ converge simplement vers la fonction $f$ sur $]0,1[$ et de plus, la fonction $f$ est continue sur $]0,1[$. Enfin

\begin{center}
$\sum_{n=0}^{+\infty}\int_{0}^{1}|f_n(t)|\;dt=\sum_{n=0}^{+\infty}\frac{1}{(n+1)^2}<+\infty$.
\end{center}

D'après un théorème d'intégration terme à terme, $\int_{0}^{1}f(t)\;dt=\sum_{n=0}^{+\infty}\int_{0}^{1}f_n(t)\;dt=\frac{\pi^2}{6}$.
\fincorrection
\correction{005719}
La fonction $f~:~t\mapsto\frac{t-1}{\ln t}$ est continue sur $]0,1[$, prolongeable par continuité en $0$ et $1$ et donc est intégrable sur $]0,1[$.

Soit $x\in]0,1[$. Chacune des deux fonctions $t\mapsto\frac{t}{\ln t}$  et $t\mapsto\frac{1}{\ln t}$  se prolonge par continuité en $0$ et est ainsi intégrable sur $]0,x]$. On peut donc écrire 

\begin{center}
$\int_{0}^{x}\frac{t-1}{\ln t}\;dt=\int_{0}^{x}\frac{t}{\ln t}\;dt -\int_{0}^{x}\frac{1}{\ln t}\;dt$.
\end{center}

Dans la première intégrale, on pose $u = t^2$ et on obtient $\int_{0}^{x}\frac{t}{\ln t}\;dt =\int_{0}^{x}\frac{2t}{\ln(t^2)}\;dt=\int_{0}^{x^2}\frac{1}{\ln u}\;du$ et donc

\begin{center}
$\int_{0}^{x}\frac{t-1}{\ln t}\;dt=\int_{0}^{x^2}\frac{1}{\ln t}\;dt -\int_{0}^{x}\frac{1}{\ln t}\;dt=\int_{x}^{x^2}\frac{1}{\ln t}\;dt$.
\end{center}

On note alors que, puisque $x\in]0,1[$, $x^2 < x$. Pour $t\in[x^2,x]$, on a $t\ln t < 0$ et donc $\frac{x}{t\ln t}\leqslant\frac{t}{t\ln t}=\frac{1}{\ln t}\leqslant\frac{x^2}\ln t$ puis par croissance de l'intégrale, $\int_{x^2}^{x}\frac{x}{t\ln t}\;dt\leqslant\int_{x^2}^{x}\frac{1}{\ln t}\;dt\leqslant\int_{x^2}^{x}\frac{x^2}{t\ln t}\;dt$ et donc

\begin{center}
$x^2\int_{x}^{x^2}\frac{1}{t\ln t}\;dt\leqslant\int_{x}^{x^2}\frac{1}{\ln t}\;dtx\leqslant x\int_{x}^{x^2}\frac{1}{t\ln t}\;dt$
\end{center}

Maintenant, $\int_{x}^{x^2}\frac{1}{t\ln t}\;dt =\ln|\ln(x^2)|-\ln|\ln x|=\ln2$ et on a montré que, pour tout réel $x$ de $]0,1[$, 

\begin{center}
$x^2\ln2\leqslant\int_{0}^{x}\frac{t-1}{\ln t}\;dt\leqslant x\ln2$
\end{center}

Quand $x$ tend vers $1$, on obtient
 
 \begin{center}
 \shadowbox{
$\int_{0}^{1}\frac{t-1}{\ln t}\;dt =\ln2$.
}
\end{center}
\fincorrection
\correction{005720}
\begin{enumerate}
 \item  La fonction $t\mapsto e^{-t^2}$ est continue, positive et intégrable sur $[0,+\infty[$. De plus, quand $t$ tend $+\infty$,  

\begin{center}
$e^{-t^2}\sim\left(1+\frac{1}{t^2}\right)=\frac{d}{dt}\left(-\frac{1}{2t}e^{-t^2}\right)$.
\end{center}

D'après un théorème de sommation des relations de comparaison, quand $x$ tend vers $+\infty$,   

\begin{center}
$\int_{x}^{+\infty}e^{-t^2}\;dt\sim\int_{x}^{+\infty}\left(-\frac{1}{2t}e^{-t^2}\right)'\;dt =\frac{1}{2x}e^{-x^2}$,
\end{center}

et donc 

\begin{center}
\shadowbox{
$e^{x^2}\int_{x}^{+\infty}e^{-t^2}\;dt\underset{x\rightarrow+\infty}{\sim}\frac{1}{2x}$.
}
\end{center}

\item  Pour $a> 0$ fixé, $\int_{a}^{+\infty}\frac{\cos x}{x}\;dx$ converge (se montre en intégrant par parties (voir exercice \ref{ex:rou3})) puis

\begin{align*}\ensuremath
\int_{a}^{+\infty}\frac{\cos x}{x}\;dx&=-\int_{1}^{a}\frac{\cos x}{x}\;dx+\int_{1}^{+\infty}\frac{\cos x}{x}\;dx\underset{a\rightarrow0}{=}-\int_{1}^{a}\frac{\cos x}{x}\;dx + O(1)\\
 &\underset{a\rightarrow0}{=}-\int_{1}^{a}\frac{1}{x}\;dx+\int_{1}^{a}\frac{1-\cos x}{x}\;dx+O(1)\underset{a\rightarrow0}{=}-\ln a+\int_{1}^{a}\frac{1-\cos x}{x}\;dx+O(1) .
\end{align*}

Maintenant, $\frac{1-\cos x}{x}\underset{x\rightarrow0}{\sim}\frac{x}{2}$ et en particulier, $\frac{1-\cos x}{x}$ tend vers $0$ quand $x$ tend vers $0$. Par suite, la fonction $x\mapsto\frac{1-\cos x}{x}$ est continue sur $]0,1]$ et  se prolonge par continuité en $0$. Cette fonction est donc intégrable sur $]0,1]$ et en particulier, $\int_{1}^{a}\frac{1-\cos x}{x}\;dx$ a une limite réelle quand $a$ tend vers $0$. On en déduit que $\int_{a}^{+\infty}\frac{\cos x}{x}\;dx\underset{a\rightarrow0}{=}-\ln a + O(1)$ et finalement

\begin{center}
\shadowbox{
$\int_{a}^{+\infty}\frac{\cos x}{x}\;dx\underset{a\rightarrow0}{\sim}-\ln a$.
}
\end{center}

\item  Soit $a>0$.

\begin{center}
$\left|\int_{0}^{1}\frac{1}{x^3+a^2}\;dx -\frac{1}{a^2}\right| =\left|\int_{0}^{1}\left(\frac{1}{x^3+a^2}-\frac{1}{a^2}\right)\;dx\right| =\int_{0}^{1}\frac{x^3}{(x^3+a^2)a^2}\;dx\leqslant\frac{1}{a^4}$
\end{center}

Donc, $\int_{0}^{1}\frac{1}{x^3+a^2}\;dx\underset{a\rightarrow+\infty}{=}\frac{1}{a^2}+o\left(\frac{1}{a^2}\right)$ ou encore

\begin{center}
\shadowbox{
$\int_{0}^{1}\frac{1}{x^3+a^2}\;dx\underset{a\rightarrow+\infty}{\sim}\frac{1}{a^2}$.
}
\end{center}
\end{enumerate}
\fincorrection
\correction{005721}
\textbullet~\textbf{Domaine de définition.} Soit $x\in\Rr$.

Si $x < 0$, la fonction $t\mapsto\frac{1}{\ln t}$  n'est pas définie sur $[x,0[\subset[x,x^2]$ et $f(x)$ n'est pas défini.

Si $0< x < 1$, $[x^2,x]\subset]0,1[$. Donc la fonction $t\mapsto\frac{1}{\ln t}$ est continue sur $[x^2,x]$. Dans ce cas, $f(x)$ existe et est de plus strictement positif car $\ln t < 0$ pour tout $t$ de $]0,1[$.
 

Si $x > 1$, $[x,x^2]\subset]1,+\infty[$. Donc la fonction $t\mapsto\frac{1}{\ln t}$ est continue sur $[x,x^2]$. Dans ce cas aussi, $f(x)$ existe et est strictement positif.

Enfin, $f(0)$ et $f(1)$ n'ont pas de sens.

\begin{center}
\shadowbox{
$f$ est définie sur $D =]0,1[\cup]1,+\infty[$ et strictement positive sur $D$.
}
\end{center}

 \textbullet~\textbf{Dérivabilité.} Soit $I$ l'un des deux intervalles $]0,1[$ ou $]1,+\infty[$. La fonction $t\mapsto\frac{1}{\ln t}$ est continue sur $I$. Soit $F$ une primitive de cette fonction sur $I$.
 
Si $x\in]0,1[$, on a $[x^2,x]\subset]0,1[$ et donc $f(x)=F(x^2)-F(x)$. De même, si $x\in]1,+\infty[$.

On en déduit que $f$ est de classe $C^1$ sur $D$. De plus, pour $x\in D$, 

\begin{center}
$f'(x) =2xF'(x^2) - F'(x) =\frac{2x}{\ln(x^2)}-\frac{1}{\ln x}=\frac{x-1}{\ln x}$.
\end{center}

\textbullet~\textbf{Variations.} $f'$ est strictement positive sur $]0,1[\cup]1,+\infty[$ et donc $f$ est strictement croissante sur $]0,1[$ et sur $]1,+\infty[$ (mais pas nécessairement sur $D$).
 
 

 \textbullet~\textbf{Etude en $0$.}
Soit $x\in]0,1[$. On a $0<x^2 < x < 1$ et de plus la fonction $t\mapsto\frac{1}{\ln t}$ est décroissante sur $[x^2,x]\subset]0,1[$ en tant qu'inverse d'une fonction strictement négative et strictement croissante sur $]0,1[$. Donc, $\frac{x-x^2}{\ln x}\leqslant\int_{x^2}^{x}\frac{1}{\ln t}\;dt\leqslant\frac{x-x^2}{\ln(x^2)}$ puis

\begin{center}
$\forall x\in]0,1[$,  $\frac{x^2-x}{2\ln x}\leqslant f(x)\leqslant\frac{x^2-x}{\ln x}$.
\end{center}

On en déduit que $\lim_{x \rightarrow 0^+}f(x) = 0$ et on peut prolonger $f$ par continuité en $0$ en posant $f(0) = 0$ (on note encore $f$ le prolongement).

Quand $x$ tend vers $0$ par valeurs supérieures, $f'(x) =\frac{x-1}{\ln x}$  tend vers 0. Ainsi,

- $f$ est continue sur $[0,1[$,

- $f$ est de classe $C^1$ sur $]0,1[$,

- $f'$ a une limite réelle quand $x$ tend vers $0$ à savoir $0$.

D'après un théorème classique d'analyse, $f$ est de classe $C^1$ sur $[0,1[$ et $f'(0) = 0$.

 \textbullet~\textbf{Etude en 1.} On a vu à l'exercice \ref{ex:rou7bis} que $\lim_{x \rightarrow 1}f(x)=\ln2$ (la limite à droite en $1$ se traite de manière analogue). On prolonge $f$ par continuité en $1$ en posant $f(1) =\ln2$ (on note encore $f$ le prolongement obtenu).
 

Ensuite quand $x$ tend vers $1$, $f'(x)$ tend vers $1$. Donc $f$ est de classe $C^1$ sur $\Rr^+$ et $f'(1)=1$.

En particulier, $f$ est continue sur $\Rr^+$ et d'après plus haut f est strictement croissante sur $\Rr^+$.

 \textbullet~\textbf{Etude en $+\infty$.} Pour  $x > 1$, $f(x)\geqslant{x^2-x}{\ln x}$. Donc $f(x)$ et $\frac{f(x)}{x}$ tendent vers $+\infty$  quand $x$ tend vers $+\infty$.
La courbe représentative de $f$ admet en $+\infty$ une branche parabolique de direction $(Oy)$.

 \textbullet~\textbf{Convexité.} Pour $x\in D$, $f''(x)=\frac{\ln x-\frac{x-1}{x}}{\ln^2x}$.
 

En $1$, en posant $x = 1+h$ où $h$ tend vers $0$, on obtient

\begin{center}
$f''(1+h)=\frac{(1+h)\ln(1+h)-h}{(1+h)\ln^2(1+h)}=\frac{(1+h)\left(h-\frac{h^2}{2}+o(h^2)\right)-h}{h^2+o(h^2)}=\frac{1}{2}+o(1)$.
\end{center}

$f$ est donc de classe $C^2$ sur $]0,+\infty[$ et $f''(1) =\frac{1}{2}$.

Pour $x\neq 1$, $f''(x)$  est du signe de $g(x)=\ln x - 1 +\frac{1}{x}$ dont la dérivée est $g'(x)=\frac{1}{x}-\frac{1}{x^2}=\frac{x-1}{x^2}$. La fonction $g$ est stictement décroissante sur $]0,1]$ et strictement croissante sur $[1,+\infty[$. Donc pour $x\neq1$, $g(x)> g(1) = 0$. On en déduit que pour tout $x\in]0,+\infty[$, $f''(x) > 0$ et donc que $f$ est strictement convexe sur $\Rr^+$.

\textbullet~\textbf{Graphe.}

$$\includegraphics{../images/img005721-1}$$

\fincorrection
\correction{005722}
La fonction $f~:~x\mapsto\frac{(-1)^{E(x)}}{x}$ est continue par morceaux sur $[1,+\infty[$ et donc localement intégrable sur $[1,+\infty[$.

Soient $X$ un réel élément de $[2,+\infty[$ et $n=E(X)$.

\begin{center}
$\int_{1}^{X}\frac{(-1)^{E(x)}}{x}\;dx=\sum_{k=1}^{n-1}\int_{k}^{k+1}\frac{(-1)^{E(x)}}{x}\;dx+\int_{n}^{X}\frac{(-1)^{E(x)}}{x}\;dx =\sum_{k=1}^{n-1}(-1)^k\ln\left(1+\frac{1}{k}\right)+\int_{n}^{X}\frac{(-1)^{E(x)}}{x}\;dx$.
\end{center}

Or, $\left|\int_{n}^{X}\frac{(-1)^{E(x)}}{x}\;dx\right|\leqslant\frac{X-n}{n}\leqslant\frac{1}{E(X)}$. Cette dernière expression tend vers $0$ quand le réel $X$ tend vers $+\infty$ et donc $\lim_{X \rightarrow +\infty}\int_{n}^{X}\frac{(-1)^{E(x)}}{x}\;dx=0$.

D'autre part, la suite $\left((-1)^k\ln\left(1+\frac{1}{k}\right)\right)_{k\geqslant1}$ est de signe alternée et sa valeur absolue tend vers $0$ en décroissant. La série de terme général $(-1)^k\ln\left(1+\frac{1}{k}\right)$, $k\geqslant1$, converge en vertu du critère spécial aux séries alternées ou encore, quand le réel $X$ tend vers $+\infty$, $\sum_{k=1}^{n-1}(-1)^k\ln\left(1+\frac{1}{k}\right)$ a une limite réelle.

Il en est de même de $\int_{1}^{X}\frac{(-1)^{E(x)}}{x}\;dx$ et l'intégrale $\int_{1}^{+\infty}\frac{(-1)^{E(x)}}{x}\;dx$ converge. De plus 

\begin{center}
$\int_{1}^{+\infty}\frac{(-1)^{E(x)}}{x}\;dx=\sum_{n=1}^{+\infty}(-1)^n\ln\left(1+\frac{1}{n}\right)$.
\end{center}

\textbf{Calcul.} Puisque la série converge, on a $\sum_{k=1}^{+\infty}(-1)^k\ln\left(1+\frac{1}{k}\right)=\lim_{n \rightarrow +\infty}\sum_{k=1}^{2n}(-1)^k\ln\left(1+\frac{1}{k}\right)$. Pour $n\in\Nn^*$,

\begin{align*}
\sum_{k=1}^{2n}(-1)^k\ln\left(1+\frac{1}{k}\right)&=\sum_{k=1}^{n}\left(-\ln\left(1+\frac{1}{2k-1}\right)+\ln\left(1+\frac{1}{2k}\right)\right)=\sum_{k=1}^{n}\ln\left(\frac{(2k-1)(2k+1)}{(2k)^2}\right)\\
 &=\ln\left(\frac{(1\times3\times\ldots\times(2n-1))^2\times(2n+1)}{(2\times4\times\ldots\times(2n))^2}\right) =\ln\left(\frac{1}{2^{4n}}\times\left(\frac{(2n)!}{(n!)^2}\right)^2\times(2n+1)\right) .
\end{align*}
 

D'après la formule de \textsc{Stirling},   

\begin{center}
$\frac{1}{2^{4n}}\times\left(\frac{(2n)!}{(n!)^2}\right)^2\times(2n+1)\underset{n\rightarrow+\infty}{\sim}\frac{1}{2^{4n}}\times\frac{\left(\frac{2n}{e}\right)^{4n}(\sqrt{4\pi n})^2}{\left(\frac{n}{e}\right)^{4n}(\sqrt{2\pi n})^4}\times(2n)=\frac{2}{\pi}$.
\end{center}

Donc $\sum_{n=1}^{+\infty}(-1)^n\ln\left(1+\frac{1}{n}\right)=\ln\left(\frac{2}{\pi}\right)$ et on a montré que

\begin{center}
\shadowbox{
$\int_{1}^{+\infty}\frac{(-1)^{E(x)}}{x}\;dx=\sum_{n=1}^{+\infty}(-1)^n\ln\left(1+\frac{1}{n}\right)=\ln\left(\frac{2}{\pi}\right)$.
}
\end{center}
\fincorrection
\correction{005723}
\begin{enumerate}
 \item  Puisque $f$ est continue, positive et décroissante sur $[1,+\infty[$, pour $x\geqslant2$ on a

\begin{center}
$0\leqslant xf(x) = 2\left(x-\frac{x}{2}\right)f(x)\leqslant2\int_{x/2}^{x}f(t)\;dt =2\left(/dint{x/2}{+\infty}f(t)\;dt-\int_{x}^{+\infty}f(t)\;dt\right)$
\end{center}

Cette dernière expression tend vers $0$ quand $x$ tend vers $+\infty$ car $f$ est intégrable sur $[1,+\infty[$. Donc si $f$ est continue, positive, décroissante et intégrable sur $[1,+\infty[$ alors $f(x) \underset{x\rightarrow+\infty}{=}o\left(\frac{1}{x}\right)$.

\item  La fonction $x\mapsto x(f(x) - f(x+1))$ est continue et positive sur $[1,+\infty[$.

Soit $X\geqslant 1$.

\begin{align*}\ensuremath
\int_{1}^{X}x(f(x) - f(x+1))\;dx&=\int_{1}^{X}xf(x)\;dx -\int_{2}^{X+1}(x-1)f(x)\;dx =\int_{1}^{X}xf(x)\;dx -  \int_{2}^{X+1}xf(x)\;dx +\int_{2}^{X+1}f(x)\;dx\\
 &=\int_{1}^{2}xf(x)\;dx -\int_{X}^{X+1}xf(x)\;dx +\int_{2}^{X+1} f(x)\;dx.
\end{align*}

Maintenant $0\leqslant\int_{X}^{X+1}xf(x)\;dx\leqslant(X+1-X)(X+1)f(X)\leqslant2Xf(X)$. D'après 1), cette dernière expression tend vers $0$ quand $X$ tend vers $+\infty$. Donc, quand $X$ tend vers $+\infty$,  $\int_{1}^{X}x(f(x) - f(x+1))\;dx$ tend vers $\int_{1}^{2}xf(x)\;dx+\int_{2}^{+\infty} f(x)\;dx$.

Puisque la fonction $x\mapsto x(f(x) - f(x+1))$ est continue et positive sur $[1,+\infty[$, on sait que $x\mapsto x(f(x) - f(x+1))$ est intégrable sur $[1,+\infty[$ si et seulement si la fonction $X\mapsto\int_{1}^{X}x(f(x) - f(x+1))\;dx$ a une limite réelle quand $X$ tend vers $+\infty$. Donc la fonction $x\mapsto x(f(x) - f(x+1))$ est intégrable sur $[1,+\infty[$ et

\begin{center}
\shadowbox{
$\int_{1}^{+\infty}x(f(x) - f(x+1))\;dx =\int_{1}^{2}xf(x)\;dx +\int_{2}^{+\infty}f(x)\;dx$.
}
\end{center}
\end{enumerate}
\fincorrection
\correction{005724}
\begin{enumerate}
 \item  Puisque $f$ est de classe $C^1$ sur $\Rr^+$, pour $x\geqslant0$,  $\int_{0}^{x}f'(t)\;dt =f(x)-f(0)$. Donc l'intégrale $\int_{0}^{+\infty}f'(t)\;dt$ converge en $+\infty$ si et seulement si $f$ a une limite réelle $\ell$ quand $x$ tend vers $+\infty$.

Si de plus l'intégrale $\int_{0}^{+\infty}f(t)\;dt$ converge, il est exclus d'avoir $\ell\neq0$ et réciproquement si $\ell= 0$ alors $\int_{0}^{x}f'(t)\;dt$ tend vers $-f(0)$ quand $x$ tend vers $+\infty$.
Donc l'intégrale $\int_{0}^{+\infty}f'(t)\;dt$ converge si et seulement si $\lim_{x \rightarrow +\infty}f(x)=0$. 

\item  
  \begin{enumerate} 
  \item Soit $x\geqslant0$. D'après la formule de \textsc{Taylor}-\textsc{Lagrange}, il existe un réel $\theta_x\in]x,x+1[$

\begin{center}
$f(x+1) = f(x) +(x+1-x)f'(x) +\frac{1}{2}f''(\theta_x)$.
\end{center}

ce qui s'écrit encore $f'(x) = f(x+1) -f(x) -\frac{1}{2}f''(\theta_x)$. Quand $x$ tend vers $+\infty$, $f(x+1)-f(x)$ tend vers $0$ et d'autre part, $\theta_x$ tend vers $+\infty$. Ainsi, si $f$ et $f''$ ont une limite réelle quand $x$ tend vers $+\infty$, $f'$ a également une limite réelle et de plus $\lim_{x \rightarrow +\infty}f'(x) = -\frac{1}{2}\lim_{x \rightarrow +\infty}f''(x)$.

Ensuite, puisque pour $x\geqslant0$, $\int_{0}^{x}f'(t)\;dt = f(x) - f(0)$ et $\int_{0}^{x}f''(t)\;dt = f'(x) - f'(0)$, les intégrales $\int_{0}^{+\infty}f'(t)\;dt$ et $\int_{0}^{+\infty}f''(t)\;dt$  convergent et d'après 1), $\lim_{x \rightarrow +\infty}f'(x) = 0$ ($=\lim_{x \rightarrow +\infty}f''(x)$).

  \item Soit $F~:~x\mapsto\int_{0}^{x}f(t)\;dt$. $F$ est de classe $C^3$ sur $\Rr^+$. De plus, $F(x) =\int_{0}^{x}f(t)\;dt$ tend vers $\int_{0}^{+\infty}f(t)\;dt$ et 
$F''(x) =f'(x) = f'(0) +\int_{0}^{x}f''(t)\;dt$ tend vers $f'(0) +\int_{0}^{+\infty}f''(t)\;dt$. Donc $F$ et $F''$ ont des limites réelles en $+\infty$. D'après a), $f =F'$ tend vers $0$ quand $x$ tend vers $+\infty$.
  \end{enumerate}
\end{enumerate}
\fincorrection
\correction{005725}
L'inégalité $|ff''|\leqslant\frac{1}{2}(f^2+f''^2)$ montre que la fonction $ff''$ est intégrable sur $\Rr$ puis, pour $X$ et $Y$ tels que $X\leqslant Y$, une intégration par parties fournit

\begin{center}
$\int_{X}^{Y}f'^2(x)\;dx=\left[f(x)f'(x)\right]_X^Y-\int_{X}^{Y}f(x)f''(x)\;dx$.
\end{center}

Puisque la fonction $f'^2$ est positive, l'intégrabilité de $f'^2$ sur $\Rr$ équivaut à l'existence d'une limite réelle quand $X$ tend vers $+\infty$ et $Y$ tend vers $-\infty$  de $\int_{X}^{Y}f'^2(x)\;dx$ et puisque la fonction $ff''$ est intégrable sur $\Rr$, l'existence de cette limite équivaut, d'après l'égalité précédente, à l'existence d'une limite réelle en $+\infty$ et $-\infty$ pour  la fonction $ff'$.

Si $f'^2$ n'est pas intégrable sur $\Rr^+$ alors $\int_{0}^{+\infty}f'^2(x)dx= +\infty$ et donc $\lim_{x \rightarrow +\infty}f(x)f'(x) = +\infty$. En particulier, pour $x$ suffisament grand, $f(x)f'(x)\geqslant1$ puis par intégration $\frac{1}{2}(f^2(x)-f^2(0))\geqslant x$ contredisant l'intégrabilité de la fonction $f^2$ sur $\Rr$. Donc la fonction $f'^2$ est intégrable sur $\Rr^+$ et la fonction $ff'$ a une limite réelle quand $x$ tend vers $+\infty$.

De même la fonction $f'^2$ est intégrable sur $\Rr^-$ et la fonction $ff'$ a une limite réelle quand $x$ tend vers $-\infty$.

Si cette limite est un réel non nul $\ell$, supposons par exemple $\ell> 0$. Pour $x$ suffisament grand, on a $f(x)f'(x)\geqslant\ell$ puis par intégration $\frac{1}{2}(f^2(x)-f^2(0))\geqslant\ell x$ contredisant de nouveau l'intégrabilité de la fonction $f^2$. Donc la fonction $ff'$ tend vers $0$ en $+\infty$ et de même en $-\infty$.

Finalement, la fonction $f'^2$ est intégrable sur $\Rr$ et $\int_{-\infty}^{+\infty}f'^2(x)\;dx =-\int_{-\infty}^{+\infty}f(x)f''(x)\;dx$.

D'après l'inégalité de \textsc{Cauchy}-\textsc{Schwarz}, on a

\begin{center}
$\left(\int_{-\infty}^{+\infty}f'^2(x)\;dx\right)^2=\left(-\int_{-\infty}^{+\infty}f(x)f''(x)\;dx\right)^2\leqslant\left(\int_{-\infty}^{+\infty}f^2(x)\;dx\right)^2\left(\int_{-\infty}^{+\infty}f''^2(x)\;dx\right)^2$.
\end{center}

Puisque les fonctions $f$ et $f''$ sont continues sur $\Rr$, on a l'égalité si et seulement si la famille $(f,f'')$ est liée.

Donc nécessairement, ou bien $f$ est du type $x\mapsto A\ch(\omega x)+ B\sh(\omega x)$, $\omega$ réel non nul, qui est intégrable sur $\Rr$ si et seulement si $A = B = 0$, ou bien $f$ est affine et nulle encore une fois, ou bien $f$ est du type $x\mapsto A\cos(\omega x)+ B\sin(\omega x)$ et nulle encore une fois.

Donc, on a l'égalité si et seulement si $f$ est nulle.

\fincorrection
\correction{005925}
\begin{enumerate}
\item Pour $n\in\mathbb{N}$, on pose~:
\begin{eqnarray*}
u_{n} = \int_{n\pi}^{(n+1)\pi} \frac{|\sin t|}{t} \, dt.
\end{eqnarray*}
Pour $n\geq 0$, on a $\frac{1}{(n+1)\pi}\leq \frac{1}{t}$, donc
$$
\frac{1}{(n+1)\pi} \,\int_{n\pi}^{(n+1)\pi}|\sin t|\,dt
\,\,\leq\,\, \int_{n\pi}^{(n+1)\pi} \frac{|\sin t|}{t}\,dt.
$$
Or $|\sin t| = (-1)^n \sin t$ sur $[n\pi, (n+1)\pi]$. Ainsi
$$\int_{n\pi}^{(n+1)\pi} |\sin t| \,dt = (-1)^{n} \left[-\cos t
\right]_{n\pi}^{(n+1)\pi} = 2.$$
Il en d\'ecoule que $
\frac{2}{(n+1)\pi} \leq u_{n}$ et $$\int_{0}^{+\infty} \frac{|\sin
t|}{t}\, dt =  \sum_{n=0}^{+\infty} u_{n} \geq \frac{2}{\pi}
\sum_{n=0}^{+\infty} \frac{1}{n+1}$$ est divergente.


\item \emph{Deuxi\`eme formule de la moyenne.}
 D'apr\`es l'\'enonc\'e, $F$ est une primitive de $f$ et est positive et d\'ecroissante.
 Puisque la fonction $g$ admet des primitives,
 la fonction $G(y) := \int_{a}^{y} g(x)\,dx$ est la primitive de
$g$ s'annulant en $a$.
 D'apr\`es le th\'eor\`eme des valeurs interm\'ediaires, pour montrer qu'il existe
 $y\in [a,b]$ tel que~:
 \begin{equation*}
\int_{a}^{b}F(x)g(x)\,dx = F(a)\int_{a}^{y}g(x)\,dx,
 \end{equation*}
il suffit de montrer que
$$
F(a) \,\min_{y\in [a, b]}G(y) \,\leq\,
\int_{a}^{b}F(x)g(x)\,dx\,\leq\, F(a) \max_{y\in[a,b]} G(y).
$$
Par int\'egration par parties, on obtient~:
$$
\int_{a}^{b}F(x)g(x)\,dx = F(b)G(b)-F(a)G(a) - \int_{a}^{b} f(x) G(x)\,dx.
$$
avec $F(a)G(a)=0$.
Comme $f$ est n\'egative sur $[a, b]$, on a~:
\begin{eqnarray*}
\min_{y\in[a,b]} G(y)\,\int_{a}^{b} -f(x)\,dx \,\leq\, -
\int_{a}^{b} f(x) G(x)\,dx \,\leq\, \max_{y\in[a,b]} G(y)
\int_{a}^{b} -f(x)\,dx,\\ \Leftrightarrow \,\min_{y\in[a,b]} G(y)
\left(F(a)- F(b)\right)\,\leq\, - \int_{a}^{b} f(x) G(x)\,dx
\,\leq\, \max_{y\in[a,b]} G(y) \left(F(a) - F(b)\right).
\end{eqnarray*}
On en d\'eduit l'encadrement suivant~:
\begin{eqnarray*}
F(b)\left(G(b) - \min_{y\in[a,b]} G(y)\right) +
F(a)\min_{y\in[a,b]} G(y)\,\leq\,  \int_{a}^{b}F(x)g(x)\,dx\\
\leq\, F(b)\left(G(b) - \max_{y\in[a,b]} G(y))\right) +
F(a)\max_{y\in[a,b]} G(y).
\end{eqnarray*}
Les in\'egalit\'es $G(b) - \min_{y\in[a,b]} G(y)\,\geq \,0$ et
$G(b) - \max_{y\in[a,b]} G(y)\,\leq\,0$ et la positivit\'e de $F$
permettent de conclure.


\item D'apr\`es le crit\`ere de Cauchy (voir la proposition des rappels), pour
montrer que $\int_{0}^{+\infty} \frac{\sin t}{t}\, dt$ est
convergente, il suffit de montrer que $\int_{x}^{x'} \frac{\sin
t}{t}\,dt$ tend vers $0$ lorsque $x$ et  $x'$ tendent vers
$+\infty$. D'apr\`es la formule de la moyenne appliqu\'ee \`a
$F(t) = \frac{1}{t}$ et $g(t) = \sin t$, il vient~:
\begin{equation*}
\int_{x}^{x'} \frac{\sin t}{t}\,dt = \frac{1}{x} \int_{x}^{y} \sin
t \,dt
\end{equation*}
pour un certain $y \in [x, x']$. On en d\'eduit que
\begin{equation*}
\left|\int_{x}^{x'} \frac{\sin t}{t}\,dt \right| = \frac{1}{x}
|\cos y - \cos x| \leq \frac{2}{x}.
\end{equation*}
Ainsi $\lim_{x, x'\rightarrow +\infty}\int_{x}^{x'} \frac{\sin
t}{t}\,dt = 0$ et $\int_{0}^{+\infty} \frac{\sin t}{t}\, dt$ est
convergente.



\item 
  \begin{enumerate}
  \item Posons pour $t > 0$, $U(t) = \frac{e^{-\lambda t}}{t}$.
On a $u(t) = U'(t) = -\frac{e^{-\lambda t}}{t^2}\left(\lambda t +
1 \right)< 0$. Ainsi $U$ est positive et d\'ecroissante sur $]0,
+\infty[$. D'apr\`es la deuxi\`eme formule de la moyenne, pour $0<
x\leq y$, il vient~:
\begin{equation*}
\left|\int_{x}^{y} f(t, \lambda)\, dt \right| = \left|\int_{x}^{y}
U(t) \sin t\, dt  \right| = \left| \frac{e^{-\lambda x}}{x}
\int_{x}^{y'} \sin t \, dt \right|,
\end{equation*}
pour un certain $y'\in [x, y]$. On en d\'eduit que
\begin{equation*}
\left|\int_{x}^{y} f(t, \lambda) \,dt \right| \,\leq\, \frac{2
e^{-\lambda x}}{x}.
\end{equation*}


  \item  On remarque que, pour tout $x>0$, la fonction $t \mapsto f(t,
\lambda)$ est continue sur $[0, x]$, donc Riemann-int\'egrable sur
cet intervalle. D'apr\`es le crit\`ere de Cauchy et la question
$4.a)$, les int\'egrales g\'en\'eralis\'ees $\int_{0}^{+\infty}
f(t, \lambda)\,dt$ sont convergentes. Soit $F(\lambda) =
\int_{0}^{+\infty} f(t, \lambda)\,dt$. Pour montrer que les
int\'egrales g\'en\'eralis\'ees $\int_{0}^{+\infty} f(t,
\lambda)\,dt$ convergent uniform\'ement en $\lambda\geq 0$, il
faut montrer que pour tout $\varepsilon> 0$, il existe $x_{0}>0$
tel que pour tout $x> x_{0}$ et pour tout $\lambda\geq 0$,
\begin{equation*}
\left|F(\lambda) - \int_{0}^{x} f(t, \lambda)\,dt \right| \,\leq\,
\varepsilon.
\end{equation*}
Or, d'apr\`es la question $4.a)$,
\begin{equation*}
\left|F(\lambda) - \int_{0}^{x} f(t, \lambda)\,dt \right| \,\leq\,
\frac{2 e^{-\lambda x}}{x} \,\leq\, \frac{2 e^{-\lambda
x_{0}}}{x_{0}} \,\leq\, \frac{2}{x_{0}}.
\end{equation*}
Ainsi pour $x_{0} > \frac{2}{\varepsilon}$, on a l'in\'egalit\'e
d\'esir\'ee, et ce ind\'ependamment de la valeur de $\lambda$.
Posons $F_{n}(\lambda) = \int_{0}^{n}f(t, \lambda) \,dt$.
D'apr\`es ce qui pr\'ec\`ede,
$$
\sup_{\lambda\in [0,+\infty[} |F(\lambda) - F_{n}(\lambda)| \leq
\frac{2}{n},
$$
i.e. $F_{n}$ converge uniform\'ement vers $F$ sur $[0, +\infty[$.
Comme les fonction $F_{n}$ sont continues, il en d\'ecoule que $F$
est continue. On peut aussi revenir \`a la d\'efinition de
continuit\'e~:
 pour montrer que la fonction $\lambda \mapsto F(\lambda)$
est continue en un point $\lambda_{0} \in [0, +\infty[$, il faut
montrer que pour tout $\varepsilon>0$, il existe un voisinage de
${\lambda}_{0}$ tel que $|F(\lambda) - F({\lambda}_{0})| <
\varepsilon$ pour tout $\lambda $ dans ce voisinage. Soit
$\varepsilon>0$ fix\'e, et posons $x_{\varepsilon} =
\frac{6}{\varepsilon}$. On a~:
\begin{eqnarray*}
|F(\lambda) - F({\lambda}_{0})| &=&
\left|\int_{0}^{x_{\varepsilon}} \left(f(t, \lambda) - f(t,
{\lambda}_{0}\right)\,dt + \int_{x_{\varepsilon}}^{+\infty} f(t,
\lambda)\, dt - \int_{x_{\varepsilon}}^{+\infty} f(t,
{\lambda}_{0})\,dt
\right|\\
&\leq & \left|\int_{0}^{x_{\varepsilon}} \left(f(t, \lambda) -
f(t, {\lambda}_{0}\right)\,dt\right| + \left|
\int_{x_{\varepsilon}}^{+\infty} f(t, \lambda)\, dt \right| +
\left| \int_{x_{\varepsilon}}^{+\infty} f(t,
{\lambda}_{0})\,dt \right|\\
& \leq & \left|\int_{0}^{x_{\varepsilon}} \left(f(t, \lambda) -
f(t, {\lambda}_{0}\right)\,dt \right|+ \frac{2}{3}\varepsilon.
\end{eqnarray*}
Pour conclure, il suffit de trouver un voisinage de ${\lambda}_0$
tel que  $\left|\int_{0}^{x_{\varepsilon}} \left(f(t, \lambda) -
f(t, {\lambda}_{0}\right)\,dt\right| < \frac{\varepsilon}{3}$ pour
tout $\lambda $ dans ce voisinage. L'existence d'un tel voisinage
est garantie par la continuit\'e de la fonction $\lambda \mapsto
\int_{0}^{x_{\varepsilon}}f(t, \lambda)\,dt$ donn\'ee par le
th\'eor\`eme de continuit\'e d'une int\'egrale d\'ependant d'un param\`etre (voir les rappels).
On peut \'egalement d\'eterminer
l'existence de ce voisinage \`a la main de la fa\c{c}on suivante~:
\begin{eqnarray*}
\left|\int_{0}^{x_{\varepsilon}} \left(f(t, \lambda) - f(t,
{\lambda}_{0})\right)\,dt\right| &\,\leq\,&
\int_{0}^{x_{\varepsilon}}
\left|\left(f(t, \lambda) - f(t, {\lambda}_{0}\right)\right|\,dt\\
&\,\leq \,& \sup_{t\in [0, x_{\varepsilon}]} \left| e^{-\lambda t}
- e^{-{\lambda}_0 t}\right| \int_{0}^{x_{\varepsilon}} \left|
\frac{\sin t}{t}
\right|\,dt,\\
& \leq & x_{\varepsilon} \sup_{t\in [0, x_{\varepsilon}]} \left|
e^{-\lambda t} - e^{-{\lambda}_0 t}\right|,
\end{eqnarray*}
o\`u l'on a utiliser l'in\'egalit\'e $|\sin t| \leq |t|$. On a~:
\begin{eqnarray*}
\left| e^{-\lambda t} - e^{-{\lambda}_0 t}\right|& =&
\left|e^{-\frac{(\lambda + {\lambda}_0)t}{2}}
\left(e^{-\frac{({\lambda}_{0} - \lambda)t}{2}} -
e^{\frac{({\lambda}_{0} - \lambda)t}{2}}\right) \right| \\
&\leq\,& 2 \sinh\left(\frac{|\lambda -
{\lambda}_0|t}{2}\right)\,\leq\, 2 \sinh\left(\frac{|\lambda -
{\lambda}_0|x_{\varepsilon}}{2}\right),
\end{eqnarray*}
car la fonction $\sinh$ est croissante.  Ainsi le voisinage de
${\lambda}_0$ d\'etermin\'e par $|\lambda-{\lambda_0}| <
\frac{\varepsilon}{3}\text{argsinh}\left(\frac{\varepsilon^2}{36}\right)
$ convient.

  \item Pour $x \in [0, +\infty[$ et $\lambda \in ]0, +\infty[$,
posons
$$
\tilde{F}(x, \lambda) = \int_{0}^{x} f(t, \lambda) \,dt.
$$
D'apr\`es le th\'eor\`eme de d\'erivabilit\'e d'une int\'egrale d\'ependant d'un
param\`etre (voir les rappels), la fonction $\tilde{F}$ est
d\'erivable par rapport \`a la deuxi\`eme variable et sa
d\'eriv\'ee partielle vaut~:
$$
\frac{\partial \tilde{F}}{\partial \lambda}(x, \lambda) =
\int_{0}^{x} \frac{\partial f}{\partial \lambda}(t, \lambda) \,dt.
$$
Lorsque $x\rightarrow +\infty$, $\tilde{F}(x, \lambda)$ tend vers
$F(\lambda)$. D'apr\`es la question $4.b)$ cette convergence est
uniforme pour $\lambda \geq 0$. D'autre part, lorsque $x$ tend
vers $+\infty$, $ \frac{\partial \tilde{F}}{\partial \lambda}(x,
\lambda)$ tend vers $F'(\lambda):= -\int_{0}^{+\infty} e^{-\lambda
t} \sin t \,dt$. On peut montrer comme dans la question $4.b)$ que
cette convergence est uniforme pour $\lambda > 0$ (attention il
faut exclure $\lambda = 0$ ici). Il en d\'ecoule que
$\lim_{h\rightarrow 0}\frac{F(\lambda + h) - F(\lambda)}{h} =
F'(\lambda)$ (\'ecrivez l'argument! On pourra soit utiliser le
dernier th\'eor\`eme des rappels, soit le montrer \`a la main...).


  \item  Soit $x > 0$. On a~:
\begin{eqnarray*}
- \int_{0}^{x} e^{-\lambda t} \sin t\,dt &=& \left[
\frac{e^{-\lambda t}}{\lambda} \sin t\right]_{0}^{x} -
\int_{0}^{x} \frac{e^{-\lambda t}}{\lambda} \cos t \,dt\\ & =&
\frac{e^{-\lambda x}}{\lambda} \sin x + \left[ \frac{e^{-\lambda
t}}{\lambda^2} \cos t\right]_{0}^{x} + \int_{0}^{x}
\frac{e^{-\lambda t}}{\lambda^2} \sin t\,dt \\ &=&
\frac{e^{-\lambda x}}{\lambda} \sin x + \frac{e^{-\lambda
x}}{\lambda^2} \cos x - \frac{1}{\lambda^2} + \int_{0}^{x}
\frac{e^{-\lambda t}}{\lambda^2} \sin t\,dt.
\end{eqnarray*}
Ainsi $$-\left(1 + \frac{1}{\lambda^2}\right)\int_{0}^{x}
e^{-\lambda t} \sin t\,dt = \frac{e^{-\lambda x}}{\lambda} \sin x
+ \frac{e^{-\lambda x}}{\lambda^2} \cos x - \frac{1}{\lambda^2}.$$
On en d\'eduit que
$$
F'(\lambda) = \lim_{x\rightarrow +\infty} -\int_{0}^{x}
e^{-\lambda t} \sin t\,dt  = \frac{-1}{1 + \lambda^2}.
$$


  \item De la question pr\'ec\'edente, il d\'ecoule que
$$
F(\lambda ) = - \arctan \lambda + C,
$$
o\`u $C$ est une constante r\'eelle. Montrons que $F(\lambda)$
tend vers $0$ lorsque $\lambda$ tend vers $+\infty$. Soit
$\varepsilon > 0$ et $x> \frac{4}{\varepsilon}$. On a~:
\begin{eqnarray*}
|F(\lambda)| & =  &\left| \int_{0}^{x} e^{-\lambda t} \frac{ \sin
t }{t}\, dt + \int_{x}^{+\infty} e^{-\lambda t} \frac{ \sin t
}{t}\, dt \right|\\
& \leq &  \int_{0}^{x} e^{-\lambda t} \frac{| \sin t| }{t}\, dt +
\left| \int_{x}^{+\infty} e^{-\lambda t} \frac{ \sin t
}{t}\, dt \right|\\
& \leq &  \int_{0}^{x} e^{-\lambda t} \,dt +
\frac{\varepsilon}{2},
\end{eqnarray*}
o\`u on a utiliser que $|\sin t| \leq t$ et la question $4. b)$.
Ainsi
\begin{eqnarray*}
|F(\lambda)| \,\leq \, \frac{1 - e^{-\lambda x_{0}}}{\lambda} +
\frac{\varepsilon}{2}.
\end{eqnarray*}
Comme $\lim_{\lambda \rightarrow +\infty}\frac{1 - e^{-\lambda
x_{0}}}{\lambda} = 0$, il existe ${\lambda}_{0} > 0$ tel que pour
tout $\lambda > {\lambda}_{0}$, $\frac{1 - e^{-\lambda
x_{0}}}{\lambda} < \frac{\varepsilon}{2}$. On en d\'eduit que pour
$\lambda > {\lambda}_{0}$, $|F(\lambda)| < \varepsilon$,
c'est-\`a-dire $\lim_{\lambda \rightarrow +\infty}F(\lambda) = 0$.
Alors $C = \lim_{\lambda \rightarrow +\infty} \arctan \lambda =
\frac{\pi}{2}$. Ainsi
$$
F(\lambda) = -\arctan \lambda + \frac{\pi}{2} \quad \text{et}\quad
\int_{0}^{+\infty} \frac{\sin t}{t}\,dt = F(0) = \frac{\pi}{2}.
$$
  \end{enumerate}
\end{enumerate}
\fincorrection
\nocorrection
\nocorrection
\nocorrection
\nocorrection
\nocorrection
\nocorrection
\nocorrection
\nocorrection
\nocorrection
\correction{004238}
$I_n = \frac\pi{\sqrt 3}(2-\sqrt 3\,)^n$.
\fincorrection
\correction{004239}
Couper en intervalles de $k\pi/n$.
         On obtient $I_n = \frac\pi{\sqrt2}$ pour tout $n\ge 1$.
\fincorrection
\correction{004240}
$f$ est paire, $\pi$-périodique. $f'(x) = 0$ pour $0 \le x \le \frac \pi2$
$ \Rightarrow  f(x) = f(\pi/4) = \frac \pi4$.
\fincorrection
\correction{004244}
Comparaison entre $ \int_{t=0}^1 \frac{d t}{(1+t)^2}$ et
son approximation des trapèzes. Découper et intégrer deux fois par
parties, $u_n\to \frac38$ lorsque $n\to\infty$.
\fincorrection
\nocorrection
\correction{004247}
$I = \Bigl[ f'(t)(1+\cos t) \Bigr]_0^{2\pi} +
      \int_{t=0}^{2\pi} f''(t)(1+\cos t)\,d t \ge 0$.
\fincorrection
\nocorrection
\nocorrection
\nocorrection
\correction{004251}
\begin{enumerate}
  \item formule de Taylor-intégrale.
  \item 
\end{enumerate}
 \fincorrection
\nocorrection
\nocorrection
\correction{004254}
$H' = f(2F-f^2) = fK$ et $K' = 2f(1-f')$ donc $H$ est croissante et positive.
\fincorrection
\nocorrection
\nocorrection
\nocorrection
\nocorrection
\correction{004259}
\begin{enumerate}
  \item 
  \item Soit $F(x) =  \int_{t=a}^{x}f(t)\,d t$ et $G = F^{-1}$.
Alors $\frac1n \sum_{k=0}^{n-1}f(x_k) = \frac 1n \sum_{k=0}^{n-1} f\circ G(\frac kn)
\to  \int_{t=a}^{b}f^2(t)\,d t\Bigm/ \int_{t=a}^{b}f(t)\,d t$ lorsque $n\to\infty$.
\end{enumerate}
 \fincorrection
\correction{004260}
On a $f=e^g$ avec $g$ de classe $\mathcal{C}^1$ par le théorème de relèvement d'où
$I(f) = \frac{g(2\pi)-g(0)}{2i\pi}\in\Z$.
\fincorrection
\correction{004261}
\begin{enumerate}
  \item Il existe toujours une unique fonction $g$ de classe $\mathcal{C}^2$
    telle que $g'' = f$, $g(a)=g'(a)=0$~: $g(x) =  \int_{t=a}^x(x-t)f(t)\,d t$
    (Taylor-Intégral).
  \item Soient $\lambda,\mu\in \R$ tels que $f_1\ :\ x \mapsto f(x) - \lambda-\mu x$
    vérifie $ \int_{x=a}^b f_1(x)\,d x =  \int_{x=a}^b xf_1(x)\,d x = 0$.
    On trouve $$\left\{\begin{array}{ll}(b-a)\lambda +(b^2-a^2)/2\mu &= - \int_{x=a}^bf(x)\,d x\cr
                      (b^2-a^2)/2\lambda + (b^3-a^3)/3\mu &= - \int_{x=a}^bxf(x)\,d x\cr\end{array}\right.$$
    et ce système a pour déterminant $(b-a)^4/12\ne 0$ donc $\lambda,\mu$
    existent et sont uniques. Soit $g_1\in F$ telle que $g_1''=f_1$~:
    $ \int_{x=a}^bg_1''(x)g''(x)\,d x = 0$ pour tout $g\in F$, en particulier
    pour $g=g_1$ donc $g_1''=f_1=0$ et $f(x) = \lambda + \mu x$.
\end{enumerate}
 \fincorrection
\correction{004262}
Soit $g=f-a$. On a $0\le g \le b-a$ et $ \int_0^1 g = -a$ d'où
$ \int_0^1g^2 \le (b-a) \int_0^1g = -a(b-a)$ et 
$ \int_0^1f^2 =  \int_0^1g^2 +2a \int_0^1g + a^2\le -ab$.
\fincorrection
\correction{005445}
\begin{enumerate}
\item  $f$ est continue sur le segment $[0,1]$ et est donc bornée sur ce segment. Soit $M$ un majorant de $|f|$ sur $[0,1]$. Pour $n\in\Nn$,

$$|u_n|\leq\int_{0}^{1}t^n|f(t)|\;dt\leq M\int_{0}^{1}t^n\;dt=\frac{M}{n+1},$$

et comme $\lim_{n\rightarrow +\infty}\frac{M}{n+1}=0$, on a montré que $\lim_{n\rightarrow +\infty}u_n=0$.

Soit $n\in\Nn$. Puisque $f$ est de classe $C^1$ sur $[0,1]$, on peut effectuer une intégration par parties qui fournit

$$u_n=\left[\frac{t^{n+1}}{n+1}f(t)\right]_{0}^{1}-\frac{1}{n+1}\int_{0}^{1}t^{n+1}f'(t)\;dt=\frac{f(1)}{n+1}-\frac{1}{n+1}\int_{0}^{1}t^{n+1}f'(t)\;dt.$$

Puisque $f'$ est continue sur $[0,1]$, $\lim_{n\rightarrow +\infty}\int_{0}^{1}t^{n+1}f'(t)\;dt=0$ ou encore $-\frac{1}{n+1}\int_{0}^{1}t^{n+1}f'(t)\;dt=o(\frac{1}{n})$. D'autre part, puisque $f(1)\neq0$, $\frac{f(1)}{n+1}\sim\frac{f(1)}{n}$ ou encore $\frac{f(1)}{n+1}=\frac{f(1)}{n}+o(\frac{1}{n})$. Finalement, $u_n=\frac{f(1)}{n}+o(\frac{1}{n})$, ou encore

$$u_n\sim\frac{f(1)}{n}.$$

\item  Puisque $f$ est de classe $C^1$ sur $[0,1]$ et que $f(1)=0$, une intégration par parties fournit

$u_n=-\frac{1}{n+1}\int_{0}^{1}t^{n+1}f'(t)\;dt$. Puisque $f'$ est de classe $C^1$ sur $[0,1]$ et que $f'(1)\neq0$, le 1) appliqué à $f'$ fournit

$$u_n=-\frac{1}{n+1}\int_{0}^{1}t^{n+1}f'(t)\;dt\sim-\frac{1}{n}\frac{f'(1)}{n}=-\frac{f'(1)}{n^2}.$$

Par exemple, $\int_{0}^{1}t^n\sin\frac{\pi t}{2}\;dt\sim\frac{1}{n}$ et $\int_{0}^{1}t^n\cos\frac{\pi t}{2}\;dt\sim\frac{\pi}{2n^2}$
\end{enumerate}

\fincorrection
\correction{005448}
\begin{enumerate}
\item  Puisque $f$ est de classe $C^1$ sur $[a,b]$, on peut effectuer une intégration par parties qui fournit pour $\lambda>0$~:

$$\left|\int_{a}^{b}f(t)\sin(\lambda t)\;dt\right|=\left|\frac{1}{\lambda}(-\left[\cos(\lambda t)f(t)\right]_{a}^{b}+\int_{a}^{b}f'(t)\cos(\lambda t)\;dt)\right|\leq\frac{1}{\lambda}(|f(a)|+|f(b)|+\int_{a}^{b}|f'(t)|\;dt).$$ 

Cette dernière expression tend vers $0$ quand $\lambda$ tend vers $+\infty$, et donc $\int_{a}^{b}f(t)\sin(\lambda t)\;dt$ tend vers $0$ quand $\lambda$ tend vers $+\infty$.

\item  Si $f$ est simplement supposée continue par morceaux, on ne peut donc plus effectuer une intégration par parties.

Le résultat est clair si $f=1$, car pour $\lambda>0$, $\left|\int_{a}^{b}\sin(\lambda t)\;dt\right|=...\leq\frac{2}{\lambda}$.

Le résultat s'étend aux fonctions constantes par linéarité de l'intégrale puis aux fonctions constantes par morceaux par additivité par rapport à l'intervalle d'intégration, c'est-à-dire aux fonctions en escaliers.

Soit alors $f$ une fonction continue par morceaux sur $[a,b]$. 

Soit $\varepsilon>0$. On sait qu'il existe une fonction en escaliers $g$ sur $[a,b]$ telle que $\forall x\in[a,b],\;|f(x)-g(x)|<\frac{\varepsilon}{2(b-a)}$.

Pour $\lambda>0$, on a alors 

\begin{align*}\ensuremath
\left|\int_{a}^{b}f(t)\sin(\lambda t)\;dt\right|&=\left|\int_{a}^{b}(f(t)-g(t))\sin(\lambda t)\;dt+\int_{a}^{b}g(t)\sin(\lambda t)\;dt\right|\\
 &\leq\int_{a}^{b}|f(t)-g(t)|\;dt+\left|\int_{a}^{b}g(t)\sin(\lambda  t)\;dt\right|\leq(b-a)\frac{\varepsilon}{2(b-a)}+\left|\int_{a}^{b}g(t)\sin(\lambda t)\;dt\right|\\
 &=\frac{\varepsilon}{2}+\left|\int_{a}^{b}g(t)\sin(\lambda t)\;dt\right|.
\end{align*}

Maintenant, le résultat étant établi pour les fonctions en esacliers, 

$$\exists A>0/\;\forall\lambda\in\Rr,\;
(\lambda>A\Rightarrow\left|\int_{a}^{b}g(t)\sin(\lambda t)\;dt\right|<\frac{\varepsilon}{2}).$$

Pour $\lambda>A$, on a alors $\left|\int_{a}^{b}f(t)\sin(\lambda t)\;dt\right|<\frac{\varepsilon}{2}+\frac{\varepsilon}{2}=\varepsilon$. On a montré que 

$$\forall\varepsilon>0,\exists A>0/\;\forall\lambda\in\Rr,\;(\lambda>A\Rightarrow\left|\int_{a}^{b}f(t)\sin(\lambda t)\;dt\right|<\varepsilon),$$

et donc que $\int_{a}^{b}f(t)\sin(\lambda t)\;dt$ tend vers $0$ quand $\lambda$ tend vers $+\infty$.
\end{enumerate}

\fincorrection
\correction{005452}
Pour $t\in[0,1]$,

\begin{align*}\ensuremath
f^2(t)&=\left(\int_{0}^{t}f'(u)\;du\right)^2\leq(\int_{0}^{t}{f'}^2(u)\;du)(\int_{0}^{t}1\;du)\quad(\textsc{Cauchy}-\textsc{Schwarz})\\
 &=t\int_{0}^{t}{f'}^2(u)\;du\leq t\int_{0}^{1}{f'}^2(u)\;du,
\end{align*}
 
et donc, par croissance de l'intégrale,

$$\int_{0}^{1}f^2(t)\;dt\leq\int_{0}^{1}t(\int_{0}^{1}{f'}^2(u)\;du)\;dt=(\int_{0}^{1}{f'}^2(u)\;du)\int_{0}^{1}t\;dt=\frac{1}{2}\int_{0}^{1}{f'}^2(u)\;du.$$
\fincorrection
\correction{005454}
Puisque $f$ est continue et strictement croissante sur $[0,a]$, $f$ réalise une bijection de $[0,a]$ sur $f([0,a])=[0,f(a)]$.

Soit $x\in[0,a]$. Pour $y\in[0,f(a)]$, posons $g(y)=\int_{0}^{x}f(t)\;dt+\int_{0}^{y}f^{-1}(t)\;dt-xy$. Puisque $f$ est continue sur $[0,a]$, on sait que $f^{-1}$ est continue sur $[0,f(a)]$ et donc la fonction $y\mapsto\int_{0}^{y}f^{-1}(t)\;dt$ est définie et de classe $C^1$ sur $[0,f(a)]$. Donc $g$ est de classe $C^1$ sur 
$[0,f(a)]$ et pour $y\in[0,f(a)]$, $g'(y)=f^{-1}(y)-x$.

Or, $f$ étant strictement croissante sur $[0,a]$, $g'(y)>0\Leftrightarrow f^{-1}(y)>x\Leftrightarrow y>f(x)$. Par suite, $g'$ est strictement négative sur $[0,f(x)[$ et strictement positive sur $]f(x),f(a)]$, et $g$ est strictement décroissante sur $[0,f(x)]$ et strictement croissante sur $[f(x),f(a)]$. $g$ admet en $y=f(x)$ un minimum global égal à
$g(f(x))=\int_{0}^{x}f(t)\;dt+\int_{0}^{f(x)}f^{-1}(t)\;dt-xf(x)$. Notons $h(x)$ cette expression.

$f$ est continue sur $[0,a]$. Donc, $x\mapsto\int_{0}^{x}f(t)\;dt$ est de classe $C^1$ sur $[0,a]$. Ensuite $f$ est de classe $C^1$ sur $[0,a]$ à valeurs dans $[0,f(a)]$ et $y\mapsto\int_{0}^{y}f^{-1}(t)\;dt$ est de classe $C^1$ sur $[0,f(a)]$ (puisque $f^{-1}$ est continue sur $[0,f(a)]$). On en déduit que $x\mapsto\int_{0}^{f(x)}f^{-1}(t)\;dt$ est de classe $C^1$ sur $[0,a]$. Il en est de même de $h$ et pour $x\in[0,a]$,

$$h'(x)=f(x)+f'(x)f^{-1}(f(x))-f(x)-xf'(x)=0.$$

$h$ est donc constante sur $[0,a]$ et pour $x\in[0,a]$, $h(x)=h(0)=0$.

On a montré que 

$$\forall(x,y)\in[0,a]\times[0,f(a)],\;\int_{0}^{x}f(t)\;dt+\int_{0}^{y}f^{-1}(t)\;dt-xy\geq
\int_{0}^{x}f(t)\;dt+\int_{0}^{f(x)}f^{-1}(t)\;dt-xf(x)=0.$$
\fincorrection
\correction{005455}
Soit, pour $x\in[0,1]$, $g(x)=f(x)-x$. g est continue sur $[0,1]$ et 
$$\int_{0}^{1}g(x)\;dx=\int_{0}^{1}f(x)\;dx-\int_{0}^{1}x\;dx=\frac{1}{2}-\frac{1}{2}=0.$$

Si $g$ est de signe constant, $g$ étant de plus continue sur $[0,1]$ et d'intégrale nulle sur $[0,1]$, on sait que $g$ est nulle. Sinon, $g$ change de signe sur $[0,1]$ et le théorème des valeurs intermédiaires montre que $g$ s'annule au moins une fois. Dans tous les cas, $g$ s'annule au moins une fois sur $[0,1]$ ou encore, $f$ admet au moins un point fixe dans $[0,1]$.
\fincorrection
\correction{005456}
D'après l'inégalité de \textsc{Cauchy}-\textsc{Schwarz},

$$\left(\int_{0}^{1}f(t)\;dt\right)\left(\int_{0}^{1}g(t)\;dt\right)=\left(\int_{0}^{1}(\sqrt{f(t)})^2\;dt\right)\left(\int_{0}^{1}(\sqrt{g(t)})^2\;dt\right)\geq\left(\int_{0}^{1}\sqrt{f(t)}\sqrt{g(t)}\;dt\right)^2\geq\left(\int_{0}^{1}1\;dt\right)^2=1.$$
\fincorrection
\correction{005461}
$f$ est continue sur $\Rr$ et admet donc des primitives sur $\Rr$. Soit $F$ une primitive donnée de $f$ sur $\Rr$. Notons $(*)$ la relation~:

$$\forall(x,y)\in\Rr^2,\;f(x)f(y)=F(x+y)-F(x-y).$$

Pour $x=y=0$, on obtient $forall x\in\Rr$, $f(0)=0$. Puis $x=0$ fournit $\forall y\in\Rr,\;F(y)-F(-y)=0$. $F$ est donc nécessairement paire et sa dérivée $f$ est nécessairement impaire.

La fonction nulle est solution du problème. Soit $f$ une éventuelle solution non nulle. Il existe alors un réel $y_0$ tel que $f(y_0)\neq0$. Pour tout réel $x$, on a alors 

$$f(x)=\frac{1}{f(y_0)}\int_{x-y_0}^{x+y_0}f(t)\;dt=\frac{1}{f(y_0)}(F(x+y_0)-F(x-y_0)).$$

$f$ est continue sur $\Rr$ et donc $F$ est de classe $C^1$ sur $\Rr$. Il en est de même de la fonction $x\mapsto\frac{1}{f(y_0)}(F(x+y_0)-F(x-y_0))$ et donc de $f$. Mais alors, $F$ est de classe $C^2$ sur $\Rr$ et donc $f$ l'est aussi ($f$ est en fait de classe $C^\infty$ par récurrence).

En dérivant $(*)$ à $y$ fixé, on obtient $f'(x)f(y)=f(x+y)-f(x-y)$ $(**)$, mais en dérivant à $x$ fixé, on obtient aussi $f(x)f'(y)=f(x+y)+f(x-y)$ $(***)$. En redérivant $(**)$ à $y$ fixé, on obtient $f''(x)f(y)=f'(x+y)-f'(x-y)$ et en dérivant $(***)$ à $x$ fixé, on obtient $f(x)f''(y)=f'(x+y)-f'(x-y)$. Mais alors,

$$\forall(x,y)\in\Rr^2,\;f''(x)f(y)=f(x)f''(y),$$

et en particulier,

$$\forall x\in\Rr,\;f''(x)-\frac{f''(y_0)}{f(y_0)}f(x)=0.$$

On a montré que si $f$ est solution du problème, il existe un réel $\lambda$ tel que $f$ est solution de l'équation différentielle $y''-\lambda y=0$ $(E)$.

\begin{itemize}
\item[- si $\lambda>0$], en posant $k=\sqrt{\lambda}$, $(E)$ s'écrit $y''-k^2y=0$. Les solutions de $(E)$ sont les fonctions de la forme $x\mapsto A\sh(kx)+B\ch(kx)$, $(A,B)\in\Rr^2$ et les solutions impaires de $(E)$ sont les fonctions de la forme $x\mapsto A\sh(kx)$, $A\in\Rr$.

Réciproquement, soit $k$ un réel strictement positif. Pour $A\in\Rr^*$ (on sait que la fonction nulle est solution) et $x\in\Rr$, posons $f(x)=A\sh(kx)$. Alors

$$\int_{x-y}^{x+y}f(t)\;dt=\frac{A}{k}(\ch(k(x+y))-\ch(k(x-y)))\frac{2A}{k}\sh(kx)\sh(ky)=\frac{2}{kA}f(x)f(y).$$

$f$ est solution si et seulement si $\frac{2}{kA}=1$ ou encore $A=\frac{2}{k}$.

\item[- si $\lambda<0$], en posant $k=\sqrt{-\lambda}$, $(E)$ s'écrit $y''+k^2y=0$. Les solutions de $(E)$ sont les fonctions de la forme $x\mapsto A\sin(kx)+B\cos(kx)$, $(A,B)\in\Rr^2$ et les solutions impaires de $(E)$ sont les fonctions de la forme $x\mapsto A\sin(kx)$, $A\in\Rr$.

Réciproquement, soit $k$ un réel strictement positif. Pour $A\in\Rr^*$ et $x\in\Rr$, posons $f(x)=A\sin(kx)$. Alors

$$\int_{x-y}^{x+y}f(t)\;dt=\frac{A}{k}(\cos(k(x-y))-\cos(k(x+y)))=\frac{2A}{k}\sin(kx)\sin(ky)=\frac{2}{kA}f(x)f(y).$$

$f$ est solution si et seulement si $\frac{2}{kA}=1$ ou encore $A=\frac{2}{k}$.

\item[- si $\lambda=0$], $(E)$ s'écrit $y''=0$. Les solutions de $(E)$ sont les fonctions affines et les solutions impaires de $(E)$ sont les fonctions de la forme $x\mapsto Ax$, $A\in\Rr$.

Réciproquement, si $f(x)=Ax$

$$\int_{x-y}^{x+y}f(t)\;dt=\frac{A}{2}((x+y)^2-(x-y)^2)=2Axy=\frac{2}{A}f(x)f(y),$$

et $f$ est solution si et seulement si $A=2$.

\end{itemize}

Les solutions sont la fonction nulle, la fonction $x\mapsto2x$, les fonctions $x\mapsto\frac{2}{k}\sin(kx)$, $k>0$ et les fonctions $x\mapsto\frac{2}{k}\sh(kx)$, $k>0$.
\fincorrection
\correction{005462}
Soit $F$ une primitive de $f$ sur $[a,b]$. $F$ est de classe $C^2$ sur le segment $[a,b]$ et l'inégalité de \textsc{Taylor}-\textsc{Lagrange} permet d'écrire

$$|F(\frac{a+b}{2})-F(a)-\frac{b-a}{2}F'(a)|\leq\frac{1}{2}\frac{(b-a)^2}{4}\sup\{|F''(x)|,\;x\in[a,b]\}.$$

Mais $F'(a)=f(a)=0$ et $F''=f'$. Donc,

$$|F(\frac{a+b}{2})-F(a)|\leq\frac{1}{2}M\frac{(b-a)^2}{4}.$$

De même, puisque $F'(b)=f(b)=0$,

$$|F(\frac{a+b}{2})-F(b)|\leq\frac{1}{2}M\frac{(b-a)^2}{4}.$$

Mais alors,

$$\left|\int_{a}^{b}f(t)\;dt\right|=|F(b)-F(a)|\leq|F(b)-F(\frac{a+b}{2})|+|F(\frac{a+b}{2})-F(a)|\leq\frac{1}{2}M\frac{(b-a)^2}{4}+\frac{1}{2}M\frac{(b-a)^2}{4}=M\frac{(b-a)^2}{4}.$$
\fincorrection
\correction{005463}
Si $\int_{0}^{1}f(t)\;dt\geq0$,

\begin{align*}\ensuremath
\left|\int_{0}^{1}f(t)\;dt\right|=\int_{0}^{1}|f(t)|\;dt&\Leftrightarrow\int_{0}^{1}f(t)\;dt=\int_{0}^{1}|f(t)|\;dt\Leftrightarrow
\int_{0}^{1}(|f(t)|-f(t))\;dt=0\\
 &\Leftrightarrow|f|-f=0\;(\mbox{fonction continue positive d'intégrale nulle})\\
 &\Leftrightarrow f=|f|\Leftrightarrow f\geq0.
\end{align*}

Si $\int_{0}^{1}f(t)\;dt\leq0$, alors $\int_{0}^{1}-f(t)\;dt\geq0$ et d'après ce qui précède, $f$ est solution si et seulement si $-f=|-f|$ ou encore $f\leq0$.

En résumé, $f$ est solution si et seulement si $f$ est de signe constant sur $[0,1]$.
\fincorrection
\nocorrection
\nocorrection
\nocorrection
\nocorrection
\nocorrection
\nocorrection
\nocorrection
\correction{003427}
Ils sont égaux.
\fincorrection
\nocorrection
\correction{003438}
Les deux membres sont $n$-linéaires alternés.
         On le vérifie sur la base du déterminant.
\fincorrection
\correction{005635}
\textbf{1ère solution.} 

\begin{align*}\ensuremath
\text{det}B&=\sum_{\sigma\in S_n}^{}\varepsilon(\sigma)b_{\sigma(1),1}...b_{\sigma(n),n}=\sum_{\sigma\in S_n}^{}\varepsilon(\sigma)(-1)^{1+2+...+n+\sigma(1)+...+\sigma(n)}a_{\sigma(1),1}...a_{\sigma(n),n}\\
 &=\sum_{\sigma\in S_n}^{}\varepsilon(\sigma)(-1)^{2(1+2+...+n)}a_{\sigma(1),1}...a_{\sigma(n),n}=\sum_{\sigma\in S_n}^{}\varepsilon(\sigma)a_{\sigma(1),1}...a_{\sigma(n),n}\\
 &=\text{det}A.
\end{align*}

\textbf{2ème solution.} On multiplie les lignes numéros $2$, $4$,... de $B$ par $-1$ puis les colonnes numéros $2$, $4$,... de la matrice obtenue par $-1$. On obtient la matrice $A$ qui se déduit donc de la matrice $B$ par multiplication des lignes ou des colonnes par un nombre pair de $-1$ (puisqu'il y a autant de lignes portant un numéro pair que de colonnes portant un numéro pair). Par suite, $\text{det}(B)=\text{det}(A)$.
\fincorrection
\correction{005636}
Soient $C\in \mathcal{M}_q(\Kk)$ et $D\in\mathcal{M}_{p,q}(\Kk)$. Soit $\begin{array}[t]{cccc}
\varphi~:&(\mathcal{M}_{p,1}(\Kk))^p&\rightarrow&\Kk\\
 &(C_1,\ldots,C_p)&\mapsto&\text{det}\left(
 \begin{array}{cc}
 X&D\\
 0&C
 \end{array}
 \right)
 \end{array}$ où $X=(C_1\ldots C_p)\in\mathcal{M}_p(\Kk)$.
 

\textbullet~$\varphi$ est linéaire par rapport à chacune des colonnes $C_1$,\ldots, $C_p$.

\textbullet~Si il existe $(i,j)\in\llbracket1,p\rrbracket^2$ tel que $i\neq j$ et $C_i=C_j$, alors $\varphi(C_1,\ldots,C_p)=0$.

Ainsi, $\varphi$ est une forme $p$-linéaire alternée sur l'espace $\mathcal{M}_{p,1}(\Kk)$ qui est de dimension $p$. On sait alors qu'il existe $\lambda\in\Kk$ tel que $\varphi=\lambda\;\text{det}_{\mathcal{B}_0}$ (où $\text{det}_{\mathcal{B}_0}$ désigne la forme déterminant dans la base canonique de $\mathcal{M}_{p,1}(\Kk)$) ou encore il existe $\lambda\in\Kk$ indépendant de $(C_1,\ldots,C_p)$ tel que $\forall (C_1,\ldots,C_p)\in(\mathcal{M}_{p,1}(\Kk))^p$, $f(C_1,\ldots,C_p)=\lambda\;\text{det}_{\mathcal{B}_0}(C_1,\ldots,C_p)$ ou enfin il existe $\lambda\in\Kk$ indépendant de $X$ tel que $\forall X\in\mathcal{M}_{p}(\Kk)$, $\text{det}\left(
 \begin{array}{cc}
 X&D\\
 0&C
 \end{array}
 \right)=\lambda\;\text{det}(X)$. Pour $X=I_p$, on obtient $\lambda=\text{det}\left(
 \begin{array}{cc}
 I_p&D\\
 0&C
 \end{array}
 \right)$ et donc
 
 \begin{center}
 $\forall B\in\mathcal{M}_p(\Kk)$, $\text{det}\left(
 \begin{array}{cc}
 B&D\\
 0&C
 \end{array}
 \right)=\text{det}(B)\times\text{det}\left(
 \begin{array}{cc}
 I_p&D\\
 0&C
 \end{array}
 \right)$.
 \end{center}
 

De même, l'application $Y\mapsto\text{det}\left(
 \begin{array}{cc}
 I_p&D\\
 0&Y
 \end{array}
 \right)$ est une forme $q$-linéaire alternée des lignes de $Y$ et donc il existe $\mu\in\Kk$ tel que $\forall Y\in\mathcal{M}_q(\Kk)$, $\text{det}\left(
 \begin{array}{cc}
 I_p&D\\
 0&Y
 \end{array}
 \right)=\mu\;\text{det}(Y)$ puis $Y=I_q$ fournit $\mu=\text{det}\left(
 \begin{array}{cc}
 I_p&D\\
 0&I_q
 \end{array}
 \right)$ et donc 
 
 \begin{center}
 $\forall B\in\mathcal{M}_p(\Kk)$, $\forall C\in\mathcal{M}_q(\Kk)$, $\forall D\in\mathcal{M}_{p,q}(\Kk)$, $\text{det}\left(
 \begin{array}{cc}
 B&D\\
 0&C
 \end{array}
 \right)=\text{det}(B)\times\text{det}(C)\times\text{det}\left(
 \begin{array}{cc}
 I_p&D\\
 0&I_q
 \end{array}
 \right)=\text{det}(B)\times\text{det}(C)$,
 \end{center}
 

(en supposant acquise la valeur d'un déterminant triangulaire qui peut s'obtenir en revenant à la définition d'un déterminant et indépendamment de tout calcul par blocs).

 \begin{center}
 \shadowbox{
 $\forall (B,C,D)\in\mathcal{M}_p(\Kk)\times\mathcal{M}_q(\Kk)\times\mathcal{M}_{p,q}(\Kk)$, $\text{det}\left(
 \begin{array}{cc}
 B&D\\
 0&C
 \end{array}
 \right)=\text{det}(B)\times\text{det}(C)$.
 }
 \end{center}

\fincorrection
\nocorrection
\nocorrection
\nocorrection
\nocorrection
\nocorrection
\nocorrection
\nocorrection
\nocorrection
\nocorrection
\nocorrection
\nocorrection
\nocorrection
\nocorrection
\nocorrection
\nocorrection
\nocorrection
\nocorrection
\nocorrection
\nocorrection
\nocorrection
\nocorrection
\nocorrection
\nocorrection
\nocorrection
\nocorrection
\nocorrection
\nocorrection
\nocorrection
\nocorrection
\correction{001143}


Commençons par un travail préparatoire : le calcul du déterminant de taille $(n-1)\times (n-1)$ :
{\footnotesize$$\Gamma_k = 
\begin{array}{|cccc|cccc|} 
x  &        &        &   &&&&\\
-1 & x      &        &   &&&&\\
   & \ddots & \ddots &   &&&&\\
   &        & -1     & x &&&&\\ 
   \hline
&&&&  -1 & x      &       &   \\
&&&&     &\ddots & \ddots      &   \\
&&&&     & &\ddots & x  \\
&&&&     &       &     & -1 \\
\end{array}
$$}
où le bloc en haut à gauche est de taille $k\times k$.

On développe, en commençant par la première ligne, puis encore une fois par la première ligne,...
pour trouver que 
$$\Gamma_k = x^k\times (-1)^{n-1-k}$$

Autre méthode : on retrouve le même résultat en utilisant les déterminant par blocs :
$$\begin{array}{|c|c|} 
A & B \\
\hline
(0) & C \\ 
\end{array}
= \det A \times \det C$$

\bigskip


Revenons à l'exercice !

Contrairement à l'habitude on développe par rapport à la colonne qui a le moins de $0$.
En développant par rapport à la dernière colonne on obtient :
{\footnotesize
\begin{align*}
 \Delta_{n}
& =
   \begin{vmatrix}
   x &  0    &        & a_{0}   \\
    -1 &\ddots &\ddots  &\vdots  \\
      &\ddots &x      & a_{n-2} \\
    0 &       & -1      & x+a_{n-1}
   \end{vmatrix} \\
 & = (-1)^{n-1} a_0 
   \begin{vmatrix}
   -1 &  x    &        &    \\
     & -1 & x  &  \\
      & & \ddots & \ddots       \\
     &       &       & -1
   \end{vmatrix}
+ (-1)^{n} a_1   
\begin{vmatrix}
   x & &&&\\
   & -1 &  x    &        &    \\
   &  & -1 & x  &  \\
    &  & & \ddots & \ddots       \\
    & &       &       & -1
   \end{vmatrix} 
\\
& \quad + \cdots +
(-1)^{2n-3} a_{n-2}   
\begin{vmatrix}
  x &      &    & &      \\
    -1 &\ddots &  & & \\
     &\ddots &\ddots & &  \\
      && -1 & x    &   \\
    &&&& -1 \\
   \end{vmatrix}
+ (-1)^{2n-2}(x+a_{n-1})
\begin{vmatrix}
  x &      &    &       \\
    -1 &\ddots &  &  \\
     &\ddots &\ddots &   \\
    &  & -1 & x       \\
   \end{vmatrix} \\
 & = \sum_{k=0}^{n-2} (-1)^{n-1+k} a_ k \times \Gamma_k \quad  + \quad (-1)^{2n-2}(x+a_{n-1})\Gamma_{n-1} \\
 & = \sum_{k=0}^{n-2} (-1)^{n-1+k} a_ k \times x^k\times (-1)^{n-1-k} \quad + \quad (x+a_{n-1})x^{n-1}\\
 & = a_0+a_1x+a_2x^2+\cdots + a_{n-1}x^{n-1} + x^n\\
\end{align*}
}
\fincorrection
\nocorrection
\correction{001145}
\begin{enumerate}
  \item En développant par rapport à la première colonne on trouve la relation suivante :

$$\Delta_n = a \Delta_{n-1} + (-1)^{n-1}(n-1) 
\left\vert
\begin{matrix}
       0   &    0   & \cdots & 0      & n-1 \\
       a   &    0   & \ddots & \vdots & \vdots \\
    \vdots & \ddots & \ddots & 0      & 3 \\
       0   & \cdots &   a    & 0      & 2 \\
       0  & \cdots &   0    & a      & 1
\end{matrix}
\right\vert
$$
Notons $\delta$ ce dernier déterminant (dont la matrice est de taille $n-1\times n-1$). 
On le calcule en développant par rapport à la première ligne 
$$\delta  = (-1)^{n-2}(n-1)
\left\vert
\begin{matrix}
       a   &    0   & \cdots      & 0 \\
       0   &    a   & \ddots      & \vdots \\
    \vdots & \ddots & \ddots      & 0  \\
       0   & \cdots &   0         & a
\end{matrix}\right\vert
= (-1)^{n-2}(n-1) a^{n-2}.
$$
Donc 
$$\Delta_n = a \Delta_{n-1} - a^{n-2}(n-1)^2.$$

  \item Prouvons la formule 
$$\Delta_n=a^n-a^{n-2}\sum_{i=1}^{n-1}{i^2}$$
par récurrence sur $n\ge 2$.

\begin{itemize}
  \item \textbf{Initialisation.} Pour $n=2$,
$\Delta_2=\begin{vmatrix}a&1\\1&a\end{vmatrix}=a^2-1$
donc la formule est vraie.

  \item \textbf{Hérédité.} Supposons la formule vraie vraie au rang $n-1$,
c'est-à-dire $\Delta_{n-1}=a^{n-1}-a^{n-3}\sum_{i=1}^{n-2}{i^2}$.
Calculons $\Delta_n$ :
\begin{align*}
\Delta_n 
  & = a \Delta_{n-1} - a^{n-2}(n-1)^2  \quad \text{ par la première question } \\
  & = a\Big(a^{n-1}-a^{n-3}\sum_{i=1}^{n-2}{i^2} \Big) - a^{n-2}(n-1)^2 \quad \text{ par l'hypothèse de récurrence} \\
  & = a^n -  a^{n-2}\sum_{i=1}^{n-2}{i^2} - a^{n-2}(n-1)^2 \\
  & = a^n-a^{n-2}\sum_{i=1}^{n-1}{i^2}
\end{align*}
La formule est donc vraie au rang $n$.

  \item \textbf{Conclusion.} Par le principe de récurrence la formule est vraie
pour tout entier $n\ge 2$.
\end{itemize}

\end{enumerate}
\fincorrection
\nocorrection
\nocorrection
\nocorrection
\nocorrection
\nocorrection
\nocorrection
\nocorrection
\correction{002453}
Notons $V_n$ le déterminant à calculer
et $C_1,\ldots,C_n$ les colonnes de la matrice correspondante.

Nous allons faire les opérations suivantes sur les colonnes
en partant de la dernière colonne.
$C_n$ est remplacée par $C_n-t_n C_{n-1}$,
puis $C_{n-1}$ est remplacée par $C_{n-1}-t_n C_{n-2}$,...
jusqu'à $C_2$ qui est remplacée par $C_2-t_nC_1$.
On obtient donc

$$V_n=\begin{vmatrix}
1 & t_1 & t_1^2 & \ldots & t_1^{n-1} \\
1 & t_2 & t_2^2 & \ldots & t_2^{n-1} \\\
\ldots&\ldots&\ldots& \ldots & \ldots \\
1 & t_n & t_n^2 & \ldots & t_n^{n-1}
\end{vmatrix} 
= 
\begin{vmatrix}
1 & t_1-t_n & t_1^2-t_1t_n & \ldots & t_1^{n-1}-t_1^{n-2}t_n \\
1 & t_2-t_n & t_2^2-t_2t_n & \ldots & t_2^{n-1}-t_2^{n-2}t_n \\\
\ldots&\ldots&\ldots& \ldots & \ldots \\
1 & 0 & 0 & \ldots & 0
\end{vmatrix}
$$

On développe par rapport à la dernière ligne et on écrit $t_i^k-t_i^{k-1}t_n=t_i^{k-1}(t_i-t_n)$ pour obtenir :
$$V_n = (-1)^{n-1}\begin{vmatrix}
 t_1-t_n & t_1(t_1-t_n) & \ldots & t_1^{n-2}(t_1-t_n) \\
 t_2-t_n & t_2(t_2-t_n) & \ldots & t_2^{n-2}(t_2-t_n) \\\
\ldots&\ldots& \ldots & \ldots \\
 t_{n-1}-t_n & \ldots & \ldots & \ldots
\end{vmatrix}$$

Nous utilisons maintenant la linéarité du déterminant par rapport à chacune des lignes :
on factorise la première ligne par $t_1-t_n$ ; la second par $t_2-t_n$,...
On obtient 
$$V_n = (-1)^{n-1}(t_1-t_n)(t_2-t_n)\cdots(t_{n-1}-t_n)
\begin{vmatrix}
1 & t_1 & t_1^2 & \ldots & t_1^{n-2} \\
1 & t_2 & t_2^2 & \ldots & t_2^{n-2} \\\
\ldots&\ldots&\ldots& \ldots & \ldots \\
1 & t_{n-1} & t_{n-1}^2 & \ldots & t_{n-1}^{n-2}
\end{vmatrix} 
$$
Donc $$V_n = V_{n-1}\prod_{j=1}^{n-1}(t_n-t_j).$$

Si maintenant on suppose la formule connue pour $V_{n-1}$
c'est-à-dire $V_{n-1}(t_1,\ldots,t_{n-1})
= \prod_{1 \le i < j \le n-1} (t_j - t_i)$

Alors on obtient par récurrence que
$$V_n(t_1,\ldots,t_{n-1},t_n)= V_{n-1}(t_1,\ldots,t_{n-1})\prod_{j=1}^{n-1}(t_n-t_j) =  \prod_{1 \le i < j \le n} (t_j - t_i).$$

\fincorrection
\nocorrection
\nocorrection
\nocorrection
\nocorrection
\nocorrection
\correction{002578}
{\it Soient $a,b,c$ des r\'eels v\'erifiant $a^2+b^2+c^2=1$ et
$P$ la matrice r\'eelle $3\times3$ suivante} :
$$P=\begin{pmatrix}a^2&ab&ac \\  ab&b^2&bc \\  ac&bc&c^2\end{pmatrix}$$
\begin{enumerate} 
   \item Calculons le d\'eterminant de $P$.
$$\det P=\begin{vmatrix}a^2&ab&ac \\  ab&b^2&bc \\  ac&bc&c^2\end{vmatrix}=abc\begin{vmatrix}a&a&a \\  b&b&b \\  c&c&c\end{vmatrix}=0.$$
  \item D\'eterminons les sous-espaces vectoriels de $\R^3$, $\ker P$ et $\Im P$.
$$\ker P=\left\{(x,y,z)\in\R^3,\ \begin{pmatrix}a^2&ab&ac \\  ab&b^2&bc \\  ac&bc&c^2\end{pmatrix}
\begin{pmatrix}x \\  y \\  z\end{pmatrix}=\begin{pmatrix}0 \\  0 \\  0\end{pmatrix}\right\},$$
on a 
$$(x,y,z)\in\ker P\iff\left\{\begin{align*}a(ax+by+cz)&=0 \\  b(ax+by+cz)&=0 \\  c(ax+by+cz)&=0 \\ \end{align*}\right.$$
Or, $a,b$ et $c$ ne sont pas simultan\'ement nuls car $a^2+b^2+c^2=1$, ainsi
$$\ker P=\{(x,y,z)\in\R^3,\ ax+by+cz=0\},$$
c'est le plan vectoriel d'\'equation $ax+by+cz=0$.

L'image de $P$ est le sous-espace de $\R^3$ engendr\'e par les vecteurs colonnes de la matrice $P$. Sachant que 
$\dim \ker P+\dim\Im P=\dim \R^3=3$, on sait que la dimension de l'image de $P$ est \'egale \`a $1$, c'est-\`a-dire que l'image est une droite vectorielle. En effet, les vecteurs colonnes de $P$ sont les vecteurs
$$\begin{pmatrix}a^2 \\  ab \\  ac\end{pmatrix},\begin{pmatrix}ab \\  b^2 \\  bc\end{pmatrix}, \begin{pmatrix}ac \\  bc \\  c^2\end{pmatrix}$$
c'est-\`a-dire
$$a\begin{pmatrix}a \\  b \\  c\end{pmatrix},b\begin{pmatrix}a \\  b \\  c\end{pmatrix}, c\begin{pmatrix}a \\  b \\  c\end{pmatrix}.$$
Le sous-espace $\Im P$ est donc la droite vectorielle engendr\'ee par le vecteur $\begin{pmatrix}a \\  b \\  c\end{pmatrix}$.

  \item Soit $Q=I-P$, calculons $P^2$, $PQ$, $QP$ et $Q^2$.
  \begin{align*}
  P^2&=\begin{pmatrix}a^2&ab&ac \\  ab&b^2&bc \\  ac&bc&c^2\end{pmatrix}\begin{pmatrix}a^2&ab&ac \\  ab&b^2&bc \\  ac&bc&c^2\end{pmatrix} \\  &=
\begin{pmatrix}a^4+a^2b^2+a^2c^2 & a^3b+ab^3+abc^2 & a^3c+ab^2c+ac^3  \\  
a^3b+ab^3+abc^2 & a^2b^2+b^4+b^2c^2 & a^2bc+b^3c+bc^3 \\ 
a^3c+ab^2c+ac^3 & a^2bc+b^3c+bc^3 & a^2c^2+b^2c^2+c^4\end{pmatrix}
 \\ &=\begin{pmatrix}a^2(a^2+b^2+c^2)&ab(a^2+b^2+c^2)&ac(a^2+b^2+c^2) \\ 
 ab(a^2+b^2+c^2)&b^2(a^2+b^2+c^2)&bc(a^2+b^2+c^2)
  \\  ac(a^2+b^2+c^2)&bc(a^2+b^2+c^2)&c^2(a^2+b^2+c^2)\end{pmatrix} \\ &=
 \begin{pmatrix}a^2&ab&ac \\  ab&b^2&bc \\  ac&bc&c^2\end{pmatrix}=P.  
 \end{align*}
Car $a^2+b^2+c^2=1$.


Si $Q=I-P$, on a
$$PQ=P(I-P)=PI-P^2=P-P=0,$$
$$QP=(I-P)P=IP-P^2=P-P=0$$ et
$$Q^2=(I-P)(I-P)=I^2-IP-PI+P^2=I-P-P+P=I-P=Q.$$


  \item Caract\'erisons g\'eom\'etriquement $P$ et $Q$.


Nous avons vu que le noyau de $P$ \'etait \'egal au plan vectoriel d'\'equation $ax+by+cz=0$ et que  son image de \'etait la droite vectorielle engendr\'ee par le vecteur $(a,b,c)$. 
Par ailleurs, on a $P^2=P$, \'egalit\'e qui caract\'erise les projecteurs, l'endomorphisme de matrice $P$ est donc la projection sur $\Im P$ suivant la direction $\ker P$.

Soit $X\in \R^3$, on a 
$$QX=0\iff IX-PX=0\iff PX=X\iff X\in \Im P,$$ 
ainsi $\ker Q=\Im P$. D'autre part,
$$Q=I-P=\begin{pmatrix}1-a^2&-ab&-ac \\  -ab&1-b^2&-bc \\  -ac&-bc&1-c^2\end{pmatrix}=
\begin{pmatrix}b^2+c^2&-ab&-ac \\  -ab&a^2+c^2&-bc \\  -ac&-bc&a^2+b^2\end{pmatrix}.$$
On a $\dim\Im Q=2$ et les vecteurs colonnes de $Q$ v\'erifient l'\'equation $ax+by+cz=0$, ainsi $\Im Q=\ker P$. L'\'egalit\'e $Q^2=Q$ prouve que $Q$ est \'egalement un projecteur, c'est la projection sur $\Im Q$ dirig\'ee par $\ker Q$.

  \end{enumerate}
\fincorrection
\correction{002582}
 Soit $a\in\R$ et $A$ la matrice suivante
$$A=\begin{pmatrix}1&0&a \\  0&a&1 \\  a&1&0\end{pmatrix}$$
\begin{enumerate}
 \item Calculons le d\'eterminant de $A$ et d\'eterminons 
pour quelles valeurs de $a$ la matrice est inversible.

$$\det A=\begin{vmatrix}1&0&a \\  0&a&1 \\  a&1&0\end{vmatrix}=\begin{vmatrix}a&1 \\ 1&0\end{vmatrix}+a\begin{vmatrix}0&a \\  a&1\end{vmatrix}=-1-a^3.$$
La matrice $A$ est inversible si et seulement si son d\'eterminant est non nul, c'est-\`a-dire si et seulement si $1+a^3\neq 0$, ce qui \'equivaut \`a $a\neq-1$ car $a\in \R$.
 \item Calculons $A^{-1}$ lorsque $A$ est inversible, c'est-\`a-dire $a\neq-1$. Pour cela nous allons d\'eterminer la comatrice $\tilde A$ de $A$. On a 
$$\tilde A=\begin{pmatrix}-1&a&-a^2 \\  a&-a^2&-1 \\  -a^2&-1&a\end{pmatrix},$$
on remarque que $\tilde A=^t\!\!\!\tilde A$ et on a bien $A ^t\!\!\tilde A=^t\!\!\!\tilde AA=(-1-a^3)I_3$ d'o\`u
$$A^{-1}={\frac{1}{(-1-a^3)}}\tilde A={\frac{1}{-1-a^3}}\begin{pmatrix}-1&a&-a^2 \\  a&-a^2&-1 \\  -a^2&-1&a\end{pmatrix}.$$

\end{enumerate}
\fincorrection
\correction{002753}
\begin{enumerate}
  \item L'aire $\mathcal{A}$ du parallélogramme construit sur les vecteurs $\vec{u} = \left(\begin{array}{c}a\\c\end{array}\right)$ et 
$\vec{v} = \left(\begin{array}{c}b\\d\end{array}\right)$ est la valeur absolue du déterminant $\begin{vmatrix} a & b \\ c & d \end{vmatrix}$
donc $\mathcal{A} =  |ad-bc|$.
Ici on trouve 
$\mathcal{A} = \text{abs} \begin{vmatrix} 2 & 1 \\ 3 & 4 \end{vmatrix} = + 5 $
où $\text{abs}$ désigne la fonction valeur absolue.

  \item Le volume  du  parall\'el\'epip\`ede construit sur trois vecteurs de $\Rr^3$ est la valeur absolue du déterminant
de la matrice formée des trois vecteurs.
Ici $$\mathcal{V} = \text{abs} \begin{vmatrix} 1 & 0 & 1 \\ 2 & 1 & 1 \\ 0 & 3 & 1\end{vmatrix} 
= \text{abs} \Big(+1 \begin{vmatrix} 1 & 1 \\ 3 & 1\end{vmatrix}  + 1  \begin{vmatrix} 2 & 1 \\ 0 & 3 \end{vmatrix}
\Big)= 4$$
où l'on a développé par rapport à la première ligne.

  \item Si un parall\'el\'epip\`ede est construit sur trois vecteurs de $\Rr^3$
dont les coefficients sont des entiers alors le volume correspond au déterminant d'une matrice à coefficients
entiers. C'est donc un entier.

\end{enumerate}
\fincorrection
\nocorrection
\nocorrection
\nocorrection
\nocorrection
\nocorrection
\nocorrection
\nocorrection
 \correction{003451}
 Notation : $\varepsilon_n = \begin{cases}1 &\text{si}n \equiv 0 \text{ ou  }1 (\mathrm{mod}\, 4) \cr
                                   -1 &\text{sinon.}\end{cases}$
 
 \begin{enumerate}
   \item $(b-a)^2(a+b+2x)(a+b-2x)$.
   \item $(a+b+c)^3$
   \item $2abc(a-b)(b-c)(c-a)$
   \item $(a-b)(b-c)(c-a)(ab+ac+bc)$.
   \item $-(a^3-b^3)^2$.
   \item  $\frac {\alpha^{n+1} - \beta^{n+1}}{\alpha - \beta}$
              où $\alpha \ne \beta$ sont les racines de $X^2 - aX + bc = 0$.
              \par
              $(n+1)\left(\frac a2\right)^n$ si $\alpha = \beta$.
   \item $a^{n-3}(a-b)\bigl(a^2 + ab -2(n-2)b^2\bigr)$.
   \item 1.
   \item $a_1 a_2 \dots a_n\left(1+\frac{b_1}{a_1}+ \dots
              + \frac{b_n}{a_n}\right)$.
   \item 0
   \item $\varepsilon_n \frac {n^{n-1}(n+1)}2$.
   \item $(-1)^{n-1}(n-1)2^{n-2}$.
 \end{enumerate}
  \fincorrection
 \correction{003452}
 \begin{enumerate}
   \item $-x(1-x)(2-x)\dots(n-1-x)$.
   \item $(x-a_1) \dots (x-a_n)( x + a_1 + \dots + a_n )$.
   \item $z(y-z)(x-y) \dots (a-b)$. 
   \item $\frac { V(a,b,c)V(x,y,z) }{ (a+x) \dots (c+z) }$.
 \end{enumerate}
  \fincorrection
\correction{003453}
3.$\sin\alpha - \sin\beta - \sin(\alpha-\beta)$.

 \fincorrection
\correction{003454}
\begin{enumerate}
  \item Développer.
  \item 
  \begin{enumerate}
    \item $\begin{cases} D(a-b,0,c-b) = (a-b)^n \cr
                       D(a-c,b-c,0) = (a-c)^n \cr\end{cases}  \Rightarrow 
               D(a,b,c) = \frac { c(a-b)^n - b(a-c)^n } {c-b}$.
    \item $\det( (a-b)I + bU ) = (a-b)^n + nb(a-b)^{n-1}$.
  \end{enumerate}
\end{enumerate}
 \fincorrection
\nocorrection
\nocorrection
\correction{003457}
\begin{enumerate}
  \item $M^2 = \begin{pmatrix}n      &0      &\dots     &0     \cr
                              0      &0      &          &n     \cr
                              \vdots &       &\dots           \cr
                              0      &n      &          &0     \cr \end{pmatrix}
                \Rightarrow  D^2 = \varepsilon_{n-1}n^n$.
  \item
  \item $n^{n/2}\exp\left({i\frac \pi4 (n-1)(3n+2)}\right)$.
\end{enumerate}
Avec la notation : $\varepsilon_n = \begin{cases}1 &\text{si}n \equiv 0 \text{ ou  }1 (\mathrm{mod}\, 4) \cr
                                   -1 &\text{sinon.}\end{cases}$
 \fincorrection
\correction{003458}
Polynômes de Tchebychev
         $ \Rightarrow  D = 2^{(n-1)(n-2)/2} V(\cos\alpha_1, \dots, \cos\alpha_n )$.
\fincorrection
\correction{003459}
$M = ( x_i^{j-1} ) \times ( C_k^{i-1} y_j^{k-i+1} )
  \Rightarrow  \\
  \det M = 
  \begin{cases}0 & \text{ si } k < n-1 \cr
                   
                     \varepsilon_n C_{n-1}^0 C_{n-1}^1 \dots C_{n-1}^{n-1}
		    V(x_1, \dots, x_n ) V(y_1, \dots, y_n )& \text{ si } k = n-1. \end{cases}$
		    
Avec la notation : $\varepsilon_n = \begin{cases}1 &\text{si}n \equiv 0 \text{ ou  }1 (\mathrm{mod}\, 4) \cr
                                   -1 &\text{sinon.}\end{cases}$		    
\fincorrection
\correction{003460}
$A = \left( \frac {i^{j-1}}{(j-1)!} \right)
         \times \Bigl( P^{(i-1)}(j) \Bigr)
          \Rightarrow  \det A = \varepsilon_n \bigl(a_{n-1}(n-1)!\bigr)^n$.
	Avec la notation : $\varepsilon_n = \begin{cases}1 &\text{si}n \equiv 0 \text{ ou  }1 (\mathrm{mod}\, 4) \cr
                                   -1 &\text{sinon.}\end{cases}$  
\fincorrection
\correction{005362}
Soit $(a,b,c)\in\Rr^3$. Notons $\Delta$ le déterminant de l'énoncé. Pour $x$ réel, on pose $D(x)=\left|
\begin{array}{ccc}
-2x&x+b&x+c\\
b+x&-2b&b+c\\
c+x&c+b&-2c
\end{array}\right|$ (de sorte que $\Delta=D(a)$)). $D$ est un polynôme de degré inférieur ou égal à $2$. Le coefficient de $x^2$ vaut
 
$$-(-2c)+(b+c)+(b+c)-(-2b)=4(b+c).$$
Puis,

$$D(-b)=\left|
\begin{array}{ccc}
2b&0&-b+c\\
0&-2b&b+c\\
c-b&c+b&-2c
\end{array}\right|=2b(4bc-(b+c)^2)+2b(c-b)^2=0,$$
et par symétrie des rôles de $b$ et $c$, $D(-c)=0$. De ce qui précède, on déduit que si $b\neq c$, $D(x)=4(b+c)(x+b)(x+c)$ (même si $b+c=0$ car alors $D$ est un polynôme de degré infèrieur ou égal à $1$ admettant au moins deux racines distinctes et est donc le polynôme nul).
Ainsi, si $b\neq c$ (ou par symétrie des roles, si $a\neq b$ ou $a\neq c$), on a~:~$\Delta=4(b+c)(a+b)(a+c)$. Un seul cas n'est pas encore étudié à savoir le cas où $a=b=c$. Dans ce cas, 

$$D(a)=\left|
\begin{array}{ccc}
-2a&2a&2a\\
2a&-2a&2a\\
2a&2a&-2a
\end{array}\right|=8a^3\left|
\begin{array}{ccc}
-1&1&1\\
1&-1&1\\
1&1&-1
\end{array}\right|=32a^3=4(a+a)(a+a)(a+a),$$
ce qui démontre l'identité proposée dans tous les cas (on pouvait aussi conclure en constatant que, pour $a$ et $b$ fixés, la fonction $\Delta$ est une fonction continue de $c$ et on obtient la valeur de $\Delta$ pour $c=b$ en faisant tendre $c$ vers $b$ dans l'expression de $\Delta$ déjà connue pour $c\neq b$).

\begin{center}
\shadowbox{
$\Delta=4(a+b)(a+c)(b+c)$.
}
\end{center}
\fincorrection
\correction{005363}
Soit $P=\left|
\begin{array}{cccc}
X&a&b&c\\
a&X&c&b\\
b&c&X&a\\
c&b&a&X
\end{array}\right|$. $P$ est un polynôme unitaire de degré $4$.
En remplaçant $C_1$ par $C_1+C_2+C_3+C_4$ et par linéarité par rapport à la première colonne, on voit que $P$ est divisible par $(X+a+b+c)$. Mais aussi, en remplaçant $C_1$ par $C_1-C_2-C_3+C_4$ ou $C_1-C_2+C_3-C_4$ ou $C_1+C_2-C_3-C_4$, on voit que $P$ est divisible par $(X-a-b+c)$ ou $(X-a+b-c)$ ou $(X+a-b-c)$.
\textbf{1er cas.} Si les quatre nombres $-a-b-c$, $-a+b+c$, $a-b+c$ et $a+b-c$ sont deux à deux distincts, $P$ est unitaire de degré $4$ et divisible par les quatre facteurs de degré $1$ précédents, ceux-ci étant deux à deux premiers entre eux. Dans ce cas, $P=(X+a+b+c)(X+a+b-c)(X+a-b+c)(X-a+b+c)$.
\textbf{2ème cas.} Deux au moins des quatre nombres $-a-b-c$, $-a+b+c$, $a-b+c$ et $a+b-c$ sont égaux. Notons alors que $-a-b-c=a+b-c\Leftrightarrow b=-a$ et que $-a+b+c=a-b+c\Leftrightarrow a=b$. Par symétrie des roles, deux des quatre nombres $-a-b-c$, $-a+b+c$, $a-b+c$ et $a+b-c$ sont égaux si et seulement si deux des trois nombres $|a|$, $|b|$ ou $|c|$ sont égaux. On conclut dans ce cas que l'expression de $P$ précédemment trouvée reste valable par continuité par rapport à $a$, $b$ ou $c$.

\begin{center}
\shadowbox{
$P=(X+a+b+c)(X+a+b-c)(X+a-b+c)(X-a+b+c)$.
}
\end{center}
\fincorrection
\correction{005364}

\begin{enumerate}
 \item  Pour $n\geq2$, posons $\Delta_n=\left|
\begin{array}{ccccc}
0&1&2&\ldots&n-1\\
1&0&1& &n-2\\
2&1&0&\ddots&\vdots\\
\vdots& &\ddots&\ddots&1\\
n-1&n-2&\ldots&1&0
\end{array}
\right|$. Tout d'abord, on fait apparaître beaucoup de $1$.
Pour cela, on effectue les transformations $C_1\leftarrow C_1-C_2$ puis $C_2\leftarrow C_2-C_3$ puis \ldots puis $C_{n-1}=C_{n-1}-C_n$. On obtient

$$\Delta_n=\mbox{det}(C_1-C_2,C_2-C_3,...,C_{n-1}-C_n,C_n)=\left|
\begin{array}{ccccc}
-1&-1&\ldots&-1&n-1\\
1&-1& &\vdots&n-2\\
1&1&-1&\vdots&\vdots\\
\vdots& &\ddots&-1&1\\
1&1&\ldots&1&0
\end{array}
\right|.$$
On fait alors apparaître un déterminant triangulaire en constatant que $\mbox{det}(L_1,L_2,...,L_n)=\mbox{det}(L_1,L_2+L_1,...,L_{n-1}+L_1,L_n+L_1)$. On obtient

$$\Delta_n=\left|
\begin{array}{ccccc}
-1&\times&\ldots&\ldots&\times\\
0&-2&\ddots& &\vdots\\
0&0&\ddots&\ddots&\vdots\\
\vdots& &\ddots&-2&\times\\
0&\ldots&\ldots&0&n-1
\end{array}
\right|=(1-n)(-2)^{n-2}.$$.

\begin{center}
\shadowbox{
$\forall n\geq2,\;\Delta_n=(1-n)(-2)^{n-2}$.
}
\end{center}
 \item  $\forall(i,j)\in\llbracket1,n\rrbracket^2,\;\sin(a_i+a_j)=\sin a_i\cos a_j+\cos a_i\sin a_j$ et donc si on pose $C=\left(
\begin{array}{c}
\cos a_1\\
\cos a_2\\
\vdots\\
\cos a_n
\end{array}
\right)$ et $S=\left(
\begin{array}{c}
\sin a_1\\
\sin a_2\\
\vdots\\
\sin a_n
\end{array}
\right)$,
on a $\forall j\in\llbracket1,n\rrbracket,\;C_j=\cos a_jS+\sin a_jC$. En particulier, $\mbox{Vect}(C_1,...,C_n)\subset\mbox{Vect}(C,S)$ et le rang de la matrice proposée est inférieur ou égal à $2$. Donc,

\begin{center}
\shadowbox{
$\forall n\geq3,\;\mbox{det}(\sin(a_i+a_j))_{1\leq i,j\leq n}=0.$
}
\end{center}
Si $n=2$, $\mbox{det}(\sin(a_i+a_j))_{1\leq i,j\leq 2}=\sin(2a_1)\sin(2a_2)-\sin^2(a_1+a_2)$.
 \item  L'exercice n'a de sens que si le format $n$ est pair. Posons $n=2p$ où $p$ est un entier naturel non nul.

\begin{align*}\ensuremath
\Delta_n&=\left|\begin{array}{cccccc}
a&0&\ldots&\ldots&0&b\\
0&\ddots&0&0& &0\\
\vdots&0&a&b&0&\vdots\\
\vdots&0&b&a&0&\vdots\\
0& &0&0&\ddots&0\\
b&0&\ldots&\ldots&0&a
\end{array}
\right|
=\left|\begin{array}{cccccc}
a+b&0&\ldots&\ldots&0&b\\
0&\ddots&0&0& &0\\
\vdots&0&a+b&b&0&\vdots\\
\vdots&0&b+a&a&0&\vdots\\
0& &0&0&\ddots&0\\
b+a&0&\ldots&\ldots&0&a
\end{array}
\right|\;(\mbox{pour}\;1\leq j\leq p,\;C_j\leftarrow C_j+C_{2p+1-j})\\
 &=(a+b)^p\left|\begin{array}{cccccc}
1&0&\ldots&\ldots&0&b\\
0&\ddots&0&0& &0\\
\vdots&0&1&b&0&\vdots\\
\vdots&0&1&a&0&\vdots\\
0& &0&0&\ddots&0\\
1&0&\ldots&\ldots&0&a
\end{array}
\right|\;(\mbox{par linéarité par rapport aux colonnes}\;C_1,\;C_2,...,\;C_p)
\end{align*}

$$
=(a+b)^p\left|\begin{array}{cccccc}
1&0&\ldots&\ldots&0&b\\
0&\ddots&0&0& &0\\
\vdots&0&1&b&0&\vdots\\
 & &\ddots&a-b&0&\vdots\\
\vdots& & &\ddots&\ddots&0\\
0&0&\ldots&\ldots&0&a-b
\end{array}
\right| (\mbox{pour}\;p+1\leq i\leq2p,\;L_i\leftarrow L_i-L_{2p+1-i}).
$$
et $\Delta_n=(a+b)^p(a-b)^p=(a^2-b^2)^p$.

\begin{center}
\shadowbox{
$\forall p\in\Nn^*,\;\Delta_{2p}=(a^2-b^2)^p$.
}
\end{center}
 \item  On retranche à la première colonne la somme de toutes les autres et on obtient

$$D_n=\left|\begin{array}{ccccc}
1&1&\ldots& &1\\
1&1&0&\ldots&0\\
\vdots&0&\ddots&\ddots&\vdots\\
\vdots&\vdots&\ddots&1&0\\
1&0&\ldots&0&1
\end{array}
\right|=\left|\begin{array}{ccccc}
-(n-2)&1&\ldots& &1\\
0&1&0&\ldots&0\\
\vdots&\ddots&\ddots&\ddots&\vdots\\
\vdots& &\ddots&1&0\\
0&\ldots&\ldots&0&1
\end{array}
\right|=-(n-2).$$
 \item  Pour $1\leq i\leq p$,

$$L_{i+1}-L_i=(C_{n+i}^0-C_{n+i-1}^0,C_{n+i}^1-C_{n+i-1}^1,...,C_{n+i}^p-C_{n+i-1}^p)=(0,C_{n+i-1}^0,C_{n+i-1}^1,...,C_{n+i-1}^{p-1}).$$
On remplace alors dans cet ordre $L_p$ par $L_p-L_{p-1}$ puis $L_{p-1}$ par $L_{p-1}-L_{p-2}$ puis ... puis $L_2$ par $L_2-L_1$ pour obtenir, avec des notations évidentes

$$\mbox{det}(A_p)=\left|
\begin{array}{cc}
1& \\
0&A_{p-1}
\end{array}\right|=\mbox{det}(A_{p-1}).$$
Par suite, $\mbox{det}(A_p)=\mbox{det}(A_{p-1})=...=\mbox{det}(A_1)=1$.
 \item  En développant suivant la dernière ligne, on obtient~:

$$D_n=(a_{n-1}-X)(-X)^{n-1}+\sum_{k=0}^{n-2}(-1)^{n+k+1}a_k\Delta_k,$$
où $\Delta_k=\left|
\begin{array}{ccc|ccc}
-X&1&0& & & \\
0&\ddots&1& & & \\
0&0&-X& & & \\
\hline
0& &0&1&0&0\\
 & & &-X&\ddots&0\\
0& &0&0&-X&1\\
\end{array}
\right|=(-1)^{k}X^k$ et donc 

\begin{center}
\shadowbox{
$\forall n\geq2,\;D_n=(-1)^n\left(X^n-\sum_{k=0}^{n-1}a_kX^k\right).$
}
\end{center} 
\end{enumerate}
\fincorrection
\correction{005365}
Si deux des $b_j$ sont égaux, $\mbox{det}(A)$ est nul car deux de ses colonnes sont égales. On suppose dorénavant que les $b_j$ sont deux à deux distincts.
Soient $\lambda_1$,..., $\lambda_n$, $n$ nombres complexes tels que $\lambda_n\neq0$. On a

$$\mbox{det}A=\frac{1}{\lambda_n}\mbox{det}(C_1,...,C_{n-1},\sum_{j=1}^{n}\lambda_jC_j)=\mbox{det}B,$$
où la dernière colonne de $B$ est de la forme $(R(a_i))_{1\leq i\leq n}$ avec $R=\sum_{j=1}^{n}\frac{\lambda_j}{X+b_j}$.
On prend $R=\frac{(X-a_1)...(X-a_{n-1})}{(X+b_1)...(X+b_n)}$. $R$ ainsi définie est irréductible (car $\forall(i,j)\in\llbracket1,n\rrbracket^2,\;a_i\neq-b_j$). Les pôles de $R$ sont simples et la partie entière de $R$ est nulle. La décomposition en éléments simples de $R$ a bien la forme espérée.
Pour ce choix de $R$, puisque $R(a_1)=...=R(a_{n-1})=0$, on obtient en développant suivant la dernière colonne 

$$\Delta_n=\frac{1}{\lambda_n}R(a_n)\Delta_{n-1},$$
avec 

$$\lambda_n=\lim_{z\rightarrow -b_n}(z+bn)R(z)=\frac{(-b_n-a_1)...(-b_n-a_{n-1})}{(-b_n+b_1)...(-b_n+b_{n-1})}
=\frac{(a_1+b_n)...(a_{n-1}+b_n)}{(b_n-b_1)...(b_n-b_{n-1})}.$$
Donc 

$$\forall n\geq2,\;\Delta_n=\frac{(a_n-a_1)...(a_{n}-a_{n-1})(b_n-b_1)...(b_n-b_{n-1})}
{(a_n+b_1)(a_n+b_2)...(a_n+b_n)..(a_2+b_n)(a_1+b_n)}\Delta_{n-1}.$$
En réitérant et compte tenu de $\Delta_1=1$, on obtient

\begin{center}
\shadowbox{
$\Delta_n=\frac{\prod_{1\leq i<j\leq n}^{}(a_j-a_i)\prod_{1\leq i<j\leq n}^{}(b_j-b_i)}{\prod_{1\leq i,j\leq n}^{}(a_i+b_j)}=\frac{\mbox{Van}(a_1,...,a_n)\mbox{Van}(b_1,...,b_n)}{\prod_{1\leq i,j\leq n}^{}(a_i+b_j)}.$
}
\end{center}
Dans le cas particulier où $\forall i\in\llbracket1,n\rrbracket,\;a_i=b_i=i$, en notant $H_n$ le déterminant (de \textsc{Hilbert}) à calculer~: $H_n=\frac{\mbox{Van}(1,2,...,n)^2}{\prod_{1\leq i,j\leq n}^{}(i+j)}$. Mais, 

$$\prod_{1\leq i,j\leq n}^{}(i+j)=\prod_{i=1}^{n}\left(\prod_{j=1}^{n}(i+j)\right)=\prod_{i=1}^{n}\frac{(n+i)!}{i!}
=\frac{\prod_{k=1}^{2n}k!}{\left(\prod_{k=1}^{n}k!\right)^2},$$  
et d'autre part,

$$\mbox{Van}(1,2,...,n)=\prod_{1\leq i<j\leq n}^{}(j-i)=\prod_{i=1}^{n-1}\left(\prod_{j=i+1}^{n}(j-i)\right)=\prod_{i=1}^{n-1}(n-i)!=\frac{1}{n!}\prod_{k=1}^{n}k!.$$
Donc,

\begin{center}
\shadowbox{
$\forall n\geq 1,\;H_n=\frac{\left(\prod_{k=1}^{n}k!\right)^3}{n!^2\times\prod_{k=1}^{2n}k!}$.
}
\end{center}

\fincorrection
\correction{005366}
On procède par récurrence sur $n\geq1$.
\textbullet~Pour $n=1$, c'est clair.
\textbullet~Soit $n\geq1$. Supposons que tout déterminant $\Delta_n$ de format $n$ et du type de l'énoncé soit divisible par $2^{n-1}$. Soit $\Delta_{n+1}$ un déterminant de format $n+1$, du type de l'énoncé.
Si tous les coefficients $a_{i,j}$ de $\Delta_{n+1}$ sont égaux à $1$, puisque $n+1\geq2$, $\Delta_{n+1}$ a deux colonnes égales et est donc nul. Dans ce cas, $\Delta_{n+1}$ est bien divisible par $2^n$.
Sinon, on va changer petit à petit tous les $-1$ en $1$.
Soit $(i,j)$ un couple d'indices tel que $a_{i,j}=-1$ et $\Delta_{n+1}'$ le déterminant dont tous les coefficients sont égaux à ceux de $\Delta_{n+1}$ sauf le coefficient ligne $i$ et colonne $j$ qui est égal à $1$.

$$\Delta_{n+1}-\Delta_{n+1}'=\mbox{det}(C_1,...,C_j,...,C_n)-\mbox{det}(C_1,...,C_j',...,C_n)=\mbox{det}(C_1,...,C_j-C_j',...,C_n),$$ 

où  $C_j-C_j'=\left(
\begin{array}{c}
0\\
\vdots\\
0\\
-2\\
0\\
\vdots\\
0
\end{array}
\right)$ ($-2$ en ligne $i$). En développant ce dernier déterminant suivant sa $j$-ème colonne, on obtient~:

$$\Delta_{n+1}-\Delta_{n+1}'=-2\Delta_n,$$ 
où $\Delta_n$ est un déterminant de format $n$ et du type de l'énoncé. Par hypothèse de récurrence, $\Delta_n$ est divisible par $2^{n-1}$ et donc $\Delta_{n+1}-\Delta_{n+1}'$ est divisible par $2^n$. Ainsi, en changeant les $-1$ en $1$ les uns après les autres, on obtient

\begin{center}
$\Delta_{n+1}\equiv\left|
\begin{array}{ccc}
1&\ldots&1\\
\vdots& &\vdots\\
1&\ldots&1
\end{array}\right|\;(\text{mod}\;2^n)$.
\end{center}
Ce dernier déterminant étant nul, le résultat est démontré par récurrence.
\fincorrection
\correction{005369}

\textbf{1ère solution.}

\begin{align*}\ensuremath
\mbox{det}B&=\sum_{\sigma\in S_n}\varepsilon(\sigma)(-1)^{1+\sigma(1)+2+\sigma(2)+...+n+\sigma(n)}a_{\sigma(1),1}a_{\sigma(2),2}...a_{\sigma(n),n}\\
 &=\sum_{\sigma\in S_n}\varepsilon(\sigma)a_{\sigma(1),1}a_{\sigma(2),2}...a_{\sigma(n),n}\;(\mbox{car}\;1+\sigma(1)+2+\sigma(2)+...+n+\sigma(n)=2(1+2+...+n)\in2\Nn)\\
 &=\mbox{det}A
\end{align*}
\textbf{2ème solution.} On multiplie par $-1$ les lignes $2$, $4$, $6$... puis les colonnes $2$, $4$, $6$...On obtient 
$\mbox{det}B=(-1)^{2p}\mbox{det}A=\mbox{det}A$ (où $p$ est le nombre de lignes ou de colonnes portant un numéro pair).
\fincorrection
\correction{005371}
\textbullet~Soit $(k,l)\in\llbracket1,n\rrbracket^2$. Le coefficient ligne $k$, colonne $l$ de $P^2$ est

$$\alpha_{k,l}=\sum_{u=1}^{n}\omega^{(k-1)(u-1)}\omega^{(u-1)(l-1)}=\sum_{u=1}^{n}\omega^{(k+l-2)(u-1)}=\sum_{u=0}^{n-1}(\omega^{k+l-2})^u.$$
 

Or, $\omega^{k+l-2}=1\Leftrightarrow k+l-2\in n\Zz$. Mais, $0\leq k+l-2\leq 2n-2<2n$ et donc, $k+l-2\in n\Zz\Leftrightarrow k+l-2\in\{0,n\}\Leftrightarrow k+l=2\;\mbox{ou}\;k+l=n+2$. Dans ce cas, $\alpha_{k,l}=n$. Sinon, 

$$\alpha_{k,l}=\frac{1-(\omega^{k+l-2})^n}{1-\omega^{k+l-2}}=\frac{1-1}{1-\omega^{k+l-2}}=0.$$
Ainsi, $P^2=n\left(
\begin{array}{ccccc}
1&0&\ldots&\ldots&0\\
0& & &0&1\\
\vdots& &0&1&0\\
 & & & & \\
0&1&0&\ldots&0
\end{array}
\right)$.
\textbullet~Soit $(k,l)\in\llbracket1,n\rrbracket^2$. Le coefficient ligne $k$, colonne $l$ de $P\overline{P}$ est

$$\beta_{k,l}=\sum_{u=1}^{n}\omega^{(k-1)(u-1)}\omega^{-(u-1)(l-1)}=\sum_{u=1}^{n}(\omega^{k-l})^{u-1}.$$
Or, $\omega^{k-l}=1\Leftrightarrow k-l\in n\Zz$. Mais, $-n<-(n-1)\leq k-l\leq n-1<n$ et donc $k-l\in n\Zz\Leftrightarrow k=l$. Dans ce cas, $\beta_{k,l}=n$. Sinon, $\beta_{k,l}=0$. Ainsi, $P\overline{P}=nI_n$ (ce qui montre que $P\in GL_n(\Cc)$ et $P^{-1}=\frac{1}{n}\overline{P}$).
Calculons enfin $PA$. Il faut d'abord écrire proprement les coefficients de $A$. Le coefficient ligne $k$, colonne $l$ de $A$ peut s'écrire $a_{l-k+1}$ si l'on adopte la convention commode $a_{n+1}=a_1$, $a_{n+2}=a_2$ et plus généralement pour tout entier relatif $k$, $a_{n+k}=a_k$.
Avec cette convention d'écriture, le coefficient ligne $k$, colonne $l$ de $PA$ vaut

$$\sum_{u=1}^{n}\omega^{(k-1)(u-1)}a_{l-u+1}=\sum_{v=l-n+1}^{l}\omega^{(k-1)(l-v)}a_v.$$
Puis on réordonne cette somme pour qu'elle commence par $a_1$.

\begin{align*}\ensuremath
\sum_{v=l-n+1}^{l}\omega^{(k-1)(l-v)}a_v&=\sum_{v=1}^{l}\omega^{(k-1)(l-v)}a_v
+\sum_{v=l-n+1}^{0}\omega^{(k-1)(l-v)}a_v\\
 &=\sum_{v=1}^{l}\omega^{(k-1)(l-v)}a_v+\sum_{w=l+1}^{n}\omega^{(k-1)(l-w+n)}a_{w+n}\;(\mbox{en posant}\;w=v+n)\\
 &=\sum_{v=1}^{l}\omega^{(k-1)(l-v)}a_v+\sum_{w=l+1}^{n}\omega^{(k-1)(l-w)}a_{w}\\
 &=\sum_{v=1}^{n}\omega^{(k-1)(l-v)}a_v=\omega^{(k-1)(l-1)}\sum_{v=1}^{n}\omega^{(k-1)(1-v)}a_v
\end{align*}
(le point clé du calcul précédent est que les suites $(a_k)$ et $(\omega^k)$ ont la même période $n$ ce qui s'est traduit par
$\omega^{(k-1)(l-w+n)}a_{w+n}=\omega^{(k-1)(l-v)}a_v$).
Pour $k$ élément de $\llbracket1,n\rrbracket$, posons alors $S_k=\sum_{v=1}^{n}\omega^{(k-1)(1-v)}a_v$. On a montré que $PA=(\omega^{(k-1)(l-1)}S_k)_{1\leq k,l\leq n}$.

Par linéarité par rapport à chaque colonne, on a alors

\begin{center}
$\mbox{det}(PA)=\mbox{det}(\omega^{(k-1)(l-1)}S_k)_{1\leq k,l\leq n}=\left(\prod_{k=1}^{n}S_k\right)\times\mbox{det}(\omega^{(k-1)(l-1)})_{1\leq k,l\leq n}=\left(\prod_{k=1}^{n}S_k\right)\times\mbox{det}P$.
\end{center}
Donc $(\mbox{det}P)(\mbox{det}A)=\left(\prod_{k=1}^{n}S_k\right)\mbox{det}P$ et finalement, puisque $\mbox{det}P\neq0$,

\begin{center}
\shadowbox{
$\mbox{det}A=\prod_{k=1}^{n}\left(\sum_{v=1}^{n}\omega^{(k-1)(1-v)}a_v\right).$
}
\end{center}
Par exemple, pour $n=3$, $\mbox{det}A=(a_1+a_2+a_3)(a_1+ja_2+j^2a_3)(a_1+j^2a_2+ja_3)$.
\fincorrection
\correction{005372}
On a toujours $A\times{^t}\mbox{com}A=(\mbox{det}A)I_n$ et donc 

$$(\mbox{det}A)(\mbox{det}(\mbox{com}A))=(\mbox{det}A)(\mbox{det}(^t\mbox{com}A))=\mbox{det}(\mbox{det}A\;I_n)=(\mbox{det}A)^n.$$
\textbullet~Si $\mbox{det}A\neq0$, on obtient $\mbox{det}(\mbox{com}A)=(\mbox{det}A)^{n-1}$.
\textbullet~Si $\mbox{det}A=0$, alors $A{^t\mbox{com}A}=0$ et $\mbox{com}A$ n'est pas inversible car sinon, $A=0$ puis $\mbox{com}A=0$ ce qui est absurde. Donc, $\mbox{det}(\mbox{com}A)=0$. Ainsi, dans tous les cas,

\begin{center}
\shadowbox{
$\forall A\in M_n(\Cc),\;\mbox{det}(\mbox{com}A)=(\mbox{det}A)^{n-1}$.
}
\end{center}
\textbullet~Si $\mbox{rg}A=n$, alors  $\mbox{com}A\in GL_n(\Kk)$ (car $\mbox{det}(\mbox{com}A)\neq0$) et $\mbox{rg}(\mbox{com}A)=n$.
\textbullet~Si $\mbox{rg}A\leq n-2$, alors tous les mineurs de format $n-1$ sont nuls et $\mbox{com}A=0$. Dans ce cas, $\mbox{rg}(\mbox{com}A)=0$.
\textbullet~Si $\mbox{rg}A=n-1$, il existe un mineur de format $n-1$ non nul et $\mbox{com}A\neq0$. Dans ce cas, $1\leq\mbox{rg}(\mbox{com}A)\leq n-1$.
Plus précisément, 
$$A{^t\mbox{com}A}=0\Rightarrow\mbox{com}A{^tA}=0\Rightarrow\mbox{Im}({^tA})\subset\mbox{Ker}(\mbox{com}A)
\Rightarrow\mbox{dim}(\mbox{Ker}(\mbox{com}A))\geq\mbox{rg}({^tA})=\mbox{rg}A=n-1\Rightarrow\mbox{rg}(\mbox{com}A)\leq1,$$
et finalement si $\mbox{rg}A=n-1$, $\mbox{rg}(\mbox{com}A)=1$.

\fincorrection
\correction{005373}
\begin{align*}\ensuremath
(\mbox{det}A)'&=\left(\sum_{\sigma\in S_n}^{}\varepsilon(\sigma)a_{\sigma(1),1}a_{\sigma(2),2}...a_{\sigma(n),n}\right)'=\sum_{\sigma\in S_n}^{}\varepsilon(\sigma)\left(\sum_{k=1}^{n}a_{\sigma(1),1}...a_{\sigma(k),k}'...a_{\sigma(n),n}\right)\\
 &=\sum_{k=1}^{n}\sum_{\sigma\in S_n}^{}\varepsilon(\sigma)a_{\sigma(1),1}...a_{\sigma(k),k}'...a_{\sigma(n),n}
=\sum_{k=1}^{n}\mbox{det}(C_1,...,C_k',...,C_n)
\end{align*}
\textbf{Applications.}
\begin{enumerate}
 \item  Soit $\Delta_n(x)=\left|\begin{array}{ccccc}
x+1&1&\ldots& &1\\
1&x+1&\ddots& &\vdots\\
\vdots&\ddots&\ddots&\ddots&\vdots\\
\vdots& &\ddots&\ddots&1\\
1&\ldots&\ldots&1&x+1
\end{array}
\right|$. $\Delta_n$ est un polynôme dont la dérivée est d'après ce qui précède, $\Delta_{n}'=\sum_{k=1}^{n}\delta_k$ où $\delta_k$ est le déterminant déduit de $\Delta_n$ en remplaçant sa $k$-ème colonne par le $k$-ème vecteur de la base canonique de $M_{n,1}(\Kk)$. En développant $\delta_k$ par rapport à sa $k$-ème colonne, on obtient $\delta_k=\Delta_{n-1}$ et donc $\Delta_{n}'=n\Delta_{n-1}$.
Ensuite, on a déjà $\Delta_1=X+1$ puis $\Delta_2=(X+1)^2-1=X^2+2X$ ...
Montrons par récurrence que pour $n\geq1$, $\Delta_n=X^n+nX^{n-1}$.
C'est vrai pour $n=1$ puis, si pour $n\geq1$, $\Delta_n=X^n+nX^{n-1}$ alors $\Delta_{n+1}'=(n+1)X^n+(n+1)nX^{n-1}$ et, par intégration,  $\Delta_{n+1}=X^{n+1}+(n+1)X^n+\Delta_{n+1}(0)$. Mais, puisque $n\geq1$, on a $n+1\geq2$ et $\Delta_{n+1}(0)$ est un déterminant ayant au moins deux colonnes identiques. Par suite, $\Delta_{n+1}(0)=0$ ce qui montre que $\Delta_{n+1}=X^{n+1}+(n+1)X^n$. Le résultat est démontré par récurrence.

\begin{center}
\shadowbox{
$\forall n\in\Nn^*,\;\Delta_n=x^n+nx^{n-1}$.
}
\end{center}
 \item  Soit $\Delta_n(x)=\left|\begin{array}{ccccc}
x+a_1&x&\ldots& &x\\
x&x+a_2&\ddots& &\vdots\\
\vdots&\ddots&\ddots&\ddots&\vdots\\
\vdots& &\ddots&\ddots&x\\
x&\ldots&\ldots&x&x+a_n
\end{array}
\right|$. $\Delta_n=\mbox{det}(a_1e_1+xC,...,a_ne_n+xC)$ où $e_k$ est le $k$-ème vecteur de la base canonique de $M_{n,1}(\Kk)$ et $C$ est la colonne dont toutes les composantes sont égales à $1$. Par linéarité par rapport à chaque colonne, $\Delta_n$ est somme de $2^n$ déterminants mais dès que $C$ apparait deux fois, le déterminant correspondant est nul. Donc, $\Delta_n=\mbox{det}(a_1e_1,...,a_ne_n)+\sum_{}^{}\mbox{det}(a_1e_1,...,xC,...,a_ne_n)$. Ceci montre que $\Delta_n$ est un polynôme de degré inférieur ou égal à $1$.
La formule de \textsc{Taylor} fournit alors : $\Delta_n=\Delta_n(0)+X\Delta_{n}'(0)$. Immédiatement, $\Delta_n(0)=\prod_{k=1}^{n}a_k=\sigma_n$ puis $\Delta_{n}'(0)=\sum_{k=1}^{n}\mbox{det}(a_1e_1,...,C,...,a_ne_n)=\sum_{k=1}^{n}\prod_{i\neq k}^{}a_i=\sigma_{n-1}$. Donc, $\Delta_n=\sigma_n+X\sigma_{n-1}$.
\end{enumerate}
\fincorrection
\correction{005374}
\begin{enumerate}
 \item  Pour le premier déterminant, on retranche la première colonne à chacune des autres et on obtient un déterminant triangulaire inférieur dont la valeur est $(-1)^{n-1}$. Pour le deuxième, on ajoute à la première colonne la somme de toutes les autres, puis on met $(n-1)$ en facteurs de la première colonne et on tombe sur le premier déterminant. Le deuxième déterminant vaut donc $(-1)^{n-1}(n-1)$.
 \item  Pour $(i,j)$ élément de $\llbracket1,n\rrbracket^2$, $(i+j-1)^2=j^2+2(i-1)j+(i-1)^2$. Donc, 

$$\forall j\in\{1,...,n\},\;C_j=j^2(1)_{1\leq i\leq n}+2j(i-1)_{1\leq i\leq n}+((i-1)^2)_{1\leq i\leq n}.$$
Les colonnes de la matrice sont donc éléments de $\mbox{Vect}((1)_{1\leq i\leq n},(i-1)_{1\leq i\leq n},((i-1)2)_{1\leq i\leq n})$ qui est de dimension inférieure ou égale à $3$ et la matrice proposée est de rang infèrieur ou égal à $3$. Donc, si $n\geq4$, $\Delta_n=0$. Il reste ensuite à calculer $\Delta_1=1$ puis $\Delta_2=\left|
\begin{array}{cc}
1&4\\
4&9
\end{array}
\right|=-7$ puis $\Delta_3=\left|
\begin{array}{ccc}
1&4&9\\
4&9&16\\
9&16&25
\end{array}
\right|=(225-256)-4(100-144)+9(64-81)=-31+176-153=-8$.
 \item 

$$\Delta_n=\mbox{det}(C_1,...,C_n)=\mbox{det}(C_1+...+C_n,C_2,...,C_n)=(a+(n-1)b)
\left|
\begin{array}{ccccc}
1&b&\ldots&\ldots&b\\
1&a&\ddots& &\vdots\\
\vdots&b&\ddots&\ddots&\vdots\\
\vdots&\vdots&\ddots&\ddots&b\\
1&b&\ldots&b&a
\end{array}
\right|,$$
par linéarité par rapport à la première colonne. Puis, aux lignes numéros $2$,..., $n$, on retranche la première ligne pour obtenir~:

$$\Delta_n=(a+(n-1)b)\left|
\begin{array}{ccccc}
1&b&\ldots&\ldots&b\\
0&a-b&0&\ldots&0\\
\vdots&0&\ddots&\ddots&\vdots\\
\vdots&\vdots&\ddots&\ddots&0\\
0&0&\ldots&0&a-b
\end{array}
\right|=(a+(n-1)b)(a-b)^{n-1}.$$
 \item  Par $n$ linéarité, $D_n$ est somme de $2^n$ déterminants. Mais dans cette somme, un déterminant est nul dès qu'il contient au moins deux colonnes de $x$. Ainsi, en posant $\Delta_n=\mbox{det}(C_1+xC,...,C_n+xC)$ où $C_k=\left(
\begin{array}{c}
a_1\\
\vdots\\
a_{k-1}\\
b\\
a_{k+1}\\
\vdots\\
a_n
\end{array}
\right)$ et $C=(1)_{1\leq i\leq n}$,
on obtient~:

$$\Delta_n=\mbox{det}(C_1,...,C_n)+\sum_{k=1}^{n}\mbox{det}(C_1,...,C_{k-1},xC,C_{k+1},...,C_n),$$
ce qui montre que $\Delta_n$ est un polynôme de degré infèrieur ou égal à $1$. Posons $\Delta_n=AX+B$ et $P=\prod_{k=1}^{n}(a_k-X)$. Quand $x=-b$ ou $x=-c$, le déterminant proposé est triangulaire et se calcule donc immédiatement. Donc~:
\textbf{1er cas.} Si $b\neq c$. $\Delta_n(-b)=P(b)$ et $\Delta_n(-c)=P(c)$ fournit le système $\left\{
\begin{array}{l}
-bA+B=P(b)\\
-cA+B=P(c)
\end{array}
\right.$ et donc $A=-\frac{P(c)-P(b)}{c-b}$ et $B=\frac{cP(b)-bP(c)}{c-b}$. Ainsi,

\begin{center}
\shadowbox{
si $b\neq c$, $\Delta_n=-\frac{P(c)-P(b)}{c-b}x+\frac{cP(b)-bP(c)}{c-b}\;\mbox{où}\;P=\prod_{k=1}^{n}(a_k-X).$
}
\end{center}
\textbf{2ème cas.} Si $b=c$, l'expression obtenue en fixant $x$ et $b$ est clairement une fonction continue de $c$ car polynômiale en $c$. On obtient donc la valeur de $\Delta_n$ quand $b=c$ en faisant tendre $c$ vers $b$ dans l'expression déjà connue de $\Delta_n$ pour $b\neq c$. Maintenant, quand $b$ tend vers $c$, $-\frac{P(c)-P(b)}{c-b}$ tend vers $-P'(b)$ et 
$$\frac{cP(b)-bP(c)}{c-b}=\frac{c(P(b)-P(c))+(c-b)P(c)}{c-b},$$

tend vers $-bP'(b)+P(b)$.

\begin{center}
\shadowbox{
si $b=c$, $\Delta_n=-xP'(b)+P(b)-bP'(b)\;\mbox{où}\;P=\prod_{k=1}^{n}(a_k-X).$
}
\end{center}
 \item  $\Delta_2=3$ et $\Delta_3=\left|
\begin{array}{ccc}
2&1&0\\
1&2&1\\
0&1&2
\end{array}
\right|=2\times3-2=4$. Puis, pour $n\geq4$, on obtient en développant suivant la première colonne~:

$$\Delta_n=2\Delta_{n-1}-\Delta_{n-2}.$$
D'où, pour $n\geq4$, $\Delta_n-\Delta_{n-1}=\Delta_{n-1}-\Delta_{n-2}$ et la suite $(\Delta_n-\Delta_{n-1})_{n\geq3}$ est constante. Par suite, pour $n\geq3$, $\Delta_n-\Delta_{n-1}=\Delta_3-\Delta_2=1$ et donc la suite $(\Delta_n)_{n\geq2}$ est arithmétique de raison $1$. On en déduit que, pour $n\geq2$, $\Delta_n=\Delta_2+(n-2)\times1=n+1$ (on pouvait aussi résoudre l'équation caractéristique de la récurrence double).

\begin{center}
\shadowbox{
$\forall n\geq2,\;\Delta_n=n+1.$
}
\end{center}
\end{enumerate}
\fincorrection
\correction{005376}
$\left|
\begin{array}{ccc}
1&2&3\\
0&1&-2\\
2&1&0
\end{array}
\right|=2+2(-7)=-12\neq0$ et le sytème est de \textsc{Cramer} en $x_1$, $x_2$ et $x_4$. On note aussi que le système est homogène de rang $3$ et donc que l'ensemble des solutions $F$ est un sous-espace vectoriel de $\Rr^5$ de dimension $5-3=2$.

\begin{align*}\ensuremath
\left\{
\begin{array}{l}
x_1+2x_2-x_3+3x_4+x_5=0\\
x_2+x_3-2x_4+2x_5=0\\
2x_1+x_2-5x_3-4x_5=0
\end{array}
\right.&\Leftrightarrow\left\{
\begin{array}{l}
x_1+2x_2+3x_4=x_3-x_5\\
x_2-2x_4=-x_3-2x_5\\
2x_1+x_2=5x_3+4x_5
\end{array}
\right.\Leftrightarrow\left\{
\begin{array}{l}
x_2=-2x_1+5x_3+4x_5\\
x_4=\frac{1}{2}((-2x_1+5x_3+4x_5)+x_3+2x_5)\\
x_1+2x_2+3x_4=x_3-x_5
\end{array}
\right.
\\
 &\Leftrightarrow\left\{
\begin{array}{l}
x_2=-2x_1+5x_3+4x_5\\
x_4=-x_1+3x_3+3x_5\\
x_1+2(-2x_1+5x_3+4x_5)+3(-x_1+3x_3+3x_5)=x_3-x_5
\end{array}
\right.
\\
 &\Leftrightarrow
 \left\{
\begin{array}{l}
x_1=3x_3+3x_5
x_2=-x_3-2x_5\\
x_4=0
\end{array}
\right.
\end{align*}

L'ensemble des solutions est $F=\{(3x_3+3x_5,-x_3-2x_5,x_3,0,x_5),\;(x_3,x_5)\in\Rr^2\}=\mbox{Vect}(e_1,e_2)$ où $e_1=(3,-1,1,0,0)$ et $e_2=(3,-2,0,0,1)$ et, puisque $\mbox{dim}F=2$, une base de $F$ est $(e_1,e_2)$.
\fincorrection
\correction{005637}
Soit $n$ un entier naturel non nul. On note $L_0$, $L_1$,\ldots, $L_n$ les lignes  du déterminant $\text{Van}(x_0,\ldots,x_n)$

A la ligne numéro $n$ du déterminant $\text{Van}(x_0,\ldots,x_n)$,  on ajoute une combinaison linéaire des lignes précédentes du type $L_{n}\leftarrow L_{n}+\sum_{i=0}^{n-1}\lambda_iL_i$. La valeur du déterminant n'a pas changé mais sa dernière ligne s'écrit maintenant $(P(x_0),...,P(x_n))$ où $P$ est un polynôme unitaire de degré $n$. On choisit alors pour
$P$ (le choix des $\lambda_i$ équivaut au choix de $P$) le polynôme $P=\prod_{i=0}^{n-1}(X-x_i)$ (qui est bien unitaire de degré $n$).
La dernière ligne s'écrit alors $(0,...,0,P(x_{n+1}))$ et en développant ce déterminant suivant cette dernière ligne, on obtient la relation de récurrence :

\begin{center}
$\forall n\in\Nn^*,\;\text{Van}(x_0,\ldots,x_n)=P(x_n)\text{Van}(x_0,\ldots,x_{n-1})=\prod_{i=0}^{n-1}(x_n-x_i)\text{Van}(x_0,\ldots,x_{n-1})$.
\end{center}

En tenant compte de $\text{Van}(x_0)=1$, on obtient donc par récurrence

\begin{center}
\shadowbox{
$\forall n\in\Nn^*,\;\forall(x_i)_{0\leqslant i\leqslant n}\in\Kk^n,\;\text{Van}(x_i)_{0\leqslant i\leqslant n-1}=\prod_{0\leqslant i<j\leqslant n-1}^{}(x_j-x_i)$.
}
\end{center}
En particulier, $\text{Van}(x_i)_{0\leqslant i\leqslant n-1}\neq 0$ si et seulement si les $x_i$ sont deux à deux distincts.
\fincorrection
\correction{005638}
Si deux des $a_i$ sont égaux ou deux des $b_j$ sont égaux, $C_n$ est nul car $C_n$ a soit deux lignes identiques, soit deux colonnes identiques.

On suppose dorénavant que les $a_i$ sont deux à deux distincts de même que les $b_j$ (et toujours que les sommes $a_i+b_j$ sont toutes non nulles).

Soit $n\in\Nn^*$. On note $L_1$,\ldots, $L_{n+1}$ les lignes de $C_{n+1}$.

On effectue sur $C_{n+1}$ la transformation $L_{n+1}\leftarrow\sum_{i=1}^{n+1}\lambda_iL_i$ avec $\lambda_{n+1}\neq 0$.

On obtient $C_{n+1}=\frac{1}{\lambda_{n+1}}D_{n+1}$ où $D_{n+1}$ est le  déterminant obtenu en remplaçant la dernière ligne de $C_{n+1}$ par la ligne $(R(b_1),...,R(b_{n+1}))$ avec $R=\sum_{i=1}^{n+1}\frac{\lambda_i}{X+a_i}$. On prend $R=\frac{(X-b_1)\ldots(X-b_n)}{(X+a_1)\ldots(X+a_{n+1})}$. 

\textbullet~Puisque les $-a_i$ sont distincts des $b_j$, $R$ est  irréductible. 

\textbullet~Puisque les $a_i$ sont deux à deux distincts, les pôles de $R$ sont simples.

\textbullet~Puisque $\text{deg}((X-b_1)...(X-b_n))<\text{deg}((X+a_1)...(X+a_{n+1}))$, la partie entière de $R$ est nulle.

$R$ admet donc effectivement une décomposition en éléments simples de la forme $R=\sum_{i=1}^{n+1}\frac{\lambda_i}{X+a_i}$ où $\lambda_{n+1}\neq 0$.

Avec ce choix des $\lambda_i$, la dernière ligne de $D_{n+1}$ s'écrit $(0,...,0,R(b_{n+1}))$ et en développant $D_{n+1}$ suivant sa dernière ligne, on obtient la relation de récurrence : 

\begin{center}
$\forall n\in\Nn^*,\;C_{n+1}=\frac{1}{\lambda_{n+1}}R(b_{n+1})C_n$.
\end{center}

Calculons $\lambda_{n+1}$. Puisque $-a_{n+1}$ est un pôle simple de $R$, 

\begin{center}
$\lambda_{n+1}=\lim{x \rightarrow -a_{n+1}}(x+a_{n+1})R(x)=\frac{(-a_{n+1}-b_1)\ldots(-a_{n+1}-b_n)}{(-a_{n+1}+a_1)\ldots(-a_{n+1}+a_n)}=\frac{(a_{n+1}+b_1)\ldots(a_{n+1}+b_n)}{(a_{n+1}-a_1)\ldots(a_{n+1}-a_n)}$.
\end{center}

On en déduit que

\begin{center}
$\frac{1}{\lambda_{n+1}} R(b_{n+1}) =\frac{(a_{n+1}-a_1)\ldots(a_{n+1}-a_n)}{(a_{n+1}+b_1)\ldots(a_{n+1}+b_n)}\frac{(b_{n+1}-b_1)\ldots(b_{n+1}-b_n)}{(b_{n+1}+a_1)\ldots(b_{n+1}+a_n)}$
\end{center}

puis la relation de récurrence

\begin{center}
$\forall n\in\Nn^*,\;C_{n+1}=\frac{\prod_{i=1}^{n}(a_{n+1}-a_i)\prod_{i=1}^{n}(b_{n+1}-b_i)}{\prod_{i=n+1\;\text{ou}\;j=n+1}(a_{i}+b_j)}C_n$.
\end{center}

En tenant compte de $C_1=\frac{1}{a_1+b_1}$, on obtient par récurrence

\begin{center}
\shadowbox{
$\text{det}\left(\frac{1}{a_i+b_j}\right)_{1\leqslant i,j\leqslant n}=\frac{\prod_{1\leqslant i<j\leqslant n}^{}(a_j-a_i)\prod_{1\leqslant i<j\leqslant n}^{}(b_j-b_i)}{\prod_{1\leqslant i,j\leqslant n}^{}(a_i+b_j)}=\frac{\text{Van}(a_i)_{1\leqslant i\leqslant n}\times\text{Van}(b_j)_{1\leqslant j\leqslant n}}{\prod_{1\leqslant i,j\leqslant n}^{}(a_i+b_j)}$.   
}
\end{center}

(y compris dans les cas particuliers analysés en début d'exercice).

Calcul du déterminant de \textsc{Hilbert}. On est dans le cas particulier où $\forall i\in\llbracket1,n\rrbracket$, $a_i = b_i = i$.
D'abord

\begin{center}
$\text{Van}(1,...,n)=\prod_{j=2}^{n}\left(\prod_{i=1}^{j-1}(j-i)\right)=\prod_{j=2}^{n}(j-1)! =\prod_{j=1}^{n-1}i!$.
\end{center}

Puis $\prod_{1\leqslant i,j\leqslant n}^{}(i+j)=\prod_{i=1}^{n}\left(\prod_{j=1}^{n}(i+j)\right)=\prod_{i=1}^{n}\frac{(i+n)!}{i!}=$   et donc

\begin{center}
\shadowbox{
$\forall n\in\Nn^*,\;H_n=\frac{\left(\prod_{i=1}^{n}i!\right)^4}{n!^2\prod_{i=1}^{2n}i!}$.
}
\end{center}  
\fincorrection
\correction{005640}
On note $C_1$,\ldots, $C_n$ les colonnes du déterminant de l'énoncé puis on pose  $C=(\cos(a_i))_{1\leqslant i\leqslant n}$ et $S=(\sin(a_i))_{1\leqslant i\leqslant n}$.

Pour tout $j\in\llbracket1,n\rrbracket$, $Cj =\sin(a_j)C+\cos(a_j)S$. Ainsi, les colonnes de la matrice proposée sont dans $\text{Vect}(C,S)$ qui est un espace de dimension au plus deux et donc,

\begin{center}
\shadowbox{
si $n\geqslant3$, $\text{det}(\sin(a_i+a_j))_{1\leqslant i,j\leqslant n}= 0$.
}
\end{center}

Si $n=2$, on a $\left|
\begin{array}{cc}
\sin(2a_1)&\sin(a_1+a_2)\\
\sin(a_1+a_2)&\sin(2a_2)
\end{array}
\right|=\sin(2a_1)\sin(2a_2)-\sin^2(a_1+a_2)$.
\fincorrection
\correction{005641}
Soient les vecteurs colonnes $A=(a_i)_{1\leqslant i\leqslant n}$ et $U =(1)_{1\leqslant i\leqslant n}$.

$\forall j\in\llbracket1,n\rrbracket$, $C_j=A +b_jU$. Les colonnes de la matrice proposée sont dans un espace de dimension au plus deux et donc,

\begin{center}
\shadowbox{
si $n\geqslant3$, $\text{det}(a_i+b_j)_{1\leqslant i,j\leqslant n}= 0$.
}
\end{center}

Si $n=2$, on a $\left|
\begin{array}{cc}
a_1+b_1&a_1+b_2\\
a_2+b_1&a_2+b_2
\end{array}
\right|=(a_1+b_1)(a_2+b_2)-(a_1+b_2)(a_2+b_1)=a_1b_2+a_2b_1-a_1b_1-a_2b_2=(a_2-a_1)(b_1-b_2)$.
\fincorrection
\correction{005642}
Pour tout $j\in\llbracket1,n\rrbracket$,

\begin{center}
$C_j =((a+i+j)^2)_{1\leqslant i\leqslant n}= j^2(1)_{1\leqslant i\leqslant n}+ 2(a+j)(i)_{1\leqslant i\leqslant n}+ (i^2)_{1\leqslant i\leqslant n}$.
\end{center}

Les colonnes de la matrice proposée sont dans un espace de dimension au plus trois et donc,

\begin{center}
\shadowbox{
si $n\geqslant 4$, $\text{det}((a+i+j)^2)_{1\leqslant i,j\leqslant n}= 0$.
}
\end{center}

Le calcul est aisé pour $n\in\{1,2,3\}$.
\fincorrection
\correction{005643}
$\frac{x_j-x_i}{j-i}$ est déjà un rationnel strictement positif.

Posons $P_i= 1$ si $i = 1$, et si $i\geqslant2$, $P_i=\frac{X(X-1)\ldots(X-(i-2))}{(i-1)!}$.

Puisque, pour $i\in\llbracket1,n\rrbracket$, $\text{deg}(P_i)=i-1$, on sait que la famille $(P_i)_{1\leqslant i\leqslant n}$ est une base de $\Qq_{n-1}[X]$.
De plus, pour $i\geqslant 2$, $P_i -\frac{X^{i-1}}{(i-1)!}$ est de degré $i-2$ et est donc combinaison linéaire de $P_1$, $P_2$,..., $P_{i-2}$ ou encore, pour $2\leqslant i\leqslant n$, la ligne numéro $i$ du déterminant $\text{det}\left(C_{x_j}^{i-1}\right)_{1\leqslant i,j\leqslant n}$ est somme de la ligne $\left(\frac{x_j^{i-1}}{(i-1)!}\right)_{1\leqslant j\leqslant n}$ et d'une combinaison linéaire des lignes qui la précède. En partant de la dernière ligne et en remontant jusqu'à la deuxième, on retranche la combinaison linéaire correspondante des lignes précedentes sans changer la valeur du déterminant. On obtient par linéarité par rapport à chaque ligne

\begin{center}
$\text{det}\left(C_{x_j}^{i-1}\right)_{1\leqslant i,j\leqslant n}=\frac{1}{\prod_{i=1}^{n}(i-1)!}\text{Van}(x_1,...,x_n)=\frac{\prod_{1\leqslant i<j\leqslant n}^{}(x_j-x_i)}{\prod_{1\leqslant i<j\leqslant n}^{}(j-i)}.$
\end{center}

Finalement,

\begin{center}
\shadowbox{
$\prod_{1\leqslant i<j\leqslant n}^{}\frac{x_j-x_i}{j-i}=\text{det}\left(C_{x_j}^{i-1}\right)_{1\leqslant i,j\leqslant n}\in\Nn^*$.
}
\end{center}
\fincorrection
\correction{005644}
Le coefficient ligne $j$, colonne $k$, $(j,k)\in\llbracket1,n\rrbracket^2$, de la matrice $A$ vaut $a_{k-j}$ avec la convention : si $-(n-1)\leqslant u\leqslant -1$, $a_u = a_{n+u}$.

Le coefficient ligne $j$, colonne $k$, $(j,k)\in\llbracket1,n\rrbracket^2$, de la matrice  $A\Omega$ vaut 

\begin{align*}\ensuremath
\sum_{u=1}^{n}a_{u-j}\omega^{(u-1)(k-1)}&=\sum_{v=-(j-1)}^{n-j}a_v\omega^{(v+j-1)(k-1)}=\sum_{v=-(j-1)}^{-1}a_v\omega^{(v+j-1)(k-1)}+\sum_{v=0}^{n-j}a_v\omega^{(v+j-1)(k-1)}\\
 &=\sum_{v=-(j-1)}^{-1}a_{v+n}\omega^{(v+n+j-1)(k-1)}+\sum_{u=0}^{n-j}a_u\omega^{(u+j-1)(k-1)}\;(\text{car}\;a_{v+n}=a_v\;\text{et}\;\omega^n = 1)\\
 &=\sum_{u=n-j+1}^{n-1}a_{u}\omega^{(u+j-1)(k-1)}+\sum_{u=0}^{n-j}a_u\omega^{(u+j-1)(k-1)}=\sum_{u=0}^{n-1}a_u\omega^{(u+j-1)(k-1)}\\
 &=\omega^{(j-1)(k-1)}\sum_{u=0}^{n-1}a_u\omega^{u(k-1)}.
\end{align*}

Pour $k\in\llbracket1,n\rrbracket$, posons $S_k =\sum_{u=0}^{n-1}a_u\omega^{u(k-1)}$. Le coefficient ligne $j$, colonne $k$ de $A\Omega$ vaut donc $\omega^{(j-1)(k-1)}S_k$. Par passage au détereminant, on en déduit que :

\begin{center}
$\text{det}(A\Omega)=\text{det}\left(\omega^{(j-1)(k-1)}S_k\right)_{1\leqslant j,k\leqslant n}=\left(\prod_{k=1}^{n}S_k\right)\text{det}(\omega^{(j-1)(k-1)})_{1\leqslant j,k\leqslant n}$
\end{center}

($S_k$ est en facteur de la colonne $k$) ou encore $(\text{det}A)(\text{det}\Omega)=\left(\prod_{k=1}^{n}S_k\right)(\text{det}\Omega)$. Enfin, $\Omega$ est la matrice de \textsc{Vandermonde} des racines $n$-èmes de l'unité et est donc inversible puisque celles-ci sont deux à deux distinctes. Par suite $\text{det}\Omega\neq0$ et après simplification on obtient

\begin{center}
\shadowbox{
$\text{det}A=\prod_{k=1}^{n}S_k$ où $S_k=\sum_{u=0}^{n-1}a_u\omega^{u(k-1)}$.
}
\end{center}

Par exemple, $\left|
\begin{array}{ccc}
a&b&c\\
c&a&b\\
b&c&a
\end{array}
\right|=S_1S_2S_3 =(a+b+c)(a+jb+j^2c)(a+j^2b+jc)$ où $j =e^{2i\pi/3}$.

Un calcul bien plus simple sera fourni dans la planche \og Réduction \fg.
\fincorrection
\correction{005645}
\begin{enumerate}
 \item  $d=\sum_{\sigma\in  S_n}^{}\varepsilon(\sigma)a_{\sigma(1),1}...a_{\sigma(n),n}$ est dérivable sur $\Rr$ en tant que combinaison linéaire de produits de fonctions dérivables sur $\Rr$ et de plus

\begin{align*}\ensuremath
d'&=\sum_{\sigma\in  S_n}^{}\varepsilon(\sigma)(a_{\sigma(1),1}...a_{\sigma(n),n})'=\sum_{\sigma\in  S_n}^{}\varepsilon(\sigma)\sum_{i=1}^{n}a_{\sigma(1),1}...a_{\sigma(i),i}'\ldots a_{\sigma(n),n}=\sum_{i=1}^{n}\sum_{\sigma\in  S_n}^{}\varepsilon(\sigma)a_{\sigma(1),1}...a_{\sigma(i),i}'\ldots a_{\sigma(n),n}\\
   &=\sum_{i=1}^{n}\text{det}(C_1,...,C_i',...,C_n)\;(\text{où}\;C_1,...,C_n\;\text{sont les colonnes de la matrice}).
\end{align*}

\item  \textbf{1 ère solution.} D'après ce qui précède, la fonction $d_n$ est dérivable sur $\Rr$ et pour $n\geqslant 2$ et $x$ réel, on a

\begin{align*}\ensuremath
d_n'(x)&=\sum_{i=1}^{n}\left|
\begin{array}{ccccccccc}
x+1&1&\ldots&1&0&1&\ldots&\ldots&1\\
1&\ddots&\ddots&\vdots&\vdots&\vdots\\
\vdots&\ddots&\ddots&1&\vdots\\
 & &\ddots&x+1&0&\vdots\\
\vdots& & &1&1&1\\
 & & &\vdots&0&x+1&\ddots\\
 & & & &\vdots&1&\ddots&\ddots&\vdots\\
\vdots& & &\vdots&\vdots&\vdots&\ddots&\ddots&1\\
1&\ldots&\ldots&1&0&1&\ldots&1&x+1
\end{array}
\right|(\text{la colonne particulière est la colonne}\;i)\\
 &=\sum_{i=1}^{n}d_{n-1}(x)(\text{en développant le}\;i\text{-ème déterminant par rapport à sa}\;i\text{-ème colonne})\\
 &=nd_{n-1}(x).
\end{align*}

En résumé, $\forall n\geqslant 2$, $\forall x\in\Rr$, $d_n(x)=nd_{n-1}(x)$. D'autre part $\forall x\in\Rr$, $d_1(x)=x+1$ et $\forall n\geqslant 2$, $d_n(0) = 0$ (déterminant ayant deux colonnes identiques).

Montrons alors par récurrence que 
$\forall n\geqslant 1$, $\forall x\in\Rr$, $d_n(x) =x^n+nx^{n-1}$.

\textbullet~C'est vrai pour $n=1$.

\textbullet~Soit $n\geqslant1$. Supposons que $\forall n\geqslant 1$, $\forall x\in\Rr$, $d_n(x) =x^n+nx^{n-1}$. Alors, pour $x\in\Rr$,

\begin{center}
$d_{n+1}(x)=d_{n+1}(0)+\int_{0}^{x}d_{n+1}'(t)\;dt=(n+1)\int_{0}^{x}d_n(t\;dt)=x^{n+1}+(n+1)x^n$.
\end{center}

On a montré que

\begin{center}
\shadowbox{
$\forall n\geqslant 1$, $\forall x\in\Rr$, $d_n(x)=x^n+nx^{n-1}$.
}
\end{center}

\textbf{2 ème solution.} $d_n$ est clairement un polynôme de degré $n$ unitaire.
Pour $n\geqslant 2$, puisque dn(0) = 0 et que $d_n'= nd_{n-1}$, $0$ est racine de $d_n$, $d_n'$, ..., $d_n^{(n-2)}$ et est donc racine d'ordre $n-1$ au moins de $d_n$. 
Enfin, $d_n(-n)=0$ car la somme des colonnes du déterminant obtenu est nulle.
Finalement $\forall n\geqslant2$, $\forall x\in\Rr$, $d_n(x) = x^{n-1}(x+n)$ ce qui reste vrai pour $n=1$.

Une variante peut être obtenue avec des connaissances sur la réduction.
\end{enumerate}
\fincorrection
\correction{005646}
On effectue sur la matrice $\left(
\begin{array}{cc}
A&-B\\
B&A
\end{array}
\right)$ les transformations : $\forall j\in\llbracket1,n\rrbracket$, $C_j\leftarrow C_j+iC_{n+j}$ (où $i^2 = -1$) sans modifier la valeur du déterminant. On obtient $\text{det}\left(
\begin{array}{cc}
A&-B\\
B&A
\end{array}
\right)=\text{det}\left(
\begin{array}{cc}
A-iB&-B\\
B+iA&A
\end{array}
\right)$.

Puis en effectuant les transformations : $\forall j\in\llbracket n+1,2n\rrbracket$, $L_j\leftarrow L_j - iL_{j-n}$, on obtient 

\begin{center}
$\text{det}\left(
\begin{array}{cc}
A&-B\\
B&A
\end{array}
\right)=\text{det}\left(
\begin{array}{cc}
A-iB&-B\\
B+iA&A
\end{array}
\right)=\text{det}\left(
\begin{array}{cc}
A-iB&-B\\
0&A+iB
\end{array}
\right)=\text{det}(A+iB)\times\text{det}(A-iB)$.
\end{center}

Comme les matrices $A$ et $B$ sont réelles, $\text{det}(A-iB)=\overline{\text{det}(A+iB)}$ et donc

\begin{center}
$\text{det}\left(
\begin{array}{cc}
A&-B\\
B&A
\end{array}
\right)=|\text{det}(A+iB)|^2\in\Rr^+$.
\end{center}
\fincorrection
\correction{005647}
Si $D$ est inversible, un calcul par blocs fournit
 
 
\begin{center}
$\left(
\begin{array}{cc}
A&B\\
C&D
\end{array}
\right)\left(
\begin{array}{cc}
D&0\\
-C&D^{-1}
\end{array}
\right)=\left(
\begin{array}{cc}
AD-BC&BD^{-1}\\
CD-DC&I
\end{array}
\right)=\left(
\begin{array}{cc}
AD-BC&BD^{-1}\\
0&I
\end{array}
\right)$ (car $C$ et $D$ commutent)
\end{center}

et donc, puisque

\begin{align*}\ensuremath
\text{det}\left(\left(
\begin{array}{cc}
A&B\\
C&D
\end{array}
\right)\left(
\begin{array}{cc}
D&0\\
-C&D^{-1}
\end{array}
\right)\right)&=\text{det}\left(
\begin{array}{cc}
A&B\\
C&D
\end{array}
\right)\text{det}\left(
\begin{array}{cc}
D&0\\
-C&D^{-1}
\end{array}
\right)=\text{det}\left(
\begin{array}{cc}
A&B\\
C&D
\end{array}
\right)\times\text{det}D\times\text{det}D^{-1}\\
 &=\text{det}\left(
\begin{array}{cc}
A&B\\
C&D
\end{array}
\right)
\end{align*}

et que $\text{det}\left(
\begin{array}{cc}
AD-BC&BD^{-1}\\
0&I
\end{array}
\right)=\text{det}(AD-BC)$, on a bien $\text{det}\left(
\begin{array}{cc}
A&B\\
C&D
\end{array}
\right)=\text{det}(AD-BC)$ (si $C$ et $D$ commutent).

Si $D$ n'est pas inversible, $\text{det}(D-xI)$ est un polynôme en $x$ de degré $n$ et donc ne s'annule qu'un nombre fini de fois. Par suite, la matrice $D-xI$ est inversible sauf peut-être pour un nombre fini de valeurs de $x$. D'autre part, pour toute valeur de $x$, les matrices $C$ et $D-xI$ commutent et d'après ce qui précède, pour toutes valeurs de $x$ sauf peut-être pour un nombre fini, on a

\begin{center}
$\text{det}\text{det}\left(
\begin{array}{cc}
A&B\\
C&D
\end{array}
\right)=\text{det}(A(D-xI)-BC)$.
\end{center}

Ces deux expressions sont encore des polynômes en $x$ qui coïncident donc en une infinité de valeurs de $x$ et sont donc égaux. Ces deux polynômes prennent en particulier la même valeur en $0$ et on a montré que

\begin{center}
si $C$ et $D$ commutent, $\text{det}\left(
\begin{array}{cc}
A&B\\
C&D
\end{array}
\right)=\text{det}(AD-BC)$.
\end{center}
\fincorrection
\correction{005649}
En développant suivant la dernière colonne, on obtient

\begin{center}
$\text{det}(A-xI_n)=\left|
\begin{array}{ccccc}
-x&0&\ldots&0&a_0\\
1&\ddots&\ddots&\vdots&a_1\\
0&\ddots&\ddots&0&\vdots\\
\vdots&\ddots&\ddots&-x&a_{n-1}\\
0&\ldots&0&1&a_n-x
\end{array}
\right|= (-x)^n(a_n-x)+\sum_{k=0}^{n-1}(-1)^{n-k+1}a_k\Delta_k$
\end{center}

où $\Delta_k=\left|
\begin{array}{cccccccc}
-x&0&\ldots&0&\times&\ldots&\ldots&\times\\
\times&\ddots&\ddots&\vdots&\vdots& & &\vdots\\
\vdots&\ddots&\ddots&0&\vdots& &  &\vdots\\
\times&\ldots&\times&-x&\times&\ldots&\ldots&\times\\
0&\ldots&\ldots&0&1&\times&\ldots&\times\\
\vdots& & &\vdots&0&\ddots &\ddots&\vdots\\
\vdots& & &\vdots&\vdots&\ddots&\ddots&\times\\
0&\ldots&\ldots&0&0&\ldots&0&1\\
\end{array}
\right|= (-x)^k1^{n-k}=(-x)^k$ (déterminant par blocs)
		

Finalement, 

\begin{align*}\ensuremath
\text{det}(A-xI_n)=(-x)^n(a_n-x) +\sum_{k=0}^{n-1}(-1)^{n-k+1}a_k(-x)^k=(-1)^{n+1}\left(x^{n+1}-\sum_{k=0}^{n}a_kx^k\right).
\end{align*}
\fincorrection
\correction{005650}
\begin{enumerate}
 \item  Sans modifier la valeur de $\text{det}A$, on effectue les transformations :$\forall j\in\llbracket1,n\rrbracket$, $C_j\leftarrow C_j+C_{2n+1-j}$.

On obtient alors par linéarité du déterminant par rapport à chacune des $n$ premières colonnes

\begin{center}
$\text{det}A =(a+b)^p\left|
\begin{array}{cccccccc}
1&0&\ldots&0&0&\ldots&0&b\\
0&\ddots&\ddots&\vdots&\vdots& & &0\\
\vdots&\ddots&\ddots&0&0& & &\vdots\\
0&\ldots&0&1&b&0&\ldots&0\\
0&\ldots&0&1&a&0&\ldots&0\\
\vdots& & &0&0&\ddots&\ddots&\vdots\\
0& & &\vdots&\vdots&\ddots&\ddots&0\\
1&0&\ldots&0&0&\ldots&0&a 
\end{array}
\right|$.
\end{center}

On effectue ensuite les transformations : $\forall i\in\llbracket n+1, 2n\rrbracket$, $L_i\leftarrow L_i - L_{2n+1-i}$ et par linéarité du déterminant par rapport aux $n$ dernières lignes, on obtient

\begin{center}
$\text{det}A = (a+b)^n(a-b)^n= (a^2-b^2)^n$.
\end{center}

\item  Ce déterminant a deux colonnes égales et est donc nul.

\item  On retranche la première colonne à toutes les autres et on obtient un déterminant triangulaire : $D_n=(-1)^{n-1}$.

Pour le deuxième déterminant, on ajoute les $n-1$ dernières colonnes à la première puis on met $n-1$ en facteur de la première colonne et on retombe sur le déterminant précédent. On obtient :
$D_n=(-1)^{n-1}(n-1)$.

\item  On ajoute les $n-1$ dernières colonnes à la première puis on met $a+(n-1)b$ en facteur de la première colonne. On obtient

\begin{center}
$D_n=(a+(n-1)b)\left|
\begin{array}{ccccc}
1&b&\ldots&\ldots&b\\
\vdots&a&\ddots& &\vdots\\
 &b&\ddots&\ddots&\vdots\\
\vdots&\vdots&\ddots&\ddots&b\\
1&b&\ldots&b&a
\end{array}
\right|$.
\end{center}

On retranche ensuite la première ligne à toutes les autres et on obtient

\begin{center}
$D_n=(a+(n-1)b)\left|
\begin{array}{ccccc}
1&b&\ldots&\ldots&b\\
0&a-b&0&\ldots&0\\
\vdots&0&\ddots&\ddots&\vdots\\
\vdots&\vdots&\ddots&\ddots&0\\
0&0&\ldots&0&a-b
\end{array}
\right|=(a+(n-1)b)(a-b)^{n-1}$.
\end{center}

\begin{center}
\shadowbox{
$\left|\begin{array}{cccc}
a&b&\ldots&b\\
b&\ddots&\ddots&\vdots\\
\vdots&\ddots&\ddots&b\\
b&\ldots&b&a
\end{array}
\right|=(a+(n-1)b)(a-b)^{n-1}$.
}
\end{center}
\end{enumerate}
\fincorrection
\correction{006885}
\begin{enumerate}
  \item  Le déterminant de la matrice $\begin{pmatrix} a & b \\ c & d \end{pmatrix}$
est $\begin{vmatrix} a & b \\ c & d \end{vmatrix} =  ad-bc$.
Donc $\begin{vmatrix} 7 & 11 \\ -8 & 4 \end{vmatrix} = 7 \times 4 - 11 \times (-8) = 116$.

  \item Nous allons voir différentes méthodes pour calculer les déterminants.

\textbf{Première méthode.} \emph{Règle de Sarrus.} 
Pour le matrice $3\times 3$ il existe une formule qui permet de calculer directement le déterminant.
$$\begin{vmatrix} 
a_{11} & a_{12} & a_{13} \\   
a_{21} & a_{22} & a_{23} \\  
a_{31} & a_{32} & a_{33} \\ 
  \end{vmatrix}
= a_{11}a_{22}a_{33} + a_{12}a_{23}a_{31} + a_{21}a_{32}a_{13}
- a_{13}a_{22}a_{31} - a_{11} a_{32}a_{23}  - a_{12}a_{21}a_{33}$$
Donc 
$$\begin{vmatrix}
    1 & 0 & 6 \\
    3 & 4 & 15\\
    5 & 6 & 21
  \end{vmatrix}
= 1\times 4 \times 21 + 0 \times 15 \times 5 + 3\times 6 \times 6
- 5\times 4\times 6 -6\times 15 \times 1 -3 \times 0 \times 21 = -18$$


Attention ! La règle de Sarrus ne s'applique qu'aux matrices $3\times 3$.

  \item \textbf{Deuxième méthode.} \emph{Se ramener à une matrice diagonale ou triangulaire.}


Si dans une matrice on change un ligne $L_i$ en $L_i-\lambda L_j$ alors le déterminant reste le même.
Même chose avec les colonnes. 

$$ \begin{array}{l|ccc|}
    _{L_1} & 1 & 0 & 2 \\
    _{L_2} & 3 & 4 & 5\\
    _{L_3} & 5 & 6  & 7
  \end{array}
= \begin{array}{l|ccc|}
    & 1 & 0 & 2  \\
    _{L_2 \leftarrow L_2-3L_1} & 0 & 4 & -1 \\
    _{L_3 \leftarrow L_3-5L_1} & 0 & 6  & -3 \\
  \end{array}
= \begin{array}{l|ccc|}
    & 1 & 0 & 2 \\
    & 0 & 4 & -1 \\
    _{L_3 \leftarrow L_3-\frac32 L_2} & 0 & 0  & -\frac32 \\
  \end{array} = 1\times 4 \times (-\tfrac32) = -6
$$
On a utilisé le fait  que le déterminant d'une matrice diagonale (ou triangulaire) est le produit
des coefficients sur la diagonale.


  \item \textbf{Troisième méthode.} \emph{Développement par rapport à une ligne ou une colonne.}
Nous allons développer par rapport à la deuxième colonne. 
$$\begin{vmatrix}
    1 & 0 & -1 \\
    2 & 3 & 5  \\
    4 & 1 & 3
  \end{vmatrix}
= (-0) \times \begin{vmatrix}
    2 & 5  \\
    4 & 3
  \end{vmatrix}
+ (+3)\times \begin{vmatrix}
    1 & -1 \\
    4 & 3
  \end{vmatrix}
+ (-1) \times \begin{vmatrix}
    1 & -1 \\
    2 & 5  \\
  \end{vmatrix}
= 0 + 3 \times 7 - 1  \times 7 = 14$$


Bien souvent on commence par simplifier la matrice en faisant apparaître un maximum de $0$
par les opérations élémentaires sur les lignes et les colonnes. Puis on développe en choisissant 
la ligne ou la colonne qui a le plus de $0$.


  \item 
On fait apparaître des $0$ sur la première colonne puis on développe par rapport à cette colonne.
$$\Delta = \begin{array}{l|cccc|} 
_{L_1} & 0 & 1 & 2 & 3 \\ _{L_2} & 1 & 2 & 3 & 0 \\ _{L_3} & 2 & 3 & 0 & 1 \\ _{L_4} & 3 & 0 & 1 & 2 \end{array}
=  \begin{array}{l|cccc|}& 0 & 1 & 2 & 3\\  & 1 & 2 & 3 & 0\\ _{L_3 \leftarrow L_3 -2 L_2} & 0 & -1 & -6 & 1\\ _{L_4\leftarrow L_4 -3 L_2} & 0 & -6 & -8 & 2\end{array}
= - \begin{array}{|ccc|}  1 & 2 & 3\\  -1 & -6 & 1\\  -6 & -8 & 2\end{array}
$$
Pour calculer le déterminant $3\times 3$ on fait apparaître des $0$ sur la première colonne, puis on la développe.
$$-\Delta = \begin{array}{l|ccc|} _{L_1} & 1 & 2 & 3\\ _{L_2} & -1 & -6 & 1\\ _{L_3} & -6 & -8 & 2\end{array}
= \begin{array}{l|ccc|}  & 1 & 2 & 3\\ _{L_2\leftarrow L_2+L_1} & 0 & -4 & 4\\ _{L_3\leftarrow L_3+6L_1} & 0 & 4 & 20\end{array}
=1\begin{vmatrix}-4 & 4 \\ 4 & 20\end{vmatrix} = -96
$$
Donc $\Delta=96$.


  \item La matrice a déjà beaucoup de $0$ mais on peut en faire apparaître davantage sur la dernière colonne, puis on développe par rapport à la dernière colonne.
$$
\Delta'= \begin{array}{l|cccc|} 
_{L_1} & 0 & 1 & 1 & 0\\ _{L_2} & 1 & 0 & 0 & 1\\ _{L_3} & 1 & 1 & 0 & 1\\  _{L_4} & 1 & 1 & 1 &0  
\end{array}
= \begin{array}{l|cccc|}  & 0 & 1 & 1 & 0\\  & 1 & 0 & 0 & 1\\ _{L_3 \leftarrow L_3-L_2} & 0 & 1 & 0 & 0 \\  & 1 & 1 & 1 &0  
\end{array}
= \begin{array}{|ccc|}  0 & 1 & 1 \\  0 & 1 & 0  \\   1 & 1 & 1  
\end{array}
$$
On développe ce dernier déterminant par rapport à la première colonne :
$$ \Delta'=\begin{array}{|ccc|}  0 & 1 & 1 \\  0 & 1 & 0  \\   1 & 1 & 1  
\end{array} = 1 \times \begin{vmatrix} 1 & 1 \\ 1 & 0 \end{vmatrix} = -1$$

  \item Toujours la même méthode, on fait apparaître des $0$ sur la première colonne, puis on développe par rapport à cette colonne.
$$
\Delta''= \begin{array}{l|cccc|} 
_{L_1} & 1 & 2 & 1 & 2\\  _{L_2} &1 & 3 & 1 & 3\\  _{L_3} &2 & 1 & 0& 6\\  _{L_4} &1 & 1& 1&7\end{array}
=\begin{array}{l|cccc|}  & 1 & 2 & 1 & 2\\  _{L_2\leftarrow L_2-L_1} &0 & 1 & 0 & 1\\  _{L_3\leftarrow L_3-2L_1} & 0 & -3 & -2 & 2 \\ 
 _{L_4\leftarrow L_4-L_1} & 0 & -1 & 0 & 5 \end{array}
= \begin{array}{|ccc|} 1 & 0 & 1\\-3 & -2 & 2 \\ -1 & 0 & 5 \end{array}
$$
On développe par rapport à la deuxième colonne :
$$\Delta''= -2 \times \begin{vmatrix}  1 & 1 \\ -1 & 5 \end{vmatrix} = -12$$
\end{enumerate}
\fincorrection
\correction{006886}
\begin{enumerate}
  \item Par la règle de Sarrus :
$$\Delta_1 = \begin{vmatrix}
a&b&c\\c&a&b\\b&c&a
\end{vmatrix} 
= a^3+b^3+c^3-3abc.$$

  \item On développe par rapport à la seconde ligne qui ne contient qu'un coefficient non nul et on calcule le déterminant
$3\times 3$ par la règle de Sarrus :
$$\Delta_2 = \begin{vmatrix}
1&0&0&1 \\ 0&1&0&0 \\ 1&0&1&1 \\ 2&3&1&1
\end{vmatrix}
= +1 \begin{vmatrix}
1&0&1 \\ 1&1&1 \\ 2&1&1
\end{vmatrix}
= -1.$$

  \item
$$\Delta_3 = 
\begin{array}{l|cccc|} 
_{L_1} & -1 & 1 & 1 & 1 \\ _{L_2} & 1 & -1 & 1 & 1\\ _{L_3} & 1 & 1 & -1& 1\\ _{L_4} & 1 & 1& 1&-1
\end{array}
= \begin{array}{l|cccc|} 
 & -1 & 1 & 1 & 1 \\ _{L_2\leftarrow L_2+L_1} & 0 & 0 & 2 & 2\\
 _{L_3\leftarrow L_3+L_1} & 0 & 2 & 0 & 2 \\  _{L_4\leftarrow L_4+L_1} &0 & 2 & 2 & 0
\end{array}
$$
On développe par rapport à la première colonne :
$$\Delta_3 = (-1) \times \begin{array}{|ccc|} 
  0 & 2 & 2\\  2 & 0 & 2 \\  2 & 2 & 0
\end{array} = -16$$

  \item Le déterminant est linéaire par rapport à chacune de ses lignes
et aussi chacune de ses colonnes. 
Par exemple les coefficients de la première ligne sont tous des multiples
de $5$ donc 
$$\Delta_4 = 
\begin{vmatrix}
10 & 0 & -5 & 15 \\ -2 & 7 & 3 & 0 \\ 8 & 14 & 0 & 2 \\ 0 & -21 & 1 & -1
\end{vmatrix}
= 5 \times \begin{vmatrix}
2 & 0 & -1 & 3 \\ -2 & 7 & 3 & 0 \\ 8 & 14 & 0 & 2 \\ 0 & -21 & 1 & -1
\end{vmatrix}
$$
On fait la même chose avec la troisième ligne :
$$\Delta_4 = 5 \times 2 \times \begin{vmatrix}
2 & 0 & -1 & 3 \\ -2 & 7 & 3 & 0 \\ 4 & 7 & 0 & 1 \\ 0 & -21 & 1 & -1
\end{vmatrix}
$$
Et enfin les coefficients la première colonne sont des multiples de $2$ et 
ceux de la troisième colonne sont des multiples de $7$ donc :
$$\Delta_4 = 5 \times 2 \times 2 \times\begin{vmatrix}
1 & 0 & -1 & 3 \\ -1 & 7 & 3 & 0 \\ 2 & 7 & 0 & 1 \\ 0 & -21 & 1 & -1
\end{vmatrix}
= 5 \times 2 \times 2 \times 7  \times
\begin{vmatrix}
1 & 0 & -1 & 3 \\ -1 & 1 & 3 & 0 \\ 2 & 1 & 0 & 1 \\ 0 & -3 & 1 & -1
\end{vmatrix}
$$

Les coefficients sont plus raisonnables !
On fait $L_2\leftarrow L_2+L_1$ et $L_3\leftarrow L_3-2L_1$
pour obtenir :
$$\Delta_4 = 140 \times 
\begin{vmatrix}
1 & 0 & -1 & 3 \\ 0 & 1 & 2 & 3 \\ 0 & 1 & 2 & -5 \\ 0 & -3 & 1 & -1
\end{vmatrix}
=140 \times \begin{vmatrix}
 1 & 2 & 3 \\  1 & 2 & -5 \\ -3 & 1 & -1
\end{vmatrix}
= 140 \times 56 = 7840$$


  \item 
$$\Delta_5 =\begin{array}{l|cccc|} 
_{L_1} &a&a&b&0 \\ _{L_2} & a&a&0&b \\ _{L_3} & c&0&a&a \\ _{L_4} &0&c&a&a
\end{array}
= \begin{array}{l|cccc|} 
 &a&a&b&0 \\ _{L_2\leftarrow L_2-L_1} & 0&0&-b&b \\  & c&0&a&a \\ _{L_4\leftarrow L_4-L_3} & -c&c&0&0
\end{array}
$$
On fait ensuite les opérations suivantes sur les colonnes :
$C_2 \leftarrow C_2+C_1$ et $C_3 \leftarrow C_3-C_4$ pour obtenir 
une dernière ligne facile à développer :
$$\Delta_5 
= \begin{array}{|cccc|} 
 a&2a&b&0 \\ 0&0&-2b&b \\  c&c&0&a \\ -c&0&0&0
\end{array}
= +c \times \begin{array}{|ccc|} 
 2a&b&0 \\ 0&-2b&b \\  c&0&a \\
 \end{array} = bc(bc-4a^2)
$$
  \item
On fait d'abord les opérations $C_1 \leftarrow C_1-C_3$ et $C_2 \leftarrow C_2-C_4$
et on développe par rapport à la première ligne :
$$\Delta_6 = \begin{vmatrix}
1&0&3&0&0 \\ 0&1&0&3&0 \\ a&0&a&0&3 \\ b&a&0&a&0 \\ 0&b&0&0&a  
\end{vmatrix}
=
\begin{vmatrix}
-2&0&3&0&0 \\ 0&-2&0&3&0 \\ 0&0&a&0&3 \\ b&0&0&a&0 \\ 0&b&0&0&a  
\end{vmatrix}
= (-2) \times \begin{vmatrix}
-2&0&3&0 \\ 0&a&0&3 \\ 0&0&a&0 \\ b&0&0&a  
\end{vmatrix}
+ 3  \times \begin{vmatrix}
0&-2&3&0 \\ 0&0&0&3 \\ b&0&a&0 \\ 0&b&0&a  
\end{vmatrix}
$$
Le premier déterminant à calculer se développe par rapport à la deuxième colonne et le second déterminant
par rapport à la première colonne :
$$\Delta_6 = (-2)\times a \times 
\begin{vmatrix}
-2&3&0 \\  0&a&0 \\ b&0&a  
\end{vmatrix}
+ 3  \times b \times 
\begin{vmatrix}
-2&3&0 \\ 0&0&3 \\ b&0&a  
\end{vmatrix}=4a^3+27b^2$$

  \item Nous allons permuter des lignes et des colonnes pour se ramener à une matrice diagonale par blocs.
Souvenons-nous que lorsque l'on échange deux lignes (ou deux colonnes) alors le déterminant change de signe.
Nous allons rassembler les zéros.
On commence par échanger les colonnes $C_1$ et $C_3$ : $C_1\leftrightarrow C_3$ :
$$\Delta_7=
\begin{vmatrix}
1&0&0&1&0 \\ 0&-4&3&0&0 \\ -3&0&0&-3&-2 \\ 0&1&7&0&0 \\ 4&0&0&7&1  
\end{vmatrix}=
- \begin{vmatrix}
0&0&1&1&0 \\ 3&-4&0&0&0 \\ 0&0&-3&-3&-2 \\ 7&1&0&0&0 \\ 0&0&4&7&1  
\end{vmatrix}$$
Puis on échange les lignes $L_1$ et $L_4$ : $L_1\leftrightarrow L_4$ :
$$\Delta_7=+\begin{array}{|ccccc|} 
7&1&0&0&0 \\ 3&-4&0&0&0 \\  0&0&-3&-3&-2 \\ 0&0&1&1&0 \\ 0&0&4&7&1  
\end{array}$$
Notre matrice est sous la forme d'une matrice diagonale par blocs
 et son déterminant est le produit des déterminants.
$$\Delta_7=\begin{array}{|cc|ccc|} 
7&1&0&0&0 \\ 3&-4&0&0&0 \\ \hline 0&0&-3&-3&-2 \\ 0&0&1&1&0 \\ 0&0&4&7&1  
\end{array}
= \begin{vmatrix}
 7&1 \\  3&-4 \\  
  \end{vmatrix} \times
\begin{vmatrix}
 -3&-3&-2 \\  1&1&0 \\  4&7&1  
  \end{vmatrix}
= (-31)\times (-6) = 186
$$

\end{enumerate}
\fincorrection
\correction{006887}
\begin{enumerate}
  \item On retire la première colonne à toutes les autres colonnes 
$$
\Delta_1 =   \begin{vmatrix}
     a_1   &a_2   &\cdots&a_n    \\
     a_1   &a_1   &\ddots&\vdots \\
     \vdots&\ddots&\ddots&a_2    \\
     a_1   &\cdots&a_1   &a_1
   \end{vmatrix}
 =  \begin{vmatrix}
     a_1   &a_2-a_1  & a_3-a_1 &\cdots&a_n -a_1   \\
     a_1   &0  & &\ddots&\vdots \\
     \vdots&\vdots &\ddots&\ddots&a_2-a_1    \\
     a_1   &0&\cdots&0   & 0
   \end{vmatrix}$$
On développe par rapport à la dernière ligne :
$$\Delta_1 = (-1)^{n-1}a_1 \begin{vmatrix}
     a_2-a_1  &  &\cdots&a_n -a_1   \\
     0  &\ddots &&\vdots \\
     \vdots &\ddots&\ddots&    \\
     0&\cdots&0   & a_2-a_1
   \end{vmatrix}
= (-1)^{n-1}a_1(a_2-a_1)^{n-1}
$$
Où l'on a reconnu le déterminant d'un matrice triangulaire supérieure.
Donc 
$$\Delta_1 = a_1(a_1-a_2)^{n-1}.$$


  \item 
On va transformer la matrice correspondante en une matrice triangulaire supérieure,
on commence par remplacer la ligne $L_2$ par $L_2-L_1$ (on ne note que les coefficients non nuls) : 
$$\Delta_2 =   \begin{vmatrix}
  1 &      &    &   &+1 \\
  1 &1      && &  \\
   &1      &1  & &\\
   & &\ddots &\ddots & \\
  &  &       &1      &1
\end{vmatrix}=
 \begin{vmatrix}
  1 &      &    &   &+1 \\
  0 &1      && &  -1\\
   &1      &1  & &\\
   & &\ddots &\ddots & \\
  &  &       &1      &1
\end{vmatrix}
$$
Puis on remplace la ligne $L_3$ par $L_3-L_2$ (attention il s'agit de la nouvelle ligne $L_2$) et on 
continue ainsi de suite jusqu'à $L_{n-1} \leftarrow L_{n-1}-L_{n-2}$ ($n$ est la taille de la matrice sous-jacente) :
$$\Delta_2 =   \begin{vmatrix}
  1 &      &    &   &&+1 \\
  0 &1      && &  &-1\\
    & 0      &1  & & &+1\\
    & & 1      &1  & && \\
   & &&\ddots &\ddots & \\
  &  &       &&1      &1                 
               \end{vmatrix}
= \cdots = 
\begin{vmatrix}
  1 &      &    &   &&+1 \\
  0 &1      && &  &-1\\
    & 0      &1  & & &+1\\
    & & \ddots      &\ddots  & &\vdots& \\
   & &&0 &1 & (-1)^{n} \\
  &  &       &&1      &1                 
               \end{vmatrix}
$$
On fait attention pour le dernier remplacement $L_n \leftarrow L_n-L_{n-1}$ légèrement différent et qui conduit au déterminant d'une matrice triangulaire : :
$$\Delta_2 = \begin{vmatrix}
  1 &      &    &   &&+1 \\
  0 &1      && &  &-1\\
    & 0      &1  & & &+1\\
    & & \ddots      &\ddots  & &\vdots& \\
   & && &1 & (-1)^{n}\\
  &  &       &&0      &1 - (-1)^{n}                 
               \end{vmatrix}
= 1 - (-1)^{n}.
$$

En conclusion $\Delta_2 = \begin{cases}
                           0 & \text{ si $n$ est pair} \\
                           2 & \text{ si $n$ est impair} \\                           
                          \end{cases}$

  \item
On retire la colonne $C_1$ aux autres colonnes $C_i$ pour faire apparaître des $0$ :
$$\Delta_3 = \begin{vmatrix}
    a+b  &    a   & \cdots &  a       \\
     a   &   a+b  & \ddots & \vdots   \\
  \vdots & \ddots & \ddots &  a       \\
     a   & \cdots &    a   & a+b 
\end{vmatrix}
= \begin{vmatrix}
    a+b  &    -b   & & \cdots &  -b       \\
     a   &   b     &   0 & \cdots        &  0   \\
     a   &   0 &  \ddots & \ddots & \vdots   \\
  \vdots & \ddots & \ddots &  b & 0       \\
     a   &  0 & \cdots &    0   & b 
\end{vmatrix}$$
On remplace ensuite $L_1$ par $L_1+L_2+L_3+\cdots +L_n$
(ou ce qui revient au même : faites les opérations 
$L_1 \leftarrow L_1+L_2$ puis $L_1 \leftarrow L_1+L_3$,\ldots
chacune de ces opérations fait apparaître un $0$ sur la première ligne)
pour obtenir une matrice triangulaire inférieure :
$$\Delta_3 = \begin{vmatrix}
   na+b  &    0   & & \cdots &  0       \\
     a   &   b     &   0 & \cdots        &  0   \\
     a   &   0 &  \ddots & \ddots & \vdots   \\
  \vdots & \ddots & \ddots &  b & 0       \\
     a   &  0 & \cdots &    0   & b 
\end{vmatrix}=(na+b)b^{n-1}.$$
\end{enumerate}
\fincorrection
\correction{001163}
\begin{enumerate}
  \item  Remarquons que comme le système est homogène (c'est-à-dire les coefficients du second membre sont
nuls) alors $(0,0,0)$ est une solution du système. Voyons s'il y en a d'autres. 
Nous faisons semblant de ne pas voir que la seconde ligne implique $x=y$ et 
que le système est en fait très simple à résoudre. 
Nous allons appliquer le pivot de Gauss en faisant les opérations suivantes sur les lignes 
$L_2 \leftarrow L_2-L_1$ et $L_3 \leftarrow L_3-L_1$ :
$$  \left\{
    \begin{array}{rcrcrcl}
         x &+&    y &-&    z &=&  0  \\
         x &-&    y & &      &=&  0  \\
         x &+&    4y &+&    z &=&  0  
    \end{array}
  \right.
\iff\left\{
    \begin{array}{rcrcrcl}
         x &+&    y &-&    z &=&  0  \\
           &-&   2y &+&    z &=&  0  \\
           & &   3y &+&    2z &=&  0  
    \end{array}
  \right.
$$
On fait maintenant $L_3 \leftarrow 2L_3+3L_2$ pour obtenir :
$$\left\{
    \begin{array}{rcrcrcl}
         x &+&    y &-&    z &=&  0  \\
           &-&   2y &+&    z &=&  0  \\
           & &    &&       7z &=&  0  
    \end{array}
  \right.
$$

En partant de la dernière ligne on trouve $z=0$, puis en remontant $y=0$, puis $x=0$.
Conclusion l'unique solution de ce système est $(0,0,0)$.

  \item
On applique le pivot de Gauss $L_2 \leftarrow L_2-L_1$ et $L_3 \leftarrow L_3-L_1$ :
$$ \left\{
    \begin{array}{rcrcrcl}
       x &+&  y &+& 2z &=& 5  \\
       x &-&  y &-&  z &=& 1  \\
       x & &    &+&  z &=& 3    
    \end{array}
  \right. 
\iff
 \left\{
    \begin{array}{rcrcrcl}
       x &+&  y &+&  2z &=& 5  \\
        &-&  2y &-&  3z &=& -4  \\
        &-&  y  &-&  z  &=& -2    
    \end{array}
  \right. 
$$
Puis $L_3\leftarrow 2L_3-L_2$ pour obtenir un système équivalent qui est triangulaire donc facile à résoudre :
$$ \left\{
    \begin{array}{rcrcrcl}
       x &+&  y &+&  2z &=& 5  \\
        &-&  2y &-&  3z &=& -4  \\
        &&    &&  z  &=& 0    
    \end{array}
  \right. 
\iff \left\{
    \begin{array}{rcl}
        x &=& 3 \\
        y &=& 2 \\
        z &=& 0    
    \end{array} \right.
$$

On n'oublie pas de vérifier que c'est une solution du système initial.

  \item  
On fait les opérations $L_2\leftarrow 3L_2+L_1$ et $L_3\leftarrow 3L_3-L_1$ pour obtenir :
$$\left\{
    \begin{array}{rcrcrcl}
      3x &-& y  &+& 2z &=& a \\
      -x &+& 2y &-& 3z &=& b \\
       x &+& 2y &+&  z &=& c 
    \end{array}
  \right.
\iff
\left\{
    \begin{array}{rcrcrcl}
      3x &-& y  &+& 2z &=& a \\
         & & 5y &-& 7z &=& 3b+a \\
         & & 7y &+&  z &=& 3c-a 
    \end{array}
  \right.
  $$
Puis on fait $L_3\leftarrow 5L_3-7L_2$, ce qui donne un système triangulaire :
$$\left\{
    \begin{array}{rcrcrcl}
      3x &-& y  &+& 2z &=& a \\
         & & 5y &-& 7z &=& 3b+a \\
         & &    & & 54z &=& 5(3c-a)-7(3b+a) 
    \end{array}
  \right.
  $$
En partant de la fin on en déduit : $z= \frac{1}{54}(-12a-21b+15c)$ puis en remontant cela donne
$$\left\{
    \begin{array}{rcl}
         x &=& \frac{1}{18}(8a+5b-c) \\
         y &=& \frac{1}{18}(-2a+b+7c) \\
         z &=& \frac{1}{18}(-4a-7b+5c)
    \end{array}
  \right.
  $$

\end{enumerate}

\fincorrection
\nocorrection
\nocorrection
\correction{001166}
$(S_1)$~: solution unique si $m^2 \not=4$, impossible 
sinon. $(S_2)$~: solution unique si  $m^2\not=1/2$, infinité sinon.
\fincorrection
\correction{001167}
$(S_1)$~: $a=b$ ou $b=c$ ou $c=a$.\\
$(S_2)$~: $2abc+bc+ca+ab=1$.
\fincorrection
\correction{001168}
$(S_1)$~: solution unique quels que soient $b_1,b_2,b_3,b_4$.\\
$(S_2)$~: solutions si $b_2=b_1+b_3$.\\
$(S_3)$~: solutions si $b_1+b_2-2b_4=0$ et $2b_1-b_3-2b_4=0$.\\
$(S_4)$~: solutions si $b_2=-2b_1$ et $b_3=-b_1$ et $b_4=3b_1$.
\fincorrection
\correction{001169}
\begin{enumerate}
  \item On commence par simplifier le système en effectuant les opérations suivantes sur les lignes :
$L_1 \leftarrow L_1 -L_4$, $L_2 \leftarrow L_2 -L_4$, $L_3 \leftarrow L_3 -L_4$ : 
$$ (S) \iff \left\{\begin{array}{rcrcrcccl}
    \lambda x &&  &&  &-& \lambda t &=& a-d\\ 
     && \lambda y&  &&-& \lambda t &=&b-d \\
    &&&&  \lambda z&-& \lambda t &=&c-d \\ 
    x&+&y&+&z&+&(1+\lambda)t &=&d
\end{array}\right.$$


  \item Traitons le cas particulier $\lambda = 0$.
Si $\lambda=0$ alors le système n'a des solutions que si $a=b=c=d$.
Les solutions sont alors les $(x,y,z,t)$ qui vérifie $x+y+z+t=d$.
(C'est un espace de dimension $3$ dans $\Rr^4$.)

  \item Si $\lambda\neq 0$ alors on peut faire l'opération suivante sur la dernière ligne : 
$L_4 \leftarrow L_4 - \frac 1\lambda L_1 - \frac1\lambda L_2 - \frac1\lambda L_3$ pour obtenir :
$$ (S) \iff \left\{\begin{array}{rcrcrcccl}
    \lambda x &&  &&  &-& \lambda t &=& a-d\\ 
     && \lambda y&  &&-& \lambda t &=&b-d \\
    &&&&  \lambda z&-& \lambda t &=&c-d \\ 
    &&&&&&(\lambda+4)t &=&d - \frac 1\lambda(a+b+c-3d)
\end{array}\right.$$



  \item Cas particulier $\lambda=-4$. La dernière ligne devient
$0=a+b+c+d$. 
Donc si $a+b+c+d\neq 0$ alors il n'y a pas de solutions.

Si $\lambda=-4$ et $a+b+c+d=0$ alors existe une infinité de solutions :
$$\left\{\Big(t-\frac{a-d}{4}, t-\frac{b-d}{4}, t-\frac{c-d}{4}, t\Big) \mid t\in\Rr \right\}.$$

  \item Cas général : $\lambda\neq 0$ et $\lambda\neq -4$.
On calcule d'abord $t= \frac{1}{\lambda+4}\left(d - \frac 1\lambda(a+b+c-3d)\right)$ 
et en remplaçant par la valeur de $t$ obtenue on en 
déduit les valeurs pour 
$x=t+\frac1\lambda(a-d),y=t+\frac1\lambda(b-d),z=t+\frac1\lambda(c-d)$. Il existe donc une solution unique :
$$\left(\frac{(\lambda+3)a-b-c-d}{\lambda(\lambda+4)}, \frac{(\lambda+3)b-a-c-d}{\lambda(\lambda+4)}, \frac{(\lambda+3)c-a-b-d}{\lambda(\lambda+4)},\frac{(\lambda+3)d-a-b-c}{\lambda(\lambda+4)}\right).$$

\end{enumerate}
\fincorrection
\correction{001170}
Pas de solution si $\lambda^2+\lambda-2\not=0$
($\lambda\not=1$ et $\lambda\not=-2$).
Si $\lambda=1$, pas de solution si $a+1\not=0$, infinité de solutions
sinon.
 Si $\lambda=-2$, solution unique.
\fincorrection
\nocorrection
\nocorrection
\nocorrection
\nocorrection
\nocorrection
\nocorrection
\nocorrection
\correction{001178}
On commence par simplifier le système :
\begin{itemize}
  \item on place la ligne $L_3$ en première position pour le pivot de Gauss ;
  \item on réordonne les variables dans l'ordre : $y,t,x,z$ pour profiter des lignes déjà simples.
\end{itemize}
$$
  \left\{
    \begin{array}{rcrcrcrcl}
y &+& t &+& x &+& z &=&  0  \\
3y &+& 3t &&  &+& z &=&  0  \\
-y &-& t &+& 2x &+& z &=&  0  \\
 &&  && 3x &+& 2z &=&  0  \\
    \end{array}\right.
$$
On commence le pivot de Gauss avec les opération $L_2 \leftarrow L_2-3L_1$ et
$L_3\leftarrow L_3+L_1$ pour obtenir :
$$
  \left\{
    \begin{array}{rcrcrcrcl}
y &+& t &+& x &+& z &=&  0  \\
 &&  &-& 3x &-& 2z &=&  0  \\
 &&  && 3x &+& 2z &=&  0  \\
 &&  && 3x &+& 2z &=&  0  \\
    \end{array}\right.
$$
Les $3$ dernières lignes sont identiques, on se ramène donc à un système avec $2$ équations et $4$ inconnues :
$$
  \left\{
    \begin{array}{rcrcrcrcl}
y &+& t &+& x &+& z &=&  0  \\
 &&  && 3x &+& 2z &=&  0  \\
    \end{array}\right.
$$
Nous choisissons $x$ et $y$ comme paramètres, alors $z=-\frac32x$
et $t = -x-y-z = \frac12x-y$.
Les solutions du système sont donc les 
$$\bigg\{ \big(x,y,z=-\frac32x,t=\frac12x-y\big) \mid x,y \in \Rr\bigg\}$$
\fincorrection
\nocorrection
\nocorrection
\nocorrection
\nocorrection
\nocorrection
\nocorrection
\nocorrection
\nocorrection
\nocorrection
\nocorrection
\nocorrection
\nocorrection
\nocorrection
\nocorrection
\nocorrection
\nocorrection
\nocorrection
\nocorrection
\correction{002478}
$S_1$: si $a\neq0$, le syst{\`e}me est {\'e}quivalent {\`a}
$y=-1-2\cosh a\ x$ et $z=x+2\cosh a$

si $a=0$, il est {\'e}quivalent {\`a} $x+y+z=1$

\smallskip

$S_2$: On peut soustraire {\`a} chaque ligne la ligne pr{\'e}c{\'e}dente, puis
2 fois la pr{\'e}c{\'e}dente, etc\dots On obtient ainsi un syst{\`e}me triangulaire
cram{\'e}rien et apr{\`e}s bien des calculs la solution $x_k=(-1)^{k+1}C_n^k$.

Voici une solution plus astucieuse.
Soit $P(X)=x_1X+x_2X^2+\dots+x_nX^n$ et $T$ l'op{\'e}rateur $Q(x)\mapsto XQ'(X)$.
Le syst{\`e}me peut s'{\'e}crire $P(1)=1,\ (TP)(1)=0,\ (T^2P)(1)=0,\ldots,\ (T^{n-1}P)(1)=0$.
On en d{\'e}duit que $P'(1)=P''(1)=\ldots=P^{(n-1)}(1)=0$, donc $P$ est de la forme
$P(X)=1+\lambda(1-X)^n$, et $\lambda=-1$ car le terme constant de $P$ est nul.
Donc $x_k=(-1)^{k+1}C_n^k$.

\fincorrection
\nocorrection
\nocorrection
\nocorrection
\nocorrection
\nocorrection
\nocorrection
\nocorrection
\nocorrection
\nocorrection
\correction{002768}
\begin{enumerate}
  \item 
  \begin{enumerate}
    \item \textbf{Par substitution.} La première équation s'écrit aussi $y=1-2x$.
On remplace maintenant $y$ dans la deuxième équation 
$$3x+7y=-2 \implies 3x+7(1-2x)=-2 \implies 11x=9 \implies  x=\frac{9}{11}.$$
On en déduit $y$ : $y=1-2x=1-2\frac{9}{11}=-\frac{7}{11}$.
La solution de ce système est donc le couple $(\frac{9}{11},-\frac{7}{11})$.

N'oubliez pas de vérifier que votre solution fonctionne !

    \item \textbf{Par le pivot de Gauss.} On garde la ligne $L_1$ et on remplace la ligne 
$L_2$ par $2L_2-3L_1$ :
$$
\left\{\begin{array}{ccccc}2x & + & y & = & 1\\ 3x & + & 7y & = & -2  \end{array}\right.
\iff 
\left\{\begin{array}{ccccc}2x & + & y & = & 1\\  &  & 11y & = & -7  \end{array}\right.
$$
On obtient un système triangulaire : on en déduit $y=-\frac{7}{11}$ et alors la première ligne permet d'obtenir
$x=\frac{9}{11}$.

    \item \textbf{Par les matrices.} En terme matriciel le système s'écrit

$$AX=Y \quad \text{  avec } A = \begin{pmatrix} 2 & 1 \\ 3 & 7 \\ \end{pmatrix}
\quad X = \begin{pmatrix} x \\ y \end{pmatrix} \quad Y = \begin{pmatrix} 1 \\ -2 \end{pmatrix}$$

On trouve la solution du système en inversant la matrice :
$$X = A^{-1} Y.$$

L'inverse d'une matrice $2\times 2$ se calcule ainsi
$$\text{ si } A = \begin{pmatrix} a & b \\ c & d \\ \end{pmatrix} \quad \text{ alors } \quad
A^{-1} = \frac{1}{ad-bc}\begin{pmatrix} d & -b \\ -c & a \\ \end{pmatrix}$$
Il faut bien sûr que le déterminant $\det A = \begin{vmatrix}  a & b \\ c & d \\ \end{vmatrix}=ad-bc$ 
soit différent de $0$.

Ici on trouve 
$$A^{-1} = \frac{1}{11}\begin{pmatrix} 7 & -1 \\ -3 & 2 \\ \end{pmatrix} \quad \text{ et } \quad
X = A^{-1} \begin{pmatrix} 1 \\ -2 \end{pmatrix} =  \frac{1}{11}\begin{pmatrix} 9 \\ -7 \end{pmatrix}$$

    \item \textbf{Par les formules de Cramer.}
Les formules de Cramer pour un système de deux équations sont les suivantes si le déterminant vérifie $ad-bc\neq 0$ :
$$
\left\{\begin{array}{ccccc}ax & + & by & = & e\\ cx & + & dy & = & f  \end{array}\right.
\implies x = \frac{\begin{vmatrix}  e & b \\ f & d\end{vmatrix}}{\begin{vmatrix}  a & b \\ c & d \end{vmatrix}}
\quad \text{ et } \quad y = \frac{\begin{vmatrix}  a & e \\ c & f\end{vmatrix}}{\begin{vmatrix}  a & b \\ c & d \end{vmatrix}}$$

Ce qui donne ici :
$$x = \frac{\begin{vmatrix}  1 & 1 \\ -2 & 7\end{vmatrix}}{\begin{vmatrix}  2 & 1 \\ 3 & 7 \end{vmatrix}} = \frac{9}{11}
 \quad \text{ et } \quad y = \frac{\begin{vmatrix}  2 & 1 \\ 3 & -2 \end{vmatrix}}{\begin{vmatrix}  2 & 1 \\ 3 & 7 \end{vmatrix}} = -\frac{7}{11}$$

  \end{enumerate}


  \item
  \begin{enumerate}
     \item Avant tout on regarde s'il existe une solution unique, c'est le cas si et seulement si le déterminant 
est non nul. Pour le premier système le déterminant est 
$\begin{vmatrix}  a & 1 \\ a^2+1 & 2a\end{vmatrix}=a^2-1$
donc il y a une unique solution si et seulement si $a\neq \pm 1$.

Bien sûr toutes les méthodes conduisent au même résultat !
Par exemple par substitution, en écrivant la première ligne $y=2-ax$,
la deuxième ligne devient $(a^2+1)x+2a(2-ax)=1$. On en déduit
que si $a\neq\pm1$ alors $x=\frac{4a-1}{a^2-1}$ puis $y= \frac{-2a^2+a-2}{a^2-1}$.

Traitons maintenant les cas particuliers.
Si $a=1$ alors le système devient :
$\left\{\begin{array}{ccccc} x & + & y & = & 2\\ 2x & + & 2y & = & 1 \end{array} \right.$
Mais on ne peut avoir en même temps $x+y=2$ et $x+y=\frac12$.
Donc il n'y a pas de solution.

Si $a=-1$ alors le système devient :
$\left\{\begin{array}{ccccc} -x & + & y & = & 2\\ 2x & - & 2y & = & 1 \end{array} \right.$
et il n'y a pas de solution.

     \item Ici le déterminant est $\begin{vmatrix}  a+1 & a-1 \\ a-1 & a+1 \end{vmatrix}=(a+1)^2-(a-1)^2=4a$.

Si $a\neq 0$ alors on trouve la solution unique $(x,y)$.
Par exemple avec la formule de Cramer
$$x=\frac{\begin{vmatrix}  1 & a-1 \\ 1 & a+1\end{vmatrix}}{4a} = \frac{1}{2a}
\text{ et } 
y = \frac{\begin{vmatrix}  a+1 & 1 \\ a-1 & 1 \end{vmatrix}}{4a} = \frac{1}{2a}.$$

Si $a=0$ il n'y a pas de solution.
  \end{enumerate}
  
\end{enumerate}
\fincorrection
\nocorrection
\correction{003408}
Système de Cramer ssi $m\ne 0,\pm2$; compatible ssi $m\ne2$.
\fincorrection
\correction{003409}
Système de Cramer ssi $m\ne 0,\pm1,\pm i$; compatible ssi $m\ne 0,\pm i$.
\fincorrection
\correction{003410}
Système de Cramer ssi $m\ne 1,\pm2i$; compatible ssi $m\ne 1$.
\fincorrection
\correction{003411}
Si $m \ne 0,-2$ alors système de Cramer ; sinon, système incompatible.
\fincorrection
\correction{003412}
Système de Cramer ssi $a,b,c$ sont distincts. Sinon, il y a des solutions ssi $d\in\{a,b,c\}$.
\fincorrection
\correction{003413}
Système compatible ssi $3a+2b+2c+d=0$.
\fincorrection
\correction{003414}
Système de Cramer.
\fincorrection
\correction{003415}
Système de Cramer ssi $\cos\alpha$, $\cos\beta$, $\cos\gamma$
         sont distincts. Sinon, il y a des solutions ssi les seconds membres
         correspondants sont égaux.
\fincorrection
\correction{003416}
CN d'existence de solution : $p+q+r = 0$. C'est une CNS si la liste
$(a,b,c)$ comporte au plus un zéro.
\fincorrection
\correction{003417}
\begin{enumerate}
  \item Pour éviter d'avoir à diviser par $a$ on réordonne nos lignes puis on applique la méthode du pivot :
{\small
$$\left\{\begin{array}{ccccccll}
x &+&  by &+& az &=& 1 & _{L_1} \cr
x &+& aby &+&  z &=& b & _{L_2} \cr
ax &+&  by &+&  z &=& 1 & _{L_3} \cr
\end{array}\right.
\iff
\left\{\begin{array}{ccccccll}
x &+&  by &+& az &=& 1 &_{L_1} \cr
  & & b(a-1)y &+&  (1-a)z &=& b-1 & _{L_2 \leftarrow L_2-L_1} \cr
  &&  b(1-a)y &+&  (1-a^2)z &=& 1-a & _{L_3 \leftarrow L_3-aL_1} \cr
\end{array}\right.
$$
}
On fait ensuite $L_3 \leftarrow L_3+L_2$ pour obtenir un système triangulaire équivalent au système initial :
$$\left\{\begin{array}{ccccccl}
x &+&  by &+& az &=& 1  \cr
  & & b(a-1)y &+&  (1-a)z &=& b-1  \cr
  &&   &&  (2-a-a^2)z &=& b-a  \cr
\end{array}\right.
$$

  \item Nous allons maintenant discuter de l'existence des solutions.
Remarquons d'abord que $2-a-a^2= -(a-1)(a+2)$.
Donc si $a\neq 1$ et $a\neq -2$ alors $2-a-a^2\neq 0$
donc $z=\frac{a-b}{(a-1)(a+2)}$. On a donc trouvé la valeur de $z$.
La deuxième ligne du système triangulaire est 
$b(a-1)y +  (1-a)z = b-1$ on sait déjà $a-1\neq 0$.
Si $b\neq 0$ alors, en reportant la valeur de $z$ obtenue, on trouve la valeur 
$y= \frac{b-1 - (1-a)z}{b(a-1)}$.
Puis avec la première ligne on en déduit aussi $x=1-by-az$.

Donc si $a\neq 1$ et $a \neq -2$ et $b\neq 0$ alors il existe une unique solution $(x,y,z)$.

  \item Il faut maintenant s'occuper des cas particuliers.
  \begin{enumerate}

     \item Si $a=1$ alors notre système triangulaire devient :
$$\left\{\begin{array}{ccccccc}
x &+&  by &+&z&=& 1 \cr
  & & & &  0 &=& b-1 \cr
  &&   &&  0 &=& b-1 \cr
\end{array}\right.
$$ 
Si $b\neq 1$ il n'y a pas de solution.
Si $a=1$ et $b=1$
alors il ne reste plus que l'équation $x+y+z=1$.
On choisit par exemple $y,z$ comme paramètres, l'ensemble des solutions 
est $$\big\{(1-y-z,y,z) \mid y,z \in \Rr \big\}.$$

  \item Si $a=-2$ alors le système triangulaire devient :
$$\left\{\begin{array}{ccccccc}
x &+&  by &-& 2z &=& 1 \cr
  & & -3by &+&  3z &=& b-1 \cr
  &&   &&  0 &=& b+2 \cr
\end{array}\right.
$$
Donc si $b\neq -2$ il n'y a pas de solution. 
Si $a=-2$ et $b=-2$ alors le système est
$$\left\{\begin{array}{ccccccc}
x &-&  2y &-& 2z &=& 1 \cr
  & & 2y &+&  z &=& -1 \cr
\end{array}\right.
$$
Si l'on choisit $y$ comme paramètre alors il y a une infinité de solutions
$$\big\{(-1-2y,y,-1-2y) \mid y \in \Rr \big\}.$$

    \item Enfin si $b=0$
alors la deuxième et troisième ligne du système triangulaire sont :
$(1-a)z =-1 $ et $(2-a-a^2)z =-a$. Donc $z=\frac{-1}{1-a}=\frac{-a}{2-a-a^2}$
(le sous-cas $b=0$ et $a=1$ n'a pas de solution).
Dans tous les cas il n'y a pas de solution.

    \item Conclusion :
\begin{itemize}
  \item Si $a\neq 1$ et $a\neq -2$ et $b\neq 0$, c'est un système de Cramer : il admet une unique solution.
  \item Si $a=1$ et $b\neq 1$ il n'y a pas de solution (le système n'est pas compatible).
  \item Si $a=1$ et $b=1$ il y a une infinité de solutions (qui forment un plan dans $\Rr^3$).
  \item Si $a=-2$ et $b\neq -2$ il n'y a pas de solution. 
  \item Si $a=-2$ et $b=-2$ il y a une infinité de solutions (qui forment une droite dans $\Rr^3$).
  \item Si $b=0$ il n'y a pas de solution.
\end{itemize}

  \end{enumerate}

\end{enumerate}
% $a\ne 1,-2$ et $b\ne0 \iff$ Cramer. \par
%          \leavevmode\vbox{\halign{$#$\hfil$\,  \Rightarrow \,$&#\hfil\cr
%          b=0  &incompatible.            \cr
%          a=1  &des solutions ssi $b=1$. \cr
% 	 a=-2 &des solutions ssi $b=-2$.\cr}}
\fincorrection
\correction{003418}
Décomposition en éléments simples de
$F = \dfrac x{X+a} + \dfrac y{X+2a} + \dfrac z{X+3a}$ avec
$F(1) = F(2) = F(3) = 1$ donc une solution unique si $a \ne 0$.
\fincorrection
\nocorrection
\nocorrection
\nocorrection
\nocorrection
\nocorrection
\correction{003424}
Notons $P(x)=ax^3+bx^2+cx+d$ un polynôme de degré $\le 3$.
\begin{enumerate}
  \item Tout d'abord calculons l'intégrale :
$$\int_{2}^4 P(x)\; dx = \left[a\frac{x^4}{4} + b\frac{x^3}{3}+c\frac{x^2}{2} +dx \right]_2^4
= 60a+\frac{56}{3}b+6c+2d.$$

  \item D'autre part
$$\alpha P(2) + \beta P(3) + \gamma P(4)
=  \alpha\big(8a+4b+2c+d\big)+\beta\big(27a+9b+3c+d\big)+\gamma\big(64a+16b+4c+d\big).$$
Donc 
$$\alpha P(2) + \beta P(3) + \gamma P(4)
= (8\alpha+27\beta+64\gamma)a + (4\alpha+9\beta+16\gamma)b+(2\alpha+3\beta+4\gamma)c+(\alpha+\beta+\gamma)d.$$

  \item Pour avoir l'égalité $\int_{2}^4 P(x)\; dx =\alpha P(2) + \beta P(3) + \gamma P(4)$
quelque soit les coefficients $a,b,c,d$ il faut et il suffit que 
$$(8\alpha+27\beta+64\gamma)a + (4\alpha+9\beta+16\gamma)b+(2\alpha+3\beta+4\gamma)c+(\alpha+\beta+\gamma)d =
60a+\frac{56}{3}b+6c+2d$$ 
ce qui équivaut à 
$$\left\{\begin{array}{rcrcrcl}
 \alpha & + &  \beta & + &  \gamma & = & 2 \\
 2\alpha & + & 3\beta & + & 4\gamma & = & 6 \\   
 4\alpha & + & 9\beta & + &  16\gamma & = & \frac{56}{3} \\ 
 8\alpha & + &  27\beta & + &  64\gamma & = & 60 \\     
         \end{array}
 \right.$$

De façon surprenante ce système à $3$ inconnues et $4$ équations a une solution unique :
$$\alpha = \frac 13,\quad \beta = \frac 43,\quad \gamma = \frac 13.$$
\end{enumerate}
\fincorrection
\correction{003425}
2. $$\begin{pmatrix} 3 &2 &6             \cr
                    4 &5 &\llap{$1$}2   \cr
                    2 &2 &5             \cr \end{pmatrix}^{-1} =
          \begin{pmatrix} -1 &         -2 & \hfill  6 \cr
                    -4 &         -3 & \hfill 12 \cr
              \hfill 2 &  \hfill  2 &        -7 \cr \end{pmatrix}  \Rightarrow 
          \left\{
          \begin{array}{lll} 
          x &=& 2^{-2} 3^6    \cr
                  y &=& 2^{-3}3^{12}  \cr
                  z &=& 2^2 3^{-7}.   \cr 
                  \end{array}\right.$$
 \fincorrection
\correction{003431}
Si $\mathrm{rg}(A) = n$, $\mathrm{rg}(\text{com}(A)) = n$.\quad
	     Si $\mathrm{rg}(A) = n-1$, $\mathrm{rg}(\text{com}(A)) = 1$.\quad
	     Si $\mathrm{rg}(A) \le n-2$, $\mathrm{rg}(\text{com}(A)) = 0$.
\fincorrection
\nocorrection
\correction{003461}
\begin{enumerate}
  \item 3.
  \item 4.
  \item 2.
  \item 3.
\end{enumerate}
 \fincorrection
\nocorrection
\correction{003463}
$\mathrm{rg}=3$ si $\lambda \ne 2$ et $\lambda \ne -25$. \par
\leavevmode\vbox{
\halign{$\lambda = #$\hfil &$ \Rightarrow  \mathrm{rg} = #$\hfil&\ $: #$.\hfil \cr
        2   &2 &11L_1 = 5L_2 + 9L_3    \cr
        -25 &2 &L_1 + 2L_2 + 9L_3 = 0  \cr }}
\fincorrection
\correction{003464}
$\mathrm{rg} = 3$ si $a \ne \frac 13$ ou $b \ne -3$, $\mathrm{rg}=2$ sinon.
\fincorrection
\correction{003465}
$\mathrm{rg} ABC \le 2  \Rightarrow  x = 13$.
         $M = \begin{pmatrix} 1&0 \cr 0&1 \cr \frac35 &-\frac 15 \cr \end{pmatrix}
              \begin{pmatrix} 1&0 \cr 0&1 \cr \end{pmatrix}
              \begin{pmatrix} 1&1&2 \cr -2&13&1 \end{pmatrix}$.
\fincorrection
\correction{003466}
Les colonnes de~$A$ engendrent les $n-p$ derniers vecteurs de la
base canonique.
\fincorrection
\correction{003467}
Échange des lignes $i$ et $j$.
\fincorrection
\nocorrection
\nocorrection
\correction{003470}
$A = \begin{pmatrix} 1 &x_1 \cr \vdots &\vdots \cr 1 &x_n \cr \end{pmatrix}
     \begin{pmatrix} y_1 &\dots &y_n \cr 1 &\dots &1 \cr \end{pmatrix}  \Rightarrow 
     \mathrm{rg}(A) = \begin{cases} 2 &\text{ si } (x_i) \text{ et } (y_j) \text{ ne sont pas constantes }\cr
                      1 &\text{ ou } 0 \text{ sinon.}\cr \end{cases}$

$B = \begin{pmatrix} 1 &x_1 &x_1^2 \cr \vdots &\vdots &\vdots \cr 1 &x_n &x_n^2 \cr \end{pmatrix}
     \begin{pmatrix} y_1^2 &\dots &y_n^2 \cr 2y_1 &\dots &2y_n \cr 1 &\dots &1 \cr \end{pmatrix}
      \Rightarrow 
     \mathrm{rg}(B) = \begin{cases} 3 &\text{ si } \mathrm{Card}\,(x_i)\ge 3 \text{ et } \mathrm{Card}\,(y_j) \ge 3 \cr
                      2 &\text{ si } \min\bigl(\mathrm{Card}\,(x_i),\mathrm{Card}\,(y_j)\bigr) = 2 \cr
                      1 &\text{ ou } 0 \text{ sinon.} \cr \end{cases}$
\fincorrection
\nocorrection
\correction{003472}
4. $A = \begin{pmatrix}1 & i \cr i &-1 \cr \end{pmatrix}$.
 \fincorrection
\nocorrection
\correction{003474}
$M = \Re\left[\begin{pmatrix} e^{i\theta} \cr \vdots \cr e^{ni\theta} \cr \end{pmatrix}
              \Bigl(\begin{matrix} 1 & \dots & e^{(n-1)i\theta} \cr \end{matrix}\Bigr)\right]
  \Rightarrow  \mathrm{rg} M \le 2$.

Le premier mineur $2\times2$ vaut $-\sin^2\theta
 \Rightarrow  \mathrm{rg} M = 2$ si $\theta \not \equiv 0 (\mathrm{mod}\, \pi)$.
Sinon, $\mathrm{rg} M = 1$.
\fincorrection
\nocorrection
\correction{003476}
$E$ est un sev et un idéal à gauche de $\mathcal{M}_n(\R)$.
	     Il est isomorphe à $\mathcal{L}({H,\R^n})$ où $H$ est un supplémentaire de
	     $\Im A$ dans $\R^n$. $\dim E = n(n-\mathrm{rg}(A))$.
\fincorrection
\nocorrection
\nocorrection
\correction{003479}
2 ou 0.
\fincorrection
\nocorrection
\nocorrection
\nocorrection
\correction{003483}
\begin{enumerate}
  \item 
  \item $B$ admet $r$ lignes indépendantes d'indices $i_1,\dots,i_r$
    et $C$ admet $r$ colonnes indépendantes d'indices $j_1,\dots,j_r$.
    Soient $B'$ et $C'$ les sous matrices carrées associées dans $B$ et $C$.
    Alors la sous-matrice de~$A$ d'indices $i_1,\dots,i_r$ pour les lignes
    et $j_1,\dots,j_r$ pour les colonnes est $B'C'$, de rang~$r$.
    Donc $\mathrm{rg}(A)\ge r$ et l'inégalité inverse est bien connue.
  \item Soient $i_1,\dots,i_r$ $r$ indices tels que les
    lignes associées dans~$A$ sont linéairement indépendantes,
    et $B\in\mathcal{M}_{r,n}(\R)$ la sous-matrice correpondante. Par construction, $\mathrm{rg}(B)=r$.
    Chaque ligne de~$A$
    étant combinaison linéaire des lignes de~$B$, il existe $C\in\mathcal{M}_{r,n}(\R)$
    telle que $A=BC$. Et on a $r=\mathrm{nb.lignes}(C)\ge\mathrm{rg}(C)\ge \mathrm{rg}(A)=r$.
  \item 
  \item Comprendre dans cette question que $B,C$ ne sont pas
    forcément les matrices construites en {\bf 2}.
    Notons $\mathrm{vect}(X)$ l'espace vectoriel engendré par les colonnes
    d'une matrice~$X$. De $A=BC={}^tC{}^tB$ on tire
    $\mathrm{vect}(A)\subset\mathrm{vect}(B)$ et $\mathrm{vect}(A)\subset\mathrm{vect}({}^tC)$,
    et tous ces espaces sont de dimension~$r$, donc ils sont égaux.
    On en déduit qu'il existe une matrice~$P\in GL_r(\R)$ telle
    que $B={}^tCP$ d'où $CB=C{}^tCP$. $\mathrm{rg}(C{}^tC)=\mathrm{rg}(C)=r$ et $P$
    est inversible donc $\mathrm{rg}(CB)=r$.
\end{enumerate}
 \fincorrection
\nocorrection
\correction{005367}
$\Delta=\mbox{det}M=\mbox{Van}(1,2,...,n)\neq0$ et le système est de \textsc{Cramer}. Les formules de \textsc{Cramer} fournissent alors pour $k\in\llbracket1,n\rrbracket$, $x_k=\frac{\Delta_k}{\Delta}$ où 

$$\Delta_k=\mbox{Van}(1,...,k-1,0,k+1,...,n)=(-1)^{k+1}\left|
\begin{array}{cccccc}
1&\ldots&k-1&k+1&\ldots&n\\
1& &(k-1)^2&(k+1)^2& &n^2\\
\vdots& &\vdots&\vdots& &\vdots\\
 & & & & & \\
\vdots& &\vdots&\vdots& &\vdots\\
1& &(k-1)^{n-1}&(k+1)^{n-1}& &n^{n-1}
\end{array}
\right|.$$
(en développant par rapport à la $k$-ème colonne). Par linéarité par rapport à chaque colonne, on a alors

\begin{align*}\ensuremath
\Delta_k&=(-1)^{k+1}1\times2...\times(k-1)\times(k+1)\times...\times n\times\mbox{Van}(1,2,...,k-1,k+1,...,n)\\
 &=(-1)^{k+1}\frac{n!}{k}\frac{\mbox{Van}(1,2,...,n)}{(k-(k-1))...(k-1)((k+1)-k)....(n-k)}=(-1)^{k+1}\frac{n!}{k!(n-k)!}\Delta,
\end{align*}
et donc,

\begin{center}
\shadowbox{
$\forall k\in\llbracket1,n\rrbracket,\;x_k=(-1)^{k+1}C_n^k$.
}
\end{center}

\fincorrection
\correction{005375}
$m$ est un paramètre réel
\begin{enumerate}
\item  $\mbox{det}S=2(m(m-5)-6)+(3(m-5)-3)+7(6-m)=2m^2-14m+12=2(m-1)(m-6)$. Le système est de \textsc{Cramer} si et seulement si $m\in\{1,6\}$. 
Si $m\notin\{1,6\}$, les formules de \textsc{Cramer} fournissent alors~:

$$\begin{array}{l}
x=\frac{1}{2(m-1)(m-6)}\left|
\begin{array}{ccc}
4&3&1\\
5&m&2\\
7&3&m-5
\end{array}
\right|
=\frac{2(m-6)(2m-9)}{2(m-1)(m-6)}=\frac{2m-9}{m-1}\\
y=\frac{1}{2(m-1)(m-6)}\left|
\begin{array}{ccc}
2&4&1\\
-1&5&2\\
7&7&m-5
\end{array}
\right|
=\frac{14(m-6)}{2(m-1)(m-6)}=\frac{7}{m-1}\\
z=\frac{1}{2(m-1)(m-6)}\left|
\begin{array}{ccc}
2&3&4\\
-1&m&5\\
7&3&m7
\end{array}
\right|  =\frac{-14(m-6)}{2(m-1)(m-6)}=-\frac{7}{m-1}\\
\end{array}
$$

Si $m\in\{1,6\}$, $\mbox{det}S=0$. Un déterminant principal est $\left|
\begin{array}{cc}
2&1\\
-1&2
\end{array}
\right|=5\neq0$. On peut choisir les deux premières équations comme équations principales et $x$ et $z$ comme inconnues principales. Le système des deux premières équations équivaut à $\left\{
\begin{array}{l}
x=\frac{3+(m-6)y}{5}\\
z=\frac{14-(2m+3)y}{5}
\end{array}
\right.$.

La dernière équation fournit alors une condition nécessaire et suffisante de compatibilité (les termes en y disparaissent automatiquement pour $m\in\{1,6\}$ et donc pas la peine de les calculer).

\begin{align*}\ensuremath
7x+2y+(m-5)z=7&\Leftrightarrow7\frac{3+(m-6)y}{5}+3y+(m-5)\frac{14-(2m+3)y}{5}=7\Leftrightarrow 21+14(m-5)-35=0\\
 &\Leftrightarrow 14(m-6)= 0\Leftrightarrow m = 6.
\end{align*}

Si $m=1$, le système n'a pas de solution et si $m=6$, l'ensemble des solutions est $\{(\frac{3}{5},y,-\frac{y}{5}),\;y\in\Rr\}$.

\item  $\mbox{det}S=2(-8m-4+2)-(4m+1)+5(2m+2m+1)=0$. Le système n'est jamais de \textsc{Cramer}. Un déterminant principal est $\left|
\begin{array}{cc}
2&1\\
1&2
\end{array}
\right|=3\neq 0$. On peut choisir les deux premières équations comme équations principales et $x$ et $z$ comme inconnues principales. Le système des deux premières équations équivaut à $\left\{
\begin{array}{l}
x=\frac{6m-4-(4m+1)y}{3}\\
z=\frac{-3m+8+(5m+2)y}{3}
\end{array}
\right.$. La dernière équation fournit alors une condition nécessaire et suffisante de compatibilité.

\begin{align*}\ensuremath
5x-y+4z=3m-2&\Leftrightarrow5\frac{6m-4-(4m+1)y}{3}-y+4\frac{-3m+8+(5m+2)y}{3}=3m-2\\
 &\Leftrightarrow5(6m-4)+4(-3m+8)-3(3m-2)=0\Leftrightarrow9(m+2)=0\Leftrightarrow m=-2.
\end{align*}

Si $m\neq-2$, le système n'a pas de solution. Si $m=-2$, l'ensemble des solutions est $\{(\frac{-16+7y}{3},y,\frac{14-8y}{3}),\;y\in\Rr\}$.

\item  $\left|
\begin{array}{ccc}
1&1&1\\
1&m&1\\
m&-1&-m
\end{array}\right|=-2m^2+2m=-2m(m-1)$. Le système est de \textsc{Cramer} en $x$, $y$ et $z$ si et seulement si $m\in\{0,1\}$.

Si $m\notin\{0,1\}$, les formules de \textsc{Cramer} fournissent~:

$$\begin{array}{l}
x=\frac{1}{-2m(m-1)}\left|
\begin{array}{ccc}
3-t&1&1\\
m+2+mt&m&1\\
-1+t&-1&-m
\end{array}\right|=\frac{(2m^2-2m)t+(-2m^2+2m)}{-2m(m-1)}=-t+1\\
y=\frac{1}{-2m(m-1)}\left|
\begin{array}{ccc}
1&3-t&1\\
1&m+2+mt&1\\
m&-1+t&-m
\end{array}\right|=\frac{(-2m^2-2m)+(-2m^2+2m)}{-2m(m-1)}=\frac{m+1}{m-1}t+1\\
z=\frac{1}{-2m(m-1)}\left|
\begin{array}{ccc}
1&1&3-t\\
1&m&m+2+mt\\
m&-1&-1+t
\end{array}\right|=\frac{(2m^2+2m)t+(-2m^2+2m)}{-2m(m-1)}=-\frac{m+1}{m-1}t+1.
\end{array}$$

Dans ce cas, l'ensemble des solutions est $\{(-t+1,\frac{m+1}{m-1}t+1,-\frac{m+1}{m-1}t+1,t),\;t\in\Rr\}$.

Si $m=0$, le système s'écrit $\left\{
\begin{array}{l}
x+y+z+t=3\\
x+z=2\\
y+t=-1
\end{array}
\right.\Leftrightarrow \left\{
\begin{array}{l}
z=2-x\\
t=-1-y
\end{array}
\right.$. Dans ce cas, l'ensemble des solutions est $\{(x,y,2-x,1-y),\;(x,y)\in\Rr^2\}$.

Si $m=1$, le système s'écrit $\left\{
\begin{array}{l}
x+y+z+t=3\\
x+y+z-t=3\\
x-y-z-t=-1
\end{array}\right.\Leftrightarrow\left\{
\begin{array}{l}
t=0\\
x+y+z=3\\
x-y-z=-1
\end{array}\right.\Leftrightarrow\left\{
\begin{array}{l}
t=0\\
x=1\\
z=2-y
\end{array}\right.$. Dans ce cas, l'ensemble de solutions est $\{(1,y,2-y,0),\;z\in\Rr\}$.

\item  
\begin{align*}\ensuremath
\mbox{det}(S)&=\left|
\begin{array}{cccc}
1&2&3&m\\
2&1&m&3\\
3&m&1&2\\
m&3&2&1
\end{array}\right|=\left|
\begin{array}{cccc}
m+6&2&3&m\\
m+6&1&m&3\\
m+6&m&1&2\\
m+6&3&2&1
\end{array}\right|=(m+6)\left|
\begin{array}{cccc}
1&2&3&m\\
1&1&m&3\\
1&m&1&2\\
1&3&2&1
\end{array}\right|\\
 &=(m+6)\left|
\begin{array}{cccc}
1&2&3&m\\
0&-1&m-3&3-m\\
0&m-2&-2&2-m\\
0&1&-1&1-m
\end{array}\right|=(m+6)\left|
\begin{array}{ccc}
-1&m-3&3-m\\
m-2&-2&2-m\\
1&-1&1-m
\end{array}\right|\\
 &=(m+6)\left|
\begin{array}{ccc}
-1&m-3&0\\
m-2&-2&-m\\
1&-1&-m
\end{array}\right|=-m(m+6)\left|
\begin{array}{ccc}
-1&m-3&0\\
m-2&-2&1\\
1&-1&1
\end{array}\right|\\
 &=-m(m+6)\left|
\begin{array}{ccc}
-1&m-3&0\\
m-3&-1&0\\
1&-1&1
\end{array}\right|
=-m(m+6)\left|
\begin{array}{cc}
-1&m-3\\
m-3&-1
\end{array}\right|=m(m-2)(m-4)(m+6).
\end{align*}

Le système est de \textsc{Cramer} si et seulement si $m\notin\{0,2,4,-6\}$. Dans ce cas~:

\begin{align*}\ensuremath
m(m-2)(m-4)(m+6)x&=\left|
\begin{array}{cccc}
m-1&2&3&m\\
1&1&m&3\\
0&m&1&2\\
0&3&2&1
\end{array}\right|=\left|
\begin{array}{cccc}
0&2-(m-1)&3-m(m-1)&m-3(m-1)\\
1&1&m&3\\
0&m&1&2\\
0&3&2&1
\end{array}\right|\\
 &=-\left|
\begin{array}{ccc}
3-m&-m^2+m+3&-2m+3\\
m&1&2\\
3&2&1
\end{array}\right|=-\left|
\begin{array}{ccc}
5m-6&-m^2+5m-3&-2m+3\\
m-6&-3&2\\
0&0&1
\end{array}\right|\\
 &=-[-3(5m-6)-(m-6)(-m^2+5m-3)]\\
 &=-m^3+11m^2-18m=-m(m-2)(m-9).
\end{align*}

et $x=-\frac{m-9}{(m-4)(m+6)}$.

\begin{align*}\ensuremath
m(m-2)(m-4)(m+6)y&=\left|
\begin{array}{cccc}
1&m-1&3&m\\
2&1&m&3\\
3&0&1&2\\
m&0&2&1
\end{array}\right|=\left|
\begin{array}{cccc}
-2m+3&0&-m^2+m+3&-2m+3\\
2&1&m&3\\
3&0&1&2\\
m&0&2&1
\end{array}\right|\\
 &=\left|
\begin{array}{ccc}
-2m+3&-m^2+m+3&-2m+3\\
3&1&2\\
m&2&1
\end{array}\right|\\
 &=\left|
\begin{array}{ccc}
3m^2-5m-6&-m^2+m+3&2m^2-4m-3\\
0&1&0\\
m-6&2&-3
\end{array}\right|\\
 &=-3(3m^2-5m-6)-(m-6)(2m^2-4m-3)\\
 &=-2m^3+7m^2-6m=-m(2m-3)(m-2)
\end{align*}

et $y=-\frac{2m-3}{(m-4)(m-6)}$.

\begin{align*}\ensuremath
m(m-2)(m-4)(m+6)z&=\left|
\begin{array}{cccc}
1&2&m-1&m\\
2&1&1&3\\
3&m&0&2\\
m&3&0&1
\end{array}\right|=\left|
\begin{array}{cccc}
-2m+3&-m+3&0&-2m+3\\
2&1&1&3\\
3&m&0&2\\
m&3&0&1
\end{array}\right|\\
 &=-\left|
\begin{array}{ccc}
-2m+3&-m+3&-2m+3\\
3&m&2\\
m&3&1
\end{array}\right|\\
 &=-(-2m+3)(m-6)+3(5m-6)-m(2m^2-5m+6)=-2m^3+7m^2-6m\\
 &=-m(2m-3)(m-2),
\end{align*}

et $z=-\frac{2m-3}{(m-4)(m-6)}$.

\begin{align*}\ensuremath
m(m-2)(m-4)(m+6)t&=\left|
\begin{array}{cccc}
1&2&3&m-1\\
2&1&m&1\\
3&m&1&0\\
m&3&2&0
\end{array}\right|=\left|
\begin{array}{cccc}
-2m+3&-m+3&-m^2+m+3&0\\
2&1&m&1\\
3&m&1&0\\
m&3&2&0
\end{array}\right|\\
 &=\left|
\begin{array}{ccc}
-2m+3&-m+3&-m^2+m+3\\
3&m&1\\
m&3&2
\end{array}\right|\\
 &=(-2m+3)(2m-3)-3(3m^2-5m-3)+m(m^3-m^2-4m+3)\\
 &=m^4-m^3-17m^2+30m=m(m-2)(m^2+m-15)
\end{align*}

et $t=\frac{m^2+m-15}{(m-4)(m-6)}$.

Si $m = 0$, le système s'écrit 

\begin{align*}\ensuremath
\left\{
\begin{array}{l}
x+2y+3z=-1\\
2x+y+3t=1\\
3x+z+2t=0\\
3y+2z+t=0
\end{array}
\right.&\Leftrightarrow\left\{
\begin{array}{l}
x+y+z+t=\;(E_1+E_2)\\
2x+y+3t=1\\
x+y+z+t=0(E_3+E_4)\\
3y+2z+t=0
\end{array}
\right.\Leftrightarrow\left\{
\begin{array}{l}
t=-x-y-z\\
-x-2y-3z=1\\
-x+2y+z=0
\end{array}
\right.
\\
 &\Leftrightarrow\left\{
\begin{array}{l}
z=x-2y\\
-x-2y-3(x-2y)=1\\
t=-x-y-z
\end{array}
\right.\Leftrightarrow\left\{
\begin{array}{l}
y=x+\frac{1}{4}\
z=-x-\frac{1}{2}\\
t=-x+\frac{1}{4}
\end{array}
\right.
\end{align*}

D'où l'ensemble de solutions~:~$\{(x,x+\frac{1}{4},-x-\frac{1}{4};-x+\frac{1}{2}),\;x\in\Rr\}$.

Si $m=2$, on obtient pour ensemble de solutions~:~$\{(x,-x-\frac{5}{8},x+\frac{1}{2};-x-\frac{1}{8}),\;x\in\Rr\}$.

Si $m=4$ ou $m=-6$, on voit en résolvant que le système est incompatible.

\item  $\left|
\begin{array}{ccc}
m&1&1\\
-1&-1&m\\
1&-1&-m
\end{array}
\right|=m(2m)+(-m+1)+(m+1)=2(m^2+1)\neq 0$ ($m$ désignant un paramètre réel).

Le système formé des équations $1$, $2$ et $4$ est donc de \textsc{Cramer}. Les formules de \textsc{Cramer} fournissent alors~:

$$x=\frac{2m^2-m-1}{m^2+1},\;y=3-m\;\mbox{et}\;z=\frac{3m-1}{m^2+1}.$$

La troisième équation fournit alors une condition nécessaire et suffisante de compatibilité~:

\begin{align*}\ensuremath
-m\frac{2m^2-m-1}{m^2+1}&+3-m+m\frac{3m-1}{m^2+1}=-m\\
 &\Leftrightarrow-m(2m^2-m-1)+(3-m)(m^2+1)+m(3m-1)=-m(m^2+1)\\
 &\Leftrightarrow-2m^3+7m^2+3= 0
\end{align*}

Le système est compatible si et seulement si $m$ est l'une des trois racines de l'équation $-2X^3+7X^2+3=0$.

\item  $\mbox{det}S=\frac{1}{abc}\left|
\begin{array}{ccc}
a&b&c\\
a^2&b^2&c^2\\
1&1&1
\end{array}
\right|=\frac{1}{abc}\left|
\begin{array}{ccc}
1&1&1\\
a&b&c\\
a^2&b^2&c^2
\end{array}
\right|=\frac{\mbox{Van}(a,b,c)}{abc}$.

Si $a$, $b$ et $c$ sont deux à deux distincts, le système est de \textsc{Cramer}. On obtient~:

$$x=\frac{abc}{mbc}\frac{\mbox{Van}(m,b,c)}{\mbox{Van}(a,b,c)}=\frac{a(b-m)(c-m)}{m(b-a)(c-a)},$$

puis, par symétrie des rôles, $y=\frac{b(a-m)(c-m)}{m(a-b)(c-b)}$ et $z=\frac{c(a-m)(b-m)}{m(a-c)(b-c)}$.

Si $a=b\neq c$ (ou $a=c\neq b$ ou $b=c\neq a$), le système s'écrit~:

$$\left\{
\begin{array}{l}
x+y=1-z\\
ax+ay+cz=m\\
\frac{1}{a}x+\frac{1}{a}y+\frac{1}{c}z=\frac{1}{m}
\end{array}
\right.\Leftrightarrow\left\{
\begin{array}{l}
x+y=1-z\\
a(1-z)+cz=m\\
\frac{1}{a}(1-z)+\frac{1}{c}z=\frac{1}{m}
\end{array}
\right.\Leftrightarrow\left\{
\begin{array}{l}
x+y=1-z\\
z=\frac{m-a}{c-a}\\
(\frac{1}{c}-\frac{1}{a})\frac{m-a}{c-a}=\frac{1}{m}-\frac{1}{a}
\end{array}
\right..$$

Le système est compatible si et seulement si $(m-a)(m-c)=0$ ou encore ($m=a$ ou $m=c$). Dans ce cas, l'ensemble des solutions est~:~$\{(x,\frac{m-c}{a-c}-x;\frac{m-a}{c-a}),\;x\in\Rr\}$.

Si $a=b=c$, le système s'écrit~:~$x+y+z=1=\frac{m}{a}=\frac{a}{m}$. Le système est compatible si et seulement si $m=a=b=c$ et dans ce cas l'ensemble des solutions est~:~$\{(x,y,1-x-y),\;(x,y)\in\Rr^2\}$.

\item  

\begin{align*}\ensuremath
\mbox{det}S&=\left|
\begin{array}{ccc}
(b+c)^2&b^2&c^2\\
a^2&(a+c)^2&c^2\\
a^2&b^2&(a+b)^2
\end{array}
\right|=\left|
\begin{array}{ccc}
(b+c)^2&b^2&c^2\\
a^2-(b+c)^2&(a+c)^2-b^2&0\\
0&b^2-(a+c)^2&(a+b)^2-c^2
\end{array}
\right|\\
 &=(a+b+c)^2\left|
\begin{array}{ccc}
(b+c)^2&b^2&c^2\\
a-b-c&a+c-b&0\\
0&b-a-c&a+b-c
\end{array}
\right|=(a+b+c)^2\left|
\begin{array}{ccc}
2bc&b^2&c^2\\
-2c&a+c-b&0\\
-2(b-c)&b-a-c&a+b-c
\end{array}
\right|
\\
 &=2(a+b+c)^2(c^2(-c(b-a-c)+(b-c)(a+c-b))+(a+b-c)(bc(a+c-b)+b^2c))\\
 &=2(a+b+c)^2(c^2b(a-b+c)+(a+b-c)bc(a+c))\\
 &=2bc(a+b+c)^2(a^2+ab+ac)=2abc(a+b+c)^3.
\end{align*}

Si $abc(a+b+c)\neq0$, le système est de \textsc{Cramer} et on obtient après calcul~:

$$x=\frac{(a-b+c)(a+b-c)}{2abc(a+b+c)},\;y=\frac{(a-b-c)(a+b-c)}{2abc(a+b+c)}\;\mbox{et}\;z=\frac{(a-b+c)(a-b-c)}{2abc(a+b+c)}.$$ 

Si $a=0$ (ou $b=0$ ou $c=0$), le système s'écrit~:

$$\left\{
\begin{array}{l}
(b+c)^2x+b^2y+c^2z=1\\
c^2(y+z)=1\\
b^2(y+z)=1
\end{array}
\right.
.$$

Donc,

Si (($a=0$ et $b^2\neq c^2$) ou ($b=0$ et $a^2\neq c^2$) ou ($c=0$ et $a^2\neq b^2$)), le système n'a pas de solution.

Si $a=0$ et $b=c\neq 0$, l'ensemble des solutions est $\{(0,y,-\frac{y}{b^2}),\;y\in\Rr\}$ (résultats analogues pour les cas ($b=0$ et $a=c\neq 0$) et ($c=0$ et $a=b\neq0$)).

Si $a=b=c=0$, il n'y a pas de solution.
 
Si $a=0$ et $c=-b\neq0$, l'ensemble des solutions est $\{(x,y-\frac{y}{b^2}),\;(x,y)\in\Rr^2\}$ (résultats analogues pour ($b=0$ et $c=-a\neq0$) et ($c=0$ et $b=-a\neq0$).

Si $abc\neq0$ et $a+b+c=0$, le système équivaut à l'équation $a^2x+b^2y+c^2z=1$. L'ensemble  des solutions est $\{(x,y,\frac{1-a^2x-b^2y}{c^2}),\;(x,y)\in\Rr^2\}$.

\item  

\begin{align*}\ensuremath
\mbox{det}S&=\left|
\begin{array}{ccc}
a&b&c\\
c&a&b\\
b&c&a
\end{array}
\right|=(a+b+c)\left|
\begin{array}{ccc}
1&b&c\\
1&a&b\\
1&c&a
\end{array}
\right|
=(a+b+c)\left|
\begin{array}{ccc}
1&b&c\\
0&a-b&b-c\\
0&c-b&a-c
\end{array}
\right|\\
 &=(a+b+c)((a-b)(a-c)+(b-c)^2)=(a+b+c)(a^2+b^2+c^2-ab-ac-bc)\\
 &=(a+b+c)(a+jb+j^2c)(a+j^2b+jc)
\end{align*}

Si $\mbox{det}S\neq0$, les formules de \textsc{Cramer} fournissent :

$$x\mbox{det}S=\left|
\begin{array}{ccc}
p&b&c\\
q&a&b\\
r&c&a
\end{array}
\right|=p(a^2-bc)+q(c^2-ab)+r(b^2-ac).$$

Je n'ai pas envie de finir.

\item  Soit $P=X^3-X-1$. $P$ et $P'=3X^2-1$ n'ont pas de racines communes dans $\Cc$ car $\frac{1}{\sqrt{3}}$ et $-\frac{1}{\sqrt{3}}$ ne sont pas racines de $P$ et donc les racines de $P$ sont simples ou encore, $a$, $b$ et $c$ sont deux à deux distincts.

Ainsi, $\mbox{det}S=\mbox{Van}(a,b,c)\neq0$ et le système est de \textsc{Cramer}.

$$(b-a)(c-a)(c-b)x=\left|
\begin{array}{ccc}
0&1&1\\
2&b&c\\
3&b^2&c^2
\end{array}
\right|=-2(c^2-b^2)+3(c-b)=(c-b)(3-2(b+c))=(c-b)(3+2a),$$

(car $a+b+c=0$) et $x=\frac{3+2a}{(b-a)(c-a)}$.

$$(b-a)(c-a)(c-b)y=\left|
\begin{array}{ccc}
1&0&1\\
a&2&c\\
a^2&3&c^2
\end{array}
\right|=2(c^2-a^2)-3(c-a)=(c-a)(2(a+c)-3)=-(c-a)(3+2b),$$

et $y=-\frac{3+2b}{(b-a)(c-a)}$.

$$(b-a)(c-a)(c-b)z=\left|
\begin{array}{ccc}
1&1&0\\
a&b&2\\
a^2&b^2&3
\end{array}
\right|=-2(b^2-a^2)+3(b-a)=(b-a)(3+2c),$$

et $z=\frac{3+2c}{(c-a)(c-b)}$ (difficile d'aller plus loin).
\end{enumerate}
\fincorrection
\correction{005378}
Soit $D_n$ le déterminant du système pour $n\geq3$.

En développant ce déterminant suivant sa première colonne, on obtient la relation de récurrence~:
 
$$\forall n\geq5,\;D_n=D_{n-1}-D_{n-2},$$

ce qui fournit aisément par récurrence, en tenant compte de $D_3=D_4=-1$~:

$$\forall k\geq1,\;D_{3k}=D_{3k+1}=(-1)^k\;\mbox{et}\;D_{3k+2}=0.$$

Pour $n$ élément de $3\Nn^*\cup(1+3\Nn^*)$, le système est de \textsc{Cramer} et homogène et admet donc une et une seule solution à savoir la solution nulle.

Pour $n=3k+2$, puisque $D_n=0$ mais que le mineur de format $n-1$ constitué des $n-1$ premières lignes et colonnes est $D_{n-1}$ et est donc non nul, le système est homogène de rang $n-1$ et l'ensemble des solutions est un sous-espace vectoriel de $\Rr^n$ de dimension $1$. On trouve aisément $\mathcal{S}=\{\lambda(1,-1,0,1,-1,0...,1,-1),\;;\lambda\in\Rr\}$.
\fincorrection
\correction{005381}
Soit $(x_1,...,x_n)\in\Rr^n$ et $F=\sum_{k=1}^{n}\frac{x_k}{X+b_k}$.

La fraction rationnelle $F$ s'écrit, après réduction au même dénominateur~:

$$F=\frac{P}{Q}\;\mbox{où}\;Q=\prod_{k=1}^{n}(X+b_k)\;\mbox{et}\;P\mbox{est un polynôme de degré infèrieur ou égal à}\;n-1.$$

Maintenant,

$$(x_1,...,x_n)\;\mbox{solution de}\;(S)\Leftrightarrow\forall k\in\{1,...,n\},\;F(a_k)=1\Leftrightarrow\forall k\in\{1,...,n\},\;(Q-P)(a_k)=0.$$

Par suite, puisque les $a_k$ sont deux à deux distincts, $Q-P$ est divisible par $\prod_{k=1}^{n}(X-a_k)$. Mais, $Q$ est unitaire de degré $n$ et $P$ est de degré infèrieur ou égal à $n-1$, et donc $Q-P$ est unitaire de degré $n$ ce qui montre que $Q-P=\prod_{k=1}^{n}(X-a_k)$ ou encore que 

$$P=\prod_{k=1}^{n}(X+b_k)-\prod_{k=1}^{n}(X-a_k).$$

Réciproquement, si $F=\frac{\prod_{k=1}^{n}(X+b_k)-\prod_{k=1}^{n}(X-a_k)}{\prod_{k=1}^{n}(X+b_k)}$, alors $\forall k\in\{1,...,n\},\;F(a_k)=1$.

En résumé, 

\begin{align*}\ensuremath
(x_1,...,x_n)\;\mbox{solution de}\;(S)&\Leftrightarrow\sum_{k=1}^{n}\frac{x_k}{X+b_k}=  \frac{\prod_{k=1}^{n}(X+b_k)-\prod_{k=1}^{n}(X-a_k)}{\prod_{k=1}^{n}(X+b_k)}\\
 &\Leftrightarrow\forall i\in\{1,...,n\},\;x_i=\lim_{x\rightarrow -b_i}(x+b_i)\frac{\prod_{k=1}^{n}(x+b_k)-\prod_{k=1}^{n}(x-a_k)}{\prod_{k=1}^{n}(x+b_k)}\\&\Leftrightarrow\forall i\in\{1,...,n\},\;x_i=\frac{\prod_{k=1}^{n}(b_i+a_k)}{\prod_{k=1}^{n}(b_k-b_i)}
\end{align*}
\fincorrection
\correction{005639}
 Le déterminant du système est $\Delta=\text{Van}(1,\ldots,n)\neq0$. Le système proposé est donc un système de \textsc{Cramer}.

Les formules de \textsc{Cramer} donnent : $\forall j\in\llbracket1,n\rrbracket$, $x_j=\frac{\Delta_j}{\Delta}$   où

\begin{align*}\ensuremath
\Delta_j&=\left|
\begin{array}{ccccccc}
1&\ldots&1&1&1&\ldots&1\\
1& &j-1&0&j+1& &n\\
\vdots& &\vdots&\vdots&\vdots& &\vdots\\
 & \\
 & \\
\vdots& &\vdots&\vdots&\vdots& &\vdots\\
1&\ldots&(j-1)^{n-1}&0&(j+1)^{n-1}&\ldots&n^{n-1}
\end{array}
\right|\\
 &= (-1)^{j+1}\left|
\begin{array}{cccccc}
1&\ldots&j-1&j+1&\ldots&n\\
\vdots& &\vdots&\vdots& &\vdots\\
 & \\
  & \\
\vdots& &\vdots&\vdots& &\vdots\\
1&\ldots&(j-1)^{n-1}&(j+1)^{n-1}&\ldots&n^{n-1}
\end{array}
\right|(\text{en développant suivant la}\;j\text{-ème colonne})\\
  &= (-1)^{j+1}1...(j-1)(j+1)...n\left|
\begin{array}{cccccc}
1&\ldots&1&1&\ldots&1\\
1& &j-1&j+1& &n\\
\vdots& &\vdots&\vdots& &\vdots\\
  & \\
\vdots& &\vdots&\vdots& &\vdots\\
1&\ldots&(j-1)^{n-2}&(j+1)^{n-2}&\ldots&n^{n-2}
\end{array}
\right|(\text{par}\;n-\text{linéarité})\\
 & = (-1)^{j+1}\frac{n!}{j}\text{Van}(1,...,(j-1),(j+1),...,n) = (-1)^{j+1}\frac{n!}{j}\frac{\text{Van}(1,...,n)}{(j-1)\ldots(j-(j-1))((j+1)-j)\ldots(n-j)}\\
  &= (-1)^{j+1}\frac{n!}{j!(n-j)!}\text{Van}(1,...,n)=(-1)^{j+1}\dbinom{n}{j}\text{Van}(1,...,n).
\end{align*}

Finalement,

\begin{center}
\shadowbox{
$\forall j\in\llbracket1,n\rrbracket$, $x_j =(-1)^{j+1}\dbinom{n}{j}$.
}
\end{center}
\fincorrection
\correction{007026}
Pour les trois premiers systèmes, on élimine les dénominateurs en multipliant, et les termes de degré deux se simplifient ce qui conduit à des systèmes linéaires. Les ensembles de solutions sont: 
\[ 
\{ (4,8)\}; \quad 
\{ (7,26)\}; \quad
\{ (19,32)\}.
\]
Pour les trois suivants, on effectue des changements de variable simples du type $X=\frac{1}{x}$, ce qui conduit à des systèmes linéaires. Il faut alors vérifier la compatibilité avec le domaine de résolution des équations, le cas échéant. Le second système aboutit en effet à un système linéaire auxiliaire qui a une unique solution dont une composante est nulle, ce qui ne convient pas. Le troisième système aboutit à un système linéaire ayant une infinité de solutions. L'ensemble de solutions du système initial est alors  une hyperbole privée d'un point. Plus précisément, les ensembles de solutions sont :
\[ 
\left\{ \left(-\frac{1}{2},\frac{1}{7}\right)\right\}; \quad 
\emptyset; \quad
\left\{ \left(t,\frac{3t}{2t-1}\right), t\in \R\setminus \left\{0,\frac{1}{2}\right\}\right\}.
\]

Pour les deux systèmes restants, on effectue également des changements de variable : $X=\frac{1}{4(x-2)}$ et $Y=\frac{1}{15(y-1)}$ pour le premier, et $X=\frac{1}{x-y}$, $Y=\frac{1}{x+y}$ pour le dernier, ce qui conduit à la résolution d'un deuxième système linéaire pour conclure. Les ensembles de solutions sont :

\[ 
\left\{ \left(\frac{17}{8},\frac{16}{5}\right)\right\}; \quad 
\left\{ \left(\frac{4}{15},\frac{1}{15}\right)\right\}.
\]
\fincorrection
\correction{007027}
$17$ et $12$.
\fincorrection
\correction{007028}
Si $a$ et $b$ sont les mesures des deux autres côtés, on obtient par Pythagore $a^2+b^2=13^2=169$, et $ab=60$. On en déduit que $(a+b)^2=289=17^2$ et $(a-b)^2=49=7^2$, et d'où on tire un système linéaire. Les solutions sont $12$ et $5$.
\fincorrection
\nocorrection
\correction{007030}
Les solutions réelles sont $\left(\sqrt 2,\frac{1}{\sqrt 2}\right)$ et $\left(-\sqrt 2,-\frac{1}{\sqrt 2}\right)$.
\fincorrection
\correction{007031}
On trouve une unique solution $\left(\frac{3}{8},\frac{16}{3\sqrt 3},\frac{1}{\sqrt 3}\right)$.
\fincorrection
\nocorrection
\correction{007033}
Si $n$ est impair, le système possède une solution unique.
Si $n$ est pair alors le système possède une solution si, et seulement si, $\sum_{k=1}^n a_k=0$.

Cet exercice peut également se résoudre en considérant les symétries centrales dont les centres ont pour affixes les $a_k$, et la composition de toutes ces isométries. Si $n$ est impair, cette composée est une symétrie centrale. Sinon, c'est une translation.
\fincorrection
\correction{007034}
Prendre le logarithme des équations puis changer de variable. Solution unique : $\left(\frac{1}{2},\frac{\sqrt{2}}{5}\right)$.
\fincorrection
\correction{001151}
\begin{enumerate}
  \item Dans le cas $n=2$, $n=4$ les matrice suivantes conviennent :
$$ J =
\begin{pmatrix}
0&-1\\
1&0\\
\end{pmatrix}, \qquad
J' =
\begin{pmatrix}
J&(0)\\
(0)&J\\
\end{pmatrix}.
$$
  \item Supposons qu'un tel morphisme existe.
Soit $J$ sa matrice pour une base fixée. Alors $J^2 = -I_n$ o\`u
$I_n$ est la matrice identité de taille $n$. En termes de
déterminant nous avons : $\det (J^2)  =  \det I_n $, ce qui
s'écrit $(\det J)^2 = (-1)^n$. Donc $n$ est pair car $(\det J)^2$
est positif.
 \end{enumerate}
\fincorrection
\nocorrection
\correction{003443}
$27a^4 = 256b^3$.
\fincorrection
\nocorrection
\nocorrection
\nocorrection
\correction{003447}
ctrex: $A=\Z$, $P(0)=0$ et $P(2)=1$.
\fincorrection
\correction{003448}
On se place dans $\Z/p\Z$ et on considère $J = (\delta_{i,i+1\bmod p})$.
On a $J^p = I$ et $A = a_0J^0 + \dots a_{p-1}J^{p-1}$ donc
$A^p = (a_0^p+\dots + a_{p-1}^p)I$ (car on est en caractéristique $p$).

On en déduit $\det(A) = \det(A)^p = (a_0^p+\dots + a_{p-1}^p)^p = a_0+\dots + a_{p-1}$.

Autre méthode en restant dans~$\Z$~:
$\det(A) = \sum_{\sigma\in S_p} \varepsilon(\sigma)a_{1,\sigma(1)}\dots a_{p,\sigma(p)}
         = \sum_{\sigma\in S_p} \varepsilon(\sigma)a_{\sigma(1)-1\bmod p}\dots a_{\sigma(p)-p\bmod p}$.
Notons $x(\sigma) = \varepsilon(\sigma)a_{\sigma(1)-1\bmod p}\dots a_{\sigma(p)-p\bmod p}$ et
$c$ le cycle $(1,2,\dots,p)$. Alors $x(\sigma) = x(c^{-k}\circ\sigma\circ c^k)$
pour tout $k\in\Z$. Le nombre de permutations distinctes que l'on obtient
à $\sigma$ fixé en faisant varier $k$ est égal à $1$ si $\sigma$ et $c$ commutent,
et à $p$ sinon, d'après la relation~: $\mathrm{Card}\,(\text{orbite})\times\mathrm{Card}\,(\text{stabilisateur}) = \mathrm{Card}\,({<c>}) =p$.
De plus, $c$ et $\sigma$ commutent si et seulement si $\sigma\in{<c>}$ (facile),
d'où $\det(A)\equiv \sum_{k=0}^{p-1}\varepsilon(c^k)a_{k}^p \equiv a_0+\dots+a_{p-1}\bmod p$.
\fincorrection
\correction{003449}
$\det(M) = \sum_{\sigma\in S_n}\varepsilon(\sigma)a_{1\sigma(1)}\dots a_{n\sigma(n)}$.
Soit $\sigma\in S_n$ telle que $\sigma\ne\sigma^{-1}$. Alors les termes associés
à $\sigma$ et $\sigma^{-1}$ sont égaux car $M$ est symétrique, donc la somme
de ces deux termes est paire. Soit $\sigma\in S_n$ telle que $\sigma=\sigma^{-1}$.
Alors comme $n$ est impair, il existe $i\in{[[1,n]]}$ tel que $\sigma(i)=i$
donc le terme associé à~$\sigma$ est pair.
\fincorrection
\nocorrection
\correction{005377}
Si $z_k$ est l'affixe complexe de $M_k$ et $a_k$ est l'affixe complexe de $A_k$, le problème posé équivaut au système~:
$$\forall k\in\{1,...,n-1\},\;z_k+z_{k+1}=2a_k\;\mbox{et}\;z_n+z_1=2a_n.$$

Le déterminant de ce système vaut~:

\begin{align*}\ensuremath
\left|
\begin{array}{ccccc}
1&1&0&\ldots&0\\
0&\ddots&\ddots&\ddots&\vdots\\
\vdots&\ddots& &\ddots&0\\
0& &\ddots&\ddots&1\\
1&0&\ldots&0&1
\end{array}
\right|
&=1.1^{n-1}+(-1)^{n+1}.1^{n-1}\;(\mbox{en développant suivant la première colonne})\\
 &=1+(-1)^{n+1}.
\end{align*}

Si $n$ est impair, $\mbox{det}S=2\neq 0$ et le système admet une et une seule solution.

On obtient $z_2=2a_1-z_1$, $z_3=2a_2-2a_1+z_1$,..., $z_n=2a_{n-1}-2a_{n-2}+...+2a_2-2a_1+z_1$ et enfin :

$$2a_{n-1}-2a_{n-2}+...+2a_2-2a_1+z_1+z_1=2a_n,$$

et donc $z_1=a_1-a_2+...-a_{n-1}+a_n$ puis $z_2=a_1+a_2-a_3+...+a_{n-1}-a_n$ puis $z_3=-a_1+a_2+a_3-a_4...+a_n$ ... puis $z_n=-a_1+a_2-a_3+...+a_{n-1}+a_n$.

Si $n$ est pair, $\mbox{det}S=0$ mais le mineur formé des $n-1$ premières lignes et $n-1$ dernières colonnes est non nul. Donc, le système est de rang $n-1$, les $n-1$ premières équations et $n-1$ dernières inconnues peuvent être choisies pour équations et inconnues principales.

On résout les $n-1$ premières équations constituant un sytème de \textsc{Cramer} en $z_2$,...,$z_n$. On obtient 

$$z_2=2a_1-z1,\;z_3=2a_2-2a_1+z_1,...,\;z_n=2a_{n-1}-2a_{n-2}+...-2a_2+2a_1-z_1.$$

La dernière équation fournit alors une condition nécessaire et suffisante de compatibilité~:

$$2a_{n-1}-2a_{n-2}+...-2a_2+2a_1-z_1+z_1=2a_n\Leftrightarrow a_1+a_3...=a_2+a_4+...$$

Cette dernière condition se traduit géométriquement par le fait que les systèmes de points $(A_1,A_3,...)$ et
$(A_2,A_4,...)$ ont même isobarycentre.

En résumé, si $n$ est pair et si les systèmes de points $(A_1,A_3,...)$ et $(A_2,A_4,...)$ n'ont pas même isobarycentre, le problème n'a pas de solutions.

Si $n$ est pair et si les systèmes de points $(A_1,A_3,...)$ et $(A_2,A_4,...)$ ont même isobarycentre, le problème a une infinité de solutions~:~$M_1$ est un point quelconque puis on construit les symétriques successifs par rapport aux points $A_1$, $A_2$ ...
\fincorrection
\nocorrection
\nocorrection
\nocorrection
\nocorrection
\correction{001157}
Soit $$ A =
\begin{pmatrix}
0&a& b\\
-a&0&c\\
-b&-c&0\\
\end{pmatrix}, \qquad
B =
\begin{pmatrix}
0&a\\
-a&0\\
\end{pmatrix}.
$$
Alors $\det A = 0$, mais $\det B = a^2$ est non nul si $a\neq 0$.

\fincorrection
\nocorrection
\nocorrection
\nocorrection
\nocorrection
\nocorrection
\nocorrection
\nocorrection
\nocorrection
\correction{003429}
\begin{enumerate}
  \item 
  \item $A = \begin{pmatrix} 1        &\lambda \cr
                            \lambda  &1       \cr \end{pmatrix}$,\quad
             $B = \begin{pmatrix} 0        &\lambda \cr
                            -\lambda &0       \cr \end{pmatrix}$,\quad
             $|\lambda| > \frac 12$.
\end{enumerate}
 \fincorrection
\nocorrection
\correction{003434}
3. On complète par $\begin{pmatrix} 0 & 0 & 1 & 0 \cr 0 & 0 & 0 & 1 \cr \end{pmatrix}
               \Rightarrow \text{base} =
               \left\{\,
               \left(\begin{smallmatrix} -1 \cr 2 \cr -1 \cr 0 \cr \end{smallmatrix}\right),
               \left(\begin{smallmatrix} -2 \cr 3 \cr 0 \cr -1 \cr \end{smallmatrix}\right)
               \,\right\}$.
 \fincorrection
\correction{003435}
Développer le produit.

Un seul coeff. non nul par ligne et colonne, ou une ligne nulle.
\fincorrection
\correction{003436}
\begin{enumerate}
  \item 
  \item %\def \dddd#1#2{\text{det}\bigl(\vec #1,\vec #2\bigr)}\relax
    Si $d_{ab} \ne 0$, prendre 
    $$\left\{
    \begin{array}{lllll}
    \vec c &=& -\frac {d_{bc}}{d_{ab}}\vec a &+& \frac {d_{ac}}{d_{ab}}\vec b \cr
    \vec d &=& \phantom{-}\frac {d_{db}}{d_{ab}}\vec a
                           &+& \frac {d_{ad}}{d_{ab}}\vec b.\cr
\end{array}\right. $$
\end{enumerate}
 \fincorrection
\correction{003437}
Si $(\vec a, \vec b, \vec c)$ est une base, décomposer $\vec d$.
         Si $\vec a = \lambda \vec b + \mu \vec c$, on obtient $\vec 0 = \vec 0$.
\fincorrection
\nocorrection
\nocorrection
\correction{003441}
2. (b) $(I + E_{ij})^k = I + kE_{ij}$.
              Calculer le pgcd d'une ligne par opérations élémentaires à l'aide
              de Bézout. Ce pgcd vaut 1 sinon $M \notin SL_n(\Z)$.
  \fincorrection
\correction{003442}
$(\det A)^n$.
\fincorrection
\correction{005368}
En remplaçant les colonnes $C_1$,..., $C_n$ par respectivement $C_1+iC_{n+1}$,..., $C_n+iC_{2n}$, on obtient~:

$$\mbox{det}C=\mbox{det}
\left(
\begin{array}{cc}
A+iB&B\\
-B+iA&A
\end{array}
\right),$$ 
puis en remplaçant les lignes $L_{n+1}$,..., $L_{2n}$ de la nouvelle matrice par respectivement $L_{n+1}-iL_1$,..., $L_{2n}-iL_n$, on obtient~:

$$\mbox{det}C=\mbox{det}
\left(
\begin{array}{cc}
A+iB&B\\
0&A-iB
\end{array}
\right)=\mbox{det}(A+iB)\mbox{det}(A-iB)=|\mbox{det}(A+iB)|^2\in\Rr^+.$$
\fincorrection
\correction{005370}
On suppose $n\geq2$. La matrice nulle est solution du problème.
Soit $A$ un élément de $M_n(\Cc)$ tel que $\forall B\in M_n(\Cc),\;\mbox{det}(A+B)=\mbox{det}A+\mbox{det}B$. En particulier, $2\mbox{det}A=\mbox{det}(2A)=2^n\mbox{det}A$ et donc $\mbox{det}A=0$ car $n\geq2$. Ainsi, $A\notin GL_n(\Cc)$.
Si $A\neq 0$, il existe une certaine colonne $C_j$ qui n'est pas nulle. Puisque la colonne $-Cj$ n'est pas nulle, on peut compléter la famille libre $(-C_j)$ en une base $(C_1',...,-C_j,...,C_n')$ de $M_{n,1}(\Cc)$. La matrice $B$ dont les colonnes sont justement $C_1$',...,$-C_j$,...,$C_n'$ est alors inversible de sorte que $\mbox{det}A+\mbox{det}B=\mbox{det}B\neq 0$. Mais, $A+B$ a une colonne nulle et donc $\mbox{det}(A+B)=0\neq\mbox{det}A+\mbox{det}B$.
Ainsi, seule la matrice nulle peut donc être solution du problème .

\begin{center}
\shadowbox{
$\forall A\in M_n(\Cc),\;(\forall M\in M_n(\Cc),\;\text{det}(A+M)=\text{det}(A)+\text{det}(M))\Leftrightarrow A=0$.
}
\end{center}
\fincorrection
\correction{005379}
$(1)\Rightarrow(2)$. Montrons par récurrence sur $n\geq1$ que~:~$(\forall(a_1,...,a_n)\in E^n/\;(\mbox{det}(f_i(a_j))_{1\leq i,j\leq n}= 0)\Rightarrow((f_1,...,f_n)$ liée).

Pour $n=1$,

$$(\forall a_1\in E/\;\mbox{det}(f_i(a_j))_{1\leq i,j\leq1}=0)\Rightarrow(\forall a_1/\;f_1(a_1)=0)\Rightarrow(f_1=0)\Rightarrow(f_1)\;\mbox{liée}.$$

Soit $n\geq2$. Supposons que $(\forall(a_1,...,a_{n-1})\in E^{n-1}/\;\mbox{det}(f_i(a_j))_{1\leq i,j\leq n-1}=0)\Rightarrow(f_1,...,f_{n-1})$ liée.

Soient $f_1$,..., $f_n$ $n$ fonctions telles que $\forall(a_1,...,a_n)\in E^n/\;\mbox{det}(f_i(a_j))_{1\leq i,j\leq n}=0$.

Si $(f_1,...,f_{n-1})$ est liée alors $(f_1,...,f_n)$ est liée en tant que sur famille d'une famille liée. Si $(f_1,...,f_{n-1})$ est libre, par hypothèse de récurrence, il existe $a_1$,...,$a_{n-1}$ $n-1$ éléments de $E$ tels que $\mbox{det}(f_i(a_j))_{1\leq i,j\leq n-1}\neq 0$. Mais, par hypothèse, on a~:

$$\forall x\in E,\;\mbox{det}(f_i(a_1),...,f_i(a_{n-1}),f_i(x))_{1\leq i\leq n}=0.$$

En développant ce déterminant suivant sa dernière colonne, on obtient une égalité du type $\sum_{i=1}^{n}\lambda_if_i(x)=0$ où les $\lambda_i$ sont indépendants de $x$ ou encore une égalité du type $\sum_{i=1}^{n}\lambda_if_i=0$ avec $\lambda_n=\mbox{det}(f_i(a_j))_{1\leq i,j\leq n-1}\neq0$ ce qui montre encore que $(f_1,...,f_n)$ est liée.

$(2)\Rightarrow(1)$. On suppose que $\exists(a_1,...,a_n)\in E^n/\;\mbox{det}(f_i(a_j))_{1\leq i,j\leq n}\neq 0)$. Montrons que $(f_1,...,f_n)$ est libre.

Soit $(\lambda_1,...,\lambda_n)\in\Cc^n$ tel que $\sum_{i=1}^{n}\lambda_if_i=0$. En particulier~:~$\forall j\in\{1,...,n\},\;\sum_{i=1}^{n}\lambda_if_i(a_j)=0$. Les $n$ égalités précédentes fournissent un système d'équations linéaires en les $\lambda_i$ à $n$ inconnues, $n$ équations, de déterminant non nul et homogène ou encore un système de \textsc{Cramer} homogène dont on sait qu'il admet pour unique solution $(\lambda_1,...,\lambda_n)=(0,...,0)$. On a montré que $(f_1,...,f_n)$ est libre.
\fincorrection
\correction{005380}
Soit $A_n$ la matrice de l'énoncé.

En développant $\mbox{det}A_n$ suivant sa première colonne puis en développant le déterminant de format $n-1$ obtenu suivant sa première ligne, on obtient $\mbox{det}A_n=-\mbox{det}A_{n-2}$ pour $n\geq3$.

Par suite, pour $p\geq1$, $\mbox{det}A_{2p}=(-1)^{p-1}\mbox{det}A_2=(-1)^p\neq0$ et pour $p\geq1$, $A_{2p}$ est inversible.

On a aussi, pour $p\geq1$, $\mbox{det}A_{2p+1}=(-1)^{p-1}\mbox{det}A_3=0$ et, pour $p\geq1$, $A_{2p+1}$ n'est pas inversible. Finalement, $A_n$ est inversible si et seulement si $n$ est pair.

Dorénavant, on pose $n=2p$ ($p\geq1$).

Pour $X=(x_i)_{1\leq i\leq n}$ et $Y=(y_i)_{1\leq i\leq n}$ vecteurs colonnes donnés, on a~:

$$AX=Y\Leftrightarrow\left\{
\begin{array}{l}
x_2=y_1\\
\forall i\in\{2,...,2p-1\},\;x_{i-1}+x_{i+1}=y_i\\
x_{2p-1}=y_{2p}
\end{array}
\right..$$
 
Ce système se résoud en $x_2=y_1$ puis, par récurrence, pour $k\leq p$, $x_{2k}=y_{2k-1}-y_{2k-3}+...+(-1)^{k-1}y_1$ et aussi $x_{2p-1}=y_{2p}$, puis, par récurrence, pour $k\leq p$, $x_{2k-1}=y_{2k}-y_{2k+2}+...+(-1)^{p-k}y_{2p}$. D'où l'inverse de $A$ quand $n=8$ par exemple~:

$$A^{-1}=
\left(
\begin{array}{cccccccc}
0&1&0&-1&0&1&0&-1\\
1&0&0&0&0&0&0&0\\
0&0&0&1&0&-1&0&1\\
-1&0&1&0&0&0&0&0\\
0&0&0&0&0&1&0&-1\\
1&0&-1&0&1&0&0&0\\
0&0&0&0&0&0&0&1\\
-1&0&1&0&-1&0&1&0
\end{array}
\right).$$
\fincorrection
\correction{005614}
On a toujours $A{^t}(\text{com}A)=(\text{det}A)I_n$. Par passage au déterminant et puisqu'une matrice a même déterminant que sa transposée, on obtient

\begin{center}
$(\text{det}A)(\text{det}(\text{com}A)) = (\text{det}A)^n$.
\end{center}

\textbullet~Si $\text{det}A$ n'est pas nul, on en déduit $\text{det}(\text{com}A)=(\text{det}A)^{n-1}$.

\textbullet~Si $\text{det}A$ est nul, on a $A{^t}(\text{com}A) = 0$ et donc ${^t}\text{com}A$ est soit nulle, soit diviseur de zéro, et donc dans tous les cas non inversible. Il en est de même de $\text{com}A$ et donc $\text{det}(\text{com}A) = 0=(\text{det}A)^{n-1}$. Finalement 

\begin{center}
\shadowbox{
$\forall A\in\mathcal{M}_(\Rr),\;\text{det}(\text{com}A)=(\text{det}A)^{n-1}$.
}
\end{center}
\fincorrection
\correction{005615}
\textbullet~Si $A$ est de rang $n$, c'est-à-dire inversible, l'égalité $(\text{com}A)\times\frac{1}{\text{det}A}{^t}A =I_n$ montre que $\text{com}A$ est inversible et donc de rang $n$.

Dans ce qui suit, le lien entre le rang d'une matrice et la nullité des différents mineurs est hors programme. On suppose maintenant $\text{rg}(A)\leqslant n-1$.

\textbullet~Si $\text{rg}A\leqslant n-2$. Montrons que tous les mineurs de format $n-1$ extraits de $A$ sont nuls.

Soient $j_1$,\ldots, $j_{n-1}$, $n-1$ numéros de colonnes deux à deux distincts puis $A'\in\mathcal{M}_{n,n-1}(\Kk)$ dont les colonnes sont $C_{j_1}$,\ldots, $C_{j_{n-1}}$. Puisque $A$ est de rang au plus $n-2$, la famille des colonnes de $A'$ est liée et donc $A'$ est de rang au plus $n-2$. Il en est de même de la matrice ${^t}A'\in\mathcal{n-1,n}(\Kk)$ et donc toute matrice $A''$ obtenue en supprimant l'une des colonnes de $A'$ est carrée, de format $n-1$, non inversible. Son déterminant est donc nul.

Ainsi, tout déterminant obtenu en supprimant une ligne et une colonne de $\text{det}(A)$ est nul ou encore tous les mineurs de format $n-1$ extraits de $A$ sont nuls. Finalement, si $\text{rg}A\leqslant n-2$, $\text{com}A= 0$.

\textbullet~Il reste à étudier le cas où $\text{rg}A=n-1$ et donc $\text{dim}KerA = 1$.

L'égalité $\text{det}A=0$ impose $A{^t}(\text{com}A) = 0$. Mais alors $\text{Im}({^t}(\text{com}A))\subset\text{Ker}A$ et en particulier $\text{rg}(\text{com}A)=\text{rg}({^t}(\text{com}A))\leqslant\text{dim}(\text{Ker}A)=1$. Ainsi, si $\text{rg}(A)=n-1$ alors $\text{rg}(\text{com}A)\in\{0,1\}$.

Montrons que l'un au moins des mineurs de format $n-1$ extraits de $A$ est non nul ce qui montrera que $\text{rg}(\text{com}A)=1$.

Puisque $\text{rg}A=n-1$, il existe $n-1$ colonnes $C_{j_1}$,\ldots, $C_{j_{n-1}}$ de $A$ constituant une famille libre. La matrice $A'\in\mathcal{M}_{n,n-1}(\Kk)$ constituée par ces colonnes est de rang $n-1$. Il en est de même de sa transposée. Mais alors, il existe $n-1$ colonnes de ${^t}A'$ linéairement indépendantes. La matrice $A''$ constituée de ces $n-1$ colonnes est carrée de format $n-1$ et de rang $n-1$. $A''$ est donc inversible et il en est de même de ${^t}A''$. Le déterminant de ${^t}A''$ est un mineur de format $n-1$ extrait de $A$ et non nul.

En résumé,

\begin{center}
\shadowbox{
$\forall A\in\mathcal{M}_n(\Rr)$, $\text{rg}(\text{com}A) =\left\{
\begin{array}{l}
n\;\text{si}\;\text{rg}(A) = n\\
1\;\text{si}\;\text{rg}(A) = n-1\\
0\;\text{si}\;\text{rg}(A)\leqslant n-2
\end{array}
\right.$.
}
\end{center}
\fincorrection
\correction{005616}
Si $\text{rg} M\leqslant n-1$, l'égalité $M=\text{com}M$ entraîne $M{^t}M=M{^t}(\text{com}M)=(\text{det}M)I_n=0$ et donc $M = 0$. En effet, 

\begin{align*}\ensuremath
M{^t}M=0&\Rightarrow\forall X\in\mathcal{M}_{n,1}(\Rr),\;M{^t}MX=0\Rightarrow\forall X\in\mathcal{M}_{n,1}(\Rr),\;{^t}XM{^t}MX = 0\Rightarrow\forall X\in\mathcal{M}_{n,1}(\Rr),\;\left\|{^t}MX\right\|^2=0\\
 &\Rightarrow\forall X\in\mathcal{M}_{n,1}(\Rr),\;{^t}MX = 0\Rightarrow{^t}M=0\Rightarrow M=0.
\end{align*}

En résumé, si $M$ est solution, $M=0$ ou $M$ est inversible.

Dans le deuxième cas, d'après l'exercice \ref{ex:rou17}, on doit avoir $\text{det}M=(\text{det}M)^{n-1}$ et donc, puisque $\text{det}M\neq0$, $\text{det}M\in\{-1,1\}$ (et même $\text{det}M=1$ si $n$ est impair) car $\text{det}M$ est réel.

\textbullet~Si $\text{det}M=-1$, on doit avoir $M{^t}M=-I_n$ mais ceci est impossible car le coefficient ligne $1$, colonne $1$, de la matrice $M{^t}M$ vaut $m_{1,1}^2+...+m_{1,n}^2\neq-1$.

\textbullet~Il reste le cas où $\text{det}M=1$, l'égalité $M=\text{com}M$ entraîne $M{^t}M=I_n$ c'est-à-dire $M$ est orthogonale positive.

Réciproquement, si $M$ est orthogonale positive, ${^t}M=M^{-1}=\frac{1}{\text{det}M}{^t}(\text{com}M)={^t}\text{com}M$ et donc $M=\text{com}M$.

Finalement ,

\begin{center}
\shadowbox{
$\mathcal{S}=\{0\}\cup O_n^+(\Rr)$.
}
\end{center}
\fincorrection
\correction{005648}
$A = 0$ convient.

Réciproquement, on a tout d'abord $\text{det}(A+A)=\text{det}A+\text{det}A$ ou encore $(2^n-2)\text{det}A=0$ et, puique $n\geqslant2$, $\text{det}A = 0$. Donc,

\begin{center}
$A\notin GL_n(\Kk)$ et $A$ vérifie : $\forall M\in M_n(\Kk)$, $\text{det}(A+M)=\text{det}M$.
\end{center}

Supposons $A\neq 0$. Il existe donc une colonne $C_j\neq0$.

La colonne $-C_j$ n'est pas nulle et d'après le théorème de la base incomplète, on peut construire une matrice $M$ inversible dont la $j$-ème colonne est $-C_j$. Puisque $M$ est inversible, $\text{det}M\neq0$ et puisque la $j$-ème colonne de la matrice $A+M$ est nulle, $\text{det}(A+M)= 0$. Pour cette matrice $M$, on a $\text{det}(A+M)\neq\text{det}A +\text{det}M$ et $A$ n'est pas solution du problème. Finalement

\begin{center}
\shadowbox{
$(\forall M\in M_n(\Kk)$, $\text{det}(A+M)=\text{det}A+\text{det}M)/lra A = 0$.
}
\end{center}
\fincorrection
\nocorrection
\nocorrection
\nocorrection
\nocorrection
\nocorrection
\nocorrection
\nocorrection
\nocorrection
\nocorrection
\nocorrection
\nocorrection
\nocorrection
\nocorrection
\nocorrection
\correction{001611}
  \begin{enumerate}
  \item 
    Soit $(v,w)\in\mathrm{Com}$, $(\lambda,\mu)\in\R^{2}$. $u(\lambda v+\mu
    w)=\lambda uv+\mu uw=\lambda vu+\mu wu=(\lambda v+\mu w)u$. Donc
    $\mathrm{Com}$ est un sous espace vectoriel de $\mathcal{L}(E,E)$.

  \item
    Soit $x\in E_{\lambda}$. $u(v(x))=uv(x)=vu(x)=v(\lambda x)=\lambda
    v(x)$ donc $v(x)\in E_{\lambda}$.

  \item
    Chaque valeur propre est de multiplicité 1 donc chaque espace propre
    est de dimension 1. Ainsi, si $x\in E_{\lambda}\setminus\{0\}$,
    $E_{\lambda}=\R x$. Comme $v(x)\in E_{\lambda}$,
    $\exists\alpha\in\R, v(x)=\alpha x$. Donc $x$ est un vecteur propre
    de $v$.

  \item
    Soit $(e_{1},...,e_{n})$ une base de vecteurs propres de $u$. C'est
    aussi une base de vecteurs propres pour tout élément de $\mathrm{Com}$. Tout
    élément de $\mathrm{Com}$ est donc représenté par une matrice diagonale dans
    $(e_{1},...,e_{n})$. Réciproquement, tout endomorphisme représenté
    dans cette base par une matrice diagonale commute avec $u$. Donc
    $$
    \mathrm{Com}=\{v\in\mathcal{L}(E,E), \exists (\alpha_{1},...,\alpha_{n})\in\R^{n},
    M_{v/(e_{1},...,e_{n})}=
      \begin{pmatrix}
        \alpha_{1}& & \\ &\ddots& \\ & & \alpha_{n}
      \end{pmatrix}
      \}$$
    On en déduit que $\mathrm{Com}$ est de dimension $n$.

  \item
    $u\,u^{i}=u(u\cdots u)=(u\cdots u)u=u^{i}\, u$. Donc $\forall
    i\in\{0,...,n-1\}, u^{i}\in\mathrm{Com}$. Ainsi
    $\mathrm{Vect}(\mathrm{id},u,...,u^{n-1})\subset\mathrm{Com}$.

  \item
    Soit $x_{k}\in E_{\lambda_{k}}\setminus\{0\}$.
    $u^{i}(x)=\lambda_{k}^{i}\,x$. Donc $(\sum \alpha_{i}u^{i})x
    =\sum\alpha_{i}u^{i}(x) =(\sum\alpha_{i}\lambda_{k}^{i})x=0$. Donc
    $\forall k\in\{1,...,n\}, \sum\alpha_{i}\lambda_{k}^{i}=0$.

  \item
    Le déterminent du système $(*)$ est non nul. Il s'agit donc d'un
    système de Cramer~: il n'a qu'une solution,
    $\alpha_{0}=...=\alpha_{n-1}=0$. La famille $(\mathrm{id}, u, ..., u^{n-1})$
    est donc libre.

  \item
    On a $\dim\mathrm{Vect}(\mathrm{id}, u, ..., u^{n-1})=n=\dim\mathrm{Com}$ et $\mathrm{Vect}(\mathrm{id}, u,
    ..., u^{n-1})\subset\mathrm{Com}$ donc $\mathrm{Vect}(\mathrm{id}, u, ..., u^{n-1})=\mathrm{Com}$

  \end{enumerate}
\fincorrection
\nocorrection
\nocorrection
\nocorrection
\nocorrection
\nocorrection
\nocorrection
\nocorrection
\nocorrection
\correction{002479}
On a une suite récurrente à trois termes reliant les composantes  $v_i$ du vecteur
propre. On calcule le terme général de la suite en résolvant le 
polynôme caractéristique. Les deux constantes sont identifiées en 
écrivant que $v_0=v_{n+1}=0$. On trouve $n+1$ valeurs propres distinctes:
\[
\lambda_k=b+2c\left(\frac{a}{c}\right)^{1/2}cos\left(\frac{2 k \pi}{n+1}\right)
\mbox{~~~pour~~~}k=1,...n
\]
avec le vecteur propre $v^k$ associé, de composantes
\[
v_j^k=\left(\frac{a}{c}\right)^{j/2} sin\left(\frac{2 k j \pi}{n+1}\right)
\mbox{~~~pour~~~}j=1,...n
\]
\fincorrection
\correction{002570}
 Soit $A$ une matrice carr\'ee d'ordre $n$. On suppose que $A$ est inversible et que $\lambda\in\R$ est une valeur propre de $A$.
\begin{enumerate}
  \item {\it D\'emontrons que $\lambda\neq 0$}.
 Si $\lambda=0$ est valeur propre de $A$, alors $\ker A\neq\{0\}$, donc $A$ n'est pas injective et sa matrice ne peut pas \^etre inversible. Par
cons\'equent, $\lambda\neq 0$. 

  \item {\it D\'emontrons que si $\vec x$ est un vecteur propre de $A$ pour la valeur propre $\lambda$ alors il est vecteur propre de $A^{-1}$ de valeur propre $\lambda^{-1}$ }.

Comme $A$ est inversible, on a $A\vec x=\lambda\vec x\iff A^{-1}(A\vec x)=A^{-1}(\lambda \vec x)\iff\vec x=\lambda A^{-1}\vec x $,
d'o\`u $A^{-1}\vec x=\lambda^{-1}\vec x$. Ce qui prouve que $\vec x$ est vecteur propre de $A^{-1}$ de valeur propre $\lambda^{-1}$.
\end{enumerate}
\fincorrection
\correction{002571}
Soit $f$ un endomorphisme de $E$ v\'erifiant $f^2={\rm mathrm{Id}}_E$.
\begin {enumerate}
  \item {\it D\'emontrons que les seules valeurs propres possibles de $f$ sont $1$ et $-1$}.

Si $\lambda$ est une valeur propre de $f$, il existe un vecteur non nul $\vec x\in E$ tel que $f(\vec x)=\lambda\vec x$. On a donc
$$f^2(\vec x)=f(\lambda\vec x)=\lambda f(\vec x)=\lambda^2\vec x.$$
Mais, $f^2={\rm mathrm{Id}}_E$ donc si $\vec x$ est un vecteur propre associ\'e \`a la valeur propre $\lambda$ on a
$$\vec x=f^2(\vec x)=\lambda^2\vec x,$$
d'o\`u $\lambda^2=1$, c'est-\`a-dire (dans $\R$ ou $\C$), $\lambda=1$ ou $\lambda=-1$. ce qui prouve que les seules valeurs propres possibles de $f$
 sont $1$ et $-1$.
  \item {\it  V\'erifions que pour tout $\vec x\in E$, on a}
$$f(\vec x-f(\vec x))=-(\vec x-f(\vec x))\ \ {\hbox{et}}\ \ f(\vec x+f(\vec x))=(\vec x+f(\vec x))$$
Soit $\vec x\in E$, on a
$$f(\vec x-f(\vec x))=f(\vec x)-f^2(\vec x)=f(\vec x)-\vec x=-(\vec x-f(\vec x))$$ et
$$f(\vec x+f(\vec x))=f(\vec x)+f^2(\vec x)=f(\vec x)+\vec x$$

{\it Nous allons en d\'eduire que $f$ admet toujours une valeur propre.}

Supposons que $1$ ne soit pas valeur propre de $f$, alors, $\vec x=f(\vec x)\Rightarrow \vec x=\vec 0$. Or, pour tout $\vec x\in E$, on a
$f(\vec x+f(\vec x))=f(\vec x)+\vec x$, donc pour tout $\vec x\in E$, on a $f(\vec x)+\vec x=\vec 0$, c'est-\`a-dire, $f(\vec x)=-\vec x$.
Ce qui prouve que $-1$ est valeur propre de $f$. On a m\^eme dans ce cas $f=-{\rm mathrm{Id}_E}$. 

Si $-1$ n'est pas valeur propre de $f$, on montre par un raisonnement
analogue que pour tout $\vec x\in E$ on a $f(\vec x)-\vec x=\vec 0$. Ce qui prouve que $1$ est valeur propre de $f$, et dans ce cas $f={\rm mathrm{Id}_E}$.

  \item  {\it D\'emontrons que si $1$ et $-1$ sont valeurs propres, alors $E$ est somme directe des sous-espaces propres correspondants.}

Supposons maintenant que $1$ et $-1$ sont valeurs propres de $f$. Ce sont alors les seules et on a, pour tout $\vec x\in E$,
$$\vec x={\frac{1}{2}}(\vec x+f(\vec x))+{\frac{1}{2}}(\vec x-f(\vec x))$$
Et, quelque soit $\vec x\in E$, $f(\vec x-f(\vec x))=-(\vec x-f(\vec x))$ et $f(\vec x+f(\vec x))=(\vec x+f(\vec x))$, c'est-\`a-dire 
$\vec x+f(\vec x)$ est dans le sous-espace propre associ\'e \`a la valeur propre $1$ et $\vec x-f(\vec x)$ est dans le sous-espace propre associ\'e \`a
la valeur propre $-1$. Par ailleurs on sait que les sous-espaces propres sont en somme directe ( on peut le v\'erifier \'egalement puisque leur intersection
est l'ensemble des vecteurs $\vec x$ tels que $\vec x=-\vec x$, donc r\'eduite au vecteur nul). par cons\'equent $E$ est bien somme directe des sous-espaces
propres correspondants aux valeurs propres $1$ et $-1$.

  \item {\it Traduisons g\'eom\'etriquement le cas $n=2$}.

Rappelons que si il n'y a qu'une valeur propre, $f$ est l'identit\'e ou son oppos\'ee. Dans le cas o\`u $1$ et $-1$ sont valeur propres, leurs sous-espaces
propres sont des droites vectorielles. Soit $u$ un vecteur
propre tel que $f(u)=u$ et $v$ un vecteur propre tel que $f(v)=-v$, alors si $w=au+bv$, $f(w)=au-bv$.
\end {enumerate}
\fincorrection
\correction{002579}
{\it Soit $E$ un espace vectoriel sur un corps $K$ ($K=\R$ ou $\C$), et $u$ un endomorphisme de $E$. On suppose $u$ nilpotent, c'est-\`a-dire qu'il existe un entier strictement positif $n$ tel que $u^n=0$.}

\begin{enumerate}  \item Montrons que $u$ n'est pas inversible. 

On a : $0=\det u^n=(\det u)^n$, d'o\`u $\det u=0$, ce qui prouve que $u$ n'est pas inversible.

  \item D\'eterminons les valeurs propres de $u$ et les sous-espaces propres associ\'es.

Soit $\lambda$ une valeur propre de $u$, il existe alors un vecteur $x\in E$ non nul tel que $u(x)=\lambda x$. Or, $u(x)=\lambda x\Rightarrow u^n(x)=\lambda^n x$. Mais, $u^n(x)=0$ et $x\neq 0$, d'o\`u $\lambda^n=0$ et donc $\lambda=0$. La seule valeur propre possible de $u$ est donc $0$ et c'est une valeur propre car, comme $u$ n'est pas inversible, le noyau de $u$ n'est pas r\'eduit \`a $\{0\}$. L'endomorphisme $u$ admet donc $0$ comme unique valeur propre, le sous-espace propre associ\'e est $\ker u$.
 \end{enumerate}
\fincorrection
\correction{002580}
{\it Soit $M$ la matrice de $\R^4$ suivante}
$$M=\begin{pmatrix}0&1&0&0 \\  2&0&-1&0 \\  0&7&0&6 \\  0&0&3&0\end{pmatrix}$$

\begin{enumerate}  
\item D\'eterminons les valeurs propres de $M$ et ses sous-espaces propres.

Les valeurs propres de $M$ sont les r\'eels $\lambda$ tels que $\det (M-\lambda I)=0$.
$$\det(M-\lambda I)=\begin{vmatrix}-\lambda&1&0&0 \\  2&-\lambda&-1&0 \\  0&7&-\lambda&6 \\  0&0&3&-\lambda\end{vmatrix}=\lambda^4-13\lambda^2+36=(\lambda^2-4)(\lambda^2-9)=(\lambda-2)(\lambda+2)(\lambda-3)(\lambda+3).$$
Les valeurs propres de $M$ sont donc $2,-2,3$ et $-3$. Notons $E_2$, $E_{-2}$, $E_3$ et $E_{-3}$ les sous-espaces propres associ\'es.

\begin{align*}
E_2&=\{X\in\R^4,MX=2X\} \\ &=\left\{(x,y,z,t)\in \R^4,y=2x,2x-z=2y,7y+6t=2z,3z=2t\right\} 
\end{align*}
$${\hbox{or}}\ \left\{\begin{align*}y&=2x \\  2x-z&=2y \\  7y+6t&=2z \\  3z&=2t\end{align*}\right.\iff
\left\{\begin{align*}y&=2x \\  2x-z&=4x \\  14x+9z&=2z \\  3z&=2t\end{align*}\right.\iff
\left\{\begin{align*}y&=2x \\  z&=-2x \\  t&=-3x\end{align*}\right.$$
ainsi, $E_2$ est la droite vectorielle engendr\'ee par le vecteur $u_1=(1,2,-2,-3)$.

\begin{align*}
E_{-2}&=\{X\in\R^4,MX=-2X\} \\ &=\left\{(x,y,z,t)\in \R^4,y=-2x,2x-z=-2y,7y+6t=-2z,3z=-2t\right\} 
\end{align*}
$${\hbox{or}}\ \left\{\begin{align*}y&=-2x \\  2x-z&=-2y \\  7y+6t&=-2z \\  3z&=-2t\end{align*}\right.\iff
\left\{\begin{align*}y&=-2x \\  2x-z&=4x \\  -14x-9z&=2z \\  3z&=-2t\end{align*}\right.\iff
\left\{\begin{align*}y&=-2x \\  z&=-2x \\  t&=3x\end{align*}\right.$$
ainsi, $E_{-2}$ est la droite vectorielle engendr\'ee par le vecteur $u_2=(1,-2,-2,3)$.

\begin{align*}
E_3&=\{X\in\R^4,MX=3X\} \\  &=\left\{(x,y,z,t)\in \R^4,y=3x,2x-z=3y,7y+6t=3z,3z=3t\right\}
\end{align*}
$${\hbox{or}}\ \left\{\begin{align*}y&=3x \\  2x-z&=3y \\  7y+6t&=3z \\  3z&=3t\end{align*}\right.\iff
\left\{\begin{align*}y&=3x \\  2x-z&=9x \\  21x+6t&=3z \\  z&=t\end{align*}\right.\iff
\left\{\begin{align*}y&=3x \\  z&=-7x \\  t&=-7x\end{align*}\right.$$
ainsi, $E_3$ est la droite vectorielle engendr\'ee par le vecteur $u_3=(1,3,-7,-7)$.

\begin{align*}
E_{-3}&=\{X\in\R^4,MX=-3X\} \\ &=\left\{(x,y,z,t)\in \R^4,y=-3x,2x-z=-3y,7y+6t=-3z,3z=-3t\right\}
\end{align*}
$${\hbox{or}}\ \left\{\begin{align*}y&=-3x \\  2x-z&=-3y \\  7y+6t&=-3z \\  3z&=-3t\end{align*}\right.\iff
\left\{\begin{align*}y&=-3x \\  2x-z&=9x \\  -21x-6z&=-3z \\  z&=-t\end{align*}\right.\iff
\left\{\begin{align*}y&=-3x \\  z&=-7x \\  t&=7x\end{align*}\right.$$
ainsi, $E_{-3}$ est la droite vectorielle engendr\'ee par le vecteur $u_4=(1,-3,-7,7)$.


  \item Montrons que $M$ est diagonalisable. 

La matrice $M$ admet quatre valeurs propres distinctes, ce qui prouve que les quatres vecteurs propres correspondants sont lin\'eairement ind\'ependants. En effet, les vecteurs $u_1,u_2,u_3$ et $u_4$ d\'etermin\'es en $1)$ forment une base de $\R^4$. L'endomorphisme dont la matrice est $M$ dans la base canonique de $\R^4$ est repr\'esent\'e par une matrice diagonale dans la base
$(u_1,u_2,u_3,u_4)$ puisque $Mu_1=2u_1, Mu_2=-2u_2$, 
$Mu_3=3u_3$ et $Mu_4=-3u_4$.

  \item D\'eterminons une base de vecteurs propres et $P$ la matrice de passage.

Une base de vecteurs propres a \'et\'e d\'etermin\'ee dans les questions pr\'ec\'edentes. C'est la base $(u_1,u_2,u_3,u_4)$ et la matrice de passage est la matrice 
$$P=\begin{pmatrix}1&1&1&1 \\  2&-2&3&-3 \\  -2&-2&-7&-7 \\  -3&3&-7&7\end{pmatrix}$$
  \item On a $D=P^{-1}MP$, pour $k\in\N$ exprimons $M^k$ en fonction de $D^k$, puis calculons $M^k$.

On a $$D=\begin{pmatrix}2&0&0&0 \\  0&-2&0&0 \\  0&0&3&0 \\  0&0&0&-3\end{pmatrix}\ {\hbox{donc}}\ 
D^k=\begin{pmatrix}2^k&0&0&0 \\  0&(-2)^k&0&0 \\  0&0&3^k&0 \\  0&0&0&(-3)^k\end{pmatrix}.$$

Mais, $M=PDP^{-1}$, d'o\`u, pour $k\in\N$, $M^k=(PDP^{-1})^k=PD^kP^{-1}$. 

Pour calculer $M^k$, il faut donc d\'eterminer la matrice $P^{-1}$ qui exprime les coordonn\'ees des vecteurs de la base canonique de $\R^4$ dans la base $(u_1,u_2,u_3,u_4)$.

On r\'esout le syst\`eme, et on a :
$$\begin{align*}u_1&=i+2j-2k-3l \\  u_2&=i-2j-2k+3l \\  u_3&=i+3j-7k-7t \\  u_4&=i-3j-7k+7l\end{align*}\iff
\left\{\begin{align*}i&={\frac{1}{10}}(7u_1+7u_2-2u_3-2u_4) \\  
j&={\frac{1}{10}}(7u_1-7u_2-3u_3+3u_4) \\  
k&={\frac{1}{10}}(u_1+u_2-u_3-u_4) \\  
l&={\frac{1}{10}}(3u_1-3u_2-2u_3+2u_4)\end{align*}\right.$$
d'o\`u
$$P^{-1}={\frac{1}{10}}\begin{pmatrix}7&7&1&3 \\  7&-7&1&-3 \\  -2&-3&-1&-2 \\  -2&3&-1&2\end{pmatrix}$$
et 
$$M^k=PD^kP^{-1}={\frac{1}{10}}\begin{pmatrix}1&1&1&1 \\  2&-2&3&-3 \\  -2&-2&-7&-7 \\  -3&3&7&7\end{pmatrix}
\begin{pmatrix}2^k&0&0&0 \\  0&(-2)^k&0&0 \\  0&0&3^k&0 \\  0&0&0&(-3)^k\end{pmatrix}
\begin{pmatrix}7&7&1&3 \\  7&-7&1&-3 \\  -2&-3&-1&-2 \\  -2&3&-1&2\end{pmatrix}.$$
 \end{enumerate}

\fincorrection
\correction{002594}
 Soit $A$ la matrice suivante 
$$A=\begin{pmatrix}3&0&-1 \\  2&4&2 \\  -1&0&3\end{pmatrix}$$ 
\begin{enumerate}
 \item {\it On d\'etermine et on factorise le polyn\^ome caract\'eristique de $A$.}

Soit $P_A$ le polyn\^ome caract\'eristique de $A$, on a 
\begin{align*}
P_A(X)=\begin{vmatrix}3-X&0&-1 \\  2&4-X&2 \\  -1&0&3-X\end{vmatrix}&=(4-X)\begin{vmatrix}3-X&-1 \\ -1&3-X\end{vmatrix} \\ &=(4-X)(X^2-6X+8) \\ &=(4-X)(X-4)(X-2) \\ &=(2-X)(4-X)^2
\end{align*}
 \item {\it On d\'emontre que $A$ est diagonalisable et on d\'etermine une matrice $D$ diagonale et une matrice $P$ inversible telles  $A=PDP^{-1}$.}

Le polyn\^ome $P_A$ admet deux racines, donc la matrice $A$ admet deux valeurs propres, $\lambda_1=2$, valeur propre simple et $\lambda_2=4$, valeur propre double. D\'eterminons les sous-espaces propres associ\'es.

Notons $E_1=\{\vec V=(x,y,z)/ A\vec V=2\vec V\}$, on r\'esout alors le syst\`eme
$$\left\{\begin{align*}3x-z&=2x \\  2x+4y+2z&=2y \\  -x+3z&=2z\end{align*}\right.\iff\left\{\begin{align*}z&=x \\  y&=-2x\end{align*}\right.$$
Le sous-espace propre $E_1$ associ\'e \`a la valeur propre $2$ est une droite vectorielle, dont un vecteur directeur est $\vec e_1=(1,-2,1)$.

Notons $E_2=\{\vec V=(x,y,z)/ A\vec V=4\vec V\}$, on r\'esout alors le syst\`eme
$$\left\{\begin{align*}3x-z&=4x \\  2x+4y+2z&=4y \\  -x+3z&=4z\end{align*}\right.\iff z=-x$$
Le sous-espace propre $E_2$ associ\'e \`a la valeur propre $4$ est le plan vectoriel, d'\'equation 

$z=-x$ dont une base est donn\'ee, par exemple par les vecteurs $\vec e_2=(0,1,0)$ et 

$\vec e_3=(1,0,-1)$. Remarquons que l'on pouvait lire directement sur la matrice $A$, le fait que le vecteur $\vec e_2$ est vecteur propre associ\'e \`a la valeur propre $4$.

Les dimensions des sous-espaces propres sont \'egales aux multiplicit\'es des valeurs propres correspondantes, par cons\'equent, l'espace $\R^3$ admet une base de vecteurs propres et la matrice $A$ est diagonalisable.

Notons $P$ la matrice de passage, on a
$$P=\begin{pmatrix}1&0&1 \\ -2&1&0 \\  1&0&-1\end{pmatrix}$$
et, si $D$ est la matrice diagonale
$$D=\begin{pmatrix}2&0&0 \\ 0&4&0 \\  0&0&4\end{pmatrix},$$
on a la relation
$$A=PDP^{-1}.$$
 \item {\it On donne en le justifiant, mais sans calculs, le polyn\^ome minimal de $A$.}

La matrice $A$ est diagonalisable, donc son polyn\^ome minimal n'a que des racines simples, par ailleurs les racines du polyn\^ome minimal sont exactement les valeurs propres de $A$ et le polyn\^ome minimal est un polyn\^ome unitaire qui divise le polyn\^ome caract\'eristique. On a donc
$$Q_A(X)=(X-2)(X-4).$$
 \item  {\it On calcule $A^n$ pour $n\in\N$.}

On a vu, dans la question 2), que $A=PDP^{-1}$, on a donc, pour $n\in\N$, $A^n=P^{-1}D^nP$, or
$$D^n=\begin{pmatrix}2^n&0&0 \\ 0&4^n&0 \\  0&0&4^n\end{pmatrix},$$
il nous reste \`a calculer $P^{-1}$. On sait que $P^{-1}={\frac{1}{\det P}}^t\!\!\tilde P$, d'o\`u
$$\det P=-2,\ \tilde P=\begin{pmatrix}-1&-2&-1 \\ 0&-2&0 \\  -1&-2&1\end{pmatrix}\ {\hbox{et}}\ 
P^{-1}=-{\frac{1}{2}}\begin{pmatrix}-1&0&-1 \\ -2&-2&-2 \\  -1&0&1\end{pmatrix}.$$
On a donc
\begin{align*}A^n&=-{\frac{1}{2}}\begin{pmatrix}1&0&1 \\ -2&1&0 \\  1&0&-1\end{pmatrix}\begin{pmatrix}2^n&0&0 \\ 0&4^n&0 \\  0&0&4^n\end{pmatrix}\begin{pmatrix}-1&0&-1 \\ -2&-2&-2 \\  -1&0&1\end{pmatrix} \\ &={\frac{1}{2}}
\begin{pmatrix}2^n+4^n&0&2^n-4^n \\ 2(4^n-2^n)&2.4^n&2(4^n-2^n) \\  2^n-4^n&0&2^n+4^n\end{pmatrix}.
\end{align*}
\end{enumerate}
\fincorrection
\correction{002595}
 Soit $A$ la matrice suivante
$$A=\begin{pmatrix}1&1 \\  2&1\end{pmatrix}$$
\begin{enumerate}
 \item {\it On calcule le polyn\^ome caract\'eristique et on d\'etermine les valeurs propres de $A$.}

Le polyn\^ome caract\'eristique $P_A(X)$ est \'egal \`a
$$P_A(X)=\begin{vmatrix}1-X&1 \\  2&1-X\end{vmatrix}=(1-X)^2-2=X^2-2X-1.$$
Calculons ses racines, le discriminant r\'eduit de ce polyn\^ome du second degr\'e est \'egal \`a $\Delta'=(-1)^2-(-1)=2$, les racines sont donc
$$\lambda_1=1+\sqrt2\ \ {\hbox{et}}\ \ \lambda_2=1-\sqrt2,$$
ce sont les valeurs propres de $A$.
 \item  On note $\lambda_1>\lambda_2$ les valeurs propres de $A$, $E_1$ et $E_2$ les sous-espaces propres associ\'es. {\it On d\'etermine une base $(\vec{\varepsilon_1},\vec{\varepsilon_2})$ de $\R^2$ telle que 
$\vec{\varepsilon_1}\in E_1$, $\vec{\varepsilon_2}\in E_2$, les deux vecteurs ayant des coordonn\'ees de la forme $(1,y)$.}

On cherche $\vec\varepsilon_1=\begin{pmatrix}1 \\  y\end{pmatrix}$ tel que $A.\vec\varepsilon_1=(1+\sqrt2)\vec\varepsilon_1,$
on calcule donc $y$ tel que
$$\begin{pmatrix}1&1 \\  2&1\end{pmatrix}\begin{pmatrix}1 \\  y\end{pmatrix}=(1+\sqrt2)\begin{pmatrix}1 \\  y\end{pmatrix}$$
ce qui \'equivaut \`a
$$\left\{\begin{align*}1+y&=1+\sqrt2 \\  2+y&=(1+\sqrt2)y\end{align*}\right.$$
d'o\`u $y=\sqrt2$ et $\vec\varepsilon_1=\begin{pmatrix}1 \\  \sqrt2\end{pmatrix}$.

On cherche $\vec\varepsilon_2=\begin{pmatrix}1 \\  y\end{pmatrix}$ tel que $A.\vec\varepsilon_2=(1-\sqrt2)\vec\varepsilon_2,$
on calcule donc $y$ tel que
$$\begin{pmatrix}1&1 \\  2&1\end{pmatrix}\begin{pmatrix}1 \\  y\end{pmatrix}=(1-\sqrt2)\begin{pmatrix}1 \\  y\end{pmatrix}$$
ce qui \'equivaut \`a
$$\left\{\begin{align*}1+y&=1-\sqrt2 \\  2+y&=(1-\sqrt2)y\end{align*}\right.$$
d'o\`u $y=-\sqrt2$ et $\vec\varepsilon_2=\begin{pmatrix}1 \\  -\sqrt2\end{pmatrix}$.
 \item  Soit $\vec x$ un vecteur de $\R^2$, on note $(\alpha,\beta)$ ses coordonn\'ees dans la base
$(\vec{\varepsilon_1},\vec{\varepsilon_2})$.{\it On d\'emontre que, pour $n\in\N$, on a
$$A^n\vec x=\alpha\lambda_1^n\vec{\varepsilon_1}+\beta\lambda_2^n\vec{\varepsilon_2}.$$}

On a $\vec x=\alpha\vec\varepsilon_1+\beta\vec\varepsilon_2$, d'o\`u, par lin\'earit\'e
$A\vec x=\alpha A\vec\varepsilon_1+\beta A\vec\varepsilon_2$ et $A^n\vec x=\alpha A^n\vec\varepsilon_1+\beta A^n\vec\varepsilon_2$. Or, on montre, par r\'ecurrence sur $n$, que $A^n\vec\varepsilon_1=\lambda_1^n\vec\varepsilon_1$ et de m\^eme $A^n\vec\varepsilon_2=\lambda_2^n\vec\varepsilon_2$. Pour $n=1$, c'est la d\'efinition des vecteurs propres. Soit $n$ fix\'e, tel que $A^n\vec\varepsilon_1=\lambda_1^n\vec\varepsilon_1$, on a alors
$A^{n+1}\vec\varepsilon_1=A.A^n\vec\varepsilon_1=\lambda_1^nA\vec\varepsilon_1=\lambda_1^{n+1}\vec\varepsilon_1$. Ainsi, pour tout $n\in\N$, on a $A^n\vec\varepsilon_1=\lambda_1^n\vec\varepsilon_1$, et, de m\^eme, $A^n\vec\varepsilon_2=\lambda_2^n\vec\varepsilon_2$. D'o\`u le r\'esultat.
 \item  Notons $A^n\vec x=\begin{pmatrix}a_n \\  b_n\end{pmatrix}$ dans la base canonique de $\R^2$. {\it On exprime $a_n$ et $b_n$ en fonction de $\alpha$, $\beta$, $\lambda_1$ et $\lambda_2$ et on en d\'eduit que, si $\alpha\neq 0$, la suite ${\frac{b_n}{a_n}}$ tend vers $\sqrt2$ quand $n$ tend vers $+\infty$.}

D'apr\`es la question pr\'ec\'edente et les vecteurs $\vec \varepsilon_1$ et $\vec\varepsilon_2$ obtenus en 2) on a
$$A^n\vec x=\alpha\lambda_1^n\begin{pmatrix}1 \\ \sqrt2\end{pmatrix}+\beta\lambda_2^n\begin{pmatrix}1 \\ -\sqrt2\end{pmatrix}=\begin{pmatrix}a_n \\  b_n\end{pmatrix}$$
d'o\`u
$$\left\{\begin{align*}a_n&=\alpha\lambda_1^n+\beta\lambda_2^n \\  b_n &=\sqrt2(\alpha\lambda_1^n-\beta\lambda_2^n)\end{align*}\right.$$
On suppose $\alpha\neq 0$, pour $n$ assez grand, on a 
$${\frac{b_n}{a_n}}=\sqrt2{\frac{\alpha\lambda_1^n-\beta\lambda_2^n}{\alpha\lambda_1^n+\beta\lambda_2^n}},$$
or, 
$$|\lambda_1|=|1+\sqrt2|>1\Longrightarrow \lim_{n\rightarrow+\infty}\lambda_1^n=+\infty,$$
et
$$|\lambda_2|=|1-\sqrt2|<1\Longrightarrow \lim_{n\rightarrow+\infty}\lambda_2^n=0.$$
D'o\`u l'\'equivalence
$${\frac{b_n}{a_n}}=\sqrt2{\frac{\alpha\lambda_1^n-\beta\lambda_2^n}{\alpha\lambda_1^n+\beta\lambda_2^n}}
\sim\sqrt2{\frac{\alpha\lambda_1^n}{\alpha\lambda_1^n}}=\sqrt2.$$
On a donc bien
$$\lim_{n\rightarrow+\infty}{\frac{b_n}{a_n}}=\sqrt2.$$
 \item {\it On explique, sans calcul, comment obtenir, \`a partir des questions pr\'ec\'edentes, une approximation de $\sqrt2$ par une suite  de nombres rationnels.}

La matrice $A$ est \`a coefficients entiers, aussi, pour tout $n\in\N$, la matrice $A^n$ est \`a coefficients entiers. Si l'on choisit un vecteur $\vec x$ \`a coordonn\'ees enti\`eres dans la base canonique de $\R^2$, alors les coordonn\'ees $a_n$ et $b_n$ du vecteur $A^n\vec x$ sont des entiers et elles nous fournissent une suite $\displaystyle{\frac{b_n}{a_n}}$ de nombres rationnels qui tend vers $\sqrt2$.
\end{enumerate}
\fincorrection
\correction{002596}
Soit $P(X)$ un polyn\^ome de $\C[X]$, soit $A$ une matrice de $M_n(\C)$. On note $B$ la matrice : $B=P(A)\in M_n(\C)$.
\begin{enumerate}
 \item {\it On d\'emontre que si $\vec x$ est un vecteur propre de $A$ de valeur propre $\lambda$, alors $\vec x$ est un vecteur propre de $B$ de valeur propre $P(\lambda)$.}

Soit $\vec x\neq0$ tel que $A\vec x=\lambda\vec x$, notons $P(X)=\sum_{k=0}^da_kX^k$, on a
$$P(A)=\sum_{k=0}^da_kA^k=a_0I_n+a_1 A+\dots+a_d A^d,$$
o\`u $I_n$ d\'esigne la matrice unit\'e.

Or, pour $k\in\N$, on a $A^k\vec x=\lambda^k\vec x$, d'o\`u
$$B\vec x=P(A)\vec x=\sum_{k=0}^da_kA^k\vec x=\left(\sum_{k=0}^da_k\lambda^k\right)\vec x=P(\lambda)\vec x,$$
ce qui prouve que $\vec x$ est un vecteur propre de la matrice $B=P(A)$ pour la valeur propre $P(\lambda)$.

 \item  Le but de cette question est de d\'emontrer que les valeurs propres de $B$ sont toutes de la forme $P(\lambda)$, avec $\lambda$ valeur propre de $A$.

 Soit $\mu\in\C$, on d\'ecompose le polyn\^ome $P(X)-\mu$ en produit de facteurs de degr\'e $1$ :
$$P(X)-\mu=a(X-\alpha_1)\cdots(X-\alpha_r).$$
   \begin{enumerate}
   \item {\it On d\'emontre que} 
$\det(B-\mu I_n)=a^n\det(A-\alpha_1I_n)\cdots\det(A-\alpha_rI_n).$

Compte tenu de la d\'ecomposition du polyn\^ome $P(X)-\mu$, on a 
$$P(A)-\mu I_n=aI_n(A-\alpha_1I_n)\cdots(A-\alpha_rI_n)$$
d'o\`u
$$\det(B-\mu I_n)=a^n\det(A-\alpha_1I_n)\cdots\det(A-\alpha_rI_n)$$
car le d\'eterminant est une forme multilin\'eaire (d'o\`u le $a^n$) et le d\'eterminant d'un produit de matrices est \'egal au produit de leurs d\'eterminants.
   \item {\it On en d\'eduit que si $\mu$ est valeur propre de $B$, alors il existe une valeur propre $\lambda$ de $A$ telle que $\mu=P(\lambda)$.}

Si $\mu$ est une valeur propre de $B$, alors, par d\'efinition, $\det(B-\mu I_n)=0$, ainsi, compte tenu de la question pr\'ec\'edente, il existe un $\alpha_i$, $1\leq i\leq r$, tel que $\det(A-\alpha_iI_n)=0$, c'est-\`a-dire que l'un des $\alpha_i$, $1\leq i\leq r$, est valeur propre de $A$. Or, pour $1\leq i\leq r$, $P(\alpha_i)-\mu=0$ donc si $\mu$ est une valeur propre de $B$, on a $\mu=P(\alpha_i)$ o\`u $\alpha_i$ est une valeur propre de $A$.

  \end{enumerate}
 \item On note $S_A$ l'ensemble des valeurs propres de $A$. {\it On d\'emontre que }
$$S_B=\{P(\lambda)/\  \lambda\in S_A\}.$$

Soit $\lambda$ une valeur propre de $A$, on a d\'emontr\'e en 1) que $P(\lambda)$ est une valeur propre de $B$, ainsi $\{P(\lambda)/\  \lambda\in S_A\}\subset S_B$. R\'eciproquement, si $\mu$ est une valeur propre de $B$ alors, d'apr\`es 2), il existe une valeur propre $\lambda$ de $A$ telle que $\mu=P(\lambda)$, ainsi on a $S_B\subset\{P(\lambda)/\  \lambda\in S_A\}$, d'o\`u l'\'egalit\'e des deux ensembles.
 \item  Soient $\lambda_1,\dots,\lambda_r$ les valeurs propres de $A$ et soit $Q(X)$ le polyn\^ome
$$Q(X)=(X-\lambda_1)\cdots(X-\lambda_r),$$
on note $C$ la matrice $C=Q(A)$. 
   \begin{enumerate}
   \item {\it On d\'emontre que $S_C=\{0\}$}.

D'apr\`es la question pr\'ec\'edente, on a $S_C=\{Q(\lambda)/\  \lambda\in S_A\}.$ Or, par d\'efinition du polyn\^ome $Q(X)$, on a $Q(\lambda)=0$ pour toute valeur propre $\lambda$ de $A$, ainsi, $S_C=\{0\}$.
   \item {\it On en d\'eduit que le polyn\^ome caract\'eristique de $C$ est $(-1)^nX^n$ et que $C^n=0$.}

Les valeurs propres de $C$ sont les racines de son polyn\^ome caract\'eristique, or $C$ admet une unique valeur propre : $0$, ainsi $P_C(X)=(-1)^nX^n$. Par ailleurs, d'apr\`es le th\'eor\`eme de Cayley-Hamilton, on a $P_C(C)=0$, ainsi $(-1)^nC^n=0$, donc $C^n=0$.
  \end{enumerate}
\end{enumerate}
\fincorrection
\nocorrection
\nocorrection
\nocorrection
\nocorrection
\nocorrection
\nocorrection
\nocorrection
\nocorrection
\correction{003379}
2. $D^{-1} = BC^{-1}A + I_p$.
 \fincorrection
\correction{003499}
\begin{enumerate}
  \item $0$ et les racines de $6\lambda^2 -6n\lambda -n(n-1)(2n-1) = 0$.
  \item $\sin\alpha+\sin2\alpha,\quad -\sin\alpha,\quad -\sin2\alpha$.
\end{enumerate}
 \fincorrection
\correction{003500}
\begin{enumerate}
  \item $\mathrm{rg}(A) = 2  \Rightarrow  0$ est valeur propre d'ordre au moins $n-2$.
    $E_0 = \{ a_1x_1 + \dots + a_{n-1}x_{n-1} = x_n = 0 \}$.\par
    vp $\lambda \ne 0$ :
    $\lambda^2 - a_n\lambda - (a_1^2 + \dots + a_{n-1}^2) = 0$.
    Il y a deux racines distinctes,
    $E_\lambda = \text{vect}((a_1,\dots,a_{n-1},\lambda))$.
  \item $A$ est diagonale. vp = 0 et $a_n$.
\end{enumerate}
 \fincorrection
\correction{003501}
\begin{enumerate}
  \item $D_n = 2\cos\theta D_{n-1} - D_{n-2}  \Rightarrow 
              D_n = \frac{\sin(n+1)\theta}{\sin\theta}$.
  \item $-2\cos\left(\frac{k\pi}{n+1}\right)$, $1 \le k \le n$.
\end{enumerate}
 \fincorrection
\correction{003502}
Soit $P_n(x)$ le polynôme caractéristique de~$x$ et $Q_n(x)$ celui de la
matrice obtenue à partir de~$A$ en rempla\c cant le premier~$1$ par $2$.
On a les relations de récurrence~:
$$P_n(x) = (1-x)Q_{n-1}(x) - Q_{n-2}(x),
  \qquad
  Q_n(x) = (2-x)Q_{n-1}(x) - Q_{n-2}(x).$$
D'où pour $x\notin\{0,4\}$~:
$$P_n(x) = \frac{(1-\alpha)(1-\alpha^{2n})}{\alpha^n(1+\alpha)},  \qquad\text{avec }
x = 2 - \alpha - \frac1\alpha.$$
Les valeurs propres de~$A$ autres que $0$ et $4$ sont les réels $x_k=2(1-\cos(k\pi/n))$
avec $0< k< n$ et $0$ est aussi valeur propre (somme des colonnes nulle) donc
il n'y en a pas d'autres.
 \fincorrection
\correction{003506}
$\lambda = 0$ :
	     $E_0 = \{\vec x$ tq $x_1+\dots+x_q+x_{n-q+1}+\dots+x_n = 0\}$,\par
	     $\lambda = 2\min(p,q)$ : $E_\lambda = \text{vect}
	     ((\underbrace{1,\dots,1}_p,0,\dots,0,\underbrace{1,\dots,1}_p))$.
\fincorrection
\nocorrection
\nocorrection
\nocorrection
\correction{003512}
$u(X^k) = -kX^k + (k-2n)X^{k+1}  \Rightarrow $ la matrice de $u$ est triangulaire
         inférieure. $\text{Spec}(u) = \{0,-1,\dots,-2n\}$.\par
         $\lambda = -k$ : Résoudre l'équation différentielle
	         $ \Rightarrow  P = cX^k(X-1)^{2n-k}$.
\fincorrection
\correction{003513}
$\alpha^3 : (X-\beta)(X-\gamma)$,  \quad
         $\beta^3  : (X-\alpha)(X-\gamma)$, \quad
         $\gamma^3 : (X-\alpha)(X-\beta)$.
\fincorrection
\correction{003514}
$\lambda = 1\phantom-$ : $P = Q\bigl((X-1)^2\bigr)$. \par
         $\lambda =-1$          : $P = (X-1)Q\bigl((X-1)^2\bigr)$.
\fincorrection
\correction{003515}
$\lambda = 1$ : $P = aX+b$.
\fincorrection
\correction{003516}
$M = \begin{pmatrix}0 &-2a &-a^2 &\dots	&-a^n \cr
			     &2   &-2a	&	&(0)  \cr
			     &	  &3	&\ddots       \cr
			     &	  &	&\ddots &-na  \cr
			  (0)&	  &	&	&n+1  \cr\end{pmatrix}$.\par
	     $\mathrm{Ker} f = \{\text{polynômes constants}\}$,
	     $\Im f = \{\text{polynômes divisibles par } X-a \}$.\par
	     Valeurs propres : $0,2,3,\dots,n+1$.
	Pour $2\le k\le n+1$, $E_k = \text{vect}\bigl((X-a)^{k-1}\bigr)$.     
\fincorrection
\correction{005655}
Soit $P=aX^3+bX^2+cX+d\in\Rr_3[X]$

\begin{center} 
$AP - (X^4-X)P= (X-1)P = aX^4+(b-a)X^3+(c-b)X^2+(d-c)X-d = a(X^4-X)+(b-a)X^3+(c-b)X^2+(a+d-c)X-d$.
\end{center}

et donc $AP=(X^4-X)(P+a)+(b-a)X^3+(c-b)X^2+(a+d-c)X-d$ et donc $f(P) = (b-a)X^3+(c-b)X^2+(a+d-c)X-d$. Par suite, $f$ est un endomorphisme de $E$ et la matrice de $f$ dans la base canonique $(1,X,X^2,X^3)$ de $E$ est

\begin{center}
$A=\left(
\begin{array}{cccc}
-1&0&0&0\\
1&-1&0&1\\
0&1&-1&0\\
0&0&1&-1
\end{array}
\right)$.
\end{center}

puis

\begin{align*}\ensuremath
\chi_A&=\left|
\begin{array}{cccc}
-1-X&0&0&0\\
1&-1-X&0&1\\
0&1&-1-X&0\\
0&0&1&-1-X
\end{array}
\right|= (-1-X)\left|
\begin{array}{ccc}
-1-X&0&1\\
1&-1-X&0\\
0&1&-1-X
\end{array}
\right|\\
 &= -(X+1)(-(X+1)^3+1) =X(X+1)(X^2+3X+3).
\end{align*}

 
$A$ admet quatre valeurs propres simples dans $\Cc$, deux réelles 0 et -1 et deux non réelles $-1+j$ et $-1+j^2$. $\chi_f$ n'est pas scindé sur $\Rr$ et donc $f$ n'est pas diagonalisable.

\textbullet~Soit $P\in E$. $P\in\text{Ker}f\Leftrightarrow b-a = c-b = a+d-c = -d = 0\Leftrightarrow a=b=c\;\text{et}\;d=0$. $\text{Ker}f=\text{Vect}(X^3+X^2+X)$.

\textbullet~Soit $P\in E$. $P\in\text{Ker}(f+Id)\Leftrightarrow b=c=a+d = 0\Leftrightarrow b=c=0\;\text{et}\;d=-a$. $\text{Ker}(f+Id)= \text{Vect}(X^3-1)$.

\textbullet~$\text{rg}(f) = 3$ et immédiatement $\text{Im}f =\text{Vect}(X-1,X^2-X,X^3-X^2)$.

Si $\Kk=\Cc$, on peut continuer :

$P\in\text{Ker}(f+(1-j)Id)\Leftrightarrow b-ja = c-jb = a + d - jc = - jd = 0\Leftrightarrow b = ja,\;c = j^2a\;\text{et}\;d=0$.

Donc  
$\text{Ker}(f+(1-j)Id)=\text{Vect}(X^3+jX^2+j^2X)$ et en conjuguant $\text{Ker}(f+(1-j^2)Id) =\text{Vect}(X^3+j^2X^2+jX)$.

\textbf{Remarque.} $B =X(X-1)(X-j)(X-j^2)$ et on a trouvé pour base de vecteurs propres les quatre polynômes de \textsc{Lagrange} $X^3-1 = (X-1)(X-j)(X-j^2)$ puis $X^3+X^2+X = X(X-j)(X-j^2)$ puis $X^3+jX^2+j^2X =X(X-1)(X-j^2)$ et enfin $X^3+j^2X^2+jX = X(X-1)(X-j )$. C'est une généralité. On peut montrer que si $E =\Cc_n[X]$ et si $B$ a $n+1$ racines deux à deux distinctes dans $\Cc$ alors $f$ est diagonalisable et une base de vecteurs propres est fournie par les polynômes de \textsc{Lagrange} associés aux racines de $B$ et ceci pour un polynôme $A$ quelconque.
\fincorrection
\correction{005660}
\textbf{1er cas.} Supposons $\alpha=\beta= 0$ et donc $uv = vu$. Puisque $E$ est un $\Cc$-espace de dimension finie non nulle, $u$ admet au moins une valeur propre que l'on note $\lambda$. Le sous-espace propre $E_\lambda$ correspondant n'est pas réduit à $\{0\}$, est stable par $u$ et d'autre part stable par $v$ car $u$ et $v$ commutent. On note $u'$ et $v'$ les restrictions de $u$ et $v$ au sous-espace $E_\lambda$. $u'$ et $v'$ sont des endomorphismes de $E_\lambda$. De nouveau, $E_\lambda$ est un $\Cc$-espace de dimension finie non nulle et donc $v'$ admet au moins un vecteur propre $x_0$. Par construction, $x_0$ est un vecteur propre commun à $u$ et $v$.

\textbf{2ème cas.} Supposons par exemple $\alpha\neq0$. 

\begin{align*}\ensuremath
uv - vu =\alpha u+\mu v&\Leftrightarrow(\alpha u+\beta v)\circ\frac{1}{\alpha}v -\frac{1}{\alpha}v\circ(\alpha u+\beta v) =\alpha u+\beta v\\
 &\Leftrightarrow fg - gf = f\;\text{en posant}\;f =\alpha u+\beta v\;\text{et}\;g =\frac{1}{\alpha}v.
\end{align*}

On va chercher un vecteur propre commun à $u$ et $v$ dans le noyau de $f$. Montrons tout d'abord que $\text{Ker}f$ n'est pas nul (on sait montrer que $f$ est en fait nilpotent (exercice \ref{ex:rou9}) mais on peut montrer directement une propriété un peu moins forte).

Si $f$ est inversible, l'égalité $fg - gf = f$ fournit $(g+Id)\circ f = f\circ g$ et donc $g+Id = f\circ g\circ f^{-1}$. Par suite, $g$ et $g+Id$ ont même polynôme caractéristique ou encore, si $\lambda$ est valeur propre de $g$ alors $\lambda+1$ est encore valeur propre de $g$. Mais alors $\lambda+2$, $\lambda+3$... sont aussi valeurs propres de $g$ et $g$ a une infinité de valeurs propres deux à deux distinctes. Ceci est exclu et donc $\text{Ker}f$ n'est pas réduit à $\{0\}$.

 
Maintenant, si $x$ est un vecteur de $\text{Ker}f$, on a $f(g(x))= g(f(x))+f(x) = 0$ et $g(x)$ est dans $\text{Ker}f$. Donc $g$ laisse $\text{Ker}f$ stable et sa restriction à $\text{Ker}f$ est un endomorphisme de $\text{Ker}f$ qui admet au moins une valeur propre et donc au moins un vecteur propre. Ce vecteur est bien un vecteur propre commun à $f$ et $g$. 

Enfin si $x$ est vecteur propre commun à $f$ et $g$ alors $x$ est vecteur propre de $v =\frac{1}{\alpha}g$ et de $u=\frac{1}{\alpha}(f-\beta v)$. $x$ est un vecteur propre commun à $u$ et $v$.
\fincorrection
\correction{005674}
\begin{enumerate}
 \item  Soit $f\in C^0(\Rr,\Rr)$. Soit $F$ une primitive de $f$ sur $\Rr$. Pour tout $x\in\Rr^*$, on a $(\varphi(f))(x)=\frac{F(x)-F(0)}{x-0}$.

$F$ est continue sur $\Rr$ donc $\varphi(f)$ est continue sur $\Rr^*$. De plus, $F$ étant dérivable en $0$

\begin{center} 
$\displaystyle\lim_{\substack{x\rightarrow0\\
x\neq0}}(\varphi(f))(x) =\displaystyle\lim_{\substack{x\rightarrow0\\
x\neq0}}\frac{F(x)-F(0)}{x-0}= F'(0) = f(0) =(\varphi(f))(0)$.
\end{center}

Finalement $\varphi(f)$ est continue sur $\Rr$. Ainsi, $\varphi$ est une application de $E$ dans $E$. La linéarité de $\varphi$ est claire et finalement

\begin{center}
\shadowbox{
$\varphi\in\mathcal{L}(C^0(\Rr,\Rr))$.
}
\end{center}

\item  Si $f$ est dans $\text{Ker}(\varphi)$ alors $f(0) = 0$ et pour tout $x$ non nul, $\int_{0}^{x}f(t)\;dt=0$. Par dérivation on obtient $\forall x\in\Rr^*$, $f(x) = 0$ ce qui reste vrai pour $x=0$ et donc $f = 0$. Finalement $\text{Ker}(\varphi)=\{0\}$ et $\varphi$ est injective.

$\varphi$ n'est pas surjective car pour toute $f\in E$, $\varphi(f)$ est de classe $C^1$ sur $\Rr^*$. Mais alors par exemple, l'application $g~:~x\mapsto|x-1|$ est dans $E$ mais n'est pas dans $\text{Im}(\varphi)$.

\begin{center}
\shadowbox{
$\varphi$ est injective et n'est pas surjective.
}
\end{center}

\item  On cherche $\lambda\in\Rr$ et $f$ continue sur $\Rr$ et non nulle telle que $\forall x\in\Rr,\;(\varphi(f))(x) =\lambda f(x)$. D'après la question précédente, $0$ n'est pas valeur propre de $\varphi$ et donc nécessairement $\lambda\neq0$.

Pour $x = 0$, nécessairement $f(0) =\lambda f(0)$ et donc ou bien  $\lambda= 1$ ou bien $f(0) = 0$.

On doit avoir pour tout $x\in\Rr^*$, $f(x)=\frac{1}{\lambda x}\int_{0}^{x}f(t)\;dt$. $f$ est nécessairement dérivable sur $\Rr^*$. Pour tout $x\in\Rr^*$, on a $\int_{0}^{x}f(t)\;dt=\lambda xf(x)$ et par dérivation, on obtient pour $x\in\Rr^*$,

\begin{center}
$f(x) =\lambda(xf'(x)+f(x))$.
\end{center}

Soit $I$ l'un des deux intervalles $]-\infty,0[$ ou $]0,+\infty[$.

\begin{align*}\ensuremath
\forall x\in I,\;f(x) =\lambda(xf'(x)+f(x))&\Rightarrow\forall x\in I,\;f'(x)+\frac{\lambda-1}{\lambda x}f(x) = 0\\
 &\Rightarrow\forall x\in I,\;e^{\frac{(\lambda-1)\ln|x|}{\lambda}}f'(x)+\frac{\lambda-1}{\lambda x}e^{\frac{(\lambda-1)\ln|x|}{\lambda}}f(x) = 0\\
  &\Rightarrow\forall x\in I,\;\left(|x|^{\frac{\lambda-1}{\lambda}}f\right)'(x)=0\\
  &\Rightarrow\exists K\in\Rr/\;\forall x\in I,\;|x|^{\frac{\lambda-1}{\lambda}}f(x)=K\Rightarrow\exists K\in\Rr/\;\forall x\in I,\;f(x)=K|x|^{\frac{1-\lambda}{\lambda}}.
\end{align*}
				   

\textbf{1er cas.} Si $\lambda\in]-\infty,0[\cup]1,+\infty[$ alors $\frac{1-\lambda}{\lambda}< 0$ et donc $\lim_{x \rightarrow 0}|x|^{\frac{1-\lambda}{\lambda}}=+\infty$. La fonction $x\mapsto K|x|^{\frac{1-\lambda}{\lambda}}$ ne peut donc être la restriction à $I$ d'une fonction continue sur $\Rr$ que dans le cas $K = 0$. Ceci fournit $f_{/]-\infty,0[}= 0$, $f_{/]0,+\infty[}= 0$ et $f(0)=0$ par continuité en $0$. Dons $f$ est nécessairement nulle  et $\lambda$ n'est pas valeur propre de $\varphi$ dans ce cas.

\textbf{2ème cas.} Si $\lambda= 1$, les restriction de $f$ à $]-\infty,0[$ ou $]0,+\infty[$ sont constantes et donc, par continuité de $f$ en $0$, $f$ est constante sur $\Rr$. Réciproquement, les fonctions constantes $f$ vérifient bien $\varphi(f) = f$. Ainsi, $1$ est valeur propre de $f$ et le sous-espace propre associé est constitué des fonctions constantes.

\textbf{3ème cas.} Si $\lambda\in]0,1[$, nécessairement $\exists(K_1,K_2)\in\Rr^2/\;\forall x\in\Rr,\;f(x)=\left\{
\begin{array}{l}
K_1x^{\frac{1}{\lambda}-1}\;\text{si}\;x\geq0\\
K_2(-x)^{\frac{1}{\lambda}-1}\;\text{si}\;x < 0
\end{array}
\right.$. $f$ ainsi définie est bien continue sur $\Rr$. Calculons alors $\varphi(f)$.

 
$(\varphi(f))(0)=f(0)=0$ puis si $x > 0$,

\begin{center}
$(\varphi(f))(x)=\frac{1}{x}\int_{0}^{x}K_1t^{\frac{1-\lambda}{\lambda}}\;dt=\frac{\lambda K_1}{x}x^{\frac{1}{\lambda}}=\lambda K_1 x^{\frac{1}{\lambda}-1}=\lambda f(x)$
\end{center}

et de même si $x < 0$. Enfin, $(\varphi(f))(0)=0=\lambda f(0)$. Finalement $\varphi(f) =\lambda f$. $\lambda$ est donc valeur propre de $\varphi$ ($K_1 = K_2 = 1$ fournit une fonction non nulle) et le sous-espace propre associé à $\lambda$ est de dimension $2$. Une base de ce sous-espace est $(f_1,f_2)$ où $\forall x\in\Rr$, $f_1(x)=\left\{
\begin{array}{l}
x^{\frac{1}{\lambda}-1}\;\text{si}\;x\geq0\\
0\;\text{si}\;x < 0
\end{array}
\right.$ et $f_2(x) =\left\{
\begin{array}{l}
0\;\text{si}\;x\geq0\\
(-x)^{\frac{1}{\lambda}-1}\;\text{si}\;x < 0
\end{array}
\right.$.

Finalement

\begin{center}
\shadowbox{
$\text{Sp}(\varphi)= ]0,1]$.
}
\end{center}
\end{enumerate}
\fincorrection
\nocorrection
\nocorrection
\nocorrection
\nocorrection
\nocorrection
\nocorrection
\nocorrection
\nocorrection
\nocorrection
\nocorrection
\nocorrection
\nocorrection
\nocorrection
\nocorrection
\nocorrection
\nocorrection
\nocorrection
\nocorrection
\nocorrection
\nocorrection
\nocorrection
\nocorrection
\correction{001635}

  $\chi_{A}=(-1-X)(2-X)^{2}$. Donc $A$ est diagonalisable ssi $\dim
  \ker(A-2I)=2$. Or $\mathrm{rg}(A-2I)=2$, donc $\dim \ker(A-2I)=1$ donc $A$
  n'est pas diagonalisable. Cependant, $\chi_{A}$ est scindé sur $\R$
  donc $A$ est triangularisable sur $\R$.

  $\Big(\begin{smallmatrix}x\\y\\z\end{smallmatrix}\Big)\in\ker(A+I)%
  \Leftrightarrow\Big\{\begin{smallmatrix}%
    4x+2y+4z=0\\-x+4y-z=0\\-2x-y-2z=0%
  \end{smallmatrix}
  \Leftrightarrow\big\{\begin{smallmatrix}%
    x+z=0\\y=0%
  \end{smallmatrix}$ donc $\ker(A+I)=\R\Big(\begin{smallmatrix}%
    1\\0\\-1%
  \end{smallmatrix}\Big)$

De m\^eme,
 
  $\Big(\begin{smallmatrix}x\\y\\z\end{smallmatrix}\Big)\in\ker(A-2I)%
  \Leftrightarrow\Big\{\begin{smallmatrix}%
    x+2y+4z=0\\-x+y-z=0\\-2x-y-5z=0%
  \end{smallmatrix}
  \Leftrightarrow\big\{\begin{smallmatrix}%
    x=2y\\z=-y%
  \end{smallmatrix}$ donc $\ker(A+I)=\R\Big(\begin{smallmatrix}%
    2\\1\\-1%
  \end{smallmatrix}\Big)$

On sait que $\ker((A-2I)^{2})$ est de dimension 2, et que $\Big(\begin{smallmatrix}%
    2\\1\\-1%
  \end{smallmatrix}\Big)\in\ker(A-2I)\subset\ker((A-2I)^{2})$.
 On cherche donc un deuxième vecteur dans $\ker((A-2I)^{2})$, linéairement
 indépendant de $\Big(\begin{smallmatrix}%
    2\\1\\-1%
  \end{smallmatrix}\Big)$.

$(A-2I)^{2}=\Big(\begin{smallmatrix}%
    -9&0&-18\\0&0&0\\9&0&18%
  \end{smallmatrix}\Big)$
donc $\Big(\begin{smallmatrix}%
    0\\1\\0%
  \end{smallmatrix}\Big)$
convient. De plus~: $A\Big(\begin{smallmatrix}%
    0\\1\\0%
  \end{smallmatrix}\Big)=\Big(\begin{smallmatrix}%
    2\\3\\-1%
  \end{smallmatrix}\Big)=\Big(\begin{smallmatrix}%
    2\\1\\-1%
  \end{smallmatrix}\Big)+2\Big(\begin{smallmatrix}%
    0\\1\\0%
  \end{smallmatrix}\Big)$.

Donc en posant $P=\Big(\begin{smallmatrix}%
    1&2&0\\0&1&1\\-1&-1&0%
  \end{smallmatrix}\Big)$, on obtient
$P^{-1}AP=\Big(\begin{smallmatrix}%
    -1&0&0\\0&2&1\\0&0&2%
  \end{smallmatrix}\Big)$.
\fincorrection
\correction{001636}
  On a $A^{3}=A$, donc $P=X^{3}-X=(X-1)(X+1)X$ est un polyn\^ome
  annulateur de $A$. Il s'agit d'un poyn\^ome scindé à racine simples
  donc $A$ est diagonalisable. Les valeurs propres de $A$ sont des
  racines de $P$ donc $\mathrm{Sp}(A)\subset\{0,1,-1\}$. On a $\mathrm{rg} A=2$
  donc $0$ est valeur propre de multiplicité 2. La résolution de système
  $(A+I)X=0$ montre que $\ker (A+I)=\R\Big(\begin{smallmatrix}%
    1\\1\\-1\\-1%
  \end{smallmatrix}\Big)$,
  donc $-1$ est valeur propre de multiplicité 1 donc $1$ est
  nécessairement valeur propre de multiplicité 1~: on en déduit que
  $\chi_{A}=X^{2}(X-1)(X+1)$.
\fincorrection
\nocorrection
\nocorrection
\nocorrection
\nocorrection
\correction{001641}
\begin{enumerate}
  \item 
$A$ est triangulaire inférieure donc ses valeurs sont ses coefficients
diagonaux~: 1, 2 et 3. $A$ a trois valeurs propres distinctes donc $A$
est diagonalisable.

\item
$\chi_{B}=-(X-1)(X+1)^{2}$. $B+I=
\begin{pmatrix}
  3&-2&1\\
  3&-2&1\\
  -1&2&1
\end{pmatrix}
$, donc $\mathrm{rg}(B+I)=2$, $\dim (\ker B+I)=1<2$ donc $B$ n'est pas
diagonalisable.

$\chi_{B}(B)=0$ donc $B(B^{2}+B-I)=I$, soit $B^{-1}=B^{2}+B-I=
\begin{pmatrix}
  -2& 2& 1\\
  -1& 1& 1\\
   3&-2& 0 
\end{pmatrix}
$.
\end{enumerate}
\fincorrection
\correction{001642}
  \begin{enumerate}
  \item 
    ${}^t{A}=A$ donc $A$ est diagonalisable dans une base
    orthonormée.

  \item
    Par exemple~: $P=
    \begin{pmatrix}
      1& 1& 1\\
      1&-1& 0\\
      1& 0&-1
    \end{pmatrix}
$, $P^{-1}AP=
\begin{pmatrix}
  -1& 0& 0\\
   0& 2& 0\\
   0& 0& 2
 \end{pmatrix}
$.

 \item
   $Q=
   \begin{pmatrix}
     1/\sqrt{3} & 1/\sqrt{2} & 1/\sqrt{6} \\
     1/\sqrt{3} &-1/\sqrt{2} & 1/\sqrt{6} \\
     1/\sqrt{3} &          0 &-2/\sqrt{6} 
   \end{pmatrix}
$, $Q^{-1}AQ=\begin{pmatrix}
  -1& 0& 0\\
   0& 2& 0\\
   0& 0& 2
 \end{pmatrix}
$ et ${}^t{Q}=Q^{-1}$
  \end{enumerate}
\fincorrection
\nocorrection
\nocorrection
\nocorrection
\nocorrection
\nocorrection
\nocorrection
\nocorrection
\nocorrection
\nocorrection
\nocorrection
\nocorrection
\nocorrection
\nocorrection
\nocorrection
\nocorrection
\nocorrection
\nocorrection
\nocorrection
\nocorrection
\nocorrection
\nocorrection
\nocorrection
\nocorrection
\nocorrection
\nocorrection
\nocorrection
\nocorrection
\nocorrection
\nocorrection
\correction{001672}
 $\mathrm{tr} a=\mathrm{tr} A=-1,$ $\det a=\det A=-6$
$$P_a(X)=X^2-\mathrm{tr} X+\det a=X^2+X-6=(X-2)(X+3).$$ Donc le spectre est
$\{2,-3\}$, il est de taille 2 comme l'espace est de dimension 2.
D'apr\`es le cours, $a$ est diagonalisable et les espaces propres
de dimension 1. L'espace propre associ\'e \`a la valeur propre 2
est l'ensemble des $(x,y)$ tels que $7x-10y=2x$ ou $x=2y.$ On peut
prendre $\vec{f}_1=(2,1)$ pour base de cet espace propre. L'espace
propre associ\'e \`a la valeur propre $-3$ est l'ensemble des
$(x,y)$ tels que $7x-10y=-3x$ ou $x=y.$ On peut prendre
$\vec{f}_2=(1,1)$ pour base de cet espace propre. Alors si
$f=(\vec{f}_1,\vec{f}_2)$ on a

$$P=[\mathrm{id}_E]_f^e=\left[\begin{array}{rr}2&1\\1&1\end{array}\right],\ \
P^{-1}=[\mathrm{id}_E]_e^f=\left[\begin{array}{rr}
1&-1\\-1&2\end{array}\right],\ \
D=[a]_f^f=\left[\begin{array}{rr}2&0\\0&-3\end{array}\right].$$

$$D^{50}=[a^{50}]_f^f=\left[\begin{array}{rr}2^{50}&0\\0&(-3)^{50}\end{array}\right],
\ A^{50}=[a^{50}]_e^e=PD^{50}P^{-1}=\left[\begin{array}{rr}
2.2^{50}-(-3)^{50}&-2.2^{50}+2.(-3)^{50}\\2^{50}-(-3)^{50}&-2^{50}+2.(-3)^{50}
\end{array}\right]$$
Donc $\lim_{n\infty}\frac{1}{3^{2n}}[a^{2n}]_f^f=L=
\left[\begin{array}{rr}0&0\\0&1\end{array}\right],$ et
$\lim_{n\infty}\frac{1}{3^{2n}}[a^{2n}]_e^e=PLP^{-1}=\left[\begin{array}{rr}
-1&2\\-1&2\end{array}\right].$
\fincorrection
\correction{001673}
 Si
$X=(x_{ij})_{1\leq i,j\leq n}\in F,$ il est clair que
$X=\sum_{1\leq i,j\leq n}x_{ij}F_{ij}.$ C'est donc une famille
g\'en\'eratrice. Elle est ind\'ependante, car si $\sum_{1\leq
i,j\leq n}x_{ij}F_{ij}$ est la matrice nulle, cela implique que
$x_{ij} =0$ pour tous $i$ et $j.$ C'est donc une base de $F$. Elle
est de taille $n^2$, donc $F$ est de dimension $n^2.$ Ensuite, si
$D=\mathrm{diag}(d_1,\ldots,d_n)$ et si $X=(x_{ij})_{1\leq i,j\leq
n}$ alors le coefficient $(i,j)$ de la matrice $\Phi(X)=\alpha
XD+\beta DX$ est $(\alpha d_j+\beta d_i)x_{ij}.$ Donc $\Phi
(F_{ij})=(\alpha d_j+\beta d_i)F_{ij},$ ce qui est dire que
$F_{ij}$ est un vecteur propre de $\Phi$ pour la valeur propre
$\alpha d_j+\beta d_i.$ L'espace $F$ admet donc une base de
vecteurs propres de $\Phi.$ D'apr\`es le cours, cela entra\^ine
que $\Phi$ est diagonalisable. Si on le repr\'esente dans la base
de vecteurs propres, le d\'eterminant de $\Phi$ est donc le
produit des \'el\'ements diagonaux, c'est \`a dire $\det
\Phi=\prod_{i=1}^n\prod_{j=1}^n(\alpha d_j+\beta d_i).$ Plus
g\'en\'eralement
 $\det (\Phi-\lambda \mathrm{id}_F)=
\prod_{i=1}^n\prod_{j=1}^n(\alpha d_j+\beta d_i-\lambda).$
\fincorrection
\correction{001674}
 Notons $D_n=\det B.$
Alors $D_1=2\cos \theta=\frac{\sin 2\theta}{\sin \theta}$ et
$D_2=4\cos^2\theta -1=\frac{\sin 3\theta}{\sin \theta}.$ Si $n>2$,
d\'eveloppons $D_n$ par rapport \`a la derni\`ere ligne, en
recommencant encore une fois avec un des d\'eterminants d'ordre
$n-1$ obtenus. On obtient $D_n=2\cos \theta D_{n-1}-D_{n-2}.$
Faisons l'hypoth\`ese de r\'ecurrence que $D_k=\frac{\sin
(k+1)\theta}{\sin \theta}$ pour $k<n.$ On a vu que c'est vrai pour
$k=1$ et 2. Alors par des identit\'es trigonom\'etriques
classiques $D_n=\frac{2\cos\theta \sin n\theta}{\sin \theta}
-\frac{\sin (n-1)\theta}{\sin \theta}= \frac{\sin
(n+1)\theta}{\sin \theta},$ et la r\'ecurrence est \'etendue.
Puisque $\sin x=0 \Leftrightarrow $ il existe un entier relatif
$k$ tel que $x=k\pi$ alors $D_n=0$ si et seulement si il existe
$k=1,2,\ldots,n$ tel que $\theta=\frac{k\pi}{n+1}$ les autres
valeurs de $k$ \'etant exclues car $0<\theta<\pi.$ Par
d\'efinition de $P_A$ on a $P_A(-2\cos \theta)=D_n=\frac{\sin
(n+1)\theta}{\sin \theta}$ qui s'annule pour les $n$ nombres
distincts $-2\cos \frac{k\pi}{n+1},$ $k=1,2,\ldots,n$ qui sont
n\'ecessairement toutes les valeurs propres de $A.$ Les valeurs
propres de $2I_n+ A$ sont donc $2-2\cos
\frac{k\pi}{n+1}=4\sin^2\frac{k\pi}{2n+2}>0.$ Le spectre de $2I_n-
A$ est le m\^eme.

\fincorrection
\nocorrection
\nocorrection
\correction{002563}
Soit $M$ la matrice r\'eelle $3\times3$ suivante :
$$M=\begin{pmatrix}0&2&-1 \\ 3&-2&0 \\ -2&2&1\end{pmatrix}$$
\begin{enumerate}  \item {\it D\'eterminons les valeurs propres de $M$.}

Ce sont les racines du polyn\^ome caract\'eristique
\begin{align}P_M(X)=\begin{vmatrix}-X&2&-1 \\ 3&-2-X&0 \\ -2&2&1-X\end{vmatrix}
&=-1\begin{vmatrix}3&-2-X \\ -2&2\end{vmatrix}
+(1-X)\begin{vmatrix}-X&2 \\ 3&-2-2X\end{vmatrix}\\
&=(1-X)(X^2+2X-8)\\
&=(1-X)(X+4)(X-2).
\end{align}
La matrice $M$ admet donc trois valeurs propres distinctes qui sont : $1,2,$ et $-4$.

  \item {\it Montrons que $M$ est diagonalisable.}

Nous venons de voir que $M$, matrice r\'eelle $3\times 3$, admet trois valeurs propres r\'eelles distinctes, 
cela prouve que $M$ est diagonalisable.

  \item {\it D\'eterminons une base de vecteurs propres et $P$ la matrice de passage.}

Les trois sous-espaces propres distincts sont de dimension $1$, il suffit de d\'eterminer un vecteur propre pour
chacune des valeurs propres.

$\lambda=1$ : Le vecteur $\vec u$ de coordonn\'ees $(x,y,z)$ est un vecteur propre pour la valeur propre $1$
 si et seulement si 
$$\left\{\begin{align*}2y-z&=x \\  3x-2y&=y \\  -2x+2y+z&=z\end{align*}\right.\iff
\left\{\begin{align*}-x+2y-z&=0 \\  3x-3y&=0 \\  -2x+2y&=0\end{align*}\right.\iff
\left\{\begin{align*}x&=y \\  x&=z\end{align*}\right.$$
Le sous-espace propre associ\'e \`a la valeur propre $\lambda=1$ est la droite vectorielle engendr\'ee par 
le vecteur $\vec{e_1}$ de coordonn\'ees $(1,1,1)$.


$\lambda=2$ : Le vecteur $\vec u$ de coordonn\'ees $(x,y,z)$ est un vecteur propre pour la valeur propre $2$
 si et seulement si 
$$\left\{\begin{align*}-2x+2y-z&=0 \\  3x-4y&=0 \\  -2x+2y-z&=0\end{align*}\right.\iff
\left\{\begin{align*}3x-4y&=0 \\  -2x+2y-z&=0\end{align*}\right.$$
Le sous-espace propre associ\'e \`a la valeur propre $\lambda=2$ est la droite vectorielle engendr\'ee par 
le vecteur $\vec{e_2}$ de coordonn\'ees $(4,3,-2)$.


$\lambda=-4$ : Le vecteur $\vec u$ de coordonn\'ees $(x,y,z)$ est un vecteur propre pour la valeur propre $-4$
 si et seulement si 
$$\left\{\begin{align*}-4x+2y-z&=0 \\  3x+2y&=0 \\  -2x+2y+5z&=0\end{align*}\right.\iff
\left\{\begin{align*}x-z&=0 \\  2y+3x&=0\end{align*}\right.$$
Le sous-espace propre associ\'e \`a la valeur propre $\lambda=-4$ est la droite vectorielle engendr\'ee par 
le vecteur $\vec{e_3}$ de coordonn\'ees $(2,-3,2)$.

Les vecteurs $\vec{e_1}, \vec{e_2}$ et $\vec{e_3}$ forment une base de $E$ compos\'ee de vecteurs propres, la matrice 
de passage $P$ est \'egale \`a
$$P=\begin{pmatrix}1&4&2 \\ 1&3&-3 \\ 1&-2&2\end{pmatrix}$$

  \item {\it Exprimons $M^k$ en fonction de $D^k$, puis calculons $M^k$.}

On a $$D=P^{-1}MP=\begin{pmatrix}1&0&0 \\ 0&2&0 \\ 0&0&-4\end{pmatrix}$$
pour $k\in\N$, on a 
$$D^k=\begin{pmatrix}1&0&0 \\ 0&2^k&0 \\ 0&0&(-4)^k\end{pmatrix},$$
et $M^k=PD^kP^{-1}.$

Calculons donc la matrice $P^{-1}$ : on a $\displaystyle P^{-1}={\frac{1}{\det P}}(\rm{com} P)^t$. Or
$$\det P=\begin{vmatrix}1&4&2 \\ 1&3&-3 \\ 1&-2&2\end{vmatrix}=\begin{vmatrix}1&6&2 \\ 1&0&-3 \\ 1&0&2\end{vmatrix}=
-6\begin{vmatrix}1&-3 \\ 1&2\end{vmatrix}=-30,$$ et
$${\rm com}P=\begin{pmatrix}0&-5&-5 \\ -12&0&6 \\ -18&5&-1\end{pmatrix}$$ d'o\`u
$$P^{-1}=-{\frac{1}{30}}\begin{pmatrix}0&-12&-18 \\ -5&0&5 \\ -5&6&-1\end{pmatrix}.$$
On a donc 
$$M^k=PD^kP^{-1}=-{\frac{1}{30}}
\begin{pmatrix}-5.2^{k+2}-10(-4)^k&-12+12(-4)^k&-18+5.2^{k+2}-2(-4)^k \\ 
-15.2^{k}-15(-4)^k&-12-18(-4)^k&-18+5.2^{k+1}+3(-4)^k \\ 
5.2^{k+1}-10(-4)^k&-12+12(-4)^k&-18-5.2^{k+1}-2(-4)^k\end{pmatrix}$$
 \end{enumerate}
\fincorrection
\correction{002566}
 Soit $$A=\begin{pmatrix}1&0&0 \\ 0&1&0 \\ 1&-1&2\end{pmatrix}$$
{\it D\'emontrons que $A$ est diagonalisable et trouvons une matrice $P$ telle que $P^{-1}AP$ soit diagonale.}

Commen\c cons par calculer le polyn\^ome caract\'eristique de $A$ :
$$P_A(X)=\begin{vmatrix}1-X&0&0 \\  0&1-X&0 \\ 1&-1&2-X\end{vmatrix}=(1-X)^2(2-X)$$
Les racines du polyn\^ome caract\'eristique sont les r\'eels $1$ avec la multiplicit\'e $2$, et $2$ avec la multiplicit\'e $1$.

D\'eterminons les sous-espaces propres associ\'es : Soit $E_1$ le sous-espace propre associ\'e \`a la valeur propre double $1$.

$E_1=\{V(x,y,z)\in\R^3/\ A.V=V\}$, 
$$V\in E_1\iff \left\{ \begin{align*}x&=x \\  y&=y \\  x-y+z&=0\end{align*}\right.\iff x-y+z=0$$
$E_1$ est donc un plan vectoriel, dont les vecteurs $e_1=(1,1,0)$ et $e_2=(0,1,1)$ forment une base.

Soit $E_2$ le sous-espace propre associ\'e \`a la valeur propre simple $2$.

$E_2=\{V(x,y,z)\in\R^3/\ A.V=2V\}$, 
$$V\in E_2\iff \left\{ \begin{align*}x&=2x \\  y&=2y \\  x-y+2z&=2z\end{align*}\right.\iff x=0, y=0$$
$E_2$ est donc une droite vectorielle, dont le vecteur $e_3=(0,0,1)$ est une base.

Les dimensions des sous-espaces propres sont \'egales \`a la multiplicit\'e des valeurs propres correspondantes, la matrice $A$ est donc diagonalisable.
Dans la base $(e_1, e_2, e_3)$ l'endomorphisme repr\'esent\'e par $A$ (dans la base canonique) a pour matrice.
$$D=\begin{pmatrix}1&0&0 \\  0&1&0 \\  0&0&2\end{pmatrix}$$
la matrice de passage $$P=\begin{pmatrix}1&0&0 \\  1&1&0 \\  0&1&1\end{pmatrix}$$ v\'erifie $P^{-1}AP=D$.
\fincorrection
\correction{002567}
Soit $$A=\begin{pmatrix}1&1&-1 \\ 0&1&0 \\ 1&0&1\end{pmatrix}$$
{\it Factorisons le polyn\^ome caract\'eristique de $A$}. 
$$P_A(X)=\begin{vmatrix}1-X&1&-1 \\  0&1-X&0 \\ 1&0&1-X\end{vmatrix}=(1-X)^3+(1-X)=(1-X)((1-X)^2+1)=(1-X)(X^2-2X+2)$$
factorisons maintenant le polyn\^ome $X^2-2X+2$, le discriminant r\'eduit $\Delta'=1-2=-1$, ce polyn\^ome n'admet donc pas de racines r\'eelles, mais 
deux racines complexes conjugu\'ees qui sont : $1+i$ et $1-i$. On a $P_A(X)=(1-X)(1-i-X)(1+i-X)$.

La matrice $A$ n'est pas diagonalisable dans $\R$ car son polyn\^ome caract\'eristique n'a pas toutes ses racines dans $\R$, elle est diagonalisable dans
$\C$ car c'est une matrice $3\times 3$ qui admet trois valeurs propres distinctes.

\fincorrection
\correction{002568}
 Soit  $$A=\begin{pmatrix}a&c \\  c&d\end{pmatrix}\in M_2(\R)$$

{\it D\'emontrons que $A$ est diagonalisable dans $\R$}.

Le polyn\^ome caract\'eristique $P_A(X)$ est \'egal \`a
$$P_A(X)=\begin{vmatrix}a-X&c \\  c& d-x\end{vmatrix}=(a-X)(d-X)-c^2=X^2-(a+d)X+ad-c^2,$$
d\'eterminons ses racines : calculons le discriminant :
\begin{align*}\Delta&=(a+d)^2-4(ad-c^2) \\  &=a^2+d^2+2ad-4ad+4c^2 \\ &=a^2+d^2-2ad+4c^2 \\ &=(a-d)^2+4c^2\geq 0\end{align*}
On a $\Delta=0\iff a-d=0\ {\hbox{et}}\ c=0$, mais, si $c=0$, la matrice $A$ est d\'ej\`a diagonale. Sinon $\Delta>0$ et le polyn\^ome caract\'eristique admet
deux racines r\'eelles distinctes, ce qui prouve que la matrice est toujours diagonalisable dans $\R$.

\fincorrection
\nocorrection
\nocorrection
\nocorrection
\correction{002575}
{\it On suppose qu'une population $x$ de lapins et une population $y$ de loups sont gouvern\'ees par le syst\`eme suivant d'\'equations diff\'erentielles} :
$$(S)\ \ \left\{\begin{align*}x'&= 4x-2y \\  y'&= x+y\end{align*}\right.$$
\begin {enumerate}
  \item On diagonalise la matrice 
$$A=\begin{pmatrix}4&-2 \\ 1&1 \\ \end{pmatrix}.$$
Pour cela on d\'etermine ses valeurs propres :
$$\det(A-\lambda I)=\begin{vmatrix}4-\lambda &-2 \\ 1&1-\lambda \\ \end{vmatrix}=(4-\lambda)(1-\lambda)+2=\lambda^2-5\lambda+6=(\lambda-2)(\lambda-3).$$
Ainsi, la matrice $A$ admet deux valeurs propres distinctes, qui sont $\lambda_1=2$ et $\lambda_2=3$. Elle est diagonalisable. D\'eterminons une base de vecteurs propres :
$$\begin{pmatrix}4&-2 \\ 1&1 \\ \end{pmatrix}\begin{pmatrix}x \\  y\end{pmatrix}=
\begin{pmatrix}2x \\  2y\end{pmatrix}\iff x=y,$$
d'o\`u le vecteur propre $u_1=(1,1)$ associ\'e \`a la valeur propre $\lambda_1=2$ . 
$$\begin{pmatrix}4&-2 \\ 1&1 \\ \end{pmatrix}\begin{pmatrix}x \\  y\end{pmatrix}=
\begin{pmatrix}3x \\  3y\end{pmatrix}\iff x=2y,$$
d'o\`u le vecteur propre $u_2=(2,1)$ associ\'e \`a la valeur propre $\lambda_2=3$ .
Dans la base $(u_1, u_2)$, la matrice s'\'ecrit
$$A'=\begin{pmatrix}2&0 \\ 0&3 \\ \end{pmatrix}.$$
On a $A=PA'P^{-1}$ o\`u
$$P=\begin{pmatrix}1&2 \\ 1&1 \\  \end{pmatrix}\ {\hbox{et}}\ P^{-1}=\begin{pmatrix}-1&2 \\ 1&-1 \\ \end{pmatrix}.$$
  \item Exprimons le syst\`eme $(S)$ et ses solutions dans une base de vecteurs propres de $A$.

Dans la base $(u_1, u_2)$, le syst\`eme $(S)$ devient
$$(S')\ \ \left\{\begin{align*}X'&= 2X \\  Y'&= 3Y\end{align*}\right.$$
Ses solutions sont les fonctions
$$\left\{\begin{align*}X(t)&= X(0)e^{2t} \\  Y(t)&=Y(0)e^{3t}\end{align*}\right.$$

  \item  Pour repr\'esenter graphiquement les trajectoires de $(S)$ dans le rep\`ere $(Oxy)$, on trace d'abord le rep\`ere $(O, u_1,u_2)$ dans le rep\`ere $(0xy)$, puis, on trace les courbes $$Y={\frac{Y(0)}{X(0)}}X^{3/2}$$
dans le rep\`ere $(O, u_1,u_2)$ (ou $OXY$).

  \item On voit sur le dessin que si $Y(0)$ est strictement positif, alors la population des lapins, $x(t)$ tend vers $+\infty$ quand $t$ tend vers $+\infty$. Si $Y(0)$ est strictement n\'egatif alors la populations des lapins s'\'eteint dans la mesure ou $x(t)$ dans ce cas tendrait vers $-\infty$.
\end {enumerate}

\fincorrection
\correction{002576}
 {\it Soit $u$ l'endomorphisme de $\R^3$, dont la matrice dans la base canonique est}
$$A=\begin{pmatrix}3&2&-2 \\  -1&0&1 \\ 1&1&0 \\ \end{pmatrix}$$ 
\begin {enumerate}
  \item Calculons les valeurs propres de $A$. 
$$\det(A-\lambda I)=\begin{vmatrix}3-\lambda&2&-2 \\  -1&-\lambda&1 \\ 1&1&-\lambda \\ \end{vmatrix}=(\lambda-1)^3.$$ 
La matrice $A$ admet une valeur propre triple qui est $\lambda=1$, elle ne peut pas \^etre diagonalisable sinon son sous-espace propre serait de dimension $3$ or, $A\neq I$.
  \item Calculons $(A-I)^2$.
$$(A-I)^2=\begin{pmatrix}2&2&-2 \\  -1&-1&1 \\ 1&1&-1 \\ \end{pmatrix}^2=\begin{pmatrix}0&0&0 \\  0&0&0 \\ 0&0&0 \\ \end{pmatrix}$$

 Montrons que $A^n=nA+(1-n)I$ en utilisant la formule du bin\^ome de Newton.
$$A^n=(A-I+I)^n=\sum_{k=0}^{n}C_n^k(A-I)^kI^{n-k}=C_n^0I^n+C_n^1(A-I)=I+n(A-I)=nA+(1-n)I.$$
Car, pour $k\geq2$, on a $(A-I)^k=0$.
  \item  {\it Soient $P(X)=(X-1)^2$ et $Q\in\R[X]$}. 

Exprimons le reste de la division euclidienne de $Q$ par $P$ en fonction de $Q(1)$ et $Q'(1)$, o\`u $Q'$ est le polyn\^ome d\'eriv\'e de $Q$. 

Il existe des polyn\^omes $S$ et $R$, avec $d^{\circ}R<d^{\circ}P$ ou $R=0$, tels que 
$$Q(X)=S(X)(X-1)^2+R(X).$$
Notons $R(X)=aX+b$ ($R(X)$ est de degr\'e $1$ car $P$ est de degr\'e 2) et d\'erivons, on obtient
$$Q'(X)=S'(X)(X-1)^2+2(X-1)S(X)+a,$$
on a donc $Q(1)=R(1)=a+b$ et $Q'(1)=a$, c'est-\`a-dire $a=Q'(1)$ et $b=Q(1)-Q'(1)$ d'o\`u
$$R(X)=Q'(1)X+(Q(1)-Q'(1)).$$
D'apr\`es la question $2)$, on remarque que $P(A)=0$, en choisissant le polyn\^ome $Q(X)=X^n$ on a
$Q(1)=1$ et $Q'(1)=n$, donc 
$$Q(A)=A^n=R(A)=Q'(1)A+(Q(1)-Q'(1))I=nA+(1-n)I.$$ 
  \item
  \begin {enumerate}
    \item  Montrons que l'image de $\R^3$ par l'endomorphisme $(A-I)$ est un sous-espace vectoriel de dimension $1$.
$$\forall(X,Y,Z)\in Im (A-I), \exists (x,y,z)\in\R^3,\ (A-I)\begin{pmatrix}x \\  y \\  z\end{pmatrix}
=\begin{pmatrix}X \\  Y \\  Z\end{pmatrix},$$
c'est-\`a-dire
$$\begin{pmatrix}X \\  Y \\  Z\end{pmatrix}=(x+y-z)\begin{pmatrix}2 \\  -1 \\  1\end{pmatrix}.$$
Ce qui prouve que $Im(A-I)$ est la droite vectorielle engendr\'ee par le vecteur 
$\varepsilon_2=(2,-1,1)$.

    \item D\'eterminons un vecteur $\varepsilon_3$ tel que $u(\varepsilon_3)=\varepsilon_2+\varepsilon_3$.
On pose $\varepsilon_3=(x,y,z)$,
$$u(\varepsilon_3)=\varepsilon_2+\varepsilon_3\iff A\begin{pmatrix}x \\  y \\  z\end{pmatrix}
=\begin{pmatrix}2 \\  -1 \\  1\end{pmatrix}+\begin{pmatrix}x \\  y \\  z\end{pmatrix},$$
c'est-\`a-dire
$$\left\{\begin{align*}3x+2y-2z &= x+2 \\  -x+z &=y-1 \\  x+y &=z+1\end{align*}\right.\iff
\left\{\begin{align*}2(x+y-z) &= 2 \\ -1(x+y-z) &=-1 \\  (x+y-z) &=+1\end{align*}\right.\iff x+y-z=1.$$
On prends, par exemple $\varepsilon_3=(1,0,0)$.

D\'eterminons un vecteur propre $\varepsilon_1$ de $u$ non colin\'eaire \`a $\varepsilon_2$.
$$A\begin{pmatrix}x \\  y \\  z\end{pmatrix}
=\begin{pmatrix}x \\  y \\  z\end{pmatrix}\iff\left\{\begin{align*}3x+2y-2z &= x \\  -x+z &=y \\  x+y &=z\end{align*}\right.\iff x+y-z=0.$$
On peut prendre le vecteur $\varepsilon_1=(0,1,1)$ qui n'est pas colin\'eaire \`a $\varepsilon_2$.

    \item   Ecrivons la matrice de $u$ dans la base $(\varepsilon_1,\varepsilon_2,\varepsilon_3)$, ainsi que les matrices de passage.

On a $u(\varepsilon_1)=\varepsilon_1,u(\varepsilon_2)=\varepsilon_2$ et $u(\varepsilon_3)=\varepsilon_2+\varepsilon_3$ d'o\`u la matrice de $u$ dans la base $(\varepsilon_1,\varepsilon_2,\varepsilon_3)$

$$A'=\begin{pmatrix}1&0&0 \\  0&1&1 \\ 0&0&1 \\ \end{pmatrix}.$$
La matrice de passage $P$ est \'egale \`a 
$$P=\begin{pmatrix}0&2&1 \\  1&-1&0 \\ 1&1&0 \\ \end{pmatrix}$$ et son inverse
$$P^{-1}=\begin{pmatrix}0&1/2&1/2 \\  0&-1/2&1/2 \\ 1&1&-1 \\ \end{pmatrix}.$$
    \item Pour retrouver $A^n$, on \'ecrit $A'=I+N$, o\`u
$$N=\begin{pmatrix}0&0&0 \\  0&0&1 \\ 0&0&0 \\ \end{pmatrix},$$
et $N^2=0$. Par ailleurs, on a $A=PA'P^{-1}$, d'o\`u 
$$A^n=PA'^nP^{-1}=P(I+N)^nP^{-1}=P(I+nN)P^{-1}=P\begin{pmatrix}1&0&0 \\  0&1&n \\ 0&0&1 \\ \end{pmatrix}P^{-1}$$
$$=\begin{pmatrix}2n+1&2n&-2n \\  -n&1-n&n \\  n&n&1-n \\ \end{pmatrix}=n\begin{pmatrix}3&2&-2 \\  -1&0&1 \\ 1&1&0 \\ \end{pmatrix}+(1-n)\begin{pmatrix}1&0&0 \\  0&1&0 \\ 0&0&1 \\ \end{pmatrix}=nA+(1-n)I.$$    
  \end {enumerate}
\end {enumerate}
\fincorrection
\correction{002577}
{\it Soient $M$ et $A$ deux matrices de ${\cal{M}}_n(\R)$ telles que 
$MA=AM$.
On suppose que $M$ admet $n$ valeurs propres distinctes.}
  \begin {enumerate}
  \item {\it Soit $x$ un vecteur propre de $M$ de valeur propre $\lambda$.} 

Montrons que $MAx=\lambda Ax$.

On a $Mx=\lambda x$, donc $AMx=A\lambda x=\lambda Ax$. Mais, $AM=MA$, donc
$MAx=AMx=\lambda Ax$. Ce qui prouve que le vecteur $Ax$ est un vecteur propre de $M$ pour la valeur propre $\lambda$, et comme les valeurs propres de $M$ sont suppos\'ees distinctes, les sous-espaces propres sont de dimension $1$, donc $Ax$ est colin\'eaire \`a $x$. Ainsi, il existe un r\'eel $\mu$ tel que $Ax=\mu x$, donc $x$ est un vecteur propre de $A$.

  \item On note maintenant $\lambda_1,\cdots,\lambda_n$ les valeurs propres de $M$ et 
$\mu_1,\cdots,\mu_n$ celles de $A$.
  \begin {enumerate}
    \item  Montrons l'\'egalit\'e suivante :
$$\begin{vmatrix} &\lambda_1 &\cdots &\lambda_1^{n-1} \\ 
\vdots& & &\vdots \\  1 &\lambda_n &\cdots & \lambda_n^{n-1}\end{vmatrix}=
\prod_{1\leq i<j\leq n}(\lambda_j-\lambda_i).$$
Il s'agit du d\'eterminant de Vandermonde. Notons le $V(\lambda_1,\cdots,\lambda_n)$. La d\'emonstration se fait par r\'ecurrence sur $n$.
Pour $n=2$, c'est \'evident. Supposons le r\'esultat vrai pour $n-1$. Dans $V(\lambda_1,\cdots,\lambda_n)$, retranchons \`a chaque colonne $\lambda_1$ fois la pr\'ec\'edente ( en commen\c cant par la derni\`ere colonne). On obtient 
$$\begin{vmatrix}
&0&0&\cdots&0 \\ 
1& \lambda_2-\lambda_1&\lambda_2^2-\lambda_1\lambda_2&\cdots&\lambda_2^{n-1}-\lambda_1\lambda_2^{n-2} \\  \vdots&\vdots&\vdots& &\vdots \\  1&\lambda_n-\lambda_1&\lambda_n^2-\lambda_1\lambda_n&\cdots&\lambda_n^{n-1}-\lambda_1\lambda_n^{n-2}
\end{vmatrix}
=
\begin{vmatrix}
\lambda_2-\lambda_1&\lambda_2^2-\lambda_1\lambda_2&\cdots&\lambda_2^{n-1}-\lambda_1\lambda_2^{n-2} \\ 
\vdots&\vdots& &\vdots \\ \lambda_n-\lambda_1&\lambda_n^2-\lambda_1\lambda_n&\cdots&\lambda_n^{n-1}-\lambda_1\lambda_n^{n-2}
\end{vmatrix}.$$
On factorise alors chaque ligne par $(\lambda_i-\lambda_1)$ et on obtient
$$\begin{align*}
V(\lambda_1,\cdots,\lambda_n)&=(\lambda_2-\lambda_1)\cdots(\lambda_n-\lambda_1)
\begin{vmatrix}1 &\lambda_2 &\cdots &\lambda_2^{n-2} \\ 
\vdots& & &\vdots \\  1 &\lambda_n &\cdots & \lambda_n^{n-2}\end{vmatrix} \\ 
&= \prod_{i=2}^n(\lambda_i-\lambda_1)V(\lambda_2,\cdots,\lambda_n)=\prod_{1\leq i<j\leq n}(\lambda_i-\lambda_j)
\end{align*}$$
car $V(\lambda_2,\cdots,\lambda_n)=\prod_{2\leq i<j\leq n}(\lambda_j-\lambda_i)$ 
par hypoth\`ese de r\'ecurrence.


Ce d\'eterminant est le d\'eterminant du syst\`eme suivant, 
$$\left\{\begin{align*}\mu_1&=\alpha_0+\alpha_1\lambda_1+\cdots+\alpha_{n-1}\lambda_1^{n-1} \\ 
\vdots \\  \mu_n &=\alpha_0+\alpha_1\lambda_n+\cdots+\alpha_{n-1}\lambda_n^{n-1} \\ \end{align*} \right.$$
or $V(\lambda_1,\cdots,\lambda_n)\neq0$ puisque les $\lambda_i$ sont suppos\'es distincts, c'est donc un syst\`eme de Cramer, il admet donc une unique solution $(\alpha_0,\cdots,\alpha_{n-1})\in\R^n.$
    \item {\it Soient $M'$ et $A'$ les matrices diagonales suivantes} :
$$M'=\begin{pmatrix}
\lambda_1 & 0 &\cdots & 0 \\  
0 &\ddots  & & 0 \\ \vdots & &  & \vdots \\  
0 & \cdots & 0 &\lambda_n \\ 
\end{pmatrix},\ \ 
A'=\begin{pmatrix}
\mu_1 & 0 &\cdots & 0 \\  
0 &\ddots  & & 0 \\
\vdots & &  & \vdots \\  
0 & \cdots & 0 &\mu_n \\ 
\end{pmatrix}$$
Montrons qu'il existe des r\'eels $\alpha_0,\cdots,\alpha_{n-1}$ tels que
$$A'=\sum_{k=0}^{n-1}\alpha_k M'^k.$$

Compte tenu des matrices $A'$ et $M'$ l'existence de r\'eels tels que $$A'=\sum_{k=0}^{n-1}\alpha_k M'^k$$ est \'equivalente \`a l'existence d'une solution pour le syst\`eme pr\'ec\'edent, d'o\`u le r\'esultat.

On en d\'eduit qu'il existe des r\'eels $\alpha_0,\cdots,\alpha_{n-1}$ tels que
$$A=\sum_{k=0}^{n-1}\alpha_k M^k.$$ 
La matrice $M$ admet $n$ vecteurs propres lin\'eairement ind\'ependants qui sont \'egalement vecteurs propres de la matrice $A$. Par cons\'equent il existe une m\^eme matrice de passage $P$ telle que $M=PM'P^{-1}$ et $A=PA'P^{-1}$, d'o\`u l'\'egalit\'e 
$$A=\sum_{k=0}^{n-1}\alpha_k M^k.$$   
  \end {enumerate}
\end {enumerate}
\fincorrection
\correction{002581}
Soit $u$ l'endomorphisme de $\R^3$ dont la matrice dans la base canonique est  
$$A=\begin{pmatrix}-3&-2&-2 \\  2&1&2 \\  3&3&2\end{pmatrix}$$ 
\begin{enumerate}
 \item D\'eterminons le polyn\^ome caract\'eristique de $A$.

On a 
$$\begin{align*}P_A(X)&=\begin{vmatrix}-3-X&-2&-2 \\ 2&1-X&2 \\ 3&3&2-X\end{vmatrix}
=\begin{vmatrix}-3-X&0&-2 \\ 2&-1-X&2 \\ 3&1+X&2-X\end{vmatrix} \\ 
&=\begin{vmatrix}-3-X&0&-2 \\ 5&0&4-X \\ 3&1+X&2-X\end{vmatrix}
=-(1+X)\begin{vmatrix}-3-X&-2 \\ 5&4-X\end{vmatrix} \\ 
&=-(1+X)[(X-4)(X+3)+10]=-(1+X)(X^2-X-2)=-(1+X)^2(X-2)\end{align*}$$
 \item D\'emontrons que les valeurs propres de $A$ sont $-1$ et $2$ et d\'eterminons les sous-espaces propres associ\'es. 

Les valeurs propres de $A$ sont les racines du polyn\^ome caract\'eristique, ce sont donc bien les r\'eels $-1$ et $2$.

Les sous-espaces propres associ\'es sont les ensembles 
$$E_{-1}=\left\{(x,y,z)\in\R^3,\ A\begin{pmatrix}x \\  y \\  z\end{pmatrix}=-\begin{pmatrix}x \\  y \\  z\end{pmatrix}\right\}=\ker(A+I_3)$$
et
$$E_{2}=\left\{(x,y,z)\in\R^3,\ A\begin{pmatrix}x \\  y \\  z\end{pmatrix}=2\begin{pmatrix}x \\  y \\  z\end{pmatrix}\right\}=\ker(A-2I_3)$$
On a 
$$(x,y,z)\in E_{-1}\iff\left\{\begin{align*}-3x-2y-2z&=-x \\  2x+y+2z&=-y \\  3x+3y+2z&=-z\end{align*}\right.
\iff\left\{\begin{align*}-2x-2y-2z&=0 \\  2x+2y+2z&=0 \\  3x+3y+3z&=0\end{align*}\right.$$
Le sous-espace caract\'eristique $E_{-1}$ associ\'e \`a la valeur propre $-1$ est donc le plan vectoriel d'\'equation $x+y+z=0$, il est de dimension $2$, \'egale \`a la multiplicit\'e de la racine $-1$.

On a 
$$(x,y,z)\in E_{2}\iff\left\{\begin{align*}-3x-2y-2z&=2x \\  2x+y+2z&=2y \\  3x+3y+2z&=2z\end{align*}\right.
\iff\left\{\begin{align*}-5x-2y-2z&=0 \\  2x-y+2z&=0 \\  3x+3y&=0\end{align*}\right.$$
ce qui \'equivaut \`a $y=-x$ et $2z=-3x$.

Le sous-espace caract\'eristique $E_{2}$ associ\'e \`a la valeur propre $2$ est donc la droite vectorielle engendr\'ee par le vecteur $(2,-2,-3)$ , il est de dimension $1$, \'egale \`a la multiplicit\'e de la racine $2$.
 \item D\'emontrons que $A$ est diagonalisable et donnons une base de $\R^3$ dans laquelle la matrice de $u$ est diagonale.

La question pr\'ec\'edente et les r\'esultats obtenus sur les dimensions des sous-espaces propres permettent d'affirmer que la matrice $A$ est diagonalisable. Une base de $\R^3$ obtenue \`a partir de bases des sous-espaces propres est une base de vecteurs propres dans laquelle la matrice de $u$ est diagonale. Par exemple dans la base form\'ee des vecteurs  $u_1=(1,-1,0)$, $u_2=(1,0,-1)$ et 
$u_3=(2,-2,-3)$, la matrice de $u$ est la matrice $D$ qui s'\'ecrit
$$D=\begin{pmatrix}-1&0&0 \\  0&-1&0 \\  0&0&2\end{pmatrix}$$
 \item  Trouvons une matrice $P$ telle que $P^{-1}AP$ soit diagonale.

La matrice cherch\'ee $P$ est la matrice de passage exprimant la base de vecteurs propres $(u_1,u_2,u_3)$ dans la base canonique.
 C'est donc la matrice
$$P=\begin{pmatrix}1&1&2 \\  -1&0&-2 \\  0&-1&-3\end{pmatrix}.$$
On a $P^{-1}AP=D$.
\end{enumerate}

\fincorrection
\correction{002583}
 Soit $A=\begin{pmatrix}1&2&0&0 \\ 0&1&2&0 \\ 0&0&1&2 \\ 0&0&0&1\end{pmatrix}$.
Expliquons sans calcul pourquoi la matrice $A$ n'est pas diagonalisable.

On remarque que le polyn\^ome caract\'eristique de $A$ est \'egal \`a $(1-X)^4$. Ainsi la matrice $A$ admet-elle une unique valeur propre : $\lambda=1$, si elle \'etait diagonalisable, il existerait une matrice $P$ inversible telle que $A=PI_4P^{-1}$ alors $A=I_4$, or ce n'est pas le cas, par cons\'equent la matrice $A$ n'est pas diagonalisable.
\fincorrection
\correction{002584}
Soit $A$ une matrice $2\times 2$ \`a coefficients r\'eels. On suppose que dans chaque colonne de $A$ la somme des coefficients est \'egale \`a $1$. 
\begin{enumerate}
 \item Soient $(x_1,x_2)$, $(y_1,y_2)$ deux vecteurs de $\R^2$, on suppose que  
$$A\begin{pmatrix}x_1 \\  x_2\end{pmatrix}=\begin{pmatrix}y_1 \\  y_2\end{pmatrix}$$ 
montrons qu'alors 
$$y_1+y_2=x_1+x_2.$$

Compte tenu des hypoth\`eses, la matrice $A$ est de la forme
$$\begin{pmatrix}a&b \\ 1-a&1-b\end{pmatrix},$$
o\`u $a$ et $b$ sont des r\'eels. On a alors 
$$\begin{pmatrix}a&b \\ 1-a&1-b\end{pmatrix}\begin{pmatrix}x_1 \\  x_2\end{pmatrix}=\begin{pmatrix}y_1 \\  y_2\end{pmatrix}\iff
\left\{\begin{align*}ax_1+bx_2&=y_1 \\  (1-a)x_1+(1-b)x_2&=y_2\end{align*}\right.$$
ce qui implique $y_1+y_2=x_1+x_2$.
 \item Montrons que le vecteur $\varepsilon=(1,-1)$ est un vecteur propre de $A$.

Si $A\varepsilon=\begin{pmatrix}y_1 \\  y_2\end{pmatrix}$, alors $y_1+y_2=0$ donc $y_2=-y_1$ et $A\varepsilon=y_1\varepsilon$, ce qui prouve que $\varepsilon$ est un vecteur propre. On peut aussi le voir de la mani\`ere suivante 
$$A\varepsilon=\begin{pmatrix}a&b \\ 1-a&1-b\end{pmatrix}\begin{pmatrix}1 \\  -1\end{pmatrix}=
\begin{pmatrix}a-b \\  b-a\end{pmatrix}=(a-b)\varepsilon.$$
On note $\lambda=(a-b)$ sa valeur propre.
 \item Montrons que si $v$ est un vecteur propre de $A$ non colin\'eaire \`a $\varepsilon$, alors la valeur propre associ\'ee \`a $v$ est \'egale \`a $1$.

Soit $v=(x_1,x_2)$ un vecteur propre de $A$ non colin\'eaire \`a $\varepsilon$, notons $\mu$ sa valeur propre, on a $Av=\mu v$, et, d'apr\`es la question $1)$, on a $$x_1+x_2=\mu x_1+\mu x_2=\mu(x_1+x_2)$$
ce qui implique $\mu=1$ car $v$ est suppos\'e non colin\'eaire \`a $\varepsilon$ donc $x_1+x_2\neq0$. 

 \item  Soit $e_1=(1,0)$. Montrons que la matrice, dans la base $(e_1,\varepsilon)$, de l'endomorphisme associ\'e \`a $A$  est de la forme
$$\begin{pmatrix}1&0  \\ \alpha&\lambda \end{pmatrix},$$
o\`u $\alpha\in\R$.

Pour cela on \'ecrit $Ae_1$ et $A\varepsilon$ dans la base $(e_1,\varepsilon)$. 
On a d'une part $A\varepsilon=\lambda\varepsilon$ et, d'autre part,
$$\begin{pmatrix}a&b \\ 1-a&1-b\end{pmatrix}\begin{pmatrix}1 \\  0\end{pmatrix}=\begin{pmatrix}a \\  1-a\end{pmatrix}
=\begin{pmatrix}1 \\  0\end{pmatrix}+(a-1)\begin{pmatrix}1 \\  -1\end{pmatrix}.$$ 
D'o\`u la matrice dans la base $(e_1,\varepsilon)$
$$\begin{pmatrix}1&0 \\  \alpha&\lambda\end{pmatrix}$$
o\`u $\alpha=a-1$ et $\lambda=a-b$.

On en d\'eduit que si $\lambda\neq 1$, alors $A$ est diagonalisable sur $\R$.
Le polyn\^ome caract\'eristique de $A$ est \'egal \`a $(1-X)(\lambda-X)$, ainsi, si $\lambda\neq 1$, il admet deux racines distinctes ce qui prouve que $A$ est diagonalisable.
\end{enumerate}
\fincorrection
\correction{002586}

\centerline{\bf I}

Soit $\alpha\in\R$ et $A_\alpha\in M_3(\R)$ la matrice suivante
$$A_\alpha=\begin{pmatrix}-1 & 0 & \alpha+1 \\  1&-2&0 \\ -1&1&\alpha\end{pmatrix}$$
{\it Premi\`ere partie} :

\begin{enumerate} 
 \item{\it Factorisons le polyn\^ome caract\'eristique $P_{A_\alpha}(X)$ en produit de facteurs du premier degr\'e.}

On a
\begin{align*}
P_{A_\alpha}(X)
&=\begin{vmatrix}-1-X & 0 & \alpha+1 \\  1&-2-X&0 \\ -1&1&\alpha-X\end{vmatrix}
=\begin{vmatrix}-1-X & 0 & \alpha+1 \\  -1-X&-2-X&0 \\ 0&1&\alpha-X\end{vmatrix} \\ 
&=\begin{vmatrix}-1-X & 0 & \alpha+1 \\  0&-2-X&-\alpha-1 \\ -1&1&\alpha-X\end{vmatrix} \\ 
&=(-1-X)[(-2-X)(\alpha-X)+\alpha+1] \\ 
&=-(X+1)[X^2+(2-\alpha)X+1-\alpha].
\end{align*}
Factorisons le polyn\^ome $X^2+(2-\alpha)X+1-\alpha$, son discriminant est \'egal \`a
$$\Delta=(2-\alpha)^2-4(1-\alpha)=\alpha^2.$$
On a donc $\sqrt{\Delta}=|\alpha|$, ce qui nous donne les deux racines 
$$\lambda_1={\frac{\alpha-2-\alpha}{2}}=-1\ \ {\hbox{et}}\ \ 
\lambda_2={\frac{\alpha-2+\alpha}{2}}=\alpha-1.$$
Le polyn\^ome caract\'eristique $P_{A_\alpha}(X)$ se factorise donc en
$$P_{A_\alpha}(X)=-(X+1)^2(X-\alpha+1).$$ 

  \item {\it D\'eterminons selon la valeur du param\`etre $\alpha$ les valeurs propres distinctes de $A_\alpha$ et leur multiplicit\'e.}

Les valeurs propres de $A_\alpha$ sont les racines du polyn\^ome caract\'eristique $P_{A_\alpha}$, ainsi,

- si $\alpha=0$, la matrice $A_\alpha$ admet une valeur propre triple $\lambda=-1$,

- si $\alpha\neq0$, la matrice $A_\alpha$ admet deux valeurs propres distinctes
$\lambda_1=-1$ valeur propre double et $\lambda_2=\alpha-1$, valeur propre simple.
 

  \item {\it D\'eterminons les valeurs de $\alpha$ pour lesquelles la matrice $A_\alpha$ est diagonalisable.}

Il est clair que dans le cas $\alpha=0$, la matrice n'est pas diagonalisable, en effet si elle l'\'etait, il existerait une matrice inversible $P$ telle que $A_\alpha=P(-I)P^{-1}=-I$, ce qui n'est pas le cas.

Si $\alpha\neq0$, la matrice $A_\alpha$ est diagonalisable si le sous-espace propre associ\'e \`a la valeur propre $-1$ est de dimension $2$. D\'eterminons ce sous-espace propre.

$$E_{-1}=\ker(A_\alpha+I)=\left\{(x,y,z)\in\R^3,\ \begin{pmatrix}-1 & 0 & \alpha+1 \\  1&-2&0 \\ -1&1&\alpha\end{pmatrix}\begin{pmatrix}x \\  y \\  z\end{pmatrix}=\begin{pmatrix}-x \\  -y \\  -z\end{pmatrix}\right\}$$
ainsi,
$$(x,y,z)\in E_{-1}\iff\left\{\begin{align*}-x+(\alpha+1)z&=-x \\  x-2y&=-y \\  -x+y+\alpha z&=-z\end{align*}\right.
\iff\left\{\begin{align*}(\alpha+1)z&=0 \\  x-y=0\end{align*}\right.$$
Il faut distinguer les cas $\alpha=-1$ et $\alpha\neq -1$.

- Si $\alpha=-1$, le sous-espace $E_{-1}$ est le plan vectoriel d'\'equation $x=y$, dans ce cas la matrice $A_\alpha$ est diagonalisable.

- Si $\alpha\neq -1$, le sous-espace $E_{-1}$ est la droite vectorielle engendr\'ee par le vecteur
$(1,1,0)$, dans ce cas la matrice $A_\alpha$ n'est pas diagonalisable.

  \item {\it D\'eterminons selon la valeur de $\alpha$ le polyn\^ome minimal de $A_\alpha$.}

 Notons $Q_A$ le polyn\^ome minimal de $A_\alpha$. On sait que la matrice $A_\alpha$ est diagonalisable sur $\R$ si et seulement si son polyn\^ome minimal a toutes ses racines dans $\R$ et que celles-ci sont simples. Or, nous venons de d\'emontrer que $A_\alpha$ est diagonalisable sur $\R$ si et seulement $\alpha=-1$, on a donc 

- Si $\alpha=-1$, $A_\alpha$ est diagonalisable, donc $Q_A(X)=(X+1)(X-\alpha+1)=(X+1)(X+2)$.

- Si $\alpha\neq -1$, $A_\alpha$ n'est pas diagonalisable, donc $Q_A(X)=P_A(X)=(X+1)^2(X-\alpha+1)$.
\end{enumerate}

{\it Seconde partie} : 

On suppose d\'esormais que $\alpha=0$, on note $A=A_0$ et $f$ l'endomorphisme de $\R^3$ associ\'e \`a la matrice $A$. On a donc 
$$A=A_0=\begin{pmatrix}-1 & 0 & 1 \\  1&-2&0 \\ -1&1&0\end{pmatrix}$$
et $P_A(X)=-(X+1)^3$.

\begin{enumerate}  
 \item {\it D\'eterminons les sous-espaces propres et caract\'eristiques de $A$.}

La matrice $A$ admet une unique valeur propre $\lambda=-1$ de multiplicit\'e $3$, le sous-espace propre associ\'e est l'espace $E_{-1}=\ker(A+I)$, et on a
$$(x,y,z)\in E_{-1}\iff\left\{\begin{align*}-x+z&=-x \\  x-2y&=-y \\  -x+y&=-z\end{align*}\right.
\iff\left\{\begin{align*}z&=0 \\  x=y\end{align*}\right.$$
Le sous-espace $E_{-1}$ est la droite vectorielle engendr\'ee par le vecteur $(1,1,0)$.

Le sous-espace caract\'eristique de $A$, associ\'e \`a l'unique valeur propre $\lambda=-1$, est le sous-espace $N_{-1}=\ker(A+I)^3$, or, compte tenu du th\'eor\`eme de Hamilton-Cayley, on sait que $P_A(A)=0$, ainsi, la matrice $(A+I)^3$ est la matrice nulle, ce qui implique $N_{-1}=\R^3$, c'est donc l'espace tout entier.

  \item {\it D\'emontrons que $f$ admet un plan stable.}

La matrice de $f$ n'est pas diagonalisable, mais comme son polyn\^ome caract\'eristique se factorise sur $\R$, elle est trigonalisable, ce qui prouve qu'elle admet un plan stable, le plan engendr\'e par les deux premiers vecteurs d'une base de trigonalisation.

Par ailleurs, on a $E_{-1}=\ker(A+I)\subset\ker(A+I)^2\subset\ker(A+I)^3=\R^3$, le sous-espace 
$\ker(A+I)^2$ est clairement stable par $A$ car  pour tout $v\in\ker(A+I)^2$, $Av\in\ker(A+I)^2$, en effet 
$$(A+I)^2Av=A(A+I)^2v=0.$$
D\'emontrons que ce sous-espace est un plan. On a
$$A^2=\begin{pmatrix}-1 & 0 & 1 \\  1&-2&0 \\ -1&1&0\end{pmatrix}\begin{pmatrix}-1 & 0 & 1 \\  1&-2&0 \\ -1&1&0\end{pmatrix}=\begin{pmatrix}-1 & 1 & 1 \\  -1&1&1 \\ 0&0&0\end{pmatrix},$$
donc $\ker(A+I)^2=\{(x,y,z)\in\R^3,\ -x+y+z=0\}$, c'est bien un plan vectoriel.

  \item {\it D\'emontrons qu'il existe une base de $\R^3$ dans laquelle la matrice de $f$ est }
$$B=\begin{pmatrix}-1 & 1 & 0 \\  0&-1&1 \\ 0&0&-1\end{pmatrix}$$
{\it et trouvons une matrice $P$ inversible telle que} $A=PBP^{-1}$.

Nous cherchons des vecteurs $e_1,e_2,e_3$ tels que $Ae_1=e_1$, $Ae_2=e_1-e_2$ et $Ae_3=e_2-e_3$.
Le vecteur $e_1$ appartient \`a $E_1=\ker(A+I)$, nous choisirons $e_2\in \ker(A+I)^2$ tel que $(e_1,e_2)$ soit une base de $\ker(A+I)^2$. Remarquons que si l'on cherche $e_2=(x,y,z)$ tel que
$Ae_2=e_1-e_2$, on obtient le syst\`eme 
$$\begin{pmatrix}-1 & 0 & 1 \\  1&-2&0 \\ -1&1&0\end{pmatrix}\begin{pmatrix}x \\  y \\  z\end{pmatrix}=\begin{pmatrix}1-x \\  1-y \\ -z\end{pmatrix}\iff
\left\{\begin{align*}-x+z&=1-x \\  x-2y&=1-y \\  -x+y&=-z\end{align*}\right.
\iff\left\{\begin{align*}z&=1 \\  x-y&=1\end{align*}\right.$$
ce qui nous donne bien un vecteur de $\ker(A+I)^2$.
Ainsi, les vecteurs $e_1=(1,1,0)$ et $e_2=(1,0,1)$ conviennent. Il nous reste \`a chercher un vecteur $e_3$ tel que $Ae_3=e_2-e_3$, c'est-\`a-dire
$$\begin{pmatrix}-1 & 0 & 1 \\  1&-2&0 \\ -1&1&0\end{pmatrix}\begin{pmatrix}x \\  y \\  z\end{pmatrix}=\begin{pmatrix}1-x \\  -y \\ 1-z\end{pmatrix}\iff
\left\{\begin{align*}-x+z&=1-x \\  x-2y&=-y \\  -x+y&=1-z\end{align*}\right.
\iff\left\{\begin{align*}z&=1 \\  x&=y\end{align*}\right.$$
Le vecteur $e_3=(0,0,1)$ convient. On obtient alors la matrice $P$ suivante qui est inversible et v\'erifie $A=PBP^{-1}$,
$$P=\begin{pmatrix}1 & 1 & 0 \\  1&0&0 \\ 0&1&1\end{pmatrix}$$  

  \item{\it D\'ecomposition de Dunford de $B$}

On a 
$$B=\begin{pmatrix}-1 & 1 & 0 \\  0&-1&1 \\ 0&0&-1\end{pmatrix}=
\begin{pmatrix}-1 & 0& 0 \\  0&-1&0 \\ 0&0&-1\end{pmatrix}+
\begin{pmatrix}0& 1 & 0 \\  0&0&1 \\ 0&0&0\end{pmatrix}$$
et il est clair que les deux matrices commutent car l'une est \'egale à $-I$. Or, il existe un unique couple de matrice $D$ et $N$, $D$ diagonalisable et $N$ nilpotente, telles que $B=D+N$ et $DN=ND$. C'est donc l\`a la d\'ecomposition de Dunford, $B=D+N$ avec
$$D=\begin{pmatrix}-1 & 0& 0 \\  0&-1&0 \\ 0&0&-1\end{pmatrix}\ \ {\hbox{et}}\ \ 
N=\begin{pmatrix}0& 1 & 0 \\  0&0&1 \\ 0&0&0\end{pmatrix}.$$

  \item{\it Pour $t\in\R$, calculons $\exp tB$ et exprimons $\exp tA$ \`a l'aide de $P$ et $\exp tB$.}

Remarquons tout d'abord que $N^2=0$ donc pour tout $t\in\R$, $(tN)^2=0$ et l'expenentielle est \'egale \`a $\exp(tN)=I+tN$, par ailleurs $ND=DN$, donc  pour tout $t\in\R$, les matrices $tN$ et $tD$ commutent \'egalement, $(tN)(tD)=(tD)(tN)$, on a donc
$$\exp(tB)=\exp(tD+tN)=\exp(tD)\exp(tN)=\exp(-tI)\exp(-tN)=e^{-t}(I+tN).$$
D'o\`u
$$\exp(tB)=e^{-t}\begin{pmatrix}1 & t & 0 \\  0&1&t \\ 0&0&1\end{pmatrix}.$$
Pour d\'eterminer l'exponentielle de la matrice $tA$, on \'ecrit
$$\exp(tA)=\exp(t(PBP^{-1}))=\exp(P(tA)P^{-1})=P\exp(tB)P^{-1}.$$
  \item {\it Solutions des syst\`emes diff\'erentiels $Y'=BY$ et $X'=AX$.}

La solution g\'en\'erale du syst\`eme $Y'=BY$ s'\'ecrit 
$$S(t)=\exp(tB)v$$
o\`u $v=(a,b,c)$ est un vecteur de $\R^3$. La solution $S: \R\mapsto \R^3$  s'\'ecrit donc
$$S(t)=\begin{pmatrix}x(t) \\  y(t) \\  z(t)\end{pmatrix}
=e^{-t}\begin{pmatrix}1 & t & 0 \\  0&1&t \\ 0&0&1\end{pmatrix}\begin{pmatrix}a \\  b \\  c\end{pmatrix}
=e^{-t}\begin{pmatrix}a+bt \\  b+ct \\  c\end{pmatrix}$$
Pour obtenir la solution du syst\`eme $X'=AX$, on \'ecrit
$$X'=AX\iff X'=(PBP^{-1})X\iff P^{-1}X'=(BP^{-1})X\iff (P^{-1}X)'=B(P^{-1}X)$$
ainsi, en notant $Y=P^{-1}X$ ou encore $X=PY$, les solutions du syst\`eme $X'=AX$ sont les $PS(t)$ o\`u $P$ est la matrice v\'erifiant
$A=PBP^{-1}$ et $S$ une solution du syst\`eme $Y'=BY$.

La solution g\'en\'erale du syst\`eme $X'=AX$ s'\'ecrit donc 
$$X(t)=\begin{pmatrix}x(t) \\  y(t) \\  z(t)\end{pmatrix}
=\begin{pmatrix}1 & 1 & 0 \\  1&0&0 \\ 0&1&1\end{pmatrix}
\begin{pmatrix}e^{-t}(a+bt) \\  e^{-t}(b+ct) \\  e^{-t}c\end{pmatrix}
=e^{-t}\begin{pmatrix}a+b+(c+b)t \\  a+bt \\  b+c+ct\end{pmatrix}$$
o\`u $(a,b,c)\in\R^3$. 

\end{enumerate}


\centerline{\bf II}

Soient $K$ un corps,  $N\in M_n(K)$ une matrice nilpotente et $A$ une matrice telle que 

$AN=NA$.

\begin{enumerate}  
 \item {\it D\'eterminons les valeurs propres de $N$.}

La matrice $N$ \'etant nilpotente, il existe un entier naturel $m$ tel que $N^n=0$, on a donc 
$\det N^m=(\det N)^m=0$ donc $\det N=0$, l'endomophisme de matrice $N$ n'est pas bijectif ce qui prouve que $0$ est valeur propre de $N$, c'est la seule, en effet si $\lambda$ est une autre valeur propre et $x\neq 0$ un vecteur propre associ\'e \`a $\lambda$ on a 
$$Nx=\lambda x\Rightarrow N^m x=\lambda^m x$$
d'o\`u $\lambda^mx=0$, mais $x\neq 0$ donc $\lambda=0$. Ainsi la matrice $N$ admet une unique valeur propre $\lambda=0$ de multiplicit\'e $n$.

  \item{\it D\'emontrons que $N$ est trigonalisable.}

Le polyn\^ome caract\'eristique de $N$ admet une unique racine $0\in K$, toutes ses racines sont donc dans $K$, ce qui prouve que la matrice $N$ est trigonalisable. Elle est semblable \`a une matrice triangulaire n'ayant que des $0$ sur la diagonale.

  \item{\it D\'emontrons que $\det (I+N)=1$.}

Compte tenu de ce qui pr\'ec\`ede, la matrice $N+I$ est une matrice triangulaire n'ayant que des $1$ sur la diagonale, or le d\'eterminant d'une matrice triangulaire est le produit de ses termes diagonaux, ainsi on a bien $\det(N+I)=1$.

  \item {\it On suppose $A$ inversible. D\'emontrons que les matrices $AN$ et $NA^{-1}$  sont nilpotentes.}

Comme les matrices $A$ et $N$ commutent, pour tout $k\in \N$, on a $(AN)^k=A^kN^k$ donc pour $k=m$, 
$(AN)^m=A^mN^m=A.0=0$ ce qui prouve que $AN$ est nilpotente. De m\^eme $NA^{-1}=A^{-1}N$ et $NA^{-1}$ est nilpotente.

{\it On en d\'eduit que $$\det(A+N)=\det A.$$}

L'\'egalit\'e $AN=NA$ implique $N=ANA^{-1}$ ainsi, on a
$$\det(N+A)=\det(ANA^{-1}+A)=\det(A(NA^{-1}+I))=\det A\det(NA^{-1}+I)=\det A.$$

  \item{\it On suppose $A$ non inversible. En exprimant $(A+N)^k$ pour $k\in\N$, d\'emontrons que $\det(A+N)=0.$}

Comme les matrices $A$ et $N$ commutent, on peut utiliser la formule du bin\^ome de Newton pour calculer les puissances de $A+N$. Soit $m$ tel que $N^m=0$ et, pour tout $k<m, N^k\neq 0$ on a alors
$$(A+N)^m=\sum_{k=0}^mC_m^kA^kN^{m-k}=\sum_{k=1}^mC_m^kA^kN^{m-k}=A\sum_{k=1}^mC_m^kA^{k-1}N^{m-k}$$
ainsi
$$\det((A+N)^m)=\det A.\det\sum_{k=1}^mC_m^kA^{k-1}N^{m-k}=0$$
car $\det A=0$. Or, $\det((A+N)^m)=(\det(A+N))^m$, on a donc bien $\det(A+N)=0$.

\end{enumerate}
\fincorrection
\correction{002587}
{\it Soit $\displaystyle A=\begin{pmatrix}a&c \\  c&d\end{pmatrix}\in M_2(\R)$,
 on montre que $A$ est diagonalisable sur $\R$.}

Le polyn\^ome caract\'eristique $P_A(X)$ est \'egal \`a
$$P_A(X)=\begin{vmatrix}a-X&c \\  c& d-x\end{vmatrix}=(a-X)(d-X)-c^2=X^2-(a+d)X+ad-c^2,$$
d\'eterminons ses racines : calculons le discriminant :
\begin{align*}
\Delta&=(a+d)^2-4(ad-c^2) \\  &=a^2+d^2+2ad-4ad+4c^2 \\ &=a^2+d^2-2ad+4c^2 \\ &=(a-d)^2+4c^2\geq 0
\end{align*}
On a $\Delta=0\iff a-d=0\ {\hbox{et}}\ c=0$, mais, si $c=0$, la matrice $A$ est d\'ej\`a diagonale. Sinon $\Delta>0$ et le polyn\^ome caract\'eristique admet
deux racines r\'eelles distinctes, ce qui prouve que la matrice est toujours diagonalisable dans $\R$.

\fincorrection
\correction{002591}
 Soit $a\in\R$, notons $A$ la matrice suivante
$$A=\begin{pmatrix}0&1 \\  -a&1+a\end{pmatrix}.$$
On d\'efinit une suite $(u_n)_{n\in\N}$, par la donn\'ee de $u_0$ et $u_1$ et la relation de r\'ecurrence suivante, pour $n\in\N$
$$u_{n+2}=(1+a)u_{n+1}-au_n$$
\begin{enumerate}
 \item {\it Pour quelles valeurs de $a$ la matrice $A$ est-elle diagonalisable ?}

Calculons le polyn\^ome caract\'eristique $P_A(X)$ :
$$P_A(X)=\begin{vmatrix}-X&1 \\  -a& 1+a-X\end{vmatrix}=-X(1+a-X)+a=X^2-(1+a)X+a.$$
La matrice $A$ est diagonalisable sur $\R$ si le polyn\^ome $P_A$ admet deux racines distinctes dans $\R$. En effet, si $P_A$ admet une racine double $r$ et $A$ diagonalisable, alors l'endomorphisme de matrice $A$ est \'egal \`a $r\mathrm{Id}_E$, ce qui n'est pas le cas. Calculons donc le discriminant du polyn\^ome caract\'eristique.
$$\Delta=(1+a)^2-4a=1+a^2+2a-4a=1+a^2-2a=(1-a)^2.$$
Ainsi la matrice $A$ est diagonalisable pour tout $a\neq 1$.
 \item {\it Lorsque $A$ est diagonalisable, calculons $A^n$ pour $n\in \N$}.

Lorsque $A$ est diagonalisable, il existe une matrice inversible $P$ et une matrice diagonale $D$ telles que $A=PDP^{-1}$, ainsi pour tout $n\in\N$, on a $A^n=PD^nP^{-1}$. D\'eterminons les matrices $P$ et $D$. Pour cela calculons les deux valeurs propres de $A$, ce sont les racines du polyn\^ome $P_A$, on a donc 
$$\lambda_1={\frac{1+a+1-a}{2}}=1\ {\hbox{et}}\ \lambda_2={\frac{1+a-1+a}{2}}=a.$$
D\'eterminons maintenant des vecteurs propres associ\'es aux valeurs propres $1$ et $a$. On cherche des vecteurs $\vec{e_1}$ et $\vec{e_2}$ tels que 
$A\vec{e_1}=\vec{e_1}$ et $A\vec{e_2}=a\vec{e_2}$.
$$\begin{pmatrix}0&1 \\  -a&1+a\end{pmatrix}\begin{pmatrix}x \\  y\end{pmatrix}=\begin{pmatrix}x \\  y\end{pmatrix}\iff y=x$$ et
$$\begin{pmatrix}0&1 \\  -a&1+a\end{pmatrix}\begin{pmatrix}x \\  y\end{pmatrix}=a\begin{pmatrix}x \\  y\end{pmatrix}\iff y=ax$$
ainsi on peut choisir $\vec e_1=(1,1)$ et $\vec e_2=(1,a)$. On a alors
$$P=\begin{pmatrix}1&1 \\  1&a\end{pmatrix},\  D=\begin{pmatrix}1&0 \\  0&a\end{pmatrix},\ 
P^{-1}={\frac{1} {a-1}}\begin{pmatrix}a&-1 \\  -1&1\end{pmatrix}.$$
D'o\`u, pour tout $n\in\N$, 
$$A^n=PD^nP^{-1}=P\begin{pmatrix}1&0 \\  0&a^n\end{pmatrix}P^{-1}={\frac{1} {a-1}}\begin{pmatrix}a-a^n&a^n-1 \\  a-a^{n+1}&a^{n+1}-1\end{pmatrix}$$
 \item {\it On suppose $A$ diagonalisable. On note $U_n$ le vecteur $U_{n}=\begin{pmatrix}u_{n} \\  u_{n+1}\end{pmatrix}$,  on exprime $U_{n+1}$ en fonction de $U_n$ et de $A$, puis $U_n$ en fonction de $U_0$ et de $A$.}

On a, par d\'efinition, pour tout $n\in\N$, $u_{n+2}=(1+a)u_{n+1}-au_n$, ainsi, 
$$U_{n+1}=\begin{pmatrix}u_{n+1} \\  u_{n+2}\end{pmatrix}=\begin{pmatrix}0&1 \\  -a&1+a\end{pmatrix}\begin{pmatrix}u_{n} \\  u_{n+1}\end{pmatrix}=AU_n.$$
On a donc $U_1=AU_0$, montrons par r\'ecurrence sur $n$, que pour tout $n\in\N$, $U_n=A^nU_0$. C'est vrai pour $n=0$, $U_0=A^0U_0=IU_0=U_0$ et pour $n=1$. Soit $n$ fix\'e pour lequel on suppose $U_n=A^nU_0$, on a alors $U^{n+1}=AU_n=A.A^nU_0=A^{n+1}U_0$, le r\'esultat est donc vrai pour tout entier naturel $n$.

La matrice $A$ \'etant suppos\'ee diagonalisable, on a donc, pour $n\in\N$,
$$U_n=A^nU_0=PD^nP^{-1}U_0={\frac {1}{a-1}}\begin{pmatrix}a-a^n&a^n-1 \\  a-a^{n+1}&a^{n+1}-1\end{pmatrix}\begin{pmatrix}u_{0} \\  u_{1}\end{pmatrix},$$
ainsi on peut exprimer pour $n\in\N$, le terme g\'en\'eral de la suite $u_n$ en fonction des premiers termes $u_0$ et $u_1$, on a 
$$u_n={\frac{1} {a-1}}\left((a-a^n)u_0+(a^n-1)u_1\right).$$

\end{enumerate}
\fincorrection
\correction{002592}
 Soit $A$ la matrice de $M_3(\R)$ suivante :
$$A=\begin{pmatrix}0&1&0 \\  -4&4&0 \\ -2&1&2\end{pmatrix}.$$ 
\begin{enumerate}
 \item {\it La matrice $A$ est-elle diagonalisable ?}

Calculons son polyn\^ome caract\'eristique 
$$P_A(X)=\begin{vmatrix}-X&1&0 \\  -4&4-X&0 \\ -2&1&2-X\end{vmatrix}=(2-X)(X^2-4X+4)=(2-X)^3.$$
la matrice $A$ admet une unique valeur propre $2$, si elle \'etait diagonalisable, elle serait semblable \`a la matrice $2.I_3$, elle serait donc \'egale \`a $2I_3$ ce qui n'est pas le cas, elle n'est donc pas diagonalisable. 

 \item  {\it Calculons $(A-2I_3)^2$, puis $(A-2I_3)^n$ pour tout $n\in\N$}.

On a 
$$(A-2I_3)^2=\begin{pmatrix}-2&1&0 \\  -4&2&0 \\ -2&1&0\end{pmatrix}\begin{pmatrix}-2&1&0 \\  -4&2&0 \\ -2&1&0\end{pmatrix}=\begin{pmatrix}0&0&0 \\  0&0&0 \\ 0&0&0\end{pmatrix},$$
ainsi, $(A-2I_3)^0=I$, $(A-2I_3)^1=\begin{pmatrix}-2&1&0 \\  -4&2&0 \\ -2&1&0\end{pmatrix}$, et,  pour tout $n\geq2$, on a $(A-2I_3)^n=0$.

{\it On en d\'eduit $A^n$}

Notons $B=A-2I_3$, on a $A=A-2I_3+2I_3=B+2I_3$ avec $B^n=0$ pour $n\geq2$. Par ailleurs, les matrices $B$ et $2I_3$ commutent, ainsi
$$A^n=(B+2I_3)^n=\sum^n_{k=0}C_n^kB^k(2I_3)^{n-k}$$
o\`u les $C_n^k$ sont les coefficients du bin\^ome de Newton :
$$C_n^k={\frac{n!}{k!(n-k)!}}.$$ 
Or, pour $k\geq2$, on a $B^k=0$ d'o\`u pour $n\geq2$,
\begin{align*}
A^n&=C_n^0B^0(2I_3)^n+C_n^1B^1(2I_3)^{n-1} \\ 
&=2^nI_3+2^{n-1}nB \\ 
&=2^nI_3+2^{n-1}n(A-2I_3) \\ 
&=2^n(1-n)I_3+2^{n-1}nA.
\end{align*}

\end{enumerate}
\fincorrection
\correction{002593}
 Soit $f$ l'endomorphisme de $\R^4$ dont la matrice dans la base canonique  est 
$$A=\begin{pmatrix}-8&-3&-3&1 \\  6&3&2&-1 \\ 26&7&10&-2 \\ 0&0&0&2\end{pmatrix}.$$ 

\begin{enumerate}
 \item {\it D\'emontrons que $1$ et $2$ sont des valeurs propres de $f$}.

Pour cela montrons que $\det(A-I)=0$ et $\det(A-2I)=0$. On a
$$\det(A-I)=\begin{vmatrix}-9&-3&-3&1 \\  6&2&2&-1 \\ 26&7&9&-2 \\ 0&0&0&1\end{vmatrix}=
-\begin{vmatrix}-9&-3&-3 \\  6&2&2 \\ 26&7&9\end{vmatrix}=-\begin{vmatrix}-9&-3&0 \\  6&2&0 \\ 26&7&2\end{vmatrix}=-2
\begin{vmatrix}-9&-3 \\  6&2\end{vmatrix}=0.$$
Et
$$\det(A-2I)=\begin{vmatrix}-10&-3&-3&1 \\  6&1&2&-1 \\ 26&7&8&-2 \\ 0&0&0&0\end{vmatrix}=0.$$ Ainsi, les r\'eels $1$ et $2$ sont bien valeurs propres de la matrice $A$.
 \item {\it D\'eterminons des vecteurs propres de $f$ associ\'es aux valeurs propres $1$ et $2$}.

Soit $\vec e=(x,y,z,t)$ tel que $A\vec e=\vec e$, on r\'esout alors le syst\`eme 
$$\left\{{\begin{align*}-9x-3y-3z+t&=0 \\ 6x+2y+2z-t&=0 \\ 26x+7y+9z-2t&=0 \\  t&=0\end{align*}}\right.\iff
\left\{{\begin{align*}3x+y+z&=0 \\ 26x+7y+9z&=0 \\  t&=0\end{align*}}\right.\iff
\left\{{\begin{align*}x&=-2y \\  z&=5y \\  t&=0\end{align*}}\right.,$$
ce syst\`eme repr\'esente une droite vectorielle engendr\'ee, par exemple, par le vecteur 

$\vec e=(-2,1,5,0)$.

Soit $\vec u=(x,y,z,t)$ tel que $A\vec u=2\vec u$, on r\'esout 
$$\left\{{\begin{align*}-10x-3y-3z+t&=0 \\ 6x+y+2z-t&=0 \\ 26x+7y+8z-2t&=0\end{align*}}\right.\iff
\left\{{\begin{align*}10x+3y+3z&=0 \\ 6x+y+2z&=0 \\  t&=0\end{align*}}\right.\iff
\left\{{\begin{align*}3y&=-2x \\  3z&=-8x \\  t&=0\end{align*}}\right.,$$
ce syst\`eme repr\'esente une droite vectorielle engendr\'ee, par exemple, par le vecteur

$\vec u=(3,-2,-8,0).$ 
 \item {\it  On consid\`ere le vecteur $\vec u$ pr\'ec\'edent et on determine des vecteurs $\vec v$ et $\vec w$ tels que }
$$f(\vec v)=2\vec v+\vec u\ {\hbox{et}}\ f(\vec w)=2\vec w+\vec v.$$
Pour d\'eterminer le vecteur $\vec v=(x,y,z,t)$, on r\'esout le syst\`eme
$$\left\{{\begin{align*}-10x-3y-3z+t&=3 \\ 6x+y+2z-t&=-2 \\ 26x+7y+8z-2t&=-8\end{align*}}\right.\iff
\left\{{\begin{align*}10x+3y+3z&=-3 \\ 6x+y+2z&=-2 \\  t&=0\end{align*}}\right.,$$
le vecteur $\vec v=(0,0,-1,0)$ convient.
Pour d\'eterminer le vecteur $\vec w=(x,y,z,t)$, on r\'esout le syst\`eme
$$\left\{{\begin{align*}-10x-3y-3z+t&=0 \\ 6x+y+2z-t&=0 \\ 26x+7y+8z-2t&=-1\end{align*}}\right.\iff
\left\{{\begin{align*}10x+3y+3z&=-1 \\ 6x+y+2z&=-1 \\  t&=-1\end{align*}}\right.,$$
le vecteur $\vec w=(1/2,0,-2,-1)$ convient.
 \item {\it Les vecteurs $\vec e$, $\vec u$, $\vec v$ et $\vec w$ sont ceux d\'efinis pr\'ec\'edemment. On d\'emontre que $(\vec e,\vec u,\vec v,\vec w)$ est une base de $\R^4$ et on  donne la matrice de $f$ dans cette base.}

La matrice $M$ des vecteurs $\vec e, \vec u, \vec v, \vec w$ dans la base cannonique est de rang $4$ car son d\'eterminant est non nul, en effet
$$\det M=\begin{vmatrix}-2&3&0&1/2 \\  1&-2&0&0 \\ 5&-8&-1&-2 \\ 0&0&0&-1\end{vmatrix}=
\begin{vmatrix}-2&3&0 \\  1&-2&0 \\ 5&-8&-1\end{vmatrix}=-\begin{vmatrix}-2&3 \\  1&-2\end{vmatrix}=-1 .$$
Compte tenu des d\'efinitions des vecteurs $\vec e, \vec u, \vec v, \vec w$, la matrice $B$ de l'endomorphisme $f$ dans la base $(\vec e,\vec u,\vec v,\vec w)$ s'\'ecrit
$$B=\begin{pmatrix}1&0&0&0 \\ 0&2&1&0 \\ 0&0&2&1 \\ 0&0&0&2\end{pmatrix}.$$
  \item {\it La matrice $A$ est-elle diagonalisable ?}

D'apr\`es la question pr\'ec\'edente, les valeurs propres de $f$ sont $1$, valeur propre simple, et $2$ de multiplicit\'e $3$. Nous avons vu dans le b) que le sous-espace propre associ\'e \`a la valeur propre $2$ est de dimension $1\neq 3$, ainsi, la matrice $A$ n'est pas diagonalisable.
\end{enumerate}
\fincorrection
\correction{002599}
Soit $m\in\R$, et $A$ la matrice
$$A=\begin{pmatrix}1+m&1+m&1 \\  -m&-m&-1 \\  m&m-1&0\end{pmatrix}$$
\begin{enumerate}
 \item  {\it Factorisons le polyn\^ome caract\'eristique de $A$ et montrons que les valeurs propres de $A$ sont $-1$ et $1$.} (1,5 points)
\begin{align*}
P_A(X)&=\begin{vmatrix}1+m-X&1+m&1 \\  -m&-m-X&-1 \\  m&m-1&-X\end{vmatrix} \\ 
&=\begin{vmatrix}1+m-X&1+m&1 \\  -m&-m-X&-1 \\  0&-X-1&-X-1\end{vmatrix} \\ 
&=\begin{vmatrix}1+m-X&m&1 \\  -m&1-m-X&-1 \\  0&0&-X-1\end{vmatrix} \\ 
&=(-X-1)\begin{vmatrix}1+m-X&m \\  -m&1-m-X\end{vmatrix} \\ 
&=-(1+X)\left((1-X)^2-m^2+m^2\right) \\ 
&=-(1+X)(1-X)^2
\end{align*}
Ainsi, les valeurs propres de la matrice $A$ sont $-1$, valeur propre simple, et $1$, valeur propre double.

 \item {\it Pour quelles valeurs de $m$ la matrice est-elle diagonalisable ? } (1,5 points)

La matrice $A$ est diagonalisable si et seulement si le sous-espace propre associ\'e \`a la valeur propre $1$ est de dimension $2$. D\'eterminons donc ce sous-espace propre $E_1=\{\vec u=(x,y,z)\in\R^3,\ A\vec u=\vec u\}$.
$$\begin{align*}
A\vec u=\vec u
&\iff\left\{{\begin{align*}mx+(1+m)y+z&=0 \\  -mx-(1+m)y-z&=0 \\  mx+(m-1)y-z&=0\end{align*}} \right. \\ 
&\iff\left\{{\begin{align*}mx+(1+m)y+z&=0 \\  mx+(m-1)y-z&=0\end{align*}}\right. \\ 
&\iff\left\{{\begin{align*}2y+2z&=0 \\  2mx+2my&=0\end{align*}}\right. \\ 
&\iff\left\{{\begin{align*}y+z&=0 \\  m(x+y)&=0\end{align*}}\right. \\ 
\end{align*}$$
Ainsi, l'espace $E_1$ est de dimension $2$ si et seulement si $m=0$, c'est alors le plan d'\'equation $y+z=0$, sinon c'est une droite, intersection des deux plans $y+z=0$ et $x+y=0$.

{\it D\'eterminons suivant les valeurs de $m$ le polyn\^ome minimal de $A$.} (1 point)

Si $m=0$, la matrice $A$ est diagonalisable, son polyn\^ome minimal n'a que des racines simples, il est \'egal \`a 
$$Q(X)=(X-1)(X+1).$$
Si $m\neq 0$, la matrice $A$ n'est pas diagonalisable, son polyn\^ome minimal ne peut pas avoir uniquement des racines simples, il est donc \'egal \`a
$$Q(X)=(X-1)^2(X+1).$$

\end{enumerate}
\fincorrection
\correction{002600}
\begin{enumerate}
 \item  {\it Donnons un exemple de matrice dans $M_2(\R)$, diagonalisable sur $\C$ mais non diagonalisable sur $\R$}. (2 points)

Soit $A$ la matrice $A=\begin{pmatrix}0&-1 \\  1&0\end{pmatrix}$. Son polyn\^ome caract\'eristique est \'egal \`a
$$P_A(X)=\begin{vmatrix}-X&-1 \\  1&-X\end{vmatrix}=X^2+1.$$
Le polyn\^ome caract\'eristique de $A$ admet deux racines complexes conjugu\'ees distinctes $i$ et $-i$ elle est donc diagonalisable sur $\C$ mais elle ne l'est pas sur $\R$.

 \item {\it Donnons un exemple de matrice dans $M_2(\R)$ non diagonalisable, ni sur $\C$, ni sur $\R$.} (2 points) 
 
Soit $A$ la matrice $A=\begin{pmatrix}1&1 \\  0&1\end{pmatrix}$. Son polyn\^ome caract\'eristique est \'egal \`a
$$P_A(X)=\begin{vmatrix}1-X&1 \\  0&1-X\end{vmatrix}=(1-X)^2.$$
Le polyn\^ome caract\'eristique de $A$ admet une racine double $1$, la matrice $A$ admet l'unique valeur propre $1$, or, elle n'est pas \'egale \`a l'identit\'e, par cons\'equent, elle n'est diagonalisable, ni 
sur $\C$, ni sur $\R$.

\end{enumerate}
\fincorrection
\correction{002601}
Soit $A$ la matrice suivante :
$$A=\begin{pmatrix}0&1 \\ 1&0\end{pmatrix}$$
\begin{enumerate}
 \item {\it Diagonalisons la matrice $A$.} (2 points)

Son polyn\^ome caract\'eristique est \'egal \`a
$$P_A(X)=\begin{vmatrix}-X&1 \\  1&-X\end{vmatrix}=X^2-1=(X-1)(X+1).$$
La matrice $A$ admet deux valeurs propres distinctes, elle est donc diagonalisable.

{\it D\'eterminons une base de vecteurs propres de $A$}. 

Soit $\vec u=(x,y)\in\R^2$,
$$A\vec u=\vec u\iff x=y\ {\hbox{et}}\ A\vec u=-\vec u\iff x=-y.$$
Notons $\vec u_1=(1,1)$ et $\vec u_2=(-1,1)$, le vecteur $\vec u_1$ est un vecteur propre associ\'e \`a la valeur propre $1$ et le vecteur $\vec u_2$ est un vecteur propre associ\'e \`a la valeur propre $-1$, ils sont lin\'eairement ind\'ependants, ils forment donc une base de $\R^2$. Ainsi, on a $A=PDP^{-1}$, o\`u
$$P=\begin{pmatrix}1&-1 \\  1&1\end{pmatrix}\ {\hbox{et}}\ D=\begin{pmatrix}1&0 \\  0&-1\end{pmatrix}$$

 \item  {\it Exprimons les solutions du syst\`eme diff\'erentiel $X'=AX$ dans une base de vecteurs propres et tra\c cons ses trajectoires.} (3 points)

Soit $Y$ tel que $PY=X$, on a alors 
$$X'=AX\iff PY'=APY\iff Y'=P^{-1}APY\iff Y'=DY.$$
Les solutions du syst\`eme diff\'erentiel $X'=AX$ dans la base de vecteurs propres $(\vec u_1,\vec u_2)$ sont les solutions du syst\`eme $Y'=DY$. Si $Y=(x,y)$, on a $x'(t)=x(t)$ et $y'(t)=-y(t)$, ainsi, les solutions du syst\`eme sont $x(t)=ae^t$ et $y(t)=be^{-t}$ o\`u $a$ et $b$ sont des constantes r\'eelles arbitraires. Les trajectoires, exprim\'ees dans la base de vecteurs propres $(\vec u_1,\vec u_2)$,  sont donc les courbes d'\'equation $y=c/x$ avec $c\in\R$, ce sont des branches d'hyperboles.

\end{enumerate}
\fincorrection
\correction{002603}
{\it Soit $a\in\R$ et $A$ la matrice suivante}
$$A=\begin{pmatrix}1&a&0 \\  a&0&1 \\  0&1&a\end{pmatrix}.$$
\begin{enumerate}
 \item {\it Calculons le d\'eterminant de $A$ et d\'eterminons 
pour quelles valeurs de $a$ la matrice est inversible.}

On développe le déterminant par rapport à la première colonne, on obtient
$$\det A=\begin{vmatrix}1&a&0 \\  a&0&1 \\  0&1&a\end{vmatrix}=\begin{vmatrix}0&1 \\  1&a\end{vmatrix}-a\begin{vmatrix}a&0 \\  1&a\end{vmatrix}=-1-a^3.$$
La matrice $A$ est inversible si et seulement si son déterminant est non nul.
$$\det A\neq 0\iff 1+a^3\neq 0\iff a\neq -1.$$
 \item {\it Calculons $A^{-1}$ lorsque $A$ est inversible.}

On suppose $a\neq-1$, on a $\displaystyle A^{-1}={\frac{1}{\det A}}\,^t\!\!\tilde A$, où $\tilde A$ est la comatrice de $A$ et $^t\!\!\tilde A$ la transposée de $\tilde A$. On a 
$$\tilde A=\begin{pmatrix}-1&-a^2&a \\  -a^2&a&-1 \\  a&-1&-a^2\end{pmatrix}=^t\!\!\tilde A.$$
D'où $\displaystyle A^{-1}={\frac{1}{1+a^3}}\begin{pmatrix}1&a^2&-a \\  a^2&-a&1 \\  -a&1&a^2\end{pmatrix}.$
\end{enumerate}
\fincorrection
\correction{002604}
{\it Soit $\theta\in \R$, on consid\`ere l'endomorphisme $f$ de $\R^3$ dont la matrice dans la base canonique est la suivante}
$$A=\begin{pmatrix}\cos\theta&-\sin\theta&0 \\  \sin\theta&\cos\theta&0 \\ 0&0&1 \\ \end{pmatrix}.$$
\begin{enumerate}
 \item  {\it Déterminons la nature géométrique de cet endomorphisme.}

Notons $(\vec i,\vec j,\vec k)$ la base canonique de $\R^3$, la matrice $A$ est la matrice de la rotation d'axe $\R\vec k$ d'angle $\theta$. 

On peut ajouter que les vecteurs colinéaires à $\vec k$ sont fixes. Un vecteur de coordonnées $(x,y,z)$ est envoyé sur le vecteur $(x\cos\theta-y\sin\theta,x\sin\theta+y\cos\theta,z)$, sa composante dans le plan engendré par $\vec i$ et $\vec j$ subit la rotation plane d'angle $\theta$.
 \item {\it Démontrons que, pour tout $\theta\in\R\setminus\pi\Z$, la matrice $A$ admet une unique valeur propre réelle et déterminons son sous-espace propre associé}. 

Calculons le polyn\^ome caractéristique de la matrice $A$.
\begin{align*}
 P_A(X)=\begin{vmatrix}\cos\theta-X&-\sin\theta&0 \\  \sin\theta&\cos\theta-X&0 \\ 0&0&1-X \\ \end{vmatrix} &=[(\cos\theta-X)^2+\sin^2\theta](1-X) \\  &=(1-X)(X^2-2X\cos\theta+1
\end{align*}
Cherchons les racines du polynôme $X^2-2X\cos\theta+1$, pour cela on calcule son discrimminant réduit
$$\Delta'=\cos^2\theta-1=-\sin^2\theta<0,$$
en effet, si $\theta\in\R\setminus\pi\Z$, alors $\sin\theta\neq 0$, donc le polyn\^ome $P_A$ n'admet qu'une racine réelle $\lambda=1$. Son sous-espace propre associé est de dimension $1$, c'est l'axe $\R\vec k$ de la rotation. 

{\it Cas où $\theta\in\pi\Z$ }

On distingue les cas $\theta=n\pi$ avec $n$ pair ou impair :

- Si $\theta=2n\pi,\ n\in\Z$, alors $\displaystyle A=\begin{pmatrix}1&0&0 \\  0&1&0 \\ 0&0&1 \\ \end{pmatrix}$, c'est la matrice de l'identité.

- Si $\theta=(2n+1)\pi,\ n\in\Z$, alors $\displaystyle A=\begin{pmatrix}-1&0&0 \\  0&-1&0 \\ 0&0&1 \\ \end{pmatrix}$, c'est la matrice de la symétrie orthogonale par rapport à l'axe $\R\vec k$. Elle admet deux valeurs propres, la valeur propre $1$ dont le sous-espace propre est l'axe $\R\vec k$ et la valeur propre $-1$ dont le sous-espace propre est le plan $\R\vec i+\R\vec j$.

\end{enumerate}
\fincorrection
\correction{002605}
{\it Soit $u$ l'endomorphisme de $\R^3$ dont la matrice dans la base canonique est}  
$$A=\begin{pmatrix}-4&-2&-2 \\  2&0&2 \\  3&3&1\end{pmatrix}.$$ 
\begin{enumerate}
 \item  {\it D\'eterminons et factorisons le polyn\^ome caract\'eristique de $A$.}

Par opérations sur les colonnes puis les lignes, on a
$$P_A(X)=\begin{vmatrix}-4-X&-2&-2 \\  2&-X&2 \\  3&3&1-X\end{vmatrix}=\begin{vmatrix}-4-X&0&-2 \\  2&-X-2&2 \\  3&2+X&1-X\end{vmatrix},$$ d'où
$$P_A(X)=\begin{vmatrix}-4-X&0&-2 \\  2&-X-2&2 \\  5&0&3-X\end{vmatrix}$$
et, en développant par rapport à la deuxième colonne
$$P_A(X)=-(X+2)[(-4-X)(3-X)+10]=-(X+2)(X^2+X-2)=-(X+2)^2(X-1).$$
 \item {\it D\'emontrons que les valeurs propres de $A$ sont $1$ et $-2$ et d\'eterminons les sous-espaces propres associ\'es}.

Les valeurs propres de $A$ sont les racines du polyn\^ome caracteristique, c'est-à-dire, $1$, valeur propre simple et, $-2$, valeur propre double.

Notons $E_1$ le sous-espace propre associé à la valeur propre $1$, 
$$E_1=\{\vec u=(x,y,z),\ A.\vec u=\vec u\}.$$
Ainsi
$$\vec u=(x,y,z)\in E_1\iff\left\{\begin{align*}-&5x-2y-2z=0 \\ &2x-y+2z=0 \\ &3x+3y=0\end{align*}\right.\iff
\left\{\begin{align*}&y=-x \\ &3x+2z=0\end{align*}\right.$$
Le sous-espace propre $E_1$ est donc une droite vectorielle dont un vecteur directeur est donné, par exemple, par $\vec e_1=(-2,2,3)$.

Notons $E_{-2}$ le sous-espace propre associé à la valeur propre $-2$, 
$$E_{-2}=\{\vec u=(x,y,z),\ A.\vec u=-2\vec u\}.$$
Ainsi
$$\vec u=(x,y,z)\in E_{-2}\iff\begin{align*}-&2x-2y-2z=0 \\ &2x+2y+2z=0 \\ &3x+3y+3z=0\end{align*}\iff
x+y+z=0$$
Le sous-espace propre $E_{-2}$ est donc le plan vectoriel d'équation $x+y+z=0$, dont une base est donnée, par exemple, par $\vec e_2=(1,-1,0)$ et $\vec e_3=(1,0,-1)$.
 \item {\it D\'emontrons que $A$ est diagonalisable et donnons une base de $\R^3$ dans laquelle la matrice de $u$ est diagonale.}

Les sous-espaces propres associés aux valeurs propres sont de dimension la multiplicité de la valeur propre correspondante, ce qui prouve que la matrice $A$ est diagonalisable. Dans la base $(\vec e_1,\vec e_2,\vec e_3)$la matrice de l'endomorphisme associé à $A$ est diagonale, elle s'écrit
$$D=\begin{pmatrix}1&0&0 \\  0&-2&0 \\  0&0&-2\end{pmatrix}.$$ 
 \item {\it Trouvons une matrice $P$ telle que $P^{-1}AP$ soit diagonale.}

La matrice de changement de base qui exprime la base $(\vec e_1,\vec e_2,\vec e_3)$ des vecteurs propres, trouvés ci-dessus, dans la base canonique est la matrice $P$ cherchée
$$P=\begin{pmatrix}-2&1&1 \\  2&-1&0 \\  3&0&-1\end{pmatrix}\ \ {\hbox{et}}\ \ P^{-1}={\frac{1}{3}}\begin{pmatrix}1&1&1 \\  2&-1&2 \\  3&3&0\end{pmatrix},$$
elle est inversible et on a $P^{-1}AP=D$. (Le calcul de $P^{-1}$ n'était pas demandé, ni nécessaire). 
\end{enumerate}
\fincorrection
\correction{002606}
{\it Soit $u$ l'endomorphisme de $\R^3$, dont la matrice dans la base canonique est}
$$A=\begin{pmatrix}3&2&-2 \\  -1&0&1 \\ 1&1&0 \\ \end{pmatrix}.$$ 
\begin{enumerate}
 \item  {\it Calculons les valeurs propres de $A$ et voyons si l'endomorphisme $u$ est diagonalisable.} 

En opérant sur les colonnes et les lignes du déterminant, on obtient
$$P_A(X)=\begin{vmatrix}3-X&2&-2 \\  -1&-X&1 \\  1&1&-X\end{vmatrix}=\begin{vmatrix}3-X&0&-2 \\  -1&1-X&1 \\  1&1-X&-X\end{vmatrix},$$ d'où
$$P_A(X)=\begin{vmatrix}3-X&0&-2 \\  -1&1-X&1 \\  2&0&-X-1\end{vmatrix}$$
et, en développant par rapport à la deuxième colonne
$$P_A(X)=(1-X)[(3-X)(-1-X)+4]=(1-X)(X^2-2X+1)=(1-X)^3.$$
Ainsi, la matrice $A$ admet $1$ comme valeur propre triple. Elle n'est donc pas diagonalisable, sinon elle serait égale à $I=I_3$, la matrice identité.
 \item {\it Calculons $(A-I)^2$ et démontrons que $A^n=nA+(1-n)I$.}
 
On calcule d'abord la matrice $A-I$,
$$A-I=\begin{pmatrix}2&2&-2 \\  -1&-1&1 \\ 1&1&-1 \\ \end{pmatrix},$$
puis la matrice $(A-I)^2$,
$$(A-I)^2=\begin{pmatrix}2&2&-2 \\  -1&-1&1 \\ 1&1&-1 \\ \end{pmatrix}
\begin{pmatrix}2&2&-2 \\  -1&-1&1 \\ 1&1&-1 \\ \end{pmatrix}=
\begin{pmatrix}0&0&0 \\  0&0&0 \\ 0&0&0 \\ \end{pmatrix},$$ 
c'est donc la matrice nulle.

Nous allons donner deux méthodes pour démontrer que $A^n=nA+(1-n)I$.

{\it Première méthode} :  En utilisant le bin\^ome de Newton. On écrit $A^n=(A-I+I)^n$, or, les matrices $A-I$ et $I$ commutent, on a donc
$$(A-I+I)^n=\sum_{k=0}^nC_n^k(A-I)^kI^{(n-k)}=C_n^0I+C_n^1(A-I)=I+n(A-I)=nA+(1-n)I.$$
{\it Deuxième méthode} : Par récurrence sur $n$. Le résultat est vrai pour $n=0$ et $n=1$. Fixons $n$ arbitrairement pour lequel on suppose que $A^n=nA+(1-n)I$, on a alors
$$A^{n+1}=A(nA+(1-n)I)=nA^2+(1-n)A,$$
sachant que $(A-I)^2=0$, on en déduit que $A^2=2A-I$ ainsi
$$A^{n+1}=n(2A-I)+(1-n)A=(n+1)A-nI=(n+1)A+(1-(n+1))I.$$
L'égalité est donc vraie pour tout $n\in\N$.
\end{enumerate}
\fincorrection
\correction{002607}
{\it Soit $A$ la matrice 
$$A=\begin{pmatrix}1&0&0 \\  -1&2&1 \\  0&0&2\end{pmatrix}$$ 
et $f$ l'endomorphisme de $\R^3$ associ\'e.}
\begin{enumerate}
 \item  {\it Déterminons les valeurs propres de $A$.}

Calculons les racines du polynôme caractéristique $P_A(X)$ :
$$P_A(X)=\begin{vmatrix}1-X&0&0 \\  -1&2-X&1 \\  0&0&2-X\end{vmatrix}=(2-X)^2(1-X).$$
Les valeurs propres de $A$ sont $\lambda_1=1$, valeur propre simple et $\lambda_2=2$, valeur propre double.

 \item {\it Déterminons, sans calculs, des vecteurs $\vec u$ et $\vec v$ tels que $f(\vec u)=2\vec u$ et $f(\vec v)=2\vec v+\vec u$.}

Si l'on note $(\vec e_1,\vec e_2, \vec e_3)$, la base dans laquelle est exprimée la matrice $A$ de l'endomorphisme $f$, on remarque que 
$$f(\vec e_2)=2\vec e_2\quad{\hbox{et}}\quad f(\vec e_3)=\vec e_2+2\vec e_3.$$
Ainsi, les vecteurs $\vec u=\vec e_2$ et $\vec v=\vec e_3$ répondent-ils à la question.
 \item  {\it Soit $\vec e$ tel que $f(\vec e)=\vec e$. Démontrons que $(\vec e,\vec u,\vec v)$ est une base de $\R^3$ et écrivons la matrice de $f$ dans cette base.}

Notons $\vec e=(x,y,z)$ alors
$$f(\vec e)=\vec e\iff\left\{\begin{align*}x&=x \\  -x+2y+z&=y \\  2z&=z\end{align*}\right.
\iff\left\{\begin{align*}z&=0 \\   x&=y\end{align*}\right.$$
Le vecteur $\vec e=(1,1,0)$ convient. Les vecteurs $\vec e$, $\vec u$ et $\vec v$ sont linéairement indépendants, ils forment donc une base de $\R^3$. La matrice de $f$ dans cette base s'écrit
$$B=\begin{pmatrix}1&0&0 \\  0&2&1 \\  0&0&2\end{pmatrix}$$

 \item  {\it La matrice $A$ est-elle diagonalisable ?} 

Le sous-espace propre associé à la valeur propre $2$ est l'ensemble des vecteurs $(x,y,z)$ tels que 
$$\left\{\begin{align*}x&=2x \\  -x+2y+z&=2y \\  2z&=2z\end{align*}\right.\iff\left\{\begin{align*}x&=0 \\   z&=0\end{align*}\right.$$
C'est une droite vectorielle, sa dimension n'est donc pas égale à la multiplicité de la valeur propre $2$ comme racine du polynôme caractéristique, la matrice $A$ n'est pas diagonalisable.
\end{enumerate}
\fincorrection
\correction{002608}
{\it Soit $A$ la matrice 
$$A=\begin{pmatrix}1&0&0 \\  1&-1&0 \\  -1&2&-1\end{pmatrix}$$ 
et $f$ l'endomorphisme de $\R^3$ associ\'e.}
\begin{enumerate}
 \item {\it  Déterminons et factorisons le polyn\^ome caract\'eristique de $A$.} 

On note $P_A(X)$ le polynôme caractéristique, on a
$$P_A(X)=\begin{vmatrix}1-X&0&0 \\  1&-1-X&0 \\  -1&2&-1-X\end{vmatrix}=(-1-X)\begin{vmatrix}1-X&0 \\  1&-1-X\end{vmatrix}=(-1-X)^2(1-X).$$
La matrice $A$ admet deux valeurs propres, $1$, valeur propre simple, et $-1$, valeur propre double.

 \item {\it D\'eterminons les sous-espaces propres et les sous-espaces caract\'eristiques de $A$.}

La valeur propre $1$ est simple, le sous-espace propre associé est égal au sous-espace caractéristique, c'est l'ensemble
$$E_1=N_1=\{\vec u=(x,y,z)\in\R^3,\ A\vec u =\vec u\}.$$
On a 
$$\vec u\in E_1\iff \left\{\begin{align*}x&=x \\  x-y&=y \\  -x+2y-z&=z\end{align*}\right.\iff\left\{\begin{align*}x&=2y \\  z&=0\end{align*}\right.$$ 
L'espace $E_1$ est une droite vectorielle dont un vecteur directeur $\vec e_1$ est donné, par exemple, par $\vec e_1=(2,1,0)$.

Le sous-espace propre associé à la valeur propre $-1$ est défini par
$$E_{-1}=\{\vec u=(x,y,z)\in\R^3,\ A\vec u =-\vec u\}.$$
On a 
$$\vec u\in E_{-1}\iff \left\{\begin{align*}x&=-x \\  x-y&=-y \\  -x+2y-z&=-z\end{align*}\right.
\iff\left\{ \begin{align*}x&=0 \\  y&=0\end{align*}\right.$$ 
L'espace $E_{-1}$ est une droite vectorielle dont un vecteur directeur $\vec e_2$ est donné, par exemple, par $\vec e_2=(0,0,1)$. La dimension de $E_{-1}$ n'est pas égale à la multiplicité de la racine, la matrice n'est pas diagonalisable. Déterminons le sous-espace caractéristique associé à la valeur propre $-1$. Pour cela calculons la matrice $(A+I)^2$.
$$(A+I)^2=\begin{pmatrix}2&0&0 \\  1&0&0 \\  -1&2&0\end{pmatrix}^2=\begin{pmatrix}4&0&0 \\  2&0&0 \\  0&0&0\end{pmatrix}$$
$$N_{-1}=\ker(A+I)^2=\{\vec u=(x,y,z)\in\R^3,\ x=0\}$$
Le sous-espace caractéristique $N_{-1}$ est le plan vectoriel engendré par les vecteurs 

$e_2=(0,0,1)$ et $e_3=(0,1,0)$.
 \item  {\it D\'emontrons qu'il existe une base de $\R^3$ dans laquelle la matrice de $f$ est 
$$B=\begin{pmatrix}1 & 0 & 0 \\  0&-1&2 \\ 0&0&-1\end{pmatrix}$$
et trouvons une matrice $P$ inversible telle que $ AP=PB$ (ou $A=PBP^{-1}$).} 

On considère les vecteurs $\vec e_1=(2,1,0)$ et $\vec e_2=(0,0,1)$ et on cherche un vecteur $\vec e\in N_{-1}$ tel que 
$f(\vec e)=2\vec e_2-\vec e$. Notons $\vec e=(0,y,z)$,
$$f(\vec e)=2\vec e_2-\vec e\iff y=1,$$
le vecteur $\vec e=\vec e_3=(0,1,0)$ convient, on pouvait le voir directement sur la deuxième colonne de la matrice $A$.
Ainsi, dans la base $(\vec e_1,\vec e_2,\vec e_3)$ avec $\vec e_1=(2,1,0)$, $\vec e_2=(0,0,1)$ et $\vec e_3=(0,1,0)$, la matrice de $f$ s'écrit
$$B=\begin{pmatrix}1 & 0 & 0 \\  0&-1&2 \\ 0&0&-1\end{pmatrix}.$$
La matrice $P$ cherchée est la matrice de passage qui exprime la base $(\vec e_1,\vec e_2,\vec e_3)$ dans la base canonique $(\vec i,\vec j,\vec k)$. On a
$$P=\begin{pmatrix}2 & 0 & 0 \\  1&0&1 \\ 0&1&0\end{pmatrix}$$
et $AP=PB$ ou $A=PBP^{-1}$. On peut calculer $P^{-1}$, c'est la matrice qui exprime les vecteurs $\vec i,\vec j$ et $\vec k$ dans la base 
$(\vec e_1,\vec e_2,\vec e_3)$,
$$P^{-1}=\begin{pmatrix}1/2 & 0 & 0 \\  0&0&1 \\ -1/2&1&0\end{pmatrix}.$$

 \item  {\it Ecrivons la d\'ecomposition de Dunford de $B$.}  

On a
$$B=\begin{pmatrix}1 & 0 & 0 \\  0&-1&2 \\ 0&0&-1\end{pmatrix}=\underbrace{\begin{pmatrix}1 & 0 & 0 \\  0&-1&0 \\ 0&0&-1\end{pmatrix}}_D+\underbrace{\begin{pmatrix}0 & 0 & 0 \\  0&0&2 \\ 0&0&0\end{pmatrix}}_N.$$
Il est clair que la matrice $D$ est diagonalisable puisque diagonale, on vérifie facilement que $N^2=0$, c'est-à-dire que la matrice $N$ est nilpotente et que les deux matrices commutent, $DN=ND$. Ainsi la décomposition $B=D+N$ est bien la décomposition de Dunford de la matrice $B$.
 \item  {\it Pour $t\in\R$, calculons $\exp tB$}.

 On utilise la décomposition de Dunford de la matrice $tB$, $tB=tD+tN$, on a donc 
 $$\exp(tB)=\exp(tD+tN)=\exp(tD).\exp(tN)$$
car les matrices commutent, par ailleurs, comme $D$ est diagonale, on a
$$\exp(tD)=\begin{pmatrix}e^t & 0 & 0 \\  0&e^{-t}&0 \\ 0&0&e^{-t}\end{pmatrix},$$ 
et comme $N^2=0$, on a
$$\exp(tN)=I+tN=\begin{pmatrix}1 & 0 & 0 \\  0&1&2t \\ 0&0&1\end{pmatrix}.$$
D'où
$$\exp(tB)=\begin{pmatrix}e^t & 0 & 0 \\  0&e^{-t}&0 \\ 0&0&e^{-t}\end{pmatrix}\begin{pmatrix}1 & 0 & 0 \\  0&1&2t \\ 0&0&1\end{pmatrix}=\begin{pmatrix}e^t & 0 & 0 \\  0&e^{-t}&2te^{-t} \\ 0&0&e^{-t}\end{pmatrix}.$$
 \item {\it Donnons les solutions des syst\`emes diff\'erentiels $y'=By$ et $x'=Ax$, où $x$ et $y$ désignent des fonctions réelles à valeurs dans $\R^3$.} 

Les solutions du système différentiel $y'(t)=B.y(t)$ sont les fonctions $y(t)=\exp(tB).V$ où $V$ est un vecteur quelconque de $\R^3$. Donc
$$y(t)=\begin{pmatrix}e^t & 0 & 0 \\  0&e^{-t}&2te^{-t} \\ 0&0&e^{-t}\end{pmatrix}\begin{pmatrix}a \\  b \\  c\end{pmatrix}=\begin{pmatrix}ae^t  \\  be^{-t}+2cte^{-t} \\  ce^{-t}\end{pmatrix},$$
$(a,b,c)\in\R^3$.

Pour trouver les solutions du système différentiel $x'(t)=A.x(t)$, on utilise le fait suivant
$$P.y\quad{\hbox{est solution de}}\quad x'=A.x\iff y\quad{\hbox{est solution de}}\quad y'=(P^{-1}AP).y=B.y.$$
Ainsi, les solutions du système $x'=A.x$ s'écrivent
$$x(t)=P.\exp(tB).V$$
où $V$ est un vecteur quelconque de $\R^3$. On remarque qu'il n'est pas utile de calculer la martrice $P^{-1}$. C'est-à-dire
$$x(t)=\begin{pmatrix}2 & 0 & 0 \\  1&0&1 \\ 0&1&0\end{pmatrix}\begin{pmatrix}ae^t  \\  be^{-t}+2cte^{-t} \\  ce^{-t}\end{pmatrix}=\begin{pmatrix}2ae^t  \\  ae^{t}+ce^{-t} \\  be^{-t}+2cte^{-t}\end{pmatrix},$$
$(a,b,c)\in\R^3$.
\end{enumerate}
\fincorrection
\correction{002609}
 {\it Soit $a\in\R$ et $A$ la matrice
$$A=\begin{pmatrix}0&1&0 \\  0&a&0 \\  0&a-2&2\end{pmatrix}.$$}
\begin{enumerate}
 \item {\it Pour quelles valeurs de $a$ la matrice $A$ est-elle diagonalisable ?}

Déterminons le polynôme caractéristique de la matrice $A$.
$$P_A(X)=\det (A-XI)=\begin{vmatrix}-X&1&0 \\  0&a-X&0 \\ 0&a-2&2-X\end{vmatrix}=-X(a-X)(2-X).$$
Ce polynôme admet trois racines $0,a$ et $2$. Ainsi, si $a\notin\{0,2\}$ la matrice est diagonalisable. Examinons les cas $a=0$ et $a=2$.

Si $a=0$, la valeur propre $0$ est valeur propre double, on a
$$A=\begin{pmatrix}0&1&0 \\  0&0&0 \\  0&-2&2\end{pmatrix}.$$
Le sous-espace propre associé à $0$ est égal à $\ker A=\{\vec u=(x,y,z),\ A\vec u=\vec 0\}$,
$$\vec u\in\ker A\iff\left\{\begin{align*}y&=0 \\ -2y+2z&=0\end{align*}\right.\iff y=z=0.$$
Le sous-espace propre associé à la valeur propre double $0$ est une droite vectorielle, la droite engendrée par $(1,0,0)$, la matrice n'est donc pas diagonalisable.

Si $a=2$, la valeur propre $2$ est double, on a
$$A=\begin{pmatrix}0&1&0 \\  0&2&0 \\  0&0&2\end{pmatrix}.$$
Le sous-espace propre associé à $2$ est égal à $E_2=\{\vec u=(x,y,z),\ A\vec u=2\vec u\}$,
$$\vec u\in E_2\iff\left\{\begin{align*}y&=2x \\ 2y&=2y \\ 2z&=2z\end{align*}\right.\iff y=2x.$$
Le sous-espace propre associé à la valeur propre double $2$ est un plan vectoriel, le plan d'équation $y=2x$, la matrice est donc diagonalisable.

Ainsi la matrice $A$ est diagonalisable si et seulement si $a\neq 0$.

{\it Lorsque $A$ est diagonalisable, déterminons une base de vecteurs propres de $A$.}

On a $a\neq 0$ et on distingue les cas $a\neq 2$ et $a=2$.

Si $a\neq 2$, les sous-espaces propres associés aux valeurs propres $0$ et $2$ sont lisibles sur la matrice, on a
$$E_0=\ker A=\R.(1,0,0)\quad{\hbox{et}}\quad E_2=\R.(0,0,1),$$
On détermine $E_a=\{\vec u=(x,y,z),\ A\vec u=a\vec u\}$.
$$\vec u\in E_a\iff\left\{\begin{align*}y&=ax \\  ay&=ay \\  (a-2)y+2z&=az\end{align*}\right.\iff\left\{\begin{align*}y&=ax \\  (a-2)y&=(a-2)z\end{align*}\right.\iff\left\{\begin{align*}y&=ax \\  y&=z\end{align*}\right.$$
C'est la droite vectorielle engendrée par le vecteur $\vec e=(1,a,a)$. Ainsi, une base de vecteurs propres est donnée par les vecteurs $(1,0,0), (0,0,1)$ et $(1,a,a)$.

Si $a=2$, nous avons vu que le sous-espace propre associé à la valeur propre double $2$ est le plan d'équation $y=2x$. 
Ainsi, une base de vecteurs propres est donnée par les vecteurs $(1,0,0), (0,0,1)$ et $(1,2,0)$.

 \item {\it Soit $E$ l'espace vectoriel des solutions du système $x'=Ax$, où $x$ est une fonction de la variable réelle $t$ à valeur dans $\R^3$.}
   \begin{enumerate}
     \item {\it Lorsque $A$ est diagonalisable, donnons une base de $E$ en fonction des vecteurs propres et des valeurs propres de $A$ et écrivons la solution générale du système.}

Si $\lambda_1, \lambda_2, \lambda_3$ sont les valeurs propres de $A$ et $\vec e_1,\vec e_2$ et $\vec e_3$ les vecteurs propres associés, on sait qu'une base de l'espace des solutions du système différentiel $x'=Ax$ est donnée par 
$$e^{\lambda_1 t}\vec e_1,\ e^{\lambda_2 t}\vec e_2,\ e^{\lambda_3 t}\vec e_3.$$
Ainsi, si $a\neq 2$ cette base est donnée par
$$(1,0,0),\ e^{2t}(0,0,1),\ e^{at}(1,a,a)$$
et si $a=2$, elle est donnée par
$$(1,0,0),\ e^{2t}(0,0,1),\ e^{2t}(1,2,0).$$
     \item  {\it Lorsque $A$ n'est pas diagonalisable, intégrons directement le système $X'=AX$}.

Lorsque $A$ n'est pas diagonalisable, $a=0$ et 
$$A=\begin{pmatrix}0&1&0 \\  0&0&0 \\  0&-2&2\end{pmatrix}.$$
Le système $X'=AX$ est équivalent à
$$\begin{pmatrix}x'(t) \\  y'(t) \\  z'(t)\end{pmatrix}=\begin{pmatrix}0&1&0 \\  0&0&0 \\  0&-2&2\end{pmatrix}
\begin{pmatrix}x(t) \\  y(t) \\  z(t)\end{pmatrix}\iff
\left\{\begin{align*}x'&=y \\  y'&=0 \\  z'&=-2y+2z\end{align*}\right.$$
Si $y'=0$, alors $y(t)=\alpha$, $\alpha\in\R$. Ainsi, si $x'=\alpha$, $x(t)=\alpha t+\beta$, $\beta\in\R$ et la troisième équation devient
$$z'=2z-2\alpha$$
et sa solution s'écrit $z(t)=\gamma e^{2t}+\alpha$, $\gamma\in\R$. D'où la solution générale du système
$$X(t)=\begin{pmatrix}x(t) \\  y(t) \\  z(t)\end{pmatrix}=\begin{pmatrix}\alpha t+\beta \\  \alpha \\  \gamma e^{2t}+\alpha\end{pmatrix}=\alpha\begin{pmatrix}t \\ 1 \\ 1\end{pmatrix}+\beta\begin{pmatrix}1 \\ 0 \\ 0\end{pmatrix}+\gamma\begin{pmatrix}0 \\ 0 \\  e^{2t}\end{pmatrix}$$
$(\alpha,\beta,\gamma)\in\R^3$.
    \end{enumerate}
 \item {\it Soit $E_0$ l'ensemble des éléments $s$ de $E$ tels que $\lim_{t\rightarrow+\infty}s(t)=\vec0$. Démontrons que $E_0$ est un sous-espace vectoriel de $E$}.

Pour démontrer que $E_0$ est un sous-espace vectoriel de $E$, il suffit de démontrer que $0_E\in E_0$ et que $E_0$ est stable par combinaison linéaire. La fonction nulle est clairement dans $E_0$, par ailleurs, si $s_1$ et $s_2$ sont dans $E_0$ et si $a_1,a_2\in\R$, on a
$$\lim_{t\rightarrow+\infty}a_1s_1(t)+a_2s_2(t)=\vec0$$
ce qui prouve que $a_1s_1+a_2s_2\in E_0$.

{\it Déterminons sa dimension en fonction de $a$.}

Nous avons, dans tous les cas, une base de l'espace des solutions, donc l'écriture de la solution générale, on regarde alors les solutions de $E$ qui sont dans $E_0$.

{\it 1er cas} : $a\neq 2$ et $a\neq 0$, la solution générale s'écrit
$$s(t)=\alpha(1,0,0)+\beta(0,0,e^{2t})+\gamma(e^{at},ae^{at},ae^{at}),\ (\alpha,\beta,\gamma)\in\R^3.$$
Ainsi, si $a>0$,  $s(t)\in E_0\iff s=\vec 0$, $\dim E_0=0$, 

et si $a<0$, $s(t)\in E_0\iff \alpha=\beta=0$, $\dim E_0=1$.

{\it 2ème cas} : $a=2$, la solution générale s'écrit
$$s(t)=\alpha(1,0,0)+\beta(0,0,e^{2t})+\gamma(e^{2t},2e^{2t},0),\ (\alpha,\beta,\gamma)\in\R^3.$$
Ainsi, $s(t)\in E_0\iff s=\vec 0$, $\dim E_0=0$.

{\it 3ème cas} : $a=0$, la solution générale s'écrit
$$s(t)=\alpha(1,0,0)+\beta(0,0,e^{2t})+\gamma(t,1,1),\ (\alpha,\beta,\gamma)\in\R^3.$$
Ainsi, $s(t)\in E_0\iff s=\vec 0$, $\dim E_0=0$.
 \item {\it Soit $F$ l'ensemble des éléments $s$ de $E$ bornés sur $[0,+\infty[$. Démontrons que $F$ est un sous-espace vectoriel de $E$.}

Pour démontrer que $F$ est un sous-espace vectoriel de $E$, il suffit de démontrer que $0_E\in F$ et que $F$ est stable par combinaison linéaire. La fonction nulle est clairement dans $F$, par ailleurs, si $s_1$ et $s_2$ sont dans $E_0$ et si $a_1,a_2\in\R$, la fonction $a_1s_1(t)+a_2s_2(t)$ est bornée sur $[0,+\infty[$, donc dans $F$.

{\it Déterminons sa dimension en fonction de $a$.}

Comme dans le cas précédent, suivant la forme de la solution générale, on a

si $a\geq0$, les seules solutions bornées sur $[0,+\infty[$ sont de la forme 

$s(t)=\alpha(1,0,0),\ \alpha\in R$, ainsi $\dim F=1$.

si $a<0$, les seules solutions bornées sur $[0,+\infty[$ sont de la forme 

$s(t)=\alpha(1,0,0)+\gamma(e^{at},ae^{at},ae^{at})$, $\ (\alpha,\gamma)\in\R^2$, $\dim F=2$.

\end{enumerate}
\fincorrection
\correction{002610}

\centerline{\bf I}

Soit $\alpha\in\R$ et $A_\alpha\in M_3(\R)$ la matrice suivante
$$A_\alpha=\begin{pmatrix}-1 & 0 & \alpha+1 \\  1&-2&0 \\ -1&1&\alpha\end{pmatrix}$$

\begin{enumerate}  
  \item{\it Factorisons le polyn\^ome caract\'eristique $P_{A_\alpha}(X)$ en produit de facteurs du premier degr\'e.}

On a
\begin{align*}P_{A_\alpha}(X)
&=\begin{vmatrix}-1-X & 0 & \alpha+1 \\  1&-2-X&0 \\ -1&1&\alpha-X\end{vmatrix}
=\begin{vmatrix}-1-X & 0 & \alpha+1 \\  -1-X&-2-X&0 \\ 0&1&\alpha-X\end{vmatrix} \\ 
&=\begin{vmatrix}-1-X & 0 & \alpha+1 \\  0&-2-X&-\alpha-1 \\ 0&1&\alpha-X\end{vmatrix} \\ 
&=(-1-X)[(-2-X)(\alpha-X)+\alpha+1] \\ 
&=-(X+1)[X^2+(2-\alpha)X+1-\alpha].
\end{align*}
Factorisons le polyn\^ome $X^2+(2-\alpha)X+1-\alpha$, son discriminant est \'egal \`a
$$\Delta=(2-\alpha)^2-4(1-\alpha)=\alpha^2.$$
On a donc $\sqrt{\Delta}=|\alpha|$, ce qui nous donne les deux racines 
$$\lambda_1={\frac{\alpha-2-\alpha}{2}}=-1\ \ {\hbox{et}}\ \ 
\lambda_2={\frac{\alpha-2+\alpha}{2}}=\alpha-1.$$
Le polyn\^ome caract\'eristique $P_{A_\alpha}(X)$ se factorise donc en
$$P_{A_\alpha}(X)=-(X+1)^2(X-\alpha+1).$$ 

  \item {\it D\'eterminons selon la valeur du param\`etre $\alpha$ les valeurs propres distinctes de $A_\alpha$ et leur multiplicit\'e.}

Les valeurs propres de $A_\alpha$ sont les racines du polyn\^ome caract\'eristique $P_{A_\alpha}$, ainsi,

- si $\alpha=0$, la matrice $A_\alpha$ admet une valeur propre triple $\lambda=-1$,

- si $\alpha\neq0$, la matrice $A_\alpha$ admet deux valeurs propres distinctes
$\lambda_1=-1$ valeur propre double et $\lambda_2=\alpha-1$, valeur propre simple.
 

  \item {\it D\'eterminons les valeurs de $\alpha$ pour lesquelles la matrice $A_\alpha$ est diagonalisable.}

Il est clair que dans le cas $\alpha=0$, la matrice n'est pas diagonalisable, en effet si elle l'\'etait, il existerait une matrice inversible $P$ telle que $A_\alpha=P(-I)P^{-1}=-I$, ce qui n'est pas le cas.

Si $\alpha\neq0$, la matrice $A_\alpha$ est diagonalisable si le sous-espace propre associ\'e \`a la valeur propre $-1$ est de dimension $2$. D\'eterminons ce sous-espace propre.

$$E_{-1}=\ker(A_\alpha+I)=\left\{(x,y,z)\in\R^3,\ \begin{pmatrix}-1 & 0 & \alpha+1 \\  1&-2&0 \\ -1&1&\alpha\end{pmatrix}\begin{pmatrix}x \\  y \\  z\end{pmatrix}= \begin{pmatrix}-x \\  -y \\  -z\end{pmatrix}\right\}$$
ainsi,
$$(x,y,z)\in E_{-1}\iff\left\{\begin{align*}-x+(\alpha+1)z&=-x \\  x-2y&=-y \\  -x+y+\alpha z&=-z\end{align*}\right.
\iff\left\{\begin{align*}(\alpha+1)z&=0 \\  x-y=0\end{align*}\right.$$
Il faut distinguer les cas $\alpha=-1$ et $\alpha\neq -1$.

- Si $\alpha=-1$, le sous-espace $E_{-1}$ est le plan vectoriel d'\'equation $x=y$, dans ce cas la matrice $A_\alpha$ est diagonalisable.

- Si $\alpha\neq -1$, le sous-espace $E_{-1}$ est la droite vectorielle engendr\'ee par le vecteur
$(1,1,0)$, dans ce cas la matrice $A_\alpha$ n'est pas diagonalisable.

  \item {\it D\'eterminons selon la valeur de $\alpha$ le polyn\^ome minimal de $A_\alpha$.}

 Notons $Q_A$ le polyn\^ome minimal de $A_\alpha$. On sait que la matrice $A_\alpha$ est diagonalisable sur $\R$ si et seulement si son polyn\^ome minimal a toutes ses racines dans $\R$ et que celles-ci sont simples. Or, nous venons de d\'emontrer que $A_\alpha$ est diagonalisable sur $\R$ si et seulement $\alpha=-1$, on a donc 

- Si $\alpha=-1$, $A_\alpha$ est diagonalisable, donc $Q_A(X)=(X+1)(X-\alpha+1)=(X+1)(X+2)$.

- Si $\alpha\neq -1$, on doit distinguer les cas $\alpha=0$ et $\alpha\neq0$, en effet, 

\indent - si $\alpha\neq 0$,  $A_\alpha$ n'est pas diagonalisable donc les racines de $Q_A$ ne sont pas toutes les deux simples, et $P_A$ admet deux racines distinctes, donc
$$Q_A(X)=-P_A(X)=(X+1)^2(X-\alpha+1).$$

\indent- si $\alpha=0$, on a $P_A(X)=-(X+1)^3$ et $A_0$ n'est pas diagonalisable, le polyn\^ome minimal peut donc \^etre \'egal \`a $(X+1)^2$ ou \`a $(X+1)^3$, or
$$(A+I)^2=\begin{pmatrix}0 & 0 & 1 \\  1&-1&0 \\ -1&1&1\end{pmatrix}\begin{pmatrix}0 & 0 & 1 \\  1&-1&0 \\ -1&1&1\end{pmatrix}=\begin{pmatrix}-1 & 1 & 1 \\  -1&1&1 \\ 0&0&0\end{pmatrix}\neq0$$
donc le polyn\^ome minimal $Q_A$ est \'egal \`a $(X+1)^3$.
\end{enumerate}


{\centerline{II}}  

On suppose d\'esormais que $\alpha=0$, on note $A=A_0$ et $f$ l'endomorphisme de $\R^3$ associ\'e \`a la matrice $A$. On a donc 
$$A=A_0=\begin{pmatrix}-1 & 0 & 1 \\  1&-2&0 \\ -1&1&0\end{pmatrix}$$
et $P_A(X)=-(X+1)^3$.

\begin{enumerate} 
  \item {\it D\'eterminons les sous-espaces propres et caract\'eristiques de $A$.}

La matrice $A$ admet une unique valeur propre $\lambda=-1$ de multiplicit\'e $3$, le sous-espace propre associ\'e est l'espace $E_{-1}=\ker(A+I)$, et on a
$$(x,y,z)\in E_{-1}\iff\left\{\begin{align*}-x+z&=-x \\  x-2y&=-y \\  -x+y&=-z\end{align*}\right.
\iff\left\{\begin{align*}z&=0 \\  x&=y\end{align*}\right.$$
Le sous-espace $E_{-1}$ est la droite vectorielle engendr\'ee par le vecteur $(1,1,0)$.

Le sous-espace caract\'eristique de $A$, associ\'e \`a l'unique valeur propre $\lambda=-1$, est le sous-espace $N_{-1}=\ker(A+I)^3$, or, compte tenu du th\'eor\`eme de Hamilton-Cayley, on sait que $P_A(A)=0$, ainsi, la matrice $(A+I)^3$ est la matrice nulle, ce qui implique $N_{-1}=\R^3$, c'est donc l'espace tout entier.

  \item {\it D\'emontrons que le sous-espace vectoriel $\ker(A+I)^2$ est un plan stable par $f$.}

On a $E_{-1}=\ker(A+I)\subset\ker(A+I)^2\subset\ker(A+I)^3=\R^3$, le sous-espace 
$\ker(A+I)^2$ est clairement stable par $A$ car  pour tout $v\in\ker(A+I)^2$, $Av\in\ker(A+I)^2$, en effet 
$$(A+I)^2Av=A(A+I)^2v=0.$$
D\'emontrons que ce sous-espace est un plan. On a
$$(A+I)^2=\begin{pmatrix}0 & 0 & 1 \\  1&-1&0 \\ -1&1&1\end{pmatrix}\begin{pmatrix}0 & 0 & 1 \\  1&-1&0 \\ -1&1&1\end{pmatrix}=\begin{pmatrix}-1 & 1 & 1 \\  -1&1&1 \\ 0&0&0\end{pmatrix}$$
donc $\ker(A+I)^2=\{(x,y,z)\in\R^3,\ -x+y+z=0\}$, c'est bien un plan vectoriel.

  \item {\it D\'emontrons qu'il existe une base de $\R^3$ dans laquelle la matrice de $f$ est }
$$B=\begin{pmatrix}-1 & 1 & 0 \\  0&-1&1 \\ 0&0&-1\end{pmatrix}$$
{\it et trouvons une matrice $P$ inversible telle que} $A=PBP^{-1}$.

Nous cherchons des vecteurs $e_1,e_2,e_3$ tels que $Ae_1=e_1$, $Ae_2=e_1-e_2$ et $Ae_3=e_2-e_3$.
Le vecteur $e_1$ appartient \`a $E_1=\ker(A+I)$, et $\ker(A+I)$ est la droite d'\'equations :

$\{z=0,x-y=0\}$ ; nous choisirons $e_2\in \ker(A+I)^2$ tel que $(e_1,e_2)$ soit une base de $\ker(A+I)^2$. Remarquons que si l'on cherche $e_2=(x,y,z)$ tel que
$Ae_2=e_1-e_2$, on obtient le syst\`eme 
$$\begin{pmatrix}-1 & 0 & 1 \\  1&-2&0 \\ -1&1&0\end{pmatrix}\begin{pmatrix}x \\  y \\  z\end{pmatrix}=\begin{pmatrix}1-x \\  1-y \\ -z\end{pmatrix}\iff
\left\{\begin{align*}-x+z&=1-x \\  x-2y&=1-y \\  -x+y&=-z\end{align*}\right.
\iff\left\{\begin{align*}z&=1 \\  x-y&=1\end{align*}\right.$$
ce qui nous donne bien un vecteur de $\ker(A+I)^2$.
Ainsi, les vecteurs $e_1=(1,1,0)$ et $e_2=(1,0,1)$ conviennent. Il nous reste \`a chercher un vecteur $e_3$ tel que $Ae_3=e_2-e_3$, c'est-\`a-dire
$$\begin{pmatrix}-1 & 0 & 1 \\  1&-2&0 \\ -1&1&0\end{pmatrix}\begin{pmatrix}x \\  y \\  z\end{pmatrix}=\begin{pmatrix}1-x \\  -y \\ 1-z\end{pmatrix}\iff
\left\{\begin{align*}-x+z&=1-x \\  x-2y&=-y \\  -x+y&=1-z\end{align*}\right.
\iff\left\{\begin{align*}z&=1 \\  x&=y\end{align*}\right.$$
Le vecteur $e_3=(0,0,1)$ convient. On obtient alors la matrice $P$ suivante qui est inversible et v\'erifie $A=PBP^{-1}$,
$$P=\begin{pmatrix}1 & 1 & 0 \\  1&0&0 \\ 0&1&1\end{pmatrix}$$  

  \item{\it D\'ecomposition de Dunford de $B$}

On a 
$$B=\begin{pmatrix}-1 & 1 & 0 \\  0&-1&1 \\ 0&0&-1\end{pmatrix}=
\begin{pmatrix}-1 & 0& 0 \\  0&-1&0 \\ 0&0&-1\end{pmatrix}+
\begin{pmatrix}0& 1 & 0 \\  0&0&1 \\ 0&0&0\end{pmatrix}$$
et il est clair que les deux matrices commutent car l'une est \'egale à $-I$. Or, il existe un unique couple de matrices $D$ et $N$, $D$ diagonalisable et $N$ nilpotente, telles que $B=D+N$ et $DN=ND$. Or si
$$D=\begin{pmatrix}-1 & 0& 0 \\  0&-1&0 \\ 0&0&-1\end{pmatrix}\ \ {\hbox{et}}\ \ 
N=\begin{pmatrix}0& 1 & 0 \\  0&0&1 \\ 0&0&0\end{pmatrix},$$
On a 
$$N^2=\begin{pmatrix}0& 1 & 0 \\  0&0&1 \\ 0&0&0\end{pmatrix}\begin{pmatrix}0& 1 & 0 \\  0&0&1 \\ 0&0&0\end{pmatrix}=\begin{pmatrix}0& 0 & 1 \\  0&0&0 \\ 0&0&0\end{pmatrix}$$
et $N^3=0$. La d\'ecomposition $B=D+N$ est donc bien la d\'ecomposition de Dunford de la matrice $B$.

  \item{\it Pour $t\in\R$, calculons $\exp tB$ et exprimons $\exp tA$ \`a l'aide de $P$ et $\exp tB$.}

On a $N^3=0$ donc pour tout $t\in\R$, $(tN)^3=0$ et l'exponentielle est \'egale \`a $$\exp(tN)=I+tN+(t^2/2)N^2,$$ par ailleurs $ND=DN$, donc  pour tout $t\in\R$, les matrices $tN$ et $tD$ commutent \'egalement, $(tN)(tD)=(tD)(tN)$, on a donc
$$\exp(tB)=\exp(tD+tN)=\exp(tD)\exp(tN)=\exp(-tI)\exp(tN)=e^{-t}I\left(I+tN+{\frac{t^2}{2}}N^2\right)$$
D'o\`u
$$\exp(tB)=e^{-t}\begin{pmatrix}1 & t & {\frac{t^2}{2}} \\  0&1&t \\ 0&0&1\end{pmatrix}.$$
Pour d\'eterminer l'exponentielle de la matrice $tA$, on \'ecrit
$$\exp(tA)=\exp(t(PBP^{-1}))=\exp(P(tA)P^{-1})=P\exp(tB)P^{-1}.$$
  \item {\it Solutions des syst\`emes diff\'erentiels $Y'=BY$ et $X'=AX$.}

La solution g\'en\'erale du syst\`eme $Y'=BY$ s'\'ecrit 
$$S(t)=\exp(tB)v$$
o\`u $v=(a,b,c)$ est un vecteur de $\R^3$. La solution $S: \R\mapsto \R^3$  s'\'ecrit donc
$$S(t)=\begin{pmatrix}x(t) \\  y(t) \\  z(t)\end{pmatrix}
=e^{-t}\begin{pmatrix}1 & t & {\frac{t^2}{2}} \\  0&1&t \\ 0&0&1\end{pmatrix}\begin{pmatrix}a \\  b \\  c\end{pmatrix}
=e^{-t}\begin{pmatrix}a+bt+c{\frac{t^2}{2}} \\  b+ct \\  c\end{pmatrix}$$
Pour obtenir la solution du syst\`eme $X'=AX$, on \'ecrit
$$X'=AX\iff X'=(PBP^{-1})X\iff P^{-1}X'=(BP^{-1})X\iff (P^{-1}X)'=B(P^{-1}X)$$
ainsi, en notant $Y=P^{-1}X$ ou encore $X=PY$, les solutions du syst\`eme $X'=AX$ sont les $PS(t)$ o\`u $P$ est la matrice v\'erifiant
$A=PBP^{-1}$ et $S$ une solution du syst\`eme $Y'=BY$.

La solution g\'en\'erale du syst\`eme $X'=AX$ s'\'ecrit donc 
$$X(t)=\begin{pmatrix}x(t) \\  y(t) \\  z(t)\end{pmatrix}
=\begin{pmatrix}1 & 1 & 0 \\  1&0&0 \\ 0&1&1\end{pmatrix}
\begin{pmatrix}e^{-t}(a+bt+c{\frac{t^2}{2}}) \\  e^{-t}(b+ct) \\  e^{-t}c\end{pmatrix}
=e^{-t}\begin{pmatrix}(a+b)+(c+b)t+c{\frac{t^2}{2}} \\  a+bt+c{\frac{t^2}{2}} \\  (b+c)+ct\end{pmatrix}$$
o\`u $(a,b,c)\in\R^3$. 

\end{enumerate}


\centerline{III}  

On suppose, dans cette partie, que $\alpha=-1$, on note $A=A_{-1}$.
$$A=\begin{pmatrix}-1 & 0 & 0 \\  1&-2&0 \\ -1&1&-1\end{pmatrix}$$

\begin{enumerate}  
\item{\it  V\'erifions que la matrice $A$ est diagonalisable.}

Nous avons vu dans la partie $I) 3)$ que lorsque $\alpha=-1$, la matrice est diagonalisable, en effet, dans ce cas, elle admet deux valeurs propres : $-1$, valeur propre double et $-2$, valeur propre simple. Le sous-espace propre associé à la valeur propre $-1$ étant le plan d'équation $x=y$.
 
  \item {\it Diagonalisons la matrice $A$.}

Pour cela déterminons une base de vecteurs propres. Le plan $x=y$ est engendré par les vecteurs $u(1,1,0)$ et $v(0,0,1)$, déterminons un vecteur directeur de la droite $E_{-2}$ :
$$E_{-2}=\ker(A_\alpha+2I)=\left\{(x,y,z)\in\R^3,\ \begin{pmatrix}-1 & 0 & 0 \\  1&-2&0 \\ -1&1&-1\end{pmatrix}
\begin{pmatrix}x \\  y \\  z\end{pmatrix}=
\begin{pmatrix}-2x \\  -2y \\  -2z\end{pmatrix}\right\}$$
ainsi,
$$(x,y,z)\in E_{-2}\iff\left\{\begin{align*}-x&=-2x \\  x-2y&=-2y \\  -x+y-z&=-2z\end{align*}\right.
\iff\left\{\begin{align*}x&=0 \\  y&=-z\end{align*}\right.$$
Le sous-espace propre $E_{-2}$ est la droite vectorielle engendrée par le vecteur $w(0,1,-1)$. Ainsi, dans la base $(u,v,w)$ la matrice est diagonale et s'écrit
$$D=\begin{pmatrix}-1 & 0 & 0 \\  0&-1&0 \\ 0&0&-2\end{pmatrix},$$ 
on a $A=PDP^{-1}$, où 
$$P=\begin{pmatrix}1 & 0 & 0 \\  1&0&1 \\ 0&1&-1\end{pmatrix}$$

  \item {\it Donnons les solutions du syst\`eme diff\'erentiel $X'=A.X$.}

 Si on note $X=PY$, les solutions du syst\`eme $X'=AX$ sont les $PS(t)$ o\`u $S$ une solution du syst\`eme $Y'=DY$. Ainsi, la solution générale du système $X'=A.X$ s'écrit
 $$X(t)=\begin{pmatrix}x(t) \\  y(t) \\  z(t)\end{pmatrix}
=\begin{pmatrix}1 & 0 & 0 \\  1&0&1 \\ 0&1&-1\end{pmatrix}
\begin{pmatrix}ae^{-t} \\  be^{-t} \\  ce^{-2t}\end{pmatrix}
=\begin{pmatrix}ae^{-t} \\  ae^{-t}+be^{-t} \\  be^{-t}+ce^{-2t}\end{pmatrix}$$
o\`u $(a,b,c)\in\R^3$. 
\end{enumerate}


\centerline{IV}

On suppose, dans cette partie, que $\alpha=1$, on note $A=A_{1}$.
$$A=\begin{pmatrix}-1 & 0 & 2 \\  1&-2&0 \\ -1&1&1\end{pmatrix}$$

\begin{enumerate}  
\item {\it D\'eterminons les sous-espaces propres et caract\'eristiques de $A$.}

La matrice $A$ admet deux valeurs propres : $-1$, valeur propre double et $0$, valeur propre simple. Le sous-espace propre $E_{-1}$, associé à la valeur propre $-1$, est la droite vectorielle engendr\'ee par le vecteur
$(1,1,0)$. Déterminons le sous-espace $E_0$ :
$$E_{0}=\ker(A)=\left\{(x,y,z)\in\R^3,\ \begin{pmatrix}-1 & 0 & 2 \\  1&-2&0 \\ -1&1&1\end{pmatrix}\begin{pmatrix}x \\  y \\  z\end{pmatrix}=\begin{pmatrix}0 \\  0 \\  0\end{pmatrix}\right\}$$
ainsi,
$$(x,y,z)\in E_{0}\iff\left\{\begin{align*}-x+2z&=0 \\  x-2y&=0 \\  -x+y+z&=0\end{align*}\right.
\iff\left\{\begin{align*}x&=2z \\  x&=2y\end{align*}\right.$$
Le sous-espace propre $E_0$ est la droite vectorielle engendrée par $(2,1,1)$. Le sous-espace caractéristique associé à la valeur propre $0$ est le sous-espace propre $E_0$. Déterminons le sous-espace caractéristique associé à la valeur propre double $-1$ c'est l'espace vectoriel $\ker(A+I)^2$. On a 
$$(A+I)^2=\begin{pmatrix}0 & 0 & 2 \\  1&-1&0 \\ -1&1&2\end{pmatrix}^2=\begin{pmatrix}-2 & 2 & 4 \\  -1&1&2 \\ -1&1&2\end{pmatrix}.$$
Ainsi, $\ker(A+I)^2$ est le plan vectoriel d'équation $-x+y+2z=0$, engendré par les vecteurs $(1,1,0)\in E_{-1}$ et $(2,0,1)$. 



  \item {\it Trigonalisons la matrice $A$.}

Dans la base $(u,v,w)$ où $u(2,1,1)$, $v(1,1,0)$ et $w(2,0,1)$ la matrice est triangulaire, 
il existe $\lambda$ tel que $A.w=\lambda v-w.$
$$A.w=\begin{pmatrix}-1 & 0 & 2 \\  1&-2&0 \\ -1&1&1\end{pmatrix}.
\begin{pmatrix}2 \\  0 \\  1\end{pmatrix}=
\begin{pmatrix}0 \\  2 \\  -1\end{pmatrix}=
2\begin{pmatrix}1 \\  1 \\  0\end{pmatrix}-\begin{pmatrix}0 \\  2 \\  1\end{pmatrix}$$
Ainsi $A=PTP^{-1}$ où 
$$T=\begin{pmatrix}0 & 0 & 0 \\  0&-1&2 \\ 0&0&-1\end{pmatrix}\ {\hbox{et}}\ P=
\begin{pmatrix}2 & 1 & 2 \\  1&1&0 \\ 1&0&1\end{pmatrix}.$$

\end{enumerate}


\fincorrection
\correction{002611}
 Soit $A$ une matrice $2\times 2$ \`a coefficients r\'eels.
$$A=\begin{pmatrix}a&b \\  c&d\end{pmatrix}$$
On suppose $a+c=b+d=1$ et $a-b\neq 1$. 
\begin{enumerate}
 \item Soient $(x_1,x_2)$, $(y_1,y_2)$ deux vecteurs de $\R^2$, tels que  
$$A\begin{pmatrix}x_1 \\  x_2\end{pmatrix}=\begin{pmatrix}y_1 \\  y_2\end{pmatrix}$$ 
{\it On montre que $y_1+y_2=x_1+x_2.$}

 On a
$$A\begin{pmatrix}x_1 \\  x_2\end{pmatrix}=\begin{pmatrix}a&b \\  c&d\end{pmatrix}\begin{pmatrix}x_1 \\  x_2\end{pmatrix}=\begin{pmatrix}ax_1+bx_2 \\  cx_1+dx_2\end{pmatrix}=\begin{pmatrix}y_1 \\  y_2\end{pmatrix},$$
d'où $y_1+y_2=ax_1+bx_2+cx_1+dx_2=(a+c)x_1+(b+d)x_2=x_1+x_2.$
 \item {\it Soit le vecteur $\vec x=(1,-1)$, vérifions que $\vec x$ est un vecteur propre de $A$, et déterminons sa valeur propre.}
$$A.\vec x=\begin{pmatrix}a&b \\  c&d\end{pmatrix}\begin{pmatrix}1 \\  -1\end{pmatrix}=\begin{pmatrix}a-b \\  c-d\end{pmatrix},$$
or $c-d=(1-a)-(1-b)=-(a-b)$, car $a+b=c+d=1$. Ainsi,
$$\begin{pmatrix}a&b \\  c&d\end{pmatrix}\begin{pmatrix}1 \\  -1\end{pmatrix}=
\begin{pmatrix}a-b \\  -(a-b)\end{pmatrix}=(a-b)\begin{pmatrix}1 \\  -1\end{pmatrix}.$$
Ainsi, le vecteur $\vec x$ est un vecteur propre de $A$ pour la valeur propre $a-b$.
 \item  {\it Déterminons le polynôme caractéristique de $A$ et calculons ses racines.}

Tout d'abord, compte tenu de l'hypothèse
$a+b=c+d=1$, nous écrirons
$$A=\begin{pmatrix}a&b \\  1-a&1-b\end{pmatrix}.$$
D'où 
$$P_A(X)=\begin{vmatrix}a-X&b \\  1-a&1-b-X\end{vmatrix}=(a-X)(1-b-X)-b(1-a)=X^2-(a-b+1)X+(a-b).$$
On sait, d'après la question précédente que $a-b$ est racine de ce polyn\^ome, or, le produit des racines est égal à $a-b$ et la somme à $a-b+1$, ainsi la seconde racine est égale é $1$.

 \item {\it Déterminons un vecteur propre, $\vec y$, de $A$ non colinéaire à $\vec x$ et exprimons la matrice de l'endomorphisme défini par $A$ dans la base $(\vec x,\vec y)$.}

Un vecteur propre non colinéaire à $\vec x$ est vecteur propre pour la valeur propre $1$. Ainsi, si on note $\vec y=(y_1,y_2)$, on a
$$A\vec y=\vec y\iff\begin{pmatrix}a&b \\  1-a&1-b\end{pmatrix}\begin{pmatrix}y_1 \\  y_2\end{pmatrix}=
\begin{pmatrix}y_1 \\  y_2\end{pmatrix},$$
ce qui équivaut à
$$\left\{\begin{align*}ay_1+by_2&=y_1 \\  (1-a)y_1+(1-b)y_2&=y_2\end{align*}\right.\iff (a-1)y_1+by_2=0.$$
Le vecteur $\vec y=(b,1-a)$ est un vecteur propre de $A$ pour la valeur propre $1$.
\end{enumerate}
\fincorrection
\correction{002612}
Soit $E$ un espace vectoriel de dimension $3$. On note ${\cal B}=(\vec e_1,\vec e_2,\vec e_3)$ une base de $E$, si $\vec u$ est un vecteur de $E$ on note $(x,y,z)$  ses coordonnées dans la base $\cal B$. Soit 
$f$ une application linéaire de $E$ dans lui-même, définie par
\begin{align*}
f:E &\longrightarrow E \\  \begin{pmatrix}x \\  y \\  z\end{pmatrix}&\longmapsto \begin{pmatrix}x' \\  y' \\  z'\end{pmatrix}=\begin{pmatrix}-x+y-z \\  2x+2z \\  4x-2y+4z\end{pmatrix}
\end{align*}
\begin{enumerate}
 \item  {\it Ecrivons la matrice $A$ de $f$ dans la base $\cal B$}.

On a $f(\vec e_1)=(-1,2,4)$, $f(\vec e_2)=(1,0,-2)$ et $f(\vec e_3)=(-1,2,4)$. D'où la matrice
$$A=\begin{pmatrix}-1&1&-1 \\ 2&0&2 \\ 4&-2&4\end{pmatrix}.$$
 \item  {\it Déterminons les sous-espaces $\ker f$ et $\Im f$.}

Le sous-espace vectoriel $\Im f$ est engendré par les vecteurs $f(\vec e_1)$, $f(\vec e_2)$ et $f(\vec e_3)=f(\vec e_1)$, c'est donc le plan vectoriel engendré par les vecteurs $f(\vec e_1)=(-1,2,4)$ et $f(\vec e_2)=(1,0,-2)$ qui sont clairement linéairement indépendants.

Pour le noyau, on a $\ker f=\{\vec u\in E,\ f(\vec u)=\vec 0\}$, ainsi,
$$\vec u=(x,y,z)\in\ker f\iff \left\{\begin{align*}-x+y-z=0 \\  2(x+z)=0 \\  4x-2y+4z=0\end{align*}\right.
\iff\left\{\begin{align*}x+z=0 \\  y=0\end{align*}\right.$$
C'est donc la droite vectorielle engendrée par le vecteur $\vec v=(1,0,-1)$.
 \item Soient $\vec u_1=(1,0,-1)$, $\vec u_2=(1,2,0)$ et $\vec u_3=(0,1,1)$. {\it Démontrons que $(\vec u_1,\vec u_2,\vec u_3)$ est une base de $E$.}

Pour cela nous allons vérifier que le déterminant de leurs coordonnées est non nul,
$$\begin{vmatrix}1&1&0 \\ 0&2&1 \\ -1&0&1\end{vmatrix}=\begin{vmatrix}2&1 \\ 0&1\end{vmatrix}-\begin{vmatrix}1&0 \\ 2&1\end{vmatrix}=2-1=1\neq0.$$
Ainsi, les trois vecteurs $\vec u_1,\vec u_2,\vec u_3$ sont linéairement indépendants, ils forment donc une base de $E$, car $E$ est de dimension $3$.
 \item  {\it Calculons $f(\vec u_1)$, $f(\vec u_2)$ et $f(\vec u_3)$ et déterminons la matrice $B$ de $f$ dans la base $(\vec u_1,\vec u_2,\vec u_3)$.}

On a $f(\vec u_1)=\vec 0$,

$f(\vec u_2)=\begin{pmatrix}-1+2 \\  2 \\  4-4 \\ \end{pmatrix}=\begin{pmatrix}1 \\ 2 \\ 0\end{pmatrix}=\vec u_2$.

$f(\vec u_3)=\begin{pmatrix}1-1 \\  2 \\  -4+2 \\\end{pmatrix}=\begin{pmatrix}0 \\ 2 \\ 2\end{pmatrix}=2\vec u_3$.

Ainsi la matrice $B$ de $f$ dans la base $(\vec u_1,\vec u_2,\vec u_3)$ s'écrit
$$B=\begin{pmatrix}0&0&0 \\ 0&1&0 \\ 0&0&2\end{pmatrix}.$$
 \item  {\it Déterminons les valeurs propres de $f$ et, pour chacune, un vecteur propre.}

D'après la question précédente, les valeurs propres de $f$ sont $0,1$ et $2$, et les vecteurs propres sont $\vec u_1$ pour la valeur propre $0$, $\vec u_2$ pour la valeur propre $1$ et $\vec u_3$ pour la valeur propre $2$.
\end{enumerate}
\fincorrection
\correction{002613}
 Soit $E$ un espace vectoriel de dimension $n$. On cherche à déterminer une matrice $A\in{\cal M}_n(\R)$ telle que $A^2=-I_n$, où $I_n$ désigne la matrice identité d'ordre $n$. On notera $f$ l'endomorphisme de $E$ de matrice $A$ dans la base canonique.
\begin{enumerate}
 \item {\it Démontrons que l'existence d'une telle matrice implique la parité de $n$.}

Supposons qu'il existe $A\in{\cal M}_n(\R)$ telle que $A^2=-I_n$, on a alors
$$\det(A^2)=(\det A)^2=(-1)^n,$$
ce qui implique $n$ pair, car un carré est toujours positif.
 \item  On suppose maintenant que $n=4$.
   \begin{enumerate}
   \item {\it Démontrons que pour tout $\vec x\in E$, $\vec x\neq 0$, les vecteurs $\vec x$ et $f(\vec x)$ sont linéairement indépendants.}

Soit $\vec x\in E$, on suppose $\vec x\neq 0$, supposons qu'il existe des réels $a,b$ tels que $a\vec x+bf(\vec x)=\vec0$, on a alors
$$a\vec x+bf(\vec x)=\vec0\Longrightarrow f(a\vec x+bf(\vec x))=\vec0\Longrightarrow af(\vec x)-b\vec x=\vec 0,$$
car $f^2=-\mathrm{id}_E$. Or,
$$\left\{\begin{align*}a\vec x+bf(\vec x)=\vec0 \\  af(\vec x)-b\vec x=\vec 0\end{align*}\right.\Longrightarrow(a^2+b^2)\vec x=\vec0,$$
ce qui implique $a^2+b^2=0$ car $\vec x\neq \vec 0$, et, donc $a=b=0$. Ce qui prouve que les vecteurs $\vec x$ et $f(\vec x)$ sont linéairement indépendants.
   \item  Soit $\vec x_1\neq0$, on note $F$ le sous-espace vectoriel de $E$ engendré par les vecteurs $\vec x_1$ et $f(\vec x_1)$.
            \begin{enumerate}
               \item {\it Démontrons que $F$ est stable par $f$.}

Soit $\vec x\in F$, il existe $(a,b)\in\R^2$ tel que $\vec x=a\vec x_1+bf(\vec x_1)$, d'où
$$f(\vec x)=f(a\vec x_1+bf(\vec x_1))=af(\vec x_1)+bf^2(\vec x)=af(\vec x_1)-b\vec x_1\in F.$$
D'où la stabilité de $F$ par $f$.
               \item  Soit $\vec x_2\in E$, on suppose que $\vec x_2\not\notin F$. 
	       
	       {\it Démontrons que ${\cal B}=(\vec x_1,f(\vec x_1),\vec x_2,f(\vec x_2))$ est une base de $E$.}
La dimension de $E$ étant égale à $4$, il suffit de démontrer que les vecteurs sont linéairement indépendants. Supposons qu'il existe $(a_1,b_1,a_2,b_2)\in\R^4$ tel que $$a_1\vec x_1+b_1f(\vec x_1)+a_2\vec x_2+b_2f(\vec x_2)=\vec0,$$
on a alors,
$$a_2\vec x_2+b_2f(\vec x_2)\in F,$$
et, comme $F$ est stable par $f$, 
$$f(a_2\vec x_2+b_2f(\vec x_2))=a_2f(\vec x_2)-b_2\vec x_2\in F.$$
Ce qui implique 
$$(a_2^2+b_2^2)\vec x_2\in F\quad{\hbox{d'où}}\quad a_2^2+b_2^2=0 $$
car on a supposé $\vec x_2\not\notin F$. On a donc $a_2=b_2=0$ et, par conséquent, $a_1\vec x_1+b_1f(\vec x_1)=0$, or les vecteurs $\vec x_1$ et $f(\vec x_1)$ sont linéairement indépendants, ce qui implique $a_1=b_1=0$. D'où l'indépendance des vecteurs $\vec x_1,f(\vec x_1),\vec x_2,f(\vec x_2)$.

             \end{enumerate}
   \item {\it Ecrivons la matrice $A$ de $f$ dans la base $\cal B$.}

On calcule les images des vecteurs de la base $\cal B$. On a $f(\vec x_1)=f(\vec x_1)$, $f(f(\vec x_1))=-\vec x_1$, $f(\vec x_2)=f(\vec x_2)$, $f(f(\vec x_2))=-\vec x_2$. D'où la matrice
$$A=\begin{pmatrix}0&-1&0&0 \\  1&0&0&0 \\ 0&0&0&-1 \\ 0&0&1&0\end{pmatrix}$$

   \item {\it Calculons $\det f$ et $\det(\lambda\mathrm{id}_E-f)$ pour $\lambda\in\R$.}

On a, en développant par blocs,
$$\det f=\det A=\begin{vmatrix}0&-1&0&0 \\  1&0&0&0 \\ 0&0&0&-1 \\ 0&0&1&0\end{vmatrix}=
\begin{vmatrix}0&-1 \\  1&0\end{vmatrix}.\begin{vmatrix}0&-1 \\  1&0\end{vmatrix}=1.$$
De même,
$$\det (\lambda\mathrm{id}_E-f)=\det (\lambda I_4-A)=\begin{vmatrix}\lambda&-1&0&0 \\  1&\lambda&0&0 \\ 0&0&\lambda&-1 \\ 0&0&1&\lambda\end{vmatrix}=\begin{vmatrix}\lambda&-1 \\  1&\lambda\end{vmatrix}.\begin{vmatrix}\lambda&-1 \\  1&\lambda\end{vmatrix}=(\lambda^2+1)^2.$$
  \item {\it L'endomorphisme $f$ admet-il des valeurs propres réelles ?}

Les valeurs propres réelles de $f$ sont les réels $\lambda$ qui annulent $\det (\lambda\mathrm{id}_E-f)$, ce sont donc les réels $\lambda$ tels que $\lambda^2+1=0$. Ainsi, $f$ n'admet pas de valeurs propres réelles.
  \end{enumerate}
\end{enumerate}
\fincorrection
\correction{002614}
Soient ${a\in \R}$, ${b\in \R}$ et $A$ la matrice 

$$\begin{pmatrix} 1&a&0\cr
0&1&b\cr
0&0&2\cr
\end{pmatrix} $$

\begin{enumerate}
 \item  {\it Donnons les valeurs de $a$ et de $b$ pour lesquelles la décomposition de Dunford de $A$ est}



$$A=\begin{pmatrix}
1 & 0 & 0 \cr
0 & 1 & 0 \cr
0 & 0 & 2 \cr
\end{pmatrix} + 
\begin{pmatrix}
 0 & a & 0 \cr
0 & 0 & b \cr
0 & 0 & 0 \cr
\end{pmatrix} 
$$
Notons $D=\begin{pmatrix}
1 & 0 & 0 \cr
0 & 1 & 0 \cr
0 & 0 & 2 \cr
\end{pmatrix}$ et $N=\begin{pmatrix}
 0 & a & 0 \cr
0 & 0 & b \cr
0 & 0 & 0 \cr
\end{pmatrix}$.
Cette décomposition de $A=D+N$ est sa décomposition de Dunford si et seulement si $N$ est nilpotente (il est clair que $D$ est diagonale) et si $ND=DN$.

Vérifions que $N$ est nilpotente :
$$N^2=\begin{pmatrix}
 0 & a & 0 \cr
0 & 0 & b \cr
0 & 0 & 0 \cr
\end{pmatrix}.\begin{pmatrix}
 0 & a & 0 \cr
0 & 0 & b \cr
0 & 0 & 0 \cr
\end{pmatrix}=\begin{pmatrix}
 0 & 0 & ab \cr
0 & 0 & 0 \cr
0 & 0 & 0 \cr
\end{pmatrix},\quad N^3=\begin{pmatrix}
 0 & 0 & 0 \cr
0 & 0 & 0 \cr
0 & 0 & 0 \cr
\end{pmatrix}$$
ainsi la matrice $N$ est bien nilpotente quelques soient les valeurs de $a$ et $b$. Déterminons pour quelles valeurs de $a$ et $b$ les matrices commutent.

$$N.D=\begin{pmatrix}
 0 & a & 0 \cr
0 & 0 & b \cr
0 & 0 & 0 \cr
\end{pmatrix}.\begin{pmatrix}
1 & 0 & 0 \cr
0 & 1 & 0 \cr
0 & 0 & 2 \cr
\end{pmatrix}=\begin{pmatrix}
0 & a & 0 \cr
0 & 0 & 2b \cr
0 & 0 & 0 \cr
\end{pmatrix}$$ et 
$$D.N=\begin{pmatrix}
1 & 0 & 0 \cr
0 & 1 & 0 \cr
0 & 0 & 2 \cr
\end{pmatrix}.\begin{pmatrix}
 0 & a & 0 \cr
0 & 0 & b \cr
0 & 0 & 0 \cr
\end{pmatrix}=\begin{pmatrix}
 0 & a & 0 \cr
0 & 0 & b \cr
0 & 0 & 0 \cr
\end{pmatrix}.$$
Ainsi, $ND=DN$ si et seulement si $b=2b$, c'est-à-dire si $b=0$. Le paramètre $a$ peut prendre n'importe quelle valeur.
 

 \item  On suppose dans la suite que $b=1$ et $a\neq 0$.
$$A=\begin{pmatrix} 1&a&0\cr
0&1&1\cr
0&0&2\cr
\end{pmatrix}$$

 \begin{enumerate}
\item {\it Déterminons les sous espaces propres et les sous espaces caractéristiques de $A$}.
 
Commençons par déterminer les valeurs propres de $A$, ce qui est immédiat car $A$ est sous forme triangulaire. Elle admet donc deux valeurs propres, $1$ valeur propre double et $2$ valeur propre simple.
 
Notons $E_1$ et $E_2$ les sous-espaces propres de $A$.

$$E_{1}=\{\vec u=(x,y,z)\in\R^3,\ A\vec u =\vec u\}.$$
On a 
$$\vec u\in E_{1}\iff \left\{\begin{align*} x+ay&=x\cr y+z&=y\cr 2z&=z\end{align*}\right.\iff\left\{\begin{align*}y&=0\cr z&=0\end{align*}\right.$$ 
L'espace $E_1$ est donc la droite vectorielle engendrée par le vecteur $(1,0,0)$, ce sous-espace propre associé à la valeur propre double $1$ est de dimension $1$, la matrice n'est pas diagonalisable.
 
$$E_{2}=\{\vec u=(x,y,z)\in\R^3,\ A\vec u =2\vec u\}.$$
On a 
$$\vec u\in E_{2}\iff \left\{\begin{align*}x+ay&=2x\cr y+z&=2y\cr 2z&=2z\end{align*}\right.\iff\left\{\begin{align*}x&=ay\cr y&=z\end{align*}\right.$$ 
L'espace $E_2$ est donc la droite vectorielle engendrée par le vecteur $(a,1,1)$. La valeur propre $2$ étant simple, le sous-espace caractéristique $N_2$ associé est égal à l'espace $E_2$.
 
Déterminons le sous-espace caractéristique $N_1$ associé à la valeur propre $1$. On a 

$N_1=\ker(A-I)^2$. Calculons la matrice $(A-I)^2$.
$$(A-I)^2=\begin{pmatrix} 0&a&0\cr
0&0&1\cr
0&0&1\cr
\end{pmatrix}\begin{pmatrix} 0&a&0\cr
0&0&1\cr
0&0&1\cr
\end{pmatrix}=\begin{pmatrix} 0&0&a\cr
0&0&1\cr
0&0&1\cr
\end{pmatrix}.$$
Ainsi, le noyau de $(A-I)^2$ est le plan engendré par les vecteurs $\vec e_1=(1,0,0)$ et $\vec e_2=(0,1,0)$. 
 

\item {\it Déterminons $D$ diagonalisable et $N$ nilpotente telles que $D$ commute avec $N$ et} 
$$A=D+N.$$
Notons $\vec e_3=(a,1,1)$, dans la base $(\vec e_1,\vec e_2,\vec e_3)$, la matrice associée à l'endomorphisme représenté par $A$ s'écrit 
$$B=\begin{pmatrix} 1&a&0\cr
0&1&0\cr
0&0&2\cr
\end{pmatrix}=\underbrace{\begin{pmatrix} 1&0&0\cr
0&1&0\cr
0&0&2\cr
\end{pmatrix}}_\Delta+\underbrace{\begin{pmatrix} 0&a&0\cr
0&0&0\cr
0&0&0\cr
\end{pmatrix}}_M .$$
Par construction, c'est la décomposition de Dunford de $B$ et on a $A=PBP^{-1}$ avec 
$$P=\begin{pmatrix} 1&0&a\cr
0&1&1\cr
0&0&1\cr
\end{pmatrix}\ {\hbox{et}}\ P^{-1}=\begin{pmatrix} 1&0&-a\cr
0&1&-1\cr
0&0&1\cr
\end{pmatrix}.$$
Les matrices $D=P\Delta P^{-1}$ et $N=PMP^{-1}$ vérifient, $N$ nilpotente, $D$ diagonalisable et $ND=DN$. Calculons les
$$D=P\Delta P^{-1}=\begin{pmatrix} 1&0&a\cr
0&1&1\cr
0&0&1\cr
\end{pmatrix}\begin{pmatrix} 1&0&0\cr
0&1&0\cr
0&0&2\cr
\end{pmatrix}\begin{pmatrix} 1&0&-a\cr
0&1&-1\cr
0&0&1\cr
\end{pmatrix}=\begin{pmatrix} 1&0&a\cr
0&1&1\cr
0&0&2\cr
\end{pmatrix}$$
$$N=PMP^{-1}=\begin{pmatrix} 1&0&a\cr
0&1&1\cr
0&0&1\cr
\end{pmatrix}\begin{pmatrix} 0&a&0\cr
0&0&0\cr
0&0&0\cr
\end{pmatrix}\begin{pmatrix} 1&0&-a\cr
0&1&-1\cr
0&0&1\cr
\end{pmatrix}=\begin{pmatrix} 0&a&-a\cr
0&0&0\cr
0&0&0\cr
\end{pmatrix}=A-D.$$
Ainsi
$$A=\begin{pmatrix} 1&a&0\cr
0&1&1\cr
0&0&2\cr
\end{pmatrix}=N+D=\begin{pmatrix} 0&a&-a\cr
0&0&0\cr
0&0&0\cr
\end{pmatrix}+\begin{pmatrix} 1&0&a\cr
0&1&1\cr
0&0&2\cr
\end{pmatrix}.$$
 \end{enumerate}
 
\item Soit le système différentiel suivant:

$${\cal E}: \ \ \left\{\begin{align*} x_1 '(t)&= x_1(t)+2x_2(t)\cr
x_2'(t)&= x_2(t)+x_3(t)\cr
x_3'(t)&= 2x_3(t)\cr
\end{align*}\right.$$

{\it Déterminons les solutions de ${\cal E}$}.
 
Remarquons que, si l'on note $X=(x_1,x_2,x_3)$, le système $\cal E$ s'écrit $X'=AX$ avec
$$A=\begin{pmatrix} 1&2&0\cr
0&1&1\cr
0&0&2\cr
\end{pmatrix} $$
qui correspond à la matrice $A$ précédente avec $a=2$ et $b=1$. La solution générale du système s'écrit 
$$X(t)=\exp(tA)V$$ 
où $V=(a,b,c)$ est un vecteur de $\R^3$.
 
Par ailleurs $X=PY$ est solution de $X'=AX\iff Y$ est solution de $Y'=\underbrace{P^{-1}AP}_B Y$. La solution générale du système $Y'=BY$ s'écrit $Y=\exp(tB)V$ où $V\in\R^3$. Calculons donc l'exponentielle de la matrice $tB$ pour $t\in\R$. On a vu dans la question précédente que $B=\Delta + M$ avec $\Delta M=M\Delta$, ainsi
$$\begin{align*}\exp(tB)&=\exp(t\Delta).\exp(tM)\cr
&=\begin{pmatrix} e^t&0&0\cr
0&e^t&0\cr
0&0&e^{2t}\cr
\end{pmatrix}.(I+tM)\quad {\hbox{car}}\ M^2=O\cr
&=\begin{pmatrix} e^t&0&0\cr
0&e^t&0\cr
0&0&e^{2t}\cr
\end{pmatrix}\begin{pmatrix} 1&2t&0\cr
0&1&0\cr
0&0&1\cr
\end{pmatrix}=\begin{pmatrix} e^t&2te^t&0\cr
0&e^t&0\cr
0&0&e^{2t}\cr
\end{pmatrix}.
\end{align*}$$
La solution générale du système $\cal E$ s'écrit donc $X=P\exp(tB).V$ où $V=(a,b,c)\in\R^3$. C'est-à-dire
$$X(t)=\begin{pmatrix} 1&0&2\cr
0&1&1\cr
0&0&1\cr
\end{pmatrix}\begin{pmatrix} e^t&2te^t&0\cr
0&e^t&0\cr
0&0&e^{2t}\cr
\end{pmatrix}\begin{pmatrix} a\cr
b\cr
c\cr
\end{pmatrix}=\begin{pmatrix} e^t&2te^t&2e^{2t}\cr
0&e^t&e^{2t}\cr
0&0&e^{2t}\cr
\end{pmatrix}\begin{pmatrix} a\cr
b\cr
c\cr
\end{pmatrix}$$ 
où encore
$$\left\{\begin{align*}x_1(t)&=(a+2bt)e^t+2ce^{2t}\cr
x_2(t)&=be^t+ce^{2t}\cr
x_3(t)&=ce^{2t}\end{align*}\right.$$
\end{enumerate}
\fincorrection
\correction{002615}
{\bf Questions préliminaires} : 
 \begin{enumerate}
\item[(a)] Soient  $E$ un espace vectoriel réel de dimension $n$ et  $u$ un endomorphisme de $E$. Soit $P\in\R[X]$ un polynôme. Soit $\lambda$ une valeur propre de $u$ et $\vec x$ un vecteur propre associé à $\lambda$. 
 
{\it Démontrons que $\vec x$ est vecteur propre de l'endomorphisme $P(u)$ pour la valeur propre $P(\lambda)$.}
 
On a $u(\vec x)=\lambda \vec x$, et, par récurrence sur $n$, pour tout $n\in\N$, $u^n(\vec x)=\lambda^n\vec x$. Notons 

$P(X)=a_0+a_1X+a_2X^2+\dots+a_dX^d$, l'endomorphisme $P(u)$ vérifie
\begin{align*}
P(u)(\vec x)&=(a_0\mathrm{id}_E+a_1u+a_2u^2+\dots+a_du^d)(\vec x)\\
  &=a_0\vec x+a_1\lambda\vec x+a_2\lambda^2\vec x+\dots+a_d\lambda^d\vec x\\
  &=P(\lambda) \vec x \\
\end{align*}

ce qui prouve que le vecteur $\vec x$ est vecteur propre de l'endomorphisme $P(u)$ pour la valeur propre $P(\lambda)$.


\item[(b)] 
 Théorème de Hamilton-Cayley. Soient  $E$ un espace vectoriel réel de dimension $n$ et  $u$ un endomorphisme de $E$. Soit $P$ le polynôme caractéristique de $u$, alors $P(u)=0$ (le zéro étant celui de l'ensemble des endomorphisme de $E$)
 
{\it Version matricielle} : Si $A\in{\cal M}_n(\C)$ est une matrice et $P_A$ son polynôme caractéristique, alors $P_A(A)=0$ (le zéro étant celui de ${\cal M}_n(\C)$).
\end{enumerate}


\bigskip 
Soit $$A=\begin{pmatrix} 1&0&0\cr-9& 1&9\cr 9&0&-8\end{pmatrix}.$$
 \begin{enumerate}
\item {\it Déterminons les valeurs propres de $A$.} 
 
Pour cela calculons son polynôme caractéristique :
$$P_A(X)=\begin{vmatrix}1-X&0&0\cr -9&1-X&9\cr9&0&-8-X\end{vmatrix}=(1-X)^2(-8-X).$$
La matrice $A$ admet deux valeurs propres, $1$ valeur propre double et $-8$ valeur propre simple. 
 
{\it Déterminons une base de vecteurs propres de $A$.}
 
Soit $\vec u=(x,y,z)$ un vecteur  propre associé à la valeur propre $1$, on a
$$A.\vec u=\vec u\iff \left\{\begin{align*}x&=x\cr-9x+ y+9z&=y\cr 9x-8z&=z\end{align*}\right.\iff x=z$$
Ainsi, le sous-espace propre associé à la valeur propre $1$ est un plan vectoriel dont les vecteurs $\vec e_1=(0,1,0)$ et $\vec e_2=(1,0,1)$ forment une base.
 
Soit $\vec u=(x,y,z)$ un vecteur  propre associé à la valeur propre $-8$, on a
$$A.\vec u=-8\vec u\iff \left\{\begin{align*}x&=-8x\cr-9x+ y+9z&=-8y\cr 9x-8z&=-8z\end{align*}\right.\iff \left\{\begin{align*}x&=0\cr z&=-y\end{align*}\right.$$
Ainsi, le sous-espace propre associé à la valeur propre $-8$ est une droite vectorielle engendrée par le vecteur $\vec e_3=(0,1,-1)$. 
 
Les vecteurs $\vec e_1,\vec e_2,\vec e_3$ ci-dessus forment une base de $E$ composée de vecteurs propres de $A$.

{\it Diagonalisons $A$.}
 
Dans cette base, la matrice s'écrit $D=\begin{pmatrix}1&0&0\cr0&1&0\cr0&0&-8\end{pmatrix}$ et on a
$A=PDP^{-1}$ avec 

$P=\begin{pmatrix}0&1&0\cr1&0&1\cr0&1&-1\end{pmatrix}$.
 
 \item  On cherche à déterminer une matrice $B$ telle que $B^3=A$.
 \begin{enumerate}
 
\item {\it Démontrons que si $\lambda$ est une valeur propre de $B$ alors $\lambda^3$ est une valeur propre de $A$.}
 
On considère le polynôme $P(X)=X^3$, on applique la question préliminaire a). Si $\lambda$ est une valeur propre de $B$, alors $P(\lambda)=\lambda^3$ est une valeur propre de $A=P(B)=B^3$. 

 
\item {\it Déterminons les valeurs propres de $B$ et leur multiplicité.}
 
Si $\lambda_1, \lambda_2,\lambda_3$ sont les valeurs propres de $B$ (elles existent toujours dans $\C$) et $\vec u_1,\vec u_2,\vec u_3$ les vecteurs propres associés, alors ces vecteurs sont aussi vecteurs propres de $A$ pour les valeurs propres $\lambda_1^3, \lambda_2^3,\lambda_3^3$. Sachant que les valeurs propres de $A$ sont $1$,$1$ et $-8$, on a $\lambda_1^3=\lambda_2^3=1$ et $\lambda_3^3=-8$. Ainsi les valeurs propres de $B$ sont $1$ de multiplicité $2$ et $-2$ de multiplicité $1$.
 
\item {\it Ecrivons le polynôme caractéristique de $B$.}
 
Compte tenu de la question précédente, on a
$$P_B(X)=(1-X)^2(-2-X).$$
 
\item {\it Déterminons $B$ telle que $B^3=A$.} 
 
On a $P_B(X)=(1-X)^2(-2-X)=-X^3+3X-2$, or, d'après le théorème de Cayley-Hamilton, on a $P_B(B)=0$, c'est-à-dire $-B^3+3B-2I=0$, par conséquent
$$A=B^3=3B-2I.$$
Ainsi, $B=1/3(A+2I)$, d'où,
$$B=1/3\begin{pmatrix}3&0&0\cr-9&3&9\cr9&0&-6\end{pmatrix}=\begin{pmatrix}1&0&0\cr-3&1&3\cr3&0&-2\end{pmatrix}.$$

\end{enumerate}
\end{enumerate}
\fincorrection
 \correction{003496}
 \begin{enumerate}
\item $     P=\begin{pmatrix} 1 &-5 \cr 1 & 2 \cr\end{pmatrix}$,\,
$    D= \begin{pmatrix} 6 & 0 \cr 0&-1 \cr\end{pmatrix}$
\item $    P= \begin{pmatrix} 5 & 1 \cr -4 & 1 \cr\end{pmatrix}$,\,
$     D=\begin{pmatrix}-2 & 0 \cr 0& 7 \cr\end{pmatrix}$
\item $ P=\begin{pmatrix} 3 & 1 \cr -8 &-1 \cr\end{pmatrix}$,\,
      $ \begin{pmatrix}-3 & 0 \cr 0& 2 \cr\end{pmatrix}$
\item   $P=\begin{pmatrix} 2 &-2 \cr 1 & 1 \cr\end{pmatrix}$,\,
    $ D=\begin{pmatrix} 6 & 0 \cr 0 & 2 \cr\end{pmatrix}$

\end{enumerate}
 \fincorrection
 \correction{003497}
 \begin{enumerate}
 \item  $P=\begin{pmatrix}  1 &6    &6    \cr -5 &12   &12   \cr -3 &9-\sqrt{57} &9+\sqrt{57} \cr\end{pmatrix}$,\,
     $D=\begin{pmatrix}-2 & 0&0 \cr 0&-\frac{3-\sqrt{57}}2  & 0 \cr 0 &0&-\frac{3+\sqrt{57}}2  \cr\end{pmatrix}$
 
 \item $ P=    \begin{pmatrix} 1 & \phantom-3 &-1 \cr -2 & 4 & 0 \cr 1 & 1 & 1 \cr\end{pmatrix}$,\,
$     D=\begin{pmatrix} 0 & 0&0 \cr 0& 2 & 0 \cr 0 &0&-2 \cr\end{pmatrix}$

 \item   $P=\begin{pmatrix}  1 &-2+\sqrt{3} &-2-\sqrt{3}  \cr 1 & 1 & 1 \cr 1 & 1-\sqrt{3}  & 1+\sqrt{3}  \cr\end{pmatrix}$,\,
   $ D= \begin{pmatrix} 6 & 0&0 \cr 0&\sqrt{3}    & 0 \cr 0 &0&-\sqrt{3}    \cr \end{pmatrix}$

 \item $     P=\begin{pmatrix} -7 & 5+3\sqrt{5} & 5-3\sqrt{5} \cr 11 &-4\sqrt{5} & 4\sqrt{5} \cr 2 & 5+5\sqrt{5} & 5-5\sqrt{5} \cr\end{pmatrix}$,\,
     $D=\begin{pmatrix} -3 & 0&0 \cr 0&\frac{5+\sqrt{5}}2  & 0 \cr 0 &0&\frac{5-\sqrt{5}}2     \cr \end{pmatrix}$
     
 \item $     P=\begin{pmatrix} 2 & 1 & 1 \cr -5 & 1 & 1 \cr 2 &-2 & 1 \cr\end{pmatrix}$,\,
$    D= \begin{pmatrix}-1 & 0&0 \cr 0&-3 & 0 \cr 0 &0& 6 \cr\end{pmatrix}$

 \item $     P=\begin{pmatrix} 1  & 1  & 1 \cr  i  &-i  & 1 \cr 0  & 0  & 1 \cr\end{pmatrix}$,\,
$    D= \begin{pmatrix}1+i & 0&0 \cr 0&1-i & 0 \cr 0 &0& 2 \cr\end{pmatrix}$

 \item $     P=\begin{pmatrix} 1 & 1 & 1 \cr  1 &-1 & 0 \cr  1 & 0 &-1 \cr\end{pmatrix}$,\,
$    D= \begin{pmatrix} 0 & 0&0 \cr 0& 3 & 0 \cr 0 &0& 3 \cr\end{pmatrix}$

 \item $     P=\begin{pmatrix} 1 & 2 & 1 \cr 1 & 0 & 3 \cr 0 &-1 & 2 \cr\end{pmatrix}$,\,
$     D=\begin{pmatrix} 0 & 0&0 \cr 0& 0 & 0 \cr 0 &0& 2 \cr\end{pmatrix}$

\item $P= \begin{pmatrix}-4 &-1  &-2 \cr -3 &-1 &-1 \cr  4 & 2 & 1 \cr\end{pmatrix}$,\,
$     D=\begin{pmatrix} 0 & 0&0 \cr 0&-1  & 0 \cr 0 &0& 2 \cr\end{pmatrix}$

 \item $     P=\begin{pmatrix}-1 & 2 & 3 \cr -1 & 1  & 0 \cr 1 & 0 & 2 \cr\end{pmatrix}$,\,
$ D=  \begin{pmatrix} 0 & 0&0 \cr 0& 2 & 0 \cr 0 &0& 2 \cr\end{pmatrix}$

 \end{enumerate}
 \fincorrection
 \correction{003498}
 \begin{enumerate}
\item    $P=  \begin{pmatrix}-1 &-1 & 1 &-1 \cr -1 & 1 & 3 & 3 \cr 1 & 1 & 3 &-3 \cr 1 &-1 & 1 & 1 \cr\end{pmatrix}$,\,
  $D= \begin{pmatrix} 1&0&0&0 \cr 0&-1 &0&0\cr 0&0 & 3 &0\cr 0&0&0&-3 \cr\end{pmatrix}$
\item  $ P=  \begin{pmatrix} 1 & 1 & 1 & 1 \cr 1 & 0 & 0 &-1 \cr 0 & 1 & 0 &-1 \cr 0 & 0 & 1 &-1 \cr\end{pmatrix}$,\,
   $D=  \begin{pmatrix} 2 &0&0&0 \cr 0& 2 &0&0\cr 0&0 & 2 &0\cr 0&0&0&-2 \cr\end{pmatrix}$
\item   $P=   \begin{pmatrix} 0 & 1 & 1 & 0 \cr  1 & 0 & 0 &-1 \cr  1 & 0 & 1 & 0 \cr 0 & 1 & 0 & 1 \cr\end{pmatrix}$,\,
    $ D=\begin{pmatrix} 2 &0&0&0 \cr 0& 2 &0&0\cr 0&0&-2 &0\cr 0&0&0&-2 \cr\end{pmatrix}$
 
\item     $ P=\begin{pmatrix} 1  &  0 & 30 & 18 \cr 0  &  0 & 15 &-99 \cr 0  &  0 & 21 & 99 \cr 0  &  1 &-11 & 11 \cr\end{pmatrix}$,\,
    $D= \begin{pmatrix}-5  &0&0&0 \cr 0&  2 &0&0\cr 0&0& -4 &0\cr 0&0&0&-16 \cr\end{pmatrix}$
\item      $P=\begin{pmatrix} 1 & 0 & 7 &-8 \cr 3 & 1 &12 & 9 \cr 0 & 0 & 1 & 0 \cr 0 & 0 & 1 & 6 \cr\end{pmatrix}$,\,
    $D= \begin{pmatrix} 2 &0&0&0 \cr 0& 1 &0&0\cr 0&0& 3 &0\cr 0&0&0&-1 \cr\end{pmatrix}$
\item $1$ est valeur propre quadruple, non diagonalisable.
\item $0$ est valeur propre quadruple, non diagonalisable.
\item $0$ est vp double, $\mathrm{rg} A = 2$. Autres vp : $\frac{-3\pm3\sqrt{13}}2$, diagonalisable.
\end{enumerate}
 \fincorrection
\correction{003503}
$n$ pair : $P = \begin{pmatrix}1&       &  &   &       &1  \cr
                                     &\ddots &  &   &\cdots     \cr
                                     &       &1 &1              \cr
                                     &       &1 &-1             \cr
                                     &\cdots &  &   &\ddots     \cr
                                   1 &       &  &   &       &-1 \cr \end{pmatrix}$,
                    $D = \text{diag}(1,\dots,1,-1,\dots,-1)$.\par
       $n$ impair : $P = \begin{pmatrix}1&       &  &  &   &       &1  \cr
                                     &\ddots &  &  &   &\cdots     \cr
                                     &       &1 &  &1              \cr
                                     &       &  &1 &               \cr
                                     &       &1 &  &-1             \cr
                                     &\cdots &  &  &   &\ddots     \cr
                                   1 &       &  &  &   &       &-1 \cr \end{pmatrix}$,
                    $D = \text{diag}(1,\dots,1,-1,\dots,-1)$.
\fincorrection
\correction{003504}
$P = \bigl(\omega^{(i-1)(1-j)}\bigr)$,
         $D = \text{diag}(1,\omega,\dots,\omega^{n-1})$
         avec $\omega = \exp(2i\pi/n)$.
\fincorrection
\correction{003505}
$P = \begin{pmatrix}1& 1& 1& 1 \cr
                        1& 1&-1&-1 \cr
                        1&-1& 1&-1 \cr
                        1&-1&-1& 1 \cr\end{pmatrix}$,
         $D = \text{diag}\,(a+b+c+e,a-b-c+e,-a+b-c+e,-a-b+c+e)$.
\fincorrection
\nocorrection
\correction{003510}
\begin{enumerate}
  \item $M = \begin{pmatrix} 0 &1 &-2 &       &(0)      \cr
                              &2 &2  &\ddots           \cr
                              &  &6  &\ddots &-n(n-1)  \cr
                              &  &   &\ddots &n        \cr
                           (0)&  &   &       &n(n+1)   \cr \end{pmatrix}$.

\end{enumerate}
 \fincorrection
\correction{003517}
Oui ssi $\mathrm{tr}(A)\ne 0$ ou $A=0$.
\fincorrection
\correction{005654}
Soit $P$ un élément de $\Rr_{2n}[X]$. $f(P)$ est un polynôme de degré inférieur ou égal à $2n+1$ et de plus, si $a$ est le coefficient de $X^{2n}$ dans $P$, le coefficient de $X^{2n+1}$ dans $f(P)$ est $2na-2na = 0$. Donc $f(P)$ est un élément de $\Rr_{2n}[X]$. La linéarité de $f$ étant claire, $f$ est bien un endomorphisme de $\Rr_{2n}[X]$.

Cherchons maintenant $P$ polynôme non nul et $\lambda$ réel tels que $f(P) =\lambda P$ ce qui équivaut à

\begin{center}
$\frac{P'}{P}=\frac{2nX+\lambda}{X^2-1}=\frac{1}{2}\left(\frac{2n+\lambda}{X-1}-\frac{-2n+\lambda}{X+1}\right)$.
\end{center}

En identifiant à la décomposition en éléments simples classique
de $\frac{P'}{P}$ (à savoir si $P=K(X-z_1)^{\alpha_1}\ldots(X-z_k)^{\alpha_k}$ avec $K\neq0$ et les $z_i$ deux à deux distincts, alors $\frac{P'}{P}=\sum_{i=1}^{k}\frac{\alpha_i}{X-z_i}$), on voit que nécessairement $P$ ne peut admettre pour racines dans $\Cc$ que $-1$ et $1$ et d'autre part que $P$ est de degré $\frac{1}{2}(2n+\lambda+2n-\lambda)=2n$. $P$ est donc nécessairement  de la forme

\begin{center}
$P=aP_k$ avec $a\in\Rr^*$ et $P_k =(X-1)^k(X+1)^{2n-k}$ avec $k\in\llbracket0,2n\rrbracket$.
\end{center}

Réciproquement, chaque $P_k$ est non nul et  vérifie 

\begin{center}
$\frac{P_k'}{P_k}=\frac{k}{X-1}+\frac{2n-k}{X+1}=\frac{1}{2}\left(\frac{2n+(2k-2n)}{X-1}+\frac{2n-(2k-2n)}{X+1}\right)$.
\end{center}

  
Donc, pour chaque $k\in\llbracket0,2n\rrbracket$, $P_k$ est vecteur propre de $f$ associé à la valeur propre $\lambda_k = 2(k-n)$.

Ainsi, $f$ admet $2n+1$ valeurs propres, nécessairement simples car $\text{dim}(\Rr_{2n}[X]=2n+1$). $f$ est donc diagonalisable et les sous espaces propres de $f$ sont les droites $\text{Vect}(P_k)$, $0\leqslant k\leqslant 2n$.
\fincorrection
\correction{005663}
$\text{det}M=\text{det}\left(
\begin{array}{cc}
A&4A\\
A&A
\end{array}
\right)=\text{det}\left(
\begin{array}{cc}
-3A&4A\\
0&A
\end{array}
\right)$  ($\forall k\in\llbracket1,n\rrbracket$, $L_k\leftarrow L_k-L_{n+k}$) et donc $\text{det}M=\text{det}(A)\text{det}(-3A) = (-3)^n(\text{det}A)^2$.

\begin{center}
\shadowbox{
$\text{det}M=(-3)^n(\text{det}A)^2$.
}
\end{center}

L'idée de l'étude de $M$ qui suit vient de l'étude de la matrice de format $2$, $B=\left(
\begin{array}{cc}
1&4\\
1&1
\end{array}
\right)$.

Une diagonalisation rapide amène à $B=\left(
\begin{array}{cc}
-2&2\\
1&1
\end{array}
\right)\times\left(
\begin{array}{cc}
-1&0\\
0&3
\end{array}
\right)\times\frac{1}{4}\left(
\begin{array}{cc}
-1&2\\
1&2
\end{array}
\right)$. Soit alors $P$ la matrice de format $2n$ définie par blocs par $P=\left(
\begin{array}{cc}
-2I_n&2I_{n}\\
I_n&I_n
\end{array}
\right)$.
Un calcul par blocs montre que $P$ est inversible et que $P^{-1}=\frac{1}{4}\left(
\begin{array}{cc}
-I_n&2I_{n}\\
I_n&2I_n
\end{array}
\right)$ puis que

\begin{center} 
$P^{-1}MP=\frac{1}{4}\left(
\begin{array}{cc}
-I_n&2I_{n}\\
I_n&2I_n
\end{array}
\right)\left(
\begin{array}{cc}
A&4A\\
A&A
\end{array}
\right)\left(
\begin{array}{cc}
-2I_n&2I_{n}\\
I_n&I_n
\end{array}
\right)=\frac{1}{4}\left(
\begin{array}{cc}
A&-2A\\
3A&6A
\end{array}
\right)\left(
\begin{array}{cc}
-2I_n&2I_{n}\\
I_n&I_n
\end{array}
\right)
=\left(
\begin{array}{cc}
-A&0\\
0&3A
\end{array}
\right)$.
\end{center}

On pose $N=\left(
\begin{array}{cc}
-A&0\\
0&3A
\end{array}
\right)$. Puisque les matrices $M$ et $N$ sont semblables, $M$ et $N$ ont même polynôme caractéristique et de plus $M$ est diagonalisable si et seulement si $N$ l'est.

Cherchons les vecteurs propres $Z$ de $N$ sous la forme $Z=\left(
\begin{array}{c}
X\\
Y
\end{array}
\right)$ où $X$ et $Y$ sont des vecteurs colonnes de format $n$. Sois $\lambda\in\Cc$.

\begin{center}
$NZ=\lambda Z\Leftrightarrow\left(
\begin{array}{cc}
-A&0\\
0&3A
\end{array}
\right)\left(
\begin{array}{c}
X\\
Y
\end{array}
\right)=\lambda\left(
\begin{array}{c}
X\\
Y
\end{array}
\right)\Leftrightarrow -AX =\lambda X\;\text{et}\;3AY =\lambda Y$.
\end{center}

Par suite

\begin{align*}\ensuremath
Z\;\text{est vecteur propre de}\;N\;\text{associé à}\;\lambda\Leftrightarrow(X\neq0\;\text{ou}\;Y\neq0)\;\text{et}\; 
(X\in\text{Ker}(A+\lambda I)\;\text{et}\;Y\in\text{Ker}\left(A-\frac{\lambda}{3}I\right)).
\end{align*}

Une discussion suivant $\lambda$ s'en suit :

\textbf{1er cas.} Si $-\lambda$ et $\frac{\lambda}{3}$ ne sont pas valeurs propres de $A$ alors $\lambda$ n'est pas valeur propre de $M$.

\textbf{2ème cas.} Si $-\lambda$ est dans $\text{Sp}A$ et $\frac{\lambda}{3}$ n'y est pas, alors $\lambda$ est valeur propre de $M$.
Le sous-espace propre associé est l'ensemble des $P\left(
\begin{array}{c}
X\\
0
\end{array}
\right)=\left(
\begin{array}{c}
-2X\\
X
\end{array}
\right)$ où $X$ décrit $\text{Ker}(A+\lambda I)$. La dimension de $E_\lambda$ est alors $\text{dim}(\text{Ker}(A+\lambda I)$.

\textbf{3ème cas.} Si $-\lambda$ n'est pas dans $\text{Sp}A$ et $\frac{\lambda}{3}$ y est, alors $\lambda$ est valeur propre de $M$.
Le sous-espace propre associé est l'ensemble des $P\left(
\begin{array}{c}
0\\
Y
\end{array}
\right)=\left(
\begin{array}{c}
2Y\\
Y
\end{array}
\right)$ où $Y$ décrit $\text{Ker}\left(A-\frac{\lambda}{3}I\right)$. La dimension de $E_\lambda$ est alors $\text{dim}\left(\text{Ker}\left(A-\frac{\lambda}{3}I\right)\right)$.

\textbf{4ème cas.} Si $-\lambda$ est dans $\text{Sp}A$ et $\frac{\lambda}{3}$  aussi, alors $\lambda$ est valeur propre de $M$.
Le sous-espace propre associé est l'ensemble des $P\left(
\begin{array}{c}
X\\
Y
\end{array}
\right)=\left(
\begin{array}{c}
-2X+2Y\\
X+Y
\end{array}
\right)$ où $X$ décrit $\text{Ker}(A+\lambda I)$  et $Y$ décrit $\text{Ker}\left(A-\frac{\lambda}{3}I\right)$. La dimension de $E_\lambda$ est alors $\text{dim}(\text{Ker}(A+\lambda I)) +\text{dim}\left(\text{Ker}\left(A-\frac{\lambda}{3}I\right)\right)$.

Dans tous les cas, $\text{dim}(E_\lambda(M))=\text{dim}(E_{-\lambda}(A))+\text{dim}(E_{\lambda/3}(A))$ (et en particulier $\text{dim}(\text{Ker}M)=2\text{dim}(\text{Ker}A)$).
Comme les applications $\lambda\mapsto -\lambda$ et $\lambda\mapsto\frac{\lambda}{3}$ sont des bijections de $\Cc$ sur lui-même,

\begin{align*}\ensuremath
A\;\text{est diagonalisable}&\Leftrightarrow\sum_{\lambda\in\Cc}^{}\text{dim}(E_{\lambda}(A)) = n\\
 &\Leftrightarrow\sum_{\lambda\in\Cc}^{}\text{dim}(E_{\lambda}(A))+\sum_{\lambda\in\Cc}^{}\text{dim}(E_{\lambda}(A))= 2n\\
 &\Leftrightarrow\sum_{\lambda\in\Cc}^{}\text{dim}(E_{-\lambda}(A))+\sum_{\lambda\in\Cc}^{}\text{dim}(E_{\lambda/3}(A))= 2n\\
 &\Leftrightarrow M\;\text{est diagonalisable}.
\end{align*}
\fincorrection
\correction{005667}
Soit $f$ l'endomorphisme de $\Rr^n$ de matrice $A$ dans la base canonique $(e_1,\ldots,e_n)$ de $\Rr^n$. $\forall i\in\llbracket1,n\rrbracket$, $f(e_{i})=e_{n+1-i}$ et donc $\forall i\in\llbracket1,n\rrbracket$, $f^2(e_{i})=e_i$. Donc $f$ est une symétrie distincte de l'identité et en particulier $\text{Sp}A=\{-1,1\}$ et $f$ est diagonalisable. On en déduit que $A$ est diagonalisable dans $\Rr$.
\fincorrection
\correction{005673}
$A$ est de format $2$ et donc, soit a deux valeurs propres distinctes et est dans ce cas diagonalisable dans $\Cc$, soit a une valeur propre double $\lambda$ non nulle car $\text{Tr}A = 2\lambda \neq 0$.

Dans ce dernier cas, $A^2$ est diagonalisable et est donc est semblable à $\text{diag}(\lambda^2,\lambda^2) =\lambda^2I$. Par suite, $A^2 =\lambda^2I$. Ainsi, $A$ annule le polynôme $X^2-\lambda^2=(X-\lambda)(X+\lambda)$ qui est scindé sur $\Rr$ à racines simples. Dans ce cas aussi, $A$ est diagonalisable.

\fincorrection
\correction{005675}
Trouvons un polynôme scindé à racines simples annulant $f$.

Le polynôme $P =X(X-\lambda)(X-\mu) = X^3 - (\lambda+\mu)X^2 +\lambda\mu X$ est annulateur de $f$. En effet,

\begin{align*}\ensuremath
P(f)&=f^3 -(\lambda+\mu)f^2 +\lambda\mu f = (\lambda^3 - (\lambda+\mu)\lambda^2 +(\lambda\mu)\lambda)u + (\mu^3 - (\lambda+\mu)\mu^2 +(\lambda\mu)\mu)v\\
 &= P(\lambda)u + P(\mu)v=0 .
\end{align*}

\textbullet~Si $\lambda$ et $\mu$ sont distincts et non nuls, $P$ est un polynôme scindé à racines simples annulateur de $f$ et donc $f$ est diagonalisable.

\textbullet~Si $\lambda=\mu=0$, alors $f=0$ et donc $f$ est diagonalisable.

\textbullet~Si par exemple $\lambda\neq0$ et $\mu=0$, $f^2 =\lambda^2u =\lambda f$ et et le polynôme $P=X(X-\lambda)$ est scindé à racines simples et annulateur de $f$. Dans ce cas aussi $f$ est diagonalisable.

\textbullet~Enfin si $\lambda=\mu\neq0$, $f^2 =\lambda^2(u+v) =\lambda f$ et de nouveau $P=X(X-\lambda)$ est scindé à racines simples et annulateur de $f$.

Dans tous les cas, $f$ est diagonalisable.
\fincorrection
\correction{005682}
(Si les $a_k$ sont réels, la matrice $A$ est symétrique réelle et les redoublants savent que la matrice $A$ est diagonalisable.)

Si tous les $a_k$, $1\leqslant k\leqslant n-1$, sont nuls la matrice $A$ est diagonalisable car diagonale.

On suppose dorénavant que l'un au moins des $a_k$, $1\leqslant k\leqslant n-1$, est non nul. Dans ce cas, $\text{rg}A = 2$.

$0$ est valeur propre d'ordre $n-2$ au moins. Soient $\lambda$ et $\mu$ les deux dernières valeurs propres. On a

\begin{center}
$\lambda +\mu=\text{Tr}A=a_n$ et $\lambda^2 +\mu^2=\text{Tr}(A^2)=\sum_{k=1}^{n-1}a_k^2+\sum_{k=1}^{n}a_k^2 = 2\sum_{k=1}^{n-1}a_k^2 +a_n^2$.
\end{center}

$\lambda$ et $\mu$ sont solutions du système $\left\{
\begin{array}{l}
\lambda+\mu=a_n\\
\lambda^2+\mu^2=2\sum_{k=1}^{n-1}a_k^2 +a_n^2
\end{array}
\right.$ qui équivaut au système $\left\{
\begin{array}{l}
\lambda+\mu=a_n\\
\lambda^2+\mu^2=-\sum_{k=1}^{n-1}a_k^2
\end{array}
\right.$ $(S)$.

On a alors les situations suivantes :

\textbullet~Si $\lambda$ et $\mu$ sont distincts et non nuls, $A$ est diagonalisable car l'ordre de multiplicité de chaque valeur propre est égale à la dimension du sous-espace propre correspondant.

\textbullet~Si $\lambda$ ou $\mu$ est nul, $A$ n'est  pas diagonalisable car l'ordre de multiplicité de la valeur propre $0$ est différent de $n-2$, la dimension du noyau de $A$.

\textbullet~Si $\lambda=\mu\neq 0$, $A$ est diagonalisable si et seulement si $\text{rg}(A-\lambda I) = n-2$ mais on peut noter que si $\lambda$ n'est pas nul, on a toujours $\text{rg}(A-\lambda I)= n-1$ en considérant la matrice extraite formée des n-1 premières lignes et colonnes.

En résumé, la matrice $A$ est diagonalisable si et seulement si le système $(S)$ admet deux solutions distinctes et non nulles.

Mais $\lambda$ et $\mu$ sont solutions du système $(S)$ si et seulement si $\lambda$ et $\mu$ sont les racines de l'équation $(E)$ : $X^2 - a_nX- \sum_{k=1}^{n-1}a_k^2 = 0$. Par suite, $A$ est diagonalisable si et seulement si $\sum_{k=1}^{n-1}a_k^2 = 0$ et $\Delta= a_n^2+4\sum_{k=1}^{n-1}a_k^2\neq0$.
\fincorrection
\nocorrection
\correction{003525}
3. $P = \begin{pmatrix} -1 &\phantom-3 & 3 \cr 2 & 4 & 0 \cr -1 & 1 & -3 \cr\end{pmatrix}$,
    $D = \begin{pmatrix} 0 & 2 & -2 \cr \end{pmatrix}$.
 \fincorrection
\correction{003526}
\begin{enumerate}
  \item 1 si $C \ne 0$, 0 si $C = 0$.
  \item $\dim(E_0) \ge n-1  \Rightarrow  X^{n-1}$ divise $\chi_M  \Rightarrow 
             \chi_M = (-1)^n(X^n - (a_1^2+\dots+a_n^2)X^{n-1})$.
  \item Oui.
\end{enumerate}
 \fincorrection
\correction{003527}
$\mathrm{rg} A = 1$ donc $\dim\mathrm{Ker} A = n-1$ et $0$ est valeur propre
             d'ordre au moins $n-1$. La somme des valeurs propres est
             $\mathrm{tr} A = n$ donc la drenière valeur propre est $n$ et le
             sous-espace propre associé est de dimension~$1$. Donc $A$
             est diagonalisable.
\fincorrection
\correction{003528}
\begin{enumerate}
  \item La fonction $f_n : x  \mapsto \frac{P_n(x)}{x^n}$ croît strictement
             de $-\infty$ à 1 quand $x$ varie de $0$ à $+\infty$.
  \item $\chi_A(x) = (-1)^n\Bigl(x^n - \sum_{k=1}^n kx^{n-k}\Bigr)$.
\end{enumerate}
 \fincorrection
\correction{003529}
Soit $M = (x_iy_j)$ : $M$ est de rang inférieur ou égal à $1$, donc
$0$ est valeur propre de $M$ d'ordre au moins $n-1$.
Comme $\mathrm{tr}(M) = x_1y_1 + \dots + x_ny_n$, le polynôme caractéristique de
$M$ est $\chi_M(x) = (-x)^{n-1}(x_1y_1 + \dots + x_ny_n - x)$, et le déterminant demandé
est $\Delta_n = \chi_M(-1) = x_1y_1 + \dots + x_ny_n + 1$.
\fincorrection
\correction{003530}
\begin{enumerate}
  \item $\det(M+(t))$ est une fonction affine de~$t$.
  \item $|\lambda+a| = k|\lambda+b|$ et $\lambda=x+iy  \Rightarrow 
    (1-k^2)(x^2+y^2)+\dots = 0$, équation d'un cercle si $|a|\ne |b|$.
\end{enumerate}
 \fincorrection
\correction{003531}
\begin{enumerate}
  \item $a_1\dots a_n + b_1a_2\dots a_n + a_1b_2a_3\dots a_n + \dots + a_1\dots a_{n-1}b_n$.
  \item 
  \item $\frac{\chi_A(t)}{\prod_{i=1}^n(a_i-t)} = 1 + \sum_{i=1}^n\frac{b_i}{a_i-t}$ change de signe
    entre deux $a_i$ successifs et dans l'un des intervalles $]-\infty,a_1[$ ou $]a_n,+\infty[$
    donc $\chi_A$ admet $n$ racines distinctes.
  \item Oui. Supposons par exemple $a_1=\dots=a_p <a_{p+1}<\dots<a_n$~:
    La question précédente met en évidence $n-p$ racines simples de $\chi_A$ entre
    les $a_i$ et $\pm\infty$, et $a_1$ est aussi racine d'ordre $p-1$
    de $\chi_A$. Or les $p$ premières lignes de $A-a_1I$ sont égales
    donc $\mathrm{rg}(A-a_1I)\le n-p+1$ et $\dim(\mathrm{Ker}(A-a_1I))\ge p-1$ d'où la
    diagonalisabilité. Le cas où il y a plusieurs groupes de $a_i$ égaux
    se traite de même.
\end{enumerate}
 \fincorrection
\correction{003532}
$\chi_B(X) = \frac{(-X)^n}{\det(A)}\chi_A(1/X)$,\quad$\chi_C(X^2) = \chi_A(X)\chi_A(-X)$.
\fincorrection
\nocorrection
\nocorrection
\correction{003535}
$(a)\Leftrightarrow(b)$~: thm du rang.

$(c)\Leftrightarrow(d)$~: immédiat.

$(c) \Rightarrow (b)$~: si $AX=XB$ alors pour tout polynôme $P$ on a $P(A)X = XP(B)$.

$\overline{(c)} \Rightarrow \overline{(b)}$~: prendre $U$ vecteur propre de $A$,
$V$ vecteur propre de ${}^tB$ associés à la même valeur propre et $X=U^tV$.
\fincorrection
\nocorrection
\nocorrection
\correction{003538}
Somme des valeurs propres $= n$.
\fincorrection
\nocorrection
\correction{003540}
Soit $ K=\Z/p\Z$. Il faut en fait prouver que pour toute matrice $A\in\mathcal{M}_n(K)$
on a $\mathrm{tr}(A^p)=\mathrm{tr}(A)$. Remarquer qu'on n'a pas forcément $A^p=A$ dans $\mathcal{M}_n(K)$,
c'est faux, entre autres, si $A$ est nilpotente d'indice 2.
Soit $X$ une indéterminée sur $ K$. On a dans l'anneau $\mathcal{M}_n(K[X])$~:
$(A-XI_n)^p = A^p-X^pI_n$, d'où, en prenant les déterminants~:
$\chi_{A^p}(X^p) = \chi_A(X)^p = \chi_A(X^p)$ et on égale les coefficients de $X^{(n-1)p}$.
\fincorrection
\correction{005656}
Si $p = q$, le résultat est connu : $\chi_{AB}=\chi_{BA}$.

Supposons par exemple $p < q$. On se ramène au cas de matrices carrées en complétant. 
Soient $A'=\left(
\begin{array}{c}
A\\
0_{q-p,q}
\end{array}
\right)$ et $B'=\left(
\begin{array}{cc}
B&0_{q,q-p}
\end{array}
\right)$. $A'$ et $B'$ sont des matrices carrées de format $q$ et $A'B'$ et $B'A'$ ont même polynôme caractéristique.

Un calcul par blocs donne $B'A'=BA$ et $A'B'=\left(
\begin{array}{c}
A\\
0_{q-p,q}
\end{array}
\right)\left(
\begin{array}{cc}
B&0_{q,q-p}
\end{array}
\right)=\left(
\begin{array}{cc}
AB&0_{p,q-p}\\
0_{q-p,p}&0_{q-p,q-p}
\end{array}
\right)$. Donc $\chi_{BA}=(-X)^{q-p}\chi_{AB}$ ou encore, avec une écriture plus symétrique, $(-X)^p\chi_{BA}= (-X)^q\chi_{AB}$ ce qui vrai dans tous les cas.

\begin{center}
\shadowbox{
$\forall A\in\mathcal{M}_{p,q}(\Kk)$, $\forall B\in\mathcal{M}_{q,p}(\Kk)$, $(-X)^p\chi_{BA}= (-X)^q\chi_{AB}$.
}
\end{center}
\fincorrection
\correction{005657}
Si $u$ est inversible,

\begin{center}
$\text{det}(u+v)=\text{det}u\Leftrightarrow\text{det}u\times\text{det}(Id+u^{-1}v) =\text{det}u\Leftrightarrow\text{det}(Id+u^{-1}v)=1$.
\end{center}

$u$ et $v$ commutent et donc $u^{-1}$ et $v$ également car $uv=vu\Rightarrow u^{-1}uvu^{-1}= u^{-1}vuu^{-1}\Rightarrow vu^{-1}=u^{-1}v$. Mais alors, puisque $v$ est nilpotent, l'endomorphisme $w = u^{-1}v$ l'est également car $(u^{-1}v)^p =u^{-p}v^p)$.

Il reste donc à calculer $\text{det}(Id+w)$ où $w$ est un endomorphisme nilpotent. On remarque que $\text{det}(Id+w)=\chi_w(-1)$. Il est connu que $0$ est l'unique valeur propre d'un endomorphisme nilpotent et donc $\chi_w=(-X)^n$ puis

\begin{center}
$\text{det}(Id+w)=\chi_w(-1)=(-(-1))^n = 1$.
\end{center}

Le résultat est donc démontré dans le cas où $u$ est inversible. Si $u$ n'est pas inversible, $u+xId$ est inversible sauf pour un nombre fini de valeurs de $x$ et commute toujours avec $v$. Donc, pour tout $x$ sauf peut-être pour un nombre fini, $\text{det}(u+xId+v)=\text{det}(u+xId)$. Ces deux polynômes coïncident en une infinité de valeurs de $x$ et sont donc égaux. Ils prennent en particulier la même valeur en $0$ ce qui refournit $\text{det}(u+v)=\text{det}u$.
\fincorrection
\correction{005666}
Soit $A\in\mathcal{M}_n(\Rr)$. $A$ est antisymétrique si et seulement si ${^t}A =-A$. Dans ce cas

\begin{center}
$\chi_A =\text{det}(A-XI) =\text{det}({^t}(A-XI))=\text{det}(-A-XI)=(-1)^n\text{det}(A+XI) =(-1)^n\chi_A(-X)$
\end{center}

Ainsi, $\chi_A$ a la parité de $n$.
\fincorrection
\correction{005678}
Soit $(\lambda_1,...,\lambda_n)$ la famille des valeurs propres de $A$. On a donc $\chi_A=(\lambda_1-X)\ldots(\lambda_n-X)$.

\begin{align*}\ensuremath
\chi_A(B)\;\text{inversible}&\Leftrightarrow(B-\lambda_1I)...(B-\lambda_nI)\;\text{inversible}\\ 
&\Leftrightarrow\forall k\in\llbracket1,n\rrbracket,\;B-\lambda_kI\;\text{inversible}\;(\text{car}\;\text{det}((B-\lambda_1I)...(B-\lambda_nI))=\text{det}(B-\lambda_1I)\times...\times\text{det}(B-\lambda_nI))\\
 &\Leftrightarrow\forall k\in\llbracket1,n\rrbracket,\;\lambda_k\;\text{n'est pas valeur propre de}\;B\\
 &\Leftrightarrow\text{Sp}A\cap\text{Sp}B =\varnothing.
\end{align*}
\fincorrection
\correction{005679}
Si $P$ et $\chi_f$ sont premiers entre eux, d'après le théorème de \textsc{Bézout}, il existe deux polynômes $U$ et $V$ tels que $UP+V\chi_f = 1$. En prenant la valeur en $f$ et puisque que $\chi_f(f) = 0$, on obtient $P(f)\circ U(f) =U(f)\circ P(f) = Id$. $P(f)$ est donc un automorphisme de $E$.

Réciproquement, si $P$ et $\chi_f$ ne sont pas premiers entre eux, $P$ et $\chi_f$ ont une racine commune $\lambda$ dans $\Cc$. Soit $A$ est la matrice de $f$ dans une base donnée (si $\Kk$ n'est pas $\Cc$ l'utilisation de la matrice est indispensable). On a $P(A)=(A-\lambda I)Q(A)$ pour un certain polynôme $Q$. La matrice $A-\lambda I$ n'est pas inversible car $\lambda$ est valeur propre de $A$ et donc $P(A)$ n'est pas inversible ($\text{det}(P(A))=\text{det}(A-\lambda I)\text{det}Q(A) = 0$).
\fincorrection
\nocorrection
\nocorrection
\nocorrection
\nocorrection
\nocorrection
\nocorrection
\nocorrection
\nocorrection
\nocorrection
\nocorrection
\nocorrection
\nocorrection
\nocorrection
\nocorrection
\correction{003612}
\begin{enumerate}
  \item Valeurs propres : $1,j,j^2$.
    sev stables : $\{\vec 0\}$, vect$\{\vec e_3\}$, vect$\{\vec e_1,\vec e_2\}$
    et $\R^3$.
  \item {$\left.\vcenter{\halign{&$#$\hfil\cr
    AB = BA                     & \Rightarrow  \phi_B(\vec e_3) = \lambda\vec e_3       \cr
    {}^t\!A{}^tB = {}^tB{}^t\!A & \Rightarrow  \phi_{{}^tB}(\vec e_3) = \lambda\vec e_3 \cr}
    }\right\rbrace  \Rightarrow  B = 
    \begin{pmatrix} a+\mu &a       &0       \cr
                                     -3a   &-2a+\mu &0       \cr
                                      0    &0       &\lambda \cr \end{pmatrix}$.}
\end{enumerate}
 \fincorrection
\nocorrection
\correction{003614}
Soit $\varphi(x,y,z) = x+2y+3z$. $f$ conserve la surface de
niveau $\varphi=1$ donc par linéarité $f\circ\varphi=\varphi$ et
$\varphi$ est vecteur propre de ${}^tf$.
\fincorrection
\nocorrection
\nocorrection
\correction{003617}
Si $\chi_u$ est irréductible, pour $x\ne 0$ le polynôme minimal de~$x$
en~$u$ est égal à~$\chi_u$ donc le sous-espace cyclique engendré par~$x$ est égal à~$E$
et il n'y a pas de sous-espace stable non trivial.

Si seuls $\{0\}$ et $E$ sont stables, soit $x\ne 0$. Le sous-espace cyclique
engendré par~$x$ est égal à~$E$ donc l'annulateur minimal de~$u$ en~$x$
est égal à~$\chi_u$. Soit~$P$ un diviseur non trivial de~$\chi_u$ et
$y = P(u)(x)$~: l'annulateur minimal de~$u$ en~$y$ est $\chi_u/P$, absurde.
\fincorrection
\correction{003618}
Soit~$F$ un hyperplan de~$E$,
$<\!e\!>$ un supplémentaire stable et~$H$ un supplémentaire
de~$<\!e\!>$ stable. Si $K$ est un sev de~$H$, alors $K$ admet un supplémentaire
$K'$ dans~$E$ stable et $H\cap K'$ est un sev de~$H$ stable, en somme directe
avec~$K$.
$K'\not\subset H$ car $K\subset H$ et $K\oplus K' = E$ donc $K'+H = E$
et $\dim(H\cap K') = \dim(H)+\dim(K')-\dim(E) = \dim(H)-\dim(K)$
soit $K\oplus(H\cap K') = H$. $f_{|H}$ vérifie la même propriété que~$f$ et
on obtient par récurrence que $f$ est diagonalisable.

Réciproquement, soit $f$ diagonalisable, $F$ un sev de~$E$ et $(e_1,\dots,e_n)$
une base propre pour~$f$. On montre que $F$ admet un supplémentaire stable
par récurrence sur $\mathrm{codim}(F)$~: si $F=E$ alors $\{0\}$ convient
et si $F\ne E$ alors il existe $i$ tel que $e_i\notin F$ d'où $F\oplus{<\!e_i\!>}$
est un sur-espace strict de~$F$, admettant un supplémentaire~$G$ stable, d'où
$G\oplus{<\!e_i\!>}$ est supplémentaire de~$F$ stable.

\fincorrection
\correction{003619}
$\mathrm{Spec}(f) = \{0,1,2\}$ donc $f$ est diagonalisable et chaque sous-espace propre est de dimension~$1$.
Comme la restriction d'un diagonalisable à un sous-espace stable est encore diagonalisable,
les sous-espaces stables par~$f$ sont les huit sous-sommes de~$E_0\oplus E_1\oplus E_2$.
\fincorrection
\correction{005683}
\begin{enumerate}
 \item  $\chi_A=\left|
\begin{array}{ccc}
1-X&1&-1\\
1&1-X&1\\
1&1&1-X
\end{array}
\right|=(1-X)(X^2-2X)-(2-X)+(2-X) = -X(X-1)(X-2)$.

On est dans le cas d'une matrice diagonalisable avec 3 valeurs propres simples.

\textbf{Recherche des droites stables.} Dans chacun des cas, les droites stables sont les droites engendrées par des vecteurs propres. On obtient immédiatement les 3 droites stables : $E_0=\text{Vect}(e1)$ où $e_1=(1,-1,0)$, $E_1=\text{Vect}(e_2)$ où $e_2=(1,-1,-1)$ et $E_2=\text{Vect}(e_3)$ où $e_3=(0,1,1)$.

\textbf{Recherche des plans stables.} Soit $P$ un plan stable par $f$. La restriction de $f$ à $P$ est un endomorphisme de $P$ et on sait de plus que le polynôme caractéristique de $f_{/P}$ divise celui de $f$. $f_{/P}$ est diagonalisable car $f$ l'est  car on dispose d'un polynôme scindé à racines simples annulant $f$ et donc $f_{/P}$. On en déduit que $P$ est engendré par deux vecteurs propres indépendants de $f_{/P}$ qui sont encore vecteurs propres de $f$. On obtient trois plans stables : $P_1 =\text{Vect}(e_2,e_3)$, $P_2 =\text{Vect}(e_1,e_3)$ et $P_3 =\text{Vect}(e_1,e_2)$.

\item  $\chi_A=\left|
\begin{array}{ccc}
2-X&2&1\\
1&3-X&1\\
1&2&2-X
\end{array}
\right|=(2-X)(X^2-5X+4)-(-2X+2)+(X-1) = (1-X)((X-2)(X-4)-2-1) = (1-X)(X^2+6X-5) = -(X-1)^2(X-5)$. Puis $E_1$ est le plan d'équation $x+2y+z=0$ et $E_5=\text{Vect}((1,1,1))$.

On est toujours dans le cas diagonalisable mais avec une valeur propre double.

Les droites stables sont $E_5=\text{Vect}((1,1,1))$ et n'importe quelle droite contenue dans $E_1$. Une telle droite est engendrée par un vecteur de la forme $(x,y,-x-2y)$  avec $(x,y)\neq(0,0)$.

\textbf{Recherche des plans stables.} Soit $P$ un plan stable par $f$. $f$ est diagonalisable et donc $f_{/P}$ est un endomorphisme diagonalisable de $P$. Par suite, $P$ est engendré par deux vecteurs propres indépendants de $f$. On retrouve le plan propre de $f$ d'équation $x+2y+z=0$ et les plans engendrés par $(1,1,1)$ et un vecteur quelconque non nul du plan d'équation $x+2y+z=0$. L'équation générale d'un tel plan est $(-a-3b)x+(2a+2b)y+(b-a)z = 0$ où $(a,b)\neq(0,0)$.

\item  $\chi_A=\left|
\begin{array}{ccc}
6-X&-6&5\\
-4&-1-X&10\\
7&-6&4-X
\end{array}
\right|=(6-X)(X^2-3X+56)+4(6X+6)+7(5X-55)=-X^3+9X^2-15X-25=-(X+1)(X^2-10X+25)=-(X+1)(X-5)^2$.

$E_{-1}=\text{Vect}(10,15,4)$ et $E_5=\text{Vect}((1,1,1))$. On est dans le cas où $A$ admet une valeur propre simple et une double mais n'est pas diagonalisable. Les droites stables par $f$ sont les deux droites propres.

\textbf{Recherche des plans stables.} Soit $P$ un plan stable par $f$. Le polynôme caractéristique de $f_{/P}$ est unitaire et divise celui de $f$. Ce polynôme caractéristique est donc soit $(X-1)(X-5)$ soit $(X-5)^2$.

Dans le premier cas, $f_{/P}$ est diagonalisable et $P$ est nécessairement le plan $\text{Vect}((10,15,4))+\text{Vect}((1,1,1))$  c'est-à-dire le plan d'équation $11x -6y -5z =0$.

Dans le deuxième cas, $\chi_{f_{/P}}=(X-5)^2$ et $5$ est l'unique valeur propre de $f_{/P}$.  Le théorème de \textsc{Cayley}-\textsc{Hamilton} montre que $(f_{/P}-5Id)^2 = 0$ et donc $P$ est contenu dans $\text{Ker}(f-5Id)^2$. $\text{Ker}(f-5Id)^2$ est le plan d'équation $x=z$ qui est bien sûr stable par $f$ car $(f-5Id)^2$ commute avec $f$.
\end{enumerate}
\fincorrection
\correction{005686}
$F$ est stable par $f$ et donc $f_{/F}$ est un endomorphisme de $F$. $f$ est diagonalisable et donc il existe un polynôme $P$, scindé à racines simples, tel que $P(f)=0$. Mais alors $P(f_{/F})=0$ et on a trouvé un polynôme scindé à racines simples annulateur de $f_{/F}$. Donc $f_{/F}$ est diagonalisable.
\fincorrection
\nocorrection
\nocorrection
\nocorrection
\nocorrection
\nocorrection
\nocorrection
\correction{003578}
\begin{enumerate}
  \item 1 est valeur propre double, $d_1 = 1$.
  \item $\begin{pmatrix} 1\cr 1 \cr 1\end{pmatrix}$, $\begin{pmatrix} 2 \cr 1 \cr 2\end{pmatrix}$.
  \item $\begin{pmatrix} 1 \cr 0 \cr 0\end{pmatrix}$.
  \item $\begin{pmatrix} 1 &0 &6 \cr 0 &0 &-4 \cr 0 &0 &1 \cr\end{pmatrix}$.
  \item $X = \begin{pmatrix} (6\alpha t + \gamma)e^t           + 2\beta \cr
                            (6\alpha t + \gamma + 3\alpha)e^t + \beta  \cr
                            (6\alpha t + \gamma -  \alpha)e^t + 2\beta \cr \end{pmatrix}$.
\end{enumerate}
 \fincorrection
\nocorrection
\nocorrection
\correction{003622}
$0$ est valeur propre, se placer dans un hyperplan stable et récurer.
\fincorrection
\correction{003623}
Non. Prendre $\mathrm{mat}(u_n) = \left(\begin{smallmatrix}0&1\cr0&1/n\cr\end{smallmatrix}\right)$.
\fincorrection
\correction{003624}
Trigonaliser fortement $M$.
\fincorrection
\correction{005677}
Montrons le résultat par récurrence sur $n =\text{dim} E\geqslant 1$.

\textbullet~Si $n = 1$, c'est clair.

\textbullet~Soit $n\geqslant 1$. Supposons que deux endomorphismes d'un $\Cc$-espace de dimension $n$ qui commutent soient simultanément trigonalisables.

Soient $f$ et $g$ deux endomorphismes d'un $\Cc$-espace vectoriel de dimension $n+1$ tels que $fg = gf$.

$f$ et $g$ ont au moins un vecteur propre en commun. En effet, $f$ admet au moins une valeur propre $\lambda$. Soit $E_\lambda$ le sous-espace propre de $f$ associé à $\lambda$. $g$ commute avec $f$ et donc laisse stable $E_\lambda$. La restriction de $g$ à $E_\lambda$ est un endomorphisme de $E_\lambda$ qui est de dimension finie non nulle. Cette restriction admet donc une valeur propre et donc un vecteur propre. Ce vecteur est un vecteur propre commun à $f$ et $g$.

Commençons à construire une base de trigonalisation simultanée de $f$ et $g$. Soit $x$ un vecteur propre commun à $f$ et $g$. On complète la famille libre $(x)$ en une base $\mathcal{B}= (x,...)$ de $E$. Dans la base $\mathcal{B}$, les matrices $M$ et $N$ de $f$ et $g$ s'écrivent respectivement  $M=\left(
\begin{array}{cc}
\lambda&\times\\
0&M_1
\end{array}
\right)$ et  $N=\left(
\begin{array}{cc}
\mu&\times\\
0&N_1
\end{array}
\right)$ où $M_1$ et $N_1$ sont de format $n$. Un calcul par blocs montre que $M_1$ et $N_1$ commutent ou encore si $f_1$ et $g_1$ sont les endomorphismes de $\Cc^n$ de matrices $M_1$ et $N_1$ dans la base canonique de $\Cc^n$, $f_1$ et $g_1$ commutent. Par hypothèse de récurrence, $f_1$ et $g_1$ sont simultanément trigonalisables. Donc il existe une matrice inversible de format $n$  $P_1$ et deux matrices triangulaires supérieures de format
$n$ $T_1$ et $T_1'$ telles que $P_1^{-1}M_1P_1=T_1$ et $P_1^{-1}N_1P_1 = T_1'$.

Soit $P=\left(
\begin{array}{cc}
1&0\\
0&P_1
\end{array}
\right)$. $P$ est inversible de format $n+1$ car $\text{det}P =\text{det}P_1\neq0$ et un calcul par blocs montre que $P^{-1}MP$ et $P^{-1}NP$ sont triangulaires supérieures.

$P$ est donc la matrice de passage de la base $\mathcal{B}$ à une base de trigonalisation simultanée de $f$ et $g$.
\fincorrection
\nocorrection
\nocorrection
\nocorrection
\nocorrection
\nocorrection
\nocorrection
\nocorrection
\correction{002589}
{\it On consid\`ere la matrice suivante
$$A=\begin{pmatrix}1&-1&0 \\  1&0&-1 \\  -1&0&2\end{pmatrix}$$
et $f$ l'endomorphisme de $\R^3$ associ\'e.}
\begin{enumerate}
 \item {\it Factorisons le polyn\^ome caract\'eristique de $A$.}

On a 
\begin{align*}P_A(X)&=\begin{vmatrix}1-X&-1&0 \\  1&-X&-1 \\  -1&0&2-X\end{vmatrix} \\ 
&=(1-X)\begin{vmatrix}-X&-1 \\  0&2-X\end{vmatrix}+\begin{vmatrix}1&-1 \\  -1&2-X\end{vmatrix} \\ 
&=(1-X)(X^2-2X)+(1-X) \\ 
&=(1-X)(X^2-2X+1)=(1-X)^3.
\end{align*}
Le polyn\^ome caract\'eristique de $A$ admet une valeur propre triple $\lambda=1$.
 \item {\it D\'eterminons les sous-espaces propres et caract\'eristiques de $A$.}

La matrice $A$ admet une unique valeur propre $\lambda=1$ de multiplicit\'e $3$, le sous-espace propre associ\'e est l'espace $E_{1}=\ker(A-I)$, et on a
$$(x,y,z)\in E_{1}\iff\left\{\begin{align*}x-y&=x \\  x-z&=y \\  -x+2z&=z\end{align*}\right.
\iff\left\{\begin{align*}y&=0 \\  x&=z\end{align*}\right.$$
Le sous-espace $E_{1}$ est la droite vectorielle engendr\'ee par le vecteur $(1,0,1)$.

Le sous-espace caract\'eristique de $A$, associ\'e \`a l'unique valeur propre $\lambda=1$, est le sous-espace $N_{1}=\ker(A-I)^3$, or, compte tenu du th\'eor\`eme de Hamilton-Cayley, on sait que $P_A(A)=0$, ainsi, la matrice $(A-I)^3$ est la matrice nulle, ce qui implique $N_{1}=\R^3$, c'est donc l'espace tout entier.

 \item {\it D\'emontrons qu'il existe une base de $\R^3$ dans laquelle la matrice de $f$ s'\'ecrit}
$$B=\begin{pmatrix}1&1&0 \\  0&1&1 \\  0&0&1\end{pmatrix}.$$

Nous cherchons des vecteurs $e_1,e_2,e_3$ tels que $Ae_1=e_1$, $Ae_2=e_1+e_2$ et $Ae_3=e_2+e_3$.
Le vecteur $e_1$ appartient \`a $E_1=\ker(A-I)$, et $\ker(A-I)$ est la droite d'\'equations :

$\{y=0,x=z\}$. On d\'etermine $e_2=(x,y,z)$ tel que
$Ae_2=e_1+e_2$, on obtient le syst\`eme 
$$\begin{pmatrix}1 & -1 & 0 \\  1&0&-1 \\ -1&0&2\end{pmatrix}\begin{pmatrix}x \\  y \\  z\end{pmatrix}=\begin{pmatrix}1+x \\  1+y \\  z\end{pmatrix}\iff
\left\{\begin{align*}x-y&=1+x \\  x-z&=y \\  -x+2z&=1+z\end{align*}\right.
\iff\left\{\begin{align*}y&=-1 \\  x-z&=-1\end{align*}\right.$$
Ainsi, les vecteurs $e_1=(1,0,1)$ et $e_2=(-1,-1,0)$ conviennent. Il nous reste \`a chercher un vecteur $e_3$ tel que $Ae_3=e_2+e_3$, c'est-\`a-dire
$$\begin{pmatrix}1 & -1 & 0 \\  1&0&-1 \\ -1&0&2\end{pmatrix}\begin{pmatrix}x \\  y \\  z\end{pmatrix}=\begin{pmatrix}-1+x \\  -1+y \\  z\end{pmatrix}\iff
\left\{\begin{align*}x-y&=x-1 \\  x-z&=y-1 \\  -x+2z&=z\end{align*}\right.
\iff\left\{\begin{align*}y&=1 \\  x&=z\end{align*}\right.$$
Le vecteur $e_3=(0,1,0)$ convient. On obtient alors la matrice $P$ suivante qui est inversible et v\'erifie $A=PBP^{-1}$,
$$P=\begin{pmatrix}1 & -1 & 0 \\  0&-1&1 \\ 1&0&0\end{pmatrix}$$  

 \item  {\it D\'ecomposition de Dunford de $B$}

On a 
$$B=\begin{pmatrix}1 & 1 & 0 \\  0&1&1 \\ 0&0&1\end{pmatrix}=
\begin{pmatrix}1 & 0& 0 \\  0&1&0 \\ 0&0&1\end{pmatrix}+
\begin{pmatrix}0& 1 & 0 \\  0&0&1 \\ 0&0&0\end{pmatrix}$$
et il est clair que les deux matrices commutent car l'une est \'egale \`a $I$. 
Or, il existe un unique couple de matrices $D$ et $N$, $D$ diagonalisable
et $N$ nilpotente, telles que $B=D+N$ et $DN=ND$. Or si
$$D=\begin{pmatrix}1 & 0& 0 \\  0&1&0 \\ 0&0&1\end{pmatrix}\ \ {\hbox{et}}\ \ 
N=\begin{pmatrix}0& 1 & 0 \\  0&0&1 \\ 0&0&0\end{pmatrix},$$
On a 
$$N^2=\begin{pmatrix}0& 1 & 0 \\  0&0&1 \\ 0&0&0\end{pmatrix}\begin{pmatrix}0& 1 & 0 \\  0&0&1 \\ 0&0&0\end{pmatrix}=\begin{pmatrix}0& 0 & 1 \\  0&0&0 \\ 0&0&0\end{pmatrix}$$
et $N^3=0$. La d\'ecomposition $B=D+N$ est donc bien la d\'ecomposition de Dunford de la matrice $B$.
\end{enumerate}
\fincorrection
\correction{002602}
 Soit $A$ la matrice 
$$A=\begin{pmatrix}3&2&4 \\  -1&3&-1 \\  -2&-1&-3\end{pmatrix}$$ 
et $f$ l'endomorphisme de $\R^3$ associ\'e.
\begin{enumerate}
 \item  {\it Factorisons le polyn\^ome caract\'eristique de $A$.} (1 point)

On a 
\begin{align*}
P_A(X)&=\begin{vmatrix}3-X&2&4 \\  -1&3-X&-1 \\  -2&-1&-3-X\end{vmatrix} \\ 
&=\begin{vmatrix}-1-X&2&4 \\  0&3-X&-1 \\  1+X&-1&-3-X\end{vmatrix} \\ 
&=\begin{vmatrix}-1-X&2&4 \\ 0&3-X&-1 \\  0&1&1-X\end{vmatrix} \\ 
&=(-1-X)(X^2-4X+4)=-(X+1)(X-2)^2
\end{align*}
Les valeurs propres de la matrice $A$ sont $\lambda_1=-1$, valeur propre simple et $\lambda_2=2$, valeur propre double.
 \item  {\it D\'eterminons les sous-espaces propres et caract\'eristiques de $A$.} (2 points)
 
Le sous-espace propre associ\'e \`a la valeur propre $-1$ est le sous-espace vectoriel $E_{-1}$ d\'efini par
$$E_{-1}=\{\vec u\in\R^3,\ A\vec u=-u\}.$$
Soit $\vec u=(x,y,z)\in\R^3$, 
$$\vec u\in E_{-1}\iff \left\{\begin{align*}4x+2y+4z&=0 \\  -x+4y-z&=0 \\  -2x-y-2z&=0\end{align*}\right.
\iff\left\{ \begin{align*}2x+y+2z&=0 \\  x-4y+z&=0 \end{align*}\right.$$
L'espace $E_{-1}$ est une droite vectorielle dont un vecteur directeur est donn\'e par 
$$\vec u_1=(1,0,-1).$$
Le sous-espace propre associ\'e \`a la valeur propre $2$ est le sous-espace vectoriel $E_{2}$ d\'efini par
$$E_{2}=\{\vec u\in\R^3,\ A\vec u=2u\}.$$
Soit $\vec u=(x,y,z)\in\R^3$, 
$$\vec u\in E_{2}\iff \left\{\begin{align*}x+2y+4z&=0 \\  -x+y-z&=0 \\  -2x-y-5z&=0\end{align*}\right.
\iff \left\{\begin{align*}x+2y+4z&=0 \\  x-y+z&=0 \end{align*}\right.$$
L'espace $E_{2}$ est une droite vectorielle dont un vecteur directeur est donn\'e par 
$$\vec u_2=(2,1,-1).$$
Le sous-espace $E_2$ \'etant de dimension $1$, la matrice $A$ n'est pas diagonalisable.

Le sous-espace caract\'eristique $N_{-1}$ associ\'e \`a la valeur propre $-1$ est \'egal au sous-espace propre $E_{-1}$. Le sous-espace caract\'eristique $N_{2}$ associ\'e \`a la valeur propre $2$ est \'egal \`a
$$N_2=\ker(A-2I)^2.$$ 
D\'eterminons-le 
$$(A-2I)^2=\begin{pmatrix}1&2&4 \\  -1&1&-1 \\  -2&-1&-5\end{pmatrix}
\begin{pmatrix}1&2&4 \\  -1&1&-1 \\  -2&-1&-5\end{pmatrix}=\begin{pmatrix}-9&0&-18 \\  0&0&0 \\  9&0&18\end{pmatrix}.$$
D'o\`u $\ker(A-2I)^2=\{(x,y,z)\R^3,\ x+2z=0\}$, c'est le plan vectoriel d'\'equation $x+2z=0$.

 \item  {\it D\'emontrons qu'il existe une base de $\R^3$ dans laquelle la matrice de $f$ est} ( 1 point)
$$B=\begin{pmatrix}-1 &0& 0 \\  0&2&1 \\ 0&0&2\end{pmatrix}.$$
Les vecteurs $\vec u_1$ et $\vec u_2$, vecteurs propres associ\'es aux valeurs propres $-1$ et $2$ conviennent pour les deux premiers vecteurs de la base recherch\'ee, on va alors chercher un vecteur 
$\vec u_3\in\ker(A-2I)^2$ tel que $A\vec u_3=\vec u_2+2\vec u_3$, notons $\vec u_3=(-2z,y,z)$, on d\'etermine $y$ et $z$ tels que 
$$\begin{pmatrix}2y-2z \\ 3y+z \\ -y+z\end{pmatrix})=\begin{pmatrix}2 \\ 1 \\ -1\end{pmatrix}+\begin{pmatrix}-4z \\ 2y \\ 2z\end{pmatrix}$$
on obtient
$y+z=1$, le vecteur $\vec u_3=(0,1,0)$ convient.

{\it Trouvons une matrice $P$ inversible telle que $A=PBP^{-1}$.} (1 point)

La matrice de passage $P$ qui exprime la base $(\vec u_1, \vec u_2, \vec u_3)$ dans la base canonique de $\R^3$ r\'epond \`a la question,
$$P=\begin{pmatrix}1 &2&0 \\  0&1&1 \\ -1&-1&0\end{pmatrix}\ {\hbox{et}}\ 
P^{-1}=\begin{pmatrix}-1 &0&-2 \\  1&0&1 \\ -1&1&-1\end{pmatrix}.$$
 \item  {\it Ecrivons la d\'ecomposition de Dunford de $B$.} (1 point)

On a 
$$B=\begin{pmatrix}-1 &0& 0 \\  0&2&1 \\ 0&0&2\end{pmatrix}=\underbrace{\begin{pmatrix}-1 &0& 0 \\  0&2&0 \\ 0&0&2\end{pmatrix}}_D+
\underbrace{\begin{pmatrix}0&0& 0 \\  0&0&1 \\ 0&0&0\end{pmatrix},}_N$$
la matrice $D$ est diagonale, la matrice $N$ est nilpotente, $N^2=0$, et $ND=DN$, c'est donc bien la d\'ecomposition de Dunford de la matrice $B$.
 \item {\it Calculons $\exp B$}. (1 point)

Compte tenu de la question pr\'ec\'edente, on a $B=N+D$, avec $DN=ND$, ainsi $\exp B=\exp D\exp N$, or
$$\exp D=\begin{pmatrix}e^{-1} &0& 0 \\  0&e^2&0 \\ 0&0&e^2\end{pmatrix}\ {\hbox{et}}\ 
\exp N=I+N=\begin{pmatrix}1 &0& 0 \\  0&1&1 \\ 0&0&1\end{pmatrix}.$$
D'o\`u
$$\exp B=\begin{pmatrix}e^{-1} &0& 0 \\  0&e^2&e^2 \\ 0&0&e^2\end{pmatrix}.$$
\end{enumerate}
\fincorrection
\nocorrection
\nocorrection
\correction{003583}
On se ramène à $\lambda=0$ en rempla\c cant $f$ par~$f-\lambda\mathrm{id}$.
$\Im f$ est de dimension~$1$ stable par~$f$ donc $f_{|\Im f}$ est une homothétie, c'est l'application nulle vu $\mathrm{Spec}(f)$. On en déduit~$\Im f\subset\mathrm{Ker} f$.
Soit~$e_2\in\Im f\setminus\{0\}$, $e_3$ un antécédent de~$e_2$ par~$f$ et $e_1\in\mathrm{Ker} f$ indépendant de~$e_2$.
Alors~${\cal B} = (e_1,e_2,e_3)$ convient.
\fincorrection
\correction{005670}
Posons $\chi_f =\prod_{k=1}^{p}(\lambda_k-X)^{\alpha_k}$  où $\lambda_1$,..., $\lambda_p$ sont les valeurs propres deux à deux distinctes de $f$.

Soit $E_k'=\text{Ker}(f-\lambda_kId)^{\alpha_k}$ le sous-espace caractéristique de $f$ associé à la valeur propre $\lambda_k$, $1\leqslant k\leqslant p$. D'après le théorème de décomposition des noyaux, $E = E_1'\oplus...\oplus E_p'$. De plus, si $f_k$ est la restriction de $f$ à $E_k'$ alors $f_k$ est un endomorphisme de $E_k'$ (car $f$ et $(f-\lambda_kId)^{\alpha_k}$ commutent).

On note que $(f_k-\lambda_k)^{\alpha_k} = 0$ et donc $\lambda_k$ est l'unique valeur propre de $f_k$ car toute valeur propre de $f_k$ est racine du polynôme annulateur $(X-\lambda_k)^{\alpha_k}$.

\textbf{Existence de $d$ et $n$.} On définit $d$ par ses restrictions $d_k$ aux $E_k'$, $1\leqslant k\leqslant p$ : $d_k$ est l'homothétie de rapport $\lambda_k$. Puis on définit $n$ par $n = f-d$.

$d$ est diagonalisable car toute base de $E$ adaptée à la décomposition $E = E_1'\oplus...\oplus E_p'$ est une base de vecteurs propres de $d$. De plus,  $f= d+n$.

Soit $n_k$ la restriction de $n$ à $E_k'$. On a $n_k = f_k-\lambda_kId_{E_k'}$ et par définition de $E_k'$,  $n_k^{\alpha_k}= 0$. Mais alors, si on pose $\alpha=\text{Max}\{\alpha_1,...,\alpha_p\}$, on a $n_k^\alpha= 0$ pour tout $k$ de $\{1,...,p\}$ et donc $n^\alpha= 0$. Ainsi, $n$ est nilpotent.
Enfin, pour tout $k\in\llbracket1,p\rrbracket$, $n_k$ commute avec $d_k$ car $d_k$ est une homothétie et donc $nd = dn$.

\textbf{Unicité de $d$ et $n$.} Supposons que $f=d+n$ avec $d$ diagonalisable, $n$ nilpotent et $nd=dn$.

$d$ commute avec $n$ et donc avec $f$ car $df= d^2+dn = d^2+nd = fd$. Mais alors, $n=f-d$ commute également avec $f$. $d$ et $n$ laissent donc stables les sous-espaces caractéristiques $E_k'$, $1\leqslant k\leqslant p$ de $f$. Pour $k\in\llbracket1,p\rrbracket$, on note $d_k$ et $n_k$ les restrictions de $d$ et $n$ à $E_k'$.

Soient $k\in\llbracket1,p\rrbracket$ puis $\mu$ une valeur propre de $d_k$. D'après l'exercice \ref{ex:rou7},

\begin{center}
$\text{det}(f_k-\mu Id)=\text{det}(d_k-\mu Id+n)=\text{det}(d_k-\mu Id)=0$,
\end{center}

car $d_k-\mu Id$ n'est pas inversible. On en déduit que $f_k-\mu Id$ n'est pas inversible et donc que $\mu$ est valeur propre de $f_k$. Puisque $\lambda_k$ est l'unique valeur propre de $f_k$, on a donc $\mu=\lambda_k$. Ainsi, $\lambda_k$ est l'unique valeur propre de $d_k$ et puisque $d_k$ est diagonalisable (voir l'exercice \ref{ex:rou36}), on a nécessairement $d_k=\lambda_kId_{E_k'}$ puis $n_k=f_k-\lambda_k Id_{E_k'}$. Ceci montre l'unicité de $d$ et $n$.

\fincorrection
\nocorrection
\nocorrection
\nocorrection
\nocorrection
\nocorrection
\nocorrection
\nocorrection
\correction{001710}
Soit $A=\Big(\begin{smallmatrix}%
  1&1&0\\0&1&1\\1&0&1%
\end{smallmatrix}\Big)$.
$\chi_{A}=(2-X)(\omega-X)(\bar\omega-X)$ donc $A$ est diagonalisable sur
$\C$.

$\ker(A-2I)=\C\Big(\begin{smallmatrix}%
  1\\1\\1%
\end{smallmatrix}\Big)$  

$\Big(\begin{smallmatrix}x\\y\\z\end{smallmatrix}\Big)\in\ker(A-\omega I)%
\Leftrightarrow\Big\{\begin{smallmatrix}%
  (1-\omega)x+y=0\\(1-\omega)y+z=0\\(1-\omega)z+x=0%
\end{smallmatrix}
\Leftrightarrow\big\{\begin{smallmatrix}%
  y=(\omega-1)x\\z=(\omega-1)^{2}x%
\end{smallmatrix} 
\Leftrightarrow\big\{\begin{smallmatrix}%
  y=\omega^{2}x\\z=\omega^{4}x%
\end{smallmatrix}$ 
donc $\ker(A-\omega I)=\C\Big(\begin{smallmatrix}1\\\omega^{2}\\\omega^{4}\end{smallmatrix}\Big)$

On en déduit que $\ker(A-\bar\omega I)=\C\Big(\begin{smallmatrix}1\\\bar\omega^{2}\\\bar\omega^{4}\end{smallmatrix}\Big)$.

Ainsi en posant $P=\Big(\begin{smallmatrix}
  1&1&1\\1&\omega^{2}&\bar\omega^{2}\\1&\omega^{4}&\bar\omega^{4}\end{smallmatrix}\Big)$ on obtient
$P^{-1}AP=\Big(\begin{smallmatrix}
  2&0&0\\
  0&\omega&0\\
  0&0&\bar\omega\end{smallmatrix}\Big)$

On en déduit que les solutions sont les suites de la forme%
$\Big(\begin{smallmatrix}
  x_{n}\\
  y_{n}\\
  z_{n}  
\end{smallmatrix}\Big)=
P\Big(\begin{smallmatrix}
  2^n&0&0\\
  0&\omega^{n}&0\\
  0&0&\bar\omega^{n}\end{smallmatrix}\Big)
\Big(\begin{smallmatrix}
  a\\b\\c\end{smallmatrix}\Big)
=
P\Big(\begin{smallmatrix}2^na\\\omega^n b\\\bar\omega^nc\end{smallmatrix}\Big)
$ soit~:
$\Big\{\begin{smallmatrix}
  x_{n}=&2^na&+\omega^{n}b  &+\bar\omega^{n}c\\
  y_{n}=&2^na&+\omega^{n+2}b&+\bar\omega^{n+2}c\\
  z_{n}=&2^na&+\omega^{n+4}b&+\bar\omega^{n+4}c\\
\end{smallmatrix}$ où $a,b,c$ sont trois complexes.

En résolvant le système%
$\Big\{\begin{smallmatrix}
  a+b+c&=2\\
  a+b\omega^{2}+c\bar\omega^{2}&=1\\
  a+b\omega^{4}+c\bar\omega^{4}&=1\\
\end{smallmatrix}$
on obtient la solution particulière cherchée,  c'est la solution
associée à $a=4/3, b=c=1/3$ :
$$\left\{\begin{array}{llll}
  x_{n}=&2^{n+2}/3 &+&2/3 \cos(n\pi/3)\\
  y_{n}=&2^{n+2}/3 &+&2/3 \cos((n+2)\pi/3)\\
  z_{n}=&2^{n+2}/3 &+&2/3 \cos((n+4)\pi/3)\\
\end{array}\right..$$ 

\fincorrection
\nocorrection
\correction{001712}
  \begin{enumerate}
  \item 
    $A{}^t{A}=(a^{2}+b^{2}+c^{2}+d^{2})\mathrm{id}$. Ainsi $\det
    A*\det{}^t A=(\det A)^{2}=(a^{2}+b^{2}+c^{2}+d^{2})^{4}$ et donc
    $\det A=\pm(a^{2}+b^{2}+c^{2}+d^{2})^{2}$.
    
  \item Dans l'expression $\det A =\sum_{\sigma\in
     S_{4}}\epsilon(\sigma)\alpha_{1\sigma(1)}...\alpha_{4\sigma(4)}$ où
   les coefficients de $A$ sont notés $\alpha_{ij}$, le seul terme en
   $a^{4}$ est obtenu pour $\sigma=\mathrm{id}$, soit $\epsilon(\sigma)=+1$. On
   en déduit que $\det A=(a^{2}+b^{2}+c^{2}+d^{2})^{2}$. Pour obtenir le
   polynôme caractéristique de $A$, on remplace $a$ par $a-X$ dans $A$,
   et on calcule le déterminant. On a donc
   $\chi_{A}=((a-X)^{2}+b^{2}+c^{2}+d^{2})^{2}$

 \item
   $\forall\lambda\in\R, \chi_{A}(\lambda)>0$ car $(b,c,d)\neq(0,0,0)$.
   Donc $A$ n'a pas de valeur propre réelle, donc $A$ n'est ni
   diagonalisable ni triangularisable sur $\R$.

 \item
   $A(i\sqrt{3},1,1,1)=(1-i\sqrt{3}(i\sqrt{3},1,1,1))$ et
   $A(-1,i\sqrt{3},-1,1)=(1-i\sqrt{3}(-1,i\sqrt{3},-1,1))$. Pour la
   seconde valeur propre, qui est le conjugué de $1-i\sqrt{3}$, on
   utilise les vecteurs conjugués. Ainsi, en posant $ P=
 \begin{pmatrix}
   i\sqrt{3} & -1         &-i\sqrt{3} & -1 \\ 
   1         &  i\sqrt{3} & 1         & -i\sqrt{3}\\
   1         & -1         & 1         & -1        \\
   1         &  1         & 1         &  1
 \end{pmatrix}
   $ on a $P^{-1}AP=
   \begin{pmatrix}
     2\bar\omega &0          &0      &0\\
     0           &2\bar\omega&0      &0\\  
     0           &0          &2\omega&0\\
     0           &0          &0      &2\omega
\end{pmatrix}=D
$.

\item
  Soit $X_{n}=(u_{n},v_{n},w_{n},h_{n})$. On a $\forall n \in\N,
  X_{n+1}=AX_{n}$, d'où $\forall n\in\N, X_{n}=A^{n}X_{0}$. Or
  $A^{n}X_{0}=PD^{n}P^{-1}X_{0}$. On en déduit que $X_{n}=
  \begin{pmatrix}
   2^{n}\bar\omega^{n}i\sqrt{3}
 & -2^{n}\bar\omega^{n}         
 &-2^{n}\omega^{n}i\sqrt{3} 
 & -2^{n}\omega^{n} \\ 
   2^{n}\bar\omega^{n}         
 &  2^{n}\bar\omega^{n}i\sqrt{3} 
 & 2^{n}\omega^{n}        
 & -2^{n}\omega^{n}i\sqrt{3}\\
   2^{n}\bar\omega^{n}         
 & -2^{n}\bar\omega^{n}         
 & 2^{n}\omega^{n}        
 & -2^{n}\omega^{n}        \\
   2^{n}\bar\omega^{n}         
 &  2^{n}\bar\omega^{n}         
 & 2^{n}\omega^{n}        
 &  2^{n}\omega^{n}
  \end{pmatrix}
$. Posons $Y_{0}=P^{-1}X_{0}$. En résolvant le système $PX_{0}=Y_{0}$,
on obtient $Y_{0}=(1/2i\sqrt{3},0,-1/2i\sqrt{3},0)$, et finalement~:
$$
X_{n}=1/2i\sqrt{3}
\begin{pmatrix}
  2^{n}(\bar\omega^{n}+\omega^{n})i\sqrt{3}\\
  2^{n}(\bar\omega^{n}-\omega^{n})\\
  2^{n}(\bar\omega^{n}-\omega^{n})\\
  2^{n}(\bar\omega^{n}-\omega^{n})
\end{pmatrix}=2^{n}
\begin{pmatrix}
  \cos \frac{n\pi}{3}\\
  -\sin \frac{n\pi}{3}\\
  -\sin \frac{n\pi}{3}\\
  -\sin \frac{n\pi}{3}
\end{pmatrix}
$$
\end{enumerate}
\fincorrection
\nocorrection
\nocorrection
\nocorrection
\nocorrection
\nocorrection
\nocorrection
\nocorrection
\nocorrection
\nocorrection
\nocorrection
\nocorrection
\nocorrection
\nocorrection
\nocorrection
\nocorrection
\nocorrection
\correction{002590}
 {\it La suite de Fibonacci $0,1,1,2,3,5,8,13,...$ est la suite $(F_n)_{n\geq 0}$ d\'efinie par la relation de r\'ecurrence $F_{n+1}=F_{n}+F_{n-1}$ pour $n\geq1$, avec $F_0=0$ et $F_1=1$.}
\begin{enumerate}
 \item {\it D\'eterminons une matrice $A\in M_2(\R)$ telle que, pour tout $n\geq1$,
$$\begin{pmatrix}F_{n+1} \\  F_n\end{pmatrix}=A^n\begin{pmatrix}F_{1} \\  F_0\end{pmatrix}.$$}

On a $\displaystyle \begin{pmatrix}F_{n+1} \\  F_n\end{pmatrix}=\begin{pmatrix}F_{n}+F_{n-1} \\  F_n\end{pmatrix}= 
\begin{pmatrix}1&1 \\  1&0\end{pmatrix}\begin{pmatrix}F_{n} \\  F_{n-1}\end{pmatrix}$.

Notons, pour $n\geq1$, $X_n$ le vecteur $\displaystyle\begin{pmatrix}F_{n} \\  F_{n-1}\end{pmatrix}$ et $A$ la matrice $\displaystyle\begin{pmatrix}1&1 \\  1&0\end{pmatrix}$. Nous allons d\'emontrer, par r\'ecurrence sur $n$ que pour tout $n\geq1$, on a $X_n=A^nX_0$, c'est clairement vrai pour $n=1$, supposons que ce soit vrai pour un $n$ arbitrairement fix\'e, on a alors
$$X_{n+1}=AX_n=AA^nX_0=A^{n+1}X_0,$$
d'o\`u le r\'esultat.

 \item {\it Montrons que $A$ admet deux valeurs propres r\'eelles distinctes que l'on note $\lambda_1$ et $\lambda_2$ avec $\lambda_1<\lambda_2$.}

On a
$$P_A(X)=\begin{vmatrix}1-X&1 \\ 1&-X\end{vmatrix}=X^2-X-1.$$
Le discriminant $\Delta=5>0$, le polyn\^ome caract\'eristique admet deux racines distinctes
$$\lambda_1={\frac{1-\sqrt{5}}{2}}<\lambda_2={\frac{1+\sqrt{5}}{2}}.$$

 \item {\it Trouvons des vecteurs propres $\varepsilon_1$ et $\varepsilon_2$ associ\'es aux valeurs propres $\lambda_1$ et $\lambda_2$, sous la forme 
$\displaystyle\begin{pmatrix}\alpha \\  1\end{pmatrix}$, avec $\alpha\in\R$.}

Posons $\varepsilon_1=\begin{pmatrix}\alpha \\  1\end{pmatrix}$ et calculons $\alpha$ tel que
$A\varepsilon_1=\lambda_1\varepsilon_1$, c'est-\`a-dire
$$\begin{pmatrix}1&1 \\  1&0\end{pmatrix}\begin{pmatrix}\alpha \\  1\end{pmatrix}=\lambda_1\begin{pmatrix}\alpha \\  1\end{pmatrix}$$
ce qui \'equivaut \`a
$$\left\{\begin{align*}\alpha+1&=\lambda_1\alpha \\ \alpha&=\lambda_1\end{align*}\right.\iff\alpha=\lambda_1$$
car $\lambda_1^2-\lambda_1-1=0$, on a donc 
$\displaystyle \varepsilon_1=\begin{pmatrix}\lambda_1 \\ 1\end{pmatrix}$ et, de m\^eme,
$\displaystyle \varepsilon_2=\begin{pmatrix}\lambda_2 \\ 1\end{pmatrix}$

 \item{\it D\'eterminons les coordonn\'ees du vecteur $\displaystyle\begin{pmatrix}F_1 \\  F_0\end{pmatrix}$ dans la base 
$(\varepsilon_1,\varepsilon_2)$, on les note $x_1$ et $x_2$.}

On cherche $x_1$ et $x_2$ tels que 
$$\begin{pmatrix}F_1 \\  F_0\end{pmatrix}=\begin{pmatrix}1 \\  0\end{pmatrix}
=x_1\varepsilon_1+x_2\varepsilon_2=x_1\begin{pmatrix}\lambda_1 \\ 1\end{pmatrix}+
x_2\begin{pmatrix}\lambda_2 \\ 1\end{pmatrix},$$
$x_1$ et $x_2$ sont donc solutions du syst\`eme
$$\left\{\begin{align*}x_1\lambda_1+x_2\lambda_2&=1 \\  x_1+x_2=0\end{align*}\right.\iff
\left\{\begin{align*}x_1(\lambda_1-\lambda_2)&=1 \\  x_2=-x_1\end{align*}\right.\iff
\left\{\begin{align*}x_1&={\frac{1}{\lambda_1-\lambda_2}}={\frac{-1}{\sqrt{5}}} \\  
x_2&=-{\frac{1}{\lambda_1-\lambda_2}}={\frac{1}{\sqrt{5}}}\end{align*}\right.$$

 \item {\it Montrons que $\displaystyle\begin{pmatrix}F_{n+1} \\  F_n\end{pmatrix}
=\lambda_1^nx_1\varepsilon_1+\lambda_2^nx_2\varepsilon_2$.}

Les vecteurs $\varepsilon_1$ et $\varepsilon_2$ \'etant vecteurs propres de $A$, on a 
$A\varepsilon_1=\lambda_1\varepsilon_1$ et $A\varepsilon_2=\lambda_2\varepsilon_2$, ainsi par r\'ecurrence, on a, pour tout $n\geq1$, 
$A^n\varepsilon_1=\lambda_1^n\varepsilon_1$ et $A^n\varepsilon_2=\lambda_2^n \varepsilon_2$. Ainsi, 
$$\begin{pmatrix}F_{n+1} \\  F_n\end{pmatrix}
=A^n\begin{pmatrix}F_{1} \\  F_0\end{pmatrix}=A^n(x_1\varepsilon_1+x_2\varepsilon_2)
=x_1A^n\varepsilon_1+x_2A^n\varepsilon_2=\lambda_1^nx_1\varepsilon_1+\lambda_2^nx_2\varepsilon_2.$$

{\it On en d\'eduit que 
$$F_n={\frac{\lambda_1^n}{\lambda_1-\lambda_2}}-{\frac{\lambda_2^n}{\lambda_1-\lambda_2}}\ .$$}

On a montr\'e que $\displaystyle \varepsilon_1=\begin{pmatrix}\lambda_1 \\ 1\end{pmatrix}$ et $\displaystyle \varepsilon_2=\begin{pmatrix}\lambda_2 \\ 1\end{pmatrix}$, on a donc, 
$$\begin{pmatrix}F_{n+1} \\  F_n\end{pmatrix}
=\lambda_1^n\left({\frac{1}{\lambda_1-\lambda_2}}\right)\begin{pmatrix}\lambda_1 \\ 1\end{pmatrix}+
\lambda_2^n\left({\frac{-1}{\lambda_1-\lambda_2}}\right)\begin{pmatrix}\lambda_2 \\ 1\end{pmatrix},$$
d'o\`u le r\'esultat
$$F_n={\frac{\lambda_1^n}{\lambda_1-\lambda_2}}-{\frac{\lambda_2^n}{\lambda_1-\lambda_2}}\ .$$

 \item {\it Donnons un \'equivalent de $F_n$ lorsque $n$ tend vers $+\infty$.}

On remarque que $|\lambda_1|<1$ et $|\lambda_2|>1$ ainsi, lorsque $n$ tend vers l'infini, $\lambda_1^n$ tend vers $0$ et $\lambda_2^n$ tend vers $+\infty$. On a donc
$$\lim_{n\rightarrow+\infty} (\lambda_2-\lambda_1){\frac{F_n}{\lambda_2^n}}=
\lim_{n\rightarrow+\infty}-{\frac{\lambda_1^n}{\lambda_2^n}}+1=1.$$
Ce qui prouve que $\displaystyle{\frac{\lambda_2^n}{\lambda_2-\lambda_1}}$ est un \'equivalent de $F_n$ lorsque $n$ tend vers $+\infty$.

\end{enumerate}
\fincorrection
\correction{002597}
Soit $A$ la matrice 
$$A=\begin{pmatrix}2&0&-1 \\  -1&1&1 \\  0&-1&0\end{pmatrix}$$ 
et $f$ l'endomorphisme de $\R^3$ associ\'e.
\begin{enumerate}
 \item  {\it Factorisons le polyn\^ome caract\'eristique de $A$}.

Le polyn\^ome caract\'eristique de $A$ est le polyn\^ome 
$$\begin{align*}P_A(X)=\begin{vmatrix}2-X&0&-1 \\ -1&1-X&1 \\  0&-1&-X\end{vmatrix}
&=(2-X)\begin{vmatrix}1-X&1 \\ -1&-X\end{vmatrix}-\begin{vmatrix}-1&1-X \\  0&-1\end{vmatrix} \\ 
&=(2-X)(X^2-X+1)-1 \\ 
&=-X^3+3X^2-3X+1=(1-X)^3\end{align*}$$

 \item {\it D\'eterminons les sous-espaces propres et caract\'eristiques de $A$}.

La matrice $A$ admet une unique valeur propre $\lambda=1$, comme $A\neq I$, elle n'est pas diagonalisable. Son sous-espace caract\'eristique est \'egal \`a $\ker(A-I_3)^3=\R^3$. En effet, d'apr\`es le th\'eor\`eme de Hamilton-Cayley, on a $P_A(A)=0$, c'est-\`a-dire $(A-I_3)^3=0$. Son sous-espace propre est \'egal \`a $\ker(A-I_3)$.
$$\ker(A-I_3)=\{\vec u\in\R^3/\ A\vec u=\vec u\}=\{(x,y,z)\in\R^3/\ x=z=-y\}.$$
C'est la droite vectorielle engendr\'ee par le vecteur $\vec u_1=(1,-1,1)$.
 \item {\it D\'emontrons qu'il existe une base de $\R^3$ dans laquelle la matrice de $f$ est 
$$B=\begin{pmatrix}1 & 1 & 0 \\  0&1&1 \\ 0&0&1\end{pmatrix}$$
et trouver une matrice $P$ inversible telle que $A=PBP^{-1}$}.

Le vecteur $\vec u_1$ v\'erifie $A\vec u_1=u_1$, on cherche un vecteur $\vec u_2=(x,y,z)$ tel que 
$A\vec u_2=\vec u_1+\vec u_2$.
$$\begin{align*}A\vec u_2=\vec u_1+\vec u_2&\iff\begin{pmatrix}2&0&-1 \\  -1&1&1 \\  0&-1&0\end{pmatrix}
\begin{pmatrix}x \\  y \\  z\end{pmatrix}=\begin{pmatrix}1+x \\  -1+y \\  1+z\end{pmatrix} \\ 
&\iff\begin{pmatrix}2x-z \\  -x+y+z \\  -y\end{pmatrix}=\begin{pmatrix}1+x \\  -1+y \\  1+z\end{pmatrix}\end{align*}$$
On obtient donc $z-x=z+y=-1$, le vecteur $\vec u_2=(1,-1,0)$ convient. On cherche alors un vecteur $\vec u_3=(x,y,z)$ tel que $A\vec u_3=\vec u_2+\vec u_3$.
$$\begin{align*}A\vec u_3=\vec u_2+\vec u_3\iff
\begin{pmatrix}2x-z \\  -x+y+z \\  -y\end{pmatrix}=\begin{pmatrix}1+x \\ -1+y \\  z\end{pmatrix}\end{align*}$$
on obtient alors $x+y=x-z=1$. Le vecteur $\vec u_3=(1,0,0)$ convient.
Ainsi, dans la base $(\vec u_1, \vec u_2, \vec u_3)$ la matrice de $f$ est \'egale \`a $B$. La matrice $P$ cherch\'ee est 
$$P=\begin{pmatrix}1& 1 & 1 \\ -1&-1&0 \\ 1&0&0\end{pmatrix},$$
on a bien $A=PBP^{-1}$ et $B=P^{-1}AP$.
 \item  {\it Ecrivons la d\'ecomposition de Dunford de $B$}.

On a 
$$B=\begin{pmatrix}1 & 1 & 0 \\  0&1&1 \\ 0&0&1\end{pmatrix}=\begin{pmatrix}1 & 0 & 0 \\  0&1&0 \\ 0&0&1\end{pmatrix}+
\begin{pmatrix}0& 1 & 0 \\  0&0&1 \\ 0&0&0\end{pmatrix}=I_3+N$$
Si $N=\begin{pmatrix}0& 1 & 0 \\  0&0&1 \\ 0&0&0\end{pmatrix}$ alors $N^2=\begin{pmatrix}0& 0 & 1 \\  0&0&0 \\ 0&0&0\end{pmatrix}$ et $N^3=0$. La matrice $I_3$ est diagonale, la matrice $N$ est nilpotente, les matrices $I_3$ et $N$ commutent, c'est donc bien la d\'ecomposition de Dunford de $B$.

 \item {\it Pour $t\in\R$, calculons} $\exp tB$. 

On a, pour $t\in\R$, $tB=tI_3+tN$, ainsi $\exp tB=\exp(tI_3+tN)=(\exp tI_3)(\exp tN)$ car les matrices $tI_3$ et $tN$ commutent. Par ailleurs, $\exp tI_3=e^tI_3$ et $\exp tN=I+tN+{\frac{t^2}{2}}N^2$. On a donc
$$\exp tB=e^t\begin{pmatrix}1&t&t^2/2 \\ 0&1&t \\ 0&0&1\end{pmatrix}.$$

 \item {\it Donnons les solutions des syst\`emes diff\'erentiels $Y'=BY$ et $X'=AX$.}

Int\'egrons le syst\`eme $Y'=BY$, sa solution g\'en\'erale s'\'ecrit
$$Y(t)=(\exp tB) Y_0,$$
o\`u $Y_0$ est un vecteur de $\R^3$. 

Int\'egrons alors le syst\`eme $X'=AX$. Remarquons que si $PY$ est solution de $X'=AX$, on a 
$$(PY)'=A(PY)\iff PY'=APY\iff Y'=P^{-1}APY\iff Y'=BY$$
ainsi $Y$ est solution de $Y'=BY$, la solution g\'en\'erale du syst\`eme $X'=AX$ s\'ecrit donc
$$\begin{align*}X(t)=P(\exp tB)Y_0&=e^t\begin{pmatrix}1& 1 &1 \\  -1&-1&0 \\ 1&0&0\end{pmatrix}\begin{pmatrix}1&t&t^2/2 \\ 0&1&t \\ 0&0&1\end{pmatrix}Y_0 \\  &=e^t
\begin{pmatrix}1& t+1& t^2/2+t+1 \\  -1&-t-1&-t^2/2-t \\  1&t&t^2/2\end{pmatrix}\begin{pmatrix}a \\  b \\  c\end{pmatrix}\end{align*}$$
o\`u $Y_0=(a,b,c)$ est un vecteur de $\R^3$. Ou encore
$$\left\{\begin{align*}
x(t)&=e^t(a+b(t+1)+c(t^2/2+t+1)) \\  y(t)&=e^t(-a-b(t+1)-c(t^2/2+t)) \\  z(t)&=e^t(a+bt+ct^2/2)
\end{align*}\right.$$ ou encore

$$\begin{pmatrix}x(t) \\  y(t) \\  z(t)\end{pmatrix}=a\begin{pmatrix}e^t \\  -e^t \\  e^t\end{pmatrix}+
b\begin{pmatrix}e^t(t+1) \\  -e^t(t+1) \\  te^t\end{pmatrix}+c\begin{pmatrix}e^t(t^2/2+t+1) \\  -e^t(t^2/2+t) \\  e^t t^2/2\end{pmatrix}$$
avec $(a,b,c)\in\R^3$.
\end{enumerate}
\fincorrection
\correction{002598}
\begin{enumerate}
 \item  On note $(\vec e_1, \vec e_2, \vec e_3)$ la base canonique de $\R^3$. Soit $A$ la matrice 
$$A=\begin{pmatrix}1&0&0 \\  0&2&0 \\  0&0&3\end{pmatrix}.$$
{\it Donnons sans calcul les valeurs propres de $A$ et une base de vecteurs propres.}

La matrice $A$ est diagonale dans la base canonique de $\R^3$, on a $A\vec e_1=\vec e_1$, $A\vec e_2=2\vec e_2$ et $A\vec e_3=3\vec e_3$. Les valeurs propres de $A$ sont les r\'eels $1$, $2$ et $3$ et les sous-espaces propres associ\'es sont les droites vectorielles engendr\'ees respectivement par $\vec e_1$, $\vec e_2$ et $\vec e_3$.

 \item On cherche \`a d\'eterminer, s'il en existe, les matrices $B$ telles que $\exp B=A$.
   \begin{enumerate}
   \item  {\it Montrons que si $A=\exp B$, alors $AB=BA$.}

On suppose qu'il existe $B$ telle que $A=\exp B$. On a alors, par d\'efinition, 
$$A=\sum_{k=0}^{+\infty}{\frac{1}{k!}}B^k,\ 
{\hbox{d'o\`u}}\  AB=BA=\sum_{k=0}^{+\infty}{\frac{1}{k!}}B^{k+1}$$
   \item {\it On en d\'eduit que la base $(\vec e_1,\vec e_2, \vec e_3)$ est une base de vecteurs propres de  B.}

On a $(BA)\vec e_1=B(A\vec e_1)=B\vec e_1$, mais $BA=AB$, on a donc $B\vec e_1=(AB)\vec e_1=A(B\vec e_1)$. Ce qui prouve que $B\vec e_1$ est un vecteur propre de $A$ associ\'e \`a la valeur propre $1$, il est donc colin\' eaire \`a $\vec e_1$ , ainsi, $\vec e_1$ est bien un vecteur propre de $B$. De m\^eme, $BA\vec e_2=2B\vec e_2=AB\vec e_2$ donc $B\vec e_2$ est colin\'eaire \`a $\vec e_2$ et $\vec e_2$ est un vecteur propre de $B$. Et aussi, $BA\vec e_3=3B\vec e_3=AB\vec e_3$ d'o\`u $B\vec e_3$ colin\'eaire \`a $\vec e_3$ et $\vec e_3$ vecteur propre de $B$.
   \item {\it D\'eterminons les matrices  $B\in M_3(\R)$ telles que $\exp B=A$.}

Les vecteurs $\vec e_1, \vec e_2, \vec e_3$ \'etant vecteurs propres de $B$, la matrice $B$ est diagonale dans la base canonique, il existe donc des r\'eels $\lambda_1,\lambda_2$ et $\lambda_3$ tels que 
$$B=\begin{pmatrix}\lambda_1&0&0 \\  0&\lambda_2&0 \\  0&0&\lambda_3\end{pmatrix},\ {\hbox{ainsi}}\ \exp B=
\begin{pmatrix}e^{\lambda_1}&0&0 \\  0& e^{\lambda_2}&0 \\  0&0&e^{\lambda_3}\end{pmatrix}=\begin{pmatrix}1&0&0 \\  0&2&0 \\  0&0&3\end{pmatrix}$$
ce qui implique $e^{\lambda_1}=1$, $e^{\lambda_2}=2$ et $e^{\lambda_3}=3$ et donc $\lambda_1=\ln 1=0$, $\lambda_2=\ln2$ et $\lambda_3=\ln3$ d'o\`u l'existence d'une unique matrice $B$ telle que $\exp B=A$, c'est la matrice
$$B=\begin{pmatrix}0&0&0 \\  0&\ln2&0 \\  0&0&\ln3\end{pmatrix}$$ 
  \end{enumerate}
 \item Soit la matrice $C$, 
$$C=\begin{pmatrix}0&1&0 \\  0&0&1 \\  0&0&0\end{pmatrix}.$$
{\it Montrons qu'il n'existe pas de matrice $D\in M_3(\R)$ telle que $C=\exp D$.}

Quelque soit la matrice $D$, la matrice $\exp D$ est inversible d'inverse $\exp (-D)$, or, il est clair que la matrice $C$ n'est pas inversible, son d\'eterminant est nul, ainsi, il n'existe pas de matrice $D$ telle que $C=\exp D$.

 \item  {\it Calculons le polyn\^ome caract\'eristique et le polyn\^ome minimal de $C$.}

Soit $P_C(X)$ le polyn\^ome caract\'eristique de $C$, on a 
$$P_C(X)=\begin{vmatrix}-X&1&0 \\  0&-X&1 \\  0&0&-X\end{vmatrix}=-X\begin{vmatrix}-X&1 \\ 0&-X\end{vmatrix}=-X^3.$$
Le polyn\^ome minimal $Q_C$ de $C$ est unitaire, divise son polyn\^ome caract\'eristique $P_C$ et s'annule en $C$, or
$C^2=\begin{pmatrix}0&0&1 \\  0&0&0 \\  0&0&0\end{pmatrix}\neq 0$, mais $C^3=0$, d'o\`u $Q_C(X)=X^3$.
 \item Supposons qu'il existe une matrice $E\in M_3(\R)$ telle que $E^2=C$. Notons $Q_E(X)$ son polyn\^ome minimal et $Q_C(X)$ le polyn\^ome minimal de $C$.

    \begin{enumerate}
   \item  {\it Montrons que $Q_E(X)$ divise $Q_C(X^2)$.}

On a $Q_C(C)=Q_C(E^2)=0$, ainsi le polyn\^ome $S(X)=Q_C(X^2)$ s'annule en $E$, ce qui prouve que $Q_E(X)$, polyn\^ome minimal de $E$, divise $Q_C(X^2)$.
   \item {\it On en d\'eduit que $E^3=0$ et que $C^2=0$.}

Le polyn\^ome minimal de $E$ divise $Q_C(X^2)=-X^6$, or son degr\'e est inf\'erieur ou \'egal \`a $3$, par ailleurs on suppose $E^2=C$ donc $E^2\neq 0$ ainsi, on a bien $Q_E(X)=X^3$ et $E^3=0$, or, $E^3=0$ implique $E^4=0$ et, comme $E^4=C^2$, on a $C^2=0$.

   \item {\it On d\'eduit de ce qui pr\'ec\`ede qu'il n'existe pas de matrice $E$ telle que $E^2=C$.}

Si une telle matrice existe, alors on a vu qu'elle v\'erifie $E^3=0$, ainsi on a $E^4=0$, or $E^4=C^2\neq0$, d'o\`u  la conclusion.

  \end{enumerate}
 \item Soient $F$ et $G$ des matrices de $M_3(\R)$ telles que $F=\exp G$. {\it D\'emontrons que pour tout $n\in \N^*$, il existe une matrice $H$ telle que $H^n=F$}.

Soit $F=\exp G$ et $n\in\N^*$, posons $H=\exp{\frac{1}{n}}G$, on a $$H^n=\left(\exp{\frac{1}{n}}G\right)^n=\exp{\frac{n}{n}}G=\exp G=F.$$
\end{enumerate}
\fincorrection
\correction{003585}
\begin{enumerate}
  \item $A^{2k}=I$, $A^{2k+1}=A$.
\end{enumerate}
 \fincorrection
\correction{003586}
3. $\alpha_n = -\frac13 + \frac {2^n}4 + \frac {(-2)^n}{12}$,
             $\beta_n  = \frac {2^n - (-2)^n}4$,
             $\gamma_n = \frac43 - \frac {2^n}2 + \frac {(-2)^n}{6}$.
 \fincorrection
\correction{003587}
2. $P = \begin{pmatrix} 1 &1 &1 \cr 1 & 2 & 3 \cr 1 &4 &9 \cr \end{pmatrix}$,
             $D = \begin{pmatrix} 1 &2 &3 \end{pmatrix}$.\par
             $2u_n = (6-6.2^n+2.3^n)u_0 + (-5+8.2^n-3.3^n)u_1 + (1-2.2^n+3^n)u_2$.
 \fincorrection
\correction{003588}
\begin{enumerate}
  \item Le polynôme minimal de $f$ est de degré supérieur ou égal à~$n$ et n'a
pas de diviseurs non triviaux. Donc $\dim E = 1$ et $f$ est une homothétie
si $ K=\C$. Si $ K=\R$ on peut aussi avoir $\dim E = 2$ et $f$ n'a pas
de valeurs propres réelles.
\end{enumerate}
 \fincorrection
\correction{003589}
\begin{enumerate}
  \item Diagonaliser ${}^tM  \Rightarrow  y_n -\frac 32x_n = \text{cste}$.
  \item $y_n-x_n = 2^n(y_0-x_0)$ donc si $y_0\ne x_0$ alors $M_n \to\infty$
sinon la suite est constante.
  \item $\frac32$ si $y_0\ne x_0$.
\end{enumerate}
 \fincorrection
\nocorrection
\correction{003591}
2. $X = \frac 12\begin{pmatrix} a+b &b-a \cr b-a &a+b \cr \end{pmatrix}$,
             $Y = \frac 1b\begin{pmatrix} 1 &1 \cr 1 &1 \cr \end{pmatrix}$ ou l'inverse.
 \fincorrection
\correction{003592}
$M = \pm\begin{pmatrix} 3 &0 &0 \cr 1/5 &\pm2 &0 \cr 7/30 &\pm1/3 &\pm1 \cr \end{pmatrix}$
      ou $M = \pm\begin{pmatrix} 3 &0 &0 \cr 1   &\mp2 &0 \cr 1/2  &\mp1   &\pm1 \cr \end{pmatrix}$.
\fincorrection
\correction{003593}
$A=PDP^{-1}$ avec $P = \begin{pmatrix}1 &1 &0 \cr
                                          2 &0 &1 \cr
                                          3 &-4&4 \cr\end{pmatrix}$
         et $D = \begin{pmatrix} 0 &0 &0 \cr
                           0 &1 &0 \cr
                           0 &0 &1 \cr\end{pmatrix}$.
         On prend $B = PMP^{-1}$ avec $M =  \begin{pmatrix} 0 &0 &0\cr
                                                      0 &0 &1\cr
                                                      0 &1 &0 \cr\end{pmatrix}$.
\fincorrection
\correction{003594}
3 valeurs propres distinctes, $M$ est diagonalisable et son
	     commutant est l'ensemble des polynômes en $M$ :
	     $aI + bM + cM^2$, $a,b,c\in K$.
\fincorrection
\correction{003595}
\begin{enumerate}
  \item Par similitude on se ramène aux cas~:
    $A=\left(\begin{smallmatrix}\lambda&0\cr0&\lambda\cr\end{smallmatrix}\right)$, $C(A) = \mathcal{M}_2(\C)$ ou
    $A=\left(\begin{smallmatrix}\lambda&0\cr0&\mu\cr\end{smallmatrix}\right)$, $C(A) = \C[A]$ ou
    $A=\left(\begin{smallmatrix}\lambda&1\cr0&\lambda\cr\end{smallmatrix}\right)$, $C(A) = \C[A]$.
  \item Si $A$ est diagonalisable de valeurs propres $\lambda_i$
    avec les multiplicités $n_i$ alors ${\dim(C(A)) = \sum n_i^2 \ge n}$.
    
    Dans le cas général, soit $(A_k)$ une suite de matrices diagonalisables
    convergeant vers $A$ et $(C_k^1,\dots,C_k^n)$ une suite de $n$-uplets
    de matrices commutant avec $A_k$ telles que $(C_k^1,\dots,C_k^n)$
    est une famille orthonormale pour un produit scalaire quelconque choisi
    sur $\mathcal{M}_n(\C)$. Par compacité il existe une sous-suite convergente, donc
    $n$ matrices $C_\infty^i$ formant une famille orthonormale et commutant
    avec~$A$ d'où $\dim(C(A))\ge n$.
\end{enumerate}
 \fincorrection
\correction{003596}
Soit $d_n(i,j)$ le nombre de chemins de longueur~$n$ allant du sommet
$i$ au sommet $j$. $j$ admet trois voisins $k_1,k_2,k_3$ et l'on a~:
$d_n(i,j) = d_{n-1}(i,k_1) + d_{n-1}(i,k_2) + d_{n-1}(i,k_3)$.
On numérote les sommets de $0$ à $7$ de sorte que les voisins du sommet~$i$
sont les sommets $i+1 \bmod 8$, $i+2 \bmod 8$ et $i+4 \bmod 8$.
Le vecteur $d_n = (d_n(0,0),\dots,d_n(0,7))$ vérifie la relation de récurrence
$d_n = A d_{n-1}$ où $A$ est la matrice suivante ($.$ désigne $0$)~:
$$A = \begin{pmatrix}. &1 &1 &. &1 &. &. &.\cr
               1 &. &. &1 &. &1 &. &.\cr
               1 &. &. &1 &. &. &1 &.\cr
               . &1 &1 &. &. &. &. &1\cr
               1 &. &. &. &. &1 &1 &.\cr
               . &1 &. &. &1 &. &. &1\cr
               . &. &1 &. &1 &. &. &1\cr
               . &. &. &1 &. &1 &1 &.\cr\end{pmatrix}
 = \begin{pmatrix}B &I_4\\ I_4&B\\\end{pmatrix}
 =  P\begin{pmatrix}B+I_4 &0\\ 0&B-I_4\end{pmatrix}P^{-1}$$
avec $$B = \begin{pmatrix}. &1 &1 &.\cr
                   1 &. &. &1\cr
                   1 &. &. &1\cr
                   . &1 &1 &.\cr\end{pmatrix}
\text{ et } P = \begin{pmatrix}I_4 &I_4\\ I_4&-I_4\\ \end{pmatrix}.$$
De même,
$$B\pm I_4 = \begin{pmatrix}C\pm I_2&I_2\\ I_2&C\pm I_2\\\end{pmatrix}
= Q\begin{pmatrix}C\pm I_2+I_2 &0\cr 0&C\pm I_2 - I_2\end{pmatrix}Q^{-1}$$
et enfin,
$$C\pm I_2\pm I_2 = \begin{pmatrix}\pm I_1\pm I_1&I_1\cr I_1&\pm I_1\pm I_1\cr\end{pmatrix}
= R\begin{pmatrix}\pm I_1\pm I_1+I_1 &0\cr 0&\pm I_1\pm I_1 - I_1\end{pmatrix}R^{-1}.$$
Donc $A$ est diagonalisable de valeurs propres $-3,-1,1,3$ et on
peut certainement terminer les calculs pour obtenir $d_n = A^nd_0$.
\fincorrection
\correction{005651}
\textbf{1ère solution.} $A =2J-I_3$ où $J=\left(
\begin{array}{ccc}
1&1&1\\
1&1&1\\
1&1&1
\end{array}
\right)$. On a $J^2 = 3J$ et plus généralement $\forall k\in\Nn^*$, $J^k =3^{k-1}J$.

Soit $n\in\Nn^*$. Puisque les matrices $2J$ et $-I$ commutent, la formule du binôme de \textsc{Newton} permet d'écrire

\begin{align*}\ensuremath
A^n&=(2J-I)^n =(-I)^n+\sum_{k=1}^{n}\dbinom{n}{k}(2J)^k(-I)^{n-k}=(-1)^nI+\left(\sum_{k=1}^{n}\dbinom{n}{k}2^k3^{k-1}(-1)^{n-k}\right)J\\
 &=(-1)^nI+\frac{1}{3}\left(\sum_{k=1}^{n}\dbinom{n}{k}6^k(-1)^{n-k}\right)J=(-1)^nI+\frac{1}{3}((6-1)^n-(-1)^n)J\\
 &=\frac{1}{3}\left(\begin{array}{ccc}
5^n+2(-1)^n&5^n-(-1)^n&5^n-(-1)^n\\
5^n-(-1)^n&5^n+2(-1)^n&5^n-(-1)^n\\
5^n-(-1)^n&5^n-(-1)^n&5^n+2(-1)^n
\end{array}
\right),
\end{align*}

ce qui reste vrai quand $n=0$.

Soit de nouveau $n\in\Nn^*$.

\begin{align*}\ensuremath
((-1)^nI+\frac{1}{3}(5^n-(-1)^n)J)&\times((-1)^{-n}I+\frac{1}{3}(5^{-n}-(-1)^{-n})J)\\
 &=I+\frac{1}{3}((-5)^n-1+(-5)^{-n}-1)J+\frac{1}{9}(1-(-5)^n-(-5)^{-n}+1)J^2\\
 &=I+\frac{1}{3}((-5)^n-1+(-5)^{-n}-1)J+\frac{3}{9}(1-(-5)^n-(-5)^{-n}+1)J=I,
\end{align*}

et donc

\begin{center}
$
A^{-n}=\frac{1}{3}\left(\begin{array}{ccc}
5^{-n}+2(-1)^{-n}&5^{-n}-(-1)^{-n}&5^{-n}-(-1)^{-n}\\
5^{-n}-(-1)^{-n}&5^{-n}+2(-1)^{-n}&5^{-n}-(-1)^{-n}\\
5^{-n}-(-1)^{-n}&5^{-n}-(-1)^{-n}&5^{-n}+2(-1)^{-n}
\end{array}
\right)$.
\end{center}

Finalement

\begin{center}
\shadowbox{
$\forall n\in\Zz$, $A^n=\frac{1}{3}\left(\begin{array}{ccc}
5^n+2(-1)^n&5^n-(-1)^n&5^n-(-1)^n\\
5^n-(-1)^n&5^n+2(-1)^n&5^n-(-1)^n\\
5^n-(-1)^n&5^n-(-1)^n&5^n+2(-1)^n
\end{array}
\right)$.
}
\end{center}

\textbf{2ème solution.} Puisque $\text{rg}(A+I)=1$, $\text{dim}(\text{Ker}(A+I))=2$ et $-1$ est valeur propre de $A$ d'ordre au moins $2$. La troisième valeur propre $\lambda$ est fournie par la trace : $\lambda-1-1=3$ et donc $\lambda=5$. Par suite, $\chi_A=-(X+1)^2(X-5)$.

De plus, $\left(\begin{array}{c}
x\\
y\\
z
\end{array}
\right)\in E_{-1}\Leftrightarrow x+y+z=0$ et donc $E_{-1}=\text{Vect}(e_1,e_2)$ où $e_1=\left(\begin{array}{c}
1\\
-1\\
0
\end{array}
\right)$ et $e_2=\left(\begin{array}{c}
1\\
0\\
-1
\end{array}
\right)$.

De même,  $\left(\begin{array}{c}
x\\
y\\
z
\end{array}
\right)\in E_{-1}\Leftrightarrow x=y=z$ et $E_5=\text{Vect}(e_3)$ où $e_3=\left(\begin{array}{c}
1\\
1\\
1
\end{array}
\right)$.

On pose $P=\left(
\begin{array}{ccc}
1&1&1\\
-1&0&1\\
0&-1&1
\end{array}
\right)$ et $D=\text{diag}(-1,-1,5)$ et on a $A =PDP^{-1}$.

Calcul de $P^{-1}$. Soit $(i,j,k)$ la base canonique de $\Rr^3$.

\begin{center}
$\left\{
\begin{array}{l}
e_1=i-j\\
e_2=i-k\\
e_3=i+j+k
\end{array}
\right.\Leftrightarrow\left\{
\begin{array}{l}
j=i-e_1\\
k=i-e_2\\
e_3=i+i-e_1+i-e_2
\end{array}
\right.\Leftrightarrow\left\{
\begin{array}{l}
i=\frac{1}{3}(e_1+e_2+e_3)\\
\rule[-4mm]{0mm}{10mm}j=\frac{1}{3}(-2e_1+e_2+e_3)\\
k=\frac{1}{3}(e_1-2e_2+e_3)
\end{array}
\right.$
\end{center}

et donc $P^{-1}=\frac{1}{3}\left(
\begin{array}{ccc}
1&-2&1\\
1&1&-2\\
1&1&1
\end{array}
\right)$. Soit alors $n\in\Zz$.

\begin{align*}\ensuremath 
A^n&= PD^nP^{-1}=\frac{1}{3}\left(
\begin{array}{ccc}
1&1&1\\
-1&0&1\\
0&-1&1
\end{array}
\right)\left(
\begin{array}{ccc}
(-1)^n&0&0\\
0&(-1)^n&0\\
0&0&5^n
\end{array}
\right)\left(
\begin{array}{ccc}
1&-2&1\\
1&1&-2\\
1&1&1
\end{array}
\right)\\
 &=\frac{1}{3}\left(
\begin{array}{ccc}
(-1)^n&(-1)^n&5^n\\
-(-1)^n&0&5^n\\
0&-(-1)^n&5^n
\end{array}
\right)\left(
\begin{array}{ccc}
1&-2&1\\
1&1&-2\\
1&1&1
\end{array}
\right)=\frac{1}{3}\left(\begin{array}{ccc}
5^n+2(-1)^n&5^n-(-1)^n&5^n-(-1)^n\\
5^n-(-1)^n&5^n+2(-1)^n&5^n-(-1)^n\\
5^n-(-1)^n&5^n-(-1)^n&5^n+2(-1)^n
\end{array}
\right),
\end{align*}

et on retrouve le résultat obtenu plus haut, le calcul ayant été mené directement avec $n$ entier relatif.

\textbf{3ème solution.} Soit $n\in\Nn^*$. La division euclidienne de $X^n$ par $\chi_A$ fournit trois réels $a_n$, $b_n$ et $c_n$ et un polynôme $Q$ tels que $X^n=\chi_AQ+a_nX^2+b_nX+c_n$. En prenant les valeurs des membres en $5$, puis la valeur des deux membres ainsi que de leurs dérivées en $-1$ , on obtient

\begin{center}
$\left\{
\begin{array}{l}
25a_n+5b_n+c_n=5^n\\
a_n-b_n+c_n=(-1)^n\\
-2a_n+b_n=n(-1)^{n-1}
\end{array}
\right.\Leftrightarrow\left\{
\begin{array}{l}
b_n=2a_n-n(-1)^{n}\\
35a_n+c_n=5n(-1)^{n}+5^n\\
-a_n+c_n=-(n-1)(-1)^n
\end{array}
\right.\Leftrightarrow\left\{
\begin{array}{l}
a_n=\frac{1}{36}(5^n+(6n-1)(-1)^n)\\
\rule[-4mm]{0mm}{10mm}c_n=\frac{1}{36}(5^n+(-30n+35)(-1)^n)\\
b_n=\frac{1}{36}(2\times5^n+(-24n-2)(-1)^n)
\end{array}
\right.$.
\end{center}
 

Le théorème de \textsc{Cayley}-\textsc{Hamilton} fournit alors

\begin{align*}\ensuremath
A^n&=\frac{1}{36}\left((5^n+(6n-1)(-1)^n)A^2+2(5^n-(12n+1)(-1)^n)A+(5^n+(-30n+35)(-1)^n)I\right)
\\
     &=\frac{1}{36}\left((5^n+(6n-1)(-1)^n)
\left(
\begin{array}{ccc}
9&8&8\\
8&9&8\\
8&8&9
\end{array}
\right)+2(5^n-(12n+1)(-1)^n)\left(
\begin{array}{ccc}
1&2&2\\
2&1&2\\
2&2&1
\end{array}
\right)\right.\\
 &\left.+(5^n+(-30n+35)(-1)^n)\left(
\begin{array}{ccc}
1&0&0\\
0&1&0\\
0&0&1
\end{array}
\right)\right)\\
     &=\frac{1}{36}\left(
\begin{array}{ccc}
12\times5^n+24(-1)^n&12\times5^n-12(-1)^n&12\times5^n-12(-1)^n\\
12\times5^n-12(-1)^n&12\times5^n+24(-1)^n&12\times5^n-12(-1)^n\\
12\times5^n-12(-1)^n&12\times5^n-12(-1)^n&12\times5^n+24(-1)^n
\end{array}
\right)\\
 &=\frac{1}{3}\left(
\begin{array}{ccc}
5^n+2(-1)^n&5^n-(-1)^n&5^n-(-1)^n\\
5^n-(-1)^n&5^n+2(-1)^n&5^n-(-1)^n\\
5^n-(-1)^n&5^n-(-1)^n&5^n+2(-1)^n
\end{array}
\right).
\end{align*}

On retrouve encore une fois le même résultat mais pour $n\in\Nn^*$ uniquement.
\fincorrection
\correction{005652}
Soit $X\in\mathcal{M}_3(\Rr)$. Si $X^2 = A$ alors $AX= X^3 = XA$ et donc $X$ et $A$ commutent.

$A$ admet trois valeurs propres réelles et simples à savoir $1$, $3$ et $4$. Donc $A$ est diagonalisable dans $\Rr$ et les sous espaces propres de $A$ sont des droites. $X$ commute avec $A$ et donc laisse stable les trois droites propres de $A$.

Ainsi une base de $\mathcal{M}_{3,1}(\Rr)$ formée de vecteurs propres de $A$ est également une base de vecteurs propres de $X$ ou encore, si $P$ est une matrice réelle inversible telle que $P^{-1}AP$ soit une matrice diagonale $D_0$ alors pour la même matrice $P$, 
$P^{-1}XP$ est une matrice diagonale $D$. De plus

\begin{center}
$X^2 =A\Leftrightarrow PD^2P^{-1}= PD_0P^{-1}\Leftrightarrow D^2 = D_0\Leftrightarrow D =\text{diag}(\pm\sqrt{3},\pm2, \pm1)$
\end{center}

ce qui fournit huit solutions deux à opposées. On peut prendre $P=\left(
\begin{array}{ccc}
2&0&0\\
-16&1&0\\
5&0&1
\end{array}
\right)$ puis $P^{-1}=\left(
\begin{array}{ccc}
1/2&0&0\\
8&1&0\\
-5/2&0&1
\end{array}
\right)$. 
D'où les solutions 

\begin{align*}\ensuremath
\left(
\begin{array}{ccc}
2&0&0\\
-16&1&0\\
5&0&1
\end{array}
\right)\left(
\begin{array}{ccc}
\sqrt{3}\varepsilon_1&0&0\\
0&2\varepsilon&0\\
0&0&\varepsilon_3
\end{array}
\right)\left(
\begin{array}{ccc}
1/2&0&0\\
8&1&0\\
-5/2&0&1
\end{array}
\right)&=\left(
\begin{array}{ccc}
2\sqrt{3}\varepsilon_1&0&0\\
-16\sqrt{3}\varepsilon_1&2\varepsilon_2&0\\
5\sqrt{3}\varepsilon_1&0&\varepsilon_3
\end{array}
\right)\left(
\begin{array}{ccc}
1/2&0&0\\
8&1&0\\
-5/2&0&1
\end{array}
\right)\\
 &=\left(
\begin{array}{ccc}
\sqrt{3}\varepsilon_1&0&0\\
-8\sqrt{3}\varepsilon_1+16\varepsilon_2&2\varepsilon_2&0\\
5(\sqrt{3}\varepsilon_1-\varepsilon_3)/2&0&\varepsilon_3
\end{array}
\right).
\end{align*}

où $(\varepsilon_1, \varepsilon_2,\varepsilon_3)\in\{-1,1\}^3$.
\fincorrection
\correction{005653}
\begin{enumerate}
 \item  $\chi_A=-(2+X)((3-X)(-1-X)+4)=-(X+2)(X^2-2X+1)=-(X+2)(X-1)^2$. 

$A$ diagonalisable $\Rightarrow\text{dim}(\text{Ker}(A-I))=2\Rightarrow\text{rg}(A-I) = 1$ ce qui n'est pas . Donc $A$ n'est pas diagonalisable. De plus, $E_{-2}=\text{Vect}(e_1)$ où $e_1=\left(
\begin{array}{c}
0\\
0\\
1
\end{array}
\right)$ et $E_1=\text{Vect}(e_2)$ où $e_2=\left(
\begin{array}{c}
1\\
-2\\
-4
\end{array}
\right)$.

\item  $(A-I)^2=\left(
\begin{array}{ccc}
2&1&0\\
-4&-2&0\\
4&8&-3
\end{array}
\right)\left(
\begin{array}{ccc}
2&1&0\\
-4&-2&0\\
4&8&-3
\end{array}
\right)=\left(
\begin{array}{ccc}
0&0&0\\
0&0&0\\
-36&-36&9
\end{array}
\right)$ et donc $\text{Ker}(A-I)^2$ est le plan d'équation $4x+4y-z=0$.

\item  On note $f$ l'endomorphisme de $\Rr^3$ dont la matrice dans la base canonique de $\Rr^3$ est $A$. Le théorème de \textsc{Cayley}-\textsc{Hamilton} et le théorème de décomposition des noyaux permettent d'affirmer

\begin{center}
$\mathcal{M}_3(\Rr)=\text{Ker}(A+2I)\oplus\text{Ker}(A-I)^2$.
\end{center}

De plus,  chacun des sous-espaces $\text{Ker}(A+2I)$ et $\text{Ker}(A-I)^2$ étant stables par $f$, la matrice de $f$ dans toute base adaptée à cette décomposition est diagonale par blocs. Enfin, $\text{Ker}(A-I)$ est une droite vectorielle contenue dans le plan $\text{Ker}(A-I)^2$ et en choisissant une base de $\text{Ker}(A-I)^2$ dont l'un des deux vecteurs est dans $\text{Ker}(A-I)$, la matrice de $f$ aura la forme voulue.

On a déjà choisi $e_1=\left(
\begin{array}{c}
0\\
0\\
1
\end{array}
\right)$ et $e_2=\left(
\begin{array}{c}
1\\
-2\\
-4
\end{array}
\right)$ puis on prend $e_3=\left(
\begin{array}{c}
1\\
-1\\
0
\end{array}
\right)$. On note $P=\left(
\begin{array}{ccc}
0&1&1\\
0&-2&-1\\
1&-4&0
\end{array}
\right)$. $P$ est inversible d'inverse $P^{-1}=\left(
\begin{array}{ccc}
-4&-4&1\\
-1&-1&0\\
2&1&0
\end{array}
\right)$. On peut déjà affirmer que $P^{-1}AP$ est de la forme $\left(
\begin{array}{ccc}
-2&0&0\\
0&1&\times\\
0&0&1
\end{array}
\right)$.
Plus précisément 

\begin{center}
$Ae_3-e_3=\left(
\begin{array}{ccc}
3&1&0\\
-4&-1&0\\
4&8&-2
\end{array}
\right)\left(
\begin{array}{c}
1\\
-1\\
0
\end{array}
\right)-\left(
\begin{array}{c}
1\\
-1\\
0
\end{array}
\right)=\left(
\begin{array}{c}
1\\
-2\\
-4
\end{array}
\right)=e_2$
\end{center}

 et donc $Ae_3=e_2+e_3$ puis
 
 \begin{center}
 \shadowbox{
 $A=PTP^{-1}$ où $P=\left(
\begin{array}{ccc}
0&1&1\\
0&-2&-1\\
1&-4&0
\end{array}
\right)$, $P^{-1}=\left(
\begin{array}{ccc}
-4&-4&1\\
-1&-1&0\\
2&1&0
\end{array}
\right)$ et $T=\left(
\begin{array}{ccc}
-2&0&0\\
0&1&1\\
0&0&1
\end{array}
\right)$.
}
\end{center}

\item  Soit $n\in\Nn$. Posons $T = D+N$ où $D=\text{diag}(-2,1,1)$ et $N =E_{2,3}$. On a $ND = DN$ et $N^2 = 0$. Puisque les matrices $D$ et $N$ commutent, la formule du binôme de \textsc{Newton} permet d'écrire

\begin{align*}\ensuremath 
T^n&= D^n+nD^{n-1}N=\text{diag}((-2)^n,1,1) + n\text{diag}((-2)^{n-1},1,1)E_{2,3}=\text{diag}((-2)^n,1,1) +nE_{2,3}\\
 &=\left(
 \begin{array}{ccc}
 (-2)^n&0&0\\
 0&1&n\\
 0&0&1
 \end{array}
 \right).
\end{align*}

Puis

\begin{align*}\ensuremath
A^n&=PT^nP^{-1}=\left(
\begin{array}{ccc}
0&1&1\\
0&-2&-1\\
1&-4&0
\end{array}
\right)\left(
 \begin{array}{ccc}
 (-2)^n&0&0\\
 0&1&n\\
 0&0&1
 \end{array}
 \right)\left(
\begin{array}{ccc}
-4&-4&1\\
-1&-1&0\\
2&1&0
\end{array}
\right)\\
 &= \left(
\begin{array}{ccc}
0&1&n+1\\
0&-2&-2n-1\\
(-2)^n&-4&-4n
\end{array}
\right)\left(
\begin{array}{ccc}
-4&-4&1\\
-1&-1&0\\
2&1&0
\end{array}
\right)=\left(
\begin{array}{ccc}
2n+1&n&0\\
-4n&-2n+1&0\\
-4(-2)^n-8n+4&-4(-2)^n-4n+4&(-2)^n
\end{array}
\right).
\end{align*}

\begin{center}
\shadowbox{
$\forall n\in\Nn$, $A^n=\left(
\begin{array}{ccc}
2n+1&n&0\\
-4n&-2n+1&0\\
-4(-2)^n-8n+4&-4(-2)^n-4n+4&(-2)^n
\end{array}
\right)$.
}
\end{center}
\end{enumerate}
\fincorrection
\correction{005664}
$\chi_A=\left|
\begin{array}{cccc}
-X&b&\ldots&b\\
a&\ddots&\ddots&\vdots\\
\vdots&\ddots&\ddots&b\\
a&\ldots&a&-X
\end{array}
\right|$. Soit $f(x) =\left|
\begin{array}{cccc}
-X+x&b+x&\ldots&b+x\\
a+x&\ddots&\ddots&\vdots\\
\vdots&\ddots&\ddots&b+x\\
a+x&\ldots&a+x&-X+x
\end{array}
\right|$.

$f$ est un polynôme en $x$. Par $n$ linéarité du déterminant, $f(x)$ est somme de $2^n$ déterminants dont $2^n-(n+1)$ sont nuls car contiennent deux colonnes de $x$. Les déterminants restant contiennent au plus une colonne de x et sont donc de degré inférieur ou égal à $1$ en $x$. $f$ est donc une fonction affine. Il existe donc deux nombres $A$ et $B$ tels que $\forall x\in \Cc$, $f(x)=Ax+B$. Les égalités $f(-a)=(-X-a)^n$ et $f(-b)=(-X-b)^n$ fournissent $\left\{
\begin{array}{l}
-aA+B=(-X-a)^n\\
-bA+B=(-X-b)^n
\end{array}
\right.$ et comme $a\neq b$,  les formules de \textsc{Cramer} fournissent

\begin{center}
$\chi_A= f(0)=B=\frac{1}{b-a}(b(-X-a)^n-a(-X-b)^n)$.
\end{center}

Soit $\lambda\in\Cc$.

\begin{center}
$\lambda$ valeur propre de $A\Rightarrow \ch_A(\lambda) = 0\Rightarrow\left(\frac{\lambda+a}{\lambda+b}\right)^n=\frac{a}{b}\Rightarrow\left|\frac{\lambda+a}{\lambda+b}\right|=\left|\frac{a}{b}\right|^{1/n}$.
\end{center}

Soient $M$ le point du plan d'affixe $\lambda$, $A$ le point du plan d'affixe $-a$ et $B$ le point du plan d'affixe $-b$ puis $k =\left|\frac{a}{b}\right|^{1/n}$. $k$ est un réel strictement positif et distinct de $1$. On peut donc poser $I =\text{bar}(A(1),B(-k))$ et $J =\text{bar}(A(1),B(k))$.

\begin{align*}\ensuremath
\lambda\;\text{valeur propre de}\;A&\Rightarrow MA = kMB \Rightarrow MA^2 - k^2MB^2 = 0\Rightarrow (\overrightarrow{MA}-k\overrightarrow{MB})(\overrightarrow{MA}+k\overrightarrow{MB}) = 0\\ 
 &\Rightarrow(1-k)\overrightarrow{MI}.(1+k)\overrightarrow{MJ}= 0 \Rightarrow\overrightarrow{MI}.\overrightarrow{MJ}= 0\\
  &\Rightarrow M\;\text{est sur le cercle de diamètre}\;[I,J]\;(\text{cercles d'\textsc{Appolonius} (de Perga)}).
\end{align*}

\fincorrection
\correction{005665}
 \begin{enumerate}
 \item  Les hypothèses fournissent $AU= U$ où $U=\left(
\begin{array}{c}
1\\
\vdots\\
1
\end{array}
\right)$ et donc $1$ est valeur propre de $A$.

\item  
  \begin{enumerate}
  \item  Soient $\lambda$ une valeur propre de $A$ et $X =\left(
\begin{array}{c}
x_1\\
\vdots\\
x_n
\end{array}
\right)\in M_{n,1}(\Cc)$ un vecteur propre associé.

\begin{align*}\ensuremath
AX =\lambda X&\Rightarrow\forall i\in\llbracket1,n\rrbracket,\;\sum_{j=1}^{n}a_{i,j}x_j=\lambda x_i\Rightarrow\forall i\in\llbracket1,n\rrbracket,\;|\lambda x_i|\leqslant\sum_{j=1}^{n}|a_{i,j}||x_j|\\
 &\Rightarrow\forall i\in\llbracket1,n\rrbracket,\;|\lambda x_i|\leqslant\left(\sum_{j=1}^{n}a_{i,j}\right)\text{Max}\{|x_j|,\;1\leqslant j\leqslant n\}\\
 &\Rightarrow\forall i\in\llbracket1,n\rrbracket,\;|\lambda| |x_i|\leqslant\text{Max}\{|x_j|,\;1\leqslant j\leqslant n\}.
\end{align*}

On choisit alors pour $i$ un indice $i_0$ tel que $|x_{i_0}|=\text{Max}\{|x_j|,\;1\leqslant j\leqslant n\}$. Puisque $X$ est non nul, on a $|x_{i_0}|> 0$. On obtient

\begin{center}
$|\lambda| |x_{i_0}|\leqslant|x_{i_0}|$ et donc $|\lambda|\leqslant 1$ puisque $|x_{i_0}| > 0$.
\end{center}

  \item Plus précisément,

\begin{center}
$|\lambda-a_{i_0,i_0}||x_{i_0}|=\left|\sum_{j\neq i_0}^{}a_{i_0,j}x_j\right|\leqslant\sum_{j\neq i_0}^{}a_{i_0,j}|x_j|\leqslant \left(\sum_{j\neq i_0}^{}a_{i_0,j}\right)|x_{i_0}|=(1-a_{i_0,i_0})|x_{i_0}|$
\end{center}

et donc $\forall\lambda\in\text{Sp}A,\;|\lambda-a_{i_0,i_0}|\leqslant1-a_{i_0,i_0}$ ce qui signifie que les valeurs propres de $A$ appartiennent au disque de centre $a_{i_0,i_0}$ et de rayon $1-a_{i_0,i_0}$. Ce disque est tangent intérieurement au cercle de centre $(1,0)$ et de rayon $1$ en le point $(1,0)$.


$$\includegraphics{../images/img005665-1}$$

  \end{enumerate}
\end{enumerate}
\fincorrection
\correction{005668}
\begin{enumerate}
 \item  $J^n = I$. $J$ annule le polynôme $X^n-1$ qui est à racines simples dans $\Cc$ et donc $J$ est diagonalisable dans $\Cc$.

Les valeurs propres de $J$ sont à choisir parmi les racines $n$-èmes de $1$ dans $\Cc$. On pose $\omega=e^{2i\pi/n}$. Vérifions que $\forall k\in\llbracket0,n-1\rrbracket$, $\omega^k$ est valeur propre de $J$.

Soient $k\in\llbracket0,n-1\rrbracket$ et $X =(x_j)_{1\leqslant j\leqslant n}$ un élément de $\mathcal{M}_{n,1}(\Cc)$.

\begin{center}
$JX =\omega^kX\Leftrightarrow\left\{
\begin{array}{l}
x_2=\omega^kx_1\\
x_3=\omega^kx_2\\
\vdots\\
x_n=\omega^kx_{n-1}\\
x_1=\omega^kx_n
\end{array}
\right.\Leftrightarrow\left\{
\begin{array}{l}
x_2=\omega^kx_1\\
x_3=(\omega^k)^2x_1\\
\vdots\\
x_n=(\omega^k)^{n-1}x_{1}\\
x_1=(\omega^k)^nx_1
\end{array}
\right.\Leftrightarrow\left\{
\begin{array}{l}
x_2=\omega^kx_1\\
x_3=(\omega^k)^2x_1\\
\vdots\\
x_n=(\omega^k)^{n-1}x_{1}
\end{array}
\right.
$
\end{center}

et donc

\begin{center}
$JX =\omega^kX\Leftrightarrow X\in\text{Vect}(U_k)\;\text{où}\;U_k=\left(
\begin{array}{c}
1\\
\omega^k\\
(\omega^k)^2\\
\vdots\\
(\omega^k)^{n-1}
\end{array}
\right)$.
\end{center}

Donc $\forall k\in\llbracket0,n-1\rrbracket$, $\omega^k$ est valeur propre de $J$. Les valeurs propres de $J$ sont les $n$ racines $n$-èmes de $1$. Ces valeurs propres sont toutes simples. Le sous espace propre associé à $\omega^k$, $0\leqslant k\leqslant n-1$, est la droite vectorielle $D_k=\text{Vect}(U_k)$.

Soit $P$ la matrice de \textsc{Vandermonde} des racines $n$-èmes de l'unité c'est-à-dire $P=(\omega^{(j-1)(k-1)})_{0\leqslant j,k\leqslant n-1}$ puis $D =\text{diag}(1,\omega,...,\omega^{n-1})$, alors on a déjà vu que $P^{-1}=\frac{1}{n}\overline{P}$ (exercice \ref{ex:rou16}) et on a 

\begin{center}
\shadowbox{
$J =PDP^{-1}$ avec $D =\text{diag}(\omega^{j})_{1\leqslant j\leqslant n}$, $P=\left(\omega^{(j-1)(k-1)}\right)_{1\leqslant j,k\leqslant n}$ et $P^{-1}=\frac{1}{n}\overline{P}$ avec $\omega=e^{2i\pi/n}$.
}
\end{center}

\textbf{Remarque.} La seule connaissance de $D$ suffit pour le 2).

\item  Soit $A$ la matrice de l'énoncé.

\begin{center}
$A= a_0I + a_1J + a_2J^2+...+a_{n-1}J^{n-1}= Q(J)$ où $Q = a_0+a_1X+...+a_{n-1}X^{n-1}$.
\end{center}

D'après 1), $A = P\times Q(D)\times P^{-1}$ et donc $A$ est semblable à la matrice $\text{diag}(Q(1),Q(\omega),...,Q(\omega^{n-1}))$. Par suite, $A$ a même déterminant que la matrice $\text{diag}(Q(1),Q(\omega),...,Q(\omega^{n-1}))$. D'où la valeur du déterminant circulant de l'énoncé :

\begin{center}
\shadowbox{
$\left|
\begin{array}{ccccc}
a_0&a_1&\ldots&a_{n-2}&a_{n-1}\\
a_{n-1}&a_0&a_1& &a_{n-2}\\
\vdots& &\ddots&\ddots&\vdots\\
a_2& &\ddots&a_0&a_1\\
a_1&a_2&\ldots&a_{n-1}&a_0
\end{array}
\right|=\prod_{k=0}^{n-1}\left(\sum_{j=0}^{n-1}e^{2i(j-1)(k-1)\pi/n}a_j\right)$.
}
\end{center}
\end{enumerate}
\fincorrection
\correction{005669}
\begin{enumerate}
 \item  Soit $\sigma\in S_n$.

\begin{center}
$\text{det}(P_\sigma)=\sum_{\sigma'\in S_n}^{}\varepsilon(\sigma')p_{\sigma'(1),1}\ldots p_{\sigma'(n),n}=\sum_{\sigma'\in S_n}^{}\varepsilon(\sigma')\delta_{\sigma'(1),\sigma(1)}\ldots \delta_{\sigma'(n),\sigma(n)}=\varepsilon(\sigma)$,
\end{center}

car $\delta_{\sigma'(1),\sigma(1)}\ldots \delta_{\sigma'(n),\sigma(n)}\neq0\Leftrightarrow\forall i\in\llbracket1,n\rrbracket,\;\sigma'(i)=\sigma(i)\Leftrightarrow \sigma'=\sigma$.

\begin{center}
\shadowbox{
$\forall\sigma\in S_n$, $\text{det}(P_\sigma)=\varepsilon(\sigma)$.
}
\end{center}

\item   
  \begin{enumerate}
  \item Soit $(\sigma,\sigma')\in S_n^2$. Soit $(i,j)\in\llbracket1,n\rrbracket^2$. Le coefficient ligne $i$, colonne $j$, de la matrice $P_\sigma\times P_{\sigma'}$ vaut

\begin{center}
$\sum_{k=1}^{n}\delta_{i,\sigma(k)}\delta_{k,\sigma'(j)}$.
\end{center}

Dans cette somme, si $k\neq\sigma'(j)$, le terme correspondant est nul et quand $k=\sigma'(j)$, le terme correspondant vaut $\delta_{i,\sigma(\sigma'(j))}$. Finalement, le coefficient ligne $i$, colonne $j$, de la matrice $P_\sigma\times P_{\sigma'}$ vaut $\delta_{i,\sigma(\sigma'(j))}$ qui est encore le coefficient ligne $i$, colonne $j$, de la matrice $P_{\sigma\circ\sigma'}$.

\begin{center}
\shadowbox{
$\forall(\sigma,\sigma')\in S_n^2$, $P_\sigma\times P_{\sigma'}=P_{\sigma\circ\sigma'}$.
}
\end{center}

  \item Montrons que $G$ est un sous-groupe du groupe $(GL_n(\Rr),\times)$. $G$ contient $I_n=P_{Id}$ et d'autre part, $G$ est contenu dans $GL_n(\Rr)$ d'après 1).

\begin{center}
\shadowbox{
$(G,\times)$ est un sous-groupe de $(GL_n(\Rr),\times)$.
}
\end{center}
  \end{enumerate}

\item  Le coefficient ligne $i$, colonne $j$, de la matrice $AP_\sigma$ vaut

\begin{center}
$\sum_{k=1}^{n}a_{i,k}\delta_{k,\sigma(j)}=a_{i,\sigma(j)}$.
\end{center}

Par suite, si $C_1$,\ldots, $C_n$ désignent les colonnes de la matrice $A$, la matrice $AP_\sigma$ est la matrice dont les colonnes sont $C_{\sigma(1)}$,\ldots, $C_{\sigma(n)}$.

\begin{center}
\shadowbox{
Si $A=(C_1\ldots C_n)$, $AP_\sigma=(C_{\sigma(1)}\ldots C_{\sigma(n)})$.
}
\end{center}

\item  Commençons par trouver le polynôme caractéristique d'un cycle $c$ de longueur $\ell$ ($1\leqslant\ell\leqslant n$). Soit $f_c$ l'endomorphisme de $E=\Rr^n$ de matrice $P_c$ dans la base canonique de $\Rr^n$. Il existe une base de $E$ dans laquelle la matrice de $f_c$ est $\left(
\begin{array}{cc}
J_\ell&0_{\ell,n-\ell}\\
0_{n-\ell,\ell}&I_{n-\ell}
\end{array}
\right)$ où la matrice $J_\ell$ est la matice de l'exercice \ref{ex:rou18}. Le polynôme caractéristique $\chi_{P_c}$ de $P_c$ est donc $(-1)^n(X-1)^{n-\ell}(X^\ell-1)$ (voir exercice \ref{ex:rou18}).

Soit maintenant $\sigma\in S_n$. On note $f_\sigma$ l'endomorphisme de $E=\Rr^n$ de matrice $P_\sigma$ dans la base canonique de $\Rr^n$. $\sigma$ se décompose de manière unique à l'ordre près des facteurs en produit de cycles à supports disjoints, ces cycles commutant deux à deux.

Posons donc $\sigma= c_1\circ...\circ c_p$, $p\geqslant1$, où les $c_i$, $1\leqslant i\leqslant p$, sont des cycles à supports disjoints, et notons $\ell_i$ la longueur du cycle $c_i$, $1\leqslant i\leqslant p$. Il existe une base de $E$ dans laquelle la matrice de $f_\sigma$ est $\left(
\begin{array}{cccc}
J_{\ell_1}&0&\ldots&0\\
0&\ddots&\ddots&\vdots\\
\vdots&\ddots&J_{\ell_p}&0\\
0&\ldots&0&I_k
\end{array}
\right)$ où $k =n -\ell_1-...-\ell_p$ est le nombre de points fixes de $\sigma$.

Le polynôme caratéristique cherché est donc $\chi_{P_\sigma}= (-1)^n(X^{\ell_1}-1)\ldots(X^{\ell_p}-1)(X-1)^{n-\ell_1-\ldots-\ell_p}$. On en déduit immédiatement les valeurs propres de $P_\sigma$.
\end{enumerate}
\fincorrection
\correction{005671}
On cherche une matrice $A$ de format $4$ dont le polynôme caractéristique est $X^4-3X^3+X^2-1$. La matrice compagnon $A=\left(
\begin{array}{cccc}
0&0&0&1\\
1&0&0&0\\
0&1&0&-1\\
0&0&1&3
\end{array}
\right)$ convient (voir l'exercice \ref{exo:rou15}) et le théorème de \textsc{Cayley}-\textsc{Hamilton} montre que $A^4-3A^3+A^2-I_4= 0$.
\fincorrection
\correction{005672}
Soit $A$ la matrice de l'énoncé. $\text{det}A$ est le produit des valeurs propres de $A$.

\textbullet~Si $b = 0$, $\text{det}A = a^n$.

\textbullet~Si $b\neq0$, $\text{rg}(A-(a-b)I) = 1$ ou encore $\text{dim}(\text{Ker}(A-(a-b)I))= n-1$. Par suite, $a-b$ est valeur propre d'ordre $n-1$ au moins. On obtient la valeur propre manquante $\lambda$ par la trace de $A$ : $(n-1)(a-b) +\lambda= na$ et donc $\lambda = a + (n-1)b$. Finalement $\text{det}A = (a-b)^{n-1}(a+(n-1)b)$ ce qui reste vrai quand $b = 0$.

\begin{center}
\shadowbox{
$\left|
\begin{array}{cccc}
a&b&\ldots&b\\
b&a&\ddots&\vdots\\
\vdots&\ddots&\ddots&b\\
b&\ldots&b&a
\end{array}
\right|=(a-b)^{n-1}(a+(n-1)b)$.
}
\end{center}
\fincorrection
\correction{005676}
Posons $N=\left(
\begin{array}{ccc}
0&1&0\\
0&0&1\\
0&0&0
\end{array}
\right)$. On a $N^2=E_{1,3}$ et $N^3 = 0$. Si $X\in\mathcal{M}_3(\Cc)$ est une matrice carrée vérifiant $X^2 = N$, alors $X^6=0$. Donc $X$ est nilpotente et, puisque $X$ est de format $3$, on sait que $X^3 = 0$. Mais alors $N^2 =X^4 = 0$ ce qui n'est pas . L'équation proposée n'a pas de solution.

\fincorrection
\correction{005680}
$\text{rg}(M_{a,b}-I) =1$, si $a = b = 0$, $2$ si l'un des deux nombres $a$ ou $b$ est nul et l'autre pas et $3$ si $a$ et $b$ ne sont pas nuls. Donc $M_{0,0}$ n'est semblable à aucune des trois autres matrices et de même pour $M_{1,1}$.

Il reste à savoir si les matrices $M_{1,0}$ et $M_{0,1}$ sont semblables.

$(M_{1,0}-I)^2 = (E_{1,2}+E_{2,3})^2 =E_{1,3}\neq0$ et $(M_{0,1}-I)^2 = (E_{1,2}+E_{3,4})^2 = 0$. Donc les matrices $M_{1,0}$ et $M_{0,1}$ ne sont pas semblables.
\fincorrection
\correction{005684}
Soit $A =\left(
\begin{array}{ccc}
1&3&-7\\
2&6&-14\\
1&3&-7
\end{array}
\right)$. $A$ est de rang $1$ et donc admet deux valeurs propres égales à $0$ . $\text{Tr}A = 0$ et donc la troisième valeur propre est encore $0$. Donc $\chi_A = -X^3$. $A$ est nilpotente et le calcul donne $A^2 = 0$. Ainsi, si $X$ est une matrice telle que $X^2 = A$ alors $X$ est nilpotente et donc $X^3=0$.

\textbf{Réduction de $A$.} $A^2 = 0$. Donc $\text{ImA}\subset\text{Ker}A$. Soit $e_3$ un vecteur non dans $\text{Ker}A$ puis $e_2=Ae_3$. $(e_2)$ est une base de $\text{Im}A$ que l'on complète en $(e_1,e_2)$ base de $\text{Ker}A$.

$(e_1,e_2,e_3)$ est une base de $\mathcal{M}_{3,1}(\Cc)$ car si $ae_1 + be_2 + ce_3 = 0$ alors $A(ae_1 + be_2 + ce_3) = 0$ c'est-à-dire $c e_2 = 0$ et donc $c = 0$. Puis $a = b = 0$ car la famille $(e_1,e_2)$ est libre.

Si $P$ est la matrice de passage de la base canonique de $\mathcal{M}_{3,1}(\Cc)$ à la base $(e_1,e_2,e_3)$ alors $P^{-1}AP=\left(
\begin{array}{ccc}
0&0&0\\
0&0&1\\
0&0&0
\end{array}
\right)$.
On voit peut prendre $P=\left(
\begin{array}{ccc}
3&-7&0\\
-1&-14&0\\
0&-7&1
\end{array}
\right)$.

Si $X^2 = A$, $X$ commute avec $A$ et donc $X$ laisse stable $\text{Im}A$ et $\text{Ker}A$. On en déduit que $Xe_2$ est colinéaire à $e_2$ et $Xe_1$ est dans $\text{Vect}(e_1,e_2)$. Donc $P^{-1}XP$ est de la forme $\left(
\begin{array}{ccc}
a&0&d\\
b&c&e\\
0&0&f
\end{array}
\right)$. De plus, $X$ est nilpotente de polynôme caractéristique $(a-\lambda)(c-\lambda)(f-\lambda)$. On a donc nécessairement $a =c =f=0$. $P^{-1}XP$ est de la forme $\left(
\begin{array}{ccc}
0&0&b\\
a&0&c\\
0&0&0
\end{array}
\right)$.

Enfin, $X^2=A\Leftrightarrow\left(
\begin{array}{ccc}
0&0&b\\
a&0&c\\
0&0&0
\end{array}
\right)^2=\left(
\begin{array}{ccc}
0&0&0\\
0&0&1\\
0&0&0
\end{array}
\right)\Leftrightarrow ab = 1$.

Les matrices $X$ solutions sont les matrices de la forme $P\left(
\begin{array}{ccc}
0&0&\frac{1}{a}\\
a&0&b\rule{0mm}{4mm}\\
0&0&0
\end{array}
\right)P^{-1}$ où $a$ est non nul et $b$ quelconque.

On trouve $P^{-1}=\frac{1}{49}\left(
\begin{array}{ccc}
14&-7&0\\
-1&-3&0\\
-7&-21&49
\end{array}
\right)$  puis

\begin{align*}\ensuremath
X&=\frac{1}{49}\left(
\begin{array}{ccc}
3&-7&0\\
-1&-14&0\\
0&-7&1
\end{array}
\right)\left(
\begin{array}{ccc}
0&0&\frac{1}{a}\\
a&0&b\rule{0mm}{4mm}\\
0&0&0
\end{array}
\right)\left(
\begin{array}{ccc}
14&-7&0\\
-1&-3&0\\
-7&-21&49
\end{array}
\right)\\
 &=\frac{1}{49}\left(
\begin{array}{ccc}
-7a&0&\frac{3}{a}-7b\\
\rule[-4mm]{0mm}{10mm}-14&0&-\frac{1}{a}-14b\\
-7a&0&-7b
\end{array}
\right)\left(
\begin{array}{ccc}
14&-7&0\\
-1&-3&0\\
-7&-21&49
\end{array}
\right)\\
&=\left(
\begin{array}{ccc}
-2a-\frac{3}{7a}+b&a-\frac{9}{7a}+3b&\frac{3}{a}-7b\\
\rule[-4mm]{0mm}{10mm}-4a+\frac{1}{7a}+2b1&2a+\frac{3}{7a}+6b&-\frac{1}{a}-14b\\
-2a+b&a+3b&-7b
\end{array}
\right),\;(a,b)\in\Cc^*\times\Cc.
\end{align*}
\fincorrection
\correction{005685}
Soit $A=\left(
\begin{array}{ccc}
1&0&-1\\
1&2&1\\
2&2&3
\end{array}
\right)$. $\chi_A=\left|
\begin{array}{ccc}
1-X&0&-1\\
1&2-X&1\\
2&2&3-X
\end{array}
\right|=(1-X)(X^2-5X+4)-(-2+2X)=(1-X)(X^2-5X+4+2)=-(X-1)(X-2)(X-3)$.

$A$ est à valeurs propres réelles et simples. $A$ est diagonalisable dans $\Rr$ et les sous-espaces propres sont des droites.

Si $M$ est une matrice qui commute avec $A$, $M$ laisse stable ces droites et donc si $P$ est une matrice inversible telle que $P^{-1}AP$ soit diagonale alors la matrice $P^{-1}MP$ est diagonale. Réciproquement une telle matrice commute avec $A$.

\begin{center}
$C(A)=\{P\text{diag}(a,b,c)P^{-1},\;(a,b,c)\in\Cc^3\}$.
\end{center}

On trouve $C(A)=\left\{\left(\begin{array}{ccc}
2b-c&-a+2b-c&\frac{a-c}{2}\\
-b+c&a-b+c&(-a+c)/2\\
2c-2b&-2b+c&c
\end{array}
\right),\;(a,b,c)\in\Cc^3\right\}$. On peut vérifier que $C(A)=\text{Vect}(I,A,A^2)$.
\fincorrection
\nocorrection
\nocorrection
\nocorrection
\nocorrection
\nocorrection
\correction{001579}
Soit $u$ un endomorphisme d'un espace vectoriel $E$ de dimension finie,
alors le polynôme caractéristiaque de $u$ est aussi un polynôme
annulateur de $u$.

Preuve si $\chi_{u}$ est scindé à racines simples~: $u$ est alors
diagonalisable et il existe donc une base $B$ dans laquelle $Mat_{B}(u)=
\begin{pmatrix}
  \lambda_{1}\\
  &\ddots\\
  &      &\lambda_{n}
\end{pmatrix}
$. Alors $Mat_{B}(\chi_{u}(u))=\begin{pmatrix}
  \chi_{u}(\lambda_{1})\\
  &\ddots\\
  &      &\chi_{u}(\lambda_{n})
\end{pmatrix}
$. Et comme $\forall i\in\{1,\dots,n\}$ on a $\chi_{u}(\lambda_{i})=0$,
on en déduit que $\chi_{u}(u)=0$. 
\fincorrection
\nocorrection
\nocorrection
\nocorrection
\nocorrection
\nocorrection
\nocorrection
\nocorrection
\nocorrection
\nocorrection
\nocorrection
\nocorrection
\nocorrection
\nocorrection
\nocorrection
\nocorrection
\nocorrection
\nocorrection
\correction{002569}
 Soit $A$ la matrice suivante
$$A=\begin{pmatrix}0&1&1 \\  1&0&1 \\  1&1&0\end{pmatrix}$$
{\it Calculons $A^2$ et v\'erifions que $A^2=A+2I_3$}. 
On a 
$$A^2=\begin{pmatrix}0&1&1 \\  1&0&1 \\  1&1&0\end{pmatrix}\begin{pmatrix}0&1&1 \\  1&0&1 \\  1&1&0\end{pmatrix}=\begin{pmatrix}2&1&1 \\  1&2&1 \\  1&1&2\end{pmatrix}=A+2I_3.$$
On a donc $A^2-A=2I_3$, c'est-\`a-dire $A(A-I_3)=2I_3$, ou encore $A.{\frac{1}{2}}(A-I_3)=I_3$. Ce qui prouve que $A$ est in versible et que son inverse 
est $A^{-1}={\frac{1}{2}}(A-I_3)$
\fincorrection
\correction{002588}
{\it Soit $N$ une matrice nilpotente, il existe $q\in\N$ tel que $N^q=0$. Montrons que la matrice $I-N$ est inversible et exprimons son inverse en fonction de $N$.}

On remarque que 
$(I-N)(I+N+N^2+\dots+N^{q-1})=I-N^q=I.$
Ainsi, la matrice $I-N$ est inversible, et son inverse est 
$(I-N)^{-1}=I+N+N^2+\dots+N^{q-1}.$
\fincorrection
\correction{003518}
\begin{enumerate}
  \item $(-1)^n(X^n-a_nX^{n-1}-\dots-a_1)$.
  \item \'Etude de $x \mapsto (x^n-a_nx^{n-1}-\dots-a_1)/x^n$.
  \item Inégalité triangulaire.
  \item Expression générale de $A^k$.
\end{enumerate}
 \fincorrection
\correction{003541}
$a=b$ ou $a,b$ non nuls.
\fincorrection
\nocorrection
\nocorrection
\correction{003544}
\begin{enumerate}
  \item 
  \item $(A-xI)(\,{}^t\!A-xI) = (x^2-2x+4)I$, $\chi_A(x) = x^2-2x+4$.
  \item $\,{}^t\!A = 2I-A$ donc $(A-xI)((2-x)I-A) = (x^2-2x+4)I$.
             En prenant pour $x$ une des racines du polynôme $x^2-2x+4$, on obtient
             un polynôme scindé à racines simples annulant $A$.
\end{enumerate}
 \fincorrection
\nocorrection
\correction{003546}
$A$ est diagonalisable car $A^2=I$. $e^A = (\ch1)I + (\sh1)A$.
\fincorrection
\correction{003547}
Si $\Im u\subset \mathrm{Ker} u$ alors $u^2 = 0$ donc $0$ est l'unique
         valeur propre de $u$ et $u\ne 0$ donc $u$ n'est pas diagonalisable.\par
         Si $\Im u \not \subset \mathrm{Ker} u$ alors $\Im u \cap \mathrm{Ker} u = \{\vec 0\}$
         et donc $\Im u + \mathrm{Ker} u = E$. Or $\Im u$ et $\mathrm{Ker} u$ sont des sous-espaces
         propres de $u$ donc $u$ est diagonalisable.
\fincorrection
\nocorrection
\correction{003549}
\begin{enumerate}
  \item Polynôme annulateur simple.
  \item Non, ctrex = $B$ nilpotent.
\end{enumerate}
 \fincorrection
\correction{003550}
$\text{spec}(p) \subset \{-1,0,1\}$. $p$ est diagonalisable si et
	     seulement s'il annule un polynôme scindé à racines simples.
\fincorrection
\correction{003551}
$A$ est $\C$-diagonalisable et les valeurs propres sont $\alpha > 0$
et $\beta,\overline \beta$ avec la même multiplicité.
\fincorrection
\nocorrection
\correction{003553}
$A$ est diagonalisable et a $n$ valeurs propres distinctes,
	     sinon il existerait un polynôme annulateur de degré inférieur ou
	     égal à $n-1$. Ces racines sont les $n$ racines $n$-èmes de 1 et
	     leur somme est nulle.
\fincorrection
\correction{003554}
$A$ est $\C$-diagonalisable (polynôme annulateur à racines simples)
$ \Rightarrow  \dim(E_1) + \dim(E_{-1}) = n$.
Les dimensions sont conservées sur $\R$.
\fincorrection
\nocorrection
\correction{003556}
Soit $P$ un polynôme tel que $P(\lambda) = 1$ et $P(\mu) = 0$ pour
toutes les autres valeurs propres, $\mu$, de~$f$. Alors $p_\lambda = P(f)$.
\fincorrection
\correction{003557}
  3. $\mathrm{Spec}(u_k)\subset \{i,-i\}$ d'après la relation $u_k^2 = -\mathrm{id}_E$.
    Si le spectre était réduit à un élément alors $u_k$ serait scalaire car
    diagonalisable, mais ceci est incompatible avec la relation d'anticommutation
    entre~$u_k$ et $u_\ell$. Donc $\mathrm{Spec}(u_k)=\{i,-i\}$.
    
  4. $u_\ell$ avec $\ell\ne k$ échange les sous-espaces propres de~$u_k$
    donc ils ont même dimension~$n/2$.
 \fincorrection
\nocorrection
\nocorrection
\correction{003560}
\begin{enumerate}
  \item Calcul Maple~:
$h=\begin{pmatrix}c+4&b&a\cr 0&c+2&b\cr 0&0&c\cr\end{pmatrix}$,
$v=ku$.
  \item 
  \begin{enumerate}
    \item 
    \item 
    \item $u^k\circ h - h\circ u^k = -2ku^k$,
$P(u)\circ h - h \circ P(u) = -2u\circ P'(u)$.
    \item Si $P(u) = 0$ alors $u\circ P'(u) = 0$ donc $P$ (polynôme minimal)
divise $XP'$ ce qui implique $P(X) = X^k$ pour un certain~$k$.
  \end{enumerate} 
\end{enumerate}
 \fincorrection
\nocorrection
\nocorrection
\correction{003563}
Aucun polynôme constant ne convient. Si $P$ est non constant et $\alpha$
est une racine de~$P$ alors en considérant $A=\alpha I_n$ on obtient une
première condition nécessaire~: $n\alpha \in\Z$. Si $P$ a une autre racine $\beta$
alors en prenant $A = \mathrm{diag}(\alpha,\dots,\alpha,\beta)$ on obtient
une deuxième condition nécessaire~: $\beta-\alpha\in\Z$. Ainsi les polynômes~$P$
cherchés ont la propriété suivante~: {\it $\deg(P)\ge 1$ et il existe $u\in\Z$ tel que
toutes les racines de~$P$ sont congrues à $u/n$ modulo~$1$.}
Cette condition est clairement suffisante.
\fincorrection
\correction{003564}
On écrit $C = PJQ$ où $P,Q$ sont inversibles et $J$ est la matrice
canonique de rang~$r$. Alors $(P^{-1}AP)J = J(QBQ^{-1})$ donc $P^{-1}AP$
et $QBQ^{-1}$ sont triangulaires par blocs avec le même bloc diagonal
$r\times r$, ce qui prouve que $\chi_A$ et $\chi_B$ ont un facteur de
degré~$r$ en commun.
\fincorrection
\correction{003565}
Le polyn{\^o}me s'{\'e}crit $(X^{2}+1)(X^{2}+X+1)$. Il n'a donc pas de racine r{\'e}elle. Or tout {\'e}l{\'e}ment de $M_{5}(\R)$
poss{\`e}de au moins une valeur propre et cette valeur propre devrait {\^e}tre {\'e}galement racine du polyn{\^o}me minimal. Par cons{\'e}quent
 $X^{4}+X^{3}+2X^{2}+X+1$ ne peut pas {\^e}tre le polyn{\^o}me minimal d'une matrice de $M_{5}(\R)$.
\fincorrection
\correction{003566}
\begin{enumerate}
  \item Que c'est un isomorphisme (et réciproquement).
  \item Soit $Q(X) = P(X)/X$. On a $u\circ Q(u)=0$ et $X, Q$ sont
    premiers entre eux, d'où $E=\mathrm{Ker} u \oplus \mathrm{Ker} Q(u)$ et
    $\Im u \subset \mathrm{Ker} Q(u)$. On conclut avec le théorème du rang.
  \item Même méthode.
\end{enumerate}
 \fincorrection
\correction{005658}
\textbullet~Si $A$ est nilpotente, pour tout $k\in\llbracket1,n\rrbracket$, $A^k$ est nilpotente et donc $0$ est l'unique valeur propre dans $\Cc$ de $A^k$. Par suite, $\forall k\in\llbracket1,n\rrbracket$, $\text{Tr}(A^k) = 0$.

\textbullet~
Réciproquement , supposons que $\forall k\in\llbracket1,n\rrbracket$, $\text{Tr}(A^k) = 0$ et  montrons alors que toutes les valeurs propres de $A$ dans $\Cc$ sont nulles. Ceci montrera que le polynôme caractéristique de $A$ est $(-X)^n$ et donc que $A$ est nilpotente d'après le théorème de \textsc{Cayley}-\textsc{Hamilton}.

Soient $\lambda_1$,..., $\lambda_n$ les $n$ valeurs propres (distinctes ou confondues) de $A$ dans $\Cc$. Pour $k\in\llbracket1,n\rrbracket$, on pose $S_k =\lambda_1^k+...+\lambda_n^k$. Il s'agit de montrer que : $(\forall k\in\llbracket1,n\rrbracket,\;S_k =0)\Rightarrow(\forall j\in\llbracket1,n\rrbracket,\;\lambda_j =0)$.

\textbf{1ère solution.}
Les $S_k$, $1\leqslant k\leqslant n$, sont tous nuls et par combinaisons linéaires de ces égalités, on en déduit que pour tout polynôme $P$ de degré inférieur ou égal à $n$ et s'annulant en $0$, on a $P(\lambda_k) = 0$  (1). Il s'agit alors de bien choisir le polynôme $P$.

Soit $i\in\llbracket1,n\rrbracket$. Soient $\mu_1$,..., $\mu_p$ les valeurs propres deux à deux distinctes de $A$ ($1\leqslant p\leqslant n$). On prend $P =X\prod_{j\neq i}^{}(X-\mu_j)$ si $p\geqslant2$ et $P = X$ si $p = 1$. $P$ est bien un polynôme de degré inférieur ou égal à $n$ et s'annule en $0$. L'égalité $P(\lambda_i)=0$ fournit $\lambda_i = 0$ ce qu'il fallait démontrer.

\textbf{2ème solution.} Pour ceux qui savent que les sommes de \textsc{Newton} $S_k$ sont liées aux fonctions élémentaires en les $\lambda_i$  $\sigma_1$,..., $\sigma_n$ par les formules de \textsc{Newton} :

\begin{center}
$\forall k\leqslant n$, $S_k -\sigma_1S_{k-1}+...+(-1)^{k-1}\sigma_{k-1}S_1+ (-1)^kk\sigma_k = 0$.
\end{center}

Par suite, si tous les $S_k$, $1\leqslant k\leqslant n$, sont nuls alors immédiatement tous les $\sigma_k$, $1\leqslant k\leqslant n$, sont nuls et donc les $\lambda_i$ sont nuls car tous racines de l'équation $x^n = 0$.
\fincorrection
\correction{005659}
Soit $k\in\Nn^*$.

\begin{align*}\ensuremath
f^kg-fg^k&=f^kg-f^{k-1}gf+f^{k-1}gf-f^{k-2}gf^2+f^{k-2}gf^2-\ldots-fgf^{k-1}+fgf^{k-1}-gf^k\\
 &=\sum_{i=0}^{k-1}(f^{k-i}gf^i-f^{k-i-1}gf^{i+1})=\sum_{i=0}^{k-1}f^{k-i-1}(fg-gf)f^i=\sum_{i=0}^{k-1}f^{k-i-1}ff^i\\
 &=kf^k.
\end{align*}

Ainsi,

\begin{center}
\shadowbox{
si $fg-gf=f$, alors $\forall k\in\Nn$, $f^kg-gf^k=kf^k$\quad$(*)$.
}
\end{center}

\textbf{1ère solution.} Soit $\begin{array}[t]{cccc}
\varphi~:&\mathcal{L}(E)&\rightarrow&\mathcal{L}(E)\\
 &h&\mapsto&hg-gh
\end{array}$. $\varphi$ est un endomorphisme de $\mathcal{L}(E)$ et $\forall k\in\Nn^*$, $\varphi(f^k)=kf^k$. Si, pour $k\in\Nn^*$ donné, $f^k$ n'est pas nul, $f^k$ est valeur propre de $\varphi$ associé à la valeur propre $k$. Par suite, si aucun des $f^k$ n'est nul, $\varphi$ admet une infinité de valeurs propres deux à deux distinctes. Ceci est impossible car $\text{dim}(\mathcal{L}(E))<+\infty$. Donc, $f$ est nilpotent.

\textbf{2ème solution.} Les égalités $(*)$ peuvent s'écrire $P(f)g-gP(f)=fP'(f)$, $(**)$, quand $P$ est un polynôme de la forme $X^k$, $k\in\Nn$. Par linéarité, les égalités $(**)$ sont vraies pour tout polynôme $P$.

En particulier, l'égalité $(**)$ est vraie quand $P$ est $Q_f$ le polynôme minimal de $f$ et donc

\begin{center}
$fQ_f'(f)=Q_f(f)g-gQ_f(f)=0$.
\end{center}

Le polynôme $XQ_f'$ est donc un polynôme annulateur de $f$ et on en déduit que le polynôme $Q_f$ divise le polynôme $XQ_f'$. Plus précisément, si $p\in\Nn^*$ est le degré de $Q_f$, les polynômes $pQ_f$ ayant mêmes degrés et mêmes coefficients dominants, on en déduit que $pQ_f=XQ_f'$ ou encore que

\begin{center}
$\frac{Q_f'}{Q_f}=\frac{p}{X}$.
\end{center}

Par identification à la décomposition en éléments simples usuelles de $\frac{Q_f'}{Q_f}$, on en déduit que $Q_f=X^p$. En particulier, $f^p=0$ et encore une fois $f$ est nilpotent.
\fincorrection
\nocorrection
\nocorrection
\nocorrection
\correction{001702}
\begin{enumerate}
  \item
    $u\big(\sum_{i=1}^{N}\alpha_{i}x_{i}\big)=\sum_{i=1}^{N}\alpha_{i}u(u^{i}(x_{0}))=\sum_{i=1}^{N}\alpha_{i} x_{i+1}$. Donc $\forall x\in F,\ u(x)\in F$.
  \item
    Si à un rang $k$, $x_{k+1}$ est une combinaison linéaire des $x_{i}$
    pour $i\leq k$~: $x_{k+1}=\sum_{i=0}^{k}a_{i}x_{i}$. On en déduit que
    $x_{k+2}=\sum_{i=0}^{k}a_{i}x_{i+1}$, et donc que $x_{k+2}\in
    Vect(x_{1},\dots,x_{k+1})\subset Vect(x_{0},\dots, x_{k})$, et par
    récurrence, on obtient finalement que $\forall p>k, x_{p}\in
    Vect(x_{0},\dots,x_{k})$. On en déduit que le rang de la famille
    $\{x_{0},\dots,x_{m}\}$, est strictement croissant avec $m$ puis
    éventuellement constant à partir d'un certain rang. Comme $E$ est de
    dimension finie $n$, on en déduit que ce rang est constant à partir
    d'un rang $k\leq n$~: la famille $(x_{0},\dots,x_{k})$ est alors
    libre, et $x_{k+1}$ est une combinaison linéaire de
    $(x_{0},\dots,x_{k})$. 
  \item
    $x_{k+1}-\sum_{i=0}^{k}a_{i}x_{i}=u^{k+1}(x_{0})-\sum_{i=0}^{k}a_{i}u^{i}(x_{0})=0$ donc $P_{0}(u)(x_{0})=0$.
  \item
    Si $x\in F$ alors $x=\sum_{i=0}^{N}\alpha_{i}u^{i}(x_{0})$. En posant
    $P=\sum_{i=0}^{N}\alpha_{i} X^{i}$, on a $x=P(u)(x_{0})$.
  \item
    Soit $P=QP_{0}+R$ la division euclidienne de $P$ par $P_{0}$, alors
    $\deg(R)<\deg(P_{0})=k+1$. Notons $R=\sum_{i=0}^{k}r_{i} X^{i}$. On a
    $x=P(u)(x_{0})=Q(u)P_{0}(u)(x_{0})+R(u)(x_{0})=R(u)(x_{0})$ 
  \item
    La famille $(x_{0},\dots,x_{k})$ est donc libre et génératrice dans
    $F$~: c'est une base.
  \item
    La matrice de $u_{|_{F}}$ dans cette base est la matrice compagnon
    associée au polynôme $P_{0}$, et $\chi_{u_{|_{F}}}=P_{0}$.
  \item
    On choisit un vecteur $y\in E\setminus F$, et on recommence le même
    travail avec ce vecteur, et on continue ainsi jusqu'à avoir obtenu
    une base de tout l'espace. La matrice de $u$ dans la base finale est
    alors du type demandé.
  \end{enumerate}
\fincorrection
\nocorrection
\correction{003520}
spec$(T) = {]-1,1]}$.
\fincorrection
\correction{003521}
2. $0 < \lambda \le 1$ : $f(x) = Cx^{1/\lambda-1}$.
 \fincorrection
\correction{003522}
$1/k$, $k\ge 1$.
\fincorrection
\correction{003523}
$\lambda = \frac1{(\pi/2+k\pi)^2}$ : $u(x) = C\sin(\pi/2+k\pi)x$.
\fincorrection
\nocorrection
\correction{003580}
\begin{enumerate}
  \item $A \sim \text{diag}(1,\alpha,\alpha^{-1})$ où $\alpha$ est une
             racine primitive $7^{\text{ème}}$ de $1$,\par
             $A \sim \text{diag}(\alpha,\alpha^{10},\alpha^{-11})$ où $\alpha$
             est une racine primitive $37^{\text{ème}}$ de $1$.
  \item pas de solution.
  \item vp${} = 0$ ou $1$.
\end{enumerate}
 \fincorrection
\correction{003584}
Soit $f$ un endomorphisme d'un ev $E$ ayant $A$ pour matrice.
On doit trouver $g\in GL(E)$ tel que $f\circ g = 2g\circ f$. Construction
de~$g$ par récurrence sur $n=\dim E$.

$n\le1$~: on a $f=0$ donc $g=\mathrm{id}_E$ convient.

$0,\dots,n-1 \Rightarrow n$~: $f$ est non surjectif donc l'hypothèse de récurrence
s'applique à~$f_{|\Im(f)}$. Soit $g_1\in GL(\Im(f))$ tel que
$f(g_1(x)) = 2g_1(f(x))$ pour tout~$x\in\Im(f)$. Soit
$E = H \oplus I \oplus K \oplus L$
avec $H=\Im(f)\cap\mathrm{Ker}(f)$, $H \oplus I = \Im(f)$ et $H \oplus K = \mathrm{Ker}(f)$.
La restriction de~$f$ à $I\oplus L$ induit un isomorphisme sur~$\Im(f)$,
on note $\varphi$ l'isomorphisme réciproque. Soit~$g\in\mathcal{L}(E)$ définie par~:

$$g(h+i+k+\ell) = g_1(h+i) + k + 2\varphi(g_1(f(\ell))).$$

On vérifie facilement que $f\circ g = 2g\circ f$ et il reste à prouver
que $g$ est injective. Si $x=h+i+k+\ell\in\mathrm{Ker} g$ alors
$g(f(x)) = g_1(f(i+\ell)) = 0$ donc $i+\ell\in\mathrm{Ker} f = H\oplus K$ soit
$i=\ell=0$. Il reste $g_1(h)+k=0$ ce qui implique $h=k=0$ car $g_1(h)\in \Im f
= H\oplus I$.

Remarque~: la démonstration passe à tout corps de caractéristique différente de~$2$.
\fincorrection
\correction{003597}
\begin{enumerate}
  \item 
  \item Par récurrence pour $P = X^k$, puis par linéarité.
  \item $A=0$.
\end{enumerate}
 \fincorrection
\correction{003598}
S'inspirer du cas $n=1$. Soit $P = \begin{pmatrix}I&I\cr I&-I\cr\end{pmatrix}$ :
	     $P^{-1}AP = \begin{pmatrix}2M&0\cr 0&0\cr\end{pmatrix}$ est diagonalisable,
	     donc $A$ aussi.
\fincorrection
\correction{003599}
$E_\lambda(M) = \left\{ \begin{pmatrix}\lambda Y\cr Y\hfil\end{pmatrix} \text{ tq }
         AY = \lambda^2Y \right\}$.
\fincorrection
\nocorrection
\correction{003601}
Calcul du polynôme caractéristique de~$B$ par opérations en blocs.
On obtient $$\chi_B(x) = \det(x^2I-2xA-A^2) = (-1)^n\chi_A\Bigl(\frac x{1+\sqrt2}\Bigr)\chi_A\Bigl(\frac x{1-\sqrt2}\Bigr)$$
donc $$\mathrm{Spec}(B) = \{(1+\sqrt2)\lambda,\ \lambda\in\mathrm{Spec}(A)\} \cup \{(1-\sqrt2)\lambda,\ \lambda\in\mathrm{Spec}(A)\}.$$
\fincorrection
\correction{003602}
%Diagonalisation par blocs.
En prenant $P=\begin{pmatrix}I_2&I_2\cr-I_2&I_2\cr\end{pmatrix}$
on trouve $P^{-1}MP = \begin{pmatrix}a^2-ab&ab-b^2&0&0\cr ab-b^2&a^2-ab&0&0\cr0&0&a^2+ab&b^2+ab\cr0&0&b^2+ab&a^2+ab\cr\end{pmatrix}
= \begin{pmatrix}M_1&0\cr0&M_2\cr\end{pmatrix}$.

En prenant $P_1=\begin{pmatrix}1&1\cr-1&1\cr\end{pmatrix}$ on a
$P_1^{-1}M_1P_1 = \begin{pmatrix}(a-b)^2&0\cr0&a^2-b^2\cr\end{pmatrix}$ et 
$P_1^{-1}M_2P_1 = \begin{pmatrix}a^2-b^2&0\cr0&(a+b)^2\cr\end{pmatrix}$.

Ainsi, $\mathrm{Spec}{A} = \{(a+b)^2,(a-b)^2,(a+b)(a-b)\}$, donc l'ensemble cherché
est la boule unité ouverte pour $\|\ \|_1$.
\fincorrection
\nocorrection
\nocorrection
\correction{003605}
Si $P(0)\ne 0$ alors $f$ est bijective. Si $P(0) = 0$ alors
$f^2\circ\text{qqch} = -P'(0)f \Rightarrow  \mathrm{Ker} f^2 = \mathrm{Ker} f$.
\fincorrection
\nocorrection
\correction{003607}
Soit $\mu$ le polynôme minimal de~$u$ et ${\cal D}$ l'ensemble des
diviseurs unitaires de~$\mu$. Pour $P\in K[X]$ et $d=P\wedge\mu$ on
a facilement $\mathrm{Ker}(P(u)) = \mathrm{Ker}(d(u))$ et $\Im(P(u)) = \Im(d(u))$.
Ceci montre déjà que $\cal K$ et $\cal I$ sont finis.

De plus, si $d\in{\cal D}$ alors l'annulateur
minimal de $u_{|\Im(d(u))}$ est $\mu/d$ donc l'application $d \mapsto\Im(d(u))$
est injective sur $\cal D$ et $\mathrm{Card}\,({\cal I}) = \mathrm{Card}\,({\cal D})$.
De même, l'annulateur minimal de $u_{|\mathrm{Ker}(d(u))}$ est $d$ car
$\mathrm{Ker}(d(u))\supset \Im(\frac\mu d(u))$ et $d$ est l'annulateur minimal
de $u_{|\Im(\frac\mu d(u))}$ donc l'application $d \mapsto\mathrm{Ker}(d(u))$
est injective sur $\cal D$ et $\mathrm{Card}\,({\cal K}) = \mathrm{Card}\,({\cal D})$.
\fincorrection
\correction{003608}
En appliquant le théorème du rang à $f_{|\mathrm{Ker} f^2}$, on a~:
$\dim(\mathrm{Ker} f^2) = \dim(\mathrm{Ker} f) + \dim(f(\mathrm{Ker} f^2))$, et $f(\mathrm{Ker} f^2)\subset \mathrm{Ker} f$,
donc $f(\mathrm{Ker} f^2)= \mathrm{Ker} f$. Soit $G_i = \mathrm{Ker} g^i$. Montrons que $g(G_{i+1}) = G_i$
pour tout~$i\in{[[0,k]]}$~: si $x\in G_{i+1}$ alors $g^i(g(x)) = g^{i+1}(x) = 0$
donc $g(x)\in G_i$. Réciproquement, si $y\in G_i$ alors $y\in G_k = f(G_{2k})$,
donc $y$ a un antécédant $x$ par $f$, cet antécédant appartient à $G_{i+k}$,
et $y = g(g^{k-1}(x)) \in g(G_{i+1})$. On en déduit, avec le théorème du rang
appliqué à $g_{|G_{i+1}}$, que $\dim(G_{i+1}) = \dim(G_i) + \dim(\mathrm{Ker} g)$
pour tout~$i\in{[[0,k]]}$, d'où $d = \dim(G_k) = \dim(G_0) + k\dim(\mathrm{Ker} g) = k\dim(\mathrm{Ker} g)$.
\fincorrection
\correction{005661}
\begin{enumerate}
 \item  \textbullet~$E$ contient $I_2$ et est inclus dans $GL_2(\Rr)$.

\textbullet~Si $A$ et $B$ sont dans $E$ alors $AB$ est à coefficients entiers et $\text{det}(AB) =  \text{det}A\text{det}B  = 1$. Donc $AB$ est dans $E$. 

\textbullet~Si $A$ est dans $E$, $\text{det}(A^{-1})= 1$ et en particulier $A^{-1}=\frac{1}{\text{det}A}{^t}\text{com}(A)$ est à coefficients entiers. On en déduit que $A^{-1}$ est dans $E$.

Finalement

\begin{center}
\shadowbox{
$E$ est un sous-groupe de $GL_2(\Rr)$.
}
\end{center}

\item  Soit $A$ un élément de $E$ tel qu'il existe un entier naturel non nul $p$ tel que $A^p =I_2$.

$A$ est diagonalisable dans $\Cc$ car annule le polynôme à racines simples $X^p-1$.

$A$ admet deux valeurs propres distinctes ou confondues qui sont des racines $p$-èmes de $1$ dans $\Cc$ et puisque $A$ est réelle, on obtient les cas suivants :

\textbf{1er cas.} Si $\text{Sp}A =(1,1)$, puisque $A$ est diagonalisable, $A$ est semblable à $I_2$ et par suite $A = I_2$. Dans ce cas, $A^{12}= I_2$.

\textbf{2ème cas.} Si $\text{Sp}A = (-1,-1)$, $A = -I_2$ et $A^{12}=I_2$.

\textbf{3ème cas.} Si $\text{Sp}A = (1,-1)$ alors $A$ est semblable à $\text{diag}(1,-1)$ et donc $A^2 = I_2$ puis encore une fois $A^{12}= I_2$.

\textbf{4ème cas.} Si $\text{Sp}A =(e^{i\theta},e^{-i\theta})$. Dans ce cas $\text{Tr}A = 2cos\theta$ est un entier ce qui impose $2\cos\theta\in\{-2,-1,0,1,2\}$. Les cas $\cos\theta= 1$ et $\cos\theta=-1$ ont déjà été étudié.

\textbullet~Si $\cos\theta= 0$, $\text{Sp}A=(i,-i)$ et $A$ est semblable à $\text{diag}(i,-i)$. Donc $A^4 =I_2$ puis $A^{12}= I_2$.

\textbullet~Si $\cos\theta=\pm\frac{1}{2}$, $\text{Sp}A=(j,j^2)$ ou $\text{Sp}A=(-j,-j^2)$. Dans le premier cas, $A^3 = I_2$ et dans le deuxième $A^6 = I_2$.

Dans tous les cas $A^{12}= I_2$.
\end{enumerate}
\fincorrection
\correction{005662}
On montre le résultat par récurrence sur $n\in\Nn^*$ le format de $A$.

\textbullet~C'est clair pour $n = 1$.

\textbullet~Soit $n\geqslant1$. Supposons que toute matrice de format $n$ et de trace nulle soit semblable à une matrice de diagonale nulle.

Soient $A$ une matrice carrée de format $n+1$ et de trace nulle puis $f$ l'endomorphisme de $\Kk^{n+1}$ de matrice $A$ dans la base canonique $(e_1,...,e_{n+1})$ de $\Kk^{n+1}$.

Si $f$ est une homothétie de rapport noté $k$, alors $0=\text{Tr}(f)=k(n+1)$ et donc $k = 0$ puis $f = 0$ puis $A = 0$. Dans ce cas, $A$ est effectivement semblable à une matrice de diagonale nulle.

Sinon $f$ n'est pas une homothétie et on sait qu'il existe un vecteur $u$ de $E$ tel que la famille $(u,f(u))$ soit libre (voir exercice \ref{ex:rou25}). On complète la famille libre $(u,f(u))$ en une base de $E$. Le coefficient ligne 1, colonne 1, de la matrice de $f$ dans cette base est nul. Plus précisément, $A$ est semblable à une matrice de la forme $\left(
\begin{array}{cccccc}
0&\times&\ldots& &\ldots&\times\\
1& & & & \\
0& & &  & \\
\vdots& & &A'& & \\
\vdots& \\
0&
\end{array}
\right)$.

Puis $\text{Tr}A'=\text{Tr}A= 0$ et par hypothèse de récurrence, $A'$ est semblable à une matrice $A_1$ de diagonale nulle ou encore il existe $A_1$ matrice carrée de format $n$ et de diagonale nulle et $Q\in GL_n(\Kk)$ telle que $Q^{-1}A'Q = A_1$.

Mais alors, si on pose $P=\left(
\begin{array}{cccc}
1&0&\ldots&0\\
0& & & \\
\vdots& &Q& \\
0& & & 
\end{array}
\right)$, $P$ est inversible car $\text{det}(P)=1\times\text{det}(Q)\neq0$ et un calcul par blocs montre que $P^{-1}=\left(
\begin{array}{cccc}
1&0&\ldots&0\\
0& & & \\
\vdots& &Q^{-1}& \\
0& & & 
\end{array}
\right)$ puis que $P^{-1}AP=\left(
\begin{array}{cccccc}
0&\times&\ldots& &\ldots&\times\\
\times& & & & \\
\vdots& & &  & \\
 & & &A_1& & \\
\vdots& \\
\times&
\end{array}
\right)$ est de diagonale nulle.
\fincorrection
\correction{005681}
Soit $B$ la matrice de l'énoncé. $\text{rg}B = 1$ et si $A$ existe, nécessairement $\text{rg}A = n-1$ (exercice \ref{ex:rou18bis}).

Une matrice de rang $1$ admet l'écriture générale $U{^t}V$ où $U$ et $V$ sont des vecteurs colonnes non nuls. Ici $U=\left(
\begin{array}{c}
1\\
2\\
\vdots\\
n
\end{array}
\right)$  et $V =\left(
\begin{array}{c}
1\\
0\\
\vdots\\
0
\end{array}
\right)$.

Si $A$ existe, $A$ doit déjà vérifier $A{^t}tB={^t}BA = 0$ ou encore $AV{^t}U= 0$ (1) et $V{^t}UA = 0$ (2). En multipliant les deux membres de l'égalité (1) par $U$ à droite puis en simplifiant par le réel non nul ${^t}UU=\|U\|_2^2$, on obtient $AV = 0$. Ceci montre que la première colonne de $A$ est nulle (les $n-1$ dernières devant alors former une famille libre).

De même, en multipliant les deux membres de l'égalité (2) par ${^t}V$ à gauche, on obtient ${^t}UA = 0$ et donc les colonnes de la matrice $A$ sont orthogonales à $U$ (pour le produit scalaire usuel) ce qui invite franchement à considérer la matrice $A=\left(
\begin{array}{ccccc}
0&-2&\ldots&\ldots&-n\\
0&1&0&\ldots&0\\
\vdots&\ddots&\ddots&\ddots&\vdots\\
\vdots& &\ddots&\ddots&0\\
0&\ldots&\ldots&0&1
 \end{array}\right)$ qui convient.
\fincorrection
\correction{005687}
Soit $P=X^3+X^2+X=X(X-j)(X-j^2)$. $P$ est à racines simples dans $\Cc$ et annulateur de $A$. Donc $A$ est diagonalisable dans $\Cc$ et ses valeurs propres sont à choisir dans $\{0,j,j^2\}$. Le polynôme caractéristique de $A$ est de la forme $(-1)^nX^\alpha(X-j)^\beta(X-j^2)^\gamma$ avec $\alpha+\beta+\gamma=n$. De plus, $A$ est réelle et on sait que $j$ et $j^2=\overline{j}$ ont même ordre de multiplicité ou encore $\gamma=\beta$.

Puisque $A$ est diagonalisable, l'ordre de multiplicité de chaque valeur propre est égale à la dimension du sous-espace propre correspondant et donc

\begin{center}
$\text{rg}(A)=n-\text{dim}(\text{Ker}A)=n-\alpha=2\beta$.
\end{center}

On a montré que $\text{rg}A$ est un entier pair.
\fincorrection
\nocorrection
\nocorrection
\correction{002564}
 Soit $E$ un espace vectoriel sur un corps $K$ ($K=\R$ ou $\C$), on appelle {\it projecteur}
un endomorphisme $p$ de $E$ v\'erifiant $p\circ p=p$. Soit $p$ un projecteur.

\begin{enumerate}  \item {\it Montrons que $\mathrm{Id}_E-p$ est un projecteur et calculons $p\circ(\mathrm{Id}_E-p)$ et $(\mathrm{Id}_E-p)\circ p$.}

On a $(\mathrm{Id}_E-p)\circ(\mathrm{Id}_E-p)=\mathrm{Id}_E-p-p+p^2=\mathrm{Id}_E-p$, car $p^2=p$, ce qui prouve que $\mathrm{Id}_E-p$ est un projecteur.

Par ailleurs, on a
$$p\circ(\mathrm{Id}_E-p)=p-p^2=p-p=0=(\mathrm{Id}_E-p)\circ p$$
donc pour tout $\vec x\in E$, on a $p(\vec x-p(\vec x))=\vec 0$.


  \item {\it Montrons que pour tout $\vec x\in \Im p$, on a $p(\vec x)=\vec x$.}

Soit $\vec x\in \Im p$, il existe $\vec y\in E$ tel que $\vec x=p(\vec y)$, on a donc 
$p(\vec x)=p^2(\vec y)=p(\vec y)=\vec x$.

  \item{\it On en d\'eduit que $\Im p$ et $\ker p$ sont suppl\'ementaires.}

Soit $\vec x\in E$, on peut \'ecrire $\vec x=p(\vec x)+\vec x -p(\vec x)$, consid\'erons $\vec x -p(\vec x)$, on a
$p(\vec x-p(\vec x))=0$

ce qui prouve que $\vec x -p(\vec x)\in \ker p$. Ainsi tout \'el\'ement de $E$ s'\'ecrit comme somme d'un 
\'el\'ement de $\Im p$, $p(\vec x)$, et d'un \'el\'ement de $\ker p$, $\vec x -p(\vec x)$, il nous reste \`a 
d\'emontrer que la somme est directe.

Soit $\vec x\in \Im p\cap\ker p$, on a, d'une part $p(\vec x)=\vec x$ d'apr\`es la question $2)$ car $\vec x\in \Im p$
et, d'autre part $p(\vec x)=\vec 0$ car $\vec x\in \ker p$, d'o\`u $\vec x=\vec 0$. On a donc 
$$E=\Im p\oplus\ker p.$$
(Sachant que $\dim E=\dim\ker p+\dim\Im p$, on pouvait se contenter de d\'emontrer que $\Im p\cap \ker p={\vec 0}$, ici nous avons
explicitement la d\'ecomposition.)


  \item{\it Montrons que le rang de $p$ est \'egal \`a la trace de $p$.}

Notons $n$ la dimension de $E$ et consid\'erons une base de $E$ de la forme
$$(\vec{e_1},\cdots,\vec{e_k},\vec {e_{k+1}},\cdots,\vec{e_n})$$
o\`u $(\vec{e_1},\cdots,\vec{e_k})$ est une base de $\Im p$ et $(\vec {e_{k+1}},\cdots,\vec{e_n})$ une base de $\ker p$.
dans une telle base, la matrice de $p$ s'\'ecrit
$$M=\begin{pmatrix}I_k&0 \\ 0&0\end{pmatrix}$$
o\`u $I_k$ d\'esigne la matrice identit\'e $k\times k$, et les $0$ des blocs de z\'eros. Le rang de $p$ est \'egal \`a la 
dimension de $\Im p$ c'est-\`a-dire ici \`a $k$ et on a bien $k=\mathrm{Tr} M=\mathrm{Tr} p.$
\end{enumerate}
\fincorrection
\nocorrection
\nocorrection
\nocorrection
\nocorrection
\nocorrection
\nocorrection
\nocorrection
\nocorrection
\nocorrection
\nocorrection
\nocorrection
\nocorrection
\nocorrection
\nocorrection
\nocorrection
\nocorrection
\nocorrection
\nocorrection
\nocorrection
\nocorrection
\nocorrection
\nocorrection
\nocorrection
\nocorrection
\nocorrection
\nocorrection
\nocorrection
\nocorrection
\nocorrection
\correction{001554}
  \begin{enumerate}
  \item 
    $<u_{k}(x),a>=k<x,a><a,a>+<x,a>=(k+1)<x,a>$ donc
    $x=\frac{-k}{k+1}<u_{k}(x),a>a+u_{k}(x)$. On en déduit que $u_{k}$
    est inversible, et que $u_{k}^{-1}=u_{\frac{-k}{k+1}}$.

  \item
    L'adjoint d'un endomorphisme $u$ est l'unique endomorphisme $v$ qui
    satisfait~: $\forall(x,y)\in E^{2}, <u(x),y>=<x,u(y)>$. Or
    $<u_{k}(x),y>=k<x,a><y,a>+<x,y>=<x,u_{k}(y)>$. Donc $u_{k}$ est égal
    à son adjoint.

  \item
    Si $u_{k}$ est orthogonal, on doit avoir $\Vert u_{k}(a)\Vert=\Vert
    a\Vert=1$, soit $|k+1|=1$. Ainsi $k=0$ ou $k=-2$.

    Pour $k=0$, $u_{k}=\mathrm{id}$ est bien orthogonal. Pour $k=-2$,
    $u_{-2}^{-1}=u_{\frac{-2}{-2+1}}=u_{-2}={}^t u_{-2}$. Donc
    $u_{-2}$ est bien orthogonal. Il s'agit de la symétrie orthogonale
    par rapport à l'hyperplan $\{a\}^{\bot}$


  \item
    Si $k=0$, 1 est la seule valeur propre et $E_{1}=E$

    Si $k\neq 1$, $\forall x\in\{a\}^{\bot}, u_{k}(x)=x$ donc $1$ est
    valeur propre de multiplicité au moins $n-1$. De plus
    $u_{k}(a)=(k+1)a$ donc $(k+1)$ est valeur propre. Finalement, $1$
    est valeur propre de multiplicité exactement $n-1$, avec pour espace
    propre $\{a\}^{\bot}$, et $k+1$ est valeur propre simple avec espace
    propre $\R a$.
   \end{enumerate}
\fincorrection
\nocorrection
\nocorrection
\nocorrection
\nocorrection
\nocorrection
\nocorrection
\nocorrection
\nocorrection
\nocorrection
\nocorrection
\nocorrection
\nocorrection
\nocorrection
\nocorrection
\nocorrection
\nocorrection
\nocorrection
\nocorrection
\nocorrection
\nocorrection
\nocorrection
\nocorrection
\nocorrection
\correction{003568}
\begin{enumerate}
  \item 
  \item 
  \begin{enumerate}
    \item Pour $p\in K[X]$ on a $P(\Phi_u) = v \mapsto v\circ P(u)$ donc $u$ et $\Phi_u$ ont
    mêmes polynômes annulateurs.
    \item $(\lambda\in\mathrm{Spec}(\Phi_u)) \Leftrightarrow (\exists\ v\ne 0\text{ tq }v\circ(u-\lambda\mathrm{id}_E) = 0)
    \Leftrightarrow (u-\lambda\mathrm{id}_E\text{ n'est pas surjectif}) \Leftrightarrow (\lambda\in\mathrm{Spec}(u))$.
    Ainsi $\Phi_u$ et $u$ ont même spectre.
    Si $\lambda\in\mathrm{Spec}(u)$ et $v\in\mathcal{L}(E)$ on a~:
    $(\Phi_u(v) = \lambda v) \Leftrightarrow (\Im(u-\lambda\mathrm{id}_E) \subset\mathrm{Ker} v)$ donc
    $\mathrm{Ker}(\Phi_u-\lambda\mathrm{id}_{\mathcal{L}(E)})$ est isomorphe à $\mathcal{L}({H,E})$ où $H$ est un supplémentaire
    de~$\Im(u-\lambda\mathrm{id}_E)$.
    On en déduit~:
    $\dim(\mathrm{Ker}(\Phi_u-\lambda\mathrm{id}_{\mathcal{L}(E)})) = \dim(E)\dim(\mathrm{Ker}(u-\lambda\mathrm{id}_E)$.
  \end{enumerate}
\end{enumerate}
 \fincorrection
\nocorrection
\correction{003570}
$\lambda = 1$ : Dir$(p) \subset \mathrm{Ker} f$,\quad $\Im f \subset \text{Base}(p)$.

$\lambda = 0$ : $f\bigl(\text{Base}(p)\bigr) \subset \text{Dir}(p)$.
\fincorrection
\nocorrection
\correction{003572}
\begin{enumerate}
  \item Pour $P\in K[X]$ on a $P(u)\circ v - v\circ P(u) = P'(u)$.
\end{enumerate}
 \fincorrection
\nocorrection
\correction{003574}
Supposons qu'il existe $g\in\mathcal{L}(E)$ tel que $\mathrm{rg}(f\circ g - g\circ f) = 1$.
Alors il existe $\ell\in E^*$ et $a\in E$ tous deux non nuls tels que~:
$$\forall\ x\in E,\ f(g(x)) - g(f(x)) = \ell(x)a.$$
D'où par récurrence sur~$k$~:
$$\forall\ x\in E,\ f^k(g(x)) - g(f^k(x)) = \ell(x)f^{k-1}(a) + \ell(f(x))f^{k-2}(a) + \dots + \ell(f^{k-1}(x))a.$$

Comme $\chi_f$ est irréductible, le sous-espace $f$-monogène
engendré par~$a$ est égal à~$E$, soit~: $(a,f(a),\dots,f^{n-1}(a))$ est
une base de~$E$ avec $n=\dim E$ et $f^n(a) = \alpha_0a + \dots + \alpha_{n-1}f^{n-1}(a)$.
Alors $\mu_f(f) = f^n-\alpha_{n-1}f^{n-1} - \dots - \alpha_0f^0 = 0$ et~:
$$\forall\ x\in E,\ 0 = \mu_f(f)(g(x)) - g(\mu_f(f)(x)) = \ell(x)f^{n-1}(a) + \dots + \ell(f^{n-1}(x) - \dots - \alpha_1x)a.$$
Ceci implique $\ell(x) = 0$ pour tout~$x$, en contradiction avec l'hypothèse
$\mathrm{rg}(f\circ g - g\circ f) = 1$.
\fincorrection
\correction{003575}
\begin{enumerate}
  \item Oui, les applications $u \mapsto p\circ u$ et $u \mapsto u\circ p$ le sont
    (ce sont des projecteurs) et elles commutent.
  \item Soit $\cal B$ une base de~$E$ obtenue par concaténation d'une
    base de~$\mathrm{Ker} p$ et d'une base de~$\Im p$.

    Si $\mathrm{mat}_{\cal B}(u) = \left(\begin{smallmatrix}A&B\cr C&D\cr\end{smallmatrix}\right)$
    alors $\mathrm{mat}_{\cal B}(\varphi(u)) = \left(\begin{smallmatrix}A&B/2\cr C/2&D\cr\end{smallmatrix}\right)$,
    d'où $\mathrm{Spec}(\varphi) \subset \{0,\frac12,1\}$ et $d_0 = (n-r)^2$,
    $d_1=r^2$ et $d_{1/2} = 2r(n-r)$.
\end{enumerate}
 \fincorrection
\correction{003576}
Si $D$ est diagonalisable alors les applications $X \mapsto DX$ et $X \mapsto XD$
le sont (annulateur scindé à racines simples) et elles commutent, donc elles sont
simmultanément diagonalisables et leur différence, $\phi_D$, est aussi diagonalisable.

Pour la réciproque, on commence par constater que si $P$ est un polynôme
quelconque, alors~:
$$\forall\ X\in\mathcal{M}_n(\C),\ P(\phi_D)(X) = \sum_{k=0}^{\deg(P)} (-1)^kD^kX\frac{P^{(k)}(D)}{k!}
= \sum_{k=0}^{\deg(P)} (-1)^k\frac{P^{(k)}(D)}{k!}XD^k.$$
(formule du binôme pour $P=X^m$ et linéarité de chaque membre par rapport à~$P$
pour $P$ quelconque).

Supposons $\phi_D$ diagonalisable, prenons~$P$ annulateur scindé à racines simples de~$\phi_D$,
$X=U^tV$ où $U$ est un vecteur propre de~$D$
associé à une certaine valeur propre~$\lambda$ et $V$ un vecteur arbitraire. Donc~:
$$0 = \sum_{k=0}^{\deg(P)} (-1)^k\lambda^kU^tV\frac{P^{(k)}(D)}{k!}
= U^tV\sum_{k=0}^{\deg(P)} (-1)^k\lambda^k\frac{P^{(k)}(D)}{k!}
= U^tVP(D-\lambda I).$$
Comme $U\ne 0$, ceci implique ${}^tVP(D-\lambda I) = 0$ pour tout~$V$,
donc $P(D-\lambda I) = 0$. Ainsi $D-\lambda I$ est diagonalisable et $D$ itou.
\fincorrection
\nocorrection
\correction{003626}
\begin{enumerate}
  \item 
  \item $( (-2,0,1),\ (0,3,-2),\ (1,-2,1) )$.
\end{enumerate}
 \fincorrection
\correction{003627}
Base ssi $n$ est impair, $2\vec e_1 = (1,1,-1,1,-1,\dots,1,-1)$ et
         les autres vecteurs s'obtiennent par rotation :\par
         $2\vec e_2 = (-1,1,1,-1,1,-1,\dots,1)$.
\fincorrection
\nocorrection
\correction{003629}
2. $\phi_i^* = (1-2d_i(X-x_i))P_i^2$,
             $\psi_i^* = (X-x_i)P_i^2$.
\fincorrection
\correction{003630}
  2. $\frac18(9-15X^2,\ 75X-105X^3,\ -15+45X^2,\ -105X+175X^3)$.
\fincorrection
\correction{003631}
$M = \begin{pmatrix} 1   &1   &1   &b-a         \cr
                        a   &b   &c   &(b^2-a^2)/2 \cr
                        a^2 &b^2 &c^2 &(b^3-a^3)/3 \cr
                        a^3 &b^3 &c^3 &(b^4-a^4)/4 \cr \end{pmatrix}$ et
         $\det(M) = (b-a)^4(c-a)(c-b)\frac{2c-a-b}{12}$,
         donc la famille est libre si et seulement si $c \ne \frac {a+b}2$.
\fincorrection
\correction{003632}
2. terme dominant $ \Rightarrow  P_n^*(Q_i) = 1$,
             donc $P_n^* = \sum_{i=0}^n \frac{f_i}{Q_i(i)}
                         = \sum_{i=0}^n \frac{(-1)^{n-i}f_i}{i!\,(n-i)!}$.

3. $P_k^* = \sum_{i=0}^k \frac{(-1)^{k-i}f_i}{i!\,(k-i)!}$.

 \fincorrection
\correction{003633}
2. $P_0^* = \frac{f_a}{(b-a)^2}$,
             $P_1^* = \frac{f_a+f_b-4f_c}{(b-a)^2}$,
             $P_2^* = \frac{f_b}{(b-a)^2}$.
 \fincorrection
\nocorrection
\nocorrection
\nocorrection
\nocorrection
\nocorrection
\correction{003639}
1. $\frac{(k+1)(k+2)}2$.

3. $(\delta_A,\delta_B,\delta_C,\delta_{A'},\delta_{B'},\delta_{C'})$
             où $A',B',C'$ sont les milieux du triangle $ABC$.
             
5. $ \int_{}^{}\!\! \int_T^{} f(x,y)\,d xd y =
              \frac{f(A')+f(B')+f(C')}6$.
 \fincorrection
\correction{003640}
3. Rmq : coefficients de Fourrier :
    $\alpha_p = \frac 2{2n+1}\sum_{k=0}^{2n} f(a_k)\cos(pa_k)$ et
    $\beta_p = \frac 2{2n+1}\sum_{k=0}^{2n} f(a_k)\sin(pa_k)$.
 \fincorrection
\nocorrection
\correction{003642}
2. $\left(\frac{\varphi^{n-1} - {\bar\varphi}^{n-1}}{\varphi-\bar\varphi},
               \frac{\varphi^n - {\bar\varphi}^n}{\varphi-\bar\varphi}\right)$
             avec $\varphi=\frac{1+\sqrt5}2$, $\bar\varphi=\frac{1-\sqrt5}2$.
 \fincorrection
\nocorrection
\nocorrection
\nocorrection
\nocorrection
\nocorrection
\nocorrection
\nocorrection
\correction{003650}
\leavevmode\hbox{\vtop{\halign{&$#$\hfil\quad\cr
{\vec e}_i \leftrightarrow {\vec e}_j                   : &e_i^* \leftrightarrow e_j^*. \cr
{\vec e}_i \leftarrow \alpha{\vec e}_i              : &e_i^* \leftarrow e_i^*/\alpha. \cr
{\vec e}_i \leftarrow {\vec e_i} + \alpha{\vec e}_j : &e_j^* \leftarrow e_j^* - \alpha e_i^*. \cr
}}}
\fincorrection
\nocorrection
\nocorrection
\nocorrection
\nocorrection
\nocorrection
\nocorrection
\nocorrection
\nocorrection
\nocorrection
\correction{001311}
\begin{enumerate}
  \item Oui.
  \item Non. Le seul élément qui peut être l'élément neutre
  est $1$ qui n'appartient pas à l'ensemble.
  \item Non. $0$ n'a pas d'inverse.
  \item Oui.
\end{enumerate}
\fincorrection
\nocorrection
\nocorrection
\correction{001314}
Le premier ensemble n'est pas un groupe car, par exemple, la
matrice $\begin{pmatrix}2&0\\0&2\\ \end{pmatrix}$ ne peut avoir
pour inverse que $\begin{pmatrix} \frac{1}{2}&0\\0&\frac{1}{2}\\
\end{pmatrix}$ qui n'appartient pas à l'ensemble.

Notons $G = \{ M \in \mathcal{M}_{2}(\Zz) : \det M = 1 \}$ et
montrons que $G$ est un sous-groupe de $Gl(2,\Rr)$.
\begin{itemize}
  \item la matrice identité appartient à $G$.
  \item si $A,B \in G$ alors $AB \in \mathcal{M}_{2}(\Zz)$ et
$\det AB = \det A \times \det B = 1\times 1 =1$, et donc $AB \in
G$.
  \item Si $A = \begin{pmatrix} a&b\\c&d\\ \end{pmatrix}$ ($a,b,c,d \in \Zz$) alors
$\frac{1}{\det A}\begin{pmatrix} d&-b\\-c&a\\ \end{pmatrix}
=\begin{pmatrix} d&-b\\-c&a\\ \end{pmatrix}$ appartient à $G$ et
est l'inverse de  $A$.
\end{itemize}
\fincorrection
\nocorrection
\nocorrection
\nocorrection
\nocorrection
\nocorrection
\nocorrection
\nocorrection
\correction{001322}
\begin{enumerate}
    \item L'ensemble  $G$  des matrices  $\begin{pmatrix}
a & c \cr b & d \cr \end{pmatrix}$  avec  $a,b,c,d\in {\R}$   tels
que  $ad-bc\not = 0 $  et  $a^2-b^2-c^2-d^2 \leq 1$  n'est pas un
sous-groupe de  $Gl_2({\R})$. En effet les deux matrices
$\begin{pmatrix} 1 & 1\cr 0 & 1/2 \cr \end{pmatrix}$  et
$\begin{pmatrix} 1 & 0 \cr 1 & 1/2\cr \end{pmatrix}$
appartiennent \`a  $G$  et leur produit  $\begin{pmatrix} 2 & 1/2
\cr 1/2 & 1/4\cr \end{pmatrix}$ n'appartient pas \`a  $G$.
    \item L'ensemble  $H$  des matrices  $\begin{pmatrix}
a & b \cr 0 & a^{-1} \cr \end{pmatrix}$  avec  $a\in {\R}^*$  et
$b\in {\R}$  est un sous groupe de  $Gl_2({\R})$. En effet,

- $I_2$  \' el\' ement neutre de  $Gl_2({\R})$  appartient \`a
$H$.

- Soient  $M=\begin{pmatrix} a & b \cr 0 & a^{-1} \cr
\end{pmatrix}$  et  $M'=\begin{pmatrix} c & d \cr 0 & c^{-1} \cr
\end{pmatrix}$ deux \' el\' ements de  $H$  alors
$MM'=\begin{pmatrix} ac & ad+bc^{-1} \cr 0 & (ac)^{-1} \cr
\end{pmatrix}$  donc le produit de deux \' el\' ements de  $H$
appartient \`a  $H$.

- Soit  $M=\begin{pmatrix} a & b \cr 0 & a^{-1} \cr
\end{pmatrix}$. Alors  $M^{-1} =\begin{pmatrix} a^{-1} &- b \cr 0
& a \cr \end{pmatrix}$  appartient \`a  $H$.
    \item Soit  $K_M$  l'ensemble des matrices  $\begin{pmatrix}
a & c \cr b & d \cr \end{pmatrix}$  avec  $a,b,c,d\in {\R}$  tels
que   $ad-bc\not = 0 $  et  $ a\leq M$. Nous allons montrer, en
raisonnant par l'absurde, qu'il n'existe pas de valeur  $M\in
{\R}$  telle que  $K_M$  forme un sous-groupe de  $Gl_2({\R})$.

Soit  $M\in {\R}$  tel que  $K_M$  forme un sous-groupe de
$Gl_2({\R})$. Alors  $I_2$  appartient \`a  $K_M$  donc  $M\geq
1$. Ainsi, les matrices  $A=\begin{pmatrix} 1 & 1 \cr 0 & 1 \cr
\end{pmatrix}$  et, pour tout  $n\in {\N}$,  $A_n=\begin{pmatrix}
1 & 1 \cr n & 1 \cr \end{pmatrix}$  appartiennent \`a  $K_n$  donc
le produit  $AA_n=\begin{pmatrix} 1+n & 0 \cr 0 & 1 \cr
\end{pmatrix}$  appartient \`a  $K_n$. En cons\' equence, pour
tout  $n\in {\N}$, on a : $ 1+n\leq M$, ce qui est absurde.
\end{enumerate}
\fincorrection
\correction{001323}
\begin{itemize}
    \item[$\bullet$] Si $H \subset K$ alors $H\cup K = K$, qui est un sous-groupe
de H. M\^eme chose si $K \subset H$.
    \item[$\bullet$] R\'eciproquement, supposons que
$H\cup K$ est un sous-groupe de $G$. Par l'absurde supposons que
$H \not\subset K$ et $K \not\subset H$. Alors il existe $x\in
H\setminus K$ et $y \in K\setminus H$. Comme $x,y \in H\cup K$ et
que $H\cup K$ est un groupe alors $x.y \in H\cup K$. Donc $x.y \in
H$ ou $x.y \in K$. Par exemple supposons $x.y \in H$ alors comme
$x\in H$, $x^{-1}\in H$ et donc comme $H$ est un groupe
$x^{-1}.x.y \in H$ et donc $y\in H$. Ce qui est en contradiction
avec l'hypoth\`ese $ y \in K\setminus H$. En conclusion, parmi les
sous-groupes $H,K$ l'un est inclus dans l'autre.
\end{itemize}
\fincorrection
\nocorrection
\nocorrection
\correction{001326}
Soit $G = \langle a,b \rangle$, tout élément $g$ de $G$ S'écrit $g
= a^{\alpha_1}b^{\beta_1}  a^{\alpha_2}b^{\beta_2}\ldots
a^{\alpha_n}b^{\beta_n}$ avec $\alpha_i,\beta_i \in \Zz$. Si $h
\in \langle a \rangle \cap \langle b \rangle$, alors en
particulier $h \in \langle a \rangle$ et $h = a^\mu$ avec $\mu \in
\Zz$, donc $h$ commute avec $a^{\alpha_i}$ pour tout $\alpha_i$
dans $\Zz$ (en effet $a^{\alpha_i}a^{\mu} = a^{\alpha_i+\mu}=
a^\mu a^{\alpha_i}$. De même $h \in \langle b \rangle$ donc $h$
s'écrit également $h=b^\nu$ ($\nu \in \Zz$) et $h$ commute avec
$b^{\beta_i}$. Donc $hg = (h a^{\alpha_1})b^{\beta_1}\ldots =
(a^{\alpha_1}h)b^{\beta_1}\ldots = a^{\alpha_1}(h
b^{\beta_1})\ldots =  a^{\alpha_1}(b^{\beta_1}h)\ldots =\cdots$
Finalement $hg = a^{\alpha_1}b^{\beta_1}\ldots
a^{\alpha_n}b^{\beta_n}h = gh$. Ainsi $h$ commute  avec tout
élément de $G$ et appartient ainsi au centre de $G$.
\fincorrection
\nocorrection
\nocorrection
\nocorrection
\nocorrection
\nocorrection
\nocorrection
\nocorrection
\nocorrection
\correction{001343}
Soit $f:(\Zz,+) \longrightarrow (\Zz,+)$ un morphisme de groupe.
Comme tout morphisme $f$ v\'erifie $f(0)=0$. Notons $a = f(1)$.
Alors
$$f(2) = f(1+1) = f(1)+f(1) = a+a = 2.a.$$
De m\^eme, pour $n \ge 0$ :
$$f(n) = f(1+\cdots+1)= f(1)+\cdots+f(1) = n.f(1) = n.a.$$
Enfin comme
$$0= f(0) = f(1+(-1)) = f(1) + f(-1) = a + f(-1),$$
alors $f(-1) = -a$ et pour tout $n \in \Zz$ :
$$f(n) = n.a.$$
Donc tous les morphisme sont de la forme $n \mapsto n.a$, avec
$a\in\Zz$.

\bigskip

Un morphisme $n \mapsto n.a$ est injectif si et seulement si $a
\not= 0$, et surjectif si et seulement si $n = \pm 1$.
\fincorrection
\nocorrection
\correction{001345}
\begin{align*}
f : (\Rr,+)& \longrightarrow (\Cc^*,\times) \\
x& \mapsto e^{ix}\\
\end{align*}

V\'erifions que $f$ est un morphisme de groupe. Soit $x,y \in
\Rr$, alors
$$f(x+y) = e^{i(x+y)}=e^{ix}e^{iy}=f(x)\times f(y),$$
et
$$f(x^{-1})= e^{i(-x)} = \frac{1}{e^{ix}}={f(x)}^{-1}.$$
Donc $f$ est un morphisme de groupe.

Montrons que $f$ n'est pas injective en prouvant que le noyau
n'est pas r\'eduit \`a $0$ :
$$ \text{Ker}\, f = \left\lbrace x\in\Rr \text{ tels que } f(x)=1 \right\rbrace
   = \left\lbrace x\in\Rr \text{ tels que } e^{ix}=1 \right\rbrace
= \left\lbrace x = 0+2k\pi,\ \ k\in\Zz \right\rbrace. $$

Enfin
$$\text{Im}\, f =  \left\lbrace y\in\Cc^*, y = e^{ix} \right\rbrace$$
est l'ensemble des complexes de module $1$, c'est-\`a-dire le
cercle de centre $0$ et de rayon $1$.
\fincorrection
\nocorrection
\nocorrection
\nocorrection
\nocorrection
\nocorrection
\nocorrection
\nocorrection
\nocorrection
\correction{001354}
Soit $\phi : \Cc^* \longrightarrow \Rr^*$ un morphisme entre les
deux groupes multiplicatifs $\Cc^*$ et $\Rr^*$. Notons $a =
\phi(i) \in \Rr^*$. Alors $\phi(-1)=\phi(i^2)=\phi(i)^2=a^2$, de
même $1=\phi(1)=\phi( (-1)^2 )=\phi(-1)^2=a^4$ ; donc $a^4 = 1$ et
nécessairement $a^2 = 1$. Le morphisme $\phi$ n'est pas injectif
car $\phi(1)= \phi(-1) = 1$, \emph{a fortiori} $\phi$ n'est pas un
isomorphisme.
\fincorrection
\nocorrection
\correction{001386}
Soit $x\not= e$ un \'el\'ement de $G$, soit $H = \{ e,x,x^2,\ldots
\}$ le sous-groupe engendr\'e par $x$. $H$ est un sous-groupe de
$G$ donc $\mathrm{Card} H$ divise $\mathrm{Card} G=p$ qui un nombre premier. En
cons\'equent $\mathrm{Card} H = 1$ ou $p$ mais $H \not= \{ e \}$ donc
$\mathrm{Card} H =p$ et $H=G$.

Nous venons de montrer que $G$ est engendr\'e par $x$ donc $G$
est cyclique, de plus le raisonnement est valide quelque soit
$x\not= e$ alors tout \'el\'ement de $G\setminus \{ e \}$ est un
g\'en\'erateur de $G$.
\fincorrection
\correction{001387}
\begin{enumerate}
    \item $H\cap H'$ est un sous-groupe  de $H$ donc
$\mathrm{Card} H\cap H'$  divise $\mathrm{Card} H = p$. Or $p$ est premier donc
$\mathrm{Card} H\cap H' = 1$ ou $p$. Mais $H\cap H'\not= H$ donc $\mathrm{Card}
H\cap H' \not= p$ et donc $ H\cap H' = \{ e \}$.
    \item Soit $E$ l'ensemble des \'el\'ements
d'ordre $p$ que l'on suppose non vide. Notons que pour $x\in E$ le
sous-groupe $H_x$ engendr\'e par $x$ est d'ordre $p$ et de plus
tout $z\in H_x\setminus \{ e \}$ est d'ordre $p$ car $H_x$ est
cyclique et $p$ est premier. Donc $H_x$ contient $p-1$ \'el\'ement
d'ordre $p$.

Si $E$ ne contient qu'un seule \'el\'ement $x$ alors $E = H_x
\setminus\{ e \}$ et donc $E$ contienet $p-1$ \'el\'ements.

Sinon, soit $x,y \in  E$ avec $x\not=y$. Alors d'apr\`es la
premi\`ere question  $H_x \cap H_y = \{ e \}$. Donc $E$ se
d\'ecompose en une union disjointe de $H_x \setminus \{ e \}$.
Donc $\mathrm{Card} E$ est multiple de $p-1$.
\end{enumerate}
\fincorrection
\nocorrection
\correction{001389}
\begin{enumerate}
    \item Notons d'abord que pour $x\in G$ $x^2=e$ et donc
$x^{-1}=x$. Soit maintenant $x,y \in G$. Alors $xy\in G$ et
$(xy)^2=e$ donc $xy=(xy)^{-1}$ et par suite $xy=y^{-1}x^{-1}=yx$
car $x$ et $y$ sont d'ordre $2$. Le produit de deux \'el\'ements
quelconques de $G$ commute donc $G$ est commutatif.
    \item Notons $E$ l'ensemble des \'el\'ements d'ordre $2$.
$$E = \{ x\in G\ / \ x^2=e \text{ et } x \not=e \} =
\{ x\in G\ / \ x=x^{-1} \text{ et } x \not=e \}.$$ Par l'absurde
supposons que $H$ est l'ensemble vide. Alors quelque soit
$x\not=e$ dans $G$ $x\not= x^{-1}$. Donc nous pouvons d\'ecomposer
$G\setminus \{ e \}$ en deux ensembles disjoints $F = \{
x_1,\ldots,x_n\}$ et $F' =\{{x_1}^{-1}, \ldots,{x_n}^{-1} \}$ qui
sont de m\^eme cardinal $n$. Donc le cardinal de $G$ est $2n+1$
(le $+1$ provient de l'\'el\'ement neutre). Ce qui contredit
l'hypoth\`ese << $G$ d'orde pair >>.
\end{enumerate}
\fincorrection
\nocorrection
\nocorrection
\nocorrection
\nocorrection
\nocorrection
\nocorrection
\nocorrection
\nocorrection
\nocorrection
\nocorrection
\nocorrection
\correction{002966}
\begin{enumerate}
 \item Non, $a$ n'est pas r{\'e}gulier.
 \item Oui, $G \approx \Z/3\Z$.
 \item Non, il n'y a pas d'{\'e}l{\'e}ment neutre.
\end{enumerate}

\fincorrection
\nocorrection
\nocorrection
\nocorrection
\nocorrection
\nocorrection
\nocorrection
\nocorrection
\nocorrection
\nocorrection
\nocorrection
\nocorrection
\nocorrection
\nocorrection
\nocorrection
\nocorrection
\nocorrection
\nocorrection
\nocorrection
\nocorrection
\nocorrection
\nocorrection
\nocorrection
\nocorrection
\nocorrection
\nocorrection
\nocorrection
\nocorrection
\nocorrection
\nocorrection
\nocorrection
\nocorrection
\correction{007360}
\begin{enumerate}
 \item Dans $\mathcal{S}_7$, $(1,2)(1,3)=(132)$ et par conjugaison $(1,2)(2,3)(1,2)=(1,3)$.
 \item L'ensemble des transpositions de $\mathcal{S}_7$ ne contient pas l'identité : ce n'est donc pas un sous-groupe de $\mathcal{S}_7$. Le premier calcul montre qu'il n'est pas non plus stable par produit.
\end{enumerate}
\fincorrection
\nocorrection
\nocorrection
\correction{007363}
\begin{enumerate}
\item Les ordres possibles des éléments d'un groupe d'ordre $6$ sont, par le théorème de Lagrange les diviseurs de $6$ c'est à dire $1,2,3$ et $6$.
\item Les éléments inversibles de $(\Z/7\Z,\times)$ sont les classes des entiers premiers à $7$ c'est à dire toutes les classes non nulles.
\item la table de multiplication de $((\Z/7\Z)^\star,\times)$ est


\begin{center}
\begin{tabular}{|c|c|c|c|c|c|c|} 
\hline
$\times$ &1&2&3&4&5&6\\
\hline
1&1&2&3&4&5&6\\
2&2&4&6&1&3&5\\
3&3&6&2&5&1&4\\
4&4&1&5&2&6&3\\
5&5&3&1&6&4&2\\
6&6&5&4&3&2&1\\
\hline
\end{tabular}
\end{center}


\item On calcule l'ordre multiplicatif de $2$. Comme $2^3=1$ l'ordre de $2$ est $3$. Comme $3^2=2$ et $3^3=6$, $3$ est d'ordre 6. C'est donc un générateur
 de $((\Z/7\Z)^\star,\times)$.
\item On déduit de la question précédente que $((\Z/7\Z)^\star,\times)$ est un groupe cyclique d'ordre $6$. Il a donc $\phi(6)=\phi(2)\phi(3)=2$ générateurs.
\item Dans le groupe $\Z/2\Z\times \Z/4\Z$ tous les éléments $g$ vérifient $4g=0$. Il n'y a donc aucun élément d'ordre $8$. Ce groupe n'est donc pas cyclique.
\item Par le théorème chinois, puisque $4$ et $5$ sont premiers entre eux, 
$\Z/20\Z$ est isomorphe à $\Z/4\Z\times \Z/5\Z$.
Le groupe  $(\Z/20\Z)^\star$ est donc isomorphe au produit 
$(\Z/4\Z)^\star\times (\Z/5\Z)^\star$. Comme $(\Z/4\Z)^\star$ est un groupe d'ordre $2$, il est isomorphe à $\Z/2\Z$. Comme $5$ est premier, $(\Z/5\Z)^\star$ est un groupe cyclique d'ordre $4$ donc isomorphe à $\Z/4\Z$.
Par conséquent, le groupe  $(\Z/20\Z)^\star$ n'est pas cyclique.
\end{enumerate}
\fincorrection
\correction{007364}
\begin{enumerate}
 \item Comme $\alpha$ et $\beta$ sont à support disjoint, ils commutent $\alpha\beta=\beta\alpha$
 \item L'ordre de $\alpha$ donner comme produit de cycles à support disjoint est le $ppcm$ des ordres des cycles qui le composent, donc $ppcm(3,2)=6$. Celui de la transposition $\beta$ est $2$.
 \item $\alpha^{-1}=(175)(43)$ et $\beta^{-1}=(26)$.
 \item Montrons que $S=\{\alpha^i\circ\beta^j, 0\leq i \leq 5, 0\leq j\leq 1\}$ est un sous-groupe de $\mathcal{S}_7$.


*stabilité par produit : Soit $\alpha^i\circ\beta^j$ et $\alpha^k\circ\beta^l$ deux éléments de $S$ où $0\leq i \leq 5, 0\leq j\leq 1$ et $0\leq k \leq 5, 0\leq l\leq 1$.
\begin{eqnarray*}
 (\alpha^i\circ\beta^j)\circ (\alpha^k\circ\beta^l)&=&\alpha^i\circ\beta^j\circ \alpha^k\circ\beta^l\\
&=&\alpha^i\alpha^k\beta^j\beta^l=\alpha^{i+k}\beta^{j+l}
\end{eqnarray*}
Comme $\alpha^6=Id$, en considérant le reste $0\leq r\leq 5$ de la division euclidienne de $i+k$ par $6$ et celui $0\leq s\leq 1$ de la division euclidienne de $j+l$ par $2$
on obtient $(\alpha^i\circ\beta^j)\circ (\alpha^k\circ\beta^l)=\alpha^r\beta^s$ qui appartient donc à $S$.


*stabilité par inversion :
Soit $\alpha^i\circ\beta^j$ dans $S$ où $0\leq i \leq 5, 0\leq j\leq 1$.
Comme $\alpha^6=Id$, $(\alpha^i)^{-1}=\alpha^{6-i}$. Comme $\beta^2=Id$, $\beta^{-1}=\beta$.
$$(\alpha^i\circ\beta^j)^{-1}=(\beta^j)^{-1}(\alpha^i)^{-1}=\beta^j\alpha^{6-i}=
\alpha^{6-i}\beta^j$$ qui appartient donc à $S$.

 \item Soit $\alpha^i\circ\beta^j$ et $\alpha^k\circ\beta^l$ deux éléments de $S$ où $0\leq i \leq 5, 0\leq j\leq 1$ et $0\leq k \leq 5, 0\leq l\leq 1$.
Si $\alpha^i\circ\beta^j=\alpha^k\circ\beta^l$, alors $\alpha^i\alpha^{6-k}=\beta^l\beta^j$. Comme $\alpha$ et $\beta$ sont à support disjoint, ceci implique que $\alpha^i\alpha^{6-k}=Id$ et $\beta^l\beta^j=Id$,
soit en particulier $l=j$. On en déduit que $i=k$.
Donc tous les éléments de la liste $\{\alpha^i\circ\beta^j, 0\leq i \leq 5, 0\leq j\leq 1\}$ sont distincts.
L'ordre de $S$ est donc $6\times 2=12$.
 \item Soit $H$ un sous-groupe de $\mathcal{S}_7$ qui contient $\alpha$ et $\beta$.
Par stabilité par produit, il contient $\alpha^i$, $\beta^j$ et $\alpha^i\circ\beta^j$ pour tout $0\leq i \leq 5, 0\leq j\leq 1$. Il 
 contient donc $S$.
 \item On peut déduire des questions 4. et 6. précédentes que le sous-groupe de $\mathcal{S}_7$ engendré par $\alpha$ et $\beta$ est $S$.
 \item Comme $\alpha$ et $\beta$ sont à support disjoint, $\alpha\circ\beta=(157)(43)(26)$ est une écriture en produit de cycles à support disjoint. L'ordre de $\alpha\circ\beta$ est donc $ppcm(3,2,2)=6$.
 \item Une écriture en produit de cycles à support disjoint de $\alpha^i\circ\beta^j$ montre que son ordre est inférieur à $ppcm(3,2)=6$.
Par conséquent, le sous-groupe $S$ qui n'a pas d'éléments d'ordre 12 n'est pas cyclique.
\end{enumerate}
\fincorrection
\correction{007365}
\begin{enumerate}
 \item On vérifie d'abord que 
$$A^{-1}=\begin{pmatrix}0&1/2\\1&0\end{pmatrix} 
\quad \text{ et } \quad B^{-1}=\begin{pmatrix}1&-4\\0&1\end{pmatrix}$$

Calculer les produits de matrices $ADA^{-1}$ et $BDB^{-1}$ dans $GL(2,\R)$ où
$$ADA^{-1}=\begin{pmatrix}0&1\\2&0\end{pmatrix}
\begin{pmatrix}2&0\\0&3 \end{pmatrix}
\begin{pmatrix}0&1/2\\1&0\end{pmatrix} =\begin{pmatrix}3&0\\0&2 \end{pmatrix}$$
$$BDB^{-1}=\begin{pmatrix}1&4\\0&1\end{pmatrix}
\begin{pmatrix}2&0\\0&3 \end{pmatrix}
\begin{pmatrix}1&-4\\0&1\end{pmatrix}=\begin{pmatrix}2&4\\0&3 \end{pmatrix}$$
 \item Comme $D\in\mathcal{D}$ et $B\in GL(2,\R)$ mais $BDB^{-1}$ n'appartient pas à $\mathcal D$, le sous-groupe $\mathcal{D}$ de $GL(2,\R)$ des matrices diagonales inversibles n'est pas distingué dans $GL(2,\R)$.
 \item $$AB=\begin{pmatrix}0&1\\2&8\end{pmatrix}.$$
 \item $trace(AB)=0+8=8$ et $trace(A)+trace(B)=0+2=2\not=8$.
Par conséquent, $trace$ n'est pas un morphisme de groupes de $GL(2,\R)$ dans $(\R,+)$.
\end{enumerate}
\fincorrection
\nocorrection
\nocorrection
\nocorrection
\correction{001336}
Notons $G$ l'ensemble des \'el\'ements d'ordre fini de $H$.
Montrons que $G$ est un sous-groupe de $H$.
\begin{itemize}
    \item $G \subset H$ et $0 \in G$.
    \item Si $x\in G$ alors $(-x)+(-x)+\cdots+(-x)=-(x+x+\cdots+x) = 0$. Donc
$-x \in G$.
    \item Si $x,y \in G$ alors $(x+y)+\cdots+(x+y)=(x+\cdots+x)+(y+\cdots+y)
=0+0=0$. Donc $x+y\in G$.
\end{itemize}
Nous venons de montrer que $G$ est un sous-groupe de $H$. De plus
comme $H$ est commutatif alors $G$ l'est aussi !
\fincorrection
\correction{001337}
\begin{enumerate}
    \item La matrice  $\begin{pmatrix}
0& 1 \cr 1 & 0 \cr \end{pmatrix}$  est d'ordre  $2$. La matrice
$\begin{pmatrix} 1 & 0 \cr 0 & 2 \cr \end{pmatrix}$  n'est pas
d'ordre fini puisque, pour tout  $n\in {\N}$ : $\begin{pmatrix} 1
& 0 \cr 0 & 2 \cr \end{pmatrix}^n=\begin{pmatrix} 1 & 0 \cr 0 &
2^n \cr \end{pmatrix}\not = \begin{pmatrix} 1 & 0 \cr 0 & 1 \cr
\end{pmatrix}$.
    \item Notons  $e_G$  et  $e_H$  les \' el\' ements neutres respectifs
de  $G$  et de  $H$. Soit  $g$  un \' el\' ement de  $G$  d'ordre
$n$.

-  Alors  $\varphi (g)^n=\varphi (g^n)=\varphi (e_G)=e_H$. Donc
$\varphi (g)$  est d'ordre inf\' erieur ou \' egal \`a  $n$, ordre
de  $g$.

- Supposons  $\varphi $  injectif et  $\varphi (g)$  d'ordre
strictement inf\' erieur \`a  $n$, c'est \`a dire qu'il existe
$p<n$  tel que : $\varphi (g)^p=e_H$. Alors  $\varphi (g^p)=e_H$
donc, puisque  $\varphi$  est injectif et  $\varphi (e_G)=e_H$, on
a aussi : $g^p=e_G$, ce qui est impossible puisque l'ordre de  $g$
est  $n$.
    \item Raisonnons par l'absurde : Soit  $G$  un groupe fini.
Supposons qu'il existe dans  $G$  un \' el\' ement  $g$  n'\'
etant pas d'ordre fini. Comme  $G$  est un groupe, on peut
consid\' erer  $X=\{ g^k \ k\in {\N}\} $. Or, pour  $i\not = j \ :
\ g^i \not = g^j$. En effet, supposons  $i<j$. Si  $g^i  = g^j$
alors  $g^{j-i}=e_G$  et  $g$  est d'ordre inf\' erieur ou \' egal
\`a  $j-i$, donc fini, ce qui est impossible. $X$  est donc un
ensemble infini.
 $G$  contient un ensemble infini donc est infini, ce qui est absurde, donc  $g$  ne
 peut \^etre que d'ordre fini.
\end{enumerate}
\fincorrection
\nocorrection
\nocorrection
\correction{001340}
Rappelons d'abord que pour $x$ un élément d'ordre $n$, alors
$$x^q=e \Longrightarrow n | q.$$
\begin{itemize}
  \item Si $n$ est pair alors $ord(x^2)=n/2$ : en effet
$(x^2)^\frac{n}{2} = x^n=e$ et pour $p \geq 1$ tel que $(x^2)^p =
e$ alors $x^{2p}= e$ et $n|2p$ donc $p \geq \frac n2$. Donc $n/2$
est le plus petit des entiers $q$ (non nul) tel que $x^q=e$ et par
conséquent $n/2$ est l'ordre de $x$.
  \item Si $n$ est impair alors $ord(x)=n$. Tout d'abord
$(x^2)^n=(x^n)^2=e$ et pour $p$ tel que $(x^2)^p=e$ alors $n|2p$
mais $2$ et $n$ sont premiers entre eux donc d'après le théorème
de Gauss, $n|p$ et en particulier $p \geq n$.
\end{itemize}
\fincorrection
\correction{001341}

\begin{enumerate}
  \item Déjà $(xy)^{mn} = x^{mn}y^{mn}= (x^m)^n(y^n)^m=e.e=e$. Soit $p$
tel que $(xy)^p = e$, alors $e= (xy)^{mp} = x^{mp}y^{mp}=y^{mp}$,
et donc $mp$ est divisible par l'ordre de $y$ , c'est-à-dire $n$.
Comme $m$ et $n$ sont premiers entre eux alors d'après le théorème
de Gauss $n$ divise $p$. Un raisonnement semblable à partir de
$(xy)^{np}=e$ conduit à : $m$ divise $p$. Finalement $m|p$ et
$n|p$ donc $mn|p$ car $m$ et $n$ sont premiers entre eux.

Voici un contre exemple dans le cas o\`u $m$ et $n$ ne sont pas
premiers entre eux : dans le groupe $\Zz/12\Zz$ : $\bar{2}$ est
d'ordre $6$, $\bar{4}$ est d'ordre $3$, mais $\bar{2}+\bar{4} =
\bar{6}$ est d'ordre $2 \not= 3\times 6$.
  \item $A$ est d'ordre $4$, $B$ est d'ordre $3$, $(AB)^n =
\begin{pmatrix} 1&n\\0&1\\ \end{pmatrix}$ n'est jamais la matrice identité pout $n\geq 1$.
\end{enumerate}
\fincorrection
\correction{001342}
Par l'absurde supposons que $(\Qq,+)$ est engendré par un seul
élément $\frac pq$ ($p$ et $q$ premiers entre eux) alors tout
élément de $\Qq$ s'écrit $n\frac pq$ avec $n\in \Zz$. Il s'ensuit
que $\frac p{2q}$ (qui appartient à $\Qq$) doit s'écrire $n\frac
pq$, mais alors $2n=1$ avec $n \in \Zz$ ce qui est impossible.
Conclusion $(\Qq,+)$ n'est pas monogène.
\fincorrection
\nocorrection
\nocorrection
\nocorrection
\nocorrection
\nocorrection
\nocorrection
\nocorrection
\correction{002981}
\begin{enumerate}
 \item $x \mapsto ax$, $a\in\Q$.
 \item $x \mapsto 0$.
 \item  $x \mapsto 1$.
\end{enumerate}

\fincorrection
\nocorrection
\nocorrection
\nocorrection
\nocorrection
\nocorrection
\nocorrection
\nocorrection
\nocorrection
\nocorrection
\nocorrection
\nocorrection
\nocorrection
\nocorrection
\nocorrection
\nocorrection
\nocorrection
\nocorrection
\nocorrection
\nocorrection
\nocorrection
\correction{003007}
\begin{enumerate}
  \item
  \item $A$ est int{\`e}gre car $\{0\}$ est premier et si $a\in A\setminus\{0\}$
alors $a\times a\in(a^2)$ qui est premier donc $a^2$ divise $a$ d'o{\`u} $a$ est inversible.
\end{enumerate}
\fincorrection
\nocorrection
\nocorrection
\correction{003010}
\begin{enumerate}
  \item 
  \item $1$.
  \item $x+y = (x+y)^2 = x^2+y^2+xy+yx = x+y+xy+yx  \Rightarrow  xy+yx = 0$.\par
             Pour $y = 1$ : $x+x = 0  \Rightarrow  1 = -1$.\par
             Pour $y$ quelconque : $xy = -yx = yx$.
  \item Antisym{\'e}trie : si $x = ay$, alors $xy = ay^2 = ay = x$.\par
             Donc $(x\le y)$ et $(y\le x)  \Rightarrow  xy = x = y$.
\end{enumerate}
\fincorrection
\nocorrection
\correction{003012}
Si $(1-ab)c = 1 = c(1-ab)$ alors $abc  = c-1 = cab$
donc $babca = bca-ba = bcaba$ soit $ba(1+bca) = bca = (1+bca)ba$ donc
$1+bca$ est inverse de $1-ba$.
\fincorrection
\correction{003013}
\begin{enumerate}
  \item 
  \item 
  \item Remarque : la r{\'e}ciproque fausse : $A = \Z[X]$, $I=(X)$, $J=(X+4)$.
  \item $114\Z$.
\end{enumerate}
\fincorrection
\nocorrection
\nocorrection
\nocorrection
\nocorrection
\nocorrection
\nocorrection
\nocorrection
\nocorrection
\correction{003022}
$f(x_1,\dots,x_n) = a_1x_1 + \dots + a_nx_n$.\par
$f$ est multiplicative sur la base canonique $ \Rightarrow $ $a_ia_j = 0$ pour
$i\ne j$.\par
$f(1,\dots,1) = 1  \Rightarrow  $ un des $a_i$ vaut $1$, et les autres $0$.\par
conclusion : $f$ = fct coordonn{\'e}e.
\fincorrection
\correction{003023}
\begin{enumerate}
  \item 
  \item idem \ref{morphismes} : les projections + la valeur de
             stationnement.
  \item 
\end{enumerate}
\fincorrection
\correction{003024}
\begin{enumerate}
  \item $\pm1,\pm i$.
  \item On a : $1+i = 0\times2 + (1+i) = 1\times2 + (i-1)$.
  \item
\end{enumerate}
\fincorrection
\nocorrection
\nocorrection
\nocorrection
\nocorrection
\nocorrection
\nocorrection
\nocorrection
\nocorrection
\nocorrection
\nocorrection
\nocorrection
\nocorrection
\nocorrection
\nocorrection
\nocorrection
\nocorrection
\nocorrection
\nocorrection
\nocorrection
\nocorrection
\nocorrection
\nocorrection
\nocorrection
\nocorrection
\nocorrection
\nocorrection
\nocorrection
\nocorrection
\nocorrection
\nocorrection
\nocorrection
\nocorrection
\nocorrection
\nocorrection
\nocorrection
\nocorrection
\correction{007369}
\begin{enumerate}
 \item On peut affirmer que l'ordre de $a$ est un diviseur de $m$.
 \item Il s'agit de l'idéal $pgcd(18,28)\Z=252\Z$.
 \item
 \item
 \item $\Z/4\Z$ et $\Z/2\Z\times\Z/2\Z$ sont deux groupes d'ordre $4$, mais le second dont tous les éléments sont d'ordre 1 ou 2 n'est pas cyclique.
 \item Comme $(1,2)(1,3)(1,2)=(2,3)$, $(1,2)$ et $(1,3)$ ne commutent pas. Par conséquent, le groupe $\mathcal{S}_4$ n'est pas cyclique.
 \end{enumerate} 
\fincorrection
\correction{007370}
\begin{enumerate}
 \item $(-7387)-(-1601)=-5786=(-2)2893$ donc, l'entier $-1601$ est un représentant de la classe $[-7387]_{ 2893}$ de $\Z/2893\Z$.
 \item Par le petit théorème de Fermat, $11^{12}=1[13]$.
 On effectue la division euclidienne de $329$ par $12$.
 $329-27\times 12=5$.
 Par conséquent, $11 ^{329}=11^5[13]$.
 $11^2=121=4[13]$, $11^4=4^2=3[13]$, d'où 
 $11 ^{329}=11^5=33=7[13]$
 \item Puisque $\text{pgcd}(51,131)=1$, la classe $[51]$ est inversible dans $\Z/131\Z$. En utilisant l'algorithme d'Euclide on trouve que $1=131\times (-7)+51\times 18$. Alors, dans $\Z/131\Z$ on a $51^{-1}=18$. 
\end{enumerate}
\fincorrection
\correction{007371}
\begin{enumerate}
  \item $ X^3+X+1=(X^2+X+1)(X-1)+(X+2)$ et $X^2+X+1=(X+2)(X-1)+3$.
Une relation de Bezout entre $X^2+X+1$ et $X^3+X+1$ dans $\R[X]$ est 
$$3=(X^2+X+1)-(X+2)(X-1)=(X^2+X+1)-(X-1)[(X^3+X+1)-(X^2+X+1)(X-1)]$$
soit 
$$3=(-X+1)(X^3+X+1)+(X^2-2X+2)(X^2+X+1).$$
 \item Puisque $X^3+X+1$ et $X^2+X+1$ sont premiers entre eux, la classe du polynôme $X^2+X+1$ inversible dans l'anneau quotient $\R[X]/(X^3+X+1)$.
Son inverse est la classe de $(X^2-2X+2)/3$, comme le donne la relation de Bezout.
\end{enumerate}
\fincorrection
\correction{007372}
\begin{enumerate}
 \item 
 \item Comme $\phi(12)=\phi(3)\phi(2^2)=2\times 2=4$ il y a quatre éléments inversibles dans l'anneau $(\Z/12\Z)$.
Il s'agit de $1,-1=11,5,-5=7$.
Comme $5^5=25=1[12]$, tous les éléments sont d'ordre $1$ ou $2$.
Le groupe $(\Z/12\Z)^\times$ est donc isomorphe à $\Z/2\Z\times\Z/2\Z$.
On peut aussi remarquer que par le théorème chinois,
$(\Z/12\Z)^\times = (\Z/4\Z)^\times\times(\Z/3\Z)^\times= \Z/2\Z\times\Z/2\Z$.
\end{enumerate}
\fincorrection
\correction{007373}
 Comme il s'agit d'un groupe d'ordre $5$, il est cyclique donc commutatif.
La table se complète donc en 
$$\begin{array}{|c|c|c|c|c|c|}\hline
 * & a & b 	& c &	d & e\\ \hline
 a & 	 & e	& d &	 & c \\ \hline
 b & e &	&a  &	 & d \\ \hline
 c &	d &a 	 &  &  &b \\ \hline
 d &	 &	 &  &d &e \\ \hline
 e & c & d	 &b &	 e & a \\ \hline
\end{array}
$$
Comme $de=e$, $d$ est l'élément neutre.
La table se complète donc en 
$$\begin{array}{|c|c|c|c|c|c|}\hline
 * & a & b 	& c &	d & e\\ \hline
 a & 	 & e	& d &	 a & c \\ \hline
 b & e &	&a  &	b & d \\ \hline
 c &	d &a 	 &  & c &b \\ \hline
 d &	a &b	 & c &d &e \\ \hline
 e & c & d	 &b &	 e & a \\ \hline
\end{array}
$$
Par la propriété de carré latin, 
on peut terminer le tableau
$$\begin{array}{|c|c|c|c|c|c|}\hline
 * & a & b 	& c &	d & e\\ \hline
 a & 	b & e	& d &	 a & c \\ \hline
 b & e &	c&a  &	b & d \\ \hline
 c &	d &a 	 & e & c &b \\ \hline
 d &	a &b	 & c &d &e \\ \hline
 e & c & d	 &b &	 e & a \\ \hline
\end{array}
$$
\fincorrection
\correction{007374}
\begin{enumerate}
 \item Dans $\mathcal{S}_7$, $(1,2)(1,3)=(132)$ et par conjugaison $(1,2)(2,3)(1,2)=(1,3)$.
 \item L'ensemble des transpositions de $\mathcal{S}_7$ ne contient pas l'identité : ce n'est donc pas un sous-groupe de $\mathcal{S}_7$. Le premier calcul montre qu'il n'est pas non plus stable par produit.
\end{enumerate} 
\fincorrection
\nocorrection
\nocorrection
\nocorrection
\nocorrection
\nocorrection
\nocorrection
\nocorrection
\nocorrection
\nocorrection
\nocorrection
\nocorrection
\nocorrection
\nocorrection
\nocorrection
\nocorrection
\nocorrection
\correction{003026}
$K = \{0,1,a,b\}$ et $\{1,a,b\}$ est un groupe multiplicatif
$ \Rightarrow  b=a^2$, $a^3=1$.\par
%
$\vtop{\offinterlineskip
\halign{\strut\ $#$\hfil\ \vrule\ &&$#$\hfil\ \cr
+   &0   &1   &a   &a^2 \cr
\omit\hfil\vrule height 2pt\ \cr
\noalign{\hrule}
\omit\hfil\vrule height 2pt\ \cr
0   &0   &1   &a   &a^2 \cr
1   &1   &0   &a^2 &a   \cr
a   &a   &a^2 &0   &1   \cr
a^2 &a^2 &a   &1   &0   \cr}}$  \qquad
%
$\vtop{\offinterlineskip
\halign{\strut\ $#$\hfil\ \vrule\ &&$#$\hfil\ \cr
\times&1   &a   &a^2 \cr
\omit\hfil\vrule height 2pt\ \cr
\noalign{\hrule}
\omit\hfil\vrule height 2pt\ \cr
1     &1   &a   &a^2 \cr
a     &a   &a^2 &1   \cr
a^2   &a^2 &1   &a   \cr}}$
\fincorrection
\nocorrection
\nocorrection
\nocorrection
\nocorrection
\nocorrection
\nocorrection
\nocorrection
\nocorrection
\nocorrection
\nocorrection
\nocorrection
\nocorrection
\nocorrection
\nocorrection
\nocorrection
\nocorrection
\nocorrection
\nocorrection
\nocorrection
\nocorrection
\nocorrection
\nocorrection
\nocorrection
\nocorrection
\nocorrection
\nocorrection
\nocorrection
\correction{001412}
\begin{enumerate}
    \item $|S_n| = n!$ donc $|S_3| = 3! =6$. Montrons plus g\'en\'eralement
qu'il n'existe pas d'\'el\'ement d'ordre $n!$ dans $S_n$ ($n\ge
3$). Par l'absurde soit $\alpha$ un tel \'el\'ement. Alors par
hypoth\`ese $S_n$ est engendr\'e par $\alpha$ et donc $S_n$ est un
groupe commutatif. Mais $(1,2)(2,3) \not= (2,3)(1,2)$ ce qui est
absurde. En conclusion il n'existe pas d'\'el\'ements d'ordre $6$.
    \item Explicitons $S_3$ :
$$S_3 = \left\lbrace  id ; \tau_1=(1,2) ; \tau_2=(2,3) ;
\tau_3=(1,3) ; \sigma_1 = (1,2,3) ; \sigma_2 = \sigma_1^{-1}=
(3,2,1) \right\rbrace.$$ Remarquons

Les sous-groupes d'ordre $2$ sont de la forme $\{ id ; \tau \}$
avec $\tau^2 = id$. Les seuls \'el\'ements d'ordre $2$ sont les
transpositions et donc se sont les groupes $\{ id ; (1,2) \}, $\{
id ; (1,3) \}, $\{ id ; (2,3) \}$.

Les sous-groupes d'ordre trois sont de la forme $\{
id,\sigma,\sigma^2 \}$ avec $\sigma^2=\sigma^{-1}$. Et donc le
seul sous-groupe d'ordre $3$ est $\{ id ; (1,2,3) ; (3,2,1) \}$.

    \item Les sous-groupes de $S_3$ ont un ordre qui divise
$|S_3| = 6$. Donc un sous-groupe peut-\^etre  d'ordre $1,2,3$ ou
$6$. L'unique sous-groupe d'ordre $1$ est $\{ id \}$, et l'unique
sous-groupe d'ordre $6$ est $S_3$. Les sous-groupes d'ordre $2$ et
$3$ ont \'et\'es donn\'es \`a la question pr\'ec\'edente.
\end{enumerate}
\fincorrection
\nocorrection
\nocorrection
\nocorrection
\nocorrection
\nocorrection
\correction{001418}
\begin{enumerate}
  \item $\sigma = (1,3)(2,7,9,5)= (2,7,9,5)(1,3)$ et
$\sigma^k = (1,3)^k(2,7,9,5)^k$. Les transpositions sont d'ordre
$2$ donc $(1,3)^k = \mathrm{id}$ si $k \equiv 0 (\mod 2)$ et $(1,3)^k =
(1,3)$ si $k \equiv 1 (\mod 2)$. Le cycle $(2,7,9,5)$ est d'ordre
$4$, et $(2,7,9,5)^k$ est respectivement égale à $\mathrm{id}, (2,7,9,5),
(2,9)(7,5), (5,9,7,2)$ si $k$ est respectivement congru à
$0,1,2,3$ modulo $4$. Le calcul de $\sigma^k$ donne donc $\mathrm{id},
(1,3)(2,7,9,5), (2,9)(7,5)$ ou $(1,3)(5,9,7,2)$ selon que $k$ est
congru à $0,1,2$ ou $3$ modulo $4$.
  \item L'écriture de
$\varphi = (10,3,4,1)(8,7)(4,7)(5,6)(2,6)(2,9)$ est une
décomposition en produit de cycles mais ils ne sont pas à supports
disjoints. \'Ecrivons $\varphi$ sous la forme :
$$\phi = \begin{pmatrix}
1&2&3&4&5&6&7&8&9&10\\
10&9&4&8&6&2&1&7&5&3\\
\end{pmatrix}$$
Ce qui se décompose $\varphi = (1,10,3,4,8,7)(2,9,5,6) =
(2,9,5,6)(1,10,3,4,8,7)$. Le calcul de $\varphi^k =
(1,10,3,4,8,7)^k(2,9,5,6)^k$ est similaire au calcul précédent
(selon $k (\mod 12)$ )

\end{enumerate}
\fincorrection
\nocorrection
\nocorrection
\nocorrection
\nocorrection
\nocorrection
\nocorrection
\nocorrection
\correction{001426}
\begin{enumerate}
  \item $\mathcal{S}_{N}$ est l'ensemble des permutations de l'ensemble
$\{ 1,2,\ldots,N \}$. Dans $\mathcal{S}_{n+2}$ notons $\tau$ la
permutation $(n+1,n+2)$. Nous définissons une application $\phi :
\mathcal{S}_{n} \longrightarrow \mathcal{S}_{n+2}$ par les
relations
$$ \phi(\sigma) = \sigma \text{ si } \epsilon(\sigma) = +1 \quad ;
\quad \phi(\sigma) =  \sigma \circ \tau \text{ sinon ;}
$$
o\`u $\epsilon$ désigne la signature. Alors $\phi$ est un
morphisme de groupe, de plus quelque soit $\sigma \in
\mathcal{S}_{n}$ alors $\epsilon (\phi(\sigma))=+1$ (si
$\epsilon(\sigma)= +1$ c'est clair, sinon $\epsilon
(\phi(\sigma))=\epsilon(\sigma)\times \epsilon(\tau) = (-1)\times
(-1) = +1$). Donc $\phi (\mathcal{S}_{n})$ est un sous-groupe de
$\mathcal{A}_{n+2}$.

Enfin $\phi$ est injective : en effet soit $\sigma$ tel que
$\phi(\sigma)=\mathrm{id}$. Soit $\epsilon(\sigma) = +1 $ et alors
$\phi(\sigma)= \sigma = \mathrm{id}$ ; soit $\epsilon(\sigma) = -1$ et
alors $\phi(\sigma)= \sigma \circ \tau$, pour $j \in
\{1,2,\ldots,n\}$
 $j = \phi(\sigma)(j)= \sigma \circ \tau(j) = \sigma(j)$, et donc
quelque soit $j\in \{1,2,\ldots,n\}$ $\sigma(j)=j$ et donc $\sigma
= \mathrm{id}$. On vient de démontrer que la composée de deux permutations
à supports disjoint est l'identité si et seulement si les
permutations sont déjà l'identité !


Notons encore $\phi : \mathcal{S}_{n} \longrightarrow
\phi(\mathcal{S}_{n} )$ le morphisme induit par $\phi$. Il est
injectif et surjectif, donc $\mathcal{S}_{n}$ est isomorphe
$\phi(\mathcal{S}_{n})$ qui est un sous-groupe de
$\mathcal{A}_{n+2}$.

  \item $\mathcal{A}_5$ est  de cardinal $5 !/2=60$, et comme $24 = \mathrm{Card} \mathcal{S}_4$ ne divise
pas $60$    alors $\mathcal{A}_5$  n'a pas de sous-groupe d'ordre
$24$.
  \item C'est un peu plus délicat car $\mathrm{Card} \mathcal{S}_5 = 5 ! = 120$ divise
$\mathrm{Card} \mathcal{A}_6 = 6 !/2 = 360$ donc l'argument ci-dessus
n'est pas valide. Cependant s'il existe un isomorphisme entre
$\mathcal{S}_5$ et un sous-groupe de $\mathcal{A}_6$ alors un
cycle d'ordre $5$ de $\mathcal{S}_5$ est envoyé sur une
permutation $\sigma \in \mathcal{A}_6$ d'ordre $5$.

Décomposons $\sigma$ en produit de cycles à supports disjoints,
$\sigma = \sigma_1 \circ \sigma_2 \circ \cdots$. Comme les cycles
$\sigma_i$ sont à supports disjoints, il y a au plus trois cycles
(de longueur $\geq 2$) dans la décomposition (car dans
$\mathcal{A}_6$ on peut permuter au plus $6$ éléments).
\begin{itemize}
  \item Le cas $\sigma = \sigma_1$ n'est pas possible car alors $\sigma_1$ serait un cycle d'ordre $5$ et donc
de signature $-1$ dans $\mathcal{A}_6$.
  \item Si $\sigma = \sigma_1 \circ \sigma_2$ alors les longueurs de $\sigma_1$
  et $\sigma_2$ sont $(4,2)$ ou $(2,2)$, et l'ordre de leur composée
  $\sigma_1 \circ \sigma_2$ est donc $4$ ou $2$ mais pas $5$.
  \item Si $\sigma = \sigma_1 \circ \sigma_2 \circ \sigma_3$ alors
les $\sigma_i$ sont des transpositions, et la signature de
$\sigma$ est alors $-1$ ce qui contredit $\sigma \in
\mathcal{A}_6$.
\end{itemize}
\end{enumerate}
\fincorrection
\nocorrection
\nocorrection
\nocorrection
\correction{003076}
\begin{enumerate}
 \item $(a\ b)\circ(c\ d)$.
 \item  
\end{enumerate}

\fincorrection
\correction{003077}
\begin{enumerate}
 \item $(1\ 2)\circ(i\ j)$.
 \item 
\end{enumerate}
\fincorrection
\correction{003078}
$\epsilon(\sigma) = (-1)^{n+c+f}$.
\fincorrection
\nocorrection
\nocorrection
\nocorrection
\correction{003082}
Compter les inversions ou r{\'e}currence :
         $\epsilon(\sigma) = (-1)^{n(n-1)/2}$.
\fincorrection
\nocorrection
\correction{003084}
Les puissances de $\sigma$.
\fincorrection
\correction{003085}
Conjugaison : $\tau = (1\ 2\ 3\ 4\ 5)^x \circ (6\ 7\ 8\ 9\ 10)^y$,\par
         ou $\tau = (1\ 6)\circ(2\ 7)\circ(3\ 8)\circ(4\ 9)\circ(5\ 10)\circ
                    (1\ 2\ 3\ 4\ 5)^x \circ (6\ 7\ 8\ 9\ 10)^y$.\par
         $ \Rightarrow $ 50 {\'e}l{\'e}ments.
\fincorrection
\nocorrection
\correction{003087}
30.
\fincorrection
\nocorrection
\correction{003089}
$$C_{26}^3
  \times \frac{2C_{23}^3\times2C_{20}^3}{2!}
  \times 4!\,C_{17}^5
  \times \frac{5!\,C_{12}^6\times5!\,C_{6}^6}{2!}
  = 10\,372\,722\,765\,601\,996\,800\,000.$$
\fincorrection
\correction{005353}
\begin{enumerate}
\item  Les inversions de $\sigma$ sont :
$\sigma=(3\;10\;7\;1\;2\;6\;4\;5\;12\;8\;9\;11)$. 

$$\begin{array}{l}
\{1,4\},\;\{1,5\},\;\{2,3\},\;\{2,4\},\;\{2,5\},\;\{2,6\},\;\{2,7\},\;\{2,8\},\;\{2,10\},\;\{2,11\},\;\{3,4\},\;\{3,5\},\;\{3,6\},\;\{3,7\},\\
\{3,8\},\;\{6,7\},\;\{6,8\},\;\{9,10\},\;\{9,11\},\;\{9,12\}.
\end{array}$$

Au total, il y a $2+8+5+2+3=20$ inversions. $\sigma$ est donc une permutation paire (de signature $1$).

\item  $\tau_{11,12}\circ\sigma=(3\;10\;7\;1\;2\;6\;4\;5\;11\;8\;9\;12)$. 

Puis, $\tau_{9,11}\circ\tau_{11,12}\circ\sigma=(3\;10\;7\;1\;2\;6\;4\;5\;9\;8\;11\;12)$. 

Puis, $\tau_{10,8}\circ\tau_{9,11}\circ\tau_{11,12}\circ\sigma=(3\;8\;7\;1\;2\;6\;4\;5\;9\;10\;11\;12)$. 

Puis, 
$\tau_{8,5}\circ\tau_{10,8}\circ\tau_{9,11}\circ\tau_{11,12}\circ\sigma=(3\;5\;7\;1\;2\;6\;4\;8\;9\;10\;11\;12)$. 

Puis, $\tau_{7,4}\circ\tau_{8,5}\circ\tau_{10,8}\circ\tau_{9,11}\circ\tau_{11,12}\circ\sigma=(3\;5\;4\;1\;2\;6\;7\;8\;9\;10\;11\;12)$. 

Puis, $\tau_{5,2}\circ\tau_{7,4}\circ\tau_{8,5}\circ\tau_{10,8}\circ\tau_{9,11}\circ\tau_{11,12}\circ\sigma=(3\;2\;4\;1\;5\;6\;7\;8\;9\;10\;11\;12)$.

Puis, $\tau_{1,4}\circ\tau_{5,2}\circ\tau_{7,4}\circ\tau_{8,5}\circ\tau_{10,8}\circ\tau_{9,11}\circ\tau_{11,12}\circ\sigma=(3\;2\;1\;4\;5\;6\;7\;8\;9\;10\;11\;12)=\tau_{1,3}$.

Par suite,

$$\sigma=\tau_{11,12}\circ\tau_{9,11}\circ\tau_{10,8}\circ\tau_{8,5}\circ\tau_{7,4}\circ\tau_{5,2}\circ\tau_{1,4}\circ\tau_{1,3}.$$
\item  $O(1)=\{1,3,4,7\}=O(3)=O(4)=O(7)$, puis $O(2)=\{2,5,8,10\}$ puis $O(6)=\{6\}$ et $O(9)=\{9,11,12\}=O(11)=O(12)$. $\sigma$ a $4$ orbites, deux de cardinal $4$, une de cardinal $3$ et un singleton (correspondant à un point fixe).

\item  $\sigma$ est donc le produit commutatif des cycles $c_1=\left(\begin{array}{cccc}
1&3&4&7\\
3&7&7&4
\end{array}
\right)$, $c_2=\left(\begin{array}{cccc}
2&5&8&10\\
10&2&5&8
\end{array}
\right)$ et

$c_3=\left(\begin{array}{ccc}
9&11&12\\
12&9&11
\end{array}
\right)$.

On a $c_1^4=c_2^4=Id$ et $c_3^3=Id$. Or, $2005=4.1001+1$. Donc, $c_1^{2005}=c_1(c_1^4)^{1001}=c_1$, et de même $c_2^{2005}=c_2$. Puis, $c_3^{2005}=(c_3^{3})^{668}c_3=c_3$. Puisque $c_1$, $c_2$ et $c_3$ commutent, 

$$\sigma^{2005}=c_1^{2005}c_2^{2005}c_3^{2005}=c_1c_2c_3=\sigma=(3\;10\;7\;1\;2\;6\;4\;5\;12\;8\;9\;11).$$
\end{enumerate}

\fincorrection
\correction{005354}
$(S_n,\circ)$ est engendré par les transpositions. Il suffit donc de montrer que pour $2\leq i<j\leq n$, la transposition $\tau_{i,j}$ est produit des $\tau_{1,k}$, $2\leq k\leq n$.

Mais $\tau_{1,i}\circ\tau_{1,j}\circ\tau_{1,i}=(i 1 j)(j i 1)(i 1 j)=(1 i j)=\tau_{i,j}$ ce qu'il fallait démontrer.
\fincorrection
\correction{005355}
Les éléments de $A_n$ sont les produits pairs de transpositions. Il suffit donc de vérifier qu'un produit de deux transpositions est un produit de cycles de longueur $3$.

Soient $i$, $j$ et $k$ trois éléments deux à deux distincts de $\{1,...,n\}$. $\tau_{i,k}\circ\tau_{i,j}$ est le $3$-cycle~:~$i\rightarrow j$ $j\rightarrow k$ $k\rightarrow i$, ce qui montre qu'un $3$-cycle est pair et que le produit de deux transpositions dont les supports ont en commun un singleton est un $3$-cycle.

Le cas $\tau_{i,j}\circ\tau_{i,j}=Id=(2 3 1)(3 1 2)$ est immédiat. Il reste à étudier le produit de deux transpositions à supports disjoints.

Soient $i$, $j$, $k$ et $l$ quatre éléments de deux à deux distincts de $\{1,...,n\}$.

$$\tau_{i,j}\circ\tau_{k,l}=(j i k l)(i j l k)=(j i l k)=(j k i l)(l j i k).$$ 

Donc, $\tau_{i,j}\circ\tau_{k,l}$ est un bien un produit de $3$-cycles ce qui achève la démonstration.
\fincorrection
\correction{005356}
D'après l'exercice \ref{exo:suprou2ter}, il suffit de montrer que pour $2\leq i\leq n$, $\tau_{1,i}$ peut s'écrire en utilisant uniquement $\tau=\tau_{1,2}$ et $c=(2\;3\;...\;n\;1)$. On note que $c^n=Id$.

Tout d'abord, pour $1\leq i\leq n-1$, étudions $\sigma=c^{i-1}\circ\tau\circ c^{n-i+1}$.

Soit $k\in\{1,...,n\}$. 

\begin{align*}\ensuremath
\tau\circ c^{n-i+1}(k)\neq c^{n-i+1}(k)&\leq c^{n-i+1}(k)\in\{1,2\}\Leftrightarrow k\in\{c^{-n+i-1}(1),c^{-n+i-1}(2)\}\Leftrightarrow k\in\{c^{i-1}(1),c^{i-1}(2)\}\\
 &\Leftrightarrow k\in\{i,i+1\}.
\end{align*}

Donc, si $k\notin\{i,i+1\}$, 

$$\sigma(k)=c^{i-1}(k)(\tau\circ c^{n-i+1}(k))=c^{i-1}(c^{n-i+1}(k))=c^n(k)=k,$$

et la restriction de $\sigma$ à $\{1,...,n\}\setminus\{i,i+1\}$ est l'identité de cet ensemble. Comme $\sigma$ n'est pas l'identité puisque $\sigma(i)\neq i$, $\sigma$ est donc nécessairement la transposition $\tau_{i,i+1}$.

On a montré que $\forall i\in\{1,...,n-1\},\;c^{i-1}\circ\tau \circ c^{n-i+1}=\tau_{i,i+1}$.

Vérifions maintenant que les $\tau_{1,i}$ s'écrivent à l'aide des $\tau_{j,j+1}$. D'après l'exercice \ref{exo:suprou2ter}, $\tau_{i,j}=\tau_{1,i}\circ\tau_{1,j}\circ\tau_{1,i}$, et donc bien sûr, plus généralement, $\tau_{i,j}=\tau_{k,i}\circ\tau_{k,j}\circ\tau_{k,i}$.

Par suite, $\tau_{1,i}=\tau_{1,2}\circ\tau_{2,i}\circ\tau_{1,2}$ puis, $\tau_{2,i}=\tau_{2,3}\circ\tau_{3,i}\circ\tau_{2,3}$, puis, $\tau_{3,i}=\tau_{3,4}\circ\tau_{4,i}\circ\tau_{3,4}$ ... et $\tau_{i-2,i}=\tau_{i-2,i-1}\circ\tau_{i-1,i}\circ\tau_{i-2,i-1}$. Finalement,

$$\tau_{1,i}=\tau_{1,2}\circ\tau_{2,3}\circ...\circ\tau_{i-2,i-1}\tau_{i-1,i}\circ\tau_{i-2,i-1}\circ...\circ\tau_{2,3}\circ\tau_{1,2},$$

ce qui achève la démonstration.
\fincorrection
\correction{005357}
Soit $(G,\times)$ un groupe. Pour $x$ élément de $G$, on considère 
$\begin{array}[t]{cccc}
f_x~:&G&\rightarrow&G\\
 &y&\mapsto&xy
\end{array}$. $f_x$ est une application de $G$ vers $G$ et de plus, clairement $f_x\circ f_{x^{-1}}=f_{x^{-1}}\circ f_x=Id_G$. Donc, pour tout élément $x$ de $G$, $f_x$ est une permutation de $G$.

Soit alors $\begin{array}[t]{cccc}
\varphi~:&(G,\times)&\rightarrow&(S_G,\circ)\\
 &x&\mapsto&f_x
\end{array}$. D'après ce qui précède, $\varphi$ est une application. De plus, $\varphi$ est de plus un morphisme de groupes. En effet, pour $(x,x',y)\in G^3$, on a~:

$$\varphi((xx'))(y)=f_{xx'}(y)=xx'y=f_x(f_x'(y)) = f_x\circ f_{x'}(y)=(\varphi(x)\circ\varphi(x'))(y),$$

et donc $\forall(x,x')\in G^2,\;\varphi(xx')=\varphi(x)o\varphi(x')$.

Enfin, $\varphi$ est injectif car, pour $x$ élément de $G$~:

$$\varphi(x)= Id\Rightarrow\forall y\in G,\;xy=y\Rightarrow xe=e\Rightarrow x=e.$$

Donc, $\mbox{Ker}\varphi=\{e\}$, et $\varphi$ est injectif.

$\varphi$ est ainsi un isomorphisme de groupes de $(G,\times)$ sur $(f(G),\circ)$ qui est un sous groupe de $(S_G,\circ)$. $(G,\times)$ est bien isomorphe à un sous groupe de $(S_G,\circ)$.
\fincorrection
\correction{005358}
Montrons d'abord par récurrence sur $l\geq 2$ que la signature d'un cycle de longueur $l$ est $(-1)^{l-1}$.

C'est connu pour $l=2$ (signature d'une transposition).

Soit $l\geq 2$. Supposons que tout cycle de longueur $l$ ait pour signature $(-1)^{l-1}$. Soit $c$ un cycle de longueur $l+1$.

On note $\{x_1,x_2,...,x_{l+1}\}$ le support de $c$ et on suppose que, pour $1\leq i\leq l$, $c(x_i)=x_{i+1}$ et que $c(x_{l+1})=x_1$.

Montrons alors que $\tau_{x_1,x_{l+1}}\circ c$ est un cycle de longueur $l$. $\tau_{x_1,x_{l+1}}\circ c$ fixe déjà $x_{l+1}$ puis, si $1\leq i\leq l-1$, $\tau_{x_1,x_{l+1}}\circ c(x_i)=\tau_{x_1,x_{l+1}}(x_{i+1})=x_{i+1}$ (car $x_{i+1}$ n'est ni $x_1$, ni $x_{l+1}$), et enfin $\tau_{x_1,x_{l+1}}\circ c(x_l)=\tau_{x_1,x_{l+1}}(x_{l+1})=x_1$. 
$\tau_{x_1,x_{l+1}}\circ c$ est donc bien un cycle de longueur $l$. Par hypothèse de récurrence, $\tau_{x_1,x_{l+1}}\circ c$ a pour signature $(-1)^{l-1}$ et donc, $c$ a pour signature $(-1)^{(l+1)-1}$.

Montrons maintenant que si $\sigma$ est une permutation quelconque de $\{1,...,n\}$ ayant $k$ orbites la signature de $\sigma$ est $(-1)^{n-k}$.

Si $\sigma$ est l'identité, $\sigma$ a $n$ orbites et le résultat est clair.

Si $\sigma$ n'est pas l'identité, on décompose $\sigma$ en produit de cycles à supports disjoints.

Posons $\sigma=c_1...c_p$ où $p$ désigne le nombre d'orbites de $\sigma$ non réduites à un singleton et donc $k-p$ est le nombre de points fixes de $\sigma$. Si $l_i$ est la longueur de $c_i$, on a donc $n=l_1+...+l_p+(k-p)$ ou encore $n-k=l_1+...+l_p-p$.

Mais alors,

$$\varepsilon(\sigma)=\prod_{i=1}^{p}\varepsilon(c_i)=\prod_{i=1}^{p}(-1)^{l_i-1}=(-1)^{l_1+...+l_p-p}=(-1)^{n-k}.$$
\fincorrection
\correction{005359}
\begin{enumerate}
\item  
\begin{enumerate}
\item Soient $\sigma$ et $\sigma'$ deux éléments de $S_n$. Soit $(i,j)\in\{1,...,n\}^2$. Le coefficient ligne $i$, colonne $j$ de $P_\sigma P_{\sigma'}$ vaut

$$\sum_{k=1}^{n}\delta_{i,\sigma(k)}\delta_{k,\sigma'(j)}=\delta_{i,\sigma(\sigma'(j))},$$

et est donc aussi le coefficient ligne $i$, colonne $j$ de la matrice $P_{\sigma\circ\sigma'}$. Par suite,

$$\forall(\sigma,\sigma')\in(S_n)^2,\;P_\sigma\times P_{\sigma'}=P_{\sigma\circ\sigma'}.$$

\item Soit $\sigma\in S_n$. D'après a), $P_\sigma P_{\sigma^{-1}}=P_{\sigma\circ\sigma^{-1}}=P_{Id}=I_n=P_{\sigma^{-1}}P_\sigma$. On en déduit que toute matrice $P_\sigma$ est inversible, d'inverse $P_{\sigma^{-1}}$. Par suite, $G\subset GL_n(\Rr)$ (et clairement, $G\neq\emptyset$).

Soit alors $(\sigma,\sigma')\in(S_n)^2$.

$$P_\sigma P_{\sigma'}^{-1}=P_\sigma P_{{\sigma'}^{-1}}=P_{\sigma\circ{\sigma'}^{-1}}\in G.$$

On a montré que $G$ est un sous-groupe de $(GL_n(\Rr),\times)$.

Soit $\begin{array}[t]{cccc}
\varphi~:&S_n&\rightarrow&G\\
 &\sigma&\mapsto&P_\sigma
\end{array}$. D'après a), $\varphi$ est un morphisme de groupes. $\varphi$ est clairement surjectif. Il reste à vérifier que $\varphi$ est injectif.

Soit $\sigma\in S_n$.

\begin{align*}\ensuremath
\sigma\in\mbox{Ker}\varphi\Rightarrow P_\sigma=I_n&\Rightarrow\forall(i,j)\in\{1,...,n\}^2,\;\delta_{i,\sigma(j)}=\delta_{i,j}\\
 &\Rightarrow \forall i\in\{1,...,n\},\;\delta_{i,\sigma(i)}=1\Rightarrow\forall i\in\{1,...,n\},\;\sigma(i)=i\\
 &\Rightarrow\sigma=Id.
\end{align*}

Puisque le noyau du morphisme $\varphi$ est réduit à $\{Id\}$, $\varphi$ est injectif.

Ainsi, $\varphi$ est un isomorphisme du groupe $(S_n,\circ)$ sur le groupe $(G,\times)$ et on a montré que $(G,\times)$ est un sous-groupe de $(GL_n(\Rr),\times)$, isomorphe à $(S_n,\circ)$.
\end{enumerate}

\item  Soit $(i,j)\in\{1,...,n\}^2$. Le coefficient ligne $i$, colonne $j$ de $AP_\sigma$ vaut~:

$$\sum_{k=1}^{n}a_{i,k}\delta_{k,\sigma(j)}=a_{i,\sigma(j)}.$$

Ainsi, l'élément ligne $i$, colonne $j$, de $AP_{\sigma}$ est l'élément ligne $i$, colonne $\sigma(j)$, de $A$, ou encore, si $j$ est un élément donné de $\{1,...,n\}$, la $j$-ème colonne de $AP_\sigma$ est la $\sigma(j)$-ème colonne de $A$. Ainsi, si on note $C_1$,...,$C_n$ les colonnes de $A$ (et donc $A=(C_1,...,C_n)$), alors $AP_\sigma=(C_{\sigma(1)},...,C_{\sigma(n)})$. En clair, multiplier $A$ par $P_\sigma$ à droite a pour effet d'appliquer la permutation $\sigma$ aux colonnes de $A$ (puisque $P_\sigma$ est inversible, on retrouve le fait que permuter les colonnes de $A$ ne modifie pas le rang de $A$).

De même, le coefficient ligne $i$, colonne $j$, de $P_\sigma A$ vaut 

$$\sum_{k=1}^{n}\delta_{i,\sigma(k)}a_{k,j}=\sum_{k=1}^{n}\delta_{\sigma^{-1}(i),k}a_{k,j}=a_{\sigma^{-1}(i),j},$$

(on a utilisé $\sigma(k)=i\Leftrightarrow k=\sigma^{-1}(i)$) et multiplier $A$ par $P_\sigma$ à gauche a pour effet d'appliquer la permutation $\sigma^{-1}$ aux lignes de $A$.
\end{enumerate}
\fincorrection
\correction{005360}
$G=\{A_1,...,A_p\}$ est déjà une partie non vide de $GL_n(\Rr)$, stable pour $\times$. Il reste à vérifier que $G$ est stable pour le passage à l'inverse.

Soient $i\in\{1,...,n\}$, puis $\begin{array}[t]{cccc}
\varphi_i~:&G&\rightarrow&G\\
 &A&\mapsto&A_iA
\end{array}$. Puisque $G$ est stable pour le produit, $\varphi_i$ est une application de $G$ dans $G$.
 
Montrons que $\varphi_i$ est injective. Soit $(A,B)\in G$.

$$\varphi_(A)=\varphi_i(B)\Rightarrow A_iA=A_iB\Rightarrow A_i^{-1}A_iA=A_i^{-1}A_iB\Rightarrow A=B.$$

Donc, $\varphi_i$ est une application injective de l'ensemble \textbf{fini} $G$ dans lui-même. On sait alors que $\varphi_i$ est une permutation de $G$.

Par $\varphi_i$, $A_i$ a un antécédent $A$ dans $G$. $A_iA=A_i$ fournit $A_i^{-1}A_iA=A_i^{-1}A_i$ puis $A=I\in G$. Ainsi, $G$ contient la matrice $I$. Ensuite, $I$ a un antécédent par $\varphi_i$ dans $G$. Donc, il existe $B\in G$ telle que $A_iB=I$. Mais alors $A_i^{-1}=B\in G$.

$G$ est bien stable pour le passage à l'inverse et est donc un sous-groupe de $(GL_n(\Rr),\times)$.
\fincorrection
\correction{005361}
Pour $(x_1,...,x_n)\in E$, on pose $\varphi((x_1,...,x_n))=x_1+...+x_n$. $\varphi$ est une forme linéaire non nulle sur $E$ et $H$ est le noyau de $\varphi$. $H$ est donc bien un hyperplan de $E$.

Il est clair que, pour $(\sigma,\sigma')\in S_n^2$, $f_\sigma\circ f_{\sigma'}=f_{\sigma\circ\sigma'}$.
$(\mathcal{L}(E),+,.)$ est un espace vectoriel et donc, $p$ est bien un endomorphisme de $E$.

$$p^2=\frac{1}{n!^2}\left(\sum_{\sigma\in S_n}^{}f_\sigma\right)^2=\sum_{(\sigma,\sigma')\in(S_n)^2}^{}f_\sigma\circ f_{\sigma'}.$$

Mais, $(S_n,\circ)$ est un groupe fini. Par suite, l'application $\begin{array}[t]{ccc}
S_n&\rightarrow&S_n\\
\sigma&\mapsto&\sigma\circ\sigma'
\end{array}$, injective (même démarche que dans l'exercice \ref{exo:suprou8ter}), est une permutation de $S_n$. On en déduit que, pour $\sigma'$ donnée, $\sum_{\sigma\in S_n}^{}f_{\sigma\circ\sigma'}=\sum_{\sigma\in S_n}^{}f_{\sigma}$. Ainsi, en posant $q=n!p$.

$$p^2=\frac{1}{n!^2}\sum_{\sigma'\in S_n}^{}(\sum_{\sigma\in S_n}^{}f_{\sigma\circ\sigma'})=\frac{1}{n!^2}\sum_{\sigma'\in S_n}^{}q=\frac{1}{n!^2}.n!q=\frac{1}{n!}q=p.$$

$p$ est donc une projection. Déterminons alors l'image et le noyau de $p$. Soit $i\in\{1,...,n\}$. 

$$p(e_i)=\frac{1}{n!}\sum_{\sigma\in S_n}^{}f_\sigma(e_i)=\frac{1}{n!}\sum_{\sigma\in S_n}^{}e_{\sigma(i)}.$$

Maintenant, il y a (bien sûr) autant de permuations $\sigma$ telles que $\sigma(i)=1$, que de permutations $\sigma$ telles que $\sigma(i)=2$,... ou de permutations $\sigma$ telles que $\sigma(i)=n$, à savoir $\frac{n!}{n}=(n-1)!$. Donc,

$$\forall i\in\{1,...,n\},\;p(e_i)=\frac{1}{n!}\frac{n!}{n}\sum_{k=1}^{n}e_k=\frac{1}{n}\sum_{k=1}^{n}e_k.$$

Posons $u=\frac{1}{n}\sum_{k=1}^{n}e_k$. D'après ce qui précède,

$$\mbox{Im}p=\mbox{Vect}(p(e_1),...,p(e_n))=\mbox{Vect}(u).$$

Ensuite, si $x=x_1e_1+...+x_ne_n$ est un élément de $E$,

$$p(x)=0\Leftrightarrow\sum_{k=1}^{n}x_kp(e_k)=0\Leftrightarrow(\sum_{k=1}^{n}x_k)u=0\Leftrightarrow\sum_{k=1}^{n}x_k=0\Leftrightarrow x\in H.$$

Ainsi, $p$ est la projection sur $\mbox{Vect}(u)$ parallèlement à $H$.
\fincorrection
\correction{002960}
\begin{enumerate}
 \item Commutative, associative, 0 = {\'e}lt neutre, tout {\'e}lt $\ne$ 1 est r{\'e}gulier,
    seul 0 est sym{\'e}trisable.
 \item Tout {\'e}lt est sym{\'e}trisable et $x^{-1} = \frac x{x-1}$.
\end{enumerate}
\fincorrection
\correction{002961}
$\exists\ b \in E$ tq $a*b*a = a$. Alors $b*a$ est neutre {\`a} droite et $a*b$ est
neutre {\`a} gauche.
\fincorrection
\correction{002962}
\begin{enumerate}
  \item Associative, commutative, $\{e\}$ = {\'e}lement neutre,
               $A$ est sym{\'e}trisable $\iff A = \{a\}$ avec $a$ sym{\'e}trisable.
  \item Oui.
\end{enumerate}
\fincorrection
\correction{002963}
\begin{enumerate}
  \item Non commutative, associative, $(1,0)$ = {\'e}lt neutre,\par
               $(a,b)$ est r{\'e}gulier $\iff a \ne 0$. \par
               $(a,b)$ est inversible $\iff a = \pm1$.
  \item
  \item 
  \item 
\end{enumerate}
\fincorrection
\nocorrection
\nocorrection
\nocorrection
\nocorrection
\nocorrection
\nocorrection
\nocorrection
\nocorrection
\nocorrection
\nocorrection
\nocorrection
\correction{002996}
\begin{enumerate}
  \item 
  \item Soit $x \in G$ : $\exists\ u,v \in \Z$ tq $ua + vb = 1
              \Rightarrow  x = (x^{ua})(x^{vb})$.
\end{enumerate}
\fincorrection
\nocorrection
\nocorrection
\nocorrection
\nocorrection
\nocorrection
\nocorrection
\correction{003003}
Soir $\{e_1,\dots,e_p\}$ une partie g{\'e}n{\'e}ratrice de cardinal
minimal. Alors les $2^p$ {\'e}l{\'e}ments $e_1^{\alpha_1}\dots e_p^{\alpha_p}$
avec $\alpha_i\in\{0,1\}$ sont distincts (sinon un des $e_i$ appartient
au groupe engendr{\'e} par les autres) donc $n\ge 2^p$.
\fincorrection
\correction{003004}
Si $a\in G$ est d'ordre infini alors il engendre un sous-groupe
isomorphe {\`a}~$\Z$, qui a une infinit{\'e} de sous-groupes~; c'est exclu.
Donc tous les sous-groupes monog{\`e}nes de~$G$ sont finis, et $G$ est
la r{\'e}union de ces sous-groupes.
\fincorrection
\nocorrection
\nocorrection
\nocorrection
\nocorrection
\nocorrection
\correction{007375}
\begin{enumerate}
 \item C'est la multiplication. Puisque $\Z/17\Z$ est un corps fini, 
par un théorème du cours, le groupe de ses inversibles est un groupe cyclique d'ordre $16$.
 \item Il en a $\phi(91)=\phi(13\times17)=\phi(13)\times\phi(17)=12\times 16= 192$.
 \item $s=(15)(273)(469)$.
 \item \begin{eqnarray*}
(19)\sigma&=&\left(\begin{array}{ccccccccc} 
1&2&3&4&5&6&7&8&9\\
1&8&6&7&2&3&9&4&5\end{array}\right)\\
(28)(19)\sigma&=&\left(\begin{array}{ccccccccc} 
1&2&3&4&5&6&7&8&9\\
1&2&6&7&8&3&9&4&5\end{array}\right)\\
(36)(28)(19)\sigma&=&\left(\begin{array}{ccccccccc} 
1&2&3&4&5&6&7&8&9\\
1&2&3&7&8&6&9&4&5\end{array}\right)\\
(47)(36)(28)(19)\sigma&=&\left(\begin{array}{ccccccccc} 
1&2&3&4&5&6&7&8&9\\
1&2&3&4&8&6&9&7&5\end{array}\right)\\
(58)(47)(36)(28)(19)\sigma&=&\left(\begin{array}{ccccccccc} 
1&2&3&4&5&6&7&8&9\\
1&2&3&4&5&6&9&7&8\end{array}\right)\\
(79)(58)(47)(36)(28)(19)\sigma&=&\left(\begin{array}{ccccccccc} 
1&2&3&4&5&6&7&8&9\\
1&2&3&4&5&6&7&9&8\end{array}\right)=(89)\\
\end{eqnarray*}
Donc,
$$\sigma=(19)(28)(36)(47)(58)(79)(89).$$
\item Non, il change dans la transposition de $t_1$ par $t_2$.

\end{enumerate}
\fincorrection
\correction{007376}
\begin{enumerate}
 \item D'après le cours, il en a $\phi(24)=\phi(2^3\times3)
=(2^3-2^2)(2)=8$.
 \item Par le théorème chinois, comme $3$ et $8$ sont premiers entre eux,
$$\Z/24\Z \simeq \Z/3\times\Z/8\Z.$$
 \item Par la question précédente, $(\Z/24\Z)^\times\simeq (\Z/3\times\Z/8\Z)^\times
\simeq (\Z/3\Z)^\times\times (\Z/8\Z)^\times\simeq \Z/2\Z\times (\Z/8\Z)^\times.$
Pour déterminer le groupe $(\Z/8\Z)^\times$, on observe que c'est un groupe à $\phi(8)=4$ éléments.
$(\Z/8\Z)^\times=\{1,3,-1,-3\}$. Comme $3^2=1$, tous les éléments sont d'ordre $2$.
Par conséquent, tous les éléments de $(\Z/24\Z)^\times$ sont d'ordre 2, et ce groupe n'est pas cyclique.
\end{enumerate}
\fincorrection
\correction{007377}
\begin{enumerate}
 \item $221=13\times17$
 \item Par le théorème chinois,
\begin{eqnarray*}
 (X-3)(X-5)=0\pmod{221}&\iff& \left\{\begin{array}{c}(X-3)(X-5)=0\pmod{13} \\ (X-3)(X-5)=0\pmod{17}.\end{array}\right.
\end{eqnarray*}
Puisque $13$ et $17$ sont premiers, par le lemme d'Euclide,
\begin{eqnarray*}
\lefteqn{\left\{\begin{array}{c}(X-3)(X-5)=0\pmod{13} \\ (X-3)(X-5)=0\pmod{17}.\end{array}\right.}\\
&\iff&\left\{\begin{array}{c} (X=3\pmod{13}\quad\text{ou}\quad X=5\pmod{13})\\
\quad \text{ et } \quad (X=3\pmod{17}\quad\text{ou}\quad X=5\pmod{17})\end{array}\right.\\
&\iff& \left\{\begin{array}{ccc}
    X=3\pmod{13}&\quad \text{ et } \quad &X=3\pmod{17}\\ &\quad\text{ou}\quad&\\
 X=3\pmod{13}&\quad \text{ et } \quad &X=5\pmod{17}\\ &\quad\text{ou}\quad&\\
 X=5\pmod{13}&\quad \text{ et } \quad &X=3\pmod{17}\\ &\quad\text{ou}\quad&\\
 X=5\pmod{13}&\quad \text{ et } \quad &X=5\pmod{17}\\ 
    \end{array}\right.\\
&\iff&X=3\pmod{221}\quad\text{ou}\quad X=107\pmod{221}\\
&&\quad\text{ou}\quad X=122\pmod{221}\quad\text{ou}\quad X=5\pmod{221}.
\end{eqnarray*}
Les solutions entières de l'équation
$(X-3)(X-5)=0\pmod{221}$ sont les entiers congrus à $3$, $107$, $122$ ou $5$ modulo $221$.
\end{enumerate}
\fincorrection
\correction{007378}
\begin{enumerate}
 \item Comme $\zeta$ est d'ordre $8$, $\zeta^3$ est d'ordre 
$8/\pgcd(8,3)=8$. Le sous-groupe engendré par $\zeta^3$ est donc tout le groupe des racines de l'unité d'ordre $8$.

\includegraphics[scale=0.7]{images/img-mour-130}
\item $z^11=1$ donc l'ordre de $z$ est un diviseur de $11$. Mais $z\not=1$. Donc $z$ est d'ordre $11$.
Comme $z^8\not=1$, $z^8$ est aussi d'ordre $11$.
$z^8=\exp(16\times 8 i \pi/11)=\exp(64\times 2i \pi/11)$.
Comme $64=5\times 11+9$, l'argument de $z^8$ est $18\pi/11$.
\end{enumerate}
\fincorrection
\correction{007379}
\begin{enumerate}
 \item 
$2^3<7<3^2$. Donc, $2<\sqrt{7}<3$.
Soit $P(x)=x^2-7$.
$$\begin{array}{c|ccc}
 &1&0&-7\\ \hline \textrm{facteur }2&1&2&-3\\
 &1&4&\\ &1&&
\end{array}$$
Donc
$P(2+y)=y^2+4y-3$. Par conséquent,
$Q(z):=100P(2+z/10)=z^2+40z-300$.
Avec $300/40\sim 7$
On trouve $Q(6)<0$ et $Q(7)>0$.
Donc $$2,6<\sqrt{7}<2.7.$$

 \item $Q(z)=z^2+40z-300$.
$$\begin{array}{c|ccc}
 &1&40&-300\\ \hline \textrm{facteur }6&1&46&-24\\
 &1&52&\\ &1&&
\end{array}$$
Donc, $Q(6+u)=u^2+52u-24$. Par conséquent,
$R(v):=100Q(6+v/10)=v^2+520u-2400$.
Comme $2400/520\sim 5$, on trouve $R(5)>0$ et $R(4)<0$.
Donc $$2,64<\sqrt{7}<2.65.$$

\end{enumerate}
\fincorrection
\nocorrection
\nocorrection
\nocorrection
\nocorrection
\nocorrection
\nocorrection
\nocorrection
\nocorrection
\nocorrection
\nocorrection
\correction{007390}
\begin{enumerate}
	\item
	\item $60^2=-53[281]$.
$53^2=-1[281]$.
Donc, $c=53$ est une racine de $-1$ modulo $281$.
	\item 
$281=5\times(53+i)+(16-5i)$.
$53-i=(3+i)(16-5i)$.
Donc, $\pgcd(281, c+i)=16-5i$.
En calculant $N(16-5i)$ on trouve $$16^2+5^2=281.$$
\end{enumerate}
\fincorrection
\nocorrection
\nocorrection
\nocorrection
\nocorrection
\nocorrection
\correction{007396}
\begin{enumerate}
 \item Si $G$ est un groupe fini et $H$ un sous-groupe de $G$, alors $\text{card}(G)=\text{card}(G/H)\text{card}(H)$. En particulier, l'ordre d'un élément d'un groupe fini divise l'ordre du groupe.
 \item L'ordre de $a^p$ est $\frac{k}{k\wedge p}$. 
 \item On note que l'ordre de $\mathcal{S}_3$ est $6$. Un groupe fini d'ordre $6$ est cyclique si est seulement s'il possède un élément d'ordre $6$. Il y a trois éléments d'ordre 2, deux éléments d'ordre 3, et un élément d'ordre 1 dans $\mathcal{S}_3$. Alors le groupe n'est pas cyclique.
Noter aussi qu'il n'est pas commutatif.
 \item Soit $D\in k[X]$ un polynôme non nul. Pour tout polynôme $A$ de $k[X]$ il existe un unique couple $(Q,R)\in (k[X])^2$
tel que $A=DQ+R$ et $\text{deg} (R)< \text{deg} (D)$.
 \item Soit $A$ et $B$ deux entiers de Gauss avec $B\neq 0$. Alors il existe deux entiers de Gauss $Q$ et $R$ tels que 
$A=BQ+R$ et $N(R)<N(B)$.
\end{enumerate} 
\fincorrection
\correction{007397}
\begin{enumerate}
 \item Puisque $\text{pgcd}(51,131)=1$, la classe $[51]$ est inversible dans $\Z/131\Z$. En utilisant l'algorithme d'Euclide on trouve que $1=131\times (-7)+51\times 18$. Alors, dans $\Z/131\Z$ on a $51^{-1}=18$ et $92\times 51^{-1}=-39\times 18=-47=84$. 
 \item Un élément $a$ dans $\Z/16\Z$ est un diviseur de 0 s'il existe un élément non nul $b$ tel que $ab=0 \text{ mod } 16$. Alors les diviseurs de 0 sont tous les éléments $a$ dans $\Z/16\Z$ tels que $\text{pgcd}(a,16)\neq 1$. On trouve: $2, 4, 6, 8, 10, 12, 14$.
\end{enumerate} 
\fincorrection
\correction{007398}
\begin{enumerate}
 \item Si le polynôme $P=X^4+X+1$ est réductible, il est divisible par un polynôme de degré~ 1 (donc $X$ ou $X+1$) ou par un polynôme irréductible de degré 2 (donc $X^2+X+1$). On vérifie que les polynômes $X, X+1$ et $X^2+X+1$ ne divisent pas $P$, alors $P$ est irréductible. 
 \item On effectue la division euclidienne de $3X^5+X^2+X+7$ par $X^4+X+1$. Puisque $3X^5+X^2+X+7=X^5+X^2+X+1=X(X^4+X+1)+1=1$ dans $A$, la classe de $3X^5+X^2+X+7$ n'est pas nulle. \\
 L'anneau $\mathbb{F}_2[X]/P$ est un corps si et seulement si le polynôme $P$ est irréductible dans $\mathbb{F}_2[X]$. Comme $X^4+X+1$ est irréductible, $A$ est un corps. $A$ est de cardinal $2^4=16$. 
 \item On a une relation $\alpha^4+\alpha+1=0$ dans $A$, alors $\alpha^4=\alpha+1$. Le groupe des inversibles $A^{\times}$ de $A$ contient $15$ éléments, alors le théorème de Lagrange implique que l'élément inversible $\alpha$ vérifie $\alpha^{15}=1$.
 \item Puisque $\alpha^4+\alpha+1=0$ et $\alpha^{15}-1=0$, les polynômes $X^{15}-1$ et $X^4+X+1$ représentent la même classe (classe nulle) dans $A$. Alors $X^{15}-1$ est bien un multiple de $X^4+X+1$ dans $\mathbb{F}_2[X]$. 
\end{enumerate} 
\fincorrection
\correction{007399}
\begin{enumerate}
 \item Soit $C$ le code donné. L'alphabet de $C$ est l'ensemble $\{0,1\}$. La longueur des mots est 15, alors la longueur de C est 15. 
La matrice génératrice donnée est exactement la matrice du code engendré par le polynôme $g=1+X^3+X^4.$ Puisque $g$ est un diviseur unitaire de $X^{15}-1$ de degré $4$, 
$$X^{15}-1=(X^4+X^3+1)(X^{11}+X^{10}+X^9+X^8+X^6+X^4+X^3+1)$$
le code engendré par $g$ est cyclique, de dimension $11=15-4$. Alors $C$ est cyclique de dimension 11 engendré par $g$. Il contient $2^{11}$ mots. 
 \item Oui, voir au-dessus.
 \item On trouve le polynôme de contrôle: 
$$h=\frac{X^{15}-1}{g}=X^{11}+X^{10}+X^9+X^8+X^6+X^4+X^3+1$$
Alors la matrice de contrôle est
$$H=\left(\begin{array}{ccccccccccccccc}
1&1&1&1&0&1&0&1&1&0&0&1&0&0&0\\
0&1&1&1&1&0&1&0&1&1&0&0&1&0&0\\
0&0&1&1&1&1&0&1&0&1&1&0&0&1&0\\
0&0&0&1&1&1&1&0&1&0&1&1&0&0&1\\
\end{array}\right)$$
Deux colonnes quelconques de $H$ sont toujours linéairement indépendantes, donc la distance du code est au moins~ $3$. Puisque le code contient un mot de poids 3, alors la distance est en fait $3$. Alors $C$ peut corriger une erreur et en détecter deux.
 \item Un mot $m$ appartient à $C$ ssi $H\cdot \mathrm{tr} m=\vec{0}$. \\
Le produit de $H$ et $m=(0,0,0,1,0,1,1,1,1,1,1,0,0,1,1)$ est la colonne obtenue comme somme des colonnes étoilées
$$\left(\begin{array}{ccccccccccccccc}
0&0&0&*&0&*&*&*&*&*&*&0&0&*&*\\ \hline
1&1&1&1&0&1&0&1&1&0&0&1&0&0&0\\
0&1&1&1&1&0&1&0&1&1&0&0&1&0&0\\
0&0&1&1&1&1&0&1&0&1&1&0&0&1&0\\
0&0&0&1&1&1&1&0&1&0&1&1&0&0&1\\
\end{array}\right)$$
soit $(0,0,0,0)$. Alors $m$ est un mot de code.
\end{enumerate} 
\fincorrection
\correction{007400}
\begin{enumerate}
 \item On vérifie que les premiers entre $2$ et $\sqrt{613}$, c'est-à-dire, $2, 3, 5, 7, 11, 13, 17, 19, 23$  ne divisent pas $613$. Alors le nombre $613$ est premier. 
 \item Puisque $613=1 \text{ mod }4$, d'après le Théorème de deux carrés, $613$ peut s'écrire comme somme de deux carrés.

 \item $35^2=1225=-1 \text{ mod } 613$.
 \item $$\frac{613}{35+i}=\frac{613\times(35-i)}{(35+i)\times(35-i)}=\frac{613\times(35-i)}{1225+1}=\frac{35}{2}-\frac{i}{2}.$$
Un entier de Gauss le plus proche à $\frac{35}{2}-\frac{i}{2}$ est $17$. Alors 
$$613=17\times(35+i)+(18-17i)$$
 \item On a trouvé que $35$ est une racine de $-1$ modulo $p=613$. Trouvons $\text{pgcd(613,35+i)}$.
Puisque $613=17\times(35+i)+(18-17i)$, $\text{pgcd}(613,35+i)=\text{pgcd}(35+i,18-17i)=18-17i$. Donc 
$$613=N(18-17i)=18^2+17^2$$
\end{enumerate} 
\fincorrection
\correction{007401}
\begin{enumerate}
 \item Pour tout $P\in k[X]$ et $D\in k[X]^*$, il existe un unique
couple $(Q,\,R)$ dans $k[X]$ tel que $P=QD+R$ et $\deg(R)<\deg(D)$.
 \item Pour tout $a\in \Z[i]$ et $d\in \Z[i]^*$, il existe un
couple $(q,\,r)$ dans $k[X]$ tel que $a=qd+r$ et $|r|<|d|$.
 \item L'ordre de $a^k$ est $\frac{n}{n\wedge k}$. En effet, soient
$d=n\wedge k$, $n'=n/d$ et $k'=k/d$, de sorte que $n'\wedge k'=1$.
Pour tout $m$, on a $a^{km}=1\iff n|km \iff n'| k'm\iff n'|m$. Donc
l'ordre de $a^k$ est $n'$.
 \item Soit $G$ d'ordre 13 et $x$ un élément non trivial de $G$. Par
le théorème de Lagrange, $x$ ne peut qu'être d'ordre 13. Mais
$\langle x\rangle$ contient déjà 13 éléments, donc $G=\langle
x\rangle$. Ainsi $G$ est abélien car monogène.
 \item Par exemple 7, puisque 7 est congru à 3 modulo 4 ou bien pusique ni 7-$1^2$=6, ni $7-2^2=3$, ni avec $n\in\N, n\geq 3$, $7-n^2<0$ ne sont des carrés.
\end{enumerate} 
\fincorrection
\correction{007402}
\begin{enumerate}
 \item Non, car $\Z/3\Z\times \Z/3\Z$ ne contient aucun élément d'ordre 9.
 \item Non, car ils n'ont pas le même ordre.
 \item Les groupes $(\mathbb{F}_7)^\star$ et $\Z/2\Z\times \Z/3\Z$ sont-ils isomorphes ? Par le théorème sur le groupe des inversibles d'un corps fini et par le théorème chinois, ils sont tous les deux cycliques d'ordre $6$, donc isomorphes. 
\end{enumerate} 
\fincorrection
\correction{007403}
\begin{enumerate}
 \item Soit $P=X^2+X+1$. Le polynôme $P$ est sans racine, donc s'il
est réductible, il se factorise en deux polynômes de degré 2. Or
$(X^2+X+1)^2=X^4+X^2+1$, qui est différent de $P$. Donc $P$ est
irréductible.
 \item L'idéal $(X^4+X^3+1)$ est premier car $X^4+X^3+1$ est
irréductible, donc il est maximal car $\mathbb{F}_2[X]$ est principal, comme
tout anneau de polynômes sur un corps. Or le quotient par un idéal
maximal est un corps. Donc $A$ est un corps. De plus, il est
isomorphe au polynôme de $\mathbb{F}_2[X]$ de degré au plus 3, donc $A$
possède 16 éléments.
 \item Puisque $A$ est un corps, $A^*$ contient 15 éléments, donc
l'ordre de tout élément de $A^*$ est un diviseur de 15. En
particulier, $\alpha^{15}=1$.
 \item 
$$\begin{array}{l}
\alpha^1=\alpha\\
\alpha^2=\alpha^2\\
\alpha^3=\alpha^3\\
\alpha^4=\alpha^3+1\\
\alpha^5=\alpha^4+\alpha=\alpha^3+\alpha+1\\
\end{array}$$
$$\begin{array}{l}
\alpha^6=\alpha^4+\alpha^2+\alpha=\alpha^3+\alpha^2+\alpha+1\\
\alpha^7=\alpha^4+\alpha^3+\alpha^2+\alpha=\alpha^2+\alpha+1\\
\alpha^8=\alpha^3+\alpha^2+\alpha\\
\alpha^9=\alpha^4+\alpha^3+\alpha^2=\alpha^2+1\\
\alpha^{10}=\alpha^3+\alpha\\
\end{array}$$
$$\begin{array}{l}
\alpha^{11}=\alpha^4+\alpha^2=\alpha^3+\alpha^2+1\\
\alpha^{12}=\alpha^4+\alpha^3+\alpha=\alpha+1\\
\alpha^{13}=\alpha^2+\alpha\\
\alpha^{14}=\alpha^3+\alpha^2\\
\alpha^{15}=1
\end{array}$$

 \item $\alpha^7+\alpha^8+\alpha^9=\alpha^7(1+\alpha+\alpha^2)=\alpha^{7+7}=\alpha^{14}=\alpha^3+\alpha^2$.
 \item $(1+\alpha+\alpha^3)^{-1}=(\alpha^5)^{-1}=\alpha^{10}$.
\end{enumerate} 
\fincorrection
\correction{007404}
\begin{enumerate}
 \item Le groupe $\mathbb{F}_{19}^\times$ est d'ordre 18, donc les ordres
possibles des éléments sont les diviseurs de 18 (et ce sont
exactement ceux-là car $\mathbb{F}_{19}^\times$ est cyclique), c'est-à-dire
1,2,3,6,9,18. Or $2^1=2$, $2^2=4$, $2^3=8$, $2^6=7$, $2^9=-1$, donc
2 est bien un générateur de $\mathbb{F}_{19}^\times$.
 \item $C=2^3=8$.
 \item $D=2^7=14$. $C^d=8^7=2^{21}=2^3=8$. $mC^d=11\times8=12$.
Finalement, $(M_1,M_2)=(14,12)$.
 \item Il utilise la formule : $m=M_2(M_1^c)^{-1}$. On vérifie en effet
que $12\times 14^{-3}=12\times(-5)^{-3}=12\times (-7)=-8=11$.
 \item Le message envoyé est $3\times8^{-3}=3\times 2^{-9}=-3=16$.
Pour crypter $m=16$, Alice a choisi une clé privée $k$ et a envoyé le
message $(2^k,\,16\times 8^k)$. Or on sait que $2^3=8$ et que le
logarithme discret est un isomorphisme, donc Alice a choisi cette
fois la clé privée 3. Et effectivement, on retrouve bien
$(2^3,\,16\times 8^3)=(8,3)$.
\end{enumerate} 
\fincorrection
\nocorrection
\nocorrection
\nocorrection
\nocorrection
\nocorrection
\nocorrection
\nocorrection
\nocorrection
\nocorrection
\nocorrection
\nocorrection
\nocorrection
\nocorrection
\nocorrection
\nocorrection
\nocorrection
\nocorrection
\nocorrection
\nocorrection
\nocorrection
\nocorrection
\nocorrection
\nocorrection
\nocorrection
\nocorrection
\nocorrection
\nocorrection
\nocorrection
\nocorrection
\nocorrection
\nocorrection
\nocorrection
\nocorrection
\nocorrection
\correction{003660}
\begin{enumerate}
  \item $a > 0, b = c, d > 0, ad - bc > 0$.
  \item $a - b > 0 \text{ et } a + (n-1)b > 0$.
  \item 
\end{enumerate}
 \fincorrection
\correction{003661}
$\left( \frac1{\sqrt2}, X\sqrt{\frac32}, (3X^2-1)\sqrt{\frac58} + (5X^3-3X)\sqrt{\frac78}\, \right)$.
\fincorrection
\correction{003662}
$\frac 1{\sqrt5}, \frac {X-2}{\sqrt{10}}, \frac {X^2-4X+2}{\sqrt{14}}$.
\fincorrection
\correction{003663}
\begin{enumerate}
  \item 
  \item \'Elever au carré.
  \item   
  \begin{enumerate}
    \item $(\vec x\mid \vec u) = 1
                  \Leftrightarrow (i(\vec x) \mid \vec u-i(\vec x)) = 0$ :
                  sphère passant par $\vec 0$.
    \item Hyperplan ne passant pas par $\vec 0$.
    \item $\|\vec x-\vec a\,\|^2 = R^2 \Leftrightarrow
                  \left\|\vec x -\frac{\vec a}{\|\vec a\,\|^2-R^2}\right\|^2 =
                  \frac{R^2}{(\|\vec a\,\|^2-R^2)^2}$ :
                  sphère ne passant pas par $\vec 0$.
  \end{enumerate}    
\end{enumerate}
 \fincorrection
\correction{003664}
\begin{enumerate}
  \item \'Elever au carré.
  \item 
\end{enumerate}
 \fincorrection
\correction{003665}
\begin{enumerate}
  \item 
  \item $\frac1{\sqrt{n+1}}$.
\end{enumerate}
 \fincorrection
\correction{003666}
\begin{enumerate}
  \item $\biggl(\frac 1{\sqrt6}(1,-2,1,0), \frac 1{\sqrt{30}}(2,-1,-4,3) \biggr)$
  \item $\frac 1{10}\begin{pmatrix}3 &-4 &-1 &2  \cr
                          -4 &7  &-2 &-1 \cr
                          -1 &-2 &7 &-4  \cr
                           2 &-1 &-4 &3  \cr\end{pmatrix}$.
  \item $\sqrt{\frac 7{10}}$.
\end{enumerate}
 \fincorrection
\correction{003667}
$\frac1{\Sigma a_i^2}\Bigl( I- (a_ia_j) \Bigr)$.
\fincorrection
\nocorrection
\nocorrection
\nocorrection
\correction{003671}
{\def\IIII{(\Im p\cap\Im q)}%
\def \orthooo{\buildrel\bot\over\oplus}%
Si $p\circ q = q\circ p$ :
Soient $x\in \IIII^\bot \cap \Im p$ et $y\in \IIII^\bot \cap \Im q$.

Alors $p\circ q(x) = q(x) \in \Im p\cap \Im q$,
donc $(q(x)\mid y)=(x\mid y) = 0$.
Si
$A = \IIII^\bot \cap \Im p$ et $B = \IIII^\bot \cap \Im q$ sont orthogonaux :
Alors $\Im p = \IIII\orthooo A$, $\Im q = \IIII \orthooo B$, et
$E = \IIII \orthooo A \orthooo B \orthooo (\Im p^\bot \cap \Im q^\bot)$.
Par décomposition, on obtient $p\circ q = q \circ p ={} $
la projection orthogonale sur $\Im p \cap \Im q$.}
\fincorrection
\correction{003672}
\begin{enumerate}
  \item  $\sum_{i=1}^n (\vec e_j\mid \vec e_i)^2 = 1  \Rightarrow {}$
        famille orthonormée et vect$(\vec e_i)^\perp$ = $\{\vec 0\}$.
  \item 
\end{enumerate}
 \fincorrection
\nocorrection
\nocorrection
\correction{003675}
sphère de centre $-\frac {\gamma \vec a}{\beta \|\vec a\,\|^2}$.
\fincorrection
\nocorrection
\nocorrection
\nocorrection
\correction{003679}
Soit $X$ la matrice de $\vec e_n$ dans ${\cal B}$.
         On a $GX = \begin{pmatrix}0\cr\vdots\cr0\cr\lambda\cr\end{pmatrix}$ et ${}^tXGX = \lambda x_p = 1$.
	 On applique alors les formules de Cramer.
\fincorrection
\nocorrection
\correction{003681}
Non, $\|\vec u_1+\vec u_2+\vec u_3\|^2 < 0$.
\fincorrection
\nocorrection
\correction{003683}
\begin{enumerate}
  \item 
  \item 
  \item $ \int_{t=0}^1 t^kt^x\,d t = \frac 1{k+x+1}$.
  \item $\Phi$ a pour pôles au plus simples $-1,-2,\dots,-n-1$ et
             pour racines $0,1,\dots,n-1$. Comme $\Phi(x) \to 0$ lorsque $x\to\infty$,
             on a donc $\Phi(x) = \lambda\frac{x(x-1)\dots(x-n+1)}{(x+1)\dots(x+n+1)}$.
  \item $a_k =$ résidu de $\Phi$ en $-k-1 =
             (-1)^{n+k}\lambda\frac{(n+k)!}{(k!)^2(n-k)!}$.
  \item
\end{enumerate}
 \fincorrection
\correction{003684}
$x^2 + (x+y)^2 + (x+2y)^2 = (\sqrt3(x-y))^2 + (\sqrt2y)^2$.
\fincorrection
\nocorrection
\nocorrection
\correction{003687}
$f \in F^\perp  \Rightarrow  xf\perp f$.
\fincorrection
\correction{003688}
\begin{enumerate}
  \item 
  \item $30X^2-36X+9$.
\end{enumerate}
 \fincorrection
\correction{003689}
$P_a(t) = \frac{3}{8}(3 -5t^2 -5a^2 +15a^2t^2)
                   +\frac{5at}8(15-21t^2 -21a^2 +35a^2t^2)$,\par
$8\|P_a\|^2 = 9 + 45a^2 - 165a^4 + 175a^6$ est maximal pour $a=\pm1  \Rightarrow  \|P_a\| = 2\sqrt2$.
\fincorrection
\correction{003690}
\begin{enumerate}
  \item 
  \item $\pi(f)(t) = f(0)\frac{\sh(1-t)}{\sh(1)} + f(1)\frac{\sh(t)}{\sh(1)}.$
  \item L'inf est atteint pour la fonction~$f\in W$ telle
    que $f(0) = \alpha$ et $f(1) = \beta$, soit
    $f(t) = \alpha\frac{\sh(1-t)}{\sh(1)} + \beta\frac{\sh(t)}{\sh(1)}$
    et $inf=\frac{(\alpha^2+\beta^2)\ch(1) - 2\alpha\beta}{\sh(1)}$.
\end{enumerate}
 \fincorrection
\correction{003691}
\begin{enumerate}
  \item Le sous-espace vectoriel engendré a un orthogonal nul.
  \item N'importe quelle famille génératrice convient (équivalence des normes).
  \item $1=\|y_i\|^2 = \|y_i\|^4 + \sum_{j\ne i}(y_i\mid y_j)^2 \Rightarrow \forall\ j\ne i,\ (y_i\mid y_j) = 0$.
  \item Par polarisation on a~:
    $\forall\ x,y,\ \sum_{j\in I}(x\mid y_j)(y\mid y_j) = A(x\mid y)$
    donc $\sum_{j\in I}(x\mid y_j)y_j-Ax\in E^\bot$.
\end{enumerate}
 \fincorrection
\correction{003692}
Soient $x\in\mathrm{Ker}(u-\mathrm{id})$ et $y=u(z)-z\in\Im(u-\mathrm{id})$.
On a $y = u(z+\lambda x)-(z+\lambda x)$ d'où~:
$$\|z+\lambda x\|^2\ge \|u(z+\lambda x)\|^2 = \|z+\lambda x\|^2 + 2\lambda(x\mid y) + 2(z\mid y) + \|y\|^2.$$
En faisant tendre $\lambda$ vers $\pm\infty$ on obtient $(x\mid y) = 0$
et on conclut avec le théorème du rang.
\fincorrection
\correction{003693}
$f$ linéaire et $f = x \mapsto\|x\|^2$ conviennent et l'ensemble~$\cal E$
des fonctions $f$ vérifiant la propriété est stable par combinaison
linéaire donc toute fonction de la forme $x \mapsto\ell(x) + a\|x\|^2$
avec $\ell\in E^*$ et $a\in\R$ convient. On montre que ce sont les
seules~:
Soit $f\in{\cal E}$ l'on décompose en sa partie paire $f_p$
et sa partie impaire~$f_i$. Alors $f_p,f_i\in{\cal E}$.

Soient $x,y\in E$ avec $\|x\|=\|y\|$ et $x\perp y$.
On a $f_i(x\pm y) = f_i(x)\pm f_i(y)$ et $f_i(2x) = f_i(x+y) + f_i(x-y) = 2f_i(x)$.
Ensuite, $f_i(2x) + f_i(x) -f_i(y) = f_i(2x+y) + f_i(x-2y) = f_i(3x-y) = f_i(3x) - f_i(y)$
d'où $f_i(3x) = 3f_i(x)$ et de proche en proche $f_i(kx) = kf_i(x)$ pour $k\in\N$
puis pour $k\in\Z,\Q,\R$ successivement vu la continuité de~$f$.
En prenant une base $(e_1,\dots,e_n)$ orthonormale on a
$f(x_1e_1+\dots+x_ne_n) = x_1f(e_1)+\dots+x_nf(e_n)$ pour tous $x_1,\dots,x_n$
réels donc $f_i$ est linéaire.

Soient à présent $x,y\in E$ avec $\|x\|=\|y\|$
alors $f_p(x+y) + f_p(x-y) = f(2x)$ et $f_p(x+y) + f_p(y-x) = f_p(2y)$
d'où $f_p(2x) = f_p(2y)$. Ainsi $f_p$ est constante sur les sphères
de centre~$0$. On écrit $f_p(x) = \varphi(\|x\|^2)$ avec $\varphi : {\R^+} \to \R$
prolongée à $\R$ par imparité ($f_p(0) = 0$ de manière évidente)
et on a $\varphi(a^2+b^2) = f_p(ae_1+be_2) = f_p(ae_1)+f_p(be_2) = \varphi(a^2)+\varphi(b^2)$
d'où l'on conclut que $\varphi$ est linéaire.
\fincorrection
\correction{005482}
Posons $\varphi~:~(A,B)\mapsto\mbox{Tr}({^t}AB)$. Montrons que $\varphi$ est un produit scalaire sur $\mathcal{M}_n(\Rr)$. 
\textbf{1ère solution.} \textbullet~$\varphi$ est symétrique. En effet, pour $(A,B)\in(\mathcal{M}_n(\Rr))^2$,

$$\varphi(A,B)=\mbox{Tr}(^{t}AB)=\mbox{Tr}({^t}({^t}AB))=\mbox{Tr}({^t}BA)=\varphi(B,A).$$

\textbullet~$\varphi$ est bilinéaire par linéarité de la trace et de la transposition.
\textbullet~Si $A=(a_{i,j})_{1\leq i,j\leq n}\in\mathcal{M}_n(\Rr)\setminus\{0\}$, alors

$$\varphi(A,A)=\sum_{i=1}^{n}\left(\sum_{j=1}^{n}a_{i,j}a_{i,j}\right)=\sum_{i,j}^{}a_{i,j}^2>0$$
car au moins un des réels de cette somme est strictement positif. $\varphi$ est donc définie, positive.

\textbf{2ème solution.} Posons $A=(a_{i,j})$ et $B=(b_{i,j})$. On a

\begin{center}
$\text{Tr}({^t}AB)=\sum_{j=1}^{n}\left(\sum_{i=1}^{n}a_{i,j}b_{i,j}\right)=\sum_{1\leq i,j\leq n}^{}a_{i,j}b_{i,j}$.
\end{center}
Ainsi, $\varphi$ est le produit scalaire canonique sur $\mathcal{M}_n(\Rr)$ et en particulier, $\varphi$ est un produit scalaire sur $\mathcal{M}_n(\Rr)$.
$N$ n'est autre que la norme associée au produit scalaire $\varphi$ (et en particulier, $N$ est une norme).
Soit $(A,B)\in(\mathcal{M}_n(\Rr))^2$.

\begin{align*}\ensuremath
N(AB)^2&=\sum_{i,j}^{}\left(\sum_{k=1}^{n}a_{i,k}b_{k,j}\right)^2\\
 &\leq\sum_{i,j}^{}\left(\sum_{k=1}^{n}a_{i,k}^2\right)\left(\sum_{l=1}^{n}b_{l,j}^2\right)\;(\mbox{d'après l'inégalité de \textsc{Cauchy}-\textsc{Schwarz}})\\
 &=\sum_{i,j,k,l}^{}a_{i,k}^2b_{l,j}^2=\left(\sum_{i,k}^{}a_{i,k}^2\right)\left(\sum_{l,j}^{}b_{l,j}^2\right)=N(A)^2N(B)^2,
\end{align*}
et donc,

\begin{center}
\shadowbox{
$\forall(A,B)\in(\mathcal{M}_n(\Rr))^2,\;N(AB)\leq N(A)N(B)$.
}
\end{center}
\fincorrection
\correction{005483}
\begin{enumerate}
 \item  Soit $(x,y,z)\in\Rr^3$.

\begin{align*}\ensuremath f(x+z,y)+f(x-z,y)&=\frac{1}{4}(||x+z+y||^2+||x-z+y||^2-||x+z-y||^2-||x-z-y||^2)\\
 &=\frac{1}{4}\left(2(||x+y||^2+||z||^2)-2(||x-y||^2+||z||^2)\right)=2f(x,y).
\end{align*}
 \item  $2f(x,y)= f(x+x,y)+f(x-x,y)=f(2x,y)+f(0,y)$ mais $f(0,y)=(||y||^2-||-y||^2)=0$ (définition d'une norme).
 \item  \textbullet~Montrons par récurrence que $\forall n\in\Nn,\;f(nx,y)=nf(x,y)$.
C'est clair pour $n=0$ et $n=1$.
Soit $n\geq0$. Si l'égalité est vraie pour $n$ et $n+1$ alors d'après 1), 

$$f((n+2)x,y)+f(nx,y)=f((n+1)x+x,y)+f((n+1)x-x,y)=2f((n+1)x,y),$$
et donc, par hypothèse de récurrence,

$$f((n+2)x,y)=2f((n+1)x,y)-f(nx,y)=2(n+1)f(x,y)-nf(x,y)=(n+2)f(x,y).$$
Le résultat est démontré par récurrence.
\textbullet~Soit $n\in\Nn^*$, $f(x,y)=f\left(n\times\frac{1}{n}.x,y\right)=nf\left(\frac{1}{n}x,y\right)$ et donc $f\left(\frac{1}{n}x,y\right)=\frac{1}{n}f(x,y)$.
\textbullet~Soit alors $r=\frac{p}{q}$, $p\in\Nn$, $q\in\Nn^*$, $f(rx,y)=\frac{1}{q}f(px,y)=p\frac{1}{q}f(x,y)=rf(x,y)$ et donc, pour tout rationnel positif $r$, $f(rx,y)=rf(x,y)$.
Enfin, si $r\leq0$, $f(rx,y)+f(-rx,y)=2f(0,y)=0$ (d'après 1)) et donc= $f(-rx,y)=-f(-rx,y)=rf(x,y)$.

\begin{center}
$\forall(x,y)\in E^2,\;\forall r\in\Qq,\;f(rx,y)=rf(x,y)$.
\end{center}
 \item  On pose $x=\frac{1}{2}(u+v)$ et $y=\frac{1}{2}(u-v)$.

$$f(u,w)+f(v,w)=f(x+y,w)+f(x-y,w)=2f(x,w)=2f\left(\frac{1}{2}(u+v),w\right)=f(u+v,w).$$
 \item  f est symétrique (définition d'une norme) et linéaire par rapport à sa première variable (d'après 3) et 4)).
Donc f est bilinéaire.
 \item  f est une forme bilinéaire symétrique.
Pour $x\in E$, $f(x,x)=\frac{1}{4}(||x+x||^2+||x-x||^2)=\frac{1}{4}||2x||^2=||x||^2$ (définition d'une norme) ce qui montre tout à la fois que $f$ est définie positive et donc un produit scalaire, et que $||\;||$ est la norme associée. $||\;||$ est donc une norme euclidienne.
\end{enumerate}
\fincorrection
\correction{005485}
Soit $A$ un éventuel polynôme solution c'est à dire tel que $\forall P\in\Rr[X],\;\int_{0}^{1}P(t)A(t)\;dt=P(0)$.

$P=1$ fournit $\int_{0}^{1}A(t)\;dt=1$ et donc nécessairement $A\neq0$. $P=XA$ fournit $\int_{0}^{1}tA^2(t)\;dt=P(0)=0$.
Mais alors, $\forall t\in[0,1],\;tA^2(t)=0$ (fonction continue positive d'intégrale nulle) puis $A=0$ (polynôme ayant une infinité de racines deux à deux distinctes). $A$ n'existe pas.
\fincorrection
\correction{005495}
Soit $\begin{array}[t]{cccc}
\varphi~:&E&\rightarrow&\Rr^n\\
 &x&\mapsto&(x|e_1,...,x|e_n)
\end{array}$. $\varphi$ est clairement linéaire et $\mbox{Ker}\varphi$ est $(e_1,...,e_n)^\bot=E^\bot=\{0\}$.
Comme $E$ et $\Rr^n$ ont mêmes dimensions finies, $\varphi$ est un isomorphisme d'espaces vectoriels. En particulier, pour tout $n$-uplet $(a_1,...,a_n)$ de réels, il existe un unique vecteur $x$ tel que $\forall i\in\llbracket1,n\rrbracket,\;x|e_i=a_i$.
\fincorrection
\correction{005496}
\textbf{1ère solution.}
Montrons par récurrence que sur $n=\mbox{dim}(E)$ que, si $(x_i)_{1\leq i\leq p}$ est obtusangle, $p\leq n+1$.
\textbullet~Pour $n=1$, une famille obtusangle ne peut contenir au moins trois vecteurs car si elle contient les vecteurs $x_1$ et $x_2$ verifiant $x_1.x_2<0$, un vecteur $x_3$ quelconque est soit nul (auquel cas $x_3.x_1=0$), soit de même sens que $x_1$ (auquel cas $x_1.x_3>0$) soit de même sens que $x_2$ (auquel cas $x_2.x_3>0$). Donc $p\leq2$.
\textbullet~Soit $n\geq 1$. Suppososons que toute famille obtusangle d'un espace de dimension $n$ a un cardinal inférieur ou égal à $n+1$.
Soit $(x_i)_{1\leq i\leq p}$ une famille obtusangle d'un espace $E$ de dimension $n+1$. Si $p=1$, il n'y a plus rien à dire. Supposons $p\geq2$. $x_p$ n'est pas nul et $H=x_p^\bot$ est un hyperplan de $E$ et donc est de dimension $n$.
Soit, pour $1\leq i\leq p-1$, $y_i=x_i-\frac{(x_i|x_p)}{||x_p||^2}x_p$ le projeté orthogonal de $x_i$ sur $H$.
Vérifions que la famille $(y_i)_{1\leq i\leq p-1}$ est une famille obtusangle. Soit $(i,j)\in\llbracket1,p-1\rrbracket$ tel que $i\neq j$.
 
$$y_i.y_j=x_i.x_j-\frac{(x_i|x_p)(x_j|x_p)}{||x_p||^2}-\frac{(x_j|x_p)(x_i|x_p)}{||x_p||^2}+\frac{(x_i|x_p)(x_j|x_p)(x_p|x_p)}{||x_p||^4}= x_i|x_j-\frac{(x_i|x_p)(x_j|x_p)}{||x_p||^2}<0.$$
Mais alors, par hypothèse de récurrence, $p-1\leq 1+\mbox{dim}H=n+1$ et donc $p\leq n+2$. 
Le résultat est démontré par récurrence.
\textbf{2ème solution.}
Montrons que si la famille $(x_i)_{1\leq i\leq p}$ est obtusangle, la famille $(x_i)_{1\leq i\leq p-1}$ est libre.
Supposons par l'absurde, qu'il existe une famille de scalaires $(\lambda_i)_{1\leq i\leq p-1}$ non tous nuls tels que $\sum_{i=1}^{p-1}\lambda_ix_i=0\;(*)$.
Quite à multiplier les deux membres de $(*)$ par $-1$, on peut supposer qu'il existe au moins un réel $\lambda_i>0$. Soit $I$ l'ensemble des indices $i$ tels que $\lambda_i>0$ et $J$ l'ensemble des indices $i$ tels que $\lambda_i\leq0$ (éventuellement $J$ est vide). $I$ et $J$ sont disjoints.
(*) s'écrit $\sum_{i\in I}^{}\lambda_ix_i=-\sum_{i\in J}^{}\lambda_ix_i$ (si $J$ est vide, le second membre est nul).
On a 

$$0\leq\left\|\sum_{i\in I}^{}\lambda_ix_i\right\|^2=\left(\sum_{i\in I}^{}\lambda_ix_i\right).\left(-\sum_{i\in  J}^{}\lambda_ix_i\right)=\sum_{(i,j)\in I\times J}^{}\lambda_i(-\lambda_j)x_i.x_j\leq 0.$$
Donc, $\left\|\sum_{i\in I}^{}\lambda_ix_i\right\|^2=0$ puis $\sum_{i\in I}^{}\lambda_ix_i=0$.
Mais, en faisant le produit scalaire avec $x_p$, on obtient $\left(\sum_{i\in I}^{}\lambda_ix_i\right).x_p=\sum_{i\in I}^{}\lambda_i(x_i.x_p)<0$ ce qui est une contradiction.
La famille $(x_i)_{1\leq i\leq p-1}$ est donc libre. Mais alors son cardinal $p-1$ est inférieur ou égal à la dimension $n$ et donc $p\leq n+1$.
\fincorrection
\correction{005629}
\begin{enumerate}
 \item  Soit $a\in\Cc$. Soient $(\lambda,\mu)\in\Cc^2$ et $(P,Q)\in E^2$.

\begin{center}
$\varphi_a(\lambda P+\mu Q)=(\lambda P+\mu Q)(a)=\lambda P(a)+\mu Q(a)=\lambda\varphi_a(P)+\mu\varphi_a(Q)$.
\end{center}

Donc, $\varphi_a$ est une forme linéaire sur $E$.

\item  On a déjà $\text{card}\left(\varphi_{a_j}\right)_{0\leqslant j\leqslant n}=n+1=\text{dim}(E)=\text{dim}(E^*)<+\infty$. Il suffit donc de vérifier que la famille $\left(\varphi_{a_j}\right)_{0\leqslant j\leqslant n}$ est libre.

Pour $k\in\llbracket0,n\rrbracket$, on pose $P_k=\prod_{j\neq k}^{}\frac{X-a_j}{a_k-a_j}$. Chaque $P_k$ est un élément de $E$ et de plus 

\begin{center}
$\forall(j,k)\in\llbracket0,n\rrbracket^2$, $\varphi_{a_j}(P_k)=\delta_{j,k}=\left\{
\begin{array}{l}
1\;\text{si}\;j\neq k\\
0\;\text{si}\;j=k
\end{array}
\right.$\quad$(*)$.
\end{center}

Soit alors $(\lambda_0,\ldots,\lambda_n)\in\Cc^{n+1}$.

\begin{align*}\ensuremath
\sum_{j=0}^{n}\lambda_j\varphi_j=0&\Rightarrow\forall P\in E,\;\sum_{j=0}^{n}\lambda_j\varphi_j(P)=0\Rightarrow\forall k\in \llbracket0,n\rrbracket,\;\sum_{j=0}^{n}\lambda_j\varphi_j(P_k)=0\Rightarrow\forall k\in \llbracket0,n\rrbracket,\;\sum_{j=0}^{n}\lambda_j\delta_{j,k}=0\\
 &\Rightarrow\forall k\in \llbracket0,n\rrbracket,\;\lambda_k=0.
\end{align*}

Ceci montre que la famille $\left(\varphi_{a_j}\right)_{0\leqslant j\leqslant n}$ est libre et donc une base de $E^*$. Les égalités $(*)$ montrent alors que la préduale de la base $\left(\varphi_{a_j}\right)_{0\leqslant j\leqslant n}$ de $E^*$ est la famille $\left(P_k\right)_{0\leqslant k\leqslant n}$.

\item  Pour $P\in E$, posons $\varphi(P)=\int_{0}^{1}P(t)\;dt$. $\varphi$ est une forme linéaire sur $E$ et donc, puisque la famille $\left(\varphi_{a_j}\right)_{0\leqslant j\leqslant n}$ est une base de $E^*$, il existe $(\lambda_0,\ldots,\lambda_n)\in\Cc^{n+1}$ tel que $\varphi=\sum_{j=0}^{n}\lambda_j\varphi_{a_j}$ ou encore il existe $(\lambda_0,\ldots,\lambda_n)\in\Cc^{n+1}$ tel que pour tout $P\in E$, $\int_{0}^{1}P(t)\;dt=\lambda_0P(a_0)+\ldots+\lambda_nP(a_n)$ (les $\lambda_j$ étant indépendants de $P$).

En appliquant cette dernière égalité au polynôme $P_k$, $0\leqslant k\leqslant n$, on obtient $\lambda_k=\int_{0}^{1}P_k(t)\;dt=\int_{0}^{1}\prod_{j\neq k}^{}\frac{t-a_j}{a_k-a_j}\;dt$.

\begin{center}
\shadowbox{
$\forall P\in\Cc_n[X],\;\int_{0}^{1}P(t)\;dt=\sum_{k=0}^{n}\lambda_kP(a_k)$ où $\forall k\in\llbracket0,n\rrbracket$, $\lambda_k=\int_{0}^{1}\prod_{j\neq k}^{}\frac{t-a_j}{a_k-a_j}\;dt$.
}
\end{center}
\end{enumerate}
\fincorrection
\correction{005630}
 Les quatre applications $\varphi_1$, $\varphi_2$, $\psi_1$ et $\psi_2$ sont effectivement des formes linéaires sur $E$.

Cherchons tout d'abord la future  base préduale de la famille $(\varphi_1,\varphi_2,\psi_1,\psi_2)$. On note $(P_0,P_1,P_2,P_3)$ cette future base.

\textbullet~On doit avoir $\varphi_1(P_2)=\varphi_2(P_2)=\psi_2(P_2)=0$ et $\psi_1(P_2)=1$. Ainsi, $P_2$ s'annule en $0$ et en $1$ et de plus $P_2'(1) = 0$. Donc $P_2$ admet $0$ pour racine d'ordre $1$ au moins et $1$ pour racine d'ordre 2 au moins. Puisque $P_2$ est de degré inférieur ou égal à $3$, il existe une constante $a$ telle que $P_2=aX(X-1)^2=aX^3-2aX^2+aX$ puis $P_2'(0) = 1$ fournit $a = 1$ puis $P_2 =X(X-1)^2$.

\textbullet~De même, il existe une constante $a$ telle que $P_3 = aX^2(X-1)=aX^3-aX^2$ et $1 = P_3'(1) = 3a - 2a$ fournit $P_3 = X^2(X-1)$.

\textbullet~$P_0$ admet $1$ pour racine double et donc il existe deux constantes $a$ et $b$ telles que $P_0 =(aX+b)(X-1)^2$ puis les égalités $P_0(0) =1$ et $P_0'(0) = 0$ fournissent 
$b = 1$ et $a - 2b = 0$. Par suite, $P_0 =(2X+1)(X-1)^2$.

\textbullet~$P_1$ admet $0$ pour racine double et il existe deux constantes $a$ et $b$ telles que 
$P_1 = (aX+b)X^2$ puis les égalités $P_1(1) = 1$ et $P_1'(1) = 0$ fournissent $a+b = 1$ et $3a+2b = 0$ et donc $P_1 = (-2X + 3)X^2$.

\begin{center}
\shadowbox{
$P_0=(2X+1)(X-1)^2$, $P_1=(-2X+3)X^2$, $P_2=X(X-1)^2$ et $P_3=X^2(X-1)$.
}
\end{center}

Montrons alors que $(\varphi_0,\varphi_1,\varphi_2,\varphi_3)$ est une base de $E^*$.
Cette famille est libre car si  $a\varphi_1+ b\varphi_2 + c\psi_1 + d\psi_2 = 0$, on obtient en appliquant successivement à $P_0$, $P_1$, $P_2$ et $P_3$, $a=b=c=d=0$. Mais alors, la famille   $(\varphi_1,\varphi_2,\psi_1,\psi_2)$ est une famille libre de $E^*$ de cardinal $4$et donc une base de $E^*$. Sa préduale est $(P_0,P_1,P_2,P_3)$.
\fincorrection
\correction{005631}
\textbf{1 ère solution.} On utilise le fait qu'une réunion de deux sous-espaces vectoriels est un sous-espace vectoriel si et seulement si l'un des deux contient l'autre. Donc

\begin{center}
$\varphi\psi= 0\Rightarrow\text{Ker}\varphi\cup\text{Ker}\psi = E\Rightarrow\text{Ker}\psi \subset\text{Ker}\varphi=\text{Ker}\varphi\cup\text{Ker}\psi= E\;\text{ou}\;\text{Ker}\varphi\subset\text{Ker}\psi=\text{Ker}\varphi\cup\text{Ker}\psi = E\Rightarrow \varphi= 0\;\text{ou}\;\psi = 0$.
\end{center}

\textbf{2ème solution.} Supposons que $\varphi\psi= 0$ et qu'il existe $x$ et $y$ tels que $\varphi(x)\neq 0$ (et donc $\psi(x) = 0$) et $\psi(y)\neq 0$ (et donc $\varphi(y) = 0$). Alors $0 = \varphi(x+y)\psi(x+y)=(\varphi(x)+\varphi(y))(\psi(x)+\psi(y))=\varphi(x)\psi(y)$ ce qui est une contradiction.

\begin{center}
\shadowbox{
$\forall(\varphi,\psi)\in(E^*)^2,\;(\forall x\in E,\;\varphi(x)\psi(x)=0)\Rightarrow\varphi=0\;\text{ou}\;\psi=0$.
}
\end{center}
\fincorrection
\correction{005632}
\begin{enumerate}
 \item  Soit $\varphi\in E^*$. 

\textbullet~$\Rightarrow/$ Supposons qu'il existe $(\lambda_1,\ldots,\lambda_n)\in\Kk^n$ tel que $\varphi=\lambda_1\varphi_1+...+\lambda_n\varphi_n$.

Soit $x\in\displaystyle\bigcap_{i=1}^{n}\text{Ker}\varphi_i$. Alors $\varphi(x)=\lambda_1\varphi_1(x)+\ldots+\lambda_n\varphi_n(x)=0+\ldots+0=0$ et donc $x\in\text{Ker}\varphi$. On a montré que 
$\displaystyle\bigcap_{i=1}^{n}
\text{Ker}\varphi_i\subset\text{Ker}\varphi
$.

\textbullet~$\Leftarrow/$ Supposons tout d'abord la famille $(\varphi_1,\ldots,\varphi_n)$ libre. On complète éventuellement la famille libre $(\varphi_1,\ldots,\varphi_n)$ de $E^*$ en une base $(\varphi_1,\ldots,\varphi_n,\varphi_{n+1},\ldots,\varphi_p)$ de $E^*$ et on note $(e_1,\ldots,e_n,e_{n+1},\ldots,e_p)$ la préduale de la base $(\varphi_1,\ldots,\varphi_p)$.

Soit $x=\sum_{i=1}^{p}x_ie_i$ un élément de $E$.

\begin{center}
$x\in\displaystyle\bigcap_{i=1}^{n}\text{Ker}\varphi_i\Leftrightarrow\forall i\in\llbracket1,n\rrbracket,\;\varphi_i(x)=0\Leftrightarrow\forall i\in\llbracket1,n\rrbracket,\;x_i=0\Leftrightarrow x\in\text{Vect}(e_{n+1},\ldots,e_p)$
\end{center}

(avec la convention usuelle $\text{Vect}(\varnothing)=\{0\}$ dans le cas $p=n$). Donc $\displaystyle\bigcap_{i=1}^{n}\text{Ker}\varphi_i=\text{Vect}(e_{n+1},\ldots,e_p)$.

Soit alors $\varphi\in E^*$. Posons $\varphi=\sum_{i=1}^{p}\lambda_i\varphi_i$.

\begin{align*}\ensuremath
\displaystyle\bigcap_{i=1}^{n}\text{Ker}\varphi_i\subset\text{Ker}\varphi&\Rightarrow\text{Vect}(e_{n+1},\ldots,e_p)\subset\text{Ker}\varphi\Rightarrow\forall j\in\llbracket n+1,p\rrbracket,\;\varphi(e_j)=0\\
 &\Rightarrow\forall j\in\llbracket n+1,p\rrbracket,\;\lambda_j=0\Rightarrow\varphi=\sum_{i=1}^{n}\lambda_i\varphi_i.
\end{align*}

Le résultat est donc démontré dans le cas où la famille $(\varphi_1,\ldots,\varphi_n)$ est libre.

Si tous les $\varphi_i$, $1\leqslant i\leqslant n$, sont nuls alors $\displaystyle\bigcap_{i=1}^{n}\text{Ker}\varphi_i=E$ puis $\text{Ker}\varphi=E$ et donc $\varphi=0$. Dans ce cas aussi, $\varphi$ est combinaison linéaire des $\varphi_i$, $1\leqslant i\leqslant n$.

Si les $\varphi_i$, $1\leqslant i\leqslant n$, ne sont pas tous nuls et si la famille $(\varphi_1,\ldots,\varphi_n)$ est liée, on extrait de la famille $(\varphi_1,\ldots,\varphi_n)$ génératrice de $\text{Vect}(\varphi_1,\ldots,\varphi_n)$ une base $\left(\varphi_{i_1},\ldots,\varphi_{i_m}\right)$ de $\text{Vect}(\varphi_1,\ldots,\varphi_n)$.

On a $\displaystyle\bigcap_{i=1}^{n}\text{Ker}\varphi_i\subset\displaystyle\bigcap_{k=1}^{m}\text{Ker}\varphi_{i_k}$ mais d'autre part, tout $\varphi_i$, $1\leqslant i\leqslant n$, étant combinaison linéaire des $\varphi_{i_k}$, $1\leqslant k\leqslant m$, chaque $\text{Ker}\varphi_i$, $1\leqslant i\leqslant n$, contient $\displaystyle\bigcap_{k=1}^{m}\text{Ker}\varphi_{i_k}$ et donc $\displaystyle\bigcap_{k=1}^{m}\text{Ker}\varphi_{i_k}\subset\displaystyle\bigcap_{i=1}^{n}\text{Ker}\varphi_i$. Finalement, $\displaystyle\bigcap_{k=1}^{m}\text{Ker}\varphi_{i_k}=\displaystyle\bigcap_{i=1}^{n}\text{Ker}\varphi_i\subset\text{Ker}\varphi$. D'après l'étude du cas où la famille est libre, $\varphi$ est combinaison linéaire des $\varphi_{i_k}$, $1\leqslant k\leqslant m$ et donc des $\varphi_i$, $1\leqslant i\leqslant n$. La réciproque est démontrée dans tous les cas.

\item  Soit $\varphi$ une forme linéaire sur $\Rr^3$ telle que $P=\text{Ker}\varphi$ (en particulier $\varphi$ n'est pas nulle).
Soient $\varphi_1$ la forme linéaire $(x,y,z)\mapsto x+y+z$ et $\varphi_2$ la forme linéaire $(x,y,z)\mapsto2x+3z$. Alors la famille $(\varphi_1,\varphi_2)$ est une famille libre du dual de $\Rr^3$ et $D=\text{Ker}\varphi_1\cap\text{Ker}\varphi_2$. D'après 1)

\begin{center}
$D\subset P\Leftrightarrow\exists(a,b)\in\Rr^2\setminus\{(0,0)\}/\;\varphi= a\varphi_1+b\varphi_2$ (théorie des faisceaux),
\end{center}

puis

\begin{center}
$u\in P\Leftrightarrow a\varphi_1(u)+b\varphi_2(u) = 0\Leftrightarrow 3a + 5b = 0$.
\end{center}

Une équation de $P$ est donc $5(x+y+z)-3(2x+3z) = 0$ ou encore $-x + 5y -4z = 0$.
\end{enumerate}
\fincorrection
\correction{005633}
Soit $\begin{array}[t]{cccc}f~:&E&\rightarrow&\Kk^n\\
 &x&\mapsto&(\varphi_1(x),...,\varphi_n(x))
\end{array}$. Il s'agit de démontrer que la famille $(\varphi_1,\ldots,\varphi_n)$ est liée si et seulement si $\text{Ker}(f)\neq\{0\}$.

\textbullet~Si la famille $(\varphi_1,\ldots,\varphi_n)$ est libre, c'est une base de $E^*$ (car $\text{dim}(E^*)=n$). Notons $(u_1,\ldots,u_n)$ sa préduale et notons $(e_1,\ldots,e_n)$ la base canonique de $\Kk^n$. Pour $1\leqslant i\leqslant n$, on a $f(u_i)=e_i$. Ainsi, l'image par $f$ d'une base de $E$ est une base de $\Kk^n$ et on sait alors que $f$ est un isomorphisme. En particulier, $\text{Ker}(f)=\{0\}$.

\textbullet~Si les $\varphi_i$ sont tous nuls, tout vecteur non nul $x$ annule chaque $\varphi_i$. Supposons alors que la famille $(\varphi_1,\ldots,\varphi_n)$ est liée et que les $\varphi_i$ ne sont pas tous nuls. On extrait de la famille $(\varphi_1,\ldots,\varphi_n)$ une base $\left(\varphi_{i_1},\ldots,\varphi_{i_m}\right)$ (avec $1\leqslant m<n$) de $\text{Vect}(\varphi_1,\ldots,\varphi_n)$. On complète la famille libre $\left(\varphi_{i_1},\ldots,\varphi_{i_m}\right)$ en une base $\left(\varphi_{i_1},\ldots,\varphi_{i_m},\psi_1,\ldots,\psi_{n-m}\right)$ de $E^*$. On note $(e_1,\ldots,e_m,e_{m+1},\ldots,e_n)$ sa préduale. Les formes linéaires $\varphi_{i_1}$,\ldots, $\varphi_{i_m}$ s'annulent toutes en $e_n$ et donc chaque $\varphi_i$ s'annule en $e_n$ puisque chaque $\varphi_i$ est combinaison linéaire des $\varphi_{i_k}$, $1\leqslant i\leqslant m$. Le vecteur $e_n$ est donc un vecteur non nul $x$ tel que $\forall i\in\llbracket1,n\rrbracket$, $\varphi_i(x)=0$.
\fincorrection
\correction{005634}
La matrice de la famille $(f_1,f_2,f_3,f_4)$ dans la base canonique du dual de $\Rr^4$ est $A =\left(
\begin{array}{cccc}
1&1&1&0\\
2&1&0&1\\
-1&m&1&-3\\
-2&1&m+4&-m
\end{array}
\right)$. La matrice $A$ a même rang que la matrice $\left(
\begin{array}{cccc}
1&0&0&0\\
2&-1&-2&1\\
-1&m+1&2&-3\\
-2&3&m+6&-m
\end{array}
\right)$ (pour $2\leqslant j\leqslant 3$, $C_j\leftarrow C_j-C_1$) puis que la matrice $\left(
\begin{array}{cccc}
1&0&0&0\\
2&-1&0&0\\
-1&m+1&-2m&m-2\\
-2&3&m&-m+3
\end{array}
\right)
$ ($C_3\leftarrow C_3-2C_2$ et $C_4\leftarrow C_4+C_2$)

\textbullet~Si $m = 0$, $A$ a même rang que la matrice $\left(
\begin{array}{cccc}
1&0&0\\
2&-1&0\\
-1&1&-2\\
-2&3&3
\end{array}
\right)
$et donc $\text{rg}(A)=3$.

\textbullet~Si $m\neq0$, $A$ a même rang que la matrice $\left(
\begin{array}{cccc}
1&0&0&0\\
2&-1&0&0\\
-1&m+1&-2&m-2\\
-2&3&1&-m+3
\end{array}
\right)
$ ($C_3\leftarrow\frac{1}{m}C_3$) puis que la matrice $\left(
\begin{array}{cccc}
1&0&0&0\\
2&-1&0&0\\
-1&m+1&-2&0\\
-2&3&1&-m+4
\end{array}
\right)
$ ($C_4\leftarrow2C_4+(m-2)C_3$)

Donc, si $m =4$, $\text{rg}(A)= 3$ et si $m$ n'est ni $0$ ni $4$, $\text{rg}(A)=4$.

\begin{center}
\shadowbox{
Si $m\notin\{0,4\}$, $\text{rg}(f_1,f_2,f_3,f_4)=4$ et si $m\in\{0,4\}$, $\text{rg}(f_1,f_2,f_3,f_4)=3$.
}
\end{center}
\fincorrection
\correction{005773}
\begin{enumerate}
 \item  \textbullet~Soient $P$ et $Q$ deux polynômes. La fonction $t\mapsto P(t)Q(t)e^{-t}$ est continue sur $[0,+\infty[$ et est négligeable en $+\infty$ devant $\frac{1}{t^2}$ d'après un théorème de croissances comparées. Donc la fonction $t\mapsto P(t)Q(t)e^{-t}$ est intégrable sur $[0,+\infty[$ et $\varphi(P,Q)$ existe dans $\Rr$.

\textbullet~La symétrie, la bilinéarité et la positivité de l'application $\varphi$ sont claires. De plus, pour $P\in E$,

\begin{align*}\ensuremath
\varphi(P,P)=0&\Rightarrow\int_{0}^{+\infty}P^2(t)e^{-t}\;dt=0\\
 &\Rightarrow\forall t\in[0,+\infty[,\;P^2(t)e^{-t}=0\;(\text{fonction continue positive d'intégrale nulle})\\
 &\Rightarrow\forall t\in[0,+\infty[,\;P(t)=0\Rightarrow P=0\;(\text{polynôme ayant une infinité de racines}).
\end{align*}

Ainsi, la forme $\varphi$ est définie et finalement

\begin{center}
\shadowbox{
l'application $\varphi$ est un produit scalaire sur $E$.
}
\end{center}

\item 
  \begin{enumerate}
  \item Soit $n\in\Nn$. La formule de \textsc{Leibniz} permet d'écrire

\begin{center}
$(X^ne^{-X})^{(n)}e^X=\left(\sum_{k=0}^{n}\dbinom{n}{k}(X^n)^{(n-k)}(e^{-X})^{(k)}\right)e^X=\sum_{k=0}^{n}(-1)^{k}\dbinom{n}{k}\frac{n!}{k!}X^k$.
\end{center}

En particulier, $\forall n\in\Nn$, $\text{deg}(h_n)=n$ (et $\text{dom}(h_n)=(-1)^n$) et on sait que

\begin{center}
\shadowbox{
la famille $(h_n)_{n\in\Nn}$ est une base de $\Rr[X]$.
}
\end{center}

  \item Soient $P\in E$ et $n\in\Nn^*$. Soit $A>0$. Les deux fonctions $t\mapsto(t^ne^{-t})^{(n-1)}$ et $P$ sont de classe $C^1$ sur le segment $[0,A]$. On peut donc effectuer une intégration par parties et on obtient

\begin{center}
$\int_{0}^{A}P(t)h_n(t)e^{-t}\;dt=\int_{0}^{A}P(t)(t^ne^{-t})^{(n)}\;dt=\left[P(t)(t^ne^{-t})^{(n-1)}\right]_0^A-\int_{0}^{A}P'(t)(t^ne^{-t})^{(n-1)}\;dt$
\end{center}

Maintenant, $(t^ne^{-t})^{(n-1)}$ peut s'écrire $Q(t)e^{-t}$ où $Q$ est un polynôme et donc $P(t)(t^ne^{-t})^{(n-1)}(t)$ tend vers $0$ quand $t$ tend vers $+\infty$ d'après un théorème de croissances comparées. D'autre part, la formule de \textsc{Leibniz} montre que le polynôme $Q$ a
une valuation au moins égale à $1$. On en déduit que la fonction $t\mapsto P(t)(t^ne^{-t})^{(n-1)}(t)$ s'annule en $0$. En faisant tendre $A$ vers $+\infty$, on obtient

\begin{center}
$\int_{0}^{+\infty}P(t)h_n(t)e^{-t}\;dt=-\int_{0}^{+\infty}P'(t)(t^ne^{-t})^{(n-1)}\;dt$.
\end{center}

De manière générale, pour $0\leqslant k\leqslant n$, les remarques précédentes s'appliquent à la fonction $t\mapsto P^{(k)}(t)(t^ne^{-t})^{(n-k)}$ et par récurrence on obtient

\begin{center}
$\forall k\in\llbracket0,n\rrbracket$, $\int_{0}^{+\infty}P(t)h_n(t)e^{-t}\;dt=(-1)^k\int_{0}^{+\infty}P^{(k)}(t)(t^ne^{-t})^{(n-k)}\;dt$.
\end{center}

En particulier, pour $k=n$ on obtient $\int_{0}^{+\infty}P(t)h_n(t)e^{-t}\;dt=(-1)^n\int_{0}^{+\infty}P^{(n)}(t)t^ne^{-t}\;dt$. Cette égalité reste vraie quand $n=0$ et on a montré que

\begin{center}
$\forall P\in\Rr[X]$, $\forall n\in\Nn$, $\varphi(P,h_n)=\int_{0}^{+\infty}P(t)h_n(t)e^{-t}\;dt=(-1)^n\int_{0}^{+\infty}P^{(n)}(t)t^ne^{-t}\;dt$.
\end{center}

En particulier, si $n\in\Nn^*$ et $\text{deg}(P)<n$, on a $P^{(n)}=0$ et donc $\varphi(P,h_n)=0$. Ainsi, $\forall n\in\Nn^*$, $h_n\in(\Rr_{n-1}[X])^\bot$. Puisque $\forall n\in\Nn$, $\text{deg}(h_n)=n$, on en déduit en particulier que $\forall n\in\Nn^*$, $\forall k\in\llbracket0,n-1\rrbracket$, $\varphi(hn,h_k)=0$ et on a montré que

\begin{center}
\shadowbox{
la famille $(h_n)_{n\in\Nn}$ est une base orthogonale de l'espace préhilbertien $(\Rr[X],\varphi)$.
}
\end{center}

  \item  Soit $n\in\Nn$. Puisque $\text{deg}(h_n)=n$ et $\text{dom}(h_n)=(-1)^n$, on a $h_n^{(n)}=(-1)^nn!$. La question précédente fournit alors

\begin{center}
$\|h_n\|^2=(-1)^n\int_{0}^{+\infty}h_n^{(n)}(t)t^ne^{-t}\;dt=n!\int_{0}^{+\infty}t^ne^{-t}\;dt=n!\Gamma(n+1)=n!^2$,
\end{center}

et donc $\|h_n\|=n!$. Par suite,

\begin{center}
\shadowbox{
la famille $\left(\frac{1}{n!}h_n\right)_{n\in\Nn}$ est une base orthonormale de l'espace préhilbertien $(\Rr[X],\varphi)$.
}
\end{center}
  \end{enumerate}
\end{enumerate}
\fincorrection
\correction{005774}
\begin{enumerate}
 \item  \textbullet~Soit $(P,Q)\in E^2$. L'application $t\mapsto\frac{P(t)Q(t)}{\sqrt{1-t^2}}$ est continue sur $]-1,1[$. Ensuite, l'application $t\mapsto\frac{P(t)Q(t)}{\sqrt{1+t}}$ est bornée au voisinage de $1$ car continue en $1$ et donc quand $t$ tend vers $1$, $\frac{P(t)Q(t)}{\sqrt{1-t^2}}=\frac{P(t)Q(t)}{\sqrt{1+t}}\times\frac{1}{\sqrt{1-t}}=O\left(\frac{1}{\sqrt{1-t}}\right)$. Puisque $\frac{1}{2}<1$, on en déduit que l'application $t\mapsto\frac{P(t)Q(t)}{\sqrt{1-t^2}}$ est intégrable sur un voisinage de $1$ à gauche. De même, quand $t$ tend vers $1$, $\frac{P(t)Q(t)}{\sqrt{1-t^2}}=O\left(\frac{1}{\sqrt{1+t}}\right)$ et l'application $t\mapsto\frac{P(t)Q(t)}{\sqrt{1-t^2}}$ est intégrable sur un voisinage de $-1$ à droite. Finalement, l'application $t\mapsto\frac{P(t)Q(t)}{\sqrt{1-t^2}}$ est intégrable sur $]-1,1[$ et $\varphi(P,Q)$ existe.

\textbullet~La symétrie, la bilinéarité et la positivité de $\varphi$ sont claires. De plus, pour $P\in E$,

\begin{align*}\ensuremath
\varphi(P,P)=0&\Rightarrow\int_{-1}^{1}\frac{P^2(t)}{\sqrt{1-t^2}}\;dt=0\\
 &\Rightarrow\forall t\in]-1,1[,\;\frac{P^2(t)}{\sqrt{1-t^2}}=0\;(\text{fonction continue, positive, d'intégrale nulle})\\
 &\Rightarrow\forall t\in]-1,1[,P(t)=0\Rightarrow P=0\;(\text{polynôme ayant une infinité de racines}).
\end{align*}

Ainsi, l'application $\varphi$ est définie et finalement

\begin{center}
\shadowbox{
l'application $\varphi$ est un produit scalaire sur $E$.
}
\end{center}

\item 
  \begin{enumerate}
  \item Soit $(n,p)\in\Nn^2$. En posant $t=\cos\theta$, on obtient

\begin{center}
$\varphi(T_n,T_p)=\int_{-1}^{1}\frac{T_n(t)T_p(t)}{\sqrt{1-t^2}}\;dt=\int_{\pi}^{0}\frac{T_n(\cos\theta)T_p(\cos\theta)}{\sqrt{1-\cos^2\theta}}\;(-\sin\theta d\theta)=\int_{0}^{\pi}\cos(n\theta)\cos(p\theta)\;d\theta$.
\end{center}

Si de plus, $n\neq p$, 

\begin{center}
$\varphi(T_n,T_p)=\frac{1}{2}\int_{0}^{\pi}(\cos((n+p)\theta)+\cos((n-p)\theta))\;d\theta=\frac{1}{2}\left[\frac{\sin((n+p)\theta)}{n+p}+\frac{\sin((n-p)\theta)}{n-p}\right]_0^\pi=0$.
\end{center}

Ainsi, la famille $(T_n)_{n\in\Nn}$ est orthogonale. De plus, on sait que $\forall n\in\Nn$, $\text{deg}(T_n)=n$ et on a donc montré que

\begin{center}
\shadowbox{
la famille $(T_n)_{n\in\Nn}$ est une base orthogonale de l'espace préhilbertien $(E,\varphi)$.
}
\end{center}

  \item Soit $n\in\Nn$. Quand $p=n$, la formule précédente fournit

\begin{center}
$\|T_n\|^2=\frac{1}{2}\int_{0}^{\pi}(1+\cos(2n\theta))\;d\theta=\left\{
\begin{array}{l}
\pi\;\text{si}\;n=0\\
\frac{\pi}{2}\;\text{si}\;n\geqslant1
\end{array}
\right.$,
\end{center}

et donc

\begin{center}
\shadowbox{
$\forall n\in\Nn$, $\|T_n\|=\left\{
\begin{array}{l}
\rule[-2mm]{0mm}{0mm}\sqrt{\pi}\;\text{si}\;n=0\\
\sqrt{\frac{\pi}{2}}\;\text{si}\;n\geqslant1
\end{array}
\right.$.
}
\end{center}
  \end{enumerate}
\end{enumerate}
\fincorrection
\correction{005775}
\begin{enumerate}
 \item  Montrons que $E$ est un sous-espace de $(\Rr^\Nn,+,.)$. La suite nulle est élément de $E$. Soient $(u,v)\in E^2$ et $(\lambda,\mu)\in\Rr^2$.

\begin{center}
$0\leqslant(\lambda u+\mu v)^2=\lambda^2u^2+2\lambda\mu uv+\mu^2v^2\leqslant\lambda^2u^2+\lambda\mu(u^2+v^2)+\mu^2v^2=(\lambda^2+\lambda\mu)u^2+(\lambda\mu+\mu^2)v^2$.
\end{center}

Par hypothèse, la série de terme général $(\lambda^2+\lambda\mu)u_n^2+(\lambda\mu+\mu^2)v_n^2$ converge et on en déduit que la suite $\lambda u+\mu v$ est de carré sommable. On a montré que

\begin{center}
\shadowbox{
$E$ est un sous-espace vectoriel de l'espace vectoriel $(\Rr^\Nn,+,.)$.
}
\end{center}

\item  \textbullet~Soient $u$ et $v$ deux éléments de $E$. Pour tout entier naturel $n$,

\begin{center}
$|u_nv_n\leqslant\frac{1}{2}(u_n^2+v_n^2)$.
\end{center}

Ainsi, la série de terme général $u_nv_n$ est absolument convergente et donc convergente. Ceci montre que $\varphi(u,v)$ existe dans $\Rr$.

\textbullet~La symétrie, la bilinéarité et la positivité de $\varphi$ sont claires. De plus, pour $u\in E$,

\begin{center}
$\varphi(u,u)=0\Rightarrow\sum_{n=0}^{+\infty}u_n^2=0\Rightarrow\forall n\in\Nn,\;u_n^2=0\Rightarrow u=0$.
\end{center}

En résumé, l'application $\varphi$ est une forme bilinéaire, symétrique, définie, positive et donc

\begin{center}
\shadowbox{
l'application $\varphi$ est un produit scalaire sur $E$.
}
\end{center}
\end{enumerate}
\fincorrection
\correction{005776}
Soit $(A,B)\in(\mathcal{M}_n(\Rr))^2$. 

\begin{center}
$\Phi(A,B)=\text{Tr}({^t}A\times B) =\sum_{1\leqslant i,j\leqslant n}^{}a_{i,j}b_{i,j}$.
\end{center}

L'application $\Phi$ n'est autre que produit scalaire canonique de $\mathcal{M}_n(\Rr)$ et en particulier est un produit scalaire. La base canonique de $\mathcal{M}_n(\Rr)$ (constituée des matrices élémentaires) est orthonormée pour ce produit scalaire.

L'application $\Phi$ n'est pas un produit scalaire sur $\mathcal{M}_n(\Cc)$. Par exemple, si $A=iE_{1,1}\neq0$ alors ${^t}AA=-E_{1,1}$ puis $\text{Tr}({^t}AA)=-1<0$.
\fincorrection
\correction{005777}
Soit $N$ une norme sur $E$ vérifiant $\forall(x,y)\in E^2$ $(N(x+y))^2+(N(x-y))^2 = 2((N(x))^2+(N(y))^2)$.

Il faut montrer que la norme $N$ est associée à un produit scalaire $B$. Si $B$ existe, $B$ est nécessairement défini par

\begin{center}
$\forall(x,y)\in E^2$, $B(x,y) =\frac{1}{4}((N(x+y))^2-(N(x-y))^2)$.
\end{center}

Réciproquement, 

\textbullet~Pour tout $x\in E$, $B(x,x) =\frac{1}{4}((N(2x))^2-(N(0))^2)=\frac{1}{4}(4(N(x))^2-0) =(N(x))2$ et donc $\forall x\in E$, $B(x,x)\geqslant0$ puis $B(x,x)=0\Leftrightarrow x=0$. De plus, $\forall x\in E$, $N(x)=\sqrt{B(x,x)}$.

\textbullet~$\forall (x,y)\in E^2$, $B(y,x)=\frac{1}{4}((N(y+x))^2-(N(y-x))^2)=\frac{1}{4}((N(x+y))^2-(N(x-y))^2)=B(x,y)$.

\textbullet~Vérifions alors que l'application $B$ est bilinéaire.

1) Montrons que $\forall(x,y,z)\in E^3$, $B(x+y,z)+B(x-y,z)=2B(x,z)$.

\begin{align*}\ensuremath
B(x+y,z) + B(x-y,z)&=\frac{1}{4}((N(x+y+z))^2-(N(x+y-z))^2 +(N(x-y+z))^2-(N(x-y-z))^2)\\
 &=\frac{1}{4}((N(x+y+z))^2+(N(x-y+z))^2)-((N(x+y-z))^2+(N(x-y-z))^2)\\
 &=\frac{1}{4}(2(N(x+z))^2 +(N(y))^2)-2((N(x-z))^2 +(N(y))^2)\;(\text{par hypothèse sur}\;N)\\
 &=\frac{2}{4}((N(x+z))^2-(N(x-z))^2) =2B(x,z).
\end{align*}

2) Montrons que $\forall(x,z)\in E^2$, $B(2x,z) = 2B(x,z)$. Tout d'abord, $B(0,z)=\frac{1}{4}((N(z))^2-(N(-z))^2)=0$ puis d'après 1)

\begin{align*}\ensuremath
B(2x,z)&= B(x+x,z)+B(x-x,z) = 2B(x,z).
\end{align*}

3) Montrons que $\forall(x,y,z)\in E^3$, $B(x,z) + B(y,z) = B(x+y,z)$.

\begin{align*}\ensuremath
B(x,z) + B(y,z)&= B\left(\frac{x+y}{2}+\frac{x-y}{2},z\right)+ B\left(\frac{x+y}{2}-\frac{x-y}{2},z\right)\\
 &=2B\left(\frac{x+y}{2},z\right)\;(\text{d'après 1)})\\
 &= B(x+y,z)\;(\text{d'après 2)}).
\end{align*}

4) Montrons que $\forall n\in\Nn$, $\forall(x,y)\in E^2$, $B(nx,y) = nB(x,y)$.

\textbullet~~C'est vrai pour $n = 0$ et $n=1$.

\textbullet~Soit $n\geqslant0$. Supposons que $\forall(x,y)\in E^2$, $B(nx,y) = nB(x,y)$ et $B((n+1)x,y) = (n+1)B(x,y)$. Alors

\begin{center}
$B((n+2)x,y) + B(nx,y) = B((n+2)x+nx,y) = B(2(n+1)x,y) = 2B((n+1)x,y)$,
\end{center}

et donc, par hypothèse de récurrence, $B((n+2)x,y)=2(n+1)B(x,y)-nB(x,y) =(n+2)B(x,y)$.

5) Montrons que $\forall n\in\Zz$, $\forall(x,y)\in E^2$, $B(nx,y) = nB(x,y)$. Le résultat est acquis pour $n\geqslant 0$. Pour $n\in\Nn$,

\begin{center}
$B(nx,y) + B(-nx,y) = B(0,y) = 0$ et donc $B(-nx,y) =- B(nx,y)=-nB(x,y)$,
\end{center}

6) Montrons que $\forall n\in\Nn^*$, $\forall(x,y)\in E^2$, $B\left(\frac{1}{n}x,y\right) =\frac{1}{n}B(x,y)$.

\begin{center}
$B(x,y) =B\left(\frac{1}{n}nx,y\right) =nB\left(\frac{1}{n}x,y\right)$ et donc $B\left(\frac{1}{n}x,y\right) =\frac{1}{n}B(x,y)$.
\end{center}

7) Montrons que $\forall r\in\Qq$, $\forall(x,y)\in E^2$, $B(rx,y) = rB(x,y)$. Soient $(p,q)\in\Zz\times\Nn^*$ puis $r=\frac{p}{q}$.

\begin{center}
$B(rx,y) = B\left(\frac{p}{q}x,y\right) =pB\left(\frac{1}{q}x,y\right)=\frac{p}{q}B(x,y) = rB(x,y)$.
\end{center}

8) Montrons que $\forall\lambda\in\Rr$, $\forall(x,y)\in E^2$, $B(\lambda x,y) =\lambda B(x,y)$. Soit $\lambda$ un réel. Puisque $\Qq$ est dense dans $\Rr$, il existe une

suite de rationnels $(r_n)_{n\in\Nn}$ convergente de limite $\lambda$.

Maintenant, l'application $\begin{array}[t]{cccc}N~:&(E,N)&\rightarrow&(\Rr,|\;|)\\
 &x&\mapsto&N(x)
 \end{array}$ est continue sur $E$ car $1$-Lipschitzienne sur $E$. Donc
 
\begin{center}
$B(\lambda x,y) =B(\lim_{n \rightarrow +\infty}r_nx,y)=\lim_{n \rightarrow +\infty}B(r_nx,y)=\lim_{n \rightarrow +\infty}r_nB(x,y)=\lambda B(x,y)$.
\end{center}

Finalement, l'application $B$ est une forme bilinéaire symétrique définie positive et donc un produit scalaire. Puisque $\forall x\in E$, $N(x)=\sqrt{B(x,x)}$, $N$ est la norme associée à ce produit scalaire. On a montré que

\begin{center}
\shadowbox{
toute norme vérifiant l'identité du parallélogramme est une norme hilbertienne.
}
\end{center}
\fincorrection
\correction{005778}
Soit $i\in\llbracket1,n\rrbracket$.

\begin{center}
$1=\|e_i\|^2=\sum_{j\ge1}(e_i|e_j)^2 =1 +\sum_{j\neq i}^{}(e_i|e_j)^2$
\end{center}

et donc $\sum_{j\neq i}^{}(e_i|e_j)^2= 0$.On en déduit que $\forall j\neq i$, $(e_i|e_j)=0$. Ainsi, pour tout couple d'indices $(i,j)$ tel que $i\neq j$, on a  $e_i|e_j = 0$. Par suite

\begin{center}
la famille $(e_i)_{1\leqslant i\leqslant n}$ est une famille orthonormale.
\end{center}

Il reste à vérifier que si $F =\text{Vect}(e_1,...,e_n)$ alors $F = E$.

Soit $x$ un vecteur de $E$. $F$ est un sous-espace vectoriel de $E$ de dimension finie. On peut donc définir le projeté orthogonal $p_F(x)$ de $x$ sur $F$. On sait que

\begin{center}
$p_F(x) =\sum_{i=1}^{n}(x|e_i)e_i$.
\end{center}

On en déduit que $\|p_F(x)\|^2 =\sum_{i=1}^{n}(x|e_i)^2 =\|x\|^2$. D'après le théorème de \textsc{Pythagore},

\begin{center}
$\|x-p_F(x)\|^2 =\|x\|2 -\|p_F(x)\|^2 = 0$,
\end{center}

et donc $x = p_F(x)$ ce qui montre que $x\in F$. Donc $F = E$ et finalement

\begin{center}
\shadowbox{
la famille $(e_i)_{1\leqslant i\leqslant n}$ est une base orthonormée de $E$.
}
\end{center}

\fincorrection
\correction{005788}
\textbf{1 ère solution.} Soit $p\geqslant2$. Montrons que si la famille $(x_1,...,x_{p})$ est obtusangle alors la famille $(x_1,...,x_{p-1})$ est libre. 

Soit $(x_1,...,x_{p})$ une famille obtusangle. Supposons que la famille $(x_1,...,x_{p-1})$ soit liée.

Il existe donc $(\lambda_1,...,\lambda_{p-1})\in\Rr^{p-1}\setminus\{(0,\ldots,0)\}$ tel que $\sum_{k=1}^{p-1}\lambda_kx_k = 0$.

Quite à multiplier les deux membres de l'égalité par $-1$, on peut supposer que l'un des $\lambda_i$ au moins est strictement positif. On pose $I =\{k\in\llbracket1,p-1\rrbracket/\;\lambda_k > 0\}$ et $J =\{k\in\llbracket1,p-1\rrbracket/\;\lambda_k\leqslant 0\}$ (éventuellement $J$ est vide).

Si $J$ est vide, il reste $\sum_{i\in I}^{}\lambda_ix_i = 0$ et si $J$ est non vide,

\begin{center}
$\left\|\sum_{i\in I}^{}\lambda_ix_i\right\|^2= -\left(\sum_{i\in I}^{}\lambda_ix_i\right)\left(\sum_{j\in J}^{}\lambda_jx_j\right)= -\sum_{(i,j)\in I\times J}^{}\lambda_i\lambda_j\left(x_i|x_j\right)\leqslant0$ (car  $\forall(i,j)\in I\times J$, $\left(x_i|x_j\right)< 0$ et $\lambda_i\lambda_j\leqslant0$).
\end{center}

Ainsi, dans tous les cas, $\sum_{i\in I}^{}\lambda_ix_i = 0$. Mais ceci est impossible car  $\left(\sum_{i\in I}^{}\lambda_ix_i\right)|x_{p}=\sum_{i\in I}^{}\lambda_i\left(x_i|x_{p}\right)< 0$.

On a montré que la famille $(x_1,\ldots,x_{p-1})$ est libre et on en déduit que $p-1\leqslant n$ ou encore $p\leqslant n+1$.

\textbf{2ème solution.} Montrons par récurrence sur $n =\text{dim}E_n\geqslant1$ que tout famille obtusangle de $E_n$ a un cardinal inférieur ou égal à $n+1$.

\textbullet~Pour $n = 1$. Soient $x_1$, $x_2$ et $x_3$ trois vecteurs de $E_1$. On peut identifier ces vecteurs à des réels. Deux des trois réels $x_1$, $x_2$ ou $x_3$ ont même signe et on ne peut donc avoir $x_1x_2 < 0$ et $x_1x_3 < 0$ et $x_2x_3 < 0$.

Une famille obtusangle de $E_1$ a donc un cardinal inférieur ou égal à $2$.

\textbullet~Soit $n\geqslant1$. Supposons que toute famille obtusangle d'un espace euclidien de dimension $n$ a un cardinal inférieur ou égal à $n+1$. Soit $(x_1,...,x_p)$ une famille obtusangle de $E_{n+1}$.

Si $p = 1$ alors $p\leqslant n+2$. Supposons dorénavant $p\geqslant2$.

On va construire à partir de cette famille une famille obtusangle de cardinal $p-1$ d'un espace euclidien de dimension $n$.

Soit $F=x_p^\bot$. Puisque la famille $(x_1,...,x_p)$ est obtusangle, le vecteur $x_p$ n'est pas nul et $F$ est un espace euclidien de dimension $n$.

On note $y_1$, $y_2$,..., $y_{p-1}$ les projetés orthogonaux des vecteurs $x_1$, ... , $x_{p-1}$ sur $F$. On sait que

\begin{center}
$\forall i\in\llbracket1,p-1\rrbracket$, $y_i =x_i -\frac{\left(x_i|x_p\right)}{\|x_p\|^2}x_p$.
\end{center}

Soit $(i,j)\in\llbracket1,p-1\rrbracket$ tel que $i\neq j$.

\begin{center}
$\left(y_i|y_j\right)=\left(x_i|x_j\right)-2\frac{\left(x_i|x_p\right)\left(x_j|x_p\right)}{\|x_p\|^2}+\frac{\left(x_i|x_p\right)\left(x_j|x_p\right)\|x_p\|^2}{\|x_p\|^4}= \left(x_i|x_j\right)-\frac{\left(x_i|x_p\right)\left(x_j|x_p\right)}{\|x_p\|^2}< 0$.
\end{center}

Ainsi, la famille $(y_i)_{1\leqslant i\leqslant p-1}$ est une famille obtusangle d'un espace euclidien de dimension $n$ et par hypothèse de récurrence $p-1\leqslant n+1$ et donc $p\leqslant n+2$. Le résultat est démontré par récurrence.
\fincorrection
\nocorrection
\nocorrection
\nocorrection
\nocorrection
\nocorrection
\nocorrection
\nocorrection
\nocorrection
\correction{003713}
\begin{enumerate}
  \item Oui ssi $\lambda^2 + \mu^2 \le 1$.
  \item Non, disc $=-1$.
  \item Oui, $= \frac {5x^2}{12} + 3\left(y+\frac x3\right)^2 + 4z^2
                   + \left(t+\frac x2\right)^2$.
\end{enumerate}
 \fincorrection
\correction{003714}
Spec$(A) = \{6,3,3\} \Rightarrow $ oui.
\fincorrection
\correction{003715}
$(n-1,0)$.
\fincorrection
\nocorrection
\correction{003717}
Récurrence,
$q(x_1,\dots,x_n) = \sum_{i<k} \frac{i+1}{2i}y_i^2
                  + \frac12\sum_{i\ge k} x_i^2
                  + \frac1{2k}\Bigl(\sum_{i\ge k} x_i\Bigr)^2$.
\fincorrection
\nocorrection
\nocorrection
\correction{003720}
\begin{enumerate}
  \item 
  \item 
  \item Si $a\in\mathrm{Ker} f$, $\mathrm{Ker}\tilde\varphi = E$.\par
Si $a\notin\mathrm{Ker} f$ et $q(a) = 0$, $\mathrm{Ker}\tilde\varphi = a^\bot$.\par
Si $q(a)\ne 0$, $\mathrm{Ker}\tilde\varphi = \mathrm{Ker}(f) \oplus {<a>}$.
\end{enumerate}
 \fincorrection
\nocorrection
\nocorrection
\correction{003723}
\begin{enumerate}
  \item 
  \item $(d-1,0)$ ou $(d,0)$.
\end{enumerate}
 \fincorrection
\correction{003724}
Récurrence sur~$n$.

Soit $(e_1,\dots,e_n)$ la base dans laquelle $A$ est la matrice de~$q$.
$\Delta_{n-1}(A) \ne 0$ donc il existe des coefficients
$\alpha_1,\dots,\alpha_{n-1}$ tels que $u_n = e_n-\sum_{i<n}\alpha_i e_i$
soit $q$-orthogonal à~$e_1,\dots,e_{n-1}$. Alors $A$ a mêmes
mineurs principaux que la matrice de $q$ dans la base
$(e_1,\dots,e_{n-1},u_n)$.
\fincorrection
\correction{003725}
Pour $A$ à diagonale fortement dominante, récurrence sur $n$.

Soit $(e_1,\dots,e_n)$ la base dans laquelle $A$ est la matrice de $q$.
$\Delta_{n-1}(A) \ne 0$ donc il existe des coefficients
$\alpha_1,\dots,\alpha_{n-1}$ tels que 
$u_n = e_n-\sum_{i<n}\alpha_i e_i$ soit $q$-orthogonal à $e_1,\dots,e_{n-1}$ et il faut montrer que
$q(u_n) > 0$ ce qui résulte de $|\alpha_i| \le 1$ en considérant la $i$-ème ligne
de $A$.
\fincorrection
\correction{003726}
\begin{enumerate}
  \item 
  \item 
  \item $GL_n^+(\R)$ l'est.
\end{enumerate}
 \fincorrection
\nocorrection
\correction{003728}
\begin{enumerate}
  \item Il n'y a pas de solution pour $n=2$ donc pas non plus pour $n>2$.
  \item Pour $n=3$ on trouve $A={}^tCC$ où $C$ est une colonne sans zéros,
pour $n\ge 3$ on obtient le même résultat en considérant les blocs $3\times3$
centrés sur la diagonale.
\end{enumerate}
 \fincorrection
\correction{003729}
Soit $P$ orthogonale diagonalisant~$S$~: ${}^t\!PSP=\mathrm{Diag}(\lambda_1,\dots,\lambda_n)$
et $Y={}^t\!PX$.

On a~: $q(X) = \begin{vmatrix}0&y_1&\dots&y_n\cr y_1&\lambda_1\cr\vdots&&\ddots\cr y_n&&&\lambda_n\cr\end{vmatrix}
=-\lambda_1\dots\lambda_n\Bigl(\frac{y_1^2}{\lambda_1} +\dots +\frac{y_n^2}{\lambda_n}\Bigr)$.
\fincorrection
\correction{003730}
Pour $a=0$, $q$ est définie positive.
Pour $a\ne 0$ prendre une base orthonormale commen\c cant par~$a$~;
la matrice de~$q$ dans cette base est $\mathrm{Diag}(\alpha+\beta\|a\|^2,\alpha,\dots,\alpha)$.
\fincorrection
\correction{003731}
Soit $(E_{ij})$ la base canonique de~$\mathcal{M}_2(\C)$~:
$E_{12}^2=0$ donc $q(E_{12}) = 0$ et si $A$ est une matrice quelconque
de rang~$1$, $A$ est équivalente à~$E_{12}$ d'où $q(A) = 0$.
Si $A=0$ on a aussi $q(A)=0$ et si $A$ est inversible alors toute matrice
est multiple de~$A$ donc $q(A)\ne 0$, en particulier $q(I)=1$
car $q^2(I)=q(I)$. On en déduit $q(A) = 0 \Leftrightarrow\det(A)=0$.

Pour $A$ quelconque, les applications~: $z \mapsto\det(A-zI)$ et $z \mapsto q(A-zI)$
sont polynomiales de degré~$2$, avec le même coefficient de~$z^2$
et les mêmes racines, donc sont égales d'où $q=\det$.

\emph{Remarque :} le même raisonnement est applicable sur un corps quelconque en se limitant
aux matrices triangulaires, et toute matrice est produit de triangulaires
(algorithme du pivot de Gauss).
\fincorrection
\correction{003732}
Quitte à remplacer $E$ par $\mathrm{vect}(x_1,\dots,y_n)$, on peut supposer $E$
de dimension finie $p$. Soit $\cal B$ une base de $E$, et $X,Y$ et $F$
les matrices de $(x_1,\dots,x_n)$, $(y_1,\dots,y_n)$ et $\phi$ dans $\cal B$.
On doit prouver $\det({}^tXFY)^2 \le \det({}^tXFX)\det({}^tYFY)$.
Comme $F$ est symétrique positive, elle est de la forme $F={}^tMM$ pour une
certaine matrice carrée $M$, donc en rempla\c cant $X$ et $Y$ par
$MX$ et $MY$, il suffit de prouver $\det({}^tXY)^2 \le \det({}^tXX)\det({}^tYY)$
pour toutes matrices $X,Y$ réelles rectangulaires de même taille.

En projetant chaque colonne de $Y$ sur le sous-espace vectoriel engendré par les colonnes de $X$,
on peut décomposer $Y = XA + B$ où $A$ est une matrice carrée et $B$ une matrice
rectangulaire de même taille que $X$ telle que ${}^tXB = 0$. Il reste à prouver:
$\det({}^tXXA)^2 \le \det({}^tXX)\det({}^tA{}^tXXA+{}^tBB)$, soit:
$\det({}^tA{}^tXXA) \le \det({}^tA{}^tXXA+{}^tBB)$.

On pose $U={}^tA{}^tXXA$ et $V = {}^tBB$: $U$ et $V$ sont des matrices réelles
symétriques positives de même taille, à priori quelquonques.
Si $U$ est inversible, on écrit $U={}^tPP$ avec $P$ inversible
et on est rammené à prouver que $1\le\det(I+{}^tP^{-1}VP^{-1}) = \det(I+W)$,
avec $W$ symétrique positive, ce qui résulte du fait que toutes les valeurs
propres de $I+W$ sont supérieures ou égales à $1$.
Si $U$ n'est pas inversible, on remplace $U$ par $U+\varepsilon I$ avec
$\varepsilon >0$, puis on fait tendre $\varepsilon$ vers $0^+$.

Remarque: il y a peut-être plus simple?
\fincorrection
\correction{005807}
\begin{enumerate}
 \item  Si la matrice de $f$ dans la base canonique de $\Rr^2$ est $A=\left(
\begin{array}{cc}
a&c\\
b&d
\end{array}
\right)$, 

\begin{center}
$Q(f) =\lambda(a^2+2bc+d^2)+\mu(ad-bc)$.
\end{center}

$Q$ est un polynôme homogène de degré $2$ en les coordonnées de $f$ dans la base canonique de $\mathcal{L}(\Rr^2)$ et donc $Q$ est une forme quadratique sur $\mathcal{L}(\Rr^2)$.

\item  \textbullet~Si $\lambda=\mu= 0$, $r = 0$ et $s = (0,0)$. Si $\lambda=0$ et $\mu\neq0$, 

\begin{center}
$Q(f)=\frac{\mu}{4}(a+d)^2-\frac{\mu}{4}(a-d)^2-\frac{mu}{4}(b+c)^2+\frac{\mu}{4}(b-c)^2$,
\end{center}

et donc $r=4$ et $s=(2,2)$.

\textbullet~Si $\lambda\neq0$,

\begin{align*}\ensuremath
Q(f)&=\lambda a^2+\mu ad + (2\lambda-\mu)bc +\lambda d^2=\lambda\left(a+\frac{\mu}{2\lambda}d\right)^2 +(2\lambda-\mu)bc +\left(\lambda-\frac{\mu^2}{4\lambda}\right)d^2\\
 &=\lambda\left(a+\frac{\mu}{2\lambda}d\right)^2+\left(\lambda-\frac{\mu^2}{4\lambda}\right)d^2+  \frac{2\lambda-\mu}{4}(b+c)^2-\frac{2\lambda-\mu}{4}(b-c)^2.
\end{align*}

Maintenant, les quatre formes linéaires $(a,b,c,d)\mapsto a+\frac{\mu}{2\lambda}d$, $(a,b,c,d)\mapsto d$, $(a,b,c,d)\mapsto b+c$ et $(a,b,c,d)\mapsto b-c$ sont linéairement indépendantes. Donc 

- si $\mu=2\lambda\;(\neq0)$, $r=1$,

- si $\mu=-2\lambda\;(\neq2\lambda)$, $r=3$,

- si $|\mu|\neq2|\lambda|\;(\neq0)$, $r=4$.

En particulier, si $\lambda=1$ et $\mu=0$, alors $r=4$ et $s=(3,1)$ et si $\lambda= 0$ et $\mu= 1$,  $r=4$ et $s = (2,2)$.
\end{enumerate}
\fincorrection
\correction{005808}
Dans le cas où $E$ est de dimension finie, la signature de $Q$ permet de conclure immédiatement. Supposons donc que $E$ n'est pas de dimension fine.

Par hypothèse, il existe un vecteur non nul $x_0$ tel que $Q(x_0) = 0$. Supposons $Q$ de signe constant. Ouite à remplacer $Q$ par $-Q$, on supposera que $Q$ est positive.
D?après l'inégalité de \textsc{Cauchy}-\textsc{Schwarz} (valable pour les formes quadratiques positives)

\begin{center}
$\forall y\in E,\;|\varphi(x_0,y)|\leqslant \sqrt{Q(x_0)}\sqrt{Q(y)}= 0$.
\end{center}

Donc $\forall y\in E,\;\varphi(x_0,y) = 0$ et $x_0$ est dans le noyau de $\varphi$. Puisque $x_0\neq0$, on en déduit que $\varphi$ est dégénérée. En résumé, si $Q$ est de signe constant, $\varphi$ est dégénérée ou encore si $\varphi$ est non dégénérée, $Q$ n'est pas de signe constant.
\fincorrection
\correction{005809}
\begin{enumerate}
 \item  Pour tout $(x_1,...,x_n)\in\Rr^n$,

\begin{center}
$Q(x_1,...,x_n)=\sum_{1\leqslant i,j\leqslant n}^{}\left(\int_{a}^{b}f_i(t)f_j(t)\;dt\right)x_ix_j=\int_{a}^{b}\left(\sum_{1\leqslant i,j\leqslant n}^{}x_ix_jf_i(t)f_j(t)\right)\;dt=\int_{a}^{b}\left(\sum_{i=1}^{n}x_if_i(t)\right)^2\;dt\geqslant0$.
\end{center}

Donc Q est une forme quadratique positive.

\item  De plus, pour tout $(x_1,\ldots,x_n)\in\Rr^n$, $Q((x_1,...,x_n)) = 0\Leftrightarrow\sum_{i=1}^{n}x_if_i  =0$ (fonction continue positive d'intégrale nulle). Donc

\begin{align*}\ensuremath
Q\;\text{définie}&\Leftrightarrow\forall(x_1,...,x_n)\in\Rr^n,\;[Q((x_1,...,x_n)) = 0\Rightarrow(x_1,...,x_n) = 0]\\
 &\Leftrightarrow\forall(x_1,...,x_n)\in\Rr^n,\;[\sum_{i=1}^{n}x_if_i=0\Rightarrow(x_1,...,x_n)= 0]\\
 &(f_1,...,f_n)\;\text{libre}.
\end{align*}

\item  Dans le cas particulier envisagé, la matrice de $Q$ dans la base canonique de $\Rr^n$ est la matrice de \textsc{Hilbert} $H_n=\left(\frac{1}{i+j-1}\right)_{1\leqslant i,j\leqslant n}$.
\end{enumerate}
\fincorrection
\correction{005810}
Posons $X=\left(
\begin{array}{c}
x_1\\
\vdots\\
x_n
\end{array}
\right)$ et $A =\left(\begin{array}{cccc}
0&x_1&\ldots&x_n\\
x_1& & & \\
\vdots& &S& \\
x_n& & & 
\end{array}
\right)=\left(
\begin{array}{cc}
0&{^t}X\\
X&S
\end{array}\right)$.

Un calcul par blocs fournit $\left(
\begin{array}{cc}
0&{^t}X\\
X&S
\end{array}\right)\left(
\begin{array}{cc}
1&0\\
0&S^{-1}
\end{array}\right)=\left(
\begin{array}{cc}
0&{^t}XS^{-1}\\
X&I_n
\end{array}\right)$  puis

\begin{center}
$\left(
\begin{array}{cc}
0&{^t}X\\
X&S
\end{array}\right)\left(
\begin{array}{cc}
1&0\\
0&S^{-1}
\end{array}\right)\left(
\begin{array}{cc}
-1&0\\
X&I_n
\end{array}\right)=\left(
\begin{array}{cc}
0&{^t}XS^{-1}\\
X&I_n
\end{array}\right)\left(
\begin{array}{cc}
-1&0\\
X&I_n
\end{array}\right)=\left(
\begin{array}{cc}
{^t}XS^{-1}X&{^t}XS^{-1}\\
0&I_n
\end{array}\right)$.
\end{center}

On en déduit que $\text{det}(A)\times\text{det}(S^{-1})\times(-1)={^t}XS^{-1}X$ puis que $Q(X)=-\text{det}(A)={^t}X((\text{det}(S))S^{-1})X ={^t}XS'X$ où $S'=(\text{det}(S))S^{-1}$.

Maintenant, la matrice $S$ est définie positive et donc ses valeurs propres sont des réels strictement positifs. Les valeurs propres de la matrice $S'$ sont les $\frac{\text{det}(S)}{\lambda}$ où $\lambda$ décrit le spectre de $S$ et donc la matrice $S'$ est aussi une matrice symétrique définie positive. $Q$ est donc une forme quadratique définie positive.
\fincorrection
\nocorrection
\nocorrection
\correction{005792}
Le produit scalaire usuel de $\mathcal{M}_3(\Rr)$ est défini par

\begin{center}
$\forall(A,B)\in(\mathcal{M}_3(\Rr))^2$, $\left(A|B\right)=\text{Tr}({^t}AB) =\sum_{1\leqslant i,j\leqslant n}^{}a_{i,j}b_{i,j}$.
\end{center}

\begin{enumerate}
 \item  Déterminons l'orthogonal de $\mathcal{A}_3(\Rr)$ dans $\mathcal{M}_3(\Rr)$. Soit $(A,B)\in\mathcal{S}_(\Rr)\times\mathcal{A}_3(\Rr)$.

\begin{center}
$\left(A|B\right)=\text{Tr}({^t}AB) =\text{Tr}(AB) =\text{Tr}(BA) = -\text{Tr}({^t}BA) = -\left(B|A\right)$.
\end{center}

et donc $\left(A|B\right)= 0$. Donc $\mathcal{S}_3(\Rr)\subset\left(\mathcal{A}_3(\Rr)\right)^\bot$ et comme de plus, $\text{dim}\left(\mathcal{S}_3(\Rr)\right)=\text{dim}\left((\left(\mathcal{A}_3(\Rr)\right)^\bot\right)$, on a montré que

\begin{center}
\shadowbox{
$\left(\mathcal{A}_3(\Rr)\right)^\bot=\mathcal{S}_3(\Rr)$.
}
\end{center}

\item  Ainsi, la projection orthogonale de $M$ sur $\mathcal{A}_3(\Rr)$ est exactement la partie antisymétrique $p_a(M)$ de $M$ et la distance cherchée est la norme de $M-p_a(M)=p_s(M)$ avec 

\begin{center}
$p_s(M)=\frac{1}{2}\left(\left(
\begin{array}{ccc}
0&1&0\\
0&0&1\\
0&0&0
\end{array}
\right)+
\left(
\begin{array}{ccc}
0&0&0\\
1&0&0\\
0&1&0
\end{array}
\right)\right)=\frac{1}{2}\left(
\begin{array}{ccc}
0&1&0\\
1&0&1\\
0&1&0
\end{array}
\right)$.
\end{center}

Par suite,

\begin{center}
$d\left(M,\mathcal{A}_3(R)\right)=\|p_s(M)\|=\frac{1}{2}\sqrt{1^2+1^2+1^2+1^2}=1$.
\end{center}
\end{enumerate}
\fincorrection
\nocorrection
\nocorrection
\nocorrection
\nocorrection
\nocorrection
\nocorrection
\nocorrection
\nocorrection
\nocorrection
\nocorrection
\nocorrection
\nocorrection
\nocorrection
\nocorrection
\nocorrection
\nocorrection
\nocorrection
\nocorrection
\correction{003735}
$\lambda = 0$, $f = \mathrm{id}_E$ et $\lambda = -\frac2{\|\vec v\,\|^2}$,
         $f =$ la symétrie par rapport à vect$(\vec v)$.
\fincorrection
\nocorrection
\nocorrection
\nocorrection
\nocorrection
\correction{003740}
$\varphi(A) = \mathrm{tr}(A\,{}^t\!A)$ donc pour toute matrice $P$
telle que $P\,{}^t\!P$ soit scalaire (non nulle) on a $\varphi(P^{-1}AP) = \varphi(A)$.
Ces matrices sont les matrices de la forme $P = \lambda M$ avec $M$ orthogonale
(matrices de similitude).

Réciproquement, soit $P$ telle que $\varphi(P^{-1}AP) = \varphi(A)$ et
$Q=P\,{}^t\!P$.

On a par polarisation~:
$\forall\ A,B,\ \mathrm{tr}(AQ\,{}^t\!B\,{}^t\!Q^{-1}) = \mathrm{tr}(A\,{}^t\!B)$ donc pour $B=QC$~:
$\forall\ A,C\ \mathrm{tr}(AQ{}^tC) = \mathrm{tr}(A{}^tCQ)$ ce qui implique~:
$\forall\ C,\ Q{}^tC = {}^tCQ$ et donc que $Q$ est scalaire.
\fincorrection
\nocorrection
\nocorrection
\nocorrection
\nocorrection
\nocorrection
\nocorrection
\nocorrection
\nocorrection
\nocorrection
\correction{003750}
\begin{enumerate}
  \item 
  \item On raisonne par récurrence sur $d = \dim(F) = \dim(G)$.
Pour $d=0$ il n'y a rien à prouver.

Pour $d \ge 1$ on considère $a \in F$ et $b \in G$ unitaires tels que
$(a\mid b)$ soit maximal.
Soient $F_1$ l'orthogonal de $a$ dans $F$ et $G_1$ l'orthogonal de $b$ dans $G$
(sous-espaces vectoriels de dimensions égales à $d-1$).
Le choix de $a,b$ fait que $F_1$ est orthogonal à $b$ et $G_1$ est orthogonal
à $a$, donc $F_1$ et $G_1$ sont tous deux inclus dans l'orthogonal de
$\mathrm{vect}(a,b)$. On peut trouver un endomorphisme de cet orthogonal qui
échange $F_1$ et $G_1$, que l'on complète par la symétrie orthogonale dans
$\mathrm{vect}(a,b)$ qui échange $a$ et $b$.
\end{enumerate}
 \fincorrection
\correction{003751}
\begin{enumerate}
  \item 
  \item 
  \item oui pour $\phi_P$.

Pour $\psi_P : \forall\ A,B,\ \mathrm{tr}({}^tAB) = \mathrm{tr}({}^tP^tA^t(P^{-1})P^{-1}BP )
      =\mathrm{tr}(P^tP^tA^t(P^{-1})P^{-1}B )$.

Donc $P^tP^tA^t(P^{-1})P^{-1} = A$, donc $P^tP$ est scalaire, donc $P$ est une
matrice de similitude.
\end{enumerate}
 \fincorrection
\correction{003752}
$A = 0$ ou $A\in{\cal O}^+(n)$.
\fincorrection
\correction{003753}
\begin{enumerate}
  \item $A = P-I$, $P \in {\cal O}(n)$.
  \item Hadamard.
\end{enumerate}
 \fincorrection
\nocorrection
\correction{003755}
\begin{enumerate}
  \item 
  \item Tout $f\in G$ vérifie $\det(f)\in\{-1,1\}$. Réciproquement,
soit $f\in U(E)\ \text{ tq }\det(f)\in\{-1,1\}$ et $F = \mathrm{Ker}(f-\mathrm{id})$.
On montre que $f$ est composée de réflexions par récurrence sur
$p=\mathrm{codim} F$.

$p=0  \Rightarrow  f=\mathrm{id}$. $p=1 \Rightarrow f$ est une réflexion car $F$ est un hyperplan et
$F^\perp$ est stable par~$f$.

$0,\dots,p-1 \Rightarrow p$~: soit $(e_1,\dots,e_n)$ une BON de~$E$ telle que
$(e_{p+1},\dots,e_n)$ est une BON de~$F$ et $e_1$ est vecteur propre de~$f$.
Soit $e'_1 = f(e_1) = \lambda e_1$, et
$\sigma,\sigma'$ deux réflexions telles que $\sigma(e_1)=e_2$ et
$\sigma'(e_2)=e'_1$. Alors $g=\sigma\circ\sigma'\circ f \in U(E)$,
$\det(g) = \det(f) \in\{-1,1\}$ et $\mathrm{codim}(\mathrm{Ker}(g-\mathrm{id})) < p$
donc $g$ est composée de réflexions et $f$ aussi.
\end{enumerate}
 \fincorrection
\nocorrection
\nocorrection
\nocorrection
\correction{003759}
$f = \mathrm{id}-r$ où $r(x_1,\dots,x_n) = (x_n,x_1,\dots,x_{n-1})$.
Donc $f^*\circ f = 2\mathrm{id}-r-r^{-1}$ a pour valeurs propres les nombres
$2-2\cos(2k\pi/n)$, $k\in{[[0,2n-1]]}$ et
$\|f\| = \begin{cases}2&\text{ si } n \text{ est pair}\cr2\cos(\pi/2n)&\text{ si } n \text{ est impair.}\cr \end{cases}$.
\fincorrection
\correction{003760}
Oui. Pour $n=1$ il y a égalité. Pour~$n=2$ cela résulte de la densité
de~$\mathbb{U}\cap\Q[i]$ dans~$\mathbb{U}$ (démonstration ci-dessous). Pour~$n$ quelconque,
il suffit de voir qu'une réflexion
quelconque est limite de réflexions à coefficients rationnels
(approcher un vecteur non nul normal à l'hyperplan de réflexion
par une suite de vecteurs rationnels).

Densité de~$\mathbb{U}\cap\Q[i]$ dans~$\mathbb{U}$~: pour $p\in\N^*$ on considère $z_p = \frac{(p^2-1)+2ip}{p^2+1}$.
On a $z_p\in \mathbb{U}\cap\Q[i]$ et $z_p \to 1$ lorsque $p\to\infty$.
Si~$z\in\mathbb{U}$ alors~$d(z,\mathbb{U}\cap\Q[i])\le d(z,\{z_p^k,\ k\in\Z\})\le \frac12|1-z_p|$.

\fincorrection
\correction{003761}
$^t\!AA = I \Leftrightarrow a^2+b^2+c^2=1$. Dans ce cas $\det(A) = \mathrm{tr}(A) = 1$ donc $f$ est un quart de tour.
 L'axe du quart de tour est engendré par $\begin{pmatrix}a\cr b\cr c\cr\end{pmatrix}$.
\fincorrection
\correction{005486}
\begin{enumerate}
 \item  Soit $\mathcal{B}$ une base orthonormée de $E$ et $M=\mbox{Mat}_{\mathcal{B}}(x_1,...,x_n)$ ($M$ est une matrice de format $(p,n)$).
Puisque $\mathcal{B}$ est orthonormée, le produit scalaire usuel des colonnes $C_i$ et $C_j$ est encore $x_i|x_j$.
Donc, $\forall(i,j)\in\llbracket1,n\rrbracket^2,\;{^t}C_iC_j=x_i|x_j$ ou encore 

\begin{center}
\shadowbox{
$G={^t}MM$.
}
\end{center}
Il s'agit alors de montrer que $\mbox{rg}(M)=\mbox{rg}({^t}MM)$. Ceci provient du fait que $M$ et ${^t}MM$ ont même noyau. En effet, pour $X\in\mathcal{M}_{n,1}(\Rr)$, 

$$X\in\mbox{Ker}M\Rightarrow MX=0\Rightarrow{^t}M\times MX=0\Rightarrow({^t}MM)X=0\Rightarrow X\in\mbox{Ker}({^t}MM)$$
et

\begin{align*}\ensuremath
X\in\mbox{Ker}({^t}MM)&\Rightarrow{^t}MMX=0\Rightarrow{^t}X{^t}MMX=0\Rightarrow{^t}(MX)MX=0\Rightarrow||MX||^2=0\Rightarrow MX=0\\
 &\Rightarrow X\in\mbox{Ker}M.
\end{align*}
Ainsi, $\text{Ker}(M)=\text{Ker}({^t}MM)=\text{Ker}(G(x_1,\ldots,x_n))$. Mais alors, d'après le théorème du rang, $\text{rg}(x_1,\ldots,x_n)=\text{rg}(M)=\text{rg}(G(x_1,\ldots,x_n))$.

\begin{center}
\shadowbox{
$\text{rg}(G(x_1,\ldots,x_n))=\text{rg}(x_1,\ldots,x_n)$.
}
\end{center}
 \item  Si la famille $(x_1,...,x_n)$ est liée, $\mbox{rg}(G)=\mbox{rg}(x_1,...,x_n)<n$, et donc, puisque $G$ est une matrice carrée de format $n$, $\gamma(x_1,...,x_n)=\mbox{det}(G)=0$.
Si la famille $(x_1, ... ,x_n)$ est libre, $(x_1,...,x_n)$ engendre un espace $F$ de dimension $n$. Soient $\mathcal{B}$ une base orthonormée de $F$ et $M$ la matrice de la famille $(x_1,...,x_n)$ dans $\mathcal{B}$. D'après 1), on a $G={^t}MM$ et d'autre part, $M$ est une matrice carrée. Par suite,

$$\gamma(x_1,...,x_n)=\mbox{det}({^t}MM)=\mbox{det}({^t}M)\mbox{det}(M)=(\mbox{det}M)^2>0.$$

 \item  On écrit $x=x-p_F(x)+p_F(x)$. La première colonne de $\gamma(x,x_1,...,x_n)$ s'écrit~:

$$\left(
\begin{array}{c}
||x||^2\\
x|x_1\\
x|x_2\\
\vdots\\
x|x_n
\end{array}
\right)=\left(
\begin{array}{c}
||x-p_F(x)+p_F(x)||^2\\
x-p_F(x)+p_F(x)|x_1\\
x-p_F(x)+p_F(x)|x_2\\
\vdots\\
x-p_F(x)+p_F(x)|x_n
\end{array}
\right)=\left(
\begin{array}{c}
||x-p_F(x)||^2\\
0|x_1\\
0|x_2\\
\vdots\\
0|x_n
\end{array}
\right)+\left(
\begin{array}{c}
||p_F(x)||^2\\
p_F(x)|x_1\\
p_F(x)|x_2\\
\vdots\\
p_F(x)|x_n
\end{array}
\right).$$  
(en 1ère ligne, c'est le théorème de \textsc{Pythagore} et dans les suivantes, $x-p_F(x)\in F^\bot$). Par linéarité par rapport à la première colonne, $\gamma(x,x_1,...,x_n)$ est somme de deux déterminants. Le deuxième est  $\gamma(p_F(x),x_1,...,x_n)$ et est nul car la famille $(p_F(x),x_1,...,x_n)$ est liée. On développe le premier suivant sa première colonne et on obtient~:

$$\gamma(x,x_1,...,x_n)=||x-p_F(x)||^2\gamma(x_1,...,x_n),$$
ce qui fournit la formule désirée.

\begin{center}
\shadowbox{
$\forall x\in E,\;d(x,F)=\|x-p_F(x)\|=\sqrt{\frac{\gamma(x,x_1,\ldots,x_n)}{\gamma(x_1,\ldots,x_n)}}$.
}
\end{center}
\end{enumerate}
\fincorrection
\correction{005488}
Un vecteur engendrant $D$ est $\overrightarrow{u}=(2,1,3)$. Pour $(x,y,z)\in\Rr^3$,

$$p((x,y,z))=\frac{(x,y,z)|(2,1,3)}{||(2,1,3)||^2}(2,1,3)=\frac{2x+y+3z}{14}(2,1,3).$$ 
On en déduit que $\mbox{Mat}_{\mathcal{B}}p=P=\frac{1}{14}\left(
\begin{array}{ccc}
4&2&6\\
2&1&3\\
6&3&9
\end{array}
\right)$, puis $\mbox{Mat}_{\mathcal{B}}s=2P-I=\frac{1}{7}\left(
\begin{array}{ccc}
-3&2&6\\
2&-6&3\\
6&3&2
\end{array}
\right)$.
Plus généralement, la matrice de la projection orthogonale sur le vecteur unitaire $(a,b,c)$ dans la base canonique orthonormée est $P=\left(
\begin{array}{ccc}
a^2&ab&ac\\
ab&b^2&bc\\
ac&bc&c^2
\end{array}
\right)$ et la matrice de la projection orthogonale sur le plan $ax+by+cz=0$ dans la base canonique orthonormée est $I-P=
\left(
\begin{array}{ccc}
1-a^2&-ab&-ac\\
-ab&1-b^2&-bc\\
-ac&-bc&1-c^2
\end{array}
\right)$.

\fincorrection
\correction{005494}
\textbf{1ère solution.}

\begin{align*}\ensuremath
\int_{0}^{1}(x^4-ax-b)^2\;dx&=\frac{1}{9}+\frac{1}{3}a^2+b^2-\frac{1}{3}a-\frac{2}{5}b+ab=\frac{1}{3}\left(a+\frac{1}{2}(3b-1)\right)^2-\frac{1}{12}(3b-1)^2+b^2-\frac{2}{5}b+\frac{1}{9}\\
 &=\frac{1}{3}\left(a+\frac{1}{2}(3b-1)\right)^2+\frac{1}{4}b^2+\frac{1}{10}b+\frac{1}{36}=\frac{1}{3}\left(a+\frac{1}{2}(3b-1)\right)^2+\frac{1}{4}\left(b+\frac{1}{5}\right)^2+\frac{4}{225}\geq\frac{4}{225},
\end{align*}
avec égalité si et seulement si $a+\frac{1}{2}(3b-1)=b+\frac{1}{5}=0$ ou encore $b=-\frac{1}{5}$ et $a=\frac{4}{5}$.

\begin{center}
\shadowbox{
$\int_{0}^{1}(x^4-ax-b)^2\;dx$ est minimum pour $a=\frac{4}{5}$ et $b=-\frac{1}{5}$ et ce minimum vaut $\frac{4}{225}$.
}
\end{center}
\textbf{2ème solution.} $(P,Q)\mapsto\int_{0}^{1}P(t)Q(t)\;dt$ est un produit scalaire sur $\Rr_4[X]$ et $\int_{0}^{1}(x^4-ax-b)^2dx$ est, pour ce produit scalaire, le carré de la distance du polynôme $X^4$ au polynôme de degré inférieur ou égal à $1$, $aX+b$. On doit calculer $\mbox{Inf}\left\{\int_{0}^{1}(x^4-ax-b)^2\;dx,\;(a,b)\in\Rr^2\right\}$ qui est le carré de la distance de $X^4$ à $F=\Rr_1[X]$. On sait que cette borne inférieure est un minimum, atteint une et une seule fois quand $aX+b$ est la projection orthogonale de $X^4$ sur $F$.
Trouvons une base orthonormale de $F$. L'orthonormalisée $(P_0,P_1)$ de $(1,X)$ convient.
$||1||^2=\int_{0}^{1}1\;dt=1$ et $P_0=1$. Puis $X-(X|P_0)P_0=X-\int_{0}^{1}t\;dt=X-\frac{1}{2}$, et comme 
$||X-(X|P_0)P_0||^2=\int_{0}^{1}\left(t-\frac{1}{2}\right)^2\;dt=\frac{1}{3}-\frac{1}{2}+\frac{1}{4}=\frac{1}{12}$, $P_1=2\sqrt{3}\left(X-\frac{1}{2}\right)=\sqrt{3}(2X-1)$.
La projection orthogonale de $X^4$ sur $F$ est alors $(X^4|P_0)P_0+(X^4|P_1)P_1$ avec $(X^4|P_0)=\int_{0}^{1}t^4\;dt=\frac{1}{5}$ et $(X^4|P_1)=\sqrt{3}\int_{0}^{1}t^4(2t-1)\;dt=\sqrt{3}(\frac{1}{3}-\frac{1}{5})=\frac{2\sqrt{3}}{15}$. Donc, la projection orthogonale de $X^4$ sur $F$ est $\frac{1}{5}+\frac{2\sqrt{3}}{15}\sqrt{3}(2X-1)=\frac{1}{5}(4X-1)$.
Le minimum cherché est alors $\int_{0}^{1}\left(t^4-\frac{1}{5}(4t-1)\right)^2\;dt=...=\frac{4}{225}$.
\fincorrection
\correction{005504}
Notons $P$ le plan d'équation $x+y=0$ dans la base $\mathcal{B}=(i,j,k)$. $P$ est le plan de vecteur normal $n=i+j$.

\begin{enumerate}
 \item  Soit $s$ la symétrie orthogonale par rapport au plan $P'$ d'équation $x-y+z=0$. $s(P)$ est le plan de vecteur normal $s(n)$. 
Or, le vecteur $n$ est dans $P'$  et donc $s(n)=n$ puis $s(P)=P$.

\begin{center}
\shadowbox{
$s(P)$ est le plan $P$.
}
\end{center}
 \item  Notons $\sigma$ la symétrie orthogonale par rapport au vecteur $u=(1,1,1)$. $\sigma(P)$ est le plan de vecteur normal 

\begin{center}
$\sigma(n)=2\frac{n.u}{\|u\|^2}u-n=2\frac{2}{3}(1,1,1)-(1,1,0)=\frac{1}{3}(1,1,4)$.
\end{center}

\begin{center}
\shadowbox{
$\sigma(P)$ est le plan d'équation $x+y+4z=0$.
}
\end{center}
 \item  Notons $r$ la rotation d'angle $\frac{\pi}{4}$ autour du vecteur unitaire $u=\frac{1}{\sqrt{3}}(1,1,1)$. $r(P)$ est le plan de vecteur normal 
\begin{align*}\ensuremath
r(n)&=\left(\cos\frac{\pi}{4}\right)n+\left(1-\cos\frac{\pi}{4}\right)(n.u)u+\left(\sin\frac{\pi}{4}\right)u\wedge n\\
 &=\frac{1}{\sqrt{2}}(1,1,0)+\left(1-\frac{1}{\sqrt{2}}\right)\times\frac{2}{3}(1,1,1)+\frac{1}{\sqrt{2}}\times\frac{1}{\sqrt{3}}(-1,1,0)\\
 &=\frac{1}{3\sqrt{2}}(3+2(\sqrt{2}-1)-\sqrt{3},3+2(\sqrt{2}-1)+\sqrt{3},2(\sqrt{2}-1))=\frac{1}{3\sqrt{2}}(1+2\sqrt{2}-\sqrt{3},1+2\sqrt{2}+\sqrt{3},2(\sqrt{2}-1)).
\end{align*}

\begin{center}
\shadowbox{
$r(P)$ est le plan d'équation $(1+2\sqrt{2}-\sqrt{3})x+(1+2\sqrt{2}+\sqrt{3})y+2(\sqrt{2}-1)z=0$.
}
\end{center}
\end{enumerate}
\fincorrection
\correction{005509}
Notons $p$ la projection sur $(P)$ parallèlement à $(\Delta)$.
\textbullet~Déterminons un repère de $(D)$. 
\begin{center}$\left\{
\begin{array}{l}
x+y+z=1\\
x-2y-z=0
\end{array}
\right.\Leftrightarrow
\left\{
\begin{array}{l}
y+z=-x+1\\
2y+z=x
\end{array}
\right.
\Leftrightarrow
\left\{
\begin{array}{l}
y=2x-1\\
z=-3x+2
\end{array}
\right.
$
\end{center}
$(D)$ est la droite de repère $(A,\overrightarrow{u})$ où $A(0,-1,2)$ et $\overrightarrow{u}(1,2,-3)$.
\textbullet~$(\Delta)$ est dirigée par le vecteur $\overrightarrow{u'}(1,3,2)$. $\overrightarrow{u}$ n'est pas colinéaire à $\overrightarrow{u'}$ et donc $(D)$ n'est pas parallèle à $(\Delta)$. On en déduit que $p(D)$ est une droite. 
Plus précisément, $p(D)$ est la droite intersection du plan $(P)$ et du plan $(P')$ contenant $(D)$ et parallèle à $(\Delta)$. Déterminons une équation de $(P')$. Un repère de $(P')$ est $(A,\overrightarrow{u},\overrightarrow{u'})$. Donc

\begin{center}
$M(x,y,z)\in(P')\Leftrightarrow\left|
\begin{array}{ccc}
x&1&1\\
y+1&2&3\\
z-2&-3&2
\end{array}
\right|=0\Leftrightarrow 13x-5(y+1)+(z-2)=0\Leftrightarrow13x-5y+z=7$.
\end{center}

Finalement

\begin{center}
\shadowbox{
$p(D)$ est la droite dont un système d'équations cartésiennes est $\left\{
\begin{array}{l}
13x-5y+z=7\\
x+3y+2z=6
\end{array}
\right.$
}
\end{center}
\fincorrection
\correction{005520}
Notons $p$ la projection orthogonale sur $(P)$.
Un repère de $(D)$ est $\left(A,\overrightarrow{u}\right)$ où $A(0,-1,2)$ et $\overrightarrow{u}(1,2,-3)$. Un vecteur normal à $(P)$ est $\overrightarrow{n}(1,3,2)$. $\overrightarrow{u}$ et $\overrightarrow{n}$ ne sont pas colinéaires et donc $p(D)$ est une droite du plan $(P)$.
Plus précisément, $p(D)$ est l'intersection du plan $(P)$ et du plan $(P')$ contenant $(D)$ et perpendiculaire à $(P)$. Un repère de $(P')$ est $\left(A,\overrightarrow{u},\overrightarrow{n}\right)$. Donc

\begin{align*}\ensuremath
M(x,y,z)\in(P')&\Leftrightarrow\left|
\begin{array}{ccc}
x&1&1\\
y+1&2&3\\
z-2&-3&2
\end{array}
\right|=0\Leftrightarrow13x-5(y+1)+(z-2)=0\Leftrightarrow13x-5y+z=7.
\end{align*}

\begin{center}
\shadowbox{
La projetée orthogonale de $(D)$ sur $(P)$ est la droite d'équations $\left\{
\begin{array}{l}
13x-5y+z=7\\
x+3y+2z=6
\end{array}
\right.$.
}
\end{center}
\fincorrection
\nocorrection
\nocorrection
\nocorrection
\nocorrection
\nocorrection
\nocorrection
\nocorrection
\nocorrection
\nocorrection
\nocorrection
\nocorrection
\nocorrection
\nocorrection
\nocorrection
\nocorrection
\nocorrection
\correction{005484}
La famille $(V_1,V_2)$ est clairement libre et donc une base de $F$. Son orthonormalisée $(e_1,e_2)$ est une base orthonormée de $F$.
$||V_1||=\sqrt{1+4+1+1}=\sqrt{7}$ et $e_1=\frac{1}{\sqrt{7}}V_1=\frac{1}{\sqrt{7}}(1,2,-1,1)$.
$(V_2|e_1)=\frac{1}{\sqrt{7}}(0+6-1-1)=\frac{4}{\sqrt{7}}$  puis $V_2-(V_2|e_1)e_1=(0,3,1,-1)-\frac{4}{7} (1,2,-1,1)=\frac{1}{7}(-4,13,11,-11)$ puis $e_2=\frac{1}{\sqrt{427}}(-4,13,11,-11)$.
Une base orthonormée de $F$ est $(e_1,e_2)$ où $e_1=\frac{1}{\sqrt{7}}(1,2,-1,1)$ et $e_2=\frac{1}{\sqrt{427}}(-4,13,11,-11)$.
Soit $(x,y,z,t)\in\Rr^4$.

$$(x,y,z,t)\in F^\bot\Leftrightarrow(x,y,z,t)\in(V_1,V_2)^\bot\Leftrightarrow\left\{
\begin{array}{l}
x+2y-z+t=0\\
3y+z-t=0
\end{array}
\right..$$
\fincorrection
\correction{005500}
\begin{enumerate}
 \item  La symétrie, la bilinéarité et la positivité sont claires. Soit alors $P\in\Rr_n[X]$. 

\begin{align*}\ensuremath
P|P=0&\Rightarrow\int_{0}^{1}P^2(t)\;dt=0\\
 &\Rightarrow\forall t\in[0,1],\;P^2(t)=0\;(\mbox{fonction continue, positive, d'intégrale nulle})\\
 &\Rightarrow P=0\;(\mbox{polynôme ayant une infinité de racines}).
\end{align*}
Ainsi, l'application $(P,Q)\mapsto\int_{0}^{1}P(t)Q(t)\;dt$ est un produit scalaire sur $\Rr_n[X]$.
 \item  Pour vérifier que la famille $\left(\frac{L_p}{||L_p||}\right)_{0\leq p\leq n}$ est l'orthonormalisée de \textsc{Schmidt} de la base canonique de $E$, nous allons vérifier que
  \begin{enumerate}
  \item $\forall p\in\llbracket0,n\rrbracket,\;\mbox{Vect}(L_0,L_1,...,L_p)=\mbox{Vect}(1,X,...,X^p)$,\rule[-4mm]{0mm}{0mm}
  \item la famille $\left(\frac{L_p}{||L_p||}\right)_{0\leq p\leq n}$ est orthonormale,
  \item $\forall p\in\llbracket0,n\rrbracket,\;L_p|X^p>0$.\rule{0mm}{5mm}
  \end{enumerate}
Pour a), on note que $L_p$ est un polynôme de degré $p$ (et de coefficient dominant $\frac{(2p)!}{p!}$). Par suite, $(L_0,L_1,...,L_p)$ est une base de $\Rr_p[X]$, ou encore, $\forall p\in\llbracket0,n\rrbracket,\;\mbox{Vect}(L_0,L_1,...,L_p)=\mbox{Vect}(1,X,...,X^p)$.
Soit $p\in\llbracket0,n\rrbracket$. Soit $P$ un polynôme de degré inférieur ou égal à $p$. Si $p\geq1$, une intégration par parties fournit :

\begin{align*}\ensuremath
L_p|P&=\int_{-1}^{1}((t^2-1)^p)^{(p)}P(t)\;dt=\left[((t^2-1)^p)^{(p-1)}P(t)\right]_{-1}^1-\int_{-1}^{1}((t^2-1)^p)^{(p-1)}P'(t)\;dt\\
 &=-\int_{-1}^{1}((t^2-1)^p)^{(p-1)}P'(t)\;dt.
\end{align*}
En effet, $1$ et $-1$ sont racines d'ordre $p$ de $(t^2-1)^p$ et donc d'ordre $p-k$ de $((t^2-1)^p)^{(k)}$ pour $0\leq k\leq p$ et en particulier, racines de chaque $((t^2-1)^p)^{(k)}$ pour $0\leq k\leq p-1$.
En réitérant, on obtient pour tout $k\in\llbracket0,p\rrbracket$, $L_p|P=(-1)^k\int_{-1}^{1}((t^2-1)^p)^{(p-k)}P^{(k)}(t)\;dt$ et pour $k=p$, on obtient enfin $L_p|P=(-1)^p\int_{-1}^{1}(t^2-1)^pP^{(p)}(t)\;dt$, cette formule restant vraie pour $p=0$.
Soient $p$ et $q$ deux entiers tels que $0\leq q<p\leq n$. D'après ce qui précède, $L_p|L_q=(-1)^p\int_{-1}^{1}(t^2-1)^pL_q^{(p)}(t)\;dt=0$ car $q=\mbox{deg}(L_q)<q$. Ainsi, la famille $(L_p)_{0\leq p\leq n}$ est donc une famille orthogonale de $n+1$ polynômes tous non nuls et est par suite est une base orthogonale de $\Rr_n[X]$. On en déduit que la famille $\left(\frac{L_p}{||L_p||}\right)_{0\leq p\leq n}$ est une base orthonormale de $\Rr_n[X]$.
Enfin, $L_p|X^p=(-1)^p\int_{-1}^{1}(t^2-1)^p(t^p)^{(p)}\;dt=p!\int_{-1}^{1}(1-t^2)^p\;dt>0$. On a montré que 

\begin{center}
\shadowbox{
la famille $\left(\frac{L_p}{||L_p||}\right)_{0\leq p\leq n}$ est l'orthonormalisée de la base canonique de $\Rr_n[X]$.
}
\end{center}
Calculons $\|L_p\|$. On note que $L_p\in(L_0,...,L_{p-1})^\bot=(\Rr_{p-1}[X])^\bot$. Par suite,

\begin{align*}\ensuremath
||L_p||^2&=L_p|L_p=L_p|\mbox{dom}(L_p)X^p\;(\mbox{car}\;L_p\in(\Rr_{p-1}[X])^\bot)\\
 &=\frac{(2p)!}{p!}L_p|X^p=\frac{(2p)!}{p!}p!\int_{-1}^{1}(1-t^2)^p\;dt=2(2p)!\int_{0}^{1}(1-t^2)^p\;dt\\
 &=2(2p)!\int_{\pi/2}^{0}(1-\cos^2u)^p(-\sin u)\;du=2(2p)!\int_{0}^{\pi/2}\sin^{2p+1}u\;du\\
 &=2(2p)!W_{2p+1}\;(\mbox{intégrales de \textsc{Wallis}})\\
 &=2(2p)!\frac{2^{2p}(p!)^2}{(2p+1)!}\;(\mbox{à revoir})\\
 &=\frac{2}{2p+1}2^{2p}(p!)^2.
\end{align*}
Donc, $\forall p\in\llbracket0,n\rrbracket,\;\|L_p\|=\sqrt{\frac{2}{2p+1}}2^pp!$. On en déduit que la famille $\left(
\sqrt{\frac{2p+1}{2}}\frac{1}{2^pp!}((X^2-1)^p)^{(p)}\right)_{0\leq p\leq n}$ est une base orthonormale de $\Rr_n[X]$ (pour le produit scalaire considéré).
\end{enumerate}
\fincorrection
\correction{005772}
Soit $n\in\Nn$. Posons $\ell_n=(X^2-1)^n$ de sorte que $L_n=\ell_n^{(n)}$. $L_n$ est un polynôme de degré $n$ car $\ell_n$ est de degré $2n$.

 
\begin{enumerate}
 \item 
   \begin{enumerate}
   \item Soient $n\in\Nn^*$ et $P\in E$. Une intégration par parties fournit

\begin{center}
$\left(L_n|P\right)=\int_{-1}^{1}L_n(x)P(x)\;dx =\int_{-1}^{1}(\ell_n)^{(n)}(x) P(x)\;dx =\left[(\ell_n)^{(n-1)}(x) P(x)\right]_{-1}^1-\int_{-1}^{1}(\ell_n)^{(n-1)}(x) P'(x)\;dx$.
\end{center}

Maintenant, $-1$ et $1$ sont racines d'ordre $n$ du polynôme $\ell_n$ et donc, pour tout $k\in\llbracket0,n\rrbracket$, $-1$ et $1$ sont racines d'ordre $n-k$ de $\ell_n^{(k)}$ et en particulier racines de $(\ell_n)^{(k}$ pour $k\in\llbracket0,n-1\rrbracket$. Donc

\begin{center}
$\left(L_n|P\right)=-\int_{-1}^{1}(\ell_n)^{(n-1)}(x) P'(x)\;dx$.
\end{center}

Plus généralement, si pour un entier $k\in\llbracket0,n-1\rrbracket$, $\left(L_n|P\right)= (-1)^{k}\int_{-1}^{1}(\ell_n)^{(n-k)}(x) P^{(k)}(x)\;dx$ alors

\begin{align*}\ensuremath
\left(L_n|P\right)&= (-1)^{k}\left(\left[(\ell_n^{n-k-1}(x)P^{(k)}(x)\right]_{-1}^{1}-\int_{-1}^{1}(\ell_n)^{(n-k-1)}(x)P^{(k+1)}(x)\;dx\right)\\
 &=(-1)^{k+1}\int_{-1}^{1}(\ell_n)^{(n-k-1)}(x) P^{(k+1)}(x)\;dx.
\end{align*}

On a montré par récurrence que pour tout entier $k\in\llbracket0,n\rrbracket$, 
$\left(L_n|P\right)=(-1)^{k}\int_{-1}^{1}(\ell_n)^{(n-k)}(x) P^{(k)}(x)\;dx$. En particulier 

\begin{center}
$\left(L_n|P\right)=(-1)^{n}\int_{-1}^{1}\ell_n(x)P^{(n)}(x)\;dx=\int_{-1}^{1}(1-x^2)^nP^{(n)}(x)\;dx$/quad$(*)$.
\end{center}

Cette dernière égalité reste vraie pour $n=0$ et on a montré que

\begin{center}
\shadowbox{
$\forall n\in\Nn$, $\forall P\in\Rr[X]$, $\left(L_n|P\right)=\int_{-1}^{1}\left(1-x^2\right)^nP^{(n)}(x)\;dx$.
}
\end{center}

Soient alors $n$ et $p$ deux entiers naturels tels que $0\leqslant p<n$. Puisque $\text{deg}(L_p)=p<n$, on a $\left(L_n|L_p\right)=0$. On a montré que

\begin{center}
\shadowbox{
La famille $(L_k)_{0\leqslant k\leqslant n}$ est une base orthogonale de l'espace $(\Rr[X],\;|\;)$.
}
\end{center}

  \item
On applique maintenant la formule $(*)$ dans le cas particulier $P =L_n$. On obtient 

\begin{align*}\ensuremath
\|L_n\|^2&=\int_{-1}^{1}\left(1-x^2\right)^nL_n^{(n)}(x)\;dx=2\times(2n)!\int_{0}^{1}(1-x^2)^n\;dx=2\times(2n)!\int_{\pi/2}^{0}(1-\cos^2t)^n(-\sin t)\;dt\\
 &=2\times(2n)!\int_{0}^{\pi/2}\sin^{2n+1}t\;dt=2\times(2n)!W_{2n+1}\;(\text{intégrales de \textsc{Wallis}}).
\end{align*}

On \og sait \fg que  $\forall n\in\Nn$, $W_{2n+1}=\frac{2n}{2n+1}W_{2n-1}=\frac{(2n)\times(2n-2)\times\ldots\times2}{(2n+1)\times(2n-1)\times\ldots\times3}W_1 =\frac{2^{2n}n!^2}{(2n+1)!}$. On obtient alors 

\begin{center}
$\|L_n\|^2=\frac{2^{2n}n!^2}{(2n+1)!}\times2\times(2n)!=\frac{2^{2n+1}n!^2}{2n+1}$,
\end{center}

\begin{center}
\shadowbox{
$\forall n\in\Nn$, $\|L_n\|=\sqrt{\frac{2}{2n+1}}\frac{2^nn!}{2n+1}$.
}
\end{center}

On en déduit que la famille $\left(\sqrt{\frac{2n+1}{2}}\frac{1}{2^nn!}L_n\right)_{n\in\Nn}$ est une base orthonormale de $(\Rr[X],\;|\;)$. Pour $n\in\Nn$, on pose $P_n=\sqrt{\frac{2n+1}{2}}\frac{1}{2^nn!}\left((X^2-1)^n\right)^{(n)}$.
  \end{enumerate}

\item  La famille $(P_n)_{n\in\Nn}$ est une base de orthonormée de $\Rr[X]$. Chaque $P_n$, $n\in\Nn$, est de degré $n$ et donc, $\forall n\in\Nn$, $\text{Vect}(P_0,...,P_n) =\text{Vect}(1,X,...,X^n)$ et de plus, pour $n\in\Nn$

\begin{center}
$P_n|X^n=\frac{1}{\text{dom}}\left((P_n)|\text{dom}(P_n)X^n\right)=\frac{1}{\text{dom}(P_n)}\left(P_n|P_n\right)$
\end{center}

car $P_n\in(P_0,\ldots,P_{n-1})^\bot=(1,X,\ldots,X^{n-1})^\bot=(\Rr_{n-1}[X])^\bot$. Ceci montre que $P_n|X^n>0$. 

\begin{center}
\shadowbox{
\begin{tabular}{c}
L'orthonormalisée de la base canonique de $\Rr[X]$ est la famille des polynômes de \textsc{Lengendre}\\
$\left(\sqrt{\frac{2n+1}{2}}\frac{1}{2^nn!}\left((X^2-1)^n\right)^{(n)}\right)_{n\in\Nn}$.
\end{tabular}
}
\end{center}
\end{enumerate}
\fincorrection
\correction{005779}
\begin{enumerate}
 \item  L'existence, la bilinéarité, la symétrie et la positivité sont immédiates. Soit $P\in\Rr[X]$. 

\begin{align*}\ensuremath
\Phi(P,P) = 0&\Rightarrow\int_{0}^{1}f(t)P^2(t)\;dt = 0\\
 &\Rightarrow \forall t\in[0,1],\;f(t)P^2(t)= 0\; (\text{fonction continue positive d'intégrale nulle}).
\end{align*}

Maintenant, la fonction $f$ est continue, positive sur $[0,1]$ et n'est pas nulle. Donc la fonction $f$ est strictement positive sur un intervalle ouvert non vide inclus dans le segment $[0,1]$. Par suite, le polynôme $P$ a une infinité de racines et finalement $P = 0$.

\begin{center}
\shadowbox{
L'application $\Phi$ est un produit scalaire sur $\Rr[X]$.
}
\end{center}

\item  L'orthonormalisée de la base canonique de $\Rr[X]$ répond à la question.

\item  Soit $n$ un entier naturel non nul. Le polynôme $P_n\in(P_0,...,P_{n-1})^\bot=(\Rr_{n-1}[X])^\bot$.
Soit $p$ le nombre de racines réelles d'ordre impair du polynôme $P_n$. Soient $a_1$,..., $a_p$ ces racines (deux à deux distinctes, réelles et d'ordre impair) dans le cas où $p\geqslant1$. Si $p\geqslant1$, on pose $Q= (X-a_1)...(X-a_p)$ et si $p=0$, on pose $Q = 1$.

Si $p < n$, le polynôme $Q$ est orthogonal à $P_n$ car de degré strictement plus petit que le de gré de $P_n$. D'autre part , au vu de la définition de $Q$, la fonction $t\mapsto f(t)P_n(t)Q(t)$ est continue sur $[0,1]$, de signe constant sur $[0,1]$, d(intégrale nulle sur $[0,1]$. La fonction $t\mapsto f(t)P_n(t)Q(t)$ est donc nulle. On en déduit que le polynôme $P_n$ est le polynôme nul ce qui n'est pas. Donc $p = n$ ce qui signifie que le polynôme $P_n$ a $n$ racines réelles simples.
\end{enumerate}
\fincorrection
\correction{005789}
Si la famille $(x_1,...,x_n)$ est liée, l'inégalité est vraie.

Si la famille $(x_1,...,x_n)$ est libre, on peut considérer $B_0 =(e_1,...,e_n)$ l'orthonormalisée de \textsc{Schmidt} de la famille $(x_1,...,x_n)$. Les bases $B_0$ et $B$ sont des bases orthonormées de $E$ et donc

\begin{align*}\ensuremath
\left|\text{det}_B(x_1,...,x_n)\right|&=\left|\text{det}_{B_0}(x_1,...,x_n)\right|
=\text{abs}\left(\left|\begin{array}{cccc}(x_1|e_1)&\times&\ldots&\times\\
 0&\ddots&\ddots&\vdots\\
 \vdots&\ddots&\ddots&\times\\
 0&\ldots&0&(x_n|e_n)\end{array}\right|\right)\\
 &=\prod_{k=1}^{n}|(x_k|e_k)|\leqslant\prod_{k=1}^{n}\|x_k\|\|e_k\|\;(\text{d'après l'inégalité de \textsc{Cauchy}-\textsc{Schwarz}})\\
 &=\prod_{k=1}^{n}\|x_k\|.
\end{align*}

\begin{center}
\shadowbox{
$\forall(x_1,\ldots,x_n)\in E^n$, $\left|\text{det}_B(x_1,...,x_n)\right|\leqslant\prod_{k=1}^{n}\|x_k\|$ (inégalité de \textsc{Hadamard}).
}
\end{center}

Ensuite, 

- si la famille $(x_1,...,x_n)$ est liée, on a l'égalité si et seulement si l'un des vecteurs $x_k$ est nul 

- si la famille $(x_1,...,x_n)$ est libre, on a l'égalité si et seulement si $\forall k\in\llbracket1,n\rrbracket$, $|(x_k|e_k)|=\|x_k\|\|e_k\|$. Les cas d'égalité

de l'inégalité de \textsc{Cauchy}-\textsc{Schwarz} étant connus, on a l'égalité si et seulement si $\forall k\in\llbracket1,n\rrbracket$, $x_k$ est colinéaire à $e_k$ 

ou encore si et seulement si la famille $(x_1,...,x_n)$ est orthogonale.

En résumé, l'inégalité de \textsc{Hadamard} est une égalité si et seulement si la famille $(x_1,...,x_n)$ est orthogonale libre ou si l'un des vecteurs est nul.
\fincorrection
\correction{005790}
C'est l'exercice \ref{exo:sperou4}.
\fincorrection
\correction{005806}
\begin{enumerate}
 \item  \textbf{1ère solution.} La matrice de la forme quadratique $Q$ dans la base canonique de $\Rr^3$ est $A=\left(\begin{array}{ccc}
2&\frac{3}{2}&-2\\
\frac{3}{2}&-2&\frac{7}{2}\\
-2&\frac{7}{2}&-6
\end{array}
\right)$.

Le polynôme caractéristique  de $A$ est 

\begin{align*}\ensuremath
\chi_A&=\left|\begin{array}{ccc}
2-X&\frac{3}{2}&-2\\
\frac{3}{2}&-2-X&\frac{7}{2}\\
-2&\frac{7}{2}&-6-X
\end{array}
\right|= (2-X)\left(X^2+8X-\frac{1}{4}\right) -\frac{3}{2}\left(-\frac{3}{2}X-2\right)-2\left(-2X+\frac{5}{4}\right) \\
 &= -X^3-6X^2+\frac{45}{2}X=-X\left(X^2+6X-\frac{45}{2}\right).
\end{align*}

Puisque $A$ est symétrique réelle, on sait que les valeurs propres de $A$ sont réelles. $\chi_A$ admet pour racines $0$ et deux réels non nuls de signes contraires (puisque leur produit vaut $-\frac{45}{2}$). Par suite, le rang et la signature de $Q$ sont 

\begin{center}
\shadowbox{
$r=2$ et $s = (1,1)$.
}
\end{center}

\textbf{2ème solution.} On effectue une réduction de \textsc{Gauss}.

\begin{align*}\ensuremath
Q((x,y,z))&= 2x^2-2y^2-6z^2+3xy-4xz+7yz=2x^2+x(3y-4z)-2y^2+7yz-6z^2\\
 &= 2\left(x+\frac{3}{4}y-z\right)^2-2\left(\frac{3}{4}y-z\right)^2-2y^2+7yz-6z^2 =2\left(x+\frac{3}{4}y-z\right)^2-\frac{25}{8}y^2+10yz-8z^2\\
 &= 2\left(x+\frac{3}{4}y-z\right)^2-\frac{25}{8}\left(y-\frac{8}{5}z\right)^2.
\end{align*}

Les formes linéaires $(x,y,z)\mapsto x+\frac{3}{4}y-z$ et $(x,y,z)\mapsto y-\frac{8}{5}z$ étant linéairement indépendantes, on retrouve le fait que $Q$ est de rang $r=2$ et de signature $s=(1,1)$. La forme quadratique $Q$ est dégénérée et n'est ni positive ni négative.

\item  La matrice de $Q$ dans la base canonique $(i,j,k)$ est $A=\left(
\begin{array}{ccc}
3&-1&-1\\
-1&3&-1\\
-1&-1&3
\end{array}
\right)$. Le nombre $4$ est valeur propre de $A$ et puisque $A$ est diagonalisable, $4$ est valeur propre d'ordre $\text{dim}(\text{Ker}(A-4I_3))= 3 -\text{rg}(A-4I_3)=2$.
La dernière valeur propre $\lambda$ est fournie par $4+4+\lambda=\text{Tr}(A)= 9$ et $? = 1$. Ainsi, $\text{Sp}(A) = (1,4,4)$.

Les trois valeurs propres de $A$ sont strictement positives et donc la forme quadratique $Q$ est de rang $3$ et de signature $(3,0)$.

\begin{center}
\shadowbox{
$Q$ est définie positive.
}
\end{center}

\item  Effectuons une réduction de \textsc{Gauss}.

\begin{center}
$Q((x,y,z,t))=xy+yz+zt+tx=(x+z)(y+t) =\frac{1}{4}(x+y+z+t)^2-\frac{1}{4}(x-y+z-t)^2$.
\end{center}

Puisque les deux formes linéaires $(x,y,z,t)\mapsto x+y+z+t$  et $(x,y,z,t)\mapsto x-y+z-t$ sont linéairement indépendantes, la forme quadratique $Q$ est de rang $r=2$ et de signature $s=(1,1)$.

\item  Effectuons une réduction de \textsc{Gauss}.

\begin{align*}\ensuremath
Q((x,y,z,t))&=x^2+(4+\lambda)y^2+(1+4\lambda)z^2+\lambda t^2+4xy+2xz+4(1-\lambda)yz+2\lambda yt+(1-4\lambda)zt\\
 &=(x+2y+z)^2 +\lambda y^2+4\lambda z^2+\lambda t^2 -4\lambda yz+2\lambda yt+(1-4\lambda)zt\\
 &= (x+2y+z)^2+\lambda(y-2z+t)^2+zt =(x+2y+z)^2+\lambda(y-2z+t)^2+\frac{1}{4}(z+t)^2-\frac{1}{4}(z-t)^2.
\end{align*}

Si $\lambda< 0$, la forme quadratique $Q$ est de rang $4$ et de signature $(2,2)$.

Si $\lambda = 0$, la forme quadratique $Q$ est de rang $3$ et de signature $(2,1)$.

Si $\lambda> 0$, la forme quadratique $Q$ est de rang $4$ et de signature $(3,1)$.

\item  \textbf{1ère solution.} La matrice de la forme quadratique $Q$ dans la base canonique est $A=\frac{1}{2}\left(
\begin{array}{ccccc}
0&1&1&1&1\\
1&0&1&1&1\\
1&1&0&1&1\\
1&1&1&0&1\\
0&1&1&1&1
\end{array}
\right)$. Les valeurs propres de $A$ sont $-\frac{1}{2}$ qui est d'ordre $4$ et $2$ qui est valeur propre simple.Donc, la signature de la forme quadratique $Q$ est 

\begin{center}
\shadowbox{
$s=(1,4)$.
}
\end{center}

\textbf{2ème solution.}Effectuons une réduction de \textsc{Gauss}.

\begin{align*}\ensuremath
Q((x_1,...,x_5))&= x_1x_2+x_1(x_3+x_4+x_5)+x_2(x_3+x_4+x_5)+x_3x_4+x_3x_5+x_4x_5\\
 &= (x_1+x_3+x_4+x_5)(x_2+ x_3+x_4+x_5) - (x_3+x_4+x_5)^2+x_3x_4+x_3x_5+x_4x_5\\
 &=\frac{1}{4}(x_1+x_2+2x_3+2x_4+2x_5)^2-\frac{1}{4}(x_1-x_2)^2-x_3^2-x_4^2-x_5^2-x_3x_4-x_3x_5-x_4x_5\\
 &=\frac{1}{4}(x_1+x_2+2x_3+2x_4+2x_5)^2-\frac{1}{4}(x_1-x_2)^2-\left(x_3+\frac{1}{2}x_4+\frac{1}{2}x_5\right)^2-\frac{3}{4}x_4^2-\frac{1}{2}x_4x_5-\frac{3}{4}x_5^2\\
 &=\frac{1}{4}(x_1+x_2+2x_3+2x_4+2x_5)^2-\frac{1}{4}(x_1-x_2)^2-\left(x_3+\frac{1}{2}x_4+\frac{1}{2}x_5\right)^2 -\frac{3}{4}\left(x_4-\frac{1}{3}x_5\right)^2 -\frac{5}{6}x_5^2,
\end{align*}

et on retrouve $s = (1,4)$.

\item  $Q(x_1,...,x_n)=(x_1+...+x_n)^2$ et donc

\begin{center}
\shadowbox{
$r = 1$ et $s = (1,0)$.
}
\end{center}

\item  Pour $n\geqslant 2$, $Q((x_1,...,x_n)) =\left(\sum_{i=1}^{n}ix_i\right)\left(\sum_{j=1}^{n}x_j\right)=\frac{1}{4}\left(\sum_{i=1}^{n}(i+1)x_i\right)^2-\frac{1}{4}\left(\sum_{i=1}^{n}(i-1)x_i\right)^2$. Donc 

\begin{center}
\shadowbox{
$r = 2$ et $s = (1,1)$
}
\end{center}

car les deux formes linéaires $(x_1,...,x_n)\mapsto \sum_{i=1}^{n}(i+1)x_i$ et $(x_1,...,x_n)\mapsto\sum_{i=1}^{n}(i-1)x_i$ sont indépendantes pour $n\geqslant2$.

\item  Puisque la matrice de $Q$ dans la base canonique est $\left(
\begin{array}{cccccc}
1&1&1&\ldots&1&1\\
1&2&2&\ldots&2&2\\
1&2&3&\ldots&3&3\\
\vdots&\vdots&\vdots&\ddots&\vdots&\vdots\\
1&2&3& &n-1&n-1\\
1&2&3&\ldots&n-1&n
\end{array}
\right)$

\begin{align*}
Q((x1,...,xn))&=\sum_{1\leqslant i,j\leqslant n}^{}x_ix_j+\sum_{2\leqslant i,j\leqslant n}^{}x_ix_j+ ... +\sum_{n-1\leqslant i,j\leqslant n}^{}x_ix_j + x_n^2\\
  &= (x_1+...+x_n)^2+(x_2+...+x_n)^2+...+(x_{n-1}+x_n)^2 + x_n^2.
\end{align*}

\begin{center}
\shadowbox{
$Q$ est donc définie positive.
}
\end{center}
\end{enumerate}
\fincorrection
\correction{005811}
\begin{enumerate}
 \item  (Quand $x^2$ et $y^2$ ont les mêmes coefficients, penser à faire une rotation d'angle $\frac{\pi}{4}$) En posant $x=\frac{1}{\sqrt{2}}(X+Y)$ et $y =\frac{1}{\sqrt{2}}(X-Y)$, on obtient

\begin{center}
$x^2+10xy+y^2=\frac{1}{2}(X+Y)^2+5(X+Y)(X-Y)+\frac{1}{2}(X-Y)^2=6X^2-4Y^2$.
\end{center}

Ainsi, si on note $(i,j)$ la base canonique de $\Rr^2$ puis $e_1=\frac{1}{\sqrt{2}}(i+j)$ et $e_2=\frac{1}{\sqrt{2}}(i-j)$, on a

\begin{center}
$x^2+10xy+y^2=Q(xi+yj)=Q(Xe_1+Ye_2)=6X^2-4Y^2$.
\end{center}

\item  La matrice de $Q$ dans la base canonique $(i,j)$ de $\Rr^2$ est $A=\left(
\begin{array}{cc}
6&2\\
2&9
\end{array}
\right)$. Les deux nombres $5$ et $10$ ont une somme égale à $15 =\text{Tr}(A)$ et un produit égal à $50=\text{det}(A)$ et sont donc les valeurs propres de $A$. On sait alors que dans une base orthonormée $(e_1,e_2)$ de vecteurs propres de $A$ associée à la famille de valeurs propres $(5,10)$, on a $Q(Xe_1+Ye_2)=5X^2+10Y^2$. Déterminons une telle base.

$(A-5I_2)\left(
\begin{array}{c}
x\\
y
\end{array}
\right)=\left(
\begin{array}{c}
0\\
0
\end{array}
\right)\Leftrightarrow x+2y = 0$ et donc $\text{Ker}(A-5I_2) =\text{Vect}(e_1)$ où $e_1=\frac{1}{\sqrt{5}}(2,-1)$  puis $\text{Ker}(A-10I_2) =(\text{Ker}(A-5I_2))^\bot=\text{Vect}(e_2)$ où $e_2=\frac{1}{\sqrt{5}}(1,2)$.

Donc, si $e_1=\frac{1}{\sqrt{5}}(2,-1)$, $e_2=\frac{1}{\sqrt{5}}(1,2)$ et $u=xi+yj=Xe_1+Ye_2$, alors $q(u)=6x2+4xy+9y2= 5X^2 + 10Y^2$. De plus, $x=\frac{1}{\sqrt{5}}(2X+Y)$ et $y=\frac{1}{\sqrt{5}}(-X+2Y)$.

\item  La matrice de $Q$ dans la base canonique est $A=\left(
\begin{array}{ccc}
4&\sqrt{6}&5\sqrt{2}\\
\sqrt{6}&9&\sqrt{3}\\
5\sqrt{2}&\sqrt{3}&-1
\end{array}
\right)$.

\begin{align*}\ensuremath
\chi_A&=\left|
\begin{array}{ccc}
4-X&\sqrt{6}&5\sqrt{2}\\
\sqrt{6}&9-X&\sqrt{3}\\
5\sqrt{2}&\sqrt{3}&-1-X
\end{array}
\right|= (4-X)(X^2-8X-12)-\sqrt{6}(-\sqrt{6}X-6\sqrt{6})+5\sqrt{2}(5\sqrt{2}X-42\sqrt{2})\\
 &= -X^3+12X^2+36X-432 =-(X-6)(X+6)(X-12).
\end{align*}

Ensuite,

\begin{align*}\ensuremath
\left(
\begin{array}{c}
x\\
y\\
z
\end{array}
\right)\in\text{Ker}(A-6I_3)&\Leftrightarrow\left\{
\begin{array}{l}
-2x+\sqrt{6}y+5\sqrt{2}z=0\\
\sqrt{6}x+3y+\sqrt{3}z=0\\
5\sqrt{2}x+\sqrt{3}y-7z=0
\end{array}
\right.\Leftrightarrow\left\{
\begin{array}{l}
z=-\sqrt{2}x-\sqrt{3}y\\
-2x+\sqrt{6}y+5\sqrt{2}(-\sqrt{2}x-\sqrt{3}y)=0\\
5\sqrt{2}x+\sqrt{3}y-7(-\sqrt{2}x-\sqrt{3}y)=0
\end{array}
\right.\\
 &\Leftrightarrow\left\{
\begin{array}{l}
z=-\sqrt{2}x-\sqrt{3}y\\
-12x-4\sqrt{6}y=0\\
12\sqrt{2}x+8\sqrt{3}y=0
\end{array}
\right.\Leftrightarrow\left\{
\begin{array}{l}
y=-\sqrt{\frac{3}{2}}x\\
z=-\sqrt{2}x-\sqrt{3}(-\sqrt{\frac{3}{2}}x)
\end{array}
\right.\Leftrightarrow\left\{
\begin{array}{l}
y=-\sqrt{\frac{3}{2}}x\\
z=\frac{\sqrt{2}}{2}x
\end{array}
\right.\Leftrightarrow\left\{
\begin{array}{l}
x=\sqrt{2}z\\
y=-\sqrt{3}{z}
\end{array}
\right.
\end{align*}

et $\text{Ker}(A-6I_3) =\text{Vect}(e_1)$ où $e_1=\frac{1}{\sqrt{6}}(\sqrt{2},-\sqrt{3},1)$. De même,

   
\begin{align*}\ensuremath
\left(
\begin{array}{c}
x\\
y\\
z
\end{array}
\right)\in\text{Ker}(A+6I_3)&\Leftrightarrow\left\{
\begin{array}{l}
10x+\sqrt{6}y+5\sqrt{2}z=0\\
\sqrt{6}x+15y+\sqrt{3}z=0\\
5\sqrt{2}x+\sqrt{3}y+5z=0
\end{array}
\right.\Leftrightarrow\left\{
\begin{array}{l}
z=-\sqrt{2}x-5\sqrt{3}y\\
10x+\sqrt{6}y+5\sqrt{2}(-\sqrt{2}x-5\sqrt{3}y)=0\\
5\sqrt{2}x+\sqrt{3}y+5(-\sqrt{2}x-5\sqrt{3}y)=0
\end{array}
\right.\\
 &\Leftrightarrow\left\{
\begin{array}{l}
z=-\sqrt{2}x-\sqrt{3}y\\
y=0
\end{array}
\right.\Leftrightarrow\left\{
\begin{array}{l}
y=0\\
z=-\sqrt{2}x
\end{array}
\right.
\end{align*}

 
et $\text{Ker}(A+6I_3) =\text{Vect}(e_2)$ où $e_2=\frac{1}{\sqrt{3}}(1,0,-\sqrt{2})$.

Enfin $\text{Ker}(A-12I_3) =\text{Vect}(e_3)$ où 

\begin{center}
$e_3 =e_1\wedge e_2=\frac{1}{3\sqrt{2}}\left(
\begin{array}{c}
\sqrt{2}\\
-\sqrt{3}\\
1
\end{array}
\right)\wedge\left(
\begin{array}{c}
1\\
0\\
-\sqrt{2}
\end{array}
\right)=\frac{1}{3\sqrt{2}}\left(
\begin{array}{c}
\sqrt{6}\\
3\\
\sqrt{3}
\end{array}
\right)$.
\end{center}

Ainsi, si on pose $P=\left(
\begin{array}{ccc}
\frac{1}{\sqrt{3}}&\frac{1}{\sqrt{3}}&\frac{1}{\sqrt{3}}\\
-\frac{1}{\sqrt{2}}&0&\frac{1}{\sqrt{2}}\\
\frac{1}{\sqrt{6}}&-\frac{2}{\sqrt{6}}&\frac{1}{\sqrt{6}}
\end{array}\right)$ alors $PA{^t}P=\text{diag}(6,-6,12)$ ou encore

\begin{center}
$Q(Xe_1+Ye_2+Ze_3) =6X^2-6Y^2+12Z^2$ où $\text{Mat}_{(i,j,k)}(e_1,e_2,e_3)=\left(
\begin{array}{ccc}
\frac{1}{\sqrt{3}}&\frac{1}{\sqrt{3}}&\frac{1}{\sqrt{3}}\\
-\frac{1}{\sqrt{2}}&0&\frac{1}{\sqrt{2}}\\
\frac{1}{\sqrt{6}}&-\frac{2}{\sqrt{6}}&\frac{1}{\sqrt{6}}
\end{array}\right)$.
\end{center} 
\end{enumerate}
\fincorrection
\correction{005812}
\begin{enumerate}
 \item  Soit $P$ un élément de $E$. D'après un théorème de croissances comparées, $P(k)P(-k)e^{-k}\underset{k\rightarrow+\infty}{=}o\left(\frac{1}{k^2}\right)$ et donc $Q(P)$ existe.

Pour tout élément $P$ de $E$, $Q(P) = B(P,P)$ où $B$ est la forme bilinéaire symétrique définie sur $E$ par

\begin{center}
$\forall(P_1,P_2)\in E^2$, $B(P_1,P_2)=\frac{1}{2}\left(\sum_{k=0}^{+\infty}(P_1(k)P_2(-k)+ P_1(-k)P_2(k))e^{-k}\right)$
\end{center}

et donc $Q$ est une forme quadratique sur $E$.

\item  Soit $F$ le sous-espace vectoriel de $E$ dont les éléments sont les polynômes pairs et $G$ le sous-espace vectoriel de $E$ dont les éléments sont les polynômes impairs. $F$ et $G$ sont supplémentaires dans $E$.

Soit $P$ est un polynôme pair et non nul. Tout d'abord , $Q(P)\sum_{k=0}^{+\infty}(P(k))^2e^{-k}\geqslant0$. De plus, comme $P$ ne peut admettre tout entier naturel pour racine, on a plus précisément $Q(P)> 0$. De même, si $P$ est impair et non nul, $Q(P) < 0$.

Ainsi, la restriction de $Q$ à $F$ (resp. $G$) est définie positive (resp.négative). Enfin, si $P_1$ est pair et $P_2$ est impair, on a

\begin{center}
$B(P_1,P_2)=\sum_{k=0}^{+\infty}P_1(k)P_2(-k)e^{-k}=-\sum_{k=0}^{+\infty}P_1(k)P_2(k)e^{-k}=-\sum_{k=0}^{+\infty}P_1(-k)P_2(k)e^{-k}=-B(P_2,P_1)=-B(P_1,P_2)$,
\end{center}

et donc $B(P_1,P_2) = 0$ ($F$ et $G$ sont orthogonaux pour $B$).

Il existe un base de $F$ dans laquelle $Q_{/F}$ est combinaison linéaire à coefficients strictement positifs de carrés de formes linéaires linéairement indépendantes en nombre égal à $\text{dim}(F)$ et de même il existe un base de $G$ dans laquelle $Q_{/G}$ est combinaison linéaire à coefficients strictement négatifs de carrés de formes linéaires linéairement indépendantes en nombre égal à $\text{dim}(G)$. Maintenant, si $P$ est un polynôme quelconque de parties paire et impaire $P_1$ et $P_2$ respectivement,

\begin{center}
$Q(P)=Q(P_1+P_2) =Q(P_1)+2B(P_1,P_2)+Q(P_2) = Q_{/F}(P_1)+Q_{/G}(P_2)$.
\end{center}

Donc la réunion des deux bases ci-dessus est une base de $E$ dans laquelle $Q$ est combinaison linéaire de carrés de formes linéaires linéairement indépendantes dans laquelle $\text{dim}(F)=E\left(\frac{n}{2}\right)+1$ coefficients sont strictement positifs et 
$\text{dim}(G)=E\left(\frac{n+1}{2}\right)$ sont strictement négatifs. Finalement,

\begin{center}
\shadowbox{
$Q$ est donc non dégénérée de signature $s=\left(E\left(\frac{n}{2}\right)+1,E\left(\frac{n+1}{2}\right)\right)$.
}

\end{center}
\end{enumerate}
\fincorrection
\nocorrection
\correction{003695}
$\vec x = -\frac{\vec a\wedge\vec b}{\|\vec a\|^2} + \lambda \vec a$.
\fincorrection
\correction{003696}
$p = [\vec a,\vec b,\vec c]  \Rightarrow 
 \begin{cases} \vec u\wedge\vec v = p\vec b \cr
         \vec v\wedge\vec w = p\vec c \cr
         \vec w\wedge\vec u = p\vec a \cr\end{cases}$
et $[\vec u,\vec v,\vec w] = p^2$.

si $[\vec u,\vec v,\vec w] < 0$ : pas de solutions.

si $[\vec u,\vec v,\vec w] = 0$ et $\mathrm{rg}(\vec u,\vec v,\vec w) > 1$ :
   pas de solutions.

si $[\vec u,\vec v,\vec w] = 0$ et $\mathrm{rg}(\vec u,\vec v,\vec w) \le 1$ :
   une infinité de solutions.

si $[\vec u,\vec v,\vec w] > 0$ : 2 solutions.
\fincorrection
\correction{003697}
3. $G = \mathrm{tr}(F)I - {}^tF$.
\fincorrection
\nocorrection
\correction{003699}
$abc\sqrt{ 1 - \cos^2\alpha - \cos^2\beta - \cos^2\gamma
              + 2\cos\alpha\cos\beta\cos\gamma }$

$=abc\sqrt{ (\cos\gamma - \cos(\alpha+\beta))(\cos(\alpha-\beta) - \cos\gamma)}$.
\fincorrection
\nocorrection
\correction{003701}
\begin{enumerate}
  \item $f(\vec x) = (\vec x|\vec u)\vec u
                        + \cos\alpha(\vec u\wedge\vec x)\wedge\vec u
                        + \sin\alpha(\vec u\wedge\vec x)$.
  \item $M = (\cos\alpha)I
                + (1-\cos\alpha)\begin{pmatrix}a^2 &ab &ac \cr
                                         ab &b^2 &bc \cr
                                         ac &bc &c^2 \cr \end{pmatrix}
                + \sin\alpha\begin{pmatrix}0  &-c &b \cr
                                     c  &0  &-a \cr
                                     -b &a  &0 \cr\end{pmatrix}$.
\end{enumerate}
 \fincorrection
\nocorrection
\correction{003703}
Pour $f,g \in {\cal O}^+(\R^3)$, $f$ et $g$ ont la même matrice
         réduite dans une base orthonormée convenable, donc sont conjugués
         dans ${\cal O}(\R^3)$. $h$ n'est pas toujours positif car les bases
         peuvent ne pas avoir même orientation (ex~: deux rotations inverses).
         Pour $f,g \in {\cal O}^-(\R^3)$, considérer $-f$, $-g$.
\fincorrection
\correction{003704}
 3. $\pi/2,\pi/2,\pi$ \par $-\pi/2,\alpha,\pi/2$.
\fincorrection
\nocorrection
\correction{003706}
\begin{enumerate}
  \item 
  \item matrice dans une {\it bond}.
  \item $f = 2(h+\mathrm{id})^{-1} - \mathrm{id}$.
\end{enumerate}
 \fincorrection
\correction{003707}
\begin{enumerate}
  \item $f(\vec x) = \vec u\wedge\vec x$ avec
    $\vec u = (\alpha,\beta,\gamma)$.
  \item $\vec y = \frac{\vec x
                     + (\vec u\mid\vec x)\vec u-\vec u\wedge\vec x}
                     {1+\|\vec u\,\|^2}$.
  \item axe dirigé par $\vec u$,
             $\cos\theta = \frac{1-\|\vec u\,\|^2}{1+\|\vec u\,\|^2}$,
             $\sin\theta = \frac{-2\|\vec u\,\|}{1+\|\vec u\,\|^2}$.
\end{enumerate}
 \fincorrection
\correction{003708}
\begin{enumerate}
  \item 
  \item $g(\vec x) = (\cos\alpha)\vec x
                        + (1-\cos\alpha)\frac{(\vec a\mid\vec x)\vec a}{\alpha^2}
                        + \frac{\sin\alpha(\vec a\wedge\vec x)}\alpha$.
\end{enumerate}
 \fincorrection
\correction{003709}
\begin{enumerate}
  \item $\frac 17 \begin{pmatrix} 6 &3 &-2 \cr -2 &6 &3 \cr 3 &-2 &6 \cr\end{pmatrix}$.
  \item Rotation autour de $(1,1,1)$ d'angle $-\Arccos(\frac {11}{14})$.
\end{enumerate}
 \fincorrection
\correction{003710}
\begin{enumerate}
\item
On a
\[{}^tMM = I \Leftrightarrow (a^2+b^2+c^2=1 \text{ et } ab+bc+ca=0),\]
 et 
\[\det(M) = 1 \Leftrightarrow a^3+b^3+c^3-3abc=1.\]
En remarquant que 
\[ (a^3+b^3+b^3-3abc)  = (a+b+c)(a^2+b^2+c^2-ab-bc-ca), \] 
on en déduit (a)$\Leftrightarrow$ (b)$\Leftrightarrow$
($a+b+c = 1$ et $ab+ac+bc = 0$), d'où le polynôme $P$.

Or, $P$ a trois racines réelles si et seulement si
 $0 \le \lambda \le \frac 4{27}$(étude de la fonction associée), d'où l'équivalence avec la condition (c).
\item On voit sur la matrice que $(1,1,1)$ est propre pour la valeur propre $a+b+c$. Si $M \in SO(3)$, c'est donc (l'identité ou) une matrice de rotation suivant l'axe $\R (1,1,1)$. L'angle $\theta$ de la rotation vérifie alors $\mathrm{Tr}(M) = 2\cos(\theta)+1$ d'où $|\theta| = \Arccos\frac{3a-1}{2}$.
\item L'ensemble est celui de toutes les rotations d'axe $(1,1,1)$. C'est un groupe isomorphe à $\R/2\pi\Z$.
\end{enumerate}
\fincorrection
\correction{003711}
\begin{enumerate}
  \item rotation autour de $(1,0,1)$ d'angle $-\arccos(1/3)$.
  \item rotation autour de $(-3,1,1)$ d'angle $-\arccos(7/18)$.
  \item demi-tour autour de $(-1,-2,1)$.
  \item rotation autour de $(0,1,1)$ d'angle $2\pi/3$.
  \item rotation autour de $(-2-\sqrt3,1+\sqrt2,\sqrt2 - \sqrt3)$
             d'angle $\arccos( \frac {\sqrt6 - \sqrt2 + 1}{2\sqrt6} )$.
  \item symétrie par rapport à $x=y+z$.
  \item symétrie par rapport à $3x=2y-z$.
  \item symétrie par rapport à $x+2y-z=0$.
  \item symétrie-rotation autour de $(1,-3,1)$ d'angle $-\arccos(5/6)$.
  \item symétrie-rotation autour de $(1,-1,0)$ d'angle $\pi/3$.
  \item projection sur $2x+2y+z=0$ puis rotation d'angle $\arccos(3/5)$.
\end{enumerate}
\fincorrection
\correction{003712}
$M=\begin{pmatrix}0 &-1 &0\cr 0 &0 &-1\cr 1 &0 &0\cr\end{pmatrix}$, $\theta = \frac{2\pi}3$.
\fincorrection
\correction{005491}
\begin{align*}\ensuremath
[u\wedge v,v\wedge w,w\wedge u]&=((u\wedge v)\wedge(v\wedge w))|(w\wedge u)=(((u\wedge v)|w)v-((u\wedge v)|v)w)|(w\wedge u)\\
 &=(((u\wedge v)|w)v)|(w\wedge u)=((u\wedge v)w)\times(v|(w\wedge u))=[u,v,w][w,u,v]\\
 &=[u,v,w]^2.
\end{align*}
\fincorrection
\correction{005493}
$(u\wedge v)|(w\wedge s)=[u,v,w\wedge s]=[w\wedge s,u,v]=((w\wedge s)\wedge u)|v=((u|w)s-(u|s)w)|v=(u|w)(v|s)-(u|s)(v|w)$.
De même, $(u\wedge v)\wedge(w\wedge s)=((u\wedge v)|s)w-((u\wedge v)|w)s=[u,v,s]w-[u,v,w]s$.
\fincorrection
\correction{005793}
Posons $k =\frac{1}{\sqrt{2}}(e_1+e_2)$. Soit $u=xe_1+ye_2+ze_3\in E$. On sait que

\begin{align*}\ensuremath
r(u)&=(\cos\theta)u+(1-\cos\theta)\left(u|k\right)k+(\sin\theta)k\wedge u=\frac{1}{2}u +\frac{1}{4}(u|(e_1+e_2))(e_1+e_2)+\frac{\sqrt{6}}{4}(e_1+e_2)\wedge u\\
 &\text{\og}=\text{\fg}\;\frac{1}{2}\left(
 \begin{array}{c}
 x\\
 y\\
 z
 \end{array}
 \right)+\frac{1}{4}(x+y)\left(
 \begin{array}{c}
 1\\
 1\\
 0
 \end{array}
 \right)+\frac{\sqrt{6}}{4}\left(
 \begin{array}{c}
 1\\
 1\\
 0
 \end{array}
 \right)\wedge\left(
 \begin{array}{c}
 x\\
 y\\
 z
 \end{array}
 \right)\\
  &=\frac{1}{2}\left(
 \begin{array}{c}
 x\\
 y\\
 z
 \end{array}
 \right)+\frac{1}{4}(x+y)\left(
 \begin{array}{c}
 1\\
 1\\
 0
 \end{array}
 \right)+\frac{\sqrt{6}}{4}\left(
 \begin{array}{c}
 z\\
 -z\\
 -x+y
 \end{array}
 \right)\\
 &=\left(
 \begin{array}{c}
 \frac{3}{4}x+\frac{1}{4}y+\frac{\sqrt{6}}{4}z\\
\rule[-4mm]{0mm}{11mm}\frac{1}{4}x+\frac{3}{4}y-\frac{\sqrt{6}}{4}z\\
-\frac{\sqrt{6}}{4}x+\frac{\sqrt{6}}{4}y+\frac{1}{2}z
 \end{array}
 \right)=\left(
 \begin{array}{cccc}
 \frac{3}{4}&\frac{1}{4}&\frac{\sqrt{6}}{4}\\
\rule[-4mm]{0mm}{11mm}\frac{1}{4}&\frac{3}{4}&-\frac{\sqrt{6}}{4}\\
-\frac{\sqrt{6}}{4}&\frac{\sqrt{6}}{4}&\frac{1}{2}
 \end{array}
 \right)\left(
 \begin{array}{c}
x\\
y\\
z
 \end{array}
 \right)
\end{align*}

et la matrice cherchée est

\begin{center}
\shadowbox{
$M=\left(
 \begin{array}{cccc}
 \frac{3}{4}&\frac{1}{4}&\frac{\sqrt{6}}{4}\\
\rule[-4mm]{0mm}{11mm}\frac{1}{4}&\frac{3}{4}&-\frac{\sqrt{6}}{4}\\
-\frac{\sqrt{6}}{4}&\frac{\sqrt{6}}{4}&\frac{1}{2}
 \end{array}
 \right)
 $.
 }
 \end{center}
\fincorrection
\correction{005797}
Soit $(i,j,k)$ une base orthonormée directe de $\Rr^3$ euclidien orienté. Posons $u = f(i)$, $v = f(j)$ et $w = f(k)$. Nécessairement, $u\wedge v =f(i)\wedge f(j) = f(i\wedge j) = f(k) = w$ et de même $v\wedge w = u$ et $w\wedge u = v$.

\textbf{1er cas.} Si l'un des vecteurs $u$ ou $v$ ou $w$ est nul alors $u = v = w = 0$ et donc $f = 0$. Réciproquement, l'application nulle convient.

\textbf{2ème cas.} \textbullet~Si les trois vecteurs $u$, $v$ et $w$ sont non nuls alors $u\wedge v\neq 0$ et donc la famille $(u,v)$ est libre. Mais alors la famille $(u,v,w)$ est une base directe de $\Rr^3$.

Ensuite $w = u\wedge v$ est orthogonal à $u$ et $v$ et $v = w\wedge u$ est orthogonal à $u$. On en déduit que la famille $(u,v,w)$ est une base orthogonale directe de $\Rr^3$.

Enfin, puisque $u$ et $v$ sont orthogonaux, $\|w\|=\|u\wedge v\| =\|u\|\|v\|$ et de même $\|u\|=\|v\|\|w\|$ et $\|v\|=\|u\|\|w\|$. Puis $\|u\|\|v\|\|w\|= (\|u\|\|v\|\|w\|)^2$ et donc, puisque les vecteurs $u$, $v$ et $w$ sont non nuls, $\|u\|\|v\|\|w\| = 1$. Les égalités $\|u\|\|v\|\|w\| = 1$ et $\|u\|=\|v\|\|w\|$ fournissent $\|u\|^2 = 1$ et de même $\|v\|^2 =\|w\|^2 = 1$.

Finalement, la famille $(u,v,w)$ est une base orthonormée directe.

En résumé, l'image par $f$ d'une certaine base orthonormée directe de $\Rr^3$ est une base orthonormée directe de $\Rr^3$ et on sait que $f$ est un automorphisme orthogonal positif de $\Rr^3$ c'est-à-dire une rotation de $\Rr^3$.

\textbullet~Réciproquement, si $f$ est la rotation d'angle $\theta$ autour du vecteur unitaire $e_3$. On considère $e_1$ et $e_2$ deux vecteurs de $\Rr^3$ tels que la famille $(e_1,e_2,e_3)$ soit une base orthonormée directe.

Pour vérifier que $f$ est solution, par linéarité, il suffit de vérifier les $9$ égalités : $\forall(i,j)\in\{1,2,3\}^2$, $f(e_i\wedge e_j)=f(e_i)\wedge f(e_j)$. Pour vérifier ces $9$ égalités, il suffit se réduisent en fin de compte d'en vérifier $2$ :

\begin{center}
$f(e_1)\wedge f(e_2)= e_3 = f(e_3) = f(e_1\wedge e_2)$ et $f(e3)\wedge f(e1) = e_3\wedge f(e_1) = f(e_2) = f(e_3\wedge e_1)$.
\end{center}

Les endomorphismes cherchés sont donc l'application nulle et les rotations de $\Rr^3$.
\fincorrection
\correction{005801}
Si $a = 0$, $f = 0$ et il n'y a plus rien à dire.

Si $a\neq0$, puisque $f(a) = 0$, $0$ est valeur propre de $f$ et $\text{Vect}(a)\subset E_0(f)$.
D'autre part, si $x$ est orthogonal à $a$, d'après la formule du double produit vectoriel

\begin{center}
$f(x) = (a.x)a -\|a\|^2x = -\|a\|^2x$.
\end{center}

Donc le réel non nul $-\|a\|^2$ est valeur propre de $f$ et $a^\bot\subset E_{-\|a\|^2}$. Maintenant,  $\text{dim}\text{Vect}(a) +\text{dim}a^\bot= 3$ et donc $\text{Sp}(f)=(0,-\|a\|^2,-\|a\|^2)$ puis $E_0(f) = \text{Vect}(a)$ et  $E_{-\|a\|^2}= a^\bot$. On en déduit aussi que $f$ est diagonalisable. On peut noter que, puisque $f$ est diagonalisable et que les sous-espaces propres sont orthogonaux, $f$ est un endomorphisme symétrique.
\fincorrection
\correction{005803}
Les deux formes linéaires considérées sont indépendantes et donc $P$ est un plan. Une base de $P$ est par exemple $(i,j) = ((1,-1,0,0)(1,0,2,-3))$. On orthonormalise la base $(i,j)$.

On prend $e_1= \frac{1}{\sqrt{2}}(1,-1,0,0)$ puis $e_2'= j - (j|e_1)e_1 = (1,0,2,-3) -\frac{1}{2}(1,-1,0,0) = \frac{1}{2}(1,1,4,-6)$
puis $e_2=\frac{1}{3\sqrt{6}}(1,1,4,-6)$.

\begin{center}
Une base orthonormée de $P$ est $(e_1,e_2)$ où $e_1= \frac{1}{\sqrt{2}}(1,-1,0,0)$ et $e_2=\frac{1}{3\sqrt{6}}(1,1,4,-6)$.
\end{center}

\begin{enumerate}
 \item  Le projeté orthogonal de $u = (x,y,z,t)$ sur $P$ est

\begin{align*}\ensuremath 
p_P(u)&=(u|e_1)e_1 + (u|e_2)e_2  =\frac{1}{2}(x-y) (1,-1,0,0) +\frac{1}{54}(x+y+4z-6t) (1,1,4,-6)\\
 &= \frac{1}{27}(14x-13y+2z-3t,-13x+14y+2z-3t,2x+2y+8z-12t,-3x-3y-12z+18t).
\end{align*}

La matrice dans la base canonique de la projection orthogonale sur $P$ est

\begin{center}
\shadowbox{
$M=\frac{1}{27}\left(
\begin{array}{cccc}
14&-13&2&-3\\
-13&14&2&-3\\
2&2&8&-12\\
-3&-3&-12&18
\end{array}\right)$.
}
\end{center}

La matrice dans la base canonique de la symétrie orthogonale par rapport à $P$ est

\begin{center}
\shadowbox{
$S = 2M - I_4= \frac{1}{27}\left(
\begin{array}{cccc}
1&-26&4&-6\\
-26&1&4&-6\\
4&4&-11&-24\\
-6&-6&-24&9
\end{array}\right)$.
}
\end{center}

\item  La distance de u = (x,y,z,t) à P est

\begin{align*}\ensuremath
\|u -p_P(u)\|&=\frac{1}{27}\|(14x+13y-2z+3t,13x+14y-2z+3t,-2x-2y+19z+12t,3x+3y+12z+9t)\|\\
 &=\frac{1}{27}\sqrt{(14x+13y-2z+3t)^2+(13x+14y-2z+3t)^2+(-2x-2y+19z+12t)^2+(3x+3y+12z+9t)^2}.
\end{align*}  
\end{enumerate}
\fincorrection
\nocorrection
\nocorrection
\nocorrection
\nocorrection
\correction{003766}
\begin{enumerate}
  \item $P = \frac13\begin{pmatrix}2 &-1 &2 \cr 2 &2 &-1 \cr -1 &2 &2 \cr\end{pmatrix}$,
             $D = \text{Diag}(3\ 6\ 9)$.
  \item $P = \frac13\begin{pmatrix}2 &-1 &2 \cr 2 &2 &-1 \cr -1 &2 &2 \cr\end{pmatrix}$,
             $D = \text{Diag}(3\ 3\ 2)$.
\end{enumerate}
 \fincorrection
\correction{003767}
Si tous les $a_i$ sont nuls, $M = 0$.\par
         Sinon, $M = C^tC  \Rightarrow  E_0 = C^\perp$ et $E_\nu = \text{vect}(C)$
         avec $\nu = \|C\|^2$.
\fincorrection
\correction{003768}

         $M = 5\begin{pmatrix} \frac12 &0         &0         & \frac12 \cr
                         0         & \frac12 & \frac12 &0         \cr
                         0         &\frac12   &\frac12   &0         \cr
                         \frac12   &0         &0         &\frac12   \cr\end{pmatrix}
              -\begin{pmatrix}\frac12   &0         &0         &-\frac12  \cr
                         0         &\frac12   &-\frac12  &0         \cr
                         0         &-\frac12  &\frac12   &0         \cr
                         -\frac12&0           &0         &\frac12   \cr\end{pmatrix}$.
\fincorrection
\correction{003769}
\begin{enumerate}
  \item 
  \item $u$ est autoadjoint pour $<,>$.
  \item $P_0 = 1$, $P_2 = x$, $P_6 = 3x^2-1$, $P_{12} = 5x^3 - 3x$.
\end{enumerate}
 \fincorrection
\correction{003770}
\begin{enumerate}
  \item 
  \item 
  \item $\lambda_k = k(k+1)$.
\end{enumerate}
 \fincorrection
\nocorrection
\nocorrection
\correction{003773}
\begin{enumerate}
  \item 
  \item 
  \item $(p\circ q)_{|\Im p} = (p\circ q \circ p)_{|\Im p}$ est
             diagonalisable et $(p\circ q)_{|\mathrm{Ker} q + (\mathrm{Ker} p\cap \Im q)} = 0$
             donc tout vecteur de $E$ est somme de vecteurs propres pour
             $p\circ q$.
\end{enumerate}
 \fincorrection
\nocorrection
\nocorrection
\nocorrection
\nocorrection
\correction{003778}
\begin{enumerate}
  \item 
  \item Récurrence : pour $n=1$ c'est évident.\par
             $n-1\to n$ : $A = \begin{pmatrix} A' &C' \cr {}^tC' &\alpha \cr\end{pmatrix}$
             avec $A' = {}^tB'B'$.\par
             On cherche $B = \begin{pmatrix} B' &X' \cr 0 &x \cr\end{pmatrix}$ d'où :
             $X' = {}^tB'^{-1}C'$ et $x^2 = \alpha - {}^tX'X'
                 = \frac{\det A}{\det A'} > 0$.
\end{enumerate}
 \fincorrection
\nocorrection
\nocorrection
\nocorrection
\correction{003782}
\begin{enumerate}
  \item Soit $(\vec u_1,\dots,\vec u_n)$ une base propre pour $u$.
             On prend $\vec x = \vec u_1 + \dots + \vec u_n$.
  \item On norme $\vec x$ et on le complète en une base orthonormée.
             La matrice de $u$ dans cette base est symétrique, de trace nulle,
             et la diagonale commence par 0. On termine par récurrence.
\end{enumerate}
 \fincorrection
\nocorrection
\correction{003784}
$ABX = \lambda X  \Rightarrow  {}^tX^tBABX = \lambda {}^tXBX$.
\fincorrection
\correction{003785}
Se ramener au cas où $A$ est diagonale.
\fincorrection
\correction{003786}
Il existe $P$ inversible telle que $A = {}^tPP$ et $B = {}^tPB'P$
         avec $B'$ symétrique définie positive.

         Alors $A+B = {}^tP(I+B')P$ et
         $\det(I+B') = \prod(1+\beta_i) \ge 1 + \prod\beta_i  \Rightarrow $ cqfd.
\fincorrection
\correction{003787}
Soit ${\cal B}$ une BON fixée, $M = \text{Mat}_{\cal B}(f)$,
         ${\cal B}'$ la BON cherchée et $P$ la matrice de passage de $\cal B$
         à ${\cal B}'$. On veut que ${}^tM'M'$ soit diagonale avec
         $M' = {}^tPMP$, cad ${}^tP^tMMP$ diagonale.
\fincorrection
\nocorrection
\nocorrection
\correction{003790}
Soit
$(\vec h_i)$ une base diagonale pour $h$,
$H_i = \text{vect}\{\vec h_1, \dots, \vec h_i\}$
et $(\vec f_i)$, $F_i$ idem pour $f$.

Pour $\vec x \in F_k \cap H_{k-1}^\perp$,
$\lambda_k\|\vec x\,\|^2 + (\vec x\mid\vec x_0)^2 \le
 (h(\vec x)\mid\vec x) + (\vec x\mid\vec x_0)^2 =
 (f(\vec x)\mid\vec x) \le \mu_k\|\vec x\,\|^2$.

Pour $\vec x \in H_{k+1} \cap F_{k-1}^\perp \cap \vec x_0^\perp$,
$\mu_k\|\vec x\,\|^2 \le (f(\vec x)\mid\vec x) = (h(\vec x)\mid\vec x)
 \le \lambda_{k+1}\|\vec x\,\|^2$.
\fincorrection
\correction{003791}
\begin{enumerate}
  \item Si $f(x) + f^*(x) = 0$ alors
$f(x) \in \Im f \cap \Im f^* = \Im f \cap (\mathrm{Ker} f)^\bot
                             = \Im f \cap (\Im f)^\bot$
donc $f(x) = f^*(x) = 0$ et
$x\in \mathrm{Ker} f \cap \mathrm{Ker} f^* = \mathrm{Ker} f \cap (\mathrm{Ker} f)^\bot$.
  \item $f^2 = 0  \Rightarrow  \Im f \subset \mathrm{Ker} f$.\par
$f+f^* \in GL(E)  \Rightarrow  \Im f + \Im f^* = \Im f + (\mathrm{Ker} f)^\bot = E                  \Rightarrow  \dim \Im f \ge \dim \mathrm{Ker} f$.
\end{enumerate}
 \fincorrection
\nocorrection
\nocorrection
\nocorrection
\nocorrection
\correction{003796}
\begin{enumerate}
  \item $((u-u^*)(x)\mid x) = 0$.
  \item Orthodiagonaliser et appliquer l'inégalité de Cauchy-Schwarz.
\end{enumerate}
 \fincorrection
\correction{003797}
Soit $K = \sup\{\|u_0+\dots+u_n\|\}$ et~$x\in H$.
On note $v_{p,q} = \sum_{n=p}^q u_n$ pour $p\le q$.
La série $\sum (u_n(x)\mid x)$ est convergente (termes positifs, sommes
partielles majorées) donc elle vérifie le critère de Cauchy~:
$(v_{p,q}(x)\mid x) \to 0$ lorsque $p,q\to\infty$.

Comme $v_{p,q}$ est positif, il vérifie l'inégalité de Cauchy-Schwarz~:
$$|(v_{p,q}(x)\mid y)|^2 \le (v_{p,q}(x)\mid x)(v_{p,q}(y)\mid y) \le 2K\|y\|^2(v_{p,q}(x)\mid x).$$
En particulier pour $y=v_{p,q}(x)$ on obtient~: $\|v_{p,q}(x)\|^2 \le 2K(v_{p,q}(x)\mid x)$
donc la série $\sum u_n(x)$ est de Cauchy.


\emph{Remarque :} exemple où $\sum u_n$ ne converge pas dans~${\cal L}_c(H)$~:
$H = \ell^2(\N)$ et $u_n =$ projection orthogonale sur $<\!e_n\!>$ où $e_n(p) = \delta_{n,p}$.
$\sum u_n$ converge simplement et non uniformément vers l'identité.
\fincorrection
\correction{003798}
${}^t\!AA$ est $\R$-diagonalisable donc annule un polynôme
$P$ scindé à racines simples. $A$ annule le polynôme $P(X^3)$, donc est $\C$-diagonalisable
si $0$ n'est pas racine de~$P$ ce que l'on peut imposer si $A$ est inversible.

Si $A$ n'est pas inversible, soit $P(X) = XQ(X)$ avec $Q(0)\ne 0$.

On a $\R^n = \mathrm{Ker}(A^3)\oplus \mathrm{Ker}(Q(A^3))$ et $\mathrm{Ker}(A^3) = \mathrm{Ker}({}^t\!AA) = \mathrm{Ker}(A)$
donc $AQ(A^3)=0$ et $A$ est encore $\C$-diagonalisable.

Contre-exemple pour la $\R$-diagonalisabilité~: prendre une
rotation d'angle $2\pi/3$ dans le plan.
\fincorrection
\correction{003799}
\begin{enumerate}
  \item On se place dans une base propre pour~$u$, soient
    $U,V,W$ les matrices correspondantes avec $U = \mathrm{diag}(\lambda_i)$.
    On doit donc résoudre $(\lambda_i+\lambda_j)W_{ij} = V_{ij}$ d'où
    l'existence, l'unicité et la symétrie de~$w$.
  \item 
\begin{verbatim}
> A := matrix([[4,1,1],[1,4,-1],[1,-1,4]]); B := matrix([[0,0,-1],[0,0,1],[-1,1,3]]);
> eigenvals(A); eigenvects(A);
> P := transpose(matrix([[1, 0, 1], [1, 1, 0],[-1, 1, 1]]));
> A1 := evalm(P^(-1)&*A&*P); B1 := evalm(P^(-1)&*B&*P);
> C1 := matrix(3,3);
> for i from 1 to 3 do for j from 1 to 3 do C1[i,j] := B1[i,j]/(A1[i,i]+A1[j,j]) od od;
> C := evalm(P&*C1&*P^(-1)); evalm(A&*C+C&*A-B);
\end{verbatim}
$ \Rightarrow  C=\frac1{140}\begin{pmatrix}11&-11&-33\cr -11&11&33\cr -33&33&69\cr\end{pmatrix}$.
  \item Si $v$ est défini positif~: on a $(v(x)\mid x) = 2(u(x)\mid w(x))$ donc si $\lambda$ est
    une valeur propre de~$w$ et~$x$ est un vecteur propre associé, on a
    $\lambda = \frac{(v(x)\mid x)}{2(u(x)\mid x)} > 0$ d'où $w$ est défini positif.
    
    Cas $w$ défini positif et $v$ non positif~:
    $U=\begin{pmatrix}1&0\cr0&2\cr\end{pmatrix}$, $W=\begin{pmatrix}1&1\cr1&1+x\end{pmatrix}$,
    $V=\begin{pmatrix}2&3\cr3&4x+4\cr\end{pmatrix}$ avec $0<x<\frac18$.
\end{enumerate}
 \fincorrection
\correction{003800}
\begin{enumerate}
  \item c'est un endomorphisme autoadjoint positif de déterminant~$1$.
  \item $X^n\det\Bigl(\frac{\mathrm{id}}X+s^*\circ s\Bigr) = \det(\mathrm{id}+Xs^*\circ s)
             =\det(s^*\circ(\mathrm{id}+Xs^*\circ s)\circ s) = \det(s^*\circ s + X\mathrm{id})$.
  \item $s^*\circ s$ est diagonalisable avec des valeurs propres $(\lambda_i)$
    réelles positives deux à deux inverses pour la même multiplicité.
    $P^2(1) = \prod_{1\le i\le n}(1+\lambda_i)\Bigl(1+\frac1{\lambda_i}\Bigr)$
    et $(1+x)\Bigl(1+\frac1x\Bigr)\ge 4$ pour tout $x>0$ avec égalité ssi
    $x=1$.
    
    Si $P(1) = 2^n$ alors toutes les valeurs propres de~$s^*\circ s$ valent~$1$
    et $s^*\circ s$ est diagonalisable donc $s^*\circ s = \mathrm{id}$ et $s$ est une
    symétrie orthogonale. La réciproque est immédiate.
  \item Se ramener au cas $A_4=I$ puis calculer $\det A$ par pivotage.
\end{enumerate}
 \fincorrection
\correction{003801}
\begin{enumerate}
  \item Que c'est un espace préhilbertien.
  \item $g_x(t) = \min(t(1-x),x(1-t))$
  \item On note $g_i = g_{x_i}$~: $(g_1,\dots,g_n)$ est libre par
    considération des points anguleux, donc engendre un espace vectoriel $G$ de dimension~$n$.
    Soit $f\in P$~: $f = f_0+f_1$ avec $f_0\in G$ et $f_1\in G^\bot$.
    Alors $\varphi(f) = \varphi(f_0) + \|f_1\|^2$ donc $\varphi$ est minimale
    en~$f$ ssi $\varphi_{|G}$ est minimale en~$f_0$ et $f_1=0$.
    Désormais on suppose $f_1=0$ et~$f\in G$.
    
    L'application~: $$ u : G \to {\R^n}, f \mapsto {(f(x_1),\dots,f(x_n)) = ((f\mid g_1),\dots,(f\mid g_n))}$$
    est un isomorphisme linéaire. Soit $v$ l'endormophisme autoadjoint défini
    positif de $\R^n$ (pour le produit scalaire canonique) tel que~:
    $\forall\ t\in\R^n, (t\mid v(t)) = \|u^{-1}(t)\|^2$.

    On a donc en notant $\alpha = (\alpha_1,\dots,\alpha_n)$ et $\beta = (\mathrm{id}+v)^{-1}(\alpha)$~:
    \begin{align*}\forall\ t\in\R^n,\ \varphi(u^{-1}(t))
      &= (t\mid v(t)) + (t-\alpha\mid t-\alpha)\\
      &= (t\mid(\mathrm{id}+v)(t)) - 2(t\mid \alpha) + (\alpha\mid\alpha)\\
      &= (t-\beta\mid(\mathrm{id}+v)(t-\beta)) + (\alpha\mid\alpha-\beta).\\ \end{align*}
    $\mathrm{id}+v$ est autoadjoint défini positif donc le minimum de~$\varphi$
    est atteint pour $f=u^{-1}(\beta)$ (solution unique) et vaut
    $(\alpha\mid\alpha-\beta)$.
\end{enumerate}
 \fincorrection
\correction{003802}
\begin{enumerate}
  \item $p$ est un projecteur orthogonal $\Leftrightarrow$ $p$ est un projecteur et $p=p^*$ $\Leftrightarrow$ $p^*$ est un projecteur orthogonal.
  \item $p$ et $p^*$ commutent donc $\mathrm{Ker} p$ et $\Im p$ sont stables par
    $p$ et par $p^*$, d'où $p^*_{|\mathrm{Ker} p} = (p_{|\mathrm{Ker} p})^* = 0_{\mathrm{Ker} p}$
    et $p^*_{|\Im p} = (p_{|\Im p})^* = \mathrm{id}_{\Im p}$. Ainsi $p=p^*$ ce
    qui implique $\mathrm{Ker} p \bot \Im p$.
\end{enumerate}
 \fincorrection
\correction{003803}
\begin{enumerate}
  \item 
  \item On a pour~$f,g\in E$~: $u\circ v(f) = g \Leftrightarrow g$ est $\mathcal{C}^2$,
$g(0)=g'(1)=0$ et $g''=-f$.
En particulier $u\circ v$ est injectif, $0$ n'est pas valeur propre de~$u\circ v$.

Pour $\lambda\in\R^*$ et $f\in E$ on a $u\circ v(f) = \lambda f$
si et seulement si $f$ est de la forme $x \mapsto ae^{\alpha x} + be^{-\alpha x}$
avec $\alpha^2 = -\frac1\lambda$ et $a+b=a\alpha e^\alpha - b\alpha e^{-\alpha} = 0$.
On obtient $f\ne 0$ en prenant $a\ne 0$, $b=-a$ et $\alpha=i\pi(\frac12+k)$, $k\in\Z$.
Donc $\mathrm{Spec}(u\circ v) = \Bigl\{\frac{1}{\pi^2(\frac 12+k)^2},\ k\in\Z\Bigr\}$.
\end{enumerate}
 \fincorrection
\correction{003804}
\begin{enumerate}
  \item $u_1+\dots+u_p$ est l'endomorphisme autoadjoint associé à
         $q_1+\dots+q_p$.
  \item $\Im(u_1)+\dots+\Im(u_p) \supset\Im(u_1+\dots+u_p) = E$ et la somme
         des dimensions est égale à~$\dim E$ donc la somme des sous-espaces est directe.
  \item On a $\mathrm{Ker}(u_1) = \{x\in E\text{ tq }x = u_2(x) + \dots + u_p(x)\}
\subset \Im(u_2+\dots+u_p) = \Im(u_2)\oplus\dots\oplus\Im(u_p)$ et les deux
termes extrêmes ont même dimension, d'où $\mathrm{Ker}(u_1) = \Im(u_2)\oplus\dots\oplus\Im(u_p)$.
Comme $u_1$ est autoadjoint, $\Im(u_1) \perp \mathrm{Ker}(u_1)$ ce qui prouve
l'orthogonalité de la somme. De plus $\Im(u_1)\subset \mathrm{Ker}(u_j)$ pour $j\ge 1$
donc $q_1(x) = \|x\|^2$ pour tout~$x\in\Im(u_1)$. En appliquant {\bf 1)} à
$\Im(u_1)$ on obtient $u_1(x)=x$ pour tout~$x\in\Im(u_1)$ ce qui prouve
que~$u_1$ est un projecteur, et c'est un projecteur orthogonal car autoadjoint.
\end{enumerate}
 \fincorrection
\correction{003805}
$A=P^{-1}DP \Rightarrow {}^tA = ({}^tPP)A(P^{-1}{}^tP^{-1})$.\par

$S$ définie positive $ \Rightarrow  \exists\ P\in GL_n(\R)$ tq $S={}^tPP$,
donc ${}^tA = SA S^{-1}  \Rightarrow  {}^tA = {}^tPM{}^tP^{-1}$ avec $M = PAP^{-1}$,
d'où ${}^tM = M$ est diagonale.
\fincorrection
\correction{003806}
Pour $A$ symétrique réelle on a $\max(\mathrm{Sp}(A)) = \sup\{(x\mid Ax)/\|x\|^2,\ x\in\R^n\setminus\{0\}\}$
donc $f$ est la borne supérieure des fontions affines $t \mapsto((x\mid Ax) + t(x\mid Bx))/\|x\|^2$
lorsque $x$ décrit $\R^n\setminus\{0\}$. En tant que sup de fonctions convexes,
c'est une fonction convexe.
\fincorrection
\nocorrection
\nocorrection
\correction{003825}
\begin{enumerate}
  \item $(f(\vec x)\mid \vec y) = -(f(\vec y)\mid \vec x)$ et
             $(f(i\vec x)\mid \vec y) = -(f(\vec y)\mid i\vec x)$.
  \item 
  \item 
\end{enumerate}
 \fincorrection
\nocorrection
\correction{003827}
$a=\frac{20}{\pi^3}- \frac{320}{\pi^5}$,
         $b=\frac{240}{\pi^4}-\frac{12}{\pi^2}$,
         $m=\frac\pi2-\frac8\pi+\frac{160}{\pi^3}-\frac{1280}{\pi^5}$.
\fincorrection
\correction{003828}
\begin{enumerate}
  \item $\frac1{16}$.
  \item $\frac14$.
\end{enumerate}
 \fincorrection
\nocorrection
\nocorrection
\nocorrection
\correction{003832}
\begin{enumerate}
  \item 
  \item 
  \item 
  \item 
  \item $\frac1{12}$.
\end{enumerate}
 \fincorrection
\nocorrection
\correction{003834}
\begin{enumerate}
  \item 
  \item 
  \item $\alpha = \sqrt{1-\frac2{\sqrt7}}$,
             $\beta  = \sqrt{1+\frac2{\sqrt7}}$.
\end{enumerate}
 \fincorrection
\correction{003835}
\begin{enumerate}
  \item 
  \item Soit $P_0 = Q_0'$. Par IPP on obtient $Q_0$ est orthogonal à la
famille $(jX^{j-1}-X^j)_{j\ge 1}$ qui est une base de $\R[X]$ donc $Q_0 = 0 = P_0$
et $ \int_{t=0}^{+\infty} e^{-t}t^0P_0(t)\,d t \ne \delta_{0,0}$.
\end{enumerate}
 \fincorrection
\correction{003836}
$f = {}$ cste.
\fincorrection
\nocorrection
\correction{003838}
Soit $\cal B$ une base orthonormée de $E$ et $P$ la matrice de
passage de $\cal B$ à $(\vec e_i)$.

Le premier membre vaut ${}^tXPG^{-1}{}^tPX = {}^t XX$.
\fincorrection
\nocorrection
\correction{003840}
\begin{enumerate}
  \item 
  \item 
  \item Soit $g \in H^\bot$ non nulle. Les formes linéaires :
             $f  \mapsto  \int_{0}^{1/2} f$ et $f  \mapsto  \int_0^1 fg$ sont
             nulles sur $H$, donc proportionnelles, ce qui est impossible
             pour $g$ continue.
\end{enumerate}
 \fincorrection
\correction{003841}
\begin{enumerate}
  \item $u \ge 0$ et $u^{-1}(0)$ est d'intérieur vide.
  \item Il existe $\alpha,\beta > 0$ tels que $\alpha u \le v \le \beta u$.
\end{enumerate}
 \fincorrection
\correction{003842}
\begin{enumerate}
  \item $(a_n)$ est partout dense.
  \item 
  \item 
  \item Si les $a_n$ sont distincts, on choisit pour tout~$n$
une fonction $f_n$ comprise entre $0$ et~$1$ valant alternativement
$1$ et~$-1$ en $a_0,\dots,a_n$. Alors la suite $(f_n)$ est de Cauchy
mais ne converge pas car si $f_n\to f$ alors $f^2 \equiv 1$, absurde.
\end{enumerate}
 \fincorrection
\correction{003843}
$\sqrt{1-|(u\mid v)|^2} = d(v,\C u) \le d(v,\C w) + |(v\mid w)|d(w,\C u)$.
\fincorrection
\correction{003844}
\begin{enumerate}
  \item $A=UT\,^tU  \Rightarrow  A\overline A = UT\overline T\,{}^t\overline{U}$ est semblable à $T\overline T$.
  \item $A\overline A$ est à valeurs propres positives distinctes. Soit $U$ unitaire trigonalisant
$A\overline A$ et $T = U^{-1}A\,{}^tU^{-1}$. Donc $T\overline T$ est triangulaire supérieure
à valeurs propres réelles distinctes. On montre que ceci implique $T$ triangulaire
supérieure par récurrence sur~$n$.

$T = \begin{pmatrix}t&X\cr Y&Z\cr\end{pmatrix}
 \Rightarrow  T\overline T = \begin{pmatrix}|t|^2 + X\overline Y &t\overline X + X\overline Z\cr 
                            \overline tY + Z\overline Y &Y\overline X + Z\overline Z\cr\end{pmatrix}
                 = \begin{pmatrix}\text{réel}&*\cr 0&\text{tr.sup à vp réelles}\cr\end{pmatrix}$.


Donc $X\overline Y = \overline XY$ et $Z\overline Y = -\overline tY$, d'où
$(Y\overline X - X\overline YI)Y = 0 = (Z\overline Z - |t|^2I)Y$.

Par hypothèse $Y\overline X + Z\overline Z - (|t|^2 + X\overline Y)I$ est inversible
donc $Y=0$ et on est ramené au cas~$n-1$.
  \item $A\overline A$ est à valeurs propres positives~: ???

Solution de Pierre Février (MP$^*$ Neuilly sur Seine) :

{\bf lemme :} Si $\lambda \in \text{Sp} (A\overline{A})$ alors il existe $W\ne 0$ et $\alpha \in \R$ 
tels que $A\overline{W}=\alpha W$ et $\alpha ^2=\lambda$.

Soit $V \in E_{\lambda}(A\overline{A}),\ V\ne 0$. Si $A\overline{V}=-\sqrt{\lambda}\,V$ on a le 
r\'esultat voulu, sinon on pose $W=A\overline{V}+\sqrt{\lambda}\,V$. On a alors :
$$\overline{A}W=\overline{A}A\overline{V}+\sqrt{\lambda}\,\overline{A}V=\lambda \overline{V}+
\sqrt{\lambda}\,\overline{A}V=\sqrt{\lambda}\,\overline{W}$$
On peut s'arranger pour que le vecteur pr\'ec\'edent soit unitaire et construire $U$ matrice unitaire 
de premi\`ere colonne $W$. 

On a alors $^t\overline UA\overline U=\begin{pmatrix}
\alpha  & x & x & x\cr
 0      &   &   &  \cr
\vdots  &   & B &  \cr
 0      &   &   &  \cr\end{pmatrix}$, 
puis
$^t\overline UA\overline AU=\begin{pmatrix}
\alpha ^2 & y & y             & y\cr
0         &   &               &  \cr
\vdots    &   & B \overline B &  \cr
0         &   &               &  \cr\end{pmatrix}$.
On en d\'eduit le r\'esultat par r\'ecurrence.
\end{enumerate}
 \fincorrection
\correction{003807}
$M = X^tY - Y^tX$.

Soit $Z \in \mathcal{M}_{n,1}(\R)$ tq $(I + aM)Z = 0$.
Donc $Z \in \text{vect}(X,Y)$ : $Z = \lambda X + \mu Y$.

on remplace :
$$\left\{\begin{array}{lllll}
 (1-a{}^tYX)\lambda &-& a{}^tYY\mu      &=&0 \cr
 a{}^tXX\lambda     &+& (1+a{}^tYX)\mu  &=&0 \cr
\end{array}\right.$$

CNS $\iff a^2({}^tXX\,^tYY - ({}^tXY)^2) + 1 \ne 0$.
\fincorrection
\correction{003808}
On a ${}^tAA - {}^tCC = I_n$.

Soit $X$ tel que $AX = 0$. Donc ${}^tXX = -{}^t(CX)(CX)$, donc $X = 0$.
\fincorrection
\correction{003809}
$A^tA = (a^2+b^2+c^2+d^2)I$.
\fincorrection
\correction{003810}
trigonaliser $A$ dans une base orthonormée.
\fincorrection
\correction{003811}
$h = f\circ f^*$ : $h^2 = \text{id}$ et $h \ge 0  \Rightarrow  h = \text{id}$.
\fincorrection
\nocorrection
\correction{003813}
$a=b=\pm c$.
\fincorrection
\correction{003814}
Ils sont égaux (décomposer $A$ en symétrique $+$ antisymétrique).
\fincorrection
\correction{003815}
\begin{enumerate}
  \item $\mathrm{spec}(M) = \{j,j^2\}  \Rightarrow $ on prend comme base orthonormale
$a = \begin{pmatrix}1\cr0\cr\end{pmatrix}$ et $b = \frac2{\sqrt3}(f(a)+\frac12a) =
\begin{pmatrix}\sqrt3\cr-2\sqrt3\cr\end{pmatrix}$.
  \item $M$ est une matrice de rotation ssi $\mathrm{spec}(M)\subset \mathbb{U}\setminus\{\pm1\}$
ou $M=\pm I$.
  \item $M$ est la matrice d'une application orthogonale ssi $\mathrm{spec}(M)\subset \mathbb{U}$
et $M$ est $\C$-diagonalisable (alors $M$ est $\R$-semblable à une
matrice diagonale par blocs dont les blocs sont des matrices de rotation).
\end{enumerate}
 \fincorrection
\nocorrection
\correction{003817}
\begin{enumerate}
  \item Inégalité de Cauchy-Schwarz.
  \item Il existe $P$ orthogonale de même taille que $A$ telle
que $D = {}^tPAP$ est diagonale positive.

Alors $\begin{pmatrix}{}^tP&0\cr0&I\cr\end{pmatrix}U\begin{pmatrix}P&0\cr0&I\cr\end{pmatrix}
= \begin{pmatrix}D&{}^tPC\cr{}^tCP&B\cr\end{pmatrix}$ est symétrique positive
donc si $d_{ii} = 0$ alors la ligne $i$ de ${}^tPC$ est nulle.
Ainsi $\begin{pmatrix}{}^tP&0\cr0&I\cr\end{pmatrix}U'\begin{pmatrix}P&0\cr0&I\cr\end{pmatrix}
= \begin{pmatrix}D&{}^tPC\cr0&0\cr\end{pmatrix}$ est, après renumérotation éventuelle des
lignes et colonnes,
de la forme $U'' = \begin{pmatrix}D'&C'\cr0&0\cr\end{pmatrix}$ où $D'$ est diagonale inversible
et $U'$ est semblable à $U''$.
Enfin $U''$ est diagonalisable~:
$\begin{pmatrix}I&D'^{-1}C'\cr0&I\cr\end{pmatrix}U''\begin{pmatrix}I&-D'^{-1}C'\cr0&I\cr\end{pmatrix} = \begin{pmatrix}D'&0\cr0&0\cr\end{pmatrix}$.
\end{enumerate}
 \fincorrection
\correction{003818}
\begin{enumerate}
  \item Si $^tUMV = D$ est diagonale alors $^tMM = VD^2{}^tV$. Inversement,
    comme $^tMM$ est symétrique définie positive, il existe $D$ diagonale inversible
    et $V$ orthogonale telles que $^tMM = VD^2{}^tV$. On pose $M = UD^tV$ ce
    qui définit $U$ puisque $D^tV$ est inversible et on a
    $VD^2{}^tV = {}^tMM = VD^tUUD^tV$ d'où $^tUU = I$.
  \item $M$ est limite de matrices $M_k$ inversibles que l'on peut décomposer
    sous la forme $M_k = U_kD_k{}^tV_k$ avec $U_k$ et $V_k$ orthogonales et $D_k$ diagonale.
    Comme $O(n)$ est compact on peut supposer, quitte à extraire des sous-suites,
    que les suites $(U_k)$ et $(V_k)$ convergent vers $U,V$ orthogonales
    d'où $^tUMV = \lim_{k\to\infty} {}^tU_kM_kV_k = \lim_{k\to\infty}D_k = D$ diagonale.
  \item En diagonalisant $^tMM$ on trouve
    $V = \begin{pmatrix}\frac1{\sqrt2} & \phantom-\frac1{\sqrt6} & \frac1{\sqrt3} \cr
                   0              & \frac2{\sqrt6} &-\frac1{\sqrt3} \cr
                  -\frac1{\sqrt2} & \frac1{\sqrt6} & \frac1{\sqrt3} \cr\end{pmatrix}$,
    $D=\begin{pmatrix}\sqrt3 &0 &0\cr 0&\sqrt3 &0\cr 0&0&\phantom00\cr\end{pmatrix}$.
    Comme $D$ n'est pas inversible il faut ruser pour trouver $U$. On donne
    des coefficients indéterminés à $U$ et on écrit que $^tUMV = D$
    ce qui donne $U=\begin{pmatrix}a &b+\scriptstyle\sqrt2 &c\cr
                             -a-\frac3{\sqrt6} &-b-\frac1{\sqrt2} &\phantom0-c\cr
                              a &b &c\cr\end{pmatrix}$
    avec $a,b,c\in\R$. On choisit alors $a,b,c$ de sorte que $U\in O(3)$
    d'où, par exemple, $c=\frac1{\sqrt3}$, $a=-\frac1{\sqrt6}$, $b=-\frac1{\sqrt2}$
    et $U=\begin{pmatrix}-\frac1{\sqrt6} &\frac1{\sqrt2}  &  \frac1{\sqrt3} \cr
                    \frac2{\sqrt6} & 0              & -\frac1{\sqrt3} \cr
                   -\frac1{\sqrt6} &-\frac1{\sqrt2} &  \frac1{\sqrt3} \cr\end{pmatrix}$.
\end{enumerate}
 \fincorrection
\correction{003819}
\begin{enumerate}
  \item $\det(A)^2 = (-1)^n.$
  \item $A$ est $\C$-diagonalisable (annulateur simple) et ses valeurs propres
sont $i,-i$ avec la même multiplicité ($A$ est réelle). La matrice $A'$ donnée a
les mêmes propriétés donc $A$ et $A'$ sont $\C$-semblables à la même matrice
diagonale, et donc $\C$-semblables l'une à l'autre. Comme la $\C$-similitude
entre matrices réelles est équivalente à la $\R$-similitude (résultat bien connu),
$A$ et $A'$ sont $\R$-semblables.

  \item Soit $e_1$ unitaire et $e'_1 = Ae_1$. Alors $e'_1$ est unitaire et
$Ae'_1 = -e_1$ d'où $(e_1\mid e'_1) = (Ae_1\mid Ae'_1) = -(e_1\mid e'_1) = 0$
donc $(e_1,e'_1)$ est une famille orthonormale. Si $F_1$ est le sous-espace vectoriel engendré
par $(e_1,e'_1)$ alors $F_1^\bot$ est stable par $A$ donc on peut construire
par récurrence une base orthonormale $(e_1,\dots,e_{n/2},e'_1,\dots,e'_{n/2})$
telle que $Ae_i = e'_i$ et $Ae'_i = -e_i$.
\end{enumerate}
 \fincorrection
\correction{003820}
On remplace $A$ par $A + b_nI$ et $B$ par $B-b_nI$ ce
qui ne modifie pas $C$. Maintenant les valeurs propres de~$B$ sont
positves donc pour tou~$x\in\C^n$ on a $(Ax\mid x) \le (Cx\mid x)$.
Soit $(x_1,\dots,x_n)$ une base orthonormale propre pour~$A$
et $(y_1,\dots,y_n)$ une base orthonormale propre pour~$C$.
Si $z\in\mathrm{vect}(x_1,\dots x_i)$ alors $(Az\mid z) \ge a_i\|z\|^2$
et si $z\in\mathrm{vect}(y_i,\dots y_n)$ alors $(Az\mid z) \le (Cz\mid z) \le c_i\|z\|^2$.
Or $\mathrm{vect}(x_1,\dots x_i)$ et $\mathrm{vect}(y_i,\dots y_n)$ ont une intersection
non triviale (la somme des dimensions est égale à $n+1$) donc il existe
$z\ne 0$ tel que $a_i\|z\|^2 \le c_i\|z\|^2$ d'où $a_i\le c_i$.
\fincorrection
\correction{003821}
\begin{enumerate}
  \item Prendre $A$ supérieur ou égal à la plus petite des valeurs propres
     de~$-M$.
  \item Surjectivité de~$\phi$~: $\Im\Phi$ est un sous-espace vectoriel de~$S_n(\R)$ contenant
    $S_n^{++}(\R)$ donc contenant $\mathrm{vect}(S_n^{++}(\R)) = S_n(\R)$
    d'après la question précédente. On en déduit que $\phi$ est un isomorphisme
    grâce au théorème du rang.
    
    Si $M\in S_n^+(\R)$ alors $M = \lim_{p\to\infty}(M+I_n/p)$ donc
    $M\in\overline{S_n^{++}(\R)}$. Réciproquement, si $M\in\overline{S_n^{++}(\R)}$
    alors $M = \lim_{p\to\infty}(M_p)$ avec $M_p$ définie positive, donc pour
    tout $x\in\R^n$ on a ${}^txMx = \lim_{p\to\infty}({}^txM_px)\ge 0$, c'est-à-dire
    $M\in S_n^+(\R)$.
    Ainsi~: $\overline{S_n^{++}(\R)} = S_n^{+}(\R)$. Comme $\phi$ est continue
    (car linéaire en dimension finie) on en déduit $\phi(S_n^{+}(\R))\subset S_n^{+}(\R)$.
    De plus, $\phi(S_n^{++}(\R)) = S_n^{++}(\R) \Rightarrow S_n^{++}(\R) = \phi^{-1}(S_n^{++}(\R))$
    donc par continuité de $\phi^{-1}$~: $\phi^{-1}(S_n^{+}(\R))\subset S_n^{+}(\R)$,
    d'où $S_n^{+}(\R)\subset \phi(S_n^{+}(\R))$.\smallskip
  \item Soit $M\in S_2(\R)$ de valeurs propres $a,b$ avec $a\le b$,
    et soient $a'\le b'$ les valeurs propres de~$\phi(M)$.
    Pour tout $\lambda>b$ on a $M+\lambda I_2\in S_2^{++}(\R)$ donc
    $\phi(M)+\lambda I_2\in S_2^{++}(\R)$ c'est-à-dire $\lambda > b'$.
    Ceci prouve que $b'\le b$ et on montre l'égalité en considérant~$\phi^{-1}$.
    De même, en considérant $-M$ on montre que $a'=a$. Finalement
    $\chi_M = (X-a)(X-b) = \chi_{\phi(M)}$. De plus, $\det(M) = ab = \det(\phi(M))$.
    
    Remarque~: soient $A = \left(\begin{smallmatrix}1&0\cr0&0\cr\end{smallmatrix}\right)$,
                      $B = \left(\begin{smallmatrix}0&0\cr0&1\cr\end{smallmatrix}\right)$,
                      $C = \left(\begin{smallmatrix}0&1\cr1&0\cr\end{smallmatrix}\right)$,
    et $A'=\phi(A)$, $B'=\phi(B)$, $C'=\phi(C)$. On sait que $A'$ est
    orthodiagonalisable avec pour valeurs propres $0$ et $1$, donc il existe
    $P\in O(2)$ telle que $A' = {}^tPAP$. $A'+B'=\phi(I_2)=I_2$ d'où
    $B' = I_2-A' = {}^tPBP$. Posons $C' = {}^tP\left(\begin{smallmatrix}u&v\cr v&w\cr\end{smallmatrix}\right)P$.
    $0=\mathrm{tr}(C) = \mathrm{tr}(C') = u+w$ et $-1=\det(C) = \det(C') = uw-v^2$
    donc $w=-u$ et $u^2+v^2=1$. De plus,
    $-1=\det(A+C) = -u-u^2-v^2$ d'où $u=0$ et $v=\pm 1$.
    
    Si $v=1$ alors $C' = {}^tPCP$ et par linéarité, $\phi(M)={}^tPMP$
    pour toute $M\in S_2(\R)$.
    Si $v=-1$ on trouve de même $\phi(M) = {}^tQMQ$ avec $Q=P\left(\begin{smallmatrix}1&0\cr0&-1\cr\end{smallmatrix}\right)\in O(2)$.
    Réciproquement, toute aplication de la forme $M  \mapsto{}^tPMP$ avec $P\in O(2)$
    vérifie les hypothèses de la question. Les fonctions $\phi$
    linéaires vérifiant la seule condition $\phi(S_2^{++}(\R)) = S_2^{++}(\R)$
    sont les fonctions de la forme $M \mapsto{}^tPMP$ avec $P\in GL_2(\R)$
    (écrire $\phi(I_2) = {}^tTT$ puis considérer $M \mapsto{}^tT^{-1}\phi(M)T^{-1}$).
    
    Généralisation en dimension quelconque~?
\end{enumerate}
 \fincorrection
\correction{003822}
$M = (^tMM)^{-2}$ est symétrique définie positive, donc diagonalisable
en base orthonormale. En examinant la forme diagonale on trouve $M=I$.
\fincorrection
\correction{005489}
\begin{enumerate}
 \item  $\|C_1\|=\|C_2\|=\frac{1}{3}\sqrt{4+4+1}=1$ et $C_1|C_2=\frac{1}{9}(-2+4-2)=0$. Enfin, 

$$C_1\wedge C_2=\frac{1}{9}\left(
\begin{array}{c}
-2\\
2\\
1
\end{array}
\right)\wedge\left(
\begin{array}{c}
1\\
2\\
-2
\end{array}
\right)
=\frac{1}{9}\left(
\begin{array}{c}
-6\\
-3\\
-6
\end{array}
\right)=-\frac{1}{3}\left(
\begin{array}{c}
2\\
1\\
2
\end{array}
\right)=C_3.$$ 
Donc, $A\in O_3^+(\Rr)$ et $f$ est une rotation (distincte de l'identité).
\textbf{Axe de $f$.} Soit $X\in\mathcal{M}_{3,1}(\Rr)$.

$$AX=X\Leftrightarrow\left\{
\begin{array}{l}
-x-y-2z=0\\
-2x-5y-z=0\\
-x+2y-5z=0
\end{array}
\right.\Leftrightarrow\left\{
\begin{array}{l}
z=-2x-5y\\
3x+9y=0\\
9x+27y=0
\end{array}
\right.\Leftrightarrow\left\{
\begin{array}{l}
x=-3y\\
z=y
\end{array}
\right..$$ 
L'axe $D$ de $f$ est $\mbox{Vect}(\overrightarrow{u})$ où $\overrightarrow{u}=(-3,1,1)$. $D$ est dorénavant orienté par $\overrightarrow{u}$.
\textbf{Angle de $f$.} Le vecteur $\overrightarrow{v}=\frac{1}{\sqrt{2}}(0,1,-1)$ est un vecteur unitaire orthogonal à l'axe. Donc,

$$\cos\theta=\overrightarrow{v}.f(\overrightarrow{v})=\frac{1}{\sqrt{2}}(0,1,-1).\frac{1}{\sqrt{2}}\frac{-1}{3}(-1,1,-4)=-\frac{1}{6}\times5
=-\frac{5}{6},$$
et donc, $\theta=\pm\Arccos(-\frac{5}{6})\;(2\pi)$. (Si on sait que $\text{Tr}(A)=2\cos\theta+1$, c'est plus court : $2\cos\theta+1=\frac{2}{3}-\frac{2}{3}-\frac{2}{3}$ fournit $\cos\theta=-\frac{5}{6}$).
Le signe de $\sin\theta$ est le signe de $[\overrightarrow{i},f(\overrightarrow{i}),\overrightarrow{u}]=\left|
\begin{array}{ccc}
1&\frac{2}{3}&-3\\
0&-\frac{2}{3}&1\\
0&-\frac{1}{3}&1\\
\end{array}
\right|=-\frac{1}{3}<0$. Donc, 

\begin{center}
\shadowbox{
$f$ est la rotation d'angle $-\Arccos(-\frac{5}{6})$ autour de $u=(-3,1,1)$.
}
\end{center}
 \item  $||C_1||=||C_2||=\frac{1}{4}\sqrt{9+1+6}=1$ et $C_1|C_2=\frac{1}{16}(3+3-6)=0$. Enfin, 

$$C_1\wedge C_2=\frac{1}{16}\left(
\begin{array}{c}
3\\
1\\
-\sqrt{6}
\end{array}
\right)\wedge\left(
\begin{array}{c}
1\\
3\\
\sqrt{6}
\end{array}
\right)
=\frac{1}{16}\left(
\begin{array}{c}
4\sqrt{6}\\
-4\sqrt{6}\\
8
\end{array}
\right)=\frac{1}{4}\left(
\begin{array}{c}
\sqrt{6}\\
-\sqrt{6}\\
2
\end{array}
\right)
=C_3.$$
Donc, $A\in O_3^+(\Rr)$ et $f$ est une rotation.
\textbf{Axe de $f$.} Soit $X\in\mathcal{M}_{3,1}(\Rr)$.

$$AX=X\Leftrightarrow\left\{
\begin{array}{l}
-x+y+\sqrt{6}z=0\\
x-y-\sqrt{6}z=0\\
-\sqrt{6}x+\sqrt{6}y-2z=0
\end{array}
\right.\Leftrightarrow x-y=\sqrt{6}z=\frac{2}{\sqrt{6}}z\Leftrightarrow x=y\;\mbox{et}\;z=0.$$ 
L'axe $D$ de $f$ est $\mbox{Vect}(\overrightarrow{u})$ où $\overrightarrow{u}=(1,1,0)$. $D$ est dorénavant orienté par $\overrightarrow{u}$.
\textbf{Angle de $f$.} $\overrightarrow{k}=[0,0,1)$ est un vecteur unitaire orthogonal à $\overrightarrow{u}$. Par suite,

$$\cos\theta=\overrightarrow{k}.f(\overrightarrow{k})=(0,0,1).\frac{1}{4}(\sqrt{6},-\sqrt{6},2)=\frac{1}{2},$$
et donc $\cos\theta=\pm\frac{\pi}{3}\;(2\pi)$. Le signe de $\sin\theta$ est le signe de $\left[\overrightarrow{i},f(\overrightarrow{i}),\overrightarrow{u}\right]
=\left|
\begin{array}{ccc}
1&3/4&1\\
0&1/4&1\\
0&-\sqrt{6}/4&0
\end{array}
\right|=\frac{1}{\sqrt{6}}>0$. Donc, 

\begin{center}
\shadowbox{
$f$ est la rotation d'angle $\frac{\pi}{3}$ autour de $\overrightarrow{u}=(1,1,0)$.
}
\end{center}
 \item  $||C_1||=||C_2||=\frac{1}{9}\sqrt{64+16+1}=1$ et $C_1|C_2=\frac{1}{81}(8-16+8)=0$. Enfin,

$$C_1\wedge C_2=\frac{1}{81}\left(
\begin{array}{c}
8\\
-4\\
1
\end{array}
\right)\wedge\left(
\begin{array}{c}
1\\
4\\
8
\end{array}
\right)=\frac{1}{81}\left(
\begin{array}{c}
-36\\
-63\\
36
\end{array}
\right)=-\frac{1}{9}\left(
\begin{array}{c}
4\\
7\\
-4
\end{array}
\right)=-C_3.$$
Donc, $A\in O_3^-(\Rr)$. $A$ n'est pas symétrique, et donc $f$ n'est pas une réflexion. $f$ est donc la composée commutative $s\circ r$ d'une rotation d'angle $\theta$ autour d'un certain vecteur unitaire $\overrightarrow{u}$ et de la réflexion de plan $\overrightarrow{u}^\bot$ où $\overrightarrow{u}$ et $\theta$ sont à déterminer.
\textbf{Axe de $r$.} L'axe de $r$ est $Ker(f+Id_E)$ (car $f\neq-Id_E$).

$$AX=-X\Leftrightarrow\left\{
\begin{array}{l}
17x+y+4z=0\\
-4x+13y+7z=0\\
x+8y+5z=0
\end{array}
\right.\Leftrightarrow\left\{
\begin{array}{l}
y=-17x-4z\\
-225x-45z=0\\
-135x-27z=0
\end{array}
\right.\Leftrightarrow\left\{
\begin{array}{l}
z=-5x\\
y=3x
\end{array}\right.$$
$\mbox{Ker}(f+Id_E)=\mbox{Vect}(\overrightarrow{u})=D$ où $u=(1,3,-5)$. $D$ est dorénavant orienté par $\overrightarrow{u}$.
$s$ est la réflexion par rapport au plan $P=u^\bot$ dont une équation est $x+3y-5z=0$.
On écrit alors la matrice $S$ de $s$ dans la base de départ. On calcule $S^{-1}A=SA$ qui est la matrice de $r$ et on termine comme en 1) et 2).
\end{enumerate}
\fincorrection
\correction{005490}
Soit $f$ l'endomorphisme de $\Rr^3$ de matrice $M$ dans la base canonique de $\Rr^3$.

\begin{align*}\ensuremath
f\;\mbox{est une rotation}&\Leftrightarrow M\in O_3^+(\Rr)\Leftrightarrow||C_1||=||C_2||=||C_3||=1\;\mbox{et}\;C_1|C_2=C_1|C_3=C_2|C_3=0\;\mbox{et}\;\mbox{det}M=1\\
 &\Leftrightarrow a^2+b^2+c^2=1\;\mbox{et}\;ab+bc+ca=0\;\mbox{et}\;a^3+b^3+c^3-3abc=1.
\end{align*}
Posons $\sigma_1=a+b+c$, $\sigma_2=ab+bc+ca$ et $\sigma_3=abc$. On a $a^2+b^2+c^2=(a+b+c)^2-2(ab+ac+bc)=\sigma_1^2-2\sigma_2$.
Ensuite,

$$\sigma_1^3=(a+b+c)^3=a^3+b^3+c^3+3(a^2b+ba^2+a^2c+ca^2+b^2c+c^2b)+6abc,$$
et 

$$\sigma_1(\sigma_1^2-2\sigma_2)=(a+b+c)(a^2+b^2+c^2)=a^3+b^3+c^3+(a^2b+b^2a+a^2c+c^2a+b^2c+c^2b).$$
Donc, 

$$\sigma_1^3-3\sigma_1(\sigma_1^2-2\sigma_2)=-2(a^3+b^3+c^3)+6\sigma_3$$ 
et finalement, $a^3+b^3+c^3=\sigma_1^3-3\sigma_1\sigma_2+3\sigma_3$. 

\begin{align*}\ensuremath
M\in O_3^+(\Rr)&\Leftrightarrow\sigma_2=0\;\mbox{et}\;\sigma_1^2-2\sigma_2=1\;\mbox{et}\;\sigma_1^3-3\sigma_1\sigma_2=1\\
 &\Leftrightarrow\sigma_2=0\;\text{et}\;\sigma_1=1\\
 &\Leftrightarrow a,\;b\;\mbox{et}\;c\;\mbox{sont les solutions réelles d'une équation du type}\;x^3-x^2+k=0\;(\mbox{où}\;k=-\sigma_3).
\end{align*}
Posons $P(x)=x^3-x^2+k$ et donc $P'(x)=3x^2-2x=x(2x-3)$.
Sur $]-\infty,0]$, $P$ est strictement croissante, strictement décroissante sur $\left[0,\frac{3}{2}\right]$ et strictement croissante sur $\left[\frac{3}{2},+\infty\right[$. $P$ admet donc au plus une racine dans chacun de ces trois intervalles.
\textbf{1er cas.} Si $P(0)=k>0$ et $P\left(\frac{2}{3}\right)=k-\frac{4}{27}<0$ ou ce qui revient au même, $0<k<\frac{4}{27}$, P admet trois racines réelles deux à deux distinctes ($P$ étant d'autre part continue sur $\Rr$), nécessairement toutes simples.
\textbf{2ème cas.} Si $k\in\left\{0,\frac{4}{27}\right\}$, $P$ et $P'$ ont une racine réelle commune (à savoir $0$ ou $\frac{4}{27}$) et $P$ admet une racine réelle d'ordre au moins $2$. La troisième racine est alors nécessairement réelle.
\textbf{3ème cas.} Si $k<0$ ou $k>\frac{4}{27}$, P admet une racine réelle exactement. Celle-ci est nécessairement simple au vu du 2ème cas et donc $P$ admet deux autres racines non réelles.
En résumé, $P$ a toutes ses racines réelles si et seulement si $0\leq k\leq\frac{4}{27}$ et donc, $f$ est une rotation si et seulement si $a$, $b$ et $c$ sont les solutions d'une équation du type $x^3-x^2+k=0$ où $0\leq k\leq\frac{4}{27}$.
\fincorrection
\correction{005780}
\begin{enumerate}
 \item  Soit $F =\text{Vect}(x_1,...,x_n)$ et $m =\text{dim}F$. Soit $\mathcal{B}= (e_i)_{1\leqslant i\leqslant m}$ une base orthonormée de $F$ puis $M$ la matrice de la famille $(x_j)_{1\leqslant j\leqslant n}$ dans la base $\mathcal{B}$ .$M$ est une matrice rectangulaire de format $(m,n)$.

Soit $(i,j)\in\llbracket1,m\rrbracket\times\llbracket1,n\rrbracket$. Puisque la base $\mathcal{B}$ est orthonormée, le coefficient ligne $i$, colonne $j$ de la matrice ${^t}MM$ est

\begin{center}
$\sum_{k=1}^{m}m_{k,i}m_{k,j}=(x_i|x_j)$,
\end{center}

et on a donc

\begin{center}
$G(x_1,x_2,...,x_n) ={^t}MM$.
\end{center}

Puisque $\text{rg}(x_1,...,x_n) =\text{rg}M$, il s'agit de vérifier que $\text{rg}({^t}MM)=\text{rg}M$. Pour cela, montrons que les matrices $M$ et ${^t}MM$ ont même noyau.

Soit $X\in\mathcal{M}_{n,1}(\Rr)$. $X\in\text{Ker}M\Rightarrow MX = 0\Rightarrow{^t}MMX = 0\Rightarrow X\in\text{Ker}({^t}MM)$ et aussi

\begin{center}
$X\in\text{Ker}({^t}MM)\Rightarrow{^t}MMX = 0\Rightarrow{^t}X{^t}MMX = 0\Rightarrow{^t}(MX)MX = 0\Rightarrow\|MX\|_2^2 = 0\Rightarrow MX = 0\Rightarrow X\in\text{Ker}M$.
\end{center}

Finalement, $\text{Ker}({^t}MM) =\text{Ker}M$ et donc, d'après le théorème du rang, $\text{rg}(x_1,...,x_n) =\text{rg}M =\text{rg}({^t}MM) =\text{rg}(G(x_1,x_2,...,x_n))$.

\begin{center}
\shadowbox{
$\text{rg}(G(x_1,x_2,...,x_n))=\text{rg}(x_1,\ldots,x_n)$.
}
\end{center}

\item  D'après 1), 

\begin{align*}\ensuremath
(x_1,...,x_n)\;\text{liée}&\Leftrightarrow \text{rg}(x_1,x_2,...,x_n)<\Leftrightarrow \text{rg}G(x_1,x_2,...,x_n)<n\Leftrightarrow G(x_1,x_2,...,x_n)\notin\mathcal{GL}_n(\Rr)\\
 &\Leftrightarrow\gamma(x_1,x_2,...,x_n) = 0.
\end{align*}

De plus, quand la famille $(x_1,x_2,...,x_n)$ libre, avec les notations de la question 1), on a $m = n$ et la matrice $M$ est une matrice carrée. On peut donc écrire

\begin{center} 
$\gamma(x_1,x_2,...,x_n) =\text{det}({^t}MM)=\text{det}({^t}M)\times\text{det}(M)= (\text{det}M)^2 > 0$.
\end{center}

\begin{center}
\shadowbox{
\begin{tabular}{c}
$(x_1,...,x_n)\;\text{liée}\Leftrightarrow\gamma(x_1,\ldots,x_n)=0$\\
$(x_1,...,x_n)\;\text{libre}\Leftrightarrow\gamma(x_1,\ldots,x_n)>0$.
\end{tabular}
}
\end{center}

\item  \textbf{1ère solution.} Soit $x$ un vecteur de $E$ et $p_F(x)$ son projeté orthogonal sur $F$. Dans la première colonne de $\gamma(x,x_1,\ldots,x_n)$, le théorème de \textsc{Pythagore} permet d'écrire (puisque $x-p_F(x)\in F^\bot$)

\begin{align*}\ensuremath
\left(
\begin{array}{c}
(x|x)\\
(x|x_1)\\
\vdots\\
(x|x_n)
\end{array}
\right)&=\left(
\begin{array}{c}
\|x-p_F(x)+p_F(x)\|^2\\
(x-p_F(x)+p_F(x)|x_1)\\
\vdots\\
(x-p_F(x)+p_F(x)|x_n)
\end{array}
\right)=\left(
\begin{array}{c}
\|x-p_F(x)\|^2+\|p_F(x)\|^2\\
(p_F(x)|x_1)\\
\vdots\\
(p_F(x)|x_n)
\end{array}
\right)\\
 &=\left(
\begin{array}{c}
\|x-p_F(x)\|^2\\
0\\
\vdots\\
0
\end{array}
\right)+\left(
\begin{array}{c}
(p_F(x)|p_F(x))\\
(p_F(x)|x_1)\\
\vdots\\
(p_F(x)|x_n)
\end{array}
\right)
\end{align*}

Après avoir remplacé aussi en première ligne les $(x|x_i)$ par $(p_F(x)|x_i)$, on obtient par linéarité par rapport à la première colonne

\begin{center}
$\gamma(x,x_1,x_2,...,x_n) =\gamma(x-p_F(x),x_1,x_2,...,x_n) +\gamma(p_F(x),x_1,x_2,...,x_n)$
\end{center}

Maintenant, $p_F(x)$ est dans $F$ et donc la famille $(p_F(x),x_1,x_2,...,x_n)$ est liée puis d'après la question 2) $\gamma(p_F(x),x_1,x_2,...,x_n) = 0$. Il reste $\gamma(x,x_1,x_2,...,x_n) =\gamma(x-p_F(x),x_1,x_2,...,x_n)$ et en développant suivant la première colonne, on obtient

\begin{center}
$\forall x\in E,\;\gamma(x,x_1,\ldots,x_n)=\gamma(x- p_F(x),x_1,x_2,...,x_n) =\|x- p_F(x)\|^2\gamma(x_1,x_2,...,x_n)$.
\end{center}

Finalement 

\begin{center}
\shadowbox{
$\|x- p_F(x)\|=\sqrt{\frac{\gamma(x,x_1,x_2,...,x_n)}{\gamma(x_1,x_2,...,x_n)}}$.
}
\end{center}

\textbf{2ème solution.} Posons $p_F(x) =\sum_{i=1}^{n}\lambda_ix_i$ puis  $d =\|x-p_F(x)\|$ de sorte que

\begin{center}
$d^2 =(x-p_F(x))|(x-p_F(x)) =(x-p_F(x))|x =\|x\|^2 -(x|p_F(x))$.
\end{center}

D'autre part, pour chaque $i\in\llbracket1,n\rrbracket$, $x|x_i=(x-p_F(x)|x_I)+(p_F(x)|x_i)=(p_F(x)|x_i)$. Par suite, les $n+1$ réels $d^2$, $\lambda_1$,..., $\lambda_n$ sont solutions du système d'équations linéaires
 
\begin{center}
$\left\{
\begin{array}{ll}
d^2+&\lambda_1(x|x_1)+\ldots+\lambda_n(x|x_n)=\|x\|^2\\
 &\lambda_1(x_1|x_1)+\ldots+\lambda_n(x_1|x_n)=(x|x_1)\\
 &\rule{1,5cm}{0mm}\vdots\\
 &\lambda_1(x_n|x_1)+\ldots+\lambda_n(x_n|x_n)=(x|x_n)
\end{array}
\right.$
\end{center}

Le déterminant de ce système vaut $\gamma(x_1,x_2,...,x_n)>0$ et le système est de \textsc{Cramer}. Le déterminant associé à $d^2$ est $\gamma(x,x_1,x_2,...,x_n)$ et les formules de \textsc{Cramer} refournissent

\begin{center}
$d^2=\frac{\gamma(x,x_1,\ldots,x_n)}{\gamma(x_1,\ldots,x_n)}$.
\end{center}
\end{enumerate}
\fincorrection
\correction{005786}
La matrice $H_n$ est symétrique réelle. Soit $X=(x_i)_{1\leqslant i\leqslant n}\in\mathcal{M}_{n,1}(\Rr)$.

\begin{align*}\ensuremath
{^t}X H_n X&=\sum_{1\leqslant i,j\leqslant n}^{}\frac{x_ix_j}{i+j-1}=\sum_{1\leqslant i,j\leqslant n}^{}x_ix_j\int_{0}^{1}t^{i+j-2}\;dt =\int_{0}^{1}\left(\sum_{1\leqslant i,j\leqslant n}^{}x_ix_j t^{i+j-2}\right)dt\\
 &=\int_{0}^{1}\left(\sum_{i=1}^{n}x_it^{i-1}\right)^2dt\geqslant0 .
\end{align*}

De plus, si $X\neq0$, le polynôme $\sum_{i=1}^{n}x_iY^{i-1}$ n'est pas le polynôme nul et donc, puisqu'un polynôme non nul admet un nombre fini de racines, la fonction $t\mapsto\left(\sum_{i=1}^{n}x_it^{i-1}\right)^2$. Ainsi, la fonction $t\mapsto\left(\sum_{i=1}^{n}x_it^{i-1}\right)^2$ est continue positive et non nulle sur $[0,1]$ et on en déduit que $\int_{0}^{1}\left(\sum_{i=1}^{n}x_it^{i-1}\right)^2\;dt>0$. On a montré que $\forall X\in\mathcal{M}_{n,1}(\Rr)\setminus\{0\}$, ${^t}X H_n X > 0$ et donc que

\begin{center}
\shadowbox{
la matrice $H_n$ est symétrique définie positive.
}
\end{center}
\fincorrection
\correction{005787}
\begin{enumerate}
 \item  ${^t}S={^t}({^t}AA)={^t}A{^t}({^t}A) ={^t}AA = S$. Donc $S\in\mathcal{S}_n(\Rr)$.

Soit $X\in\mathcal{M}_{n,1}(\Rr)$, ${^t}XSX ={^t}X{^t}AAX ={^t}(AX)AX =\|AX\|_2^2\geqslant0$. Donc $S\in\mathcal{S}_n^+(\Rr)$.

\begin{center}
\shadowbox{
$\forall A\in\mathcal{M}_n(\Rr)$, ${^t}AA\in\mathcal{S}_n^+(\Rr)$.
}
\end{center}

\item  Soit $S\in\mathcal{S}_n^+(\Rr)$. D'après le théorème spectral, il existe $P$ dans $O_n(\Rr)$ et $D$ dans $\mathcal{D}_n(\Rr)$ telles que $S = PD{^t}P$.

Posons $D =\text{diag}(\lambda_1,...,\lambda_n)$. Puisque $S$ est dans $\mathcal{S}_n^+(\Rr)$, $D$ est dans $\mathcal{D}_n^+(\Rr)$ et on peut poser $D'=\text{diag}(\sqrt{\lambda_1},...,\sqrt{\lambda_n})$ de sorte que $D'^2 = D$.
On peut alors écrire

\begin{center}
$S = PD{^t}P =PD'D'{^t}P ={^t}(D{^t}P) D'{^t}P$,
\end{center}

et la matrice $A = D'{t}P$ convient.

\begin{center}
\shadowbox{
$\forall S\in\mathcal{S}_n^+(\Rr)$, $\exists A\in\mathcal{M}_n(\Rr)/$ $S={^t}AA$.
}
\end{center}

On a aussi ${^t}(-A)(-A) = S$ et comme en général  $-A\neq A$, on n'a pas l'unicité de la matrice $A$.

\item  

\begin{align*}\ensuremath
S\;\text{définie positive}&\Leftrightarrow \forall X\in\mathcal{M}_{n,1}(\Rr)\setminus\{0\},\;{^t}XSX> 0\Leftrightarrow \forall X\in\mathcal{M}_{n,1}(\Rr)\setminus\{0\},\;\|AX\|_2^2 > 0\\
 &\Leftrightarrow \forall X\in\mathcal{M}_{n,1}(\Rr)\setminus\{0\},\;AX\neq0
\Leftrightarrow\text{Ker} A =\{0\}\Leftrightarrow A\in\mathcal{GL}_n(\Rr).
\end{align*}

\item  Montrons que les matrices $A$ et $S$ ont même noyau. Soit $X\in\mathcal{M}_{n,1}(\Rr)$. 

\begin{center}
$X\in\text{Ker}A\Rightarrow AX = 0\Rightarrow{^t}AAX = 0\Rightarrow SX = 0\Rightarrow X\in\text{Ker}S$,
\end{center}

et

\begin{center}
$X\in\text{Ker}S\Rightarrow{^t}AAX = 0\Rightarrow{^t}X{^t}AAX = 0\Rightarrow{^t}(AX)AX = 0\Rightarrow\|AX\|_2^2 = 0\Rightarrow AX = 0\Rightarrow X\in\text{Ker}A$.
\end{center}

Ainsi, $\text{Ker}({^t}AA)=\text{Ker}(A)$ et en particulier, grâce au théorème du rang, on a montré que

\begin{center}
\shadowbox{
$\forall A\in\mathcal{M}_n(\Rr)$, $\text{rg}({^t}AA)=\text{rg}(A)$.
}
\end{center}

\item  Soit $S\in\mathcal{S}_n^+(\Rr)$.

\textbf{Existence.} D'après le théorème spectral, il existe $P_0\in O_n(\Rr)$ et $D_0\in\mathcal{D}_n^+(\Rr)$ telles que $S = P_0D_0{^t}P_0$.

Posons $D_0 =\text{diag}(\lambda_1,...,\lambda_n)$ où les $\lambda_i$, $1\leqslant i\leqslant n$, sont des réels positifs puis $\Delta_0 =\text{diag}(\sqrt{\lambda_1},...,\sqrt{\lambda_n})$ et enfin 
$R =P_0\Delta_0{^t}P_0$. La matrice $R$ est orthogonalement semblable à une matrice de $\mathcal{D}_n^+(\Rr)$ et est donc un élément de $\mathcal{S}_n^+(\Rr)$. Puis

\begin{center}
$R^2 = P_0\Delta_0^2{^t}P_0 = P_0D_0{^t}P_0 = S$.
\end{center}

\textbf{Unicité.} Soit $M$ un élément de $\mathcal{S}_n^+(\Rr)$ telle que $M^2 = S$.

$M$ est diagonalisable d'après le théorème spectral et donc $\mathcal{M}_{n,1}(\Rr)=\underset{\lambda\in\text{Sp}(M)}{\oplus}E_M(\lambda)$. Mais si $\lambda$ est une valeur propre de $M$, $\text{Ker}(M-\lambda I_n)\subset\text{Ker}(M^2-\lambda^2I_n)=\text{Ker}(S-\lambda^2I_n)$. De plus, les valeurs propres de $M$ étant positive, les $\lambda^2$, $\lambda\in\text{Sp}(M)$, sont deux à deux distincts ou encore les $\text{Ker}(S-\lambda^2I_n)$, $\lambda\in\text{Sp}(M)$, sont deux à deux distincts.

Ceci montre que pour chaque $\lambda\in\text{Sp}(M)$, $\text{Ker}(M-\lambda I_n)=\text{Ker}(S-\lambda^2I_n)$ et que les $\lambda^2$, $\lambda\in\text{Sp}(M)$, sont toutes les valeurs propres de $S$.

Ainsi, nécessairement la matrice ${^t}P_0MP_0$ est une matrice diagonale $D$. L'égalité $M^2=S$ fournit $D^2=D_0$ puis $D=\Delta_0$ (car $D\in\mathcal{D}_n^+(\Rr)$) et finalement $M = R$.

\begin{center}
\shadowbox{
$\forall S\in\mathcal{S}_n^+(\Rr)$, $\exists!R\in\mathcal{S}_n^+(\Rr)/\;R^2=S$.
}
\end{center}
\end{enumerate}
\fincorrection
\correction{005791}
Soit $A =(a_{i,j})_{1\leqslant i,j\leqslant n}$ une matrice orthogonale. On pose $X =\left(
\begin{array}{c}
1\\
\vdots\\
1
\end{array}
\right)\in\mathcal{M}_{n,1}(\Rr)$.

D'après l'inégalité de \textsc{Cauchy}-\textsc{Schwarz}

\begin{align*}\ensuremath
\left|\sum_{1\leqslant i,j\leqslant n}^{}a_{i,j}\right|&=\left|\sum_{1\leqslant i,j\leqslant n}^{}1\times a_{i,j}\times1\right|=\left|{^t}XAX\right| =\left|\left(AX|X\right)\right|\\
 &\leqslant\|AX\|\|X\|\;(\text{d'après l'inégalité de \textsc{Cauchy}-\textsc{Schwarz}})\\
 &=\|X\|^2\;(\text{puisque la matrice}\;A\;\text{est orthogonale})\\
 &=n.
\end{align*}
	        

On a l'égalité si et seulement si la famille $(X,AX)$ est liée ce qui équivaut à $X$ vecteur propre de $A$.

On sait que les valeurs propres (réelles) de $A$ ne peuvent être que $1$ ou $-1$. Donc, 

\begin{center}
égalité $\Leftrightarrow AX = X\;\text{ou}\;AX = -X\Leftrightarrow \forall i\in\llbracket1,n\rrbracket,\;\left|\sum_{j=1}^{n}a_{i,j}\right|= 1$
\end{center}

Il paraît difficile d'améliorer ce résultat dans le cas général. Supposons de plus que $\forall(i,j)\in\llbracket1,n\rrbracket^2$, $a_{i,j}\geqslant 0$. Soit $i\in\llbracket1,n\rrbracket$. Puisque tous les $a_{i,j}$ sont éléments de $[0,1]$,

\begin{center}
$1=\sum_{j=1}^{n}a_{i,j}\geqslant\sum_{j=1}^{n}a_{i,j}^2=1$.
\end{center}

L'inégalité écrite est donc une égalité et on en déduit que chaque inégalité $a_{i,j}\geqslant a_{i,j}^2$, $1\leqslant j\leqslant n$, est une égalité. Par suite, $\forall(i,j)\in\llbracket1,n\rrbracket^2$, $a_{i,j}\in\{0,1\}$. Ceci montre que la matrice $A$ est une matrice de permutation qui réciproquement convient.

\fincorrection
\correction{005794}
La matrice $A$ est symétrique réelle positive. Donc ses valeurs propres $\lambda_1$,..., $\lambda_n$ sont des réels positifs. De plus,

\begin{center}
$\text{det}A =\lambda_1...\lambda_n$ et $\text{det}(I_n+A) =\chi_A(-1) =(1+\lambda_1)...(1+\lambda_n)$.
\end{center}

L'inégalité à démontrer équivaut donc à :

\begin{center}
$\forall(\lambda_1,...,\lambda_n)\in(\Rr^+)^n$, $1+\sqrt[n]{\prod_{k=1}^{n}\lambda_k}\leqslant\sqrt[n]{\prod_{k=1}^{n}(1+\lambda_k)}$.
\end{center}

Soit donc $(\lambda_1,...,\lambda_n)\in(\Rr^+)^n$. Si l'un des $\lambda_k$ est nul, l'inégalité est immédiate.

Supposons dorénavant tous les $\lambda_k$ strictement positifs. L'inégalité à démontrer s'écrit

\begin{center}
$\ln\left(1+\text{exp}\left(\frac{1}{n}\left(\ln(\lambda_1)+...+\ln(\lambda_n)\right)\right)\right)\leqslant\frac{1}{n}\left(\ln(1+\text{exp}(\ln(\lambda_1)))+ ... +\ln(1+\text{exp}(\ln(\lambda_n)))\right)$\quad$(*)$
\end{center}

ou encore $f\left(\frac{1}{n}(x_1+...+x_n)\right)\leqslant\frac{1}{n}(f(x_1)+...+f(x_n))$ où $\forall x\in\Rr$, $f(x) =\ln(1+e^x)$ et $\forall k\in\llbracket1,n\rrbracket$, $x_k =\ln(\lambda_k)$.

L'inégalité à démontrer est une inégalité de convexité. La fonction $f$ est deux fois dérivable sur $\Rr$ et pour tout réel $x$,

\begin{center}
$f'(x)=\frac{e^x}{e^x+1}=1-\frac{1}{e^x+1}$ puis $f''(x) =\frac{e^x}{(e^x+1)^2}\geqslant0$.
\end{center}

La fonction $f$ est donc convexe sur $\Rr$ ce qui démontre l'inégalité $(*)$.

\begin{center}
\shadowbox{
$\forall A\in\mathcal{S}_n^+(\Rr)$, $1 +\sqrt[n]{\text{det}(A)}\leqslant\sqrt[n]{\text{det}(I_n+A)}$.
}
\end{center}
\fincorrection
\correction{005795}
Soit $A$ une matrice orthogonale à coefficients entiers. Puisque les colonnes ou les lignes de $A$ sont unitaires, on trouve par ligne ou par colonne un et un seul coefficient de valeur absolue égale à $1$, les autres coefficients étant nuls. $A$ est donc obtenue en multipliant chaque coefficient d'une matrice de permutation par $1$ ou $-1$. Réciproquement, une telle matrice est orthogonale à coefficients entiers.

Il y a $n!$ matrices de permutation et pour chaque matrice de permutation $2^n$ façons d'attribuer un signe $+$ ou $-$ à chaque coefficient égal à $1$. Donc

\begin{center}
\shadowbox{
$\text{card}(O_n(\Rr)\cap\mathcal{M}_n(\Zz)) = 2^nn!$.
}
\end{center}
\fincorrection
\correction{005798}
Puisque les matrices $S_1 ={^t}AA$ et $S_2 = A{^t}A$ sont symétriques réelles, ces deux matrices sont à valeurs propres réelles. On sait d'autre part que si $M$ et $N$ sont deux matrices quelconques alors les matrices $MN$ et $NM$ ont même polynôme caractéristique.

Notons alors $(\lambda_i)_{1\leqslant i\leqslant n}$ la famille des valeurs propres des matrices $S_1$ et $S_2$ et posons $D =\text{Diag}(\lambda_1,...\lambda_n)$. D'après le théorème spectral, il existe deux matrices orthogonales $P_1$ et $P_2$ telles que $S_1 =P_1D{^t}P_1$ et $S_2 = P_2D{^t}P_2$. Mais alors 

\begin{center}
$S_2=P_2({^t}P_1S_1P_1){^t}P_2= (P_2{^t}P_1)S_1{^t}(P_2{^t}P_1)$.
\end{center}

Comme la matrice $P_2{^t}P_1$ est orthogonale, on a montré que les matrices $S_1$ et $S_2$ sont orthogonalement semblables.

\begin{center}
\shadowbox{
$\forall A\in\mathcal{M}_n(\Rr)$, les matrices ${^t}AA$ et $A{^t}A$ sont orthogonalement semblables.
}
\end{center}
\fincorrection
\correction{005799}
\textbf{Remarque.} Il faut prendre garde au fait que le produit de deux matrices symétriques n'est pas nécessairement symétrique. Plus précisément, si $A$ et $B$ sont deux matrices symétriques alors 

\begin{center}
$AB\in\mathcal{S}_n(\Rr)\Leftrightarrow{^t}(AB) = AB\Leftrightarrow{^t}B{^t}A = AB\Leftrightarrow BA = AB$
\end{center}

et le produit de deux matrices symétriques est symétrique si et seulement si ces deux matrices commutent. Donc au départ, rien n'impose que les valeurs propres de $AB$ soient toutes réelles .

Soient $A$ et $B$ deux matrices symétriques réelles positives. D'après l'exercice \ref{ex:rou2}, il existe deux matrices carrées $M$ et $N$ telles que $A ={^t}MM$ et $B={^t}NN$. On a alors $AB ={^t}MM{^t}NN$. La  matrice $AB$ a même polynôme caractéristique que la matrice $N({^t}MM{^t}N={^t}(M{^t}N)M{^t}N$. D'après l'exercice \ref{ex:rou2}, cette dernière matrice est symétrique positive et a donc des valeurs propres réelles positives. On a montré que les valeurs propres de la matrice $AB$ sont réelles et positives.

\begin{center}
\shadowbox{
$\forall(A,B)\in(\mathcal{S}_n^+(\Rr)$, $\text{Sp}(AB)\subset\Rr^+$.
}
\end{center}
\fincorrection
\correction{005800}
Soient $A$ et $B$ deux matrices symétriques réelles positives.

\textbf{1er cas.} Supposons qu'aucune des deux matrices $A$ ou $B$ n'est inversible, alors $\text{det}A +\text{det}B = 0$.

D'autre part, la matrice $A+B$ est symétrique car $(\mathcal{S}_n(R),+,.)$ est un espace vectoriel et ses valeurs propres sont donc réelles. De plus, pour $X$ vecteur colonne donné, ${^t}X(A+B)X= {^t}XAX+{^t}XBX\geqslant 0$.

La matrice $A+B$ est donc symétrique réelle positive. Par suite, les valeurs propres de la matrice $A+B$ sont des réels positifs et puisque $\text{det}(A+B)$ est le produit de ces valeurs propres, on a $\text{det}(A+B)\geqslant 0 =\text{det}A +\text{det}B$.

\textbf{2ème cas.}
Sinon, une des deux matrices $A$ ou $B$ est inversible (et donc automatiquement définie positive).
Supposons par exemple $A$ définie positive.

D'après l'exercice \ref{ex:rou2}, il existe une matrice inversible $M$ telle que $A ={^t}MM$. On peut alors écrire $A + B ={^t}MM + B ={^t}M(I_n+{^t}(M^{-1}BM^{-1})M$ et donc

\begin{center}
$\text{det}(A+B) =(\text{det}M)^2\text{det}(I_n+{^t}(M^{-1})BM^{-1}= (\text{det}M)^2\text{det}(I_n+C)$
\end{center}

 où $C ={^t}M^{-1}BM^{-1}$. La matrice $C$ est symétrique, positive car pour tout vecteur colonne $X$,
 
 \begin{center}
 ${^t}XCX={^t}X{^t}(M^{-1})BM^{-1} X ={^t}(M^{-1}X)B(M^{-1}X)\geqslant 0$
 \end{center}
 
 
et ses valeurs propres $\lambda_1$,..., $\lambda_n$ sont des réels positifs. Les valeurs propres de la matrice $I_n + C$ sont les réels $1 +\lambda_i$, $1\leqslant i\leqslant n$  et donc 

\begin{center}
$\text{det}(I_n+C) = (1+\lambda_1)...(1+\lambda_n)\geqslant 1 +\lambda_1...\lambda_n = 1 +\text{det}C$.
\end{center}

Maintenant, $\text{det}A = (\text{det}M)^2$ puis $\text{det}B=(\text{det}M)^2\text{det}C$ et donc

\begin{center}
$\text{det}A +\text{det}B =(\text{det}M)^2(1+\text{det}C)\leqslant (\text{det}M)^2\text{det}(I_n+C) = \text{det}(A+B)$.
\end{center}

On a montré que

\begin{center}
\shadowbox{
$\forall(A,B)\in(\mathcal{S}_n^+(\Rr)$, $\text{det}A+\text{det}B\leqslant\text{det}(A+B)$.
}
\end{center}
\fincorrection
\correction{005804}
\textbf{1ère solution.} (n'utilisant pas les valeurs propres)
Soient $A$ la matrice de l'énoncé puis $X =(x_i)_{1\leqslant i\leqslant n}$ un élément de $\mathcal{M}_{n,1}(\Rr)$.

\begin{align*}\ensuremath
{^t}XAX&= (n-1)\sum_{i=1}^{n}x_i^2 -\sum_{i\neq j}^{}x_ix_j=\sum_{i\neq j}^{}(x_i^2- x_ix_j)=\frac{1}{2}\left(\sum_{i\neq j}^{}x_i^2-2 x_ix_j+x_j^2\right)\\
 &=\sum_{i\neq j}(x_i-x_j)^2\geqslant 0
\end{align*}

et donc la matrice $A$ est positive. De plus, si $X=\left(1\right)_{1\leqslant i\leqslant n}\neq0$, ${^t}XAX= 0$ et donc la matrice $A$ n'est pas définie.

\textbf{2ème solution.} La matrice $A$ est symétrique réelle. Donc ses valeurs propres sont réelles et $A$ est diagonalisable. Par suite, la dimension de chacun de des sous-espaces propres de $A$ est égale à l'ordre de multiplicité de la valeur propre correspondante.

On note alors que $\text{rg}(A - nI_n) = 1$ et donc $n$ est valeur propre de $A$ d'ordre $n-1$. Soit $\lambda$ la valeur propre manquante.
\begin{center}
$(n-1)n+\lambda=\text{Tr}A = n(n-1)$.
\end{center}

Donc $\lambda= 0$. Ainsi, $\text{Sp}(A)\subset\Rr^+$ et donc la matrice $A$ est positive mais $0$ est valeur propre de $A$ et donc la matrice $A$ n'est pas définie.

\begin{center}
\shadowbox{
La matrice $A$ est positive et non définie.
}
\end{center}
\fincorrection
\correction{005805}
Pour la première question, une simple observation suffit :
les matrices $I$ et $-I$ sont dans $O_n(\Rr)$,
mais pas la matrice nulle qui est leur milieu.



Soient $A$ et $B$ deux matrices orthogonales distinctes. Montrons que pour tout réel $\lambda\in]0,1[$, la matrice $(1-\lambda)A +\lambda B$ n'est pas orthogonale.

Supposons par l'absurde qu'il existe $\lambda\in]0,1[$ tel que la matrice $(1-\lambda)A +\lambda B$ soit orthogonale.

Pour $j\in\llbracket1,n\rrbracket$, on note respectivement $A_j$, $B_j$ et $C_j$ la $j$-ème colonne de matrice $A$, de la matrice $B$ et de la matrice $(1-\lambda)A +\lambda B$. Ces trois matrices étant orthogonales, pour tout $j\in\llbracket1,n\rrbracket$,

\begin{center}
$1 =\|C_j\|\leqslant(1-\lambda)\|A_j\|+\lambda\|B_j\|=(1 -\lambda)+\lambda= 1$,
\end{center}

et donc $\|C_j\| =(1-\lambda)\|A_j\|+\lambda\|B_j\|$. On est dans un cas d'égalité de l'inégalité de \textsc{Minkowski}. Puisque $\lambda\in]0,1[$, les colonnes $(-\lambda)A_j$ et $\lambda B_j$ ne sont pas nulles et donc sont colinéaires et de même sens. Puisque les réels $1-\lambda$ et $\lambda$ sont strictement positifs, il en est de même des colonnes $A_j$ et $B_j$ et puisque ces colonnes sont des vecteurs unitaires, ces colonnes sont en fin de compte égales. En résumé, si il existe $\lambda\in]0,1[$ tel que la matrice $(1-\lambda)A +\lambda B$ soit orthogonale, alors $A = B$. Ceci est une contradiction et on a montré que

\begin{center}
\shadowbox{
$O_n(\Rr)$ n'est pas convexe.
}
\end{center}
\fincorrection
\correction{005813}
$A$ est la matrice d'un produit scalaire $\varphi$ dans une certaine base $\mathcal{B}$ fixée de $\Rr^n$. Soit $\mathcal{B}'$ l'orthonormalisée de \textsc{Schmidt} de la base $\mathcal{B}$ pour le produit scalaire $\varphi$ et $T$ la matrice de passage de la base $\mathcal{B}'$ à la base $\mathcal{B}$. La matrice $T$ est triangulaire de même que la matrice ${^t}T$.

Puisque la base $\mathcal{B}'$ est orthonormée pour le produit scalaire $\varphi$, la matrice de $\varphi$ dans la base $\mathcal{B}'$ est $I_n$. D'après les formules de changement de base, $A= {^t}T(\text{Mat}_{\mathcal{B}'}\varphi)T={^t}TT$.
\fincorrection
\correction{005814}
 Puisque la matrice $A$ est définie positive, il existe d'après le l'exercice \ref{ex:rou8} une matrice triangulaire supèrieure inversible $T$ telle que $A={^t}TT$. Posons alors $T= (t_{i,j})_{1\leqslant i,j\leqslant n}$.

\begin{center}
$\text{det}(A)= (\text{det}(T))^2=t_{1,1}^2 ...t_{n,n}^2$
\end{center}

Mais pour $i\in\llbracket1,n\rrbracket$, $a_{i,i}=\sum_{k=1}^{n}t_{k,i}^2\geqslant t_{i,i}^2$ et donc $\text{det}(A)\leqslant\prod_{i=1}^{n}a_{i,i}$.

\textbf{Remarque.} On a montré au passage que les coefficients diagonaux $a_{i,i}$ de $A$ étaient des réels strictement positifs.
\fincorrection
\correction{005487}
Je vous laisse vérifier la linéarité.
Si $x$ est colinéaire à $a$, $f(x)=0$ et les vecteurs de $\mbox{Vect}(a)\setminus\{0\}$ sont des vecteurs non nuls colinéaires à leur image.
Si $x$ n'est pas colinéaire à $a$, $a\wedge x$ est un vecteur non nul orthogonal à $a$ et il en est de même de $f(x)=a\wedge(a\wedge x)$. Donc, si $x$ est colinéaire à $f(x)$, $x$ est nécessairement orthogonal à $a$.
Réciproquement, si $x$ est un vecteur non nul orthogonal à $a$, $f(x)=(a.x)a-\|a\|^2x=-||a||^2x$ et $x$ est colinéaire à $f(x)$. Les vecteurs non nuls colinéaires à leur image sont les vecteurs non nuls de $\mbox{Vect}(a)$ et de $a^\bot$.
\fincorrection
\correction{005492}
Si la famille $(x_i)_{1\leq i\leq n}$ est une famille liée, l'inégalité est claire et de plus, on a l'égalité si et seulement si l'un des vecteurs est nuls.
Si la famille $(x_i)_{1\leq i\leq n}$ est une famille libre et donc une base de E, considérons $\mathcal{B}'=(e_1,...,e_n)$ son orthonormalisée de \textsc{Schmidt}. On a 

$$\left|\mbox{det}_{\mathcal{B}}(x_i)_{1\leq i\leq n}\right|=\left|\mbox{det}_{\mathcal{B}'}(x_i)_{1\leq i\leq n}\times\mbox{det}_{\mathcal{B}}(\mathcal{B}')\right|=|\mbox{det}_{\mathcal{B}'}(x_i)_{1\leq i\leq n}|,$$
car $\mbox{det}_{\mathcal{B}}\mathcal{B}'$ est le déterminant d'une d'une base orthonormée dans une autre et vaut donc $1$ ou $-1$.
Maintenant, la matrice de la famille $(x_i)_{1\leq i\leq n}$ dans $\mathcal{B}'$ est triangulaire supérieure et son déterminant est le produit des coefficients diagonaux à savoir les nombres $x_i|e_i$ (puisque $\mathcal{B}'$ est orthonormée). Donc

$$|\mbox{det}_{\mathcal{B}}(x_i)_{1\leq i\leq n}|=|\mbox{det}_{\mathcal{B}'}(x_i)_{1\leq i\leq n}|=\left|\prod_{i=1}^{n}(x_i|e_i)\right|\leq\prod_{i=1}^{n}\|x_i\|\times\|e_i\|=\prod_{i=1}^{n}||x_i||,$$
d'après l'inégalité de \textsc{Cauchy}-\textsc{Schwarz}. De plus, on a l'égalité si et seulement si, pour tout $i$,  $|x_i|e_i|=\|x_i\|\times\|e_i\|$ ou encore si et seulement si, pour tout $i$, $x_i$ est colinéaire à $e_i$ ou enfin si et seulement si la famille $(x_i)_{1\leq i\leq n}$ est orthogonale.
\fincorrection
\correction{005497}
L'application $(P,Q)\mapsto\int_{0}^{1}P(t)Q(t)\;dt$ est un produit scalaire sur $E=\Rr_3[X]$. Déterminons une base orthonormée de $E$. Pour cela, déterminons $(Q_0,Q_1,Q_2,Q_3)$ l'orthonormalisée de la base canonique $(P_0,P_1,P_2,P_3)=(1,X,X^2,X^3)$.
\textbullet~$||P_0||^2=\int_{-1}^{1}1^2\;dt=2$ et on prend \shadowbox{$Q_0=\frac{1}{\sqrt{2}}$.}
\textbullet~$P_1|Q_0=\frac{1}{\sqrt{2}}\int_{-1}^{1}t\;dt=0$ puis $P_1-(P_1|Q_0)Q_0=X$ puis $||P_1-(P_1|Q_0)Q_0||^2=\int_{-1}^{1}t^2\;dt=\frac{2}{3}$ et \shadowbox{$Q_1=\sqrt{\frac{3}{2}}X$.}
\textbullet~$P_2|Q_0=\frac{1}{\sqrt{2}}\int_{-1}^{1}t^2\;dt=\frac{\sqrt{2}}{3}$ et $P_2|Q_1=0$. Donc, $P_2-(P_2|Q_0)Q_0-(P_2|Q_1)Q_1=X^2-\frac{1}{3}$,
puis $||P_2-(P_2|Q_0)Q_0-(P_2|Q_1)Q_1||^2=\int_{-1}^{1}\left(t^2-\frac{1}{3}\right)^2\;dt=2\left(\frac{1}{5}-\frac{2}{9}+\frac{1}{9}\right)=
\frac{8}{45}$ et \shadowbox{$Q_2=\frac{\sqrt{5}}{2\sqrt{2}}(3X^2-1)$.}
\textbullet~$P_3|Q_0=P_3|Q_2=0$ et $P_3|Q_1=\sqrt{\frac{3}{2}}\int_{-1}^{1}t^4\;dt=\frac{\sqrt{6}}{5}$ et 
$P_3-(P_3|Q_0)Q_0-(P_3|Q_1)Q_1-(P_3|Q_2)Q_2=X^3-\frac{3}{5}X$,
puis $\left\|X^3-\frac{3}{5}X\right\|^2=\int_{-1}^{1}\left(t^3-\frac{3}{5}t\right)^2\;dt=2\left(\frac{1}{7}-\frac{6}{25}+\frac{3}{25}\right)=2\frac{25-21}{175}=\frac{8}{175}$, et \shadowbox{$Q_3=\frac{\sqrt{7}}{2\sqrt{2}}(5X^3-3X)$.}
Une base orthonormée de $E$ est $(Q_0,Q_1,Q_2,Q_3)$ où $Q_0=\frac{1}{\sqrt{2}}$, $Q_1=\frac{\sqrt{3}}{\sqrt{2}}X$, $Q_2=\frac{\sqrt{5}}{2\sqrt{2}}(3X^2-1)$ et $Q_3=\frac{\sqrt{7}}{2\sqrt{2}}(5X^3-3X)$.
Soit alors $P$ un élément quelconque de $E=\Rr_3[X]$ tel que $\int_{-1}^{1}P^2(t)\;dt=1$. Posons $P=aQ_0+bQ_1+cQ_2+dQ_3$.
Puisque $(Q_0,Q_1,Q_2,Q_3)$ est une base orthonormée de $E$, $\int_{-1}^{1}P^2(t)\;dt=||P||^2=a^2+b^2+c^2+d^2=1$. Maintenant, pour $x\in[-1,1]$, en posant $M_i=\mbox{Max}\{|Q_i(x)|,\;x\in[-1,1]\}$, on a~:

\begin{align*}\ensuremath
|P(x)|&\leq|a|\times|Q_0(x)|+|b|\times|Q_1(x)|+|c|\times|Q_2(x)|+|d|\times|Q_3(x)|\leq|a|M_0+|b|M_1+|c|M_2+|d|M_3\\
 &\leq\sqrt{a^2+b^2+c^2+d^2}\sqrt{M_0^2+M_1^2+M_2^2+M_3^2}=\sqrt{M_0^2+M_1^2+M_2^2+M_3^2}.
\end{align*}
Une étude brève montre alors que chaque $|P_i|$ atteint son maximum sur $[-1,1]$ en $1$ (et $-1$) et donc 

$$\sqrt{M_0^2+M_1^2+M_2^2+M_3^2}=\sqrt{\frac{1}{2}+\frac{3}{2}+\frac{5}{2}+\frac{7}{2}}=2\sqrt{2}.$$
Ainsi, $\forall x\in[-1,1],\;|P(x)|\leq2\sqrt{2}$ et donc $\mbox{Max}\{|P(x)|,\;x\in[-1,1]\}\leq2\sqrt{2}$.
Etudions les cas d'égalité. Soit $P\in\Rr_3[X]$ un polynôme éventuel tel que $\mbox{Max}\{|P(x)|,\;x\in[-1,1]\}\leq2\sqrt{2}$.
Soit $x_0\in[-1,1]$ tel que $\mbox{Max}\{|P(x)|,\;x\in[-1,1]\}=|P(x_0)|$. Alors :

\begin{align*}\ensuremath
2\sqrt{2}&=|P(x_0)|\leq|a|\times|Q_0(x_0)|+|b|\times|Q_1(x_0)|+|c|\times|Q_2(x_0)|+|d|\times|Q_3(x_0)|\leq|a|M_0+|b|M_1+|c|M_2+|d|M_3\\
 &\leq\sqrt{M_0^2+M_1^2+M_2^2+M_3^2}=2\sqrt{2}.
\end{align*}
Chacune de ces inégalités est donc une égalité. La dernière (\textsc{Cauchy}-\textsc{Schwarz}) est une égalité si et seulement si $(|a|,|b|,|c|,|d|)$ est colinéaire à $(1,\sqrt{3},\sqrt{5},\sqrt{7})$ ou encore si et seulement si $P$ est de la forme $\lambda(\pm Q_0\pm\sqrt{3}Q_1\pm\sqrt{5}Q_2\pm\sqrt{7}Q_3)$ où $\lambda^2(1+3+5+7)=1$ et donc $\lambda=\pm\frac{1}{4}$, ce qui ne laisse plus que $16$ polynômes possibles. L'avant-dernière inégalité est une égalité si et seulement si $x_0\in\{-1,1\}$ (clair). La première inégalité est une égalité si et seulement si 

$$|aQ_0(1)+bQ_1(1)+cQ_2(1)+dQ_3(1)|=|a|Q_0(1)+|b|Q_1(1)+|c|Q_2(1)+|d|Q_3(1),$$
ce qui équivaut au fait que $a$, $b$, $c$ et $d$ aient même signe et $P$ est l'un des deux polynômes
 
\begin{align*}\ensuremath
\pm\frac{1}{4}(Q_0+\sqrt{3}Q_1+\sqrt{5}Q_2+\sqrt{7}Q_3)&=\pm\frac{1}{4\sqrt{2}}\left(1+3X+\frac{5}{2}(3X^2-1)
+\frac{7}{2}(5X^3-3X)\right)\\
 &=\pm\frac{1}{8\sqrt{2}}(35X^3+15X2-15X-3)
\end{align*}
\fincorrection
\correction{005498}
Si $x$ est colinéaire à $k$, $r(x)=x$, et si $x\in k^\bot,\;r(x)=(\cos\theta)x+(\sin\theta)k\wedge x$.
Soit $x\in E$. On écrit $x=x_1+x_2$ où $x_1\in k^\bot$ et $x_2\in\mbox{Vect}(k)$. On a $x_2=(x.k)k$ (car $k$ est unitaire) et $x_1=x-(x.k)k$. Par suite,

\begin{align*}\ensuremath
r(x)&=r(x_1)+r(x_2)=(\cos\theta)x_1+(\sin\theta)k\wedge x_1+x_2=(\cos\theta)(x-(x.k)k)+(\sin\theta)k\wedge x+(x.k)k\\
 &=(\cos\theta)x+(1-\cos\theta)(x.k)k+sin\theta(k\wedge x)=(\cos\theta)x+2\sin^2\left(\frac{\theta}{2}\right)(x.k)k+sin\theta(k\wedge x)
\end{align*}
\textbf{Application.} Si $k=\frac{1}{\sqrt{2}}(e_1+e_2)$ et $\theta=\frac{\pi}{3}$, pour tout vecteur $x$, on a~:

$$r(x)=\frac{1}{2}x+\frac{1}{2}(x.k)k+\frac{\sqrt{3}}{2}(k\wedge x),$$
puis,
$r(e_1)=\frac{1}{2}e_1+\frac{1}{4}(e_1+e_2)-\frac{\sqrt{3}}{2\sqrt{2}}e_3=\frac{1}{4}(3e_1+e_2-\sqrt{6}e_3)$

$r(e_2)=\frac{1}{2}e_2+\frac{1}{4}(e_1+e_2)+\frac{\sqrt{3}}{2\sqrt{2}}e_3=\frac{1}{4}(e_1+3e_2+\sqrt{6}e_3)$

$r(e_3)=\frac{1}{2}e_3+\frac{\sqrt{3}}{2\sqrt{2}}(-e_2+e_1)=\frac{1}{4}(\sqrt{6}e_1-\sqrt{6}e_2+2e_3)$.

\begin{center}
\shadowbox{
La matrice cherchée est $\frac{1}{4}\left(
\begin{array}{ccc}
3&1&\sqrt{6}\\
1&3&-\sqrt{6}\\
-\sqrt{6}&\sqrt{6}&2
\end{array}
\right)$.
}
\end{center}
\fincorrection
\correction{005499}
L'application $(f,g)\mapsto\int_{0}^{1}f(t)g(t)\;dt$ est un produit scalaire sur $C^0([0,1],\Rr)$. D'après l'inégalité de \textsc{Cauchy}-\textsc{Schwarz},
 
\begin{align*}\ensuremath
I_nI_{n+2}&=\int_{0}^{1}f^n(t)\;dt\int_{0}^{1}f^{n+2}(t)\;dt=\int_{0}^{1}\left(\sqrt{(f(t))^n}\right)^2\;dt\int_{0}^{1}\left(\sqrt{(f(t))^{n+2}}\right)^2\;dt\\
 &\geq\left(\int_{0}^{1}\sqrt{(f(t))^n}\sqrt{(f(t))^{n+2}}\;dt\right)^2=\left(\int_{0}^{1}f^{n+1}(t)\;dt\right)^2=I_{n+1}^2
\end{align*}
Maintenant, comme $f$ est continue et strictement positive sur $[0,1]$, $I_n$ est strictement positif pour tout naturel $n$. On en déduit que $\forall n\in\Nn,\;\frac{I_{n+1}}{I_n}\leq\frac{I_{n+2}}{I_{n+1}}$ et donc que

\begin{center}
\shadowbox{
la suite $\left(\frac{I_{n+1}}{I_n}\right)_{n\in\Nn}$ est définie et croissante.
}
\end{center}
\fincorrection
\correction{005796}
\begin{enumerate}
 \item  Soit $f$ un endomorphisme du $\Rr$-espace vectoriel $\Cc$. Pour tout nombre complexe $z$

\begin{align*}\ensuremath
f(z)&= f((\text{Re}(z)).1 +(\text{Im}(z)).i) =(\text{Re}(z))f(1) +(\text{Im}(z))f(i) =\frac{1}{2}(z+\overline{z})f(1)+\frac{1}{2i}(z-\overline{z})f(i)\\
 &= \frac{f(1)-if(i)}{2}z +\frac{f(1)+if(i)}{2}\overline{z},
\end{align*}

et on peut prendre $a =\frac{f(1)-if(i)}{2}$ et $b =\frac{f(1)+if(i)}{2}$. (Réciproquement pour $a$ et $b$ complexes donnés, l'application $f$ ainsi définie est $\Rr$-linéaire et on a donc l'écriture générale complexe d'un endomorphisme du plan).

\item  

\begin{center}
$\text{Tr}(f) =\text{Re}(f(1)) +\text{Im}(f(i)) =\text{Re}(a+b)+\text{Im}(i(a-b)) =\text{Re}(a+b)+\text{Re}(a-b) = 2\text{Re}(a)$
\end{center}

et

\begin{align*}\ensuremath
\text{det}(f)&=\text{Re}(a+b)\text{Im}(i(a-b))-\text{Im}(a+b)\text{Re}(i(a-b))=\text{Re}(a+b)\text{Re}(a-b)+\text{Im}(a+b)\text{Im}(a-b)\\
 &=(\text{Re}(a))^2 -(\text{Re}(b))^2+(\text{Im}(a))^2-(\text{Im}(b))^2=|a|^2 - |b|^2.
\end{align*}

\begin{center}
\shadowbox{
$\text{Tr}(f)=2\text{Re}(a)$ et $\text{det}(f)=|a|^2 - |b|^2$.
}
\end{center}

\item  Soient $z$ et $z'$ deux nombres complexes. On rappelle que

\begin{center}
$z|z'= (\text{Re}z)(\text{Re}z')+(\text{Im}z)(\text{Im}z') =\frac{1}{4}(z+\overline{z})(z'+\overline{z'})-\frac{1}{4}(z-\overline{z})(z'-\overline{z'})=\frac{1}{2}(\overline{z}z'+z\overline{z'}) =\text{Re}(\overline{z}z')$.
\end{center}

et au passage si on oriente le plan de sorte que la base orthonormée $(1,i)$ soit directe,

\begin{center}
$[z,z']= (\text{Re}z)(\text{Im}z')+(\text{Im}z)(\text{Re}z') =\frac{1}{4i}(z+\overline{z})(z'-\overline{z'})-\frac{1}{4i}(z-\overline{z})(z'+\overline{z'})=\frac{1}{2i}(\overline{z}z'-z\overline{z'}) =\text{Im}(\overline{z}z')$.
\end{center}

Notons $M$ la matrice de $f$ dans la base $(1,i)$. Puisque la base $(1,i)$ est orthonormée,

\begin{center}
$f = f^*\Leftrightarrow M ={^t}M\Leftrightarrow\text{Im}(a+b) =\text{Re}(i(a-b))\Leftrightarrow\text{Im}(a+b) =-\text{Im}(a-b)\Leftrightarrow2\text{Im}a = 0\Leftrightarrow a\in\Rr$.
\end{center}

\begin{center}
\shadowbox{
$f=f^*\Leftrightarrow a\in\Rr$.
}
\end{center}
\end{enumerate}
\fincorrection
\correction{005802}
Il s'agit de montrer qu'un endomorphisme d'un espace euclidien $E$ qui conserve l'orthogonalité est une similitude.

On peut raisonner sur une base orthonormée de $E$ que l'on note  $\mathcal{B}=(e_i)_{1\leqslant i\leqslant n}$. Par hypothèse, la famille $(f(e_i))_{1\leqslant i\leqslant n}$ est orthogonale. De plus, pour $i\neq j$, $(e_i+e_j)|(e_i-e_j)=\|e_i\|^2-\|e_j\|^2= 0$ et donc $f(e_i+e_j)|f(e_i-e_j) = 0$ ce qui fournit $\|f(e_i)\| =\|f(e_j)\|$. Soit $k$ la valeur commune des normes des $f(e_i)$, $1\leqslant i\leqslant n$.

Si $k = 0$, tous les $f(e_i)$ sont nuls et donc $f$ est nulle.

Si $k\neq0$, l'image par l'endomorphisme $\frac{1}{k}f$ de la base othonormée $\mathcal{B}$ est une base orthonormée. Donc l'endomorphisme $\frac{1}{k}f$ est un automorphisme orthogonal de $E$ et donc l'endomorphisme $\frac{1}{k}f$ conserve la norme.

Dans tous les cas, on a trouvé un réel positif $k$ tel que $\forall x\in E$, $\|f(x)\| = k\|x\|$.
\fincorrection
\correction{002658}
a) Par r{\'e}currence et int{\'e}gration par parties, on montre que
$I_{p,q}={p!q!\over(p+q+1)!}$

b) En d{\'e}veloppant $(1-x)^n$ {\`a} l'aide de la formule du bin{\^o}me, on obtient apr{\`e}s
int{\'e}gration
$I_{n,n}=\sum_{k=0}^n {(-1)^kC_n^k\over n+k+1}$.

$D_n$ {\'e}tant le ppmc des d{\'e}nominateurs dans l'expression pr{\'e}c{\'e}dente, nous en d{\'e}duisons que
$I_{n,n}={a\over D_n}$ o{\`u} $a$ est un entier. Comme $I_{n,n}>0$, $a\geq1$ et donc
$I_{n,n}\geq{1\over D_n}$. En utilisant le r{\'e}sultat de la question a), nous en d{\'e}duisons
l'in{\'e}galit{\'e} demand{\'e}e.

c) Soit $D_n=\prod_{i=1}^k p_i^{\alpha_i}$ la d{\'e}composition en
facteurs premiers de $D_n$. Pour tout $i$ compris entre 1 et $k$,
$p_i^{\alpha_i}$ divise un des nombres $n+1,\ n+2,\dots,\ 2n+1$.
Par cons{\'e}quent, $p_i^{\alpha_i}\leq2n+1$. De plus, les $p_i$ {\'e}tant deux {\`a} deux
distincts et inf{\'e}rieurs ou {\'e}gaux {\`a} $2n+1$, $k\leq\pi(2n+1)$. D'o{\`u} la majoration 
demand{\'e}e.

\fincorrection
\correction{005291}
Soit $n$ un entier naturel.

$$n(n+1)(n+2)(n+3)+1=n^4+6n^3+11n^2+6n+1=(n^2+3n+1)^2,$$

avec $n^2+3n+1$ entier naturel.
\fincorrection
\correction{005292}
\begin{enumerate}
\item  Soit $n$ un entier relatif.

Si $n$ est pair, $n$ et $5n^3$ sont pairs de même que $5n^3+n$ et $2$ divise $5n^3+n$.

Si $n$ est impair, $n$ et $5n^3$ sont impairs et de nouveau $5n^3+n$ est pair. Finalement~:~$\forall n\in\Zz,\;2|(5n^3+n)$.

Si $n$ est multiple de $3$, $n$ et $5n^3$ sont multiples de $3$ de même que $5n^3+n$.

Si $n$ est de la forme $3p+1$, alors 

$$5n^2+1=5(3p+1)^2+1=45p^2+30p+6=3(9p^2+10p+2)$$

et $5n^2+1$ est divisible par $3$. Il en est de même de $5n^3+n=n(5n^2+1)$.

Si $n$ est de la forme $3p+2$, $5n^2+1=5(3p+2)^2+1=45p^2+60p+21=3(9p^2+20p+7)$ et $5n^2+1$ est divisible par $3$. Il en est de même de $5n^3+n=n(5n^2+1)$.

Finalement, $\forall n\in\Zz,\;3|(5n^3+n)$.

Enfin, $5n^3+n$ est divisible par $2$ et $3$ et donc par $2\times 3=6$. On a montré que~:~$\forall n\in\Zz,\;6|(5n^3+n)$. (Tout ceci s'exprime beaucoup mieux à l'aide de congruences. Par exemple~:~si $n\equiv1\;(3)$, $5n^2+1\equiv5.1^2+1=6\equiv0\;(3)$)

\item  $4^{2^n}$ signifie $(...((4^2)^2)^2...)^2$. Etudions la suite de ces élévations au carré successives modulo $7$. $4^{2^0}=4$ est dans $4+7\Zz$. $4^{2^1}=16$ est dans $2+7\Zz$. $4^{2^2}=16^2=(7k+2)^2=4+7k'$ est dans $4+7\Zz$...
Montrons par récurrence sur $p$ entier naturel que~:~$\forall p\in\Nn$, $4^{2^{2p}}$ est dans $4+7\Zz$ et $4^{2^{2p+1}}$ est dans $2+7\Zz$.

C'est vrai pour $p=0$.

Soit $p\geq0$. Si il existe deux entiers relatifs $k_{2p}$ et $k_{2p+1}$ tels que $4^{2^{2p}}=4+7k_{2p}$ et $4^{2^{2p+1}}=2+7k_{2p+1}$, alors~:

$$4^{2^{2p+2}}=(4^{2^{2p+1}})^2=(2+7k_{2p+1})^2=4+7(4k_{2p+1}+7k_{2p+1}^2)\in4+7\Zz,$$

puis 

$$4^{2^{2p+3}}=(4^{2^{2p+2}})^2=(4+7k_{2p+2})^2=16+28k_{2p+2}+49k_{2p+2}^2=2+7(2+4k_{2p+2}+7k_{2p+2}^2)\in2+7\Zz.$$

On a montré par récurrence que si $n$ est pair, $4^{2^n}$ est dans $4+7\Zz$ et si $n$ est impair, $4^{2^n}$ est dans $2+7\Zz$.

Ensuite $2^{2^0}=2$ est dans $2+7\Zz$ puis, pour $n\geq1$, $2^{2^n}=2^{2.2^{n-1}}=4^{2^{n-1}}$ est dans $4+7\Zz$ si $n-1$ est pair ou encore si $n$ est impair et est dans $2+7\Zz$ si $n$ est pair. Ainsi, que $n$ soit pair ou impair, $4^{2^n}+2^{2^n}+1$ est dans $(4+2)+1+7\Zz=7+7\Zz=7\Zz$ et on a montré que~:

$$\forall n\in\Nn,\;7|4^{2^n}+2^{2^n}+1.$$
\end{enumerate}
\fincorrection
\correction{005293}
Soient $m$, $n$ et $p$ trois entiers naturels et $r_1$, $r_2$ et $r_3$ les restes des divisions euclidiennes de $m$, $n$ et $p$ par $8$. Alors,

$$m^2+n^2+p^2=(8q_1+r_1)^2+(8q_2+r_2)^2+(8q_3+r_3)^2\in r_1^2+r_2^2+r_3^2+8\Zz.$$

Donc $m^2+n^2+p^2$ est dans $7+8\Zz$ si et seulement si $r_1^2+r_2^2+r_3^2$ est dans $7+8\Zz$.

Comme $r_1$, $r_2$ et $r_3$ sont des entiers entre $0$ et $7$, il suffit de vérifier que les sommes de trois carrés d'entiers compris au sens large entre $0$ et $7$ ne sont pas dans $7+8\Zz$.

Or, $0^2=0\in8\Zz$, $1^2=1\in1+8\Zz$, $2^2=4\in4+8\Zz$, $3^2=9\in1+8\Zz$, $4^2=16\in8\Zz$, $5^2=25\in1+8\Zz$, $6^2=36\in4+8\Zz$ et $7^2=49\in1+8\Zz$. Donc, les carrés des entiers de $0$ à $7$ sont dans $8\Zz$ ou $1+8\Zz$ ou $4+8\Zz$. Enfin,

$$\begin{array}{llll}
0+0+0=0\in8\Zz,\quad&0+0+1=1\in1+8\Zz,\quad&0+0+4=4\in4+8\Zz,&0+1+1=2\in2+8\Zz,\\
0+1+4=5\in5+8\Zz&0+4+4=8\in8\Zz,&1+1+1=3\in3+8\Zz,&1+1+4=6\in6+8\Zz,\\
1+4+4=9\in1+8\Zz,&4+4+4=12\in4+8\Zz.
\end{array}
$$

Aucune de ces sommes n'est dans $7+8\Zz$ et on a montré qu'un entier de la forme $8n+7$ n'est pas la somme de trois carrés.
\fincorrection
\correction{005294}
Soit $n\in\Nn^*$. En développant $(1+\sqrt{2})^n$ par la formule du binôme de \textsc{Newton} et en séparant les termes où $\sqrt{2}$ apparaît à un exposant pair des termes où $\sqrt{2}$ apparaît à un exposant impair, on écrit $(1+\sqrt{2})^n$ sous la forme $a_n+b_n\sqrt{2}$ où $a_n$ et $b_n$ sont des entiers naturels non nuls.

Mais alors $(1-\sqrt{2})^n=a_n-b_n\sqrt{2}$ et donc 

$$(-1)^n=(1+\sqrt{2})^n(1-\sqrt{2})^n=(a_n+b_n\sqrt{2})(a_n-b_n\sqrt{2})=a_n^2-2b_n^2$$
ou finalement,
 
$$((-1)^na_n)a_n+(2(-1)^{n+1}b_n)b_n=1$$
 
où $(-1)^na_n=u$ et $2(-1)^{n+1}b_n=v$ sont des entiers relatifs. Le théorème de \textsc{Bezout} permet d'affirmer que $a_n$ et $b_n$ sont premiers entre eux.

\fincorrection
\correction{005295}
Posons $(1+\sqrt{3})^n=a_n+b_n\sqrt{3}$ où $a_n$ et $b_n$ sont des entiers naturels. On a alors $(1-\sqrt{3})^n=a_n-b_n\sqrt{3}$ et donc 

$$(1+\sqrt{3})^{2n+1}+(1-\sqrt{3})^{2n+1}=2a_{2n+1}\in\Nn.$$

Mais de plus, $-1<1-\sqrt{3}< 0$ et donc, puisque $2n+1$ est impair, $-1<(1-\sqrt{3})^{2n+1}<0$. Par suite, 

$$2a_{2n+1}<(1+\sqrt{3})^{2n+1}<2a_{2n+1}+1,$$

ce qui montre que $E((1+\sqrt{3})^{2n+1})=2a_{2n+1}=(1+\sqrt{3})^{2n+1}+(1-\sqrt{3})^{2n+1}$ et montre déjà que $E((1+\sqrt{3})^{2n+1})$ est un entier pair. Mais on en veut plus~:

\begin{align*}\ensuremath
(1+\sqrt{3})^{2n+1}+(1-\sqrt{3})^{2n+1}&=(1+\sqrt{3})((1+\sqrt{3})^2)^n+(1-\sqrt{3})((1-\sqrt{3})^2)^n\\
 &=(1+\sqrt{3})(4+2\sqrt{3})^n+(1-\sqrt{3})(4-2\sqrt{3})^n\\
 &=2^n((1+\sqrt{3})(2+\sqrt{3})^n+(1-\sqrt{3})(2-\sqrt{3})^n)
\end{align*}

Montrons enfin que $(1+\sqrt{3})(2+\sqrt{3})^n+(1-\sqrt{3})(2-\sqrt{3})^n$ est un entier, pair. Mais, $(1+\sqrt{3})(2+\sqrt{3})^n$ est de la forme $A+B\sqrt{3}$ où $A$ et $B$ sont des entiers naturels et donc, puisque $(1-\sqrt{3})(2-\sqrt{3})^n=A-B\sqrt{3}$, on a finalement $(1+\sqrt{3})(2+\sqrt{3})^n+(1-\sqrt{3})(2-\sqrt{3})^n=2A$ où $A$ est un entier.

Donc, $(1+\sqrt{3})(2+\sqrt{3})^n+(1-\sqrt{3})(2-\sqrt{3})^n$ est un entier pair, ou encore $(1+\sqrt{3})^{2n+1}+(1-\sqrt{3})^{2n+1}=E((1+\sqrt{3})^{2n+1})$ est un entier divisible par $2^{n+1}$.
\fincorrection
\correction{005296}
Soit $n$ un entier naturel non nul. On note $\sigma(n)$ la somme de ses chiffres en base $10$ (voir l'exercice \ref{exo:suprou19}). Si $n=a_0+10a_1+...+10^ka_k$ où $k\in\Nn$, $0\leq a_i\leq9$ pour $0\leq i\leq k$ et $a_k\neq0$, alors

$$\sigma(n)=a_0+...+a_k\leq9(k+1)\leq 9(E(\log n)+1)\leq9(\log n+1).$$

Donc, 

$$A=\sigma(4444^{4444})\leq9(\log(4444^{4444})+1)\leq9(4444\log(10^5)+1)=9(4444.5+1)=9.22221=199989.$$

Puis, $B=\sigma(A)\leq1+5.9=46$, puis $\sigma(B)\leq\sigma(39)=12$. Donc, $1\leq\sigma(B)\leq12$.

D'autre part, on sait que modulo 9~:~$\sigma(B)\equiv B\equiv A=4444^{4444}$.
Enfin, $4444^{4444}=(9.443+7)^{4444}\equiv7^{4444}\;(9)$. De plus, $7\equiv-2\;(9)$ puis $7^2\equiv4\;(9)$ puis $7^3\equiv 28\equiv1\;(9)$ et donc $7^{4444}=(7^3)^{1481}.7\equiv(1^3)^{1481}.7\equiv7\;(9)$. Finalement, $1\leq\sigma(B)\leq12$ et $C\equiv7\;(9)$ ce qui impose $C=7$.
\fincorrection
\correction{005297}
On a trois possibilités~:~$p\in3\Zz$, $p\in3\Zz+1$ ou $p\in3\Zz-1$.

Dans les deux derniers cas, $p^2\in1+3\Zz$ et $8p^2+1\in9+3\Zz=3\Zz$. Mais alors, $8p^2+1$ est premier et multiple de $3$ ce qui impose $8p^2+1=3$. Cette dernière égalité est impossible.

Il ne reste donc que le cas où $p$ est premier et multiple de $3$, c'est-à-dire $p=3$ (en résumé, 
$p$ et $8p^2+1$ premiers impliquent $p=3$). Dans ce cas, $8p^2+1=73$ et $8p^2-1=71$ sont effectivement premiers.
\fincorrection
\correction{005298}
\begin{enumerate}
\item  Pour $1\leq k\leq n$, $kC_n^k=nC_{n-1}^{k-1}$. Donc, si $k$ et $n$ sont premiers entre eux, puisque $n$ divise $kC_n^k$, le théorème de \textsc{Gauss} permet d'affirmer que $n$ divise $C_n^k$.

\item  De même, $(n+1)C_{2n}^{n-1}=nC_{2n}^{n}$ montre que $(n+1)$ divise $nC_{2n}^{n}$ et, puisque $n$ et $(n+1)$ sont premiers entre eux (d'après \textsc{Bezout} puisque $(n+1)-n=1$), $(n+1)$ divise $C_{2n}^{n}$ d'après le théorème de \textsc{Gauss}.
\end{enumerate}
\fincorrection
\correction{005299}
\begin{enumerate}
\item  Posons $d=x\wedge y$ et $m=x\vee y$. $d$ divise $m=105=3.5.7$ mais, puisque $d$ divise $x$ et $y$, $d$ divise aussi $x+y=56=2^3.7$. Donc, $d$ divise $105\wedge56=7$ et nécessairement $d=1$ ou $d=7$.

\begin{itemize}
\item[1er cas.] $d=1$ fournit, puisque $m=105$, $xy=md=105$. $x$ et $y$ sont donc les solutions de l'équation $X^2-56X+105=0$ qui n'admet pas de solutions entières.

\item[2ème cas.] $d=7$ fournit $xy=7.105=735$. $x$ et $y$ sont donc les solutions de l'équation $X^2-56X+735=0$ qui admet les solutions $21$ et $35$.
\end{itemize}

Réciproquement, $21+35=56$ et $21\vee35=3.5.7=105$. $\mathcal{S}=\{(21,35),(35,21)\}$.

\item  On pose $x=dx'$ et $y=dy'$ avec $x'$ et $y'$ premiers entre eux et $d=x\wedge y$. Le système s'écrit $\left\{
\begin{array}{l}
x'-y'=1\\
dx'y'=72
\end{array}
\right.
$ ou encore $\left\{
\begin{array}{l}
x'=y'+1\\
d(y'+1)y'=72
\end{array}
\right.
$. En particulier, $y'$ et $y'+1$ sont deux diviseurs consécutifs de $72$. $72=2^3.3^2$ admet $4.3=12$ diviseurs à savoir $1$, $2$, $3$, $4$, $6$, $8$, $9$, $12$, $18$, $24$, $36$ et $72$. Donc $y'$ est élément de $\{1,2,3,8\}$.

\begin{itemize}
\item[1er cas.] $y'=1$ fournit $d=\frac{72}{1.2}=36$ puis $y=36.1=36$ et $x=y+d=72$. Réciproquement, $72-36=36=36\wedge72$ et $36\vee72=72$.

\item[2ème cas.] $y'=2$ fournit $d=12$, $y=24$, $x=36$ qui réciproquement conviennent.

\item[3ème cas.] $y'=3$ fournit $d=6$, $y=18$, $x=24$ qui réciproquement conviennent.

\item[4ème cas.] $y'= 8$ fournit $d=1$, $y=8$, $x=9$ qui réciproquement conviennent.
\end{itemize}

$$\mathcal{S}=\{(9,8),(24,18),(36,24),(72,36)\}.$$

\item  $d$ divise $m$ et donc $d$ divise $243=3^5$ et $d\in\{1,3,9,27,81,243\}$. On pose alors $x=dx'$, $y=dy'$ avec $x'$ et $y'$ premiers entre eux.

\begin{itemize}
\item[1er cas.] Si $d=1$ on a $x'y'-1=243$ ou encore $x'y'=244$ ce qui fournit les possibilités (en n'oubliant pas que $x'$ et $y'$ sont premiers entre eux)~:

$x'=1$, $y'=244$ puis $x=1$ et $y=244$, 

$x'=4$, $y'=61$ puis $x=4$ et $y=61$,

$x'=61$, $y'=4$ puis $x=61$ et $y=4$,
 
$x'=244$, $y'=1$ puis $x=244$ et $y=1$ qui réciproquement conviennent.

\item[2ème cas.] Si $d=3$, on a $x'y'=81+1=82$ ce qui fournit les possibilités~:

$x'=1$, $y'=82$ puis $x=3$ et $y=246$,
 
$x'=2$, $y'=41$ puis $x=6$ et $y=123$,

$x'=41$, $y'=2$ puis $x=123$ et $y=6$,
 
$x'=82$, $y'=1$ puis $x=246$ et $y=3$ qui réciproquement conviennent.

\item[3ème cas.] Si $d=9$ on a $x'y'=27+1=28$ ce qui fournit les possibilités~:

$x'=1$, $y'=28$ puis $x=9$ et $y=252$,
 
$x'=4$, $y'=7$ puis $x=36$ et $y=63$,

$x'=7$, $y'=4$ puis $x=63$ et $y=36$,
 
$x'=28$, $y'=1$ puis $x=252$ et $y=9$ qui réciproquement conviennent.

\item[4ème cas.] Si $d=27$ on a $x'y'=9+1=10$ ce qui fournit les possibilités~:

$x'=1$, $y'=10$ puis $x=27$ et $y=270$,
 
$x'=2$, $y'=5$ puis $x=54$ et $y=135$,

$x'=5$, $y'=2$ puis $x=135$ et $y=54$,
 
$x'=10$, $y'=1$ puis $x=270$ et $y=27$ qui réciproquement conviennent.

\item[5ème cas.] Si $d=81$, on a $x'y'=3+1=4$ ce qui fournit les possibilités~:

$x'=1$, $y'=4$ puis $x=81$ et $y=324$,
 
$x'=4$, $y'=1$ puis $x=324$ et $y=81$ qui réciproquement conviennent.

\item[6ème cas.] Si $d=243$, on a $x'y'=1+1=2$ ce qui fournit les possibilités~:

$x'=1$, $y'=2$ puis $x=243$ et $y=486$,
 
$x'=2$, $y'=1$ puis $x=486$ et $y=243$ qui réciproquement conviennent.
\end{itemize}
\end{enumerate}
\fincorrection
\correction{005300}
Soit $n$ un entier supérieur ou égal à $2$.

$$(n-2)^2+(n-1)^2+n^2+(n+1)^2+(n+2)^2=5n^2+10=5(n^2+2).$$

$5(n^2+2)$ devant être un carré parfait, $n^2+2$ doit encore être divisible par $5$ mais si $n$ est dans $5\Zz$, $n^2+2$ est dans $2+5\Zz$, si $n$ est dans $\pm1+5\Zz$, $n^2+2$ est dans $3+5\Zz$ et si $n$ est dans $\pm2+5\Zz$, $n^2+2$ est dans $1+5\Zz$ et $n^2+2$ n'est jamais divisible par $5$. Une somme de cinq carrés d'entiers consécutifs n'est donc pas un carré parfait.
\fincorrection
\correction{005301}
Soient $n$ et $m$ deux entiers naturels tels que $n<m$. Posons $m=n+k$ avec $k>0$. On note que 

$$F_m=2^{2^{n+k}}+1=(2^{2^n})^{2^k}+1=(F_n-1)^{2^k}+1.$$

En développant l'expression précédente par la formule du binôme de \textsc{Newton} et en tenant compte du fait que $2^k$ est pair puisque $k$ est strictement positif, on obtient une expression de la forme $q.F_n+1+1=q.F_n+2$.

Le P.G.C.D. de $F_n$ et $F_m$ doit encore diviser $F_m-q.F_n=2$ et vaut donc $1$ ou $2$. Enfin, puisque $2^n$ et $2^m$ sont strictement positifs, $F_n$ et $F_m$ sont impairs et leur P.G.C.D. vaut donc $1$ (ce résultat redémontre l'existence d'une infinité de nombres premiers).
\fincorrection
\correction{005302}
\begin{enumerate}
\item  Soit, pour $n$ entier naturel non nul donné, $v_n=u_{n+1}u_{n-1}-u_n^2$. Alors,

$$v_{n+1}=u_{n+2}u_n-u_{n+1}^2=(u_n+u_{n+1})u_n-u_{n+1}(u_{n-1}+u_n)=u_n^2-u_{n+1}u_{n-1}=-v_n.$$

La suite $v$ est donc une suite géométrique de raison $-1$ et on a~:

$$\forall n\in\Nn^*,\;v_n=(-1)^{n-1}v_1=(-1)^n.$$

Cette égalité s'écrit encore $((-1)^nu_{n-1})u_{n+1}+((-1)^{n+1}u_n)u_n=1$ et le théorème de \textsc{Bezout} permet d'affirmer que pour tout entier naturel $n$, les entiers $u_n$ et $u_{n+1}$ sont premiers entre eux (il est clair par récurrence que la suite $u$ est à valeurs entières).

\item  Pour $m=1$ et $n$ entier naturel quelconque~:

$$u_{n+m}=u_{n+1}=u_{n+1}u_1+u_nu_0=u_{n+1}u_m+u_{m-1}u_n.$$

Pour $m=2$ et $n$ entier naturel quelconque~:

$$u_{n+m}=u_{n+2}=u_{n+1}+u_n=u_{n+1}u_2+u_nu_1=u_{n+1}u_m+u_{m-1}u_n.$$

Soit $m\geq1$. Supposons que pour tout entier naturel $n$, on a $u_{n+m}=u_{n+1}u_m+u_{m-1}u_n$ et $u_{n+m+1}=u_{n+1}u_{m+1}+u_mu_n$. Alors, pour tout entier naturel $n$,

\begin{align*}\ensuremath
u_{n+m+2}&=u_{n+m+1}+u_{n+m}=u_{n+1}u_{m+1}+u_mu_n+u_{n+1}u_m+u_{m-1}u_n\;(\mbox{par hypothèse de récurrence})\\
 &=u_{n+1}(u_{m+1}+u_m)+u_n(u_m+u_{m-1})=u_{n+1}u_{m+2}+u_nu_{m+1}.
\end{align*}

ce qui démontre l'égalité proposée par récurrence.

Soient $n$ et $m$ deux entiers naturels tels que $n\geq m$. La division euclidienne de $n$ par $m$ s'écrit $n=mq+r$ avec $q$ et $r$ entiers tels que $0\leq r\leq m-1$.

Or, $u_{m+r}=u_mu_{r+1}+u_{m-1}u_r$. Par suite, un diviseur commun à $u_m$ et $u_r$ divise encore $u_m$ et $u_{m+r}$ et réciproquement un diviseur commun à $u_m$ et $u_{m+r}$ divise $u_{m-1}u_r$. Mais, $u_m$ et $u_{m-1}$ sont premiers entre eux et, d'après le théorème de \textsc{Gauss}, un diviseur commun à $u_m$ et $u_{m+r}$ divise $u_r$. Les diviseurs communs à $u_m$ et $u_r$ sont encore les diviseurs communs à $u_m$ et $u_{m+r}$ et donc~:

$$u_m\wedge u_r=u_m\wedge u_{m+r}.$$
 
Puis, par récurrence 

$$u_m\wedge u_r=u_m\wedge u_{m+r}=u_m\wedge u_{2m+r}=...=u_m\wedge u_{qm+r}=u_m\wedge u_n.$$

Ainsi, les algorithmes d'\textsc{Euclide} appliqués d'une part à $u_m$ et $u_n$ et d'autre part à $m$ et $n$ s'effectuent en parallèle et en particulier, $u_m\wedge u_n=u_{m\wedge n}.$
\end{enumerate}
\fincorrection
\correction{005303}
\begin{enumerate}
\item  Posons $d=x\wedge y\wedge z$ puis $x=dx'$, $y=dy'$ et $z=dz'$ où $x'\wedge y'\wedge z'=1$.

$$x^2+y^2=z^2\Leftrightarrow d^2({x'}^2+d^2{y'}^2)=d^2{z'}2\Leftrightarrow {x'}^2+{y'}^2={z'}^2,$$

avec $x'\wedge y'\wedge z'=1$, ce qui montre que l'on peut se ramener au cas où $x$, $y$ et $z$ sont premiers entre eux.

Supposons donc $x$, $y$ et $z$ premiers entre eux (dans leur ensemble). Soit $p$ un nombre premier. Si $p$ divise $x$ et $y$ alors $p$ divise $x^2+y^2=z^2$ et donc $p$ est également un facteur premier de $z$ contredisant le fait que $x$, $y$ et $z$ sont premiers entre eux. Donc, $x$ et $y$ sont premiers entre eux.

Si $p$ divise $x$ et $z$ alors $p$ divise $z^2-x^2=y^2$ et donc $p$ est également un facteur premier de $y$, contredisant le fait que $x$, $y$ et $z$ sont premiers entre eux. Donc, $x$ et $z$ sont premiers entre eux. De même, $y$ et $z$ sont premiers entre eux. Finalement, $x$, $y$ et $z$ sont premiers entre eux deux à deux.

\item  Puisque $x$, $y$ et $z$ sont deux à deux premiers entre eux, parmi les nombres $x$, $y$ et $z$, il y a au plus un nombre pair. Mais si ces trois nombres sont impairs, $x^2+y^2=z^2$ est pair en tant que somme de deux nombres impairs contredisant le fait que $z$ est impair. Ainsi, parmi les nombres $x$, $y$ et $z$, il y a exactement un nombre pair et deux nombres impairs.

Si $x$ et $y$ sont impairs, alors d'une part, $z$ est pair et $z^2$ est dans $4\Zz$ et d'autre part $x^2$ et $y^2$ sont dans $1+4\Zz$. Mais alors, $x^2+y^2$ est dans $2+4\Zz$ excluant ainsi l'égalité $x^2+y^2=z^2$. Donc, $z$ est impair et l'un des deux nombres $x$ ou $y$ est pair. Supposons, quite à permuter les lettres $x$ et $y$, que $x$ est impair et $y$ est pair.

Posons alors $y=2y'$ puis $X=\frac{z+x}{2}$ et $Z=\frac{z-x}{2}$ (puisque $x$ et $z$ sont impairs, $X$ et $Z$ sont des entiers). 

\item  On a

$$x^2+y^2=z^2\Leftrightarrow4{y'}^2=(z+x)(z-x)\Leftrightarrow{y'}^2=XZ.$$

Un diviseur commun à $X$ et $Z$ divise encore $z=Z+X$ et $x=Z-X$ et est donc égal à $\pm1$ puisque $x$ et $z$ sont premiers  entre eux. $X$ et $Z$ sont des entiers premiers entre eux.

Le produit des deux entiers $X$ et $Z$ est un carré parfait et ces entiers sont premiers entre eux. Donc, un facteur premier de $X$ n'apparaît pas dans $Z$ et apparaît donc dans $X$ à un exposant pair ce qui montre que $X$ est un carré parfait. De même, $Z$ est un carré parfait.
 
\item  Donc, il existe deux entiers relatifs $u$ et $v$ tels que $X=u^2$ et $Z=v^2$. Mais alors, $z=Z+X=u^2+v^2$ et $x=Z-X=u^2-v^2$. Enfin, $y^2=z^2-x^2=(u^2+v^2)^2-(u^2-v^2)^2=4u^2v^2$ et donc, $y=2uv$ quite à remplacer $u$ par $-u$.

En résumé, si $x^2+y^2=z^2$ alors il existe $(d,u,v)\in\Nn^*\times\Zz\times\Zz$ tel que $x=d(u^2-v^2)$, $y=2duv$ et $z=d(u^2+v^2)$ ou bien $x=2duv$, $y=d(u^2-v^2)$ et $z=d(u^2+v^2)$.

Réciproquement, 

$$(d(u^2-v^2))^2+(2duv)^2=d^2(u^4+2u^2v^2+v^4)=(d(u^2+v^2))^2,$$

et on a trouvé tous les triplets Pythagoriciens. Par exemple, $d=1$, $u=2$ et $v=1$ fournissent le triplet $(3,4,5)$. $d=2$, $u=2$ et $v=1$ fournissent le triplet $(6,8,10)$ et $d=1$, $u=3$ et $v=2$ fournissent le triplet $(5,12,13)$.
\end{enumerate}
\fincorrection
\correction{005304}
Soient $x$ et $y$ deux entiers naturels tels que $3x^3+xy+4y^3=349$. On a $4y^3\leq 3x^3+xy+4y^3=349$ et donc 

$$y\leq\sqrt[3]{\frac{349}{4}}=4,4...$$

Donc, $y\in\{0,1,2,3,4\}$. De même, $3x^3\leq 3x^3+xy+4y^3=349$ et donc 

$$x\leq\sqrt[3]{\frac{349}{3}}=4,8...$$

Donc, $x\in\{0,1,2,3,4\}$ ce qui ne laisse plus que $5.5=25$ couples candidats. Ensuite,

$y=0$ donne $3x^3=349$ qui ne fournit pas de solutions.

$y=1$ donne $3x^3+x-345=0$, équation dont aucun des entiers de $0$ à $4$ n'est solution.

$y=2$ donne $3x^3+2x-317=0$, équation dont aucun des entiers de $0$ à $4$ n'est solution.

$y=3$ donne $3x^3+3x-241=0$, équation dont aucun des entiers de $0$ à $4$ n'est solution.

$y=4$ donne $3x^3+4x-93=0$ dont seul $x=3$ est solution.

$$\mathcal{S}=\{(3,4)\}.$$

\fincorrection
\correction{005305}
Si $x\geq5$ et $5\leq k\leq x$, alors $k!$ est divisible par $2.5=10$. D'autre part, $1!+2!+3!+4!=33$ et le chiffre des unités de $\sum_{k=1}^{x}k!$ est $3$. $\sum_{k=1}^{x}k!$ n'est donc pas un carré parfait car le chiffre des unités (en base 10) d'un carré parfait est à choisir parmi $0$, $1$, $4$, $5$, $6$, $9$. Donc, $x\leq4$.
Ensuite, $1!=1=1^2$ puis $1!+2!=1+2=3$ n'est pas un carré parfait, puis $1!+2!+3!=9=3^2$ puis $1!+2!+3!+4!=33$ n'est pas un carré parfait.

$$\mathcal{S}=\{(1,1),(3,3)\}.$$
\fincorrection
\correction{005306}
\begin{align*}\ensuremath
n&=9+8(10+10^2+...+10^{p-1})+4(10^p+...+10^{2p-1})=9+80\frac{10^{p-1}-1}{10-1}+4.10^p\frac{10^p-1}{10-1}\\
 &=\frac{1}{9}(81+80(10^{p-1}-1)+4.10^p(10^p-1))=\frac{1}{9}(4.10^{2p}+4.10^p+1)=\left(\frac{2.10^p+1}{3}\right)^2,
\end{align*}

(ce qui montre déjà que $n$ est le carré d'un rationnel). Maintenant, 

$$2.10^p+1=2(9+1)^p+1=2.\sum_{k=0}^{p}C_p^k9^k+1=3+2\sum_{k=1}^{p}C_p^k3^{2k}=3(1+2\sum_{k=1}^{p}C_p^k3^{2k-1}),$$

et $2.10^p+1$ est un entier divisible par $3$. Finalement, $n=\left(\frac{2.10^p+1}{3}\right)^2$ est bien le carré d'un entier.
\fincorrection
\correction{005307}
Pour $k\in\Nn$, posons $a_k=11...1$ ($k+1$ chiffres $1$ en base $10$).

Soit $n$ un entier naturel quelconque.

La division euclidienne de $a_k$ par $n$ s'écrit~:~$a_k=n.q_k+r_k$ où $q_k$ et $r_k$ sont des entiers naturels tels que $0\leq r_k\leq n-1$.

Les $n+1$ entiers $r_0$,..., $r_n$ sont à choisir parmi les $n$ entiers $0$, $1$,..., $n-1$. Les $n+1$ restes considérés ne peuvent donc être deux à deux distincts. Par suite,

$$\exists(k,l)\in\Nn^2/\;0\leq k<l\leq n\;\mbox{et}\;r_k=r_l.$$

Mais alors, $a_l-a_k=(q_l-q_k)n$ est un multiple de $n$. Comme $a_l-a_k=11...10...0$ ($l-k$ chiffres $1$ et $k+1$ chiffres $0$), on a montré que tout entier naturel admet un multiple de la forme $11...10...0=11...1.10^K$. Si de plus $n$ est impair, non divisible par $5$, alors $n$ est premier à $2$ et à $5$ et donc à $10^K$. D'après le théorème de \textsc{Gauss}, $n$ divise $11...1$.
\fincorrection
\correction{005308}
\begin{enumerate}
\item  $u_n^2=(2^{n+1}+1)^2=2^{2n+2}+2^{n+2}+1=10...010...01_2$ ($n-1$ puis $n+1$ chiffres $0$)
\item 
\begin{align*}\ensuremath
u_n^3&=(2^{n+1}+1)^3=2^{3n+3}+3.2^{2n+2}+3.2^{n+1}+1=2^{3n+3}+(2+1).2^{2n+2}+(2+1).2^{n+1}+1\\
 &=2^{3n+3}+2^{2n+3}+2^{2n+2}+2^{n+2}+2^{n+1}+1=10...0110...0110...01_2
\end{align*}

($n-1$ puis $n-1$ puis $n$ chiffres $0$)

\item  

\begin{align*}\ensuremath
u_n^3-u_n^2+u_n&=2^{3n+3}+3.2^{2n+2}+3.2^{n+1}+1-2^{2n+2}-2^{n+2}-1+2^{n+1}+1=2^{3n+3}+2^{2n+3}+2^{n+2}+1\\
 &=10...010...010...01
\end{align*}

($n-1$ puis $n$ puis $n+1$ chiffres $0$)
\end{enumerate}
\fincorrection
\correction{005309}
\begin{enumerate}
\item  Soit $n\in\Nn^*$. Posons $n=\sum_{k=0}^{p}a_k10^k$, où $p\in\Nn$, et $\forall k\in\{0,...,p\},\;a_k\in\{0,...,9\}$, et $a_p\neq0$. Le nombre de chiffres de $n$ est alors $p+1$. L'entier $p$ vérifie $10^p\leq n<10^{p+1}$ ou encore $p\leq \log n<p+1$. Par suite, $p=E(\log n)$. Ainsi, le nombre de chiffres de $n$ en base $10$ est $E(\log n)+1$.

\item  Pour $n\in\Nn^*$, posons $u_n=\frac{\sigma(n+1)}{\sigma(n)}$
\begin{enumerate}
\item Soit $n\in\Nn^*$. Posons $n=a_p10^p+...+10a_1+a_0=\overline{a_p...a_1a_0}_{10}$. Si au moins un des chiffres de $n$ n'est pas $9$, on note $k$ le plus petit indice tel que $a_k\neq9$. Alors, $0\leq k\leq p-1$ et $n=\overline{a_p...a_k9...9}_{10}$ et $n+1=\overline{a_p...a_{k+1}(a_k+1)0...0}_{10}$. Dans ce cas, si $k=0$,

$$\frac{\sigma(n+1)}{\sigma(n)}=\frac{\sigma(n)+1}{\sigma(n)}=1+\frac{1}{\sigma(n)}\leq1+1=2.$$

Si $1\leq k\leq p-1$,

$$\frac{\sigma(n+1)}{\sigma(n)}=\frac{a_p+...+a_k+1}{a_p+...+a_k+9k}\leq\frac{a_p+...+a_k+1}{a_p+...+a_k+1}=1\leq2.$$

Sinon, tous les chiffres de $n$ sont égaux à $9$, et dans ce cas,

$$\frac{\sigma(n+1)}{\sigma(n)}=\frac{1}{9(p+1)}\leq2.$$

Ainsi, pour tout entier naturel non nul $n$, on a $u_n\leq2$. La suite $u$ est donc bornée.

 

Pour $p\in\Nn^*$, $u_{10^p-1}=\frac{\sigma(10^p)}{\sigma(10^p-1)}=\frac{1}{9p}$. La suite extraite $(u_{10^p-1})_{p\in\Nn}$ converge et a pour limite $0$.

Pour $p\in\Nn^*$, $u_{10^p}=\frac{\sigma(10^p+1)}{\sigma(10^p)}=\frac{2}{1}=2$. La suite extraite $(u_{10^p})_{p\in\Nn}$ converge et a pour limite $2\neq0$.

On en déduit que la suite $u$ diverge.

\item Avec les notations du a), $1\leq\sigma(n)\leq9(p+1)=9(E(\log n)+1)\leq9(\log n+1)$.

\item Soit $n\in\Nn^*$. $1\leq\sqrt[n]{\sigma(n)}\leq\sqrt[n]{9(\log n+1)}=\mbox{exp}
(\frac{1}{n}(\ln9+\ln(1+\frac{\ln n}{\ln 10})$. Les deux membres de cet encadrement tendent vers $1$ et donc la suite $(\sqrt[n]{\sigma(n)})_{n\geq1}$ converge et $\lim_{n\rightarrow +\infty}\sqrt[n]{\sigma(n)}=1$.
\end{enumerate}
\end{enumerate}
\fincorrection
\correction{005310}
\begin{enumerate}
\item  (Formule de \textsc{Legendre}) Soit $n$ un entier naturel supérieur ou égal à $2$.

Si $p$ est un nombre premier qui divise $n!=1.2...n$, alors $p$ est un facteur premier de l'un des entiers $2$,..., $n$ et en particulier, $p\leq n$. Réciproquement, il est clair que si $p$ est un nombre premier tel que $p\leq n$, $p$ divise $n!$. Les facteurs premiers de $n!$ sont donc les nombres premiers inférieurs ou égaux à $n$.

Soit donc $p$ un nombre premier tel que $p\leq n$. Pour trouver l'exposant de $p$ dans la décomposition primaire de $n!$, on compte $1$ pour chaque multiple de $p$ inférieur ou égal à $n$, on rajoute $1$ pour chaque multiple de $p^2$ inférieur ou égal à $n$, on rajoute encore $1$ pour chaque multiple de $p^3$ inférieur ou égal à $n$... et on s'arrête quand l'exposant $k$ vérifie $p^k>n$.

$$n\geq p^k\Leftrightarrow\ln n\geq k\ln p\Leftrightarrow k\leq\frac{\ln n}{\ln p},$$

(car $\ln p>0$). Donc, si $k\geq E(\frac{\ln n}{\ln p})+1$, alors $p^k>n$.

Dit autrement, l'exposant de $p$ est la somme du nombre de multiples de $p$ inférieurs ou égaux à $n$, du nombre de multiples de $p^2$ inférieurs ou égaux à $n$, du nombre de multiple de $p^3$ inférieurs ou égaux à $n$... et du nombre de multiples de $p^{E(\ln n/\ln p)}$.

Soit $k$ un entier tel que $1\leq k\leq E(\frac{\ln n}{\ln p})$ et $K$ un entier naturel.

$$1\leq K.p^k\leq n\Leftrightarrow \frac{1}{p^k}\leq K\leq\frac{n}{p^k}\Leftrightarrow 1\leq K\leq E(\frac{n}{p^k}).$$

Il y a donc $E(\frac{n}{p^k})$ multiples de $p^k$ compris au sens large entre $1$ et $n$. On a montré que l'exposant de $p$ dans la décomposition de $n!$ en facteurs premiers est

$$E(\frac{n}{p})+E(\frac{n}{p^2})+E(\frac{n}{p^3})+...$$

\item  L'exposant de $5$ dans la décomposition primaire de $1000!$ est

$$E(\frac{1000}{5})+E(\frac{1000}{5^2})+E(\frac{1000}{5^3})+E(\frac{1000}{5^4})=200+40+8+1=249.$$

L'exposant de $2$ est évidemment supérieur (il y a déjà au moins $500$ nombres pairs entre $1$ et $1000$). Donc, la plus grande puissance de $10$ divisant $1000!$ est encore la plus grande puissance de $5$ divisant $1000!$, à savoir $249$. L'écriture en base $10$ de $1000!$ se termine par $249$ zéros.
\end{enumerate}
\fincorrection
\correction{005311}
(Petit théorème de \textsc{Fermat}) Soit $p$ un nombre premier.
\begin{enumerate}
\item  Soit $p$ un nombre premier et $k$ un entier tel que $1\leq k\leq p-1$. On a $kC_p^k=pC_{p-1}^{k-1}$. Donc, $p$ divise $kC_p^k$. Mais, $p$ est premier et donc $p$ est premier à tous les entiers compris entre $1$ et $p-1$ au sens large. D'après le théorème de \textsc{Gauss}, $p$ divise $C_p^k$.
\item  Soit $p$ un nombre premier. Montrons par récurrence que $\forall a\in\Nn^*$, $a^p\equiv a\;(p)$.

C'est clair pour $a=1$.

Soit $a\geq1$. Supposons que $a^p\equiv a\;(p)$. On a alors 

\begin{align*}\ensuremath
(a+1)^p&=\sum_{k=0}^{p}C_p^ka^k=a^p+1+\sum_{k=1}^{p-1}C_p^ka^k\\
 &\equiv a^p+1\;(p)\quad(\mbox{d'après 1)})\\
 &\equiv a+1\;(p)\quad(\mbox{par hypothèse de récurrence})
\end{align*}

On a montré par récurrence que $\forall a\in\Nn^*,\;a^p\equiv a\;(p)$.

\end{enumerate}
\fincorrection
\correction{005312}
Soit $p$ un entier naturel supérieur ou égal à $2$. 

Supposons que $(p-1)!\equiv-1\;(p)$. Il existe donc un entier relatif $a$ tel que $(p-1)!=-1+ap$ $(*)$.

Soit $k\in\{1,...,p-1\}$. L'égalité $(*)$ s'écrit encore $k(-\prod_{j\neq k}^{}j)+ap=1$. Le théorème de \textsc{Bezout} permet alors d'affirmer que $k$ et $p$ sont premiers entre eux. Ainsi, $p$ est premier avec tous les entiers naturels éléments de $\{1,...,p-1\}$ et donc, $p$ est un nombre premier.
\fincorrection
\nocorrection
\nocorrection
\nocorrection
\nocorrection
\nocorrection
\nocorrection
\nocorrection
\nocorrection
\nocorrection
\correction{003151}
\begin{enumerate}
  \item $x=\dot{\overline{25}}$, $y=\dot{\overline{32}}$.
  \item $x=\dot{\overline{15}}$ ou $\dot{\overline{16}}$.
\end{enumerate}
\fincorrection
\correction{003152}
\begin{enumerate}
  \item $\dot0$, $\pm\dot1$, $\pm\dot5$.
  \item
  \item 
\end{enumerate}
\fincorrection
\nocorrection
\nocorrection
\correction{003155}
{\'E}tudier le m{\^e}me produit dans $\Z/n\Z$.
\fincorrection
\nocorrection
\correction{003157}
\begin{enumerate}
  \item 
  \item 3.
  \item 
  \item 
  \item 
  \item $\dot{\overline{11}}$, $\dot{\overline{27}}$.
\end{enumerate}
\fincorrection
\nocorrection
\nocorrection
\correction{003160}
Pour $1\le k< p$~: $k!\,C_{p+k}^k
=(p+1)\dots(p+k) \equiv k! (\mathrm{mod}\,{p})$ donc $C_{p+k}^k \equiv 1(\mathrm{mod}\,{p})$.
De plus $C_p^k \equiv 0(\mathrm{mod}\, p)$ d'o{\`u} $C_p^kC_{p+k}^k\equiv C_p^k(\mathrm{mod}\,{p^2})$.

Ensuite $(p-1)!\,C_{2p}^p = 2(p+1)\dots(p+p-1)\equiv2(p-1)!+2p\sum_{i=1}^{p-1}\frac{(p-1)!}i (\mathrm{mod}\,{p^2})
\equiv2(p-1)!\Bigl(1+p\sum_{i=1}^{p-1}i'\Bigr) (\mathrm{mod}\,{p^2})$ o{\`u} $i'$ d{\'e}signe
l'inverse de~$i$ modulo~$p$. L'application $x \mapsto x^{-1}$ est une permutation
de $(\Z/p\Z)^*$ donc $\sum_{i=1}^{p-1}i'\equiv \frac{p(p-1)}2(\mathrm{mod}\, p)\equiv 0(\mathrm{mod}\, p)$,
d'o{\`u} $C_p^pC_{2p}^p\equiv 2(\mathrm{mod}\,{p^2})$.

Enfin $\sum_{k=0}^pC_p^kC_{p+k}^k \equiv 1 + \sum_{k=1}^{p-1}C_p^k + 2
(\mathrm{mod}\,{p^2})\equiv 2^p+1(\mathrm{mod}\,{p^2})$.
\fincorrection
\correction{003161}
L'{\'e}quation caract{\'e}ristique, $X^3=4(X^2+X+1)$ admet trois racines distinctes
dans~$\Z/11Z$~: $1,6,8$. Donc $x_n$ est de la forme~: $x_n = a + 6^nb + 8^nc$ avec $a,b,c\in\Z/11\Z$.
On a $6^{10} \equiv 8^{10} \equiv 1 (\mathrm{mod}\,{11})$, donc $(x_n)$ est p{\'e}riodique de p{\'e}riode divisant~$10$.
La plus petite p{\'e}riode est $1$ si $b=c=0$, $10$ sinon car les suites $(6^n)$ et $(8^n)$
ont $10$ comme plus petite p{\'e}riode modulo~$11$ et l'on a~:
$8(x_{n+1}-x_n) - 5(x_{n+2}-x_{n+1}) = 7\cdotp8^nc$ et
$7(x_{n+2}-x_{n+1}) - (x_{n+1}-x_n) = 7\cdotp6^nb$.
\fincorrection
\correction{003162}
\begin{enumerate}
  \item 
  \item 
  \begin{enumerate}
    \item Le nombre de solutions de l'{\'e}quation $x^q = \dot 1$ est inf{\'e}rieur
    ou {\'e}gal {\`a}~$q < p-1$.
    \item $\dot 0 = a^{3q}-\dot 1 = (a^q - \dot 1)(a^{2q} + a^q + \dot 1)$
    donc $a^{2q}$ est racine de $x^2 + x + \dot 1 = \dot 0$, de discriminant
    $-\dot 3$.
  \end{enumerate}
  \item Il existe $x\in\Z/p\Z$ solution de $x^2 + x + \dot 1 = \dot 0$,
    et un tel $x$ est d'ordre multiplicatif $3$. Par le th{\'e}or{\`e}me de Lagrange, on en d{\'e}duit
    $3\mid p-1$.
\end{enumerate}
\fincorrection
\correction{003163}
Regrouper $x$ et $n-x$.
\fincorrection
\nocorrection
\correction{004565}
\begin{enumerate}
  \item $R=1$.
  \item $R=1$.
  \item $R=1$.
  \item $R=\frac1e$.
  \item $R=\frac1{\sqrt b}$.
  \item $R=1$.
  \item $R=1$.
  \item $R=1$.
  \item $R=\frac13$.
  \item $R=1$.
  \item $R=1$.
  \item $R=\sqrt2-1$.
  \item $R=\frac{(k-1)^{k-1}}{k^k}$.
  \item $R=0$.
  \item $R=\frac12$, $2t\le 1+t^2\le2$.
  \item $R=1$, $a_n\sim\frac{\ln n}{n^2}$.
  \item $R=1$.
\end{enumerate}
 \fincorrection
\nocorrection
\correction{004567}
La suite $\Bigl(\cos\Bigl(\frac{2k\pi}5+\alpha\Bigr)\Bigr)_{k\in\N}$ est périodique de
période~$5$, donc prend au plus cinq valeurs distinctes.
soit $a$ celle de plus grande valeur absolue. Alors $R=\frac1{|a|}$.
\fincorrection
\correction{004568}
$\sum_{k=1}^n k^{-\alpha}\sim
\begin{cases}
\frac{n^{1-\alpha}}{1-\alpha} & \text{ si } \alpha<1,\cr
\ln(n) & \text{ si } \alpha=1,\cr
\zeta(\alpha) & \text{ si } \alpha>1.\cr
\end{cases}$
Dans les trois cas, on obtient $R=1$.

Il y convergence en~$x=1$ si et seulement si $\alpha<0$ et il y a divergence grossière en~$x=-1$
lorsque $\alpha>1$ vu les équivalents. Pour $\alpha\le 1$ et $x=-1$ il y a convergence (CSA).
\fincorrection
\correction{004569}
\begin{enumerate}
  \item $(a_n)$ est bornée et $(na_n)$ ne l'est pas, donc $R_a = 1$. $|b_n|\sim|a_n|$ donc $R_b=1$.
  \item Il y a doute seulement pour $x=\pm 1$. Le critère de convergence d'Abel (hors programme)
    s'applique, $\sum a_nx^n$ converge si $x=\pm1$.
    $b_n = a_n -\frac16a_n^3 +  O(n^{-5/3})$ et le critère d'Abel s'applique aussi
    à $\sum a_n^3x^n$ (linéariser le $\cos^3$), il y a aussi convergence pour $x=\pm 1$.
    
    Résolution conforme au programme : regrouper par paquets de six termes.
\end{enumerate}
 \fincorrection
\correction{004570}
\begin{enumerate}
  \item $R' = R^2$.
  \item $R'=\infty$.
  \item $R' = eR$.
\end{enumerate}
 \fincorrection
\correction{004571}
$\min(\sqrt R,\sqrt {R'}\,)$.
\fincorrection
\correction{004572}
\begin{enumerate}
  \item Série produit de $a(z)$ et $\frac1{z-\rho}  \Rightarrow 
             b_n = \sum_{k=0}^\infty a_{k+n+1}\rho^k$.
  \item Si $a(\rho) \ne 0$ : $b(z)$ converge pour $|z|<\rho$ et tend vers
             l'infini pour $z\to\rho^-  \Rightarrow  R = \rho$.\par
             Si $a(\rho) = 0$ : $\forall\ r > \rho$, $|a_p| \le \frac M{r^p} \Rightarrow 
             |b_n| \le \frac M{r^n(r-\rho)}  \Rightarrow  R=\infty$.
\end{enumerate}
 \fincorrection
\correction{004573}
\begin{enumerate}
  \item $]-1,2[$.
  \item Pour $0\le k\le 4^n$, on a $|a_k| \le C_{4^n}^{4^n/2}\Bigm/2^{4^n}$
(atteint pour $k=4^n/2$).

Donc $a_n\to0$ lorsque $n\to\infty$ et si $x>1$ alors
$a_{3*4^n/2}x^{3*4^n/2}$
\hbox to 0pt{\hskip5mm $\not{}$\hss}
$\to 0$ lorsque $n\to\infty$.
\end{enumerate}
 \fincorrection
\correction{005745}
\begin{enumerate}
 \item  Soit $z\neq0$. Pour $n >e^{1/|z|}$, on a $|z|\ln n > 1$ et donc la suite $((\ln n)^nz^n)$ ne tend pas vers 0 quand $n$ tend vers $+\infty$. Ainsi, pour tout nombre complexe non nul $z$, la série proposée diverge grossièrement. 

\begin{center}
\shadowbox{
$R = 0$.
}
\end{center}

\item   Soit $z\neq0$. Pour $n >\frac{1}{|z|^2}$, on a $|z|\sqrt{n}> 1$ et donc la suite $((\sqrt{n})^nz^n)$ ne tend pas vers 0 quand $n$ tend vers $+\infty$. Pour tout nombre complexe non nul $z$, la série proposée diverge grossièrement. 

\begin{center}
\shadowbox{
$R = 0$.
}
\end{center}

\item  D'après la formule de \textsc{Stirling}

\begin{center}
$(\ln(n!))^2\underset{n\rightarrow+\infty}{\sim}\ln^2\left(\left(\frac{n}{e}\right)^n\sqrt{2\pi n}\right)=\left(\left(n+\frac{1}{2}\right)\ln n -n +\ln(\sqrt{2\pi)}\right)^2\underset{n\rightarrow+\infty}{\sim}n^2\ln^2n$.
\end{center}

La série entière proposée a même rayon de convergence que la série entière associée à la suite $(n^2\ln^2n)$. Comme $\lim_{n \rightarrow +\infty}\frac{(n+1)^2\ln^2(n+1)}{n^2\ln^2n}=1$, la règle de d'\textsc{Alembert} permet d'affirmer que

\begin{center}
\shadowbox{
$R = 1$.
}
\end{center}

\item 

\begin{center}
$n^4\ln\left(\frac{1}{2}\left(\ch\frac{1}{n}+\cos\frac{1}{n}\right)\right))\underset{n\rightarrow+\infty}{=}n^4\ln\left(1+\frac{1}{24n^4}+o\left(\frac{1}{n^4}\right)\right)\underset{n\rightarrow+\infty}{=}\frac{1}{24}+o(1) .$
\end{center}

Donc $\left(\frac{1}{2}\left(\ch\frac{1}{n}+\cos\frac{1}{n}\right)\right)^{n^4}\underset{n\rightarrow+\infty}{\sim}e^{1/24}$ et

\begin{center}
\shadowbox{
$R = 1$.
}
\end{center}

\item  Pour $n\in\Nn^*$, posons $a_n=\frac{C_{2n+2}^{n+1}}{(n+1)^{n+1}}$.

\begin{align*}\ensuremath
\frac{a_{n+1}}{a_n}&=\frac{(2n+2)!}{(2n)!}\times\frac{n!^2}{(n+1)!^2}\times\frac{n^n}{(n+1)^{n+1}}=\frac{(2n+2)(2n+1)n^n}{(n+1)^2(n+1)^{n+1}}=\frac{1}{n+1}\times\frac{4n+2}{n+1}\times\left(1+\frac{1}{n}\right)^{-n}\\
 &\underset{n\rightarrow+\infty}{\sim}\frac{4}{ne}.
\end{align*}

et donc $\lim_{n \rightarrow +\infty}\left|\frac{a_{n+1}}{a_n}\right|=0$. D'après la règle de d'\textsc{Alembert},

\begin{center}
\shadowbox{
$R = +\infty$.
}
\end{center}

\item  On a vu que $\ln(n!)\underset{n\rightarrow+\infty}{\sim}n\ln n$. Donc la série entière proposée a même rayon de convergence que la série entière associée à la suite $\left(\frac{(n\ln n)^a}{n!^b}\right)$. Puis

\begin{center}
$\frac{((n+1)\ln(n+1))^a/(n+1)!^b}{(n\ln n)^a/n!^b}\underset{n\rightarrow+\infty}{\sim}\frac{1}{n^b}$
\end{center}

et donc, d'après la règle de d'\textsc{Alembert} 

\begin{center}
\shadowbox{
si $b >0$, $R =+\infty$, si $b = 0$, $R = 1$ et si $b < 0$, $R = 0$.
}
\end{center}

\item  Si $a=0$, $R=+\infty$. On suppose $a\neq0$.

\textbullet~Si $b> 1$,  $\frac{a^n}{1+b^n}\underset{n\rightarrow+\infty}{\sim}\left(\frac{a}{b}\right)^n$ et donc $R =\frac{b}{a}$.

\textbullet~Si $b= 1$,   $\frac{a^n}{1+b^n}=\frac{a^n}{2}$ et $R =a$.

\textbullet~Si $0\leqslant b< 1$,  $\frac{a^n}{1+b^n}\underset{n\rightarrow+\infty}{\sim}a^n$ et $R =a$.

Dans tous les cas 

\begin{center}
\shadowbox{
$R=\frac{\text{Max}(1,b)}{a}$ si $a>0$ et $R=+\infty$ si $a=0$.
}
\end{center}
\end{enumerate}
\fincorrection
\nocorrection
\nocorrection
\correction{004581}
\begin{enumerate}
  \item $-1+\sqrt x\Argth\sqrt x$ pour $0\le x < 1$
             et $-1-\sqrt{-x}\Arctan\sqrt{-x}$ pour $-1\le x \le 0$.
  \item $\frac{x+x^2}{(1-x)^3}$.
  \item $\frac{x(1+4x+x^2)}{(1-x)^4}$.
  \item $\frac{2(1-x^2)\ln(1-x) + x^2+2x}{4x^3}$
             (décomposer en éléments simples).
  \item $-\frac12(x+(x^2+1)\Arctan x)$ (décomposer en éléments simples).
  \item $-1+\frac u4\Argth u -\frac u2\Arctan u$, $u=\sqrt[4]x$.
  \item $\frac1{2(1-x)} + \frac5{2\sqrt x}\Argth\sqrt x$
             pour $0\le x < 1$
             et $\frac1{2(1-x)} + \frac5{2\sqrt{-x}}\Arctan\sqrt{-x}$
             pour $-1 < x \le 0$.
  \item $-\frac12\ln(1-2x\ch a+x^2)$.
  \item $1-\frac{5\cos2\theta-4}{(5-4\cos2\theta)^2}$ (linéariser).
  \item $\frac{2x-1}{(1-x)^2} - \frac{2\ln(1-x)}x$.
  \item $\ch\sqrt x$ pour $x \ge 0$ et $\cos\sqrt{-x}$ pour $x \le 0$.
  \item $\frac{e^{x^2}}2
             - \frac{e^{x^2\cos2\theta}}2\cos(x^2\sin2\theta)$.
  \item $(x+15x^2+25x^3+10x^4+x^5)e^x$.
  \item $\frac{e^x+2e^{-x/2}\cos(x\sqrt3/2)}3$, ($f''' = f$).
  \item $\frac{1-\sqrt{1-4x}-2x}{2x\sqrt{1-4x}}$.
  \item $\frac{x^2-1}2$.
  \item $-\frac{\ln(1-x)}{1-x}$.
\end{enumerate}
 \fincorrection
\correction{004582}
$R=\sqrt2-1$, $\Sigma = \frac{1-x}{1-2x-x^2}$.
\fincorrection
\correction{004583}
\begin{enumerate}
  \item 
  \item S'il existe $\lambda\in\mathrm{Sp}(A)$ tel que $|\lambda|\ge 1$
    et si $x$ est un vecteur propre associé alors $kA^kx = k\lambda^kx\not\rightarrow0$
    donc la série diverge.

    Si toutes les valeurs propres de~$A$ sont de module $<1$, comme
    $kA^k = \sum_\lambda \lambda^kP_\lambda(k)$ où les $P_\lambda$
    sont des polynômes à coefficients matriciels, la série converge absolument.
  \item $S = \sum_{k=0}^\infty (k+1)A^{k+1} = AS + \sum_{k=0}^\infty A^{k+1}
                = AS + A(I-A)^{-1}$ donc $S = A(I-A)^{-2}$ est inversible ssi
             $A$ l'est.
\end{enumerate}
 \fincorrection
\correction{004584}
\begin{enumerate}
  \item $\chi_A(\lambda) = -\lambda^3 + 2\lambda^2 + \lambda - 1$.
    $\chi_A(-1)>0$, $\chi_A(0) < 0$, $\chi_A(1)>0$, $\chi_A(2)>0$, $\chi_A(3)<0$
    donc $\chi_A$ admet une racine dans chacun des intervalles
    $]-1,0[$, $]0,1[$ et $]2,3[$.
  \item Cayley-Hamilton~: $t_n = 2t_{n-1}+t_{n-2}-t_{n-3}$.
  \item Soient $-1<\alpha<0<\beta<1<2<\gamma<3$ les valeurs propres de~$A$.
    On a $t_nz^n = (\alpha z)^n + (\beta z)^n + (\gamma z)^n$ donc la série $\sum_{n=0}^\infty t_nz^n$
    converge si et seulement si $|\gamma z|<1$ et vaut~:
    $$\frac1{1-\alpha z} + \frac1{1-\beta z} + \frac 1{1-\gamma z}
      = \frac1z\frac{\chi'}{\chi}\Bigl(\frac1z\Bigr)
      = \frac{-z^2-4z+3}{z^3-z^2-2z+1}.$$
\end{enumerate}
 \fincorrection
\correction{004585}
$=- \int_{t=0}^{1} \frac{t^3}{1+t^3}\,d t = \frac{\ln 2}{\strut 3} +\frac\pi{\strut 3\sqrt3}-1$.
\fincorrection
\correction{004586}
$R=1$. On décompose $P$ sous la forme :
	     $P = a_0 + a_1(X+1) + a_2(X+1)(X+2) + \dots + a_p(X+1)\dots(X+p)$.\par
	     Alors $\sum_{n=0}^\infty P(n)x^n = \frac{a_0}{1-x} +
	     \frac{a_1}{(1-x)^2} + \dots + \frac{p!\,a_p}{(1-x)^{p+1}}$.
\fincorrection
\correction{004587}
$\sum_{n=1}^\infty \frac{\sin n\theta}{2^n} =
		\frac{2\sin\theta}{5-4\cos\theta}$,\quad
	     $\sum_{n=1}^\infty \frac{\cos n\theta}{n2^n} =
		\ln2 - \frac12\ln(5-4\cos\theta)$.
\fincorrection
\correction{004588}
$f(t) = \sum_{n=0}^\infty u_nt^n = \frac{e^{-t}}{1-t\strut} = \sum_{n=0}^\infty\sum_{k=0}^n\frac{(-1)^k}{\strut k!}t^n$
donc $u_n\to\frac1{\strut e}$ lorsque $n\to\infty$.
\fincorrection
\correction{005746}
\begin{enumerate}
 \item  La règle de d'\textsc{Alembert} montre que la série proposée a un rayon de convergence égal à $1$.

\textbf{1ère solution.} Pour $x\in]-1,1[$, on pose $f(x)=\sum_{n=2}^{+\infty}\frac{1}{n(n-1)}x^n$.
$f$ est dérivable sur $]-1,1[$ et pour $x$ dans $]-1,1[$, 

\begin{center}
$f'(x) =\sum_{n=2}^{+\infty}\frac{1}{n-1}x^{n-1}=\sum_{n=1}^{+\infty}\frac{x^n}{n}= -\ln(1-x)$.
\end{center}

Puis, pour $x\in]-1,1[$, $f(x) = f(0) +\int_{0}^{x}f'(t) dt =(1-x)\ln(1-x) + x$.

\textbf{2ème solution.} Pour $x\in]-1,1[$, 

\begin{center}
$f(x) =\sum_{n=2}^{+\infty}\left(\frac{1}{n-1}-\frac{1}{n}\right)x^n =x\sum_{n=2}^{+\infty}\frac{x^{n-1}}{n-1}-\sum_{n=2}^{+\infty}\frac{x^n}{n}=-x\ln(1-x) +\ln(1-x) + x$.
\end{center}

\begin{center}
\shadowbox{
$\forall x\in]-1,1[$, $\sum_{n=2}^{+\infty}\frac{x^n}{n(n-1)}=-x\ln(1-x) +\ln(1-x) + x$.
}
\end{center}

\item  La règle de d'\textsc{Alembert} montre que la série proposée a un rayon égal à $1$.
Pour $x\in]-1,1[\setminus\{0\}$

\begin{center}
$\sum_{n=0}^{+\infty}\frac{3n}{n+2}x^n =3\left(\sum_{n=0}^{+\infty}x^n - 2\sum_{n=0}^{+\infty}\frac{x^n}{n+2}\right) = 3\left(\frac{1}{1-x}- \frac{2}{x^2}\sum_{n=2}^{+\infty}\frac{x^n}{n}\right) = 3\left(\frac{1}{1-x}+\frac{2}{x^2}(x+\ln(1-x))\right)$
\end{center}

\begin{center}
\shadowbox{
$\forall x\in]-1,1[$, $\sum_{n=0}^{+\infty}\frac{3n}{n+2}x^n=\left\{
\begin{array}{l}3\left(\frac{1}{1-x}+\frac{2}{x^2}(x+\ln(1-x))\right)\;\text{si}\;x\in]-1,1[\setminus\{0\}\\
\rule{0mm}{4mm}0\;\text{si}\;x=0
\end{array}
\right.$.
}
\end{center}

\item  La règle de d'\textsc{Alembert} montre que la série proposée a un rayon égal à $1$.

\textbullet~Soit $x\in]0,1[$.

\begin{align*}\ensuremath
\sum_{n=0}^{+\infty}\frac{x^n}{2n+1}&=\sum_{n=0}^{+\infty}\frac{\left(\sqrt{x}\right)^{2n}}{2n+1}=\frac{1}{\sqrt{x}}\sum_{n=0}^{+\infty}\frac{\left(\sqrt{x}\right)^{2n+1}}{2n+1}=\frac{1}{\sqrt{x}}(\ln(1+\sqrt{x}) -\ln(1-\sqrt{x}))\\
 &=\frac{\Argth(\sqrt{x})}{\sqrt{x}}.
\end{align*}	

\textbullet~Soit $x\in]-1,0[$.

\begin{align*}\ensuremath
\sum_{n=0}^{+\infty}\frac{x^n}{2n+1}&=\sum_{n=0}^{+\infty}(-1)^n\frac{(-x)^n}{2n+1}=\frac{1}{\sqrt{-x}}\sum_{n=0}^{+\infty}(-1)^n\frac{\left(\sqrt{-x}\right)^{2n+1}}{2n+1}=\frac{\Arctan(\sqrt{-x})}{\sqrt{-x}}
\end{align*}

\begin{center}
\shadowbox{
$\forall x\in]-1,1[$, $\sum_{n=0}^{+\infty}\frac{x^n}{2n+1}=\left\{
\begin{array}{l}
\frac{\Argth(\sqrt{x})}{\sqrt{x}}\;\text{si}\;x\in]0,1[\\
\rule[-2mm]{0mm}{4mm}1\;\text{si}\;x=0\\
\frac{\Arctan(\sqrt{-x})}{\sqrt{-x}}\;\text{si}\;x\in]-1,0[
\end{array}
\right.$.
}
\end{center}

\item  La règle de d'\textsc{Alembert} montre que la série proposée a un rayon égal à $+\infty$. Pour $x$ réel, 

\begin{center}
$f(x) =\frac{1}{2}\sum_{n=0}^{+\infty}\frac{2n+1-1}{(2n+1)!}x^n=\frac{1}{2}\sum_{n=0}^{+\infty}\frac{1}{(2n)!}x^n-\frac{1}{2}\sum_{n=0}^{+\infty}\frac{1}{(2n+1)!}x^n$.
\end{center}

\textbullet~Si $x > 0$,  

\begin{center}
$f(x) =\frac{1}{2}\sum_{n=0}^{+\infty}\frac{1}{(2n)!}\left(\sqrt{x}\right)^{2n}-\frac{1}{2\sqrt{x}}\sum_{n=0}^{+\infty}\frac{1}{(2n+1)!}\left(\sqrt{x}\right)^{2n+1}=\frac{1}{2}\left(\ch(\sqrt{x})-\frac{1}{\sqrt{x}}\sh(\sqrt{x})\right)$.
\end{center}

\textbullet~Si $x<0$,  

\begin{center}
$f(x) =\frac{1}{2}\sum_{n=0}^{+\infty}(-1)^n\frac{1}{(2n)!}\left(\sqrt{-x}\right)^{2n}-\frac{1}{2\sqrt{-x}}\sum_{n=0}^{+\infty}(-1)^n\frac{1}{(2n+1)!}\left(\sqrt{-x}\right)^{2n+1}=\frac{1}{2}\left(\cos(\sqrt{-x})-\frac{1}{\sqrt{-x}}\sin(\sqrt{-x})\right)$.
\end{center}

\begin{center}
\shadowbox{
$\forall x\in\Rr$, $\sum_{n=0}^{+\infty}\frac{nx^n}{(2n+1)!}=\left\{
\begin{array}{l}
\frac{1}{2}\left(\ch(\sqrt{x})-\frac{1}{\sqrt{x}}\sh(\sqrt{x})\right)\;\text{si}\;x>0\\
\rule[-2mm]{0mm}{4mm}0\;\text{si}\;x=0\\
\frac{1}{2}\left(\cos(\sqrt{-x})-\frac{1}{\sqrt{-x}}\sin(\sqrt{-x})\right)\;\text{si}\;x<0
\end{array}
\right.$.
}
\end{center}

\item  Immédiatement $R = +\infty$ et 

\begin{center}
\shadowbox{
$\forall x\in\Rr$, $\sum_{n=0}^{+\infty}\frac{x^{4n}}{(4n)!}=\frac{1}{2}(\cos x +\ch x)$.
}
\end{center}

\item  $\ch n\underset{n\rightarrow+\infty}{\sim}\frac{e^n}{2}$ et donc $R=\frac{1}{e}$.
Pour $x$ dans $\left]-\frac{1}{e},\frac{1}{e}\right[$,

\begin{align*}\ensuremath 
\sum_{n=0}^{+\infty}\left(\ch n\right)x^n=\frac{1}{2}\left(\sum_{n=0}^{+\infty}(ex)^n+\sum_{n=0}^{+\infty}\left(\frac{x}{e}\right)^n\right)=\frac{1}{2}\left(\frac{1}{1-ex}+\frac{1}{1-\frac{x}{e}}\right)=\frac{1}{2}\frac{2-\left(e+\frac{1}{e}\right)x}{x^2-\left(e+\frac{1}{e}\right)x+1}\\
 &=\frac{1-x\ch1}{x^2-2x\ch1+1}.
\end{align*}

\begin{center}
\shadowbox{
$\forall x\in\left]-\frac{1}{e},\frac{1}{e}\right[$, $\sum_{n=0}^{+\infty}\left(\ch n\right)x^n=\frac{1-x\ch1}{x^2-2x\ch1+1}$.
}
\end{center}

\item  La série proposée est le produit de \textsc{Cauchy} des séries entières  $\sum_{n=0}^{+\infty}x^n$ et  $\sum_{n=1}^{+\infty}\frac{x^n}{n}$ qui sont toutes deux de rayon $1$. Donc $R\geqslant1$. Mais d'autre part, pour tout entier naturel non nul $n$, $a_n=\sum_{k=1}^{n}\frac{1}{k}\geqslant1$ et $R\leqslant1$. Finalement $R = 1$. De plus, pour $x$ dans $]-1,1[$, $f(x)=\frac{1}{1-x}\times-\ln(1-x)=\frac{\ln(1-x)}{x-1}$.

\begin{center}
\shadowbox{
$\forall x\in\left]-1,1\right[$, $\sum_{n=1}^{+\infty}\left(\sum_{k=1}^{n}\frac{1}{k}\right)x^n=\frac{\ln(1-x)}{x-1}$.
}
\end{center}

\item  La règle de d'\textsc{Alembert} montre que le rayon de convergence est égal à $+\infty$.

Pour $n$ entier naturel donné, $\frac{n^2+4n-1}{n!(n+2)}=\frac{n^3+5n^2+3n-1}{(n+2)!}$ puis 

\begin{align*}\ensuremath
n^3+5n^2+3n-1&= (n+2)(n+1)n+2n^2+n-1 = (n+2)(n+1)n+2(n+2)(n+1) - 5n - 5\\
 &= (n+2)(n+1)n + 2(n+2)(n+1) - 5(n+2) + 5
\end{align*}

Donc, pour tout réel $x$,

\begin{center}
$f(x) =\sum_{n=0}^{+\infty}\frac{(n+2)(n+1)n}{(n+2)!}x^n +2\sum_{n=0}^{+\infty}\frac{(n+2)(n+1)}{(n+2)!}x^n- 5\sum_{n=0}^{+\infty}\frac{n+2}{(n+2)!}x^n+ 5\sum_{n=0}^{+\infty}\frac{1}{(n+2)!}x^n.$
\end{center}

Ensuite $f(0) = -\frac{1}{2}$  et pour $x\neq 0$,

\begin{align*}
f(x)&=\sum_{n=1}^{+\infty}\frac{1}{(n-1)!}x^n +2\sum_{n=0}^{+\infty}\frac{1}{n!}x^n- 5\sum_{n=0}^{+\infty}\frac{1}{(n+1)!}x^n+ 5\sum_{n=0}^{+\infty}\frac{1}{(n+2)!}x^n\\
 &= xe^x + 2e^x -5\frac{e^x-1}{x}+ 5\frac{e^x-1-x}{x^2}= \frac{e^x(x^3+2x^2-5x+5) -5x}{x^2}.
\end{align*}

\begin{center}
\shadowbox{
$\forall x\in\Rr$, $\sum_{n=0}^{+\infty}\frac{n^2+4n-1}{n!(n+2)}x^n=\left\{
\begin{array}{l}
\frac{e^x(x^3+2x^2-5x+5) -5x}{x^2}\;\text{si}\;x\in\Rr^*\\
-\frac{1}{2}\;\text{si}\;x=0
\end{array}
\right.$.
}
\end{center}

\item  Pour $n\in\Nn^*$,  $\frac{1}{n}\leqslant a_n=n^{(-1)^n}\leqslant n$ et donc $R = 1$. Pour $x$ dans $]-1,1[$, $f(x)=\sum_{k=0}^{+\infty}\frac{x^{2k+1}}{2k+1}+\sum_{k=1}^{+\infty}(2k)x^k$. Puis 

\begin{center}
$\sum_{k=1}^{+\infty}(2k)x^k= 2x\sum_{k=1}^{+\infty}kx^{k-1}= 2x\left(\sum_{k=0}^{+\infty}x^k\right)'= 2x\left(\frac{1}{1-x}\right)'=\frac{2x}{(1-x)^2}$.
\end{center}

\begin{center}
\shadowbox{
$\forall x\in]-1,1[$, $\sum_{n=1}^{+\infty}n^{(-1)^n}x^n=\Argth x+\frac{2x}{(1-x)^2}$.
}
\end{center}

\item  $R = 1$. Pour $x$ réel non nul dans $]-1,1[$, $f(x)=-\frac{1}{x}\sum_{n=1}^{+\infty}(-1)^{n-1}\frac{(x^4)^n}{n}= -\frac{\ln(1+x^4)}{4x}$ et sinon $f(0) = 0$.

\begin{center}
\shadowbox{
$\forall x\in]-1,1[$, $\sum_{n=1}^{+\infty}(-1)^n\frac{x^{4n-1}}{4n}=\left\{
\begin{array}{l}
-\frac{\ln(1+x^4)}{4x}\;\text{si}\;x\neq0\\
\rule{0mm}{4mm}0\;\text{si}\;x=0
\end{array}
\right.$.
}
\end{center}

\item  La règle de d'\textsc{Alembert} fournit $R=\frac{1}{2}$. Pour $x$ dans $\left]-\frac{1}{2},\frac{1}{2}\right[$,

\begin{align*}\ensuremath
\sum_{n=0}^{+\infty}(n^2+1)2^{n+1}x^n&=2\left(\sum_{n=0}^{+\infty}(n+2)(n+1)(2x)^n -3\sum_{n=0}^{+\infty}(n+1)(2x)^n +2\sum_{n=0}^{+\infty}(2x)^n\right)\\
  &= 2\left(\left(\sum_{n=0}^{+\infty}(2x)^n\right)''- 3\left(\sum_{n=0}^{+\infty}(2x)^n\right)'+2\sum_{n=0}^{+\infty}(2x)^n\right)
 = 2\left(2\frac{1}{1-2x}  - 3\frac{2}{(1-2x)^2}+\frac{4}{(1-2x)^3}\right)\\
 &=2\frac{2(1-2x)^2-6(1-2x)+8}{(1-2x)^3}=2\frac{8x^2+4x}{(1-2x)^3}.
\end{align*}

\item  Pour $x = 1$, la suite $((-1)^{n+1}nx^{2n+1})$ n'est pas bornée et donc $R\geqslant 1$. Mais la série converge si $|x| < 1$ et $R\leqslant 1$. Finalement $R=1$.

Pour $x$ dans $]-1,1[$,

\begin{align*}\ensuremath
\sum_{n=0}^{+\infty}(-1)^{n+1}nx^{2n+1}&=\frac{1}{2}\left(\sum_{n=0}^{+\infty}(-1)^{n+1}(2n+2)x^{2n+1}-2\sum_{n=0}^{+\infty}(-1)^{n+1}x^{2n+1}\right)\\ 
 &=\frac{1}{2}\left(\left(\sum_{n=0}^{+\infty}(-1)^{n+1}x^{2n+2}\right)'+2x\sum_{n=0}^{+\infty}(-x^2)^n\right)=\frac{1}{2}\left(\left(\frac{-x^2}{1+x^2}\right)'+\frac{2x}{1+x^2}\right)\\
 &= -\frac{x(1+x^2)-x^3}{1+x^2)^2}+\frac{x}{1+x^2}=\frac{x^3}{1+x^2)^2}.
\end{align*}

\begin{center}
\shadowbox{
$\forall x\in]-1,1[$, $\sum_{n=0}^{+\infty}(-1)^{n+1}nx^{2n+1}=\frac{x^3}{1+x^2)^2}$.
}
\end{center}

\item  \textbf{1ère solution.} Les racines de l'équation caractéristique $z^2 - z - 1 = 0$ sont $\alpha=\frac{1+\sqrt{5}}{2}$ et $\beta=\frac{1-\sqrt{5}}{2}$. On sait qu'il existe deux nombres réels $\lambda$ et $\mu$ tels que pour tout entier naturel $n$,

\begin{center}
$a_n =\lambda\left(\frac{1+\sqrt{5}}{2}\right)^n+\mu\left(\frac{1-\sqrt{5}}{2}\right)^n$.
\end{center}

Les égalités $n=0$ et $n=1$ fournissent

\begin{center}
$\left\{
\begin{array}{l}
\rule[-2mm]{0mm}{0mm}\lambda+\mu=1\\
\frac{1+\sqrt{5}}{2}\lambda+\frac{1-\sqrt{5}}{2}\mu=1
\end{array}
\right.\Leftrightarrow\left\{
\begin{array}{l}
\rule[-2mm]{0mm}{0mm}\lambda+\mu=1\\
\lambda-\mu=\frac{1}{\sqrt{5}}
\end{array}
\right.\Leftrightarrow
\left\{
\begin{array}{l}
\lambda=\frac{1}{2}\left(1+\frac{1}{\sqrt{5}}\right)\\
\mu=\frac{1}{2}\left(1-\frac{1}{\sqrt{5}}\right)\end{array}
\right.\Leftrightarrow\left\{
\begin{array}{l}
\lambda=\frac{1}{\sqrt{5}}\frac{1+\sqrt{5}}{2}\\
\mu=-\frac{1}{\sqrt{5}}\frac{1-\sqrt{5}}{2}
\end{array}
\right.$.
\end{center}

Finalement, pour tout entier naturel $n$, $a_n =\frac{1}{\sqrt{5}}\left(\frac{1+\sqrt{5}}{2}\right)^{n+1}-\frac{1}{\sqrt{5}}\left(\frac{1-\sqrt{5}}{2}\right)^{n+1}$.

Les séries entières respectivement associées aux suites $\left(\left(\frac{1+\sqrt{5}}{2}\right)^{n+1}\right)$ et $\left(\left(\frac{1-\sqrt{5}}{2}\right)^{n+1}\right)$ ont pour rayons respectifs $\left|\frac{1}{(1+\sqrt{5})/2}\right|=\frac{\sqrt{5}-1}{2}$ et $\left|\frac{1}{(1-\sqrt{5})/2}\right|=\frac{\sqrt{5}+1}{2}$. Ces rayons étant distincts, la série proposée a pour rayon

\begin{center}
$R=\text{Min}\left\{\frac{\sqrt{5}-1}{2},\frac{\sqrt{5}+1}{2}\right\}=\frac{\sqrt{5}-1}{2}$.
\end{center}

Pour $x$ dans $\left]-\frac{\sqrt{5}-1}{2},\frac{\sqrt{5}-1}{2}\right[$, on a

\begin{align*}\ensuremath
\sum_{n=0}^{+\infty}a_nx^n&=\frac{\alpha}{\sqrt{5}}\sum_{n=0}^{+\infty}(\alpha x)^n -\frac{\beta}{\sqrt{5}}\sum_{n=0}^{+\infty}(\beta x)^n
=\frac{1}{\sqrt{5}}\left(\frac{\alpha}{1-\alpha x}-\frac{\beta}{1-\beta x}\right)
=\frac{\alpha-\beta}{\sqrt{5}}\times\frac{1}{\alpha\beta x^2-(\alpha+\beta)x+1}\\
 &=\frac{1}{1-x-x^2}.
\end{align*}

\textbf{2ème solution.} Supposons à priori le rayon $R$ de la série proposée strictement positif. Pour $x$ dans $]-R,R[$, on a

\begin{align*}\ensuremath
f(x)&= 1+ x +\sum_{n=2}^{+\infty}a_nx^n= 1 + x +\sum_{n=0}^{+\infty}a_{n+2}x^{n+2}=1+x+\sum_{n=0}^{+\infty}(a_{n+1}+a_n)x^{n+2}\\
 &= 1 + x + x\sum_{n=0}^{+\infty}a_{n+1}x^{n+1}+ x^2\sum_{n=0}^{+\infty}a_nx^n\;(\text{les deux séries ont même rayon})\\ 
 &=1+x+x(f(x)-1)+ x^2 f(x).
\end{align*}

Donc, nécessairement $\forall x\in]-R,R[$, $f(x) =\frac{1}{1-x-x^2}$.

Réciproquement, la fraction rationnelle ci-dessus n'admet pas $0$ pour pôle et est donc développable en série entière. Le rayon de convergence de la série obtenue est le minimum des modules des pôles de $f$ à savoir $R=\frac{\sqrt{5}-1}{2}$. Notons $\sum_{n=0}^{+\infty}b_nx^n$ ce développement. Pour tout $x$ de $]-R,R[$, on a $\left(\sum_{n=0}^{+\infty}b_nx^n\right)(1-x-x^2)=1$ et donc $\sum_{n=0}^{+\infty}b_nx^n -\sum_{n=0}^{+\infty}b_nx^{n+1}-\sum_{n=0}^{+\infty}b_nx^{n+2}= 1$ ce qui s'écrit encore  $\sum_{n=0}^{+\infty}b_nx^n-\sum_{n=1}^{+\infty}b_{n-1}x_n -\sum_{n=2}^{+\infty}b_{n-2}x^n = 1$. Finalement 

\begin{center}
$\forall x\in]-R,R[$, $b_0 +(b_1-b_0)x+\sum_{n=2}^{+\infty}(bn-b_{n-1}-b_{n-2})x^n = 1$.
\end{center}

Par unicité des coefficients d'un développement en série entière, on a alors $b_0=b_1= 1$ et $\forall n\geqslant2$, $b_n = b_{n-1}+b{n-2}$. On en déduit alors par récurrence que $\forall n\in\Nn$, $b_n = a_n$.

\begin{center}
\shadowbox{
$\forall x\in\left]-\frac{\sqrt{5}-1}{2},\frac{\sqrt{5}-1}{2}\right[$, $\sum_{n=0}^{+\infty}a_nx^n=\frac{1}{1-x-x^2}$.
}
\end{center}

\textbf{Remarque.} En généralisant le travail précédent, on peut montrer que les suites associées aux développements en série entière des fractions rationnelles sont justement les suites vérifiant des relations de récurrence linéaire.

\item  Pour tout entier naturel $n$, $1\leqslant a_n\leqslant n+1$. Donc $R = 1$.

On remarque que pour tout entier naturel $n$, $a_n =\sum_{k+5l=n}^{}1$. La série entière proposée est donc le produit de \textsc{Cauchy} des séries $\sum_{k=0}^{+\infty}x^k$ et $\sum_{l=0}^{+\infty}x^{5l}$. Pour $x$ dans $]-1,1[$, on a donc

\begin{center}
$f(x) =\left(\sum_{k=0}^{+\infty}x^k\right)\left(\sum_{l=0}^{+\infty}x^{5l}\right)=\frac{1}{1-x}\times\frac{1}{(1-x)^5}$.
\end{center}

\textbf{Remarque.} De combien de façons peut -on payer $100$ euros avec des pièces de $1$, $2$, $5$, $10$, $20$ et $50$ centimes d'euros, des pièces de $1$ et $2$ euros et des billets de $10$ et $20$ euros ? Soit $N$ le nombre de solutions. $N$ est le nombre de solutions en nombres entiers a,b,... de l'équation

\begin{center}
$a +2b +5c +10d +20e +50f +100g +200h+500k+1000i+2000j=10000$
\end{center}

et est donc le coefficient de $x^{10000}$ du développement en série entière de 

\begin{center}
$\frac{1}{(1-x)(1-x^2)(1-x^5)(1-x^{10})(1-x^{20})(1-x^{50})(1-x^{100})(1-x^{200})(1-x^{500})(1-x^{1000})(1-x^{2000})}$,
\end{center}

La remarque est néanmoins anecdotique et il semble bien préférable de dénombrer à la main le nombre de solutions. Les 
exercices \ref{ex:rou19x} et \ref{ex:rou20x} de cette planche font bien mieux comprendre à quel point les séries entières sont un outil intéressant pour les dénombrements.
\end{enumerate}
\fincorrection
\correction{005752}
Pour tout entier naturel non nul, $|a_n|\leqslant\frac{1}{n}$ et donc $R\geqslant1$. Mais si $x > 1$, la suite $\left(\frac{1}{n}\cos\left(\frac{2n\pi}{3}\right)x^n\right)_{n\geqslant1}$ n'est pas bornée comme on le voit en considérant la suite extraite des termes d'indices multiples de $3$ et donc $R = 1$. Pour $x$ dans $]-1,1[$, $f(x)=\text{Re}\left(\sum_{n=1}^{+\infty}\frac{(jx)^n}{n}\right)$. Le problème est alors de ne pouvoir écrire $-\ln(1-jx)$. Il faut s'y prendre autrement.

$f$ est donc dérivable sur $]-1,1[$ et pour $x$ dans $]-1;1[$,

\begin{center}
$f'(x)=\sum_{n=1}^{+\infty}\cos\left(\frac{2n\pi}{3}\right)x^{n-1}=\text{Re}\left(\sum_{n=0}^{+\infty}j^nx^{n-1}\right) =\text{Re}\left(\frac{j}{1-jx}\right)=\text{Re}\left(\frac{j(1-j^2x)}{x^2+x+1}\right) = -\frac{1}{2}\frac{2x+1}{x^2+x+1}$.
\end{center}

Par suite, pour $x\in]-1,1[$, $f(x) = f(0)+\int_{0}^{x}f'(t)\;dt= -\frac{1}{2}\ln(x^2+x+1)$.

\begin{center}
\shadowbox{
$\forall x\in]-1,1[$, $\sum_{n=1}^{+\infty}\frac{1}{n}\cos\left(\frac{2n\pi}{3}\right)x^{n}=-\frac{1}{2}\ln(x^2+x+1)$.
}
\end{center}
\fincorrection
\correction{005753}
Le rayon de la série considérée est égal $1$. Soit $x\in]-1,1[$.

\begin{center}
$f(x) =\frac{1}{2}\sum_{n=0}^{+\infty}\left(\frac{1}{2n-1}-\frac{1}{2n+1}\right)x^n =\frac{1}{2}\left(-1 +\sum_{n=1}^{+\infty}\frac{x^n}{2n-1}-\sum_{n=0}^{+\infty}\frac{x^n}{2n+1}\right)$.
\end{center}

\textbullet~Si $x$ est dans $]0,1[$,

\begin{align*}\ensuremath
f(x)&=\frac{1}{2}\left(-1 +\sum_{n=0}^{+\infty}\frac{x^{n+1}}{2n+1}-\sum_{n=0}^{+\infty}\frac{x^n}{2n+1}\right)=\frac{1}{2}\left(-1+\left(\sqrt{x}-\frac{1}{\sqrt{x}}\right)\sum_{n=0}^{+\infty}\frac{\left(\sqrt{x}\right)^{2n+1}}{2n+1}\right)\\
 &=\frac{1}{2}\left(-1+\left(\sqrt{x}-\frac{1}{\sqrt{x}}\right)\ln\left(\frac{1+\sqrt{x}}{1-\sqrt{x}}\right)\right).
\end{align*}

\textbullet~Si $x$ est dans $]-1,0[$,

\begin{align*}\ensuremath
f(x)&=\frac{1}{2}\left(-1 +\sum_{n=0}^{+\infty}\frac{x^{n+1}}{2n+1}-\sum_{n=0}^{+\infty}\frac{x^n}{2n+1}\right)=\frac{1}{2}\left(-1-\left(\sqrt{-x}+\frac{1}{\sqrt{-x}}\right)\sum_{n=0}^{+\infty}(-1)^n\frac{\left(\sqrt{-x}\right)^{2n+1}}{2n+1}\right)\\
 &=\frac{1}{2}\left(-1-\left(\sqrt{-x}+\frac{1}{\sqrt{-x}}\right)\Arctan(\sqrt{-x})\right).
\end{align*}

\textbullet~$f(0)=-1$.

Maintenant, la somme est en fait définie sur $[-1,1]$ car les séries numériques de termes généraux $\frac{1}{4n^2-1}$ et $\frac{(-1)^n}{4n^2-1}$ convergent. Vérifions que la somme est continue sur $[-1,1]$.

Pour $x$ dans $[-1,1]$ et $n\in\Nn^*$, $\left|\frac{x^n}{4n^2-1}\right|\leqslant\frac{1}{4n^2-1}$ qui est le  terme général d'une série numérique convergente. La série entière considérée converge donc normalement sur $[-1,1]$. On en déduit que cette somme est continue sur $[-1,1]$. Donc 

\begin{align*}\ensuremath
 \sum_{n=0}^{+\infty}\frac{1}{4n^2-1}&=f(1)=\displaystyle\lim_{\substack{x\rightarrow1\\ x<1}}f(x)
=\displaystyle\lim_{\substack{x\rightarrow1\\ x<1}}\frac{1}{2}\left(-1+\left(\sqrt{x}-\frac{1}{\sqrt{-x}} \right)\ln\left(1+\sqrt{x}\right) +\frac{1}{\sqrt{x}}(1+\sqrt{x})(1-\sqrt{x})\ln(1-\sqrt{x})\right)\\
 &=-\frac{1}{2}
\end{align*} 
			

\textbf{Remarque.} $\sum_{n=0}^{+\infty}\frac{1}{4n^2-1}=\lim_{n \rightarrow +\infty}\left(\frac{1}{2}\sum_{k=0}^{n}\left(\frac{1}{2k-1}-\frac{1}{2k+1}\right)\right)=\lim_{n \rightarrow +\infty}\left(\frac{1}{2}\left(-1-\frac{1}{2n+1}\right)\right)=-\frac{1}{2}$ (série télescopique).

On a aussi

\begin{align*}\ensuremath
 \sum_{n=0}^{+\infty}\frac{(-1)^n}{4n^2-1}&=f(-1)=\displaystyle\lim_{\substack{x\rightarrow-1\\ x>-1}}f(x)
=\displaystyle\lim_{\substack{x\rightarrow-1\\ x>-1}}\frac{1}{2}\left(-1-\left(\sqrt{-x}+\frac{1}{\sqrt{-x}} \right)\Arctan(\sqrt{-x})\right)\\
 &=\frac{1}{2}\left(-1-2\Arctan1\right)=-\frac{\pi+2}{4}.
\end{align*} 

\fincorrection
\correction{005754}
Pour tout entier naturel $n$, $|a_n|\geqslant\frac{1}{2n+1}$ et donc la série proposée ne converge pas absolument.

Pour tout entier naturel $n$,

\begin{align*}\ensuremath
|u_n| - |u_{n+1}|&=\frac{1}{2n+1}\sum_{k=0}^{n}\frac{1}{4k+1}-\frac{1}{2n+3}\sum_{k=0}^{n+1}\frac{1}{4k+1}=\left(\frac{1}{2n+1}-\frac{1}{2n+3}\right)\sum_{k=0}^{n}\frac{1}{4k+1}-\frac{1}{2n+3}\times\frac{1}{4n+5}\\
 &=\frac{2}{(2n+1)(2n+3)}\sum_{k=0}^{n}\frac{1}{4k+1}-\frac{1}{(2n+3)(4n+5)}\geqslant\frac{1}{(2n+3)(2n+1)}-\frac{1}{(2n+3)(4n+5)}> 0.
\end{align*}

 
La suite $(|u_n|)_{n\in\Nn}$ est donc décroissante. De plus, pour tout entier naturel non nul $n$,

\begin{center}
$\sum_{k=0}^{n}\frac{1}{4k+1}\leqslant\sum_{k=1}^{4n+1}\frac{1}{k}\leqslant1+\sum_{k=2}^{4n+1}\int_{k-1}^{k}\frac{1}{t}\;dt=1 +\ln(4n+1)$
\end{center}

et donc $|u_n|\leqslant\frac{1 +\ln(4n+1)}{2n+1}$\rule[-4mm]{0mm}{0mm}. On en déduit que $\lim_{n \rightarrow +\infty}u_n=0$. Finalement, la série proposée converge en vertu du critère spécial aux séries alternées.

Considérons la série entière $\sum_{n=0}^{+\infty}u_n x^{2n+1}$.  La série de terme général $a_n$ converge et donc $R\geqslant1$ mais puisque la série de terme général $|a_n|$ diverge et donc $R\leqslant1$. Finalement, $R=1$. Pour $x\in]-1,1[$, posons $f(x)=\sum_{n=0}^{+\infty}u_n x^{2n+1}$. Pour $x$ dans $]-1,1[$, 

\begin{align*}\ensuremath
f'(x)&=\sum_{n=0}^{+\infty}(-1)^n\left(\sum_{k=0}^{n}\frac{1}{4k+1}\right)x^{2n}=\sum_{n=0}^{+\infty}\left(\sum_{k=0}^{n}\frac{1}{4k+1}\right)(-x^2)^n\\
 &=\left(\sum_{n=0}^{+\infty}\frac{(-x^2)^n}{4n+1}\right)\left(\sum_{n=0}^{+\infty}(-x^2)^n\right)\;(\text{produit de \textsc{Cauchy} de deux séries numériques absolument convergentes})
\end{align*}

Donc, pour $x$ dans $]0,1[$, $f'(x)  = g(x)h(x)$ où $h(x) =\sum_{n=0}^{+\infty}(-x^2)^n =\frac{1}{1+x^2}$ puis 

\begin{center}
$g(x) =\frac{1}{\sqrt{x}}\sum_{n=0}^{+\infty}(-1)^n\frac{1}{4n+1}(\sqrt{x})^{4n+1}$.
\end{center}

Maintenant, en posant $k(X)=\sum_{n=0}^{+\infty}(-1)^n X^{4n+1}$ pour $X$ dans $]-1,1[$, $k'(X)=\sum_{n=0}^{+\infty}(-1)^nX^{4n}=\frac{1}{X^4+1}$. Ensuite, en posant $\omega=e^{i\pi/4}$, par réalité et parité

\begin{center}
$\frac{1}{X^4+1}=\frac{a}{X-\omega}+\frac{\overline{a}}{X-\overline{\omega}}-\frac{a}{X+\omega}-\frac{\overline{a}}{X+\overline{\omega}}$
\end{center}

où $a=\frac{1}{4\omega^3}=-\frac{\omega}{4}$. Il vient alors

\begin{align*}\ensuremath
\frac{1}{X^4+1}&=-\frac{1}{4}\left(\frac{\omega}{X-\omega}+\frac{\overline{\omega}}{X-\overline{\omega}}-\frac{\omega}{X+\omega}-\frac{\overline{\omega}}{X+\overline{\omega}}\right)=\frac{1}{4}\left(-\frac{X\sqrt{2}-2}{X^2-\sqrt{2}X+1}+\frac{X\sqrt{2}+2}{X^2+\sqrt{2}X+1}\right)\\
 &=\frac{1}{4\sqrt{2}}\left(\frac{2X+2\sqrt{2}}{X^2+\sqrt{2}X+1}-\frac{2X-2\sqrt{2}}{X^2-\sqrt{2}X+1}\right)\\
 &=\frac{1}{4\sqrt{2}}\left(\frac{2X+\sqrt{2}}{X^2+\sqrt{2}X+1}+\frac{\sqrt{2}}{\left(X+\frac{1}{\sqrt{2}}\right)^2+\left(\frac{1}{\sqrt{2}}\right)^2}-\frac{2X-\sqrt{2}}{X^2-\sqrt{2}X+1}+\frac{\sqrt{2}}{\left(X-\frac{1}{\sqrt{2}}\right)^2+\left(\frac{1}{\sqrt{2}}\right)^2}\right)
\end{align*}
 

En tenant compte de $k(0) = 0$, on obtient donc pour $X\in]-1,1[$,

\begin{center}
$k(X) =\frac{1}{4\sqrt{2}}\left(\ln(X^2+X\sqrt{2}+1)-\ln(X^2-X\sqrt{2}+1))+2\left(\Arctan(X\sqrt{2}+1)+\Arctan(X\sqrt{2}-1)\right)\right)$.
\end{center}

Ensuite, pour tout réel $x\in]0,1[$, $f'(x) =\frac{1}{\sqrt{x}}k\left(\sqrt{x}\right)\frac{1}{1+x^2}=\frac{1}{\sqrt{x}}k'\left(\sqrt{x}\right)k\left(\sqrt{x}\right)$ et donc 

\begin{align*}\ensuremath
f(x)&=f(0)+\left(k\left(\sqrt{x}\right)^2 -k(0)^2\right)=k\left(\sqrt{x}\right)^2\\
 &=\frac{1}{32}\left(\ln(X^2+X\sqrt{2}+1)-\ln(X^2-X\sqrt{2}+1))+2\left(\Arctan(X\sqrt{2}+1)+\Arctan(X\sqrt{2}-1)\right)\right)^2.
\end{align*}

Quand $x$ tend vers $1$, $f(x)$ tend vers 

\begin{center}
$\frac{1}{32}
\left(\ln\left(\frac{2+\sqrt{2}}{2-\sqrt{2}}\right)+ 2 (\Arctan (\sqrt{2}+1)+\Arctan(\sqrt{2}-1))\right)^2=  \frac{1}{32}\left(\ln(3+2\sqrt{2}) +\pi\right)^2$.
\end{center}

(car $\Arctan (\sqrt{2}+1)+\Arctan(\sqrt{2}-1)=\Arctan (\sqrt{2}+1)+\Arctan\left(\frac{1}{\sqrt{2}+1}\right)=\frac{\pi}{2}$).

Enfin, pour $x$ dans $[0,1]$ et $n$ dans $\Nn$, $|u_n|x^n - |u_{n+1}|x^{n+1}\geqslant(|u_n| - |u_{n+1}|)x^n\geqslant 0$ et la série numérique de terme général $u_nx^n$ est alternée. D'après une majoration classique du reste à l'ordre $n$ d'une telle série, pour tout entier naturel $n$ et tout réel $x$ de $[0,1]$,

\begin{center}
$|R_n(x)|=\left|\sum_{k=n+1}^{+\infty}u_kx^k\right|\leqslant\left|u_{n+1}x^{n+1}\right|\leqslant |u_{n+1}|$,
\end{center}

et donc $\underset{x\in[0,1]}{\text{Sup}}|R_n(x)|\leqslant|a_{n+1}|\underset{n\rightarrow+\infty}{\rightarrow}0$. La convergence est uniforme sur $[0,1]$ et on en déduit que la somme est continue sur $[0,1]$. En particulier 

\begin{center}
$\sum_{n=0}^{+\infty}u_n=f(1)=\displaystyle\lim_{\substack{x\rightarrow1\\ x<1}}f(x)=\frac{1}{32}\left(\ln(3+2\sqrt{2}) +\pi\right)^2$.
\end{center}

\begin{center}
\shadowbox{
$\sum_{n=0}^{+\infty}\left(\frac{(-1)^n}{2n+1}\sum_{k=0}^{n}\frac{1}{4k+1}\right)=\frac{1}{32}\left(\ln(3+2\sqrt{2}) +\pi\right)^2$.
}
\end{center}
\fincorrection
\correction{005757}
Pour tout entier naturel $n$, $a_{n+1}+b_{n+1}= 2(a_n + b_n)$ et $3a_{n+1}+ 2b_{n+1}= 3a_n+2b_n$ (rappel : ces combinaisons linéaires sont fournies par les vecteurs propres de ${^t}A$ si on ne les devine pas). On en déduit que pour tout entier naturel $n$, $a_n+b_n = 2^n(a_0+b_0) = 2^n$ et 
$3a_n+2b_n = 3a_0+2b_0 = 3$. Finalement,

\begin{center}
$\forall n\in\Nn$, $a_n =3-2^{n+1}$ et $b_n = 3(2^n-1)$.
\end{center}

Les deux séries proposées sont alors clairement de rayons infini et pour tout réel $x$, $f(x) =3e^x- 2e^{2x}$ et $g(x) =3(e^{2x}-e^x)$. (On peut avoir d'autres idées de résolution, plus astucieuses, mais au bout du compte moins performantes).
\fincorrection
\correction{005758}
Pour $n\geqslant1$, posons $a_n =\frac{1}{nC_{2n}^n}$. Pour $n\in\Nn^*$,

\begin{center}
$\left|\frac{a_{n+1}}{a_n}\right|=\frac{n}{n+1}\times\frac{(2n)!}{(2n+2)!}\times\frac{(n+1)!^2}{n!^2}=\frac{n}{2(2n+1)}$\quad$(*)$.
\end{center}

Par suite, $\left|\frac{a_{n+1}}{a_n}\right|\underset{n\rightarrow+\infty}{\rightarrow}\frac{1}{4}$ et d'après la règle de d'\textsc{Alembert}, le rayon de la série entière considérée est $R = 4$.

Pour $x\in\left]-4,4\right[$, posons $f(x)=\sum_{n=1}^{+\infty}a_nx^n$.

Les relations $(*)$ s'écrivent encore $\forall n\in\Nn^*$, $4(n+1)a_{n+1}-2a_{n+1}=na_n$.

Soit $x\in]-4,4[$. On multiplie les deux membres de l'égalité précédente par $x^{n+1}$ et on somme sur $n$. On obtient

\begin{center}
$4x\sum_{n=1}^{+\infty}(n+1)a_{n+1}x^n-2\sum_{n=1}^{+\infty}a_{n+1}x^{n+1}=x^2\sum_{n=1}^{+\infty}na_nx^{n-1}$,
\end{center}

ou encore $x^2f'(x) =4x(f'(x)-a_1)-2(f(x)-a_1x)$ ou encore $x(x-4)f'(x)+2f(x)=-x$\quad$(E)$. Soit $I$ l'un des deux intervalles $]-4,0[$ ou $]0,4[$.Sur $I$, l'équation $(E)$ s'écrit :

\begin{center}
$f'(x)+\frac{1}{2}\left(\frac{1}{x}-\frac{1}{x-4}\right)f(x)=-\frac{1}{x-4}$.
\end{center}

Une primitive sur $I$ de la fonction $a~:~x\mapsto\frac{1}{2}\left(\frac{1}{x}-\frac{1}{x-4}\right)$ est la fonction $A~:~x\mapsto\frac{1}{2}(\ln|x-4|-\ln|x|)=\ln\sqrt{\frac{|x-4|}{|x|}}$. 

\begin{align*}\ensuremath
f\;\text{solution de}\;(E)\;\text{sur}\;I&\Leftrightarrow\forall x\in I,\;f'(x)+\frac{1}{2}\left(\frac{1}{x}-\frac{1}{x-4}\right)f(x)=-\frac{1}{x-4}\\
 &\Leftrightarrow\forall x\in I,\;e^{A(x)}f'(x)+a(x)e^{A(x)}f(x)=\frac{1}{4-x}\sqrt{\frac{|x-4|}{|x|}}\\
 &\Leftrightarrow\forall x\in I,\;(e^Af)'(x)=\frac{1}{\sqrt{|x(x-4)|}}\quad(*).
\end{align*}

Déterminons une primitive de la fonction $x\mapsto\frac{1}{\sqrt{|x(x-4)|}}$ sur $I$.

\textbullet~Si $I=]0,4[$, $\frac{1}{\sqrt{|x(x-4)|}}=\frac{1}{\sqrt{x(4-x)}}=\frac{1}{\sqrt{4-(x-2)^2}}$ et une primitive de la fonction $x\mapsto\frac{1}{\sqrt{|x(x-4)|}}$ sur $I$ est la fonction $x\mapsto\Arcsin\left(\frac{x-2}{2}\right)$. Puis

\begin{align*}\ensuremath
f\;\text{solution de}\;(E)\;\text{sur}\;I&\Leftrightarrow\exists C\in\Rr/\;\forall x\in I,\;e^{A(x)}f(x)=\Arcsin\left(\frac{x-2}{2}\right)+C\\
 &\Leftrightarrow\exists C\in\Rr/\;\forall x\in I,\;f(x)=\sqrt{\frac{x}{4-x}}\left(\Arcsin\left(\frac{x-2}{2}\right)+C\right).
\end{align*}

\textbullet~Si $I=]-4,0[$, $\frac{1}{\sqrt{|x(x-4)|}}=\frac{1}{\sqrt{x(x-4)}}=\frac{1}{\sqrt{(2-x)^2-4}}$ et une primitive de la fonction $x\mapsto\frac{1}{\sqrt{|x(x-4)|}}$ sur $I$ est la fonction $x\mapsto-\Argch\left(\frac{2-x}{2}\right)$. Puis

\begin{align*}\ensuremath
f\;\text{solution de}\;(E)\;\text{sur}\;I&\Leftrightarrow\exists C'\in\Rr/\;\forall x\in I,\;e^{A(x)}f(x)=\Argch\left(\frac{2-x}{2}\right)+C'\\
 &\Leftrightarrow\exists C'\in\Rr/\;\forall x\in I,\;f(x)=\sqrt{\frac{x}{x-4}}\left(-\Argch\left(\frac{2-x}{2}\right)+C'\right).
\end{align*}

$f$ doit être définie, continue et dérivable sur $]-4,4[$ et en particulier dérivable en $0$. Ceci impose $\lim_{x \rightarrow 0^+}\Arcsin\left(\frac{x-2}{2}\right)+C=0$ (car sinon $f(x)\underset{0^+}{\sim}C\sqrt{x}$) et donc $C=\frac{\pi}{2}$. Pour $x\in]0,4[$, on a alors $f(x)=\sqrt{\frac{x}{4-x}}\left(\frac{\pi}{2}-\Arcsin\left(\frac{2-x}{2}\right)\right)=\sqrt{\frac{x}{4-x}}\Arccos\left(\frac{2-x}{2}\right)$ ce qui reste vrai pour $x=0$ par continuité..

De même, $\lim_{x \rightarrow 0^-}-\Argch\left(\frac{2-x}{2}\right)+C'=0$ et donc $C'=0$. On a montré que

\begin{center}
\shadowbox{
$\forall x\in]-4,4[,\;\sum_{n=1}^{+\infty}\frac{1}{nC_{2n}^n}x^n=\left\{
\begin{array}{l}
\sqrt{\frac{x}{4-x}}\Arccos\left(\frac{2-x}{2}\right)\;\text{si}\;x\in[0,4[\\
-\sqrt{\frac{x}{x-4}}\Argch\left(\frac{2-x}{2}\right)\;\text{si}\;x\in]-4,0]
\end{array}
\right.
$.
}
\end{center}
\fincorrection
\correction{005763}
On a $I_0 = 0$, $I_1 = 1$ et $I_2 = 2$ (l'identité et la transposition $\tau_{1,2}$).

Soit $n\in\Nn^*$. Il y a $I_{n+1}$ involutions $\sigma$ de $\llbracket1,n+2\rrbracket$ vérifiant $\sigma(n+2) =n+2$ car la restriction d'une telle permutation à $\llbracket1,n+1\rrbracket$ est une involution de $\llbracket1,n+1\rrbracket$ et réciproquement.

Si $\sigma(n+2) = k\in\llbracket1,n+1\rrbracket$, nécessairement $\sigma(k) =n+2$ puis la restriction de $\sigma$ à $\llbracket1,n+2\rrbracket\setminus\{k,n+2\}$ est une involution et réciproquement  Il y a $I_n$ involutions de $\llbracket1,n+2\rrbracket\setminus\{k,n+2\}$ et $n+1$ choix possibles de $k$ et donc $(n+1)I_n$ involutions de $\llbracket1,n+2\rrbracket$ telles que $\sigma(n+2)\neq n+2$. En résumé, 

\begin{center}
$\forall n\in\Nn^*$, $I_{n+2}= I_{n+1}+(n+1)I_n$.
\end{center}

Le rayon $R$ de la série entière associée à la suite $\left(\frac{I_n}{n!}\right)_{n\in\Nn^*}$ est supérieur ou égal à $1$ car $\forall n\in\Nn^*$, $\frac{I_n}{n!}\leqslant1$. Pour $x$ dans $]-R,R[$, posons $f(x)=\sum_{n=1}^{+\infty}\frac{I_n}{n!}x^n$. $f$ est dérivable sur $]-R,R[$ et pour $x\in]-R,R[$

\begin{align*}\ensuremath
f'(x)&=\sum_{n=1}^{+\infty}\frac{I_n}{(n-1)!}x^{n-1}=1+2x+\sum_{n=1}^{+\infty}\frac{I_{n+2}}{(n+1)!}x^{n+1}=1+2x+\sum_{n=1}^{+\infty}\frac{I_{n+1}+(n+1)I_n}{(n+1)!}x^{n+1}\\
 &=1+2x+\sum_{n=2}^{+\infty}\frac{I_n}{n!}x^n+x\sum_{n=1}^{+\infty}\frac{I_n}{n!}x^n\\
  &=1+2x+f(x)-x+xf(x)=1+x+(x+1)f(x).
\end{align*}

Donc, pour $x\in]-R,R[$, $f'(x)+(x+1)f(x)=x+1$ ou encore $e^{\frac{x^2}{2}+x}f'(x)+(x+1)e^{\frac{x^2}{2}+x}f(x)=(x+1)e^{\frac{x^2}{2}+x}$. Par suite, pour $x\in]-R,R[$,

\begin{center}
$e^{\frac{x^2}{2}+x}f(x)-f(0)=\int_{0}^{x}(t+1)e^{\frac{t^2}{2}+t}\;dt=e^{\frac{x^2}{2}+x}-1$,
\end{center}

et puisque $f(0)=0$, $\forall x\in]-R,R[$, $f(x)=e^{\frac{x^2}{2}+x}-1$.

Réciproquement, la fonction précédente est développable en série entière sur $\Rr$ en vertu de théorèmes généraux ($=e^{\frac{x^2}{2}}\times e^x$) et les coefficients de ce développement vérifient les relations définissant $\frac{I_n}{n!}$ de manière unique. Donc, ces coefficients sont les $\frac{I_n}{n!}$ ce qui montre que $R = +\infty$.

\begin{center}
\shadowbox{
$\forall x\in\Rr$, $\sum_{n=1}^{+\infty}\frac{I_n}{n!}x^n=e^{\frac{x^2}{2}+x}-1$.
}
\end{center}
\fincorrection
\correction{004574}
\begin{enumerate}
  \item $= \ln(1-x^3) - \ln(1-x)
              = \sum_{n=0}^\infty\left( \frac{x^{3n+1}}{3n+1}
              + \frac{x^{3n+2}}{3n+2} - 2\frac{x^{3n+3}}{3n+3}\right)$.
  \item Factoriser : $-\ln 6 + \left(\frac56 + \ln6\right)x
             - \sum_{n=2}^\infty\left(\frac{n+1}{2^n}
             + \frac{2n+1}{3^n}\right)\frac{x^n}{n(n-1)}$.
  \item Dériver le ln :
             $\sum_{n=1}^\infty {\binom{-1/2}{n-1}}\frac{x^{2n}}{2n-1}$.
  \item $\sum_{n=0}^\infty -\frac{2n+ 5 + 3(-1)^n}4 x^n$.
  \item $\frac15\sum_{n=0}^\infty \Bigl(1+2{\sqrt2\,}^n
              (2\cos(3n\pi/4)-\sin(3n\pi/4))\Bigr)x^n$.
  \item Intégrer : $\sum_{n=0}^\infty \frac{n+1}{4\sqrt 2}
             \Bigl((-\sqrt2-1)^{n+2} - (\sqrt2-1)^{n+2}\Bigr)x^n$.
  \item $=\frac{1-x}{\sqrt{1-x^2}}
              =\sum_{n=0}^\infty {\binom{-1/2}{n}}(-1)^n(x^{2n}-x^{2n+1})$.
  \item Dériver : $\frac\pi4 - \sum_{n=1}^\infty
                      \frac{\sin(n\pi/4)}{n{\sqrt2\,}^n}(-1)^nx^n$.
  \item Dériver : $\frac\pi3 + \sum_{n=1}^\infty
             (-1)^{n-1}\frac{\sin(n\pi/6)}{n2^n}x^n$.
  \item Dériver, factoriser : $\sum_{n=1}^\infty
             -\frac{2^n+2^{-n}}{n^2}x^n$.
  \item Linéariser : $1 + \sum_{n=1}^\infty \frac{(-1)^{n+1}4^n}{(2n)!}
             \Bigl(x^{2n-1} + \frac{(2n^2+3n-1)}{(2n+1)(2n+2)}x^{2n}\Bigr)$.
  \item Dériver : $\sum_{n=0}^\infty
             \frac{(-1)^n(2^{2n+1}-1)}{n!\,(2n+1)}x^{2n+1}$.
  \item $y' = -4xy+1$ : $\sum_{n=0}^\infty
             (-1)^n\frac{8^nn!}{(2n+1)!}x^{2n+1}$.
  \item $2x(1-x)y'+(1-2x)y = 1$ : $\sum_{n=0}^\infty
             \frac{4^n}{(2n+1)C_{2n}^n}x^n$.
  \item $(1-x^2)y'' - xy' + \frac y9 = 0$ : $\sum_{n=0}^\infty
              \frac{4^nC_{3n}^n}{(2n+1)3^{3n+1}}x^{2n+1}$.
\end{enumerate}
 \fincorrection
\correction{004575}
$=\frac12\ln(e^a-x) + \frac12\ln(e^{-a}-x)
=\sum_{k=1}^\infty \frac{\ch(ka)}kx^k$.
\fincorrection
\correction{004576}
$\frac{e^{x^2}}{1-x} = (1+x)\frac{e^{x^2}}{1-x^2}
          = \sum_{n=0}^\infty \left(1+\frac1{1!} + \dots
          + \frac1{n!}\right)(x^{2n}+x^{2n+1})$.
\fincorrection
\correction{004577}
$f(\sh y) = e^{y/2}$ d'où l'équation différentielle~:
$(1+x^2)f''(x) + xf'(x) = \frac14f(x)$.

En posant $f(x) = \sum_{n=0}^\infty a_n x^n$ on obtient
$4(k+1)(k+2)a_{k+2} = -(2k+1)(2k-1)a_k$ avec $a_0 = f(0) = 1$ et $a_1 = f'(0) = \frac12$,
d'où $a_{2p} = \frac{(-1)^{p+1}C_{4p-2}^{2p-1}}{p2^{4p}}$ si $p\ge 1$
et $a_{2p+1} = \frac{(-1)^pC_{4p}^{2p}}{2^{4p+1}(2p+1)}$ si $p\ge 0$.

Le rayon de convergence de la série correspondante est 1, ce qui
valide la méthode (avec le théorème d'unicité de Cauchy-Lipschitz).
\fincorrection
\correction{004578}
$a_n = a_{n-1} + a_{n-2}  \Rightarrow 
          a_n = \frac1{\sqrt5}\left(\left(\frac{1+\sqrt5}2\right)^n
              - \left(\frac{1-\sqrt5}2\right)^n\right)$.
\fincorrection
\correction{004579}
Coefficient de $x^n$ dans $(\sum x^k)(\sum (k+1)x^k)(\sum (k+1)^2x^k)
         = \frac{1+x}{(1-x)^6}$\par $ \Rightarrow  c_n = {\binom{n+5}{5}} + {\binom{n+4}{5}}
         = \frac{(2n+5)(n+4)(n+3)(n+2)(n+1)}{120}$.
\fincorrection
\nocorrection
\correction{004589}
\begin{enumerate}
  \item Pour $|x| < \frac1q$ : $\ln f(x) = \sum_{n=1}^\infty \ln(1-q^nx)
             = -\sum_{n=1}^\infty\sum_{k=1}^\infty \frac{q^{kn}x^k}k
             = -\sum_{k=1}^\infty \frac{q^kx^k}{k(1-q^k)}$,\par
             $f = e^{\ln f}$ est DSE par composition.
  \item $a_n = \frac{q^{n(n+1)/2}}{(q-1)\dots(q^n-1)}$, $R = \infty$.
\end{enumerate}
 \fincorrection
\correction{004590}
$|f^{(k)}(0)| = \sum_{n=0}^\infty n^{2k}e^{-n} \ge k^{2k}e^{-k}$,
donc $\frac{|f^{(k)}(0)|}{k!} \ge k^ke^{-k}  \Rightarrow  R=0$.
\fincorrection
\correction{004591}
Il y a dérivation terme à terme facilement et indéfiniment.

DSE au voisinage de~$0$~: on envisage de permuter les $\Sigma$ dans~:
$f_\alpha(x) = \sum_{n=1}^\infty\sum_{p=0}^\infty e^{-n^\alpha}\frac{(inx)^p}{p!}$,
ce qui est légitime si la série $\sum_{n=1}^\infty e^{-n^\alpha}e^{n|x|}$ converge.
On en déduit qu'une condition suffisante pour que~$f$ soit DSE au voisinage de~$0$
est $\alpha \ge 1$ (avec convergence si $x\in{]-1,1[}$ pour $\alpha = 1$ et pour
tout~$x\in\R$ si $\alpha > 1$).

Cas $\alpha < 1$~: $|f^{(k)}(0)| = \sum_{n=1}^\infty e^{-n^\alpha}n^k
\ge e^{-N^\alpha}N^k$ avec $N=\lfloor k^{1/\alpha}\rfloor$ donc pour $r>0$
fixé et $k$ tendant vers l'infini on~a
$\ln\Bigl(\Bigl|\frac{f^{(k)}(0)r^k}{k!}\Bigr|\Bigr) \sim
 \Bigl(\frac1\alpha-1\Bigr)k\ln(k)$
et la série de terme général $\frac{f^{(k)}(0)r^k}{k!}$ diverge grossièrement.

DSE au voisinage de~$a\ne 0$~: même raisonnement en écrivant
$f(x) = \sum_{n=1}^\infty\sum_{p=0}^\infty e^{-n^\alpha}e^{ina}\frac{(in(x-a))^p}{p!}$.
En conclusion, $f$ est analytique sur~$\R$ si et seulement si~$\alpha\ge 1$.
\fincorrection
\correction{004592}
\begin{enumerate}
  \item     Pour $x\ne 0$ la série comporte un nombre fini de termes non nuls
    au voisinage de~$x$, donc est $\mathcal{C}^\infty$ au voisinage de~$x$.
    On a $|f^{(k)}(x)| = \Bigl|\sum_{n=0}^\infty a_n\lambda_n^{k-n}\varphi_n^{(k)}(\lambda_nx)\Bigr|
                    \le\sum_{n=0}^\infty |a_n|\lambda_n^{k-n}M_n
                    \le \text{cste}(k) + \sum_{n=k+1}^\infty|a_n|M_n/\lambda_n$
    en supposant $\lambda_n\ge 1$ pour $n\ge k$, donc $f^{(k)}$ est bornée
    sur~$\R$. Ceci implique que $f$ est $\mathcal{C}^\infty$ en $0$ et on a le
    développemment limité~: $f(x) = \sum_{n=0}^k a_nx^n +  o(x^k)$ car
    $\phi \equiv 1$ au voisinage de~$0$ donc $f^{(k)}(0) = k!\,a_k$.
  \item $\psi(x) = \exp\Bigl(\frac1{(1-x)(x-2)}\Bigr)$ sur $]1,2[$, $\psi(x) = 0$ ailleurs.
\end{enumerate}
 \fincorrection
\correction{005747}
Dans chaque question, on note $f$ la fonction considérée.

\begin{enumerate}
 \item  $f$ est développable en série entière à l'origine en tant que fraction rationelle n'admettant pas $0$ pour pôle. Le rayon du développement est le minimum des modules des pôles de $f$ à savoir $1$. Pour $x$ dans $]-1,1[$,

\begin{center}
$f(x)=\frac{1}{x-1}-\frac{1}{x-2}=\frac{1}{1-x}-\frac{1}{2}\times\frac{1}{1-\frac{x}{2}}=\sum_{n=0}^{+\infty}x^n -\frac{1}{2}\sum_{n=0}^{+\infty}\frac{x^n}{2^n}=\sum_{n=0}^{+\infty}\left(1-\frac{1}{2^{n+1}}\right)x^n$.
\end{center}

\item  $f$ est développable en série entière à l'origine en tant que fraction rationelle n'admettant pas $0$ pour pôle.

\textbf{1er cas.} Si $|t| < 1$, soit $\theta=\Arccos t$. On a donc $\theta\in]0,\pi[$ et $t=\cos(\theta)$. Pour tout réel $x$, on a

\begin{center}
$x^2 -2tx +1 =x^2 -2x\cos(\theta)+1=(x-e^{i\theta})(x-e^{-i\theta})$,
\end{center}

avec $e^{i\theta}\neq e^{-i\theta}$. Les pôles sont de modules $1$ et le rayon du développement est donc égal à $1$. Pour $x$ dans $]-1,1[$,

\begin{align*}\ensuremath
 \frac{1}{x^2-2x\cos(\theta)+1}&=\frac{1}{2i\sin(\theta)}\left(\frac{1}{x-e^{i\theta}}-\frac{1}{x-e^{-i\theta}}\right) =\frac{1}{2i\sin(\theta)}\left(-\frac{e^{-i\theta}}{1-xe^{-i\theta}}+\frac{e^{i\theta}}{1-xe^{i\theta}}\right)\\
  &=\frac{1}{2i\sin(\theta)}\left(e^{i\theta}\sum_{n=0}^{+\infty}e^{in\theta}x^n-e^{-i\theta}\sum_{n=0}^{+\infty}e^{-in\theta}x^n\right) =\sum_{n=0}^{+\infty}\frac{\sin((n+1)\theta)}{\sin\theta}x^n.
\end{align*}

\begin{center}
\shadowbox{
$\forall t\in]-1,1[$, $\forall x\in]-1,1[$, $\frac{1}{x^2-2xt+1}=\sum_{n=0}^{+\infty}\frac{\sin((n+1)\theta)}{\sin\theta}x^n$ où $\theta=\Arccos t$.
}
\end{center}

\textbf{2ème cas.} Si $t > 1$, on peut poser $t =\ch(\theta)$ où $\theta$ est un certain réel positif ou nul. Plus précisément, $\theta=\Argch t=\ln(t+\sqrt{t^2-1})\in]0,+\infty[$. Pour tout réel $x$, on a

\begin{center}
$x^2 -2tx +1 =x^2 -2x\ch(\theta)+1=(x-e^{\theta})(x-e^{-\theta})$,
\end{center}

avec $e^\theta\neq e^{-\theta}$. Le minimum des modules des pôles de $f$ est $e^{-\theta}=\frac{1}{t+\sqrt{t^2-1}}= t-\sqrt{t^2-1}$. Le rayon du développement est donc $R=t-\sqrt{t^2-1}$. Pour $x\in]-R,R[$,

\begin{align*}\ensuremath
\frac{1}{x^2-2x\ch(\theta)+1}&=\frac{1}{2\sh(\theta)}\left(\frac{1}{x-e^{\theta}}-\frac{1}{x-e^{-\theta}}\right) =\frac{1}{2\sh(\theta)}\left(-\frac{e^{i\theta}}{1-xe^{-\theta}}+\frac{e^{\theta}}{1-xe^{\theta}}\right)\\
  &=\frac{1}{2\sh(\theta)}\left(e^{\theta}\sum_{n=0}^{+\infty}e^{n\theta}x^n-e^{-\theta}\sum_{n=0}^{+\infty}e^{-n\theta}x^n\right) =\sum_{n=0}^{+\infty}\frac{\sh((n+1)\theta)}{\sh\theta}x^n.
\end{align*}

\textbf{3ème cas.} Si $t < -1$, on applique ce qui précède à $-t$ et $-x$.

\textbf{4ème cas.} Si $t = 1$, pour $x\in]-1,1[$,

\begin{center}
$\frac{1}{x^2-2xt+1}=\frac{1}{(1-x)^2}=\left(\frac{1}{1-x}\right)'=\sum_{n=0}^{+\infty}(n+1)x^n$
\end{center}

Si $t = -1$, en remplaçant $x$ par $-x$,  on obtient pour $x\in]-1,1[$, $\frac{1}{(1+x)^2}=\sum_{n=0}^{+\infty}(-1)^n(n+1)x^n$.

\item  Pour tout réel $x$, $x^2 -5x+6=(x-2)(x-3)$ et donc si $x < 2$, $x^2-5x+6>0$. Pour $x\in ]-2,2[$,

\begin{center}
$\ln(x^2-5x+6)=\ln(2-x)+\ln(3-x)=\ln(6)+\ln\left(1-\frac{x}{2}\right)+\ln\left(1-\frac{x}{3}\right)$,
\end{center}

et puisque pour $x$ dans $]-2,2[$, $\frac{x}{2}$ et $\frac{x}{3}$ sont dans $]-1,1[$, 

\begin{center}
$\ln(x^2-5x+6)=\ln(6)-\sum_{n=1}^{+\infty}\left(\frac{1}{2^n}+\frac {1}{3^n}\right)\frac{x^n}{n}$,
\end{center}

et en particulier la fonction $f$ est développable en série entière et le rayon du développement est $2$ clairement.

\item  Si $\cos a = 0$, la fonction $f$ est définie et dérivable sur $\mathcal{D}=\Rr$ et si $\cos a\neq 0$, $f$ est définie et dérivable sur $\mathcal{D}=\left]-\infty,\frac{1}{\cos a}\right[\cup\left]\frac{1}{\cos a},+\infty\right[$. Pour $x\in\mathcal{D}$, 

\begin{center}
$f'(x) =\sin a\times\frac{1}{(1-x\cos a)^2}\times\frac{1}{1+\left(\frac{x\sin a}{1-x\cos a}\right)^2}=\frac{\sin a}{x^2-2x\cos a+1}$.
\end{center}

D'après 2), la fonction $f'$ est dans tous les cas développable en série entière, le rayon du développement est $1$ et pour $x$ dans $]-1,1[$

\begin{center}
$f'(x) =\sum_{n=0}^{+\infty}\frac{\sin((n+1)a)}{\sin a}x^n$.
\end{center}

On sait alors que la fonction $f$ est développable en série entière, que le développement a même rayon de convergence et s'obtient en intégrant terme à terme. Donc pour $x$ dans $]-1,1[$,

\begin{center}
$f(x)= f(0)+\int_{0}^{x}f'(t)\;dt=\sum_{n=0}^{+\infty}\frac{\sin(n+1)a)}{\sin a}x^{n+1}$.
\end{center}

\item  La fonction $f$ est développable en série entière en tant que fraction rationnelle n'admettant pas $0$ pour pôle. Le rayon est le minimum des modules des pôles de $f$ à savoir $1$.

\begin{center}
$\frac{1}{(x-1)(x-2)\ldots(x-p)}=\sum_{k=1}^{p}\frac{\lambda_k}{x-k}$
\end{center}

avec $\lambda_k=(-1)^{p-k}\frac{1}{(k-1)!(p-k)!}= (-1)^{p-k}\frac{k}{p!}C_p^k$.
Par suite, pour $x$ dans $]-1,1[$

\begin{align*}\ensuremath
f(x)&=\sum_{k=1}^{p}(-1)^{p-k}\frac{k}{p!}C_p^k\left(-\frac{1}{k}\right)\frac{1}{1-\frac{x}{k}}=\frac{(-1)^p}{p!}\sum_{k=1}^{p} (-1)^{k+1}C_p^k\left(\sum_{n=0}^{+\infty}\frac{x^n}{k^n}\right)\\
 &=\frac{(-1)^p}{p!}\sum_{n=0}^{+\infty}\left(\sum_{k=1}^{p}(-1)^{k+1}\frac{C_p^k}{k^n}\right)x^n.
\end{align*}

\item  La fonction $f$ est deux fois dérivable sur $]-1,1[$ et pour x dans $]-1,1[$, $f'(x) =2\frac{1}{\sqrt{1-x^2}}\Arcsin x$ puis

\begin{center}
$f''(x) =2x\frac{1}{(1-x^2)^{3/2}}\Arcsin x+\frac{2}{1-x^2}=\frac{x}{1-x^2}f'(x)+\frac{2}{1-x^2}$.
\end{center}

Donc, pour $x$ dans $]-1,1[$,

\begin{center}
$(1-x^2)f''(x)-xf'(x)=2$ \quad$(1)$\quad et $f(0)=f'(0) = 0$\quad$(2)$.
\end{center}

On admettra que ces égalités déterminent la fonction $f$ de manière unique.

Soit $\sum_{n=0}^{+\infty}a_nx^n$ une série entière de rayon $R$ supposé à priori strictement positif. Pour $x\in]-R,R[$, on pose $g(x)=\sum_{n=0}^{+\infty}a_nx^n$.

\begin{align*}\ensuremath
g\;\text{est solution de}\;(1)\;\text{sur}\;]-R,R[&\Leftrightarrow\forall x\in]-R,R[,\;(1-x^2)  \sum_{n=2}^{+\infty}n(n-1)a_nx^{n-2} -x\sum_{n=1}^{+\infty}na_nx^{n-1}=2\\
 &\Leftrightarrow \forall x\in]-R,R[,\;\sum_{n=2}^{+\infty}n(n-1)a_nx^{n-2}-\sum_{n=0}^{+\infty}n(n-1)a_nx^n+\sum_{n=0}^{+\infty}na_nx^n= 2\\
 &\Leftrightarrow \forall x\in]-R,R[,\;\sum_{n=2}^{+\infty}n(n-1)a_nx^{n-2}-\sum_{n=0}^{+\infty}n^2a_nx^n = 2\\
 &\Leftrightarrow \forall x\in]-R,R[,\;\sum_{n=0}^{+\infty}(n+2)(n+1)a_{n+2}x^{n}-\sum_{n=0}^{+\infty}n^2a_nx^n = 2\\
 &\Leftrightarrow\forall x\in]-R,R[,\;\sum_{n=0}^{+\infty}((n+2)(n+1)a_{n+2}-n^2a_n)x^{n}= 2\\
 &\Leftrightarrow a_2=1\;\text{et}\;\forall n\in\Nn^*,\;a_{n+2}=\frac{n^2}{(n+2)(n+1)}a_n\;(\text{par unicité des coefficients d'un DES}).
\end{align*}

En résumé, la fonction $g$ est solution de $(1)$ et $(2)$ sur $]-R,R[$ si et seulement si $a_0=a_1=0$ et $a_2=1$ et $\forall n\in\Nn^*,\;a_{n+2}=\frac{n^2}{(n+2)(n+1)}a_n$\quad(3) puis

\begin{align*}\ensuremath
(3)&\Leftrightarrow\forall n\in\Nn,\; a_{2n+1}= 0\;\text{et}\;a_0=0,\;a_2 =1\;\text{et}\;\forall n\geqslant2,\;a_{2n}= \frac{\left((2n-2)\times\ldots\times4\times2\right)^2}{(2n)\times(2n-1)\ldots\times4\times3}a_2\\
 &\Leftrightarrow a_0 = 0\;\text{et}\;\forall n\in\Nn,\;a_{2n+1}= 0\;\text{et}\;\forall n\in\Nn^*,\;a_{2n}=\frac{2^{2n-1}((n-1)!)^2}{(2n)!}
\end{align*}

En résumé, sous l'hypothèse $R > 0$, la fonction $g$ est solution de (1) et (2) sur $]-R,R[$ si et seulement si $\forall x\in]-R,R[$, $g(x)=\sum_{n=1}^{+\infty}\frac{2^{2n-1}}{n^2C_{2n}^n}x^{2n}$.

Réciproquement, calculons le rayon de la série entière précédente. Pour $x$ réel non nul,

\begin{center} 
$\left|\frac{2^{2n+1}(n!)^2x^{2n+2}}{(2n+2)!}\times\frac{(2n)!}{2^{2n-1}((n-1)!)^2x^{2n}}\right| = \frac{4x^2n^2}{(2n+2)(2n+1}\underset{n\rightarrow+\infty}{\rightarrow}x^2$.
\end{center}

D'après la règle de d'\textsc{Alembert}, la série proposée converge absolument pour $|x| < 1$ et diverge grossièrement pour $|x|> 1$. Le rayon de la série proposée est donc $1>0$ ce qui valide les calculs précédents.

Par unicité de la solution de (1) et (2) sur $]-1,1[$, $f$ est développable en série entière et

\begin{center}
\shadowbox{
$\forall x\in]-1,1[$, $\Arcsin^2x=\sum_{n=1}^{+\infty}\frac{2^{2n-1}}{n^2C_{2n}^n}x^{2n}$.
}
\end{center}

\item  Pour tout réel $x$, $\cos(x^2) =\sum_{n=0}^{+\infty}(-1)^n\frac{x^{4n}}{(2n)!}$ (le rayon est infini).
On sait alors que la fonction $f$ est développable en série entière, que le rayon du développement est encore infini et que l'on peut intégrer terme à terme pour obtenir (en tenant compte de $f(0)=0$) 

\begin{center}
\shadowbox{
$\forall x\in\Rr$, $\int_{0}^{x}=\sum_{n=0}^{+\infty}(-1)^n\frac{x^{4n+1}}{(4n+1)\times(2n)!}$.
}
\end{center}

\item  Les zéros du polynôme $t^4+t^2 + 1$ sont  $j$, $j^2$, $-j$ et $-j^2$. Donc la fonction $t\mapsto\frac{1}{t^4+t^2+1}$ est développable en série entière en tant que fraction rationnelle n'admettant pas zéro pour pôle et que le rayon de la série obtenue est $1$. Puis pour $t$ dans $]-1,1[$, 

\begin{center}
$\frac{1}{t^4+t^2+1}=\frac{1-t^2}{1-t^6}=(1-t^2)\sum_{n=0}^{+\infty}t^{6n}=\sum_{n=0}^{+\infty}t^{6n}-\sum_{n=0}^{+\infty}t^{6n+2}=1-t^2+t^6-t^8+t^{12}-t^{14}+\ldots$.
\end{center}

La fonction $t\mapsto\frac{1}{t^4+t^2+1}$ est continue sur $]-\infty,0]$ et négligeable devant $\frac{1}{t^2}$ quand $t$ tend vers $-\infty$. La fonction $t\mapsto\frac{1}{t^4+t^2+1}$ est donc intégrable sur $]-\infty,0]$.

Par intégration terme à terme licite, on obtient pour $x$ dans $]-1,1[$,

\begin{center}
$f(x)=\int_{-\infty}^{0}\frac{1}{t^4+t^2+1}\;dt+\int_{0}^{x}\frac{1}{t^4+t^2+1}\;dt=\int_{-\infty}^{0}\frac{1}{t^4+t^2+1}\;dt+\sum_{n=0}^{+\infty}\left(\frac{t^{6n+1}}{6n+1}-\frac{t^{6n+3}}{6n+3}\right)$.
\end{center}

Calcul de $I=\int_{-\infty}^{0}\frac{1}{t^4+t^2+1}\;dt$. Par parité et réalité,

\begin{center}
$\frac{1}{t^4+t^2+1}=\frac{a}{t-j}+\frac{\overline{a}}{t-j^2}-\frac{a}{t+j}-\frac{\overline{a}}{t+j^2}$,
\end{center}

avec $a=\frac{1}{4j^3+2j}=\frac{1}{2(2+j)}=\frac{2+j^2}{2(2+j)(2+j^2)}=\frac{1-j}{6}$.
Puis

\begin{align*}\ensuremath
\frac{1}{t^4+t^2+1}&=\frac{1}{6}\left(\frac{1-j}{t-j}+\frac{1-j^2}{t-j^2}-\frac{1-j}{t+j}-\frac{1-j^2}{t+j^2}\right)=\frac{1}{6}\left(\frac{3t+3}{t^2+t+1}+\frac{-3t+3}{t^2-t+1}\right)\\
 &=\frac{1}{4}\left(\frac{2t+1}{t^2+t+1}+\frac{1}{t^2+t+1}-\frac{2t-1}{t^2-t+1}+\frac{1}{t^2-t+1}\right)\\
 &=\frac{1}{4}\left(\frac{2t+1}{t^2+t+1}+\frac{1}{\left(t+\frac{1}{2}\right)^2+\left(\frac{\sqrt{3}}{2}\right)^2}-\frac{2t-1}{t^2-t+1}+\frac{1}{\left(t-\frac{1}{2}\right)^2+\left(\frac{\sqrt{3}}{2}\right)^2}\right).
\end{align*}

Par suite,

\begin{align*}\ensuremath
\int_{-\infty}^{0}\frac{1}{t^4+t^2+1}\;dt&=\frac{1}{4}\left[\ln\left(\frac{t^2+t+1}{t^2-t+1}\right)+\frac{2}{\sqrt{3}}\left(\Arctan\frac{2t+1}{\sqrt{3}}+\Arctan\frac{2t-1}{\sqrt{3}}\right)\right]_{-\infty}^0=\frac{1}{2\sqrt{3}}\left(\frac{\pi}{2}+\frac{\pi}{2}\right)=\frac{\pi}{2\sqrt{3}}.
\end{align*}

En résumé,

\begin{center}
\shadowbox{
$\forall x\in]-1,1[$, $\int_{-\infty}^{x}\frac{1}{t^4+t^2+1}\;dt=\frac{\pi}{2\sqrt{3}}+\sum_{n=0}^{+\infty}\left(\frac{t^{6n+1}}{6n+1}-\frac{t^{6n+3}}{6n+3}\right)$.
}
\end{center}

\item  $f$ est développable en série entière sur $\Rr$ en tant que produit de fonctions développables en série entière sur $\Rr$. Pour $x$ réel,

\begin{align*}\ensuremath
\cos x\ch x&=\frac{1}{4}\left(e^{(1+i)x}+e^{(1-i)x}+e^{(-1+i)x}+e^{(-1-i)x}\right)=\frac{1}{4}\sum_{n=0}^{+\infty}\left((1+i)^n+(1-i)^n+(-1+i)^n+(-1-i)^n\right)\frac{x^n}{n!}\\
 &=\frac{1}{4}\sum_{n=0}^{+\infty}\left((\sqrt{2}e^{i\pi/4})^n+(\sqrt{2}e^{-i\pi/4})^n+(\sqrt{2}e^{3i\pi/4})^n+(\sqrt{2}e^{-3i\pi/4})^n\right)\frac{x^n}{n!}\\
  &=\frac{1}{2}\sum_{n=0}^{+\infty}(\sqrt{2})^n\left(\cos\left(\frac{n\pi}{4}\right)+\cos\left(\frac{3n\pi}{4}\right)\right)\frac{x^n}{n!}
=\frac{1}{4}\sum_{n=0}^{+\infty}(\sqrt{2})^n\cos\left(\frac{n\pi}{2}\right)\cos\left(\frac{n\pi}{4}\right)\frac{x^n}{n!}\\
 &=\frac{1}{4}\sum_{p=0}^{+\infty}(\sqrt{2})^{2p}(-1)^p\cos\left(\frac{p\pi}{2}\right)\frac{x^{2p}}{(2p)!}=\frac{1}{4}\sum_{k=0}^{+\infty}(\sqrt{2})^{4k}(-1)^{2k}\cos\left(\frac{2k\pi}{2}\right)\frac{x^{4k}}{(4k)!}=\sum_{k=0}^{+\infty}(-1)^k2^{2k-2}\frac{x^{4k}}{(4k)!}.
\end{align*}

\begin{center}
\shadowbox{
$\forall x\in\Rr$, $\cos x\ch x=\sum_{k=0}^{+\infty}(-1)^k2^{2k-2}\frac{x^{4k}}{(4k)!}$.
}
\end{center}
\end{enumerate}
\fincorrection
\correction{005748}
Pour $x$ réel non nul, $f(x) =\sum_{n=0}^{+\infty}(-1)^n\frac{x^{2n}}{(2n+1)!}$ ce qui reste vrai pour $x=0$. La fonction $f$ est donc développable en série entière sur $\Rr$ et en particulier, la fonction $f$ est  de classe $C^\infty$ sur $\Rr$.
\fincorrection
\correction{005756}
Pour $x$ réel, on sait que $F(x)=e^{-x^2}\int_{0}^{x}e^{t^2}\;dt=\left(\sum_{n=0}^{+\infty}\sum_{n=0}^{+\infty}(-1)^n\frac{x^{2n}}{n!}\right)\left(\sum_{n=0}^{+\infty}\frac{x^{2n+1}}{n!(2n+1)}\right)$.

La fonction $F$ est impaire donc les coefficients d'indices pairs sont nuls. D'autre part, pour $n\in\Nn$, le coefficient de $x^{2n+1}$ du produit de \text{Cauchy} des deux séries précédentes vaut

\begin{center}
$\sum_{k=0}^{n}\frac{1}{k!(2k+1)}\times\frac{(-1)^{n-k}}{(n-k)!}$.
\end{center}

La méthode choisie fournit classiquement une expression compliquée des coefficients.

On peut aussi obtenir $F$ comme solution d'une équation différentielle linéaire du premier ordre. $F$ est dérivable sur $\Rr$ et pour tout réel $x$, $F'(x) =-2xe^{-x^2}\int_{0}^{x}e^{t^2}\;dt+1=-2xF(x)+1$.

$F$ est uniquement déterminée par les conditions $F'+ 2xF= 1$ et $F(0) = 0$\quad(*). $F$ est développable en série entière sur $\Rr$ d'après le début de l'exercice et impaire. Pour $x$ réel, posons donc $F(x)=\sum_{n=0}^{+\infty}a_nx^{2n+1}$.

\begin{align*}\ensuremath
(*)&\Leftrightarrow\forall x\in\Rr,\;\sum_{n=0}^{+\infty}(2n+1)a_nx^{2n}+ 2\sum_{n=0}^{+\infty}a_nx^{2n+2}=1\Leftrightarrow\forall x\in\Rr,\;a_0+\sum_{n=1}^{+\infty}((2n+1)a_n+ 2a_{n-1})x^{2n}= 1\\
 &\Leftrightarrow a_0 = 1\;\text{et}\;\forall n\geqslant1,\;(2n+1)a_n+ 2a_{n-1}=0\Leftrightarrow a_0 = 1\;\text{et}\;\forall n\geqslant1,\;a_n =-\frac{2}{2n+1}a_{n-1}\\
 &a_0 = 1\;\text{et}\;\forall n\geqslant1,\;a_n =\frac{(-1)^n2^n}{(2n+1)(2n-1)\ldots1}a_0\\
 &\Leftrightarrow\forall n\in\Nn,\;\frac{(-1)^n2^{2n}n!}{(2n+1)!}.
\end{align*}

On a montré que pour tout réel $x$, $F(x)=\sum_{n=0}^{+\infty}\frac{(-1)^n2^{2n}n!}{(2n+1)!}x^{2n+1}$. Par unicité des coefficients d'une série entière, $\forall n\in\Nn$, on obtient en particulier,

\begin{center}
$\forall n\in\Nn$, $\sum_{k=0}^{n}\frac{1}{k!(2k+1)}\times\frac{(-1)^{n-k}}{(n-k)!}=\frac{(-1)^n2^{2n}n!}{(2n+1)!}$.
\end{center}
\fincorrection
\correction{005761}
\begin{enumerate}
 \item  La fonction $f$ est de classe $C^\infty$ sur $\left]-\frac{\pi}{2},\frac{\pi}{2}\right[$ en tant que quotient de fonctions de classe $C^\infty$ sur $\left]-\frac{\pi}{2},\frac{\pi}{2}\right[$ dont le dénominateur ne s'annule pas sur $\left]-\frac{\pi}{2},\frac{\pi}{2}\right[$ et de plus $f'= 1 + f^2$.

Montrons par récurrence que pour tout naturel $n$, il existe un polynôme $P_n$ à coefficients entiers naturels tel que $f^{(n)}= P_n\circ f$ (ou encore $\forall x\in\left]-\frac{\pi}{2},\frac{\pi}{2}\right[$, $\tan^{(n)}(x) = P_n(\tan x)$).

\textbullet~C'est vrai pour $n=0$ avec $P_0=X$ et pour $n =1$ avec $P_1= 1+X^2$.

\textbullet~Soit $n\geqslant1$. Supposons que pour tout $k\in\llbracket0,n\rrbracket$, il existe un polynôme $P_k$ à coefficients entiers naturels tel que $f^{(k)}= P_k\circ f$. D'après la formule de \textsc{Leibniz},

\begin{center}
$f^{(n+1)}=(1+f^2)^{(n)}= (f^2)^{(n)}=\sum_{k=0}^{n}\dbinom{n}{k}f^{(k)}f^{(n-k)}=\left(\sum_{k=0}^{n}\dbinom{n}{k}P_kP_{n-k}\right)\circ f$
\end{center}

et le polynôme $P_{n+1}=\sum_{k=0}^{n}\dbinom{n}{k}P_kP_{n-k}$ est un polynôme à coefficients entiers naturels tel que $\tan^{(n+1)}=P_{n+1}\circ f$.

\textbf{Remarque.} On aurait pu aussi dériver l'égalité $f^{(n)}=P_n\circ f$ pour obtenir $f^{(n+1)}= f'\times P_n'\circ f = (P_1\times P_n')\circ f$ mais on a déjà dans l'idée une relation de récurrence sur les coefficients du développement de $\tan$ qui n'est pas fournie par cette dernière égalité.

\item  Soient $x\in\left[0,\frac{\pi}{2}\right[$ et $n\in\Nn$. La formule de \textsc{Taylor}-\textsc{Laplace} à l'ordre $n$ en $0$ fournit

\begin{center}
$f(x)=\sum_{k=0}^{n}\frac{f^{(k)}(0)}{k!}x^k+\int_{0}^{x}\frac{(x-t)^n}{n!}f^{(n+1)}(t)\;dt $.
\end{center}

Le 1) montre que pour tout réel $t$ de $\left[0,\frac{\pi}{2}\right[$ et tout entier naturel $k$, $f^{(k)}(t)=P_{k}(\tan t)\geqslant 0$.

Donc, d'une part $\frac{f^{(k)}(0)}{k!}x^k\geqslant0$ et d'autre part,

\begin{center}
$\sum_{k=0}^{n}\frac{f^{(k)}(0)}{k!}x^k=f(x)-\int_{0}^{x}\frac{(x-t)^n}{n!}f^{(n+1)}(t)\;dt$.
\end{center}

La suite des sommes partielles de la série de terme général $\frac{f^{(k)}(0)}{k!}x^k\geqslant0$ est majorée et donc la série de terme général $\frac{f^{(k)}(0)}{k!}x^k$ converge.

Ainsi, la série de \textsc{Taylor} de $f$ à l'origine converge pour tout réel $x$ de $\left[0,\frac{\pi}{2}\right[$. Son rayon de convergence $R$ est donc supérieur ou égal à $\frac{\pi}{2}$ (et donc la série de terme général $\frac{f^{(k)}(0)}{k!}x^k$ converge aussi pour $x\in\left]-\frac{\pi}{2},0\right]$). Il n'y a par contre aucune raison pour le moment pour que sa somme soit $f$.

\item  Pour $n$ entier naturel donné, posons $a_n =\frac{f^{(n)}(0)}{n!}$ puis pour $x$ dans $\left]-\frac{\pi}{2},\frac{\pi}{2}\right[$, posons $g(x) =\sum_{n=0}^{+\infty}a_nx_n$.

On a vu que $\forall n\in\Nn^*$, $P_{n+1}=\sum_{k=0}^{n}P_kP_{n-k}$. On divise les deux membres de ces égalités par $n!$ et on prend la valeur en $0$ ($=\tan0$). On obtient

\begin{center}
$\forall n\in\Nn^*$, $(n+1)a_{n+1}=a_ka_{n-k}$ et aussi $a_0 = 0$ et $a_1 = 1$.
\end{center}

Donc, pour $x\in\left]-\frac{\pi}{2},\frac{\pi}{2}\right[$,

\begin{align*}\ensuremath
g'(x)&=\sum_{n=0}^{+\infty}(n+1)a_{n+1}x^n=1+\sum_{n=1}^{+\infty}\left(\sum_{k=0}^{n}a_ka_{n-k}\right)x^n =1+\sum_{n=0}^{+\infty}\left(\sum_{k=0}^{n}a_ka_{n-k}\right)x^n=1 +\left(\sum_{n=0}^{+\infty}a_nx^n\right)^2\\
 &= 1 + g^2(x).
\end{align*}

De plus, $g(0) = a_0 = 0$.

Pour $x\in\left]-\frac{\pi}{2},\frac{\pi}{2}\right[$, posons alors $h(x) =\Arctan(g(x))$. La fonction $h$ est dérivable sur $\left]-\frac{\pi}{2},\frac{\pi}{2}\right[$ et pour $x\in\left]-\frac{\pi}{2},\frac{\pi}{2}\right[$

\begin{center}
$h'(x)=\frac{g'(x)}{1+g^2(x)}= 1$ puis $h(x) =h(0)+(x-0) = x$.
\end{center}

Ainsi, pour tout $x\in\left]-\frac{\pi}{2},\frac{\pi}{2}\right[$, $g(x) =\tan x = f(x)$. Ceci montre déjà que $f$ est développable en série entière sur $\left]-\frac{\pi}{2},\frac{\pi}{2}\right[$. Mais quand $x$ tend vers $\frac{\pi}{2}$  par valeurs inférieures, $g(x) = f(x)$ tend vers $+\infty$ et donc $R\leqslant\frac{\pi}{2}$ puis $R=\frac{\pi}{2}$.

En résumé, la fonction tangente est développable en série entière sur $\left]-\frac{\pi}{2},\frac{\pi}{2}\right[$ et pour $x\in\left]-\frac{\pi}{2},\frac{\pi}{2}\right[$, $\tan x =\sum_{n=0}^{+\infty}a_nx^n$ où
$a_0 = 0$, $a_1 = 1$ et $\forall n\in\Nn^*$, $(n+1)a_{n+1}=\sum_{k=0}^{n}a_ka_{n-k}$. De plus, $\forall n\in\Nn$, $a_{2n}= 0$ puisque la fonction tangente est impaire.

\item  $a_0 = a_2 = a_4 = a_6 = 0$ puis $a_1 = 1$.

$3a_3 = a_0a_2 + a_1^2 + a_2a_0 = 1$ et  donc $a_3 =\frac{1}{3}$.

$5a_5 =2a_1a_3=\frac{2}{3}$ et donc $a_5 =\frac{2}{15}$.

$7a_7 =2a_1a_5 + a_3^2=\frac{4}{15}+\frac{1}{9}=\frac{51}{135}=\frac{17}{45}$ et $a_7=\frac{17}{315}$.

\begin{center}
\shadowbox{
$\forall x\in\left]-\frac{\pi}{2},\frac{\pi}{2}\right[$, $\tan x = x +\frac{x^3}{2}+\frac{2x^5}{15}+\frac{17x^7}{315}+ ...$
}
\end{center}

\item  Pour tout  réel $x$, $\tanh(x) =\frac{1}{x}\tan(ix)$ et donc pour $x\in\left]-\frac{\pi}{2},\frac{\pi}{2}\right[$,

\begin{center}
$\tanh(x) =\frac{1}{i}\sum_{n=0}^{+\infty}a_{2n+1}(ix)^{2n+1}=\sum_{n=0}^{+\infty}(-1)^na_{2n+1}x^{2n+1}$.
\end{center}

Cette série entière a aussi pour rayon de convergence $\frac{\pi}{2}$.
\end{enumerate}
\fincorrection
\correction{005762}
Soit $x\in\Rr$. La fonction $t\mapsto e^{-t^2}\sin(tx)$ est continue sur $[0,+\infty[$, négligeable devant $\frac{1}{t^2}$ quand $t$ tend vers $+\infty$ et est donc intégrable sur $[0,+\infty[$. La fonction $F$ est donc définie sur $\Rr$ et impaire.

Soit $x\in\Rr$. Pour tout réel $t$, posons $f(t)=e^{-t^2}\sin(tx)$. Pour $t\in\Rr$, on a

\begin{center}
$e^{-t^2}\sin(tx) =\sum_{n=0}^{+\infty}(-1)^n\frac{x^{2n+1}}{(2n+1)!}t^{2n+1}e^{-t^2}$.
\end{center}

Pour $n\in\Nn$ et $t\in\Rr$, posons $f_n(t)=(-1)^n\frac{x^{2n+1}}{(2n+1)!}t^{2n+1}e^{-t^2}$.

\textbullet~Chaque fonction $f_n$, $n\in\Nn$, est continue puis intégrable sur $[0,+\infty[$ car négligeable devant $\frac{1}{t^2}$ quand $t$ tend vers $+\infty$.

\textbullet~La série de fonctions de terme général $f_n$, $n\in\Nn$, converge simplement vers la fonction $f$ sur $[0,+\infty[$.

\textbullet~Ensuite, $\sum_{n=0}^{+\infty}\int_{0}^{+\infty}|f_n(t)|\;dt=\sum_{n=0}^{+\infty}\frac{|x|^{2n+1}}{(2n+1)!}\int_{0}^{+\infty}t^{2n+1}e^{-t^2}\;dt$. Pour $n\in\Nn$, posons $I_n=\int_{0}^{+\infty}t^{2n+1}e^{-t^2}\;dt$.

Soit $n\in\Nn^*$. Soit $A$ un réel strictement positif. Les deux fonctions $t\mapsto t^{2n}$ et $t\mapsto-\frac{1}{2}e^{-t^2}$ sont de classe $C^1$ sur le segment $[0,A]$. On peut donc effectuer une intégration par parties et on obtient

\begin{align*}\ensuremath
\int_{0}^{A}t^{2n+1}e^{-t^2}\;dt&=\int_{0}^{A}t^{2n}\times te^{-t^2}\;dt=\left[-\frac{1}{2}t^{2n}e^{-t^2}\right]_0^A+n\int_{0}^{A}t^{2n-1}e^{-t^2}\;dt\\
 &=-\frac{1}{2}A^{2n}e^{-A^2}+n\int_{0}^{A}t^{2n-1}e^{-t^2}\;dt.
\end{align*}

Quand $A$ tend vers $+\infty$, on obtient $I_{n}=nI_{n-1}$. En tenant compte, de $I_0=\int_{0}^{+\infty}te^{-t^2}\;dt=\frac{1}{2}$ on a donc $\forall n\in\Nn$, $I_n=\frac{n!}{2}$ puis

\begin{center}
$\sum_{n=0}^{+\infty}\int_{0}^{+\infty}|f_n(t)|\;dt=\sum_{n=0}^{+\infty}\frac{n!|x|^{2n+1}}{(2n+1)!}$.
\end{center}

Soient $n\in\Nn$ et $x\in\Rr$. $\left|\frac{\frac{(n+1)!|x|^{2n+3}}{(2n+3)!}}{\frac{n!|x|^{2n+1}}{(2n+1)!}}\right|=\frac{(n+1)x^2}{(2n+3)(2n+2)}$ et donc $\lim_{n \rightarrow +\infty}\frac{\frac{(n+1)!|x|^{2n+3}}{(2n+3)!}}{\frac{n!|x|^{2n+1}}{(2n+1)!}}=0$. D'après la règle de d'\textsc{Alembert}, la série numérique de terme général $\frac{n!|x|^{2n+1}}{(2n+1)!}$ converge.

En résumé, pour tout réel $x$,

\textbullet~Chaque fonction $f_n$, $n\in\Nn$, est continue puis intégrable sur $[0,+\infty[$ car négligeable devant $\frac{1}{t^2}$ quand $t$ tend vers $+\infty$.

\textbullet~La série de fonctions de terme général $f_n$, $n\in\Nn$, converge simplement vers la fonction $f$ sur $[0,+\infty[$.

\textbullet~$\sum_{n=0}^{+\infty}\int_{0}^{+\infty}|f_n(t)|\;dt<+\infty$.

D'après un théorème d'intégration terme à terme, pour tout réel $x$,

\begin{center}
$\int_{0}^{+\infty}e^{-t^2}\sin(tx)\;dt=\sum_{n=0}^{+\infty}\int_{0}^{+\infty}f_n(t)\;dt=\sum_{n=0}^{+\infty}(-1)^n\frac{n!x^{2n+1}}{2(2n+1)!}$.
\end{center}

\begin{center}
\shadowbox{
$\forall x\in\Rr$, $\int_{0}^{+\infty}e^{-t^2}\sin(tx)\;dt=\sum_{n=0}^{+\infty}(-1)^n\frac{n!x^{2n+1}}{2(2n+1)!}$.
}
\end{center}

$F$ est dérivable sur $\Rr$ et pour tout réel $x$,

\begin{center}
$F'(x)=\sum_{n=0}^{+\infty}(-1)^n\frac{n!x^{2n}}{2(2n)!}=\frac{1}{2}-\frac{x}{2}\sum_{n=1}^{+\infty}(-1)^{n-1}\frac{(n-1)!x^{2n-1}}{(2(2n-1)!}=\frac{1}{2}-\frac{x}{2}F(x)$.
\end{center}

Par suite, pour tout réel $x$, $e^{x^2/4}F'(x) +\frac{x}{2}e^{x^2/4}F(x)=\frac{e^{x^2/4}}{2}$ et donc

\begin{center}
$F(x)=F(0)+\frac{e^{-x^2/4}}{2}\int_{0}^{x}e^{t^2/4}\;dt =\frac{e^{-x^2/4}}{2}\int_{0}^{x}e^{t^2/4}\;dt$.
\end{center}

\begin{center}
\shadowbox{
$\forall x\in\Rr$, $\int_{0}^{+\infty}e^{-t^2}\sin(tx)\;dt=\frac{e^{-x^2/4}}{2}\int_{0}^{x}e^{t^2/4}\;dt$.
}
\end{center}

\fincorrection
\nocorrection
\nocorrection
\nocorrection
\nocorrection
\nocorrection
\correction{004593}
\begin{enumerate}
  \item $R=1$.
  \item $x=-1  \Rightarrow $ cv (série alternée), $x=1  \Rightarrow $ dv.
  \item $f$ est croissante sur $[0,1[$ donc $L$ existe dans $[0,+\infty]$.

             $L = \sup\limits_{[0,1[} f(x)
              \ge \sup\limits_{[0,1[} \sum_{n=1}^N x^n\sin\frac1{\sqrt n}
              \ge \sum_{n=1}^N \sin\frac1{\sqrt n}  \Rightarrow  L = +\infty$.
\end{enumerate}
 \fincorrection
\correction{004594}
\begin{enumerate}
  \item Fonction croissante. $\lim_{x\to1^-} A(x) \ge \sum_{n=0}^N a_n$.
  \item Dém de type Césaro.
\end{enumerate}
 \fincorrection
\correction{004595}
Continuité radiale.
\fincorrection
\correction{004596}
$\frac{c_n}{b_n} = a_0 + a_1\frac{b_{n-1}}{b_n} + \dots + a_n\frac{b_0}{b_n}
                  = \sum_{k=0}^\infty a_ku_{n,k}$
et le théorème de convergence dominée s'applique.
\fincorrection
\correction{004597}
$y = \sum_{n=1}^\infty \frac{5^{n-2}x^n}{\prod_{2\le k\le n}(3k+2)}$
         $(R=\infty)$.\par
         $N = 8  \Rightarrow  0.409954 \le y(1) \le 0.409973$.
\fincorrection
\correction{004598}
\begin{enumerate}
  \item $\tan^{(n)} =
             \sum_{k=0}^{n-1} C_{n-1}^k \tan^{(k)}\tan^{(n-1-k)}$.
  \item Pour $0 \le x < \pi/2$ la série est à termes positifs et les
             sommes partielles sont majorées par $\tan x$.
             Pour $-\pi/2 < x \le 0$, il y a convergence absolue.
  \item 
  \item Si $R > \pi/2$, $f$ aurait une limite finie en $\pi/2$.
\end{enumerate}
 \fincorrection
\correction{004599}
\begin{enumerate}
  \item Produit de deux séries $ \Rightarrow  R \ge 1$.
          Lorsque $x\to1^-$   $f(x) \to +\infty  \Rightarrow  R = 1$.
  \item $(1-x^2)y' = xy+1  \Rightarrow  (n+2)a_{n+2} = (n+1)a_n  \Rightarrow 
              a_{2k} = a_0\frac{C_{2k}^k}{4^k}$,
             $a_{2k+1} = a_1\frac{4^k}{(2k+1)C_{2k}^k}$.\par
             $a_0 = 0$, $a_1 = 1  \Rightarrow 
             y = \sum_{k=0}^\infty \frac{4^kx^{2k+1}}{(2k+1)C_{2k}^k}$.
  \item $\Arcsin^2x = 2 \int_{t=0}^x f(t)\,d t
              = \sum_{k=1}^\infty \frac{2^{2k-1}x^{2k}}{k^2C_{2k}^k}$.
\end{enumerate}
 \fincorrection
\correction{004600}
\begin{enumerate}
  \item $R=4$.
  \item $y = 4\sqrt{\frac x{(4-x)^3}}\left(\sqrt{\frac{4-x}x}
                  - \Arctan\sqrt{\frac{4-x}x} + c\right)$.
             $f(x) = 1 + \frac x2 + \text{o}(x)  \Rightarrow  c = \frac\pi2$.
             $ \Rightarrow \sum_{n=0}^\infty \frac 1{\strut C_{2n}^n}
              = \frac 43 + \frac{2\pi}{9\sqrt 3}$.
\end{enumerate}
 \fincorrection
\correction{004601}
\begin{enumerate}
  \item $R=\sqrt2$.
  \item Stirling $ \Rightarrow  a_n{\sqrt2\,}^{2n+1} \sim \frac2{\sqrt{\pi n}}
              \Rightarrow $ DV.
  \item $(x^2-2)y' + xy + 2 = 0  \Rightarrow 
             f(x) = \frac{2\Arcsin (x/\sqrt2\,)}{\sqrt{2-x^2}}$.
\end{enumerate}
 \fincorrection
\correction{004602}
\begin{enumerate}
  \item 
  \item $f'(x) = e^xf(x)$.
\end{enumerate}
 \fincorrection
\correction{004603}
\begin{enumerate}
  \item $a_n \le n!$ par récurrence.
  \item $2f' = f^2  \Rightarrow  f(x) = \frac1{1-x/2}$.
  \item $a_n = n!\,2^{-n}$.
\end{enumerate}
 \fincorrection
\correction{004604}
$Z(x) = \sum_{n,p\ge 1} \frac{x^n}{p^{2n}} = \sum_{p\ge 1}\frac{x}{p^2-x}$.

$Z'(x) = \sum_{p\ge 1}\frac{p^2}{(p^2-x)^2} = \sum_{p\ge 1}\frac{1}{p^2-x} + \sum_{p\ge 1}\frac{x}{(p^2-x)^2}$.

$Z^2(x) = \sum_{p,q\ge 1}\frac{x^2}{(p^2-x)(q^2-x)}
  = \sum_{p\ne q}\frac{x^2}{q^2-p^2}\Bigl(\frac{1}{p^2-x} - \frac{1}{q^2-x}\Bigr) + \sum_{p\ge1}\frac{x^2}{(p^2-x)^2}$

$Z^2(x)-xZ'(x) + Z(x) = 2\sum_{p\ne q}\frac{x^2}{(q^2-p^2)(p^2-x)}$.

A $p$ fixé, $\sum_{q\ne p}\frac1{q^2-p^2} = \frac{1}{2p}\sum_{q\ne p}\Bigl(\frac1{q-p}-\frac1{q+p}\Bigr) = \frac1{2p}\Bigl(\frac1p + \frac1{2p}\Bigr) = \frac{3}{4p^2}$.

Donc $Z^2(x)-xZ'(x) + Z(x) = \frac32\sum_{p\ge 1}\frac{x^2}{p^2(p^2-x)} = \frac32(Z(x)-x\zeta(2))$.

Rmq~: $2Z(x^2) = 1-\pi x\mathrm{cotan}(\pi x)$ (Euler).
\fincorrection
\correction{004605}
$\alpha=\pi/4$, $\beta=\pi/2$.
\fincorrection
\nocorrection
\correction{004607}
\begin{enumerate}
  \item $t^t = \exp(t\ln t) = \sum_{k=0}^\infty \frac{t^k\ln^k t}{k!}$.
  \item 0.78343
\end{enumerate}
 \fincorrection
\correction{004608}
$=\sum_{n=1}^\infty  \int_{t=0}^1 -\frac{t^n\ln t}n\,d t
          = \sum_{n=1}^\infty \frac1{n(n+1)^2} = 2-\frac{\pi^2}6$.
\fincorrection
\correction{004609}
Développer en série entière $\ln(1-t^2)$.
$I = \frac{\pi^2}2 - 4\ln 2$.
\fincorrection
\nocorrection
\nocorrection
\nocorrection
\correction{005751}
On a déjà vu que $W_n\underset{n\rightarrow+\infty}{\sim}\sqrt{\frac{\pi}{2n}}$ et la règle de d'\textsc{Alembert} fournit $R = 1$. Soit $x\in]-1,1[$.

Pour tout $t\in\left[0,\frac{\pi}{2}\right]$ et tout entier naturel $n$, $\left|x^n\cos^nt\right|\leqslant|x|^n$. Comme la série numérique de terme général $|x|^n$, $n\in\Nn$, converge, la série de fonctions de terme général $t\mapsto x^ncos^nt$ est normalement et donc uniformément convergente sur le segment $\left[0,\frac{\pi}{2}\right]$. D'après le théorème d'intégration terme à terme sur un segment,

\begin{align*}\ensuremath
\sum_{n=0}^{+\infty}W_nx^n&=\sum_{n=0}^{+\infty}x^n\int_{0}^{\pi/2}\cos^nt\;dt  =\int_{0}^{\pi/2}\left(\sum_{n=0}^{+\infty}x^n\cos^nt\right)dt =\int_{0}^{\pi/2}\frac{1}{1-x\cos t}\;dt\\
 &=\int_{0}^{1}\frac{1}{1-x\frac{1-u^2}{1+u^2}}\frac{2du}{1+u^2}\;(\text{en posant}\;u=\tan\left(\frac{t}{2}\right))\\
 &=2\int_{0}^{1}\frac{1}{(1+x)u^2+(1-x)}\;du=2\times\frac{1}{1+x}\times\frac{1}{\sqrt{\frac{1-x}{1+x}}}\left[\Arctan\left(\frac{u}{\sqrt{\frac{1-x}{1+x}}}\right)\right]_0^1\\
 &= \frac{2}{\sqrt{1-x^2}}\Arctan\sqrt{\frac{x+1}{x-1}}.
\end{align*}

\begin{center}
\shadowbox{
$\forall x\in]-1,1[$, $\sum_{n=0}^{+\infty}W_nx^n=\frac{2}{\sqrt{1-x^2}}\Arctan\sqrt{\frac{x+1}{x-1}}$.
}
\end{center}
\fincorrection
\nocorrection
\correction{004614}
\begin{enumerate}
  \item Calcul.
  \item     Soit $0<r_0<d$ et $R(\theta)$ le rayon de la série de Taylor de~$f$ en $r_0e^{i\theta}$.
    Le cercle de centre $0$ et de rayon $r_0$ est recouvert par les disques ouverts
    $D(r_0e^{i\theta}, \frac12R(\theta))$, $\theta$ variant de $0$ à~$2\pi$,
    donc on peut en extraire un recouvrement fini~; soit $\rho$ le rayon minimum
    des disques extraits. Alors pour $0\le\theta\le2\pi$ on a $R(\theta) \ge \rho$
    (cf. analycité de la somme d'une série entière dans le disque ouvert de convergence).
    
    D'après la première question on a~:
    $\Bigl|\frac{f^{(n)}(r_0e^{i\theta})}{n!}\Bigr| \le \frac M{\rho^n}$
    où $M$ majore $|f|$ sur $\overline D(0,r_0+\rho)$ d'où pour $|r-r_0|<\rho$~:
    $$ \int_{\theta=0}^{2\pi}\frac{f(re^{i\theta})}{e^{in\theta}}\, d\theta
      =  \int_{\theta=0}^{2\pi} \sum_{k=0}^\infty
           \frac{f^{(k)}(r_0e^{i\theta})}{k!} (r-r_0)^k e^{i(k-n)\theta}\, d\theta
      = \sum_{k=0}^\infty (r-r_0)^k  \int_{\theta=0}^{2\pi}
           \frac{f^{(k)}(r_0e^{i\theta})}{k!} e^{i(k-n)\theta}\, d\theta.$$
    ce qui démontre l'analycité de $\varphi = r \mapsto  \int_{\theta=0}^{2\pi}\frac{f(re^{i\theta})}{e^{in\theta}}\, d\theta$
    sur $]0,d[$.
    
    Enfin, $\varphi(r) = a_nr^n$ au voisinage de~$0$ d'où $\varphi(r) = a_nr^n$
    sur~$[0,d[$ par prolongement analytique.
  \item $g(z)= \frac1{2\pi} \int_{\theta=0}^{2\pi}\frac{f(re^{i\theta})}{re^{i\theta}-z}re^{i\theta}\, d\theta
                  = \frac1{2\pi} \int_{\theta=0}^{2\pi}f(re^{i\theta})\sum_{k=0}^\infty \frac{z^k}{r^ke^{ik\theta}}\, d\theta
                  = \sum_{k=0}^\infty \frac{z^k}{2\pi} \int_{\theta=0}^{2\pi}\frac{f(re^{i\theta})}{r^ke^{ik\theta}}\, d\theta
                  = \sum_{k=0}^\infty a_kz^k$.

             Le rayon est au moins égal à~$r$ car $f$ est bornée sur $\overline D(0,r)$.
  \item résulte de la question {3.}.
  \item D'après la question {1.}, $|a_n| \le \|f\|_\infty/r^n$ pour tout $r>0$ donc $a_n = 0$ si $n\ge 1$.
  \item $1/P$ est analytique bornée sur~$\C$.
  \item On peut passer à la limite uniforme (ou dominée) dans la question {3.}.
  \item $f(z) = \sum_{n=0}^\infty a_nz^n  \Rightarrow  f\circ g(z) = \sum_{n=0}^\infty a_ng^n(z)$
      et il y a convergence localement uniforme.
\end{enumerate}
 \fincorrection
\nocorrection
\correction{004616}
$2 \Rightarrow 1$ : évident.\par
         {$1 \Rightarrow 2$ : }Soit $a > 0$ et $M = \sup(|f(z)|e^{-a|z|})$.

         $a_n = \frac1{2\pi} \int_{\theta=0}^{2\pi}
                \frac{f(Re^{i\theta})}{R^ne^{in\theta}}\,d\theta
           \Rightarrow  |a_n| \le M\frac{e^{aR}}{R^n} \le M\inf\limits_{R>0}
                \frac{e^{aR}}{R^n}
          = M\left(\frac{ea}n\right)^n$.

          Donc $\sqrt[n]{n!\,\|a_n\|} \le \sqrt[n]{n!}\frac{ea}n
               \to a$ lorsque $n\to\infty$. CQFD
\fincorrection
\correction{004617}
\begin{enumerate}
  \item $f(\overline z) = \overline{f(z)}$.
  \item $\Im(f)$ est de signe constant sur le connexe $D\cap \Omega ^+$
    et $f(z)\sim z$ au voisinage de~$0$.
  \item Intégrer terme à terme.
  \item $\Im (f(re^{i\theta})) \sin\theta \ge 0$ par la question {2.}.
\end{enumerate}
\fincorrection
\nocorrection
\correction{004619}
On pose, sous réserve de convergence,
$f(t) = \sum_{n=1}^\infty z_nt^n$. Alors~:
$$f(t)
=\sum_{n=1}^p z_nt^n + \sum_{j=1}^pe^{i\omega}p_j\sum_{n=p+1}^\infty z_{n-j}t^n
=\sum_{n=1}^p z_nt^n + \sum_{j=1}^pe^{i\omega}p_jt^j\Bigl(f(t)-\sum_{n=1}^{p-j} z_nt^n\Bigr)
$$
soit~:
$$\Bigl(1-\sum_{j=1}^pe^{i\omega}p_jt^j\Bigr)f(t) = P(t)f(t)
= \sum_{n=1}^p z_nt^n - \sum_{j=1}^pe^{i\omega}p_jt^j\sum_{n=1}^{p-j} z_nt^n
=Q(t),$$
donc $f(t) = Q(t)/P(t)$. Réciproquement, soit $Q(t)/P(t) = \sum_{n=1}^\infty a_nt^n$~:
en remontant les calculs précédents on voit que $(a_n)$ vérifie la même
relation de récurrence que~$(z_n)$ avec les mêmes premiers termes d'où
$z_n = a_n$ pour tout~$n$.

Si $|t|< 1$ alors $\Bigl|\sum_{j=1}^pe^{i\omega}p_jt^j\Bigr|< 1$
donc $P$ n'a pas de racine dans le disque unité ouvert.
Si $P$ n'a pas non plus de racine sur le cercle unité alors le développement
en série entière de $Q(t)/P(t)$ a un rayon $>1$ et $z_n\to0$ lorsque $n\to\infty$.
Si $P$ admet des racines dans~$\mathbb{U}$ on peut déja dire que la suite $(z_n)$
est bornée par $\max(|z_1|,\dots,|z_p|)$ puis\dots~?
\fincorrection
\correction{004620}
\begin{enumerate}
  \item 
  \item Complétude~: soit $(f_k)$ une suite d'éléments de~$E$ de Cauchy,
    $f_k(z) = \sum_{n\in\N}a_{n,k}z^n$. On a, à $k$ et $n$ fixés, par convergence dominée~:
    $$\frac1{2\pi} \int_{\theta=0}^{2\pi}f_k(e^{i\theta})e^{-in\theta}\, d\theta
      = \lim_{r\to 1^-}\frac1{2\pi r^n} \int_{\theta=0}^{2\pi}f_k(re^{i\theta})e^{-in\theta}\, d\theta
      = a_{n,k}.$$
    La suite $(f_k)$
    converge uniformément sur~$\overline D$ vers une fonction $\varphi : {\overline D} \to \C$
    continue. On note~:
    $$a_n = \frac1{2\pi} \int_{\theta=0}^{2\pi}\varphi(e^{i\theta})e^{-in\theta}\, d\theta
          = \lim_{k\to\infty} a_{n,k}.$$
    La suite $(a_n)$ est bornée, donc le rayon de convergence de $\sum_{n\in\N}a_nz^n$
    est supérieur ou égal à~$1$. Pour $z\in D$ fixé on a alors lorsque $k\to\infty$ :
    \begin{align*} f_k(z)
        &= \sum_{n\in\N}a_{n,k}z^n
         = \frac1{2\pi} \int_{\theta=0}^{2\pi}\Bigl(\sum_{n\in\N}f_k(e^{i\theta})e^{-in\theta}z^n\Bigr)\, d\theta
         = \frac1{2\pi} \int_{\theta=0}^{2\pi}\frac{f_k(e^{i\theta})}{1-ze^{-i\theta}}\, d\theta\cr
        &\to \frac1{2\pi} \int_{\theta=0}^{2\pi}\frac{\varphi(e^{i\theta})}{1-ze^{-i\theta}}\, d\theta
         = \frac1{2\pi} \int_{\theta=0}^{2\pi}\Bigl(\sum_{n\in\N}\varphi(e^{i\theta})e^{-in\theta}z^n\Bigr)\, d\theta
         = \sum_{n\in\N}a_nz^n\cr
    \end{align*}
    ce qui prouve que $\varphi\in E$. Enfin on a $\|f_k-\varphi\| \to 0$ lorsque $k\to\infty$
    par convergence uniforme, d'où $\varphi = \lim_{k\to\infty}f_k$ dans~$E$.
  \item Soit $f\in E$ et $f_n(z) = f\Bigl(\frac{nz}{n+1}\Bigr)$.
    Comme $f$ est uniformément continue, $f_n$ converge uniformément vers~$f$
    sur $\overline D$. Soit $\varepsilon>0$ et $n$ tel que $\|f-f_n\|_\infty \le \varepsilon$.
    Comme $f_n$ est développable en série entière avec un rayon au moins égal à $1+\frac1n$,
    son développement converge uniformément vers~$f_n$ sur $\overline D$ donc
    il existe $P\in\C[X]$ tel que $\|f_n-P\|_\infty \le \varepsilon$.
\end{enumerate}
 \fincorrection
\correction{004621}
\begin{enumerate}
  \item $\frac z{(1-z)^2} = \sum_{n=1}^\infty nz^n$ ($R=1$).

    Pour $z,t\in \mathring D(0,1)$ on a
    $\frac{z}{(1-z)^2} - \frac{t}{(1-t)^2} = \frac{(z-t)(1-zt)}{(1-z)^2(1-t)^2}$,
    quantité nulle si et seulement si $z=t$, d'où l'injectivité de $z \mapsto\frac z{(1-z)^2}$.
  \item
  \begin{enumerate}
    \item $f(z)\in\R \Leftrightarrow f(z) = \overline{f(z)} = f(\overline z) \Leftrightarrow z=\overline z \Leftrightarrow z\in\R$.

Par injectivité, on en déduit que $\Im(f(z))$ garde un signe constant sur chaque demi-disque
limité par $]-1,1[$, et comme $f(z) = z + o_{z\to0}(z)$, ce signe est celui de~$\Im z$.
    \item $ \int_{t=0}^{\pi }\Im(f(re^{it}))\sin nt\,d t = \frac{\pi a_nr^n}2$.

On a $|\sin(nt)|\le n\sin(t)$ pour $0\le t\le\pi$ par récurrence,
donc $\frac{\pi |a_n|r^n}2\le n \int_{t=0}^{\pi }\Im(f(re^{it}))\sin t\,d t = \frac{n\pi a_1r}2$.
On en déduit $|a_n|r^n\le n|a_1|r$ et on conclut $|a_n|\le n$ en faisant tendre $r$ vers~$1$.
  \end{enumerate}
\end{enumerate}
 \fincorrection
\correction{005749}
Soit $R > 0$. Notons $D_R$ le disque fermé de centre $0$ et de rayon $R$.
Soient $z\in D_R$ et $n$ un entier naturel.

\begin{center}
$|P_n(z)| =\left|e^z - (e^z-P_n(z))\right|\geqslant|e^z|-|e^z-P_n(z)|\geqslant e^{-R}-|e^z-P_n(z)|$.
\end{center}

On sait que la suite de polynômes $(P_n)_{n\in\Nn}$ converge uniformément vers la fonction exponentielle sur $D_R$.Donc il existe un entier $n_0$ tel que pour tout $z\in D_R$ et tout entier $n\geqslant n_0$, $\left|e^z-P_n(z)\right|\leqslant\frac{1}{2}e^{-R}$. Pour $n\geqslant n_0$ et $z\in D_R$, $\left|P_n(z)\right|\geqslant\frac{1}{2}e^{-R}>0$ et $P_n$ ne s'annule pas dans $D_R$.
\fincorrection
\correction{005750}
On cherche une série entière $\sum_{n=0}^{+\infty}b_nx^n$ de rayon $R$ strictement positif telle que $\left(\sum_{n=0}^{+\infty}a_nx^n\right)\left(\sum_{n=0}^{+\infty}b_n^xn\right) = 1$ pour $x$ élément d'un certain intervalle ouvert non vide de centre $0$.

Cette égalité impose à la suite $(b_n)$ de vérifier le système d'équations $\left\{
\begin{array}{l}
a_0b_0=1\\
a_0b_1+a_1b_0=0\\
a_0b_2+a_1b_1+a_2b_0=0\\
\rule{1.5cm}{0mm}\vdots\\
a_0b_n+a_1b_{n-1}+\ldots+a_{n-1}b_1+a_nb_0=1\\
\rule{2.5cm}{0mm}\vdots
\end{array}
\right.$

\begin{enumerate}
 \item  Montrons par récurrence que $\forall n\in\Nn$, $b_n$ existe et est unique.

\textbullet~Puisque $a_0 = 1$, $a_0b_0=1\Leftrightarrow b_0=1$. Ceci montre l'existence et l'unicité de $b_0$.

\textbullet~Soit $n\in\Nn$. Supposons avoir démontré l'existence et l'unicité de $b_0$, $b_1$, \ldots, $b_n$. 

Alors $a_0b_{n+1}+a_1b_n+\ldots+a_nb_1+a_{n+1}b_0=0\Leftrightarrow b_{n+1}=-a_1b_n-\ldots-a_nb_1-a_{n+1}b_0$. Ceci montre l'existence

et l'unicité de $b_{n+1}$.

On a montré par récurrence que la suite $(b_n)$ existe et est unique.

\item  Il faut alors vérifier que la série entière associée à la suite $(b_n)_{n\in\Nn}$ a un rayon de convergence strictement positif .

Soit $R > 0$ le rayon de la série associée à la suite $(a_n)_{n\in\Nn}$ et soit $r$ un réel tel que $0 < r < R$.

On sait que la suite $(a_nr^n)_{n\in\Nn}$ est bornée et il existe $M > 0$ tel que pour tout entier naturel $n$, $|a_n|\leqslant\frac{M}{r^n}$.

$b_0 = 1$ puis $|b_1| = |-a_1b_0|\leqslant\frac{M}{r}$ puis $|b_2|=|-a_2b_0-a_1b_1|\leqslant\frac{M}{r^2}+\frac{M}{r}\times\frac{M}{r}=\frac{M(M+1)}{r^2}$ puis

\begin{center}
$|b_3| =|-a_3b_0-a_2b_1-a_1b_2|\leqslant\frac{M}{r^3}+\frac{M}{r^2}\times+\frac{M}{r}+\frac{M}{r}\times\frac{M(M+1)}{r^2}=\frac{M(M^2+2M+1)}{r^3}=\frac{M(M+1)^2}{r^3}$.
\end{center}

Montrons alors par récurrence que $\forall n\in\Nn^*$, $b_n|\leqslant\frac{M(M+1)^{n-1}}{r^n}$.

\textbullet~C'est vrai pour $n=1$.

\textbullet~Soit $n\geqslant1$, supposons que $\forall k\in\llbracket1,n\rrbracket$, $|b_k|\leqslant\frac{M(M+1)^{k-1}}{r^k}$. Alors 

\begin{align*}\ensuremath
\left|b_{n+1}\right|&\leqslant|-a_{n+1}b_0| +|-a_nb_1| + ...+ |-a_1b_n|\leqslant\frac{M}{r^{n+1}}+\sum_{k=1}^{n}\frac{M(M+1)^{k-1}}{r^k}\times\frac{M}{r^{n+1-k}}\\
 &=\frac{M}{r^{n+1}}\left(1+ M\sum_{k=1}^{n}(M+1)^{k-1}\right)=\frac{M}{r^{n+1}}\left(1+ M\frac{(M+1)^n-1}{(M+1)-1}\right)=\frac{M(M+1)^n}{r^{n+1}}.
\end{align*}

On a montré par récurrence que pour tout entier naturel non nul $n$, $|b_n|\leqslant\frac{M(M+1)^n}{r^{n+1}}$. En particulier , le rayon $R'$ de la série entière associée à la suite $(b_n)_{n\in\Nn}$ vérifie $R'\geqslant\frac{r}{M+1}> 0$. Ceci valide les calculs initiaux sur $]-\rho,\rho[$ où $\rho=\text{Min}(R,R')>0$ et donc l'inverse d'une fonction $f$ développable en série entière à l'origine et telle que $f(0)\neq0$ est développable en série entière à l'origine.
\end{enumerate}
\fincorrection
\correction{005755}
Posons $\text{Sp}_\Cc A=\left(\lambda_1,\ldots,\lambda_p\right)$. On sait que pour tout entier naturel $n$, $\text{Tr}(A^n) =\lambda_1^n+...+\lambda_p^n$.

Soit $\lambda$ un nombre complexe. 

\textbullet~Si $\lambda= 0$, la série entière associée à la suite $(\lambda^n)_{n\in\Nn}$ est de rayon infini et pour tout nombre complexe $z$, $\sum_{n=0}^{+\infty}\lambda^nz^n=1=\frac{1}{1-\lambda z}$.

\textbullet~Si $\lambda\neq0$, la série entière associée à la suite $(\lambda^n)$ est de rayon $\frac{1}{|\lambda|}$ et pour $|z| <\frac{1}{|\lambda|}$,  $\sum_{n=0}^{+\infty}\lambda^nz^n=\frac{1}{1-\lambda z}$.

Soit $\rho=\text{Max}\left(|\lambda_1|,\ldots,|\lambda_p|\right)$ ($\rho$ est le rayon spectral de la matrice $A$) et $R=\frac{1}{\rho}$  si $\rho\neq0$ et $R = +\infty$ si $\rho= 0$.

Pour $|z| < R$,  

\begin{align*}\ensuremath
\sum_{n=0}^{+\infty}\text{Tr}(A^n)z^n&=\sum_{n=0}^{+\infty}\left(\sum_{k=1}^{p}(\lambda_kz)^n\right)
=\sum_{k=1}^{p}\left(\sum_{n=0}^{+\infty}(\lambda_kz)^n\right)\;(\text{somme de}\;p\;\text{séries convergentes})\\
 &=\sum_{k=1}^{p}\frac{1}{1-\lambda_kz}.
\end{align*}

Il est alors clair que $R$ est le rayon de convergence de la série entière proposée (développement en série entière d'une fraction rationnelle).

Si de plus, $0< |z| < R$,  $\sum_{n=0}^{+\infty}\text{Tr}(A^n)z^n=\frac{1}{z}\left(\sum_{k=1}^{p}\frac{1}{\frac{1}{z}-\lambda_k}\right)=\frac{\chi_A'\left(\frac{1}{z}\right)}{\frac{1}{z}\chi_A\left(\frac{1}{z}\right)}$ (décomposition usuelle de $\frac{P'}{P}$). 
\fincorrection
\correction{005759}
\begin{enumerate}
 \item  Soient $A$ et $B$ les sommes des séries entières associées aux suites $a$ et $b$ sur $]-1,1[$. La fonction $B$ est strictement positive sur $]0,1[$ et en particulier ne s'annule pas sur $]0,1[$.

\textbullet~La suite $a$ est positive donc la fonction $A$ est croissante sur $[0,1[$ et admet ainsi une limite réelle ou infinie quand $x$ tend vers $1$ par valeurs inférieures. De plus, pour $N$ entier naturel donné et $x\in[0,1[$, on a $\sum_{n=0}^{+\infty}a_nx^n\geqslant\sum_{n=0}^{N}a_nx^n$ et donc

\begin{center}
$\forall N\in\Nn$, $\displaystyle\lim_{\substack{x\rightarrow1,\;x<1}}A(x)\geqslant\displaystyle\lim_{\substack{x\rightarrow1,\;x<1}}\sum_{n=0}^{N}a_nx^n =\sum_{n=0}^{N}a_n$.
\end{center}

Puisque la série de terme général positif $a_n$ diverge, quand $N$ tend tend vers $+\infty$, on obtient $\displaystyle\lim_{\substack{x\rightarrow1,\;x<1}}A(x)\geqslant+\infty$ et donc $\displaystyle\lim_{\substack{x\rightarrow1,\;x<1}}A(x)=+\infty$. Il en est de même pour $B$ car la série de terme général $b_n$ diverge quelque soit la valeur de $k$.

\textbullet~On veut alors montrer que $A-k B\underset{x\rightarrow1^-}{=}o(B)$.

Soit $\varepsilon> 0$. Par hypothèse, $a_n-kb_n\underset{n\rightarrow+\infty}o(b_n)$ et donc il existe un entier naturel $N$ tel que pour $n\geqslant N$, $|a_n-kb_n|\leqslant\frac{\varepsilon}{2} b_n$.

Soit $x\in[0,1[$.

\begin{center}
$|A(x)-kB(x)|\leqslant \sum_{n=0}^{+\infty}|a_n-kb_n|x^n\leqslant \sum_{n=0}^{N}|a_n-kb_n|x^n+   \frac{\varepsilon}{2}\sum_{n=N+1}^{+\infty}b_nx^n\leqslant\sum_{n=0}^{N}|a_n-kb_n|+\frac{\varepsilon}{2}B(x)$.
\end{center}

Maintenant, $B(x)$ tend vers $+\infty$ quand $x$ tend vers $1$ par valeurs inférieures. Donc il existe $\alpha\in]0,1[$ tel que pour $x\in]1-\alpha,1[$,  $B(x) >\frac{2}{\varepsilon}\sum_{n=0}^{N}|a_n-kb_n|$. Pour $x\in]1-\alpha,1[$, on a alors $|A(x)-kB(x)| <\frac{\varepsilon}{2}B(x)+\frac{\varepsilon}{2}B(x)=\varepsilon B(x)$.

On a montré que $\forall\varepsilon>0$, $\exists\alpha\in]0,1[/$ $\forall x\in]1-\alpha,1[$, $|A(x)-kB(x)|<\varepsilon B(x)$ et donc $\lim_{x \rightarrow 1^-}\frac{A(x)}{B(x)}=k$.

\item  
  \begin{enumerate}
  \item La série entière proposée \og vérifie \fg les hypothèses du 1) et de plus , $\ln n\underset{n\rightarrow+\infty}{\sim}1+\frac{1}{2}+...+\frac{1}{n}$. Donc 

\begin{center}
$f(x)\underset{x\rightarrow1^-}{\sim}\sum_{n=1}^{+\infty}\left(\sum_{k=1}^{n}\frac{1}{k}\right)x^n=\left(\sum_{n=0}^{+\infty}x^n\right)\left(\sum_{n=1}^{+\infty}\frac{x^n}{n}\right) =\frac{\ln(1-x)}{x-1}$.
\end{center}

\begin{center}
\shadowbox{
$\sum_{n=1}^{+\infty}(\ln n) x^n\underset{x\rightarrow1^-}{\sim}\frac{\ln(1-x)}{x-1}$.
}
\end{center}

  \item Soit $p\geqslant2$. $n^{p-1}\underset{n\rightarrow+\infty}{\sim}(n+1)(n+2)...(n+p-1)$. Comme les deux suites $(n^{p-1})$ et $((n+1)(n+2)...(n+p-1))$ vérifient les hypothèses du 1)

\begin{center}
$\sum_{n=0}^{+\infty}n^{p-1}x^n\underset{x\rightarrow1^-}{\sim}\sum_{n=0}^{+\infty}(n+p-1)...(n+1)x^n =\left(\sum_{n=0}^{+\infty}x^n\right)^{(p-1)}=\left(\frac{1}{1-x}\right)^{(p-1)}=\frac{(p-1)!}{(1-x)^p}$.
\end{center}

Par suite,

\begin{center}
\shadowbox{
$\lim_{x \rightarrow 1^-}(1-x)^p\sum_{n=1}^{+\infty}n^{p-1}x^n=(p-1)!$.
}
\end{center}
  \end{enumerate}
\end{enumerate}
\fincorrection
\correction{005760}
Supposons qu'il existe un entier naturel $p$ tel que $a_p = a_{p+1}$. Le développement limité à l'ordre $1$ de $f^{(p)}$ en $0$ s'écrit $f^{(p)}(x)\underset{x\rightarrow0}{=}f^{(p)}(0)+xf^{(p+1)}(0) +o(x) = a_p(1+x) + o(x)$ et on en déduit 

\begin{align*}\ensuremath
|f^{(p)}(x)|&\geqslant|a_p(1+x)|-|o(x)|=1+x-|o(x)|\geqslant 1+x -\frac{x}{2}\;(\text{sur un voisinage pointé de}\;0\;\text{à droite})\\
 &=1 +\frac{x}{2}>1\;(\text{sur un voisinage pointé de}\;0\;\text{à droite}).
\end{align*}

Donc si deux termes consécutifs sont égaux, $f$ nne vérifie pas les conditions de l'énoncé ou encore si $f$ vérifie les conditions de l'énoncé, alors $\forall p\in\Nn$, $a_{p+1}= - a_p$ puis $a_p = (-1)^pa_0$. Mais alors, nécessairement pour tout réel $x$, $f(x) =e^{-x}$ ou pour tout réel $x$, $f(x) = -e^{-x}$.

Réciproquement, ces deux fonctions sont clairement solutions du problème posé.
\fincorrection
\correction{005764}
\begin{enumerate}
 \item  Soient $n\geqslant2$ puis $k\in\llbracket1,n-1\rrbracket$. On met une parenthèse autour de $X_1...X_k$ et une autour de $X_{k+1}...X_n$. Ensuite, pour chacun des $a_k$ parenthésages de $X1...X_k$, il y a $a_{n-k}$ parenthésages possibles de $X_{k+1}...X_n$. Finalement, en faisant varier $k$ de $1$ à $n-1$, on a montré que 

\begin{center}
\shadowbox{
$\forall n\geqslant2$, $a_n =\sum_{k=1}^{n-1}a_ka_{n-k}$.
}
\end{center}

\item  On suppose momentanément le rayon $R$ de la série entière associé à la suite $(a_n)_{n\in\Nn^*}$ strictement positif. On pose conventionnellement $a_0=0$. Pour $x\in]-R,R[$,

\begin{center}
$f^2(x)=\left(\sum_{n=0}^{+\infty}a_nx^n\right)^2 =\sum_{n=0}^{+\infty}\left(\sum_{k=0}^{n}a_ka_{n-k}\right)x^n =\sum_{n=2}^{+\infty}\left(\sum_{k=1}^{n-1}a_ka_{n-k}\right)x^n=\sum_{n=2}^{+\infty}a_nx^n= f(x) - x$,
\end{center}

et donc

\begin{center}
\shadowbox{
$\forall x\in]-R,R[$, $f^2(x) = f(x) - x$.
}
\end{center}

\item  Nécessairement, pour tout $x$ de $]-R,R[$, $f(x)=\frac{1}{2}(1+\sqrt{1+4x})$ $(I)$ ou $f(x)=  \frac{1}{2}(1-\sqrt{1-4x})$ $(II)$. Ainsi, pour chaque $x\in]-R,R[$, on doit choisir l'une de ces deux expressions. Puisque $f(0)=0$, il faut choisir l'expression $(II)$ quand $x=0$. 

Pour $x\in\left]-\frac{1}{4},\frac{1}{4}\right[$, posons $g(x)=\frac{1}{2}(1-\sqrt{1-4x})$. $g$ est développable en série entière sur $\left]-\frac{1}{4},\frac{1}{4}\right[$ en vertu de théorèmes généraux. Notons $(b_n)_{n\in\Nn}$ la suite des coefficients du développement. Puisque $g(0)=0$, on a $b_0=0=a_0$ et puisque $g'(0)=1$, on a $b_1=1=a_1$. Enfin, la fonction $g$ vérifie $\forall x\in\left]-\frac{1}{4},\frac{1}{4}\right[$, $g^2(x)=g(x)-x$ et donc $\forall n\geqslant2$, $b_n =\sum_{k=1}^{n-1}b_kb_{n-k}$. On en déduit par récurrence que pour tout entier naturel $n$, $b_n=a_n$ et donc $\forall x\in\left]-\frac{1}{4},\frac{1}{4}\right[$, $f(x)=g(x)$.

\begin{center}
\shadowbox{
$\forall x\in\left]-\frac{1}{4},\frac{1}{4}\right[$, $f(x)=\frac{1}{2}(1-\sqrt{1-4x})$.
}
\end{center}

\item  Pour connaitre les $a_n$, il reste à développer la fonction $g$ en série entière. Pour $x\in\left]-\frac{1}{4},\frac{1}{4}\right[$, 

\begin{center}
$g(x) =\frac{1}{2}(1-(1-4x)^{1/2})=\frac{1}{2}\left(1-\sum_{n=0}^{+\infty}C_{1/2}^n(-4x)^n\right)=\sum_{n=1}^{+\infty}(-1)^{n-1}C_{1/2}^n2^{2n-1}x^n$.
\end{center}

Enfin, pour $n\in\Nn^*$,

\begin{align*}\ensuremath
(-1)^{n-1}C_{1/2}^n2^{2n-1}&=(-1)^{n-1}\frac{\frac{1}{2}\times\left(\frac{1}{2}-1\right)\times\ldots\times\left(\frac{1}{2}-(n-1)\right)}{n!}2^{2n-1}=\frac{2^{n-1}}{n!}\times1\times3\times....\times(2n-3)\\
 &=\frac{2^{n-1}}{n!}\times\frac{1\times2\times3\times....\times(2n-3)\times(2n-2)}{2\times4\times\ldots\times(2n-2)}=\frac{(2n-2)!}{n!(n-1)!}=\frac{C_{2n-2}^{n-1}}{n}.
\end{align*}

Donc

\begin{center}
\shadowbox{
$\forall n\in\Nn^*$, $a_n =\frac{C_{2n-2}^{n-1}}{n}$.
}
\end{center}
\end{enumerate}
\fincorrection
\nocorrection
\nocorrection
\correction{001951}
\begin{enumerate}

	\item $a_0 = \pi$ et pour $n\ne 0$, $a_n = \frac{2}{\pi n^2} ((-1)^n - 1)$ (nul quand $n$ est pair), $b_n=0$.

	\item Par la formule de Parseval 
	\[
	\frac{2 \pi^2}{3} 
	= \frac{1}{\pi} \int_{-\pi}^\pi |x|^2 \, dx
	= \frac{\pi^2}{2} + \frac{16}{\pi^2} \sum_{p=0}^\infty \frac{1}{(2p+1)^4}
	\]
	donc $\sum_{p=0}^\infty \frac{1}{(2p+1)^2} = \frac{\pi^4}{96}$.
	
	\item Comme
	\[
	\sum_{n=1}^\infty \frac{1}{n^4} = \sum_{p=1}^\infty \frac{1}{(2p)^2} + \sum_{p=0}^\infty \frac{1}{(2p+1)^2}
	= \frac{1}{16} \sum_{n=1}^\infty \frac{1}{n^4} + \frac{\pi^4}{96}
	\]
	il vient $\displaystyle \sum_{n=1}^\infty \frac{1}{n^4} = \frac{15}{16} \frac{\pi^4}{96}
	= \frac{\pi^4}{90}$.
	
	\item C'est le théorème de Dirichlet. Appliqué en $x=0$, cela implique que 
	$\displaystyle \sum_{p=0}^\infty \frac{1}{(2p+1)^2} = \frac{\pi}{4} \frac{\pi}{2} = \frac{\pi^2}{8}$. 
	On en déduit que $\displaystyle \sum_{p=0}^\infty \frac{1}{(2p+1)^2} = \frac{\pi^2}{6}$.
	

\end{enumerate}

\fincorrection
\nocorrection
\nocorrection
\correction{004622}
\begin{enumerate}
  \item $a_0 = \pi$, $a_{2p} = 0$, $a_{2p+1} = \frac4{\pi(2p+1)^2}$,
             $b_n = 0$.
  \item $a_n = 0$, $b_n = \frac2n$.
  \item $a_0 = \frac{8\pi^2}3$, $a_n = \frac4{n^2}$,
             $b_n = -\frac{4\pi}n$.
  \item $a_0 = \frac 2\pi$, $a_{2p} = \frac{-2}{\pi(4p^2-1)}$,
             $a_{2p+1} = 0$, $b_1 = \frac 12$, $b_p = 0$.
  \item $a_{2p} = \frac{24}{\pi(4p^2-1)(4p^2-9)}$, $a_{2p+1} = 0$,
             $b_p = 0$.
\end{enumerate}
 \fincorrection
\correction{004623}
$I_{n+1}-I_{n-1} =  \int_{t=0}^{\pi/2} 2\cos(nt)\,d t
                          = \frac2n\sin(n\pi/2)$.

Donc $I_{2p} = 2\Bigl(1 - \frac13 + \dots+\frac{(-1)^{p-1}}{2p-1}\Bigr)$,
     $I_{2p+1} = \frac\pi2$.

$b_n = 0$ (parité), $a_{2p} = 0$ (symétrie par rapport à $(\pi/2,0)$),
$a_{2p+1} = -\frac{4}{(2p+1)\pi}I_{2p+1} = -\frac{2}{2p+1}$.
\fincorrection
\correction{004624}
$$\left\{\begin{array}{lll}2a_k &= &(a_{k-1}+a_{k+1})\cos\alpha\cr
       a_0 &= &a_1\cos\alpha+2\hfill\cr\end{array}\right.
       \Rightarrow  a_k = A\tan(\frac\pi4-\frac\alpha2)^k
                                          + B\mathrm{cotan}(\frac\pi4-\frac\alpha2)^k.$$

Comme $a_k\to 0$ (lorsque $k\to\infty$) on a $B=0$ d'où $A=\frac2{\sin\alpha}$.
Finalement, $f(t) = \frac2{\sin\alpha}\Bigl(\frac12+\sum_{k=1}^\infty
\Bigl(\frac{1-\sin\alpha}{\cos\alpha}\Bigr)^k\cos(kt)\Bigr)$.
\fincorrection
\correction{004625}
\begin{enumerate}
  \item $\sum_{n=0}^\infty a^n\cos nx = \Re\Bigl(\sum_{n=0}^\infty (ae^{ix})^n\Bigr)
    =\Re\Bigl(\frac1{1-ae^{ix}}\Bigr) = g(x)$.
  \item Il y a convergence normale.
  \item \begin{align*} h(x)
    &=  \int_{t=0}^{2\pi}\sum_{n=0}^\infty a^n\cos(nx-nt)f(t)\,d t\cr
    &= \sum_{n=0}^\infty \int_{t=0}^{2\pi} a^n\cos(nx-nt)f(t)\,d t\cr
    &= \sum_{n=0}^\infty\Bigl(\cos(nx) \int_{t=0}^{2\pi} a^n\cos(nt)f(t)\,d t + \sin(nx) \int_{t=0}^{2\pi} a^n\sin(nt)f(t)\,d t\Bigr)\cr
    &= \sum_{n=0}^\infty(\pi a^na_n(f)\cos(nx) + \pi a^nb_n(f)\sin(nx)).\cr
    \end{align*}
    Il y a convergence normale car $|a|<1$ et les coefficients de Fourier
    de~$f$ sont bornés.
    On en déduit que $h$ est continue, puis que les coefficients de Fourier de~$h$
    sont $a^na_n(f)$ et $a^nb_n(f)$.
  \item Les coefficients de Fourier des deux membres doivent être égaux, ce qui
    donne~: $a_n(f)= \frac1{n^2(1-\pi\lambda a^n)}$ et $b_n(f) = 0$ si pour
    tout~$n\in\N^*$ on a $\pi\lambda a^n\ne 1$ (sinon il n'y a pas de solution),
    et $a_0(f) = 0$ si $2\pi\lambda\ne1$, $a_0(f)$ quelconque sinon.
    Réciproquement, en posant ${f(x) = [a_0/2] + \sum_{n=1}^\infty\frac{\cos nx}{n^2(1-\pi\lambda a^n)}}$
    on définit $f$, $2\pi$-périodique continue (la série converge normalement),
    solution de l'équation par égalité des coefficients de Fourier de chaque membre.
\end{enumerate}
 \fincorrection
\correction{004626}
\begin{enumerate}
  \item $f(x) = \frac 12 + \sum_{n=1}^\infty \frac{\cos nx}{2^n}
              = \frac3{2(5-4\cos x)}$.
  \item $\frac\pi3$.
\end{enumerate}
 \fincorrection
\correction{004627}
\begin{enumerate}
  \item 
  \item $g(x) = \frac1{\sh a}\bigl(e^af(e^{ix}) - e^{-ix}f(e^{-ix})\bigr)
                   = \frac1{\sh a}\Bigl(1 + 2\sum_{k=1}^\infty
                      (-1)^ke^{-ka}\cos kx \Bigr)$.
\end{enumerate}
 \fincorrection
\nocorrection
\correction{004629}
\begin{enumerate}
  \item $c_k(h) = c_k(f)c_k(g)$.
  \item 
\end{enumerate}
 \fincorrection
\correction{004630}
$c_k = \frac1{2\pi} \int_{x=0}^{2\pi}\sum_{n=-\infty}^{+\infty}e^{-(x-2n\pi)^2-ikx}\,d x
              = \frac1{2\pi}\sum_{n=-\infty}^{+\infty} \int_{x=0}^{2\pi}e^{-(x-2n\pi)^2-ikx}\,d x
              = \frac1{2\pi} \int_{x=-\infty}^{+\infty}e^{-x^2-ikx}\,d x
              = \frac{e^{-k^2/4}}{2\sqrt\pi}$

$\Bigl($calculer $ \int_{x=-\infty}^{+\infty}e^{-x^2-i\xi x}\,d x = \sqrt\pi e^{-\xi^2/4}$
par équation différentielle$\Bigr)$.
\fincorrection
\correction{004631}
\begin{enumerate}
  \item $a_0 = \frac{2\sh\pi}\pi$,
             $a_n = \frac{2(-1)^n\sh\pi}{\pi(1+n^2)}$, $b_n = -na_n$.
  \item $S = \frac{\pi-\tanh\pi}{2\tanh\pi}$,
             $S' = \frac{\pi-\sh\pi}{2\sh\pi}$.
\end{enumerate}
 \fincorrection
\correction{004632}
\begin{enumerate}
  \item Si $a\ne 0$~: $S_f(x) = \sum_{n=-\infty}^{+\infty} \frac{e^{2\pi a}-1}{\strut2\pi(a-in)}e^{inx}
    = \frac{e^{2\pi a}-1}{2\pi a} + \sum_{n=1}^{+\infty}\frac{e^{2\pi a}-1}{\pi(a^2+n^2)}(a\cos(nx) - n\sin(nx))$.
  \item On peut supposer $a > 0$ car $I(-a) = -I(a)$ et $I(0) = 0$.
    On envisage d'intégrer terme à terme la
    relation~: $$\frac{e^{-u}}{1-e^{-u}}\sin(au) = \sum_{n=1}^{\infty}e^{-nu}\sin(au).$$
    On coupe l'intégrale $\int_0^{+\infty}$ en $\int_0^{\pi/a}+\int_{\pi/a}^{+\infty}$~:
    sur~$[0,\pi/a]$ le sinus est positif et le théorème de convergence monotone
    s'applique. Sur~$[\pi/a,+\infty[$ le théorème d'intégration terme
    à terme s'applique (série des normes~$1$ convergente) car
    $\int_{\pi/a}^{+\infty}|e^{-nu}\sin(au)|\,d u \le \int_{\pi/a}^{+\infty}e^{-nu}\,d u 
    = e^{-n\pi/a}/n$. Ainsi,
    $$I(a) = \sum_{n=1}^\infty  \int_{u=0}^{+\infty}e^{-nu}\sin(au)\,d u
           = \sum_{n=1}^\infty  \frac{a}{n^2+a^2}.$$
  \item Déjà fait, $ \int_{u=0}^{+\infty}e^{-u}\sin(au)\,d u =  \frac{a}{a^2+1}$.
    Il doit y avoir une autre méthode pour la question précédente ?!
  \item En comparant avec~{\bf 1)} pour~$x=0$ on obtient~:
    $I(a) = \frac\pi{\strut2}\frac{e^{2a\pi}+1}{e^{2a\pi}-1}-\frac1{\strut2a}$ pour $a>0$.
\end{enumerate}
 \fincorrection
\correction{004633}
\begin{enumerate}
  \item $S_f(x) = \sum_{n=-\infty}^{+\infty} \frac{e^{2\pi a}-1}{\strut2\pi(a-in)}e^{inx}$.
  \item $\sum_{n\in\Z}\frac{(e^{2\pi a}-1)^2}{4\pi^2(a^2+n^2)} = \frac1{\strut2\pi} \int_{t=0}^{2\pi}|f(t)|^2\,d t
    = \frac{e^{4a\pi}-1}{4a\pi}$
    donc $\sum_{n\ge1}\frac{a}{a^2+n^2} = \frac\pi{\strut2}\frac{e^{2a\pi}+1}{e^{2a\pi}-1}-\frac1{\strut2a}$.
    
    $\sum_{n\ge1}\frac{1}{a^2+n^2} = \frac1{\strut2a}\Bigl(\frac\pi{\tanh(a\pi)}-\frac1a\Bigr)
    \to \frac{\pi^2}{\strut6}$ lorsque $a\to0 $ et il y a convergence dominée.
  \item $\frac\pi{\strut2}$.
\end{enumerate}
 \fincorrection
\correction{004634}
\begin{enumerate}
  \item $a_n = -\frac{4}{4n^2-1}$,
         $b_n = -\frac{32n}{\pi(4n^2-1)^2}$.
  \item 
  \item $\frac{\pi-2}4$.
\end{enumerate}
 \fincorrection
\correction{004635}
\begin{enumerate}
  \item $\cos ax = \frac{\sin \pi a}{\pi a} +
              \frac{2a\sin\pi a}\pi\sum_{n=1}^\infty
              \frac{(-1)^n\cos nx}{a^2-n^2}$.
  \item $g'(t) = \sum_{n=1}^\infty \frac{2t}{t^2-n^2}
                    = \frac{\pi\cos\pi t}{\sin\pi t} - \frac 1t  \Rightarrow 
              g(t) = \ln\left(\lambda\frac{\sin\pi t}t\right)$ et $g(0) = 0  \Rightarrow 
              g(t) = \ln\left(\frac{\sin\pi t}{\pi t}\right)$.
\end{enumerate}
 \fincorrection
\correction{004636}
\begin{enumerate}
  \item $f(x) = \sum_{n=1}^\infty \frac{\sin nx}n$.
  \item 
  \item Le premier membre vaut $f(1) = \frac{\pi-1}2$ et le second
             $\frac1{4\pi} \int_{t=0}^{2\pi} (f(t+1)-f(t-1))^2\,d t
             = \frac{\pi-1}2$.
\end{enumerate}
\fincorrection
\correction{004637}
\begin{enumerate}
  \item $a_{2p+1} = b_{2p+1} = 0$.
  \item $a_{2p} = b_{2p} = 0$.
\end{enumerate}
 \fincorrection
\nocorrection
\correction{004639}
$a'_k = kb_k$, $b'_k = -ka_k$ $+$ inégalité de Bessel.
\fincorrection
\correction{004640}
\begin{enumerate}
  \item 
  \item $S_g'(x) = \frac{f(2\pi)-f(0)}{2\pi}
                    + \sum_{n=1}^\infty \bigl(nb_n - \frac{f(0) - f(2\pi)}{\pi}\bigr)
                    \cos nx - na_n\sin nx$.
\end{enumerate}
 \fincorrection
\nocorrection
\correction{004642}
\begin{align*}
\pi a_k &=  \int_{t=0}^{2\pi}f(t)\cos(kt)\,d t\cr
        &= \sum_{i=0}^{k-1} \int_{t=0}^{2\pi/k}f(t+2i\pi/k)\cos(kt)\,d t\cr
        &= \sum_{i=0}^{k-1} \int_{t=0}^{\pi/2k}\Bigl(
                 f(t+2i\pi/k) - f(t+2i\pi/k+\pi/k)
                -f(t+2(i+1)\pi/k-\pi/k) + f(t+2(i+1)\pi/k)\Bigr)\cos(kt)\,d t.\cr\end{align*}
\fincorrection
\correction{005782}
\begin{enumerate}
 \item  La fonction $f$ est continue par morceaux sur $\Rr$ et $2\pi$-périodique. On peut donc calculer ses coefficients de \textsc{Fourier}.

$$\includegraphics{../images/img005782-1}$$


Puisque $f$ est paire, $\forall n\in\Nn^*$, $b_n(f)=0$ puis pour $n\in\Nn$, $a_n(f)=\frac{2}{\pi}\int_{0}^{\pi}\left(1-\frac{2x}{\pi}\right)\cos(nx)\;dx$.

Par suite, $a_0(f)=0$ puis pour $n\in\Nn^*$,

\begin{center}
$a_n(f)=\frac{2}{\pi}\left(\left[\left(1-\frac{2x}{\pi}\right)\frac{\sin(nx)}{n}\right]_0^\pi+\frac{2}{n\pi}\int_{0}^{\pi}\sin(nx)\;dx\right)=\frac{4}{n\pi^2}\left[\frac{-\cos(nx)}{n}\right]_0^\pi=\frac{4(1-(-1)^n)}{n^2\pi^2}$.
\end{center} 

La fonction $f$ est $2\pi$-périodique, continue sur $\Rr$ et de classe $C^1$ par morceaux sur $\Rr$. D'après le théorème de \textsc{Dirichlet}, la série de \textsc{Fourier} de $f$ converge vers $f$ sur $\Rr$. Par suite, pour tout réel $x$, 

\begin{center}
$f(x)=\frac{a_0(f)}{2}+\sum_{n=1}^{+\infty}(a_n(f)\cos(nx)+b_n(f)\sin(nx))=\frac{4}{\pi^2}\sum_{n=1}^{+\infty}\frac{1-(-1)^n}{n^2}\cos(nx)=\frac{8}{\pi^2}\sum_{p=0}^{+\infty}\frac{\cos((2p+1)x)}{(2p+1)^2}$.
\end{center}

\begin{center}
\shadowbox{
$\forall x\in\Rr$, $f(x)=\frac{8}{\pi^2}\sum_{n=0}^{+\infty}\frac{\cos((2n+1)x)}{(2n+1)^2}$.
}
\end{center}

L'égalité $f(0)=1$ fournit $\sum_{n=0}^{+\infty}\frac{1}{(2n+1)^2}=\frac{\pi^2}{8}$. Ensuite, si $S=\sum_{n=1}^{+\infty}\frac{1}{n^2}$, on a

\begin{center}
$S=\sum_{n=0}^{+\infty}\frac{1}{(2n+1)^2}+\sum_{n=1}^{+\infty}\frac{1}{(2n)^2}=\frac{\pi^2}{8}+\frac{S}{4}$,
\end{center}

et donc $S=\frac{4}{3}\times\frac{\pi^2}{8}=\frac{\pi^2}{6}$.

D'autre part, puisque $f$ est continue par morceaux sur $\Rr$ et $2\pi$-périodique, la formule de \textsc{Parseval} fournit $\frac{(a_0(f))^2}{2}+\sum_{n=1}^{+\infty}((a_n(f))^2+(b_n(f))^2)=\frac{1}{\pi}\int_{-\pi}^{\pi}(f(x))^2\;dx$ et donc

\begin{center}
$\frac{64}{\pi^4}\sum_{n=0}^{+\infty}\frac{1}{(2n+1)^4}=\frac{2}{\pi}\int_{0}^{\pi}\left(1-\frac{2x}{\pi}\right)^2\;dx=\left[-\frac{1}{3}\left(1-\frac{2x}{\pi}\right)^3\right]_0^\pi=\frac{2}{3}$
\end{center}

et donc $\sum_{n=0}^{+\infty}\frac{1}{(2n+1)^4}=\frac{2}{3}\times\frac{\pi^4}{64}=\frac{\pi^4}{96}$. Enfin, si on pose $S=\frac{1}{n^4}$,

\begin{center}
$S=\sum_{n=0}^{+\infty}\frac{1}{(2n+1)^4}+\sum_{n=1}^{+\infty}\frac{1}{(2n)^4}=\frac{\pi^4}{96}+\frac{S}{16}$,
\end{center}

et donc $S=\frac{16}{15}\times\frac{\pi^4}{96}=\frac{\pi^4}{90}$.

\begin{center}
\shadowbox{
$\sum_{n=0}^{+\infty}\frac{1}{(2n+1)^2}=\frac{\pi^2}{8}$, $\sum_{n=1}^{+\infty}\frac{1}{n^2}=\frac{\pi^2}{6}$ et $\sum_{n=1}^{+\infty}\frac{1}{n^4}=\frac{\pi^4}{90}$.
}
\end{center}

\item  La fonction $f$ est continue par morceaux sur $\Rr$ et $2\pi$-périodique. On peut donc calculer ses coefficients de \textsc{Fourier}.

$$\includegraphics{../images/img005782-2}$$



Puisque $f$ est impaire, $\forall n\in\Nn$, $a_n(f)=0$ puis pour $n\in\Nn^*$, 

\begin{align*}\ensuremath
b_n(f)&=\frac{2}{\pi}\int_{0}^{\pi}x(\pi-x)\sin(nx)\;dx=\frac{2}{\pi}\left(\left[x(\pi-x)\frac{-\cos(nx)}{n}\right]_0^\pi+\frac{1}{n}\int_{0}^{\pi}(\pi-2x)\cos(nx)\;dx\right)\\
 &=\frac{2}{n\pi}\left(\left[(\pi-2x)\frac{\sin(nx)}{n}\right]_0^\pi+\frac{2}{n}\int_{0}^{\pi}\sin(nx)\;dx\right)=\frac{4}{n^2\pi}\left[-\frac{\cos(nx)}{n}\right]_0^\pi=\frac{4(1-(-1)^n)}{n^3\pi}.
\end{align*} 

La fonction $f$ est $2\pi$-périodique, continue sur $\Rr$ et de classe $C^1$ par morceaux sur $\Rr$. D'après le théorème de \textsc{Dirichlet}, la série de \textsc{Fourier} de $f$ converge vers $f$ sur $\Rr$. Par suite, pour tout réel $x$, 

\begin{center}
$f(x)=\frac{a_0(f)}{2}+\sum_{n=1}^{+\infty}(a_n(f)\cos(nx)+b_n(f)\sin(nx))=\frac{4}{\pi}\sum_{n=1}^{+\infty}\frac{1-(-1)^n}{n^3}\sin(nx)=\frac{8}{\pi}\sum_{p=0}^{+\infty}\frac{\sin((2p+1)x)}{(2p+1)^3}$.
\end{center}

\begin{center}
\shadowbox{
$\forall x\in\Rr$, $f(x)=\frac{8}{\pi}\sum_{n=0}^{+\infty}\frac{\sin((2n+1)x)}{(2n+1)^3}$.
}
\end{center}

L'égalité $f\left(\frac{\pi}{2}\right)=\frac{\pi^2}{4}$ fournit $\sum_{n=0}^{+\infty}(-1)^n\frac{1}{(2n+1)^3}=\frac{\pi^3}{32}$. Ensuite,  puisque $f$ est continue par morceaux sur $\Rr$ et $2\pi$-périodique, la formule de \textsc{Parseval} fournit $\frac{(a_0(f))^2}{2}+\sum_{n=1}^{+\infty}((a_n(f))^2+(b_n(f))^2)=\frac{1}{\pi}\int_{-\pi}^{\pi}(f(x))^2\;dx$ et donc

\begin{center}
$\frac{64}{\pi^2}\sum_{n=0}^{+\infty}\frac{1}{(2n+1)^6}=\frac{2}{\pi}\int_{0}^{\pi}x^2(\pi-x)^2\;dx=\frac{2}{\pi}\left[\pi^2\frac{x^3}{3}-2\pi\frac{x^4}{4}+\frac{x^5}{5}\right]_0^\pi=2\pi^4\left(\frac{1}{3}-\frac{1}{2}+\frac{1}{5}\right)=\frac{\pi^4}{15}$
\end{center}

et donc $\sum_{n=0}^{+\infty}\frac{1}{(2n+1)^6}=\frac{\pi^2}{64}\times\frac{\pi^4}{15}=\frac{\pi^6}{960}$. Enfin, si on pose $S=\frac{1}{n^6}$,

\begin{center}
$S=\sum_{n=0}^{+\infty}\frac{1}{(2n+1)^6}+\sum_{n=1}^{+\infty}\frac{1}{(2n)^6}=\frac{\pi^6}{960}+\frac{S}{64}$,
\end{center}

et donc $S=\frac{64}{63}\times\frac{\pi^6}{960}=\frac{\pi^6}{945}$.

\begin{center}
\shadowbox{
$\sum_{n=0}^{+\infty}(-1)^n\frac{1}{(2n+1)^3}=\frac{\pi^3}{32}$, $\sum_{n=1}^{+\infty}\frac{1}{(2n+1)^6}=\frac{\pi^6}{960}$ et $\sum_{n=1}^{+\infty}\frac{1}{n^6}=\frac{\pi^6}{945}$.
}
\end{center}

\item   La fonction $f$ est continue par morceaux sur $\Rr$ et $2\pi$-périodique. On peut donc calculer ses coefficients de \textsc{Fourier}.

$$\includegraphics{../images/img005782-3}$$


La fonction $f$ a mêmes coefficients de \textsc{Fourier} que la fonction $g$ définie sur $\Rr$, impaire et $2\pi$-périodique telle que $\forall x\in\left]0,\frac{\pi}{2}\right[$, $g(x)=0$. Donc $\forall n\in\Nn$, $a_n(f)=0$ puis pour $n\in\Nn^*$,

\begin{align*}\ensuremath
b_n(f)&=\frac{2}{\pi}\int_{0}^{\pi}\sin\left(\frac{x}{2}\right)\sin(nx)\;dx=\frac{1}{\pi}\int_{0}^{\pi}\left(\cos\left(\left(n-\frac{1}{2}\right)x\right)-\cos\left(\left(n+\frac{1}{2}\right)x\right)\right)dx\\
 &=\frac{1}{\pi}\left[\frac{\sin\left(\left(n-\frac{1}{2}\right)x\right)}{n-\frac{1}{2}}-\frac{\sin\left(\left(n+\frac{1}{2}\right)x\right)}{n+\frac{1}{2}}\right]_0^\pi=\frac{1}{\pi}\left(-\frac{(-1)^n}{n-\frac{1}{2}}-\frac{(-1)^n}{n+\frac{1}{2}}\right)\\
 &=-\frac{(-1)^n}{\pi}\frac{2n}{n^2-\frac{1}{4}}=-\frac{(-1)^n}{\pi}\frac{8n}{4n^2-1}.
\end{align*}

La fonction $f$ est $2\pi$-périodique et de classe $C^1$ par morceaux sur $\Rr$. D'après le théorème de \textsc{Dirichlet}, la série de \textsc{Fourier} de $f$ converge en tout réel $x$ et a pour somme $\frac{1}{2}(f(x^+)+f(x^-)$. En particulier,

\begin{center}
\shadowbox{
$\forall x\in]-\pi,\pi[$, $\sin\left(\frac{x}{2}\right)=-\frac{8}{\pi}\sum_{n=1}^{+\infty}(-1)^n\frac{n}{4n^2-1}\sin(nx)$.
}
\end{center}

L'égalité $f\left(\frac{\pi}{2}\right)=\frac{1}{\sqrt{2}}$ fournit 

\begin{center}
$\frac{1}{\sqrt{2}}=-\frac{8}{\pi}\sum_{n=0}^{+\infty}(-1)^n\frac{n}{4n^2-1}\sin\left(n\frac{\pi}{2}\right)=\frac{8}{\pi}\sum_{p=0}^{+\infty}\frac{2p+1}{4(2p+1)^2-1}\sin\left((2p+1)\frac{\pi}{2}\right)=\frac{8}{\pi}\sum_{p=0}^{+\infty}(-1)^p\frac{2p+1}{16p^2+1p+3}$,
\end{center}

\begin{center}
\shadowbox{
$\sum_{n=0}^{+\infty}(-1)^n\frac{2n+1}{16n^2+16n+3}=\frac{\pi}{8\sqrt{2}}$.
}
\end{center}

\item  $f$ est $2\pi$-périodique, continue par morceaux sur $\Rr$ et paire. Pour $n\in\Nn^*$, $b_n(f)=0$ puis pour $n\in\Nn$, $a_n(f)=\frac{1}{\pi}\int_{-\pi}^{\pi}\ch(\lambda x)\cos(nx)\;dx$.

\textbf{1ère solution.} Soit $n\in\Nn$.

\begin{align*}\ensuremath
a_n(f)&=\frac{1}{\pi}\text{Re}\left(\int_{-\pi}^{\pi}\ch(\lambda x)e^{inx}\;dx\right)=\frac{1}{2\pi}\text{Re}\left(\int_{-\pi}^{\pi}e^{(\lambda+in)x}\;dx+\int_{-\pi}^{\pi}e^{(-\lambda+in)x}\;dx\right)\\
 &=\frac{1}{2\pi}\text{Re}\left(\frac{e^{(\lambda+in)\pi}-e^{-(\lambda+in)\pi}}{\lambda+in}+\frac{e^{(-\lambda+in)\pi}-e^{-(-\lambda+in)\pi}}{-\lambda+in}\right)=\frac{(-1)^n}{2\pi}\text{Re}\left(\frac{2\sh(\lambda\pi)}{\lambda+in}+\frac{-2\sh(\lambda\pi)}{-\lambda+in}\right)\\
 &=\frac{(-1)^n\sh(\lambda\pi)}{\pi}\text{Re}\left(\frac{\lambda-in}{\lambda^2+n^2}+\frac{\lambda+in}{\lambda^2+n^2}\right)=\frac{2\lambda\sh(\lambda\pi)}{\pi}\times\frac{(-1)^n}{n^2+\lambda^2}
\end{align*}

\textbf{2ème solution.} Une double intégration par parties fournit

\begin{align*}\ensuremath
a_n(f)&=\frac{1}{\pi}\left(\left[\frac{\sh(\lambda x)}{\lambda}\cos(nx)\right]_{-\pi}^{\pi}+\frac{n}{\lambda}\int_{-\pi}^{\pi}\sh(\lambda x)\sin(nx)\;dx\right)=\frac{1}{\pi}\left(\frac{2(-1)^n\sh(\lambda\pi)}{\lambda}+\frac{n}{\lambda}\int_{-\pi}^{\pi}\sh(\lambda x)\sin(nx)\;dx\right)\\
 &=\frac{1}{\pi}\left(\frac{2(-1)^n\sh(\lambda\pi)}{\lambda}+\frac{n}{\lambda}\left(\left[\frac{\ch(\lambda x)}{\lambda}\sin(nx)\right]_{-\pi}^{\pi}-\frac{n}{\lambda}\int_{-\pi}^{\pi}\ch(\lambda x)\cos(nx)\;dx\right)\right)\\
 &=\frac{2(-1)^n\sh(\lambda\pi)}{\lambda\pi}-\frac{n^2}{\lambda^2}a_n(f),
\end{align*}

et donc $\forall n\in\Nn$, $a_n(f)=\frac{2(-1)^n\sh(\lambda\pi)}{\lambda\pi}\times\frac{\lambda^2}{n^2+\lambda^2}=\frac{2\lambda\sh(\lambda\pi)}{\pi}\times\frac{(-1)^n}{n^2+\lambda^2}$.

La fonction $f$ est $2\pi$-périodique, continue sur $\Rr$ et de classe $C^1$ par morceaux sur $\Rr$. D'après le théorème de \textsc{Dirichlet}, la série de \textsc{Fourier} de $f$ converge vers $f$ sur $\Rr$. On en déduit que

\begin{center}
\shadowbox{
$\forall x\in\Rr$, $f(x)=\frac{\sh(\lambda\pi)}{\lambda\pi}+\frac{2\lambda\sh(\lambda\pi)}{\pi}\sum_{n=1}^{+\infty}\frac{(-1)^n}{n^2+\lambda^2}\cos(nx)$.
}
\end{center}

L'égalité $f(0)=1$ fournit $1=\frac{\sh(\lambda\pi)}{\lambda\pi}+\frac{2\lambda\sh(\lambda\pi)}{\pi}\sum_{n=1}^{+\infty}\frac{(-1)^n}{n^2+\lambda^2}$ et donc

\begin{center}
$\sum_{n=1}^{+\infty}\frac{(-1)^n}{n^2+\lambda^2}=\frac{\pi}{2\lambda\sh(\lambda\pi)}\left(1-\frac{\sh(\lambda\pi)}{\lambda\pi}\right)=\frac{\pi\left(\sh(\lambda\pi)-\pi\lambda\right)}{2\lambda^2\pi\sh(\lambda\pi)}$
\end{center}

et l'égalité $f(\pi)=\ch(\lambda\pi)$ fournit 

\begin{center}
$\sum_{n=1}^{+\infty}\frac{1}{n^2+\lambda^2}=\frac{\pi}{2\lambda\sh(\lambda\pi)}\left(\ch(\lambda\pi)-\frac{\sh(\lambda\pi)}{\lambda\pi}\right)=\frac{\lambda\pi\ch(\lambda\pi)-\sh(\lambda\pi)}{2\lambda^2\sh(\lambda\pi)}$
\end{center}

\begin{center}
\shadowbox{
$\forall\lambda>0$, $\sum_{n=1}^{+\infty}\frac{(-1)^n}{n^2+\lambda^2}=\frac{\pi}{2\lambda\sh(\lambda\pi)}$ et $\sum_{n=1}^{+\infty}\frac{1}{n^2+\lambda^2}=\frac{\lambda\pi\ch(\lambda\pi)-\sh(\lambda\pi)}{2\lambda^2\sh(\lambda\pi)}$.
}
\end{center}

La fonction $f$ est $2\pi$-périodique, continue par morceaux sur $\Rr$. L'égalité de \textsc{Parseval} s'écrit $\frac{(a_0(f))^2}{2}+\sum_{n=1}^{+\infty}((a_n(f))^2+(b_n(f))^2)=\frac{1}{\pi}\int_{-\pi}^{\pi}(f(x))^2\;dx$ avec

\begin{center}
$\frac{1}{\pi}\int_{-\pi}^{\pi}(f(x))^2\;dx=\frac{1}{\pi}\int_{-\pi}^{\pi}\ch^2(\lambda x)\;dx=\frac{1}{\pi}\int_{-\pi}^{\pi}\frac{\ch(2\lambda x)+1}{2}\;dx=1+\frac{\sh(2\lambda\pi)}{2\pi}$,
\end{center}

et donc $1+\frac{\sh(2\lambda\pi)}{2\pi}=\frac{2\sh^2(\lambda\pi)}{\pi^2\lambda^2}+\frac{4\lambda^2\sh^2(\lambda\pi)}{\pi^2}\sum_{n=1}^{+\infty}\frac{1}{(\lambda^2+n^2)^2}$ puis

\begin{center}
$\sum_{n=1}^{+\infty}\frac{1}{(\lambda^2+n^2)^2}=\frac{\pi^2}{4\lambda^2\sh^2(\lambda\pi)}\left(1+\frac{\sh(2\lambda\pi)}{2\pi}-\frac{2\sh^2(\lambda\pi)}{\pi^2\lambda^2}\right)=\frac{2\pi^2\lambda^2+\pi\lambda\sh(2\lambda\pi)-4\lambda^2\sh^2(\lambda\pi)}{8\lambda^4\sh^2(\lambda\pi)}$.
\end{center}

\begin{center}
\shadowbox{
$\forall\lambda>0$, $\sum_{n=1}^{+\infty}\frac{1}{(\lambda^2+n^2)^2}=\frac{\pi^2\lambda^2+\pi\lambda\ch(\lambda\pi)\sh(\lambda\pi)-2\lambda^2\sh^2(\lambda\pi)}{4\lambda^2\sh^2(\lambda\pi)}$.
}
\end{center}

\item  La fonction $f$ est continue par morceaux sur $\Rr$ et $2\pi$-périodique. On peut donc calculer ses coefficients de \textsc{Fourier}.

$$\includegraphics{../images/img005782-4}$$


Soit $n\in\Nn$.

\begin{align*}\ensuremath
a_n(f)&=\frac{1}{\pi}\int_{-\pi}^{\pi}\text{sup}(\sin x,0)\cos(nx)\;dx=\frac{1}{\pi}\int_{0}^{\pi}\sin x\cos(nx)\;dx=\frac{1}{2\pi}\int_{0}^{\pi}\sin((n+1)x)-\sin((n-1)x)\;dx\\
 &=\left\{
 \begin{array}{l}
 \frac{1}{2\pi}\int_{0}^{\pi}\sin(2x)\;dx\;\text{si}\;n=1\\
 \frac{1}{2\pi}\left[-\frac{\cos((n+1)x)}{n+1}+\frac{\cos((n-1)x)}{n-1}\right]_0^\pi\;\text{si}\;n\neq1
 \end{array}
 \right.=\left\{
 \begin{array}{l}
 \frac{1}{2\pi}\left[-\frac{\cos(2x)}{2}\right]_0^\pi\;\text{si}\;n=1\\
 \frac{1}{2\pi}\left(-\frac{(-1)^{n+1}-1}{n+1}+\frac{(-1)^{n-1}-1}{n-1}\right)\;\text{si}\;n\neq1
 \end{array}
 \right.\\
 &=\left\{
 \begin{array}{l}
\rule[-2mm]{0mm}{0mm}0\;\text{si}\;n=1\\
-\frac{1+(-1)^n}{\pi}\frac{1}{n^2-1}\;\text{si}\;n\neq1
 \end{array}
 \right.
\end{align*}

Soit $n\in\Nn^*$.

\begin{align*}\ensuremath
b_n(f)&=\frac{1}{\pi}\int_{0}^{\pi}\sin x\sin(nx)\;dx=\frac{1}{2\pi}\int_{0}^{\pi}(\cos((n-1)x)-\cos((n+1)x))\;dx=\left\{
 \begin{array}{l}
 \frac{1}{2}\;\text{si}\;n=1\\
\rule{0mm}{4mm} 0\;\text{si}\;n\neq1
 \end{array}
 \right.\end{align*}

La fonction $f$ est $2\pi$-périodique, continue sur $\Rr$ et de classe $C^1$ par morceaux sur $\Rr$. D'après le théorème de \textsc{Dirichlet}, la série de \textsc{Fourier} de $f$ converge vers $f$ sur $\Rr$. On en déduit que pour tout réel $x$

\begin{center}
$\text{sup}(\sin x,0)=\frac{1}{\pi}+\frac{\sin x}{2}-\frac{1}{\pi}\sum_{n=2}^{+\infty}\frac{1+(-1)^n}{n^2-1}\cos(nx)=\frac{1}{\pi}+\frac{\sin x}{2}-\frac{2}{\pi}\sum_{p=1}^{+\infty}\frac{1}{4p^2-1}\cos(2px)$.
\end{center}

\begin{center}
\shadowbox{
$\forall x\in\Rr$, $\text{sup}(\sin x,0)=\frac{1}{\pi}+\frac{\sin x}{2}-\frac{2}{\pi}\sum_{n=1}^{+\infty}\frac{1}{4n^2-1}\cos(2nx)$.
}
\end{center}

L'égalité $f(0)=0$ fournit $\frac{1}{\pi}-\frac{2}{\pi}\sum_{n=1}^{+\infty}\frac{1}{4n^2-1}=0$ et donc

\begin{center}
\shadowbox{
$\sum_{n=1}^{+\infty}\frac{1}{4n^2-1}=\frac{1}{2}$.
}
\end{center}

\textbf{Remarque.} $\sum_{n=1}^{+\infty}\frac{1}{4n^2-1}=\lim_{N \rightarrow +\infty}\frac{1}{2}\sum_{n=1}^{N}\left(\frac{1}{2n-1}-\frac{1}{2n+1}\right)=\lim_{N \rightarrow +\infty}\frac{1}{2}\left(1-\frac{1}{2N+1}\right)=\frac{1}{2}$
\end{enumerate}
\fincorrection
\correction{005783}
\begin{enumerate}
 \item 
   \begin{enumerate}
   \item Soit $a\in\Cc\setminus[-1,1]$. Pour tout réel $t$, $a-\cos t\neq0$ et

\begin{center}
$\frac{1}{a-\cos t}=\frac{2}{2a-e^{it}-e^{-it}}=\frac{-2e^{it}}{(e^{it})^2-2ae^{it}+1}$.
\end{center}

L'équation $z^2-2az+1=0$ admet deux solutions non  nulles inverses l'une de l'autre. On note $b$ la solution de plus petit module de sorte que $|b|\leqslant1$.

On ne peut avoir $|b|=1$ car alors il existe $\theta\in\Rr$ tel que $b=e^{i\theta}$. On en déduit que $2a=b+\frac{1}{b}=2\cos\theta\in[-2,2]$ puis que $a\in[-1,1]$ ce qui n'est pas. Donc $|b|\neq1$. Plus précisément, puisque $|b|\leqslant\left|\frac{1}{b}\right|$, on a $|b|<1$ et $\left|\frac{1}{b}\right|$. En particulier, $b\neq\frac{1}{b}$.

Ensuite, pour $|t|<|b|$, on a

\begin{align*}\ensuremath
\frac{1}{a-\cos t}&=\frac{-2e^{it}}{(e^{it}-b)\left(e^{it}-\frac{1}{b}\right)}=\frac{2}{\frac{1}{b}-b}\left(\frac{b}{e^{it}-b}-\frac{1/b}{e^{it}-\frac{1}{b}}\right)=\frac{2b}{1-b^2}\left(\frac{be^{-it}}{1-be^{-it}}+\frac{1}{1-be^{it}}\right)\\
 &=\frac{2b}{1-b^2}\left(be^{-it}\sum_{n=0}^{+\infty}b^ne^{-int}+\sum_{n=0}^{+\infty}b^ne^{int}\right)\;(\text{car}\;|be^{it}|=|be^{-it}|=|b|<1)\\
 &=\frac{2b}{1-b^2}\left(\sum_{n=0}^{+\infty}b^{n+1}e^{-i(n+1)t}+\sum_{n=0}^{+\infty}b^ne^{int}\right)=\frac{2b}{1-b^2}\left(1+\sum_{n=1}^{+\infty}b^ne^{int}+\sum_{n=1}^{+\infty}b^{n}e^{-int}\right)\\
 &=\frac{2b}{1-b^2}\left(1+2\sum_{n=1}^{+\infty}b^n\cos(nt)\right).
\end{align*}

\begin{center}
\shadowbox{
$\forall t\in\Rr$, $\frac{1}{a-\cos t}=\frac{2b}{1-b^2}\left(1+2\sum_{n=1}^{+\infty}b^n\cos(nt)\right)$.
}
\end{center}

   \item Pour tout réel $t\in[-\pi,\pi]$ et tout entier naturel non nul $n$, on a $|b^n\cos(nt)|\leqslant|b|^n$. Comme la série numérique de terme général $|b|^n$ converge, on en déduit que la série de fonctions de terme général $t\mapsto b^n\cos(nt)$, $n\in\Nn$, converge normalement et donc uniformément sur le segment $[-\pi,\pi]$.

On sait alors que la série obtenue est la série de \textsc{Fourier} de $f$.

  \end{enumerate}

\item  Puisque la fonction $f$ est paire, pour tout entier naturel $n$, $a_n(f)=\frac{2}{\pi}\int_{0}^{\pi}\frac{\cos(nt)}{a-\cos t}\;dt$. Donc, pour tout entier naturel $n$ (y compris pour $n=0$),

\begin{center}
$\int_{0}^{\pi}\frac{\cos(nt)}{a-\cos t}\;dt=\frac{\pi a_n(f)}{2}=\frac{2b^{n+1}\pi}{1-b^2}$
\end{center}

Finalement,

\begin{center}
\shadowbox{
$\forall n\in\Nn$, $\int_{0}^{\pi}\frac{\cos(nt)}{a-\cos t}\;dt=\frac{2b^{n+1}\pi}{1-b^2}$.
}
\end{center}
\end{enumerate}
\fincorrection
\correction{005784}
\begin{enumerate}
 \item  Soit $\alpha\in\Cc\setminus\Zz$. La fonction $f$ est $2\pi$-périodique, continue sur $\Rr$ et de classe $C^1$ par morceaux sur $\Rr$. Donc la série de \textsc{Fourier} de $f$ converge vers $f$ sur $\Rr$ d'après le théorème de \textsc{Dirichlet}.

Puisque $f$ est paire, $\forall n\in\Nn^*$, $b_n(f)=0$ puis pour $n\in\Nn$,

\begin{align*}\ensuremath
a_n(f)&=\frac{2}{\pi}\int_{0}^{\pi}\cos(\alpha x)\cos(nx)\;dx=\frac{1}{\pi}\int_{0}^{\pi}\left(\cos((n+\alpha)x)+\cos((n-\alpha)x)\right)dx\\
 &=\frac{1}{\pi}\left[\frac{\sin((\alpha+n)x)}{\alpha+n}+\frac{\sin((\alpha-n)x)}{\alpha-n}\right]_0^\pi\;(\text{car}\;\alpha\notin\Zz)\\
 &=\frac{1}{\pi}\left(\frac{\sin((\alpha+n)\pi)}{\alpha+n}+\frac{\sin((\alpha-n)\pi)}{\alpha-n}\right)=(-1)^n\frac{2\alpha\sin(\alpha\pi)}{\pi(\alpha^2-n^2)}
\end{align*}

Finalement,

\begin{center}
\shadowbox{
$\forall\alpha\in\Cc\setminus\Zz$, $\forall x\in[-\pi,\pi]$, $\cos(\alpha x)=\frac{\sin(\alpha\pi)}{\alpha\pi}+\frac{\sin(\alpha\pi)}{\pi}\sum_{n=1}^{+\infty}(-1)^n\frac{2\alpha}{\alpha^2-n^2}\cos(nx)$.
}
\end{center}

\item  Soit $z\in\Cc\setminus\Zz$.

On prend $\alpha=z$ et $x=0$ dans la formule précédente et on obtient $1=\frac{\sin(\pi z)}{\pi z}+\frac{\sin(\pi z)}{\pi}\sum_{n=1}^{+\infty}(-1)^n\frac{2z}{z^2-n^2}$ $(*)$. Maintenant,

\begin{center}
$\sin(\pi z)=0\Leftrightarrow\frac{1}{2i}(e^{i\pi z}-e^{-i\pi z})=0\Leftrightarrow e^{i\pi z}=e^{-i\pi z}\Leftrightarrow e^{2i\pi z}=1\Leftrightarrow 2i\pi z\in2i\pi\Zz\Leftrightarrow z\in\Zz$.
\end{center}

Puisque $z\in\Cc\setminus\Zz$, $\sin(\pi z)\neq 0$ et l'égalité $(*)$ peut s'écrire $\frac{\pi}{\sin(\pi z)}=\frac{1}{z}+\sum_{n=1}^{+\infty}(-1)^n\frac{2z}{z^2-n^2}$.

De même, en prenant $\alpha=z$ et $x=\pi$, on obtient $\cos(\pi z)=\frac{\sin(\pi z)}{\pi z}+\frac{\sin(\pi z)}{\pi}\sum_{n=1}^{+\infty}\frac{2z}{z^2-n^2}$ et donc $\pi\mathrm{cotan}(\pi z)=\frac{1}{z}+\sum_{n=1}^{+\infty}\frac{2z}{z^2-n^2}$.

\begin{center}
\shadowbox{
$\frac{\pi}{\sin(\pi z)}=\frac{1}{z}+\sum_{n=1}^{+\infty}(-1)^n\frac{2z}{z^2-n^2}$ et $\pi\mathrm{cotan}(\pi z)=\frac{1}{z}+\sum_{n=1}^{+\infty}\frac{2z}{z^2-n^2}$.
}
\end{center}
\end{enumerate}
\fincorrection
\correction{005785}
La fonction $f$ est $1$-périodique, continue par morceaux sur $\Rr$. On peut donc calculer ses coefficients de \textsc{Fourier}. 

$$\includegraphics{../images/img005785-1}$$



La fonction $f$ a mêmes coefficients de \textsc{Fourier} que la fonction $g~:~x\mapsto\left\{
\begin{array}{l}
f(x)\;\text{si}\;x\notin\Zz\\
0\;\text{si}\;x\in\Zz
\end{array}
\right.$ qui est impaire. Donc, $\forall n\in\Nn$, $a_n(f)=0$ puis pour $n\in\Nn^*$

\begin{align*}\ensuremath
b_n(f)&=\frac{2}{1}\int_{0}^{1}f(t)\sin\left(\frac{2n\pi t}{1}\right)\;dt=\int_{0}^{1}(2t-1)\sin(2n\pi t)\;dt\\
 &=\left[-\frac{(2t-1)\cos(2n\pi t)}{2n\pi}\right]_0^1+\frac{1}{n\pi}\int_{0}^{1}\cos(2n\pi t)\;dt=\left(-\frac{1}{2n\pi}-\frac{1}{2n\pi}\right)+0\\
 &=-\frac{1}{n\pi}.
\end{align*}

La fonction $f$ est de plus de classe $C^1$ par morceaux sur $\Rr$ et d'après le théorème de \textsc{Dirichlet}, en tout réel $x$, la série de \textsc{Fourier} de $f$ converge et a pour pour somme 
$\frac{1}{2}(f(x^+)+f(x^-))$. En particulier,

\begin{center}
\shadowbox{
$\forall x\in\Rr\setminus\Zz$, $f(x)=x-E(x)-\frac{1}{2}=-\sum_{n=1}^{+\infty}\frac{\sin(2n\pi x)}{n\pi}$. 
}
\end{center}

\item  Soit $p\in\Nn^*$. Pour $n\in\Nn^*$,

\begin{align*}\ensuremath
b_n(f_p)&=2\int_{0}^{1}f(pt)\sin(2n\pi t)\;dt=2\int_{0}^{p}f(u)\sin\left(2n\pi \frac{u}{p}\right)\frac{du}{p}\\
 &=\left[-\frac{(2t-1)\cos(2n\pi t)}{2n\pi}\right]_0^1+\frac{1}{n\pi}\int_{0}^{1}\cos(2n\pi t)\;dt=\left(-\frac{1}{2n\pi}-\frac{1}{2n\pi}\right)+0\\
 &=-\frac{1}{n\pi}.
\end{align*}

\textbf{Remarque.} Soient $p\in\Nn^*$ et $x\in[0,1]\setminus\left\{
\frac{k}{p},\;k\in\llbracket0,p\rrbracket\right\}$. Alors $px\notin\Zz$ et donc

\begin{center}
$f_p(x)=f(px)=-\sum_{n=1}^{+\infty}\frac{\sin(2np\pi x)}{n\pi}=\sum_{k=1}^{+\infty}b_{k,p}\sin(2k\pi x)$
\end{center}

où $\forall k\in\Nn^*$, $b_{k,p}=\left\{
\begin{array}{l}
0\;\text{si}\;k\notin p\Zz\\
-\frac{1}{\frac{k}{p}\pi}\;\text{si}\;k\in p\Zz
\end{array}
\right.$ mais malheureusement, on ne peut pas récupérer ces coefficients car la série obtenue ne converge pas normalement.

\begin{center}
\shadowbox{
$\forall(p,q)\in(\Nn^*)^2$, $\int_{0}^{1}f_q(x)f_q(x)\;dx=\frac{\left(\text{PGCD}(p,q)\right)^2}{12pq}$.
}
\end{center}
\fincorrection
\correction{004650}
\begin{enumerate}
  \item $\frac\pi{n+1}$.
  \item Somme de Riemman : $\ell =  \int_{t=0}^\pi \frac{\sin t}t\,d t$.
\end{enumerate}
 \fincorrection
\correction{004651}
$\sup\limits_{[\alpha,\beta]}|a_n \cos nx + b_n\sin nx|
          = \sqrt{a_n^2 + b_n^2}$ pour $n \ge \frac{\beta-\alpha}{2\pi}$.
\fincorrection
\correction{004652}
S'il y a convergence uniforme :
$\|a_n\sin nx + \dots + a_p\sin px\|_\infty \to 0$ lorsque $n,p\to\infty$.

On prend $x = \frac\pi{2p}$ : $0 \le \frac{a_p}p(n+\dots+p)
                                  \le \frac1p(na_n + \dots + pa_p)
                               \to 0$ lorsque $n,p\to\infty$.
$n = [p/2]  \Rightarrow $ cqfd.

Si $na_n \to 0$ : Soit $x \in {]0,\pi]}$ et $n$ tel que
$\frac 1n \le x\le \frac1{n-1}$.

Transformation d'Abel :
$|a_n\sin nx + \dots + a_p\sin px| \le \frac{2a_n}{\sin x/2} \le \frac{2na_n}\pi$,

et $|a_k\sin kx + \dots + a_{n-1}\sin (n-1)x| \le (ka_k + \dots + (n-1)a_{n-1})x
\le \frac{n-k}{n-1}a_k\le 2a_k$.

\fincorrection
\nocorrection
\correction{004654}
\begin{enumerate}
  \item $R(n) = \frac an + S(n)$ avec $\deg S \le -2$.
             Donc $f(x) = R(0) + 2ia\sum_{n=1}^\infty \frac{\sin nx}n
                        + \sum_{n=1}^\infty (S(n)e^{inx} + S(-n)e^{-inx})$
             converge pour tout $x \in \R$.
  \item $R(n) = \frac {a_1}n + \dots + \frac{a_k}{n^k} + S(n)$
             avec $\deg S \le -k-1  \Rightarrow 
             f(x) = R(0) + a_1f_1(x) + \dots + a_kf_k(x) + g(x)$
             avec $f_p(x) = \sum_{n=1}^\infty
             \frac{e^{inx}+(-1)^pe^{-inx}}{n^p}$
             et $g$ de classe $\mathcal{C}^k$ sur $\R\setminus2\pi\Z$.
             $f_p' = if_{p-1}$ et $f_1$ est $\mathcal{C}^\infty$ sur $\R\setminus2\pi\Z$
             donc $f_p$ aussi.
\end{enumerate}
 \fincorrection
\correction{004655}
\begin{enumerate}
  \item $f(x) = \frac{\pi^2}3 + 4\sum_{n=1}^\infty
             \frac{\cos nx}{n^2}$.
  \item Soit $P(n) = an^2 + bn + c$.
             Alors $f(x) - 4ag(x) = \frac{\pi^2}3
             + 4\sum_{n=1}^\infty \frac{bn+c}{an^4+bn^3+cn^2}\cos nx$.
\end{enumerate}
 \fincorrection
\correction{004656}
\begin{enumerate}
  \item $k_n(x)= \frac{1-\cos((n+1)x)}{(n+1)(1-\cos x)}
                    = \frac{\sin^2((n+1)x/2)}{(n+1)\sin^2(x/2)}$.
  \item 
\end{enumerate}
 \fincorrection
\nocorrection
\nocorrection
\correction{004643}
\begin{enumerate}
  \item Immédiat. La fonction prolongée est $\mathcal{C}^1$ sur~$\R$ et $\mathcal{C}^2$ par morceaux.
  \item On décompose $f$ en série de Fourier~:
    $f(x) = - \sum_{n=1}^\infty \frac{c_n}{n^2\pi^2}\sin(n\pi x)$ avec
    $c_n = 2 \int_{u=0}^1 f''(u)\sin(n\pi u)\,d u$.
    En appliquant l'inégalité de Cauchy-Schwarz on obtient~:
    $\|f\|_\infty^2\le\Bigl(\sum_{n=1}^\infty\frac1{n^4\pi^4}\Bigr)\Bigr(\sum_{n=1}^\infty c_n^2\Bigr)
    = \frac{2\zeta(4)}{\pi^4}\|f''\|_2^2 = \frac{\|f''\|_2^2}{45}$.
    
    Autre démonstration sans utiliser les séries de Fourier~: pour $x\in{[0,1]}$
    on a
    \begin{align*}
    f(x) &{}=  \int_{t=0}^x f'(t)\,d t = xf'(x) - \int_{t=0}^xtf'(t)\,d t\cr
    f(x) &{}=  \int_{t=1}^x f'(t)\,d t = (x-1)f'(x) - \int_{t=1}^x(t-1)f'(t)\,d t\cr
    f(x) &= (1-x)f(x) + xf(x) =  \int_{t=0}^x t(x-1)f''(t)\,d t +  \int_{t=x}^1 x(t-1)f''(t)\,d t\cr
         &=  \int_{t=0}^1\varphi(x,t)f''(t)\,d t.\text{ avec }
          \varphi(x,t) = xt - \min(x,t).\cr\end{align*}
    On en déduit $|f(x)|^2\le \|f''\|_2^2 \int_{t=0}^1\varphi(x,t)^2\,d t
    = \frac{x^2(x-1)^2}3\|f''\|_2^2\le \frac{\|f''\|_2^2}{48}$.
\end{enumerate}
 \fincorrection
\correction{004644}
Parseval pour $f$ et $f'$. \'Egalité ssi $f(x) = a\cos x + b\sin x$.
\fincorrection
\correction{004645}
\begin{enumerate}
  \item Développer $f$, $f'$ et $g'$ en séries de Fourier et
appliquer l'inégalité $2|\overline a\,b| \le |a|^2 + |b|^2$.
Il y a égalité si et seulement si $f'$ et $g'$ sont CL de cos et sin.
  \item On paramètre par une abscisse curviligne~ :
$x=f(t)$, $y=g(t)$ $ \Rightarrow {\cal A} =  \int_0^{2\pi}fg' \le  \int_0^{2\pi} \frac{f'^2+g'^2}2 = \pi$.
\end{enumerate}
 \fincorrection
\nocorrection
\nocorrection
\nocorrection
\correction{004649}
On pose $g(t) = f(a|t|/\pi)$ pour $t\in{[-\pi,\pi]}$, prolongée par $2\pi$-périodicité.
Alors $g$ est paire, continue, et tous ses coefficients de Fourier sont nuls donc $g=0$.
\fincorrection
\nocorrection
\nocorrection
\nocorrection
\correction{004659}
\begin{enumerate}
  \item $|\sin x| = \frac2\pi -\frac 4\pi\sum_{n=1}^\infty
             \frac{\cos 2nx}{4n^2-1}$.
  \item 
\end{enumerate}
 \fincorrection
\correction{004660}
$|\sin x| = \frac2\pi -\frac 4\pi\sum_{n=1}^\infty
         \frac{\cos 2nx}{4n^2-1}
          \Rightarrow  y = \frac2\pi -\frac4\pi\sum_{n=1}^\infty
         \frac{\cos 2nx}{(4n^2-1)(16n^4-4n^2+1)}$.
         \par
         Cette série converge et définit une fonction de classe $\mathcal{C}^4$ solution
         de l'équation.

         Unicité : les solutions de l'équation homogène sont combinaison de
         $e^{jx}$, $e^{-jx}$, $e^{j^2x}$ et $e^{-j^2x}$ donc non
         $\pi$-périodiques.
\fincorrection
\correction{004661}
$k \notin\Z$ : $y = \sum_{k=1}^\infty \frac{\cos nx}{n^2(k^2-n^2)}
                        + a\cos kx + b\sin kx$.\par
         $k\in\Z$ : remplacer $\frac{\cos kx}{k^2(k^2-k^2)}$ par
                        $\frac{x\cos kx}{2k^3}$.
\fincorrection
\correction{004662}
\begin{enumerate}
  \item Développer $f$ en série de Fourier.
  \item Densité des polynômes trigonométriques dans $\mathcal{C}^0$.
  \item $f(t) = \sin^2(\pi t)$, $\alpha = \frac1\pi$, $x = 0$ :
             $S_n = \sin^21 + \dots + \sin^2n \sim \frac n2$.\par
             Transformation d'Abel : $\sum_{n=1}^N \frac{\sin^2 n}n =
             -\sin^21 + \sum_{n=2}^{N-1} \frac{S_k}{k(k-1)} + \frac{S_N}N
             \to +\infty$ lorsque $N\to\infty$.
    Remarque : on a un raisonnement plus simple en écrivant $2\sin^2(n) = 1-\cos(2n)$.
\end{enumerate}
 \fincorrection
\correction{004663}
Intégration à $x'-x$ constant~: $a_k =  \int_{y=-1}^1 e^{-|y|^\beta}(1-|y|)e^{-2ik\pi y}\,d y$
est le $2k$-ème coefficient de Fourier de la fonction $f$, $2$-périodique,
telle que ${f(y) = e^{-|y|^\beta}(1-|y|)}$ si $-1\le y\le 1$
donc le $k$-ème coefficient de Fourier de $g$, $1$-périodique, telle
que ${g(y) = \frac12(f(y) + f(y+1))}$.
Soit $g_n$ la $n$-ème somme partielle de la série de Fourier de~$g$,
$g_n(y) = \sum_{|k|\le n}a_ke^{2ik\pi y}$.

On a par convergence normale de la série de Fourier de~$g$~:
$\sum_{|k|>n\text{ ou }|\ell|>n} a_ka_\ell = g^2(0) - g_n^2(0)$.

\begin{align*}
g(0)-g_n(0) 
&=  \int_{y=-{\frac 1 2}}^{\frac 1 2} \frac{g(0)-g(y)}{\sin\pi y}\sin((2n+1)\pi y)\,d y\cr
&= 2 \int_{y=0}^{\frac 1 2} \frac{g(0)-g(y)}{\sin\pi y}\sin((2n+1)\pi y)\,d y\cr
&= 2\Bigl[-\frac{g(0)-g(y)}{\sin\pi y}\frac{\cos((2n+1)\pi y)}{(2n+1)\pi}\Bigr]_{y=0}^{\frac 1 2}
+2 \int_{y=0}^{\frac 1 2} \frac{d}{d y}\Bigl(\frac{g(0)-g(y)}{\sin\pi y}\Bigr)\frac{\cos((2n+1)\pi y)}{(2n+1)\pi}\,d y\cr
&=\frac{1-e^{-1}}{(2n+1)\pi^2} +  \int_{y=0}^{\frac 1 2}\mathrm{fct continue}(y)\frac{\cos((2n+1)\pi y)}{(2n+1)\pi}\,d y\cr
&\sim\frac{1-e^{-1}}{(2n+1)\pi^2}.\cr
\end{align*}
$g(0)+g_n(0) \to 2g(0) = 1$ (lorsque $n\to\infty$) d'où
$\sum_{|k|>n\text{ ou }|\ell|>n} a_ka_\ell\sim\frac{1-e^{-1}}{(2n+1)\pi^2}$.
\fincorrection
\correction{004664}
\begin{enumerate}
  \item Si $f=0$ alors $c_n = 0$ pour tout~$n$ par intégration
    terme à terme. Donc une fonction $f\in E$ possède un unique développement
    trigonométrique et $\|f\|$ est bien défini. Alors $E$ est isomorphe
    à $\ell^1(\Z)$ qui est un espace vectoriel normé complet.
  \item Produit de convolution de deux $\Z$-suites sommables.
  \item 
  \begin{enumerate}
    \item $z_0 = \varphi(x \mapsto e^{2i\pi x})$. On a $|z_0| = 1$ car la suite
    $(\varphi(x \mapsto e^{2in\pi x})_{n\in\Z})$ est bornée.
    \item
  \end{enumerate}
\end{enumerate}
 \fincorrection
\correction{004665}
$u_0+\dots+u_n = 1 + (-1)^n \int_{t=0}^1\frac{\cos((n+1)t^2)}{\cos(t^2)}\,d t
=1 + (-1)^n \int_{u=0}^1\frac{\cos((n+1)u)}{2\cos(u)\sqrt u}\,d u\to 1$ lorsque $n\to\infty$.
\fincorrection
\correction{004666}
\begin{enumerate}
  \item Si~$f$ est $\mathcal{C}^1$ $2\pi$-périodique alors $f'$ est continue,
    $2\pi$-périodique de moyenne nulle donc $D(E_0^1)\subset E_0$.
    Ré\-ci\-pro\-que\-ment, si $g\in E_0$ alors toutes les primitives de~$g$ sont
    $2\pi$-périodiques de classe~$\mathcal{C}^1$ et il y en a exactement une qui a
    une valeur moyenne nulle (une seule possibilité de régler la constante).
  \item Non (et ceci quelle que soit la norme) car le spectre de~$D$ n'est par borné.
  \item $\|f\|_2^2 = \sum_{k\in\Z^*}|c_k(f)|^2 \le \sum_{k\in\Z^*}|ikc_k(f)|^2 = \|f'\|^2$.
  \item Idem, $\|\kern-1.2pt|D^{-1}_{|E_n}\|\kern-1.2pt| = \frac1{n+1}$.
\end{enumerate}
 \fincorrection
\correction{004667}
On a $c_0(g) = c_1(g) = c_{-1}(g) = 0$, donc $g$ est orthogonale à tout polynôme trigonométrique de degré
inférieur ou égal à~$1$. Si $g$ est de signe constant sur $]0,2\pi[$, on contredit $c_0(g) = 0$. Donc
$g$ a au moins une racine $a\in{]0,2\pi[}$. Si $g$ n'a pas d'autre racine
dans $]0,2\pi[$ alors $g$ est de signes constants opposés sur $]0,a[$ et $]a,2\pi[$.
Mais alors $g(t)(\cos(t-a/2)-\cos(a/2))$ est de signe constant sur la réunion de
ces intervalles, c'est absurde. Donc $g$ a une deuxième racine dans $]0,2\pi[$,
par exemple $b\in{]a,2\pi[}$. Si $g$ n'a pas d'autre racine sur $]0,2\pi[$
alors $g$ est de signes constants sur $]0,a[$, $]a,b[$ et $]b,2\pi[$ et les
signes alternent. On obtient une nouvelle contradiction car alors
$g(t)(\cos(t-(a+b)/2)-\cos((b-a)/2))$ est de signe constant sur la réunion de
ces intervalles. Ainsi $g$ admet une troisième racine, par exemple $c\in{]b,2\pi[}$.
Enfin, si l'on suppose que $g$ n'a pas d'autre racine sur $]0,2\pi[$ alors on
a $g(t) > 0$ sur $]0,a[$ et $]b,c[$ et $g(t) < 0$ sur $]a,b[$ et $]c,2\pi[$ ou
l'inverse. Dans les deux cas, on en déduit que $g(0) = g(2\pi) = 0$.
\fincorrection
\correction{005781}
\begin{enumerate}
 \item  \textbullet~Puisque $f$ est impaire, $f(0)=0$. Puisque $f$ est impaire et $2\pi$-périodique, $-f(\pi)=f(-\pi)=f(\pi)$ et donc $f(\pi)=0$. Puisque $f$ est $2\pi$-périodique, pour $k\in\Zz$, $f(2k\pi)=f(0)=0$ et $f((2k+1)\pi)=f(\pi)=0$. Finalement, $\forall k\in\Zz$, $f(k\pi)=0$.

Soit $x\in]-\pi,0[$. Puisque $f$ est impaire, $f(x)=-f(-x)=-\sin\left(-\frac{x}{2}\right)=\sin\left(\frac{x}{2}\right)$ et donc $\forall x\in]-\pi,\pi[$, $f(x)=\sin\left(\frac{x}{2}\right)$.

Soit $x\in\Rr\setminus\pi\Zz$. Il existe $k\in\Zz$ tel que $-\pi<x-2k\pi<\pi$ et puisque $f$ est $2\pi$-périodique, $f(x)=f(x-2k\pi)=\sin\left(\frac{x-2k\pi}{2}\right)=(-1)^k\sin\left(\frac{x}{2}\right)$. De plus, $-\pi<x-2k\pi<\pi\Rightarrow k<\frac{x+\pi}{2\pi}<k+1$ et $k=E\left(\frac{x+\pi}{2\pi}\right)$.

\begin{center}
\shadowbox{
$\forall x\in\Rr,\;f(x)=\left\{
\begin{array}{l}
0\;\text{si}\;x\in\pi\Zz\\
(-1)^k\sin\left(\frac{x}{2}\right)\;\text{où}\;k=E\left(\frac{x+\pi}{2\pi}\right)\;\text{si}\;x\notin\pi\Zz
\end{array}
\right.$.
}
\end{center}

\item  \textbullet~Soit $x\in[-\pi,0]$. Puisque $f$ est paire, $f(x)=f(-x)=\sin\left(-\frac{x}{2}\right)=\sin\left(\left|\frac{x}{2}\right|\right)$ et donc $\forall x\in[-\pi,\pi]$, $f(x)=\sin\left(\left|\frac{x}{2}\right|\right)$.

Soit $x\in\Rr$. Il existe $k\in\Zz$ tel que $-\pi<x-2k\pi\leqslant\pi$ et puisque $f$ est $2\pi$-périodique, $f(x)=f(x-2k\pi)=\sin\left(\left|\frac{x-2k\pi}{2}\right|\right)$. 

\begin{center}
\shadowbox{
$\forall x\in\Rr,\;f(x)=\sin\left(\left|\frac{x}{2}-k\pi\right|\right)\;\text{où}\;k=E\left(\frac{x+\pi}{2\pi}\right)$.
}
\end{center}
\end{enumerate}
\fincorrection
\correction{002683}
\vfill
{\leftskip=-1cm A\qquad  1.\par} 
$$\begin{aligned}
\qquad &{d\over dt}\sum_{m=1}^n (-1)^m (z-a)^m
                \phi^{(n-m)}(t) f^{(m)}\left(a+t(z-a)\right) \cr
        &=\sum_{m=1}^n (-1)^m (z-a)^m \Big[
                \phi^{(n-m+1)}(t)f^{(m)}\left(a+t(z-a)\right)  \cr
        &  \mkern260mu   +\phi^{(n-m)}(t) (z-a) f^{(m+1)}\left(a+t(z-a)\right)  \Big]  \cr
        &=\sum_{m=1}^n (-1)^m 
                (z-a)^m \phi^{(n-m+1)}(t)f^{(m)}\left(a+t(z-a)\right)   \cr
        &  \mkern220mu  + \sum_{m=1}^n (-1)^m(z-a)^{m+1} \phi^{(n-m)}(t)f^{(m+1)}\left(a+t(z-a)\right)   \cr
\end{aligned}
$$
Effectuons le changement d'indice de sommation $m=m'+1$ dans la deuxième somme; tous les
termes s'éliminent deux à deux, à l'exception du premier terme de la première somme
et du dernier terme de la deuxième somme, d'où le résultat demandé. 

\medskip
2.a. Plus généralement, on a le résultat suivant: si une fonction $f$ est nulle
en zéro et de classe $C^{n+1}$ sur un intervalle $I$ contenant zéro, $g(t)=f(t)/t$ est
prolongeable par continuité en zéro et son prolongement est de classe $C^n$.
Gardons nous de déduire fallacieusement ce résultat de l'existence d'un développement
limité d'ordre $n$ de $g$.
On montre à l'aide d'un développement limité d'ordre 1 de $f$ que $g$ est
prolongeable par continuité en $0$ en posant $\tilde g(0)=f'(0)$. Par ailleurs, $g$ est de
classe $C^{n+1}$ sur $I\backslash\{0\}$. 
Faisons l'hypothèse de
récurrence $g^{(k-1)}(t)={f^{(k)}(0) \over k} + {f^{(k+1)}(0) \over k+1} {t \over 1!}
+ \dots + {f^{(n+1)}(0) \over n+1} {t^{n-k+1} \over (n-k+1)!}+o(t^{n-k+1})$
$(1\leq k \leq n,\ t\neq0)$.
En dérivant $k$ fois l'identité $f(t)=tg(t)$, on obtient que
$\forall t\neq0 \ \ g^{(k)}(t)={ f^{(k)}(t) - k g^{(k-1)}(t) \over t }$.
Or $f^{(k)}(t)=f^{(k)}(0)+f^{(k+1)}(0){t\over1!}+\dots
+f^{(n+1)}(0){t^{n+1-k}\over(n-k)!}+o(t^{n+1-k})$.
En substituant ces développements limités dans l'identité précédente, on montre que
l'hypothèse de récurrence est vérifiée au rang suivant, donc pour tout entier
$k$ de 1 à $n+1$.
 Faisons l'hypothèse de récurrence que $g^{(k)}(0)$ existe
et est égale à $f^{(k+1)}(0)/(k+1)$ ($0\leq k\leq n$). Du développement limité
de $g^{(k)}$ (tronqué à l'ordre 1) et de l'hypothèse de
récurrence, il résulte que $g^{(k)}$ est continue et dérivable en zéro et que, si
$k<n$, $g^{(k+1)}(0)=f^{(k+2)}(0)/(k+2)$, ce qui prouve par récurrence que $g$ est
$n$ fois continument dérivable sur $I$.

Par conséquent, $t\mapsto{e^t-1 \over t}$ est prolongeable par continuité en une
fonction indéfiniment dérivable sur $\Rr$. Cette fonction ne s'annulant jamais,
son inverse est également définie et indéfiniment dérivable sur $\Rr$.



La fonction $t \mapsto {t \over e^t-1} + {t\over2}$ est paire car
$$\begin{aligned}
{-t \over e^{-t}-1} + {-t\over2}
&= {-t e^t \over 1-e^t} - {t\over2}  \cr
&= {t e^t - t + t \over e^t-1} - {t\over2}  \cr
&= {t \over e^t-1} + {t\over2}  \cr
\end{aligned}
$$
Donc  le développement limité de ${t \over e^t-1}$, dont l'existence est garantie par
le fait que la fonction est indéfiniment dérivable, est de la forme demandée par
l'énoncé.

Par conséquent
\begin{eqnarray*}
{t \over e^t-1} (e^{zt}-1) &=&
\Big( 1-{t\over2}+ {b_1 t^2 \over 2!} + {b_2 t^4 \over 4!} + \dots 
+  {b_N t^{2N} \over (2N)!} + o(t^n) \Big) 
\\
\times \Big( zt + {z^2 t^2 \over 2!} + \dots + {z^n t^n \over n!} + o(t^n)   \Big)   
 \end{eqnarray*}
et $\phi_n(z)/n!$ est le coefficient de $t^n$ dans ce développement:
$$\phi_n(z)/n!= {z^n\over n!} - {1\over2} {z^{n-1}\over(n-1)!}
+ {b_1\over2!}{z^{n-2}\over(n-2)!} + {b_2\over4!}{z^{n-4}\over(n-4)!}
+ \dots + {b_N\over(2N)!}{z^{n-2N}\over(n-2N)!},$$
d'où l'expression de $\phi_n$ demandée.

\medskip
2.b. 
$$\begin{aligned}
\phi_n(z+1)-\phi_n(z)&={d^n \over dt^n}_{| t=0}
                      \left(t{e^{zt}-1 \over e^t-1} - t{e^{(z+1)t}-1 \over e^t-1} \right) \cr
                     &={d^n \over dt^n}_{| t=0} \left(t e^{zt} \right) \cr
\end{aligned}
$$
Comme $t \mapsto t e^{zt}$ est de classe $C^\infty$ et que son développement limité
d'ordre $n$ en $t=0$ est $\sum_{k=0}^{n-1} {z^k\over k!} t^{k+1} + o(t^n)$, il vient
$\phi_n(z+1)-\phi_n(z)=nz^{n-1}$.


\medskip
3. $(i)$ est obtenue en dérivant autant de fois que nécessaire l'identité précédente
et en donnant à $z$ la valeur zéro.
   $(ii)$, $(iii)$, $(iv)$, $(v)$ et $(vi)$ sont des conséquences immédiates de 2.a.

4.a. On applique la question 1 au polynôme $\phi_{2n}$ de degré 2n et on
intègre entre $0$ et $1$. Il vient
\[
\sum_{m=1}^{2n} (-1)^m (z-a)^m \Big[ \phi_{2n}^{(2n-m)}(1) f^{(m)}(z)
- \phi_{2n}^{(2n-m)}(0) f^{(m)}(a)  \Big]
\]\[
 = -(2n)! \bigl( f(z)-f(a) \bigr) + (z-a)^{2n+1} \int_0^1 \phi_{2n}(t)f\Big( a+(z-a)t \Big)dt
\]
en tenant compte du fait que $\phi_{2n}^{(2n)}=(2n)!$

On obtient l'égalité demandée en substituant aux  dérivées itérées
de $\phi_{2n}$ les expressions déterminées dans la question 3.      

\medskip
4.b. Appliquons la question précédente en remplaçant $f$ par une
primitive de $F$ et $z$ par $\omega$. Il vient
\begin{eqnarray*}
0 &=& \int_a^{a+\omega} F(t)dt - {\omega\over2} (F(a+\omega)+F(a)) 
+ \sum_{m=1}^{n-1} b_m {(z-a)^{2m} \over (2m)!} \left[ F^{(2m-1)}(a+\omega) - F^{(2m-1)}(a) \right] 
\\
&& - {\omega^{2n+1}\over (2n)!} \int_0^1 \phi_{2n}(t)F^{(2n)}\Big( a+(z-a)t \Big)dt  
 \end{eqnarray*}
Lorsqu'on somme les égalités obtenues en remplaçant $a$ successivement
par lui-m{\^e}me, $a+\omega$, \dots, $a+(r-1)\omega$, on obtient le résultat demandé,
certains termes se simplifiant deux à deux.

\bigskip
{\leftskip=-1cm B\qquad 1. On a pour tout $x>0$ fixé \par}
$$u_k(x)=\ln \left( 1+{x\over k} \right)
                        + x \ln \left( 1-{k\over k+1} \right)
                        ={x\over k}-{x\over k+1}+O({1\over k^2})
                        =O({1\over k^2})$$
Donc la série $\sum_{k\geq1}u_k(x)$ converge.

\medskip
2. 
$$\begin{aligned}
\ln(x+1)+\sum_{k=1}^n &u_k(x+1)  \cr
        &=\ln(x+1) + \sum_{k=2}^{n+1} \ln(x+k)
                        + \sum_{k=1}^n \Big[ -\ln(k)
                                + (x+1) \big(\ln(k)-\ln(k+1)\big) \Big] \cr
        &=\sum_{k=1}^{n+1} \ln(x+k)                                
                        + \sum_{k=1}^n  x (\ln(k)-\ln(k+1)) 
                        - \sum_{k=1}^n \ln(k+1)        \cr
        &=\sum_{k=1}^n u_k(x) + \ln(x+n+1) - \ln(n+1)  \cr
\end{aligned}
$$
Lorsque $n$ tend vers $+\infty$,
$\ln(x+n+1) - \ln(n+1)=\ln \left( {x+n+1 \over n+1} \right)$ tend vers zéro;
on obtient donc par passage à la limite l'égalité souhaitée.

\medskip
3. Pour tout $k\geq1$ entier, $u_k(1)=0$, donc $G(1)=0$ et
on prouve aisément par récurrence à l'aide de la question précédente
que pour tout entier strictement positif $n$, $G(n+1)=\sum_{k=1}^n \ln(k)$,
égalité de laquelle on déduit immédiatement le résultat demandé.

\medskip
4. immédiat

\medskip
5. C'est une application directe de la question A.4.b.\par
$T_{p,n}(x,y)=-{1\over(2p)!}\int_0^1 \phi_{2p}(t)
                        \sum_{m=0}^{n-1} f^{(2p)}(m+t) dt
=-{1\over(2p)!} \int_0^n \phi_{(2p)} \big( t-E(t) \big) f^{(2p)}(t)dt$


\medskip
6. L'intégrande dans l'expression de $T_{p,n}(x,y)$ est majorée en valeur
absolue par le produit de la borne supérieure de la fonction continue
$\phi_{2p}$ sur le segment $[0,1]$ et de la valeur absolue de $f^{(2p)}$.

On prouve aisément par récurrence que
$f^{(m)}(t)=(-1)^{m-1}\left({1\over(y+t)^m}-{1\over(x+t)^m}\right)
= O \left( {1\over t^{m+1}} \right)$ (quand $t \to +\infty$).
Donc l'intégrale $\int_0^{+\infty} \phi_{(2p)} \big( t-E(t) \big) f^{(2p)}(t)dt$
est absolument convergente, ce qui prouve que $T_{p,n}(x,y)$ admet une limite
finie lorsque $n$ tend vers $+\infty$.



\medskip
7. D'après les questions 4 et 5,
$$\begin{aligned}
G(y)-G(x)&=\sum_{k=1}^{+\infty}\bigl[ \ln(y+k) - \ln(x+k)
                        + (y-x)\ln\left({k\over k+1}\right) \bigr] + \ln y - \ln x \cr
        &=\lim_{n\to +\infty}\Bigl\{\sum_{k=0}^n [\ln(y+k) - \ln(x+k)]
                 + (y-x)\ln{1\over n+1} \Bigr\}                          \cr 
        &\mkern-30mu =\lim_{n\to +\infty} \Biggl\{
                 (y+n)\ln(y+n) - (y+n) - y\ln y + y - (x+n)\ln(x+n)             \cr
        &\mkern40mu  + (x+n) + x\ln x - x
         + {1\over2} \bigl( \ln y - \ln x + \ln(y+n) - \ln(x+n) \bigr)          \cr
        &\mkern10mu +\sum_{h=1}^{p-1} {b_h\over(2h)!}
                \big( f^{(2h-1)}(n) - {1\over y^{2h-1}} + {1\over x^{2h-1}}\big) 
         + T_{p,n}(x,y) + (y-x)\ln{1\over n+1} \Biggr\}                  \cr
\end{aligned}
$$

Or $(y+n)\ln(y+n) - (x+n)\ln(x+n) = (y-x)\ln n + y - x + o(1)$
quand $n\to +\infty$. On obtient donc après simplifications
$G(y)-G(x)=g(x)-g(y)+R_n(x,y)$, ce qu'il fallait démontrer.

\medskip
8. D'après la question 5 et l'expression des dérivées successives de $f$ donnée dans
la question 6, $T_{p,n}(x,y)$ est majoré en valeur absolue par le produit d'une constante
et de l'intégrale
$\int_0^n \left| {1\over (y+t)^{2p}} - {1\over (x+t)^{2p}} \right| dt$.
$|R_p(x,y)|$ est majoré de la m{\^e}me façon en remplaçant la borne finale d'intégration
$n$ par $+\infty$. L'argument de la valeur absolue gardant un signe constant, l'intégrale
majorante est égale à
${1\over 2p-1} \left[ \ \left| {1\over (y+t)^{2p-1}}
                                - {1\over (x+t)^{2p-1}} \right| \ \right]_0^{+\infty}$
et on obtient ainsi l'estimée souhaitée.


\medskip
9. On a $g(m)=m\ln m - m - {1\over2}\ln m + o(1)$ et $G(m)=-\ln(m-1)!=-\ln m! + \ln m$ pour
$m$ entier, donc
le résultat demandé découle immédiatement de la formule de Stirling
$m! \sim \sqrt{2\pi} m^{m+{1\over2}} e^{-m}$.

\medskip
10. Le résultat demandé est obtenu à partir de l'égalité de la question 7 par passage à la
limite. On fait tendre $x$ vers $+\infty$ par valeurs entières et on tient compte de l'estimée
obtenue dans la question 8.

\medskip
11. En calculant les premiers termes du développement limité de la question A.2.a,
on trouve $b_1=1/6,\ b_2=-1/30,\ b_3=1/42$.
Des questions 3 et 10, il résulte que
$$\begin{aligned}
\ln(m!) &= -G(m) + \ln m        \cr
        &= m\ln m - m + {1\over2}\ln m + {1\over2}\ln2\pi
        + {1\over12} {1\over m} - {1\over360} {1\over m^3} + {1\over1260} {1\over m^5}
        + O({1\over m^7})  \cr
\end{aligned}
$$
\fincorrection
\correction{004503}
\begin{enumerate}
  \item 
  \item $\|f_n\|_\infty = f_n\left(\frac1{n+1}\right)
             \sim en^{\alpha-1}$.
\end{enumerate}
 \fincorrection
\nocorrection
\nocorrection
\correction{004506}
\begin{enumerate}
  \item 
  \item Intégrale constante = 1.
\end{enumerate}
 \fincorrection
\correction{004507}
\begin{enumerate}
  \item $e^{-x}$.
  \item 
  \item 
  \item 
\end{enumerate}
 \fincorrection
\correction{004508}
CVU sur tout compact par encadrement du logarithme.
\fincorrection
\nocorrection
\nocorrection
\correction{004511}
\begin{enumerate}
  \item $y_n = (n+1)\left(e^x - e^{nx/(n+1)}\right)$.
  \item $y = xe^x$.
  \item 
\end{enumerate}
 \fincorrection
\correction{004512}
$\bigl(|f_n(x)|\bigr)$ décro\^\i t donc tend vers $L$. On extrait une
sous suite $(f_{\varphi(n)})$ convergeant vers $\ell  \Rightarrow  |\ell| = L$.

La sous-suite $(f_{\varphi(n)+1})$ converge vers
$f(\ell)  \Rightarrow  |f(\ell)| = L  \Rightarrow  L = 0$.
\fincorrection
\correction{004513}
\begin{enumerate}
  \item $\ell(t) = \begin{cases} 0 \text{ si } t=0 \cr
                       \frac {1+\sqrt{1+4t}}2 \text{ si } t >0.\cr\end{cases}$
  \item 
  \item Accroissements finis.
  \item 
\end{enumerate}
 \fincorrection
\nocorrection
\correction{004515}
$f_n(t) \to \sqrt t$ par valeurs croisantes, il y a convergence uniforme.
\fincorrection
\correction{004516}
$P_n\left(t+\frac32\right) = \frac 12 + \frac t4
          \sum_{k=0}^n {\binom{-1/2}{k}}\left(\frac{t^2}4-1\right)^k$.
\fincorrection
\correction{004522}
\begin{enumerate}
  \item Polynôme de Lagrange.
  \item 
\end{enumerate}
 \fincorrection
\correction{004523}
\begin{enumerate}
  \item Il y a convergence simple vers la fonction nulle en $0$ et $1$ et {\'e}gale {\`a} $1/2$ ailleurs. La convergence
est uniforme sur tout $[a,b]\subset {]0,1[}$.
  \item La question pr{\'e}c{\'e}dente donne le r{\'e}sultat pour $1/2$, il suffit alors d'utiliser le th{\'e}or{\`e}me de
Weierstrass et les nombres dyadiques.
\end{enumerate}
 \fincorrection
\nocorrection
\nocorrection
\correction{005726}
\begin{enumerate}
 \item  Pour tout entier naturel $n$, $f_n$ est définie sur $\Rr$ et impaire.

\textbf{Convergence simple sur $\Rr$.} Soit $x\in\Rr$.

\textbullet~Si $x =0$, pour tout entier naturel $n$, $f_n(x)= 0$ et donc $\lim_{n \rightarrow +\infty}f_n(x)=0$. 

\textbullet~Si $x\neq0$, $f_n(x)\underset{n\rightarrow+\infty}{\sim}\frac{1}{nx}$ et de nouveau $\lim_{n \rightarrow +\infty}f_n(x)=0$. 

\begin{center}
\shadowbox{
La suite de fonctions $(f_n)_{n\in\Nn}$ converge simplement sur $\Rr$ vers la fonction nulle.
}
\end{center}

\textbf{Convergence uniforme sur $\Rr$.}
On peut noter tout de suite que pour tout $n\in\Nn^*$, $f_n\left(\frac{1}{n}\right)=\frac{1}{2}$ et donc $\|f_n\|_\infty\geqslant\frac{1}{2}$. On en déduit que $\|f_n\|_\infty$ ne tend pas vers $0$ quand $n$ tend vers $+\infty$.

\begin{center}
\shadowbox{
La suite de fonctions $(f_n)_{n\in\Nn}$ ne converge pas uniformément sur $\Rr$ vers la fonction nulle.
}
\end{center}

Si on n'a pas remarqué ce qui précède, on étudie la fonction $f_n$ sur $\Rr^+$ ($f_n$ étant impaire) dans le but de déterminer $\underset{x\in\Rr}{\text{sup}}|f_n(x) - 0|$.

Soit $n\in\Nn^*$. La fonction $f_n$ est dérivable sur $\Rr^+$ et pour tout réel positif $x$, $f_n'(x) = n\frac{(1+n^2x^2)-x(n^2x)}{(1+n^2x)^2}=\frac{n(1-n^2x^2)}{(1+n^2x)^2}$. Par suite, la fonction $f_n$ est croissante sur $\left[0,\frac{1}{n}\right]$ et décroissante sur $\left[\frac{1}{n},+\infty\right[$.

Puisque la fonction $f_n$ est positive sur $\Rr^+$,  $\underset{x\in\Rr}{\text{sup}}|f_n(x) - 0|= f_n\left(\frac{1}{n}\right)=\frac{1}{2}$ qui ne tend pas vers $0$ quand $n$ tend vers l'infini.

\textbf{Convergence uniforme et localement uniforme sur $]0,+\infty[$.}
La suite de fonctions $(f_n)_{n\in\Nn}$ ne converge toujours pas uniformément vers la fonction nulle sur $]0,+\infty[$ car pour $n\geqslant1$, $\underset{x\in\Rr}{\text{sup}}|f_n(x) - 0|=\frac{1}{2}$.

 
Soit $a$ un réel strictement positif fixé. Soit $n>\frac{1}{a}$. On a $0<\frac{1}{n}<a$ et donc  la fonction $f_n$ est décroissante sur $[a,+\infty[$. Par suite, pour tout réel $x$ de $[a,+\infty[$, $0\leqslant f_n(x)\leqslant f_n(a)$.

Donc $\underset{x\in[a,+\infty[}{\text{sup}}|f_n(x) - 0|= fn(a)$ pour $n>\frac{1}{a}$. On en déduit que  $\lim_{n \rightarrow +\infty}\underset{x\in[a,+\infty[}{\text{sup}}|f_n(x) - 0|=0$.
Donc la suite de fonctions $(f_n)_{n\in\Nn}$ converge uniformément vers la fonction nulle sur tout intervalle de la forme $[a,+\infty[$ où $a > 0$ et en particulier converge localement uniformément vers la fonction nulle sur $]0,+\infty[$ mais ne converge pas uniformément vers la fonction nulle sur $]0,+\infty[$. 

\item  \textbf{Convergence simple sur $\Rr$.} Soit $x\in\Rr$. On sait que $e^x =\lim_{n \rightarrow +\infty}\sum_{k=0}^{n}\frac{x^k}{k!}$ et donc la suite $(f_n)_{n\in\Nn}$ converge simplement sur $\Rr$ vers la fonction constante $f~:~x\mapsto1$.

\textbf{Convergence uniforme sur $\Rr$ et $\Rr^+$.} $\lim_{x \rightarrow -\infty}|f_n(x)-f(x)| = +\infty$. Par suite, pour tout entier naturel $n$, la fonction $|f_n-f|$ n'est pas bornée sur $\Rr$. La suite de fonctions $(f_n)_{n\in\Nn}$ ne converge donc pas uniformément vers $f$ sur $\Rr$.

$\lim_{x \rightarrow +\infty}|f_n(x)-f(x)| = 1$ et donc $\underset{x\in[0,+\infty[}{\text{sup}}|f_n(x)-f(x)|\geqslant1$. La suite de fonctions $(f_n)_{n\in\Nn}$ ne converge donc pas uniformément vers $f$ sur $\Rr^+$.

\textbf{Convergence localement uniforme sur $\Rr$.} Soit $[a,b]$ un segment de $\Rr$.

 
Pour $n\in\Nn^*$, posons $g_n= f_n- f$. La fonction $g_n$ est dérivable sur $\Rr$ et pour $x\in\Rr$

\begin{center}
$g_n'(x)=e^{-x}\left(-\sum_{k=0}^{n}\frac{x^k}{k!}+\sum_{k=0}^{n-1}\frac{x^k}{k!}\right)=-\frac{e^{-x}x^n}{n!}$.
\end{center}

Si $n$ est pair, la fonction $g_n$ est décroissante sur $\Rr$ et s'annule en $0$.

Si $n$ est impair, la fonction $g_n$ est croissante sur $\Rr^-$, décroissante sur $\Rr^+$ et s'annule en $0$.

Dans les deux cas, si $x\in[a,b]$, $|g_n(x)|\leqslant\text{Max}\{|g_n(a)|,|g_n(b)|\}$ avec égalité effectivement obtenue pour $x=a$ ou $x=b$. Donc 

\begin{center}
$\underset{x\in[a,b]}{\text{sup}}|gn(x)|=\text{Max}\{|g_n(a)|,|g_n(b)|\}=\frac{g_n(a)+g_n(b)+|g_n(a)-g_n(b)|}{2}$.
\end{center}

Cette dernière expression tend vers $0$ quand $n$ tend vers $+\infty$. On en déduit que la suite de fonctions $(f_n)_{n\in\Nn}$ converge uniformément vers $f$ sur tout segment $[a,b]$ contenu dans $\Rr$ ou encore 

\begin{center}
\shadowbox{
la suite de fonctions $(f_n)_{n\in\Nn}$ converge localement uniformément vers la fonction $f~:~x\mapsto1$ sur $\Rr$.
}
\end{center}

\item  Pour $x$ réel et $n$ entier naturel, on pose $f_n(x) =n(1-x)^n\sin\left(\frac{\pi}{2}x\right)$.

\textbf{Convergence simple.} Soit $x$ réel fixé. $\sin\left(\frac{\pi}{2}x\right)=0\Leftrightarrow x\in 2\Zz$. Dans ce cas, $\lim_{n \rightarrow +\infty}f_n(x) =0$.

Si $x\notin 2\Zz$, la suite $(f_n(x))_{n\in\Nn}$ converge $\Leftrightarrow$ la suite $(n(1-x)^n)_{n\in\Nn}$ converge $\Leftrightarrow |1-x| < 1 \Leftrightarrow 0 < x < 2$. Dans ce cas, $\lim_{n \rightarrow +\infty}f_n(x)=0$.

\begin{center}
\shadowbox{
La suite de fonctions $(f_n)_{n\in\Nn}$ converge simplement  vers la fonction nulle sur $[0,2]\cup 2\Zz$.
}
\end{center}

\textbf{Convergence uniforme sur $[0,2]$.} Soit  $n$ un entier naturel non nul fixé.

\begin{center}
$\underset{x\in[0,2]}{\text{sup}}|f_n(x)-0|\geqslant\left|f_n\left(\frac{1}{n}\right)\right|=n\left(1-\frac{1}{n}\right)^n\sin\left(\frac{\pi}{2n}\right)$.
\end{center}

Cette dernière expression est équivalente à $\frac{\pi}{2e}$ en $+\infty$ et en particulier ne tend pas vers $0$ quand $n$ tend vers $+\infty$. 

\begin{center}
\shadowbox{
La suite de fonctions $(f_n)_{n\in\Nn}$ ne converge pas uniformément  vers la fonction nulle sur $[0,2]$.
}
\end{center}

$$\includegraphics{../images/img005726-1}$$


\begin{center}
\shadowbox{
La suite de fonctions $(f_n)_{n\in\Nn}$ ne converge pas uniformément  vers la fonction nulle sur $[0,2]$.
}
\end{center}
\end{enumerate}
\fincorrection
\correction{005727}
\textbf{Convergence simple sur $\Rr^+$.} Soit $x$ un réel positif fixé. Pour $n> x$, $f_n(x)=\left(1-\frac{x}{n}\right)^n$ et donc

\begin{center}
$f_n(x)\underset{n\rightarrow+\infty}{=}\left(1-\frac{x}{n}\right)^n\underset{n\rightarrow+\infty}{=}\text{exp}\left(n\ln\left(1-\frac{x}{n}\right)\right)\underset{n\rightarrow+\infty}{=}\text{exp}(-x+o(1)$.
\end{center}

Donc la suite de fonctions $(f_n)_{n\in\Nn^*}$ converge simplement sur $\Rr^+$ vers la fonction $f~:~x\mapsto e^{-x}$.

\textbf{Convergence uniforme sur $\Rr^+$.}
Pour $x$ réel positif et $n$ entier naturel non nul, posons $g_n(x) = f(x)-f_n(x) =\left\{
\begin{array}{l}
e^{-x}-\left(1-\frac{x}{n}\right)^n\;\text{si}\;x\in[0,n]\\
e^{-x}\;\text{si}\;x>n
\end{array}
\right.$. Déterminons la borne supérieure de la fonction $|g_n|$ sur $[0,+\infty[$.

La fonction $g_n$ est définie et continue sur $R^+$. Pour $x\geqslant n$, $0<g_n(x)\leqslant e^{-n} = g_n(n)$.

Etudions la fonction $g_n$ sur $[0,n]$. Pour $x\in[0,n]$, $g_n'(x)=-e^{-x}+\left(1-\frac{x}{n}\right)^{n-1}$.($g_n'(n)$ est la dérivée à gauche de la fonction $g_n$ en $n$, mais on peut montrer qu'en fait la fonction $g_n$ est dérivable en $n$ pour $n > 1$).

La fonction $g_n$ est continue sur le segment $[0,n]$ et admet donc sur $[0,n]$ un minimum et un maximum.

\textbullet~La fonction $g_n$ a un minimum égal à $0$ atteint en $0$. En effet, on sait que pour tout réel $u$, $e^u\geqslant 1+u$ (inégalité de convexité) et donc pour tout réel $x$ de $[0,n]$,  $e^{-x/n}\geqslant 1-\frac{x}{n}\geqslant 0$. Après élévation des deux membres de cette inégalité, par croissance de $t\mapsto t^n$ sur $\Rr^+$, on obtient $e^{-x}\geqslant\left(1-\frac{x}{n}\right)^n$ ou encore $g_n(x)\geqslant 0 = g_n(0)$.

\textbullet~Pour $0 < x\leqslant n$, les inégalités précédentes sont strictes et la fonction ${g_n}_{/[0,n]}$ admet son maximum dans $]0,n]$. De plus, $g_n'(n)=-e^{-n}< 0$ et puisque la fonction $g_n$ est de classe $C^1$ sur $[0,n]$, sa dérivée $g_n'$ est strictement négative sur un voisinage à gauche de $n$. La fonction $g_n$ est alors strictement décroissante sur ce voisinage et la fonction $g_n$ admet nécessairement son maximum sur $\Rr^+$ en un certain point $x_n$ de $]0,n[$. En un tel point, puisque l'intervalle $]0,n[$ est ouvert, on sait que la dérivée de la fonction $g_n$ s'annule. L'égalité $g_n'(x_n) = 0$ fournit $\left(1-\frac{x_n}{n}\right)^{n-1}=e^{-x_n}$ et donc 

\begin{center}
$g_n(x_n) =e^{-x_n}-\left(1-\frac{x_n}{n}\right)^n =\left(1-\left(1-\frac{x_n}{n}\right)\right)e^{-x_n}=  \frac{x_ne^{-x_n}}{n}$.
\end{center}

En résumé, pour tout réel positif $x$, $0\leqslant g_n(x)\leqslant\frac{x_ne^{-x_n}}{n}$ où $x_n$ est un certain réel de $]0,n[$.

Pour $u$ réel positif, posons $h(u) = ue^{-u}$. La fonction $h$ est dérivable sur $/mbr^+$ et pour $u\geqslant0$, $h'(u) = (1-u)e^{-u}$. Par suite, la fonction $h$ admet un maximum en $1$ égal à $\frac{1}{e}$. On a montré que

\begin{center}
$\forall x\in[0,+\infty[$, $\forall n\in\Nn^*$, $0\leqslant g_n(x)\leqslant\frac{1}{ne}$
\end{center}

ou encore
$\forall n\in\Nn^*$, $\text{sup}\{|g_n(x)|,\;x\geqslant 0\}\leqslant\frac{1}{ne}$. Ainsi, $\lim_{n \rightarrow +\infty}\text{sup}\{|g_n(x)|,\;x\geqslant 0\}=0$ et on a montré que
 
 
\begin{center}
\shadowbox{
la suite de fonctions $(f_n)_{n\in\Nn^*}$ converge uniformément sur $\Rr^+$ vers la fonction $x\mapsto e^{-x}$.
}
\end{center}

\item  \textbf{Existence de $I=\int_{0}^{+\infty}e^{-x^2}\;dx$.} La fonction $x\mapsto e^{-x^2}$ est continue sur $[0,+\infty[$ et négligeable devant $\frac{1}{x^2}$ en $+\infty$. Donc la fonction $x\mapsto e^{-x^2}$  est intégrable sur $[0,+\infty[$. Par suite, $I$ existe dans $\Rr$.

On est alors en droit d'espérer que $I=\lim_{n \rightarrow +\infty}\int_{0}^{+\infty}f_n(x^2)\;dx$.

La fonction $x\mapsto f_n(x^2)$ est continue sur $[0,+\infty[$ et nulle sur $[\sqrt{n},+\infty[$. Donc la fonction $x\mapsto f_n(x^2)$ est intégrable sur $[0,+\infty[$. Pour $n\in\Nn^*$, posons $I_n=\int_{0}^{+\infty}f_n(x^2)\;dx=\int_{0}^{\sqrt{n}}\left(1-\frac{x^2}{n}\right)^n\;dx$.

Montrons que $I_n$ tend vers $I$ quand $n$ tend vers $+\infty$.

\begin{center}
$|I-I_n|\leqslant\int_{0}^{\sqrt{n}}|f(x^2)-f_n(x^2)|\;dx+\int_{\sqrt{n}}^{+\infty}e^{-x^2}\;dx\leqslant\sqrt{n}\times\frac{1}{ne}+\int_{\sqrt{n}}^{+\infty}e^{-x^2}\;dx=\frac{1}{e\sqrt{n}}+\int_{\sqrt{n}}^{+\infty}e^{-x^2}\;dx$.
\end{center}

Puisque la fonction $x\mapsto e^{-x^2}$ est intégrable sur $[0,+\infty[$, cette dernière expression tend vers $0$ quand $n$ tend vers $+\infty$ et donc $\lim_{n \rightarrow +\infty}I_n=I$.

\textbf{Calcul de la limite de $I_n$.} Soit $n\in\Nn^*$. Les changements de variables $x = u\sqrt{n}$  puis $u =\cos v$ fournissent

\begin{center}
$I_n =\int_{0}^{\sqrt{n}}\left(1-\frac{x^2}{n}\right)^n\;dx=\sqrt{n}\int_{0}^{1}(1-u^2)^n\;du =\sqrt{n}\int_{0}^{\pi/2}\sin^{2n+1}v\;dv =\sqrt{n}W_{2n+1}$
\end{center}

où $W_n$ est la $n$-ème intégrale de \textsc{Wallis}. On a déjà vu (exercice classique, voir fiches de Maths Sup) que $W_n\underset{n\rightarrow+\infty}{\sim}\sqrt{\frac{\pi}{2n}}$ et donc 

\begin{center}
$I_n\underset{n\rightarrow+\infty}{\sim}\sqrt{n}\times\sqrt{\frac{\pi}{2(2n+1)}}\underset{n\rightarrow+\infty}{\sim}\frac{\sqrt{\pi}}{2}$.
\end{center}

Finalement, $I_n$ tend vers $\frac{\sqrt{\pi}}{2}$ quand $n$ tend vers $+\infty$ et donc 

\begin{center}
\shadowbox{
$\int_{0}^{+\infty}e^{-x^2}\;dx=\frac{\sqrt{\pi}}{2}$.
}
\end{center}

Vous pouvez voir différents calculs de l'intégrale de \textsc{Gauss} dans \og Grands classiques de concours : intégration \fg.
\fincorrection
\correction{005729}
Posons $f=\lim_{n \rightarrow +\infty}P_n$.

Le critère de \textsc{Cauchy} de convergence uniforme (appliqué à $\varepsilon=1$) permet d'écrire

\begin{center}
$\exists N\in\Nn/\;\forall n\geqslant N,\;\forall m\geqslant N,\;\forall x\in\Rr,\; |P_n(x) - P_m(x)|\leqslant1$.
\end{center}

Pour $n\geqslant N$, les polynômes $P_N-P_n$  sont bornés sur $\Rr$ et donc constants. Par suite, pour chaque $n\geqslant N$, il existe $a_n\in\Rr$ tel que $P_N-P_n = a_n$\quad$(*)$.
Puisque la suite $(P_n)$ converge simplement sur $\Rr$, La suite $(a_n)=(P_N(0) - P_n(0))$ converge  vers un réel que l'on note $a$. On fait alors tendre $n$ tend vers $+\infty$ dans l'égalité $(*)$ et on obtient 

\begin{center}
$f = P_N - a$
\end{center}

On a montré que $f$ est un polynôme.

\fincorrection
\correction{005732}
\begin{enumerate}
 \item  \textbf{Convergence simple.} Chaque fonction $f_n$, $n\in\Nn$, est définie sur $\Rr$. Soit $x\in\Rr$.

\textbullet~Si $x<0$, $f_n(x)\underset{n\rightarrow+\infty}{\rightarrow}+\infty$ et la série de terme général $f_n(x)$, $n\in\Nn$, diverge grossièrement.

\textbullet~Si $x=0$, puisque $\forall n\in\Nn$, $f_n(x)=f_n(0)=0$, la série de terme général $f_n(x)$, $n\in\Nn$, converge.

\textbullet~Si $x>0$, $n^2f_n(x)=x^2e^{-x\sqrt{n}+3\ln n}\underset{n\rightarrow+\infty}{\rightarrow}0$ et donc $f_n(x)\underset{n\rightarrow+\infty}{=}o\left(\frac{1}{n^2}\right)$. Dans ce cas aussi, la série de terme général $f_n(x)$, $n\in\Nn$, converge.

\begin{center}
\shadowbox{
La série de fonctions de terme général $f_n$, $n\in\Nn$, converge simplement sur $\Rr^+$.
}
\end{center}

\textbf{Convergence normale.} La fonction $f_0$ est la fonction nulle. Soit $n\in\Nn^*$. La fonction $f_n$ est dérivable sur $\Rr^+$ et pour tout réel positif $x$,

\begin{center}
$f_n'(x)=n(2x-x^2\sqrt{n})e^{-x\sqrt{n}}=nx(2-x\sqrt{n})e^{-x\sqrt{n}}$.
\end{center}

La fonction $f_n$ est positive sur $[0,+\infty[$, croissante sur $\left[0,\frac{2}{\sqrt{n}}\right]$ et décroissante sur $\left[\frac{2}{\sqrt{n}},+\infty\right[$. On en déduit que

\begin{center}
$\|f_n\|_\infty=\underset{x\in[0,+\infty[}{\text{sup}}|f_n(t)|=f_n\left(\frac{2}{\sqrt{n}}\right)=4e^{-2}$.
\end{center}

Par suite, la série numérique de terme général $\|f_n\|_\infty$, $n\in\Nn$, diverge grossièrement et donc

\begin{center}
\shadowbox{
La série de fonctions de terme général $f_n$, $n\in\Nn$, ne converge pas normalement sur $\Rr^+$.
}
\end{center}

Soit $a>0$. Pour $n\geqslant\frac{4}{a^2}$, on a $\frac{2}{\sqrt{n}}\leqslant a$ et donc la fonction $f_n$ est décroissante sur $[a,+\infty[$. Soit donc $n$ un entier supérieur ou égal à $\frac{4}{a^2}$. Pour tout réel $t$ supérieur ou égal à $a$, on a $|f_n(t)|=f_n(t)\leqslant f_n(a)$ et donc $\underset{x\in[a,+\infty[}{\text{sup}}|f_n(t)|=f_n(a)$.

Comme la série numérique de terme général $f_n(a)$, $n\in\Nn$, converge, la série de fonctions de terme général $f_n$, $n\in\Nn$, converge normalement et donc uniformément sur $[a,+\infty[$.

\begin{center}
\shadowbox{
Pour tout $a>0$, la série de fonctions de terme général $f_n$, $n\in\Nn$, converge normalement et uniformément sur $[a,+\infty[$.
}
\end{center}

\textbf{Convergence uniforme sur $[0,+\infty[$.} Pour $n\in\Nn$ et $t\in\Rr^+$,

\begin{center}
$|R_n(t)|=\sum_{k=n+1}^{+\infty}f_k(t)\geqslant f_{n+1}(t)$,
\end{center}

et donc $\underset{t\in[0,+\infty[}{\text{sup}}|R_n(t)|\geqslant\underset{t\in[0,+\infty[}{\text{sup}}|f_{n+1}(t)|4e^{-2}$. Par suite, $\underset{t\in[0,+\infty[}{\text{sup}}|R_n(t)|$ ne tend pas vers $0$ quand $n$ tend vers $+\infty$ et donc

\begin{center}
\shadowbox{
la série de fonctions de terme général $f_n$, $n\in\Nn$, ne converge pas uniformément sur $\Rr^+$.
}
\end{center}

\item  \textbf{Convergence simple.} Chaque fonction $f_n$, $n\in\Nn^*$, est définie sur $]0,+\infty[$. Soit $x\in]0,+\infty[$. Puisque $f_n(x)\underset{n\rightarrow+\infty}{\sim}\frac{1}{n^3x^2}>0$, la série numérique de terme général $f_n(x)$ converge. Donc

\begin{center}
\shadowbox{
la série de fonctions de terme général $f_n$, $n\in\Nn^*$, converge simplement sur $]0,+\infty[$.
}
\end{center}

\textbf{Convergence normale.} Soit $n\in\Nn^*$. La fonction $f_n$ est décroissante et positive sur $]0,+\infty[$. Donc $\underset{x\in]0,+\infty}{\text{sup}}|f_n(x)|=f_n(0)=\frac{1}{n}$.Puisque la série numérique de terme général $\frac{1}{n}$, $n\in\Nn^*$, diverge

\begin{center}
\shadowbox{
la série de fonctions de terme général $f_n$, $n\in\Nn^*$, ne converge pas normalement sur $\Rr^+$.
}
\end{center}

Soit $a>0$. Pour $n\in\Nn^*$, la fonction $f_n$ est décroissante et positive sur $5a,+\infty[$ et donc $\underset{x\in[a,+\infty}{\text{sup}}|f_n(x)|=f_n(a)$.

Comme la série numérique de terme général $f_n(a)$, $n\in\Nn^*$, converge, la série de fonctions de terme général $f_n$, $n\in\Nn$, converge normalement et donc uniformément sur $[a,+\infty[$.

\begin{center}
\shadowbox{
Pour tout $a>0$, la série de fonctions de terme général $f_n$, $n\in\Nn^*$, converge normalement et uniformément sur $[a,+\infty[$.
}
\end{center}

\item  \textbf{Convergence simple.} Chaque fonction $f_n$, $n\in\Nn$, est définie sur $\Rr$ et impaire. Soit $x\in\Rr^+$.

\textbullet~Si $x=0$, pour tout entier naturel $n$, $f_n(x)=f_n(0)=0$. Dans ce cas, la série numérique de terme général $f_n(x)$ converge.

\textbullet~Si $x>0$, la suite $\left(\frac{x}{(x^2+1)^n}\right)_{n\in\Nn}$ est une suite géométrique de premier $x>0$ et de raison $\frac{1}{x^2+1}\in]0,1[$. On en déduit que la suite $\left(\frac{x}{(x^2+1)^n}\right)_{n\in\Nn}$ est positive décroissante de limite nulle. Par suite, la série numérique de terme général $f_n(x)$ converge en vertu du critère spécial aux séries alternées.

\textbullet~Si $x<0$, puisque pour tout entier naturel $n$, $f_n(x)=-f_n(-x)$, la série numérique de terme général $f_n(x)$ converge.

Finalement

\begin{center}
\shadowbox{
la série de fonctions de terme général $f_n$, $n\in\Nn$, converge simplement sur $\Rr$.
}
\end{center}

\textbf{Convergence normale.} La fonction $f_0$ n'est pas bornée sur $\Rr$ et donc la série de fonctions de terme général $f_n$, $n\in\Nn$, n'est pas normalement convergente sur $\Rr$.

Analysons la convergence normale de la série de fonctions de terme général $f_n$, $n\geqslant1$, sur $\Rr$.

Soit $n\in\Nn^*$. La fonction $g_n=(-1)^nf_n$ est dérivable sur $\Rr$ et pour tout réel $x$,

\begin{center}
$g_n'(x)=\frac{1}{(1+x^2)^n}+x\times\frac{-2nx}{(1+x^2)^{n+1}}=\frac{1-(2n-1)x^2}{(1+x^2)^{n+1}}$.
\end{center}

La fonction $g_n$ est positive sur $\Rr^+$, croissante sur $\left[0,\frac{1}{\sqrt{2n-1}}\right]$ et décroissante sur $\left[\frac{1}{\sqrt{2n-1}},+\infty\right[$. Puisque la fonction $g_n$ est impaire, on en déduit que

\begin{center}
$\|f_n\|_\infty=\underset{x\in\Rr}{\text{sup}}|f_n(x)|=g_n\left(\frac{1}{\sqrt{2n-1}}\right)=\frac{1}{\sqrt{2n-1}}\times\frac{1}{\left(1+\frac{1}{2n-1}\right)^{n+1}}=\frac{1}{\sqrt{2n-1}}\left(1-\frac{1}{2n}\right)^{-(n+1)}$.
\end{center}

Mais $\left(1-\frac{1}{2n}\right)^{-(n+1)}=\text{exp}\left(-(n+1)\ln\left(1-\frac{1}{2n}\right)\right)\underset{n\rightarrow+\infty}{=}\text{exp}\left(\frac{1}{2}+o(1)\right)$ et donc 

\begin{center}
$\|f_n\|_\infty=\frac{1}{\sqrt{2n-1}}\left(1-\frac{1}{2n}\right)^{-(n+1)}\underset{n\rightarrow+\infty}{\sim}\frac{1}{e\sqrt{2}\times\sqrt{n}}>0$.
\end{center}

Par suite, la série numérique de terme général $\|f_n\|_\infty$, $n\in\Nn^*$, diverge et donc

\begin{center}
\shadowbox{
la série de fonctions de terme général $f_n$, $n\in\Nn^*$, ne converge pas normalement sur $\Rr$.
}
\end{center}

\textbf{Convergence uniforme sur $\Rr$.} Soit $n\in\Nn$. Pour $x\in\Rr^+$, puisque la suite $\left(\frac{x}{(1+x^2)^n}\right)_{n\in\Nn}$ est positive décroissante et de limite nulle, d'après une majoration classique du reste à l'ordre $n$ d'une série alternée,

\begin{center}
$R_n(x)|=\left|\sum_{k=n+1}^{+\infty}(-1)^k\frac{x}{(1+x^2)^k}\right|\leqslant\left|(-1)^{n+1}\frac{x}{(1+x^2)^{n+1}}\right|=\frac{x}{(1+x^2)^{n+1}}=g_{n+1}(x)\leqslant g_{n+1}\left(\frac{1}{\sqrt{2n+1}}\right)$,
\end{center}

cette inégalité restant valable pour $x<0$ par parité. Donc $\underset{x\in\Rr}{\text{sup}}|R_n(x)|\leqslant g_{n+1}\left(\frac{1}{\sqrt{2n+1}}\right)$. D'après ci-dessus, 

$g_{n+1}\left(\frac{1}{\sqrt{2n+1}}\right)$ tend vers $0$ quand $n$ tend vers $+\infty$ et il en est de même de $\underset{x\in\Rr}{\text{sup}}|R_n(x)|$. On a montré que

\begin{center}
\shadowbox{
la série de fonctions de terme général $f_n$, $n\in\Nn$, converge uniformément sur $\Rr$.
}
\end{center}
\end{enumerate}
\fincorrection
\correction{005733}
Soit $n\in\Nn$.

\begin{center}
$\sum_{k=0}^{n}\frac{(-1)^k}{1+ka}=\sum_{k=0}^{n}(-1)^k\int_{0}^{1}t^{ka}\;dt=\int_{0}^{1}\left(\sum_{k=0}^{n}(-t^a)^k\right)\;dt=\int_{0}^{1}\frac{1}{1+t^a}\;dt+(-1)^n\int_{0}^{1}\frac{t^{(n+1)a}}{1+t^a}\;dt$,
\end{center}

avec $\left|(-1)^n\int_{0}^{1}\frac{t^{(n+1)a}}{1+t^a}\;dt\right|\leqslant\int_{0}^{1}t^{(n+1)a}\;dt=\frac{1}{1+(n+1)a}$. Par suite, $\lim_{n \rightarrow +\infty}(-1)^n\int_{0}^{1}\frac{t^{(n+1)a}}{1+t^a}\;dt=0$. On en déduit que la série de terme général $\frac{(-1)^k}{1+ka}$, $k\geqslant 0$, converge et que

\begin{center}
\shadowbox{
$\sum_{k=0}^{+\infty}\frac{(-1)^k}{1+ka}=\int_{0}^{1}\frac{1}{1+t^a}\;dt$.
}
\end{center}
\fincorrection
\nocorrection
\nocorrection
\nocorrection
\correction{004518}
\begin{enumerate}
  \item 
  \item oui, par convolution.
\end{enumerate}
 \fincorrection
\nocorrection
\nocorrection
\correction{004521}
$|g_n(f_n(x))-g(f(x))| \le |g_n(f_n(x))-g(f_n(x))|
                                  + |g(f_n(x))-g(f(x))|$
         et $g$ est uniformément continue.
\fincorrection
\nocorrection
\correction{004527}
Prendre une subdivision régulière de $[a,b]$ et encadrer $f_n$ par
les cordes associées.
\fincorrection
\correction{005728}
\begin{enumerate}
 \item 
 \begin{enumerate}
 \item Soit $n\in\Nn^*$.

\textbullet~Si $\forall x\in[0,1]$, $f(x)= 1$,

\begin{center}
$B_n(f)=\sum_{k=0}^{n}\dbinom{n}{k}X^k(1-X)^{n-k}= (X+(1-X))^n = 1$.
\end{center}

\textbullet~Si $\forall x\in[0,1]$, $f(x)=x$,

\begin{align*}\ensuremath
B_n(f)&=\sum_{k=0}^{n}\frac{k}{n}\dbinom{n}{k}X^k(1-X)^{n-k}=\sum_{k=1}^{n}\dbinom{n-1}{k-1}X^k(1-X)^{n-k}=X\sum_{k=1}^{n}\dbinom{n-1}{k-1}X^{k-1}(1-X)^{(n-1)-(k-1)}\\
 &= X\sum_{k=0}^{n-1}\dbinom{n-1}{k}X^k(1-X)^{n-1-k}= X.
\end{align*}
	

\textbullet~Si $\forall x\in[0,1]$, $f(x)=x(x-1)$, alors $B_n(f)=\sum_{k=0}^{n}\dbinom{n}{k}\frac{k}{n}\left(\frac{k}{n}-1\right)X^k(1-X)^{n-k}$ et donc $B_1(f)=0$. Pour $n\geqslant2$ et $k\in\llbracket1,n-1\rrbracket$

\begin{center}
$\frac{k}{n}\left(\frac{k}{n}-1\right)\dbinom{n}{k}=-\frac{1}{n^2}k(n-k)\frac{n!}{k!(n-k)!}=-\frac{n-1}{n}\frac{(n-2)!}{(k-1)(n-k-1)!}=-\frac{n-1}{n}\dbinom{n-2}{k-1}$.
\end{center}

Par suite,

\begin{align*}\ensuremath
B_n(f)&=-\frac{n-1}{n}\sum_{k=1}^{n-1}\dbinom{n-2}{k-1}X^k(1-X)^{n-k}  = -\frac{n-1}{n}X(1-X)   \sum_{k=1}^{n-1}X^{k-1}(1-X)^{(n-2)-(k-1)}\\
 &=-\frac{n-1}{n}X(1-X)\sum_{k=0}^{n-2}\dbinom{n-2}{k}X^k(1-X)^{n-2-k}=-\frac{n-1}{n} X(1-X).
\end{align*}

ce qui reste vrai pour n = 1.

 \item D'après la question précédente

\begin{align*}\ensuremath
\sum_{k=0}^{n}\dbinom{n}{k}(k-nX)^2X^k(1-X)^{n-k}&=\sum_{k=0}^{n}\dbinom{n}{k}k^2X^k(1-X)^{n-k}- 2nX\sum_{k=0}^{n}\dbinom{n}{k}kX^k(1-X)^{n-k}+n^2X^2\sum_{k=0}^{n}\dbinom{n}{k}X^k(1-X)^{n-k}\\
 &=\sum_{k=0}^{n}\dbinom{n}{k}k(k-n)X^k(1-X)^{n-k}-n(2X-1)\sum_{k=0}^{n}\dbinom{n}{k}kX^k(1-X)^{n-k}\\
  &+n^2X^2\sum_{k=0}^{n}\dbinom{n}{k}X^k(1-X)^{n-k}\\
  &=n^2\sum_{k=0}^{n}\frac{k}{n}\left(\frac{k}{n}-1\right)\dbinom{n}{k}X^k(1-X)^{n-k}-n^2(2X-1)\sum_{k=0}^{n}\dbinom{n}{k}\frac{k}{n}X^k(1-X)^{n-k}+n^2X^2\\
 &= -n(n-1)X(1-X) -n^2(2X-1)X + n^2X^2 = -nX^2 + nX = nX(1-X).
\end{align*}
  \end{enumerate}

\item  Soit $\varepsilon> 0$. Soient $n$ un entier naturel non nul et $\alpha$ un réel strictement positif donné. Soit $x$ un réel de $[0,1]$.

Notons $A$ (resp. $B$) l'ensemble des entiers $k\in\llbracket0,n\rrbracket$ tels que $\left|x-\frac{k}{n}\right|<\alpha$ (resp. $\left|x-\frac{k}{n}\right|\geqslant\alpha$). (Si $A$ ou $B$ sont vides, les sommes ci-dessous correspondantes sont nulles).

\begin{align*}\ensuremath
|f(x) - B_n(f)(x)|&=\left|\sum_{k=0}^{n}\dbinom{n}{k}\left(f(x)-f\left(\frac{k}{n}\right)\right)x^k(1-x)^{n-k}\right|\\
 &\leqslant\sum_{k\in A}^{}\dbinom{n}{k}\left|f(x)-f\left(\frac{k}{n}\right)\right|x^k(1-x)^{n-k}+\sum_{k\in B}^{}\dbinom{n}{k}\left|f(x)-f\left(\frac{k}{n}\right)\right|x^k(1-x)^{n-k}
\end{align*}

$f$ est continue sur le segment $[0,1]$ et donc est uniformément continue sur ce segment d'après le théorème de \textsc{Heine}. Par suite, il existe $\alpha>0$ tel que si $x$ et $y$ sont deux réels de $[0,1]$ tels que $|x-y|<\alpha$ alors $|f(x)-f(y)| <\frac{\varepsilon}{2}$. $\alpha$ est ainsi dorénavant fixé. Pour ce choix de $\alpha$, 

\begin{center}
$\sum_{k\in A}^{}\dbinom{n}{k}\left|f(x)-f\left(\frac{k}{n}\right)\right|x^k(1-x)^{n-k}\leqslant \frac{\varepsilon}{2}\sum_{k\in A}^{}\dbinom{n}{k}x^k(1-x)^{n-k}\leqslant \frac{\varepsilon}{2}\sum_{k=0}^{n}\dbinom{n}{k}x^k(1-x)^{n-k} =\frac{\varepsilon}{2}.$
\end{center}

Ensuite, la fonction $f$ est continue sur le segment $[0,1]$ et donc est bornée sur ce segment. Soit $M$ un majorant de la fonction $|f|$ sur $[0,1]$.

\begin{center}
$\sum_{k\in B}^{}\dbinom{n}{k}\left|f(x)-f\left(\frac{k}{n}\right)\right|x^k(1-x)^{n-k}\leqslant2M\sum_{k\in B}^{}\dbinom{n}{k}x^k(1-x)^{n-k}$
\end{center}

Mais si $k\in B$, l'inégalité $\left|x-\frac{k}{n}\right|\geqslant\alpha$ fournit $1\leqslant\frac{1}{\alpha^2n^2}(k-nx)^2$ et donc

\begin{align*}\ensuremath
\sum_{k\in B}^{}\dbinom{n}{k}x^k(1-x)^{n-k}&\leqslant 1\leqslant\frac{1}{\alpha^2n^2}\sum_{k\in B}^{}\dbinom{n}{k}(k-nx)^2x^k(1-x)^{n-k}\leqslant\frac{1}{\alpha^2n^2}\sum_{k=0}^{n}\dbinom{n}{k}(k-nx)^2x^k(1-x)^{n-k}\\
 &=\frac{1}{\alpha^2n^2}\times nx(1-x)=\frac{1}{\alpha^2n}\left(\frac{1}{4}-\left(x-\frac{1}{2}\right)^2\right)\leqslant\frac{1}{4\alpha^2n}.
\end{align*}

En résumé, pour tout réel $x\in[0,1]$

\begin{center}
$|f(x) - B_n(f)(x)|\leqslant\frac{\varepsilon}{2}+2M\times\frac{1}{4\alpha^2n}=\frac{\varepsilon}{2}+\frac{M}{2\alpha^2n}$.
\end{center}

Maintenant, puisque $\lim_{n \rightarrow +\infty}\frac{M}{2\alpha^2n}=0$, il existe un entier naturel non nul $N$ tel que pour $n\geqslant N$, $\frac{M}{2\alpha^2n}<\frac{\varepsilon}{2}$. Pour $n\geqslant N$, on a $|f(x) - B_n(f)(x)|<\frac{\varepsilon}{2}+\frac{\varepsilon}{2}=\varepsilon$. On a montré que

\begin{center}
$\forall\varepsilon>0$, $\exists N\in\Nn^*/\;\forall n\in\Nn^*,\;\forall x\in[0,1],\;(n\geqslant N\Rightarrow|f(x)-(B_n(f))(x)|<\varepsilon$,
\end{center}

et donc que

\begin{center}
\shadowbox{
la suite de polynômes $(B_n(f))_{n\in\Nn^*}$ converge uniformément sur $[0,1]$ vers $f$.
}
\end{center}

\item  La question 2) montre le théorème de \textsc{Weierstrass} dans le cas du segment $[0,1]$.

Soient $[a,b]$ un segment quelconque et $f$ un application continue sur $[a,b]$.

Pour $x\in[0,1]$, posons $g(x)= f(a+(b-a)x)$. La fonction $g$ est continue sur $[0,1]$ et donc il existe une suite de polynômes $(P_n)$ convergeant uniformément vers $g$ sur $[0,1]$.
Pour $n\in\Nn$, posons $Q_n =P_n\left(\frac{X-a}{b-a}\right)$.

Soit $\varepsilon>0$. $\exists N\geqslant1$ tel que $\forall n\geqslant N$, $\forall y\in[0,1]$, $|g(y)-P_n(y)| <\varepsilon$.

Soient $x\in[a,b]$ et $n\geqslant N$. Le réel $y=\frac{x-a}{b-a}$ est dans $[0,1]$ et

\begin{center}
$|f(x)-Q_n(x)|=\left|f(a+(b-a)y)-Q_n(a+(b-a)y)\right|=|g(y)-P_n(y)|<\varepsilon$.
\end{center}

Ceci démontre que la suite de polynômes $(Q_n)_{n\in\Nn}$ converge uniformément vers la fonction $f$ sur $[a,b]$.
\end{enumerate}
\fincorrection
\correction{005730}
\begin{enumerate}
 \item  Pour $x\in]-1,1[$ et $n$ entier naturel non nul, posons $f_n(x)=\frac{x^n\sin(nx)}{n}$. 

Soit $x\in]-1,1[$. Pour $n$ entier naturel non nul, $|f_n(x)|\leqslant|x|^n$. Or, la série géométrique de terme général $|x|^n$, $n\geqslant1$, est convergente et donc la série numérique de terme général $f_n(x)$ est absolument convergente et en particulier convergente. On en déduit que $f(x)$ existe.

\begin{center}
\shadowbox{
$f$ est définie sur $]-1,1[$.
}
\end{center}

Soit $a\in]0,1[$. Chaque $f_n$, $n\geqslant1$, est de classe $C^1$ sur $[-a,a]$ et pour $x\in[-a,a]$,

\begin{center}
$f_n'(x) =x^{n-1}\sin(nx)+x^n\cos(nx)$.
\end{center}

Pour $x\in[-a,a]$ et $n\in\Nn^*$, 

\begin{center}
$|f_n'(x)|\leqslant a^{n-1}+a^n\leqslant2a^{n-1}$.
\end{center} 

Puisque la série numérique de terme général $2a^{n-1}$, $n\geqslant1$, converge, la série de fonctions de terme général $f_n'$, $n\geqslant1$, est normalement et donc uniformément sur $[-a,a]$.

En résumé,

\textbullet~la série de fonctions de terme général $f_n$, $n\geqslant1$, converge simplement vers $f$ sur $[-a,a]$,

\textbullet~chaque fonction $f_n$, $n\geqslant1$, est de classe $C^1$ sur $[-a,a]$,

\textbullet~la série de fonctions de terme général $f_n'$ converge uniformément sur $[-a,a]$.

D'après un corollaire du théorème de dérivation terme à terme, $f$ est de classe $C^1$ sur $[-a,a]$ pour tout réel $a$ de $]0,1[$  et donc sur $]-1,1[$ et sa dérivée s'obtient par dérivation terme à terme. 

\begin{center}
\shadowbox{
$f$ est de classe $C^1$ sur $]-1,1[$ et $\forall x\in]-1,1[$, $f'(x)=\sum_{n=1}^{+\infty}(x^{n-1}\sin(nx)+x^n\cos(nx))$.
}
\end{center}

\item  Ainsi, pour $x\in]-1,1[$

\begin{align*}\ensuremath
f'(x)&=\sum_{n=1}^{+\infty}(x^{n-1}\sin(nx)+x^n\cos(nx))=\text{Im}\left(\sum_{n=1}^{+\infty}x^{n-1}e^{inx}\right)+\text{Re}\left(\sum_{n=1}^{+\infty}x^{n}e^{inx}\right)\\
 &=\text{Im}\left(\frac{e^{ix}}{1-xe^{ix}}\right)+\text{Re}\left(\frac{xe^{ix}}{1-xe^{ix}}\right)=\text{Im}\left(\frac{e^{ix}(1-xe^{-ix})}{x^2-2x\cos x+1}\right) +\text{Re}\left(\frac{xe^{ix}(1-xe^{-ix})}{x^2-2x\cos x+1}\right)\\
  &=\frac{\sin x+x\cos x-x^2}{x^2-2x\cos x+1}  .
\end{align*}

Mais, pour $x\in]-1,1[$,

\begin{center}
$\left(\frac{x\sin x}{1-x\cos x}\right)'=\frac{(\sin x+x\cos x)(1-x\cos x)-x\sin x(-\cos x+x\sin x)}{(1-x\cos x)^2}=\frac{\sin x+x\cos x-x^2}{(1-x\cos x)^2}$.
\end{center}

et donc

\begin{align*}\ensuremath
\left(\Arctan\left(\frac{x\sin x}{1-x\cos x}\right)\right)'&=\frac{\sin x+x\cos x-x^2}{(1-x\cos x)^2}\times\frac{1}{1+\left(\frac{x\sin x}{1-x\cos x}\right)^2}=\frac{\sin x+x\cos x-x^2}{(1-x\cos x)^2+x^2\sin^2x}\\
 &=\frac{\sin x+x\cos x-x^2}{x^2-2x\cos x+1}= f'(x).
\end{align*}

Finalement, pour $x\in]-1,1[$,

\begin{center}
$f(x)=f(0)+\int_{0}^{x}f'(t)\;dt=0+\Arctan\left(\frac{x\sin x}{1-x\cos x}\right)-\Arctan(0)=\Arctan\left(\frac{x\sin x}{1-x\cos x}\right)$.
\end{center}

\begin{center}
\shadowbox{
$\forall x\in]-1,1[$, $\sum_{n=1}^{+\infty}\frac{x^n\sin(nx)}{n}=\Arctan\left(\frac{x\sin x}{1-x\cos x}\right)$.
}
\end{center}
\end{enumerate}
\fincorrection
\correction{005731}
\begin{enumerate}
 \item  Pour $n$ entier naturel non nul, on note $f_n$ la fonction $x\mapsto\frac{(-1)^n}{\ln(nx)}$. Pour tout réel $x$, $f(x)$ existe si et seulement si chaque $f_n(x)$, $n\in\Nn^*$, existe et la série numérique de terme général $f_n(x)$, $n\in\Nn^*$, converge.

Pour $n\in\Nn^*$ et $x\in\Rr$, $f_n(x)$ existe si et seulement si $x>0$ et $x\neq \frac{1}{n}$.

Soit donc $x\in D =]0,+\infty[\setminus\left\{\frac{1}{p},\;p\in\Nn^*\right\}$.

Pour $n>\frac{1}{x}$, on a $\ln(nx)>0$. On en déduit que la suite $\left(\frac{1}{\ln(nx)}\right)_{n\in\Nn^*}$ est positive et décroissante à partir d'un certain et tend vers $0$ quand $n$ tend vers $+\infty$. Ainsi, la série numérique de terme général $f_n(x)$ converge en vertu du critère spécial aux séries alternées et donc $f(x)$ existe. 

\begin{center}
\shadowbox{
Le domaine de définition de $f$ est $D=]0,+\infty[\setminus\left\{\frac{1}{n},\;n\in\Nn^*\right\}$.
}
\end{center}

\item  \textbf{Limite de $f$ en $+\infty$.}
Soit $x > 1$. Donc $f(x)$ existe. Pour tout entier naturel non nul $n$, $\ln(nx) > 0$. On en déduit que la suite $\left(\frac{1}{\ln(nx)}\right)_{n\in\Nn^*}$ est décroissante. On sait alors que la valeur absolue de $f(x)$ est majorée par la valeur absolue du premier terme de la série. Ainsi

\begin{center}
$\forall x>1$, $|f(x)|\leqslant\left|\frac{(-1)^0}{\ln(x)}\right|=\frac{1}{\ln x}$,
\end{center}

et en particulier 

\begin{center}
\shadowbox{
$\lim_{x \rightarrow +\infty}f(x) = 0$.
}
\end{center}

On peut noter de plus que pour $x>1$, $f(x)$ est du signe du premier terme de la série à savoir $\frac{1}{\ln(x)}$ et donc $\forall x\in]1,+\infty[$, $f(x) > 0$.

\textbf{Convergence uniforme sur$]1,+\infty[$.} D'après une majoration classique du reste à l'ordre $n$ alternée d'une série alternée, pour $x > 1$ et $n$ naturel non nul,

\begin{center}
$|R_n(x)| =\left|\sum_{k=n+1}^{+\infty}\frac{(-1)^{k-1}}{\ln(kx)}\right|\leqslant\left|\frac{(-1)^{n}}{\ln((n+1)x)}\right|=\frac{1}{\ln((n+1)x}\leqslant\frac{1}{n+1}$.
\end{center}

Donc, pour tout entier naturel non nul, $\underset{x\in]1,+\infty[}{\text{sup}}|R_n(x)|\leqslant\frac{1}{\ln(n+1)}$ et donc $\lim_{n \rightarrow +\infty}\underset{x\in]1,+\infty[}{\text{sup}}|R_n(x)|=0$. La série de fonctions de terme général $f_n$ converge uniformément vers sa somme sur $]1,+\infty[$.

\textbf{Continuité sur $]1,+\infty[$.} Chaque fonction $f_n$, $n\in\Nn^*$ est continue sur $]1,+\infty[$ et donc $f$ est donc continue sur $]1,+\infty[$ en tant que limite uniforme sur $]1,+\infty[$ d'une suite de fonctions continues sur $]1,+\infty[$.

\begin{center}
\shadowbox{
$f$ est continue sur $]1,+\infty[$.
}
\end{center}

\textbf{Limite en $1$ à droite.} Soit $n\geqslant2$. Quand $x$ tend vers $1$ par valeurs supérieures, $f_n(x)$ tend vers $\ell_n=\frac{(-1)^{n-1}}{\ln(n)}$. Puisque la série de fonctions de terme général $f_n$, $n\geqslant2$, converge uniformément vers sa somme sur $]1,+\infty[$, le théorème d'interversion des limites permet d'affirmer que la série numérique de terme général $\ell_n$, $n\geqslant2$ converge et que la fonction $x\mapsto f(x) -\frac{1}{\ln(x)}=\sum_{n=2}^{+\infty}f_n(x)$ tend vers le réel $\sum_{n=2}^{+\infty}\frac{(-1)^{n-1}}{\ln(n)}$ quand $x$ tend vers $1$ par valeurs supérieures ou encore 

\begin{center}
\shadowbox{
$f(x)\underset{x\rightarrow1^+}{=}\frac{1}{\ln x}+O(1)$ et en particulier, $\displaystyle\lim_{\substack{x\rightarrow1\\
x>1}}f(x) = +\infty$.
}
\end{center}

\item  La série de fonctions de terme général $f_n$, $n\geqslant1$, converge simplement vers la fonction $f$ sur $]1,+\infty[$. De plus chaque fonction $f_n$ est de classe $C^1$ sur $]1,+\infty[$ et pour $n\in\Nn^*$ et $x>1$, 

\begin{center}
$f_n'(x) =\frac{(-1)^n}{x\ln^2(nx)}$.
\end{center}

Il reste à vérifier la convergence uniforme de la série de fonctions de terme général $f_n'$ sur $]1,+\infty[$.

Soit $x > 1$. La série de terme général $f_n'(x)$ est alternée car son terme général est alterné en signe et sa valeur absolue à savoir $\frac{1}{x\ln^2(nx)}$ tend vers zéro quand $n$ tend vers $+\infty$ en décroissant. Donc, d'après une majoration classique du reste à l'ordre $n$ d'une série alternée,

\begin{center}
$|R_n(x)|=\left|\sum_{k=n+1}^{+\infty}\frac{(-1)^k}{x\ln^2(kx)}\right|\leqslant\left|\frac{(-1)^{n+1}}{x\ln^2((n+1)x)}\right|=\frac{1}{x\ln^2((n+1)x)}\leqslant\frac{1}{\ln^2(n+1)}$.
\end{center}

Par suite, $\underset{x\in]1,+\infty[}{\text{sup}}|R_n(x)|\leqslant\frac{1}{\ln^2(n+1)}$ et donc $\lim_{n \rightarrow +\infty}\underset{x\in]1,+\infty[}{\text{sup}}|R_n(x)|=0$. Ainsi, la série de fonctions de terme général $f_n'$, $n\geqslant1$, converge uniformément sur $]1,+\infty[$.

En résumé,

\textbullet~la série de fonctions de terme général $f_n$, $n\geqslant1$, converge simplement vers $f$ sur $]1,+\infty[$,

\textbullet~chaque fonction $f_n$, $n\geqslant1$, est de classe $C^1$ sur $]1,+\infty[$,

\textbullet~la série de fonctions de terme général $f_n'$ converge uniformément sur $]1,+\infty[$.

D'après un corollaire du théorème de dérivation terme à terme, $f$ est de classe $C^1$ sur $]1,+\infty[$ et sa dérivée s'obtient par dérivation terme à terme. 

\begin{center}
\shadowbox{
$f$ est de classe $C^1$ sur $]1,+\infty[$ et $\forall x>1$, $f'(x) =\sum_{n=1}^{+\infty}\frac{(-1)^n}{x\ln^2(nx)}$.
}
\end{center}

Pour $x >1$, puisque la série de somme $f'(x)$ est alternée, $f'(x)$ est du signe du premier terme de la somme à savoir $-\frac{1}{x\ln^2x}$. Par suite, $\forall x\in]-1,1[$, $f'(x)\leqslant0$ et $f$ est donc strictement décroissante sur $]1,+\infty[$.

\begin{center}
\shadowbox{
La fonction $f$ est décroissante sur $]1,+\infty[$.
}
\end{center}
\end{enumerate}
\fincorrection
\correction{005734}
\begin{enumerate}
 \item  \textbf{Convergence simple.} Soit $t\in\Rr$. Pour tout entier naturel non nul $n$, $1+\frac{t^2}{n(1+t^2)}\geqslant1>0$ et donc $f_n(t)$ existe. Ensuite, $\ln\left(1+\frac{t^2}{n(1+t^2)}\right)>0$ et donc la suite numérique $(f_n(t))_{n\in\Nn^*}$ est alternée en signe. De plus, $|f_n(t)|=\ln\left(1+\frac{t^2}{n(1+t^2)}\right)$ et la suite $(|f_n(t)|)_{n\in\Nn^*}$ tend vers $0$ en décroissant.

On en déduit que la série de terme général $f_n(t)$, $n\geqslant1$, converge en vertu du critère spécial aux séries alternées.

\begin{center}
\shadowbox{
La série de fonctions de terme général $f_n$, $n\geqslant1$, converge simplement sur $\Rr$.
}
\end{center}

On pose alors $f=\sum_{n=1}^{+\infty}f_n$.

\textbf{Convergence uniforme.} Soit $n\in\Nn^*$. D'après une majoration classique du reste à l'ordre $n$ d'une série alternée, pour tout réel $t$ on a

\begin{align*}\ensuremath
|R_n(t)|&=\left|\sum_{k=n+1}^{+\infty}f_k(t)\right|\leqslant|f_{n+1}(t)|=\ln\left(1+\frac{t^2}{(n+1)(1+t^2)}\right)=\ln\left(1+\frac{t^2+1-1}{(n+1)(1+t^2)}\right)=\ln\left(1+\frac{1}{n+1}-\frac{1}{(n+1)(1+t^2)}\right)\\
 &\leqslant\ln\left(1+\frac{1}{n+1}\right),
\end{align*}

et donc, $\forall n\in\Nn^*$, $\underset{t\in\Rr}{\text{sup}}|R_n(t)|\leqslant\ln\left(1+\frac{1}{n+1}\right)$. Comme $\lim_{n \rightarrow +\infty}\ln\left(1+\frac{1}{n+1}\right)=0$, on a encore $\lim_{n \rightarrow +\infty}\underset{t\in\Rr}{\text{sup}}|R_n(t)|=0$ et on a montré que

\begin{center}
\shadowbox{
La série de fonctions de terme général $f_n$, $n\geqslant1$, converge uniformément vers $f$ sur $\Rr$.
}
\end{center}

\textbf{Continuité.} Puisque chaque fonction $f_n$, $n\geqslant1$, est continue sur $\Rr$, la fonction $f$ est continue sur $\Rr$ en tant que limite uniforme sur $\Rr$ d'une suite de fonctions continues sur $\Rr$.

\begin{center}
\shadowbox{
$f$ est continue sur $\Rr$.
}
\end{center}

\item  D'après le théorème d'interversion des limites, $f$ a une limite réelle en $+\infty$ et

\begin{center}
$\lim_{t \rightarrow +\infty}f(t)=\sum_{n=1}^{+\infty}\lim_{t \rightarrow +\infty}f_n(t)=\sum_{n=1}^{+\infty}(-1)^n\ln\left(1+\frac{1}{n}\right)=\ln\left(\frac{2}{\pi}\right)$ (voir l'exercice \ref{ex:rou4}, 5)).
\end{center}

\begin{center}
\shadowbox{
$\lim_{t \rightarrow +\infty}f_n(t)=\ln\left(\frac{2}{\pi}\right)$.
}
\end{center}
\end{enumerate}
\fincorrection
\correction{005735}
\textbf{Domaine de définition.} Soit $t\in\Rr$. Pour chaque $n\in\Nn^*$, $f_n(t)$ existe et de plus $f_n(t)=\frac{\Arctan(nt)}{n^2}\underset{n\rightarrow+\infty}{=}O\left(\frac{1}{n^2}\right)$. Donc la série numérique de terme général $f_n(t)$, $n\geqslant1$, converge absolument et en particulier converge. On a montré que

\begin{center}
\shadowbox{
$f$ est définie sur $\Rr$.
}
\end{center}

\textbf{Parité.} Pour tout réel $t$,

\begin{center}
$f(-t)=\sum_{n=1}^{+\infty}\frac{\Arctan(-nt)}{n^2}=-\sum_{n=1}^{+\infty}\frac{\Arctan(nt)}{n^2}=-f(t)$.
\end{center}

\begin{center}
\shadowbox{
$f$ est impaire.
}
\end{center}

\textbf{Convergence normale.} Pour tout réel $t$ et tout entier naturel non nul $n$, $|f_n(t)|\leqslant\frac{\pi}{2n^2}$ et donc pour tout entier naturel non nul $n$,

\begin{center}
$\underset{t\in\Rr}{\text{sup}}|f_n(t)|\leqslant\frac{\pi}{2n^2}$.
\end{center}

\rule[-4mm]{0mm}{0mm}Comme la série numérique de terme général $\frac{\pi}{2n^2}$, $n\geqslant1$, converge, la série de fonctions de terme général $f_n$ converge normalement et donc uniformément vers $f$ sur $\Rr$.

\textbf{Limite de $f$ en $+\infty$.} Puisque la série de fonctions de terme général $f_n$, $n\geqslant1$, converge uniformément vers $f$ sur $\Rr$ et que chaque fonction $f_n$ a une limite réelle quand $t$ tend vers $+\infty$ à savoir $\ell_n=\frac{\pi}{2n^2}$, le théorème d'interversion des limites permet d'affirmer que $f$ a une limite réelle en $+\infty$ et que

\begin{center}
$\lim_{t \rightarrow +\infty}f(t)=\sum_{n=1}^{+\infty}\ell_n=\frac{\pi}{2}\sum_{n=1}^{+\infty}\frac{1}{n^2}=\frac{\pi^3}{12}$.
\end{center}

\begin{center}
\shadowbox{
$\lim_{t \rightarrow +\infty}f(t)=\frac{\pi^3}{12}$ et $\lim_{t \rightarrow -\infty}f(t)=-\frac{\pi^3}{12}$.
}
\end{center}

\textbf{Continuité.} Puisque chaque fonction $f_n$, $n\in\Nn^*$, est continue sur $\Rr$ et que la série de fonctions de terme général $f_n$ converge uniformément vers $f$ sur $\Rr$, la fonction $f$ est continue sur $\Rr$ en tant que limite uniforme sur $\Rr$ d'une suite de fonctions continues sur $\Rr$.

\begin{center}
\shadowbox{
$f$ est continue sur $\Rr$.
}
\end{center}

\textbf{Dérivation.} Soit $a>0$. Chaque fonction $f_n$, $n\geqslant1$, est de classe $C^1$ sur $[a,+\infty[$ et pour $n\in\Nn^*$ et $t\geqslant a$,

\begin{center}
$f_n'(t)=\frac{n}{n^2(1+n^2t^2)}=\frac{1}{n(1+n^2t^2)}$.
\end{center}

Pour $n\in\Nn^*$, on a alors $\underset{t\in[a,+\infty[}{\text{sup}}|f_n'(t)|=f_n'(a)=\frac{1}{n(1+n^2a^2)}$. Puisque $\frac{1}{n(1+n^2a^2)}\underset{n\rightarrow+\infty}{\sim}\frac{1}{a^2n^3}>0$, la série de terme général $\frac{1}{n(1+n^2a^2)}$ converge et par suite, la série de fonctions de terme général $f_n'$, $n\geqslant1$, converge normalement et donc uniformément sur $[a,+\infty[$.

En résumé,

\textbullet~la série de fonctions de terme général $f_n$, $n\geqslant 1$, converge simplement vers $f$ sur $[a,+\infty[$,

\textbullet~chaque fonction $f_n$ est de classe $C^1$ sur $[a,+\infty[$,

\textbullet~la série de fonctions de terme général $f_n'$ converge uniformément sur $[a,+\infty[$.

D'après un corollaire du théorème de dérivation terme à terme, $f$ est de classe $C^1$ sur $[a,+\infty[$ et sa dérivée s'obtient par dérivation terme à terme. Ceci étant vrai pour tout $a>0$, $f$ est de classe $C^1$ sur $]0,+\infty[$ et puisque $f$ est impaire

\begin{center}
\shadowbox{
$f$ est de classe $C^1$ sur $\Rr^*$ et $\forall t\in\Rr^*$, $f'(t)=\sum_{n=1}^{+\infty}\frac{1}{n(1+n^2t^2)}$.
}
\end{center}

\textbf{Dérivabilité en $0$.} La fonction $f'$ est décroissante sur $]0,+\infty[$. Donc la fonction $f'$ admet une limite en $0^+$ élément de $]-\infty,+\infty]$. Pour $t>0$ et $N\in\Nn^*$, on a $f'(t)\geqslant\sum_{n=1}^{N}\frac{1}{n(1+n^2t^2)}$ et quand $t$ tend vers $0$, on obtient

\begin{center}
$\displaystyle\lim_{\substack{t\rightarrow0\\ t>0}}f'(t)\geqslant\sum_{n=1}^{N}\frac{1}{n}$.
\end{center}

Cette inégalité étant vraie pour tout entier naturel non nul $N$, quand $N$ tend vers $+\infty$ on obtient

\begin{center}
$\displaystyle\lim_{\substack{t\rightarrow0\\ t>0}}f'(t)\geqslant\sum_{n=1}^{+\infty}\frac{1}{n}=+\infty$.
\end{center}

On a montré que $\displaystyle\lim_{\substack{t\rightarrow0\\ t>0}}f'(t)=+\infty$.

En résumé, $f$ est de classe $C^0$ sur $[0,+\infty[$, de classe $C^1$ sur $]0,+\infty[$ et $f'(t)$ tend vers $+\infty$ quand $t$ tend vers $0$ par valeurs supérieures. D'après un corollaire du théorème des accroissements finis, on sait que $f$ n'est pas dérivable en $0$ à droite et que sa courbe représentative admet $[Oy)$ pour demi-tangente en $(0,0)$. Puisque $f$ est impaire, $f$ n'est pas dérivable en $0$ et sa courbe représentative admet $(Oy)$ pour tangente en $(0,0)$.

\textbf{Allure du graphe.}

$$\includegraphics{../images/img005735-1}$$

\fincorrection
\correction{005738}
\begin{enumerate}
 \item  Soit $x\in[0,+\infty[$. Pour $n> x^2$, $f_n(x)=\text{exp}\left(n\ln\left(1-\frac{x^2}{n}\right)\right)$ et donc $f_n(x)\underset{n\rightarrow+\infty}{=}\text{exp}(-x^2+o(1))$. Donc la suite $(f_n)_{n\in\Nn^*}$ converge simplement sur $\Rr^+$ vers la fonction $f~:~x\mapsto e^{-x^2}$.

\item  Chaque fonction $f_n$, $n\in\Nn^*$, est continue par morceaux sur $[0,+\infty[$ et nulle au voisinage de $+\infty$. Donc chaque fonction $f_n$, $n\in\Nn^*$, est intégrable sur $[0,+\infty[$.

La fonction $f$ est continue sur $[0,+\infty[$ et négligeable devant $\frac{1}{x^2}$ quand $x$ tend vers $+\infty$. Donc la fonction $f$ est intégrable sur $[0,+\infty[$.

Soit $n\in\Nn^*$. Par convexité de la fonction exponentielle, $\forall u\in\Rr$, $1+u\leqslant e^u$. Par suite, $\forall x\in\left[0,\sqrt{n}\right]$, $0\leqslant1-\frac{x^2}{n}\leqslant e^{-x^2/n}$ puis par croissance de la fonction $t\mapsto t^n$ sur $\Rr^+$, $0\leqslant f_n(x)=\left(1-\frac{x^2}{n}\right)^n\leqslant e^{-x^2}=f(x)$. D'autre part, pour $x>\sqrt{n}$, $f_n(x)=0\leqslant f(x)$. Finalement 

\begin{center}.
$\forall n\in\Nn^*$, $\forall x\in[0,+\infty[$, $|f_n(x)|\leqslant f(x)$.
\end{center}

En résumé,

\textbullet~chaque fonction $f_n$, $n\in\Nn^*$, est continue par morceaux et intégrable sur $[0,+\infty[$,

\textbullet~la suite de fonctions $(f_n)$ converge simplement vers la fonction $f$ sur $[0,+\infty[$ et la fonction $f$ est continue sur $[0,+\infty[$.

\textbullet~$\forall n\in\Nn^*$, $|f_n|\leqslant f$, la fonction $f$ étant intégrable sur $[0,+\infty[$.

D'après le théorème de convergence dominée, la suite $\left(\int_{0}^{+\infty}f_n(x)\;dx\right)_{n\in\Nn^*}$ converge vers $\int_{0}^{+\infty}f(x)\;dx$. Ainsi,

\begin{center}
$\int_{0}^{+\infty}e^{-x^2}\;dx=\lim_{n \rightarrow +\infty}\int_{0}^{\sqrt{n}}\left(1-\frac{x^2}{n}\right)^n\;dx$.
\end{center}

Soit $n\in\Nn^*$. En posant $t=\Arccos\left(\frac{x}{\sqrt{n}}\right)$ et donc $\frac{x^2}{n}=\cos^2t$ et $dx=-\sqrt{n}\sin t\;dt$, on obtient

\begin{center}
$\int_{0}^{\sqrt{n}}\left(1-\frac{x^2}{n}\right)^n\;dx=\int_{\pi/2}^{0}(1-\cos^2t)^n\times(-\sqrt{n}\sin t)\;dt=\sqrt{n}\int_{0}^{\pi/2}\sin^{2n+1}t\;dt=\sqrt{n}W_{2n+1}$,
\end{center}

où $W_n$ est la $n$-ème intégrale de \textsc{Wallis}. Classiquement, $W_n\underset{n\rightarrow+\infty}{\sim}\sqrt{\frac{\pi}{2n}}$ (voir Exercices Maths Sup) et donc

\begin{center}
$\frac{W_{2n+1}}{\sqrt{n}}\underset{n\rightarrow+\infty}{\sim}\sqrt{n}\sqrt{\frac{\pi}{2(2n+1)}}\underset{n\rightarrow+\infty}{\sim}\frac{\sqrt{\pi}}{2}$.
\end{center}

On a montré que

\begin{center}
\shadowbox{
$\int_{0}^{+\infty}e^{-x^2}\;dx=\frac{\sqrt{\pi}}{2}$.
}
\end{center}
\end{enumerate}
\fincorrection
\correction{005739}
Pour $x\in]0,1]$, $x^{-x}=e^{-x\ln(x)}$ et donc $\lim_{x \rightarrow 0^+}x^{-x}=1$. Donc si on pose $\forall x\in[0,1]$, $f(x)=\left\{
\begin{array}{l}
x^{-x}\;\text{si}\;x\in]0,1]\\
1\;\text{si}\;x=0
\end{array}
\right.$, $f$ est une fonction continue sur le segment $[0,1]$ et donc intégrable sur le segment $[0,1]$.

Pour $x\in]0,1]$, $x^{-x}=e^{-x\ln(x)}=\sum_{n=0}^{+\infty}\frac{(-x\ln(x))^n}{n!}$. Posons alors $\forall x\in[0,1]$, $f_0(x)=1$ puis $\forall n\in\Nn^*$,  $\forall x\in[0,1]$, $f_n(x)=\left\{
\begin{array}{l}
\frac{(-x\ln(x))^n}{n!}\;\text{si}\;x\in]0,1]\\
\rule{0mm}{4mm}0\;\text{si}\;x=0
\end{array}
\right.$. La fonction $f_0$ est continue sur $[0,1]$ et pour $n\in\Nn^*$, puisque $-x\ln(x)\underset{x\rightarrow0^+}{\rightarrow}0$, la fonction $f_n$ est continue sur $[0,1]$. En résumé, chaque fonction $f_n$, $n\in\Nn$, est continue sur $[0,1]$. De plus,

\begin{center}
$\forall x\in[0,1]$, $f(x)=\sum_{n=0}^{+\infty}f_n(x)$.
\end{center}

Vérifions alors que la série de fonctions de terme général $f_n$ converge normalement et donc uniformément vers $f$ sur le segment $[0,1]$. Pour $x\in[0,1]$, posons $g(x)=\left\{
\begin{array}{l}
-x\ln x\;\text{si}\;x\in]0,1]\\
0\;\text{si}\;x=0
\end{array}
\right.$. La fonction $g$ est continue sur le segment $[0,1]$ et admet donc un maximum $M$ sur ce segment. Pour $x\in[0,1]$, on a $0\leqslant g(x)\leqslant M$ (on peut montrer que $M=g\left(\frac{1}{e}\right)=\frac{1}{e}$). Mais alors $\forall n\in\Nn$, $\forall x\in]0,1]$, $|f_n(x)|=\frac{(g(x))^n}{n!}\leqslant\frac{M^n}{n!}$ ce qui reste vrai pour $x=0$. Comme la série numérique de terme général $\frac{M^n}{n!}$ converge, on a montré que la série de fonctions de terme général $f_n$ converge normalement et donc uniformément vers $f$ sur le segment $[0,1]$.

D'après le théorème d'intégration terme à terme sur un segment, la série numérique de terme général $\int_{0}^{1}f_n(x)\;dx$, converge et

\begin{center}
$\int_{0}^{1}f(x)\;dx=\sum_{n=0}^{+\infty}\int_{0}^{1}f_n(x)\;dx$\quad$(*)$.
\end{center}

Pour $n\in\Nn$, posons $I_n=\int_{0}^{1}f_n(x)\;dx$. Soit $n\in\Nn$. En posant $u=-\ln(x)$ puis $v=(n+1)u$, on obtient

\begin{align*}\ensuremath
I_n&=\frac{1}{n!}\int_{0}^{1}(-x\ln x)^n\;dx=\frac{1}{n!}\int_{+\infty}^{0}(ue^{-u})^n\times(-e^{-u}\;du)=\frac{1}{n!}\int_{0}^{+\infty}u^ne^{-(n+1)u}\;du\\
 &=\frac{1}{n!(n+1)^{n+1}}\int_{0}^{+\infty}v^ne^{-v}\;dv=\frac{\Gamma(n+1)}{n!(n+1)^{n+1}}=\frac{1}{(n+1)^{n+1}}.
\end{align*}

L'égalité $(*)$ s'écrit alors $\int_{0}^{1}x^{-x}\;dx=\sum_{n=0}^{+\infty}\frac{1}{(n+1)^{n+1}}=\sum_{n=1}^{+\infty}\frac{1}{n^n}$.

\textbf{Remarque.} Pour calculer $I_n=\int_{0}^{1}\frac{(-x\ln x)^n}{n!}\;dx$, on peut aussi s'intéresser plus généralement à $J_{n,p}=\int_{0}^{1}\frac{x^n(-\ln x)^p}{n!}\;dx$ que l'on calcule par récurrence grâce à une intégration par parties.

Le travail qui précède permet encore d'écrire

\begin{center}
$\int_{0}^{1}x^x\;dx=\sum_{n=0}^{+\infty}\int_{0}^{1}\frac{(x\ln x)^n}{n!}\;dx=\sum_{n=1}^{+\infty}\frac{(-1)^n}{n^n}$.
\end{center}

\begin{center}
\shadowbox{
$\int_{0}^{1}x^{-x}\;dx =\sum_{n=1}^{+\infty}\frac{1}{n^n}$ et $\int_{0}^{1}x^{x}\;dx =\sum_{n=1}^{+\infty}\frac{(-1)^n}{n^n}$.
}
\end{center}
\fincorrection
\correction{005740}
Pour $x>0$, posons $f(x)=\frac{x^2}{e^x-1}$. $f$ est continue sur $]0,+\infty[$.

Ensuite, pour tout réel strictement positif $x$, on a $0<e^{-x}<1$ et donc

\begin{center}
$\frac{x^2}{e^x-1}=\frac{x^2e^{-x}}{1-e^{-x}}=x^2e^{-x}\sum_{n=0}^{+\infty}e^{-nx}=\sum_{n=0}^{+\infty}x^2e^{-(n+1)x}=\sum_{n=1}^{+\infty}x^2e^{-nx}$.
\end{center}

Pour $n\in\Nn^*$ et $x>0$, posons $f_n(x)=x^2e^{-nx}$. Chaque fonction $f_n$, $n\in\Nn^*$, est continue et intégrable sur $[0,+\infty[$ car négligeable devant $\frac{1}{x^2}$ quand $x$ tend vers $+\infty$. En particulier, chaque fonction $f_n$, $n\in\Nn^*$, est intégrable sur $]0,+\infty[$. De plus, pour $n\in\Nn^*$,

\begin{center}
$\int_{0}^{+\infty}|f_n(x)|\;dx=\int_{0}^{+\infty}x^2e^{-nx}\;dx=\frac{1}{n^3}\int_{0}^{+\infty}u^2e^{-u}\;du=\frac{\Gamma(3)}{n^3}=\frac{2}{n^3}$,
\end{center}

qui est le terme général d'une série numérique convergente.

En résumé,

\textbullet~chaque fonction $f_n$, $n\in\Nn^*$, est continue et intégrable sur $]0,+\infty[$,

\textbullet~la série de fonctions de terme général $f_n$, $n\in\Nn^*$, converge simplement vers la fonction $f$ sur $]0,+\infty[$ et la fonction $f$

est continue sur $]0,+\infty[$.

\textbullet~$\sum_{n=1}^{+\infty}\int_{0}^{+\infty}|f_n(x)|\;dx<+\infty$.

D'après un théorème d'intégration terme à terme, $f$ est intégrable sur $]0,+\infty[$, la série numérique de terme général $\int_{0}^{1}f_n(x)\;dx$, $n\in\Nn^*$, converge et

\begin{center}
$\int_{0}^{+\infty}f(x)\;dx=\sum_{n=1}^{+\infty}\int_{0}^{+\infty}f_n(x)\;dx=\sum_{n=1}^{+\infty}\frac{2}{n^3}$.
\end{center}

On a montré que

\begin{center}
\shadowbox{
$\int_{0}^{+\infty}\frac{x^2}{e^x-1}\;dx=2\sum_{n=1}^{+\infty}\frac{1}{n^3}$.
}
\end{center}
\fincorrection
\correction{005741}
C'est presque le même exercice que l'exercice \ref{ex:rou3bis}. Pour tout réel $x>0$,

\begin{center}
$\frac{x}{\sh x}=\frac{2xe^{-x}}{1-e^{-2x}}=2xe^{-x}\sum_{n=0}^{+\infty}e^{-2nx}=\sum_{n=0}^{+\infty}2xe^{-(2n+1)x}$,
\end{center}

puis avec la même démarche que dans l'exercice précédent

\begin{align*}\ensuremath
\int_{0}^{+\infty}\frac{x}{\sh x}\;dx&=\sum_{n=0}^{+\infty}\int_{0}^{+\infty}2xe^{-(2n+1)x}\;dx=\sum_{n=0}^{+\infty}\frac{2}{(2n+1)^2}\int_{0}^{+\infty}ue^{-u}\;du=\sum_{n=0}^{+\infty}\frac{2\Gamma(2)}{(2n+1)^2}\\
 &=2\sum_{n=0}^{+\infty}\frac{1}{(2n+1)^2}=2\left(\sum_{n=1}^{+\infty}\frac{1}{n^2}-\sum_{n=1}^{+\infty}\frac{1}{(2n)^2}\right)=2\left(1-\frac{1}{4}\right)\frac{\pi^2}{6}=\frac{\pi^2}{4}.
\end{align*}

\begin{center}
\shadowbox{
$\int_{0}^{+\infty}\frac{x}{\sh x}\;dx=\frac{\pi^2}{4}$.
}
\end{center}
\fincorrection
\correction{005742}
Ici, le plus simple est peut-être de ne pas utiliser de théorème d'intégration terme à terme. La fonction $f~:~x\mapsto\frac{\ln x}{1+x^2}$ est continue sur $]0,1]$. De plus, quand $x$ tend vers $0$, $f(x)\underset{x\rightarrow0^+}{\sim}\ln x\underset{x\rightarrow0^+}{=}o\left(\frac{1}{\sqrt{x}}\right)$. On en déduit que $f$ est intégrable sur $]0,1]$.

Soit $n\in\Nn$.

\begin{center}
$\frac{\ln x}{1+x^2}=\sum_{k=0}^{n}(-1)^kx^{2k}\ln x+\frac{(-1)^{n+1}x^{2n+2}\ln x}{1+x^2}$.
\end{center}

Maintenant, chacune des fonctions $f_k~:~x\mapsto(-1)^kx^{2k}\ln x$, $0\leqslant k\leqslant n$, est intégrable sur $]0,1]$ car continue sur $]0,1]$ et négligeable devant $\frac{1}{\sqrt{x}}$ quand $x$ tend vers $0$. On en déduit encore que la fonction $g_n~:~x\mapsto\frac{(-1)^{n+1}x^{2n+2}\ln x}{1+x^2}$ est intégrable sur $]0,1]$ car $g_n=f-\sum_{k=0}^{n}f_k$. On a donc

\begin{center}
$\forall n\in\Nn$, $\int_{0}^{1}\frac{\ln x}{1+x^2}\;dx=\sum_{k=0}^{n}(-1)^k\int_{0}^{1}x^{2k}\ln x\;dx+\int_{0}^{1}\frac{(-1)^{n+1}x^{2n+2}\ln x}{1+x^2}\;dx$.
\end{center}

La fonction $h~:~x\mapsto\frac{x\ln x}{1+x^2}\;dx$ est continue sur $]0,1]$ et prolongeable par continuité en $0$. On en déduit que la fonction $h$ est bornée sur $]0,1]$. Soit $M$ un majorant de la fonction $|h|$ sur $]0,1]$. Pour tout entier naturel $n$, on a alors

\begin{center}
$\left|\int_{0}^{1}\frac{(-1)^{n+1}x^{2n+2}\ln x}{1+x^2}\;dx\right|\leqslant \int_{0}^{1}x^{2n}\left|\frac{x\ln x}{1+x^2}\right|\;dx\leqslant M\int_{0}^{1}x^{2n}\;dx=\frac{M}{2n+1}$.
\end{center}

En particulier, $\lim_{n \rightarrow +\infty}\int_{0}^{1}\frac{(-1)^{n+1}x^{2n+2}\ln x}{1+x^2}\;dx=0$. Ceci montre que la série numérique de terme général $(-1)^k\int_{0}^{1}x^{2k}\ln x\;dx$, $k\in\Nn$, converge et que

\begin{center}
$\int_{0}^{1}\frac{\ln x}{1+x^2}\;dx=\sum_{k=0}^{+\infty}(-1)^k\int_{0}^{1}x^{2k}\ln x\;dx$.
\end{center}

Soient $n\in\Nn$ et $\varepsilon\in]0,1[$. Les deux fonctions $x\mapsto\frac{x^{2n+1}}{2n+1}$ et $x\mapsto\ln x$ sont de classe $C^1$ sur le segment $[\varepsilon,1]$. On peut donc effectuer une intégration par parties et on obtient

\begin{center}
$\int_{\varepsilon}^{1}x^{2n}\ln x\;dx=\left[\frac{x^{2n+1}}{2n+1}\ln x\right]_\varepsilon^1-\frac{1}{2n+1}\int_{\varepsilon}^{1}x^{2n}\;dx=-\frac{\varepsilon^{2n+1}}{2n+1}\ln \varepsilon-\frac{1}{(2n+1)^2}(1-\varepsilon^{2n+1})$.
\end{center}

Quand $\varepsilon$ tend vers $0$, on obtient $\int_{0}^{1}x^{2n}\ln x\;dx=-\frac{1}{(2n+1)^2}$. Par suite,

\begin{center}
$\int_{0}^{1}\frac{\ln x}{1+x^2}\;dx=-\sum_{n=0}^{+\infty}\frac{(-1)^n}{2n+1}=-\frac{\pi}{4}$.
\end{center}

Vérifions maintenant l'intégrabilité de la fonction $f$ sur $]0,+\infty[$. La fonction $f$ est continue sur $]0,+\infty[$ et on sait déjà que $f$ est intégrable sur $]0,1]$. De plus, $x^{3/2}f(x)\underset{x\rightarrow+\infty}{\sim}\frac{\ln x}{\sqrt{x}}\underset{x\rightarrow+\infty}{\rightarrow}0$ et donc $f(x)\underset{x\rightarrow+\infty}{=}o\left(\frac{1}{x^{3/2}}\right)$. Ceci montre que la fonction $f$ est intégrable sur $[1,+\infty[$ et finalement sur $]0,+\infty[$.

Pour calculer $I=\int_{0}^{+\infty}\frac{\ln x}{1+x^2}\;dx$, la méthode précédente ne marche plus du tout  car pour $x>1$, $x^{n}$ tend vers $+\infty$ quand $n$ tend vers $+\infty$. C'est une toute autre idée qui permet d'aller au bout. On pose $u=\frac{1}{x}$ et on obtient

\begin{center}
$I=\int_{0}^{+\infty}\frac{\ln x}{1+x^2}\;dx=\int_{+\infty}^{0}\frac{\ln\left(\frac{1}{u}\right)}{1+\frac{1}{u^2}}\times\frac{-du}{u^2}=-\int_{0}^{+\infty}\frac{\ln u}{1+u^2}\;du=-I$,
\end{center}

et donc $I=0$.

\begin{center}
\shadowbox{
$\int_{0}^{1}\frac{\ln x}{1+x^2}\;dx=-\frac{\pi}{4}$ et $\int_{0}^{+\infty}\frac{\ln x}{1+x^2}\;dx=0$.
}
\end{center}
\fincorrection
\correction{005743}
\begin{enumerate}
 \item  Soit $x\in[0,1[$. Pour tout réel $t$ de $[0,x]$, on a $\frac{1}{1-t}=\sum_{n=0}^{+\infty}t^n$. Maintenant, pour tout réel $t\in[0,x]$ et tout entier naturel $n$, on a $|t|^n\leqslant x^n$. Puisque la série numérique de terme général $x^n$ converge, on en déduit que la série de fonctions de terme général $t\mapsto t^n$ converge normalement et donc uniformément sur le segment $[0,x]$. D'après le théorème d'intégration terme à terme sur un segment, on peut affirmer que

\begin{center}
$-\ln(1-x)=\int_{0}^{x}\frac{1}{1-t}\;dt=\sum_{n=0}^{+\infty}\int_{0}^{x}t^n\;dt=\sum_{n=0}^{+\infty}\frac{x^{n+1}}{n+1}=\sum_{n=1}^{+\infty}\frac{x^n}{n}$.
\end{center}

\begin{center}
\shadowbox{
$\forall t \in[0,1[$, $-\ln(1-t) =\sum_{n=1}^{+\infty}\frac{t^n}{n}$.
}
\end{center}

\item  Par suite, pour $t\in]0,1[$,

\begin{center}
$\frac{\ln(t)\ln(1-t)}{t}=-\sum_{n=1}^{+\infty}\frac{t^{n-1}\ln t}{n}$.
\end{center}

Pour $t\in]0,1[$, posons $f(t)=\frac{\ln(t)\ln(1-t)}{t}$ puis pour $t\in]0,1]$ et $n\in\Nn^*$, posons $f_n(t)=-\frac{t^{n-1}\ln t}{n}$.

Soit $n\in\Nn^*$. La fonction $f_n$ est continue sur $]0,1]$ et négligeable devant $\frac{1}{\sqrt{t}}$ quand $t$ tend vers $0$. La fonction $f_n$ est donc intégrable sur $]0,1]$. En particulier, la fonction $f_n$ est donc intégrable sur $]0,1[$. Calculons alors $\int_{0}^{1}f_n(t)\;dt$.

Soit $a\in]0,1[$. Les deux fonctions $t\mapsto\frac{t^n}{n}$ et $t\mapsto-\ln t$ sont de classe $C^1$ sur le segment $[a,1]$. On peut donc effectuer une intégration par parties et on obtient

\begin{center}
$\int_{a}^{1}t^{n-1}(-\ln t)\;dt=\left[-\frac{t^n\ln t}{n}\right]_a^1+\frac{1}{n}\int_{a}^{1}t^{n-1}\;dt=\frac{a^n\ln a}{n}+\frac{1}{n^2}(1-a^n)$.
\end{center}

Quand $a$ tend vers $0$, on obtient $\int_{0}^{1}-t^{n-1}\ln t\;dt=\frac{1}{n^2}$ et donc $\int_{0}^{1}f_n(t)\;dt=\frac{1}{n^3}$. Puisque la fonction $f_n$ est positive sur $]0,1[$, on a encore $\int_{0}^{1}|f_n(t)|\;dt=\frac{1}{n^3}$. On en déduit que la série numérique de terme général $\int_{0}^{1}|f_n(t)|\;dt$ converge.

En résumé,

\textbullet~chaque fonction $f_n$ est continue par morceaux et intégrable sur $]0,1[$,

\textbullet~la séries de fonctions de terme général $f_n$, $n\in\Nn^*$, converge simplement vers la fonction $f$ sur $]0,1[$ et la fonction $f$

est continue sur $]0,1[$,

\textbullet~$\sum_{n=1}^{+\infty}\int_{0}^{1}|f_n(t)|\;dt<+\infty$.

D'après un théorème d'intégration terme à terme,

\begin{center}
$\int_{0}^{1}\frac{\ln(t)\ln(1-t)}{t}\;dt=\sum_{n=1}^{+\infty}\int_{0}^{1}\frac{-t^{n-1}\ln t}{n}\;dt=\sum_{n=1}^{+\infty}\frac{1}{n^3}$.
\end{center}

\begin{center}
\shadowbox{
$\int_{0}^{1}\frac{\ln(t)\ln(1-t)}{t}\;dt =\sum_{n=1}^{+\infty}\frac{1}{n^3}$.
}
\end{center}
\end{enumerate}
\fincorrection
\correction{005744}
\textbf{Existence de l'intégrale.} Soit $x\in\Rr$. La fonction $f~:~t\mapsto\frac{\cos(xt)}{\ch t}$ est continue sur $[0,+\infty[$. De plus, pour tout réel positif $t$, $|f(t)|\leqslant\frac{1}{\ch t}$ et donc $|f(t)|\underset{t\rightarrow+\infty}{\sim}\frac{2}{e^t}\underset{t\rightarrow+\infty}{=}o\left(\frac{1}{t^2}\right)$. On en déduit que la fonction $f$ est intégrable sur  $[0,+\infty[$.

\begin{center}
Pour tout réel $x$, $\int_{0}^{+\infty}\frac{\cos(xt)}{\ch t}\;dt$ existe.
\end{center}

\textbf{Convergence de la série.} Soit $x\in\Rr$. Pour $n\in\Nn$, posons $u_n(x)=\frac{2n+1}{(2n+1)^2+x^2}$. Pour $n\in\Nn$,

\begin{align*}\ensuremath
u_n(x)-u_{n+1}(x)&=\frac{2n+1}{(2n+1)^2+x^2}-\frac{2n+3}{(2n+3)^2+x^2}=\frac{(2n+1)((2n+3)^2+x^2)-(2n+3)((2n+1)^2+x^2)}{((2n+1)^2+x^2)((2n+3)^2+x^2)}\\
 &=\frac{2(2n+1)(2n+3)-2x^2}{((2n+1)^2+x^2)((2n+3)^2+x^2)}.
\end{align*}

Puisque le numérateur de cette dernière expression tend vers $+\infty$ quand $n$ tend vers $+\infty$, cette expression est positive pour $n$ grand. On en déduit que la suite $(u_n(x))$ décroît à partir d'un certain rang. D'autre part, $\lim_{n \rightarrow +\infty}u_n(x)=0$.

On en déduit que la série de terme général $(-1)^nu_n(x)$ converge en vertu du critère spécial aux séries alternées.

\begin{center}
Pour tout réel $x$, la série de terme général $(-1)^n\frac{2n+1}{(2n+1)^2+x^2}$ converge.
\end{center}

\textbf{Egalité de l'intégrale et de la somme de la série.} Soit $n\in\Nn$. Pour $t\in]0,+\infty[$, on a $e^{-t}\in]0,1[$ et donc

\begin{align*}\ensuremath
\frac{\cos(xt)}{\ch t}&=\frac{2\cos(xt)e^{-t}}{1+e^{-2t}}=2\cos(xt)e^{-t}\sum_{k=0}^{n}(-1)^ke^{-2kt}+(-1)^{n+1}\frac{\cos(xt)e^{-(2n+3)t}}{1+e^{-2t}}\\
 &=2\sum_{k=0}^{n}(-1)^k\cos(xt)e^{-(2k+1)t}+(-1)^{n+1}\frac{\cos(xt)e^{-(2n+3)t}}{1+e^{-2t}}.
\end{align*}

Maintenant, pour chaque $k\in\Nn$, la fonction $t\mapsto\cos(xt)e^{-(2k+1)t}$ est intégrable sur $[0,+\infty[$ car continue sur $[0,+\infty[$ et négligeable devant $\frac{1}{t^2}$ quand $t$ tend vers $+\infty$. On en déduit encore que la fonction $t\mapsto (-1)^{n+1}\frac{\cos(xt)e^{-(2n+3)t}}{1+e^{-2t}}$ est intégrable sur $[0,+\infty[$ puis que 

\begin{center}
$\forall n\in\Nn$, $\int_{0}^{+\infty}\frac{\cos(xt)}{\ch t}\;dt=2\sum_{k=0}^{n}(-1)^k\int_{0}^{+\infty}\cos(xt)e^{-(2k+1)t}\;dt+(-1)^{n+1}\int_{0}^{+\infty}\frac{\cos(xt)e^{-(2n+3)t}}{1+e^{-2t}}\;dt$.
\end{center}

Ensuite, $\left|\int_{0}^{+\infty}(-1)^{n+1}\frac{\cos(xt)e^{-(2n+3)t}}{1+e^{-2t}}\;dt\right|\leqslant\int_{0}^{+\infty}e^{-(2n+3)t}\;dt=\frac{1}{2n+3}$, et donc $\lim_{n \rightarrow +\infty}(-1)^{n+1}\frac{\cos(xt)e^{-(2n+3)t}}{1+e^{-2t}}\;dt=0$ puis

\begin{center}
$\int_{0}^{+\infty}\frac{\cos(xt)}{\ch t}\;dt=2\sum_{n=0}^{+\infty}(-1)^n\int_{0}^{+\infty}\cos(xt)e^{-(2n+1)t}\;dt$.
\end{center}

Soit $n\in\Nn$.

\begin{align*}\ensuremath
\int_{0}^{+\infty}\cos(xt)e^{-(2n+1)t}\;dt&=\text{Re}\left(\int_{0}^{+\infty}e^{ixt}e^{-(2n+1)t}\;dt\right)=\text{Re}\left(\int_{0}^{+\infty}e^{(-(2n+1)+ix)t}\;dt\right)\\
 &=\text{Re}\left(\left[\frac{e^{(-(2n+1)+ix)t}}{-(2n+1)+ix}\right]_0^{+\infty}\right)=\text{Re}\left(\frac{1}{(2n+1)-ix}\left(1-\lim_{t \rightarrow +\infty}e^{(-(2n+1)+ix)t}\right)\right)\\
 &=\text{Re}\left(\frac{1}{(2n+1)-ix}\right)\;(\text{car}\;\left|e^{(-(2n+1)+ix)t}\right|=e^{-(2n+1)t}\underset{n\rightarrow+\infty}{\rightarrow}0)\\
 &=\text{Re}\left(\frac{2n+1+ix}{(2n+1)^2+x^2}\right)=\frac{2n+1}{(2n+1)^2+x^2}.
\end{align*}

On a enfin montré que

\begin{center}
\shadowbox{
$\int_{0}^{+\infty}\frac{\cos(xt)}{\ch t}\;dt = 2\sum_{n=0}^{+\infty}(-1)^n\frac{2n+1}{(2n+1)^2+x^2}$.
}
\end{center}

\fincorrection
\correction{005864}
Soit $a\in\Rr$. Pour $n\in\Nn^*$, on pose $A_n=\left(
\begin{array}{cc}
1&- \frac{a}{n}\\
 \frac{a}{n}&1
\end{array}
\right)$.

Soit $n\in\Nn^*$. On peut écrire $A_n=\sqrt{1+ \frac{a^2}{n^2}}\left(
\begin{array}{cc}
 \frac{1}{\sqrt{1+ \frac{a^2}{n^2}}}&- \frac{a/n}{\sqrt{1+ \frac{a^2}{n^2}}}\\
 \frac{a/n}{\sqrt{1+ \frac{a^2}{n^2}}}& \frac{1}{\sqrt{1+ \frac{a^2}{n^2}}}
\end{array}
\right)$. Les sommes des carrés des deux nombres $ \frac{1}{\sqrt{1+ \frac{a^2}{n^2}}}$ et $ \frac{a/n}{\sqrt{1+ \frac{a^2}{n^2}}}$ est égale à $1$. Donc il existe un réel $\theta_n\in]-\pi,\pi]$ tel que $\cos(\theta_n)= \frac{1}{\sqrt{1+ \frac{a^2}{n^2}}}$ et $\sin(\theta_n)= \frac{a/n}{\sqrt{1+ \frac{a^2}{n^2}}}$. De plus, $\cos(\theta_n)>0$ et $\sin(\theta_n)>0$ et donc on peut prendre 

\begin{center}
$\theta_n=\Arctan\left( \frac{a}{n}\right)\in\left]0, \frac{\pi}{2}\right[$.
\end{center}

Pour $n\in\Nn^*$, on a alors

\begin{center}
$A_n^n=\left(\sqrt{1+ \frac{a^2}{n^2}}\right)^n\left(
\begin{array}{cc}
\cos(\theta_n)&-\sin(\theta_n)\\
\sin(\theta_n)&\cos(\theta_n)
\end{array}
\right)^n=\left(1+ \frac{a^2}{n^2}\right)^{n/2}\left(
\begin{array}{cc}
\cos(n\theta_n)&-\sin(n\theta_n)\\
\sin(n\theta_n)&\cos(n\theta_n)
\end{array}
\right)$.
\end{center}

Maintenant, $\left(1+ \frac{a^2}{n^2}\right)^{n/2}=\text{exp}\left( \frac{n}{2}\ln\left(1+ \frac{a^2}{n^2}\right)\right)\underset{n\rightarrow+\infty}{=}\text{exp}\left( \frac{a^2}{2n}+o\left( \frac{1}{n}\right)\right)\underset{n\rightarrow+\infty}{=}1+o(1)$.

D'autre part, $n\theta_n=n\Arctan\left( \frac{a}{n}\right)\underset{n\rightarrow+\infty}{=}n\times \frac{a}{n}=a$. Donc

\begin{center}
$\lim_{n \rightarrow +\infty}A_n^n=1.\left(
\begin{array}{cc}
\cos(a)&-\sin(a)\\
\sin(a)&\cos(a)
\end{array}
\right)=\left(
\begin{array}{cc}
\cos(a)&-\sin(a)\\
\sin(a)&\cos(a)
\end{array}
\right)$.
\end{center}

\begin{center}
\shadowbox{
$\forall a>0$, $\lim_{n \rightarrow +\infty}\left(
\begin{array}{cc}
1&- \frac{a}{n}\\
 \frac{a}{n}&1
\end{array}
\right)^n=\left(
\begin{array}{cc}
\cos(a)&-\sin(a)\\
\sin(a)&\cos(a)
\end{array}
\right)$.
}
\end{center}
\fincorrection
\correction{005865}
Soit $A\in\mathcal{M}_p(\Cc)$. 

\textbf{(3) $\Rightarrow$ (2).} On sait que si la série de terme général $A^n$, $n\in\Nn$, converge, alors $\lim_{n \rightarrow +\infty}A^n=0$.

\textbf{(2) $\Rightarrow$ (1).} Supposons $\lim_{n \rightarrow +\infty}A^n=0$. Soit $\lambda\in\Cc$ une valeur propre de $A$ et $X\in\mathcal{M}_{p,1}(\Cc)\setminus\{0\}$ un vecteur propre associé. Pour tout entier naturel $n$, $A^nX=\lambda^nX$. Puisque $\lim_{n \rightarrow +\infty}A^n=0$, on a encore $\lim_{n \rightarrow +\infty}A^nX=0$ puis $\lim_{n \rightarrow +\infty}\lambda^nX=0$ et donc $\lim_{n \rightarrow +\infty}\lambda^n=0$.

Ainsi, si $\lim_{n \rightarrow +\infty}A^n=0$ alors $\text{Sp}(A)\subset B_o(0,1)$.

\textbf{(1) $\Rightarrow$ (3)}. Soit $A\in\mathcal{M}_p(\Cc)$ telle que $\text{Sp}(A)\subset B_o(0,1)$. On sait (voir exercice \ref{ex:rou20} : décomposition de \textsc{Dunford}) qu'il existe deux matrices $D$ et $N$ telles que 

1) $A=D+N$ 

2) $D$ diagonalisable

3) $N$ nilpotente

4) $DN=ND$. 

De plus, les valeurs propres de $D$ sont les valeurs propres de $A$.

On note $k$ l'indice de nilpotence de $N$. Puisque les matrices $D$ et $N$ convergent, la formule du binôme de \textsc{Newton} permet d'écrire pour $n\geqslant k$

\begin{center}
$A^n=(D+N)^n=\sum_{j=0}^{n}\dbinom{n}{j}D^{n-j}N^j=\sum_{j=0}^{k}\dbinom{n}{j}D^{n-j}N^j$.
\end{center}

Il existe une matrice $P\in\mathcal{GL}_p(\Cc)$ et une matrice diagonale $\Delta$ tel que $D=P\Delta P^{-1}$. Mais alors, $\forall j\in\llbracket0,k\rrbracket$, $\forall n\geqslant j$, $\dbinom{n}{j}D^{n-j}N^j=P\times\dbinom{n}{j}\Delta^{n-j}\times PN^j$.

Soit $j\in\llbracket0,k\rrbracket$. Vérifions tout d'abord que la série de terme général $\dbinom{n}{j}\Delta^{n-j}$, $n\geqslant j$ converge. Posons $\Delta=\text{diag}(\lambda_1,\ldots,\lambda_p)$. Alors $\forall n\geqslant j$, $\dbinom{n}{j}\Delta^{n-j}=\text{diag}\left(\dbinom{n}{j}\lambda_1^{n-j},\ldots,\dbinom{n}{j}\lambda_p^{n-j}\right)$. Maintenant, si $\lambda$ est une valeur propre de $\Delta$ (et donc de $A$), $\dbinom{n}{j}\lambda^{n-j}= \frac{n(n-1)\ldots(n-j+1)}{j!}\lambda^{n-j}\underset{n\rightarrow+\infty}{\sim}n^j\lambda^{n-j}=\underset{n\rightarrow+\infty}o\left( \frac{1}{n^2}\right)$ car $|\lambda|<1$ et donc la série de terme général $\dbinom{n}{j}\lambda^{n-j}$, $n\geqslant j$, converge.

Ainsi, la série de terme général $\dbinom{n}{j}\Delta^{n-j}$ converge. D'autre part, l'application $M\mapsto P\times M\times PN^j$ est continue sur $\mathcal{M}_p(\Cc)$ en tant qu'endomorphisme d'un espace de dimension finie. On en déduit que la série de terme général $P\times\dbinom{n}{j}\Delta^{n-j}\times PN^j$ converge.

Finalement, pour chaque $j\in\llbracket0,k\rrbracket$, la série de terme général $P\times\dbinom{n}{j}\Delta^{n-j}\times PN^j$ converge et donc la série de terme général $A^n$ converge car est somme de $j+1$ séries convergentes.
\fincorrection
\correction{005866}
$\chi_A=\left|
\begin{array}{cc}
4/3-X&-5/6\\
5/3&-7/6-X
\end{array}
\right|=X^2- \frac{1}{6}X- \frac{1}{6}=\left(X- \frac{1}{2}\right)\left(X+ \frac{1}{3}\right)$. Par suite, $A=PDP^{-1}$ où $D=\text{diag}\left( \frac{1}{2},- \frac{1}{3}\right)$, $P=\left(
\begin{array}{cc}
1&1\\
1&2
\end{array}
\right)$ et donc $P^{-1}=\left(
\begin{array}{cc}
2&-1\\
-1&1
\end{array}
\right)$.

Soit $n\in\Nn$.

\begin{center}
$\sum_{k=0}^{n}A^k=P\left(\sum_{k=0}^{n}D^k\right)P^{-1}=P\;\text{diag}\left(\sum_{k=0}^{n}\left( \frac{1}{2}\right)^k,\sum_{k=0}^{n}\left(- \frac{1}{3}\right)^k\right)P^{-1}$.
\end{center}

Puisque $ \frac{1}{2}$ et $- \frac{1}{3}$ sont dans $]-1,1[$, les séries numériques de termes généraux respectifs $\left( \frac{1}{2}\right)^k$ et $\left(- \frac{1}{3}\right)^k$ convergent. Il en est de même de la série de terme général $D^k$. Maintenant, l'application $M\mapsto PMP^{-1}$, converge est continue car linéaire sur un espace de dimension finie et on en déduit que la série de terme général $A^k$ converge. De plus,

\begin{align*}\ensuremath
\sum_{n=0}^{+\infty}A^n&=\sum_{n=0}^{+\infty}PD^nP^{-1}=P\left(\sum_{n=0}^{+\infty}D^n\right)P^{-1}\;(\text{par continuité de l'application}\;M\mapsto PMP^{-1})\\
  &=P\;\text{diag}\left(\sum_{n=0}^{+\infty}\left( \frac{1}{2}\right)^n,\sum_{n=0}^{+\infty}\left(- \frac{1}{3}\right)^n\right)P^{-1}=P\;\text{diag}\left( \frac{1}{1- \frac{1}{2}}, \frac{1}{1+ \frac{1}{3}}\right)P^{-1}\\
  &=\left(
\begin{array}{cc}
1&1\\
1&2
\end{array}
\right)\left(
\begin{array}{cc}
2&0\\
0& \frac{3}{4}
\end{array}
\right)\left(
\begin{array}{cc}
2&-1\\
-1&1
\end{array}
\right)=\left(
\begin{array}{cc}
\rule[-4mm]{0mm}{0mm}2& \frac{3}{4}\\
2& \frac{3}{2}
\end{array}
\right)\left(
\begin{array}{cc}
2&-1\\
-1&1
\end{array}
\right)=\left(
\begin{array}{cc}
\rule[-4mm]{0mm}{0mm} \frac{13}{4}&- \frac{5}{4}\\
 \frac{5}{2}&- \frac{1}{2}
\end{array}
\right).
\end{align*}

\begin{center}
\shadowbox{
$\sum_{n=0}^{+\infty}A^n=\left(
\begin{array}{cc}
\rule[-4mm]{0mm}{0mm} \frac{13}{4}&- \frac{5}{4}\\
 \frac{5}{2}&- \frac{1}{2}
\end{array}
\right)$.
}
\end{center}

\textbf{Remarque.} D'après l'exercice suivant, la matrice obtenue est $(I-A)^{-1}$.
\fincorrection
\correction{005867}
Soit $A\in\mathcal{M}_p(\Cc)$ telle que $\|A\|<1$. Pour tout entier naturel $n$, on a $\|A^n\|\leqslant\|A\|^n$. Puisque $\|A\|<1$, la série numérique de terme général $\|A\|^n$, $n\in\Nn$, converge. Il en est de même de la série de terme général $\|A^n\|$ et donc la série de terme général $A^n$, $n\in\Nn$, converge absolument. 
Puisque $\mathcal{M}_p(\Cc)$ est complet en tant que $\Cc$ espace de dimension finie, on en déduit que la série de terme général $A^n$, $n\in\Nn$, converge. De plus,

\begin{align*}\ensuremath
(I-A)\sum_{n=0}^{+\infty}A^n&=(I-A)\lim_{n \rightarrow +\infty}\left(\sum_{k=0}^{n}A^k\right)=\lim_{n \rightarrow +\infty}\left((I-A)\sum_{k=0}^{n}A^k\right)\;(\text{par continuité de l'application}\;M\mapsto(I-A)M)\\
 &=\lim_{n \rightarrow +\infty}(I-A^{n+1})=I\;(\lim_{n \rightarrow +\infty}A^{n+1}=0\;\text{car}\;\forall n\in\Nn,\;\|A^{n+1}\|\leqslant\|A\|^{n+1}).
\end{align*}

Ainsi, la matrice $I-A$ est inversible à droite et donc inversible et de plus, $(I-A)^{-1}=\sum_{n=0}^{+\infty}A^n$. On en déduit encore

\begin{center}
$\|(I-A)^{-1}-(I+A)\|=\left\|\sum_{n=2}^{+\infty}A^n\right\|\leqslant\sum_{n=2}^{+\infty}\|A\|^n= \frac{\|A\|^2}{1-\|A\|}$.
\end{center}
\fincorrection
\correction{005868}
Soit $A\in\mathcal{M}_n(\Cc)$. On sait que d'une part $\text{det}(\text{exp}(A))\neq0$ et d'autre part $\text{exp}(A)=\lim_{p \rightarrow +\infty}\left(\sum_{k=0}^{p} \frac{A^k}{k!}\right)$. Par continuité du déterminant, on a donc $\lim_{p \rightarrow +\infty}\text{det}\left(\sum_{k=0}^{p} \frac{A^k}{k!}\right)=\text{det}(\text{exp}(A))\neq0$. Par suite, il existe $p_0\in\Nn$ tel que $\forall p\geqslant p_0$, $\text{det}\left(\sum_{k=0}^{p} \frac{A^k}{k!}\right)\neq0$ et donc tel que $\sum_{k=0}^{p} \frac{A^k}{k!}\in GL_n(\Rr)$.

\fincorrection
\correction{005869}
\begin{enumerate}
 \item  $\chi_A=\left|
\begin{array}{ccc}
3-X&2&2\\
1&-X&1\\
-1&1&-X
\end{array}
\right|=(3-X)(X^2-1)-(-2X-2)-(2X+2)=-(X+1)(X-1)(X-3)$.

Soit $n\in\Nn$. La division euclidienne de $X^n$ par $\chi_A$ s'écrit $X^n=Q_n\times\chi_A+a_nX^2+b_nX+c_n$ où $Q_n\in\Rr[X]$ et $(a_n,b_n,c_n)\in\Rr^3$. En évaluant les deux membres de cette égalité en $-1$, $1$ et $3$, on obtient

\begin{center}
$\left\{
\begin{array}{l}
a_n-b_n+c_n=(-1)^n\\
a_n+b_n+c_n=1\\
9a_n+3b_n+c_n=3^n
\end{array}
\right.
\Rightarrow\left\{
\begin{array}{l}
b_n= \frac{1}{2}(1-(-1)^n)\\
a_n+c_n= \frac{1}{2}(1+(-1)^n)\\
8a_n+ \frac{3}{2}(1-(-1)^n)+ \frac{1}{2}(1+(-1)^n)=3^n
\end{array}
\right.\Rightarrow\left\{
\begin{array}{l}
a_n= \frac{1}{8}(3^n-2+(-1)^n)\\
\rule[-4mm]{0mm}{10mm}b_n= \frac{1}{2}(1-(-1)^n)\\
c_n= \frac{1}{8}(-3^n+6+3(-1)^n)
\end{array}
\right.$.
\end{center}

Le théorème de \textsc{Cayley}-\textsc{Hamilton} fournit alors 

\begin{center}
$\forall n\in\Nn$, $A^n= \frac{1}{8}((3^n-2+(-1)^n)A^2+4(1-(-1)^n)A+(-3^n+6+3(-1)^n)I_3)$.
\end{center}

Maintenant,

\begin{center}
$A^2=\left(
\begin{array}{ccc}
3&2&2\\
1&0&1\\
-1&1&0
\end{array}
\right)\left(
\begin{array}{ccc}
3&2&2\\
1&0&1\\
-1&1&0
\end{array}
\right)=\left(
\begin{array}{ccc}
9&8&8\\
2&3&2\\
-2&-2&-1
\end{array}
\right)$
\end{center}

et donc, pour tout réel $t$,

\begin{align*}\ensuremath
\text{exp}(tA)&=\sum_{n=0}^{+\infty} \frac{t^n}{n!}A^n=\sum_{n=0}^{+\infty} \frac{t^n}{n!}. \frac{1}{8}((3^n-2+(-1)^n)A^2+4(1-(-1)^n)A+(-3^n+6+3(-1)^n)I_3)\\
 &= \frac{e^{3t}-2e^t+e^{-t}}{8}\left(
\begin{array}{ccc}
9&8&8\\
2&3&2\\
-2&-2&-1
\end{array}
\right)+ \frac{4(e^t-e^{-t})}{8}\left(
\begin{array}{ccc}
3&2&2\\
1&0&1\\
-1&1&0
\end{array}
\right)+ \frac{-e^{3t}+6e^t+3e^{-t}}{8}\left(
\begin{array}{ccc}
1&0&0\\
0&1&0\\
0&0&1
\end{array}
\right)\\
&= \frac{1}{8}\left(
\begin{array}{ccc}
8e^{3t}&8e^{3t}-8e^t&8e^{3t}-8e^t\\
2e^{3t}-2e^{-t}&2e^{3t}+6e^{-t}&2e^{3t}-2e^{-t}\\
-2e^{3t}+2e^{-t}&-2e^{3t}+8e^t-6e^{-t}&2e^{3t}+8e^t+2e^{-t}
\end{array}
\right)\\
 &= \frac{1}{4}\left(
\begin{array}{ccc}
4e^{3t}&4e^{3t}-4e^t&4e^{3t}-4e^t\\
e^{3t}-e^{-t}&e^{3t}+3e^{-t}&e^{3t}-e^{-t}\\
-e^{3t}+e^{-t}&-e^{3t}+4e^t-3e^{-t}&e^{3t}+4e^t+e^{-t}
\end{array}
\right).
\end{align*}

\begin{center}
\shadowbox{
$\forall t\in\Rr$, $\text{exp}(tA)= \frac{1}{4}\left(
\begin{array}{ccc}
4e^{3t}&4e^{3t}-4e^t&4e^{3t}-4e^t\\
e^{3t}-e^{-t}&e^{3t}+3e^{-t}&e^{3t}-e^{-t}\\
-e^{3t}+e^{-t}&-e^{3t}+4e^t-3e^{-t}&e^{3t}+4e^t+e^{-t}
\end{array}
\right)$.
}
\end{center}

\item  $\chi_A=\left|
\begin{array}{ccc}
4-X&1&1\\
6&4-X&2\\
-10&-4&-2-X
\end{array}
\right|=(4-X)(X^2-2X)-6(-X+2)-10(X-2)=(X-2)[-X(X-4)+6-10]=-(X-2)(X^2-4X+4)=-(X-2)^3$. On est dans la situation où $A$ a une unique valeur propre. D'après le théorème de \textsc{Cayley}-\textsc{Hamilton}, $(A-2I_3)^3=0$ et donc pour tout réel $t$,

\begin{align*}\ensuremath
\text{exp}(tA)&=\text{exp}(t(A-2I_3)+2tI_3)=\text{exp}(t(A-2I_3))\times\text{exp}(2tI_3)\;(\text{car les matrices}\;t(A-2I_3)\;\text{et}\;2tI_3\;\text{commutent})\\
 &=\left(I_3+t(A-2I_3)+ \frac{t^2}{2}(A-2I_3)^2\right)\times e^{2t}I_3\\
 &=e^{2t}\left(
\begin{array}{ccc}
1&0&0\\
0&1&0\\
0&0&1
\end{array}
\right)+te^{2t}\left(
\begin{array}{ccc}
2&1&1\\
6&2&2\\
-10&-4&-4
\end{array}
\right)+ \frac{t^2e^{2t}}{2}\left(
\begin{array}{ccc}
2&1&1\\
6&2&2\\
-10&-4&-4
\end{array}
\right)\left(
\begin{array}{ccc}
2&1&1\\
6&2&2\\
-10&-4&-4
\end{array}
\right)
\\
 &=e^{2t}\left(
\begin{array}{ccc}
1&0&0\\
0&1&0\\
0&0&1
\end{array}
\right)+te^{2t}\left(
\begin{array}{ccc}
2&1&1\\
6&2&2\\
-10&-4&-4
\end{array}
\right)+ \frac{t^2e^{2t}}{2}\left(
\begin{array}{ccc}
0&0&0\\
4&2&2\\
-4&-2&-2
\end{array}
\right)\\
 &=\left(
\begin{array}{ccc}
(2t+1)e^{2t}&te^{2t}&te^{2t}\\
(2t^2+6t)e^{2t}&(t^2+2t+1)e^{2t}&(t^2+2t)e^{2t}\\
(-2t^2-10t)e^{2t}&(-t^2-4t)e^{2t}&(-t^2-4t)e^{2t}
\end{array}
\right).
\end{align*}

\begin{center}
\shadowbox{
$\forall t\in\Rr$, $\text{exp}(tA)=\left(
\begin{array}{ccc}
(2t+1)e^{2t}&te^{2t}&te^{2t}\\
(2t^2+6t)e^{2t}&(t^2+2t+1)e^{2t}&(t^2+2t)e^{2t}\\
(-2t^2-10t)e^{2t}&(-t^2-4t)e^{2t}&(-t^2-4t)e^{2t}
\end{array}
\right)$.
}
\end{center}
\end{enumerate}
\fincorrection
\correction{005870}
$\;$

$\chi_A=\left|
\begin{array}{ccc}
-X&1/2&-2\\
1/2&-X&0\\
0&0&-1/2-X
\end{array}
\right|=-X\left(X^2+ \frac{1}{2}X\right)- \frac{1}{2}\left(- \frac{1}{2}X- \frac{1}{4}\right)=-X^2\left(X+ \frac{1}{2}\right)+ \frac{1}{4}\left(X+ \frac{1}{2}\right)=-\left(X+ \frac{1}{2}\right)^2\left(X- \frac{1}{2}\right)$. 

Soit $n\in\Nn$. La division euclidienne de $X^n$ par $\chi_A$ s'écrit $X^n=Q_n\chi_A+a_nX^2+b_nX+c_n$ où $Q_n\in\Rr[X]$ et $(a_n,b_n,c_n)\in\Rr^3$.

On évalue les deux membres de cette égalité en $ \frac{1}{2}$ et $- \frac{1}{2}$ et on obtient $ \frac{a_n}{4}+ \frac{b_n}{2}+c_n=\left( \frac{1}{2}\right)^n$ et $ \frac{a_n}{4}- \frac{b_n}{2}+c_n=\left(- \frac{1}{2}\right)^n$.

Puis en dérivant les deux membres de l'égalité et en évaluant en $- \frac{1}{2}$, on obtient $-a_n+b_n=n\left(- \frac{1}{2}\right)^{n-1}=-2n\left(- \frac{1}{2}\right)^{n}$. Maintenant,

\begin{align*}\ensuremath
\left\{
\begin{array}{l}
 \frac{a_n}{4}+ \frac{b_n}{2}+c_n=\left( \frac{1}{2}\right)^n\\
 \frac{a_n}{4}- \frac{b_n}{2}+c_n=\left(- \frac{1}{2}\right)^n\\
-a_n+b_n=-2n\left(- \frac{1}{2}\right)^{n}
\end{array}
\right.&\Leftrightarrow\left\{
\begin{array}{l}
b_n=\left( \frac{1}{2}\right)^n-\left(- \frac{1}{2}\right)^n\\
 \frac{a_n}{2}+2c_n=\left( \frac{1}{2}\right)^n+\left(- \frac{1}{2}\right)^n\\
-a_n+\left( \frac{1}{2}\right)^n-\left(- \frac{1}{2}\right)^n=-2n\left(- \frac{1}{2}\right)^{n}
\end{array}
\right.
\\
 &\Leftrightarrow\left\{
\begin{array}{l}
a_n=\left( \frac{1}{2}\right)^n+(2n-1)\left(- \frac{1}{2}\right)^n\\
b_n=\left( \frac{1}{2}\right)^n-\left(- \frac{1}{2}\right)^n\\
c_n= \frac{1}{4}\left( \frac{1}{2}\right)^n- \frac{2n-3}{4}\left(- \frac{1}{2}\right)^n
\end{array}
\right.
\end{align*}

Donc $\forall n\in\Nn$, $A^n=\left(\left( \frac{1}{2}\right)^n+(2n-1)\left(- \frac{1}{2}\right)^n\right)A^2+\left(\left( \frac{1}{2}\right)^n-\left(- \frac{1}{2}\right)^n\right)A+\left( \frac{1}{4}\left( \frac{1}{2}\right)^n- \frac{2n-3}{4}\left(- \frac{1}{2}\right)^n\right)I_3$ avec $A^2=\left(
\begin{array}{ccc}
0&1/2&-2\\
1/2&0&0\\
0&0&-1/2
\end{array}
\right)\left(
\begin{array}{ccc}
0&1/2&-2\\
1/2&0&0\\
0&0&-1/2
\end{array}
\right)=\left(
\begin{array}{ccc}
1/4&0&1\\
0&1/4&-1\\
0&0&1/4
\end{array}
\right)$. On en déduit que pour $|t|<2$,

\begin{align*}\ensuremath
\ln(I_3+tA)&=\sum_{n=1}^{+\infty} \frac{(-1)^{n-1}t^n}{n}A^n\\
 &=\left(\sum_{n=1}^{+\infty} \frac{(-1)^{n-1}t^n}{n}\left(\left( \frac{1}{2}\right)^n+(2n-1)\left(- \frac{1}{2}\right)^n\right)\right)A^2+\left(\sum_{n=1}^{+\infty} \frac{(-1)^{n-1}t^n}{n}\left(\left( \frac{1}{2}\right)^n-\left(- \frac{1}{2}\right)^n\right)\right)A\\
  &+\left(\sum_{n=1}^{+\infty} \frac{(-1)^{n-1}t^n}{n} \frac{1}{4}\left( \frac{1}{2}\right)^n- \frac{2n-3}{4}\left(- \frac{1}{2}\right)^n\right)I_3.
\end{align*}

et donc

\begin{align*}\ensuremath
\ln(I_3+tA)&=\left(\sum_{n=1}^{+\infty} \frac{(-1)^{n-1}(t/2)^n}{n}-2\sum_{n=1}^{+\infty}(t/2)^n-\sum_{n=1}^{+\infty} \frac{(-1)^{n-1}(-t/2)^n}{n}\right)A^2\\
 &+\left(\sum_{n=1}^{+\infty} \frac{(-1)^{n-1}(t/2)^n}{n}-\sum_{n=1}^{+\infty} \frac{(-1)^{n-1}(-t/2)^n}{n}\right)A\\
  &+ \frac{1}{4}\left(\sum_{n=1}^{+\infty} \frac{(-1)^{n-1}(t/2)^n}{n}+2\sum_{n=1}^{+\infty}(t/2)^n+3\sum_{n=1}^{+\infty} \frac{(-1)^{n-1}(-t/2)^n}{n}\right)I_3\\
  &=\left(\ln\left(1+ \frac{t}{2}\right)-2\left( \frac{1}{1- \frac{t}{2}}-1\right)-\ln\left(1- \frac{t}{2}\right)\right)A^2+\left(\ln\left(1+ \frac{t}{2}\right)-\ln\left(1- \frac{t}{2}\right)\right)A\\
  &+ \frac{1}{4}\left(\ln\left(1+ \frac{t}{2}\right)+2\left( \frac{1}{1- \frac{t}{2}}-1\right)+3\ln\left(1- \frac{t}{2}\right)\right)I_3\\
  &=\left(\ln\left( \frac{2+t}{2-t}\right)- \frac{2t}{2-t}\right)\left(
\begin{array}{ccc}
1/4&0&1\\
0&1/4&-1\\
0&0&1/4
\end{array}
\right)+\ln\left( \frac{2+t}{2-t}\right)\left(
\begin{array}{ccc}
0&1/2&-2\\
1/2&0&0\\
0&0&-1/2
\end{array}
\right)\\
  &+ \frac{1}{4}\left(\ln\left(1+ \frac{t}{2}\right)+ \frac{2t}{2-t}+3\ln\left(1- \frac{t}{2}\right)\right)\left(
\begin{array}{ccc}
1&0&0\\
0&1&0\\
0&0&1
\end{array}
\right)\\
 &=\left(
 \begin{array}{ccc}
 \frac{1}{2}\ln\left(1- \frac{t^2}{4}\right)& \frac{1}{2}\ln\left( \frac{2+t}{2-t}\right)&-\ln\left( \frac{2+t}{2-t}\right)- \frac{2t}{2-t}\\
 \frac{1}{2}\ln\left( \frac{2+t}{2-t}\right)& \frac{1}{2}\ln\left(1- \frac{t^2}{4}\right)&-\ln\left( \frac{2+t}{2-t}\right)+ \frac{2t}{2-t}\\
 0&0&\ln\left(1- \frac{t}{2}\right)
 \end{array}
 \right).
\end{align*}

\begin{center}
\shadowbox{
$\forall t\in]-2,2[$, $\ln(I_3+tA)=\left(
 \begin{array}{ccc}
 \frac{1}{2}\ln\left(1- \frac{t^2}{4}\right)& \frac{1}{2}\ln\left( \frac{2+t}{2-t}\right)&-\ln\left( \frac{2+t}{2-t}\right)- \frac{2t}{2-t}\\
 \frac{1}{2}\ln\left( \frac{2+t}{2-t}\right)& \frac{1}{2}\ln\left(1- \frac{t^2}{4}\right)&-\ln\left( \frac{2+t}{2-t}\right)+ \frac{2t}{2-t}\\
 0&0&\ln\left(1- \frac{t}{2}\right)
 \end{array}
 \right)$.
}
\end{center}
\fincorrection
\correction{005871}
\begin{enumerate}
 \item
   \begin{enumerate}
   \item  Soit $\overrightarrow{\omega}\in\Rr^3$. $f_{\overrightarrow{\omega}}$ est un endomorphisme de $\Rr^3$ par bilinéarité du produit vectoriel. De plus, pour $(\overrightarrow{x},\overrightarrow{y})\in(\Rr^3)^2$,

\begin{center}
$f_{\overrightarrow{\omega}}\left(\overrightarrow{x}\right).\overrightarrow{y}=\left(\overrightarrow{\omega}\wedge\overrightarrow{x}\right).\overrightarrow{y}=\left[\overrightarrow{\omega},\overrightarrow{x},\overrightarrow{y}\right]=-\left[\overrightarrow{\omega},\overrightarrow{y},\overrightarrow{x}\right]=-\left(\overrightarrow{\omega}\wedge\overrightarrow{y}\right).\overrightarrow{x}=-\overrightarrow{x}.f_{\overrightarrow{\omega}}\left(\overrightarrow{y}\right)$.
\end{center}

Donc,

\begin{center}
\shadowbox{
$\forall\overrightarrow{x}\in\Rr^3$, $f_{\overrightarrow{\omega}}\in\mathcal{A}(\Rr^3)$.
}
\end{center}

   \item Soit $\begin{array}[t]{cccc}
\varphi~:&\Rr^3&\rightarrow&\mathcal{A}(\Rr^3)\\
 &\overrightarrow{\omega}&\mapsto&f_{\overrightarrow{\omega}}
\end{array}$. 

\textbullet~Vérifions que $\varphi\in\mathcal{L}(\Rr^3,\mathcal{A}(\Rr^3))$. Soient $(\lambda_1,\lambda_2)\in\Rr^2$ et $(\overrightarrow{\omega}_1,\overrightarrow{\omega}_2)\in(\Rr^3)^2$. Pour tout $\overrightarrow{x}\in\Rr^3$,

\begin{align*}\ensuremath
\left(\varphi(\lambda_1\overrightarrow{\omega}_1+\lambda_2\overrightarrow{\omega}_2)\right)\left(\overrightarrow{x}\right)&=f_{\lambda_1\overrightarrow{\omega}_1+\lambda_2\overrightarrow{\omega}_2}\left(\overrightarrow{x}\right)=\left(\lambda_1\overrightarrow{\omega}_1+\lambda_2\overrightarrow{\omega}_2\right)\wedge=\lambda_1\left(\overrightarrow{\omega}_1\wedge\overrightarrow{x}\right)+\lambda_2\left(\overrightarrow{\omega}_2\wedge\overrightarrow{x}\right)\\
 &=\lambda_1f_{\overrightarrow{\omega}_1}\left(\overrightarrow{x}\right)+\lambda_2f_{\overrightarrow{\omega}_2}\left(\overrightarrow{x}\right)=\left((\lambda_1\varphi(\overrightarrow{\omega}_1)+\lambda_2\varphi(\overrightarrow{\omega}_2)\right)\left(\overrightarrow{x}\right)
\end{align*}

et donc $\varphi(\lambda_1\overrightarrow{\omega}_1+\lambda_2\overrightarrow{\omega}_2)=\lambda_1\varphi(\overrightarrow{\omega}_1)+\lambda_2\varphi(\overrightarrow{\omega}_2)$. On a montré que $\varphi\in\mathcal{L}(\Rr^3,\mathcal{A}(\Rr^3))$.

\textbullet~Vérifions que $\varphi$ est injective. Soit $\omega\in\Rr^3$.

\begin{center}
$\overrightarrow{\omega}\in\text{Ker}(\varphi)\Rightarrow f_{\overrightarrow{\omega}}=0\Rightarrow\forall\overrightarrow{x}\in\Rr^3$, $\overrightarrow{\omega}\wedge\overrightarrow{x}=\overrightarrow{0}$.
\end{center}

On applique alors ce dernier résultat à deux vecteurs non colinéaires $\overrightarrow{u}$ et $\overrightarrow{v}$. On obtient $\overrightarrow{\omega}\wedge\overrightarrow{u}=\overrightarrow{\omega}\wedge\overrightarrow{v}=\overrightarrow{0}$ et donc $\overrightarrow{x}\in\text{Vect}\left(\overrightarrow{u},\overrightarrow{v}\right)=\left\{\overrightarrow{0}\right\}$. On a montré que $\varphi$ est injective.

\textbullet~Enfin, $\text{dim}\left(\mathcal{A}_3(\Rr)\right)= \frac{3\times(3-1)}{2}=3=\text{dim}(\Rr^3)<+\infty$. On en déduit que $\varphi$ est un isomorphisme de $\Rr^3$ sur $\mathcal{A}(\Rr^3)$. En particulier,

\begin{center}
\shadowbox{
$\forall f\in\mathcal{A}(\Rr^3)$, $\exists\overrightarrow{\omega}\in\Rr^3/\;f=f_{\overrightarrow{\omega}}$.
}
\end{center}
   \end{enumerate}

\item  Soit $\overrightarrow{\omega}\in\Rr^3$. Si $\overrightarrow{\omega}=\overrightarrow{0}$, alors $f_{\overrightarrow{\omega}}=0$ et donc $\text{exp}\left(f_{\overrightarrow{\omega}}\right)=Id_{\Rr^3}$.

On suppose dorénavant $\overrightarrow{\omega}\neq\overrightarrow{0}$. On pose $\overrightarrow{e_3}= \frac{1}{\left\|\overrightarrow{\omega}\right\|}\overrightarrow{\omega}$ puis on complète la famille orthonormale $\left(\overrightarrow{e_3}\right)$ en une base orthonormale directe $\left(\overrightarrow{e_1},\overrightarrow{e_2},\overrightarrow{e_3}\right)$ (en particulier $\overrightarrow{e_3}\wedge\overrightarrow{e_1}=\overrightarrow{e_2}$).

\textbullet~Puisque $\overrightarrow{e_3}$ est colinéaire à $\overrightarrow{\omega}$, $f_{\overrightarrow{\omega}}\left(\overrightarrow{e_3}\right)=\overrightarrow{0}$. On en déduit que

\begin{center}
$\text{exp}\left(f_{\overrightarrow{\omega}}\right)\left(\overrightarrow{e_3}\right)=Id\left(\overrightarrow{e_3}\right)+f_{\overrightarrow{\omega}}\left(\overrightarrow{e_3}\right)+ \frac{1}{2}f_{\overrightarrow{\omega}}^2\left(\overrightarrow{e_3}\right)+ \frac{1}{6}f_{\overrightarrow{\omega}}^3\left(\overrightarrow{e_3}\right)+\ldots=\overrightarrow{e_3}$.
\end{center}

\textbullet~D'autre part, $f_{\overrightarrow{\omega}}\left(\overrightarrow{e_1}\right)=\overrightarrow{\omega}\wedge\overrightarrow{e_1}=\left\|\overrightarrow{\omega}\right\|\overrightarrow{e_3}\wedge\overrightarrow{e_1}=\left\|\overrightarrow{\omega}\right\|\overrightarrow{e_2}$ et de même $f_{\overrightarrow{\omega}}\left(\overrightarrow{e_2}\right)=-\left\|\overrightarrow{\omega}\right\|\overrightarrow{e_1}$.

On en déduit que $f_{\overrightarrow{\omega}}^2\left(\overrightarrow{e_1}\right)=-\left\|\overrightarrow{\omega}\right\|^2\overrightarrow{e_1}$ et donc que 

\begin{center}
$\forall n\in\Nn$, $f_{\overrightarrow{\omega}}^{2n}\left(\overrightarrow{e_1}\right)=(-1)^n\left\|\overrightarrow{\omega}\right\|^{2n}\overrightarrow{e_1}$\quad puis\quad$f_{\overrightarrow{\omega}}^{2n+1}\left(\overrightarrow{e_1}\right)=(-1)^n\left\|\overrightarrow{\omega}\right\|^{2n+1}\overrightarrow{e_2}$.
\end{center}

Par suite,

\begin{align*}\ensuremath
\text{exp}\left(f_{\overrightarrow{\omega}}\right)\left(\overrightarrow{e_1}\right)&=\sum_{n=0}^{+\infty} \frac{1}{n!}f_{\overrightarrow{\omega}}^{n}\left(\overrightarrow{e_1}\right)\\
 &=\left(\sum_{n=0}^{+\infty}(-1)^n \frac{\left\|\overrightarrow{\omega}\right\|^{2n}}{(2n)!}\right)\overrightarrow{e_1}+\left(\sum_{n=0}^{+\infty}(-1)^n \frac{\left\|\overrightarrow{\omega}\right\|^{2n+1}}{(2n+1)!}\right)\overrightarrow{e_2}\;(\text{somme de deux séries convergentes})\\
 &=\cos\left(\left\|\overrightarrow{\omega}\right\|\right)\overrightarrow{e_1}+\sin\left(\left\|\overrightarrow{\omega}\right\|\right)\overrightarrow{e_2}.
\end{align*}

De même,

\begin{center}
$\forall n\in\Nn$, $f_{\overrightarrow{\omega}}^{2n}\left(\overrightarrow{e_2}\right)=(-1)^n\left\|\overrightarrow{\omega}\right\|^{2n}\overrightarrow{e_2}$\quad puis\quad$f_{\overrightarrow{\omega}}^{2n+1}\left(\overrightarrow{e_2}\right)=-(-1)^n\left\|\overrightarrow{\omega}\right\|^{2n+1}\overrightarrow{e_1}$.
\end{center}

et donc

\begin{align*}\ensuremath
\text{exp}\left(f_{\overrightarrow{\omega}}\right)\left(\overrightarrow{e_2}\right)&=\sum_{n=0}^{+\infty} \frac{1}{n!}f_{\overrightarrow{\omega}}^{n}\left(\overrightarrow{e_2}\right)=\left(\sum_{n=0}^{+\infty}(-1)^n \frac{\left\|\overrightarrow{\omega}\right\|^{2n}}{(2n)!}\right)\overrightarrow{e_2}-\left(\sum_{n=0}^{+\infty}(-1)^n \frac{\left\|\overrightarrow{\omega}\right\|^{2n+1}}{(2n+1)!}\right)\overrightarrow{e_1}\\
 &=-\sin\left(\left\|\overrightarrow{\omega}\right\|\right)\overrightarrow{e_1}+\cos\left(\left\|\overrightarrow{\omega}\right\|\right)\overrightarrow{e_2}.
\end{align*}

Ainsi, la matrice de $\text{exp}\left(f_{\overrightarrow{\omega}}\right)$ dans la base orthonormée directe $\left(\overrightarrow{e_1},\overrightarrow{e_2},\overrightarrow{e_3}\right)$ est $\left(
\begin{array}{ccc}
\cos\left(\left\|\overrightarrow{\omega}\right\|\right)&-\sin\left(\left\|\overrightarrow{\omega}\right\|\right)&0\\
\sin\left(\left\|\overrightarrow{\omega}\right\|\right)&\cos\left(\left\|\overrightarrow{\omega}\right\|\right)&0\\
0&0&1
\end{array}
\right)$ et $\text{exp}(f_{\overrightarrow{\omega}})$ est la rotation d'angle $\left\|\overrightarrow{\omega}\right\|$ autour de $\overrightarrow{\omega}$.
\end{enumerate}
\fincorrection
\correction{005872}
On munit $\mathcal{M}_n(\Cc)$ d'une norme sous-multiplicative notée $\|\;\|$. Soit $A\in\mathcal{M}_n(\Cc)$. Soit $p\in\Nn^*$.

\begin{center}
$\left\|\sum_{k=0}^{p} \frac{A^k}{k!}-\left(I+ \frac{A}{p}\right)^p\right\| =\left\|\sum_{k=0}^{p}\left( \frac{1}{k!}- \frac{C_p^k}{p^k}\right)A^k\right\|\leqslant\sum_{k=0}^{p}\left| \frac{1}{k!}- \frac{C_p^k}{p^k}\right|\|A\|^k $.
\end{center}

Maintenant, $\forall k\in\llbracket1,p\rrbracket$, $ \frac{1}{k!}- \frac{C_p^k}{p^k}= \frac{1}{k!}\left(1-  \frac{\overbrace{p\times(p-1)\times\ldots\times(p-k+1)}^{k}}{\underbrace{p\times p\times\ldots\times p}_{k}}\right)\geqslant 0$. Donc,

\begin{center}
$\sum_{k=0}^{p}\left| \frac{1}{k!}- \frac{C_p^k}{p^k}\right|\|A\|^k=\sum_{k=0}^{p} \frac{\|A\|^k}{k!}-\left(1+ \frac{\|A\|^p}{p}\right)^p\underset{n\rightarrow+\infty}{\rightarrow}e^{\|A\|}-e^{\|A\|}= 0$.
\end{center}

On en déduit que $\sum_{k=0}^{p} \frac{A^k}{k!}-\left(I+ \frac{A}{p}\right)^p$ tend vers $0$ quand $p$ tend vers $+\infty$ et puisque $\sum_{k=0}^{p} \frac{A^k}{k!}$ tend vers $\text{exp}(A)$ quand $p$ tend vers $+\infty$, il en est de même de $\left(I+ \frac{A}{p}\right)^p$.

\begin{center}
\shadowbox{
$\forall A\in\mathcal{M}_n(\Cc)$, $\text{exp}(A)=\lim_{p \rightarrow +\infty}\left(I_n+ \frac{A}{p}\right)^p$.
}
\end{center}
\fincorrection
\correction{005873}
Soit $A\in\mathcal{M}_n(\Rr)$. Soit $k\in\Nn$. Puisque $\chi_A$ est de degré $n$, la division euclidienne de $X^k$ par $\chi_A$ s'écrit

\begin{center}
$X^k=Q_k\times\chi_A+a_{n-1}^{(k)}X^{n-1}+\ldots+a_1^{(k)}X+a_0^{(k)}$ où $Q_k\in\Rr[\Cc]$ et $(a_0^{(k)},\ldots,a_{n-1}^{(k)})\in\Cc^n$.
\end{center}

Le théorème de \textsc{Cayley}-\textsc{Hamilton} montre alors que $A^k=a_{n-1}^{(k)}A^{n-1}+\ldots+a_1^{(k)}A+a_0^{(k)}I_n$.

Ainsi, $\forall k\in\Nn$, $A^k\in\text{Vect}(A^{n-1},\ldots,A,I_n)$ puis $\forall p\in\Nn$, $\sum_{k=0}^{p} \frac{A^k}{k!}\in\text{Vect}(A^{n-1},A^{n-2},\ldots,A,I_n)$. Enfin, puisque

$\text{Vect}(A^{n-1},A^{n-2},\ldots,A,I_n)$ est un fermé de $\mathcal{M}_n(\Cc)$ en tant que sous-espace vectoriel d'un espace vectoriel de dimension finie, $\text{exp}(A)=\lim_{p \rightarrow +\infty}\sum_{k=0}^{p} \frac{A^k}{k!}\in\text{Vect}(A^{n-1},A^{n-2},\ldots,A,I_n)$.

On a montré que

\begin{center}
\shadowbox{
$\forall A\in\mathcal{M}_n(\Cc)$, $\text{exp}(A)\in\Cc_{n-1}[A]$.
}
\end{center}
\fincorrection
\correction{004528}
\begin{enumerate}
  \item 
  \item $f(0) = 0$, $f(\pi) = \pi\sh 1$,
             $f\left(\frac\pi2\right) = \frac\pi2(e-\cos 1)$.
\end{enumerate}
 \fincorrection
\correction{004529}
\begin{enumerate}
  \item $\R^*$.
  \item CSI~: $\frac{\sqrt\pi}{2a} =  \int_{x=0}^\infty e^{-a^2x^2}\,d x \le f(a) \le  \int_{x=0}^\infty e^{-a^2x^2}\,d x + 1 = \frac{\sqrt\pi}{2a} + 1$. Donc $af(a)\to \frac{\sqrt\pi}2$ lorsque $a\to0^+$.
  \item TCM~: $f(a)\to 1$ lorsque $a\to+\infty$.
\end{enumerate}
 \fincorrection
\nocorrection
\nocorrection
\correction{004532}
\begin{enumerate}
  \item 
  \item $f(x) = \frac{e^x}{(e^x-1)^2}$.
\end{enumerate}
 \fincorrection
\nocorrection
\correction{004534}
\begin{enumerate}
  \item 
  \item $xg(x) - g(x+1) = \frac1e$.
  \item CSA $ \Rightarrow  g' < 0$. $g(x)\to +\infty$ lorsque $x\to0^+$,
                               $g(x)\to 0$ lorsque $x\to+\infty$.
  \item $g(x) \sim \frac1x$ en $0^+$ et $g(x) \sim \frac1{ex}$ en
             $+\infty$.
\end{enumerate}
 \fincorrection
\nocorrection
\correction{004536}
\begin{enumerate}
  \item 
  \item CSA $ \Rightarrow  |R_n(x)| \le |u_{n+1}(x)| \le
             \ln\left(1+\frac1{n+1}\right)$.
  \item Non, $\|u_n\|_\infty = \ln\left(1+\frac1n\right)$.
\end{enumerate}
 \fincorrection
\correction{004537}
\begin{enumerate}
  \item 
  \item $f(x+1) = xf(x)-1$.
  \item 
\end{enumerate}
 \fincorrection
\correction{004538}
\begin{enumerate}
  \item CVU sur tout $[a,b]$.
  \item 
  \item $f(x+1) = f(x) + \frac\pi2 - \Arctan x$.
  \item $f(x+1) - f(x) \sim \frac 1x$ donc la suite $(f(n))$ diverge et
    $f$ est croissante $ \Rightarrow  \lim = +\infty$.
\end{enumerate}
 \fincorrection
\correction{004539}
$\frac1{t+1} = \sum_{n=0}^\infty (-1)^nt^n$.
\fincorrection
\correction{004540}
\begin{enumerate}
  \item $\frac{f_n(x)}{f_{n+1}(x)} = 1 - \frac{x(x+1)}{2n^2}
         +  o\left(\frac1{n^2}\right)$
         donc la série $\sum\ln f_n(x)$ est convergente pour tout $x\notin -\N^*$.
  \item 
  \item $\frac{f_n'(x)}{f_n(x)} \to -\gamma + \sum_{k=1}^\infty \frac x{k(k+x)}$ lorsque $n\to\infty$.
\end{enumerate}
 \fincorrection
\correction{004541}
Poser $t = xu$ puis intégrer deux fois par parties :
         $f_n(x) = 1 -  \int_{u=0}^1 (1-u)^{n+1}x\sin(xu)\,d u$
         donc $(f_n)$ converge simplement vers la fonction constante $1$, et
         la convergence est uniforme sur tout intervalle borné.
\fincorrection
\nocorrection
\correction{004543}
\begin{enumerate}
  \item cva si $|\cos x| < 1$, scv si $\cos x = 1$, dv si $\cos x = -1$.
  \item TCM en regroupant les termes deux par deux.
  \item $ \int_{x=0}^{\pi/2} -\ln(1-\cos x)\,d x$.
\end{enumerate}
 \fincorrection
\correction{004544}
\begin{enumerate}
  \item $F_n(X) = \sum_{k=0}^{n-1}\frac{e^{2ik\pi/n}}{X+n(1-e^{2ik\pi/n})}$.
  \item     $F_n(2x) - F_n(-2x) = \sum_{k=0}^{n-1}\frac{4xe^{2ik\pi/n}}{4x^2-n^2(1-e^{2ik\pi/n})^2}
                        = \sum_{k=0}^{n-1}\frac{x}{x^2e^{-2ik\pi/n}+n^2\sin(k\pi/n)^2}$.

    Supposons $n$ impair, et regroupons les termes conjugués obtenus pour $k$ et $n-k$~:

    $F_n(2x) - F_n(-2x) = \frac1x + \sum_{k=1}^{(n-1)/2}\Bigl(\underbrace{
    \frac{x}{x^2e^{-2ik\pi/n}+n^2\sin(k\pi/n)^2} + \frac{x}{x^2e^{2ik\pi/n}+n^2\sin(k\pi/n)^2}
    }_{=u(k,n,x)}\Bigr)$.

    On transforme la somme en série de $k=1$ à $k=\infty$ en posant $u(k,n,x) = 0$
    si $k>(n-1)/2$, puis on passe à la limite, sous réserve de justification,
    dans cette série pour $n\to\infty$, ce qui donne la formule demandée.
    
    Justification de l'interversion limite-série~: en utilisant $\sin(t)\ge \frac{2t}{\pi}$
    pour $0\le t\le \frac{\pi}2$ on a $|u(k,n,x)| \le \frac{2|x|}{4k^2-x^2}$ pour tout
    $k\ge |x/2|$, donc il y a convergence normale par rapport à $n$, à $x$ fixé.
  \item $\sum_{k=1}^\infty\frac{2}{x^2+k^2\pi^2} = \frac{\coth(x)}x-\frac1{x^2}$
    est normalement convergente sur $\R$, on peut passer à la limite pour $x\to 0$.
\end{enumerate}
 \fincorrection
\correction{004545}
\begin{enumerate}
  \item Transformation d'Abel.
  \item $f(x) = \Arctan\left(\frac{x\sin x}{1-x\cos x}\right)$.
  \item $\frac{\pi-1}2$.
\end{enumerate}
 \fincorrection
\correction{004546}
\begin{enumerate}
  \item 
  \item $\zeta(x) - \frac 1{x-1} =
    \sum_{n=1}^\infty \left(\frac 1{\strut n^x} -  \int_{t=n}^{n+1} \frac{d t}{\strut t^x} \right)$.

    A $n$ fixé, $\frac 1{n^x} -  \int_{t=n}^{n+1} \frac{d t}{t^x}\to
                 \frac 1n -  \int_{t=n}^{n+1} \frac{d t}t$
   lorsque $x\to1^+$ et la convergence est monotone donc lorsque $x\to1^+$

    $\zeta(x) - \frac 1{x-1} \to
    \sum_{n=1}^\infty \left(\frac 1n -  \int_{t=n}^{n+1} \frac{d t}t\right) = \gamma$.
  \item $\sum_{n=1}^\infty \frac{(-1)^n\ln n}n = \eta'(1) = \gamma\ln 2 - \frac12\ln(2)^2$.
\end{enumerate}
 \fincorrection
\nocorrection
\correction{004548}
Pour $k$ fixé et $x\in{[0,1[}$ on a
$0\le f(x) \le\text{polynôme}(x) + \sum_{n=0}^\infty x^k = \text{polynôme}(x) + \frac1{1-x^k}$
et $\frac1{1-x^k}\sim\frac1{k(1-x)}$ au voisinage de~$1$ donc
$0\le f(x) \le \frac2{k(1-x)}$ pour $x$ suffisament proche de~$1$.
\fincorrection
\correction{004549}
\begin{enumerate}
  \item 
  \item soit $y\in{[c,d]}$ et $x_n = f_n^{-1}(y)$. La suite $(x_n)$ admet au plus
         une valeur d'adhérence, $x = f^{-1}(y)$.
  \item 
\end{enumerate}
 \fincorrection
\correction{004550}
Oui~: $|\sup f_n - \sup f| \le \|f_n-f\|_\infty$.
\fincorrection
\correction{004551}
\begin{enumerate}
  \item Il y a convergence normale sur tout intervalle $[a,+\infty[$ avec $a>0$.
    Il n'y a pas convergence normale au voisinage de $0$ car
    $\sup\Bigl\{\frac{xe^{-nx}}{\ln n},\ x\ge 0\Bigr\} = \frac1{en\ln n}$ atteint pour $x=\frac1n$ et
    $\sum\frac1{n\ln n}$ diverge (série de Bertrand).
    Par contre il y a convergence uniforme sur $[0,+\infty[$ car
    $$0\le \sum_{k=n}^\infty f_k(x) \le \frac1{\ln n}\sum_{k=n}^\infty xe^{-kx}
       = \frac{xe^{-nx}}{\ln n(1-e^{-x})} \le \frac{\sup\{t/(1-e^{-t}),\ t\ge 0\}}{\ln n}.$$
  \item 
  \item Lorsque $x\to0^+$, $\frac{S(x)-S(0)}x = \sum_{n=2}^\infty\frac{e^{-nx}}{\ln n}
    \to \sum_{n=2}^\infty\frac1{\ln n} = +\infty$ par convergence monotone.
  \item 
\end{enumerate}
 \fincorrection
\correction{004552}
Comparaison série-intégrale, $f(x)\to \ln(2)$ lorsque $x\to{1^-}$.
\fincorrection
\correction{004553}
\begin{enumerate}
  \item $-1<t<1$.
  \item Pour $0\le t<1$ et $n\ge 2$ on a~:
\begin{align*}
(1-t)\frac{t^n}{1-t^n}
&= \frac{t^n}{1+t+\dots+t^{n-1}}\cr
&= \frac{t^n}{n} + \frac{t^n((1-t) + (1-t^2) + \dots + (1-t^{n-1}))}{n(1+t+\dots+t^{n-1})}\cr
&= \frac{t^n}{n} + \frac{(t^n-t^{n+1})((n-1) + (n-2)t + \dots + t^{n-2})}{n(1+t+\dots+t^{n-1})}\cr
\end{align*}
d'où $0 \le (1-t)\frac{t^n}{1-t^n} - \frac{t^n}{n} \le \frac{n-1}{n}(t^n-t^{n+1}) \le t^n-t^{n+1}$
(vrai aussi si $n=1$) et en sommant~:
$$0\le (1-t)S(t) + \ln(1-t) \le 1.$$
\end{enumerate}
 \fincorrection
\correction{004554}
\begin{enumerate}
  \item La série converge normalement et $\phi$ est continue.
  \item $\phi$ est $1$-lipschitzienne, mais on ne peut rien en déduire pour $f$~:

pour $N$ fixé et $0<h\le \frac1{2.4^N}$, on a $|f(h)-f(0)| = f(h) \ge \sum_{n=1}^N3^nh = \frac{3^{N+1}-3}2h$
donc $f$ n'est pas lipschitzienne au voisinage de~$0$.
  \item D'après ce qui précède, le taux d'accroissement de~$f$ en $0$ est arbitrairement
grand, donc $f$ n'est pas dérivable en~$0$. On montre de même que $f$ n'est pas dérivable
en $x\in\R$.
\end{enumerate}
 \fincorrection
\correction{004555}
On suppose $h$ réel.
La série converge localement normalement sur $\R^*$ donc $f$ est définie sur~$\R$
et continue sur $\R^*$. Continuité en~$0$~: on pose $A_n = \sum_{k=n}^\infty a_k$
et $\varphi(t) = \frac{\sin^2(t)}{t^2}$ si $t\ne 0$, $\varphi(0) = 1$ ($\varphi$ est $\mathcal{C}^\infty$
sur~$\R$ comme somme d'une série entière de rayon infini). Pour $h\ne 0$ on a~:
$$f(h) = \sum_{n=1}^\infty(A_n-A_{n+1})\varphi(nh)
       = A_1\varphi(h) + \sum_{n=2}^\infty A_n(\varphi(nh)-\varphi((n-1)h))
       = A_1\varphi(h) + \sum_{n=2}^\infty A_n \int_{t=(n-1)h}^{nh}\varphi'(t)\,d t.$$
Cette dernière série est uniformément convergente sur~$\R$ car $A_n \to 0$ (lorsque $n\to\infty$) et
$ \int_{t=0}^{+\infty}|\varphi'(t)|\,d t$ est convergente.
\fincorrection
\correction{004556}
\begin{enumerate}
  \item Soit $k=\lfloor n/2\pi\rfloor$.
    On a $F_n(x) = \frac{2k\pi}n \int_{t=0}^{2\pi}f(x+t)f(t)\,d t + \frac1n \int_{t=2k\pi}^nf(x+t)f(t)\,d t
    \to \int_{t=0}^{2\pi}f(x+t)f(t)\,d t$  lorsque $n\to\infty$.
  \item Uniforme.
  \item Cauchy-Schwarz.
\end{enumerate}
 \fincorrection
\correction{004557}
$g = x \mapsto f(\tan(x/2))$ est limite uniforme de polynômes trigonométriques.
\fincorrection
\correction{004558}
\begin{enumerate}
  \item CSA~: $0\le f(x)\le \frac1{\sqrt{1+x^2}}$ donc $f(x)\to 0$ lorsque $x\to+\infty$.
  \item {$xf(x)$}
$= \sum_{p=0}^\infty \frac{x}{\sqrt{(2p+1)^2+x^2}} - \frac{x}{\sqrt{(2p+2)^2+x^2}}$\par
$= \sum_{p=0}^\infty \int_{t=2p+1}^{2p+2}\frac{xt}{(t^2+x^2)^{3/2}}\,d t$\par
$= \sum_{p=0}^\infty \int_{u=(2p+1)/x}^{(2p+2)/x}\frac{u}{(u^2+1)^{3/2}}\,d u$.\par


On a $ \int_{u=0}^\infty\frac u{(u^2+1)^{3/2}}\,d u = 1 = a + b$ avec~:

$a = \sum_{p=0}^\infty \int_{u=(2p)/x}^{(2p+1)/x}\frac{u}{(u^2+1)^{3/2}}\,d u$ et
$b = \sum_{p=0}^\infty \int_{u=(2p+1)/x}^{(2p+2)/x}\frac{u}{(u^2+1)^{3/2}}\,d u = xf(x)$.

$h$ : $u \mapsto\frac{u}{(u^2+1)^{3/2}}$ est croissante sur~$[0,\sqrt{\frac12}]$ et
décroissante sur $[\sqrt{\frac12},+\infty[$ donc $|a-b| \le \frac{3\|h\|_\infty}x$,
et $xf(x)\to \frac12$ lorsque $x\to+\infty$.
\end{enumerate}
 \fincorrection
\correction{004559}
On a
$ \int_{t=x}^{+\infty}\frac{d t}{\sh t} \le xS_1(x)
\le \frac x{\sh x} +  \int_{t=x}^{+\infty}\frac{d t}{\sh t}$
et $\frac1{\sh t}=\frac1t + O(t)$ donc
$ \int_{t=x}^{+\infty}\frac{d t}{\sh t}=-\ln(x)+ O(1)$.
On en déduit $S_1(x)\sim-\frac{\ln x}x$.

La même méthode ne marche pas pour $S_2$ car le terme résiduel,
$\frac x{\sh^2(x)}$ n'est pas négligeable devant
$ \int_{t=x}^{+\infty}\frac{d t}{\sh^2(t)}$.
Par contre, on peut remarquer que la série
$\sum_{n=1}^{\infty}{\frac{x^2}{\sh^{2}{(nx)}}}$ est normalement
convergente sur~$\R$, d'où
$S_2(x)\sim\frac{\zeta(2)}{x^2}$.
\fincorrection
\correction{004560}
\begin{enumerate}
  \item Lorsque $n\to\infty$, $S_n(t) = \Im\Bigl(\frac{e^{ix}-t^{n}e^{i(n+1)x}}{1-te^{ix}}\Bigr) \to
              \Im\Bigl(\frac{e^{ix}}{1-te^{ix}}\Bigr) = 
              \frac{\sin x}{1-2t\cos x + t^2}$ pour $-1<t<1$.
  \item $ \int_{t=0}^1 S_n(t)\,d t = \sum_{p=1}^n \frac{\sin(px)}{p}$.
             \par
             $ \int_{t=0}^1 S(t)\,d t
	      = (t-\cos x = u\sin x) =  \int_{u=-\cot x}^{\tan x/2} \frac{d u}{1+ u^2} = \frac{\pi-x}2$.
  \item TCD~: $|S_n(t)|\le \frac2{\sin x}$ intégrable par rapport à~$t$ sur $[0,1]$. On en déduit $\sum_{p=1}^\infty \frac{\sin(px)}{p} = \frac{\pi-x}2$.
\end{enumerate}
 \fincorrection
\correction{004561}
\begin{enumerate}
  \item $R$ est trivialement un $\R$-espace vectoriel. Le théorème de décomposition en éléments
    simples donne une base de~$R$ en se limitant aux éléments simples n'ayant pas de pôle dans~$[0,1]$.

    $R_{m,n}$ n'est pas un espace vectoriel. Par exemple $\frac1{X+1}$ et $\frac1{X+2}$ appartiennent
    à~$R_{0,1}$ mais pas leur somme.
  \item Soit $(f_k)$ une suite d'éléments de~$R_{m,n}$ telle que $\|g-f_k\|\to d$ quand $k\to\infty$.
    On note $f_k=P_k/Q_k$ avec $P_k\in\R_m[X]$, $Q_k\in\R_n[X]$ et $\|Q_k\|=1$.
    On a $\|P_k\| \le \|g-f_k\| + \|g\|$ donc les suites $(P_k)$ et $(Q_k)$
    sont bornées dans $\R_m[X]$ et $\R_n[X]$. Quitte à prendre une sous-suite, on se ramène
    au cas $P_k\to P\in\R_m[X]$ et $Q_k\to Q\in\R_n[X]$ (quand $k\to\infty$) avec de plus
    $\|Q\|=1$.
    
    Si $Q$ n'a pas de racine dans~$[0,1]$, il existe $\alpha>0$ tel
    que $|Q(x)| \ge \alpha$ pour tout~$x\in[0,1]$, donc $|Q_k(x)| \ge \frac12\alpha$
    pour tout~$x\in[0,1]$ et tout~$k$ assez grand. On en déduit que la suite $(P_k/Q_k)$
    converge uniformément vers $P/Q$ sur~$[0,1]$ et que $r_0=P/Q$ convient.
    
    Si $Q$ admet dans~$[0,1]$ des racines $a_1,\dots,a_p$ de multiplicités $\alpha_1,\dots,\alpha_p$,
    on note $Q^0 = \prod_i(X-a_i)^{\alpha_i}$ et $Q^1 = Q/Q^0$.
    Soit $M = \max\{\|g-f_k\|,\ k\in\N\}$. Pour tous $x\in[0,1]$ et $k\in\N$ on a
    $|g(x)Q_k(x)-P_k(x)|\le M|Q_k(x)|$ donc à la limite, $|g(x)Q(x)-P(x)|\le M|Q(x)|$
    pour tout~$x\in[0,1]$. Ceci implique que $Q^0$ divise~$P$, on note $P^1=P/Q^0$.
    Alors pour tout~$x\in[0,1]$ et $k\in\N$ on a $|g(x)Q^0(x)-P_k(x)Q^0(x)/Q_k(x)|\le \|g-f_k\||Q^0(x)|$,
    d'où $|g(x)Q^0(x)-P^1(x)Q^0(x)/Q^1(x)|\le d|Q^0(x)|$ et finalement $r_0=P^1/Q^1$
    convient.
\end{enumerate}
 \fincorrection
\correction{004562}
\begin{enumerate}
  \item 
  \item Soit $P_n$ le polynôme de Lagrange défini par $P_n(x_i) = f_n(x_i)$
    et $\deg P_n < p$.
    Les coordonnées de $P_n$ dans la base de Lagrange
    forment des suites convergentes donc la suite $(P_n)$ est uniformément
    convergente sur $[a,b]$. Quant à la suite $(P_n^{(p)})$, c'est la suite nulle.
    Donc on peut remplacer $f_n$ par $f_n - P_n$ dans l'énoncé, ce qui revient
    à supposer que $f_n(x_i) = 0$ pour tous $n$ et $i$. Soit $f$ la fonction
    définie par $f(x_i) = 0$ et $f^{(p)} = g$~: $f$ existe (prendre une primitive
    $p$-ème arbitraire de~$g$ et lui soustraire un polynôme de Lagrange approprié)
    et est unique (la différence entre deux solutions est polynomiale de degré $<p$
    et s'annule en $p$ points distincts). On remplace maintenant $f_n$ par $f_n-f$,
    et on est rammené à montrer que~: si $f_n(x_i) = 0$ pour tous $n$ et $i$
    et si $(f_n^{(p)})$ converge uniformément vers la fonction nulle, alors
    $(f_n)$ converge uniformément vers la fonction nulle. Ceci résulte du lemme
    suivant~:
    
    {\it Il existe une fonction $\varphi_p$ bornée sur $[a,b]^2$,
    indépendante de~$n$, telle que $f_n(x) =  \int_{t=a}^b \varphi_p(x,t)f_n^{(p)}(t)\,d t$.}
    
    Démonstration. On écrit la formule de Taylor-intégrale pour $f_n$ entre
    $x$ et $y$~:
    $$f_n(y) = f_n(x) + (y-x)f_n'(x) + \dots + \frac{(y-x)^{p-1}}{(p-1)!}f_n^{(p-1)}(x)
             +  \int_{t=x}^y \frac{(y-t)^{p-1}}{(p-1)!}f_n^{(p)}(t)\,d t.$$
    L'intégrale peut être étendue à l'intervalle $[a,b]$ sous la forme
    $ \int_{t=a}^b u_p(x,y,t)f_n^{(p)}(t)\,d t$ en posant
    $$u_p(x,y,t) = \begin{cases}(y-t)^{p-1}/(p-1)!  & \text{ si } x < t < y~;\cr
                          -(y-t)^{p-1}/(p-1)! & \text{ si } y < t < x~;\cr
                          0                   & \text{ sinon}.\cr\end{cases}$$
    En prenant successivement $y=x_1$, \ldots, $y=x_n$, on obtient un
    système linéaire en $f_n(x)$, \dots, $f_n^{(p-1)}(x)$ de la forme~:
    $$\left\{\begin{array}{lll} f_n(x) + (x_1-x)f_n'(x) + \dots + \frac{(x_1-x)^{p-1}}{(p-1)!}f_n^{(p-1)}(x) &= &- \int_{t=a}^b u_p(x,x_1,t)f_n^{(p)}(t)\,d t\cr
             \noalign{\smallskip}\vdots\cr
             f_n(x) + (x_p-x)f_n'(x) + \dots + \frac{(x_p-x)^{p-1}}{(p-1)!}f_n^{(p-1)}(x) &= &- \int_{t=a}^b u_p(x,x_p,t)f_n^{(p)}(t)\,d t\cr\end{array}\right.$$
    La matrice $M$ de ce système est la matrice de Vandermonde de $x_1-x,\dots,x_p-x$,
    inversible. On en déduit, avec les
    formules de Cramer, une expression de $f_n(x)$ à l'aide des intégrales du second membre,
    de la forme voulue. Le facteur $\varphi_p$ est borné car le dénominateur
    est $\det(M) = \prod_{i<j}(x_j-x_i)$, indépendant de~$x$.
\end{enumerate}
 \fincorrection
\correction{004563}
Développer en séries sous l'intégrale, multiplier, permuter
avec l'intégrale puis simplifier.
\fincorrection
\correction{005736}
Soit $x>0$. Pour $n\in\Nn^*$, $n^2e^{-x\sqrt{n}}=e^{-x\sqrt{n}+2\ln n}\underset{n\rightarrow+\infty}{=}o(1)$ d'après un théorème de croissances comparées. On en déduit que $e^{-x\sqrt{n}}\underset{n\rightarrow+\infty}{=}o\left(\frac{1}{n^2}\right)$ et donc que la série de terme général $e^{-x\sqrt{n}}$ converge. Ainsi, $f$ est bien définie sur $]0,+\infty[$.

Soit $x\in]0,+\infty[$. La fonction $t\mapsto e^{-x\sqrt{t}}$ est décroissante sur $[0,+\infty[$. Donc, $\forall k\in\Nn$, $\int_{k}^{k+1}e^{-x\sqrt{t}}\;dt\leqslant e^{-x\sqrt{k}}$ et $\forall k\in\Nn^*$, $e^{-x\sqrt{k}}\leqslant\int_{k-1}^{k}e^{-x\sqrt{t}}\;dt$. En sommant ces inégalités, on obtient

\begin{center}
$\forall x\in]0,+\infty[$, $\int_{0}^{+\infty}e^{-x\sqrt{t}}\;dt\leqslant f(x)\leqslant1+\int_{0}^{+\infty}e^{-x\sqrt{t}}\;dt$\quad$(*)$.
\end{center}

Soit $x\in]0,+\infty[$. En posant $u=x\sqrt{t}$ et donc $t=\frac{u^2}{x^2}$ puis $dt=\frac{2u}{x^2}\;du$, on obtient

\begin{center}
$\int_{0}^{+\infty}e^{-x\sqrt{t}}\;dt=\frac{2}{x^2}\int_{0}^{+\infty}ue^{-u}\;du=\frac{2}{x^2}\times\Gamma(2)=\frac{2}{x^2}$.
\end{center}

L'encadrement $(*)$ s'écrit alors

\begin{center}
$\forall x\in]0,+\infty[$, $\frac{2}{x^2}\leqslant f(x)\leqslant1+\frac{2}{x^2}$.
\end{center}

Comme $\displaystyle\lim_{\substack{x\rightarrow0\\
x>0}}\frac{2}{x^2}=+\infty$, on a montré que

\begin{center}
\shadowbox{
$\sum_{n=0}^{+\infty}e^{-x\sqrt{n}}\underset{x\rightarrow0,\;x>0}{\sim}\frac{2}{x^2}$.
}
\end{center}
\fincorrection
\correction{005737}
Soit $x\in]-1,1[$. Pour $n\in\Nn^*$, $\left|x^{n^2}\right|=|x|^{n^2}\leqslant|x|^n$. Puisque la série numérique de terme général $|x|^n$ converge, on en déduit que la série de terme général $x^{n^2}$ est absolument convergente et en particulier convergente. Donc, $f$ est bien définie sur $]-1,1[$.

Soit $x\in]0,1[$. La fonction $t\mapsto x^{t^2}=e^{t^2\ln x}$ est décroissante sur $[0,+\infty[$. Donc, $\forall k\in\Nn^*$, $\int_{k}^{k+1}x^{t^2}\;dt\leqslant x^{k^2}\leqslant\int_{k-1}^{k}x^{t^2}\;dt$. En sommant ces inégalités, on obtient

\begin{center}
$\forall x\in]0,1[$, $\int_{1}^{+\infty}x^{t^2}\;dt\leqslant f(x)\leqslant\int_{0}^{+\infty}x^{t^2}\;dt$\quad$(*)$.
\end{center}

Soit $x\in]0,1[$. En posant $u=t\sqrt{-\ln x}$, on obtient

\begin{center}
$\int_{0}^{+\infty}x^{t^2}\;dt=\int_{0}^{+\infty}e^{t^2\ln x}\;dt=\int_{0}^{+\infty}e^{-(t\sqrt{-\ln x})^2}\;dt=\frac{1}{\sqrt{-\ln x}}\int_{0}^{+\infty}e^{-u^2}\;du=\frac{\sqrt{\pi}}{2\sqrt{-\ln x}}$.
\end{center}

L'encadrement $(*)$ s'écrit alors

\begin{center}
$\forall x\in]0,1[$, $\frac{\sqrt{\pi}}{2\sqrt{-\ln x}}-\int_{0}^{1}x^{t^2}\;dt\leqslant f(x)\leqslant\frac{\sqrt{\pi}}{2\sqrt{-\ln x}}$.
\end{center}

Comme $\displaystyle\lim_{\substack{x\rightarrow1\\
x<1}}\frac{\sqrt{\pi}}{2\sqrt{-\ln x}}=+\infty$, on a montré que

\begin{center}
\shadowbox{
$\sum_{n=1}^{+\infty}x^{n^2}\underset{x\rightarrow1,\;x<1}{\sim}\frac{\sqrt{\pi}}{2\sqrt{-\ln x}}$.
}
\end{center}

\fincorrection
\correction{001784}

\begin{enumerate} 
\item 
Avec $f(x)=x^4$, d'o\`u $y= x^4-x$, on obtient
$\frac{x^2y}{x+y} =x^2-\frac 1x$ d'o\`u
\[
\mathrm{lim}_{\begin{smallmatrix}(x,y)\to (0,0)\\ x+y\ne 0\end{smallmatrix}} \frac{x^2y}{x+y}
\]
n'existe pas.

\item
$\mathrm{lim}_{\begin{smallmatrix}(x,y,z)\to (0,0,0)\\x=y=z\ne 0\end{smallmatrix}} \frac{xyz+z^3}{2x^3+yz^2} =\tfrac 23$ et
$\mathrm{lim}_{\begin{smallmatrix}(x,y,z)\to (0,0,0)\\x\ne 0, y=z=0\end{smallmatrix}} \frac{xyz+z^3}{2x^3+yz^2} =0$.
Il s'ensuit que
\[
\mathrm{lim}_{\begin{smallmatrix}(x,y,z)\to (0,0,0)\\2x^3+yz^2 \ne 0\end{smallmatrix}} \frac{xyz+z^3}{2x^3+yz^2}
\]
n'existe pas.

\item
Sur $\R\setminus \{0\}$, la fonction $f$ d\'efinie par
$f(x)=\frac{|x|}{x^2}=\frac{1}{|x|}$ tend vers $+\infty$ 
quand $x$ tend vers z\'ero d'o\`u 
\[
\mathrm{lim}_{\begin{smallmatrix}(x,y)\to (0,0)\\(x,y)\ne (0,0)\end{smallmatrix}} \frac{|x|+|y|}{x^2+y^2}
\]
n'existe pas en tant que limite finie.


\item
D'une part,
$\mathrm{lim}_{\begin{smallmatrix}(x,y)\to (0,0)\\ x \ne 0, y=0\end{smallmatrix}} \frac{x^4y}{x^2-y^2} = 0$.
D'autre part, vue l'indication, 
avec $x^2-y^2=h(y)$, un calcul imm\'ediat donne
\[
 \frac{x^4y}{x^2-y^2}=  \frac{y^5+2y^3h(y) + (h(y))^2y}{h(y)}=
\frac {y^5}{h(y)} + 2 y^3 + h(y)y .
\]
Avec $h(y)=y^6$, l'expression  $\frac{x^4y}{x^2-y^2}$ tend donc vers
$+\infty$ quand $y$ tend vers zero 
d'o\`u
\[
\mathrm{lim}_{\begin{smallmatrix}(x,y)\to (0,0)\\ x \ne \pm y\end{smallmatrix}} \frac{x^4y}{x^2-y^2}
\]
n'existe pas.

\item
Le long de la demi-droite $x>0,y=0,z=0$,
la limite existe et vaut z\'ero et
le long de la demi-droite $x=y=z>0$
la limite existe et vaut $1/3$ d'o\`u
\[
\mathrm{lim}_{\begin{smallmatrix}(x,y,z)\to (0,0,0)\\(x,y,z) \ne(0,0,0)\end{smallmatrix}} \frac{xy+yz}{x^2+2y^2+3z^2}
\]
n'existe pas.
\end{enumerate}

\fincorrection
\correction{001785}
\[ 
\lim_{y\to 0}f(x,y) =  \lim_{y\to 0}\frac{y^2}{y^2 + (1-y/x)^2} 
=\frac{\lim_{y\to 0} y^2}{\lim_{y\to 0} y^2 + (1-y/x)^2}
=\frac{0}{1} =0
\]
De m\^eme $\lim_{x\to 0}f(x,y) =0$ d'o\`u \eqref{non}.
D'autre part,
$f(x,x) =  \frac{x^4}{x^4} =1$ d'o\`u
$\lim_{x\to 0}f(x,x) = 1$ et
$\lim_{(x,y)\to (0,0)} f(x,y)$ ne peut pas exister.

\fincorrection
\nocorrection
\correction{001787}

\begin{enumerate}  
\item
$\lim_{(x,y)\to (0,0),y=0} \frac{x}{x^2+y^2}$ n'existe pas
d'o\`u $\lim_{(x,y)\to (0,0)} \frac{x}{x^2+y^2}  $ n'existe pas.

\item
$\frac{(x+2y)^3}{x^2+y^2} =r(\cos \varphi +2 \sin \varphi)^3$ d'o\`u
$\left|\frac{(x+2y)^3}{x^2+y^2}\right| \leq 27 r$ et
\[
\lim_{(x,y)\to (0,0)} \frac{(x+2y)^3}{x^2+y^2} =0
\]
car $\lim_{(x,y)\to (0,0)} r=0$.

\item  
$\lim_{(x,y)\to (1,0)} {\sqrt{x^2+y^2}} =1 \ne 0$ et
$\lim_{(x,y)\to (1,0)} {\log (x+e^y)}=\log 2$  d'o\`u
\[
\lim_{(x,y)\to (1,0)} \frac{\log (x+e^y)}{\sqrt{x^2+y^2}}=\log 2. 
\]

\item 
$\mathrm{lim}_{\begin{smallmatrix}(x,y)\to (0,0)\\y=2\end{smallmatrix}} \frac{x^4+y^3-xy}{x^4+y^2} = 1$ tandis que
$\mathrm{lim}_{\begin{smallmatrix}(x,y)\to (0,0)\\x=0\end{smallmatrix}} \frac{x^4+y^3-xy}{x^4+y^2} =0$ 
 d'o\`u
\[
\lim_{(x,y)\to (0,0)} \frac{x^4+y^3-xy}{x^4+y^2} 
\]
n'existe pas.
\end{enumerate}
\fincorrection
\correction{001788} 
\begin{enumerate} 
\item Supposons $x+y+z \ne 0$. Alors
\[
\frac{xyz}{x+y+z} =\frac{xy(h(x,y)-x-y)}{h(x,y)}
=xy -\frac{xy(x+y)}{h(x,y)}
\]
 d'o\`u, avec
\[h(x,y)=(x+y)^4,
\]
nous obtenons
\[
\frac{xyz}{x+y+z} 
=xy -\frac{xy}{(x+y)^3} .
\]
Il s'ensuit que
\[
\mathrm{lim}_{\begin{smallmatrix}(x,y,z)\to (0,0,0)\\x+y+z=(x+y)^4\\
x\ne 0, y\ne  0, z \ne 0\end{smallmatrix}}\frac{xyz}{x+y+z} 
\]
n'existe pas, au moins non pas en tant que limite finie.
D'autre part,
\[
\mathrm{lim}_{\begin{smallmatrix}(x,y,z)\to (0,0,0)\\
x+z\ne 0,y= 0\end{smallmatrix}}\frac{xyz}{x+y+z} =0.
\]
Par cons\'equent,
\[
\mathrm{lim}_{\begin{smallmatrix}(x,y,z)\to (0,0,0)\\
x+y+z\ne 0\end{smallmatrix}}\frac{xyz}{x+y+z}
\]
ne peut pas exister.

\item La limite
\[
\mathrm{lim}_{\begin{smallmatrix}(x,y,z) \to (0,0,0)\\ x\ne \pm y, z = 0\end{smallmatrix}}f(x,y,z)= 
\mathrm{lim}_{\begin{smallmatrix}(x,y) \to (0,0)\\
x \ne y\end{smallmatrix}}\frac 1{x-y}
\]
n'existe pas car $\mathrm{lim}_{\begin{smallmatrix}(x,y) \to (0,0)\\
y=x-x^2\end{smallmatrix}}\frac 1{x-y}$ n'existe pas.
Par cons\'equent,
\[
\mathrm{lim}_{\begin{smallmatrix}(x,y,z) \to (0,0,0)\\x^2-y^2+z^2\ne 0\end{smallmatrix}}f(x,y,z)
= 
\mathrm{lim}_{\begin{smallmatrix}(x,y,z) \to (0,0,0)\\x^2-y^2+z^2\ne 0\end{smallmatrix}}\frac{x+y}{x^2-y^2+z^2}
\]
ne peut pas exister.
\end{enumerate}

\fincorrection
\nocorrection
\nocorrection
\nocorrection
\nocorrection
\correction{002621} 
Des calculs \'el\'ementaires donnent

\begin{enumerate}
  \item $u_1=(\frac{1}{2}, \cos 1),\ u_2=(\frac{16}{15}, \cos \frac 12),\ 
\ldots \ ,
 u_{10}=(\frac{400}{143}, \cos \frac 1{10}),\ 
\ldots$
  \item 
  $u_1=(\frac{1}{2}\arctan 1,   \sin(\frac{\pi}{4\mathrm e     })),\ 
   u_2=(\frac{4}{5}\arctan 2,  \sin(\frac{\pi}{4\mathrm e^{1/2}})),\ \\
   u_3=(\frac{9}{10}\arctan 3, \sin(\frac{\pi}{4\mathrm e^{1/3}})),\ \ldots \ ,
u_{10}=(\frac{100}{101}\arctan (10), \sin(\frac{\pi}{4\mathrm e^{1/10}})),\ \ldots \ 
$
  \item 
  $u_1=(\sinh 1, 0), \ 
   u_2=(\sinh 2, \frac{\ln 2}{2}), \ 
   u_3=(\sinh 3, \frac{\ln 3}{3}), \ \ldots \ ,\\ 
u_{10}=(\sinh 10, \frac{\ln 10}{10}),\ 
\ldots $
  \item 
  $u_1=a^n   (\cos   (\alpha), \sin  (\alpha) ),\ 
   u_2=a^2   (\cos  (2\alpha), \sin  (2\alpha)),  \\
   u_3=a^3   (\cos  (3\alpha), \sin  (3\alpha)),\   \ldots \ , 
u_{10}=a^{10}(\cos (10\alpha), \sin (10\alpha)),\ \ldots
$
\end{enumerate}

Les limites pouvu qu'elles existent se calculent ainsi :

\begin{enumerate}
\item 
$
\lim_{n \to \infty}\frac{4n^2}{n^2+4n+3}
=\lim_{n \to \infty}\frac{4}{1+\frac 4n+ \frac 3{n^2}}=4,
\
\lim_{n \to \infty}\cos(1/n) = \cos (0)=1
$
d'o\`u
\[
\lim_{n \to \infty}u_n=\lim_{n \to \infty}
\left(\frac{4n^2}{n^2+4n+3}, \cos\frac1n\right)= (4,0) .
\]
\item 
$
\lim_{n \to \infty}\frac{n^2}{n^2+1}
=\lim_{n \to \infty}\frac{1}{1+1/n^2}=1,
\
\lim_{n \to \infty}\arctan n = \pi/2,
\\
\lim_{n \to \infty}
\frac{n^2\arctan n}{n^2+1}= \pi/2
$
mais
$
\lim_{n \to \infty}\sin(\frac{\pi}{4}\exp(-{\frac 1n}))
$
n'existe pas d'o\`u 
\[
\lim_{n \to \infty}u_n
\]
n'existe pas.


  \item $
\lim_{n \to \infty}\frac{\ln n}{n} = 0$ tandis que
$\lim_{n \to \infty}\sinh n$ n'existe pas en tant que limite finie
car 
\[
\lim_{x \to +\infty}\sinh x = +\infty
\]
d'o\`u
\[
\lim u_n=\lim\left(\sinh n, \frac{\ln n}{n}\right)
\]
n'existe pas.
\item $\lim_{n \to \infty} (\cos  (n\alpha), \sin  (n\alpha))$
n'existe pas tandis que
pour que $\lim_{n \to \infty}a^n$ existe il faut et il suffit que $a\leq 1$ et,
s'il en est ainsi, $\lim_{n \to \infty}a^n =0$ si $a< 1$
et $\lim_{n \to \infty}a^n =1$ si $a= 1$.
Par cons\'equent: 
Pour que 
\[
\lim_{n \to \infty}u_n=\lim_{n \to \infty}a^n(\cos  (n\alpha), \sin  (n\alpha))
\]
existe il faut et il suffit que $a<1$, et la limite vaut alors z\'ero.
\end{enumerate}
\fincorrection
\nocorrection
\correction{005553}
On note $f$ la fonction considérée.
\begin{enumerate}
 \item  Pour $x\neq0$, $f(x,-x+x^3)=\frac{x(-x+x^3)}{x-x+x^3}\underset{x\rightarrow0^+}{\sim}-\frac{1}{x}$. Quand $x$ tend vers $0$, $-x+x^3$ tend vers $0$ puis
$\displaystyle\lim_{\substack{(x,y)\rightarrow(0,0)\\x>0,\;y=-x+x^3}}f(x,y)=-\infty$. $f$ n'a de limite réelle en $(0,0)$.
 \item  Pour $x\neq0$, $f(x,0)=\frac{x\times0}{x^2+0^2}=0$ puis $\displaystyle\lim_{\substack{(x,y)\rightarrow(0,0)\\y=0}}f(x,y)=0$. Mais aussi, pour $x\neq0$, $f(x,x)=\frac{x\times x}{x^2+x^2}=\frac{1}{2}$ puis $\displaystyle\lim_{\substack{(x,y)\rightarrow(0,0)\\x=y}}f(x,y)=\frac{1}{2}$.
Donc si $f$ a une limite réelle, cette limite doit être égale à $0$ et à $\frac{1}{2}$ ce qui est impossible. $f$ n'a pas de limite réelle en $(0,0)$.
 \item  Pour tout $(x,y)\in\Rr^2$, $x^2-2|xy|+y^2=(|x|-|y|)^2\geqslant0$ et donc $|xy|\leqslant\frac{1}{2}(x^2+y^2)$.
Par suite, pour $(x,y)\neq(0,0)$,

\begin{center}
$|f(x,y)|=\frac{x^2y^2}{x^2+y^2}\leqslant\frac{(x^2+y^2)^2}{4(x^2+y^2)}=\frac{1}{4}(x^2+y^2)$.
\end{center}
Comme $\lim_{(x,y)\rightarrow (0,0)}\frac{1}{4}(x^2+y^2)=0$, on a aussi $\lim_{(x,y)\rightarrow (0,0)}f(x,y)=0$.

 \item   $\lim_{(x,y)\rightarrow (0,0)}\frac{\sin y}{y}=1$ et $\lim_{(x,y)\rightarrow (0,0)}(1+x^2+y^2)=1$. Donc $\lim_{(x,y)\rightarrow (0,0)}f(x,y)=1$.
 \item  Pour $(x,y)\in\Rr^2$, $|x^3+y^3|=|x+y|(x^2+xy+y^2)\leqslant\frac{3}{2}|x+y|(x^2+y^2)$ et donc pour $(x,y)\neq(0,0)$, 

\begin{center}
$|f(x,y)|=\frac{|x^3+y^3|}{x^2+y^2}\leqslant\frac{3}{2}|x+y|$.
\end{center}
Comme $\lim_{(x,y)\rightarrow (0,0)}\frac{3}{2}|x+y|=0$, on a aussi $\lim_{(x,y)\rightarrow (0,0)}f(x,y)=0$.
 \item  Pour $(x,y)\in\Rr^2$, $|x^4+y^4|=(x^2+y^2)^2-2x^2y^2\leqslant(x^2+y^2)^2+2\times\left(\frac{1}{2}(x^2+y^2)\right)^2=\frac{3}{2}(x^2+y^2)^2$ et donc pour $(x,y)\neq(0,0)$, 

\begin{center}
$|f(x,y)|=\frac{|x^4+y^4|}{x^2+y^2}\leqslant\frac{3}{2}(x^2+y^2)$.
\end{center}
Comme $\lim_{(x,y)\rightarrow (0,0)}\frac{3}{2}(x^2+y^2)=0$, on a aussi $\lim_{(x,y)\rightarrow (0,0)}f(x,y)=0$.
\end{enumerate}
\fincorrection
\correction{005887}
\begin{enumerate}
 \item  $f$ est définie sur $\Rr^2\setminus\{(0,0)\}$.

Pour $x\neq0$, $f(x,0)=0$. Quand $x$ tend vers $0$, le couple $(x,0)$ tend vers le couple $(0,0)$ et $f(x,0)$ tend vers $0$. Donc, si $f$ a une limite réelle en $0$, cette limite est nécessairement $0$.

Pour $x\neq0$, $f(x,x)= \frac{1}{2}$. Quand $x$ tend vers $0$, le couple $(x,x)$ tend vers $(0,0)$ et $f(x,x)$ tend vers $ \frac{1}{2}\neq0$. Donc $f$ n'a pas de limite réelle en $(0,0)$.

\item  $f$ est définie sur $\Rr^2\setminus\{(0,0)\}$.

Pour $(x,y)\neq(0,0)$, $|f(x,y)|= \frac{x^2y^2}{x^2+y^2}= \frac{|xy|}{x^2+y^2}\times |xy|\leqslant \frac{1}{2}|xy|$. Comme $ \frac{1}{2}|xy|$ tend vers $0$ quand le couple $(x,y)$ tend vers le couple  $(0,0)$, il en est de même de $f$. $f(x,y)$ tend vers $0$ quand $(x,y)$ tend vers $(0,0)$.

\item  $f$ est définie sur $\Rr^2\setminus\{(0,0)\}$.

Pour $y\neq0$, $f(0,y)= \frac{y^3}{y^4}= \frac{1}{y}$. Quand $y$ tend vers $0$ par valeurs supérieures, le couple $(0,y)$ tend vers le couple $(0,0)$ et $f(0,y)$ tend vers $+\infty$. Donc $f$ n'a pas de limite réelle en $(0,0)$.

\item  $f$ est définie sur $\Rr^2\setminus\{(0,0)\}$.

Pour $x\neq0$, $f(x,x)= \frac{\sqrt{2x^2}}{2|x|\sqrt{|x|}}= \frac{1}{\sqrt{2|x|}}$.Quand $x$ tend vers $0$, le couple $(x,x)$ tend vers le couple $(0,0)$ et $f(x,x)$ tend vers $+\infty$. Donc $f$ n'a pas de limite réelle en $(0,0)$.

\item  $f$ est définie sur $\Rr^2\setminus\{(x,-x),\;x\in\Rr\}$.

Pour $x\neq0$, $f(x,-x+x^3)= \frac{(x+x^2-x^3)(-x+(-x+x^2)^2)}{x^3}\underset{x\rightarrow0}{\sim}- \frac{1}{x}$. Quand $x$ tend vers $0$ par valeurs supérieures, le couple $(x,-x+x^3)$ tend vers $(0,0$ et $f(x,-x+x^3)$ tend vers $-\infty$. Donc $f$ n'a pas de limite réelle en $(0,0)$.

\item  $f$ est définie sur $\Rr^2\setminus\{(x,0),\;x\in\Rr\}$.

$ \frac{1-\cos\sqrt{|xy|}}{|y|}\underset{(x,y)\rightarrow(0,0)}{\sim} \frac{(\sqrt{|xy|})^2}{2|y|}= \frac{|x|}{2}$ et donc $f$ tend vers $0$ quand $(x,y)$ tend vers $(0,0)$.

\item  $f$ est définie sur $\Rr^3$ privé du cône de révolution d'équation $x^2-y^2+z^2=0$.

$f(x,0,0)= \frac{1}{x}$ qui tend vers $+\infty$ quand $x$ tend vers $0$ par valeurs supérieures. Donc $f$ n'a pas de limite réelle en $(0,0,0)$.

\item  $f(2+h,-2+k,l)= \frac{h+k}{h^2-k^2+l^2+4h+4k}=g(h,k,l)$. $g(h,0,0)$ tend vers $ \frac{1}{4}$ quand $h$ tend vers $0$ et $g(0,0,l)$ tend vers $0\neq \frac{1}{4}$ quand $l$ tend vers $0$. Donc, $f$ n'a pas de limite réelle quand $(x,y,z)$ tend vers $(2,-2,0)$.
\end{enumerate}
\fincorrection
\nocorrection
\nocorrection
\nocorrection
\nocorrection
\nocorrection
\nocorrection
\correction{005554}
Déterminons tout d'abord $F(x,y)$ pour $(x,y)\in\Rr^2$.
\textbullet~Pour $y\in\Rr$, $F(x,y)=\text{Max}\left\{f_{0,y}(-1),f_{0,y}(1)\right\}=\text{Max}\left\{y,-y\right\}=|y|$.
\textbullet~Si $x\neq0$, $F(x,y)=\text{Max}\left\{f_{x,y}(-1),f_{x,y}\left(-\frac{y}{2x}\right),f_{x,y}(1)\right\}=\text{Max}\left\{x+y,x-y,-\frac{y^2}{4x}\right\}=\text{Max}\left\{x+|y|,-\frac{y^2}{4x}\right\}$.
Plus précisément, si $x>0$, on a $x+|y|>0$ et $-\frac{y^2}{4x}\leqslant0$. Donc $F(x,y)=x+|y|$ ce qui reste vrai quand $x=0$.
Si $x<0$, $x+|y|-\left(-\frac{y^2}{4x}\right)=\frac{4x^2+4x|y|+y^2}{4x}=\frac{(2x+|y|)^2}{4x}<0$ et donc $F(x,y)=-\frac{y^2}{4x}$.

\begin{center}
\shadowbox{
$\forall(x,y)\in\Rr^2,\;F(x,y)=\left\{
\begin{array}{l}
x+|y|\;\text{si}\;x\geqslant0\\
-\frac{y^2}{4x}\;\text{si}\;x<0
\end{array}
\right.$.
}
\end{center}
En vertu de théorèmes généraux, $F$ est continue sur $\{(x,y)\in\Rr^2,\;x>0\}$ et $\{(x,y)\in\Rr^2,\;x<0\}$.
Soit $y_0\neq0$. $\displaystyle\lim_{\substack{(x,y)\rightarrow(0,y_0)\\
x<0,\;y=y_0}}F(x,y)=+\infty\neq|y_0|=F(0,y_0)$ et donc $F$ n'est pas continue en $(0,y_0)$.
Enfin, $\displaystyle\lim_{\substack{(x,y)\rightarrow(0,0)\\
x<0,\;y=\sqrt{-x}}}F(x,y)=\frac{1}{4}\neq 0=F(0,0)$ et donc $F$ n'est pas continue en $(0,0)$.

\begin{center}
\shadowbox{
$F$ est continue sur $\Rr^2\setminus\{(0,y),\;y\in\Rr\}$ et est discontinue en tout $(0,y)$, $y\in\Rr$.
}
\end{center}
\fincorrection
\correction{005900}
\textbullet~Pour tout $x\in E$, $\|f(x)\|= \frac{\|x\|}{1+\|x\|}< \frac{\|x\|+1}{\|x\|+1}=1$. Donc $f$ est bien une application de $E$ dans $B$.

\textbullet~Si $y=0$, pour $x\in E$, $f(x)=y\Leftrightarrow \frac{1}{1+\|x\|}x=0\Leftrightarrow x=0$.

Soit alors $y\in B\setminus\{0\}$. Pour $x\in E$,

\begin{center}
$f(x)=y\Rightarrow x=(1+\|x\|)y\Rightarrow\exists\lambda\in\Kk/\;x=\lambda y$.
\end{center}

Donc un éventuel antécédent de $y$ est nécessairement de la forme $\lambda y$, $\lambda\in\Rr$. Réciproquement, pour $\lambda\in\Rr$, $f(\lambda y)= \frac{\lambda}{1+|\lambda|\|y\|}y$ et donc

\begin{align*}\ensuremath
f(\lambda y)=y&\Leftrightarrow \frac{\lambda}{1+|\lambda|\|y\|}=1\Leftrightarrow\lambda=1+|\lambda|\|y\|\\
 &\Leftrightarrow(\lambda\geqslant0\;\text{et}\;(1-\|y\|)\lambda=1)\;\text{ou}\;(\lambda<0\;\text{et}\;(1+\|y\|)\lambda=1)\\
 &\Leftrightarrow\lambda= \frac{1}{1-\|y\|}\;(\text{car}\;\|y\|<1).
\end{align*}

Dans tous les cas, $y$ admet un antécédent par $f$ et un seul à savoir $x= \frac{1}{1-\|y\|}y$. Ainsi,

\begin{center}
$f$ est bijective et $\forall x\in B$, $f^{-1}(x)= \frac{1}{1-\|x\|}x$.
\end{center}

\textbullet~On sait que l'application $x\mapsto\|x\|$ est continue sur $\Rr^2$. Donc l'application $x\mapsto \frac{1}{1+\|x\|}$ est continue sur $\Rr^2$ en tant qu'inverse d'une fonction continue sur $\Rr^2$ à valeurs dans $\Rr$, ne s'annulant pas sur $\Rr^2$. L'application $x\mapsto \frac{1}{1-\|x\|}$ est continue sur $B$ pour les mêmes raisons. Donc les applications $f$ et $f^{-1}$ sont continues sur $\Rr^2$ et $B$ respectivement et on a montré que

\begin{center}
\shadowbox{
l'application $\begin{array}[t]{cccc}
f~:&E&\rightarrow&B\\
 &x&\mapsto& \frac{x}{1+\|x\|}
\end{array}$ est un homéomorphisme.
}
\end{center}
\fincorrection
\correction{005901}
\textbf{1ère solution.} Pour $x=(x_1,\ldots,x_n)\in\Rr^n$, $f(x)=\sqrt{\sum_{i=1}^{n}x_i^2}$. $f$ est de classe $C^1$ sur $\Rr^n\setminus\{0\}$ en vertu de théorèmes généraux et pour tout $x=(x_1,\ldots,x_n)\in\Rr^n\setminus\{0\}$ et tout $i\in\llbracket1,n\rrbracket$

\begin{center}
$ \frac{\partial f}{\partial x_i}(x)= \frac{x_i}{\sqrt{\sum_{i=1}^{n}x_i^2}}= \frac{x_i}{\|x\|_2}$.
\end{center}

On en déduit que $f$ est différentiable sur $\Rr^n\setminus\{0\}$ et pour $x\in\Rr^n\setminus\{0\}$ et $h\in\Rr^n$

\begin{center}
$df_x(h)=\sum_{i=1}^{n} \frac{\partial f}{\partial x_i}(x)h_i= \frac{1}{\|x\|_2}\sum_{i=1}^{n}x_ih_i= \frac{x|h}{\|x\|_2}$.
\end{center}

\begin{center}
\shadowbox{
$\forall x\in\Rr^n\setminus\{0\}$, $\forall h\in\Rr^n$, $df_x(h)= \frac{x|h}{\|x\|_2}$.
}
\end{center}

\textbf{2 ème solution.} Soit $x\in\Rr^n\setminus\{0\}$. Pour $h\in\Rr^n$,

\begin{center}
$\|x+h\|_2-\|x\|_2= \frac{\left(\|x+h\|_2-\|x\|_2\right)\left(\|x+h\|_2+\|x\|_2\right)}{\|x+h\|_2+\|x\|_2}= \frac{2(x|h)+\|h\|_2^2}{\|x+h\|_2+\|x\|_2}$,
\end{center}

puis

\begin{center}
$\|x+h\|_2-\|x\|_2- \frac{x|h}{\|x\|_2}= \frac{2(x|h)+\|h\|_2^2}{\|x+h\|_2+\|x\|_2}- \frac{x|h}{\|x\|_2}= \frac{-\left(\|x+h\|_2-\|x\|_2\right)(x|h)+\|x\|_2\|h\|_2^2}{\left(\|x+h\|_2+\|x\|_2\right)\|x\|_2}$.
\end{center}

Maintenant, on sait que l'application $x\mapsto\|x\|_2$ est continue sur $\Rr^n$. On en déduit que $ \frac{1}{\left(\|x+h\|_2+\|x\|_2\right)\|x\|_2}\underset{h\rightarrow0}{\sim} \frac{1}{2\|x\|_2^2}$ et aussi que $\|x+h\|_2-\|x\|_2$ tend vers $0$ quand $h$ tend vers $0$. Ensuite, puisque $\left|(x|h)\right|\leqslant\|x\|_2\|h\|_2$ (inégalité de \textsc{Cauchy}-\textsc{Schwarz}), on a $x|h\underset{h\rightarrow0}{=}O(\|h\|_2)$ puis $\left(\|x+h\|_2-\|x\|_2\right)\left(x|h\right)\underset{h\rightarrow0}{=}o(\|h\|_2)$.

Finalement, $ \frac{-\left(\|x+h\|_2-\|x\|_2\right)(x|h)+\|x\|_2\|h\|_2^2}{\left(\|x+h\|_2+\|x\|_2\right)\|x\|_2}\underset{h\rightarrow0}{=}o(\|h\|_2)$ et donc

\begin{center}
$\|x+h\|_2\underset{h\rightarrow0}{=}\|x\|_2+ \frac{x|h}{\|x\|_2}+o(\|h\|_2)$.
\end{center}

Puisque l'application $h\mapsto \frac{x|h}{\|x\|_2}$ est linéaire, on a redémontré que $f$ est différentiable en tout $x$ de $\Rr^n\setminus\{0\}$ et que $\forall x\in\Rr^n\setminus\{0\}$, $\forall h\in\Rr^n$, $df_x(h)= \frac{x|h}{\|x\|_2}$.

Soit $L$ une application linéaire de $\Rr^n$ dans $\Rr$ c'est-à-dire une forme linéaire.

\begin{center}
$ \frac{1}{\|h\|_2}\left(\|0+h\|_2-\|0\|_2-L(h)\right)=1-L\left( \frac{h}{\|h\|_2}\right)$.
\end{center}

Supposons que cette expression tende vers $0$ quand $h$ tend vers $0$. Pour $u$ vecteur non nul donné et $t$ réel non nul, l'expression $1-L\left( \frac{tu}{\|tu\|_2}\right)=1- \frac{t}{|t|}L\left( \frac{u}{\|u\|_2}\right)$ tend donc vers $0$ quand $t$ tend vers $0$. Mais si $t$ tend vers $0$ par valeurs supérieures, on obtient $L(u)=\|u\|_2$ et si $t$ tend vers $0$ par valeurs inférieures, on obtient $L(u)=-\|u\|_2$ ce qui est impossible car $u\neq0$. Donc $f$ n'est pas différentiable en $0$.
\fincorrection
\nocorrection
\nocorrection
\correction{001800}
Puisque $\left|\frac{x^{2}-y^{2}}{x^{2} + y^{2}}\right|$ reste born\'e,
\[
\mathrm{lim}_{(x,y)\to (0,0)}f(x,y)= \mathrm{lim}_{(x,y)\to (0,0)}xy\frac{x^{2}-y^{2}}{x^{2} + y^{2}}
= \mathrm{lim}_{(x,y)\to (0,0)}xy =0
\]
d'o\`u $f$ est continue en $(0,0)$.
De m\^eme,
\begin{align*}
\mathrm{lim}_{(x,y)\to (0,0)}\frac{f(x,y)}x&= \mathrm{lim}_{(x,y)\to (0,0)}y\frac{x^{2}-y^{2}}{x^{2} + y^{2}}
= \mathrm{lim}_{(x,y)\to (0,0)}y =0
\\
\mathrm{lim}_{(x,y)\to (0,0)}\frac{f(x,y)}y&= \mathrm{lim}_{(x,y)\to (0,0)}x\frac{x^{2}-y^{2}}{x^{2} + y^{2}}
= \mathrm{lim}_{(x,y)\to (0,0)}x =0
\end{align*}
d'o\`u les d\'eriv\'ees partielles
$\frac{\partial f}{\partial x}(0,0)$ et $\frac{\partial f}{\partial y}(0,0)$
existent, et
$\frac{\partial f}{\partial x}(0,0)=0$ et $\frac{\partial f}{\partial y}(0,0)=0$.
En plus, en dehors de l'origine,
\begin{align*}
\frac{\partial f}{\partial x}&= \frac{f(x,y)}x + xy \frac{\partial}{\partial x}
\frac{x^{2}-y^{2}}{x^{2} + y^{2}} = \frac{f(x,y)}x +  4\frac{x^{2}y^{3}}{(x^{2} + y^{2})^2}
\\
\frac{\partial f}{\partial y}&= \frac{f(x,y)}y +  xy \frac{\partial}{\partial y}
\frac{x^{2}-y^{2}}{x^{2} + y^{2}}  = \frac{f(x,y)}x +  4\frac{x^{3}y^{2}}{(x^{2} + y^{2})^2} .
\end{align*}
Puisque
\[
\mathrm{lim}_{(x,y) \to (0,0)}\frac{x^{2}y^{3}}{(x^{2} + y^{2})^2}=0,\ 
\mathrm{lim}_{(x,y) \to (0,0)}\frac{x^{3}y^{2}}{(x^{2} + y^{2})^2}=0,
\]
il s'ensuit que
\[
\mathrm{lim}_{(u,v) \to (0,0)}\frac{\partial f}{\partial x}(u,v) =0,\ 
\mathrm{lim}_{(u,v) \to (0,0)}\frac{\partial f}{\partial y}(u,v) =0,
\]
d'o\`u les d\'eriv\'ees partielles $\frac{\partial f}{\partial x}$ et
$\frac{\partial f}{\partial y}$ sont continues en $(0,0)$.

\fincorrection
\correction{001801}
\begin{align*}
\frac{\partial g}{\partial x}(x,y)&= f'(x+y)
\\
\frac{\partial g}{\partial y}(x,y)&= f'(x+y)
\\
\frac{\partial h}{\partial x}(x,y)&= 2x f'(x^{2} + y^{2})
\\
\frac{\partial h}{\partial y}(x,y)&= 2y f'(x^{2} + y^{2})
\\
\frac{\partial k}{\partial x}(x,y)&= yf'(xy)
\\
\frac{\partial k}{\partial y}(x,y)&= xf'(xy)
\end{align*}
\fincorrection
\nocorrection
\correction{001803}
Il est \'evident que, en tout point tel que  $\left|x\right|<\left|y\right|$
ou   $\left|x\right|>\left|y\right|$,
la fonction est continue et les d\'eriv\'ees
partielles existent.

Soit $x \ne 0$. Alors $f$ n'est ni continue en $(x,x)$ ni en $(x,-x)$.
Car
\begin{align*}
\mathrm{lim}_{\begin{smallmatrix} (u,v) \to (x,x)\\ |u|>|v| \end{smallmatrix}}f(u,v)
&= \mathrm{lim}_{u \to x}u = x \ne 0,
\\
\mathrm{lim}_{(u,u) \to (x,x)} f(u,u) &=0,
\\
\mathrm{lim}_{\begin{smallmatrix} (u,v) \to (x,-x)\\ |u|>|v| \end{smallmatrix}}f(u,v)
&= \mathrm{lim}_{u \to x}u = x \ne 0,
\\
\mathrm{lim}_{(u,-u) \to (x,-x)} f(u,u) &=0.
\end{align*}
Par contre, $f$ est continue en $(0,0)$. Car
\[
\mathrm{lim}_{(u,v) \to (0,0)}f(u,v)=0
\]
puisque
\[
\begin{array}{lll}
 f(u, v)=u &\mathrm{si}&
\left|u\right|>\left|v\right|,\\ 
f(u, v)=v &\mathrm{ si }&\left|u\right|<\left|v\right|,\\ 
f(u, v)= 0 &\mathrm{ si }&
\left|u\right| = \left|v\right|, 
\end{array}
\]
et puisque alors $\mathrm{lim}_{u \to 0}u=0$ et $\mathrm{lim}_{v \to 0}v=0$.

Soit  $(x,y)$ un point o\`u 
$\left|x\right| = \left|y\right|$.
Il reste \`a \'etudier les d\'eriv\'ees partielles en un tel point $(x,y)$.
Soit $x \ne 0$.
Alors la fonction $h$ de la variable $t$ d\'efinie par
\[
h(t)=f(x+t,y)=\begin{cases} x+t,\quad & |x+t|>|y|\\ y,\quad & |x+t|<|y| \end{cases}
\]
n'est pas d\'erivable en $t=0$ donc la d\'eriv\'ee partielle
$\frac{\partial f}{\partial x}(x,y)$ n'existe pas.
De m\^eme,
la fonction $k$ de la variable $t$ d\'efinie par
\[
k(t)=f(x,y+t)=\begin{cases} x,\quad & |x|>|y+t|,\\ y+t,\quad & |x|<|y+t|, \end{cases}
\]
n'est pas d\'erivable en $t=0$ donc la d\'eriv\'ee partielle
$\frac{\partial f}{\partial y}(x,y)$ n'existe pas.
Enfin soit $x = 0$.
Alors la fonction $h$ de la variable $t$ d\'efinie par
\[
h(t)=f(t,0)=t
\]
est d\'erivable en $t=0$ donc la d\'eriv\'ee partielle
$\frac{\partial f}{\partial x}(0,0)$ existe.
De m\^eme,
la fonction $k$ de la variable $t$ d\'efinie par
\[
k(t)=f(0,t)=t
\]
est d\'erivable en $t=0$ donc la d\'eriv\'ee partielle
$\frac{\partial f}{\partial y}(0,0)$ existe.
\fincorrection
\nocorrection
\nocorrection
\nocorrection
\nocorrection
\nocorrection
\nocorrection
\nocorrection
\nocorrection
\nocorrection
\nocorrection
\nocorrection
\nocorrection
\nocorrection
\nocorrection
\nocorrection
\nocorrection
\nocorrection
\nocorrection
\nocorrection
\correction{002628}
\begin{enumerate} 
\item Le plan tangent \`a la surface d'\'equation 
$z^2=19-x^2-y^2$ au point $(x_0,y_0,z_0)$ 
est donn\'e par l'\'equation
\[
2z_0(z-z_0)= -2x_0(x-x_0) -2y_0(y-y_0)
\]
d'o\`u, au point $(1,3,3)$, cette \'equation s'\'ecrit
\[
6(z-3)= -2(x-1) -6(y-3)  
\]
ou
\[
x+3y+3z= 19
\] 
\item Soit $f$ la fonction d\'efinie par 
$f(x,y)=\sin(\pi xy)\exp(2x^2y-1)$.
Les d\'eriv\'ees partielles de $f$ sont
\begin{align*}
\frac{\partial f}{\partial x}&=
\pi y\cos (\pi xy)\exp(2x^2y-1)+4xy\sin(\pi xy)\exp(2x^2y-1)
\\
\frac{\partial f}{\partial y}&=
\pi x\cos (\pi xy)\exp(2x^2y-1)+2x^2\sin(\pi xy)\exp(2x^2y-1)
\end{align*}
d'o\`u
\[
\frac{\partial f}{\partial x}(1,1/2)= 2,
\quad
\frac{\partial f}{\partial y}(1,1/2)= 2.
\]
Le plan tangent \`a la surface d'\'equation 
$z=\sin(\pi xy)\exp(2x^2y-1)$ au point $(x_0,y_0,z_0)$ 
est donn\'e par l'\'equation
\[
z-z_0=\frac{\partial f}{\partial x}(x_0,y_0) (x-x_0) 
+\frac{\partial f}{\partial y}(x_0,y_0)
(y-y_0)
\]
d'o\`u, au point $(1,1/2,1)$, cette \'equation s'\'ecrit
\[
z-1= 2  (x-1) +2(y-1/2) 
\]
ou
\[
 2x+2y-z =2.
\]
\end{enumerate}

\fincorrection
\correction{002629}
\begin{enumerate} 
\item L'\'equation d'un plan tangent doit \^etre une \'equation lin\'eaire !

\item La confusion est exactement celle \`a \'eviter suivant les indications
donn\'ees.


\item D'apr\`es \eqref{tang1}, le plan tangent \`a la surface d'\'equation
$z=f(x,y)=x^4-y^2$ au point $(x_0,y_0,z_0)=(2,3,7)$ est donn\'e 
par l'\'equation 
\[
z-7=\frac{\partial f}{\partial x}(2,3) (x-2) 
+\frac{\partial f}{\partial y}(2,3)
(y-3)
\]
c.a.d.
\[
z-7=32 (x-2) -6 (y-3) .
\]
\end{enumerate}
\fincorrection
\correction{002630}
Suivant l'indication,
le plan tangent \`a la surface d'\'equation 
$z=4x^2 +y^2$ au point $(x_0,y_0,z_0)$ 
est donn\'e par l'\'equation
\begin{align*}
z&=z_0 +8x_0 (x-x_0) + 2y_0 (y-y_0) 
\\
&= 8x_0x +2y_0 y +z_0 - 8x^2_0 -  2y^2_0= 
8x_0x +2y_0 y  - z_0 
\end{align*}
d'o\`u par
\begin{equation}
z-8x_0x -2y_0 y  = z_0.
\label{tang2}
\end{equation}
Pour que
ce plan soit parall\`ele au plan d'\'equation $x+2y+z=6$ il faut et il 
suffit que
$(1,2)=(-8x_0, -2y_0)$ d'o\`u que $x_0=-1/8$ et $y_0=-1$.
Par cons\'equent, 
le point cherch\'e sur le parabolo\"\i de $z=4x^2 +y^2$
est le point
$(-1/8,-1,17/16)$.
De m\^eme, 
pour que
le plan \eqref{tang2} 
soit parall\`ele au plan d'\'equation $3x+5y-2z=3$ il faut et il 
suffit que
$(3/2,5/2)=(8x_0, 2y_0)$ d'o\`u que $x_0=3/16$ et $y_0=5/4$,
et le point cherch\'e sur le parabolo\"\i de $z=4x^2 +y^2$
est alors le point
$(3/16,5/4,9/64+25/16)$ =$(3/16,5/4,109/64)$.
\fincorrection
\correction{002631}
\begin{enumerate} 
\item 
Le vecteur normal du c\^one $C$
au point $(x_0,y_0,z_0)$ de $C$
est le vecteur
$(x_0,y_0,-z_0)$ et le plan tangent au c\^one $C$
en ce point est donn\'e par l'\'equation
\[
x_0x + y_0y -z_0z =0
\]
car l'origine appartient \`a ce plan.

\item L'intersection du c\^one $C$ avec 
le plan vertical d'\'equation $y=ax$ o\`u $a\in \R$ 
est constitu\'ee des points
$x(1,a, \pm \sqrt {1+a^2})$ o\`u $x \in \R$, c.a.d. des deux droites
\[
{\cal D}_1=\{x(1,a, \sqrt {1+a^2}); x \in \R\},
\quad 
{\cal D}_2=\{x(1,a, - \sqrt {1+a^2}); x \in \R\}.
\]
L'intersection du demi-c\^one $C^+$ avec 
ce plan vertical est donc constitu\'ee des 
deux demi-droites
\begin{align*}
{}&{\cal D}^+_1=\{x(1,a, \sqrt {1+a^2}); x \in \R, x \geq 0\}
\\
{}&{\cal D}^+_2=\{x(-1,-a, \sqrt {1+a^2}); x \in \R, x \geq 0\}.
\end{align*} 
\item  Le vecteur normal en un point quelconque 
$x(1,a, \sqrt {1+a^2})$ de ${\cal D}_1$ 
respectivement $x(1,a, -\sqrt {1+a^2})$ de ${\cal D}_2$ 
est le vecteur
$x(1,a, -\sqrt {1+a^2})$ 
respectivement $x(1,a, \sqrt {1+a^2})$ 
d'o\`u 
la direction et donc
le plan tangent au c\^one $C$ 
sont le m\^eme en tout point de ${\cal D}_1\setminus \{(0,0,0)\}$ 
respectivement  ${\cal D}_2\setminus \{(0,0,0)\}$.
\end{enumerate}
\fincorrection
\correction{002632}
\begin{enumerate} 
\item  La forme \eqref{tang3} de l'\'equation du plan tangent au graphe
$z=x^2-2y^3$ de la fonction $f$ au point $(x_0,y_0,z_0)$ 
nous donne l'\'equation
\[
z-z_0 = 2x_0(x-x_0)-6y^2_0(y-y_0) = 2x_0x -6y^2_0y -2x^2_0+6y^3_0
\]
\item Au point $(2,1,2)$, 
ce plan tangent 
est ainsi  donn\'e par l'\'equation
\[
4x-6y-z =0.
\]
Pour que ce plan soit parall\`ele au  plan tangent  au point 
$(x_1,y_1,z_1)$ distinct de $(x_0,y_0,z_0)$  il faut et il suffit que
$(4,6,-1)=(2x_1, 6y_1^2,-1)$ et
$y_1 \ne 1$, c.a.d. que
$(x_1,y_1,z_1)=(2, -1,6)$.
\end{enumerate}
\fincorrection
\correction{002633}
\begin{enumerate}  
\item $\mathrm{lim}
_{\begin{smallmatrix} (x,y) \to (0,0)\\ (x,y) \ne (0,0) \end{smallmatrix}}
\frac{xy^2}{x^2+y^2}=\lim_{r\to 0}r \cos \varphi \sin^2 \varphi
$ 
existe et vaut z\'ero puisque $\cos \varphi \sin^2 \varphi$ est born\'e.
Par cons\'equent $f$ est continue \`a l'origine et donc partout.
Il est \'evident que la fonction $f$ est diff\'erentiable
en chaque point distinct de l'origine.
Soit $v=(a,b)\in\R^2$ non nul. Alors
\[
D_vf(0,0) = \left .\frac{d}{dt}\left(t\frac{ab^2}{a^2+b^2}\right)\right|_{t=0}=\frac{ab^2}{a^2+b^2}
\]
existe d'o\`u
$\frac{\partial f}{\partial x}(0,0)=0$ et
$\frac{\partial f}{\partial y}(0,0)=0$;
puisqu'il existe une d\'eriv\'ee directionnelle non nulle,
la fonction $f$ ne peut pas \^etre diff\'erentiable en $(0,0)$.

\item L'association $\R^2 \to \R,\ v \mapsto D_vf(0,0)$
n'est \'evidemment pas lin\'eaire
et les droites appartenant \`a la famille  des droites passant par l'origine
et de vecteurs directeurs $(v, D_vf(0,0))\in\R^3$ ne forment pas un plan. 

\item Dans $\R^2\setminus\{ (0,0)\}$, 
\begin{align*}
\frac{\partial f}{\partial x} &=\frac{y^2(x^2+y^2)-2x^2y^2}{(x^2+y^2)^2}
=\frac{y^4-x^2y^2}{(x^2+y^2)^2}=\sin^4 \varphi -\sin^2 \varphi \cos^2 \varphi
\\
\frac{\partial f}{\partial y} &=\frac{2xy(x^2+y^2) -2xy^3}{(x^2+y^2)^2}
=\frac{2x^3y}{(x^2+y^2)^2} = 2\sin \varphi \cos^3 \varphi
\end{align*}
d'o\`u, en coordonn\'ees polaires,
\[
D_vf(x,y)=D_vf(r \cos \varphi,r \sin \varphi)=
a(\sin^4 \varphi -\sin^2 \varphi \cos^2 \varphi)
+2b\sin \varphi \cos^3 \varphi 
\]
et, $\varphi$ \'etant fix\'e, 
\[ 
\mathrm{lim}_{r \to 0}D_vf(r \cos \varphi,r \sin \varphi)=
a(\sin^4 \varphi -\sin^2 \varphi \cos^2 \varphi)
+2b\sin \varphi \cos^3 \varphi .
\]
Par cons\'equent,
$D_vf(x,y)$ n'est pas continu en $(x,y)$ sauf peut-\^etre si $a=0$.
Par exemple, avec $\sin \varphi =1$, on trouve
\[ 
\mathrm{lim}_{r \to 0}D_vf(0,r)=a 
\]
et $a \ne \frac{ab^2}{a^2+b^2}$ sauf si $a=0$.
Si $a=0$, la d\'eriv\'ee directionnelle $D_v$ est la d\'eriv\'ee partielle
$\frac{\partial f}{\partial y}$ et, $\varphi$ \'etant fix\'e,
\[ 
\mathrm{lim}_{r \to 0}D_vf(r \cos \varphi,r \sin \varphi)=
2b\sin \varphi \cos^3 \varphi 
\]
ce qui n'est pas  nul si $\sin \varphi \cos \varphi$ ne l'est pas.
Puisque $\frac{\partial f}{\partial y}(0,0)=0$, la d\'eriv\'ee partielle
$\frac{\partial f}{\partial y}$ n'est continue en $(0,0)$ non plus.
\end{enumerate}
\fincorrection
\correction{002634} 
\begin{align*}
\frac{\partial f}{\partial x}&= -\cos x \cos y\exp[-\sin x \cos y ]
\\
\frac{\partial f}{\partial y}&= \sin x \sin y\exp[-\sin x \cos y ]
\end{align*}
etc. d'o\`u, avec $\frac{\partial f}{\partial x}(0,0)=-1$ et
$\frac{\partial f}{\partial y}(0,0)=0$,
\[
f(x,y)
=
f(0,0)+\frac{\partial f}{\partial x}(0,0)x+\frac{\partial f}{\partial y}(0,0)y
+\ldots
=1-x +\ldots
\]
Avec $x=0,0184$ on trouve, pour
$\exp[\sin(3.16)\cos(0.02))]$,
la valeur approch\'ee $1-0,0184 =0,9816$. 

N.B. On peut faire mieux si n\'ecessaire : Avec
\begin{align*}
\frac{\partial^2 f}{\partial x^2}&= (\sin x \cos y
+\cos^2 x \cos^2 y)\exp[-\sin x \cos y ]
\\
\frac{\partial^2 f}{\partial x \partial y}&
= (\cos x \sin y+\cos x \cos y \sin x \sin y)\exp[-\sin x \cos y ]
\\
\frac{\partial^2 f}{\partial y^2}&=(\sin x \cos y
+\sin^2 x \sin^2 y)\exp[-\sin x \cos y ]
\end{align*}
on trouve
\[
f(x,y)
=
f(0,0)+\frac{\partial f}{\partial x}(0,0)x+\frac{\partial f}{\partial y}(0,0)y
+\ldots
=1-x++\tfrac 12 x^2 +\ldots
\]
etc.

De m\^eme,
\begin{align*}
\frac{\partial h}{\partial x}&= \frac 1 {2(1+(\sqrt{4+x}-2\exp(y))^2)\sqrt{4+x}}
\\
\frac{\partial h}{\partial y}&= -\frac {2\exp(y)} {1+(\sqrt{4+x}-2\exp(y))^2}
\end{align*}
etc. d'o\`u,  avec $\frac{\partial h}{\partial x}(0,0)=\tfrac 14$ et
$\frac{\partial h}{\partial y}(0,0)=-2$,
\[
h(x,y)
=
h(0,0)+\frac{\partial h}{\partial x}(0,0)x+\frac{\partial h}{\partial y}(0,0)y
+\ldots
=\tfrac 14 x-2y +\ldots .
\]
Avec $x=0,03$ et $y=0,01$ on trouve, pour
$ \arctan[\sqrt{4.03}-2\exp(0.01)]$,
la valeur approch\'ee $0,0075-0,02 =-0,00125$. 

\fincorrection
\nocorrection
\nocorrection
\correction{005555}
\textbullet~Pour $(x,y)\in\Rr^2$, $x^2+y^2=0\Leftrightarrow x=y=0$ et donc $f$ est définie sur $\Rr^2$.
\textbullet~$f$ est de classe $C^\infty$ sur $\Rr^2\setminus\{(0,0)\}$ en tant que quotient de fonctions de classe $C^\infty$ sur $\Rr^2\setminus\{(0,0)\}$ dont le dénominateur ne s'annule pas sur $\Rr^2\setminus\{(0,0)\}$.
 

\textbullet~Pour $(x,y)\neq(0,0)$, $|f(x,y)|\leqslant\frac{|xy|(x^2+y^2)}{x^2+y^2}=|xy|$. Comme $\lim_{(x,y)\rightarrow (0,0)}|xy|=0$, on en déduit que $\displaystyle\lim_{\substack{(x,y)\rightarrow(0,0)\\(x,y)\neq(0,0)}}f(x,y)=0=f(0,0)$. Ainsi, $f$ est continue en $(0,0)$ et donc sur $\Rr^2$.
\textbullet~\textbf{Existence de $\frac{\partial f}{\partial x}(0,0)$.} Pour $x\neq0$,

\begin{center}
$\frac{f(x,0)-f(0,0)}{x-0}=\frac{x\times0\times(x^2-0^2)}{x\times(x^2+0^2)}=0$,
\end{center}
et donc $\lim_{x\rightarrow 0}\frac{f(x,0)-f(0,0)}{x-0}=0$. Ainsi, $f$ admet une dérivée partielle par rapport à sa première variable en $(0,0)$ et $\frac{\partial f}{\partial x}(0,0)=0$.
\textbullet~Pour $(x,y)\neq(0,0)$, $\frac{\partial f}{\partial x}(x,y)=y\frac{(3x^2-y^2)(x^2+y^2)-(x^3-y^2x)(2x)}{(x^2+y^2)^2}=\frac{y(x^4+4x^2y^2-y^4)}{(x^2+y^2)^2}$.

Finalement, $f$ admet sur $\Rr^2$ une dérivée partielle par rapport à sa première variable définie par 

\begin{center}
$\forall(x,y)\in\Rr^2$, $\frac{\partial f}{\partial x}(x,y)=\left\{
\begin{array}{l}
\rule[-4mm]{0mm}{0mm}0\;\text{si}\;(x,y)=(0,0)\\
\frac{y(x^4+4x^2y^2-y^4)}{(x^2+y^2)^2}\;\text{si}\;(x,y)\neq(0,0)
\end{array}
\right.$.
\end{center}
\textbullet~Pour $(x,y)\in\Rr^2$, $f(y,x)=-f(x,y)$. Par suite, $\forall(x,y)\in\Rr^2$, $\frac{\partial f}{\partial y}(x,y)=-\frac{\partial f}{\partial x}(y,x)$.
En effet, pour $(x_0,y_0)$ donné dans $\Rr^2$

\begin{center}
$\frac{f(x_0,y)-f(x_0,y_0)}{y-y_0}=\frac{-f(y,x_0)+f(y_0,x_0)}{y-y_0}=-\frac{f(y,x_0)-f(y_0,x_0)}{y-y_0}\underset{y\rightarrow y_0}{\rightarrow}-\frac{\partial f}{\partial x}(y_0,x_0)$.
\end{center}
Donc, $f$ admet sur $\Rr^2$ une dérivée partielle par rapport à sa deuxième variable définie par 

\begin{center}
$\forall(x,y)\in\Rr^2$, $\frac{\partial f}{\partial y}(x,y)=-\frac{\partial f}{\partial x}(y,x)=\left\{
\begin{array}{l}
\rule[-4mm]{0mm}{0mm}0\;\text{si}\;(x,y)=(0,0)\\
\frac{x(x^4-4x^2y^2-y^4)}{(x^2+y^2)^2}\;\text{si}\;(x,y)\neq(0,0)
\end{array}
\right.$.
\end{center}
\textbullet~\textbf{Continuité de $\frac{\partial f}{\partial x}$ et $\frac{\partial f}{\partial y}$ en $(0,0)$.} Pour $(x,y)\neq(0,0)$,

\begin{center}
$\left|\frac{\partial f}{\partial x}(x,y)-\frac{\partial f}{\partial x}(0,0)\right|=\frac{|y(x^4+4x^2y^2-y^4)|}{(x^2+y^2)^2}\leqslant\frac{|y|(x^4+4x^2y^2+y^4)}{(x^2+y^2)^2}\leqslant\frac{|y|(2x^4+4x^2y^2+2y^4)}{(x^2+y^2)^2}=2|y|$.
\end{center}
Comme $2|y|$ tend vers $0$ quand $(x,y)$ tend vers $(0,0)$, $\left|\frac{\partial f}{\partial x}(x,y)-\frac{\partial f}{\partial x}(0,0)\right|$ tend vers $0$ quand $(x,y)$ tend vers $(0,0)$. On en déduit que l'application $\frac{\partial f}{\partial x}$ est continue en $(0,0)$ et donc sur $\Rr^2$.
Enfin, puisque $\forall(x,y)\in\Rr^2$, $\frac{\partial f}{\partial y}(x,y)=-\frac{\partial f}{\partial x}(y,x)$, $\frac{\partial f}{\partial y}$ est continue sur $\Rr^2$. $f$ est donc au moins de classe $C^1$ sur $\Rr^2$.
\textbullet~Pour $x\neq0$, $\frac{\frac{\partial f}{\partial y}(x,0)-\frac{\partial f}{\partial y}(0,0)}{x-0}=\frac{x^4}{x^4}=1$ et donc $\lim_{x\rightarrow 0}\frac{\frac{\partial f}{\partial y}(x,0)-\frac{\partial f}{\partial y}(0,0)}{x-0}=1$. Donc $\frac{\partial^2f}{\partial y\partial x}(0,0)$ existe et $\frac{\partial^2f}{\partial y\partial x}(0,0)=1$.
Pour $y\neq0$, $\frac{\frac{\partial f}{\partial x}(y,0)-\frac{\partial f}{\partial x}(0,0)}{y-0}=-\frac{y^4}{y^4}=-1$ et donc $\lim_{y\rightarrow 0}\frac{\frac{\partial f}{\partial x}(y,0)-\frac{\partial f}{\partial x}(0,0)}{y-0}=-1$. Donc $\frac{\partial^2f}{\partial x\partial y}(0,0)$ existe et $\frac{\partial^2f}{\partial x\partial y}(0,0)=-1$.
$\frac{\partial^2f}{\partial y\partial x}(0,0)\neq\frac{\partial^2f}{\partial x\partial y}(0,0)$ et donc $f$ n'est pas de classe $C^2$ sur $\Rr^2$ d'après le théorème de \textsc{Schwarz}.

\begin{center}
\shadowbox{
$f$ est de classe $C^1$ exactement sur $\Rr^2$.
}
\end{center}
\fincorrection
\correction{005888}
\textbullet~$f$ est définie sur $\Rr^2$.

\textbullet~$f$ est de classe $C^\infty$ sur $\Rr^2\setminus\{(0,0)\}$ en tant que fraction rationnelle dont le dénominateur ne s'annule pas sur $\Rr^2\setminus\{(0,0)\}$.

\textbullet~\textbf{Continuité en $(0,0)$.} Pour $(x,y)\neq(0,0)$,

\begin{center}
$|f(x,y)-f(0,0)|= \frac{|xy||x^2-y^2|}{x^2+y^2}\leqslant|xy|\times \frac{x^2+y^2}{x^2+y^2}=|xy|$.
\end{center}

Comme $|xy|$ tend vers $0$ quand le couple $(x,y)$ tend vers le couple $(0,0)$, on a donc $\displaystyle\lim_{\substack{(x,y)\rightarrow(0,0)\\(x,y)\neq(0,0)}}f(x,y)=f(0,0)$. On en déduit que $f$ est continue en $(0,0)$ et finalement $f$ est continue sur $\Rr^2$.

\begin{center}
$f$ est de classe $C^0$ au moins sur $\Rr^2$.
\end{center}

\textbullet~\textbf{Dérivées partielles d'ordre $1$ sur $\Rr^2\setminus\{(0,0)\}$.} $f$ est de classe $C^1$ au moins sur $\Rr^2\setminus\{(0,0)\}$ et pour $(x,y)\neq(0,0)$,

\begin{center}
$ \frac{\partial f}{\partial x}(x,y)=y \frac{(3x^2-y^2)(x^2+y^2)-(x^3-xy^2)(2x)}{(x^2+y^2)^2}= \frac{y(x^4+4x^2y^2-y^4)}{(x^2+y^2)^2}$,
\end{center}

D'autre part, pour $(x,y)\neq(0,0)$ $f(x,y)=-f(y,x)$. Donc pour $(x,y)\neq(0,0)$,

\begin{center}
$ \frac{\partial f}{\partial y}(x,y)=- \frac{\partial f}{\partial x}(y,x)= \frac{x(x^4-4x^2y^2-y^4)}{(x^2+y^2)^2}$.
\end{center}

\textbullet~\textbf{Existence de $ \frac{\partial f}{\partial x}(0,0)$ et $ \frac{\partial f}{\partial y}(0,0)$.} Pour $x\neq0$,

\begin{center}
$ \frac{f(x,0)-f(0,0)}{x-0}= \frac{0-0}{x}=0$,
\end{center}

et donc $\lim_{x \rightarrow 0} \frac{f(x,0)-f(0,0)}{x-0}=0$. Ainsi, $ \frac{\partial f}{\partial x}(0,0)$ existe et $ \frac{\partial f}{\partial x}(0,0)=0$. De même, $ \frac{\partial f}{\partial y}(0,0)=0$. Ainsi, $f$ admet des dérivées partielles premières sur $\Rr^2$ définies par

\begin{center}
$\forall(x,y)\in\Rr^2$, $ \frac{\partial f}{\partial x}(x,y)=\left\{
\begin{array}{l}
 \frac{y(x^4+4x^2y^2-y^4)}{(x^2+y^2)^2}\;\text{si}\;(x,y)\neq(0,0)\\
0\;\text{si}\;(x,y)=(0,0)
\end{array}
\right.$ et $ \frac{\partial f}{\partial y}(x,y)=\left\{
\begin{array}{l}
 \frac{x(x^4-4x^2y^2-y^4)}{(x^2+y^2)^2}\;\text{si}\;(x,y)\neq(0,0)\\
0\;\text{si}\;(x,y)=(0,0)
\end{array}
\right.$.
\end{center}

\textbullet~\textbf{Continuité de $ \frac{\partial f}{\partial x}$ et $ \frac{\partial f}{\partial y}$ en $(0,0)$.} Pour $(x,y)\neq(0,0)$,

\begin{center}
$\left| \frac{\partial f}{\partial x}(x,y)- \frac{\partial f}{\partial x}(0,0)\right|= \frac{|y||x^4+4x^2y^2-y^4|}{(x^2+y^2)^2}\leqslant|y| \frac{x^4+4x^2y^2+y^4}{(x^2+y^2)^2}\leqslant|y| \frac{2x^4+4x^2y^2+2y^4}{(x^2+y^2)^2}=2|y|$.
\end{center}

Comme $2|y|$ tend vers $0$ quand $(x,y)$ tend vers $(0,0)$, on en déduit que $\left| \frac{\partial f}{\partial x}(x,y)- \frac{\partial f}{\partial x}(0,0)\right|$ tend vers $0$ quand $(x,y)$ tend vers $(0,0)$. Donc la fonction $ \frac{\partial f}{\partial x}$ est continue en $(0,0)$ et finalement sur $\Rr^2$. Il en est de même de la fonction $ \frac{\partial f}{\partial y}$ et on a montré que

\begin{center}
\shadowbox{
$f$ est au moins de classe $C^1$ sur $\Rr^2$.
}
\end{center}
\fincorrection
\correction{005894}
\begin{enumerate}
 \item  $f$ est de classe $C^1$ sur $\Rr^2$ qui est un ouvert de $\Rr^2$. Donc si $f$ admet un extremum local en un point $(x_0,y_0)$ de $\Rr^2$, $(x_0,y_0)$ est un point critique de $f$.

\begin{center}
$df_{(x,y)}=0\Leftrightarrow\left\{
\begin{array}{l}
3x^2+6xy-15=0\\
3x^2-12=0
\end{array}
\right.\Leftrightarrow\left\{
\begin{array}{l}
x=2\\
y= \frac{1}{4}
\end{array}
\right.\;\text{ou}\;\left\{
\begin{array}{l}
x=-2\\
y=- \frac{1}{4}
\end{array}
\right.$.
\end{center}

Réciproquement, $r=6x+6y$, $t=0$ et $s=6x$ puis $rt-s^2=-36x^2$. Ainsi, $(rt-s^2)\left(2, \frac{1}{4}\right)=(rt-s^2)\left(-2,- \frac{1}{4}\right)=-144<0$ et $f$ n'admet pas d'extremum local en $\left(2, \frac{1}{4}\right)$ ou $\left(-2,- \frac{1}{4}\right)$.

\begin{center}
\shadowbox{
$f$ n'admet pas d'extremum local sur $\Rr^2$.
}
\end{center}

\item  La fonction $f$ est de classe $C^1$ sur $\Rr^2$ en tant que polynôme à plusieurs variables. Donc, si $f$ admet un extremum local en $(x_0,y_0)\in\Rr^2$, $(x_0,y_0)$ est un point critique de $f$. Soit $(x,y)\in\Rr^2$.

\begin{align*}\ensuremath
\left\{
\begin{array}{l}
 \frac{\partial f}{\partial x}(x,y)=0\\
\rule{0mm}{7mm} \frac{\partial f}{\partial y}(x,y)=0
\end{array}
\right.&\Leftrightarrow\left\{
\begin{array}{l}
-4(x-y)+4x^3=0\\
4(x-y)+4y^3=0
\end{array}
\right.\Leftrightarrow\left\{
\begin{array}{l}
x^3+y^3=0\\
-4(x-y)+4x^3=0
\end{array}
\right.\Leftrightarrow\left\{
\begin{array}{l}
y=-x\\
x^3-2x=0
\end{array}
\right.\\
 &\Leftrightarrow(x,y)\in\left\{(0,0),\left(\sqrt{2},\sqrt{2}\right),\left(-\sqrt{2},-\sqrt{2}\right)\right\}.
\end{align*}

Réciproquement, $f$ est plus précisément de classe $C^2$ sur $\Rr^2$ et

\begin{center}
$r(x,y)t(x,y)-s^2(x,y)=(-4+12x^2)(-4+12y^2)-(4)^2=-48x^2-48y^2+144x^2y^2=48(3x^2y^2-x^2-y^2)$
\end{center}

\textbullet~$(rt-s^2)\left(\sqrt{2},\sqrt{2}\right)=48(12-2-2)>0$. Donc $f$ admet un extremum local en $\left(\sqrt{2},\sqrt{2}\right)$. Plus précisément, puisque $r\left(\sqrt{2},\sqrt{2}\right)=2\times12-4=20>0$, $f$ admet un minimum local en $\left(\sqrt{2},\sqrt{2}\right)$. De plus, pour $(x,y)\in\Rr^2$,

\begin{align*}\ensuremath
f(x,y)-f\left(\sqrt{2},\sqrt{2}\right)&=-2(x-y)^2+x^4+y^4-8=x^4+y^4-2x^2-2y^2+4xy+8\\
 &\geqslant x^4+y^4-2x^2-2y^2-2(x^2+y^2)+8=(x^4-4x^2+4)+(y^4-4y^2+4)=(x^2-2)^2+(y^2-2)^2\\
 &\geqslant0.
\end{align*}

et $f\left(\sqrt{2},\sqrt{2}\right)$ est un minimum global.

\textbullet~Pour tout $(x,y)\in\Rr^2$, $f(-x,-y)=f(x,y)$ et donc $f$ admet aussi un minimum global en $\left(-\sqrt{2},-\sqrt{2}\right)$ égal à $8$.

\textbullet~$f(0,0)=0$. Pour $x\neq0$, $f(x,x)=2x^4>0$ et donc $f$ prend des valeurs strictement supérieures à $f(0,0)$ dans tout voisinage de $(0,0)$. Pour $x\in\left]-\sqrt{2},\sqrt{2}\right[\setminus\{0\}$, $f(x,0)=x^4-2x^2=x^2(x^2-2)<0$ et $f$ prend des valeurs strictement inférieures à $f(0,0)$ dans tout voisinage de $(0,0)$. Finalement, $f$ n'admet pas d'extremum local en $(0,0)$.

\begin{center}
\shadowbox{
$f$ admet un minimum global égal à $8$, atteint en $\left(\sqrt{2},\sqrt{2}\right)$ et $\left(-\sqrt{2},-\sqrt{2}\right)$.
}
\end{center}
\end{enumerate}
\fincorrection
\correction{005895}
On munit $\mathcal{M}_n(\Rr)$ d'une norme sous-multiplicative $\|\;\|$. Soit $A\in GL_n(\Rr)$. On sait que $GL_n(\Rr)$ est un ouvert de $\mathcal{M}_n(\Rr)$ et donc pour $H\in\mathcal{M}_n(\Rr)$ de norme suffisamment petite, $A+H\in GL_n(\Rr)$. Pour un tel $H$

\begin{center}
$(A+H)^{-1}-A^{-1}=(A+H)^{-1}(I_n-(A+H)A^{-1})=-(A+H)^{-1}HA^{-1}$
\end{center}

puis

\begin{align*}\ensuremath
(A+H)^{-1}-A^{-1}+A^{-1}HA^{-1}&=-(A+H)^{-1}HA^{-1}+A^{-1}HA^{-1}=(A+H)^{-1}(-HA^{-1}+(A+H)A^{-1}HA^{-1})\\
 &=(A+H)^{-1}HA^{-1}HA^{-1}.
\end{align*}

Par suite, $\left\|f(A+H)-f(A)+A^{-1}HA^{-1}\right\|=\left\|(A+H)^{-1}-A^{-1}+A^{-1}HA^{-1}\right\|\leqslant\left\|(A+H)^{-1}\right\|\left\|A^{-1}\right\|^2\left\|H\right\|^2$. 

Maintenant, la formule $M^{-1}= \frac{1}{\text{det}(M)}{^t}(\text{com}(M))$, valable pour tout $M\in GL_n(\Rr)$, et la continuité du déterminant montre que l'application $M\mapsto M^{-1}$ est continue sur l'ouvert $GL_n(\Rr)$. On en déduit que $\left\|(A+H)^{-1}\right\|$ tend vers $\left\|A^{-1}\right\|$ quand $H$ tend vers $0$. Par suite,

\begin{center}
$\lim_{H \rightarrow 0}\left\|(A+H)^{-1}\right\|\left\|A^{-1}\right\|^2\left\|H\right\|=0$ et donc $\lim_{H \rightarrow 0} \frac{1}{\|H\|}\left\|(A+H)^{-1}-A^{-1}+A^{-1}HA^{-1}\right\|=0$.
\end{center}

Comme l'application $H\mapsto -A^{-1}HA^{-1}$ est linéaire, c'est la différentielle de $f$ en $A$.

\begin{center}
\shadowbox{
$\forall A\in GL_n(\Rr)$, $\forall H\in\mathcal{M}_n(\Rr)$, $df_A(H)=-A^{-1}HA^{-1}$.
}
\end{center}
\fincorrection
\correction{005896}
Pour tout complexe $z$ tel que $|z|\leqslant1$,

\begin{center}
$|\sin(z)|=\left|\sum_{n=0}^{+\infty}(-1)^n \frac{z^{2n+1}}{(2n+1)!}\right|\leqslant\sum_{n=0}^{+\infty} \frac{|z|^{2n+1}}{(2n+1)!}=\sh(|z|)\leqslant\sh1$,
\end{center}

l'égalité étant obtenue effectivement pour $z=i$ car $|\sin(i)|=\left| \frac{e^{i^2}-e^{-i^2}}{2i}\right|= \frac{e-e^{-1}}{2}=\sh(1)$.

\begin{center}
\shadowbox{
$\text{Max}\{|\sin z|,\;z\in\Cc,\;|z|\leqslant1\}=\sh(1)$.
}
\end{center}
\fincorrection
\correction{005899}
On munit $(\Rr^3)^2$ de la norme définie par $\forall(x,y)\in(\Rr^3)^2$, $\|(x,y)\|=\text{Max}\{\|h\|_2,\|k\|_2\}$.

\textbullet~Soit $(a,b)\in(\Rr^3)^2$. Pour $(h,k)\in(\Rr^3)^2$,

\begin{center}
$f((a,b)+(h,h))=(a+h).(b+k)=a.b+a.h+b.k+h.k$,
\end{center}

et donc $f((a,b)+(h,h))-f((a,b))=(a.h+b.k)+h.k$. Maintenant l'application $L~:~(h,k)\mapsto a.h+b.k$ est linéaire et de plus, pour $(h,k)\neq(0,0)$,

\begin{center}
$|f((a,b)+(h,h))-f((a,b))-L((h,k))|=|h.k|\leqslant\|h\|_2\|k\|_2\leqslant\|(h,k)\|^2$,
\end{center}

et donc $ \frac{1}{\|(h,k)\|}|f((a,b)+(h,h))-f((a,b))-L((h,k))|\leqslant\|(h,k)\|$ puis

\begin{center}
$\lim_{(h,k) \rightarrow (0,0)} \frac{1}{\|(h,k)\|}|f((a,b)+(h,h))-f((a,b))-L((h,k))|=0$.
\end{center}

Puisque l'application $(h,k)\mapsto a.h+b.k$ est linéaire, on en déduit que $f$ est différentiable en $(a,b)$ et que $\forall(h,k)\in(\Rr^3)^2$, $df_{(a,b)}(h,k)=a.h+b.k$.

La démarche est analogue pour le produit vectoriel :

\begin{center}
$ \frac{1}{\|(h,k)\|}\|(a+h)\wedge(b+k)-a\wedge b-a\wedge h-b\wedge k\|_2= \frac{\|h\wedge k\|_2}{\|(h,k)\|}\leqslant \frac{\|h\|_2\|k\|_2}{\|(h,k)\|}\leqslant\|(h,k)\|$.
\end{center}

Puisque l'application $(h,k)\mapsto a\wedge h+b\wedge k$ est linéaire, on en déduit que $g$ est différentiable en $(a,b)$ et que $\forall(h,k)\in(\Rr^3)^2$, $dg_{(a,b)}(h,k)=a\wedge h+b\wedge k$.
\fincorrection
\correction{002622}

\noindent
\begin{enumerate}  

\item $D_f=\R^2$. 
\begin{align*}
\frac{\partial f}{\partial x}&= 2x\exp(xy)+x^2y\exp(xy)
\\
\frac{\partial f}{\partial y}&= x^3\exp(xy)
\end{align*}
\noindent
\item $D_f=\{(x,y); x > 0 \ \text{ou} \ y \ne 0\}\subseteq \R^2$.
\begin{align*}
\frac{\partial f}{\partial x}&= 
\frac{1+\frac{x}{\sqrt{x^2+y^2}}}{x+\sqrt{x^2+y^2}}= 
\frac 1{\sqrt{x^2+y^2}}
\\
\frac{\partial f}{\partial y}&=  
\frac{\frac{y}{\sqrt{x^2+y^2}}}{x+\sqrt{x^2+y^2}}
=\frac{y}{x\sqrt{x^2+y^2}+x^2+y^2}
\end{align*}

\item $D_f=\R^2$. 
\begin{align*}
\frac{\partial f}{\partial x}&= 2\sin x \cos x
\\
\frac{\partial f}{\partial y}&= -2\sin y \cos y 
\end{align*}

\item $D_f=\{(x,y,z); z \ne 0\}\subseteq \R^3$.
\begin{align*}
\frac{\partial f}{\partial x}&= 2xy^2\sqrt{z}
\\
\frac{\partial f}{\partial y}&= 2x^2y\sqrt{z} 
\\
\frac{\partial f}{\partial z}&= \frac {x^2y^2}{2\sqrt{z}}  
\end{align*}
\end{enumerate}
\fincorrection
\correction{002623}
\begin{enumerate}  
\item $\frac{\partial f}{\partial x}=\cos y + y\exp x$,\ 
$\frac{\partial f}{\partial y}=-x\sin y + \exp x $.

\item $D_vf(0,0)=\cos \theta \,\frac{\partial f}{\partial x}(0,0)
+\sin \theta \,\frac{\partial f}{\partial y}(0,0)
=\cos \theta +\sin \theta$.
Cette d\'eriv\'ee directionnelle de $f$ est maximale
quand $\sin \theta = \cos \theta = \frac{\sqrt{2}}{2}$, c.a.d. quand
  $\theta =\frac \pi 4$,
et minimale quand $\sin \theta = \cos \theta = -\frac{\sqrt{2}}{2}$, 
c.a.d. quand  $\theta =\frac 54 \pi$.

Signification g\'eom\'etrique:
Le plan engendr\'e par le vecteur $(\cos \theta,\sin \theta ,0)$ 
et l'axe des $z$ rencontre le graphe $z=f(x,y)$ en une courbe.
Cette courbe est de pente maximale en valeur absolue pour
$\cos \theta=\sin \theta=\frac{\sqrt{2}}{2}$ et
$\cos \theta=\sin \theta=-\frac{\sqrt{2}}{2}$ (m\^eme plan).
Les deux signes s'expliquent par les deux orientations possibles de cette courbe
(sens du param\'etrage).
\end{enumerate}
\fincorrection
\correction{002624}
\begin{enumerate}  
\item  $f(x,y)=(x^2+y^2)^x= \mathrm e^{x \log (x^2+y^2)}=
\mathrm e^{2r\cos \varphi\log r}$. Puisque $\cos \varphi$
est born\'e,
 $\lim_{\begin{smallmatrix} r \to 0\\ r >0
\end{smallmatrix}} 2r\cos \varphi\log r =0
$
d'o\`u
\[
\mathrm{lim}_{\begin{smallmatrix} (x,y) \to (0,0)\\ (x,y) \ne (0,0)
\end{smallmatrix}} f(x,y)=
\mathrm e^{\mathrm{lim}_{\begin{smallmatrix} r \to 0\\ r >0
\end{smallmatrix}} 2r\cos \varphi\log r}=\mathrm e^0=1,
\]
car la fonction exponentielle est continue.

\item  Dans $\R^2 \setminus \{(0,0)\}$ les d\'eriv\'ees partielles
par rapport aux variables $x$ et $y$ se calculent ainsi:
\begin{align*}\frac{\partial f}{\partial x}&
= \left(\ln (x^2+y^2)+\frac {2x^2}{x^2+y^2}\right)(x^2+y^2)^x
\\
\frac{\partial f}{\partial y}&
= \left(\frac {2xy}{x^2+y^2}\right)(x^2+y^2)^x
\end{align*}

\item Pour que la d\'eriv\'ee partielle
$\frac{\partial f}{\partial x}(0,0)$ existe, il faut et il suffit que
\[
\mathrm{lim}_{\begin{smallmatrix} x \to 0\\ x \ne 0
\end{smallmatrix}}\frac {f(x,0) -1}x =
\mathrm{lim}_{\begin{smallmatrix} x \to 0\\ x \ne 0
\end{smallmatrix}}\frac {(x^2)^x-1}x = 
\mathrm{lim}_{\begin{smallmatrix} x \to 0\\ x > 0
\end{smallmatrix}}\frac {\mathrm e^{2x \log x} -1}x
\]
existe. Si $x>0$,
\[
\frac {\mathrm e^{2x \log x} -1}x = 2 \log x + \varepsilon(x)
\]
o\`u $\mathrm{lim}_{\begin{smallmatrix} x \to 0\\ x > 0
\end{smallmatrix}}\varepsilon(x)=0$. Par cons\'equent,
 la d\'eriv\'ee partielle
$\frac{\partial f}{\partial x}(0,0)$ n'existe pas.
D'autre part,
\[
\frac{\partial f}{\partial y}(0,0)=
\mathrm{lim}_{\begin{smallmatrix} y \to 0\\ y \ne 0
\end{smallmatrix}}\frac {f(0,y) -1}y =
\mathrm{lim}_{\begin{smallmatrix} y \to 0\\ y \ne 0
\end{smallmatrix}}\frac {(y^2)^0-1}y = 0
\]
existe.
\end{enumerate}
\fincorrection
\correction{002625}
\begin{enumerate}  
\item Puisque $f(x,y)=\frac{x^2y +3y^3}{x^2+y^2}
=r(\cos^2 \varphi \sin \varphi + 3 \sin^3 \varphi)$, il s'ensuit que
\[
\mathrm{lim}_{\begin{smallmatrix} (x,y) \to (0,0)\\ (x,y) \ne (0,0)
\end{smallmatrix}} f(x,y) = 0
\]
car $\cos^2 \varphi \sin \varphi + 3 \sin^3 \varphi$ reste born\'e.
Par cons\'equent la fonction $f$ est continue en $(0,0)$.

\item Les d\'eriv\'ees partielles
\begin{align*}
\frac{\partial f}{\partial x}(0,0)&=
\mathrm{lim}_{\begin{smallmatrix} x \to 0\\ x \ne 0
\end{smallmatrix}}\frac {f(x,0)}x =
\mathrm{lim}_{\begin{smallmatrix} x \to 0\\ x \ne 0
\end{smallmatrix}}\frac{0}{x^2} = 0
\\
\frac{\partial f}{\partial y}(0,0)&=
\mathrm{lim}_{\begin{smallmatrix} y \to 0\\ y \ne 0
\end{smallmatrix}}\frac {f(0,y)}y =
\mathrm{lim}_{\begin{smallmatrix} y \to 0\\ y \ne 0
\end{smallmatrix}}\frac{3y^3}{y^2} = 3
\end{align*}
existent.

\item Puisque $f(x,x)=\frac{4x^3}{2x^2}= 2x$, la d\'eriv\'ee
directionnelle $D_vf(0,0)$ suivant le vecteur $v=(1,1)$ est non nulle.
Par cons\'equent, la fonction $f$ n'est pas diff\'erentiable en $(0,0)$.

\item 
\begin{align*}
\frac{\partial f}{\partial x}&=
\frac{\partial }{\partial x}\frac{x^2y +3y^3}{x^2+y^2}
=\frac{2xy(x^2+y^2)-2x(x^2y +3y^3)}{(x^2+y^2)^2}=
-4\frac{xy^3}{(x^2+y^2)^2}
\\
\frac{\partial f}{\partial y}&=
\frac{\partial }{\partial y}\frac{x^2y +3y^3}{x^2+y^2}
=\frac{(x^2+3y^2)(x^2+y^2)-2y(x^2y +3y^3)}{(x^2+y^2)^2}=
\frac{x^4 +8x^2y^2 +3y^4}{(x^2+y^2)^2}
\end{align*}

\item D'apr\`es \eqref{eqt}, cette \'equation s'\'ecrit
\[
z-2 =\frac{\partial f}{\partial x}(1,1)(x-1)+
\frac{\partial f}{\partial y}(1,1)(y-1)
=1-x+3(y-1)
\]
d'où $z=3y-x$.
\item La fonction $F\colon \R^2 \to \R^2$ s'\'ecrit 
$F(x,y)=\left(\frac{x^2y +3y^3}{x^2+y^2}, \frac{y^2x +3x^3}{x^2+y^2}\right)$
et sa matrice jacobienne
\[
\mathrm{J}_F(1,1)=
\left[
\begin{matrix} 
\frac{\partial f}{\partial x}(1,1)
&
\frac{\partial f}{\partial y}(1,1)
\\
\frac{\partial f}{\partial y}(1,1)
&
\frac{\partial f}{\partial x}(1,1)
\end{matrix}
\right]
= \left[
\begin{matrix} -1 & 3\\ 3 & -1 \end{matrix}
\right]
\]
au point $(1,1)$ est inversible.
Par cons\'equent, la fonction $F$ admet une r\'eciproque locale au
voisinage du point $(1,1)$.
Au point $(2,2)$,
\[
\mathrm{J}_F(2,2)=
\left[
\begin{matrix} 
\frac{\partial f}{\partial x}(2,2)
&
\frac{\partial f}{\partial y}(2,2)
\\
\frac{\partial f}{\partial y}(2,2)
&
\frac{\partial f}{\partial x}(2,2)
\end{matrix}
\right]
= \left[
\begin{matrix} -1 & 3\\ 3 & -1 \end{matrix}
\right]
\]
d'o\`u la fonction $F$ admet \'egalement une r\'eciproque locale au
voisinage du point $(2,2)$.
\end{enumerate}
\fincorrection
\correction{002626}
\begin{align*}
\frac{\partial g}{\partial x}
&=
\frac{\partial f}{\partial x}-\frac{\partial f}{\partial z}
\\
\frac{\partial g}{\partial y}
&=
\frac{\partial f}{\partial y}-\frac{\partial f}{\partial x}
\\
\frac{\partial g}{\partial z}
&=
\frac{\partial f}{\partial z}-\frac{\partial f}{\partial y}
\end{align*}
d'o\`u \eqref{ex3}.
\fincorrection
\correction{002627}
\begin{enumerate} 
\item $g(f(x,y))= xy^2 \sin^2 (xy) \cos x \exp(y^2)$

\item \begin{align*}
\frac{\partial (g \circ f)}{\partial x}
&= 
y^2 \sin (xy) \exp(y^2)
(
2xy \cos x  \cos (xy) 
-
x \sin x\sin (xy)
+
\cos x \sin (xy)
)
\\
\frac{\partial (g\circ f)}{\partial y}
&=
2xy \cos x   \sin (xy)\exp(y^2)
(xy  \cos (xy) +(1+y^2)  \sin (xy) )
\end{align*}
\item Calculons d'abord
\begin{align*}
\frac{\partial w}{\partial x}
&=y\sin(xy)\exp(y^2) +xy^2\cos(xy)\exp(y^2)
\\
&=y\exp(y^2)(\sin(xy) +xy\cos(xy))
\\
\frac{\partial w}{\partial y}
&=x\sin(xy)\exp(y^2) +x^2y\cos(xy)\exp(y^2) +2xy^2\sin(xy)\exp(y^2)
\\
&=
x\exp(y^2)(\sin(xy) +xy\cos(xy) +2y^2\sin(xy))
\\
&=
x\exp(y^2)((1+2y^2)\sin(xy) +xy\cos(xy)) .
\end{align*}
Ainsi la matrice jacobienne  $\mathrm J_f$ de $f$ s'\'ecrit
\begin{align*}
\mathrm{J_f}&=
\left[
\begin{matrix} 
\frac{\partial u}{\partial x}
&
\frac{\partial u}{\partial y}
\\
\frac{\partial v}{\partial x}
&
\frac{\partial v}{\partial y}
\\
\frac{\partial w}{\partial x}
&
\frac{\partial w}{\partial y}
\end{matrix}
\right]
\\
&=
\left[
\begin{matrix} 
y \cos (xy)
&
x \cos (xy)
\\
-y\sin x
&
\cos x
\\
y\exp(y^2)(\sin(xy) +xy\cos(xy))
&
x\exp(y^2)((1+2y^2)\sin(xy) +xy\cos(xy))
\end{matrix}
\right] .
\end{align*}
De m\^eme, la matrice jacobienne $\mathrm J_g$ de $g$ est :
\begin{align*}
\mathrm J_g
&=
\left[
\frac{\partial g}{\partial u}, \frac{\partial g}{\partial v},\frac{\partial g}{\partial w}
\right]=[vw,uw,uv]
\\
&=\left [
xy^2 \sin (xy) \cos x \exp(y^2),xy \sin^2 (xy) \exp(y^2),y \sin (xy) \cos x 
\right]
\end{align*}
\item La matrice jacobienne $\mathrm J_{g\circ f}$ de la fonction
compos\'ee  $g\circ f$ s'\'ecrit comme produit matricielle
\begin{align*}
\mathrm J_{g\circ f}&=\mathrm J_g \circ \mathrm J_f =\left[
\frac{\partial g}{\partial u}, \frac{\partial g}{\partial v},\frac{\partial g}{\partial w}
\right]
\left[
\begin{matrix} 
\frac{\partial u}{\partial x}
&
\frac{\partial u}{\partial y}
\\
\frac{\partial v}{\partial x}
&
\frac{\partial v}{\partial y}
\\
\frac{\partial w}{\partial x}
&
\frac{\partial w}{\partial y}
\end{matrix}
\right]
\end{align*}
d'o\`u
\begin{align*}
\frac{\partial (g \circ f)}{\partial x}
&= 
(xy^2 \sin (xy) \cos x \exp(y^2)) y \cos (xy)
\\
&\quad
-(xy \sin^2 (xy) \exp(y^2)) y\sin x
\\
&\quad
+(y \sin (xy) \cos x)
y\exp(y^2)(\sin(xy) +xy\cos(xy))
\\
&=
xy^3 \cos x \sin (xy) \cos (xy)  \exp(y^2))
-
xy^2 \sin x\sin^2 (xy) \exp(y^2)
\\
&\quad
+y^2 \cos x \sin^2 (xy) \exp(y^2)
+xy^3 \cos x \sin (xy) \cos(xy)\exp(y^2)
\\&=
y^2 \sin (xy) \exp(y^2)
(
2xy \cos x  \cos (xy) 
-
x \sin x\sin (xy)
+
\cos x \sin (xy)
)
\\
\frac{\partial (g\circ f)}{\partial y}
&=
(xy^2 \sin (xy) \cos x \exp(y^2))
x\cos (xy)
\\
&\quad
+
(xy \sin^2 (xy) \exp(y^2))
\cos x
\\
&\quad
+(y \sin (xy) \cos x)
x\exp(y^2)((1+2y^2)\sin(xy) +xy\cos(xy))
\\
&=
x^2y^2  \cos x\sin (xy) 
\cos (xy)\exp(y^2) +xy \cos x\sin^2 (xy) \exp(y^2)
\\
&\quad
+xy(1+2y^2)  \cos x \sin^2 (xy)
\exp(y^2)
+x^2y^2 \cos x\sin (xy)\cos(xy)\exp(y^2)
\\
&=
2 x^2y^2  \cos x\sin (xy) 
\cos (xy)\exp(y^2) +2xy(1+y^2) \cos x\sin^2 (xy) \exp(y^2)
\\
&=
2xy \cos x   \sin (xy)\exp(y^2)
(xy  \cos (xy) +(1+y^2)  \sin (xy) ) .
\end{align*}
\end{enumerate}
\fincorrection
\nocorrection
\correction{004137}
$\lim = 3$.
\fincorrection
\correction{004138}
\begin{enumerate}
  \item 
  \item $\frac{\partial f}{\partial x} = \frac{\pm 1}{1+x^2}$, $+$ si $y > x$, $-$ si $y < x$.\par
             $\frac{\partial f}{\partial y} = \frac{\pm 1}{1+y^2}$, $-$ si $y > x$, $+$ si $y < x$.
  \item Pour $y \ge x$, $f(x,y) = \frac\pi2 + g(x,y)$,\quad
             pour $y \le x$, $f(x,y) = \frac\pi2 - g(x,y)$.
\end{enumerate}
 \fincorrection
\nocorrection
\nocorrection
\correction{004141}
$\frac{\partial^2 f}{\partial x \partial y}= \frac{\partial u}{\partial x}\frac{\partial u}{\partial y}\frac{\partial^2 g}{\partial u^2} + \left(\frac{\partial u}{\partial x}\frac{\partial v}{\partial y}+\frac{\partial u}{\partial y}\frac{\partial v}{\partial x}\right)\frac{\partial^2 g}{\partial u \partial v}
         + \frac{\partial v}{\partial x}\frac{\partial v}{\partial y}\frac{\partial^2 g}{\partial v^2} + \frac{\partial^2 u}{\partial x \partial y}\frac{\partial g}{\partial u} + \frac{\partial^2 u}{\partial x \partial y}\frac{\partial g}{\partial v}$.
\fincorrection
\nocorrection
\correction{004143}
\begin{enumerate}
  \item 
  \item $\Delta f = \frac{\partial^2 g}{\partial \rho^2} + \frac 1\rho \frac{\partial g}{\partial \rho} + \frac 1{\rho^2} \frac{\partial^2 g}{\partial \theta^2}$.
\end{enumerate}
 \fincorrection
\nocorrection
\correction{004145}
$\Delta F = \frac{n-1}r f'(r) + f''(r)$.
\fincorrection
\correction{004146}
$\frac{\partial g}{\partial x}(0,y) = - yf'(\pi/2)$, $\frac{\partial  g}{\partial y}(x,0) = xf'(0)$.
\fincorrection
\correction{004147}
\begin{enumerate}
  \item $f$ est homogène de degré~$2$, $\frac{\partial f}{\partial x}$ et $\frac{\partial f}{\partial y}$ sont homogènes de degré~$1$, donc
    ces trois fonctions tendent vers~$0$ en~$(0,0)$. Ainsi $f$ est de classe~$\mathcal{C}^1$.
  \item $\frac{\partial^2 f}{\partial x \partial y}(0,0)=0$, $\frac{\partial^2 f}{\partial y \partial x}(0,0)=1$.
\end{enumerate}
 \fincorrection
\nocorrection
\nocorrection
\nocorrection
\correction{004151}
$\frac{\partial^n(fg)}{\partial x^k\partial y^{n-k}} =
\sum_{i=0}^k\sum_{j=0}^{n-k} C_k^iC_{n-k}^j
\frac{\partial^{i+j}f}{\partial x^i\partial y^j}
\frac{\partial^{n-i-j}g}{\partial x^{k-i}\partial y^{n-k-j}}$.
\fincorrection
\correction{004152}
\begin{enumerate}
  \item pas de solution.
  \item $f(x,y) = \frac {y-1}{x+y+1} + {}$ cste.
  \item $f(x,y) = \frac{xy}{x+y} + {}$ cste.
  \item $f(x,y) = x^2+y^2+\frac xy + {}$ cste.
\end{enumerate}
 \fincorrection
\correction{004153}
$f(y) = \lambda e^{-y}$,
         $F(x,y) = \lambda\left({\frac{x^2}2 - (y+1)e^{-y}}\right)$.
 \fincorrection
\correction{004154}
$f(x,y) = g(x^2+y^2)$.
 \fincorrection
\correction{004155}
$f(y) = \frac yk$, $g(z) = kz+\ell$,
         $F(x,y) = \left(x^2+\frac{y^2}2\right)z + \frac {\ell y^2}{2k} + C$.
\fincorrection
\correction{004156}
\begin{enumerate}
  \item $f(x,y) = -\frac {\ln x}{xy} - \frac 1x$.
  \item $y = \frac {\ln x}{\lambda x - 1}$.
\end{enumerate}
 \fincorrection
\nocorrection
\nocorrection
\nocorrection
\correction{004160}
$\det(J_f) = \prod_{i<j}(x_i-x_j)$.
\fincorrection
\nocorrection
\correction{004162}
$f(\R^3) = \{ (u,v,w) \text{ tq } u+v > 0,\ ue^{2w}-v > 0\}$.
\fincorrection
\nocorrection
\correction{004164}
$y = -\frac x2 + \frac {9x^2}{16} + o (x^2)  \Rightarrow {}$
         tangente de pente $-\frac 12$, au dessus.
\fincorrection
\correction{004165}
\begin{enumerate}
  \item 
  \item $\varphi(x) = 1 + x -\frac23x^3 + o (x^3)$.
\end{enumerate}
 \fincorrection
\correction{004166}
$\varphi'(1) = -\frac13$, $\varphi''(1) = -\frac8{27}$.
\fincorrection
\nocorrection
\nocorrection
\correction{004169}
$x=1 - \frac{\lambda+\mu}5 + \frac{\lambda^2+\lambda\mu}{25} + o(\lambda^2+\mu^2)$.
\fincorrection
\correction{004170}
$x=\sh(a),\ y=\sh(b)  \Rightarrow  u=\sh(a+b),\ v=\sh(a+b)+\ch(a+b)$.
donc $f(\R^2)$ est inclus dans l'hyperbole d'équation $v^2-2uv=1$
et on doit avoir $v\ge u$ ce qui donne la branche supérieure.
\fincorrection
\nocorrection
\nocorrection
\correction{004173}
\begin{enumerate}
  \item $f(x,y) =  \int_{t=0}^1 (x\frac{\partial f}{\partial x}(tx,ty) + y\frac{\partial f}{\partial y}(tx,ty))\,d t$.
  \item 
  \item 
\end{enumerate}
 \fincorrection
\correction{004174}
\begin{enumerate}
  \item 
  \begin{enumerate}
    \item thm des trois cordes pour $u \mapsto f(x+uh)$ sur $[-1,1]$.
    \item Soit $\overline{B}_\infty(x,r)\subset U$ et $y$ tel
que $0<\|x-y\|_\infty \le r$. On note $t = \|x-y\|_\infty/r$ et $h=(y-x)/t$.
Alors $y = x+th$ et $\|x\pm h\| = r$ donc~:
$$|f(x)-f(y)| \le t\max(|f(x+h)-f(x)|,|f(x-h)-f(x)|) \le Mt$$
car la restriction de~$f$ aux côtés de $\overline{B}_\infty(x,r)$ est
bornée (fonction convexe en dimension~1).
  \end{enumerate}  
  \item La surface représentative de~$f$ est au dessus de ses plans tangents.
  \item 
\end{enumerate}
 \fincorrection
\nocorrection
\correction{004176}
\begin{enumerate}
  \item $d f_M(H) = \sum_{k=0}^\infty \frac1{k!}(M^{k-1}H + \dots + HM^{k-1})$.
  \item 
  \item les matrices $M_k = \begin{pmatrix}0&1/k\cr0&2i\pi\cr\end{pmatrix}$ sont toutes
semblables à~$M_\infty$ et ont même exponentielle $I$.
  \item $M_k=\begin{pmatrix}0&2\pi+1/k\cr-4\pi^2/(2\pi+1/k)&0\cr\end{pmatrix}$.
\end{enumerate}
 \fincorrection
\nocorrection
\nocorrection
\nocorrection
\nocorrection
\correction{004181}
\begin{enumerate}
  \item $xy/(x^2+y^2)$.
  \item $f$ est lipschitzienne pour $\|\ \|_\infty$.
\end{enumerate}
 \fincorrection
\correction{004182}
\begin{enumerate}
  \item $2x^4y^2 \le (x^4+y^2)^2$.
  \item $g_\theta(r)\sim r^2$.
  \item $f(x,x^2)=-x^4$. Donc $(0,0)$ n'est pas minimum local de~$f$.
\end{enumerate}
 \fincorrection
\nocorrection
\correction{004184}
\begin{enumerate}
  \item 
  \item $\|f(x)\|\to+\infty$ lorsque $\|x\|\to\infty$ donc $x \mapsto\|f(x)-a\|^2$
    admet un minimum sur~$\R^n$. En ce point on a pour tout $h\in\R^n$~:
    $(f(x)-a\mid d f_x(h)) = 0$ et $d f_x$ est surjective
    (linéaire injective en dimension finie) donc $f(x)=a$.
\end{enumerate}
 \fincorrection
\correction{004185}
\begin{enumerate}
  \item Si $u$ atteint son maximum en $(x,y)\in\Omega$ alors $d^2u(x,y)$ est
    négative, contradiction avec $\Delta u(x,y)>0$.
  \item Soit $u_p(x,y) = u(x,y) + (x^2+y^2)/p$~: $\Delta u_p = \Delta u + 2/p >0$
    donc $u_p$ relève du cas précédent.

    \def\maxl{\max\limits}
    On a $\maxl_\Omega u  + \frac Mp \ge \maxl_\Omega u_p = \maxl_{\overline\Omega\setminus\Omega} u_p
    \ge \maxl_{\overline\Omega\setminus\Omega} u$ et on passe à la limite.
  \item Soit $u_1(x,y) = u(x,y) + \alpha\ln(x^2+y^2)$ où $\alpha$ est
    tel que $M_1(r_1) = M_1(r_2)$, avec
    $M_1(r) = \max\limits_{x^2+y^2=r^2} (u(x,y))$.
    On a $\Delta u_1\ge 0$ d'où $M_1(r) \le M_1(r_1) = M_1(r_2)$ c'est-à-dire~:
    $$M(r) \le M(r_1) + \alpha\ln(r_1/r) = M(r_2)-\alpha\ln(r/r_2)
    = \frac{M(r_1)\ln(r_2/r)+M(r_2)\ln (r/r_1)}{\ln (r_2/r_1)}.$$
\end{enumerate}
 \fincorrection
\correction{004186}
\begin{enumerate}
  \item On pose $g(r,\theta) = f(r\cos\theta,r\sin\theta)$,
    $h(r) =  \int_{\theta=0}^{2\pi}f(r\cos\theta,r\sin\theta)\,d\theta$ et l'on a
    $0=\Delta f = \frac{\partial^2 g}{\partial r^2} + \frac1r\frac{\partial g}{\partial r} + \frac1{r^2}\frac{\partial^2 g}{\partial \theta^2}$
    d'où~:
    $$0=h''(r) + \frac1rh'(r) + \frac1{r^2}\Bigl[\frac{\partial g}{\partial \theta}(r,\theta)\Bigr]_{\theta=0}^{2\pi}$$
    Le crochet est nul par $2\pi$-périodicité de $g$ donc $h''(r) + \frac1rh'(r) = 0$
    soit $h'(r) = \frac Kr$ et $K=0$ par continuité de $h'$ en $0$.
  \item $\pi r^2f(0,0)$.
\end{enumerate}
 \fincorrection
\correction{004187}
\begin{enumerate}
  \item $\det(J_\varphi(x,y)) =  f'(x)f'(y) - g'(x)g'(y) > 0$ donc le théorème
    d'inversion locale s'applique, il suffit de vérifier l'injectivité de~$\varphi$.
    Si $\varphi(x,y) = \varphi(u,v)$ alors~:
    \begin{align*}&|x-u| \le |f(x)-f(u)| = |g(v)-g(y)| \le |v-y|\\
               &|v-y| \le |f(v)-f(y)| = |g(x)-g(u)| \le |x-u|\\ 
\end{align*}
    d'où $|x-u| \le |x-u|$ et il y a inégalité stricte si $v\ne y$ ce qui est
    absurde donc $v=y$ et de même $u=x$.
  \item On a $\varphi(x,y) = (u,v)$ si et seulement si $x = f^{-1}(u-g(y))$ et
    $y = f^{-1}(v - g(f^{-1}(u-g(y)))) = h(y)$. $h$ est $k^2$-lipschitzienne
    donc le théorème du point fixe s'applique.
\end{enumerate}
 \fincorrection
\correction{004188}
Sinon il existe $a\in\R^2$ telle que $g$~: $x \mapsto f(x)-(a\mid x)$
n'a pas de point critique, donc pas de minimum ni de maximum.
On a aussi $|g(x)|/\|x\|\to +\infty$ lorsque $\|x\|\to\infty$ d'où $\sup(g) = +\infty$
et $\inf(g) = -\infty$.
Considérons pour $r> 0$ $E_r = \{x\in\R^2\text{ tq }\|x\|\ge r\}$~: $g(E_r)$ est une partie connexe
de~$\R$ donc un intervalle, et $\sup(g(E_r)) = \sup(g) = +\infty$,
$\inf(g(E_r)) = \inf(g) = -\infty$, d'où $g(E_r) = \R$. Ainsi il existe des $x$ de normes arbitrairement
grandes tels que $g(x) = 0$ en contradiction avec la
propriété $|g(x)|/\|x\| \to +\infty$ lorsque $\|x\|\to\infty$.

Remarque~: l'hypothèse $f$ de classe $\mathcal{C}^2$ est surabondante, la classe $\mathcal{C}^1$
suffit à conclure.
\fincorrection
\correction{004189}
Remarques~:\par
-- la transformation $f \mapsto g$ est appelée {\it transformation de Legendre.}
On notera $g = f^*$ ci-dessous.\par
-- l'hypothèse \og\ $H_f$ est définie positive en tout point \fg\ implique que $f$ est convexe.

{\bf Étude d'un cas particulier~:} $f(x) = \alpha\|x\|^2 + \beta(x\mid a) + \gamma$
avec $\alpha \in\R^{+*}$, $\beta,\gamma\in\R$ et $a\in\R^n$.

Alors $(x\mid y) - f(x) =
- \alpha\Bigl\|x - \frac{y-\beta a}{2\alpha}\Bigr\|^2
+ \alpha\Bigl\|\frac{y-\beta a}{2\alpha}\Bigr\|^2 - \gamma$
d'où $f^*(y) = \alpha\Bigl\|\frac{y-\beta a}{2\alpha}\Bigr\|^2 - \gamma
= \alpha^*\|y\|^2 + \beta^*(y\mid a) + \gamma^*$
avec $\alpha^* = 1/4\alpha$, $\beta^* = -\beta/2\alpha$ et
$\gamma^* = \beta^2\|a\|^2/4\alpha - \gamma$.
Ainsi, $f^*$ a la même forme que $f$, et on vérifie immédiatement que $f^{**} = f$.

{\bf Cas général~:} on montre que $f^*$ est bien définie, vérifie les mêmes hypothèses
que~$f$ et que l'on a $f^{**} = f$.

{\bf 1.} Bonne définition de~$f^*$~: à $y$ fixé on a $(x\mid y) -f(x) \to -\infty$ lorsque $\|x\|\to\infty$
donc le sup existe et est un max, atteint en un point $x$ tel que $\nabla f(x) = y$.
Ce point $x$ est unique~: en effet, si $h\in\R^n\setminus\{0\}$ alors la fonction
$\R\ni t \mapsto(h\mid\nabla f(x+th))$ est strictement croissante (définie-positivité
de~$H_f$) ce qui implique $\nabla f(x+h) \ne \nabla f(x)$.
Ainsi, $$f^*(y) = (y\mid y^*) - f(y^*) \text{ avec } \nabla f(y^*) = y.$$

{\bf 2.} $f^*$ est $\mathcal{C}^2$ et $H_{f^*}$ est définie-positive~:
d'après ce qui précède, la fonction $\nabla f$ est un $\mathcal{C}^1$ difféomorphisme de~$\R^n$~;
sa différentielle est l'endomorphisme de~$\R^n$ de matrice $H_f$ dans la base
canonique de~$\R^n$. Donc $f^*$ est de classe $\mathcal{C}^1$ et pour $y,h\in\R^n$~:
$$d (f^*)_y(h) = (\nabla f^*(y)\mid h) = (h\mid y^*) + (y\mid d y^*(h)) - (\nabla f(y^*)\mid d y^*(h))
= (h\mid y^*)$$
puisque $\nabla f(y^*) = y$. On en déduit~: $\nabla f^*(y) = y^*$,
puis $H_{f^*}(y) = (H_f(y^*))^{-1}$, matrice symétrique définie positive.

{\bf 3.} $f^*(y)/\|y\|\to +\infty$ lorsque $\|y\|\to\infty$~:
soit $a>1$ et $M_a = \sup\{f(x),\ \|x\|\le a\}$.
Pour $y\in\R^n\setminus\{0\}$ et $x=ay/\|y\|$ on a
$f^*(y) \ge (x\mid y) - f(x) \ge a\|y\|-M_a\ge (a-1)\|y\|$ si $\|y\|$ est assez grand.

{\bf 4.} $f^{**} = f$~: car pour $x\in\R^n$ on a $f^{**}(x) = (x\mid y) - f^*(y)$ où $y$ est
défini par $\nabla f^*(y) = x$, c'est-à-dire $y^* = x$ et donc
$f^{**}(x) = (y^*\mid y) - f^*(y) = f(y^*) = f(x)$.
\fincorrection
\correction{004190}
\begin{enumerate}
  \item ???
  \item On a pour $a,b\in\R$~: $|\varphi(a+b) - \varphi(a) - b\varphi'(a)| \le \frac12\|\varphi''\|_\infty b^2$.

Donc pour $f,h\in E$~: $|T(f+h) - T(f) - (h\mid \varphi'\circ f)| \le \frac12\|\varphi''\|_\infty \|h\|^2$,
ce qui prouve que $T$ est différentiable en~$f$ de différentielle $h \mapsto(h\mid\varphi'\circ f)$.
On en déduit alors que~$T$ est continue en~$f$.
\end{enumerate}
 \fincorrection
\correction{005556}
\begin{enumerate}
 \item  Posons $\Delta=\{(x,y)/\;y\neq0\}$. $f$ est continue sur $\Rr^2\setminus\Delta$ en vertu de théorèmes généraux. Soit $x_0\in\Rr$.

\begin{center}
$|f(x,y)-f(x_0,0)|=\left\{
\begin{array}{l}
0\;\text{si}\;y=0\\
y^2\left|\sin\left(\frac{x}{y}\right)\right|\;\text{si}\;y\neq0
\end{array}
\right.\leqslant y^2$.
\end{center}
Comme $\displaystyle\lim_{(x,y)\rightarrow(x_0,0)}y^2=0$, $\displaystyle\lim_{(x,y)\rightarrow(x_0,0)}|f(x,y)-f(x_0,0)|=0$ et donc $f$ est continue en $(x_0,0)$. Finalement,

\begin{center}
\shadowbox{
$f$ est continue sur $\Rr^2$.
}
\end{center}
 \item  \textbullet~$f$ est de classe $C^2$ sur $\Rr^2\setminus\Delta$. En particulier, d'après le théorème de \textsc{Schwarz}, $\frac{\partial^2f}{\partial x\partial y}=\frac{\partial^2f}{\partial y\partial x}$ sur $\Delta$. pour $(x,y)\in\Rr^2\setminus\Delta$,

\begin{center}
$\frac{\partial f}{\partial x}(x,y)=y\cos\left(\frac{x}{y}\right)$ et $\frac{\partial f}{\partial y}(x,y)=2y\sin\left(\frac{x}{y}\right)-x\cos\left(\frac{x}{y}\right)$,
\end{center}
puis

\begin{center}
$\frac{\partial^2f}{\partial x^2}(x,y)=-\sin\left(\frac{x}{y}\right)$, $\frac{\partial^2f}{\partial x\partial y}(x,y)=\cos\left(\frac{x}{y}\right)-\frac{x}{y}\sin\left(\frac{x}{y}\right)$,
\end{center}
et enfin

\begin{center}
$\frac{\partial^2f}{\partial y^2}(x,y)=2\sin\left(\frac{x}{y}\right)-2\frac{x}{y}\cos\left(\frac{x}{y}\right)-\frac{x^2}{y^2}\sin\left(\frac{x}{y}\right)$.
\end{center}
\textbullet~\textbf{Existence de $\frac{\partial f}{\partial x}(x_0,0)$.} Pour $x\neq x_0$,  $\frac{f(x,0)-f(x_0,0)}{x-x_0}=0$ et donc$\frac{f(x,0)-f(x_0,0)}{x-x_0}\underset{x\rightarrow x_0}{\rightarrow}0$. On en déduit que $\frac{\partial f}{\partial x}(x_0,0)$ existe et $\frac{\partial f}{\partial x}(x_0,0)=0$. En résumé, $f$ admet une dérivée partielle par rapport à sa première variable sur $\Rr^2$ définie par

\begin{center}
\shadowbox{
$\forall(x,y)\in\Rr^2,\;\frac{\partial f}{\partial x}(x,y)=\left\{
\begin{array}{l}
0\;\text{si}\;y=0\\
y\cos\left(\frac{x}{y}\right)\;\text{si}\;y\neq0
\end{array}
\right.$.
}
\end{center}
\textbullet~\textbf{Existence de $\frac{\partial f}{\partial y}(x_0,0)$.} Soit $x_0\in\Rr$. Pour $y\neq0$,  

\begin{center}
$\left|\frac{f(x_0,y)-f(x_0,0)}{y-0}\right|=\left\{
\begin{array}{l}
0\;\text{si}\;y=0\\
y\left|\sin\left(\frac{x_0}{y}\right)\right|\;\text{si}\;y\neq0
\end{array}
\right.\leqslant|y|.
$
\end{center} 

et donc$\frac{f(x_0,y)-f(x_0,0)}{y-0}\underset{y\rightarrow0}{\rightarrow}0$. On en déduit que $\frac{\partial f}{\partial y}(x_0,0)$ existe et $\frac{\partial f}{\partial y}(x_0,0)=0$. En résumé, $f$ admet une dérivée partielle par rapport à sa deuxième variable sur $\Rr^2$ définie par

\begin{center}
\shadowbox{
$\forall(x,y)\in\Rr^2,\;\frac{\partial f}{\partial y}(x,y)=\left\{
\begin{array}{l}
0\;\text{si}\;y=0\\
2y\sin\left(\frac{x}{y}\right)-x\cos\left(\frac{x}{y}\right)\;\text{si}\;y\neq0
\end{array}
\right.$.
}
\end{center}
\textbullet~\textbf{Existence de $\frac{\partial^2f}{\partial x\partial y}(0,0)$.} Pour $x\neq0$,

\begin{center}
$\frac{\frac{\partial f}{\partial y}(x,0)-\frac{\partial f}{\partial y}(0,0)}{x-0}=0$
\end{center}
et donc $\frac{\frac{\partial f}{\partial y}(x,0)-\frac{\partial f}{\partial y}(0,0)}{x-0}$ tend vers $0$ quand $x$ tend vers $0$. On en déduit que $\frac{\partial^2f}{\partial x\partial y}(0,0)$ existe et

\begin{center}
\shadowbox{
$\frac{\partial^2f}{\partial x\partial y}(0,0)=0$.
}
\end{center}
\textbullet~\textbf{Existence de $\frac{\partial^2f}{\partial y\partial x}(0,0)$.} Pour $y\neq0$,

\begin{center}
$\frac{\frac{\partial f}{\partial x}(0,y)-\frac{\partial f}{\partial x}(0,0)}{y-0}=\frac{y\cos\left(\frac{0}{y}\right)}{y}=1$
\end{center}
et donc $\frac{\frac{\partial f}{\partial x}(0,y)-\frac{\partial f}{\partial x}(0,0)}{y-0}$ tend vers $1$ quand $y$ tend vers $0$. On en déduit que $\frac{\partial^2f}{\partial y\partial x}(0,0)$ existe et

\begin{center}
\shadowbox{
$\frac{\partial^2f}{\partial y\partial x}(0,0)=1$.
}
\end{center}
\end{enumerate}
\fincorrection
\correction{005557}
Pour $(x,y)\in\Rr^2$, $\frac{\cos(2x)}{\ch(2y)}\in[-1,1]$. Plus précisément, quand $x$ décrit $\Rr$, $\frac{\cos(2x)}{\ch(2\times0)}$ décrit $[-1,1]$ et donc quand $(x,y)$ décrit $\Rr^2$, $\frac{\cos(2x)}{\ch(2y)}$ décrit $[-1,1]$. 
On suppose déjà que $f$ est de classe $C^2$ sur $[-1,1]$. L'application $g$ est alors de classe $C^2$ sur $\Rr^2$ et pour $(x,y)\in\Rr^2$,

\begin{center}
$\frac{\partial g}{\partial x}(x,y)=-\frac{2\sin(2x)}{\ch(2y)}f'\left(\frac{\cos2x}{\ch2y}\right)$ puis $\frac{\partial^2g}{\partial x^2}(x,y)=-\frac{4\cos(2x)}{\ch(2y)}f'\left(\frac{\cos2x}{\ch2y}\right)+\frac{4\sin^2(2x)}{\ch^2(2y)}f''\left(\frac{\cos2x}{\ch2y}\right)$.
\end{center}
Ensuite,

\begin{center}
$\frac{\partial g}{\partial y}(x,y)=-\frac{2\cos(2x)\sh(2y)}{\ch^2(2y)}f'\left(\frac{\cos2x}{\ch2y}\right)$
\end{center}
puis 
\begin{center}
$\frac{\partial^2g}{\partial y^2}(x,y)=-\frac{4\cos(2x)}{\ch(2y)}f'\left(\frac{\cos2x}{\ch2y}\right)
-2\cos(2x)\sh(2y)\frac{-4\sh(2y)}{\ch^3(2y)}f'\left(\frac{\cos2x}{\ch2y}\right)
+\frac{4\cos^2(2x)\sh^2(2y)}{\ch^4(2y)}f''\left(\frac{\cos2x}{\ch2y}\right)$.
\end{center}
Mais alors

\begin{align*}\ensuremath
\Delta g(x,y)&=\frac{-8\cos(2x)\ch^2(2y)+8\cos(2x)\sh^2(2y)}{\ch^3(2y)}f'\left(\frac{\cos2x}{\ch2y}\right)+\frac{4\sin^2(2x)\ch^2(2y)+4\cos^2(2x)\sh^2(2y)}{\ch^4(2y)}f''\left(\frac{\cos2x}{\ch2y}\right)\\
 &=\frac{-8\cos(2x)}{\ch^3(2y)}f'\left(\frac{\cos2x}{\ch2y}\right)+\frac{4(1-\cos^2(2x))\ch^2(2y)+4\cos^2(2x)(\ch^2(2y)-1)}{\ch^4(2y)}f''\left(\frac{\cos2x}{\ch2y}\right)\\
 &=\frac{-8\cos(2x)}{\ch^3(2y)}f'\left(\frac{\cos2x}{\ch2y}\right)+\frac{4\ch^2(2y)-4\cos^2(2x)}{\ch^4(2y)}f''\left(\frac{\cos2x}{\ch2y}\right)\\
 &=\frac{4}{\ch^2(2y)}\left(-2\frac{\cos(2x)}{\ch(2y)}f'\left(\frac{\cos2x}{\ch2y}\right)+\left(1-\frac{\cos^2(2x)}{\ch^2(2y)}\right)f''\left(\frac{\cos2x}{\ch2y}\right)\right).
\end{align*}
Par suite,

\begin{align*}\ensuremath
\Delta g=0&\Leftrightarrow\forall(x,y)\in\Rr^2,\;-2\frac{\cos(2x)}{\ch(2y)}f'\left(\frac{\cos2x}{\ch2y}\right)+\left(1-\frac{\cos^2(2x)}{\ch^2(2y)}\right)f''\left(\frac{\cos2x}{\ch2y}\right)=0\\
 &\Leftrightarrow\forall t\in[-1,1],\;-2tf'(t)+(1-t^2)f''(t)=0\Leftrightarrow\forall t\in[-1,1],((1-t^2)f')'(t)=0\\
 &\Leftrightarrow\exists \lambda\in\Rr,\;\forall t\in[-1,1],\;(1-t^2)f'(t)=\lambda.
\end{align*}
Le choix $\lambda\neq0$ ne fournit pas de solution sur $[-1,1]$. Donc $\lambda=0$ puis $f'=0$ puis $f$ constante ce qui est exclu. Donc, on ne peut pas poursuivre sur $[-1,1]$.
On cherche dorénavant $f$ de classe $C^2$ sur $]-1,1[$ de sorte que $g$ est de classe $C^2$ sur $\Rr^2\setminus\left\{\left(\frac{k\pi}{2},0\right),\;k\in\Zz\right\}$.

\begin{align*}\ensuremath
f\;\text{solution}&\Leftrightarrow\exists \lambda\in\Rr^*,\;\forall t\in]-1,1[,\;(1-t^2)f'(t)=\lambda\Leftrightarrow\exists \lambda\in\Rr^*/\;\forall t\in]-1,1[,\;f'(t)=\frac{\lambda}{1-t^2}\\
 &\Leftrightarrow\exists(\lambda,\mu)\in\Rr^*\times\Rr/\;\forall t\in]-1,1[,\;f(t)=\lambda\Argth t+\mu.
\end{align*}
\fincorrection
\correction{005893}
On dérive par rapport à $\lambda$ les deux membres de l'égalité $f(\lambda x) =\lambda^rf(x)$ et on obtient

\begin{center}
$\forall x=(x_1,...,x_n)\in\Rr^n$, $\forall\lambda>0$, $\sum_{i=1}^{n}x_i \frac{\partial f}{\partial x_i}(\lambda x)=r\lambda^{r-1}f(x)$,
\end{center}

et pour $\lambda=1$, on obtient

\begin{center}
$\forall x = (x_1,...,x_n)\in \Rr^n$ $\sum_{i=1}^{n}x_i \frac{\partial f}{\partial x_i}(x) = rf(x)$.
\end{center}
\fincorrection
\correction{002635}
\begin{align*}
\mathrm{d} \left(\ln(xy)\right)&= \frac{\mathrm d(xy)}{xy}
= \frac{x\mathrm dy+y\mathrm dx}{xy}= \frac{\mathrm dy}{y}
+\frac{\mathrm dx}{x} ;
\\
\mathrm{d} \left(xyz(1+\sinh(yz))\right)&=(1+\sinh(yz))\mathrm{d}(xyz)
+xyz\mathrm{d}(\sinh(yz))
\\
&=
yz(1+\sinh(yz))\mathrm{d}x
+
xz(1+\sinh(yz))\mathrm{d}y
+
xy(1+\sinh(yz))\mathrm{d}z
\\
&\quad
+xyz\cosh(yz)\mathrm{d}(yz)
\\
&=
yz(1+\sinh(yz))\mathrm{d}x
+
xz(1+\sinh(yz))\mathrm{d}y
+
xy(1+\sinh(yz))\mathrm{d}z
\\
&\quad
+xyz^2\cosh(yz)\mathrm{d}y
+xy^2z\cosh(yz)\mathrm{d}z
\\
&=
yz(1+\sinh(yz))\mathrm{d}x
\\
&\quad
+
xz(1+\sinh(yz) +yz\cosh(yz))
\mathrm{d}y
\\
&\quad +
xy(1+\sinh(yz)+ yz\cosh(yz))\mathrm{d}z ;
\end{align*}

\begin{align*}
\mathrm{d} \left(\sin(x^2y) \mathrm e^{x-y}\right)&=
\left(\cos(x^2y)\mathrm e^{x-y}\right)\mathrm{d}(x^2y)
+\sin(x^2y)\mathrm e^{x-y}\mathrm{d}(x-y)
\\
&=
x^2\cos(x^2y)\mathrm e^{x-y}\mathrm{d}y
+
2xy\cos(x^2y)\mathrm e^{x-y}\mathrm{d}x
\\
&\quad
+\sin(x^2y)\mathrm e^{x-y}\mathrm{d}x-
\sin(x^2y)\mathrm e^{x-y}
\mathrm{d}y
\\
&=
\left(x^2\cos(x^2y)x-
\sin(x^2y)\right)\mathrm e^{x-y}\mathrm{d}y
\\
&\quad +
\left(2xy\cos(x^2y)+\sin(x^2y)\right) \mathrm e^{x-y}\mathrm{d}x .
\end{align*}
\fincorrection
\correction{002636}
\begin{enumerate} 
\item La forme diff\'erentielle
$x^2y^2 \mathrm{d} x + x^3 y\mathrm{d} y$  de degr\'e 1 n'est pas ferm\'ee
car la forme diff\'erentielle de degr\'e 2
\[
d(x^2y^2 \mathrm{d} x + x^3 y\mathrm{d} y) 
= 2x^2y \mathrm{d} y\mathrm{d} x + 3x^2 y\mathrm{d} x\mathrm{d} y
= x^2 y \mathrm{d} x\mathrm{d} y 
\]
est non nulle. Par cons\'equent, une fonction $g\colon \R^2\to\R$ du 
type cherch\'e
ne peut pas exister.

\item Une fonction $b$ du 
type cherch\'e doit satisfaire \`a l'\'equation diff\'erentielle
partielle
\[
2x^2y -\frac{\partial b}{\partial x} =0
\]
d'o\`u $b(x,y)= \frac 23 x^3 y +k(y)$ o\`u $k$ est une fonction
de la variable $y$. Une fonction $g$ correspondante doit alors satisfaire
aux \'equations diff\'erentielles partielles
\[
\frac{\partial g}{\partial x}=x^2 y^2,
\
\frac{\partial g}{\partial y}=\frac 23 x^3 y +k(y).
\]
Il s'ensuit que $g$ est de la forme
$g(x,y)=\frac 13 x^3 y^2 +K(y)$ o\`u $K$ est une fonction de la varriable $y$.
\end{enumerate}
\fincorrection
\correction{002637}
\begin{enumerate}  
\item Un calcul imm\'ediat donne $\mathrm d u+\mathrm d v =\mathrm d g$.

\item Par cons\'equent, $g=u+v+c$ o\`u la constante $c$ est d\'etermin\'ee par
la condition
\[
3=g(1,1)= u(1,1)+v(1,1)+c=1+1+c
\]
d'o\`u $c=1$.

\item  Un calcul direct montre que l'application r\'eciproque 
\[
k \colon \R_{>0}\times\R_{>0} \longrightarrow \R_{>0}\times\R_{>0}
\]
de $h$ 
est donn\'ee par la formule
\[
k(u,v)= (x(u,v),y(u,v))= 
\left(\left(\frac{u^2}{v}\right)^{1/3},\left(\frac{v^2}{u}\right)^{1/3}\right).
\]
\item $\mathrm d (g \circ k)=\mathrm d (u \circ k)+ \mathrm d (v \circ k)
=\mathrm d u +\mathrm d v$ car $u(k(u,v))=u$ et $v(k(u,v))=v$.

\item  Un calcul imm\'ediat donne 
$$J_h=\left [\begin{matrix} 2xy & x^2\\ y^2 &2xy \end{matrix}\right ], \quad 
J_k=\left [\begin{matrix} 
\frac 23 (uv)^{-1/3} & -\frac {u^{2/3}}{3v^{4/3}}
\\  -\frac {v^{2/3}}{3u^{4/3}} &\frac 23 (uv)^{-1/3} 
\end{matrix}
\right ]$$ d'o\`u $J_h(x,y)J_{k}(h(x,y))=I_2$.
\end{enumerate}
\fincorrection
\correction{002638}
\[
\mathrm d\sin(xyz)= yz\cos(xyz)\mathrm dx
+zx\cos(xyz)\mathrm dy +xy\cos(xyz)\mathrm dz
\]
d'o\`u la matrice hessienne
\[
\left[ \begin{smallmatrix}
-y^2z^2\sin(xyz)          &z\cos(xyz) -xyz^2\sin(xyz)&y\cos(xyz)-xy^2z\sin(xyz)
\\
z\cos(xyz)-xyz^2\sin(xyz)&-x^2z^2\sin(xyz)           &x\cos(xyz)-x^2yz\sin(xyz)
\\
y\cos(xyz)-xy^2z\sin(xyz)&y\cos(xyz) -x^2yz\sin(xyz)&-x^2y^2\sin(xyz)
\end{smallmatrix}
\right] .
\]
De m\^eme
\begin{align*}
\mathrm d(\sin^2(y/x))&= -{2y}{x^{-2}}\sin(y/x)\cos(y/x)\mathrm dx 
+ 2x^{-1} \sin(y/x)\cos(y/x) \mathrm dy
\\
&= \sin(2y/x)\left( -\frac y {x^2} dx + \frac 1x dy\right)
\end{align*}
d'o\`u la matrice hessienne
\[
\left[ \begin{matrix}
{2y}{x^{-3}}\sin(2y/x)+ {2y^2}{x^{-4}}\cos(2y/x) 
&
-{x^{-2}}\sin(2y/x)  -{2y}{x^{-3}}\cos(2y/x)
\\
-x^{-2} \sin(2y/x)  -2yx^{-3} \cos(2y/x)&
2x^{-2} \cos(2y/x)
\end{matrix}
\right] .
\]
\fincorrection
\correction{002639}
\begin{enumerate}
\item $\frac{\partial }{\partial r}\left (r \frac{\partial F}{\partial r}\right)=
\frac{\partial F}{\partial r} +r \frac{\partial^2 F}{\partial r^2}$

\item $\frac{\partial F}{\partial r}=
\frac{\partial f}{\partial x}\frac xr+ \frac{\partial f}{\partial y}\frac yr$

\item $\frac{\partial F}{\partial \theta}=
\frac{\partial f}{\partial x}\frac{\partial x}{\partial \theta}
+ \frac{\partial f}{\partial y}\frac{\partial y}{\partial \theta}
= -y \frac{\partial f}{\partial x}+x\frac{\partial f}{\partial y}
$

\item En prenant la somme des trois \'equations suivantes
\begin{align*}
r^2 \frac{\partial^2 F}{\partial r^2}&=
x^2 \frac{\partial^2 f}{\partial x^2}
+2xy \frac{\partial^2 f}{\partial x \partial y}
+
y^2 \frac{\partial^2 f}{\partial y^2}
\\
r\frac{\partial F}{\partial r}&=
x\frac{\partial f}{\partial x}+ y\frac{\partial f}{\partial y}
\\
\frac{\partial^2 F}{\partial \theta^2}&=
x^2 \frac{\partial^2 f}{\partial y^2}
-2xy \frac{\partial^2 f}{\partial x \partial y}
+
y^2 \frac{\partial^2 f}{\partial x^2}
-
x \frac{\partial f}{\partial x}
-
y\frac{\partial f}{\partial y}
\end{align*}
on trouve le r\'esultat cherch\'e.
\end{enumerate}
\fincorrection
\correction{002640}
\begin{enumerate}  
\item Avec $\frac{\partial }{\partial u}
=1/2(\frac{\partial}{\partial x}+\frac{\partial}{\partial t})$ et
$\frac{\partial }{\partial v}
=1/2(-\frac{\partial}{\partial x}+\frac{\partial}{\partial t})$
nous obtenons les identit\'es
\begin{align*}
\frac{\partial F}{\partial u}&= 
\frac 12\frac{\partial f}{\partial x}+\frac 12\frac{\partial f}{\partial t}
\\
\frac{\partial F}{\partial v}&= 
-\frac 12\frac{\partial f}{\partial x}+\frac 12\frac{\partial f}{\partial t}
\\
\frac{\partial^2 F}{\partial u \partial v}&= 
-\frac 14\frac{\partial^2 f}{\partial x^2}
+\frac 14\frac{\partial^2 f}{\partial t^2}
\end{align*}
d'o\`u pour que $f$ satisfasse \`a l'\'equation \eqref{eq:ondes} il faut et il suffit que $F$ satisfasse \`a l'\'equation
 \eqref{eq:ondesbis}.

\item Supposons que $F$ satisfasse \`a l'\'equation \eqref{eq:ondesbis}.
Alors la fonction $\frac{\partial F}{\partial u}$ est une fonction
disons $h_1$ seulement de la variable $u$
et la fonction $\frac{\partial F}{\partial v}$ est une fonction
disons $h_2$ seulement de la variable $v$.
Par cons\'equent,
$F(u,v)= g_1(u)+g_2(v)$ o\`u $g'_1=h_1$ et $g'_2=h_2$.

\item La solution g\'en\'erale de
\eqref{eq:ondes} s'\'ecrit alors
\[
f(x,t)=g_1(u)+g_2(v)=g_1(x+t)+g_2(t-x) .
\]
La fonction $g_1$ d\'ecrit
une onde qui se d\'eplace vers la droite et 
la fonction $g_1$ d\'ecrit
une onde qui se d\'eplace vers la gauche.

Enfin, pour trouver la solution unique
satisfaisant aux condition initiales \eqref{in} nous constatons que
les conditions initiales entra\^ \i nent les identit\'es
\begin{align*}
f(x,0)&=g_1(x)+g_2(-x)=\sin x
\\
\frac {\partial f}{\partial x}(x,0)&=g'_1(x)-g'_2(-x)=\cos x
\\
\frac {\partial f}{\partial t}(x,0)&=g'_1(x)+g'_2(-x)=-\cos x
\end{align*}
d'o\`u $g'_1=0$ et $g'_2(-x)=-\cos x$, c.a.d. 
$g_2(x)=\sin (-x)$. Par cons\'equent,
la solution unique cherch\'ee $f$ s'\'ecrit
\[
f(x,t)= \sin(x-t).
\]
\end{enumerate}
\fincorrection
\correction{005889}
On pose $D=\{(x,0),\;x\in\Rr\}$ puis $\Omega=\Rr^2\setminus D$.

\textbullet~$f$ est définie sur $\Rr^2$.

\textbullet~$f$ est de classe $C^1$ sur $\Omega$ en vertu de théorèmes généraux et pour $(x,y)\in\Omega$,

\begin{center}
$ \frac{\partial f}{\partial x}(x,y)=y\cos\left( \frac{x}{y}\right)$ et $ \frac{\partial f}{\partial y}(x,y)=2y\sin\left( \frac{x}{y}\right)-x\cos\left( \frac{x}{y}\right)$.
\end{center}

\textbullet~Etudions la continuité de $f$ en $(0,0)$. Pour $(x,y)\neq(0,0)$,

\begin{center}
$|f(x,y)-f(0,0)|=\left\{
\begin{array}{l}
y^2\left|\sin\left( \frac{x}{y}\right)\right|\;\text{si}\;y\neq0\\
\rule{0mm}{5mm}0\;\text{si}\;y=0
\end{array}
\right.\leqslant\left\{
\begin{array}{l}
y^2\;\text{si}\;y\neq0\\
\rule{0mm}{5mm}0\;\text{si}\;y=0
\end{array}
\right.\leqslant y^2$.
\end{center}

Comme $y^2$ tend vers $0$ quand $(x,y)$ tend vers $0$, $\displaystyle\lim_{\substack{(x,y)\rightarrow(0,0)\\(x,y)\neq(0,0)}}f(x,y)=f(0,0)$ et donc $f$ est continue en $(0,0)$ puis

\begin{center}
$f$ est continue sur $\Rr^2$.
\end{center}

\textbullet~Etudions l'existence et la valeur éventuelle de $ \frac{\partial f}{\partial x}(x_0,0)$, $x_0$ réel donné. Pour $x\neq x_0$,

\begin{center}
$ \frac{f(x,0)-f(x_0,0)}{x-x_0}= \frac{0-0}{x-x_0}=0$.
\end{center}

Donc $ \frac{f(x,x_0)-f(x_0,0)}{x-x_0}$ tend vers $0$ quand $x$ tend vers $x_0$. On en déduit que $ \frac{\partial f}{\partial x}(x_0,0)$ existe et $ \frac{\partial f}{\partial x}(x_0,0)=0$. Finalement, la fonction $ \frac{\partial f}{\partial x}$ est définie sur $\Rr^2$ par

\begin{center}
$\forall (x,y)\in\Rr^2$, $ \frac{\partial f}{\partial x}(x,y)=\left\{
\begin{array}{l}
y\cos\left( \frac{x}{y}\right)\;\text{si}\;y\neq0\\
\rule{0mm}{5mm}0\;\text{si}\;y=0
\end{array}
\right.$.
\end{center}

\textbullet~Etudions l'existence et la valeur éventuelle de $ \frac{\partial f}{\partial y}(x_0,0)$, $x_0$ réel donné. Pour $y\neq 0$,

\begin{center}
$ \frac{f(x_0,y)-f(x_0,0)}{y-0}= \frac{y^2\sin\left( \frac{x_0}{y}\right)}{y}=y\sin\left( \frac{x_0}{y}\right)$.
\end{center}

On en déduit que $\left| \frac{f(x_0,y)-f(x_0,0)}{y-0}\right|\leqslant|y|$ puis que $ \frac{f(x_0,y)-f(x_0,0)}{y-0}$ tend vers $0$ quand $y$ tend vers $0$. Par suite, $ \frac{\partial f}{\partial y}(x_0,0)$ existe et $ \frac{\partial f}{\partial y}(x_0,0)=0$. Finalement, la fonction $ \frac{\partial f}{\partial y}$ est définie sur $\Rr^2$ par

\begin{center}
$\forall (x,y)\in\Rr^2$, $ \frac{\partial f}{\partial y}(x,y)=\left\{
\begin{array}{l}
2y\sin\left( \frac{x}{y}\right)-x\cos\left( \frac{x}{y}\right)\;\text{si}\;y\neq0\\
\rule{0mm}{5mm}0\;\text{si}\;y=0
\end{array}
\right.$.
\end{center}

\textbullet~Etudions la continuité de $ \frac{\partial f}{\partial x}$ en $(x_0,0)$, $x_0$ réel donné. Pour $(x,y)\in\Rr^2$,

\begin{center}
$\left| \frac{\partial f}{\partial x}(x,y)- \frac{\partial f}{\partial x}(x_0,0)\right|=\left\{
\begin{array}{l}
|y|\left|\cos\left( \frac{x}{y}\right)\right|\;\text{si}\;y\neq0\\
\rule{0mm}{5mm}0\;\text{si}\;y=0
\end{array}
\right.\leqslant|y|$.
\end{center}

Quand $(x,y)$ tend vers $(0,0)$, $|y|$ tend vers $0$ et donc $ \frac{\partial f}{\partial x}(x,y)$ tend vers $ \frac{\partial f}{\partial x}(x_0,0)$ quand $(x,y)$ tend vers $(x_0,0)$. La fonction $ \frac{\partial f}{\partial x}$ est donc continue en $(x_0,0)$ et finalement

\begin{center}
la fonction $ \frac{\partial f}{\partial x}$ est continue sur $\Rr^2$.
\end{center}

\textbullet~Etudions la continuité de $ \frac{\partial f}{\partial y}$ en $(x_0,0)$, $x_0$ réel donné. Supposons tout d'abord $x_0=0$. Pour $(x,y)\in\Rr^2$,

\begin{center}
$\left| \frac{\partial f}{\partial y}(x,y)- \frac{\partial f}{\partial y}(0,0)\right|=\left\{
\begin{array}{l}
\left|2y\sin\left( \frac{x}{y}\right)-x\cos\left( \frac{x}{y}\right)\right|\;\text{si}\;y\neq0\\
\rule{0mm}{5mm}0\;\text{si}\;y=0
\end{array}
\right.\leqslant2|y|+|x|$.
\end{center}

Quand $(x,y)$ tend vers $(0,0)$, $|x|+2|y|$ tend vers $0$ et donc $ \frac{\partial f}{\partial y}(x,y)$ tend vers $ \frac{\partial f}{\partial y}(0,0)$ quand $(x,y)$ tend vers $(0,0)$. 

Supposons maintenant $x_0\neq0$. Pour $y\neq0$, $ \frac{\partial f}{\partial y}(x_0,y)=2y\sin\left( \frac{x_0}{y}\right)-x_0\cos\left( \frac{x_0}{y}\right)$. Quand $y$ tend vers $0$, $2y\sin\left( \frac{x_0}{y}\right)$ tend vers $0$ car $\left|2y\sin\left( \frac{x_0}{y}\right)\right|$ et $x_0\cos\left( \frac{x_0}{y}\right)$ n'a pas de limite réelle car $x_0\neq0$. Donc $ \frac{\partial f}{\partial y}(x_0,y)$ n'a pas de limite quand $y$ tend vers $0$ et la fonction $ \frac{\partial f}{\partial y}$ n'est pas continue en $(x_0,0)$ si $x_0\neq0$. On a montré que

\begin{center}
$f$ est de classe $C^1$ sur $\Omega\cup\{(0,0)\}$.
\end{center}

\textbullet~Etudions l'existence et la valeur éventuelle de $ \frac{\partial^2f}{\partial x\partial y}(0,0)$. Pour $x\neq0$,

\begin{center} 
$ \frac{ \frac{\partial f}{\partial y}(x,0)- \frac{\partial f}{\partial y}(0,0)}{x-0}= \frac{0-0}{x}=0$.
\end{center}

Donc $ \frac{ \frac{\partial f}{\partial y}(x,0)- \frac{\partial f}{\partial y}(0,0)}{x-0}$ tend vers $0$ quand $x$ tend vers $0$. On en déduit que $ \frac{\partial^2}{\partial x\partial y}(0,0)$ existe et $ \frac{\partial^2}{\partial x\partial y}(0,0)=0$. 

\textbullet~Etudions l'existence et la valeur éventuelle de $ \frac{\partial^2f}{\partial y\partial x}(0,0)$. Pour $y\neq0$,

\begin{center} 
$ \frac{ \frac{\partial f}{\partial x}(0,y)- \frac{\partial f}{\partial x}(0,0)}{y-0}= \frac{y\cos\left( \frac{0}{y}\right)}{y}=1$.
\end{center}

Donc $ \frac{ \frac{\partial f}{\partial x}(0,y)- \frac{\partial f}{\partial x}(0,0)}{y-0}$ tend vers $1$ quand $y$ tend vers $0$. On en déduit que $ \frac{\partial^2}{\partial y\partial x}(0,0)$ existe et $ \frac{\partial^2}{\partial y\partial x}(0,0)=1$. On a montré que $ \frac{\partial^2}{\partial x\partial y}(0,0)$ et $ \frac{\partial^2}{\partial y\partial x}(0,0)$ existent et sont différents. D'après le théorème de \textsc{Schwarz}, $f$ n'est pas de classe $C^2$ sur $\Omega\cup\{(0,0)\}$.
\fincorrection
\correction{005904}
Puisque la fonction $\ch$ ne s'annule pas sur $\Rr$, $g$ est de classe $C^2$ sur $\Rr^2$ et pour $(x,y)\in\Rr^2$,

\begin{center}
$ \frac{\partial g}{\partial x}(x,y)=-2 \frac{\sin(2x)}{\ch(2y)}f'\left( \frac{\cos(2x)}{\ch(2y)}\right)$
\end{center}

puis

\begin{align*}\ensuremath
 \frac{\partial^2g}{\partial x^2}(x,y)&=-4 \frac{\cos(2x)}{\ch(2y)}f'\left( \frac{\cos(2x)}{\ch(2y)}\right)+4 \frac{\sin^2(2x)}{\ch^2(2y)}f''\left( \frac{\cos(2x)}{\ch(2y)}\right)\\
 &=-4 \frac{\cos(2x)}{\ch(2y)}f'\left( \frac{\cos(2x)}{\ch(2y)}\right)+4 \frac{1-\cos^2(2x)}{\ch^2(2y)}f''\left( \frac{\cos(2x)}{\ch(2y)}\right).
\end{align*}

De même,

\begin{center}
$ \frac{\partial g}{\partial y}(x,y)=-2 \frac{\cos(2x)\sh(2y)}{\ch^2(2y)}f'\left( \frac{\cos(2x)}{\ch(2y)}\right)$
\end{center}

puis

\begin{align*}\ensuremath
 \frac{\partial^2g}{\partial y^2}(x,y)&=-2\cos(2x) \frac{2\ch^3(2y)-4\sh^2(2y)\ch(2y)}{\ch^4(2y)}f'\left( \frac{\cos(2x)}{\ch(2y)}\right)+4 \frac{\cos^2(2x)\sh^2(2y)}{\ch^4(2y)}f''\left( \frac{\cos(2x)}{\ch(2y)}\right)\\
 &=-4 \frac{\cos(2x)}{\ch^3(2y)}(-\ch^2(2y)+2)f'\left( \frac{\cos(2x)}{\ch(2y)}\right)+4 \frac{\cos^2(2x)(\ch^2(2y)-1)}{\ch^4(2y)}f''\left( \frac{\cos(2x)}{\ch(2y)}\right).
\end{align*}

Donc, pour tout $(x,y)\in\Rr^2$,

\begin{align*}\ensuremath
 \frac{\ch^2(2y)}{4}\Delta g(x,y)&=-2 \frac{\cos(2x)}{\ch(2y)}f'\left( \frac{\cos(2x)}{\ch(2y)}\right)+\left(1- \frac{\cos^2(2x)}{\ch^2(2y)}\right)f''\left( \frac{\cos(2x)}{\ch(2y)}\right).
\end{align*}

Maintenant, pour $(x,y)\in\Rr^2$, $-1\leqslant  \frac{\cos(2x)}{\ch(2y)}\leqslant1$ et d'autre part, l'expression  $ \frac{\cos(2x)}{\ch(2\times)}=\cos(2x)$ décrit $[-1,1]$ quand $x$ décrit $\Rr$. Donc $\left\{ \frac{\cos(2x)}{\ch(2\times)},\;(x,y)\in\Rr^2\right\}=[-1,1]$. Par suite,

\begin{center}
$\forall(x,y)\;\Rr^2,\;\Delta g(x,y)=0\Leftrightarrow\forall t\in[-1,1],\;(1-t^2)f''(t)-2tf'(t)=0$.
\end{center}

On cherche une application $f$ de classe $C^2$ sur $]-1,1[$. Or $\left| \frac{\cos(2x)}{\ch(2y)}\right|=1\Leftrightarrow|\cos(2x)|=\ch(2y)\Leftrightarrow|\cos(2x)|=\ch(2y)=1\Leftrightarrow y=0\;\text{et}\;x\in \frac{\pi}{2}\Zz$. Donc

\begin{align*}\ensuremath
\forall(x,y)\;\Rr^2\setminus\left\{\left( \frac{k\pi}{2},0\right),\;k\in\Zz\right\},\;\Delta g(x,y)=0&\Leftrightarrow\forall t\in]-1,1[,\;(1-t^2)f''(t)-2tf'(t)=0\\
 &\Leftrightarrow\forall t\in]-1,1[,\;((1-t^2)f')'(t)=0\Leftrightarrow\exists\lambda\in\Rr/\;\forall t\in]-1,1[,\;f'(t)= \frac{\lambda}{1-t^2}\\
 &\Leftrightarrow\exists(\lambda,\mu)\in\Rr^2/\;\forall t\in]-1,1[,\;f(t)=\lambda\Argth t+\mu.
\end{align*}

De plus, $f$ n'est pas constante si et seulement si $\mu=0$.

\begin{center}
\shadowbox{
L'application $t\mapsto\Argth t$ convient.
}
\end{center}
\fincorrection
\correction{005905}
Soit $(x,y)\in\Rr^2$. La matrice jacobienne de $f$ en $(x,y)$ s'écrit $\left(
\begin{array}{cc}
c(x,y)&-s(x,y)\\
s(x,y)&c(x,y)
\end{array}
\right)$ où $c$ et $s$ sont deux fonctions de classe $C^1$ sur $\Rr^2$ telle que $c^2+s^2=1$ $(*)$. Il s'agit dans un premier temps de vérifier que les fonctions $c$ et $s$ sont constantes sur $\Rr^2$.

Puisque $f$ est de classe $C^2$ sur $\Rr^2$, d'après le théorème de \textsc{Schwarz}, $ \frac{\partial^2f}{\partial x\partial y}= \frac{\partial^2f}{\partial x\partial y}$. Ceci s'écrit encore $ \frac{\partial}{\partial y}\left(
\begin{array}{c}
c\\
s
\end{array}
\right)= \frac{\partial}{\partial x}\left(
\begin{array}{c}
-s\\
c
\end{array}
\right)$ ou enfin

\begin{center}
$\forall(x,y)\in\Rr^2$, $\left(
\begin{array}{c}
 \frac{\partial c}{\partial y}(x,y)\\
\rule{0mm}{6mm} \frac{\partial s}{\partial y}(x,y)
\end{array}
\right)=\left(
\begin{array}{c}
- \frac{\partial s}{\partial x}(x,y)\\
\rule{0mm}{6mm} \frac{\partial c}{\partial x}(x,y)
\end{array}
\right)$ $(**)$.
\end{center}

En dérivant $(*)$ par rapport à $x$ ou à $y$, on obtient les égalités $c \frac{\partial c}{\partial x}+s \frac{\partial s}{\partial x}=0$ et $c \frac{\partial c}{\partial y}+s \frac{\partial s}{\partial y}=0$. Ceci montre que les deux vecteurs $\left(
\begin{array}{c}
 \frac{\partial c}{\partial x}\\
\rule{0mm}{6mm} \frac{\partial s}{\partial x}
\end{array}
\right)$ et $\left(
\begin{array}{c}
 \frac{\partial c}{\partial y}\\
\rule{0mm}{6mm} \frac{\partial s}{\partial y}
\end{array}
\right)$ sont orthogonaux au vecteur non nul $\left(
\begin{array}{c}
c\\
s
\end{array}
\right)$ et sont donc colinéaires. Mais l'égalité $(**)$ montre que les deux vecteurs $\left(
\begin{array}{c}
 \frac{\partial c}{\partial x}\\
\rule{0mm}{6mm} \frac{\partial s}{\partial x}
\end{array}
\right)$ et $\left(
\begin{array}{c}
 \frac{\partial c}{\partial y}\\
\rule{0mm}{6mm} \frac{\partial s}{\partial y}
\end{array}
\right)$ sont aussi orthogonaux l'un à l'autre. Finalement, pour tout $(x,y)\in\Rr^2$, les deux vecteurs $\left(
\begin{array}{c}
 \frac{\partial c}{\partial x}(x,y)\\
\rule{0mm}{6mm} \frac{\partial s}{\partial x}(x,y)
\end{array}
\right)$ et $\left(
\begin{array}{c}
 \frac{\partial c}{\partial y}(x,y)\\
\rule{0mm}{6mm} \frac{\partial s}{\partial y}(x,y)
\end{array}
\right)$ sont nuls. On en déduit que les deux applications $c$ et $s$ sont constantes sur $\Rr^2$ et donc, il existe $\theta$ dans $\Rr$ tel que pour tout $(x,y)\in\Rr^2$, la matrice jacobienne de $f$ en $(x,y)$ est $\left(
\begin{array}{cc}
\cos(\theta)&-\sin(\theta)\\
\sin(\theta)&\cos(\theta)
\end{array}
\right)$.

Soit $g$ la rotation d'angle $\theta$ prenant la même valeur que $f$ en $(0,0)$. $f$ et $g$ ont mêmes différentielles en tout point et coïncident en un point. Donc $f=g$ et $f$ est une rotation affine.

\begin{center}
\shadowbox{
\begin{tabular}{c}
Soit $f~:~\Rr^2\rightarrow\Rr^2$ de classe $C^2$ dont la différentielle en tout point est une rotation.\\
Alors $f$ est une rotation affine.
\end{tabular}
}
\end{center}
\fincorrection
\correction{002641}
\begin{enumerate}  
\item $df=(2x-y)\mathrm dx + (2y-x)\mathrm dy$ et
$\mathrm{Hess}_f=\left[\begin{matrix} 2 & -1\\ -1 & 2 \end{matrix}\right]$ 
d'o\`u
\begin{align*}
(u,v)\mathrm{Hess}_f(0,0)\left[\begin{matrix} u \\ v \end{matrix}\right]
&=
(u,v)
\left[\begin{matrix} 2 & -1\\ -1 & 2 \end{matrix} \right]
\left[\begin{matrix} u \\ v \end{matrix}\right]
\\
&=
u(2u-v)+v(2v-u)=2(u^2-uv+v^2)
\\
&=
2\left(\left(u-\tfrac v2\right)^2 + \tfrac 34 v^2\right) .
\end{align*}
Par cons\'equent la forme hessienne au point $(0,0)$
est positive et ce point pr\'esente donc un minimum local.

\item $f(x,y)=x^2+2xy+y^2+6= (x+y)^2+6$ d'o\`u 
le point $(0,0)$
pr\'esente un minimum local.

\item 
$df=(3x^2+2x+2y^2 +3y)\mathrm dx + (4xy-4y^3 +3x+2y)\mathrm dy$ et
\[
\mathrm{Hess}_f=\left[\begin{matrix} 
6x +2 & 4y+3\\ 4y+3 &  -12y^2 +4x +2
\end{matrix}\right]
\] 
d'o\`u
\begin{align*}
(u,v)\mathrm{Hess}_f(0,0)
\left[\begin{matrix} u \\ v \end{matrix}\right]&=
(u,v)\left[\begin{matrix} 
2 & 3\\ 3 &  2
\end{matrix}\right]
\left[\begin{matrix} u \\ v \end{matrix}\right]
\\
&=
(2u+3v)u+(3u+2v)v=2(u^2+3uv+v^2)
\\
&=
2\left(\left(u+\tfrac {3v}2\right)^2 - \tfrac 54 v^2\right) .
\end{align*}
Par cons\'equent la forme hessienne au point $(0,0)$
est non d\'eg\'en\'er\'ee et 
ind\'efinie et ce point pr\'esente  un point selle.
\end{enumerate}
\fincorrection
\correction{002642} 
Puisque
$df=\cos x \,\mathrm dx + (2y-2)\mathrm dy$, les points critiques sont les points
$((k+1/2)\pi,1)$ ($k\in \mathbb Z)$.
En plus,
$\mathrm{Hess}_f=\left[\begin{matrix} - \sin x & 0\\ 0 & 2 \end{matrix}\right]$ 
et 
\[
-\sin ((k+1/2)\pi)=(-1)^{k+1}
\]
d'o\`u 
$\mathrm{Hess}_f((k+1/2)\pi,1)
=\left[\begin{matrix} (-1)^{k+1} & 0\\ 0 & 2 \end{matrix}\right]$.
Par cons\'equent,  si
$k$ est impaire,
le point $((k+1/2)\pi,1)$ pr\'esente un minimum local et,
 si $k$ est paire, le point $((k+1/2)\pi,1)$ pr\'esente un point selle. 
\fincorrection
\correction{002643}
\begin{enumerate}  
\item $dF=f(y) f'(x)\mathrm dx + f(x) f'(y)\mathrm dy$ et
\begin{equation}
\mathrm{Hess}_f(x,y)=\left[\begin{matrix} 
f(y) f''(x) &  f'(x) f'(y)\\  f'(x) f'(y) & f(x)  f''(y)
\end{matrix}\right]
\label{Hess}
\end{equation} 
d'o\`u
$\mathrm{Hess}_f(0,0)= (f'(0))^2\left[\begin{matrix} 
0 &  1\\  1 & 0
\end{matrix}\right]$
et
\[
(u,v)\mathrm{Hess}_f(0,0)
\left[\begin{matrix} u \\ v \end{matrix}\right]
=(f'(0))^2
(u,v)\left[\begin{matrix} 
0 & 1\\ 1 &  0
\end{matrix}\right]
\left[\begin{matrix} u \\ v \end{matrix}\right]
=2(f'(0))^2uv
\]
Par cons\'equent la forme hessienne au point $(0,0)$
est non d\'eg\'en\'er\'ee et 
ind\'efinie et ce point ne peut pas pr\'esenter 
un extremum relatif.
En effet, 
le point $(0,0)$ est critique mais  un point selle.

\item D'apr\`es la partie (1.) et la p\'eriodicit\'e, les points de la forme
\begin{equation}
(x,y)=(k,l) \in \R^2,\ k,l \in \mathbb Z,
\label{selle}
\end{equation} 
pr\'esentent des points selle.
\'Egalement d'apr\`es la partie (1.),
\begin{align*}
\frac{\partial f}{\partial x}&=f(y) f'(x)=2 \pi \sin (2\pi y) \cos (2\pi x)
\\
\frac{\partial f}{\partial y}&=f(x) f'(y)=2 \pi \sin (2 \pi x) \cos (2 \pi y).
\end{align*}
Par cons\'equent,
pour que le point $(x,y)$  soit critique
il faut et il suffit qu'il soit de la forme
\[
(k,l),\ (k+\tfrac 12,l),\ (k,l+\tfrac 12),\ (k+\tfrac 12,l+\tfrac 12),\  
k,l \in \mathbb Z,
\]
ou
\[
 (k+\tfrac 14,l+\tfrac 14),
\
 (k+\tfrac 14,l+\tfrac 34),
\
 (k+\tfrac 34,l+\tfrac 14),
\
 (k+\tfrac 34,l+\tfrac 34), \ k,l \in \mathbb Z.
\]
D'apr\`es la p\'eriodicit\'e, il suffit d'examiner les huit points
\[
(0,0),\ (\tfrac 12,0),\ (0,\tfrac 12),\ (\tfrac 12,\tfrac 12),
\
 (\tfrac 14,\tfrac 14),
\
 (\tfrac 14,\tfrac 34),
\
 (\tfrac 34,\tfrac 14),
\
 (\tfrac 34,\tfrac 34)
\]
et, d'apr\`es (1.), l'origine pr\'esente un point selle.
D'apr\`es \eqref{Hess},
\begin{align*}
\mathrm{Hess}_f(0,\tfrac 12)&=\left[\begin{matrix} 
f(\tfrac 12) f''(0) &  f'(0) f'(\tfrac 12)\\  f'(0) f'(\tfrac 12) & f(0)  f''(\tfrac 12)
\end{matrix}\right]=16 \pi^2\left[\begin{matrix} 
0 &  -1\\  -1 & 0
\end{matrix}\right]
\\
\mathrm{Hess}_f(\tfrac 12,0)&=\left[\begin{matrix} 
f(0) f''(\tfrac 12) &  f'(\tfrac 12) f'(0)\\  f'(\tfrac 12) f'(0) & f(\tfrac 12)  f''(0)
\end{matrix}\right]=16 \pi^2\left[\begin{matrix} 
0 &  -1\\  -1 & 0
\end{matrix}\right]
\\
\mathrm{Hess}_f(\tfrac 12,\tfrac 12)&=\left[\begin{matrix} 
f(\tfrac 12) f''(\tfrac 12) &  f'(\tfrac 12) f'(\tfrac 12)\\  f'(\tfrac 12) f'(\tfrac 12) & f(\tfrac 12)  f''(\tfrac 12)
\end{matrix}\right]=16 \pi^2\left[\begin{matrix} 
0 &  1\\  1 & 0
\end{matrix}\right]
\end{align*}
d'o\`u les points $(0,\frac 12)$, $(\tfrac 12,0)$ et $(\tfrac 12,\tfrac 12)$
pr\'esentent des points selle.
Il est g\'eo\-m\'etri\-quement 
\'evident que le comportement de la fonction sin entra\^\i ne
que les points $(\tfrac 14,\tfrac 14)$ et  $(\tfrac 34,\tfrac 34)$
pr\'esentent des maxima et que les points
$(\tfrac 14,\tfrac 34)$ et  $(\tfrac 34,\tfrac 14)$
pr\'esentent des minima.
\end{enumerate}
\fincorrection
\correction{002644} 
Soit $f\colon \R^3 \longrightarrow \R$ la fonction d\'efinie par
\[f(x,y,z)=\sin(\pi xy)+\sin(\pi yz) -1.\]
Ses d\'eriv\'ees partielles sont
\[\frac{\partial f}{\partial x} =\pi y \cos(\pi xy),\ 
\frac{\partial f}{\partial y} =\pi (x \cos(\pi xy) + z \cos(\pi yz)),
\ \frac{\partial f}{\partial y} =\pi y \cos(\pi yz)\]
et, apr\`es simplification, au point  $(1,\frac{1}{6}, 1)$, 
l'\'equation \eqref{tang4}  du plan tangent \`a la surface de niveau
en discussion devient
\[(x-1)+12(y-1/6)+(z-1)=0 .\]
Ainsi,  en ce point, le vecteur $(1,12,1)$ est perpendiculaire \`a la surface. 
\fincorrection
\correction{002645}
\begin{enumerate} 
\item 
Puisque 
$\frac{\partial f}{\partial x}= ye^{x}+2e^y\cos(2x)$ et 
$\frac{\partial f}{\partial y}=e^{x}+e^y\sin(2x)$ il s'ensuit que
$\frac{\partial f}{\partial x}(0,0)=2$ et 
$\frac{\partial f}{\partial y}(0,0)=1$.
Par cons\'equent,
il existe une fonction 
$h$ de la variable $x$ d\'efinie au voisinage de $0$ telle que
$h(0)=0$ et telle que,
pour qu'au voisinage de $(0,0)$
les coordonn\'ees $x$ et $y$ du point $(x,y)$
satisfassent \`a l'\'equation $ye^{x}+e^y\sin(2x)=0$
il faut et il suffit que
$y=h(x)$; 
de m\^eme il existe une fonction $k$ 
de la variable $y$ d\'efinie au voisinage de $0$  telle que
$h(0)=0$ et telle que,
pour qu'au voisinage de $(0,0)$
les coordonn\'ees $x$ et $y$ du point $(x,y)$
satisfassent \`a  l'\'equation $ye^{x}+e^y\sin(2x)=0$
il faut et il suffit que $x=k(y)$. 
En plus,
\[
h'(0)=-\frac{\frac{\partial f}{\partial x}(0,0)}
{\frac{\partial f}{\partial y}(0,0)}= -2,
\ 
k'(0)= -\frac{\frac{\partial f}{\partial y}(0,0)}
{\frac{\partial f}{\partial x}(0,0)} =-\tfrac 12.
\]

\item Puisque le point $(0,0)$ appartient \`a la courbe $\cal C$,
en $0$, les fonctions $h$ et $k$ prennent les valeurs
$h(0)=0$ et $k(0)=0$. Par cons\'equent,
\[
\mathrm{lim}_{\begin{smallmatrix}
(x,y)\to(0,0), (x,y)\ne 0\\
ye^{x}+e^y\sin(2x)=0
\end{smallmatrix}}y/x =h'(0) = -2.
\]
\end{enumerate}
\fincorrection
\correction{002646}
\begin{enumerate}
  \item  Puisque $\frac{\partial f}{\partial y} = -2y$ et
\[
\frac{\partial f}{\partial x} =  (x+1)^2 + 2x(x+1)=(x+1)(3x+1)=3x^2+4x+1, 
\] 
les points stationnaires de $f$ sont les points
$(-1,0)$ et $(-1/3,0)$.
En plus,
\[
\mathrm{Hess}_f(x,y)=\left[\begin{matrix} 
6x+4 & 0\\  0 & -2
\end{matrix}\right]
\] 
d'o\`u
$\mathrm{Hess}_f(-1,0)= \left[\begin{matrix} 
-2 &  0\\  0 & -2
\end{matrix}\right]$
et
$\mathrm{Hess}_f(-1/3,0)= \left[\begin{matrix} 
2 &  0\\  0 & -2
\end{matrix}\right]$.
Par cons\'equent la forme hessienne au point $(-1,0)$
est d\'efinie n\'egative et ce point pr\'esente  un maximum local;
de m\^eme, la forme hessienne au point $(-1/3,0)$
est non d\'eg\'en\'er\'ee  et ind\'efinie 
et ce point pr\'esente  un point selle.

\item  La courbe $y= \sqrt{x} \hskip 2pt (x+1)$ pour $x\geq 0$
passe par les points $(0,0)$, $(\tfrac 13,\tfrac 43 \sqrt 3)$, 
$(1,2)$, et $(2,3 \sqrt 2)$;
elle a une tangente verticale \`a l'origine,
le point $(\tfrac 13,\tfrac 43 \sqrt 3)$
est un point d'inflexion, la pente en ce point vaut $\sqrt 3$,
et c'est la pente minimale de la courbe.
Ces faits se d\'eduisent
des expressions
$y'= \tfrac 32 \sqrt x + \tfrac 12 (\sqrt x)^{-1}$ et
$y''= \tfrac 34 x^{-\tfrac 12} - \tfrac 14 x^{-\tfrac 32}$.
La courbe constitu\'ee des points tels que $f(x,y)=0$ et $x \geq 0$
s'obtient par r\'eflexion de la courbe
$y= \sqrt{x} \hskip 2pt (x+1)$ pour $x\geq 0$
par rapport \`a l'axe des $x$.

\item Dans la boule ouverte 
\[
\{(x,y,z);(x+1)^2+y^2+x^2 <1\} \subseteq \R^3, 
\]
le graphe
$z=f(x,y)$ de la fonction $f$ ne rencontre le plan des $x$ et $y$ qu'au point
$(-1,0)$. Par cons\'equent, l'intersection 
$D \cap \cal C$ du disque
\[
D=\{(x,y); (x+1)^2+y^2<1\}
\]
avec $\cal C$ ne consiste qu'au point $(-1,0)$.


\item Voir l'indication de l'exercice pr\'ec\'edent.

\item Quel que soit le point 
$(x_0,y_0)$ de ${\cal C}$
distinct de  $(-1,0)$,
d'apr\`es (1.),
\[
\left(\frac{\partial f}{\partial x}(x_0,y_0),\frac{\partial f}{\partial y}(x_0,y_0) \right) 
\ne (0,0).
\] 
L'assertion est donc une cons\'equence imm\'ediate du th\'eor\`eme des 
fonctions implicites. 
\end{enumerate}
\fincorrection
\nocorrection
\correction{004191}
\begin{enumerate}

  \item $(0,0)$ : non extrémal\par
             $(1,1)$ : maximum local
  \item $(0,0)$ : non extrémal\par
             $\pm(\sqrt2, -\sqrt2)$ : minimum absolu
  \item $(-2/15,-1/5)$ : maximum local \par
             $(x,0)$ : max. local pour $x<-1/3$, min. local pour $x > -1/3$\par
             $(0,y)$ : max. local pour $y<-1/2$, min. local pour $y > -1/2$
  \item $(2^{-1/3},4^{-1/3})$ : maximum absolu
  \item prendre le log. $(1,1)$ : maximum absolu
  \item $(-1,-1)$ : non extrémal
  \item $(1,0)$ : minimum absolu\par
             $(e^{-2},0)$ : non extrémal
  \item $= MA + MB  \Rightarrow $ minimum absolu sur $[A,B]$
  \item $M \in \text{med}(A,B),\ \overline{(\vec{MA}, \vec{MB})} = \pm \frac {2\pi}3$\par
             $M = A,B$ : minimum absolu\par
             $M = O$
\end{enumerate}
 \fincorrection
\correction{004192}
\begin{enumerate}
  \item 
  \item 
  \begin{enumerate}
    \item isobarycentre de $ABC$.
    \item Point de Fermat ou $A,B,C$.
    \item$A,B,C$.
  \end{enumerate}
\end{enumerate}
 \fincorrection
\correction{004193}
\begin{enumerate}
  \item $4S=\sqrt{2a^2b^2 + 2a^2c^2 + 2b^2c^2 - a^4 - b^4 - c^4}$.
  \item max = $\frac4{\sqrt3}$.
\end{enumerate}
 \fincorrection
\correction{004194}
Existence d'un maximum par compacité.
Soient $x,y$ les coordonnées d'un point $M$ dans un repère orthonormé du plan et
$(u,v,w)$ les coordonnées barycentriques de $M$ par rapport à $A,B,C$ (avec $u+v+w=1$).
$u,v,w$ sont des fonctions affines de $x,y,z$ et $(AB)$ a pour équation
barycentrique $w=0$ d'où $d(M,AB) = \alpha |w|$ pour un certain réel $\alpha>0$.
De même pour $d(M,AC)$ et $d(M,BC)$ et $f(M) = \alpha\beta\gamma|u||v||w|$.
Lorsque $M$ varie dans le triangle, $(u,v,w)$ décrit tous les triplets de réels positifs de
somme~$1$ et on cherche le maximum du produit $uvw$, il est atteint quand $u,v,w$
sont égaux, c'est-à-dire au centre de gravité du triangle.
\fincorrection
\nocorrection
\correction{004196}
Il existe une base orthonormale de~$E$ et un réel $\lambda$
tels que $f(x) = \lambda x_1$ et $g(x) = \lambda x_1e^{-x_1^2}e^{-x_2^2-\dots-x_n^2}$.
Donc $g$ est maximale/minimale pour $x_1=\pm1/\sqrt2$, $x_2=\dots=x_n=0$.
\fincorrection
\correction{004197}
Le minimum demandé existe car $\varphi(M,N,P)\to+\infty$ quand l'un au
moins des points $M,N,P$ tend vers l'infini sur sa droite personnelle.
Soient $D_1\cap D_2 = \{A\}$, $D_2\cap D_3 = \{B\}$, $D_3\cap D_1 = \{C\}$
et $A'$, $B'$, $C'$ les milieux de $[B,C]$, $[C,A]$ et $[A,B]$.
Déjà on a $\varphi(B',C',A') = \frac32a$ et $\varphi(A,N,P) \ge 2AP \ge a\sqrt3$
donc le minimum n'est pas atteint lorsque l'un des points $M,N,P$ est confondu
avec l'un des points $A,B,C$, ni non plus si l'un des points $M,N,P$
est hors du triangle $ABC$.
Pour $N,P$ fixés hors de $D_1$, on fait varier $M$ sur $D_1$~: $M = A + t\vec{AB}$ avec $t\in\R$
et on considère $f(t) = \varphi(M,N,P)$.
Alors $f'(t) = \Bigl(\vec{AB} \mid \frac{\vec{MN}}{MN} + \frac{\vec{MP}}{MP}\Bigr)$
donc $f(t)$ est minimal lorsque $D_1$ est la bissectrice extérieure des
demi-droites $[MN)$ et $[MP)$.

Soit $(M,N,P)$ un triplet réalisant le minimum de~$\varphi$ et $\alpha,\beta,\gamma$
les angles du triangle $MNP$ en $P,M$ et $N$. Les angles du triangle
$AMN$ sont $\pi/3$, $(\pi-\beta)/2$ et $(\pi-\gamma)/2$ d'où
$2\pi/3 = \beta+\gamma = \pi-\alpha$ et donc $\alpha = \pi/3 = \beta = \gamma$.
On en déduit que $(MP)$ est paralèle à $(AB)$, $(MN)$ à $(BC)$ et $(NP)$ à $(AC)$
puis que $(M,N,P) = (B',C',A')$.
\fincorrection
\correction{004198}
\begin{enumerate}
  \item On fixe $A_i\in D_i$ et $\vec u_i$ un vecteur directeur
    de~$D_i$. Soit $M_i = A_i + x_i\vec u_i$. Alors

    \begin{align*}f(M_1,M_2,M_3) &= f(A_1,A_2,A_3)\\
                    &+ 2\Bigl(
                        (\vec{A_1A_2}\mid x_2\vec u_2 - x_1\vec u_1)
                      + (\vec{A_2A_3}\mid x_3\vec u_3 - x_2\vec u_2)
                      + (\vec{A_3A_1}\mid x_1\vec u_1 - x_3\vec u_3)
                     \Bigr)\\
                    &+ \Bigl(
                        \|x_2\vec u_2 - x_1\vec u_1\|^2
                      + \|x_3\vec u_3 - x_2\vec u_2\|^2
                      + \|x_1\vec u_1 - x_3\vec u_3\|^2
                     \Bigr)\\
                    &= a + b(x_1,x_2,x_3) + c(x_1,x_2,x_3).
\end{align*}

    $b$ est une forme linéaire et $c$ est une forme quadratique positive,
    et même définie positive car $\vec u_1,\vec u_2,\vec u_3$ sont deux
    à deux non colinéaires. Il en résulte que
    $f(M_1,M_2,M_3) \to +\infty$ lorsque $|x_1|+|x_2|+|x_3|\to\infty$, donc par
    continuité, $f$ admet un minimum.
    
    Choisissons alors $A_1,A_2,A_3$ de sorte que $f(A_1,A_2,A_3)$ soit égal
    à ce minimum. On a alors $b=0$ car $(A_1,A_2,A_3)$ est point critique
    de~$f$, d'où $f(M_1,M_2,M_3) > f(A_1,A_2,A_3)$ si $(x_1,x_2,x_3)\ne(0,0,0)$
    vu la définie-positivité de~$c$. Ceci prouve l'unicité du triplet
    où $f$ atteint son minimum.
  \item On soup\c conne fortement le triplet constitué des milieux
    des côtés. En notant $A_1,A_2,A_3$ ces milieux, il suffit de
    vérifier que la forme linéaire $b$ de la réponse précédente
    est nulle, et c'est clairement le cas après regroupement autour
    de~$x_1$, $x_2$, $x_3$.
\end{enumerate}
 \fincorrection
\correction{004199}
On param{\`e}tre le chemin en coordonn{\'e}ees sph{\'e}riques par $t\mapsto (\theta(t),\phi(t))$.

La longueur du chemin est ${ \int_{t=0}^{1}\sqrt{\phi'^{2}(t)+\sin^{2}(\phi(t))\theta'^{2}(t)}d t\ge \Bigl| \int_{t=0}^{1}\phi'(t)d t\Bigr|}$ avec {\'e}galit{\'e} si et 
seulement si $\theta'=0$ et $\phi'$ est de signe constant. On trouve donc les m{\'e}ridiens.
\fincorrection
\correction{004200}
Soit $y\in B$ orthogonal {\`a} $x$. La fonction $g : \theta \mapsto f(x\cos \theta +y\sin \theta)$ admet un extr{\'e}mum en $0$, donc
$g'(0)=0$, soit $\nabla f(x) \perp y$. Si $x=0$ on a donc $\nabla f(0) = 0$.
Sinon, $\nabla f(x) = \lambda x$ avec $\lambda\in\R$,
et en faisant un d{\'e}veloppement limit{\'e} de $f(x-tx)$ on voit
que $\lambda \ge 0$.
\fincorrection
\correction{005558}
\begin{enumerate}
 \item  $f$ est de classe $C^1$ sur $\Rr^2$ qui est un ouvert de $\Rr^2$. Donc si $f$ admet un extremum local en un point $(x_0,y_0)$ de $\Rr^2$, $(x_0,y_0)$ est un point critique de $f$. Or, pour $(x,y)\in\Rr^2$,

\begin{center}
$\left\{
\begin{array}{l}
\frac{\partial f}{\partial x}(x,y)=0\\
\rule{0mm}{7mm}\frac{\partial f}{\partial x}(x,y)=0
\end{array}
\right.\Leftrightarrow\left\{
\begin{array}{l}
2x+y+2=0\\
x+2y+3=0
\end{array}
\right.\Leftrightarrow\left\{
\begin{array}{l}
x=-\frac{1}{3}\\
\rule{0mm}{7mm}y=-\frac{4}{3}
\end{array}
\right.$.
\end{center}
Donc si $f$ admet un extremum local, c'est nécessairement en $\left(-\frac{1}{3},\frac{4}{3}\right)$ avec $f\left(-\frac{1}{3},\frac{4}{3}\right)=-\frac{7}{3}$. D'autre part,

\begin{align*}\ensuremath
f(x,y)&=x^2+xy+y^2+2x+3y=\left(x+\frac{y}{2}+1\right)^2-\left(\frac{y}{2}+1\right)^2+y^2+3y=\left(x+\frac{y}{2}+1\right)^2+\frac{3y^2}{4}+2y-1\\
 &=\left(x+\frac{y}{2}+1\right)^2+\frac{3}{4}\left(y+\frac{4}{3}\right)^2-\frac{7}{3}\geqslant-\frac{7}{3}=f\left(-\frac{1}{3},\frac{4}{3}\right).
\end{align*}
Donc $f$ admet un minimum local en $\left(-\frac{1}{3},\frac{4}{3}\right)$ égal à $-\frac{7}{3}$ et ce minimum local est un minimum global. D'autre part, $f$ n'admet pas de maximum local.
 \item   $f$ est de classe $C^1$ sur $\Rr^2$ qui est un ouvert de $\Rr^2$. Donc si $f$ admet un extremum local en un point $(x_0,y_0)$ de $\Rr^2$, $(x_0,y_0)$ est un point critique de $f$. Or, pour $(x,y)\in\Rr^2$,

\begin{center}
$\left\{
\begin{array}{l}
\frac{\partial f}{\partial x}(x,y)=0\\
\rule{0mm}{7mm}\frac{\partial f}{\partial x}(x,y)=0
\end{array}
\right.\Leftrightarrow\left\{
\begin{array}{l}
4x^3-4y=0\\
4y^3-4x=0
\end{array}
\right.\Leftrightarrow\left\{
\begin{array}{l}
y=x^3\\
x^9-x=0
\end{array}
\right.\Leftrightarrow(x,y)\in\{(0,0),(1,1),(-1,-1)$.
\end{center}
Les points critiques de $f$ sont $(0,0)$, $(1,1)$ et $(-1,-1)$.
Maintenant, pour $(x,y)\in\Rr^2$, $f(-x,-y)=f(x,y)$. Ceci permet de restreindre l'étude aux deux points $(0,0)$  et $(1,1)$.
\textbullet~Pour $x\in\Rr$, $f(x,0)=x^4>0$ sur $\Rr^*$ et $f(x,x)=-4x^2+2x^4=2x^2(-2+x^2)<0$ sur $]-\sqrt{2},0[\cup]0,\sqrt{2}[$. Donc $f$ change de signe dans tous voisinage de $(0,0)$ et puisque $f(0,0)=0$, $f$ n'admet pas d'extremum local en $(0,0)$.
\textbullet~Pour $(h,k)\in\Rr^2$,

\begin{align*}\ensuremath
f(1+h,1+k)-f(1,1)&=(1+h)^4+(1+k)^4-4(1+h)(1+k)+2=6h^2+6k^2-4hk+4h^3+4k^3+h^4+k^4\\
 &\geqslant6h^2+6k^2-2(h^2+k^2)+4h^3+4k^3+h^4+k^4=4h^2+4h^3+h^4+4k^2+4k^3+k^4\\
 &=h^2(2h^2+1)^2+k^2(2k^2+1)^2\geqslant0.
\end{align*}
$f$ admet donc un minimum global en $(1,1)$ (et en $(-1,-1)$) égal à $-2$.
\end{enumerate}
\fincorrection
\correction{005559}
Soit $M$ un point intérieur au triangle $ABC$. On pose $a=BC$, $b=CA$ et $c=AB$. On note $x$, $y$, $z$ et $\mathcal{A}$ les aires respectives des triangles $MBC$, $MCA$, $MAB$ et $ABC$.
On a

\begin{center}
$d(M,(BC))d(M,(CA))d(M(AB))=\frac{2\text{aire}(MBC)}{a}\frac{2\text{aire}(MCA)}{b}\frac{2\text{aire}(MAB)}{c}=\frac{8xyz}{abc}=\frac{8}{abc}xy(\mathcal{A}-x-y)$.
\end{center}
On doit donc déterminer le maximum de la fonction $f(x,y)=xy(\mathcal{A}-x-y)$ quand $(x,y)$ décrit le triangle ouvert $T=\{(x,y)\in\Rr^2,\;x>0,\;y>0,\;x+y<\mathcal{A}\}$. On admet que $f$ admet un maximum global sur le triangle fermé $T'=\{(x,y)\in\Rr^2,\;x\geqslant0,\;y\geqslant0,\;x+y\leqslant\mathcal{A}\}$ (cela résulte d'un théorème de math Spé : \og une fonction numérique continue sur un compact admet un minimum et un maximum \fg). Ce maximum est atteint dans l'intérieur $T$ de $T'$ car $f$ est nulle au bord de $T'$ et strictement positive à l'intérieur de $T'$.

Puisque $f$ est de classe $C^1$ sur $T$ qui est un ouvert de $\Rr^2$, $f$ atteint son maximum sur $T$ en un point critique de $f$. Or, pour $(x,y)\in T^2$,

\begin{align*}\ensuremath
\left\{
\begin{array}{l}
\frac{\partial f}{\partial x}(x,y)=0\\
\rule{0mm}{6mm}\frac{\partial f}{\partial y}(x,y)=0
\end{array}
\right.&\Leftrightarrow\left\{
\begin{array}{l}
y(\mathcal{A}-x-y)-xy=0\\
y(\mathcal{A}-x-y)-xy=0
\end{array}
\right.\Leftrightarrow\left\{
\begin{array}{l}
y(\mathcal{A}-2x-y)=0\\
x(\mathcal{A}-x-2y)=0
\end{array}
\right.\\
 &\Leftrightarrow\left\{
\begin{array}{l}
2x+y=\mathcal{A}\\
x+2y=\mathcal{A}
\end{array}
\right.\Leftrightarrow x=y=\frac{\mathcal{A}}{3}.
\end{align*}
Le maximum cherché est donc égal à $\frac{8}{abc}\times\frac{\mathcal{A}}{3}\times\frac{\mathcal{A}}{3}\times\left(\mathcal{A}-\frac{\mathcal{A}}{3}-\frac{\mathcal{A}}{3}\right)=\frac{8\mathcal{A}^3}{27abc}$.
(On peut montrer que ce maximum est obtenu quand $M$ est le centre de gravité du triangle $ABC$).
\fincorrection
\correction{005560}
Soient $\mathcal{R}$ un repère orthonormé de $\Rr^2$ muni de sa structure euclidienne canonique puis $M$, $A$ et $B$ les points de coordonnées respectives $(x,y)$, $(0,a)$ et $(a,0)$ dans $\mathcal{R}$.
Pour $(x,y)\in\Rr^2$, $f(x,y)=MA+MB\geqslant AB=a\sqrt{2}$ avec égalité si et seulement si $M\in[AB]$. Donc

\begin{center}
\shadowbox{
Le minimum de $f$ sur $\Rr^2$ existe et vaut $a\sqrt{2}$.
}
\end{center}
\fincorrection
\correction{005902}
On pose $BC=a$, $CA=b$ et $AB=c$ et on note $\mathcal{A}$ l'aire du triangle $ABC$. Soit $M$ un point intérieur au triangle $ABC$. On note $I$, $J$ et $K$ les projetés orthogonaux de $M$ sur les droites $(BC)$, $(CA)$ et $(AB)$ respectivement. On pose $u=\text{aire de}\;MBC$, $v=\text{aire de}\;MCA$ et $w=\text{aire de}\;MAB$. On a

\begin{center}
$d(M,(BC))\times d(M,(CA))\times d(M,(AB))=MI\times MJ\times MK= \frac{2u}{a}\times \frac{2v}{b}\times \frac{2w}{c}= \frac{8}{abc}uv(\mathcal{A}-u-v)$.
\end{center}

Il s'agit alors de trouver le maximum de la fonction $f~:~(u,v)\mapsto uv(\mathcal{A}-u-v)$ sur le domaine 

\begin{center}
$T=\left\{(u,v)\in\Rr^2/\;u\geqslant0,\;v\geqslant0\;\text{et}\;u+v\leqslant\mathcal{A}\right\}$.
\end{center}

$T$ est un compact de $\Rr^2$. En effet :

- $\forall(u,v)\in T^2$, $\|(u,v)\|_1=u+v\leqslant\mathcal{A}$ et donc $T$ est bornée.

- Les applications $\varphi_1~:~(u,v)\mapsto u$, $\varphi_2~:~(u,v)\mapsto v$ et $\varphi_3~:~(u,v)\mapsto u+v$ sont continues sur $\Rr^2$ en tant que formes 

linéaires sur un espace de dimension finie. Donc les ensembles $P_1=\{(u,v)\in\Rr^2/\;u\geqslant0\}=\varphi_1^{-1}([0,+\infty[)$,

$P_2=\{(u,v)\in\Rr^2/\;v\geqslant0\}=\varphi_2^{-1}([0,+\infty[)$ et $P_3=\{(u,v)\in\Rr^2/\;u+v\leqslant0\}=\varphi_3^{-1}(]-\infty,0])$ sont des fermés de $\Rr^2$

en tant qu'images réciproques de fermés par des applications continues. On en déduit que $T=P_1\cap P_2\cap P_3$ est un

fermé de $\Rr^2$ en tant qu'intersection de fermés de $\Rr^2$.

Puisque $T$ est un fermé borné de $\Rr^2$, $T$ est un compact de $\Rr^2$ puisque $\Rr^2$ est de dimension finie et d'après le théorème de \textsc{Borel}-\textsc{Lebesgue}.

$f$ est continue sur le compact $T$ à valeurs dans $\Rr$ en tant que polynôme à plusieurs variables et donc $f$ admet un maximum sur $T$.

Pour tout $(u,v)$ appartenant à la frontière de $T$, on a $f(u,v)=0$. Comme $f$ est strictement positive sur $\overset{\circ}{T}=\{(u,v)\in\Rr^2/\;u>0,\;v>0\;\text{et}\;u+v<0\}$, $f$ admet son maximum dans $\overset{\circ}{T}$. Puisque $f$ est de classe $C^1$ sur $\overset{\circ}{T}$ qui est un ouvert de $\Rr^2$, si $f$ admet un maximum en $(u_0,v_0)\in\overset{\circ}{T}$, $(u_0,v_0)$ est nécessairement un point critique de $f$. Soit $(u,v)\in\overset{\circ}{T}$.

\begin{center}
$\left\{
\begin{array}{l}
 \frac{\partial f}{\partial u}(u,v)=0\\
\rule{0mm}{6mm} \frac{\partial f}{\partial v}(u,v)=0
\end{array}
\right.\Leftrightarrow\left\{
\begin{array}{l}
v(\mathcal{A}-2u-v)=0\\
u(\mathcal{A}-u-2v)=0
\end{array}
\right.\Leftrightarrow\left\{
\begin{array}{l}
2u+v=\mathcal{A}\\
u+2v=\mathcal{A}
\end{array}
\right.\Leftrightarrow u=v= \frac{\mathcal{A}}{3}$.
\end{center}

Puisque $f$ admet un point critique et un seul à savoir $(u_0,v_0)=\left( \frac{\mathcal{A}}{3}, \frac{\mathcal{A}}{3}\right)$, $f$ admet son maximum en ce point et ce maximum vaut $f(u_0,v_0)= \frac{\mathcal{A}^3}{27}$. Le maximum du produit des distances d'un point $M$ intérieur au triangle $ABC$ aux cotés de ce triangle est donc $ \frac{8\mathcal{A}^3}{27abc}$.

\textbf{Remarque.} On peut démontrer que pour tout point $M$ intérieur au triangle $ABC$, on a $M=\text{bar}\left((A,\text{aire de}\;MBC),(B,\text{aire de}\;MAC),(C,\text{aire de}\;MAB)\right)$. Si maintenant $M$ est le point en lequel on réalise le maximum, les trois aires sont égales et donc le maximum est atteint en $G$ l'isobarycentre du triangle $ABC$.

\fincorrection
\correction{005903}
Soient $A$ et $B$ les points du plan de coordonnées respectives $(0,a)$ et $(a,0)$ dans un certain repère $\mathcal{R}$ orthonormé. Soit $M$ un point du plan de coordonnées $(x,y)$ dans $\mathcal{R}$. Pour $(x,y)\in\Rr^2$,

\begin{center}
$f(x,y)=\left\|\overrightarrow{MA}\right\|_2+\left\|\overrightarrow{MB}\right\|_2=MA+MB\geqslant AB$ avec égalité si et seulement si $M\in[AB]$.
\end{center}

Donc $f$ admet un minimum global égal à $AB=a\sqrt{2}$ atteint en tout couple $(x,y)$ de la forme $(\lambda a,(1-\lambda)a)$, $\lambda\in[0,1]$.
\fincorrection
\nocorrection
\correction{005890}
Soit $(x,y,z,t)\in\Rr^4$.

\begin{align*}\ensuremath
\varphi(x,y)=(z,t)&\Leftrightarrow\left\{
\begin{array}{l}
e^x-e^y=z\\
x+y=t
\end{array}
\right.\Leftrightarrow\left\{
\begin{array}{l}
y=t-x\\
e^x-e^{t-x}=z
\end{array}
\right.\Leftrightarrow\left\{
\begin{array}{l}
y=t-x\\
(e^x)^2-ze^x-e^t=0
\end{array}
\right.\\
 &\Leftrightarrow\left\{
\begin{array}{l}
y=t-x\\
e^x=z-\sqrt{z^2+4e^t}\;\text{ou}\;e^x=z+\sqrt{z^2+4e^t}
\end{array}
\right.\\
&\Leftrightarrow\left\{
\begin{array}{l}
e^x=z+\sqrt{z^2+4e^t}\\
y=t-x
\end{array}
\right.\;(\text{car}\;z-\sqrt{z^2+4e^t}<z-\sqrt{z^2}=z-|z|\leqslant0)\\
&\Leftrightarrow\left\{
\begin{array}{l}
x=\ln(z+\sqrt{z^2+4e^t})\\
y=t-\ln(z+\sqrt{z^2+4e^t})
\end{array}
\right.\;(\text{car}\;z+\sqrt{z^2+4e^t}>z+\sqrt{z^2}=z+|z|\geqslant0).
\end{align*}

Ainsi, tout élément $(z,t)\in\Rr^2$ a un antécédent et un seul dans $\Rr^2$ par $\varphi$ et donc $\varphi$ est une bijection de $\Rr^2$ sur lui-même.

La fonction $\varphi$ est de classe $C^1$ sur $\Rr^2$ de jacobien $J_\varphi(x,y)=\left|
\begin{array}{cc}
e^x&-e^y\\
1&1
\end{array}
\right|=e^x+e^y$. Le jacobien de $\varphi$ ne s'annule pas sur $\Rr^2$. En résumé, $\varphi$ est une bijection de $\Rr^2$ sur lui-même, de classe $C^1$ sur $\Rr^2$ et le jacobien de $\varphi$ ne s'annule pas sur $\Rr^2$. On sait alors que

\begin{center}
\shadowbox{
$\varphi$ est un $C^1$-difféomorphisme de $\Rr^2$ sur lui-même.
}
\end{center}
\fincorrection
\correction{005891}
Soit $n\in\Nn$. Soit $x\in\Rr$. La fonction $f_x:y\mapsto y^{2n+1}+y-x$ est continue et strictement croissante sur $\Rr$ en tant que somme de fonctions continues et strictement croissantes sur $\Rr$. Donc la fonction $f_x$ réalise une bijection de $\Rr$ sur $]\lim_{y \rightarrow -\infty}f_x(y),\lim_{y \rightarrow +\infty}f_x(y)[=\Rr$. En particulier, l'équation $f_x(y)=0$ a une et une seule solution dans $\Rr$ que l'on note $\varphi(x)$.

La fonction $f:(x,y)\mapsto y^{2n+1}+y-x$ est de classe $C^1$ sur $\Rr^2$ qui est un ouvert de $\Rr^2$ et de plus, $\forall(x,y)\in\Rr^2$, $ \frac{\partial f}{\partial y}(x,y)=(2n+1)y^{2n}+1\neq0$. D'après le théorème des fonctions implicites, la fonction $\varphi$ implicitement définie par l'égalité $f(x,y)=0$ est dérivable en tout réel $x$ et de plus, en dérivant l'égalité $\forall x\in\Rr$, $(\varphi(x))^{2n+1}+\varphi(x)-x=0$, on obtient $\forall x\in\Rr$, $(2n+1)\varphi'(x)(\varphi(x))^{2n}+\varphi'(x)-1=0$ et donc

\begin{center}
$\forall x\in\Rr$, $\varphi'(x)= \frac{1}{(2n+1)(\varphi(x))^{2n}+1}$.
\end{center}

Montrons par récurrence que $\forall p\in\Nn^*$, la fonction $\varphi$ est $p$ fois dérivable sur $\Rr$.

- C'est vrai pour $p=1$.

- Soit $p\geqslant1$. Supposons que la fonction $\varphi$ soit $p$ fois dérivable sur $\Rr$. Alors la fonction $\varphi'= \frac{1}{(2n+1)\varphi^{2n}+1}$ est 

$p$ fois dérivable sur $\Rr$ en tant qu'inverse d'une fonction $p$ fois dérivable sur $\Rr$ ne s'annulant pas sur $\Rr$. On en déduit

que la fonction $\varphi$ est $p+1$ fois dérivable sur $\Rr$.

On a montré par récurrence que $\forall p\in\Nn^*$, la fonction $\varphi$ est $p$ fois dérivable sur $\Rr$ et donc que

\begin{center}
\shadowbox{
la fonction $\varphi$ est de classe $C^\infty$ sur $\Rr$.
}
\end{center}

Calculons maintenant $I=\int_{0}^{2}\varphi(t)\;dt$. On note tout d'abord que, puisque $0^{2n+1}+0-0=0$, on a $\varphi(0)=0$ et puisque $1^{2n+1}+1-2=0$, on a $\varphi(2)=1$. 

Maintenant, pour tout réel $x$ de $[0,2]$, on a $\varphi'(x)(\varphi(x))^{2n+1}+\varphi'(x)\varphi(x)-x\varphi'(x)=0$ (en multipliant par $\varphi'(x)$ les deux membres de l'égalité définissant $\varphi(x)$) et en intégrant sur le segment $[0,2]$, on obtient

\begin{center}
$\int_{0}^{2}\varphi'(x)(\varphi(x))^{2n+1}\;dx+\int_{0}^{2}\varphi'(x)\varphi(x)\;dx-\int_{0}^{2}x\varphi'(x)\;dx=0$ $(*)$.
\end{center}

Or, $\int_{0}^{2}\varphi'(x)(\varphi(x))^{2n+1}\;dx=\left[ \frac{(\varphi(x))^{2n+2}}{2n+2}\right]_0^2= \frac{1}{2n+2}$. De même, $\int_{0}^{2}\varphi'(x)\varphi(x)\;dx=\left[ \frac{(\varphi(x))^{2}}{2}\right]_0^2= \frac{1}{2}$ et donc $\int_{0}^{2}\varphi'(x)(\varphi(x))^{2n+1}\;dx+\int_{0}^{2}\varphi'(x)\varphi(x)\;dx= \frac{1}{2n+2}+ \frac{1}{2}= \frac{n+2}{2n+2}$. D'autre part, puisque les deux fonctions $x\mapsto x$ et $x\mapsto\varphi(x)$ sont de classe $C^1$ sur le segment $[0,2]$, on peut effectuer une intégration par parties qui fournit

\begin{center}
$-\int_{0}^{2}x\varphi'(x)\;dx=\left[-x\varphi(x)\right]_0^2+\int_{0}^{2}\varphi(x)\;dx=-2+I$.
\end{center}

L'égalité $(*)$ s'écrit donc $ \frac{n+2}{2n+2}-2+I=0$ et on obtient $I= \frac{3n+2}{2n+2}$.

\begin{center}
\shadowbox{
$\int_{0}^{2}\varphi(x)\;dx= \frac{3n+2}{2n+2}$.
}
\end{center}
\fincorrection
\correction{005892}
Soit $x\in\Rr$. La fonction $f_x:y\mapsto e^{x+y}+y-1$ est continue et strictement croissante sur $\Rr$ en tant que somme de fonctions continues et strictement croissantes sur $\Rr$. Donc la fonction $f_x$ réalise une bijection de $\Rr$ sur $]\lim_{y \rightarrow -\infty}f_x(y),\lim_{y \rightarrow +\infty}f_x(y)[=\Rr$. En particulier, l'équation $f_x(y)=0$ a une et une seule solution dans $\Rr$ que l'on note $\varphi(x)$.

La fonction $f:(x,y)\mapsto e^{x+y}+y-1$ est de classe $C^1$ sur $\Rr^2$ qui est un ouvert de $\Rr^2$ et de plus, $\forall(x,y)\in\Rr^2$, $ \frac{\partial f}{\partial y}(x,y)=e^{x+y}+1\neq0$. D'après le théorème des fonctions implicites, la fonction $\varphi$ implicitement définie par l'égalité $f(x,y)=0$ est dérivable en tout réel $x$ et de plus, en dérivant l'égalité $\forall x\in\Rr$, $e^{x+\varphi(x)}+\varphi(x)-1=0$, on obtient $\forall x\in\Rr$, $(1+\varphi'(x))e^{x+\varphi(x)}+\varphi'(x)=0$ ou encore

\begin{center}
$\forall x\in\Rr$, $\varphi'(x)=- \frac{e^{x+\varphi(x)}}{e^{x+\varphi(x)}+1}$ $(*)$.
\end{center}

On en déduit par récurrence que $\varphi$ est de classe $C^\infty$ sur $\Rr$ et en particulier admet en $0$ un développement limité d'ordre $3$. Déterminons ce développement limité.

\textbf{1ère solution.} Puisque $e^{0+0}+0-1=0$, on a $\varphi(0)=0$. L'égalité $(*)$ fournit alors $\varphi'(0)=- \frac{1}{2}$ et on peut poser $\varphi(x)\underset{x\rightarrow0}{=}- \frac{1}{2}x+ax^2+bx^3+o(x^3)$. On obtient

\begin{align*}\ensuremath
e^{x+\varphi(x)}&\underset{x\rightarrow0}{=}e^{\frac{x}{2}+ax^2+bx^3+o(x^3)}\\
 &\underset{x\rightarrow0}{=}1+\left( \frac{x}{2}+ax^2+bx^3\right)+ \frac{1}{2}\left( \frac{x}{2}+ax^2\right)^2+ \frac{1}{6}\left( \frac{x}{2}\right)^3+o(x^3)\\
  &\underset{x\rightarrow0}{=}1+ \frac{x}{2}+\left(a+ \frac{1}{8}\right)x^2+\left(b+ \frac{a}{2}+ \frac{1}{48}\right)x^3+o(x^3).
\end{align*}

L'égalité $e^{x+\varphi(x)}+\varphi(x)-1=0$ fournit alors $a+ \frac{1}{8}+a=0$ et $b+ \frac{a}{2}+ \frac{1}{48}+b=0$ ou encore $a=- \frac{1}{16}$ et $b= \frac{1}{192}$.

\textbf{2ème solution.} On a déjà $\varphi(0)=0$ et $\varphi'(0)=0$. En dérivant l'égalité $(*)$, on obtient

\begin{center}
$\varphi''(x)=- \frac{(1+\varphi'(x))e^{x+\varphi(x)}(e^{x+\varphi(x)}+1)-(1+\varphi'(x))e^{x+\varphi(x)}(e^{x+\varphi(x)})}{\left(e^{x+\varphi(x)}+1\right)^2}=- \frac{(1+\varphi'(x))e^{x+\varphi(x)}}{\left(e^{x+\varphi(x)}+1\right)^2}$,
\end{center}

et donc $ \frac{\varphi''(0)}{2}=- \frac{ \frac{1}{2}}{2\times2^2}=- \frac{1}{16}$. De même,

\begin{center}
$\varphi^{(3)}(x)=-\varphi''(x) \frac{e^{x+\varphi(x)}}{\left(e^{x+\varphi(x)}+1\right)^2}-(1+\varphi'(x))e^{x+\varphi(x)} \frac{(1+\varphi'(x))}{\left(e^{x+\varphi(x)}+1\right)^2}+(1+\varphi'(x))e^{x+\varphi(x)} \frac{2(1+\varphi'(x))e^{x+\varphi(x)}}{\left(e^{x+\varphi(x)}+1\right)^3}$,
\end{center}

et donc $ \frac{\varphi^{(3)}(0)}{6}= \frac{1}{6}\left( \frac{1}{8}\times \frac{1}{4}- \frac{1}{2}\times \frac{1/2}{4}+ \frac{1}{2}\times \frac{1}{8}\right)= \frac{1}{192}$. La formule de \textsc{Taylor}-\textsc{Young} refournit alors

\begin{center}
\shadowbox{
$\varphi(x)\underset{x\rightarrow0}{=}- \frac{x}{2}- \frac{x^2}{16}+ \frac{x^3}{384}+o(x^3)$.
}
\end{center}
\fincorrection
\nocorrection
\nocorrection
\nocorrection
\nocorrection
\nocorrection
\nocorrection
\nocorrection
\correction{005561}
 Soit $\varphi$ une application de classe $C^2$ sur $\Rr$ puis $f$ l'application définie sur $U$ par $\forall(x,y)\in U$, $f(x,y)=\varphi\left(\frac{y}{x}\right)$ vérifie :

\begin{center}
$\frac{\partial^2f}{\partial x^2}-\frac{\partial^2f}{\partial y^2}=\frac{y}{x^3}$.
\end{center}

\begin{align*}\ensuremath
\frac{\partial^2f}{\partial x^2}-\frac{\partial^2f}{\partial y^2}&=\frac{\partial^2}{\partial x^2}\left(\varphi\left(\frac{y}{x}\right)\right)-\frac{\partial^2}{\partial y^2}\left(\varphi\left(\frac{y}{x}\right)\right)\\
 &=\frac{\partial}{\partial x}\left(-\frac{y}{x^2}\varphi'\left(\frac{y}{x}\right)\right)-\frac{\partial}{\partial y}\left(\frac{1}{x}\varphi'\left(\frac{y}{x}\right)\right)=\frac{2y}{x^3}\varphi'\left(\frac{y}{x}\right)+\frac{y^2}{x^4}\varphi''\left(\frac{y}{x}\right)-\frac{1}{x^2}\varphi''\left(\frac{y}{x}\right)\\
 &=\frac{1}{x^2}\left(\frac{2y}{x}\varphi'\left(\frac{y}{x}\right)+\left(\frac{y^2}{x^2}-1\right)\varphi''\left(\frac{y}{x}\right)\right).
\end{align*}
Puis, quand $(x,y)$ décrit $U$, $\frac{y}{x}$ décrit $\Rr$ (car $\frac{y}{1}$ décrit déjà $\Rr$)

\begin{align*}\ensuremath
\forall(x,y)\in U,\;\frac{\partial^2f}{\partial x^2}(x,y)-\frac{\partial^2f}{\partial y^2}(x,y)=\frac{y}{x^3}&\Leftrightarrow\forall(x,y)\in U,\;\frac{2y}{x}\varphi'\left(\frac{y}{x}\right)+\left(\frac{y^2}{x^2}-1\right)\varphi''\left(\frac{y}{x}\right)=\frac{y}{x}\\
 &\Leftrightarrow\forall t\in\Rr,\;2t\varphi'(t)+(t^2-1)\varphi''(t)=t\\
 &\Leftrightarrow\exists\lambda\in\Rr/\;\forall t\in\Rr,\;(t^2-1)\varphi'(t)=\frac{t^2}{2}+\lambda\quad(*)
\end{align*}
Maintenant, $\frac{t^2}{2}+\lambda$ ne s'annule pas en $\pm1$, l'égalité $(*)$ fournit une fonction $\varphi$ telle que $\varphi'$ n'a pas une limite réelle en $\pm1$. Une telle solution n'est pas de classe $C^2$ sur $\Rr$. Donc nécessairement $\lambda=-\frac{1}{2}$ puis

\begin{align*}\ensuremath
\forall(x,y)\in U,\;\frac{\partial^2f}{\partial x^2}(x,y)-\frac{\partial^2f}{\partial y^2}(x,y)=\frac{y}{x^3}&\Leftrightarrow\forall t\in\Rr,\;(t^2-1)\varphi'(t)=\frac{t^2-1}{2}\Leftrightarrow\forall t\in\Rr\setminus\{-1,1\},\;\varphi'(t)=\frac{1}{2}\\
 &\Leftrightarrow\forall t\in\Rr,\;\varphi'(t)=\frac{1}{2}\;(\text{par continuité de}\;\varphi'\;\text{en}\;\pm1)\\
 &\Leftrightarrow\exists\lambda\in\Rr/\;\forall t\in\Rr,\;\varphi(t)=\frac{t}{2}+\lambda.
\end{align*}
\fincorrection
\correction{005562}
\begin{enumerate}
 \item  $\left\{
\begin{array}{l}
u=x+y\\
v=x+2y
\end{array}
\right.\Leftrightarrow\left\{
\begin{array}{l}
x=2u-v\\
y=-u+v
\end{array}
\right.$. L'application $(x,y)\mapsto(u,v)$ est un $C^1$-difféomorphisme de $\Rr^2$ sur lui-même.
Pour $(u,v)\in\Rr^2$, posons alors $g(u,v)=f(2u-v,u+v)=f(x,y)$ de sorte que $\forall(x,y)\in\Rr^2$, $f(x,y)=g(x+y,x+2y)=g(u,v)$. $f$ est de classe $C^1$ sur $\Rr^2$ si et seulement si $g$ est de classe $C^1$ sur $\Rr^2$ et

\begin{align*}\ensuremath
 2\frac{\partial f}{\partial x}(x,y)-\frac{\partial f}{\partial y}(x,y)&=2\frac{\partial}{\partial x}\left(g(x+y,x+2y)\right)-\frac{\partial}{\partial y}\left(g(x+y,x+2y)\right)\\
  &=2\left(\frac{\partial u}{\partial x}\times\frac{\partial g}{\partial u}(u,v)+\frac{\partial v}{\partial x}\times\frac{\partial g}{\partial v}(u,v)\right)-\left(\frac{\partial u}{\partial y}\times\frac{\partial g}{\partial u}(u,v)+\frac{\partial v}{\partial y}\times\frac{\partial g}{\partial v}(u,v)\right)\\
  &=2\left(\frac{\partial g}{\partial u}(u,v)+\frac{\partial g}{\partial v}(u,v)\right)-\left(\frac{\partial g}{\partial u}(u,v)+2\frac{\partial g}{\partial v}(u,v)\right)=\frac{\partial g}{\partial u}(u,v).
 \end{align*}
 
 
 Par suite,
 
 \begin{align*}\ensuremath
 \forall(x,y)\in\Rr^2,\;2\frac{\partial f}{\partial x}(x,y)-\frac{\partial f}{\partial y}(x,y)=0&\Leftrightarrow\forall(u,v)\in\Rr^2,\;\frac{\partial g}{\partial u}(u,v)=0\\
  &\Leftrightarrow\exists F~:~\Rr\rightarrow\Rr\;\text{de classe}\;C^1\;\text{telle que}\;\forall (u,v)\in\Rr^2,\;g(u,v)=F(v)\\
  &\Leftrightarrow\exists F~:~\Rr\rightarrow\Rr\;\text{de classe}\;C^1\;\text{telle que}\;\forall (x,y)\in\Rr^2,\;f(x,y)=F(x+2y).
 \end{align*}
 \item  On pose $r=\sqrt{x^2+y^2}$ et $\theta=\Arctan\left(\frac{y}{x}\right)$ de sorte que $x=r\cos\theta$ et $y=r\sin\theta$. On pose $f(x,y)=f(r\cos\theta,r\sin\theta)=g(r,\theta)$. On sait que $\frac{\partial r}{\partial x}=\cos\theta$, $\frac{\partial r}{\partial y}=\sin\theta$, $\frac{\partial \theta}{\partial x}=-\frac{\sin\theta}{r}$, $\frac{\partial \theta}{\partial y}=\frac{\cos\theta}{r}$
\begin{align*}\ensuremath
x\frac{\partial f}{\partial x}+ y\frac{\partial f}{\partial y}=r\cos\theta\left(\cos\theta\frac{\partial g}{\partial r}-\frac{\sin\theta}{r}\frac{\partial g}{\partial \theta}\right)+r\sin\theta\left(\sin\theta\frac{\partial g}{\partial r}+\frac{\cos\theta}{r}\frac{\partial g}{\partial \theta}\right)=r\frac{\partial g}{\partial r},
\end{align*}
puis

\begin{align*}\ensuremath
\forall(x,y)\in D,\;x\frac{\partial f}{\partial x}(x,y)&+ y\frac{\partial f}{\partial y}(x,y)=\sqrt{x^2+y^2}\Leftrightarrow\forall r>0,\;r\frac{\partial g}{\partial r}(r,\theta)=r\Leftrightarrow\forall r>0,\;\frac{\partial g}{\partial r}(r,\theta)=1\\
 &\Leftrightarrow\exists \varphi\;\text{de classe}\;C^1\;\text{sur}\;\left]-\frac{\pi}{2},\frac{\pi}{2}\right[/\;\forall(r,\theta)\in]0,+\infty[\times\left]-\frac{\pi}{2},\frac{\pi}{2}\right[,\;g(r,\theta)=r+\varphi(\theta)\\
 &\Leftrightarrow\exists \varphi\;\text{de classe}\;C^1\;\text{sur}\;\left]-\frac{\pi}{2},\frac{\pi}{2}\right[/\;\forall(x,y)\in D,\;f(x,y)=\sqrt{x^2+y^2}+\varphi\left(\Arctan\frac{y}{x}\right)\\
  &\Leftrightarrow\exists \psi\;\text{de classe}\;C^1\;\text{sur}\;\Rr/\;\forall(x,y)\in D,\;f(x,y)=\sqrt{x^2+y^2}+\psi\left(\frac{y}{x}\right).
\end{align*}
\end{enumerate}
\fincorrection
\nocorrection
\nocorrection
\nocorrection
\nocorrection
\nocorrection
\nocorrection
\nocorrection
\nocorrection
\nocorrection
\nocorrection
\nocorrection
\nocorrection
\nocorrection
\nocorrection
\nocorrection
\nocorrection
\nocorrection
\nocorrection
\nocorrection
\nocorrection
\nocorrection
\nocorrection
\nocorrection
\correction{004381}
\begin{enumerate}
  \item $\frac1{30}$.
  \item $0$.
  \item $\frac\pi4 ab(a^2+b^2)$.
  \item $\frac{96}{35}$.
  \item $\frac\pi2$.
  \item $\pi(1-\ln 2)$.
  \item $\frac56$.
  \item $2(\ln 2-1)$.
  \item $\frac{3\pi}2$.
  \item $\frac{65\pi}{48}$.
  \item $\frac{2\sqrt2}3 a^3$.
  \item $\frac7{45}$.
  \item $\pi(1-\frac1e )$.
  \item $\frac{(e^{2p}-1)^2}3 \quad (x=u^2v,y=uv^2)$.
\end{enumerate}
 \fincorrection
\correction{004382}
Poser $u=x, v=x+y$. On obtient $I=\frac{2}{1701}$.
\fincorrection
\correction{004383}
symétrie $+$ passage en polaires.
             $I = \frac34\pi-\frac{11}6$.
\fincorrection
\correction{004384}
\begin{enumerate}
  \item $\frac12 \ln\left(\frac{32}{27}\right)$.
  \item $2\pi a^2\Arcsin\frac Ra -2\pi R\sqrt{a^2-R^2}$.
  \item $\frac1{720}$.
  \item $\frac34 -\ln2$.
  \item $\frac{\pi R^2a^2}4 (a^2+3R^2)$.
  \item $\frac\pi2 (1-\ln 2)$.
  \item $\frac{4\pi}{15} abc (a^2+b^2)$.
\end{enumerate}
 \fincorrection
\correction{004385}
\begin{enumerate}
  \item Intégrer en $z$ d'abord :
             $\frac1{(1+x^2z^2)(1+y^2z^2)} =
              \frac 1{x^2-y^2}\Bigl(\frac{x^2}{1+x^2z^2} - \frac{y^2}{1+y^2z^2}\Bigr)$.
             On obtient $I = \pi\ln 2$.
  \item Intégrer $I$ en $x$ et $y$ d'abord. On obtient
    $I =  \int_{z=0}^{+\infty} \Bigl(\frac{\Arctan z}z\Bigr)^2 d z$.
\end{enumerate}
 \fincorrection
\correction{004386}
$I=\frac{\pi\ln2}8$.
\fincorrection
\correction{004387}
$\frac12\pi^2Rr^2(4R^2+3r^2)$.
\fincorrection
\correction{004388}
$2\pi b\left(a^2-\frac {b^2}3\right)$.
\fincorrection
\correction{004389}
\begin{enumerate}
  \item 
  \item 
  \item Fubini $ \Rightarrow  I=\frac{\pi^2}8$.
\end{enumerate}
 \fincorrection
\nocorrection
\correction{004391}
\begin{enumerate}
  \item $2A = \left( \int_{t=0}^1 \frac{d t}{1+t^2}\right)^2  \Rightarrow  A = \frac{\pi^2}{32}$.
  \item 
  \item 
  \item $B+C = \frac D2 $, $B-C = -D$.
  \item $C = -\frac{3\pi^2}{32}$, $D = -\frac{\pi^2}8$.
\end{enumerate}
 \fincorrection
\correction{004392}
\begin{enumerate}
  \item $A= \frac\pi6 -\frac{\sqrt3}4 \ln3$.
  \item $4ab\arctan \frac ba$.
\end{enumerate}
 \fincorrection
\correction{004393}
Formule de Green : ${\cal A} = \frac{3\pi a^2}8$.
\fincorrection
\correction{004394}
Formule de Green. $A = \frac43(p_2-p_1)(q_2-q_1)$.
\fincorrection
\correction{004395}
Formule de Green. $A = \pi a^2$.
\fincorrection
\correction{004396}
\begin{enumerate}
  \item $V = \frac{4\pi}3 (1-\sqrt{1-a^2}^3\,)$.
  \item $V = \frac{2\pi}3 (2-\sqrt{2})$.
  \item $V = 2\pi^2Rr^2$.
\end{enumerate}
 \fincorrection
\correction{004397}
$V=\frac{4\pi p^3}{3\lambda^4}$.
\fincorrection
\correction{004398}
$\frac{\pi a^3}{12\sqrt2} (3\ln(1+\sqrt2\,) - \sqrt2\,)$.
\fincorrection
\correction{004399}
hauteur${} = \alpha R$ avec $\alpha^3-3\alpha+1 = 0  \Rightarrow  \alpha \approx 0.347$.
\fincorrection
\nocorrection
\nocorrection
\correction{004402}
C'est manifestement vrai pour $\psi\equiv 1$ et aussi pour $\psi(t)=t$.
De manière générale, si $\psi(t) = t^k$ avec $k\in\N$ alors pour $n\ge k$,
$(x_1+\dots+x_n)^k$ est une somme de $n^k$ monômes parmi lesquels il y a
$n(n-1)\dots(n-k+1)$ monômes où chaque variable apparaît avec l'exposant $0$ ou $1$.
On a alors~:
$$\Lambda_n(\psi) = \frac{n(n-1)\dots(n-k+1)}{n^k}\Bigl( \int_{x=0}^1xf(x)\,d x\Bigr)^k\Bigl( \int_{x=0}^1f(x)\,d x\Bigr)^{n-k}
+\Bigl(1-\frac{n(n-1)\dots(n-k+1)}{n^k}\Bigr) O(1),$$
ce qui prouve que $\Lambda_n(\psi) \to\psi\Bigl( \int_{x=0}^1xf(x)\,d x\Bigr)$ (lorsque $n\to\infty$) lorsque
$\psi(t) = t^k$. Par linéarité, cette relation est encore vraie pour tout $\psi$
polynôme. On conclut pour $\psi$ continue quelconque avec le théorème de
Stone-Weierstrass.
\fincorrection
\correction{005908}
\begin{enumerate}
 \item  Représentons le domaine $D=\{(x,y)\in\Rr^2/\;x\leqslant1,\;y\leqslant1,\;x+y\geqslant1\}$.

$$\includegraphics{../images/img005908-1}$$



\begin{align*}\ensuremath
I&=\displaystyle\iint_{D}(x+y)\;dxdy=\int_{0}^{1}\left(\int_{1-x}^{1}(x+y)dy\right)dx\;(\text{ou aussi}\;\int_{0}^{1}\left(\int_{0}^{1-y}(x+y)dx\right)dy)\\
 &=\int_{0}^{1}\left[xy+ \frac{y^2}{2}\right]_{y=1-x}^{y=1}dx=\int_{0}^{1}\left(x+ \frac{1}{2}-x(1-x)- \frac{(1-x)^2}{2}\right)dx\\
 &=\int_{0}^{1}\left( \frac{x^2}{2}+x\right)dx= \frac{1}{6}+ \frac{1}{2}= \frac{2}{3}.
\end{align*}

\begin{center}
\shadowbox{
$\displaystyle\iint_{D}(x+y)\;dxdy= \frac{2}{3}$.
}
\end{center}

\item  Si on pose pour $(x,y)\in/mbr^2$, $f(x,y)=|x+y|$ alors pour tout $(x,y)\in\Rr^2$, $f(-x,-y)=f(x,y)$ ou encore $f$ prend les mêmes valeurs en deux points symétriques par rapport à $O$. Puisque le point $O$ est centre de symétrie de $[-1,1]^2$, on en déduit que

\begin{align*}\ensuremath
I&=\displaystyle\iint_{-1\leqslant x,y\leqslant1,\;x+y\geqslant0}f(x,y)\;dxdy+\displaystyle\iint_{-1\leqslant x,y\leqslant1,\;x+y\leqslant0}f(x,y)\;dxdy\\
 &=2\iint_{-1\leqslant x,y\leqslant1,\;x+y\geqslant0}(x+y)\;dxdy=2\int_{-1}^{1}\left(\int_{-x}^{1}(x+y)\;dy\right)dx\\
 &=2\int_{-1}^{1}\left[xy+ \frac{y^2}{2}\right]_{y=-x}^{y=1}dx=2\int_{-1}^{1}\left(x+ \frac{1}{2}+x^2- \frac{x^2}{2}\right)dx\\
 &=2\left( \frac{1}{2}\times \frac{2}{3}+ \frac{1}{2}\times2\right)= \frac{8}{3}.
\end{align*}

\begin{center}
\shadowbox{
$\displaystyle\iint_{[-1,1]^2}|x+y|\;dxdy= \frac{8}{3}$.
}
\end{center}

\item  Représentons le domaine $D=\{(x,y)\in\Rr^2/\;0\leqslant x\leqslant1,\;0\leqslant y\leqslant1,\;x^2\leqslant y\leqslant\sqrt{x}\}$.

$$\includegraphics{../images/img005908-2}$$


\begin{align*}\ensuremath
I&=\int_{0}^{1}\left(\int_{x^2}^{\sqrt{x}}y\;dy\right)x\;dx=\int_{0}^{1}x\left[ \frac{y^2}{2}\right]_{y=x^2}^{y=\sqrt{x}}\;dx= \frac{1}{2}\int_{0}^{1}x(x-x^4)\;dx= \frac{1}{2}\left( \frac{1}{3}- \frac{1}{6}\right)\\
&= \frac{1}{12}.
\end{align*}

\item  En passant en polaires, on obtient

\begin{align*}\ensuremath
I&=\iint_{x^2+y^2\leqslant1} \frac{1}{1+x^2+y^2}\;dxdy=\iint_{0\leqslant r\leqslant1,\;0\leqslant\theta\leqslant2\pi} \frac{1}{1+r^2}\;rdrd\theta\\
 &=\left(\int_{0}^{1} \frac{r}{1+r^2}\;dr\right)\times\left(\int_{0}^{2\pi}d\theta\right)\;(\text{intégrales indépendantes})\\
 &=2\pi\times\left[ \frac{1}{2}\ln(1+r^2)\right]_0^1=\pi\ln2.
\end{align*}

\begin{center}
\shadowbox{
$\displaystyle\iint_{x^2+y^2\leqslant1} \frac{1}{1+x^2+y^2}\;dxdy=\pi\ln2$.
}
\end{center}

\item  Posons $D=\{(x,y)\in\Rr^2/\;x\leqslant x^2+y^2\leqslant1\}$. Puisque $x\leqslant x^2+y^2\Leftrightarrow\left(x- \frac{1}{2}\right)^2+y^2\geqslant \frac{1}{4}$, $D$ est l'intersection de l'intérieur du disque de centre $O$ et de rayon $1$, bord compris, et de l'extérieur du disque de centre $\left( \frac{1}{2},0\right)$ et de rayon $ \frac{1}{2}$, bord compris. Soit $M$ un point du plan. On note $(r,\theta)$ un couple de coordonnées polaires de $M$ tel que $r\geqslant0$ et $\theta\in[0,2\pi]$.

\begin{center}
$M\in D\Leftrightarrow r\cos\theta\leqslant r^2\leqslant1\Leftrightarrow r=0\;\text{ou}\;(0<r\leqslant1\;\text{et}\;r\geqslant\cos\theta$.
\end{center}

$$\includegraphics{../images/img005908-3}$$


En passant en polaires, on obtient

\begin{align*}\ensuremath
I&=2\iint_{x\leqslant x^2+y^2\leqslant1,\;y\geqslant0} \frac{1}{(1+x^2+y^2)^2}dxdy=2\left(\int_{0}^{\pi/2}\left(\int_{\cos\theta}^{1} \frac{r}{(1+r^2)^2}dr\right)d\theta+\int_{\pi/2}^{\pi}\left(\int_{0}^{1} \frac{r}{(1+r^2)^2}dr\right)d\theta\right)\\
 &=2\left(\int_{0}^{\pi/2}\left[- \frac{1}{2(1+r^2)}dr\right]_{\cos\theta}^{1}d\theta+\int_{\pi/2}^{\pi}\left[- \frac{1}{2(1+r^2)}dr\right]_{0}^{1}d\theta\right)\\
 &=\int_{0}^{\pi/2}\left( \frac{1}{1+\cos^2\theta}- \frac{1}{2}\right)d\theta+\int_{\pi/2}^{\pi} \frac{1}{2}\;d\theta=\int_{0}^{\pi/2} \frac{1}{1+\cos^2\theta}d\theta\\
 &=\int_{0}^{\pi/2} \frac{1}{ \frac{1}{\cos^2\theta}+1} \frac{d\theta}{\cos^2\theta}=\int_{0}^{\pi/2} \frac{1}{2+\tan^2\theta}d(\tan\theta)=\int_{0}^{+\infty} \frac{1}{t^2+2}\;dt=\left[ \frac{1}{\sqrt{2}}\Arctan\left( \frac{t}{\sqrt{2}}\right)\right]_0^{+\infty}= \frac{\pi}{2\sqrt{2}}.
\end{align*}

\begin{center}
\shadowbox{
$\iint_{x\leqslant x^2+y^2\leqslant1} \frac{1}{(1+x^2+y^2)^2}dxdy= \frac{\pi}{2\sqrt{2}}$.
}
\end{center}

\item  

\begin{align*}\ensuremath
I&=\iiint_{0\leqslant x\leqslant y\leqslant z\leqslant1}xyzdxdydz=\int_{0}^{1}\left(\int_{x}^{1}\left(\int_{y}^{1}zdz\right)ydy\right)xdx=\int_{0}^{1}\left(\int_{x}^{1} \frac{1}{2}(1-y^2)ydy\right)xdx\\
 &= \frac{1}{2}\int_{0}^{1}\left(\int_{x}^{1}(y-y^3)dy\right)xdx= \frac{1}{2}\int_{0}^{1}\left[ \frac{y^2}{2}- \frac{y^4}{4}\right]_x^1xdx= \frac{1}{2}\int_{0}^{1}\left( \frac{1}{4}- \frac{x^2}{2}+ \frac{x^4}{4}\right)xdx\\
 &= \frac{1}{8}\int_{0}^{1}(x^5-2x^3+x)dx= \frac{1}{8}\left( \frac{1}{6}- \frac{1}{2}+ \frac{1}{2}\right)= \frac{1}{48}.
\end{align*}

\begin{center}
\shadowbox{
$\iiint_{0\leqslant x\leqslant y\leqslant z\leqslant1}xyzdxdydz= \frac{1}{48}$.
}
\end{center}

\item  En sommant par tranches, on obtient

\begin{align*}\ensuremath
I&=\displaystyle\iiint_{\sqrt{x}+\sqrt{y}+\sqrt{z}\leqslant1}zdxdydz=\int_{0}^{1}\left(\iint_{\sqrt{x}+\sqrt{y}\leqslant1-\sqrt{z}}dxdy\right)zdz\\
 &=\int_{0}^{1}\left(\iint_{\sqrt{u}+\sqrt{v}\leqslant1}(1-\sqrt{z})^4dudv\right)zdz\;(\text{en posant}\;x=(1-\sqrt{z})^2u\;\text{et}\;y=(1-\sqrt{z})^2v)\\
 &=\mathcal{A}(D)\times\int_{0}^{1}z(1-\sqrt{z})^4dz\;\text{où}\;D=\{(u,v)\in\Rr^2/\;\sqrt{u}+\sqrt{v}\leqslant1\}.
\end{align*}

Maintenant,

\begin{center}
$\mathcal{A}(D)=\int_{0}^{1}\left(\int_{0}^{(1-\sqrt{u})^2}dv\right)du=\int_{0}^{1}(1-2\sqrt{u}+u)du=1- \frac{4}{3}+ \frac{1}{2}= \frac{1}{6}$
\end{center}

et

\begin{center}
$\int_{0}^{1}z(1-\sqrt{z})^4dz=\int_{0}^{1}(z-4z^{3/2}+6z^2-4z^{5/2}+z^3)\;dz= \frac{1}{2}- \frac{8}{5}+2- \frac{8}{7}+ \frac{1}{4}= \frac{1}{140}$.
\end{center}

Finalement

\begin{center}
\shadowbox{
$\displaystyle\iiint_{\sqrt{x}+\sqrt{y}+\sqrt{z}\leqslant1}zdxdydz= \frac{1}{840}$.
}
\end{center}

\end{enumerate}
\fincorrection
\correction{005910}
\ 
$$\includegraphics{../images/img005910-1}$$


L'aire du domaine considéré $D=\{(x,y)\in\Rr^2/\;2p_1x\leqslant y^2\leqslant 2p_2x\;\text{et}\;2q_2y\leqslant x^2\leqslant2q_2x\}$ est

\begin{center}
$\mathcal{A}=\displaystyle\iint_{D}dxdy$.
\end{center}

Pour $(x,y)\in D^2$, posons $p= \frac{y^2}{2x}$ et $q= \frac{x^2}{2y}$ ou encore considérons l'application $\begin{array}[t]{cccc}
\varphi~:&D&\rightarrow&[p_1,p_2]\times[q_1,q_2]\\
 &(x,y)&\mapsto&\left( \frac{y^2}{2x}, \frac{x^2}{2y}\right)
 \end{array}$ et vérifions que $\varphi$ est un $C^1$-difféomorphisme.
 
 

 

\textbullet~Pour chaque $(x,y)\in D^2$, on $2p_1x\leqslant y^2\leqslant 2p_2x$ et $2q_1y\leqslant x^2\leqslant2q_2y$ ou encore $p_1\leqslant \frac{y^2}{2x}\leqslant p_2$ et $q_1\leqslant \frac{x^2}{2y}\leqslant q_2$. Donc $\varphi$ est bien une application.

\textbullet~Soit $(p,q)\in[p_1,p_2]\times[q_1,q_2]$. Pour $(x,y)\in(]0,+\infty[)^2$,

\begin{center}
$\left\{
\begin{array}{l}
 \frac{y^2}{2x}=p\\
\rule{0mm}{7mm} \frac{x^2}{2y}=q
\end{array}
\right.\Leftrightarrow\left\{
\begin{array}{l}
y= \frac{x^2}{2q}\\
\rule{0mm}{7mm} \frac{(x^2/2q)^2}{2x}=p
\end{array}
\right.\Leftrightarrow\left\{
\begin{array}{l}
x=\sqrt[3]{8pq^2}\\
y=\sqrt[3]{8p^2q}
\end{array}
\right.$
\end{center}

Donc, l'équation $\varphi(x,y)=(p,q)$ a exactement une solution $(x_0,y_0)$ dans $]0,+\infty[^2$. De plus, puisque $ \frac{y_0^2}{2x_0}=p\in[p_1,p_2]$ et $ \frac{x_0^2}{2y_0}=q\in[q_1,q_2]$, on a $2p_1x_0\leqslant y_0^2\leqslant 2p_2x_0$ et $2q_1y_0\leqslant x_0^2\leqslant 2q_2y_0$ et donc $(x_0,y_0)\in D^2$. Donc $\varphi$ est une bijection.

\textbullet~$\varphi$ est de classe $C^1$ sur $D$ et pour $(x,y)\in D^2$,

\begin{center}
$ \frac{D(p,q)}{D(x,y)}=J(\varphi)(x,y)=\left|
\begin{array}{cc}
- \frac{y^2}{2x^2}& \frac{y}{x}\\
 \frac{x}{y}&- \frac{x^2}{2y^2}
\end{array}
\right|= \frac{1}{4}-1=- \frac{3}{4}\neq0$.
\end{center}

Ainsi, $\varphi$ est une bijection de $D$ sur $[p_1,p_2]\times[q_1,q_2]$, de classe $C^1$ sur $D$ et son jacobien ne s'annule pas sur $D$. On sait alors que $\varphi$ est un $C^1$-difféomorphisme de $D$ sur $[p_1,p_2]\times[q_1,q_2]$.

Posons alors $(p,q)=\varphi(x,y)$ dans $\displaystyle\iint_{D}dxdy$. On obtient

\begin{center}
$\mathcal{A}=\displaystyle\iint_{D}dxdy=\displaystyle\iint_{[p_1,p_2]\times[q_1,q_2]}\left| \frac{D(x,y)}{D(p,q)}\right|dpdq=\displaystyle \frac{4}{3}\iint_{[p_1,p_2]\times[q_1,q_2]}dpdq= \frac{4}{3}(p_2-p_1)(q_2-q_1)$.
\end{center}

\begin{center}
\shadowbox{
$\mathcal{A}= \frac{4}{3}(p_2-p_1)(q_2-q_1)$.
}
\end{center}

\fincorrection
\correction{005911}
Pour $n\in\Nn^*$ et $R\geqslant0$, posons $B_n(R)=\{(x_1,\ldots,x_n)\in\Rr^n/\;x_1^2+\ldots+x_n^2\leqslant R^2\}$ et notons $V_n(R)$ le volume de $B_n(R)$. Par définition,

\begin{center}
$V_n(R)=\displaystyle\int\ldots\iint_{x_1^2+\ldots+x_n^2\leqslant R^2}dx_1\ldots dx_n$.
\end{center}

En posant $x_1=Ry_1$, \ldots, $x_n=Ry_n$, on a $ \frac{D(x_1,\ldots,x_n)}{D(y_1,\ldots,y_n)}=R^n$ (quand $R>0$) puis

\begin{center}
$V_n(R)=\displaystyle\int\ldots\iint_{x_1^2+\ldots+x_n^2\leqslant R^2}dx_1\ldots dx_n=R^n\displaystyle\int\ldots\iint_{y_1^2+\ldots+y_n^2\leqslant1}dy_1\ldots dy_n=R^nV_n(1)$.
\end{center}

ce qui reste vrai quand $R=0$. Pour $n\geqslant2$, on peut alors écrire

\begin{align*}\ensuremath
V_n(1)&=\displaystyle\int_{-1}^1\left(\displaystyle\int\ldots\iint_{x_1^2+\ldots+x_{n-1}^2\leqslant 1-x_n^2}dx_1\ldots dx_{n-1}\right)dx_n=\int_{-1}^{1}V_{n-1}\left(\sqrt{1-x_n^2}\right)\;dx_n\\
 &=\int_{-1}^{1}(1-x_n^2)^{(n-1)/2}V_{n-1}(1)\;dx_n=I_{n}V_{n-1}(1)
\end{align*}

où $I_{n}=\int_{-1}^{1}(1-x^2)^{(n-1)/2}\;dx$. Pour calculer $I_n$, on pose $x=\cos\theta$. On obtient

\begin{center}
$I_n=\int_{\pi}^{0}(1-\cos^2\theta)^{(n-1)/2}(-\sin\theta)\;d\theta=\int_{0}^{\pi}\sin^{n}\theta\;d\theta=2\int_{0}^{\pi/2}\sin^n\theta\;d\theta=2W_n$ (intégrales de \textsc{Wallis}).
\end{center}

Finalement,

\begin{center}
$V_1(1)=2$ et $\forall n\geqslant2$, $V_n(1)=2W_nV_{n-1}(1)$.
\end{center}

On en déduit que pour $n\geqslant2$,

\begin{center}
$V_n(1)=(2W_n)(2W_{n-1})\ldots(2W_2)V_1(1)=2^n\prod_{k=2}^{n}W_k=2^n\prod_{k=1}^{n}W_k$,
\end{center} 

 ce qui reste vrai pour $n=1$. Maintenant, il est bien connu que la suite $((n+1)W_{n+1}W_n)_{n\in\Nn}$ est constante et plus précisément que $\forall n\in\Nn$, $(n+1)W_{n+1}W_n=W_1W_0= \frac{\pi}{2}$. Donc, pour $p\in\Nn^*$,

\begin{align*}\ensuremath
V_{2p}(1)&=2^{2p}\prod_{k=1}^{2p}W_k=2^{2p}\prod_{k=1}^{p}(W_{2k-1}W_{2k})=2^{2p}\prod_{k=1}^{p} \frac{\pi}{2(2k)}= \frac{\pi^p}{p!},
\end{align*}

et de même

\begin{align*}\ensuremath
V_{2p+1}(1)&=2^{2p+1}\prod_{k=2}^{2p+1}W_k=2^{2p+1}\prod_{k=1}^{p}(W_{2k}W_{2k+1})=2^{2p+1}\prod_{k=1}^{p} \frac{\pi}{2(2k+1)}\\
 &= \frac{\pi^p2^{p+1}}{3\times5\times\ldots\times(2p+1)}= \frac{\pi^p2^{p+1}(2p)\times(2p-2)\times\ldots\times2}{(2p+1)!}= \frac{\pi^p2^{2p+1}p!}{(2p+1)!}.
\end{align*}

\begin{center}
\shadowbox{
$\forall p\in\Nn^*$, $\forall R>0$, $V_{2p}(R)= \frac{\pi^pR^{2p}}{p!}$ et $V_{2p-1}(R)= \frac{\pi^p2^{2p+1}p!R^{2p+1}}{(2p+1)!}$.
}
\end{center}

En particulier, $V_1(R)=2R$, $V_2(R)=\pi R^2$ et $V_3(R)= \frac{4}{3}\pi R^3$.
\fincorrection
\correction{005912}
 \textbf{1ère solution.} $V=\displaystyle\iiint_{x^2+ \frac{1}{2}y^2+ \frac{3}{4}z^2+xz\leqslant1}dxdydz$. Or $x^2+ \frac{1}{2}y^2+ \frac{3}{4}z^2+xz=\left(x+ \frac{z}{2}\right)^2+ \frac{y^2}{2}+ \frac{z^2}{2}$. On pose donc $u=x+ \frac{z}{2}$, $v= \frac{y}{\sqrt{2}}$ et $w= \frac{z}{\sqrt{2}}$.

\begin{center}
$ \frac{D(u,v,w)}{D(x,y,z)}=\left|
\begin{array}{ccc}
1&0& \frac{1}{2}\\
0& \frac{1}{\sqrt{2}}&0\\
0&0& \frac{1}{\sqrt{2}}
\end{array}
\right|= \frac{1}{2}$.
\end{center}

On en déduit que $ \frac{D(x,y,z)}{D(u,v,w)}=2$ puis que

\begin{center}
$V=\displaystyle\iiint_{x^2+ \frac{1}{2}y^2+ \frac{3}{4}z^2+xz\leqslant1}dxdydz=\displaystyle\iiint_{u^2+v^2+w^2\leqslant1}\left| \frac{D(x,y,z)}{D(u,v,w)}\right|dudvdw=2\times \frac{4\pi}{3}= \frac{8\pi}{3}$.
\end{center}

\textbf{2ème solution.} Supposons savoir que le volume délimité par l'ellipsoïde d'équation $ \frac{X^2}{a^2}+ \frac{Y^2}{b^2}+ \frac{Z^2}{c^2}=1$ est $ \frac{4}{3}\pi abc$. La matrice de la forme quadratique $(x,y,z)\mapsto x^2+ \frac{1}{2}y^2+ \frac{3}{4}z^2+xz$ dans la base canonique orthonormée de $\Rr^3$ est $A=\left(
\begin{array}{ccc}
1&0& \frac{1}{2}\\
0& \frac{1}{2}&0\\
 \frac{1}{2}&0& \frac{3}{4}
\end{array}
\right)$. On sait que cette matrice a $3$ valeurs propres strictement positives $\lambda= \frac{1}{a^2}$, $\mu= \frac{1}{b^2}$ et $\nu= \frac{1}{c^2}$ puis qu'il existe une base orthonormée dans laquelle l'ellipsoïde a pour équation $ \frac{X^2}{a^2}+ \frac{Y^2}{b^2}+ \frac{Z^2}{c^2}=1$. Le volume de l'ellipsoïde est alors

\begin{center}
$V= \frac{4}{3}\pi abc= \frac{4}{3} \frac{\pi}{\sqrt{\lambda\mu\nu}}= \frac{4}{3} \frac{\pi}{\sqrt{\text{det}(A)}}= \frac{4}{3} \frac{\pi}{\sqrt{ \frac{1}{4}}}= \frac{8\pi}{3}$
\end{center}

\begin{center}
\shadowbox{
$V= \frac{8\pi}{3}$.
}
\end{center}
\fincorrection
\correction{005914}
On pose déjà $x=ua$ et $y=vb$ de sorte que $ \frac{D(u,v)}{D(x,y)}=ab$. On obtient

\begin{center}
$I=\displaystyle\iint_{\frac{x^2}{a^2}+\frac{y^2}{b^2}\leqslant1}(x^2-y^2)\;dxdy=ab\displaystyle\iint_{u^2+v^2\leqslant1}(a^2u^2-b^2v^2)\;dudv$.
\end{center}

Ensuite,

\begin{align*}\ensuremath
\displaystyle\iint_{u^2+v^2\leqslant1}u^2\;dudv&=\displaystyle\iint_{u^2+v^2\leqslant1}v^2\;dudv= \frac{1}{2}\displaystyle\iint_{u^2+v^2\leqslant1}(u^2+v^2)\;dudv= \frac{1}{2}\int_{r=0}^{r=1}\int_{\theta=0}^{2\pi}r^2\times rdrd\theta\\
 &= \frac{1}{2}\int_{0}^{1}r^3dr\int_{0}^{2\pi}d\theta= \frac{1}{2}\times \frac{1}{4}\times2\pi= \frac{\pi}{4},
\end{align*}

et donc

\begin{center}
\shadowbox{
$I= \frac{\pi ab(a^2-b^2)}{4}$.
}
\end{center}
\fincorrection
\correction{004313}
$= \int_{t=x}^{2x} \left({\frac1t - \frac{5t}6 +  o(t)}\right)d t \to \ln 2$.
\fincorrection
\correction{004314}

\fincorrection
\correction{004315}
Formule de la moyenne sur $[3,x]$ et $[x,x+x^2]  \Rightarrow  \lim = 0$.
\fincorrection
\correction{004316}
$\ln 2$.
\fincorrection
\correction{004317}
$\sum_{n=1}^N (-1)^n \int_0^{+\infty} \frac{dt}{(1+t^3)^n} =
	       \int_0^{+\infty}\Bigl(1-\frac{(-1)^N}{(1+t^3)^N}\Bigr)\frac{dt}{2+t^3}
	      \to  \int_0^{+\infty}\frac{dt}{2+t^3} =
	      \frac{\pi2^{5/3}}{3\sqrt3}$ lorsque $N\to\infty$.
\fincorrection
\correction{004318}
\begin{enumerate}
  \item 
  \item $I(x) =  \int_{t=x}^{+\infty} \frac{\sin(t-x)}t\,d t
                   = \cos x  \int_{t=x}^{+\infty} \frac{\sin t}t\,d t
                   - \sin x  \int_{t=x}^{+\infty} \frac{\cos t}t\,d t
                   \to 0$ lorsque $x\to+\infty$.
\end{enumerate}
 \fincorrection
\correction{004319}
$\frac\pi2f(0)$.
\fincorrection
\correction{004320}
$t=ux$ puis intégration par parties $ \Rightarrow \sim \frac 1{x^2}$.
\fincorrection
\correction{004321}
$f'(x) = \frac{2\ln x}{a(1+x^2)}  \Rightarrow  f(1+h) = \frac1{2a}(h^2-h^3) +  o(h^3)$.
\fincorrection
\correction{004322}
\begin{enumerate}
  \item 
  \item Soit $\varepsilon > 0$ : Pour $x$ assez petit,
             $\bigl|f(t)^x-1-x\ln(f(t))\bigr| \le \varepsilon x$ car $\ln f$
             est borné sur $[a,b]$.\par
             Donc $\left| \int_{t=a}^b f(t)^xd t - 1 - x \int_{t=a}^b\ln(f(t))\,d t\right| \le \varepsilon x$,
             et   $\left|\ln\left({ \int_{t=a}^b f(t)^xd t}\right) - x \int_{t=a}^b\ln(f(t))\,d t\right| \le 2\varepsilon x$.
\end{enumerate}
 \fincorrection
\nocorrection
\nocorrection
\correction{004325}
\begin{enumerate}
  \item Couper en $ \int_{t=0}^{1-\varepsilon} +  \int_{t=1-\varepsilon}^1$
  \item $=\left[\frac{t\ln(1+t^n)}n\right]_{t=0}^1 - \frac1n \int_{t=0}^1 \ln(1+t^n)\,d t \sim \frac{\ln 2}n$.
  \item $\frac1n \int_{t=0}^1  \frac{d t}{\sqrt{1+t}+1} = \frac{2\sqrt2-2+2\ln(2\sqrt2-2)}n$.
\end{enumerate}
 \fincorrection
\correction{004326}
$1-\frac{\ln 2}n +  o\Bigl(\frac1n\Bigr)$.
\fincorrection
\correction{004327}
$1+\frac1n \int_{t=0}^1\frac{d t}{\sqrt{1+t}+1} +  o\Bigl(\frac1n\Bigr)
        = 1+\frac{2\sqrt2-2+2\ln(2\sqrt2-2)}n +  o\Bigl(\frac1n\Bigr)$.
\fincorrection
\correction{004328}
$u=t^n  \Rightarrow  \sim \frac 1n  \int_{u=1}^e \frac{\sqrt{1+u}}u\,d u$.
\fincorrection
\nocorrection
\correction{004330}
$f(0)$.
\fincorrection
\correction{004331}
Soit $f_n(x) = (1-x/n)^n$ si $0\le x \le n$ et $f_n(x) = 0$ si
$x>n$. Alors $f_n(x)$ converge simplement vers $e^{-x}$ et il y a
convergence dominée.
\fincorrection
\correction{004332}
$f(x) = \cos x$.
\fincorrection
\correction{004333}
\begin{enumerate}
  \item $I_n+I_{n+2} = \frac1{n+1}$.
  \item $I_{2k} = \frac1{2k-1} - \frac1{2k-3} + \dots + \frac{(-1)^{k-1}}1 + (-1)^k\frac\pi4$,\par
         $I_{2k+1} = \frac1{2k} - \frac1{2k-2} + \dots + \frac{(-1)^{k-1}}2 - (-1)^k\ln\sqrt2$.
  \item $\frac11-\frac13+\frac15 - \dots = \frac\pi4$ et $\frac11-\frac12+\frac13-\dots=\ln 2$.
\end{enumerate}
 \fincorrection
\correction{004334}
$ \int_0^1\frac{t^n-t^{2n}}{1-t}\,d t = \frac1{n+1} + \dots + \frac1{2n}
              \to \ln 2$ lorsque $n\to\infty$.
\fincorrection
\correction{004335}
\begin{enumerate}
  \item 
  \item $g(c)$.
\end{enumerate}
 \fincorrection
\correction{004336}
$\Phi(x) =  \int_{t=0}^x f^2(t)\,d t  \Rightarrow  \Phi'\Phi^2 \to \ell^2
          \Rightarrow  \Phi^3 \sim 3\ell^2 x  \Rightarrow  f = \sqrt{\Phi'} \sim \sqrt[3]{\frac \ell{3x}}$.
\fincorrection
\nocorrection
\correction{004338}
\begin{enumerate}
  \item 
  \item 
  \item 
  \item Pour $\alpha=0$ on a $h(x) =  \int_{t=0}^x\frac{\sin t}t\,d t$, quantité
    bornée car l'intégrale converge en~$+\infty$.

    Pour $\alpha=1$ on a $h(x) = \cos x \int_{t=0}^x \frac{\cos t\sin t}t\,d t
                               + \sin x \int_{t=0}^x \frac{\sin^2 t}t\,d t$,
    quantité non bornée car la deuxième intégrale diverge en~$+\infty$.

    Pour $0<\alpha<1$, développer le $\cos(x-t)$ puis linéariser les produits
    obtenus. On obtient quatre intégrales convergentes, donc $h$ est bornée.
\end{enumerate}
 \fincorrection
\correction{004339}
\begin{enumerate}
  \item $\frac{\ln(1+a)}a$.
  \item $=-\varphi'(a) = \frac{\ln(1+a)}{a^2} - \frac1{a(1+a)}$.
\end{enumerate}
 \fincorrection
\nocorrection
\nocorrection
\correction{004342}
$f'(x) = \frac\pi{x+1}$, $f(x) = \pi\ln(x+1)$.
\fincorrection
\correction{004343}
\begin{enumerate}
  \item 
  \item $I'(x) = \frac\pi{x+1}$, $I(x) = \pi\ln\left(\frac{x+1}2\right)$.
\end{enumerate}
 \fincorrection
\correction{004344}
\begin{enumerate}
  \item $g'(x) =  \int_{t=0}^1 (-2x)e^{-x^2(1+t^2)}\,d t = -2e^{-x^2}\sqrt{f(x)} = -f'(x)$.
  \item 
  \item 
\end{enumerate}
 \fincorrection
\correction{004345}
$u=\frac at  \Rightarrow  I = -\frac12\frac{d I}{d a}  \Rightarrow  I = \frac{\sqrt\pi}2e^{-2a}$.
\fincorrection
\correction{004346}
\begin{enumerate}
  \item 
  \item $I'(x) = -2xI(x)$.
  \item $I(x) = \frac{\sqrt\pi}2 e^{-x^2}$.
\end{enumerate}
 \fincorrection
\nocorrection
\nocorrection
\nocorrection
\correction{004350}
\begin{enumerate}
  \item $g(x,y) =  \int_{u=1}^y f(ux)\,d u$.
  \item $\frac{\partial g}{\partial x} =
             \frac1x \Bigl( yf(xy)-f(x) -  \int_{u=1}^y f(ux)\,d u\Bigr)
             \to \frac{y^2-1}2f'(0)$ lorsque $x\to0$.
\end{enumerate}
 \fincorrection
\correction{004351}
\begin{enumerate}
  \item $-2$ de $-\infty$ à $-\frac\pi2$,\quad
             $2\sin x$ de $-\frac\pi2$ à $\frac\pi2$,\quad
             $2$ de $\frac\pi2$ à $+\infty$.
  \item 
  \item 
  \item 
  \item 
  \item
\end{enumerate}
 \fincorrection
\correction{004352}
$2.40 < x < 2.41$.
\fincorrection
\correction{004353}
\begin{enumerate}
  \item 
  \item $a=\alpha$, $b=\alpha^2$.
  \item comparaison série-intégrale $ \Rightarrow  I(\alpha)\to\frac\pi2$ lorsque $\alpha\to+\infty$.
\end{enumerate}
 \fincorrection
\correction{004354}
$\ln\Gamma$ est convexe, encadrer $\ln\Gamma(x)$ par les cordes
passant par $(\lfloor x\rfloor,\ln\Gamma(\lfloor x\rfloor))$.
\fincorrection
\nocorrection
\correction{004356}
\begin{enumerate}
  \item $f'(a) = -\frac a2f(a)  \Rightarrow  f(a) = \frac{\sqrt\pi}2\exp(-a^2/4)$.
  \item $g'(a) = f(a)  \Rightarrow  g(a) \to \frac\pi2$ lorsque $a\to+\infty$.
\end{enumerate}
 \fincorrection
\nocorrection
\correction{004358}
\begin{enumerate}
  \item IPP.
  \item Soit $F(x) =  \int_{t=x}^{+\infty}f(t)\,d t$.
    On a $\varphi(a) = F(0) -a \int_{t=0}^{+\infty} e^{-at}F(t)\,d t$.

    Soit $\varepsilon > 0$ et $A$ tel que $x>A  \Rightarrow  |F(x)| \le \varepsilon$.
    On a~:
    $$\left|a \int_{t=0}^{+\infty} e^{-at}F(t)\,d t\right| \le
      a\sup{|F(t)|, t\in{[0,A]}} + \varepsilon e^{-aA} \le 2\varepsilon$$
    pour $a$ suffisament petit.
\end{enumerate}
 \fincorrection
\correction{004359}
$v_n\to 1$ (lorsque $n\to\infty$) par convergence dominée.
$v_n-1 =  \int_{x=0}^1\frac{x^n}{1+x^n}\, d x = \frac1n \int_{u=0}^1\frac{u^{-1/n}}{1+u}\,d u
\sim \frac{\ln 2}n$ donc la série diverge.
\fincorrection
\correction{004360}
$u_n \to 1$ (lorsque $n\to\infty$) par convergence dominée.

$u_n-1 =  \int_{x=0}^1\Bigl(\frac{1+x^{n-2}}{1+x^n}-1\Bigr)\,d x
       = \frac1n \int_{u=0}^1\frac{u^{1-1/n}}{1+u}(u^{-2/n}-1)\,d u$.

On a $0\le u^{-2/n}-1=\exp\Bigl(-\frac{2\ln(u)}n\Bigr)-1\le -\frac{2\ln(u)}n\exp\Bigl(-\frac{2\ln(u)}n\Bigr)$

d'où $0\le u_n\le\frac2{n^2} \int_{u=0}^1\frac{u^{1-3/n}(-\ln u)}{1+u}\,d u
= O\Bigl(\frac1{n^2}\Bigr)$.
\fincorrection
\correction{004361}
$I(\alpha)$ est définie pour tout~$\alpha>1$.
\smallskip
$I(2) = (x=e^u) =  \int_{u=-\infty}^{+\infty} \frac{ue^{2u}}{(1+e^{2u})^2}\,d u
      =  \int_{u=-\infty}^{+\infty} \frac{u}{(e^u+e^{-u})^2}\,d u
      = 0$ (parité).
\smallskip
$I(3) =  \int_{u=-\infty}^{+\infty} \frac{ue^{-u}}{(e^u+e^{-u})^3}\,d u
      =  \int_{u=0}^{+\infty} \frac{-u(e^u-e^{-u})}{(e^u+e^{-u})^3}\,d u
      = \Bigl[\frac{u}{2(e^u+e^{-u})^2}\Bigr]_{u=0}^{+\infty}
      -  \int_{u=0}^{+\infty} \frac{d u}{2(e^u+e^{-u})^2}$
\smallskip
$\phantom{I(3)}
      = - \int_{u=0}^{+\infty} \frac{e^{2u}\,d u}{2(1+e^{2u})^2}
      = \Bigl[\frac{1}{4(1+e^{2u})}\Bigr]_{u=0}^{+\infty}
      = -\frac18$.
\smallskip
$I(\alpha)\to 0$ (lorsque $\alpha\to+\infty$) par convergence dominée.
\fincorrection
\correction{004362}
\begin{enumerate}
  \item $D_f = {]0,1[}$. $f$ est convexe sur~$]0,1[$ par intégration de
    l'inégalité de convexité pour $x \mapsto t^{-x}$ et $f(x)\to +\infty$ (lorsque $x\to0$ ou $x\to 1$)
    par convergence monotone donc $f$ décroît puis recroît.
  \item 
  \item En~$0$~: $\frac1{\strut t^x(1+t)} = \frac1{\strut t^{x+1}} - \frac1{\strut t^{x+1}(1+t)}$
    donc $f(x) =   \int_{t=0}^1\frac{d t}{t^x(1+t)} + \frac1{\strut x}
    -  \int_{t=1}^{+\infty} \frac{d t}{\strut t^{x+1}(1+t)} \sim \frac1{\strut x}.$
    
    En~$1$~: $\frac1{\strut t^x(1+t)} = \frac1{\strut t^x} - \frac{t^{1-x}}{\strut 1+t}$
    donc $f(x) = \frac1{\strut 1-x} -  \int_{t=0}^1\frac{t^{1-x}}{\strut 1+t}\,d t
    +  \int_{t=1}^{+\infty}\frac{d t}{t^x(1+t)} \sim \frac1{\strut 1-x}$.
  \item $f(1/n) =_{(t=u^n)}  \int_{u=0}^{+\infty} \frac{nu^{n-1}\,d u}{u(1+u^n)}
    =_{(v=1/u)}  \int_{v=0}^{+\infty} \frac{n\,d v}{1+v^n} = \frac\pi{\strut\sin(\pi/n)}$
    (formule bien connue\dots)
\end{enumerate}
 \fincorrection
\correction{004363}
Si $\alpha\in2\pi\,\Z$ alors $u_n=0$ pour tout~$n$.

Sinon, $u_n = \Im\Bigl( \int_{t=0}^1 \sum_{k=1}^nt^{n+k-1}e^{ik\alpha}\,d t\Bigr)
=\Im\Bigl( \int_{t=0}^1\frac{t^ne^{i\alpha}-t^{2n}e^{i(n+1)\alpha}}{1-te^{i\alpha}}\,d t\Bigr)
\to 0$ (lorsque $n\to\infty$) par convergence dominée.
\fincorrection
\correction{004364}
$I_a =  \int_{x=-a}^a|f(x)|\,d x
\le  \int_{x=-a}^a \int_{t=-\infty}^{+\infty} \frac{e^{-|x+t|}}{\sqrt{|t|}(1+|t|)}\,d t\,d x
=  \int_{t=-\infty}^{+\infty} \int_{x=-a}^a \frac{e^{-|x+t|}}{\sqrt{|t|}(1+|t|)}\,d x\,d t$.
\smallskip
On a $ \int_{x=-a}^a e^{-|x+t|}\,d x = 
\begin{cases}2-2e^{-a}\ch t & \text{ si } |t|<a\cr
       2e^{-|t|}\sh a & \text{ si } |t|\ge a,\cr\end{cases}$
\smallskip
donc $I_a \le  \int_{t=0}^a \frac{2\,d t}{(1+t)\sqrt t} +
               \int_{t=a}^{+\infty}\frac{4e^{-t}\sh a}{(1+a)\sqrt a}\,d t
          = 4\Arctan(\sqrt a) + \frac{4e^{-a}\sh a}{(1+a)\sqrt a}
          \le \text{cste}$.
\fincorrection
\correction{004365}
\begin{enumerate}
  \item $s(x) =   \int_{t=0}^{+\infty} \sum_{k=1}^\infty\sin(t)e^{-kxt}\,d t$.

On a $|\sin(t)e^{-kxt}| \le te^{-kxt}$ et $ \int_{t=0}^{+\infty}te^{-kxt}\,d t = \frac 1{k^2}$
donc $\sum_{k=1}^{\infty} \int_{t=0}^{+\infty}|\sin(t)e^{-kxt}|\,d t$ converge
ce qui légitime l'interversion intégrale-série.
D'où $s(x) = \sum_{k=1}^\infty \int_{t=0}^{+\infty}\sin(t)e^{-kxt}\,d t
           = \sum_{k=1}^{\infty}\frac1{k^2x^2+1}$.
  \item Sachant (?) que $ \int_{t=0}^{+\infty}\frac{\sin t}t\,d t = \frac\pi2$, on obtient~:

\begin{align*}
xs(x)-\frac\pi2
&=  \int_{t=0}^{+\infty}\Bigl(\frac{x\sin t}{e^{xt}-1}-\frac{\sin t}{t}\Bigr)\,d t\cr
&=  \int_{u=0}^{+\infty}\Bigl(\frac{1}{e^{u}-1}-\frac{1}{u}\Bigr)\sin\Bigl(\frac ux\Bigr)\,d u\cr
&= -x\Bigr[\underbrace{\Bigl(\frac{1}{e^{u}-1}-\frac{1}{u}\Bigr)}_{\to{\textstyle\frac12} \text{ si } u\to0^+}\cos\Bigl(\frac ux\Bigr)\Bigr]_{u=0}^{+\infty}
   +x \int_{u=0}^{+\infty}\underbrace{\Bigl(\frac{-e^u}{(e^{u}-1)^2}+\frac{1}{u^2}\Bigr)}_{\to{\textstyle\frac1{12}\text{ si } u\to0^+ }}\cos\Bigl(\frac ux\Bigr)\,d u\cr
&= x(\text{quantité bornée})\to 0 \text{ si } u\to0^+.\cr
\end{align*}
\end{enumerate}
 \fincorrection
\correction{004366}
Fonction de $x$ de classe $\mathcal{C}^\infty$ sur $]0,+\infty[$, décroissante
de limite $\pi/2$ en $0^+$ et $0$ en $+\infty$. Demi-tangente verticale en~$0^+$,
\'Equivalente à $1/x$ en $+\infty$ (par IPP). \'Equation différentielle~:
$f(x) + f''(x) = 1/x$.
\fincorrection
\correction{004367}
\begin{enumerate}
  \item $[0,+\infty[$.
  \item 
  \item $f(x) - f'(x) =  \int_{t=0}^{+\infty}e^{-t^2x}\,d t = (u=t\sqrt x) = \frac{\sqrt\pi}{2\sqrt x}$.
  \item     $f'(x)
    = \frac{-1}{\sqrt x} \int_{u=0}^{+\infty}\frac{u^2e^{-u^2}}{x+u^2}\,d u
    = \frac{-1}{x\sqrt x} \int_{u=0}^{+\infty}\frac{u^2e^{-u^2}}{1+u^2/x}\,d u
    \sim\frac{-1}{x\sqrt x} \int_{u=0}^{+\infty}u^2e^{-u^2}\,d u
    =\frac{-\sqrt\pi}{4x\sqrt x}$.
  \item 
  \item
\end{enumerate}
 \fincorrection
\correction{004368}
\begin{enumerate}
  \item 
  \item Théorème de Fubini~:
   $ \int_{x=0}^{+\infty} f(x)\,d x =  \int_{\theta=0}^\pi \int_{x=0}^{+\infty}\Re(e^{(-\alpha+i\sin\theta)x})\,d x\, d\theta
   =  \int_{\theta=0}^\pi\frac{\alpha\, d\theta}{\alpha^2+\sin^2\theta}=\frac\pi{\sqrt{1+\alpha^2}}$

   (couper en $\theta=\pi/2$ et poser $u=\tan\theta$).
\end{enumerate}
 \fincorrection
\correction{004369}
$I'(x) =  \int_{t=0}^{+\infty}(\cos t-e^{-t})e^{-xt}\, d t
=\frac{x}{1+x^2}-\frac1{x+1}$
donc $I(x) = \ln\Bigl(\frac{1+x^2}{(1+x)^2}\Bigr) + \text{cste}$
et $I(x)\to0$ (pour $x\to+\infty$) d'où $\text{cste}=0$.
Alors $I(x)\to0$ pour $x\to0^+$.
\fincorrection
\nocorrection
\correction{004371}
$D_f = {]0,+\infty[}$. Il y a domination locale, donc $f$ est continue.

De même, pour $x>0$ on a $f'(x) =  \int_{t=0}^{+\infty}\frac{-\ln(t)t^{x+1}}{(t^{x+1}+t+1)^2}\,d t$.
En coupant l'intégrale en~$1$ et en posant $u=1/t$ dans l'intégrale sur $[1,+\infty[$
il vient~: $f'(x) =  \int_{t=0}^1\ln(t)t^{x+1}\Bigl(\frac1{(t+t^{x+1}+t^{x+2})^2} - \frac1{(t^{x+1}+t+1)^2}\Bigr)\,d t < 0$
car $\ln(t)<0$ et $t^{x+2}<1$ si $t\in{]0,1[}$. Donc $f$ est strictement décroissante sur~$]0,+\infty[$.

$ \int_{t=1}^{+\infty}\frac{d t }{t^{x+1}+t}
=  \int_{t=1}^{+\infty}\sum_{k=0}^\infty\frac{(-1)^k}{t^{(k+1)x+1}}\,d t =$
(domination du reste avec le CSA)
$= \sum_{k=0}^\infty\frac{(-1)^k}{(k+1)x}= \frac{\ln 2}x$.

$\Bigl|f(x)- \int_{t=1}^{+\infty}\frac{1}{t^{x+1}+t}\Bigr|
= \int_{t=0}^1\frac{d t}{t^{x+1}+t+1} +  \int_{t=1}^{+\infty}\frac{d t}{(t^{x+1}+t)(t^{x+1}+t+1)}
\le  \int_{t=0}^1\frac{d t}{t+1} +  \int_{t=1}^{+\infty}\frac{d t}{t(t+1)} = 2\ln 2$ 
donc $f(x) = \frac{\ln 2}x + O_{x\to0^+}(1)$.

Pour $x\to+\infty$, on a avec le TCM séparément sur $[0,1]$ et sur $[1,+\infty[$~:
lorsque $x\to+\infty$, $f(x)\to \int_{t=0}^1\frac{d t}{t+1} = \ln 2$.
\fincorrection
\correction{004372}
Pour $\lambda\ne 0$~:
$I_n = \Bigl[\frac{\exp(\lambda n\sin^2(x))}{2\lambda n\cos(x)}\Bigr]_{x=0}^\alpha
- \int_{x=0}^{\alpha}\frac{\sin(x)}{2\lambda n\cos^2(x)}\exp(\lambda n\sin^2(x))\,d x
= \frac{\exp(\lambda n\sin^2(\alpha))}{2\lambda n\cos(\alpha)} - \frac1{2\lambda n}
- \frac{J_n}{2\lambda n}$
avec $0\le J_n\le \frac{I_n}{\cos^2(\alpha)}$.
Donc $I_n\sim \frac{\exp(\lambda n\sin^2(\alpha))}{2\lambda n\cos(\alpha)}$ si $\lambda > 0$
et $I_n\sim - \frac1{2\lambda n}$ si $\lambda < 0$.
\fincorrection
\correction{004373}
\begin{enumerate}
  \item Pour $0\le t\le 1$ on a $t(1-t)(n-1)! \le t(1-t)\dots(n-t)\le n!$
    d'où $\frac1{6n}\le|a_n|\le1$ et $R=1$.
  \item $(-1)^na_n= \int_{t=0}^1t(1-t)(1-t/2)\dots(1-t/n)\,d t$.
    Pour $0\le x\le\frac12$ on a $x\le-\ln(1-x)\le x+x^2$ (étude
    de fonction) donc pour $k\ge 2$ et $0\le t\le 1$~:
    $e^{-t/k-t^2/k^2}\le 1-t/k\le e^{-t/k}$ d'où~:
 
    $$b_n =  \int_{t=0}^1t(1-t)e^{-t(H_n-1)-t^2K_n}\,d t\le (-1)^na_n
    \le  \int_{t=0}^1te^{-tH_n}\,d t = c_n$$
    avec $H_n = 1+1/2+\dots+1/n$ et $K_n = 1/2^2+\dots+1/n^2$.

    \'Equivalent du majorant~:
    $$c_n = \frac{1-(1+H_n)e^{-H_n}}{H_n^2}\sim\frac1{H_n^2}.$$

    \'Equivalent du minorant~:
    \begin{align*}b_n&\ge  \int_{t=0}^1t(1-t)(1-t^2K_n)e^{-t(H_n-1)}\,d t\cr
    &=  \int_{t=0}^1te^{-t(H_n-1)}\,d t -  \int_{t=0}^1t^2(1+t(1-t)K_n)e^{-t(H_n-1)}\,d t\cr
    &\ge  \int_{t=0}^1te^{-t(H_n-1)}\,d t - (1+{\textstyle\frac14}K_n) \int_{t=0}^1t^2e^{-t(H_n-1)}\,d t\cr
    &\ge \frac{1-H_ne^{1-H_n}}{(H_n-1)^2} - (1+{\textstyle\frac14}K_n)\frac{2-(H_n^2+1)e^{1-H_n}}{(H_n-1)^3}\cr
    &\sim\frac1{H_n^2}.\cr\end{align*}
    Finalement, $a_n\sim\frac{(-1)^n}{H_n^2}\sim\frac{(-1)^n}{\ln^2n}$.
\end{enumerate}
 \fincorrection
\correction{004374}
\begin{enumerate}
  \item $\sum u_n(t)$ converge pour $|t|<1$.
  \item $P_n(t) = t^{n+1}\sin(nx) - t^n\sin((n+1)x) + \sin(x)$,
$Q(t) = t^2 - 2t\cos(x) + 1 = (t-\cos x)^2 + \sin^2x\ge \sin^2x$.
  \item Pour $|t|<1$ on a $S_n(t)\to\frac{\sin t}{Q(t)}$ lorsque $n\to\infty$
et il y a convergence dominée vu la minoration de~$Q$ donc l'intégrale
suit~:
$ \int_{t=0}^1 S_n(t)\,d t\to \int_{t=0}^1 \frac{\sin x\,d t}{t^2-2t\cos x + 1}
= (t-\cos x = u\sin x) =  \int_{u=-\cot x}^{\tan(x/2)}\frac{d u}{u^2+1}
=  \frac{\pi-x}2$.
  \item $\frac{\sin(nx)}n =  \int_{t=0}^1 u_n(t)\,d t$ d'où
$\sum_{n=1}^\infty \frac{\sin(nx)}n = \frac{\pi-x}2$.
\end{enumerate}
 \fincorrection
\correction{004375}
\begin{enumerate}
  \item 
  \item $\sin^2(nt) = \frac{1-\cos(2nt)}2$, donc il suffit d'étudier
$I_n =  \int_{t=0}^{n\pi}\cos(2nt)f(t)\,d t$.

Posons $I_{n,p} =  \int_{t=0}^{\min(n,p)\pi}\cos(2nt)f(t)\,d t$~:
on a $|I_n-I_{n,p}| \le
 \int_{t=p\pi}^{+\infty}|f(t)|\,d t$, quantité indépendante de~$n$ et
tendant vers $0$ quand $p\to\infty$ donc le théorème d'interversion
des limites s'applique~:

$\lim_{n\to\infty}I_n
= \lim_{n\to\infty}\lim_{p\to\infty}I_{n,p}
= \lim_{p\to\infty}\lim_{n\to\infty}I_{n,p} = 0$.
On en déduit $u_n \to \frac12 \int_{t=0}^{+\infty} f(t)\,d t$ lorsque $n\to\infty$.
\end{enumerate}
 \fincorrection
\correction{004376}
\begin{enumerate}
  \item $0\le f_n(x)\le x^n$ et $f_n(1) = 1$ donc  lorsque $n\to\infty$ $f_n(x)\to \begin{cases}0 & \text{ si }  x<1 \cr 1 & \text{ si }x=1.\cr\end{cases}$
  \item
  \begin{enumerate}  
    \item   Non, la continuité n'est pas conservée.
    \item  Oui, il y a décroissance évidente.
  \end{enumerate}
  \item Changement de variable $u=\Bigl(\frac{1+x^{n-1}}2\Bigr)^n$~:
    $J_n=\frac2{n(n-1)} \int_{u=1/2^n}^1 (2u^{1/n}-1)^{2/(n-1)}u^{1/n}\,d u$
    et l'intégrale tend vers~$1$ quand $n\to\infty$ par convergence dominée.
\end{enumerate}
 \fincorrection
\correction{004377}
Il y a convergence si et seulement si $x>-1$.
$f'(x) =  \int_{t=0}^1 (1-t)t^x\,d t = \frac1{x+1} - \frac1{x+2}$,
donc $f(x) = \ln\Bigl(\frac{x+1}{x+2}\Bigr) + C$ et $f(x)\to 0$ (lorsque $x\to+\infty$) d'où $C=0$.
\fincorrection
\correction{004378}
\begin{enumerate}
  \item $]1,+\infty[$.
  \item 
  \item $I'(a) = - \int_{x=0}^{+\infty}\sh xe^{-ax}\,d x = \frac12\Bigl(\frac1{a-1}-\frac1{a+1}\Bigr)$.
D'où $I(a) = \frac12\ln\Bigl(\frac{a-1}{a+1}\Bigr) +$ cste et $I(a)\to 0$ lorsque $a\to+\infty$
donc la constante est nulle.
\end{enumerate}
 \fincorrection
\correction{004379}
\begin{enumerate}
  \item Intégrer par parties.
  \item Intégrer deux fois par parties.
  \item Pour $0<u<v$~: $ \int_{t=u}^{v}e^{-it^2/2}\,d t =
    \Bigl[\frac{e^{-it^2/2}}{-it}\Bigr]_{t=u}^v -
     \int_{t=u}^v\frac{e^{-it^2/2}}{it^2}\,d t
    \to \frac{e^{-iu^2/2}}{iu}- \int_{t=u}^{+\infty}\frac{e^{-it^2/2}}{it^2}\,d t$
     lorsque $v\to+\infty$.

    Ainsi $ \int_{t=u}^{+\infty}e^{-it^2/2}\,d t$ converge et de même pour
    $ \int_{t=-\infty}^{-u}e^{-it^2/2}\,d t$.
  \item On pose $f(t) = f(0) + t\varphi(t)$ avec $\varphi$ de
    classe $\mathcal{C}^1$. Il vient~:
    \begin{align*}g(x)\sqrt x &= f(0) \int_{u=a\sqrt x}^{b\sqrt x}e^{-iu^2/2}\,d u
      -\frac1{i\sqrt x}\Bigl[e^{-iu^2/2}\varphi(u/\sqrt x)\Bigr]_{u=a\sqrt
    x}^{b\sqrt x} + \frac1{ix} \int_{u=a\sqrt x}^{b\sqrt
    x}e^{-iu^2/2}\varphi'(u/\sqrt x)\,d u\cr
    &\to f(0).I\cr\end{align*}
    lorsque $x\to+\infty$.
\end{enumerate}
 \fincorrection
\nocorrection
\correction{005453}
Pour $x$ réel donné, la fonction $t\mapsto|t-x|f(t)$ est continue sur $[a,b]$ et donc $F(x)$ existe. Pour $x\leq a$,
$F(x)=\int_{a}^{b}(t-x)f(t)\;dt=-x\int_{a}^{b}f(t)\;dt+\int_{a}^{b}tf(t)\;dt$. $F$ est donc de classe $C^1$ sur $]-\infty,a]$ en tant que fonction affine et, pour $x<a$, $F'(x)=-\int_{a}^{b}f(t)\;dt$ (en particulier $F'_g(a)=-\int_{a}^{b}f(t)\;dt$)..

De même, pour $x\geq b$, $F(x)=x\int_{a}^{b}f(t)dt-\int_{a}^{b}tf(t)\;dt$. $F$ est donc de classe $C^1$ sur $[b,+\infty[$ en tant que fonction affine et, pour $x\geq b$, $F'(x)=\int_{a}^{b}f(t)\;dt$ (en particulier $F'_d(b)=\int_{a}^{b}f(t)\;dt$).

Enfin, si $a\leq x\leq b$, 

$$F(x)=\int_{a}^{x}(x-t)f(t)\;dt+\int_{x}^{b}(t-x)f(t)\;dt=x(\int_{a}^{x}f(t)\;dt-\int_{x}^{b}f(t)\;dt)-\int_{a}^{x}tf(t)\;dt+\int_{x}^{b}tf(t)\;dt.$$

$F$ est donc de classe $C^1$ sur $[a,b]$ et, pour $a\leq x\leq b$, 

\begin{align*}\ensuremath
F'(x)&=\int_{a}^{x}f(t)\;dt-\int_{x}^{b}f(t)\;dt+x(f(x)-(-f(x)))-xf(x)-xf(x)\\
 &=\int_{a}^{x}f(t)dt-\int_{x}^{b}f(t)\;dt.
\end{align*}

(et en particulier, $F'_d(a)=-\int_{a}^{b}f(t)\;dt=F'_g(a)$ et $F'_g(b)=\int_{a}^{b}f(t)\;dt=F'_d(b)$).

$F$ est continue $]-\infty,a]$, $[a,b]$ et $[b,+\infty[$ et donc sur $\Rr$. $F$ est de classe $C^1$ sur $]-\infty,a]$, $[a,b]$ et $[b,+\infty[$. De plus, $F'_g(a)=F'_d(a)$ et $F'_g(b)=F'_d(b)$. $F$ est donc de classe $C^1$ sur $\Rr$.
\fincorrection
\correction{005460}
Pour $t\in\Rr$, posons $g(t)=\frac{1}{\sqrt{t^4+t^2+1}}$. $g$ est continue sur $\Rr$ et admet donc des primitives sur $\Rr$. Soit $G$ une primitive de $g$ sur $\Rr$.

\textbf{Définition, dérivabilté, dérivée.}

Puisque $g$ est continue sur $\Rr$, $F$ est définie sur $\Rr$ et pour tout réel $x$, $F(x)=G(2x)-G(x)$. $G$ est de classe $C^1$ sur $\Rr$ et donc $F$ est de classe $C^1$ sur $\Rr$ et pour tout réel $x$,

$$F'(x)=2G'(2x)-G'(x)=2g(2x)-g(x)=\frac{2}{\sqrt{16x^4-4x^2+1}}-\frac{1}{\sqrt{x^4+x^2+1}}.$$

\textbf{Parité.}

Soit $x\in\Rr$. En posant $t=-u$ et donc $dt=-du$, on obtient, en notant que $g$ est paire

$$F(-x)=\int_{-x}^{-2x}g(t)\;dt=\int_{x}^{2x}g(-u).-du=-\int_{x}^{2x}g(u)\;du=-F(x).$$

$F$ est donc impaire.

\textbf{Variations.}

Pour $x$ réel,

\begin{align*}\ensuremath
\mbox{sgn}(F'(x))&=\mbox{sgn}(\frac{2}{\sqrt{16x^4-4x^2+1}}-\frac{1}{\sqrt{x^4+x^2+1}})=\mbox{sgn}(2\sqrt{x^4+x^2+1}-\sqrt{16x^4-4x^2+1})\\
 &=\mbox{sgn}(4(x^4+x^2+1)-(16x^4+4x^2+1))\;(\mbox{par croissance de}\;t\mapsto t^2\;\mbox{sur}\;\Rr^+)\\
  &=\mbox{sgn}(-12x^4+3)=\mbox{sgn}(1-4x^4)=\mbox{sgn}(1-2x^2).
\end{align*}

$F$ est donc strictement croissante sur $[-\frac{1}{\sqrt{2}},\frac{1}{\sqrt{2}}]$ et strictement décroissante sur $]-\infty,-\frac{1}{\sqrt{2}}]$ et sur $[\frac{1}{\sqrt{2}},+\infty[$.

\textbf{Etude en} $\bf{+{\infty}}$.

Pour $x>0$, $0\leq F(x)\leq\int_{x}^{2x}\frac{1}{\sqrt{x^4}}\;dt=\frac{2x-x}{x^2}=\frac{1}{x}$. Comme $\frac{1}{x}$ tend vers $0$ quand $x$ tend vers $+\infty$, le théorème des gendarmes permet d'affirmer que $\lim_{x\rightarrow +\infty}F(x)=0$.

\textbf{Graphe.}

$$\includegraphics{../images/img005460-1}$$

\fincorrection
\correction{005464}
\begin{enumerate}
\item  Si $x>1$, $[x,x^2]\subset]1,+\infty[$ et $t\mapsto\frac{1}{\ln t}$ est continue sur $]1,+\infty[$. Par suite, $\int_{x}^{x^2}\frac{dt}{\ln t}$ existe. De plus,

$$x\int_{x}^{x^2}\frac{1}{t\ln t}\;dt\leq\int_{x}^{x^2}\leq\int_{x}^{x^2}\frac{dt}{\ln t}=\int_{x}^{x^2}t\frac{1}{t\ln t}\;dt\leq x^2\int_{x}^{x^2}\frac{1}{t\ln t}\;dt.$$

Mais, 

$$\int_{x}^{x^2}\frac{1}{t\ln t}\;dt=\left[\ln|\ln t|\right]_{x}^{x^2}=\ln|\ln(x^2)|-\ln|\ln(x)|=\ln\left|\frac{2\ln x}{\ln x}\right|=\ln2.$$

Donc, $\forall x>1,\;x\ln2\leq F(x)\leq x^2\ln2$. On en déduit que $\lim_{x\rightarrow 1,\;x>1}F(x)=\ln2$.

Si $0<x<1$, on a $x^2<x$ puis $[x^2,x]\subset]0,1[$. Donc, $t\mapsto\frac{1}{\ln t}$ est continue sur $[x^2,x]$ et $F(x)=-\int_{x^2}^{x}\frac{1}{\ln t}\;dt$ existe.

Pour $t\in[x^2,x]$, on a $t\ln t<0$ et $x^2\leq t\leq x$. Par suite, 

$$x\frac{1}{t\ln t}\leq t\frac{1}{t\ln t}=\frac{1}{\ln t}\leq x^2\frac{1}{t\ln t},$$

puis, $\int_{x^2}^{x}x\frac{1}{t\ln t}\;dt\leq\int_{x^2}^{x}\frac{1}{\ln t}\;dt\leq\int_{x^2}^{x}x^2\frac{1}{t\ln t}\;dt$, et finalement,

$$x^2\ln2=\int_{x}^{x^2}x^2\frac{1}{t\ln t}\;dt\leq F(x)=\int_{x}^{x^2}\frac{1}{\ln t}\;dt\leq\int_{x}^{x^2}x\frac{1}{t\ln t}\;dt=x\ln 2.$$

On obtient alors $\lim_{x\rightarrow 1,\;x<1}F(x)=\ln2$ et finalement, $\lim_{x\rightarrow 1}F(x)=\ln2$. On en déduit que $F$ se prolonge par continuité en $1$ en posant $F(1)=\ln 2$ (on note encore $F$ le prolongement obtenu).

\item  On a déjà vu que $F$ est définie (au moins) sur $]0,+\infty[$ ($F$ désignant le prolongement). Il ne parait pas encore possible de donner un sens à $F(0)$ et encore moins à $F(x)$ quand $x<0$, car alors $[x,0]$ est un intervalle de longueur non nulle contenu dans $[x,x^2]$, sur lequel la fonction $t\mapsto\frac{1}{\ln t}$ n'est même pas définie.

$$D_F=]0,+\infty[.$$

Pour $t\in]0,1[\cup]1,+\infty[$, posons $g(t)=\frac{1}{\ln t}$ et notons $G$ une primitive de $g$ sur cet ensemble. Alors, pour $x$ dans $]0,1[\cup]1,+\infty[$, $F(x)=G(x^2)-G(x)$. On en déduit que $F$ est dérivable (et même de classe $C^1$) sur $]0,1[\cup]1,+\infty[$ et que pour $x$ dans $]0,1[\cup]1,+\infty[$,

$$F'(x)=2xg(x^2)-g(x)=\frac{2x}{\ln(x^2)}-\frac{1}{\ln x}=\frac{x-1}{\ln x}.$$

Maintenant, quand $x$ tend vers $1$, $\frac{x-1}{\ln x}$ tend vers $1$. Ainsi, $F$ est continue sur $]0,+\infty[$, de classe $C^1$ sur $]0,1[\cup]1,+\infty[$ et $F'$ a une limite réelle en $1$. Un théorème classique d'analyse permet d'affirmer que  $F$ est de classe $C^1$ sur $D_F$ et en particulier, dérivable en $1$ avec $F'(1)=1$.

$$\forall x\in]0,+\infty[,\;F'(x)=\left\{
\begin{array}{l}
\frac{x-1}{\ln x}\;\mbox{si}\;x\neq1\\
1\;\mbox{si}\;x=1.
\end{array}
\right..$$

Si $x>1$, $x-1>0$ et $\ln x>0$ et si $0<x<1$, $x-1<0$ et $\ln x<0$. Dans tous les cas ($0<x<1$, $x=1$, $x>1$)  $F'(x)>0$. $F$ est strictement croissante sur $]0,+\infty[$.

On a vu que $\forall x>1,\;F(x)>x\ln2$ et donc $\lim_{x\rightarrow +\infty}F(x)=+\infty$. Plus précisément, pour $x>1$,

$$\frac{F(x)}{x}=\frac{1}{x}\int_{x}^{x^2}\frac{1}{\ln t}\;dt\geq\frac{x^2-x}{x\ln x}=\frac{x-1}{\ln x}.$$

Comme $\frac{x-1}{\ln x}$ tend vers $+\infty$ quand $x$ tend vers $+\infty$, on en déduit que $\frac{F(x)}{x}$ tend vers $+\infty$ quand $x$ tend vers $+\infty$ et donc que la courbe représentative de $F$ admet en $+\infty$ une branche parabolique de direction $(Oy)$.

Pour $x\in]0,1[$ et $t\in[x^2,x]$, on a $2\ln x=\ln(x^2)\leq\ln t\leq\ln x<0$ et donc $\frac{1}{\ln x}\leq\frac{1}{\ln t}\leq\frac{1}{2\ln x}$, puis $(x-x^2)\frac{1}{\ln x}\leq\int_{x^2}^{x}\frac{1}{\ln t}\;dt\leq(x-x^2)\frac{1}{2\ln x}$ et finalement,

$$\forall x\in]0,1[,\;\frac{x-x^2}{-2\ln x}\leq F(x)\leq\frac{x-x^2}{-\ln x}.$$

On obtient déjà $\lim_{x\rightarrow 0}F(x)=0$. On peut prolonger $F$ par continuité en $0$ en posant $F(0)=0$. Ensuite, $\frac{F(x)-F(0)}{x-0}=\frac{F(x)}{x}$ est compris entre $\frac{1-x}{-2\ln x}$ et $\frac{1-x}{-\ln x}$. Comme ces deux expressions tendent vers $0$ quand $x$ tend vers $0$, on en déduit que $\frac{F(x)-F(0)}{x-0}$ tend vers $0$ quand $x$ tend vers $0$. $F$ est donc dérivable en $0$ et $F'(0)=0$.

$$\includegraphics{../images/img005464-1}$$

\end{enumerate}

\fincorrection
\correction{005472}
Notons $D$ le domaine de définition de $f$.

Si $x\in D$, $-x\in D$ et $f(-x)=-f(x)$. $f$ est donc impaire.

Si $x\in D$, $x+2\pi\in D$ et $f(x+2\pi)=f(x)$. $f$ est donc $2\pi$-périodique.

On étudiera donc $f$ sur $[0,\pi]$.

Soient $x\in[0,\pi]$ et $t\in[-1,1]$. $t^2-2t\cos x+1=(t-\cos x)^2+\sin^x\geq0$ avec égalité si et seulement si $\sin x=0$ et $t-\cos x=0$.

Ainsi, si $x\in]0,\pi[$, $\forall t\in]-1,1[,\;t^2-2t\cos x+1\neq0$. On en déduit que la fraction rationnelle $t\mapsto
\frac{\sin t}{1-2t\cos x+t^2}$ est continue sur $[-1,1]$, et donc que $f(x)$ existe.

Si $x=0$, $\forall t\in[-1,1[$, $\frac{\sin x}{t^2-2t\cos x+1}=\frac{0}{(t-1)^2}=0$. On prut prolonger cette fonction par continuité en $1$ et consisérer que $f(0)=\int_{-1}^{1}0\;dt=0$. De même, on peut considérer que $f(\pi)=0$.

Ainsi, $f$ est définie sur $[0,\pi]$ et donc, par parité et $2\pi$-périodicité, sur $\Rr$.

Soit $x\in]0,\pi[$.Calculons $f(x)$.

\begin{align*}\ensuremath
f(x)&=\int_{-1}^{1}\frac{\sin x}{(t-\cos x)^2+\sin^2x}\;dt=\left[\Arctan\frac{t-\cos x}{\sin x}\right]_{-1}^{1}
=\Arctan\frac{1-\cos x}{\sin x}+\Arctan\frac{1+\cos x}{\sin x}\\
 &=\Arctan\frac{2\sin^2(x/2)}{2\sin(x/2)\cos(x/2)}+\Arctan\frac{2\cos^2(x/2)}{2\sin(x/2)\cos(x/2)}
=\Arctan(\tan(x/2))+\Arctan(\frac{1}{\tan(x/2)})\\
 &=\frac{\pi}{2}\;(\mbox{car}\;\tan(x/2)>0\;\mbox{pour}\;x\in]0,\pi[).
\end{align*}

Ce calcul achève l'étude de $f$. En voici le graphe~:

$$\includegraphics{../images/img005472-1}$$

\fincorrection
\correction{005473}
Soit $x\in\Rr$. La fonction $t\mapsto\mbox{Max}(x,t)=\frac{1}{2}(x+t+|x-t|)$ est continue sur $[0,1]$ en vertu de théorèmes généraux. Par suite, $\int_{0}^{1}\mbox{Max}(x,t)\;dt$ existe. 

Si $x\leq0$, alors $\forall t\in[0,1]$, $x\leq t$ et donc $\mbox{Max}(x,t)=t$. Par suite, $f(x)=\int_{0}^{1}t\;dt=\frac{1}{2}$.

Si $x\geq1$, alors $\forall t\in[0,1]$, $t\leq x$ et donc $\mbox{Max}(x,t)=x$. Par suite, $f(x)=\int_{0}^{1}x\;dt=x$.

Si $0<x<1$,

$$f(x)=\int_{0}^{x}x\;dt+\int_{x}^{1}t\;dt=x^2+\frac{1}{2}(1-x^2)=\frac{1}{2}(1+x^2).$$

En résumé, $\forall x\in\Rr,\;f(x)=\left\{
\begin{array}{l}
\frac{1}{2}\;\mbox{si}\;x\leq0\\
\frac{1}{2}(1+x^2)\;\mbox{si}\;0<x<1\\
x\;\mbox{si}\;x\geq1
\end{array}
\right.$.

$f$ est déjà continue sur $]-\infty,0]$, $[1,+\infty[$ et $]0,1[$. De plus, $f(0^+)=\frac{1}{2}=f(0)$ et $f(1^-)=1=f(1)$. $f$ est ainsi continue à droite en $0$ et continue à gauche en $1$ et donc sur $\Rr$.

$f$ est de classe $C^1$ sur $]-\infty,0]$, $[1,+\infty[$ et $]0,1[$. De plus, $\lim_{x\rightarrow 0,\;x>0}f'(x)=\lim_{x\rightarrow 0,\;x>0}x=0$. $f$ est donc continue sur $[0,1[$ de classe $C^1$ sur $]0,1[$ et $f'$ a une limite réelle quand $x$ tend vers $0$. D'après un théorème classique d'analyse, $f$ est de classe $C^1$ sur $[0,1[$ et en particulier, $f$ est dérivable à droite en $0$ et ${f'}_{d}(0)=0$. Comme d'autre part, $f$ est dérivable à gauche en $0$ et que ${f'}_{g}(0)=0={f'}_{d}(0)$, $f$ est dérivable en $0$ et $f'(0)=0$.

L'étude en $1$ montre que $f$ est dérivable en $1$ et que $f'(1)=1$. Le graphe de $f$ est le suivant~:

$$\includegraphics{../images/img005473-1}$$

\fincorrection
\correction{005765}
\begin{enumerate}
 \item  Soit $n\in\Nn^*$. Soient $a$ et $A$ deux réels tels que $0<a<A$. On considère $\begin{array}[t]{cccc}
F_n~:&[a,A]\times\Rr&\rightarrow&\Rr\\
 &(x,t)&\mapsto&\frac{1}{(t^2+x^2)^n}
\end{array}$.

\textbullet~Pour chaque $x$ de $[a,A]$, la fonction $t\mapsto F_n(x,t)$ est continue par morceaux et intégrable sur $[0,+\infty[$ car $F_n(x,t)\underset{t\rightarrow+\infty}{\sim}\frac{1}{t^{2n}}>0$ avec $2n>1$.

\textbullet~La fonction $F_n$ est admet sur $[a,A]\times[0,+\infty[$ une dérivée partielle par rapport à sa première variable $x$ définie par :

\begin{center}
$\forall(x,t)\in[a,A]\times[0,+\infty[$, $\frac{\partial F_n}{\partial x}(x,t)=\frac{-2nx}{(t^2+x^2)^{n+1}}$.
\end{center}

De plus,

- pour chaque $x\in[a,A]$, la fonction $t\mapsto\frac{\partial F_n}{\partial x}(x,t)$ est continue par morceaux sur $[0,+\infty[$,\rule[-4mm]{0mm}{0mm}

-pour chaque $t\in[0,+\infty[$, la fonction $x\mapsto\frac{\partial F_n}{\partial x}(x,t)$ est continue sur $[a,A]$,

-pour chaque $(x,t)\in[a,A]\times[0,+\infty[$,

\begin{center}
$\left|\frac{\partial F_n}{\partial x}(x,t)\right|=\frac{2nx}{(t^2+x^2)^{n+1}}\leqslant\frac{2nA}{(t^2+a^2)^{n+1}}=\varphi(t)$,
\end{center}

où la fonction $\varphi$ est continue par morceaux et intégrable sur $[0,+\infty[$ car négligeable devant $\frac{1}{t^2}$ quand $t$ tend vers $+\infty$.

D'après le théorème de dérivation des intégrales à paramètres (théorème de \textsc{Leibniz}), la fonction $I_n$ est de classe $C^1$ sur $[a,A]$ et sa dérivée s'obtient par dérivation sous le signe somme. Ceci étant vrai pour tous réels $a$ et $A$ tels que $0<a<A$, on a montré que la fonction $I_n$ est de classe $C^1$ sur $]0,+\infty[$ et que

\begin{center}
$\forall x>0$, $I_n'(x)=-2nx\int_{0}^{+\infty}\frac{1}{(t^2+x^2)^{n+1}}\;dt=-2nxI_{n+1}(x)$.
\end{center}

\begin{center}
\shadowbox{
$\forall n\in\Nn^*$, $I_n'(x)=-2nxI_{n+1}(x)$.
}
\end{center}

\item  Pour $x>0$, on a $I_1(x)=\left[\frac{1}{x}\Arctan\left(\frac{t}{x}\right)\right]_0^{+\infty}=\frac{\pi}{2x}$. Ensuite, $I_2(x)=-\frac{1}{2x}I_1'(x)=\frac{\pi}{4x^3}$ puis $I_3(x)=-\frac{1}{4x}I_2'(x)=\frac{3\pi}{16x^5}$ et donc $I_3(1)=\frac{3\pi}{16}$.

\begin{center}
\shadowbox{
$\int_{0}^{+\infty}\frac{1}{(t^2+1)^3}\;dt=\frac{3\pi}{16}$.
}
\end{center}
\end{enumerate}
\fincorrection
\correction{005766}
\begin{enumerate}
 \item 
   \begin{enumerate}
   \item \textbf{Parité de $F$.} Soit $x$ un réel du domaine de définition de $F$. En posant $t=\theta+\pi$, on obtient

\begin{align*}\ensuremath
F(x)&=\int_{-\pi}^{\pi}\ln(x^2-2x\cos\theta+1)\;d\theta=\int_{0}^{2\pi}\ln(x^2+2x\cos t+1)\;dt\\
 &=\int_{-\pi}^{\pi}\ln(x^2+2x\cos t+1)\;dt\;(\text{par}\;2\pi\text{-périodicité})\\
 &=\int_{-\pi}^{\pi}\ln((-x)^2-2(-x)\cos t+1)\;dt=F(-x).
\end{align*}

Ainsi, pour tout réel $x$, $F(x)$ existe si et seulement si $F(-x)$ existe et de plus $F(x)=F(-x)$.

\begin{center}
\shadowbox{
$F$ est paire.
}
\end{center}

\textbf{Définition de $F$.} Soit $x\in[0,+\infty[$. Pour tout réel $\theta\in[-\pi,\pi]$,

\begin{center}
$x^2-2x\cos\theta+1=(x-\cos\theta)^2+(\sin\theta)^2=\left|x-e^{i\theta}\right|^2\geqslant0$.
\end{center}

De plus, $\left|x-e^{i\theta}\right|=0\Leftrightarrow e^{i\theta}=x\Leftrightarrow x=1\;\text{et}\;\theta=0$. Par suite,

\textbullet~si $x\neq1$, la fonction $\theta\mapsto x^2-2x\cos\theta+1$ est continue sur le segment $[0,\pi]$ et donc intégrable sur ce segment.

\textbullet~si $x=1$, pour tout réel $\theta\in[-\pi,\pi]$ on a $x^2-2x\cos\theta+1=2-2\cos\theta=4\sin^2\frac{\theta}{2}$. La fonction $\theta\mapsto\ln\left(4\sin^2\frac{\theta}{2}\right)$ 

est continue sur $[-\pi,0[\cup]0,\pi]$ et quand $\theta$ tend vers $0$

\begin{center}
$\ln\left(4\sin^2\frac{\theta}{2}\right)=2\ln2+2\ln\left|\sin\frac{\theta}{2}\right|\sim2\ln\left|\frac{\theta}{2}\right|\sim2\ln|\theta|=o\left(\frac{1}{\sqrt{|\theta|}}\right)$.
\end{center}

On en déduit que la fonction $\theta\mapsto\ln\left(4\sin^2\frac{\theta}{2}\right)$ est intégrable sur $[-\pi,\pi]$ et donc que $F(1)$ existe.

Finalement, $F$ est définie sur $[0,+\infty[$ et par parité

\begin{center}
\shadowbox{
$F$ est définie sur $\Rr$.
}
\end{center}

\text{Remarque.} Par parité de la fonction $\theta\mapsto\ln(x^2-2x\cos\theta+1)$, pour tout réel $x$, on a encore $F(x)=2\int_{0}^{\pi}\ln(x^2-2x\cos\theta+1)\;d\theta$.

\textbf{Continuité de $F$.} Soit $A>1$. Soit $\begin{array}[t]{cccc}
\Phi~:&[0,A]\times]0,\pi]&\rightarrow&\Rr\\
 &(x,\theta)&\mapsto&\ln(x^2-2x\cos\theta+1)
\end{array}$.

\textbullet~Pour chaque $x\in[0,A]$, la fonction $\theta\mapsto\Phi(x,\theta)$ est continue par morceaux sur $]0,\pi]$.

\textbullet~Pour chaque $\theta\in]0,\pi]$, la fonction $x\mapsto\Phi(x,\theta)$ est continue par morceaux sur $[0,A]$.

\textbullet~Pour chaque $(x,\theta)\in\Rr^+\times]0,\pi]$, puisque $x^2-2x\cos\theta+1=(x-\cos\theta)^2+(\sin\theta)^2$,

\begin{align*}\ensuremath
|\Phi(x,\theta)|&\leqslant\text{Max}\{\left|\ln(0^2-0\cos\theta+1)\right|,\left|\ln(\cos^2\theta-2\cos\theta\cos\theta+1)\right|,\left|\ln(A^2-2A\cos\theta+1)\right|\}\\
 &=\text{Max}\{2\left|\ln(|\sin\theta|)\right|,\left|\ln(A^2-2A\cos\theta+1)\right|\}=\varphi(\theta).
\end{align*}

On a vu que la fonction $f_1~:~\theta\mapsto2\left|\ln(|\sin\theta|)\right|$ est intégrable sur $]0,\pi]$ et d'autre part, la fonction $f_2~:~\theta\left|\ln(A^2-2A\cos\theta+1)\right|$ est intégrable sur $[0,\pi]$ et donc sur $]0,\pi]$ car continue sur $[0,\pi]$. Puisque $\varphi=\frac{1}{2}(f_1+f_2+|f_1-f_2|)$, on en déduit que la fonction $\varphi$ est continue par morceaux et intégrable sur $]0,\pi]$.

D'après le théorème de continuité des intégrales à paramètres, la fonction $F$ est continue sur $[0,A]$ et ceci pour tout $A>1$. Par suite, la fonction $F$ est continue sur $\Rr^+$ puis par parité,

\begin{center}
\shadowbox{
la fonction $F$ est continue sur $\Rr$.
}
\end{center}
 

\textbf{Dérivabilité de $F$.} Soient $A\in]0,1[$ puis  $\begin{array}[t]{cccc}
\Phi~:&[-A,A]\times[0,\pi]&\rightarrow&\Rr\\
 &(x,\theta)&\mapsto&\ln(x^2-2x\cos\theta+1)
\end{array}$.

\textbullet~Pour chaque $x\in[-A,A]$, la fonction $\theta\mapsto \Phi(x,\theta)$ est continue sur le segment $[0,\pi]$ et donc intégrable sur ce segment.

\textbullet~La fonction $\Phi$ admet sur $[-A,A]\times[0,\pi]$ une dérivée partielle par rapport à sa première variable $x$ définie par

\begin{center}
$\forall(x,\theta)\in[-A,A]\times[0,\pi]$, $\frac{\partial \Phi}{\partial x}(x,\theta)=\frac{2x-2\cos\theta}{x^2-2x\cos\theta+1}$.
\end{center}

De plus,

- pour chaque $x\in[-A,A]$, la fonction $\theta\mapsto\frac{\partial \Phi}{\partial x}(x,\theta)$ est continue par morceaux sur $[0,\pi]$,

- pour chaque $\theta\in[0,\pi]$, la fonction $x\mapsto\frac{\partial \Phi}{\partial x}(x,\theta)$ est continue sur $[-A,A]$,

- pour chaque $(x,\theta)\in[-A,A]\times[0,\pi]$,

\begin{center}
$\left|\frac{\partial \Phi}{\partial x}(x,\theta)\right|=\frac{2|x-\cos\theta|}{|x-e^{i\theta}|^2}\leqslant\frac{4}{|A-1|^2}=\varphi(\theta)$.
\end{center}

La dernière inégalité écrite est claire géométriquement :


$$\includegraphics{../images/img005766-1}$$


De plus, la fonction constante $\varphi$ est intégrable sur le segment $[0,\pi]$.

D'après le théorème de dérivation des intégrales à paramètres, la fonction $F$ est de classe $C^1$ sur $[-A,A]$ et sa dérivée s'obtient par dérivation sous le signe somme. Ceci étant vrai pour tout $A\in]0,1[$, $F$ est de classe $C^1$ sur $]-1,1[$ et $\forall x\in]-1,1[$, \rule[-4mm]{0mm}{0mm}$F'(x)=4\int_{0}^{\pi}\frac{x-\cos\theta}{x^2-2x\cos\theta+1}\;d\theta$. La démarche est analogue sur $]-\infty,-1[$ et sur $]1,+\infty[$ et finalement $F$ est de classe $C^1$ sur $\Rr\setminus\{-1,1\}$ et

\begin{center}
\shadowbox{
$\forall x\in\Rr\setminus\{-1,1\}$, $F'(x)=4\int_{0}^{\pi}\frac{x-\cos\theta}{x^2-2x\cos\theta+1}\;d\theta$.
}
\end{center}

   \item \textbf{Calcul de $F'(x)$.} Soit $x\in\Rr\setminus\{-1,1\}$. On pose $t=\tan\left(\frac{\theta}{2}\right)$. On a donc $\cos\theta=\frac{1-t^2}{1+t^2}$ et $d\theta=\frac{2dt}{1+t^2}$. On obtient

\begin{align*}\ensuremath
F'(x)&=4\int_{0}^{\pi}\frac{x-\cos\theta}{x^2-2x\cos\theta+1}\;d\theta=8\int_{0}^{+\infty}\frac{x-\frac{1-t^2}{1+t^2}}{x^2-2x\frac{1-t^2}{1+t^2}+1}\frac{dt}{1+t^2}\\
 &=8\int_{0}^{+\infty}\frac{(1+t^2)x-(1-t^2)}{((1+t^2)x^2-2x(1-t^2)+(1+t^2))(1+t^2)}\;dt\\
 &=8\int_{0}^{+\infty}\frac{(x+1)t^2+(x-1)}{((x+1)^2t^2+(x-1)^2)(1+t^2)}\;dt
\end{align*}

Pour tout réel $t$,

\begin{center}
$\left(t^2+\left(\frac{x-1}{x+1}\right)^2\right)(t^2+1)=\left(t-i\frac{x-1}{x+1}\right)\left(t+i\frac{x-1}{x+1}\right)(t-i)(t+i)$.
\end{center}

De plus, $\pm\frac{x-1}{x+1}=\pm1\Leftrightarrow\frac{x-1}{x+1}=-1\Leftrightarrow x=0$.

\textbullet~$F'(0)=4\int_{0}^{\pi}(-\cos\theta)\;d\theta=0$.

\textbullet~Si $x\neq0$, les pôles de la fraction rationnelle $\frac{(x+1)t^2+(x-1)}{((x+1)^2t^2+(x-1)^2)(1+t^2)}$ sont simples et par parité, la décomposition en éléments simples de cette fraction s'écrit

\begin{center}
$\frac{(x+1)t^2+(x-1)}{((x+1)^2t^2+(x-1)^2)(1+t^2)}=\frac{a}{t-i\frac{x-1}{x+1}}-\frac{a}{t+i\frac{x-1}{x+1}}+\frac{b}{t-i}-\frac{b}{t+i}$,
\end{center}

avec

\begin{align*}\ensuremath
a&=\frac{-(x+1)\left(\frac{x-1}{x+1}\right)^2+(x-1)}{(x+1)^2\left(2i\frac{x-1}{x+1}\right)\left(1-\left(\frac{x-1}{x+1}\right)^2\right)}=\frac{-(x+1)(x-1)^2+(x-1)(x+1)^2}{2i(x+1)(x-1)((x+1)^2-(x-1)^2)}\\
 &=\frac{2(x^2-1)}{2i(x^2-1)(4x)}=\frac{1}{4ix},
\end{align*}

et

\begin{center}
$b=\frac{-(x+1)+(x-1)}{2i(-(x+1)^2+(x-1)^2)}=\frac{1}{4ix}$.
\end{center}

Donc

\begin{align*}\ensuremath
8\frac{(x+1)t^2+(x-1)}{((x+1)^2t^2+(x-1)^2)(1+t^2)}&=\frac{2}{ix}\left(\frac{1}{t-i\frac{x-1}{x+1}}-\frac{1}{t+i\frac{x-1}{x+1}}+\frac{1}{t-i}-\frac{1}{t+i}\right)\\
 &=\frac{2}{ix}\left(\frac{2i\frac{x-1}{x+1}}{t^2+\left(\frac{x-1}{x+1}\right)^2}+\frac{2i}{t^2+1}\right)=\frac{4}{x}\left(\frac{x^2-1}{(x+1)^2t^2+(x-1)^2}+\frac{1}{t^2+1}\right)
\end{align*}

Ensuite, en notant $\varepsilon$ le signe de $\frac{x-1}{x+1}$

\begin{align*}\ensuremath
F'(x)&=\frac{4}{x}\int_{0}^{+\infty}\left(\frac{x^2-1}{(x+1)^2t^2+(x-1)^2}+\frac{1}{t^2+1}\right)\;dt\\
 &=\frac{4}{x}\left[\frac{x^2-1}{(x+1)^2}\frac{1}{\frac{x-1}{x+1}}\Arctan\left(\frac{t}{\frac{x-1}{x+1}}\right)\right]_0^{+\infty}=\frac{4}{x}\left(\varepsilon+1\right)\frac{\pi}{2}
\end{align*}

Par suite, si $x\in]-1,1[$, $F'(x)=0$ et si $x\in]-\infty,-1[\cup]1,+\infty[$, $F'(x)=\frac{4\pi}{x}$.

\begin{center}
\shadowbox{
$\forall x\in\Rr\setminus\{-1,1\}$, $F'(x)=\left\{
\begin{array}{l}
\rule[-2mm]{0mm}{0mm}0\;\text{si}\;x\in]-1,1[\\
\frac{4\pi}{x}\;\text{si}\;x\in]-\infty,-1[\cup]1,+\infty[
\end{array}
\right.$.
}
\end{center}

   \item Soit $x>1$.

\begin{center}
$F(x)-4\pi\ln(x)=\int_{-\pi}^{\pi}\ln(x^2-2x\cos\theta+1)\;d\theta-\int_{-\pi}^{\pi}\ln(x^2)\;d\theta=\int_{-\pi}^{\pi}\ln\left(1-\frac{2}{x}\cos\theta+\frac{1}{x^2}\right)\;d\theta=F\left(\frac{1}{x}\right)$.
\end{center}

Par suite, $\lim_{x \rightarrow +\infty}(F(x)-4\pi\ln(x))=\lim_{x \rightarrow +\infty}F\left(\frac{1}{x}\right)=\lim_{y \rightarrow 0}F(y)=F(0)=0$ par continuité de $F$ en $0$.

\begin{center}
\shadowbox{
$\lim_{x \rightarrow +\infty}(F(x)-4\pi\ln(x))=0$.
}
\end{center}

   \item \textbullet~$F$ est continue sur $[-1,1]$, dérivable sur $]-1,1[$ de dérivée nulle sur $]-1,1[$. Donc la fonction $F$ est constante sur l'intervalle $[-1,1]$. Par suite, pour tout réel $x\in[-1,1]$, $F(x)=F(0)=0$.

\textbullet~$F$ est dérivable sur $]1,+\infty[$ et $\forall x>1$, $F'(x)=\frac{4\pi}{x}$. Donc il existe $C\in\Rr$ tel que $\forall x>1$, $F(x)=4\pi\ln x+C$ avec $C=\lim_{x \rightarrow +\infty}(F(x)-4\pi\ln x)=0$. Donc $\forall x>1$, $F(x)=4\pi\ln x$.

\textbullet~Si $x<-1$, $F(x)=F(-x)=4\pi\ln(-x)=4\pi\ln|x|$.

\begin{center}
\shadowbox{
$\forall x\in\Rr$, $\int_{-\pi}^{\pi}\ln(x^2-2x\cos\theta+1)\;d\theta=\left\{
\begin{array}{l}
0\;\text{si}\;x\in[-1,1]\\
4\pi\ln(|x|)\;\text{si}\;x\in]-\infty,-1[\cup]1,+\infty[
\end{array}
\right.$.
}
\end{center}
  \end{enumerate}

\item  
  \begin{enumerate}
  \item Soit $x\in]-1,1[$. Pour $\theta\in[-\pi,\pi]$, on pose $f(\theta)=\ln(x^2-2x\cos\theta+1)$. Puisque $\forall\theta\in[-\pi,\pi]$, $x^2-2x\cos\theta+1>0$ (voir 1)), $f$ est dérivable sur $[-\pi,\pi]$ et pour $\theta\in[-\pi,\pi]$,

\begin{align*}\ensuremath
f'(\theta)&=\frac{2x\sin\theta}{x^2-2x\cos\theta+1}=\frac{1}{i}\left(\frac{e^{i\theta}}{x-e^{i\theta}}-\frac{e^{-i\theta}}{x-e^{-i\theta}}\right)=\frac{1}{i}\left(-\frac{1}{1-xe^{-i\theta}}+\frac{1}{1-xe^{i\theta}}\right)\\
 &=\frac{1}{i}\left(\sum_{n=0}^{+\infty}x^ne^{in\theta}-\sum_{n=0}^{+\infty}x^ne^{-in\theta}\right)\;(\text{car}\;|xe^{i\theta}|=|xe^{-i\theta}|=|x|<1)\\
 &=2\sum_{n=1}^{+\infty}\sin(n\theta)x^n.
\end{align*}

Soit $\theta\in[-\pi,\pi]$. $I$ désigne l'intervalle $[0,\theta]$ ou $[\theta,0]$ suivant que $\theta$ soit positif ou négatif.

Pour $n\in\Nn^*$ et $t\in I$, posons $g_n(t)=2\sin(nt)x^n$. Pour tout $n\in\Nn^*$ et tout $t\in I$, on a $|f_n(t)|\leqslant|x|^n$. Comme $|x|^n$ est le terme général d'une série numérique convergente, la série de fonctions de terme général $f_n$, $n\in\Nn^*$, converge normalement et donc uniformément sur le segment $I$. D'après le théorème d'intégration terme à terme sur un segment, on peut écrire

\begin{align*}\ensuremath
f(\theta)&=f(0)+\int_{0}^{\theta}f'(t)\;dt=2\ln(1-x)+\sum_{n=1}^{+\infty}2x^n\int_{0}^{\theta}\sin(nt)\;dt\\
 &=2\left(-\sum_{n=1}^{+\infty}\frac{x^n}{n}+\sum_{n=1}^{+\infty}(1-\cos(n\theta))\frac{x^n}{n}\right)=-2\sum_{n=1}^{+\infty}\frac{\cos(n\theta)}{n}x^n.
\end{align*}

\begin{center}
\shadowbox{
$\forall x\in]-1,1[$, $\forall\theta\in[-\pi,\pi]$, $\ln(x^2-2x\cos\theta+1)=-2\sum_{n=1}^{+\infty}\frac{\cos(n\theta)}{n}x^n$.
}
\end{center}

Soit $x\in]-1,1[$. Pour $n\in\Nn^*$ et $\theta\in[-\pi,\pi]$, $\left|\frac{\cos(n\theta)}{n}x^n\right|\leqslant|x|^n$. Comme précédemment, on peut intégrer terme à terme et on obtient

\begin{center}
$F(x)=-2\sum_{n=1}^{+\infty}\frac{x^n}{n}\int_{-\pi}^{\pi}\cos(n\theta)\;d\theta=0$.
\end{center}

\begin{center}
\shadowbox{
$\forall x\in]-1,1[$, $\int_{-\pi}^{\pi}\ln(x^2-2x\cos\theta+1)\;d\theta=0$.
}
\end{center}

\  \item Soit $x\in\Rr^*$.

\begin{align*}\ensuremath
F\left(\frac{1}{x}\right)&=\int_{-\pi}^{\pi}\ln\left(1-\frac{2}{x}\cos\theta+\frac{1}{x^2}\right)\;d\theta=\int_{-\pi}^{\pi}(\ln(1-2x\cos\theta+x^2)-\ln(x^2))\;d\theta=-4\pi\ln|x|+F(x).
\end{align*}

\begin{center}
\shadowbox{
$\forall x\in\Rr^*$, $F\left(\frac{1}{x}\right)=-4\pi\ln|x|+F(x)$.
}
\end{center}

Soit $x>1$. Puisque $\frac{1}{x}\in]0,1[$, $F(x)=4\pi\ln x+F\left(\frac{1}{x}\right)=4\pi\ln x$. On retrouve alors les résultats du 1).
  \end{enumerate}
\end{enumerate}
\fincorrection
\correction{005767}
\begin{enumerate}
 \item  Soit $A>0$. Soit $\begin{array}[t]{cccc}
\Phi~:&[-A,A]\times[0,1]&\rightarrow&\Rr\\
 &(x,t)&\mapsto&\frac{e^{-x^2(1+t^2)}}{1+t^2}
\end{array}$.

\textbullet~Pour chaque $x$ de $[-A,A]$, la fonction $t\mapsto F(x,t)$ est continue sur le segment $[0,1]$ et donc intégrable sur ce segment.

\textbullet~La fonction $\Phi$ admet sur $[-A,A]\times[0,1]$ une dérivée partielle par rapport à sa première variable $x$ définie par :

\begin{center}
$\forall(x,t)\in[-A,A]\times[0,1]$, $\frac{\partial \Phi}{\partial x}(x,t)=-2xe^{-x^2(1+t^2)}$.
\end{center}

De plus,

- pour chaque $x\in[-A,A]$, la fonction $t\mapsto \frac{\partial \Phi}{\partial x}(x,t)$ est continue par morceaux sur le segment $[0,1]$,

- pour chaque $t\in[0,1]$, la fonction $x\mapsto \frac{\partial \Phi}{\partial x}(x,t)$ est continue par morceaux sur $\Rr$,

- pour chaque $(x,t)\in[-A,A]\times[0,1]$, $\left|\frac{\partial \Phi}{\partial x}(x,t)\right|\leqslant 2A=\varphi(t)$, la fonction $\varphi$ étant continue et donc intégrable sur le

segment $[0,1]$.

D'après le théorème de dérivation des intégrales à paramètres (théorème de \textsc{Leibniz}), la fonction $F$ est de classe $C^1$ sur $[-A,A]$ et sa dérivée s'obtient en dérivant sous le signe somme. Ceci étant vrai pour tout $A>0$, $F$ est de classe $C^1$ sur $\Rr$ et

\begin{center}
\shadowbox{
$\forall x\in\Rr$, $F'(x)=-2x\int_{0}^{1}e^{-x^2(1+t^2)}\;dt$.
}
\end{center}

\item  La fonction $x\mapsto e^{-x^2}$ est continue sur $\Rr$. On en déduit que la fonction $x\mapsto\int_{0}^{x}e^{-t^2}\;dt$ est de classe $C^1$ sur $\Rr$. Il en est de même de la fonction $G$ et pour tout réel $x$,

\begin{center}
\shadowbox{
$G'(x)=2e^{-x^2}\int_{0}^{x}e^{-t^2}\;dt$.
}
\end{center}

\item  Soit $x\in\Rr^*$. En posant $u=xt$, on obtient

\begin{center}
$F'(x)=-2x\int_{0}^{1}e^{-x^2(1+t^2)}\;dt=-2e^{-x^2}\int_{0}^{1}e^{-(xt)^2}\;xdt=-e^{-x^2}\int_{0}^{x}e^{-u^2}\;du=-G'(x)$,
\end{center}

cette égalité restant vraie quand $x=0$ par continuité des fonctions $F'$ et $G'$ sur $\Rr$.

Ainsi, $F'+G'=0$ et donc $\forall x\in\Rr$, $F(x)+G(x)=F(0)+G(0)=\int_{0}^{1}\frac{1}{1+t^2}\;dt=\frac{\pi}{4}$.

\begin{center}
\shadowbox{
$\forall x\in\Rr$, $F(x)+G(x)=\frac{\pi}{4}$.
}
\end{center}

\item  Pour $x\in\Rr$,

\begin{center}
$|F(x)|=\int_{0}^{1}\frac{e^{-x^2(1+t^2)}}{1+t^2}\;dt\leqslant e^{-x^2}\int_{0}^{1}\frac{1}{1+t^2}\;dt=\frac{\pi}{4e^{x^2}}$,
\end{center}

et puisque $\lim_{x \rightarrow +\infty}\frac{\pi}{4e^{x^2}}=0$, on a montré que

\begin{center}
\shadowbox{
$\lim_{x \rightarrow +\infty}F(x)=0$.
}
\end{center}

\item  Pour $x>0$, on a $\int_{0}^{x}e^{-t^2}\;dt\geqslant0$ et donc d'après la question 3),

\begin{center}
$\int_{0}^{x}e^{-t^2}\;dt=\sqrt{G(x)}=\sqrt{\frac{\pi}{2}-F(x)}$.
\end{center}

La question 4) permet alors d'affirmer que $\lim_{x \rightarrow +\infty}G(x)=\frac{\sqrt{\pi}}{2}$ et donc que

\begin{center}
\shadowbox{
$\int_{0}^{+\infty}e^{-t^2}\;dt=\frac{\sqrt{\pi}}{2}$.
}
\end{center}
\end{enumerate}
\fincorrection
\correction{005768}
Soit $x\in\Rr$. La fonction $t\mapsto e^{-t^2}\ch(tx)$ est continue sur $[0,+\infty[$. Quand $t$ tend vers $+\infty$, $e^{-t^2}\ch(tx)=\frac{1}{2}(e^{-t^2+tx}+e^{-t^2-tx})=o\left(\frac{1}{t^2}\right)$ d'après un théorème de croissances comparées et donc la fonction $t\mapsto e^{-t^2}\ch(tx)$ est intégrable sur $[0,+\infty[$. Pour $x\in\Rr$, on peut poser $f(x)=\int_{0}^{+\infty}e^{-t^2}\ch(tx)\;dt$.

\textbf{Calcul de $f(x)$.} Soit $A>0$. On pose $\begin{array}[t]{cccc}
\Phi~:&[-A,A]\times[0,+\infty[&\rightarrow&\Rr\\
 &(x,t)&\mapsto&e^{-t^2}\ch(tx)
\end{array}$.

\textbullet~Pour chaque $x\in[-A,A]$, la fonction $t\mapsto e^{-t^2}\ch(tx)$ est continue par morceaux et intégrable sur $[0,+\infty[$.

\textbullet~La fonction $\Phi$ admet sur $[-A,A]\times[0,+\infty[$ une dérivée partielle par rapport à sa première variable définie par :

\begin{center}
$\forall(x,t)\in[-A,A]\times[0,+\infty[$, $\frac{\partial \Phi}{\partial x}(x,t)=te^{-t^2}\sh(tx)$.
\end{center}

De plus,

- pour chaque $x\in[-A,A]$, la fonction $t\mapsto\frac{\partial \Phi}{\partial x}(x,t)$ est continue par morceaux sur $[0,+\infty[$,

-pour chaque $t\in[0,+\infty[$, la fonction $x\mapsto\frac{\partial \Phi}{\partial x}(x,t)$ est continue sur $[-A,A]$,

- pour chaque $(x,t)\in[-A,A]\times[0,+\infty[$,

\begin{center}
$\left|\frac{\partial \Phi}{\partial x}(x,t)\right|=te^{-t^2}|\sh(tx)|\leqslant te^{-t^2}\sh(t|A|)=\varphi(t)$.
\end{center}

La fonction $\varphi$ est continue par morceaux sur $[0,+\infty[$ et intégrable sur $[0,+\infty[$ car négligeable devant $\frac{1}{t^2}$ quand $t$ tend vers $+\infty$.

D'après le théorème de dérivation des intégrales à paramètres (théorème de \textsc{Leibniz}), la fonction $f$ est de classe $C^1$ sur $[-A,A]$ et sa dérivée s'obtient par dérivation sous le signe somme. Ceci étant vrai pour tout réel $A>0$, la fonction $f$ est de classe $C^1$ sur $\Rr$ et

\begin{center}
$\forall x\in\Rr$, $f'(x)=\int_{0}^{+\infty}te^{-t^2}\sh(tx)\;dt$.
\end{center}

Soit $x\in\Rr$. On effectue maintenant une intégration par parties. Soit $A>0$. Les deux fonctions $t\mapsto te^{-t^2}$ et $t\mapsto\sh(tx)$ sont de classe $C^1$ sur le segment $[0,A]$. On peut donc effectuer une intégration par parties et on obtient

\begin{center}
$\int_{0}^{A}te^{-t^2}\sh(tx)\;dt=\left[-\frac{1}{2}e^{-t^2}\sh(tx)\right]_0^A+\frac{x}{2}\int_{0}^{A}e^{-t^2}\ch(tx)\;dt=-\frac{1}{2}e^{-A^2}\sh(tA)+\frac{x}{2}\int_{0}^{A}e^{-t^2}\ch(tx)\;dt$.
\end{center}

Quand $A$ tend vers $+\infty$, on obtient $f'(x)=\int_{0}^{+\infty}te^{-t^2}\sh(tx)\;dt=\frac{x}{2}\int_{0}^{+\infty}e^{-t^2}\ch(tx)\;dt=\frac{x}{2}f(x)$.

Ensuite, pour tout réel $x$, $e^{-x^2/4}f'(x)-\frac{x}{2}e^{-x^2/4}f(x)=0$ ou encore $(e^{-x^2/4}f)'(x)=0$. On en déduit que $\forall x\in\Rr$, $e^{-x^2/4}f(x)=e^0f(0)=\int_{0}^{+\infty}e^{-t^2}\;dt=\frac{\sqrt{\pi}}{2}$ et donc que $\forall x\in\Rr$, $f(x)=\frac{\sqrt{\pi}}{2}e^{x^2/4}$.

\begin{center}
\shadowbox{
$\forall x\in\Rr$, $\int_{0}^{+\infty}e^{-t^2}\ch(tx)\;dt=\frac{\sqrt{\pi}}{2}e^{x^2/4}$.
}
\end{center}
\fincorrection
\correction{005769}
Soit $x\in\Rr$. La fonction $t\mapsto \frac{t-1}{\ln t}t^x$ est continue sur $]0,1[$. 

\textbf{Etude en $1$.} $\frac{t-1}{\ln t}t^x\underset{t\rightarrow1}{\sim}1\times1=1$ et donc la fonction $t\mapsto \frac{t-1}{\ln t}t^x$ se prolonge par continuité en $1$. On en déduit que la fonction $t\mapsto \frac{t-1}{\ln t}t^x$ est intégrable sur un voisinage de $1$ à gauche.

\textbf{Etude en $0$.} $\frac{t-1}{\ln t}t^x\underset{t\rightarrow0}{\sim}-\frac{t^x}{\ln t}>0$.

-si $x>-1$, $-\frac{t^x}{\ln t}\underset{t\rightarrow0}{=}o(t^x)$ et puisque $x>-1$, la fonction $t\mapsto \frac{t-1}{\ln t}t^x$ est intégrable sur un voisinage de $0$ à droite.

- si $x\leqslant-1$, la fonction $t\mapsto-\frac{t^x}{\ln t}$ domine la fonction $t\mapsto-\frac{1}{t\ln t}$ quand $t$ tend vers $0$ par valeurs supérieures. Puisque la fonction $-\frac{1}{t\ln t}$ est positive et que $\int_{x}^{1/2}-\frac{1}{t\ln t}\;dt=\ln|\ln(x)|-\ln|\ln(1/2)|\underset{x\rightarrow0}{\rightarrow}+\infty$, la fonction $t\mapsto-\frac{1}{t\ln t}$ n'est pas intégrable sur un voisinage de $0$. Il en est de même de la fonction $t\mapsto\frac{t-1}{\ln t}t^x$.

Finalement, la fonction $t\mapsto\frac{t-1}{\ln t}t^x$ est intégrable sur $]0,1[$ si et seulement si $x>-1$. Pour $x>-1$, on peut poser $f(x)=\int_{0}^{1}\frac{t-1}{\ln t}t^x\;dt$.

\textbf{Calcul de $f(x)$.} Soit $a>-1$. On pose $\begin{array}[t]{cccc}
\Phi~:&[a,+\infty[\times]0,1[&\rightarrow&\Rr\\
 &(x,t)&\mapsto&\frac{t-1}{\ln t}t^x
\end{array}$.

\textbullet~Pour chaque $x\in[a,+\infty[$, la fonction $t\mapsto \frac{t-1}{\ln t}t^x$ est continue par morceaux et intégrable sur $]0,1[$.

\textbullet~La fonction $\Phi$ admet sur $[a,+\infty[\times]0,1[$ une dérivée partielle par rapport à sa première variable définie par :

\begin{center}
$\forall(x,t)\in[a,+\infty[\times]0,1[$, $\frac{\partial \Phi}{\partial x}(x,t)=(t-1)t^x=t^{x+1}-t^x$.
\end{center}

De plus,

- pour chaque $x\in[a,+\infty[$, la fonction $t\mapsto\frac{\partial \Phi}{\partial x}(x,t)$ est continue par morceaux sur $]0,1[$,

-pour chaque $t\in]0,1[$, la fonction $x\mapsto\frac{\partial \Phi}{\partial x}(x,t)$ est continue sur $[a,+\infty[$,

- pour chaque $(x,t)\in[a,+\infty[\times]0,1[$,

\begin{center}
$\left|\frac{\partial \Phi}{\partial x}(x,t)\right|=(1-t)t^a=\varphi(t)$.
\end{center}

La fonction $\varphi$ est continue par morceaux sur $]0,1[$ et intégrable sur $]0,1[$ car $a>-1$.

D'après le théorème de dérivation des intégrales à paramètres (théorème de \textsc{Leibniz}), la fonction $f$ est de classe $C^1$ sur $[a,+\infty[$ et sa dérivée s'obtient par dérivation sous le signe somme. Ceci étant vrai pour tout réel $a>-1$, la fonction $f$ est de classe $C^1$ sur $]-1,+\infty[$ et

\begin{center}
$\forall x>-1$, $f'(x)=\int_{0}^{1}(t-1)t^x\;dt=\left[\frac{t^{x+2}}{x+2}-\frac{t^{x+1}}{x+1}\right]_0^1=\frac{1}{x+2}-\frac{1}{x+1}$.
\end{center}

Par suite, il existe $C\in\Rr$ tel que $\forall x>-1$, $f(x)=\ln\left(\frac{x+2}{x+1}\right)+C$\quad$(*)$. Pour déterminer la constante $C$, on peut utiliser le résultat de l'exercice \ref{ex:rou7bis} : $\int_{0}^{1}\frac{t-1}{\ln t}\;dt=\ln2$. On peut aussi obtenir directement la constante $C$ sans aucun calcul d'intégrale. Pour cela, déterminons $\lim_{x \rightarrow +\infty}f(x)$.

La fonction $g~:~t\mapsto\frac{t-1}{\ln t}$ est continue sur le segment $]0,1[$, prolongeable par continuité en $0$ et en $1$ en posant $g(0)=0$ et $g(1)=1$. On en déduit que cette fonction est bornée sur l'intervalle $]0,1[$ (car son prolongement est une fonction continue sur un segment).

Soit $M$ un majorant de la fonction $|g|$ sur $]0,1[$. Pour $x>-1$, on a

\begin{center}
$|g(x)|\leqslant M\int_{0}^{1}t^x\;dt=\frac{M}{x+1}$.
\end{center}

Ceci montre que $\lim_{x \rightarrow +\infty}f(x)=0$ et en passant à la limite quand $x$ tend vers $+\infty$ dans l'égalité $(*)$, on obtient $C=0$. On a donc montré que

\begin{center}
\shadowbox{
$\forall x>-1$, $\int_{0}^{1}\frac{t-1}{\ln t}t^x\;dt=\ln\left(\frac{x+2}{x+1}\right)$.
}
\end{center}

On retrouve en particulier $\int_{0}^{1}\frac{t-1}{\ln t}\;dt=\ln2$.
\fincorrection
\correction{005770}
\textbf{Existence de $\int_{0}^{+\infty}\frac{e^{-tx}}{1+t^2}\;dt$.} Soit $x\geqslant0$. La fonction $t\mapsto\frac{e^{-tx}}{1+t^2}$ est continue sur $[0,+\infty[$ et est dominée par $\frac{1}{t^2}$ quand $t$ tend vers $+\infty$. Cette fonction est donc intégrable sur $[0,+\infty[$. Donc $\int_{0}^{+\infty}\frac{e^{-tx}}{1+t^2}\;dt$ existe pour tout réel positif $x$ et on pose $\forall x\geqslant 0$, $f(x)=\int_{0}^{+\infty}\frac{e^{-tx}}{1+t^2}\;dt$.

\textbf{Continuité de $f$ sur $[0,+\infty[$.} Soit $\begin{array}[t]{cccc}
\Phi~:&[0,+\infty[\times[0,+\infty[&\rightarrow&\Rr\\
 &(x,t)&\mapsto&\frac{e^{-tx}}{1+t^2}
 \end{array}$.
 

\textbullet~Pour chaque $x\in[0,+\infty[$, la fonction $t\mapsto\Phi(x,t)$ est continue par morceaux sur $[0,+\infty[$.
 

\textbullet~Pour chaque $t\in[0,+\infty[$, la fonction $x\mapsto\Phi(x,t)$ est continue sur $[0,+\infty[$.
 

\textbullet~Pour chaque $(x,t)\in[0,+\infty[\times[0,+\infty[$,

\begin{center}
$|\Phi(x,t)|=\frac{e^{-tx}}{1+t^2}\;dt\leqslant\frac{1}{1+t^2}=\varphi_0(t)$.
\end{center}

De plus, la fonction $\varphi_0$ est continue et intégrable sur $[0,+\infty[$ car équivalente à $\frac{1}{t^2}$ quand $t$ tend vers $+\infty$.

D'après le théorème de continuité des intégrales à paramètres, $f$ est continue sur $[0,+\infty[$. 

\textbf{Dérivée seconde de $f$.} Soit $a>0$. On pose $\begin{array}[t]{cccc}
\Phi~:&[0,+\infty[\times[a,+\infty[&\rightarrow&\Rr\\
 &(x,t)&\mapsto&\frac{e^{-tx}}{1+t^2}
 \end{array}$.
 

En plus de ce qui précède, $\Phi$ admet sur $[a,+\infty[\times[0,+\infty[$ des dérivées partielles d'ordre $1$ et $2$ définies par
 
\begin{center}
$\forall(x,t)\in[a,+\infty[\times[0,+\infty[$, $\frac{\partial\Phi}{\partial x}(x,t)=-\frac{te^{-tx}}{1+t^2}$ et $\frac{\partial^2\Phi}{\partial x^2}(x,t)=\frac{t^2e^{-tx}}{1+t^2}$.
\end{center}

\textbullet~Pour chaque $x\in[a,+\infty[$, les fonctions $t\mapsto\frac{\partial\Phi}{\partial x}(x,t)$ et $t\mapsto\frac{\partial^2\Phi}{\partial x^2}(x,t)$ sont continues par morceaux sur $[0,+\infty[$.

\textbullet~Pour chaque $t\in[0,+\infty[$, les fonctions $x\mapsto\frac{\partial\Phi}{\partial x}(x,t)$ et $x\mapsto\frac{\partial^2\Phi}{\partial x^2}(x,t)$ sont continues sur $[a,+\infty[$.

\textbullet~Pour chaque $(x,t)\in[a,+\infty[\times[0,+\infty[$,

\begin{center}
$\left|\frac{\partial\Phi}{\partial x}(x,t)\right|=\frac{te^{-tx}}{1+t^2}\leqslant\frac{te^{-at}}{1+t^2}=\varphi_1(t)$ et $\left|\frac{\partial^2\Phi}{\partial x}(x,t)\right|=\frac{t^2e^{-tx}}{1+t^2}\leqslant\frac{t^2e^{-at}}{1+t^2}=\varphi_2(t)$.
\end{center}

De plus, les fonctions $\varphi_1$ et $\varphi_2$ sont continues par morceaux et intégrables sur $[0,+\infty[$ car négligeables devant $\frac{1}{t^2}$ quand $t$ tend vers $+\infty$.

D'après une généralisation du théorème de dérivation des intégrales à paramètres, $f$ est de classe $C^2$ sur $[a,+\infty[$ et ses dérivées premières et secondes s'obtiennent par dérivation sous le signe somme. Ceci étant vrai pour tout $a>0$, $f$ est de classe $C^2$ sur $]0,+\infty[$ et

\begin{center}
\shadowbox{
$\forall x>0$, $f'(x)=\int_{0}^{+\infty}\frac{te^{-tx}}{1+t^2}\;dt$ et $f''(x)=\int_{0}^{+\infty}\frac{t^2e^{-tx}}{1+t^2}\;dt$.
}
\end{center}

\textbf{Equation différentielle vérifiée par $f$.} Pour $x>0$,

\begin{center}
$f''(x)+f(x)=\int_{0}^{+\infty}\frac{t^2e^{-tx}}{1+t^2}\;dt+\int_{0}^{+\infty}\frac{e^{-tx}}{1+t^2}\;dt=\int_{0}^{+\infty}e^{-tx}\;dt=\left[-\frac{e^{-tx}}{x}\right]_0^{+\infty}=\frac{1}{x}$.
\end{center}

\begin{center}
\shadowbox{
$\forall x>0$, $f(x)+f''(x)=\frac{1}{x}$.
}
\end{center}

\textbf{Existence de $\int_{0}^{+\infty}\frac{\sin t}{x+t}\;dt$.} Si $x=0$, l'exercice \ref{ex:rou3}, 1)  montre que $\int_{0}^{+\infty}\frac{\sin t}{t}\;dt$ est une intégrale convergente.

Si $x>0$, une intégration par parties fournit pour $A>0$

\begin{center}
$\int_{0}^{A}\frac{\sin t}{t+x}\;dt=-\frac{\cos A}{A+x}+\frac{1}{x}-\int_{0}^{A}\frac{\cos t}{(t+x)^2}\;dt$.
\end{center}

La fonction $t\mapsto\frac{\cos t}{(t+x)^2}$ est continue sur $[0,+\infty[$ et est dominée par $\frac{1}{t^2}$ quand $t$ tend vers $+\infty$. Cette fonction est donc intégrable sur $[0,+\infty[$. On en déduit que la fonction $A\mapsto\int_{0}^{A}\frac{\cos t}{(t+x)^2}\;dt$ a une limite réelle quand $A$ tend vers $+\infty$ et il en est de même de la fonction $A\mapsto\int_{0}^{A}\frac{\sin t}{t+x}\;dt$. Ainsi, pour chaque $x\in[0,+\infty[$, $\int_{0}^{+\infty}\frac{\sin t}{t+x}\;dt$ est une intégrale convergente. Pour $x\geqslant 0$, on peut donc poser $g(x)=\int_{0}^{+\infty}\frac{\sin t}{t+x}\;dt$.

\textbf{Equation différentielle vérifiée par $g$.} Pour $x>0$, on pose $u=x+t$. on obtient

\begin{center}
$g(x)=\int_{0}^{+\infty}\frac{\sin t}{t+x}\;dt=\int_{x}^{+\infty}\frac{\sin(u-x)}{u}\;du=\cos x\int_{x}^{+\infty}\frac{\sin u}{u}\;du-\sin x\int_{x}^{+\infty}\frac{\cos u}{u}\;du$.
\end{center}

(car toutes les intégrales considérées sont convergentes). Maintenant, les fonctions $c~:~u\mapsto\frac{\cos u}{u}$ et $s~:~u\mapsto\frac{\sin u}{u}$ sont continues sur $]0,+\infty[$ et admettent donc des primitives sur $]0,+\infty[$. On note $C$ (respectivement $S$) une primitive de la fonction $c$ (respectivement $s$) sur $]0,+\infty[$). Pour tout réel $x>0$, $\int_{x}^{+\infty}\frac{\cos u}{u}\;du=\lim_{t \rightarrow +\infty}C(t)-C(x)$. On en déduit que la fonction $x\mapsto\int_{x}^{+\infty}\frac{\cos u}{u}\;du$ est de classe $C^1$ sur $]0,+\infty[$, de dérivée $-c$. De même, la fonction $x\mapsto\int_{x}^{+\infty}\frac{\sin u}{u}\;du$ est de classe $C^1$ sur $]0,+\infty[$, de dérivée $-s$. Mais alors la fonction $g$ est de classe $C^1$ sur $]0,+\infty[$ et pour tout réel $x>0$,

\begin{align*}\ensuremath
g'(x)&=-\sin x\int_{x}^{+\infty}\frac{\sin u}{u}\;du-\frac{\sin x\cos x}{x}-\cos x\int_{x}^{+\infty}\frac{\cos u}{u}\;du+\frac{\cos x\sin x}{x}\\
 &=-\cos x\int_{x}^{+\infty}\frac{\cos u}{u}\;du-\sin x\int_{x}^{+\infty}\frac{\sin u}{u}\;du.
\end{align*}

La fonction $g'$ est encore de classe $C^1$ sur $]0,+\infty[$ et pour tout réel $x>0$,

\begin{align*}\ensuremath
g'(x)&=\sin x\int_{x}^{+\infty}\frac{\cos u}{u}\;du+\frac{\cos^2 x}{x}-\cos x\int_{x}^{+\infty}\frac{\sin u}{u}\;du+\frac{\sin^2 x}{x}\\
 &=\frac{1}{x}-g(x).
\end{align*}

\begin{center}
\shadowbox{
$\forall x>0$, $g''(x)+g(x)=\frac{1}{x}$.
}
\end{center}

\textbf{Egalité de $f$ et $g$ sur $]0,+\infty[$.} Pour tout réel $x>0$, $(f-g)''(x)+(f-g)(x)=0$. Donc il existe deux réels $\lambda$ et $\mu$ tels que $\forall x>0$, $(f-g)(x)=\lambda\cos x+\mu\sin x=A\cos(x+\varphi)$ pour $A=\sqrt{\lambda^2+\mu^2}$ et pour un certain $\varphi$.

Maintenant, pour $x>0$, $|f(x)|=\int_{0}^{+\infty}\frac{e^{-tx}}{1+t^2}\;dt\leqslant\int_{0}^{+\infty}e^{-tx}\;dt=\frac{1}{x}$ et on en déduit que $\lim_{x \rightarrow +\infty}\frac{1}{x}=0$.

Ensuite, $|g(x)\leqslant\left|\int_{x}^{+\infty}\frac{\sin u}{u}\;du\right|+\left|\int_{x}^{+\infty}\frac{\cos u}{u}\;du\right|$. Puisque les intégrales $\int_{1}^{+\infty}\frac{\sin u}{u}\;du$ et $\int_{1}^{+\infty}\frac{\cos u}{u}\;du$ sont des intégrales convergentes, on a $\lim_{x \rightarrow +\infty}\int_{x}^{+\infty}\frac{\sin u}{u}\;du=\lim_{x \rightarrow +\infty}\int_{x}^{+\infty}\frac{\cos u}{u}\;du=0$ et donc aussi $\lim_{x \rightarrow +\infty}g(x)=0$.

Finalement, $\lim_{x \rightarrow +\infty}(f(x)-g(x))=0$ ce qui impose $A=0$ et donc $\forall x>0$, $f(x)=g(x)$.

\begin{center}
\shadowbox{
$\forall x>0$, $\int_{0}^{+\infty}\frac{e^{-tx}}{1+t^2}\;dt=\int_{0}^{+\infty}\frac{\sin t}{t+x}\;dt$.
}
\end{center}

\textbf{Continuité de $g$ en $0$ et valeur de $\int_{0}^{+\infty}\frac{\sin t}{t}\;dt$.} Pour $x>0$,

\begin{align*}\ensuremath
g(x)&=\cos x\int_{x}^{+\infty}\frac{\sin u}{u}\;du-\sin x\int_{x}^{+\infty}\frac{\cos u}{u}\;du\\
 &=\cos x\int_{x}^{+\infty}\frac{\sin u}{u}\;du-\sin x\int_{1}^{+\infty}\frac{\cos u}{u}\;du+\sin x\int_{x}^{1}\frac{1-\cos u}{u}\;du-\sin x\ln x.
\end{align*}

Quand $x$ tend vers $0$, $\sin x\ln x\sim x\ln x$ et donc $\lim_{x \rightarrow 0}\sin x\ln x=0$. Ensuite, la fonction $u\mapsto\frac{1-\cos u}{u}$ est intégrable sur $]0,1]$ car continue sur $]0,1]$ et prolongeable par continuité en $0$. On en déduit que $\lim_{x \rightarrow 0}\sin x\int_{x}^{1}\frac{1-\cos u}{u}\;du=0\times\int_{0}^{1}\frac{1-\cos u}{u}\;du=0$. Il reste

\begin{center}
$\displaystyle\lim_{\substack{x\rightarrow0\\ x>0}}g(x)=\int_{0}^{+\infty}\frac{\sin t}{t}\;dt=g(0)$.
\end{center}

La fonction $g$ est donc continue en $0$. Puisque la fonction $f$ est également continue en $0$, on en déduit que

\begin{center}
$g(0)=\displaystyle\lim_{\substack{x\rightarrow0\\ x>0}}g(x)=\displaystyle\lim_{\substack{x\rightarrow0\\ x>0}}f(x)=f(0)=\int_{0}^{+\infty}\frac{1}{1+t^2}\;dt=\frac{\pi}{2}$.
\end{center}

\begin{center}
\shadowbox{
$\forall x\geqslant0$, $\int_{0}^{+\infty}\frac{e^{-tx}}{1+t^2}\;dt=\int_{0}^{+\infty}\frac{\sin t}{t+x}\;dt$ et en particulier, $\int_{0}^{+\infty}\frac{\sin t}{t}\;dt=\frac{\pi}{2}$.
}
\end{center}
\fincorrection
\correction{005771}
\begin{enumerate}
 \item
\textbullet~Soit $x\in\Rr$. La fonction $t\mapsto f(x-t)g(t)$ est continue sur le segment $[0,T]$ et donc intégrable sur ce segment. Par suite, $f*g(x)$ existe.

\textbullet~Soit $x\in\Rr$. $f*g(x+T)=\int_{0}^{T}f(x+T-t)g(t)\;dt=\int_{0}^{T}f(x-t)g(t)\;dt=f*g(x)$. Donc la fonction $f*g$ est $T$-périodique.

\textbullet~Les fonction $f$ et $g$ sont continues sur $\Rr$ et $T$-périodiques. Ces fonctions sont en particulier bornées sur $\Rr$. On note $M_1$ et $M_2$ des majorants sur $\Rr$ des fonctions $|f|$ et $|g|$ respectivement.

Soit $\begin{array}[t]{cccc}
\Phi~:&\Rr\times[0,T]&\rightarrow&\Rr\\
 &(x,t)&\mapsto&f(x-t)g(t)
\end{array}$.

- Pour chaque $x\in\Rr$, la fonction $t\mapsto\Phi(x,t)$ est continue par morceaux sur $[0,T]$.

- Pour chaque $t\in[0,T]$, la fonction $x\mapsto\Phi(x,t)$ est continue sur $\Rr$.
 

- Pour chaque $(x,t)\in\Rr\times[0,T]$, $|\Phi(x,t)|\leqslant M_1M_2=\varphi(t)$. De plus, la fonction $\varphi$ est continue et donc intégrable sur le segment $[0,T]$.

D'après le théorème de continuité des intégrales à paramètres, la fonction $f*g$ est continue sur $\Rr$.

\item  Soient $f$ et $g$ deux éléments de $E$. Soit $x\in\Rr$. En posant $u=x-t$, on obtient

\begin{align*}\ensuremath
f*g(x)&=\int_{0}^{T}f(x-t)g(t)\;dt=\int_{x}^{x-T}f(u)g(u-t)\;(-du)=\int_{x-T}^{x}g(u-t)f(u)\;du\\
 &=\int_{0}^{T}g(u-t)f(u)\;du\;(\text{car la fonction}\;u\mapsto g(u-t)f(u)\;\text{est}\;T\text{-périodique})\\
 &=g*f(x).
\end{align*}
\end{enumerate}
\fincorrection
\nocorrection
\nocorrection
\correction{000847}

% Correction : Nicolas Bonnotte

Le but de l'exercice est de résoudre l'équation
\begin{equation}\tag{E}\label{eq}
y'(x)- \frac{y(x)}{x}-y(x)^2=-9x^2.
\end{equation}

\begin{enumerate}
    \item Trouvons $a \in ]0, \infty[ $ tel que $y_0(x)=ax$ soit une solution
particuli\`ere. Puisque
\[ y_0'(x) - \frac{y_0(x)}{x}-y_0(x)^2 = - a^2 x^2,\]
$y_0$ est solution si et seulement si $a = \pm 3$. On choisit $a=3$.

    \item Si $z$ est une fonction $\mathcal{C}^1$ ne s'annulant pas, on pose $y(x) = 3x-1/{z(x)}$. Alors $y$ est solution si et seulement si
\[ \frac{z'(x)}{z(x)^2} + \frac{1}{xz(x)} - \frac{1}{z(x)^2} + \frac{6x}{z(x)} = 0.\]
En multipliant par $z(x)^2$, on obtient que $y$ est solution de \eqref{eq} ssi $z$ vérifie
\begin{equation}\tag{E$_1$}\label{eq1}
z'(x) + \left(6x+ \frac{1}{x}\right)z(x) =1.
\end{equation}


    \item On résout \eqref{eq1} sur $]0,\infty[$. Une primitive de $x \mapsto 6x + 1/x$ est $x \mapsto 3x^2 + \ln(x)$, donc les solutions de l'équation homogène sont les $x \mapsto A \exp(-3x^2 - \ln(x))$. On cherche une solution particulière de \eqref{eq1} sous la forme $z_p(x) = \alpha(x) \exp(-3x^2-\ln(x))$ ; alors $z_p$ est solution si
 $\alpha'(x) \exp(-3x^2 - \ln(x)) = 1$, c'est-à-dire si $\alpha'(x) = x\exp(3x^2)$, par exemple si $\alpha(x) = \exp(3x^2)/6$. Les solutions de \eqref{eq1} sont donc les 
 \[ z(x) = \frac{1+A\exp(-3x^2)}{6x}, \qquad \text{ avec } A \in \R.\]


   \item  On va maintenant en déduire les solutions de \eqref{eq} d\'efinies sur $]0,\infty[$. 
    
    Soit $y$ une solution $\mathcal{C}^1$ définie sur $]0,\infty[$. On suppose dans un premier temps que $y(x) > 3x$ sur l'intervalle ouvert $I \subset ]0,\infty[$, pris aussi grand que possible.  Alors $y(x) = 3x - 1/{z_I(x)}$ pour une certaine fonction $z_I < 0$ qui est $\mathcal{C}^1$ sur $I$. D'après la question précédente, on a nécessairement $z_I(x) = [1+A_I\exp(-3x^2)]/6x$ pour une certaine constante $A_I \in \R$. Puisque $z_I < 0$, cela impose $A_I<0$, mais du coup $I \neq ]0,+\infty[$ car $1 > A_I\exp(-3x^2)$ si $x$ est assez grand.
  
  Dans tous les cas, il existe donc un intervalle ouvert $J$ tel que $y(x) < 3x$ sur $J$. On suppose encore que $J$ est aussi grand que possible. Sur $J$, $y(x) = 3x - 1/{z_J(x)}$ pour une certaine fonction $z_J > 0$ qui est $\mathcal{C}^1$ sur $J$. Encore d'après la question précédente, $z_J(x) = [1+A_J\exp(-3x^2)]/6x$ pour une certaine constante $A_J$.
 Puisque l'intervalle ouvert $J = ]a,b[$ a été supposé maximal, et puisque $y$ est supposée définie sur $]0,+\infty[$, si $a > 0$ on a $y(a) = 3a$ et de même si $b < \infty$, $y(b) = 3b$, car sinon par continuité de $y$ on aurait encore $y(x) < 3x$ sur $]a-\epsilon,b+\epsilon[$ pour un petit $\epsilon > 0$. Cela n'est possible respectivement que si $z_J(x) \rightarrow +\infty$ lorsque $x \rightarrow a$ ou $z_J(x) \rightarrow +\infty$ lorsque $x \rightarrow b$. Or on a dit que $z_J = [1+A_J\exp(-3x^2)]/6x$, cela n'est donc pas possible du tout (sauf précisément si respectivement $a =0$ et $b=0$).
  
  Donc soit $y(x) = 3x$ sur $]0,+\infty[$, soit $y(x) < 3x$ sur $]0,+\infty[$. Dans ce dernier cas, $z(x) = 1/(3x-y(x))$ est définie sur $]0,+\infty[$ et s'écrit $z(x) = [1+A\exp(-3x^2)]/6x$. Puisque $z > 0$, nécessairement $A \geq -1$. Donc si $y$ est solution, alors     
    \[ y(x) = 3x \qquad \text{ou} \qquad y(x) =  3x - \frac{6x}{1+A\exp(-3x^2)} \quad \text{avec } A \geq -1.\]
    
 Réciproquement, si $y$ est ainsi définie, alors $y$ est définie et $\mathcal{C}^1$ sur $]0,\infty[$, et on peut vérifier que c'est bien une solution.
\end{enumerate}
\fincorrection
\nocorrection
\nocorrection
\nocorrection
\correction{000851}
Les primitives de la fonction $a(x)=2x$ sont les fonctions $A(x)=
x^2/2 + k $ o\`u $k \in \R$ est une constante r\'eelle quelconque.
Donc les solutions de l'\'equation homog\`ene associ\'ee \`a $E$
sont toutes les fonctions d\'efinies sur $\R$ du type : $ y(x)=
ce^{-x^2}$ o\`u $c\in\R$ est une constante arbitraire. On cherche
maintenant une solution particuli\`ere de $E$ sous la forme
$y_p(x)=c(x)e^{-x^2}$ (m\'ethode de la variation de la constante).
On a :\\ $y_p^\prime(x)+ 2xy_p(x)= c^\prime(x)e^{-x^2}$. Donc
$y_p$ est solution de $E$ si et seulement si :
$c^\prime(x)=xe^{x^2}$ pour tout $x\in\R$. On choisit la fonction
$c$ parmi les primitives de la fonction $xe^{x^2}$, par exemple :
$c(x)=1/2e^{x^2}$. Donc la fonction $y_p$
telle que $y_p(x)=1/2e^{x^2}e^{-x^2}=1/2$ est solution de $E$.\\
Par cons\'equent les solutions de $E$ sont toutes les fonctions de
la forme :
  $$ y(x)= ce^{-x^2} + \frac{1}{2} \; c \in\R .$$
Pour $y$ solution de $E_1$, la condition $y(0)=1$ \'equivaut \`a :
$c=1/2$.

\fincorrection
\nocorrection
\correction{000853}
% Correction : Sylvie Ruette
Si $z=x+\lambda y$, alors $\dot z=x(1+3\lambda)+y(1-\lambda)$.
Pour que $\dot z$ soit proportionnel à $z$, il suffit que $\lambda(1+3\lambda)
=1-\lambda$, autrement dit $3\lambda^2+2\lambda-1=0$. Il y a deux
solutions : $a=-1$ et $b=1/3$. On pose $u=x-y$ et $v=x+\frac 13 y$.
On a alors : $\dot u=-2u$ et $\dot v=2v$. On en déduit que les solutions
$u$ et $v$ sont de la forme 
$u(t)=C_1 e^{-2t}$ et $v(t)=C_2e^{2t}$, avec $C_1,C_2\in\mathbb{R}$.
De plus,
$$
\left\{\begin{array}{rcl}
u&=&x-y\\
v&=&x+\frac 13 y
\end{array}\right.
\Longleftrightarrow
\left\{\begin{array}{rcl}
x&=&\frac 14(u+3v)\\
y&=&\frac 34(v-u)
\end{array}\right.
$$
Donc les solutions du système différentiel sont les couples de fonctions 
$(x,y)$ de la forme 
$$x(t)=D_1 e^{-2t}+3D_2e^{2t},\ y(t)=-3D_1 e^{-2t}+3D_2e^{2t},\ 
D_1,D_2\in\mathbb{R}.$$

On veut $x(0)=2$ et $y(0)=-2$. Ceci est équivalent à
$$
\left\{\begin{array}{rcrcr}
D_1&+&3D_2&=&2\\
-3D_1&+&3D_2&=&-2
\end{array}\right.\Longleftrightarrow
\left\{\begin{array}{rcr}
D_1&=&1\\D_2&=&1/3
\end{array}\right.
$$
Conclusion : il y a une unique solution, donnée par
$$x(t)=e^{-2t}+e^{2t},\ y(t)=-3e^{-2t}+e^{2t}.$$

\fincorrection
\nocorrection
\nocorrection
\nocorrection
\nocorrection
\nocorrection
\correction{005476}
 Les équations différentielles à résoudre dans cet exercice sont toutes linéaires du premier ordre.
On note $(E)$ l'équation différentielle proposée et $(E_H)$ l'équation homogène associée.

\begin{enumerate}

\item  Les fonctions $x\mapsto\frac{1}{x\ln x}$ et $x\mapsto\frac{1}{\ln x}$ sont continues sur $I$ et on sait
que les solutions de $(E)$ sur $I$ sont de la forme
$f_0+\lambda f_1$ où $f_0$ est une solution particulière de $(E)$ et $f_1$ est une solution particulière non nulle de
$(E_H)$.

Soit $f$ une fonction dérivable sur $I$.

\begin{align*}\ensuremath
f\;\mbox{solution de}\;(E)\;\mbox{sur}\;I&\Leftrightarrow\forall x\in I,\;x\ln xf'(x)+f(x)=x\Leftrightarrow\forall x\in I,\;\ln
xf'(x)+\frac{1}{x}f(x)=1\\
 &\Leftrightarrow\forall x\in I,\;(\ln x.f)'(x)=1\Leftrightarrow\exists \lambda\in\Rr/\;\forall x\in I,\;f(x)=\frac{x+\lambda}{\ln x}
\end{align*}

\item  Les fonctions $x\mapsto\frac{1}{x}$ et $x\mapsto1+\frac{1}{x^2}$ sont continues sur $I$ et on sait
que les solutions de $(E)$ sur $I$ sont de la forme
$f_0+\lambda f_1$ où $f_0$ est une solution particulière de $(E)$ et $f_1$ est une solution particulière non nulle de
$(E_H)$.

Soit $f$ une fonction dérivable sur $I$.

\begin{align*}\ensuremath
f\;\mbox{solution de}\;(E)\;\mbox{sur}\;I&\Leftrightarrow\forall x\in I,\;x(xf'(x)+f(x)-x)=1
\Leftrightarrow\forall x\in I,\;(xf)'(x)=x+\frac{1}{x}\\
 &\Leftrightarrow\exists\lambda\in\Rr/\;\forall x\in I,\;
xf(x)=\frac{x^2}{2}+\ln(-x)+\lambda\Leftrightarrow\exists\lambda\in\Rr/\;\forall x\in I,\;f(x)=\frac{x}{2}+\frac{\ln
(-x)+\lambda}{x}.
\end{align*}

Les solutions de $(E)$ sur  $I$ sont les fonctions de la forme $x\mapsto\frac{x}{2}+\frac{\ln(-x)+\lambda}{x}$,
$\lambda\in\Rr$.

\item  Les fonctions $x\mapsto\frac{1}{2x}$ et $x\mapsto\frac{x^3}{2}$ sont continues sur $I=]-\infty,0[$ et on
saitque les solutions de $(E)$ sur $I$ sont de la forme
$f_0+\lambda f_1$ où $f_0$ est une solution particulière de $(E)$ et $f_1$ est une solution particulière non nulle de
$(E_H)$.

Soit $f$ une fonction dérivable sur $I$.

\begin{align*}\ensuremath
f\;\mbox{solution de}\;(E)\;\mbox{sur}\;I&\Leftrightarrow\forall x\in I,\;f'(x)+\frac{1}{2x}f(x)=\frac{x^3}{2}\\
 &\forall x\in I,\;e^{\ln|x|/2}f'(x)+\frac{1}{2x}e^{\ln|x|/2}f(x)=\frac{x^3}{2}e^{\ln|x|/2}\Leftrightarrow\forall x\in
I,\;(\sqrt{-x}f)'(x)=-\frac{1}{2}(-x)^{7/2}\\
 &\Leftrightarrow\exists\lambda\in\Rr/\;\forall x\in I,\;\sqrt{-x}f(x)=\frac{1}{9}(-x)^{9/2}+\lambda\\
 &\Leftrightarrow\exists\lambda\in\Rr/\;\forall x\in I,\;f(x)=\frac{1}{9}x^4+\frac{\lambda}{\sqrt{-x}}
\end{align*}

Les solutions de $(E)$ sur  $I$ sont les fonctions de la forme $x\mapsto\frac{x^4}{9}+\frac{\lambda}{\sqrt{-x}}$,
$\lambda\in\Rr$.

\item  Les fonctions $x\mapsto2$ et $x\mapsto x^2-3x$ sont continues sur $\Rr$ et on sait
que les solutions de $(E)$ sur $\Rr$ sont de la forme
$f_0+\lambda f_1$ où $f_0$ est une solution particulière de $(E)$ et $f_1$ est une solution particulière non nulle de
$(E_H)$.

Soit $f$ une fonction dérivable sur $\Rr$.

\begin{align*}\ensuremath
f\;\mbox{solution de}\;(E)\;\mbox{sur}\;\Rr&\Leftrightarrow\forall x\in\Rr,\;f'(x)+2f(x)=x^2-3x\\
 &\Leftrightarrow\forall x\in\Rr,\;e^{2x}f'(x)+2e^{2x}f(x)=(x^2-3x)e^{2x}\Leftrightarrow\forall
x\in\Rr,\;(e^{2x}f)'(x)=(x^2-3x)e^{2x}
\end{align*}

Recherche d'une primitive sur ${\Rr}$ de la fonction $x\mapsto(x^2-3x)e^{2x}$.
\begin{itemize}
\item[\textbf{1ère méthode}.] Deux intégrations par parties fournissent~:

\begin{align*}\ensuremath
\int_{}^{}(x^2-3x)e^{2x}\;dx&=\frac{1}{2}(x^2-3x)e^{2x}-\frac{1}{2}\int_{}^{}(2x-3)e^{2x}\;dx\\
 &=\frac{1}{2}(x^2-3x)e^{2x}-\frac{1}{4}(2x-3)e^{2x}+\frac{1}{2}\int_{}^{}e^{2x}\;dx\\
 &=\frac{1}{4}(2x^2-8x+3)e^{2x}+\frac{1}{4}e^{2x}+C=\frac{1}{2}(x^2-4x+2)e^{2x}+C
\end{align*}

\item[\textbf{2ème méthode}.] Cherchons les primitives de $x\mapsto(x^2-3x)e^{2x}$ sous la forme
$x\mapsto(ax^2+bx+c)e^{2x}$.

$$((ax^2+bx+c)e^{2x})'=(2(ax^2+bx+c)+(2ax+b))e^{2x}=(2ax^2+2(a+b)x+b+2c)e^{2x}.$$

Donc,

$$((ax^2+bx+c)e^{2x})'=(x^2-3x)e^{2x}\Leftrightarrow
\left\{
\begin{array}{l}
2a=1\\
2(a+b)=-3\\
b+2c=0
\end{array}
\right.\Leftrightarrow\left\{
\begin{array}{l}
a=\frac{1}{2}\\
b=-2\\
c=1
\end{array}
\right..$$

\end{itemize}

\textbf{Résolution de} $\bf{(E)}$.

\begin{align*}\ensuremath
f\;\mbox{solution de}\;(E)\;\mbox{sur}\;\Rr&\Leftrightarrow\forall
x\in\Rr,\;(e^{2x}f)'(x)=(x^2-3x)e^{2x}\Leftrightarrow\exists\lambda\in\Rr/\;\forall
x\in\Rr,\;e^{2x}f(x)=(\frac{x^2}{2}-2x+1)e^{2x}+\lambda\\
 &\Leftrightarrow\exists\lambda\in\Rr/\;\forall
x\in\Rr,\;f(x)=\frac{x^2}{2}-2x+1+\lambda e^{-2x}.
\end{align*}

Les solutions de $(E)$ sur  $\Rr$ sont les fonctions de la
forme$x\mapsto\frac{x^2}{2}-2x+1+\lambda e^{-2x}$,$\lambda\in\Rr$.

\item  Les fonctions $x\mapsto1$ et $x\mapsto\frac{1}{1+2e^x}$ sont continues sur $\Rr$ et on sait
que les solutions de $(E)$ sur $\Rr$ sont de la forme
$f_0+\lambda f_1$ où $f_0$ est une solution particulière de $(E)$ et $f_1$ est une solution particulière non nulle de
$(E_H)$.

Soit $f$ une fonction dérivable sur $\Rr$.

\begin{align*}\ensuremath
f\;\mbox{solution de}\;(E)\;\mbox{sur}\;\Rr&\Leftrightarrow\forall x\in\Rr,\;f'(x)+f(x)=\frac{1}{1+2e^x}
\Leftrightarrow\forall x\in\Rr,\;e^xf'(x)+e^xf(x)=\frac{e^x}{1+2e^x}\\
 &\Leftrightarrow\exists\lambda\in\Rr/\;\forall x\in\Rr,\;e^xf(x)=\frac{1}{2}\ln(1+2e^x)+\lambda\\
 &\Leftrightarrow\exists\lambda\in\Rr/\;\forall x\in\Rr,\;f(x)=(\frac{1}{2}\ln(1+2e^x)+\lambda)e^{-x}
\end{align*}

Les solutions de $(E)$ sur  $\Rr$ sont les fonctions de la forme
$x\mapsto(\frac{1}{2}\ln(1+2e^x)+\lambda)e^{-x}$,$\lambda\in\Rr$.

\item  Les fonctions $x\mapsto-\frac{\cos x}{\sin x}$ et $x\mapsto-\frac{1}{\sin x}$ sont continues sur $I=]0,\pi[$
et on sait que les solutions de $(E)$ sur $\Rr$ sont de la forme
$f_0+\lambda f_1$ où $f_0$ est une solution particulière de $(E)$ et $f_1$ est une solution particulière non nulle de
$(E_H)$.

Mais $x\mapsto\sin x$ est une solution non nulle de $(E_H)$ sur $I$ et $x\mapsto\cos x$ est une solution de $(E)$ sur
$]0,\pi[$.

Les solutions de $(E)$ sur  $]0,\pi[$ sont les fonctions de la forme $x\mapsto\lambda\sin x+\cos x$, $\lambda\in\Rr$.

\end{enumerate}
\fincorrection
\correction{005477}
L'équation différentielle à résoudre dans cet exercice est linéaire du premier ordre. On note $(E)$ l'équation
différentielle proposée et $(E_H)$ l'équation homogène associée.

Soit $I$ l'un des deux intervalles $]-1,1[$ ou $]1,+\infty[$. Les fonctions $x\mapsto\frac{-2x}{1-x^2}$ et
$x\mapsto\frac{x^2}{1-x^2}$ sont continues sur $I$ et on sait que les solutions de $(E)$ sur $I$ sont de la forme
$f_0+\lambda f_1$ où $f_0$ est une solution particulière de $(E)$ et $f_1$ est une solution particulière non nulle de
$(E_H)$.

Résolution de $(E)$ sur $I$. Soit $f$ une fonction dérivable sur $I$.

\begin{align*}\ensuremath
f\;\mbox{solution de}\;(E)\;\mbox{sur}\;I&\Leftrightarrow\forall x\in I,\;(1-x^2)f'(x)-2xf(x)=x^2\\
 &\Leftrightarrow\forall x\in I,\;((1-x^2)f)'(x)=x^2\Leftrightarrow\exists\lambda\in\Rr/\;\forall x\in
I,\;(1-x^2)f(x)=\frac{x^3}{3}+\lambda\\ &\Leftrightarrow\exists \lambda\in\Rr/\;\forall x\in 
I,\;f(x)=\frac{x^3+\lambda}{3(1-x^2)},\end{align*}

(en renommant $\lambda$ la constante $3\lambda$).

Si $I=]-1,+\infty[$.

Soit $f$ une éventuelle solution de $(E)$ sur $I$. Les restrictions de $f$ à $]-1,1[$ et $]1,+\infty[$ sont
encore solution de $(E)$ et donc de la forme précédente. Par suite, nécessairement, il existe deux constantes
$\lambda_1$ et$\lambda_2$ telles que, pour $-1<x<1$, $f(x)=\frac{x^3+\lambda_1}{3(1-x^2)}$ et pour $x>1$,
$f(x)=\frac{x^3+\lambda_2}{3(1-x^2)}$. Enfin, l'équation impose $f(1)=-\frac{1}{2}$.

En résumé, une éventuelle solution de $(E)$ sur $I$ est nécessairement de la forme~:

$$\forall x>-1,\;f(x)=\left\{
\begin{array}{l}
\frac{x^3+\lambda_1}{3(1-x^2)}\;\mbox{si}\;-1<x<1\\
-\frac{1}{2}\;\mbox{si}\;x=1\\
\frac{x^3+\lambda_2}{3(1-x^2)}\;\mbox{si}\;x>1
\end{array}
\right..$$

Réciproquement, $f$ ainsi définie, est dérivable sur $]-1,1[$ et solution de $(E)$ sur $]-1,1[$, dérivable sur
$]1,+\infty[$ et solution de $(E)$ sur $]1,+\infty[$ et, si $f$ est dérivable en $1$, $f$ vérifie encore $(E)$ pour
$x=1$. Donc, $f$ est solution de $(E)$ sur $]-1,+\infty[$ si et seulement si $f$ est dérivable en $1$.

Pour $-1<x<1$,

$$\frac{f(x)-f(1)}{x-1}=\frac{\frac{x^3+\lambda_1}{3(1-x^2)}+\frac{1}{2}}{x-1}=\frac{2x^3+2\lambda_1+3(1-x^2)}
{6(1-x^2)(x-1)}$$

Quand $x$ tend vers $1$ par valeurs inférieures, le dénominateur de la fraction tend vers $0$ et le numérateur tend vers
$2(1+\lambda_1)$. Donc, si $\lambda_1\neq-1$, $f$ n'est pas dérivable à gauche en $1$. De même, si $\lambda_2$ n'est pas
$-1$, $f$ n'est pas dérivable à droite en $-1$. Ainsi, si $f$ est solution de $(E)$ sur $I$, nécessairement
$\lambda_1=\lambda_2=-1$. Dans ce cas, pour $x\in]-1,+\infty[\setminus\{1\}$,

$$f(x)=\frac{x^3-1}{3(1-x^2)}=\frac{(x-1)(x^2+x+1)}{3(1-x)(1+x)}=-\frac{x^2+x+1}{3(x+1)},$$

ce qui reste vrai pour $x=1$. Ainsi, si $f$ est une solution de $(E)$ sur $]-1,+\infty[$, nécessairement
pour $x>-1$, $f(x)=-\frac{x^2+x+1}{3(x+1)}$. Réciproquement, $f$ ainsi définie est dérivable sur $]-1,+\infty[$ et en
particulier en $1$. $f$ est donc solution de $(E)$ sur $]-1,+\infty[$.

Sur $]-1,+\infty[$, $(E)$ admet une et une seule solution à savoir la fonction $x\mapsto-\frac{x^2+x+1}{3(x+1)}$.

Si $I=\Rr$, soit $f$ une éventuelle solution de $(E)$ sur $\Rr$. La restriction de $f$ à $]-1,+\infty[$ est
nécessairement la fonction précédente. Mais cette fonction tend vers $-\infty$ quand $x$ tend vers $-1$ par valeurs
supérieures. Donc $f$ ne peut être continue sur $\Rr$ et $(E)$ n'a pas de solution sur $\Rr$.
\fincorrection
\correction{005478}
\textbf{Résolution} de $(E)$ sur $]0,+\infty[$.

Soit $f$ une fonction dérivable sur $]0,+\infty[$.

\begin{align*}\ensuremath
f\;\mbox{solution de}\;(E)\;&\mbox{sur}\;]0,+\infty[\Leftrightarrow\forall x\in]0,+\infty[,\;|x|f'(x)+(x-1)f(x)=x^3\\
 &\Leftrightarrow\forall x\in]0,+\infty[,\;xf'(x)+(x-1)f(x)=x^3\Leftrightarrow\forall x\in]0,+\infty[,\;f'(x)+(1-\frac{1}{x})f(x)=x^2\\
 &\Leftrightarrow\forall x\in]0,+\infty[,\;e^{x-\ln
x}f'(x)+(1-\frac{1}{x})e^{x-\ln x}f(x)=e^{x-\ln x}x^2\\
 &\Leftrightarrow\forall x\in]0,+\infty[,\;(\frac{e^x}{x}f)'(x)=xe^x=((x-1)e^x)'\\
 &\Leftrightarrow\exists\lambda\in\Rr/\;\forall
x\in]0,+\infty[,\;f(x)=xe^{-x}((x-1)e^x+\lambda)=x^2-x+\lambda xe^{-x}
\end{align*}

Les solutions de $(E)$ sur $]0,+\infty[$ sont les fonctions de la forme $x\mapsto x^2-x+\lambda xe^{-x}$,
$\lambda\in\Rr$.

\textbf{Résolution} de $(E)$ sur $]-\infty,0[$.

Soit $f$ une fonction dérivable sur $]-\infty,0[$.

\begin{align*}\ensuremath
f\;\mbox{solution de}\;(E)\;\mbox{sur}\;]-\infty,0[&\Leftrightarrow\forall x\in]0,+\infty[,\;
-xf'(x)+(x-1)f(x)=x^3\\
 &\Leftrightarrow\forall x\in]-\infty,0[,\;f'(x)+(-1+\frac{1}{x})f(x)=-x^2\\
 &\Leftrightarrow\forall x\in]-\infty,0[,\;e^{-x+\ln|x|}f'(x)+
(-1+\frac{1}{x})e^{-x+\ln|x|}f(x)=-e^{-x+\ln|x|}x^2\\
 &\Leftrightarrow\forall x\in]-\infty,0[,\;(-xe^{-x}y)'=x^3e^{-x}\;(*)
\end{align*}

Déterminons une primitive de la fonction $x\mapsto-x^3e^{-x}$ de la forme $(ax^3+bx^2+cx+d)e^{-x}$.

$$((ax^3+bx^2+cx+d)e^{-x})'=(-(ax^3+bx^2+cx+d)+(3ax^2+2bx+c))e^{-x}=(-ax^3+(3a-b)x^2+(2b-c)x+c-d)e^{-x},$$

et

$$((ax^3+bx^2+cx+d)e^{-x})'=x^3e^{-x}\Leftrightarrow
\left\{
\begin{array}{l}
-a=-1\\
3a-b=0\\
2b-c=0\\
c-d=0
\end{array}
\right.
\Leftrightarrow
\left\{
\begin{array}{l}
a=1\\
b=3\\
c=6=d
\end{array}
\right..$$

Par suite,

\begin{align*}\ensuremath
(*)&\Leftrightarrow\exists\lambda\in\Rr/\;\forall x\in]-\infty,0[,\;xe^{-x}f(x)=(x^3+3x^2+6x+6)e^{-x}+\lambda\\
 &\Leftrightarrow\exists\lambda\in\Rr/\;\forall x\in]-\infty,0[,\;
f(x)=x^2+3x+6+\frac{\lambda e^x+6}{x}.
\end{align*}

Les solutions de $(E)$ sur $]-\infty,0[$ sont les fonctions de la forme $x\mapsto x^2+3x+6+\frac{\lambda e^x+6}{x}$,
$\lambda\in\Rr$.

On peut montrer que l'équation admet une et une seule solution sur $\Rr$ en \og~recollant~\fg~les expressions
précédentes, mais en ce début d'année, on manque encore d'outils.
\fincorrection
\correction{005481}
On sait que les solutions sur $\Rr$ de l'équation proposée sont les fonctions de la forme~:

$$g~:~x\mapsto\lambda e^{-ax}+e^{-ax}\int_{0}^{x}e^{at}f(t)\;dt,\;\lambda\in\Rr.$$
Dans ce cas, pour $x\in\Rr$, $g(x+T)=\lambda e^{-a(x+T)}+e^{-a(x+T)}\int_{0}^{x+T}e^{at}f(t)\;dt$. Or,

\begin{align*}\ensuremath
\int_{0}^{x+T}e^{at}f(t)\;dt&=\int_{0}^{x}e^{at}f(t)\;dt+\int_{x}^{x+T}e^{at}f(t)\;dt=\int_{0}^{x}e^{at}f(t)\;dt+\int_{0}^{T}
e^{a(u+T)}f(u+T)\;du\\
 &=\int_{0}^{x}e^{at}f(t)\;dt+e^{aT}\int_{0}^{T}e^{au}f(u)du.
\end{align*}

Donc,

\begin{align*}\ensuremath
g(x+T)&=\lambda e^{-a(x+T)}+e^{-a(x+T)}\int_{0}^{x}e^{at}f(t)\;dt+e^{-ax}\int_{0}^{T}e^{au}f(u)\;du\\
 &=\lambda e^{-a(x+T)}+e^{-a(x+T)}\int_{0}^{T}e^{at}f(t)\;dt+g(x)-\lambda e^{-ax}.
\end{align*}

Par suite,

\begin{align*}\ensuremath
g\;\mbox{est}\;T\mbox{-périodique}&\Leftrightarrow\forall
x\in\Rr,\;\lambda e^{-a(x+T)}+e^{-a(x+T)}\int_{0}^{T}e^{at}f(t)\;dt-\lambda e^{-ax}=0\\
 &\Leftrightarrow\lambda(1-e^{-aT})=e^{-aT}\int_{0}^{T}e^{at}f(t)\;dt\Leftrightarrow\lambda=\frac{e^{-aT}}{1-e^{-aT}}\int_{0}^{T}e^{at}f(t)\;dt
\end{align*}

($e^{-aT}\neq1$ car $a\neq0$ et $T\neq0$). D'où l'existence et l'unicité d'une solution $T$-périodique~:

$$\forall x\in\Rr,\;g(x)=\frac{e^{-aT}}{1-e^{-aT}}\int_{0}^{T}e^{at}f(t)\;dte^{-ax}+e^{-ax}\int_{0}^{x}e^{at}f(t)\;dt.$$

\fincorrection
\correction{006991}\
\begin{enumerate}
\item Il s'agit d'une équation différentielle linéaire d'ordre 1, à coefficients constants, avec second membre.\\ 

On commence par résoudre l'équation homogène associée $y'+2y=0$: les solutions sont les $y(x)=\lambda e^{-2x}$, $\lambda\in\R$.\\

Il suffit ensuite de trouver une solution particulière de $(E_1)$. Le second membre étant polynomial de degré 2, on cherche une solution particulière de la m\^eme forme: \\
$y_0(x)=ax^2+bx+c$ est solution de $(E_1)$\\
\ \ \ $\Longleftrightarrow \forall x\in\R,\ y_0'(x)+2y_0(x)=x^2$ \\
\ \ \  $\Longleftrightarrow \forall x\in\R,\ 2ax^2+(2a+2b)x+b+2c=x^2$ 

\noindent Ainsi, en identifiant les coefficients, on voit que $y_0(x)=\frac{1}{2}x^2-\frac{1}{2}x+\frac{1}{4}$ convient.

Les solutions de $(E_1)$ sont obtenues en faisant la somme de cette solution particulière et des solutions 
de l'équation homogène: 
$$y(x)=\frac{1}{2}x^2-\frac{1}{2}x+\frac{1}{4}+\lambda e^{-2x}\quad (x\in\R)$$
où $\lambda$ est un paramètre réel.
\item Il s'agit d'une équation différentielle linéaire d'ordre 1, à coefficients constants, avec second membre.\\ 


Les solutions de l'équation homogène associée $y'+y=0$ sont les $y(x)=\lambda e^{-x}$, $\lambda\in\R$.\\

Il suffit ensuite de trouver une solution particulière de $(E_2)$. Le second membre est cette fois une fonction trigonométrique, on cherche une solution particulière sous la forme d'une combinaison linéaire de $\cos$ et $\sin$:\\
 $y_0(x)=a\cos x+b\sin x$ est solution de $(E_2)$\\
$\Longleftrightarrow \forall x\in\R,\ y_0'(x)+y_0(x)=2\sin x$ \\
$\Longleftrightarrow \forall x\in\R,\ (a+b)\cos x+(-a+b)\sin x=2\sin x$ 

\noindent Ainsi, en identifiant les coefficients, on voit que $y_0(x)=-\cos x+\sin x$ convient.

Les solutions de $(E_2)$ sont obtenues en faisant la somme de cette solution particulière et des solutions 
de l'équation homogène:
$$y(x)=-\cos x+\sin x+\lambda e^{-x}\quad (x\in\R)$$
où $\lambda$ est un paramètre réel.
\item Les solutions de l'équation homogène associée $y'-y=0$ sont les $y(x)=\lambda e^{x}$, $\lambda\in\R$. On remarque que le second membre est le produit d'une fonction exponentielle par une fonction polynomiale de degré $d=1$: or la fonction exponentielle du second membre est la m\^eme ($e^x$) que celle qui appara\^it dans les solutions de l'équation homogène. On cherche donc une solution particulière sous la forme d'un produit de $e^x$ par une fonction polynomiale de degré $d+1=2$:\\
$y_0(x)=(ax^2+bx+c)e^x$ est solution de $(E_3)$ \\
$\Longleftrightarrow \forall x\in\R,\ y_0'(x)-y_0(x)=(x+1)e^x$ \\
$\Longleftrightarrow \forall x\in\R,\ (2ax+b)e^x=(x+1)e^x$ 

\noindent Ainsi, en identifiant les coefficients, on voit que $y_0(x)=(\frac{1}{2}x^2+x)e^x$ convient.

Les solutions de $(E_3)$ sont obtenues en faisant la somme de cette solution particulière et des solutions de l'équation homogène: 
$$y(x)=(\frac{1}{2}x^2+x+\lambda)e^{x}\quad (x\in\R)$$
où $\lambda$ est un paramètre réel.
\item Les solutions de l'équation homogène associée $y'+y=0$ sont les $y(x)=\lambda e^{-x}$, $\lambda\in\R$. On remarque que le second membre est la somme d'une fonction polynomiale de degré 1, d'une fonction exponentielle (différente de $e^{-x}$) et d'une fonction trigonométrique. D'après le principe de superposition, on cherche donc une solution particulière sous la forme d'une telle somme:\\
$y_0(x)=ax+b+\mu e^x+\alpha\cos x+\beta\sin x$ est solution de $(E_4)$ \\
$\Longleftrightarrow \forall x\in\R,\ y_0'(x)+y_0(x)=x-e^x+\cos x$ \\
$\Longleftrightarrow \forall x\in\R,\ ax+a+b+2\mu e^x+(\alpha+\beta)\cos x+(-\alpha +\beta)\sin x=x-e^x+\cos x$ 

\noindent Ainsi, en identifiant les coefficients, on voit que $$y_0(x)=x-1-\frac{1}{2}e^x+\frac{1}{2}\cos x+\frac{1}{2}\sin x$$ convient.

Les solutions de $(E_4)$ sont obtenues en faisant la somme de cette solution particulière et des solutions de l'équation homogène:
$$y(x)=x-1-\frac{1}{2}e^x+\frac{1}{2}\cos x+\frac{1}{2}\sin x+\lambda e^{-x}\quad (x\in\R)$$
où $\lambda$ est un paramètre réel.
\end{enumerate}
\fincorrection
\correction{006992}
Une fonction $f:[0;1]\to\R$ convient si et seulement si 
\begin{itemize}
\item $f$ est dérivable
\item $f$ est solution de $y'+y=c$ 
\item $f$ vérifie $f(0)+f(1)=c$ (où $c$ est un réel quelconque)
\end{itemize}
Or les solutions de l'équation différentielle $y'+y=c$ sont exactement les 
$f:x\mapsto \lambda e^{-x}+c$, où $\lambda\in\R$ (en effet, on voit facilement que 
la fonction constante égale à $c$ est une solution particulière de $y'+y=c$). \'Evidemment ces fonctions sont dérivables, et $f(0)+f(1)=\lambda(1+e^{-1})+2c$, donc la troisième condition est satisfaite si et seulement si $-\lambda(1+e^{-1})=c$.

Ainsi les solutions du problème sont exactement les 
$$f(x)=\lambda(e^{-x}-1-e^{-1})$$
pour $\lambda\in\R$.
\fincorrection
\correction{006993}\
\begin{enumerate}
\item Comme le coefficient de $y'$ ne s'annule pas, on peut réécrire l'équation sous la forme
$$y'+\frac{2x}{x^2+1}y=\frac{3x^2+1}{x^2+1}$$
  \begin{enumerate}
    \item Les solutions de l'équation homogène associée sont les $y(x)=\lambda e^{A(x)}$, 
    où $A$ est une primitive de $a(x)=-\frac{2x}{x^2+1}$ et $\lambda\in\R$. 
    Puisque $a(x)$ est de la forme $-\frac{u'}{u}$ avec $u>0$, on peut choisir $A(x)=-\ln(u(x))$ 
    où $u(x)=x^2+1$. Les solutions sont donc les 
    $\displaystyle y(x)=\lambda e^{-\ln(x^2+1)}=\frac{\lambda}{x^2+1}$.
    
    \item Il suffit ensuite de trouver une solution particulière de l'équation 
    avec second membre: on remarque  que $y_0(x)=x$ convient.

    \item Les solutions sont obtenues en faisant la somme:  
    $$y(x)=x+\frac{\lambda}{x^2+1}\quad (x\in\R)$$
    où $\lambda$ est un paramètre réel.
    
    \item $y(0)=3$ si et seulement si $\lambda=3$. La solution cherchée est donc
    $y(x)=x+\frac{3}{x^2+1}$.
  \end{enumerate}
  
  Voici les courbes intégrales pour $\lambda=-1,0,\ldots,5$.
  
  \begin{center}
  \begin{tikzpicture}[scale=1]

 \draw[->,>=latex,thick,gray] (-3,0)--(4,0);
 \draw[->,>=latex,thick,gray] (0,-3)--(0,+5.5);

\foreach \k in {-1,0,1,2,3,4,5}{
 \draw [thick, color=red!60,samples=100,smooth, domain=-3:4] plot(\x,{\x + \k/(1+\x*\x)});
}
\draw [very thick, color=red,samples=100,smooth, domain=-3:4] plot(\x,{\x + 3/(1+\x*\x)});

 \node[below right,gray] at (0,0) {$0$};
\node[below,gray] at (1,0) {$1$};
\node[left,gray] at (0,1) {$1$};
\end{tikzpicture}
\end{center}
  
\item On commence par remarquer que $y_0(x)=\cos x$ est une solution particulière. Pour l'équation homogène: sur l'intervalle considéré, le coefficient de $y'$ ne s'annule pas, et l'équation se réécrit
$$y'-\frac{\cos x}{\sin x}y=0$$
Les solutions sont les $y(x)=\lambda e^{A(x)}$, où $\lambda\in\R$ et $A$ est une primitive de $a(x)=\frac{\cos x}{\sin x}$. Puisque $a(x)$ est de la forme $\frac{u'}{u}$ avec $u>0$, on peut choisir $A(x)=\ln(u(x))$ avec $u(x)=\sin x$. Les solutions de l'équation sont donc les $\displaystyle y(x)=\lambda e^{\ln(\sin x)}=\lambda\sin x$.

Finalement, les solutions de l'équation sont les 
$$y(x)=\cos x+\lambda\sin x$$
où $\lambda$ est un paramètre réel.

  \item On a 
  $$y(\frac\pi4) = 1 
  \iff \cos\frac\pi4 +\lambda\sin \frac\pi4 = 1
  \iff \frac{\sqrt 2}{2}(1+\lambda) = 1
  \iff \lambda = \frac{2}{\sqrt 2}-1$$
  La solution cherchée est $y(x)=\cos x+\left(\frac{2}{\sqrt 2}-1\right)\sin x$
\end{enumerate}

  Voici les courbes intégrales pour $\lambda=-2,-1,0,\ldots,4$
  et $\frac{2}{\sqrt 2}-1$ (en gras).
\begin{center}
\begin{tikzpicture}[scale=1]

 \draw[->,>=latex,thick,gray] (-0.5,0)--(4,0);
 \draw[->,>=latex,thick,gray] (0,-2.5)--(0,+4.5);

\foreach \k in {-2,-1,0,1,2,3,4}{
 \draw [thick, color=red!60,samples=100,smooth, domain=0.0:3.14] plot(\x,{cos(\x r) + \k*sin(\x r)});
}
\def\k{0.4142}
 \draw [ultra thick, color=red,samples=100,smooth, domain=0.0:3.14] plot(\x,{cos(\x r) + \k*sin(\x r)});

 \node[below right,gray] at (0,0) {$0$};
\node[below,gray] at (1,0) {$1$};
\node[below right,gray] at (3.14,0) {$\pi$};
\node[left,gray] at (0,1) {$1$};
\end{tikzpicture}
\end{center}

\fincorrection
\correction{006994}\ 
\begin{enumerate}
  \item $y' - (2x - \frac{1}{x})y = 1$ sur $]0;+\infty[$
  \begin{enumerate}
    \item  \textbf{Résolution de l'équation homogène $y' - (2x - \frac{1}{x})y = 0$.}
    
    Une primitive de $a(x) = 2x - \frac1x$ est $A(x) = x^2 - \ln x$,
    donc les solutions de l'équation homogène sont les $y(x) = \lambda \exp(x^2 - \ln x) = \lambda \frac1x\exp(x^2)$,
    pour $\lambda$ une constante réelle quelconque.
    
    
    \item \textbf{Recherche d'une solution particulière.}
    
    Nous allons utiliser la méthode de variation de la constante pour trouver une solution particulière 
    à l'équation $y' - (2x - \frac{1}{x})y = 1$.
    On cherche une telle solution sous la forme $y_0(x) = \lambda(x) \frac1x\exp(x^2)$ 
    où $x \mapsto \lambda(x)$ est maintenant une fonction.
    
    On calcule d'abord 
    $$y_0'(x) = \lambda'(x) \frac1x\exp(x^2) + \lambda(x) \left(-\frac{1}{x^2}+2\right) \exp(x^2)$$
    
    Maintenant :
    \begin{align*}
          & y_0 \quad \text{ est  solution de }  y' - (2x + \frac{1}{x})y = 1 \\
    \iff & y_0' - (2x - \frac{1}{x})y_0 = 1 \\
    \iff& \lambda'(x) x\exp(x^2) + \lambda(x) \left(-\frac{1}{x^2}+2\right) \exp(x^2)
        - (2x - \frac{1}{x})\lambda(x) \frac1x\exp(x^2) = 1 \\
    \iff& \lambda'(x) \frac1x\exp(x^2) = 1 \qquad \text{cela doit se simplifier !}\\
    \iff& \lambda'(x) = x\exp(-x^2)
    \end{align*}
    
    Ainsi on peut prendre $\lambda(x) = -\frac12\exp(x^2)$, ce qui fournit la solution particulière :
    $$y_0(x) = \lambda(x) \frac1x\exp(x^2) = -\frac12\exp(-x^2)\frac1x\exp(x^2) = -\frac1{2x}$$
    
    Pour se rassurer, on n'oublie pas de vérifier que c'est bien une solution !
    
    \item \textbf{Solution générale.}    
    
    L'ensemble des solutions s'obtient par la somme de la solution particulière avec les solutions de l'équation
    homogène. Autrement dit, les solutions sont les :
    $$y(x) = -\frac1{2x} + \lambda \frac1x\exp(x^2)\qquad (\lambda\in\Rr).$$
    
  \end{enumerate}  

  \item $y'-y = x^k \exp(x)$ sur $\R$, avec $k \in \Nn$
  \begin{enumerate}
    \item  \textbf{Résolution de l'équation homogène $y' - y = 0$.}
    
    Les solutions de l'équation homogène sont les $y(x) = \lambda \exp(x)$, $\lambda \in \Rr$.
        
    \item \textbf{Recherche d'une solution particulière.}
    
    On cherche une solution particulière sous la forme $y_0(x) = \lambda(x) \exp(x)$ 
    où $x \mapsto \lambda(x)$ est maintenant une fonction.
    
    Comme $y_0'(x) = \lambda'(x)\exp(x) + \lambda(x) \exp(x)$
    alors

    \begin{align*}
          & y_0 \quad \text{ est  solution de }  y' - y = x^k \exp(x) \\
    \iff& \lambda'(x)\exp(x) + \lambda(x) \exp(x) - \lambda(x) \exp(x) = x^k \exp(x)\\
    \iff& \lambda'(x)\exp(x) = x^k \exp(x)\\
    \iff& \lambda'(x) = x^k
    \end{align*}
    
    On fixe $\lambda(x) = \frac{x^{k+1}}{k+1}$, ce qui conduit à la solution particulière :
    $$y_0(x) = \frac{x^{k+1}}{k+1} \exp(x)$$
    
    \item \textbf{Solution générale.}    
    
    L'ensemble des solutions est formé des 
    $$y(x) =\frac{x^{k+1}}{k+1} \exp(x) + \lambda \exp(x)\qquad (\lambda\in\Rr).$$
    
  \end{enumerate}  
  
  
\item $x(1+\ln^2(x))y'+2\ln(x)y=1$ sur $]0;+\infty[$

Le coefficient de $y'$ ne s'annule pas sur $]0;+\infty[$, l'équation peut donc se mettre sous la forme
$$y'+\frac{2\ln x}{x(1+\ln^2(x))}y=\frac{1}{x(1+\ln^2(x))}$$
  \begin{enumerate}
  
    \item Les solutions de l'équation homogène associée sont les $y(x)=\lambda e^{A(x)}$, 
    où $A$ est une primitive de $a(x)=-\frac{2\ln x}{x(1+\ln^2(x))}$ et $\lambda\in\R$. 
    On peut donc choisir $A(x)=-\ln(u(x))$ avec $u(x)=1+\ln^2(x)$. Les solutions 
    de l'équation sont les $\displaystyle y(x)=\lambda e^{-\ln(1+\ln^2(x))}=\frac{\lambda}{1+\ln^2(x)}$.
    
    \item Utilisons la méthode de variation de la constante pour trouver une solution 
    particulière de l'équation avec second membre. On cherche $y_0(x)=\frac{\lambda(x)}{1+\ln^2(x)}$, 
    avec $\lambda$ une fonction dérivable. Or $z(x)=\frac{1}{1+\ln^2(x)}$ est solution 
    de l'équation homogène et $y_0(x)=\lambda(x)z(x)$:
    
    \begin{align*}
          & y_0 \quad \text{ est  solution} \\
    \iff&  y_0'+\frac{2\ln x}{x(1+\ln^2(x))}y_0  =  \frac{1}{x(1+\ln^2(x))}\\
    \iff& \lambda'(x)z(x)+\lambda(x)\underbrace{\left[z'(x)+\frac{2\ln x}{x(1+\ln^2(x))}z(x)\right]}_{=0} = \frac{1}{x(1+\ln^2(x))} \\
    \iff& \frac{\lambda'(x)}{1+\ln^2(x)} = \frac{1}{x(1+\ln^2(x))}\\
    \iff& \lambda'(x)=\frac{1}{x}
    \end{align*}    
    On peut donc choisir $\lambda(x)=\ln x$, ce qui donne la solution particulière 
$y_0(x)=\frac{\ln x}{1+\ln^2(x)}$.
     

    \item Les solutions sont obtenues en faisant la somme de cette solution particulière et 
    des solutions de l'équation homogène: ce sont les 
$$y(x)=\frac{\ln x +\lambda}{1+\ln^2(x)}\qquad (\lambda\in\R)$$
où $\lambda$ est un paramètre réel.
  \end{enumerate}
Remarque: le choix d'une primitive de $\lambda'$ se fait à constante additive près. 
Si on avait choisi par exemple $\lambda(x)=\ln x + 1$, la solution particulière aurait été différente, mais les solutions de l'équation avec second membre auraient été les 
$$y(x)=\frac{\ln x +1+\lambda}{1+\ln^2(x)}\quad (x\in\R)$$
Quitte à poser $\lambda_0'=1+\lambda$, ce sont évidemment les m\^emes que celles trouvées précédemment!
\end{enumerate}
\fincorrection
\correction{006995}\ 
\begin{enumerate}
\item L'équation différentielle $y'-e^xe^y=0$ est à variables séparées: 
en effet, en divisant par $e^{y}$, on obtient $-y'e^{-y}=-e^x$. 
Le terme de gauche est la dérivée de $e^{-y}$ ($y$ est une fonction de $x$), 
celui de droite est la dérivée de $x\mapsto -e^x$: 
$$\frac{\dd e^{-y}}{\dd x} = \frac{\dd (-e^x)}{\dd x}$$
Les dérivées étant égales, cela implique que les deux fonctions sont égales à 
une constante additive près: ainsi $y$ est solution sur $I$ si et seulement si 
elle est dérivable sur $I$ et $\exists\,c\in\R,\ \forall x\in I$, $e^{-y}=-e^x+c$. 
\`A $c$ fixé, cette égalité n'est possible que si $-e^x+c>0$, 
c'est-à-dire si $c>0$ et $x<\ln c$. On obtient ainsi les solutions:
$$y_c(x)=-\ln(c-e^x)\quad (\text{pour } x\in I_c=]-\infty;\ln c[)$$
où $c$ est un paramètre réel strictement positif.

Pour que l'une des courbes intégrales passe par l'origine, il faut qu'il existe $c>0$ tel que $0\in I_c$ et $y_c(0)=0$: autrement dit, $c>1$ et $c-1=1$. Il s'agit donc de $y_2:x\mapsto-\ln(2-e^x)$, la courbe intégrale cherchée est son graphe, au-dessus de l'intervalle $I_2=]-\infty;\ln 2[$. Sa tangente en l'origine a pour pente $y_2'(0)=e^0e^{y(0)}=1$, c'est la première bissectrice. Comme par construction $y_2'$ est à valeurs strictement positives, la fonction $y_2$ est strictement croissante.

\begin{tikzpicture}[scale=1]

 \draw[->,>=latex,thick,gray] (-3,0)--(2,0);
 \draw[->,>=latex,thick,gray] (0,-1.5)--(0,+3);

 \draw [very thick, color=red,samples=200,smooth, domain=-3:0.66] plot(\x,{-ln(2-exp(\x))})
node[right]{$y_2(x)$};

\draw[very thick, blue] (-1,-1)--(+1,+1);

 \node[below right,gray] at (0,0) {$0$};
\node[below,gray] at (1,0) {$1$};
\node[left,gray] at (0,1) {$1$};


\begin{scope} [xshift=6cm]
 \draw[->,>=latex,thick,gray] (-3,0)--(2,0);
 \draw[->,>=latex,thick,gray] (0,-1.5)--(0,+3);

 \draw [very thick, color=red,samples=200,smooth, domain=-3:0.95] plot(\x,{-\x - ln(1-\x)})
node[right]{$y_1(x)$};

\draw[very thick, blue] (-2,0)--(+1.5,0);

 \node[below right,gray] at (0,0) {$0$};
\node[below,gray] at (1,0) {$1$};
\node[left,gray] at (0,1) {$1$};
\end{scope}

\end{tikzpicture}

\item Posons $z(x)=x+y(x)$: $z$ a le m\^eme domaine de définition que $y$ et est dérivable 
si et seulement si $y$ l'est. En remplaçant $y(x)$ par $z(x)-x$ dans l'équation 
différentielle $y'-e^xe^y=-1$, on obtient $z'-e^z=0$, c'est-à-dire $z'e^{-z}=1$. 
Il s'agit de nouveau d'une équation à variables séparées: en intégrant cette égalité, 
on obtient que $z$ est solution sur $J$ si et seulement si elle est dérivable sur 
$J$ et $\exists\,c\in\R,\ \forall x\in J$, $e^{-z}=-x+c$. 
\`A $c$ fixé, cette égalité n'est possible que si $c>x$. On obtient ainsi les solutions:
$$y_c(x)=z_c(x)-x=-x-\ln(c-x)\quad (\text{pour } x\in J_c=]-\infty;c[)$$
où $c$ est un paramètre réel.

Pour que l'une des courbes intégrales passe par l'origine, 
il faut qu'il existe $c\in\R$ tel que $0\in J_c$ et $y_c(0)=0$: 
autrement dit, $c>0$ et $c=1$. Il s'agit donc de $y_1:x\mapsto-x-\ln(1-x)$, 
la courbe intégrale cherchée est son graphe, au-dessus de l'intervalle $J_1=]-\infty;1[$. 
Sa tangente en l'origine a pour pente $y_1'(0)=e^0e^{y(0)}-1=0$: elle est horizontale. 
\end{enumerate}
\fincorrection
\correction{006996}\
\begin{enumerate}
\item $x^2y'-y = 0$ $(E_1)$


Pour se ramener à l'étude d'une équation différentielle de la forme $y'+ay=b$, 
on résout d'abord sur les intervalles où le coefficient de $y'$ ne s'annule pas: 
on se place donc sur $]-\infty;0[$ ou $]0;+\infty[$.

\begin{enumerate}
 \item \textbf{Résolution sur $]-\infty;0[$ ou $]0;+\infty[$.}
 
Sur chacun de ces intervalles, l'équation différentielle se réécrit 
$$y'-\frac{1}{x^2}y=0$$ 
qui est une équation différentielle linéaire homogène d'ordre 1 à coefficients non constants. Ses solutions sont de la forme $y(x)=\lambda e^{-1/x}$ (en effet, sur $]-\infty;0[$ ou $]0;+\infty[$, une primitive de $\frac{1}{x^2}$ est $\frac{-1}{x}$).

  \item \textbf{Recollement en $0$.}
  
Une solution $y$ de $(E_1)$ sur $\R$ doit \^etre solution sur $]-\infty;0[$ et $]0;+\infty[$, il existe donc $\lambda_+,\ \lambda_-\in\R$ tels que 
$$y(x)=\left\lbrace\begin{array}{l}
\lambda_+e^{-1/x}\ \mathrm{si}\ x>0\\
\lambda_-e^{-1/x}\ \mathrm{si}\ x<0
\end{array}\right.$$
Il reste à voir si l'on peut recoller les deux expressions pour obtenir une solution sur $\R$: autrement dit, pour quels choix de $\lambda_+,\ \lambda_-$ la fonction $y$ se prolonge-t-elle en 0 en une fonction dérivable vérifiant $(E_1)$? 
\begin{itemize}
\item $e^{-1/x}\xrightarrow[x\to\,0^-]{}+\infty$ et $e^{-1/x}\xrightarrow[x\to\,0^+]{}0$, donc $y$ est prolongeable par continuité en 0 si et seulement si \fbox{$\lambda_-=0$}. On peut alors poser \fbox{$y(0)=0$}, quel que soit le choix de $\lambda_+$.
\item Pour voir si la fonction ainsi prolongée est dérivable en 0, on étudie son taux d'accroissement:
$$\left\{\begin{array}{l}
\text{pour}\ x>0,\ \frac{y(x)-y(0)}{x-0}=\frac{\lambda_+e^{-1/x}}{x}=-\lambda_+(\frac{-1}{x})e^{-1/x}\xrightarrow[x\to\,0^+]{}0\\
\ \\
\text{pour}\ x<0,\ \frac{y(x)-y(0)}{x-0}=0\xrightarrow[x\to\,0^-]{}0
\end{array}\right.$$
Ainsi la fonction $y$ est dérivable en 0 et \fbox{$y'(0)=0$}.
\item Par construction, l'équation différentielle $(E_1)$ est satisfaite sur $\R^*$. 
Vérifions qu'elle est également satisfaite au point $x=0$: $0^2 \cdot y'(0)-y(0)=-y(0)=0$.
\end{itemize}

  \item \textbf{Conclusion.}
  
Finalement, les solutions sur $\R$ sont exactement les fonctions suivantes: 
$$y(x)=\left\lbrace\begin{array}{l}
\lambda e^{-1/x}\quad \text{si}\ x>0\\
0\qquad \text{si}\ x\le 0
\end{array}\right. \qquad (\lambda\in\R)$$
  \end{enumerate}

\item $xy'+y-1 = 0$ $(E_2)$

Pour se ramener à l'étude d'une équation différentielle de la forme 
$y'+ay=b$, on résout d'abord sur les intervalles où le coefficient 
de $y'$ ne s'annule pas: on se place donc sur $I=]-\infty;0[$ ou $I=]0;+\infty[$.

  \begin{enumerate}
  \item \textbf{Résolution sur $I$.}
  
  Sur l'intervalle $I$, l'équation différentielle se réécrit 
$$y'+\frac{1}{x}y=\frac{1}{x}$$ 
qui est une équation différentielle linéaire d'ordre 1 à 
coefficients non constants.

  \begin{itemize}
    \item Pour l'équation homogène $y'+\frac{1}{x}y=0$
 une primitive de $-\frac{1}{x}$ sur $I$, est $-\ln|x|$.
Les solutions de l'équation homogène sont donc les $\lambda e^{-\ln|x|}=\lambda\frac{1}{|x|}$. 
Quitte à changer $\lambda$ en $-\lambda$ si $I=]-\infty;0[$, on peut écrire les solutions de 
l'équation homogène sous la forme $y(x)=\lambda \frac1x$.
    
    \item Pour trouver les solutions de 
l'équation avec second membre, on applique la méthode de variation de la constante 
en cherchant $y(x)=\lambda(x) \frac 1x$: en remplaçant, on voit que $y$ est solution sur $I$ 
si et seulement si $\lambda'(x)=1$. En intégrant, on obtient $\lambda(x)=x$. 
Une solution particulière en donc $y_0(x) = 1$.
    
    \item Sur $I$ les solutions de $(E_2)$ sont les $y(x)=1+\frac{\lambda}{x}$ où $\lambda$ est un paramètre réel.
  \end{itemize}

 
  \item \textbf{Recollement en $0$.}
  
Une solution $y$ de $(E_2)$ sur $\R$ doit \^etre solution sur $]-\infty;0[$ et $]0;+\infty[$, 
il existe donc $\lambda_+,\ \lambda_-\in\R$ tels que 
$$y(x)=\left\lbrace\begin{array}{l}
1+\frac{\lambda_+}{x}\quad \text{si}\ x>0\\
1+\frac{\lambda_-}{x}\quad \text{si}\ x<0
\end{array}\right.$$
Il reste à voir si l'on peut recoller les deux expressions pour obtenir 
une solution sur $\R$. On voit tout de suite que $y$ a une limite (finie) 
en $0$ si et seulement si \fbox{$\lambda_+=\lambda_-=0$}. 
Dans ce cas, on peut alors poser \fbox{$y(0)=1$} et $y$ est la fonction constante 
égale à 1, qui est bien s\^ur dérivable sur $\R$. De plus, $(E_2)$ est bien satisfaite au point $x=0$.

  \item \textbf{Conclusion.}
  
Finalement, $(E_2)$ admet sur $\R$ une unique solution, qui est la fonction constante égale à $1$.
  \end{enumerate}
\end{enumerate}
\fincorrection
\correction{007001}
\begin{enumerate}
  \item \textbf{\'Equation de Bernoulli}
  \begin{enumerate}
    \item On suppose qu'une solution $y$ ne s'annule pas.
    On divise l'équation $y'+a(x)y+b(x)y^n = 0$ par $y^n$, ce qui donne
    $$\frac{y'}{y^n} +a(x) \frac{1}{y^{n-1}} + b(x) = 0.$$
    On pose $z(x) = \frac{1}{y^{n-1}}$ et donc $z'(x) = (1-n)\frac{y'}{y^{n}}$.
    L'équation de Bernoulli devient une équation différentielle linéaire :
    $$\tfrac{1}{1-n} z' + a(x) z +b(x) = 0$$
    
    
    \item \'Equation $xy'+y-xy^3 = 0$.
    
    Cherchons les solutions $y$ qui ne s'annulent pas. On peut alors diviser par $y^3$ pour obtenir :
    $$x \frac{y'}{y^3} + \frac{1}{y^2} - x=0$$
    On pose $z(x) = \frac{1}{y^2(x)}$, et donc $z'(x) = -2\frac{y'(x)}{y(x)^3}$.
    L'équation différentielle s'exprime alors $\frac{-1}{2} xz'+z -x = 0$, c'est-à-dire :
    $$xz'-2z = -2x.$$
    Les solutions  sur $\Rr$ de cette dernière équation sont les 
    $$z(x) =\begin{cases} \lambda_+ x^2+2x\ \text{si}\ x\ge 0 \cr\lambda_- x^2+2x\ 
    \text{si}\ x< 0 \end{cases} ,\quad\lambda_+,\lambda_- \in \Rr$$
    
    Comme on a posé $z(x) = \frac{1}{y^2(x)}$, on se retreint à un intervalle $I$ sur lequel $z(x)>0$: nécessairement $0\notin I$, donc on considère $z(x)=\lambda x^2+2x$, qui est strictement positif sur $I_\lambda$ où
$$I_\lambda=\begin{cases}
]0;+\infty[\quad\text{si}\ \lambda=0\\
\left]0;-\frac{2}{\lambda}\right[\quad\text{si}\ \lambda<0\\
\left]-\infty;-\frac{2}{\lambda}\right[\ \mathrm{ou}\ \left]0;+\infty\right[\quad\text{si}\ \lambda>0              
            \end{cases}
$$
On a $(y(x))^2 = \frac{1}{z(x)}$ pour tout $x\in I_\lambda$ et donc  $y(x) = \epsilon(x)\frac{1}{\sqrt{z(x)}}$, 
où $\epsilon(x)=\pm 1$. Or $y$ est continue sur l'intervalle $I_\lambda$, et 
ne s'annule pas par hypothèse: d'après le théorème des valeurs intermédiaires, $y$ 
ne peut pas prendre à la fois des valeurs strictement positives et des valeurs 
strictement négatives, donc $\epsilon(x)$ est soit constant égal à $1$, soit constant égal à $-1$.
    Ainsi les solutions cherchées sont les :
    $$y(x) = \frac{1}{\sqrt{\lambda x^2 + 2x}}\ \mathrm{ou}\ y(x) = \frac{-1}{\sqrt{\lambda x^2 + 2x}}\quad\mathrm{sur}\ I_\lambda\qquad (\lambda\in \Rr)$$
    
    Noter que la solution nulle est aussi solution.
  \end{enumerate}
  \item \textbf{\'Equation de Riccati}
  \begin{enumerate}
    \item Soit $y_0$ une solution de $y'+ a(x)y +b(x)y^2 = c(x)$.
    Posons $u(x) = y(x)-y_0(x)$, donc $y = u +y_0$. L'équation devient :
    $$u'+y_0' + a(x) (u+y_0)  + b(x) (u^2+2uy_0 + y_0^2) = c(x)$$
    Comme $y_0$ est une solution particulière alors
    $$y_0'+ a(x)y_0 +b(x)y_0^2 = c(x)$$
    Et donc l'équation se simplifie en :
    $$u' + \big(a(x)+2y_0(x)b(x)\big) u + b(x) u^2 = 0$$
    qui est une équation du type Bernoulli.
    
    \item \'Equation $x^2(y'+y^2) = xy-1$.
    
    \begin{itemize}
      \item Après division par $x^2$ c'est bien une équation de Riccati sur $I=]-\infty;0[$ ou $I=]0;+\infty[$.
    
      \item $y_0 = \frac{1}{x}$ est bien une solution particulière.
      
      \item On a $u(x) = y(x)-y_0(x)$ et donc $y = u + \frac1x$.
      L'équation $x^2(y'+y^2) = xy-1$ devient 
      $$x^2 \left( u'-\frac{1}{x^2} + u^2 + 2\frac{u}{x} + \frac{1}{x^2} \right) = x\left(u+\frac1x\right)-1$$
      qui se simplifie en
      $$x^2 \left( u' + u^2 + 2\frac{u}{x}  \right) = xu$$
      ce qui correspond à l'équation de Bernoulli :
      $$u'+\frac{1}{x}u+u^2 = 0.$$
      
      \item Si $u$ ne s'annule pas, en divisant par $u^2$, cette équation devient 
      $\frac{u'}{u^2} + \frac{1}{x}\frac{1}{u} + 1 =0$.
      On pose $z(x) = \frac{1}{u}$, l'équation devient
      $-z'+\frac1x z+1 =0$.
      Ses solutions sur $I$ sont les $z(x) = \lambda x + x\ln |x|$, $\lambda \in \Rr$.
      Ainsi $u(x) = \frac{1}{z(x)} = \frac1{\lambda x + x\ln |x|}$ mais il y a aussi la solution nulle $u(x)= 0$.
      
      \item Conclusion. Comme $y = u + \frac1x$, on obtient alors des solutions de l'équation de départ sur $]-\infty;0[$ et $]0;+\infty[$:
      $$y(x) = \frac1x \qquad \text{ ou } \qquad y(x) = \frac1x +  \frac1{\lambda x + x\ln |x|} \qquad (\lambda \in \Rr).$$
    \end{itemize}

  \end{enumerate}
  \end{enumerate}  
\fincorrection
\correction{007002}
\ \begin{enumerate}
\item 
Notons $A(x) = \int_0^xe^{t^2}\,\dd t$, une primitive de $e^{x^2}$.
On ne sait pas expliciter cette primitive.
Les solutions de $y'+e^{x^2}y=0$ s'écrivent 
$f(x)=\lambda e^{-A(x)}$. 

Si $x\ge 1$, on a par positivité de l'intégrale $A(x) = \int_0^xe^{t^2}\,\dd t\ge 0$ et 
comme $e^{t^2} \ge 1$ alors 
$$A(x) = \int_0^xe^{t^2}\,\dd t \ge \int_0^x 1 \,\dd t = x$$
En conséquence :
$$0 \le e^{-A(x)} \le e^{-x}$$
Ainsi, 
$$0\le |f(x)|\le |\lambda|e^{-x}$$ et $f(x)\xrightarrow[x\to +\infty]{}0$.

\item Supposons que $y$ vérifie sur $\R$ l'équation, et 
posons $u(x)=y(x)^2+e^{-x^2}y'(x)^2$. La fonction $u$ est à valeurs positives, 
dérivable, et 
$$u'(x) = 2y'(x)y(x)+e^{-x^2}2y''(x)y'(x)-2xe^{-x^2}y'(x)^2$$
en utilisant que $e^{-x^2}y''(x) = -y(x)$ 
(car $y$ est solution de l'équation différentielle) on obtient :
$$u'(x)=-2xe^{-x^2}y'^2.$$ 

Ainsi la fonction $u$ est croissante sur $]-\infty;0[$ 
et décroissante sur $]0;+\infty[$, donc pour tout $x\in\R$, $0 \le u(x)\le u(0)$. 
Or $y^2(x)\le u(x)$ par construction, donc
$$\forall x\in\R,\qquad |y(x)|\le \sqrt{u(0)}=\sqrt{y(0)^2+y'(0)^2}$$
\end{enumerate}
\fincorrection
\nocorrection
\nocorrection
\nocorrection
\nocorrection
\correction{000863}
 $ y^{\prime\prime}- 3y^\prime + 2 y = e^x $. Le polyn\^ome caract\'eristique est
$f(r)= (r-1)(r-2)$ et les solutions de l'\'equation homog\`ene
sont donc toutes les fonctions :
$$ y(x) = c_1 e^x +c_2e^{2x}  \hbox{ avec } c_1, c_2 \in \R. $$
On cherche une solution particuli\`ere de la forme $y_p(x)=
P(x)e^x$, on est dans la situation $(\imath\imath)$ la condition
$(*)$ sur $P$ est : $P^{\prime\prime} -P^\prime = 1$, et $P(x)=-x$
convient. Les solutions de l'\'equation sont donc les fonctions :
$$ y(x) = (c_1-x)e^x +c_2e^{2x} \hbox{ avec } c_1, c_2 \in \R. $$
\fincorrection
\correction{000864}
 $ y^{\prime\prime}- y =-6\cos x + 2x\sin x $. Ici $f(r) =(r-1)(r+1)$ et l'\'equation
homog\`ene a pour solutions :
$$ y(x) = c_1 e^x +c_2e^{-x}  \hbox{ avec } c_1, c_2 \in \R .$$
On remarque que la fonction $3\cos x$ v\'erifie l'\'equation :
$y^{\prime\prime}- y =-6\cos x $, il nous reste donc \`a chercher
une solution $y_1$ de l'\'equation $y^{\prime\prime}-y=2x\sin x$,
car $y_p(x)=3\cos x+y_1(x)$ sera une solution de l'\'equation
consid\'r\'ee. Pour cela, on remarque que $2x\sin x = \mathop{\mathrm{Im}}\nolimits
(2xe^{ix})$ et on utilise la m\'ethode d\'ecrite plus haut pour
trouver une solution $z_1$ de l'\'equation : $y^{\prime\prime}- y
=2xe^{ix}$. On cherche $z_1$ sous la forme $P(x)e^{ix}$ o\`u $P$
est un polyn\^ome de degr\'e 1 car $f(i)=-2\not =0$. On a
$f^\prime(i)= 2i$, la condition $(*)$ sur $P$ est donc : $
2iP^\prime(x)-2P(x) = 2x $ ce qui donne apr\`es identification
$P(x) = -x-i$. Alors $y_1(x)=\mathop{\mathrm{Im}}\nolimits((-x+i)e^{ix})=-x\sin x-\cos x$.
Les solutions  sont par cons\'equent les fonctions :
$$y(x) = c_1 e^x +c_2e^{-x}+2\cos x -x\sin x  \hbox{ avec } c_1, c_2 \in \R. $$
Autre m\'ethode pour trouver une solution de $ y^{\prime\prime}- y
= 2x\sin x $  : On la cherche de la forme $y_1(x) = A(x)\sin x +
B(x)\cos x $ o\`u $A,B$ sont des polyn\^omes de degr\'e 1 car $i$
n'est pas racine de l'\'equation caract\'eristique ({\sl danger} :
pour un second membre du type $Q(x)\sin (\beta x)e^{\alpha x}$ la
discussion porte sur $\alpha+i\beta$ et non sur $\alpha$ ou
$\beta$...). On calcule $y_1^\prime$, $y_1^{\prime\prime}$ et on
applique l'\'equation \'etudi\'ee \`a $y_1$ \ldots on obtient la
condition :
$$ (A^{\prime\prime}-A-2B^\prime)\sin x +(B^{\prime\prime}-B-2A^\prime)= 2x\sin x$$
qui sera r\'ealis\'ee si : $ \left\lbrace \begin{array}{c}
                    A^{\prime\prime}-A-2B^\prime = 2x \\
                    B^{\prime\prime}-B-2A^\prime = 0
                \end{array} \right. . $ \\
On \'ecrit : $A(x)= ax+b$ et $B(x)=cx+d$, apr\`es identification
on obtient : $a=d=-1$, $b=c=0$, ce qui d\'etermine $y_1$.
\fincorrection
\correction{000865}
% $4y^{\prime\prime}+4y^\prime + 5y = \sin x e^{-x/2} $. L'\'equation
%caract\'eristique a 2 racines complexes $r_1 = -1/2+i$ et
%$r_2=\overline{r_1}$ et les solutions de l'\'equation homog\`ene
%sont :
%$$ y(x) = e^{-x/2}(c_1\cos x +c_2\sin x) \hbox{ avec } c_1,c_2\in \R $$
%On a $\sin x e^{-x/2} =\mathop{\mathrm{Im}}\nolimits(e^{(-1/2+i)x})$, on commence donc par
%chercher une solution $z_p$ de l'\'equation  avec le nouveau
%second membre $e^{(-1/2+i)x}$.Comme $-1/2+i$ est racine de
%l'\'equation caract\'eristique, on cherchera $
%z_p(x)=P(x)e^{(-1/2+i)x}$ avec $P$ de degr\'e 1. Par cons\'equent
%la condition $(*)$ sur $P$ :
%$$ 4P^{\prime\prime}+f^\prime(-1/2+i)P^\prime+f(-1/2+i)P = 1$$
%s'\'ecrit ici : $8iP^\prime =1$ ( $P^{\prime\prime} = 0$,
%$f(-1/2+i)=0$ et $f^\prime(-1/2+i)=8i$), on peut donc prendre
%$P(x)=-i/8x$ et $z_p(x)=-i/8xe^{(-1/2+i)x}$, par cons\'equent sa
%partie imaginaire $y_p(x)=\mathop{\mathrm{Im}}\nolimits(-i/8xe^{(-1/2+i)x})= 1/8 x\sin
%xe^{-x/2}$ est une solution de notre \'equation. Les solutions
%sont donc toutes les fonctions de la forme :
%$$y(x) = e^{-x/2}(c_1\cos x +(c_2+1/8x)\sin x) \hbox{ avec } c_1,c_2\in \R.$$

La solution générale est de la forme
$$y(x) = K_1 \cos(x) e^{-x/2} + K_2 \sin(x) e^{-x/2} - \frac18 x e^{-x/2} \cos(x)$$
($K_1$ et $K_2$ constantes réelles) et les conditions initiales donnent $K_1=0$, $K_2=1/8$.
\fincorrection
\correction{000866}
\begin{enumerate}
    \item Le polyn\^ome caract\'eristique associ\'e \`a $E$ est : $p(x) = x^2+2x+4$ ; son discriminant
est $\Delta = -12$ et il a pour racines les 2 nombres complexes $
-1+i\sqrt{3}$ et $-1-i\sqrt3$. Les solutions de l'\'equation
homog\`ene sont donc toutes fonctions :
$$ y(x) = e^{-x}(a\cos{\sqrt3 x} + b\sin{\sqrt3 x})$$ obtenues lorsque $a,b$ d\'ecrivent $\R$.
    \item Le second membre est de la forme $ e^{\lambda x}Q(x) $ avec $\lambda = 1$ et $Q(x)=x$. On
cherchera une solution de l'\'equation sous la forme : $ y_p(x)=
R(x)e^x$ avec $R$ polyn\^ome de degr\'e \'egal \`a celui de $Q$
puisque $p(1) \not =0$. On pose donc $R(x) = ax+b$. On a
 $$ y_p^{\prime\prime}(x) + 2y_p^{\prime}(x) +4y_p(x) = (7ax+7b+4a)e^x.$$
  Donc $y_p$ est solution
si et seulement si $ 7ax + 7a+4b = x$. On trouve apr\`es
identification des coefficients :
 $$ a=\frac{1}{7} \qquad \hbox{et}\qquad b=\frac{-4}{49}.$$
La fonction $y_p(x)=\frac{1}{7}(x-\frac{4}{7})e^x$ est donc
solution de $E$ et la forme g\'en\'erale des solutions de $E$ est
:
$$ y(x)= e^{-x}(a\cos{\sqrt3 x} + b\sin{\sqrt3 x}) +\frac{1}{7}(x-\frac{4}{7})e^x \; ;\; a,b \in \R.$$

    \item  Soit $h$ une solution de $E$. Les conditions $h(0)=1$, $h(1)=0$ sont r\'ealis\'ees ssi
$$a=\frac{53}{49}\qquad \hbox{et} \qquad b=-\frac{53\cos\sqrt3+3e^2}{49\sin\sqrt3}.$$
    \item
    \begin{enumerate}
        \item  On a : $g^\prime(x) =e^xf^\prime(e^x)$ et
$g^{\prime\prime}(x)=e^xf^\prime(e^x)+e^{2x}f^{\prime\prime}(e^x)$
d'o\`u pour tout $x\in\R$ :
$$g^{\prime\prime}(x)+2g^\prime(x)+4g(x)=e^{2x}f^{\prime\prime}(e^x)+2e^xf^\prime(e^x)+4f(e^x)= e^x\log{e^x}=xe^x$$
donc $g$ est solution de $E$.
        \item  R\'eciproquement pour $f(t)= g(\log t)$ o\`u $g$ est une
solution de $E$ on montre que $f$ est 2 fois d\'erivable et
v\'erifie l'\'equation donn\'ee en 4. Donc les fonctions $f$
recherch\'ees sont de la forme :
$$\frac{1}{t}(a\cos{(\sqrt3\log t)} + b\sin{(\sqrt3\log t)}) +\frac{t}{7}(\log t-\frac{4}{7}) \; ;\; a,b \in \R.$$
    \end{enumerate}
\end{enumerate}
\fincorrection
\nocorrection
\nocorrection
\nocorrection
\nocorrection
\nocorrection
\correction{000872}
\begin{enumerate}
    \item  L'\'equation caract\'eristique $r^2-4r+4=0$ a une racine (double) $r=2$
donc les solutions de l'\'equation homog\`ene sont les fonctions :
             $$ y(x) = (c_1x+c_2)e^{2x} \hbox{ o\`u } c_1,c_2 \in\R. $$
    \item Pour $d(x) = e^{-2x}$ on peut chercher une solution particuli\`ere de la forme : $ y_1(x)
= ae^{-2x} $ car $-2$ n'est pas racine de l'\'equation
caract\'eristique. On a $y_1'(x)= -2e^{-2x}$ et
$y_1''(x)=4ae^{-2x}$. Par cons\'equent $y_1$ est solution si et
seulement si :
             $$\forall x\in\R\quad  (4a -4(-2a)+4a)e^{-2x} = e^{-2x} $$
donc si et seulement si $ a =\frac{1}{16}$. \\ Pour $d(x) =e^{2x}$
on cherche une solution de la forme $y_2(x)=ax^2e^{2x}$, car $2$
est racine double de l'\'equation caract\'eristique. On a $y_2'(x)
= (2ax+2ax^2)e^{2x}$ et
$y_2''(x)=(2a+4ax+4ax+4ax^2)e^{2x}=(4ax^2+8ax+2a)e^{2x}$. Alors
$y_2$ est solution si et seulement si
$$\forall x\in\R\quad (4ax^2+8ax+2a-4(2ax+2ax^2)+4ax^2)e^{2x} =e^{2x} $$
donc si et seulement si $a=\frac{1}{2}$.
    \item  On d\'eduit du principe de superposition que la
fonction $$
y_p(x)=\frac{1}{4}(y_1(x)+y_2(x))=\frac{1}{64}e^{-2x}+\frac{1}{8}x^2e^{2x}$$
est solution de l'\'equation pour le second membre donn\'e dans
cette question, et la forme g\'en\'erale des solutions est alors :
$$ y(x)=(c_1x+c_2)e^{2x}+\frac{1}{64}e^{-2x}+\frac{1}{8}x^2e^{2x} \hbox{ o\`u }c_1,c_2\in\R.$$
\end{enumerate}
\fincorrection
\nocorrection
\nocorrection
\nocorrection
\nocorrection
\nocorrection
\nocorrection
\nocorrection
\correction{000880}
R\'{e}ponse : $\left( \lambda x+\mu \right)
e^{-x}+\frac{e^{x}}{25}\left[ \left( 3x-4\right) \cos x-\left(
4x-2\right) \sin x\right] +\left( \sin x-x\cos x\right) e^{-x}$.
\fincorrection
\correction{000881}
R\'{e}ponse : $\frac{1}{2}\left( -x\cos x+\sin x\right) +\lambda
\cos x+\mu \sinh x.$
\fincorrection
\nocorrection
\nocorrection
\correction{000884}
R\'{e}ponse : $x\rightarrow \frac{\lambda x+\mu
}{\sqrt{1+x^{2}}},\left( \lambda ,\mu \right) \in {\Rr}^{2}.$
\fincorrection
\correction{000885}
R\'{e}ponse : $x\rightarrow \lambda x\sinh x+\mu x\cosh x,\left(
\lambda ,\mu \right) \in {\Rr}^{2}.$
\fincorrection
\correction{005479}
\begin{enumerate}
\item  L'équation caractéristique de l'équation homogène $y''-2y'+2y=0$ est $r^2-2r+2=0$ dont les racines sont $1-i$
et $1+i$. Les solutions de l'équation homogène sont les fonctions de la forme $x\mapsto e^x(\lambda\cos x+\mu\sin
x),\;(\lambda,\mu)\in\Rr^2$. L'équation avec second membre s'écrit

$$y''-2y'+2y=
\frac{x}{4}(e^{(1+i)x}+e^{(-1+i)x}+e^{(1-i)x}+e^{(-1-i)x}).$$

On applique alors le principe de superposition des solutions.

Recherche d'une solution particulière de l'équation $y''-2y'+2y=xe^{(1+i)x}$.

$1+i$ est racine simple de l'équation caractéristique et donc
l'équation précédente admet une solution particulière de la forme $f~:~x\mapsto(ax^2+bx)e^{(1+i)x}$. D'après la formule
de \textsc{Leibniz}~:

\begin{align*}\ensuremath
f''-2f'+2f&=(((1+i)^2(ax^2+bx)+2(1+i)(2ax+b)+2a)\\
 &\;-2((1+i)(ax^2+bx)+(2ax+b))+2(ax^2+bx))e^{(1+i)x}\\
 &=(2(1+i)(2ax+b)+2a-2((2ax+b)))e^{(1+i)x}=(2i(2ax+b)+2a)e^{(1+i)x}\\
 &=(4iax+2a+2ib)e^{(1+i)x}
\end{align*}

puis,

$$f''-2f'+2f=xe^{(1+i)x}\Leftrightarrow4ia=1\;\mbox{et}\;2ib+2a=0\Leftrightarrow a=-\frac{i}{4}\;\mbox{et}\;b=\frac{1}{4}.$$

Une solution particulière de l'équation $y''-2y'+2y=xe^{(1+i)x}$ est $x\mapsto\frac{1}{4}(-ix^2+x)e^{(1+i)x}$. Par
conjugaison, une solution particulière de l'équation $y''-2y'+2y=xe^{(1-i)x}$ est
$x\mapsto\frac{1}{4}(ix^2+x)e^{(1-i)x}$.

Recherche d'une solution particulière de l'équation $y''-2y'+2y=xe^{(-1+i)x}$.

$-1+i$ n'est pas racine de l'équation caractéristique et donc
l'équation précédente admet une solution particulière de la forme $f~:~x\mapsto(ax+b)e^{(-1+i)x}$. D'après la formule
de \textsc{Leibniz}~:

\begin{align*}\ensuremath
f''-2f'+2f&=(((-1+i)^2(ax+b)+2(-1+i)a)-2((-1+i)(ax+b)+a)+2(ax+b))e^{(-1+i)x}\\
 &=((ax+b)(-2i-2(-1+i)+2)+2(-1+i)a-2a)e^{(-1+i)x}\\
 &=(4(1-i)(ax+b)-2(2-i)a)e^{(1+i)x}=(4(1-i)ax-2(2-i)a+4(1-i)b)e^{(1+i)x}
\end{align*}

puis,

\begin{align*}\ensuremath
f''-2f'+2f=xe^{(-1+i)x}&\Leftrightarrow4(1-i)a=1\;\mbox{et}\;4(1-i)b-2(2-i)a=0\\
 &\Leftrightarrow
a=\frac{1+i}{8}\;\mbox{et}\;b=\frac{(2-i)(1+i)}{16(1-i)}=\frac{(3+i)(1+i)}{32}=\frac{1+2i}{16}.
\end{align*}

Une solution particulière de l'équation $y''-2y'+2y=xe^{(-1+i)x}$ est $x\mapsto\frac{1}{16}(2(1+i)x+1+2i)e^{(-1+i)x}$.
Par conjugaison, une solution particulière de l'équation $y''-2y'+2y=xe^{(-1-i)x}$ est
$x\mapsto\frac{1}{16}(2(1-i)x+1-2i)e^{(-1-i)x}$.

Une solution particulière de l'équation $y''-2y'+2y=x\cos x\ch x$ est donc

\begin{align*}\ensuremath
\frac{1}{4}&(2\mbox{Re}(\frac{1}{4}(-ix^2+x)e^{(1+i)x}+\frac{1}{16}(2(1+i)x+1+2i)e^{(-1+i)x})\\
 &=\frac{1}{32}\mbox{Re}(4(-ix^2+x)(\cos x+i\sin x)e^x+(2x+1+2(x+1)i)(\cos x+i\sin x)e^{-x}\\
 &=\frac{1}{32}(4(x\cos x+x^2\sin x)e^x+((2x+1)\cos x-2(x+1)\sin x)e^{-x})
\end{align*}

Les solutions sur $\Rr$ de l'équation proposée sont les fonctions de la forme $x\mapsto(\frac{1}{8}(x\cos x+x^2\sin
x)+\lambda\cos x+\mu\sin x)e^x+((2x+1)\cos x-2(x+1)\sin x)e^{-x})$, $(\lambda,\mu)\in\Rr^2$.

\item  L'équation caractéristique de l'équation homogène $y''+6y'+9y=0$ est $r^2+6r+9=0$ qui admet la racine double
$r=-3$. Les solutions de l'équation homogène sont les fonctions de la forme $x\mapsto
e^{-3x}(\lambda x+\mu),\;(\lambda,\mu)\in\Rr^2$.

$2$ n'est pas racine de l'équation caractéristique et donc
l'équation proposée admet une solution particulière de la forme $f~:~x\mapsto(ax^2+bx+c)e^{2x}$. D'après la formule
de \textsc{Leibniz}~:

\begin{align*}\ensuremath
f''+6f'+9f&=((4(ax^2+bx+c)+4(2ax+b)+2a)+6(2(ax^2+bx+c)+(2ax+b))+9(ax^2+bx+c))e^{2x}\\
 &=(25(ax^2+bx+c)+10(2ax+b)+2a)e^{2x}=(25ax^2+(20a+25b)x+2a+10b+25c)e^{2x}
\end{align*}

puis,

$$f''+6f'+9f=x^2e^{2x}\Leftrightarrow25a=1\;\mbox{et}\;20a+25b=0\;\mbox{et}\;2a+10b+25c=0\Leftrightarrow
a=\frac{1}{25}\;\mbox{et}\;b=-\frac{4}{125}\;\mbox{et}\;c=\frac{6}{625}.$$
Une solution particulière de l'équation $y''+6y'+9y=x^2e^{2x}$ est $x\mapsto\frac{1}{625}(25x^2-20x+6)e^{2x}$.

Les solutions sur $\Rr$ de l'équation proposée sont les fonctions de la forme
$x\mapsto\frac{1}{625}(25x^2-20x+6)e^{2x}+(\lambda x+\mu)e^{-3x}$, $(\lambda,\mu)\in\Rr^2$.

\item  L'équation caractéristique de l'équation homogène $y''-2y'+y=0$ est $r^2-2r+1=0$ qui admet la racine double
$r=1$. Les solutions de l'équation homogène sont les fonctions de la forme $x\mapsto
e^{x}(\lambda x+\mu),\;(\lambda,\mu)\in\Rr^2$.

Le second membre s'écrit $\frac{1}{2}(e^{x}+e^{-x})$. Appliquons le principe de superposition des solutions.

Recherche d'une solution particulière de l'équation $y''-2y'+y=e^{x}$.

$1$ est racine double de l'équation caractéristique et donc
l'équation proposée admet une solution particulière de la forme $f~:~x\mapsto ax^2e^{x}$. D'après la formule
de \textsc{Leibniz}~:

\begin{align*}\ensuremath
f''-2f'+f&=((ax^2+2(2ax)+2a)-2(ax^2+(2ax))+ax^2)e^{2x}=2ae^{x}
\end{align*}

puis,

$$f''-2f'+f=e^{x}\Leftrightarrow2a=1\Leftrightarrow a=\frac{1}{2}.$$
Une solution particulière de l'équation $y''-2y'+y=e^{x}$ est $x\mapsto\frac{1}{2}x^2e^{x}$.

Recherche d'une solution particulière de l'équation $y''-2y'+y=e^{-x}$.

$-1$ n'est pas racine de l'équation caractéristique et donc
l'équation proposée admet une solution particulière de la forme $f~:~x\mapsto ae^{-x}$.

\begin{align*}\ensuremath
f''-2f'+f&=(a+2a+a)e^{-x}=4ae^{-x}
\end{align*}

puis,

$$f''-2f'+f=e^{-x}\Leftrightarrow a=\frac{1}{4}.$$
Une solution particulière de l'équation $y''-2y'+y=e^{-x}$ est $x\mapsto\frac{1}{4}e^{-x}$.

Les solutions sur $\Rr$ de l'équation proposée sont les fonctions de la forme
$x\mapsto(\frac{x^2}{4}+\lambda x+\mu)e^x+\frac{1}{8}e^{-x}$, $(\lambda,\mu)\in\Rr^2$.

\item  Soit $k\in\Rr$. L'équation caractéristique de l'équation homogène $y''-2ky'+(1+k^2)y=0$ est $r^2-2kr+1+k^2=0$
dont le discriminant réduit vaut $-1=i^2$. Cette équation admet donc pour racines $k+i$ et $k-i$. Les solutions de
l'équation homogène sont les fonctions de la forme $x\mapsto e^{kx}(\lambda\cos x+\mu\sin x),\;(\lambda,\mu)\in\Rr^2$.

Le second membre s'écrit $\mbox{Im}(e^{(1+i)x})$. Résolvons donc l'équation $y''-2y'+y=e^{(1+i)x}$.

Si $k\neq1$, $1+i$ n'est pas racine de l'équation caractéristique et donc
l'équation proposée admet une solution particulière de la forme $f~:~x\mapsto ae^{(1+i)x}$. Or,

\begin{align*}\ensuremath
f''-2kf'+(1+k^2)f&=a((1+i)^2-2k(1+i)+1+k^2)e^{(1+i)x}=((k-1)^2-2(k-1)i)ae^{(1+i)x}
\end{align*}

et donc,

$$f''-2kf'+(1+k^2)f=e^{(1+i)x}\Leftrightarrow a=\frac{1}{k-1}\frac{1}{k-1-2i}=\frac{k-1+2i}{(k-1)(k^2-2k+5)}.$$
Une solution particulière de l'équation $y''-2y'+y=e^{(1+i)x}$ est $x\mapsto\frac{k-1-2i}{(k-1)(k^2-2k+5)}e^{(1+i)x}$
et une solution particulière de l'équation $y''-2y'+y=e^{x}\sin x$ est

$$\frac{1}{(k-1)(k^2-2k+5)}\mbox{Im}((k-1-2i)(\cos x+i\sin x)e^x=\frac{1}{(k-1)(k^2-2k+5)}(-2\cos x+(k-1)\sin x)e^x.$$

Si $k\neq1$, les solutions de l'équation proposée sont les fonctions de la forme

$$x\mapsto\frac{1}{(k-1)(k^2-2k+5)}(-2\cos x+(k-1)\sin x)e^x+(\lambda\cos x+\mu\sin x)e^{kx},\;(\lambda,\mu)\in\Rr.$$

\end{enumerate}
\fincorrection
\correction{005480}
\begin{enumerate}
\item  Supposons $y$ deux fois dérivable sur $]0,+\infty[$. La fonction $t\mapsto e^t$ est deux fois dérivable sur
$\Rr$ à valeurs dans $]0,+\infty[$ et la fonction $x\mapsto y(x)$ est deux fois dérivable sur $]0,+\infty[$. Donc,
puisque pour tout réel $t$, $z(t)=y(e^t)$, la fonction $z$ est deux fois dérivable sur $\Rr$ en tant que composée de
fonctions deux fois dérivables.
Réciproquement, supposons que $z$ est deux fois dérivable sur $\Rr$. La fonction $x\mapsto\ln x$ est deux fois
dérivable sur $]0,+\infty[$ à valeurs dans $\Rr$ et la fonction $t\mapsto z(t)$ est deux fois dérivable sur $\Rr$.
Donc, puisque pour tout réel strictement positif $x$, $y(x)=z(\ln x)$, la fonction $y$ est deux fois dérivable sur
$]0,+\infty[$.

\item  Pour $t$ réel, posons donc $x=e^t$ puis, $z(t)=y(x)=y(e^t)$. Alors, $z'(t)=e^ty'(e^t)=xy'(x)$ puis
$z''(t)=e^ty'(e^t)+(e^t)^2y''(e^t)=xy'(x)+x^2y''(x)$. Donc, $xy'(x)=z'(t)$ et $x^2y''(x)=z''(t)-xy'(x)=z''(t)-z'(t)$.

Par suite,

$$ax^2y''(x)+bxy'(x)+cy(x)=a(z''(t)-z'(t))+bz'(t)+cz(t)=az''(t)+(b-a)z'(t)+cz(t).$$

Donc,

$$\forall x>0,\;ax^2y''(x)+bxy'(x)+cy(x)=0\Leftrightarrow\forall t\in\Rr,\;az''(t)+(b-a)z'(t)+cz(t)=0.$$

\item  On applique le 2) avec $a=1$, $b=-1$ et $c=1$. L'équation à résoudre sur $\Rr$ est alors $z''-2z'+z=0$. Les
solutions de cette équation sur $\Rr$ sont les fonctions de la forme $t\mapsto(\lambda t+\mu)e^t$,
$(\lambda,\mu)\in\Rr^2$. Les solutions sur $]0,+\infty[$ de l'équation initiale sont donc les fonctions de la forme
$x\mapsto\lambda x\ln x+\mu x$, $(\lambda,\mu)\in\Rr^2$.

\end{enumerate}
\fincorrection
\correction{006997}
\begin{enumerate}
\item Il s'agit d'une équation homogène du second ordre. L'équation caractéristique associée est $r^2-3r+2=0$, qui admet deux solutions: $r=2$ et $r=1$. Les solutions sont donc les fonctions définies sur $\R$ par $y(x)=\lambda e^{2x}+\mu e^x$ ($\lambda,\mu\in\R$).
\item L'équation caractéristique associée est $r^2+2r+2=0$, qui admet deux solutions: $r=-1+i$ et $r=-1-i$. On sait alors que les solutions sont donc les fonctions définies sur $\R$ par $y(x)=e^{-x}(A\cos x+ B\sin x)$ ($A, B\in\R$). Remarquons que, en utilisant l'expression des fonctions $\cos$ et $\sin$ à l'aide d'exponentielles, ces solutions peuvent aussi s'écrire sous la forme $\lambda e^{(-1+i)x}+\mu e^{(-1-i)x}$ $(\lambda,\mu\in\R)$.
\item L'équation caractéristique est $r^2-2r+1=0$, dont 1 est racine double. 
Les solutions de l'équation homogène sont donc de la forme $(\lambda x+\mu)e^x$. 
\item Les solutions de l'équation homogène sont les $\lambda\cos x+\mu \sin x$. 
Le second membre peut en fait se réécrire $\cos^2 x=1+\cos(2x)$: 
d'après le principe de superposition, on cherche une solution particulière sous 
la forme $a+b\cos(2x)+c\sin(2x)$. En remplaçant, on trouve qu'une telle fonction est solution si $a=1$, $b=-\frac{1}{3}$, $c=0$. Les solutions générales sont donc les $\lambda\cos x+\mu \sin x-\frac{1}{3}\cos(2x)+1$.
\end{enumerate}
\fincorrection
\correction{006998}
L'équation caractéristique associée à l'équation homogène est $r^2-4r+4=0$, 
pour laquelle $r=2$ est racine double. Les solutions de l'équation homogène 
sont donc les $(\lambda x+\mu)e^{2x}$. 

Lorsque $d(x)=e^{-2x}$, on cherche une solution particulière sous la forme $ae^{-2x}$, 
qui convient si $a=\frac{1}{16}$.

Lorsque $d(x)=e^{2x}$, comme 2 est la racine double de l'équation caractéristique, 
on cherche une solution comme le produit de $e^{2x}$ par un polyn\^ome de degré 2. 
Comme on sait déjà que $(\lambda x+\mu)e^{2x}$ est solution de l'équation homogène, 
il est inutile de faire intervenir des termes de degré $1$ et $0$: on cherche 
donc une solution de la forme $ax^2e^{2x}$, qui convient si et seulement si $a=\frac{1}{2}$. 

Puisque $\frac{1}{2}\ch(2x) =\frac{1}{4}(e^{2x}+e^{-2x})$, les solutions générales sont obtenues sous 
la forme $y(x)=\frac{1}{64}e^{-2x}+\frac{1}{8}x^2e^{2x}+(\lambda x+\mu)e^{2x}$.
\fincorrection
\correction{006999}
Les solutions de l'équation homogène sont les $\lambda\cos x+\mu\sin x$. 
En posant $y_1(x)=\cos x$ et $y_2(x)=\sin x$, on va chercher les solutions sous 
la forme $\lambda y_1+\mu y_2$, vérifiant
\begin{eqnarray*}
\left\{\begin{array}{l}
\lambda'y_1+\mu'y_2=0\\
\lambda'y'_1+\mu'y'_2=\mathrm{cotan}\, x
\end{array}\right.
&\Longleftrightarrow&
\left\{\begin{array}{l}
\lambda'\cos x+\mu'\sin x=0\\
\lambda'(-\sin x)+\mu'\cos x=\mathrm{cotan}\, x
\end{array}\right.\\
\ &\Longleftrightarrow&
\left\{\begin{array}{l}
\lambda'(x)=\frac{\begin{array}{|cc|}0&\sin x\\\mathrm{cotan}\, x&\cos x \end{array}}{\begin{array}{|cc|}\cos x&\sin x\\-\sin x&\cos x \end{array}}\\
\ \\
\mu'(x)=\frac{\begin{array}{|cc|}\cos x&0\\-\sin x&\mathrm{cotan}\, x \end{array}}{\begin{array}{|cc|}\cos x&\sin x\\-\sin x&\cos x \end{array}}
\end{array}\right.
\end{eqnarray*}
d'après les formules de Cramer, où $\begin{array}{|cc|}\cos x&\sin x\\-\sin x&\cos x \end{array}=1$. On obtient donc 
$$\left\{\begin{array}{l}
\lambda'(x)=-\cos x\\
\mu'(x)=\frac{\cos^2x}{\sin x}
\end{array}\right.
$$
ce qui donne une primitive $\lambda(x)=-\sin x$. 

Pour $\mu$, on cherche à primitiver 
$\frac{\cos^2x}{\sin x}$ à l'aide du changement de variable $t=\cos x$ (et donc $\dd t= -\sin t \,\dd x$), 
on calcule une primitive
\begin{eqnarray*}
\int \frac{\cos^2x}{\sin x}\,\dd x&=&-\int\frac{t^2}{1-t^2}\,\dd t=t-\int\frac{1}{1-t^2}\,\dd t\\ 
 &=&t+\frac{1}{2}\ln(1-t) -\frac{1}{2}\ln(1-t)= \cos x + \frac{1}{2}\ln(1-\cos x) - \frac{1}{2}\ln(1-\cos x)
\end{eqnarray*}
En remplaçant, les solutions générales sont les 
$$y(x)=c_1\cos x+c_2\sin x +  (-\sin x)\cos x+\left(\cos x + \frac{1}{2}\ln(1-\cos x) - \frac{1}{2}\ln(1-\cos x)\right)\sin x$$
qui se simplifie $y(x)=c_1\cos x+c_2\sin x + \frac{1}{2}\sin x\ln\frac{1-\cos x}{1+\cos x}$, $c_1,c_2\in\Rr$.
\fincorrection
\correction{007000}
\begin{enumerate}
\item Puisqu'on cherche $y$ fonction de $x\in]0;+\infty[$, et que l'application $t\mapsto e^t$ 
est bijective de $\R$ sur $]0;+\infty[$, on peut poser $x=e^t$ et $z(t)=y(e^t)$. 
On a alors $t= \ln x$ et $y(x) = z(\ln x)$.
Ce qui donne :
\begin{eqnarray*}
y(x)&=&z(\ln x)=z(t)\\
y'(x)&=&\frac{1}{x}z'(\ln x)=e^{-t}z'(t)\\
y''(x)&=&-\frac{1}{x^2}z'(\ln x)+\frac{1}{x^2}z''(\ln x)=-e^{-2t}z'(t)+e^{-2t}z''(t)
\end{eqnarray*}
En remplaçant, on obtient donc que 
$$\forall x\in]0;+\infty[,\ x^2y''+xy'+y=0\Longleftrightarrow\forall t\in\R,\ z''(t)+z(t)=0$$
autrement dit, $z(t)=\lambda\cos t+\mu\sin t$ où $\lambda,\mu\in\R$. 
Finalement, les solutions de l'équation de départ sont de la forme 
$$y(x) = z(\ln x) = \lambda\cos(\ln x)+\mu\sin(\ln x)$$ où $\lambda,\mu\in\R$.

\item L'application $t\mapsto \tan t$ étant bijective de 
$]-\frac{\pi}{2};\frac{\pi}{2}[$ sur $\R$, on peut poser 
$x=\tan t$ et $z(t)=y(\tan t)$. On a alors $t = \arctan x$ et ainsi :
\begin{eqnarray*}
y(x)&=&z(\arctan x)=z(t)\\
y'(x)&=&\frac{1}{1+x^2}z'(\arctan x)\\
y''(x)&=&\frac{1}{(1+x^2)^2}\big(z''(\arctan x)-2xz'(\arctan x)\big)
\end{eqnarray*}
En remplaçant, on obtient que $z$ est solution de l'équation différentielle $z''+mz=0$. Pour résoudre cette équation, on doit distinguer trois cas:
\begin{itemize}
  \item $m<0$: alors $z(t)=\lambda e^{\sqrt{-m}t}+\mu e^{-\sqrt{-m}t}$ et donc  
$$y(x)=\lambda e^{\sqrt{-m}\arctan x}+\mu e^{-\sqrt{-m}\arctan x},$$
  
  \item $m=0$: $z''=0$ et donc $z(t)=\lambda t+\mu$ et $y(x)=\lambda\arctan x+\mu$,
  
  \item $m>0$: alors $z(t)=\lambda\cos(\sqrt{m}t)+\mu\sin(\sqrt{m}t)$ et donc 
$$y(x)=\lambda\cos(\sqrt{m}\arctan x)+\mu\sin(\sqrt{m}\arctan x)$$
où $\lambda,\mu\in\R$.
\end{itemize}

\end{enumerate}
\fincorrection
\correction{007003}
\ \begin{enumerate}
\item En faisant le changement de variable $x=e^t$ (donc $t = \ln x$) et en posant 
$z(t)=y(e^t)$ (donc $y(x) = z(\ln x)$, l'équation $x^2y''+y=0$ devient $z''-z'+z=0$, 
dont les solutions sont les 
$z(t)=e^{t/2}\cdot\left( \lambda\cos(\frac{\sqrt{3}}{2}t)+\mu\sin(\frac{\sqrt{3}}{2}t) \right)$,
$\lambda,\mu\in\Rr$. 
Autrement dit, 
$$y(x)=\sqrt{x}\left(\lambda\cos(\frac{\sqrt{3}}{2}\ln x)+\mu\sin(\frac{\sqrt{3}}{2}\ln x)\right)$$
\item Supposons que $f$ convienne: par hypothèse, $f$ est de classe $\mathcal{C}^1$, 
donc $x\mapsto f(\frac{1}{x})$ est de classe $\mathcal{C}^1$ sur $\R^*$ et 
par conséquent $f'$ aussi. Ainsi $f$ est nécessairement de classe 
$\mathcal{C}^2$ sur $\R^*$ (en fait, en itérant le raisonnement, on montrerait 
facilement que $f$ est $\mathcal{C}^\infty$ sur $\R^*$). 


En dérivant l'équation $f'(x)=f\left(\frac{1}{x}\right)$, 
on obtient 
$$f''(x)=-\frac{1}{x^2}f'\left(\frac{1}{x}\right)$$
et en réutilisant l'équation :
$$f''(x)=-\frac{1}{x^2}f'\left(\frac{1}{x}\right)= -\frac{1}{x^2} f(x).$$
Ainsi on obtient que $f$ est solution de $x^2y''+y=0$ sur $\R^*$. 
Nécessairement, il existe $\lambda,\mu\in\Rr$ tels que 
$$\forall x>0,\ f(x)=\sqrt{x}\left(\lambda\cos(\frac{\sqrt{3}}{2}\ln x)
+\mu\sin(\frac{\sqrt{3}}{2}\ln x)\right)$$


Par hypothèse, $f$ est de classe $\mathcal{C}^1$ sur $\R$, en particulier 
elle se prolonge en $0$ de façon $\mathcal{C}^1$. Cherchons à quelle condition 
sur $\lambda,\mu$ cela est possible. Déjà, 
$f(x)=\sqrt{x}\left(\lambda\cos(\frac{\sqrt{3}}{2}\ln x)+\mu\sin(\frac{\sqrt{3}}{2}\ln x)\right)
\xrightarrow[x\to 0^+]{}0$ pour tous $\lambda,\mu$ ; donc $f(0)=0$. Mais
$$\frac{f(x)-f(0)}{x-0}= \frac1{\sqrt{x}}
\left(\lambda\cos(\frac{\sqrt{3}}{2}\ln x)+\mu\sin(\frac{\sqrt{3}}{2}\ln x)\right)$$ 
n'a pas de limite en $0$ si $\lambda\neq0$ ou $\mu\neq0$. En effet, pour 
$x_n=e^{\frac{2}{\sqrt{3}}(-2n\pi)}$, on a $x_n\to 0$ 
mais $\frac{f(x_n)-f(0)}{x_n-0}=\frac{\lambda}{\sqrt{x_n}}$ qui admet une limite finie seulement si $\lambda = 0$.
De même avec $x_n'=e^{\frac{2}{\sqrt{3}}(-2n\pi+\frac\pi2)}$ qui donne
$\frac{f(x_n')-f(0)}{x_n'-0}=\frac{\mu}{\sqrt{x_n'}}$ et implique donc $\mu=0$.


Par conséquent, la seule possibilité est $\lambda=\mu=0$. Ainsi $f$ est la fonction nulle, sur $[0,+\infty[$.
Le même raisonnement s'applique sur $]-\infty,0]$. La fonction est donc nécessairement nulle sur $\Rr$.
Réciproquement, la fonction constante nulle est bien solution du problème initial. 
\end{enumerate}
\fincorrection
\correction{004054}
\begin{enumerate}
  \item $y = 2 + \frac\lambda{x+2}$.
  \item $y = \frac{C+\sin x}x$.
  \item $y = -\cos x + \frac{\sin x + \lambda}{1+x}$.
  \item $y = \lambda x - \frac 1{3x^2}$.
  \item $y = \lambda x^{4/3} - x$.
  \item $y = \frac{\sin x - \cos x}2 + \frac{3\sin 2x - 6\cos 2x}5 + \lambda e^{-x}$.
  \item $y = \frac{\Argch(1-2x)+\lambda}{2\sqrt{x^2-x}}$ pour $x < 0$\par
               $y = \frac{\arcsin(2x-1)+\mu}{2\sqrt{x-x^2}}$ pour $0 < x < 1$\par
               $y = \frac{-\Argch(2x-1)+\nu}{2\sqrt{x^2-x}}$ pour $1 < x$.
  \item  $y = \frac {x-1}{2x} \arctan x +
       \frac {x+1}{2x} \left(\ln\frac {|x+1|}{\sqrt{x^2+1}} + \lambda \right)$.
  \item $y = \frac{x}{1-x^2}\Bigl( (1+x)\ln x + 1 + \lambda x\Bigr)$.
\end{enumerate}
 \fincorrection
\correction{004055}
\begin{enumerate}
  \item $y = (x + a\cos x + b\sin x) e^x$.
  \item $y = (ax+b)e^{2x} + 2xe^x$.
  \item $y = e^{2x}(a\cos3x + b\sin3x) + 2\cos2x + \sin 2x$.
  \item $y = \sin x\ln\Bigl|\tan \frac x2\Bigr| + \lambda \cos x + \mu\sin x$
               (variation de la constante avec sin).
  \item $y = (\lambda + \ln|x|)e^{-x} + \mu e^{-2x}$.
  \item $y = \lambda\cos x + \mu\sin x + \sum_{n=0}^\infty (-1)^n P^{(2n)}(x)$.
  \item $y = a\cos x + b\sin x$ avec $a^2 + b^2 = 1$ ou $y = \pm1$.
\end{enumerate}
 \fincorrection
\correction{004056}
\begin{enumerate}
  \item $y = -e^x + \lambda e^{e^x} + \mu e^{-e^x}$.
  \item $y = \lambda e^{x^2} + \mu e^{2x^2} + \frac{2x^2+3}{16}$.
  \item $y=\lambda x^2 + \mu\ln x$.
  \item $y = ax + bx^2 + 1 - 2x\sin x$.
  \item $y = x^2\ln|x+1| + \lambda \left(x^2\ln\left|\frac x{x+1}\right| + x-\frac12\right) + \mu x^2$.
  \item $y = \frac{-1 + a\ch x + b\sh x}{x^2}$.
  \item $y = \lambda\sqrt{x^2+3} + \mu x - 1$.
  \item $y^{(4)} - 2y'' + y = 0  \Rightarrow  y = a(\ch x -x\sh x) + b(x\ch x -\sh x)$.
\end{enumerate}
 \fincorrection
\correction{004057}
\begin{enumerate}
  \item $2n(2n-3)a_n = -9a_{n-3}
                   \Rightarrow  y = \begin{cases} a_0\cos(x^{3/2})  &\text{ si }x \ge 0\cr
                                  a_0\ch(|x|^{3/2}) &\text{ si }x \le 0.\cr\end{cases}$\par
             Solution générale :
             $y = \begin{cases} a\cos u + b\sin u &\text{ si } x \ge 0\cr
                          a\ch u + b\sh u   &\text{ si } x \le 0,\cr\end{cases}$ avec $u = |x|^{3/2}$.
  \item $n(n+1)a_n = a_{n-2}  \Rightarrow  y = a_0\frac{\sh x}x$.\par
             Solution générale : $y = \frac{a\ch x + b\sh x}x$.
  \item $(2n+1)(2n+2)a_{n+1} = a_n  \Rightarrow  y = a_0\ch(\sqrt x\,)$.\par
             Solution générale sur $\R^+$ : $u = a\ch(\sqrt x\,) + b\sh(\sqrt x\,)$.
  \item $n(n-1)a_n + (n+1)a_{n-2} = 0
                   \Rightarrow  y = a_0(1-x^2)e^{-x^2/2} + a_1z$.\par
             Solution générale : $y = (1-x^2)e^{-x^2/2}(a+bF(x)) + bx$
             avec $F(x) =  \int_{t=0}^x e^{t^2/2}\,d t$.
  \item $(n+2)(n+3)a_n = a_{n-2}  \Rightarrow  y = \frac{x-\sh x}{x^3}$.\par
             Solution générale : $y = \frac{a\ch x + b\sh x + x}{x^3}$.
  \item $na_{n+1} = (n+1)a_n  \Rightarrow  y = \frac{\lambda x}{(1-x)^2}$.\par
             Solution générale : $y = \frac {ax + b(1+x\ln|x|)}{(1-x)^2}$.
\end{enumerate}
 \fincorrection
\correction{004058}
$|\sin x| = \frac2\pi -\frac 4\pi\sum_{n=1}^\infty \frac{\cos 2nx}{4n^2-1}
 \Rightarrow  y = \frac2\pi -\frac4\pi\sum_{n=1}^\infty \frac{\cos 2nx}{(4n^2-1)(16n^4-4n^2+1)}$.
\par
Cette série converge et définit une fonction de classe $\mathcal{C}^4$ solution de l'équation.

Unicité : les solutions de l'équation homogène sont combinaison de
$e^{jx}$, $e^{-jx}$, $e^{j^2x}$ et $e^{-j^2x}$ donc non $\pi$-périodiques.
\fincorrection
\correction{004059}
$k \notin\Z$ : $y = \sum_{k=1}^\infty \frac{\cos nx}{n^2(k^2-n^2)} + a\cos kx + b\sin kx$.\par
         $k\in\Z$ : remplacer $\frac{\cos kx}{k^2(k^2-k^2)}$ par $\frac{x\sin kx}{2k^3}$.
\fincorrection
\correction{004060}
$y = x+1 + \lambda e^x$ ou $y = x - 1 + \lambda e^{-x}$.
\fincorrection
\correction{004061}
$$y = \begin{cases}e^x-1 &\text{ si } x < 0\cr
1-\cos x + \sin x &\text{ si } 0 \le x < \frac {3\pi}2\cr
e^{3\pi/2-x}-1    &\text{ si } x \ge \frac {3\pi}2.\cr\end{cases}$$
\fincorrection
\correction{004062}
$4xz''+2z'+z = \frac{-1}{\strut y^2}\Bigl(4xy''+2y'-y-\frac{8xy'^2}{\strut y}\Bigr)
=\frac{2}{\strut y^3}(y^2+4xy'^2)$
donc $\frac{y'}{\strut y} = \pm\frac1{\sqrt{-4x}}$
et $y = \lambda\exp(\pm\sqrt{-x})$ pour $x<0$.

Résolution sans indication~: on pose $x=\varepsilon t^2$ et $y(x) = z(t)$
d'où $\frac{d^2z}{d t^2}+\varepsilon z = 0$.
\fincorrection
\correction{004063}
\begin{enumerate}
  \item $g(x) = be^{-ax} +  \int_{t=0}^x e^{a(t-x)}f(t)\,d t$.
  \item $ \int_{x=0}^Xg(x)\,d x
    = \frac{b}{\strut a}(1-e^{-aX}) +  \int_{t=0}^X \int_{x=t}^Xe^{a(t-x)}f(t)\,d x\,d t
    = \frac{b}{\strut a}(1-e^{-aX}) +  \int_{t=0}^X\frac{1-e^{a(t-X)}}{\strut a}f(t)\,d t$

    $\phantom{ \int_{x=0}^Xg(x)\,d x}\to
     \frac{b}{\strut a}+ \frac{1}{\strut a} \int_{t=0}^{+\infty}f(t)\,d t$ lorsque $X\to+\infty$.

    Donc l'intégrale de~$g$ converge. On montre la convergence
    absolue par majoration élémentaire.
\end{enumerate}
 \fincorrection
\correction{004064}
\begin{enumerate}
  \item $x = 2\alpha e^t + (2\gamma t + 2\beta-\gamma) e^{2t}$, \quad
             $y = (\gamma t + \beta) e^{2t}$, \quad
             $z = \alpha e^t + (\gamma t + \beta) e^{2t}$.
  \item $y = -1 + \lambda e^{\alpha t} + \mu e^{\beta t}$,\quad
             $z = -1 + \lambda(1+\alpha) e^{\alpha t} + \mu(1+\beta) e^{\beta t}$,\quad
             $\alpha = \frac{-3+\sqrt5}2$,
             $\beta  = \frac{-3-\sqrt5}2$.
  \item $y = \frac{-3\cos t-13\sin t}{25} + (at+b)e^{2t}$,\quad
             $z = \frac{-4\cos t- 3\sin t}{25} + (at+a+b)e^{2t}$.
  \item $x = (a+bt+ct^2)e^t$,\quad
             $y = \left({a+\frac{b-c}2 + (b+c)t+ct^2}\right)e^t$,\quad
             $z = \left({a-\frac{b+c}2 + (b-c)t+ct^2}\right)e^t$.
  \item $x = -(b+c)e^t + (a+b+c)e^{2t}$,\par
             $y = \frac12(-a+5b+3c) - 2(b+c)e^t + \frac12(a+b+c)e^{2t}$,\par
             $z = \frac12(a-5b-3c)  + 3(b+c)e^t - \frac12(a+b+c)e^{2t}$.
  \item $x = (at^2 + (a+b+\frac12)t + a+b+c)e^t$,\par
             $y = (at^2 + (b-a+\frac12)t + a+c)e^t - \frac12e^{-t}$,\par
             $z = (-at^2 + (a-b-\frac12)t -c)e^t + \frac12e^{-t}$.
\end{enumerate}
 \fincorrection
\correction{004065}
$y = (t^2+1)x' -tx -2t^2+1  \Rightarrow  (t^2+1)x'' + 2tx' -2x = 6t$.\par
Résolution par DSE $ \Rightarrow  x = a(1+t\arctan t) + bt + t\ln(1+t^2)$,
                        $y = a\arctan t + b + 1 + \ln(1+t^2)$.
\fincorrection
\nocorrection
\nocorrection
\correction{004068}
Dériver deux fois. $f(x) = \frac{\sin x-x\cos x}2 + \lambda\sh x + \mu\cos x$.
\fincorrection
\correction{004069}
\begin{enumerate}
  \item $z'+ \frac z{\tanh x} = 0$.
  \item $y = \frac {ax+b}{\sh x}$.
  \item 
  \item 
  \item 
\end{enumerate}
 \fincorrection
\correction{004070}
\begin{enumerate}
  \item spectre ${}= \C$, $f_\lambda(t) = e^{-t^2/2}e^{\lambda t}$.
  \item Pour $\lambda \ne 0$, $\Phi^2(f) = \lambda^2 f \Leftrightarrow
             f = af_\lambda + bf_{-\lambda}$.\par
             Pour $\lambda = 0$, $\Phi^2(f) = 0 \Leftrightarrow f(t) = (at+b)e^{-t^2/2}$.
  \item $\Phi^2(y)=-2y \Leftrightarrow y = e^{-t^2/2}\bigl( a\cos(t\sqrt2\,) + b\sin(t\sqrt2\,)\bigr)$.
\end{enumerate}
 \fincorrection
\correction{004071}
\begin{enumerate}
  \item $\lambda = n^2$ : $P(X) = aX^n$.
  \item $\lambda > 0$ : $f(x) = \alpha x^{\sqrt\lambda} + \beta x^{-\sqrt\lambda}$.\par
            $\lambda = 0$ : $f(x) = \alpha + \beta\ln x$.\par
            $\lambda < 0$ : $f(x) = \alpha\cos\bigl(\sqrt{-\lambda}\ln x\bigr)
                                  + \beta \sin\bigl(\sqrt{-\lambda}\ln x\bigr)$.
  \item $\lambda \in \R$ : $f(x) = \alpha\exp\lambda\bigl( \arcsin\sqrt x  - \sqrt{x(1-x)} \bigr)$.
\end{enumerate}
 \fincorrection
\correction{004072}
\begin{enumerate}
  \item $\lambda = 2k$, $P = \alpha(X-1)^{n-k}(X+1)^{n+k}$ pour $-n \le k \le n$.
  \item $\lambda = 0$, $P = \alpha X^{2n}$.
  \item $\lambda = 0$, $P = \alpha (X^2+1)^n$.
\end{enumerate}
\fincorrection
\correction{004073}
$y =  \int_{t=0}^x g(t)\,d t  \Rightarrow  \frac{y'}y = \frac{2\pm\sqrt2}x
           \Rightarrow  g(x) = \alpha x^{1\pm\sqrt2}$.
          Continuité en $0  \Rightarrow  g(x) = \alpha x^{1+\sqrt2}$.
\fincorrection
\correction{004074}
\'Etudier $e^{-A}(y-z)$, $A'=a$.
\fincorrection
\correction{004075}
Point de concours : $\Bigl(x_0-\frac1{a(x_0)}, -\frac{b(x_0)}{a(x_0)}\Bigr)$.
\fincorrection
\nocorrection
\correction{004077}
\begin{enumerate}
  \item 
  \item Soient $y_0$ une solution particulière et $y_1$ une solution
             non nulle de l'équation homogène : $y_1(x) = e^{-A(x)}$ avec
             $A'=a$. Alors $y_0(x+T) = y_0(x) + \alpha y_1(x)$, et pour
             une solution $y$ quelconque, $y = y_0 + \lambda y_1$ :
             $y(x+T)-y(x) = (\alpha + \lambda(e^{-I}-1))y_1(x)$ où
             $I =  \int_{t=0}^T a(t)\,d t$.
\end{enumerate}
\fincorrection
\correction{004078}
$\lambda \ne 0 :\quad f(x) = a\sin(\alpha x),\quad
 \alpha \equiv \frac \pi4 (\mathrm{mod}\, \pi)$

$\lambda = 0 : \quad f = 0$.
\fincorrection
\correction{004079}
$f(1) = e^{-k} \int_{t=0}^1 g(t)e^{kt}\,d t$.

Avec Cauchy-Schwarz on obtient
 $ \int_{t=0}^1 (f'(t)+kf(t))^2\,d t \ge
            \frac{2k}{1-e^{-2k}}f(1)^2 = 9\frac{2k}{1-e^{-2k}}$.

Il y a égalité pour $f(t) = 3\frac{e^{kt}-e^{-kt}}{e^{k}-e^{-k}}$.
\fincorrection
\correction{004080}
\begin{enumerate}
  \item $u_n-u(t_n) =  \frac1{2\pi} \int_{t=t_n}^{t_n+2\pi} \int_{x=t_n}^t u'(x)\,d x\,d t$
et on majore l'intégrale interne par Cauchy-Schwarz.
  \item $w+w''=u'$ donc $w(t) =  \int_{x=t_n}^t\sin(t-x)u'(x)\,d x + \alpha\cos t + \beta\sin t$
puis 
\begin{align*} \int_{t=t_n}^{t_n+2\pi}w(t)\,d t\ 
&=  \int_{t=t_n}^{t_n+2\pi} \int_{x=t_n}^t\sin(t-x)u'(x)\,d x\,d t\\
&=  \int_{x=t_n}^{t_n+2\pi} \int_{t=x}^{t_n+2\pi}\sin(t-x)u'(x)\,d t\,d x\\
&=  \int_{x=t_n}^{t_n+2\pi} u'(x)(\cos(t_n-t)-1)\,d x\\
&\to 0,\\ \text{ lorsque }n\to\infty\\ \end{align*}
et
$ \int_{t=t_n}^{t_n+2\pi}v(t)\,d t
= w(t_n)-w(t_n+2\pi)
= - \int_{x=t_n}^{t_n+2\pi}\sin(t-x)u'(x)\,d x$.
\end{enumerate}
\fincorrection
\correction{004081}
\begin{enumerate}
  \item Formule de Duhamel : $y(t) = - \int_{x=0}^t e^{t-x}f(x)\,d x + \lambda e^t$.

    Par convergence dominée, l'intégrale tend vers $0$ quand $t$ tend vers $-\infty$
    donc toutes les solutions de~$(E)$ sont bornées au voisinage de $-\infty$.

    Pour $t\ge 0$ on a $y(t) = e^t\Bigl(\lambda- \int_{x=0}^t e^{-x}f(x)\,d x\Bigr)$
    donc il y a au plus une valeur de $\lambda$ telle que $y$ soit éventuellement
    bornée au voisinage de $+\infty$, c'est $\lambda =  \int_{x=0}^{+\infty} e^{-x}f(x)\,d x$.

    Pour ce choix on a~:
    $|y(t)| = \Bigl| \int_{x=t}^{+\infty}e^{t-x}f(x)\,d x\Bigr|
    \le \int_{x=t}^{+\infty}|f(x)|\,d x \to 0$ lorsque $t\to+\infty$.
  \item \begin{align*} \int_{t=a}^b|F(t)|\,d t
    &\le  \int_{t=a}^b \int_{x=t}^{+\infty}e^{t-x}|f(x)|\,d xd t\\
    &\le  \int_{x=a}^b \int_{t=a}^xe^{t-x}|f(x)|\,d td x
       +  \int_{x=b}^{+\infty} \int_{t=a}^be^{t-x}|f(x)|\,d td x\\
    &\le  \int_{x=a}^b(1-e^{a-x})|f(x)|\,d x
       +  \int_{x=b}^{+\infty}(e^{b-x}-e^{a-x})|f(x)|\,d td x\\
    &\le  \int_{x=-\infty}^{+\infty}|f(x)|\,d x\\ \end{align*}
    donc $F$ est intégrable. $F' = F-f$ est aussi intégrable et
    $ \int_{t=-\infty}^{+\infty}F'(t)\,d t = \Bigl[F(t)\Bigr]_{t=-\infty}^{+\infty} = 0$
    d'où $ \int_{-\infty}^{+\infty} F = \int_{-\infty}^{+\infty} f$.
\end{enumerate}
\fincorrection
\correction{004082}
Poser $F(x) =  \int_{t=0}^x f(t)\,d t$ et $G(x) =  \int_{t=0}^x g(t)\,d t$
puis résoudre~: $$\left\{\begin{array}{lll} F(x) &= &x-1+G'(x)\\ G(x) &= &x-1 + F'(x)\\ F(0) &= &G(0) = 0.\hfill\\ \end{array}\right.$$

On trouve $f(x) = g(x) = 1$.
\fincorrection
\correction{004083}
$y'$ étant bornée, $y$ admet une limite finie en tout point fini
donc la solution non prolongeable est définie sur~$\R$.

Une solution $y$ est de classe $\mathcal{C}^\infty$ vu l'équation et
$y'' = (1+\sin(x+y))\cos(x+y)$ est du signe de $\cos(x+y)$.
En un point $(x_0,y_0)$ tel que $x_0+y_0\equiv \frac\pi2(\mathrm{mod}\,{\pi})$
on a $y'' = 0$ et $\frac{d}{d x}(x+y)\ne 0$ donc $y''$ change
de signe et il y a inflexion.
En un point $(x_0,y_0)$ tel que $x_0+y_0\equiv \frac{3\pi}2(\mathrm{mod}\,{\pi})$
on a $x+y=$cste (car $y=\text{cste}-x$ est solution et il y a unicité)
donc il n'y a pas inflexion.
\fincorrection
\correction{004084}
\begin{enumerate}
  \item thm de Cauchy-Lipschitz linéaire.
  \item $x \mapsto A(-x)$ et $x \mapsto -B(-x)$ sont solutions de $(E)$ et
    vérifient les bonnes conditions initiales.
  \item Résoudre $u''(x) + ku(x) = 2d\cos(x)u(x)$ par la formule de Duhamel.
\end{enumerate}
 \fincorrection
\correction{004085}
\begin{enumerate}
  \item 
  \item Oui, $\mathrm{Ker} u = \{0\}$.
  \item Oui~: si $g\in E$ alors $f = t \mapsto \int_{s=-\infty}^t e^{s-t}g(s)\,d s$
    appartient~à~$E$ et $f+f'=g$.
\end{enumerate}
 \fincorrection
\correction{004086}
\begin{enumerate}
  \item 
  \item $u_p(x) = - \int_{t=0}^x pf(t)\sh\Bigl(\frac{x-t}p\Bigr)\,d t + \frac{\sh(x/p)}{\sh(1/p)} \int_{t=0}^1 pf(t)\sh\Bigl(\frac{1-t}p\Bigr)\,d t$.
  \item TCD : lorsque $p\to\infty$ $u_p(x)\to - \int_{t=0}^x (x-t)f(t)\,d t + x \int_{t=0}^1 (1-t)f(t)\,d t =  \int_{t=0}^x t(1-x)f(t)\,d t +  \int_{t=x}^1 x(1-t)f(t)\,d t$
    (primitive deuxième de~$-f$ s'annulant en~$0$ et~$1$).
\end{enumerate}
 \fincorrection
\correction{004087}
\begin{enumerate}
  \item Il suffit de démontrer que les solutions de~$({\cal E})$ sont développables en série entière.
    La méthode des coefficients indéterminés donne $n(n-1)a_n = (\lambda-1)a_{n-4}$ si $n\ge 4$
    et $a_2 = a_3 = 0$ d'où $a_{4k} = \frac{(\lambda - 1)^ka_0}{\prod_{i=1}^k4i(4i-1)}$,
    $a_{4k+1} = \frac{(\lambda - 1)^ka_1}{\prod_{i=1}^k4i(4i+1)}$, $a_{4k+2} = a_{4k+3} = 0$.
    On obtient une série de rayon infini pour tout choix de~$a_0,a_1$ donc les solutions DSE
    forment un espace vectoriel de dimension~$2$ et on a ainsi trouvé toutes les solutions.
  \item On doit avoir $H''(x) - 2xH'(x) + ((2-\lambda)x^2 - 1)H(x) = 0$.
    Si $H$ est une fonction polynomiale non nulle, en examinant les termes de
    plus haut degré on obtient une contradiction. Donc il n'existe pas de
    telle solution.
\end{enumerate}
 \fincorrection
\correction{004088}
Considérer $h(t) = a +  \int_{u=0}^t f(u)g(u)\,d u$ et
résoudre l'inéquation différentielle $h'(t) \le g(t)h(t)$ par la formule
de Duhamel.
\fincorrection
\correction{004089}
$f(x) =  \int_{t=0}^x g(t)\sin(x-t)\,d t + \lambda\cos x + \mu\sin x$
         avec $g = f+f''$.
\fincorrection
\correction{004090}
On pose $\varphi(t) = f''(t)+f'(t)+f(t)$.

$f(t) = e^{-t/2} \left[ -\frac 2{\sqrt3}
         \int_{u=0}^t \varphi(u)e^{u/2} \sin\left({\frac {\sqrt3\,(u-t)}2 }\right), d u
        + A\cos\left(\frac {\sqrt3\,u}2 \right) + B\sin\left(\frac {\sqrt3\,u}2\right)
        \right ]$.
\fincorrection
\nocorrection
\correction{004092}
\begin{enumerate}
  \item $y =  \int_{t=\alpha}^x b(t)e^{A(t)-A(x)}\,d t + y(\alpha)e^{-A(x)}$
             avec $A' = a$ et $A(\alpha) = 0$.\par
             Comme $a \ge 1$, on a $A(x) \ge x-\alpha$ et $A(t)-A(x) \le t-x$ pour $t \le x$.\par
             Donc $|y| \le  \int_{t=\alpha}^z |b(t)|e^{t-x}\,d t +
                            \int_{t=z}^x |b(t)|e^{t-x}\,d t + |y(\alpha)|e^{\alpha-x}
                       \le \|b\|_{\infty} e^{z-x} +
                           \sup\limits_{[z,+\infty[}|b| + |y(\alpha)|e^{\alpha-x}$.\par
             On choisit $z$ tel que $z\to+\infty$ et $x-z\to+\infty  \Rightarrow $ cqfd.
  \item Comme $A(t)-A(x) \le t-x$ pour $t \le x$, l'intégrale
             $ \int_{t=-\infty}^x b(t)e^{A(t)-A(x)}\,d t$ converge et fournit
             une solution nulle en $-\infty$.
             Comme $e^{-A(x)}\to +\infty$ lorsque $x\to-\infty$, c'est la seule.
\end{enumerate}
 \fincorrection
\correction{004093}
\begin{enumerate}
  \item Sinon, la convexité de $y$ est incorrecte.
  \item S'il existe $x$ tel que $z(x) = z'(x) = 0$, alors $z=0$ ce qui
             est absurde.\par
             S'il existe $x$ tel que $z(x) = 0 \ne z'(x)$, alors par convexité,
             $z$ ne peut s'annuler ailleurs.
\end{enumerate}
 \fincorrection
\nocorrection
\correction{004095}
Si $y_1$ et $y_2$ sont deux solutions bornées alors $y_1''$ et
$y_2''$ sont intégrables donc $y_i'(t) \to 0$ lorsque $t\to+\infty$ et $W_{y_1,y_2}(t)= 0$
pour tout~$t$.
\fincorrection
\correction{004096}
\begin{enumerate}
  \item On suppose $y\ne 0$ sinon $y$ n'a pas de zéros consécutifs.
    Comme $y(x_0) = 0$, on a $y'(x_0)\ne 0$ sinon $y=0$. Ceci implique que chaque
    zéro de~$y$ est isolé, donc la notion de zéros consécutifs est pertinente.
    Enfin, $y'(x_0)$ et $y'(x_1)$ sont de signes opposés sinon il existe un
    autre zéro dans~$]x_0,x_1[$.
    
  \item $W' = (q-r)yz$. $W(x_1)-W(x_0) = y'(x_0)z(x_0) - y'(x_1)z(x_1)$ (non simplifiable).
  \item Si $z$ ne s'annule pas dans $]x_0,x_1[$ alors $W'$ est de signe constant
    sur cet intervalle. L'examen des différents cas possibles de signe apporte une
    contradiction entre les signes de~$W'$ et de~$W(x_1)-W(x_0)$ si $z(x_0)\ne 0$
    ou $z(x_1) \ne 0$.
  \item On prend $r=q$, $z=u$. Si $u(x_0)\ne 0$
    alors $u$ admet un zéro dans~$]x_0,x_1[$ et en permutant les rôles de~$u$ et~$y$,
    le prochain zéro éventuel de~$u$ vient après $y_1$.
    Sinon, $u = \frac{u'(x_0)}{y'(x_0)}y$.
\end{enumerate}
 \fincorrection
\correction{004097}
\begin{enumerate}
  \item 
  \item 
    \begin{enumerate}
     \item Wronskien.
     \item $\left(\frac zy\right)' = \frac {z'y-zy'}{y^2}$ est de signe constant
$ \Rightarrow  \frac zy$ est monotone.

$\frac zy$ admet des limites infinies en $u$ et $v$. TVI
   \end{enumerate}
\end{enumerate}
 \fincorrection
\correction{004098}
$z' = -\frac\lambda{t^2}\sin(t-a)y(t)$ donc si $y$ ne s'annule pas
sur $]a,a+\pi[$, alors $z$ est strictement monotone sur $[a,a+\pi]$.
Mais $z(a+\pi)-z(a) = y(a+\pi)+y(a)  \Rightarrow $ contradiction de signe.
\fincorrection
\correction{004099}
\begin{enumerate}
  \item L'ensemble des zéros est localement fini d'après Cauchy-Lipchitz.

             Si $y$ ne s'annule pas sur $[a,+\infty[$, par exemple $y > 0$,
             alors $y$ est concave positive donc minorée, donc $y''\to-\infty$
             ce qui implique $y',y\to-\infty$, contradiction.
  \item 
  \item Soit $b_n = \frac\pi{e^{a_n/2}}$.
             Alors $b_{n+1} \le 2\ln\left(\frac {b_n}{b_{n+1}}\right) \le b_n$ et
             $b_n \to 0$ donc $b_n \sim b_{n+1} \sim 2\left(\frac{b_n}{b_{n+1}}-1\right)$.\par
             Alors $\frac 1{b_{n+1}} - \frac 1{b_n} \to \frac 12$,
             $b_n \sim \frac 1{2n}$ et $a_n \sim 2\ln n$.
\end{enumerate}
 \fincorrection
\correction{004100}
Les formes linéaires $y \mapsto y(a)$ et $y \mapsto y(b)$ sont linéairement
indépendantes sur l'espace des solutions de l'équation homogène.
\fincorrection
\correction{004101}
On se ramène au cas $z=0$. Soit $x$ tel que $y(x) < 0$ et $y'(x) = 0$.
         Alors $y''(x) < 0$, donc $y$ n'est pas minimale en $x$. Donc $y$ n'a
         pas de minimum local sur $]a,+\infty[$.
\fincorrection
\correction{004102}
Si $x_i(0)>0$ pour tout~$i$ on obtient une contradiction
en considérant le plus petit $t$ tel qu'il existe $i$ avec $x_i(t)<0$.

Cas général~: dépendance continue de la solution par rapport aux conditions
initiales\dots
\fincorrection
\correction{004103}
$f'(x) = \frac {x+i}{2(x^2+1)} f(x)
  \Rightarrow  f(x) = \sqrt\pi (x^2+1)^{-1/4}\exp\left(\frac i2\arctan x\right)$.
\fincorrection
\correction{004104}
\begin{enumerate}
  \item Soit $f$ non identiquement nulle vérifiant $f''=\lambda\Delta f$ avec $\lambda > 0$~:
    sur tout intervalle où $f$ est strictement positive, $f$ est strictement convexe
    donc ne peut pas s'annuler aux deux bords~; idem quand $f$ est strictement
    négative, il y a contradiction. Le cas $\lambda=0$ est trivial.
  \item $(f\mid g) =  \int_{t=a}^bf'(t)g'(t)\,d t = - \int_{t=a}^bf''(t)g(t)\,d t  = - \int_{t=a}^bf(t)g''(t)\,d t$.
  \item
  \begin{enumerate}
     \item   
     \item
  Si $f_\lambda$ a un nombre fini de zéros, soit $x_n$ le
    dernier et $A=\max(x_n,2)$. Sur $[A,+\infty[$, $f$ est de signe constant, $\varepsilon$,
    et on a $f_\lambda''-\lambda f_\lambda = \lambda(\Delta-1)f_\lambda = \varphi$
    d'où $f_\lambda(x) =  \int_{t=A}^x\sin((x-t)\sqrt{-\lambda}\,)\varphi(t)\,d t + \alpha\cos(x\sqrt{-\lambda}\,) + \beta\sin(x\sqrt{-\lambda}\,)$.
    En particulier $f_\lambda(A)+f_\lambda\Bigl(A+\frac{\pi}{\sqrt{-\lambda}}\Bigr)
    =  \int_{t=A}^{A+\pi/\sqrt{-\lambda}}\sin((x-t)\sqrt{-\lambda}\,)\varphi(t)\,d t$ est du
    signe de $-\varepsilon$, absurde.
    
    Si l'ensemble des zéros de~$f$ admet un point d'accumulation~$x$ on
    a $f_\lambda(x) = f_\lambda'(x) = 0$ d'où $f_\lambda=0$, absurde.
  \end{enumerate}
\end{enumerate}
 \fincorrection
\correction{004105}
\begin{enumerate}
  \item Après le prolongement indiqué on peut appliquer le relation
de Parseval à $f$ et $f'$ sachant que $c_0(f) = 0$ par imparité et
$|c_n(f)| = |c_n(f')|/n\le |c_n(f')|$ pour $n\ne0$.
  \item $x''(t) + q(t)x(t) = 0  \Rightarrow   \int_0^\pi x'^2 = \Bigl[xx'\Bigr]_0^\pi -  \int_0^\pi xx''
=  \int_0^\pi qx^2  \Rightarrow  x'=0  \Rightarrow  x=0$.

Rmq: il n'est pas nécessaire d'avoir $q$ de classe $\mathcal{C}^1$.
  \item Il existe $x_0$ de classe $\mathcal{C}^2$ vérifiant l'équation différentielle.
Par différence avec $x_0$ on se ramène au cas $f = 0$ et il faut montrer
que l'application $\varphi : {\cal S} \to {\R^2}, x \mapsto {(x(0),x(\pi))}$ est bijective,
en notant $\cal S$ l'espace des solutions de l'équation homogène $x'' + qx = 0$.
Or $\varphi$ est linéaire et est injective d'après la question précédente, c'est
donc une bijection car $\dim{\cal S} = 2$.

Remarque~: l'hypothèse $f$ de classe $\mathcal{C}^1$ est inutile, continue suffit.
\end{enumerate}
 \fincorrection
\correction{004106}
$x$ et $1/x$ sont solution $ \Rightarrow $ $x'^2+qx^2=0$ donc une condition nécessaire
est~: $q(t)\le 0$ et $q=-x'^2/x^2$ est de classe $\mathcal{C}^1$ .
Réciproquement, supposons $q$ négative de classe $\mathcal{C}^1$ et soit $r(t)=\sqrt{-q(t)}$.
Si $x$ est solution de $x'=r(t)x$ alors sur tout intervalle~$I$ où $q$ ne
s'annule pas on a $x'' = r(t)x' + r'(t)x$ donc
$$x''+p(t)x'+q(t)x = (r(t)+p(t))x'+(r'(t)+q(t))x = (r(t)p(t)+r'(t))x$$
donc une deuxième condition nécessaire est~: $p(t)q(t)=-\frac12q'(t)$.
Ces deux conditions sont suffisantes si $q$ est strictement négative.
\fincorrection
\correction{004107}
La suite $(X_k)$ de fonctions définie par $X_k(t) = X_0$,
$X_{k+1}(t) = X_0 +  \int_{u=0}^t A(u)X_k(u)\,d u$ converge localement
uniformément vers $X$ et $X_k(t)$ est clairement à composantes positives
pour $t\ge 0$.
\fincorrection
\correction{004108}
Pour $n=1$ et $A(t)=a>0$ on trouve après les incantations usuelles
(équation homogène, variation de la constante et mise en forme de l'intégrale) que
$X\ :\ t \mapsto \int_{u=0}^1F(tu)u^{a-1}\,d u$ est l'unique solution
prolongeable en $0$ et qu'elle est de classe $\mathcal{C}^\infty$ sur~$\R$.
Pour $n=1$ et $A$ non constante, on trouve de même~:
$$X(t) =  \int_{u=0}^1F(tu)u^{A(0)-1}\exp\Bigl( \int_{v=t}^{tu}\frac{A(v)-A(0)}v\,d v\Bigr)\,d u$$
et l'on voit que $X$ est $\mathcal{C}^\infty$ en écrivant
$\frac{A(v)-A(0)}v =  \int_{w=0}^1A'(vw)\,d w$.

Pour $n$ quelconque et $A$ constante~: alors la fonction $X\ :\ t \mapsto \int_{u=0}^1u^{A(0)-I}F(tu)\,d u$ est
l'unique solution prolongeable en $0$,
en convenant que $u^{A(0)-I} = \exp((A(0)-I)\ln(u))$ (l'intégrale converge en~$0$
car $u^{A(0)-I} = O(u^{\alpha-1}\ln(u)^n)$ pour tout $\alpha>0$
minorant les parties réelles des valeurs propres de~$A(0)$).

Pour $A$ non constante, on met l'équation sous forme intégrale~:
$$tX'(t) + A(t)X(t) = F(t)\Leftrightarrow X(t) =  \int_{u=0}^1u^{A(0)-I}\{F(tu)-(A(tu)-A(0))X(tu)\}\,d u .$$
Soit $a>0$ à choisir.
Posons $E = \mathcal{C}([-a,a],\C^n)$ et pour $X\in E$~: $$\Phi(X) = t \mapsto \int_{u=0}^1u^{A(0)-I}\{F(tu)-(A(tu)-A(0))X(tu)\}\,d u.$$
On a facilement~: $\Phi(X)\in E$ si $X\in E$ et $\Phi$ est contractante sur~$E$ pour $\|\ \|_\infty$
si $a$ est choisi suffisament petit. Donc l'équation $tX'(t) + A(t)X(t) = F(t)$
admet une solution (unique) définie au voisinage de $0$, et cette solution est prolongeable
en une solution sur~$\R$ car le théorème de Cauchy-Lipschitz s'applique
en dehors de~$0$. Par ailleurs, on a~:
$$\forall\ t\in\R,\ X(t) =  \int_{u=0}^1u^{A(0)-I}\{F(tu)-(A(tu)-A(0))X(tu)\}\,d u,$$
ce qui montre par récurrence sur~$k\in\N$ que $X$ est de classe $\mathcal{C}^k$ sur~$\R$.
\fincorrection
\correction{004109}
\begin{enumerate}
  \item Poser $F(t) =  \int_{u=t_0}^t f(u)\,d u$ et résoudre l'inéquation
    différentielle $F'(t) \le g(t) + kF(t)$ par la formule de Duhamel.
  \item \begin{align*}
    M' = AM
    & \Rightarrow  \|\kern-1.2pt|M'(t)\|\kern-1.2pt| \le K\|\kern-1.2pt|M(t)\|\kern-1.2pt|\\
    & \Rightarrow  \|\kern-1.2pt|M(t)-I\|\kern-1.2pt| \le K \int_{u=t0}^t\|\kern-1.2pt|M(u)\|\kern-1.2pt|\,d u\\
    & \Rightarrow  \|\kern-1.2pt|M(t)\|\kern-1.2pt| \le 1 + K \int_{u=t0}^t\|\kern-1.2pt|M(u)\|\kern-1.2pt|\,d u\\
    & \Rightarrow  \|\kern-1.2pt|M(t)\|\kern-1.2pt| \le 1 + K \int_{u=t0}^te^{K(t-u)}\,d u = e^{K(t-t_0)}.\\
    \noalign{}
    (M-N)' &= (A-B)M + B(M-N)\\
    & \Rightarrow  \|\kern-1.2pt|(M-N)'(t)\|\kern-1.2pt| \le \eta e^{K(t-t_0)} + (K+\eta)\|\kern-1.2pt|(M-N)(t)\|\kern-1.2pt|\\
    & \Rightarrow  \|\kern-1.2pt|(M-N)(t)\|\kern-1.2pt| \le \frac\eta K (e^{K(t-t_0)}-1) + (K+\eta) \int_{u=t_0}^t\|\kern-1.2pt|(M-N)(u)\|\kern-1.2pt|\,d u\\
    & \Rightarrow  \|\kern-1.2pt|(M-N)(t)\|\kern-1.2pt| \le \underbrace{\frac\eta K (e^{K(t-t_0)}-1) + \frac{(K+\eta)\eta}K \int_{u=t_0}^te^{(K+\eta)(t-u)}(e^{K(u-t_0)}-1)\,d u}_{\textstyle=e^{K(t-t_0)}(e^{\eta(t-t_0)}-1)}\\
    \end{align*}
  \item $X_0(t) = M_0(t)\alpha$ et $Y_0(t) = N_0(t)\alpha$, d'où
    $\|X_0(t)-Y_0(t)\| \le e^{K(t-t_0)}(e^{\eta(t-t_0)}-1)\|\alpha\|$.
\end{enumerate}
 \fincorrection
\correction{005874}
Dans tout l'exercice, on note $(E)$ l'équation différentielle considérée et $(E_H)$ l'équation homogène associée.

\begin{enumerate}
 \item  Les solutions de $(E)$ sur $\Rr$ forment un $\Rr$-espace affine de direction l'espace des solutions de $(E_H)$ sur $\Rr$ qui est de dimension $1$. La fonction $x\mapsto1$ est une solution de $(E)$ sur $\Rr$ et la fonction $x\mapsto e^{-x}$ est une solution non nulle de $(E_H)$ sur $\Rr$. Donc

\begin{center}
\shadowbox{
$\mathcal{S}_\Rr=\left\{x\mapsto1+\lambda e^{-x},\;\lambda\in\Rr\right\}$.
}
\end{center}

\item  Les solutions de $(E_H)$ sur $\Rr$ sont les fonctions de la forme $x\mapsto \lambda e^{x/2}$. Déterminons maintenant une solution particulière de $(E)$ sur $\Rr$.

\textbf{1ère solution.} Il existe une solution particulière de $(E)$ sur $\Rr$ de la forme $x\mapsto a\cos x+b\sin x$, $(a,b)\in\Rr^2$. Soit $f$ une telle fonction. Alors, pour tout réel $x$,

\begin{center}
$2f'(x)-f(x)=2(-a\sin x+b\cos x)-(a\cos x+b\sin x)=(-a+2b)\cos x+(-2a-b)\sin x$.
\end{center}

Par suite,

\begin{align*}\ensuremath
f\;\text{solution de}\;(E)\;\text{sur}\;\Rr&\Leftarrow\forall x\in\Rr,\;2f'(x)-f(x)=\cos x\Leftarrow
\left\{
\begin{array}{l}
-a+2b=1\\
-2a-b=0
\end{array}
\right.
\\
 &\Leftarrow a=- \frac{1}{5}\;\text{et}\;b= \frac{2}{5}.
\end{align*}

\begin{center}
\shadowbox{
$\mathcal{S}_\Rr=\left\{x\mapsto \frac{1}{5}(-\cos x+2\sin x)+\lambda e^{x/2},\;\lambda\in\Rr\right\}$.
}
\end{center}

\textbf{2ème solution.} Par la méthode de variation de la constante, il existe une solution particulière de $(E)$ sur $\Rr$ de la forme $x\mapsto\lambda(x)e^{x/2}$ où $\lambda$ est une fonction dérivable sur $\Rr$. Soit $f$ une telle fonction.

\begin{align*}\ensuremath
f\;\text{solution de}\;(E)\;\text{sur}\Rr&\Leftarrow\forall x\in\Rr,\;2\left(\lambda'(x)e^{x/2}+ \frac{1}{2}\lambda(x)e^{x/2}\right)-2\lambda(x)e^{x/2}=\cos(x)\\
 &\Leftarrow\forall x\in\Rr,\;\lambda'(x)= \frac{1}{2}e^{-x/2}\cos x.
\end{align*}

Or, 

\begin{align*}\ensuremath
\int_{}^{} \frac{1}{2}e^{-x/2}\cos x\;dx&= \frac{1}{2}\text{Re}\left(\int_{}^{}e^{(-\frac{1}{2}+i)x}\;dx\right)= \frac{1}{2}\text{Re}\left( \frac{e^{(-\frac{1}{2}+i)x}}{- \frac{1}{2}+i}\right)+C= \frac{1}{5}e^{-x/2}\text{Re}\left((\cos x+i\sin x)(-1-2i)\right)+C\\
 &= \frac{1}{5}e^{-x/2}(-\cos x+2\sin x)+C.
\end{align*}

Par suite, on peut prendre $\lambda(x)= \frac{1}{5}e^{-x/2}(-\cos x+2\sin x)$ ce qui fournit la solution particulière $f_0(x)= \frac{1}{5}(-\cos x+2\sin x)$.

\item  Puisque les fonctions $x\mapsto-2$ et $x\mapsto xe^{2x}$ sont continues sur $\Rr$, l'ensemble des solutions de $(E)$ sur $\Rr$ est un $\Rr$-espace affine de dimension $1$. Soit $f$ une fonction dérivable sur $\Rr$.

\begin{align*}\ensuremath
f\;\text{solution de}\;(E)\;\text{sur}\;\Rr&\Leftrightarrow\forall x\in\Rr,\;f'(x)-2f(x)=xe^{2x}\Leftrightarrow\forall x\in\Rr,\;e^{-2x}f'(x)-2e^{-2x}f(x)=x\Leftrightarrow\forall x\in\Rr,\;(e^{-2x}f)'(x)=x\\
 &\Leftrightarrow\exists\lambda\in\Rr/\;\forall x\in\Rr,\;e^{-2x}f(x)= \frac{x^2}{2}+\lambda\Leftrightarrow\exists\lambda\in\Rr/\;\forall x\in\Rr,\;f(x)=\left( \frac{x^2}{2}+\lambda\right)e^{2x}.
\end{align*}

\begin{center}
\shadowbox{
$\mathcal{S}_\Rr=\left\{x\mapsto\left( \frac{x^2}{2}+\lambda\right)e^{2x},\;\lambda\in\Rr\right\}$.
}
\end{center}

\item  L'équation caractéristique $(E_c)$ associée à l'équation homogène $y''-4y'+4y =0$ est $z^2-4z+4=0$ et admet $z_0=2$ pour racine double. On sait que les solutions de $(E_H)$ sur $\Rr$ sont les fonctions de la forme $x\mapsto(\lambda x+\mu)e^{2x}$, $(\lambda,\mu)\in\Rr^2$.

Puisque $2$ est racine double de l'équation caractéristique, l'équation $y''-4y'+4y=e^{2x}$ admet une solution particulière $f_0$ de la forme : $\forall x\in\Rr$, $f_0(x)=ax^2e^{2x}$, $a\in\Rr$. La formule de \textsc{Leibniz} fournit pour tout réel $x$,

\begin{center}
$f_0''(x)-4f_0'(x)+4f_0(x)=a(4x^2+8x+2)e^{2x}-4a(2x^2+2x)e^{2x}+4ax^2e^{2x}=2ae^{2x}$,
\end{center}

et $f_0$ est solution de $(E)$ sur $\Rr$ si et seulement si $a= \frac{1}{2}$.

\begin{center}
\shadowbox{
$\mathcal{S}_\Rr=\left\{x\mapsto\left( \frac{x^2}{2}+\lambda x+\mu\right)e^{2x},\;(\lambda,\mu)\in\Rr^2\right\}$.
}
\end{center}

\item  L'équation caractéristique $(E_c)$ associée à l'équation homogène $y''+4y =0$ est $z^2+4=0$ et admet deux racines non réelles conjuguées $z_1=2i$ et $z_2=\overline{z_1}=-2i$. On sait que les solutions de $(E_H)$ sur $\Rr$ sont les fonctions de la forme $x\mapsto\lambda \cos(2x)+\mu\sin(2x)$, $(\lambda,\mu)\in\Rr^2$.

Une solution réelle de l'équation $y''+4y =\cos(2x)$ est la partie réelle d'une solution de l'équation $y''+4y=e^{2ix}$. Puisque le nombre $2i$ est racine simple de $(E_c)$, cette dernière équation admet une solution de la forme $f_1~:~x\mapsto axe^{2ix}$, $a\in\Cc$. La formule de \textsc{Leibniz} fournit pour tout réel $x$,

\begin{center}
$f_1''(x)+4f_1(x)=a((-4x+4i)e^{2ix}+4xe^{2ix})=4iae^{2ix}$.
\end{center}

et $f_1$ est solution de $y''+4y=e^{2ix}$ si et seulement si $a= \frac{1}{4i}$. On obtient $f_1(x)= \frac{1}{4i}xe^{2ix}= \frac{1}{4}x(-i\cos(2x)+\sin(2x))$ ce qui fournit une solution particulière de $(E)$ sur $\Rr$ : $\forall x\in\Rr$, $f_0(x)= \frac{1}{4}x\sin(2x)$.

\begin{center}
\shadowbox{
$\mathcal{S}_\Rr=\left\{x\mapsto \frac{1}{4}x\sin(2x)+\lambda\cos(2x)+\mu\sin(2x),\;(\lambda,\mu)\in\Rr^2\right\}$.
}
\end{center}

\item  L'équation caractéristique $(E_c)$ associée à l'équation $(E_H)$ est $z^2+2z+2=0$ et admet deux racines non réelles conjuguées $z_1=-1+i$ et $z_2=\overline{z_1}=-1-i$. On sait que les solutions de $(E_H)$ sur $\Rr$ sont les fonctions de la forme $x\mapsto(\lambda \cos(x)+\mu\sin(x))e^{-x}$, $(\lambda,\mu)\in\Rr^2$.

Pour tout réel $x$, $\cos(x)\ch(x)=\text{Re}\left(e^{ix}\ch(x)\right)= \frac{1}{2}\text{Re}\left(e^{(1+i)x}+e^{(-1+i)x}\right)$. Notons $(E_1)$ l'équation $y''+2y'+2y=e^{(1+i)x}$ et $(E_2)$ l'équation $y''+2y'+2y=e^{(-1+i)x}$. Si $f_1$ est une solution de $(E_1)$ et $f_2$ est une solution de $(E_2)$ alors $f_0= \frac{1}{2}\text{Re}(f_1+f_2)$ est une solution de $(E)$ sur $\Rr$ d'après le principe de superposition des solutions.

\textbullet~$(E_1)$ admet une solution particulière de la forme $f_1~:~x\mapsto ae^{(1+i)x}$, $a\in\Cc$. Pour tout réel $x$,

\begin{center}
$f_1''(x)+2f_1'(x)+2f_1(x)=a((1+i)^2+2(1+i)+2)e^{(1+i)x}=a(4+4i)e^{(1+i)x}$
\end{center}

et $f_1$ est solution de $(E_1)$ sur $\Rr$ si et seulement si $a= \frac{1}{4+4i}= \frac{1-i}{8}$. On obtient $f_1(x)= \frac{1-i}{8}e^{(1+i)x}$.

\textbullet~$(E_2)$ admet une solution particulière de la forme $f_2~:~x\mapsto axe^{(-1+i)x}$, $a\in\Cc$. La formule de \textsc{Leibniz} fournit pour tout réel $x$,

\begin{center}
$f_2''(x)+2f_2'(x)+2f_2(x)=a(((-1+i)^2x+2(-1+i))+2((-1+i)x+1)+2x)e^{(-1+i)x}=2iae^{(-1+i)x}$
\end{center}

et $f_2$ est solution de $(E_2)$ sur $\Rr$ si et seulement si $a= \frac{1}{2i}=- \frac{i}{2}$. On obtient $f_2(x)=- \frac{i}{2}e^{(-1+i)x}$.

\textbullet~Une solution particulière $f_0$ de $(E)$ sur $\Rr$ est donc définie pour tout réel $x$ par

\begin{align*}\ensuremath
f_0(x)&= \frac{1}{2}\text{Re}\left( \frac{1-i}{8}e^{(1+i)x}- \frac{i}{2}e^{(-1+i)x}\right)= \frac{1}{2}\text{Re}\left( \frac{1}{8}(1-i)(\cos(x)+i\sin(x))e^{x}- \frac{i}{2}(\cos(x)+i\sin(x))e^{-x}\right)\\
 &= \frac{1}{2}\left( \frac{1}{8}(\cos(x)+\sin(x))e^{x}+ \frac{1}{2}\sin(x)e^{-x}\right)
\end{align*}

\begin{center}
\shadowbox{
$\mathcal{S}_\Rr=\left\{x\mapsto \frac{1}{16}(\cos(x)+\sin(x))e^{x}+ \frac{1}{4}\sin(x)e^{-x}+(\lambda \cos(x)+\mu\sin(x))e^{-x},\;(\lambda,\mu)\in\Rr^2\right\}$.
}
\end{center}
\end{enumerate}
\fincorrection
\correction{005875}
\begin{enumerate}
 \item   Posons $g=f'+\alpha f$. La fonction $g$ est continue sur $\Rr$ et la fonction $f$ est solution sur $\Rr$ de l'équation différentielle $y'+\alpha y=g$. De plus, $\lim_{x \rightarrow +\infty}g(x)=\ell$. Ensuite,

\begin{align*}\ensuremath
f'+\alpha f=g&\Rightarrow \forall x\in\Rr,\;e^{\alpha x}f'(x)+\alpha e^{\alpha x}f(x)=e^{\alpha x}g(x)\Rightarrow\forall x\in\Rr,\;(e^{\alpha x}f)'(x)=e^{\alpha x}g(x)\\
 &\Rightarrow\forall x\in\Rr,\;e^{\alpha x}f(x)=f(0)+\int_{0}^{x}e^{\alpha t}g(t)\;dt\Rightarrow\forall x\in\Rr,\;f(x)=f(0)e^{-\alpha x}+e^{-\alpha x}\int_{0}^{x}e^{\alpha t}g(t)\;dt.
\end{align*}

Puisque $\text{Re}(\alpha)>0$ et que $\left|e^{-\alpha x}\right|=e^{-\text{Re}(\alpha)x}$, $\lim_{x \rightarrow +\infty}f(0)e^{-\alpha x}=0$. Vérifions alors que $\lim_{x \rightarrow +\infty}e^{-\alpha x}\int_{0}^{x}e^{\alpha t}g(t)\;dt= \frac{\ell}{\alpha}$ sachant que $\lim_{x \rightarrow +\infty}g(x)=\ell$.

On suppose tout d'abord $\ell=0$. Soit $\varepsilon>0$. Il existe $A_1>0$ tel que $\forall t\geqslant A_1$, $|g(t)|\leqslant \frac{\varepsilon}{2}$. Pour $x\geqslant A_1$,

\begin{align*}\ensuremath
\left|e^{-\alpha x}\int_{0}^{x}e^{\alpha t}g(t)\;dt\right|&\leqslant e^{-\text{Re}(\alpha)x}\left|\int_{0}^{A_1}e^{\alpha t}g(t)\;dt\right|+e^{-\text{Re}(\alpha)x}\int_{A_1}^{x}e^{\text{Re}(\alpha)t}|g(t)|\;dt\\
 &\leqslant e^{-\text{Re}(\alpha)x}\left|\int_{0}^{A_1}e^{\alpha t}g(t)\;dt\right|+e^{-\text{Re}(\alpha)x}\int_{A_1}^{x}e^{\text{Re}(\alpha)t}\times \frac{\varepsilon}{2}\;dt\\
  &=e^{-\text{Re}(\alpha)x}\left|\int_{0}^{A_1}e^{\alpha t}g(t)\;dt\right|+ \frac{\varepsilon}{2}\left(1-e^{-\text{Re}(\alpha)(x-A_1)}\right)\\
  &\leqslant e^{-\text{Re}(\alpha)x}\left|\int_{0}^{A_1}e^{\alpha t}g(t)\;dt\right|+ \frac{\varepsilon}{2}.
\end{align*}

Maintenant $\lim_{x \rightarrow +\infty}e^{-\text{Re}(\alpha)x}\left|\int_{0}^{A_1}e^{\alpha t}g(t)\;dt\right|=0$ et donc il existe $A\geqslant A_1$ tel que $\forall x> A$, $e^{-\text{Re}(\alpha)x}\left|\int_{0}^{A_1}e^{\alpha t}g(t)\;dt\right|< \frac{\varepsilon}{2}$. Pour $x>A$, on a $\left|e^{-\alpha x}\int_{0}^{x}e^{\alpha t}g(t)\;dt\right|< \frac{\varepsilon}{2}+ \frac{\varepsilon}{2}=\varepsilon$. On a ainsi montré que $\lim_{x \rightarrow +\infty}f(x)=0= \frac{\ell}{\alpha}$.

On revient maintenant au cas général $\ell$ quelconque.

\begin{align*}\ensuremath
f'+\alpha f\underset{x\rightarrow+\infty}{\rightarrow}\ell&\Rightarrow f'+\alpha f-\ell
\underset{x\rightarrow+\infty}{\rightarrow}0\Rightarrow\left(f- \frac{\ell}{\alpha}\right)'+
\alpha\left(f- \frac{\ell}{\alpha}\right)\underset{x\rightarrow+\infty}{\rightarrow}0\\
 &\Rightarrow f- \frac{\ell}{\alpha}\underset{x\rightarrow+\infty}{\rightarrow}0\Rightarrow f\underset{x\rightarrow+\infty}{\rightarrow} \frac{\ell}{\alpha}.
 \end{align*}

\begin{center}
\shadowbox{
$\forall f\in C^1(\Rr,\Rr)$, $\forall \alpha\in\Cc$ tel que $\text{Re}(\alpha)>0$, $\lim_{x \rightarrow +\infty}(f'(x)+\alpha f(x))=\ell\Rightarrow\lim_{x \rightarrow +\infty}f(x)= \frac{\ell}{\alpha}$.
}
\end{center}

\item  $f''+f'+f=\left(f'-jf\right)'-j^2\left(f'-jf\right)$. D'après 1), comme $\text{Re}(-j^2)=\text{Re}(-j)= \frac{1}{2}>0$,

\begin{center}
$f''+f'+f\underset{x\rightarrow+\infty}{\rightarrow}0\Rightarrow\left(f'-jf\right)'-j^2\left(f'-jf\right)\underset{x\rightarrow+\infty}{\rightarrow}0\Rightarrow f'-jf\underset{x\rightarrow+\infty}{\rightarrow}0\Rightarrow f\underset{x\rightarrow+\infty}{\rightarrow}0$.
\end{center}

\begin{center}
\shadowbox{
$\forall f\in C^2(\Rr,\Rr)$, $\lim_{x \rightarrow +\infty}(f''(x)+f'(x)+f(x))=0\Rightarrow\lim_{x \rightarrow +\infty}f(x)=0$.
}
\end{center}

\item  Montrons le résultat par récurrence sur $n$.

\textbullet~Pour $n=1$, c'est le 1) dans le cas particulier $\ell=0$ (si $P=X-\alpha$, $P(D)(f)=f'-\alpha f$ avec $\text{Re}(-\alpha)>0$).

\textbullet~Soit $n\in\Nn^*$. Supposons le résultat acquis pour $n$. Soit $P$ un polynôme de degré $n+1$ dont les racines ont des parties réelles strictement négatives et tel que $\lim_{x \rightarrow +\infty}(P(D))(f)(x)=0$. Soit $\alpha$ une racine de $P$. $P$ s'écrit $P=(X-\alpha)Q$ où $Q$ est un polynôme dont les racines ont toutes une partie réelle strictement négative. Puisque

\begin{center}
$P(D)(f)=((D-\alpha Id)\circ(Q(D))(f)=(Q(D)(f))'-\alpha(Q(D)(f))\underset{+\infty}{\rightarrow}0$,
\end{center}

on en déduit que $Q(D)(f)\underset{+\infty}{\rightarrow}0$ d'après le cas $n=1$ puis que $f\underset{+\infty}{\rightarrow}0$ par hypothèse de récurrence.

Le résultat est démontré par récurrence.
\end{enumerate}
\fincorrection
\correction{005876}
On pose $g=f+f''$. Par hypothèse, la fonction $g$ est une application continue et positive sur $\Rr$ et de plus, la fonction $f$ est solution sur $\Rr$ de l'équation différentielle $y''+y=g$ sur $\Rr$. Résolvons cette équation différentielle, notée $(E)$, sur $\Rr$.

Les solutions de l'équation homogène associée sont les fonctions de la forme $x\mapsto\lambda\cos x+\mu\sin x$, $(\lambda,\mu)\in\Rr^2$. D'après la méthode de variation des constantes, il existe une solution particulière de $(E)$ sur $\Rr$ de la forme $f_0~:~x\mapsto\lambda(x)\cos(x)+\mu(x)\sin(x)$ où de plus les fonctions $\lambda$ et $\mu$ sont solutions du système

\begin{center}
$\left\{
\begin{array}{l}
\lambda'\cos(x)+\mu'\sin(x)=0\\
-\lambda'\sin(x)+\mu'\cos(x)=g
\end{array}
\right.$.
\end{center}

Les formules de \textsc{Cramer} fournissent $\forall x\in\Rr$, $\lambda'(x)=-g(x)\sin(x)$ et $\mu'(x)=g(x)\cos(x)$. On peut alors prendre $\forall x\in\Rr$, $\lambda(x)=-\int_{0}^{x}g(t)\sin(t)\;dt$ et $\mu(x)=\int_{0}^{x}g(t)\cos(t)\;dt$ puis 

\begin{center}
$\forall x\in\Rr$, $f_0(x)=-\cos(x)\int_{0}^{x}g(t)\sin(t)\;dt+\sin(x)\int_{0}^{x}g(t)\cos(t)\;dt=\int_{0}^{x}g(t)\sin(x-t)\;dt$.
\end{center}

Ainsi, les solutions de $(E)$ sur $\Rr$ sont les fonctions de la forme $x\mapsto\lambda\cos(x)+\mu\sin(x)+\int_{0}^{x}g(t)\sin(x-t)\;dt$, $(\lambda,\mu)\in\Rr^2$. La fonction $f$ est l'une de ces solutions. Par suite, il existe $(\lambda_0,\mu_0)\in\Rr^2$ tel que $\forall x\in\Rr$, $f(x)=\lambda_0\cos(x)+\mu_0\sin(x)+\int_{0}^{x}g(t)\sin(x-t)\;dt$ et donc pour tout réel $x$,

\begin{align*}\ensuremath
f(x)+f(x+\pi)&=\int_{0}^{x+\pi}g(t)\sin(x+\pi-t)\;dt+\int_{0}^{x}g(t)\sin(x-t)\;dt=-\int_{0}^{x+\pi}g(t)\sin(x-t)\;dt+\int_{0}^{x}g(t)\sin(x-t)\;dt\\
 &=\int_{x}^{x+\pi}g(t)\sin(t-x)\;dt=\int_{0}^{\pi}g(u+x)\sin(u)\;du\geqslant0.
\end{align*}

On a montré que si $\forall x\in\Rr$, $f(x)+f''(x)\geqslant 0$, alors $\forall x\in\Rr$, $f(x)+f(x+\pi)\geqslant 0$.
\fincorrection
\correction{005877}
Dans tout l'exercice, on note $(E)$ l'équation proposée et $(E_H)$ l'équation homogène associée.

\begin{enumerate}
 \item  On note $J$ l'un des deux intervalles $]-\infty,0[$ ou $]0,+\infty[$. Sur $J$, l'équation $(E)$ s'écrit encore $y'- \frac{2}{x}y=0$. Comme la fonction $x\mapsto- \frac{2}{x}$ est continue sur $J$, les solutions de $(E)$ sur $J$ constituent un $\Rr$-espace vectoriel de dimension $1$. Enfin, la fonction $x\mapsto x^2$ est une solution non nulle de $(E)$ sur $J$ et donc $\mathcal{S}_J=\left\{
x\mapsto \lambda x^2,\;\lambda\in\Rr\right\}$.

Soit $f$ une solution de $(E)$ sur $\Rr$. Nécessairement, $\exists(\lambda_1,\lambda_2)\in\Rr^2/\;\forall x\in\Rr,\;f(x)=\left\{
\begin{array}{l}
\lambda_1x^2\;\text{si}\;x>0\\
0\;\text{si}\;x=0\\
\lambda_2x^2\;\text{si}\;x<0
\end{array}
\right.=\left\{
\begin{array}{l}
\lambda_1x^2\;\text{si}\;x\geqslant0\\
\lambda_2x^2\;\text{si}\;x<0
\end{array}
\right.$.

Réciproquement, une telle fonction $f$ est définie sur $\Rr$, dérivable sur $\Rr^*$, solution de $(E)$ sur $\Rr^*$ et vérifie encore l'équation $(E)$ en $0$ si de plus elle est dérivable en $0$. Donc, une telle fonction est solution de $(E)$ sur $\Rr$ si et seulement si elle est dérivable en $0$.

Il est géométriquement clair que $f$ est dérivable en $0$ pour tout choix de $\lambda_1$ et $\lambda_2$ et donc $f$ est solution de $(E)$ sur $\Rr$ pour tout choix de $\lambda_1$ et $\lambda_2$.

\begin{center}
\shadowbox{
$\mathcal{S}_\Rr=\left\{
x\mapsto\left\{
\begin{array}{l}
\lambda_1x^2\;\text{si}\;x\geqslant0\\
\lambda_2x^2\;\text{si}\;x<0
\end{array}
\right.,\;(\lambda_1,\lambda_2)\in\Rr^2\right\}$.
}
\end{center}

On note que $\mathcal{S}_\Rr$ est un $\Rr$-espace vectoriel de dimension $2$. En effet, pour toute solution $f$ de $(E)$ sur $\Rr$, $f(x)=\lambda_1\left\{\begin{array}{l}
x^2\;\text{si}\;x\geqslant0\\
0\;\text{si}\;x<0
\end{array}
\right.+\lambda_2\left\{
\begin{array}{l}
0\;\text{si}\;x\geqslant0\\
x^2\;\text{si}\;x<0
\end{array}
\right.=\lambda_1f_1(x)+\lambda_2f_2(x)$. Donc $\mathcal{S}_\Rr=\text{Vect}(f_1,f_2)$ avec $(f_1,f_2)$ clairement libre.

Un exemple de graphe de solution est donné à la page suivante.

$$\includegraphics{../images/img005877-1}$$



\item  L'ensemble des solutions sur $]-\infty,0[$ ou $]0,+\infty[$ est $\{x\mapsto\lambda x,\;\lambda\in\Rr\}$.

Soit $f$ une solution de $(E)$ sur $\Rr$. Nécessairement, $\exists(\lambda_1,\lambda_2)\in\Rr^2/\;\forall x\in\Rr,\;f(x)=\left\{
\begin{array}{l}
\lambda_1x\;\text{si}\;x\geqslant0\\
\lambda_2x\;\text{si}\;x<0
\end{array}
\right.$.

Réciproquement, une telle fonction $f$ est solution de l'équation $(E)$ sur $\Rr$ si et seulement si elle est dérivable en $0$.

Il est géométriquement clair que $f$ est dérivable en $0$ si et seulement si $\lambda_1=\lambda_2$ et donc

\begin{center}
\shadowbox{
$\mathcal{S}_\Rr=\{x\mapsto\lambda x,\;\lambda\in\Rr\}$.
}
\end{center}

Dans ce cas, $\mathcal{S}_\Rr$ est un $\Rr$-espace vectoriel de dimension $1$.

\item  L'ensemble des solutions sur $]-\infty,0[$ ou $]0,+\infty[$ est $\{x\mapsto \frac{\lambda}{x},\;\lambda\in\Rr\}$.

Soit $f$ une solution de $(E)$ sur $\Rr$. Nécessairement, $\exists(\lambda_1,\lambda_2)\in\Rr^2/\;\forall x\in\Rr,\;f(x)=\left\{
\begin{array}{l}
\lambda_1x\;\text{si}\;x>0\\
0\;\text{si}\;x=0\\
\lambda_2x\;\text{si}\;x<0
\end{array}
\right.$.

Réciproquement, une telle fonction $f$ est solution de l'équation $(E)$ sur $\Rr$ si et seulement si elle est dérivable en $0$.

Il est géométriquement clair que $f$ est dérivable en $0$ si et seulement si $\lambda_1=\lambda_2=0$ et donc

\begin{center}
\shadowbox{
$\mathcal{S}_\Rr=\{0\}$.
}
\end{center}

Dans ce cas, $\mathcal{S}_\Rr$ est un $\Rr$-espace vectoriel de dimension $0$.

\item  Soit $f$ une fonction dérivable sur $]0,+\infty[$.

\begin{align*}\ensuremath
f\;\text{solution de}\;(E)\;\text{sur}\;J&\Leftrightarrow\forall x\in J,\;xf'(x)-2f(x)=x^3\Leftrightarrow\forall x\in J,\; \frac{1}{x^2}f'(x)- \frac{2}{x^3}f(x)=1\\
 &\Leftrightarrow\forall x\in J,\;\left( \frac{1}{x^2}f\right)'(x)=1\Leftrightarrow\exists\lambda\in\Rr/\;\forall x\in J,\; \frac{f(x)}{x^2}=x+\lambda\Leftrightarrow\exists\lambda\in\Rr/\;\forall x\in J,\; \frac{f(x)}{x^2}=x+\lambda\\
  &\Leftrightarrow\exists\lambda\in\Rr/\;\forall x\in J,\;f(x)=x^3+\lambda x^2.
\end{align*}

\begin{center}
\shadowbox{
$\mathcal{S}_{]0,+\infty[}=\{x\mapsto x^3+\lambda x^2,\;\lambda\in\Rr\}$.
}
\end{center}

\item  Si $f$ est une solution sur $\Rr$ de l'équation $x^2y'+2xy = 1$ alors $0^2\times f'(0)+0\times f(0)=1$ ce qui est impossible. Donc

\begin{center}
\shadowbox{
$\mathcal{S}_{\Rr}=\varnothing$.
}
\end{center}

\item  \textbullet~\textbf{Résolution sur $]-\infty,0[$, $]0,1[$ et $]1,+\infty[$.} Soit $I$ l'un des trois intervalles $]-\infty,0[$, $]0,1[$ ou $]1,+\infty[$. Sur $I$, l'équation $(E)$ s'écrit encore $y'+ \frac{1}{2x}y= \frac{1}{2x(1-x)}$. Puisque les fonctions $x\mapsto \frac{1}{2x}$ et $x\mapsto \frac{1}{2x(1-x)}$ sont continues sur $I$, les solutions de $(E)$ sur $I$ constituent un $\Rr$-espace affine de dimension $1$. Soit $f$ une fonction dérivable sur $I$. Pour $x\in I$, on note $\varepsilon$ le signe de $x$ sur $I$.

\begin{align*}\ensuremath
f\;\text{solution de}\;(E)\;\text{sur}\;I&\Leftrightarrow\forall x\in I,\;f'(x)+ \frac{1}{2x}f(x)= \frac{1}{2x(1-x)}\Leftrightarrow\forall x\in I,\;e^{\frac{1}{2}\ln(\varepsilon x)}f'(x)+ \frac{1}{2x}e^{\frac{1}{2}\ln(\varepsilon x)}f(x)= \frac{e^{\frac{1}{2}\ln(\varepsilon x)}}{2x(1-x)}\\
 &\Leftrightarrow\forall x\in I,\;(\sqrt{\varepsilon x}\;f)'(x)= \frac{\sqrt{\varepsilon x}}{2x(1-x)}\Leftrightarrow\forall x\in I,\;(\sqrt{\varepsilon x}\;f)'(x)= \frac{
 \varepsilon}{2\sqrt{\varepsilon x}(1-x)}.
\end{align*}

Déterminons alors les primitives de la fonction $x\mapsto \frac{
 \varepsilon}{2\sqrt{\varepsilon x}(1-x)}$ sur $I$. En posant $u=\sqrt{\varepsilon x}$ et donc $x=\varepsilon u^2$ puis $du=2\varepsilon udu$.
 
\begin{center}
$\int_{}^{} \frac{
 \varepsilon}{2\sqrt{\varepsilon x}(1-x)}\;dx=\int_{}^{} \frac{\varepsilon}{2u(1-\varepsilon u^2)}2\varepsilon u\;du=\int_{}^{} \frac{1}{1-\varepsilon u^2}\;du$.
\end{center}

-\textit{Résolution sur $]-\infty,0[$.}

Dans ce cas, $\varepsilon =-1$ puis $\int_{}^{} \frac{
 \varepsilon}{2\sqrt{\varepsilon x}(1-x)}\;dx=\int_{}^{} \frac{1}{1+u^2}\;du=\Arctan(u)+\lambda=\Arctan(\sqrt{-x})+\lambda$. Par suite,
 
 \begin{align*}\ensuremath
 f\;\text{solution de}\;(E)\;\text{sur}\;]-\infty,0[&\Leftrightarrow\exists\lambda\in\Rr/\;\forall x\in]-\infty,0[,\;\sqrt{-x}f(x)=\Arctan(\sqrt{-x})+\lambda\\
  &\Leftrightarrow\exists\lambda\in\Rr/\;\forall x\in]-\infty,0[,\;f(x)= \frac{\Arctan(\sqrt{-x})+\lambda}{\sqrt{-x}}.
 \end{align*}

\begin{center}
\shadowbox{
$\mathcal{S}_{]-\infty,0[}=\left\{x\mapsto \frac{\Arctan(\sqrt{-x})+\lambda}{\sqrt{-x}},\;\lambda\in\Rr\right\}$.
}
\end{center}

-\textit{Résolution sur $]0,1[$ et $]1,+\infty[$.}

Dans ce cas, $\varepsilon =1$ puis $\int_{}^{} \frac{
 \varepsilon}{2\sqrt{\varepsilon x}(1-x)}\;dx=\int_{}^{} \frac{1}{1-u^2}\;du=\left\{
 \begin{array}{l}
  \frac{1}{2}\ln\left(
  \frac{\sqrt{x}+1}{\sqrt{x}-1}\right)\;\text{si}\;x>1
 \\
\rule{0mm}{5mm} \Argth(\sqrt{x})\;\text{si}\;x\in]0,1[
 \end{array}
 \right.+\lambda$. Par suite,

-\textit{Résolution sur $]0,1[$ et $]1,+\infty[$.}

\begin{center}
\shadowbox{
$\mathcal{S}_{]0,1[}=\left\{x\mapsto \frac{\Argth\left(\sqrt{x}\right)+\lambda}{\sqrt{x}},\;\lambda\in\Rr\right\}$ et $\mathcal{S}_{]1,+\infty[}=\left\{x\mapsto \frac{ \frac{1}{2}\ln\left(
  \frac{\sqrt{x}+1}{\sqrt{x}-1}\right)+\lambda}{\sqrt{x}},\;\lambda\in\Rr\right\}$.
}
\end{center}

-\textit{Résolution sur $]0,+\infty[$ et sur $\Rr$.} Si $f$ est une solution de $(E)$ sur $]0,+\infty[$ ou sur $\Rr$, alors $0\times f'(1)+0\times f(1)=1$ ce qui

est impossible. Donc

\begin{center}
\shadowbox{
$\mathcal{S}_{]0,+\infty[}=\varnothing$ et $\mathcal{S}_{\Rr}=\varnothing$.
}
\end{center}

-\textit{Résolution sur $]-\infty,1[$.} Si $f$ est une solution de $(E)$ sur $]-\infty,1[$, alors il existe nécessairement $(\lambda_1,\lambda_2)\in\Rr^2$ tel que

$\forall x\in]-\infty,1[$, $f(x)=\left\{
\begin{array}{l}
 \frac{\Arctan(\sqrt{-x})+\lambda}{\sqrt{-x}}\;\text{si}\;x<0\\
1\;\text{si}\;x=0\\
 \frac{\Argth\left(\sqrt{x}\right)+\lambda}{\sqrt{x}}\;\text{si}\;0<x<1
\end{array}
\right.$. Réciproquement une telle fonction est solution si et seulement

si elle est dérivable en $0$.

Quand $x$ tend vers $0$ par valeurs inférieures,

\begin{center}
$f(x)= \frac{1}{\sqrt{-x}}\left(\lambda_1+\sqrt{-x}- \frac{\left(\sqrt{-x}\right)^3}{3}+o\left((\sqrt{-x})^3\right)\right)= \frac{\lambda_1}{\sqrt{-x}}+1+ \frac{x}{3}+o(x)= \frac{\lambda_1}{\sqrt{-x}}+f(0)+ \frac{x}{3}+o(x)$.
\end{center}

et quand $x$ tend vers $0$ par valeurs supérieures,

\begin{center}
$f(x)= \frac{1}{\sqrt{x}}\left(\lambda_2+\sqrt{x}+ \frac{\left(\sqrt{x}\right)^3}{3}+o\left((\sqrt{x})^3\right)\right)= \frac{\lambda_2}{\sqrt{x}}+1+ \frac{x}{3}+o(x)= \frac{\lambda_2}{\sqrt{x}}+f(0)+ \frac{x}{3}+o(x)$.
\end{center}

Par suite, $f$ est dérivable à droite et à gauche en $0$ si et seulement si $\lambda_1=\lambda_2=0$ et dans ce cas, quand $x$ tend vers $0$,

$f(x)=f(0)+ \frac{x}{3}+o(x)$ ce qui montre que $f$ est dérivable en $0$. En résumé, $f$ est solution de $(E)$ sur $]-\infty,1[$ si et

seulement si $\lambda_1=\lambda_2=0$.

\begin{center}
\shadowbox{
$\mathcal{S}_{]-\infty,1[}=\left\{
x\mapsto\left\{
\begin{array}{l}
 \frac{\Arctan\left(\sqrt{-x}\right)}{\sqrt{-x}}\;\text{si}\;x<0\\
\rule[-2mm]{0mm}{5mm}1\;\text{si}\;x=0\\
 \frac{\Argth\left(\sqrt{x}\right)}{\sqrt{x}}\;\text{si}\;0<x<1
\end{array}
\right.
\right\}$.
}
\end{center}

\item  \textbf{Résolution de $(E)$ sur $]-\infty,0[$ et sur $]0,+\infty[$.} Soit $I$ l'un des deux intervalles $]-\infty,0[$ ou $]0,+\infty[$. On note $\varepsilon$ le signe de $x$ sur $I$. Sur $I$, $(E)$ s'écrit encore $y'+\varepsilon\left(1- \frac{1}{x}\right)y=\varepsilon x^2$. Puisque les deux fonctions $x\mapsto\varepsilon\left(1- \frac{1}{x}\right)$ et $x\mapsto\varepsilon x^2$ sont continues sur $I$, les solutions de $(E)$ sur $I$ constituent un $\Rr$-espace affine de dimension $1$. Soit $f$ une fonction dérivable sur $I$.

\begin{align*}\ensuremath
f\;\text{solution de}\;(E)\;\text{sur}\;I&\Leftrightarrow\forall x\in I,\;f'(x)+\varepsilon\left(1- \frac{1}{x}\right)f(x)=\varepsilon x^2\\
 &\Leftrightarrow\forall x\in I,\;e^{\varepsilon x-\varepsilon\ln(\varepsilon x)}f'(x)+\varepsilon\left(1- \frac{1}{x}\right)e^{\varepsilon x-\varepsilon\ln(\varepsilon x)}f(x)=\varepsilon x^2e^{\varepsilon x-\varepsilon\ln(\varepsilon x)}\\
 &\Leftrightarrow\forall x\in I,\;\left((\varepsilon x)^{-\varepsilon}e^{\varepsilon x}f\right)'(x)= x^{-\varepsilon}x^2e^{\varepsilon x}
\end{align*}

\textbullet~Si $I=]0,+\infty[$, $\varepsilon=1$ et 

\begin{align*}\ensuremath
f\;\text{solution de}\;(E)\;\text{sur}\;]0,+\infty[&\Leftrightarrow\forall x\in ]0,+\infty[,\;\left( \frac{e^x}{x}f\right)'(x)=xe^{x}
\Leftrightarrow\exists\lambda\in\Rr/\;\forall x\in]0,+\infty[,\; \frac{e^x}{x}f(x)=(x-1)e^x+\lambda\\
 &\Leftrightarrow\exists\lambda\in\Rr/\;\forall x\in]0,+\infty[,\;f(x)=x^2-x+\lambda xe^{-x}.
\end{align*}

\begin{center}
\shadowbox{
$\mathcal{S}_{]0,+\infty[}=\left\{x\mapsto x^2-x+\lambda xe^{-x},\;\lambda\in\Rr\right\}$.
}
\end{center}

\textbullet~Si $I=]-\infty,0[$, $\varepsilon=-1$ et 

\begin{align*}\ensuremath
f\;\text{solution de}\;(E)\;\text{sur}\;]0,+\infty[&\Leftrightarrow\forall x\in ]0,+\infty[,\;\left(-xe^{-x}f\right)'(x)=x^3e^{-x}
\Leftrightarrow\forall x\in ]0,+\infty[,\;\left(xe^{-x}f\right)'(x)=-x^3e^{-x}.
\end{align*}

Or, $\int_{}^{}-x^3e^{-x}\;dx=x^3e^{-x}-3\int_{}^{}x^2e^{-x}\;dx=(x^3+3x^2)e^{-x}-6\int_{}^{}xe^{-x}\;dx=(x^3+3x^2+6x+6)e^{-x}+\lambda$ et donc

\begin{align*}\ensuremath
f\;\text{solution de}\;(E)\;\text{sur}\;]-\infty,0[&\Leftrightarrow\exists\lambda\in\Rr/\;\forall x\in]-\infty,0[,\;xe^{-x}f(x)=(x^3+3x^2+6x+6)e^{-x}+\lambda\\
 &\Leftrightarrow\exists\lambda\in\Rr/\;\forall x\in]-\infty,0[,\;f(x)=x^2+3x+6+ \frac{6+\lambda e^x}{x}.
\end{align*}

\begin{center}
\shadowbox{
$\mathcal{S}_{]0,+\infty[}=\left\{x\mapsto x^2+3x+6+ \frac{6+\lambda e^x}{x},\;\lambda\in\Rr\right\}$.
}
\end{center}

\textbf{Résolution de $(E)$ sur $\Rr$.} Si $f$ est une solution de $(E)$ sur $\Rr$, nécessairement il existe $(\lambda_1,\lambda_2)\in\Rr^2$ tel que $\forall x\in\Rr$, $f(x)=\left\{
\begin{array}{l}
x^2-x+\lambda_1 xe^{-x}\;\text{si}\;x>0\\
0\;\text{si}\;x=0\\
x^2+3x+6+ \frac{6+\lambda_2 e^x}{x}\;\text{si}\;x<0
\end{array}
\right.=\left\{
\begin{array}{l}
x^2-x+\lambda_1 xe^{-x}\;\text{si}\;x\geqslant0\\
x^2+3x+6+ \frac{6+\lambda_2 e^x}{x}\;\text{si}\;x<0
\end{array}
\right.$. Réciproquement, une telle fonction est solution sur $\Rr$ si et seulement si elle est dérivable en $0$.

Quand $x$ tend vers $0$ par valeurs supérieures, $f(x)=-x+o(x)+\lambda_1 x(1+o(1))=(\lambda_1-1)x+o(x)$. Par suite, $f$ est dérivable à droite en $0$ pour tout choix de $\lambda_1$ et $f_d'(0)=\lambda_1-1$.

Quand $x$ tend vers $0$ par valeurs inférieures,

\begin{align*}\ensuremath
f(x)&=6+3x+o(x)+ \frac{6+\lambda_2\left(1+x+ \frac{x^2}{2}+o(x^2)\right)}{x}= \frac{6+\lambda_2}{x}+6+\lambda_2+\left(3+ \frac{\lambda_2}{2}\right)x+o(x)
\end{align*}

Par suite, $f$ est dérivable à gauche en $0$ si et seulement si $\lambda_2=-6$. Dans ce cas, quand $x$ tend vers $0$ par valeurs inférieures, $f(x)=o(x)$ et $f_g'(0)=0$. Maintenant, $f$ est dérivable en $0$ si et seulement si $f$ est dérivable à droite et à gauche en $0$ et $f_d'(0)=f_g'(0)$. Ceci équivaut à $\lambda_2=-6$ et $\lambda_1=1$.

\begin{center}
\shadowbox{
$\mathcal{S}_{\Rr}=\left\{x\mapsto \left\{
\begin{array}{l}
x^2-x+xe^{-x}\;\text{si}\;x\geqslant0\\
x^2+3x+6+ \frac{6(1-e^x)}{x}\;\text{si}\;x<0
\end{array}
\right.\right\}$.
}
\end{center}
\end{enumerate}
\fincorrection
\correction{005878}
\textbullet~Pour $n\in\Nn$, posons $a_n=(-1)^{n-1} \frac{C_{2n}^n}{2n-1}$. La suite $(a_n)_{n\in\Nn}$ ne s'annule pas et pour $n\in\Nn$,

\begin{align*}\ensuremath
 \frac{a_{n+1}}{a_n}&=- \frac{(2n+2)!}{(2n)!}\times \frac{n!^2}{(n+1)!^2}\times \frac{2n-1}{2n+1}=- \frac{(2n+2)(2n+1)}{(n+1)^2}\times \frac{2n-1}{2n+1}=- \frac{2(2n-1)}{n+1}.
\end{align*}

Par suite, $\left| \frac{a_{n+1}}{a_n}\right|\underset{n\rightarrow+\infty}{\rightarrow}4$ et d'après la règle de d'\textsc{Alembert}, $R_a= \frac{1}{4}$. Pour $x$ tel que la série converge, on pose $f(x)=\sum_{n=0}^{+\infty}(-1)^{n-1} \frac{C_{2n}^n}{2n-1}x^n$.

\textbullet~Soit $x\in\left]- \frac{1}{4}, \frac{1}{4}\right[$. Pour $n\in\Nn$, on a $(n+1)a_{n+1}+4na_n=2a_n$. Pour chaque $n\in\Nn$, on multiplie les deux membres de cette égalité par $x^n$ puis on somme sur $n$. On obtient $\sum_{n=0}^{+\infty}(n+1)a_{n+1}x^{n}+4x\sum_{n=1}^{+\infty}na_nx^{n-1}=2\sum_{n=0}^{+\infty}a_nx^n$ ou encore $(1+4x)f'(x)=2f(x)$. De plus $f(0)=a_0=1$. Mais alors

\begin{align*}\ensuremath
\forall x\in\left]- \frac{1}{4}, \frac{1}{4}\right[,\;(1+4x)f'(x)=2f(x)&\Rightarrow\forall x\in\left]- \frac{1}{4}, \frac{1}{4}\right[,\;e^{-\frac{1}{2}\ln(1+4x)}f'(x)- \frac{2}{1+4x}e^{-\frac{1}{2}\ln(1+4x)}f(x)=0\\
 &\Rightarrow\forall x\in\left]- \frac{1}{4}, \frac{1}{4}\right[,\;\left( \frac{f}{\sqrt{1+4x}}\right)'(x)=0\Rightarrow\forall x\in\left]- \frac{1}{4}, \frac{1}{4}\right[,\; \frac{f(x)}{\sqrt{1+4x}}= \frac{f(0)}{\sqrt{1+0}}\\
 &\Rightarrow\forall x\in\left]- \frac{1}{4}, \frac{1}{4}\right[,\;f(x)=\sqrt{1+4x}.
\end{align*}

\begin{center}
\shadowbox{
$\forall x\in\left]- \frac{1}{4}, \frac{1}{4}\right[$, $\sum_{n=0}^{+\infty}(-1)^n \frac{C_{2n}^n}{2n-1}x^n=\sqrt{1+4x}$.
}
\end{center}

\textbullet~Pour $n\in\Nn$, posons $u_n= \frac{C_{2n}^n}{(2n-1)4^n}$. La suite $u$ est strictement positive à partir du rang $1$ et pour $n\geqslant1$, 

\begin{center}
$ \frac{u_{n+1}}{u_n}= \frac{2(2n-1)}{4(n+1)}= \frac{2n-1}{2n+2}<1$.
\end{center}

Ainsi, la suite $u$ est décroissante à partir du rang $1$. De plus, d'après la formule de \textsc{Stirling},

\begin{center}
$u_n= \frac{(2n)!}{n!^2(2n-1)4^n}\underset{n\rightarrow+\infty}{\sim} \frac{\left( \frac{2n}{e}\right)^{2n}\sqrt{4\pi n}}{\left( \frac{n}{e}\right)^{2n}(2\pi n)(2n)4^n}= \frac{1}{2\sqrt{\pi}n^{3/2}}$.
\end{center}

Par suite, $u_n\underset{n\rightarrow+\infty}{\rightarrow}0$. En résumé, la suite $u$ est positive et décroissante à partir du rang $1$ et $\lim_{n \rightarrow +\infty}u_n=0$. On en déduit que la série numérique de terme général $(-1)^{n-1} \frac{C_{2n}^n}{(2n-1)4^n}=(-1)^{n-1}u_n$ converge en vertu du critère spécial aux séries alternées (théorème de \textsc{Leibniz}).

\textbullet~La fonction $f$ est donc définie en $ \frac{1}{4}$. Vérifions que $f$ est continue en $ \frac{1}{4}$.

Pour  $x\in\left]0, \frac{1}{4}\right]$ et $n\in\Nn$, posons $f_n(x)=a_nx_n=(-1)^{n-1} \frac{C_{2n}^n}{2n-1}x^n$. Pour chaque $x$ de $\left]0, \frac{1}{4}\right]$, la suite $(f_n(x))_{n\in\Nn}$ ne s'annule pas et la suite $((-1)^{n-1}f_n(x))_{n\in\Nn}$ est positive à partir du rang $1$. Ensuite, pour $n\geqslant1$ et $x\in\left]0, \frac{1}{4}\right]$,

\begin{align*}\ensuremath
\left| \frac{f_{n+1}(x)}{f_n(x)}\right|&=\left| \frac{a_{n+1}}{a_n}\right|x= \frac{2(2n-1)}{n+1}x\leqslant \frac{2(2n-1)}{n+1}\times \frac{1}{4}= \frac{4n-2}{4n+4}<1
\end{align*}

On en déduit que pour chaque $x$ de $\left]0, \frac{1}{4}\right]$, la suite numérique $(|f_n(x)|)_{n\in\Nn}$ décroît à partir du rang $1$. D'après une majoration classique du reste à l'ordre $n$ d'une série alternée, pour $n\geqslant1$ et $x\in\left]0, \frac{1}{4}\right]$

\begin{center}
$\left|\sum_{k=n+1}^{+\infty}f_k(x)\right|\leqslant|f_{n+1}(x)|=|a_{n+1}|x^{n+1}\leqslant \frac{|a_{n+1}|}{4^{n+1}}=u_{n+1}$,
\end{center}

et donc $\forall n\in\Nn$, $\text{Sup}\left\{\left|\sum_{k=n+1}^{+\infty}f_k(x)\right|,\;x\in\left]0, \frac{1}{4}\right]\right\}\leqslant u_{n+1}$. Puisque la suite $(u_n)$ tend vers $0$ quand $n$ tend vers $+\infty$, on a montré que la série de fonction de terme général $f_n$, $n\in\Nn$, converge uniformément vers $f$ sur $\left]0, \frac{1}{4}\right]$. Puisque chaque fonction $f_n$ est continue sur $\left]0, \frac{1}{4}\right]$, $f$ est continue sur $\left]0, \frac{1}{4}\right]$ et en particulier en $ \frac{1}{4}$.

Mais alors

\begin{center}
$\sum_{n=0}^{+\infty} \frac{(-1)^{n-1}C_{2n}^n}{(2n-1)4^n}=f\left( \frac{1}{4}\right)=\displaystyle\lim_{\substack{x\rightarrow\frac{1}{4}\\
x<\frac{1}{4}}}f(x)=\displaystyle\lim_{\substack{x\rightarrow\frac{1}{4}\\
x<\frac{1}{4}}}\sqrt{1+4x}=\sqrt{1+4\times \frac{1}{4}}=\sqrt{2}$.
\end{center}

\begin{center}
\shadowbox{
$\sum_{n=0}^{+\infty} \frac{(-1)^{n-1}C_{2n}^n}{(2n-1)4^n}=\sqrt{2}$.
}
\end{center}
\fincorrection
\correction{005879}
\begin{enumerate}
 \item  Posons $A=\left(
\begin{array}{cc}
4&-2\\
1&1
\end{array}
\right)$.

$\chi_A=\lambda^2-5\lambda+6=(\lambda-2)(\lambda-3)$ puis $A=PDP^{-1}$ où $D=\text{diag}(2,3)$ et $P=\left(
\begin{array}{cc}
1&2\\
1&1
\end{array}
\right)$.

Posons $X=\left(
\begin{array}{c}
x\\
y
\end{array}
\right)$ puis $X_1=P^{-1}X=\left(
\begin{array}{c}
x_1\\
y_1
\end{array}
\right)$.

\begin{align*}\ensuremath
\left\{
\begin{array}{l}
x'=4x-2y\\
y'=x+y
\end{array}
\right.&\Leftrightarrow X'=AX\Leftrightarrow X'=PDP^{-1}X\Leftrightarrow P^{-1}X'=DP^{-1}X\Leftrightarrow(P^{-1}X)'=D(P^{-1}X)\Leftrightarrow X_1'=DX_1\\
 &\Leftrightarrow\left\{
\begin{array}{l}
x_1'=2x_1\\
y_1'=3y_1
\end{array}
\right.\Leftrightarrow\exists(a,b)\in\Rr/\;\forall t\in\Rr,\;\left\{
\begin{array}{l}
x_1(t)=ae^{2t}\\
y_1(t)=be^{3t}
\end{array}
\right.\\
&\Leftrightarrow\exists(a,b)\in\Rr/\;\forall t\in\Rr,\;\left(
\begin{array}{l}
x(t)\\
y(t)
\end{array}
\right)=\left(
\begin{array}{cc}
1&2\\
1&1
\end{array}
\right)\left(
\begin{array}{l}
ae^{2t}\\
be^{3t}
\end{array}
\right)\\
 &\Leftrightarrow\exists(a,b)\in\Rr/\;\forall t\in\Rr,\;\left(
\begin{array}{l}
x(t)\\
y(t)
\end{array}
\right)=\left(
\begin{array}{l}
ae^{2t}+2be^{3t}\\
ae^{2t}+be^{3t}
\end{array}
\right)\\
\end{align*}

\begin{center}
\shadowbox{
$\mathcal{S}=\left\{t\mapsto \left(
\begin{array}{l}
ae^{2t}+2be^{3t}\\
ae^{2t}+be^{3t}
\end{array}
\right),\;(a,b)\in\Rr^2\right\}$.
}
\end{center}

\item  Puisque la fonction $t\mapsto\left(
\begin{array}{c}
 \frac{1}{\cos t}\\
\rule{0mm}{5mm}0
\end{array}
\right)$ est continue sur $\left]- \frac{\pi}{2}, \frac{\pi}{2}\right[$, les solutions réelles sur $\left]- \frac{\pi}{2}, \frac{\pi}{2}\right[$ du système proposé constituent un $\Rr$-espace affine de dimension $2$.

\textbf{Résolution du système homogène associé.} Posons $A=\left(
\begin{array}{cc}
1&-1\\
2&-1
\end{array}
\right)$. $\chi_A=\lambda^2+1=(\lambda-i)(\lambda+i)$ et en particulier $A$ est diagonalisable dans $\Cc$. Un vecteur propre de $A$ associé à la valeur propre $i$ est $\left(
\begin{array}{c}
1\\
1-i
\end{array}
\right)$ et un vecteur propre de $A$ associé à la valeur propre $-i$ est $\left(
\begin{array}{c}
1\\
1+i
\end{array}
\right)$. On sait alors que les solutions complexes sur $\Rr$ du système homogène associé sont les fonctions de la forme $X~:~t\mapsto ae^{it}\left(
\begin{array}{c}
1\\
1-i
\end{array}
\right)+be^{-it}\left(
\begin{array}{c}
1\\
1+i
\end{array}
\right)$, $(a,b)\in\Cc^2$.

Déterminons alors les solutions réelles du système homogène.

\begin{align*}\ensuremath
X\;\text{réelle}&\Leftrightarrow\forall t\in\Rr,\;ae^{it}\left(
\begin{array}{c}
1\\
1-i
\end{array}
\right)+be^{-it}\left(
\begin{array}{c}
1\\
1+i
\end{array}
\right)=\overline{a}e^{-it}\left(
\begin{array}{c}
1\\
1+i
\end{array}
\right)+\overline{b}e^{it}\left(
\begin{array}{c}
1\\
1-i
\end{array}
\right)\\
 &\Leftrightarrow b=\overline{a}\;(\text{car la famille de fonctions}\;(e^{it},e^{-it})\;\text{est libre}.)
\end{align*}

Les solutions réelles sur $\Rr$ du système homogène sont les fonctions de la forme $X~:~t\mapsto ae^{it}\left(
\begin{array}{c}
1\\
1-i
\end{array}
\right)+\overline{a}e^{-it}\left(
\begin{array}{c}
1\\
1+i
\end{array}
\right)=2\text{Re}\left(ae^{it}\left(
\begin{array}{c}
1\\
1-i
\end{array}
\right)\right)$, $a\in\Cc$. En posant $a=\lambda+i\mu$, $(\lambda,\mu)\in\Rr^2$, 

\begin{center}
$2\text{Re}\left(ae^{it}\left(
\begin{array}{c}
1\\
1-i
\end{array}
\right)\right)=2\text{Re}\left(\left(
\begin{array}{c}
(\lambda+i\mu)(\cos t+i\sin t)\\
(\lambda+i\mu)(1-i)(\cos t+i\sin t)
\end{array}
\right)\right)=2\left(
\begin{array}{c}
\lambda\cos t-\mu\sin t\\
\lambda(\cos t+\sin t)+\mu(\cos t-\sin t)
\end{array}
\right)$.
\end{center}

Maintenant, le couple $(\lambda,\mu)$ décrit $\Rr^2$ si et seulement si le couple $(2\lambda,2\mu)$ décrit $\Rr^2$ et en renommant les constantes $\lambda$ et $\mu$, on obtient les solutions réelles du système homogène : $t\mapsto\lambda\left(\begin{array}{c}
\cos t\\
\cos t+\sin t
\end{array}
\right)+\mu\left(\begin{array}{c}
-\sin t\\
\cos t-\sin t
\end{array}
\right)$, $(\lambda,\mu)\in\Rr^2$.

\textbf{Résolution du système.} D'après la méthode de variation de la constante, il existe une solution particulière du système de la forme $t\mapsto\lambda(t)\left(
\begin{array}{c}
\cos t\\
\cos t+\sin t
\end{array}
\right)+\mu(t)
\left(
\begin{array}{c}
-\sin t\\
\cos t-\sin t
\end{array}
\right)$ où $\lambda$ et $\mu$ sont deux fonctions dérivables sur $\left]- \frac{\pi}{2}, \frac{\pi}{2}\right[$ telles que pour tout réel $t$ de $\left]- \frac{\pi}{2}, \frac{\pi}{2}\right[$, $\lambda'(t)\left(
\begin{array}{c}
\cos t\\
\cos t+\sin t
\end{array}
\right)+\mu'(t)
\left(
\begin{array}{c}
-\sin t\\
\cos t-\sin t
\end{array}
\right)=\left(
\begin{array}{c}
 \frac{1}{\cos t}\\
\rule{0mm}{5mm}0
\end{array}
\right)$. Les formules de \textsc{Cramer} fournissent $\lambda'(t)= \frac{1}{\cos t}(\cos t-\sin t)=1- \frac{\sin t}{\cos t}$ et $\mu'(t)=- \frac{1}{\cos t}(\cos t+\sin t)=-1- \frac{\sin t}{\cos t}$. On peut prendre $\lambda(t)=t+\ln(\cos t)$ et $\mu(t)=-t+\ln(\cos t)$ et on obtient la solution particulière 
\begin{center}
$X(t)=(t+\ln(\cos t))\left(
\begin{array}{c}
\cos t\\
\cos t+\sin t
\end{array}
\right)+(-t+\ln(\cos t))\left(
\begin{array}{c}
-\sin t\\
\cos t-\sin t
\end{array}
\right)=\left(
\begin{array}{c}
t(\cos t+\sin t)+\ln(\cos t)(\cos t-\sin t)\\
2t\sin t+2\cos t\ln(\cos t)
\end{array}.
\right)$.
\end{center}

\begin{center}
\shadowbox{
$\mathcal{S}_{\left]- \frac{\pi}{2}, \frac{\pi}{2}\right[}=\left\{t\mapsto \left(
\begin{array}{l}
t(\cos t+\sin t)+\ln(\cos t)(\cos t-\sin t)+\lambda\cos t-\mu\sin t\\
2t\sin t+2\cos t\ln(\cos t)+\lambda(\cos t+\sin t)+\mu(\cos t-\sin t)
\end{array}
\right),\;(\lambda,\mu)\in\Rr^2\right\}$.
}
\end{center}

\item  Puisque la fonction $t\mapsto\left(
\begin{array}{c}
e^t\\
t
\end{array}
\right)$ est continue sur $\Rr$, les solutions sur $\Rr$ du système proposé constituent un $\Rr$-espace affine de dimension $2$.

\textbf{Résolution du système homogène associé.} Posons $A=\left(
\begin{array}{cc}
5&-2\\
-1&6
\end{array}
\right)$. $\chi_A=\lambda^2-11\lambda+28=(\lambda-4)(\lambda-7)$. Un vecteur propre de ${^t}A$ associé à la valeur propre $4$ est $\left(
\begin{array}{c}
1\\
1
\end{array}
\right)$ et un vecteur propre de ${^t}A$ associé à la valeur propre $7$ est $\left(
\begin{array}{c}
1\\
-2
\end{array}
\right)$. Ces vecteurs fournissent des combinaisons linéaires intéressantes des équations :

\begin{align*}\ensuremath
\left\{
\begin{array}{l}
x'=5x-2y+e^t\\
y'=-x+6y+t
\end{array}
\right.&\Leftrightarrow\left\{
\begin{array}{l}
(x+y)'=4(x+y)+e^t+t\\
(x-2y)'=7(x-2y)+e^t-2t
\end{array}
\right.\\
 &\Leftrightarrow\exists(\lambda,\mu)\in\Rr^2/\;\forall t\in\Rr,\;\left\{
\begin{array}{l}
x(t)+y(t)=- \frac{e^t}{3}- \frac{t}{4}- \frac{1}{16}+\lambda e^{4t}\\
x(t)-2y(t)=- \frac{e^t}{6}+ \frac{2t}{7}+ \frac{2}{49}+\mu e^{7t}
\end{array}
\right.\\
 &\Leftrightarrow\exists(\lambda,\mu)\in\Rr^2/\;\forall t\in\Rr,\;\left\{
\begin{array}{l}
x(t)=- \frac{5e^t}{6}- \frac{3t}{14}- \frac{33}{392}++2\lambda e^{4t}+\mu e^{4t}\\
y(t)=- \frac{e^t}{6} \frac{15t}{28}- \frac{81}{784}+\lambda e^{4t}-\mu e^{7t}
\end{array}
\right.
\end{align*}  

\item  Posons $A=\left(
\begin{array}{ccc}  
5&1&-1\\
2&4&-2\\
1&-1&1
\end{array}
\right)$.

\begin{align*}\ensuremath
\chi_A&=\left|
\begin{array}{ccc}  
5-\lambda&1&-1\\
2&4-\lambda&-2\\
1&-1&1-\lambda
\end{array}
\right|=(5-\lambda)(\lambda^2-5\lambda+2)-2(-\lambda)+(-\lambda+2)=-\lambda^3+10\lambda^2-26\lambda+12\\
 &=-(\lambda-6)(\lambda^2-4\lambda+2)=-(\lambda-6)(\lambda-2+\sqrt{2})(\lambda-2-\sqrt{2}),
\end{align*}

et en particulier $A$ est diagonalisable dans $\Rr$.

\begin{center}
$(x,y,z)\in\text{Ker}(A-6I)\Leftrightarrow\left\{
\begin{array}{l}
-x+y-z=0\\
2x-2y-2z=0\\
x-y-5z=0
\end{array}
\right.\Leftrightarrow z=0\;\text{et}\;x=y$.
\end{center}

$\text{Ker}(A-6I)$ est la droite vectorielle engendrée par le vecteur $(1,1,0)$.

\begin{align*}\ensuremath
(x,y,z)\in\text{Ker}(A-(2+\sqrt{2})I)&\Leftrightarrow\left\{
\begin{array}{l}
(3-\sqrt{2})x+y-z=0\\
2x+(2-\sqrt{2})y-2z=0\\
x-y-(1+\sqrt{2})z=0
\end{array}
\right.\Leftrightarrow\left\{
\begin{array}{l}
z=(3-\sqrt{2})x+y\\
2x+(2-\sqrt{2})y-2((3-\sqrt{2})x+y)=0\\
x-y-(1+\sqrt{2})((3-\sqrt{2})x+y)=0
\end{array}
\right.
\\
 &\Leftrightarrow\left\{
\begin{array}{l}
z=(3-\sqrt{2})x+y\\
(-4+2\sqrt{2})x-\sqrt{2}y=0\\
-2\sqrt{2}x-(2+\sqrt{2})y=0
\end{array}
\right.\Leftrightarrow\left\{
\begin{array}{l}
z=(3-\sqrt{2})x+y\\
y=(-2\sqrt{2}+2)x
\end{array}
\right.\\
 &\Leftrightarrow\left\{
\begin{array}{l}
y=(-2\sqrt{2}+2)x\\
z=(5-3\sqrt{2})x
\end{array}
\right..
\end{align*} $\text{Ker}(A-(2+\sqrt{2})I)$ est la droite vectorielle engendrée par le vecteur $(1,2-2\sqrt{2},5-3\sqrt{2})$. Un calcul conjugué montre alors que $\text{Ker}(A-(2-\sqrt{2})I)$ est la droite vectorielle engendrée par le vecteur $(1,2+2\sqrt{2},5+3\sqrt{2})$.

On sait alors que les solutions du système homogène $t\mapsto ae^{6t}\left(
\begin{array}{c}
1\\
1\\
0
\end{array}
\right)+be^{(2+\sqrt{2})t}\left(
\begin{array}{c}
1\\
2-2\sqrt{2}\\
5-3\sqrt{2})
\end{array}
\right)+ce^{(2-\sqrt{2})t}\left(
\begin{array}{c}
1\\
2+2\sqrt{2}\\
5+3\sqrt{2})
\end{array}
\right)$, $(a,b,c)\in\Rr^3$.

\item  Posons $A=\left(
\begin{array}{ccc}
2&1&0\\
-1&0&0\\
1&1&1
\end{array}
\right)$. $\chi_A=\left|
\begin{array}{ccc}
2-\lambda&1&0\\
-1&-\lambda&0\\
1&1&1-\lambda
\end{array}
\right|=(1-\lambda)(\lambda^2-2\lambda+1)=(1-\lambda)^3$. Le théorème de \textsc{Cayley}-\textsc{Hamilton} permet alors d'affirmer que $(A-I)^3=0$. 

On sait que les solutions du système $X'=AX$ sont les fonctions de la forme $t\mapsto e^{tA}X_0$ où $X_0\in\mathcal{M}_{3,1}(\Rr)$. Or, pour $t\in\Rr$,

\begin{align*}\ensuremath
e^{tA}&=e^{t(A-I)}\times e^{tI}\;(\text{car les matrices}\;t(A-I)\;\text{et}\;tI\;\text{commutent})\\
 &=\left(\sum_{n=0}^{+\infty} \frac{t^n}{n!}(A-I)^n\right)\times e^tI=e^t\left(\sum_{n=0}^{2} \frac{t^n}{n!}(A-I)^n\right)\\
 &=e^t\left(\left(
 \begin{array}{ccc}
 1&0&0\\
 0&1&0\\
 0&0&1
 \end{array}
 \right)+t\left(
 \begin{array}{ccc}
 1&1&0\\
 -1&-1&0\\
 1&1&0
 \end{array}
 \right)+ \frac{t^2}{2}\left(
 \begin{array}{ccc}
 1&1&0\\
 -1&-1&0\\
 1&1&0
 \end{array}
 \right)\left(
 \begin{array}{ccc}
 1&1&0\\
 -1&-1&0\\
 1&1&0
 \end{array}
 \right)\right)\\
 &=e^t\left(
 \begin{array}{ccc}
 1+t&t&0\\
 -t&1-t&0\\
 t&t&1
 \end{array}
 \right).
\end{align*}

Les solutions du système sont les fonctions de la forme $t\mapsto e^{tA}X_0=e^t\left(
 \begin{array}{ccc}
 1+t&t&0\\
 -t&1-t&0\\
 t&t&1
 \end{array}
 \right)\left(
 \begin{array}{c}
 a\\
 b\\
 c
 \end{array}
 \right)=\left(
 \begin{array}{c}
 (a+(a+b)t)e^t\\
 (b-(a+b)t)e^t\\
 ((a+b)t+c)e^t
 \end{array}
 \right)$, $(a,b,c)\in\Rr^3$. Maintenant,
 
\begin{center}
$\left(
\begin{array}{c}
x(0)\\
y(0)\\
z(0)
\end{array}
\right)=\left(
\begin{array}{c}
0\\
1\\
-1
\end{array}
\right)\Leftrightarrow\left\{
\begin{array}{l}
a=0\\
b=1\\
c=-1
\end{array}
\right.$.
\end{center}

La solution cherchée est $t\mapsto\left(
 \begin{array}{c}
 te^t\\
 (1-t)e^t\\
 (t-1)e^t
 \end{array}
 \right)$.
\end{enumerate}
\fincorrection
\correction{005880}
Soit $A\in\mathcal{S}_n^+(\Rr)$. Pour $t\in\Rr$, posons $g(t)=\|X(t)\|_2^2=\left(X(t)|X(t)\right)$. La fonction $g$ est dérivable sur $\Rr$ et pour tout réel $t$

\begin{center}
$g'(t)=2\left(X(t)|X'(t)\right)=2\left(X(t)|AX(t)\right)=2{^t}X(t)AX(t)\geqslant0$.
\end{center}

Ainsi, la fonction $g$ est croissante sur $\Rr$ et il en est de même de la fonction $\sqrt{g}~:~t\mapsto\|X(t)\|_2$.
\fincorrection
\correction{005881}
\begin{enumerate}
 \item  Puisque les fonctions $t\mapsto\left(
\begin{array}{cc}
- \frac{1}{2t}& \frac{1}{2t^2}\\
\rule{0mm}{6mm} \frac{1}{2}& \frac{1}{2t}
\end{array}
\right)$ et $t\mapsto\left(
\begin{array}{c}
2t\\
t^2
\end{array}
\right)$ sont continues sur $]0,+\infty[$, l'ensemble des solutions sur $]0,+\infty[$ du système proposé est un $\Rr$-espace affine de dimension $2$.

\textbf{Résolution du système homogène associé.} Le couple de fonctions $(x,y)=(1,t)$ est solution du système homogène associé sur $]0,+\infty[$. Pour chaque réel strictement positif $t$, les deux vecteurs $\left(
\begin{array}{c}
1\\
t
\end{array}
\right)$ et $\left(
\begin{array}{c}
0\\
1
\end{array}
\right)$ constituent une base de $\mathcal{M}_2(\Rr)$ car $\left|
\begin{array}{cc}
1&0\\
t&1
\end{array}
\right|=1\neq0$. Cherchons alors les solutions du système homogène sous la forme $t\mapsto\alpha(t)\left(
\begin{array}{c}
1\\
t
\end{array}
\right)+\beta(t)\left(
\begin{array}{c}
0\\
1
\end{array}
\right)=\left(
\begin{array}{c}
\alpha(t)\\
t\alpha(t)+\beta(t)
\end{array}
\right)$.

\begin{align*}\ensuremath
\left\{
\begin{array}{l}
x'=- \frac{1}{2t}x+ \frac{1}{2t^2}y\\
y'= \frac{1}{2}x+ \frac{1}{2t}y
\end{array}
\right.&\Leftrightarrow\left\{
\begin{array}{l}
\alpha'=- \frac{1}{2t}\alpha+ \frac{1}{2t^2}(t\alpha+\beta)\\
\rule{0mm}{6mm}t\alpha'+\alpha+\beta'= \frac{1}{2}\alpha+ \frac{1}{2t}(t\alpha+\beta)
\end{array}
\right.\Leftrightarrow\left\{
\begin{array}{l}
\alpha'= \frac{\beta}{2t^2}\\
t\alpha'+\beta'= \frac{\beta}{2t}
\end{array}
\right.\\
 &\Leftrightarrow\left\{
\begin{array}{l}
\alpha'= \frac{\beta}{2t^2}\\
 \frac{\beta}{2t}+\beta'= \frac{\beta}{2t}
\end{array}
\right.\Leftrightarrow\left\{
\begin{array}{l}
\beta'=0\\
\alpha'= \frac{\beta}{2t^2}
\end{array}
\right.\\
 &\Leftrightarrow\exists(\lambda,\mu)\in\Rr^2/\;\forall t\in]0,+\infty[,\;\left\{
\begin{array}{l}
\beta(t)=\lambda\\
\alpha'(t)= \frac{\lambda}{2t^2}
\end{array}
\right.\\
 &\Leftrightarrow\exists(\lambda,\mu)\in\Rr^2/\;\forall t\in]0,+\infty[,\;\left\{
\begin{array}{l}
\beta(t)=\lambda\\
\alpha(t)=- \frac{\lambda}{2t}+\mu
\end{array}
\right.\\
 &\Leftrightarrow\exists(\lambda,\mu)\in\Rr^2/\;\forall t\in]0,+\infty[,\;\left\{
\begin{array}{l}
x(t)=- \frac{\lambda}{2t}+\mu\\
\rule{0mm}{6mm}y(t)= \frac{\lambda}{2}+\mu t
\end{array}
\right.
\end{align*}

Maintenant, pour tout réel strictement positif $t$, $w(t)=\left|
\begin{array}{cc}
- \frac{1}{2t}&1\\
\rule{0mm}{6mm} \frac{1}{2}&t
\end{array}
\right|=-1\neq0$ et donc les deux fonctions $t\mapsto\left(
\begin{array}{c}
- \frac{1}{2t}\\
\rule{0mm}{6mm} \frac{1}{2}
\end{array}
\right)$ et $t\mapsto\left(
\begin{array}{c}
1\\
t
\end{array}
\right)$ sont deux solutions indépendantes du système homogène sur $]0,+\infty[$. Les solutions sur $]0,+\infty[$ du système homogène sont les fonctions de la forme $t\mapsto\lambda\left(
\begin{array}{c}
- \frac{1}{2t}\\
\rule{0mm}{6mm} \frac{1}{2}
\end{array}
\right)
+\mu\left(
\begin{array}{c}
1\\
t
\end{array}
\right)$.

\textbf{Recherche d'une solution particulière du système par la méthode de variations des constantes.}

Il existe une solution particulière du système de la forme $t\mapsto\lambda(t)\left(
\begin{array}{c}
- \frac{1}{2t}\\
\rule{0mm}{6mm} \frac{1}{2}
\end{array}
\right)
+\mu(t)\left(
\begin{array}{c}
1\\
t
\end{array}
\right)$ où $\lambda$ et $\mu$ sont deux fonctions dérivables sur $]0,+\infty[$ telles que pour tout réel strictement positif $t$, $\lambda'(t)\left(
\begin{array}{c}
- \frac{1}{2t}\\
\rule{0mm}{6mm} \frac{1}{2}
\end{array}
\right)
+\mu'(t)\left(
\begin{array}{c}
1\\
t
\end{array}
\right)=\left(
\begin{array}{c}
2t\\
t^2
\end{array}
\right)$. Les formules de \textsc{Cramer} fournissent $\lambda'(t)= \frac{1}{-1}\left|
\begin{array}{cc}
2t&1\\
t^2&t
\end{array}
\right|=-t^2$ et $\mu'(t)= \frac{1}{-1}\left|
\begin{array}{cc}
- \frac{1}{2t}&2t\\
\rule{0mm}{6mm} \frac{1}{2}&t^2
\end{array}
\right|= \frac{3t}{2}$. On peut prendre $\lambda(t)=- \frac{t^3}{3}$ et $\mu(t)= \frac{3t^2}{4}$ et on obtient la solution particulière $X(t)=- \frac{t^3}{3}\left(
\begin{array}{c}
- \frac{1}{2t}\\
\rule{0mm}{6mm} \frac{1}{2}
\end{array}
\right)
+ \frac{3t^2}{4}\left(
\begin{array}{c}
1\\
t
\end{array}
\right)=\left(
\begin{array}{c}
11t^2/12\\
7t^3/12
\end{array}
\right)$

\begin{center}
\shadowbox{
$\mathcal{S}_{]0,+\infty[}=\left\{
t\mapsto\left(
\begin{array}{l}
 \frac{11t^2}{12}- \frac{\lambda}{2t}+\mu\\
\rule{0mm}{7mm} \frac{7t^3}{12}+ \frac{\lambda}{2}+\mu t
\end{array}
\right),\;(\lambda,\mu)\in\Rr^2\right\}$.
}
\end{center}

\item  Puisque les fonctions $t\mapsto \frac{1}{t^2+1}\left(
\begin{array}{cc}
t&-1\\
\rule{0mm}{6mm}1&t
\end{array}
\right)$ et $t\mapsto \frac{1}{t^2+1}\left(
\begin{array}{c}
2t^2-1\\
3t
\end{array}
\right)$ sont continues sur $\Rr$, l'ensemble des solutions sur $\Rr$ du système proposé est un $\Rr$-espace affine de dimension $2$.

\textbf{Résolution du système homogène associé.} Les couples de fonctions $X_1=(x,y)=(t,-1)$ et $(x,y)=(1,t)$ sont solutions du système homogène associé sur $\Rr$. De plus, pour chaque réel $t$, $w(t)=\left|
\begin{array}{cc}
t&1\\
-1&t
\end{array}
\right|=t^2+1\neq0$. Le couple de fonctions $(X_1,X_2)$ est donc un système fondamental de solutions sur $\Rr$ du système homogène $X'=AX$. Les fonctions solutions du système homogène $X'=AX$ sont les fonctions de la forme $t\mapsto\lambda\left(
\begin{array}{c}
t\\
-1
\end{array}
\right)+\mu
\left(
\begin{array}{c}
1\\
t
\end{array}
\right)$, $(\lambda,\mu)\in\Rr^2$.

\textbf{Recherche d'une solution particulière du système par la méthode de variation de la constante.}

Il existe une solution particulière du système de la forme $t\mapsto\lambda(t)\left(
\begin{array}{c}
t\\
-1
\end{array}
\right)
+\mu(t)\left(
\begin{array}{c}
1\\
t
\end{array}
\right)$ où $\lambda$ et $\mu$ sont deux fonctions dérivables sur $\Rr$ telles que pour tout réel $t$, $\lambda'(t)\left(
\begin{array}{c}
t\\
-1
\end{array}
\right)
+\mu'(t)\left(
\begin{array}{c}
1\\
t
\end{array}
\right)= \frac{1}{t^2+1}\left(
\begin{array}{c}
2t^2-1\\
3t
\end{array}
\right)$. Les formules de \textsc{Cramer} fournissent $\lambda'(t)= \frac{1}{t^2+1}\left|
\begin{array}{cc}
(2t^2-1)/(t^2+1)&1\\
3t/(t^2+1)&t
\end{array}
\right|= \frac{2t^3+2t}{(t^2+1)^2}= \frac{2t}{t^2+1}$ et $\mu'(t)= \frac{1}{t^2+1}\left|
\begin{array}{cc}
t&(2t^2-1)/(t^2+1)\\
-1&3t/(t^2+1)
\end{array}
\right|= \frac{5t^2-1}{(t^2+1)^2}$. On peut déjà prendre $\lambda(t)= \frac{1}{2}\ln(t^2+1)$. Ensuite, 
$\int_{}^{} \frac{5t^2-1}{(t^2+1)^2}\;dt=5\int_{}^{} \frac{1}{t^2+1}\;dt-6\int_{}^{} \frac{1}{(t^2+1)^2}\;dt$ puis

\begin{center}
$\int_{}^{} \frac{1}{t^2+1}\;dt= \frac{t}{t^2+1}-\int_{}^{}t\times \frac{-2t}{(t^2+1)^2}\;dt= \frac{t}{t^2+1}+2\int_{}^{} \frac{t^2+1-1}{(t^2+1)^2}\;dt= \frac{t}{t^2+1}+2\int_{}^{} \frac{1}{t^2+1}\;dt-2\int_{}^{} \frac{1}{(t^2+1)^2}\;dt$,
\end{center}

et donc $\int_{}^{} \frac{1}{(t^2+1)^2}\;dt= \frac{1}{2}\left( \frac{t}{t^2+1}+\Arctan t\right)+C$. On peut prendre $\mu(t)= \frac{2t}{t^2+1}-3\Arctan t$.

\begin{center}
\shadowbox{
$\mathcal{S}_{\Rr}=\left\{
t\mapsto\left(
\begin{array}{l}
 \frac{t}{2}\ln(1+t^2)+ \frac{2t}{t^2+1}-3\Arctan t+\lambda t+\mu\\
- \frac{t}{2}\ln(1+t^2)+ \frac{2t^2}{t^2+1}-3t\Arctan t-\lambda+\mu t
\end{array}
\right),\;(\lambda,\mu)\in\Rr^2\right\}$.
}
\end{center}

\item  Si de plus $y= \frac{1}{x}$, le système s'écrit $\left\{
\begin{array}{l}
\sh(2t)x'=\ch(2t)x- \frac{1}{x}\\
-\sh(2t) \frac{x'}{x^2}=-x+\ch(2t) \frac{1}{x}
\end{array}
\right.$ ou encore $\left\{
\begin{array}{l}
\sh(2t)x'=\ch(2t)x- \frac{1}{x}\\
\sh(2t)x'=x^3-\ch(2t)x
\end{array}
\right.$. On obtient $x^3-\ch(2t)x=\ch(2t)x- \frac{1}{x}$ ou encore $x^4-2\ch(2t)x^2+1=0$. Ensuite,

\begin{center}
$x^4-2\ch(2t)x^2+1=(x^2-\ch(2t))^2-\sh^2(2t)=(x^2-e^{2t})(x^2-e^{-2t})=(x-e^{t})(x+e^{t})(x-e^{-t})(x+e^{-t})$.
\end{center}

Ainsi, nécessairement $(x,y)\in\left\{(e^{t},e^{-t}),(e^{-t},e^t),(-e^t,-e^{-t}),(-e^{-t},e^t)\right\}$. Réciproquement, si $(x,y)=(e^t,e^{-t})$, 

\begin{center}
$\ch(2t)x-y= \frac{1}{2}(e^{3t}+e^{-t})-e^{-t}= \frac{1}{2}(e^{3t}-e^{-t})= \frac{1}{2}(e^{2t}-e^{-2t})e^t=\sh(2t)e^t=\sh(2t)x'$
\end{center}

et

\begin{center}
$-x+\ch(2t)y=-e^t+ \frac{1}{2}(e^{t}+e^{-3t})= \frac{1}{2}(-e^{t}+e^{-3t})=- \frac{1}{2}(e^{2t}-e^{-2t})e^{-t}=-\sh(2t)e^{-t}=\sh(2t)y'$.
\end{center}

Donc le couple $X_1=(x,y)=(e^t,e^{-t})$ est une solution non nulle du système. De même, si $(x,y)=(e^{-t},e^t)$,

\begin{center}
$\ch(2t)x-y= \frac{1}{2}(e^{t}+e^{-3t})-e^{t}= \frac{1}{2}(-e^{t}-e^{-3t})=- \frac{1}{2}(e^{2t}-e^{-2t})e^{-t}=-\sh(2t)e^{-t}=\sh(2t)x'$
\end{center}

et

\begin{center}
$-x+\ch(2t)y=-e^{-t}+ \frac{1}{2}(e^{3t}+e^{-t})= \frac{1}{2}(e^{3t}-e^{-t})= \frac{1}{2}(e^{2t}-e^{-2t})e^{t}=\sh(2t)e^{t}=\sh(2t)y'$.
\end{center}

Donc le couple $X_2=(x,y)=(e^{-t},e^{t})$ est une solution non nulle du système. Enfin, $w(t)=\left|
\begin{array}{cc}
e^t&e^{-t}\\
e^{-t}&e^t
\end{array}
\right|=e^{2t}-e^{-2t}=2\sh(2t)\neq0$ et le couple $(X_1,X_2)$ est un système fondamental de solutions sur $]0,+\infty[$.

\begin{center}
\shadowbox{
$\mathcal{S}_{]0,+\infty[}=\left\{
t\mapsto\left(
\begin{array}{c}
\lambda e^t+\mu e^{-t}\\
\lambda e^{-t}+\mu e^t
\end{array}
\right),\;(\lambda,\mu)\in\Rr^2\right\}$.
}
\end{center}
\end{enumerate}
\fincorrection
\correction{005882}
\begin{enumerate}
 \item  Sur $I=\left]- \frac{1}{2},+\infty\right[$, $(E)$ s'écrit $y''+ \frac{4x-2}{2x+1}y'- \frac{8}{2x+1}y=0$. Puisque les deux fonctions $x\mapsto \frac{4x-2}{2x+1}$ et $x\mapsto- \frac{8}{2x+1}$ sont continues sur $I$, les solutions de $(E)$ sur $I$ forment un$\Rr$-espace vectoriel de dimension $2$.

\textbf{Recherche d'une solution polynomiale non nulle de $(E)$.} Soit $P$ un éventuel polynôme non nul solution de $(E)$. On note $n$ son degré. Le polynôme $Q=(2X+1)P''+(4X-2)P'-8P$ est de degré au plus $n$. De plus, le coefficient de $X^n$ dans $Q$ est $(4n-8)\text{dom}(P)$. Si $P$ est solution de $(E)$, on a nécessairement $(4n-8)\text{dom}(P)=0$ et donc $n=2$.

Posons alors $P=aX^2+bX+c$.

\begin{center}
$(2X+1)P''+(4X-2)P'-8P=(2X+1)(2a)+(4X-2)(2aX+b)-8(aX^2+bX+c)=-4bX+2a-2b-8c$
\end{center}

Par suite, $P$ est solution de $(E)$ sur $I$ si et seulement si $-4b=2a-2b-8c=0$ ce qui équivaut à $b=0$ et $a=4c$. La fonction $f_1~:~x\mapsto4x^2+1$ est donc une solution non nulle de $(E)$ sur $I$.

\textbf{Recherche d'une solution particulière de la forme $f_\alpha~:~x\mapsto e^{\alpha x}$, $\alpha\in\Cc$.}

\begin{center}
$(2x+1)(e^{\alpha x})''+(4x-2)(e^{\alpha x})'-8e^{\alpha x}=\left(\alpha^2(2x+1)+\alpha(4x-2)-8\right)e^{\alpha x}=\left(2\alpha(\alpha+2)x+\alpha^2-2\alpha-8\right)e^{\alpha x}$
\end{center}

Par suite, $f_{\alpha}$ est solution de $(E)$ sur $I$ si et seulement si $2\alpha(\alpha+2)=\alpha^2-2\alpha-8=0$ ce qui équivaut à $\alpha=-2$. Ainsi, la fonction $f_2~:~x\mapsto e^{-2x}$ est solution de $(E)$ sur $I$.

\textbf{Résolution de $(E)$ sur $\left]- \frac{1}{2},+\infty\right[$.} Vérifions que le couple $(f_1,f_2)$ est un système fondamental de solution de $(E)$ sur $\left]- \frac{1}{2},+\infty\right[$. Pour $x>- \frac{1}{2}$,

\begin{center}
$w(x)=\left|
\begin{array}{cc}
4x^2+1&e^{-2x}\\
8x&-2e^{-2x}
\end{array}
\right|
=(-8x^2-8x-2)e^{-2x}=-2(2x+1)^2e^{-2x}\neq0$.
\end{center}

Donc le couple $(f_1,f_2)$ est un système fondamental de solution de $(E)$ sur $\left]- \frac{1}{2},+\infty\right[$ et

\begin{center}
\shadowbox{
$\mathcal{S}_{\left]-\frac{1}{2},+\infty\right[}=\left\{x\mapsto\lambda(4x^2+1)+\mu e^{-2x},\;(\lambda,\mu)\in\Rr^2\right\}$.
}
\end{center}

\textbf{Résolution de $(E)$ sur $\Rr$.} On a aussi $\mathcal{S}_{\left]-\infty,-\frac{1}{2}\right[}=\left\{x\mapsto\lambda(4x^2+1)+\mu e^{-2x},\;(\lambda,\mu)\in\Rr^2\right\}$. Soit $f$ une solution de $(E)$ sur $\Rr$. Nécessairement, il existe $(\lambda_1,\lambda_2,\mu_1,\mu_2)\in\Rr^4$ tel que $\forall x\in\Rr$, $f(x)=\left\{
\begin{array}{l}
\lambda_1(4x^2+1)+\mu_1 e^{-2x}\;\text{si}\;x\leqslant- \frac{1}{2}\\
\rule{0mm}{6mm}\lambda_2(4x^2+1)+\mu_2 e^{-2x}\;\text{si}\;x>- \frac{1}{2}
\end{array}
\right.$ (par continuité à gauche en $- \frac{1}{2}$).

$f$ ainsi définie est deux fois dérivables sur $\left]-\infty,- \frac{1}{2}\right[$ et sur $\left]- \frac{1}{2},+\infty\right[$, solution de $(E)$ sur chacun de ces deux intervalles et vérifie encore $(E)$ en $x=- \frac{1}{2}$ si de plus $f$ est deux fois dérivable en $- \frac{1}{2}$. 

En résumé, $f$ est solution de $(E)$ sur $\Rr$ si et seulement si $f$ est deux fois dérivables en $- \frac{1}{2}$.

$f$ est déjà deux fois dérivable à droite et à gauche en $- \frac{1}{2}$. De plus, en posant $h=x+ \frac{1}{2}$ ou encore $x=- \frac{1}{2}+h$, on obtient quand $x$ tend vers $- \frac{1}{2}$ par valeurs inférieures

\begin{center}
$f(x)=\lambda_1(2-4h+4h^2)+\mu_1ee^{-2h}=(2\lambda_1+e\mu_1)+(-4\lambda_1-2e\mu_1)h+(4\lambda_1+2e\mu_1)h^2+o(h^2)$,
\end{center}

et de même quand $x$ tend vers $- \frac{1}{2}$ par valeurs supérieures, $f(x)=(2\lambda_2+e\mu_2)+(-4\lambda_2-2e\mu_2)h+(4\lambda_2+2e\mu_2)h^2+o(h^2)$. Par suite, $f$ est deux fois dérivable en $- \frac{1}{2}$ si et seulement si $2\lambda_1+e\mu_1=2\lambda_2+e\mu_2$ ou encore $\mu_2= \frac{2}{e}(\lambda_1+\lambda_2)+\mu_1$.

Ainsi, les solutions de $(E)$ sur $\Rr$ sont les fonctions de la forme $x\mapsto\left\{
\begin{array}{l}
a(4x^2+1)+b e^{-2x}\;\text{si}\;x\leqslant- \frac{1}{2}\\
\rule{0mm}{6mm}c(4x^2+1)+\left( \frac{2}{e}(a+c)-b\right)e^{-2x}\;\text{si}\;x>- \frac{1}{2}
\end{array}
\right.$, $(a,b,c)\in\Rr^3$. Ainsi, l'espace des solutions sur $\Rr$ est de dimension $3$ et une base de cet espace est par exemple $(f_1,f_2,f_3)$ où $f_1~:~x\mapsto\left\{
\begin{array}{l}
4x^2+1\;\text{si}\;x\leqslant- \frac{1}{2}\\
\rule{0mm}{6mm} \frac{2}{e}e^{-2x}\;\text{si}\;x>- \frac{1}{2}
\end{array}
\right.$, $f_2~:~x\mapsto\left\{
\begin{array}{l}
e^{-2x}\;\text{si}\;x\leqslant- \frac{1}{2}\\
-e^{-2x}\;\text{si}\;x>- \frac{1}{2}
\end{array}
\right.$ et $f_3~:~x\mapsto\left\{
\begin{array}{l}
0\;\text{si}\;x\leqslant- \frac{1}{2}\\
4x^2+1+ \frac{2}{e}e^{-2x}\;\text{si}\;x>- \frac{1}{2}
\end{array}
\right.$.

\item  Sur $I=]0,+\infty[$, l'équation $(E)$ s'écrit $y''- \frac{2}{x+1}y'+ \frac{2}{x(x+1)}y=0$. Puisque les deux fonctions $x\mapsto- \frac{2}{x+1}$ et $x\mapsto \frac{2}{x(x+1)}$ sont continues sur $I$, les solutions de $(E)$ sur $I$ forment un $\Rr$-espace vectoriel de dimension $2$.

La fonction $f_1~:~x\mapsto x$ est solution de $(E)$ sur $I$. Posons alors $y=f_1z$. Puisque la fonction $f_1$ ne s'annule pas sur $I$, la fonction $y$ est deux fois dérivables sur $I$ si et seulement si la fonction $z$ est deux fois dérivables sur $I$. De plus, d'après la formule de \textsc{Leibniz},

\begin{align*}\ensuremath
(x^2+x)y''-2y'+2y&=(x^2+x)(f_1''z+2f_1'z'+f_1z'')-2x(f_1'z+f_1z')+2f_1z\\
 &=(x^2+x)f_1z''+(2(x^2+x)f_1'-2xf_1)z'+((x^2+x)f_1''-2f_1'+2f_1)z\\
 &=(x^3+x^2)z''+2xz'.
\end{align*}

Par suite,

\begin{align*}\ensuremath
y\;\text{solution de}\;(E)\;\text{sur}\;I&\Leftrightarrow\forall x\in I,\;(x^3+x^2)z''(x)+2xz'(x)=0\\
 &\Leftrightarrow\forall x\in I,\;z''(x)+ \frac{2}{x(x+1)}z'(x)=0\Leftrightarrow\forall x\in I,\;\left(e^{2\ln|x|-2\ln|x+1|}z'\right)'(x)=0\\
  &\Leftrightarrow\exists\lambda\in\Rr/\;\forall x\in I,\;z'(x)=\lambda\left( \frac{x+1}{x}\right)^2\Leftrightarrow\exists(\lambda,\mu)\in\Rr^2/\;\forall x\in I,\;z(x)=\lambda\left(x+2\ln|x|- \frac{1}{x}\right)+\mu\\
  &\Leftrightarrow\exists(\lambda,\mu)\in\Rr^2/\;\forall x\in I,\;y(x)=\lambda(x^2+2x\ln|x|-1)+\mu x.
\end{align*}

\item  Cherchons les solutions développables en série entière. Soit $f(x)=\sum_{n=0}^{+\infty}a_nx_n$ une série entière dont le rayon $R$ est supposé à priori strictement positif. Pour $x\in]-R,R[$,

\begin{align*}\ensuremath
4xf''(x)-2f'(x)+9x^2f(x)&=4x\sum_{n=2}^{+\infty}n(n-1)a_nx^{n-2}-2\sum_{n=1}^{+\infty}na_nx^{n-1}+9x^2\sum_{n=0}^{+\infty}a_nx^n\\
 &=4\sum_{n=1}^{+\infty}n(n-1)a_nx^{n-1}-2\sum_{n=1}^{+\infty}na_nx^{n-1}+9\sum_{n=0}^{+\infty}a_nx^{n+2}\\
  &=\sum_{n=1}^{+\infty}2n(2n-3)a_nx^{n-1}+9\sum_{n=0}^{+\infty}a_nx^{n+2}=\sum_{n=1}^{+\infty}2n(2n-3)a_nx^{n-1}+9\sum_{n=3}^{+\infty}a_{n-3}x^{n-1}\\
  &=-a_1+4a_2x+\sum_{n=3}^{+\infty}(2n(2n-3)a_n+9a_{n-3})x^{n-1}
\end{align*}

Par suite, $f$ est solution de $(E)$ sur $]-R,R[$ si et seulement si $a_1=a_2=0$ et $\forall n\geqslant3$, $2n(2n-3)a_n+9a_{n-3}=0$ ce qui s'écrit encore 

\begin{center}
$a_1=a_2=0$ et $\forall n\geqslant3$, $a_n=- \frac{9}{2n(2n-3)}a_{n-3}$.
\end{center}

Les conditions $a_1=0$ et $\forall n\geqslant3$, $a_n=- \frac{9}{2n(2n-3)}a_{n-3}$ sont équivalentes à $\forall p\in\Nn$, $a_{3p+1}=0$ et les conditions $a_2=0$ et $\forall n\geqslant3$, $a_n=- \frac{9}{2n(2n-3)}a_{n-3}$ sont équivalentes à $\forall p\in\Nn$, $a_{3p+2}=0$.

Enfin les conditions $\forall p\in\Nn^*$, $a_{3p}=- \frac{9}{6p(6p-3)}a_{3p-3}=- \frac{1}{2p(2p-1)}a_{3(p-1)}$ sont équivalentes pour $p\geqslant1$ à

\begin{center}
$a_{3p}=- \frac{1}{2p(2p-1)}\times- \frac{1}{(2p-2)(2p-3)}\times\ldots\times- \frac{1}{2\times1}a_0= \frac{(-1)^p}{(2p)!}a_0$.
\end{center}

En résumé, sous l'hypothèse $R>0$, $f$ est solution de $(E)$ sur $]-R,R[$ si et seulement si $\forall x\in]-R,R[$, $f(x)=\sum_{p=0}^{+\infty} \frac{(-1)^p}{(2p)!}x^{3p}$.

Réciproquement, puisque pour tout réel $x$, $\lim_{x \rightarrow +\infty} \frac{(-1)^p}{(2p)!}x^{3p}=0$ d'après un théorème de croissances comparées, $R=+\infty$ pour tout choix de $a_0$ ce qui valide les calculs précédents sur $\Rr$.

Les solutions de $(E)$ développables en série entière sont les fonctions de la forme $x\mapsto\lambda\sum_{n=0}^{+\infty} \frac{(-1)^n}{(2n)!}x^{3n}$, $x\in\Rr$. Ensuite, pour $x>0$,

\begin{center}
$\sum_{n=0}^{+\infty} \frac{(-1)^n}{(2n)!}x^{3n}=\sum_{n=0}^{+\infty} \frac{(-1)^n}{(2n)!}(x^{3/2})^{2n}=\cos\left(x^{3/2}\right)$.
\end{center}

Donc la fonction $x\mapsto\cos\left(x^{3/2}\right)$ est une solution de $(E)$ sur $]0,+\infty[$. La forme de cette solution nous invite à changer de variable en posant $t=x^{3/2}$. Plus précisément, pour $x>0$, posons $y(x)=z(x^{3/2})=z(t)$. Puisque l'application $\varphi~:~x\mapsto x^{3/2}$ est un $C^2$-difféomorphisme de $]0,+\infty[$ sur lui-même, la fonction $y$ est deux fois dérivables sur $]0,+\infty[$ si et seulement si la fonction est deux fois dérivable sur $]0,+\infty[$.

Pour $x>0$, on a $y(x)=z(x^{3/2})$ puis $y'(x)= \frac{3}{2}x^{1/2}z'(x^{3/2})$ puis $y''(x)= \frac{3}{4}x^{-1/2}z'(x^{3/2})+ \frac{9}{4}xz''(x^{3/2})$ et donc

\begin{align*}\ensuremath
4xy''(x)- 2y'(x)+ 9x^2y(x)&=4x\left( \frac{3}{4}x^{-1/2}z'(x^{3/2})+ \frac{9}{4}xz''(x^{3/2})\right)-2\left( \frac{3}{2}x^{1/2}z'(x^{3/2})\right)+9x^2z(x^{3/2})\\
 &=9x^2(z''(x^{3/2})+z(x^{3/2})).
\end{align*}

Par suite,

\begin{align*}\ensuremath
y\;\text{solution de}\;(E)\;\text{sur}\;]0,+\infty[&\Leftrightarrow\forall x>0,\;9x^2(z''(x^{3/2})+z(x^{3/2}))=0\Leftrightarrow\forall t>0,\;z''(t)+z(t)=0\\
 &\Leftrightarrow\exists(\lambda,\mu)\in\Rr^2/\;\forall t>0,\;z(t)=\lambda\cos t+\mu\sin t\\
 &\Leftrightarrow\exists(\lambda,\mu)\in\Rr^2/\;\forall x>0,\;y(x)=\lambda\cos(x^{3/2})+\mu\sin(x^{3/2}).
\end{align*}

\begin{center}
\shadowbox{
$\mathcal{S}_{]0,+\infty[}=\left\{x\mapsto\lambda\cos(x^{3/2})+\mu\sin(x^{3/2}),\;(\lambda,\mu)\in\Rr^2\right\}$.
}
\end{center}

\item  Puisque les fonctions $x\mapsto- \frac{2}{1+x}$, $x\mapsto \frac{1-x}{1+x}$ et $x\mapsto \frac{xe^{-x}}{1+x}$ sont continues sur $]-1,+\infty[$, les solutions de $(E)$ sur $]-1,+\infty[$ constituent un $\Rr$-espace affine de dimension $2$.

\textbf{Résolution de l'équation homogène.} 

La fonction $f_1~:~x\mapsto e^x$ est solution sur $]-1,+\infty[$ de l'équation $(1+x)y''-2y'+(1-x)y=0$. Posons alors $y=f_1z$. Puisque la fonction $f_1$ ne s'annule pas sur $]-1,+\infty[$, la fonction $y$ est deux fois dérivable sur $]-1,+\infty[$ si et seulement si la fonction $z$ est deux fois dérivable sur $]-1,+\infty[$. De plus, la formule de \textsc{Leibniz} permet d'écrire pour $x>-1$

\begin{align*}\ensuremath
(1+x)y''(x)-2y'(x)+(1-x)y(x)&=(1+x)(f_1''z(x)+2f_1'(x)z'(x)+f_1(x)z''(x))-2(f_1'(x)z(x)+f_1(x)z'(x))\\
 &+(1-x)f_1(x)z(x)\\
 &=(1+x)f_1(x)z''(x)+(2(1+x)f_1'(x)-2f_1(x))z'(x)=((1+x)z''(x)+2xz'(x))e^x.
\end{align*}

Par suite,

\begin{align*}\ensuremath
y\;\text{solution de}\;(E_H)\;\text{sur}\;]-1,+\infty[&\Leftrightarrow\forall x>-1,\;(1+x)z''(x)+2xz'(x)=0\Leftrightarrow\forall x>-1,\;z''(x)+\left(2- \frac{2}{1+x}\right)z'(x)=0\\
 &\Leftrightarrow\forall x>-1,\;e^{2x-2\ln(1+x)}z''(x)+\left(2- \frac{2}{1+x}\right)e^{2x-2\ln(1+x)}z'(x)=0\\
 &\Leftrightarrow\forall x>-1,\;\left( \frac{e^{2x}}{(x+1)^2}z'\right)'(x)=0\Leftrightarrow\exists\lambda\in\Rr/\;\forall x>-1,\;z'(x)=\lambda(x+1)^2e^{-2x}.
\end{align*}

Maintenant

\begin{align*}\ensuremath
\int_{}^{}(x+1)^2e^{-2x}\;dx&=- \frac{1}{2}(x+1)^2e^{-2x}+\int_{}^{}(x+1)e^{-2x}\;dx=- \frac{1}{2}(x+1)^2e^{-2x}- \frac{1}{2}(x+1)e^{-2x}+ \frac{1}{2}\int_{}^{}e^{-2x}\;dx\\
 &=\left(- \frac{1}{2}(x+1)^2- \frac{1}{2}(x+1)- \frac{1}{4}\right)e^{-2x}+C=\left(- \frac{x^2}{4}- \frac{3x}{2}- \frac{5}{4}\right)e^{-2x}+C=- \frac{1}{4}(2x^2+6x+5)e^{-2x}+C
\end{align*}

On en déduit que

\begin{align*}\ensuremath
y\;\text{solution de}\;(E_H)\;\text{sur}\;]-1,+\infty[&\Leftrightarrow\exists\lambda\in\Rr/\;\forall x>-1,\;z'(x)=\lambda(x+1)^2e^{-2x}\\
 &\Leftrightarrow\exists(\lambda,\mu)\in\Rr^2/\;\forall x>-1,\;z(x)=- \frac{\lambda}{4}(2x^2+6x+5)e^{-2x}+\mu\\
 &\Leftrightarrow\exists(\lambda,\mu)\in\Rr^2/\;\forall x>-1,\;y(x)=- \frac{\lambda}{4}(2x^2+6x+5)e^{-x}+\mu e^x.
\end{align*}

Maintenant, $\lambda$ décrit $\Rr$ si et seulement si $- \frac{\lambda}{4}$ décrit $\Rr$ et en renommant la constante $\lambda$, les solutions de $(E_H)$ sur $]-1,+\infty[$ sont les fonctions de la forme $x\mapsto\lambda(2x^2+6x+5)e^{-x}+\mu e^x$, $(\lambda,\mu)\in\Rr^2$.

\textbf{Recherche d'une solution particulière de $(E)$.} Au vu du second membre, on peut chercher une solution particulière de la forme $f_0~:~x\mapsto(ax+b)e^{-x}$, $(a,b)\in\Rr^2$.

\begin{align*}\ensuremath
(1+x)((ax+b)e^{-x})''-2((ax+b)e^{-x})'+(1-x)(ax+b)e^{-x}&=((1+x)((ax+b)-2a)-2(-(ax+b)+a)\\
 &+(1-x)(ax+b))e^{-x}\\
 &=(2bx+(4b-4a))e^{-x}.
\end{align*}

Par suite, $f_0$ est solution de $(E)$ sur $]-1,+\infty[$ si et seulement si $2b=1$ et $4b-4a=0$ ce qui équivaut à $a=b= \frac{1}{2}$. Une solution de $(E)$ sur $]-1,+\infty[$ est $x\mapsto \frac{x+1}{2}e^{-x}$.

\begin{center}
\shadowbox{
$\mathcal{S}_{]-1,+\infty[}=\left\{x\mapsto\lambda(2x^2+6x+5)e^{-x}+\mu e^x+ \frac{x+1}{2}e^{-x},\;(\lambda,\mu)\in\Rr^2\right\}$.
}
\end{center}

\item  Puisque la fonction $x\mapsto \frac{e^{-2x}}{\sqrt{x^2+1}}$ est continue sur $\Rr$, les solutions de $(E)$ sur $\Rr$ constituent un $\Rr$-espace affine de dimension $2$.

L'équation caractéristique de l'équation homogène est $z^2+4z+4=0$. Puisque cette équation admet $-2$ pour racine double, les solutions de l'équation homogène associée sont les fonctions de la forme $x\mapsto\lambda e^{-2x}+\mu xe^{-2x}$, $(\lambda,\mu)\in\Rr^2$.

D'après la méthode de variation de la constante, il existe une solution particulière de $(E)$ sur $\Rr$ de la forme $x\mapsto\lambda(x)e^{-2x}+\mu(x)xe^{-2x}$ où $\lambda$ et $\mu$ sont deux fonctions dérivables sur $\Rr$ telles que 

\begin{center}
$\left\{
\begin{array}{l}
\lambda'(x)e^{-2x}+\mu'(x)xe^{-2x}=0\\
-2\lambda'(x)e^{-2x}+\mu'(x)(-2x+1)e^{-2x}= \frac{e^{-2x}}{\sqrt{x^2+1}}
\end{array}
\right.$.
\end{center}

Les formules de \textsc{Cramer} fournissent $\lambda'(x)= \frac{1}{e^{-4x}}\left|
\begin{array}{cc}
0&xe^{-2x}\\
 \frac{e^{-2x}}{\sqrt{x^2+1}}&(-2x+1)e^{-2x}
\end{array}
\right|=- \frac{x}{\sqrt{x^2+1}}$ et 

$\mu'(x)= \frac{1}{e^{-4x}}\left|
\begin{array}{cc}
e^{-2x}&0\\
-2e^{-2x}& \frac{e^{-2x}}{\sqrt{x^2+1}}
\end{array}
\right|= \frac{1}{\sqrt{x^2+1}}$. On peut prendre $\lambda(x)=-\sqrt{x^2+1}$ et $\mu(x)=\Argsh(x)=\ln\left(x+\sqrt{x^2+1}\right)$ puis $f_0(x)=\left(-\sqrt{x^2+1}+x\ln\left(x+\sqrt{x^2+1}\right)\right)e^{-2x}$.

\begin{center}
\shadowbox{
$\mathcal{S}_\Rr=\left\{x\mapsto\left(\lambda+\mu x+\left(-\sqrt{x^2+1}+x\ln\left(x+\sqrt{x^2+1}\right)\right)\right)e^{-2x},\;(\lambda,\mu)\in\Rr^2\right\}$.
}
\end{center}
\end{enumerate}
\fincorrection
\correction{005883}
Soit $f$ une éventuelle solution. $f$ est dérivable sur $\Rr$ et pour tout réel $x$, $f'(x)=-f(-x)+e^{x}$. On en déduit que $f'$ est dérivable sur $\Rr$ ou encore que $f$ est deux fois dérivable sur $\Rr$. En dérivant l'égalité initiale, on obtient pour tout réel $x$

\begin{center}
$f''(x)=f'(-x)+e^{x}=-f(x)+e^{-x}+e^x$,
\end{center}

et donc $f$ est solution sur $\Rr$ de l'équation différentielle $y''+y=2\ch(x)$. Par suite, il existe $(\lambda,\mu)\in\Rr^2$ tel que $\forall x\in\Rr$, $f(x)=\ch(x)+\lambda\cos(x)+\mu\sin(x)$.

Réciproquement, soit $f$ une telle fonction. $f$ est dérivable sur $\Rr$ et pour tout réel $x$,

\begin{center}
$f'(x)+f(-x)=(\sh(x)-\lambda\sin(x)+\mu\cos(x))+(\ch(x)+\lambda\cos(x)-\mu\sin(x))=e^x+(\lambda+\mu)(\cos(x)-\sin(x))$,
\end{center}

et $f$ est solution si et seulement si $\lambda+\mu=0$.

Les fonctions solutions sont les fonctions de la forme $x\mapsto\ch(x)+\lambda(\cos(x)-\sin(x))$, $\lambda\in\Rr$.
\fincorrection
\correction{005884}
Soit $f$ une éventuelle solution. $f$ est dérivable sur $]0,+\infty[$ et pour tout réel $x>0$, $f'(x)=f\left( \frac{3}{16x}\right)$. On en déduit que $f'$ est dérivable sur $]0,+\infty[$ ou encore que $f$ est deux fois dérivable sur $]0,+\infty[$. En dérivant l'égalité initiale, on obtient pour tout réel $x$

\begin{center}
$f''(x)=- \frac{3}{16x^2}f'\left( \frac{3}{16x}\right)=- \frac{3}{16x^2}f\left( \frac{3/16}{(3/16)/x}\right)=- \frac{3}{16x^2}f(x)$,
\end{center}

et donc $f$ est solution sur $\Rr$ de l'équation différentielle $x^2y''+ \frac{3}{16}y=0$ $(E)$. Les solutions de $(E)$ sur $]0,+\infty[$ constituent un $\Rr$-espace vectoriel de dimension $2$. Cherchons une solution particulière de $(E)$ sur $]0,+\infty[$ de la forme $g_\alpha~:~x\mapsto x^\alpha$, $\alpha\in\Rr$.

\begin{align*}\ensuremath
f_\alpha\;\text{solution de}\;(E)\;\text{sur}\;]0,+\infty[&\Leftrightarrow\forall x>0,\;x^2\alpha(\alpha-1)x^{\alpha-2}+ \frac{3}{16}x^\alpha=0\Leftrightarrow\alpha^2-\alpha+ \frac{3}{16}=0\\
 &\Leftrightarrow\alpha= \frac{1}{4}\;\text{ou}\;\alpha= \frac{3}{4}.
\end{align*}

Les deux fonctions $f_1~:~x\mapsto x^{1/4}$ et $f_2~:~x\mapsto x^{3/4}$ sont solutions de $(E)$ sur $]0,+\infty[$. Le wronskien de ces solutions est $w(x)=\left|
\begin{array}{cc}
x^{1/4}&x^{3/4}\\
 \frac{1}{4}x^{-3/4}& \frac{3}{4}x^{-1/4}
\end{array}
\right|= \frac{1}{2}\neq0$ et donc $(f_1,f_2)$ est un système fondamental de solutions de $(E)$ sur $]0,+\infty[$. Ainsi, si $f$ est solution du problème, nécessairement $\exists(\lambda_1,\lambda_2)\in\Rr^2$ tel que $\forall x>0$, $f(x)=\lambda_1x^{1/4}+\lambda_2x^{3/4}$.

Réciproquement, soit $f$ une telle fonction.

Pour tout réel $x>0$, $f'(x)= \frac{\lambda_1}{4}x^{-3/4}+ \frac{3\lambda_2}{4}x^{-1/4}$ et $f\left( \frac{3}{16x}\right)= \frac{3^{1/4}\lambda_1}{2}x^{-1/4}+ \frac{3^{3/4}\lambda_2}{8}x^{-3/4}$. Donc

\begin{align*}\ensuremath
f\;\text{solution}&\Leftrightarrow\forall x>0,\; \frac{\lambda_1}{4}x^{-3/4}+ \frac{3\lambda_2}{4}x^{-1/4}= \frac{3^{1/4}\lambda_1}{2}x^{-1/4}+ \frac{3^{3/4}\lambda_2}{8}x^{-3/4}\\
 &\Leftrightarrow\forall x>0,\; \frac{\lambda_1}{4}x^{-3/4}+ \frac{3\lambda_2}{4}x^{-1/4}= \frac{3^{1/4}\lambda_1}{2}x^{-1/4}+ \frac{3^{3/4}\lambda_2}{8}x^{-3/4}\\
 &\Leftrightarrow\forall x>0,\; \frac{\lambda_1}{4}x^{1/4}+ \frac{3\lambda_2}{4}x^{3/4}= \frac{3^{1/4}\lambda_1}{2}x^{3/4}+ \frac{3^{3/4}\lambda_2}{8}x^{1/4}\;(\text{après multiplication par}\;x)\\
 &\Leftrightarrow \frac{\lambda_1}{4}= \frac{3^{3/4}\lambda_2}{8}\;\text{et}\; \frac{3\lambda_2}{4}= \frac{3^{1/4}\lambda_1}{2}\Leftrightarrow\lambda_2= \frac{2}{3^{3/4}}\lambda_1.
\end{align*}

Les fonctions solutions sont les fonctions de la forme $x\mapsto\lambda\left(x^{1/4}+2\left( \frac{x}{3}\right)^{3/4}\right)$.
\fincorrection
\correction{005885}
La fonction nulle est solution. Dorénavant, $f$ est une éventuelle solution non nulle.  Il existe donc $x_0\in\Rr$ tel que $f(x_0)\neq0$. 

\textbullet~L'égalité $f(x_0)f(0)=\int_{x_0-0}^{x_0+0}f(t)\;dt=0$ fournit $f(0)=0$.

\textbullet~Pour tout réel $y$, $\int_{-y}^{y}f(t)\;dt=f(0)f(y)=0$. Maintenant, la fonction $f$ est continue sur $\Rr$ et donc la fonction $y\mapsto \int_{-y}^{y}f(t)\;dt$ est de classe $C^1$ sur $\Rr$. En dérivant, on obtient pour tout réel $y$, $f(y)+f(-y)=0$ et donc $f$ est impaire.

\textbullet~Pour tout réel $y$, on a alors $f(y)= \frac{1}{f(x_0)}\int_{x_0-y}^{x_0+y}f(t)\;dt$. Puisque $f$ est continue sur $\Rr$, la fonction $x\mapsto \frac{1}{f(x_0)}\int_{x_0-y}^{x_0+y}f(t)\;dt$ est de classe $C^1$ sur $\Rr$ et il en est de même de $f$. Mais alors la fonction $x\mapsto \frac{1}{f(x_0)}\int_{x_0-y}^{x_0+y}f(t)\;dt$ est de classe $C^2$ sur $\Rr$ et il en est de même de $f$.

\begin{center}
$f$ est de classe $C^2$ sur $\Rr$.
\end{center}

\textbullet~En dérivant à $y$ fixé ou $x$ fixé l'égalité $f(x)f(y) =\int_{x-y}^{x+y}f(t)\;dt$, on obtient pour tout $(x,y)\in\Rr^2$, $f'(x)f(y)=f(x+y)-f(x-y)$ et $f(x)f'(y)=f(x+y)+f(x-y)$. En dérivant la première égalité à $y$ fixé et la deuxième à $x$ fixé, on obtient pour tout $(x,y)\in\Rr^2$, $f''(x)f(y)=f'(x+y)-f'(x-y)=f(x)f''(y)$. En particulier, pour tout réel $x$, $f''(x)f(x_0)-f(x)f''(x_0)=0$ ou encore $f''(x)-kf(x)=0$ où $f= \frac{f''(x_0)}{f(x_0}$.

\begin{center}
$f$ est solution d'une équation différentielle du type $y''-ky=0$.
\end{center}

\textbullet~$f$ est donc de l'un des types suivants : $x\mapsto\lambda\cos(\omega x)+\mu\sin(\omega x)$, $(\lambda,\mu)\in\Rr^2$ et $\omega>0$ ou $x\mapsto ax+b$, $(a,b)\in\Rr^2$ ou $x\mapsto \lambda\ch(\omega x)+\mu\sh(\omega x)$, $(\lambda,\mu)\in\Rr^2$ et $\omega>0$ (suivant que $k>0$, $k=0$ ou $k<0$). De plus, $f$ étant impaire, $f$ est nécessairement de l'un des types suivants :

\begin{center}
$x\mapsto\lambda\sin(\omega x)$, $\lambda\in\Rr^*$ et $\omega>0$ ou $x\mapsto ax$, $a\in\Rr*$ ou $x\mapsto \lambda\sh(\omega x)$, $\lambda\in\Rr$ et $\omega>0$.
\end{center}

Réciproquement,

- si $\forall x\in\Rr$, $f(x)=ax$, $a\in\Rr^*$ alors $f(x)f(y)=a^2xy$ et $\int_{x-y}^{x+y}f(t)\;dt= \frac{a}{2}((x+y)^2-(x-y)^2)=2axy$. Donc $f$ est solution si

et seulement si $a=2$. On obtient la fonction solution $x\mapsto2x$.

- si $\forall x\in\Rr$, $f(x)=\lambda\sin(\omega x)$, $\lambda\in\Rr^*$ et $\omega>0$, alors $f(x)f(y)=\lambda^2\sin(\omega x)\sin(\omega y)$ et 

$\int_{x-y}^{x+y}f(t)\;dt= \frac{\lambda}{\omega}(\cos(\omega(x-y))-\cos(\omega(x+y)))= \frac{2\lambda}{\omega}\sin(\omega x)\sin(\omega y)$. Donc $f$ est solution si et seulement si $\lambda= \frac{2}{\omega}$.

On obtient les fonctions solutions $x\mapsto \frac{2}{\omega}\sin(\omega x)$, $\omega>0$.

- si $\forall x\in\Rr$, $f(x)=\lambda\sh(\omega x)$, $\lambda\in\Rr^*$ et $\omega>0$, alors $f(x)f(y)=\lambda^2\sh(\omega x)\sh(\omega y)$ et 

$\int_{x-y}^{x+y}f(t)\;dt= \frac{\lambda}{\omega}(\ch(\omega(x+y))-\ch(\omega(x-y)))= \frac{2\lambda}{\omega}\sh(\omega x)\sh(\omega y)$. Donc $f$ est solution si et seulement si $\lambda= \frac{2}{\omega}$.

On obtient les fonctions solutions $x\mapsto \frac{2}{\omega}\sh(\omega x)$, $\omega>0$.
\fincorrection
\correction{005886}
\textbf{Existence de $F(x)=\int_{0}^{+\infty} \frac{e^{-tx}}{1+t^2}\;dt$.} Soit $x>0$. La fonction $t\mapsto \frac{e^{-tx}}{1+t^2}\;dt$ est continue sur $[0,+\infty[$ et est dominée par $ \frac{1}{t^2}$ quand $t$ tend vers $+\infty$. Donc la fonction $t\mapsto \frac{e^{-tx}}{1+t^2}\;dt$ est intégrable sur $[0,+\infty[$. On en déduit que

\begin{center}
$F$ est définie sur $]0,+\infty[$.
\end{center}

\textbf{Dérivées de $F$.} Soient $a>0$ puis $\begin{array}[t]{cccc}
\Phi~:&[a,+\infty[\times[0,+\infty[&\rightarrow&\Rr\\
 &(x,t)&\mapsto& \frac{e^{-tx}}{1+t^2}
 \end{array}$.
 
 

\textbullet~Pour tout réel $x\in[a,+\infty[$, la fonction $t\mapsto\Phi(x,t)$ est continue et intégrable sur $[0,+\infty[$.

\textbullet~La fonction $\Phi$ admet sur $[a,+\infty[\times[0,+\infty[$ des dérivées partielles d'ordre $1$ et $2$ par rapport à sa première variable $x$ et pour tout $(x,t)\in[a,+\infty[\times[0,+\infty[$

\begin{center}
$ \frac{\partial \Phi}{\partial x}(x,t)=- \frac{te^{-tx}}{1+t^2}$ et $ \frac{\partial^2\Phi}{\partial x^2}(x,t)= \frac{t^2e^{-tx}}{1+t^2}$.
\end{center}

De plus

- pour tout $x\in[a,+\infty[$, les fonctions $t\mapsto \frac{\partial \Phi}{\partial x}(x,t)$ et $t\mapsto \frac{\partial^2\Phi}{\partial x^2}(x,t)$ sont continues par morceaux sur $[0,+\infty[$.

- pour tout $t\in[0,+\infty[$, les fonctions $x\mapsto \frac{\partial \Phi}{\partial x}(x,t)$ et $x\mapsto \frac{\partial^2\Phi}{\partial x^2}(x,t)$ sont continues sur $[0,+\infty[$.

- pour tout $(x,t)\in[a,+\infty[\times[0,+\infty[$, $\left| \frac{\partial \Phi}{\partial x}(x,t)\right|\leqslant \frac{te^{-ta}}{1+t^2}=\varphi_1(t)$ et  $\left| \frac{\partial^2\Phi}{\partial x^2}(x,t)\right|\leqslant \frac{t^2e^{-ta}}{1+t^2}=\varphi_2(t)$ où les fonctions 

$\varphi_1$ et $\varphi_2$ sont continues par morceaux et intégrables sur $[0,+\infty[$ car sont dominées en $+\infty$ par $ \frac{1}{t^2}$.
 
 
 

 
 
D'après le théorème de dérivation des intégrales à paramètres (théorème de \textsc{Leibniz}), $F$ est deux fois dérivable sur $[a,+\infty[$ et les dérivées de $F$ s'obtiennent par dérivation sous le signe somme. Ceci étant vrai pour tout réel $a>0$, $F$ est deux fois dérivable sur $]0,+\infty[$ et pour $x>0$, $F''(x)=\int_{0}^{+\infty} \frac{\partial^2F}{\partial x^2} \frac{e^{-tx}}{1+t^2}\;dt=\int_{0}^{+\infty} \frac{t^2e^{-tx}}{1+t^2}\;dt$.

\begin{center}
$\forall x>0$, $F''(x)=\int_{0}^{+\infty} \frac{t^2e^{-tx}}{1+t^2}\;dt$.
\end{center}

\textbf{Equation différentielle vérifiée par $F$.} Pour $x>0$, $F''(x)+F(x)=\int_{0}^{+\infty}e^{-tx}\;dt=\left[- \frac{1}{x}e^{-tx}\right]_0^{+\infty}= \frac{1}{x}$.

\begin{center}
$\forall x>0$, $F''(x)+F(x)= \frac{1}{x}$.
\end{center}

\textbf{Existence de $G(x)=\int_{0}^{+\infty} \frac{\sin t}{t+x}\;dt$.} Soit $a>0$. Montrons l'existence de $\int_{a}^{+\infty} \frac{\sin u}{u}\;du$ et $\int_{a}^{+\infty} \frac{\cos u}{u}\;du$.

Soit $A>a$. Une intégration par parties fournit $\int_{a}^{A} \frac{\sin u}{u}\;du=- \frac{\cos a}{a}+ \frac{\cos A}{A}-\int_{a}^{A} \frac{\cos u}{u^2}\;du$. Puisque $\left| \frac{\cos A}{A}\right|\leqslant \frac{1}{A}$, on a $\lim_{A \rightarrow +\infty} \frac{\cos A}{A}=0$. D'autre part, puisque $\forall u\geqslant a$, $\left| \frac{\cos u}{u^2}\right|\leqslant \frac{1}{u^2}$, la fonction $u\mapsto \frac{\cos u}{u^2}$ est intégrable sur $[a,+\infty[$ et en particulier, $\int_{a}^{A} \frac{\cos u}{u^2}\;du$ a une limite quand $A$ tend vers $+\infty$. On en déduit que $\int_{a}^{A} \frac{\sin u}{u}\;du$ a une limite quand $A$ tend vers $+\infty$ ou encore que l'intégrale $\int_{a}^{+\infty} \frac{\sin u}{u}\;du$ converge en $+\infty$. De même, $\int_{a}^{+\infty} \frac{\cos u}{u}\;du$ converge en $+\infty$. 

Mais alors, pour $x>0$,

\begin{center}
$\cos x\int_{x}^{+\infty} \frac{\sin u}{u}\;du-\sin x\int_{x}^{+\infty} \frac{\cos u}{u}\;du=\int_{x}^{+\infty} \frac{\sin(u-x)}{u}\;du=\int_{0}^{+\infty} \frac{\sin t}{t+x}\;dt=G(x)$ existe.
\end{center}

\begin{center}
$G$ est définie sur $]0,+\infty[$.
\end{center}

\textbf{Equation différentielle vérifiée par $G$.}  Puisque la fonction $u\mapsto \frac{\sin u}{u}$ est continue sur $]0,+\infty[$, la fonction $x\mapsto\int_{x}^{+\infty} \frac{\sin u}{u}\;du=\int_{1}^{+\infty} \frac{\sin u}{u}\;du-\int_{1}^{x} \frac{\sin u}{u}\;du$ est de classe $C^1$ sur $]0,+\infty[$. De même, la fonction $x\mapsto\int_{x}^{+\infty} \frac{\cos u}{u}\;du$ est de classe $C^1$ sur $]0,+\infty[$ puis $G$ est de classe $C^1$ sur $]0,+\infty[$. De plus, pour tout réel $x>0$,

\begin{center}
$G'(x)=-\sin x\int_{x}^{+\infty} \frac{\sin u}{u}\;du- \frac{\cos x\sin x}{x}-\cos x\int_{x}^{+\infty} \frac{\cos u}{u}\;du+ \frac{\cos x\sin x}{x}=-\sin x\int_{x}^{+\infty} \frac{\sin u}{u}\;du-\cos x\int_{x}^{+\infty} \frac{\cos u}{u}\;du$,
\end{center}

puis en redérivant

\begin{center}
$G''(x)=-\cos x\int_{x}^{+\infty} \frac{\sin u}{u}\;du+ \frac{\sin^2x}{x}+\sin x\int_{x}^{+\infty} \frac{\cos u}{u}\;du+ \frac{\cos^2x}{x}=-G(x)+ \frac{1}{x}$.
\end{center}

\begin{center}
$\forall x>0$, $G''(x)+G(x)= \frac{1}{x}$.
\end{center}

\textbf{Limites de $F$ et $G$ en $+\infty$.} Puisque $\int_{1}^{+\infty} \frac{\sin u}{u}\;du$ et $\int_{1}^{+\infty} \frac{\cos u}{u}\;du$ sont deux intégrales convergentes, $\lim_{x \rightarrow +\infty}\int_{x}^{+\infty} \frac{\sin u}{u}\;du=\lim_{x \rightarrow +\infty}\int_{x}^{+\infty} \frac{\cos u}{u}\;du=0$. Puisque les fonctions sinus et cosinus sont bornées sur $\Rr$. On en déduit que $\lim_{x \rightarrow +\infty}G(x)=0$.

Pour tout réel $x>0$, $|F(x)|=\int_{0}^{+\infty} \frac{e^{-tx}}{1+t^2}\;dt\leqslant\int_{0}^{+\infty}e^{-tx}\;dt= \frac{1}{x}$ et donc $\lim_{x \rightarrow +\infty}F(x)=0$.

\begin{center}
$\lim_{x \rightarrow +\infty}F(x)=\lim_{x \rightarrow +\infty}G(x)=0$.
\end{center}

\textbf{Egalité de $F$ et $G$.} D'après ce qui précède, $(F-G)''+(F-G)=0$ et donc il existe $(\lambda,\mu)\in\Rr^2$ tel que pour tout $x>0$, $F(x)-G(x)=\lambda\cos x+\mu\sin x$. Si $(\lambda,\mu)\neq(0,0)$, alors $\lambda\cos x+\mu\sin x=\sqrt{\lambda^2+\mu^2}\cos(x-x_0)$ ne tend pas vers $0$ quand $x$ tend vers $+\infty$. Puisque $\lim_{x \rightarrow +\infty}F(x)-G(x)=0$, on a nécessairement $\lambda=\mu=0$ et donc $F-G=0$. On a montré que

\begin{center}
\shadowbox{
$\forall x>0$, $\int_{0}^{+\infty} \frac{e^{-tx}}{1+t^2}\;dt=\int_{0}^{+\infty} \frac{\sin t}{t+x}\;dt$.
}
\end{center}

\textbf{Remarque.} On peut montrer que l'égalité persiste quand $x=0$ (par continuité) et on obtient $\int_{0}^{+\infty} \frac{\sin t}{t}\;dt= \frac{\pi}{2}$.

\fincorrection
\correction{004110}
\begin{enumerate}
  \item $y = -1 + \frac1{1-\lambda e^x}$ ou $y = -1$.
  \item $y \equiv 2\arctan(\lambda e^{-\cos x}) - \frac\pi2 (\mathrm{mod}\,\pi)$.
  \item $y = \pm\sqrt{1+(\sqrt x+\lambda)^2}$ ou $y = \pm1$.
  \item $y = -\ln\bigl(1-x(1-1/e)\bigr)$.
  \item $y=\left(\lambda+\frac x2\right)\left|\lambda+\frac x2\right|$ ou $y = 0$.
\end{enumerate}
 \fincorrection
\correction{004111}
\begin{enumerate}
  \item $y = \frac{1-\lambda^2x^2}{2\lambda}$, $\lambda > 0$.
  \item $y = -x\pm\sqrt{2x^2-\lambda}$ ou $y = x(-1\pm\sqrt2)$.
  \item $y = \frac {-1 \pm \sqrt{1+4\lambda^2x^2}}{2\lambda}$
               ou $y = \pm x$ ou $y=0$.
  \item $y = -x \pm \sqrt{\lambda+3x^2}$ et $y = x(-1\pm\sqrt3)$.
\end{enumerate}
 \fincorrection
\correction{004112}
\begin{enumerate}
  \item $y = \frac{\pm1}{\sqrt{\lambda x^2+2x}}$ ou $y = 0$.
  \item $y = \pm\sqrt[4]{x^2+\frac\lambda{x^2}}$.
  \item $y = ( (\sqrt x + 2)^2 + \lambda e^{\sqrt x}\, )^2$.
  \item $y = \frac 1x\left(\frac2{\lambda-x}\right)^2$ ou $y = 0$.
  \item $y = \pm\frac{\sqrt2\,x^3}{\sqrt{2\lambda-x^4}}$ ou $y=0$.
\end{enumerate}
 \fincorrection
\correction{004113}
$y=\frac1x+\frac1{x\ln|x|+\lambda x}$ ou $y = \frac1x$.
\fincorrection
\correction{004114}
$y = \lambda\pm\sqrt{\lambda^2+1-x^2}$ ou $y=0$.
 \fincorrection
\correction{004115}
\begin{enumerate}
  \item $\gamma_1(t) = (-x(t),y(t))$ et $\gamma_2(t) = (x(-t),-y(-t))$
sont aussi solutions de~$(S)$.

Par ailleurs, la théorie de Cauchy-Lipschitz s'applique, en particulier
s'il existe $t_0$ tel que $x(t_0) = 0$ alors $x(t) = 0$ pour tout $t$.
De même s'il existe $t_0$ tel que $x(t_0) = y(t_0) = 0$ alors
$x(t)=y(t)=0$ pour tout~$t$.

Pour $\lambda,\mu$ non nuls et $x$ ne s'annulant pas, $t \mapsto(\lambda x(\mu t), \lambda y(\mu t))$
est solution de~$(S)$ si et seulement si $\mu=\lambda$.

L'ensemble des trajectoires maximales
est donc stable par les symétries par rapport aux deux axes et par les
homothéties de centre~$(0,0)$. De plus toute
trajectoire maximale qui touche l'axe des $x$ est symétrique par
rapport à cet axe.
  \item $x'(t_0) = 0 \Leftrightarrow x(t_0) = 0$ ou $y(t_0) = 0$. Dans le premier cas on
a $x(t) = 0$ pour tout $t$ et $y(t)$ est arbitraire (solution de $y' = y^2$).
Dans le second cas $x(t_0)$ est arbitraire (Cauchy-Lipschitz) donc
l'ensemble des points à tangente verticale est la réunion des deux
axes privé de $(0,0)$ (où il n'y a pas de tangente).

$y'(t_0) = 0 \Leftrightarrow x(t_0) = \pm y(t_0)$, quantité arbitraire, donc
l'ensemble des points à tangente horizontale est la réunion des deux
bissectrices des axes, privée de $(0,0)$.

Solutions particulières~: $x(t)=0, y(t) = \frac{1}{\lambda-t}$.
  \item En supposant $\Phi$ de classe $\mathcal{C}^1$ on obtient l'équation
$\frac2nx\Phi\Phi' = \Phi^2-x^2$ soit $\frac2nx\psi' = \psi-x^2$ avec $\psi=\Phi^2$.
On obtient $\psi(x) = |x|^n\Bigl(\lambda + \frac{n}{(n-1)x}\Bigr)$ si $n\ne 1$
et $\psi(x) = |x|(\lambda  - \ln|x|)$ si $n=1$.

Une courbe intégrale (en fait une trajectoire) qui
ne touche aucun des deux axes vérifie l'hypothèse $y = $
fonction de $x$ car $x'$ ne peut s'annuler donc $x$ est une fonction
injective de~$t$. Une trajectoire qui touche l'axe des $y$ est incluse
dans cet axe (déjà vu) et une trajectoire qui touche l'axe des $x$ en
dehors de $(0,0)$ le traverse ($y' \ne 0$), donc est réunion de
sous-arcs localement d'un seul côté de l'axe des $x$,
de la forme $y = \Phi(x)=\pm\sqrt{\psi(x)}$.
\end{enumerate}
 \fincorrection
\correction{004116}
Méthode d'Euler : \\
{\small \vbox{\hrule
	     \vtop{\offinterlineskip
		   \halign{\vrule \ \strut$#$\hfil\ \vrule&&\quad$#$\hfil\cr
	     x &0.0   &0.1   &0.2   &0.3   &0.4   &0.5	 &0.6	&0.7   &0.8   &0.9   &1.0\phantom{00}	\ \vrule\cr\noalign{\hrule}
	     y &1.000 &0.900 &0.810 &0.729 &0.656 &0.590 &0.531 &0.478 &0.430 &0.387 &0.348		\ \vrule\cr\noalign{\hrule}
	     z &0.000 &0.100 &0.180 &0.243 &0.292 &0.328 &0.354 &0.372 &0.383 &0.387 &0.387		\ \vrule\cr\noalign{\hrule}}}}
}
	     \par
	     Solution théorique : $y = e^{-x}$, $z = xe^{-x}$.
\fincorrection
\correction{004117}
\begin{enumerate}
  \item $y=4\arctan((\sqrt2-1) e^x)$.
  \item $y=2a-(\lambda x+\mu)^{2/3}$.
  \item $y=\exp\left(-\frac {x^2}2 + \lambda x + \mu\right)$.
\end{enumerate}
 \fincorrection
\nocorrection
\correction{004119}
Si $a > 0$, $y'(0) = a^3$,
si $a=0$, $y''(0) = 0$, $y'''(0) = 0$, $y^{IV}(0) = 6$.
Donc $y$ est croissante au voisinage de $0$.

Si $y'>0$ sur $]0,\gamma[$, alors $y(\gamma) > 0$ donc $y'(\gamma)>0$ et
$y'>0$ sur $[\gamma,\gamma+\varepsilon[$ donc, par connexité, $y'>0$ sur
$]0,\beta[$.

$y'\ge y^3  \Rightarrow  1 \le \frac{y'}{y^3}  \Rightarrow  x \le \frac1{2a^2}-\frac1{2y^2} \le \frac1{2a^2}$.
\fincorrection
\correction{004120}
\begin{enumerate}
  \item Régionnement.
  \item idem.
  \item Pour $x<0$, $y'<-e^y  \Rightarrow  -y'e^{-y}>1  \Rightarrow  x > e^{-y} + C > C$.
\end{enumerate}
 \fincorrection
\nocorrection
\correction{004122}
Supposons $t>t_{0}$ tel que $x^{2}(t)-t\ge0$. On peut alors poser $t_{1}=\min \{ t>t_{0}\,|\,x^{2}(t)-t>0\}$. On a alors 
$x^{2}(t_{1})-t_{1}=0$. Si $x(t_{1})=\sqrt{t_{1}}$, on {\'e}tudie la fonction $y(t)=x(t)-\sqrt t$. On a $y'(t_{1})=-\frac{1}{2\sqrt {t_{1}}}<0$. Cela
contredit le fait que, pour tout $t\in [t_{0},t_{1}[$, $y(t)<0$. De m{\^e}me si $x(t_{1})=-\sqrt{t_{1}}$, on {\'e}tudie la fonction $z(t)=x(t)+\sqrt t$ et on aboutit
{\`a} une contradiction. Par cons{\'e}quent la courbe int{\'e}grale reste dans $D_{0}$. Si la solution maximale ({\`a} droite) est d{\'e}finie sur $[t_{0},\beta[$, avec 
$\beta \in \R$, alors pour tout $t\in [t_{0},\beta[$, $-\beta\le x'(t)\le 0$. On en d{\'e}duit que $x'$ est int{\'e}grable sur $[t_{0},\beta[$ et donc que $x(t)$ admet
une limite finie quand $t$ tend vers $\beta$. On prolonge la fonction en $\beta$ et la fonction prolong{\'e}e v{\'e}rifie $(E)$ sur $[t_{0},\beta]$ ce qui est impossible.
On en d{\'e}duit que $\beta =+\infty$. On a, pour tout $t\ge t_{0}$, $x'(t)<0$, donc $x$ est d{\'e}croissante. Si $x(t)$ a une limite $\ell \in \R$ en $+\infty$
alors $x'(t)\mathop{\sim}\limits_{t\to +\infty}-t$, ce qui est impossible. Par cons{\'e}quent $x(t)\to -\infty$ (lorsque $t\to +\infty$). En particulier, pour $t$ assez
grand, $x(t)\le 0$. En d{\'e}rivant $(E)$ on a $x''(t)=2x(t)(x^{2}(t)-t)-1$. Si, {\`a} partir d'un certain rang, pour tout $t$, $x''(t)\ge 0$ alors $x'$ est croissante et 
major{\'e}e. Elle ne peut tendre que vers 
$0$ car sinon $x'(t) \sim \ell$, puis $x(t)\sim \ell t$ et $x'(t)\sim \ell ^{2} t^{2}$. Sinon il existe $t_{1}$ tel que $x''(t_{1})<0$. S'il existe $t_{2}>t_{1}$ tel
que $x''(t_{2})=0$ (avec $t_{2}$ minimal) alors $x'''(t_{2})=2x+\frac{1}{2x^{2}}$ qui est n{\'e}gatif pour $t$ assez grand. Ceci est impossible et donc
dans ce cas $x''(t)$ reste n{\'e}gatif lorsque $t$ tend vers $+\infty$, on a alors $0<t-x^{2}(t)\le \frac{-1}{2x(t)}\cdotp$ Par cons{\'e}quent 
$x^{2}(t)-t\to 0$ lorsque $t\to +\infty$, on en d{\'e}duit que $x(t)\mathop{\sim}\limits_{t\to +\infty}-\sqrt t$.
\fincorrection
\nocorrection
\nocorrection
\nocorrection
\correction{004126}
Soient $y,z$ deux solutions distinctes. D'après Cauchy-Lipschitz,
$y'(a) \ne z'(a)$, donc par exemple $y'(a) > z'(a)$.
Soit $c > a$ maximal tel que : $\forall\ x \in {]a,c[}$, $y(x) > z(x)$.
Donc $y-z$ est strictement positive convexe sur $[a,c]$, et s'annule en $a$ et $c$,
ce qui est impossible.
\fincorrection
\nocorrection
\nocorrection
\nocorrection
\correction{004130}
\begin{enumerate}
  \item Sinon $d(0,f'(\R)) > 0$ et $f$ ne peut pas être minorée.
  \item Supposons que pour tout~$a\in\R^p$ on a~$\nabla f(a)\ne 0$.

On considère l'équation différentielle autonome~:
$x' =  \frac{\nabla f(x)}{\|\nabla f(x)\|}$. Pour $x(0)$ donné il existe
une solution maximale, et elle est définie sur~$\R$ car $x'$ est bornée.
Alors la fonction~: $t \mapsto f(x(t))$ est $\mathcal{C}^1$ minorée sur~$\R$ donc il
esiste une suite de réels $(t_n)$ telle que
$\frac{d}{d t}(f(x(t_n))) = \|\nabla f(x(t_n))\| \to 0$ lorsque $n\to\infty$.
\end{enumerate}
 \fincorrection
\correction{004131}
$f_u$ : $p \mapsto p + u\wedge p$ est linéaire injective car
$f_u(p) = 0  \Rightarrow  p\perp p$, donc bijective. L'application $u  \mapsto f_u$
de $\R^3$ dans $\mathcal{L}(\R^3)$ est $\mathcal{C}^\infty$, donc il en est de même
de l'application inverse~: $u \mapsto(f_u)^{-1}$ et l'équation différentielle
donnée équivaut à $u' = (f_u)^{-1}(-u\wedge(u_3e_3))$ qui relève de la théorie
de Cauchy-Lipschitz~: il existe une unique solution maximale
définie sur un intervalle ouvert~$I$.
D'après l'équation différentielle, $(u' \mid u) = 0$ d'où $\|u\|$ est constant.
Alors $u' = (f_u)^{-1}(-u\wedge(u_3e_3))$ est borné, donc $u$ admet une limite
finie en tout point fini par uniforme continuité, ceci prouve que $I=\R$.
\fincorrection
\correction{004132}
\begin{enumerate}
  \item $f_{n+1}(x) - f_n(x) =  \int_{t=0}^x (f_n-f_{n-1})(t-t^2)\,d t$
    donc par récurrence $f_{n+1}-f_n \ge 0$.

    De plus $f_{n+1}(x) - f_n(x) \le x\|f_n-f_{n-1}\|_\infty$
    d'où $f_{n+2}(x) - f_{n+1}(x) \le \|f_n-f_{n-1}\|_\infty \int_{t=0}^x(t-t^2)\,d t
                                  \le \frac16\|f_n-f_{n-1}\|_\infty$
    et $\|f_{n+2} - f_{n+1}\|_\infty \le \frac16\|f_n-f_{n-1}\|_\infty$
    ce qui prouve que la série télescopique $\sum(f_{n+1}-f_n)$ est normalement
    convergente.
  \item Par passage à la limite uniforme sous le signe intégral
    on a $f(x) = 1 +  \int_{t=0}^x f(t-t^2)\,d t$ d'où $f$ est $\mathcal{C}^1$
    et $f'(x) = f(x-x^2)$ ce qui entraine le caractère $\mathcal{C}^\infty$ de $f$ par
    récurrence.
    $f'(0) = f'(1) = f(0) = 1$.
  \item $f'$ est positive d'après l'équation différentielle vérifiée par~$f$
    et $f''(x) = (1-2x)f'(x-x^2)$ est du signe de $1-2x$, c'est-à-dire que $f$
    est convexe sur $[0,\frac12]$ et concave sur $[\frac12,1]$.
  \item $1+x = f_1(x) \le f(x) = f_1(x) + \sum_{k=1}^\infty(f_{k+1}(x) - f_k(x))$.
    De plus, $f'(x) = f(x-x^2) \le f(x)$ d'où $x \mapsto f(x)e^{-x}$ est décroissante
    et vaut $1$ en $0$ ce qui prouve que $f(x) \le e^x$.
\end{enumerate}
 \fincorrection
\correction{004133}
\begin{enumerate}
  \item Qu'il en existe et qu'il y en a une unique maximale, son intervalle de définition est ouvert.
  \item Soit $(]\alpha,\beta[,x)$ une solution maximale. Si $t_0\in{]\alpha,\beta[}$ est tel que $x(t_0)=a$
    alors $x'(t_0)>0$ donc $x(t)-a$ est du signe de $t-t_0$ au voisinage de~$t_0$.
    Ceci montre que $t_0$ (éventuel) est unique, et en particulier $t_0 < 0$.
    De même, il existe au plus un réel $t_1$ tel que $x(t_1) = b$ et $t_1 < 0$.
    Par ailleurs l'existence de l'un des deux réels $t_0$ ou $t_1$ exclut l'autre.
    Enfin, $a\le x(t)\le b$ pour tout $t\in{[0,\beta[}$ donc d'après le théorème
    des bouts on a $\beta=+\infty$.
  \item Soit $\varphi : {[a,b]} \to {[a,b]},y \mapsto {x(T).}$ Comme deux courbes intégrales
    maximales distinctes n'ont aucun point commun, $\varphi$ est injective et
    par disjonction de cas on montre que $\varphi$ est strictement croissante et
    satisfait à la propriété des valeurs intermédiaires. En particulier $\varphi$
    est continue et $\varphi(y)-y$ prend une valeur positive en~$a$, négative en~$b$
    donc s'annule pour un certain $y\in{[a,b]}$. Pour cet~$y$, la solution
    correspondante est $T$-périodique.
\end{enumerate}
 \fincorrection
\correction{004134}
Remarque~: la seule continuité de~$f$ implique l'existence d'une solution
maximale à condition initiale donnée (thm. de Cauchy-Arzela, HP), mais pas son unicité.

thm des bouts~: supposons $y$ solution, définie sur $[t_0,\alpha[$ avec $\alpha < \sup J$.

On a $\frac{d}{d t}(\|y\|^2) = 2(y'\mid y) = 2(f(t,y)\mid y) \le 2a\|y\|^2 + 2b$,
ce que l'on écrit $z' = 2az + 2b-c$ avec $z=\|y\|^2$ et $c$ fonction continue positive.
Donc $z(t) = \exp(2A(t)-2A(t_0))z(t_0) +  \int_{s=t_0}^t\exp(2A(t)-2A(s))(2b(s)-c(s))\,d s$
où $A$ est une primitive de $a$ sur $J$. On en déduit que $z$ est majorée
sur $[t_0,\alpha[$ car $A$ et $b$ sont continues sur $[t_0,\alpha]$ et $c\ge 0$, donc
$\|y'\| = \|f(t,y)\|$ est aussi majorée, et $ \int_{s=t_0}^\alpha y'(s)d s$
est absolument convergente. Ainsi $y$ admet une limite finie en $\alpha^{-}$, et
l'on peut prolonger $y$ au dela de $\alpha$ avec le thm de Cauchy-Arzela~; $y$ n'est
pas maximale.
\fincorrection
\correction{004135}
\begin{enumerate}
  \item $\frac{d}{d t}f(x,y) = x'\frac{\partial f}{\partial x} + y'\frac{\partial f}{\partial y}$ donc $f$ convient
si $\frac{\partial f}{\partial x} = y(x-1)$ et $\frac{\partial f}{\partial y} = x(y-1)$ (condition suffisante).
Il n'existe pas de telle fonction (thm. de Schwarz), mais on peut accepter
$f$ telle que $\frac{\partial f}{\partial x} = \lambda(x,y)y(x-1)$ et $\frac{\partial f}{\partial y} = \lambda(x,y)x(y-1)$
où $\lambda$ est une fonction bien choisie (appelée {\it facteur intégrant}).
On voit immédiatement que $\lambda(x,y)=\frac1{xy}$ convient,
d'où $f(x,y) = x+y-\ln(xy)$.
  \item D'après le théorème d'unicité de Cauchy-Lipschitz,
s'il existe $t_0$ tel que $x(t_0)=0$ alors $x(t)=0$ pour
tout~$t$, et de même pour~$y$. Ainsi, si on fixe une
condition initiale $x(0)>0$, $y(0)>0$ alors $x(t)>0$ et
${y(t)>0}$ pour tout~$t$. 
De plus, par le même raisonnement,
si $(x(0),y(0))\ne(1,1)$ alors $(x(t),y(t))\ne(1,1)$ pour
tout~$t$. Désormais on suppose ces conditions satisfaites.
Soit $k=f(x(0),y(0))=x(0)+y(0)-\ln(x(0)y(0))$.
Par étude de fonction, on voit que $k\ne 2$ et
la courbe $C_k$ d'équation $f(x,y)=k$ est une courbe fermée de classe
$\mathcal{C}^1$ entourant le point $(1,1)$. Le point $M_t=(x(t),y(t))$
se déplace sur $C_k$ avec une vitesse numérique
${d s/d t = \sqrt{x^2(1-y)^2 + y^2(x-1)^2}\ge\alpha_k>0}$
où $\alpha_k$ ne dépend que de~$k$. On en déduit qu'une
abscisse curviligne de~$M_t$ décrit~$\R$ quand $t$ décrit~$\R$.
En particulier il existe $t_0>0$ tel que $s(t_0)-s(0)=\mathrm{longueur}(C_k)$
ce qui implique $M_{t_0}=M_0$ et le mouvement est $t_0$-périodique.
\end{enumerate}
 \fincorrection
\correction{004201}
$f(x,y) = (3x-2y)e^{4(x-y)}$.
\fincorrection
\correction{004202}
$f(x,y) = \frac{a(x-y)}2 + g(x+y)$.
\fincorrection
\correction{004203}
$f(x,y) = g(xy)$.
\fincorrection
\correction{004204}
$f(x,y) = \ln|xy| + g\Bigl(\frac yx\Bigr)$.
\fincorrection
\correction{004205}
$f(x,y) = g(\theta)$.
\fincorrection
\correction{004206}
\begin{enumerate}
  \item $g(\rho,\theta) = \lambda(\rho)e^{-2\theta}$.
  \item $g$ est $2\pi$-périodique, donc $\lambda = 0$, et $f = 0$.
\end{enumerate}
 \fincorrection
\correction{004207}
$f(x,y) = g\Bigl(\frac{1+y^2}x\Bigr)$.
\fincorrection
\nocorrection
\correction{004209}
$f(x,y) = \rho^\alpha A(\theta) + \rho^{1-\alpha} B(\theta)$.
\fincorrection
\correction{004210}
\begin{enumerate}
  \item $  (a + b\alpha + c\alpha^2)\frac{\partial^2 g}{\partial u^2}
              + (2a + b(\alpha+\beta) + 2c\alpha\beta)\frac{\partial^2 g}{\partial u \partial v}
              + (a + b\beta + c\beta^2)\frac{\partial^2 g}{\partial v^2} = 0$.
  \item 
\end{enumerate}
 \fincorrection
\correction{004211}
$2u\frac{\partial^2 g}{\partial u^2} + \frac{\partial g}{\partial u} = 0  \Rightarrow 
         f(x,y) = A\Bigl(\frac xy\Bigr)\sqrt{xy} + B\Bigl(\frac xy\Bigr)$.
\fincorrection
\correction{004212}
$2tg'(t) + (1+t^2)g''(t) = t  \Rightarrow 
          g(t) = \frac t2 + \lambda\arctan t + \mu$.
\fincorrection
\correction{004213}
\begin{enumerate}
  \item $\frac{\partial^2 g}{\partial x^2} = 4\frac{\partial^2 f}{\partial x^2} + 8\frac{\partial^2 f}{\partial x \partial y} + 4\frac{\partial^2 f}{\partial y^2}$\par
             $\frac{\partial^2 g}{\partial y^2} =  \frac{\partial^2 f}{\partial x^2} - 2\frac{\partial^2 f}{\partial x \partial y} +  \frac{\partial^2 f}{\partial y^2}$\par
             $\frac{\partial^2 g}{\partial x \partial y} = 2\frac{\partial^2 f}{\partial x^2} - 2\frac{\partial^2 f}{\partial y^2}$.
  \item $f(x,y) = \frac{xy}{16} + h(x) + k(y)$.
\end{enumerate}
 \fincorrection
\correction{004214}
$2u \frac{\partial^2 g}{\partial u \partial v} - \frac{\partial g}{\partial v} = 0  \Rightarrow f(x,y) = g\Bigl(\frac xy\Bigr)\sqrt{xy} + h(xy)$.
\fincorrection
\correction{004215}
$(1-t^2)f'' - 2tf' = 0
 \Rightarrow  f(t) = \lambda\ln\left|\frac {1+t}{1-t} \right| + \mu$.
\fincorrection
\correction{004216}
$\Delta f = 4(x^2+y^2)\Delta F$.
\fincorrection
\correction{004217}
\begin{enumerate}
  \item $\frac{\partial^2 g}{\partial u \partial v} = uv\frac{\partial^2 f}{\partial x^2} + (u+v)\frac{\partial^2 f}{\partial x \partial y} + \frac{\partial^2 f}{\partial y^2} + \frac{\partial f}{\partial x}$.
  \item $f(x,y) = \frac{xy}2 + h(u) + k(v)$ avec $u+v=y$, $uv=x$.
\end{enumerate}
 \fincorrection
\correction{004218}
$f(r) = A\cos r + B\sin r$.
\fincorrection
\correction{004219}
\begin{enumerate}
  \item Il existe $G$ telle que $\frac{\partial G}{\partial u} = g+u\frac{\partial g}{\partial u} = \frac{\partial (ug)}{\partial u}$ et
                                     $\frac{\partial G}{\partial v} = g+v\frac{\partial g}{\partial v} = \frac{\partial(vg)}{\partial v}$.\par
             Donc $G = ug + \varphi(v) = vg + \psi(u)$, d'où
             $g = \frac{\varphi(v) - \psi(u)}{v-u}$.
             La réciproque est immédiate.
  \item $f(x,y) = x\Bigl(\varphi\Bigl(y-\frac1x\Bigr)
                     + \psi\Bigl(y+\frac1x\Bigr)\Bigr)$.
\end{enumerate}
 \fincorrection
\correction{005898}
\begin{enumerate}
 \item  Soit $f$ une application de classe $C^1$ sur $\Rr^2$. Posons $f(x,y)=g(u,v)$ où $u = x+y$ et $v = x+2y$. L'application $(x,y)\mapsto(x+y,x+2y)=(u,v)$ est un automorphisme de $\Rr^2$ et en particulier un $C^1$-difféormorphisme de $\Rr^2$ sur lui-même.

\begin{center}
$ \frac{\partial f}{\partial x}= \frac{\partial}{\partial x}(g(u,v))= \frac{\partial u}{\partial x}
\times \frac{\partial g}{\partial u}+ \frac{\partial v}{\partial x}
\times \frac{\partial g}{\partial v}= \frac{\partial g}{\partial u}+ \frac{\partial g}{\partial v}$
\end{center}

De même, $ \frac{\partial f}{\partial y}= \frac{\partial g}{\partial u}+2 \frac{\partial g}{\partial v}$ et donc

\begin{center}
$2 \frac{\partial f}{\partial x}- \frac{\partial f}{\partial y}=2 \frac{\partial g}{\partial u}+2 \frac{\partial g}{\partial v}- \frac{\partial g}{\partial u}-2 \frac{\partial g}{\partial v}= \frac{\partial g}{\partial u}$.
\end{center}

Par suite, $2 \frac{\partial f}{\partial x}- \frac{\partial f}{\partial y}=0\Leftrightarrow \frac{\partial g}{\partial u}=0\Leftrightarrow\exists h\in C^1(\Rr,\Rr)/\;\forall(u,v)\in\Rr^2,\;g(u,v)=h(v)\Leftrightarrow\exists h\in C^1(\Rr,\Rr)/\;\forall(x,y)\in\Rr^2,\;f(x,y)=h(x+2y)$.

\begin{center}
\shadowbox{
Les solutions sont les $(x,y)\mapsto h(x+2y)$ où $h\in C^1(\Rr,\Rr)$.
}
\end{center}

Par exemple, la fonction $(x,y)\mapsto\cos\sqrt{(x+2y)^2+1}$ est solution.

\item  Soit $f$ une application de classe $C^1$ sur $\Rr^2\setminus\{(0,0)\}$. Posons $f(x,y)=g(r,\theta)$ où $x=r\cos\theta$ et $y=r\sin\theta$. L'application $(r,\theta)\mapsto(r\cos\theta,r\sin\theta)=(x,y)$ est un $C^1$-difféormorphisme de $]0,+\infty[\times[0,2\pi[$ sur $\Rr^2\setminus\{(0,0)\}$. De plus,

\begin{center}
$ \frac{\partial g}{\partial r}= \frac{\partial}{\partial r}(f(x,y))= \frac{\partial x}{\partial r} \frac{\partial f}{\partial x}+ \frac{\partial y}{\partial r} \frac{\partial f}{\partial y}=\cos\theta \frac{\partial f}{\partial x}+\sin\theta \frac{\partial f}{\partial y}$,
\end{center}

et

\begin{center}
$ \frac{\partial g}{\partial \theta}= \frac{\partial}{\partial \theta}(f(x,y))= \frac{\partial x}{\partial \theta} \frac{\partial f}{\partial x}+ \frac{\partial y}{\partial \theta} \frac{\partial f}{\partial y}=-r\sin\theta \frac{\partial f}{\partial x}+r\cos\theta \frac{\partial f}{\partial y}=x \frac{\partial f}{\partial y}-y \frac{\partial f}{\partial x}$.
\end{center}

Donc

\begin{align*}\ensuremath
 \frac{\partial f}{\partial y}-y \frac{\partial f}{\partial x}=0&\Leftrightarrow \frac{\partial g}{\partial \theta}=0\Leftrightarrow\exists h_1\in C^1(]0,+\infty[,\Rr)/\;\forall(r,\theta)\in]0,+\infty[\times[0,2\pi[,\;g(r,\theta)=h_1(r)\\
 &\Leftrightarrow\exists h_1\in C^1(]0,+\infty[,\Rr)/\;\forall(x,y)\in\Rr^2\setminus\{(0,0)\},\;f(x,y)=h_1\left(\sqrt{x^2+y^2}\right)\\
 &\Leftrightarrow\exists h\in C^1(]0,+\infty[,\Rr)/\;\forall(x,y)\in\Rr^2\setminus\{(0,0)\},\;f(x,y)=h(x^2+y^2).
\end{align*}

\begin{center}
\shadowbox{
Les solutions sont les $(x,y)\mapsto h(x^2+y^2)$ où $h\in C^1(]0,+\infty[,\Rr)$.
}
\end{center}

\item  Soit $f$ une fonction de classe $C^2$ sur $]0,+\infty[\times\Rr$. D'après le théorème de \textsc{Schwarz}, $ \frac{\partial^2f}{\partial x\partial y}= \frac{\partial^2f}{\partial y\partial x}$.

Soit $\begin{array}[t]{cccc}
\varphi~:&]0,+\infty[\times\Rr&\rightarrow&]0,+\infty[\times\Rr\\
 &(u,v)&\mapsto&(u,uv)=(x,y)
\end{array}
$. Donc si on pose $f(x,y)=g(u,v)$, on a $g=f\circ\varphi$.

Soit $(x,y,u,v)\in]0,+\infty[\times\Rr\times]0,+\infty[\times\Rr$.

\begin{center}
$\varphi(u,v)=(x,y)\Leftrightarrow\left\{
\begin{array}{l}
u=x\\
uv=y
\end{array}
\right.\left\{
\begin{array}{l}
u=x\\
v= \frac{y}{x}
\end{array}
\right.$.
\end{center}

Ainsi, $\varphi$ est une bijection de $]0,+\infty[$ sur lui-même et sa réciproque est l'application 

\begin{center}
$\begin{array}[t]{cccc}
\varphi^{-1}~:&]0,+\infty[\times\Rr&\rightarrow&]0,+\infty[\times\Rr\\
 &(x,y)&\mapsto&\left(x, \frac{y}{x}\right)=(u,v)
\end{array}
$.
\end{center}

De plus, $\varphi$ est de classe $C^2$ sur $]0,+\infty[\times\Rr$ et son jacobien

\begin{center}
$J_\varphi(u,v)=\left|
\begin{array}{cc}
1&0\\
v&u
\end{array}
\right|=u$
\end{center}

ne s'annule pas sur $]0,+\infty[\times\Rr$. On sait alors que $\varphi$ est un $C^2$-difféomorphisme de $]0,+\infty[\times\Rr$ sur lui-même.

Puisque $g=f\circ\varphi$ et que $\varphi$ est un $C^2$-difféomorphisme de $]0,+\infty[\times\Rr$ sur lui-même, $f$ est de classe $C^2$ sur $]0,+\infty[\times\Rr$ si et seulement si $g$ est de classe $C^2$ sur $]0,+\infty[\times\Rr$.

\textbullet~$ \frac{\partial f}{\partial x}= \frac{\partial u}{\partial x} \frac{\partial g}{\partial u}+ \frac{\partial v}{\partial x} \frac{\partial g}{\partial v}= \frac{\partial g}{\partial u}- \frac{y}{x^2} \frac{\partial g}{\partial v}$.

\textbullet~$ \frac{\partial f}{\partial y}= \frac{\partial u}{\partial y} \frac{\partial g}{\partial u}+ \frac{\partial v}{\partial y} \frac{\partial g}{\partial v}= \frac{1}{x} \frac{\partial g}{\partial v}$.

\textbullet~$ \frac{\partial^2f}{\partial x^2}= \frac{\partial}{\partial x}\left( \frac{\partial g}{\partial u}- \frac{y}{x^2} \frac{\partial g}{\partial v}\right)=\left( \frac{\partial^2g}{\partial u^2}- \frac{y}{x^2} \frac{\partial^2g}{\partial u\partial v}\right)+\left( \frac{2y}{x^3} \frac{\partial g}{\partial v}- \frac{y}{x^2} \frac{\partial^2g}{\partial u\partial v}+ \frac{y^2}{x^4} \frac{\partial^2g}{\partial v^2}\right)= \frac{\partial^2g}{\partial u^2}- \frac{2y}{x^2} \frac{\partial^2g}{\partial u\partial v}+ \frac{y^2}{x^4} \frac{\partial^2g}{\partial v^2}+ \frac{2y}{x^3} \frac{\partial g}{\partial v}$.

\textbullet~$ \frac{\partial^2f}{\partial x\partial y}= \frac{\partial}{\partial y}\left( \frac{\partial g}{\partial u}- \frac{y}{x^2} \frac{\partial g}{\partial v}\right)= \frac{1}{x} \frac{\partial^2g}{\partial v\partial  v}- \frac{1}{x^2} \frac{\partial g}{\partial v}- \frac{y}{x^3} \frac{\partial^2g}{\partial v^2}$.

\textbullet~$ \frac{\partial^2f}{\partial y^2}= \frac{\partial}{\partial y}\left( \frac{1}{x} \frac{\partial g}{\partial v}\right)= \frac{1}{x^2} \frac{\partial^2g}{\partial v^2}$.

Ensuite,

\begin{align*}\ensuremath
x^2 \frac{\partial^2f}{\partial x^2}+2xy \frac{\partial^2f}{\partial x\partial y}+y^2 \frac{\partial^2f}{\partial y^2}&=x^2 \frac{\partial^2g}{\partial u^2}-2y \frac{\partial^2g}{\partial u\partial v}+ \frac{y^2}{x^2} \frac{\partial^2g}{\partial v^2}+ \frac{2y}{x} \frac{\partial g}{\partial v}+2y \frac{\partial^2g}{\partial v\partial  v}- \frac{2y}{x} \frac{\partial g}{\partial v}- \frac{2y^2}{x^2} \frac{\partial^2g}{\partial v^2}+ \frac{y^2}{x^2} \frac{\partial^2g}{\partial v^2}\\
 &=x^2 \frac{\partial^2g}{\partial u^2}.
\end{align*}

Ainsi,

\begin{align*}\ensuremath
\forall(x,y)\in]0,+\infty[\times\Rr,&\;x^2 \frac{\partial^2f}{\partial x^2}(x,y)+2xy \frac{\partial^2f}{\partial x\partial y}(x,y)+y^2 \frac{\partial^2f}{\partial y^2}(x,y)=0\Leftrightarrow\forall(u,v)\in]0,+\infty[\times\Rr,\; \frac{\partial^2g}{\partial u^2}(u,v)=0\\
 &\Leftrightarrow\exists h\in C^2(\Rr,\Rr)/\;\forall(u,v)\in]0,+\infty[\times\Rr,\; \frac{\partial g}{\partial u}(u,v)=h(v)\\
  &\exists(h,k)\in(C^2(\Rr,\Rr))^2/\;\forall(u,v)\in]0,+\infty[\times\Rr,\;g(u,v)=uh(v)+k(v)\\
  &\exists(h,k)\in(C^2(\Rr,\Rr))^2/\;\forall(x,y)\in]0,+\infty[\times\Rr,\;f(x,y)=xh(xy)+k(xy).
\end{align*}

Les fonctions solutions sont les $(x,y)\mapsto xh(xy)+k(xy)$ où $h$ et $k$ sont deux fonctions de classe $C^2$ sur $\Rr$.
\end{enumerate}
\fincorrection
\nocorrection
\nocorrection
\nocorrection
\nocorrection
\nocorrection
\nocorrection
\nocorrection
\nocorrection
\nocorrection
\nocorrection
\nocorrection
\nocorrection
\nocorrection
\nocorrection
\nocorrection
\nocorrection
\nocorrection
\nocorrection
\nocorrection
\nocorrection
\nocorrection
\nocorrection
\nocorrection
\nocorrection
\nocorrection
\nocorrection
\nocorrection
\correction{002616}
\begin{enumerate} 

\item Le graphe est bien un parabolo\" \i de de r\'evolution
ayant l'origine pour sommet, d'axe de r\'evolution l'axe des $z$,
et dont la concavit\'e tourne vers le haut.
Les lignes de niveau sont les cercles
$x ^2+y^2=z_0,\ z_0 = c$, $c>0$ \'etant une constante;
pour $c=0$ c'est le somment, c.a.d. l'origine.

\item 
Le graphe de la fonction $f$ est  un parabolo\" \i de de r\'evolution
ayant le point $(0,0,25)$  pour sommet et plafonn\'e par
le plan des $x$ et $y$,
d'axe de r\'evolution l'axe des $z$, et
dont la concavit\'e tourne vers le bas. 
Les lignes des niveau sont les cercles
$x^2+y^2=25-z_0,\ z_0 = c$, $c<25$ \'etant une constante
qui d\'eg\'en\`erent en un point, le sommet,
pour $c=25$.

Le graphe de la fonction $g$ est un demi-c\^one
de r\'evolution
ayant le point $(0,0,5)$  pour sommet et plafonn\'e par
le plan des $x$ et $y$,
d'axe de r\'evolution l'axe des $z$, et
dont la concavit\'e tourne vers le bas. 
Les lignes des niveau sont les cercles
$x^2+y^2=(5-z_0)^2$, $z_0=c$ \'etant une constante
telle que $0 \leq c \leq 5$
qui d\'eg\'en\`erent en un point, le sommet,
pour $c=5$.

\item Dans $\R^3$ avec coordonn\'ees $(x,y,z)$,
avec $f(x)=(y,z)$,
le graphe en discussion est une h\'elice sur le c\^one de r\'evolution
$y2+z^2=x^2$.

\item Le support de cette courbe param\'etr\'ee
est une spirale planaire qui rencontre l'origine 
et dont la pente \`a l'origine vaut z\'ero.

\item Pour que $f(x,y,z)= \exp (x+y^2-z^2)$ soit constant 
il faut et il suffit que $x+y^2-z^2$ soit constant. Les surfaces de niveau
en discussion sont donc les surfaces
$x+y^2-z^2=c$. Ce sont des parabolo\"\i des hyperboliques.

\item Le graphe de l'application $f$ 
en discussion est une surface dans $\R^4$, et la dimension
4 est trop grande pour
repr\'esenter, sur une feuille de papier, ce graphe plong\'e dans
$\R^4$.
L'application $f$ est un champs de vecteurs dans le plan cependant.
De fa\c con g\'en\'erale, on peut repr\'esenter graphiquement
le champ de vecteurs $X\colon U \to \R^2$ dans l'ouvert $U$ du plan en dessinant,
au point $(u_1,u_2)$ de $U$, le vecteur $X(u_1,u_2)= (x_1(u_1,u_2),x_2(u_1,u_2))$.

N.B. Quand on repr\'esente une surface dans l'espace de dimension 3 ordinaire
par un dessin sur une feuille de papier, en v\'erit\'e on ne dessine qu'une projection de
l'espace de dimension 3 sur un plan.

\end{enumerate}

\fincorrection
\correction{002617}
\begin{enumerate}
\item La partie $A_1$ est ouverte.
Car la courbe $x^2y^2 = 1$ a quatre branches,
les deux branches de $xy = 1$ et les deux branches de
$xy = -1$; ces quatre branches coupent le plan en cinq parties dont une
contient l'origine.
La courbe $x^2y^2 = 1$ \'etant une partie ferm\'ee,
le compl\'ementaire est un ouvert qui est r\'eunion de cinq ouverts.
La partie $A_1$ est la r\'eunion des quatre parties
qui ne contiennent pas l'origine.
Puisque $A_1$ est ouvert, $A_1$ co\"\i ncide avec son int\'erieur.
L'adh\'erence de $A_1$ est la r\'eunion de $A_1$ avec les quatre branches
de la courbe $x^2y^2 = 1$.

\item La partie $A_2$ est le demi-cercle de rayon 1 ayant l'origine pour
centre constitu\'e des angles $0<\varphi<\pi$ en radians
et ce n'est ni ouvert ni ferm\'e.
La partie $A_2$ n'est pas ouverte car aucun disque de rayon positif n'est 
dans $A_2$; elle n'est pas ferm\'ee car les points $(\pm 1,0)$ sont
des points d'adh\'erence qui n'appartiennent pas \`a $A_2$.
L'adh\'erence de $A_2$ est la partie
\[
\{ (x,y)\in\R^2 | x^2+y^2=1, y\geq 0\}
\] 
du plan.
\end{enumerate}

\fincorrection
\correction{002618} 
\begin{enumerate}
\item
Soient $q_1$ un point de $B_1$ et $q_2$ un point de $B_2$,  soient
$d_1$ resp. $d_2$ la distance
de $q_1$ au bord de $B_1$ resp. la distance
de $q_2$ au bord de $B_2$, et soit $0<d\leq \min(d_1,d_2)$.
Alors la boule ouverte dans $\R^{n+m}$  centr\'ee en $(q_1,q_2)$
et de rayon $d$ est dans 
$B_1\times B_2$.

\item
Soient $p$ un point de $A$,  $q$ un point de $B$, soit
$B_1$ un disque ouvert dans $A$ contenant $p$, et
soit $B_2$ un intervalle ouvert dans $B$ contenant $q$.
D'apr\`es (1.), $B_1\times B_2$ 
est un ouvert de $\R^3$ tel que 
$B_1\times B_2 \subseteq A \times B$
et $(p,q)$ appartient \`a $B_1\times B_2$.
Par cons\'equent, $A\times B$
est un ouvert de $\R^3$. 
\end{enumerate}
\fincorrection
\correction{002619}
\begin{enumerate} 
\item La r\'eunion $\cup_n A_n$ d'une suite 
de parties ouvertes $A_n$ de $\R^2$
est bien une partie ouverte
de $\R^2$. En effet, soit $q$ un point de $\cup_n A_n$. Alors il existe
$n_0$ tel que $q$ appartienne \`a $A_{n_0}$. Puisque 
$A_{n_0}$ est ouvert, il existe un disque ouvert $D$
dans $A_{n_0}$ tel que $q$ appartienne \`a $D$. 
Par cons\'equent, il existe un disque ouvert $D$
dans $\cup_n A_n$ tel que $q$ appartienne \`a $D$. 

L'intersection $\cap_n A_n$  d'une suite 
de parties ouvertes $A_n$ de $\R^2$ n'est pas n\'ecessairement
ouverte. Par exemple, dans $\R$, l'intersection des intervalles ouverts
$]-1/n,1/n[$ est la partie $\{0\}\subseteq \R$
qui n'est pas ouverte.

\item  La r\'eunion $\cup_n B_n$ d'une suite 
de parties ferm\'ees $B_n$ de $\R^2$
n'est pas n\'ecessairement
une partie  ferm\'ee de $\R^2$.
Car le compl\'ementaire $\mathcal C(\cup_n B_n)$ de $\cup_n B_n$
est l'intersection $\cap_n \mathcal CB_n$ des compl\'ementaires
et c'est donc l'intersection d'une suite $\left (\mathcal C B_n\right)$
de parties ouvertes  de $\R^2$ qui, d'apr\`es (1.), n'est pas 
 n\'ecessairement
une partie ouverte de $\R^2$.
De m\^eme,
l'intersection $\cap_n B_n$ d'une suite 
de parties ferm\'ees $B_n$ de $\R^2$
est bien une partie ferm\'ee
de $\R^2$. 
Car le compl\'ementaire $\mathcal C(\cap_n B_n)$ de $\cap_n B_n$
est la r\'eunion $\cup_n \mathcal CB_n$ des compl\'ementaires
et c'est donc la r\'eunion d'une suite  $\left(\mathcal C B_n\right)$
de parties ouvertes de $\R^2$ qui, d'apr\`es (1.),
est une partie ouverte de $\R^2$.
\end{enumerate}
\fincorrection
\correction{002620}
La partie $A$ du plan n'est pas ouverte puisqu'elle
ne contient aucun disque ouvert.
Cette partie $A$ n'est pas ferm\'ee non plus: L'origine est un point 
d'adh\'erence: Quel que soit le disque ouvert $B$ centr\'e \`a l'origine,
il existe un point de $B$ qui appartient \`a $A$.
Mais l'origine n'appartient pas \`a $A$ d'o\`u $A$ n'est pas ferm\'e.
L'adh\'erence $\overline A$ de $A$ est la r\'eunion
$A \cup (\{0\}\times [-1,1])$. Car quelle que soit
la suite $(x_n)$ de points de $A \cup (\{0\}\times [-1,1])$
telle que cette suite converge dans le plan,
la limite appartient \`a $A \cup (\{0\}\times [-1,1])$.
\fincorrection
\nocorrection
\correction{005844}
\begin{enumerate}
 \item  Soit $A$ une partie de $E$. $\overline{A}$ est fermé et donc  $\overline{(\overline{A})}=\overline{A}$. $\overset{\circ}{A}$ est ouvert et donc $\overset{\circ}{\overset{\circ}{A}}=\overset{\circ}{A}$.

\item  Soient $A$ et $B$ deux parties de $E$ telles que  $A\subset B$.

\textbullet~Pour tout $x\in E$, $x\in\overline{A}\Rightarrow \forall V\in\mathcal{V}(x),\; V\cap A\neq\varnothing\Rightarrow \forall V\in\mathcal{V}(x),\;V\cap B\neq\varnothing\Rightarrow x\in\overline{B}$. Donc $\overline{A}\subset\overline{B}$.

 
\textbullet~Pour tout $x\in E$, $x\in\overset{\circ}{A}\Rightarrow A\in\mathcal{V}(x)\Rightarrow B\in\mathcal{V}(x)\Rightarrow x\in\overset{\circ}{B}$. Donc $\overset{\circ}{A}\subset\overset{\circ}{B}$.

\item  Soient $A$ et $B$ deux parties de $E$.

$\overline{A}\cup\overline{B}$ est une partie fermée de $E$ contenant $A\cup B$. Donc  $\overline{A\cup B}\subset\overline{A}\cup\overline{B}$ (puisque $\overline{A\cup B}$ est le plus petit fermé de $E$ au sens de l'inclusion contenant $A\cup B$).

Réciproquement, $A\subset A\cup B$ et $B\subset A\cup B\Rightarrow\overline{A}\subset\overline{A\cup B}$ et  $\overline{B}\subset\overline{A\cup B}\Rightarrow\overline{A}\cup\overline{B}\subset\overline{A\cup B}$.

Finalement $\overline{A\cup B}=\overline{A}\cup\overline{B}$.

$\overset{\circ}{A}\cap\overset{\circ}{B}$ est un ouvert contenu dans $A\cap B$ et donc $\overset{\circ}{A}\cap\overset{\circ}{B}\subset\overset{\circ}{A\cap B}$.

Réciproquement ,  $A\cap B\subset A$ et $A\cap B\subset B\Rightarrow\overset{\circ}{A\cap B}\subset\overset{\circ}{A}$ et $\overset{\circ}{A\cap B}\subset\overset{\circ}{B}\Rightarrow\overset{\circ}{A\cap B}\subset\overset{\circ}{A}\cap\overset{\circ}{B}$.

Finalement, $\overset{\circ}{A\cap B}=\overset{\circ}{A}\cap\overset{\circ}{B}$.

\item  $\overline{A}\cap\overline{B}$ est un fermé contenant $A\cap B$ et donc  $\overline{A\cap B}\subset\overline{A}\cap\overline{B}$.

On n'a pas nécessairement l'égalité. Si $A = [0,1[$ et $B = ]1,2]$, $A\cap B =\varnothing$ puis $\overline{A\cap B}=\varnothing$ mais $\overline{A}\cap\overline{B}= [0,1]\cap[1,2]=\{1\}\neq\varnothing$.

$\overset{\circ}{A}\cup\overset{\circ}{B}$ est un ouvert contenu dans $A\cup B$ et donc $\overset{\circ}{A}\cup\overset{\circ}{B}\subset\overset{\circ}{A\cup B}$.

On n'a pas nécessairement l'égalité. Si $A = [0,1]$ et $B =[1,2]$, $A\cup B =[0,2]$ puis $\overset{\circ}{A\cup B}=]0,2[$ mais $\overset{\circ}{A}\cup\overset{\circ}{B}=]0,1[\cup]1,2[\neq]0,2[$.

\item  Soient $A$ et $B$ deux parties de $E$. Soit $x\in E$.

\begin{align*}\ensuremath
x\in\overset{\circ}{A\setminus B}&\Leftrightarrow A\setminus B\in\mathcal{V}(x)\Leftrightarrow\exists\mathcal{B}\;\text{boule ouverte de centre}\;x\;\text{telle que}\;\mathcal{B}\subset A\setminus B\\
 &\Leftrightarrow \exists\mathcal{B}\;\text{boule ouverte de centre}\;x\;\text{telle que}\;\mathcal{B}\subset A\;\text{et}\;\mathcal{B}\subset{^c}B\Leftrightarrow A\in\mathcal{V}(x)\;\text{et}\;{^c}B\in\mathcal{V}(x)\\
  &\Leftrightarrow x\in\overset{\circ}{A}\;\text{et}\;x\in\overset{\circ}{({^c}B)}\Leftrightarrow x\in\overset{\circ}{A}\;\text{et}\;x\in{^c}(\overline{B})\Leftrightarrow x\in\overset{\circ}{A}\cap{^c}(\overline{B})\Leftrightarrow x\in\overset{\circ}{A}\setminus\overline{B}.
 \end{align*}
 

Donc $\overset{\circ}{A\setminus B}=\overset{\circ}{A}\setminus\overline{B}$.

\item  Soit $A$ une partie de $E$. $\overset{\circ}{\overline{\overset{\circ}{A}}}\subset\overline{\overset{\circ}{A}}\Rightarrow\overline{\overset{\circ}{\overline{\overset{\circ}{A}}}}\subset\overline{\overline{\overset{\circ}{A}}}=\overline{\overset{\circ}{A}}$. D'autre part $\overset{\circ}{A}\subset\overline{\overset{\circ}{A}}\Rightarrow\overset{\circ}{\overset{\circ}{A}}=\overset{\circ}{A}\subset\overset{\circ}{\overline{\overset{\circ}{A}}}\Rightarrow\overline{\overset{\circ}{A}}\subset\overline{\overset{\circ}{\overline{\overset{\circ}{A}}}}$. Finalement, $\overline{\overset{\circ}{A}}=\overline{\overset{\circ}{\overline{\overset{\circ}{A}}}}$.

$\overset{\circ}{\overline{A}}\subset\overline{\overset{\circ}{\overline{A}}}\Rightarrow\overset{\circ}{\overset{\circ}{\overline{A}}}=\overset{\circ}{\overline{A}}\subset\overset{\circ}{\overline{\overset{\circ}{\overline{A}}}}$. D'autre part  $\overset{\circ}{\overline{A}}\subset\overline{A}\Rightarrow\overline{\overset{\circ}{\overline{A}}}\subset\overline{\overline{A}}=\overline{A}\Rightarrow\overset{\circ}{\overline{\overset{\circ}{\overline{A}}}}\subset\overset{\circ}{\overline{A}}$. Finalement, $\overset{\circ}{\overline{\overset{\circ}{\overline{A}}}}=\overset{\circ}{\overline{A}}$.
\end{enumerate}
\fincorrection
\correction{005845}
 L'exercice \ref{ex:rou6} montre que l'on ne peut pas faire mieux.

Soit $A = ([0,1[\cup]1,2])\cup\{3\}\cup(\Qq\cap[4,5])$.

\textbullet~$\overset{\circ}{A}= ]0,1[\cup]1,2[$.

\textbullet~$\overline{\overset{\circ}{A}}= [0,2]$.

\textbullet~$\overset{\circ}{\overline{\overset{\circ}{A}}}= ]0,2[$.

\textbullet~$\overline{A}= [0,2] ]\cup\{3\}\cup[4,5]$

\textbullet~$\overset{\circ}{\overline{A}}= ]0,2[\cup]4,5[$.

\textbullet~$\overline{\overset{\circ}{\overline{A}}}=[0,2]\cup[4,5]$.

Les $7$ ensembles considérés sont deux à deux distincts.
\fincorrection
\correction{005846}
Soit $f\in E$. Pour $n\in\Nn^*$, soit $g_n$ l'application définie par $\forall x\in[0,1]$, $g_n(x) = f(x)+ \frac{1}{n}\left|x- \frac{1}{2}\right|$.

Chaque fonction $g_n$ est continue sur $[0,1]$ mais non dérivable en $ \frac{1}{2}$ ou encore $\forall n\in\Nn^*$, $g_n\in E\setminus D$. De plus, $\forall n\in\Nn^*$ $\|f-g_n\|_\infty= \frac{1}{2n}$. On en déduit que la suite $(g_n)_{n\geqslant1}$ tend vers $f$ dans l'espace vectoriel normé $(E,\|\;\|_\infty)$.

$f$ est donc limite d'une suite d'éléments de ${^c}D$ et donc est dans l'adhérence de ${^c}D$. Ceci montre que $\overline{{^c}D}=E$ ou encore ${^c}(\overset{\circ}{D})=E$ ou enfin $\overset{\circ}{D}=\varnothing$.

Enfin, puisque $P\subset D$, on a aussi $\overset{\circ}{P}=\varnothing$.
\fincorrection
\correction{005848}
\begin{enumerate}
 \item  Soit $x\in E$. Puisque $D$ est dense dans $E$, il existe une suite $(d_n)_{n\in\Nn}$ d'éléments de $D$ convergeant vers $x$ et puisque $f$ et $g$ sont continues et coincident sur $D$ et donc en $x$

\begin{center}
$f(x) = f\left(\lim_{n \rightarrow +\infty}d_n\right)=\lim_{n \rightarrow +\infty} f(d_n) =\lim_{n \rightarrow +\infty}g(d_n) = g\left(\lim_{n \rightarrow +\infty}d_n\right) = g(x)$.
\end{center}

On a montré que $f = g$.

\item  Soit $f\in\Rr^\Rr$. On suppose que $\forall(x,y)\in\Rr^2$ $f(x+y) = f(x)+f(y)$. Soit $a = f(1)$.

\textbullet~$x=y=0$ fournit $f(0) = 0 = a\times0$.

\textbullet~Soit $n\in \Nn^*$ et $x\in\Rr$. $f(nx) = f(x+...+x) = f(x)+...+f(x)  = nf(x)$. Ceci reste vrai pour $n=0$.
 
\textbullet~En particulier $x = 1$ fournit pour tout entier naturel non nul $n$, $f(n) = nf(1) = an$ puis $x = \frac{1}{n}$   fournit 

$nf\left( \frac{1}{n}\right) = f(1) = a$ et donc $f\left( \frac{1}{n}\right) = \frac{a}{n}$.

\textbullet~Ensuite, $\forall(p,q)\in(\Nn\times\Nn^*)^2$, $f\left( \frac{p}{q}\right) = pf\left( \frac{1}{q}\right)=a \frac{p}{q}$.

\textbullet~Soit $x\in\Rr$. L'égalité $f(x)+f(-x)=f(0)=0$ fournit $f(-x)=-f(x)$.

\textbullet~En particulier, $\forall(p,q)\in(\Nn^*)^2$, $f\left(- \frac{p}{q}\right) = -f\left( \frac{p}{q}\right)= - a \frac{p}{q}$.

En résumé, si $f$ est morphisme du groupe $(\Rr,+)$ dans lui-même, $\forall r\in\Qq$, $f(r) = ar$ où $a = f(1)$.

Si de plus $f$ est continue sur $\Rr$, les deux applications $f~:~x\mapsto f(x)$ et $g~:~x\mapsto ax$ sont continues sur $\Rr$ et coïncident sur $\Qq$ qui est dense dans $\Rr$. D'après le 1), $f = g$ ou encore $\forall x\in\Rr$, $f(x)=ax$ où $a=f(1)$.

Réciproquement, toute application linéaire $x\mapsto ax$ est en particulier un morphisme du groupe $(\Rr,+)$ dans lui-même, continu sur $\Rr$. 

\begin{center}
\shadowbox{
Les morphismes continus du groupe $(\Rr,+)$ dans lui-même sont les applications linéaires $x\mapsto ax$, $a\in\Rr$.
}
\end{center}
\end{enumerate}
\fincorrection
\correction{005849}
Soit $u = (u_n)_{n\in\Nn}$ une suite bornée de l'espace normé $(E,\|\;\|)$ ayant une unique valeur d'adhérence que l'on note $\ell$. Montrons que la suite $u$ converge vers $\ell$.

Supposons par l'absurde que la suite $u$ ne converge pas vers $\ell$. Donc

\begin{center}
$\exists\varepsilon>0/\;\forall n_0\in\Nn,\;\exists n\geqslant n_0/\;\|u_n-\ell\|\geqslant\varepsilon$\quad$(*)$.
\end{center}

$\varepsilon$ est ainsi dorénavant fixé.

En appliquant $(*)$ à $n_0 = 0$, il existe un rang $\varphi(0)\geqslant n_0 = 0$ tel que $\|u_{\varphi(0)}-\ell\|\geqslant\varepsilon$.

Puis en prenant $n_0=\varphi(0)+1$, il existe un rang $\varphi(1)>\varphi(0)$ tel que $\|u_{\varphi(1)}-\ell\|\geqslant\varepsilon$ ... et on construit ainsi par récurrence une suite extraite $(u_{\varphi(n)})_{n\in\Nn}$ telle que $\forall n\in\Nn$, $\|u_{\varphi(n)}-\ell\|\geqslant\varepsilon$.

Maintenant, la suite $u$ est bornée et il en est de même de la suite $(u_{\varphi(n)})$. Puisque $E$ est de dimension finie, le théorème de \textsc{Bolzano}-\textsc{Weierstrass} permet d'affirmer qu'il existe une suite $(u_{\psi(n)})_{n\in\Nn}$ extraite de $(u_{\varphi(n)})$ et donc de $u$ convergeant vers un certain $\ell'\in E$. $\ell'$ est donc une valeur d'adhérence de la suite $u$. Mais quand $n$ tend vers $+\infty$ dans l'inégalité $\|u_{\psi(n)}-\ell\|\geqslant\varepsilon$, on obtient $\|\ell'-\ell\|\geqslant\varepsilon$ et donc $\ell\neq\ell'$. Ceci constitue une contradiction et donc $u$ converge vers $\ell$.
\fincorrection
\nocorrection
\nocorrection
\nocorrection
\nocorrection
\nocorrection
\nocorrection
\nocorrection
\correction{004828}
\begin{enumerate}
  \item 
  \item $d(x_n,a)$ d{\'e}croit, donc tend vers $d$.  Il existe une
    sous-suite $(x_{n_k})$ convergeant vers $\ell$ et $d(\ell,a) = d$.
    La suite $(f(x_{n_k}))$ converge vers $f(\ell)$ et on a
    $d(f(\ell),a) = d$, donc $\ell = a$. Il y a une seule valeur
    d'adh{\'e}rence, donc la suite converge.
\end{enumerate}
 \fincorrection
\correction{004829}
$C$ est stable par~$f_n$ qui est $\Bigl(1-\frac1n\Bigr)$-lipschitzienne.
Donc il existe $x_n\in C$ tel que $f_n(x_n) = x_n$~; toute valeur d'adh{\'e}rence de $(x_n)$
est point fixe de~$f$.
\fincorrection
\nocorrection
\nocorrection
\nocorrection
\correction{004833}
\begin{enumerate}
  \item 
  \item 
  \item Soit $\ell = \lim_{n\to\infty} \delta(K_n)$.
Il existe $x_n,y_n \in K_n$ tels que $d(x_n,y_n) = \delta(K_n)$.
Apr{\`e}s extraction de sous-suites, on peut supposer que $x_n \to x$ et
$y_n \to y$. Pour $\varepsilon > 0$, $(\overline B(x,\varepsilon) \cap K_n)$
et $(\overline B(y,\varepsilon) \cap K_n)$ forment des suites d{\'e}croissantes de
compacts non vides, donc $\overline B(x,\varepsilon) \cap K$ et $\overline B(y,\varepsilon) \cap K$
sont non vides. Par cons{\'e}quent, $\delta(K) \ge \ell - 2\varepsilon$.
\end{enumerate}
 \fincorrection
\nocorrection
\correction{004835}
$U_{i,n} = \{x \in E \text{ tq } \overline B(x,1/n) \subset O_i\}$ est ouvert
et les $U_{i,n}$ recouvrent $E$. On extrait un recouvrement fini
$ \Rightarrow  r = \min(1/n)$.
\fincorrection
\correction{004836}
Soit $(u^n)$ une suite de suites {\'e}l{\'e}ments de $A$ : $u^n = (u^n_k)$.
On peut trouver une sous-suite $(u^{n_{p_0}})$ telle que $(u^{n_{p_0}}_0)$ converge
vers $u_0 \in {[0,1]}$, puis une sous-suite $(u^{n_{p_1}})$ telle que
$(u^{n_{p_1}}_0,u^{n_{p_1}}_1)$ converge vers $(u_0,u_1) \in {[0,1]}^2$, etc.
Alors la suite $(u^{n_{p_k}})_k$ converge dans $A$ vers $(u_0,u_1,\dots)$.
\fincorrection
\nocorrection
\nocorrection
\correction{004839}
On choisit $a\in K$ et on consid{\`e}re pour $n\ge1$
la fonction $f_n$~: $x \mapsto\frac1na + \Bigl(1-\frac1n\Bigr)f(x)$.
$f_n$ est une $\Bigl(1-\frac1n\Bigr)$-contraction de~$K$ donc admet un point
fixe $x_n$. Si $x$ est une valeur d'adh{\'e}rence de la suite $(x_n)$ alors
$f(x) = x$.
\fincorrection
\nocorrection
\nocorrection
\nocorrection
\nocorrection
\nocorrection
\correction{005850}
Pour $\alpha\in]0,\pi[$, posons $f(\alpha) =\underset{n\in\Zz}{\text{Sup}}|\sin(n\alpha)|=\underset{n\in\Nn}{\text{Sup}}|\sin(n\alpha)|$.

\textbullet~Tout d'abord $\forall\alpha\in]0,\pi[$, $\forall n\in\Nn$, $|\sin(n(\pi-\alpha))| = |\sin(n\alpha)|$ et donc $\forall\alpha\in]0,\pi[$, $f(\pi-\alpha) = f(\alpha)$.

On en déduit que $\underset{\alpha\in]0,\pi[}{\text{Inf}}f(\alpha) =\underset{\alpha\in\left]0, \frac{\pi}{2}\right]}{\text{Inf}}f(\alpha)$.

\textbullet~$f\left( \frac{\pi}{3}\right)=\text{sup}\left\{0, \frac{\sqrt{3}}{2}\right\}= \frac{\sqrt{3}}{2}$.

\textbullet~Ensuite, si $\alpha\in\left[ \frac{\pi}{3}, \frac{\pi}{2}\right]$, $f(\alpha)\geqslant\sin(\alpha)\geqslant\sin\left( \frac{\pi}{3}\right)= \frac{\sqrt{3}}{2}=f\left( \frac{\pi}{3}\right)$.
Par suite $\underset{\alpha\in\left]0, \frac{\pi}{2}\right]}{\text{Inf}}f(\alpha)=\underset{\alpha\in\left]0, \frac{\pi}{3}\right]}{\text{Inf}}f(\alpha)$.

\textbullet~Soit alors $\alpha\in\left]0, \frac{\pi}{3}\right]$. Montrons qu'il existe un entier naturel non nul $n_0$ tel que $n_0\alpha\in\left[ \frac{\pi}{3}, \frac{2\pi}{3}\right]$.

Il existe un unique entier naturel $n_1$ tel que $n_1\alpha\leqslant \frac{\pi}{3}<(n_1+1)\alpha$ à savoir $n_1= E\left( \frac{\pi}{3\alpha}\right)$.

Mais alors, $ \frac{\pi}{3}<(n_1+1)\alpha= n\alpha+\alpha\leqslant \frac{\pi}{3}+ \frac{\pi}{3}= \frac{2\pi}{3}$ et l'entier $n_0=n_1+1$ convient.

Ceci montre que $f(\alpha)\geqslant\sin \frac{\pi}{3}= \frac{\sqrt{3}}{2}=f\left( \frac{\pi}{3}\right)$.

Finalement $\forall\alpha\in]0,\pi[$, $f(\alpha)\geqslant f\left( \frac{\pi}{3}\right)$ et donc 
$\underset{\alpha\in]0,\pi[}{\text{Inf}}\left\{\underset{n\in\Zz}{\text{Sup}}|\sin(n\alpha)|\right\}=\underset{\alpha\in]0,\pi[}{\text{Min}}\left\{\underset{n\in\Zz}{\text{Sup}}|\sin(n\alpha)|\right\}=f\left( \frac{\pi}{3}\right)= \frac{\sqrt{3}}{2}$.

\begin{center}
\shadowbox{
$\underset{\alpha\in]0,\pi[}{\text{Inf}}\left\{\underset{n\in\Zz}{\text{Sup}}|\sin(n\alpha)|\right\}= \frac{\sqrt{3}}{2}$.
}
\end{center}
\fincorrection
\correction{004717}
\begin{enumerate}
  \item $A \setminus ]a-\frac 1n, a+\frac 1n[$ est fini.
  \item 
\end{enumerate}
 \fincorrection
\nocorrection
\correction{004719}
Soient $x<y$ : 
         Il existe $a \in \N$ tel que $n\ge a  \Rightarrow  \sqrt{n+1} - \sqrt{n} < y-x$.\par
         Il existe $b \in \N$ tel que $\sqrt a - \sqrt b < x$.\par
         Alors il existe $c \ge a$ tel que $x < \sqrt c - \sqrt b < y$.\par
\fincorrection
\correction{004720}
$n+p\sqrt2 > 1  \Rightarrow  n>0,\ p>0$, donc $A\cap]1,+\infty[$
         admet un plus petit {\'e}l{\'e}ment : $3+2\sqrt2$.
\fincorrection
\correction{004721}
On construit un ensemble de type Cantor dont les trous ont pour
         longueur $1$, $\frac 1{2^a}$, $\frac 1{4^a}$, \dots,
         et on r{\'e}partit les $x_k$ de part et d'autre des trous en fonction
         de l'{\'e}criture d{\'e}cimale de $k$ (0 $\rightarrow$ {\`a} gauche,
         1 $\rightarrow$ {\`a} droite).
\fincorrection
\nocorrection
\correction{004723}
L'ensemble des valeurs d'adh{\'e}rence est un intervalle constitu{\'e} de
         points fixes de $f$ $ \Rightarrow $ la suite $(u_n)$ a une seule valeur
         d'adh{\'e}rence.
\fincorrection
\nocorrection
\nocorrection
\nocorrection
\correction{004727}
Soit $\ell = \liminf \frac{u_n}n$ et $\varepsilon > 0$.
         Il existe $p$ tel que $(\ell-\varepsilon)p \le u_p \le (\ell+\varepsilon)p$.
         \par
         Alors pour $n\in\N$ et $0\le k < p$ :
         $u_k + (\ell-\varepsilon)np \le u_{np+k} \le u_k + (\ell+\varepsilon)np$.
\fincorrection
\correction{004728}
\begin{enumerate}
  \item Il est suppos{\'e} connu (et {\`a} savoir d{\'e}montrer) le fait suivant~:
    {\it si $G$ est un sous-groupe de~$\R$, alors soit $G$ est monog{\`e}ne, soit $\overline G=\R$.}
    Dans le cas de la question, le groupe $G$ des p{\'e}riodes de~$f$ contient $1$ et~$\sqrt2$
    donc n'est pas monog{\`e}ne car~$\sqrt2\notin\Q$ (la d{\'e}monstration a {\'e}t{\'e} demand{\'e}e {\`a} l'{\'e}l{\`e}ve).
    De plus $G$ est ferm{\'e} par continuit{\'e} de~$f$, d'o{\`u} $f$ est constante.
  \item D'apr{\`e}s la premi{\`e}re question, pour tout~$y\in\R\setminus\Q$
    l'application $x \mapsto f(x,y)$ est constante et il en va de m{\^e}me si $y\in\Q$ par
    continuit{\'e} de~$f$. Donc $f$ est de la forme $(x,y) \mapsto g(y)$ o{\`u} $g$ est
    $1$-p{\'e}riodique. R{\'e}ciproquement, toute fonction~$f$ de cette forme convient.
\end{enumerate}
 \fincorrection
\correction{005843}
\textbf{1ère solution.} \textbullet~Montrons qu'entre deux réels distincts, il existe un rationnel.

Soient $x$ et $y$ deux réels tels que $x < y$. Soient $d= y -x$ puis $n$ un entier naturel non nul tel que  $ \frac{1}{n}<d$ (par exemple, $n=E\left( \frac{1}{d}\right)+1$). Soient enfin $k = E(nx)$ et $r = \frac{k+1}{n}$. $r$ est un rationnel et de plus

\begin{center}
$x = \frac{nx}{n}< \frac{k+1}{n}=r\leqslant \frac{nx+1}{n}=x+ \frac{1}{n}<x+d=x+y-x = y$.
\end{center}

En  résumé, $\forall(x,y)\in\Rr^2$, $(x<y\Rightarrow\exists r\in\Qq/\;x<r<y)$. Ceci montre que $\Qq$ est dense dans $\Rr$.

\textbf{2ème solution.} On sait que tout réel est limite d'une suite de décimaux et en particulier tout réel est limite d'une suite de rationnels. Donc $\Qq$ est dense dans $\Rr$.

\begin{center}
\shadowbox{
$\Qq$ est dense dans $\Rr$.
}
\end{center}
\fincorrection
\correction{005851}
Soit $f$ une application uniformément continue sur $\Rr$. 
$\exists\alpha> 0/\;\forall(x,y)\in\Rr^2,\;(|x-y|\leqslant\alpha\Rightarrow|f(x)-f(y)|\leqslant1)$.

Soit $x\in\Rr^+$ (le travail est analogue si $x\in\Rr^-$). 

Pour $n\in\Nn$

\begin{center}
$|x-n\alpha|\leqslant\alpha\Leftrightarrow-\alpha\leqslant x-n\alpha\leqslant\alpha\Leftrightarrow \frac{x}{\alpha}-1\leqslant n\leqslant \frac{x}{\alpha}+1\Leftarrow n = E\left( \frac{x}{\alpha}\right).$
\end{center}

On pose $n_0=E\left( \frac{x}{\alpha}\right)$.

\begin{align*}\ensuremath
|f(x)| &\leqslant|f(x)-f(x-\alpha)|+|f(x-\alpha)-f(x-2\alpha)|+\ldots+|f(x-(n_0-1)\alpha)-f(x-n_0\alpha)|+|f(x-n_0\alpha)-f(0)|+|f(0)|\\
 &\leqslant n_0+1+|f(0)|\;(\text{car}\;|x-n_0\alpha-0|\leqslant\alpha)\\
 &\leqslant \frac{x}{\alpha}+2+|f(0)| .
\end{align*}

Ainsi, $\forall x\in\Rr^+$, $|f(x)|\leqslant \frac{x}{\alpha}+2+|f(0)|$. Par symétrie des calculs, $\forall x\in\Rr^-$, $|f(x)|\leqslant \frac{-x}{\alpha}+2+|f(0)|$ et donc $\forall x\in\Rr$, $|f(x)|\leqslant \frac{|x|}{\alpha}+2+|f(0)|$.

\begin{center}
\shadowbox{
$f$ uniformément continue sur $\Rr\Rightarrow\exists(a,b)\in\Rr^2/\;\forall x\in\Rr,\;|f(x)\leqslant a|x|+b$.
}
\end{center}
\fincorrection
\correction{004729}
Soit $a \in \mathring{\overline U} \cap \mathring{\overline V}$ :
Il existe $B(a,r) \subset \overline U\cap \overline V$.

Soit $b \in B(a,r) \cap U$ : Il existe $B(b,r') \subset B(a,r) \cap U$.
Donc $b \notin \overline V$, contradiction.
\fincorrection
\nocorrection
\correction{004731}
$\mathring{\overline U} \subset \overline U  \Rightarrow  \overline{\mathring{\overline U}} \subset \overline U$.\par
$U \subset \overline U  \Rightarrow  U \subset \mathring{\overline U}  \Rightarrow  \overline U \subset \overline{\mathring{\overline U}}$.
\fincorrection
\correction{004732}
Soit $a$ int{\'e}rieur {\`a} $\text{Fr}(U)$ : Il existe $B(a,r) \subset \overline U\setminus U$.
$B(a,r) \cap U = \varnothing  \Rightarrow  a \notin \overline U$, contradiction.
\fincorrection
\nocorrection
\correction{004734}
Par passage {\`a} la limite, $\delta(A) = \delta(\overline A) \ge \delta(\text{Fr}(A))$.\par
Soient $x,y \in A$ distincts et $D$ la droite passant par $x$ et $y$.
$D$ coupe $A$ suivant un ensemble born{\'e} dont les extr{\'e}mit{\'e}s appartiennent
{\`a} $\text{Fr}(A)$. Donc $\delta(\text{Fr}(A)) \ge \delta(A)$.
\fincorrection
\nocorrection
\correction{004736}
Si $f$ est continue : 
soit $x \in \overline A$ : $x = \lim a_n  \Rightarrow  f(x) = \lim f(a_n) \in \overline{f(A)}$.\par
soit $x \in f^{-1}(B)^\circ$ : $f(x) \in \mathring B  \Rightarrow  \exists\ B(f(x),r) \subset B$,
$\exists\ \delta > 0 \text{ tq } f(B(x,\delta)) \subset B(f(x),r)$\par
$ \Rightarrow  B(x,\delta) \subset f^{-1}(B)$.

si $f(\overline A) \subset \overline{f(A)}$ : 
soit $B \subset F$ ferm{\'e} et $A = f^{-1}(B)$ :
$f(\overline A) \subset B$ donc $\overline A \subset A$.

si $f^{-1}(\mathring B) \subset f^{-1}(B)^\circ$ : 
soit $B \subset F$ ouvert et $A = f^{-1}(B)$ :
$\mathring A \supset f^{-1}(\mathring B) = A$.
\fincorrection
\nocorrection
\nocorrection
\nocorrection
\nocorrection
\correction{004741}
\begin{enumerate}
  \item Pour $x,x'\in\R^n$ et $y\in A$ on a $d_A(x)\le \|x-y\| \le \|x-x'\| + \|x'-y\|$.
    En prenant la borne inf{\'e}rieure sur~$y$ on obtient $d_A(x) \le \|x-x'\| + d_A(x')$.
    Par sym{\'e}trie on a aussi $d_A(x') \le \|x-x'\| + d_A(x)$ d'o{\`u}
    $|d_A(x)-d_A(x')\le\|x-x'\|$.
  \item On sait que $\overline A = \{x\in\R^n\text{ tq }d_A(x)=0\}$.
    Donc $d_A = d_B \Rightarrow \overline A = \overline B$ et la r{\'e}ciproque r{\'e}sulte
    de la propri{\'e}t{\'e} facile $d_A = d_{\overline A}$.
  \item On note~: $M = \sup\{|d_A(y)-d_B(y)|,\ y\in\R^n\}$,
    $\alpha = \sup\{d_B(x), x\in A\}$ et $\beta=\sup\{d_A(x),\ x\in B\}$.
    Par restriction de $y$ {\`a} $A\cup B$ on obtient $M \ge \max(\alpha,\beta)$.
    Par ailleurs, pour $y\in\R^n$, $a\in A$ et $b\in B$ on a
    $\|y-a\|-\|y-b\|\le\|a-b\|$ d'o{\`u} $d_A(y)-\|y-b\|\le d_A(b)$ puis
    $d_A(y)-d_B(y)\le \beta$. Par sym{\'e}trie on a aussi
    $d_B(y)-d_A(y)\le \alpha$ donc $|d_A(y)-d_B(y)|\le\max(\alpha,\beta)$
    et finalement $M\le\max(\alpha,\beta)$.
\end{enumerate}
 \fincorrection
\nocorrection
\nocorrection
\correction{004744}
\begin{enumerate}
  \item Pour~$\theta\in\R$, la demi-droite d'origine~$O$ et d'angle
    polaire~$\theta$ coupe~$K$ selon un intervalle non trivial ($K$ est convexe
    et $O$ est interieur {\`a}~$K$), ferm{\'e} born{\'e} ($K$ est compact).
    On note~$f(\theta)$ la longueur de cet intervalle, ce qui d{\'e}finit
    $f : \R \to {\R^{+*}}$ $2\pi$-p{\'e}riodique telle que $K = \{M(\rho,\theta)\text{ tq } 0\le \rho\le f(\theta)\}$.
    
    Continuit{\'e} de~$f$~: soit~$\theta_0\in\R$ et $\varepsilon>0$. Soit
    $M(\rho_0,\theta_0)$ tel que $f(\theta_0)-\varepsilon<\rho_0<f(\theta_0)$.
    Donc $M\in\mathring K$ et il existe~$\alpha>0$ tel que la boule de
    centre~$M$ et de rayon~$\alpha$ est incluse dans~$K$ (faire un dessin).
    Ainsi, pour tout~$\theta$
    suffisament proche de~$\theta_0$ on a~$f(\theta)\ge OM > f(\theta_0)-\varepsilon$.
    Consid{\'e}rons alors une hypoth{\'e}tique suite~$(\theta_k)$ de r{\'e}els convergeant
    vers~$\theta_0$ telle que la suite $(f(\theta_k))$ ne converge pas vers~$f(\theta_0)$.
    Comme la suite $(f(\theta_k))$ est born{\'e}e on peut, quitte {\`a} extraire une
    sous-suite, supposer qu'elle converge vers un r{\'e}el~$\ell$ et le raisonnement
    pr{\'e}c{\'e}dent montre que~$\ell>f(\theta_0)$. Si $M_k$ d{\'e}signe le point de~$K$
    {\`a} la distance~$f(\theta_k)$ dans la direction d'angle polaire~$\theta_k$
    alors la suite $(M_k)$ converge vers le point $M(\ell,\theta_0)$ qui doit
    appartenir {\`a}~$K$ par compacit{\'e}, mais qui contredit la d{\'e}finition de~$f(\theta_0)$.
  \item Si $g$ ne s'annule pas alors $ \int_{x=0}^{\pi}g(x)\sin(x)\,d x \ne 0$.
    Si $g$ ne s'annule qu'en~$\alpha\in{[0,\pi]}$ alors $g$ est de signe constant
    sur~$[0,\alpha]$ et sur~$[\alpha,\pi]$, les signes sont oppos{\'e}s, et on obtient encore
    une contradiction en consid{\'e}rant $ \int_{x=0}^{\pi}g(x)\sin(x-\alpha)\,d x$
    qui vaut~$0$ (d{\'e}velopper le sinus).
  \item On choisit $O=G$. On a 
    $\iint_{M\in K}\vec{OM} = \vec 0$, soit
    ${ \int_{\theta=0}^{2\pi} f^3(\theta)\cos\theta\, d\theta =
      \int_{\theta=0}^{2\pi} f^3(\theta)\sin\theta\, d\theta = 0}$, soit encore~:
    $ \int_{\theta=0}^{\pi} (f^3(\theta)-f^3(\theta+\pi))\cos\theta\, d\theta =
      \int_{\theta=0}^{\pi} (f^3(\theta)-f^3(\theta+\pi))\sin\theta\, d\theta = 0$.
    D'apr{\`e}s la question pr{\'e}c{\'e}dente, il existe $\alpha\ne\beta\in{]0,\pi[}$ tels que
    $f(\alpha) = f(\alpha+\pi)$ et $f(\beta) = f(\beta+\pi)$, ce qui prouve qu'il y
    a au moins deux diam{\`e}tres de~$K$ dont~$O=G$ est le milieu. On prouve
    l'existence d'un troisi{\`e}me diam{\`e}tre en d{\'e}calant l'origine des angles
    polaires de fa\c con {\`a} avoir $f(0) = f(\pi)$, ce qui est possible vu l'existence
    de~$\alpha,\beta$.
\end{enumerate}
\fincorrection
\correction{004745}
Pour $\|x\| \le 1$ et $\|y\| \le 1$ on a $\|f(x)-f(y)\| = \|x-y\|$.

Pour $\|x\| \le 1< \|y\|$ on a
$\|f(x)-f(y)\| \le \|x-y\| + \Bigr\|y - \frac y{\|y\|}\Bigr\|
= \|x-y\| + \|y\| - 1 \le \|x-y\| + \|y\| - \|x\| \le 2\|x-y\|$.

Pour $1 < \|x\| \le \|y\|$ on a $\|f(x) - f(y)\| \le
\Bigl\|\frac x{\|x\|} - \frac y{\|x\|}\Bigr\| + 
\Bigl\|\frac y{\|x\|} - \frac y{\|y\|}\Bigr\| \le
\frac{\|x-y\| + \|y\| - \|x\|}{\|x\|} \le \frac{2\|x-y\|}{\|x\|}$.


Remarque~: dans le cas o{\`u} la norme est euclidienne, $f(x)$ est le projet{\'e}
de $x$ sur la boule unit{\'e}, c'est-{\`a}-dire le point de la boule unit{\'e} le plus
proche de~$x$. Dans ce cas, $f$ est $1$-lipschitzienne. Dans le cas d'une norme
non euclidienne on peut avoir $\|f(x)-f(y)\| > \|x-y\|$, par exemple
avec $x=(1,1)$ et $y=(\frac12,\frac32)$ dans $\R^2$ pour $\|\ \|_\infty$.
\fincorrection
\nocorrection
\nocorrection
\nocorrection
\nocorrection
\nocorrection
\nocorrection
\nocorrection
\nocorrection
\nocorrection
\correction{004755}
On construit $(s_k)$ de proche en proche de sorte que pour tout~$n$ fix{\'e}
la suite $(y_n^{s_k})$ soit convergente vers $z_n$.
Comme $\sum_{n\le N}(y_n^{s_k})^2$ est born{\'e}e ind{\'e}pendamment de~$N$ et~$k$ la s{\'e}rie
$\sum_nz_n^2$ a ses sommes partielles born{\'e}es donc converge.
On a alors $(x\mid y^{s_k})\xrightarrow[k\to\infty]{}(x\mid z)$ pour toute suite~$x$
{\`a} support fini, puis pour toute suite de carr{\'e} sommable par interversion de
limites.
\fincorrection
\correction{004756}
\begin{enumerate}
  \item Soit $(e_1,\dots,e_p)$ une base de~$E$. On remplace la norme sur $E$
         par la norme infinie associ{\'e}e {\`a} $(e_1,\dots,e_p)$. Alors
         $\|\hskip-1pt|u^n\|\hskip-1pt| \le \sum_{i=1}^p\|u^n(e_i)\|$.
  \item Trigonaliser fortement $u$ (ou son prolongement au complexifi{\'e} de~$E$).
         Comme $(u^n)$ est born{\'e}, les valeurs propres de~$u$ sont de module inf{\'e}rieur
         ou {\'e}gal {\`a}~$1$, et pour celles de module $1$ le bloc triangulaire associ{\'e}
         est en fait diagonal.
         On trouve $\frac1{n+1}\sum_{i=0}^n u^i \xrightarrow[n\to\infty]{}$ projection sur
         $\mathrm{Ker}(u-\mathrm{id})$ parall{\`e}lement {\`a} $\Im(u-\mathrm{id})$.
\end{enumerate}
 \fincorrection
\nocorrection
\nocorrection
\nocorrection
\correction{004760}
\begin{enumerate}
  \item 
  \item 
  \item 
  \item $(P_n)$ converge vers $0$ pour $a\in{]-2,2[}$ et vers $1$ pour $a=2$.
    La suite est non born{\'e}e si $|a| > 2$~; elle est born{\'e}e divergente pour $a=-2$.
  \item $(X/b)^n$ converge vers $1$ pour $N_b$ et vers $0$ pour $N_a$.
\end{enumerate}
 \fincorrection
\nocorrection
\correction{004762}
\begin{enumerate}
  \item 
  \item 
  \item $f(x)= \int_{t=0}^x\sin(x-t)(f(t)+f''(((t))\,d t$,
    $f''(x) = (f(x)+f''(x))-f(x)$.
\end{enumerate}
 \fincorrection
\correction{004763}
$(AB)_{ij} = \sum_{k=1}^n A_{ik}B_{kj}  \Rightarrow 
 (AB)_{ij}^2 \le \sum_{k=1}^n A_{ik}^2 \times \sum_{k=1}^n B_{kj}^2$.
\fincorrection
\correction{004764}
Si $A$ est une matrice de rang $r>0$ telle que $p(A) = 0$
alors pour toute matrice $M$ de rang $<r$ on peut trouver $P$ et $Q$
telles que $M=PAQ$ d'o{\`u} $P(M)=0$. Donc $p$ est nulle sur toute matrice
de rang~$1$ et par in{\'e}galit{\'e} triangulaire sur tout matrice.
\fincorrection
\nocorrection
\nocorrection
\nocorrection
\nocorrection
\nocorrection
\nocorrection
\correction{004771}
\begin{enumerate}
  \item $2^n\sum_{i=1}^n\|u_i\|^2$.
  \item Supposons qu'il existe une norme euclidienne $\|\ \|$ et deux r{\'e}els
    $\alpha,\beta>0$ tels que $\alpha\|u\|_\infty\le\|u\|\le\beta\|u\|_\infty$
    pour tout $u\in\mathcal{C}([0,1],\R)$.
    On pose $u(x) = 1-2|x|$ pour $x\in[\frac{-1}2,\frac12]$ et $u(x) = 0$
    sinon. Soit $n\in\N$ et pour $1\le i\le n$~: $u_i(x) = u((n+1)x-i)$.
    Alors $\sum_\sigma\bigl\|\sum_{i=1}^n\sigma(i)u_i\bigr\|^2 \le 2^n\beta^2$
    et $2^n\sum_{i=1}^n\|u_i\|^2 \ge 2^nn\alpha^2$ donc ces deux sommes ne peuvent
    rester {\'e}gales quand $n\to\infty$.
  \item M{\^e}me construction. On trouve
    $$\sum_\sigma\bigl\|\sum_{i=1}^n\sigma(i)u_i\bigr\|^2 \le 2^n\beta^2\|u\|_p^2\Bigl(\frac{n}{n+1}\Bigr)^{2/p}$$
    et $2^n\sum_{i=1}^n\|u_i\|^2 \ge 2^n\alpha^2\|u\|_p^2\frac{n}{(n+1)^{2/p}}$.
\end{enumerate}
 \fincorrection
\nocorrection
\correction{004773}
\begin{enumerate}
  \item 
  \item On prouve la convexit{\'e} de~$B$. Soient $x,y\in B$, $t\in{[0,1]}$
    et $z=(1-t)x + ty$. On a $N^2(z) \le 2t^2 + 2(1-t)^2$, d'o{\`u} $N(z)\le 1$
    si $t=\frac12$. Ceci prouve d{\'e}j{\`a} que $B$ est stable par milieu, et on en
    d{\'e}duit par r{\'e}currence sur~$n\in\N$ que $z\in B$ si $t$ est de la
    forme $a/2^n$ avec $a\in{[[0,2^n]]}$.
    
    Si $t$ n'est pas de cette forme, on {\'e}crit $t$ comme barycentre
    de deux nombres dyadiques~$t=u\frac{a}{2^n} + (1-u)\frac{b}{2^n}$ en
    faisant en sorte que $u$ soit arbitrairement proche de $\frac12$.
    Si c'est possible,
    on obtient que $z$ est barycentre de deux {\'e}l{\'e}ments de~$B$ avec les
    coefficients $u$ et $1-u$, d'o{\`u} $N^2(z)\le 2u^2+2(1-u)^2\xrightarrow[u\to1/2]{}$.
    Reste donc {\`a} choisir $n$, $a$, $b$~: pour $n$ donn{\'e},
    on choisit $a = {[2^n]}-n$ et $b={[2^n]}+n$. C'est possible car
    $[2^nt]\sim 2^nt$ et on est dans le cas $0<t<1$
    donc on a bien $a,b\in {[[0,2^n]]}$ si $n$ est suffisament grand.
    Il vient $u=\frac{b-2^nt}{b-a}$, quantit{\'e} comprise entre $\frac{n-1}{2n}$
    et $\frac12$ et donc qui tend bien vers~$\frac12$.
    
    Remarque~: la condition $(iii)$ est aussi n{\'e}cessaire, donc une norme
    est une application v{\'e}rifiant $(i)$, $(ii)$ et $(iii)$.
\end{enumerate}
 \fincorrection
\nocorrection
\nocorrection
\nocorrection
\nocorrection
\nocorrection
\correction{004779}
\begin{enumerate}
  \item On proc{\`e}de par r{\'e}currence sur~$n$. Pour $n=1$, l'application ${[0,+\infty[}\ni\lambda_1 \mapsto \|x-\lambda_1a_1\|$
est continue, constante si $a_1 = 0$ et tend vers $+\infty$ quand $\lambda_1\to+\infty$ si $a_1\ne0$.
Dans les deux cas elle admet un minimum.

Pour $n\ge 2$, soit $C'=\Bigl\{\sum_{i=1}^{n-1}\lambda_{i}a_{i},\lambda_{i}\ge 0\Bigr\}$.
Soit pour $\lambda_n\in{[0,+\infty[}$~ : $\varphi(\lambda_n) = d(x-\lambda_na_n,C')$.
La distance {\`a} $C'$ est 1-lipschitzienne donc
$\varphi(\lambda_n) \ge d(-\lambda_na_n,C') - \|x\| = \lambda_nd(-a_n,C') - \|x\|\xrightarrow[\lambda_n\to+\infty]{}+\infty$.
{\'E}tant continue, $\varphi$ admet un minimum sur $[0,+\infty[$ et on
applique l'hypoth{\`e}se de r{\'e}currence {\`a} $x-\lambda_na_n$.
  \item 
\end{enumerate}
 \fincorrection
\correction{004780}
On proc{\`e}de par r{\'e}currence sur~$n=\dim E$. Pour $n=1$, en confondant $E$ et $\R$,
$C$ est un intervalle dense, c'est $\R$.
Pour $n\ge 2$, soit $E = H\oplus{<a>}$ o{\`u} $H$ est un hyperplan de~$E$.
On montre ci-dessous que $C' = C \cap H$ est une partie de~$H$ convexe et dense,
donc {\'e}gale {\`a}~$H$, d'o{\`u} $H\subset C$ et ce pour tout~$H$. Ainsi $C=E$.

Densit{\'e} de $C'$~: soit $x\in H$, et $(y_k)$, $(z_k)$ des suites d'{\'e}l{\'e}ments de~$C$
convergeant respectivement vers $x+a$ et $x-a$. On {\'e}crit $y_k = y'_k + \lambda_ka$
et $z_k = z'_k + \mu_ka$ avec $y'_k,z'_k\in H$ et $\lambda_k,\mu_k\in\R$.
Par {\'e}quivalence des normes en dimension finie, on a $\lambda_k\xrightarrow[k\to\infty]{}1$
et $\mu_k\xrightarrow[k\to\infty]{}-1$, donc le point $x_k = \frac{\lambda_kz_k-\mu_ky_k}{\lambda_k-\mu_k}$
est bien d{\'e}fini et appartient {\`a}~$C'$ pour $k$ assez grand, et converge vers~$x$.

Remarque~: Si $E$ est de dimension infinie, alors il contient des hyperplans non ferm{\'e}s,
donc des parties strictes, convexes denses.
\fincorrection
\nocorrection
\correction{004782}
\begin{enumerate}
  \item $\overline P = P$,
             $\mathring P = \{\text{fonctions strictement positives}\}$.
  \item $\overline P = P$,
             $\mathring P = \varnothing$.
\end{enumerate}
 \fincorrection
\correction{004783}
$\overline F = F$, $\mathring F = \varnothing$.
\fincorrection
\correction{004784}
Oui pour $\pm\mathrm{id}_{\R^n}$, non pour les autres (les sym{\'e}tries non triviales).

$\mathrm{id}_{\R^n}$ est isol{\'e} car si $u\ne\mathrm{id}_{\R^n}$ et $u^2=\mathrm{id}_{\R^n}$ alors $-1$
est valeur propre de~$u$ et $\|u-\mathrm{id}_{\R^n}\|\ge 2$. De m{\^e}me pour~$-\mathrm{id}_{\R^n}$.

Si $u$ est une sym{\'e}trie non triviale, soit $(e_1,\dots,e_n)$ une base propre
de~$u$ avec $u(e_1)=e_1$ et $u(e_2)=-e_2$. Pour $p\in\N^*$ soit~$u_p\in\mathcal{L}(\R^n)$
d{\'e}finie par $u_p(e_1)=e_1+e_2/p$ et $u_p(e_i)=u(e_i)$ pour~$i\ge 2$.
On a $u_p^2=\mathrm{id}_{\R^n}$, $u_p\ne u$ et $u_p\xrightarrow[p\to\infty]{}u$.
\fincorrection
\nocorrection
\nocorrection
\correction{004787}
\begin{enumerate}
  \item $GL_n(\C) = \det^{-1}(\C\setminus\{0\})$ est ouvert.
    Il est dense car $A\in\mathcal{M}_n(\C)$ quelconque est limite des matrices
    $A - \frac1p I$ inversibles pour presque tout entier $p$ ($A$ a un
    nombre fini de valeurs propres).
  \item Toute matrice triangulaire est limite de matrices triangulaires
    {\`a} coefficients diagonaux distincts.
  \item $\begin{pmatrix}\lambda &0\cr 0&\lambda\cr\end{pmatrix}
             = \lim_{p\to\infty}\begin{pmatrix}\lambda &1/p\cr 0&\lambda\cr\end{pmatrix}$
    donc une matrice triangulaire {\`a} valeurs propres non distinctes est
    limite de matrices non diagonalisables.
    Par conjugaison, la fronti{\`e}re de $D_n(\C)$ contient
    l'ensemble des matrices ayant au moins une valeur propre multiple.
    
    R{\'e}ciproquement, soit $(A_k)$ une suite de matrices non diagonalisables
    convergeant vers une matrice~$A$. Les matrices $A_k$ ont toutes au moins une valeur
    propre multiple, et ces valeurs propres sont born{\'e}es
    (car si $\lambda$ est une valeur propre de~$M$ alors $|\lambda| \le \| M\|$ en prenant une norme
    sur $\mathcal{M}_n(\C)$ subordonn{\'e}e {\`a} une norme sur $\C^n$) donc on peut trouver
    une suite $(z_k)$ de complexes convergeant vers un complexe $z$ telle
    que $\chi'_{A_k}(z_k) = 0$. A la limite on a $\chi_A(z) = 0$ ce qui prouve
    que $A$ a au moins une valeur propre multiple.
    
    Conclusion~: la fronti{\`e}re de~$D_n(\C)$ est exactement l'ensemble des
    matrices diagonalisables ayant au moins une valeur propre multiple et l'int{\'e}rieur
    de~$D_n(\C)$ est l'ensemble des matrices {\`a} valeurs propres distinctes.
\end{enumerate}
 \fincorrection
\correction{004788}
Si $N$ est nilpotente, on peut se ramener au cas o{\`u} $N$
est triangulaire sup{\'e}rieure stricte.

Soit alors $P=\mathrm{diag}(1,\alpha,\dots,\alpha^{n-1})$ avec $\alpha\in\C^*$.
Le coefficient g{\'e}n{\'e}ral de $P^{-1}NP$ est $\alpha^{j-i}N_{ij}\xrightarrow[\alpha\to0]{}0$.

R{\'e}ciproquement, s'il existe une suite $(N_k)$ de matrices semblables
{\`a}~$N$ convergeant vers la matrice nulle, alors par continuit{\'e} du
polyn{\^o}me caract{\'e}ristique, on a $\chi_N=(-X)^n$ et $N$ est nilpotente.
\fincorrection
\correction{004789}
\begin{enumerate}
  \item $\Omega$ est ouvert~: si~$P$ a $n$ racines distinctes $a_1<a_2<\dots< a_n$
    on choisit $b_0 < a_1 < b_1 < a_2 <\dots < a_n < b_n$. La suite $(P(b_0),\dots,P(b_n))$
    est constitu{\'e}e de termes non nuls de signes altern{\'e}s, il en est de m{\^e}me
    pour la suite $(Q(b_0),\dots,Q(b_n))$ o{\`u}~$Q$ est un polyn{\^o}me unitaire arbitraire
    suffisament proche de~$P$ (pour une norme quelconque).
    
    $\Omega$ n'est pas ferm{\'e} car $\varnothing\ne\Omega\ne\R^n$ et $\R^n$ est connexe.
  \item D{\'e}j{\`a}, si l'on munit $\R^n$ et $\R_n[X]$ de normes convenables,
    $f$ est une isom{\'e}trie bicontinue donc $f(\overline\Omega) = \overline{f(\Omega)}$.
    Montrons que $\overline{f(\Omega)}$ est l'ensemble $\cal S$ des
    polyn{\^o}mes de~$\R_n[X]$ unitaires et scind{\'e}s sur~$\R$.

    Si $P = X^n + a_{n-1}X^{n-1} + \dots + a_0$, notons~$M(P)$ la matrice compagne de~$P$.
    
    Si $P\in\cal S$ alors $M(P)$ est $\R$-trigonalisable, donc limite de matrices
    {\`a} valeurs propres r{\'e}elles distinctes. Les polyn{\^o}mes caract{\'e}ristiques de
    ces matrices, au signe pr{\`e}s, appartienent {\`a}~$f(\Omega)$ et convergent vers~$P$
    d'o{\`u} ${\cal S} \subset \overline{f(\Omega)}$.
    
    Si $(P_k)$ est une suite de polyn{\^o}mes de~$f(\Omega)$ convergeant vers~$P$
    alors il existe une suite $(O_k)$ de matrices orthogonales telle que
    ${}^tO_kM(P_k)O_k$ est triangulaire sup{\'e}rieure (m{\'e}thode de Schmidt). Quitte
    {\`a} extraire une sous-suite, on peut supposer que $O_k$ converge vers une
    matrice orthogonale~$O$ et donc ${}^tOM(P)O$ est aussi triangulaire
    sup{\'e}rieure ce qui implique que~$P\in\cal S$.
\end{enumerate}
 \fincorrection
\nocorrection
\correction{004791}
Remarquer que la restriction de $f$ {\`a} toute partie compacte est uniform{\'e}ment
continue.
\fincorrection
\nocorrection
\nocorrection
\nocorrection
\nocorrection
\nocorrection
\nocorrection
\correction{004798}
\begin{enumerate}
  \item 
  \item Soit $F$ est un tel ferm{\'e}, et $a\in F$.
On prend $f(x) = x + d(x,F)$ si $0\le x\le a$
et $f(x) = x-d(x,F)$ si $a\le x\le 1$.
\end{enumerate}
 \fincorrection
\correction{004799}
\begin{enumerate}
  \item Pour simplifier, on suppose $z=0$ (sinon, se placer dans la base $(1,X-z,\dots,(X-z)^d)$
    et invoquer l'{\'e}quivalence des normes en dimension finie).

    Soit $P_n(x) = a_{n,0} + a_{n,1}x + \dots + a_{n,d}x^d$. La suite $(P_n)$ {\'e}tant convergente
    est born{\'e}e donc il existe $M\in\R$ tel que $|a_{n,k}|\le M$ pour
    tous $n,k$. De plus, $a_{n,0}\xrightarrow[n\to\infty]{} a_0=0$ et
    $a_{n,1}\xrightarrow[n\to\infty]{}a_1\ne 0$.

    Posons alors $Q_n(x) = -\frac{\strut a_{n,0} + a_{n,2}x^2 + \dots + a_{n,d}x^d}{a_{n,1}}$
    (bien d{\'e}fini si $n$ est assez grand).
    On va montrer que $Q_n$ v{\'e}rifie les hypoth{\`e}ses du th{\'e}or{\`e}me du point fixe sur
    $\overline{B(0,\delta)}$ pour tout~$n$ assez grand si $\delta$ est choisi assez petit,
    ce qui implique l'existence et l'unicit{\'e} d'une racine pour $P_n$ dans $\overline{B(0,\delta)}$.
    
    $Q_n(\overline{B(0,\delta)})\subset\overline{B(0,\delta)}$~?
    Soit $x\in\overline{B(0,\delta)}$~: on a $$|Q_n(x)| \le \frac{|a_{n,0}| + M(\delta^2+\dots+\delta^d)}{|a_{n,1}|}
    \xrightarrow[n\to\infty]{}\frac{M(\delta^2+\dots+\delta^d)}{|a_1|}.$$
    On choisit $\delta>0$ tel que $\frac{M(\delta+\dots+\delta^{d-1})}{|a_1}\le\frac12$.
    Il existe alors $N_1\in\N$ tel que $\frac{|a_{n,0}| + M(\delta^2+\dots+\delta^d)}{|a_{n,1}|}\le\delta$
    pour tout~$n\ge N_1$.
    
    $Q_n$ est contractante sur $\overline{B(0,\delta)}$~?
    Soient $x,y\in\overline{B(0,\delta)}$. On a~:
    $$\begin{aligned}|Q_n(x)-Q_n(y)|
    &\le\frac{|a_{n,2}||x^2-y^2|+\dots+|a_{n,d}||x^d-y^d|}{|a_n,1|}\cr
    &\le|x-y|\frac{|a_{n,2}||x+y|+\dots+|a_{n,d}||x^{d-1}+\dots+y^{d-1}|}{|a_{n,1}|}\cr
    &\le|x-y|\frac{M(2\delta+\dots+d\delta^{d-1})}{|a_{n,1}|}.\cr\end{aligned}$$
    Quitte {\`a} diminuer $\delta$ on peut imposer $\frac{M(2\delta+\dots+d\delta^{d-1})}{|a_1|}\le\frac12$
    et donc $\frac{M(2\delta+\dots+d\delta^{d-1})}{|a_{n,1}|}\le \frac23$ pour tout~$n\ge N_2$
    et $Q_n$ est $\frac23$-lipschitzienne.
  \item Voir r{\'e}ponse pr{\'e}c{\'e}dente. Y a-t-il une r{\'e}ponse plus simple pour~{1)}~?
  \item Si $z$ est z{\'e}ro d'ordre~$k$ de~$P$ alors il existe $\delta>0$ tel
    que pour tout~$n$ assez grand, $P_n$ a exactement $k$ racines compt{\'e}es avec leur ordre de
    multiplicit{\'e} dans~$\overline{B(0,\delta)}$. Ceci est une cons{\'e}quence du {\it th{\'e}or{\`e}me des
    r{\'e}sidus\/} largement hors programme\dots
\end{enumerate}
 \fincorrection
\correction{004800}
Si $f$ est constante c'est {\'e}vident. Sinon, on a facilement $|f(z)|\xrightarrow[|z|\to\infty]{}\infty$.
Consid{\'e}rons un ferm{\'e} $F$ et une suite $(f(z_n))$ d'{\'e}l{\'e}ments de~$f(F)$ convergeant
vers $Z\in\C$. D'apr{\`e}s la remarque, la suite $(z_n)$ est born{\'e}e, elle
admet une valeur d'adh{\'e}rence $z\in F$ et $Z=f(z)\in f(F)$.

Remarque~: ce r{\'e}sultat est faux pour une fonction polynomiale sur $\C^p$ avec $p\ge 2$,
prendre par exemple $f(x,y) = x$ sur $\C^2$ et $F = \{(x,y)\in\C^2$ tq $xy=1\}$.
\fincorrection
\correction{004801}
On suppose $P$ non constant, sans quoi le r{\'e}sultat est trivial.
Soit $S(X) = \sup(|P(x)|,\ x\in X)$.
On a par inclusion et continuit{\'e}~: $S(\mathrm{Fr}(U)) \le S(\overline U) = S(U)$.
Soit $x\in\overline U$ tel que $|P(x)| = S(\overline U)$. On d{\'e}montre
par l'absurde que $x\in\mathrm{Fr}(U)$, ce qui entra{\^\i}ne l'{\'e}galit{\'e} demand{\'e}e.
Supposons donc $x\in U$ et soit $n=\deg(P)$. Alors pour $\rho>0$ suffisament petit,
et $\theta\in\R$, on a $x+\rho e^{i\theta}\in U$ et~:
$$P(x+\rho e^{i\theta}) = P(x) + \rho e^{i\theta}P'(X) + \dots + \frac{\rho^n e^{in\theta}}{n!}P^{(n)}(x).$$
avec $P^{(n)}(x) = P^{(n)} \ne 0$. On en d{\'e}duit~:
$$2\pi |P(x)| = \Bigl| \int_{\theta=0}^{2\pi} P(x+\rho e^{i\theta})\,d \theta\Bigr|
\le  \int_{\theta=0}^{2\pi} |P(x+\rho e^{i\theta})|\,d \theta \le 2\pi\, S(U) = 2\pi|P(x)|.$$
On en d{\'e}duit que les in{\'e}galit{\'e}s sont des {\'e}galit{\'e}s, et en particulier que
la quantit{\'e} $|P(x+\rho e^{i\theta})|$ est ind{\'e}pendante de $\rho$ et $\theta$.
Il y a contradiction car $|P(x+\rho e^{i\theta})|^2$ est un polyn{\^o}me de
degr{\'e} $2n$ en $\rho$.
\fincorrection
\correction{004802}
\begin{enumerate}
  \item $f(rx) = rf(x)$ pour tout $r\in\N$ par r{\'e}currence, puis pour
    tout $r\in\Z$ par diff{\'e}rence, pour tout $r\in\Q$ par quotient et enfin
    pour tout $r\in\R$ par densit{\'e}. Dans le cas de $\C$-ev $f$ est $\R$-lin{\'e}aire
    mais pas forc{\'e}ment $\C$-lin{\'e}aire, ctrex~: $z \mapsto\overline z$ de $\C$
    dans $\C$.
  \item $\|f_{n+1}(x) - f_n(x)\|\le M2^{-n-1}$ donc la s{\'e}rie t{\'e}lescopique
    $\sum(f_{n+1}(x) - f_n(x))$ est uniform{\'e}ment convergente.
  \item $\|f_n(x+y) - f_n(x) - f_n(y)\|\le M2^{-n}$ donc
     $\|g(x+y) - g(x) - g(y)\|\le 0$ et $g$ est continue (limite uniforme des $f_n$)
     d'o{\`u} $g$ est lin{\'e}aire continue.
     $\|f(x) - g(x)\| = \|\sum_{k=0}^\infty(f_{k}(x) - f_{k+1}(x))\|
                        \le 2M$
     donc $f-g$ est born{\'e}e.
     Si $h$ est une application lin{\'e}aire telle que $f-h$ est born{\'e}e alors $g-h$ est aussi
     born{\'e}e ce qui entra{\^\i}ne $g=h$ par lin{\'e}arit{\'e}.
\end{enumerate}
 \fincorrection
\nocorrection
\nocorrection
\correction{004805}
\begin{enumerate}
  \item 
  \item Si $c\ne 0$ alors $P_c$ est continue pour toutes les normes $N_p$
    et $\|\!|P_c|\!\|_{N_p}=|c|^{-p}e^{|c|}$.
    Par contre $P_0$ n'est continue que pour~$N_0$ car si $p>0$ alors
    $N_p(x \mapsto e^{-n|x|})=\frac{p^pe^{-p}}{(n+1)^p}\xrightarrow[n\to\infty]{}0$
    donc la suite $(x \mapsto e^{-n|x|})$ converge vers la fonction nulle
    pour $N_p$, mais $P_0(x \mapsto e^{-n|x|})=1\rlap{\hskip9pt/}\to0$.
  \item Si $p<q$ alors $N_p(x \mapsto e^{-n|x|})/N_q(x \mapsto e^{-n|x|}) \to \infty$ losrque $n\to \infty$.
\end{enumerate}
 \fincorrection
\nocorrection
\nocorrection
\nocorrection
\nocorrection
\correction{004810}
Prendre une base.
\fincorrection
\correction{004811}
$\frac1{1-X^n} = \frac 1n \sum_{k=0}^{n-1} \frac{1}{1-\omega_k X}$,
$\omega_k = e^{2ik\pi/n}$.
Donc $1 = \sum_{k=0}^{n-1} \frac{1-X^n}{n(1-\omega_k X)}$.

Il s'agit de polyn{\^o}mes, donc on peut remplacer $X$ par $u$, d'o{\`u} :
$(\mathrm{id}_E - u^n)^{-1} = \frac1n \sum_{k=0}^{n-1}\bigl(\mathrm{id}_E - e^{2ik\pi/n}u\bigr)^{-1}$.
\fincorrection
\correction{004812}
$\|X+iY\|^2 = \|X\|^2 + \|Y\|^2$, $\|A(X+iY)\|^2 = \|AX\|^2+\|AY\|^2
         \le |||A|||_\R^2(\|X^2\|+\|Y\|^2)$ donc $|||A|||_\C \le |||A|||_\R$.
\fincorrection
\correction{004813}
\begin{enumerate}
  \item Non : $\|(x^2+1)^n\| = 2^n$.
  \item Oui : $\|\psi\| = \|A\|$.
  \item $\|\phi\| = e$, $\|\psi(x^n)\|/\|x^n\| \text{ si }m n  \Rightarrow  \psi$
             est discontinue.
\end{enumerate}
 \fincorrection
\correction{004814}
\begin{enumerate}
  \item $nv^{n-1}$.
  \item Si $u$ et $v$ sont continus,
             $n\|v^{n-1}\| \le 2\|u\|\,\|v^n\| \le 2\|u\|\,\|v^{n-1}\|\,\|v\|$.\par
             S'il existe $k$ tel que $v^k = 0$, on peut remonter jusqu'{\`a} $v = 0$, absurde.
             Sinon, on a aussi une contradiction.
\end{enumerate}
 \fincorrection
\correction{004815}
\begin{enumerate}
  \item 
  \item $\|P\| = N(P) + N(P') + N(P'') + \dots$ o{\`u} $N$ est une
norme quelconque sur~$F$.
\end{enumerate}
 \fincorrection
\correction{004816}
Les formes lin{\'e}aires $P \mapsto P(0)$, $P \mapsto P(1)$ et $P \mapsto P(2)$
constituent une base de $E_2^*$ donc engendrent les formes lin{\'e}aires
$P \mapsto P'(1)$, $P \mapsto P'(2)$ et $P \mapsto P'(3)$. Apr{\`e}s calculs, on trouve~:
$$\forall\ P\in E_2,\quad
\left\{\begin{aligned}2P'(1) &= &P(2)  & &      &-&P(0)\cr
      2P'(2) &= &3P(2) &-&4P(1) &+&P(0)\cr
      2P'(3) &= &5P(2) &-&8P(1) &+&3P(0).\cr
\end{aligned}\right.$$
En notant $P(0) = a$, $P(1) = b$ et $P(2) = c$ on doit
donc chercher~:
$$\|\hskip-1pt|\varphi\|\hskip-1pt|
=\frac12\sup\{|c-a| + 4|3c-4b+a| + 9|5c-8b+3a|,\text{ tq }|a|+|b|+|c|\le1\}.$$
La fonction $f$~: $(a,b,c) \mapsto|c-a| + 4|3c-4b+a| + 9|5c-8b+3a|$ est convexe donc
son maximum sur l'{\it icosa{\`e}dre\/} $I = \{(a,b,c)\text{ tq }|a|+|b|+|c|\le 1\}$
est atteint en l'un des sommets $(\pm1,0,0)$, $(0,\pm1,0)$, $(0,0,\pm1)$.

Finalement,~$\|\hskip-1pt|\varphi\|\hskip-1pt| = \frac12f(0,1,0) = 44$.
\fincorrection
\nocorrection
\correction{004818}
\begin{enumerate}
  \item 
  \item Si $x \notin \mathrm{Ker} f$ :
             $\forall\ y \in \mathrm{Ker} f,\ |f(x)| = |f(x-y)| \le \|f\|\,\|x-y\|$
             donc $|f(x)| \le \|f\|d(x,\mathrm{Ker} f)$.\par
             Soit $z \in E : z = \alpha x + y$ avec $y \in \mathrm{Ker} f$.
             Alors $|f(z)| = |\alpha||f(x)|$ et $\|z\| \ge |\alpha|d(x,\mathrm{Ker} f)$

             donc $\frac{|f(z)|}{\|z\|} \le \frac{|f(x)|}{d(x,\mathrm{Ker} f)}$.
\end{enumerate}
 \fincorrection
\correction{004819}
\begin{enumerate}
  \item $f(x_1)+f(x_2) \le \|\kern-1.2pt|f\|\kern-1.2pt|\;\|x_1+x_2\|
                            \le \|\kern-1.2pt|f\|\kern-1.2pt|\;(\|x_1-\varepsilon\| + \|x_2+\varepsilon\|)$.
  \item Prendre $\alpha$ compris entre le sup du premier membre et l'inf du troisi{\`e}me.
    Le sup et l'inf sont dans cet ordre d'apr{\`e}s la question pr{\'e}c{\'e}dente.
  \item   
  Rmq~: $\varphi$ est mal d{\'e}finie, il faut ajouter "$\varphi$ est lin{\'e}aire".
    On a {\'e}videmment $\|\kern-1.2pt|\varphi\|\kern-1.2pt| \ge \|\kern-1.2pt|f\|\kern-1.2pt|$
    puisque $\varphi$ prolonge~$f$, et il reste {\`a} montrer~:
    $$\forall x\in F,\ \forall\ t\in\R,\ |f(x)+t\alpha| \le \|\kern-1.2pt|f\|\kern-1.2pt|\;\|x+t\varepsilon\|.$$
    Pour~$t=0$ c'est un fait connu. Pour $t>0$, cela r{\'e}sulte de l'encadrement de~$\alpha$
    en prenant $x_1=-x/t$ et $x_2=x/t$. Pour~$t<0$, prendre $x_1=x/t$ et $x_2=-x/t$.
  \item Si $u^k = (u^k_n)_{n\in\N}\in E$ et $(u^k)_{k\in\N}$ est une suite de Cauchy,
    alors pour tout~$n\in\N$ la suite r{\'e}elle $(u^k_n)_{k\in\N}$ est de Cauchy dans~$\R$
    donc converge vers un r{\'e}el~$\ell_n$. De plus la suite $(u^k)_{k\in\N}$ est born{\'e}e dans~$E$ donc
    la suite~$\ell$ ainsi mise en {\'e}vidence est sommable (les sommes partielles de~$\sum|\ell_n|$
    sont major{\'e}es), et on montre que $\|u^k-\ell\|\xrightarrow[k\to\infty]{}0$ par interversion
    de limites.
  \item Prendre~$F_n = \{u\in E\text{ tq }u_k=0\text{ si }k\ge n\}$.
  \item D'apr{\`e}s la question~{\bf 3)} on peut construire $f_n$, forme
    lin{\'e}aire sur~$F+F_n$ telle que $f_{n+1}$ prolonge $f_n$ et a m{\^e}me norme que~$f_n$
    (donc $\|\kern-1.2pt|f_n\|\kern-1.2pt| = \|\kern-1.2pt|f\|\kern-1.2pt|$).
    Soit~$G = \cup_{n\in\N}(F+F_n)$ et $g$ la forme lin{\'e}aire sur~$G$ co{\"\i} ncidant avec
    chaque $f_n$ sur~$F+F_n$. $G$ est dense dans~$E$
    donc on peut prolonger~$g$ en $\varphi : E \to \R$ par uniforme continuit{\'e}. Il est alors
    clair que~$\varphi$ est une forme lin{\'e}aire prolongeant~$f$ et a m{\^e}me norme
    que~$f$.
\end{enumerate}
 \fincorrection
\nocorrection
\correction{004821}
\begin{enumerate}
  \item Si $\nu$ est subordonn{\'e}e {\`a} $\|\ \|$~:
    on a $|\lambda| \le \nu(f^p)^{1/p}$ pour toute valeur propre $\lambda$
    et tout $p\ge 1$, donc il suffit de prouver que la suite $(x_p = \nu(f^p)^{1/p})$
    est convergente. Soit $\ell = \inf\{x_p,\ p\ge 1\}$, $\varepsilon > 0$
    et $p\ge 1$ tel que $x_p \le \ell+\varepsilon$. Pour $n> p$ on note
    $n-1=pq+r$ la division euclidienne de $n-1$ par $p$ et l'on a~:
    $$\nu(f^n) = \nu((f^p)^q\circ f^{r+1})
               \le \nu(f^p)^q\nu(f^{r+1})$$
    d'o{\`u}~:
    $$\ell \le x_n \le x_p^{pq/n}x_{r+1}^{(r+1)/n}\le (\ell+\varepsilon)^{pq/n}\max(x_1,\dots,x_p)^{(r+1)/n}.$$
    Le majorant tend vers $\ell+\varepsilon$ quand $n$ tend vers l'infini
    donc pour $n$ assez grand on a $\ell \le x_n \le \ell+2\varepsilon$
    ce qui prouve la convergence demand{\'e}e.
    
    Dans le cas o{\`u} $\nu$ est une norme quelconque sur $\mathcal{L}(E)$, il existe
    une norme subordonn{\'e}e $\mu$ et deux r{\'e}els $a,b>0$ tels que $a\mu \le \nu\le b\mu$
    et donc les suites $(\nu(f^p)^{1/p})$ et  $(\mu(f^p)^{1/p})$ ont m{\^e}me
    limite par le th{\'e}or{\`e}me des gens d'armes.
    Remarque~: il r{\'e}sulte de ceci que $\lim_{p\to\infty}(\nu(f^p)^{1/p})$
    est ind{\'e}pendant de $\nu$.
  \item Consid{\'e}rer la matrice de $f^p$ dans une base propre pour~$f$.
  \item On sait que $f^p = \sum_{\lambda\in\mathrm{spec}(f)} \lambda^pP_\lambda(p)$
    o{\`u} $P_\lambda$ est un polyn{\^o}me. D'o{\`u} $\rho(f) \le \nu(f^p)^{1/p}\le \rho(f)+ o(1)$
    et donc $\nu(f^p)^{1/p}\xrightarrow[p\to\infty]{} \rho(f)$ (thm du rayon spectral).
\end{enumerate}
 \fincorrection
\correction{004822}
Si $u(\mathring B(0,1))$ est ouvert alors il engendre $\R^m$ donc
$u$ est surjective.

Si $u$ est surjective, soit $A=u(\mathring B(0,1))$. $A$ est convexe,
born{\'e}, sym{\'e}trique par rapport {\`a}~$0$ et la r{\'e}union des homoth{\'e}tiques
de~$A$ est {\'e}gale {\`a}~$\R^m$~; la jauge associ{\'e}e {\`a}~$A$ est une norme
sur~$\R^m$ {\'e}quivalente {\`a} l'une des normes usuelles donc
$A$ contient une boule de centre $0$ et, par homoth{\'e}tie-translation, tout
ouvert de~$\R^n$ a une image ouverte dans~$\R^m$.
\fincorrection
\nocorrection
\nocorrection
\correction{004825}
$E\setminus B$ est connexe par arcs et contient au moins un point $a\in A$.
Soit $x\in E\setminus B$ et $\varphi : {[0,1]} \to {E\setminus B}$ un arc
continu joignant $a$ {\`a} $x$ dans $E\setminus B$.
Alors $\varphi^{-1}(A) = \varphi^{-1}(A\cup B)$ est non vide, relativement ouvert
et relativement ferm{\'e} dans $[0,1]$, donc c'est $[0,1]$ ce qui prouve
que $x\in A$.
\fincorrection
\correction{004826}
Le sens $H$ est ferm{\'e} $ \Rightarrow $ $E\setminus H$ n'est pas connexe (par arcs) est {\'e}vident.
R{\'e}ciproquement, si $H$ n'est pas ferm{\'e} alors $\overline H = E$.
Soient $a,b\in E\setminus H$ et $(x_n)$ une suite d'{\'e}l{\'e}ments de~$H$ telle que
$x_0 = 0$ et $x_n\xrightarrow[n\to\infty]{}a-b$. On d{\'e}finit un arc continu $\varphi : {[0,1]} \to {E\setminus H}$
reliant $a$ {\`a} $b$ par~: $\varphi$ est affine sur $[\frac1{n+2},\frac1{n+1}]$,
$\varphi(\frac1{n+1})=b+x_n$ et $\varphi(0)=a$.
\fincorrection
\correction{005839}
\textbf{Cas de la boule fermée.}  Soit $B=\{u\in E/\;\|u\|\leqslant 1\}$. Soient $(x,y)\in B^2$ et $\lambda\in[0,1]$.

\begin{center}
$\|\lambda x+(1-\lambda)y\|\leqslant \lambda\|x\| + (1-\lambda)\|y\|\leqslant\lambda+1-\lambda= 1$.
\end{center}

Ainsi, $\forall(x,y)\in B^2$, $\forall\lambda\in[0,1]$, $\lambda x+(1-\lambda)y\in B$ et donc $B$ est convexe.

\textbf{Cas de la boule ouverte.} Soit $B=\{u\in E/\;\|u\|< 1\}$. Soient $(x,y)\in B^2$ et $\lambda\in[0,1]$.

Puisque $0\leqslant \lambda\leqslant 1$ et $0\leqslant\|x\|<1$, on en déduit que $\lambda\|x\|<1$. Comme $(1-\lambda)\|y\|\leqslant1$ (et même $<1$) et donc

\begin{center}
$\|\lambda x+(1-\lambda)y\|\leqslant \lambda\|x\| + (1-\lambda)\|y\|<1$.
\end{center} 

\begin{center}
\shadowbox{
La boule unité fermée (ou ouverte) de l'espace vectoriel normé $(E,\|\;\|)$ est un convexe de l'espace vectoriel $E$.
}
\end{center}
\fincorrection
\correction{005840}
\begin{enumerate}
 \item  Puisque $p>0$ et $q>0$, $1= \frac{1}{p}+ \frac{1}{q}> \frac{1}{p}$ et donc $p>1$. De même, $q>1$. D'autre part, $q= \frac{p}{p-1}$.

  \begin{enumerate}
  \item L'inégalité est immédiate quand $y=0$. Soit $y > 0$ fixé.

Pour $x\geqslant0$, on pose $f(x)= \frac{x^p}{p}+ \frac{y^q}{q} -xy$. Puisque $p > 1$, la fonction $f$ est dérivable sur $[0,+\infty[$ et $\forall x\geqslant0$, $f'(x)=x^{p-1}-y$. $f$ admet donc un minimum en $x_0=y^{1/(p-1)}$ égal à 

\begin{center}
$f\left(y^{1/(p-1)}\right)= \frac{y^{p(p-1)}}{p}+ \frac{y^{p/(p-1)}}{q}-y^{1/(p-1)}y=y^{p/(p-1}\left( \frac{1}{p}+ \frac{1}{q}-1\right)= 0$.
\end{center}

Finalement, $f$ est positive sur $[0,+\infty[$ et donc

\begin{center}
\shadowbox{
$\forall x\geqslant0$, $\forall y\geqslant0$, $xy\leqslant \frac{x^p}{p}+ \frac{y^q}{q}$.
}
\end{center}

  \item Posons $A=\sum_{k=1}^{n}|a_k|^p$ et $B =\sum_{k=1}^{n}|b_k|^q$.

Si $A$ (ou $B$) est nul, tous les $a_k$ (ou tous les $b_k$) sont nuls et l'inégalité est vraie.

On suppose dorénavant que $A > 0$ et $B > 0$. D'après la question a),

\begin{center}
$\sum_{k=1}^{n} \frac{|a_k|}{A^{1/p}}\times \frac{|b_k|}{B^{1/q}}\leqslant\sum_{k=1}^{n}\left( \frac{|a_k|^p}{pA}+ \frac{|b_k|^q}{qB}\right) = \frac{1}{pA}\sum_{k=1}^{n}|a_k|^p+ \frac{1}{qB}\sum_{k=1}^{n}|b_k|^q= \frac{1}{pA}\times A+ \frac{1}{qB}\times B= \frac{1}{p}+ \frac{1}{q}=1$,
\end{center}

et donc $\sum_{k=1}^{n}|a_k||b_k|\leqslant A^{1/p}B^{1/q}=\left(\sum_{k=1}^{n}|a_k|^p\right)^{1/p}\left(\sum_{k=1}^{n}|b_k|^q\right)^{1/q}$. Comme $\left|\sum_{k=1}^{n}a_kb_k\right|\leqslant\sum_{k=1}^{n}|a_k||b_k|$, on a montré que

\begin{center}
\shadowbox{
$\forall((a_k)_{1\leqslant k\leqslant n},(b_k)_{1\leqslant k\leqslant n})\in(\Rr^n)^2$, $\sum_{k=1}^{n}|a_kb_k|\leqslant\left(\sum_{k=1}^{n}|a_k|^p\right)^{1/p}\left(\sum_{k=1}^{n}|b_k|^q\right)^{1/q}$  (Inégalité de \textsc{Hölder}).
}
\end{center}

\textbf{Remarque.} Quand $p=q=2$, on a bien $ \frac{1}{p}+ \frac{1}{q}=1$ et l'inégalité de \textsc{Hölder} s'écrit 

\begin{center}
$\sum_{k=1}^{n}|a_kb_k|\leqslant\left(\sum_{k=1}^{n}|a_k|^2\right)^{1/2}\left(\sum_{k=1}^{n}|b_k|^2\right)^{1/2}$ (inégalité de \textsc{Cauchy}-\textsc{Schwarz}).
\end{center}

  \item Soit $((a_k)_{1\leqslant k\leqslant n},(b_k)_{1\leqslant k\leqslant n})\in(\Rr^n)^2$. D'après l'inégalité de \textsc{Hölder}, on a

\begin{align*}
\sum_{k=1}^{n}(|a_k|+|b_k|)^p&=\sum_{k=1}^{n}|a_k|(|a_k|+|b_k|)^{p-1}+\sum_{k=1}^{n}|b_k|(|a_k|+|b_k|)^{p-1}\\
 &\Leftrightarrow\left(\sum_{k=1}^{n}|a_k|^p\right)^{1/p}\left(\sum_{k=1}^{n}(|a_k|+|b_k|)^{(p-1)q}\right)^{1/q}\left(\sum_{k=1}^{n}|b_k|^p\right)^{1/p}\left(\sum_{k=1}^{n}(|a_k|+|b_k|)^{(p-1)q}\right)^{1/q}\\
 &=\left(\left(\sum_{k=1}^{n}|a_k|^p\right)^{1/p}+\left(\sum_{k=1}^{n}|b_k|^p\right)^{1/p}\right)\left(\sum_{k=1}^{n}(|a_k|+|b_k|)^{p}\right)^{1-\frac{1}{p}}.
\end{align*}

Si  $\sum_{k=1}^{n}(|a_k|+|b_k|)^p = 0$, tous les $a_k$ et les $b_k$ sont nuls et l'inégalité est claire.

Sinon $\sum_{k=1}^{n}(|a_k|+|b_k|)^p > 0$ et après simplification des deux membres de l'inégalité précédente par le réel strictement positif $\sum{k=1}^{n}(|a_k|+|b_k|)^p$,  on obtient $\left(\sum_{k=1}^{n}|a_k+b_k|^p\right)^{1/p}\leqslant\left(\sum_{k=1}^{n}|a_k|^p\right)^{1/p}+\left(\sum_{k=1}^{n}|b_k|^p\right)^{1/p}$

\begin{center}
\shadowbox{
$\forall((a_k)_{1\leqslant k\leqslant n},(b_k)_{1\leqslant k\leqslant n})\in(\Rr^n)^2$, $\left(\sum_{k=1}^{n}|a_k+b_k|^p\right)^{1/p}\leqslant\left(\sum_{k=1}^{n}|a_k|^p\right)^{1/p}+\left(\sum_{k=1}^{n}|b_k|^p\right)^{1/p}$  (Inégalité de \textsc{Minkowski}).
}
\end{center}
  \end{enumerate}

\item  
  \begin{enumerate}
  \item On sait déjà que $N_1$ est une norme sur $\Rr^n$. Soit $\alpha> 1$.

(1) $N_\alpha$ est bien une application de $\Rr^n$ dans $\Rr^+$.

(2) Soit $x=(x_k)_{1\leqslant k\leqslant n}\in\Rr^n$. $N_\alpha(x) = 0\Rightarrow \forall k\in\llbracket1,n\rrbracket,\;|x_k| = 0\Rightarrow x = 0$.

(3) Soient $\lambda\in\Rr$ et $x=(x_k)_{1\leqslant k\leqslant n}\in\Rr^n$. $N_\alpha(\lambda x)= \left(\sum_{k=1}^{n}|\lambda x_k|^\alpha\right)^{1/\alpha}=\left(|\lambda|^\alpha\right)^{1/\alpha} N_\alpha(x)=|\lambda|N_\alpha(x)$.

(4) L'inégalité triangulaire est l'inégalité de \textsc{Minkowski}.

\begin{center}
\shadowbox{
$\forall \alpha\in\Rr^n$, $N_\alpha$ est une norme sur $\Rr^n$.
}
\end{center}

  \item Quelques \og boules unités \fg~dans $\Rr^2$.

$$\includegraphics{../images/img005840-1}$$



\textbf{Remarque.} Toute boule unité est symétrique par rapport à $O$ puisque $\forall x\in E$, $N(x) = N(-x)$ et donc

\begin{center}
$\forall x\in E$, $N(x)\leqslant1\Leftrightarrow N(-x)\leqslant1$.
\end{center}

 \item Soient $\alpha> 0$ et $x\in E$. On a

\begin{center}
$N_\infty(x)\leqslant N_\alpha(x)\leqslant n^{1/\alpha}N_\infty(x)$,
\end{center}

et le théorème des gendarmes fournit $\lim_{\alpha \rightarrow +\infty} N_\alpha(x) = N_\infty(x)$.

\begin{center}
\shadowbox{
$\forall x\in E$, $\lim_{\alpha \rightarrow +\infty} N_\alpha(x) = N_\infty(x)$.
}
\end{center}

 \item Soient $\alpha\in]0,1[$ puis $B=\{x\in\Rr^n/\;N_\alpha(x)\leqslant1\}$.
Les vecteurs $x=(1,0,0,\ldots,0)$ et $y=(0,1,0,\ldots,0)$ sont des éléments de $B$. Le milieu du segment $[xy]$ est $z= \frac{1}{2}(1,1,0,\ldots,0)$.

\begin{center}
$N_\alpha(z)= \frac{1}{2}(1^\alpha+1^\alpha)^{1/\alpha}=2^{\frac{1}{\alpha}-1}> 1$ car  $ \frac{1}{\alpha}-1 > 0$
\end{center}

et donc $z\notin B$. Ainsi, $B$ n'est pas convexe et donc $N_\alpha$ n'est pas une norme d'après l'exercice \ref{ex:rou1bis}.

On peut remarquer que pour $n = 1$, les $N_\alpha$ coïncident toutes avec la valeur absolue.
  \end{enumerate}
\end{enumerate}
\fincorrection
\correction{005841}
\textbullet~Il est connu que $N$ est une norme sur $E$.

\textbullet~Montrons que $N'$ est une norme sur $E$.

(1) $N'$ est une application de $E$ dans $\Rr^+$ car pour $f$ dans $E$, $f'$ est continue sur le segment $[0,1]$ et donc $f'$ est intégrable

sur le segment $[0,1]$.

(2) Soit $f\in E$. Si $N'(f) = 0$ alors  $f(0)=0\;\text{et}\;f'= 0$ (fonction continue positive d'intégrale nulle). Par suite, $f$ est un

polynôme de degré inférieur ou égal à $0$ tel que $f(0)=0$ et on en déduit que $f = 0$.

(3) $\forall f\in E$, $\forall \lambda\in\Rr$, $N'(\lambda f) =|\lambda f(0)|+\int_{0}^{1}|\lambda f'(t)|\;dt=|\lambda|\left(|f(0)|+\int_{0}^{1}|f'(t)|\;dt\right)=|\lambda| N'(f)$.

(4) Soit $(f,g)\in E^2$.

\begin{center}
$N'(f+g)\leqslant|f(0)|+|g(0)|+\int_{0}^{1}|f'(t)|\;dt +\int_{0}^{1}|g'(t)|dt= N'(f)+N'(g)$.
\end{center}

Donc $N'$ est une norme sur $E$.

\textbullet~Montrons que $N''$ est une norme sur $E$. On note que $\forall f\in E$, $N''(f)=|f(0)|+N'(f')$ et tout est immédiat.

\begin{center}
\shadowbox{
$N$, $N'$ et $N''$ sont des normes sur $E$.
}
\end{center}

\textbullet~Soit $f\in E$ et $t\in[0,1]$. Puisque la fonction $f'$ est continue sur $[0,1]$

\begin{center}
$|f(t)| = |f(0) +\int_{0}^{t}f'(u)\;du|\leqslant|f(0)|+\int_{0}^{t}|f'(u)|du\leqslant |f(0)|+\int_{0}^{1}|f'(u)|\;du =N'(f)$,
\end{center}

et donc $N(f) =\int_{0}^{1}|f(t)|\;dt\leqslant\int_{0}^{1}N'(f)\;dt=N'(f)$.

Ensuite en appliquant le résultat précédent à $f'$,  on obtient 

\begin{center}
$N'(f)=|f(0)|+N(f')\leqslant|f(0)|+N'(f')=N''(f)$.
\end{center}

Finalement 

\begin{center}
\shadowbox{
$\forall f\in E$, $N(f)\leqslant N'(f)\leqslant N''(f)$.
}
\end{center}

Pour $n\in\Nn$ et $t\in[0,1]$, on pose $f_n(t) = t^n$. 

$N(f_n)=\int_{0}^{1}t^n\;dt= \frac{1}{n+1}$ et donc la suite $(f_n)_{n\in\Nn}$ tend vers $0$ dans l'espace vectoriel normé $(E,N)$.

Par contre, pour $n\geqslant1$, $N'(f_n)=n\int_{0}^{1}t^{n-1}\;dt=1$ et la suite $(f_n)_{n\in\Nn}$  ne tend pas vers $0$ dans l'espace vectoriel normé $(E,N')$. On en déduit que

\begin{center}
\shadowbox{
les normes $N$ et $N'$ ne sont pas des normes équivalentes.
}
\end{center}

De même en utilisant $f_n(t)= \frac{t^n}{n}$, on montre que les normes $N'$ et $N''$ ne sont pas équivalentes.
\fincorrection
\correction{005847}
\begin{enumerate}
 \item  Soit $x\in E$. $\left\{\|x-a\|,\;a\in A\right\}$ est une partie non vide et minorée (par $0$) de $\Rr$. $\left\{\|x-a\|,\;a\in A\right\}$ admet donc une borne inférieure dans $\Rr$. On en déduit l'existence de $d_A(x)$.

\item 
 \begin{enumerate}
 \item Soit $A$ une partie fermée et non vide de $E$. Soit $x\in E$.

\textbullet~Supposons que $x\in A$. Alors $0\leqslant f(x)=\text{Inf}\{\|x-a\|,\;a\in A\}\leqslant\|x-x\|=0$ et donc $d_A(x) = 0$.

\textbullet~Supposons que $d_A(x)=0$. Par définition d'une borne inférieure, $\forall\varepsilon>0$ $\exists a_\varepsilon\in A/\;\|x-a_\varepsilon\| <\varepsilon$.

Soit $V$ un voisinage de $x$. $V$ contient une boule ouverte de centre $x$ et de rayon $\varepsilon>0$ puis d'après ce qui précède, $V$ contient un élément de $A$. Finalement, $\forall V\in\mathcal{V}(x)$, $V\cap A\neq\varnothing$ et donc $x\in\overline{A}=A$.

\begin{center}
\shadowbox{
Si $A$ est fermée, $\forall x\in E,\;(d_A(x) = 0\Leftrightarrow x\in A)$.
}
\end{center}

 \item Posons $d=d_A(x)$. Pour chaque entier naturel $n$, il existe $a_n\in A$ tel que $d\leqslant \|x-a_n\|\leqslant d+ \frac{1}{n}$.

La suite $(a_n)_{n\in\Nn}$ est bornée. En effet, $\forall n\in\Nn^*$ $\|a_n\|\leqslant \|a_n-x\|+\|x\|\leqslant d+ \frac{1}{n}+\| x\|\leqslant d+\|x\|+1$.

Puisque $E$ est de dimension finie, d'après le théorème de \textsc{Bolzano}-\textsc{Weierstrass},  on peut extraire de la suite $(a_n)_{n\geqslant1}$ une suite $(a_{\varphi(n)})_{n\geqslant1}$ convergeant vers un certain élément $a$ de $E$.

Ensuite, puisque $A$ est fermée, on en déduit que $a\in A$. Puis, comme

\begin{center}
$\forall n\in\Nn^*$, $d\leqslant\|x-a_{\varphi(n)}\| \leqslant d+ \frac{1}{\varphi(n)}$,
\end{center}

et puisque $\varphi(n)$ tend vers l'infini quand $n$ tend vers $+\infty$, on obtient quand $n$ tend vers l'infini, $d =\lim_{n \rightarrow +\infty}\|x-a_{\varphi(n)}\|$. Maintenant on sait que l'application $y\mapsto\|y\|$ est continue sur l'espace normé $(E,\|\;\|)$ et donc 

\begin{center}
$\lim_{n \rightarrow +\infty}\left\|x-a_{\varphi(n)}\right\|=\left\|x-\lim_{n \rightarrow +\infty}a_{\varphi(n)}\right\|=\|x-a\|$.
\end{center}

On a montré qu'il existe $a\in A$ tel que $d_A(x)=\|x-a\|$.
  \end{enumerate}

\item  Soit $x\in E$.

Puisque $A\subset\overline{A}$, $d_{\overline{A}}(x)$ est un minorant de $\{\|x-a\|,\;a\in A\}$. Comme $d_A(x)$ est le plus grand des minorants de $\{\|x-a\|,\;a\in A\}$, on a donc $d_{\overline{A}}(x)\leqslant d_A(x)$.

Soit alors $\varepsilon>0$. Il existe $y\in\overline{A}$ tel que $\|x-y\|<d(x,\overline{A})+ \frac{\varepsilon}{2}$ et puis il existe $a\in A$ tel que $\||y-a\| < \frac{\varepsilon}{2}$. On en déduit que

\begin{center}
$d_A(x)\leqslant \|x-a\|\leqslant\|x-y\|+\|y-a\|<d_{\overline{A}}(x)+ \frac{\varepsilon}{2}+ \frac{\varepsilon}{2}=d_{\overline{A}}(x)+\varepsilon$.
\end{center}

Ainsi, $\forall\varepsilon>0$, $d_A(x)<d_{\overline{A}}(x)+\varepsilon$. Quand $\varepsilon$ tend vers $0$, on obtient $d_A(x)\leqslant d_{\overline{A}}(x)$.

Finalement 

\begin{center}
\shadowbox{
$\forall x\in E$, $d_A(x)=d_{\overline{A}}(x)$.
}
\end{center}

\item  Montrons que l'application $d_A$ est Lipschitzienne. Soit $(x,y)\in E^2$

Soit $a\in A$. $d_A(x)\leqslant\|x-a\|\leqslant\|x-y\|+\|y-a\|$. Donc, $\forall a\in A$, $d_A(x) -\|x-y\|\leqslant \|y-a\| $ ou encore $d_A(x) -\|x-y\|$ est un minorant de $\{\|y-a\|,\;a\in A\}$. Puisque $d_A(y)$ est le plus grand des minorants de $\{\|y-a\|,\;a\in A\}$, on a donc $d_A(x) -\|x-y\|\leqslant d_A(y)$.

En résumé, $\forall(x,y)\in E^2$, $d_A(x) - d_A(y)\leqslant\|x-y\|$. En échangeant les rôles de $x$ et $y$, on obtient $\forall(x,y)\in E^2$, $d_A(y) - d_A(x)\leqslant\|x-y\|$ et finalement

\begin{center}
$\forall(x,y)\in E^2$, $\left|d_A(x) - d_A(y)\right|\leqslant\|x-y\|$.
\end{center}

Ainsi l'application $\begin{array}[t]{cccc}d_A~:&(E,\|\;\|)&\rightarrow&(\Rr,|\;|)\\
 &x&\mapsto&d_A(x)
\end{array}$  est $1$-Lipschitzienne et en particulier $d_A$ est continue sur l'espace  vectoriel normé $(E,\|\;\|)$.

\item  Soient $A$ et $B$ deux parties fermées et non vides de $E$ telles que $d_A = d_B$.

Soit $a\in A$. $d_B(a) = d_A(a) = 0$ (d'après 2)) et donc $a\in B$ (d'après 2)). Ainsi $A\subset B$ puis, par symétrie des rôles, $B\subset A$ et finalement $A=B$.

\item  ($A$ n'est pas un sous espace vectoriel de $E$.)

Soit $f\in A$. $1\leqslant\int_{0}^{1}f(t)\;dt\leqslant\int_{0}^{1}|f(t)|\;dt\leqslant\|f\|_\infty$. Ainsi, $\forall f\in A$, $\|f\|_\infty\geqslant1$ et donc $d_A(0)\geqslant1$.

Pour $n\in\Nn^*$ et $x\in[0,1]$, on pose $f_n(x)=\left\{
\begin{array}{l}
(n+1)x\;\text{si}\;x\in\left[0, \frac{1}{n}\right]\\
1+ \frac{1}{n}\;x\in\left[ \frac{1}{n},1\right]
\end{array}
\right.$.

$$\includegraphics{../images/img005847-1}$$


Pour chaque entier naturel non nul $n$, la fonction $f_n$ est continue sur $[0,1]$ et  

\begin{center}
$\int_{0}^{1}f_n(x)\;dx = \frac{1}{2n}\left(1+ \frac{1}{n}\right)+\left(1- \frac{1}{n}\right)\left(1+ \frac{1}{n}\right)=1+ \frac{1}{2n}- \frac{1}{2n^2}\geqslant1$.
\end{center}

Donc, la suite $(f_n)_{n\geqslant1}$ est une suite d'éléments de $A$. On en déduit que $\forall n\in\Nn^*$, $d_A(0)\leqslant\|f_n\|_\infty= 1+ \frac{1}{n}$. 

En résumé, $\forall n\in\Nn^*$, $1\leqslant d_A(0)\leqslant 1+ \frac{1}{n}$ et finalement 

\begin{center}
\shadowbox{
$d_A(0)=1$.
}
\end{center}

\textbf{Remarque.} $A$ est fermée mais la distance à $A$ n'est malgré tout pas atteinte. En effet

\textbullet~Soit $(f_n)_{n\in\Nn}$ est une suite d'éléments de $A$ convergeant dans l'espace vectoriel normé $(E,\|\;\|_\infty)$ vers un certain élément $f$ de $E$.
La suite de fonctions $(f_n)_{n\in\Nn}$ converge uniformément vers $f$ sur $[0,1]$ et donc d'une part, $f(0)=\lim_{n \rightarrow +\infty} f_n(0) = 0$ et d'autre part $\int_{0}^{1}f(x)\;dx=\int_{0}^{1}\lim_{n \rightarrow +\infty}f_n(x)\;dx=\lim_{n \rightarrow +\infty}\int_{0}^{1}f_n(x)\;dx\geqslant1$. Donc $f\in A$ et on a montré que $A$ est fermée.
 

\textbullet~Supposons qu'il existe $f\in A$ telle que $\|f\|_\infty= 1$.
Alors l'encadrement $1\leqslant\int_{0}^{1}f(x)\;dx\leqslant\|f\|_\infty= 1$ fournit $\int_{0}^{1}f(x)\;dx =\|f\|_\infty= 1$ puis $\int_{0}^{1}(\|f\|_\infty-f(x))\;dx = 0$ et donc $\|f\|_\infty- f = 0$ (fonction continue positive d'intégrale nulle) ou encore $f = 1$ ce qui contredit $f(0) = 0$. On ne peut donc pas trouver $f\in A$ tel que $d_A(0) =d(0,f)$.
\end{enumerate}
\fincorrection
\correction{004841}
Soient $F_1$, $F_2$ ferm{\'e}s non vides disjoints tels que
$F_1 \cup F_2 = A$ : Alors $\text{Fr}(A) = \text{Fr}(F_1) \cup \text{Fr}(F_2)$.
\fincorrection
\nocorrection
\correction{004843}
Soient $F_1$, $F_2$ ferm{\'e}s non vides disjoints tels que
$F_1 \cup F_2 = \{\text{va de } u_n\}$. Il existe $\varepsilon > 0$ tel que
$d(F_1,F_2) > \varepsilon$. Alors, {\`a} partir d'un certain rang,
tous les termes de la suite sont dans un seul des $F_i$.
\fincorrection
\nocorrection
\nocorrection
\nocorrection
\nocorrection
\nocorrection
\correction{004848}
Soit $r = \lim r_n$ :
$\|a_n - a_{n+k}\| \le r_n - r_{n+k}$ donc la suite $(a_n)$ est de Cauchy,
et converge vers $a$.

On a $\|a_n - a\| \le r_n - r$ donc $B(a,r) \subset B_n$.

R{\'e}ciproquement, si $x \in \bigcap_n B_n$, alors $\|x-a_n\| \le r_n$ donc
$\|x-a\| \le r$.
\fincorrection
\nocorrection
\correction{004850}
Soit $a \in \mathring F$ et $B(a,r) \subset \bigcup_n F_n$ :
$B \setminus F_1$ est un ouvert non vide donc contient une boule
$B_1(a_1,r_1)$.
De m{\^e}me, $B_1 \setminus F_2$ contient une boule $B_2(a_2,r_2)$ etc.
On peut imposer $r_n \xrightarrow[n\to\infty]{} 0$, donc il existe $c \in \bigcap_n B_n$,
c.a.d. $c\in B$ mais pour tout $n$, $c \notin F_n$. Contradiction.
\fincorrection
\correction{004851}
Soit $(a_n)$ d{\'e}finie par $a_{n+1} = f(a_n)$ : les sous-suites
$(a_{2n})$ et $(a_{2n+1})$ convergent vers le point fixe de $f\circ f$.
\fincorrection
\correction{004852}
Supposons qu'il existe une famille $(\mathcal{C}_i = \mathcal{C}(a_i,R_i))_{i\in I}$ de cercles disjoints
dont la r{\'e}union est {\'e}gale au plan $P$. On note $D_i$ le disque ferm{\'e} de fronti{\`e}re $\mathcal{C}_i$.
Soit $i_0\in I$ choisi arbitrairement, $i_1$ tel que $a_{i_0}\in\mathcal{C}_{i_1}$,
$i_2$ tel que $a_{i_1}\in\mathcal{C}_{i_2}$ etc. On a $R_{i_k} < \frac12R_{i_{k-1}}$
donc la suite $(D_{i_k})$ v{\'e}rifie le th{\'e}or{\`e}me des ferm{\'e}s emboit{\'e}s, l'intersection
des $D_{i_k}$ est r{\'e}duite {\`a} un point $x$ par lequel ne passe aucun cercle $\mathcal{C}_j$.
\fincorrection
\correction{004853}
\begin{enumerate}
  \item 
  \item demi-cercle unit{\'e} $ \Rightarrow  x = 0$, $y = \frac 2\pi$.
  \item Sommes de Riemann + l'enveloppe convexe d'un compact est compacte.
  \item $\vec N\,'(t) = \vec\sigma(\vec M\,'(t))  \Rightarrow 
             \|\vec N\,'(t)\| = \|\vec M\,'(t))\|$.\par
             $ \int_{t=a}^b \vec{GN_t}\|\vec N\,'(t)\|\,dt = \vec 0
             =  \int_{t=a}^b \vec{\sigma(G)N_t}\|\vec M\,'(t)\|\,dt$,
             donc $\sigma(G) = G$.
  \item 
\end{enumerate}
 \fincorrection
\correction{004854}
\begin{enumerate}
  \item 
  \item 
  \item $\vec e_i\,'' = \vec\Omega\,'\wedge\vec e_i
             + (\vec\Omega|\vec e_i)\vec\Omega - \|\vec\Omega\|^2\vec e_i$.
\end{enumerate}
 \fincorrection
\nocorrection
\nocorrection
\correction{004857}
\begin{enumerate}
  \item Non : $f(t) = (t,t^2)$, $g(t) = (1,t)$.
  \item Non : $f(t) = \left(\frac{t^2}2, \frac{t^3}3\right)$,
                   $g(t) = \left(t,\frac{t^2}2\right)$.
\end{enumerate}
 \fincorrection
\nocorrection
\correction{005854}
\begin{enumerate}
 \item  \textbullet~Soit $P\in E$. Si on pose $P=\sum_{k=0}^{+\infty}a_kX^k$, il existe $n\in\Nn$ tel que $\forall k>n$, $a_k=0$. Donc $\|P\|_\infty=\text{Sup}\left\{\left| \frac{P^{(k)}(0)}{k!}\right|,\;k\in\Nn\right\}=\text{Max}\{|a_k|,\;0\leqslant k\leqslant n\}$ existe dans $\Rr$.

\textbullet~$\forall P\in E$, $\|P\|_\infty\geqslant0$.

\textbullet~Soit $P\in E$. $\|P\|_\infty=0\Rightarrow\forall k\in\Nn,\;|a_k|\leqslant0\Rightarrow\forall k\in\Nn,\;a_k=0\Rightarrow P=0$.

\textbullet~Soient $P\in E$ et $\lambda\in\Rr$. $\|\lambda P\|_\infty=\text{Max}\{|\lambda a_k|,\;0\leqslant k\leqslant n\}=|\lambda|\text{Max}\{|a_k|,\;0\leqslant k\leqslant n\}=|\lambda|\|P\|_\infty$.

\textbullet~Soient $P=\sum_{k\geqslant0}^{}a_kX^k$ et $Q=\sum_{k\geqslant0}^{}b_kX^k$ deux polynômes. Pour $k\in\Nn$, $|a_k+b_k|\leqslant|a_k|+|b_k|\leqslant\|P\|_\infty+\|Q\|_\infty$ et donc $\|P+Q\|_\infty\leqslant\|P\|_\infty+\|Q\|_\infty$.

\begin{center}
\shadowbox{
$\|\;\|_\infty$ est une norme sur $E$.
}
\end{center}

\item  $\forall P\in E$, $\|f(P)\|_\infty=\|P\|_\infty$ et donc $\forall P\in E\setminus\{0\}$, $ \frac{\|f(P)\|_\infty}{\|P\|_\infty}=1$. On en déduit que $\text{Sup}\left\{ \frac{\|f(P)\|_\infty}{\|P\|_\infty},\;P\in E\setminus\{0\}\right\}=1$. Ceci montre tout à la fois que $f$ est continue sur $(E,\|\;\|_\infty)$ et $|||f|||=1$.

\begin{center}
\shadowbox{
$f$ est continue sur $(E,\|\;\|_\infty)$ et $|||f|||=1$.
}
\end{center}
\end{enumerate}
\fincorrection
\correction{005855}
(La linéarité de $\Delta$ est claire et de plus $\Delta$ est un endomorphisme de $E$ car si $u$ est une suite bornée, $\Delta(u)$ l'est encore. Plus précisément,)

\begin{center}
$\forall u\in E$, $\forall n\in\Nn$, $\left|\Delta(u)_n\right|\leqslant|u_n|+|u_{n+1}|\leqslant 2\|u\|_\infty$ et donc $\forall u\in E$, $\|\Delta(u)\|_\infty\leqslant 2\|u\|_\infty$.
\end{center}

Ceci montre que $\Delta$ est continu sur $E$ et $|||\Delta|||\leqslant2$. Ensuite, si $u$ est la suite définie par $\forall n\in\Nn$, $u_n = (-1)^n$ alors $u$ est un élément non nul de $E$ tel que $\|u\|_\infty=1$ et $\|\Delta(u)\|_\infty= 2$. En résumé,

\textbullet~$\forall u\in E\setminus\{0\}$, $ \frac{\|\Delta(u)\|_\infty}{\|u\|_\infty}\leqslant2$,

\textbullet~$\exists u\in E\setminus\{0\}$, $ \frac{\|\Delta(u)\|_\infty}{\|u\|_\infty}=2$.

On en déduit que

\begin{center}
\shadowbox{
$\Delta$ est continu sur $(E,\|\;\|_\infty)$ et $|||\Delta|||=2$.
}
\end{center}

(La linéarité de $C$ est claire et $C$ est un endomorphisme de $E$ car si $u$ est bornée, $C(u)$ l'est encore. Plus précisément,)

\begin{center}
$\forall u\in E$, $\forall n\in\Nn$, $\left|(C(u))_n\right|\leqslant \frac{1}{n+1}\sum_{k=0}^{n}\|u\|_\infty =\|u\|_\infty$ et donc $\forall u\in E$, $\|C(u)\|_\infty\leqslant\|u\|_\infty$.
\end{center}

Par suite $T$ est continue sur $E$ et $|||T|||\leqslant1$. Ensuite, si $u$ est la suite définie par $\forall n\in\Nn$, $u_n = 1$ alors $u$ est un élément non nul de $E$ tel que $\|u\|_\infty=1$ et $\|C(u)\|_\infty=1$. En résumé,

\textbullet~$\forall u\in E\setminus\{0\}$, $ \frac{\|C(u)\|_\infty}{\|u\|_\infty}\leqslant1$,

\textbullet~$\exists u\in E\setminus\{0\}$, $ \frac{\|C(u)\|_\infty}{\|u\|_\infty}=1$.

On en déduit que

\begin{center}
\shadowbox{
$C$ est continu sur $(E,\|\;\|_\infty)$ et $|||C|||=1$.
}
\end{center}
\fincorrection
\correction{005856}
\begin{enumerate}
 \item  Soit $f\in E$.

\begin{align*}\ensuremath
\|Tf\|_1&=\int_{0}^{1}|Tf(x)|\;dx=\int_{0}^{1}\left|\int_{0}^{x}f(t)\;dt\right|dx\\
 &\leqslant\int_{0}^{1}\left(\int_{0}^{x}|f(t)|\;dt\right)dx\\
 &\leqslant\int_{0}^{1}\left(\int_{0}^{1}|f(t)|\;dt\right)dx=\int_{0}^{1}\|f\|_1\;dx=\|f\|_1.
\end{align*}

Ceci montre que $\forall f\in E\setminus\{0\}$, $ \frac{\|Tf\|_1}{\|f\|_1}\leqslant1$. Ceci montre que $T$ est continu sur $(E,\|\;\|_1)$ et que $|||T|||\leqslant1$.

Pour $n\in\Nn$ et $x\in[0,1]$, posons $f_n(x)=(1-x)^n$. Pour $n\in\Nn$,

\begin{center}
$\|f_n\|_1=\int_{0}^{1}(1-x)^n\;dx=\left[- \frac{(1-x)^{n+1}}{n+1}\right]_0^1= \frac{1}{n+1}$,
\end{center}

puis pour $x\in[0,1]$, $Tf_n(x)=\int_{0}^{x}(1-t)^n\;dt= \frac{1}{n+1}(1-(1-x)^{n+1})$ et donc

\begin{center}
$\|Tf_n\|_1=\int_{0}^{1}|Tf_n(x)|\;dx= \frac{1}{n+1}\int_{0}^{1}(1-(1-x)^{n+1})\;dx= \frac{1}{n+1}\left(1- \frac{1}{n+2}\right)= \frac{1}{n+2}$.
\end{center}

On en déduit que $\forall n\in\Nn$, $|||T|||\geqslant \frac{\|Tf_n\|_1}{\|f_n\|_1}= \frac{n+1}{n+2}$.

En résumé, $\forall n\in\Nn$, $ \frac{n+1}{n+2}\leqslant|||T|||\leqslant 1$ et donc $|||T|||=1$.

\begin{center}
\shadowbox{
$T$ est continu sur $(E,\|\;\|_1)$ et $|||T|||=1$.
}
\end{center}
 
 
 \item  Supposons qu'il existe $f\in E\setminus\{0\}$ tel que $\|Tf\|_1=\|f\|_1$. On en déduit que chaque inégalité écrite au début de la question 1) est une égalité et en particulier $\int_{0}^{1}\left(\int_{0}^{x}|f(t)|\;dt\right)dx=\int_{0}^{1}\left(\int_{0}^{1}|f(t)|\;dt\right)dx$ ou encore $\int_{0}^{1}\left(\int_{0}^{1}|f(t)|\;dt-\int_{0}^{x}|f(t)|\;dt\right)dx=0$. Par suite, $\forall x\in[0,1]$, $\int_{0}^{1}|f(t)|\;dt-\int_{0}^{x}|f(t)|\;dt=0$ (fonction continue, positive, d'intégrale nulle) puis en dérivant la dernière inégalité, $\forall x\in[0,1]$, $|f(x)|=0$ et finalement $f=0$. Ceci est une contradiction et donc $|||T|||$ n'est pas atteinte.
\end{enumerate}
\fincorrection
\correction{005857}
L'application $f$ est linéaire de $(E,N)$ dans $(\Rr,|\;|)$. Soit $A=(a_{i,j})_{1\leqslant i,j\leqslant n}\in E$.

\begin{align*}\ensuremath
|f(A)|&=|\text{Tr}(A)|\leqslant\sum_{i=1}^{n}|a_{i,i}|\\
 &\leqslant\sum_{i=1}^{n}\left(\sum_{j=1}^{n}|a_{i,j}|\right)\leqslant\sum_{i=1}^{n}N(A)=nN(A).
\end{align*}
 

Ceci montre déjà que $f$ est continue sur $(E,N)$ et que $|||f|||\leqslant n$. De plus, si $A=I_n\neq0$, $ \frac{|f(A)|}{N(A)}= \frac{n}{1}=n$. Donc

\begin{center}
\shadowbox{
$f$ est continue sur $(E,N)$ et $|||f|||=n$.
}
\end{center}
\fincorrection
\correction{005858}
\textbullet~$\forall A=(a_{i,j})_{1\leqslant i,j\leqslant n}\in\mathcal{M}_n(\Rr)$, $\|A\|_\infty=\text{Max}\{|a_{i,j}|,\;1\leqslant i,j\leqslant n\}$.

Soient $A=(a_{i,j})_{1\leqslant i,j\leqslant n}$ et $B=(b_{i,j})_{1\leqslant i,j\leqslant n}$. Posons $AB=(c_{i,j})_{1\leqslant i,j\leqslant n}$ où $\forall(i,j)\in\llbracket1,n\rrbracket^2$, $c_{i,j}=\sum_{k=1}^{n}a_{i,k}b_{k,j}$.

Pour $(i,j)\in\llbracket1,n\rrbracket^2$,

\begin{center}
$|c_{i,j}|\leqslant\sum_{k=1}^{n}|a_{i,k}||b_{k,j}|\leqslant\sum_{k=1}^{n}\|A\|_\infty\|B\|_\infty=n\|A\|_\infty\|B\|_\infty$,
\end{center}

et donc, $\|AB\|_\infty\leqslant n\|A\|_\infty\|B\|_\infty$. Ainsi, $\forall(A,B)\in(\mathcal{M}_n(\Cc)\setminus\{0\})^2$, $ \frac{\|AB\|_\infty}{\|A\|_\infty\|B\|_\infty}\leqslant n$.

De plus, pour $A_0=B_0=(1)_{1\leqslant i,j\leqslant n}\neq0$, $\|A_0\|_\infty=\|B_0\|_\infty=1$ puis $\|A_0B_0\|_\infty=\|nA_0\|_\infty=n$ et donc $ \frac{\|A_0B_0\|_\infty}{\|A_0\|_\infty\|B_0\|_\infty}=n$. Ceci montre que

\begin{center}
\shadowbox{
$\text{Sup}\left\{ \frac{\|AB\|_\infty}{\|A\|_\infty\|B\|_\infty},\;(A,B)\in(\mathcal{M}_n(\Cc)\setminus\{0\})^2\right\}=n$.
}
\end{center}

En particulier, $\|\;\|_\infty$ n'est pas une norme sous-multiplicative.

\textbullet~$\forall A=(a_{i,j})_{1\leqslant i,j\leqslant n}\in\mathcal{M}_n(\Rr)$, $\|A\|_1=\sum_{1\leqslant i,j\leqslant n}^{}|a_{i,j}|$. Avec les notations précédentes,

\begin{align*}\ensuremath
\|AB\|_1&=\sum_{1\leqslant i,j\leqslant n}^{}|c_{i,j}|=\sum_{1\leqslant i,j\leqslant n}^{}\left|\sum_{k=1}^{n}a_{i,k}b_{k,j}\right|\\
 &\leqslant\sum_{1\leqslant i,j\leqslant n}^{}\left(\sum_{k=1}^{n}|a_{i,k}||b_{k,j}|\right)=\sum_{1\leqslant i,j,k\leqslant n}^{}|a_{i,k}||b_{k,j}|\\
 &\sum_{1\leqslant i,j,k,l\leqslant n}^{}|a_{i,j}||b_{k,l}|=\|A\|_1\|B\|_1.
\end{align*}

Donc $\forall(A,B)\in\left(\mathcal{M}_n(\Rr)\setminus\{0\}\right)^2$, $ \frac{\|AB\|_1}{\|A\|_1\|B\|_1}\leqslant 1$.

De plus, pour $A_0=B_0=E_{1,1}$, on a $A_0B_=E_{1,1}$ et donc $ \frac{\|A_0B_0\|_1}{\|A_0\|_1\|B_0\|_1}=1$. Ceci montre que

\begin{center}
\shadowbox{
$\text{Sup}\left\{ \frac{\|AB\|_1}{\|A\|_1\|B\|_1},\;(A,B)\in(\mathcal{M}_n(\Cc)\setminus\{0\})^2\right\}=1$.
}
\end{center}

En particulier, $\|\;\|_1$ est une norme sous-multiplicative.

\textbullet~$\forall A=(a_{i,j})_{1\leqslant i,j\leqslant n}\in\mathcal{M}_n(\Rr)$, $\|A\|_2=\sqrt{\sum_{1\leqslant i,j\leqslant n}^{}a_{i,j}^2}$. Avec les notations précédentes,

\begin{align*}\ensuremath
\|AB\|_2^2&=\sum_{1\leqslant i,j\leqslant n}^{}c_{i,j}^2=\sum_{1\leqslant i,j\leqslant n}^{}\left(\sum_{k=1}^{n}a_{i,k}b_{k,j}\right)^2\\
 &\leqslant\sum_{1\leqslant i,j\leqslant n}^{}\left(\sum_{k=1}^{n}a_{i,k}^2\right)\left(\sum_{k=1}^{n}b_{k,j}^2\right)\;(\text{inégalité de \textsc{Cauchy}-\textsc{Schwarz}})\\
 &=\sum_{1\leqslant i,j\leqslant n}^{}\left(\sum_{k=1}^{n}a_{i,k}^2\right)\left(\sum_{l=1}^{n}b_{l,j}^2\right)=\sum_{1\leqslant i,j,k,l\leqslant n}^{}a_{i,k}^2b_{l,j}^2=\left(\sum_{1\leqslant i,k\leqslant n}^{}a_{i,k}^2\right)\left(\sum_{1\leqslant j,l\leqslant n}^{}b_{l,j}^2\right)=\|A\|_2\|B\|_2
\end{align*}

Donc $\forall(A,B)\in\left(\mathcal{M}_n(\Rr)\setminus\{0\}\right)^2$, $ \frac{\|AB\|_2}{\|A\|_2\|B\|_2}\leqslant 1$.

De plus, pour $A_0=B_0=E_{1,1}$, on a $A_0B_=E_{1,1}$ et donc $ \frac{\|A_0B_0\|_2}{\|A_0\|_2\|B_0\|_2}=1$. Ceci montre que

\begin{center}
\shadowbox{
$\text{Sup}\left\{ \frac{\|AB\|_2}{\|A\|_2\|B\|_2},\;(A,B)\in(\mathcal{M}_n(\Cc)\setminus\{0\})^2\right\}=1$
}
\end{center}

En particulier, $\|\;\|_2$ est une norme sous-multiplicative.
\fincorrection
\correction{005859}
Une \og norme trois barres \fg~sur $\mathcal{M}_n(\Rr)$ est nécessairement sous-multiplicative. L'exercice précédent montre qu'il existe des normes sur $\mathcal{M}_n(\Rr)$ qui ne sont pas sous-multiplicatives (par exemple $\|\;\|_\infty$). Donc une norme sur $\mathcal{M}_n(\Rr)$ n'est pas nécessairement une \og norme trois barres \fg.
\fincorrection
\correction{005860}
Soit $N$ une norme sur $\mathcal{M}_n(\Rr)$. D'après l'exercice \ref{ex:rou5}, $\|\;\|_1$ est une norme sous-multiplicative.

Puisque $\mathcal{M}_n(\Rr)$ est un espace vectoriel de dimension finie sur $\Rr$, $N$ et $\|\;\|_1$ sont des normes équivalentes. Par suite, il existe deux réels strictement positifs $\alpha$ et $\beta$ tels que $\alpha\|\;\|_1\leqslant N\leqslant\beta\|\;\|_1$.

Pour $(A,B)\in\left(\mathcal{M}_n(\Rr)\right)^2$,

\begin{center}
$N(AB)\leqslant\beta\|AB\|_1\leqslant\beta\|A\|_1\|B\|_1\leqslant \frac{\beta}{\alpha^2}N(A)N(B)$
\end{center}

et le réel $k= \frac{\beta}{\alpha^2}$ est un réel strictement positif tel que $\forall(A,B)\in(\mathcal{M}_n(\Rr))^2$, $N(AB)\leqslant kN(A)N(B)$.

\textbf{Remarque.} Le résultat précédent signifie que $N'= \frac{1}{K}N$ est une norme sous-multiplicative car pour $(A,B)\in\left(\mathcal{M}_n(\Rr)\right)^2$,

\begin{center}
$N'(AB)= \frac{1}{k^2}N(AB)\leqslant \frac{1}{k^2}N(A)N(B)= \frac{1}{k}N(A) \frac{1}{k}N(B)=N'(A)N'(B)$.
\end{center}
\fincorrection
\correction{005861}
Non, car si $A=E_{1,1}\neq0$ et $B=E_{2,2}\neq0$ alors $AB=0$ puis $N(AB)<N(A)N(B)$.
\fincorrection
\correction{005862}
\textbullet~Pour $\|\;\|_1$. Soient $A=(a_{i,j})_{1\leqslant i,j\leqslant n}\in\mathcal{M}_n(\Rr)$ puis $X=(x_i)_{1\leqslant i\leqslant n}\mathcal{M}_{n,1}(\Rr)$.

\begin{align*}\ensuremath
\|AX\|_1&=\sum_{i=1}^{n}\left|\sum_{j=1}^{n}a_{i,j}x_j\right|\\
 &\leqslant\sum_{i=1}^{n}\left(\sum_{j=1}^{n}|a_{i,j}||x_j|\right)=\sum_{j=1}^{n}|x_j|\left(\sum_{i=1}^{n}|a_{i,j}|\right)\\
 &\leqslant\left(\sum_{j=1}^{n}|x_j|\right)\text{Max}\left\{\sum_{i=1}^{n}|a_{i,j}|,\;1\leqslant j\leqslant n\right\}=\text{Max}\{\|C_j\|_1,\;1\leqslant j\leqslant n\}\times\|X\|_1,
\end{align*}

en notant $C_1$,\ldots, $C_n$ les colonnes de la matrice $A$. Donc, $\forall A\in\mathcal{M}_n(\Rr)$, $|||A|||_1\leqslant\text{Max}\{\|C_j\|_1,\;1\leqslant j\leqslant n\}$.

Soit alors $j_0\in\llbracket1,n\rrbracket$ tel que $\|C_{j_0}\|_1=\text{Max}\{\|C_j\|_1,\;1\leqslant j\leqslant n\}$. On note $X_0$ le vecteur colonne dont toutes les composantes sont nulles sauf la $j_0$-ème qui est égale à $1$. $X_0$ est un vecteur non nul tel que

\begin{center}
$\|AX_0\|_1=\sum_{i=1}^{n}|a_{i,j_0}|=\text{Max}\left\{\|C_j\|_1,\;1\leqslant j\leqslant n\right\}\times\|X_0\|_1$.
\end{center}

En résumé,

\textbf{(1)} $\forall X\in\mathcal{M}_{n,1}(\Rr)\setminus\{0\}$, $ \frac{\|AX\|_1}{\|X\|_1}\leqslant\text{Max}\left\{\|C_j\|_1,\;1\leqslant j\leqslant n\right\}$,

\textbf{(2)} $\exists X_0\in\mathcal{M}_{n,1}(\Rr)\setminus\{0\}$, $ \frac{\|AX_0\|_1}{\|X_0\|_1}=\text{Max}\left\{\|C_j\|_1,\;1\leqslant j\leqslant n\right\}$.

On en déduit que $\forall A\in\mathcal{M}_n(\Rr)$, $|||A|||_1=\text{Max}\left\{\|C_j\|_1,\;1\leqslant j\leqslant n\right\}$.

\textbullet~Pour $\|\;\|_\infty$. Soient $A=(a_{i,j})_{1\leqslant i,j\leqslant n}\in\mathcal{M}_n(\Rr)$ puis $X=(x_i)_{1\leqslant i\leqslant n}\mathcal{M}_{n,1}(\Rr)$. Pour $i\in\llbracket1,n\rrbracket$,

\begin{align*}\ensuremath
|(AX)_i|&=\left|\sum_{j=1}^{n}a_{i,j}x_j\right|\leqslant\sum_{j=1}^{n}|a_{i,j}||x_j|\leqslant\left(\sum_{j=1}^{n}|a_{i,j}|\right)\|X\|_\infty\\
 &\leqslant\text{Max}\left\{\sum_{j=1}^{n}|a_{i,j}|,\;1\leqslant i\leqslant n\right\}\|X\|_\infty=\text{Max}\{\|L_k\|_1,\;1\leqslant k\leqslant n\}\times\|X\|_\infty,
\end{align*}

en notant $L_1$,\ldots, $L_n$ les lignes de la matrice $A$. Donc, $\forall A\in\mathcal{M}_n(\Rr)$, $|||A|||_\infty\leqslant\text{Max}\{\|L_i\|_1,\;1\leqslant i\leqslant n\}$.

Soit alors $i_0\in\llbracket1,n\rrbracket$ tel que $\|L_{i_0}\|_1=\text{Max}\{\|L_i\|_1,\;1\leqslant i\leqslant n\}$. On pose $X_{0}=\left(\varepsilon_i\right)_{1\leqslant i\leqslant n}$ où $\forall j\in\llbracket 1,n\rrbracket$, $\varepsilon_j$ est un élément de $\{-1,1\}$ tel que $a_{i_0,j}=\varepsilon_j|a_{i_0,j}|$ (par exemple, $\varepsilon_j= \frac{a_{i_0,j}}{|a_{i_0,j}|}$ si $a_{i_0,j}\neq0$ et $\varepsilon_j=1$ si $a_{i_0,j}=1$).

\begin{align*}\ensuremath
\|AX_0\|_\infty&=\text{Max}\left\{\left|\sum_{j=1}^{n}a_{i,j}\varepsilon_j\right|,\;1\leqslant i\leqslant n\right\}\\
 &\geqslant\left|\sum_{j=1}^{n}a_{i_0,j}\varepsilon_j\right|=\sum_{j=1}^{n}|a_{i_0,j}|=\|L_{i_0}\|_1
=\text{Max}\{\|L_i\|_1,\;1\leqslant i\leqslant n\}\times\|X_0\|_\infty.
\end{align*}

En résumé,

\textbf{(1)} $\forall X\in\mathcal{M}_{n,1}(\Rr)\setminus\{0\}$, $ \frac{\|AX\|_\infty}{\|X\|_\infty}\leqslant\text{Max}\left\{\|L_i\|_1,\;1\leqslant i\leqslant n\right\}$,

\textbf{(2)} $\exists X_0\in\mathcal{M}_{n,1}(\Rr)\setminus\{0\}$, $ \frac{\|AX_0\|_\infty}{\|X_0\|_\infty}\geqslant\text{Max}\left\{\|L_i\|_1,\;1\leqslant i\leqslant n\right\}$.

On en déduit que $\forall A\in\mathcal{M}_n(\Rr)$, $|||A|||_\infty=\text{Max}\left\{\|L_i\|_1,\;1\leqslant j\leqslant n\right\}$.

Ainsi, en notant $C_1$,\ldots, $C_n$ et $L_1$,\ldots, $L_n$ respectivement les colonnes et les lignes d'une matrice $A$,

\begin{center}
\shadowbox{
$\forall A\in\mathcal{M}_n(\Rr)$, $|||A|||_1=\text{Max}\{\|C_j\|_1,\;1\leqslant j\leqslant n\}$ et $|||A|||_\infty=\text{Max}\{\|L_i\|_1,\;1\leqslant i\leqslant n\}$.
}
\end{center}
\fincorrection
\correction{005863}
Soit $D=\text{diag}(\lambda_i)_{1\leqslant i\leqslant n}\in\mathcal{D}_n(\Rr)$. Pour $X=(x_i)_{1\leqslant i\leqslant n}\in\mathcal{M}_{n,1}(\Rr)$,

\begin{center}
$\|DX\|_2=\sqrt{\sum_{i=1}^{n}\lambda_i^2x_i^2}\leqslant\sqrt{(\rho(D))^2\sum_{i=1}^{n}x_i^2}=\rho(D)\|X\|_2$,
\end{center}

De plus, si $\lambda$ est une valeur propre de $D$ telle que $|\lambda|=\rho(D)$ et $X_0$ est un vecteur propre associé, alors 

\begin{center}
$\|DX_0\|_2=\|\lambda X_0\|_2=|\lambda|\|X_0\|_2=\rho(D)\|X_0\|_2$.
\end{center}

En résumé

\textbf{(1)} $\forall X\in\mathcal{M}_{n,1}(\Rr)\setminus\{0\}$, $ \frac{\|DX\|_2}{\|X\|_2}\leqslant\rho(D)$,

\textbf{(2)} $\exists X_0\in\mathcal{M}_{n,1}(\Rr)\setminus\{0\}$, $ \frac{\|DX_0\|_2}{\|X_0\|_2}=\rho(D)$.

On en déduit que $\forall D\in\mathcal{D}_n(\Rr)$, $|||D|||_2=\rho(D)$.

Soit alors $A\in\mathcal{S}_n(\Rr)$. D'après le théorème spectral, il existe $P\in O_n(\Rr)$ et $D=\text{diag}(\lambda_i)_{1\leqslant i\leqslant n}\in\mathcal{D}_n(\Rr)$ tel que $A=PD{^t}P$. De plus $\rho(A)=\rho(D)$. Pour $X\in\mathcal{M}_{n,1}(\Rr)$,

\begin{align*}\ensuremath
\|AX\|_2&=\|PD{^t}PX\|_2\\
 &=\|D({^t}PX)\|_2\;(\text{car}\;P\in O_n(\Rr)\Rightarrow\forall Y\in\mathcal{M}_{n,1}(\Rr),\;\|PY\|_2=\|Y\|_2)\\
 &=\|DX'\|_2\;\text{où on a posé}\;X'={^t}PX.
\end{align*}

Maintenant l'application $X\mapsto{^t}PX=X'$ est une permutation de $\mathcal{M}_{n,1}(\Rr)$ car la matrice ${^t}P$ est inversible et donc $X$ décrit $\mathcal{M}_{n,1}(\Rr)$ si et seulement si $X'$ décrit $\mathcal{M}_{n,1}(\Rr)$. De plus, pour tout vecteur colonne $X$, $\|X'\|_2=\|{^t}PX\|_2=\|X\|_2$. On en déduit que $\left\{
 \frac{\|AX\|_2}{\|X\|_2},\;X\in\mathcal{M}_{n,1}(\Rr)\setminus\{0\}\right\}=\left\{
 \frac{\|DX'\|_2}{\|X'\|_2},\;X'\in\mathcal{M}_{n,1}(\Rr)\setminus\{0\}\right\}$ et en particulier,

\begin{center}
$|||A|||_2=|||D|||_2=\rho(D)=\rho(A)$.
\end{center}

\begin{center}
\shadowbox{
$\forall A\in\mathcal{S}_n(\Rr)$, $|||A|||_2=\text{Sup}\left\{ \frac{\|AX\|_2}{\|X\|_2},\;X\in\mathcal{M}_{n,1}(\Rr)\setminus\{0\}\right\}=\rho(A)$.
}
\end{center}

\textbf{Remarque.} L'application $A\mapsto\rho(A)$ est donc une norme sur $\mathcal{S}_n(\Rr)$ et de plus cette norme est sous-multiplicative.
\fincorrection
\nocorrection
\nocorrection
\nocorrection
\nocorrection
\nocorrection
\nocorrection
\nocorrection
\nocorrection
\nocorrection
\correction{005842}
\begin{enumerate}
 \item  Soit $\begin{array}[t]{cccc}
d~:&\mathcal{M}_n(\Rr)&\rightarrow&\Rr\\
 &M&\mapsto&\text{det}(M)
\end{array}
$. On sait que l'application $d$ est continue sur $\mathcal{M}_n(\Rr)$ (muni de n'importe quelle norme) et que $\Rr^*$ est un ouvert de $\Rr$ en tant que réunion de deux intervalles ouverts.

Par suite, $GL_n(\Rr)=d^{-1}(\Rr^*)$ est un ouvert de $\mathcal{M}_n(\Rr)$ en tant qu'image réciproque d'un ouvert par une application continue.

Soit $A\in \mathcal{M}_n(\Rr)$. Le polynôme$\text{det}(A-xI)$ n'a qu'un nombre fini de racines (éventuellement nul) donc pour $p$ entier naturel supérieur ou égal à un certain $p_0$, $\text{det}\left(A- \frac{1}{p}I\right)\neq 0$. La suite $\left(A- \frac{1}{p}I\right)_{p\geqslant p_0}$ est une suite d'éléments de $GL_n(\Rr)$ convergente de limite $A$. Ceci montre que l'adhérence de $GL_n(\Rr)$ est $\mathcal{M}_n(\Rr)$ ou encore $GL_n(\Rr)$ est dense dans $\mathcal{M}_n(\Rr)$.

\begin{center}
\shadowbox{
$GL_n(\Rr)$ est un ouvert de $\mathcal{M}_n(\Rr)$, dense dans $\mathcal{M}_n(\Rr)$.
}
\end{center}

\item  $\mathcal{M}_n(\Rr)\setminus GL_n(\Rr)$ est fermé en tant que complémentaire d'un ouvert.

Soit $n\geqslant2$. Les matrices $A_p = pE_{1,1}$, $p\in\Nn$, sont non inversibles et la suite $(A_p)_{p\in\Nn}$ est non bornée. Par suite $Mn(R)\setminus GLn(R)$  est non borné et donc non compact.

\begin{center}
\shadowbox{
$\forall n\geqslant2$, $M_n(\Rr)\setminus GL_n(\Rr)$  est fermé mais non compact.
}
\end{center}

\item  \textbullet~Montrons que $O_n(\Rr)$ est fermé. Posons $\begin{array}[t]{cccc}
g~:&\mathcal{M}_n(\Rr)&\rightarrow&(\mathcal{M}_n(\Rr))^2\\
 &M&\mapsto&(M,{^t}M)
\end{array}$, $\begin{array}[t]{cccc}h :&(\mathcal{M}_n(\Rr))^2&\rightarrow&\mathcal{M}_n(\Rr)\\
 &(M,N)&\mapsto MN\end{array}$puis $\begin{array}[t]{cccc}f~:&\mathcal{M}_n(\Rr)&\rightarrow&\mathcal{M}_n(R)\\
  &M&\mapsto&M{^t}M\end{array}$.
		
		
$g$ est continue sur $\mathcal{M}_n(\Rr)$ car linéaire sur un espace de dimension finie. $h$ est continue sur $(M_n(\Rr))^2$ car bilinéaire sur un espace de dimension finie. On en déduit que
$f = h\circ g$ est  continue sur $\mathcal{M}_n(\Rr)$.

Enfin $O_n(\Rr) =f^{-1}({I_n})$ est fermé en tant qu'image réciproque d'un fermé par une application continue.

\textbullet~Montrons que $O_n(\Rr)$ est borné. $\forall A\in O_n(\Rr)$, $\forall(i,j)\in\llbracket1,n\rrbracket^2$, $|a_{i,j}|\leqslant 1$ et donc $\forall A\in O_n(\Rr)$, $\|A\|_{\infty}\leqslant1$.

D'après le théorème de \textsc{Borel}-\text{Lebesgue}, puisque $O_n(\Rr)$ est un fermé borné de l'espace de dimension finie $\mathcal{M}_n(\Rr)$, $O_n(\Rr)$ est un compact de $\mathcal{M}_n(\Rr)$.

$O_n(\Rr)$ n'est pas convexe. En effet, les deux matrices $I_n$ et $-I_n$ sont orthogonales mais le milieu du segment joignant ces deux matrices est $0$ qui n'est pas une matrice orthogonale.

\begin{center}
\shadowbox{
$O_n(\Rr)$  est compact mais non convexe.
}
\end{center}

\item  $S_n(\Rr)$ est un sous espace vectoriel de l'espace de dimension finie $\mathcal{M}_n(\Rr)$ et est donc un fermé de $\mathcal{M}_n(\Rr)$.

\begin{center}
\shadowbox{
$S_n(\Rr)$  est fermé.
}
\end{center}

\item  Soit $A\in\mathcal{M}_n(\Rr)$ et $p$ un élément fixé de $\llbracket1,n-1\rrbracket$ (le résultat est clair si $p=0$ ou $p=n$).

$A$ est de rang inférieur ou égal à $p$ si et seulement si tous ses mineurs de format $p+1$ sont nuls (hors programme).

Soient $I$ et $J$ deux sous-ensembles donnés de $\llbracket1,n\rrbracket$ de cardinal $p+1$ et $A_{I,J}$ la matrice extraite de $A$ de format $p+1$ dont les numéros de lignes sont dans $I$ et les numéros de colonnes sont dans $J$.

Pour $I$ et $J$ donnés, l'application $A\mapsto A_{I,J}$ est continue car linéaire de $\mathcal{M}_n(\Rr)$ dans $\mathcal{M}_{p+1}(\Rr)$. Par suite, l'application 

$f_{I,J}~:~A\mapsto\text{det}(A_{I,J})$ est continue sur $\mathcal{M}_n(\Rr)$. L'ensemble des matrices $A$ telles que $\text{det}(A_{I,J}) = 0$ est donc un fermé de $\mathcal{M}_n(\Rr)$ (image réciproque du fermé $\{0\}$ de $\Rr$ par l'application continue $f_{I,J}$) et l'ensemble des matrices de rang inférieur ou égal à $p$ est un fermé de $\mathcal{M}_n(\Rr)$ en tant qu'intersection de fermés.

\item  Soit $A\in\mathcal{M}_n(\Cc)$. Posons $\text{Sp}(A)=(\lambda_i)_{1\leqslant i\leqslant n}$. On sait que toute matrice est triangulable dans $\Cc$ et donc il existe $P\in GL_n(\Cc)$ et $T\in\mathcal{T}_n(\Cc)$ avec $\forall i\in\llbracket1,n\rrbracket$, $t_{i,i}=\lambda_i$ telle que $A=PTP^{-1}$.

On munit dorénavant $\mathcal{M}_n(\Cc)$ d'une norme multiplicative notée $\|\;\|$. 
Puisque toutes les normes sont équivalentes en dimension finie, il existe un réel strictement positif $K$ telle que pour toute matrice $M$, $\|M\|\leqslant K\|M\|_\infty$.

Soit $\varepsilon> 0$. Il existe un $n$-uplet de réels $(\lambda_1,...,\lambda_n)$ tels que $\forall k\in\llbracket1,n\rrbracket$, $0\leqslant\varepsilon_k < \frac{\varepsilon}{K\|P\|\|P^{-1}\|}$ et les $\lambda_k+\varepsilon_k$ sont deux à deux distincts. (On prend $\varepsilon_1 = 0$ puis $\varepsilon_2$ dans $\left[0, \frac{\varepsilon}{K\|P\|\|P^{-1}\|}\right[$ tel que $\lambda_2+\varepsilon_2 \neq \lambda_1+\varepsilon_1$ ce qui est possible puisque $\left[0, \frac{\varepsilon}{K\|P\|\|P^{-1}\|}\right[$ est infini puis $\varepsilon_3$ dans $\left[0, \frac{\varepsilon}{K\|P\|\|P^{-1}\|}\right[$ tel que $\lambda_3+\varepsilon_3$ soit différent de $\lambda_1+\varepsilon_1$ et $\lambda_2+\varepsilon_2$ ce qui est possible puisque $\left[0, \frac{\varepsilon}{K\|P\|\|P^{-1}\|}\right[$ est infini ...)

On pose $D=\text{diag}(\lambda_i)_{1\leqslant i\leqslant n}$ puis $T' = T+D$ et enfin $A'=PT'P^{-1}$. Tout d'abord les valeurs propres de $A'$ sont deux à deux distinctes (ce sont les $\lambda_i+\varepsilon_i$, $1\leqslant i\leqslant n$) et donc $A'$ est diagonalisable. Ensuite 

\begin{center}
$\|A'-A\| =\|PDP^{-1}\|\leqslant\|P\|\|D\|\|P^{-1}\|\leqslant K\|P\|\|P^{-1}\|\|D\|_\infty <\varepsilon.$
\end{center}

En résumé, $\forall A\in\mathcal{M}_n(\Cc)$, $\forall\varepsilon>0$, $\exists A'\in\mathcal{M}_n(\Cc)/\;\|A'-A\|<\varepsilon$ et $A'$ diagonalisable. On a montré que

\begin{center}
\shadowbox{
L'ensemble des matrices complexes diagonalisables dans $\Cc$ est dense dans $\mathcal{M}_n(\Cc)$.
}
\end{center}

On ne peut remplacer $\mathcal{M}_n(\Cc)$ par $\mathcal{M}_n(\Rr)$.

Soient $A=\left(
\begin{array}{cc}
0&-1\\
1&0
\end{array}
\right)$ et $E=\left(
\begin{array}{cc}
a&c\\
b&d
\end{array}
\right)\in\mathcal{M}_2(\Rr)$.

\begin{center}
$\chi_{A+E}=\left|
\begin{array}{cc}
a-X&c-1\\
b+1&d-X
\end{array}
\right|=X^2 - (a+d)X+(ad-bc) +(b-c)+ 1$.
\end{center}

Le discriminant  de $\chi_{A+E}$ est $\Delta= (a+d)^2 - 4(ad-bc) - 4(b-c) -4$. Supposons de plus que $\|E\|_\infty\leqslant \frac{1}{4}$. Alors 

\begin{center}
$\Delta= (a+d)^2 - 4(ad-bc) - 4(b-c) -4\leqslant \frac{1}{4}+ 4\left( \frac{1}{16}+ \frac{1}{16}\right)+4\left( \frac{1}{4}+ \frac{1}{4}\right)-4=- \frac{5}{4}<0$.
\end{center}

Par suite, aucune des matrices $A+E$ avec $\|E\|_\infty\leqslant \frac{1}{4}$ n'a de valeurs propres réelles et donc aucun donc diagonalisable dans $\Rr$. On a montré que l'ensemble des matrices réelles diagonalisables dans $\Rr$ n'est pas dense dans $\mathcal{M}_n(\Rr)$.

\item  La matrice de la forme quadratique $Q~:~(x,y)\mapsto ax^2+2bxy+cy^2$ dans la base canonique est $\left(
\begin{array}{cc}
a&b\\
b&c
\end{array}
\right)$.

Les valeurs propres de cette matrice sont strictement positives si et seulement si $a+c>0$ et $ac -b^2>0$. L'application $(a,b,c)\mapsto a+c$  est continue sur $\Rr^3$ car linéaire sur $\Rr^3$ qui est de dimension finie et l'application $(a,b,c)\mapsto ac -b^2$ est continue sur $\Rr^3$ en tant que polynôme.

L'ensemble des triplets considéré est l'intersection des images réciproques par ces applications de l'ouvert $]0,+\infty[$ de $\Rr$ et est donc un ouvert de $\Rr^3$.

\item  Notons $\mathcal{S}$ l'ensemble des matrices stochastiques.

\textbullet~Vérifions que $\mathcal{S}$ est borné. Soit $A=(a_{i,j})_{1\leqslant i,j\leqslant n}\in\mathcal{S}$. $\forall(i,j)\in\llbracket1,n\rrbracket^2$, $0\leqslant a_{i,j}\leqslant1$ et donc $\|A\|_\infty\leqslant1$. Ainsi, $\forall A/in\mathcal{S}$, $\|A\|_\infty\leqslant 1$ et donc $\mathcal{S}$ est borné.

\textbullet~Vérifions que $\mathcal{S}$ est fermé. 

Soit $(i,j)\in\llbracket1,n\rrbracket^2$. L' application $f_{i,j}~:~A\mapsto a_{i,j}$ est continue sur $\mathcal{M}_n(\Rr)$ à valeurs dans $\Rr$ car linéaire sur $\mathcal{M}_n(\Rr)$ qui est de dimension finie. $[0,+\infty[$ est un fermé de $\Rr$ car son complémentaire $]-\infty,0[$ est un ouvert de $\Rr$. Par suite, $\left\{A=(a_{k,l})_{1\leqslant k,l\leqslant n}/\;a_{i,j}\geqslant0\right\}=f_{i,j}^{-1}([0,+\infty[)$ est un fermé de $\mathcal{M}_n(\Rr)$ en tant qu'image réciproque d'un fermé par une application continue.

Soit $i\in\llbracket1,n\rrbracket$. L' application $g_i~:~A\mapsto \sum_{j=1}^{n}a_{i,j}$ est continue sur $\mathcal{M}_n(\Rr)$ à valeurs dans $\Rr$ car linéaire sur $\mathcal{M}_n(\Rr)$ qui est de dimension finie. Le singleton $\{1\}$ est un fermé de $\Rr$. Par suite, $\left\{A=(a_{k,l})_{1\leqslant k,l\leqslant n}/\;\sum_{j=1}^{n}a_{i,j}=1\right\}=g_i^{-1}(\{1\})$ est un fermé de $\mathcal{M}_n(\Rr)$ en tant qu'image réciproque d'un fermé par une application continue.

$\mathcal{S}$ est donc un fermé de $\mathcal{M}_n(\Rr)$ en tant qu'intersection de fermé de $\mathcal{M}_n(\Rr)$.

En résumé, $\mathcal{S}$ est un fermé borné de l'espace $\mathcal{M}_n(\Rr)$ qui est de dimension finie et donc $\mathcal{S}$ est un compact de $\mathcal{M}_n(\Rr)$ d'après le théorème de \textsc{Borel}-\textsc{Lebesgue}.

\textbullet~Vérifions que $\mathcal{S}$ est convexe. Soient $(A,B)\in(\mathcal{S})^2$ et $\lambda\in[0,1]$. D'une part, $\forall(i,j)\in\llbracket1,n\rrbracket^2$, $(1-\lambda)a_{i,j}+\lambda b_{i,j}\geqslant0$ et d'autre part, pour $i\in\llbracket1,n\rrbracket$

\begin{center}
$\sum_{j=1}^{n}((1-\lambda)a_{i,j}+\lambda b_{i,j})=(1-\lambda)\sum_{j=1}^{n}a_{i,j}+\lambda\sum_{j=1}^{n}b_{i,j}=(1-\lambda)+\lambda=1$,
\end{center}

ce qui montre que $(1-\lambda)A+\lambda B\in\mathcal{S}$. On a montré que $\forall(A,B)\in\mathcal{S}^2$, $\forall \lambda\in[0,1]$, $(1-\lambda)A+\lambda B\in\mathcal{S}$ et donc $\mathcal{S}$ est convexe.

\begin{center}
\shadowbox{
l'ensemble des matrices stochastiques est un compact convexe de $\mathcal{M}_n(\Rr)$.
}
\end{center}

\item  Soient $A$ et $B$ deux matrices réelles diagonalisables. Soient $\begin{array}[t]{cccc}
\gamma_1~:&[0,1]&\rightarrow&\mathcal{M}_n(\Rr)\\
 &t&\mapsto&(1-t).A+t.0=(1-t)A
\end{array}
$ et 

$\begin{array}[t]{cccc}
\gamma_2~:&[0,1]&\rightarrow&\mathcal{M}_n(\Rr)\\
 &t&\mapsto &tB
 \end{array}$. Soit enfin $\begin{array}[t]{cccc}
\gamma~:&[0,1]&\rightarrow&\mathcal{M}_n(\Rr)\\
 &t&\mapsto &\left\{
 \begin{array}{l}
 \rule[-4mm]{0mm}{10mm}\gamma_1(2t)\;\text{si}\;t\in\left[0, \frac{1}{2}\right]\\
 \gamma_2(2t-1)\;\text{si}\;t\in\left[ \frac{1}{2},1\right]
 \end{array}
 \right.
 \end{array}$.
 
 
$\gamma_1$ est un chemin continu joignant la matrice $A$ à la matrice nulle et $\gamma_2$ est un chemin continu joignant la matrice nulle à la matrice $B$. Donc $\gamma$ est un chemin continu joignant la matrice $A$ à la matrice $B$. De plus, pour tout réel $t\in[0,1]$, la matrice $\gamma_1(t)=(1-t)A$ est diagonalisable (par exemple, si $A=P\text{diag}(\lambda_i)_{1\leqslant i\leqslant n}P^{-1}$ alors $(1-t)A=P\text{diag}((1-t)\lambda_i)_{1\leqslant i\leqslant n}P^{-1}$) et de même, pour tout réel $t\in[0,1]$, la matrice $\gamma_2(t)=tB$ est diagonalisable. Finalement $\gamma$ est un chemin continu joignant les deux matrices $A$ et $B$ diagonalisables dans $\Rr$, contenu dans l'ensemble des matrices diagonalisables dans $\Rr$. On a montré que

\begin{center}
\shadowbox{
l'ensemble des matrices diagonalisables dans $\Rr$ est connexe par arcs.
}
\end{center}
\end{enumerate}
\fincorrection
\correction{001956}
\begin{enumerate}
  \item 
  \begin{enumerate}
    \item Un vecteur directeur est $\overrightarrow{AB}$ dont les coordonnées sont $(x_B-x_A,y_B-y_A)=(-3,1)$.
Pour n'importe quel vecteur directeur $\vec v = (x_v,y_v)$ la pente est le réel $p=\frac{y_v}{x_v}$.
La pente est indépendante du choix du vecteur directeur. On trouve ici $p=-\frac13$.
Une équation paramétrique de la droite de vecteur directeur $\vec v$ passant par $A=(x_A,y_A)$ est donnée par
$\left\{ \begin{array}{l}
x = x_v t + x_A \\ y = y_v t+y_A \\ 
\end{array} \right..$
Donc ici pour le vecteur directeur $\overrightarrow{AB}$ on trouve l'équation paramétrique 
$\left\{ \begin{array}{l}
x = -3t + 2 \\ y = t+3 \\ 
\end{array} \right.$

Il y a plusieurs façons d'obtenir une équation cartésienne $ax+by+c=0$.


\textbf{Première méthode.} 
On sait que $A=(x_A,y_A)$ appartient à la droite donc ses coordonnées vérifient l'équation 
$ax_A+by_A+c=0$, idem avec $B$.
On en déduit le système 
$\left\{ \begin{array}{l}
2a+3b+c = 0 \\  -a+4b+c=0 \\ 
\end{array} \right..$
Les solutions s'obtiennent à une constante multiplicative près,
on peut fixer $a=1$ et on trouve alors $b=3$ et $c=-11$.
L'équation est donc $x+3y-11=0$.

    \item On trouve $\vec v = \overrightarrow{AB}=(5,-3)$, $p=-\frac35$
et $\left\{ \begin{array}{l}
x = 5t -7 \\ y = -3t-2 \\ 
\end{array} \right.$

\textbf{Deuxième méthode.} 
Pour trouver l'équation cartésienne on part de l'équation
paramétrique réécrite ainsi
$\left\{ \begin{array}{l}
\frac{x+7}{5} = t \\ -\frac{y+2}{3} = t \\ 
\end{array} \right.$
On en déduit $\frac{x+7}{5}=-\frac{y+2}{3}$ ; d'où
l'équation $3x+5y+31=0$.

    \item On trouve $\vec v = \overrightarrow{AB}=(0,3)$, la droite est donc verticale (sa pente est infinie)
une équation paramétrique est
$\left\{ \begin{array}{l}
x = 3 \\ y = 3t+6 \\ 
\end{array} \right.$.
Une équation cartésienne est simplement $(x=3)$.
  \end{enumerate}


  \item
  \begin{enumerate}
    \item Equation paramétrique
$\left\{ \begin{array}{l}
x = -3t+2 \\ y = -t+1 \\ 
\end{array} \right.$

\textbf{Troisième méthode.} 
Pour une droite d'équation cartésienne $ax+by+c=0$, on sait que $\vec n = (a,b)$ est un vecteur normal à la droite et donc 
$\vec v = (-b,a)$ est un vecteur directeur (car alors $\vec v \cdot \vec n =0$).
Réciproquement si $\vec  v = (-b,a)$ est un vecteur directeur alors une équation 
est de la forme $ax+by+c=0$ pour une certaine constante $c$ à déterminer.

Ici on nous donne le vecteur directeur $\vec v= (-3,-1)$ donc on cherche une équation
sous la forme $-x+3y+c=0$.
Pour trouver $c$, on utilise que $A$ appartient à la droite donc
$-x_A+3y_A+c=0$, ce qui conduit à $c=-1$.
Ainsi une équation de la droite est $-x+3y=1$.


    \item On trouve $2x-y+1=0$.

    \item Droite horizontale d'équation $(y=1)$.
  \end{enumerate}
  \item Voici juste les résultats :
  \begin{enumerate}
    \item $y=3x+4$,
    \item $y=-3$,
    \item $8x+4y=4$ (les droites parallèles à $8x+4y=3$ sont de la forme $8x+4y=c$).
  \end{enumerate}
\end{enumerate}



\fincorrection
\nocorrection
\nocorrection
\nocorrection
\correction{001960}
\begin{enumerate}
  \item Le point $A$ est l'intersection des droites $(AB)$ et $(AC)$.
Les coordonnées $(x,y)$ de $A$ sont donc solutions du système :
$\left\{ \begin{array}{l}
x+2y=3 \\ x+y=2 \\ 
\end{array} \right.$ donné par 
les équations des deux droites.
La seule solution est $(x,y)=(1,1)$.
On a donc $A=(1,1)$. On fait de même pour obtenir le point $B=(-1,2)$ et $C=(2,0)$.

  \item Notons $A'$ le milieu de $[BC]$ alors les coordonnées se trouvent par la formule suivante
$A'= (\frac{x_B+x_C}{2},\frac{y_B+y_C}{2}) = (\frac12,1)$.
De même on trouve $B'=(\frac32,\frac12)$ et $C'=(0,\frac32)$.

  \item 
  \begin{enumerate}
    \item Les médianes ont pour équations : 
$(AA') : (y=1)$ ; $(BB') : (3x+5y=7)$ ; $(CC') : (3x+4y=6)$.
    \item Vérifions que les trois médianes sont concourantes (ce qui est vrai quelque soit le triangle).
On calcule d'abord l'intersection $I=(AA')\cap (BB')$, les coordonnées du point $I$ d'intersection vérifient 
donc le système 
$\left\{ \begin{array}{l}
y=1 \\ 3x+5y=7 \\ 
\end{array} \right.$. On trouve $I=(\frac23,1)$.

Il ne reste plus qu'à vérifier que $I$ appartient à la droite $(CC')$ d'équation $3x+4y=6$.
En effet $3x_I+4y_I=6$ donc $I \in (CC')$.

Conclusion : les médianes sont concourantes au point $I=(\frac23,1)$.
  \end{enumerate}
\end{enumerate}

\fincorrection
\nocorrection
\nocorrection
\nocorrection
\nocorrection
\nocorrection
\nocorrection
\nocorrection
\nocorrection
\nocorrection
\nocorrection
\nocorrection
\nocorrection
\nocorrection
\nocorrection
\nocorrection
\nocorrection
\nocorrection
\nocorrection
\nocorrection
\nocorrection
\nocorrection
\nocorrection
\nocorrection
\nocorrection
\nocorrection
\nocorrection
\nocorrection
\nocorrection
\nocorrection
\nocorrection
\nocorrection
\nocorrection
\nocorrection
\nocorrection
\nocorrection
\nocorrection
\nocorrection
\nocorrection
\nocorrection
\correction{002009}
Soit $z\mapsto \alpha z + \beta$ la représentation en coordonnée complexe de la similitude directe envoyant $A$ sur $C$ et $B$ sur $D$. On a donc
\[
\left\{\begin{array}{lc}\alpha a+\beta &= c \\ \alpha b+\beta &= d\\ \end{array}\right.
\] ce qui donne $\alpha = \frac{c-d}{a-b}$ et $\beta=\frac{ad-bc}{a-b}$. D'après la condition fixée par l'énoncé, on a $\alpha\neq 1$ donc cette similitude admet un unique point fixe $\Omega$ d'affixe 
\[
\omega = \frac{\beta}{1-\alpha} = \frac{ad-bc}{a-b-c+d}.
\]
On remarque que l'expression de $\omega$ est inchangée en permutant $b$ et $c$. Cela  signifie qu'en faisant les mêmes calculs pour déterminer la représentation complexe de la similitude envoyant $A$ sur $B$ et $C$ sur $D$, on obtient le même point fixe.
\fincorrection
\nocorrection
\correction{002011}
\begin{enumerate}
\item   
  \begin{enumerate}
  \item Une équation d'un plan est $ax+by+cz+d=0$. 
Si un point appartient à un plan cela donne une condition linéaire sur $a,b,c,d$. 
Si l'on nous donne trois point cela donne un système linéaire de trois équations à trois inconnues
(car l'équation est unique à un facteur multplicatif non nul près). On trouve : 
    \begin{enumerate}
       \item $x+y+z-1=0$
       \item $3x+3y+z-7=0$
    \end{enumerate}

    \item $\vec n = \vec{u} \wedge \vec{v}$ est normal au plan. Si $\vec n=(a, b, c)$ alors une équation du plan est $ax+by+cz+d=0$.
   On trouve : 
      \begin{enumerate}
         \item $-9x+7y+12z-17=0$
         \item $17x+13y-7z-3=0$
       \end{enumerate}

 \item Trouver deux points $B, C$ de la droite $D$. Le vecteurs $\vec u = \overrightarrow{AB}$ et 
$\vec v = \overrightarrow{AC}$ sont des vecteurs directeurs de $P$. Procédé ensuite comme la question précédente.
On obtient : 
\begin{enumerate}
\item Par exemple $B=(0,-6,-3)$ et $C=(-1,0,2)$ appartiennent à $D$. On trouve l'équation
$4x-y+2z=0$. 
\item Par exemple $B=(0,-1,1)$ (pout $t=0$) et $C=(1,1,-2)$ (pour $t=1$) appartiennent à $D$. On trouve l'équation
$2x-y-1=0$. 
\end{enumerate}             

\item  Trouver un point $A$ de $D$ et deux points $B, C$ de la droite $D'$. Le vecteurs $\vec u = \overrightarrow{AB}$ et 
$\vec v = \overrightarrow{AC}$ sont des vecteurs directeurs de $P$. Puis procédé comme avant.
\end{enumerate}  

 \item  Les plans sont définis paramétriquement par
$(P)$ : $(2,2,1)+s(1,2,-1)+t(2,1,-1)$ donc
deux des vecteurs directeurs sont $\vec u =(1,2,-1)$ et $\vec v = (2,1,-1)$.
Un vecteur normal à $(P)$ est alors $\vec n = \vec{u} \wedge \vec{v}=(-1,-1,-3)$.


Pour le plan $(P')$ défini par $(1,3,1)+s'(3,3,-2)+t'(-1,1,0)$, il a 
pour vecteurs directeurs $\vec {u'} =(3,3,-2)$ et $\vec {v'} = (-1,1,0)$.
Un vecteur normal à $(P')$ est alors $\vec {n'} = \vec{u'} \wedge \vec{v'}=(2,2,6)$.

Les vecteurs normaux $\vec n$ et $\vec {n'}$ sont colinéaires donc les plans $(P)$ et $(P')$ sont parallèles
(ou confondus).

Maintenant le point $A=(2,2,1)$ appartient à $(P)$ (on a fait $s=0$ et $t=0$).
Il appartient aussi à $(P')$ (en prenant $s'=0$ et $t'=-1$).

Bilan. $(P)$ et $(P')$ sont parallèles et ont un point commun : ils sont égaux !

\end{enumerate} 

\fincorrection
\nocorrection
\nocorrection
\nocorrection
\nocorrection
\nocorrection
\nocorrection
\nocorrection
\nocorrection
\correction{002020}
\begin{enumerate}
  \item  
  \begin{enumerate}
    \item Un point $A$ appartient à un plan d'équation $ax+by+cz+d=0$ si et seulement si 
$ax_A+by_A+cz_A+d=0$.
Donc $A(1,1,1)\in P_m$ si et seulement si $ m^2+(2m-1)+m=3$.
Ce qui équivaut à $m^2+3m-4=0$. Les deux solutions sont $m=1$ et $m=-4$.
Donc $A$ appartient aux plans $P_1$ et $P_{-4}$ et pas aux autres.

    \item Un plan d'équation $ax+by+cz+d=0$ a pour vecteur normal $\vec n = (a,b,c)$.
Donc si $\vec{n}=(2,-\frac 52,-1)$ est un vecteur normal à $P_m$
une équation cartésienne est de la forme $2 x-\frac52 y-z +d=0$.
Or une équation de $P_m$ est $m^2x+(2m-1)y+mz-3=0$.
Ces deux équations sont égales à un facteur multiplicatif près $\lambda\in \Rr^*$ :
 $2 x-\frac52 y-z +d = \lambda \big(m^2x+(2m-1)y+mz-3\big)$.
On en déduit 
$2=\lambda m^2$, $-\frac52=\lambda(2m-1)$ et $-1=\lambda m$.
En divisant la première égalité par la troisième on trouve : $m=-2$.
D'où $\lambda=\frac12$. La seconde égalité est alors vérifiée.

Le seul plan ayant $\vec n$ pour vecteur normal est $P_{-2}$.


    \item Un vecteur est directeur du plan $P$ si et seulement si le produit scalaire
$\vec v\cdot \vec n=0$. Ici $\vec n = (m^2,2m-1,m)$.
Donc $\vec{v}=(1,1,1)$ est vecteur directeur si et seulement si 
$m^2+2m-1+m=0$. Ce qui équivaut à $m^2+3m-1=0$.
Les deux plans qui ont pour vecteur directeur $\vec v$ sont les plans ayant le paramètre
$m = \frac{-3\pm\sqrt{13}}{2}$.
  \end{enumerate}

  \item Nous allons prendre $3$ plans de la famille $(P_m)$, calculer leur point d'intersection
et finalement montrer que ce point appartient aux autres plans.

Prenons trois paramètre ``au hasard'' $m=0$, $m=1$, $m=-1$.
Un point qui appartient à ces trois plans doit vérifier les trois équations :
$$\left\{ \begin{array}{l}
y=-3 \\ 
x+y+z=3 \\ 
x-3y-z=3 \\ 
\end{array} \right.$$
On résout ce système pour trouver que l'intersection des trois plan $P_0$, $P_1$ et $P_{-1}$ est le point 
$Q=(0,-3,6)$.

Il ne reste plus qu'à vérifier que ce point appartient à tous les plans $P_m$: 
c'est le cas car $m^2\cdot 0 +(2m-1)\cdot (-3) +m \cdot 6 -3=0$.


\textbf{Autre méthode.}
On cherche un point $Q=(x_0,y_0,z_0)$ qui vérifie
l'égalité $m^2 x_0 +(2m-1)y_0 +m z_0 -3=0$
pour tout $m$.
En considérant que c'est une égalité polynomiale en $m$ ($x_0,y_0,z_0$ sont fixés)
on en déduit que $m^2 x_0 +(2m-1)y_0 +m z_0 -3$ est le polynôme nul : $x_0 m^2 + (2y_0+z_0)m -y_0-3=0$.
Ces coefficients sont nuls : $x_0=0$ (le coefficient de $m^2$),
$2y_0+z_0=0$ (le coefficient de $m$), $-y_0-3=0$ (le terme constant).
On trouve bien sûr le même point d'intersection de tous les plans : $Q=(0,-3,6)$.

\end{enumerate}

\fincorrection
\correction{002021}
\begin{enumerate}
  \item La distance d'un point $A=(x_0,y_0,z_0)$ à un plan $P$ d'équation 
$ax+by+cz+d=0$ est donnée par la formule :
$$d(A,P) = \frac{|ax_0+by_0+cz_0+d|}{\sqrt{a^2+b^2+c^2}}.$$
On trouve donc
  \begin{enumerate}
    \item $d(A,P)=\frac{|2\cdot 1+1\cdot 0+1\cdot 2+4|}{\sqrt{2^2+1^2+1^2}}= \frac{8}{\sqrt6}=\frac{4\sqrt6}{3}$.
    \item $d(A,P)=\frac{2}{\sqrt{42}}$.
  \end{enumerate}


  \item Trouvons d'abord une équation paramétrique de la droite $D$.
On pose par exemple $z=t$ et on exprime $x$ et $y$ en fonction de $t$.
Partant du système $\left\{ \begin{array}{l}
-2x+y-3z=1\\ x+z=1  \end{array} \right.$
on trouve $x=1-t$ et $y=3+t$.
La droite $D$ est donc l'ensemble des point $M_t=(1-t,3+t,t)$ ($t$ parcourant $\Rr$).

La distance $AM_t$ vérifie donc
$$AM_t^2 = \|\vec AM_t \|^2 = \| (1-t-1,3+t-2,t-3)\|^2 = t^2 + (t+1)^2+(t-3)^2 = 3t^2-4t+10.$$
Minimiser cette distance c'est trouver le minimum de la fonction $\delta(t)=3t^2-4t+10$.
Il est donc atteint pour $t_0$ vérifiant $\delta'(t_0)=0$, donc pour $t_0=\frac23$.
La distance entre $A$ et la droite $D$ est donc la longueur $AM_{t_0}=\sqrt{\delta(t_0)}=\sqrt{\frac{26}{3}}$.
Au passage on a obtenu la perpendiculaire à $D$ passant par $A$ c'est la droite $(AM_{t_0})$.

\bigskip


\textbf{Autre méthode.}

Il existe une formule pour calculer directement la distance.
Si $\vec v$ est un vecteur directeur de $D$  et $M_0$ un point de $D$
alors
$$d(A,D)= \frac{\|\vec v \wedge \overrightarrow{AM_0}\|}{\|\vec v\|}.$$
On a paramétré la droite $D$ par les points $M_t = (1,3,0)+t (-1,1,1)$.
Donc $M_0=(1,3,0)\in D$ et $\vec v= (-1,1,1)$ est un vecteur directeur de $D$.
On a alors $\overrightarrow{AM_0} = (0,1,-3)$ et
$\vec v \wedge \overrightarrow{AM_0} = (-4,-3,-1)$ :
on obtient : 
$$d(A,D)= \frac{\|\vec v \wedge \overrightarrow{AM_0}\|}{\|\vec v\|} = \frac{\sqrt{26}}{\sqrt{3}}.$$

\end{enumerate}

\fincorrection
\nocorrection
\nocorrection
\nocorrection
\nocorrection
\nocorrection
\nocorrection
\nocorrection
\nocorrection
\nocorrection
\nocorrection
\correction{005195}
\begin{enumerate}
\item  Soit $G$ l'isobarycentre du triangle $(ABC)$. On a donc $G=\mbox{bar}(A(1),B(1),C(1))$. Notons $A'$, $B'$ et $C'$ les milieux respectifs des côtés $[B,C]$, $[C,A]$ et $[A,B]$. D'après le théorème du barycentre partiel, 
$G=\mbox{bar}(A(1),A'(2))$. En particulier, $G$ est sur la médiane $(AA')$. De même, $G$ est sur la médiane $(BB')$ et sur la médiane $(CC')$.

Finalement, $G$ est sur les trois médianes. les trois médianes sont donc concourantes en $G$.

\item  Les droites $(BC)$ et $(CA)$ ne sont pas parallèles. Par suite, les médiatrices respectives des côtés $[B,C]$ et $[C,A]$ ne sont pas parallèles. Elles sont donc sécantes en un point que l'on note $O$. Par définition de $O$, on a $OA=OB=OC$. $O$ est donc à égale distance de $A$ et $B$ et est ainsi sur la médiatrice de $[A,B]$. Finalement, les trois médiatrices sont concourantes en $O$. De plus, $O$ étant à égale distance de $A$, $B$ et $C$, le cercle de centre $O$ et de rayon $OA$ passe par $B$ et $C$.

Réciproquement, un cercle passant par $A$, $B$ et $C$ a pour centre un point à égale distance de ces points et donc nécessairement de centre $O$ et de rayon $OA$. Ceci démontre l'existence et l'unicité du cercle circonscrit au triangle $(ABC)$~:~c'est le cercle de centre $O$ et de rayon $OA$.

\item  Les hauteurs issues de $A$ et $B$ ne sont pas parallèles (car perpendiculaires à deux droites non parallèles). Elles admettent ainsi un et un seul point d'intersection. Ceci assure l'unicité d'un point commun aux trois hauteurs.

Soit $h$ l'homthétie de centre $G$ et de rapport $-2$. Puisque $\overrightarrow{GA}=-2\overrightarrow{GA}$, on a $h(A')=A$ et de même $h(B')=B$ et $h(C')=C$.

Par $h$, l'image de la médiatrice de $[B,C]$, c'est-à-dire de la droite passant par $A'$ et perpendiculaire à $(BC)$ est la droite passant par $h(A')=A$ et perpendiculaire à $(BC)$ (car parallèle à la médiatrice de $[B,C]$). Cette droite est la hauteur issue de $A$ du triangle $(ABC)$. De même, les images des médiatrices de $[C,A]$ et $[A,B]$ sont respectivement les hauteurs issues de $B$ et $C$.

Le point $O$ est sur les trois médiatrices. Son image par $h$ est donc sur les trois hauteurs (d'où l'existence d'un point commun aux trois hauteurs). Ces trois hauteurs sont ainsi concourantes en un point noté $H$ et appelé l'orthocentre du triangle $(ABC)$. De plus, l'égalité $h(O)=H$ s'écrit $\overrightarrow{GH}=-2\overrightarrow{GO}$ ou encore $\overrightarrow{G0}+\overrightarrow{OH}=2\overrightarrow{OG}$ ou enfin,

\begin{center}
\shadowbox{$\overrightarrow{OH}=3\overrightarrow{OG}$  \textsc{Euler}.}
\end{center}

Les trois points $O$, $G$ et $H$, s'ils sont deux à deux distincts, sont en particulier alignés sur une droite appelée \textbf{droite d'}\textsc{Euler} du triangle $(ABC)$.

\item  Deux bissectrices intérieures ne sont pas parallèles (démontrez-le) et sont donc sécantes en un point $I$ à égale distance des trois côtés et à l'intérieur du triangle $(ABC)$. Ce point étant à égale distance des trois côtés est centre du cercle tangent intérieurement aux trois côtés, le cercle inscrit.

\end{enumerate}
\fincorrection
\correction{005196}
 $\;$(Notez bien l'alignement des points $G$, $H$ et $O$). 

$$\includegraphics{../images/img005196-1}$$

\begin{enumerate}
\item  On a $AB=\sqrt{3^2+1^2}=\sqrt{10}$ et $AC=\sqrt{1+2^2}=\sqrt{5}$. Par suite,

$$\cos(\widehat{BAC})=\frac{\overrightarrow{AB}.\overrightarrow{AC}}{AB.AC}=\frac{(-3)(-1)+(-1)(2)}{\sqrt{5}\sqrt{10}}=\frac{1}{5\sqrt{2}}.$$

Par suite, $\widehat{BAC}=81^\circ$ à un degré près.
\item  $\mbox{aire}(ABC)=\frac{1}{2}|\mbox{det}(\overrightarrow{AB},\overrightarrow{AC})|
=\frac{1}{2}\mbox{abs}(\left|
\begin{array}{cc}
-3&-1\\
-1&2
\end{array}
\right|)=\frac{7}{2}$.

\item  Notons $G$ l'isobarycentre du triangle $(ABC)$. $z_G=\frac{1}{3}(z_A+z_B+z_C)=\frac{1}{3}(1+2i-2+i+4i)=\frac{1}{3}(-1+7i)$, et donc \shadowbox{$G(-\frac{1}{3},\frac{7}{3}).$}

Notons $(x,y)$ les coordonnées de $\Omega$, le centre du cercle circonscrit au triangle $(ABC)$ (dans cette exercice, la lettre $O$ désigne certainement l'origine du repère).

\begin{align*}
\left\{
\begin{array}{l}
\Omega A=\Omega B\\
\Omega A=\Omega C
\end{array}
\right.&
\Rightarrow
\left\{
\begin{array}{l}
(x-1)^2+(y-2)^2=(x+2)^2+(y-1)^2\\
(x-1)^2+(y-2)^=x^2+(y-4)^2
\end{array}
\right.
\Rightarrow
\left\{
\begin{array}{l}
3x+y=0\\
2x-4y=-11
\end{array}
\right.\\
 &\Rightarrow x=-\frac{11}{14}\;\mbox{et}\;y=\frac{33}{14}\;(\mbox{d'après les formules de \textsc{Cramer}}),
\end{align*}

et donc 

\shadowbox{$\Omega(-\frac{11}{14},\frac{33}{14})$.}

Notons $(x,y)$ les coordonnées de l'orthocentre $H$ du triangle $(ABC)$.

\begin{itemize}
\item[\textbf{1ère solution.}]

\begin{align*}
\left\{
\begin{array}{l}
\overrightarrow{AH}.\overrightarrow{BC}=0\\
\overrightarrow{BH}.\overrightarrow{AC}=0
\end{array}
\right.&
\Rightarrow
\left\{
\begin{array}{l}
2(x-1)+3(y-2)=0\\
-(x+2)+2(y-1)=0
\end{array}
\right.
\Rightarrow
\left\{
\begin{array}{l}
2x+3y=8\\
-x+2y=4
\end{array}
\right.\\
 &\Rightarrow x=\frac{4}{7}\;\mbox{et}\;y=\frac{16}{7}\;(\mbox{d'après les formules de \textsc{Cramer}}),
\end{align*}

et donc, \shadowbox{$H(\frac{4}{7},\frac{16}{7})$.}

\item[\textbf{2ème solution.}] Il est bien meilleur de connaître la relation d'\textsc{Euler} $\overrightarrow{\Omega H}=3\overrightarrow{\Omega G}$ et de l'utiliser.

$$H=\Omega+3\overrightarrow{\Omega G}
=\left(
\begin{array}{c}
-\frac{11}{14}\\
\frac{33}{14}
\end{array}
\right)+3\left(
\begin{array}{c}
-\frac{1}{3}+\frac{11}{14}\\
\frac{7}{3}-\frac{33}{14}
\end{array}
\right)
=
\left(
\begin{array}{c}
\frac{4}{7}\\
\frac{16}{7}
\end{array}
\right).$$

\end{itemize}

Pour trouver le cercle circonscrit au triangle $(ABC)$, on a déjà le centre $\Omega$ et le rayon 

$$\Omega A=\sqrt{(1+\frac{11}{14})^2+(2-\frac{33}{14})^2}=\frac{1}{14}\sqrt{25^2+5^2}=\frac{5}{14}\sqrt{5^2+1}=\frac{5\sqrt{26}}{14}.$$

Il n'y a plus qu'à écrire l'équation cherchée~:

$$(x+\frac{11}{14})^2+(y-\frac{33}{14})^2=\frac{325}{98}\;\mbox{ou encore}\;
x^2+y^2+\frac{11}{7}x-\frac{33}{7}y+\frac{20}{7}=0.$$

Néanmoins, on peut trouver directement une équation de ce cercle. Les points $A$, $B$ et $C$ n'étant pas alignés, on sait que le cercle circonscrit existe et est unique.

Soient alors $(a,b,c)\in\Rr^3$ et $\mathcal{C}$ le cercle d'équation $x^2+y^2+ax+by+c=0$.

\begin{align*}
(A,B,C)\in\mathcal{C}^3&\Leftrightarrow
\left\{
\begin{array}{l}
a+2b+c=-5\\
-2a+b+c=-5\\
4b+c=-16
\end{array}
\right.\Leftrightarrow
\left\{
\begin{array}{l}
c=-4b-16
a-2b=11\\
-2a-3b=11
\end{array}\right.
\Leftrightarrow
\left\{
\begin{array}{l}
a=\frac{11}{7}\quad(\mbox{\textsc{Cramer}})\\
b=-\frac{33}{7}\\
c=\frac{20}{7}
\end{array}
\right. 
\end{align*}

\item  Les bissectrices de l'angle $A$ sont les deux droites constituées des points à égale distance des droites $(AB)$ et $(AC)$. Ces deux droites admettent pour vecteurs normaux respectifs $\vec{n}_1(1,-3)$ et $\vec{n}_2(2,1)$.

Soit $M(x,y)$ un point du plan.

\begin{align*}
d(M,(AB))=d(M,(AC))&\Leftrightarrow \frac{(\overrightarrow{AM}.\vec{n}_1)^2}{||\vec{n_1}||^2}=\frac{(\overrightarrow{AM}.\vec{n}_2)^2}{||\vec{n_2}||^2}\\
 &\Leftrightarrow\frac{((x-1)-3(y-2))^2}{10}=\frac{(2(x-1)+(y-2))^2}{5}\Leftrightarrow(x-3y+5)^2=2(2x+y-4)^2\\
 &\Leftrightarrow[(x-3y+5)+\sqrt{2}(2x+y-4)].[(x-3y+5)-\sqrt{2}(2x+y-4)]=0\\
 &\Leftrightarrow(1+2\sqrt{2})x+(-3+\sqrt{2})y+5-4\sqrt{2}=0\;\mbox{ou}\;(1-2\sqrt{2})x-(3+\sqrt{2})y+5+4\sqrt{2}=0\\
 &\Leftrightarrow y=(1+\sqrt{2})x+1-\sqrt{2}\;\mbox{ou}\;y=(1-\sqrt{2})x+1+\sqrt{2}
\end{align*}

La bissectrice intérieure $\delta_A$ de l'angle $\widehat{A}$ est la droite (pour certains, cette bissectrice est une demi-droite) passant par $A(2,1)$ et dirigée par le vecteur $\vec{u}=-\sqrt{10}.(\frac{1}{AB}\overrightarrow{AB}+\frac{1}{AC}\overrightarrow{AC})$. Ce vecteur a pour coordonnées $(3+\sqrt{2},1-2\sqrt{2})$.

Soit $M(x,y)$ un point du plan.

\begin{align*}
M\in\delta_A&\Leftrightarrow\mbox{det}(\overrightarrow{AM},\vec{u})=0\Leftrightarrow(1-2\sqrt{2})(x-1)-(3+\sqrt{2})(y-2)=0\\
 &\Leftrightarrow(1-2\sqrt{2})x-(3+\sqrt{2})y+5+4\sqrt{2}=0\Leftrightarrow y=(1-\sqrt{2})x+1+\sqrt{2}
\end{align*}

\end{enumerate}
\fincorrection
\correction{005199}
Soit $(x,y)\in\Rr^2$.

\begin{align*}
2x^2+5xy+3y^2-3x-2y-5&=2x^2+x(5y-3)+3y^2-2y-5=2(x+\frac{1}{4}(5y-3))^2-\frac{1}{8}(5y-3)^2+3y^2-2y-5\\
 &=\frac{1}{8}(4x+5y-3)^2-\frac{1}{8}y^2+\frac{14}{8}y-\frac{49}{8}\\
 &=\frac{1}{8}[(4x+5y-3)^2-(y^2-14y+49)]=\frac{1}{8}[(4x+5y-3)^2-(y-7)^2]\\
 &=\frac{1}{8}(4x+4y+4)(4x+6y-10)=(x+y+1)(2x+3y-5)
\end{align*}

Par suite,

$$\forall(x,y)\in\Rr^2,\;2x^2+5xy+3y^2-3x-2y-5=0\Leftrightarrow(x+y+1=0\;\mbox{ou}\;2x+3y-5=0.$$

$(E)$ est la réunion de la droite $(D_1)$ d'équation $x+y+1=0$ et de la droite $(D_2)$ d'équation $2x+3y-5=0$.

$$\includegraphics{../images/img005199-1}$$

La parallèle à $(D_1)$ passant par $O$ est la droite $(D_1')$ d'équation $x+y=0$ et la parallèle à $(D_2)$ passant par $O$ est la droite $(D_2')$ d'équation $2x+3y=0$. Ces droites se coupent en les quatre points $O(0,0)$, $A(-5,5)$, $B(-8,7)$ et $C(-3,2)$. L'aire de ce parallélogramme vaut $\left|\mbox{det}(\overrightarrow{OA},\overrightarrow{OC})\right|=5$.
\fincorrection
\correction{005200}
Notons $(D_1)$, $(D_2)$ et $(D_3)$ les droites d'équations respectives $y=2x+1$, $y=2x+7$ et $y=-\frac{1}{2}x$. Soit $\mathcal{C}$ un cercle.

Les droites $(D_1)$ et $(D_2)$ sont parallèles. Donc, $\mathcal{C}$ est un cercle tangent à $(D_1)$ et $(D_2)$ si et seulement si son centre est sur l'ensemble des points à égale distance de $(D_1)$ et $(D_2)$ à savoir la droite d'équation $y=2x+4$ et son rayon est la moitié de la distance de $(D_1)$ à $(D_2)$, ou encore la moitié de la distance d'un point de $(D_1)$, par exemple $(0,1)$, à $(D_2)$. Cette distance vaut $\frac{|2.0-1+7|}{\sqrt{2^2+1^2}}=\frac{6}{\sqrt{5}}$. 
Finalement, $\mathcal{C}$ est un cercle tangent à $(D_1)$ et $(D_2)$ si et seulement si son centre $\Omega$ a des coordonnées de la forme $(a,2a+4)$, $a\in\Rr$, et son rayon vaut $\frac{3}{\sqrt{5}}$.

Un cercle de centre $\Omega$ et de rayon $\frac{3}{\sqrt{5}}$ est tangent à $(D_3)$ si et seulement si la distance de $\Omega$ à $(D_3)$ est le rayon $\frac{3}{\sqrt{5}}$. Donc,

\begin{align*}
\mathcal{C}\;\mbox{solution}&\Leftrightarrow d(\Omega,(D_3))=\frac{3}{\sqrt{5}}
\Leftrightarrow\frac{|a+2(2a+4)|}{\sqrt{5}}=\frac{3}{\sqrt{5}}\Leftrightarrow|5a+8|=3\\
 &\Leftrightarrow5a+8=3\;\mbox{ou}\;5a+8=-3\Leftrightarrow a=-1\;\mbox{ou}\;a=-\frac{11}{5}
\end{align*}

On trouve deux cercles solutions, le cercle $\mathcal{C}_1$ de centre $\Omega_1(-1,2)$ et de rayon $\frac{3}{\sqrt{5}}$ et le cercle $\mathcal{C}_2$ de centre $\Omega_2(-\frac{11}{5},-\frac{2}{5})$ et de rayon $\frac{3}{\sqrt{5}}$
\fincorrection
\correction{005202}
 Pour $1\leq i\leq n$, notons $s_i$ la symétrie centrale de centre $A_i$. Le problème revient à trouver $n$ points $B_1$,..., $B_n$ tels que $B_2=s_1(B_1)$, $B_3=s_2(B_2)$,...,$B_n=s_{n-1}(B_{n-1})$, $B_1=s_n(B_n)$. Ceci équivaut à

$$\forall i\in\{2,...,n\},\;B_i=s_{i-1}\circ s_{i-2}\circ...\circ s_1(B_1)\;\mbox{et}\;B_1=s_n\circ s_{n-1}\circ...\circ s_1(B_1)\;(*).$$

Posons alors $f=s_n\circ s_{n-1}\circ...\circ s_1$. $f$ est une composée de symétries centrales. Il y a donc deux cas. Si $n$ est pair, on peut regrouper les symétries deux par deux. $f$ est alors (d'après l'exercice \ref{exo:suprou7bis}) une composée de translations et donc $f$ est une translation. Si $n$ est impair, $n-1$ est pair et donc la composée des $n-1$ premières symétries est une translation. Par suite, $f$ est la composée d'une translation et d'une symétrie centrale et est donc une symétrie centrale (d'après l'exercice \ref{exo:suprou7bis}).

Maintenant, $(*)$ a une solution si et seulement si $f$ a un point invariant.

\begin{itemize}
\item[\textbf{1er cas.}] Si $n$ est impair, $f$ étant une symétrie centrale, $f$ a un et un seul point invariant~:~son centre. Il existe donc un et un seul point $B_1$ vérfiant $B_1=s_n\circ s_{n-1}\circ...\circ s_1(B_1)$ et finalement, un et un seul $n$-uplet $(B_1,...,B_n)$ solution du problème posé.
\item[\textbf{2ème cas.}] Si $n$ est pair, $f$ est une translation. Si son vecteur est non nul, $f$ n'a pas de point invariant et le problème n'a pas de solution. Si son vecteur est nul, $f$ est l'identité et tout point est invariant par $f$.

Déterminons le vecteur de $f$. On pose $n=2p$. On a alors

$$f=s_{2p}\circ s_{2p-1}\circ...s_2\circ s_1=t_{2\overrightarrow{A_{2p-1}A_{2p}}}\circ...\circ t_{2\overrightarrow{A_{1}A_{2}}}=t_{2(\overrightarrow{A_1A_2}+...+\overrightarrow{A_{2p-1}A_{2p}})}.$$

Quand $n=2p$ est pair, le problème posé a des solutions si et seulement si $\overrightarrow{A_1A_2}+...+\overrightarrow{A_{2p-1}A_{2p}}=\overrightarrow{0}$.
\end{itemize}
\fincorrection
\correction{005203}
 Tout d'abord, pour $(x,y)\in\Rr^2$, $x^2+y^2-2x+4y+1=0\Leftrightarrow(x-1)^2+(y+2)^2=4$ et $\mathcal{C}$ est le cerlce de centre $\Omega(1,-2)$ et de rayon $2$.
\begin{enumerate}
\item  Le point $A(2,-2+\sqrt{3})$ est effectivement sur $\mathcal{C}$ car $(2-1)^2+(-2+\sqrt{3}+2)^2=1+3=4$. La tangente $(T)$ en $A$ à $\mathcal{C}$ est la droite passant par $A$ et de vecteur normal $\overrightarrow{A\Omega}$.

$$M(x,y)\in(T)\Leftrightarrow\overrightarrow{AM}.\overrightarrow{A\Omega}=0\Leftrightarrow(x-2)+\sqrt{3}(y+2-\sqrt{3})=0
\Leftrightarrow x+\sqrt{3}y-5+2\sqrt{3}=0.$$
\item  Soit $\mathcal{C}'$ le cercle de centre $(1,0)$ et de rayon $2$. Une équation de ce cercle est $x^2+y^2-2x-3=0$. Par suite,

\begin{align*}
M(x,y)\in\mathcal{C}\cap\mathcal{C}'&\Leftrightarrow
\left\{
\begin{array}{l}
x^2+y^2-2x+4y+1=0\\
x^2+y^2-2x-3=0
\end{array}
\right.\Leftrightarrow
\left\{
\begin{array}{l}
4y+4=0\quad((1)-(2))\\
x^2+y^2-2x-3=0
\end{array}
\right.
\Leftrightarrow
\left\{
\begin{array}{l}
y=-1\\
x^2-2x-2=0
\end{array}
\right.
\\
 &\Leftrightarrow\left\{
\begin{array}{l}
y=-1\\
x=1+\sqrt{3}\;\mbox{ou}\;x=1-\sqrt{3}
\end{array}
\right.
\end{align*}

Il ya donc deux points d'intersection~:~$(1+\sqrt{3},-1)$ et $(1-\sqrt{3},-1)$.
\end{enumerate}
\fincorrection
\correction{005204}
 
$$\includegraphics{../images/img005204-1}$$


Montrons tout d'abord que si $M$, $N$ et $P$ sont alignés, alors $\frac{\overline{MB}}{\overline{MC}}.\frac{\overline{NC}}{\overline{NA}}.\frac{\overline{PA}}
{\overline{PB}}=1$ $(*)$.

On suppose donc que $M$, $N$ et $P$ sont alignés et on note $(\Delta)$ la droite contenant $M$, $N$ et $P$.

\begin{itemize}
\item[\textbf{1ère solution.}] Soit $A_1$ le projeté de $A$ sur la droite $(BC)$ parallèlement à la droite $(\Delta)$. D'après le théorème de \textsc{Thales}, on a

$$\frac{\overline{NC}}{\overline{NA}}=\frac{\overline{MC}}{\overline{MA_1}}\;\mbox{et}\;
\frac{\overline{PA}}{\overline{PB}}=\frac{\overline{MA_1}}{\overline{MB}},$$

et donc,

$$\frac{\overline{MB}}{\overline{MC}}.\frac{\overline{NC}}{\overline{NA}}.\frac{\overline{PA}}
{\overline{PB}}=\frac{\overline{MB}}{\overline{MC}}.\frac{\overline{MC}}{\overline{MA_1}}.\frac{\overline{MA_1}}
{\overline{MB}}=1.$$

\item[\textbf{2ème solution.}] Soit $h_1$ l'homothétie de centre $M$ et de rapport $k_1=\frac{\overline{MB}}{\overline{MC}}$, de sorte que $h_1(C)=B$. Soit $h_2$ l'homothétie de centre $N$ et de rapport $k_2=\frac{\overline{NC}}{\overline{NA}}$, de sorte que $h_2(A)=C$.

Maintenant, le produit $k_1k_2$ peut-il être égal à $1$~?~Si c'était le cas, on aurait $\frac{\overline{MB}}{\overline{MC}}\frac{\overline{NC}}{\overline{NA}}=1$ et donc, $\frac{\overline{MB}}{\overline{MC}}=\frac{\overline{NA}}{\overline{NC}}$. La réciproque du théorème de \textsc{Thales} permettrait alors d'affirmer que $(MN)$ et $(AB)$ sont parallèles, ce qui n'est pas. Donc, $k_1k_2\neq1$ et d'après l'exercice \ref{exo:suprou7bis}, $h_1\circ h_2$ est une homothétie. Puisque $h_1\circ h_2$ transforme $A$ en $B$, son centre est sur la droite $(AB)$. Mais d'autre part, son centre est sur la droite des centres $(MN)$. Finalement, le centre de $h_1\circ h_2$ est le point d'intersection de $(MN)$ et $(AB)$, c'est-à-dire le point $P$.

Mais alors, le rapport de $h_1\circ h_2$ vaut également $\frac{\overline{PB}}{\overline{PA}}$. Ainsi, 
$\frac{\overline{MB}}{\overline{MC}}\frac{\overline{NC}}{\overline{NA}}=\frac{\overline{PB}}{\overline{PA}}$ et finalement, $\frac{\overline{MB}}{\overline{MC}}\frac{\overline{NC}}{\overline{NA}}\frac{\overline{PA}}{\overline{PB}}=1$.

\item[\textbf{3ème solution.}] On se place dans le repère $\mathcal{R}=(A,\overrightarrow{AB},\overrightarrow{AC})$. Dans ce repère, les coordonnées des différents points sont~:~$A(0,0)$, $B(1,0)$, $C(0,1)$, $M(m,1-m)$, $N(0,n)$ et $P(p,0)$ où $m$, $n$ et $p$ sont distincts de $0$ et de $1$.

Les coordonnées de $\overrightarrow{MB}$ sont $(1-m,m-1)$ et celles de $\overrightarrow{MC}$ sont $(-m,m)$. Par suite, $m\overrightarrow{MB}=(m-1)\overrightarrow{MC}$ et finalement, $\frac{\overline{MB}}{\overline{MC}}=\frac{m-1}{m}$. On trouve de même $\frac{\overline{NC}}{\overline{NA}}=\frac{n-1}{n}$ et $\frac{\overline{PA}}{\overline{PB}}=\frac{p}{p-1}$. Finalement,

$$\frac{\overline{MB}}{\overline{MC}}\frac{\overline{NC}}{\overline{NA}}\frac{\overline{PA}}{\overline{PB}}=
\frac{(m-1)(n-1)p}{mn(p-1)}.$$

Maintenant,

\begin{align*}
M,\;N\;\mbox{et}\;P\;\mbox{alignés}\;&\Leftrightarrow
\left|
\begin{array}{cc}
-m&p-m\\
m+n-1&m-1
\end{array}
\right|
\Leftrightarrow-m(m-1)-(p-m)(m+n-1)=0\\
 &\Leftrightarrow-pm-pn+p+mn=0\Leftrightarrow mn=p(m+n-1)\Leftrightarrow mn\\
 &=-p(m-1)(n-1)+pmn\Leftrightarrow p(m-1)(n-1)=mn(p-1)\\
 &\Leftrightarrow\frac{(m-1)(n-1)p}{mn(p-1)}=1
\end{align*}

\end{itemize}

Montrons maintenant que si $\frac{\overline{MB}}{\overline{MC}}\frac{\overline{NC}}{\overline{NA}}\frac{\overline{PA}}{\overline{PB}}=1$, alors les points $M$, $N$ et $P$ sont alignés. Pour cela, vérifions tout d'abord que $(MN)$ n'est pas parallèle à $(AB)$. Dans le cas contraire, le théorème de \textsc{Thales} fournirait $\frac{\overline{MB}}{\overline{MC}}\frac{\overline{NC}}{\overline{NA}}=1$ et donc  $\frac{\overline{PA}}{\overline{PB}}=1$, puis $\overline{PA}=\overline{PB}$ et finalement $\overline{AB}=0$, ce qui n'est pas.

Par suite, la droite $(MN)$ coupe la droite $(AB)$ en un point $P_1$ vérifiant d'après le début de l'exercice

$\frac{\overline{MB}}{\overline{MC}}\frac{\overline{NC}}{\overline{NA}}\frac{\overline{P_1A}}{\overline{P_1B}}=1$. On en déduit que $\frac{\overline{P_1A}}{\overline{P_1B}}=\frac{\overline{PA}}{\overline{PB}}$. Notons $k$ la valeur commune de ce rapport.

On a déjà que $k\neq1$, ou encore $1-k\neq0$. Par suite, $P_1=\mbox{bar}\{A(1),B(-k)\}=P$, ce qui montre que les points $M$, $N$ et $P$ sont alignés.
\fincorrection
\correction{005208}
\begin{enumerate}
\item  Le fait que $(D)$ et $(D')$ soient sécantes équivaut à $ab'-a'b\neq0$.

Soit $A(x_A,y_A)$ le point d'intersection de $(D)$ et $(D')$.

Si $(\Delta)$ est une droite ayant une équation de la forme $\lambda(ax+by+c)+\mu(a'x+b'y+c')=0$, $(\lambda,\mu)\neq(0,0)$ alors, puisque

$$\lambda(ax_A+by_A+c)+\mu(a'x_A+b'y_A+c')=\lambda.0+\mu.0=0,$$

le point $A$ appartient à $(\Delta)$.

Réciproquement, soit $(\Delta)$ une droite d'équation $\alpha x+\beta y+\gamma=0$, $(\alpha,\beta)\neq(0,0)$. Soit $\overrightarrow{v}$ le vecteur de coordonnées $(\alpha,\beta)$.
Puisque $ab'-a'b\neq0$, les deux vecteurs $\overrightarrow{u}(a,b)$ et $\overrightarrow{u}'(a',b')$ ne sont pas colinéaires. Mais alors, la famille $(\overrightarrow{u},\overrightarrow{u}')$ est une base du plan (vectoriel). Par suite, il existe $(\lambda,\mu)\neq(0,0)$ (car $\overrightarrow{v}\neq\overrightarrow{0}$) tel que $\overrightarrow{v}=\lambda\overrightarrow{u}+\mu\overrightarrow{u}'$, ou encore tel que $\alpha=\lambda a+\mu a'$ et $\beta=\lambda b+\mu b'$. Toute droite $(\Delta)$ admet donc une équation cartésienne de la forme 
$\lambda(ax+by)+\mu(a'x+b'y)+\gamma=0$, $(\lambda,\mu)\neq(0,0)$.

Maintenant, si $A\in(\Delta)$, alors

$$\gamma=-\lambda(ax_A+by_A)+\mu(a'x_A+b'y_A)=-\lambda(-c)-\mu(-c')=\lambda c+\mu c'.$$

Finalement, si $A\in(\Delta)$, $(\Delta)$ admet une équation de la forme $\lambda(ax+by+c)+\mu(a'x+b'y+c')=0$, $(\lambda,\mu)\neq(0,0)$.

\item  Les deux droites $(D)$ et $(D')$ considérées sont bien sécantes car $5.2-7(-3)=31\neq0$. Notons $A$ leur point d'intersection et $B$ le point de coordonnées $(1,0)$. $B$ n'est sur aucune des deux droites considérées de sorte qu'il existe une et seule droite, notée $(\Delta)$, solution du problème posé.

Puisque $(\Delta)$ passe par $A$, $(\Delta)$ a une équation de la forme $\lambda(5x+7y+1)+\mu(-3x+2y+1)=0$. Il est clair que l'on ne peut avoir $\lambda=0$ (car $(\Delta)$ n'est pas $(D')$) et après division par $\lambda$, l'équation s'écrit sous la forme $(5x+7y+1)+k(-3x+2y+1)=0$ où $k$ est un réel. 
Maintenant, $(\Delta)$ passe par $B$ si et seulement si $6-2k=0$ ou encore $k=3$.

Une équation cartésienne de $(\Delta)$ est donc $(5x+7y+1)+3(-3x+2y+1)=0$ ou encore $-4x+13y+4=0$.

\item  Soit $M(x,y)$ un point du plan.

\begin{align*}
\forall m\in\Rr,\;M\in(D_m)&\Leftrightarrow\forall m\in\Rr,\;(2m-1)x+(m+1)y-4m-1=0
\Leftrightarrow\forall m\in\Rr,\;m(2x+y-4)-x+y-1=0\\
 &\Leftrightarrow
\left\{
\begin{array}{l}
2x+y-4=0\\
-x+y-1=0
\end{array}
\right.\Leftrightarrow x=1\;\mbox{et}\;y=2
\end{align*}

Toutes les droites $(D_m)$ passent par le point $A(1,2)$.

La droite $(D_{-1})$ passe par $A$ et est parallèle à $(Oy)$. Ensuite, pour $m\neq-1$, $(D_m)$ est la droite passant par $A$ et de coefficient directeur $f(m)=\frac{-2m+1}{m+1}=-2+\frac{3}{m+1}$. Quand $m$ décrit $\Rr\setminus\{-1\}$, $f(m)$ prend toutes les valeurs réelles sauf $-2$.

La droite passant par $A$ de coefficient directeur $-2$ (et donc d'équation $y=-2x+4$) n'est pas une droite $(D_m)$. Toute autre droite passant par $A$ est une droite $(D_m)$.

\end{enumerate}
\fincorrection
\correction{005510}
\textbullet~\textbf{Repère de $(D)$.}

\begin{center}
$\left\{
\begin{array}{l}
x-z-a=0\\
y+3z+1=0
\end{array}
\right.\Leftrightarrow\left\{
\begin{array}{l}
x=a+z\\
y=-1-3z
\end{array}
\right.$.
\end{center}

$(D)$ est la droite passant par $A(a,-1,0)$ et dirigée par $u(1,-3,1)$.
 \textbullet~\textbf{Repère de $(D')$.}

\begin{center}$\left\{
\begin{array}{l}
x+2y+z-2b=0\\
3x+3y+2z-7=0
\end{array}
\right.\Leftrightarrow\left\{
\begin{array}{l}
2y+z=2b-x\\
3y+2z=7-3x
\end{array}
\right.\Leftrightarrow\left\{
\begin{array}{l}
y=4b-7+x\\
z=14-6b-3x
\end{array}
\right.$
\end{center}
$(D')$ est la droite passant par $A'(0,4b-7,-6b+14)$ et dirigée par $u'(1,1,-3)$.
\textbullet~Les vecteurs $u$ et $u'$ ne sont pas colinéaires et donc $(D)$ et $(D')$ ne sont pas parallèles.
\textbullet~Le plan $(P)$ contenant $(D)$ et parallèle à $(D')$ est le plan de repère $(A,u,u')$. Déterminons une équation de ce plan.

\begin{center}
$M(x,y,z)\in(P)\Leftrightarrow
\left|
\begin{array}{ccc}
x-a&1&1\\
y+1&-3&1\\
z&1&-3
\end{array}\right|=0\Leftrightarrow 8(x-a)+4(y+1)+4z=0\Leftrightarrow2x+y+z=2a-1$.
\end{center}
\textbullet~Enfin, $(D)$ et $(D')$ sont sécantes si et seulement si $(D')$ est contenue dans $(P)$. Comme $(D')$ est déjà parallèle à $(P)$, on a

\begin{center}
$(D)$ et $(D')$ sécantes $\Leftrightarrow A'\in(P)\Leftrightarrow(4b-7)+(-6b+14)=2a-1\Leftrightarrow b=-a+4$.
\end{center}

\begin{center}
\shadowbox{
\begin{tabular}{c}
$(D)$ et $(D')$ sont sécantes si et seulement si $b=-a+4$ et dans ce cas, une équation du plan contenant $(D)$\\
 et $(D')$ est $2x+y+z=2a-1$.
 \end{tabular}
}
\end{center}
\fincorrection
\correction{005511}
\textbullet~$(\Delta)$ est parallèle à $(D)$ si et seulement si $(\Delta)$ est dirigée par le vecteur $u(3,2,1)$ ou encore $(\Delta)$ admet un système d'équations paramétriques de la forme $\left\{
\begin{array}{l}
x=a+3\lambda\\
y=b+2\lambda\\
z=c+\lambda
\end{array}
\right.$. Ensuite, 
$(\Delta)$ est sécante à $(D_1)$ si et seulement si on peut choisir le point $(a,b,c)$ sur $(D_1)$ ou encore si et seulement si $(\Delta)$ admet un système d'équations paramétriques de la forme $\left\{
\begin{array}{l}
x=3\lambda\\
y=b+2\lambda\\
z=4+\lambda
\end{array}
\right.$.
Enfin,

\begin{center}
$(\Delta)$ et $(D_2)$ sécantes $\Leftrightarrow\exists\lambda\in\Rr/\;b+2\lambda=4+\lambda+4=0\Leftrightarrow b+2\times(-8)=0\Leftrightarrow b=16$.
\end{center}
Ceci démontre l'existence et l'unicité de $(\Delta)$ : un système d'équations paramétriques de $(\delta)$ est $\left\{
\begin{array}{l}
x=3\lambda\\
y=16+2\lambda\\
z=4+\lambda
\end{array}
\right.$. Un système d'équations cartésiennes de $(\Delta)$ est $\left\{
\begin{array}{l}
x=3(z-4)\\
y=16+2(z-4)
\end{array}
\right.$ ou encore

\begin{center}
\shadowbox{
$(\Delta)$ : $\left\{
\begin{array}{l}
x-3z+12=0\\
y-2z-8=0
\end{array}
\right.$.
}
\end{center}
\fincorrection
\correction{005512}
Notons $(\Delta)$ une éventuelle droite solution.
\textbullet~$(\Delta)$ est sécante à $(D_1)$ et $(D_2)$ si et seulement si $(\Delta)$ passe par un point de la forme $(1,0,a)$ et par un point de la forme $(b,1,0)$ ou encore si et seulement si $(\Delta)$ passe par un point de la forme $(1,0,a)$ et est dirigée par un vecteur de la forme $(b-1,1,-a)$.
Ainsi, $(\Delta)$ est sécante à $(D_1)$ et $(D_2)$ si et seulement si $(\Delta)$ admet un système d'équations paramétriques de la forme $\left\{
\begin{array}{l}
x=1+\lambda(b-1)\\
y=\lambda\\
z=a-\lambda a
\end{array}
\right.$ ou encore un système d'équations cartésiennes de la forme $\left\{
\begin{array}{l}
x-(b-1)y=1\\
ay+z=a
\end{array}
\right.$.

\textbullet~Ensuite, $(\Delta)$ et $(D_3)$ sécantes $\Leftrightarrow\exists y\in\Rr/\;\left\{
\begin{array}{l}
-(b-1)y=1\\
ay+1=a
\end{array}
\right.\Leftrightarrow b\neq 1\;\text{et}\;-\frac{a}{b-1}+1=a\Leftrightarrow b\neq0\;\text{et}\;b\neq1\;\text{et}\;a=1-\frac{1}{b}$.
En résumé, les droites sécantes à $(D_1)$, $(D_2)$ et $(D_3)$ sont les droites dont un système d'équations cartésiennes est 

\begin{center}
$\left\{
\begin{array}{l}
x-(b-1)y=1\\
\left(1-\frac{1}{b}\right)y+z=1-\frac{1}{b}
\end{array}
\right.$, $b\notin\{0,1\}$.
\end{center}
Enfin,

\begin{align*}\ensuremath
(\Delta)\;\text{et}\;(D)\;\text{sécantes}&\Leftrightarrow\exists(x,y,z)\in\Rr^3/\;\left\{
\begin{array}{l}
x-(b-1)y=1\\
\left(1-\frac{1}{b}\right)y+z=1-\frac{1}{b}\\
x=y=-6z
\end{array}
\right.\\
&\Leftrightarrow\exists(x,y,z)\in\Rr^3/\;\left\{
\begin{array}{l}
-6z+6(b-1)z=1\\
-6\left(1-\frac{1}{b}\right)z+z=1-\frac{1}{b}\\
x=y=-6z
\end{array}
\right.\\
 &\Leftrightarrow b\notin\{0,1,2\}\;\text{et}\;-6\left(1-\frac{1}{b}\right)\frac{1}{6(b-2)}+\frac{1}{6(b-2)}=1-\frac{1}{b}\\
 &\Leftrightarrow b\notin\{0,1,2\}\;\text{et}\;-6(b-1)+b=6(b-1)(b-2)\Leftrightarrow b\notin\{0,1,2\}\;\text{et}\;6b^2-13b+6=0\\
 &\Leftrightarrow b\in\left\{\frac{2}{3},\frac{3}{2}\right\}.
\end{align*}

\begin{center}
\shadowbox{
Les droites solutions sont $(\Delta_1)$ : $\left\{
\begin{array}{l}
3x+y=3\\
y-2z=1
\end{array}
\right.$ et $(\Delta_2)$ : $\left\{
\begin{array}{l}
2x-y=2\\
y+3z=1
\end{array}
\right.$.
}
\end{center}
\fincorrection
\correction{005513}
\textbullet~Déterminons le centre de gravité $G$.

\begin{center}
$G=\frac{1}{3}A+\frac{1}{3}B+\frac{1}{3}C=\frac{1}{3}(2,-2,0)+\frac{1}{3}(4,2,6)+\frac{1}{3}(-1,-3,0)=\left(\frac{5}{3},-1,2\right)$.
\end{center}
\textbullet~Déterminons le centre du cercle circonscrit $O$. Une équation du plan $(ABC)$ est $\left|
\begin{array}{ccc}
x-2&2&-3\\
y+2&4&-1\\
z&6&0
\end{array}
\right|=0$ ou encore $6(x-2)-18(y+2)+10z=0$ ou enfin $3x-9y+5z=24$. Posons alors $O(a,b,c)$.
Ensuite, $OA=OB\Leftrightarrow(a-2)^2+(b+2)^2+c^2=(a-4)^2+(b-2)^2+(c-6)^2\Leftrightarrow4a+8b+12c=48\Leftrightarrow a+2b+3c=16$ et 
$OA=OC\Leftrightarrow(a-2)^2+(b+2)^2+c^2=(a+1)^2+(b+3)^2+c^2\Leftrightarrow-6a-2b=2\Leftrightarrow3a+b=-1$.
D'où le système

\begin{align*}\ensuremath
\left\{
\begin{array}{l}
3a-9b+5c=24\\
a+2b+3c=16\\
3a+b=-1
\end{array}
\right.&\Leftrightarrow\left\{
\begin{array}{l}
b=-3a-1\\
3a-9(-3a-1)+5c=24\\
a+2(-3a-1)+3c=16
\end{array}
\right.\Leftrightarrow\left\{
\begin{array}{l}
b=-3a-1\\
6a+c=3\\
-5a+3c=18
\end{array}
\right.\\
 &\Leftrightarrow\left\{
\begin{array}{l}
b=-3a-1\\
c=3-6a\\
-5a+3(3-6a)=18
\end{array}
\right.
\Leftrightarrow\left\{
\begin{array}{l}
a=-\frac{9}{23}\\
b=\frac{4}{23}\rule{0mm}{7mm}\\
c=\frac{123}{23}\rule{0mm}{7mm}
\end{array}
\right.
\end{align*}
Donc $O\left(-\frac{9}{23},\frac{4}{23},\frac{123}{23}\right)$.
\textbullet~Déterminons l'orthocentre $H$. D'après la relation d'\textsc{Euler},

\begin{center}
$H=O+3\overrightarrow{OG}=\left(-\frac{9}{23},\frac{4}{23},\frac{123}{23}\right)+3\left(-\frac{9}{23}-\frac{5}{3},\frac{4}{23}+1,\frac{123}{23}-2\right)=\left(\frac{-151}{23},\frac{85}{23},\frac{354}{23}\right)$.
\end{center}
\textbullet~Déterminons le centre du cercle inscrit $I$. On sait que $I=\text{bar}\left\{A(a),B(b),C(c)\right\}$ où $a=BC=\sqrt{5^2+5^2+6^2}=\sqrt{86}$, $b=AC=\sqrt{3^2+1^2+0^2}=\sqrt{10}$ et $c=AB=\sqrt{2^2+4^2+6^2}=\sqrt{54}$. Donc

\begin{align*}\ensuremath
I&=\frac{\sqrt{86}}{\sqrt{86}+\sqrt{10}+\sqrt{54}}A
+\frac{\sqrt{10}}{\sqrt{86}+\sqrt{10}+\sqrt{54}}B+\frac{\sqrt{54}}{\sqrt{86}+\sqrt{10}+\sqrt{54}}C\\
 &=\left(
 \frac{2\sqrt{86}+4\sqrt{10}-\sqrt{54}}{\sqrt{86}+\sqrt{10}+\sqrt{54}},
 \frac{-2\sqrt{86}+2\sqrt{10}-3\sqrt{54}}{\sqrt{86}+\sqrt{10}+\sqrt{54}},
 \frac{6\sqrt{10}}{\sqrt{86}+\sqrt{10}+\sqrt{54}}
 \right).
\end{align*}

Dans $\Rr^3$ euclidien rapporté à un repère orthonormé, on donne $A(2,-2,0)$, $B(4,2,6)$ et 
$C(-1,-3,0)$. Déterminer l'orthocentre, le centre de gravité, les centres des cercles circonscrits et inscrits au triangle $(A,B,C)$.
\begin{center}
\shadowbox{
\begin{tabular}{c}
$G\left(\frac{5}{3},-1,2\right)$, $O\left(-\frac{9}{23},\frac{4}{23},\frac{123}{23}\right)$ et $H\left(\frac{-151}{23},\frac{85}{23},\frac{354}{23}\right)$ puis\\
$I\left(\frac{2\sqrt{86}+4\sqrt{10}-\sqrt{54}}{\sqrt{86}+\sqrt{10}+\sqrt{54}},\frac{-2\sqrt{86}+2\sqrt{10}-3\sqrt{54}}{\sqrt{86}+\sqrt{10}+\sqrt{54}},\frac{6\sqrt{10}}{\sqrt{86}+\sqrt{10}+\sqrt{54}}\right)$.\rule{0mm}{10mm}
\end{tabular}
}
\end{center}
\fincorrection
\correction{005514}
\textbullet~Déterminons un repère de $(D)$.
\begin{center}
$\left\{
\begin{array}{l}
x+y+z+1=0\\
2x+y+5z=2
\end{array}
\right.\Leftrightarrow\left\{
\begin{array}{l}
x+y=-1-z\\
2x+y=2-5z
\end{array}
\right.\Leftrightarrow\Leftrightarrow\left\{
\begin{array}{l}
x=3-4z\\
y=-4+3z
\end{array}
\right.$. 
\end{center}
Un repère de $(D)$ est $\left(A,\overrightarrow{u}\right)$ où $A(3,-4,0)$ et $\overrightarrow{u}(-4,3,1)$.
\textbullet~Soit $M(x,y,z)$ un point du plan. On sait que

\begin{center}
$d(A,(D))=\frac{\|\overrightarrow{AM}\wedge\overrightarrow{u}\|}{\|\overrightarrow{u}\|}=\frac{\sqrt{(y-3z+4)^2+(x+4z-3)^2+(3x+4y+7)^2}}{\sqrt{26}}$
\end{center}
\textbullet~Notons $\mathcal{C}$ le cylindre de révolution d'axe $(D)$ et de rayon $2$.
\begin{center}
$M(x,y,z)\in\mathcal{C}\Leftrightarrow d(A,(D))=2\Leftrightarrow(y-3z+4)^2+(x+4z-3)^2+(3x+4y+7)^2=104$
\end{center}

\begin{center}
\shadowbox{
\begin{tabular}{c}
Une équation cartésienne  du cylindre de révolution d'axe $(D)$ et de rayon $2$ est\\
$(y-3z+4)^2+(x+4z-3)^2+(3x+4y+7)^2=104$.
\end{tabular}
}
\end{center}
\fincorrection
\correction{005515}
\textbullet~Déterminons un repère de $(D)$.

\begin{center}
$\left\{
\begin{array}{l}
x+y+z+1=0\\
2x+y+5z=2
\end{array}
\right.\Leftrightarrow\left\{
\begin{array}{l}
x+y=-z-1\\
2x+y=-5z+2
\end{array}
\right.\Leftrightarrow\left\{
\begin{array}{l}
x=-4z+3\\
y=3z-4
\end{array}
\right.$
\end{center}
Un repère de $(D)$ est $\left(A,\overrightarrow{u}\right)$ où $A(3,-4,0)$ et $\overrightarrow{u}(-4,3,1)$.
\textbullet~Déterminons un repère de $(D')$.

\begin{center}
$\left\{
\begin{array}{l}
x+y+z=2\\
2x+y-5z=3
\end{array}
\right.\Leftrightarrow\left\{
\begin{array}{l}
x+y=-z+2\\
2x+y=5z+3
\end{array}
\right.\Leftrightarrow\left\{
\begin{array}{l}
x=6z+1\\
y=-7z+1
\end{array}
\right.$
\end{center}
Un repère de $(D')$ est $\left(A',\overrightarrow{u'}\right)$ où $A'(1,1,0)$ et $\overrightarrow{u'}(6,-7,1)$.
\textbullet~$\overrightarrow{u}\wedge\overrightarrow{u'}=\left(
\begin{array}{c}
-4\\
3\\
1
\end{array}
\right)\wedge\left(
\begin{array}{c}
6\\
-7\\
1
\end{array}
\right)=\left(
\begin{array}{c}
10\\
10\\
10
\end{array}
\right)\neq\overrightarrow{0}$.
\rule{0mm}{5mm}Puisque $\overrightarrow{u}$ et $\overrightarrow{u'}$ ne sont pas colinéaires, les droites $(D)$ et $(D')$ ne sont parallèles. Ceci assure l'unicité de la perpendiculaire commune  à $(D)$ et $(D')$.
\textbullet~On sait que la distance $d$ de $(D)$ à $(D')$ est donnée par

\begin{center}
$d=\frac{\text{abs}\left(\left[\overrightarrow{AA'},\overrightarrow{u},\overrightarrow{u'}\right]\right)}{\|\overrightarrow{u}\wedge\overrightarrow{u'}\|}$,
\end{center}
avec $[\overrightarrow{AA'},\overrightarrow{u},\overrightarrow{u'}]=\left|
\begin{array}{ccc}
-2&-4&6\\
5&3&-7\\
0&1&1
\end{array}
\right|=10\times(-2)+10\times5=30$ et donc $d=\frac{30}{10\sqrt{3}}=\sqrt{3}$.

\begin{center}
\shadowbox{
$d((D),(D'))=\sqrt{3}$.
}
\end{center}

\textbullet~Un système d'équations de la perpendiculaire commune est 
$\left\{
\begin{array}{l}
\left[\overrightarrow{AM},\overrightarrow{u},\overrightarrow{u}\wedge\overrightarrow{u'}\right]=0\\
\rule{0mm}{6mm}\left[\overrightarrow{A'M},\overrightarrow{u'},\overrightarrow{u}\wedge\overrightarrow{u'}\right]=0
\end{array}
\right.$. Or,

\begin{center}
$\frac{1}{10}\left[\overrightarrow{AM},\overrightarrow{u},\overrightarrow{u}\wedge\overrightarrow{u'}\right]=\left|
\begin{array}{ccc}
x-3&-4&1\\
y+4&3&1\\
z&1&1
\end{array}
\right|=2(x-3)+5(y+4)-7z=2x+5y-7z+14$,
\end{center}
et

\begin{center}
$\frac{1}{10}\left[\overrightarrow{A'M},\overrightarrow{u'},\overrightarrow{u}\wedge\overrightarrow{u'}\right]=\left|
\begin{array}{ccc}
x-1&6&1\\
y-1&-7&1\\
z&1&1
\end{array}
\right|=-8(x-1)-5(y-1)+13z=-8x-5y+13z+13$.
\end{center}
Donc

\begin{center}
\shadowbox{
\begin{tabular}{c}
un système d'équations cartésienne de la perpendiculaire commune à $(D)$ et $(D')$ est\\
$\left\{
\begin{array}{l}
2x+5y-7z=-14\\
8x+5y-13z=13
\end{array}
\right.$.
\end{tabular}
}
\end{center}
\fincorrection
\correction{005516}
$\overrightarrow{u}\in\overrightarrow{P_1}\cap\overrightarrow{P_2}\cap\overrightarrow{P_3}\Leftrightarrow\left\{
\begin{array}{l}
z-2y=0\\
2x-3z=0\\
3y-x=0
\end{array}
\right.\Leftrightarrow\left\{
\begin{array}{l}
x=3y\\
z=2y
\end{array}
\right.$.
Ainsi, les plans $(P_1)$, $(P_2)$ et $(P_3)$ sont tous trois parallèles à la droite affine $(D)$ d'équations $\left\{
\begin{array}{l}
x=3y\\
z=2y
\end{array}
\right.$. Ces plans définissent donc un prisme.
Déterminons alors l'aire d'une section droite. Le plan $(P)$ d'équation $3x+y+2z=0$ est perpendiculaire à la droite $(D)$. Son intersection avec les plans $(P_1)$, $(P_2)$ et $(P_3)$ définit donc une section droite du prisme.
\textbullet~Soit $M(x,y,z)$ un point de l'espace.

\begin{center}
$M\in(P_1)\cap(P_2)\cap(P)\Leftrightarrow
\left\{
\begin{array}{l}
z-2y=5\\
2x-3z=0\\
3x+y+2z=0
\end{array}
\right.\Leftrightarrow\left\{
\begin{array}{l}
y=\frac{z-5}{2}\\
x=\frac{3}{2}z\\
\frac{9}{2}z+\frac{z-5}{2}+2z=0
\end{array}
\right.\Leftrightarrow\left\{
\begin{array}{l}
z=\frac{5}{14}\\
y=-\frac{65}{28}\rule[-4mm]{0mm}{11mm}\\
x=\frac{15}{28}
\end{array}
\right.$
\end{center}
Notons $A\left(\frac{15}{28},-\frac{65}{28},\frac{5}{14}\right)$.
\textbullet~Soit $M(x,y,z)$ un point de l'espace.

\begin{center}
$M\in(P_1)\cap(P_3)\cap(P)\Leftrightarrow
\left\{
\begin{array}{l}
z-2y=5\\
3y-x=0\\
3x+y+2z=0
\end{array}
\right.\Leftrightarrow\left\{
\begin{array}{l}
z=2y+5\\
x=3y\\
9y+y+2(2y+5)=0
\end{array}
\right.\Leftrightarrow\left\{
\begin{array}{l}
y=-\frac{5}{7}\\
x=-\frac{15}{7}\rule[-4mm]{0mm}{11mm}\\
z=\frac{25}{7}
\end{array}
\right.$
\end{center}
Notons $B\left(-\frac{15}{7},-\frac{5}{7},\frac{25}{7}\right)$.

\textbullet~Soit $M(x,y,z)$ un point de l'espace.

\begin{center}
$M\in(P_2)\cap(P_3)\cap(P)\Leftrightarrow
\left\{
\begin{array}{l}
2x-3z=0\\
3y-x=0\\
3x+y+2z=0
\end{array}
\right.\Leftrightarrow x=y=z=0$
\end{center}
Une section droite est $OAB$ où $A\left(\frac{15}{28},-\frac{65}{28},\frac{5}{14}\right)$ et $B\left(-\frac{15}{7},-\frac{5}{7},\frac{25}{7}\right)$. De plus

\begin{align*}\ensuremath
\text{aire de}(OAB)&=\frac{1}{2}\left\|\overrightarrow{OA}\wedge\overrightarrow{OB}\right\|=\frac{1}{2}\times\frac{5}{28}\times\frac{5}{7}\left\|\left(
\begin{array}{c}
3\\
-13\\
2
\end{array}\right)
\wedge
\left(
\begin{array}{c}
-3\\
-1\\
5
\end{array}\right)\right\|=\frac{1}{2}\times\frac{5}{28}\times\frac{5}{7}\sqrt{63^2+21^2+42^2}\\
 &=\frac{1}{2}\times\frac{5}{28}\times\frac{5}{7}\times21\sqrt{3^2+1^2+2^2}=\frac{75}{4\sqrt{14}}
\end{align*}
\begin{center}
\shadowbox{
L'aire d'une section droite est $\frac{75}{4\sqrt{14}}$.
}
\end{center}
\fincorrection
\correction{005517}
Soient $(P)$ le plan d'équation $x+2y+2z=3$ et $(P')$ le plan d'équation $x+y=0$. L'angle entre $(P)$ et $(P')$ est l'angle entre les vecteurs normaux $\overrightarrow{n}(1,2,2)$ et $\overrightarrow{n'}(1,1,0)$ :

\begin{center}
$\left(\widehat{\overrightarrow{n},\overrightarrow{n'}}\right)=\Arccos\left(\frac{\overrightarrow{n}.\overrightarrow{n'}}{\|\overrightarrow{n}\|\|\overrightarrow{n'}\|}\right)=\Arccos\left(\frac{3}{3\sqrt{2}}\right)=\Arccos\left(\frac{1}{\sqrt{2}}\right)=\frac{\pi}{4}$.
\end{center}
\fincorrection
\correction{005518}
Soit $M(x,y,z)$ un point de l'espace. On a 
\begin{center}
$d(M,(P_1))=\frac{\left|4x+4y-7z-1\right|}{\sqrt{4^2+4^2+7^2}}=\frac{\left|4x+4y-7z-1\right|}{9}$ et $d(M,(P_2))=\frac{\left|8x-4y+z+7\right|}{\sqrt{8^2+4^2+1^2}}=\frac{\left|8x-4y+z+7\right|}{9}$.
\end{center}
Par suite,

\begin{align*}\ensuremath
d(M,(P_1))=d(M,(P_2)&\Leftrightarrow|4x+4y-7z-1|=|8x-4y+z+7|\Leftrightarrow(4x+4y-7z-1)^2=(8x-4y+z+7)^2\\
 &\Leftrightarrow\left((4x+4y-7z-1)-(8x-4y+z+7)\right)\left((4x+4y-7z-1)+(8x-4y+z+7)\right)=0\\
 &\Leftrightarrow(-4x+8y-8z-8)(12x-6z+6)=0\Leftrightarrow x-2y+2z+2=0\;\text{ou}\;2x-z+1=0.
\end{align*}
\begin{center}
\shadowbox{
Les plans bissecteurs de $(P_1)$ et $(P_2)$ admettent pour équation cartésienne $x-2y+2z+2=0$ et $2x-z+1=0$.
}
\end{center}
\fincorrection
\correction{005519}
\textbullet~Déterminons un repère de $(D)$.

\begin{center}
$\left\{
\begin{array}{l}
x+y-3z+4=0\\
2x-z+1=0
\end{array}
\right.\Leftrightarrow\left\{
\begin{array}{l}
y-3z=-x-4\\
z=2x+1
\end{array}
\right.\Leftrightarrow\left\{
\begin{array}{l}
y=5x-1\\
z=2x+1
\end{array}
\right.$
\end{center}
Un repère de $(D)$ est $\left(A,\overrightarrow{u}\right)$ où $A(0,-1,1)$ et $\overrightarrow{u}(1,5,2)$.
\textbullet~Puisque un système d'équations de $(D')$ est $\left\{
\begin{array}{l}
x=z-1\\
y=z-1
\end{array}
\right.$, un repère de $(D')$ est $\left(A',\overrightarrow{u'}\right)$ où $A'(-1,-1,0)$ et $\overrightarrow{u'}(1,1,1)$.
\textbullet~$\overrightarrow{u}\wedge\overrightarrow{u'}=\left(
\begin{array}{c}
1\\
5\\
2
\end{array}
\right)\wedge\left(
\begin{array}{c}
1\\
1\\
1
\end{array}
\right)=\left(
\begin{array}{c}
3\\
1\\
-4
\end{array}
\right)\neq\overrightarrow{0}$.
\rule{0mm}{5mm}Puisque $\overrightarrow{u}$ et $\overrightarrow{u'}$ ne sont pas colinéaires, les droites $(D)$ et $(D')$ ne sont parallèles. Ceci assure l'unicité de la perpendiculaire commune  à $(D)$ et $(D')$.

\textbullet~Un système d'équations de la perpendiculaire commune est 
$\left\{
\begin{array}{l}
\left[\overrightarrow{AM},\overrightarrow{u},\overrightarrow{u}\wedge\overrightarrow{u'}\right]=0\\
\rule{0mm}{6mm}\left[\overrightarrow{A'M},\overrightarrow{u'},\overrightarrow{u}\wedge\overrightarrow{u'}\right]=0
\end{array}
\right.$. Or,

\begin{center}
$\left[\overrightarrow{AM},\overrightarrow{u},\overrightarrow{u}\wedge\overrightarrow{u'}\right]=\left|
\begin{array}{ccc}
x&1&3\\
y+1&5&1\\
z-1&2&-4
\end{array}
\right|=-22x+10(y+1)-14(z-1)=-22x+10y-14z+24$,
\end{center}
et

\begin{center}
$\left[\overrightarrow{A'M},\overrightarrow{u'},\overrightarrow{u}\wedge\overrightarrow{u'}\right]=\left|
\begin{array}{ccc}
x+1&1&3\\
y+1&1&1\\
z&1&-4
\end{array}
\right|=-5(x+1)+7(y+1)-2z=-5x+7y-2z+2$.
\end{center}
Donc

\begin{center}
\shadowbox{
\begin{tabular}{c}
un système d'équations cartésienne de la perpendiculaire commune à $(D)$ et $(D')$ est\\
$\left\{
\begin{array}{l}
11x-5y+7z=12\\
5x-7y+2z=2
\end{array}
\right.$.
\end{tabular}
}
\end{center}

\fincorrection
\correction{005521}



$$\includegraphics{../images/img005521-1}$$


\textbf{Angle entre deux arêtes.} Les faces du tétraèdre $ABCD$ sont des triangles équilatéraux et donc l'angle entre deux arêtes est $60^\circ$.

$$\includegraphics{../images/img005521-2}$$

\textbf{Angle entre une  arête et une face.} C'est l'angle $\widehat{CDI}$ de la figure ci-dessus.

\begin{center}
$\widehat{CDI}=\Arccos\left(\frac{HD}{DI}\right)=\Arccos\left(\frac{a/2}{a\sqrt{3}/2}\right)=\Arccos\left(\frac{1}{\sqrt{3}}\right)=54,7\ldots^\circ$.
\end{center}
\textbf{Angle entre deux faces.} C'est l'angle $\widehat{CID}$ de la figure ci-dessus.

\begin{center}
$\widehat{CID}=\pi-2\widehat{CDI}=2\left(\frac{\pi}{2}-\Arccos\left(\frac{1}{\sqrt{3}}\right)\right)=2\Arcsin\left(\frac{1}{\sqrt{3}}\right)=70,5\ldots^\circ$.
\end{center}
\fincorrection
\correction{005522}
Déterminons un repère de $(D)$.

\begin{center}
$\left\{
\begin{array}{l}
x-y-z=0\\
x+2y-z=10
\end{array}
\right.\Leftrightarrow\left\{
\begin{array}{l}
x-z=y\\
y+2y=10
\end{array}
\right.\Leftrightarrow\left\{
\begin{array}{l}
y=\frac{10}{3}\\
z=x-\frac{10}{3}
\end{array}
\right.$.
\end{center}
Un repère de $(D)$ est $\left(A,\overrightarrow{u}\right)$ où $A\left(\frac{10}{3},\frac{10}{3},0\right)$ et $\overrightarrow{u}(1,0,1)$.
On sait alors que

\begin{center}
$d(O,(D))=\frac{\|\overrightarrow{AO}\wedge\overrightarrow{u}\|}{\|\overrightarrow{u}\|}=\frac{1}{\sqrt{2}}\times\frac{10}{3}\left\|
\left(
\begin{array}{c}
1\\
1\\
0
\end{array}
\right)\wedge\left(
\begin{array}{c}
1\\
0\\
1
\end{array}
\right)\right\|=\frac{10}{\sqrt{6}}
$.
\end{center}
\begin{center}
\shadowbox{
$d(O,(D))=\frac{10}{\sqrt{6}}$.
}
\end{center}
\fincorrection
\correction{006884}
$(D)$ est une droite de vecteur normal $\vec n = (2,-3)$. Le projeté orthogonal 
$p(M_0)$ de $M_0$ sur $(D)$ est de la forme $M_0+\lambda.\vec{n}$ 
où $\lambda$ est un réel à déterminer.
Le point $M_0+\lambda.\vec{n}$ a pour coordonnées $(x_0+2\lambda,y_0-3\lambda)$.

$$M_0+\lambda.\vec{n}\in(D)\iff 2(x_0+2\lambda)-3(y_0-3\lambda)=5
\iff\lambda=\frac{-2x_0+3y_0+5}{13}.$$

$p(M_0)$ a pour coordonnées $\big(x_0+2\frac{-2x_0+3y_0+5}{13},y_0-3\frac{-2x_0+3y_0+5}{13}\big)$ ou encore
$p(M_0)= \big(\frac{9x_0+6y_0+10}{13},\frac{6x_0+4y_0-15}{13}\big)$.

Le symétrique orthogonal $s(M_0)$ vérifie~:~$s(M_0)=M_0+2\overrightarrow{M_0p(M_0)}$
 (car $p(M_0)$ est le milieu du segment $[M_0,s(M_0)]$)
autrement dit $s(M_0)=M_0+2\lambda . \vec n$ (pour le $\lambda$ obtenu ci-dessus).

Ses coordonnées sont donc 
$s(M_0) = \big(x_0+4\frac{-2x_0+3y_0+5}{13},y_0-6\frac{-2x_0+3y_0+5}{13}\big)$ ou encore
$\big(\frac{5x_0+12y_0+20}{13},\frac{12x_0-5y_0-30}{13}\big)$.
\fincorrection
\nocorrection
\nocorrection
\nocorrection
\nocorrection
\nocorrection
\nocorrection
\nocorrection
\nocorrection
\nocorrection
\nocorrection
\nocorrection
\nocorrection
\nocorrection
\nocorrection
\nocorrection
\nocorrection
\nocorrection
\nocorrection
\nocorrection
\nocorrection
\nocorrection
\nocorrection
\nocorrection
\nocorrection
\nocorrection
\nocorrection
\nocorrection
\nocorrection
\nocorrection
\nocorrection
\nocorrection
\nocorrection
\nocorrection
\nocorrection
\nocorrection
\nocorrection
\nocorrection
\nocorrection
\nocorrection
\nocorrection
\nocorrection
\nocorrection
\nocorrection
\nocorrection
\nocorrection
\nocorrection
\nocorrection
\nocorrection
\nocorrection
\nocorrection
\nocorrection
\nocorrection
\nocorrection
\nocorrection
\nocorrection
\nocorrection
\nocorrection
\nocorrection
\correction{004859}
 $   \det(\vec{MI},\vec{MJ}) +  \det(\vec{MJ},\vec{MK}) +  \det(\vec{MK},\vec{MI}) = \det(\vec{IJ},\vec{IK})$.
\fincorrection
\correction{004860}
Le système $\begin{cases}P \cr P' \cr Q \cr\end{cases}$ doit être lié.
\fincorrection
\correction{004861}
$11x + 2y - 13z = -4$.
\fincorrection
\correction{004862}
\begin{enumerate}
  \item $a=-4$
  \item $x-5y+3z = -9$
\end{enumerate}
 \fincorrection
\correction{004863}
\begin{enumerate}
  \item Le plan passant par $A$ et $D'$ n'est pas parallèle à $D''$.
  \item On doit pouvoir s'en sortir avec un repère adéquat \dots
\end{enumerate}
 \fincorrection
\nocorrection
\correction{004865}
Repère $(A,\vec{AB},\vec{AD})$.
\fincorrection
\correction{004866}
$\left(\begin{smallmatrix} 0 \cr 4/3 \cr\end{smallmatrix}\right)$.
\fincorrection
\correction{004867}
$\dim({\cal H}) = \dim({\cal F}) + \dim({\cal G})
                - \dim({\cal F} \cap {\cal G})$ si
${\cal F} \cap {\cal G} \ne \varnothing$.

$\dim({\cal H}) = \dim({\cal F}) + \dim({\cal G})
                - \dim({\cal F} \cap {\cal G}) + 1$
sinon.
\fincorrection
\nocorrection
\correction{004869}
\begin{enumerate}
  \item $ux+vy+h = \lambda^2(ua+v) + \lambda(u+va) + h$.
             Ceci est nul pour tout $\lambda$ si et seulement si
             $ua+v=0$, $u+va=0$, $h=0$ soit
             $(u,v,a,h) = (u,-u,1,0)$ ou $(u,v,a,h) = (u,u,-1,0)$.
  \item $x-y+(\lambda-\lambda^2)(z-1) = 0$ et $x+y-(\lambda+\lambda^2)(z+1) = 0$.
  \item 
\end{enumerate}
 \fincorrection
\correction{005505}
Puisque $\left|
\begin{array}{cc}
1&2\\
2&3
\end{array}\right|=-1$, on choisit d'exprimer $x$ et $z$ en fonction de $y$.
Soit $M(x,y,z)\in\Rr^3$. D'après les formules de \textsc{Cramer}, on a
\begin{align*}\ensuremath
M\in(D)&\Leftrightarrow\left\{
\begin{array}{l}
x-y+2z+7=0\\
2x+2y+3z-5=0
\end{array}
\right.\Leftrightarrow\left\{
\begin{array}{l}
x+2z=y-7\\
-2x-3z=2y-5
\end{array}
\right.\\
 &\Leftrightarrow x=\frac{1}{1}\left|
\begin{array}{cc}
y-7&2\\
2y-5&-3
\end{array}
\right|\;\text{et}\;z=\frac{1}{1}\left|
\begin{array}{cc}
1&y-7\\
-2&2y-5
\end{array}
\right|\\
 &\Leftrightarrow\left\{\begin{array}{l}
 x=31-7y\\
 z=-19+4y
 \end{array}\right..
\end{align*}

\begin{center}
\shadowbox{
$(D)$ est la droite passant par $A(31,0,-19)$ dirigée par le vecteur $u(-7,1,4)$.
}
\end{center}
\fincorrection
\correction{005506}
Soit $M(2+\lambda,3-\lambda,7)$, $\lambda\in\Rr$, un point quelconque de $(D)$.

\begin{center}
$M\in(P)\Leftrightarrow(2+\lambda)+3(3-\lambda)-5\times7+2=0\Leftrightarrow \lambda=12$.
\end{center}
$(P)\cap(D)$ est donc un singleton. Pour $\lambda=12$, on obtient les coordonnées du point d'intersection

\begin{center}
\shadowbox{
$(P)\cap(D)=\{(14,-9,7)\}$.
}
\end{center}
\fincorrection
\correction{005507}
\textbullet~\textbf{Repère de $(D)$.}

\begin{center}
$\left\{
\begin{array}{l}
x+2=-2z\\
y=3x+z
\end{array}
\right.\Leftrightarrow\left\{
\begin{array}{l}
x=-2-2z\\
y=3(-2-2z)+z
\end{array}
\right.\Leftrightarrow\left\{
\begin{array}{l}
x=-2-2z\\
y=-6-5z
\end{array}
\right.$
\end{center}

$(D)$ est la droite passant par $A(0,-1,-1)$ et dirigée par $u(2,5,-1)$.
 \textbullet~\textbf{Repère de $(D')$.}

\begin{center}$\left\{
\begin{array}{l}
x+y+z=1\\
2x+y-z=a
\end{array}
\right.\Leftrightarrow\left\{
\begin{array}{l}
-x-y=z-1\\
2x+y=z+a
\end{array}
\right.\Leftrightarrow\left\{
\begin{array}{l}
x=2z+a-1\\
y=2-a-3z
\end{array}
\right.$
\end{center}
$(D')$ est la droite passant par $A'(a-1,2-a,0)$ et dirigée par $u'(2,-3,1)$.
\textbullet~Déjà $u$ et $u'$ ne sont pas colinéaires et donc $(D)$ et $(D')$ sont ou bien sécantes en un point et dans ce cas coplanaires ou bien non coplanaires.
\textbullet~Le plan $(P)$ contenant $(D)$ et parallèle à $(D')$ est le plan de repère $(A,u,u')$. Déterminons une équation de ce plan.

\begin{center}
$M(x,y,z)\in(P)\Leftrightarrow
\left|
\begin{array}{ccc}
x&2&2\\
y+1&5&-3\\
z+1&-1&1
\end{array}\right|=0\Leftrightarrow 2x-4(y+1)-16(z+1)=0\Leftrightarrow-x+2y+8z=-10$.
\end{center}
\textbullet~Enfin, $(D)$ et $(D')$ sont coplanaires si et seulement si $(D')$ est contenue dans $(P)$. Comme $(D')$ est déjà parallèle à $(P)$, on a

\begin{center}
$(D)$ et $(D')$ coplanaires $\Leftrightarrow A'\in(P)\Leftrightarrow-(a-1)+2(2-a)=-10\Leftrightarrow a=\frac{5}{3}$.
\end{center}

\begin{center}
\shadowbox{
\begin{tabular}{c}
$(D)$ et $(D')$ sont coplanaires si et seulement si $a=\frac{5}{3}$ et dans ce cas, une équation du plan contenant $(D)$\\
 et $(D')$ est $-x+2y+8z=-10$.
 \end{tabular}
}
\end{center}
\fincorrection
\correction{005508}
Puisque $P$ parallèle à la droite $(Oy)$, le vecteur $\overrightarrow{j}=(0,1,0)$ est dans $\overrightarrow{P}$. De même, le vecteur $\overrightarrow{AB}=(-1,3,1)$ est dans $\overrightarrow{P}$.
$P$ est donc nécessairement le plan passant par $A(0,-1,2)$ et de vecteur normal $\overrightarrow{j}\wedge\overrightarrow{AB}=(1,0,1)$. Réciproquement, ce plan convient.
Une équation de $P$ est donc $(x-0)+(z-2)=0$ ou encore $x+z=2$.

\begin{center}
\shadowbox{
Une équation du plan parallèle à la droite $(Oy)$ et passant par $A(0,-1,2)$ et $B(-1,2,3)$ est $x+z=2$.
}
\end{center}
\fincorrection
\nocorrection
\nocorrection
\correction{004872}
\begin{enumerate}
  \item $\begin{cases}2x' = x - 2y - z + 1 \cr
                     2y' = -x -z + 1      \cr
                     2z' = -x -2y +z + 1  \cr\end{cases}$
  \item $\begin{cases}2x' = -5x - 3y + 2z - 3 \cr
                     2y' = 3x + y - 2z - 1   \cr
                     z' = -3x -3y +z -3      \cr\end{cases}$
\end{enumerate}
 \fincorrection
\correction{004873}
affinité de base ${\cal P} : x + 2y + z = 2$,
            de direction vect$(\vec e_1 - \vec e_2 + 2\vec e_3)$,
            de rapport 3.
\fincorrection
\correction{004874}
si trois points sont non alignés, $ABCD$ doit être un parallélogramme.

si deux points sont distincts et $A,B,C,D$ sont alignés, on doit avoir
$A=C$, $B=D$.
\fincorrection
\nocorrection
\correction{004876}
affinité de rapport $\lambda\mu$ si $\lambda\mu \ne 1$, transvection ou $\mathrm{id}$ sinon.
\fincorrection
\nocorrection
\correction{004878}
5.     ${\cal B} = \{ 3(x+y+z) = 1 \}$,\quad
    $\vec {\cal F} = \text{vect}(\vec e_1 + \vec e_2 + \vec e_3)$,\quad
    $9\vec u = \vec e_1 + 4\vec e_2 - 5\vec e_3$.
 \fincorrection
\correction{004879}
oui ssi $\vec{P'Q'}$ est colinéaire à $\vec{PQ}$. Dans ce cas, $f$ est unique.
\fincorrection
\correction{004880}
Il existe une homothétie de centre $O$ transformant $A$ en $A'$, $B$ en $B'$,
et $C$ en $C'$, et l'homothétie de centre $G,-\frac 12$ transforme $A$ en
$\alpha$, $B$ en $\beta$, $C$ en $\gamma$.
\fincorrection
\nocorrection
\nocorrection
\nocorrection
\correction{004885}
Repère $(O, \vec{OA}, \vec{OB}, \vec{OC})  \Rightarrow $
         les plans $(ABC)$ et $(A'B'C')$ sont parallèles.
\fincorrection
\correction{004886}
$A_i^{(k)} = \text{Bar}\left(
          A_j : \frac 1{2^k}\sum_{l\equiv j(\mathrm{mod}\, n)} C_k^{|l-i|}\right)$.
\fincorrection
\correction{004887}
$\vec{ GA_i^{(k+1)} } = -\frac 1{n-1} \vec{ GA_i^{(k)} }$.
\fincorrection
\correction{004888}
$A_k = \text{Bar}(A_0 : \alpha_k, A_1 : \beta_k, A_2 : \gamma_k)$
         où $\alpha_k$, $\beta_k$, $\gamma_k$ vérifient :
         $x_k = x_{k-1} + x_{k-2} + 2x_{k-3}$.

         Les racines de l'équation caractéristique sont $2,j,j^2$, donc
         $x_k \sim \lambda2^k$ avec $\lambda = \frac{x_0+x_1+x_2}7 = \frac 17$.

         Donc $\alpha_k \sim \beta_k \sim \gamma_k$, et $A_k \to G$,
         isobarycentre de $A_0A_1A_2$.
\fincorrection
\nocorrection
\nocorrection
\nocorrection
\correction{004892}
$\frac 17$.
\fincorrection
\correction{004893}
\begin{enumerate}
  \item 
  \item $N = \text{Bar}(A:1-\alpha, B:1-\beta, C:1-\gamma)$.
  \item 
  \item homothétie de centre $G$, de rapport $-\frac 12$.
\end{enumerate}
 \fincorrection
\correction{004894}
Repère $(A,\vec{AB},\vec{AC})$.
\fincorrection
\correction{004895}
$M_1 \mapsto M_4$ est affine, et échange $A$ et $B$. $\Rightarrow$ involutive.
\fincorrection
\correction{004896}
\begin{enumerate}
  \item 
  \item $G$ et les symétriques de $A,B,C$ par rapport aux milieux des côtés opposés.
\end{enumerate}
 \fincorrection
\nocorrection
\correction{004898}
\begin{enumerate}
  \item Soit $\alpha' = \overline{(\vec{A'A},\vec{A'B})}$ :
             $\frac{A'B}{\sin(\alpha/2)} = \frac{AB}{\sin(\alpha')}$,
             $\frac{A'C}{\sin(\alpha/2)} = \frac{AC}{\sin(\pi-\alpha')}$.
  \item $I=\text{Bar}(A:a,B:b,C:c)$.
\end{enumerate}
 \fincorrection
\nocorrection
\correction{004885}
Repère $(O, \vec{OA}, \vec{OB}, \vec{OC})  \Rightarrow $
         les plans $(ABC)$ et $(A'B'C')$ sont parallèles.
\fincorrection
\correction{004886}
$A_i^{(k)} = \text{Bar}\left(
          A_j : \frac 1{2^k}\sum_{l\equiv j(\mathrm{mod}\, n)} C_k^{|l-i|}\right)$.
\fincorrection
\correction{004887}
$\vec{ GA_i^{(k+1)} } = -\frac 1{n-1} \vec{ GA_i^{(k)} }$.
\fincorrection
\correction{004888}
$A_k = \text{Bar}(A_0 : \alpha_k, A_1 : \beta_k, A_2 : \gamma_k)$
         où $\alpha_k$, $\beta_k$, $\gamma_k$ vérifient :
         $x_k = x_{k-1} + x_{k-2} + 2x_{k-3}$.

         Les racines de l'équation caractéristique sont $2,j,j^2$, donc
         $x_k \sim \lambda2^k$ avec $\lambda = \frac{x_0+x_1+x_2}7 = \frac 17$.

         Donc $\alpha_k \sim \beta_k \sim \gamma_k$, et $A_k \to G$,
         isobarycentre de $A_0A_1A_2$.
\fincorrection
\nocorrection
\nocorrection
\nocorrection
\correction{004892}
$\frac 17$.
\fincorrection
\correction{004893}
\begin{enumerate}
  \item 
  \item $N = \text{Bar}(A:1-\alpha, B:1-\beta, C:1-\gamma)$.
  \item 
  \item homothétie de centre $G$, de rapport $-\frac 12$.
\end{enumerate}
 \fincorrection
\correction{004894}
Repère $(A,\vec{AB},\vec{AC})$.
\fincorrection
\correction{004895}
$M_1 \mapsto M_4$ est affine, et échange $A$ et $B$. $\Rightarrow$ involutive.
\fincorrection
\correction{004896}
\begin{enumerate}
  \item 
  \item $G$ et les symétriques de $A,B,C$ par rapport aux milieux des côtés opposés.
\end{enumerate}
 \fincorrection
\nocorrection
\correction{004898}
\begin{enumerate}
  \item Soit $\alpha' = \overline{(\vec{A'A},\vec{A'B})}$ :
             $\frac{A'B}{\sin(\alpha/2)} = \frac{AB}{\sin(\alpha')}$,
             $\frac{A'C}{\sin(\alpha/2)} = \frac{AC}{\sin(\pi-\alpha')}$.
  \item $I=\text{Bar}(A:a,B:b,C:c)$.
\end{enumerate}
 \fincorrection
\correction{004899}
\begin{enumerate}
  \item $-\frac{\tan \gamma}{\tan \beta}$.
  \item $H = \text{Bar}(A:\tan\alpha, B:\tan\beta, C:\tan\gamma)
                = \text{Bar}\left(A:\frac a{\cos\alpha},
                                  B:\frac b{\cos\beta},
                                 C:\frac c{\cos\gamma}\right)$.
\end{enumerate}
\fincorrection
\correction{005206}
Il revient au même de démontrer que, si le plan est rapporté à un repère orthonormé, il n'existe pas de triangle équilatéral dont les sommets ont pour coordonnées des nombres entiers.

Le plan est muni d'un repère orthonormé direct. Soient $A$, $B$ et $C$ trois points deux à deux distincts, non alignés et à coordonnées entières. On sait que $\cos(\overrightarrow{AB},\overrightarrow{AC})=\frac{\overrightarrow{AB}.\overrightarrow{AC}}
{||\overrightarrow{AB}||.||\overrightarrow{AC}||}$ et $\sin(\overrightarrow{AB},\overrightarrow{AC})=\frac{\mbox{det}(\overrightarrow{AB},\overrightarrow{AC})}
{||\overrightarrow{AB}||.||\overrightarrow{AC}||}$. 

Par suite, ou bien le triangle $(ABC)$ est rectangle en $A$ (et n'est donc pas équilatéral), ou bien 

$\tan(\overrightarrow{AB},\overrightarrow{AC})=\frac{\mbox{det}(\overrightarrow{AB},\overrightarrow{AC})}
{\overrightarrow{AB}.\overrightarrow{AC}}$. Dans ce dernier cas, $\tan(\overrightarrow{AB},\overrightarrow{AC})$ est un quotient de deux nombres entiers, et est donc un rationnel. Malheureusement, pour un triangle équilatéral, la tangente de chacun de ses angles vaut $\sqrt{3}$ qui n'est pas un rationnel.

Quand le repère est orthonormé, il n'existe pas de triangle équilatéral dont les sommets sont à coordonnées entières.
\fincorrection
\correction{007059} 

~\\
%\begin{center}
\includegraphics{../images/img007059-1}
%\end{center}

D'après les hypothèses, la droite $(BC)$ est perpendiculaire à la droite $(AH)$ et $ADA'$ est rectangle en $D$. Donc, les droites $(AH)$ et $(A'D)$ sont perpendiculaires. Or, deux droites perpendiculaires à une même droite sont parallèles.  
\fincorrection  
\correction{007060}  Le point $A$ est l'orthocentre de $PQB$. L'angle est donc droit.
\fincorrection  
\correction{007061} 
Un triangle dont un des côtés est un diamètre du cercle circonscrit est rectangle. On en déduit que les trois droites sont les hauteurs de $ABC$. Elles sont donc concourantes.

\begin{center}
\includegraphics{../images/img007061-1}
\end{center}

\fincorrection  
\correction{007062} 
Prendre un deuxième point $N$ sur le cercle de telle sorte que $(AM)$ et $(BN)$ se coupent en un point $C$. On peut alors construire l'orthocentre de $ABC$. La troisième hauteur fournit une droite orthogonale à $(AB)$, coupant le cercle en deux points $P$ et $Q$. On peut alors compléter $MPQ$ en un trapèze (isocèle) $MPQR$, en utilisant les diagonales d'un tel trapèze. La droite $(MR)$ est orthogonale à $(AB)$.
\fincorrection  
\correction{007063} 
\begin{enumerate}
\item Par définition, $(IA)$ et $(IB)$ sont tangentes au cercle $\mathcal C$, donc $IA=IB$. On a de même $IA=IC$ et donc $I$ est le milieu de $[BC]$. Le triangle $ABC$ est donc un triangle d'écolier et il est rectangle en $A$.
\end{enumerate}
\fincorrection  
\nocorrection     
\nocorrection     
\nocorrection     
\correction{007067}  
L'exercice se résout assez simplement en utilisant trois triangles isocèles, mais on peut remarquer que les trois tangentes sont les trois axes radicaux, qui s'intersectent tous trois  au centre radical des trois cercles. 
\fincorrection  
\correction{007068} 

Soit $D$ le symétrique de $B$ par rapport à $C$.

\begin{center}
\includegraphics{../images/img007068-2}
\end{center}

Alors $[AC]$ est une médiane de $ABD$ et $M$ est son centre de gravité. Comme $BA=2BC=BD$, le triangle $ABD$ est isocèle en $B$. On en déduit que la médiane issue de $B$ est également la bissectrice issue de $B$. Les angles $\widehat{ABM}$ et $\widehat{MBC}$ sont donc égaux.
\fincorrection  
\nocorrection     
\nocorrection     
\correction{007071}  
%%%%%%%%%%%%
%% FIGURE %%
%%%%%%%%%%%%

Pour montrer $CA=CP$, on va montrer que le triangle $CAP$ est donc isocèle en $C$.

On a les égalités d'angles:
\begin{align*}
\widehat{CAP} &= \widehat{OAB} \text{ car les angles sont opposés par le sommet}\\
&= \pi/2 - \widehat{ABO} \\
&= \pi/2 -\widehat{OPB}  \text{ car $POB$ est isocèle en $O$}\\
&= \widehat{APC}
\end{align*}

Le triangle $CAP$ est donc isocèle en $C$, et donc $CA=CP$.
\fincorrection  
\correction{007072} 
Deuxième indication: un des angles du triangle a une mesure $\geq \pi/3$, et un autre a une mesure $\leq \pi/3$.
\fincorrection  
\nocorrection     
\nocorrection     
\nocorrection     
\correction{007076} 

Pour le cercle de centre $P$, il suffit de montrer que $P$ est le centre du cercle inscrit du triangle $MAB$. Pour cela, en notant $C$ le projeté orthogonal de $P$ sur $(MA)$, il suffit de montrer que $PC=PH$, ou de montrer que $AC=AH$. Or, on a $AC=AH = \cos(\widehat{AMO}) / OA$.

D'autres solutions sont possibles, par exemple avec des homothéties.

\fincorrection  
\nocorrection
\nocorrection
\nocorrection
\nocorrection
\nocorrection
\nocorrection
\nocorrection
\nocorrection
\nocorrection
\correction{005207}
Soit $f$ la transformation considérée.
\begin{enumerate}
\item  $f$ est la translation de vecteur $\vec{u}(3,-1)$.
\item  $\omega=2\omega+3\Leftrightarrow\omega=-3$. $f$ est l'homothétie de rapport $2$ et de centre $\Omega(-3,0)$.
\item  $\omega=i\omega+1\Leftrightarrow\omega=\frac{1}{2}(1+i)$. Comme $i=e^{i\pi/2}$, $f$ est la rotation d'angle
$\frac{\pi}{2}$ et de centre $\Omega(\frac{1}{2},\frac{1}{2})$.
\item  $\omega=(1-i)\omega+2+i\Leftrightarrow\omega=1-2i$. Comme $1-i=\sqrt{2}e^{-i\pi/4}$, $f$ est la similitude de centre
$\Omega(1,-2)$, de rapport $\sqrt{2}$ et d'angle $-\frac{\pi}{4}$.
\end{enumerate}
\fincorrection
\nocorrection
\nocorrection
\nocorrection
\nocorrection
\nocorrection
\nocorrection
\nocorrection
\nocorrection
\nocorrection
\nocorrection
\nocorrection
\nocorrection
\correction{004949}
Cercle circonscrit au triangle $A'BC$ symétrique de $ABC$ par
         rapport à $(BC)$.
\fincorrection
\correction{004950}
$xAM^2 + yBM^2 + zCM^2 = r^2 - OM^2$ avec $\mathcal{C} = \mathcal{C}(O,r)$.\par
         $xyAB^2 + xzAC^2 + yzBC^2 = xAM^2 + yBM^2 + zCM^2$.
\fincorrection
\nocorrection
\correction{004952}
Nous savons que la distance d'un point $M_0(x_0,y_0)$ à une droite $D$ d'équation $ax+by+c=0$ est donnée par la formule
$d(M_0,D) = \frac{|ax_0+by_0+c_0|}{\sqrt{a^2+b^2}}$.

Pour une droite $D_\lambda$ la formule donne :
$d(M_0,D_\lambda)  = \frac{|(1-\lambda^2)x_0 + 2\lambda y_0 - (4\lambda + 2)|}{\sqrt{(1-\lambda^2)^2+4\lambda^2}}$.

\bigskip
\textbf{Analyse.}

On cherche un point $M_0=(x_0,y_0)$ tel que pour tout $\lambda$, $d(M_0,D_\lambda)=k$ où $k\in \Rr$ est une constante.

L'égalité $d(M_0,D_\lambda)^2=k^2$ conduit à 
$$\Big((1-\lambda^2)x_0 + 2\lambda y_0 - (4\lambda + 2)\Big)^2 
= k^2 \Big( (1-\lambda^2)^2+4\lambda^2 \Big)$$
pour tout $\lambda\in \Rr$.
Nos inconnues sont $x_0,y_0,k$.
On regarde l'égalité comme une égalité de deux polynômes en la variable $\lambda$.

Pour ne pas avoir à tout développer on raffine un peu :
on identifie les termes de plus haut degré en $\lambda^4$ :
$x_0^2 \lambda^4=k^2\lambda^4$ donc $x_0^2=k^2$.

En évaluant l'égalité pour $\lambda=0$ cela donne $(x_0-2)^2=k^2$.
On en déduit $(x_0-2)^2=x_0^2$ dont la seule solution est $x_0=1$.
Ainsi $k=1$ (car $k>0$).

L'égalité pour $\lambda=+1$ donne $(2y_0-6)^2=4k^2$
et pour $\lambda=-1$ donne $(-2y_0+2)^2=4k^2$.
La seule solution est  $y_0=2$. 

\bigskip
\textbf{Synthèse.}
Vérifions que le point de coordonnées $M_0=(1,2)$ est situé à une distance $k=1$ de
toutes les droites $D_\lambda$.

Pour $(x_0,y_0)=(1,2)$, on trouve :
$d(M_0,D_\lambda)  = \frac{|(1-\lambda^2) + 4\lambda - (4\lambda + 2)|}{\sqrt{(1-\lambda^2)^2+4\lambda^2}}
= \frac{|\lambda^2+1|}{\sqrt{(\lambda^2+1)^2}} =  \frac{|\lambda^2+1|}{|\lambda^2+1|}=1$.
Donc $M_0=(1,2)$ est bien équidistant de toutes les droites $D_\lambda$.
\fincorrection
\nocorrection
\correction{004954}
Décomposer les rotations en symétries.
\fincorrection
\correction{004955}
la symétrie centrale ($\alpha + \beta + \gamma = \pi$) autour de $K$, point
de contact du cercle inscrit et de $(AC)$.
\fincorrection
\nocorrection
\correction{004957}
$x' = \frac{ay^2}{\strut x^2+y^2-ax}$, $y' = \frac{-axy}{\strut x^2+y^2-ax}$,
$M'$ est bien défini ssi $M$ n'appartient pas au cercle de diamètre $[AO]$.

Soit $D$ le demi-disque supérieur de diamètre $[AO]$, $D$ est
caractérisé par les inégalités $x^2+y^2-ax < 0$, $y>0$ d'où $x'<0$ et
$y'>0$. La réciproque se traite (péniblement) en remarquant que seuls les points
de~$D$ ont une image dans ce quart de plan et que $f$ est quasi-involutive.
\fincorrection
\correction{005197}
$(D)$ est une droite de vecteur normal $(1,3)$. Le projeté orthogonal $p(M_0)$ de $M_0$ sur $(D)$ est de la forme $M_0+\lambda.\vec{n}$ où $\lambda$ est un réel à déterminer. Le point $M_0+\lambda.\vec{n}$ a pour coordonnées $(x_0+\lambda,y_0+3\lambda)$.

$$M_0+\lambda.\vec{n}\in(D)\Leftrightarrow(x_0+\lambda)+3(y_0+3\lambda)-5=0\Leftrightarrow\lambda=\frac{-x_0-3y_0+5}{10}.$$

$p(M_0)$ a pour coordonnées $(x_0+\frac{-x_0-3y_0+5}{10},y_0+3\frac{-x_0-3y_0+5}{10})$ ou encore $(\frac{9x_0-3y_0+5}{10},\frac{-3x_0+y_0+15}{10})$.

Le symétrique orthogonal $s(M_0)$ vérifie~:~$s(M_0)=M_0+2\overrightarrow{M_0p(M_0)}$.

Ses coordonnées sont donc 
$(x_0+2(\frac{9x_0-3y_0+5}{10}-x_0),y0+2(\frac{-3x_0+y_0+15}{10}-y_0)$ ou encore

$(\frac{4x_0-3y_0+5}{5},\frac{-3x_0-4y_0+15}{5})$.

(Remarque. Si on n'avait pas déjà $p(M_0)$ on aurait cherché le symétrique sous la forme $M_0+\lambda.\vec{n}$, $\lambda$ étant entièrement déterminé par la condition~:~le milieu du segment $[M_0,s(M_0)]$ appartient à $(D)$.)
\fincorrection
\correction{005198}
 Puisque $(ABDC)$ un parallélogramme, $\overrightarrow{AD}=\overrightarrow{AB}+\overrightarrow{AC}$. Les coordonnées de $D$ dans le repère $(A,\overrightarrow{AB},\overrightarrow{AC})$ sont donc $(1,1)$.
\fincorrection
\correction{005201}
\begin{enumerate}
\item 
Soient $k$ et $k'$ deux réels non nuls, $\Omega$ et $\Omega'$ deux points (pas nécessairement distincts), puis $h$ (resp.$h'$) l'homothétie de centre $\Omega$ (resp. $\Omega'$) et de rapport $k$ (resp. $k'$).

Soient $M$ un point du plan, puis $M'=h(M)$ et $M''=h'(M')$.

\begin{align*}
M''&=\Omega'+k'\overrightarrow{\Omega'M'}=\Omega'+k'(\overrightarrow{\Omega'\Omega}+\overrightarrow{\Omega M'})
=\Omega'+k'\overrightarrow{\Omega'\Omega}+kk'\overrightarrow{\Omega M}\;(*)
\end{align*}

Chechons alors les points invariants par $h'\circ h$.

\begin{align*}
h'\circ h(M)=M&\Leftrightarrow\Omega'+k'\overrightarrow{\Omega'\Omega}+kk'\overrightarrow{\Omega M}=M
\Leftrightarrow-\overrightarrow{\Omega'M}+k'\overrightarrow{\Omega'\Omega}+kk'\overrightarrow{\Omega M}=\overrightarrow{0}\\
 &\Leftrightarrow(kk'-1)\overrightarrow{\Omega M}=(k'-1)\overrightarrow{\Omega\Omega'}\;(**)
\end{align*}

\begin{itemize}
\item[\textbf{1er cas.}] Si $kk'\neq1$, $(**)\Leftrightarrow\overrightarrow{\Omega M}=\frac{k'-1}{kk'-1}\overrightarrow{\Omega\Omega'}$, ce qui signifie que l'équation $(**)$ a une et eune seule solution que l'on note $\Omega''$, ou encore $h'\circ h$ a un et un seule point invariant, le point $\Omega''$ tel que 
$\Omega''=\Omega'+k'\overrightarrow{\Omega'\Omega}+kk'\overrightarrow{\Omega\Omega''}$.

Mais alors, l'égalié $(*)$ s'écrit pour tout point $M$

$$M''=\Omega'+k'\overrightarrow{\Omega'\Omega}+kk'\overrightarrow{\Omega M}
=\Omega'+k'\overrightarrow{\Omega'\Omega}+kk'\overrightarrow{\Omega\Omega''}+kk'\overrightarrow{\Omega''M}
=\Omega''+kk'\overrightarrow{\Omega''M}.$$

$h'\circ h$ est donc l'homothétie de rapport $kk'$ et de centre $\Omega''$. On doit noter que le centre $\Omega''$ est sur la droite $(\Omega\Omega')$.

\shadowbox{Si $kk'\neq1$,\;$h'\circ h$\; est une homothétie de rapport\;$kk'$.}

\item[\textbf{2ème cas.}] Si $kk'=1$, l'égalité $(*)$ s'écrit pour tout point $M$, $M''=\Omega'+k'\overrightarrow{\Omega'\Omega}+\overrightarrow{\Omega M}$ et donc

$$\overrightarrow{MM''}=\Omega'+k'(\Omega-\Omega')+(M-\Omega)-M=(1-k')\overrightarrow{\Omega\Omega'}.$$

Dans ce cas, $h'\circ h$ est la translation de vecteur $(1-k')\overrightarrow{\Omega\Omega'}$.
\end{itemize}

En résumé, \textbf{la comoposée de deux homothéties de rapport respectifs} $\bf{k}$ \textbf{et} $\bf{k'}$ \textbf{tous deux non nuls est une homothétie de rapport} $\bf{kk'}$ \textbf{si} $\bf{kk'}\neq1$ \textbf{et une tranlation si} $\bf{kk'=1}$ (ce résultat est à connaître).

\item  C'est un cas particulier de la question précédente. Une symétrie centrale est une homothétie de rapport $-1$. Puisque $(-1)(-1)=1$, $s'\circ s$ est une translation. Son vecteur est $\overrightarrow{\Omega s'\circ s(\Omega)}=\overrightarrow{\Omega s'(\Omega)}=2\overrightarrow{\Omega \Omega'}$.

$$\includegraphics{../images/img005201-1}$$

\textbf{La composée de deux symétries centrales est une translation.}

\item  Soit $\Omega'$ le point tel que $\overrightarrow{u}=2\overrightarrow{\Omega'\Omega}$, c'est-à-dire 
$\Omega'=\Omega-\frac{1}{2}\overrightarrow{u}$. Soit $s'$ la symétrie centrale de centre $\Omega'$. D'après 2), $s\circ s'$ est la translation de vecteur $2\overrightarrow{\Omega'\Omega}=\overrightarrow{u}$. Par suite, $s\circ t=s\circ s\circ s'=s'$.

\textbf{La composée d'une symétrie centrale et d'une translation est une symétrie centrale.}

\end{enumerate}
\fincorrection
\nocorrection
\nocorrection  
\correction{007049} 
On relie les points $A$ et $B$. On construit le milieu $I$ de $[AB]$. On a donc $AI = \frac{1}{2}AB$. On trace le cercle $\mathcal{C}$ de centre $A$ et de rayon $AI$. Ceci fait, On trace le cercle $\mathcal{C}'$ centr\'e en $I$ de diam\`etre $AB$. Nommons $C$ l'un des deux points d'intersection de $\mathcal{C}$ et $\mathcal{C}'$. Il s'agit de l'une des deux solutions possibles.

En effet, par construction $AC = \frac{1}{2} AB$, puisque $C$ appartient \`a $\mathcal{C}$. En outre $ACB$ est rectangle en $C$. En effet, $ACB$ est inscrit dans le cercle $\mathcal{C}'$ et $AB$ est un diam\`etre de $\mathcal{C}'$.
\fincorrection  
\correction{007050} 
Comme par hypoth\`ese $DA = DC = DB$, les points $A$, $B$ et $C$ sont situés sur le cercle de centre $A$ et de diam\`etre $AB$. Donc, $ACB$ est rectangle en $C$. 

\begin{center}
\begin{tikzpicture}[line cap=round,line join=round,x=1.0cm,y=1.0cm]
\clip(-1.38,-0.88) rectangle (9.2,5.2);
\fill[fill opacity=0.1] (-0.22,0.26) -- (3.88,0.26) -- (1.83,3.81) -- cycle;
\draw [shift={(1.83,3.81)},fill opacity=0.1] (0,0) -- (-60.:0.6) arc (-60.:-30.:0.6) -- cycle;
\draw [shift={(7.98,0.26)},fill opacity=0.1] (0,0) -- (150.:0.6) arc (150.:180.:0.6) -- cycle;
\draw  (-0.22,0.26)-- (3.88,0.26);
\draw  (3.88,0.26)-- (1.83,3.81);
\draw  (1.83,3.81)-- (-0.22,0.26);
\draw (7.98,0.26)-- (1.83,3.81);
\draw (7.98,0.26)-- (3.88,0.26);
\draw [shift={(3.88,0.26)}] plot[domain=0.:3.141592653589793,variable=\t]({1.*4.1*cos(\t r)+0.*4.1*sin(\t r)},{0.*4.1*cos(\t r)+1.*4.1*sin(\t r)});
\begin{scriptsize}
\draw (-0.5,0.62) node {$A$};
\draw (4.02,0.62) node {$D$};
\draw (1.98,4.18) node {$C$};
\draw (8.12,0.62) node {$B$};
\end{scriptsize}
\end{tikzpicture}
\end{center}

Puisque $ADC$ est équilatéral, $\widehat{DAC} = \widehat{BAC}= 60^{\circ}$. Comme $ACB$ est rectangle en $C$, donc $\widehat{ACB}=90^{\circ}$.

De la relation $180^{\circ} = \widehat{BAC} + \widehat{ACB} + \widehat{CBA}$ on tire $\widehat{CBA} = 30^{\circ}$.    
\fincorrection  
\nocorrection     
\correction{007052}  
\begin{enumerate}
\item 
Soit $G$ l'isobarycentre. On a
\[
G =\frac{A+B+C+D}{4} = \frac{\frac{A+B}{2} + \frac{C+D}{2}}{2}
\]
Donc si $I$ et $J$ sont les milieux de $[AB]$ et $[CD]$, on a montré que $G$ est le milieu de $[IJ]$.
On construit $I$ et $J$, puis leur milieu $G$.
\item Soit $H$ le deuxième barycentre. Par définition, 
\[
H = \frac{A+B+3C+3D}{8} = \frac{1}{4}\left(\frac{A+B}{2}\right)
+ \frac{3}{4}\left(\frac{C+D}{2}\right)
\]
On construit donc $I$ et $J$ les milieux de $[AB]$ et $[CD]$, puis on construit le barycentre de $(I,1/4)$ et $(J,3/4)$.
\item Même technique : tracer les milieux $I$ et $J$ de $[AB]$ et $[CD]$, puis tracer le milieu $K$ de $[IJ]$. L'isobarycentre du pentagone est le barycentre de $(K,4/5)$ et $(E,1/5)$.
\end{enumerate}
\fincorrection  
\nocorrection     
\correction{007054} 
\begin{enumerate}
\item Les droites sont perpendiculaires.
\item Pour les tangentes communes extérieures, le cercle de centre $O'$ et de rayon $r'-r$ intersecte le cercle de diamètre $[OO']$ en deux points $C$ et $D$. Les droites $(O'C)$ et $(O'D)$ coupent $\mathcal C'$ en deux points $A'$ et $B'$. Les tangentes extérieures sont les parallèles à $(OC)$ et $(OD)$ passant par $A'$ et $B'$. Pour les tangentes intérieures, utiliser le cercle de centre $O'$ et de rayon $r'+r$. 
\end{enumerate}
\fincorrection  
\nocorrection    
\nocorrection
\nocorrection  
\correction{007058} Tracer le lieu des points à distance $R$ des cercles et droites en présence. Leurs éventuels points d'intersection fournissent des solutions.
\fincorrection  
\nocorrection
\nocorrection
\nocorrection
\correction{007080} 

Si $n$ est impair, la composée définie dans l'indication est une symétrie centrale, et son centre (que l'on peut construire en considérant les images de plusieurs points) est un sommet du polygone. On récupère ensuite les autres sommets en appliquant les autres symétries centrales les unes après les autres.


Remarque : cet exercice peut également se résoudre à l'aide de nombres complexes, en écrivant que le milieu de deux points a pour affixe $(a+b)/2$. On obtient un système linéaire, dont on discute l'existence de solutions. 
\fincorrection  
\correction{007081} 
On complète le trapèze rectangle $ABB'A'$ en un rectangle comme conseillé, en utiisant la symétrie de centre $I$. 
\begin{center}
\includegraphics{../images/img007081-2}
\end{center}
Le résultat demandé est alors une conséquence du fait que les diagonales d'un rectangle sont égales et se coupent en leur milieu. 

%On peut aussi utiliser le théorème des milieux.
\fincorrection
\correction{007082}
\fincorrection
\correction{007083}  
Tracer des rayons de $\mathcal C$ et $\mathcal C'$ parallèles entre eux.
\fincorrection   
\nocorrection
\correction{007085} 

Analyse. Traçons comme suggéré une figure avec le carré déjà construit : on trace un carré puis on trace un triangle adéquat autour. On constate qu'un des côtés du carré, notons-le $[IJ]$, est parallèle à $[BC]$. Il y a une homothétie $h$ de centre $A$ qui envoie $[IJ]$ sur $[BC]$. Alors, l'image du carré $IJKL$ par $h$ est un carré dont un des côtés est $[BC]$. Notons $BCDE$ ce carré et traçons-le. On constate que $h(K)=D$ et $h(L)=E$, c'est-à-dire $K=h^{-1}(D)$ et $L = h^{-1}(E)$. Il ne reste plus qu'à faire la synthèse.


\fincorrection  
\correction{007086} 
\begin{enumerate}
\item Le triangle des milieux est l'image de $ABC$ par l'homothétie de centre $G$ et de rapport $-1/2$. On en déduit que $\mathcal C$ est l'image du cercle circonscrit par cette homothétie, et donc que son centre est l'image de $\Omega$ par cette homothétie : il est donc sur la droite $(G\Omega)$.

\item Montrons que $\mathcal C'$ est l'image de $\mathcal C$ par l'homothétie de centre $H$ et de rapport $1/2$. Considérons la composition de l'homothétie de centre $G$ et de rapport $-1/2$ avec l'homothétie de rapport $-1$ et de centre $J$. C'est une homothétie de rapport $1/2$ qui envoie $\mathcal C$ sur $\mathcal C'$. Comme elle envoie de plus $\Omega$ sur $J$, son centre est le point $M$ tel que $\overrightarrow{MJ}=\frac12\overrightarrow{M\Omega}$. Or on sait déjà, par exemple en considérant l'homothétie de centre $G$ et de rapport $-\frac12$, que $\overrightarrow{G\Omega}=-\frac12\overrightarrow{GH}$. On en déduit que $M=H$.
\end{enumerate}
\fincorrection  
\nocorrection
\nocorrection     
\nocorrection
\nocorrection      
\nocorrection
\nocorrection
\nocorrection
\correction{007094}     
Soit $O$ le pt d'intersection. On note $\phi$ l'homothétie qui envoie $A$ sur $B$, et $\psi$ celle qui envoie $B$ sur $C$. Alors $\phi\psi = \psi\phi$. L'image de $A$ est $C$ et l'image de $A'$ est $C'$, d'où le  parallélisme demandé.

Si les droites sont parallèles, on remplace les homothéties par des translations.
\fincorrection
\nocorrection 
\correction{007096} 
La partie linéaire de $f\circ g^{-1}$ est l'identité. Donc c'est une translation. En regardant l'image d'un des centres, on trouve que le vecteur est $(1-\lambda)O1O2$.
\fincorrection  
\correction{007097} 
Appliquer la translation au cercle. (Si on n'a pas donné le centre du cercle, commencer par construire le centre.)

Les points d'intersection des deux cercles fournissent les (ou la) solutions du problème.
\fincorrection  
\correction{007098}
Après avoir suivi l'indication, translater ce cercle de manière à ce qu'il contienne le point $A$.
\fincorrection
\correction{007099}
Commencer par tracer la bissectrice, puis $(OA)$. Ensuite, tracer un cercle quelconque tangent aux deux droites (dans le même secteur angulaire), et utiliser une homothétie. 

Note : les exercices faisant intervenir des homothéties se résolvent plus facilement en \og partant de la fin\fg, c'est-à-dire en procédant par analyse-synthèse et en faisant une figure approximative de ce que sera la solution.
\fincorrection 
\nocorrection     
\correction{007101} 
On a et $MM' = 2IJ$ et d'autre part $I$ et $J$ sont les projetés de $O$ et $O'$ sur la droite $D$, donc $IJ \leq OO'$ avec égalité ssi $(OO') // (IJ)$, dans ce cas le maximum est donc $2OO'$.
\fincorrection  
\correction{007102} 

En suivant l'indication on comprend qu'un carré inscrit dans un cercle inscrit au carré initial $ABCD$ a une aire deux fois plus petite que $ABCD$.

\begin{center}
\includegraphics{../images/img007102-2}
\end{center}


On en déduit une construction du petit carré : ses sommets sont les points d'intersection des diagonales avec le cercle inscrit.
\fincorrection  
\correction{007103}
On construit l'image du carré par une rotation d'angle $\pi/4$. Les points d'intersection des deux carrés forment un octogone régulier qui répond à la question.
\fincorrection
\correction{007104}
Voici une solution dans une configuration \og générique\fg.

Commençons par analyser le problème. On suppose que $ABCD$ est direct et que $E\in [AB]$, $F \in [BC]$ etc.
Considérons la rotation de centre $O$ (le centre du carré), et d'angle $\pi/4$. L'image du segment $[EG]$ est un segment $[E'G']$, que l'on suppose distinct de $[FH]$. En fait, on suppose pour simplifier $E'\neq F$. Considérons alors la translation de vecteur $\overrightarrow{E'F}$. Elle envoie $G'$ sur un point $G''$ appartenant à la droite $(DA)$. Si on suppose que ce point est différent de $H$, alors on a $(DA) = (G''H)$. 

Voici comment obtenir le point $G''$. On construit la perpendiculaire à $(EG)$ passant par $F$, et sur cette droite, on place le point $G''$ tel que $FG''=EG$ et $(\overrightarrow{EG},\overrightarrow{FG''}) = \pi/2$. Ceci permet de tracer la droite $(HG'')$ c'est-à-dire $(AD)$.

\begin{center}
\includegraphics{../images/img007104-1}
\end{center}

On projette ensuite les points $E$ et $G$ sur cette droite, ce qui donne $A$ et $D$. On peut ensuite terminer la construction du carré.

Il existe d'autres solutions qui utilisent le théorème de l'angle au centre.
\fincorrection 
\correction{007105} 
La rotation de centre $O$ (le centre du carré) et d'angle $\pi/2$ envoie le triangle $DAM$ sur $ABN$. On en déduit que $(DM)\bot (AN)$ et donc que $(AN)$ est une hauteur de $DMN$. On procède de même pour la deuxième hauteur.
\fincorrection  
\correction{007106} 
Tracer une figure en prolongeant les segments $[OP]$ et $[OR]$, et considérer une rotation de centre $O$ et d'angle $\pi/2$.
\fincorrection  
\nocorrection
\nocorrection
\nocorrection
\nocorrection
\nocorrection
\correction{007112}
Deux carrés sont toujours semblables, et une similitude qui n'est pas une translation admet toujours un point fixe, son centre.
\fincorrection
\nocorrection
\correction{007114} 
Les triangles $AA'I$ et $BB'C$ sont semblables, par une similitude d'angle $\pi/2$. L'image de $(AJ)$ par cette similitude est $(BI)$.
\fincorrection  
\nocorrection
\correction{007116} 
La composée est une translation, on en déduit que $\overrightarrow{HG}=\overrightarrow{EF}$.
\fincorrection  
\nocorrection
\correction{007118} 
Traçons la figure, où on a placé $I$ le milieu de $[AB]$, de telle sorte que $\frac12(OA,OB)=(OA,OI)$.

\begin{center}
\includegraphics{../images/img007118-2}
\end{center}


Les angles $(AO,\mathcal T)$ et $(AI,IO)$ sont droits.
On a d'une part :
\[ 0=(\mathcal T,\mathcal T) =  (\mathcal T,AI) +(AI,AO)+ \pi/2, \]
et d'autre part, dans le triangle $AIO$:
\[ 0=(AI,AO)+(IO,IA)+(OA,OI)=(AI,AO)+\pi/2+(OA,OI).\]
Finalement, on a donc:
\[ (\mathcal T,AB) = (\mathcal T, AI) = -(AI,AO)-\pi/2 = (OA,OI)=\frac{1}{2}(OA,OB),\]
ce qu'il fallait démontrer.$\qed$


\fincorrection  
\correction{007119} 
Par le théorème de l'angle inscrit, c'est un arc de cercle, dont le centre est sur la médiatrice de $[AB]$. 

Par le cas limite du théorème de l'angle inscrit, on sait aussi que si $\mathcal T$ est la tangente à ce cercle en $A$, alors $(\mathcal T,AB)=\alpha$.

On trace donc la droite $\mathcal T$ faisant un angle $\alpha$ avec $(AB)$ en $A$, puis la perpendiculaire à $\mathcal T$ passant par $A$. Cette droite coupe la médiatrice en un point $O$ qui est donc le centre du cercle recherché.
\fincorrection  
\correction{007120}  
Construisons un triangle $AIB$ isocèle rectangle en $I$ et le cercle de centre $I$ et de rayon $IA$. Ce cercle intersecte la médiatrice de $[AB]$ en un point $O$ qui vérifie $\widehat{AOB}=\pm \pi/4$, par le théorème de l'angle au centre. C'est donc le centre d'un octogone appuyé sur $[AB]$. En traçant le cercle de centre $O$ et de rayon $OA$, on peut terminer la construction de cet octogone.
\fincorrection  
\correction{007121}
Commençons par rappeler deux points:
\begin{enumerate}
\item dans un trapèze, deux angles non adjacents à une même base sont supplémentaires, puisque les deux bases sont parallèles.
\item un quadrilatère non croisé est inscriptible ssi les angles opposés sont supplémentaires.
\end{enumerate}

Un trapèze est isocèle ssi les angles adjacents à une même base sont égaux, donc (par le premier point ci-dessus) ssi les angles opposés sont supplémentaires, donc (par le deuxième point) ssi il est inscriptible.

\begin{center}
\includegraphics{../images/img007121-1}
\end{center}

\fincorrection
\nocorrection
\correction{007123} 
Traçons une figure. \emph{On marque dès à présent quelques égalités d'angles obtenues par le théorème de l'angle inscrit:}

\begin{center}
\includegraphics{../images/img007123-1}
\end{center}

\emph{Les égalités d'angles repérées sur la figure permettent de voir la solution, au moins dans la configuration particulière dessinée. On voit en effet que les angles $\widehat{ECD}$ et $\widehat{AEF}$ sont égaux. Attention toutefois, les angles géométriques sont trompeurs et les égalités que l'on voit sur une figure peuvent dépendre de la façon de tracer la figure. Sur la figure ci-dessous par exemple, les angles en question ne sont pas égaux mais supplémentaires.}

\begin{center}
\includegraphics{../images/img007123-2}
\end{center}


\emph{Il ne reste plus qu'à rédiger rigoureusement la solution  avec des angles de droites, en s'appuyant sur l'intuition donnée par la figure.}

Pour montrer que $(CD)$ et $(EF)$ sont parallèles, il suffit par exemple de montrer qu'elles forment le même angle avec la droite $(CA)$. Or on a la suite d'égalités d'angles de droites :

\begin{align*}
(CD,CA) &= (BD,BA) \text{ car $CDAB$ est inscriptible}\\
&= (BF,BA) \text{ car $(BD)=(BF)$}\\
&=(EF,EA) \text{ car $BFAE$ est inscriptible}\\
&=(EF,CA)  \text{ car $(EA)=(CA)$.}
\end{align*}

\fincorrection   
\correction{007124} 

Pour montrer le résultat, il suffit de montrer que $IBC$ et $JBC$ sont isocèles en $I$ et $J$. 

On commence par prouver le résultat pour $I$ :

\begin{center}
\includegraphics{../images/img007124-1}
\end{center}

Pour montrer que $BCI$ est isocèle en $I$, il suffit de montrer que $(BC,BI) = (CI,CB)$. Or, on a 
\begin{align*}
(BC,BI) &= (AC,AI) \text{ car $ABIC$ est inscriptible}\\
&= (AI,AB) \text{ car $(AI)$ est une bissectrice de $(AC)$ et $(AB)$}\\
&= (CI,CB) \text{ car $ABIC$ est inscriptible.}
\end{align*}

On remarque qu'en rédigeant avec des angles de droites, on n'a pas eu besoin (ni en fait la possibilité) de préciser si la bissectrice était intérieure ou extérieure, ce qui implique que la preuve sera la même pour $J$. Traçons juste une figure pour visualiser la deuxième situation.


\begin{center}
\includegraphics{../images/img007124-2}
\end{center}



\fincorrection  
\correction{007125}
Soit $\mathcal T$ la tangente commune  aux deux cercles.

\begin{center}
\includegraphics{../images/img007125-2}
\end{center}

Par le cas limite du théorème des angles inscrits, on a 
\[ (AB,AT) = (BT,\mathcal T)=(B'T,\mathcal T)=(A'B',A'T)\]

Comme $(AT) = (A'T)$, on en déduit que
\[ (AB,AT) = (A'B',AT),\]
et donc que $(AB)//(A'B')$. 

\underline{Autre preuve:} considérer une homothétie de centre $T$ qui envoie un cercle sur l'autre.

\fincorrection
\correction{007126} 



Soient $\mathcal C$ et $\mathcal C'$ les cercles circonscrits à $ARQ$ et $BPR$. 
Traitons le premier cas, celui où ils se coupent en $R$ et en un deuxième point $T$.

\begin{center}
\includegraphics{../images/img007126-1}
\end{center}

Il s'agit de montrer que $T, P, C, Q$ sont cocycliques.


Par le cours,  il suffit de montrer l'égalité d'angles de droites $(QT,QC)=(PT,PC)$. Or on a :
\begin{align*}
(QT,QC)&=(QT,QA) \text{ car $(QC)=(QA)$}\\
&=(RT,RA) \text{ car $AQTR$ est inscriptible} \\
&= (RT,RB) \text{ car $(RA)=(RB)$} \\
&=(PT,PB) \text{ car $PTRB$ est inscriptible} \\
&=(PT,PC) \text{ car $(PB)=(PC)$.}
\end{align*}

Attention, si on utilise des angles géométriques au lieu des angles de droites pour rédiger la solution, on peut être amené à distinguer plusieurs configurations possibles, par exemple celle-ci:
\begin{center}
\includegraphics{../images/img007126-2}
\end{center}

(Les angles $\widehat{BRT}$ et $\widehat{BPT}$ sont supplémentaires dans la première figure, et égaux dans la seconde.)

Le second cas, celui où les deux premiers cercles sont tangents en $R$, se traite à l'aide du cas limite du théorème de l'angle au centre.

\fincorrection  
\correction{007127}
Les deux parties ont la même aire.
\fincorrection
\correction{007128} 
Traçons une figure :

\begin{center}
\includegraphics{../images/img007128-1}
\end{center}


\emph{[Sur la figure, on voit que les angles $\widehat{GFD}$ et $\widehat{GED}$ sont supplémentaires, car $\widehat{GFD}=\widehat{BFA}$ et $\widehat{GED}=\widehat{GEB}+\widehat{BED} = \widehat{FBA}+\widehat{BAF}$. Il ne reste plus qu'à rédiger cette preuve un peu plus rigoureusement avec des angles de droites.]}

Montrons que $(FD,FG)=(ED,EG)$, ce qui prouve que $EDFG$ est inscriptible.

Tout d'abord, comme $(FD)=(FA)$ et $(FG)=(FB)$, on a 
\[(FD,FG)=(FA,FB).\]

Ensuite, la somme des angles du triangle $ABF$ vaut $\pi$, donc en termes d'angles de droites on a la relation 
$(FA,FB)+(AB,AF)+(BF,BA)=0$, c'est-à-dire:
\[
(FA,FB) = (AF,AB)+(BA,BF).
\]
Calculons chacun de ces deux angles. D'une part, on a :
\begin{align*}
(AF,AB) &= (AD,AB) \text{ car $(AD)=(AF)$}\\
&=(ED,EB) \text{ car $ABDE$ est inscriptible}.
\end{align*}
Et d'autre part :
\begin{align*}
(BA,BF) &= (BA,BC) \text{ car $(BF)=(BC)$} \\
&= (CB,CA) \text{ car $ABC$ est isocèle en $A$}\\
&= (EB,EA) \text{ car $ABCE$ est inscriptible.}
\end{align*}

Finalement, on obtient donc:
\begin{align*}
(FD,FG)&=(FA,FB) \\
&= (AF,AB)+(BA,BF) \\
&= (ED,EB) + (EB,EA)\\
&= (ED,EA)\\
&= (ED,EG) \text{ car $(EG)=(EA)$,}
\end{align*}
ce qu'il fallait démontrer.
\begin{center}
\includegraphics{../images/img007128-1}
\includegraphics{../images/img007128-2}
\end{center}

\fincorrection  
\correction{007129}  Traçons une figure. \emph{[Le fait de marquer toutes les égalités d'angles disponibles donne le résultat. Sur la figure, on ne marque que celles utilisées dans la rédaction proposée.]}
\begin{center}
\includegraphics{../images/img007129-1}
\end{center}

Montrons que $(PA,PC)=(DQ,BQ)$. On a:
\begin{align*}
(PA,PC)&= (PA,PQ)+(PQ,PC) \\
&= (BA,BQ)+(DQ,DC) \text{ par cocyclicité dans chaque cercle}\\
&= (BA,BQ)+(DQ,BA) \text{ car $(DC)=(BA)$}\\
&= (DQ,BQ).
\end{align*}

\fincorrection  
\correction{007130} 
\begin{enumerate}
\item Soit $\theta$ l'angle de la similitude. La partie linéaire $\overrightarrow{\phi}$ de $\phi$ est une similitude vectorielle d'angle $\theta$. On a donc $\theta =  \widehat{(\overrightarrow{AB},\overrightarrow \phi\overrightarrow{AB}}= \widehat{(\overrightarrow{AB},\overrightarrow{CD}}$.
\item On a $\theta = \widehat{(\overrightarrow{AB},\overrightarrow{CD}} = \widehat{(\overrightarrow{AQ},\overrightarrow{CQ}}$ et d'autre part $\theta:\widehat{(\overrightarrow{OA},\overrightarrow{O \phi(A)}}=\widehat{(\overrightarrow{OA},\overrightarrow{OC}}$. Donc $AQCO$ est inscriptible.
\item De même $BQDO$ est inscriptible. Donc $O$ appartient à l'intersection des cercles circonscrits à $ACQ$ et $BDQ$.
\item Si les segments sont parallèles, la similitude est une homothétie. S'ils sont de plus de même longueur, c'est une translation ou une symétrie centrale.
\end{enumerate}
\fincorrection  
\correction{007131}

Le quadrilatère $ABA'B'$ est inscriptible dans un cercle de diamètre $[AB]$. En effet, les triangles $ABA'$ et $ABB'$ sont par définition rectangles en $A'$ et $B'$, et ont même hypoténuse $[AB]$.

De même, les quadrilatères $BCB'C'$ et $CAC'A'$ sont inscriptibles dans des cercles de diamètre $[BC]$ et $[CA]$.

\begin{center}
\includegraphics{../images/img007131-1}
\end{center}

Montrons que la hauteur $(BB')$ est une bissectrice des droites $(B'C')$ et $(B'A')$. Pour cela, on montre que $(B'C',B'B)=(B'B,B'A')$.

On a :
\begin{align*}
(B'C',B'B)
&= (CC',CB) \text{ (car $BCB'C'$ est inscriptible)}\\
&= (CC',CA') \text{ (mêmes droites)}\\
&= (AC'AA') \text{ (car $ACA'C'$ est inscriptible)}\\
&= (AB,AA') \text{ (mêmes droites)}\\
&= (B'B,B'A') \text{ (car $ABA'B'$ est inscriptible)}\\
\end{align*}

\fincorrection 
\correction{007132} 



Le théorème de l'angle inscrit sur l'arc $BC$ donne l'égalité $\widehat{BAC} = \widehat{BDC}$. La considération des autres arcs $CD$, $DA$ et $AB$ donne les égalités $\widehat{DAC} = \widehat{DBC}$, $\widehat{DBA} = \widehat{DCA}$ et $\widehat{ACB} = \widehat{ADB}$.

\begin{center}
\includegraphics{../images/img007132-1}
\end{center}
 
Suivons l'énoncé. Soit $K$ le point de la diagonale $[AC]$ tel que $\widehat{ABK} = \widehat{DBC}$.

Alors,  $ABK$ et $DBC$ sont semblables (première figure) car leurs angles sont égaux :  $\widehat{BAK} = \widehat{BDC}$ et $\widehat{ABK} = \widehat{DBC}$. On a une similitude de rapport 
\[\frac{BD}{BA} = \frac{BC}{BK} = \frac{DC}{AK},\]
envoyant $ABK$ sur $DBC$.

De même, $ABD$ et $KBC$ sont semblables (deuxième figure) car leurs angles sont égaux : $\widehat{ADB} = \widehat{ACB} = \widehat{KCB}$ et $\widehat{ABD} = \widehat{KBC}$. La similitude envoyant $ABD$ sur $KBC$  a pour rapport \[\frac{BK}{BA}=\frac{BC}{BD} = \frac{KC}{AD}.\]

 


On a $AC = AK + KC$, donc $AC\cdot BD  = AK \cdot BD + KC\cdot BD$, 
or, $AK\cdot BD = AB\cdot CD$ (première similitude), et $KC \cdot BD = AD\cdot BC$ (seconde similitude), d'où le résultat:
\[ AC\cdot BD = AB\cdot CD + AD\cdot BC.\]
\fincorrection  
\nocorrection
\nocorrection    
\nocorrection
\correction{007136}
La somme des angles d'un quadrilatère convexe vaut $2\pi$ :
\begin{align*}
2\pi &=
\widehat{ABC}+\widehat{BCD}+\widehat{CDA}+\widehat{DAB}\\
&= 2\widehat{ABI}
+2\widehat{KCD}+2\widehat{CDK}+2\widehat{IAB}
\end{align*}
d'où \[\widehat{ABI}
+\widehat{KCD}+\widehat{CDK}+\widehat{IAB}=\pi,\]
autrement dit la somme des demi-angles vaut $\pi$.

On termine alors la preuve en utilisant le critère de cocyclicité.

\fincorrection
\correction{007137} 
On trouve un parallélogramme dont les angles sont les supplémentaires de ceux de $ABCD$.
\fincorrection  
\nocorrection
\correction{007139}

Si $ABC$ est rectangle, l'orthocentre coïncide avec un des sommets et la vérification de l'assertion est relativement facile. Dans la suite on suppose qu'on n'est pas dans ce cas.

Par définition, $H'$ est le symétrique de $H$ par rapport à $(AC)$ si $(AC)$ est la médiatrice de $[HH']$. C'est cela qu'on doit montrer.



D'autre part, par définition, on a $(AC) \bot (HH')$, donc $(AC)$ est la hauteur de $AHH'$ issue de $A$. 

Donc si $AHH'$ est isocèle en $A$, alors cette hauteur de $AHH'$  est aussi la médiane issue de $A$ et c'est encore la médiatrice du côté opposé à $A$ c'est-à-dire $[HH']$.

Il suffit donc de montrer que $AHH'$ est isocèle en $A$. Pour cela, il suffit de montrer que les angles adjacents à la base sont égaux, autrement dit $\widehat{AHH'} = \widehat{AH'H}$ avec des angles géométriques non orientés, ou plus précisément avec des angles orientés $(H'A,H'H)=(HH',AH)$.

Suivant la méthodologie habituelle, on marque de façon systématique les angles égaux (ou complémentaires, supplémentaires etc) sur la figure. Ceci indique la marche à suivre pour la preuve.


\begin{center}
\includegraphics{../images/img007139-2}
\end{center}

Montrons que $(H'A,H'H)=(HH',AH)$. On a :
\begin{align*}
(H'A,H'H) &= (H'A,H'B) \text{ (car $(H'H)=(HB)$}\\
&= (CA,CB) \text{ (car $ABCH'$ est inscriptible)}\\
&= (CA,AH)+(AH,CB) \text{ (par Chasles)}\\
&= (CA,AH)+\pi/2 \text{ (car $(AH)$ est une hauteur de $ABC$)}\\
&= (CA,AH)+(HH',CA)\\
&= (HH',AH) \text{ (par Chasles)}
\end{align*}


\fincorrection
\correction{007140}

Rappelons la figure :

\begin{center}
\includegraphics{../images/img007140-1}
\end{center}

\begin{enumerate}
\item On rappelle que le milieu de l'hypoténuse d'un triangle rectangle est le centre de son cercle circonscrit (une autre façon de le dire est que l'hypoténuse est un diamètre du cercle circonscrit).

Si $IO=IA$, cela signifie que $I$ est sur la médiatrice de $[OA]$. D'autre part, $AOB$ est rectangle en $O$ et $I$ est par définition sur l'hypoténuse $[AB]$. Donc $I$ est  l'intersection de l'hypoténuse et d'une médiatrice d'un autre côté, c'est donc le milieu de l'hypoténuse par la propriété rappelée plus haut. Il est donc suffisant de montrer que $IO=IA$.


\item Pour montrer que $IO=IA$, il suffit de montrer que $IOA$ est isocèle en $I$, c'est-à-dire que $(AI,AO)=(OA,OI)$. Or, on a :
\begin{align*}
(AI,AO)& = (AB,AC) \text{ (mêmes droites)}\\
&= (DB,DC) \text{ (car $ABCD$ est inscriptible}\\
&= (DO,DH) \text{ (mêmes droites)}\\
&= (DO,OH)+(OH,DH) \text{ (par Chasles)}\\
&= (DO,OH)+\pi/2 \text{ (par définition de $H$)}\\
&= (OB,OI)+\pi/2 \text{ (mêmes droites)}\\
&= (OB,OI) + (OA,OB) \text{ (car $(OA)\bot (OB)$ d'après l'énoncé)}\\
&= (OA,OI) \text{ (par Chasles)}
\end{align*}
\end{enumerate}

\fincorrection
\correction{007141}
\begin{enumerate}
\item 
Soit $\mathcal D'$  une autre droite passant par $P$ et intersectant le cercle en deux points $A'$ et $B'$. 

Les quantités $\overline{PA}\cdot\overline{PB}$ et $\overline{PA'}\cdot\overline{PB'}$ ont forcément même signe, donc il suffit de montrer que $PA\cdot PB = PA'\cdot PB'$.

Comme $ABA'B'$ est inscriptible, le théorème de l'angle inscrit donne
\[ (BA,BA')=(B'A,B'A').\]
On en déduit que les triangles  $PAB'$ et $PA'B$ ont donc deux de leurs angles égaux, donc tous leurs angles égaux, donc sont semblables. On en déduit que:
\[\frac{PA}{PB'} = \frac{PA'}{PB},
\]
donc
\[
PA\cdot PB = PA'\cdot PB'.
\]
\item On procède de la même manière, en utilisant la version tangentielle du théorème de l'angle inscrit.
\item Avec les notations de la question précédente, on a par le théorème de Pythagore $PT^2+r^2=PO^2$, d'où le résultat.
\item L'ensemble est donc $\{P\in \mathcal P,\: ||\overrightarrow{OP}||=\lambda + r^2\}$ Si $\lambda < -r^2$, l'ensemble est vide. Si $\lambda=-r^2$, c'est le centre du cercle, et si $\lambda > -r^2$, c'est un cercle de centre $O$ et de rayon $\sqrt{\lambda + r^2}$.
\item Si les deux produits sont nuls, alors $P$ coïncide avec l'un des deux points $A$ et $B$, ainsi qu'avec l'un des deux points $C$ et $D$. On en déduit que $ABCD$ est en fait un triangle et qu'il est donc inscriptible.

Si les deux produits ne sont pas nuls, on obtient 
\[\frac{PA}{PD} = \frac{PC}{PB},
\]
donc les triangles $PAD$ et $PBC$ sont semblables, donc ont mêmes angles. On conclut en utilisant la réciproque du théorème de l'angle inscrit.

\end{enumerate}
\fincorrection
\correction{007142} 
Rappelons la figure:

\begin{center}
\includegraphics{../images/img007142-1}
\end{center}

Soit $M$ un point de $(AB)$. Sa puissance par rapport aux deux cercles est la même et vaut $\overrightarrow{MA}\cdot \overrightarrow{MB}$.
 
Par ailleurs, si $M$ est de plus sur la tangente commune $(CD)$, alors sa puissance par rapport au premier cercle est $MC^2$, et celle par rapport au deuxième cercle est $MD^2$. On en déduit que $MC=MD$, donc $M$ est le milieu de $[CD]$.
\fincorrection  
\correction{007143} 
Soit $A'$ le symétrique de $A$ par rapport à $[BC]$. Alors, $BACA'$ est inscriptible, et la puissance de $H$ par rapport à son cercle circonscrit est 
\[ p_{\mathcal C}(H) = HB\cdot HC=  HA\cdot HA' = HA^2.\]

\textbf{Deuxième solution, n'utilisant pas la puissance d'un point par rapport à un cercle} : 

Les triangles $ABC$, $ABH$ et $ACH$ sont semblables car ils ont à chaque fois deux (donc trois) angles identiques. Les rapports de longueurs de côtés homologues sont donc égaux, ce qui donne
\[ \frac{AB}{AC} = \frac{HB}{HA} = \frac{HA}{HC}\]
d'où on tire $HA^2 = HB\cdot HC$, ce qu'il fallait démontrer.

\fincorrection
\correction{007144} 

Soit $G$ le projeté orthogonal de $F$ sur $\mathcal D$.

\begin{align*}
\overrightarrow{PE}\cdot \overrightarrow{PF}
&= \overrightarrow{PE}\cdot (\overrightarrow{PG}+\overrightarrow{GF})\\
&= \overrightarrow{PE}\cdot \overrightarrow{PG}
\end{align*}


Distinguons deux cas $ G=E$ et $G \neq E$. Dans le premier cas, cela signifie que le diamètre $(EF)$ est orthogonal à la droite $\mathcal D$, donc que celle-ci est tangente au cercle. Dans ce cas, $PE\cdot PG = PE^2$ est la puissance de $P$ par rapport au cercle.

Dans le deuxième cas, $EFG$ est rectangle en $G$, donc $G$ est sur le cercle de diamètre $[EF]$. On en déduit que $G$ est le second point d'intersection de $(PE)$ avec le cercle, et donc que $\overrightarrow{PE}\cdot \overrightarrow{PG}$ est là aussi la puissance de $P$ par rapport au cercle.


\begin{center}
\includegraphics{../images/img007144-2}
\end{center}

\fincorrection  
\correction{007152} 
Il existe une unique similitude directe $s$ qui envoie $[AE]$ sur $[CB]$. Son rapport est $1/l$ et son angle $-\pi/2$.

Soit $D' = s(D)$. Alors $(\overrightarrow{AD},\overrightarrow{CD'})=(\overrightarrow{AD},\overrightarrow{s(A)s(D)})=-\pi/2$ et $CD' = \frac1lAD = \frac1l$.

Cette similitude envoie donc $D$ et $F$ ssi $CF=CD'$ c'est-à-dire ssi $1/l = l-1$. Comme $l\neq 0$, cette équation est équivalente à $1=l^2-l$ . Elle admet une unique solution positive : $\frac{1+\sqrt 5}{2}$.
\fincorrection  
\correction{007153} 
Il existe une unique similitude directe qui envoie le couple $(B,D)$ sur le couple $(C,E)$. Comme $BD=CE$, c'est une isométrie donc une rotation, et comme $(\overrightarrow{BD},\overrightarrow{CE}) = \pi/2$, son angle est $\pi/2$.

Or, il n'y a qu'une rotation d'angle $\pi/2$ qui envoie $B$ sur $C$ : son centre est $I$. On en déduit que $IDE$ est rectangle isocèle en $I$.
\fincorrection  
\correction{007154} 
\begin{enumerate}
\item Pour faire la figure on peut
\begin{enumerate}
\item tracer la droite $(AB_0)$, sa perpendiculaire passant par $A$, puis les deux bissectrices de ces deux droites. 
\item Tracer les cercles de centre $A$ et de rayons $8$, $4$, $2$ et plus généralement $8\cdot 2^{-k}$. Les points de la suite $(B_n)$ sont sur les intersections de ces droites et cercles.
\end{enumerate}
\begin{center}
\includegraphics{../images/img007154-1}
\end{center}
\item  Soit $n\in\N$. On a $AB_{n+1}B_{n+2} = s(A)s(B_n)s(B_{n+1})$. 
\item Pour tout entier naturel $n$, on définit $l_n=B_nB_{n+1}$. La longueur (finie ou infinie) de la spirale est la somme de la série $\sum l_n$.

D'après la question précédente, $B_{n+1}B_{n+2} = \frac12 B_{n}B_{n+1}$. Autrement dit, la suite $(l_n)_{n\in\N}$ est géométrique de raison $\frac12$, et donc pour tout $n\in \N$ on a $l_n = \left(\frac12\right)^n\cdot l_0$. On en déduit que la série $\sum l_n$ converge absolument, et que sa somme est $l_0\cdot \frac{1}{1-\frac12} = 2l_0=2B_0B_1$.
\end{enumerate}
\fincorrection  
\correction{007155} 
Complétons la figure :

\begin{center}
\includegraphics{../images/img007155-2}
\end{center}



\begin{enumerate}
\item Suivons l'indication. Soit $s = s_2\circ s_1$. C'est une similitude directe de rapport $1$ et d'angle $\pi/2$. On vérifie également que $s(C')=C'$, donc c'est la rotation de centre $C'$ et d'angle $\pi/2$. 

Or, d'une part on a  $s(R)=A$, et d'autre part $s_1(Q) = C$ et $s_2(C) = P$, donc $s(Q) = P$. On en déduit que $\overrightarrow{AP}$ est l'image de $\overrightarrow{RQ}$ par une rotation de $\pi/2$. Les vecteurs ont donc même norme et sont orthogonaux.
\item  Les droites $(AP)$, $(BQ)$ et $(CE)$ sont les hauteurs de $PQR$ donc sont concourantes en l'orthocentre de $PQR$. Ce point est le \emph{point intérieur de Vecten} de $ABC$.
\end{enumerate}
\fincorrection  
\correction{007156} 
Soient $H_{12}$, $H_{34}$, $H_{13}$ et $H_{24}$ les projetés orthogonaux de $M$ sur les droites $(M_1M_2)$, $(M_3M_4)$, $(M_1M_3)$ et $(M_2M_4)$. Soit $s$ la similitude directe de centre $M$ qui envoie $M_2$ sur $M_3$. L'image par $s$ de la droite $(M_1M_2)$ est une droite $\mathcal D$ passant par $M_3$ et telle que 
\[ (M_1M_2,\mathcal D)=(MM_2,MM_3).\]
Or, on a $(MM_2,MM_3)=(M_1M_2,M_1M_3)$ car les quatre points sont cocycliques. On en déduit que $\mathcal D = (M_1M_3)$, et donc que $s$ envoie la droite $(M_1M_2)$ sur la droite $(M_1M_3)$.

On prouve de la même manière que $s$ envoie la droite $(M_2M_4)$ sur la droite $(M_3M_4)$.

Comme $s(M)=M$ et que les similitudes conservent les angles orientés et en particulier l'orthogonalité, on en déduit que $s(H_{12})=H_{13}$ et que $s(H_{24})=H_{34}$.

Comme une similitude conserve les rapports de longueurs, on a finalement
$\frac{MH_{13}}{MH_{34}} = \frac{MH_{12}}{MH_{24}}$ c'est-à-dire  :
\[
MH_{12} \cdot MH_{34} = MH_{13}\cdot MH_{24}.
\]

\fincorrection  
\correction{007157}
\begin{enumerate}
\item
Soit $\phi$ un déplacement préservant $\mathcal S$. C'est une translation ou une rotation.
% pas le groupe car il est un peu compliqué : c'est Z semi-direct Z/2Z.
\begin{enumerate}
\item Si $\phi$ est une translation de vecteur $\vec v$, alors $\vec v$ est horizontal (s'il avait une composante verticale, alors pour $n$ grand et $s\in \mathcal S$, le point $\phi^n(s)$ aurait une ordonnée qui n'appartiendrait pas à $[-1;1]$ ce qui est impossible. Si $\vec v $ est horizontal, on en déduit que $||\vec v||$ soit être une période de sinus, donc un multiple de $2\pi$. Réciproquement, toute translation horizontale d'un multiple de $2\pi$ est une isométries de $\mathcal S$.
\item Si $\phi$ est une rotation, alors son centre est un point de la forme $(n\pi,0)$ pour un certain $n\in \Z$, et l'angle est $\pi$, autrement dit $\phi$ est une symétrie centrale.

Réciproquement, les symétries centrales de centres les points $(n\pi,0)$ pour $n\in \Z$ sont dans $G$. La composée de deux telles symétries est une des translations trouvées plus haut.
\end{enumerate}
Déterminons maintenant les antidéplacements préservant $\mathcal A$. Un antidéplacement est soit une réflexion soit une réflexion glissée.
\begin{enumerate}
\item Si $\phi$ est une réflexion, son axe est forcément vertical par un argument similaire à celui plus haut. On en déduit ensuite que l'axe doit être une droite d'équation $x=\pi/2+n\pi$ avec $n$ un certain entier relatif.
\item Si $\phi$ est une réflexion glissée, la direction de glissement doit être horizontale (idem), et on en déduit ensuite que la longueur de translation doit être un multiple de $\pi$.
\end{enumerate}
 \item On trouve les translations horizontales de vecteur de norme multiple de $\pi$ et les symétries centrales dont le centre est de la forme $(k\pi,0)$ pour $k\in\Z$. 
 \end{enumerate}
\fincorrection  
\correction{007158} 
Comme une isométrie est affine, elle conserve les barycentres. Soit $P$ un sommet du triangle. Comme ce n'est pas un barycentre d'autres points du triangle, son image par une isométrie fixant le triangle non plus, c'est-à-dire que son image est un sommet. On en déduit que les sommets sont envoyés sur les sommets.

Par préservation du barycentre une isométrie de $\mathcal T$ fixe son isobarycentre.

Comme elle a au moins un point fixe, ça ne peut être une translation ou une symétrie glissée.
\fincorrection  
\correction{007159} 
\begin{enumerate}
\item Les trois réflexions suivant les médiatrices du triangle conviennent, de même que l'identité, et les rotations d'angles $\pm 2\pi/3$ de centre $O$. Ceci donne six isométries, trois directes et trois indirectes. On remarque en les composant entre elles que l'on n'obtient pas d'autres isométries. Cet ensemble de six isométries est donc stable par composition, et il est également stable par inverse, c'est donc un sous-groupe de $G$. Dans la suite, on va prouver que c'est $G$ tout entier.
\item Comme une isométrie est affine, elle conserve les barycentres. Soit $P$ un sommet du triangle. Comme ce n'est pas un barycentre d'autres points du triangle, son image par une isométrie fixant le triangle non plus, c'est-à-dire que son image est un sommet. On en déduit que les sommets sont envoyés sur les sommets.

\item D'après ce qui précède, une isométrie de $\mathcal T$ permute les sommets. Donc à toute isométrie $f\in G$ on peut associer la bijection dans ${\rm Bij}(\{A,B,C\})$ qui lui correspond. Ce groupe de bijections est isomorphe à $\mathfrak S_3$, en numérotant les sommets de $1$ à $3$ ($A$ est le premier sommet, $B$ le second  etc). On obtient donc une application
\[ G \to \mathfrak S_3.\]
Elle est injective car une application affine est complètement déterminée par l'image de trois points non alignés : donc préciser la permutation sur les sommets du triangle détermine complètement l'isométrie du plan. C'est un morphisme de groupes par construction : composer les isométries va composer les permutations des sommets.
\item Pour montrer que le morphisme de la question précédente est surjectif, on va utiliser la première question. À part l'identité qui est envoyée sur l'identité, les deux rotations sont envoyées sur les $3$-cycles $(123)$ et $(132)$. Les trois réflexions sont envoyées sur les trois transpositions, par exemple la réflexion suivant la médiatrice de $[BC]$ est envoyée sur la transposition $(23)$. On remarque qu'une isométrie est directe ssi la permutation associée est paire.
\item D'après la question précédente, $H$ est composé de l'identité et des deux rotations décrites plus haut. Ce sous-groupe est envoyé sur le groupe des permutations paires $\mathfrak A_3$. Il est isomorphe à $\Z/3\Z$ par l'application qui envoie $[0]$ sur l'identité, $[1]$ sur $(123)$ et $[2]$ sur $(132)$, et qui est un morphisme de groupe bijectif.
\end{enumerate}
 
\fincorrection  
\correction{007160} 

Comme une isométrie est affine, elle conserve les barycentres. Soit $P$ un sommet du triangle. Comme ce n'est pas un barycentre d'autres points du triangle, son image par une isométrie fixant le triangle non plus, c'est-à-dire que son image est un sommet. On en déduit que les sommets sont envoyés sur les sommets.

Dans la suite on suppose que $ABC$ est isocèle en $A$. Soit $f$ une isométrie de $ABC$. Comme une isométrie conserve les angles non orientés et que le triangle est isocèle en $A$ et non équilatéral, on en déduit que $B$ et $C$ sont envoyés soit sur $B$ soit sur $C$. Le point $A$ est donc fixe.

Comme une application affine est déterminée par l'image de trois points non alignés, on conclut que le triangle $ABC$ n'admet que deux isométries : l'identité, et la réflexion $\sigma$ suivante la médiatrice de $[BC]$.

Le groupe ${\rm Isom}(\mathcal T)$ est donc de cardinal deux, donc isomorphe à $\Z/2\Z$.

(Un isomorphisme est donné par $\phi : \Z/2\Z \to {\rm Isom}(\mathcal T)$ l'application qui envoie $[0]=2\Z$ sur l'identité et $[1]=1+2\Z$ sur $\sigma$. C'est une bijection, et on vérifie que c'est un isomorphisme de groupes.)
\fincorrection  
\correction{007161} 
\begin{enumerate}
\item Montrons d'abord qu'une isométrie du carré permute les sommets.
\begin{enumerate}
\item Première preuve : 
considérons une diagonale, par exemple $[AC]$. Comme $g$ est une isométrie, la distance entre $A$ et $C$ est la même qu'entre leurs images. Or, deux points du carré à distance $AC$ sont forcément deux sommets d'une diagonale du carré. L'isométrie $g$ permute donc les sommets.
\item Deuxième preuve :  
une isométrie est affine, donc préserve les barycentres, donc envoie les points extrémaux (ceux qui ne sont pas des barycentres d'autres points du carré) sur d'autres points extrémaux, et donc permute les sommets.
\end{enumerate}

Le centre $O$ étant l'isobarycentre des sommets, et une isométrie conservant les barycentres, on en déduit que $g(O)=O$. Ce n'est pas une translation ou une réflexion glissée, puisque ces dernières n'ont pas de points fixes. C'est donc une rotation (dont le centre est forcément $O$ car $O$ est fixe), ou bien une réflexion (dont l'axe contient $O$).

\item Soit $f$ une isométrie du carré. Elle envoie un sommet $P$ sur un des autres sommets.

\begin{enumerate}
\item Si c'est une réflexion, son axe est la médiatrice du segment $[Pf(P)]$, donc l'axe est forcément une diagonale, ou bien une médiatrice d'un côté du carré.

Réciproquement, on vérifie que les réflexions dont les axes sont les deux diagonales ou bien les deux médiatrices des côtés sont bien des isométries du carré.

\item Si c'est une rotation, son centre est $O$, et comme $P$ est envoyé sur un autre sommet, les angles possibles sont les multiples de $\pi/2$. Réciproquement, on vérifie que les rotations de centre $O$ et d'angles $0$, $\pi/2$, $\pi$ et $3\pi/2$ sont des isométries du carré.
\end{enumerate}
Le groupe des isométries du carré a donc huit éléments : quatre réflexions et quatre rotations (dont l'identité).

\item Le groupe $H$ est composé des quatre rotations décrites plus haut. L'application  qui à un entier $n$ associe la rotation de centre $O$ et d'angle $n\pi/2$ est un morphisme de groupes, noté $\phi : (\Z,+) \to (H,\circ)$. Il est surjectif.  Déterminons son noyau. Soit $n\in \Z$. On a :
\begin{align*}
n\in \ker(\phi)
&\Leftrightarrow \phi(n) = \operatorname{Id}\\
&\Leftrightarrow n\pi/2 \equiv 0[2\pi]\\
&\Leftrightarrow n \equiv 0[4]\\
&\Leftrightarrow n \in 4\Z.\\
\end{align*}
On en déduit $\ker(\phi) = 4\Z$. Ce morphisme de groupes passe donc au quotient par $4\Z$ et induit un morphisme injectif $\overline{\phi} = \Z/4\Z \to H$ qui est donc un isomorphisme de groupes.
\end{enumerate}

\fincorrection  
\correction{007162} 
Dans le rectangle, les points maximalement éloignés sont les sommets des diagonales. Comme une isométrie conserve les distances, on en déduit qu'une isométrie du rectangle doit envoyer une diagonale sur une autre, et donc soit permuter les sommets. Elle fixe donc le centre, qui est l'isobarycentre des sommets, et donc est une rotation ou une réflexion, car les translations et les réflexions glissées n'ont pas de points fixes.

Soit $P$ un sommet et $Q$ son image par une isométrie du rectangle. Si c'est une réflexion, son axe est donc la médiatrice de $[PQ]$, donc l'axe peut être une des deux médiatrices des côtés du rectangle (et pas une médiatrice d'une diagonale, car le rectangle n'est pas carré). Réciproquement, ces deux réflexions, notons-les $\sigma$ et $\sigma'$, sont des isométries du rectangle.

Si c'est une rotation, son angle est $0$ ou bien $\pi$, car le rectangle n'est pas carré. Réciproquement, ces deux rotations ($\operatorname{Id}$ et $-\operatorname{Id}$) conviennent.

Le groupe d'isométries du rectangle est isomorphe à $(\Z/2\Z)^2$, un isomorphisme possible étant celui qui envoie $(1,0)$ sur $\sigma$ et $(0,1)$ sur $\sigma'$ (et donc $(1,1) = (0,1)+(1,0)$ sur $\sigma\circ \sigma' = -\operatorname{Id}$.
\fincorrection  
\correction{007165} 
\begin{enumerate}
\item Il suffit de montrer que $(AU,AW)=(PU,PW)$. 
\begin{center}
\includegraphics{../images/img007165-1}
\end{center}
Or, on a :
\begin{align*}
(AU,AW)
&= (AB,AC) \text{ (mêmes droites)}\\
&= \pi/3 \text{ (car équilatéral direct)}
\end{align*}
et
\begin{align*}
(PU,PW)
&= (BC,BA) \text{ (droites parallèles donc mêmes angles)}\\
&= \pi/3\text{ (car équilatéral direct)}.
\end{align*}
Donc $(AU,AW)=(PU,PW)$ et donc $APUW$ est inscriptible.

On procède de même pour l'autre quadrilatère.
\item Pour chaque quadrilatère inscriptible, on peut en déduire  six égalités d'angles inscrits, car il y a $6=\binom{2}{4}$ façons de choisir une corde. Si $APUW$ est inscriptible, on a donc les égalités d'angles de droites:
\[(UA,UP)=(WA,WP),\quad (AU,AW)=(PU,PW),\]
et surtout:
\[ (WA,WU)=(PA,PU)\quad (AU,AP)=(WU,WP),\]
\[ (UP,UW)=(AP,AW),\quad (PW,PA)=(UW,UA).\]
Les plus intéressantes (et utiles pour la suite) sont les quatre dernières.

On procède de même pour l'autre quadrilatère.
\item Il suffit de montrer que $ABCP$ est inscriptible si et seulement si $(WU,WP)=(WV,WP)$. 
Cette dernière assertion dit en effet que les droites $(WU)$ et $(WV)$ forment le même angle avec $WP$ donc sont parallèles, et donc égales car elles ont le point $W$ en commun.

\begin{center}
\includegraphics{../images/img007165-2}
\end{center}
On a :
\begin{align*}
(WU,WP)
&= (AU,AP) \text{ (car $AUWP$ est inscriptible)}\\
&= (AB,AP) \text{ (mêmes droites car $U\in (AB)$ )}\\
&= (CB,CP) \text{ \underline{ssi $ABPC$ est inscriptible}}
\end{align*}
et d'autre part, on a 
\begin{align*}
(CB,CP)
&= (CV,CP) \text{ (mêmes droites car $V\in (BC)$ )}\\
&= (WV,WP) \text{ (car $PWCV$ est inscriptible)}
\end{align*}

\end{enumerate}

\fincorrection  
\correction{007167}
\begin{enumerate}
\item Soit $\rho$ la rotation de centre $O$ et d'angle $\pi/2$. D'après l'énoncé, on a $\rho(B)=A$ et $\rho(D)=C$. Donc $[AC]$ est l'image de $[AD]$ par $\rho$, d'où on déduit que $(AC)\bot (BD)$ et que $AC=BD$.

\item Le quadrilatère $IJKL$ est toujours un parallélogramme, même sans hypothèses sur $ABCD$ (c'est le théorème de Varignon). En effet, par le théorème de Thalès (ou simplement le théorème des milieux) : 
\[ \overrightarrow{IL} = \frac12 \overrightarrow{BD} = \overrightarrow{JK}.\]
Ceci signifie que les côtés $[IL]$ et $[JK]$ ont même longueur et sont parallèles, donc $IJKL$ est un parallélogramme.

Pour voir que c'est un carré, remarquons qu'on montre de même que  
\[ \overrightarrow{IJ} = \frac12 \overrightarrow{AC} = \overrightarrow{LK},\]
et d'après la première question, $(AC)\bot (BD)$ et que $AC=BD$, donc $IJKL$ est un parallélogramme ayant un angle droit et des cotés consécutifs de même longueur. C'est donc un carré.
\end{enumerate}
\fincorrection
\correction{007168}
\begin{enumerate}
\item Par construction, la droite $\Delta$ est la hauteur et la médiane de $ABD$. Le triangle est donc isocèle en $A$ et $\Delta$ est également sa bissectrice, ce qui montre que $\widehat{BAA_1} = \widehat{A_1AD}$ et donc que $D\in (AC)$.
\item Il suffit de prouver que $(A_2A_1,A_2C)=(DA_1,DC)$.
\begin{center}
\includegraphics{../images/img007168-2}
\end{center}
Or on a :
\begin{align*}
(A_2A_1,A_2C)
&= (A_2A,A_2C)\\
&= (BA,BC) \text{ (car $ABCA_2$ est inscriptible)}\\
&= (BA,BA_1) \text{ (mêmes droites)}\\
&= (DA,DA_1) \text{ (par réflexion suivant $\Delta$)}\\
&= (DC,DA_1) \text{ (mêmes droites)},
\end{align*}
ce qu'il fallait démontrer.
\item Remarquons déjà que $AB=AD$.  D'autre part, il suffit de montrer que
\[ \frac{AA_1}{AD} = \frac{AC}{AA_2}.\]

La question précédente montre que les triangles $ADA_1$ et $AA_2C$ sont (inversement) semblables, puisqu'ils ont deux (et donc trois) angles en commun. Les rapports entre les côtés sont donc égaux, c'est-à-dire précisément
\[ \frac{AA_1}{AD} = \frac{AC}{AA_2}.\]

Remarque : si on connaît la notion de puissance d'un point par rapport à un cercle, on peut conclure plus vite : les deux produits égaux sont la puissance de $A$ par rapport au cercle $A_1A_2CD$.
\end{enumerate}
\fincorrection
\nocorrection
\nocorrection
\nocorrection
\nocorrection
\nocorrection
\nocorrection
\nocorrection
\nocorrection
\nocorrection
\nocorrection
\nocorrection
\nocorrection
\nocorrection
\nocorrection
\nocorrection
\nocorrection
\nocorrection
\nocorrection
\nocorrection
\nocorrection
\nocorrection
\nocorrection
\nocorrection
\nocorrection
\nocorrection
\nocorrection
\nocorrection
\nocorrection
\nocorrection
\nocorrection
\nocorrection
\nocorrection
\nocorrection
\nocorrection
\nocorrection
\nocorrection
\nocorrection
\nocorrection
\nocorrection
\nocorrection
\nocorrection
\nocorrection
\nocorrection
\nocorrection
\nocorrection
\nocorrection
\nocorrection
\nocorrection
\nocorrection
\nocorrection
\nocorrection
\correction{007505}
\begin{enumerate}
    \item Comme la somme des masses $2+1$ est non nulle, le barycentre est bien défini.
    L'abscisse de $G$ est $\frac{2\times (-2)+1\times 4}{2+1}=0$
    et l'ordonnée $\frac{2\times 1+1\times 4}{2+1}=2$.
    \item $2GA^ 2+GB^2=2((-2)^2+(1-2)^2)+(4^2+(4-2)^2)=30$.
    \item Soit $M$ un point du plan $P$.
    \begin{eqnarray*}
        2MA^2+MB^2&=&2\vec{MA}^2+\vec{MB}^2=2(\vec{MG}+\vec{GA})^2+(\vec{MG}+\vec{GB})^2\\
        &=&3\vec{MG}^2+<\vec{MG},2\vec{GA}+\vec{GB}>+2\vec{GA}^ 2+\vec{GB}^2\\
        &=&3MG^2+2GA^ 2+GB^2.
    \end{eqnarray*}
    \item Soit $M$ un point du plan $P$.
    \begin{eqnarray*}
        M\in\mathcal{L}&\iff & 2MA^2+MB^2=42\iff 3MG^2+2GA^ 2+GB^2=42\\
        &\iff& 3MG^2+30=42\iff MG^2=4.
    \end{eqnarray*}
    L'ensemble $\mathcal L$ est donc le cercle de centre $G$ et de rayon $2$.
\end{enumerate}
\fincorrection
\correction{007506}
\begin{enumerate}
    \item[a.] $\ z\mapsto 3z+4$ ne représente pas une isométrie à cause du facteur $3$.
    \item[b.] $\ z\mapsto 3\overline{z}+4$ ne représente pas une isométrie à cause du facteur $3$.
    \item[c.] $\ z\mapsto \overline{z}+4$ représente une isométrie, indirecte (présence de la conjugaison). Il s'agit de la réflexion par rapport à l'axe des abscisses $\ z\mapsto \overline{z}$ composée par la translation de vecteur $4\vec{\imath}$ dans la direction de l'axe des abscisses. C'est donc une symétrie glissée. 
    \item[d.] $\ z\mapsto i\overline{z}+4$ représente une isométrie, indirecte (présence de la conjugaison). Il s'agit de la réflexion $s$ par rapport à l'axe des abscisses $\ z\mapsto \overline{z}$ composée par $z\mapsto iz$, la rotation $r$ de centre $O$ et d'angle $+\frac{\pi}{2}$
    puis par la translation $t$ de vecteur $4\vec{\imath}$.
    On peut écrire $r\circ s=(s'\circ s)\circ s=s'$ où $s'$ est la réflexion par rapport à la droite $d'$ d'équation $y=x$. On décompose ensuite $4\vec{\imath}=2(\vec{\imath}+\vec{\jmath})+2(\vec{\imath}-\vec{\jmath})$
    comme somme d'un vecteur $\vec{v}_2:=2(\vec{\imath}+\vec{\jmath})$ directeur de $d'$ et d'un vecteur $\vec{v}_1:=2(\vec{\imath}-\vec{\jmath})$ normal à $d'$.
    On en déduit
    $t\circ s'=t_{\vec{v}_2}\circ t_{\vec{v}_1}\circ s'= t_{\vec{v}_2}\circ s''$
    où $s''$ est la réflexion par rapport à la droite $d''=d'+\frac{\vec{v}_1}{2}$ parallèle à $d'$. Maintenant, $t\circ r\circ s=t\circ s'=t_{\vec{v}_2}\circ s''$. Comme $\vec{v}_2$ est parallèle à $d''$, il s'agit d'une symétrie glissée.
\end{enumerate}
\fincorrection
\correction{007507}
\begin{enumerate}
    \item Soit $A, B, C$ sont trois points distincts d'un plan affine euclidien $\mathcal{P}$.
    La somme des angles de vecteurs se réécrit par symétrie centrale par rapport à $C$ comme
    \begin{eqnarray*}
        \lefteqn{((\vec{AB}, \vec{AC}))+((\vec{BC},\vec{BA}))+((\vec{CA},\vec{CB}))}
        &&\\
        &=&((\vec{AB}, \vec{AC}))+((\vec{BC},\vec{BA}))+((-\vec{CA},-\vec{CB}))\\
        &=&((\vec{AB}, \vec{AC}))+((\vec{BC},\vec{BA}))+((\vec{AC},\vec{BC}))\\
        &=&((\vec{AB}, \vec{AC}))+((\vec{AC},\vec{BC}))+((\vec{BC},\vec{BA}))\\
        &=&((\vec{AB},\vec{BA}))\\
    \end{eqnarray*}
    c'est à dire l'angle plat.
    
    \item Soit $\mathcal{C}$ un cercle de $\mathcal{P}$ et $A$ un point de $\mathcal{C}$.
    Soit $\mathcal{C}'$ l'image par une rotation $r$ de centre $A$ du cercle $\mathcal{C}$.
    Soit $B$ l'autre point d'intersection de $\mathcal{C}$ et $\mathcal{C}'$.
    Soit $D$ le point de $\mathcal{C}$ diamétralement opposé à $A$ sur $\mathcal{C}$.
    Soit $D'=r(D)$ son image par $r$. 
    Soit $O$ le centre de $\mathcal{C}$ et $O'$ le centre de $\mathcal{C}'$.
    
    Comme $D$ appartient à $(AO)$, le point $D'=r(D)$ appartient à $r((AO))=(A'O')$.
    Donc $[A'D']$ est un diamètre de $\mathcal{C}'$.
    Comme $B$ appartient à $\mathcal{C}$ et comme $[AD]$ est un diamètre de $\mathcal{C}$, la droite $(BD)$ est orthogonale à $(AB)$.
    De même, comme $B$ appartient à $\mathcal{C}'$ et comme $[AD']$ est un diamètre de $\mathcal{C}$, la droite $(BD')$ est orthogonale à $(AB)$.
    Par conséquent, les droites $(BD)$ et $(BD')$ sont confondues et $D$, $D'$ et $B$ sont alignés.
    \item Soit $M$ un point quelconque de $\mathcal{C}$.
    Dans les triangles $BAM'$ et $BAM$
    $$((\vec{BM'}, \vec{BA}))+((\vec{AB},\vec{AM'}))+((\vec{M'A},\vec{M'B}))$$
    et $$(( \vec{BA},\vec{BM}))+((\vec{AM},\vec{AB}))+((\vec{MB},\vec{MA}))$$
    sont deux angles plats.
    Par somme,$$((\vec{BM'},\vec{BM}))+((\vec{AM},\vec{AM'}))
    +((\vec{M'A},\vec{M'B}))+((\vec{MB},\vec{MA}))$$ est un angle nul.
    
    Par propriété des angles inscrits, $((\vec{M'A},\vec{M'B}))=((\vec{D'A},\vec{D'B}))$ et $((\vec{MB},\vec{MA}))=((\vec{DB},\vec{DA}))$.
    Dans le triangle $ADD'$, on trouve que
    $$((\vec{D'A},\vec{D'B}))+((\vec{DB},\vec{DA}))$$
    et la différence d'un angle plat et de $((\vec{AD},\vec{AD'}))=((\vec{AM},\vec{AM'}))$ l'angle de la rotation.
    En conséquence,
    $((\vec{BM'},\vec{BM}))$ est un angle plat et $M$, $M'=r(M)$ et $B$ sont alignés.
\end{enumerate}
\fincorrection
\correction{007508}
\begin{enumerate}
    \item La matrice $A$ de la forme bilinéaire symétrique donnée en coordonnées par
    $$f(\begin{pmatrix}x\\y\end{pmatrix},\begin{pmatrix}x'\\y'\end{pmatrix})
    =xx'+19yy'+12xy'+12x'y$$
    est $$\begin{pmatrix}
    1&12\\12&19
    \end{pmatrix}$$
    
    \item On trouve que $e_1+2e_2$ est vecteur propre de $A$ de valeur propre $25$ et $2e_1-e_2$ est vecteur propre de $A$ de valeur propre $-5$.
    \item Comme la matrice $A$ est inversible, la conique d'équation
    $$x^2+19y^2+24xy+5y=0$$
    est une conique à centre. Comme les valeurs propres de $A$ sont de signe opposé, il s'agit d'une hyperbole.
\end{enumerate}
\fincorrection
\correction{007512}
\begin{enumerate}
    \item Comme $A$ a pour coordonnées barycentriques $(1,0,0)$ et le milieu $A'$ de $[BC]$ $(0,1,1)$, la médiane $(AA')$ de $ABC$ a pour équation barycentrique 
    $$\det\left(\begin{array}{ccc}1&0&a\\ 0&1&b\\0&1&c
    \end{array}\right)=c-b=0.$$
    De même la médiane $(BB')$ a pour équation barycentrique $a=c$ et $(CC')$
    $a=b$.
    \item Les médianes du triangle $ABC$ sont concourantes au point de coordonnées barycentriques $(1,1,1)$..
\end{enumerate}
\fincorrection
\correction{007513}
\begin{enumerate}
    \item Les éléments du groupe des isométries du plan qui conservent un carré ont tous le centre $O$ du carré comme points fixes. Ce sont donc outre l'identité, des réflexions d'axe passant par $O$ et des rotations de centre $O$.
    Il y a quatre rotations $Id,r,r^2, r^3$ où $r$ est la rotation de centre $O$ et d'angle $+\pi/2$. Il y a aussi quatre réflexions $s,rs,r^2s,r^3s$.
    \item Comme la composée des réflexions dépend de l'ordre, $(rs)s=r$ et $s(rs)=r^{-1}$, ce groupe n'est pas commutatif.
    \item Le groupe des déplacements du plan qui conservent un carré est un sous-groupe du groupe commutatif des rotation de centre $O$. Il est donc commutatif.
\end{enumerate}
\fincorrection
\correction{007514}
On construit le symétrique $s(d)$ par rapport à $O$ de la droite $d$.
Comme le symétrique de $A$ doit être sur $d$, $A$ doit être sur $s(d)$ et sur le cercle $\mathcal{C}$.

\includegraphics[scale=0.5]{images/img-mour-269-2}
\fincorrection
\nocorrection
\nocorrection
\nocorrection
\correction{002034}
\begin{enumerate}
  \item Notons $\mathcal{R}$ le repère initial $(0,\vec i,\vec j,\vec k)$.
Dire qu'un point $M$ du plan a pour coordonnées $(x,y,z)$ dans $\mathcal{R}$ 
signifie $\overrightarrow{OM} = x \vec i + y \vec j+ z\vec k$.

Si $\mathcal{R}'$ désigne un autre repère $(A, \vec u, \vec v,\vec w)$
alors le même point $M$ a pour coordonnées $(x',y',z')$ dans $\mathcal{R}'$ signifie
$\overrightarrow{AM} = x' \vec u + y'\vec v+z'\vec w$.

La formule de changement c'est simplement écrire les coordonnées de l'égalité
$\overrightarrow{OM} = \overrightarrow{OA}+\overrightarrow{AM}$.


$$\begin{pmatrix} x \\ y \\ z \end{pmatrix}
= \begin{pmatrix} -2 \\ 4 \\ 1 \end{pmatrix} + x'\vec u + y'\vec v+z'\vec w 
$$
Mais on connaît les coordonnées de 
$\vec u, \vec v,\vec w$ dans $\mathcal{R}$ :
$$\begin{pmatrix} x \\ y \\ z \end{pmatrix}
= \begin{pmatrix} -2 \\ 4 \\ 1 \end{pmatrix} +
 x' \begin{pmatrix} 1 \\ 1 \\ 1 \end{pmatrix}
+ y' \begin{pmatrix} 2 \\ 2 \\ -4 \end{pmatrix}
+ z' \begin{pmatrix} 3 \\ -1 \\ 1 \end{pmatrix}$$

D'où l'égalité de changement de repère :
$$(\mathcal{S}) \qquad \left\lbrace
\begin{array}{l}
 x = -2 + x' + 2y' + 3z' \\  
 y = 4 + x' + 2y' - z' \\ 
 z = 1 + x'- 4y' + z' \\
\end{array}
\right.$$

  \item  Dans l'équation de la droite $(D)$ 
$\left\{\begin{array}{l} y-z=3 \\ x+y=2 \end{array}\right.$ dans le repère $\mathcal{R}$
on remplace $x,y,z$ par la formule $(\mathcal{S})$  obtenue à la question précédente.

On obtient :
$$\left\{\begin{array}{l} \big(4 + x' + 2y' - z'\big)-
\big(1 + x'- 4y' + z'\big)=3 \\ 
\big(-2 + x' + 2y' + 3z'\big)+
\big(4 + x' + 2y' - z'\big) = 2 \end{array}\right..$$
Ce qui donne une équation de $(D)$ dans le repère $\mathcal{R}'$ :
$$\left\{\begin{array}{l} 
6y'-2z'=0   \\ 
2x'+4y'+2z'=0 \end{array}\right.
\quad \text{ou encore} \quad 
\left\{\begin{array}{l} 
3y'-z'=0   \\ 
x'+2y'+z'=0 \end{array}\right.$$
En particulier en faisant $(x',y',z')=(0,0,0)$ on remarque que cette droite passe par $A$.

  \item Nous avions obtenu l'égalité $(\mathcal{S})$ de changement de repère de $\mathcal{R}'$ 
vers $\mathcal{R}$ qui s'écrit :
$$\left\lbrace
\begin{array}{l}
 x + 2 =  x' + 2y' + 3z' \\  
 y - 4 =  x' + 2y' - z' \\ 
 z - 1 =  x'- 4y' + z' \\
\end{array}
\right. 
\quad \implies \quad 
\left\lbrace
\begin{array}{l}
 X =  x' + 2y' + 3z' \\  
 Y =  x' + 2y' - z' \\ 
 Z =  x'- 4y' + z' \\
\end{array}
\right. $$
Où l'on a noté $X=x+2$, $Y=y-4$, $Z=z-1$. On inverse le système 
par la méthode de Gauss pour obtenir après calculs :
$$\left\lbrace
\begin{array}{l}
 x' =  \frac{1}{12}\big(X+7Y+4Z\big) \\  
 y' =  \frac{1}{12}\big(X+Y-2Z\big) \\ 
 z' =  \frac{1}{12}\big(3X-3Y\big) \\
\end{array}
\right. $$

Donc 
$$\left\lbrace
\begin{array}{l}
 x' =  \frac{1}{12}\big(x+7y+4z -30\big) \\  
 y' =  \frac{1}{12}\big(x+y-2z\big) \\ 
 z' =  \frac{1}{12}\big(3x-3y +18\big) \\
\end{array}
\right. $$

\bigskip

Avec les matrices cela se fait ainsi :
le système $(\mathcal{S})$ devient
$$\begin{pmatrix} x+2 \\ y-4 \\ z-1 \end{pmatrix}= M \begin{pmatrix} x' \\ y' \\ z' \end{pmatrix}
\quad \text{ où } \quad 
M = \begin{pmatrix}
     1 & 2 & 3 \\ 1 & 2 & -1 \\ 1 & -4 & 1 \\
    \end{pmatrix}.
$$
Ainsi
$$\begin{pmatrix} x' \\ y' \\ z' \end{pmatrix}= M^{-1} \begin{pmatrix} x+2 \\ y-4 \\ z-1 \end{pmatrix}
\quad \text{ où } \quad 
M^{-1} = \frac{1}{12}\begin{pmatrix}
     1 & 7 & 4 \\ 1 & 1 & -2 \\ 3 & -3 & 0 \\
    \end{pmatrix}.
$$
\end{enumerate}
\fincorrection
\nocorrection
\correction{002039}
\begin{enumerate}
  \item 
Notons $P$ le plan d'équation $2x+2y-z=1$. Et soit $M_0=(x_0,y_0,z_0)$ un point quelconque.
$\vec n = (2, 2, -1)$ est un vecteur normal au plan. On cherche $p(M_0)$ appartenant au plan
sous la forme $M_0 + \lambda\cdot \vec n$.

\begin{align*}
p(M_0)\in P 
& \iff M_0 + \lambda\cdot \vec n \in P \\
& \iff (x_0,y_0,z_0)+ \lambda (2, 2, -1) \in P  \\
& \iff (x_0+2\lambda, y_0+2\lambda,z_0-\lambda) \in P \\
& \iff 2(x_0+2\lambda)+2(y_0+2\lambda)-(z_0-\lambda)=1 \\
& \iff \lambda = \frac{1-2x_0-2y_0+z_0}{9} \\
\end{align*}

En posant $\lambda_0= \frac{1-2x_0-2y_0+z_0}{9}$, 
le projeté orthogonal de $M_0$ sur $P$ est défini par 
$p(M_0)=(x_0+2\lambda_0,y_0+2\lambda_0,z_0-\lambda_0)$.
 
  \item
Notons $D$ la droite d'équation $\left\{\begin{array}{l}x+y+z=1 \\2x-z=2 \end{array}\right.$
et soit $M_0=(x_0,y_0,z_0)$ un point quelconque.

Il nous faut deux vecteurs normaux : par exemple $\vec{n_1}=(1,1,1)$
et $\vec{n_2}=(2,0,-1)$ (qui sont les vecteurs normaux aux deux plans définissant $D$).

On cherche le projeté orthogonal $\pi(M_0)$ sur la droite $D$ sous la forme
$M_0+\lambda_1 \vec{n_1}+\lambda_2\vec{n_2}$. On va déterminer $\lambda_1,\lambda_2 \in \Rr$ de sorte que ce point
appartienne à $D$.

\begin{align*}
\pi(M_0)\in D
& \iff M_0 + \lambda_1 \vec{n_1}+\lambda_2\vec{n_2}\in D \\
& \iff (x_0,y_0,z_0)+ \lambda_1 (1, 1, 1) +\lambda_2(2,0,-1) \in D  \\
& \iff (x_0+\lambda_1+2\lambda_2, y_0+\lambda_1,z_0+\lambda_1-\lambda_2) \in D \\
& \iff 
\left\{\begin{array}{l}
(x_0+\lambda_1+2\lambda_2)+(y_0+\lambda_1)+(z_0+\lambda_1-\lambda_2)=1 \\
2(x_0+\lambda_1+2\lambda_2)-(z_0+\lambda_1-\lambda_2)=2 \end{array}\right. \\
& \iff 
\left\{\begin{array}{l}
3\lambda_1+\lambda_2=1-x_0-y_0-z_0 \\
\lambda_1+5\lambda_2=2-2x_0+z_0 \\
\end{array}\right. \\
& \iff 
\lambda_1 = \frac{1}{14}\big(3-3x_0-5y_0-6z_0\big)\text{ et } \lambda_2 = \frac{1}{14}\big(5-5x_0+y_0+4z_0\big) \\
\end{align*}

Ainsi $\pi(M_0)= (x_0,y_0,z_0)+ \lambda_1 (1, 1, 1) +\lambda_2(2,0,-1)$
avec les valeurs de $\lambda_1,\lambda_2$ obtenues.
  
  \item Le principe est similaire, voici les étapes :
  \begin{enumerate} 
     \item Trouver une équation du plan. Un vecteur normal au plan est $\vec u \wedge \vec{u'}=(-1,-2,-2)$.
Donc le plan est d'équation $x+2y+2z-2=0$.
     \item Chercher le projeté d'un point $M_0=(x_0,y_0,z_0)$ sous la forme 
$M_0+\lambda \cdot \overrightarrow{AB}$. Trouver $\lambda_0$ de sorte que $M_0+\lambda_0 \cdot \overrightarrow{AB}$
appartiennent au plan. 
     \item On trouve $\overrightarrow{AB}=(-1,0,3)$ et $\lambda_0=-\frac15(x_0+2y_0+2z_0-2)$
et donc le projeté cherché est $p(M_0)=(x_0-\lambda_0,y_0,z_0+3\lambda_0)$.
   \end{enumerate}
  
\end{enumerate}


\fincorrection
\nocorrection
\nocorrection
\nocorrection
\nocorrection
\correction{004958}
Droite parallèle à $(AC)$ passant par
         $D = \text{Bar}(A:1, B:-\frac13)$.
\fincorrection
\correction{004959}
Pour $k=-1,-5$ : plan médiateur de $G = \text{Bar}(A:1,B:2,C:k)$ et
                 $I = \text{Bar}(D:1,E:1)$.\par
Pour $k\ne -1,-3,-5$ : sphère de centre $O = \text{Bar}(G:(3+k)^2, I:-4 )$.\par
Pour $k = -3$ : sphère de centre $I$.
\fincorrection
\correction{004960}
$221C_1 = (949, 149, -615)$, $93C_2 = (128, -71, 397)$.
\fincorrection
\correction{004961}
$A_P=(\frac12,\frac12,\lambda)$,
         $A_Q=(1,1-\frac\lambda2,\frac\lambda2)$,
         $A_R=(1-\frac\lambda2,1,\frac\lambda2)$,
         $A_S=(\frac{7-\lambda}{11},
               -\frac{1+3\lambda}{11},
               \frac{10\lambda-4}{11})$.
	       
        coplanaires $\Leftrightarrow 8\lambda+3\lambda^2+3\lambda^3 = 0 \Leftrightarrow \lambda=0$.
\fincorrection
\correction{004962}
\begin{enumerate}
  \item $\Omega:(1,1,1)$, $R = \sqrt5$.
  \item $(ABC) : x+y+z = 6$.\quad
             $(ABD) : 4x-2y+z = 3$.\quad
             $(ACD) : x+4y-2z = 3$.\quad
             $(BCD) : 7x+y-5z = 12$.
  \item $$I:(a,b,c) \Leftrightarrow \left\{\begin{array}{ccc} 6-a-b-c     &=& r\sqrt3    \cr
                                    4a-2b+c-3   &=& r\sqrt{21} \cr
                                    a+4b-2c-3   &=& r\sqrt{21} \cr
                                    12-7a-b+5c  &=& r\sqrt{75} \cr \end{array}\right.
              \Leftrightarrow \left\{\begin{array}{ccc} 2a &=& 9 - 2\sqrt 7\cr
                          2b &=& 6 - \sqrt7  \cr
                          2c &=& 3           \cr
                          2r &=& \sqrt{21}-2\sqrt3.\cr \end{array}\right.$$
\end{enumerate}
 \fincorrection
\correction{004963}
$H:(-18/54, 8/54, -28/54)$,  $K:(7/54, 13/54, -33/54)$.
\fincorrection
\correction{004964}
$H:(-2/19, 28/19, 35/19)$, $K:(-6/19, 40/19, 23/19)$, $d=4\sqrt{19}$.
\fincorrection
\correction{004965}
\begin{itemize}
  \item Soient $B'$, $C'$ les projetés orthogonaux de $C$, $D$ sur $(AB)$. 
Le triangle $ABC$ a pour aire $\frac12 AB \cdot CC'$.
Le triangle $ABD$ a pour aire $\frac12 AB \cdot DD'$.
Comme ces triangles ont la même aire alors $CC'=DD'$.

  \item Notons $I$ le milieu de $[C,D]$ et $I'$ son projeté orthogonal sur $(AB)$.
Par définition la droite $(II')$ est perpendiculaire à $(AB)$.
Montrons qu'elle est aussi perpendiculaire à $(CD)$.
Par Pythagore dans $I'CC'$ : $I'C^2 = I'C'^2 + CC'^2$.
Par Pythagore dans $IDD'$ : $I'D^2 = I'D'^2 + DD'^2$.
Comme $I$ est le milieu de $[C,D]$, alors par projection orthogonale, $I'$ est le milieu de $[C',D']$, donc $I'C'=I'D'$. D'autre part $CC'=DD'$ et ainsi $I'C=I'D$. Le triangle $I'CD$ étant isocèle en $I'$, alors $(II')$ est perpendiculaire à $(CD)$.

Bilan : on a montré que la perpendiculaire commune à $(AB)$ et $(CD)$ est $(II')$ et passe donc par le milieu de $[C,D]$.

\item On refait le même raisonnement en projetant cette fois sur $(CD)$, pour trouver que la perpendiculaire commune à $(AB)$ et $(CD)$ passe par le milieu $J$ de $[A,B]$. Comme la perpendiculaire commune est unique et coupe $(AB)$ et $I'$ et en $J$ alors $I'=J$. 
Bilan : la perpendiculaire commune à $(AB)$ et $(CD)$ est la droite $(IJ)$ où
$I$ est le milieu de $[C,D]$ et $J$ le milieu de $[A,B]$.

\item On considère le retournement d'axe $(IJ)$ (rotation d'angle $\pi$). Alors $A$ s'envoie sur $B$, $C$ s'envoie sur $D$. Le segment $[A,C]$ s'envoie sur $[B,D]$. Un retournement préserve les longueurs, donc $AC = BD$.

Conclusion : les deux côtés opposés $AC$ et $BD$ ont même longueur.
\end{itemize}

\fincorrection
\correction{004966}
$a/\sqrt2$.
\fincorrection
\correction{004967}
$d^2 = \frac{(x+2y-z+3)^2}6 + \frac{(3x+3z+11)^2}9$.
\fincorrection
\correction{004968}
$\begin{cases} 14x' = 13x - 2y - 3z + 4\cr
                  14y' = -2x + 10y - 6z + 8\cr
                  14z' = -3x - 6y + 5z +12.\cr\end{cases}$
\fincorrection
\correction{004969}
$D':(9/53, 40/53, -117/53)$, $E':(-14/53, 32/53, 23/53)$, $\vec u:(23, 8, -140)$.
\fincorrection
\correction{004970}
$2x-7y-z=-1$.
\fincorrection
\nocorrection
\nocorrection
\correction{004973}
$\cos\theta = 1 - \frac{1}{2\sin^2(\pi/5)} = -\frac{1}{\sqrt5}$.
\fincorrection
\correction{004974}
$x^2+y^2+z^2-10z=9$.
\fincorrection
\correction{004975}
Soient $A,B$ deux points de $S$ distincts. Intersection de $S$ avec
un plan passant par $A$ et $B$ $ \Rightarrow $ $S$ est réunion de cercles passant par
$A$ et $B$.
On considère le plan médiateur de $[A,B]$, $P$ qui coupe $S$ suivant un cercle
$\mathcal{C}$ de centre $O$.
Le plan $Q = (OAB)$ coupe $S$ suivant un cercle $\mathcal{C}'$. $\mathcal{C}$ et $\mathcal{C}'$ ont en
commun les points $C$, $D$.
$(CD)$ est médiatrice de $[A,B]$ dans $Q$ donc est un diamètre de $\mathcal{C}'$ et
passe par $O$, donc est aussi diamètre de $\mathcal{C}$.
Ainsi $\mathcal{C}$ et $\mathcal{C}'$ sont deux cercles de même centre et même rayon dans des
plans perpendiculaires.
En considérant les plans coupant $P$ et $Q$ à angle droit, on obtient que $S$
est une sphère de centre $O$.
\fincorrection
\correction{004976}
Soient $D,\vec u, \alpha$ l'axe, le vecteur, et l'angle de $g$.
Alors $f \circ g \circ f^{-1}$ est le vissage d'axe $f(D)$, de vecteur
$\vec f(\vec u)$, et d'angle $\alpha$.

On veut que ce soit $g$, donc $D$ est invariant par $f$, ce qui implique
que $D$ soit l'axe de $f$.

Réciproquement, si $f$ et $g$ ont même axe, alors ils commutent.
\fincorrection
\correction{004977}
Soit $v = \sigma_1 \circ \sigma_2$. (vissage autour de la perp. commune à
$D_1$ et $D_2$)

$\sigma_1 \circ \sigma_2 \circ \sigma_3$ est un $\frac12$-tour $ \Rightarrow $
$v \circ \sigma_3 \circ v \circ \sigma_3 = $id, donc
$\sigma_3 \circ v \circ \sigma_3^{-1} = v^{-1}$.

L'axe de $v$ est donc invariant par $\sigma_3$, donc parallèle ou
perpendiculaire à $D_3$.
Si parallèle, alors $\sigma_1 \circ \sigma_2 \circ \sigma_3$ est encore un
vissage $ \Rightarrow $ ne convient pas.
\fincorrection
\correction{004978}
$d_{AD}(D) = D,\quad d_{AC}(D) = D'$ tq $\vec{AD'} = \vec{DC}$. \par
La droite $(AD')$ est parallèle à $(DC)$ qui est perpendiculaire à $(AB)$.

Donc $f(D) = d_{AB}(D') = D''$ tq $\vec{AD''} = \vec{CD}$.

$d_{AD}(D'') = C,\quad d_{AC}(C) = C,\quad f(D'') = d_{AB}(C) = C'$
tq $\vec{AC'} = \vec{CB}$.

$d_{AD}(C') = C''$ tq $\vec{AC''} = \vec{BC},\quad d_{AC}(C'') = B,\quad
 f(C') = d_{AB}(B) = B$.

$d_{AD}(B) = B'$ tq $\vec{AB'} = \vec{BD},\quad d_{AC}(B') = B''$
tq $\vec{AB''} = \vec{DB},\quad f(B) = d_{AB}(B'') = D$.

Soit $E$ le symétrique de $C$ par rapport à $(BD)$ : $f$ est le
$\frac14$-tour autour de $(AE)$ envoyant $D$ sur $B$.
\fincorrection
\correction{004979}
$ABC$ équilatéral : 12, \quad $ABC$ isocèle : 4,\quad $ABC$ scalène : 2.
\fincorrection
\correction{004980}
\begin{enumerate}
  \item       $24$ éléments :
      \halign{ \hskip 5cm #\hfill&\quad(#)\cr
      identité &1\cr
      rotations autour de l'axe d'une face &8\cr
      symétries \% plan médiateur d'une arête &6\cr
      $\frac12$-tour autour de la perpendiculaire commune à deux arêtes opposées &3\cr
      $\frac14$-tour autour de la perp. \dots + symétrie \% plan médian &6\cr }
  \item       $48$ éléments :
      \halign{ \hskip 5cm #&\hfill#\hfill&\quad(#)\cr
      identité &&1\cr
      symétries &\% plan médian d'une face &3\cr
      &\% plan diagonal d'une face &6\cr
      &\% axe d'une face &3\cr
      &\% axe d'une arête &6\cr
      &\% centre du cube &1\cr
      rotation &$\pm 2\pi/3$ autour d'une diagonale &8\cr
      &$\pm \pi/2$ autour de l'axe d'une face &6\cr
      symétries-rotations &$\pm \pi/2$ \% axe face &6\cr
      &$\pm \pi/3$ \% diagonale &8\cr }

  \item       si les droites ne sont pas perpendiculaires :
      \halign{ \hskip 5cm #\hfill&\quad(#)\cr
      identité &1\cr
      $\frac12$-tour autour de la perpendiculaire commune &1\cr
      $\frac12$-tour autour d'une bissectrice &2\cr }
      si elles sont perpendiculaires, il y a aussi :
      \halign{ \hskip 5cm #\hfill&\quad(#)\cr
      symétrie \% plan contenant une droite et la perp. commune &2\cr
      symétrie-$\frac14$-tour autour de la perp. commune &2\cr }
\end{enumerate}
 \fincorrection
\correction{004981}
L'application $M_1 \mapsto M_4$ est affine donc est une
         homothétie-translation
         (dim 1) et le coefficient d'homothétie est strictement inférieur à 1.
\fincorrection
\correction{005501}
\textbullet~$\overrightarrow{n}=\overrightarrow{OA}\wedge\overrightarrow{OB}$ a pour coordonnées $(2,-3,-4)$. Ce vecteur n'est pas nul. Par suite, les points $O$, $A$ et $B$ ne sont pas alignés et le plan $(OAB)$ est bien défini. C'est le plan passant par $O$ et de vecteur normal $\overrightarrow{n}(2,-3,-4)$. Une équation cartésienne du plan $(OAB)$ est donc $2x-3y-4z=0$.
\textbullet~$\overrightarrow{n}'=\overrightarrow{OC}\wedge\overrightarrow{OD}$ a pour coordonnées $(4,-9,-1)$. Ce vecteur n'est pas nul. Par suite, les points $O$, $C$ et $D$ ne sont pas alignés et le plan $(OCD)$ est bien défini. C'est le plan passant par $O$ et de vecteur normal $\overrightarrow{n}'(4,-9,-1)$. Une équation cartésienne du plan $(OAB)$ est donc $4x-9y-z=0$.
\textbullet$-\overrightarrow{n}\wedge\overrightarrow{n}'$ a pour coordonnées $(33,14,6)$. Ce vecteur n'est pas nul et on sait que les plans $(OAB)$ et $(OCD)$ sont sécants en une droite, à savoir la droite passant par $O(0,0,0)$ et de vecteur directeur $(33,14,6)$. Un système d'équations cartésiennes de cette droite est 
$\left\{
\begin{array}{l}
2x-3y-4z=0\\
4x-9y-z=0
\end{array}
\right.$.
\fincorrection
\correction{005502}
Les vecteurs $(2,-3,1)$ et $(1,2,0)$ ne sont pas colinéaires, de sorte que $(P)$ est bien un plan. Trouvons alors une équation cartésienne de $(P)$

\begin{align*}\ensuremath
M(x,y,z)\in(P)&\Leftrightarrow\exists(\lambda,\mu)\in\Rr^2/\;
\left\{
\begin{array}{l}
x=1+2\lambda+\mu\\
y=-1-3\lambda+2\mu\\
z=1+\lambda
\end{array}
\right.
\Leftrightarrow\exists(\lambda,\mu)\in\Rr^2/\;
\left\{
\begin{array}{l}
\lambda=z-1\\
x=1+2(z-1)+\mu\\
y=-1-3(z-1)+2\mu
\end{array}
\right.
\\
 &\Leftrightarrow\exists(\lambda,\mu)\in\Rr^2/\;
\left\{
\begin{array}{l}
\lambda=z-1\\
\mu=x-2z+1\\
y=-1-3(z-1)+2(x-2z+1)
\end{array}
\right.
\\
 &\Leftrightarrow
-2x+y+7z-4=0
\end{align*}
Soit alors $M(2+3t,-t,1+t)$, $t\in\Rr$, un point de $(D)$

\begin{align*}\ensuremath
M\in(P)&\Leftrightarrow-2(2+3t)+(-t)+7(1+t)-4=0
\Leftrightarrow
0\times t-1=0.
\end{align*}
Ce dernier système n'a pas de solution et donc $(D)\cap(P)=\varnothing$. La droite $(D)$ est strictement parallèle au plan $(P)$.

\begin{align*}\ensuremath
M(x,y,z)\in(P)\cap(P')&\Leftrightarrow\exists(\nu,\eta)\in\Rr^2/\;\left\{
\begin{array}{l}
x=-5-\nu\\
y=3+\nu+3\eta\\
z=\nu+\eta\\
-2x+y+7z-4=0
\end{array}
\right.\\
 &\Leftrightarrow\exists(\nu,\eta)\in\Rr^2/\;\left\{
\begin{array}{l}
x=-5-\nu\\
y=3+\nu+3\eta\\
z=\nu+\eta\\
-2(-5-\nu)+(3+\nu+3\eta)+7(\nu+\eta)-4=0
\end{array}
\right.\\
 &\Leftrightarrow\exists(\nu,\eta)\in\Rr^2/\;\left\{
\begin{array}{l}
\eta=-\nu-\frac{9}{10}\\
x=-5-\nu\\
y=3+\nu+3\left(-\nu-\frac{9}{10}\right)\\
z=\nu+\left(-\nu-\frac{9}{10}\right)
\end{array}
\right.
\Leftrightarrow
\exists\nu\in\Rr/\;
\left\{
\begin{array}{l}
x=-\nu-5\\
y=-2\nu+\frac{3}{10}\\
z=-\frac{9}{10}
\end{array}
\right.
\end{align*}
$(P)$ et $(P')$ sont donc sécants en la droite passant par le point $\left(-5,\frac{3}{10},-\frac{9}{10}\right)$ et de vecteur directeur $(1,2,0)$.
\fincorrection
\correction{005503}
Soit $r$ la rotation cherchée. Notons $u$ le vecteur $\frac{1}{3}(1,2,2)$ ($u$ est unitaire) et $\theta$ l'angle de $r$. $r$ est la rotation d'angle $\theta$ autour du vecteur unitaire $u$.  On sait que pour tout vecteur $v$ de $\Rr^3$

\begin{center}
$r(v)=(\cos\theta)v+(1-\cos\theta)(v.u)u+(\sin\theta) u\wedge v\quad(*)$
\end{center}
et en particulier que $[v,r(v),u]=\sin\theta\|v\wedge u\|^2$. L'égalité $r(j)=k$ fournit

\begin{center}
$\sin\theta\|j\wedge u\|^2=\left[j,r(j),u\right]=\left[u,j,k\right]=\frac{1}{3}\left|\begin{array}{ccc}
1&0&0\\
2&1&0\\
3&0&1
\end{array}
\right|=\frac{1}{3}$.
\end{center}

Comme $u\wedge j=\frac{1}{3}(i+2j+2k)\wedge j=-\frac{2}{3}j+\frac{1}{3}k$, on a $\|j\wedge u\|^2=\frac{5}{9}$ et donc $\sin\theta=\frac{3}{5}$. L'égalité $r(j)=k$ fournit ensuite

\begin{center}
$k=(\cos\theta)j+(1-\cos\theta)\times\frac{2}{3}\times\frac{1}{3}(i+2j+2k)+\frac{3}{5}\times\frac{1}{3}(i+2j+2k)\wedge j$
\end{center}
En analysant la composante en $i$, on en déduit que $\frac{2}{9}(1-\cos\theta)-\frac{2}{5}=0$ et donc $\cos\theta=-\frac{4}{5}$. Ainsi, pour tout vecteur $v=(x,y,z)$ de $\Rr^3$, l'égalité $(*)$ s'écrit

\begin{align*}\ensuremath
r(v)&=-\frac{4}{5}(x,y,z)+\frac{9}{5}\times\frac{1}{3}\times\frac{1}{3}(x+2y+2z)(1,2,2)+\frac{3}{5}\times\frac{1}{3}
(2z-2y,2x-z,-2x+y)\\
 &=\frac{1}{5}(-4x+(x+2y+2z)+(2z-2y),-4y+2(x+2y+2z)+(2x-z),-4z+2(x+2y+2z)+(-2x+y))\\
 &=\frac{1}{5}(-3x+4z,4x+3z,5y)=\frac{1}{5}\left(
 \begin{array}{ccc}
 -3&0&4\\
 4&0&3\\
 0&5&0
 \end{array}
 \right)\left(
 \begin{array}{c}
 x\\
 y\\
 z
 \end{array}
 \right)
\end{align*}
La matrice cherchée est

\begin{center}
\shadowbox{
$\left(
 \begin{array}{ccc}
 -\frac{3}{5}&0&\frac{4}{5}\\
\rule[-4mm]{0mm}{10mm} \frac{4}{5}&0&\frac{3}{5}\\
 0&1&0
 \end{array}
 \right)$.
 }
 \end{center}
\fincorrection
\nocorrection
\nocorrection
\nocorrection
\nocorrection
\nocorrection
\nocorrection
\nocorrection
\nocorrection
\nocorrection
\nocorrection
\nocorrection
\nocorrection
\nocorrection
\nocorrection
\nocorrection
\correction{007516}
On considère l'isobarycentre $G$ des points $A,B,C,D$.
Par associativité, $G$ est le barycentre des barycentres $I$ de $(A,1),(B,1)$ 
et $J$ de $(C,1),(D,1)$ affecté des masses $(I,1+1)$ et $(J,1+1)$.
Par conséquent, $G$ est au milieu de $I$, $J$.
En particulier, les droites joignant les milieux des cotés
opposés du tétraèdre sont concourantes en $G$.
\fincorrection
\correction{007517}
Une base orthonormée du sous-espace $vect(e_1, e_1+e_2+e_3)$ est 
$f_1=e_1,f_2=\frac{e_2+e_3}{\sqrt{2}}$.
Pour la compléter en une base orthonormée de $\R^3$, il suffit de considérer $f_3=f_1\wedge f_2=\frac{e_3-e_2}{\sqrt{2}}$.
\fincorrection
\correction{005540}
On note $\mathcal{C}$ la courbe considérée.
\begin{enumerate}
 \item
        \begin{enumerate}
        \item $\mathcal{C}$ est la parabole de sommet $O$, d'axe focal $(Ox)$, de paramètre $p=\frac{1}{2}$ tournée vers
les $x$ positifs. Son foyer est le point $F\left(\frac{1}{4},0\right)$ et sa directrice est $\mathcal{D}~:~x=-\frac{1}{4}$.
        \item $\mathcal{C}$ est la parabole de sommet $O$, d'axe focal $(Ox)$, de paramètre $p=\frac{1}{2}$ tournée vers
les $x$ négatifs. Son foyer est le point $F\left(-\frac{1}{4},0\right)$ et sa directrice est $\mathcal{D}~:~x=\frac{1}{4}$.
        \item $\mathcal{C}$ est la parabole de sommet $O$, d'axe focal $(Oy)$, de paramètre $p=\frac{1}{2}$ tournée vers
les $y$ positifs. Son foyer est le point $F\left(0,\frac{1}{4}\right)$ et sa directrice est $\mathcal{D}~:~y=-\frac{1}{4}$.
        \item $\mathcal{C}$ est la parabole de sommet $O$, d'axe focal $(Oy)$, de paramètre $p=\frac{1}{2}$ tournée vers
les $y$ négatifs. Son foyer est le point $F\left(0,-\frac{1}{4}\right)$ et sa directrice est $\mathcal{D}~:~y=\frac{1}{4}$.
        \end{enumerate}

 \item  
        \begin{enumerate}
        \item $\mathcal{C}$ est une ellipse, de centre $O$ avec $a=5>3=b$ et donc d'axe focal $(Ox)$.

Ses sommets sont $A(5,0)$, $A'(-5,0)$, $B(0,3)$ et $B'(0,-3)$.

$c=\sqrt{a^2-b^2}=4$ et donc les foyers sont $F(4,0)$ et $F'(-4,0)$.

L'excentricité $e$ vaut $e=\frac{c}{a}=\frac{4}{5}$.

Les directrices ont pour équations respectives $x=\frac{a}{e}=\frac{25}{4}$ et $x=-\frac{25}{4}$.
        \item $\mathcal{C}$ est une ellipse, de centre $O$ avec $a=3<5=b$ et donc d'axe focal $(Oy)$.

Ses sommets sont $A(3,0)$, $A'(-3,0)$, $B(0,5)$ et $B'(0,-5)$.

$c=\sqrt{b^2-a^2}=4$ et donc les foyers sont $F(0,4)$ et $F'(0,-4)$.

L'excentricité $e$ vaut $e=\frac{c}{b}=\frac{4}{5}$.

Les directrices ont pour équations respectives $y=\frac{b}{e}=\frac{25}{4}$ et $y=-\frac{25}{4}$.
        \item $x^2+2y^2=1\Leftrightarrow\frac{x^2}{1^2}+\frac{y^2}{\left(\frac{1}{\sqrt{2}}\right)^2}=1$.

$\mathcal{C}$ est une ellipse, de centre $O$ avec $a=1>\frac{1}{\sqrt{2}}=b$ et donc d'axe focal $(Ox)$.

Ses sommets sont $A(1,0)$, $A'(-1,0)$, $B\left(0,\frac{1}{\sqrt{2}}\right)$ et $B'\left(0,-\frac{1}{\sqrt{2}}\right)$.

$c=\sqrt{a^2-b^2}=\frac{1}{\sqrt{2}}$ et donc les foyers sont $F\left(\frac{1}{\sqrt{2}},0\right)$ et
$F'\left(-\frac{1}{\sqrt{2}},0\right)$.

L'excentricité $e$ vaut $e=\frac{c}{a}=\frac{1}{\sqrt{2}}$.

Les directrices ont pour équations respectives $x=\frac{a}{e}=\sqrt{2}$ et $x=-\sqrt{2}$.
        \end{enumerate}
 \item 
        \begin{enumerate}
        \item $\mathcal{C}$ est une hyperbole de centre $O$ et d'axe focal $(Ox)$ avec $a=4$ et $b=3$ et donc
$c=\sqrt{a^2+b^2}=5$, puis $e=\frac{c}{a}=\frac{5}{4}$.

Les sommets sont $A(4,0)$ et $A'(-4,0)$ et les foyers sont $F(5,0)$ et $F(-5,0)$.

Les directrices sont les droites d'équations respectives $x=\frac{a}{e}=\frac{16}{5}$ et $x=-\frac{16}{5}$.

Les asymptotes sont les les droites d'équations respectives $y=\frac{3}{4}x$ et $y=-\frac{3}{4}x$.
        \item $\mathcal{C}$ est une hyperbole de centre $O$ et d'axe focal $(Oy)$ avec $a=4$ et $b=3$ et donc
$c=\sqrt{a^2+b^2}=5$, puis $e=\frac{c}{b}=\frac{5}{3}$.

Les sommets sont $B(0,3)$ et $B'(0,-3)$ et les foyers sont $F(0,5)$ et $F(0,-5)$.

Les directrices sont les droites d'équations respectives $y=\frac{b}{e}=\frac{9}{5}$ et $y=-\frac{9}{5}$.

Les asymptotes sont les droites d'équations respectives $y=\frac{3}{4}x$ et $y=-\frac{3}{4}x$.

$$\includegraphics{../images/img005540-1}$$


        \item $\mathcal{C}$ est une hyperbole de centre $O$ et d'axe focal $(Ox)$ avec $a=b=1$ et donc
$c=\sqrt{2}$, puis $e=\sqrt{2}$.

Les sommets sont $A(1,0)$ et $A'(-1,0)$ et les foyers sont $F(\sqrt{2},0)$ et $F(-\sqrt{2},0)$.

Les directrices sont les droites d'équations respectives $x=\frac{1}{\sqrt{2}}$ et $y=-\frac{1}{\sqrt{2}}$.

Les asymptotes sont les les droites d'équations respectives $y=x$ et $y=-x$.
         \end{enumerate}
\end{enumerate}
\fincorrection
\correction{005541}
\begin{enumerate}
 \item 
  \begin{enumerate}
   \item $y=x^2+x+1\Leftrightarrow y=\left(x+\frac{1}{2}\right)^2+\frac{3}{4}\Leftrightarrow\left(y-\frac{3}{4}\right)=\left(x+\frac{1}{2}\right)^2$.
$\mathcal{C}$ est la parabole de sommet $S\left(-\frac{1}{2},\frac{3}{4}\right)$, d'axe focal la droite d'équation
$x=-\frac{1}{2}$, de paramètre $p=\frac{1}{2}$ et donc de foyer
$F\left(-\frac{1}{2},\frac{3}{4}+\frac{1}{4}\right)=\left(-\frac{1}{2},1\right)$ et de directrice d'équation
$y=\frac{3}{4}-\frac{1}{4}=\frac{1}{2}$.

     \item $y^2+y-2x=0\Leftrightarrow\left(y+\frac{1}{2}\right)^2-\frac{1}{4}-2x=0\Leftrightarrow\left(y+\frac{1}{2}\right)^2=2\left(x+\frac{1}{8}\right)$.
$\mathcal{C}$ est la parabole de sommet $S\left(-\frac{1}{8},-\frac{1}{2}\right)$, d'axe focal la droite d'équation
$y=-\frac{1}{2}$, de paramètre $p=1$ et donc de foyer
$F\left(-\frac{1}{8}+\frac{1}{2},-\frac{1}{2}\right)=\left(\frac{3}{8},-\frac{1}{2}\right)$ et de directrice d'équation
$x=-\frac{1}{8}-\frac{1}{2}=-\frac{5}{8}$.
    \item $y=\sqrt{2x+3}\Leftrightarrow y^2=2\left(x+\frac{3}{2}\right)\;\mbox{et}\;y\geq0$. $\mathcal{C}$ est une demi-parabole de sommet
$S\left(-\frac{3}{2},0\right)$, d'axe focal $(Ox)$, de paramètre $p=1$ et donc de foyer $F\left(-\frac{3}{2}+\frac{1}{2},0\right)=(-1,0)$ et
de directrice d'équation $x=-\frac{3}{2}-\frac{1}{2}=-2$.
     \end{enumerate}
 \item 
  \begin{enumerate}
   \item $x^2+x+2y^2+y=0\Leftrightarrow\left(x+\frac{1}{2}\right)^2+2\left(y+\frac{1}{4}\right)^2=\frac{3}{8}
\Leftrightarrow
\frac{\left(x+\frac{1}{2}\right)^2}{\left(\sqrt{\frac{3}{8}}\right)^2}+\frac{\left(y+\frac{1}{4}\right)^2}{\left(\frac{\sqrt{3}}{4}\right)^2}=1
$.
$\mathcal{C}$ est une ellipse. Centre~:~$\left(-\frac{1}{2},-\frac{1}{4}\right)$. $a=\sqrt{\frac{3}{8}}>\frac{\sqrt{3}}{4}=b$. Axe
focal~:~$y=-\frac{1}{4}$.
Sommets~:~$A\left(-\frac{1}{2}+\sqrt{\frac{3}{8}},-\frac{1}{4}\right)$,
$A'\left(-\frac{1}{2}-\sqrt{\frac{3}{8}},-\frac{1}{4}\right)$, $B\left(-\frac{1}{2},-\frac{1}{4}+\frac{\sqrt{3}}{4}\right)$ et
$B'\left(-\frac{1}{2},-\frac{1}{4}-\frac{\sqrt{3}}{4}\right)$.
$c=\sqrt{a^2-b^2}=\frac{\sqrt{3}}{4}$. $e=\frac{c}{a}=\frac{1}{\sqrt{2}}$.
Foyers~:~$F\left(-\frac{1}{2}+\frac{\sqrt{3}}{4},-\frac{1}{4}\right)$ et $F'\left(-\frac{1}{2}-\frac{\sqrt{3}}{4},-\frac{1}{4}\right)$.
Directrices~:~$x=-\frac{1}{2}+\frac{\sqrt{3}}{2}$ et $x=-\frac{1}{2}-\frac{\sqrt{3}}{2}$.
       \item $y=-2\sqrt{-x^2+x}\Leftrightarrow y^2=4(-x^2+x)\;\mbox{et}\;y\leq0\Leftrightarrow\frac{\left(x-\frac{1}{2}\right)^2}{\left(\frac{1}{2}\right)^2}+y^2=1\;\mbox{et}\;y\leq0$.
$\mathcal{C}$ est une demi-ellipse. Centre~:~$\left(\frac{1}{2},0\right)$. $a=\frac{1}{2}<1=b$. Axe
focal~:~$x=0$.
Sommets~:~$A\left(1,0\right)$,
$A'\left(0,0\right)$ et
$B'\left(\frac{1}{2},-1\right)$.
$c=\sqrt{b^2-a^2}=\frac{\sqrt{3}}{2}$. $e=\frac{c}{b}=\frac{\sqrt{3}}{2}$.
Foyers~:~$F\left(\frac{1}{2},\frac{\sqrt{3}}{2}\right)$ et $F'\left(\frac{1}{2},-\frac{\sqrt{3}}{2}\right)$.
Directrices~:~$y=\frac{2}{\sqrt{3}}$ et $y=-\frac{2}{\sqrt{3}}$.

     \end{enumerate}

 \item  $x^2-y^2+x+y+1=0\Leftrightarrow\left(x+\frac{1}{2}\right)^2-\left(y-\frac{1}{2}\right)^2=-1$.
$\mathcal{C}$ est une hyperbole de centre
$\left(-\frac{1}{2},\frac{1}{2}\right)$ et d'axe focal la droite d'équation $x=-\frac{1}{2}$. $a=b=1$.
Sommets : $B\left(-\frac{1}{2},\frac{3}{2}\right)$ et $B'\left(-\frac{1}{2},-\frac{1}{2}\right)$.
$c=\sqrt{a^2+b^2}=\sqrt{2}$ puis $e=\frac{c}{b}=\sqrt{2}$. Foyers : $F\left(-\frac{1}{2},\frac{1}{2}+\sqrt{2}\right)$ et $F'\left(-\frac{1}{2},\frac{1}{2}-\sqrt{2}\right)$. Directrices : $y=\frac{1}{2}+\frac{1}{\sqrt{2}}$ et $y=\frac{1}{2}-\frac{1}{\sqrt{2}}$.

Asymptotes : $y=x+1$ et $y=-x$.
\end{enumerate}
\fincorrection
\correction{005542}
\begin{enumerate}
 \item  On note $\mathcal{H}$ l'hyperbole considérée. On tourne de $\frac{\pi}{4}$. Pour cela, on pose $\left\{
\begin{array}{l}
x=\frac{1}{\sqrt{2}}(X-Y)\\
y=\frac{1}{\sqrt{2}}(X+Y)
\end{array}
\right.$. On a alors

$$y=\frac{1}{x}\Leftrightarrow xy=1\Leftrightarrow \frac{1}{2}(X-Y)(X+Y)=1\Leftrightarrow\frac{X^2}{\left(\sqrt{2}\right)^2}-\frac{Y^2}{\left(\sqrt{2}\right)^2}=1.$$
Ainsi, si $\mathcal{R}$ est le repère orthonormé initial $\left(O,\overrightarrow{i},\overrightarrow{j}\right)$ et $\mathcal{R}'$ est le repère
$\left(O,\overrightarrow{I},\overrightarrow{J}\right)$ où $\overrightarrow{I}=\frac{1}{\sqrt{2}}(\overrightarrow{i}+\overrightarrow{j})$ et
$\overrightarrow{J}=\frac{1}{\sqrt{2}}(-\overrightarrow{i}+\overrightarrow{j})$, une équation de $\mathcal{H}$ dans $\mathcal{R}$ est $xy=1$ et une
équation de $\mathcal{H}$ dans $\mathcal{R}'$ est $\frac{X^2}{(\sqrt{2})^2}-\frac{Y^2}{(\sqrt{2})^2}=1$. On obtient
$a=b=\sqrt{2}$, $c=\sqrt{a^2+b^2}=2$ et $e=\frac{c}{a}=\sqrt{2}$.
Les formules de changement de repère s'écrivent $\left\{
\begin{array}{l}
x=\frac{1}{\sqrt{2}}(X-Y)\\
y=\frac{1}{\sqrt{2}}(X+Y)
\end{array}
\right.$ et les formules inverses s'écrivent

$\left\{
\begin{array}{l}
X=\frac{1}{\sqrt{2}}(x+y)\\
y=\frac{1}{\sqrt{2}}(-x+y)
\end{array}
\right.$ (dans ce qui suit, les coordonnées d'un point dans $\mathcal{R}'$ seront notées avec $\mathcal{R}'$ en indice
alors que les coordonnées dans $\mathcal{R}$ seront notées sans écrire $\mathcal{R}$ en indice).

$\shadowbox{\mbox{Centre}\;O(0,0)}$.

Asymptotes~:~bien sûr, les axes $(Ox)$ et $(Oy)$.

Axe focal~:~l'axe $(OX)$ ou encore la droite d'équation $y=x$ (dans $\mathcal{R}$).

Sommets~:~$A(\sqrt{2},0)_{\mathcal{R}'}$, $A'(-\sqrt{2},0)_{\mathcal{R}'}$ et donc
$\shadowbox{\mbox{Sommets}\;A(1,1)\;\mbox{et}\;A'(-1,-1)}$.

Foyers~:~$F(2,0)_{\mathcal{R}'}$, $F'(-2,0)_{\mathcal{R}'}$ et
donc \shadowbox{Foyers\;$F(\sqrt{2},\sqrt{2})$\; et \;$F'(-\sqrt{2},-\sqrt{2})$.}

Directrices~:~les droites d'équations $X=\pm\frac{a}{e}=\pm1$ et donc dans $\mathcal{R}$, les droites d'équations
respectives $x+y=\pm\sqrt{2}$.

%41x^2-24xy+34y^2-106x+92y+74=0
 \item  Le discriminant de cette conique vaut $41\times34-12^2=1250>0$. Il s'agit donc d'une conique du genre ellipse.
On pose $\left\{
\begin{array}{l}
x=\cos(\theta)X-\sin(\theta)Y\\
y=\sin(\theta)X+\cos(\theta)Y
\end{array}
\right.$ et on détermine $\theta$ (ou plutôt $\cos\theta$ et $\sin\theta$) de sorte que le terme en $XY$ disparaisse.
Mais, le coefficient de $XY$ dans

$$41x^2-24xy+34y^2=41(\cos(\theta)X-\sin(\theta)Y)^2-24(\cos(\theta)X-\sin(\theta)Y)(\sin(\theta)X+\cos(\theta)Y)+
34(\sin(\theta)X+\cos(\theta)Y)^2,$$
vaut

$$-82\cos\theta\sin\theta-24(\cos^2\theta-\sin^2\theta)+68\cos\theta\sin\theta=-24(\cos^2\theta-\sin^2\theta)
-14\cos\theta\sin\theta.$$
Ce coefficient est nul si et seulement si $-12\cos^2\theta+12\sin^2\theta-7\cos\theta\sin\theta=0$ ou encore, après division
par $\cos^2\theta$, $12\tan^2\theta-7\tan\theta-12=0$. On peut alors prendre $\tan\theta=\frac{4}{3}$, puis on peut prendre
$\cos\theta=\frac{1}{\sqrt{1+\tan^2\theta}}=\frac{3}{5}$ et
$\sin\theta=\cos\theta\tan\theta=\frac{3}{5}\frac{4}{3}=\frac{4}{5}$.

Posons donc $\left\{
\begin{array}{l}
x=\frac{3X-4Y}{5}\\
y=\frac{4X+3Y}{5}
\end{array}\right.$ $(*)$. On a alors

\begin{align*}\ensuremath
41x^2-24xy+34y^2-106x+92y+74&=\frac{1}{25}(41(3X-4Y)^2-24(3X-4Y)(4X+3Y)+34(4X+3Y)^2\\
 &\;-530(3X-4Y)+460(4X+3Y)+1850)\\
 &=\frac{1}{25}(625X^2+1250Y^2+250X+3500Y+1850)\\
 &=25\left(X^2+2Y^2+\frac{2}{5}X+\frac{28}{5}Y+\frac{74}{25}\right)
\end{align*}
Une équation de la courbe dans le repère défini par $(*)$ est donc
$X^2+2Y^2+\frac{2}{5}X+\frac{28}{5}Y+\frac{74}{25}=0$. Ensuite,

\begin{center}
$X^2+2Y^2+\frac{2}{5}X+\frac{28}{5}Y+\frac{74}{25}=0\Leftrightarrow\left(X+\frac{1}{5}\right)^2+2\left(Y+\frac{7}{5}\right)^2=1\Leftrightarrow\left(X+\frac{1}{5}\right)^2+\frac{\left(Y+\frac{7}{5}\right)^2}{\left(\sqrt{2}\right)^2}=1$.
\end{center}
$\mathcal{C}$ est une ellipse. On trouve $a=1$, $b=\sqrt{2}$, $c=1$ $e=\frac{1}{\sqrt{2}}$ puis
Centre $\Omega\left(1,-1\right)$. Axe focal : $3x+4y+1=0$ et axe non focal : $-4x+3y+7=0$.

Sommets : $A\left(\frac{8}{5},-\frac{1}{5}\right)$, $A'\left(\frac{2}{5},-\frac{9}{5}\right)$, $B\left(1-\frac{4\sqrt{2}}{5},-1+\frac{3\sqrt{2}}{5}\right)$ et $B'\left(1+\frac{4\sqrt{2}}{5},-1-\frac{3\sqrt{2}}{5}\right)$.

Foyers : $F\left(\frac{1}{5},-\frac{2}{5}\right)$ et $F'\left(\frac{9}{5},-\frac{8}{5}\right)$. Directrices : $4x-3y+3=0$ et $4x-3y+17=0$.
 \item  $x^2+2xy+y^2=(x+y)^2$. On pose donc $\left\{
\begin{array}{l}
X=\frac{1}{\sqrt{2}}(x-y)\\
Y=\frac{1}{\sqrt{2}}(x+y)
\end{array}
\right.$ ou encore $\left\{
\begin{array}{l}
x=\frac{1}{\sqrt{2}}(X+Y)\\
y=\frac{1}{\sqrt{2}}(-X+Y)
\end{array}
\right.$.

\begin{align*}\ensuremath
x^2+2xy+y^2+3x-2y+1=0&\Leftrightarrow 2Y^2+\frac{3}{\sqrt{2}}(X+Y)-\frac{2}{\sqrt{2}}(-X+Y)+1=0\\
 &\Leftrightarrow 2\left(Y+\frac{1}{4\sqrt{2}}\right)^2+\frac{5}{\sqrt{2}}X+\frac{15}{16}=0\Leftrightarrow\left(Y+\frac{1}{4\sqrt{2}}\right)^2=-\frac{5}{2\sqrt{2}}\left(X+\frac{3}{8\sqrt{2}}\right).
\end{align*}
$\mathcal{C}$ est une parabole de paramètre $p=\frac{5}{4\sqrt{2}}$.

Sommet : $S\left(-\frac{5}{16},\frac{1}{16}\right)$. Axe focal : $x+y+\frac{1}{4}=0$.

Foyer : $F\left(-\frac{5}{8},\frac{3}{8}\right)$. Directrice : $x-y-\frac{1}{4}=0$.
 \item  $\mathcal{C}$ est le point d'intersection des droites d'équation $x-y+1=0$ et $x+y-1=0$ c'est-à-dire le point de coordonnées $(0,1)$.

 \item  $x^2+y^2-3x-y+3=\left(x-\frac{3}{2}\right)^2+\left(y-\frac{1}{2}\right)^2+\frac{1}{2}>0$ et donc $\mathcal{C}$ est vide.
 \item  $x(x-1)+(y-2)(y-3)=0$ est une équation du cercle de diamètre $[AB]$ où $A(0,2)$ et $B(1,3)$.
 \item  Si on pose $\left\{
\begin{array}{l}
X=x+y+1\\
Y=x-y+3
\end{array}
\right.$, on effectue un changement de repère non orthonormé. Dans le nouveau repère, $\mathcal{C}$ admet pour équation cartésienne $XY=3$ et donc $\mathcal{C}$ est une hyperbole. Avec le changement de repère effectué, on obtient directement les éléments affines de cette hyperbole mais pas ses éléments métriques : hyperbole d'asymptotes les droites d'équations $x+y+1=0$ et $x-y+3=0$ et donc de centre $(-2,1)$.
Pour obtenir l'axe focal, l'excentricité, les foyers et les directrices il faut faire un changement de repère orthonormé.
 \item   Si on pose $\left\{
\begin{array}{l}
X=2x+y+1\\
Y=3x+3y
\end{array}
\right.$, $\mathcal{C}$ admet pour équation cartésienne dans le nouveau repère $Y=X^2$ et donc $\mathcal{C}$ est une parabole. Pour obtenir ces éléments métriques, il faut un changement de repère orthonormé.
\end{enumerate}
\fincorrection
\correction{005543}
Etudier les courbes dont une équation polaire (en repère orthonormé direct) est
$$\begin{array}{ccccc}
1)\;r=\frac{1}{1+2\cos\theta}&2)\;r=\frac{1}{1+\cos\theta}&3)\;r=\frac{1}{2+\cos\theta}&4)\;r=\frac{1}{1-\sin\theta}
&5)\;r=\frac{1}{2-\cos\theta}.
\end{array}$$
\begin{enumerate}
 \item  $\mathcal{C}$ est une conique d'excentricité $2$ et donc une hyperbole.

$$\includegraphics{../images/img005543-1}$$

L'axe focal est $(Ox)$. Les sommets sont les points d'intersection de $\mathcal{C}$ et $(Ox)$ c'est-à-dire les points $M(0)$ et $M(\pi)$ de coordonnées cartésiennes respectives $A'\left(\frac{1}{3},0\right)$ et $A(1,0)$. Le centre $\Omega$ est le milieu de $[AA']$ c'est-à-dire $\Omega\left(\frac{2}{3},0\right)$.
L'un des foyers est $F'=O$ et l'autre est le symétrique de $F'$ par rapport à $\Omega$ : c'est le point $F\left(\frac{4}{3},0\right)$. Puisque $a=\frac{1}{3}$ et $e=2$, les directrices sont les droites d'équation $x=x_\Omega-\frac{a}{e}=\frac{1}{2}$ et $x=\frac{5}{6}$.
Les branches infinies sont obtenues pour $\theta=\pm\frac{2\pi}{3}$. Les asymptotes sont donc les droites passant par $\Omega$ d'angle polaire $\pm\frac{2\pi}{3}$. Ce sont les droites d'équations cartésiennes $y=\pm\sqrt{3}\left(x-\frac{2}{3}\right)$.
 \item  $\mathcal{C}$ est une conique d'excentricité $1$ et donc une parabole.

$$\includegraphics{../images/img005543-2}$$

L'axe focal est $(Ox)$. Le sommet est le point $M(0)$ de coordonnées cartésiennes $S\left(\frac{1}{2},0\right)$.
Le foyer est $F=O$. Le point $K$ est le symétrique de $F$ par rapport à $S$ et a pour coordonnées $(1,0)$. La directrice a donc pour équation $x=1$.
 \item  $\mathcal{C}$ est une conique d'excentricité $e=\frac{1}{2}$ et donc une ellipse.

$$\includegraphics{../images/img005543-3}$$

L'axe focal est $(Ox)$. Les sommets sur cet axe sont $A=M(0)$ de coordonnées $\left(\frac{1}{3},0\right)$ et $A'=M(\pi)$ de coordonnées $\left(-1,0\right)$. Le centre $\Omega$ est le milieu de $[AA']$ et a pour coordonnées $\left(-\frac{1}{3},0\right)$.
L'un des foyers est $F=O$. L'autre est le symétrique de $F$ par rapport à $\Omega$ : c'est le point $F'$ de coordonnées $\left(-\frac{2}{3},0\right)$. Par suite, $c=\frac{1}{3}$ $b=\sqrt{a^2-c^2}=\sqrt{\left(\frac{2}{3}\right)^2-\left(\frac{1}{3}\right)^2}=\frac{1}{\sqrt{3}}$. D'où les sommets $B\left(-\frac{1}{3},\frac{1}{\sqrt{3}}\right)$ et $B'\left(-\frac{1}{3},-\frac{1}{\sqrt{3}}\right)$.
Les directrices sont les droites d'équations $x=x_\Omega+\frac{a}{e}=1$ et $x=-\frac{5}{3}$.
 \item  $M\left(\theta-\frac{\pi}{2}\right)=\left[r\left(\theta-\frac{\pi}{2}\right),\theta-\frac{\pi}{2}\right]=\left[\frac{1}{1+\cos\theta},\theta-\frac{\pi}{2}\right]=\text{rot}_{O,-\pi/2}\left(\left[\frac{1}{1+\cos\theta},\theta\right]\right)$. Donc $\mathcal{C}$ est l'image de la parabole d'équation polaire $r=\frac{1}{1+\cos\theta}$ par le quart de tour indirect de centre $O$.

$$\includegraphics{../images/img005543-4}$$

 \item  $M\left(\theta+\pi\right)=\left[r\left(\theta+\pi\right),\theta+\pi\right]=\left[\frac{1}{2+\cos\theta},\theta+\pi\right]=s_O\left(\left[\frac{1}{2+\cos\theta},\theta\right]\right)$. Donc $\mathcal{C}$ est l'image de l'ellipse d'équation polaire $r=\frac{1}{2+\cos\theta}$ par la symétrie centrale de centre $O$.

$$\includegraphics{../images/img005543-5}$$

\end{enumerate}
\fincorrection
\correction{005544}
Un point du plan est sur le cercle de centre $O$ et de rayon $1$ si et seulement si son affixe $z$
est de module $1$ ou encore si et seulement si il existe un réel $\theta$ tel que $z=e^{i\theta}$. Or, pour $\theta$
réel,

$$f(e^{i\theta})=\frac{1}{1+e^{i\theta}+e^{2i\theta}}=\frac{e^{-i\theta}}{e^{i\theta}+1+e^{-i\theta}}
=\overline{\left(\frac{e^{i\theta}}{1+2\cos\theta}\right)}.$$
L'ensemble cherché est donc la symétrique par rapport à $(Ox)$ de la courbe d'équation polaire
$r=\frac{1}{1+2\cos\theta}$. Cette dernière est une ellipse, symétrique par rapport à $(0x)$. Donc l'ensemble cherché
est l'ellipse d'équation polaire $r=\frac{1}{1+2\cos\theta}$ (voir l'exercice \ref{exo:routhe4}, 1)).
\fincorrection
\correction{005546}
\begin{enumerate}
 \item  Soit $M$ un point du plan.
\textbf{1er cas.} Supposons que $M\notin(AB)\cup(AC)\cup(BC)$. 

\begin{align*}\ensuremath
P,\;Q\;\text{et}\;R\;\text{alignés}\Leftrightarrow\left(\overrightarrow{PQ},\overrightarrow{PR}\right)=0\;[\pi]
\Leftrightarrow\left(\overrightarrow{PQ},\overrightarrow{PM}\right)=\left(\overrightarrow{PR},\overrightarrow{PM}\right)\;[\pi].
\end{align*}
Maintenant, puisque les triangles $MPC$ et $MQC$ sont rectangles en $P$ et $Q$ respectivement, les points $P$ et $Q$ sont sur le cercle de diamètre $[MC]$. On en déduit que $\left(\overrightarrow{PQ},\overrightarrow{PM}\right)=\left(\overrightarrow{CQ},\overrightarrow{CM}\right)\;[\pi]$. De même, $\left(\overrightarrow{PR},\overrightarrow{PM}\right)=\left(\overrightarrow{BR},\overrightarrow{BM}\right)\;[\pi]$. Par suite,

\begin{align*}\ensuremath
P,\;Q\;\text{et}\;R\;\text{alignés}&\Leftrightarrow\left(\overrightarrow{CQ},\overrightarrow{CM}\right)=\left(\overrightarrow{BR},\overrightarrow{BM}\right)\;[\pi]\\
 &\Leftrightarrow\left(\overrightarrow{CA},\overrightarrow{CM}\right)=\left(\overrightarrow{BA},\overrightarrow{BM}\right)\;[\pi]\\
 &\Leftrightarrow M\;\text{appartient au cercle circonscrit au triangle}\;ABC\;(\text{privé des points}\;A,\;B\;\text{et}\;C).
\end{align*}
\textbf{2ème cas.} Supposons par exemple que $M\in(AB)$. Dans ce cas, $M=R$. Si de plus $M$ n'est ni $A$, ni $B$, alors $M\neq P$ et $M\neq Q$ puis les droites $(MP)$ et $(MQ)$ sont perpendiculaires aux droites $(BC)$ et $(AC)$ respectivement. Si par l'absurde, les points $P$, $Q$ et $R$ sont alignés, on a $(MP)=(MQ)$ et donc $(AB)//(AC)$. Ceci est une contradiction.
Donc, si les points $P$, $Q$ et $R$ sont alignés, $M$ est l'un des trois points $A$, $B$ ou $C$. La réciproque est immédiate.
En résumant les deux cas,

\begin{center}
\shadowbox{
$P$, $Q$ et $R$ sont alignés si et seulement si $M$ est sur le cercle circonscrit au triangle $ABC$.
}
\end{center}

$$\includegraphics{../images/img005546-1}$$


 \item  \textbf{Parabole tangente aux trois côtés d'un triangle.} Commençons par rappeler une construction usuelle de la tangente en un point d'une parabole : le triangle $FMH$ est isocèle en $M$ et la tangente en $M$ à $\mathcal{P}$ est la médiatrice du segment $[FH]$. Par suite, le projeté orthogonal $P$ de $F$ sur la tangente $(T)$ est sur $5T_0)$ la tangente au sommet de la parabole $\mathcal{P}$.

$$\includegraphics{../images/img005546-2}$$

Soient $A$, $B$ et $C$ trois points non alignés. Si $\mathcal{P}$ est une parabole tangente aux droites $(BC)$, $(CA)$ et $(AB)$, les projetés orthogonaux $P$, $Q$ et $R$ de son foyer $F$ sur les droites  $(BC)$, $(CA)$ et $(AB)$ sont alignés sur la tangente au sommet de la parabole $\mathcal{P}$. D'après 1), le point $F$ est nécessairement sur le cercle circonscrit au triangle $ABC$.
Réciproquement, si $F$ est l'un des trois points $A$, $B$ ou $C$, $F$ n'est pas solution car une tangente à une parabole ne passe jamais par son foyer.
Soit donc $F$ un point du cercle circonscrit au triangle $ABC$ et distinct des points $A$, $B$ et $C$. Montrons alors qu'il existe une parabole de foyer $F$, tangente aux droites $(BC)$, $(CA)$ et $(AB)$.
On construit les projetés orthogonaux $P$, $Q$ et $R$ de $F$ sur les droites $(BC)$, $(CA)$ et $(AB)$. Ils sont alignés sur la droite de \textsc{Simson} $(T_0)$ de $F$ relativement au triangle $ABC$. La parabole de foyer $F$ et de tangente au sommet $(T_0)$ est solution du problème posé. La construction des  points de contact est fournie par le graphique de la page précédente : on construit les symétriques de $F$ par rapport aux points $P$, $Q$ et $R$ (ces symétriques sont sur la directrice) puis on remonte perpendiculairement à $(T_0)$ jusqu'à la parabole.

$$\includegraphics{../images/img005546-3}$$

\end{enumerate}
\fincorrection
\correction{005547}
On choisit un repère orthonormé dans lequel $A$ a pour coordonnées $(R,0)$ et $(\mathcal{C})$ a pour représentation paramétrique $\left\{
\begin{array}{l}
x=R\cos t\\
y=R\sin t
\end{array}
\right.$, $t\in\Rr$. Soit $P(R\cos t,R\sin t)$ un point de $(\mathcal{C})$.
La tangente $(D)$ à $(\mathcal{C})$ en $A$ est la droite d'équation $x=R$ et la tangente $(T)$ à $(\mathcal{C})$  en $P$ est la droite d'équation $x\cos t+y\sin t=R$. Quand $t\notin\pi\Zz$, $(T)$ recoupe $(D)$ en le point $S$ de coordonnées $\left(R,R\frac{1-\cos t}{\sin t}\right)$ ou encore $\left(R,R\tan\left(\frac{t}{2}\right)\right)$.

Une équation de la droite $(BS)$ est $-\tan\left(\frac{t}{2}\right)(x+R)+2y=0$. L'abscisse de $M$ est $R\cos t$ et donc

\begin{center}
$y_M=\frac{1}{2}\tan\left(\frac{t}{2}\right)(x_M+R)=\frac{1}{2}R\tan\left(\frac{t}{2}\right)(\cos t+1)=R\sin\left(\frac{t}{2}\right)\cos\left(\frac{t}{2}\right)=\frac{1}{2}R\sin t$.
\end{center}
L'ensemble des points $M$ est donc le support de l'arc $\left\{
\begin{array}{l}
x=R\cos t\\
y=\frac{1}{2}R\sin t
\end{array}
\right.$, $t\in\Rr$. C'est l'image du cercle $\mathcal{C}$ dans l'affinité de base $(AB)$, de direction $(D)$ et de rapport $\frac{1}{2}$ et donc une ellipse de grand axe $[AB]$.

$$\includegraphics{../images/img005547-1}$$

\fincorrection
\correction{005815}
Dans tout l'exercice, on pose $\mathcal{R}=(O,i,j)$ et $(\Gamma)$ l'ensemble considéré, d'équation $f(x,y)=0$.

\begin{enumerate}
 \item  Pour $(x,y)\in\Rr^2$, posons $f(x,y) =2x^2+6xy+5y^2+4x+6y+1$  et $Q((x,y)) = 2x^2+6xy+5y^2$.

 
Le discriminant de cette conique est $\Delta=2\times5 - 3^2 = 1>0$ et la courbe $(\Gamma)$ est du genre ellipse c'est-à-dire soit une ellipse, éventuellement un cercle, soit un point, soit l'ensemble vide.

\textbf{Point critique.}

\begin{center}
$\left\{
\begin{array}{l}
\rule[-5mm]{0mm}{0mm}\frac{\partial f}{\partial x}(x,y)=0\\
\frac{\partial f}{\partial y}(x,y)=0
\end{array}
\right.\Leftrightarrow
\left\{
\begin{array}{l}
4x+6y+4=0\\
6x+10y+6=0
\end{array}
\right.\Leftrightarrow x=-1\;\text{et}\;y=0$.
\end{center}

On note $\Omega$ le point de coordonnées $(-1,0)$ dans le repère $\mathcal{R}$.

\textbf{Réduction de $Q$ en base orthonormée.}
La matrice de $Q$ dans la base $(i,j)$ est $\left(
\begin{array}{cc}
2&3\\
3&5
\end{array}
\right)$. Son polynôme caractéristique est $\chi_A = X^2-7X+1$ et les valeurs propres de $A$ sont $\alpha=\frac{7+3\sqrt{5}}{2}$  et $\beta=\frac{7-3\sqrt{5}}{2}$.

$\text{Ker}(A-\alpha I_2)$ est la droite d'équation $-(1+\sqrt{5})x + 2y=0$ et est engendrée par le vecteur unitaire $e_1=\frac{1}{\sqrt{10+2\sqrt{5}}}(2,1+\sqrt{5})$. Puis $\text{Ker}(A-\alpha I_2)$ est la droite d'équation $-(1-\sqrt{5})x + 2y=0$ et est engendrée par le vecteur unitaire $e_2=\frac{1}{\sqrt{10-2\sqrt{5}}}(2,1-\sqrt{5})$.

\textbf{Equation réduite de $(\Gamma)$ dans $\mathcal{R}'=(\Omega,e_1,e_2))$.}

Les termes de degré $1$ disparaissent car $\Omega$ est l'origine de $\mathcal{R}'$ et d'autre part, $Q(xi+yj) = Q(Xe_1+Ye_2) =\alpha X^2+\beta Y^2$.

Il manque simplement la constante mais si on effectue le changement de variables $x = x_0+aX+bY$ et $y = y_0+cX+dY$, la constante est bien sûr $f((x_0,y_0))$. Donc une équation cartésienne de $(\Gamma)$ dans $\mathcal{R}'$ est $\alpha X^2+\beta Y^2+f((-1,0))= 0$ ce qui s'écrit encore

\begin{center}
$\frac{7+3\sqrt{5}}{2}X^2+\frac{7-3\sqrt{5}}{2}Y^2 = 1$.
\end{center}

\textbf{Eléménts caractéristiques de la courbe $(\Gamma)$ dans le repère $\mathcal{R}'$.}

$(\Gamma)$ est une ellipse de centre $\Omega$, d'axe focal $(\Omega Y)$ car $a=\frac{1}{\sqrt{\frac{7+3\sqrt{5}}{2}}}<\frac{1}{\sqrt{\frac{7-3\sqrt{5}}{2}}}=b$ et d'axe non focal $(\Omega X)$.

- \emph{Centre} $\Omega(0,0)_{\mathcal{R}'}$.
 
 
- \emph{Excentricité} $c^2=b^2-a^2 =\frac{2}{7-3\sqrt{5}}-\frac{2}{7+3\sqrt{5}}=3\sqrt{5}$ puis $e^2 =\left(\frac{c}{b}\right)^2=\frac{3\sqrt{5}}{7+3\sqrt{5}}=\frac{-45+21\sqrt{5}}{4}$ et 

$e=\frac{1}{2}\sqrt{-45+21\sqrt{5}}= 0,69...$
 
- \emph{Sommets} $A\left(\sqrt{\frac{7-3\sqrt{5}}{2}},0\right)_{\mathcal{R}'}$  $A'\left(-\sqrt{\frac{7-3\sqrt{5}}{2}},0\right)_{\mathcal{R}'}$  $B\left(0,\sqrt{\frac{7+3\sqrt{5}}{2}}\right)_{\mathcal{R}'}$  $B'\left(0,-\sqrt{\frac{7+3\sqrt{5}}{2}}\right)_{\mathcal{R}'}$ 

- \emph{Foyers} $\Omega F =\Omega F'=c=\sqrt{3\sqrt{5}}$  et puisque $(\Omega Y)$ est l'axe focal, $F(0,\sqrt{3\sqrt{5}})$ et $F'(0,-\sqrt{3\sqrt{5}})$. 

- \emph{Directrices} $\Omega K =\Omega K'=\frac{b}{e}=2\sqrt{\frac{2}{(7-3\sqrt{5})(-45+21\sqrt{5})}}=\frac{2}{7-3\sqrt{5}}\sqrt{\frac{2}{3\sqrt{5}}}=\frac{7+3\sqrt{5}}{\sqrt{6\sqrt{5}}}$ et donc 

$(D)$ : $Y =\frac{7+3\sqrt{5}}{\sqrt{6\sqrt{5}}}$ et $(D')$ : $Y = -\frac{7+3\sqrt{5}}{\sqrt{6\sqrt{5}}}$. 

\textbf{Eléménts caractéristiques de $(\Gamma)$ dans $\mathcal{R}$.}

Les formules de changement de repère s'écrivent $\left(\begin{array}{c}
x\\
y
\end{array}
\right)=\left(
\begin{array}{cc}
\frac{2}{\sqrt{10+2\sqrt{5}}}&\frac{2}{\sqrt{10+2\sqrt{5}}}\\
\frac{1+\sqrt{5}}{\sqrt{10+2\sqrt{5}}}&\frac{1-\sqrt{5}}{\sqrt{10-2\sqrt{5}}}
\end{array}
\right)\left(\begin{array}{c}
X\\
Y
\end{array}
\right)+\left(\begin{array}{c}
-1\\
0
\end{array}
\right)$.

- \emph{Centre} $\Omega(-1,0)_{\mathcal{R}}$ et excentricité $e =\frac{1}{2}\sqrt{-45+21\sqrt{5}}= 0,69...$

- \emph{Sommets} $A\left(
-1+\sqrt{\frac{25-11\sqrt{5}}{10}},
\sqrt{\frac{5-2\sqrt{5}}{5}}
\right)_{\mathcal{R}}$  et $A'\left(
-1-\sqrt{\frac{25-11\sqrt{5}}{10}},
-\sqrt{\frac{5-2\sqrt{5}}{5}}
\right)_{\mathcal{R}}$  puis

$B\left(
-1+\sqrt{\frac{25+11\sqrt{5}}{10}},
\sqrt{\frac{5+2\sqrt{5}}{5}}
\right)_{\mathcal{R}}$  et $B'\left(
-1-\sqrt{\frac{25+11\sqrt{5}}{10}},
\sqrt{\frac{5+2\sqrt{5}}{5}}
\right)_{\mathcal{R}}$

- \emph{Foyers} $F\left(
-1+\sqrt{\frac{3(\sqrt{5}-1)}{2}},
\sqrt{\frac{3(\sqrt{5}-1)}{2}}
\right)_{\mathcal{R}}$  et $F'\left(
-1-\sqrt{\frac{3(\sqrt{5}-1)}{2}},
-\sqrt{\frac{3(\sqrt{5}-1)}{2}}
\right)_{\mathcal{R}}$ 

-\emph{Directrices} $(D)$ :  $\frac{1}{\sqrt{10-2\sqrt{5}}}(2(x+1)+(1-\sqrt{5})y) =\frac{7+3\sqrt{5}}{\sqrt{6\sqrt{5}}}$  et 
$(D')$ :  $\frac{1}{\sqrt{10-2\sqrt{5}}}(2(x+1)+(1-\sqrt{5})y) =-\frac{7+3\sqrt{5}}{\sqrt{6\sqrt{5}}}$.

$$\includegraphics{../images/img005815-1}$$



\item  Pour $(x,y)\in\Rr^2$, posons $f(x,y) =x^2+2xy+y^2+3x-2y+1$ puis $Q((x,y)) = x^2+2xy+y^2 = (x+y)^2$. $Q$ est de rang $1$ et donc $(\Gamma)$ est du genre parabole c'est-à-dire soit une parabole, soit une réunion de deux droites parallèles éventuellement confondues, soit l'ensemble vide.

Si $(\Gamma)$ est non vide, $(\Gamma)$ est une conique est de direction asymptotique d'équation $y = - x$ (fournie par $Q(x,y)= 0$).

\textbf{1ère étude.} On étudie l'intersection de $(\Gamma)$ avec une perpendiculaire quelconque à sa direction. Soit $(D_k)$ la droite d'équation $y = x+k$, $k\in\Rr$.

L'équation aux abscisses des points d'intersection de $(\Gamma)$ et $(D_k)$ est
$x^2+2x(x+k)+(x+k)^2+3x-2(x+k)+1=0$ ou encore 

\begin{center}
$4x^2+(4k+1)x+k^2-2k+1=0$.
\end{center}

Le discriminant de cette équation est $\Delta=(4k+1)^2-16(k^2-2k+1) = 40k-15$. Puisque ce discriminant change de signe, $(\Gamma)$ est une parabole.

Le discriminant est nul pour $k=\frac{3}{8}$ ce qui founit la tangente au sommet $(T)$ : $y=x+\frac{3}{8}$ et aussi le sommet :

\begin{center}
$x_S=-\frac{4\times\frac{3}{8}+1}{2\times4}= -\frac{5}{16}$  et $y_S =x_S+\frac{3}{8}=\frac{1}{16}$.
\end{center}

Le sommet de la parabole $(\Gamma)$ est le point $S\left(-\frac{5}{16},\frac{1}{16}\right)$.

L'axe focal est la perpendiculaire à la droite $(T)$ en $S$. Une équation de l'axe focal $(\Delta)$ est $y+\frac{5}{16}=-\left(x-\frac{1}{16}\right)$ ou encore $y = -x-\frac{1}{4}$.

Pour obtenir le paramètre, le foyer et la directrice, on constate tout d'abord au vu du signe du discriminant calculé plus haut que $F = S+\frac{p}{2}\left(-\frac{1}{\sqrt{2}},\frac{1}{\sqrt{2}}\right)$    et $K= S-\frac{p}{2}\left(-\frac{1}{\sqrt{2}},\frac{1}{\sqrt{2}}\right)$.

Il ne manque plus que le paramètre $p$. Soit $M$ l'un des deux points de $(\Gamma)$ situé sur la parallèle à la tangente au sommet passant par $F$. La construction usuelle d'une parabole point par point montre que le quadrilatère $(M,F,K,H)$ est un carré. Le paramètre $p$ cherché est alors $p=FK=FM$.

 
Dans ce cas, la droite $(MK)$ est la tangente à $(\Gamma)$ en $M$ et la bissectrice de l'angle des droites $(D)$ et $(\Delta)$. Cette tangente est donc parallèle à l'un des axes de coordonnées.

L'équation générale de la tangente en un point $(x_0,y_0)$ de $(\Gamma)$ est fournie par la règle de dédoublement des termes : $xx_0+xy_0+x_0y+yy_0+\frac{3}{2}(x+ x_0)-(y+y_0)+1=0$ ou encore $x\left(x_0+y_0+\frac{3}{2}\right)+y(x_0+y_0-1)+ x_0-y_0+1=0$. Cette tangente est parallèle à l'un des axes de coordonnées si et seulement si $x_0+y_0+\frac{3}{2}=0$ ou $x_0+y_0-1=0$.

$M$ est donc sur l'une des deux droites $(\Delta_1)$ : $x+y+\frac{3}{2}=0$ ou $(\Delta_2)$ : $x+y-1=0$ qui sont toutes deux parallèles à l'axe focal $(\Delta)$ : $x+y+\frac{1}{4}=0$.

$p$ est donc aussi la distance de $(\Delta)$ à l'une quelconque de ces deux droites ou la distance d'un point quelconque de $(\Delta)$ à la droite $(\Delta_1)$. Comme le point de coordonnées $\left(-\frac{1}{4},0\right)$ est sur $(\Delta)$,

\begin{center}
$p=\frac{\left|-\frac{1}{4}+0+\frac{3}{2}\right|}{\sqrt{1^2+1^2}}=\frac{5}{4\sqrt{2}}$
\end{center}

puis $F=S+\frac{5}{8\sqrt{2}}\left(-\frac{1}{\sqrt{2}},\frac{1}{\sqrt{2}}\right)=\left(-\frac{5}{16}-\frac{5}{16},\frac{1}{16}+\frac{5}{16}\right)=\left(-\frac{5}{8},\frac{3}{8}\right)$  
et $K =S-\frac{5}{8\sqrt{2}}\left(-\frac{1}{\sqrt{2}},\frac{1}{\sqrt{2}}\right)=\left(-\frac{5}{16}+\frac{5}{16},\frac{1}{16}-\frac{5}{16}\right)=\left(0,-\frac{1}{4}\right)$ de sorte que la directrice $(D)$ a pour équation $y=x-\frac{1}{4}$.

$$\includegraphics{../images/img005815-2}$$



\textbf{2ème étude.} On pose $X =\frac{1}{\sqrt{2}}(x+y)$ et $Y=\frac{1}{\sqrt{2}}(-x+y)$ ou encore $x=\frac{1}{\sqrt{2}}(X-Y)$ et $y=\frac{1}{\sqrt{2}}(X+Y)$ ce qui correspond au changement de bases orthonormées de matrice $P=\left(
\begin{array}{cc}
\frac{1}{\sqrt{2}}&-\frac{1}{\sqrt{2}}\\
\rule{0mm}{6mm}\frac{1}{\sqrt{2}}&\frac{1}{\sqrt{2}}
\end{array}
\right)$. On note $(e_1,e_2)$ la famille de matrice $P$ dans la base $(i,j)$. Déterminons une équation de $(\Gamma)$ dans le repère $\mathcal{R}'=(O,e_1,e_2)$.

\begin{align*}\ensuremath
(x+y)^2+3x-2y+1=0&\Leftrightarrow2X^2+\frac{3}{\sqrt{2}}(X-Y)-\frac{2}{\sqrt{2}}(X+Y)+1 = 0\Leftrightarrow2X^2+\frac{1}{\sqrt{2}}X-\frac{5}{\sqrt{2}}Y+1=0\\
 &\Leftrightarrow2\left(X+\frac{1}{4\sqrt{2}}\right)^2-\frac{1}{16}-\frac{5}{\sqrt{2}}Y+1=0\Leftrightarrow\left(X+\frac{1}{4\sqrt{2}}\right)^2=\frac{5}{2\sqrt{2}}\left(Y-\frac{3}{8\sqrt{2}}\right).
\end{align*}

\textbf{Eléments de $(\Gamma)$ dans $\mathcal{R}'$.}

- $(\Gamma)$ est une parabole de \emph{sommet} $S\left(-\frac{1}{4\sqrt{2}},\frac{3}{8\sqrt{2}}\right)_{\mathcal{R}'}$.

- \emph{Paramètre} $p =\frac{5}{4\sqrt{2}}$. L'axe focal de $(\Gamma)$ est l'axe $(SY)$ et le foyer a une ordonnée strictement supérieure à $Y_S$. 

- \emph{Foyer} $F = S+\frac{5}{8\sqrt{2}}(0,1)_{\mathcal{R}'}=\left(-\frac{1}{4\sqrt{2}},\frac{1}{\sqrt{2}}\right)_{\mathcal{R}'}$.

- \emph{Directrice} $K = S-\frac{5}{8\sqrt{2}}(0,1)_{\mathcal{R}'}=\left(-\frac{1}{4\sqrt{2}},-\frac{1}{4\sqrt{2}}\right)_{\mathcal{R}'}$ et donc $(D)$ : $Y = -\frac{1}{4\sqrt{2}}$.

 

\textbf{Eléments de la parabole $(\Gamma)$ dans le repère $\mathcal{R}$.}

- \emph{Paramètre} $p =\frac{5}{4\sqrt{2}}$. \emph{Sommet} $S=\left(\frac{1}{\sqrt{2}}X_S-\frac{1}{\sqrt{2}}Y_S,\frac{1}{\sqrt{2}}X_S+\frac{1}{\sqrt{2}}Y_S\right)_{\mathcal{R}}=\left(-\frac{5}{16},\frac{1}{16}\right)_{\mathcal{R}}$.

- Le \emph{foyer} $F$ a pour coordonnées  $\left(\frac{1}{\sqrt{2}}X_F-\frac{1}{\sqrt{2}}Y_F,\frac{1}{\sqrt{2}}X_F+\frac{1}{\sqrt{2}}Y_F\right)_{\mathcal{R}}=\left(-\frac{5}{8},\frac{3}{8}\right)_{\mathcal{R}}$.

- La \emph{directrice} $(D)$ a pour équation $\frac{1}{\sqrt{2}}(-x+y)=-\frac{1}{4\sqrt{2}}$ et donc $y = x-\frac{1}{4}$.

\item  Pour $(x,y)\in\Rr^2$, posons $f(x,y) =2x^2-4xy-3x+3y+1$ puis $Q((x,y)) = 2x^2-4xy=0$.

Le discriminant de cette conique est $\Delta= 2\times0 -(-2)^2 = -4 < 0$ et la courbe est du genre hyperbole c'est-à-dire soit une hyperbole, soit une réunion de deux droites sécantes.
Dans les deux cas, les deux directions asymptotiques admettent pour équation respective $x=0$ et $x=2y$ (fourni par $Q(x,y)=0$)

\textbf{Point critique.}

\begin{center}
$\left\{
\begin{array}{l}
\rule[-5mm]{0mm}{0mm}\frac{\partial f}{\partial x}(x,y)=0\\
\frac{\partial f}{\partial y}(x,y)=0
\end{array}
\right.\Leftrightarrow
\left\{
\begin{array}{l}
4x-4y-3=0\\
-4x+3=0
\end{array}
\right.\Leftrightarrow x=\frac{3}{4}\;\text{et}\;y=0$.
\end{center}

On note $\Omega$ le point de coordonnées $\left(\frac{3}{4} ,0\right)$ dans le repère $\mathcal{R}$.

\textbf{Asymptotes.} Ce sont les droites passant par $\Omega$ de directions d'équations $x=0$ et $x=2y$. Les asymptotes sont les droites $(D_1)$ : $x =\frac{3}{4}$  et $(D_2)$ : $x-\frac{3}{4}=2y$. La réunion de ces deux droites a pour équation $\left(x-\frac{3}{4}\right)\left(x-\frac{3}{4}-2y\right)=0$ ou encore $2x^2-4xy-3x+3y+\frac{9}{2}=0$. $(\Gamma)$ n'est pas $(D_1)\cup(D_2)$ et donc $(\Gamma)$ est une hyperbole.

\textbf{Axe focal et axe transverse.}

Ce sont les deux bissectrices de la paire de droites $((D_1),(D_2))$ ou encore l'ensemble des points à égale distance de ces deux droites ou encore l'ensemble des points de coordonnées $(x,y)$ dans $\mathcal{R}$ tels que $\left(x-\frac{3}{4}\right)^2=\frac{1}{5}\left(x-2y-\frac{3}{4}\right)^2$.

Ce sont donc les droites d'équations respectives $\left(x-\frac{3}{4}\right)-\frac{1}{\sqrt{5}}\left(x-2y-\frac{3}{4}\right)=0$ et  
$\left(x-\frac{3}{4}\right)+\frac{1}{\sqrt{5}}\left(x-2y-\frac{3}{4}\right)=0$ ou encore $y=-\frac{\sqrt{5}-1}{8}(4x-3)$ et $y=\frac{\sqrt{5}+1}{8}(4x-3)$. Seule l'une de ses deux droites a une intersection non vide avec $(\Gamma)$, à savoir l'axe focal et les deux points d'intersection sont les sommets de l'hyperbole.

L'équation aux abscisses des points d'intersection de $(\Gamma)$ et $(D_1)$ est 

\begin{center}
$2x^2-4x\left(-\frac{\sqrt{5}-1}{8}(4x-3)\right)-3x+3\left(-\frac{\sqrt{5}-1}{8}(4x-3)\right)+1=0$
\end{center}

ou encore $2\sqrt{5}x^2-3\sqrt{5}x+\frac{9\sqrt{5}-1}{8}=0$ ou enfin $x^2-\frac{3}{2}x+\frac{9}{16} -\frac{1}{16\sqrt{5}}=0$ dont le discriminant vaut $\frac{1}{4\sqrt{5}}>0$.

L'axe focal est donc la droite d'équation $y=-\frac{\sqrt{5}-1}{8}(4x-3)$. Les solutions de l'équation précédente fournissent les abscisses des sommets.

$$\includegraphics{../images/img005815-3}$$



\item  Pour $(x,y)\in\Rr^2$, posons $Q(x,y)= -5x^2+6\sqrt{3}xy+y^2$. Le discriminant de $(\Gamma)$ vaut $-5 -27=-32< 0$.

La conique est du genre hyperbole et de centre $O$. Comme $O\notin(\Gamma)$, $(\Gamma)$ est plus précisément une hyperbole de centre $O$.
Les asymptotes sont fournies par l'égalité $Q(x,y)=0$ et sont donc les droites d'équations $y=(-3\sqrt{3}\pm4\sqrt{2})x$.

La matrice de $Q$ dans la base canonique est $M=\left(
\begin{array}{cc}
-5&3\sqrt{3}\\
3\sqrt{3}&1
\end{array}
\right)$. Son polynôme caractéristique est $\chi_M=X^2+4X-32 = (X-4)(X+8)$. Ensuite, $M = PD{^t}P $ où $P=\frac{1}{2}\left(
\begin{array}{cc}
\sqrt{3}&1\\
-1&\sqrt{3}
\end{array}
\right)$ et $D =\text{diag}(-8,4)$.

Les formules de changement de repère s'écrivent $\left\{
\begin{array}{l}
x=\frac{1}{2}(\sqrt{3}X+Y)\\
y=\frac{1}{2}(-X+\sqrt{3}Y)
\end{array}
\right.$ ou aussi $\left\{
\begin{array}{l}
X=\frac{1}{2}(\sqrt{3}x-y)\\
Y=\frac{1}{2}(x+\sqrt{3}y)
\end{array}
\right.$.

Dans $\mathcal{R}'(O,e_1,e_2)$, $(\Gamma)$ a pour équation $-8X^2+4Y^2-4=0$ ou encore $-\frac{X^2}{(1/\sqrt{2})^2}+Y^2 = 1$. Donc $a=\frac{1}{\sqrt{2}}$, $b=1$ et $c=\sqrt{a^2+b^2}=\sqrt{\frac{3}{2}}$. 

\textbf{Eléments de l'hyperbole dans $\mathcal{R}'$ puis $\mathcal{R}$.} L'axe focal est $(O,e_2)$ c'est-à-dire la droite d'équation $X=0$ dans $\mathcal{R}'$ ou encore $y=\sqrt{3}x$ dans $\mathcal{R}$.

- Les \emph{sommets} sont les points $B$ et $B'$ de coordonées $(1,0)$ et $(-1,0)$ dans $\mathcal{R}'$ et donc de coordonnées $\left(\frac {1}{2},\frac{\sqrt{3}}{2}\right)$ et $\left(-\frac {1}{2},-\frac{\sqrt{3}}{2}\right)$ dans $\mathcal{R}$.

- \emph{Excentricité, foyers, directrices.} $c=\sqrt{\frac{3}{2}}$ puis $e=\frac{c}{b}=\sqrt{\frac{3}{2}}$. Les foyers $F$ et $F'$ ont pour coordonnées $\left(0,\sqrt{\frac{3}{2}}\right)$ et  $\left(0,-\sqrt{\frac{3}{2}}\right)$ dans $\mathcal{R}'$ et donc $\left(\frac{\sqrt{3}}{2\sqrt{2}},\frac{3}{2\sqrt{2}}\right)$ et  $\left(-\frac{\sqrt{3}}{2\sqrt{2}},-\frac{3}{2\sqrt{2}}\right)$ dans $\mathcal{R}$.

En ce qui concerne les directrices, $K=O +\frac{1}{e}\overrightarrow{OB}=\left(0,\sqrt{\frac{2}{3}}\right)_{\mathcal{R}'}$. Les directrices sont les droites d'équations respectives $Y=\sqrt{\frac{2}{3}}$ et $Y=-\sqrt{\frac{2}{3}}$ dans $\mathcal{R}'$ ou encore d'équations respectives $x+\sqrt{3}y=\frac{4}{\sqrt{6}}$ et $x+\sqrt{3}y=-\frac{4}{\sqrt{6}}$.

$$\includegraphics{../images/img005815-4}$$


\item  L'équation proposée s'écrit $(2x+3y)^2- 2x+1 = 0$. Il s'agit d'une conique du genre parabole de direction asymptotique éventuelle $2x+3y = 0$.

Posons $X =\frac{1}{\sqrt{13}}(2x+3y)$ et $Y =\frac{1}{\sqrt{13}}(3x-2y)$ ou encore $x =\frac{1}{\sqrt{13}}(2X+3Y)$ et $y =\frac{1}{\sqrt{13}}(3X-2Y)$.

Dans $\mathcal{R}'=(O,X,Y)$, $(\Gamma)$ admet pour équation cartésienne :

\begin{align*}\ensuremath
13X^2 -\frac{2}{\sqrt{13}}(2X+3Y) + 1= 0&\Leftrightarrow 13\left(X-\frac{2}{13\sqrt{13}}\right)^2 -\frac{6}{\sqrt{13}}Y+\frac{165}{169}= 0\\
 &\Leftrightarrow\left(X-\frac{2}{13\sqrt{13}}\right)^2=\frac{6}{13\sqrt{13}}\left(Y-\frac{55}{26\sqrt{13}}\right)
\end{align*}

Ceci montre que $(\Gamma)$ est une parabole, fournit le paramètre $p =\frac{6}{13\sqrt{13}}$  puis les éléments de $(\Gamma)$ dans le repère $\mathcal{R}'$ : 

$S=\left(\frac{2}{13\sqrt{13}},\frac{55}{26\sqrt{13}}\right)_{\mathcal{R}'}$  puis $F = S +\frac{p}{2}(0,1)_{\mathcal{R}'}=\left(\frac{2}{13\sqrt{13}},\frac{61}{26\sqrt{13}}\right)_{\mathcal{R}'}$ et $K = S -\frac{p}{2}(0,1)_{\mathcal{R}'}=\left(\frac{2}{13\sqrt{13}},\frac{49}{26\sqrt{13}}\right)_{\mathcal{R}'}$  et donc $(D)$ : $Y= \frac{49}{26\sqrt{13}}$.

\textbf{Eléments de $(\Gamma)$ dans le repère $\mathcal{R}$.}

 
$S\left(\frac{173}{338},-\frac{49}{169}\right)_{\mathcal{R}}$  puis $F\left(\frac{191}{338},-\frac{55}{169}\right)_{\mathcal{R}}$   et $(D)$ :  $3x-2y=\frac{49}{26}$.

$$\includegraphics{../images/img005815-5}$$

  

\item  $(\Gamma)$ est le point d'intersection des droites d'équations respectives $x-y+1 = 0$ et $x+y-1 = 0$ à savoir le point de coordonnées $(0,1)$.

\item  L'équation s'écrit $\left(x-\frac{3}{2}\right)^2 +\left(y-\frac{1}{2}\right)^2=\frac{1}{2}$. $(\Gamma)$ est le cercle de centre $\left(\frac{3}{2},\frac{1}{2}\right)$ et de rayon $\frac{1}{\sqrt{2}}$.

\item  On reconnait une équation du cercle de diamètre $[AB]$ où $A(0,2)$ et $B(1,3)$.

\item  Si on pose $X = x+2y-4$ et $Y = x-y-1$, l'équation s'écrit $XY = 3$ ce qui montre immédiatement que la courbe est une hyperbole dont les asymptotes sont les droites d'équations respectives $x+2y-4= 0$ et $x-y-1=0$ et donc de centre le point d'intersection de ces deux droites $\Omega(2,1)$. Ce changement de repère non orthonormé ne peut pas fournir davantage et si on veut les éléments métriques de l'hyperbole, il faut revenir aux méthodes de 3) ou 4).

\item  De nouveau, si on pose $X = 2x+y-1$ et $Y = x+y$ (ou même $Y=3(x+y)$) , l'équation s'écrit $X^2 = 3Y$. Le nouveau repère est quelconque mais on peut tout de même affirmer que la courbe est une parabole de direction asymptotique $2x+y=0$.
Avec cette équation, on ne lit cependant aucun des éléments métriques de celle-ci.
\end{enumerate}
\fincorrection
\correction{005816}
Les trois courbes proposées sont des coniques propres (équation polaire d'une conique propre dans un repère dont l'axe des abscisses est l'axe focal d'origine un foyer : $r =\frac{p}{1+e\cos\theta}$ où $p = ed = eFK$).

\begin{enumerate}
 \item  $e = 2$. Il s'agit une hyperbole dont l'un des foyers est l'origine.


$$\includegraphics{../images/img005816-1}$$



L'axe focal est $(Ox)$ et donc les sommets de l'hyperbole sont les points d'intersection de la courbe avec l'axe $(Ox)$. Ce sont les points $A$ et $A'$ de coordonnées cartésiennes $\left(-\frac{2}{3},0\right)$ et $(-2,0)$ obtenus pour $\theta=\pi$ et $\theta=0$ respectivement.

Le centre est le milieu du segment $[AA']$ à savoir le point $\Omega\left(-\frac{4}{3},0\right)$.

les directions asymptotiques sont fournies par : $2\cos\theta- 1 = 0\Leftrightarrow\theta\in\pm\frac{\pi}{3}+2\pi\Zz$. Les asymptotes sont les droites d'angle polaire $\pm\frac{\pi}{3}$ passant par $\Omega$. Ce sont les droites d'équations respectives $y =\sqrt{3}\left(x+\frac{4}{3}\right)$ et $y = -\sqrt{3}\left(x+\frac{4}{3}\right)$.

L'un des foyers $F$ est l'origine. L'autre est le symétrique de $F$ par rapport à $\Omega$ à savoir  le point $F'$ de coordonnées $(-3,0)$.

Les directrices sont fournies par les points $K =\Omega+\frac{1}{e}\overrightarrow{OA}=(-1,0)$  et 
$K'=\Omega-\frac{1}{e}\overrightarrow{OA}=\left(-\frac{5}{3},0\right)$. Les directrices sont les droites $(D)$ et $(D')$ d'équations respectives $x = -1$ et $x=-\frac{5}{3}$.

\item  L'équation s'écrit $r=\frac{3}{\frac{1}{2}+\cos\theta}$. Donc $e =\frac{1}{2}$ et la courbe est une ellipse.

$$\includegraphics{../images/img005816-2}$$


Les sommets du grand axe sont les points $A(2,0)$ et $A'(-6,0)$ obtenus pour $\theta=0$ et $\theta=\pi$. Le centre est le milieu $\Omega$ de $[AA']$ de coordonnées $(-2,0)$.

Le premier foyer $F$  est l'origine et le deuxième est le symétrique du point $F$ par rapport à $\Omega$ à savoir le point $F'(-4,0)$.

Les points $K$ et $K'$ sont définis par : $K = \Omega+\frac{1}{e}\overrightarrow{OA}=(6,0)$   et 
$K'= \Omega-\frac{1}{e}\overrightarrow{OA}=(-10,0)$. Les directrices sont les droites $(D)$ et $(D')$ d'équations respectives $x = 6$ et $x = -10$.

Les sommets du petit axe sont déterminés par $b=\sqrt{a^2-c^2}=\sqrt{4^2 - 2^2}=\sqrt{12}$ puis 
$B=\Omega+b\overrightarrow{j}=\left(-2,2\sqrt{3}\right)$ et $B=\Omega-b\overrightarrow{j}=\left(-2,-2\sqrt{3}\right)$

\item  L'équation $r=\frac{2}{1+\cos\left(\theta+\frac{\pi}{2}\right)}$. On reconnait une parabole dans une présentation non traditionnelle.

Le foyer est toujours l'origine et comme la direction asymptotique est obtenue pour $\theta=\frac{\pi}{2}$ (et a donc pour angle polaire $\frac{\pi}{2}$), l'axe focal est donc la droite passant par l'origine $F$ et d'angle polaire $\frac{\pi}{2}$, c'est-à-dire l'axe des ordonnées. Le sommet est l'intersection de la courbe avec l'axe $(Oy)$ obtenue pour $\theta=-\frac{\pi}{2}$. Pour $\theta=-\frac{\pi}{2}$, on obtient $r=1$ et donc $S(-1,0)$. Puis $K =s_S(F) =(-2,0)$ et la directrice $(D)$ est la droite d'équation $y = -2$. Enfin, $p=FK= 2$.

$$\includegraphics{../images/img005816-3}$$


\textbf{Remarque.} Si on n'est pas à l'aise en polaires, on peut toujours repasser en cartésien mais c'est une très grosse perte de temps :

En 1), $r(1-2\cos\theta) = 2$ s'écrit $r -2x = 2$ puis $x^2+y^2 = (2x+2)^2$.

En 2), $r(2+\cos\theta) = 6$ s'écrit $2r + x = 6$ et donc $4(x^2+y^2) =(-x+6)^2$.

En 3), $r(1-\sin\theta) = 2$ s'écrit $r - y = 2$ puis $x^2+y^2 = (y+2)^2$ et donc $y=\frac{x^2}{4}-1$.
\end{enumerate}
\fincorrection
\correction{005817}
\begin{enumerate}
 \item  Une ellipse (privée d'un point) admet une représentation paramétrique de la forme $\left\{
\begin{array}{l}
a\frac{1-t^2}{1+t^2}\\
\rule{0mm}{6mm}b\frac{2t}{1+t^2}
\end{array}
\right.$, $t\in\Rr$, dans un repère adapté. Une branche d'hyperbole admet une représentation paramétrique de la forme 
 $\left\{
\begin{array}{l}
a\frac{1+t^2}{1-t^2}\\
b\frac{2t}{1-t^2}
\end{array}
\right.$, $t\in]-1,1[$, dans un repère adapté. Une parabole admet une représentation paramétrique de la forme 
 $\left\{
\begin{array}{l}
\frac{t^2}{2p}\\
\rule{0mm}{6mm}t
\end{array}
\right.$, dans un repère adapté ...

Réciproquement, si la courbe admet une paramétrisation du type de l'énoncé, les six polynômes $P^2$, $PQ$, $Q^2$, $PR$, $QR$ et $R^2$ sont dans $\Rr_4[X]$ qui est de dimension $5$ et donc sont linéairement dépendants. On en déduit qu'il existe $(a,b,c,d,e,f)\neq(0,0,0,0,0,0)$ tel que 
$aP^2+2bPQ+cQ^2+2dPR+2eQR+fR^2 = 0$ ou encore tel que pour tout réel $t$ tel que $R(t)\neq 0$,

\begin{center}
$a\left(\frac{P(t)}{R(t)}\right)^2 +2b\frac{P(t)}{R(t)}\times\frac{Q(t)}{R(t)}+c\left(\frac{Q(t)}{R(t)}\right)^2+2d\frac{P(t)}{R(t)}+2e\frac{Q(t)}{R(t)}+f = 0$.
\end{center}

Le support de l'arc est donc contenu dans la courbe d'équation $ax^2 + 2bxy + cy^2 + 2dx + 2ey + f = 0$ où $(a,b,c,d,e,f)\neq(0,0,0,0,0,0))$.

\item  Construction de la courbe $\left\{
\begin{array}{l}
x=\frac{2t+1}{t^2+2t-1}\\
y=\frac{t^2-1}{t^2+2t-1}
\end{array}
\right.$.

$$\includegraphics{../images/img005817-1}$$


\end{enumerate}
\fincorrection
\nocorrection
\nocorrection
\correction{004912}
Soient $P,P'$ les symétriques de $F$ par rapport aux tangentes.
         Donc $F'P = F'P' = 2a$.
         \par
         Le triangle $FPP'$ est rectangle, donc $T$ est le milieu de $[P,P']$,
         et $TF = TP = TP'$.
         \par
         Donc, $TF^2 + TF'^2 = F'P^2 = 4a^2$.
         \par
         $TF^2 + TF'^2 = 2TO^2 + OF^2 + OF'^2$ donc $T$ appartient au cercle de
         centre $O$ et de rayon $\sqrt{a^2+b^2}$.
\fincorrection
\correction{004913}
\begin{enumerate}
  \item $a^2u^2+b^2v^2 -w^2=0$.
  \item $M : \begin{pmatrix}2a\cos\theta\cr2a\sin\theta\end{pmatrix}$,
             $P : \begin{pmatrix}2a\cos\alpha\cr2a\sin\alpha\end{pmatrix}$ :
             $(MP)$ est tangente à ${\cal E}'
             \Leftrightarrow \theta \equiv \alpha \pm \frac{2\pi}3\ [2\pi]$.
\end{enumerate}
 \fincorrection
\correction{004914}
\begin{enumerate}
  \item $\frac{x^2}{(1-\alpha)^2} + \frac{y^2}{\alpha^2} = d^2$.
\end{enumerate}
 \fincorrection
\correction{004915}
$M : \begin{pmatrix}X\cr Y\end{pmatrix}  \Rightarrow  T:\begin{pmatrix}a/e\cr b^2(e-X)/d\end{pmatrix}
                                \Rightarrow  \vec{FM}\cdot\vec{FT} = 0$.
\fincorrection
\nocorrection
\nocorrection
\correction{005550}
Notons $\mathcal{C}$ l'ensemble des points considérés. Pour $x$ réel, posons $P(x)=x^3+Ax^2+Bx+C$.

\begin{align*}\ensuremath
P(x)=P(y)&\Leftrightarrow (x^3-y^3)+A(x^2-y^2)+B(x-y)=0\Leftrightarrow(x-y)((x^2+xy+y^2)+A(x+y)+B)=0\\
 &\Leftrightarrow y=x\;\mbox{ou}\;x^2+xy+y^2+A(x+y)+B=0.
\end{align*}
$\mathcal{C}$ est donc la réunion de la droite d'équation $y=x$ et de la courbe $\mathcal{E}$
d'équation $x^2+xy+y^2+A(x+y)+B=0$. Pour déterminer la nature de $\mathcal{E}$, on fait un changement de repère orthonormé en posant

$$\left\{
\begin{array}{l}
x=\frac{1}{\sqrt{2}}(X-Y)\\
y=\frac{1}{\sqrt{2}}(X+Y)
\end{array}
\right.
$$
On obtient

\begin{align*}\ensuremath
x^2+xy+y^2+A(x+y)+B=0&\Leftrightarrow\frac{1}{2}((X-Y)^2+(X-Y)(X+Y)+(X+Y)^2)+\frac{A}{\sqrt{2}}X+B=0\\
 &\Leftrightarrow3X^2+Y^2+\sqrt{2}AX+2B=0\Leftrightarrow3\left(X+\frac{A\sqrt{2}}{6}\right)^2+Y^2=\frac{A^2-12B}{6}\;(*)
\end{align*}
$\mathcal{E}$ est une ellipse si et seulement si $A^2-12B>0$ (sinon $\mathcal{E}$ est un point ou est vide). Dans ce cas,
puisque $a=\frac{1}{\sqrt{3}}<1=b$,

\begin{center}
$e=\frac{c}{b}=\frac{\sqrt{b^2-a^2}}{b}=\sqrt{1-\frac{a^2}{b^2}}=\sqrt{1-\frac{1}{3}}=\sqrt{\frac{2}{3}}$.
\end{center}
\fincorrection
\correction{005820}
On choisit un repère orthonormé $\mathcal{R}$ dans lequel le point $A$ a pour coordonnées $(1,0)$ et le point $B$ a pour coordonnées $(-1,0)$. Dans le repère $\mathcal{R}$, la droite $(\mathcal{D})$ a pour équation $x = 1$. Ensuite, il existe un réel $\theta$ tel que le point $P$ ait pour coordonnées $(\cos\theta,\sin\theta)$. La tangente $(\mathcal{T})$ a pour équation $x\cos\theta + y\sin\theta = 1$. Pour $\theta\notin\pi\Zz$, le point $S$ a pour coordonnées $\left(1,\frac{1-\cos\theta}{\sin\theta}\right)$ ou encore $\left(1,\tan\left(\frac{\theta}{2}\right)\right)$.

$$\includegraphics{../images/img005820-1}$$


La perpendiculaire à la droite $(AB)$ passant $P$ admet pour équation $x=\cos\theta$. La droite $(BS)$ admet pour équation $-\tan\left(\frac{\theta}{2}\right)(x+1)+2y=0$. Ces deux droites se coupent en le point $M$ de coordonnées $\left(\cos\theta,\frac{1}{2}\tan\left(\frac{\theta}{2}\right)(1+\cos\theta)\right)$ ou encore $\left(\cos\theta,\cos\left(\frac{\theta}{2}\right)\sin\left(\frac{\theta}{2}\right)\right)$ ou enfin $(\cos\theta,\frac{1}{2}\sin\theta)$.

L'ensemble des points $M$ est donc l'ensemble des points de coordonnées $\left(\cos\theta,\frac{1}{2}\sin\theta\right)$ quand $\theta$ décrit $\Rr\setminus\pi\Zz$ ou encore l'ellipse d'équation $x^2+4y^2 = 1$ privée des points $A$ et $B$.
\fincorrection
\correction{005822}
Posons $P = X^3+\alpha X^2+\beta X+\gamma$.

\begin{align*}\ensuremath
P(x)=P(y)&\Leftrightarrow(y^3-x^3) +\alpha(y^2-x^2)+\beta(y-x) = 0\Leftrightarrow(y-x)(x^2+xy+y^2+\alpha(x+y)+\beta) = 0\\
 &\Leftrightarrow y-x = 0\;\text{ou}\;x^2+xy+y^2+\alpha(x+y)+\beta= 0.
\end{align*}

L'ensemble cherché est la réunion de la droite $(D)$ d'équation $y=x$ et de la courbe $(\Gamma)$ d'équation $x^2+xy+y^2+\alpha(x+y)+\beta=0$. Pour étudier la courbe $(\Gamma)$ qui est du genre ellipse,
posons $x =\frac{1}{\sqrt{2}}(X+Y)$ et $y =\frac{1}{\sqrt{2}}(X-Y)$ puis notons $\mathcal{R}'$ le repère $(OXY)$.

\begin{align*}\ensuremath 
x^2+xy+y^2+\alpha(x+y)+\beta=0&\Leftrightarrow\frac{1}{2}((X+Y)^2+(X+Y)(X-Y)+(X-Y)^2)+\frac{\alpha}{\sqrt{2}}(X+Y+X-Y)+\beta=0\\
 &\Leftrightarrow\frac{1}{2}(3X^2+Y^2)+\alpha\sqrt{2}X+\beta=0\Leftrightarrow3\left(X+\frac{\alpha\sqrt{2}}{3}\right)^2+Y^2=\frac{2}{3}(\alpha^2-3\beta).
\end{align*}

$(\Gamma)$ est une ellipse si et seulement si $\alpha^2-3\beta> 0$ (sinon $(\Gamma)$ est un point ou est vide). Dans ce cas, $3\left(X+\frac{\alpha\sqrt{2}}{3}\right)^2+Y^2=\frac{2}{3}(\alpha^2-3\beta)\Leftrightarrow\frac{x^2}{a^2}+\frac{y^2}{b^2}= 1$ où $a^2=\frac{2}{9}(\alpha^2-3\beta) <  \frac{2}{3}(\alpha^2-3\beta) = b^2$. Par suite,

\begin{center}
$e=\frac{c}{b}=\frac{\sqrt{b^2-a^2}}{b}=\sqrt{1-\frac{a^2}{b^2}}=\sqrt{1-\frac{1}{3}}=\sqrt{\frac{2}{3}}$.
\end{center}
\fincorrection
\nocorrection
\nocorrection
\nocorrection
\correction{004903}
\begin{enumerate}
  \item Soit $O$ ce milieu. La tangente en $M$ est parallèle à $(FH')$,
             et passe par le milieu de $[F,H]$, donc par le milieu de $[H,H']$.
  \item Calcul d'angles. $\overline{(\vec{MO},\vec{MF})} \equiv
             \overline{(\vec{FH'},\vec{FM'})}$.
\end{enumerate}
 \fincorrection
\correction{004904}
Soit $O'$ ce centre.
         Les triangles $MPQ$ et $MAB$ sont semblables, donc $O$' est l'image de
         $O$ par l'homothétie de centre $M$ qui transforme $A$ en $P$.
         \par
         Soit $(A'B')$ la symétrique de $(AB)$ par rapport à $O$.
         D'après l'homothétie,
         $$\frac {O'M}{d(O',\Delta)} = \frac {OM}{d(O,(AB))} = (cste)
           = \frac {OM-O'M}{d(O,(AB))-d(O',\Delta)}
           = \frac {OO'}{d(O',(A'B'))}.$$
         Donc $O'$ décrit une partie d'une conique de foyer $O$ et de directrice
         $(A'B')$.
\fincorrection
\correction{004905}
Repère $(O,\frac{\vec{OA}}R,\vec j)   \Rightarrow $ parabole $\rho = \frac{R}{1+\sin\theta}$.
\fincorrection
\correction{004906}
arcs de paraboles de foyer $F$ et de directrices $\Delta,\Delta'$,
         parallèles à $D$ à la distance $2a$ de $D$.
\fincorrection
\nocorrection
\correction{004908}
$A:(t^2/2p,t)$, $B:(u^2/2p,u)$ avec $t(t+u) = -2p^2$.
         $AB$ est minimal pour $t^2=2p^2$ et vaut alors $3p\sqrt3$.
\fincorrection
\correction{004909}
\begin{enumerate}
  \item Parabole~: $y^2 = 2px \Rightarrow  x=2pt^2, y=2pt$.
         Corde~: $\begin{vmatrix}2pa^2 &2pb^2 &x\cr
                            2pa   &2pb   &y\cr
                            1     &1     &1\cr\end{vmatrix}=0$.
  \item $a^2  + ab + ac + bc + 1 = 0$.
  \item $c=-\frac{a^2+ab+1}{a+b}$.\par
$(BC)$~: $(2pa+y)b^2 + (2pa^2+2p-x)b -(ax+a^2y+y) = 0$.\par
Point fixe~: $y=-2pa$, $x=2p(a^2+1)$. 

  \item Parabole translatée de~${\cal P}$ de $(2p,0)$.
\end{enumerate}
 \fincorrection
\correction{004910}
\begin{enumerate}
  \item $M = (2pt^2,2pt) \Rightarrow 2t^2+2tt_0+1=0$. Il y a deux solutions si
         $|t_0| >\sqrt2$, une seule si $|t_0|=\sqrt2$ et aucune si
         $|t_0| <\sqrt2$.
  \item $t_1+t_2 = -t_0$, $t_1^2+t_2^2=t_0^2-1$.
         Centre~: $(4pt_0^2-2p,0)$ (1/2-droite).
\end{enumerate}
 \fincorrection
\correction{004911}
\begin{enumerate}
  \item $x_C-x_A = 2p\pm\sqrt{4p^2+8pxA}$.
  \item $x_{n+1} = x_n + \sqrt{8px_n+4p^2}+2p
                      = x_n\Bigl(1+\sqrt{\frac{8p}{x_n}}+ o\Bigl(\frac1{\sqrt{x_n}}\Bigr)\Bigr)$
    donc $\sqrt{x_{n+1}} = \sqrt{x_n} + \sqrt{2p}+ o(1)$ et $x_n\sim 2pn^2$.
\end{enumerate}
 \fincorrection
\correction{005545}
Dans un certain repère orthonormé, la parabole $\mathcal{P}$ admet une 
équation cartésienne de la forme $x^2=2py$. D'après la règle de dédoublement des termes, une équation de la tangente $\mathcal{T}_{x_0}$ en un point $(x_0,y_0)=\left(x_0,\frac{x_0^2}{2p}\right)$ de $\mathcal{P}$ est

\begin{center}
$xx_0=p(y+y_0)$.
\end{center}
Les tangentes en $M_0(x_0,y_0)$ et $M_1(x_1,y_1)$ sont perpendiculaires si et seulement si $x_0x_1+p^2=0$. L'orthoptique $\mathcal{C}$ est donc l'ensemble des points d'intersection de $\mathcal{T}_{x_0}$ et $\mathcal{T}_{-p^2/x_0}$ où $x_0$ décrit $\Rr^*$.
\begin{align*}\ensuremath
M(x,y)\in\mathcal{C}&\Leftrightarrow\exists x_0\in\Rr^*/\;\left\{
\begin{array}{l}
xx_0=p\left(y+\frac{x_0^2}{2p}\right)\\
-x\frac{p^2}{x_0}=p\left(y+\frac{p^3}{2x_0^2}\right)
\end{array}
\right.\Leftrightarrow\exists t\in\Rr^*/\;\left\{
\begin{array}{l}
tx-py=\frac{t^2}{2}\\
px+ty=-\frac{p^3}{2t}
\end{array}
\right.
\\
 &\Leftrightarrow\exists t\in\Rr^*/\;\left\{
\begin{array}{l}
x=\frac{1}{t^2+p^2}\left(\frac{t^3}{2}-\frac{p^4}{2t}\right)\\
y=\frac{1}{t^2+p^2}\left(-\frac{p^3}{2}-\frac{pt^2}{2}\right)
\end{array}
\right.\Leftrightarrow\exists t\in\Rr^*/\;\left\{
\begin{array}{l}
x=\frac{t^2-p^2}{2t}\\
\rule{0mm}{6mm}y=-\frac{p}{2}
\end{array}
\right..
\end{align*}
Maintenant, $\lim_{t\rightarrow 0^+}\frac{t^2-p^2}{2t}=-\infty$ et $\lim_{t\rightarrow +\infty}\frac{t^2-p^2}{2t}=+\infty$. Comme la fonction $t\mapsto\frac{t^2-p^2}{2t}$ est continue sur $]0,+\infty[$, quand $t$ décrit $]0,+\infty[$, $x=\frac{t^2-p^2}{2t}$ décrit $\Rr$. Finalement, l'orthoptique $\mathcal{C}$ est la droite d'équation $y=-\frac{p}{2}$ ou encore

\begin{center}
\shadowbox{
l'orthoptique d'une parabole est sa directrice.
}
\end{center}

$$\includegraphics{../images/img005545-1}$$

\fincorrection
\correction{005548}
On choisit un repère orthonormé $\mathcal{R}_1=\left(O',\overrightarrow{I},\overrightarrow{J},\overrightarrow{K}\right)$ tel que le plan d'équation $x+y+z=1$ dans $\mathcal{R}$ soit le plan d'équation $Z=0$ dans $\mathcal{R}_1$.
On prend $O'=(1,0,0)$ puis $\overrightarrow{K}=\left(\frac{1}{\sqrt{3}},\frac{1}{\sqrt{3}},\frac{1}{\sqrt{3}}\right)$, $\overrightarrow{I}=\left(\frac{1}{\sqrt{2}},-\frac{1}{\sqrt{2}},0\right)$ et enfin $\overrightarrow{J}=\overrightarrow{K}\wedge\overrightarrow{I}=\left(\frac{1}{\sqrt{6}},\frac{1}{\sqrt{6}},-\frac{2}{\sqrt{6}}\right)$. Les formules de changement de repère s'écrivent

\begin{center}
$\left\{
\begin{array}{l}
x=\frac{X}{\sqrt{2}}+\frac{Y}{\sqrt{6}}+\frac{Z}{\sqrt{3}}+1\\
y=-\frac{X}{\sqrt{2}}+\frac{Y}{\sqrt{6}}+\frac{Z}{\sqrt{3}}\\
z=-\frac{2Y}{\sqrt{6}}+\frac{Z}{\sqrt{3}}
\end{array}
\right.
$
\end{center}
Ensuite, soit $M$ un point de l'espace dont les coordonnées dans $\mathcal{R}$ sont notées $(x,y,z)$ et les coordonnées dans $\mathcal{R}_1$ sont notées $(X,Y,Z)$.

\begin{align*}\ensuremath
M\in(\Gamma)&\Leftrightarrow\left\{
\begin{array}{l}
y=x^2+x+1\\
x+y+z=1
\end{array}\right.\Leftrightarrow\left\{
\begin{array}{l}
-\frac{X}{\sqrt{2}}+\frac{Y}{\sqrt{6}}+\frac{Z}{\sqrt{3}}=\left(\frac{X}{\sqrt{2}}+\frac{Y}{\sqrt{6}}+\frac{Z}{\sqrt{3}}+1\right)^2+\left(\frac{X}{\sqrt{2}}+\frac{Y}{\sqrt{6}}+\frac{Z}{\sqrt{3}}+1\right)+1\\
Z=0
\end{array}\right.
\\
 &\Leftrightarrow\left\{
\begin{array}{l}
Z=0\\
-\frac{X}{\sqrt{2}}+\frac{Y}{\sqrt{6}}=\left(\frac{X}{\sqrt{2}}+\frac{Y}{\sqrt{6}}+1\right)^2+\frac{X}{\sqrt{2}}+\frac{Y}{\sqrt{6}}+2\end{array}\right.
\\
 &\Leftrightarrow\left\{
\begin{array}{l}
Z=0\\
\left(\frac{X}{\sqrt{2}}+\frac{Y}{\sqrt{6}}\right)^2+\frac{4X}{\sqrt{2}}+\frac{2Y}{\sqrt{6}}+3=0\end{array}\right..
\end{align*}
On travaille maintenant en dimension $2$ et on note encore $\mathcal{R}_1$ le repère $\left(O',\overrightarrow{I},\overrightarrow{J}\right)$. Une équation de $(\Gamma)$ dans $\mathcal{R}_1$ est 
$\left(\frac{X}{\sqrt{2}}+\frac{Y}{\sqrt{6}}\right)^2+\frac{4X}{\sqrt{2}}+\frac{2Y}{\sqrt{6}}+3=0$ ou encore $\frac{2}{3}\left(\frac{\sqrt{3}X}{2}+\frac{Y}{2}\right)^2+\frac{4X}{\sqrt{2}}+\frac{2Y}{\sqrt{6}}+3=0$.
On pose $\left\{
\begin{array}{l}
x'=\frac{\sqrt{3}X}{2}+\frac{Y}{2}\\
\rule{0mm}{7mm}y'=-\frac{X}{2}+\frac{\sqrt{3}Y}{2}
\end{array}
\right.$ ou encore $\left\{
\begin{array}{l}
X=\frac{\sqrt{3}x'}{2}-\frac{y'}{2}\\
\rule{0mm}{7mm}Y=\frac{x'}{2}+\frac{\sqrt{3}y'}{2}
\end{array}
\right.$ et on note $\mathcal{R}'=\left(O',\overrightarrow{i'},\overrightarrow{j'}\right)$ le nouveau repère défini par ces formules.
\begin{align*}\ensuremath
M\in(\Gamma)&\Leftrightarrow\frac{2}{3}x'^2+\frac{4}{\sqrt{2}}\left(\frac{\sqrt{3}x'}{2}-\frac{y'}{2}\right)+\frac{2}{\sqrt{6}}\left(\frac{x'}{2}+\frac{\sqrt{3}y'}{2}\right)+3=0\Leftrightarrow\frac{2}{3}x'^2+\frac{7x'}{\sqrt{6}}-\frac{y'}{\sqrt{2}}+3=0\\
 &\Leftrightarrow\frac{2}{3}\left(x'+\frac{21}{4\sqrt{6}}\right)^2=\frac{y'}{\sqrt{2}}+\frac{1}{16}\Leftrightarrow\left(x'+\frac{21}{4\sqrt{6}}\right)^2=\frac{3}{2\sqrt{2}}\left(y'+\frac{1}{8\sqrt{2}}\right).
\end{align*}
$\left(\Gamma\right)$ est une parabole de paramètre $p=\frac{3}{4\sqrt{2}}$.
Eléments de $(\Gamma)$ dans $\mathcal{R}'$ : sommet $S\left(-\frac{21}{4\sqrt{6}},-\frac{1}{8\sqrt{2}}\right)_{\mathcal{R}'}$, axe : $x'=-\frac{21}{4\sqrt{6}}$, foyer $F\left(-\frac{21}{4\sqrt{6}},\frac{1}{4\sqrt{2}}\right)_{\mathcal{R}'}$, directrice : $y'=-\frac{1}{2\sqrt{2}}$.
Eléments de $(\Gamma)$ dans $\mathcal{R}_1$ en repassant à trois coordonnées : sommet $S\left(-\frac{41}{16\sqrt{2}},-\frac{45}{16\sqrt{6}},0\right)_{\mathcal{R}_1}$, axe : $\left\{
\begin{array}{l}
\sqrt{3}X+Y=-\frac{21}{2\sqrt{6}}\\
Z=0\\
\end{array}
\right.$, foyer $F\left(-\frac{11}{4\sqrt{2}},-\frac{3}{\sqrt{6}},0\right)_{\mathcal{R}_1}$, directrice : $\left\{
\begin{array}{l}
-X+\sqrt{3}Y=-\frac{1}{\sqrt{2}}\\
Z=0
\end{array}
\right.$.
Eléments de $(\Gamma)$ dans $\mathcal{R}$ : sommet $S\left(-\frac{3}{4},\frac{13}{16},\frac{15}{16}\right)_{\mathcal{R}}$, axe : $\left\{
\begin{array}{l}
8x-4y-4z+21=0\\
x+y+z=1\\
\end{array}
\right.$, foyer $F\left(-\frac{7}{8},\frac{7}{8},10\right)_{\mathcal{R}}$, directrice : $\left\{
\begin{array}{l}
2y-2z+1=0\\
x+y+z=1
\end{array}
\right.$.
\fincorrection
\correction{005552}
On cherche l'équation d'une telle parabole $\mathcal{P}$ sous la 
forme $(ax+by)^2+2cx+2dy+e=0,\;a^2+b^2=1,\;a>0$.

$$(1,0)\in\mathcal{P}\Leftrightarrow\boxed{a^2+2c+e=0}\;\mbox{et}\;(0,2)\in\mathcal{P}\Leftrightarrow \boxed{4b^2+4d+e=0}.$$
D'après la règle de dédoublement des termes, une équation cartésienne de la tangente à
$\mathcal{P}$ en $(1,0)$ est $a^2x+aby+c(x+1)+dy+e=0$ ou encore $(a^2+c)x+(ab+d)y+c+e=0$. Cette tangente est l'axe
$(Ox)$ si et seulement si \shadowbox{$a^2+c=c+e=0\;\mbox{et}\;ab+d\neq0$.}
Une équation cartésienne de la tangente à
$\mathcal{P}$ en $(0,2)$ est $2abx+2b^2y+cx+d(y+2)+e=0$ ou encore $(2ab+c)x+(2b^2+d)y+2d+e=0$. Cette tangente est l'axe
$(Oy)$ si et seulement si \shadowbox{$2b^2+d=2d+e=0$\; et \;$2ab+c\neq0$.}
En résumé, $\mathcal{P}$ est solution si et seulement si

$$\left\{
\begin{array}{l}
c=-a^2\\
d=-2b^2\\
e=a^2=4b^2\\
a^2+b^2=1\\
ab+d\neq0\\
2ab+c\neq0\\
a>0
\end{array}
\right..$$
Maintenant, $\left(a^2=4b^2,\;a^2+b^2=1\;\mbox{et}\;a>0\right)\Leftrightarrow a=\frac{2}{\sqrt{5}}\;\mbox{et}\;b=\pm\frac{1}{\sqrt{5}}$. Le
cas $b=\frac{1}{\sqrt{5}}$ fournit $d=-\frac{2}{5}$ puis $ab+d=0$ ce qui est exclu. Donc, nécessairement
$a=\frac{2}{\sqrt{5}}$ et $b=-\frac{1}{\sqrt{5}}$ puis $c=-\frac{4}{5}$, $d=-\frac{2}{5}$ et $e=\frac{4}{5}$ qui 
sont effectivement solution du système.
On obtient ainsi une et une seule courbe du second degré solution, à savoir la courbe d'équation cartésienne

\begin{center}
\shadowbox{
$(2x-y)^2-8x-4y+4=0.$
}
\end{center}
Il reste à vérifier que cette courbe est effectivement une parabole.
On pose $\left\{
\begin{array}{l}
X=\frac{1}{\sqrt{5}}(-x-2y)\\
Y=\frac{1}{\sqrt{5}}(2x-y)
\end{array}
\right.$ ou encore $\left\{
\begin{array}{l}
x=\frac{1}{\sqrt{5}}(-X+2Y)\\
y=\frac{1}{\sqrt{5}}(-2X-Y)
\end{array}
\right.$.

\begin{align*}\ensuremath
(2x-y)^2-8x-4y+4=0&\Leftrightarrow5Y^2-\frac{8}{\sqrt{5}}(-X+2Y)-\frac{4}{\sqrt{5}}(-2X-Y)+4=0\Leftrightarrow5Y^2-\frac{12}{\sqrt{5}}Y+\frac{16}{\sqrt{5}}X+4=0\\
 &\Leftrightarrow5\left(Y-\frac{6}{5\sqrt{5}}\right)^2=-\frac{16}{\sqrt{5}}\left(X+\frac{4}{5\sqrt{5}}\right).
\end{align*}
$\mathcal{C}$ est donc effectivement une parabole.

$$\includegraphics{../images/img005552-1}$$

\fincorrection
\correction{005818}
\textbf{Solution géométrique.}

$$\includegraphics{../images/img005818-1}$$


Soit $M$ un point de l'othoptique. $M$ est sur les tangentes $(T_0)$ et $(T_1)$ à la parabole en deux points distincts $M_0$ et $M_1$ et clairement distinct du sommet $S$. Soient $I_0$ et $I_1$ les points d'intersection des droites $(T_0)$ et $(T_1)$ respectivement avec la tangente $(T)$ au sommet de la parabole. On sait (construction usuelle de la parabole par points et tangentes) que les droites $(T_0)$ et $(T_1)$ sont perpendiculaires aux droites $(FI_0)$ et $(FI_1)$ respectivement. Donc le quadrilatère $FI_0MI_1$ est un rectangle (3 angles droits connus). Par suite, le milieu de $[FM]$ qui est aussi le milieu de $[I_0I_1]$ est sur la tangente au sommet $(T)$ et $M$ est l'image d'un point de la tangente $(T)$ par l'homothétie de centre $F$ et de rapport $2$. Finalement, le point $M$ est sur la directrice de la parabole.

Réciproquement, soit $M$ un point de la directrice et $I$ le milieu de $[FM]$. On reconstruit le rectangle précédent en plaçant d'abord sur la tangente au sommet $(T)$ les points $I_0$ et $I_1$ tels que $I$ soit le milieu de $[I_0I_1]$ et tels que $I_0I_1 = MI$. Les droites $(MI_0)$ et $(MI_1)$ sont effectivement des tangentes à la parabole qui sont perpendiculaires l'une à l'autre.

\begin{center}
\shadowbox{
L'orthoptique d'une parabole est sa directrice.
}
\end{center}

\textbf{Solution analytique.} On choisit un repère orthonormé dans lequel la parabole a pour équation cartésienne $y^2=2px$. Une équation de la tangente $(T_0)$ à la parabole en un point $M_0=(x_0,y_0)=\left(\frac{y_0^2}{2p},y_0\right)$ est $-px_0+yy_0 = px_0$.

Deux tangentes $(T_0)$ et $(T_1)$ sont perpendiculaires si et seulement si $p^2 + y_0y_1 = 0$ (en particulier $y_0y_1\neq0$). Le point d'intersection des tangentes $(T_0)$ et (T1) en deux points distincts est la solution du système $\left\{
\begin{array}{l}
-px+yy_0=px_0\\
-px+yy_1=px_1
\end{array}
\right.$ et a donc pour coordonnées  $\left(\frac{p(x_0y_1-y_0x_1)}{p(y_0-y_1)},\frac{p^2(x_0-x_1)}{p(y_0-y_1)}\right)=\left(-\frac{y_0y_1}{2p},\frac{y_0+y_1}{2}\right)=\left(-\frac{p}{2},\frac{1}{2}\left(y_0-\frac{p^2}{y_0}\right)\right)$.

En résumé, un point $M(x,y)$  est sur l'orthoptique si et seulement si il existe $y_0\neq0$ tel que  
$\left\{
\begin{array}{l}
x=-\frac{p}{2}\\
y=\frac{1}{2}\left(y_0-\frac{p^2}{y_0}\right)
\end{array}
\right.$ ce qui montre déjà qu'un point de l'orthoptique est sur la droite d'équation $x = -\frac{p}{2}$  c'est-à-dire sur la directrice de la parabole. Réciproquement, la fonction $y_0\mapsto\frac{1}{2}\left(y_0-\frac{p^2}{y_0}\right)$ est continue sur $]0,+\infty[$, tend vers $-\infty$ quand $y_0$ tend vers $0$ par valeurs supérieures et tend vers $+\infty$ quand $y_0$ tend vers $+\infty$. Donc   $\frac{1}{2}\left(y_0-\frac{p^2}{y_0}\right)$ décrit $\Rr$ quand $y_0$ décrit $]0,+\infty[$, ce qui montre que l'on obtient la totalité de la directrice.

\fincorrection
\correction{005819}
\begin{enumerate}
 \item  Soit $M$ un point du plan.

\textbf{1er cas.} Supposons que $M\notin(AB)\cup(AC)\cup(BC)$. 

\begin{align*}\ensuremath
P,\;Q\;\text{et}\;R\;\text{alignés}\Leftrightarrow\left(\overrightarrow{PQ},\overrightarrow{PR}\right)=0\;[\pi]
\Leftrightarrow\left(\overrightarrow{PQ},\overrightarrow{PM}\right)=\left(\overrightarrow{PR},\overrightarrow{PM}\right)\;[\pi].
\end{align*}

Maintenant, puisque les triangles $MPC$ et $MQC$ sont rectangles en $P$ et $Q$ respectivement, les points $P$ et $Q$ sont sur le cercle de diamètre $[MC]$. On en déduit que $\left(\overrightarrow{PQ},\overrightarrow{PM}\right)=\left(\overrightarrow{CQ},\overrightarrow{CM}\right)\;[\pi]$. De même, $\left(\overrightarrow{PR},\overrightarrow{PM}\right)=\left(\overrightarrow{BR},\overrightarrow{BM}\right)\;[\pi]$. Par suite,

\begin{align*}\ensuremath
P,\;Q\;\text{et}\;R\;\text{alignés}&\Leftrightarrow\left(\overrightarrow{CQ},\overrightarrow{CM}\right)=\left(\overrightarrow{BR},\overrightarrow{BM}\right)\;[\pi]\\
 &\Leftrightarrow\left(\overrightarrow{CA},\overrightarrow{CM}\right)=\left(\overrightarrow{BA},\overrightarrow{BM}\right)\;[\pi]\\
 &\Leftrightarrow M\;\text{appartient au cercle circonscrit au triangle}\;ABC\;(\text{privé des points}\;A,\;B\;\text{et}\;C).
\end{align*}

\textbf{2ème cas.} Supposons par exemple que $M\in(AB)$. Dans ce cas, $M=R$. Si de plus $M$ n'est ni $A$, ni $B$, alors $M\neq P$ et $M\neq Q$ puis les droites $(MP)$ et $(MQ)$ sont perpendiculaires aux droites $(BC)$ et $(AC)$ respectivement. Si par l'absurde, les points $P$, $Q$ et $R$ sont alignés, on a $(MP)=(MQ)$ et donc $(AB)//(AC)$. Ceci est une contradiction.

Donc, si les points $P$, $Q$ et $R$ sont alignés, $M$ est l'un des trois points $A$, $B$ ou $C$. La réciproque est immédiate.

En résumant les deux cas,

\begin{center}
\shadowbox{
$P$, $Q$ et $R$ sont alignés si et seulement si $M$ est sur le cercle circonscrit au triangle $ABC$.
}
\end{center}

$$\includegraphics{../images/img005819-1}$$


\item  \textbf{Parabole tangente aux trois côtés d'un triangle.} Commençons par rappeler une construction usuelle de la tangente en un point d'une parabole : le triangle $FMH$ est isocèle en $M$ et la tangente en $M$ à $\mathcal{P}$ est la médiatrice du segment $[FH]$. Par suite, le projeté orthogonal $P$ de $F$ sur la tangente $(T)$ est sur $5T_0)$ la tangente au sommet de la parabole $\mathcal{P}$.

$$\includegraphics{../images/img005819-2}$$



Soient $A$, $B$ et $C$ trois points non alignés. Si $\mathcal{P}$ est une parabole tangente aux droites $(BC)$, $(CA)$ et $(AB)$, les projetés orthogonaux $P$, $Q$ et $R$ de son foyer $F$ sur les droites  $(BC)$, $(CA)$ et $(AB)$ sont alignés sur la tangente au sommet de la parabole $\mathcal{P}$. D'après 1), le point $F$ est nécessairement sur le cercle circonscrit au triangle $ABC$.

Réciproquement, si $F$ est l'un des trois points $A$, $B$ ou $C$, $F$ n'est pas solution car une tangente à une parabole ne passe jamais par son foyer.

Soit donc $F$ un point du cercle circonscrit au triangle $ABC$ et distinct des points $A$, $B$ et $C$. Montrons alors qu'il existe une parabole de foyer $F$, tangente aux droites $(BC)$, $(CA)$ et $(AB)$.

On construit les projetés orthogonaux $P$, $Q$ et $R$ de $F$ sur les droites $(BC)$, $(CA)$ et $(AB)$. Ils sont alignés sur la droite de \textsc{Simson} $(T_0)$ de $F$ relativement au triangle $ABC$. La parabole de foyer $F$ et de tangente au sommet $(T_0)$ est solution du problème posé. La construction des  points de contact est fournie par le graphique de la page précédente : on construit les symétriques de $F$ par rapport aux points $P$, $Q$ et $R$ (ces symétriques sont sur la directrice) puis on remonte perpendiculairement à $(T_0)$ jusqu'à la parabole.

$$\includegraphics{../images/img005819-3}$$

\end{enumerate}
\fincorrection
\correction{005821}
$(\Gamma)$ est l'intersection d'un cylindre parabolique de direction $(Oz)$ et d'un plan non perpendiculaire à la direction de ce cylindre. On choisit un repère orthonormé $\mathcal{R}'=(\Omega,X,Y,Z)$ dans lequel le plan d'équation $x+y+z-1= 0$ dans le repère $\mathcal{R}=(O,x,y,z)$ soit le plan $(\Omega,X,Y)$  ou encore le plan d'équation $Z=0$ dans le repère $\mathcal{R}'$.

On pose donc $Z =\frac{1}{\sqrt{3}}(x+y+z-1)$ puis par exemple $X =\frac{1}{\sqrt{2}}(x-y)$ et $Y=\frac{1}{\sqrt{6}}(x+y-2z)$ ce qui s?écrit

\begin{center}
$\left(
\begin{array}{c}
X\\
Y\\
Z
\end{array}
\right)=\left(
\begin{array}{ccc}
\frac{1}{\sqrt{2}}&-\frac{1}{\sqrt{2}}&0\\
\rule[-4mm]{0mm}{10mm}\frac{1}{\sqrt{6}}&\frac{1}{\sqrt{6}}&-\frac{2}{\sqrt{6}}\\
\frac{1}{\sqrt{3}}&\frac{1}{\sqrt{3}}&\frac{1}{\sqrt{3}}
\end{array}
\right)\left(
\begin{array}{c}
x\\
y\\
z
\end{array}
\right)+\left(
\begin{array}{c}
0\\
0\\
-\frac{1}{\sqrt{3}}
\end{array}
\right)$
\end{center}

ou encore

\begin{center}
$\left(
\begin{array}{c}
x\\
y\\
z
\end{array}
\right)=\left(
\begin{array}{ccc}
\frac{1}{\sqrt{2}}&\frac{1}{\sqrt{6}}&\frac{1}{\sqrt{3}}\\
\rule[-4mm]{0mm}{10mm}-\frac{1}{\sqrt{2}}&\frac{1}{\sqrt{6}}&\frac{1}{\sqrt{3}}\\
0&-\frac{2}{\sqrt{6}}&\frac{1}{\sqrt{3}}
\end{array}
\right)\left(
\begin{array}{c}
X\\
Y\\
Z
\end{array}
\right)+\left(
\begin{array}{c}
\frac{1}{3}\\
\rule[-4mm]{0mm}{10mm}\frac{1}{3}\\
\frac{1}{3}
\end{array}
\right)$
\end{center}

Dans le repère $\mathcal{R}'$, la courbe $(\Gamma)$ admet pour système d'équations 

\begin{center}
$\left\{
\begin{array}{l}
-\frac{1}{\sqrt{2}}X+\frac{1}{\sqrt{6}}Y+\frac{1}{\sqrt{3}}Z+\frac{1}{3}=\left(\frac{1}{\sqrt{2}}X+\frac{1}{\sqrt{6}}Y+\frac{1}{\sqrt{3}}Z+\frac{1}{3}\right)^2 +\left(\frac{1}{\sqrt{2}}X+\frac{1}{\sqrt{6}}Y+\frac{1}{\sqrt{3}}Z+\frac{1}{3}\right)+ 1\\
Z=0
\end{array}
\right.$
\end{center}

ou encore

\begin{center}
$\left\{
\begin{array}{l}
-\frac{1}{\sqrt{2}}X+\frac{1}{\sqrt{6}}Y+\frac{1}{3}=\left(\frac{1}{\sqrt{2}}X+\frac{1}{\sqrt{6}}Y+\frac{1}{3}\right)^2 +\left(\frac{1}{\sqrt{2}}X+\frac{1}{\sqrt{6}}Y+\frac{1}{3}\right)+ 1\\
Z=0
\end{array}
\right.$.
\end{center}

Continuons à deux coordonnées $X$ et $Y$ dans le plan $(\Omega,X,Y)$.

\begin{align*}\ensuremath
M(X,Y)\in(\Gamma)&\Leftrightarrow-\frac{1}{\sqrt{2}}X+\frac{1}{\sqrt{6}}Y+\frac{1}{3}=\left(\frac{1}{\sqrt{2}}X+\frac{1}{\sqrt{6}}Y+\frac{1}{3}\right)^2 +\left(\frac{1}{\sqrt{2}}X+\frac{1}{\sqrt{6}}Y+\frac{1}{3}\right)+ 1\\
 &\Leftrightarrow\frac{1}{6}(\sqrt{3}X+Y)^2 + \frac{2}{3\sqrt{6}}(\sqrt{3}X+Y)+\frac{1}{9}+\sqrt{2}X+1=0\\
 &\Leftrightarrow\frac{1}{6}(\sqrt{3}X+Y)^2+\frac{4\sqrt{2}}{3}X+\frac{2}{3\sqrt{6}}Y+\frac{10}{9}=0\\
 &\Leftrightarrow(\sqrt{3}X+Y)^2+8\sqrt{2}X+\frac{2\sqrt{6}}{3}Y+\frac{20}{3}=0.
\end{align*}

On trouve déjà une conique du genre parabole. On pose maintenant $x'=\frac{1}{2}(\sqrt{3}X+Y)$ et $y'=\frac{1}{2}(-X+\sqrt{3}Y)$ correspondant aux formules de changement de repère
$\left\{
\begin{array}{l}
x'=\frac{1}{2}(\sqrt{3}X+Y)\\
\rule{0mm}{6mm}y'=\frac{1}{2}(-X+\sqrt{3}Y)
\end{array}
\right.$    ou encore  $\left\{
\begin{array}{l}
X=\frac{1}{2}(\sqrt{3}x'-y')\\
\rule{0mm}{6mm}Y=\frac{1}{2}(x'+\sqrt{3}y')
\end{array}
\right.$.

Dans le repère $(\Omega,x',y')$, une équation de la courbe $(\Gamma)$ est

\begin{align*}\ensuremath
(\sqrt{3}X+Y)^2+8\sqrt{2}X+\frac{2\sqrt{6}}{3}Y+\frac{20}{3}=0&\Leftrightarrow4x'^2+4\sqrt{2}(\sqrt{3}x'-y')+\frac{\sqrt{6}}{3}(x'+\sqrt{3}y')+\frac{20}{3}=0\\
 &\Leftrightarrow x'^2+\frac{13}{2\sqrt{6}}x'-\frac{3}{2\sqrt{2}}y'+\frac{5}{3}=0\\
 &\Leftrightarrow\left(x'+\frac{13}{4\sqrt{6}}\right)^2-\frac{3}{2\sqrt{2}}y'+\frac{5}{3}-\frac{169}{96}=0\Leftrightarrow\left(x'+\frac{13}{4\sqrt{6}}\right)^2=\frac{3}{2\sqrt{2}}\left(y'+\frac{1}{8\sqrt{2}}\right).
\end{align*}

$(\Gamma)$ est la parabole de paramètre $p=\frac{3}{4\sqrt{2}}$ et dont les éléments caractéristiques dans le repère $(\Omega,x',y')$ sont 

$S\left(-\frac{13}{4\sqrt{6}},-\frac{1}{8\sqrt{2}}\right)$, $F = S +\frac{p}{2}(0,1)=\left(-\frac{13}{4\sqrt{6}},-\frac{1}{8\sqrt{2}}\right)+\frac{3}{8\sqrt{2}}(0,1)=\left(-\frac{13}{4\sqrt{6}},\frac{1}{4\sqrt{2}}\right)$   
puis $K= S-\frac{p}{2}(0,1)=\left(-\frac{13}{4\sqrt{6}},-\frac{1}{8\sqrt{2}}\right)-\frac{3}{8\sqrt{2}}(0,1)=\left(-\frac{13}{4\sqrt{6}},-\frac{1}{2\sqrt{2}}\right)$ et  donc $(D)$ : $y'= -\frac{1}{2\sqrt{2}}$.

On repasse maintenant dans le repère $(\Omega,X,Y)$. $S$ a pour coordonnées $\frac{1}{2}\left(
\begin{array}{cc}
\sqrt{3}&-1\\
1&\sqrt{3}
\end{array}
\right)\left(
\begin{array}{c}
-\frac{13}{4\sqrt{6}}\\
\rule{0mm}{6mm}-\frac{1}{8\sqrt{2}}
\end{array}\right)=\left(
\begin{array}{c}
-\frac{25}{16\sqrt{2}}\\
\rule{0mm}{6mm}-\frac{29}{16\sqrt{6}}
\end{array}\right)$, $F$ a pour coordonnées $\frac{1}{2}\left(
\begin{array}{cc}
\sqrt{3}&-1\\
1&\sqrt{3}
\end{array}
\right)\left(
\begin{array}{c}
-\frac{13}{4\sqrt{6}}\\
\rule{0mm}{6mm}\frac{1}{4\sqrt{2}}
\end{array}\right)=\left(
\begin{array}{c}
-\frac{7}{4\sqrt{2}}\\
\rule{0mm}{6mm}-\frac{5}{4\sqrt{6}}
\end{array}\right)$  
puis $(D)$ a pour équation $-X+\sqrt{3}Y=-\frac{1}{\sqrt{2}}$.

On revient enfin au repère $(O,x,y,z)$.

Le point $S$ a pour coordonnées $\left(
\begin{array}{ccc}
\frac{1}{\sqrt{2}}&\frac{1}{\sqrt{6}}&\frac{1}{\sqrt{3}}\\
\rule[-4mm]{0mm}{10mm}-\frac{1}{\sqrt{2}}&\frac{1}{\sqrt{6}}&\frac{1}{\sqrt{3}}\\
0&-\frac{2}{\sqrt{6}}&\frac{1}{\sqrt{3}}
\end{array}
\right)\left(
\begin{array}{c}
-\frac{25}{16\sqrt{2}}\\
\rule{0mm}{6mm}-\frac{29}{16\sqrt{6}}\\
\rule{0mm}{6mm}0
\end{array}\right)+\left(
\begin{array}{c}
\frac{1}{3}\\
\rule[-4mm]{0mm}{10mm}\frac{1}{3}\\
\frac{1}{3}
\end{array}
\right)=\left(
\begin{array}{c}
-\frac{3}{4}\\
\rule[-4mm]{0mm}{10mm}\frac{13}{16}\\
\frac{15}{16}
\end{array}
\right)$ 
puis le point $F$ a pour coordonnées   $\left(
\begin{array}{ccc}
\frac{1}{\sqrt{2}}&\frac{1}{\sqrt{6}}&\frac{1}{\sqrt{3}}\\
\rule[-4mm]{0mm}{10mm}-\frac{1}{\sqrt{2}}&\frac{1}{\sqrt{6}}&\frac{1}{\sqrt{3}}\\
0&-\frac{2}{\sqrt{6}}&\frac{1}{\sqrt{3}}
\end{array}
\right)\left(
\begin{array}{c}
-\frac{7}{4\sqrt{2}}\\
\rule{0mm}{6mm}-\frac{5}{4\sqrt{6}}\\
\rule{0mm}{6mm}0
\end{array}\right)+\left(
\begin{array}{c}
\frac{1}{3}\\
\rule[-4mm]{0mm}{10mm}\frac{1}{3}\\
\frac{1}{3}
\end{array}
\right)=\left(
\begin{array}{c}
-\frac{3}{4}\\
\rule[-4mm]{0mm}{10mm}1\\
\frac{3}{4}
\end{array}
\right)$
et enfin 

$(D)$ :  $\left\{
\begin{array}{l}
-\frac{1}{\sqrt{2}}(x-y)+\frac{\sqrt{3}}{\sqrt{6}}(x+y-2z)=-\frac{1}{\sqrt{2}}\\
x+y+z-1=0
\end{array}
\right.$

\begin{center}
\shadowbox{
\begin{tabular}{c}
$(\Gamma)$ est la parabole de paramètre $p=\frac{3}{4\sqrt{2}}$, de sommet $S\left(-\frac{3}{4},\frac{13}{16},
\frac{15}{16}\right)$, de foyer $F\left(-\frac{3}{4},1,\frac{3}{4}\right)$\\
et de directrice $(D)$ : $\left\{
\begin{array}{l}
-y+z=\frac{1}{2}\\
x+y+z-1=0
\end{array}
\right.$.
\end{tabular}
}
\end{center}
\fincorrection
\correction{005824}
On cherche une équation sous la forme $(ax+by)^2 + 2cx + 2dy + e = 0$ avec $a^2+b^2 = 1$  et $a > 0$.

\textbullet~$(1,0)\in(\mathcal{P})\Leftrightarrow a^2+2c+e = 0$   et $(0,2)\in(\mathcal{P})\Leftrightarrow4b^2+4d+e=0$.

\textbullet~D'après la règle de dédoublement des termes, la tangente en $(1,0)$ à $(\mathcal{P})$ admet pour équation cartésienne $a^2x +aby+ c(x+1) +dy + e = 0$  ou encore $(a^2+c)x + (ab+d)y + c+e = 0$. Cette tangente est l'axe $(Ox)$ si et seulement si $a^2+c = 0$ et $c+e = 0$ et $ab+d\neq 0$.

\textbullet~La tangente en $(0,2)$ à $(\mathcal{P})$ admet pour équation cartésienne $2abx + 2b^2y + cx + d(y+2) + e = 0$ ou encore $(2ab+c)x + (2b^2+d)y +2d+e = 0$. Cette tangente est l'axe $(Oy)$ si et seulement si $2b^2+d = 0$ et $2d+e = 0$ et $2ab+c\neq 0$.

En résumé, $(\mathcal{P})$ est solution si et seulement si $c=-a^2$, $d = -2b^2$, $e = a^2 = 4b^2$, $a^2+b^2 =1$, $ab+d\neq0$, $2ab+c\neq 0$ et $a > 0$.

$a > 0$, $a^2 = 4b^2$ et $a^2+b^2 =1\Leftrightarrow a =\frac{2}{\sqrt{5}}$ et $b =\pm\frac{1}{\sqrt{5}}$.

Les égalités $a=\frac{2}{\sqrt{5}}$ et $b =\frac{1}{\sqrt{5}}$ fournissent $ab+d =\frac{2}{5}-\frac{2}{5}= 0$ ce qui ne convient pas.

Il reste $a =\frac{2}{\sqrt{5}}$, $b = -\frac{2}{\sqrt{5}}$, $c = -\frac{4}{5}$ , $d= -\frac{2}{5}$  et $e  =\frac{4}{5}$  qui fournit bien une solution. On trouve une et une seule parabole. 

\begin{center}
\shadowbox{
La parabole $(\mathcal{P})$ admet pour équation cartésienne $(2x-y)^2 -8x - 4y + 4 = 0$.
}
\end{center}

$$\includegraphics{../images/img005824-1}$$


\fincorrection
\correction{004918}
Hyperbole d'excentricité $\frac1{\cos\alpha}$,
         avec $\overline{(\vec{D},\vec{\Delta})} \equiv \frac\pi2 - \alpha$.
\fincorrection
\correction{004919}
\begin{enumerate}
  \item $xy= \frac{a^2}2$.
\end{enumerate}
 \fincorrection
\nocorrection
\correction{004921}
\begin{enumerate}
  \item $MH = \frac12MF  \Rightarrow $ $IMH$, $IMN$, et $INF$ sont semblables.
\end{enumerate}
 \fincorrection
\correction{004922}
Soit $\alpha \equiv \overline{(\vec{AM},\vec{AO})}$.\par
         $\frac{OM}{\sin\alpha} = \frac{AM}{\sin(2\alpha)}
           \Rightarrow  2\cos\alpha = \frac{AM}{OM}$.\par
         Al-Khâshi $ \Rightarrow  \frac{AM^2}{OM}(OM-OA) = (OM-OA)(OM+OA)$.\par
         $OM = OA  \Rightarrow  \alpha \equiv \frac\pi4$.\par
         $OM \ne OA  \Rightarrow  OM = 2d(M,\Delta)$ où $\Delta$ est la médiatrice de
         $[O,A]$, et $M$ est du côté de $O$.
\fincorrection
\correction{004923}
$M : \begin{pmatrix}1/\cos\theta\cr-\tan\theta\end{pmatrix} \Rightarrow$ hyperbole.
\fincorrection
\correction{004924}
On se ramène à une hyperbole d'équation $xy = 1$.
	     Soient $A=\Bigl(a,\frac1a\Bigr)$, $B=\Bigl(b,\frac1b\Bigr)$, $C=\Bigl(c,\frac1c\Bigr)$.
	     Alors $H=\Bigl(-\frac1{abc},-abc\Bigr) \in \cal H$.\par
	     L'équation du cercle circonscrit à $ABC$ est :
	     $x^2+y^2+\alpha x + \beta y + \gamma = 0$ et les points communs
	     au cercle et à $\cal H$ vérifient donc :
	     $x^4 + \alpha x^3 + \gamma x^2 + \beta x + 1 = 0$.
	     On conna\^\i t 3 racines : $x=a,b,c$ donc la quatrième est
	     $q=\frac1{abc}$ ce qui prouve que $Q$ et $H$ sont symétriques par
	     rapport à $O$.
\fincorrection
\correction{005549}
Soit $\mathcal{H}$ une hyperbole. Il existe un repère orthonormé dans lequel $\mathcal{H}$ admet
une équation catésienne de la forme $\frac{x^2}{a^2}-\frac{y^2}{b^2}=1$, ($a>0$, $b>0$).
Dans ce repère, les asymptotes ont pour équations $y=\frac{b}{a}x$ et $y=-\frac{b}{a}x$. Elles sont perpendiculaires si
et seulement si $\frac{b}{a}\left(-\frac{b}{a}\right)^2=-1$ ou encore si et seulement si $a=b$. L'excentricité de $\mathcal{H}$ est alors

$$e=\frac{c}{a}=\frac{\sqrt{a^2+b^2}}{a}=\sqrt{1+\frac{b^2}{a^2}}=\sqrt{2}.$$

\begin{center}
\shadowbox{
L'excentricité de l'hyperbole équilatère vaut $\sqrt{2}$.
}
\end{center}
\fincorrection
\correction{005551}
On choisit un repère orthonormé dans lequel $P$ a pour coordonnées $(a,b)$, $a>0$, $b>0$ et l'hyperbole $\mathcal{H}$ a pour équation $xy=ab$.
Le cercle $\mathcal{C}$ de centre $P$ et de rayon $PQ$ admet pour représentation paramétrique 
$\left\{\begin{array}{l}
x=a+2\sqrt{a^2+b^2}\cos t\\
y=b+2\sqrt{a^2+b^2}\sin t
\end{array}
\right.$, $t\in\Rr$. Soit donc $M\left(a+2\sqrt{a^2+b^2}\cos t,b+2\sqrt{a^2+b^2}\sin t\right)$ un point de $\mathcal{C}$.

\begin{align*}\ensuremath
M\in\mathcal{C}&\Leftrightarrow(a+2\sqrt{a^2+b^2}\cos t)(b+2\sqrt{a^2+b^2}\sin t)=ab\\
 &\Leftrightarrow2\sqrt{a^2+b^2}(b\cos t+a\sin t)+4(a^2+b^2)\cos t\sin t=0\\
  &\Leftrightarrow2(a^2+b^2)\left(\frac{b}{\sqrt{a^2+b^2}}\cos t+\frac{a}{\sqrt{a^2+b^2}}\sin t+\sin(2t)=0\right)\\
  &\Leftrightarrow\sin(2t)+\sin(t+t_0)+0\;\text{où}\;\cos(t_0)=\frac{a}{\sqrt{a^2+b^2}}\;\text{et}\;\sin(t_0)=\frac{b}{\sqrt{a^2+b^2}}\\
  &\Leftrightarrow\exists k\in\Zz/\;2t=-t-t_0+2k\pi\;\text{ou}\;\exists k\in\Zz/\;2t=\pi+t+t_0+2k\pi\\
  &\Leftrightarrow\exists k\in\Zz/\;t=-\frac{t_0}{3}+\frac{2k\pi}{3}\;\text{ou}\;\exists k\in\Zz/\;t=\pi+t_0+2k\pi.
\end{align*}
$t=\pi+t_0+2k\pi$ fournit le point de coordonnées $(-a,-b)$ c'est-à-dire le point $Q$. Sinon, on obtient trois autres points les points $M\left(-\frac{t_0}{3}\right)$, $M\left(-\frac{t_0}{3}+\frac{2\pi}{3}\right)$ et $M\left(-\frac{t_0}{3}+\frac{4\pi}{3}\right)$. On note $A$, $B$ et $C$ ces trois points. Puisque ces trois points sont sur un cercle de centre $P$ et que $\left(\overrightarrow{PA},\overrightarrow{PB}\right)=\left(\overrightarrow{PB},\overrightarrow{PC}\right)=\left(\overrightarrow{PC},\overrightarrow{PA}\right)=\frac{2\pi}{3}$, le triangle $ABC$ est équilatéral.

$$\includegraphics{../images/img005551-1}$$

\fincorrection
\correction{005823}
On peut choisir un repère orthonormé dans lequel $(\mathcal{H})$ admet pour équation équation cartésienne $xy = ab$ et $P$ a pour coordonnées $(a,b)$ où $a$ et $b$ sont deux réels strictement positifs.

$$\includegraphics{../images/img005823-1}$$



Le cercle $(\mathcal{C})$ de centre $P$ et de rayon $PQ = 2OP$ admet la paramétrisation $\left\{
\begin{array}{l}
x=a+2\sqrt{a^2+b^2}\cos\theta\\
y=b+2\sqrt{a^2+b^2}\sin\theta
\end{array}
\right.$.

Soit $M(a+2\sqrt{a^2+b^2}\cos\theta,b+2\sqrt{a^2+b^2}\sin\theta)$, $\theta\in\Rr$, un point du cercle $\mathcal{C}$.

\begin{align*}\ensuremath
M(x,y)\in(\mathcal{H})&\Leftrightarrow (a+2\sqrt{a^2+b^2}\cos\theta)(b+2\sqrt{a^2+b^2}\sin\theta) = ab\\
 &\Leftrightarrow(a^2+b^2)\cos\theta\sin\theta + 2\sqrt{a^2+b^2}(b\cos\theta+a\sin\theta) = 0\\
 &\Leftrightarrow\sin(2\theta)+\frac{b}{\sqrt{a^2+b^2}}\cos\theta+\frac{a}{\sqrt{a^2+b^2}}\sin\theta= 0\\
 &\Leftrightarrow\sin(2\theta) +\sin(\theta+\theta_0) = 0\;\text{où}\;\theta_0 =\Arccos\left(\frac{a}{\sqrt{a^2+b^2}}\right)  \\
 &\Leftrightarrow\sin\left(\frac{3\theta+\theta_0}{2}\right)\cos\left(\frac{\theta-\theta_0}{2}\right)=0\Leftrightarrow\theta=\theta_0 +\pi\;(2\pi)\;\text{ou}\;\theta=-\frac{\theta_0}{3}\;\left(\frac{2\pi}{3}\right).
\end{align*}

Les égalités $\theta=\theta_0 +\pi\;(2\pi)$ fournissent le point $Q$.

Les égalités $\theta=-\frac{\theta_0}{3}\;\left(\frac{2\pi}{3}\right)$ fournissent trois valeurs deux à deux distinctes de $\theta$ modulo $2\pi$ et donc trois points deux à deux distincts $M_1$, $M_2$ et $M_3$ de l'hyperbole tels que les trois angles au centre $P$ du triangle $M_1M_2M_3$ soient égaux à $\frac{2\pi}{3}$. Puisque $P$ est le centre du cercle circonscrit à ce triangle, ce triangle est équilatéral de centre $P$.

\fincorrection
\correction{004925}
\begin{enumerate}
  \item vp $=1,1\pm\sqrt2$.
             er : $X^2 + (1+\sqrt2\,)Y^2 + (1-\sqrt2\,)Z^2 + \frac34 = 0  \Rightarrow $
             hyperboloïde à 2 nappes.
  \item vp $=-\frac32,-\frac32,0$.
             er : $-\frac32(X^2+Y^2)+\frac13 = 0  \Rightarrow $
             cylindre de révolution.
  \item vp $= 0,0,14$. er : $14Z^2 + 4 = 0  \Rightarrow  \varnothing$.
  \item vp $=0,-1\pm\sqrt7$.
             er : $-(1+\sqrt7\,)Y^2 + (-1+\sqrt7\,)Z^2 + 3 = 0  \Rightarrow $
             cylindre hyperbolique.
  \item vp $=0,2,3$.
             er : $2Y^2 + 3Z^2 - \frac{8X}{\sqrt6} = 0  \Rightarrow $
             paraboloïde elliptique.
  \item vp $=-\frac12,-\frac12,1$.
             er : $-\frac{X^2}2 - \frac{Y^2}2 + Z^2 + 1 = 0 \Rightarrow $
             hyperboloïde de révolution à une nappe.
  \item vp $=0,\frac92,-\frac32$.
             er : $\frac{9Y^2}2 - \frac{3Z^2}2 = 0  \Rightarrow $
             deux plans sécants.
  \item vp $=0,\pm\frac1{\sqrt2}$.
             er : $Y^2 - Z^2 =  \sqrt2  \Rightarrow $
             cylindre hyperbolique.
  \item vp $=1,1,0$.
             er : $X^2 + Y^2 =  \frac15  \Rightarrow $
             cylindre de révolution.
\end{enumerate}
 \fincorrection
\nocorrection
\nocorrection
\correction{004928}
\begin{enumerate}
  \item $\frac{(xz_0-x_0z)^2}{a^2} + \frac{(yz_0-y_0z)^2}{b^2}
              = (z-z_0)^2$.
  \item $y_0 = 0$, $\frac{x_0^2}{a^2} - \frac{z_0^2}{b^2} = 1$.
\end{enumerate}
 \fincorrection
\correction{004929}
\begin{enumerate}
  \item 
  \item 
  \item Non lorsque la valeur propre médiane est nulle
             (paraboloïde hyperbolique, cylindre hyperbolique,
	     cylindre parabolique, plans).
\end{enumerate}
 \fincorrection
\correction{004930}
\begin{enumerate}
  \item $x^2+y^2 = 1 + \lambda^2 z^2$.
  \item 
\end{enumerate}
 \fincorrection
\nocorrection
\correction{004932}
\begin{enumerate}
  \item Hyperbole équilatère d'asymptotes :
             $\begin{cases} x-z-1 = \pm y\sqrt2 \cr x+z = 1.\cr\end{cases}$
  \item $x^2+y^2 = 2z^2$, cône de révolution.
\end{enumerate}
 \fincorrection
\correction{004933}
$\lambda_1\lambda_2\lambda_3 = \frac1{a^2b^2c^2}= \frac 14  \Rightarrow 
 V = \frac{8\pi}3$.
\fincorrection
\correction{004934}
$\lambda = 0$ : $S=(0,0,0)$, cône de révolution.\par
         $\lambda = \frac43$ : $S=(\frac23,\frac23,\frac23)$, cône de révolution.
\fincorrection
\nocorrection
\correction{004936}
$\frac{\vec{MP}}{a^2} = \frac{\vec{MQ}}{b^2} = \frac{\vec{MR}}{c^2}$.
\fincorrection
\nocorrection
\nocorrection
\correction{004939}
$e<1$ : ellipsoïde de révolution, $e = 1$ : cylindre parabolique,
         $e>1$ : hyperboloïde de révolution.
\fincorrection
\correction{004940}
Ellipsoïde.
\fincorrection
\correction{004941}
$x^2\cos^2\theta + y^2 - 2xz\cos\theta\sin\theta-z^2\cos^2\theta = 0
\Rightarrow $ cône elliptique.
\fincorrection
\correction{004942}
Soit $r$ le rayon de $S$ et $h$ la distance du centre $I$ de $S$ à $P$.
         On choisit un repère tel que $P = Oxy$ et $I = (0,0,h)$.
         Soit $S'$ une sphère de centre $M(x,y,z)$ :

{Pour $z > 0$ : }
         $S$ et $S'$ extérieures $\Leftrightarrow 2(h+r)z = h^2-r^2+x^2+y^2$.\par
         $S$ à l'intérieur de $S'$ $\Leftrightarrow 2(h-r)z = h^2-r^2+x^2+y^2$ si $h > r$.\par
         $S'$ à l'intérieur de $S$ $\Leftrightarrow 2(h-r)z = h^2-r^2+x^2+y^2$ si $h < r$.\par

{Pour $z < 0$ : }
         $S$ et $S'$ extérieures $\Leftrightarrow 2(h-r)z = h^2-r^2+x^2+y^2$ si $h < r$.\par
	 $S'$ à l'intérieur de $S$ $\Leftrightarrow 2(h+r)z = h^2-r^2+x^2+y^2$ si $h < r$.\par
\fincorrection
\nocorrection
\correction{005825}
\begin{enumerate}
 \item  Pour $(x,y,z)\in\Rr^3$, on pose $Q(x,y,z)=x^2+y^2+z^2-2yz$.

Pour tout $(x,y,z)\in\Rr^3$, $x^2+y^2+z^2-2yz = x^2 + (y-z)^2 = X^2 + 2Y^2$ en posant $X = x$, $Y = \frac{1}{\sqrt{2}}(y-z)$ et $Z =\frac{1}{\sqrt{2}}(y+z)$ correspondant au changement de bases orthonormées de matrice $P =\left(
\begin{array}{ccc}
1&0&0\\
\rule[-5mm]{0mm}{11mm}0&\frac{1}{\sqrt{2}}&\frac{1}{\sqrt{2}}\\
0&-\frac{1}{\sqrt{2}}&\frac{1}{\sqrt{2}}
\end{array}
\right)$.

Notons $\mathcal{R}'= (O,e_1,e_2,e_3)$ le repère orthonormé ainsi défini. La surface $(\mathcal{S})$ admet pour équation dans $\mathcal{R}?$ $X^2+2Y^2-4X+2\sqrt{2}(Y+Z) -1 = 0$ ou encore 

\begin{center}
$(X-2)^2 +2\left(Y+\frac{1}{\sqrt{2}}\right)^2=-2\sqrt{2}\left(Z-\frac{3}{\sqrt{2}}\right)$.
\end{center}

La surface $(\mathcal{S})$ est un paraboloïde elliptique de sommet $S$ de coordonnées $\left(2,\frac{1}{\sqrt{2}},\frac{3}{\sqrt{2}}\right)$ dans $\mathcal{R}'$ et donc $\left(2,2,1\right)$ dans $\mathcal{R}$.

\item  En posant $X=\frac{1}{\sqrt{2}}(x+y)$, $Y=\frac{1}{\sqrt{2}}(-x+y)$ et $Z = z$, on obtient : $2X^2 + Z^2 = 1$.

La surface $(\mathcal{S})$ est un cylindre elliptique d'axe $(OY)$ ou encore d'axe la droite d'équations $\left\{
\begin{array}{l}
y=-x\\
z=0
\end{array}
\right.$. 

\item  Pour $(x,y,z)\in\Rr^3$, posons $Q(x,y,z)=x^2+y^2+z^2-2xy+2xz$.

La matrice de $Q$ dans la base canonique $(i,j,k)$ de $\Rr^3$ est $A=\left(
\begin{array}{ccc}
1&-1&1\\
-1&1&0\\
1&0&1
\end{array}
\right)$.

\begin{align*}\ensuremath
\chi_A&=\left|
\begin{array}{ccc}
1-X&-1&1\\
-1&1-X&0\\
1&0&1-X
\end{array}
\right|= (1-X)^3+(X-1)+(X-1) = (1-X)((1-X)^2 - 2)\\
 &= (1-X)(1+\sqrt{2}-X)(1+\sqrt{2}-X).
\end{align*}

$Q$ est de rang $3$ et de signature $(2,1)$. La surface $(\mathcal{S})$ peut être un hyperboloïde à une ou deux nappes ou un cône de révolution.

$\text{Ker}(A-I_3) =\text{Vect}(e_1)$ où $e_1=\frac{1}{\sqrt{2}}(0,1,1)$.
  

$\text{Ker}(A-(1+\sqrt{2})I_3)=\text{Vect}(e_2)$ où $e_2=\frac{1}{2}(\sqrt{2},-1,1)$ et  $\text{Ker}(A-(1-\sqrt{2})I_3)=\text{Vect}(e_3)$ où $e_3=\frac{1}{2}(\sqrt{2},1,-1)$

 
La matrice de passage correspondante est la matrice $P=\left(
\begin{array}{ccc}
0&\frac{1}{\sqrt{2}}&\frac{1}{\sqrt{2}}\\
\rule[-4mm]{0mm}{10mm}\frac{1}{\sqrt{2}}&-\frac{1}{2}&\frac{1}{2}\\
\frac{1}{\sqrt{2}}&\frac{1}{2}&-\frac{1}{2}
\end{array}
\right)$.

Déterminons une équation réduite de la surface $(\mathcal{S})$ dans le repère $(O,e_1,e_2,e_3)$.

\begin{align*}\ensuremath
X^2+&(1+\sqrt{2})Y^2+(1-\sqrt{2})Z^2+\frac{3}{\sqrt{2}}(Y+Z)-\frac{1}{2}(\sqrt{2}X-Y+Z)+\frac{1}{2}(\sqrt{2}X+Y-Z)+1 = 0\\
 &\Leftrightarrow X^2+(1+\sqrt{2})\left(Y^2+\frac{3+\sqrt{2}}{2+\sqrt{2}}Y\right)+(1-\sqrt{2})\left(Z^2-\frac{3-\sqrt{2}}{2-\sqrt{2}}Z\right)+1 = 0\\
 &\Leftrightarrow X^2+(1+\sqrt{2})\left(Y+\frac{3+\sqrt{2}}{2(2+\sqrt{2})}\right)^2-\frac{(3+\sqrt{2})^2}{8(1+\sqrt{2})}+(1-\sqrt{2})\left(Z-\frac{3-\sqrt{2}}{2(2-\sqrt{2})}\right)^2-\frac{(3-\sqrt{2})^2}{8(1-\sqrt{2})}+1 = 0\\
 &\Leftrightarrow X^2+(1+\sqrt{2})\left(Y+1-\frac{\sqrt{2}}{4}\right)^2+(1-\sqrt{2})\left(Z-1-\frac{\sqrt{2}}{4}\right)^2 =\frac{11+6\sqrt{2}}{8(1+\sqrt{2})}+\frac{11-6\sqrt{2}}{8(1-\sqrt{2})}-1\\
 &\Leftrightarrow X^2+(1+\sqrt{2})\left(Y+1-\frac{\sqrt{2}}{4}\right)^2+(1-\sqrt{2})\left(Z-1-\frac{\sqrt{2}}{4}\right)^2 =-\frac{3}{4}\\
  &\Leftrightarrow-\frac{4}{3}X^2-\frac{4(1+\sqrt{2})}{3}\left(Y+1-\frac{\sqrt{2}}{4}\right)^2+\frac{4(\sqrt{2}-1)}{3}\left(Z-1-\frac{\sqrt{2}}{4}\right)^2 =1.
\end{align*}

La surface $(\mathcal{S})$ est un hyperboloïde à deux nappes de centre de coordonnées $\left(0,-1+\frac{\sqrt{2}}{4} ,1+\frac{\sqrt{2}}{4}\right)$ dans le repère $\mathcal{R}'$.

\item  On pose $X =\frac{1}{\sqrt{5}}(x-2y)$, $Y =\frac{1}{\sqrt{5}}(2x+y)$ et $Z = z$. Dans le repère $\mathcal{R}'$ ainsi défini, la surface $(\mathcal{S})$ admet pour équation $5X^2 + 5Z^2-\frac{2}{\sqrt{5}}(X+2Y) +\frac{4}{\sqrt{5}}(-2X+Y) = 0$ ou encore 
$5\left(X-\frac{1}{\sqrt{5}}\right)^2+5Z^2 = 1$.

La surface $(\mathcal{S})$ est un cylindre de révolution d'axe la droite d'équations $\left\{
\begin{array}{l}
X=\frac{1}{\sqrt{5}}\\
Z=0
\end{array}
\right.$  dans $\mathcal{R}'$ et de rayon $\frac{1}{\sqrt{5}}$.

\item  $x^2-4x-3y-2 = 0\Leftrightarrow(x-2)^2 =3(y+2)$. La surface $(\mathcal{S})$ est un cylindre parabolique de direction $(Oz)$.

\item  Pour $(x,y,z)\in\Rr^3$, posons $Q(x,y,z)= 7x^2-2y^2+4z^2+4xy+20xz+16yz$. La matrice de $Q$ dans la base canonique $(i,j,k)$ de $\Rr^3$ est $A =\left(
\begin{array}{ccc}
7&2&10\\
2&-2&8\\
10&8&4
\end{array}
\right)$.

\begin{align*}\ensuremath
\chi_A&=\left|
\begin{array}{ccc}
7-X&2&10\\
2&-2-X&8\\
10&8&4-X
\end{array}
\right|= (7-X)(X^2-2X-72) - 2(-2X-72) + 10(10X+36)\\
 &= -X^3+9X^2+162X= -X(X+9)(X-18) .
\end{align*}

Donc $Q$ est de rang $2$ et de signature $(1,1)$.

$(x,y,z)\in\text{Ker}(A)\Leftrightarrow\left\{
\begin{array}{l}
7x+2y+10z=0\\
2x-2y+4z=0\\
10x+8y+4z=0
\end{array}
\right.\Leftrightarrow\left\{
\begin{array}{l}
x-y=-4z\\
5x+4y=-2z
\end{array}
\right.\Leftrightarrow\left\{
\begin{array}{l}
x=-2z\\
y=2z
\end{array}
\right.$ et $\text{Ker}(A)=\text{Vect}(e_1)$ où $e_1=\frac{1}{3}(-2,2,1)$.
  

$(x,y,z)\in\text{Ker}(A+9I_3)\Leftrightarrow\left\{
\begin{array}{l}
16x+2y+10z=0\\
2x+7y+8z=0\\
10x+8y+13z=0
\end{array}
\right.\Leftrightarrow\left\{
\begin{array}{l}
y=-8x-5z\\
z=-2x
\end{array}
\right.\Leftrightarrow\left\{
\begin{array}{l}
y=2x\\
z=-2x
\end{array}
\right.$  et $\text{Ker}(A+9I_3)=\text{Vect}(e_2)$ où $e_2 =\frac{1}{3}(1,2,-2)$.

$\text{Ker}(A-18I_3) =\text{Vect}(e_3)$ où $e3 =-e_1\wedge e_2 =\frac{1}{3}(2,1,2)$.

La matrice de passage du changement de bases ainsi défini est $P =\frac{1}{3}\left(
\begin{array}{ccc}
-2&1&2\\
2&2&1\\
1&-2&2
\end{array}
\right)$.

Déterminons une équation réduite de la surface $(\mathcal{S})$ dans le repère $\mathcal{R}'= (O,e_1,e_2,e_3)$

\begin{align*}\ensuremath
7x^2-&2y^2+4z^2+4xy+20xz+16yz-36x+72y-108z+36=0\\
 &\Leftrightarrow-9Y^2+18Z^2 -12(-2X+Y+2Z)+24(2X+2Y+Z)-36(X-2Y+2Z)+36 = 0\\
 &\Leftrightarrow-9Y^2+18Z^2+36X+108Y-72Z+36 = 0\Leftrightarrow-Y^2+2Z^2+4X+12Y-8Z+4=0\\
 &\Leftrightarrow4(X+8)=(Y-6)^2 - 2(Z-2)^2.
\end{align*}

La surface $(\mathcal{S})$ est un paraboloïde hyperbolique. Son point selle est le point de cordonnées $(-8,6,2)$ dans le repère $\mathcal{R}'$.

\item  La surface $(\mathcal{S})$ admet pour équation cartésienne : $x^2+y^2+z^2-xy-xz-zx-x+y = 0$.

Pour $(x,y,z)\in\Rr^3$, posons $Q(x,y,z)= x^2+y^2+z^2-xy-xz-zx$. La matrice de $Q$ dans la base canonique $(i,j,k)$ de $\Rr^3$ est $A=\left(
\begin{array}{ccc}
1&-\frac{1}{2}&-\frac{1}{2}\\
\rule[-4mm]{0mm}{10mm}-\frac{1}{2}&1&-\frac{1}{2}\\
-\frac{1}{2}&-\frac{1}{2}&1
\end{array}
\right)$. $\text{Sp}(A)=\left(\frac{3}{2},\frac{3}{2},0\right)$. Un base orthonormée $(e_1,e_2,e_3)$ de vecteurs propres est la famille de matrice $P=\left(
\begin{array}{ccc}
\frac{1}{\sqrt{2}}&\frac{1}{\sqrt{6}}&\frac{1}{\sqrt{3}}\\
\rule[-4mm]{0mm}{10mm}-\frac{1}{\sqrt{2}}&\frac{1}{\sqrt{6}}&\frac{1}{\sqrt{3}}\\
0&-\frac{2}{\sqrt{6}}&\frac{1}{\sqrt{3}}
\end{array}
\right)$.

Dans le repère $\mathcal{R}'=(O,e_1,e_2,e_3)$, la surface $(\mathcal{S})$ admet pour équation cartésienne 
$\frac{3}{2}(X^2+Y^2) -\left(\frac{1}{\sqrt{2}}X+\frac{1}{\sqrt{6}}Y+\frac{1}{\sqrt{3}}Z\right)+\left(-\frac{1}{\sqrt{2}}X+\frac{1}{\sqrt{6}}Y+\frac{1}{\sqrt{3}}Z\right) = 0$ ou encore $\frac{3}{2}(X^2+Y^2) -\sqrt{2}X = 0$ ou enfin $\left(X-\frac{\sqrt{2}}{3}\right)^2 + Y^2 =\frac{2}{9}$.

La surface $(\mathcal{S})$ est un cylindre de révolution d'axe la droite d'équation $\left\{
\begin{array}{l}
X=\frac{\sqrt{2}}{3}\\
Y=0
\end{array}
\right.$ dans le repère $\mathcal{R}'$ et de rayon $\frac{\sqrt{2}}{3}$.

\item  En posant $X=y$, $Y=\frac{1}{\sqrt{2}}(y+z)$ (et $Z=\frac{1}{\sqrt{2}}(-y+z)$, $xy + yz=0\Leftrightarrow XY=\sqrt{2}$.

La surface $(\mathcal{S})$ est un cylindre hyperbolique.

\item  Pour $(x,y,z)\in\Rr^3$, posons $Q(x,y,z)=xy+yz+zx$.

La matrice de $Q$ dans la base canonique $(i,j,k)$ de $\Rr^3$ est $\frac{1}{2}\left(
\begin{array}{ccc}
0&1&1\\
1&0&1\\
1&1&0
\end{array}
\right)$. $\text{Sp}(A)=(-\frac{1}{2},-\frac{1}{2},1)$ et donc la surface $(\mathcal{S})$ est soit un hyperboloïde à une ou deux nappes, soit un cône du second degré et dans tous les cas une surface de révolution (puisque les deux valeurs propres négatives sont égales) d'axe de direction 
$\text{Ker}(A-I_3) =\text{Vect}(1,1,1)$ et passant par le point critique $\Omega(-1,1,-1)$.

Quand on se place dans le repère $(\Omega,i,j,k)$, la surface $(\mathcal{S}$) admet pour équation $XY+YZ+ZX +2 = 0$ (car $f(-1,1,-1)=2$) puis dans le repère $(\Omega,e_1,e_2,e_3)$, $-\frac{1}{2}X^2-\frac{1}{2}Y^2+Z^2 +2 = 0$ ou encore $\frac{1}{4}X^2+\frac{1}{4}Y^2-\frac{1}{2}Z^2=1$.

La surface $(\mathcal{S})$ est un hyperboloïde de révolution à une nappe.
\end{enumerate}
\fincorrection
\correction{005826}
On cherche $(a,b,c,d,e,f,g,h,i,j)\neq (0,...,0)$ tel que la surface $(\mathcal{S})$ d'équation $ax^2+by^2+cz^2+2dxy+2eyz+2fzx+2gx+2hy+2iz+j = 0$ contienne la parabole $(\mathcal{P})$ de représentation paramétrique $\left\{
\begin{array}{l}
x=\frac{t^2}{2}\\
y=t\\
z=0
\end{array}
\right.$, $t\in\Rr$, la parabole $(\mathcal{P}')$ de représentation paramétrique $\left\{
\begin{array}{l}
x=0\\
y=t\\
z=\frac{t^2}{2}
\end{array}
\right.$, $t\in\Rr$, et le point $A(2,3,2)$.

\begin{align*}
(\mathcal{P})\subset(\mathcal{S})&\Leftrightarrow \forall t\in\Rr,\;\frac{a}{4}t^4+bt^2+dt^3+gt^2+2ht+j=0\Leftrightarrow \forall t\in\Rr,\;\frac{a}{4}t^4+dt^3+(b+g)t^2+2ht+j=0\\
 &\Leftrightarrow a=d=h=j=0\;\text{et}\;g = -b.
\end{align*}

Donc $(\mathcal{P})$ est contenue dans $(\mathcal{S})$ si et seulement si $(\mathcal{S})$ a une équation de la forme $by^2+cz^2+2eyz+2fzx-2bx+2iz = 0$ avec $(b,c,e,f,i)\neq(0,0,0,0,0)$. 

\begin{align*}
(\mathcal{P'})\subset(\mathcal{S})&\Leftrightarrow \forall t\in\Rr,\;bt^2+\frac{c}{4}t^4+et^3+it^2=0\Leftrightarrow \forall t\in\Rr,\;\frac{c}{4}t^4+et^3+(b+i)t^2=0\\
 &\Leftrightarrow c=e=0\;\text{et}\;i=-b.
\end{align*}

Donc $(\mathcal{P})$ et $(\mathcal{P}')$ sont contenues dans $(\mathcal{S})$ si et seulement si $(\mathcal{S})$ a une équation de la forme $by^2+2fzx-2bx-2bz = 0$ avec $(b,f)\neq(0,0)$. 

Enfin, $A\in(\mathcal{S})\Leftrightarrow 9b+8f-4b-4b = 0\Leftrightarrow b = -8f$ et $f\neq0$. On trouve donc une et une seule quadrique à savoir la surface $(\mathcal{S})$ d'équation $-4y2+zx+8x+8z=0$.

En posant $X=\frac{1}{\sqrt{2}}(x+z)$, $Y=y$ et $Z=\frac{1}{\sqrt{2}}(x-z)$, on obtient

\begin{align*}\ensuremath
-4y2+zx+8x+8z&=-4Y^2+\frac{1}{2}(X+Z)(X-Z)+8\sqrt{2}X\\
 &=\frac{1}{2}\left(X+8\sqrt{2}\right)^2-4Y^2+\frac{1}{2}Z^2 - 64.
\end{align*}

Dans le nouveau repère ainsi défini, une équation cartésienne de la surface $(\mathcal{S})$ est $\frac{1}{128}\left(X+8\sqrt{2}\right)^2-\frac{1}{16}Y^2+\frac{1}{128}Z^2=1$ et $(\mathcal{S})$ est  un hyperboloïde à deux nappes.
\fincorrection
\correction{005827}
Soit $(\mathcal{S})$ une surface du second degré d'équation $f(x,y,z)=0$ où $f$ est symétrique en $x$, $y$ et $z$. Soient $\sigma_1$, $\sigma_2$ et $\sigma_3$ les trois fonctions symétriques élémentaires en $x$, $y$ et $z$.

Puisque $f$ est symétrique en $x$, $y$ et $z$, $f$ est un polynôme en $\sigma_1$, $\sigma_2$ et $\sigma_3$. $f$ est d'autre part un polynôme de degré $2$ en $x$, $y$ et $z$ et donc 

\begin{center}
il existe $(a,b,c,d)\in\Rr^4$ avec $(a,b)\neq(0,0)$ tel que $f = a\sigma_1^2+b\sigma_2+c\sigma_1+d$.
\end{center}

Réciproquement, si $f$ est de la forme ci-dessus, alors $f$ est symétrique en $x$, $y$ et $z$.

Puisque $\sigma_2 = xy+yz+zx =\frac{1}{2}((x+y+z)^2-(x^2+y^2+z^2))$, $(\mathcal{S})$ admet une une équation cartésienne de la forme :

\begin{center}
$\left(a+\frac{b}{2}\right)(x+y+z)^2-b(x^2+y^2+z^2)+c(x+y+z)+d = 0$ où $(a,b)\neq(0,0)$.
\end{center}

Soit $(\mathcal{D})$ la droite passant par $O$ dirigée par $\overrightarrow{n}=\overrightarrow{i}+\overrightarrow{j}+\overrightarrow{k}$ ($\overrightarrow{n}$ est vecteur normal à tout plan d'équation $x+y+z = k$, $k\in\Rr$) et soit $r$ une rotation quelconque d'axe $(\mathcal{D})$.

Si $M$ est un point de coordonnées $(x,y,z)$ et $M'=r(M)$ a pour coordonnées $(x',y',z')$ alors $x+y+z=x'+y'+z'$ car $M$ et $M'$ sont dans un plan perpendiculaire à $(\mathcal{D})$
et $x^2+y^2+z^2=x'^2+y'^2+z'^2$ car une rotation est une isométrie et car $r(O) = O$.

Finalement, pour toute rotation $r$ d'axe $(\mathcal{D})$, $M\in(\mathcal{S})\Leftrightarrow r(M)\in(\mathcal{S})$ et donc la surface $(\mathcal{S})$ est une surface de révolution d'axe $(\mathcal{D})$.
\fincorrection
\correction{005828}
Soit $A(a,b,c)$ un point quelconque de l'espace $E_3$.

Déterminons un système d'équation du cercle $(C_A)$ d'axe $(\Delta)$ d'équations $x=y=z$ passant par $A$.

Ce cercle est par exemple l'intersection du plan passant par $A$ de vecteur normal $(1,1,1)$ et de la sphère de centre $O$ et de rayon $OA$.

Un système d'équations de $(C_A)$ est $\left\{
\begin{array}{l}
x+y+z=a+b+c\\\
x^2+y^2+z^2=a^2+b^2+c^2
\end{array}
\right.$.

Déterminons alors une équation cartésienne de la surface $\mathcal{S}$. Une condition nécessaire et suffisante pour qu'un point $M(x,y,z)$ soit un point de $(\mathcal{S})$ est $(C_M)\cap(\mathcal{D})\neq\varnothing$. Donc

\begin{align*}
M\in(\mathcal{S})&\Leftrightarrow\exists(\alpha,\beta,\gamma)\in\Rr^3/\;\left\{
\begin{array}{l}
x+y+z=\alpha+\beta+\gamma\\
x^2+y^2+z^2=\alpha^2+\beta^2+\gamma^2\\
\alpha=\gamma+2\\
\beta=2\gamma+1
\end{array}
\right.\Leftrightarrow\exists(\alpha,\beta,\gamma)\in\Rr^3/\;\left\{
\begin{array}{l}
\alpha=\gamma+2\\
\beta=2\gamma+1\\
x+y+z=\gamma+2+2\gamma+1+\gamma\\
x^2+y^2+z^2=\alpha^2+\beta^2+\gamma^2
\end{array}
\right. \\
&\Leftrightarrow\exists(\alpha,\beta,\gamma)\in\Rr^3/\;\left\{
\begin{array}{l}
\alpha=\gamma+2\\
\beta=2\gamma+1\\
\gamma=\frac{1}{4}(x+y+z-3)\\
x^2+y^2+z^2=\alpha^2+\beta^2+\gamma^2
\end{array}
\right.\\
 &\Leftrightarrow x^2+y^2+z^2 =\left(\frac{1}{4}(x+y+z-3)+2\right)^2 +\left(\frac{2}{4}(x+y+z-3)+1\right)^2 +\left(\frac{1}{4}(x+y+z-3)\right)^2\\
 &\Leftrightarrow16(x^2+y^2+z^2) =(x+y+z+5)^2+4(x+y+z-1)^2+(x+y+z-3)^2 \\
 &\Leftrightarrow16(x^2+y^2+z^2) = 6(x+y+z)^2-2(x+y+z) + 38\\
 &\Leftrightarrow5(x^2+y^2+z^2) -6(xy+yz+zx)-(x+y+z)-19 = 0.
\end{align*}

\begin{center}
\shadowbox{
Une équation cartésienne de $(\mathcal{S})$ est $5(x^2+y^2+z^2) -6(xy+yz+zx)-(x+y+z)-19 = 0$.
}
\end{center}

La matrice de la forme quadratique $(x,y,z)\mapsto5(x^2+y^2+z^2) -6(xy+yz+zx)$ dans la base canonique de $\Rr^3$ est $\left(
\begin{array}{ccc}
5&-3&-3\\
-3&5&-3\\
-3&-3&5
\end{array}
\right)$. Ses valeurs propres sont $8$, valeur propre d'ordre $2$ associée au plan d'équation $x+y+z=0$ et $-1$ valeur propre d'ordre $1$ associé à la droite d'équation. Dans le repère $\left(O,\overrightarrow{e_1},\overrightarrow{e_2},\overrightarrow{e_3}\right)$ où $\overrightarrow{e_1}=\frac{1}{\sqrt{2}}(1,-1,0)$, $\overrightarrow{e_2}=\frac{1}{\sqrt{6}}(1,1,-2)$ et $\overrightarrow{e_3}=\frac{1}{\sqrt{3}}(1,1,1)$ 

\begin{center}
$M\in(\mathcal{S})\Leftrightarrow8x'^2+8y'^2-z'^2-\sqrt{3}z'-19=0\Leftrightarrow8\left(x'-\frac{\sqrt{3}}{16}\right)^28y'^2-z'^2=19+\frac{3}{32}$.
\end{center}

La surface $(\mathcal{S})$ est un hyperboloïde à une nappe.
\fincorrection
\correction{005829}
\begin{enumerate}
 \item  On note $(\mathcal{S})$ le cône de sommet $S$ et de directrice $(\mathcal{C})$.

\begin{align*}\ensuremath
M(x,y,z)\in(\mathcal{S})\setminus\{O\}&\Leftrightarrow\exists\lambda\in\Rr^*/\;O+\lambda\overrightarrow{OM}\in(\mathcal{C})\Leftrightarrow\exists\lambda\in\Rr^*,\;\exists t\in\Rr^*/\;\left\{
\begin{array}{l}
\lambda x=t\\
\lambda y=t^2\\
\lambda z=t^3
\end{array}
\right.\\
 &\Leftrightarrow \exists\lambda\in\Rr^*,\;\exists t\in\Rr^*/\;\left\{
\begin{array}{l}
t=\lambda x\\
y=\lambda x^2\\
z=\lambda^2x^3
\end{array}
\right. \Leftrightarrow x\neq0\;\text{et}\;y\neq0\;\text{et}\;z=\left(\frac{y}{x^2}\right)^2x^3 \\
 &\Leftrightarrow x\neq0\;\text{et}\;y\neq0\;\text{et}\;z = y^2x.
\end{align*}

Si on récupère le point $O$, $M(x,y,z)\in(\mathcal{S})\Leftrightarrow(x=y=0\;\text{ou}\;xy\neq0)\;\text{et}\;z=y^2x$.

On peut noter que la surface d'équation $z = y^2x$ est la réunion du cône, sommet $O$ compris, et des axes $(Ox)$ et $(Oy)$ qui ne font pas partie du cône (à l'exception du point $O$).

\item  On note $(\mathcal{S})$ le cône de sommet $S$ et de directrice $(\mathcal{C})$.

\begin{align*}\ensuremath
M(x,y,z)\in(\mathcal{S})\setminus\{S\}&\Leftrightarrow\exists\lambda\in\Rr^*/\;S+\lambda\overrightarrow{SM}\in(\mathcal{C})\Leftrightarrow\exists\lambda\in\Rr^*/\;(1+\lambda(x-1),-1+\lambda(y+1),\lambda z)\in(\mathcal{C})\\
 &\Leftrightarrow \exists\lambda\in\Rr^*/\;\left\{
\begin{array}{l}
-1+\lambda(y+1)+\lambda z=1\\
(1+\lambda(x-1))^2+(-1+\lambda(y+1))^2=\lambda z
\end{array}
\right.\\
 &\Leftrightarrow\left(1+\frac{2}{y+z+1}(x-1)\right)^2+\left(-1+\frac{2}{y+z+1}(y+1)\right)^2=\frac{2}{y+z+1}z \\
 &\Leftrightarrow (2x+y+z-1)^2 + (y-z+1)^2 = 2(y+z+1)z\;\text{et}\;y+z+1\neq0.
\end{align*}

En résumé, $M(x,y,z)$ est dans $(\mathcal{S})$ si et seulement si $M=S$ ou $M\neq S$ et $(2x+y+z-1)^2 + (y-z+1)^2 = 2(y+z+1)z\;\text{et}\;y+z+1\neq0$.

Maintenant le point $S(1,-1,0)$ est dans le plan $(P)$ d'équation $y+z+1 = 0$ et la courbe $(\mathcal{C})$ n'a aucun point dans ce plan. Donc la surface $(\mathcal{S})$ contient un et un seul point de ce plan.

Notons alors $(\mathcal{S}')$ la surface d'équation $(2x+y+z-1)^2 + (y-z+1)^2 = 2(y+z+1)z$ et vérifions que l'intersection de $(\mathcal{S}')$ et de $(P)$ est $\{S\}$. Ceci montrera que $(\mathcal{S}')=(\mathcal{S})$.

\begin{align*}\ensuremath
M(x,y,z)\in(\mathcal{S})\cap(P)&\Leftrightarrow\left\{
\begin{array}{l}
y+z+1=0\\
(2x+y+z-1)^2 + (y-z+1)^2 = 2(y+z+1)z
\end{array}
\right.\\
 &\Leftrightarrow\left\{
\begin{array}{l}
y+z+1=0\\
2x+y+z-1=0\\
y-z+2=0
\end{array}
\right.\Leftrightarrow\left\{
\begin{array}{l}
y=-1\\
z=0\\
x=1
\end{array}
\right.\Leftrightarrow M=S
\end{align*}

Finalement $(\mathcal{S}')=(\mathcal{S})$. Une équation de $(\mathcal{S})$ est donc $(2x+y+z-1)^2 + (y-z+1)^2 = 2(y+z+1)z$ ou encore $4x^2+2y^2+4xy+4xz-2yz-4x+2y-6z+2= 0$. $(\mathcal{S})$ est donc un cône du second degré.
\end{enumerate}
\fincorrection
\correction{005830}
Notons $(C)$ le cône de sommet $S$ circonscrit à la surface $(\mathcal{S})$.

\begin{enumerate}
 \item   Ici $(\mathcal{S})$ est la sphère de centre $O$ et de rayon $3$ et  le point $S$ est extérieur à cette sphère. Donc

\begin{align*}\ensuremath
M(x,y,z)\in(C)&\Leftrightarrow M=S\;\text{ou}\;M\neq S\;\text{et}\;d(O,(SM))= 3\Leftrightarrow M=S\;\text{ou}\;M\neq S\;\text{et}\;\|\overrightarrow{SO}\wedge\overrightarrow{SM}\|= 3\|\overrightarrow{SM}\|\\
 &\Leftrightarrow\|\overrightarrow{SO}\wedge\overrightarrow{SM}\|= 3\|\overrightarrow{SM}\|\Leftrightarrow
\|(0,5,0)\wedge(x,y-5,z) \| = 3\|(x,y-5,z)\|\\
 &\Leftrightarrow(5z)^2+(5x)^2 =9(x^2+(y-5)^2+z^2)\Leftrightarrow16x^2-9(y-5)^2 +16z^2= 0.
\end{align*}

\item  Soit $M_0(x_0,y_0,z_0)$ un point de $(\mathcal{S})$ (c'est-à-dire tel que $x_0^2+x_0y_0+z_0-1 = 0$). $(\mathcal{S})$ est une surface du second degré. Une équation du plan tangent à $(\mathcal{S})$ en $M_0$ est fournie par la règle de dédoublement des termes : 

\begin{center}
$xx_0+\frac{1}{2}(y_0x+x_0y)+ \frac{1}{2}(z+z_0)-1 = 0$.
\end{center}

Ce plan tangent contient le point $S(0,0,0)$ si et seulement si $z_0 = 2$ ce qui montre déjà que la courbe de contact admet pour système d'équations  $\left\{
\begin{array}{l}
x^2+xy+z-1=0\\
z=2
\end{array}
\right.$ ou encore$\left\{
\begin{array}{l}
x^2+xy+1=0\\
z=2
\end{array}
\right.$. C'est une hyperbole du plan d'équation $z = 2$.

Le cône de sommet $S$ circonscrit à $(\mathcal{S})$ est alors le cône de sommet $S$ et de directrice $(\mathcal{C})$ d'équations $\left\{
\begin{array}{l}
x^2+xy+1=0\\
z=2
\end{array}
\right.$. On trouve la surface d'équation $4x^2+4xy+z^2 = 0$. C'est un cône du second degré.
\end{enumerate}
\fincorrection
\correction{005831}
Une équation de $(\mathcal{S})$ est encore $xy+yz+zx-\lambda x-\lambda y-\lambda z+\lambda=0$.

La matrice de la forme quadratique $Q~:~(x,y,z)\mapsto xy+yz+zx$ dans la base $(i,j,k)$ est $A= \frac{1}{2}
\left(
\begin{array}{ccc}
0&1&1\\
1&0&1\\
1&1&0
\end{array}
\right)$ est les valeurs propres de cette matrice sont $- \frac{1}{2}$, $- \frac{1}{2}$ et $1$. Le rang de $Q$ est $3$ et sa signature est $(1,2)$. La surface $(\mathcal{S})$ est à priori soit un hyperboloïde, soit un cône du second degré. Donc $(\mathcal{S})$ est un cône du second degré si et seulement si son (unique) centre de symétrie qui est aussi l'unique point critique de la fonction $f~:~(x,y,z)\mapsto x(\lambda-y)+y(\lambda-z)+z(\lambda-x)-\lambda$ appartient à $(\mathcal{S})$.

\textbf{Point critique.}

\begin{center}
$\left\{
\begin{array}{l}
 \frac{\partial f}{\partial x}(x,y,z)=0\\
 \frac{\partial f}{\partial y}(x,y,z)=0\\
 \frac{\partial f}{\partial z}(x,y,z)=0
\end{array}
\right.
\Leftrightarrow \left\{
\begin{array}{l}
y+z=\lambda\\
z+x=\lambda\\
x+y=\lambda
\end{array}
\right.\Leftrightarrow x=y=z= \frac{\lambda}{2}$.
\end{center} 

On note alors $\Omega$ le point de coordonnées $\left( \frac{\lambda}{2}, \frac{\lambda}{2}, \frac{\lambda}{2}\right)$.

$(\mathcal{S})$ est un cône $\Leftrightarrow\Omega\in(\mathcal{S})\Leftrightarrow \frac{3\lambda^2}{4}-\lambda =0\Leftrightarrow\lambda\in\left\{0, \frac{4}{3}\right\}$.

\textbullet~Si $\lambda=0$, $(\mathcal{S})$ admet pour équation $xy+yz+zx = 0$.
Dans le repère $(O,X,Y,Z)$ où $X=/dfrac{1}{\sqrt{2}}(x-y)$, $Y = \frac{1}{\sqrt{6}}(x+y-2z)$ et $Z = \frac{1}{\sqrt{3}}(x+y+z)$, $(\mathcal{S})$ admet pour équation cartésienne $- \frac{1}{2}X^2-  \frac{1}{2}Y^2+ \frac{1}{2}Z^2 = 0$ ou encore $(\mathcal{S})$ est le cône de révolution de sommet 
$O$ et de section droite le cercle d'équations $\left\{
\begin{array}{l}
Z=1\\
X^2+Y^2+Z^2=3
\end{array}
\right.$ dans $(O,X,Y,Z)$ ou encore $\left\{
\begin{array}{l}
x+y+z=\sqrt{3}\\
x^2+y^2+z^2=3
\end{array}
\right.$ dans $(O,x,y,z)$.

Puisque $(\mathcal{S})$ est un cône de révolution de sommet $O$ et d'axe la droite d'équations $x=y=z$, il est plus interessant de fournir le demi angle au sommet $\theta$. Le point $A(1,1,1)$ est sur l'axe et le point $M(2,2-1)$ est sur le cône. Donc $\theta=\Arccos\left( \frac{\left|\overrightarrow{OA}.\overrightarrow{OM}\right|}{OA\times OM}\right)=\Arccos\left( \frac{3}{3\sqrt{3}}\right)=\Arccos\left( \frac{1}{\sqrt{3}}\right)$. 

\textbullet~Si $\lambda= \frac{4}{3}$, $(\mathcal{S})$ admet pour équation $xy+yz+zx- \frac{4}{3}(x+y+z) + \frac{4}{3}= 0$ dans $(O,i,j,k)$  ou encore 
$XY+XZ+YZ = 0$ dans $(\Omega,i,j,k)$ ce qui ramène au cas précédent.
\fincorrection
\correction{005832}
Pour tout réel $t$, $(x(t))^2 +(y(t))^2= \frac{1}{4}e^{2t}((\cos t-\sin t)^2+(\cos t+\sin t)^2)= \frac{1}{2}e^{2t} = \frac{1}{2}(z(t))^2$ et le support de l'arc considéré est contenu dans le cône de révolution d'équation $z^2=2(x^2+y^2)$.
\fincorrection
\correction{005833}
\begin{enumerate}
 \item 

\begin{align*}\ensuremath
M(x,y,z)\in(\mathcal{C})&\Leftrightarrow\exists\lambda\in\Rr,\;\exists m\in(C)/\;M = m +\lambda\overrightarrow{u}\Leftrightarrow\exists\lambda\in\Rr,\;\exists t\in\Rr/\;\left\{
\begin{array}{l}
x=a\cos t+\lambda\\
y=b\sin t\\
z=a\cos t\sin t+\lambda
\end{array}
\right.\\
 &\Leftrightarrow\exists\lambda\in\Rr,\;\exists t\in\Rr/\;\left\{
\begin{array}{l}
\lambda=x-a\cos t\\
y=b\sin t\\
z=a\cos t\sin t+x-a\cos t
\end{array}
\right.\Leftrightarrow\exists t\in\Rr/\;\left\{
\begin{array}{l}
y=b\sin t\\
z=a\cos t\sin t+x-a\cos t
\end{array}
\right.\\
 &\Leftrightarrow\exists t\in\Rr/\;\left\{
\begin{array}{l}
y=b\sin t\\
z-x=a\cos t(\sin t-1)
\end{array}
\right.\Leftrightarrow\exists t\in\Rr/\;\left\{
\begin{array}{l}
y=b\sin t\\
b(z-x)=a\cos t(y-b)
\end{array}
\right.\\
 &\Leftrightarrow b^4(z-x)^2+y^2a^2(y-b)^2=a^2b^2(y-b)^2.
\end{align*}

En effet,

\textbullet~$\Rightarrow/$ s'il existe $t\in\Rr$ tel que $y=b\sin t$ et $b(z-x)=a\cos t(y-b)$ alors

\begin{align*}\ensuremath
b^4(z-x)^2+y^2a^2(y-b)^2&=b^2a^2\cos^2t(y-b)^2+b^2\sin^2ta^2(y-b)^2=a^2b^2(y-b)^2(\cos^2t+\sin^2t)\\
 &=a^2b^2(y-b)^2.
\end{align*}

\textbullet~$\Leftarrow/$ Réciproquement, si $b^4(z-x)^2+y^2a^2(y-b)^2=a^2b^2(y-b)^2$ alors $
b^4(z-x)^2=a^2(y-b)^2 (b^2-y^2)$ et donc

ou bien $y=b$, ou bien $b^2-y^2\geqslant 0$. Par suite, il existe un réel $t$ tel que $y = b\sin t=b\sin(\pi-t)$ puis

\begin{align*}\ensuremath
b^4(z-x)^2=a^2(y-b)^2 (b^2-y^2)&\Rightarrow b^4(z-x)^2=a^2(b\sin t-b)^2b^2\cos^2t\Rightarrow b(z-x) =\pm a\cos t(b\sin t-b)\\
 &\Rightarrow b(z-x) = a\cos t(y-b)\;\text{ou}\;b(z-x) = a\cos(\pi-t)(y-b)
\end{align*}

et il existe un réel $t'$ tel que $y=b\sin t'$ et $b(z-x) = a\cos t'(y-b)$.

\item  

\begin{align*}\ensuremath
M(x,y,z)\in(\mathcal{C})&\Leftrightarrow\exists\lambda\in\Rr,\;\exists m\in(C)/\;M = m +\lambda\overrightarrow{u}\Leftrightarrow\exists\lambda\in\Rr,\;\exists (X,Y,Z)\in\Rr^3/\;\left\{
\begin{array}{l}
x=X\\
y=Y+\lambda\\
z=Z+\lambda\\
Y+Z=1\\
X^2+Y^2=1
\end{array}
\right.\\
 &\Leftrightarrow\exists\lambda\in\Rr,\;\left\{
\begin{array}{l}
(y-\lambda)+(z-\lambda)=1\\
x^2+(y-\lambda)^2=1
\end{array}
\right.\Leftrightarrow x^2 +\left(y- \frac{1}{2}(y+z-1)\right)^2 = 1\\
 &\Leftrightarrow4x^2+(y-z+1)^2 = 4.
\end{align*}
\end{enumerate}
\fincorrection
\correction{005834}
La direction du cylindre est orthogonale au plan d'équation $z=x$ et est donc engendrée par le vecteur $\overrightarrow{u}(1,0,-1)$.

\begin{align*}\ensuremath
M(x,y,z)\in(\mathcal{C})&\Leftrightarrow\exists\lambda\in\Rr,\;\exists m\in(C)/\;M =m+\lambda\overrightarrow{u}\Leftrightarrow\exists\lambda\in\Rr,\;\exists(X,Y,Z)\in\Rr^3/\;\left\{
\begin{array}{l}
x=X-\lambda\\
y=Y\\
z=Z+\lambda\\
Z=X\\
2X^2+Y^2=1
\end{array}
\right. \\
 &\Leftrightarrow \exists\lambda\in\Rr/\;\left\{
\begin{array}{l}
z-\lambda=x+\lambda\\
2(x+\lambda)^2+y^2=1
\end{array}
\right.  \\
 &\Leftrightarrow 2\left(x+ \frac{1}{2}(z-x)\right)^2+y^2 =1\Leftrightarrow (x+z)^2+2y^2 = 2.
\end{align*}
\fincorrection
\correction{005835}
Un repère de $(\mathcal{D})$ est $\left(A,\overrightarrow{u}\right)$ où $A(2,1,0)$ et $\overrightarrow{u}(1,1,1)$.

\begin{align*}\ensuremath
M\in(\mathcal{C})&\Leftrightarrow d(M,(\mathcal{D})) =R\Leftrightarrow\|\overrightarrow{AM}\wedge\overrightarrow{u}\|^2 = R^2\|\overrightarrow{u}\|^2\Leftrightarrow\|(x-2,y-1,z)\wedge(1,1,1)\|^2 = R^2\|(1,1,1)\|^2\\
 &\Leftrightarrow (y-z-1)^2 + (x-z-2)^2 + (x-y-1)^2 = 3R^2 \\
 &\Leftrightarrow3x^2+3y^2+3z^2-2xy-2xz-2yz-6x+6z+6-3R^2 = 0.
\end{align*}

La droite $(Oz)$ est tangente à $(\mathcal{C})$ si et seulement si $d((Oz),(\mathcal{D})) = R$.

\begin{center}
$(Oz)$ est tangente à $(\mathcal{C})\Leftrightarrow \frac{\left[\overrightarrow{OA},\overrightarrow{k},\overrightarrow{u}\right]^2}{\left\|\overrightarrow{k}\wedge\overrightarrow{u}\right\|^2}= R^2\Leftrightarrow\left|\begin{array}{ccc}
2&0&1\\
1&0&1\\
0&1&1
\end{array}
\right|^2=\|(-1,1,0\|^2\Leftrightarrow 1 = 2R^2\Leftrightarrow R = \frac{1}{\sqrt{2}}$.
\end{center}
\fincorrection
\correction{005836}
En un point $M_0(x_0,y_0,z_0)$ de l'ellipsoïde la règle de dédoublement des termes fournit une équation du plan tangent : $xx_0+2yy_0+3zz_0 = 21$.

Ce plan est parallèle au plan d'équation $x+4y+6z=0$ si et seulement si le vecteur $(x_0,2y_0,3z_0)$ est colinéaire au vecteur $(1,4,6)$ ou encore si et seulement si $2x_0 = y_0 = z_0$.

Enfin le point $(x_0,2x_0,2x_0)$  est sur l'ellipsoïde si et seulement si $x_0^2+8x_0^2+12x_0^2=21$ ce qui équivaut à $x_0^2=1$.

Les plans cherchés sont les deux plans d'équations respectives $x+4y+6z=21$ et $x+4y+6z=-21$.
\fincorrection
\correction{005837}
Le plan tangent $(P_0)$ en $(x_0,y_0,z_0)$ tel que $x_0-8y_0z_0 = 0$ admet pour équation $(x+x_0) -8(z_0y+y_0z)=0$ ou encore $x-8z_0y-8y_0z+8y_0z_0=0$.

Un repère de $(\mathcal{D})$ est $\left(A,\overrightarrow{u}\right)$ où $A(-2,1,0)$ et $\overrightarrow{u}(4,0,-1)$.

\begin{align*}\ensuremath
(\mathcal{D})\subset(P_0)&\Leftrightarrow \forall\lambda\in\Rr,\; (-2+4\lambda)-8z_0+8y_0\lambda+8y_0z_0=0\Leftrightarrow \forall\lambda\in\Rr,\;(8y_0+4)\lambda+8y_0z_0-8z_0-2 = 0\\
 &\Leftrightarrow 8y_0+4=0\;\text{et}\;8y_0z_0-8z_0-2 = 0\Leftrightarrow y_0=- \frac{1}{2}\;\text{et}\;z_0=- \frac{1}{6}.
\end{align*}

On trouve un et un seul plan tangent contenant la droite $(\mathcal{D})$, à savoir le plan tangent à $(\mathcal{S})$ en $\left( \frac{2}{3},- \frac{1}{2},- \frac{1}{6}\right)$ d'équation $3x+4y+12z+2=0$.
\fincorrection
\correction{005838}
\begin{enumerate}
 \item  Un repère de $(\mathcal{D})$ est $\left(A,\overrightarrow{u}\right)$ où $A(0,-1,2)$ et $\overrightarrow{u}(3,3,1)$.

\begin{align*}\ensuremath
M(x,y,z)\in(\mathcal{C})&\Leftrightarrow d(M,(\mathcal{D}))=3\Leftrightarrow \left\|\overrightarrow{AM}\wedge\overrightarrow{u}\right\|^2=9\|\overrightarrow{u}\|^2\\
 &\Leftrightarrow\|(x,y+1,z-2)\wedge(3,3,1)\|^2=9\times19\Leftrightarrow(y-3z+7)^2+(x-3z+6)^2+9(x-y-1)^2=171.
\end{align*}

\item  Un repère de $(\mathcal{D})$ est $\left(A,\overrightarrow{u}\right)$ où $A(0,-1,2)$ et $\overrightarrow{u}(3,3,1)$. De plus, $S=A$.

\begin{align*}\ensuremath
M(x,y,z)\in(\mathcal{C})&\Leftrightarrow M=A\;\text{ou}\;M\neq A\;\text{et}\; \frac{\left|\overrightarrow{AM}.\overrightarrow{u}\right|}{AM\times\|\overrightarrow{u}\|}=\cos\left( \frac{\pi}{3}\right)\Leftrightarrow \left(\overrightarrow{AM}.\overrightarrow{u}\right)^2= \frac{1}{4}AM^2\|\overrightarrow{u}\|^2\\
 &\Leftrightarrow4(3x+3(y+1)+(z-2))^2=19(x^2+(y+1)^2+(z-2)^2)\\
 &\Leftrightarrow4(3x+3y+z+1)^2-19(x^2+(y+1)^2+(z-2)^2)=0\\
 &\Leftrightarrow17x^2+17y^2-15z^2+72xy+24xz+24yz+24x-14y+84z-91=0.
\end{align*}
\end{enumerate}
\fincorrection
\nocorrection
\nocorrection
\correction{004900}
\begin{enumerate}
  \item Parabole.
             axes = $\frac 15\begin{pmatrix}4\cr-3\end{pmatrix}$, $\frac 15\begin{pmatrix}3\cr4\end{pmatrix}$.
             sommet : $X = -\frac 45$, $Y = -\frac{16}5$.
  \item Ellipse.
             $\Omega = \begin{pmatrix}1\cr1\end{pmatrix}$,
             axes à $-\frac\pi3,\frac\pi6$,
             $a=2$, $b=\sqrt2$.
  \item Ellipse.
             $\Omega = \begin{pmatrix}0\cr0\end{pmatrix}$,
             axes à $-\frac\pi4,\frac\pi4$,
             $a=\sqrt2$, $b=\sqrt{\frac 23}$.
  \item 
  \item Centre $\begin{pmatrix}-2\cr 0\end{pmatrix}$.\par
             $$\left\{\begin{array}{cccccc} %????????
                &  &m &< 0  & \Rightarrow  &\text{ellipse horizontale.}   \hfill\cr
             0  &< &m &<1/2 & \Rightarrow  &\text{hyperbole verticale.}   \hfill\cr
             1/2&< &m &<1   & \Rightarrow  &\text{hyperbole horizontale.} \hfill\cr
               1&< &m &     & \Rightarrow  &\text{ellipse verticale.}     \hfill\cr  \end{array}\right.$$
\end{enumerate}
 \fincorrection
\correction{004901}
axes d'angle $-\frac12\Arctan 2\ [\pi/2]$, excentricité = $\sqrt{\frac{3\sqrt5-5}2}$.
\fincorrection
\correction{004902}
$\cal C$ : $x^2(1-e^2) + y^2 + 2e^2dx - e^2d^2 = 0$\par
         $T_u,v$ : $x(u(1-e^2)+e^2d) + vy = e^2d(d-u)$\par
         $uv'-u'v = 0  \Rightarrow  x=d$.
\fincorrection
\nocorrection
\nocorrection
\nocorrection
\nocorrection
\nocorrection
\nocorrection
\nocorrection
\nocorrection
\nocorrection
\nocorrection
\nocorrection
\nocorrection
\nocorrection
\nocorrection
\nocorrection
\nocorrection
\nocorrection
\nocorrection
\nocorrection
\nocorrection
\nocorrection
\correction{004984}
\begin{enumerate}
  \item 
  \item 
  \item $\left(\frac{y'}{x'}\right)' = \frac{2t(t^2-t-1)}{(t-1)^2}
               \Rightarrow  t = \frac{1\pm\sqrt5}2$.
  \item $\det(M',M'')=\frac32\sin4t+3\sin2t \Rightarrow t=k\pi$.
\end{enumerate}
 \fincorrection
\correction{004985}
\begin{enumerate}
  \item 
  \item $\frac{m'(t)}{m(t)} =\frac{-2(t-1)(t^2+t+1)(t^3-3t-1)}{t(2t^3+1)(t^3+2+3t)}$.
\end{enumerate}
 \fincorrection
\nocorrection
\nocorrection
\nocorrection
\nocorrection
\correction{004990}
\begin{enumerate}
  \item 
  \item Repère $(0,\vec u_\theta,\vec v_\theta)$ avec $\theta$ constant
             $ \Rightarrow $ point de concours : $X=\cos\theta$, $Y=\sin\theta$.
\end{enumerate}
 \fincorrection
\correction{004991}
Coordonnées polaires :
	     $\rho = \frac a2\Bigl(4\cos\theta-\frac1{\cos\theta}\Bigr)$
	     avec $A = (a,0)$.
\fincorrection
\correction{004992}
$$
\centerline{\includegraphics[height=4cm]{../images/img004992-1}}
$$
%\mapleplot{plot([t/(t^2-1),t,t=-20..20],x=-1..1,y=-0.5..1.5,coords=polar,axes=frame);}%

La courbe est symétrique par rapport à $Oy$. $M(\rho,\theta) = M_1(\rho_1,\theta_1)$
si et seulement si $\begin{cases}\theta_1\equiv \theta\ [2\pi]\cr \rho_1 = \rho\cr\end{cases}$
ou $\begin{cases}\theta_1\equiv \theta+\pi\ [2\pi]\cr \rho_1 = -\rho.\cr\end{cases}$

Dans le premier cas, $\rho=\rho_1 \Leftrightarrow \theta\theta_1 = -1$ pour
$\theta \ne \theta_1$ ce qui donne l'équation en $\theta$ :
$$\theta^2+2k\pi\theta+1 = 0$$
avec $k\in \Z$.
Cette équation à $k$ fixé non nul admet deux racines, ce qui donne deux
familles de points doubles.
De même, le second cas ammène deux autres familles définies par l'équation :
$$\theta^2+(2k+1)\pi\theta+1 = 0.$$
\fincorrection
\correction{005523}
\textbf{(les grands classiques)}
\begin{enumerate}
 \item  \textbf{L'astroïde.}
  \begin{enumerate}
  \item 
 
\textbf{Domaine d'étude.}

\textbullet~Pour tout réel $t$, $M(t)$ existe.

\textbullet~Pour tout réel $t$, $M(t+2\pi)=M(t)$. Par suite, la courbe complète est obtenue quand $t$ décrit un segment de
longueur

$2\pi$ comme par exemple $[-\pi,\pi]$.

\textbullet~Pour tout réel $t$,

$$M(-t)=\left(\begin{array}{c}
\cos^3(-t)\\
\sin^3(-t)
\end{array}\right)=\left(\begin{array}{c}
\cos^3t\\
-\sin^3t
\end{array}\right)
=s_{(Ox)}(M(t)).$$

On étudie et on construit la courbe pour $t\in[0,\pi]$, puis on obtient la courbe complète par réflexion d'axe $(Ox)$.

\textbullet~Pour tout réel $t$,

$$M(t+\pi)=\left(\begin{array}{c}
\cos^3(t+\pi)\\
\sin^3(t+\pi)
\end{array}\right)=\left(\begin{array}{c}
-\cos^3t\\
-\sin^3t
\end{array}\right)
=s_O(M(t)).$$

La portion de courbe obtenue quand $t$ décrit $[-\pi,0]$ est donc aussi la symétrique par rapport à $O$ de la portion de

courbe obtenue quand $t$ décrit $[0,\pi]$. Néanmoins, cette constatation ne permet pas de réduire davantage le domaine

d'éude.

\textbullet~Pour tout réel $t$,

$$M(\pi-t)=\left(\begin{array}{c}
\cos^3(\pi-t)\\
\sin^3(\pi-t)
\end{array}\right)=\left(\begin{array}{c}
-\cos^3t\\
\sin^3t
\end{array}\right)
=s_{(Oy)}(M(t)).$$

On étudie et on construit la courbe pour $t\in\left[0,\frac{\pi}{2}\right]$, puis on obtient la courbe complète par réflexion
d'axe $(Oy)$, puis

par réflexion d'axe $(Ox)$.

\textbullet~Pour tout réel $t$,

$$M\left(\frac{\pi}{2}-t\right)=\left(\begin{array}{c}
\cos^3\left(\frac{\pi}{2}-t\right)\\
\rule{0mm}{6mm}\sin^3\left(\frac{\pi}{2}-t\right)
\end{array}\right)=\left(\begin{array}{c}
\sin^3t\\
\cos^3t
\end{array}\right)
=s_{y=x}(M(t)).$$

On étudie et on construit la courbe pour $t\in\left[0,\frac{\pi}{4}\right]$, puis on obtient la courbe complète par réflexion
d'axe la droite

d'équation $y=x$, puis d'axe $(Oy)$ et enfin d'axe $(Ox)$.\rule{0mm}{5mm}

\textbf{Variations conjointes de} $\bf{x}$ \textbf{et} $\bf{y.}$ La fonction $t\mapsto x(t)$ est strictement
décroissante sur $\left[0,\frac{\pi}{4}\right]$ et la fonction $t\mapsto y(t)$ est strictement croissante sur
$\left[0,\frac{\pi}{4}\right]$.
\textbf{Etude des points singuliers.} Pour $t\in\Rr$,

$$\overrightarrow{\frac{dM}{dt}}(t)=\left(\begin{array}{c}
-3a\cos^2t\sin t\\
3a\sin^2t\cos t
\end{array}\right)=3a\cos t\sin t\left(\begin{array}{c}
-\cos t\\
\sin t
\end{array}\right).$$
Pour tout réel $t$, le vecteur $\left(\begin{array}{c}
-\cos t\\
\sin t
\end{array}\right)$ est unitaire et n'est donc pas nul. Par suite,

$$\overrightarrow{\frac{dM}{dt}}(t)=\overrightarrow{0}\Leftrightarrow 3a\cos t\sin t=0\Leftrightarrow\cos t=0\;\mbox{ou}\;\sin t=0\Leftrightarrow
t\in\frac{\pi}{2}\Zz.$$
Les points singuliers sont donc les $M\left(\frac{k\pi}{2}\right)$, $k\in\Zz$. Pour $t\notin\frac{\pi}{2}\Zz$, $M(t)$ est un
point régulier et la tangente en $M(t)$ est dirigée par le vecteur $\left(\begin{array}{c}
-\cos t\\
\sin t
\end{array}\right)$.
Etudions alors le point singulier $M(0)$. Pour $t\in\left[-\frac{\pi}{2},\frac{\pi}{2}\right]\setminus\{0\}$,

\begin{align*}\ensuremath
\frac{y(t)-y(0)}{x(t)-x(0)}&=\frac{a\sin^3t}{a\cos^3t-a}=\frac{\sin^3t}{(\cos t-1)(\cos^2t+\cos t+1)}\\
 &=\frac{8\sin^3\frac{t}{2}\cos^3\frac{t}{2}}{-2\sin^2\frac{t}{2}(\cos^2t+\cos t+1)}
=\frac{-4\sin\frac{t}{2}\cos^3\frac{t}{2}}{\cos^2t+\cos t+1},
\end{align*}
et donc, $\lim_{t\rightarrow 0}\frac{y(t)-y(0)}{x(t)-x(0)}=0$.
(Si on connaît déjà les équivalents, c'est plus court :  $\frac{\sin^3t}{(\cos t-1)(\cos^2t+\cos t+1)}\underset{x\rightarrow0}{\sim}\frac{t^3}{-\frac{t^2}{2}\times3}=-\frac{2t}{3}\rightarrow0$). La courbe admet en $M(0)$ une tangente dirigée par le vecteur
$(1,0)$. Par symétrie, la courbe admet également une tangente en $M\left(-\frac{\pi}{2}\right)$, $M\left(\frac{\pi}{2}\right)$ et $M(\pi)$,
dirigée respectivement par $(0,1)$, $(0,1)$ et $(1,0)$. Toujours par symétrie, ces quatre points sont des points de
rebroussement de première espèce. Il en résulte aussi que

$$\text{pour tout réel}\;t,\;\mbox{la tangente en}\;M(t)\;\mbox{est dirigée par le vecteur}\;(-\cos t,\sin t).$$
On en déduit la courbe.

$$\includegraphics{../images/img005523-4}$$

     \item  Soit $t\in\left]0,\frac{\pi}{2}\right[$. On a vu que la tangente $(T_t)$ en $M(t)$ est dirigée par le vecteur $(-\cos
t,\sin t)$. Une équation cartésienne de $T_t$ est donc : $-\sin t(x-a\cos^3t)-\cos t(y-a\sin^3t)=0$, ou encore

$$x\sin t+y\cos t=a\sin t\cos t\;(T_t).$$
On en déduit immédiatement que $A(t)$ a pour coordonnées $(a\cos t,0)$ et que $B(t)$ a pour coordonnées$(0,b\sin t)$
puis que

\begin{center}
\shadowbox{
$\forall t\in]0,\frac{\pi}{2}[,\;A(t)B(t)=a.$
}
\end{center}
  \end{enumerate}
 \item  \textbf{La cycloïde.}
  \begin{enumerate}
   \item La condition de roulement sans glissement se traduit par $\overline{OI}=MI$

$$\includegraphics{../images/img005523-6}$$


ou encore $x_{\Omega}=Rt$. On en déduit que 

\begin{center}
$x_M=x_{\Omega}+x_{\overrightarrow{\Omega M}}=Rt+R\cos\left(2\pi-\frac{\pi}{2}-t\right)=Rt-R\sin t=R(t-\sin t)$
\end{center} et

\begin{center}
$y_M=y_{\Omega}+y_{\overrightarrow{\Omega M}}=R+R\sin\left(2\pi-\frac{\pi}{2}-t\right)=R-R\cos t=R(1-\cos t)$.
\end{center}
   \item
\textbf{Domaine d'étude.}

\textbullet~Pour tout réel $t$, $M(t)$ existe.

\textbullet~Pour tout réel $t$, $M(t+2\pi)=M(t)+\overrightarrow{u}$ où $\overrightarrow{u}(
2\pi R,0)$. Par suite, on trace la courbe quand $t$ décrit $[0,2\pi]$ et la

courbe complète est obtenue par translations de vecteurs $k\overrightarrow{u}$, $k\in\Zz$.

\textbullet~Pour tout réel $t$, $M(-t)=(-x(t),y(t))=s_{(Oy)}(M(t))$. On trace la courbe quand $t$ décrit $[0,\pi]$, puis on complète

par réflexion d'axe $(Oy)$ puis par translations.

\textbf{Etude des points singuliers.}
Pour $t\in[0,\pi]$, $x'(t)=R(1-\cos t)=2R\sin^2\left(\frac{t}{2}\right)$ et $y'(t)=R\sin t=2R\sin\left(\frac{t}{2}\right)\cos\left(\frac{t}{2}\right)$. Le point $M(t)$ est régulier si et seulement si $t\in]0,\pi]$. Dans ce cas, la tangente en $M(t)$ est dirigée par $\left(
\begin{array}{c}
2R\sin^2(t/2)\\
2R\sin(t/2)\cos(t/2)
\end{array}
\right)$ ou encore par $\left(
\begin{array}{c}
\sin(t/2)\\
\cos(t/2)
\end{array}
\right)$. Etudions également le point singulier $M(0)$. Pour $t\in]0,\pi]$,

$$\frac{y(t)-y(0)}{x(t)-x(0)}=\frac{R(1-\cos t)}{R(t-\sin t)}\underset{t\rightarrow0}{\sim}\frac{t^2/2}{t^3/6}=\frac{3}{t}.$$
Ainsi, $\displaystyle\lim_{\substack{t\rightarrow0\\ t>0}}\frac{y(t)-y(0)}{x(t)-x(0)}=+\infty$ et la tangente en $M(0)$ est dirigée par $(0,1)$. Ainsi, dans tous les cas, la tangente en $M(t)$ est dirigée par le vecteur $\left(
\begin{array}{c}
\sin(t/2)\\
\cos(t/2)
\end{array}
\right)$. Par symétrie, $M(0)$ est un point de rebroussement de première espèce.
Sinon, $x$ et $y$ sont des fonctions croissantes sur $[0,\pi]$.

$$\includegraphics{../images/img005523-5}$$

  \end{enumerate}
 \item  \textbf{une courbe de }\textsc{Lissajous}
\textbf{Domaine d'étude.}

\textbullet~Pour tout réel $t$, $M(t)$ existe.

\textbullet~Pour tout réel $t$, $M(t+2\pi)=M(t)$ et la courbe complète est obtenue quand $t$ décrit $[-\pi,\pi]$.

\textbullet~Pour tout réel $t$,

$$M(-t)=\left(\begin{array}{c}
\sin(-2t)\\
\sin(-3t)
\end{array}\right)=\left(\begin{array}{c}
-\sin(2t)\\
-\sin(3t)
\end{array}\right)
=s_{O}(M(t)).$$

On étudie et on construit la courbe pour $t\in[0,\pi]$, puis on obtient la courbe complète par symétrie centrale de

centre $O$.

\textbullet~Pour tout réel $t$,

$$M(\pi-t)=\left(\begin{array}{c}
\sin(2\pi-2t)\\
\sin(3\pi-3t)
\end{array}\right)=\left(\begin{array}{c}
-\sin(2t)\\
\sin(3t)
\end{array}\right)
=s_{(Oy)}(M(t)).$$

On étudie et on construit la courbe pour $t\in\left[0,\frac{\pi}{2}\right]$, puis on obtient la courbe complète par réflexion d'axe $(Oy)$ puis 

par symétrie centrale de centre $O$.

\textbullet~On note aussi que $M(t+\pi)=s_{(Ox)}(M(t))$, mais cette constatation ne permet pas de réduire davantage le domaine

d'étude.

\textbf{Variations conjointes de $x$ et $y$.}
Pour $t\in\left[0,\frac{\pi}{2}\right]$, $x'(t)=2\cos(2t)$ et $y'(t)=3\cos(3t)$. On en déduit immédiatement le tableau suivant~:

$$\includegraphics{../images/img005523-1}$$

puis on en déduit la courbe.

$$\includegraphics{../images/img005523-10}$$

\textbf{Points multiples.} D'abord, tout point de l'arc est multiple, puisque la courbe est parcourue une infinité de fois. Il y a essentiellement deux \og vrais points \fg~multiples à déterminer, les autres s'en déduisent par symétrie. L'un des deux est le point de $(Ox)$ d'abscisse strictement positive obtenu pour un certain réel $t$ de $\left]0,\frac{\pi}{2}\right[$.
Soit $t\in\left]0,\frac{\pi}{2}\right[$.

$$y(t)=0\Leftrightarrow\sin(3t)=0\Leftrightarrow
3t\in\pi\Zz\Leftrightarrow
t\in\frac{\pi}{3}\Zz\Leftrightarrow t=\frac{\pi}{3}.$$
Le point de la courbe qui est sur $(Ox)$ et qui a une abscisse strictement positive est le point $M\left(\frac{\pi}{3}\right)=\left(\frac{\sqrt{3}}{2},0\right)$.
Sinon, on cherche $t_1\in\left]0,\frac{\pi}{3}\right[$ et $t_2\in\left]-\frac{2\pi}{3},-\frac{\pi}{2}\right[$ tels que $M(t_1)=M(t_2)$. 

\begin{align*}\ensuremath
M(t_1)=M(t_2)&\Rightarrow x(t_1)=x(t_2)\Leftrightarrow t_2\in t_1+\pi\Zz\;\mbox{ou}\;t_2\in\frac{\pi}{2}-t_1+\pi\Zz
\Rightarrow t_2\in\frac{\pi}{2}-t_1+\pi\Zz\\
 &\Rightarrow t_2=\frac{\pi}{2}-t_1-\pi\Rightarrow t_2=-\frac{\pi}{2}-t_1.
\end{align*}
Réciproquement, si $t_2=-\frac{\pi}{2}-t_1$, alors $x(t_1)=x(t_2)$ et donc,
 
\begin{align*}\ensuremath
M(t_1)=M(t_2)&\Leftrightarrow y\left(-\frac{\pi}{2}-t_1\right)=y(t_1)\Leftrightarrow\sin\left(3\left(-\frac{\pi}{2}-t_1\right)\right)=\sin(3t_1)\\
 &\Leftrightarrow3t_1\in-\frac{3\pi}{2}-3t_1+2\pi\Zz\;\mbox{ou}\;3t_1\in\pi+\frac{3\pi}{2}+3t_1+2\pi\Zz\Leftrightarrow
 6t_1\in-\frac{3\pi}{2}+2\pi\Zz\\
 &\Leftrightarrow t_1\in-\frac{\pi}{4}+\frac{\pi}{3}\Zz\Leftrightarrow t_1=\frac{\pi}{12}
\end{align*}
Le point $M\left(\frac{\pi}{12}\right)=\left(\frac{1}{2},\frac{\sqrt{2}}{2}\right)$ est le point multiple d'abscisse et d'ordonnée strictement positives.
 \item  \textbf{La lemniscate de } \textsc{Bernoulli}
\textbf{Domaine d'étude.}

\textbullet~Pour tout réel $t$, $M(t)$ existe.

\textbullet~Pour tout réel $t$, $M(-t)=s_O(M(t))$. On étudie et construit la courbe quand $t$ décrit $\Rr^+$ et on obtient la courbe

complète par symétrie centrale de centre $O$.

\textbullet~Pour $t>0$,

$$M\left(\frac{1}{t}\right)=\left(
\frac{\frac{1}{t}}{1+\frac{1}{t^4}},
\frac{\frac{1}{t^3}}{1+\frac{1}{t^4}}
\right)=\left(
\frac{t^3}{1+t^4},
\frac{t}{1+t^4}
\right)=s_{y=x}(M(t)).$$
On étudie et construit la courbe quand $t$ décrit $[0,1]$ et on obtient la courbe complète par réflexion d'axe la droite

d'équation $y=x$ puis par symétrie centrale de centre $O$.
\textbf{Variations conjointes de $x$ et $y$.} Les fonctions $x$ et $y$ sont dérivables sur $[0,1]$ et pour $t\in[0,1]$,

$$x'(t)=\frac{(1+t^4)-t(4t^3)}{(1+t^4)^2}=\frac{1-3t^4}{(1+t^4)^2}\;\mbox{et}\;y'(t)=\frac{3t^2(1+t^4)-t^3(4t^3)}{(1+t^4)^2}=\frac{t^2(3-t^4)}{(1+t^4)^2}.$$
On en déduit immédiatement le tableau~:

$$\includegraphics{../images/img005523-2}$$

La tangente en $M(0)$ est dirigée par le vecteur $(1,0)$. Par symétrie, la tangente en \og $M(+\infty)$ \fg~est dirigée par le vecteur $(0,1)$.

$$\includegraphics{../images/img005523-7}$$


 \item  \textbf{Les tractrices}

  \begin{enumerate}
   \item Cherchons les arcs solutions sous la forme $\left\{
\begin{array}{l}
x=f(t)+R\cos t\\
y=R\sin t
\end{array}
\right.$ où $f$ est une foncton dérivable sur un certain intervalle $I$ (de sorte que le point $M(t)$ est sur le cercle $\mathcal{C}(t)$ de centre $\left(
\begin{array}{c}
f(t)\\
0
\end{array}
\right)$ et de rayon $R$). La trajectoire cherchée est orthogonale à chaque cercle $\mathcal{C}(t)$ si et seulement si la tangente à cette trajectoire en $M(t)$ est orthogonale à la tangente au cercle $\mathcal{C}(t)$ en $M(t)$ ou encore \og si et seulement si \fg~les vecteurs $(f'(t)-R\sin t,R\cos t)$ et $(-\sin t,\cos t)$ sont orthogonaux. Cette dernière condition s'écrit $-f'(t)\sin t+R(\sin^2t+\cos^2t)=0$ ou encore $f'(t)=\frac{R}{\sin t}$ ou enfin, $f(t)=R\ln\left|\tan\frac{t}{2}\right|+C$. Les arcs solutions sont les arcs de la forme $t\mapsto\left(
\begin{array}{l}
R\left(\ln\left|\tan\frac{t}{2}\right|+\cos t\right)+C\\
R\sin t
\end{array}
\right)$, où $C\in\Rr$.

$$\includegraphics{../images/img005523-8}$$

Les courbes solutions se déduisent de la courbe $t\mapsto
\left(
\begin{array}{l}
R\left(\ln\left|\tan\frac{t}{2}\right|+\cos t\right)\\
R\sin t
\end{array}
\right)$ par translations de vecteurs colinéaires à $\overrightarrow{i}$. On peut montrer que la courbe obtenue est la trajectoire de la roue arrière d'une voiture quand celle-ci se gare en marche avant, la roue avant étant quant à elle collée au trottoir.
   \item \textbf{Domaine d'étude.} La fonction $t\mapsto M(t)$ est $2\pi$-périodique et on l'étudie donc sur $[-\pi,\pi]$. Pour $t\in[-\pi,\pi]$, $M(t)$ existe si et seulement si $t\in]-\pi,\pi[\setminus\{0\}$.
Pour $t\in]-\pi,\pi[\setminus\{0\}$, $M(-t)=s_{(Ox)}(M(t))$ puis 

$$M(\pi-t)=\left(
\begin{array}{l}
R\left(\ln\left|\tan\left(\frac{\pi}{2}-\frac{t}{2}\right)\right|+\cos(\pi-t)\right)\\
\rule{0mm}{5mm}R\sin(\pi-t)
\end{array}
\right)=\left(
\begin{array}{l}
R\left(-\ln\left|\tan\frac{t}{2}\right|-\cos t\right)\\
\rule{0mm}{5mm}R\sin(t)
\end{array}
\right)=s_{Oy}(M(t).$$
On étudie et on construit la courbe quand $t$ décrit $\left]0,\frac{\pi}{2}\right]$, et on obtient la courbe complète par réflexion d'axe $(Oy)$ puis par réflexion d'axe $(Ox)$.
\textbf{Dérivée. Etude des points singuliers.} Pour $t\in\left]0,\frac{\pi}{2}\right]$,

$$\overrightarrow{\frac{dM}{dt}}(t)=\left(
\begin{array}{l}
R\left(\frac{1}{\sin t}-\sin t\right)\\
\rule{0mm}{5mm}R\cos t
\end{array}
\right)
=\left(
\begin{array}{l}
R\frac{\cos^2t}{\sin t}\\
\rule{0mm}{5mm}R\cos t
\end{array}
\right)=R\frac{\cos t}{\sin t}\left(
\begin{array}{l}
\cos t\\
\sin t
\end{array}
\right).$$
Par suite, $\overrightarrow{\frac{dM}{dt}}(t)=\overrightarrow{0}\Leftrightarrow\frac{\cos^2t}{\sin t}=0\Leftrightarrow t=\frac{\pi}{2}$. Le point $M\left(\frac{\pi}{2}\right)$ est un point singulier.
Quand $t$ tend vers $\frac{\pi}{2}$, $y(t)-y\left(\frac{\pi}{2}\right)=R(\sin t-1)=-R\left(1-\cos\left(\frac{\pi}{2}-t\right)\right)\sim-\frac{R}{2}\left(\frac{\pi}{2}-t\right)^2$. D'autre part, posons $h=\frac{\pi}{2}-t$ ou encore $t=\frac{\pi}{2}-h$. Quand $t$ tend vers $\frac{\pi}{2}$,

$$x'(t)=R\frac{\cos^2t}{\sin t}=R\frac{\sin^2h}{\cos h}\sim Rh^2=R\left(t-\frac{\pi}{2}\right)^2+o\left(\left(t-\frac{\pi}{2}\right)^2\right),$$
et donc par intégration,

$$x(t)-x(\frac{\pi}{2})=\frac{R}{3}\left(t-\frac{\pi}{2}\right)^3+o\left(\left(t-\frac{\pi}{2}\right)^3\right)\sim\frac{R}{3}\left(t-\frac{\pi}{2}\right)^3.$$
Comme d'autre part, $y(t)-y\left(\frac{\pi}{2}\right)=-R(1-\sin t)=-R(1-\cos h)\sim-\frac{R}{2}h^2=-\frac{R}{2}\left(t-\frac{\pi}{2}\right)^2$, on en déduit que

$$\frac{y(t)-y\left(\frac{\pi}{2}\right)}{x(t)-x\left(\frac{\pi}{2}\right)}\sim\frac{-\frac{R}{2}\left(t-\frac{\pi}{2}\right)^2}{\frac{R}{3}\left(t-\frac{\pi}{2}\right)^3}=-\frac{3}{2\left(t-\frac{\pi}{2}\right)},$$
et donc $\displaystyle\lim_{\substack{t\rightarrow\frac{\pi}{2}\\
t<\frac{\pi}{2}}}\frac{y(t)-y(\frac{\pi}{2})}{x(t)-x(\frac{\pi}{2})}=+\infty$. Par symétrie d'axe $(Oy)$, la tangente en $M\left(\frac{\pi}{2}\right)$ est dirigée par $\overrightarrow{j}$ et $M\left(\frac{\pi}{2}\right)$ est un point de rebroussement de première espèce.
Sinon, $x'$ et $y'$ sont strictement positives sur $\left]0,\frac{\pi}{2}\right[$. On en déduit que $x$ et $y$ sont strictement croissantes sur cet intervalle. Quand $t$ tend vers $0$ par valeurs supérieures, $x(t)$ tend vers $-\infty$ et $y(t)$ tend vers $0$. On en déduit que la droite d'équation $x=0$ est asymptote à la courbe. D'autre part, $x$ croit de $-\infty$ à $0$ pendant que $y$ croit de $0$ à $1$.
\textbf{Courbe.}

$$\includegraphics{../images/img005523-9}$$


  \end{enumerate}
\end{enumerate}
\fincorrection
\correction{005524}
\begin{enumerate}

\item 

\textbf{Domaine d'étude.}
$M(t)$ existe si et seulement si $t\notin\{-1,1\}$. Sinon, il n'y a pas de symétrie particulière (la fonction $y$ est effectivement paire, mais $x$ n'est ni paire ni impaire).

\textbf{Dérivée.} Pour $t\in]-1,1[\setminus\{0\}$,

\begin{align*}\ensuremath
x'(t)&=x(t)(3\ln|t|-2\ln|t+1|-\ln|t-1|)'=\frac{t^3}{(t+1)^2(t-1)}(\frac{3}{t}-\frac{2}{t+1}-\frac{1}{t-1})\\
 &=\frac{t^3}{(t+1)^2(t-1)}\frac{3(t^2-1)-2(t^2-t)-(t^2+t)}{t(t+1)(t-1)}=\frac{t^2(t-3)}{(t+1)^3(t-1)^2},
\end{align*}

et

$$y'(t)=\frac{2t(t^2-1)-2t(t^2)}{(t^2-1)^2}=\frac{-2t}{(t^2-1)^2},$$

ce qui reste vrai par continuité de $x$ et $y$ en $0$.

\textbf{Etude des points singuliers.} Pour $t\in]-1,1[$, $\overrightarrow{\frac{dM}{dt}}(t)=\vec{0}\Leftrightarrow t=0$. $M(0)=(0,0)$ est l'unique point singulier. Pour $t\in]-1,1[\setminus\{0\}$,

$$\frac{y(t)-y(0)}{x(t)-x(0)}=\frac{t^2}{t^2-1}\frac{(t+1)^2(t-1)}{t^3}=\frac{t+1}{t}.$$

Par suite, $\frac{y(t)-y(0)}{x(t)-x(0)}$ tend vers $+\infty$ quand $t$ tend vers $0$ par valeurs supérieures et vers $-\infty$ quand $t$ tend vers $0$ par valeurs inférieures. La tangente en $M(0)$ est dirigée par $\vec{j}$ et d'autre part, $M(0)$ est un point de rebroussement de première espèce.

\textbf{Etude quand $t$ tend vers $\pm\infty$.} Quand $t$ tend vers $\pm\infty$, $M(t)$ tend vers le point $(1,1)$. On prolonge la courbe en posant $M(\infty)=1,1)$. On a alors

$$\frac{y(t)-y(\infty)}{x(t)-x(\infty)}=(\frac{t^2}{t^2-1}-1)(\frac{t^3}{(t+1)^2(t-1)}-1)^{-1}=\frac{1}{t^2-1}
\frac{(t+1)^2(t-1)}{-t^2+t+1}=\frac{t+1}{-t^2+t+1}\sim-\frac{1}{t}.$$

Cette expression tend donc vers $0$ quand $t$ tend vers $\pm\infty$ et la tangente en $M(\infty)$ est dirigée par $\vec{i}$.

\textbf{Etude quand $t$ tend vers 1.} Quand $t$ tend vers $1$, $x(t)\sim{1}{4(t-1)}$ et $y(t)\sim\frac{1}{2(t-1)}$. Donc, $x$ et $y$ tendent vers l'inifini et il y a branche infinie. De plus, $\frac{y(t)}{x(t)}\sim2$. Puis,

$$y(t)-2x(t)=\frac{t^2}{t^2-1}-2\frac{t^3}{(t+1)^2(t-1)}=\frac{t^2(t+1)-2t^3}{(t+1)^2(t-1)}=-\frac{t^2}{(t+1)^2}.$$

Cette dernière expression tend vers $-\frac{1}{4}$ et la droite $(\Delta)$ d'équation $y=2x-\frac{1}{4}$ est asymptote à la courbe.

\textbf{Etude quand $t$ tend vers -1.} Quand $t$ tend vers $1$, $x(t)\sim{1}{2(t+1)^2}$ et $y(t)\sim\frac{-1}{2(t+1)}$. Donc, $x$ et $y$ tendent vers l'inifini et il y a branche infinie. De plus, $\frac{y(t)}{x(t)}\sim-(t+1)$. Par suite, $\frac{y(t)}{x(t)}$ tend vers $0$ quand $t$ tend vers $-1$. La courbe admet une barnche parabolique de direction $(Ox)$.

\textbf{Variations conjointes de $x$ et $y$.} On rappelle que pour $t\in\Rr\setminus\{-1,1\}$, $x'(t)=\frac{t^2(t-3)}{(t+1)^3(t-1)^2}$ et $y'(t)=\frac{-2t}{(t^2-1)^2}$. On en déduit le tableau suivant~:

$$\includegraphics{../images/img005524-1}$$


On peut noter que la tangente en $M(3)$ est dirigée par le vecteur $\vec{j}$. Voir graphique page suivante.

$$\includegraphics{../images/img005524-2}$$



Dans la suite de cet exercice, je ne détaillerai que très peu ou pas du tout l'étude de la courbe.
\item  

$$\includegraphics{../images/img005524-3}$$


\item  

$$\includegraphics{../images/img005524-4}$$



\item  $\left\{\begin{array}{l}
x=\frac{2t}{1+t^2}\\
y=\frac{t+2}{1-t^2}
\end{array}\right.$

$$\includegraphics{../images/img005524-5}$$



\item 

$$\includegraphics{../images/img005524-6}$$


\item  $\left\{\begin{array}{l}
x=\frac{t^3}{t^2-9}\\
y=\frac{t(t-2)}{t-3}
\end{array}\right.$

$$\includegraphics{../images/img005524-7}$$


\item  $\left\{\begin{array}{l}
x=\frac{t^3}{1+3t}\\
y=\frac{3t^2}{1+3t}
\end{array}\right.$

$$\includegraphics{../images/img005524-8}$$


\item  $\left\{\begin{array}{l}
x=t^2+t^3\\
y=t^2+t^3-2t^4-2t^5
\end{array}\right.$

$$\includegraphics{../images/img005524-9}$$

\end{enumerate}
\fincorrection
\correction{005525}
\begin{enumerate}
\item  On a vu dans l'exercice \ref{exo:routhe1}, que la tangente $(T_t)$ en $M(t)$ est toujours dirigée par le vecteur $\vec{u}(t)=(-\cos t,\sin t)$. Une équation de la tangente en $M(t)$ est donc $\sin t(x-a\cos^3t)+\cos t(y-a\sin^3t)=0$ ou encore

$$x\sin t+y\cos t=a\sin t\cos t\;(T_t).$$

Soit $(t,u)\in[-\pi,\pi]^2$. 

$$(T_t)\bot(T_u)\Leftrightarrow\vec{u}(t)|\vec{u}(u)=0\Leftrightarrow\cos t\cos u+\sin t\sin u=0\Leftrightarrow\cos(t-u)=0\Leftrightarrow u\in t+\frac{\pi}{2}+\pi\Zz.$$

Il est alors clair que l'orthoptique est l'ensemble des points d'intersection des tangente $(T_t)$ et $(T_{t+\frac{\pi}{2}})$ quand $t$ décrit $\Rr$.

\begin{align*}\ensuremath
M(x,y)\;(T_t)\cap(T_{t+\frac{\pi}{2}})&\Leftrightarrow
\left\{
\begin{array}{l}
x\sin t+y\cos t=a\sin t\cos t\\
x\cos t-y\sin t=-a\sin t\cos t
\end{array}
\right.
\\
 &\Leftrightarrow
x=-\left|
\begin{array}{cc}
a\sin t\cos t&\cos t\\
-a\sin t\cos t&-\sin t
\end{array}
\right|
\;\mbox{et}
\;
y=-\left|
\begin{array}{cc}
\sin t&a\sin t\cos t\\
\cos t&-a\sin t\cos t
\end{array}
\right|
\\
 &\Leftrightarrow x=a\sin t\cos t(-\cos t+\sin t)\;\mbox{et}\;y=a\sin t\cos t(\cos t+\sin t)
\end{align*}

L'orthoptique cherchée est la courbe $t\mapsto\left(
\begin{array}{c}
a\sin t\cos t(-\cos t+\sin t)\\
a\sin t\cos t(\cos t+\sin t)
\end{array}
\right)$.

$$\includegraphics{../images/img005525-1}$$


\end{enumerate}
\fincorrection
\nocorrection
\nocorrection
\nocorrection
\correction{006981}\ 
\begin{enumerate}
\item Pour $f(x)=\sin^2x+\cos x$, le domaine de définition de $f$ est $\R$, et 
$f$ est de classe $\mathcal{C}^\infty$. On remarque que $f$ est $2\pi$-périodique et paire, 
il suffit donc de faire l'étude de $f$ sur l'intervalle $[0;\pi]$. 
\begin{itemize}
\item Variations de $f$\\
Pour $x\in [0;\pi]$, $f'(x)=2\sin x\cos x-\sin x=\sin x(2\cos x-1)$ et donc $f'(x)=0$ 
si et seulement si $x\in\{0;\frac{\pi}{3};\pi\}$. Comme $\sin x>0$ si $x\in]0;\pi[$, 
pour étudier le signe de $f'(x)$, il suffit d'étudier le signe de $(2\cos x-1)$, et on obtient 
$$\begin{array}{c|lcccr}
x&0&\ &\frac{\pi}{3}&\ &\pi \\\hline
f'(x)&0&+&0&-&0\\\hline
\ &\ &\ &\frac{5}{4}&\ &\ \\
f&\ &\nearrow &\ &\searrow &\ \\
\ &1 &\ &\ &\ &-1 \\
\end{array}$$
\item Tangentes horizontales\\
Le graphe de $f$ possède une tangente horizontale là où $f'$ s'annule, c'est-à-dire aux points de coordonnées $(0,1)$, $(\frac{\pi}{3},\frac{5}{4})$ et $(\pi,-1)$. En particulier, la tangente au point d'abscisse 0 est horizontale et a pour équation $y=1$. Pour déterminer la position de la courbe par rapport à sa tangente en ce point, on étudie le signe de $f(x)-1$ pour $x$ proche de 0:
$$f(x)-1=\sin^2x-1+\cos x=-\cos^2x+\cos x=\cos x(1-\cos x)$$
Cette expression est positive au voisinage de $0$ (et m\^eme $>0$ pour $x\not=0$ proche de $0$). La courbe est donc au-dessus de sa tangente.
\item Points particuliers\\
Le graphe de $f$ coupe l'axe des abscisses entre $0$ et $\pi$ en un unique point $x_0$, qu'on détermine en résolvant
$$f(x)=0\Longleftrightarrow 1-\cos^2x+\cos x=0 \Longleftrightarrow X^2-X-1=0\ (X=\cos x)$$
ce qui donne deux solutions pour $X$, mais une seule dans $[-1;1]$: $X=\frac{1-\sqrt{5}}{2}$ et donc 
$x_0= \Arccos (\frac{1-\sqrt{5}}{2})$.
\end{itemize}
Le graphe de $f$ est obtenu sur $[-\pi;\pi]$ par symétrie par rapport à l'axe des ordonnées, 
puis sur $\R$ par $2\pi$-périodicité.

\begin{center}
\begin{tikzpicture}[scale=1]
     \draw[->,>=latex,thick, gray] (-3.2,0)--(6.7,0) node[below,black] {$x$};
     \draw[->,>=latex,thick, gray] (0,-1.5)--(0,2.5) node[right,black] {$y$};  


  \draw[red, very thick,domain=0:pi,samples=100] plot ({\x},{sin(\x r)^2+cos(\x r)});
  \draw[gray, very thick,domain=-pi:0,samples=100] plot ({\x},{sin(\x r)^2+cos(\x r)});
  \draw[gray, very thick,domain=pi:2*pi,samples=100] plot ({\x},{sin(\x r)^2+cos(\x r)});

  \fill[blue] (2.23,0) circle (2pt) node[below left]{$x_0$};  
  \draw[blue, dashed] (0,0)--(0,1);  
  \fill[blue] (0,0) circle (2pt) node[below left]{$0$}; 
  \draw[blue, dashed] (3.14,-1)--(3.14,0);
  \fill[blue] (3.14,0) circle (2pt)  node[above]{$\pi$}; 
  \draw[blue, dashed] (1.04,1.25)--(1.04,0);  

  \fill[blue] (1.04,0) circle (2pt)  node[below]{$\frac\pi3$}; 

  \draw[blue] (0,1)--+(1,0)--+(-1,0);
  \draw[blue] (3.14,-1)--+(1,0)--+(-1,0);
  \draw[blue] (1.04,1.25)--+(1,0)--+(-1,0);
  \fill[blue] (0,1) circle (2pt); 
  \fill[blue] (3.14,-1) circle (2pt); 
  \fill[blue] (1.04,1.25) circle (2pt); 
\end{tikzpicture}
\end{center}

\begin{center}
\begin{tikzpicture}[scale=0.5]
  \draw[->,>=latex,thick, gray] (-4,0)--(17,0) node[above,black] {$x$};
  \draw[->,>=latex,thick, gray] (0,-1.5)--(0,2.5) node[right,black] {$y$};  
  \draw[red, very thick,domain=-pi:5*pi,samples=100] plot ({\x},{sin(\x r)^2+cos(\x r)})
  node[black, right] {$y = \sin^2x+\cos x$};
\end{tikzpicture}
\end{center}

\item Pour $f(x)=x+\ln(1+e^x)$, le domaine de définition de $f$ est $\R$ et $f$ est de classe $\mathcal{C}^\infty$.
\begin{itemize}
\item Variations de $f$\\
Comme $f'(x)=1+\frac{e^x}{1+e^x}$, pour tout $x$, $f'(x)>1$. En particulier $f$ est strictement croissante sur $\R$. 
\item Allure du graphe en $+\infty$\\
On a $f(x)\xrightarrow[x\to +\infty]{}+\infty$ et
$$\frac{f(x)}{x}=1+\frac{\ln\big( e^x (e^{-x}+1)\big)}{x}=1+\frac{x+\ln(e^{-x}+1)}{x}\xrightarrow[x\to +\infty]{}2$$
puis $f(x)-2x=\ln(e^{-x}+1)\xrightarrow[x\to +\infty]{}0^+$. Ainsi le graphe de $f$ a en $+\infty$ une asymptote, d'équation $y=2x$, et reste au-dessus de cette asymptote.
\item Allure du graphe en $-\infty$\\
On a $f(x)\xrightarrow[x\to -\infty]{}-\infty$ et
$$\frac{f(x)}{x}=1+\frac{\ln(1+e^x)}{x}\xrightarrow[x\to -\infty]{}1$$
puis $f(x)-x=\ln(1+e^{x})\xrightarrow[x\to -\infty]{}0^+$. Ainsi le graphe de $f$ 
a en $-\infty$ une asymptote, d'équation $y=x$, et reste au-dessus de cette asymptote.
\item Tangente au point d'abscisse $0$\\
L'équation de la tangente au graphe de $f$ au point d'abscisse $x_0$, 
et la position du graphe par rapport à cette tangente, peuvent \^etre 
obtenues simultanément à partir du développement limité de $f$ en $x_0$. 
Pour l'équation de la tangente, un développement limité à l'ordre $1$ suffit, 
mais pour avoir la position il faut pousser le développement limité à l'ordre $2$ 
(ou à l'ordre $3$ si le terme d'ordre $2$ est nul, ou plus encore...):
\begin{eqnarray*}
f(x)&=&x+\ln(1+e^x)=x+\ln\left(1+1+x+\frac{1}{2}x^2+o(x^2)\right)\\
 &=&x+\ln 2+\ln\left(1+\frac{1}{2}x+\frac{1}{4}x^2+o(x^2)\right)\\
 &=&x+\ln 2+\left(\frac{1}{2}x+\frac{1}{4}x^2\right)-\frac{1}{2}\left(\frac{1}{2}x+\frac{1}{4}x^2\right)^2+o(x^2)\\
 &=&\ln 2+\frac{3}{2}x+\frac{1}{8}x^2+o(x^2)
\end{eqnarray*}
L'équation de la tangente au point d'abscisse $0$ (donnée par le DL à l'ordre $1$) est donc 
$$y=\ln 2+\frac{3}{2}x$$ 
De plus, $f(x)-\left(\ln 2+\frac{3}{2}x\right)=\frac{1}{8}x^2+o(x^2)=\frac{1}{8}x^2(1+o(1))$
où $o(1)$ est un terme qui tend vers $0$ quand $x\to 0$. Ainsi $(1+o(1))$ a le m\^eme signe que $1$ 
pour $x$ proche de $0$, et $f(x)-\left(\ln 2+\frac{3}{2}x\right)$ est positif au voisinage de $0$: 
la courbe reste localement au-dessus de sa tangente.
\end{itemize}

\begin{center}
\begin{tikzpicture}[scale=0.8]

     \draw[->,>=latex,thick, gray] (-4,0)--(2.5,0) node[below,black] {$x$};
     \draw[->,>=latex,thick, gray] (0,-4.5)--(0,4.5) node[left,black] {$y$};  


  \draw[red, very thick,domain=-4:2,samples=100] plot ({\x},{\x + ln(1+exp(\x))})
  node[above] {$y=x+\ln(x+e^x)$};

  \draw[thick, blue] (-1,-1)--(-4,-4) node[left] {$y=x$};
    \draw[thick, blue] (0.5,1)--(2,4) node[right] {$y=2x$};
  \draw[thick, blue] (1,2.19)--(-1,-0.80) node[pos=0, left] {$y=\ln 2+\frac32 x$};

  \fill[blue] (0,0.69) circle (2pt)  node[left]{$\ln 2$}; 
  \fill[blue] (0,0) circle (2pt) node[below right]{$0$}; 
  \fill[blue] (1,0) circle (2pt) node[below]{$1$};
\end{tikzpicture}
\end{center}
\end{enumerate}
\fincorrection
\correction{006982}\
\begin{enumerate}
\item Pour transformer une équation cartésienne $y=f(x)$ en paramétrisation, il suffit de poser $x=t$ et $y=f(t)$, en faisant décrire au paramètre $t$ le domaine de définition de $f$. Ici, $f(x)=\sqrt{-x^2-3x+4}$ est bien définie pour les $x\in\R$ tels que $-x^2-3x+4\ge 0$ \textsl{i.e.} $x\in[-4;-1]$. On obtient donc la paramétrisation suivante:
$$\left\{\begin{array}{l}
x(t)=t\\
y(t)=\sqrt{-t^2-3t+4}
\end{array}\quad (t\in[-4;1])\right.$$
ce qui signifie
\begin{eqnarray*}
(x,y)\in\mathcal{C}&\Longleftrightarrow&\left\{\begin{array}{l}
x\in[-4;1]\\
y=\sqrt{-x^2-3x+4}
\end{array}\right.\\
 &\Longleftrightarrow&\exists\, t\in[-4;1]\ |\ \left\{\begin{array}{l}
x(t)=t\\
y(t)=\sqrt{-t^2-3t+4}
\end{array}\right.
\end{eqnarray*}
où $\mathcal{C}$ est la courbe étudiée.

\begin{center}
 \begin{tikzpicture}[scale=1]

     \draw[->,>=latex,thick, gray] (-4.5,0)--(2,0) node[below,black] {$x$};
     \draw[->,>=latex,thick, gray] (0,-1)--(0,6) node[right,black] {$y$};  


  \draw[red, very thick,domain=-4:1,samples=100] plot ({\x},{sqrt(-\x*\x-3*\x+4))}) ;
  
  \node[black, above right] at (-4,5.6) {$y=\sqrt{-x^2-3x+4}$};


  \draw[blue, dashed] (-4,0)--(-4,5.6); 
  \fill[blue] (-4,0) circle (2pt)  node[below]{$-4$}; 

 
  \fill[blue] (1,0) circle (2pt)  node[below]{$1$}; 
\end{tikzpicture} 
\end{center}

\item S'il est toujours possible de représenter le graphe d'une fonction comme une courbe 
paramétrée, la réciproque n'est pas vraie. Ici, la courbe considérée est le cercle de 
rayon $1$ centré au point $(3,0)$. Ce n'est donc pas un graphe de fonction, puisque 
plusieurs points de la courbe ont la m\^eme abscisse: conna\^itre $x$ ne donne pas 
$y$! Par exemple, pour $t=\pm\frac{\pi}{2}$, on obtient les deux points de la courbe $(3,-1)$ et $(3,+1)$.
\begin{center}
\begin{tikzpicture}[scale=1.2]

     \draw[->,>=latex,thick, gray] (-1.5,0)--(5,0) node[below,black] {$x$};
     \draw[->,>=latex,thick, gray] (0,-1.5)--(0,1.5) node[right,black] {$y$};  


  \draw[red, very thick,domain=-4:1,samples=100] (3,0) circle (1cm);
  \node[black, above right] at (4,0.6) {$(\cos t + 3,\sin t)$};

  \fill[blue] (0,0) circle (2pt) node[below left]{$0$}; 
  \fill[blue] (3,0) circle (2pt)  node[below]{$3$}; 
  
  \fill[blue] (3,1) circle (2pt)  node[above]{$(3,+1)$}; 
  \fill[blue] (3,-1) circle (2pt)  node[below]{$(3,-1)$};   
\end{tikzpicture}  
\end{center}


\item On constate, en utilisant la formule $\sin^2t=1-\cos^2t=-1-x(t)$, que
\begin{eqnarray*}
y(t)&=&\sin^4t+4\sin^2t+4=(-1-x(t))^2+4(-1-x(t))+4\\
 &=&x(t)^2-2x(t)+1 = (x(t)-1)^2
\end{eqnarray*}
Ainsi les points $(x,y)$ de la courbe vérifient l'équation $y=(x-1)^2$. 
De plus, lorsque le paramètre $t$ décrit $\R$, $x(t)=\cos^2t-2$ décrit l'intervalle $[-2;-1]$. Finalement, 
\begin{eqnarray*}
(x,y)\in\mathcal{C}&\Longleftrightarrow&\exists\, t\in\R\ |\ \left\{\begin{array}{l}
x(t)=\cos^2t-2\\
y(t)=\sin^4t+4\sin^2t+4
\end{array}\right.\\
\ &\Longleftrightarrow&\left\{\begin{array}{l}
x\in[-2;-1]\\
y=(x-1)^2
\end{array}\right.
\end{eqnarray*}
et la courbe est donc le graphe de la fonction
$$\begin{array}{rccc}
f:&[-2;-1]&\to&\R\\
\ &x&\mapsto&(x-1)^2
\end{array}$$

\begin{center}
\begin{tikzpicture}[scale=0.6]

     \draw[->,>=latex,thick, gray] (-4.5,0)--(3.5,0) node[below,black] {$x$};
     \draw[->,>=latex,thick, gray] (0,-1)--(0,9.5) node[right,black] {$y$};  


  \draw[red, very thick,domain=-2:-1,samples=100] plot ({\x},{(\x-1)*(\x-1)}) ;
  \draw[gray, very thick,domain=-1:3,samples=100] plot ({\x},{(\x-1)*(\x-1)}) ;  
  \node[black, left] at (-2,9) {$y=(x-1)^2$};


  \draw[blue, dashed] (-2,0)--(-2,9); 
  \fill[blue] (-2,0) circle (2pt)  node[below]{$-2$}; 
  \draw[blue, dashed] (-1,0)--(-1,4); 
  \fill[blue] (-1,0) circle (2pt)  node[below]{$-1$};
 
  \fill[blue] (1,0) circle (2pt)  node[below]{$1$}; 
\end{tikzpicture}
\end{center}

\end{enumerate}
\fincorrection
\correction{006983}\
\begin{enumerate}
\item Les expressions $x(t)=\cos^3t$ et $y(t)=\sin^3t$ sont bien définies pour tout $t\in\R$.
\begin{itemize}
\item Réduction de l'intervalle d'étude\\
Les fonctions $x$ et $y$ étant $2\pi$-périodiques, il suffit de restreindre l'étude à un intervalle de longueur $2\pi$ pour obtenir l'intégralité du support de la courbe.\\

La fonction $x$ est paire, la fonction $y$ est impaire: on fait donc l'étude sur $[0;\pi]$, puis la courbe complète sera obtenue par symétrie par rapport à l'axe $(Ox)$.\\

On constate que $x(\pi-t)=-x(t)$ et que $y(\pi-t)=y(t)$, par conséquent les points $M(\frac{\pi}{2}-t)$ et $M(\frac{\pi}{2}+t)$ sont symétriques par rapport à l'axe $(Oy)$: on restreint donc l'étude à $[0;\frac{\pi}{2}]$, puis on complète par symétrie par rapport à $(Oy)$.\\
\fbox{Finalement, on fait l'étude sur $[0;\frac{\pi}{2}]$} puis on complète
en utilisant successivement les symétries par rapport à $(Oy)$ et $(Ox)$.
\item Tableau de variations conjointes\\
Les fonctions $x$ et $y$ sont de classe $\mathcal{C}^1$. Soit $t\in[0;\frac{\pi}{2}]$:
$$\begin{array}{lcl}
x(t)=\cos^3t&\ &y(t)=\sin^3t\\
x'(t)=-3\sin t\cos^2t&\ &y'(t)=3\cos t\sin^2t\\
x'(t)<0\Longleftrightarrow t\in]0;\frac{\pi}{2}[&\ &y'(t)>0\Longleftrightarrow t\in]0;\frac{\pi}{2}[\\
x'(t)=0\Longleftrightarrow t\in\{0;\frac{\pi}{2}\} &\ &y'(t)=0\Longleftrightarrow t\in\{0;\frac{\pi}{2}\}
\end{array}$$
$$\begin{array}{c|lcr}
t&0&\ &\frac{\pi}{2}\\\hline
x'(t)&0&-&0\\\hline
\ &1 &\ &\ \\
x&\ &\searrow &\ \\
\ & &\ &0 \\\hline
\ &\ &\ &1 \\
y&\ &\nearrow &\ \\
\ &0 &\ & \\\hline
y'(t)&0&+&0\\
\end{array}$$

Cela signifie que lorsque $t$ varie de $0$ à $\frac\pi2$ la courbe va vers la gauche (car $x(t)$ décroît)
en montant (car $y(t)$ croît) du point $(1,0)$ à $(0,1)$.

\item Points particuliers\\
  \begin{itemize}
    \item $M(\frac\pi6) = \big(\frac{3\sqrt3}{8},\frac18\big) = (0.64\ldots,0.125) $ ; 
    la tangente est dirigée par $\big(x'(\frac\pi6),y'(\frac\pi6)\big) = \big(-\frac98,\frac{3\sqrt3}{8} \big) = (-1.125,0.64\ldots)$.
    \item $M(\frac\pi4) = \big(\frac{\sqrt2}{4},\frac{\sqrt2}{4}\big) = (0.35\ldots,0.35\ldots)$ ; 
    la tangente est dirigée par $\big(x'(\frac\pi4),y'(\frac\pi4)\big) = \big(-\frac{3\sqrt2}{4},\frac{3\sqrt2}{4}\big) = (-1.06\ldots,1.06\ldots)$.
    \item $M(\frac\pi3) =  \big(\frac18,\frac{3\sqrt3}{8}\big) = (0.125,0.64\ldots)$ ;
    la tangente est dirigée par $\big(x'(\frac\pi3),y'(\frac\pi3)\big) = \big(-\frac{3\sqrt3}{8},\frac98\big) = (-0.64\ldots,1.125)$. 
  \end{itemize}

\item \'Etude des points singuliers\\
Le point $M(t)$ est singulier si $x'(t)=y'(t)=0$, ce qui est le cas dans le domaine d'étude $[0;\frac{\pi}{2}]$ uniquement pour $t=0$ et $t=\frac{\pi}{2}$. Pour déterminer la tangente au point $M(0)$ (de coordonnées cartésiennes $(1,0)$), on étudie la limite en 0 de 
$$\frac{y(t)-y(0)}{x(t)-x(0)}=\frac{\sin^3t}{\cos^3t-1}$$ 
Or $\sin^3t\underset{0}{\sim}t^3$ et $\cos^3t-1=(1-\frac{t^2}{2}+o(t^2))^3-1\underset{0}{\sim}-3\frac{t^2}{2}$, donc le quotient est équivalent à $-\frac{2}{3}t$ et tend vers 0 en 0. Ainsi, $\mathcal{C}$ admet au point $M(0)$ une tangente, de pente nulle c'est-à-dire horizontale.
\end{itemize}

\begin{center}
\begin{tikzpicture}[scale=4]
     \draw[->,>=latex,thick, gray] (-1.1,0)--(1.3,0) node[below,black] {$x$};
     \draw[->,>=latex,thick, gray] (0,-1.1)--(0,1.3) node[right,black] {$y$};  



  \fill[black] (0,0) circle (0.75pt) node[below left]{$0$}; 


  \fill[blue] (1,0) circle (0.75pt)  node[below]{$M(0)$}; 
  \fill[blue] (0,1) circle (0.75pt)  node[right]{$M(\frac\pi2)$}; 
  \draw[very thick, blue] (1,0)--+(0.2,0)--+(-0.2,0);
  \draw[very thick, blue] (0,1)--+(0,0.2)--+(0,-0.2);
  
  \def\epsilon{0.2};
  \coordinate (M1) at (0.64,0.125);
  \fill[blue] (M1) circle (0.75pt)  node[right]{$M(\frac\pi6)$}; 
  \draw[very thick, blue] (M1)--+(-\epsilon*1.125,\epsilon*0.64)--+(\epsilon*1.125,-\epsilon*0.64);

  \coordinate (M2) at (0.35,0.35);
  \fill[blue] (M2) circle (0.75pt)  node[right]{$M(\frac\pi4)$}; 
  \draw[very thick, blue] (M2)--+(-\epsilon*1.06,\epsilon*1.06)--+(\epsilon*1.06,-\epsilon*1.06);


    \coordinate (M3) at (0.125,0.64);
  \fill[blue] (M3) circle (0.75pt)  node[right]{$M(\frac\pi3)$}; 
  \draw[very thick, blue] (M3)--+(\epsilon*0.64,-\epsilon*1.125)--+(-\epsilon*0.64,\epsilon*1.125);

  \draw (0,0)--(M1);
  \draw (0,0)--(M2);
  \draw (0,0)--(M3);
  
  \draw[red, very thick,domain=0:pi/2,samples=100] plot ({cos(\x r)^3},{sin(\x r)^3});
  \draw[gray, very thick,domain=pi/2:2*pi,samples=100] plot ({cos(\x r)^3},{sin(\x r)^3});

\end{tikzpicture}  
\end{center}

\item Les expressions $x(t)=t-\tanh t$ et $y(t)=\frac{1}{\ch t}$ sont bien définies pour tout $t\in\R$.
\begin{itemize} 
\item Réduction du domaine d'étude\\
Comme $x$ est impaire et $y$ paire, \fbox{on restreint l'étude à $\R^+$} puis on complète par symétrie par rapport à l'axe $(Oy)$.
\item Tableau de variations conjointes\\
Les fonctions $x$ et $y$ sont de classe $\mathcal{C}^1$. Pour $t\in\R^+$:
$$\begin{array}{lcl}
x(t)=t-\tanh t&\ &y(t)=\frac{1}{\ch t}\\
x'(t)=\tanh ^2t&\ &y'(t)=-\frac{\sh t}{\ch ^2t}\\
x'(t)>0\Longleftrightarrow t>0&\ &y'(t)<0\Longleftrightarrow t>0\\
x'(t)=0\Longleftrightarrow t=0 &\ &y'(t)=0\Longleftrightarrow t=0
\end{array}$$
$$\begin{array}{c|lcr}
t&0&\ &+\infty\\\hline
x'(t)&0&+&\ \\\hline
\ &  &\ &+\infty\ \\
x&\ &\nearrow &\ \\
\ &0 &\ & \\\hline
\ &1 &\ & \\
y&\ &\searrow &\ \\
\ & &\ & 0\\\hline
y'(t)&0& - &\ \\
\end{array}$$
Cela signifie que le courbe va vers la droite et vers le bas lorsque $t$ va de $0$ à $+\infty$.


\item \'Etude des points singuliers\\
Le seul point singulier est $M(0)$, or $\displaystyle \frac{y(t)-y(0)}{x(t)-x(0)}=\frac{1-\ch t}{t\ch t-\mathrm{sh t}}$ et 
$$\frac{1-\ch t}{t\ch t-\mathrm{sh t}}=\frac{-t^2/2+o(t^2)}{t(1+t^2/2+o(t^2))-(t+t^3/6+o(t^3))}\underset{0}{\sim}\frac{-t^2/2}{t^3/3}$$ 
et par conséquent $\frac{y(t)-y(0)}{x(t)-x(0)}\xrightarrow[t\to 0^+]{}-\infty$. Ainsi $\mathcal{C}$ possède une tangente verticale au point $M(0)$ de coordonnées cartésiennes $(0,1)$.
\item \'Etude des branches infinies\\
Comme $x(t)\xrightarrow[t\to+\infty]{}+\infty$ et $y(t)\xrightarrow[t\to+\infty]{}0$, l'axe des abscisses est asymptote à $\mathcal{C}$.
\end{itemize}

\begin{center}
\begin{tikzpicture}[scale=1.5]
     \draw[->,>=latex,thick, gray] (-3.5,0)--(3.5,0) node[below,black] {$x$};
     \draw[->,>=latex,thick, gray] (0,-.5)--(0,1.5) node[right,black] {$y$};  


  \draw[red, very thick,domain=0.2:4,samples=100] plot ({\x - tanh(\x)},{1/cosh(\x)});
  \draw[gray, very thick,domain=-0.2:-4,samples=100] plot ({\x - tanh(\x)},{1/cosh(\x)});

  \fill[black] (0,0) circle (1pt) node[below left]{$0$}; 


  \fill[black] (1,0) circle (1pt)  node[below]{$1$}; 
  
  \fill[blue] (0,1) circle (1pt)  node[right]{$M(0)$}; 
  \draw[very thick, blue] (0,1)--+(0,0.3)--+(0,-0.7);
  \draw[very thick, blue] (2.3,0)--+(1,0)--+(-1,0);
  
 

\end{tikzpicture}  
\end{center}

\item Les expressions $x(t)=t-\sin t$ et $y(t)=1-\cos t$ sont bien définies pour $t\in\R$.
\begin{itemize}
\item Réduction du domaine d'étude\\
On remarque que $x(t+2\pi)=2\pi+x(t)$ et $y(t+2\pi)=y(t)$: le point $M(t+2\pi)$ 
se déduit de $M(t)$ par translation de vecteur $2\pi\cdot\vec{i}$. 
Il suffit donc d'étudier la courbe sur l'intervalle $[-\pi;\pi]$.

La fonction $x$ étant impaire et $y$ paire, on restreint l'étude à $[0;\pi]$ 
puis on complète par symétrie par rapport à l'axe $(Oy)$. \\
\fbox{Finalement, on fait l'étude sur $[0;\pi]$} puis on complète
en utilisant successivement la symétrie par rapport à $(Oy)$, puis des translations successives 
de vecteur $2\pi\cdot\vec{i}$.

\item Tableau de variations conjointes\\
Les fonctions $x$ et $y$ sont de classe $\mathcal{C}^1$. Soit $t\in[0;\pi]$:
$$\begin{array}{lcl}
x(t)=t-\sin t&\ &y(t)=1-\cos t\\
x'(t)=1-\cos t&\ &y'(t)=\sin t\\
x'(t)>0\Longleftrightarrow t>0&\ &y'(t)>0\Longleftrightarrow 0<t<\pi\\
x'(t)=0\Longleftrightarrow t=0 &\ &y'(t)=0\Longleftrightarrow t\in\{0;\pi\}
\end{array}$$
$$\begin{array}{c|lcr}
t&0&\ &\pi\\\hline
x'(t)&0&+&2 \\\hline
\ &  &\ &\pi \\
x&\ &\nearrow &\ \\
\ &0 &\ & \\\hline
\ & &\ & 2\\
y&\ &\nearrow &\ \\
\ &0 &\ & \\\hline
y'(t)&0& + &0\ \\
\end{array}$$
La courbe va vers la droite en montant lorsque $t$ varie de $0$ à $\frac\pi2$.




\item \'Etude des points singuliers\\
Le point $M(0)$, qui est l'origine, est singulier. Pour étudier l'existence d'une tangente en ce point, on considère
$$\frac{y(t)-y(0)}{x(t)-x(0)}=\frac{1-\cos t}{t-\sin t}\underset{0}{\sim}\frac{t^2/2}{t^3/6}$$
et donc $\frac{y(t)-y(0)}{x(t)-x(0)}\xrightarrow[t\to 0^+]{}+\infty$. Par conséquent, la courbe possède une tangente de pente verticale au point $M(0)$.
\end{itemize}

\begin{center}
\begin{tikzpicture}[scale=0.7]
     \draw[->,>=latex,thick, gray] (-6.5,0)--(10,0) node[below,black] {$x$};
     \draw[->,>=latex,thick, gray] (0,-.5)--(0,2.5) node[right,black] {$y$};  

  \draw[red, very thick,domain=0:pi,samples=100] plot ({\x - sin(\x r)},{1-cos(\x r)});
  \draw[gray, very thick,domain=pi:3*pi,samples=100] plot ({\x - sin(\x r)},{1-cos(\x r)});
  \draw[gray, very thick,domain=-2*pi:0,samples=100] plot ({\x - sin(\x r)},{1-cos(\x r)});

  \fill[blue] (0,0) circle (2pt)  node[below right]{$M(0)$}; 
  \fill[blue] (3.14,2) circle (2pt)  node[above]{$M(\pi)$}; 
  \fill[black] (6.28,0) circle (2pt)  node[below]{$M(2\pi)$};
  \fill[blue] (-3.14,2) circle (2pt)  node[above]{$M(-\pi)$}; 
  \draw[very thick, blue] (0,0)--+(0,1)--+(0,-1);

\end{tikzpicture}  
\end{center}

\end{enumerate}
\fincorrection
\correction{006984}\ 
\begin{enumerate}
\item Soit $t>0$:
$$\left\{\begin{array}{l}
x(\frac{1}{t})=\frac{1}{t}\ln(\frac{1}{t})=-y(t)\\
\ \\
y(\frac{1}{t})=t\ln(\frac{1}{t})=-x(t)
\end{array}\right.$$
et par conséquent, le point $M(\frac{1}{t})$ est le symétrique 
de $M(t)$ par rapport à la droite d'équation $y=-x$. 
\fbox{On restreint l'étude à l'intervalle $]0;1]$}, 
puis on obtiendra l'intégralité de la courbe par symétrie 
par rapport à la seconde bissectrice.
\item Les fonctions $x$ et $y$ sont de classe $\mathcal{C}^1$ sur $]0;1]$.
\begin{itemize}
\item Tableau de variations conjointes\\
Pour $t\in]0;1]$: 
$$\begin{array}{lcl}
x(t)=t\ln t&\ &y(t)=\frac{\ln t}{t}\\
x'(t)=1+\ln t&\ &y'(t)=\frac{1-\ln t}{t^2}\\
x'(t)>0\Longleftrightarrow t>1/e&\ &y'(t)>0 \\
x'(t)=0\Longleftrightarrow t=1/e &\ &y'(t)\not=0
\end{array}$$
puisque $\frac{1}{e}<1<e$. On obtient donc le tableau suivant:
$$\begin{array}{c|lcccr}
t&0&\ &1/e&\ &1\\\hline
x'(t)&-\infty &-&0&+&1 \\\hline
\ &0  &\ & & & 0\\
x&\ &\searrow &\ &\nearrow & \\
\ & & &-1/e &\ & \\\hline
\ & &\ & & &0 \\
&\ & &\ &\nearrow& \\
&\ & &-e & & \\
y&\ &\nearrow &\ &\ & \\
\ &-\infty & & &\ & \\\hline
y'(t)&+\infty & + &2e^2 & +&1 \\
\end{array}$$
Il n'y a pas de point singulier.
\item \'Etude des branches infinies\\
Comme $x(t)\xrightarrow[t\to 0^+]{}0$ et $y(t)\xrightarrow[t\to 0^+]{}-\infty$, l'axe des ordonnées est asymptote à $\mathcal{C}$.
\end{itemize}
\end{enumerate}

\begin{center}
\begin{tikzpicture}[scale=1]
     \draw[->,>=latex,thick, gray] (-1,0)--(6,0) node[below,black] {$x$};
     \draw[->,>=latex,thick, gray] (0,-6)--(0,1.5) node[right,black] {$y$};  


  \draw[red, very thick,domain=0.25:1,samples=100] plot ({\x *ln(\x)},{ln(\x)/\x});
  \draw[gray, very thick,domain=1:4,samples=100] plot ({\x *ln(\x)},{ln(\x)/\x});

  \fill[black] (0,0) circle (2pt) node[above right]{$0$}; 
  \fill[black] (1,0) circle (2pt)  node[below]{$1$}; 
 % \fill[blue] (0,1) circle (2pt)  node[right]{$M(0)$}; 

  \fill[blue] (-0.37-0.02,-2.72) circle (1pt)  node[left]{$M(1/e)$}; 

  \draw[very thick, blue] (-0.37-0.02,-2.72)--+(0,1)--+(0,-1);

  \draw[very thick, gray, dashed] (-1,1)--+(5,-5);

\end{tikzpicture}
\end{center}

\fincorrection
\correction{006985}
Les expressions $x(t)=\frac{1}{t^2-t}$ et $y(t)=\frac{t}{t^2-1}$ 
sont bien définies, et de classe $\mathcal{C}^1$ en dehors de $t=0$ et $t=\pm1$. Le domaine de définition est donc

$$\mathcal{D}=\,]-\infty;-1[\ \cup\ ]-1;0[\ \cup\ ]0;1[\ \cup\ ]1;+\infty[$$
La courbe possède un point double si elle se recoupe: on cherche donc deux paramètres $t_1,t_2\in\mathcal{D}$ tels que $t_1\not=t_2$ et $M=M(t_1)=M(t_2)$, \textsl{i.e.}
\begin{eqnarray*}
\left\{\begin{array}{l}
\frac{1}{t_1^2-t_1}=\frac{1}{t_2^2-t_2}\\
\ \\
\frac{t_1}{t_1^2-1}=\frac{t_2}{t_2^2-1}
\end{array}\right.
&\Longleftrightarrow&\left\{\begin{array}{l}
t_1^2-t_1=t_2^2-t_2\\
t_1(t_2^2-1)-t_2(t_1^2-1)=0
\end{array}\right.\\
 &\Longleftrightarrow&\left\{\begin{array}{l}
(t_1-t_2)(t_1+t_2-1)=0\\
(t_2-t_1)(t_1t_2+1)=0
\end{array}\right.
\end{eqnarray*}
Comme on cherche $t_1\not=t_2$, le système obtenu est équivalent à 
$\displaystyle\left\{\begin{array}{l}t_1+t_2=1\\t_1t_2=-1\end{array}\right.$, 
autrement dit à un système du type somme-produit: 
cela signifie que $t_1$ et $t_2$ doivent \^etre les deux racines (distinctes) de $X^2-X-1$, 
c'est-à-dire $\frac{1\pm\sqrt{5}}{2}$ (qui sont bien dans $\mathcal{D}$). 
On a donc un seul point double, c'est
$$M\left(\frac{1+\sqrt{5}}{2}\right)=M\left(\frac{1-\sqrt{5}}{2}\right)$$
de coordonnées cartésiennes $(1,1)$. Pour déterminer les tangentes en ce point, on calcule le vecteur dérivé:
$$\vec{V}(t)=\left\{\begin{array}{l}
x'(t)=\frac{1-2t}{(t^2-t)^2}\\
\ \\
y'(t)=\frac{-1-t^2}{(t^2-1)^2}
\end{array}\right.$$
En remplaçant, on obtient 
$$\vec{V}(t_1)=\begin{pmatrix}x'(t_1)\\y'(t_1)\end{pmatrix}
=\begin{pmatrix}\sqrt{5}\\ \frac{-5-\sqrt{5}}{2}\end{pmatrix}\quad\text{ et }
\quad \vec{V}(t_2)=\begin{pmatrix}x'(t_2)\\y'(t_2)\end{pmatrix}
=\begin{pmatrix}-\sqrt{5}\\ \frac{-5+\sqrt{5}}{2}\end{pmatrix}$$
Les deux tangentes à la courbe au point de coordonnées $(1,1)$ sont donc dirigées respectivement par les vecteurs $\vec{V}(t_1)$ et $\vec{V}(t_2)$, dont on vérifie en faisant le produit scalaire $\vec{V}_1\cdot\vec{V}_2=x'(t_1)x'(t_2)+y'(t_1)y'(t_2)=0$ qu'ils sont orthogonaux.

\begin{center}
\begin{tikzpicture}[scale=1]
     \draw[->,>=latex,thick, gray] (-6.5,0)--(5,0) node[below,black] {$x$};
     \draw[->,>=latex,thick, gray] (0,-3.5)--(0,5) node[right,black] {$y$};  


  \draw[red, very thick,domain=0.2:0.8,samples=100] plot ({1/(\x*\x-\x)},{\x/(\x*\x-1)});
  \draw[red, very thick,domain=-0.2:-0.9,samples=100] plot ({1/(\x*\x-\x)},{\x/(\x*\x-1)});
  \draw[red, very thick,domain=1.2:10,samples=100] plot ({1/(\x*\x-\x)},{\x/(\x*\x-1)});
  \draw[red, very thick,domain=-1.2:-10,samples=100] plot ({1/(\x*\x-\x)},{\x/(\x*\x-1)});

  \fill[black] (0,0) circle (2pt) node[below left]{$0$}; 
  \fill[black] (1,0) circle (2pt)  node[below]{$1$}; 
  \fill[black] (0,1) circle (2pt)  node[left]{$1$}; 
  \draw[black, dashed] (0,1)--(1,1)--(1,0); 


 \coordinate (P) at (1,1);

  \fill[blue] (P) circle (2.5pt)  node[above=5pt]{$M$}; 

  \draw[very thick, blue] (P)--+(32:2)--+(32:-2);
  \draw[very thick, blue] (P)--+(122:2)--+(122:-2);

\end{tikzpicture}  
\end{center}

\fincorrection
\correction{006986}
Les fonctions $x$ et $y$ sont de classe $\mathcal{C}^1$ sur $\R$. 
Un point $M(t)$ de la courbe est singulier si $x'(t)=y'(t)=0$, or
$$\left\{\begin{array}{l}
x'(t)=\frac{4(t^2+1)-2t(4t-3)}{(t^2+1)^2}=\frac{-2(2t^2-3t-2)}{(t^2+1)^2}\\
\ \\
y'(t)=\frac{2(t^2+2)-2t(2t-1)}{(t^2+2)^2}=\frac{-2(t^2-t-2)}{(t^2+2)^2}
\end{array}\right.$$
Ainsi $M(t)$ est singulier si et seulement si 
$\left\{\begin{array}{l}2t^2-3t-2=0\\t^2-t-2=0\end{array}\right.$. 
Ce système admet une unique solution $t=2$, correspondant au point $M(2)$ de coordonnées $(1,\frac{1}{2})$.

Le vecteur dérivé est nul au point $M(2)$ ; 
pour obtenir l'allure de la courbe au voisinage de ce point,
il faut donc effectuer un développement limité à un 
ordre assez grand pour trouver deux termes non constants non nuls. 
Ici l'ordre 3 suffira, on pose $t=2+h$ pour simplifier 
(ainsi ``$t$ proche de 2'' devient ``$h$ proche de 0''):
\begin{eqnarray*}
x(2+h)&=&\frac{4(2+h)-3}{(2+h)^2+1}=\frac{5+4h}{5+4h+h^2}=1-\frac{1}{5}h^2\cdot\frac{1}{1+\frac{4h+h^2}{5}}\\
 &=&1-\frac{1}{5}h^2\cdot\left(1-\left(\frac{4h+h^2}{5}\right)+o\left(\frac{4h+h^2}{5}\right)\right)\\
 &=&1-\frac{1}{5}h^2\cdot\left(1-\frac{4}{5}h+o(h)\right)\\
 &=&1-\frac{1}{5}h^2+\frac{4}{25}h^3+o(h^3)
\end{eqnarray*}
\begin{eqnarray*}
y(2+h)&=&\frac{2(2+h)-1}{(2+h)^2+2}=\frac{3+2h}{6+4h+h^2}=\frac{1}{2}-\frac{1}{12}h^2\cdot\frac{1}{1+\frac{4h+h^2}{6}}\\
 &=&\frac{1}{2}-\frac{1}{12}h^2\cdot\left(1-\left(\frac{4h+h^2}{6}\right)+o\left(\frac{4h+h^2}{6}\right)\right)\\
 &=&\frac{1}{2}-\frac{1}{12}h^2\cdot\left(1-\frac{2}{3}h+o(h)\right)\\
 &=&\frac{1}{2}-\frac{1}{12}h^2+\frac{1}{18}h^3+o(h^3)
\end{eqnarray*}
On a donc le développement limité vectoriel suivant:
$$M(2+h)=\begin{pmatrix}1\\ \frac{1}{2}\end{pmatrix}
+\begin{pmatrix}-\frac{1}{5}\\-\frac{1}{12}\end{pmatrix}\cdot h^2+
\begin{pmatrix}\frac{4}{25}\\  \frac{1}{18}\end{pmatrix}\cdot h^3+o(h^3)$$
On vérifie que le terme constant du développement limité correspond bien à 
$\begin{pmatrix}x(2)\\ y(2)\end{pmatrix}$ et que le terme linéaire, 
qui vaut $\begin{pmatrix}x'(2)\\ y'(2)\end{pmatrix}\cdot h$, est nul. 
Les coefficients de $h^2$ et $h^3$ sont des vecteurs non nuls, $M(2)$ 
est donc un point de rebroussement de première espèce ($p=2$, $q=3$).
La tangente est dirigée par le premier vecteur 
non nul, coefficient de $h^k$ (avec $k\ge1$), donc ici
le coefficient de $h^2$ ; ainsi la tangente en $M(2)$ est dirigée par 
$\left(\begin{smallmatrix}-\frac{1}{5}\\ -\frac{1}{12}\end{smallmatrix}\right)$.


\begin{center}
\begin{tikzpicture}[scale=1.5]
     \draw[->,>=latex,thick, gray] (-4,0)--(3.5,0) node[below,black] {$x$};
     \draw[->,>=latex,thick, gray] (0,-1.5)--(0,2) node[right,black] {$y$};  

  \coordinate (P) at (1,1/2);
  \fill[blue] (P) circle (2.5pt)  node[above=5pt]{$M(2)$}; 
  \draw[very thick, blue] (P)--+(22.6:2)--+(22.6:-2);

  \draw[red, very thick,domain=-10:10,samples=200] plot ({(4*\x-3)/(\x*\x+1)},{(2*\x-1)/(\x*\x+2)});

  \draw[red, very thick] (0,0)--+(27:0.45)--+(26:-0.5);


  \fill[black] (0,0) circle (2pt) node[below left]{$0$}; 
  \fill[black] (1,0) circle (2pt)  node[below]{$1$}; 
  \fill[black] (0,0.5) circle (2pt)  node[left]{$\frac12$}; 
  \draw[black, dashed] (0,0.5)--(1,0.5)--(1,0); 

\end{tikzpicture}  
\end{center}

\fincorrection
\correction{006987} Les expressions $x(t)=t+\frac{4}{t}$ et $y(t)=\frac{t}{3}+2+\frac{3}{t+1}$ sont bien définies pour $t\in\mathcal{D}=\R\setminus\{-1;0\}$.
\begin{enumerate}
\item Les fonctions $x$ et $y$ sont de classe $\mathcal{C}^1$ sur $\mathcal{D}$. Soit $t\in\mathcal{D}$:
$$\begin{array}{lcl}
x(t)=t+\frac{4}{t}&\ &y(t)=\frac{t}{3}+2+\frac{3}{t+1}\\
x'(t)=1-\frac{4}{t^2}&\ &y'(t)=\frac{1}{3}-\frac{3}{(t+1)^2}\\
x'(t)>0\Longleftrightarrow |t|>2&\ &y'(t)>0\Longleftrightarrow |t+1|>3\Longleftrightarrow \left\{\begin{array}{l}t>2\\ \text{ou}\\ t<-4\end{array}\right.
\\
x'(t)=0\Longleftrightarrow t\in\{-2;2\} &\ &y'(t)=0\Longleftrightarrow t\in\{-4;2\}
\end{array}$$
$$\begin{array}{c|lcccccccccccr}
t&-\infty&\ &-4& &-2& &-1& &0& &2& &+\infty\\\hline
x'(t)& &+& &+&0&-&&-&&-&0&+&\\\hline
\ &  &\ &&&-4&&&&+\infty&&&&+\infty\\
\ &  &\ &&\nearrow&&\searrow&&&&&&&\\
x&\ & &-5&&&&-5&&&\searrow&&\nearrow&\ \\
\ &  &\nearrow &&&&&&\searrow&&&&&\\
\ &-\infty &\ & &&&&&&-\infty&&4&&\\\hline
\ & &\ & -\frac{1}{3}&&&&+\infty&&&&&&+\infty\\
\ & &\ & &\searrow&&&&\searrow&&&&&\\
y&\ &\nearrow & &&-\frac{5}{3}&&&&5&&&\nearrow&\\
\ & &\ & &&&\searrow&&&&\searrow&&&\\
\ &-\infty & & &&&&-\infty&&&&\frac{11}{3}&&\\\hline
y'(t)&& + &0 &-&&-&&-&&-&0&+&\\
\end{array}$$

\item Le tableau de variations conjointes indique:
\begin{itemize}
\item $t=-4$: tangente horizontale, au point de coordonnées $(-5,-\frac{1}{3})$;
\item $t=-2$: tangente verticale, au point de coordonnées $(-4,-\frac{5}{3})$;
\item $t=-1$: une asymptote verticale, d'équation $x=-5$;
\item $t=0$: une asymptote horizontale, d'équation $y=5$;
\item $t=2$: il y a un point singulier en $(4,\frac{11}{3})$ (voir après).
\end{itemize}
Il reste à étudier le comportement quand $t\to\pm\infty$:
$$\frac{y(t)}{x(t)}=\frac{t^3+16t^2+6t}{3(t^3+t^2+4t+4)}\xrightarrow[t\to\pm\infty]{}\frac{1}{3}$$
puis $y(t)-\frac{1}{3}x(t)\xrightarrow[t\to\pm\infty]{}2$. La courbe a donc pour asymptote, quand $t\to-\infty$ et quand $t\to+\infty$, la m\^eme droite d'équation $y=\frac{1}{3}x+2$.


\item On constate sur le tableau de variation qu'il n'y a qu'un seul point singulier, correspondant au paramètre $t=2$. Pour conna\^itre l'allure de la courbe au voisinage du point $M(2)$, on fait un développement limité de $x$ et $y$ au voisinage de $t=2$. Comme ici $x$ et $y$ sont de classe $\mathcal{C}^\infty$ et d'expressions assez simples, on peut directement appliquer la formule de Taylor-Young:
$$\left\{\begin{array}{l}
x(t)=x(2)+x'(2)\cdot(t-2)+\frac{1}{2}x''(2)\cdot(t-2)^2+\frac{1}{6}x'''(2)\cdot(t-2)^3+o((t-2)^3)\\
y(t)=y(2)+y'(2)\cdot(t-2)+\frac{1}{2}y''(2)\cdot(t-2)^2+\frac{1}{6}y'''(2)\cdot(t-2)^3+o((t-2)^3)
\end{array}\right.$$ 
On sait déjà que $x(2)=4$, $y(2)=\frac{11}{3}$ et $x'(2)=y'(2)=0$. De plus
$x''(t)=\frac{8}{t^3}$, $x'''(t)=\frac{-24}{t^4}$ et $y''(t)=\frac{6}{(t+1)^3}$, $y'''(t)=\frac{-18}{(t+1)^4}$, ce qui donne
$$M(t)=\begin{pmatrix}4 \\ \frac{11}{3}\end{pmatrix}+
\begin{pmatrix}\frac{1}{2}\\\frac{1}{9}\end{pmatrix}\cdot (t-2)^2+
\begin{pmatrix}-\frac{1}{4}\\ \frac{-1}{27}\end{pmatrix}\cdot (t-2)^3+o((t-2)^3)$$
C'est un point de rebroussement de première espèce. L'équation (sous forme paramétrée) de la tangente $T_{M(2)}$ s'obtient en tronquant le développement limité:
$$\begin{pmatrix}x\\ y\end{pmatrix}\in T_{M(2)}
\Longleftrightarrow\exists\lambda\in\R,\ \begin{pmatrix}x \\ y\end{pmatrix}
=\begin{pmatrix}4\\  \frac{11}{3}\end{pmatrix}+
\begin{pmatrix}\frac{1}{2}\\\frac{1}{9}\end{pmatrix}\cdot \lambda$$
En éliminant le paramètre $\lambda$, on récupère une équation cartésienne 
$$T_{M(2)}:\ y=\frac{11}{3}+\frac{1}{9}\cdot 2(x-4)=\frac{2}{9}x+\frac{25}{9}$$


\end{enumerate}


\begin{center}
\begin{tikzpicture}[scale=0.35]
     \draw[->,>=latex,thick, gray] (-20,0)--(21,0) node[below,black] {$x$};
     \draw[->,>=latex,thick, gray] (0,-14)--(0,18) node[right,black] {$y$};  


  \coordinate (P) at (4,11/3);
  \fill[blue] (P) circle (5pt)  node[above=5pt]{$M(2)$}; 
  \draw[very thick, blue] (P)--+(12.5:5)--+(12.5:-4);



  \fill[black] (0,0) circle (4pt) node[below left]{$0$}; 
  \fill[black] (1,0) circle (4pt)  node[below]{$1$}; 
  \fill[black] (0,1) circle (4pt)  node[left]{$1$}; 

  \draw[thick, green!80!black] (0,2)--+(18.4:21)--+(18.4:-21);
 \node[green!80!black]  at (18,7) {$y=\frac13x+2$};

  \draw[thick, green!80!black] (0,5)--+(20.5,0)--+(-20.5,0);
 \node[green!80!black]  at (-18,4) {$y=5$};

  \draw[thick, green!80!black] (-5,0)--+(0,17)--+(0,-14.5);
 \node[green!80!black]  at (-4,14) {$x=-5$};



  \coordinate (M2) at (-4,-5/3);
  \fill[blue] (M2) circle (5pt)  node[right=5pt]{$M(-2)$}; 
  \draw[very thick, blue] (M2)--+(0,3)--+(0,-3);

  \coordinate (M4) at (-5,-1/3);
  \fill[blue] (M4) circle (5pt)  node[below left=5pt]{$M(-4)$}; 
  \draw[very thick, blue] (M4)--+(3,0)--+(-3,0);

  \draw[red, very thick,domain=0.2:20,samples=100] plot ({\x+4/\x},{\x/3+2+3/(\x+1)});
  \draw[red, very thick,domain=-0.8:-0.2,samples=100] plot ({\x+4/\x},{\x/3+2+3/(\x+1)});
  \draw[red, very thick,domain=-1.2:-20,samples=100] plot ({\x+4/\x},{\x/3+2+3/(\x+1)});

\end{tikzpicture}  
\end{center}


\fincorrection
\correction{006988}

\begin{enumerate}
  \item Commençons par trouver le vecteur tangent à la courbe au point $M(t)$. 
  Les fonctions $x$ et $y$ sont de classe $\mathcal{C}^1$ sur $\R$, et 
  $x'(t)=6t$, $y'(t)=12t^2$. Si $t\not=0$, le vecteur tangent à la courbe 
  au point $M(t)$ est donc le vecteur dérivé 
  $\begin{pmatrix}6\\ 12t\end{pmatrix}$. Si $t=0$, 
  le vecteur dérivé est nul et il faut dériver encore une fois pour 
  obtenir un vecteur non nul $\begin{pmatrix}x''(0)\\ y''(0)\end{pmatrix}
  =\begin{pmatrix}6\\ 0\end{pmatrix}$, qui est donc un vecteur 
  directeur de la tangente au point $M(0)$. Finalement, pour tout $t$, 
  la tangente au point $M(t)$ est dirigée par le vecteur 
$$\vec{V}=\begin{pmatrix}6\\ 12t\end{pmatrix}$$  

  \item 
\begin{itemize}
  \item Une droite $D$ est tangente à $\mathcal{C}$ s'il existe $t\in\R$ tel 
que $M(t)\in D$ et $\vec{V}(t)$ soit un vecteur directeur de $D$.
  
  \item Une droite $D$ est orthogonale à $\mathcal{C}$ s'il existe $t'\in\R$ tel que $M(t')\in D$ 
et $\vec{V}(t')$ soit un vecteur orthogonal à la droite $D$. 
\end{itemize}


  \item On cherche donc à quelle condition $\vec{V}(t)$ et $\vec{V}(t')$ sont orthogonaux:
$$\vec{V}(t)\cdot\vec{V}(t')=0\Longleftrightarrow 36+144tt'=0\Longleftrightarrow tt'=-\frac{1}{4}$$
ce qui exclut le paramètre $t=0$. 

  \item Soit donc $t\not=0$, et $T(t)$ la tangente à $\mathcal{C}$ en $M(t)$:
$$T(t)=\big\{M(t)+\lambda\vec{V}(t)\ |\ \lambda\in\R\big\}=\left\{\begin{pmatrix}x\\ y\end{pmatrix}\ |\ \exists\lambda\in\R,\ \begin{array}{|l}x=3t^2+6\lambda\\ y=4t^3+12\lambda t\end{array}\right\}$$
En éliminant $\lambda$, on trouve que $T(t)$ a pour équation cartésienne 
$$y=4t^3+2t(x-3t^2)=2tx-2t^3$$

  \item On sait déjà que $T(t)$ et $T(-\frac{1}{4t})$ sont perpendiculaires. Il reste à voir si $T(t)$ coupe bien $\mathcal{C}$ au point $M(-\frac{1}{4t})$:
\begin{eqnarray*}
M\left(-\frac{1}{4t}\right)\in T(t)&\Longleftrightarrow& 4\left(-\frac{1}{4t}\right)^3=2t\cdot 3\left(-\frac{1}{4t}\right)^2-2t^3\\
 &\Longleftrightarrow& 32t^6-6t^2-1=0\\
 &\Longleftrightarrow& X=t^2\quad\mathrm{et}\quad X^3-\frac{3}{16}X-\frac{1}{32}=0
\end{eqnarray*}
L'étude des variations du polyn\^ome $X^3-\frac{3}{16}X-\frac{1}{32}$ montre qu'il admet $-\frac{1}{4}$ comme racine (double), il se factorise donc sous la forme $X^3-\frac{3}{16}X-\frac{1}{32}=(X+\frac{1}{4})^2(X-\frac{1}{2})$ et sa seule racine positive est $\frac{1}{2}$:
\begin{eqnarray*}
M\left(-\frac{1}{4t}\right)\in T(t)&\Longleftrightarrow& X=t^2\quad\mathrm{et}\quad X\in\left\{-\frac{1}{4};\frac{1}{2}\right\}\\
 &\Longleftrightarrow& t^2=\frac{1}{2}\Longleftrightarrow t=\pm\frac{\sqrt{2}}{2}
\end{eqnarray*}

  \item 
Ainsi $T(\frac{\sqrt{2}}{2})$ et $T(-\frac{\sqrt{2}}{2})$ sont les seules 
droites à la fois tangentes et orthogonales à $\mathcal{C}$ :

\begin{itemize}
  \item La droite 
$$T(\frac{\sqrt{2}}{2}):\ y=\sqrt{2}\,x-\frac{\sqrt{2}}{2}$$
est tangente au point $M(\frac{\sqrt{2}}{2})$ et orthogonale au point $M(-1/(4\frac{\sqrt{2}}{2}))$.
  
  \item La droite 
$$T(-\frac{\sqrt{2}}{2}):\ y=-\sqrt{2}\,x+\frac{\sqrt{2}}{2}$$
est tangente au point $M(-\frac{\sqrt{2}}{2})$ et orthogonale au point $M(1/(4\frac{\sqrt{2}}{2}))$.
\end{itemize}

\end{enumerate}

\begin{center}
\begin{tikzpicture}[scale=0.80]
     \draw[->,>=latex,thick, gray] (-2,0)--(5,0) node[below,black] {$x$};
     \draw[->,>=latex,thick, gray] (0,-4)--(0,4) node[right,black] {$y$};  


  \fill[black] (0,0) circle (2pt) node[below left]{$0$}; 
  \fill[black] (1,0) circle (2pt)  node[above]{$1$}; 
  \fill[black] (0,1) circle (2pt)  node[left]{$1$}; 


  \coordinate (M1) at (3/2,1.41);
  \fill[blue] (M1) circle (2.5pt)  node[right]{$M(+\frac{\sqrt{2}}{2})$}; 
  \draw[very thick, blue] (M1)--+(54.7:3)--+(54.7:-3);

  \coordinate (MM1) at (0.36,-0.18);
  \fill[blue] (MM1) circle (2.5pt)  node[right=5pt]{$M(-1/(4\frac{\sqrt{2}}{2}))$}; 
  \draw[very thick, blue] (MM1)--+(144.7:3)--+(144.7:-3);

%   \coordinate (M2) at (3/2,-1.41);
%   \fill[green!80!black] (M2) circle (2.5pt)  node[below left=5pt]{$M(-\frac{\sqrt{2}}{2})$}; 
%   \draw[very thick, green!80!black] (M2)--+(-54.7:3)--+(-54.7:-3);
% 
%   \coordinate (MM2) at (0.36,0.18);
%   \fill[green!80!black] (MM2) circle (2.5pt)  node[above]{$M(-\frac{\sqrt{2}}{2})$}; 
%   \draw[very thick, green!80!black] (MM2)--+(-144.7:3)--+(-144.7:-3);

  \draw[red, very thick,domain=-1:1,samples=100] plot ({3*\x*\x},{4*\x*\x*\x});

\end{tikzpicture}  \quad
\begin{tikzpicture}[scale=0.80]
     \draw[->,>=latex,thick, gray] (-2,0)--(5,0) node[below,black] {$x$};
     \draw[->,>=latex,thick, gray] (0,-4)--(0,4) node[right,black] {$y$};  


  \fill[black] (0,0) circle (2pt) node[below left]{$0$}; 
  \fill[black] (1,0) circle (2pt)  node[below]{$1$}; 
  \fill[black] (0,1) circle (2pt)  node[left]{$1$}; 


%   \coordinate (M1) at (3/2,1.41);
%   \fill[blue] (M1) circle (2.5pt)  node[right]{$M(+\frac{\sqrt{2}}{2})$}; 
%   \draw[very thick, blue] (M1)--+(54.7:3)--+(54.7:-3);
% 
%   \coordinate (MM1) at (0.36,-0.18);
%   \fill[blue] (MM1) circle (2.5pt)  node[right=5pt]{$M(-1/4\frac{\sqrt{2}}{2})$}; 
%   \draw[very thick, blue] (MM1)--+(144.7:3)--+(144.7:-3);

  \coordinate (M2) at (3/2,-1.41);
  \fill[green!80!black] (M2) circle (2.5pt)  node[right]{$M(-\frac{\sqrt{2}}{2})$}; 
  \draw[very thick, green!80!black] (M2)--+(-54.7:3)--+(-54.7:-3);

  \coordinate (MM2) at (0.36,0.18);
  \fill[green!80!black] (MM2) circle (2.5pt)  node[right=5pt]{$M(1/(4\frac{\sqrt{2}}{2}))$}; 
  \draw[very thick, green!80!black] (MM2)--+(-144.7:3)--+(-144.7:-3);

  \draw[red, very thick,domain=-1:1,samples=100] plot ({3*\x*\x},{4*\x*\x*\x});

\end{tikzpicture} 
\end{center}


\fincorrection
\nocorrection
\nocorrection
\nocorrection
\nocorrection
\nocorrection
\nocorrection
\nocorrection
\nocorrection
\nocorrection
\nocorrection
\nocorrection
\correction{004993}
 \begin{enumerate}
\item 
$$
\centerline{\includegraphics[height=4cm]{../images/img004993-1}}
$$
\item 
$$
\centerline{\includegraphics[height=4cm]{../images/img004993-2}}
$$
aire de la boucle : $2-\dfrac\pi2$                            \par
\item 
$$
\centerline{\includegraphics[height=4cm]{../images/img004993-3}}
$$
asymptotes : $y=\pm x\sqrt3-1$                                \par
la courbe traverse ses asymptotes au point de concours        \par
\item 
$$
\centerline{\includegraphics[height=4cm]{../images/img004993-4}}
$$
asymptotes : $x\pm y\sqrt3=\dfrac43$                           \par
\item 
$$
\centerline{\includegraphics[height=4cm]{../images/img004993-5}}
$$
\item 
$$
\centerline{\includegraphics[height=4cm]{../images/img004993-6}}
$$
asymptotes : $3x\pm y\sqrt3=-1$                               \par
\item 
$$
\centerline{\includegraphics[height=4cm]{../images/img004993-7}}
$$
\item 
$$
\centerline{\includegraphics[height=4cm]{../images/img004993-8}}
$$
\item 
$$
\centerline{\includegraphics[height=4cm]{../images/img004993-9}}
$$
\item 
$$
\centerline{\includegraphics[height=4cm]{../images/img004993-10}}
$$
\item   
$$
\centerline{\includegraphics[height=4cm]{../images/img004993-11}}
$$
 \end{enumerate}
\fincorrection
\correction{005205}
Notons $\mathcal{E}$ l'ensemble cherché.

Tout d'abord, pour tout réel $\theta$, $1+\sin(2\theta)\geq0$, $1-\sin(2\theta)\geq0$ puis $\sqrt{1+\sin(2\theta)}+\sqrt{1-\sin(2\theta)}>0$, car $\sin(2\theta)$ ne peut valoir simultanément $1$ et $-1$. La fonction $r\mapsto r(\theta)$ est donc définie sur $\Rr$, clairement $2\pi$-périodique.

Ainsi, 

$$M(\theta+2\pi)=[r(\theta+2\pi),\theta+2\pi]=[r(\theta),\theta+2\pi]=M(\theta).$$

On obtient donc l'ensemble complet quand $\theta$ décrit un intervalle de longueur $2\pi$ comme $[-\pi,\pi]$ par exemple.

La fonction $r\mapsto r(\theta)$ est plus paire. Par suite,

$$M(-\theta)=[r(-\theta),-\theta]=[r(\theta),-\theta]=s_{(Ox)}(M(\theta)).$$

On construit l'ensemble des points correspondant à $\theta\in[0,\pi]$ et on obtient l'ensemble complet par symétrie orthogonale d'axe $(Ox)$.

Pour $\theta\in[0,\pi]$, on a clairement $r(\pi-\theta)=r(\theta)$. Par suite,

$$M(\pi-\theta)=[r(\pi-\theta),\pi-\theta]=[r(\theta),\pi-\theta]=s_{(Oy)}(M(\theta)).$$

On construit l'ensemble des points correspondant à $\theta\in[0,\frac{\pi}{2}]$ et on obtient l'ensemble complet par symétrie orthogonale d'axe $(Oy)$ puis par symétrie orthogonale d'axe $(Ox)$.

Pour $\theta\in[0,\frac{\pi}{2}]$, on a clairement $r(\frac{\pi}{2}-\theta)=r(\theta)$. Par suite, en notant $(\Delta)$ la droite d'équation $y=x$,

$$M(\frac{\pi}{2}-\theta)=[r(\frac{\pi}{2}-\theta),\frac{\pi}{2}-\theta]=[r(\theta),\frac{\pi}{2}-\theta]
=s_{(\Delta)}(M(\theta)).$$

On construit l'ensemble des points correspondant à $\theta\in[0,\frac{\pi}{4}]$ et on obtient l'ensemble complet par symétrie orthogonale d'axe $(\Delta)$ puis par symétrie orthogonale d'axe $(Oy)$ et enfin par symétrie orthogonale d'axe $(Ox)$.

Maintenant, pour $\theta\in[0,\frac{\pi}{4}]$,

\begin{align*}
\frac{1}{\sqrt{1+\sin(2\theta)}+\sqrt{1-\sin(2\theta)}}&=
\frac{1}{\sqrt{1+\cos(\frac{\pi}{2}-2\theta)}+\sqrt{1-\cos(\frac{\pi}{2}-2\theta)}}\\
 &=\frac{1}{\sqrt{2\cos^2(\frac{\pi}{4}-\theta)}+\sqrt{2\sin^2(\frac{\pi}{4}-\theta)}}
 =\frac{1}{\sqrt{2}\cos(\frac{\pi}{4}-\theta)+\sqrt{2}\sin(\frac{\pi}{4}-\theta)}\\
 &=\frac{1}{2\cos(\frac{\pi}{4}-(\frac{\pi}{4}-\theta))}=\frac{1}{2\cos\theta}.
\end{align*}

En notant $x$ et $y$ les coordonnées d'un point $M$, on a alors 

$$M\in\mathcal{E}\Leftrightarrow r=\frac{1}{2\cos\theta}\Leftrightarrow r\cos(\theta)=\frac{1}{2}\Leftrightarrow x=\frac{1}{2}.$$

D'où le graphique~:

$$\includegraphics{../images/img005205-1}$$

\fincorrection
\correction{005530}
\begin{enumerate}
 \item  (\textbf{Lemniscate de } \textsc{Bernoulli}.) Soit $\mathcal{C}$ la courbe d'équation polaire $r=\sqrt{\cos(2\theta)}$.
\textbf{Domaine d'étude.} Notons $D$ le domaine de définition de la fonction $r~:~\theta\mapsto\sqrt{\cos(2\theta)}$.
\textbullet~$\theta\in D\Leftrightarrow\theta+2\pi\in D$ et pour $\theta\in D$, 

\begin{center}
$M(\theta+2\pi)=\left[r(\theta+2\pi),\theta+2\pi\right]=\left[r(\theta),\theta+2\pi\right]=\left[r(\theta),\theta\right]=M(\theta)$.
\end{center}
On obtient donc la courbe complète quand $\theta$ décrit un intervalle de longueur $2\pi$ comme $[-\pi,\pi]$.
\textbullet~$\theta\in D\Leftrightarrow-\theta\in D$ et pour $\theta\in D$, 

\begin{center}
$M(-\theta)=\left[r(-\theta),-\theta\right]=\left[r(\theta),-\theta\right]=s_{(Ox)}(M(\theta))$.
\end{center} On étudie et on construit la portion de courbe correspondant à $\theta\in[0,\pi]$ puis on obtient la courbe complète par réflexion d'axe $(Ox)$.
\textbullet~$\theta\in D\Leftrightarrow\pi-\theta\in D$ et pour $\theta\in D$, 

\begin{center}
$M(\pi-\theta)=\left[r(\pi-\theta),\pi-\theta\right]=\left[r(\theta),\pi-\theta\right]=s_{(Oy)}(M(\theta))$.
\end{center} On étudie et on construit la portion de courbe correspondant à $\theta\in\left[0,\frac{\pi}{2}\right]$ puis on obtient la courbe complète par réflexion d'axe $(Oy)$ puis d'axe $(Ox)$.
Pour $\theta\in\left[0,\frac{\pi}{2}\right]$, $\theta\in D\Leftrightarrow\cos\left(2\theta\right)\geqslant0\Leftrightarrow\theta\in\left[0,\frac{\pi}{4}\right]$. On étudie donc la courbe sur $\left[0,\frac{\pi}{4}\right]$.
\textbf{Variations et signe de $r$.} La fonction $r$ est strictement décroissante sur $\left[0,\frac{\pi}{4}\right]$, strictement positive sur $\left]0,\frac{\pi}{4}\right]$ et s'annule en $\frac{\pi}{4}$.
\textbf{Etude en $\frac{\pi}{4}$.} $M\left(\frac{\pi}{4}\right)=O$ et donc la tangente en $M\left(\frac{\pi}{4}\right)$ est la droite passant par $O$ et d'angle polaire $\frac{\pi}{4}$ ou encore la droite d'équation $y=x$.

\textbf{Etude en $0$.} $M(0)$ est le point de coordonnées cartésiennes $(1,0)$. Pour $\theta\in\left]-\frac{\pi}{4},\frac{\pi}{4}\right[$,

\begin{center}
$\overrightarrow{\frac{dM}{d\theta}}(\theta)=-\frac{\sin(2\theta)}{\sqrt{\cos(2\theta)}}\overrightarrow{u}_\theta+\sqrt{\cos(2\theta)}\overrightarrow{v}_\theta$ et donc $\overrightarrow{\frac{dM}{d\theta}}(0)=\overrightarrow{v}_0=\overrightarrow{j}$.
\end{center}
$M(0)$ est le point de coordonnées cartésiennes $(1,0)$ et la tangente en $M(0)$ est dirigée par $\overrightarrow{j}$

$$\includegraphics{../images/img005530-1}$$


 \item  Soit $\mathcal{C}$ la courbe d'équation polaire $r=\sin\left(\frac{2\theta}{3}\right)$.
\textbf{Domaine d'étude.} \textbullet~Pour $\theta\in \Rr$, 

\begin{center}
$M(\theta+6\pi)=\left[r(\theta+6\pi),\theta+6\pi\right]=\left[r(\theta),\theta+6\pi\right]=\left[r(\theta),\theta\right]=M(\theta)$.
\end{center}
On obtient donc la courbe complète quand $\theta$ décrit un intervalle de longueur $6\pi$ comme $[-3\pi,3\pi]$.
\textbullet~Pour $\theta\in[-3\pi,3\pi]$, 

\begin{center}
$M(-\theta)=\left[r(-\theta),-\theta\right]=\left[-r(\theta),-\theta\right]=\left[r(\theta),\pi-\theta\right]=s_{(Oy)}(M(\theta))$.
\end{center} 
On étudie et on construit la portion de courbe correspondant à $\theta\in[0,3\pi]$ puis on obtient la courbe complète par réflexion d'axe $(Oy)$.
\textbullet~Pour $\theta\in[0,3\pi]$, $M(3\pi-\theta)=\left[r(3\pi-\theta),3\pi-\theta\right]=\left[-r(\theta),3\pi-\theta\right]=\left[r(\theta),-\theta\right]=s_{(Ox)}(M(\theta))$.
 On étudie et on construit la portion de courbe correspondant à $\theta\in\left[0,\frac{3\pi}{2}\right]$ puis on obtient la courbe complète par réflexion d'axe $(Ox)$ puis d'axe $(Oy)$.
 
 
\textbullet~Pour $\theta\in\left[0,\frac{3\pi}{2}\right]$, $M\left(\frac{3\pi}{2}-\theta\right)=\left[r\left(\frac{3\pi}{2}-\theta\right),\frac{3\pi}{2}-\theta\right]=\left[r(\theta),\frac{3\pi}{2}-\theta\right]=s_{y=-x}(M(\theta))$.
 On étudie et on construit la portion de courbe correspondant à $\theta\in\left[0,\frac{3\pi}{4}\right]$ puis on obtient la courbe complète par réflexions successives d'axes la droite d'équation $y=-x$, puis d'axe $(Ox)$ et enfin d'axe $(Oy)$.
 
 
\textbullet~\textbf{Remarque.} La fonction $r$ admet $3\pi$ pour plus petite période strictement positive. Pourtant, on n'obtient pas la courbe complète quand $\theta$ décrit $[0,3\pi]$ car $3\pi$ ne fournit pas un nombre entier de tours. Plus précisément,

\begin{center}
$M(\theta+3\pi)=\left[r(\theta+3\pi),\theta+3\pi\right]=\left[r(\theta),\theta+\pi\right]=s_O(M(\theta))$.
\end{center}

\textbf{Variations et signe de $r$.} La fonction $r$ est strictement positive sur $\left]0,\frac{3\pi}{4}\right]$ et s'annule en $0$. La fonction $r$ est strictement croissante sur $\left[0,\frac{3\pi}{4}\right]$.
\textbullet~$M(0)$ est le point $O$. La tangente en $M(0)$ est la droite passant par $O$ d'angle polaire $0$ c'est-à-dire l'axe $(Ox)$.

$$\includegraphics{../images/img005530-2}$$

$$\includegraphics{../images/img005530-3}$$


 \item  Soit $\mathcal{C}$ la courbe d'équation polaire $r=ae^{b\theta}$.
L'étude est très brève. La fonction $r~:~\theta\mapsto ae^{b\theta}$ est strictement positive et strictement croissante sur $\Rr$. Tout en tournant, on ne cesse de s'écarter de l'origine~:~la courbe est une spirale.

$$\includegraphics{../images/img005530-4}$$

 \item  Soit $\mathcal{C}$ la courbe d'équation polaire $r=2\cos(\theta)+1$.

\textbf{Domaine d'étude.} \textbullet~Pour $\theta\in \Rr$, $M(\theta+2\pi)=M(\theta)$. On obtient donc la courbe complète quand $\theta$ décrit un intervalle de longueur $2\pi$ comme $[-\pi,\pi]$.
\textbullet~Pour $\theta\in[-\pi,\pi]$, $M(-\theta)=s_{(Ox)}(M(\theta))$. On étudie et on construit la portion de courbe correspondant à $\theta\in[0,\pi]$ puis on obtient la courbe complète par réflexion d'axe $(Ox)$.
\textbf{Variations et signe de $r$.} La fonction $r$ est strictement décroissante sur $[0,\pi]$. La fonction $r$ est strictement positive sur $\left[0,\frac{2\pi}{3}\right[$, strictement négative sur $\left]\frac{2\pi}{3},0\right]$ et s'annule en $\frac{2\pi}{3}$. Donc la fonction $\theta\mapsto OM(\theta)=|r(\theta)|$ est strictement décroissante sur $\left[0,\frac{2\pi}{3}\right]$ et strictement croissante sur $\left[\frac{2\pi}{3},\pi\right]$.
\textbullet~$M\left(\frac{2\pi}{3}\right)$ est le point $O$. La tangente en $M\left(\frac{2\pi}{3}\right)$ est la droite passant par $O$ d'angle polaire $\frac{2\pi}{3}$ c'est-à-dire la droite d'équation $y=-\sqrt{3}x$.
\textbullet~Par symétrie par rapport à $(Ox)$, les tangentes en $M(0)$ et $M(\pi)$ sont parrallèles à $(Oy)$.

$$\includegraphics{../images/img005530-5}$$


 \item  Soit $\mathcal{C}$ la courbe d'équation polaire $r=\tan\left(\frac{2\theta}{3}\right)$.
\textbf{Domaine d'étude.}  Notons $D$ le domaine de définition de la fonction $r~:~\theta\mapsto\tan\left(\frac{2\theta}{3}\right)$.
\textbullet~$\theta\in D\Leftrightarrow\theta+6\pi\in D$ et $M(\theta+6\pi)=M(\theta)$. On obtient donc la courbe complète quand $\theta$ décrit un intervalle de longueur $6\pi$ comme $[-3\pi,3\pi]$.
\textbullet~$\theta\in D\Leftrightarrow -\theta\in D$ et $M(-\theta)=s_{(Oy)}(M(\theta))$. On étudie et on construit la portion de courbe correspondant à $\theta\in[0,3\pi]$ puis on obtient la courbe complète par réflexion d'axe $(Oy)$.
\textbullet~$\theta\in D\Leftrightarrow 3\pi-\theta\in D$ et $M(3\pi-\theta)=s_{(Ox)}(M(\theta))$. On étudie et on construit la portion de courbe correspondant à $\theta\in\left[0,\frac{3\pi}{2}\right]$ puis on obtient la courbe complète par réflexion d'axe $(Ox)$ puis par réflexion d'axe $(Oy)$.
\textbullet~$\theta\in D\Leftrightarrow \frac{3\pi}{2}-\theta\in D$ et 

\begin{center}
$M\left(\frac{3\pi}{2}-\theta\right)=\left[-r(\theta),\frac{3\pi}{2}-\theta\right]=\left[r(\theta),\frac{\pi}{2}-\theta\right]=s_{y=x}(M(\theta))$.
\end{center}
On étudie et on construit la portion de courbe correspondant à $\theta\in\left[0,\frac{3\pi}{4}\right]$ puis on obtient la courbe complète par réflexions successives d'axe la droite d'équation $y=x$, puis d'axe $(Ox)$ et enfin d'axe $(Oy)$.
\textbullet~Pour $\theta\in\left[0,\frac{3\pi}{4}\right]$, $r(\theta)$ existe si et seulement si $\theta\neq\frac{3\pi}{4}$. On étudie donc sur $\theta\in\left[0,\frac{3\pi}{4}\right[$.

\textbf{Variations et signe de $r$.} La fonction $r$ est strictement croissante sur $\left[0,\frac{3\pi}{4}\right[$, strictement positive sur $\left]0,\frac{3\pi}{4}\right[$ et s'annule en $0$.

\textbullet~La tangente en $M(0)=O$ est la droite passant par $O$ et d'angle polaire $0$ c'est-à-dire l'axe $(Ox)$.
\textbullet~\textbf{Etude quand $\theta$ tend vers $\frac{3\pi}{4}$.} Quand $\theta$ tend vers $\frac{3\pi}{4}$ par valeurs inférieures, $r(\theta)$ tend vers $+\infty$. la courbe admet donc une direction asymptotique d'angle polaire $\frac{3\pi}{4}$ ou encore d'équation $y=-x$.
Recherchons une éventuelle droite asymptote. Pour cela, étudions $\displaystyle\lim_{\substack{\theta\rightarrow\frac{3\pi}{4}\\ \theta<\frac{3\pi}{4}}}r(\theta)\sin\left(\theta-\frac{3\pi}{4}\right)$. Posons $h=\frac{3\pi}{4}-\theta$ ou encore $\theta=\frac{3\pi}{4}-h$.
\begin{center}
$r(\theta)\sin\left(\theta-\frac{3\pi}{4}\right)=\tan\left(\frac{\pi}{2}-\frac{2h}{3}\right)\sin(-h)=-\cotan h\sin h=-\cos h\rightarrow-1$.
\end{center}
Ainsi, $\mathcal{C}$ admet une droite asymptote $(D)$ quand $\theta$ tend vers $\frac{3\pi}{4}$. De plus,

\begin{center}
$M(x,y)\in(D)\Leftrightarrow\overrightarrow{OM}.\overrightarrow{v}_{\frac{3\pi}{4}}=-1\Leftrightarrow-\frac{1}{\sqrt{2}}x-\frac{1}{\sqrt{2}}y=-1\Leftrightarrow y=-x+\sqrt{2}$.
\end{center}

$$\includegraphics{../images/img005530-6}$$

\end{enumerate}
\fincorrection
\correction{005531}
\textbf{Domaine d'étude.} Notons $D$ le domaine de définition de la fonction $r~:~\theta\mapsto\frac{2\cos\theta+1}{2\sin\theta+1}$.
$\forall \theta\in\Rr$, $\theta\in D\Leftrightarrow \theta+2\pi\in D$ et $M(\theta+2\pi)=M(\theta)$. On obtient donc la courbe complète quand $\theta$ décrit un intervalle de longueur $2\pi$ comme $[-\pi,\pi]$. Pour $\theta\in[-\pi,\pi]$, $2\sin\theta+1=0\Leftrightarrow\theta\in\left\{-\frac{5\pi}{6},-\frac{\pi}{6}\right\}$. On étudie donc la courbe sur $[-\pi,\pi]\setminus\left\{-\frac{5\pi}{6},-\frac{\pi}{6}\right\}$.
\textbf{Signe de $r$.}

$$\includegraphics{../images/img005531-1}$$


\textbf{Variations de $r$.} La fonction $r$ est dérivable sur  $[-\pi,\pi]\setminus\left\{-\frac{5\pi}{6},-\frac{\pi}{6}\right\}$ et pour $\theta\in[-\pi,\pi]\setminus\left\{-\frac{5\pi}{6},-\frac{\pi}{6}\right\}$

\begin{center}
$r'(\theta)=\frac{-2\sin\theta(2\sin\theta+1)-2\cos\theta(2\cos\theta+1)}{(2\sin\theta+1)^2}=\frac{-4-2\cos\theta-2\sin\theta}{(2\sin\theta+1)^2}=\frac{-4-2\sqrt{2}\cos\left(\theta-\frac{\pi}{4}\right)}{(2\sin\theta+1)^2}<0$.
\end{center}
La fonction $r$ est strictement décroissante sur $\left[-\pi,-\frac{5\pi}{6}\right[$, sur $\left]-\frac{5\pi}{6},-\frac{\pi}{6}\right[$ et sur $\left]-\frac{\pi}{6},\pi\right]$.
\textbf{Etude quand $\theta$ tend vers $-\frac{5\pi}{6}$.} $\displaystyle\lim_{\substack{\theta\rightarrow-\frac{5\pi}{6}\\ x<-\frac{5\pi}{6}}}r(\theta)=-\infty$ et $\displaystyle\lim_{\substack{\theta\rightarrow-\frac{5\pi}{6}\\ x>-\frac{5\pi}{6}}}r(\theta)=+\infty$. Donc la courbe $\mathcal{C}$ admet une direction asymptotique d'angle polaire $-\frac{5\pi}{6}$ ou encore d'équation $y=\frac{1}{\sqrt{3}}x$. Etudions maintenant l'existence d'une éventuelle droite asymptote et pour cela étudions $\lim_{\theta \rightarrow -\frac{5\pi}{6}}r(\theta)\sin\left(\theta+\frac{5\pi}{6}\right)$.
On pose $h=\theta+\frac{5\pi}{6}$ ou encore $\theta=-\frac{5\pi}{6}+h$ de sorte que $\theta$ tend vers $-\frac{5\pi}{6}$ si et seulement si $h$ tend vers $0$. Quand $h$ tend vers $0$
\begin{align*}\ensuremath
r(\theta)\sin\left(\theta+\frac{5\pi}{6}\right)&=\frac{2\cos\left(-\frac{5\pi}{6}+h\right)+1}{2\sin\left(-\frac{5\pi}{6}+h\right)+1}\sin h=\frac{(1-\sqrt{3}\cos h)+\sin h}{-\sqrt{3}\sin h+(1-\cos h)}\sin h\sim\frac{1-\sqrt{3}}{-\sqrt{3}h}\times h=1-\frac{1}{\sqrt{3}}.
\end{align*}
Par suite, $\mathcal{C}$ admet une droite asymptote $(D_1)$ quand $\theta$ tend vers $-\frac{5\pi}{6}$. De plus

\begin{center}
$M(x,y)\in(D_1)\Leftrightarrow \overrightarrow{OM}.\overrightarrow{v}_{-\frac{5\pi}{6}}=1-\frac{1}{\sqrt{3}}\Leftrightarrow\frac{1}{2}x-\frac{\sqrt{3}}{2}y=1-\frac{1}{\sqrt{3}}\Leftrightarrow y=\frac{1}{\sqrt{3}}x+\frac{2}{3}-\frac{2}{\sqrt{3}}$
\end{center}
\textbf{Etude quand $\theta$ tend vers $-\frac{\pi}{6}$.} $\displaystyle\lim_{\substack{\theta\rightarrow-\frac{\pi}{6}\\ x<-\frac{\pi}{6}}}r(\theta)=-\infty$ et $\displaystyle\lim_{\substack{\theta\rightarrow-\frac{\pi}{6}\\ x>-\frac{\pi}{6}}}r(\theta)=+\infty$. Donc la courbe $\mathcal{C}$ admet une direction asymptotique d'angle polaire $-\frac{\pi}{6}$ ou encore d'équation $y=-\frac{1}{\sqrt{3}}x$.
On pose ensuite $h=\theta+\frac{\pi}{6}$. Quand $h$ tend vers $0$
\begin{align*}\ensuremath
r(\theta)\sin\left(\theta+\frac{\pi}{6}\right)&=\frac{2\cos\left(-\frac{\pi}{6}+h\right)+1}{2\sin\left(-\frac{\pi}{6}+h\right)+1}\sin h=\frac{(1+\sqrt{3}\cos h)+\sin h}{\sqrt{3}\sin h+(1-\cos h)}\sin h\sim\frac{1+\sqrt{3}}{\sqrt{3}h}\times h=1+\frac{1}{\sqrt{3}}.
\end{align*}
Par suite, $\mathcal{C}$ admet une droite asymptote $(D_2)$ quand $\theta$ tend vers $-\frac{\pi}{6}$. De plus

\begin{center}
$M(x,y)\in(D_2)\Leftrightarrow \overrightarrow{OM}.\overrightarrow{v}_{-\frac{\pi}{6}}=1+\frac{1}{\sqrt{3}}\Leftrightarrow\frac{1}{2}x+\frac{\sqrt{3}}{2}y=1+\frac{1}{\sqrt{3}}\Leftrightarrow y=-\frac{1}{\sqrt{3}}x+\frac{2}{3}+\frac{2}{\sqrt{3}}$
\end{center}

\textbf{Tableau de variation de $r$.}


$$\includegraphics{../images/img005531-2}$$


\textbf{Recherche des points multiples.} Soit $(\theta_1,\theta_2)\in\left([-\pi,\pi]\setminus\left\{-\frac{5\pi}{6},-\frac{\pi}{6}\right\}\right)^2$ tel que $\theta_1<\theta_2$. On suppose de plus que $\theta_1\notin\left\{\pm\frac{2\pi}{3}\right\}$ et $\theta_1\notin\left\{\pm\frac{2\pi}{3}\right\}$ de sorte que $M(\theta_1)\neq O$ et $M(\theta_2)\neq O$. 

\begin{align*}\ensuremath
M(\theta_1)=M(\theta_2)&\Leftrightarrow(\exists k\in\Zz/\;\theta_2=\theta_1+2k\pi\;\text{et}\;r(\theta_2)=r(\theta_1))\;\text{ou}\;(\exists k\in\Zz/\;\theta_2=\theta_1+\pi+2k\pi\;\text{et}\;r(\theta_2)=-r(\theta_1))\\
 &\Leftrightarrow\theta_1\in[-\pi,0],\;\theta_2=\theta_1+\pi\;\text{et}\;r(\theta_2)=-r(\theta_1)\\
 &\Leftrightarrow\theta_1\in[-\pi,0],\;\theta_2=\theta_1+\pi\;\text{et}\;\frac{-2\cos(\theta_1)+1}{-2\sin(\theta_1)+1}=-\frac{2\cos(\theta_1)+1}{2\sin(\theta_1)+1}.
\end{align*}
Maintenant, pour $\theta\in[-\pi,0]\setminus\left\{-\frac{5\pi}{6},-\frac{2\pi}{3},-\frac{\pi}{6}\right\}$

\begin{align*}\ensuremath
\frac{-2\cos(\theta)+1}{-2\sin(\theta)+1}=-\frac{2\cos(\theta)+1}{2\sin(\theta)+1}&\Leftrightarrow-4\cos(\theta)\sin(\theta)+1=4\cos(\theta
)\sin(\theta)-1\Leftrightarrow\sin(2\theta)=\frac{1}{2}\\
 &\Leftrightarrow2\theta\in\frac{\pi}{6}+2\pi\Zz\;\text{ou}\;2\theta\in\frac{5\pi}{6}+2\pi\Zz\Leftrightarrow\theta\in\frac{\pi}{12}+\pi\Zz\;\text{ou}\;\theta\in\frac{5\pi}{12}+\pi\Zz\\
 &\Leftrightarrow\theta\in\left\{-\frac{11\pi}{12},-\frac{7\pi}{12}\right\}.
\end{align*}
Ainsi, les points doubles distincts de l'origine sont $M\left(-\frac{11\pi}{12}\right)=M\left(\frac{\pi}{12}\right)$ et $M\left(-\frac{7\pi}{12}\right)=M\left(\frac{5\pi}{12}\right)$. Sinon, $M\left(-\frac{2\pi}{3}\right)=M\left(\frac{2\pi}{3}\right)=O$.

$$\includegraphics{../images/img005531-3}$$

\fincorrection
\correction{005532}
\begin{enumerate}
 \item  \textbf{Domaine d'étude.} La fonction $r$ est $2\pi$-périodique et paire. Donc on étudie et on construit la courbe quand $\theta$ décrit $[0,\pi]$ et on obtient la courbe complète par réflexion d'axe $(Ox)$.
\textbf{Variations et signe de $r$.} La fonction $r$ est strictement décroissante sur $[0,\pi]$, strictement positive sur $]0,\pi]$ et s'annule en $\pi$.
\textbf{Etude pour $\theta=\pi$.} La tangente en $M(\pi)=O$ est la droite passant par $O$ d'angle polaire $\pi$ c'est-à-dire l'axe $(Ox)$. Par symétrie par rapport à $(Ox)$, le point $M(\pi)$ est un point de rebroussement de première espèce.

$$\includegraphics{../images/img005532-1}$$

 \item  Soient $\theta\in[-\pi,\pi]$ puis $M=O+a(1+\cos\theta)\overrightarrow{u}_{\theta}$ le point de $\mathcal{C}$ de paramètre $\theta$. 

\begin{align*}\ensuremath
\overrightarrow{\frac{dM}{d\theta}}&=-a\sin\theta\overrightarrow{u}_{\theta}+a(1+\cos\theta)\overrightarrow{v}_{\theta}=2a\cos\left(\frac{\theta}{2}\right)\left(-\sin\left(\frac{\theta}{2}\right)\overrightarrow{u}_{\theta}+\cos\left(\frac{\theta}{2}\right)\overrightarrow{v}_{\theta}\right)\\
 &=2a\cos\left(\frac{\theta}{2}\right)\left(\cos\left(\frac{\theta}{2}+\frac{\pi}{2}\right)\overrightarrow{u}_{\theta}+\sin\left(\frac{\theta}{2}+\frac{\pi}{2}\right)\overrightarrow{v}_{\theta}\right)=2a\cos\left(\frac{\theta}{2}\right)\overrightarrow{u}_{\frac{3\theta}{2}+\frac{\pi}{2}}.
\end{align*}
\textbf{Longueur $\ell$ de la cardioïde.} On a $\left\|\overrightarrow{\frac{dM}{d\theta}}\right\|=\left|2a\cos\left(\frac{\theta}{2}\right)\right|=2a\cos\left(\frac{\theta}{2}\right)$ (pour $\theta\in[-\pi,\pi]$) et donc

\begin{center}
$\ell=\int_{-\pi}^{\pi}\left\|\overrightarrow{\frac{dM}{d\theta}}\right\|d\theta=2a\int_{-\pi}^{\pi}\cos(\theta/2)\;d\theta=4a\left[\sin(\theta/2)\right]_{-\pi}^\pi=8a$.
\end{center}

\begin{center}
\shadowbox{
La cardioïde d'équation polaire $r=a(1+\cos\theta)$, $a>0$, a pour longueur $8a$.
}
\end{center}
\textbf{Développée.} Le point $M(\theta)$ est régulier si et seulement si $\theta\neq\pm\pi$. Dans ce cas, 

\begin{center}
$\frac{ds}{d\theta}=\left\|\overrightarrow{\frac{dM}{d\theta}}\right\|=2a\cos\left(\frac{\theta}{2}\right)$ et aussi $\overrightarrow{\tau}(\theta)=\overrightarrow{u}_{\frac{3\theta}{2}+\frac{\pi}{2}}$
\end{center}

En notant $\alpha(\theta)$ une mesure de l'angle $\left(\overrightarrow{i},\overrightarrow{\tau}(\theta)\right)$, on peut prendre $\alpha(\theta)=\frac{3\theta}{2}+\frac{\pi}{2}$. En notant $R(\theta)$ le rayon de courbure au point $M(\theta)$,

\begin{center}
$R(\theta)=\frac{ds}{d\alpha}=\frac{ds/d\theta}{d\alpha/d\theta}=\frac{4}{3}a\cos\left(\frac{\theta}{2}\right)$.
\end{center}
Ensuite, $\overrightarrow{n}(\theta)=r_{\pi/2}\left(\overrightarrow{\tau}(\theta)\right)=-\overrightarrow{u}_{3\theta/2}$ et donc, en notant $\Omega(\theta)$ le centre de courbure au point $M(\theta)$,

\begin{align*}\ensuremath
\Omega(\theta)&=M(\theta)+R(\theta)\overrightarrow{n}(\theta)\\
 &=O+a(1+\cos\theta)\overrightarrow{u}_\theta-\frac{4}{3}a\cos\left(\frac{\theta}{2}\right)\overrightarrow{u}_{3\theta/2}\\
  &=O+a(1+\cos\theta)\left(\cos(\theta)\overrightarrow{i}+\sin(\theta)\overrightarrow{j}\right)-\frac{4}{3}a\left(\cos\left(\frac{\theta}{2}\right)\cos\left(\frac{3\theta}{2}\right)\overrightarrow{i}+\cos\left(\frac{\theta}{2}\right)\sin\left(\frac{3\theta}{2}\right)\overrightarrow{j}\right)\\
  &=O+a\left[\left(\cos(\theta)+\cos^2(\theta)-\frac{2}{3}(\cos(\theta)+\cos(2\theta))\right)\overrightarrow{i}+\left(\sin(\theta)+\sin(\theta)\cos(\theta)-\frac{2}{3}(\sin(\theta)+\sin(2\theta))\right)\overrightarrow{j}\right]\\
    &=O+a\left[\left(\frac{2}{3}+\frac{1}{3}\cos(\theta)-\frac{1}{3}\cos^2(\theta)\right)\overrightarrow{i}+\left(\frac{1}{3}\sin(\theta)-\frac{1}{3}\sin(\theta)\cos(\theta)\right)\overrightarrow{j}\right]\\
    &=O+\frac{2a}{3}\overrightarrow{i}+\frac{a}{3}(1-\cos\theta)\overrightarrow{u}_\theta
\end{align*}
Notons $\Gamma$ la développée cherchée. On a $\Gamma=t\circ h(\mathcal{C}_1)$ où $t$ est la translation de vecteur $\frac{2a}{3}\overrightarrow{i}$, $h$ est l'homothétie de centre $O$ et de rapport $\frac{1}{3}$ et $\mathcal{C}_1$ la courbe d'équation polaire $r=a(1-\cos\theta)$.
Maintenant, en notant $r$ la fonction $\theta\mapsto a(1+\cos\theta)$ et $r_1$ la fonction $\theta\mapsto a(1-\cos\theta)$, 

\begin{center}
$[r_1(\theta+\pi),\theta+\pi)]=[a(1+\cos\theta),\theta+\pi]=s_O([r(\theta),\theta)])$.
\end{center}
La courbe $\mathcal{C}_1$ est donc la symétrique par rapport à $O$ de la courbe $\mathcal{C}$. En résumé, la développée de $\mathcal{C}$ est l'image de $\mathcal{C}$ par la transformation $t\circ h\circ s_O$ : c'est encore une cardioïde.

$$\includegraphics{../images/img005532-2}$$

\end{enumerate}
\fincorrection
\correction{005533}
Soient $(R,\theta)\in\Rr^2$ puis $M$ le point du plan dont un couple de coordonnées polaires est $[r,\theta]$.

\begin{align*}\ensuremath
M\in\mathcal{C}&\Leftrightarrow x^2(x^2+y^2)-(y-x)^2=0
\Leftrightarrow r^2\cos^2\theta\times r^2-(r\sin\theta-r\cos\theta)^2=0\\
 &\Leftrightarrow r^2[r^2\cos^2\theta-(\sin\theta-\cos\theta)^2]=0\Leftrightarrow r=0\;\text{ou}\;r^2=\left(\frac{\sin\theta-\cos\theta}{\cos\theta}\right)^2\;(\cos\theta=0\;\text{ne fournit pas de solution})\\
 &\Leftrightarrow r=0\;\text{ou}\;r=\tan\theta-1\;\text{ou}\;r=1-\tan\theta.
\end{align*}
$\mathcal{C}$ est donc la réunion de la courbe $(\mathcal{C}_1)$ d'équation polaire $r=\tan\theta-1$, $(\mathcal{C}_2)$ d'équation polaire $r=1-\tan\theta$ et $\{O\}$.
On note que le point $O$ appartient à $(\mathcal{C}_1)$ car $\theta=\frac{\pi}{4}$ fournit $r=0$. Donc $\mathcal{C}=\mathcal{C}_1\cup\mathcal{C}_2\cup\{O\}=\mathcal{C}_1\cup\mathcal{C}_2$. Ensuite, on notant $r_1$ et $r_2$ respectivement la fonction $\theta\mapsto\tan\theta-1$ et $r_2=-r_1$,

\begin{center}
$M[\theta+\pi,r_2(\theta+\pi)]=M[\theta+\pi,r_2(\theta)]=M[\theta+\pi,-r_1(\theta)]=M[\theta,r_1(\theta)]$,
\end{center}
et comme $\theta+\pi$ décrit $\Rr$ si et seulement si $\theta$ décrit $\Rr$, les courbes $\mathcal{C}_1$ et $\mathcal{C}_2$ sont une seule et même courbe.

\begin{center}
\shadowbox{
$\mathcal{C}$ est la courbe d'équation polaire $r=\tan\theta-1$.
}
\end{center}
\textbf{Construction de $\mathcal{C}$.}

$$\includegraphics{../images/img005533-1}$$

\fincorrection
\correction{005534}
\textbf{Développée.} $M(\theta)=O+ae^\theta\overrightarrow{u}_\theta$ puis 

\begin{center}
$\overrightarrow{\frac{dM}{d\theta}}=ae^\theta(\overrightarrow{u}_\theta+\overrightarrow{v}_\theta)=a\sqrt{2}e^{\theta}\left(\cos\left(\frac{\pi}{4}\right)\overrightarrow{u}_\theta+\sin\left(\frac{\pi}{4}\right)\overrightarrow{v}_\theta\right)=a\sqrt{2}e^\theta\overrightarrow{u}_{\theta+\frac{\pi}{4}}$.
\end{center}
On en déduit $\frac{ds}{d\theta}=a\sqrt{2}e^\theta$ et $\overrightarrow{\tau}(\theta)=\overrightarrow{u}_{\theta+\frac{\pi}{4}}$. On peut alors prendre $\alpha(\theta)=\theta+\frac{\pi}{4}$ et donc $\frac{d\alpha}{d\theta}=1$. Par suite

\begin{center}
$R(\theta)=\frac{ds/d\theta}{d\alpha/d\theta}=\frac{a\sqrt{2}e^\theta}{1}=a\sqrt{2}e^\theta$.
\end{center}
D'autre part, $\overrightarrow{n}(\theta)=\overrightarrow{\tau}\left(\theta+\frac{\pi}{2}\right)=\overrightarrow{u}_{\theta+\frac{3\pi}{4}}=\frac{1}{\sqrt{2}}(-\overrightarrow{u}_\theta+\overrightarrow{v}_\theta)$ et donc

\begin{center}
$\Omega(\theta)=M(\theta)+R(\theta)\overrightarrow{n}(\theta)=O+ae^\theta\overrightarrow{u}_\theta+a\sqrt{2}e^\theta.\frac{1}{\sqrt{2}}(-\overrightarrow{u}_\theta+\overrightarrow{v}_\theta)=O+ae^\theta\overrightarrow{v}_\theta=r_{O,\frac{\pi}{2}}(M(\theta))$.
\end{center}
La développée de la spirale logarithmique d'équation polaire $r=ae^{\theta}$ est l'image de cette spirale par le quart de tour direct de centre $O$.

$$\includegraphics{../images/img005534-1}$$

\fincorrection
\correction{006989}\
\begin{enumerate}
\item L'expression $r(\theta)=\frac{1}{\sqrt{\tan(2\theta)}}$ est bien définie 
sur $\mathcal{D}= ]0,\frac\pi4[$ (il faut $\tan(2\theta)$ 
bien défini et strictement positif).
\begin{itemize} 
\item Passages par l'origine\\
Puisque $r$ ne s'annule pas, la courbe ne passe pas par l'origine. Mais elle admet l'origine pour point limite: $r(\theta)\xrightarrow[\theta\to\frac{\pi}{4}^-]{}0^+$.
\item Variations et signe de la fonction $r$\\
La fonction $r$ est strictement décroissante, et strictement positive, sur $]0;\frac{\pi}{4}[$:
$$\begin{array}{c|lcr}
\theta&0&\ &\frac{\pi}{4}\\\hline
\ &+\infty &\ & \\
r&\ &\searrow &\ \\
 & &\ &0 \\
\end{array}$$
Cela signifie que la courbe tourne (dans le sens trigonométrique) en se rapprochant de l'origine.

\item Tangente à l'origine\\
Le point $M(\frac{\pi}{4})$ est à l'origine : la tangente en ce point est donc dirigée par $\vec{u}_{\frac{\pi}{4}}$, c'est la première bissectrice.

\item \'Etude des branches infinies\\
Lorsque $\theta$ tend vers 0, $r(\theta)$ tend vers $+\infty$: il y a donc une branche infinie. Pour étudier sa nature, passons en coordonnées cartésiennes:
\begin{eqnarray*}
x(\theta)&=&\frac{\cos\theta}{\sqrt{\tan(2\theta)}}\xrightarrow[\theta\to 0^+]{}+\infty\\
y(\theta)&=&\frac{\sin\theta}{\sqrt{\tan(2\theta)}}\underset{0^+}{\sim}\frac{\theta}{\sqrt{2\theta}}
\end{eqnarray*}
Ainsi $x(\theta)\xrightarrow[\theta\to 0^+]{}+\infty$, $y(\theta)\xrightarrow[\theta\to 0^+]{}0$: la droite d'équation $y=0$ est asymptote horizontale. 
\end{itemize}

\begin{center}
\begin{tikzpicture}[scale=1.2]

% Axes
     \draw[->,>=latex,thick, gray] (-1,0)--(5.5,0) node[below,black] {$x$};
     \draw[->,>=latex,thick, gray] (0,-1.5)--(0,1.5) node[right,black] {$y$};  

 % Ticks
        \foreach \x in {1,...,4}
                \draw (\x,2pt) -- (\x,-2pt)
                        node[anchor=north] {\x};
        \foreach \x in {1}
                \draw (2pt,\x) -- (-2pt,\x)
                        node[anchor=east] {\x};
        \foreach \x in {-1}
                \draw (2pt,\x) -- (-2pt,\x)
                        node[anchor=east] {\x};

  \draw[->,>=latex,thick, blue] (0,0)--+(45:1);
  \draw[->,>=latex,thick, blue] (0,0)--+(45:-1);

% Courbe

\draw [very thick, color=red, domain=0.02:pi/4, samples=100, smooth]
  plot (xy polar cs:angle=\x r, radius={1/sqrt(tan(2*\x r)});
% \draw [very thick, color=gray, domain=3.162:3.92, samples=100, smooth]
%   plot (xy polar cs:angle=\x r, radius={1/sqrt(tan(2*\x r)});

   \fill (0:0) circle (2pt) node[above left] {$\theta=\frac\pi4$};

\end{tikzpicture}  
\end{center}



\item L'expression $r(\theta)=\frac{\sin^2\theta}{\cos\theta}$ est bien définie sur $\mathcal{D}=]-\frac\pi2,\frac\pi2[$.
\begin{itemize} 
\item Réduction du domaine d'étude\\
Comme $r$ est paire, il suffit en fait de faire l'étude pour les $\theta>0$, donc 
sur $[0;\frac{\pi}{2}[$, 
puis de compléter  par réflexion d'axe $(Ox)$, en effet
$M(-\theta) = s_{(Ox)}(M(\theta))$.

On se restreint donc dans la suite à $\theta\in[0,\frac{\pi}{2}[$, puis on obtient la courbe 
complète par réflexion d'axe $(Ox)$.

\item Passages par l'origine\\
La courbe passe par l'origine si $r$ s'annule: pour $\theta\in[0,\frac{\pi}{2}[$,
$$r(\theta)=0\Longleftrightarrow\theta=0$$ 

\item Variations et signe de la fonction $r$\\
La fonction $r$ est strictement croissante sur $]0,\frac{\pi}{2}[$, 
strictement positive sur $\left]0,\frac{\pi}{2}\right]$ et s'annule en $0$:
$$\begin{array}{c|lcr}
\theta&0&\ &\frac{\pi}{2}\\\hline
\ & &\ & +\infty\\
r&\ &\nearrow &\ \\
 &0 &\ & \\
\end{array}$$
Ainsi la courbe tourne en s'éloignant de l'origine.

\item Tangente à l'origine\\
La courbe passe par l'origine en $\theta=0$. Par conséquent, la tangente en $O=M\left(0\right)$ est la droite passant par $O$ et d'angle polaire $0$, c'est-à-dire l'axe $(Ox)$. 

\item \'Etude des branches infinies\\
Lorsque $\theta$ tend vers $\frac{\pi}{2}$, $r(\theta)$ tend vers $+\infty$: il y a donc une branche infinie. Pour étudier sa nature, passons en coordonnées cartésiennes:
\begin{eqnarray*}
x(\theta)&=&\sin^2\theta\xrightarrow[\theta\to \frac{\pi}{2}^-]{}1\\
y(\theta)&=&\frac{\sin^3\theta}{\cos\theta}\xrightarrow[\theta\to \frac{\pi}{2}^-]{}+\infty
\end{eqnarray*}
Ainsi la droite d'équation $x=1$ est asymptote verticale. 
\end{itemize}


\begin{center}
\begin{tikzpicture}[scale=1.5]

% Axes
     \draw[->,>=latex,thick, gray] (-1.5,0)--(1.5,0) node[below,black] {$x$};
     \draw[->,>=latex,thick, gray] (0,-4)--(0,4) node[right,black] {$y$};  

 % Ticks
        \foreach \x in {1,...,1}
                \draw (\x,2pt) -- (\x,-2pt)
                        node[anchor=north] {\x};
        \foreach \x in {1,2,3}
                \draw (2pt,\x) -- (-2pt,\x)
                        node[anchor=east] {\x};
        \foreach \x in {-1,-2,-3}
                \draw (2pt,\x) -- (-2pt,\x)
                        node[anchor=east] {\x};

  \draw[->,>=latex,thick, blue] (0,0)--+(0:0.7);
  \draw[->,>=latex,thick, blue] (0,0)--+(0:-0.7);

% Courbe

\draw [very thick, color=red, domain=0:1.34, samples=100, smooth]
  plot (xy polar cs:angle=\x r, radius={sin(\x r)^2/(cos(\x r)});

\draw [very thick, color=gray, domain=0:-1.34, samples=100, smooth]
  plot (xy polar cs:angle=\x r, radius={sin(\x r)^2/(cos(\x r)});


   \fill (0,0) circle (1pt) node[above left] {$\theta=0$};

  \draw[thick, blue] (1,0)--+(0,4)--+(0,-4);
  \node[right] at (1,3) {$x=1$};
\end{tikzpicture}  
\end{center}


\item L'expression $r(\theta)=\sqrt{\cos(2\theta)}$ est bien définie si $\cos(2\theta)$ est positif, \textsl{i.e.} sur $\mathcal{D}=\bigcup_{k\in\Z}[-\frac{\pi}{4}+k\pi;\frac{\pi}{4}+k\pi]$.
\begin{itemize} 
\item Réduction du domaine d'étude\\
La fonction $r$ est $\pi$-périodique: on l'étudie sur un intervalle de longueur $\pi$, par exemple $[-\frac{\pi}{2};\frac{\pi}{2}]\cap\mathcal{D}=[-\frac{\pi}{4};\frac{\pi}{4}]$, puis on complète par rotation d'angle $\pi$. De plus $r$ est paire: $\theta\in D\Leftrightarrow-\theta\in D$, et pour $\theta\in D$ on a
$$M(-\theta)=\left[r(-\theta):-\theta\right]=\left[r(\theta):-\theta\right]=s_{(Ox)}(M(\theta))$$
Finalement, on étudie et on construit la portion de courbe correspondant à $\theta\in[0,\frac{\pi}{4}]$, puis on obtient la courbe complète d'abord par réflexion d'axe $(Ox)$, puis par rotation d'angle $\pi$ (symétrie centrale par rapport à l'origine).
\item Passages par l'origine\\
La courbe passe par l'origine si $r$ s'annule: pour $\theta\in[0,\frac{\pi}{4}]$,
$$r(\theta)=0\Longleftrightarrow\theta=\frac{\pi}{4}$$ 
\item Variations et signe de la fonction $r$\\
La fonction $r$ est strictement décroissante sur $[0,\frac{\pi}{4}]$, strictement positive sur $\left]0,\frac{\pi}{4}\right]$ et s'annule en $\frac{\pi}{4}$:
$$\begin{array}{c|lcr}
\theta&0&\ &\frac{\pi}{4}\\\hline
\ &1 &\ & \\
r&\ &\searrow &\ \\
 & &\ &0 \\
\end{array}$$
Ainsi la courbe tourne en se rapprochant de l'origine.

\item Tangentes\\
La courbe passe par l'origine en $\theta=\frac{\pi}{4}$, et donc la tangente en $M\left(\frac{\pi}{4}\right)$ est la droite passant par $O$ et d'angle polaire $\frac{\pi}{4}$ c'est-à-dire la première bissectrice. 

En $\theta=0$, $r(0)=1$ (et $r'(0)=0$) et la tangente est dirigée par le vecteur
$$\overrightarrow{\frac{\dd M}{\dd \theta}}(0)=r'(0)\vec{u}_0+r(0)\vec{v}_0=\vec{v}_0=\vec{j}$$
et la tangente à la courbe au point $M(0)$ de coordonnées cartésiennes $(1,0)$ est donc verticale.
\end{itemize}

\begin{center}
\begin{tikzpicture}[scale=3.5]

% Axes
     \draw[->,>=latex,thick, gray] (-1.3,0)--(1.5,0) node[below,black] {$x$};
     \draw[->,>=latex,thick, gray] (0,-0.6)--(0,0.7) node[right,black] {$y$};  

 % Ticks
        \foreach \x in {1,...,1}
                \draw (\x,2pt) -- (\x,-2pt)
                        node[anchor=north] {\x};
%       \foreach \x in {1}
%               \draw (2pt,\x) -- (-2pt,\x)
%                       node[anchor=east] {\x};
%       \foreach \x in {-1}
%               \draw (2pt,\x) -- (-2pt,\x)
%                       node[anchor=east] {\x};

  \draw[->,>=latex,thick, blue] (0,0)--+(45:0.7);
  \draw[->,>=latex,thick, blue] (0,0)--+(45:-0.7);

% Courbe

\draw [very thick, color=red, domain=0:pi/4, samples=100, smooth]
  plot (xy polar cs:angle=\x r, radius={sqrt(cos(2*\x r)});

\draw [very thick, color=gray, domain=-pi/4:0, samples=100, smooth]
  plot (xy polar cs:angle=\x r, radius={sqrt(cos(2*\x r)});

\draw [very thick, color=gray, domain=2.357:5*pi/4, samples=100, smooth]
  plot (xy polar cs:angle=\x r, radius={sqrt(cos(2*\x r)});


   \fill (0:0) circle (1pt) node[left] {$\theta=\frac\pi4$};
   \fill (1,0) circle (1pt) node[above right] {$\theta=0$};
\end{tikzpicture}  
\end{center}

\end{enumerate}
\fincorrection
\correction{006990} Les deux équations sont $2\pi$-périodiques en $\theta$, soit donc $\theta\in[0;2\pi[$, cherchons pour chaque courbe le vecteur tangent au point $M_i(\theta)$. Déjà, $\mathcal{C}_2$ ne passe pas par le p\^ole mais $\mathcal{C}_1$ oui, pour $\theta=\pi$: elle a donc en $M_1(\pi)$ une tangente horizontale (dirigée par le vecteur $\vec{u}_\pi$) et $N_1(\pi)$ est l'axe $(Oy)$. Dans tous les autres cas, la tangente à $\mathcal{C}_i$ au point $M_i(\theta)$ est dirigée par le vecteur
$$\overrightarrow{\frac{\dd OM_i}{\dd \theta}}(\theta)=r'(\theta)\vec{u}_\theta+r(\theta)\vec{v}_\theta=-\sin\theta\,\vec{u}_\theta+(a_i+\cos\theta)\vec{v}_\theta$$
où $a_1=1,a_2=3$. Le vecteur directeur de $N_i(\theta)$ est donc
$$\vec{n_i}(\theta):=(a_i+\cos\theta)\vec{u}_\theta+\sin\theta\,\vec{v}_\theta$$ 
et 
$M\in N_i(\theta)\Longleftrightarrow \exists t\in\R\ |\ \overrightarrow{OM}=\overrightarrow{OM_i(\theta)}+t\cdot\vec{n_i}(\theta)$.
Finalement, 
$$\forall \theta\not=\pi,\ N_i(\theta)
=\big\{(1+t)(a_i+\cos\theta)\vec{u}_\theta+t\sin\theta\,\vec{v}_\theta\ |\ t\in\R\big\}$$
et le résultat s'étend au cas $\theta=\pi$ pour $i=2$.
\begin{itemize}
\item Si $\theta=\pi$, $N_1(\pi)=(Oy)$ et $N_2(\pi)=\{2(1+t)\vec{u}_\pi\ |\ t\in\R\}=(Ox)$, ces deux droites s'intersectent en $O$.
\item Si $\theta\not=\pi$, $N_1(\theta)$ et $N_2(\theta)$ sont sécantes si et seulement si les vecteurs $\vec{n_1}(\theta)$ et $\vec{n_2}(\theta)$ ne sont pas colinéaires, c'est-à-dire si et seulement si $\begin{array}{|cc|}1+\cos\theta&3+\cos\theta\\ \sin\theta&\sin\theta\end{array}\not=0$ \textsl{i.e.} $\sin\theta\not=0$. Ainsi, pour $\theta\not=0$, les droites $N_1(\theta)$ et $N_2(\theta)$ sont sécantes en un point $P(\theta)$, que l'on peut déterminer: 
\begin{eqnarray*}
\lefteqn{(1+t_1)(1+\cos\theta)\vec{u}_\theta+t_1\sin\theta\,\vec{v}_\theta=(1+t_2)(3+\cos\theta)\vec{u}_\theta+t_2\sin\theta\,\vec{v}_\theta}\\
&\Longleftrightarrow& \begin{cases}(1+t_1)(1+\cos\theta)=(1+t_2)(3+\cos\theta)\cr t_1\sin\theta=t_2\sin\theta\end{cases}\\
&\Longleftrightarrow& \begin{cases}t_1=-1\cr t_1=t_2\end{cases}
\end{eqnarray*}
puisqu'ici $\sin\theta\not=0$. 
On obtient alors $\overrightarrow{OP(\theta)}=-\sin\theta\,\vec{v}_\theta$. 
\end{itemize}

La formule donnant $P(\theta)$ pour $\theta\not=\pi$ est en fait encore valable en $\theta=\pi$ puisqu'on retrouve dans ce cas $P(\pi)=O$. En coordonnées cartésiennes, on a donc
$$\forall\theta\in]0;2\pi[,\ P(\theta)=(\sin^2\theta, -\sin\theta\cos\theta)$$
Remarquons que 
\begin{eqnarray*}
(\sin^2\theta, -\sin\theta\cos\theta)&=&\left(\frac{1-\cos(2\theta)}{2},-\frac{1}{2}\sin(2\theta)\right)\\
 &=&\left(\frac{1}{2},0\right)-\frac{1}{2}\big(\cos(2\theta),\sin(2\theta)\big)
\end{eqnarray*}
Lorsque $\theta$ décrit $]0;2\pi[$, $2\theta$ décrit $]0;4\pi[$ et 
$\big(\cos(2\theta),\sin(2\theta)\big)$ décrit (deux fois, sauf en $(1,0)$ où l'on ne passe 
qu'une fois) le cercle unité. Par conséquent $\{P(\theta)\ |\ \theta\in]0;2\pi[\}$ 
est le cercle de centre $(\frac{1}{2},0)$ et de rayon $\frac{1}{2}$.

\begin{center}
\begin{tikzpicture}[scale=1.5]

% Axes
     \draw[->,>=latex,thick, gray] (-2.5,0)--(5,0) node[below,black] {$x$};
     \draw[->,>=latex,thick, gray] (0,-3.5)--(0,4) node[right,black] {$y$};  

 % Ticks
        \foreach \x in {1,...,4}
                \draw (\x,2pt) -- (\x,-2pt)
                        node[anchor=north] {\x};
        \foreach \x in {-1,-2}
                \draw (\x,2pt) -- (\x,-2pt)
                        node[anchor=north] {\x};
        \foreach \x in {1,2,3}
                \draw (2pt,\x) -- (-2pt,\x)
                        node[anchor=east] {\x};
        \foreach \x in {-1,-2,-3}
                \draw (2pt,\x) -- (-2pt,\x)
                        node[anchor=east] {\x};



% Courbe

\draw [very thick, color=red, domain=0:2*pi, samples=100, smooth]
  plot (xy polar cs:angle=\x r, radius={1+cos(\x r)});


\draw [very thick, color=red, domain=0:2*pi, samples=100, smooth]
  plot (xy polar cs:angle=\x r, radius={3+cos(\x r)});

\draw [very thick, color=green!60!black] (0.5,0) circle (0.5cm);


  \coordinate (M1) at (45:1.707);
  \fill[blue] (M1) circle (2pt)  node[above left=5pt]{$M_1(\theta)$}; 
  \draw[very thick, gray] (M1)--+(-23:3)--+(-23:-3);
  \draw[very thick, blue] (M1)--+(90-23:3)--+(90-23:-3);

  \coordinate (M2) at (45:3.707);
  \fill[blue] (M2) circle (2pt)  node[right=5pt]{$M_2(\theta)$}; 
  \draw[very thick, gray] (M2)--+(-34:3)--+(-34:-3);
  \draw[very thick, blue] (M2)--+(90-34:2)--+(90-34:-4.5);

   \fill [color=black]  (0.5,-0.5) circle (2pt) node[below right] {$P(\theta)$};

  \node at (-30:2.3) {$\mathcal{C}_1$};
  \node at (-30:4.3) {$\mathcal{C}_2$};
\end{tikzpicture}  
\end{center}

\fincorrection
\correction{004994}
\begin{enumerate}
  \item $y = ae^{bx}$.
  \item $y = \pm\sqrt{ax+b}$.
  \item $(a-x)^2 + y^2 = b^2$.
  \item $x = a\left(\ln\left|\tan\frac t2\right| + \cos t\right) + b,\quad
              y = a\sin t$.
  \item $y^2 = \frac{x^2}2 + a^2\ln|x| + b$.
\end{enumerate}
 \fincorrection
\correction{004995}
\begin{enumerate}
  \item $\rho = \frac1{a\theta+b}$.
  \item $\rho = a\theta+b$. (Spirale d'Archimède)
\end{enumerate}
 \fincorrection
\correction{004996}
$D = Ox  \Rightarrow  x_T = x - \frac{x'y}{y'}$, $x_N = x + \frac{yy'}{x'}
 \Rightarrow  2x + y\left(t-\frac1t\right) = a$ (cste).

On dérive : $2x' + y'\left(t-\frac1t\right) + y\left(1+\frac1{t^2}\right) = 0
 \Rightarrow  y'\left(t+\frac1t\right) + y\left(1+\frac1{t^2}\right) = 0$.
 
$ \Rightarrow  y = \frac\lambda t$, $x = b + \frac\lambda{2t^2}$ (Parabole)
\fincorrection
\correction{004997}
$D = Ox  \Rightarrow  x_T = x - \frac{x'y}{y'}$, $x_N = x + \frac{yy'}{x'}
 \Rightarrow  y\left(t+\frac1t\right) = a$ (cste).

$y = \frac{at}{1+t^2}$ et $x' = ty'  \Rightarrow 
x = a\left(\ln\frac t{\sqrt{1+t^2}} + \frac 1{1+t^2}\right)+b$.
\fincorrection
\correction{004998}
La tangente ne doit pas être parallèle à $Oy$, donc on peut paramétrer
$\mathcal{C}$ sous la forme : $y=f(x)$, ce qui donne l'équation :
$$|x+yy'| = |y|\sqrt{1+y'^2} \Leftrightarrow
2xyy' = y^2-x^2.$$
(équation homogène) on obtient : $y = \pm\sqrt{\lambda x - x^2}$. Les courbes
cherchées sont des arcs de cercles centrés sur $Ox$ passant par $O$.
\fincorrection
\correction{004999}
             On suppose que la droite est $Ox$ et on paramètre la courbe
             cherchée, $\mathcal{C}$, par une abscisse curviligne $s$.
	     Soient $M=(x,y) \in \mathcal{C}$,         $I=(x-R\frac{dy}{ds}, y+R\frac{dx}{ds})$
	     le centre de courbure en $M$ où $R$ est le rayon de courbure.
             On veut $|R| = \Bigl|y+R\frac{dx}{ds}\Bigr| = |y+R\cos\varphi|$ d'où~:
             $$\pm \frac{d R }{ds} = \frac{dy}{ds} -R\sin\varphi \frac{d \varphi}{ds} +  \frac{d R }{ds}\cos\varphi = \frac{d R }{ds}\cos\varphi.$$
             Ceci implique $\frac{d R }{ds} = 0$ donc $R$ est constant (cercle) ou
             $\varphi\equiv 0 \bmod \pi$ (droite horizontale). Le deuxième cas est
             exclu (courbe birégulière) donc il reste le cas d'un cercle qui
             convient s'il est tangent à $Ox$.
\fincorrection
\correction{005000}
{\def\p{\frac\varphi2}%
         $y + 2R\cos\varphi = \pm2R  \Rightarrow 
          2\frac{d R}{ d\varphi}\sin\p + R\cos\p = 0$ ou
         $2\frac{d R}{ d\varphi}\cos\p - R\sin\p = 0$.\par
          cas 1 : $R = \frac K{\sin\varphi/2}$,
                  $x = 2K\ln\left|\tan\frac\varphi4\right| - 4K\cos\p + L$,
                  $y = 4K\sin\p$.\par
          cas 2 : $R = \frac K{\cos\varphi/2}$,
                  $x = -2K\ln\left|\tan\frac{\varphi+\pi}4\right| + 4K\sin\p + L$,
                  $y = -4K\cos\p$.\par}
\fincorrection
\correction{005001}
$M + a\vec t \in Ox$ (tractrices)
	     $x = a\cos \varphi + a\ln|\tan \varphi/2| + b$, $y = a\sin\varphi$.
\fincorrection
\nocorrection
\correction{005003}
\begin{enumerate}
  \item $x =  \int_{}^s \cos\ln|t|\,d t,\quad
              y =  \int_{}^s \sin\ln|t|\,d t,\quad
              \rho = \pm\frac1{\sqrt2}e^{\theta+\pi/4}$.
  \item $x =  \int_0^s \cos\frac{u^2}2\,du,\quad
              y =  \int_0^s \sin\frac{u^2}2\,du$
             (Clothoïde ou spirale de Cornu)
  \item $x = a(\varphi\sin\varphi + \cos\varphi),\quad
              y = a(-\varphi\cos\varphi + \sin\varphi)$.
  \item $x = \ln\Bigl|\tan\Bigr(\frac\varphi2+\frac\pi4\Bigr)\Bigr|,\quad
              y = \frac1{\cos\varphi} = \ch x$.
  \item $x=\frac a4(\sin2\varphi + 2\varphi),\quad
             y = \frac a4\cos2\varphi$ (cycloïde).
\end{enumerate}
 \fincorrection
\correction{005004}
Développante de cercle :
         $\frac{ d\vec I}{d s} = \frac{d R}{d s}\vec N  \Rightarrow 
          \frac{d R}{ d\varphi} = r$.\par
         $ \Rightarrow 
            x = x_0 + r(\cos\varphi-1+\varphi\sin\varphi),\quad
            y = y_0 + r(\sin\varphi-\varphi\cos\varphi)$.
\fincorrection
\correction{005005}
$y = \frac s2\sin\varphi  \Rightarrow  s = a\sin\varphi \Rightarrow 
         x = \frac {a\sin2\varphi}4 + \frac {a\varphi}2 + b,\quad
         y = \frac {a\sin^2\varphi}2$. (cycloïde)
\fincorrection
\correction{005006}
Soit $\theta$ l'angle polaire de $\vec{OM}$ :
$\vec{MC} = \frac{d s}{ d\varphi}\vec n$ et
$\vec{MN} = \frac{d s}{ d\theta}\vec n$.

$\frac{d s}{ d\theta} = k\frac{d s}{ d\varphi}  \Rightarrow 
\varphi = \frac\theta k+b  \Rightarrow  V = a\theta+b$ avec $a = \frac1k-1$.

$\frac\rho{\rho'} = \tan(a\theta+b)  \Rightarrow  \rho = \lambda\cos(a\theta+b)^{-1/a}$
si $a \ne 0$ ou $\rho = \lambda e^{\mu\theta}$ si $a=0$.

\leavevmode\vbox{\halign{$k=#$\hfil &$ \Rightarrow $ #\hfil\cr
1       & Spirale logarithmique.\cr
\frac23 & Parabole de foyer $O$.\cr
2       & Cardioïde.\cr
\frac13 & Hyperbole de centre $O$.\cr
-1      & Lemniscate de Bernouilli.\cr}}
\fincorrection
\nocorrection
\correction{005007}
\begin{enumerate}
\item
$$
\centerline{\includegraphics[height=4cm]{../images/img005007-1}}
$$
\item
$$
\centerline{\includegraphics[height=4cm]{../images/img005007-2}}
$$
\item
$$
\centerline{\includegraphics[height=4cm]{../images/img005007-3}}
$$
asymptotes : $y=\pm x$                                  \par
$x-y \sim 1/(4x)  \Rightarrow $ aire infinie                      \par
$y'=0 \Leftrightarrow t=0,(\sqrt6\pm\sqrt2)/2$                     \par
\item
$$
\centerline{\includegraphics[height=4cm]{../images/img005007-4}}
$$
asymptote : $y=\frac x2 - \frac 3{16}$ (traversée)     \par
\item
$$
\centerline{\includegraphics[height=4cm]{../images/img005007-5}}
$$
asymptote : $y=x$                                        \par
\item
$$
\centerline{\includegraphics[height=4cm]{../images/img005007-6}}
$$
inflexions : $\tan\frac t2 = 0,\pm3$                    \par
\item
$$
\centerline{\includegraphics[height=4cm]{../images/img005007-7}}
$$
hyperbole : $(y+2)^2-x^2=4$                              \par
\item
$$
\centerline{\includegraphics[height=4cm]{../images/img005007-8}}
$$
asymptote : $x+y=-1$                                     \par
équation cartésienne : $x^3+y^3=3xy$                     \par
\item
$$
\centerline{\includegraphics[height=4cm]{../images/img005007-9}}
$$
asymptote : $x+y=e^{-1}$                                 \par
branche parabolique horizontale
\item
$$
\centerline{\includegraphics[height=4cm]{../images/img005007-10}}
$$
branche parabolique horizontale                          \par
rebroussement pour $t=1$                                 \par
\item
$$
\centerline{\includegraphics[height=4cm]{../images/img005007-11}}
$$
branche parabolique de coefficient 1                     \par
\item
$$
\centerline{\includegraphics[height=4cm]{../images/img005007-12}}
$$
 asymptote : $y=2x+\frac 32$                              \par
 point double : $t^2+t=1$, $x=y=-1$
 les tangentes sont orthogonales
\end{enumerate}
\fincorrection
\correction{005008}
\begin{enumerate}
\item 
$$
\centerline{\includegraphics[height=4cm]{../images/img005008-1}}
$$
branche parabolique horizontale
\item
$$
\centerline{\includegraphics[height=4cm]{../images/img005008-2}}
$$
branche parabolique de coefficient 1
\item
$$
\centerline{\includegraphics[height=4cm]{../images/img005008-3}}
$$
inflexion : $\cos t = \frac 23$
\item
$$
\centerline{\includegraphics[height=4cm]{../images/img005008-4}}
$$
rebroussement : $\cos^2t=\frac13$.
\item
$$
\centerline{\includegraphics[height=4cm]{../images/img005008-5}}
$$
\end{enumerate}

\fincorrection
\correction{005009}
$x=\frac{3\cos\theta-\cos3\theta}4$,
	     $y=\frac{3\sin\theta-\sin3\theta}4$.
\fincorrection
\correction{005010}
         $M = (R{\cos\theta},R{\sin\theta})$,  $S=(a,0)$ :\par
         On obtient les équations paramétriques :
         $x=\frac{R(R{\cos\theta}-a)}{R-a{\cos\theta}}$, $y=\frac{(R^2-a^2){\sin\theta}}{R-a{\cos\theta}}$.\par
         Pour $R\ne a$, il s'agit de la conique de centre $O$ et d'équation cartésienne :
         $\frac{x^2}{R^2}+\frac{y^2}{R^2-a^2} = 1$.
\fincorrection
\correction{005011}
$y_A = t  \Rightarrow  \begin{cases} 2px = t^2 + ht + h^2/2 \cr y = t + h/2\cr \end{cases}
                   \Rightarrow  \text{parabole } y^2 + h^2/4 = 2px$.
\fincorrection
\correction{005012}
$y_A = t, y_B = u  \Rightarrow  C:\left(\frac{ut}{2p},\frac{t+u}2\right)$;
         aire ${}=\frac{|u-t|^3}{8p}$.\par
         enveloppe : $M = \text{mil}(A,B)$, parabole $y^2+a^2/4 = 2px$.
\fincorrection
\correction{005013}
\begin{enumerate}
  \item $F$.
  \item $M = (t^2/2p,t)$, $M' = (t'^2/2p,t')  \Rightarrow  tt' = -p^2  \Rightarrow 
    \begin{cases} x = \frac p4\left(u^2+\frac1{u^2}\right)\cr
           y = \frac p2\left(u-\frac1u\right)\end{cases}$
    avec $u = \frac tp$. (Parabole passant par $F$)
  \item $\begin{cases} x = \frac{3p}4\left(u^2+\frac1{u^2}\right)\cr
                     y = -\frac p4\left(u-\frac1u\right)^3.\cr\end{cases}$
\end{enumerate}
 \fincorrection
\correction{005014}
\begin{enumerate}
  \item Foyer.
  \item Point d'impact : $\left(\frac{t^2}{2p},t\right)$
             Point caractéristique :
             $\left(\frac{3t^2}{2p},\frac{t(3p^2-t^2)}{2p^2}\right)$.
\end{enumerate}
 \fincorrection
\correction{005015}
$M = (t^2/2p,t)  \Rightarrow  I = (3t^2/2p+p,-t^3/p^2)$.
         Soit $P = (u^2/2p,u)$ :

         $IP=IM \Leftrightarrow (u-t)^3(u+3t) = 0  \Rightarrow  u = -3t$.

         Enveloppe : $\begin{cases}x = -3t^2/2p\cr y = 3t.\cr\end{cases}$ (Parabole)
\fincorrection
\correction{005016}
 équation polaire : $\rho = \frac p{1+e\cos\theta}  \Rightarrow 
\begin{cases}x = \frac {p(\cos\theta-\sin\theta)}{2 + e(\cos\theta-\sin\theta)}\cr
        y = \frac {p(\cos\theta+\sin\theta)}{2 + e(\cos\theta-\sin\theta)}\cr \end{cases}$
	 conique d'excentricité $\frac e{\sqrt2}$.
\fincorrection
\correction{005017}
$3x = \cos2\theta + 2\cos\theta$, $3y = \sin2\theta + 2\sin\theta$ :
         cardioïde à rebroussement en $(-1/3,0)$.
\fincorrection
\correction{005018}
$D=Ox$, rayon${}=1$ :
         $x = \theta -\cos\theta\sin\theta$, $y = \sin^2\theta$.
         pt caractéristique = projeté de $I$.
\fincorrection
\correction{005019}
$M = \begin{pmatrix}a(1+\cos t)\cr a\sin t\cr\end{pmatrix}  \Rightarrow 
          P = \begin{pmatrix}2a\cos t(1+\cos t)\cr 2a\sin t(1-\cos t)\cr\end{pmatrix}$.
          Hypocycloïde à trois rebroussements.
\fincorrection
\correction{005020}
$x = a\cos 4\theta$, $y = a\sin 4\theta$.
\fincorrection
\correction{005021}
\begin{enumerate}
  \item $x = \frac{\cos t}a(a^2 + (a^2-b^2)\sin^2t)$,
             $y = \frac{\sin t}b(b^2 - (a^2-b^2)\cos^2t)$.
  \item 
  \item Point stationnaire ssi $a^2 > 2b^2$, obtenu pour
             $\sin^2t = \frac{a^2-2b^2}{3(a^2-b^2)}$. Rebroussement de $1^{\text{ère}}$ espèce.
\end{enumerate}
 \fincorrection
\correction{005022}
$D = Ox$, $A = (0,a)$,$M = (t,0)   \Rightarrow  \begin{cases} x = 2t\cr y = t^2/a.\cr\end{cases}$ (Parabole)
\fincorrection
\nocorrection
\nocorrection
\nocorrection
\nocorrection
\nocorrection
\nocorrection
\nocorrection
\correction{005024}
\begin{enumerate}
  \item $\sh^2t$.
  \item $\theta-\tanh\frac\theta2$.
\end{enumerate}
\fincorrection
\correction{005025}
$2-\sqrt2+3\ln(1+\sqrt2)$.
\fincorrection
\correction{005026}
\begin{enumerate}
  \item $4\sqrt2+4\arccos\frac1{\sqrt3}-\pi = 4\sqrt2+4\Arctan\sqrt2-\pi$.
  \item $4$.
\end{enumerate}
 \fincorrection
\correction{005027}
$\sqrt x + \sqrt y = 1  \Rightarrow  2\sqrt{xy}=1-x-y \Rightarrow (x-y)^2=2(x+y)-1$.
La courbe est un arc de parabole d'axe la première bissectrice et tangent aux axes en $(1,0)$ et en $(0,1)$.

Longueur~: $x-y=\sh t$, $2(x+y)=\ch^2t$ $ \Rightarrow $ $x=\frac12\ch^2t+\frac14\sh t$, $y=\frac12\ch^2t-\frac14\sh t$,
$\sqrt{x'^2+y'^2}=\frac{\ch^2t}{\sqrt2}$.

$L= \int_{t=-\Argsh 1}^{\Argsh1}\frac{\ch^2td t}{\sqrt2}=\frac{\ln(1+\sqrt2)}{\sqrt2}+1$.
\fincorrection
\correction{005028}
Soit $(a_i)$ une subdivision de $[a,b]$ et $P$ la ligne brisée passant
par les points $(a_i, f(a_i))$. On montre ci-dessous que pour toute courbe
rectifiable $L$ située au dessus de $P$ et ayant même extrémités,
on a $\mathrm{long}(L) \ge \mathrm{long}(P)$ (résultat intuitivement évident~:
planter des clous aux points $(a_i,f(a_i))$ et attacher un élastique en
$(a,f(a))$ et $(b,f(b))$, passant au dessus de ces clous).
Cela étant montré, l'inégalité demandée en résulte en faisant tendre le pas
de la subdivision vers zéro.

Démonstration du thm de l'élastique~: par récurrence sur le nombre $n$ de
segments de $P$. Pour $n=1$ c'est un fait connu.
$n-1  \Rightarrow  n$~: si $L$ passe par $(a_1,f(a_1))$ alors l'hypothèse de récurrence
s'applique. Sinon, notons $D$ la demi-droite issue de $(a_0,f(a_0))$ et
passant par $(a_1,f(a_1))$. Par concavité, $P$ est en dessous de $D$. $L$
contient un point d'abscisse $a_1$ strictement au dessus de $D$, et aboutit
en $(b,f(b))$ en dessous de $D$, donc il existe un point $(u,v)$ sur $L \cap D$
avec $u > a_1$. En rempla\c cant l'arc $(a_0,f(a_0))$ -- $(u,v)$ de $L$ par le
segment correspondant on obtient une ligne $L'$ plus courte que $L$,
encore au dessus de $P$, et qui relève du premier cas.
\fincorrection
\correction{005029}
\begin{enumerate}
  \item $x=-4t^3$, $y=\frac{3+6t^2-3t^4}2$.
  \item $x=6\cos t-3\cos2t$,\quad $y=6\sin t+3\sin 2t$.
  \item $x=t+\sin t$,\quad $y=-1+\cos t$, $I_t = M_{t-\pi} + (\pi,-2)$.
  \item $x=a(\cos^3t + 3\cos t\sin^2t)$,\quad
             $y=a(\sin^3t + 3\sin t\cos^2t)$.
             $x\pm y = a(\cos t\pm \sin t)^3  \Rightarrow  $ similitude de centre $O$,
             rapport $2$, angle $\frac\pi4$.
  \item $x_I = \frac{3x^4+1}{2x^3}$,\quad $y_I = \frac{x^4+3}{2x}$.
  \item $x=\left(a-\frac {b^2}a\right)\cos^3 t$,\quad
             $y=\left(b-\frac {a^2}b\right)\sin^3 t$.
  \item $\rho = e^{\theta-\pi/2}$.
  \item $x=\frac{2+\cos\theta-\cos^2\theta}3$,\quad
             $y=\frac{\sin\theta(1-\cos\theta)}3$,\quad
             cardioïde homothétique.
\end{enumerate}
 \fincorrection
\correction{005030}
Calcul.
\fincorrection
\nocorrection
\nocorrection
\correction{005033}
\begin{enumerate}
  \item 
  \item Cercle de centre $(\frac12,0)$ et de rayon $\frac12$.
\end{enumerate}
 \fincorrection
\correction{005034}
$R=\frac12$ aux sommets principaux $(0,\pm1)$ et $R=\sqrt2$ aux sommets secondaires $(\pm1/\sqrt2,0)$.
\fincorrection
\correction{005035}
$-\frac{156}{125\sqrt2}$.
\fincorrection
\correction{005036}
$y'(0)=\lambda  \Rightarrow  I = (-\lambda-\lambda^3,1+\lambda^2)$.
\fincorrection
\nocorrection
\correction{005038}
$c_i = \frac1{\sqrt{1+a^2c^2}}\left(\frac{ac'}{1+a^2c^2} \pm c\right)$
où $c$ est la courbure en $M$ et $c' = \frac{d c}{d s}$.

Dans le repère de Frenet, les normales ont pour équations :
$X=0$, $\pm X = acY - a$, donc se coupent en $C$.
\fincorrection
\correction{005039}
\begin{enumerate}

  \item Soit $\theta$ l'angle polaire de $D$. Dans le repère
    $(O,\vec u_\theta, \vec v_\theta)$, $\cal P$ a pour équation :
    $Y = aX^2 + bX$.

    On veut que $\cal P$ soit tangente à $Oy$, soit $b = -\tan\theta$ et que
    le rayon de courbure soit $R$, soit $a = \frac1{2R\cos^3\theta}$.

    \'Equation dans $0xy$ : $x^2\sin^2\theta - 2xy\cos\theta\sin\theta
                             + y^2\cos^2\theta - 2Rx\cos^2\theta = 0$. 
  \item $O$.
\end{enumerate}
 \fincorrection
\correction{005040}
\begin{enumerate}
  \item 
  \item $x=\frac{4t^3 + a(1-t^2)}{1+t^2}$,\quad
             $y=\frac{t^4-3t^2+2at}{1+t^2}$.
\end{enumerate}
 \fincorrection
\correction{005041}
$x=a\ln\left|\tan\left(\frac t2+\frac\pi4\right)\right|$,\quad
         $y=a\cos t$.
\fincorrection
\correction{005535}
\begin{enumerate}
 \item  L'astroïde complète est obtenue quand $t$ décrit $[-\pi,\pi]$ et pour des raisons de symétrie, $L=4\int_{0}^{\pi}{2}\left\|\overrightarrow{\frac{dM}{dt}}\right\| dt$.
Or $\overrightarrow{\frac{dM}{dt}}=\left(\begin{array}{c}
-3a\sin t\cos^2t\\
3a\cos t\sin^2t
\end{array}
\right)=3a\sin t\cos t\left(\begin{array}{c}
-\cos t\\
\sin t
\end{array}
\right)$ et donc $\left\|\overrightarrow{\frac{dM}{dt}}\right\|=3a|\sin t\cos t|=\frac{3a}{2}|\sin(2t)|$ puis

\begin{center}
$L=4\int_{0}^{\pi/2}\left\|\overrightarrow{\frac{dM}{dt}}\right\| dt=6a\int_{0}^{\pi/2}\sin(2t)\;dt=6a\left[-\frac{\cos(2t)}{2}\right]_0^{\pi/2}=6a$.
\end{center}

\begin{center}
\shadowbox{
$L=6a$.
}
\end{center}

 \item   $\overrightarrow{\frac{dM}{dt}}=\left(\begin{array}{c}
R(1-\cos t)\\
R\sin t
\end{array}
\right)=2R\sin\left(\frac{t}{2}\right)\sin t\cos t\left(\begin{array}{c}
\sin(t/2)\\
\cos(t/2)
\end{array}
\right)$ et donc $\left\|\overrightarrow{\frac{dM}{dt}}\right\|=2R\left|\sin\left(\frac{t}{2}\right)\right|$ puis

\begin{center}
$L=\int_{0}^{2\pi}\left\|\overrightarrow{\frac{dM}{dt}}\right\| dt=2R\int_{0}^{2\pi}\sin\left(\frac{t}{2}\right)\;dt=4R\left[-\cos\left(\frac{t}{2}\right)\right]_0^{2\pi}=8R$.
\end{center}

\begin{center}
\shadowbox{
$L=8R$.
}
\end{center}
 \item  Une représentation paramétrique de $\Gamma$ est $\left\{
\begin{array}{l}
x=t\\
y=\frac{t^2}{2p}
\end{array}
\right.$, $0\leqslant t\leqslant a$ et donc

\begin{align*}\ensuremath
L&=\int_{0}^{a}\sqrt{(x'(t))^2+(y'(t))^2}\;dt=\int_{0}^{a}\sqrt{1+\frac{t^2}{p^2}}\;dt=p\int_{0}^{a/p}\sqrt{u^2+1}\;du\\
 &=p\left(\left[u\sqrt{u^2+1}\right]_0^{a/p}-\int_{0}^{a/p}\frac{u^2}{\sqrt{u^2+1}}\;du\right)=a\sqrt{1+\frac{a^2}{p^2}}-p\int_{0}^{a/p}\frac{u^2+1-1}{\sqrt{u^2+1}}\;du\\
 &=a\sqrt{1+\frac{a^2}{p^2}}-L+p\Argsh\left(\frac{a}{p}\right),
\end{align*}
et donc

\begin{center}
\shadowbox{
$L=\frac{1}{2}\left(a\sqrt{1+\frac{a^2}{p^2}}+p\Argsh\left(\frac{a}{p}\right)\right)$.
}
\end{center}
 \item  La cardioïde complète est obtenue quand $\theta$ décrit $[-\pi,\pi]$.

\begin{center}
$\overrightarrow{\frac{dM}{d\theta}}=a\left((-\sin\theta)\overrightarrow{u}_\theta+(1+\cos\theta)\overrightarrow{v}_\theta\right)=2a\cos\left(\frac{\theta}{2}\right)\left(-\sin\left(\frac{\theta}{2}\right)\overrightarrow{u}_\theta+\cos\left(\frac{\theta}{2}\right)\overrightarrow{v}_\theta\right)$.
\end{center}
Comme le vecteur $-\sin\left(\frac{\theta}{2}\right)\overrightarrow{u}_\theta+\cos\left(\frac{\theta}{2}\right)\overrightarrow{v}_\theta$ est unitaire, $\left\|\overrightarrow{\frac{dM}{d\theta}}\right\|=\left|2a\cos\left(\frac{\theta}{2}\right)\right|$ puis

\begin{center}
$L=\int_{-\pi}^{\pi}\left|2a\cos\left(\frac{\theta}{2}\right)\right|dt=4a\int_{0}^{\pi}\cos\left(\frac{\theta}{2}\right)dt=8a\left[\sin\left(\frac{\theta}{2}\right)\right]_0^{\pi}=8a$.
\end{center}

\begin{center}
\shadowbox{
$L=8a$.
}
\end{center}
\end{enumerate}
\fincorrection
\correction{005536}
On obtient la courbe complète quand $t$ décrit $]-\pi,0[\cup]0,\pi[$. Puisque $M(-t)=s_{(Ox)}(M(t))$ et $M(\pi-t)=s_{(Oy)}(M(t))$, on se contente d'étudier et de construire la courbe quand $t\in\left]0,\frac{\pi}{2}\right[$ puis on obtient la courbe complète par réflexions successives d'axe $(Oy)$ puis d'axe $(Ox)$.
Pour $t\in\left]0,\frac{\pi}{2}\right[$,

\begin{center}
$\overrightarrow{\frac{dM}{dt}}=\left(
\begin{array}{c}
R(-\sin t+\frac{1}{\sin t})\\
R\cos t
\end{array}
\right)=R\left(
\begin{array}{c}
\frac{\cos^2t}{\sin t}\\
\cos t
\end{array}
\right)=R\frac{\cos t}{\sin t}\left(
\begin{array}{c}
\cos t\\
\sin t
\end{array}
\right)=R\cotan t\left(
\begin{array}{c}
\cos t\\
\sin t
\end{array}
\right).$
\end{center}
Puisque $R\cotan t>0$ pour $t\in\left]0,\frac{\pi}{2}\right[$ et puisque le vecteur $\left(
\begin{array}{c}
\cos t\\
\sin t
\end{array}
\right)$ est unitaire, on a

\begin{center}
$\frac{ds}{dt}=R\cotan t$ puis $\overrightarrow{\tau}(t)=\left(
\begin{array}{c}
\cos t\\
\sin t
\end{array}
\right)$.
\end{center}
On a donc $\overrightarrow{n}(t)=\left(
\begin{array}{c}
-\sin  t\\
\cos t
\end{array}
\right)$ et d'autre part, on peut prendre $\alpha(t)=t$. En notant $\rho(t)$ le rayon de courbure au point $M(t)$,

\begin{center}
$\rho(t)=\frac{ds}{d\alpha}=\frac{ds/dt}{d\alpha/dt}=R\cotan t$,
\end{center}
puis

\begin{align*}\ensuremath
\Omega(t)&=M(t)+\rho(t)\overrightarrow{n}(t)=\left(
\begin{array}{c}
R\left(\cos t+\ln\left|\tan \frac{t}{2}\right|\right)\\
R\sin t
\end{array}
\right)+R\cotan t\left(
\begin{array}{c}
-\sin  t\\
\cos t
\end{array}
\right)\\
 &=\left(
\begin{array}{c}
R\ln\left|\tan \frac{t}{2}\right|\\
\frac{R}{\sin t}
\end{array}
\right).
\end{align*}
La développée cherchée est l'arc $t\mapsto\left(
\begin{array}{c}
R\ln\left|\tan \frac{t}{2}\right|\\
\frac{R}{\sin t}
\end{array}
\right)$, $t\in]-\pi,0[\cup]0,\pi[$ (en complétant par symétrie). Quand $t$ décrit $]0,\pi[$, on effectue alors le changement de paramètres $t\mapsto R\ln\left|\tan \frac{t}{2}\right|=u$ qui est un $C^1$-difféomorphisme de $]0,\pi[$ sur $\Rr$. On obtient $x=u$ puis

\begin{center}
$y=\frac{R}{\frac{2\tan\frac{t}{2}}{1+\tan^2\frac{t}{2}}}=\frac{R}{2}\left(\tan\frac{t}{2}+\frac{1}{\tan\frac{t}{2}}\right)=R\frac{e^{u/R}+e^{-u/R}}{2}=R\ch\left(\frac{u}{R}\right)$.
\end{center}
Le support de la développée sur $]0,\pi[$ est aussi le support de l'arc $u\mapsto\left(
\begin{array}{c}
u\\
R\ch\left(\frac{u}{R}\right)
\end{array}
\right)$, $u\in\Rr$ ou encore la chaînette d'équation cartésienne $y=R\ch\left(\frac{x}{R}\right)$.

$$\includegraphics{../images/img005536-1}$$

 \item  Quand $t$ décrit $[0,2\pi]$, on obtient une arche de cycloïde complète. Les autres arches s'en déduisent par translations de vecteurs $2k\pi R\overrightarrow{i}$. Pour $t\in[0,2\pi]$

\begin{center}
$\overrightarrow{\frac{dM}{dt}}=\left(
\begin{array}{c}
R(1-\cos t)\\
R\sin t
\end{array}
\right)=2R\sin\left(\frac{t}{2}\right)\left(
\begin{array}{c}
\sin\left(\frac{t}{2}\right)\\
\rule{0mm}{6mm}\cos\left(\frac{t}{2}\right)
\end{array}
\right).$
\end{center}
Le point $M(t)$ est régulier pour $t\in]0,2\pi[$ et pour $t\in]0,2\pi[$, $2R\sin\left(\frac{t}{2}\right)>0$. Puisque le vecteur $\left(
\begin{array}{c}
\sin\left(\frac{t}{2}\right)\\
\rule{0mm}{6mm}\cos\left(\frac{t}{2}\right)
\end{array}
\right)$ est unitaire, on a

\begin{center}
$\frac{ds}{dt}=2R\sin\left(\frac{t}{2}\right)$ et $\overrightarrow{\tau}(t)=\left(
\begin{array}{c}
\sin\left(\frac{t}{2}\right)\\
\rule{0mm}{6mm}\cos\left(\frac{t}{2}\right)
\end{array}
\right)=\left(
\begin{array}{c}
\cos\left(\frac{\pi}{2}-\frac{t}{2}\right)\\
\rule{0mm}{6mm}\sin\left(\frac{\pi}{2}-\frac{t}{2}\right)
\end{array}
\right)$.
\end{center}
On en déduit que $\overrightarrow{n}(t)=\left(
\begin{array}{c}
-\cos\left(\frac{t}{2}\right)\\
\rule{0mm}{6mm}\sin\left(\frac{t}{2}\right)
\end{array}
\right)$ et d'autre part, on peut prendre $\alpha(t)=\frac{\pi}{2}-\frac{t}{2}$. En notant $\rho(t)$ le rayon de courbure au point $M(t)$,

\begin{center}
$\rho(t)=\frac{ds}{d\alpha}=\frac{ds/dt}{d\alpha/dt}=\frac{2R\sin\left(\frac{t}{2}\right)}{-\frac{1}{2}}=-4R\sin\left(\frac{t}{2}\right)$,
\end{center}
et donc

\begin{align*}\ensuremath
\Omega(t)&=M(t)+\rho(t)\overrightarrow{n}(t)=\left(
\begin{array}{c}
R(t-\sin t)\\
R(1-\cos t)
\end{array}
\right)-4R\sin\left(\frac{t}{2}\right)\left(
\begin{array}{c}
-\cos\left(\frac{t}{2}\right)\\
\rule{0mm}{6mm}\sin\left(\frac{t}{2}\right)
\end{array}
\right)=\left(
\begin{array}{c}
R(t-\sin t)+2R\sin t\\
R(1-\cos t)-2R(1-\cos t)
\end{array}
\right)\\
 &=\left(
\begin{array}{c}
R(t+\sin t)\\
-R(1-\cos t)
\end{array}
\right).
\end{align*}
La développée cherchée est l'arc $t\mapsto\left(
\begin{array}{c}
R(t+\sin t)\\
-R(1-\cos t)
\end{array}
\right)$. Poursuivons.

\begin{align*}\ensuremath
\Omega(t+\pi)&=\left(
\begin{array}{c}
R(t+\pi-\sin t)\\
-R(1+\cos t)
\end{array}
\right)=\left(
\begin{array}{c}
R(t-\sin t)\\
R(1-\cos t)
\end{array}
\right)+\left(
\begin{array}{c}
\pi R\\
-2R
\end{array}
\right)=t_{\overrightarrow{u}}(M(t))\;\text{où}\;\overrightarrow{u}=(\pi R,-2R).
\end{align*}
Ainsi, le centre de courbure au point $M(t+\pi)$ est le translaté du point $M(t)$ dans la translation de vecteur $(\pi R,-2R)$ et donc la développée de la cycloïde est la translatée de la cycloïde par la translation de vecteur $(\pi R,-2R)$. En particulier, c'est encore une cycloïde.

$$\includegraphics{../images/img005536-2}$$

 \item  $\mathcal{C}$ est le support de la courbe paramétrée $t\mapsto M(t)=\left(
\begin{array}{c}
t\\
t^3
\end{array}
\right)$. $M(t)$ est birégulier si et seulement si $t/neq0$. Pour $t\in\Rr$, $\overrightarrow{\frac{dM}{dt}}=\left(
\begin{array}{c}
1\\
3t^2
\end{array}
\right)$. Par suite 

\begin{center}
$\frac{ds}{dt}=\sqrt{1+9t^4}$ et $\overrightarrow{\tau}(t)=\frac{1}{\sqrt{1+9t^4}}\left(
\begin{array}{c}
1\\
3t^2
\end{array}
\right)$.
\end{center}
Donc, d'une part $\overrightarrow{n}(t)=\frac{1}{\sqrt{1+9t^4}}\left(
\begin{array}{c}
-3t^2\\
1
\end{array}
\right)$ et d'autre part, puisque les coordonnées de $\overrightarrow{\tau}(t)$ sont positives, on peut prendre $\alpha(t)=\Arccos\left(\frac{1}{\sqrt{1+9t^4}}\right)$. Par suite, pour $t\neq0$

\begin{center}
$\frac{d\alpha}{dt}=-\left(-\frac{1}{2}\right)36t^3(1+9t^4)^{-3/2}\times\frac{1}{\sqrt{1-\frac{1}{1+9t^4}}}=\frac{6t}{1+9t^4}$
\end{center}
puis

\begin{center}
$R(t)=\frac{ds/dt}{d\alpha/dt}=\frac{(1+9t^4)^{3/2}}{6t}$,
\end{center}
et donc

\begin{center}
$\Omega(t)=M(t)+R(t)\overrightarrow{n}(t)=\left(
\begin{array}{c}
t\\
t^3
\end{array}
\right)+
\frac{1+9t^4}{6t}\left(
\begin{array}{c}
-3t^2\\
1
\end{array}
\right)
=\left(
\begin{array}{c}
\frac{t}{2}-\frac{9t^5}{2}\\
\rule{0mm}{6mm}\frac{5t^3}{2}+\frac{1}{6t}
\end{array}
\right)$.
\end{center}

$$\includegraphics{../images/img005536-3}$$

\fincorrection
\correction{005537}
$\mathcal{C}$ est le support de l'arc paramétré $t\mapsto\left(
\begin{array}{c}
t\\
\ln t
\end{array}
\right)$, $t>0$. 

\begin{center}
$\overrightarrow{\frac{dM}{dt}}=\left(\begin{array}{c}
1\\
1/t
\end{array}
\right)=\sqrt{1+\frac{1}{t^2}}\left(\begin{array}{c}
1/\sqrt{1+\frac{1}{t^2}}\\
1/\left(t\sqrt{1+\frac{1}{t^2}}\right)
\end{array}
\right)=\frac{\sqrt{1+t^2}}{t}\left(
\begin{array}{c}
\frac{t}{\sqrt{1+t^2}}\\
\frac{1}{\sqrt{1+t^2}}
\end{array}
\right)$.
\end{center}
Donc, $\frac{ds}{dt}=\frac{\sqrt{1+t^2}}{t}$ et on peut prendre $\alpha(t)=\Arcsin\left(\frac{1}{\sqrt{1+t^2}}\right)$ puis 

\begin{center}
$\frac{d\alpha}{dt}=-\frac{t}{(t^2+1)^{3/2}}\frac{1}{\sqrt{1-\frac{1}{1+t^2}}}=-\frac{1}{t^2+1}$,
\end{center}
et finalement

\begin{center}
$R(t)=\frac{ds/dt}{d\alpha/dt}=-\frac{1}{t}(t^2+1)^{3/2}$.
\end{center}
Pour $t>0$, posons $f(t)=|R(t)|=\frac{1}{t}(t^2+1)^{3/2}$. $f$ est dérivable sur $]0,+\infty[$ et pour $t>0$, 

\begin{center}
$f'(t)=-\frac{1}{t^2}(t^2+1)^{3/2}+3(t^2+1)^{1/2}=\frac{\sqrt{t^2+1}}{t^2}(-(t^2+1)+3t^2)=\frac{\sqrt{t^2+1}}{t^2}(2t^2-1)$.
\end{center}
$f$ admet un minimum en $t=\frac{1}{\sqrt{2}}$ égal à $\sqrt{2}\left(\frac{1}{2}+1\right)^{3/2}=\frac{3\sqrt{3}}{2}$.

\begin{center}
\shadowbox{
Le rayon de courbure minimum est $\frac{3\sqrt{3}}{2}$ et est le rayon de courbure en $M\left(\frac{1}{\sqrt{2}}\right)$, $-\frac{\pi}{2}<t<\frac{\pi}{2}$.
}
\end{center}
\fincorrection
\correction{005538}
 $\mathcal{C}$ est le support de l'arc paramétré $t\mapsto\left(
\begin{array}{c}
t\\
\ln(\cos t)
\end{array}
\right)$.

\begin{center}
$\overrightarrow{\frac{dM}{dt}}=\left(
\begin{array}{c}
1\\
-\sin t/\cos t
\end{array}
\right)=\frac{1}{\cos t}\left(
\begin{array}{c}
\cos t\\
-\sin t
\end{array}
\right)$.
\end{center}
Puisque $\frac{1}{\cos t}>0$ et que $\left(
\begin{array}{c}
\cos t\\
-\sin t
\end{array}
\right)$ est unitaire, on a successivement $\frac{ds}{dt}=\frac{1}{\cos t}$, $\overrightarrow{\tau}(t)=\left(
\begin{array}{c}
\cos t\\
-\sin t
\end{array}
\right)$, $\overrightarrow{n}(t)=\left(
\begin{array}{c}
\sin t\\
\cos t
\end{array}
\right)$, $\alpha(t)=-t$ puis

\begin{center}
$R(t)=\frac{ds/dt}{d\alpha/dt}=-\frac{ds}{dt}=-\frac{1}{\cos t}$.
\end{center}
Ensuite, si $s$ est l'abscisse curviligne d'origine $0$ orientée dans le sens des $t$ croissants,

\begin{center}
$s(t)=\int_{0}^{t}s'(u)\;du=\int_{0}^{t}\frac{1}{\cos u}\;du=\ln\left|\tan\left(\frac{t}{2}+\frac{\pi}{4}\right)\right|$.
\end{center}
Enfin,

\begin{center}
$\Omega(t)=M(t)+R(t)\overrightarrow{n}(t)=\left(
\begin{array}{c}
t\\
\ln(\cos t)
\end{array}
\right)-\frac{1}{\cos t}\left(
\begin{array}{c}
\sin t\\
\cos t
\end{array}
\right)=\left(
\begin{array}{c}
t-\tan t\\
\ln(\cos t)-1
\end{array}
\right)$.
\end{center}

$$\includegraphics{../images/img005538-1}$$

\fincorrection
\correction{005539}
Soit $\lambda\in\Rr$. $\mathcal{C}_\lambda$ est le support de l'arc paramétré $t\mapsto\left(
\begin{array}{c}
t\\
\lambda te^{-t}
\end{array}
\right)$. $\mathcal{C}_0$ est l'axe $(Ox)$ et donc $C_0$ n'est pas défini, puis $\mathcal{C}_{-\lambda}$ est la symétrique de $\mathcal{C}_\lambda$ par rapport à l'axe $(Ox)$ et donc $C_{-\lambda}$ est le symétrique de $C_\lambda$ par rapport à l'axe $(Ox)$. Dans ce qui suit, on suppose $\lambda>0$.

\begin{center}
$\overrightarrow{\frac{dM}{dt}}=\left(
\begin{array}{c}
1\\
\lambda(1-t)e^{-t}
\end{array}
\right)$.
\end{center}
Par suite $\frac{ds}{dt}=\sqrt{1+\lambda^2(1-t)^2e^{-2t}}$, $\overrightarrow{\tau}(t)=\frac{1}{\sqrt{1+\lambda^2(1-t)^2e^{-2t}}}\left(
\begin{array}{c}
1\\
\lambda(1-t)e^{-t}
\end{array}
\right)$,

$\overrightarrow{n}(t)=\frac{1}{\sqrt{1+\lambda^2(1-t)^2e^{-2t}}}\left(
\begin{array}{c}
-\lambda(1-t)e^{-t}\\
1
\end{array}
\right)$ et on peut prendre $\alpha(t)=\Arccos\left(\frac{1}{\sqrt{1+\lambda^2(1-t)^2e^{-2t}}}\right)$ (car $\overrightarrow{\tau}(t)$ a une abscisse strictement positive). Ensuite,

\begin{center}
$\frac{d\alpha}{dt}=\frac{\lambda^2((2t-2)-2(t-1)^2)e^{-2t}}{2(1+\lambda^2(1-t)^2e^{-2t})^{3/2}}\frac{1}{\sqrt{1-\frac{1}{1+\lambda^2(1-t)^2e^{-2t}}}}$
\end{center}
et donc $\frac{d\alpha}{dt}(0)=\frac{-4\lambda^2}{2(1+\lambda^2)^{3/2}}\frac{1}{\sqrt{1-\frac{1}{1+\lambda^2}}}=\frac{-2\lambda}{1+\lambda^2}$ puis $R(0)=\frac{ds/dt}{d\alpha/dt}(0)=-\frac{1}{2\lambda}(1+\lambda^2)^{3/2}$ et donc

\begin{center}
$C_\lambda=\Omega(0)=M(0)+R(0)\overrightarrow{n}(0)=O-\frac{1}{2\lambda}(1+\lambda^2)^{3/2}\frac{1}{\sqrt{1+\lambda^2}}\left(
\begin{array}{c}
-\lambda\\
1
\end{array}
\right)=\left(
\begin{array}{c}
(1+\lambda^2)/2\\
-(1+\lambda^2)/(2\lambda)
\end{array}
\right).$
\end{center}
L'ensemble des $C_\lambda$, $\lambda\in\Rr^*$, est le support de l'arc $\lambda\mapsto\left(
\begin{array}{c}
(1+\lambda^2)/2\\
-(1+\lambda^2)/(2\lambda)
\end{array}
\right)$, $\lambda\in\Rr^*$.

$$\includegraphics{../images/img005539-1}$$

\fincorrection
\correction{005042}
$M_1,M_2,M_3,M_4$ sont coplanaires si et seulement s'il existe
	     $a,b,c,d\in \R$ avec $(a,b,c)\ne(0,0,0)$ tels que le plan $P$
	     d'équation $ax+by+cz-d=0$ passe par ces points, ce qui équivaut
	     à : $t_1,t_2,t_3,t_4$ sont les racines (distinctes) du polynôme
	     $at^4+bt^3+ct^2-d$. Un tel polynôme existe si et seulement si
	     $t_1t_2t_3 + t_1t_2t_4 + t_1t_3t_4 + t_2t_3t_4 = 0$ soit :
	     $\frac1{t_1} + \frac1{t_2} + \frac1{t_3} + \frac1{t_4} = 0$
	     si aucun des $t_i$ n'est nul.
\fincorrection
\correction{005043}
\begin{enumerate}
  \item $\frac{d\vec I}{ds} = -\frac \tau c \vec B  \Rightarrow 
              \vec T_1 = \vec B$, $\frac {ds_1}{ds} = -\frac \tau c$,
              $\vec N_1 = -\vec N$, $c_1=c$.
  \item $\tau_1 = -\frac{c^2}{\tau}$.
\end{enumerate}
 \fincorrection
\correction{005044}
$c_1 = \sqrt{1+\frac{\tau^2}{c^2}}$,
         $\tau_1 = \frac{c\tau'-\tau c'}{c(c^2+\tau^2)}$.
\fincorrection
\correction{005045}
Pt caractéristique : $P = M + a(s)\vec N + b(s)\vec B$ :
         CNS ${}\Leftrightarrow \begin{cases} a = \frac1c \cr
                             \frac{ab'-a'b}{a^2+b^2} = \Arctan(b/a)'= \tau.\cr \end{cases}$
Rmq : le point caractéristique se projette sur $I$.
\fincorrection
\correction{005046}
$\vec B = \vec T + 2\vec k  \Rightarrow 
         s^2\frac{d^2\vec T}{ds^2} + s\frac{d\vec T}{ds} + \vec T = -\vec k$.

         on pose $s = e^u$ :
         $\vec{OM} = e^u\cos u\vec A + e^u\sin u\vec B - e^u\vec k$
         où $(\vec A,\vec B, \vec k)$ est orthogonale et
         $\|\vec A\| = \|\vec B\| = \|\vec k\|$.

         (spirale logarithmique relevée sur un cône)
\fincorrection
\nocorrection
\nocorrection
\nocorrection
\nocorrection
\nocorrection
\nocorrection
\nocorrection
\nocorrection
\nocorrection
\nocorrection
\nocorrection
\nocorrection
\nocorrection
\nocorrection
\nocorrection
\nocorrection
\nocorrection
\nocorrection
\nocorrection
\nocorrection
\nocorrection
\nocorrection
\nocorrection
\nocorrection
\nocorrection
\nocorrection
\nocorrection
\nocorrection
\nocorrection
\nocorrection
\nocorrection
\nocorrection
\nocorrection
\correction{005897}
\begin{enumerate}
 \item  Pour $(x,y)\in\Rr^2$, on pose $P(x,y)=2x+2y+e^{x+y}=Q(x,y)$. Les fonctions $P$ et $Q$ sont de classe $C^1$ sur $\Rr^2$ qui est un ouvert étoilé de $\Rr^2$. Donc, d'après le théorème de \textsc{Schwarz}, $\omega$ est exacte sur $\Rr^2$ si et seulement si $ \frac{\partial P}{\partial y}= \frac{\partial Q}{\partial x}$ et comme $ \frac{\partial P}{\partial y}=2+e^{x+y}= \frac{\partial Q}{\partial x}$, la forme différentielle $\omega$ est une forme différentielle exacte sur $\Rr^2$.

Soit $f$ une fonction $f$ de classe $C^1$ sur $\Rr^2$.

\begin{align*}\ensuremath
df=\omega&\Leftrightarrow\forall(x,y)\in\Rr^2,\;\left\{
\begin{array}{l}
 \frac{\partial f}{\partial x}(x,y)=2x+2y+e^{x+y}\\
\rule{0mm}{7mm} \frac{\partial f}{\partial y}(x,y)=2x+2y+e^{x+y}
\end{array}
\right.\\
 &\Leftrightarrow\exists g\in C^1(\Rr,\Rr)/\;\forall(x,y)\in\Rr^2,\;
\left\{
\begin{array}{l}
f(x,y)=x^2+2xy+e^{x+y}+g(y)\\
2x+e^{x+y}+g'(y)=2x+2y+e^{x+y}
\end{array}
\right.\\
 &\Leftrightarrow\exists\lambda\in\Rr/\;\forall(x,y)\in\Rr^2,\;
\left\{
\begin{array}{l}
f(x,y)=x^2+2xy+e^{x+y}+g(y)\\
g(y)=y^2+\lambda
\end{array}
\right.\\
 &\Leftrightarrow\exists\lambda\in\Rr/\;\forall(x,y)\in\Rr^2/\;f(x,y)=(x+y)^2+e^{x+y}+\lambda.
\end{align*}

Les primitives de $\omega$ sur $\Rr^2$ sont les fonctions de la forme $(x,y)\mapsto(x+y)^2+e^{x+y}+\lambda$, $\lambda\in\Rr$.

\textbf{Remarque.} On pouvait aussi remarquer immédiatement que si $f(x,y)=(x+y)^2+e^{x+y}$ alors $df=\omega$.

\item  La forme différentielle $\omega$ est de classe $C^1$ sur $\Omega=\{(x,y)\in\Rr^2/\;y > x\}$ qui est un ouvert étoilé de $\Rr^2$ car convexe. Donc, d'après le théorème de \textsc{Schwarz}, $\omega$ est exacte sur $\Omega$ si et seulement si $\omega$ est fermée sur $\Omega$.

$ \frac{\partial}{\partial x}\left( \frac{x}{(x-y)^2}\right)= \frac{\partial}{\partial x}\left( \frac{1}{x-y}+y \frac{1}{(x-y)^2}\right)=- \frac{1}{(x-y)^2}- \frac{2y}{(x-y)^3}=- \frac{x+y}{(x-y)^3}= \frac{x+y}{(y-x)^3}$.

$ \frac{\partial}{\partial y}\left(- \frac{y}{(x-y)^2}\right)= \frac{\partial}{\partial y}\left(- \frac{1}{y-x}-x \frac{1}{(y-x)^2}\right)= \frac{1}{(y-x)^2}+ \frac{2x}{(y-x)^3}= \frac{x+y}{(y-x)^3}= \frac{\partial}{\partial x}\left( \frac{x}{(x-y)^2}\right)$.

Donc $\omega$ est exacte sur l'ouvert $\Omega$. Soit $f$ une fonction $f$ de classe $C^1$ sur $\Rr^2$.

\begin{align*}\ensuremath
df=\omega&\Leftrightarrow\forall(x,y)\in\Omega,\;\left\{
\begin{array}{l}
 \frac{\partial f}{\partial x}(x,y)=- \frac{y}{(x-y)^2}\\
\rule{0mm}{6mm} \frac{\partial f}{\partial y}(x,y)= \frac{x}{(x-y)^2}
\end{array}
\right.\\
 &\Leftrightarrow\exists g\in C^1(\Rr,\Rr)/\;\forall(x,y)\in\Omega,\;
\left\{
\begin{array}{l}
f(x,y)= \frac{y}{x-y}+g(y)\\
 \frac{x}{(x-y)^2}+g'(y)= \frac{x}{(x-y)^2}
\end{array}
\right.\\
 &\Leftrightarrow\exists\lambda\in\Rr/\;\forall(x,y)\in\Omega,\;f(x,y)= \frac{y}{x-y}+\lambda.
\end{align*}

Les primitives de $\omega$ sur $\Omega$ sont les fonctions de la forme $(x,y)\mapsto \frac{y}{x-y}+\lambda$, $\lambda\in\Rr$.

 

\item  $\omega$ est de classe $C^1$ sur $\Rr^2\setminus\{(0,0)\}$ qui est un ouvert de $\Rr^2$ mais n'est pas étoilé. On se place dorénavant sur $\Omega=\Rr^2\setminus\{(x,0),\;x\in]-\infty,0]\}$ qui est un ouvert étoilé de $\Rr^2$. Sur $\Omega$, $\omega$ est exacte si et seulement si $\omega$ est fermée d'après le théorème de \textsc{Schwarz}.

$ \frac{\partial}{\partial x}\left( \frac{y}{x^2+y^2}-y\right)=- \frac{2xy}{(x^2+y^2)^2}= \frac{\partial}{\partial y}\left( \frac{x}{x^2+y^2}\right)$. Donc $\omega$ est exacte sur $\Omega$. Soit $f$ une application de classe $C^1$ sur $\Omega$.

\begin{align*}\ensuremath
df=\omega&\Leftrightarrow\forall(x,y)\in\Omega,\;\left\{
\begin{array}{l}
 \frac{\partial f}{\partial x}(x,y)= \frac{x}{x^2+y^2}\\
\rule{0mm}{6mm} \frac{\partial f}{\partial y}(x,y)= \frac{y}{x^2+y^2}-y
\end{array}
\right.\\
 &\Leftrightarrow\exists g\in C^1(\Rr,\Rr)/\;\forall(x,y)\in\Omega,\;\left\{
\begin{array}{l}
 \frac{\partial f}{\partial x}(x,y)= \frac{1}{2}\ln(x^2+y^2)+g(y)\\
 \frac{y}{x^2+y^2}+g'(y)= \frac{y}{x^2+y^2}-y
\end{array}
\right.\\
 &\Leftrightarrow\exists\lambda\in\Rr/\;\forall(x,y)\in\Omega,\;f(x,y)= \frac{1}{2}(\ln(x^2+y^2)-y^2)+\lambda.
\end{align*}

Les primitives de $\omega$ sur $\Omega$ sont les fonctions de la forme $(x,y)\mapsto \frac{1}{2}(\ln(x^2+y^2)-y^2)+\lambda$, $\lambda\in\Rr$.

Les fonctions précédentes sont encore des primitives de $\omega$ sur $\Rr^2\setminus\{(0,0)\}$ et donc $\omega$ est exacte sur $\Rr^2\setminus\{(0,0)\}$.

\item   $\omega$ est de classe $C^1$ sur $]0,+\infty[^2$ qui est un ouvert étoilé de $\Rr^2$. Donc $\omega$ est exacte sur $]0,+\infty[^2$ si et seulement si $\omega$ est fermée sur $]0,+\infty[^2$ d'après le théorème de \textsc{Schwarz}.

$ \frac{\partial}{\partial x}\left(
- \frac{1}{xy^2}\right)= \frac{1}{x^2y^2}$ et $ \frac{\partial}{\partial y}\left(
 \frac{1}{x^2y}\right)=- \frac{1}{x^2y^2}$. Donc $ \frac{\partial}{\partial x}\left(
- \frac{1}{xy^2}\right)\neq \frac{\partial}{\partial y}\left(
 \frac{1}{x^2y}\right)$ et $\omega$ n'est pas exacte sur $]0,+\infty[^2$.

On cherche un facteur intégrant de la forme $h~:~(x,y)\mapsto g(x^2+y^2)$ où $g$ est une fonction non nulle de classe $C^1$ sur $]0,+\infty[$.

$ \frac{\partial}{\partial x}\left(
- \frac{1}{xy^2}g(x^2+y^2)\right)= \frac{1}{x^2y^2}g(x^2+y^2)- \frac{2}{y^2}g'(x^2+y^2)$ et $ \frac{\partial}{\partial y}\left(
 \frac{1}{x^2y}g(x^2+y^2)\right)=- \frac{1}{x^2y^2}g(x^2+y^2)+ \frac{2}{x^2}g'(x^2+y^2)$.

\begin{align*}\ensuremath
h\omega\;\text{est exacte sur}\;]0,+\infty[^2&\Leftrightarrow\forall(x,y)\in]0,+\infty[^2,\; \frac{1}{x^2y^2}g(x^2+y^2)- \frac{2}{y^2}g'(x^2+y^2)=- \frac{1}{x^2y^2}g(x^2+y^2)+ \frac{2}{x^2}g'(x^2+y^2)\\
 &\Leftrightarrow\forall(x,y)\in]0,+\infty[^2,\; \frac{1}{x^2y^2}g(x^2+y^2)- \frac{x^2+y^2}{x^2y^2}g'(x^2+y^2)=0\\
 &\Leftrightarrow\forall t>0,\;-tg'(t)+g(t)=0\Leftrightarrow\exists \lambda\in\Rr/\;\forall t>0,\;g(t)=\lambda t.
\end{align*}

La forme différentielle $(x^2+y^2)\omega$ est exacte sur $]0,+\infty[^2$. De plus, 

\begin{center}
$d\left( \frac{x}{y}- \frac{y}{x}\right)=\left( \frac{1}{y}+ \frac{y}{x^2}\right)dx-\left( \frac{x}{y^2}+ \frac{1}{x}\right)dy=(x^2+y^2)\omega$.
\end{center}
\end{enumerate}
\fincorrection
\correction{005906}
\begin{enumerate}
 \item  $C$ est l'arc paramétré $t\mapsto\left( \frac{t^2-1}{2},t\right)$, $t$ variant en croissant de $-1$ à $1$.

\begin{align*}\ensuremath
\int_{C}^{}\omega&=\int_{-1}^{1}\left( \frac{(t^2-1)/2}{\left( \frac{t^2-1}{2}\right)^2+t^2}t+ \frac{t}{\left( \frac{t^2-1}{2}\right)^2+t^2}\right)dt\\
 &=0\;(\text{fonction impaire}).
\end{align*}

\begin{center}
\shadowbox{
$\int_{C}^{}\omega=2\ln2$.
}
\end{center}

\item  

\begin{align*}\ensuremath
\int_{C}^{}\omega&=\int_{0}^{2\pi}((\cos t-\sin^3t)(-\sin t)+\cos^3t(\cos t))dt=\int_{0}^{2\pi}(\cos^4t+\sin^4t-\cos t\sin t)dt\\
 &=\int_{0}^{2\pi}((\cos^2t+\sin^2t)^2-2\cos^2t\sin^2t-\cos t\sin t)dt=\int_{0}^{2\pi}\left(1- \frac{\sin(2t)}{2}- \frac{\sin^2(2t)}{2}\right)dt\\
 &=\int_{0}^{2\pi}\left(1- \frac{\sin(2t)}{2}- \frac{1}{4}(1-\cos(4t))\right)dt=2\pi\left(1- \frac{1}{4}\right)= \frac{3\pi}{2}.
\end{align*}

\begin{center}
\shadowbox{
$\int_{C}^{}\omega= \frac{3\pi}{2}$.
}
\end{center}

\item  

\begin{align*}\ensuremath
\int_{C}^{}\omega&=\int_{0}^{\pi/2}(\cos t\sin t\cos t\sin t)(-\sin t)\;dt=-\int_{0}^{\pi/2}\cos^2t\sin^3t\;dt\\
 &=\int_{0}^{\pi/2}(-\cos^2t\sin t+\cos^4t\sin t)dt=\left[ \frac{\cos^3t}{3}- \frac{\cos^5t}{5}\right]_0^{\pi/2}=- \frac{1}{3}+ \frac{1}{5}\\
 &=- \frac{2}{15}.
\end{align*}

\begin{center}
\shadowbox{
$\int_{C}^{}\omega=- \frac{2}{15}$.
}
\end{center}
\end{enumerate}
\fincorrection
\correction{005907}
\begin{enumerate}
 \item  $\omega=x^2dx+y^2dy$ est de classe $C^1$ sur $\Rr^2$ qui est un ouvert étoilé de $\Rr^2$ et est fermée car $ \frac{\partial P}{\partial y}=0= \frac{\partial Q}{\partial x}$. On en déduit que $\omega$ est exacte sur $\Rr^2$ d'après le théorème de \textsc{Schwarz}. Par suite, l'intégrale de $\omega$ le long de tout cercle parcouru une fois dans le sens trigonométrique est nulle.

\item  $\omega=y^2dx+x^2dy$ est de classe $C^1$ sur $\Rr^2$ et n'est pas fermée car $ \frac{\partial P}{\partial y}=2y\neq2x= \frac{\partial Q}{\partial x}$. On en déduit que $\omega$ n'est pas exacte sur $\Rr^2$. L'intégrale de $\omega$ le long d'un cercle parcouru une fois dans le sens trigonométrique n'est plus nécessairement nulle.

On parcourt le cercle $C$ le cercle de centre $(a,b)$ et de rayon $R>0$ une fois dans le sens trigonométrique ou encore on considère l'arc paramétré $\gamma~:~t\mapsto(a+R\cos t,b+R\sin t)$, $t$ variant en croissant de $0$ à $2\pi$.

\begin{align*}\ensuremath
\int_{\gamma}^{}\omega&=\int_{0}^{2\pi}\left((b+R\sin t)^2(-R\sin t)+(a+R\cos t)^2(R\cos t)\right)\;dt\\
 &=R\int_{0}^{2\pi}(a\cos t-b\sin t+2aR\cos^2t-2bR\sin^2t+R^2(\cos^3t-\sin^3t))\;dt\\
 &=R^2\int_{0}^{2\pi}(2a\cos^2t-2b\sin^2t+R(\cos^3t-\sin^3t))\;dt\\
 &=R^2\int_{0}^{2\pi}(a(1+\cos t)-b(1-\cos t)+R(\cos t-\sin t)(\cos^2t+\cos t\sin t+\sin^2t))\;dt\\
 &=R^2\int_{0}^{2\pi}(a-b+R(\cos t-\sin t)(1+\cos t\sin t))\;dt\\
 &=R^2\left(2\pi(b-a)+\int_{0}^{2\pi}R(\cos t-\sin t+\cos^2t\sin t-\cos t\sin^2t)\;dt\right)\\
 &=2\pi R^2(b-a).
\end{align*}
\end{enumerate}
\fincorrection
\correction{005909}

\begin{enumerate}
 \item  La forme différentielle $\omega$ est de classe $C^1$ sur $\Rr^2\setminus\{(0,0)\}$. D'après le théorème de \textsc{Schwarz}, sur tout ouvert étoilé $\Omega$ contenu dans $\Rr^2\setminus\{(0,0)\}$, la forme différentielle $\omega$ est exacte si et seulement si la forme différentielle $\omega$ est fermée.

Pour $(x,y)\in\Rr^2\setminus\{(0,0)\}$, posons $P(x,y)= \frac{e^{-y}}{x^2+y^2}(x\sin x-y\cos x)$ et $Q(x,y)= \frac{e^{-y}}{x^2+y^2}(x\cos x+y\sin x)$. Pour $(x,y)\in\Rr^2$,

\begin{align*}\ensuremath
 \frac{\partial Q}{\partial x}(x,y)&= \frac{-2xe^{-y}}{(x^2+y^2)^2}(x\cos x+y\sin x)+ \frac{e^{-y}}{x^2+y^2}(-x\sin x+\cos x+y\cos x)\\
 &= \frac{e^{-y}}{(x^2+y^2)^2}(-2x(x\cos x+y\sin x)+(x^2+y^2)(-x\sin x+\cos x+y\cos x))\\
 &= \frac{e^{-y}}{(x^2+y^2)^2}((-x^2+y^2+x^2y+y^3)\cos x+(-2xy-x^3-xy^2)\sin x),
\end{align*}

et

\begin{align*}\ensuremath
 \frac{\partial P}{\partial y}(x,y)&= \frac{-e^{-y}}{x^2+y^2}(x\sin x-y\cos x)+ \frac{-2ye^{-y}}{(x^2+y^2)^2}(x\sin x-y\cos x)+ \frac{e^{-y}}{x^2+y^2}(-\cos x)\\
 &= \frac{e^{-y}}{(x^2+y^2)^2}(-(x^2+y^2)(x\sin x-y\cos x)-2y(x\sin x-y\cos x)-(x^2+y^2)\cos x)\\
 &= \frac{e^{-y}}{(x^2+y^2)^2}((-x^2+y^2+x^2y+y^3)\cos x+(-2xy-x^3-xy^2)\sin x)\\
 &= \frac{\partial Q}{\partial x}(x,y).
\end{align*}

Finalement, la forme différentielle $\omega$ est exacte sur tout ouvert étoilé $\Omega$ contenu dans $\Rr^2\setminus\{(0,0)\}$.

On choisit $\Omega=\Rr^2\setminus\{(0,y),\;y\leqslant0\}$. $\Omega$ est un ouvert étoilé (en tout point de la forme $(0,y)$, $y>0$) de $\Rr^2$ contenant le contour fermé $\Gamma$. Puisque $\omega$ est exacte sur $\Omega$, on sait alors que $\int_{\Gamma}^{}\omega=0$.

\item  Le contour $\Gamma$ est constitué de $4$ arcs :

\textbullet~$\Gamma_1$ est l'arc $t\mapsto(t,0)$, $t$ variant en croissant de $r$ à $R$,

\textbullet~$\Gamma_2$ est l'arc $t\mapsto(R\cos t,R\sin t)$, $t$ variant en croissant de $0$ à $\pi$.

\textbullet~$\Gamma_3$ est l'arc $t\mapsto(t,0)$, $t$ variant en croissant de $-R$ à $-r$,

\textbullet~$\Gamma_4$ est l'arc $t\mapsto(r\cos t,r\sin t)$, $t$ variant en décroissant de $\pi$ à $0$.

D'après la question 1), $\int_{\Gamma_1}^{}\omega+\int_{\Gamma_2}^{}\omega+\int_{\Gamma_3}^{}\omega+\int_{\Gamma_4}^{}\omega=0$.

\begin{align*}\ensuremath
\int_{\Gamma_1}^{}\omega&=\int_{r}^{R}(P(x(t),y(t))x'(t)+Q(x(t),y(t))y'(t))\;dt=\int_{r}^{R}P(t,0)\;dt\\
 &=\int_{r}^{R} \frac{1}{t^2}\times t\sin t\;dt=\int_{r}^{R} \frac{\sin t}{t}\;dt.
\end{align*}

De même, $\int_{\Gamma_3}^{}\omega=\int_{-R}^{-r} \frac{\sin t}{t}\;dt=\int_{r}^{R} \frac{\sin t}{t}\;dt$ (puisque la fonction $x\mapsto \frac{\sin x}{x}$ est paire) et donc $\int_{\Gamma_1}^{}\omega+\int_{\Gamma_3}^{}\omega=2\int_{r}^{R} \frac{\sin x}{x}\;dx$ puis pour tout $(r,R)\in]0,+\infty[^2$ tel que $r<R$,

\begin{center}
$\int_{r}^{R} \frac{\sin x}{x}\;dx=- \frac{1}{2}\left(\int_{\Gamma_2}^{}\omega+\int_{\Gamma_4}^{}\omega\right)$.
\end{center}

Ensuite,

\begin{align*}\ensuremath
\int_{\Gamma_2}^{}\omega&=\int_{0}^{\pi}(P(R\cos t,R\sin t)(-\sin t)+Q(R\cos t,\sin t)(\cos t))\;dt\\
 &=\int_{0}^{\pi}e^{-R\sin t}((\cos t\sin(R\cos t)-\sin t\cos(R\cos t))(-\sin t)+(\cos t\cos(R\cos t)+\sin t\sin(R\cos t))(\cos t))\;dt\\
 &=\int_{0}^{\pi}e^{-R\sin t}\cos(R\cos t)\;dt.
\end{align*}

De même, $\int_{\Gamma_4}^{}\omega=\int_{\pi}^{0}e^{-r\sin t}\cos(r\cos t)\;dt=-\int_{0}^{\pi}e^{-r\sin t}\cos(r\cos t)\;dt$ et on a montré que

\begin{center}
\shadowbox{
$\forall(r,R)\in]0,+\infty[^2$, $r<R\Rightarrow\int_{r}^{R} \frac{\sin x}{x}\;dx= \frac{1}{2}\left(\int_{0}^{\pi}e^{-r\sin t}\cos(r\cos t)\;dt-\int_{0}^{\pi}e^{-R\sin t}\cos(R\cos t)\;dt\right)$.
}
\end{center}

\item  \textbullet~Etudions $\lim_{R \rightarrow +\infty}\int_{0}^{\pi}e^{-R\sin t}\cos(R\cos t)\;dt$. Pour $R>0$,

\begin{align*}\ensuremath
\left|\int_{0}^{\pi}e^{-R\sin t}\cos(R\cos t)\;dt\right|&\leqslant\int_{0}^{\pi}e^{-R\sin t}\left|\cos(R\cos t)\right|\;dt\leqslant\int_{0}^{\pi}e^{-R\sin t}\;dt=2\int_{0}^{\pi/2}e^{-R\sin t}\;dt\\
 &\leqslant2\int_{0}^{\pi/2}e^{-R(2t/\pi)}\;dt\;(\text{la fonction sinus étant concave sur}\;\left[0, \frac{\pi}{2}\right])\\
 &= \frac{\pi}{R}\left[-e^{-2Rt/\pi}\right]_0^{\pi}= \frac{\pi}{R}(1-e^{-2R})\\
 &\leqslant \frac{\pi}{R}.
\end{align*}

Comme $ \frac{\pi}{R}$ tend vers $0$ quand $R$ tend vers $+\infty$, $\lim_{R \rightarrow +\infty}\int_{0}^{\pi}e^{-R\sin t}\cos(R\cos t)\;dt=0$. On en déduit que pour tout $r>0$, l'intégrale $\int_{r}^{+\infty} \frac{\sin x}{x}\;dx$ converge en $+\infty$ et que

\begin{center}
\shadowbox{
$\forall r>0$, $\int_{r}^{+\infty} \frac{\sin x}{x}\;dx= \frac{1}{2}\int_{0}^{\pi}e^{-r\sin t}\cos(r\cos t)\;dt$.
}
\end{center}

\textbullet~Etudions maintenant $\lim_{r \rightarrow 0}\int_{0}^{\pi}e^{-r\sin t}\cos(r\cos t)\;dt$. Soit $\begin{array}[t]{cccc}
F~:&[0,+\infty[\times[0,\pi]&\rightarrow&\Rr\\
 &(r,t)&\mapsto&e^{-r\sin t}\cos(r\cos t)
\end{array}$.

- Pour tout réel $r\in[0,+\infty[$, la fonction $t\mapsto F(r,t)$ est continue par morceaux sur $[0,\pi]$.

- Pour tout réel $t\in[0,\pi]$, la fonction $r\mapsto F(r,t)$ est continue sur $[0,+\infty[$.

- Pour tout $(r,t)\in[0,+\infty[\times[0,\pi]$, $|F(r,t)|\leqslant1=\varphi(t)$ où $\varphi$ est une fonction continue par morceaux et intégrable sur le segment $[0,\pi]$.

D'après le théorème de continuité des intégrales à paramètres, la fonction $r\mapsto\int_{0}^{\pi}e^{-r\sin t}\cos(r\cos t)\;dt$ est continue sur $[0,+\infty[$. On en déduit que

\begin{center}
$\lim_{r \rightarrow 0}\int_{0}^{\pi}e^{-r\sin t}\cos(r\cos t)\;dt=\int_{0}^{\pi}e^0\cos(0)\;dt=\pi$,
\end{center}

et finalement que

\begin{center}
\shadowbox{
$\int_{0}^{+\infty} \frac{\sin x}{x}\;dx= \frac{\pi}{2}$.
}
\end{center}
\end{enumerate}
\fincorrection
\correction{005913}
Supposons tout d'abord que le support de l'arc $\gamma$ est de longueur $L=2\pi$. Puisque $\gamma$ est un arc de classe $C^1$ régulier, on peut choisir pour $\gamma$ une paramétrisation normale c'est-à-dire une paramétrisation de classe $C^1$ $t\mapsto(x(t),y(t))$, $t\in[0,2\pi]$, telle que $\forall t\in[0,2\pi]$, $x'^2(t)+y'^2(t)=1$. L'arc étant fermé, on a de plus $\gamma(0)=\gamma(2\pi)$. Cette dernière condition permet de prolonger les fonctions $x$ et $y$ en des fonctions continues sur $\Rr$ de classe $C^1$ par morceaux et $2\pi$-périodiques.

Puisque les fonctions $x'$ et $y'$ sont continues par morceaux sur $\Rr$, la formule de \textsc{Parseval} permet d'écrire

\begin{align*}\ensuremath
L=2\pi&=\int_{0}^{2\pi}1\;dt=\int_{0}^{2\pi}(x'^2(t)+y'^2(t))\;dt=\int_{0}^{2\pi}x'^2(t)\;dt+\int_{0}^{2\pi}y'^2(t)\;dt\\
 &=\pi\left( \frac{a_0^2(x')}{2}\sum_{n=1}^{+\infty}(a_n^2(x')+b_n^2(x'))+ \frac{a_0^2(y')}{2}+\sum_{n=1}^{+\infty}(a_n^2(y')+b_n^2(y'))\right)\\
 &=\pi\left(\sum_{n=1}^{+\infty}(a_n^2(x')+b_n^2(x')+a_n^2(y')+b_n^2(y'))\right) (a_0(x')= \frac{1}{\pi}\int_{0}^{2\pi}x'(t)\;dt= \frac{1}{\pi}(x(2\pi)-x(0))=0=a_0(y'))\\
 &\pi\left(\sum_{n=1}^{+\infty}n^2(a_n^2(x)+b_n^2(x)+a_n^2(y)+b_n^2(y))\right).
\end{align*}

D'autre part, d'après la formule de \textsc{Green}-\textsc{Riemann}

\begin{align*}\ensuremath
\mathcal{A}&=\int_{\gamma}^{}xdy=\int_{0}^{2\pi}x(t)y'(t)\;dt= \frac{1}{4}\int_{0}^{2\pi}((x(t)+y'(t))^2-(x(t)-y'(t))^2)\;dt\\
 &= \frac{\pi}{4}\left( \frac{a_0^2(x+y')}{2}- \frac{a_0^2(x-y')}{2}+\sum_{n=1}^{+\infty}(a_n^2(x+y')-a_n^2(x-y')+b_n^2(x+y')-b_n^2(x-y'))\right)\\
 &=\pi\left( \frac{a_0(x)a_0(y')}{2}+\sum_{n=1}^{+\infty}(a_n(x)a_n(y')+b_n(x)b_n(y'))\right)\;(\text{par linéarité des coefficients de \textsc{Fourier}})\\
 &= \frac{L}{2}\sum_{n=1}^{+\infty}n(a_n(x)b_n(y)-b_n(x)a_n(y))\\
 &\leqslant \frac{L}{2}\sum_{n=1}^{+\infty} \frac{n^2}{2}(a_n^2(x)+b_n^2(y)+b_n^2(x)+a_n^2(y))= \frac{\mathcal{L}}{2}\times \frac{\mathcal{L}}{\pi}= \frac{\mathcal{L}^2}{4\pi}.
\end{align*}

Si on a l'égalité, alors les inégalités valables pour $n\geqslant1$, 

\begin{center}
$n(a_n(x)b_n(y)-b_n(x)a_n(y))\leqslant n\times \frac{1}{2}(a_n^2(x)+b_n^2(y)+b_n^2(x)+a_n^2(y))\leqslant  \frac{n^2}{2}(a_n^2(x)+b_n^2(y)+b_n^2(x)+a_n^2(y))$,
\end{center}

sont des égalités. En particulier, pour $n\geqslant2$, on a $a_n(x)=a_n(y)=b_n(x)=b_n(y)=0$. D'autre part, quand $n=1$, $a_1(x)b_1(y)-b_1(x)a_1(y)= \frac{1}{2}(a_1^2(x)+b_1^2(y)+b_1^2(x)+a_1^2(y))$ impose $(a_1(x)-b_1(y))^2+(b_1(x)+a_1(y))^2=0$ et donc $a_1(y)=-b_1(x)$ et $b_1(y)=a_1(x)$.

D'après le théorème de \textsc{Dirichlet}, en posant $\alpha= \frac{a_0(x)}{2}$, $\beta= \frac{a_0(y)}{2}$, $a=a_1(x)$ et $b=b_1(x)$,

\begin{center}
$\forall t\in[0,2\pi]$, $\left\{
\begin{array}{l}
x(t)=\alpha+a\cos t+b\sin t=\alpha+\sqrt{a^2+b^2}\cos(t-t_0)\\
y(t)=\beta-b\cos t+a\sin t=\beta+\sqrt{a^2+b^2}\sin(t-t_0)
\end{array}
\right.$
\end{center}

où $\cos(t_0)= \frac{a}{\sqrt{a^2+b^2}}$ et $\sin(t_0)= \frac{b}{\sqrt{a^2+b^2}}$. Le support de l'arc $\gamma$ est donc un cercle. La réciproque est claire.

L'inégalité isopérimétrique est donc démontrée dans le cas où $L=2\pi$ et on a l'égalité si et seulement si le support de l'arc $\gamma$ est un cercle. Dans le cas où la longueur de la courbe $C$ est un réel strictement positif $\mathcal{L}$ quelconque, l'homothétique $(C')$ de $(C)$ dans l'homothétie de centre $O$ et de rapport $ \frac{2\pi}{\mathcal{L}}$ a une longueur $\mathcal{L}'$ égale à $2\pi$ et délimite une aire $\mathcal{A}'=\left( \frac{2\pi}{\mathcal{L}}\right)\times\mathcal{A}$.

L'inégalité $\mathcal{A}'\leqslant \frac{\mathcal{L}'^2}{2\pi}=2\pi$ s'écrit encore $\mathcal{A}\leqslant2\pi\times \frac{\mathcal{L}^2}{4\pi^2}= \frac{\mathcal{L}^2}{4\pi}$. De plus on a l'égalité si et seulement si la courbe $(C)$ est un cercle (dans ce cas, $ \frac{\mathcal{L}^2}{4\pi}= \frac{4\pi^2R^2}{4\pi}=\pi R^2=\mathcal{A}$).

\begin{center}
\shadowbox{
$\mathcal{A}\leqslant \frac{\mathcal{L}^2}{4\pi}$ avec égalité si et seulement si la courbe $(C)$ est un cercle.
}
\end{center}

(A périmètre donné, le cercle est la courbe fermée délimitant la plus grande aire)
\fincorrection
\correction{006873}
\begin{enumerate}
\item Pour $\omega_1$, on pose $P(x,y)=2xy$ et $Q(x,y)=x^2$. Comme $\omega_1$ est définie
sur l'ouvert étoilé $\Rr^2$
et que $ \frac{\partial P}{\partial y}=\frac{\partial Q}{\partial x}=2x$, le théorème de Poincaré
permet de dire que $\omega_1$ est exacte. On cherche $f$ tel que $df=\omega_1$. Ceci équivaut à résoudre le système 
$$  \left\{ \begin{array}{lll}
\frac{\partial f} {\partial x}& = &2xy \\
\frac{\partial f}{\partial y} &= & x^2\\
\end{array} \right .$$
En intégrant la première ligne par rapport à $x$, on trouve
$f(x,y)= x^2y+c(y)$. En dérivant l'expression que l'on vient d'obtenir par
rapport à $y$ et en identifiant avec la deuxième ligne
du système, on trouve $$ \frac{\partial f}{\partial y}= x^2+c'(y)=x^2.$$ Il s'ensuit que $c'(y)=0$ et 
donc que $c(y)=c \in \Rr$. Par suite, la fonction 
$f$ cherchée est :
$$f(x,y)=x^2y+c$$
o\`u $c$ est une constante réelle.

\item Pour $\omega_2$, on pose $P(x,y,z)=xy$, $Q(x,y,z)=-z$ et $R(x,y,z)=xz.$ 
On constate que 
$ \frac{\partial P}{\partial y}=x$ alors que $ \frac{\partial Q}{\partial x}=0$. La forme
$\omega_2$ n'est donc pas exacte.

\item Pour $\omega_3$, on pose $P(x,y)=2xe^{x^2-y}$ et $Q(x,y)=-2e^{x^2-y}$.
Là aussi, $ \frac{\partial P}{\partial y}\not= \frac{\partial Q}{\partial x} $
puisque $ \frac{\partial P}{\partial y}=-2xe^{x^2-y}$ alors
que $ \frac{\partial Q}{\partial x}=-4xe^{x^2-y}$ ; $\omega_3$ n'est donc pas exacte.

\item Pour $\omega_4$, posons
$P(x,y,z)=yz^2$, $Q(x,y,z)=xz^2+z$, $R(x,y,z)=2xyz+2z+y$. On constate
que 
\begin{enumerate}
\item $ \frac{\partial P}{\partial y}= \frac{\partial Q}{\partial x}= z^2$
\item $ \frac{\partial P}{\partial z}= \frac{\partial R}{\partial x}=2zy$
\item $ \frac{\partial Q}{\partial z}= \frac{\partial R}{\partial y}= 2xz+1$.
\end{enumerate}
La forme $\omega_4$ est de plus définie sur l'ouvert étoilé $\Rr^3$, elle est donc exacte
d'après le théorème de Poincaré.
Cherchons maintenant $f$ telle que $df=\omega_4$, ceci revient à résoudre le système :
$$  \left\{ \begin{array}{lll}
\frac{\partial f} {\partial x}& = &yz^2 \\
\frac{\partial f}{\partial y} &= &xz^2+z\\
\frac{\partial f}{\partial z} & = & 2xyz+2z+y \\
\end{array} \right .$$
En intégrant la première équation par rapport à $x$, on trouve
$$f(x,y,z)=xyz^2+\psi(y,z).$$
Maintenant, en dérivant l'expression obtenue successivement par $y$ et $z$
et en égalisant avec les deux dernières équations du système, on obtient 
un nouveau système 
$$  \left\{ \begin{array}{lll}
 xz^2+\frac{\partial \psi} {\partial y}& = &xz^2+z \\
 2xyz+\frac{\partial \psi}{\partial z} &= &2xyz+2z+y\\
\end{array} \right . $$
qui  équivaut à :
$$  \left\{ \begin{array}{lll}
 \frac{\partial \psi} {\partial y}& = & z \quad (1) \\
 \frac{\partial \psi}{\partial z} &= & 2z+y \quad (2)\\
\end{array} \right .$$
Finalement, en intégrant (1) par rapport à $y$, il vient 
$\psi(y,z)=zy+c(z)$. En dérivant cette expression de $\psi$ par rapport à $z$
et en égalisant avec (2), on trouve
$y+c'(z)=2z+y$, c'est-à-dire $c'(z)=2z$ donc $c(z)=z^2+c$ o\`u $c\in \Rr$.
Ainsi, la fonction $f$ telle que $\omega_4=df$ est de la forme
$$f(x,y,z)= xyz^2+zy  +z^2+c$$
o\`u $c\in \Rr$.
\end{enumerate}
\fincorrection
\correction{006874}
\begin{enumerate}
\item On vérifie que :
\begin{enumerate}
\item $dx= \cos \varphi \cos \theta dr-r\sin \varphi\cos \theta d\varphi-r\sin \theta\cos \varphi 
d\theta$

\item $dy = \cos \varphi \sin \theta dr-r\sin \varphi\sin \theta d\varphi+r\cos \theta\cos\varphi d\theta $
\item $dz=\sin \varphi dr+r\cos \varphi d\varphi.$
\end{enumerate}
Par suite, on a :
\begin{enumerate}
\item $xdx= r\cos^2 \varphi \cos^2 \theta dr-r^2\sin \varphi\cos \varphi\cos^2 \theta d\varphi-r^2\sin \theta\cos \theta\cos^2 \varphi 
d\theta$
\item $ydy = r\cos^2\varphi \sin^2 \theta dr-r^2\sin \varphi\cos \varphi\sin^2 \theta d\varphi+r^2\cos \theta\sin \theta\cos^2\varphi d\theta $
\item $zdz=r\sin^2 \varphi dr+r^2\cos \varphi \sin \varphi d\varphi.$
\end{enumerate}

\item En additionnant, on obtient $xdx+ydy+zdz=rdr$.
On en déduit que :
$$xdx+ydy+zdz=r(  \frac{\partial r}{\partial x}dx+  \frac{\partial r}{\partial y}dy+  \frac{\partial r}{\partial z}dz).$$
Ainsi
$$  \frac{\partial r}{\partial x}= \frac{x}{r}\quad\quad
  \frac{\partial r}{\partial y}=\frac{y}{r} \quad\quad
   \frac{\partial r}{\partial z}=\frac{z}{r}.$$
  \end{enumerate}
\fincorrection
\correction{006875}
\begin{enumerate}
\item Posons $P(x,y)=x^2+y^2+2x$ et $Q(x,y)=2y$. On voit facilement
que $ \frac{\partial P}{\partial y}\not=\frac{\partial Q}{\partial x}$. La forme
$\omega$ n'est donc pas exacte.

\item Comme $\omega$ est définie sur $\Rr^2$, il suffit que $\psi\omega$ soit exacte pour que $f$ existe.
Maintenant, $\psi\omega$ est exacte si et seulement si
$$ \frac{\partial (\psi(x)(x^2+y^2+2x))}{\partial y } =  \frac{\partial(\psi(x) 2y)} {\partial x}.$$
Ceci équivaut à $2y \psi(x)=2y\psi'(x)$. Ainsi, $\psi(x)=\psi'(x)$ pour tout $x$. Donc
$\psi(x)=k e^x$ avec $k$ constante. On peut choisir $k=1$.
Ainsi  $$\psi\omega=e^x(x^2+y^2+2x) dx+e^x(2y)dy.$$ On cherche ensuite $f$ telle que :
$$  \left\{ \begin{array}{lll}
\frac{\partial f} {\partial x}& = &e^x(x^2+y^2+2x) \\
\frac{\partial f}{\partial y} &= &e^x(2y) \\
\end{array} \right .$$
En intégrant la deuxième équation par rapport à $y$, on trouve 
$$f(x,y)=e^xy^2+c(x).$$
En dérivant cette expression par rapport à $x$ et en égalisant avec la première équation du système, on obtient
$$ e^xy^2+c'(x)=e^x(x^2+y^2+2x)$$ c'est-à-dire
$$c'(x)=e^x(x^2+2x).$$
Il en résulte que $c(x)=x^2 e^x+c$ et donc que
$$f(x,y)=e^x (x^2+y^2)+c$$ avec $c$ dans $\Rr$.
\end{enumerate}
\fincorrection
\correction{006876}
Au champ $\vec{V}(x,y)$ est associée la forme
$$\omega=(1+2xy) dx + (x^3-3) dy.$$
Cette forme n'est pas exacte puisque
$ \frac{\partial (1+2xy)}{\partial y}\not=  \frac{\partial(x^3-3)}{\partial x}$. Il s'ensuit
que $\vec{V(x,y)}$ n'est pas un champ de gradient.
\fincorrection
\correction{006877}
Le champ vectoriel qui dérive du potentiel $U$ est
$$\vec{\mathrm{grad}}(U)= (  \frac{\partial u}{\partial x},   \frac{\partial u}{\partial y},  \frac{\partial u}{\partial z}).$$ 
Il s'agit donc du champ vectoriel de composantes :
$$\vec{\mathrm{grad}}(U)=(1+y+yz, x+xz, xy).$$
\fincorrection
\correction{006878}
Soit $\omega=3x dx + (x+y) dy$ la forme différentielle naturellement associée à $\vec{V}(x,y)$ et considérons $x=\cos t$ et $y=\sin t$ 
comme paramétrage du cercle de centre $O$ et de rayon $1$ (avec $t\in [0;2\pi]$). Il s'ensuit que la circulation $ \int_{C} \vec{V}.\vec{dl}$
n'est autre que :
$$ \int_{C} \vec{V}.\vec{dl}= \int_{C} w
=\int_0^{2\pi} (3\cos t (-\sin t) + (\cos t +\sin t) \cos t)dt.$$
Comme $\cos^2t= \frac{\cos(2t)+1}{2}$, on obtient :
$$ \int_{C} \vec{V}.\vec{dl}= \int_{0}^{2\pi} (-2\sin t\cos t + \frac{\cos (2t)+1}{2}) dt
= [\cos^2(t)+ \frac{1}{4}\sin(2t) +\frac{t}{2}]_0^{2\pi}=\pi.$$
Remarquons que si la forme $\omega$ avait été exacte, on aurait obtenu $ \int_{C} \vec{V}.\vec{dl}=0$ comme réponse, puisque l'intégrale curviligne d'une forme exacte sur une courbe fermée
est nulle.
\fincorrection
\correction{006879}
Notons $\omega= yz dx + zx dy + xy dz$ la forme différentielle associée à $\vec{F}(x,y,z)$.
Par définition de $W$, on a $W= \int_H \vec{F}. \vec{dl}=\int_H \omega.$
D'après le paramétrage donné pour $H$, on a
\begin{align*}
W & = \int_0^{\frac{\pi}{4}} &yz dx + zx dy + xy dz \\
  & = \int_0^{\frac{\pi}{4}}& ((\sin t )t (-\sin t)+t\cos^2 t + \cos t \sin t ) dt\\
  & = \int_0^{\frac{\pi}{4}}&(t \cos(2t) +\cos t\sin t ) dt.
\end{align*} 
On a utilisé ici la formule trigonométrique : $\cos(2t)=\cos^2 t-\sin^2 t.$
En faisant une intégration par parties, on constate que
$$ \int_0^\frac{\pi}{4} t \cos(2t) dt =[ \frac{t \sin(2t)}{2}]_0^{\frac{\pi}{4}} -
\int_0^{\frac{\pi}{4}}  \frac{\sin(2t)}{2} dt . $$ On en déduit que 
$$W=  [  \frac{t \sin(2t)}{2}]_0^{\frac{\pi}{4}} +  \frac{1}{4} [\cos (2t)]_0^{\frac{\pi}{4}}
+ \frac{1}{2} [\sin^2 (t) ]_0^{\frac{\pi}{4}}
=  \frac{\pi}{8} - \frac{1}{4}+ \frac{1}{4}=  \frac{\pi}{8}.$$
Remarquons que $\omega=  yz dx + zx dy + xy dz$ est exacte. De plus, 
on vérifie aisément que $\omega= d(xyz)$. On peut alors retrouver le résultat précédent en faisant :
$$W= f(B)-f(A)$$ o\`u l'on a posé $f(x,y,z)=xyz$, 
$$B=(\cos( \frac{\pi}{4}) , \sin(\frac{\pi}{4}) ,  \frac{\pi}{4})=
( \frac{\sqrt{2}}{2}, \frac{\sqrt{2}}{2} ,  \frac{\pi}{4})$$
et $$A=(\cos(0) , \sin(0) , 0)=(1, 0, 0).$$
\fincorrection
\correction{006880}
\begin{enumerate}
\item On note $P(x,y,z)=y^2\cos x$, $Q(x,y,z)=2y\sin x+e^{2z}$ et
$R(x,y,z)=2ye^{2z}$. La forme
$\omega=Pdx+Qdy+Rdz$, naturellement associée au champ $\vec{V}(x,y,z)$,
est exacte puisqu'elle est définie sur $\Rr^3$ et 
\begin{enumerate}
\item $ \frac{\partial P}{\partial y}=  \frac{\partial Q}{\partial x}=2y\cos x$
\item $ \frac{\partial P}{\partial z}=  \frac{\partial R}{\partial x}= 0$
\item $ \frac{\partial Q}{\partial z}=  \frac{\partial R}{\partial y}=2e^{2z}.$
\end{enumerate}
Le champ $\vec{V}(x,y,z)$ est donc un champ de gradient.
\item  Cherchons $U$ tel que $\omega=dU$. Cela nous conduit à résoudre le système :
 $$  \left\{ \begin{array}{lll}
 \frac{\partial U} {\partial x}& = &y^2 \cos x \\
 \frac{\partial U}{\partial y} &= &2y\sin x+e^{2z}\\
\frac{\partial U}{\partial z} &=&2ye^{2z}\\
\end{array} \right .$$

En intégrant la première équation par rapport à $x$, on trouve :
$$U(x,y,z)=y^2\sin x+\psi(y,z).$$
Maintenant, en utilisant les deux dernières équations, on est amené à
résoudre le système suivant :
$$  \left\{ \begin{array}{lll}
\frac{\partial \psi} {\partial y}& = &e^{2z} \\
 \frac{\partial \psi}{\partial z} &= &2ye^{2z} \\
\end{array} \right .$$
Par suite, on vérifie que $\psi(y,z)=e^{2z}y+c(z)$ avec $c'(z)=0$. Donc $c(z)=c$ 
avec $c$ constante
réelle et finalement :
$$U(x,y,z)=y^2\sin x+e^{2z}y+c$$
avec $c\in \Rr$.
Par ailleurs, on veut que $U(0,0,0)=1$ ce qui donne $c=1$.

\item La circulation du champ de $A(0,1,0)$ à $B( \frac{\pi}{2},3,0)$
est 
$$ \int_{\overset{\frown}{AB}} \vec{V}.\vec{dl}=\int_{\overset{\frown}{AB}} \omega = U(B)-U(A)=
U( \frac{\pi}{2},3,0)-U(0,1,0)=11.$$
Remarquons que lorsque $\omega$ est exacte, pour calculer l'intégrale curviligne de
$\omega$ sur un chemin, il suffit de conna\^{\i}tre l'origine et l'extrémité du chemin. Autrement
dit, l'intégrale curviligne d'une forme exacte sur $\overset{\frown}{AB}$ ne dépend que de $A$
et de $B$, et non du chemin choisi pour aller de $A$ à $B$.
\end{enumerate}
\fincorrection
\correction{006881}
On rapporte le plan à un repère orthonormé direct d'origine $O$.
D'après la formule de Green-Riemann, en choisissant de prendre $P=0$ et
$Q=x^2y$ de sorte que  $ \frac{\partial Q}{\partial x}-\frac{\partial P}{\partial y}=xy$, on obtient :
$$I= \iint_{\mathcal{D}} xydxdy=\int_{T} x^2ydy$$
o\`u l'on a noté $T$ le triangle $OAB$ orienté dans le sens direct
avec $O(0,0)$, $A(1,0)$ et $B(1,1)$.
Ainsi 
$$I=\iint_{\mathcal{D}} xydxdy=\int_{\overset{\frown}{OA}} x^2ydy
+\int_{\overset{\frown}{AB}} x^2ydy+
\int_{\overset{\frown}{BO}} x^2ydy.$$
L'intégrale curviligne d'une forme différentielle sur un chemin est indépendant du paramétrage choisi 
pour ce chemin. Pour le calcul, nous choisissons de paramétrer
$\overset{\frown}{OA}$ par $x=t$ et $y=0$ avec $t$ variant de $0$ à $1$ et ainsi
$\int_{\overset{\frown}{OA}} x^2ydy=0$. 
De m\^eme, nous choisissons de paramétrer
$\overset{\frown}{BO}$ par $x=0$ et $y=t$ avec $t$ variant de $1$ à $0$ et ainsi
$\int_{\overset{\frown}{BO}} x^2ydy=0$. 
Enfin, nous choisissons de paramétrer
$\overset{\frown}{AB}$ par $x=t$ et $y=1-t$ avec $t$ allant de $1$ à $0$ et donc :
$$I= \iint_{\mathcal{D}} xydxdy=\int_{\overset{\frown}{AB}} x^2ydy
= \int_1^0  \frac{t^2(1-t) }{2}(-dt)
= \int_0^1  \frac{t^2(1-t)}{2} dt = 
   \frac{1}{24}.$$
 Remarquons qu'il n'aurait pas été plus difficile ici de calculer directement l'intégrale
 double sans utiliser la formule de Green-Riemann :
 $$\iint_{\mathcal{D}} xy dx dy=
  \int_{0^1} (\int_{0}^{1-x} xy dy) dx=  \int_0^1 x [\frac{y^2}{2}]_0^{1-x} dx
 = \int_0^1 x\frac{(1-x)^2}{2}dx =  \frac{1}{24}.$$
\fincorrection
\correction{006882}
\begin{enumerate}
\item La forme $\omega= \frac{-y}{x^2+y^2}dx+\frac{x}{x^2+y^2} dy$ est définie
sur $\Rr^2 \setminus \{(0,0)\}$.

\item Paramétrons le cercle $C$ par $x=\cos t$, $y=\sin t$ avec $t\in [0;2\pi]$.
On obtient :
\begin{align*}
  \int_{C} \omega&=& \int_{0}^{2\pi}& (-\sin t(-\sin t)+\cos t(\cos t)) dt &\\
 &=&\int_0^{2\pi} &\sin^2 t+\cos^2 t dt&\\
 &=& \int_0^{2\pi}& 1 dt& \\
 &=& 2\pi.&&
 \end{align*}

\item La forme $\omega$ n'est pas exacte, sinon son intégrale curviligne sur la courbe
fermée $C$ serait nulle et cela contredirait notre résultat de la
question précédente. Remarquons cependant que
$$
 \frac{\partial}{\partial y}( \frac{-y}{x^2+y^2})=
 \frac{\partial}{\partial x} ( \frac{x}{x^2+y^2})=  \frac{y^2-x^2}{(x^2+y^2)^2}.$$
En fait, avec cet exemple, on voit que dans le théorème de Poincaré, l'hypothèse que l'ouvert doit \^etre étoilé,
est indispensable. Ici $\Rr^2 \setminus \{(0,0)\}$ n'est pas étoilé, c'est un domaine "troué". De plus,
$\int_{C} \omega$ n'est pas nulle car le cercle entoure le "trou".
\end{enumerate}
\fincorrection
\correction{004944}
Droite orthogonale à $\vec S\wedge\vec u$,
         $\vec u \perp\vec{{\cal P}}$,
         ou $\varnothing$ ou ${\cal P}$.
\fincorrection
\nocorrection
\nocorrection
\correction{004947}
Soit $\cal T$ un torseur : on décompose ${\cal T}(A)$ en
         $\alpha\vec{AB}\wedge\vec{BC} +
          \beta \vec{AB}\wedge\vec{BD} +
          \gamma\vec{AC}\wedge\vec{CD}$, et $\vec R$ en
         $\alpha'\vec{AB} + \beta'\vec{AC} + \gamma'\vec{AD}$.
         \par   $ \Rightarrow $ famille génératrice.
\fincorrection
\nocorrection
\correction{005915}
\begin{enumerate}
 \item  Pour $(x,y)\in\Rr^2$, posons $f(x,y)=x^4-x^3+xy-y^2$ puis pour $(x,y,z)\in\Rr^3$, posons $g(x,y,z)=z-f(x,y)$. $\mathcal{S}$ est la surface d'équation $z=f(x,y)$ ou encore $g(x,y,z)=0$.

La fonction $g$ est de classe $C^1$ sur $\Rr^3$ et pour tout $(x,y,z)\in\Rr^3$,

\begin{center}
$\left(\overrightarrow{\text{grad}}\;g\right)(x,y,z)=\left(
\begin{array}{c}
- \frac{\partial f}{\partial x}(x,y)\\
\rule[-4mm]{0mm}{11mm}- \frac{\partial f}{\partial y}(x,y)\\
1
\end{array}
\right)=\left(
\begin{array}{c}
-4x^3+3x^2-y\\
-x+2y\\
1
\end{array}
\right)\neq\overrightarrow{0}$.
\end{center}

Donc, la surface $\mathcal{S}$ est régulière et en tout point $(x_0,y_0,z_0)$ de la surface $\mathcal{S}$, le vecteur gradient est un vecteur normal au plan tangent $\mathcal{P}_0$ à la surface $\mathcal{S}$ en $(x_0,y_0,z_0)$. Le plan 

\begin{align*}\ensuremath
\mathcal{P}_0\;\text{parallèle à}\;\left(O,\overrightarrow{i},\overrightarrow{j}\right)&\Leftrightarrow\left\{
\begin{array}{l}
-4x_0^3+3x_0^2-y_0=0\\
-x_0+2y_0=0
\end{array}
\right.\Leftrightarrow\left\{
\begin{array}{l}
x_0=2y_0\\
-y_0(32y_0^2-12y_0+1)=0
\end{array}
\right.\Leftrightarrow\left\{
\begin{array}{l}
x_0=2y_0\\
-y_0(32y_0^2-12y_0+1)=0
\end{array}
\right.
\\
 &\Leftrightarrow y_0=0=x_0\;\text{ou}\;(y_0= \frac{1}{4}\;\text{et}\;x_0= \frac{1}{2})\;\text{ou}\;(y_0= \frac{1}{8}\;\text{et}\;x_0= \frac{1}{4})
\end{align*}

On obtient ainsi les trois points $O(0,0,0)$, $A\left( \frac{1}{2}, \frac{1}{4},0\right)$ et $B\left( \frac{1}{4}, \frac{1}{8}, \frac{1}{256}\right)$.

\item  La fonction $f$ est de classe $C^2$ sur $\Rr^2$ et

\begin{center}
$rt-s^2= \frac{\partial^2}{\partial x^2} \frac{\partial^2}{\partial y^2}-\left( \frac{\partial^2}{\partial x\partial y}\right)^2=(12x^2-6x)(-2)-1^2=-24x^2+12x-1$
\end{center}

\textbullet~En $O$, le plan tangent est le plan $\left(O,\overrightarrow{i},\overrightarrow{j}\right)$. De plus, $(rt-s^2)(0,0)=-1<0$. Donc le point $O$ est un point selle.

\textbullet~En $A$, le plan tangent est aussi le plan $\left(O,\overrightarrow{i},\overrightarrow{j}\right)$. De plus, $(rt-s^2)\left( \frac{1}{2}, \frac{1}{4}\right)=-1<0$. Donc le point $A$ est un point selle.

\textbullet~En $B$, le plan tangent est le plan d'équation $z= \frac{1}{256}$. De plus, $(rt-s^2)\left( \frac{1}{4}, \frac{1}{8}\right)= \frac{1}{2}>0$. Donc la surface $\mathcal{S}$ a une disposition en ballon au point $B$.

\item  Il s'agit maintenant d'étudier le signe de $z=f(x,y)=x^4-x^3+xy-y^2$ sur $\Rr^2$.

\begin{center}
$f(x,y)=x^4-x^3+xy-y^2=(x^4-y^2)-x(x^2-y)=(x^2-y)(x^2+y-x)$.
\end{center}

L'intersection de la surface $\mathcal{S}$ avec le plan $\left(O,\overrightarrow{i},\overrightarrow{j}\right)$ est donc la réunion des deux paraboles d'équations respectives $y=x^2$ et $y=-x^2+x$ dans le plan $\left(O,\overrightarrow{i},\overrightarrow{j}\right)$. Représentons cette intersection ainsi que le signe de $f(x,y)\oplus\ominus$.

$$\includegraphics{../images/img005915-1}$$


\end{enumerate}
\fincorrection
\nocorrection
\correction{005047}
$4x^2 + 4y^2 - 3z^2 = a^2$.
\fincorrection
\correction{005048}
$(\Gamma)$ est l'intersection d'un cylindre hyperbolique et d'un plan.
         C'est une hyperbole dans ce plan.\par
         Pour $M(x,y,z) \in \Gamma$, on
         pose $r=\sqrt{x^2+y^2}$ et on élimine $x$ et $y$ entre les équations :
         $$\begin{cases}x^2+y^2=r^2\cr
                  x^2-y^2-4x+2=0\cr
                  x+z=1.\cr\end{cases}$$
         ce qui donne $2z^2=r^2$, donc $\Gamma$ est incluse dans l'hyperboloïde
         de révolution d'équation $2z^2=x^2+y^2$ et la surface cherchée itou.
	 La réciproque est évidente.
\fincorrection
\correction{005049}
\begin{enumerate}
  \item $\left(\sin\theta\frac{\partial f}{\partial \theta} -\rho\cos\theta\frac{\partial f}{\partial \rho}\right)x
             -\left(\cos\theta\frac{\partial f}{\partial \theta} +\rho\sin\theta\frac{\partial f}{\partial \rho}\right)y
             + \rho z = \rho f - \rho^2\frac{\partial f}{\partial \rho}$.
  \item $f-\rho\frac{\partial f}{\partial \rho} = a(\theta)
              \Rightarrow  f(\rho,\theta) = a(\theta) + b(\theta)\rho$.
  \item
$$
\centerline{\includegraphics[height=6cm]{../images/img005049-1}}
$$
  % \mapleplot{ x := r*cos(t); y := r*sin(t); z := t;    plot3d([x,y,z],r=0..2,t=0..4*Pi,style=hidden,color=black,orientation=[20,60]);   }
\end{enumerate}
 \fincorrection
\correction{005050}
$$
\centerline{\includegraphics[height=6cm]{../images/img005050-1}}
$$
% \mapleplot{%
% x := cos(u)/cosh(v); y := sin(u)/cosh(v); z := v-tanh(v);
% plot3d([x,y,z],u=0..2*Pi,v=-4..4,style=hidden,color=black,orientation=[45,70]);}
\fincorrection
\correction{005051}
La normale en $M$ est parallèle ou sécante à $Oz
         \Leftrightarrow y\frac{\partial f}{\partial x} - x\frac{\partial f}{\partial y} = 0
	 \Leftrightarrow \frac{\partial f}{\partial \theta} = 0 \Leftrightarrow f = f(\rho)$.
\fincorrection
\nocorrection
\correction{005053}
\begin{enumerate}
  \item Hyperboloïde de révolution à deux nappes.
  \item $x = 2y$, $z^2 = 1 + 5y^2$.
\end{enumerate}
 \fincorrection
\correction{005054}
\begin{enumerate}
  \item $z=x^2+y^2$.
  \item $x+y=\frac12$.
  \item $(x-y+\frac12)^2 = 2(z-y+\frac14)$.
\end{enumerate}
 \fincorrection
\correction{005055}
$x^2+y^2+z^2-2xz-2yz+2z=1$.
\fincorrection
\nocorrection
\nocorrection
\correction{005058}
\begin{enumerate}
  \item $\begin{pmatrix} (x+2y)(y+2z) \cr -(2x+y)(y+2z) \cr (2x+y)(2y+z)\cr \end{pmatrix}$
             sauf pour $M=\pm\frac2{\sqrt3}(1,-2,1)$.
  \item $(x-z)(x+y+z)=0$.
  \item segment $x=z\in {[-1,1]}$ et ellipse $x^2+z^2+xz=1$.
\end{enumerate}
 \fincorrection
\correction{005059}
$a^2y^2 = (x^2+y^2)(r^2-z^2)$.
\fincorrection
\correction{005060}
$y(x^2+(y-1)^2 + z^2) = z^2$.
\fincorrection
\correction{005061}
$x=\frac a2(1+\cos u)$,
	     $y=\frac v2(1+\cos u)$,
	     $z=\frac{\sqrt{a^2+v^2}}2\sin u$.
\fincorrection
\correction{005062}

\begin{enumerate}

  \item $x^2+y^2+z^2=a^2$.
  \item On paramètre $(\Sigma)$ par :
             $\begin{cases} x = a\cos u/\ch v\cr y = a\sin u/\sh v\cr z = a\tanh v.\cr\end{cases}$\par
        La tangente à la méridienne passant par $M(u,v)$ est dirigée par $\frac{\partial {\vec M}}{\partial v}$ et la tangente à $(\Gamma)$ passant par $M(t,mt)$ est dirigée par $\frac{\partial  {\vec M}}{\partial u} + m\frac{\partial {\vec M}}{\partial v}$. Après calculs, le cosinus de ces deux vecteurs vaut $\frac m{\sqrt{m^2+1}}$ donc est constant.
  \item Une courbe tracée sur $(\Sigma)$ est définie par la donnée de $u$ et $v$ en fonction d'un paramètre $t$. Le cosinus de l'angle entre cette courbe et une méridienne de $(\Sigma)$ vaut $\frac{u'}{\sqrt{u'^2+v'^2}}$, donc est constant si et seulement si le rapport \frac{v'}{u'}$ est constant. En notant $m$ cette constante et en prenant $u(t)=t$, on trouve les courbes déduites de $(\Gamma)$ par rotation autour de $Oz$. 
 \item 
$$\centerline{\includegraphics[height=4cm]{../images/img005062-1}}$$
%     \mapleplot{plot([1/cosh(t/3),t,t=-20..20],view=[-1..1,-1..1],coords=polar,axes=frame);}
%     \hbox{\indent équation polaire $\rho = \frac{a}{\ch m\theta}$.}
\end{enumerate}
\fincorrection
\correction{005983}
Classements possibles : sans ex-aequo, il y en a 20!.

Avec exactement $2$ ex-aequo, il y en a :
\begin{enumerate}
\item Choix des deux ex-aequo : $\binom{20}{2}=$ $190$ choix;
\item Place des ex-aequo : il y a $19$ possibilités;
\item Classements des $18$ autres personnes, une fois les ex-aequo placés : il y a $18!$ choix.
\end{enumerate}
Il y a au total : $19\binom{20}{2}(18!)$ choix possibles.
\fincorrection
\correction{005984}
- Une tenue est un triplet $(P, T, C)$ : il y a $5\times 6\times 8=240$ tenues
différentes;

- <<Il est tout en noir>> : de combien de façons différentes ? 
Réponse : de $2\times 4\times 5=40$ façons.

La probabilité de l'événement <<Il est tout en
noir>> est donc : $\frac{40}{240}=\frac{1}{6}$.

- <<Une seule pièce est noire sur les trois >> : notons les événements :
$N_{1}$ la première pièce (pantalon) est noire, $N_{2}$ la deuxième pièce (tee-shirt)
 est noire, $N_{3}$ la troisième pièce
(chaussette) est noire: l'événement est représenté par :
$(N_{1}\cap \overline{N_{2}}\cap \overline{N_{3}})\cup (\overline{N_{1}}\cap
N_{2}\cap \overline{N_{3}})\cup (\overline{N_{1}}\cap \overline{N_{2}}\cap
N_{3})$.
Ces trois événements sont disjoints, leurs probabilités
s'ajoutent. La probabilité de l'événement <<une seule pièce est noire sur les trois>> est donc : 
$0.325$.
\fincorrection
\correction{005985}
Il y a $\binom{30}{2}$ façons de choisir $2$ personnes parmi $30$ et donc $2\cdot\binom{30}{2} = 870$ bises.
\fincorrection
\correction{005986}
\begin{enumerate}
\item Une grille-réponses est une suite ordonnée de $10$ réponses,
il y a $4$ choix possibles pour chacune. Il y a donc $4^{10}$ grilles-réponses possibles.
\item L'événement $E$ <<répondre au hasard au moins $6$ fois correctement>> est 
réalisé si le candidat répond bien à $6$ ou $7$ ou $8$ ou $9$ ou $10$ questions. Notons $A_{n}$ l'événement : <<répondre au hasard
exactement $n$ fois correctement>>. Alors, $A_{n}$ est réalisé si $n$ réponses sont correctes et $10-n$
sont incorrectes : $3$ choix sont possibles pour chacune de ces dernières.
Comme il y a $\binom{10}{n}$ choix de $n$ objets parmi $10$, et donc il y a :
$\binom{10}{n} \times 3^{10-n}$ façons de réaliser $A_{n}$ et :
\begin{equation*}
P(A_n)=\frac{\binom{10}{n}\cdot 3^{10-n}}{4^{10}}
\end{equation*}
pour $n= 6, 7, 8, 9, 10$.
$P(E)=\sum_{n=6}^{10}\frac{\binom{10}{n}\cdot 3^{10-n}}{4^{10}} \simeq 1.9728\times 10^{-2}$, soit environ $2$\%.
\end{enumerate}
\fincorrection
\correction{005987}

Considérons plutôt l'événement complémentaire :
l'oiseau n'est pas touché s'il n'est touché ni par Amédée,
ni par Barnabé, ni par Charles.
Cet événement a pour probabilité : $(1-0.7)\cdot (1-0.5)\cdot (1-0.9)=0.015$.
La probabilité que l'oiseau soit touché est donc : $1-0.015=0.985$.
\fincorrection
\correction{005988}

L'univers des possibles est ici l'ensemble des combinaisons de $10$ billets
parmi les $300$ ; il y en a $\binom{300}{10}$.
Je ne gagne rien si les $10$ billets achetés se trouvent parmi les $296$
billets perdants, ceci avec la probabilité : 
\begin{equation*}
\frac{\binom{296}{10}}{\binom{300}{10}}.
\end{equation*}

La probabilité cherchée est celle de l'événement complémentaire : 
\begin{equation*}
1-\frac{\binom{296}{10}}{\binom{300}{10}}\simeq 0.127.
\end{equation*}

La probabilité est environ $12.7$\% de gagner au moins un lot.
\fincorrection
\correction{005989}
$P(A\cap B)=pq$ car les maladies sont indépendantes.
$P(A\cup B)=P(A)+P(B)-P(A\cap B)=p+q-pq$
\fincorrection
\correction{005990}
Soit $A$ : l'événement <<tirer un roi>>
et $B$ : <<tirer un pique>>.

$P(A\cap B)=\frac{1}{52};P(A)=\frac{4}{52}=\frac{1}{13};P(B)=\frac{13}{52}=\frac{1}{4}.$

Donc $P(A\cap B)=P(A)P(B)$ et donc les événements $A$ et $B$ sont indépendants.

$P(A\cup B)=P(A)+P(B)-P(A\cap B)=\frac{1}{13}+\frac{1}{4}-\frac{1}{52}=\frac{4}{13}$.
\fincorrection
\correction{005991}
Notons, pour le cas où la famille Potter comporte $2$
enfants, l'univers des possibles pour les enfants :
$\Omega =\{(G,G),(G,F),(F,G),(F,F)\}$, représente les cas possibles, 
équiprobables, d'avoir garçon-garçon, garçon-fille etc... :
Alors $P(A)=\frac{2}{4},$ $P(B)=\frac{3}{4}, P(A\cap B) = \frac{2}{4}$.
On en conclut que : $P(A\cap B)\neq P(A)P(B)$ et donc que les événements $A$ et $B$ ne sont pas indépendants.

Si maintenant la famille Potter comporte $3$ enfants :
Alors $\Omega' = \{ (a,b,c) \mid a \in \{ G,F \}, b\in \{G,F\}, c\in \{G,F\}\}$ 
représente les $2^{3}=8$ cas possibles, équiprobables.
Cette fois, $P(A)=1-P(\{(G,G,G),(F,F,F)\})=\frac{6}{8}$ ;
$P(B)=\frac{4}{8}, P(A\cap B)=P\{(F,G,G),(G,F,G),\{(G,G,F)\}=\frac{3}{8}$.
On a $P(A)P(B)=\frac{3}{8}=P(A\cap B),$ et les événements $A$ et $B$
sont indépendants

Avec $n$ enfants, on peut généraliser sans difficulté : 
$P(A)=1-\frac{2}{2^{n}},$ $P(B)=\frac{1+n}{2^{n}}$
$P(A\cap B)=\frac{n}{2^{n}}$
Un petit calcul montre que
$P(A)P(B)=P(A\cap B)$ si et seulement si $n=3$.
\fincorrection
\nocorrection
\nocorrection
\correction{006891}
Pour gagner 10 euros, il faut avoir tiré exactement une boule rouge et 4 boules blanches, 
ou 2 boules vertes et 3 boules blanches. Ces deux événements, notés $A$ et $B$, sont 
incompatibles. Ici l'univers $\Omega$ est l'ensemble des combinaisons de 5 boules, 
c'est-à-dire l'ensemble des parties à 5 éléments d'un ensemble de 20 boules. 
Les ${20\choose{5}}=\frac{20.19.18.17.16}{5!}=15504$ combinaisons sont équiprobables. 
Les éléments de l'événement $A$ sont les combinaisons formées de la boule rouge et 
d'une combinaison de 4 boules blanches. Il y en a ${16\choose4}=\frac{16.15.14.13}{4!}=1820$. 
Par conséquent, $p(A)=\frac{1820}{15504}$. Les éléments de l'événement $B$ sont les 
combinaisons formées d'une combinaison de deux boules vertes et d'une combinaison de 3 
boules blanches. Il y en a ${3\choose{2}}{16\choose{3}}=3.\frac{16.15.14}{6!}=1680$. 
Par conséquent, $p(B)=\frac{1680}{15504}$.  La probabilité de gagner 10 euros est donc 
égale à $p(A)+p(B)= \frac{3500}{
15504}\simeq 0.225...$.
\fincorrection
\correction{006892}
\begin{enumerate}
\item 
Dés non pipés signifie que la probabilité de tirer 3 (ou tout autre chiffre) sur un 
dé est $1/6$. Les tirages étant indépendants, la probabilité
d'avoir trois 3 est $(1/6)^3 =1/216$.
\item
Il y a trois manières d'obtenir deux 2 et un 1, et 216 tirages possibles, donc la probabilité cherchée est $3/216=1/72$.
\item $6/216$
\item
Un total de 9 s'obtient par l'une des additions suivantes,
$$9=6+2+1=5+3+1=5+2+2=4+4+1=4+3+2=3+3+3,$$
où on a rangé les résultats d'un tirage par ordre décroissant. On rencontre 
l'addition 6+2+1 dans $3!=6$ tirages différents. De même pour 5+3+1 et 4+3+2. 
En revanche, 5+2+2 et 4+4+1 ne correspondent qu'à 3 tirages et 3+3+3 à un seul. 
Il y a donc 6+6+6+3+3+1=25 tirages qui donnent une somme de 9. La probabilité 
que la somme soit 9 vaut donc $25/216$.
\item
Un total de 10 s'obtient par l'une des additions suivantes, 
$$10=6+3+1=6+2+2=5+4+1=5+3+2=4+4+2=4+3+3.$$ 
On rencontre les additions 6+3+1, 5+4+1, 5+3+2 dans $3!=6$ tirages différents. 
En revanche, 6+2+2, 4+4+2 et 4+3+3 ne correspondent qu'à 3 tirages. Il y a donc 
6+6+6+3+3+3=27 tirages qui donnent une somme de 9. La probabilité que la somme 
soit 9 vaut donc $27/216$.
\end{enumerate}
\fincorrection
\nocorrection
\nocorrection
\correction{006895}
\begin{enumerate}
\item 
Soit $X$ la variable aléatoire ``nombre d'erreurs commises lors de la transmission 
de 5 bits''. Alors $X$ suit une loi binomiale $B(5;0,1)$. Recevoir une majorité de 1 
alors que 00000 a été émis correspond à l'événement $X\geq 3$. Sa probabilité est

$\displaystyle
P(X=3)+P(X=4)+P(X=5)={5\choose{3}}(0,1)^3 (0,9)^2 +{5\choose{4}}(0,1)^4 (0,9)^1 +
{5\choose{5}}(0,1)^5 (0,9)^0 
=0,0081+0,00045+0,00001=0,00856.
$
\item
Recevoir une majorité de 1 alors que 11111 a été émis correspond à l'événement $X\leq 2$, 
i.e. au complémentaire du précédent. Sa probabilité est donc $1-0,00856=0,99144$. 
Par conséquent, au prix de multiplier par 5 le temps de transmission, on améliore considérablement la fiabilité.
\end{enumerate}
\fincorrection
\correction{005992}
Notons les différents événements :
$Fe$ : <<être femme>>, $Lu$ : <<porter des lunettes>>, $H$ : <<être homme>>

Alors on a $P(Fe)=0.6,$ $P(Lu/Fe)=\frac{1}{3};$ il s'agit de la 
probabilité conditionnelle probabilité de 
<<porter des lunettes>> sachant que la personne est une femme.
De même, on a $P(Lu/H)=0.5$. On cherche la probabilité
conditionnelle $P(Fe/Lu)$.
D'après la formule des probabilités totales on a :
$P(Fe/Lu)P(Lu)=P(Lu/Fe)P(Fe)$ avec $P(Lu)=P(Lu/Fe)P(Fe)+P(Lu/H)P(H)$.

Application numérique : $P(Lu)=0.4$, donc  $P(Fe/Lu)=\frac{P(Lu/Fe)P(Fe)}{P(Lu)}=0.5$.
Remarque : on peut trouver les mêmes réponses par des raisonnements 
élémentaires.
\fincorrection
\correction{005993}
C'est évidemment le même que le précédent (exercice \ref{exo:quinio11}), seul le contexte
est différent : il suffit d'adapter les calculs faits.
En pronostiquant un enfant, le présentateur a une chance sur deux
environ de ne pas se tromper.
\fincorrection
\correction{005994}
Fumeurs

Définissons les événements: $F_{n}$ <<Fumer le $n$\up{ème} jour>>, et $\overline{F_{n}}$ 
l'événement complémentaire.
Alors $\{\overline{F_{n}},F_{n}\}$ constitue un système complet d'événements, 
$P_{n}=$ $P(F_{n})$; on peut donc écrire :
$P(\overline{F_{n+1}})=P(\overline{F_{n+1}}/F_{n})P(F_{n})
+P(\overline{F_{n+1}}/\overline{F_{n}})P(\overline{F_{n}})$.

Comme $P(\overline{F_{n+1}}/F_{n})=0.9$ et $P(\overline{F_{n+1}}/\overline{F_{n}})=0.3$
$1-P_{n+1}=0.9P_{n}+0.3(1-P_{n})$, soit $P_{n+1}=-0.6P_{n}+0.7$. Notons (R)
cette relation.

Pour connaître le comportement à long terme, il faut étudier cette
suite récurrente; il y a des techniques mathématiques pour ça,
c'est le moment de s'en servir.

Cherchons la solution de l'équation <<$\ell=-0.6\ell+0.7$>>, 
la limite éventuelle satisfait nécessairement cette équation : faire un passage à la limite dans la
relation (R), ou utiliser le théorème du point fixe.

On trouve $\ell=\frac{7}{16};$ alors, la suite $Q_{n}=(P_{n}- \ell)$ vérifie : 
$Q_{n+1}= - 0.6Q_{n}$, ce qui permet de 
conclure : $Q_{n+1}=(-0.6)^{n}Q_{1}$
et comme $((-0.6)^{n})$ est une suite qui tend vers $0$, on peut dire que la
suite $(Q_{n})$ tend vers $0$ et donc que la suite $(P_n)$ tend vers $\ell=\frac{7}{16}.$

Conclusion : la probabilité $P_{n}$ pour qu'elle fume le jour $J_{n}$
tend vers $\frac{7}{16} \simeq 0.4375$.
\fincorrection
\correction{005995}
 $P_{n+1}=P(E_{n+1})=
P(E_{n+1}/E_n)P(E_n)+P(E_{n+1}/\overline{E_n})P(\overline{E_n})=\frac{1}{10}P_{n}+\frac{4}{10}Q_{n}$.
Donc
$P_{n+1}=\frac{1}{10}P_{n}+\frac{4}{10}(1-P_{n})=\frac{4}{10}-\frac{3}{10}P_{n}$.

La suite ($P_{n}-\ell)$ est géométrique, où $\ell$ est solution
de $\frac{4}{10}-\frac{3}{10}\ell=\ell$ soit $\ell=\frac{4}{13}$.
Donc $P_{n}=\frac{4}{13}+(-\frac{3}{10})^{n-1} \times (a-\frac{4}{13})$.
\fincorrection
\correction{005996}
La probabilité d'avoir Princecharmant dans la barre B est $\frac{1}{5}; $ si j'achète $n$ barres, la probabilité de n'avoir la
figurine dans aucune des $n$ barres est $(\frac{4}{5})^{n}$, puisqu'il
s'agit de $n$ événements indépendants de probabilité $\frac{4}{5}$.
Je cherche donc $n$ tel que : $1-(\frac{4}{5})^{n}\geq 0.8$. On a facilement : $n\geq 8$.

Puis, je cherche $m$ tel que : $1-(\frac{4}{5})^{m}\geq 0.9$ ; il
faut au moins $11$ barres pour que la probabilité dépasse $90$\%.
Pour la probabilité $99$\%, $n\geq 21$ .
\fincorrection
\correction{005997}
\begin{enumerate}
\item Le taux global de personnes soulagées :
$P(S)=\frac{3}{5}0.75+\frac{2}{5}0.90=0.81$.

\item Probabilité pour un patient d'avoir pris de l'aspirine sachant qu'il
est soulagé :
$P(A/S)=P(A\cap S)/P(S)=P(A)P(S/A)/P(S)=\frac{\frac{3}{5}0.75}{0.81}=55.6\%$.
\end{enumerate}
\fincorrection
\correction{005998}
\begin{enumerate}
\item Probabilité conditionnelle : si un individu a les yeux bruns d'avoir
les cheveux blonds. 
C'est $P(CB/YB)=P(YB/CB)P(CB)/P(YB)$=$P(YB\cap CB)/P(YB)=\frac{0.15}{0.4}=0.375$.

\item La probabilité de l'événement : si un individu a les cheveux
blonds d'avoir les yeux bruns.
C'est $P(YB/CB)=P(YB\cap CB)/P(CB)$=$\frac{0.15}{0.25}=0.6$.

\item La probabilité de l'événement : si un individu a les cheveux
blonds,de ne pas avoir les yeux bruns.
C'est $P(\text{non}YB/CB)=1-P(YB/CB)=0.4$.
\end{enumerate}
\fincorrection
\correction{005999}
On obtient par calcul direct ou par événement contraire
la probabilité de voler : $1-p+p(1-q)^{2}$.
\fincorrection
\correction{006000}
\begin{enumerate}
\item La probabilité pour une personne d'être malade si son test est
positif est $P(M/T^{+})=P(T^{+}/M)P(M)/P(T^{+})$
or $P(T^{+})=P(T^{+}/M)P(M)+P(T^{+}/S)P(S)=0.95\cdot 0.03+0.1\cdot
0.97=0.125\,5$. 
D'où : $P(M/T^{+})=22.7\%$.

\item La probabilité pour une personne d'être saine si son test est
positif est $P(S/T^{+})=1-P(M/T^{+})=77.3\%$.

\item La probabilité pour une personne d'être malade si son test est négatif 
est $P(M/T^{-})=0.0017$.

\item La probabilité pour une personne d'être saine si son test est négatif
est $1-P(M/T^{-})=0.998=99.8\%$.
\end{enumerate}
\fincorrection
\correction{006001}
Une manière de résoudre le problème est la suivante:
puisqu'il y a $8$ clés et que j'écarte une après l'autre les mauvaises clés, 
je considère comme ensemble de toutes les possibilités, toutes les permutations de ces huit clés : 
il y en a $8$!. Alors la solution de chaque question est basée sur le même principe:
\begin{enumerate}
\item Les permutations (fictives) qui traduisent le cas (1) sont celles qui
peuvent être représentées par une suite : 
$BMMMMMMM$, la lettre $B$ désigne la bonne, $M$ désigne une
mauvaise. Il y a $7!$ permutations de ce type.
Donc $P(A)=\frac{7!}{8!}=\frac{1}{8},$ on s'en doutait!
\item De même, les permutations (fictives) sont
celles qui peuvent être représentées par une suite:
$MBMMMMMM$: il y en a encore $7$!, et la probabilité est la même.
\item Le raisonnement permet en fait de conclure que la probabilité, avant
de commencer, d'ouvrir la porte est la même pour le premier, deuxième,..., huitième essai.
\end{enumerate}
\fincorrection
\correction{006002}
\begin{enumerate}
\item L'univers des possibles est l'ensemble des couples possibles:
il y en a $6!=720$ (imaginez les dames assises et les hommes choisissant
leur partenaire). La probabilité $P(A)$ pour que chacun des $6$ hommes
danse avec son épouse légitime est, si chacun choisit au hasard, $\frac{1}{6!}$.

\item André danse avec son épouse, les autres choisissent au hasard: il y
a $5!$ permutations pour ces derniers:
$P(B)=\frac{5!}{6!}=\frac{1}{6}$.

\item André et René dansent avec leur épouse, les $4$ autres
choisissent au hasard: il y a $4!$ permutations pour ces derniers:
$P(C)=\frac{4!}{6!}=\frac{1}{30}$.

\item André ou René dansent avec leur épouse, les $4$ autres font ce
qu'ils veulent. Considérons les événements $D_{1}:$ <<André
danse avec son épouse>> ; $D_{2}$ : <<René danse avec son épouse>>. 
Alors $D=D_{1}\cup D_{2}$ et
 $P(D_{1}\cup D_{2})=P(D_{1})+P(D_{2})-P(D_{1}\cap D_{2})=\frac{3}{10}$.
\end{enumerate}
\fincorrection
\correction{006003}
\begin{enumerate}
  \item Combien de grilles ?
Il y en a $\binom{49}{6}=13\,983\,816$

  \item Combien de grilles avec 2 nombres consécutifs ?
Ce problème peut être résolu par astuce: considérer les numéros gagnants 
comme 6 places à <<choisir>> parmi 49.
En considérant des cloisons matérialisant les numéros gagnants,
c'est un problème de points et cloisons
Par exemple:
$$\mid \bullet \bullet \left\vert {}\right\vert \bullet
\left\vert \bullet \bullet \bullet \mid \bullet \bullet \right\vert$$
les gagnants sont: 1; 4; 5; 7; 11; 14.
Dans notre cas on ne veut pas de cloisons consécutives.
Les cinq cloisons séparent les numéros en 7 bo\^{\i}tes.
Les 5 bo\^{\i}tes intérieures étant non vides, on y met 5 points,
puis $38(=49-5-6)$ dans 7 bo\^{\i}tes.
Il y a $\frac{(38-1+7)!}{38!6!}=7.\,059\,1\times
10^{6}$ séquences ne comportant pas 2 nombres consécutifs.

D'où la probabilité d'avoir une grille
comportant 2 nombres consécutifs: $0.4952$.
\end{enumerate}
\fincorrection
\correction{006004}
\begin{enumerate}
\item $u_{n+1}=P(G_{n+1})=P(G_{n+1}/Gn)P(Gn)+P(G_{n+1}/\overline{G_n})P(\overline{G_n})
=0.6u_{n}+0.3v_{n}$.  

$v_{n+1}=0.4u_{n}+0.7v_{n}$.

Donc $\left( 
\begin{array}{c}
u_{n+1} \\ 
v_{n+1}\end{array}\right) =\left( 
\begin{array}{cc}
0.6 & 0.3 \\ 
0.4 & 0.7\end{array}\right) \left( 
\begin{array}{c}
u_{n} \\ 
v_{n}\end{array}\right)$

Comme $u_{n}+v_{n}=1$, $u_{n+1}=0.6u_{n}+0.3(1-u_{n})=0.3+0.3u_{n}$.
La suite $(u_{n}-\ell)$ est géométrique, où $\ell$ est solution
de $0.3+0.3\ell=\ell$, donc $\ell=\frac{3}{7}$.
Donc $u_{n}=\frac{3}{7}+u_{1}(0.3)^{n-1}=\frac{3}{7}+0.5 (0.3)^{n-1}$.
\end{enumerate}
\fincorrection
\nocorrection
\correction{006897}
On note $A$ l'événement ``l'enfant a la maladie M'', $A^c$ son
complémentaire (événement ``l'enfant n'a pas la maladie M''), $B$ l'événement
``l'enfant a une réaction positive au test'', $B^c$ son complémentaire
(événement
``l'enfant a une réaction négative au test''). D'après
l'énoncé, $P(A)=0,01$, $P(A^c)=0,99$, $P(B^c|A^c)=0,9$ et
$P(B|A)=0,95$. Donc,  
\begin{eqnarray*}P(B)&=&P(A\cap B)+P(A^c\cap
B)=P(B|A)P(A)+P(B|A^c)P(A^c)\\
&=&P(B|A)P(A)+(1-P(B^c|A^c))P(A^c)=0,0095+0,099=0,1085.
\end{eqnarray*}
La probabilité
qu'un enfant de moins de trois mois pris au hasard et ayant une réaction 
positive soit atteint par M est donnée par $P(A|B)=\frac{P(B\cap
A)}{P(B)}=\frac{P(B|A)P(A)}{P(B)}=\frac{95}{1085}\simeq 0,088$. Un tel test serait d'une utilité discutable.
\fincorrection
\nocorrection
\nocorrection
\correction{006900}
La méthode la plus simple consiste à introduire la population
totale $N$ et à compter les daltoniens. Soit $d_H=5\%$, $d_F=0,25\%$
(taux de daltoniens chez les hommes et les femmes),
$p_H=48\%$, $p_F=52\%$ (proportions d'hommes et de femmes dans la population).
Le nombre d'hommes est $Np_H$, le nombre d'hommes daltoniens
est $Np_Hd_H$. De même, le nombre de femmes daltoniennes est
$Np_Fd_F$. La proportion de daltoniens hommes parmi les daltoniens est
donc
$$
\frac{\mbox{nombre de daltoniens hommes}}{\mbox{nombre de daltoniens}}=
\frac{N p_Hd_H}{Np_Hd_H+Np_Fd_F}=\frac{p_Hd_H}{p_Hd_H+p_Fd_F}\approx 0,95.
$$
La probabilité pour qu'un daltonien soit un homme est d'environ 95\%.

\medskip
Une formulation plus élaborée (mais strictement équivalente) 
consiste à utiliser la formule de Bayes.
Soit $H$ l'événement ``être un homme'' et $D$ l'événement ``être daltonien''.
On veut calculer $P(H|D)$. Selon la formule de Bayes,
$P(H|D)=\frac{P(D|H)P(H)}{P(D)}.$ En utilisant que
$P(D)=P(D|H)P(H)+P(D|F)P(F)$ (formule des probabilités totales), on
obtient
$$P(H|D)=\frac{P(D|H)P(H)}{P(D|H)P(H)+P(D|F)P(F)}=\frac{p_Hd_H}{p_Hd_H+p_Fd_F}
.$$
\fincorrection
\correction{006901}
Soit $U_1$ l'événement ``on tire la boule dans la première urne''
et $U_2$ l'événement ``on tire la boule dans la seconde urne''. Le choix
de l'urne étant équiprobable, on a : $P(U_1)=P(U_2)=0,5$.
Soit $B$ l'événement ``on tire une boule blanche''.
L'énoncé donne les probabilités conditionnelles suivantes : \\
$P(B|U_1)=30/40=0,75$ et $P(B|U_2)=20/40=0,5$. 

On cherche la probabilité $P(U_1|B)$.
La formule de Bayes,
appliquée à la partition $(H_1,H_2)$, nous donne :
$$
P(U_1|B)=\frac{P(B|U_1)P(U_1)}{P(B|U_1)P(U_1)+P(B|U_2)P(U_2)}=
\frac{0,75\times 0,5}{0,75\times 0,5+0,5\times 0,5}=0,6$$
(probabilité
a posteriori)

\medskip
\textit{Interprétation : avant de regarder la couleur de la boule, la probabilité d'avoir choisi la 
première urne est une probabilité a priori $P(U_1)$ soit 50 \%. 
Après avoir regardé la boule, on révise notre jugement et on considère 
$P(U_1|B)$, soit 60 \%.}
\fincorrection
\correction{006005}
\begin{enumerate}
\item On utilise une loi binomiale, loi de la variable aléatoire:
<<nombre de lettres affranchies au tarif urgent parmi 4 lettres>>
$n=5$, $p=\frac{3}{5}$. 
On obtient $P(A)=1-(\frac{2}{5})^{4}=0.9744$, 
$P(B)=\binom{4}{2}(\frac{2}{5})^{2}(\frac{3}{5})^{2}=0.3456$.

\item La loi de probabilité de $X$ est une loi binomiale, loi de la variable
aléatoire: <<nombre de lettres affranchies au tarif
urgent parmi 10 lettres>>.
$n=10$, $p=\frac{3}{5}$, son espérance est $np=6$, sa variance est $np(1-p)=\frac{12}{5}$.
\end{enumerate}
\fincorrection
\correction{006006}
On utilise une loi hypergéométrique

$P(A)=1-\frac{\binom{10}{3}}{\binom{15}{3}}=
0.736\,26$

$P(B)=\frac{\binom{5}{3}}{\binom{15}{3}}=2.\,
197\,8\times 10^{-2}$

$P(C)=\frac{\binom{5}{1}\binom{10}{2}}{\binom{15}{3}}=0.494\,51$
\fincorrection
\correction{006007}
Soit $X$ la variable aléatoire nombre de clients
qui viennent après réservation parmi 20.
La loi de $X$ est une loi binomiale de paramètres $n=20$, $p=0.75$.
Son espérance est $np=15$, son écart-type est $\sqrt{np(1-p)}=\sqrt{15\cdot 0.25}$.
La probabilité pour que $X$ soit égal à 15 est $\binom{20}{15}0.75^{15}0.25^{5}=0.202\,33$.
\fincorrection
\correction{006008}
La variable aléatoire associée à ce problème est $X$ 
<<nombre de sujets révisés parmi les $3$>> ; son support est l'ensemble
$\{0,1,2,3\}$. La loi de $X$ est une loi hypergéométrique puisque
l'événement $[X=k]$, pour $k$ compris entre $0$ et $3$, se produit si
le candidat tire $k$ sujet(s) parmi les $60$ révisés, et $3-k$ sujets
parmi les $40$ non révisés.

Alors:
\begin{enumerate}
\item Les trois sujets tirés ont été révisés : $P[X=3]=\frac{\binom{60}{3}}{\binom{100}{3}}$.
\item Deux des trois sujets tirés ont été révisés: $P[X=2]=\frac{\binom{60}{2}.\binom{40}{1}}{\binom{100}{3}}$.
\item Aucun des trois sujets: $P[X=0]=\frac{\binom{40}{3}}{\binom{100}{3}}$.
\end{enumerate}
La loi de probabilité de $X$ est donnée sur le support $\{0,1,2,3\}$ par:
\begin{equation*}
P[X=k]=\frac{\binom{60}{k}.\binom{40}{3-k}}{\binom{100}{3}}
\end{equation*}

Résultats numériques:

$k=0:P[X=0]$ $\simeq 6.\,110\times 10^{-2}$

$k=1:P[X=1] \simeq 0.289\,$

$k=2:P[X=2]\simeq 0.438$

$k=3:P[X=3] \simeq 0.212$

L'espérance est $E(X)=1.8$ (selon la formule $E(X)=np$).
\fincorrection
\correction{006009}
Puisque les réponses sont données au hasard, chaque grille-réponses est en fait la 
répétition indépendante de $20$ épreuves aléatoires (il y a $4^{20}$ grilles-réponses). 
Pour chaque question la probabilité de succès est de $\frac{1}{4}$ et
l'examinateur fait le compte des succès: la variable aléatoire $X$,
nombre de bonnes réponses, obéit à une loi binomiale donc on a
directement les résultats.
Pour toute valeur de $k$ comprise entre $0$ et $20$:
$P[X=k]=C_{20}^{k}(\frac{1}{4})^{k}(1-\frac{1}{4})^{20-k}$, ce qui
donne la loi de cette variable aléatoire.

Quelle est l'espérance d'un candidat fumiste? C'est $E(X)=np=5$
\fincorrection
\correction{006010}
Une variable aléatoire adaptée à ce problème est le nombre $X$ de personnes 
se présentant au guichet entre 10h et 11h. Compte tenu
des hypothèses, on partage l'heure en $60$ minutes. Alors $X$ suit une
loi binomiale de paramètres $n=60$ et $p=0.1$. On est dans le cas de
processus poissonnien : on peut approcher la loi de $X$ par la loi de
Poisson de paramètre $\lambda =60\times 0.1=6$.
L'espérance de $X$ est donc $E(X)=6$;

On peut alors calculer les probabilités demandées:
$P[X=k]=\frac{6^{k}e^{-6}}{k!}$. Valeurs lues dans une table ou calculées :
$P[X=3]\simeq 0.9\%;$ $P[X=4]\simeq 13.4\%;$ $P[X=5]=P[X=6]\simeq
\,16.1\%;$
$P[X=7]\,\simeq 13.8\%;P[X=8]\simeq 10.3\,\%.$

Remarque : de façon générale si le paramètre 
$\lambda$ d'une loi de Poisson est un entier $K$, on a:
$P[X=K-1]=\frac{K^{K-1}e^{-K}}{(K-1)!}=\frac{K^{K}e^{-K}}{K!}=P[X=K]\,.$

Calculons maintenant la probabilité pour que au moins $10$
personnes se présentent au guichet entre $10$h et $11$h: C'est $P[X\geq
10]=1-\sum_{k=0}^{9}\frac{6^{k}e^{-6}}{k!}\simeq 8.392\times
10^{-2}.$

\fincorrection
\correction{006011}
La probabilité $p=\frac{1}{100}$ étant faible, on peut appliquer la
loi de Poisson d'espérance $100p=1$ au nombre $X$ de centenaires pris
parmi cent personnes. On cherche donc: $P[X\geq 1]=1-P[X=0]=1-e^{-1}\simeq
63\%$.

Sur un groupe de $200$ personnes: l'espérance est 2 donc: $P[X'\geq 1]=1-e^{-2}\simeq 86\%.$
La probabilité des événements : $[X' =1]$ et $[X' =2]$
sont les mêmes et valent: $0.14$.
Ainsi, sur 200 personnes, la probabilité de trouver exactement un
centenaire vaut $0.14$, égale à la probabilité de trouver
exactement deux centenaires. Cette valeur correspond au maximum de probabilité 
pour une loi de Poisson d'espérance $2$ et se généralise.
Si $X$ obeit à une loi de Poisson d'espérance $K$, alors le maximum de
probabilité est obtenu pour les événements $[X=K-1]$ et $[X=K].$
\fincorrection
\correction{006012}
\begin{enumerate}
\item 30\% est la probabilité de l'événement Panne, noté $Pa$;
la probabilité pour une machine donnée de plus de cinq ans, d'être hors d'usage est
$P(HU)=P(HU/Pa) P(Pa)+P(HU/nonPa) P(nonPa)=0.3\cdot 0.75+0.4\cdot
0.7=0.505$.

\item La probabilité pour une machine hors d'usage de n'avoir jamais eu de
panne auparavant est 
$P(\text{non}\,Pa/HU)=P(HU/\text{non}\,Pa) P(\text{non}\,Pa)/P(HU)=0.4\cdot 0.7/
0.505\,=0.554\,46$.

\item La loi de probabilité de $X$ est une loi binomiale, $n= 10$, $p=0.3$, espérance $3$.

\item $P[X=5]=\binom{10}{5}(0.3)^{5}(0.7)^{5}=0.102\,92$
\end{enumerate}
\fincorrection
\correction{006013}
Le nombre $X$ de personnes mesurant plus de 1.90m parmi 100 obéit à
une loi de Poisson de paramètre $\frac{100}{80}$.

La probabilité qu'il y ait au moins une personne mesurant plus de 1.90m
est donc $1-P[X=0]=1-e^{-\frac{100}{80}}=1-e^{-\frac{5}{4}}=0.713\,50$.

Sur 300 personnes:
la probabilité qu'il y ait au moins une personne mesurant plus de 1.90m
est donc $1-P[Y=0]=1-e^{-\frac{300}{80}}=0.976\,48$.

\fincorrection
\nocorrection
\nocorrection
\correction{006904}
\begin{enumerate}
\item
Une lampe tirée au hasard a une probabilité de $0,2$ d'avoir une durée de vie 
inférieure à 3000 heures. Le nombre $X$ de lampes qui ont une durée de vie 
inférieure à 3000 heures dans un échantillon de taille 15 tiré au hasard est 
la somme de 15 variables de Bernoulli de paramètre $p=0,2$. Par conséquent, 
il suit une loi binomiale $B(15;0,2)$. Son espérance vaut $E(X)=15\times 0,2=3$.
\item
C'est $p(X=0)={15\choose{0}}(0,2)^0 (0,8)^{15}\sim 0,0352$.
\item
C'est $p(X=0)+p(X=1)+p(X=2)$
\begin{eqnarray*}
&=&{15\choose{0}}(0,2)^0 (0,8)^{15}+{15\choose{1}}(0,2)^1 (0,8)^{14}+
{15\choose{2}}(0,2)^2 (0,8)^{13}\\&=&0,0352+0,1319+0,2309\simeq 0,398.
\end{eqnarray*}
\end{enumerate}
\fincorrection
\correction{006905}
\begin{enumerate}
\item
La loi de $X$ est une loi binomiale de paramètres $n=23$, $p=0,75$ : 
$P(X=k)={k\choose{n}}p^k(1-p)^{n-k}$ si $0\le k\le n$.
Son espérance est $np=17,25$.
\item
$P(X\le 20)=1-P(X\in\{21, 22, 23\})\simeq 0,951$.
\end{enumerate}
\fincorrection
\correction{006906}
\begin{enumerate}
\item $Y\sim B(10,1/2)$, $X=Y/10$. 
\item $P(X>0,5)=P(Y>5)=P(Y=6,7,8,9,10)\simeq 0,377$.
\item $P(0,4\le X\le 0,6)=P(4\le Y\le 6)=P(Y=4,5,6)\simeq 0,656$
\item $P(3\le Y\le 7)\simeq 0,891$. $P(2\le Y\le 8)\simeq 0,978$. Donc $a=3$.
\item Oui. Non.
\end{enumerate}
\fincorrection
\correction{006907}
\begin{enumerate}
\item
C'est une loi de Poisson de paramètre $6$ : $P(X=n)=e^{-6}\frac{6^n}{n!}$.
La probabilité qu'il n'y ait aucun appel est $p(X=0)=e^{-6}\simeq 0,002$.
\item
Soit $Y$ la variable aléatoire ``Nombre d'appels reçus en 2 minutes''. 
Alors $Y$ suit une loi de Poisson de paramètre 4. La probabilité qu'il 
y ait entre 0 et 4 appels est $P(Y=0)+P(Y=1)+P(Y=2)+P(Y=3)+P(Y=4)\simeq
0,629$. Donc $P(Y\ge 5)=1-P(Y\le 4)\simeq 0,371$.
\end{enumerate}
\fincorrection
\correction{006908}
\begin{enumerate}
\item
$P(X=k)=0,5^k 0,5= 0,5^{k+1}$ si $0\le k\le 4$, $P(X=5)=0,5^5$.
\item $E(X)=\sum_{0\le k\le 5} kP(X=k)= \frac 14+ 2\frac 18+3\frac 1{16}+4\frac{1}{32}+5\frac{1}{32}=0,96875$.

$Y=$ nombre de garçons. Il y a exactement 1 garçon sauf s'il y a 5 filles.
$P(Y=0)=0,5^5=1/32$. $E(Y)=P(Y=1)\times 1=0,96875$.  $E(X)=E(Y)$, donc pas
efficace.
\end{enumerate}
\fincorrection
\nocorrection
\nocorrection
\nocorrection
\nocorrection
\nocorrection
\correction{006914}
\begin{enumerate}
\item 
Soit $A$ l'événement ``Charles n'a pas de chien'' et $B$ l'événement 
``Sophie n'a pas de chat''. L'énoncé donne $P(B|A)=0,9$. Par conséquent, 
la probabilité pour que le ménage n'ait aucun animal est $P(A\cap B)=P(B|A)P(A)=0,9\times 0,8=0,72$.
\item
\begin{enumerate}
\item 
$Z$ ne peut prendre que les valeurs 0, 1 et 2. L'événement $\{Z=0\}$ 
coïncide avec $A\cap B$, donc $P(Z=0)=0,72$. Il vient $P(Z=2)=1-P(Z=0)-P(Z=1)=1-0,72-0,1=0,18$.
\item
$E(Z)=0 \cdot P(Z=0)+1 \cdot P(Z=1)+2 \cdot P(Z=2)=0,1+0,36=0,46$. $E(Z^2 )=0^2 P(Z=0)+1^2 
P(Z=1)+2^2 P(Z=2)=0,1+0,72=0,82$, d'où $\text{Var}(Z)=E(Z^2 )-E(Z)^2 =0,6084$, $\sigma(Z)=\sqrt{\text{Var} (Z)}=0,78$.
\item
On calcule $P(X=0\mbox{ et } Y=0)=P(Z=0)=0,72$, $P(X=1\mbox{ et } Y=1)=P(Z=2)=0,18$, 
$
P(X=0\mbox{ et }Y=1)=P(A\cap B^c)=P(B^c|A)P(A)=(1-P(B|A))P(A)=0,1\times 0,8=0,08. 
$
On complète le tableau en utilisant le fait que le somme des probabilités des événements élémentaires vaut 1.
\begin{center}
\begin{tabular}{|c|c|c|c|}
\hline
&$Y=0$&$Y=1$&Total\\\hline
$X=0$&0,72&0,08&0,8\\\hline
$X=1$&0,02&0,18&0,2\\\hline
Total&0,74&0,26&1\\
\hline
\end{tabular}
\end{center}
La dernière ligne du tableau donne la loi de $Y$, $P(Y=0)=0,74$ et $P(Y=1)=0,26$.
\item
On constate que $P(X=1\mbox{ et } Y=1)=0,18$ n'est pas égal à 
$P(X=1)P(Y=1)=0,2\times 0,26$, donc $X$ et $Y$ ne sont pas indépendantes.
\end{enumerate}\end{enumerate}
\fincorrection
\nocorrection
\correction{006014}
La probabilité qu'une bille soit rejetée est, en notant $D$ la
variable aléatoire <<diamètre>>, $p=1-P[7.97\leq D\leq 8.03]$.
Or $P[7.97\leq D\leq 8.03]=P[-\frac{0.03}{0.02}\leq \frac{D-8}{0.02}\leq 
\frac{0.03}{0.02}]=F(1.5)-F(-1.5)=0.866\,4$.
La proportion de billes rejetées est donc $p=13.4\%$.
\fincorrection
\correction{006015}
\begin{enumerate}
\item La probabilité pour que $X$ soit inférieur à 0.36mm est :
$P[X\leq 0.36]=P[\frac{X-0.3}{0.1}\leq 0.6]=0.726,$ soit $72.6$\%.

 La probabilité pour que $X$ soit compris entre $0.25$ et $0.35\text{mm}$ est
$P[0.25\leq X\leq 0.35]=2 F(0.5)-1= 0.383$, soit $38.3$\%.

\item Pour $n = 20$, la loi de $Z=\sum X_{i}$ est une loi normale de paramètres: d'espérance $E(Z)=20 m = 6$
et de variance $\text{Var}\,Z=20 \sigma^2 = 0.2$.
\end{enumerate}
\fincorrection
\correction{006016}
Pour $n = 2000$, la loi suivie par la variable aléatoire $N$
<<nombre de plaques inutilisables parmi les 2000>> est une loi de Poisson de paramètre 2:
alors $P[N\leq 3]=0.86$.

Remarquons qu'en faisant l'approximation par une loi normale et en employant
le théorème central limite, on obtient: $P[N\leq 3]\simeq 0.76,$ et
avec correction de continuité on obtient $P[N\leq 3]\simeq 0.85.$
\fincorrection
\correction{006017}
Par des méthodes analogues on trouve que la probabilité pour que $X$ 
soit compris entre 6.3mm et 6.6 mm est
$14.3$.
\fincorrection
\correction{006018}
Si $X$ est de moyenne $m$ et d'écart-type $\sigma$ alors 
$Y=\frac{X-m}{\sigma}$ suit une loi centrée réduite.
Donc si $P[X\leq 165]$ alors $P[\frac{X-m}{\sigma }\leq \frac{165-m}{\sigma }]=0,56$.
Or on peut lire dans la table de Gauss $F(0.15)=0.5596$.

De même, si $P[X\geq 180]$ alors $P[\frac{X-m}{\sigma }\geq \frac{180-m}{\sigma }]=0.1$.
Donc $P[\frac{X-m}{\sigma }\leq \frac{180-m}{\sigma }]=0.9$
et l'on peut lire de même $F(1.28)=0.8997.$

Pour trouver $m$ et $\sigma$ il suffit de résoudre le système d'équations:
$\frac{165-m}{\sigma}=0.15$ et $\frac{180-m}{\sigma}=1.28$ d'où
$\sigma \simeq 13.27$, $m\simeq 163$ cg.
Alors,
$P[X\geq 182]=P[\frac{X-m}{\sigma }\geq \frac{182-m}{\sigma }]=1-F(1.\,43)=0.0764$.

Sur $10\,000$ personnes on estime le nombre de personnes à soigner de
l'ordre de $764$ personnes ; en fait la théorie de l'estimation donnera
une fourchette.
\fincorrection
\correction{006019}
\begin{enumerate}
\item Loi binomiale $B(365;\frac{4}{365}),$ approchée par la loi de Poisson
de paramètre 4, d'espérance et variance 4.

\item Loi binomiale $B(6;\frac{1}{2})$, d'espérance 3 et variance $\frac{3}{2}$.

\item Loi hypergéométrique.
\end{enumerate}
\fincorrection
\correction{006020}
\begin{enumerate}
\item La loi de $X$ est la loi binomiale $B(1000;0.02)$, d'espérance 20, 
d'écart-type $\sqrt{19.6}$.

\item En approchant cette loi par celle d'une loi normale de paramètre
$m=20$, écart-type $\sqrt{19.6}$.
$P[18\leq X\leq 22]=P[(17.5-20)/\sqrt{19.6}\leq (X-20)/\sqrt{19.6}\leq
(22.5-20)/\sqrt{19.6}]\simeq 0.428$.

Sans correction de continuité on trouve $P[(17-20)/\sqrt{19.6}\leq
(X-20)/\sqrt{19.6}\leq (22-20)/\sqrt{19.6}]\simeq 0.348$.

Approchée par la loi de Poisson de paramètres : espérance 20 et
variance 20, on trouve $P[18\leq X\leq 22]\simeq 0.423$.

Enfin par la vraie loi binomiale: on trouve $P[18\leq X\leq 22]\simeq 0.427$.
\end{enumerate}
\fincorrection
\correction{006021}
\begin{enumerate}
\item Soit $F$ l'événement <<la pièce est fausse>>; soit $U$ l'événement <<la pièce est un euro>>; 
soit $D$ l'événement <<la pièce est deux euros>>. Alors
$P(F)=P(F/U)P(U)+P(F/D)P(D)=2.9\%$.

\item On cherche $P(U/F)=$($P(F/U)P(U))/P(F)=51.7\%$.

\item $X$ la variable aléatoire <<nombre de pièces 
fausses parmi 1000>> obéit à une loi binomiale$B(1000;5\%)$.
Espérance: 50; écart-type: $\sigma =\sqrt{47.5}$.
En approchant cette loi par une loi normale $N(50; \sigma)$,
la probabilité pour que $X$ soit compris entre 48 et 52 est :
$P[(47.5-50)/\sigma \leq (X-50)/\sigma \leq (52.5-50)/\sigma ]\simeq 28.3\%$.
\end{enumerate}
\fincorrection
\correction{006022}
\begin{enumerate}
\item La loi de $X$ est une loi binomiale $B(180;\frac{1}{6})$
Espérance: 30; écart-type: $\sigma =\sqrt{25}=5$.

\item En approchant cette loi par une loi normale $N(30$; $\sigma)$ la
probabilité pour que $X$ soit compris entre 29 et 32 :
$P[(28.5-30)/\sigma \leq (X-30)/\sigma \leq (32.5-30)/\sigma ]\simeq 30.94\%$.
Avec la vraie loi, on trouve la probabilité pour que $X$ soit compris entre 29 et 32 est $30.86\%$.
\end{enumerate}
\fincorrection
\correction{006023}
\begin{enumerate}
\item Lorsque l'on tire un bulletin au hasard, la probabilité que ce soit
un bulletin pour A est de $0.2.$

\item Il y a suffisamment de bulletins de vote en tout pour que l'on puisse
assimiler ces tirages à des tirages avec remise; alors la loi de
probabilité de $X$ est une loi binomiale de paramètres $n=200$ et $p=0.2;$ or
 $np=40$; on peut faire l'approximation normale. L'espérance de $X$ est 
donc $m=40$ et l'écart-type: $\sqrt{40\times 0.8}=4\sqrt{2}$.

\item $P[X\geq 45]=1-P[X\leq 44]\simeq 1-F(\frac{44.5-40}{4\sqrt{2}})\simeq 21\%$,
c'est la probabilité pour que le nombre de voix pour A
soit supérieur à $45$ dans un lot de $200$ bulletins.
De même, $P[30\leq X\leq 50]\simeq F(\frac{50.5-m}{\sigma })-F(\frac{29.\text{$5$}-m}{\sigma })\simeq 93.6\%$.

\item Reprenons le calcul pour le candidat B qui n'a obtenu que $2$\% des voix.
Alors pour $n=100$ et $p=0.02$ l'approximation par une loi de Poisson d'espérance $\lambda =2$ est légitime.
On peut dire que $P[Y\geq 5]=1-P[Y\leq 4]=1-\sum_{k=0}^{4}\frac{e^{-2}2^{k}}{k!}$, de l'ordre de $5$\%.

Enfin $P[1\leq Y\leq 4]=\sum_{k=1}^{4}\frac{e^{-2}2^{k}}{k!}\simeq 0.812$.
\end{enumerate}
\fincorrection
\correction{006024}
\begin{enumerate}
\item
\begin{enumerate}
\item Pour calculer la probabilité que Monsieur A soit contrôlé
entre 60 et 80 fois dans l'année, posons le nombre de contrôles
comme une variable aléatoire. Elle obéit à une loi binomiale
$B(700;0.1)$.
On peut l'approcher par la loi normale $N(70;\sqrt{63})$.

$P[60\leq X\leq 80]=P[-10/\sqrt{63}\leq X\leq 10/\sqrt{63}]\simeq 2F(10.5/\sqrt{63})-1=0.814$.
La probabilité d'être contrôlé entre 60 et 80 fois dans l'année est $81.4$.
\item Calculons le prix que devrait payer le voyageur: $1,12 \times 700=
784$ euros. Il est perdant si l'amende dépasse ce prix.
Or l'amende est $aX$, si $a$ est l'amende fixée par la compagnie.

On cherche donc $a$ pour que : $P[aX\geq 784]\geq 0.75:$
Soit $P[aX\leq 784]\leq 0.25:$
Par lecture de table:
$a=784/64.642=12.128$
Il faut que l'amende dépasse 13 euros.
\end{enumerate}

\item Calculons le prix que devrait payer le voyageur: $1,12 \times 300=
336$ euros
Il est perdant si l'amende dépasse ce prix.
Or l'amende est $bX$, si $b$ est l'amende fixée par la compagnie.
$X$ obéit à une loi binomiale $B(300;0.5)$.
On cherche donc $b$ pour que : $P[bX\geq 336]\geq 0.75$.
Par un raisonnement analogue, on obtient cette fois le résultat:
il suffit que l'amende dépasse 2 euros 30 !
\end{enumerate}
\fincorrection
\correction{006916}
\begin{enumerate}
\item
$E(X_1)=3.5$, $\text{Var}(X_1)=\frac{35}{12}\simeq 2,92$.
\item
$X_1,X_2$ sont indépendantes de même loi.
D'où $E(S)=2E(X_1)=7$ et $\text{Var}(S)=2\text{Var}(X_1)=\frac{35}{6}$.
\end{enumerate} 
\fincorrection
\nocorrection
\nocorrection
\correction{006919}
$E(X)=2$, $E(Y)=1/2$, $E(XY)=3/2$, d'où $\text{Cov}(X,Y)=1/2$.
$\text{Var}(X)=2$, $\text{Var}(Y)=1/2$, d'où $\text{Corr}(X,Y)=\frac{\text{Cov}(X,Y)}{\sqrt{\text{Var}(X)\text{Var}(Y)}}=
1/2$.
\fincorrection
\correction{006920}
$U$ et $V$ étant à valeurs dans $\{1, \dots ,N\}$, elles admettent des moments de tout ordre, donc $X$ et $Y$ aussi, les calculs qui suivent sont donc justifiés.
$\text{Cov}(X,Y)=E(XY)-E(X)E(Y)=E(XY)$ car $E(X)=E(U)-E(V)=0$ ($U, V$ ont
même loi).
$E(XY)=E((U-V)(U+V))=E(U^2)-E(V^2)=0$ ($U, V$ ont
même loi). D'où $\text{Cov}(X,Y)=0$.
\fincorrection
\correction{006924}
Il faut que Mme Michel arrive entre 9h30 et 10h30. Intervalle de
temps d'1h sur 4h, donc proba = 1/4.
\fincorrection
\correction{006925}
\def\I1{{ \rm 1\:\!\!\! l}}
On donne 2 méthodes graphiques et un calcul direct.

On peut prendre comme modèle $\Omega=[0,1]^2$, $P$ la mesure de 
Lebesgue sur le carré, $X(x,y)=x$, $Y(x,y)=y$. On peut alors 
représenter $Z=t$ par la droite d'équation $x+y=t$ dans le carré (figure
de gauche).

Si $0\leq t\leq 1$ alors $F_Z(t)=P(X+Y\leq t)=\frac{t^2}{2}$ (figure
de gauche).
Si $1\leq t\leq 2$, alors $F_Z(t)=1-\frac{(2-t)^2}{2}$ (figure du milieu).
Si $t<0$ alors $F_Z(t)=0$ et si $t>2$ alors $F_Z(t)=1$.

$F_Z$ est continue sur $\Rr$, $C^1$ par morceaux, on peut donc dériver
pour obtenir $f_Z$ la densité de $Z$ : $f_Z(t)=t$ si $0\leq t\leq 1$,
$f_Z(t)=2-t$ si $1\leq t\leq 2$, $f_Z(t)=0$ sinon.

\medskip
On peut également calculer la densité directement, ``à la
physicienne'', en considérant $dt$ comme un tout petit accroissement
(on l'a dessiné assez gros pour rendre le dessin lisible). 
Sur la figure de droite,
on a dessiné $X+Y\in [t,t+dt]$ pour $0\leq t\leq 1$. L'aire de la
partie hachurée est à peu près $t\sqrt{2}\times\frac{dt}{\sqrt{2}}$ (longueur
$\times$ largeur), donc on  a $P(t\leq X+Y\leq t+dt)=tdt$. Or
$P(t\leq Z\leq t+dt)=f_Z(t)dt$, ce qui donne la
densité $f_Z(t)=t$ pour $0\leq t\leq 1$. On fait pareil pour $1\leq t\leq 2$.

%\bigskip
\centerline{\includegraphics{../images/img006925-1}}
%\vspace{1cm}

\medskip
Calcul direct : La densité de $X$ et $Y$ est $f(t)=\I1_{[0,1]}(t)$. Par
indépendance, $Z=X+Y$ a une densité $f_Z(t)=f*f(t)=\displaystyle
\int_{\Rr} f(t-x)f(x)dx$ (résultat du cours, qu'on peut redémontrer
par changement de variable à partir de la définition de $P_X*P_Y$).

$f(t-x)f(x)=1$ si $0\leq t-x\leq 1$ et $0\leq x\leq 1$,
$f(t-x)f(x)=0$ sinon.

Soit $D_t=\{x\in\Rr\mid 0\leq t-x\leq 1 \mbox{ et } 0\leq x\leq 1\}=
\{x\in\Rr\mid  t-1\leq x\leq t\mbox{ et }0\leq x\leq 1\}$. $D_t=
[\max (0,t-1),\min(t,1)]$. On a
$f_Z(t)=\displaystyle\int_{D_t} 1 \, dx$, alors :\\
-- si $0\leq t\leq 1$, $D_t=[0,t]$ et $f_Z(t)=t$,\\
-- si $1\leq t\leq 2$, $D_t=[t-1,1]$ et $f_Z(t)=2-t$,\\
-- si $t<0$ ou $t>2$, $D_t=\emptyset$ et $f_Z(t)=0$.
\fincorrection
\correction{006926}
\def\I1{{ \rm 1\:\!\!\! l}}
$P(Z>t)=P(X_1>t, X_2>T,\ldots, X_n>t)=P(X_1>t)P(X_2>t)\cdots P(X_n>t)$
par indépendance. Comme $P(X>t)=1-P(X\leq t)$, on a
$\displaystyle
1-F_Z(t)=\prod_{1\leq i\leq n} (1-F_{X_i}(t)).$\\
Comme $F_{X_i}(t)=1-e^{-t}$ si $t\geq 0$ et $F_{X_i}(t)=0$ sinon, on trouve
que $F_Z(t)=1-e^{-nt}$ si $t\geq 0$ et $F_Z(t)=0$ sinon. 
La densité de $Z$ est donc $\I1_{\Rr_+}ne^{-nt}$. Par conséquent, 
$Z$ est de loi exponentielle de paramètre $n$.
\fincorrection
\correction{006927}
\begin{enumerate}
\item
Inégalité de Markov : $P(X\ge 75)\le \frac{E(X)}{75}=\frac 23$.
\item
Inégalité de Bienaymé-Tchebychev : $P(|X-50|\ge 25)\le \frac{\text{Var}(X)}{25^2}=
\frac{5^2}{25^2}=0,04$. Donc $P(X\ge 75)\le 0,04$.
\end{enumerate} 
\fincorrection
\correction{006928}
Inégalité de Bienaymé-Tchebychev : 
$P(|X-10^2|\ge 10^3-10^2)\le \frac{\text{Var}(X)}{(10^3-10^2)^2}\le 
\frac{\text{Var}(X)}{(10^3)^2}=10^{-4}$. Donc $P(X\ge 10^3)\le 10^{-4}$.
\fincorrection
\nocorrection
\nocorrection
\correction{006923}
\begin{enumerate}
\item
L'énoncé suggère que le poids en grammes des paquets est une variable aléatoire 
qui suit une loi normale d'espérance 500 et d'écart-type 25. Soit $X$ la variable aléatoire correspondante, et $Y=(X-500)/25$. 
\item
 
$P(480\leq X\leq 520)=P(|Y|\leq  0,8)=0,576.$
On s'attend donc à ce que, sur 1000 paquets, il y en ait 576 dont le poids est compris entre 480g et 520g.
\item\begin{eqnarray*}
P(480\leq X\leq 490)&=&P(-0,8\leq Y\leq -0,4)\\
&=&P(0,4\leq Y\leq 0,8)\\
&=&p(Y\leq 0,8)-p(Y\leq 0,4)=0,1327.
\end{eqnarray*}
On s'attend donc à ce que, sur 1000 paquets, il y en ait 132 dont le poids est compris entre 480g et 490g.
\item
$P(450\leq X)=0,5+P(0\le Y\leq 2)=0,5+\frac 12 P(-2\le Y\leq 2)=0,5+0,4772=0,9772.$
On s'attend donc à ce que, sur 1000 paquets, il y en ait 977 dont le poids est supérieur à 450g.
\item
Il faut trouver $t$ tel que $p(|Y|<t)=0,9$. La table donne $t=1,645$, 
puis $a=25t=41$. Par conséquent, environ 90\% de la production a un poids 
compris entre $500-41=459$g et $500+41=541$g.
\end{enumerate}
\fincorrection
\correction{006954}
Voici plusieurs façons de définir une corde, qui mènent à choisir un point
au hasard dans un espace différent. Même avec la probabilité ``naturelle''
(probabilité uniforme), on obtient des résultats différents.

\begin{enumerate}
  \item La corde $[AB]$ est déterminée par sa distance au centre $O$, qui
est $r=OM$, où $M$ est le milieu de $[AB]$,
et par l'angle $\theta$ entre l'horizontale et la droite $(OM)$
(voir la figure de gauche dans (1) ci-dessous). 
Le problème étant symétrique
par rotation, l'angle $\theta$ n'intervient pas. On est donc ramené à
choisir une distance $r$ au hasard dans $[0,R]$, où $R$ est le rayon du 
cercle. La corde est à l'extérieur si $r\geq r_0$ avec $r_0=\frac{1}{2}R$ 
(voir les figures (1) ci-dessous). Donc la probabilité cherchée est $1/2$
(on a pris la probabilité uniforme sur $[0,R]$).

\medskip
\centerline{\includegraphics{../images/img006954-2}}
\medskip


  \item La corde $[AB]$ est déterminée par son centre $M$ : il suffit de
tracer la perpendiculaire à $(OM)$ passant par $M$ (figure de gauche
ci-dessus, comme pour 1). On est donc ramené à choisir un point
$M$ au hasard dans le disque de rayon $R$.
La corde passe à l'intérieur du triangle
si $|OM|\leq \frac{1}{2}R$ (comme au 1), autrement
dit si le point $M$ se trouve dans le disque de rayon $R/2$, hachuré sur
la figure (2) ci-dessous.
L'aire du domaine est $\pi(R/2)^2$ et l'aire
totale est $\pi R^2$, donc la probabilité cherchée est $1/4$
(on a pris la probabilité uniforme sur le disque de rayon $R$).

\medskip
\centerline{\includegraphics{../images/img006954-3}}
\bigskip

  \item Une corde est déterminée par 2 points $A$ et $B$ 
sur le cercle (ses extrémités). Le problème étant symétrique par rotation,
on peut fixer le point $A$ et considérer la position du point $B$.
On est donc ramené à choisir un point $B$ au hasard sur le cercle.
La corde passe à l'intérieur du triangle pour $B$ appartenant à un arc
de cercle faisant 1/3 du cercle (voir la figure (3) ci-dessus).
Donc la probabilité cherchée est $1/3$
(on a pris la probabilité uniforme sur le cercle, c'est-à-dire la longueur
d'un arc divisé par la longueur du cercle entier).

  \item  On peut déterminer la corde par son extrémité $A$ et l'angle $\theta$ 
que fait la corde avec la tangente au cercle en $A$
(voir la figure (4) ci-dessus). On est donc ramené à choisir 
un angle $\theta$ au hasard dans $[0,\pi]$.
On voit que
la corde passe dans le triangle si $\pi/3\leq \theta\leq 2\pi/3$. Ce
qui donne une probabilité de $1/3$ (on a pris la probabilité
uniforme sur $[0,\pi]$).
\end{enumerate}

\fincorrection
\correction{006028}
\begin{enumerate}
\item On obtient, sur l'échantillon, la moyenne $m_{e}=214$, l'écart-type $\sigma _{e}=55.77$.
\item La moyenne sur l'entreprise est estimée par $m_{e}$.
L'écart-type est estimé par: $\widehat{\sigma _{e}}=\sqrt{\frac{100}{99}}55.77\simeq 56.05$.
\item On en déduit, au seuil 95\%, un intervalle de confiance pour la
moyenne :
$[m_{e} - y_{\alpha }\frac{\widehat{\sigma _{e}}}{\sqrt{n}};
m_{e} + y_{\alpha }\frac{\widehat{\sigma _{e}}}{\sqrt{n}}]=[203.01;224.99]$.
Ainsi le taux moyen de cholestérol est, à un seuil de confiance $95$\%, 
située entre $203$ et $225$ cg.
\end{enumerate}
\fincorrection
\correction{006029}
Il s'agit ici d'estimer une proportion, suite à une observation qui vaut:
$f=\frac{13}{12000}\simeq 1.0833\times 10^{-3}$.

On peut utiliser une approximation par une loi normale pour 
la moyenne d'échantillon. On en déduit un intervalle de confiance pour la
proportion, au seuil 95\%:
$I_{\alpha }=[f-y_{\alpha }\sqrt{\frac{f(1-f)}{n-1}};
p+y_{\alpha }\sqrt{\frac{f(1-f)}{n-1}}]\simeq \lbrack 4.7\times 10^{-4},1.7\times 10^{-3}].$

On peut choisir $I_{\alpha }$ comme intervalle de confiance, au seuil 95\%,
de la proportion cherchée.
Par l'inégalité de Bienaymé-Tchebychev, on a l'intervalle $I=[f-a,f+a],$ avec:
$P[\left| \overline{X}-p\right| \leq a]\geq 1-(\frac{\text{Var}\, \overline{X}}{a^{2}})$ 
et $P[\left| X-p\right| \leq a]\geq 0.95$ si
$1-\frac{\text{Var}\, \overline{X}}{a^{2}} \geq 0.95,$ soit $a \geq 1.3979\times 10^{-3}$. 
On préfèrera donc la première méthode.
\fincorrection
\correction{006025}
Un intervalle dans lequel on soit <<sûr>> à $95$\%
 de trouver le nombre exact de personnes à soigner sur les $10\,000$:
$[p-y_{\alpha }\sqrt{\frac{p(1-p)}{n}};fe+y_{\alpha }\sqrt{\frac{p(1-p)}{n}}]$.
Fréquence entre 65,7\% et 94,3\%.
Donc entre 698 et 802 personnes sur 10000
\fincorrection
\correction{006026}
La loi exacte suivie par $X$ est une loi binomiale de paramètres : $n, p$.
$E(X)=0.75n$ et $\text{Var}\, X=0.25 \cdot0.75n$.
Comme $n>150$, on peut faire l'approximation par la loi normale 
d'espérance $0,75n$ et d'écart-type $\sigma =\sqrt{0.25\cdot0.75n}$.
$P[X>150]\leq 0.05$ si $P[X\leq 150]\geq 0.95$ si:
$P[\frac{X-0.75n}{\sqrt{0.25 \cdot 0.75n}}\leq \frac{150-0.75n}{\sqrt{0.25 \cdot 0.75n}}]\geq 0.95$.
Dans la table de Gauss, on lit $F(1.645)=0.95$.
On n'a plus qu'à résoudre l'inéquation: 
$\frac{150.5-0.75n}{\sqrt{0.25 \cdot 0.75n}}\geq 1.645$, dont les solutions sont:
\begin{equation*}
0\leq n\leq 187.
\end{equation*}

Ainsi, en vendant moins de $187$ billets, la compagnie ne prend qu'un risque
inférieur à $5$\% de devoir indemniser des voyageurs en surnombre.
Faisons varier les paramètres, cela ne pose aucun problème :

$N=150$, $p=0.5$. $n$ est solution de l'inéquation: $\frac{150.5-0.5n}{\sqrt{0.5.0.5n}}\geq 1.645$.
Solution : $n\leq 272$.

$N=300$, $p=0.75$.
$n$ est solution de l'inéquation: $\frac{300.5-0.75n}{\sqrt{0.25.0.75n}}\geq 1.645$.
Solution : $n\leq 381$.

$N=300$, $p=0.5$.
$n$ est solution de l'inéquation: $\frac{300.5-0.5n}{\sqrt{0.5.0.5n}}\geq 1.645$. 
Solution : $n\leq 561$.
\fincorrection
\correction{006027}
\begin{enumerate}
\item La loi de $X$ est la loi binomiale $n=30$, $p=0.2$.

\item Un intervalle de confiance au seuil 95\%, permettant d'estimer le nombre
de clients à prévoir : c'est pour la fréquence: 0.657; 0.943.
Soit entre 20 et 28 personnes.
C'est une large fouchette due à $n$ petit.
\end{enumerate}
\fincorrection
\correction{006030}
\begin{enumerate}
\item On peut estimer $m$ par la moyenne de l'échantillon: $68$ kg, et $\sigma$ par
$\sigma _{e}\sqrt{\frac{300}{299}}=7\sqrt{\frac{300}{299}} \simeq 7.0117$ kg. 
On en déduit un intervalle de confiance pour la moyenne $m$ : $I_\alpha = [67.2;68.8]$.
\item La borne supérieure de l'intervalle étant de $69$ kg, il est
raisonnable de prendre $70$ kg comme espérance de la variable poids d'un passager.
\item Le décollage est autorisé si le poids total des voyageurs et de
leurs bagages ne dépasse pas $26.2$ tonnes.
Pour chacun des $300$ passagers, notons: $X_{i}$ son poids et $Y_{i}$ le poids de ses bagages.
Faisons l'hypothèse d'indépendance entre les variables $X_{i}$ et $Y_{i}$.
Le poids total $Z=\sum_{i=1}^{300}(X_{i}+Y_{i})$ est la somme de $600$
variables aléatoires indépendantes; le théorème central
limite s'applique sous cette hypothèse.
Comme l'espérance totale est $E(Z)=300 \cdot (70+15)=25\,500$ et la variance
de $Z$ est :
$\text{Var}\, Z=300 \cdot(\text{Var}\,X_{i}+\text{Var}\,Y_{i})$.
 Alors $Z$ suit approximativement une loi
normale de moyenne $m=25\,500$, d'écart-type $\sigma =\sqrt{300\cdot(8^{2}+5^{2})}=163.4$.
Alors $Z' =\frac{Z-m}{\sigma}$ suit approximativement une loi normale
centrée réduite.
Le décollage est interdit si : $Z>26\,200$, c'est-à-dire si  $Z' >4.284$.
On lit dans la table de Gauss: pour $t=4$, $F(t)=0.999\,968=P[Z' \leq 4]$.
Le décollage est interdit pour cause de surcharge pondérale avec une
probabilité inférieure à $0.000\,04$.
\end{enumerate}
\fincorrection
\correction{006031}
\begin{enumerate}
\item L'intervalle de temps de $4$ minutes est la répétition de $240$
secondes, au cours desquelles les appels surviennent de façon indépendante, avec la probabilité d'appel de $\frac{1}{20}$; la loi de
probabilité du nombre d'appels reçus en $4$ minutes est donc une loi
binomiale, de paramètres $n=240$ et $p=\frac{1}{20}$.
\item Comme $n\geq 30$ et $np\leq 15,$ il est possible d'approcher cette loi
par une loi de Poisson de paramètre $\lambda $ estimé par $np=12$.
\item Un échantillon de taille $200$ a été réalisé pour
estimer le nombre moyen d'appels par minute; c'est un échantillon de
taille $50$ pour la variable précédente (nombre d'appels reçus
en 4 minutes) qui suit une loi de Poisson d'espérance et de variance $12$. 
Un intervalle de confiance au niveau 95\% pour la moyenne est $I_{\alpha}=[11;13]$.
\end{enumerate}
\fincorrection
\correction{006032}
Posons $H_{0}$ <<les rejets chimiques ne modifient pas le
nombre de plages atteintes par les algues>>.

Notons $p_{0}=0.1$ la proportion théorique de plages atteintes par
l'algue verte avant les rejets chimiques; $p$ la proportion théorique de
plages atteintes par l'algue verte après les rejets chimiques et $f$ la
fréquence observée dans l'échantillon.

Considérons alors la variable aléatoire $X_{i}$, $i\leq 50,$ qui a
deux modalités: $1$ si la plage est atteinte, $0$ sinon. C'est une
variable de Bernoulli, alors le nombre total de plages atteintes 
dans l'échantillon est une variable aléatoire qui, sous $H_{0}$, obéit à
une loi binomiale de paramètres $n=50$, $p_{0}=0.1$.

Sous $H_{0}$, <<$p=p_{0}=0.1$>>  la variable 
<<moyenne d'échantillon>> : 
\begin{equation*}
\overline{X}=\frac{\sum_{i=1}^{i=50}X_{i}}{n}
\end{equation*}
dont une réalisation est la fréquence observée, soit $\frac{10}{50}$, 
obéit à une loi que l'on peut approcher par une loi normale
de paramètres : moyenne $p_{0}$ et écart-type $\sqrt{\frac{p_{0}(1-p_{0})}{50}}$.

A l'aide de la formule de cours, on détermine l'intervalle de confiance
associé:
$I\simeq [0.017;0.183]$. On constate que la fréquence observée est dans la zone de rejet
(non chimique) : $0.2$ n'est pas dans
l'intervalle de confiance au seuil $95$\%.
On peut donc rejeter $H_{0}$ et conclure, au risque $0.05$, que les rejets
chimiques modifient de façon significative le nombre de plages atteintes par l'algue.
\fincorrection
\correction{006033}
Mise en oeuvre du test:
\begin{enumerate}
\item On définit un risque: 5\%.
 Pour étudier la dépendance de ces caractères
 faisons l'hypothèse $H_{0}$ : <<les deux caractères sont indépendants 
>> et voyons ce qui se passerait sous cette hypothèse.
Notons les événements:
\begin{itemize}
\item $C$ : <<avoir un cancer dans la population observée>>
\item $F$: <<être fumeur dans la population observée>>
\end{itemize}
Si les événements $F$ et $C$ sont indépendants, alors: $P(F\cap
C)=P(F)\cdot P(C)$ et de même pour les trois autres possibilités: 
$P(\overline{C}\cap F),$ $P(\overline{C}\cap \overline{F}),P(C\cap \overline{F})$, 
quantités que l'on peut donc calculer sous $H_{0}$:

$P(F)=\frac{600}{1000}$, $P(C)=\frac{500}{1000}$, $P(F) \cdot P(C)=\frac{3}{10}$,
alors l'effectif théorique correspondant à la catégorie 
<<fumeur et cancéreux>> est de $300$.
\item On en déduit le tableau théorique sous $H_{0}$ :

\begin{center}
\begin{tabular}{cccc}
\emph{Théorique} & cancer & non cancer & marge \\ 
\text{fumeur} & 300 & 300 & 600 \\ 
\text{non fumeur} & 200 & 200 & 400 \\ 
\text{marge} & 500 & 500 & 1000
\end{tabular}
\end{center}

\item On calcule alors la valeur de 
$s=\sum_{i=1}^{i=4} \frac{(O_i-T_i)^{2}}{T_i}$ : on obtient : $s=34.73$.
On a précisé le risque de \%, mais pour $\alpha =0,001$, on lit
dans la table du khi-deux à un degré de liberté :
$P[\chi ^{2}\geq 10.83]=0.001$ et le $\chi^{2}$ calculé est $34.73$!
\item On décide de rejeter $H_{0}$.
Ainsi, en rejetant l'hypothèse de l'indépendance des caractères
<<être fumeur>> et <<avoir un cancer de la gorge>>, on a moins de
une chance sur $1000$ de se tromper, puisque moins de un tableau possible
sur mille conduit à un calcul de $\chi ^{2}$ plus grand que $10.83$ ;
beaucoup moins sans doute, conduiraient à un calcul de $\chi ^{2}$ plus
grand que $34.73$.
\end{enumerate}

\fincorrection
\correction{006929}
\begin{enumerate}
\item
La meilleure estimation de $G$ est la valeur moyenne mesurée, $g_1 =1,364$. 
Pour donner un intervalle de confiance, on fait l'hypothèse que les mesures 
suivent une loi normale d'espérance $G$ et d'écart-type $\sigma$. 
La probabilité que l'intervalle $[g_1 -1,645\sigma_{pop},g_1 +1,645\sigma_{pop}]$ 
ne contienne pas $G_v$ est inférieure à $0,1$. On conclut que $[1,357;1,371]$ 
est un intervalle de confiance relatif à $G$ au seuil de 90\%.
\item
La meilleure estimation de $G$ est la valeur moyenne mesurée, $\bar{g}=\frac{1}{10}(g_1 +\cdots+g_{5})=1,365$. Soit 
$\bar\sigma^2=\frac{1}{4}\sum_{i=1}^5(g_i-\bar g)^2\simeq 2,25 .10^{-5}$.
La meilleure estimation de $\sigma$ est
$\bar\sigma\simeq 4,7 .10^{-3}$.

Comme $\frac{1}{5}(X_1 +\cdots+X_{5})$ suit une loi normale d'espérance $G$ 
et de variance $\frac{1}{25}(\text{Var}(X_1 )+\cdots+\text{Var}(X_{5}))=\frac{1}{5}\sigma^{2}$, 
la probabilité que l'intervalle 
$[\bar{g} -1,645\frac{\sigma}{\sqrt{5}},\bar{g} +1,645\frac{\sigma}{\sqrt{5}}]$ 
ne contienne par $G$ est inférieure à $0,1$. 
Si on estime $\sigma$ par $\bar\sigma$, on trouve que $[1,363;1,367]$ est un 
intervalle de confiance relatif à $G$ au seuil de 90\%.
\end{enumerate}
\fincorrection
\correction{006930}
On note $p$ la proportion inconnue. Soit $X_i$ la variable qui vaut 1 si le 
$i$-ème électeur interrogé déclare avoir l'intention de voter pour A, 
0 sinon. Les $X_i$ sont indépendantes et suivent une loi de Bernoulli 
$\mathcal{B}(p)$, donc le nombre $Z=\sum_{i=1}^{1000}$ d'électeurs 
favorables à A dans un échantillon de 1000 électeurs suit une loi 
binomiale $\mathcal{B}(p,1000)$. Comme $p$ semble de l'ordre de 0,5, 
on peut approcher $\mathcal{B}(p,1000)$ par $\mathcal{N}(1000p,\sqrt{1000p(1-p)})$. 
De nouveau, comme $n=1000$ est grand, on peut estimer $p$ par la fréquence $f=0,521$
 observée dans l'échantillon, et supposer que l'écart-type de la variable 
fréquence d'échantillon $F=Z/1000$ vaut $\sqrt{p(1-p)}/\sqrt{1000}\simeq 0,0158$. 
Comme la variable $\displaystyle \frac{F-p}{\sigma(F)}$ est normale standard, 
la probabilité que l'intervalle $[f-1,96\times 0,0158,f+1,96\times 0,0158]$ 
ne contienne pas $p$ est inférieure à 0,05. On conclut que 
$[0,49, 0,55]$ (autrement dit $52\pm 3$\%) est intervalle de confiance relatif à $p$ au seuil de 95\%.
\fincorrection
\correction{006931}
On suppose que la variable $X=$ ``durée de vie d'un appareil'' suit 
une loi normale d'espérance $m_{pop}$ et d'écart-type $\sigma_{pop}=100$h. 
Si on fait $n$ essais indépendants, alors $T=\sqrt{n}\frac{\bar{X}-m_{pop}}{\sigma_{pop}}$ 
suit une loi normale standard, donc $p(|T|>1,96)<0,05$. Au niveau de confiance 95\%,
on peut affirmer que l'intervalle 
$[\bar{x}-1,96\frac{100}{\sqrt{n}},\bar{x}+1,96\frac{100}{\sqrt{n}}]$ 
contient la durée de vie moyenne cherchée $m_{pop}$. 

La marge d'erreur n'excède pas 50h dès que $1,96\frac{100}{\sqrt{n}}<50$, i.e. si $n\geq 16$.

La marge d'erreur n'excède pas 20h dès que $1,96\frac{100}{\sqrt{n}}<20$, i.e. si $n\geq 97$.
\fincorrection
\nocorrection
\correction{006933}
Loi estimée $\mathcal{E}(\lambda)$ d'espérance
$m=1/\lambda$ et d'écart-type $\sigma=1/\lambda$. 
Estimateur de $m$ : $\bar X_n=10$. Donc on approxime $\sigma$ par $10$.
TCL : $P(\frac{|S_n-n m|}{\sigma\sqrt{n}}<c)\simeq P(|\mathcal{N}(0,1)|<c)$.
Pour avoir un intervalle de confiance à 95\%, on prend $c=1,96$.
L'encadrement $-c<\frac{\bar X_n-1/\lambda}{\sigma/\sqrt{n}}<c$ donne
$$
\frac{1}{\bar X_n+c\sigma/\sqrt{n}}<\lambda<\frac{1}{\bar X_n-c\sigma/\sqrt{n}}
$$
d'où l'intervalle de confiance $[0,091,0,111]$.
\fincorrection
\correction{006934}
On choisit de comparer les fréquences observées sur les deux échantillons.

On fait l'hypothèse $H_0$ : les deux échantillons sont des tirages de 
variables de Bernoulli indépendantes ayant la même espérance $p_{pop}$. 

Sous l'hypothèse $H_{0}$, l'effectif $n_{i}F_{i}$ suit une loi binomiale 
$\mathcal{B}(p_{pop},n_{i})$. Les échantillons étant de grande taille 
($n_1$, $n_2 \geq 30$), on peut supposer que $F_{i}$ suit une loi 
normale $\mathcal{N}(p_{pop},\sqrt{p_{pop}(1-p_{pop})/n_{i}})$. De plus, 
pour le calcul de la variance, on peut estimer $p_{pop}$ au moyen de 
l'ensemble des données disponibles, i.e.
\begin{eqnarray*}
p_{pop}\sim \frac{300+125}{700+225}=0,46.
\end{eqnarray*}
La différence $D=F_{1}-F_{2}$ suit une loi normale d'espérance nulle et de variance
$
\text{Var}(D)=\text{Var}(F_{1})+\text{Var}(F_{2})=p_{pop}(1-p_{pop})
(\frac{1}{n_{1}}+\frac{1}{n_{2}})\sim 0,2484(\frac{1}{700}+\frac{1}{225})=0,00146.
$

Par conséquent, on choisit comme fonction discriminante $T=D/\sqrt{0.00146}=D/0,0382$, 
qui suit une loi normale standard.
La table donne $p(|T|>1,96)=0,05$.
Or, sur l'échantillon mesuré, la variable $T$ prend la valeur 
$t=\frac{-0,127}{0,0382}=-3,324$. Comme $|t|>1,96$, on rejette 
$H_{0}$. On conclut, avec au plus 5\% de chances de se tromper, 
que les deux lots de boulons ne sont pas tirés de la même population.
\fincorrection
\correction{006935}
Notons $X$ le résultat d'un tirage  d'un entier entre $0$
et $9$ à l'aide de ce générateur et $p=(p_0,\ldots,p_9)$ sa loi. On cherche à tester
l'hypothèse  $H_0$ : ``$p$ est la loi $\mathcal{U}_{\{0,\ldots,9\}}$''
contre l'hypothèse $H_1$ : ``$p$ n'est pas la loi
$\mathcal{U}_{\{0,\ldots,9\}}$'' au niveau $5\%$.
Notons $N^{(i)}$ le nombre d'apparitions du chiffre $i$ sur $1000$ 
tirages de chiffres à l'aide de ce générateur et
$Z=\sum_{i=0}^{9}\frac{(N^{(i)}-100)^{2}}{100}$.\\ On rejette
l'hypothèse $H_0$ au niveau $5\%$ si la valeur observée $z_{obs}$
de $Z$ est telle que $P_{H_{0}}(Z\geq z_{obs})\leq 0,05$. Comme pour tout $i\in\{0,\ldots,9\}$,  $1000p_i=100$ est
suffisamment grand,  on peut approximer la fonction de répartition de la  loi de
$Z$ sous $H_0$ par celle du $\chi^{2}$ à $9$ degrés de
liberté. Donc, on peut approximer $P_{H_{0}}(Z\geq z_{obs})$ par
$1-F_{\chi^{2}_{9}}(z_{obs})$. A partir des valeurs observées pour les variables aléatoires
$N^{(i)}$ qui sont données dans le tableau, on obtient
$z_{obs}\simeq 21,86$. D'après la table de valeurs numériques de
la loi du $\chi^2$, %\versun{donnée à l'annexe \ref{tablechi}},
$F_{\chi^{2}_{9}}(z_{obs})$ est compris entre $0,99$ et
$0,995$. Donc, au niveau $5\%$, on rejette l'hypothèse $H_0$. On
rejette encore l'hypothèse $H_0$ au niveau $1\%$ mais pas au
niveau $0,5\%$.  
\fincorrection
\nocorrection
\nocorrection
\nocorrection
\correction{001435}
Soit $G$ sous-groupe de $\Zz/8\Zz$, alors $\mathrm{Card} G$ divise $\mathrm{Card}
\Zz/8\Zz = 8$. Donc $\mathrm{Card} G \in \{1,2,4,8\}$. De plus si $G$
contient la classe $\bar{n}$ d'un nombre impair, alors $G$
contient le sous-groupe engendr\'e par $\bar{n}$ qui est
$\Zz/8\Zz$ car alors $n$ et $8$ sont premiers entre eux, donc $G
=\Zz/8\Zz$.

\'Etude des cas. Si $\mathrm{Card} G=8$ alors $G =\Zz/8\Zz$. Si $\mathrm{Card} G =
4$ alors $G$ ne peut contenir que des classes d'entiers pairs
d'après la remarque précédente, mais comme il y a exactement $4$
classes d'entiers pairs alors $G =\{
\bar{0},\bar{2},\bar{4},\bar{6}\}$. Si  $\mathrm{Card} G = 2$ alors $G=\{
\bar{0},x\}$ et $x$ est un élément d'ordre $2$, le seul élément
d'ordre $2$ de $\Zz/8\Zz$ est $\bar{4}$. Donc $G=\{
\bar{0},\bar{4}\}$. Enfin si $\mathrm{Card} G = 1$ alors $G = \{ \bar{0}
\}$.

\fincorrection
\nocorrection
\nocorrection
\correction{001438}
La relation d'équivalence associée au quotient $\Rr^*/\Rr_+^*$ est
:
$$x\sim y \Leftrightarrow xy^{-1} >0.$$
Si $x>0$ alors $x\sim +1$ car $x(1)^{-1} >0$ (en fait $x$ est
équivalent à n'importe quel réel strictement positif) ; si $x<0$
alors $x\sim -1$ car $x(-1)^{-1} >0$, enfin $-1$ et $+1$ ne sont
pas équivalents. Il y a donc deux classes d'équivalence :
$\Rr^*/\Rr_+^* = \{ \overline{+1} , \overline{-1} \}$.

L'application $\phi : \Rr^*/\Rr_+^* \longrightarrow \Zz/2\Zz$
définie par $\phi(\overline{+1})=\tilde{0}$ et
$\phi(\overline{-1})=\tilde{1}$ est un isomorphisme entre les deux
groupes.
\fincorrection
\nocorrection
\nocorrection
\nocorrection
\correction{001442}
\begin{enumerate}
    \item Il faut montrer que pour $x\in G$ et $y\in D(G)$, $xyx^{-1} \in D(G)$.
Commen\c{c}ons par montrer ceci pour $y$ un g\'en\'erateur de
$D(G)$. Si $y=ghg^{-1}h^{-1}$ avec $g,h\in G$. Nous remarquons que
:
$$xyx^{-1} = \left( x gh x^{-1}(gh)^{-1} \right) \left( gh g^{-1}h^{-1} \right)
\left( hg x (hg)^{-1}x^{-1} \right) $$ qui est un produit
d'\'el\'ements de $D(G)$. Donc $xyx^{-1}$ est un \'el\'ement de
$D(G)$.

Soit maintenant $y$ un \'el\'ement quelconque de $D(G)$, alors il
s'\'ecrit comme produit de g\'en\'erateurs :
$$y = y_1y_2\ldots y_n,\quad \text{ avec } y_i = g_ih_ig_i^{-1}h_i^{-1}.$$
\'Ecrivons $xyx^{-1} =
(xy_1x^{-1})(xy_2x^{-1})\ldots(xy_nx^{-1})$. Chaque $xy_ix^{-1}$
appartient \`a $D(G)$. Et donc $xyx^{-1}$. Donc $D(G)$ est un
sous-groupe distingu\'e de $G$.

    \item Soit $\alpha,\beta \in G/D(G)$, alors il existe
$a,b \in G$ tels que $\overline{a}=\alpha$ et
$\overline{b}=\beta$. Nous savons que $aba^{-1}b^{-1}\in D(G)$ et
donc $\overline{aba^{-1}b^{-1}} = \epsilon$
 o\`u $\epsilon$ est l'\'el\'ement neutre de $G/D(G)$. Mais
$$\overline{aba^{-1}b^{-1}} = \overline{a}\overline{b}\overline{a^{-1}}\overline{b^{-1}}
= \overline{a}\overline{b}\overline{a}^{-1}\overline{b}^{-1} =
\alpha\beta\alpha^{-1}\beta^{-1}.$$ Donc
$\alpha\beta\alpha^{-1}\beta^{-1}=\epsilon$, autrement dit
$\alpha\beta= \beta\alpha$. Et ceci quelque soit $\alpha$ et
$\beta$, donc $G/D(G)$ est commutatif.
\bigskip
G\'en\'eralisation : si $H$ est un sous-groupe distingu\'e.
\begin{itemize}
    \item[$\bullet$] Si $D(G) \subset H$ alors $G/D(G)$ est un sous-groupe de $G/H$ donc
$G/H$ est commutatif car $G/D(G)$ l'est.
    \item[$\bullet$] Si $G/H$ est commutatif alors pour $g,h \in G$
la classe de $ghg^{-1}h{-1}$ dans $G/H$ v\'erifie :
$$\overline{ ghg^{-1}h{-1}} = \overline{g}\overline{h}\overline{g^{-1}}\overline{h^{-1}}
= \overline{g}\overline{ g^{-1}}\overline{h}\overline{h^{-1}} =
\epsilon.$$ Mais les \'el\'ements dont la classe dans $G/H$ est
l'\'el\'ement neutre sont exactement les \'el\'ements de $H$. Donc
$ghg^{-1}h^{-1}$ appartient \`a $H$. Ainsi  tous les
g\'en\'erateurs de $D(G)$ sont dans $H$ et donc $D(G) \subset H$.
\end{itemize}
\end{enumerate}
\fincorrection
\nocorrection
\nocorrection
\nocorrection
\correction{001446}
Notons $C = AB = \frac{1}{\sqrt{2}}
\begin{pmatrix}
1 & 1 \\
-1 & 1
\end{pmatrix}$.
\begin{enumerate}
    \item Un calcul donne $C^8 = I$ et pour $1\le k \le 7$, $C^k\not= I$.
Donc le groupe $H$ engendr\'e par $C$ est d'ordre $8$. Attention !
m\^eme si $A^2=I$ et $B^2=I$ on a $(AB)^2 \not= I$ car $AB \not=
BA$.
    \item Pour montrer que $H$ est distingu\'e il suffit de montrer que
$ACA^{-1}$ et $BCB^{-1}$ sont dans $H$. Mais $ACA^{-1} = ACA =
AABA = BA = (AB)^{-1} \in H$. De m\^eme $BCB^{-1} = (AB)^{-1}$.
Donc $H$ est distingu\'e dans $H$.

\bigskip
Un \'el\'ement $M$ de $G$ s'\'ecrit
$$M= A^{a_1}B^{b_1}A^{a_2}\ldots A^{a_n}B^{b_n} \quad a_i,b_i \in \Zz.$$
Mais dans $G/H$ tout terme $\overline{AB}$ ou $\overline{BA}$ vaut
$\overline{I}$ Donc $G/H = \{ \overline{I},
\overline{A},\overline{A}^2, \overline{A}^3,\ldots,
\overline{B},\overline{B}^2,\overline{B}^3,\ldots \}$ mais comme
$A^2=B^2=I$ et $AB \in H$ alors $G/H$ s'\'ecrit simplement :
$$G/H = \left\lbrace \overline{I},\overline{A} \right\rbrace.$$

\bigskip
Enfin, par la formule $|G| = |H| \times |G/H|$ nous obtenons $|G|
= 8 \times 2 = 16$.

\end{enumerate}
\fincorrection
\correction{001447}
\begin{enumerate}
  \item
  \begin{enumerate}
    \item $f((x,y)+(x',y'))=f(x+x',y+y')=3(x+x')+6(y+y')=3x+6y+3x'+6y'=f(x,y)+f(x',y')$.
    \item $\mathrm{Ker} f = \{ (x,y) ; f(x,y)=0\} = \{ (x,y) ; 3x+6y=0\}=
\{ (x,y) ; x=-2y \}= \{ (-2k,k) ; k\in \Zz\}$. Si $\mathrm{Ker} f =
p\Zz\times q\Zz$ alors $f(p,0)=0$ donc $3p=0$ soit $p=0$. De même
$f(0,q)=0$ implique $q=0$ et alors $\mathrm{Ker} f = \{ (0,0) \}$, ceci
contredit le fait que $f(-2,1)=0$.
    \item On a $f(\Zz^2) = 3\Zz$, le morphisme
$f : \Zz^2 \longrightarrow 3\Zz$ définit par passage au quotient
par le noyau un morphisme injectif $\bar{f} : \Zz^2/\mathrm{Ker} f
\longrightarrow 3\Zz$ (c'est le théorème de factorisation). De
plus comme $f$ est surjectif alors $\bar{f}$ l'est aussi. Ainsi
$\bar{f}$ est un isomorphisme entre $\Zz^2/\mathrm{Ker} f
=\Zz^2/(-2,1)\Zz$ et $3\Zz$.
  \end{enumerate}
  \item Définissons $g : \Zz^2 \longrightarrow \Zz/2\Zz \times \Zz/2\Zz$
par $g(x,y)=(\bar{x},\bar{y})$ o\`u $\bar{n}$ désigne la classe de
$n$ dans $\Zz/2\Zz$. Le noyau de $g$ est $2\Zz \times 2\Zz =
\langle (2,0) ;(0,2) \rangle = G$. Le passage au quotient par le
noyau définit l'isomorphisme $\bar{g}$ cherché.
\end{enumerate}
\fincorrection
\nocorrection
\nocorrection
\nocorrection
\correction{002102}
(I) (a) $E\not= \emptyset$ car $X\in E$. L'ensemble $A_0=\bigcap_{A\in E} A$ est
de mani\`ere \'evidente le plus petit \'el\'ement de $E$.\smallskip

(b) On a $\varphi(A_0) \subset A_0$ puisque $A_0\in E$.  On d\'eduit, par la croissance
de $\varphi$, que $\varphi(\varphi(A_0)) \subset \varphi(A_0)$, ce qui donne
$\varphi(A_0)\in E$ et donc $A_0\subset \varphi(A_0)$.\medskip

(II) (a) La croissance de $\varphi$ est imm\'ediate. \smallskip

(b) Consid\'erons la partie $A_0$ associ\'ee \`a $\varphi$. D'apr\`es le (b) du (I), on a
$X\setminus h(X\setminus g(A_0))=A_0$. Autrement dit, les parties $A_0$ et $h(X\setminus
g(A_0))$ constituent une partition de $X$. Consid\'erons l'application $f:X\rightarrow X$
d\'efinie comme \'etant $g$ sur $A_0$ et $h^{-1}$ sur $h(Y\setminus g(A_0))$. On voit sans
difficult\'e que $f$ est une bijection (noter que les images respectives des deux restrictions
pr\'ec\'edentes sont $g(A_0)$ et $Y\setminus g(A_0)$ et qu'elles constituent une partition de
$Y$).\smallskip
\fincorrection
\correction{002103}
Pour tout $x\in X$, posons $C(x)=\{ y\in X \hskip 2pt |\hskip 2pt  x\ \hbox{\rm et}\ y \
\hbox{\rm sont comparables}\}\ $ et consid\'erons $Y=\bigcap_{x\in X}C(x)$. La partie $Y$ est
totalement ordonn\'ee puisque d\`es que $y,y^\prime \in Y$, alors $y^\prime \in C(y)$ et donc
$y$ et
$y^\prime$ sont comparables. De plus, pour tout $x\notin Y$, il existe $y\in X$ tel que 
$x\notin C(y)$, c'est-\`a-dire, $y$ et $x$ non comparables.\smallskip

Il n'y a pas unicit\'e de l'ensemble $Y$ en g\'en\'eral. En effet, dans un ensemble
ordonn\'e o\`u il existe un \'el\'ement $y$ qui n'est comparable qu'\`a lui-m\^eme,
on peut prendre $Y=C(y)=\{y\}$. Il est facile de construire des ensembles ordonn\'es
poss\'edant plusieurs tels \'el\'ements $y$ (penser \`a la relation d'\'egalit\'e, dont     
le graphe est la diagonale). 
\fincorrection
\nocorrection
\nocorrection
\nocorrection
\correction{002107} 
Pour la derni\`ere question, v\'erifier par r\'ecurrence que $\displaystyle x^{\star\hskip 2pt n}=
\sum_{k=1}^n (-1)^{k-1} C_n^k x^k$.
\fincorrection
\correction{002108}
(a) D\'esignant par $b$ l'inverse \`a gauche de $a$ et par $c$ l'inverse \`a gauche de $b$,
on a $ab= (cb)(ab)= c(ba)b=cb=e$. L'\'el\'ement $b$ est donc l'inverse de $a$.
\smallskip

(b) d\'ecoule imm\'ediatement de (a).
\smallskip
\fincorrection
\correction{002109}
(a) Pour $x,y\in E$ quelconques, notons $x^\prime$ et $y^\prime$ leurs inverses \`a gauche
respectifs. Si $xy=e$, on a aussi $yx=(x^\prime x)yx=x^\prime (xy) x =x^\prime x = e$.
\smallskip

(b) Soit $f$ un \'el\'ement neutre \`a gauche. On a donc $fe=e$. D'apr\`es (a), on a aussi
$ef=e$, c'est-\`a-dire $f=e$.
\smallskip

(c) Pour tout $x\in E$, on a $xe=x(x^\prime x)=(x x^\prime) x = x$ puisque d'apr\`es (a), $x
x^\prime = e$.
\smallskip

(d) r\'esulte alors de (a), (b) et (c).
\smallskip 
\fincorrection
\nocorrection
\nocorrection
\nocorrection
\nocorrection
\correction{002114}
Pour tous $x,y\in G$, on a $xyx^{-1}y^{-1} = xyxy = (xy)(xy) = 1$ c'est-\`a-dire $xy=yx$. Donc $G$ est 
ab\'elien. Si $G$ est fini, il peut \^etre consid\'er\'e comme espace vectoriel sur le corps $\Zz/2\Zz$,             et est alors n\'ecessairement de dimension finie, ce qui donne $G$ isomorphe comme espace vectoriel \`a $(\Zz/2\Zz)^n$ et donc $|G| = 2^n$.
\fincorrection
\correction{002115}
En groupant chaque \'el\'ement $x\in G$ avec son inverse $x^{-1}$, on obtient une
partition de $G$ en sous-ensembles $\{y,y^{-1}\}$ qui ont deux \'el\'ements sauf
si $y=y^{-1}$, c'est-\`a-dire si $y^2=e$. L'\'el\'ement neutre $e$ est un tel
\'el\'ement $y$. Ce ne peut pas \^etre le seul, sinon $G$ serait d'ordre impair.
\fincorrection
\nocorrection
\nocorrection
\correction{002118}
Pour tout $h\in H$, on a $ha=k_hb$ pour un certain $k_h\in K$. En \'ecrivant 
$ha=h(ea)=hk_eb$, on obtient $k_h=hk_e$, ce qui donne $h=k_h (k_e)^{-1} \in K$.
\fincorrection
\nocorrection
\correction{002120}
(a) Supposons que $H \cup K$ soit un sous-groupe de $G$ et que $H$ ne soit pas inclus dans
$K$, c'est-\`a-dire, qu'il existe $h\in H$ tel que $h\notin K$. Montrons que $K\subset H$.
Soit $k\in K$ quelconque. On a $hk \in H\cup K$. Mais $hk\notin K$ car sinon $h=(hk)k^{-1} 
\in K$. D'o\`u $hk\in H$ et donc $k=h^{-1}(hk)\in H$.
\smallskip

(b) d\'ecoule imm\'ediatement de (a).
\fincorrection
\correction{002121}
Soit $H$ une partie finie non vide de $G$ stable par la loi de composition. Pour montrer que
$H$ est un sous-groupe, il reste \`a voir que pour tout $x\in H$, $x^{-1}\in H$. Les
puissances $x^k$ o\`u $k\in \N$ restant dans $H$, il existe $m,n\in \N$ tels que $m>n$ et
$x^m=x^n$. On a alors $x^{m-n-1} \cdot x=1$, soit $x^{-1}= x^{m-n-1}$, ce qui montre que
$x^{-1} \in H$.\smallskip

Si $H$ est infini, la propri\'et\'e pr\'ec\'edente n'est pas vraie en g\'en\'eral. Par
exemple $\N$ est une partie stable de $\Z$ pour l'addition mais n'en est pas un sous-groupe.
\fincorrection
\nocorrection
\nocorrection
\correction{002124}
Soient $a,b\in G$ d'ordre respectifs $m$ et $n$. Posons $\mu = \mathrm{ppcm}(m,n)$. 
On a $(ab)^\mu=a^\mu \cdot b^\mu = e\cdot e = e$ ($a^\mu = b^\mu = e$ r\'esultant du fait
que $m$ et $n$ divisent $\mu$). L'ordre de $ab$ divise donc $\mu$.\smallskip

Supposons que $\mathrm{pgcd}(m,n)=1$. Soit $k\in \Z$ tel que $(ab)^k=1$, soit $a^k=b^{-k}$. On en
d\'eduit que $a^{nk} = e$ et $b^{mk}=e$. D'o\`u $m|nk$ et $n|mk$. L'hypoth\`ese 
$\mathrm{pgcd}(m,n)=1$ donne alors $m|k$ et $n|k$ et donc $\mathrm{ppcm}(m,n)|k$. Cela combin\'e \`a la
premi\`ere partie montre que $ab$ est d'ordre $\mathrm{ppcm}(m,n)=mn$.
\fincorrection
\nocorrection
\nocorrection
\correction{002127} 
Etant donn\'e $a\in F$, soit $S$ une partie de $G$ contenant $a$
et engendrant $G$. Si $<S-\{ a \}> \not= G$, alors il existe un sous-groupe propre
maximal $G_i$ tel que $<S-\{ a \}> \subset G_i$. Mais alors $<S> \hskip 3pt \subset
\hskip 3pt <S-\{ a \}> <a> \hskip 3pt \subset \hskip 3pt G_i$. Contradiction, donc
$<S-\{ a \}> \not= G$.

Inversement, supposons que $a\notin F$, c'est-\`a-dire, il existe $i\in I$ tel que
$a\notin G_i$. Alors pour $S=G_i\cup \{a\}$, on a $<S>=G$ (par maximalit\'e de
$G_i$) mais $<S-\{ a \}> = G_i \not= G$.
\fincorrection
\nocorrection
\nocorrection
\correction{002130}
(a) ($\Rightarrow$) Si $HK$ est un groupe, pour tous $h\in H$ et $k\in K$, on a $(hk)^{-1} =k^{-1} h^{-1} \in HK$ et donc $kh\in (HK)^{-1} = K^{-1} H^{-1} = KH$. D'o\`u $HK\subset KH$. L'autre inclusion s'obtient similairement.

($\Leftarrow$) On v\'erifie ais\'ement en utilisant l'hypoth\`ese $HK=KH$ que $(HK)\cdot (HK) \subset HK$ et que $(HK)^{-1}\subset HK$.
\smallskip

(b) Etant donn\'es $h_0, h \in H$ et $k_0,k\in K$, on a $h_0 k_0 = hk$ si et seulement si $h_0^{-1} h = k_0 k^{-1}$. Cet \'el\'ement est n\'ecessairement dans l'intersection $H\cap K$. On a donc $h_0 k_0 = hk$ si et seulement s'il existe $u\in H\cap K$ tel que $h=h_0 u$ et $k=u^{-1} k_0$. Pour chaque \'el\'ement fix\'e $h_0k_0 \in HK$, il y a donc $|H\cap K|$ fa\c cons de l'\'ecrire $hk$ avec $(h,k)\in H\times K$. D'o\`u le r\'esultat.
\fincorrection
\correction{002131}
D'apr\`es le th\'eor\`eme de Lagrange, les sous-groupes de $S_3$ sont d'ordre $1$, $2$, $3$ ou $6$. Les sous-groupes d'ordre $1$ et $6$ sont les sous-groupes triviaux $\{1\}$ et $S_3$ respectivement.
Comme $2$ et $3$ sont premiers, les sous-groupes d'ordre $2$ et $3$ sont cycliques. Un sous-groupe d'ordre $2$ est tout sous-groupe engendr\'e par une transposition: il y en a $3$. Il existe un seul sous-groupe d'ordre $3$, celui engendr\'e par le $3$-cycle $(1 \hskip 2pt 2 \hskip 2pt 3)$.
\fincorrection
\correction{002132}
Les \'el\'ements diff\'erents de $1$ sont d'ordre $5$, $7$ ou $35$. S'il existe un \'el\'ement $g$ d'ordre $35$ (\textit{i.e.}, si le groupe est cyclique d'ordre $35$), alors $g^5$ est d'ordre $7$ et $g^7$ est d'ordre $5$. Supposons que le groupe n'est pas cyclique et qu'il n'existe pas d'\'el\'ement d'ordre $7$. Tout 
\'el\'ement  diff\'erent de $1$ serait alors d'ordre $5$ et le groupe serait r\'eunion de sous-groupes d'ordre $5$. Mais de tels sous-groupes sont soit \'egaux soit d'intersection $\{1\}$ (car $5$ est premier).
On aurait alors $35 = 4n+1$ avec $n$ le nombre de sous-groupes distincts d'ordre $5$, ce qui donne la contradiction cherch\'ee. Le raisonnement est le m\^eme s'il n'existe pas d'\'el\'ement d'ordre $5$.
\fincorrection
\correction{002133}
Si $p=2$ alors $|G|$ est d'ordre $4$: $G$ est le groupe de Klein $(\Zz/2\Zz)^2$ dont tous les \'el\'ements diff\'erents de $1$ sont d'ordre $2$. On peut donc supposer pour la suite que $p$ est impair. En proc\'edant comme dans l'exercice  \ref{ex:deb32}, on montre qu'il existe forc\'ement dans $G$ un \'el\'ement d'ordre $2$. Enfin si tous les \'el\'ements diff\'erents de $1$  \'etaient d'ordre $2$, alors d'apr\`es l'exercice \ref{ex:deb14}, l'ordre de $G$ serait une puissance de $2$. Il existe donc aussi un \'el\'ement d'ordre $p$.
\fincorrection
\correction{002134}
On a $2^{2^n} \equiv -1$ modulo $p$. On en d\'eduit que $2^{2^{n+1}} \equiv 1$ modulo $p$. Ces deux conditions donnent que l'ordre de $2$ dans $(\Zz/p\Zz)^\times$ est $2^{n+1}$. Cet ordre devant diviser l'ordre de $(\Zz/p\Zz)^\times$, c'est-\`a-dire $p-1$, on obtient le r\'esultat souhait\'e.
\fincorrection
\correction{002135}
Comme $2^n \equiv 1$ modulo $2^n-1$, l'ordre de $2$ modulo $2^n-1$, disons $m$, divise $n$. Si $m<n$, on aurait $2^m \equiv 1$ modulo $2^n-1$, c'est-\`a-dire $2^n-1$ divise $2^m-1$, ce qui n'est pas possible. L'ordre de $2$ modulo $2^n-1$ est donc $n$, et celui-ci doit diviser l'ordre de $(\Zz/(2^n-1)\Zz)^\times$, qui vaut $\varphi(2^n-1)$.
\fincorrection
\nocorrection
\nocorrection
\nocorrection
\nocorrection
\nocorrection
\nocorrection
\nocorrection
\nocorrection
\nocorrection
\nocorrection
\nocorrection
\nocorrection
\nocorrection
\nocorrection
\correction{001428}
$HK = \left\lbrace hk \ / h\in H, k\in K \right\rbrace$.
\begin{enumerate}
    \item Soit $\phi : H\times K \rightarrow HK$ d\'efinie par
$\phi(h,k) = hk$. Montrons que $\phi$ est bijective : $\phi$ est
surjective par d\'efinition de $HK$ et si $\phi(h,k)=\phi(h',k')$
alors $hk=h'k'$ et donc ${h'}^{-1}h = k'k^{-1}$ or $H\cap K =\{
e_G \}$ et donc ${h'}^{-1}h = e_G$ et donc $h=h'$, de m\^eme
$k=k'$ et donc $\phi$ est injective.

Comme $\phi$ est bijective $\mathrm{Card} H\times K = \mathrm{Card} HK$ et donc
$\mathrm{Card} HK = \mathrm{Card} H . \mathrm{Card} K$.

    \item Supposons qu'il existe deux sous-groupes $H$ et $K$ distincts et d'ordre
$p$. Montrons d'abord que $H\cap K =\{ e_G \}$. En effet $H\cap K$
est un sous-groupe de $H$ et donc le cardinal de $H\cap K$ divise
$\mathrm{Card} H = p$ avec $p$ premier. Or comme $H\not= K$ alors $H\cap K
\not= H$ et donc $\mathrm{Card} H\cap K =1$, c'est ce que nous voulions
d\'emontrer.

Maintenant d'apr\`es la premi\`ere question $HK$ est un
ensemble de cardinal $p^2$ dans le groupe  $G$ de cardinal $pq
< p^2$. Donc il ne peut exister deux sous-groupe d'ordre $p$.

Supposons maintenant que $H$ soit un sous-groupe d'ordre $p$,
c'est donc l'unique sous-groupe d'ordre $p$ d'apr\`es ce que nous
venons de d\'emontrer. Pour $g\in G$ le sous-groupe $gHg^{-1}$ est
du m\^eme ordre que $H$ (car pour $g$ fix\'e le morphisme
$\theta_g$ de $G$ dans $G$, $\theta_g(h) =ghg^{-1}$ est un
automorphisme et en particulier un biction donc $\mathrm{Card}
\theta_g(H)=\mathrm{Card} H$ ). Par cons\'equent $gHg^{-1} = H$ et donc
$H$ est un sous-groupe distingu\'e.
\end{enumerate}
\fincorrection
\nocorrection
\nocorrection
\nocorrection
\nocorrection
\nocorrection
\nocorrection
\correction{002137}
Soient $x,y\in G$ quelconques. De $(xy)^n=x^n y^n$, on d\'eduit 
$(yx)^{n-1}=x^{n-1} y^{n-1}$ puis $(yx)^{n}=yx^n y^{n-1}$ et donc $y^n x^n=yx^n
y^{n-1}$, ce qui donne $y^{n-1}x^n = x^n y^{n-1}$. Ainsi, pour tout $y\in G$,
$y^{n-1}$ commute \`a tous les \'el\'ements de la forme $x^n$ avec $x\in G$, et est
donc dans le centre de $G$, puisque l'application $x\rightarrow x^n$ est suppos\'ee
surjective. 
\fincorrection
\correction{002138}
Tout automorphisme $\varphi$ du groupe $G=\Z/2\Z \times \Z/2\Z$ permute les trois \'el\'ements
d'ordre $2$, c'est-\`a-dire l'ensemble $G^\ast$ des trois \'el\'ements non triviaux. La
correspondance qui \`a $\varphi \in \hbox{\rm Aut}(\Z/2\Z \times \Z/2\Z)$ associe sa
restriction \`a $G^\ast$ induit un morphisme $\chi: \hbox{\rm Aut}(\Z/2\Z \times
\Z/2\Z) \rightarrow S_3$. Tout morphisme $\varphi \in \hbox{\rm Aut}(\Z/2\Z \times \Z/2\Z)$
\'etant d\'etermin\'e par sa restriction \`a $G^\ast$, ce morphisme $\chi$ est injectif.
De plus, tout automorphisme lin\'eaire (pour la structure de $\Z/2\Z$-espace vectoriel de
$\Z/2\Z \times \Z/2\Z$) est un automorphisme de groupes. Il y a $6$ tels automorphismes
(autant qu'il y a de bases). L'image de $\chi$ contient donc au moins $6$ \'el\'ements.
Comme c'est un sous-groupe de $S_3$, c'est $S_3$ lui-m\^eme et $\chi$ est un isomorphisme.
\fincorrection
\correction{002139} 
Le sous-groupe $H$ est \`a la fois la classe \`a gauche et la classe \`a
droite modulo $H$ de l'\'el\'ement neutre. Si $[G:H]=2$, son compl\'ementaire $H^c$
dans $G$ est donc l'autre classe, \`a droite et \`a gauche. Classes \`a droite et
classes \`a gauche   coincident donc, soit $gH=Hg$ et donc $gHg^{-1}=H g g^{-1}=H$
pour tout $g\in G$.
\fincorrection
\correction{002140}
D'apr\`es l'hypoth\`ese, pour tout $x\in G$, il existe $z\in G$ tel que $xH \cdot
x^{-1}H = zH$. On en d\'eduit $xHx^{-1} \subset zH$. Cela entraine que $1\in zH$ et
donc que $z\in H$. D'o\`u finalement $xHx^{-1} \subset H$.
\fincorrection
\correction{002141} 
Etant donn\'es $y,z\in H$, on a $y\simeq 1$ et $z\simeq 1$. La
compatibilit\'e de la loi donne d'une part $yz \simeq 1$, soit $yz \in H$, et d'autre
part $y y^{-1} \simeq y^{-1}$ soit $y^{-1}\in H$. Cela montre que $H$ est un
sous-groupe de $G$. Pour tout $x\in G$, on a aussi $xyx^{-1} \simeq x 1 x^{-1}=1$ et
donc $xyx^{-1}\in H$. Le sous-groupe $H$ est donc distingu\'e.

De plus, pour $x,x^\prime \in G$, si $x\simeq x^\prime$, alors par compatibilit\'e de
la loi, on a $x^\prime x^{-1} \simeq x x^{-1} = 1$, c'est-\`a-dire $x^\prime x^{-1}\in
H$. R\'eciproquement, si $x^\prime x^{-1}\in H$, alors $x^\prime x^{-1} \simeq 1$,
et donc, par compatibilit\'e de la loi, $x\simeq x^\prime$.
\fincorrection
\correction{002142}
Pour tout $g\in G$, la conjugaison $c_g:G\rightarrow G$ par $g$ induit un automorphisme de
$H$ si $H$ est distingu\'e dans $G$. Si de plus $K$ est caract\'eristique dans $H$, alors
$K$ est stable par $c_g$. D'o\`u $K$ est alors distingu\'e dans $G$.
\smallskip

Le sous-ensemble $V_4$ du groupe sym\'etrique $S_4$ consistant en l'identit\'e et
les trois produits de transpositions disjointes: $(1 \hskip 2pt 2)(3 \hskip 2pt 4)$, $(1
\hskip 2pt 3)(2 \hskip 2pt 4)$ et $(1 \hskip 2pt 4)(2 \hskip 2pt 3)$ est un sous-groupe
(v\'erification imm\'ediate) qui est distingu\'e: cela r\'esulte de la formule $g (i \hskip
2pt j)(k \hskip 2pt l) g^{-1} = (g(i) \hskip 2pt g(j))(g(k) \hskip 2pt g(l))$ pour
$i,j,k,l \in \{1,2,3,4\}$ distincts. Le sous-groupe $K$ (d'ordre $2$) engendr\'e par $(1
\hskip 2pt 2)(3 \hskip 2pt 4)$ est distingu\'e dans $V_4$ (car $V_4$ est ab\'elien). Mais $K$
n'est pas distingu\'e dans $S_4$ (comme le montre encore la formule pr\'ec\'edente).
\fincorrection
\nocorrection
\nocorrection
\correction{002145}
Le groupe $\mu_{mn}$ a un \'el\'ement d'ordre $mn$. En revanche tout \'el\'ement $x\in
\mu_m\times \mu_n$ v\'erifie $x^\mu=1$ avec $\mu=\mathrm{ppcm}(m,n)$ et est donc d'ordre un diviseur
de $\mu$, lequel est $< mn$ si $m$ et $n$ ne sont pas premiers entre eux. Les groupes
$\mu_{mn}$ et $\mu_m\times \mu_n$ ne peuvent donc pas \^etre isomorphes.
\fincorrection
\nocorrection
\nocorrection
\nocorrection
\correction{002149}
Consid\'erons la surjection canonique $s:G\rightarrow G/H$. D'apr\`es l'exercice \ref{ex:le12}, $|s(K)|$
divise $\mathrm{pgcd}(|K|,|G/H|)$ qui est \'egal \`a $\mathrm{pgcd}(|H|,|G/H|)$ (puisque $|H|=|K|$) et vaut
donc $1$. Conclusion: $s(K)=\{1\}$, c'est-\`a-dire $K\subset H$. D'o\`u $K=H$ puisqu'ils ont
m\^eme ordre.
\fincorrection
\correction{002150}
On a $f(n)=f(1)^n$ pour tout entier $n>0$. Mais on a aussi $f(1/n)^n = f(1)$ pour tout
$n>0$. Cela n'est pas possible car un nombre rationnel positif $\not=0,1$ ne peut \^etre une
puissance $n$-i\`eme dans $\Q$ pour tout $n>0$. (Pour ce dernier point, noter
par exemple qu'\^etre une puissance $n$-i\`eme dans $\Q$ entra\^\i ne que tous les
exposants de la d\'ecomposition en facteurs premiers sont des multiples de $n$). 
Les deux groupes $(\Q,+)$ et $(\Q_+^\times,\times)$ ne sont donc pas isomorphes.
\fincorrection
\nocorrection
\nocorrection
\correction{002153}
On a $n=|G/H|$. Pour toute classe $aH\in G/H$, on a donc $(aH)^n=H$ c'est-\`a-dire,
$a^nH=H$ ou encore $a^n\in H$. Cela devient faux si $H$ n'est pas distingu\'e dans $G$. Par
exemple le sous-groupe $H$ de $S_3$ engendr\'e par la transposition $(1\hskip 2pt 2)$ est
d'indice $3$ dans $S_3$ et, pour $a= (2\hskip 2pt 3)$, on a $a^3=a\notin H$.
\fincorrection
\correction{002154}
Soit $H^\prime$ un sous-groupe de $G$ d'ordre $n$  et d'indice $m$. Pour tout $h\in
H^\prime$, on a $h^n=1$ et $h^m\in H$ (voir l'exercice \ref{ex:le18}). Puisque $n$ et $m$ sont
premiers en eux, on peut trouver $u,v\in \Z$ tels que $um+vn=1$. On obtient alors
$h=(h^m)^u (h^n)^v \in H$. D'o\`u $H^\prime \subset H$ et donc $H=H^\prime$ puisque
$|H|=|H^\prime|$.
\fincorrection
\nocorrection
\correction{002156}
(a) La correspondance $x\rightarrow e^{2i\pi x}$ induit un morphisme $\R\rightarrow T$,
surjectif et de noyau $\Z$. D'o\`u $\R/\Z \simeq T$. La correspondance $z\rightarrow
z/|z|$ induit l'isomorphisme $\C^\times/\R_+^\times \simeq T$. Similairement
$z\rightarrow z^2/|z|^2$ fournit l'isomorphisme $\C^\times/\R^\times \simeq T$. Les
isomorphismes $T/\mu_n \simeq T$ et $\C^\times/\mu_n \simeq \C^\times$ s'obtiennent \`a
partir  de la correspondance
$z\rightarrow z^n$.
\smallskip

(b) La correspondance $x\rightarrow e^{2i\pi x}$ induit un morphisme $\Q\rightarrow
\mu_\infty$, surjectif et de noyau $\Z$. D'o\`u $\Q/\Z \simeq \mu_\infty$. Si $G$ est un
sous-groupe fini de $\mu_\infty$, alors il existe $m\in \N$ tel que $G\subset \mu_m$. Les
sous-groupes du groupe cyclique $\mu_m$ sont les $\mu_n$ o\`u $n|m$. \smallskip

(c) Soit $G$ un sous-groupe de $\Q$ de type fini, c'est-\`a-dire engendr\'e par un nombre
fini de rationnels $p_1/q_1,\ldots,p_r/q_r$. On a alors $q_1\cdots q_r G \subset \Z$. Soit
$q$ le plus petit entier $>0$ tel que $qG\subset \Z$. Le sous-groupe $qG$ est de la forme
$a\Z$ avec $a\in \N$ premier avec $q$ (car l'existence d'un facteur commun contredirait la
minimalit\'e de $q$). On obtient $G=(a/q)\Z$. Si de plus $\Z \subset G$ alors $1\in G$ et
s'\'ecrit donc $1=ka/q$ avec $k\in \Z$, ce qui donne $ka=q$. Comme $\mathrm{pgcd}(a,q)=1$, on a
n\'ecessairement $a=1$ et donc $G = (1/q)\Z$.
\smallskip

Soit $s:\Q\rightarrow \Q/\Z$ la surjection canonique. Si $\overline G$ est un
sous-groupe de type fini de $\Q/\Z$, alors $G=s^{-1}(\overline G)$ est un sous-groupe de
$\Q$, contenant $\Z$ et de type fini (si $p_1/q_1,\ldots,p_r/q_r$ sont des ant\'ec\'edents
par $s$ de g\'en\'erateurs de $\overline G$, alors $1,p_1/q_1,\ldots,p_r/q_r$ engendrent $G$).
D'apr\`es ce qui pr\'ec\`ede, on a $G= \frac{1}{q}\Z$ et donc $\overline G =
\frac{1}{q}\Z/\Z$, qui est isomorphe \`a $\Z/q\Z$.
\smallskip

Via l'isomorphisme de la question (b), on d\'eduit les sous-groupes de $\Q/\Z$ de type fini:
ce sont les sous-groupes $\{e^{2ik\pi/q}\hskip 2pt |\hskip 2pt  k\in \Z\} = \mu_q$ avec $q$
d\'ecrivant
$\N^\times$. \smallskip

(d) On v\'erifie sans difficult\'e que pour tout nombre premier $p$, $\mu_{p^\infty}$ est
un sous-groupe de $\mu_\infty$. Il n'est pas de type fini: en effet le sous-groupe de
$\Q/\Z$ qui lui correspond par l'isomorphisme de la question (b) est engendr\'e par les
classes de rationnels $1/p^n$ modulo $\Z$, $n$ d\'ecrivant $\N$. Un tel sous-groupe $G$ n'a
pas de d\'enominateur commun, c'est-\`a-dire, il n'existe pas d'entier $q\in \Z$ tel que
$qG\subset G$. En cons\'equence il ne peut pas \^etre de type fini. 
\fincorrection
\correction{002157}
Soit $z\in \C$ quelconque et $\zeta\in \C$ une racine $n$-i\`eme de $z$. Le sous-groupe $G$
est distingu\'e dans $\C$ (puisque $\C$ est commutatif). Si $n$ est l'indice de $G$ dans
$\C$, on a donc $\zeta^n = z\in G$ (voir l'exercice \ref{ex:le18}). D'o\`u $\C\subset G$. L'inclusion
inverse est triviale.
\fincorrection
\nocorrection
\nocorrection
\correction{002160}
(a) Soit $\varphi: \mathbb{F}_p^n \rightarrow \mathbb{F}_p^m$ un morphisme de groupes. Pour tout $n\in \Z$,
on note $\overline n \in \Z/p\Z = \mathbb{F}_p$ sa classe modulo $p$. Tout \'el\'ement $\overline x
\in \mathbb{F}_p^n$ peut s'\'ecrire $\overline x=(\overline{x_1},\ldots, \overline{x_n})$ avec
$x=(x_1,\ldots,x_n) \in \Z^n$. On a alors $\varphi(\overline n \cdot \overline x)= \varphi
(\overline{nx}) = \varphi (n \overline x) = n\varphi (\overline x) = \overline n \cdot
\varphi (\overline x)$. Le morphisme $\varphi$ est donc compatible avec les lois
externes de $\mathbb{F}_p^n$ et $\mathbb{F}_p^m$. Comme il est aussi additif, c'est une application
$\mathbb{F}_p$-lin\'eaire.
\smallskip

(b) Consid\'erons l'application $V:\hbox{\rm Aut}(\Z/p\Z) \rightarrow \Z/p\Z$ qui \`a
tout automorphisme $\chi$ associe $\chi (1)$. Cette application est \`a valeurs dans   
$\Z/p\Z \setminus \{0\}$ (si $\chi \in \hbox{\rm Aut}(\Z/p\Z)$, alors $\hbox{\rm
ker}(\chi)=\{0\}$). C'est un morphisme de $\hbox{\rm Aut}(\Z/p\Z)$ muni de la composition
vers le groupe multiplicatif $\Z/p\Z \setminus \{0\} = \mathbb{F}_p^\times$: en effet si
$\chi, \chi^\prime \in \hbox{\rm Aut}(\Z/p\Z)$ et si on pose $\chi^\prime (1)=\overline c$
(classe de $c\in \Z$ modulo $p$), alors $(\chi \circ \chi^\prime) (1) = \chi(\overline c) = c
\chi (1)= \overline c \cdot \chi(1) = \chi^\prime(1) \cdot \chi(1) =  \chi(1) \cdot
\chi^\prime(1)$. Ce morphisme $V$ est de plus injectif puisque tout automorphisme $\chi$ de
$\Z/p\Z$ est d\'etermin\'e par $\chi(1)$. Enfin, pour tout $\overline a\in \Z/p\Z$ non nul, la
correspondance $\overline n \rightarrow \overline a \cdot \overline n$ induit un automorphisme
$\chi$ de $\Z/p\Z$ tel que $\chi(1)=\overline a$. L'image du morphisme $V$ est donc tout
$\mathbb{F}_p^\times$. Ce qui \'etablit l'isomorphisme demand\'e.
\smallskip

(c) D'apr\`es la question (a), il s'agit de compter le nombre d'automorphismes lin\'eaires du
$\mathbb{F}_p$-espace vectoriel $\mathbb{F}_p^n$, qui est \'egal au nombre de bases de $\mathbb{F}_p^n$,
c'est-\`a-dire $(p^n-1)(p^n-p)\cdots (p^n -p^{n-1})$.
\fincorrection
\nocorrection
\correction{002162}
Soit $G$ un groupe ab\'elien fini tel que $pG=\{0\}$. Pour tout entier $n\in \Z$ et pour tout
$g\in G$, l'\'el\'ement $ng$ ne d\'epend que de la classe de $n$ modulo $p$; on peut le
noter $\overline n \cdot g$. La correspondance $(\overline n, g)\rightarrow \overline n \cdot
g$ d\'efinit une loi externe sur le groupe additif $(\Z/p\Z)^n$ et lui conf\`ere
ainsi une structure de $\mathbb{F}_p$-espace vectoriel. Cet espace vectoriel, \'etant fini, est de
dimension finie. Il est donc isomorphe comme espace vectoriel, et en particulier comme groupe
\`a $(\Z/p\Z)^n$ pour un certain entier $n\geq 0$.
\fincorrection
\nocorrection
\nocorrection
\correction{002165} Le centre $Z(G)$ est ni trivial (car $G$ est un $p$-groupe) ni
\'egal \`a $G$ (car $G$ non ab\'elien). En utilisant l'exercice \ref{ex:le23}, on voit qu'il n'est pas 
non plus d'ordre $p^2$. Il est donc d'ordre $p$. Mais alors $G/Z(G)$ est d'ordre
$p^2$ et est donc ab\'elien (exercice \ref{ex:le24}). D'apr\`es l'exercice \ref{ex:le29}, on a
alors $D(G)\subset Z(G)$. Comme $D(G)\not=\{1\}$ (sinon $G$ serait ab\'elien), on a
$D(G) = Z(G)$.
\fincorrection
\nocorrection
\nocorrection
\nocorrection
\nocorrection
\nocorrection
\nocorrection
\nocorrection
\nocorrection
\nocorrection
\nocorrection
\nocorrection
\nocorrection
\nocorrection
\nocorrection
\nocorrection
\nocorrection
\nocorrection
\nocorrection
\nocorrection
\nocorrection
\nocorrection
\nocorrection
\nocorrection
\nocorrection
\nocorrection
\nocorrection
\nocorrection
\nocorrection
\nocorrection
\nocorrection
\correction{002167}
(a) On v\'erifie les deux formules: $(a \hskip 2pt b) (b\hskip 2pt c) =
(a\hskip 2pt b\hskip 2pt c)$ pour $a,b,c$ distincts, et $(a \hskip 2pt b) (c\hskip 2pt d) =
(a \hskip 2pt b) (b\hskip 2pt c) (b \hskip 2pt c) (c\hskip 2pt d)= (a\hskip 2pt b\hskip
2pt c)(b\hskip 2pt c\hskip 2pt d)$, pour $a,b,c,d$ distincts. On d\'eduit que toute
permutation paire, produit d'un nombre pair de transpositions, peut s'\'ecrire comme produit
de $3$-cycles. Le groupe altern\'e $A_n$ est donc engendr\'e par les $3$-cycles si $n\geq 3$.
\smallskip 

(b) On a $(1 \hskip 2pt 2\hskip 2pt j)\hskip 2pt (1\hskip 2pt 2\hskip 2pt i)\hskip 2pt
(1\hskip 2pt 2\hskip 2pt j)^{-1}=(2\hskip 2pt j\hskip 2pt i)$ pour $i,j$ distincts et
diff\'erents de $1$ et $2$, et si en plus $k$ est diff\'erent de $1,2,i,j$, on a $(1
\hskip 2pt 2\hskip 2pt k)\hskip 2pt (2\hskip 2pt j\hskip 2pt i)\hskip 2pt (1\hskip 2pt
2\hskip 2pt k)^{-1}=(k\hskip 2pt j\hskip 2pt i)$. Le groupe engendr\'e par les $3$-cycles
$(1\hskip 2pt 2\hskip 2pt i)$ o\`u $i\geq 3$ contient donc tous les $3$-cycles; d'apr\`es
(a), c'est le groupe altern\'e $A_n$.
\fincorrection
\nocorrection
\correction{002169}
Les cas $n=1$ et $n=2$ sont imm\'ediats. On peut supposer $n\geq 3$. On v\'erifie
ais\'ement la formule $(a_1 \hskip 2pt a_2\hskip 2pt\ldots\hskip 2pt a_{n-1}\hskip 2pt a_n)\ 
(a_{n-1} \hskip 2pt a_n\hskip 2pt a_{n-2} \ldots\hskip 2pt a_2\hskip 2pt a_1) = (a_1 \hskip
2pt a_n \hskip 2pt  a_{n-1})$ o\`u $a_1,\ldots,a_n$ sont les \'el\'ements d'un ensemble de
cardinal $n$. On en d\'eduit que le groupe $PC_n$ engendr\'e par les permutations circulaires
contient les $3$-cycles et donc le groupe altern\'e $A_n$ (voir exercice \ref{ex:deb67}). Les
permutations circulaires sont de signature $(-1)^{n-1}$. Si $n$ est impair, elles sont donc
paires d'o\`u $PC_n\subset A_n$ et donc finalement $PC_n = A_n$ dans ce cas. Si $n$ pair, les
permutations circulaires sont impaires, donc $PC_n\not=A_n$. L'indice de $PC_n$ dans $S_n$
devant diviser $2$ (puisque $PC_n\supset A_n$), il vaut $1$, c'est-\`a-dire $PC_n=S_n$.
\fincorrection
\correction{002170}
Supposons $\sigma \tau = \tau \sigma$. Pour tout $x\notin I$, on a
$\sigma(\tau(x))=\tau(\sigma(x))=\tau(x)$; $\tau(x)$, fix\'e par $\sigma$, n'appartient pas
\`a $I$. Cela montre que le compl\'ementaire de $I$ est invariant par $\tau$. Comme $\tau$
est injective, $I$ l'est aussi. Montrons que, sur $I$, $\tau$ est \'egal \`a une puissance de
$\sigma$. Quitte \`a renum\'eroter $\{1,\ldots,n\}$, on peut supposer que
$I=\{1,\ldots,m\}$ (o\`u $m\leq n$) et $\sigma|_I= (1\hskip 2pt 2\hskip 2pt\ldots\hskip 2pt
m)$. L'entier $\tau(1)$ est dans $I$; soit $k$ l'unique entier entre
$1$ et $m$ tel que $\tau(1)=\sigma^k(1)$. Pour tout $i\in I$, on a alors $\tau(i) = \tau
\sigma^{i-1}(1) =
\sigma^{i-1}\tau(1)= \sigma^{i-1}\sigma^{k}(1)=\sigma^{k}\sigma^{i-1}(1)=\sigma^{k}(i)$
(l'identit\'e $\tau \sigma^{i-1}= \sigma^{i-1}\tau$ utilis\'ee dans le calcul d\'ecoule
facilement de l'hypoth\`ese $\sigma \tau = \tau \sigma$). On obtient donc
$\tau|_I=(\sigma|_I)^k$. L'implication r\'eciproque est facile.
\fincorrection
\correction{002171} Un sous-groupe distingu\'e de $S_n$ qui contient une transposition
contient toute sa classe de conjugaison, c'est-\`a-dire, toutes les transpositions (cf
les indications de l'exercice \ref{ex:deb68}, ``Rappel'') et donc le groupe qu'elles engendrent, 
c'est-\`a-dire $S_n$.
\fincorrection
\correction{002172}
L'ensemble $H$ est le sous-groupe de $S_4$ fixant la paire $\{1,2\}$. Tout \'el\'ement de $H$
fixe aussi la paire $\{3,4\}$. Cela fournit un morphisme $H\rightarrow S_2 \times S_2$ qui est
clairement bijectif. D'o\`u $H\simeq S_2 \times S_2 \simeq \Z/2\Z \times \Z/2\Z$.
\medskip

\hskip 5mm On a $\sigma\in K$ si et seulement si $\sigma(1) \equiv \sigma(3)\ [\hbox{\rm mod}\ 2]$ et
$\sigma(2) \equiv \sigma(4)\ [\hbox{\rm mod}\ 2]$, c'est-\`a-dire si et seulement si
$\sigma(\{1,3\})$ est soit la paire $\{1,3\}$ soit la paire $\{2,4\}$ (auquel cas
$\sigma(\{2,4\})$ est la paire $\{2,4\}$ ou la paire $\{1,3\}$ respectivement).
Gr\^ace \`a l'identit\'e $\sigma \hskip 2pt(1\hskip 2pt 3)\hskip 2pt(2\hskip 2pt 4)\hskip
2pt \sigma^{-1} = (\sigma(1)\hskip 2pt \sigma(3))\hskip 2pt(\sigma(2)\hskip 2pt
\sigma(4))$, on voit que la condition est \'egalement \'equivalente au fait que la
conjugaison par $\sigma$ stabilise la permutation $(1\hskip 2pt 3)\hskip 2pt(2\hskip 2pt 4)$.
Autrement dit $K$ est le sous-groupe des \'el\'ements de $S_4$ commutant avec $(1\hskip 2pt
3)\hskip 2pt(2\hskip 2pt 4)$. La classe de conjugaison 2-2 ayant $3$ \'el\'ements, le
groupe $H$ est d'ordre $4!/3=8$. On peut dresser la liste de ses \'el\'ements: si $\omega =
(1\hskip 2pt 2\hskip 2pt 3\hskip 2pt 4)$ et $\tau = (1\hskip 2pt 2)\hskip 2pt (3\hskip 2pt
4)$, alors
$K=\{1,\omega,\omega^2,\omega^3,\tau,\omega \tau, \omega^2 \tau, \omega^3\tau\}$. On
v\'erifie les relations $\sigma^4=1$, $\tau^2=1$ et $\tau \sigma \tau^{-1}= \sigma^{-1}$. Le
groupe $K$ est \'egal au produit semi-direct de son sous-groupe distingu\'e
$<\omega>$ par son sous-groupe $<\tau>$ et est donc isomorphe au groupe di\'edral 
d'ordre $8$.
\fincorrection
\nocorrection 
\correction{002174} 
L'ordre d'une permutation $\omega \in S_n$ est le $\mathrm{ppcm}$ des longueurs
des cycles de la d\'ecomposition de $\omega$ en cycles \`a supports disjoints. De plus, la
somme des longueurs de ces cycles (ceux de longueur $1$ y compris) vaut $n$. Pour une
permutation d'ordre $10$ dans $S_8$, il n'y a qu'un type possible: 5-2-1. La signature vaut
alors $(-1)^{5-1} (-1)^{2-1}=-1$.
\fincorrection
\correction{002175}
(a) Un $3$-cycle $\omega$ est d'ordre $3$ et v\'erifie donc $\omega^3=1$ soit encore $\omega=
(\omega^2)^2$. Le groupe engendr\'e par tous les carr\'es de permutations dans $S_n$ contient
donc tous les $3$-cycles, et donc aussi le groupe qu'ils engendrent, c'est-\`a-dire $A_n$.
L'autre inclusion est facile puisque le carr\'e d'une permutation est toujours une 
permutation paire. 
\smallskip

(b) Si $H$ est un sous-groupe d'indice $2$ de $S_n$, il est distingu\'e. On a alors
$\sigma^2 \in H$ pour tout $\sigma \in S_n$ (cf exercice \ref{ex:le18}). D'apr\`es la question
(a), $H=A_n$.
\fincorrection
\correction{002176}
Les classes de conjugaison de $S_n$ correspondent aux types possibles d'une permutation de
$n$ \'el\'ements (cf indication exercice 3 Rappel). Pour $n=4$, on a $5$ classes:
1-1-1-1, 2-1-1, 2-2, 3-1 et 4.
\smallskip

\hskip 5mm Soit $H$ un sous-groupe distingu\'e non trivial de $S_4$. Si $H$ contient la classe 2-1-1
(transpositions), alors $H=S_4$. Si $H$ contient la classe 3-1, alors $H\supset A_4$
(cf exercice \ref{ex:deb67}) et donc $H=A_4$ ou $H=S_4$. Si $H$ contient la classe 4, alors $H=S_4$
(cf exercice \ref{ex:deb69}). Si $H$ contient la classe 2-2, alors $H\supset V_4$ (voir la correction
de l'exercice \ref{ex:le7} d\'efinition de $V_4$), ce qui donne $H=V_4$ ou bien, au vu
des cas pr\'ec\'edents, $H=A_4$ ou $H=S_4$. Les sous-groupes distingu\'es de $S_4$ sont
donc $\{1\}$,
$V_4$, $A_4$ et $S_4$.
\fincorrection
\nocorrection
\correction{002178}
Soit $H$ un sous-groupe d'indice $m$ d'un groupe $G$.  L'action de $G$ par translation \`a gauche sur l'ensemble quotient $G/H$ des classes \`a gauche modulo $H$ induit un morphisme $G\rightarrow \textrm{Per}(G/H)$ qui est non-trivial et donc est injectif puisque le noyau, distingu\'e dans $G$, ne peut \^etre trivial si $G$ est simple. L'ordre de $G$ doit donc diviser l'ordre du groupe $\textrm{Per}(G/H)$ qui vaut $m!\ $. Il faut n\'ecessairement que $|G|=m!\ $. Mais alors le morphisme pr\'ec\'edent est un isomorphisme et $G$ est isomorphe au groupe sym\'etrique $S_m$, ce qui contredit la simplicit\'e de $G$. 
\fincorrection
\nocorrection
\correction{002180}
(a) L'identit\'e $a^2b^2=(ab)^2$, par simplification \`a gauche par $a$ et \`a
droite par $b$, se r\'e\'ecrit $ab=ba$.
\smallskip

(b) La correspondance $(x,y) \rightarrow (x+y,y)$ d\'efinit un automorphisme
$\sigma$ de $\mathbb{F}_3^2$ d'ordre $3$. Identifions le groupe $<\sigma>$ au
groupe $\Z/3\Z$ et consid\'erons le produit semi-direct  $\mathbb{F}_3^2 \times \hskip -6pt
{\raise 1.4pt\hbox{${\scriptscriptstyle |}$}} \Z/3\Z$. Pour tout \'el\'ement
$((x,y),i)$, on a $((x,y),i)^2 = ((x,y)+\sigma^i(x,y), 2i)$ et $((x,y),i)^3 =
((x,y)+\sigma^i(x,y)+\sigma^{2i}(x,y), 3i) = ((0,0),0)$ puisque $({\rm Id}+\sigma^i
+\sigma^{2i})(x,y) = (3x+iy+2iy,3y)=(0,0)$. La formule $a^3b^3=(ab)^3$ est donc
satisfaite pour tous $a,b$ dans $\mathbb{F}_3^2 \times \hskip -5pt
{\raise 1.4pt\hbox{${\scriptscriptstyle |}$}} \Z/3\Z$. Mais ce produit semi-direct
n'est pas commutatif car l'action de $\Z/3\Z$ n'est pas l'action triviale.
\fincorrection
\correction{002181}
(a) Que $R$ soit une relation d'\'equivalence est imm\'ediat. La classe d'un
\'el\'ement $x\in G$ est l'ensemble $HxH$, lequel est \'egal \`a la r\'eunion
des ensembles $hxH$ o\`u $h$ d\'ecrit $H$. Ces derniers ensembles sont des classes
\`a gauche modulo $H$ et sont donc \'egaux ou disjoints.
\smallskip

(b) Pour tout $i=1,\ldots,d(x)$, $hx_iH$ est une classe \`a gauche, contenue dans
$h(HxH)H \subset HxH$, donc est de la forme $x_jH$. La formule $h\hskip 2pt \ast
\hskip 2pt x_iH= hx_iH$ d\'efinit ainsi une permutation de l'ensemble  des classes
$x_1H, \dots, x_{d(x)} H$ (la permutation r\'eciproque est celle induite par
$h^{-1}$) et donc une action de $H$ sur cet ensemble. Cette action est transitive:
pour $i,j\in \{1,\ldots,d(x)\}$, $h=x_i^{-1}x_j$ v\'erifie $h\hskip 2pt \ast
\hskip 2pt x_iH= x_jH$.

\hskip 5mm Un \'el\'ement $h\in H$ est dans le fixateur $H(x_iH)$ d'une classe
$x_iH$ si et seulement si $hx_iH=x_iH$ c'est-\`a-dire si $h\in x_iHx_i^{-1}$. D'o\`u
$H(x_iH)= H\cap x_iHx_i^{-1}$. On obtient alors $d(x)= [H:(H\cap x_iHx_i^{-1})]$ ce
qui prouve que $d(x)$ divise $|H|$ et donc aussi $|G|$.
\smallskip

(c) Si $H$ est distingu\'e dans $G$, alors classes \`a droite et classes \`a
gauche modulo $H$ coincident d'o\`u $HxH=xHH=xH$ et donc $d(x)=1$ pour tout $x\in
G$. Inversement, pour tout $x\in G$, si $d(x)=1$, alors $HxH=xH$ ce qui entraine
$Hx\subset xH$ et donc $x^{-1}Hx \subset H$.
\smallskip

(d) (i) De fa\c con g\'en\'erale, on a $d(x) \leq [G:H]$. On a ainsi $d(x) \leq p$
si $[G:H]=p$. Comme $d(x)$ divise $|G|$ et que $p$ est le plus petit premier
divisant $|G|$, n\'ecessairement $d(x)=1$ ou $d(x)=p$.
\smallskip

(ii) Si $H$ n'est pas distingu\'e alors il existe $x\in G$ avec $d(x)\not=1$ et
donc $d(x)=p$. Mais alors $\hbox{\rm card}(HxH) = d(x)\hskip 2pt |H| = p\hskip 2pt
|H| = [G:H]\hskip 2pt |H| = |G|$. C'est-\`a-dire, il n'existe qu'une seule classe
$HxH=G$, laquelle est aussi la classe de l'\'el\'ement neutre $H1H=H$, ce qui
contredit l'hypoth\`ese $[G:H]=p>1$. Conclusion: le sous-groupe $H$ est distingu\'e
dans $G$.
\fincorrection
\correction{002182}
Toute orbite ${\cal O}={\cal O}_x$ d'un \'el\'ement $x\in X$ est en bijection avec
l'ensemble $G/\cdot G(x)$ des classes \`a gauche de $G$ modulo le fixateur $G(x)$ de $G$. En
particulier, le cardinal de ${\cal O}$ divise l'ordre de $G$. De plus la somme des longueurs
des orbites est \'egale au cardinal de l'ensemble $X$.
\smallskip 

(a) Si $|G|=15$, $\hbox{\rm card}(X)=17$ et s'il n'y a pas d'orbite \`a un
seul \'el\'ement, il n'y a qu'une seule possibilit\'e: 4 orbites de longueur $3$ et une de
longueur
$5$.
\smallskip

(b) Supposons $|G|=33$ et $\hbox{\rm card}(X)=19$. Aucune somme de diviseurs $\not=1$ de $33$
n'est \'egale \`a $19$ donc n\'ecessairement il existe au moins une orbite r\'eduite \`a un
\'el\'ement.
\smallskip
\fincorrection
\correction{002183}
(a) Si $g_1^\prime, g_2^\prime$ sont dans la m\^eme classe \`a gauche de $G$ modulo $H$,
c'est-\`a-dire, si $g_1^\prime H= g_2^\prime H$ ou encore si $(g_2^\prime)^{-1}g_1^\prime
\in H$ alors $(gg_2^\prime)^{-1}(gg_1^\prime)=(g_2^\prime)^{-1}g_1^\prime
\in H$: les classes $gg_1^\prime H$ et $gg_2^\prime H$ sont \'egales. Pour tous $g,g^\prime
\in H$, la classe $gg^\prime H$ ne d\'epend donc pas du repr\'esentant choisi $g^\prime$ de la
classe $g^\prime H$; on peut la noter $g\cdot g^\prime H$. On v\'erifie sans difficult\'e
que la correspondance $(g,g^\prime H)\rightarrow g\cdot g^\prime H$ satisfait les autres
conditions de la d\'efinition d'une action de $G$ sur l'ensemble quotient $G/\cdot H$.
\smallskip

\hskip 5mm Pour $g,\gamma\in G$, on a $\gamma\cdot gH = gH$ si et seulement si $g^{-1} \gamma g \in H$ ce qui \'equivaut \`a $\gamma \in gHg^{-1}$. Le fixateur de la classe $gH$ est le
sous-groupe conjugu\'e $gHg^{-1}$ de $H$ par $g$.
\medskip

(b) Pour tout $y\in Y$ et tout $g\in G$, on a $f(g\cdot f^{-1}(y)) = g\cdot f(f^{-1}(y)) =
g\cdot y$. En appliquant $f^{-1}$, on obtient $g\cdot f^{-1}(y)=f^{-1}(g\cdot y)$, ce qui
montre que $f^{-1}$ est compatible \`a l'action de $G$. 
\medskip

(c) Soit $x\in X$ fix\'e. Pour $g\in G$, l'\'el\'ement $g\cdot x$ ne d\'epend que de la classe
\`a gauche de $g$ modulo le fixateur $G(x)$ de $x$. Cela permet de d\'efinir une application
$G/\cdot G(x) \rightarrow X$: \`a chaque classe $gG(x)$ on associe $g\cdot x$. On montre
sans difficult\'e que cette application est compatible avec l'action de $G$ (v\'erification
formelle), injective (par construction) et surjective (par l'hypoth\`ese de transitivit\'e);
c'est donc un isomorphisme de $G$-ensembles.
\medskip

(d) i) Supposons donn\'ee une application $f:G/\cdot H\rightarrow G/\cdot K$
compatible avec l'action de $G$. Pour tout $h\in H$, on a $f(hH)=f(H)=h\cdot f(H)$. Ce qui,
d'apr\`es la question (a), donne $h\in gKg^{-1}$, o\`u $g$ est un repr\'esentant de la classe
$f(H)$ dans $G/\cdot K$.
\smallskip

\hskip 5mm R\'eciproquement, supposons $H\subset gKg^{-1}$ avec $g\in G$. Consid\'erons l'application
$\varphi: G/\cdot H \rightarrow G/\cdot K$ qui \`a toute classe $\gamma H$ associe la classe
$\gamma g K$. Cette application est bien d\'efinie: en effet, si $\gamma_2^{-1}\gamma_1\in H$,
alors $(\gamma_2g)^{-1} \gamma_1g = g^{-1}(\gamma_2^{-1}\gamma_1)g \in g^{-1}Hg\subset K$;
la classe $\gamma g K$ ne d\'epend donc pas du repr\'esentant $\gamma$ de la classe $\gamma
H$. De plus $\varphi$ est compatible \`a l'action de $G$: pour tous $\gamma, \gamma^\prime \in
G$, on a $\varphi (\gamma^{\prime}\cdot\gamma H) = \varphi (\gamma^{\prime}\gamma H) =
\gamma^{\prime}\gamma gK = \gamma^{\prime}\cdot \varphi (\gamma H)$.
\smallskip

\hskip 5mm Si $f:G/\cdot H\rightarrow G/\cdot K$ est compatible avec l'action de $G$, alors son image
contient toute orbite d\`es qu'elle en contient un \'el\'ement. Comme l'action de $G$ sur sur
$G/\cdot K$ ne poss\`ede qu'une orbite, l'image de $f$ contient tout $G/\cdot K$: $f$ est
surjective.
\smallskip

\hskip 5mm D'apr\`es ce qui pr\'ec\`ede, les ensembles $G/\cdot H$ et $G/\cdot K$ sont isomorphes
comme $G$-ensembles si et seulement si $H\subset gKg^{-1}$ avec $g\in G$ et $\hbox{\rm
card}(G/\cdot H) = \hbox{\rm card}(G/\cdot K)$ ce qui \'equivaut \`a $H\subset gKg^{-1}$ et
$|H|=|K|$ ou encore \`a $H=gKg^{-1}$.
\smallskip

ii) Il suffit de r\'e\'ecrire les r\'esultats de la question pr\'ec\'edente en rempla\c cant
$G/\cdot H$ et $G/\cdot K$ par $G/\cdot G(x)$ et $G/\cdot G(y)$ qui, d'apr\`es la question (c)
sont $G$-isomorphes \`a $X$ et $Y$ respectivement (o\`u $x$ et $y$ sont des points fix\'es de
$X$ et $Y$ respectivement).
\fincorrection
\correction{002184}
(a) Pour $1\leq i,j \leq r$ quelconques et $x_i,x_j \in X_i \times X_j$, il existe $g\in G$
tel que $g\cdot x_i=x_j$ (par transitivit\'e de $G$). On a alors $g\cdot X_i = X_j$. En
particulier
$\hbox{\rm card}(X_i)= \hbox{\rm card}(g\cdot X_i)=\hbox{\rm card}(X_j)$.
\medskip

(b) Si l'action de $G$ sur $G/\cdot H$ est imprimitive, le sous-ensemble $K=\{g\in G\hskip 2pt
|\hskip 2pt g\cdot X_1 = X_1\}$, o\`u $X_1$ est par exemple celui des sous-ensembles
$X_i\subset X$ qui contient la classe neutre $H$ de $G/\cdot H$, est un sous-groupe propre de
$G$ ($K\not=G$ car
$G$ agissant transitivement, il existe $g\in G$ tel que $(g\cdot X_1) \cap X_2\not=\emptyset$)
et contenant strictement $H$ (car encore par transitivit\'e, il existe $g\in G$ tel que
$g\cdot H$ soit un \'el\'ement de $X_1$ (ce qui assure que $g\in K$) mais diff\'erent de $H$
(ce qui assure que $g\notin H$)). 

\hskip 5mm Inversement, si un tel sous-groupe $K$ de $G$ existe, la relation ``$gH \sim g^\prime H$
si $(g^\prime)^{-1} g\in K$'' est bien d\'efinie sur $G/\cdot H$ (la d\'efinition ne
d\'epend pas des repr\'esentants dans $G$ des classes $gH$ et $g^\prime H$) et est une
relation d'\'equivalence (imm\'ediat). La partition associ\'ee de $G/\cdot H$ en classes
d'\'equivalence v\'erifie les conditions de la d\'efinition d'imprimitivit\'e (pour l'action
de $G$ sur $G/\cdot H$): la partition est non triviale car $K$ est strictement contenu entre
$H$ et $K$; et si $(\gamma H)K$ est une de ces classes d'\'equivalence et
$g\in G$, alors $g\cdot (\gamma H)K$ est la classe $(g\gamma H)K$: l'action de $G$ permute
bien les classes constituant la partition de $X$.
\medskip

(c) D'apr\`es l'exercice \ref{ex:deb83}, les ensembles $X$ et $G/\cdot G(x)$ sont isomorphes comme
$G$-ensembles. L'action de $G$ sur $X$ est primitive si et seulement si celle de $G$ sur
$G/\cdot G(x)$ l'est, ce qui, d'apr\`es la question pr\'ec\'edente, \'equivaut \`a dire que le
fixateur $G(x)$ est maximal parmi les sous-groupes de $G$.
\medskip 

(d) Soient $x\in X$ et $G(x)$ son fixateur. Le sous-groupe $H$ \'etant distingu\'e dans
$G$, l'ensemble $HG(x)$ est un sous-groupe; c'est le sous-groupe engendr\'e par $H$ et
$G(x)$. De plus, l'action de $H$ sur $G$ n'\'etant pas triviale, $H$ n'est pas contenu dans
$G(x)$ et par cons\'equent $HG(x)$ contient strictement $G(x)$. D'apr\`es la question (c), il
en r\'esulte que $HG(x) = G$. On v\'erifie sans peine que l'application $H/\cdot (H\cap G(x))
\rightarrow (HG(x))/\cdot G(x)$ qui \`a toute classe $h(H\cap G(x))$ associe la classe $hG(x)$
est une bijection (ce qui g\'en\'eralise le th\'eor\`eme d'isomorphisme
$HK/K\simeq H/(H\cap K)$ qui est vrai sous l'hypoth\`ese suppl\'ementaire ``$K$
distingu\'e'' (qui assure que les ensembles $HK/K$ et $H/(H\cap K)$ sont des groupes et non
de simples ensembles comme ici)). On obtient donc que les ensembles $H/\cdot (H\cap
G(x))$ et $G/\cdot G(x)$ sont isomorphes comme $G$-ensembles (la compatibilit\'e des
actions est imm\'ediate). Or ces deux ensembles sont en bijection avec les orbites de $x$
sous $H$ et sous $G$ respectivement. Conclusion: l'action de $H$ est, comme celle de $G$,
transitive sur l'ensemble $X$. 
\fincorrection
\correction{002185}
Soit $H$ un sous-groupe primitif de $S_n$ contenant une transposition. On peut supposer que
$H$ contient la transposition $(1\hskip 2pt 2)$. Le sous-groupe engendr\'e par le fixateur
$H(1)$ et $(1\hskip 2pt 2)$ contient strictement $H(1)$. D'apr\`es
l'exercice \ref{ex:deb84} (question (c)), ce groupe est $H$.
\smallskip

\hskip 5mm Consid\'erons l'ensemble ${\cal O}$ r\'eunion de l'orbite $H(1)\cdot 2$ de $2$ sous $H(1)$ et du singleton $\{1\}$. Pour montrer que ${\cal O}$ est l'orbite de $2$ sous $H$, il suffit de
montrer que $2\in {\cal O}$ (ce qui est clair) et que ${\cal O}$ est stable sous l'action de
$H$, ou, ce qui est \'equivalent, stable sous l'action de $H(1)$ et de $(1\hskip 2pt 2)$.
L'\'el\'ement $1$ est envoy\'e sur $1\in {\cal O}$ par les \'el\'ements de $H(1)$
et sur $2\in {\cal O}$ par $(1\hskip 2pt 2)$. L'ensemble $H(1)\cdot 2$ est invariant sous
l'action de $H(1)$. Enfin, si $h\cdot 2$ d\'esigne un \'el\'ement quelconque de $H(1)\cdot
2$, alors son image par la permutation $(1\hskip 2pt 2)$ est $2$ si $h\cdot 2=1$, $1$ si
$h\cdot 2=2$ et $h\cdot 2$ si $h\cdot 2\not=1,2$; dans tous les cas, l'image est dans ${\cal
O}$. 
\smallskip

\hskip 5mm On a donc ${\cal O} = H\cdot 2 = H(1)\cdot 2 \cup \{1\}$. L'action de $H$ \'etant
transitive, cet ensemble est \'egal \`a $\{1,\ldots,n\}$ et donc $H(1)\cdot
2=\{2,\ldots,n\}$ (puisque $1\notin H(1)\cdot 2$). Cela montre que l'action de $H(1)$ sur
$\{2,\ldots,n\}$ est transitive, et donc que $H$ agit transitivement sur $\{1,\ldots,n\}$ (exercice 21).

\hskip 5mm Pour $i,j$ entiers distincts entre $1$ et $n$, choisissons alors $g\in G$ tel que $g(1)=i$
et $g(2)=j$. On a $g (1\hskip 2pt 2) g^{-1} = (g(1)\hskip 2pt g(2)) = (i\hskip 2pt j)$. Cela
montre que $H$ contient toutes les transpositions. Conclusion: $H=S_n$.
\fincorrection
\nocorrection
\nocorrection
\correction{002188}
Notons $G$ le groupe des isom\'etries de l'espace euclidien de dimension $3$ laissant
invariant l'ensemble $\{a_1,\ldots,a_4\}$ des $4$ sommets d'un t\'etra\`edre r\'egulier.
Le fixateur $G(a_4)$ agit transitivement sur $\{a_1,a_2,a_3\}$: en effet ce sous-groupe
contient la rotation d'axe la droite joignant $a_4$ au centre de gravit\'e du triangle
de sommets $a_1,a_2,a_3$, laquelle agit sur ces points comme un $3$-cycle. D'apr\`es
l'exercice \ref{ex:deb86}, le groupe $G$ agit $2$-transitivement sur
$\{a_1,\ldots,a_4\}$. De plus $G(a_4)$ contient une isom\'etrie agissant sur
$\{a_1,\ldots,a_4\}$ comme une transposition, par exemple la sym\'etrie par rapport au
plan m\'ediateur $P$ du segment $[a_1,a_2]$, laquelle
\'echange $a_1$ et $a_2$ et fixe $a_3$ et $a_4$ qui sont dans $P$. D'apr\`es
l'exercice \ref{ex:deb85}, on a $G\simeq S_4$.
\smallskip

Notons $G_+$ le sous-groupe de $G$ constitu\'e de ses isom\'etries directes. Le groupe
$G_+$ est le noyau du morphisme $\hbox{\rm det}: G_+ \rightarrow \{1,-1\}$ qui \`a tout
$g\in G$ vu comme matrice associe son d\'eterminant. Comme ce morphisme est surjectif
(la rotation et la sym\'etrie consid\'er\'ees ci-dessus sont respectivement    
directe et indirecte), $G_+$ est d'indice $2$. D'o\`u $G\simeq A_4$ puisque $A_4$ est le
seul sous-groupe de $S_4$ d'indice $2$ (cf exercice \ref{ex:deb75}). 
\fincorrection
\nocorrection
\nocorrection
\nocorrection
\nocorrection
\nocorrection
\nocorrection
\nocorrection
\correction{007729}
\begin{enumerate}
 \item Soit $G$ un groupe fini agissant sur un ensemble fini $E$.
 Alors le nombre $N$ d'orbites se calcule par 
 $$N=\frac{1}{\text{Card} G}\sum_{g\in G}\text{Card} Fix(\phi(g))=
 \frac{1}{\text{Card} G}\sum_{x\in E}\text{Card} Stab(x).$$
 \item Le groupe des déplacements du tétraède est isomorphe à $\mathcal{A}_4$.
 Il y a l'identité, les $8=4\times 2$ rotations d'ordre $3$, dont l'axe passe par un sommet et le centre de la face opposée,
 les $3$ demi-tours (d'ordre 2) dont l'axe passe par les milieux de deux arêtes opposées.
 \item On considère l'ensemble des coloriages. Il est de cardinal $c^4$, puisqu'il s'agit de choisir une couleur pour chacune des quatre faces. Le groupe des déplacements du tétraèdre agit sur cet ensemble et les façons différentes de colorier sont les orbites de cette action. Par la formule de Burnside, il suffit de déterminer le nombre de coloriages fixés par chacun des déplacements.
 Pour l'identité : $c^4$.
 Pour les rotations d'ordre $3$ : $c^2$ (si l'axe est vertical, toutes les trois faces obliques doivent avoir la même couleur et la face horizontale doit avoir une couleur.)
 Pour les demi-tours : $c^2$ (les faces se correspondent deux à deux.).
 En conclusion, le nombre de façons différentes de colorier est
 $$N=\frac{1}{\text{Card} G}\sum_{g\in G}\text{Card} Fix(\phi(g))=\frac{1}{c^4}(c^4+11c^2)=\frac{1}{c^2}(c^2+11).$$  
\end{enumerate}
\fincorrection
\nocorrection
\nocorrection
\nocorrection
\nocorrection
\nocorrection
\nocorrection
\nocorrection
\nocorrection
\nocorrection
\nocorrection
\nocorrection
\nocorrection
\nocorrection
\nocorrection
\nocorrection
\nocorrection
\nocorrection
\nocorrection
\nocorrection
\nocorrection
\nocorrection
\correction{002191}
Le sous-groupe $H\subset G$ \'etant distingu\'e, $G$ agit par conjugaison sur $H$.
Comme $G$ est un $p$-groupe, $H$ l'est aussi et les orbites non triviales de cette
action sont de longueur divisible par $p$. On d\'eduit que la r\'eunion des orbites
triviales, c'est-\`a-dire l'ensemble $H\cap Z(G)$ des points fixes, est aussi de
cardinal divisible par $p$. Comme il contient l'\'el\'ement neutre, il contient au
moins $p$ \'el\'ements et n'est donc pas r\'eduit \`a l'\'el\'ement neutre.
\fincorrection
\correction{002192} 
(a) Soit $G$ un $p$-groupe d'ordre $p^r$. Son centre $Z(G)$ est un $p$-groupe    
non trivial. Soit $x\in Z(G)\setminus\{1\}$. Si $p^\nu>0$ est son ordre, alors
$x^{p^{\nu-1}}$ est d'ordre $p$ et dans $Z(G)$; on peut donc supposer que $x$
lui-m\^eme est d'ordre $p$. Le groupe $<x>$ est distingu\'e dans
$G$ et le groupe quotient $G/<x>$ est d'ordre $p^{r-1}$. Par hypoth\`ese
de r\'ecurrence, pour tout $k\leq r$, le groupe $G/<x>$ poss\`ede un sous-groupe
distingu\'e ${\cal H}$ d'ordre $p^{k-1}$. Soit $H$ le sous-groupe image r\'eciproque
de ${\cal H}$ par la surjection canonique $G\rightarrow G/<x>$. Le sous-groupe $H$,
image r\'eciproque par un morphisme d'un sous-groupe distingu\'e, est
distingu\'e dans $G$ et ${\cal H}=H/<x>$, ce qui donne $|H| = |{\cal H}| \hskip 2pt
|<x>| = p^k$. 
\fincorrection
\correction{002193}
Comme $p$ divise $|G|$, il existe dans $G$ un \'el\'ement $s$ d'ordre $p$. Le
sous-groupe $H=\hskip 2pt <s>$, d'indice $2$, est n\'ecessairement distingu\'e dans
$G$. Il est de plus le seul sous-groupe d'ordre $p$ (cf l'exercice \ref{ex:le19}).

\hskip 5mm De fa\c con g\'en\'erale, un automorphisme $\chi$ d'un groupe cyclique
$<\zeta>$ d'ordre $p$ est d\'etermin\'e par $\chi(\zeta) = \zeta^{i_\chi}$ et cet
automorphisme est d'ordre $2$ si et seulement si $i_\chi^2 \equiv 1\ [\hbox{\rm
mod}\ p]$, c'est-\`a-dire si $\chi(\zeta) = \zeta$ ou $\chi(\zeta) = \zeta^{-1}$ ce
qui correspond aux deux automorphismes ``identit\'e'' et ``passage \`a l'inverse''
(que $p$ soit premier n'intervient pas ici; le r\'esultat est valable pour tout
entier $p\geq 1$).

\hskip 5mm Soit $t\in G$ d'ordre $2$ (qui existe car $2$ divise $|G|$). La
conjugaison par $t$ induit un automorphisme du sous-groupe distingu\'e $H$.
D'ap\`es ce qui pr\'ec\`ede, on a $tst^{-1}=s$ ou bien $tst^{-1}=s^{-1}$. Dans
le premier cas, la correspondance $(s^i,t^\varepsilon) \rightarrow
s^i \cdot t^\varepsilon$ ($i=0,1,2$ et $\varepsilon = \pm 1$) induit un morphisme
entre le produit direct $<s> \times <t>$ et $G$, lequel est injectif (car $<s>
\cap <t>=\{1\}$) et donc est bijectif (puisque les groupes de d\'epart et
d'arriv\'ee ont m\^eme ordre $2p$). Dans ce cas on a donc $G\simeq \Z/p\Z
\times \Z/2\Z \simeq \Z/2p\Z$ cyclique. Dans l'autre cas, $G$ est non
commutatif (puisque $tst^{-1}=s^{-1}\not=s$); il est engendr\'e par
$s$ et $t$ qui v\'erifient les relations $s^p=1$, $t^2=1$ et $tst^{-1} =s^{-1}$.
Dans ce cas $G$ est isomorphe au groupe di\'edral $\Z/p\Z \times \hskip -6pt
{\raise 1.4pt\hbox{${\scriptscriptstyle |}$}} \Z/2\Z$ d'ordre $2p$.
\fincorrection
\correction{002194}
(a) Le groupe $G$ n'\'etant pas ab\'elien n'est pas cyclique d'ordre $8$ et
poss\`ede au moins un \'el\'ement $a\not= 1$ qui n'est pas d'ordre $2$ (cf l'exercice \ref{ex:deb14}). 
Cet \'el\'ement est n\'ecessairement d'ordre $4$. Le sous-groupe
$H=\hskip 2pt <a>$ est distingu\'e car d'indice $2$.
\smallskip

(b) Supposons qu'il existe $b\in G\setminus H$ d'ordre $2$ et posons $K=\hskip
2pt <b>$. On a $H\cap K=\{1\}$ car $b\notin H$. Le sous-groupe
$H$ \'etant distingu\'e dans $G$, on peut \'ecrire que $HK/H \simeq K$, ce qui
donne $|HK|= |H|\hskip 2pt |K| = 8$ et donc $G=HK$. De plus, l'inclusion $K\subset
G$ est une section de la suite exacte $1 \rightarrow H \rightarrow G \rightarrow
K \rightarrow 1$. Le groupe
$G$ est donc isomorphe au produit semi-direct de $H$ par $K$. L'action sur $H$ du
g\'en\'erateur
$b$ d'ordre $2$ de $K$ est n\'ecessairement donn\'ee par le passage \`a l'inverse
(cf exercice \ref{ex:deb93}).
\smallskip

(c) Dans le cas contraire \`a (b), tous les \'el\'ements de $G\setminus H$ sont
n\'ecessairement d'ordre $4$. Les \'el\'ements de $G$ d'ordre $2$ sont donc dans
$H$, qui n'en poss\`ede qu'un: $a^2$, qu'on note $-1$.

\hskip 5mm Le centre $Z(G)$ est d'ordre diff\'erent de $1$ car $G$ est un
$2$-groupe et diff\'erent de $8$ car $G$ est non ab\'elien. Il n'est pas non plus
d'ordre $4$ car alors on aurait $G=Z(G) \cup x Z(G)$ pour un $x\in G\setminus
Z(G)$ mais alors $G$ serait ab\'elien. Le centre $Z(G)$ est donc d'ordre $2$.
D'apr\`es ce qui pr\'ec\`ede $Z(G)=\{1,-1\}$. 

\hskip 5mm Soit $b\in G\setminus H$. Alors $G$ est engendr\'e par
$a$ et $b$. D'autre part $b$ est d'ordre $4$ et $b^2$ d'ordre $2$ ce qui
entraine $b^2=-1$. La conjugaison par $b$ induit un automorphisme du sous-groupe
distingu\'e $<a>$; on a donc $bab^{-1} = a^{-1}$, le seul autre cas $bab^{-1} = a$
\'etant exclu car $G$ non ab\'elien. On obtient ensuite ais\'ement que si $ab=c$, 
on a $c^2=-1$ ($c^2=abab=aa^{-1}bb=b^2=-1$) et $ba=-ab=-c$, $bc=-cb=a$,
$ca=-ac=b$.
\fincorrection
\nocorrection 
\correction{002196}
(a) On a $\theta (g) (xH) = gx H$ $(g,x\in G)$. Le noyau de $\theta$
est l'intersection de tous les conjugu\'es $xHx^{-1}$ de $H$, c'est-\`a-dire,
d'apr\`es les th\'eor\`emes de Sylow, l'intersection de tous les $3$-Sylow de $G$.
Comme l'intersection de deux $3$-Sylow distincts est triviale, le noyau est $\not=
\{1\}$ si et seulement s'il n'existe qu'un seul $3$-Sylow, qui est alors
automatiquement distingu\'e dans $G$.

\hskip 5mm Si $H$ est non distingu\'e dans $G$, alors $\theta$ est injectif et
fournit un isomorphisme entre $G$ et un sous-groupe de $S_4$. Ce sous-groupe devant
\^etre d'ordre $12$ comme $G$, c'est n\'ecessairement $A_4$ (cf l'exercice \ref{ex:deb75}).
\smallskip

(b) Si $G$ n'est pas isomorphe \`a $A_4$, alors n\'ecessairement $H$  est
distingu\'e dans $G$ et c'est alors l'unique $3$-Sylow de $G$. Notons $1,a,a^2$ les
trois
\'el\'ement distincts du groupe cyclique $H$.

\hskip 5mm Supposons que $G$ contienne un \'el\'ement $b$ d'ordre $4$. On a
$b^4=a^3=1$. D'autre part, la conjugaison par $b$ laissant invariant le sous-groupe
distingu\'e $H=\hskip 2pt <a>$, l'\'el\'ement $bab^{-1}$ doit \^etre un
g\'en\'erateur de $<a>$, c'est-\`a-dire $a$ ou $a^{-1}$. Mais la premi\`ere
possibilit\'e est exclue car sinon $b$ serait dans le centre de $G$ et $G$ serait
ab\'elien (cf exercice \ref{ex:le23}). La seconde possibilit\'e existe bien: on prend
par exemple pour $G$ le produit semi direct $\Z/3\Z \times \hskip -6pt {\raise
1.4pt\hbox{${\scriptscriptstyle |}$}} \Z/4\Z$ o\`u l'action de $\Z/4\Z$ sur
$\Z/3\Z$ se fait \`a travers la surjection canonique $\Z/4\Z \rightarrow \Z/2\Z$,
c'est-\`a-dire, les classes de $0$ et $2$ modulo $4$ agissent comme l'identit\'e et
celles de $1$ et $3$ comme le passage \`a l'inverse.

\hskip 5mm Supposons au contraire qu'aucun \'el\'ement de $G\setminus H$ soit
d'ordre $4$. Les $2$-Sylow sont donc isomorphes au groupe de Klein $\Z/2\Z
\times \Z/2\Z$. De plus, deux quelconques $B$ et $B^\prime$ d'entre eux sont
forc\'ement d'intersection non triviale car sinon l'ensemble produit $B B^\prime$
(qui est en bijection avec $B\times B^\prime$ par $(b,b^\prime)\rightarrow
bb^\prime$) serait de cardinal $|B|\hskip 2pt |B^\prime| = 16>12$. Il y a donc
strictement moins de $3\times 3 =9$ \'el\'ements d'ordre $2$ dans $G$. Comme
$G\setminus H$ est de cardinal $9$, il existe dans $G$ un \'el\'ement $c$ 
d'ordre $\not=2$. Cet \'el\'ement ne pouvant non plus \^etre  
d'ordre $3$ ($H$ est le seul $3$-Sylow), ni d'ordre $4$ (par hypoth\`ese) est
d'ordre $6$. Le groupe $<c>$ est alors d'indice $2$ et donc distingu\'e dans $G$.
Comme $<c>$ est cyclique, il ne poss\`ede qu'un seul \'el\'ement d'ordre $2$. On
peut donc trouver dans un $2$-Sylow de $G$ un \'el\'ement $d\in G\setminus <c>$
d'ordre $2$. La conjugaison par $d$ induit un automorphisme de $<c>$ qui envoie $c$
sur un g\'en\'erateur de $<c>$, c'est-\`a-dire ou bien $c$ ou bien $c^{-1}$. Mais la
premi\`ere possibilit\'e est exclue car $G$ n'est pas ab\'elien. On a donc
$dcd^{-1}=c^{-1}$; le groupe $G$ est dans ce cas isomorphe au groupe di\'edral
$D_6$.
\smallskip

(c) Les groupes d'ordre $12$ sont 

- les groupes ab\'eliens: $\Z/3Z \times \Z/4\Z \simeq \Z/12\Z$ et $\Z/3 \times
\Z/2\Z \times \Z/2\Z \simeq \Z/6\Z \times \Z/2\Z$, et

- les groupes non ab\'eliens: $A_4$, $\Z/3\Z \times \hskip -6pt {\raise
1.4pt\hbox{${\scriptscriptstyle |}$}} \Z/4\Z$ (pour l'action donn\'ee ci-dessus)
et $D_6$.  
\fincorrection
\nocorrection 
\correction{002198} Le groupe $P$ est un $p$-sous groupe maximal de $G$ et donc
aussi de $HP$ puisque $P\subset HP$ (noter que $HP$ est un sous-groupe car $H$ est
suppos\'e distingu\'e dans $G$); $P$ est donc un $p$-Sylow de $HP$. Si $|P|=p^n$,
alors $|HP|=p^n s$ avec $p$ ne divisant pas $s$. On peut aussi \'ecrire $|H|=p^{m}
r$ avec $p$ ne divisant pas $r$; on a alors n\'ecessairement $m\leq n$ et $s$
multiple de $r$. On a aussi $HP/H
\simeq P/(H\cap P)$ ce qui donne $|H\cap P| = |P||H|/|HP| = p^m (r/s)$. On obtient donc
que $s=r$ et que $H\cap P$ est un $p$-Sylow du groupe $H$.
\smallskip

On a aussi $|G|=p^n t$ avec $p$ ne divisant pas $t$ et $t$ multiple de $s$. On en
d\'eduit $|G/H| = p^{n-m} (t/r)$. Comme $t/r$ est un entier non divisible par $p$
et que $HP/H$ est un sous-groupe de $G/H$ d'ordre $|HP/H| = p^{n-m}$, le groupe
$HP/H$ est un $p$-Sylow de $G/H$.
  
\fincorrection
\correction{002199} 
D'apr\`es les th\'eor\`emes de Sylow, le nombre de $5$-Sylow d'un groupe d'ordre
$200=5^2.2^3$ est $\equiv 1\ [\hbox{\rm mod}\ 5]$ et divise $8$. Ce ne peut \^etre
que $1$. L'unique $5$-Sylow est n\'ecessairement distingu\'e puisque ses 
conjugu\'es sont des $5$-Sylow et coincident donc avec lui. Le groupe ne peut pas
\^etre simple.
\fincorrection
\correction{002200} 
Les $p$-Sylow de $S_p$ sont d'ordre $p$ puisque $p$, \'etant premier, ne divise pas $p!/p=(p-1)!\ $.
Chaque $p$-Sylow est donc cyclique d'ordre $p$ et contient $p-1$ \'el\'ements d'ordre $p$. Les
\'el\'ements d'ordre $p$ de $S_p$ sont les $p$-cycles; il y en a $(p-1)!\ $. Il y a donc $(p-2)!$ $p$-Sylow. 
(On retrouve le th\'eor\`eme de Wilson: $(p-2)! \equiv 1 \ [\hbox{\rm mod}\ p]$ (ou $(p-1)! +1\equiv 0 \ [\hbox{\rm mod}\ p]$) si $p$ est premier).
\fincorrection
\nocorrection
\correction{002202} 
\hskip 5mm Le groupe altern\'e $A_5$ est d'ordre $60=2^2.3.5$.

\hskip 5mm Les  $5$-Sylow sont d'ordre $5$, donc cycliques; chacun est engendr\'e par un $5$-cycle et contient  $4$ $5$-cycles. Les  $5$-Sylow sont deux \`a deux d'intersection r\'eduite \`a $\{1\}$. Comme il y a $24$ $5$-cycles dans $A_5$, il y a $6$ $5$-Sylow.
(On peut aussi utiliser les th\'eor\`emes de Sylow: Le nombre de $5$-Sylow est $\equiv 1\ [\hbox{\rm mod}\ 5]$ et divise $12$; c'est donc $1$ ou $6$. Comme ce ne peut \^etre $1$ (car il y aurait alors un unique $5$-Sylow qui serait distingu\'e, ce qui est impossible car $A_5$ est simple), 
%(on peut voir aussi plus simplement que le sous-groupe 
%engendr\'e par un $5$-cycle est un $5$-Sylow et qu'il n'est pas distingu\'e), 
c'est $6$.)

\hskip 5mm Les  $3$-Sylow sont d'ordre $3$, donc cycliques; chacun est engendr\'e par un $3$-cycle et contient  $2$ $3$-cycles. Les  $3$-Sylow sont deux \`a deux d'intersection r\'eduite \`a $\{1\}$. Comme il y a $20$ $3$-cycles dans $A_5$, il y a $10$ $3$-Sylow.
(Par les th\'eor\`emes de Sylow: le nombre de $3$-Sylow est $\equiv 1\ [\hbox{\rm mod}\ 3]$ et divise $20$; c'est donc $1$, $4$ ou $10$.   Comme ci-dessus, ce ne peut \^etre $1$. Si c'etait $4$, la conjugaison de $A_5$ sur ces $3$-Sylow induirait un morphisme $A_5\rightarrow S_4$ non trivial (puisque cette action par conjugaison est transitive) et donc injectif (puisque le noyau, distingu\'e, est forc\'ement trivial). Or l'ordre de $A_5$ ne divise pas celui de $S_4$. Il y a donc 10 $3$-Sylow.)

\hskip 5mm Les  $2$-Sylow sont d'ordre $4$, donc commutatifs. Comme il n'y a pas d'\'el\'ement d'ordre $4$ dans $A_5$, chaque $2$-Sylow est isomorphe au groupe $\Zz/2\Zz \times \Zz/2\Zz$; il est engendr\'e par deux produits de deux transpositions qui commutent et contient $3$ \'el\'ements d'ordre $2$. On voit ensuite que ces trois \'el\'ements d'ordre $2$ sont les $3$ produits de deux transpositions qui commutent qu'on peut former avec quatre \'el\'ements de $\{1,\ldots,5\}$. On en d\'eduit que les  $2$-Sylow sont deux \`a deux d'intersection r\'eduite \`a $\{1\}$. Il y a $15$ \'el\'ements d'ordre $2$ dans $A_5$ et il y a $5$ $2$-Sylow.


\hskip 5mm Tout \'el\'ement de $A_5$ est d'ordre $1$, $2$, $3$ ou $5$ et est donc contenu dans un $p$-Sylow.
On a bien $6.4+10.2 + 5.3+1 = 60$.
\fincorrection
\correction{002203}
(a) Le nombre de $5$-Sylow dans un groupe $G$ d'ordre $60=2^2.3.5$ est $\equiv 1\
[\hbox{\rm mod}\ 5]$ et divise $12$. Comme $G$ est suppos\'e simple, ce ne peut
\^etre $1$; il y a donc $6$ $5$-Sylow. Le morphisme $\alpha: G\rightarrow S_6$
correspondant \`a l'action de $G$ par conjugaison sur les $5$-Sylow (une fois une
num\'erotation des $5$-Sylow de $G$ choisie) est forc\'ement injectif puisque son
noyau, \'etant un sous-groupe distingu\'e diff\'erent de $G$ (d'apr\`es
les th\'eor\`emes de Sylow, $G$ agit transitivement sur les $5$-Sylow), est
n\'ecessairement trivial. Consid\'erons ensuite le groupe $\alpha^{-1}(A_6)$. C'est
un sous-groupe distingu\'e de $G$ (comme image r\'eciproque par un morphisme du
sous-groupe distingu\'e $A_6$ de $S_6$). Si $\alpha^{-1}(A_6)=\{1\}$ alors, pour
tout $g\in G$, comme $\alpha(g^2)=\alpha(g)^2\in A_6$, on aurait $g^2=1$ et donc
$G$ ab\'elien, ce qui est absurde. On a donc $\alpha^{-1}(A_6)=G$, c'est-\`a-dire,
$\alpha (G)=H\subset A_6$. 
\smallskip

(b) Notons $\varphi:A_6 \rightarrow S_6$ le morphisme
correspondant \`a l'action de $A_6$ par translation \`a gauche sur $A_6/.H$ (une
fois une num\'erotation des \'el\'ements de $A_6/.H$ choisie). En utilisant la
simplicit\'e de $A_6$, on montre comme ci-dessus que $\varphi$ est injectif et que
$\varphi(A_6)\subset A_6$. Il en d\'ecoule que $\varphi$ est un isomorphisme
entre $A_6$ et $\varphi(A_6)= A_6$.
\smallskip

(c) Un \'el\'ement $x\in A_6$ fixe la classe neutre $H$ si et seulement si $x\in H$.
On obtient que $H$ est isomorphe, {\it via} $\varphi$, au fixateur d'un entier,
disons $6$, dans l'action de $A_6$ sur $\{1,\ldots,6\}$, c'est-\`a-dire, \`a $A_6
\cap S_5 = A_5$.
\fincorrection
\nocorrection
\nocorrection
\correction{002206}
Le nombre de $q$-Sylow d'un groupe $G$ d'ordre $p^2q$ est $\equiv 1\
[\hbox{\rm mod}\ q]$ et divise $p^2$. Ce ne peut \^etre ni $p$ ni $p^2$ car $p^2-1$
est suppos\'e non divisible par $q$; c'est donc $1$. De m\^eme le nombre de
$p$-Sylow est $\equiv 1\ [\hbox{\rm mod}\ p]$ et divise $q$ et ce ne peut \^etre
$q$ car $q-1$ est suppos\'e non divisible par $p$; c'est donc $1$. Ainsi il y a 
un unique $p$-Sylow $P$ d'ordre $p^2$, et donc ab\'elien, et un unique $q$-Sylow
$Q$ d'ordre $q$, et donc cyclique, tous deux n\'ecessairement distingu\'es. Il en
r\'esulte que tout \'el\'ement $x\in P$ commute avec tout \'el\'ement $y\in Q$:
en effet le commutateur $xyx^{-1}y^{-1} = (xyx^{-1})y^{-1}= x(yx^{-1}y^{-1})$ est
dans l'intersection $P\cap Q$ qui est le groupe trivial. Cela montre que le groupe
$PQ$ est ab\'elien; il est isomorphe au produit direct $P\times Q$ et est donc de
cardinal $|P|\hskip 2pt |Q| = p^2q=|G|$. D'o\`u finalement $G=PQ$ est ab\'elien.
\fincorrection
\correction{002207}
Soit $G$ un groupe d'ordre $p^2q$ qu'on suppose simple. On distingue deux cas:
\smallskip

\underbar{1er cas}: $p>q$. Le nombre de $p$-Sylow de $G$ est $\equiv 1\ [\hbox{\rm
mod}\ p]$ et divise $q$. Comme $G$ est simple, ce ne peut \^etre $1$ (car sinon
l'unique $p$-Sylow serait distingu\'e). Il y a donc $q$ $p$-Sylow d'ordre $p^2$,
lesquels sont conjugu\'es. L'action par conjugaison de $G$ sur ces $q$ $p$-Sylow
d\'efinit un morphisme $G \rightarrow S_q$ non trivial (car l'action est
transitive) et   donc injectif puisque le noyau, distingu\'e et $\not=G$, est
forc\'ement trivial. On en d\'eduit que $p^2q$ divise $q!$ et donc $p$ divise un
entier entre $1$ et
$q-1$, ce qui contredit l'hypoth\`ese $p>q$.
\smallskip

\underbar{2\`eme cas}: $p<q$. Le nombre de $q$-Sylow de $G$ est $\equiv 1\
[\hbox{\rm mod}\ q]$ et divise $p^2$. Comme ci-dessus, $G$ \'etant simple, ce ne
peut \^etre $1$. Ce ne peut-\^etre ni $p$ ni $p^2$. En effet, dans le cas
contraire, $p$ serait $\equiv \pm 1\ [\hbox{\rm mod}\ q]$ et donc $p\geq q-1$. Comme
$p<q$, la seule possibilit\'e est $p=q-1$ et donc $p=2$ et $q=3$. Dans ce dernier
cas, il y a $4$ $3$-Sylow d'ordre $3$ qui contiennent $8$ \'el\'ements d'ordre
$3$. Ne reste de la place que pour un seul $2$-Sylow qui devrait \^etre distingu\'e. Ce dernier cas
n'est donc lui non plus pas possible. 
\smallskip

Conclusion: il n'existe pas de groupe $G$ simple d'ordre $p^2q$.
\fincorrection
\correction{002208}
(a) Le nombre de $19$-Sylow de $G$ est $\equiv 1\
[\hbox{\rm mod}\ 19]$ et divise $21$; ce ne peut \^etre que $1$. Le groupe $G$ a
donc un unique $19$-Sylow $P$ qui est distingu\'e.
\smallskip

(b) Comme $P$ est distingu\'e dans $G$, $N=PQ$ est un sous-groupe de $G$. De
$P\cap Q=\{1\}$, on d\'eduit que $PQ/P \simeq Q$ et donc que $PQ$ est d'ordre
$7.19=133$. D'apr\`es l'exercice 15, le groupe $N$ est isomorphe au produit direct
$\Z/19\Z \times \Z/7\Z$, lequel est isomorphe au groupe cyclique $\Z/133\Z$ par le
lemme chinois.
\smallskip

(c) Le nombre de $7$-Sylow de $G$ est $\equiv 1\
[\hbox{\rm mod}\ 7]$ et divise $57$. Les seules possibilit\'es sont $1$ et $57$. Or
ce n'est pas $1$ non plus car on suppose que $Q$ n'est pas distingu\'e. Le groupe
$G$ admet donc $57$ $7$-Sylow, et donc $57$ sous-groupes cycliques d'ordre $133$
par la question pr\'ec\'edente. Ces $57$ groupes d'ordre $133$ sont bien distincts
car deux
$7$-Sylow distincts engendrent avec $P$ deux groupes cycliques d'ordre $133$
distincts puisque le $7$-Sylow est l'unique sous-groupe d'ordre $7$ du groupe
cyclique. Par cons\'equent leurs ensembles de g\'en\'erateurs sont deux \`a deux
disjoints. On obtient ainsi $57 \times \phi(133)= 57\times 6\times 18$ \'el\'ements
d'ordre
$133$ dans $G$ ($\phi$ d\'esigne ici la fonction indicatrice d'Euler), ce qui est
manifestement absurde. On peut donc conclure que $Q$ est distingu\'e dans $G$ et
que l'unique sous-groupe cyclique $N=PQ$ d'ordre $133$ l'est aussi.
\smallskip

(d) Comme $N$ est distingu\'e dans $G$, $NR$ est un sous-groupe de $G$. De
$N\cap R=\{1\}$, on d\'eduit que $NR/N \simeq R$ et donc que $NR$ est d'ordre
$133.3=399$. Ainsi $G=NR$ et l'isomorphisme pr\'ec\'edent $G/N \simeq R$ montre que
l'inclusion $R\rightarrow G$ est une section de la suite exacte $1\rightarrow
N\rightarrow G
\rightarrow R \rightarrow 1$. Le groupe $G$ est donc
isomorphe au produit semi-direct du groupe cyclique $N$ d'ordre $133$ par le groupe
cyclique $R$ d'ordre $3$. 
\fincorrection
\nocorrection
\nocorrection
\correction{007740}
\begin{enumerate}
 \item On note $O$ le centre de gravité de $\mathcal{P}_n$.
  Dans $D_{n}$, toutes les isométries, qui sont affines, conservent le centre de gravité. Il n'y a donc aucune translation ni symétrie glissée. Il y a les $n$ rotations de centre $O$ et d'angle $2k\pi/n$ avec $k = 0, \cdots, n-1$. Si $n$ est pair, il y a les $n/2$ symétries orthogonales d'axe $(OS)$ pour chaque sommet $S$ et les $n/2$ symétries d'axe médiateur d'un des $n$ côtés. Si $n$ est impair, il y a les $n$ symétries d'axe médiateur d'un des $n$ côtés.
 \item On considère l'action de $D_{n}$ sur l'ensemble fini des $n$ sommets de $\mathcal{P}_n$.
 \`A l'aide des rotations de centre $O$ et d'angle $2k\pi/n$, on montre que l'action est transitive.
 Le stabilisateur d'un sommet $S$ est le sous-groupe d'ordre $2$ engendré par la symétrie orthogonale d'axe $(OS)$ où $O$ est le centre du polygone. Par la seconde formule des classes, $D_{n}$ est d'ordre $2n$.
 La liste précédente est donc complète.
 \item Comme $n$ est impair, les $2$-Sylow sont d'ordre $2$. Ce sont les sous-groupes engendrés par les symétries orthogonales.
 Si $s$ est une symétrie d'axe $d$ et $r$ la rotation de centre $O$ et d'angle $2k\pi/n$, 
 $rsr^{-1}$ est la symétrie d'axe $r(d)$.
 Comme $n$ est impair, le groupe engendré par $r$ agit transitivement sur les axes de symétrie. Toutes les symétries sont donc conjuguées. 
 \item Les $2$-Sylow sont d'ordre $4$ dans $D_{6}$ d'ordre $2^2\times 3$.
 Le groupe engendré par deux symétries d'axe orthogonal (qui commutent) est d'ordre $4$. En conjuguant par la rotation de centre $O$ et d'angle $2\pi/6$, on obtient un autre $2$-Sylow.
 Les $2$-Sylow ne sont donc pas distingués. Le nombre de $2$-Sylow est congru à $1$ modulo $2$ et divise $12$ et n'est pas $1$.
 Il y a donc trois $2$-Sylow.
 \item La conjuguée, par une rotation d'angle $2k\pi/6$ ou par une symétrie, d'une symétrie d'axe médiateur d'un coté est une symétrie d'axe médiateur d'un coté. 
 Les symétries d'axe médiateur d'un coté et les symétries d'axe passant par un sommet (qui diffèrent car $6$ est pair) ne sont donc pas conjuguées dans $D_{6}$.
 \item Le sous-groupe engendré par une rotation d'angle $4\pi/6$ est d'ordre $3$ et c'est un $3$-Sylow.  
\end{enumerate}
\fincorrection
\correction{007741}
Les ordres des éléments de $G$ sont $3, 11$ ou $33$. Une application directe du
théorème de Sylow montre qu’on a un seul groupe d’ordre $11$ et un seul groupe d’ordre
$3$. Les éléments d’ordre $3$ et $11$ sont contenus dans ces deux groupes. On a au plus
$$1 + (3 - 1) + (11 - 1) = 1 + 2 + 10 = 13$$ éléments d’ordre 1, 3 ou $11$. Il existe donc un
élément d’ordre $33$ dans $G$ qui est donc cyclique isomorphe à $\mathbf{Z}/33\mathbf{Z}$.
\fincorrection
\nocorrection
\nocorrection
\nocorrection
\nocorrection
\nocorrection
\nocorrection
\nocorrection
\nocorrection
\nocorrection
\nocorrection
\nocorrection
\nocorrection
\nocorrection
\nocorrection
\nocorrection
\nocorrection
\nocorrection
\nocorrection
\nocorrection
\nocorrection
\nocorrection
\nocorrection
\nocorrection
\nocorrection
\nocorrection
\nocorrection
\nocorrection
\nocorrection
\nocorrection
\nocorrection
\nocorrection
\nocorrection
\nocorrection
\nocorrection
\nocorrection
\nocorrection
\nocorrection
\nocorrection
\nocorrection
\nocorrection
\nocorrection
\nocorrection
\nocorrection
\nocorrection
\nocorrection
\nocorrection
\nocorrection
\nocorrection
\nocorrection
\nocorrection
\nocorrection
\nocorrection
\nocorrection
\nocorrection
\nocorrection
\nocorrection
\nocorrection
\nocorrection
\nocorrection
\nocorrection
\nocorrection
\nocorrection
\nocorrection
\nocorrection
\nocorrection
\nocorrection
\nocorrection
\nocorrection
\nocorrection
\nocorrection
\nocorrection
\nocorrection
\nocorrection
\nocorrection
\nocorrection
\nocorrection
\nocorrection
\nocorrection
\nocorrection
\nocorrection
\nocorrection
\nocorrection
\nocorrection
\nocorrection
\correction{007738}
\begin{enumerate}
 \item Le groupe $\mathcal{A}_6$ est simple et n'a donc pas de sous-groupes distingués. De façon plus élémentaire,
 $$(1,2)(3,4)\circ (3,4)(1,5)=(1,2)(1,5)=(1,5,2)$$
 et l'ensemble proposé n'est donc pas un sous-groupe.
 \item L'intersection de deux espaces vectoriels $F$ et $G$ de dimension $3$ dans un espace vectoriel $E$ de dimension $4$ est un sous-espace vectoriel de dimension $\dim F+\dim G-\dim(F+G)=6-\dim(F+G)$ avec $3\leq \dim F+G\leq \min(\dim E, \dim F+\dim G)=4$ donc soit $3$, soit $2$. 
 L'intersection de deux plans projectifs de $\mathbb{P}^3$ est donc soit un plan projectif soit une droite projective.
 
 
 L'intersection de deux espaces vectoriels $F$ et $G$ de dimension $3$ dans un espace vectoriel $E$ de dimension $5$ est un sous-espace vectoriel de dimension $\dim F+\dim G-\dim(F+G)=6-\dim(F+G)$ avec $3\leq dim F+G
 \leq \min(\dim E, \dim F+\dim G)=5$ donc soit $3$, soit $2$, soit $1$. 
 L'intersection de deux plans projectifs de $\mathbb{P}^3$ est donc soit un plan projectif, soit une droite projective, soit un point.
 \item Comme une homographie est caractérisée par l'image d'un repère projectif, il suffit de prendre cinq points $A$, $B$, $C$, $D$ $E$ deux à deux distincts : il n'existe aucune homographie qui fixe $A$, $B$ et $C$ et qui échange $D$ et $E$ par exemple.
 
 Sinon,
 On choisit un repère projectif et des coordonnées homogènes.
 Les points de coordonnées $A[1:0], B[0,1], C[1,1]$ et $D[2,1]$ ont pour birapport $2$ alors que ceux de coordonnées $A[1:0], B[0,1], C[1,1]$ et $E[3,1]$ ont pour birapport $3$.
 Par conséquent, pour tout point $M$ les quintuplets 
 $$(A,B,C,D,M) \textrm{ et } (A,B,C,E,M)$$
 ne sont images l'un de l'autre dans aucune homographie.  
\end{enumerate}
\fincorrection
\nocorrection
\nocorrection
\correction{007744}
\begin{enumerate}
    
    \item ~
    
    \begin{center}
        \includegraphics[scale=0.5]{images/img-mour-496}
    \end{center}
    \item Non, il faut au moins avoir l'image d'un repère projectif. 
\end{enumerate}
\fincorrection
\correction{007745}
\begin{enumerate}
    \item Comme $F$ admet une droite $d$ de points fixes, $f$ fixe toutes les droites d'un plan. C'est donc en restriction à ce plan, une homothétie. Quitte à diviser par ce rapport non nul, on peut supposer que $f$ est l'identité sur ce plan.
    
    \item L'application $f$ est alors soit une dilatation, soit une transvection. Dans le premier cas, avec un repère adapté, tout point de coordonnées $(X,Y,Z)$ et son image de coordonnées $X,Y,\lambda Z)$ sont coplanaires avec le point de coordonnées $(0,0,1)$.
    Dans le second cas tout point de coordonnées $(X,Y,Z)$ et son image de coordonnées $(X,Y+Z,Z)$ sont coplanaires avec le point de coordonnées $(0,1,0)$.
    On en déduit donc l'existence d'un centre $O$.
    
    \item On choisit un repère projectif de $P^2$ composé de $M_1$ et $M_2$ sur $d$, puis $M_3=O$ et comme point unitaire $M_4=A$. On obtient un repère tel que $d$ ait pour équation $Z=0$ et $O$ pour coordonnées homogènes $[0:0:1]$ et $A$ pour coordonnées homogènes $[1:1:1]$. Notons $[x:x:z]$ des coordonées homogènes de $A'$ aligné avec $O$ et $A$ avec $x\not=0$ car $A'\not=O$ et  $z\not=0$ car $A'\not\in d$.
    On peut donc normaliser avec $A'[1:1:z]$. Comme $d$ doit être laissée fixe par l'homographie cherchée $P(f)$, quitte à normaliser, on peut supposer que $f$ est l'identité sur $vect(e_1,e_2)$. Comme $P(f)(A)=A'$, $f(e_1+e_2+e_3)=e_1+e_2+f(e_3)$ est proportionel à $e_1+e_2+ze_3$. Donc, $f(e_3)=ze_3$.
    Réciproquement, cette application $f$ satisfait les conditions requises (voir le cas de la dilatation).
    
    \item \begin{figure}[!ht] 
        \begin{center}\includegraphics[scale=0.6]{images/img-mour-497-4}\end{center}
        \caption{}
    \end{figure}
    Notons $\phi$ l'homographie de droite fixe $d$ de centre $O$ qui envoie $P$ sur $P'$. On vérifie successivement qu'elle envoie $P$ sur $P'$, $Q$ sur $Q'$, $P'$ sur $P$ et enfin $Q'$ sur $Q$. Elle coïncide avec $H$ sur un repère projectif donc partout. 
\end{enumerate}
\fincorrection
\correction{007746}
\begin{enumerate}
 \item Dire que $F$ et $F'$ ont les mêmes points fixes revient dire que $f$ et $f'$ ont les mêmes directions propres $d_1=vect(v_1)$ et $d_2=vect(v_2)$. les matrices de $f$ et $f'$ dans la base $(v_1,v_2)$ sont diagonales et donc commutent, ce qui implique que $F$ et $F'$ commutent.
 
 \item Comme $dim(E)=2$, un élément de $Gl(E)$ qui a trois directions propres distinctes est une homothétie et donc induit l'identité sur $P(E)$.
 
 \item On utilise la commutativité :
 $F(F'(A))=F \circ F'(A)=F' \circ F(A)=F'(A)$ donc $F'(A)$ est un point fixe de $F$. De même $F'(B)$ est un point fixe de $F$, $F(A')$ est un point fixe de $F'$ et $F(B')$ est un point fixe de $F'$. On en déduit $\{F'(A),F'(B)\} = \{ A,B\}$ et $\{F(A'),F(B')\} = \{A',B'\}$ puisque $F$ et $F'$ n'ont chacune que deux points fixes.
 
 \item $F$ n'a que deux points fixes (car $F\neq Id_\Delta$ puisque  $F^2\neq Id_\Delta$), comme $A'$ et $B'$ sont distincts (par hypothèse sur $F'$) et sont fixés par $F$ on a le résultat souhaité.
 
 \item $F^2$ ne possède que deux points fixes, puisque $F^2\neq Id_\Delta$.
 Sous l'hypothèse, $A'$ et $B'$ sont des points fixes de $F^2$, distincts par hypothèse sur $F'$, donc l'ensemble des points fixes de $F^2$ est $\{A',B'\}$. D'autre part $A$ et $B$ sont aussi fixés par $F^2$ puisque fixés par $F$ et $A \neq B$ par hypothèse sur $F$, donc l'ensemble des points fixes de $F^2$ est $\{A,B\}$. Enfin, par transitivité de l'égalité, $\{A',B'\}=\{A,B\}$. 
 On a ainsi obtenu une absurdité, puisqu'on a supposé que $F$ échangeait $A'$ et $B'$ et qu'on a obtenu qu'elle les fixait.
 
 \item On a donc toujours la situation envisagée en 4., et $F$ et $F'$ ont les mêmes points fixes.
\end{enumerate}
\fincorrection
\nocorrection
\nocorrection
\nocorrection
\nocorrection
\nocorrection
\nocorrection
\nocorrection
\nocorrection
\nocorrection
\correction{007756}
Soit $6$ points $A,B,\cdots F$ du plan $\R^2$ tels que $ABCDEF$ soit un hexagone régulier.

\includegraphics[scale=0.4]{images/img-mour-508-2}

Etant données les images $A'=h(A)$, $B'=h(B)$, $D'=h(D)$ et $E'=h(E)$ par une homographie $h$ de $P^2(\R)$ dans lui-même, pour construire à la règle les images des autres points, on commence par déterminer l'ancienne droite à l'infini, en déterminant les points d'intersection d'image de couple de droites parallèles.
\fincorrection
\nocorrection
\nocorrection
\nocorrection
\nocorrection
\nocorrection
\nocorrection
\nocorrection
\nocorrection
\nocorrection
\nocorrection
\nocorrection
\nocorrection
\nocorrection
\nocorrection
\nocorrection
\nocorrection
\nocorrection
\nocorrection
\nocorrection
\nocorrection
\nocorrection
\nocorrection
\nocorrection
\nocorrection
\nocorrection
\nocorrection
\nocorrection
\nocorrection
\nocorrection
\nocorrection
\nocorrection
\nocorrection
\nocorrection
\nocorrection
\nocorrection
\nocorrection
\nocorrection
\nocorrection
\nocorrection
\nocorrection
\nocorrection
\nocorrection
\nocorrection
\nocorrection
\nocorrection
\nocorrection
\nocorrection
\nocorrection
\nocorrection
\nocorrection
\nocorrection
\nocorrection
\nocorrection
\nocorrection
\nocorrection
\nocorrection
\nocorrection
\nocorrection
\nocorrection
\nocorrection
\nocorrection
\nocorrection
\nocorrection
\nocorrection
\nocorrection
\nocorrection
\nocorrection
\nocorrection
\nocorrection
\nocorrection
\nocorrection
\nocorrection
\nocorrection
\nocorrection
\nocorrection
\nocorrection
\nocorrection
\nocorrection
\nocorrection
\nocorrection
\nocorrection
\nocorrection
\nocorrection
\nocorrection
\nocorrection
\nocorrection
\nocorrection
\nocorrection
\nocorrection
\nocorrection
\correction{007723}
Le théorème de Sylow donne que ce nombre $N$ est congru à 1 modulo $5$ et divise $24$. C'est donc $1$ ou $6$.
Comme $(12)(12345)(12)=(21345)$ n'est pas dans le $5$-Sylow (d'ordre $5$) $<(12345)>$ il y a un $5$-Sylow non distingué.
 Ainsi, $N=6$.
\fincorrection
\nocorrection
\correction{007725}
\begin{enumerate}
 \item Comme le seul élément non nul de $\mathbb{F}^2$ est $1$, $SL(4,\mathbb{F}^2)=GL(4,\mathbb{F}_2)$. Le centre de $GL(4,\mathbb{F}_2)$ composé des homothéties inversibles est donc réduit à l'identité ainsi que celui de $SL(4,\mathbb{F}_2)$. Ainsi, $|PSL(4,\mathbb{F}^2)|=|GL(4,\mathbb{F}^2)|=(2^4-1)(2^4-2)(2^4-2^2)(2^4-2^3)=15\times 14\times 12\times 8=7\times 5\times 3^2\times 2^6$.
 
 \item Ces deux matrices sont d'ordre $2$.
 Elles ne sont pas conjuguées car la dimension des espaces propres de valeurs propres $1$ n'est pas la même pour les deux.
 
 \item Tous les éléments de $\mathbb{F}^\times_4$ sont des racines $3$-ième d'unité par le théorème de Lagrange. Par conséquent toutes les homothéties inversibles sont dans le centre de $SL(3,\mathbb{F}_4)$.
 Donc, 
 $|PSL(3,\mathbb{F}^4)=GL(3,\mathbb{F}_4)/(3\times 3)=(4^3-1)(4^3-4)(4^3-4^2)/9
 =63\times 60\times 48/9=7\times 5\times 3^2\times 2^6$.
 
 \item 
\begin{enumerate}
 \item $P(f^3)=F^3=F$ car $F$ est une involution.
 $g=f^3$ n'est donc pas l'identité, mais $g^2=f^6=(f^2)^3=Id$ car $f^2$ est dans le centre de $SL(3,\mathbb{F}_4)$ qui est d'ordre $3$. 
 \item Le corps $\mathbb{F}_4$ est de caractéristique $2$. $(g-Id)^2=g^2-2g+Id=0$. Par conséquent, $\dim Im(g-Id)\leq\dim Ker(g-Id)$ et $\dim Im(g-Id)+\dim Ker(g-Id)=3$. Comme $g\not=Id$, $\dim Ker(g-Id)=2$.
 \item Ainsi $g$ est un élément de $SL(3,\mathbb{F}_4)$ qui admet un plan de points fixes. C'est donc une transvection.
\end{enumerate}
\item Dans $PSL(4,\mathbb{F}_2)$, il y a au moins deux classes de conjugaison d'éléments d'ordre $2$ par la question $1$. Dans $SL(3,\mathbb{F}_4)$ comme toutes les transvections sont conjuguées, la question 4 permet d'affirmer qu'il n'y a qu'une classe de conjugaison d'éléments d'ordre $2$. Les deux groupes ne sont donc pas isomorphes.  
\end{enumerate}
\fincorrection
\nocorrection
\nocorrection
\nocorrection
\nocorrection
\nocorrection
\nocorrection
\nocorrection
\nocorrection
\nocorrection
\correction{002249}
Cours... Non, les rôles des deux opérations ne sont pas interchangeables,
puisque l'une est distributive sur l'autre.
\fincorrection
\correction{002250}
\begin{enumerate}
  \item une seule solution $x=a^{-1}(c-b)$
  \item
    pas de solution, et deux solutions. Attention, dans $\Zz/10\Zz$, on ne
    peut pas inverser $2$. Ecrire $2x=3+10k$ pour obtenir que $2|3$, et
    $2x=6+10k$ pour simplifier par $2$... dans $\Rr$.
  \end{enumerate}
\fincorrection
\correction{002251}
  \begin{enumerate}
  \item
    Ecrire $(0+a)a=a.a$ d'une part ($0$ est neutre pour $+$) et
    $(0+a).a=0.a+a.a$ (distributivité).
  \item 
    $(-1).a+a=(-1+1).a=0.a=0$ (distributivité, puis question précédente)
  \item
    Si $|A|=1$, $1=0$. Si $1=0$, $\forall a\in A, a=1.a=0.a=0$, donc
    $A=\{0\}$. 
  \end{enumerate}
\fincorrection
\correction{002252}
  \begin{enumerate}
  \item 
    Si $xy\in A^{\times}$, soit $z\in A, (xy)z=1$. Alors $x(yz)=1$ et
    $(zx)y=1$ donc $x$ et $y$ sont inversibles.

  \item
     Soit $x\in A^{\times}$, et $y\in A, xy=0$. Alors $x^{-1}xy=y=0$.
     Donc $x$ n'est pas diviseur de $0$.
  \end{enumerate}

\fincorrection
\correction{002253}
 Soit $a\in A\setminus\{0\}$. Soit $\phi_{a}:A\to A, x\mapsto ax$. Si
  $\phi_{a}(x)=\phi_{a}(y)$, alors $ax=ay$. 
Mais $ax=ay $ ssi $a(x-y)=0$, or $a\neq 0$ et $A$ est int\`egre, donc $x=y$.
%donc $a^{-1}ax=a^{-1}ay$ et $x=y$. 
Ainsi $\phi_{a}$ est injective de $A$ dans $A$. Comme $A$ est
  fini, elle est donc aussi surjective~: $\exists x\in A, \phi_{a}(x)=1$.
\fincorrection
\correction{002254}
Ce sont tous des anneaux. Montrer que $A$ est stable par addition, par
passage à l'opposé, contient $0$, est stable par multiplication et
contient $1$. Le reste (associativité et distributivité) est automatique
puisqu'il s'agit des restrictions des opérations usuelles sur $\Cc$)
  \begin{enumerate}
  \item 
    $A$ est l'ensemble des nombres dont le développement décimal s'arrête
    (``nombre fini de chiffres après la virgule'').

    Stabilité par addition~: Soit $x=10^{-n}a$ et $y=10^{-m}b$. Supposons
    par exemple que $n\geq m$. Alors $x+y=10^{-n}(a+10^{n-m}b)$ et
    $a+10^{n-m}b\in\Zz$ donc $x+y\in A$. Les autres vérifications sont
    analogues.

    Ce n'est pas un corps~: $3$ n'est pas inversible, puisque si
    $3\cdot10^{-n}a=1$, alors $3a=10^{n}$ donc $3|10^{n}$ ce qui est
    impossible. Un élément est inversible ssi il est de la forme
    $10^{-n}2^{\alpha}5^{\beta}$, $\alpha,\beta\in\Nn$.

  \item
    Stabilité par addition~: Soit $x=\frac{a}{b}\in A$ et
    $y=\frac{c}{d}\in A$, avec $\pgcd(a,b)=\pgcd(c,d)=\pgcd(p,b)=\pgcd(p,d)=1$. Alors $x+y=\frac{ad+bc}{bd}$. 

    Ce n'est pas un corps~: $p$ n'est pas inversible. Un élément est
    inversible ssi ce n'est pas un multiple de $p$.

  \item
    N'est pas un corps~: $2$ n'est pas inversible. Les seuls éléments
    inversibles sont $1,-1,i,-i$. En effet, si $z\in A^{\times}$, alors
    $|z|\geq 1$ et $|z^{-1}|\geq 1$. Donc $|z|=1$ et $z\in\{\pm1,\pm
    i\}$. Réciproquement, ces éléments sont bien tous inversibles.
  \end{enumerate}
\fincorrection
\nocorrection
\nocorrection
\nocorrection
\nocorrection
\nocorrection
\nocorrection
\correction{002300}
  $1\in I+J$ donc $\exists(x,y)\in I\times J, 1=x+y$. En multipliant
  cette égalité par $x$, on obtient $x^{2}+xy=x$. On en déduit que $xy\in
  I$, donc $\forall p\in\Nn~; x^{p}y\in I^{p}$, et donc
  $\forall(p,q)\in\Nn^{2}, x^{p}y^{q}\in I^{p}$. Par symétrie, on a aussi
  $\forall(p,q)\in\Nn^{2}, x^{p}y^{q}\in J^{q}$.

  Soit maintenant $(m,n)\in\Nn^{2}$. Notons $N=2\sup(m,n)$. Alors
  $1=1^{N}=(x+y)^{N}=\sum_{p+q=N}C^{p}_{N}x^{p}y^{q}$. Comme~:
  $(p+q=2N)\Rightarrow (p\geq n \text{ ou }q\geq m)$, tous les termes de
  cette somme sont dans $I^{n}$ ou dans $J^{m}$, et donc $1\in I^{n}+J^{m}$
\fincorrection
\correction{002301}
  \begin{enumerate}
  \item 
    $3,5,7,11$ sont deux à deux premiers entre eux, donc la solution est
    unique modulo $1155=3\cdot5\cdot 7\cdot11$.
    
    \begin{align*}
      \begin{cases}        
        x\equiv 1 \mod 3\\
        x \equiv 3 \mod 5\\
        x \equiv 4 \mod 7\\
        x\equiv 2 \mod 11
      \end{cases}
      \Leftrightarrow
      \begin{cases}        
        x\equiv 13 \mod 15\\
        x \equiv 4 \mod 7\\
        x\equiv 2 \mod 11
      \end{cases}
      \Leftrightarrow&
      \begin{cases}        
        x\equiv 88 \mod 105\\
        x\equiv 2 \mod 11
      \end{cases}\\
      \Leftrightarrow&
      \begin{cases}        
        x\equiv 508 \mod 1155
      \end{cases}
    \end{align*}
    
    \item
      Un diviseur commun de $2001$ et $2002$ divise leur différence, et
      donc $\pgcd(2001,2002)=1$. De même, $\pgcd(2002,2003)=1$, et comme
      $2\!\!\!\not| 2001$, $\pgcd(2001,2003)=1$.

      $2001,2002,2003$ sont donc deux à deux premiers entre eux, et la
      solution est donc unique modulo $2001\cdot2002\cdot2003$.
    \begin{align*}
      \begin{cases}        
        x \equiv 997 \mod 2001\\
        x \equiv 998 \mod 2002\\
        x \equiv 999 \mod 2003
      \end{cases}
      \Leftrightarrow&
      \begin{cases}        
        x \equiv -1004 \mod 2001\\
        x \equiv -1004 \mod 2002\\
        x \equiv -1004 \mod 2003
      \end{cases}\\
      \Leftrightarrow&
        x \equiv -1004 \mod (2001\cdot2002\cdot2003)
    \end{align*}
  \end{enumerate}
\fincorrection
\correction{002302}
On a $72=8\cdot9$ et $\pgcd(8,9)=1$, donc
$\Zz_{72}\simeq\Zz_{8}\times\Zz_{9}$. De même,
$\Zz_{84}\simeq\Zz_{4}\times\Zz_{3}\times\Zz_{7}$,
$\Zz_{36}\simeq\Zz_{4}\times\Zz_{9}$ et
$\Zz_{168}\simeq\Zz_{8}\times\Zz_{3}\times\Zz_{7}$. Donc
$\Zz_{72}\times\Zz_{84}\simeq
\Zz_{8}\times\Zz_{9}\times\Zz_{4}\times\Zz_{3}\times\Zz_{7}\simeq
\Zz_{4}\times\Zz_{9}\times\Zz_{8}\times\Zz_{3}\times\Zz_{7}\simeq
\Zz_{36}\times\Zz_{128}$
\fincorrection
\correction{002303}
\begin{enumerate}
  \item 
  $11,31,61$ sont premiers donc $2$ à $2$ premiers entre eux. Ainsi
  $20^{15}\equiv 1[11\cdot31\cdot61]\Leftrightarrow 
  \begin{cases}
    20^{15}\equiv1[11]\\
    20^{15}\equiv1[31]\\
    20^{15}\equiv1[61]
  \end{cases}
  $

  \begin{itemize}
  \item 
    En utilisant le petit théorème de Fermat, on obtient que, modulo
    $11$~: $20^{15}\equiv 20^{5}\equiv -2^{5}\equiv 1[11]$.

  \item
    $(20^{15})^{2}=20^{30}\equiv 1[31]$. On en déduit que
    $20^{15}\equiv\pm1[31]$. Comme $31\not\equiv1[4]$, d'après le
    théorème de Wilson, $x^{2}=-1$ n'a pas de solution modulo $31$, et
    donc $20^{15}\equiv1[31]$. $20^{2}\equiv-3[31]$ est premier

  \item
    $20^{15}\equiv(9^{2})^{15}\equiv3^{60}\equiv1[61]$
  \end{itemize}
  
\item
  $1155=11\cdot7\cdot5\cdot3$. De plus (petit théorème de Fermat)
  $2^{6754}\equiv 2^{4}\equiv5[11]$. De même,
  $2^{6754}\equiv2^{4}\equiv2[7]$, $2^{6754}\equiv2^{2}\equiv-1[5]$, et
  $2^{6754}\equiv2^{0}\equiv1[3]$. Or 
  $$
  \begin{cases}
    a\equiv5[11]\\
    a\equiv2[ 7]\\
    a\equiv4[ 5]\\
    a\equiv1[ 3]
  \end{cases}
  \Leftrightarrow
  \begin{cases}
    a\equiv5[11]\\
    a\equiv2[ 7]\\
    a\equiv4[15]
  \end{cases}
  \Leftrightarrow
  \begin{cases}
    a\equiv5[11]\\
    a\equiv-26[105]
  \end{cases}
  \Leftrightarrow
    a\equiv709[1155]\\
  $$
  Donc le reste de la division de $2^{6754}$ par $1155$ est $709$.
\end{enumerate}

\fincorrection
\correction{002304}

$13$ est premier et $100=12\cdot8+4$ donc $10^{100}\equiv
10^{4}\equiv(-3)^{4}\equiv3\equiv-10[13]$. De même
$10^{100}\equiv10^{-8}\equiv2^{8}\equiv9\equiv-10[19]$. En utilisant le
lemme chinois, on en déduit que $10^{100}\equiv-10[247]$. Comme
$\pgcd(10,247)=1$, on peut simplifier cette expression par $10$ et on a
$10^{99}\equiv -1[247]$, et donc $247|10^{99}+1$.
\fincorrection
\correction{002305}

$C=A\times B$.
\begin{align*}
(a,b)\in (A\times B)^{\times}
\Leftrightarrow&
\exists(c,d)\in A\times B, \ (a,b)(c,d)=(1,1)\\
\Leftrightarrow&
\exists(c,d)\in A\times B, \ ac=1 \text{ et } bd=1\\
\Leftrightarrow&
a\in A^{\times} \text{ et } b\in B^{\times}\\
\end{align*}
 donc $(A\times B)^{\times}=A^{\times}\times B^{\times}$.

De même, on obtient que l'ensemble $\mathcal{D}_{A\times B}$ des
diviseurs de $0$ de $A\times B$ est
$$
\mathcal{D}_{A\times
  B}=\mathcal{D}_{A}\times B\cup A\times\mathcal{D}_{B}\cup
(A\setminus\{0\})\times\{0\}\cup\{0\}\times(B\setminus\{0\}).
$$
Enfin, pour les nilpotents $Nil(A\times B)=Nil(A)\times Nil(B)$.
\fincorrection
\correction{002306}
\begin{enumerate}
  \item 
En posant $y=x+1$, on a
$\Zz_{2}[x]/(x^{3}+x^{2}+x+1)=\{0,1,x,y,x^{2},y^{2},xy,xy+1\}$.
Les tables des opérations sont les suivantes (elles sont symétriques)~:
$$
\begin{array}{c|c|c|c|c|c|c|c|c|c|c|}
\oplus    &   0 &   1 &   x &   y & x^2  & y^2  & xy   & xy+1\\\hline
0         &   0 &   1 &   x &   y & x^2  & y^2  & xy   & xy+1\\\hline
1         &     &   0 &   y &   x & y^2  & x^2  & xy+1 & xy  \\\hline
x         &     &     &   0 &   1 & xy   & xy+1 & x^2  & y^2 \\\hline
y         &     &     &     &   0 & xy+1 & xy   & y^2  & x^2 \\\hline
x^2       &     &     &     &     & 0    & 1    & x    & y   \\\hline
y^2       &     &     &     &     &      & 0    & y    & x   \\\hline
xy        &     &     &     &     &      &      & 0    & 1   \\\hline
xy+1      &     &     &     &     &      &      &      & 0   \\\hline
\end{array}
$$

$$
\begin{array}{c|c|c|c|c|c|c|c|c|c|c|}
\otimes   &   0 &   1 &   x &   y & x^2  & y^2  & xy   & xy+1\\\hline
0         &   0 &   0 &   0 &   0 & 0    & 0    & 0    & 0   \\\hline
1         &     &   1 &   x &   y & x^2  & y^2  & xy   & xy+1\\\hline
x         &     &     & x^2 &  xy & xy+1 & y^2  & y    & 1   \\\hline
y         &     &     &     & y^2 &  y   & 0    & y^2  & xy  \\\hline
x^2       &     &     &     &     &  1   & y^2  & xy   & x   \\\hline
y^2       &     &     &     &     &      & 0    & 0    & y^2 \\\hline
xy        &     &     &     &     &      &      & y^2  & y   \\\hline
xy+1      &     &     &     &     &      &      &      & x^2 \\\hline
\end{array}
$$
 
Pour $\Zz[x]/(x^{2}-1)$, $(x-1)$ et $(x+1)$ sont deux idéaux étrangers,
et le lemme chinois nous donne $\Zz[x]/(x^{2}-1)\simeq
\Zz[x]/(x-1)\times\Zz[x]/(x+1)$. Or $\Zz[x]/(x+1)\simeq\Zz$ et
$\Zz[x]/(x-1)\simeq\Zz$ donc $\Zz[x]/(x^{2}-1)\simeq\Zz\times\Zz$.


La factorisation de $(x^{8}-1)$ sur $\Qq$ est
$(x^{8}-1)=(x-1)(x+1)(x^{2}+1)(x^{4}+1)$. En utilisant le lemme chinois,
on obtient que $\Qq[x]/(x^{8}-1)\simeq \Qq[x]/(x+1)\times
\Qq[x]/(x^{2}+1)\times \Qq[x]/(x^{4}+1)$ soit~:

$$
\Qq[x]/(x^{8}-1)\simeq
\Qq\times\Qq\times\Qq[i]\times\Qq[e^{i\pi/4}].
$$

Montrons en effet que $\Qq[x]/(x^{2}+1)\simeq\Qq[i]$~: l'application
$\phi:\Qq[x]/(x^{2}+1)\to \Qq[i]$ définie par $\bar P\mapsto P(i)$ est un
morphisme d'anneau.
\begin{itemize}
\item injectivité~:
  Soit $\bar{P}\in\ker\phi$. Alors $P(i)=0$. Comme $P$ est à coefficient
  rationnels donc réels, $-i$ est aussi raine de $P$. Donc $x^{2}+1|P$.

\item surjectivité~:
  Soit $z=a+ib\in\Qq[i]$. Alors $z=\phi(ax+b)$.
\end{itemize}

\smallskip

De même pour $\Qq[x]/(x^{4}+1)\simeq\Qq[e^{i\pi/4}]$. Considérons le
morphisme $\phi~:\Qq[x]/(x^{4}+1)\to\Qq[e^{i\pi/4}]$ défini par
$\phi(\bar P)=P(e^{i\pi/4})$. $\phi$ est bien définie, c'est un morphisme
d'anneau.
\begin{itemize}
\item injectivité~:
  Soit $\bar{P}\in\ker\phi$. Alors $P(e^{i\pi/4})=0$. Par ailleurs
  $X^{4}+1$ est \emph{irréductible} dans $\Qq$~: sa factorisation sur
  $\Rr$ est $(x^{2}+\sqrt{2}x+1)(x^{2}-\sqrt{2}x+1)$, et aucun de ces
  deux polynômes, même à inversible réel près, n'est rationnel. On en
  déduit que si $(x^{4}+1)$ ne divise pas $P$, alors $\pgcd(X^{4}+1,P)=1$.
  Il existerait donc $U,V\in\Qq[x]$, $UP+V(X^{4}+1)=1$. En évaluant en
  $x=e^{i\pi/4}$, on obtient une contradiction. Donc $X^{4}+1|P$. (cf.
  exexercice \ref{ex:bar9}).

\item surjectivité~:
  Soit $z=a+be^{i\pi/4}\in\Qq[e^{i\pi/4}]$. Alors $z=\phi(ax+b)$.
\end{itemize}


\item
  On a $K[x]/(f^{n}g^{m})\simeq K[x]/(f^{m})\times K[x]/(g^{m})$. On en
  déduit que les diviseurs de $0$ sont les polynômes de la forme
  $\bar{P}$ où $P$ satisfait l'une des conditions suivantes~:
$$
\left|
\begin{array}{ll}
  f^{n}|P \text{ et } g^{m}\!\!\!\not| P & \qquad(\{0\}\times K[x]/(g^{m})\setminus\{0\})\\
  g^{m}|P \text{ et } f^{n}\!\!\!\not| P & \qquad(K[x]/(f^{n})\setminus\{0\}\times \{0\})\\
  f|P \text{ et } f^{n}\!\!\!\not| P & \qquad (\mathcal{D}_{K[x]/(f^{n})}\times K[x]/(g^{m}))\\
  g|P \text{ et } g^{m}\!\!\!\not| P & \qquad (K[x]/(f^{n})\times\mathcal{D}_{ K[x]/(g^{m})})
\end{array}\right.
$$
Les nilpotents sont donnés par les conditions
$$
\left\{
\begin{array}{ll}
  fg|P\\
  (f^{n}g^{m}\!\!\!\not| P \text{ si on veut exclure }0)
\end{array}\right.
$$

\item
   Les idéaux de $K[x]/(f^{n})$ sont les idéaux engendrés par les
   diviseurs de $f^{n}$ soit les $f^{k}$ pour $0\leq k\leq n$.

   \smallskip

   La démonstration peut se faire en toute généralité exactement de la
   même manière que dans $\Zz/n\Zz$~: Soit $\mathcal{D}$ l'ensemble des
   diviseurs de $f^{n}$ (modulo $K^{*}$). Ici, $\mathcal{D}=\{f^{k},0\leq
   k\leq n\}$. Soit $\mathcal{I}$ l'ensemble de idéaux de $K[x]/(f^{n})$.

   On a une flèche de $\mathcal{D}\to\mathcal{I}$, donnée par
   $d\mapsto(\bar{d})$.
   \begin{itemize}
   \item surjectivité
     Soit $I\in\mathcal{I}$. $I$ est principal~: notons $I=(\bar{h})$.
     Soit $d=\pgcd(f,h)$, et $h_{1}$ le polynôme déterminé par
     $h=dh_{1}$. Alors $\pgcd(f,h_{1})=0$ et $h_{1}$ est inversible dans
     le quotient. On en déduit que $(\bar{h})=(\bar{d})=I$ (or $d\in\mathcal{D}$).

   \item injectivité
     Soit $d,d'\in\mathcal{D}$ tels que $(\bar{d})=(\bar{d'})$. On a
     alors $d=h_{1}d'+h_{2}f$ donc $d'|d$. De même, $d|d'$. On en déduit
     que $d\sim d'$.
   \end{itemize}

   \smallskip
   
   Revenons à notre exercice~: les idéaux de $K[x]/(f^{n})\times
   K[x]/g^{m}$ sont donc de la forme $(f^{\alpha})\times(g^{\beta})$. En
   revenant à $K[x]/(f^{n}g^{m})$, on obtient que l'ensemble des idéaux
   est $$\{(f^{\alpha}g^{\beta}), 0\leq\alpha,\beta\leq n\}$$

\item
   Les inversibles de $K[x]/(f^{n})$ sont les (classes des) polynômes
   premiers avec $f$. Le complémentaire est donc formé des multiples de
   $f$, il y en a donc autant que de polynômes de degré $(nd-1)-d$ où $d$
   est le degré de $f$, soit $p^{(n-1)d}$. Il y a donc $p^{(n-1)d}(p-1)$
   inversibles dans $K[x]/(f^{n})$.

   On en déduit qu'il y en a $p^{(n-1)d_{f}+(m-1)d_{g}}(p-1)^{2}$ dans
   $K[x]/(f^{n}g^{m})$, où $d_{f}$ et $d_{g}$ sont les degrés respectifs
   de $f$ et $g$.

\item
   Plus généralement, si les $f_{i}$ sont des polynômes irréductibles
   distincts, dans $K[x]/(f_{1}^{n_{1}}\cdots f_{k}^{n_{k}})$ il y a
   $p^{\sum (n_{i}-1)d_{i}}(p-1)^{k}$ inversibles, où $d_{i}$ est le
   degré de $f_{i}$.
  
\end{enumerate}
\fincorrection
\correction{002307}
Pour obtenir les facteurs multiples, on utilise la remarque suivante~:
$g$ est un facteur multiple de $f$ ssi $g$ est un facteur commun à $f$ et
à $f'$ (dérivé formel de $f$).

Ainsi $\pgcd(f,f')$ est le produit de tous les facteurs multiples de $f$,
avec exposant diminué de $1$ par rapport à $f$. Ainsi $f/\pgcd(f,f')$ est
le produit de tous les facteurs irréductibles de $f$, avec exposant $1$
pour tous. Finalement, $\pgcd(\pgcd(f,f'),f/\pgcd(f,f'))$ est le produit
de tous les facteurs multiples de $f$ avec exposant $1$.
\fincorrection
\nocorrection
\correction{002309}
Soit $z=n+m\sqrt{d},z'=n'+m'\sqrt{d}\in \Z[\sqrt{d}]$. Alors
\begin{align*}
\overline{zz'}
&=\overline{(n+m\sqrt{d})(n'+m'\sqrt{d})} \\
&=\overline{(nn'+mm'd)+(nm'+n'm)\sqrt{d}} \\
&=          (nn'+mm'd)-(nm'+n'm)\sqrt{d}  \\
&=(n-m\sqrt{d})(n'-m'\sqrt{d}) \\
&=\bar{z}\, \bar{z}'
\end{align*}

Donc $\forall z,z'\in\Z[\sqrt{d}], \overline{zz'}=\bar{z}\,\bar{z}'$.

On a alors $\forall z,z'\in\Z[\sqrt{d}],\
 N(zz')=zz'\,\overline{zz'}=z\bar{z}\,z'\bar{z}'=N(z)\,N(z')$.
\fincorrection
\correction{002310}
\begin{enumerate}
\item
  \begin{itemize}
  \item {Si $z\in\Z[\sqrt{d}]$ est inversible~:} 

    Alors $zz^{-1}=1$, donc $N(z)N(z^{-1})=1$. Comme $N(z)\in\Z$ et
    $N(z^{-1})\in\Z$, on a donc $N(z)\in\{1,-1\}$.
    
  \item Si $N(z=\pm1)$~:

    Alors $z\bar{z}=\pm1$, donc $z(\pm\bar{z})=1$. Comme
    $\pm\bar{z}\in\Z[\sqrt{d}]$, $z$ est inversible.
  \end{itemize}
\item Soient $z_{1},z_{2}\in\Z[\sqrt{d}]$ tels que $z=z_{1}z_{2}$. Alors
$N(z_{1})N(z_{2})=\pm p$. Comme $\pm p$ est irréductible sur $\Z$, on en
déduit que $N(z_{1})=\pm1$ ou $N(z_{2})=\pm1$. D'après la question
précédente, on a $z_{1}\in\Z[\sqrt{d}]^{\times}$ ou
$z_{2}\in\Z[\sqrt{d}]^{\times}$~: on en déduit que $z$ est irréductible
dans $\Z[\sqrt{d}]$.

(Attention~: $p$ est premier donc irréductible dans $\Z$, mais peut être
réductible dans $\Z[\sqrt{d}]$~! cf. $2$ dans $\Z[i]$.)
\item On a $N(3)=N(2+\sqrt{-5})=9$. On peut montrer en fait que tout élément
$z$ de norme $9$ est irréductible~: si $z=z_{1}z_{2}$, alors
$N(z_{1})N(z_{2})=9$. Donc $\{N(z_{1}),N(z_{2})\}=\{1,9\}$ ou $\{3,3\}$
(dans $\Z[\sqrt{-5}]$, la norme est toujours positive). Or pour tout
$(n,m)\in\Z^{2}, n^{2}+5m^{2}\neq 3$. En effet, si $|m|\geq 1,
n^{2}+5m^{2}\geq 5$ et pour $m=0$, l'équation revient à $n^{2}=3$, qui
n'a pas de solution entière. Ainsi, $N(z_{1})=1$ ou $N(z_{2})=1$, donc
$z_{1}$ ou $z_{2}$ est inversible. $z$ n'a donc pas de factorisation non
triviale~: $z$ est irréductible dans $\Z[\sqrt{-5}]$. En particulier, $3$
et $2+\sqrt{-5}$ le sont.
\item Tout élément de $A$ de norme $9$ est irréductible. Il suffit donc de
trouver tous les éléments de norme $9$. Soit $z=n+m\sqrt{-5}\in A$. Si
$|m|\geq 2$ ou $|n|\geq 4$, alors $N(z)>9$. On cherche donc les éléments
de norme $9$ parmi les éléments $z=n+m\sqrt{-5}$ avec $|n|\leq3$ et
$|m|\leq 1$. Pour $m=0$, les seules solutions sont $n=\pm 3$, pour
$|m|=1$, les solutions sont obtenues pour $|n|=2$. Ainsi~:
$$
\forall z\in A:\ \ N(z)=9\Leftrightarrow z\in\{\pm3,\pm(2\pm\sqrt{5})\}
$$
\item On a $N(9)=81$. Donc si $9=z_{1}z_{2}$ est une factorisation de $9$ dans
$A$, $N(z_{1})N(z_{2})$ est une factorisation de $81$ (dans $\Z$), et
plus précisément on a
$\{N(z_{1}),N(z_{2})\}\in\Big\{\{1,81\},\{3,27\},\{9,9\}\Big\}$.

Si $N(z_{1})=1$ ou $N(z_{2})=1$, la factorisation est triviale.

$A$ n'a pas d'élément de norme $3$ donc la paire $\{3,27\}$ n'est pas
réalisable.

Si enfin $N(z_{1})=N(z_{2})=9$, alors
$z_{1},z_{2}\in\{\pm3,\pm(2\pm\sqrt{5})\}$. Comme
$9=3\cdot3=(2+\sqrt{-5})(2-\sqrt{-5})$, tous ces éléments sont diviseurs
de $9$.

Les diviseurs de $9$ sont donc $\{\pm1,\pm3,\pm(2\pm\sqrt{-5}),\pm9\}$. 

\medskip

Comme $N(3(2+\sqrt{-5}))=81$, le même raisonnement montre que si $d\in A$
divise $3(2+\sqrt{-5})$, alors
$d\in\{\pm1,\pm3,\pm(2\pm\sqrt{-5}),\pm3(2\pm\sqrt{-5})\}$.

Si $(2-\sqrt{-5})a=3(2+\sqrt{-5})$, alors $N(a)=9$, donc $a=\pm3$ ou
$\pm(2\pm\sqrt{-5})$. Comme $A$ est intègre, si $a=\pm3$, on obtient
$2-\sqrt{-5}=\pm(2+\sqrt{-5})$, ce qui est faux. Si $a=\pm(2+\sqrt{-5})$,
on obtient $2-\sqrt{-5}=\pm3$, ce qui est faux. Si enfin
$a=\pm(2-\sqrt{-5})$, on obtient $\pm(-1-4\sqrt{-5})=6+3\sqrt{-5})$, ce
qui est encore faux. Donc $2-\sqrt{-5}$ ne divise pas $3(2+\sqrt{-5})$
dans $A$. Tous les autres éléments de norme $9$ divisent
$3(2+\sqrt{-5})$, donc, finalement~:

Les diviseurs de $3(2+\sqrt{-5})$ sont
$\{\pm1,\pm3,\pm(2+\sqrt{-5}),\pm3(2+\sqrt{-5})\}$.

(Attention~: Le seul fait que $3$ et $2+\sqrt{-5}$ soient irréductibles
ne permet pas de conclure~! Si l'anneau n'est pas factoriel, un produit
d'irréductibles $p_{1}p_{2}$ peut avoir d'autres diviseurs (à association
près) que $p_{1}$ et $p_{2}$... cf $3\cdot3=(2+\sqrt{-5})(2-\sqrt{-5})$~!)

\item On connaît la liste des diviseurs de $3$ et de $2+\sqrt{-5}$. Les seuls
qui soient communs sont $1$ et $-1$. On en déduit que $1$ est un $\pgcd$
de $3$ et $2+\sqrt{-5}$. 

\medskip

$9$ et $3(2+\sqrt{-5})$ sont des multiples communs de $3$ et
$2+\sqrt{-5}$, donc si ces deux éléments admettent un $\mathrm{ppcm}$ $m$, on a
$m|9$ et $m|3(2+\sqrt{-5})$. On connaît la liste des diviseurs de $9$ et
$3(2+\sqrt{-5})$~: à association près, on en déduit que
$m\in\{1,3,2+\sqrt{-5}\}$. Comme $3|m$, la seule possibilité est $m=3$,
et comme $(2+\sqrt{-5})|m$, la seule possibilité est $m=2+\sqrt{-5}$. Il
y a donc contradiction~:

$3$ et $2+\sqrt{-5}$ n'ont pas de $\mathrm{ppcm}$ dans $A$.

\item 
Supposons $I$ principal~: soit $a\in A$ un générateur~: $I=(a)$. Alors
$a$ est un diviseur commun à $3$ et $2+\sqrt{-5}$, donc $a=\pm1$. (En
particulier, $I=A$). Soient $u=u_{1}+u_{2}\sqrt{-5}$ et
$v=v_{1}+v_{2}\sqrt{-5}$ deux éléments de $A$. On a~:
\begin{align*}
3u+(2+\sqrt{-5})v=1
&\Leftrightarrow (3u_{1}+2v_{1}-5v_{2})+(3u_{2}+v_{1}+2v_{2})\sqrt{-5}=1\\
&\Leftrightarrow \left\{
\begin{array}{rl}
3u_{1}+2v_{1}-5v_{2}&=1\\
3u_{2}+v_{1}+2v_{2}&=0
\end{array}\right.\\
&\Rightarrow \left\{
\begin{array}{rl}
-v_{1}+v_{2}&\equiv1 [3]\\
v_{1}-v_{2}&\equiv0 [3]
\end{array}\right.
\end{align*}
Donc $\forall u,v\in A,\ 3u+(2+\sqrt{-5})v\neq 1$. Donc $1\notin I$, ce
qui est une contradiction~: $I$ n'est pas principal.

L'anneau $A$ n'est pas principal puisqu'il a au moins un idéal non
principal. Il n'est pas non plus factoriel, puisque
$9=3\,3=(2+\sqrt{-5})(2-\sqrt{-5})$ admet deux factorisation en
irréductibles non équivalentes à association près.

\item
  \begin{itemize}
  \item 
Les diviseurs communs de $9$ et $3(2+\sqrt{-5})$ sont $\{\pm1,\pm 3,
\pm(2+\sqrt{-5})\}$. Si $9$ et $3(2+\sqrt{-5})$ admettent un $\pgcd$ $d$,
alors $d$ est dans cette liste, et divisible par tous les membre de cette
liste. Mais $3$ n'est pas divisible par $2+\sqrt{-5}$ et $2+\sqrt{-5}$ ne
divise pas $3$~: $9$ et $2+\sqrt{-5}$ n'ont pas de $\pgcd$.


  \item
Supposons que $9$ et $3(2+\sqrt{-5})$ admettent un $\mathrm{ppcm}$ $M$. Alors il
existe des éléments $a,b\in A$ tels que $M=9a=3(2+\sqrt{-5})b$. Notons
$m=3a=(2+\sqrt{-5})b$ ($A$ est intègre).

$m$ est un multiple commun de $3$ et $2+\sqrt{-5}$.

Soit $k$ un multiple commun de $3$ et $2+\sqrt{-5}$. Alors $3k$ est un
multiple commun de $9$ et $3(2+\sqrt{-5})$, donc $M|3k$~: $\exists c\in
A, 3k=Mc=3mc$. On en déduit que $k=mc$ ($A$ est intègre), donc $m|k$. On
en déduit que $m$ est un $\mathrm{ppcm}$ de $3$ et $2+\sqrt{-5}$, ce qui est
impossible.
  \end{itemize}
\end{enumerate}
\fincorrection
\correction{002311}
\begin{enumerate}
  \item  $\bar{n}$ est inversible ssi $\pgcd(n,36)=1$ (Bezout~!), i.e.
$\bar{n}\in\{\pm1, \pm5, \pm7, \pm11, \pm13, \pm17\}$. Les autres
éléments sont tous des diviseurs de $0$ puisque $\bar{n}$ divise $0$ ssi
$\pgcd(n,36)\neq1$. Enfin, $\bar{n}$ est nilpotent ssi $2|n$ et $3|n$,
donc ssi $6|n$, soit $\bar{n}\in\{0,\pm6,\pm12,18\}$.

  \item Montrons que l'ensemble $\mathcal{I}$ des idéaux de $\Zz/36\Zz$ est en
bijection avec l'ensemble $\mathcal{D}=\{1,2,3,4,6,9,12,18,36\}$  des
diviseurs (positifs) de $36$.

Considérons l'application $\phi:\mathcal{D}\to\mathcal{I}$ définie par
$\phi(d)=(\bar{d})$.

\textsl{Injectivité~:} Si $\phi(d)=\phi(d')$, alors $\exists a,b\in\Zz,
d=d'a+36b$. Comme $d|36$, on en déduit que $d|d'$. De même, on a $d'|d$,
et donc $d=d'$.

\textsl{Surjectivité~:}Soit $I\in\mathcal{I}$. $\Zz/36\Zz$ est principal,
donc $\exists a\in\Zz, I=(\bar{a})$. Soit $d=\pgcd(a,36)$. Notons
$a=da'$~: $\pgcd(a',36)=1$. On en déduit que $\bar{a}'$ est inversible
dans $\Zz/36\Zz$. Alors $\bar{d}\sim\bar{a}$ dans $\Zz/36\Zz$. On en
déduit que $I=(\bar{d})=\phi(d)$.

Finalement, il y a donc 9 idéaux dans $\Zz_{36}$~:
\begin{itemize}
\item $(\overline{ 1})=\Zz_{36}$,
\item $(\overline{ 2})=\{0,\pm2,\pm4,\pm6,\pm8,\pm10,\pm12,\pm14,\pm16,18\}$, 
\item $(\overline{ 3})=\{0,\pm3,\pm6,\pm9,\pm12,\pm15,18\}$,
\item $(\overline{ 4})=\{0,\pm4,\pm8,\pm12,\pm16\}$,
\item $(\overline{ 6})=\{0,\pm6,\pm12\}$
\item $(\overline{ 9})=\{0,\pm9,18\}$
\item $(\overline{12})=\{0,\pm12\}$
\item $(\overline{18})=\{0,18\}$
\item $(\overline{36})=\{0\}$,
\end{itemize}

  \item Si $a,b\in A^{\times}$, alors $(a b)(b^{-1} a^{-1})=1$ donc $ab\in
A^{\times}$.

Si $ab\in A^{\times}$, soit $c=(ab)^{-1}$. Alors $a(bc)=1$ donc $a\in
A^{\times}$ et $b(ac)=1$ donc $b\in A^{\times}$.
  \item On a $(6x+1)(-6x+1)=1$ dans $\Zz_{36}[x]$, donc $18x+1$ y est inversible.
  \item Soit $f$ un inversible de $\Zz_{36}[x]$. Choisissons $P\in\Zz[x]$ tel que
$\bar{P}=f$ et $Q\in\Zz[x]$ tel que $\bar{Q}=f^{-1}$.

La projection $\Zz\to\Zz_{2}$ se factorise par
$\Zz\to\Zz_{36}\to\Zz_{2}$. Ces projections sont bien définies, et sont
des morphismes d'anneaux. Notons $P_{[2]}$ la réduction de $P$ modulo
$2$~: on a alors $P_{[2]}Q_{[2]}=(PQ)_{[2]}=1$, et comme $\Zz_{2}$ est un
corps, $P_{[2]}=1$, $Q_{[2]}=1$. On en déduit que $2$ divise tous les
coefficients de $P$, sauf celui de degré $0$. De même, en considérant la
réduction modulo $3$, on obtient que $3$ divise tous les coefficients de
$P$, sauf celui de degré $0$. Finalement, $6$ divise tous les
coefficients de $P$ sauf celui de degré $0$, qui est inversible modulo
$36$~: à association (dans $\Zz_{36}$) près, $f$ est donc de la forme~:
$$
f=\sum_{i=1}^{d} 6a_{i} x^{i}+1,\qquad (a_{i})\in\Zz_{36}.
$$

Réciproquement, si $f$ est de cette forme, c'est à dire  $f=1+6xf_{1}$,
avec $f_{1}\in\Zz_{36}[x]$, alors~:
$$
 (1+6xf_{1})(1-6xf_{1})=1
$$
donc $f$ est inversible.

\end{enumerate}
\fincorrection
\correction{002312}
\begin{enumerate}
  \item Le critère d'Eisenstein avec $2$ pour module donne directement le
résultat.
  \item La réduction modulo $2$ de $Q$ est $Q_{[2]}=x^{6}+x^{2}+1$, qui n'a pas
de racine, et n'est pas divisible par $x^{2}+x+1$, le seul irréductible
de degré $2$ de $\Zz_{2}[x]$. Ainsi, $Q_{[2]}$ est soit irréductible,
auquel cas $Q$ l'est aussi sur $\Zz$, soit le produit de deux
irréductibles de degré $3$.

Si $Q_{[2]}$ n'est pas irréductible, on considère la réduction modulo $3$
de $Q$~: $Q_{[3]}=x^{6}+1=(x^{2}+1)^{3}$. $x^{2}+1$ est irréductible sur
$\Zz_{3}$, car il est de degré $2$ et n'a pas de racine. Soit $Q=RS$ une
factorisation non triviale de $Q$ sur $\Zz$. On peut supposer $R$ et $S$
unitaires. Alors, en considérant la réduction modulo $2$, on obtient que
$R_{[2]}$ et $S_{[2]}$ sont deux irréductibles de degré $3$ de
$\Zz_{2}[x]$. En particulier $\deg(R)=\deg(R_{[2]})=3$ (car $R$ est
unitaire) et $\deg(S)=\deg(S_{[2]})=3$. Cependant, la réduction modulo
$3$ de $Q$ n'admet pas de factorisation suivant deux polynômes de degré
3. C'est une contradiction~: on en déduit que $Q$ n'a pas de factorisation
non triviale.
\end{enumerate}
\fincorrection
\correction{002313}
Soit $p$ un nombre premier impair. Notons $p=2m+1$. On a
$$
 (m!)^{2}\equiv(-1)^{m+1} [p]
$$
en effet, (modulo $p$)~:
\begin{align*}
(p-1)!
&=\prod_{k=1}^{2m}k=m!\prod_{k=1}^{m}(m+k)\\
&=m!\prod_{k=1}^{m}(m+k-p)=m!\prod_{k=1}^{m}(-k)\\
&=(-1)^{m}(m!)^{2}
\end{align*}
Or, dans $\Zz_{p}[x]$, $1^{-1}=1$ et $(p-1)^{-1}=p-1$,  donc $\forall
k\in\{2,...,p-2\}$, $k^{-1}\in\{2,...,p-2\}$. Ainsi,
$\prod_{k=2}^{p-1}k\equiv1 [p]$, et donc $(p-1)!\equiv -1 [p]$. D'où le
résultat.

\begin{itemize}
\item 
  Si $p\equiv 1[4]$, $(-1)^{m+1}=-1$, et donc $m!$ est une solution de
  $x^{2}\equiv-1[p]$.
  
\item
  Si cette équation a une solution, alors $x^{2m}\equiv 1[p]$, et comme
  $x^{p-1}\equiv 1[p]$, $1\equiv (-1)^{m}[p]$. On en déduit que $m$ est
  pair, donc $p\equiv1[4]$.
\end{itemize}

\fincorrection
\correction{002314}
\begin{enumerate}
  \item \begin{align*}
%(x^{6}+x^{5}+x^{4}+x^{3}+1)&=(x^{3}+x^{2}+1)(x^{3}+x+1)+(x^{2}+x)\\
%(x^{3}+x^{2}+1)&=(x^{2}+x)x+1
f&=g(x^{3}+x+1)+(x^{2}+x)\\
g&=(x^{2}+x)x+1
\end{align*}
donc $\pgcd(f,g)=1$ et
$$
1=g-(x^{2}+x)x=g-(f-g(x^{3}+x+1))x=(x^{4}+x^{2}+x+1)g-xf
$$
  \item $f=(x^{4}+x+1)(x^{2}+x+1)$ donc $f$ n'est pas irréductible.

$g$ est de degré $3$ et n'a pas de racine, donc $g$ est irréductible.
  \item Les éléments de $A$ sont en bijection avec les polynômes de $\Zz_{2}[x]$
de degré $<\deg(g)=3$. Il y a $8$ polynômes de degré au plus $2$ sur
$\Zz_{2}$, donc $A$ a $8$ éléments.
  \item On utilise la représentation linéaire $uf+vg=1$ de $\pgcd(f,g)$ obtenue
plus haut. $uf=1+vg$, donc $\bar{u} \bar{f}=\bar{1}+\bar{0}=\bar{1}$.
Donc $(\bar{f})^{-1}=\bar{u}=\bar{x}$.
  \item Soit $f_{1}=x^{2}+x+1$ et $f_{2}=x^{4}+x+1$. Alors $f_{1}f_{2}=f$ donc
$\bar{f_{1}}\bar{f_{2}}=\bar{0}$. Pourtant, $f$ ne divise ni $f_{1}$ ni
$f_{2}$, donc $\bar{f_{1}}\neq \bar{0}$ et $\bar{f_{2}}\neq \bar{0}$~:
 $B$ n'est pas intègre, donc $B$ n'est pas un corps.
\end{enumerate}
\fincorrection
\nocorrection
\nocorrection
\nocorrection
\nocorrection
\nocorrection
\nocorrection
\nocorrection
\nocorrection
\nocorrection
\nocorrection
\nocorrection
\nocorrection
\nocorrection
\nocorrection
\nocorrection
\nocorrection
\nocorrection
\nocorrection
\nocorrection
\nocorrection
\nocorrection
\nocorrection
\nocorrection
\nocorrection
\nocorrection
\nocorrection
\nocorrection
\correction{002261}
 \begin{enumerate}
  \item
    Le polynôme $X$ n'est jamais inversible dans $A[X]$. Si $A$ n'est pas
    intègre, comme $A\subset A[X]$, $A[X]$ ne l'est pas non plus et ne
    peut pas être un corps. Si $A$ est intègre et si $X=PQ$, alors
    $\deg(P)+\deg(Q)=1$ donc $P$ ou $Q$ est une constante. Supposons par
    exemple que ce soit $P$. $P|X$ donc $P|1$ donc $P$ est inversible, et
    $Q\sim X$.
  \item
    Soit $P=X+a$ un polynôme unitaire linéaire de $A[X]$. Supposons que 
    $P=P_1P_2$. Comme $A$ estintègre, on a $\deg(P_1)+\deg(P_2)=1$, donc
    $P_1$ ou $P_2$ est une constante. Supposons que ce soit $P_1$. Alors
    $P_1|1$ et $P_1|a$. En particulier, $P_1$ est inversible, et donc 
    $P_2\sim P$.
  \item
    Les polynômes irréductibles de $\Cc[X]$ sont les polynômes de degré 1
    (théorème de Gauss).

    Les irréductibles de $\Rr[X]$ sont les polynômes de degré 1  et les
    polynômes de degré 2 sans racine réelles. En effet, soit
    $P\in\Rr[X]$. $P$ se factorise sur $\Cc[X]$ sous la forme
    $P=a\prod(X-\lambda_i)^{\nu_i}$ (avec $i\neq j\Rightarrow
    \lambda_i\neq \lambda_j$). Comme cette factorisation est unique, et
    que $P=\overline{P}$, on en déduit que si $\lambda_i$ est racine de
    $P$ avec multiplicité $\nu_i$, alors il en va de même pour
    $\overline{\lambda_i}$. Ainsi, on obtient une factorisation de $P$
    dans $\Rr[X]$: $P=a\prod_{\lambda_i\in\Rr} (X-\lambda_i)^{\nu_i}
    \prod(X^2-2\Re(\lambda_i)X+|\lambda_i|^2)^{\nu_i}$.

    $P$ est donc irréductible ssi $P$ est de la forme $P=a(X-\lambda)$
    avec $\lambda\in\Rr$ ou $P=a(X^2-2\Re(\lambda_i)X+|\lambda_i|^2)$
    avec $\lambda\notin\Rr$.

  \item
     Supposons que $K[X]$ ait un nombre fini de polynômes unitaires
     irréductibles $P_1,\dots ,P_k$. Soit alors $P=\prod_{i=1}^k P_i+1$.

     Comme $K$ est un corps, les irréductibles sont de degré au moins
     $1$, et donc $P$ n'est pas l'un des $P_i$. Comme $P$ est unitaire,
     $P$ n'est pas irréductible. En particulier, l'un au moins des $P_i$
     divise $P$. Supposons par exemple que ce soit $P_1$: $\exists Q\in
     K[X], P=P_1Q$. Alors $P_1(Q-\prod_{i=2}^k P_i)=1$. Donc $P_1$ est
     inversible, ce qui est faux.
  \end{enumerate}
\fincorrection
\correction{002262}
  \begin{enumerate}
  \item
    Supposons $(X,n)$ principal dans $\Zz[X]$: $(X,n)=(P_0)$. Alors
    $P_0|n$ donc $P_0\in\Zz$, et $P_0|X$ donc $P_0=\pm1$. Ainsi
    $(P_0)=\Zz[X]$. Or $(X,n)$ est l'ensemble des polynômes dont le terme
    constant est un multiple de $n$: en effet, si $P\in(X,n)$, $\exists
    A,B\in\Zz[X], P=AX+Bn$ donc le terme constant de $P$ est un multiple
    de $n$. Réciproquement, si le terme constant de $P=\sum p_i X^i$ est
    un multiple de $n$, $p_0=p'_0n$, alors $P=X(\sum_{i\geq
      1}p_iX^i)+p'_0n\in(X,n)$. Ainsi, $1\notin(X,n)$. Donc $(X,n)$ n'est
    pas principal.

  \item
    Si $A[X]$ est principal, soit $a\in A\setminus\{0\}$, et $I=(X,a)$.
    $A[X]$ étant principal, $\exists P_0\in A[X], I=(P_0)$. Alors $P_0|a$
    donc $P_0\in A$, et $P_0|X$ donc $P_0|1$ et $P_0$ est inversible. On
    en déduit que $I=A[X]$. En particulier $1\in I$: $\exists U,V\in
    A[X], XU+aV=1$. Le terme constant de $XU+aV$ est multiple de $a$ et
    vaut $1$. $a$ est donc inversible.

    Si $A$ est un corps, on dispose de la division euclidienne. Soit $I$
    un idéal de $A[X]$. Soit $P_0$ un élément de $I\setminus\{0\}$ de
    degré minimal. Soit $P\in I$. $\exists!(Q,R)\in A[X]^2, P=P_0Q+R$ et
    $\deg(R)<\deg(P)$. Comme $R=P-P_0Q$, on a $R\in I$, et comme
    $\deg(R)<\deg(P_0)$, on a $R=0$. Ainsi $P\in(P_0)$. On a donc
    $I\subset(P_0)\subset I$.  
  \end{enumerate}
\fincorrection
\correction{002263}
  Notons $f(x^{n})=P(x-1)$. Alors $f(1)=0\cdot P(1)=0$ et donc $(x-1)|f$.
  Notons $f=Q(x-1)$. On a alors $f(x^{n})=Q(x^{n})(x^{n}-1)$. $(x^{n}-1)$
  divise bien $f$.
\fincorrection
\correction{002264}
Notons $(Q,R)$ le quotient et le reste de cette division euclidienne:
  $(x-2)^{m}+(x-1)^{n}-1=Q(x-2)(x-1)+R$ avec $\deg(R)\leq 1$. Notons
  $R=ax+b$. En évaluant en $1$, on obtient $(-1)^{m}-1=a+b$, et en
  évaluant en $2$, $2a+b=0$. On en déduit $b=-2a$ et $a=1-(-1)^{m}$, soit
  $R=(1-(-1)^{m})(x-2)$.
\fincorrection
\correction{002265}
 \begin{enumerate}
  \item 
    Soit $P$ un polynôme de degré $d=2$ ou $3$ de $K[X]$.

    Si $P$ a une racine $a\in K$, alors $(X-a)|P$, et $P$ n'est pas
    irréductible.

    Réciproquement, si $P=AB$ avec $A,B\in K[X]$ et $A,B\notin
    K[X]^{\times}=K\setminus\{0\}$, alors $\deg(A)\geq1$, $\deg(B)\geq1$,
    et $\deg(A)+\deg(B)=d=2$ ou $3$, donc l'un au moins des deux
    polynômes $A$ et $B$ est de degré $1$. On peut supposer que c'est
    $A$. Notons $A=aX+b$. Alors $(X+a^{-1}b)|P$, et $-a^{-1}b$ est racine
    de $P$.

    Finalement $P$ a une racine ssi $P$ n'est pas irréductible.

    \item
      Irréductibles de degré $2$ de $\Zz/2\Zz$: Soit $P=aX^{2}+bX+c$ un
      polynôme de degré $2$. $a\neq0$ donc $a=1$.
      \begin{align*}
        P\text{ irréductible} &\Leftrightarrow P\text{ n'a pas de racine}\\
        &\Leftrightarrow
          \begin{cases}
            P(0)&\neq 0\\
            P(1)&\neq 0
          \end{cases}\\
        &\Leftrightarrow
          \begin{cases}
            P(0)&=1\\
            P(1)&=1
          \end{cases}\\
        &\Leftrightarrow
          \begin{cases}
            c&=1\\
            1+b+1&=1
          \end{cases}\\
        &\Leftrightarrow
          P=X^{2}+X+1
      \end{align*}
      Ainsi, il y a un seul irréductible de degré $2$, c'est $I_{2}=X^{2}+X+1$.

      Irréductibles de degré $3$ de $\Zz/2\Zz$: Soit $P=aX^{3}+bX^{2}+cX+d$ un
      polynôme de degré $2$. $a\neq0$ donc $a=1$.
      \begin{align*}
        P\text{ irréductible} &\Leftrightarrow P\text{ n'a pas de racine}\\
        &\Leftrightarrow
          \begin{cases}
            d&=1\\
            1+b+c+1&=1
          \end{cases}\\
        &\Leftrightarrow
          \begin{cases}
            d&=1\\
            (b,c)&=(1,0)\text{ ou }(b,c)=(0,1)
          \end{cases}\\
        &\Leftrightarrow
          P=X^{3}+X+1 \text{ ou } P=X^{3}+X^{2}+1
      \end{align*}
      Ainsi, il y a deux irréductibles de degré $3$ dans $\Zz/3\Zz[X]$:
      $I_{3}=X^{3}+X+1$ et $I'_{3}=X^{3}+X^{2}+1$.

    \item
      Soit $P=5X^{3}+8X^{2}+3X+15\in\Zz[X]$. Soient $A$ et $B$ deux
      polynômes tels que $P=AB$. L'application $\Zz\to\Zz/2\Zz,
      n\mapsto\bar n$ induit une application $\Zz[X]\to\Zz/2\Zz[X],
      P=\sum a_{i}X^{i}\mapsto \bar P=\sum \bar{a_{i}}X^{i}$. Cette
      application est compatible avec les opérations: en particulier
      $\overline{AB}=\bar{A}\,\bar{B}$ (pourquoi?). Ainsi on a:
      $\bar{P}=\bar{A}\bar{B}$. Or $\bar{P}=X^{3}+X+1$ est irréductible,
      donc (quitte à échanger les rôles de $A$ et $B$ on peut supposer
      que) $\bar{A}=1$ et $\bar{B}=X^{3}+X+1$. On en déduit que $B$ est
      au moins de degré $3$, d'où $\deg(A)=0$. $A\in\Zz$ et $A|P$, donc
      $A|5$, $A|8$, $A|3$, et $A|15$. On en déduit que $A=\pm1$.
      Finalement, $A=\pm1$ et $B\sim P$. $P$ est donc irréductible dans
      $\Zz[X]$.

      \bigskip
      Soit $P=X^{5}+2X^{3}+3X^{2}-6x-5\in\Zz[X]$. Soient $A$ et $B$ deux
      polynômes tels que $P=AB$. On a comme précédemment:
      $\bar{P}=\bar{A}\bar{B}$ où $\bar{P}=X^{5}+X^{2}+1$. $\bar{P}$ n'a
      pas de racine dans $\Zz/2\Zz$, donc si $\bar{P}$ est réductible, il
      doit être le produit d'un irréductible de degré $2$ et d'un
      irréductible de degré $3$. Or $\bar{P}\neq I_{2}I_{3}$ et
      $\bar{P}\neq I_{2}I'_{3}$ (faire le calcul!), donc $\bar{P}$ est
      irréductible. Le même raisonnement montre alors que $P$ est
      irréductible dans $\Zz[X]$.

    \item
      Un polynôme de degré $4$ est réductible ssi il a une racine ou est
      le produit de deux irréductibles de degré $2$. Soit
      $P=\sum_{i=0}^{4}a_{i}X^{i}\in\Zz/2\Zz[X]$, avec $a_{4}=1$.
      \begin{align*}
        P\text{ irréductible}&\Leftrightarrow 
          \begin{cases}
            P(0)\neq 0\\
            P(1)\neq 0\\
            P\neq I_{2}^{2}\\
          \end{cases}\\
        &\Leftrightarrow 
          \begin{cases}
            a_{0}= 1\\
            1+a_{3}+a_{2}+a_{1}+1=1\\
            P\neq I_{2}^{2}\\
          \end{cases}\\
        &\Leftrightarrow 
          P\in\{X^{4}+X^{3}+1,X^{4}+X+1,X^{4}+X^{3}+X^{2}+X+1\}
      \end{align*}

      \bigskip
      Un polynôme de degré $5$ est irréductible ssi il n'a pas de racine
      et l'est pas le produit d'un irréductible de degré $2$ et d'un
      irréductible de degré $3$. Tous calculs fait, on obtient la liste
      suivante: $\{%
      X^{5}            +X^{2}  +1,%
      X^{5}      +X^{3}        +1,%
      X^{5}+X^{4}+X^{3}+X^{2}  +1,%
      X^{5}+X^{4}+X^{3}      +X+1,%
      X^{5}+X^{4}      +X^{2}+X+1,%
      X^{5}      +X^{3}+X^{2}+X+1,%
      \}$.
  \end{enumerate}
\fincorrection
\correction{002266}  
\begin{enumerate}
  \item 
    On raisonne exactement comme pour l'exercice \ref{exoprec}. On peut
    réduire un peu les discussions en remarquant que puisqu'on est sur un
    corps, on peut se contenter de chercher les irréductibles
    \emph{unitaires}: on obtient les autres en multipliant les
    irréductibles unitaires par les inversibles, soit $\pm1$.
  
    Les irréductibles de degré 2 sont caractérisés par $P(0)\neq0$,
    $P(1)\neq0$ et $P(-1)\neq 0$. On obtient finalement la liste
    suivante: $\{X^{2}+1, X^{2}-X-1, -X^{2}-1, -X^{2}+X+1\}$.

    \bigskip
    Sans commentaire, on obtient la liste suivante pour les irréductibles
    de degré $3$ de $\Zz/3\Zz[X]$:
    $\{%
    \pm(X^{3}+X^{2}-X+1),%
    \pm(X^{3}-X^{2}+X+1),%
    \pm(X^{3}-X^{2}  +1),%
    \pm(X^{3}      -X+1),%
    \pm(X^{3}+X^{2}+X-1),%
    \pm(X^{3}-X^{2}-X-1)%
    \pm(X^{3}+X^{2}  -1),%
    \pm(X^{3}      -X-1),%
    \}$.
  
  \item
    $X^{2}+X+1=(X-1)^{2}$

    $X^{3}+X+2=(X+1)(X^{2}-X+2)$
    
    $X^{4}+X^{3}+X+1=(X+1)(X^{3}+1)=(X+1)^{4}$
  \end{enumerate}
\fincorrection
\correction{002267}
  On raisonne comme pour l'exercice \ref{exoprec}. Soit $P=X^{5}-6X^{3}+2X^{2}-4X+5$,
  $A,B$ deux polynômes tels que $P=AB$. En considérant la réduction
  modulo $2$, on a $\bar{P}=X^{5}+1$ donc la décomposition en facteurs
  irréductibles est $\bar{P}=(X+1)(X^{4}+X^{3}+X^{2}+X+1)$. Comme $P$ est
  unitaire, $A$ et $B$ le sont aussi, et la réduction modulo $2$ préserve
  donc le degré de $A$ et $B$. On en déduit que si $\bar{A}=X+1$, alors
  $A$ est de degré 1.

  La réduction modulo $3$ de $P$ devrait donc avoir une racine. Mais
  $P\mod 3=X^{5}-X^{2}-X-1$ n'a pas de racine dans $\Zz/3\Zz$. On en
  déduit que dans la réduction modulo $2$, la factorisation
  $\bar{P}=`\bar{A}\bar{B}$ est triviale ($\bar{A}=1$ et
  $\bar{B}=\bar{P}$ ou le contraire), puis que la factorisation $P=AB$
  elle même est triviale ($A=\pm1$ et $B=\mp P$ ou le contraire). Ainsi,
  $P$ est irréductible dans $\Zz[X]$. 

  \bigskip

  Pour $P=7X^{4}+8X^{3}+11X^{2}-24X-455$, on procède de la même façon. Si
  $P=AB$, comme $7$ est premier, l'un des polynômes $A$ ou $B$ a pour
  coefficient dominant $\pm7$ et l'autre $\mp1$.  On en déduit que les
  réductions modulo $2$ ou $3$ préservent le degré de $A$ et de $B$. Les
  décompositions en facteurs irréductibles sont les suivantes: $P\mod
  2=(X^{2}+X+1)^{2}$ et $P\mod3=(X-1)(X^{3}-X-1)$. Si la factorisation
  $P=AB$ est non triviale, alors les réductions modulo $2$ de $A$ et $B$
  sont de degré $2$, et donc $\deg(A)=\deg(B)=2$. Mais la décomposition
  modulo $3$ impose que ces degrés soient $1$ et $3$. La factorisation
  $P=AB$ est donc nécessairement triviale, et $P$ est donc irréductible.
\fincorrection
\correction{002268}
  Commençons par montrer que ces polynômes sont irréductibles sur $\Zz$.


  \paragraph{-Le cas de $f=\prod_{i=1}^{n}(X-a_{i})-1$}
  Soit $P,Q\in\Zz[X]$ tels que $f=PQ$. On peut supposer sans perte de
  généralité que $P$ et $Q$ ont des coefficients dominants positifs (i.e.
  sont unitaires).

  On a~: $\forall i, f(a_{i})= P(a_{i})Q(a_{i})=-1$ donc
  $$
   P(a_{i})=\pm1
  \qquad\text{ et }\qquad
   Q(a_{i})=\mp1
  $$
  Soit $I=\{i, P(a_{i})=-1\}$ et $J=\{1,\dots,n\}\setminus I$. On notera
  $|I|$ et $|J|$ le nombre d'éléments de $I$ et $J$.

  \textit{Supposons $I\neq\emptyset$ et $J\neq\emptyset$}: Alors
  $\prod_{i\in I}(X-a_{i})|(P+1)$ et $\prod_{i\in J}(X-a_{i})|(Q+1)$.
  Ainsi $\deg(P+1)\geq |I|$ et $\deg(Q+1)\geq |J|=n-|I|$, et comme
  $\deg(P)+\deg(Q)=n$, on en déduit que $\deg(P)=|I|$ et $\deg(Q)=|J|$,
  puis que (puisque $P$ et $Q$ sont unitaires)~:
  $$
    P=\prod_{i\in I}(X-a_{i})-1
    \qquad\text{ et }\qquad
    Q=\prod_{i\in J}(X-a_{i})-1.
  $$
  
  Ainsi $f=\prod_{k\in I\cup J}(X-a_{k})-1=(\prod_{i\in
    I}(X-a_{i})-1)(\prod_{j\in J}(X-a_{j})-1)=f-\big(\prod_{i\in
    I}(X-a_{i})+\prod_{j\in J}(X-a_{j})-2\big)$, donc $\prod_{i\in
    I}(X-a_{i})+\prod_{j\in J}(X-a_{j})-2=0_{\Zz[X]}$, ce qui est faux.

  Ainsi $I=\emptyset$ ou $J=\emptyset$. On peut supposer sans perte de
  généralité que $I=\emptyset$. Alors $\forall i\in\{1,\dots,n\},
  Q(a_{i})=-1$. Donc les $a_{i}$ sont tous racine de $Q+1$. Comme
  $\deg(Q+1)\leq n$ et $Q+1\neq 0$, on en déduit que $Q=f$, et $P=1$. $f$
  est donc bien irréductible dans $\Zz[X]$.

  \medskip
  \paragraph{-Le cas de $g=\prod_{i=1}^{n}(X-a_{i})^{2}+1$}. Supposons que $g=PQ$,
  avec $P,Q\in\Zz[X]$.On a $g(a_{i})=1=P(a_{i})Q(a_{i})$, donc
  $P(a_{i})=Q(a_{i})=\pm1$.

  Comme $g$ n'a pas de racine réelle, il en va de même de $P$ et $Q$, qui
  sont donc de signe constant (théorème des valeurs intermédiaires pour
  les fonctions continues sur $\Rr$~!). On peut donc supposer sans perte
  de généralité que $P$ et $Q$ sont positifs.

  Alors $P(a_{i})=Q(a_{i})=1$. Ainsi, tous les $a_{i}$ sont racines de
  $P-1$ et de $Q-1$. On a donc $\prod_{i=1}^{n}(X-a_{i})|P-1$ et
  $\prod_{i=1}^{n}(X-a_{i})|Q-1$.

  En particulier, si $P-1\neq0$ et $Q-1\neq0$, $\deg(P)\geq n$ et
  $\deg(Q)=2n-\deg(P)\geq n$. Ainsi $\deg(P)=\deg(Q)=n$. Comme en plus
  $P$ et $Q$ sont unitaires, on en déduit que 
  $$
    P-1=\prod_{i=1}^{n}(X-a_{i}) \qquad\text{ et }
    Q-1=\prod_{i=1}^{n}(X-a_{i}).
  $$
  On devrait donc avoir
  $(\prod_{i=1}^{n}(X-a_{i})+1)^{2}=\prod_{i=1}^{n}(X-a_{i})^{2}+1$, ce
  qui est faux ($\prod_{i=1}^{n}(X-a_{i})\neq 0_{\Zz[X]}$)~!

  Ainsi $P-1=0$ ou $Q-1=0$, et on en déduit bien que $g$ est irréductible
  dans $\Zz[X]$.


  
  \bigskip
  \paragraph{Irréductibilité dans $\Qq[X]$}
 On a le lemme suivant~:

  Si $P\in\Zz[X]$ est unitaire et irréductible dans $\Zz[X]$, alors il
  l'est aussi dans $\Qq[X]$.

  L'ingrédient de base de la démonstration est la notion de
  \emph{contenu} d'un polynôme $P\in\Zz[X]$~: c'est le $\pgcd$ de ses
  coefficients, souvent noté $c(P)$. Il satisfait la relation suivante :
  $$
    c(PQ)=c(P)c(Q).
  $$

  Supposons que $P=QR$, avec $Q,R\in\Qq[X]$, $Q$ et $R$ unitaires. En
  réduisant tous leurs coefficients de au même dénominateur, on peut
  mettre $Q$ et $R$ sous la forme~:
  $$
    Q=\frac{1}{a}Q_{1} 
    \qquad\text{ et }\qquad
    R=\frac{1}{b}R_{1} 
  $$
  avec $a,b\in\Zz$, $Q_{1},R_{1}\in\Zz[X]$ et $c(Q_{1})=1$, $c(R_{1})=1$.

  Alors $abP=Q_{1}R_{1}$, donc $c(abP)=c(Q_{1})c(R_{1})=1$. Comme
  $ab|c(abP)$, on a $ab=\pm1$, et en fait $P,Q\in\Zz[X]$. 
\fincorrection
\correction{002269}
$f$ est irréductible, donc si $f$, ne divise pas $g$, alors $f$ et $g$
 sont premiers entre eux. Ainsi,$\exists u,v\in\Qq[X], uf+vg=1$. En
 évaluant en $\alpha$, on obtient $u(\alpha)\cdot0+v(\alpha)\cdot0=1$ ce
 qui est impossible!
\fincorrection
\correction{002270}
Supposons que la fraction soit réductible. Alors, il existe
  $p,q,d\in\Zz$ tels que
  $$
  \begin{cases}
    11n+2m&=pd\\
    18n+5m&=qd\\
  \end{cases}
  $$
  On en déduit que 
  $$
  \begin{cases}
    19n&=5pd-2qd\\
    19m&=-18pd+1qd\\
  \end{cases}
  $$
  En particulier, $d|19n$ et $d|19m$. Si $d\neq19$, on a $\pgcd(n,m)\neq
  1$. Si $d=19$, alors 
  \begin{equation}\label{eq:n=5p-2q...}
  \begin{cases}
    n&=5p-2q\\
    m&=-18p+1q\\
  \end{cases}
  \end{equation}
  Réciproquement, si $\pgcd(n,m)\neq1$ ou si $n,m$ sont de la forme
  donnée par \eqref{eq:n=5p-2q...}, alors la fraction est réductible.

\fincorrection
\correction{002271}
Soit $d=\pgcd(m,n)$. Notons $n=dn'$ et $m=dm'$. Alors
  $X^{n}-1=(X^{d})^{n'}-1$. Or $(Y-1)|Y^{n'}-1$ donc
  $(X^{d}-1)|(X^{n}-1)$. De même, $(X^{d}-1)|(X^{m}-1)$, et donc
  $(X^{d}-1)|\pgcd(X^{n}-1,X^{m}-1)$.

  Par ailleurs, soit $D=\pgcd(X^{n}-1,X^{m}-1)$. Les racines de $D$ dans
  $\Cc$ sont des racines à la fois n-iéme et m-ième de $1$, qui sont
  touts simples~: elles sont donc de la forme $\omega=e^{i2\pi\alpha}$ où
  $\alpha=\frac{k}{n}=\frac{k'}{m}$. Ainsi $km'=k'n'$. On a
  $\pgcd(m',n')=1$, donc par le théorème de Gauss, on en déduit que $k'$
  est un multiple de $m'$, soit $\frac{k'}{m}=\frac{k''}{d}$, et $\omega$
  est donc une racine $d$-ième de 1. On en déduit que $D|X^{d}-1$, et
  finalement~:
  $$\pgcd(X^{n}-1,X^{m}-1)=X^{\pgcd(m,n)}-1.$$

\fincorrection
\correction{002272}
Utiliser l'algorithme d'Euclide. (on travaille dans $\Zz/2\Zz$).
  \begin{align*}
    x^{5}+x^{4}+1&=(x^{4}+x^{2}+1)(x+1)+x^{3}+x^{2}+x\\
    x^{4}+x^{2}+1&=(x^{3}+x^{2}+x)(x+1)+x^{2}+x+1\\
    x^{3}+x^{2}+x&=(x^{2}+x+1)x+0\\
  \end{align*}
  Donc $\pgcd(x^{5}+x^{4}+1,x^{4}+x^{2}+1)=x^{2}+x+1$, et
  \begin{align*}
    x^{2}+x+1&=(x^{4}+x^{2}+1)+(x^{3}+x^{2}+x)(x+1)\\
    &=(x^{4}+x^{2}+1)+\big((x^{5}+x^{4}+1)+(x^{4}+x^{2}+1)(x+1)\big)(x+1)\\
    &=(x^{4}+x^{2}+1)(1+(x+1)^{2})+(x^{5}+x^{4}+1)(x+1)\\
    &=(x^{4}+x^{2}+1)(x^{2})+(x^{5}+x^{4}+1)(x+1)\\
  \end{align*}

  \smallskip
  De même, $\pgcd(x^{5}+x^{3}+x+1,x^{4}+1)=x^{3}+1$ et
  $x^{3}+1=(x^{5}+x^{3}+x+1)+(x^{4}+1)x$.
\fincorrection
\correction{002273}
  Dans $\Zz/3\Zz$~: $\pgcd(x^{4}+1,x^{3}+x+1)=x^{2}+x-1$.

  \noindent 
  Dans $\Zz/5\Zz$~: $\pgcd(x^{4}+1,x^{3}+x+1)=1$.
\fincorrection
\correction{002274}
Sur $\Zz[X]$, $\pgcd(x^{4}+x^{3}-3x^{2}-4x-1,x^{3}+x^{2}-x-1)=1$.
\fincorrection
\correction{002275}
  \begin{enumerate}
  \item 
    $P$ est primitif, $2$ divise tous les coefficients de $P$ sauf le
    dominant, et $4$ ne divise pas le terme constant~: d'après le critère
    d'Eisenstein, on en déduit que $P$ est irréductible dans $\Zz[x]$
    (puis dans $\Qq[x]$ car il est unitaire...).

  \item
    On peut appliquer le même critère, avec $3$ cette fois.

  \item
    $f$ est primitif, et sa réduction modulo $2$ est irréductible. Donc
    $f$ est irréductible dans $\Zz[x]$.

  \item
    $f(x+1)=\sum_{k=1}^{p}C_{p}^{k}x^{k-1}$. Or $p|\frac{p!}{k!(p-k)!}$
    (car $p$ apparaît au numérateur, tandis que tous les facteurs du
    dénominateur sont $<p$~; comme $p$ est premier, ils sont donc
    premiers avec $p$). De plus $C_{p}^{1}=p$, donc $p^{2}$ ne divise pas
    le terme constant de $f(x+1)$. D'après le critère d'Eisenstein,
    $f(x+1)$ est irréductible, et donc $f$ aussi.
  \end{enumerate}
\fincorrection
\correction{002276}
 Soit $P=x^{2}-x+1$. Si $P$ a une factorisation non triviale, $P$ est
  divisible par un polynôme de degré $1$, et comme $P$ est unitaire, ce
  diviseur peut être choisi unitaire~: on en déduit que $P$ a une racine.
  On calcule $P(a+bi\sqrt{3})=(a^{2}-3b^{2}-a+1)+(2ab-b)i\sqrt{3}$. Comme
  $1/2\notin A=\Zz[i\sqrt{3}]$, $2a-1\neq 0$, donc si
  $P(a+bi\sqrt{3})=0$, alors $b=0$, et $P(a)=0$. Mais $x^{2}-x+1$ est
  primitif et se réduction modulo $2$ est irréductible, donc il est
  irréductible sur $\Zz[x]$. En particulier il n'a pas de racine dans
  $\Zz$. On en déduit que $P$ n'a pas de racine sur $A$, et est donc
  irréductible.

  \medskip
  Soit $K=\mathrm{frac}(A)=\Qq[i\sqrt{3}]$. On a
  $P(\frac{1+i\sqrt{3}}{2})=0$ donc $P$ a une racine dans $K$, donc $P$
  est réductible sur $K$.
\fincorrection
\correction{002277}
  Si $P$ a une racine $\alpha$ dans $\Zz$, alors $P(\alpha)=0$, et en
  considérant la réduction modulo $n$, $\bar{P}(\bar{\alpha})=0$, donc
  $\bar{P}$ a une racine dans $\Zz/n\Zz$ pour tout $n$.
  \begin{enumerate}
  \item 
     Si $P(0)$ et $P(1)$ sont impairs, $\bar{P}(\bar{0})=\bar{1}$ et
     $\bar{P}(\bar{1})=\bar{1}$, donc $\bar{P}$ n'a pas de racine sur
     $\Zz/2\Zz$. Donc $P$ n'a pas de racine sur $\Zz$.

  \item
     Si $n$ ne divise aucun des $P(0),\dots,P(n-1)$, alors
     $\bar{P}(\bar{0})\neq 0$,\dots, $\bar{P}(\overline{n-1})\neq 0$,
     donc $\bar{P}$ n'a pas de racine sur $\Zz/n\Zz$. Donc $P$ n'a pas de
     racine sur $\Zz$.
  \end{enumerate}
\fincorrection
\correction{002278}
  \begin{enumerate}
  \item 
    $(X-\frac{a}{b})|P$ donc $\exists Q\in\Qq[x], \,
    P=(x-\frac{a}{b})Q=(bx-a)\frac{Q}{b}$. En réduisant tous les
    coefficients de $Q$ au même dénominateur, on peut mettre $Q$ sous la
    forme~: $Q=\frac{1}{m}Q_{1}$, avec $Q_{1}\in\Zz[X]$ primitif. Alors
    $bdP=(bx-a)Q_{1}$. En considérant les contenus de ces polynômes, on a
    $c(bx-a)=\pgcd(a,b)=1$, $c(Q_{1})=1$ donc $c(bdP)=bd\,c(P)=1$. Ainsi
    $bd=\pm1$, et $(bx-a)|P$.

  \item
    On considère par exemple les cas $k=0,\dots,3$. (Pour $k=2$, on
    constate que $P(2)=0$~: on peut diviser $P$ par $(X-2)$ et déterminer
    les trois racines complexes de $P$...). On obtient que
    \begin{align*}
    (*)&&\qquad       a&|14       &\quad   (k=0),      \\
    (**)&&\qquad   (a-b)&|4       &\quad   (k=1),      \\
    (***)&&\qquad  (a-3b)&|2^{3}5 &\quad   (k=3).      \\
    \end{align*}
    Au passage On peut remarquer que si $\alpha\leq0$, $P(\alpha)<0$, donc on
    peut supposer $a>0$ et $b>0$.
    \begin{itemize}
    \item 
      Si $a=1$~: $(**)\Rightarrow b\in\{2,3,5\}$. Aucune de ces
      possibilités n'est compatible avec $(***)$.

    \item
      Si $a=2$~: $(**)\Rightarrow b\in\{1,3,4,6\}$. Comme $\pgcd(a,b)=1$,
      $4$et $6$ sont exclus. $3$ n'est pas compatible avec $(***)$. Pour
      $2$, on vérifie que $P(2)=0$.

    \item
      Si $a=7$~: $(**)\Rightarrow b\in\{3,5,9,11\}$. Mais aucune de ces
      solution ne convient.

    \item
      Si $a=14$~: $(**)\Rightarrow b\in\{10,12,16,18\}$ mais
      $\pgcd(a,b)=1$ exclu toutes ces possibilités.
    \end{itemize}
    Finalement, $2$ est la seule racine rationnelle de $P$.
  \end{enumerate}
\fincorrection
\correction{002279}
  \begin{enumerate}
  \item
    Notons $P=\sum_{i=0}^{d}a_{i}X^{i}$. Dans le calcul de $P(n+km)$, en
    développant tous les termes $(n+km)^{i}$ à l'aide du binôme, on
    obtient que $P(n+km)=\sum_{0\leq j\leq i\leq
      d}a_{i}C_{i}^{j}n^{j}(km)^{i-j} =P(n)+mN$ où $N=\sum_{0\leq j<i\leq
      d}a_{i}C_{i}^{j}n^{j}(km)^{i-j}-1\in\Zz$. Donc $m|P(n+km)$.

  \item
    Supposons qu'un tel polynôme existe~: soit $m=P(0)$. $\forall
    k\in\Zz, m|P(km)$. Comme $P(km)$ est premier, on en déduit que
    $P(km)=\pm m$. Ceci est en contradiction avec 
    $\lim_{k\to+\infty}P(km)=\pm\infty$.
  \end{enumerate}
\fincorrection
\correction{002280}
 \begin{enumerate}
  \item
    Soit $f=\sum_{i=0}^{n}a_{i}x^{i}\in\Qq[x]$. Soit
    $a_{i}=\frac{p_{i}}{q_{i}}$ le représentant irréductible de $a_{i}$.
    Soit $m=\mathrm{ppcm}(q_{0},\dots,q_{n})$. Notons $m=q_{i}m_{i}$. Alors
    $f=\frac{1}{m}\sum a_{i}m_{i}x^{i}$. En mettant en facteur
    $d=\pgcd(a_{0}m_{0},\dots,a_{n}m_{n})$, on obtient
    $f=\frac{d}{m}f_{0}$, où $f_{0}\in\Zz[x]$ est primitif.

  \item
    Notons $\alpha=\frac{p}{q}$, avec $\pgcd(p,q)=1$ et $q>0$. Soit
    $g_{1}=\alpha g$. On a $qg=pg_{1}$, donc $qc(g)=pc(g_{1})$. On en
    déduit que $q|p$, et donc que $q=1$~:$\alpha\in\Zz$.

  \item
    Soit $g\in\Qq[x]$ tel que $f=dg$. Soit $g=\frac{p}{q}g_{0}$ la
    décomposition de $g$ donnée par la question $1$. Alors $qf=pdg_{0}$
    donc $qc(f)=pc(d)c(g_{0})=p$. Donc $q|p$ et finalement $q=1$. On en
    déduit que $g=pg_{1}\in\Zz[x]$.

  \item
    $d=\pgcd_{\Qq}(f,g)=\frac{p}{q}d_{0}$. Alors $d_{0}$ est primitif et
    divise $f$ et $g$ sur $\Qq$. Donc $d_{0}$ divise $f$ et $g$ sur $\Zz$.

    Soit $h$ un diviseur commun de $f$ et $g$ dans $\Zz[x]$. On a
    $c(h)|c(f)=1$ donc $h$ est primitif. Par ailleurs, $h$ est un
    diviseur commun à $f$ et $g$ dans $\Qq[x]$, donc $h|d_{0}$ dans
    $\Qq[x]$. On en déduit que $h|d_{0}$ dans $\Zz[x]$.

    Ainsi, $d_{0}$ est bien un $\pgcd$ de $f$ et $g$ dans $\Zz[x]$.

  \item
    Soit $d=\pgcd(c(f),c(g))$, $h=\pgcd(f,g)=c(h)h_{0}$,
    $h'=\pgcd(f_{0},g_{0})$.
    
    On a $d|c(f)$, $d|c(g)$, $h'|f_{0}$ et $h'|g_{0}$ donc $dh'|f$ et
    $h'|g$, et donc $dh'|h$.

    $c(h)|c(f)$ et $c(h)|c(g)$ donc $c(h)|d$. $h|f$, donc il existe
    $f_{1}\in\Zz[x]$ tel que  $f=h_{0}c(h)f_{1}$. On a alors
    $c(h)c(f_{1})=c(f)$, et après simplification, on en déduit que
    $f_{0}=h_{0}f'_{1}$, avec $f'_{1}\in\Zz[x]$~: $h_{0}|f_{0}$. De même
    pour $g$~: $h_{0}|g_{0}$. On en déduit que $h_{0}|h'$, et donc que
    $h|dh'$.
  \end{enumerate}
\fincorrection
\correction{002281}
 Soit $K$ un corps, $A$ un anneau non trivial, et $K\xrightarrow{\phi}A$
  un morphisme d'anneaux. Soit $x\in K\setminus\{0\}$. On a
  $1=\phi(1)=\phi(xx^{-1})=\phi(x)\phi(x^{-1})\neq 0$ (car $A$ n'est pas
  l'anneau trivial). Donc $\phi(x)\neq 0$. Ainsi $\ker\phi=\{0\}$, donc
  $\phi$ est injectif.
\fincorrection
\correction{002282}
 Soit $x\in R\setminus\{0\}$. Alors
  $(x)\supset(x^{2})\supset(x^{3})\supset$ est une suite décroissante
  d'idéaux. Elle est donc stationnaire à partir d'un certain rang~:
  $\exists k\in\Nn, (x^{k})=(x^{k+1})$. En particulier, $\exists a\in R,
  k^{k+1}=ax^{k}$. Comme $A$ est intègre, on en déduit que $ax=1$, donc
  $x\in R^{\times}$.

  $R^{\times}=R\setminus\{0\}$ donc $R$ est un corps.
\fincorrection
\correction{002283}
  Soit $A$ un anneau fini, et $I$ un idéal premier. Alors $A/I$ est
  intègre, et fini (!), donc $A/I$ est un corps (voir exercice \ref{ex:bar5}). Donc $I$ est maximal.
\fincorrection
\correction{002284}
  On rappelle que le produit de deux idéaux $I$ et $J$ est l'idéal
  engendré par les produits de la forme $ab$ avec $a\in I$, $b\in J$~:
  $$
   I\cdot J=\{\sum_{i=0}^{N}a_{i}b_{i}, N\in\Nn, a_{i}\in I, b_{i}\in J\}
  $$
  \begin{itemize}
  \item Si $I$ est un idéal premier~: Soient $J$ et  $K$ deux idéaux tels
    que $J\cdot K\subset I$. Alors si $J\not\subset I$, $\exists a\in x
    \setminus I$. Soit $y\in K$. On a $xy\in J\cdot K$ donc $xy\in I$.
    Comme $I$ est premier, $x\in I$ ou $y\in I$. Mais $x\notin I$ donc
    $y\in I$. Ainsi $\forall y\in K, y\in I$~: on a montré que~:
    $J\not\subset I\Rightarrow K\subset I$. On a donc bien $J\subset I$
    ou $K\subset I$.

  \item
    Si $\forall J,K \text{ idéaux},( J\cdot K\subset I\Rightarrow
    J\subset I \text{ ou }K\subset I)$~: Soit $a,b\in A$ avec $ab\in I$.
    Alors $(a)\cdot(b)=(ab)$ donc $(a)\subset I$ ou $(b)\subset I$ et
    donc  $a\in I$ ou $b\in I$. $I$ est donc premier.
  \end{itemize}

  On a $M^{n}=M\cdot M^{n-1}$. Donc si $I$ est premier et contient
  $M^{n}$ alors $I$ contient $M$ ou $M^{n-1}$, et par une récurrence
  finie, on obtient que $I$ contient $M$. Ainsi~: $M\subset I\subsetneq
  A$. Comme $M$ est maximal on en déduit que $M=I$.
\fincorrection
\correction{002285}
 \begin{itemize}
  \item   $A[X]/(X)$~: 
    $X$ est unitaire donc on dispose de la division euclidienne par $X$.
    On vérifie (comme dans le cours) que chaque classe a un et un seul
    représentant de degré $0$. On en déduit que $A[X]/(X)$ est en
    bijection avec $A$. Il reste alors à remarquer que cette bijection
    est un morphisme d'anneaux. 

    Une autre façon de dire la même chose est de remarquer que
    l'application $\phi:A[X]\to A$, $P\mapsto P(0)$ est un morphisme
    d'anneaux. $\ker\phi=(X)$ et $\mathrm{Im}\,\phi=A$. Comme $A/\ker\phi\sim
    \mathrm{Im}\,\phi$, on a bien $A[X]/(X)\sim A$.

  \item
    On peut considérer $\phi:A[X,Y]\to A[Y]$, $P\mapsto P(0,Y)$. C'est un
    morphisme d'anneaux. En séparant les termes ne dépendant que de $Y$
    des autres, on peut mettre tout polynôme $P$ de $A[X,Y]$ sous la
    forme $P=P_{1}(Y)+XP_{2}(X,Y)$ où $P_{1}\in A[Y]$ et $P_{2}\in
    A[X,Y]$. Alors $\phi(P)=0$ ssi $P_{1}=0$, ssi $P=XP_{2}$, c'est à
    dire $P\in(X)$. Ainsi $\ker\phi=(X)$. Par ailleurs, tout polynôme $P$
    de $A[Y]$ peut être vu comme un polynôme $\tilde{P}$ de $A[X,Y]$.
    Alors $P=\phi(\tilde{P})$, donc $\mathrm{Im}\,\phi=A[Y]$. Finalement~:
     $A[X,Y]/(X)\sim A[Y]$.
    
   \item $A[X,Y]/(X,Y)$~:
     Soit $\phi:A[X,Y]\to A$, $P\mapsto P(0,0)$. $\phi$ est un morphisme
     d'anneaux, et avec les notations précédentes, pour
     $P=P_{1}(Y)+XP_{2}(X,Y)$, avec $\phi(P)=0$, on a $P_{1}(0)=0$, donc
     $Y|P_{1}(Y)$. Ainsi, $P$ est la somme de deux polynômes, l'un
     multiple de $X$, l'autre multiple de $Y$ donc $P\in(X,Y)$.
     Réciproquement, si $P\in(X,Y)$, alors $P(0,0)=0$. Donc
     $\ker\phi=(X,Y)$. $\forall a\in A \phi(a)=a$ donc $\phi$ est
     surjective. Finalement $A[X,Y]/(X,Y)\sim A$.
     
   \item $A[X_{1},\dots,X_{n}]/(X_{1},\dots,X_{n})$~:
     Soit $\phi~:A[X_{1},\dots,X_{n}]\to A$, $P\mapsto P(0)$. $\phi$ est
     un morphisme d'anneaux. En regroupant tous les termes dépendant de
     $X_{n}$, puis tous les termes restant dépendant de $X_{n-1}$, et
     ainsi de suite jusqu'aux termes dépendant seulement de $X_{1}$, et
     enfin le terme constant, tout polynôme $P\in A[X_{1},\dots,X_{n}]$
     peut se mettre sous la forme
     $P=X_{n}P_{n}+X_{n-1}P_{n-1}+\dots+X_{1}P_{1}+p_{0}$, avec $P_{i}\in
     A[X_{1},\dots,X_{i}]$ (et $p_{0}\in A$). On en déduit que
     $\ker\phi=(X_{1},\dots,X_{n})$. Par ailleurs $\forall a\in A,
     \phi(a)=a$, donc $A[X_{1},\dots,X_{n}]/(X_{1},\dots,X_{n})\sim A$.
  \end{itemize}
  
  Comme un idéal est premier (resp. maximal) ssi le quotient est intègre
  (resp. un corps), on en déduit que
  \begin{itemize}
  \item dans $A[X]$, $(X)$ est premier ssi $A$ est intègre, maximal ssi $A$ est un corps,
  \item dans $A[X,Y]$, $(X)$ est premier ssi $A$ est intègre, et n'est jamais maximal,
  \item dans $A[X_{1},\dots,X_{n}]$, $(X_{1},\dots,X_{n})$ est premier
    ssi $A$ est intègre, maximal ssi $A$ est un corps.
  \end{itemize}
\fincorrection
\correction{002286}
  Soit $\alpha=a+b\sqrt{d}\in\Zz[\sqrt{d}]$. Soit $a=mp+a'$ la division
  euclidienne de $a$ par $m$, et $b=mq+b'$ celle de $b$ par $m$. Alors
  $\alpha=m(p+q\sqrt{d})+a'+b'\sqrt{d}$. On en déduit que chaque classe
  du quotient $\Zz[\sqrt{d}]/(m)$ a un représentant dans
  $$
    \mathcal{C}=\Big\{a+b\sqrt{d},\  (a,b)\in\{0,\dots,m-1\}^{2}\Big\}
  $$
  Par ailleurs si deux éléments $a+b\sqrt{d}$ et $a'+b'\sqrt{d}$ de cet
  ensemble sont dans la même classe, alors $\exists c,d\in\Zz, \
  a+b\sqrt{d}=(a'+b'\sqrt{d})+m(c+d\sqrt{d})$. On en déduit que $a=a'+mc$
  et $b=b'+md$, et donc $a=a'$, $b=b'$.
 
  Ainsi chaque classe de $\Zz[\sqrt{d}]/(m)$ a un représentant unique
  dans $\mathcal{C}$. $\Zz[\sqrt{d}]/(m)$ et $\mathcal{C}$ sont donc en
  bijection~: en particulier, $\Zz[\sqrt{d}]/(m)$ a $m^{2}$ éléments.
  

  \emph{Remarque}~: on a
  $$
    \Zz[\sqrt{d}]\sim \Zz[X]/(X^{2}-d).
  $$
  En effet l'application $\phi:\Zz[X]/(X^{2}-d)\to\Zz[\sqrt{d}]$,
  $\bar{P}\mapsto P(\sqrt{d})$ est bien définie (si $\bar(P)=\bar{Q}$,
  alors $P(\sqrt{d})=Q(\sqrt{d})$), et c'est un morphisme d'anneaux. De
  plus, si $\phi(P)=0$, notons $P=Q(X^{2}-d)+(aX+b)$ la division
  euclidienne de $P$ par $X^{2}-d$. En évaluant en $\sqrt{d}$, on a
  $a\sqrt{d}+b=0$ donc $R=0$. On en déduit que $(X^{2}-d)|P$, i.e.
  $\bar{P}=0$. On en déduit que $\ker\phi=\{0\}$, donc $\phi$ est
  injective. Par ailleurs $\forall(a,b)\in\Zz^{2},
  \phi(a+bX)=a+b\sqrt{d}$ donc $\phi$ est surjective.

  \medskip

  Si $d$ est pair, comme $\sqrt{d}\cdot\sqrt{d}=|d|\in(2)$ alors que
  $\sqrt{d}\notin(2)$, $(2)$ n'est pas premier.

  Si $d$ est impair~: $(1+\sqrt{d})(1+\sqrt{d})=(1+d)+2\sqrt{d}\in(2)$,
  mais $(1+\sqrt{d})\notin(2)$ donc $(2)$ n'est pas premier.

  \emph{Remarque}~: $\Zz[\sqrt{d}]/(2)\sim\Zz_{2}[X]/(X^{2}+\bar{d})$.
  $(X^{2}+\bar{d})$ est $X^{2}$ ou $X^{2}+1$. Aucun de ces deux polynômes
  n'est irréductible. Donc le quotient ne saurait être intègre.
\fincorrection
\correction{002287}
  \begin{itemize}
  \item 
    Si $x\in A$ est premier~: soit $a,b\in A$ tels que $ab=x$. Alors
    $ab\in(x)$ donc $a\in(x)$ ou $b\in(x)$. On en déduit que $a\sim x$ ou
    $b\sim x$. Donc $x$ est irréductible.

  \item
    $A$ est supposé factoriel. Soit $I$ un idéal premier. Soit $x\in I$
    et $x=p_{1}\dots p_{k}$ ``la'' factorisation de $x$ en produit
    d'irréductibles. Alors $(p_{1}\cdots p_{n-1})p_{n}\in I$ donc
    $(p_{1}\cdots p_{n-1})\in I$ ou $p_{n}\in I$. si $p_{n}$ in $I$, $I$
    contient un irréductible. Sinon, $(p_{1}\cdots p_{n-2})p_{n-1}\in I$.
    Par une récurrence finie, l'un au moins des $p_{i}\in I$, donc $I$
    contient un irréductible.

  \item
    Dans $\Zz[\sqrt{-5}]$, $9\in(3)$. Pourtant
    $9=(2+\sqrt{-5})(2-\sqrt{-5})$ et $(2\pm\sqrt{-5})\notin(3)$. Donc
    $(3)$ n'est pas premier.

  \item
    $2$ est irréductible~: $2=z_{1}z_{2}$ avec $z_{i}\in\Zz[\sqrt{-5}]$,
    alors $|z_{1}|^{2}|z_{2}|^{2}=4$, donc
    $\{|z_{1}|^{2},|z_{2}|^{2}\}=\{1,4\}$ ou $\{2,2\}$. Dans le premier
    cas, on a affaire à une factorisation triviale. Le second est
    impossible, puisque l'équation $a^{2}+5b^{2}=2$ n'a pas de solution
    entière $(a,b)$.

    Par ailleurs, $(1+\sqrt{-5})(1+\sqrt{-5})=6\in(2)$, mais
    $(1\pm\sqrt{-5})\notin(2)$ donc $2$ n'est pas premier dans
    $\Zz[\sqrt{-5}]$.
  \end{itemize}
\fincorrection
\correction{002288}
 \begin{enumerate}
  \item 
    Soit $\mathcal{J}$ un idéal de $A/I$. Soit $\pi$ la projection
    canonique $A\to A/I$, et $J=\pi^{-1}(\mathcal{J})$. $J$ est un idéal
    de $A$ qui est principal donc $\exists a\in A, J=(a)$. Montrons que
    $\mathcal{J}=(\pi(a))$.

    On a $\pi(a)\in\mathcal{J}$ donc $(\pi(a))\subset\mathcal{J}$. Soit
    $\alpha\in\mathcal{J}$, et $b$ un représentant de $\alpha$, i.e.
    $b\in A$ et $\pi(b)=\alpha$. Alors $b\in J=(a)$, donc $\exists k\in
    A, b=ka$. Alors $\pi(b)=\pi(ka)=\pi(k)\pi(a)$, donc
    $\pi(b)\in(\pi(a))$. Donc $\mathcal{J}\subset(\pi(a))$.

    Finalement, $\mathcal{J}=(\pi(a))$. On en déduit que $A/I$ est
    principal.

  \item
    \begin{itemize}
    \item $\Zz/n\Zz$~:
      Soit $I$ un idéal de $\Zz/n\Zz$. $I$ est principal, donc $\exists
      a\in\Zz, I=(\bar{a})$. Or $(\bar{a}) =
      \{\alpha\bar{a},\alpha\in\Zz/n\Zz\} = \{\bar{p}\bar{a},p\in\Zz\} =
      \{\overline{pa},p\in\Zz\}$. Donc $\pi^{-1}(I)=
      \{pa+qn,(p,q)\in\Zz^{2}\}$ est l'idéal engendré sur $\Zz$ par $a$
      et $n$ donc l'idéal engendré par $d=(\pgcd(n,a))$. On en déduit que
      $I=(\bar{d})$. En particulier, $I$ est engendré par un diviseur de
      $n$.

      Soit maintenant $d_{1}$ et $d_{2}$ deux diviseurs (positifs) de $n$
      tels que $(\bar{d_{1}})=(\bar{d_{2}})$. On a
      $\pi^{-1}((d_{1}))=d_{1}\Zz=d_{2}\Zz$ donc $d_{1}=d_{2}$.

      Ainsi, les idéaux de $\Zz/n\Zz$ sont engendrés par les diviseurs de
      $n$, et deux diviseurs distincts engendrent deux idéaux distincts~:
      il y a donc autant d'idéaux dans $\Zz/n\Zz$ que de diviseurs de $n$.

    \item $\Qq[X]/(f)$~:
      On raisonne de la même manière~: la remarque clef étant si
      $I=(\bar{g})$ est un idéal de $\Qq[X]/(f)$, alors
      $\pi^{-1}(I)=(f,g)=(\pgcd(f,g))$.
    \end{itemize}

  \item
    Les idéaux maximaux sont ceux pour lesquels le quotient est un corps,
    (donc aussi ceux pour lesquels le quotient est intègre puisque
    $\Zz/n\Zz$ est fini). On a le diagramme suivant ($I=(\bar{d})$)~:
    $$    
    \xymatrix{%
       \Zz      \ar[r]^{\pi_{1}}\ar[d]^{\pi}\ar@(ur,ul)[rr]^{\pi_{2}\circ\pi_{1}} & 
       \Zz/n\Zz \ar[r]^{\pi_{2}}&(\Zz/n\Zz)/I \\ 
       \Zz/d\Zz \ar[urr]_{\sim}
     }%
     $$
     En effet, $\pi_{1}$ et $\pi_{2}$ sont des morphismes d'anneaux, et
     $\ker(\pi_{2}\circ\pi_{1})=d\Zz$. Donc $(\Zz/n\Zz)/I$ est un corps
     ssi $d$ est premier. 
     
     De même, $(\Qq[X]/(f))/I$ est un corps ssi $I=(\bar{g})$ où $g$ est
     un facteur premier de $f$.
     
  \end{enumerate}
\fincorrection
\correction{002289}
 \begin{enumerate}
  \item 
     Soit $\alpha,\beta\in\bar{J}$ et $\lambda,\mu\in A/I$. Alors
     $\exists a,b\in J$, $l,m\in A$,
     $\alpha=\pi(a),\beta=\pi(b),\lambda=\pi(l),\mu=\pi(m)$. On a donc
     $\lambda\alpha+\mu\beta=\pi(la+mb)$. Or $la+mb\in J$ (car $J$ est un
     idéal), donc $\lambda\alpha+\mu\beta\in\bar{J}$. Donc $\bar{J}$ est
     un idéal de $A/I$.

   \item
     Comme dans l'exercice \ref{ex:bar49}, on a le diagramme suivant~:
     $$    
     \xymatrix{%
       A        \ar[r]^{\pi_{1}}\ar[d]^{\pi}\ar@(ur,ul)[rr]^{\pi_{2}\circ\pi_{1}} & 
       A/I      \ar[r]^{\pi_{2}}&(A/I)/\bar{J} \\ 
       A/(I+J) \ar[urr]_{\sim}
      }%
      $$
      En effet, si $x\in\ker(\pi_{2}\circ\pi_{1})$, alors
      $\pi_{1}(x)\in\ker\pi_{2}=\bar{J}$, donc $\exists y\in A,
      \pi_{1}(x)=\pi_{1}(y)$. Alors $x-y\in\ker\pi_{1}=I$, donc $\exists
      z\in I, x=y+z$~: on a donc $x\in I+J$. Réciproquement, si $x\in
      I+J$, alors $\exists (x_{1},x_{2})\in I\times J, x=x_{1}+x_{2}$.
      Alors $\pi_{1}(x)=\pi_{1}(x_{2})\in\bar{J}$, donc
      $\pi_{2}\circ\pi_{1}(x)=0$.

      Donc $\ker(\pi_{2}\circ\pi_{1})=I+J$. Donc
      $A/(I+J)\sim(A/I)/\bar{J}$.
  \end{enumerate}
\fincorrection
\correction{002290}
  \begin{enumerate}
  \item 
    Soit $J\subset B$ un idéal premier de $B$. Soient $a,b\in A$ tels que
    $ab\in f^{-1}(J)$. Alors $f(a)f(b)=f(ab)\in J$ donc $f(a)\in J$ ou
    $f(b)\in J$. Ainsi, $a\in f^{-1}(J)$ ou $b\in f^{-1}(J)$. On en
    déduit que $f^{-1}(J)$ est premier.

    Cette proposition n'est pas vraie pour les idéaux maximaux. Par
    exemple, $A=\Zz$, $B=\Qq[X]$, $f(k)=k$, et $J=(X)$. Alors
    $f^{-1}(J)=\{0\}$ n'est pas maximal.

  \item
    Prenons $A=\Zz$, $B=\Qq$, $f(k)=k$. $f(\Zz)=\Zz$ n'est pas un idéal
    de $\Qq$ ($1\in\Zz$, $\frac{1}{2}\in\Qq$ et pourtant
    $1\times\frac{1}{2}\notin\Zz$)

    Supposons $f$ surjectif. Soit $x,y\in f(I)$, $a,b\in B$. Il existe
    $x_{0},y_{0}\in I$ tels que $x=f(x_{0})$ et $y=f(y_{0})$. De plus,
    comme $f$ est surjectif, $\exists a_{0},b_{0}\in A$ tels que
    $a=f(a_{0})$ et $b=f(b_{0})$. Alors $ax+by = f(a_{0})f(x_{0})
    +f(b_{0})f(y_{0}) = f(a_{0}x_{0}+b_{0}y_{0})$ et comme $I$ est un
    idéal, $(a_{0}x_{0}+b_{0}y_{0})\in I$, donc $(ax+by)\in f(I)$.

    $f(I)$ est donc bien un idéal de $B$.

  \item
    Soit $I$ un idéal maximal de $A$ et $J=f(I)$. Supposons $J\neq B$.
    Soit $K$ un idéal de $B$ tel que $J\subset K$. Alors $I\subset
    f^{-1}(K)$, donc $f^{-1}(K)=I$ ou $f^{-1}(K)=A$. Dans le premier cas,
    on $K=f(f^{-1}(K))=J$, dans le second cas, on a
    $K=f(f^{-1}(K))=f(A)=B$. L'idéal $J$ est donc maximal.

  \item
    $(X+2)(X+3)=X^{2}+5X$ dans $\Zz_{6}[X]$, donc $(X+\bar{2})(X+\bar{3})\in
    (X)$, mais $(X+\bar{2})\notin(X)$ et $(X+\bar{3})\notin(X)$, donc
    $r_{6}((X))$ n'est pas premier dans $\Zz_{36}[X]$.

    $(X+1)^{2}=(X^{2}+1)$ dans $\Zz_{2}[X]$, or $(X+1)\notin(X^{2}+1)$,
    donc $r_{2}((X^{2}+1))$ n'est pas premier dans $\Zz_{2}[X]$.
  \end{enumerate}
\fincorrection
\correction{002291}
  \begin{enumerate}
  \item 
    Soit $J=B\cap I$. Soit $x,y\in J$, $a,b\in B$, alors $ax+by\in B$
    puisque $B$ est un sous-anneau de $A$. $ax+by\in I$ puisque $I$ est
    un idéal. On en déduit que $J$ est un idéal.

    $B+I$ est stable par addition (car $B$ et $I$ le sont). Soit
    $\alpha=a+x\in B+I$ et $\beta= b+y\in B+I$. Alors
    $\alpha\beta=(ab)+(ay+bx+xy)\in B+I$, donc $B+I$ est stable par
    multiplication. $1\in B+I$, donc $B+I$ est un sous anneau de $A$.
    $I\subset B+I$, et $I$ est absorbant pour la multiplication dans $A$,
    donc aussi dans $B$~: $I$est un idéal de $B+I$.

  \item
    On a le diagramme (de morphismes d'anneaux) suivant~:
    $$
    \xymatrix{%
      B\ar[r]^{i} 
       \ar[d]^{{\pi_{0}}}
       \ar@(ur,ul)[rr]^{\phi}
      & B+I \ar[r]^{\pi} & (B+I)/I\\ 
      B/\ker\phi\ar[urr]_{\sim}
    }%
    $$
    Or, pour $x\in B$, on a~: $x\in\ker\phi \Leftrightarrow x=i(x)\in
    \ker\pi=I$. Donc $\ker\phi=B\cap I$, et par suite~:
    $$
    B/(B\cap I)\sim (B+I)/I.
    $$
  \end{enumerate}
\fincorrection
\correction{002292}  
\begin{enumerate}
  \item 
    Soit $P=x^{3}-x+2$. Sa réduction  $\bar{P}=x^{3}-x-1$ modulo $3$ est
    de degré $3$ et n'a pas de racine, donc $\bar{P}$ est irréductible
    dans $\Zz_{3}[x]$. Comme $P$ est primitif, on en déduit que $P$ est
    irréductible dans $\Zz[x]$, puis dans $\Qq[x]$. Comme $\Qq[x]$ est
    principal, on en déduit que $(P)$ est maximal, et donc que
    $\Qq[x]/(P)$ est un corps.

  \item
    Dans $\Qq[x]/(P)$, on a $y^{3}-y+2=0$, donc $y(y^{2}-1)=-2$ et
    finalement $y(\frac{1}{2}(1-y^{2}))=1$. Ainsi $y^{-1}=\frac{1}{2}(1-y^{2})$.

  \item 
    $1+y+y^{2}=\pi(1+x+x^{2})$. On a $\pgcd(P,1+x+x^{2})=1$, et plus
    précisément, en utilisant l'algorithme d' Euclide~:
    $13=(x+4)P-(x^{2}+3x-5)(x^{2}+x+1)$ donc
    $(y^{2}+y+1)^{-1}=\frac{-1}{13}(y^{2}+3y-5)$.
\end{enumerate}
\fincorrection
\correction{002293}
  Notons $f=\sum_{i=0}^{d} a_{i}x^{i}$. On a
  $\pgcd(a_{0},\dots,a_{d})\sim1$ et $\pi\!\!\!\not| a_{d}$. Notons
  $\bar{f}\in A/(\pi)[X]$ la réduction de $f$ modulo $\pi$. Soit $f=gh$
  une factorisation de $f$ dans $A[x]$. Alors $\bar{f}=\bar{g}\bar{h}$,
  et donc (quitte à échanger $g$ et $h$) $\bar{g}\sim 1$ et $\bar{h}\sim
  \bar{f}$. Comme $\pi\!\!\!\not| a_{d}$, on a $\deg(\bar{f})=d$, et donc
  $\deg(\bar{h})=d$ puis $\deg(h)\geq d$, et finalement $\deg(h)=d$. Par
  conséquent $\deg(g)=0$~: $g\in A$. Comme $g|f$, on a $g|c(f)\sim 1$
  donc $g\sim 1$. Ainsi, toute  factorisation de $f$ dans $A[x]$ est
  triviale~: $f$est irréductible.

\fincorrection
\correction{002294}
  \begin{enumerate}
  \item 
    Ce polynôme est unitaire donc primitif. $11$ est nombre premier qui
    divise tous les coefficients sauf le dominant. $11^{2}=121$ ne divise
    pas le coefficient de degré $0$, donc, d'après le critère
    d'Eisenstein, c'est un polynôme irréductible de $\Qq[X]$.

  \item
    $f(X,Y)=(X^2+1)Y^3+(X-1)^{2}Y^2+(X-1)$. Regardons $f$ comme un
    polynôme de $A[Y]$ avec $A=\Cc[X]$. Alors, $f$ est primitif sur $A$, et
    $(X-1)$ est un irréductible de $A$ qui divise tous les
    coefficients de $f$ sauf le dominant, et dont le carré ne divise pas
    le terme constant. D'après le critère d'Eisenstein, on en déduit que
    $f$ est irréductible dans $A[Y]=\Cc[X,Y]$.

    Dans $\Zz_{2}[X,Y]$, on a $(X^{2}+1)=(X+1)^{2}$ et
    $f=(X+1)((X+1)(Y^3+Y^2)+1)$, donc $f$ n'est pas irréductible..

  \item
    $f(X,Y)=Y^7+Y^6+7Y^4+XY^3+3X^2Y^2-5Y+X^2+X+1$. Considérons $f$ comme
    un polynôme de $A[X]$ où $A=\Qq[Y]$. Alors $f$ est primitif sur $A$.
    Soit $\pi=Y\in A$. $\pi$ est irréductible, $\pi$ ne divise pas le
    coefficient dominant de $f$, et la réduction $\bar{f}$ modulo $\pi$
    est $\bar{f}=X^{2}+X+1\in A/(\pi)[X]=\Qq[X,Y]/(Y)\simeq\Qq[X]$.
    $\bar{f}$ est donc irréductible dans $A/(\pi)$, donc d'après
    l'exercice précédent, $f$ est irréductible dans $\Qq[X,Y]$.
  \end{enumerate}

\fincorrection
\correction{002295}
  Soit $f=x^{2}+y^{2}+1\in A[x,y]$ ($A=\Cc,\Rr,\Qq,\Zz,\Zz_{2}$). Soit
  $B=A[y]$, et regardons $f$ comme un polynôme de $B[x]$. Le coefficient
  dominant de $f$ (qui est $1$) est inversible dans $B$, donc on peut
  effectuer la division euclidienne de tout polynôme par $f$~: $\forall
  g\in B[y], \exists(q,r)\in B[x]^{2}, g=qf+r$ et $\deg_{x}r\leq 1$.
  Notons $r=a(y)x+b(y), a,b\in A[y]$. De plus, pour des raisons de degré,
  le quotient et le reste de cette division sont uniques. On peut donc
  identifier $A[x,y]/(x^{2}+y^{2}+1)$ à $\{a(y)x+b(y),\ a(y),b(y)\in
  A[y]\}$. Supposons que $\bar{y}$ soit inversible dans cet quotient. Il
  existe $a,b\in A[y]$ tels que  $y(a(y)x+b(y))=\bar{1}$. On a donc
  $ya(y)=0$ et $yb(y)=1$, ce qui est impossible.
\fincorrection
\nocorrection
\correction{002297}
  Rappelons que $(a)\cdot(b)=\{\sum_{i=1}^{n}a_{i}b_{i}, n\in\Nn,
  a_{i}\in(a), b_{i}\in(b)\}=(ab)$. De plus $(ab)\subset(a)\cap(b)$ donc 
  \begin{align*}
    (ab)=(a)\cap(b)
    &\Leftrightarrow (a)\cap(b)\subset(ab)\\
    &\Leftrightarrow \forall m\in A,\ (a|m\text{ et }b|m\Rightarrow ab|m)\\
    &\Leftrightarrow \mathrm{ppcm}(a,b)\sim ab\\
    &\Leftrightarrow \mathrm{ppcm}(a,b)\sim \pgcd(a,b)\mathrm{ppcm}(a,b)\\
    &\Leftrightarrow \pgcd(a,b)\sim 1
  \end{align*}

  Si $A$ est principal, alors $\exists d\in A,\ (a,b)=(d)$. Alors
  $a\in(d)$ et $b\in(d)$ donc $d$ est un diviseur commun à $a$ et $b$. Si
  de plus $d'$ est un autre diviseur commun à $a$ et $b$, alors
  $a\in(d')$ et $b\in(d')$ et comme $(a,b)$ est le plus petit idéal
  contenant $a$ et $b$, on en déduit que $(a,b)=(d)\subset(d')$, et donc
  que $d'|d$~: finalement, $\pgcd(a,b)=d$.

\fincorrection
\correction{002298}
  \begin{enumerate}
  \item 
    $I=(5,x^{2}+3)$. On a $\pgcd(5, x^{2}+3)=1$, donc si $I$ était
    principal, on aurait $1\in I$, et donc $I=\Zz[X]$. Si $1\in I$, il
    existe $P,Q\in\Zz[x]$,  tels que $1=5P+(x^{2}+3)Q$. En considérant la
    réduction modulo $5$ de ces polynômes, on obtient
    $(x^{2}+\bar{3})\bar{Q}=\bar{1}$, ce qui est impossible pour des
    raisons de degré ($\Zz/5\Zz$ est intègre). Donc $1\notin I$, et $I$
    n'est donc pas intègre.

    \smallskip
    $x^{2}+1=(x+2)(x-2)+5$, donc $(x^{2}+1,x+2)=(x+2,5)$. Or $(x+2,5)$
    n'est pas principal pour les mêmes raisons que précédemment.

    \smallskip
    On a $(x-1)=(x^{4}-1)-x(x^{3}-1)$ donc $(x-1)\subset
    (x^{4}-1,x^{3}-1)$. Par ailleurs, $(x-1)|(x^{4}-1)$ et
    $(x-1)|(x^{3}-1)$ donc $x^{4}-1\in(x-1)$ et $x^{3}-1\in(x-1)$, donc
    $(x^{4}-1,x^{3}-1)\subset (x-1)$. Donc $(x^{4}-1,x^{3}-1)$ est
    principal.

  \item
    $I=(x,x+1)=\Zz$ car $1=(x+1) - x$. Donc $I$ n'est pas propre.

    $I=(5,x^{2}+4)$. $\Zz[X]/I\sim\Zz_{5}/(x^{2}+\bar{4})$. Mais
    $(x^{2}+\bar{4})=(x-\bar{1})(x+\bar{1})$ est réductible dans
    $\Zz_{5}[x]$, donc $\Zz_{5}/(x^{2}+\bar{4})$ n'est pas intègre~: $I$
    n'est pas premier.

    $I=(x^{2}+1,x+2)=(x+2,5)$. $\Zz[x]/I\simeq\Zz_{5}[x]/(x+\bar{2})$.
    $x+\bar{2}$ est irréductible dans $\Zz_{5}[x]$, qui est principal,
    donc $(x+\bar{2})$ est maximal, donc le quotient est un corps, et $I$
    est maximal.
    
  \end{enumerate}
\fincorrection
\correction{002299}  
\begin{enumerate}
  \item 
    Soit $a,b\in B,\ ab\in I\cap B$. Alors $ab\in I$ donc $a\in I$ ou
    $b\in I$. Comme $a,b\in B$, on a $a\in I\cap B$ ou $b\in I\cap B$.
    Donc, si $I\cap B$ est propre, $I\cap B$ est premier.

  \item
    Soit $J$ un idéal premier de $\Zz[X]$. Alors $J\cap\Zz$ est soit
    $\Zz$ soit un idéal premier de $\Zz$. Si $J\cap\Zz=\Zz$, alors $1\in
    J$, et donc $J=\Zz[X]$, ce qui est exclu. On en déduit que $J=(0)$ ou
    $J=(p)$ avec $p$ premier.

  \item
    On suppose $J\cap\Zz=(0)$ et $J\neq(0)$. Soit alors $f$ un polynôme
    de $J\setminus\{0\}$ de degré minimal. Notons $f=c(f)f_{0}$ où
    $f_{0}\in\Zz[x]$ est primitif. Comme $J$ est premier, on a $c(f)\in
    J$ ou $f_{0}\in J$. Comme $J\cap\Zz=\{0\}$, le premier cas est exclu,
    donc $f_{0}\in J$.

    Soit maintenant $g\in J$. Soit $g=f_{0}q+r$ la division euclidienne de
    $g$ par $f_{0}$ dans $\Qq$ ($q,r\in\Qq[x]$). Notons $q=\frac{a}{b}q_{0}$
    avec $q_{0}\in\Zz[x]$ primitif, et $r=\frac{a'}{b'}r_{0}$, avec
    $r_{0}\in\Qq[x]$ primitif.

    Alors $bb'g=ab'\,q_{0}f_{0}+a'b\, r_{0}$ On en déduit que
    $a'b\,r_{0}\in J$, et pour des raisons de degré, $r_{0}=0$.
    Finalement, $bb'g=ab'\,q_{0}f_{0}$, et en considérant les contenus,
    on en déduit que $bb'|ab'$, donc $b|a$, et donc $q\in\Zz[x]$. On en
    déduit que $g\in(f_{0})$, et finalement $J=(f_{0})$.

  \item
    On suppose que $J\cap\Zz=(p)$. Soit $r_{p}$ la projection
    $\Zz[x]\to\Zz_{p}[x]$. Soit $\alpha,\beta\in\Zz_{p}[x]$ tels que
    $\alpha\beta\in r_{p}(J)$. Soit $f,g$ des représentants de $\alpha$
    et $\beta$ (i.e. $r_{p}(f)=\alpha$, $r_{p}(g)=\beta$). Alors $fg\in
    r_{p}^{-1}(r_{p}(J))=J+(p)=J$. Donc $f\in J$ ou $g\in J$, et donc
    $\alpha\in r_{p}(J)$ ou $\beta\in r_{p}(J)$~: $r_{p}(J)$ est premier.
    
    $\Zz_{p}[x]$ est principal, donc il existe un polynôme $\pi$
    irréductible dans $\Zz_{p}[x]$ tel que $r_{p}(J)=(\pi)$. Soit $g$ un
    représentant de $\pi$. Alors $J=(p,g)$~: en effet, on a vu que
    $J=r_{p}^{-1}((\pi))$ et $r_{p}^{-1}((\pi))=(g)+(p)=(p,g)$.

  \item
    Supposons $J$ maximal dans $\Zz[x]$. $J$ est en particulier premier,
    donc a une des deux formes ci dessus. Supposons $J=(f)$, avec $f$
    irréductible et primitif. Soit $p$ un nombre premier ne divisant pas
    le coefficient dominant de $f$. Alors $J\subset(p,f)\subset\Zz[x]$,
    mais $(p,f)\neq\Zz[x]$. En effet, sinon, il existerait $g,h\in\Zz[x]$
    tels que $1=pg+fh$, et en considérant la réduction modulo $p$,
    $\bar{f}$ serait inversible dans $\Zz_{p}[x]$~: comme
    $\deg{\bar{f}}>0$, c'est impossible. On en déduit que $J$ n'est pas
    maximal.

    $J$ est donc de la forme $(p,g)$, avec $r_{p}(g)$ irréductible dans
    $\Zz_{p}[x]$. 
  \end{enumerate}
\fincorrection
\nocorrection
\nocorrection
\nocorrection
\nocorrection
\nocorrection
\nocorrection
\nocorrection
\nocorrection
\nocorrection
\nocorrection
\nocorrection
\nocorrection
\nocorrection
\nocorrection
\nocorrection
\nocorrection
\nocorrection
\nocorrection
\nocorrection
\nocorrection
\nocorrection
\nocorrection
\nocorrection
\nocorrection
\nocorrection
\nocorrection
\nocorrection
\nocorrection
\nocorrection
\nocorrection
\nocorrection
\nocorrection
\nocorrection
\nocorrection
\nocorrection
\nocorrection
\nocorrection
\nocorrection
\nocorrection
\nocorrection
\nocorrection
\nocorrection
\nocorrection
\nocorrection
\nocorrection
\nocorrection
\nocorrection
\nocorrection
\nocorrection
\nocorrection
\nocorrection
\nocorrection
\nocorrection
\nocorrection
\nocorrection
\nocorrection
\nocorrection
\correction{007726}
\begin{enumerate}
 \item Si $q$ et $Q$ sont équivalents, il existe un isomorphisme linéaire $u$ de $E$ tel que $Q=q\circ u$. Par conséquent $Q$ prend toutes les valeurs que prend $q$.
 \item On sait dans ce cas, que le vecteur isotrope est dans un plan hyperbolique, sur lequel $q$ est équivalente à la forme $Q(x,y)=xy$ qui prend toutes les valeurs de $k$.
 \item On procède par récurrence comme dans le cours en utilisant le fait que l'orthogonal d'un espace non singulier dan un espace non-singulier est non-singulier et on s'arrête quand la forme n'a plus de vecteurs isotropes non nuls.
 \item Si dans une telle décomposition, il y a $k$ plans hyperboliques, et dans une autre $k'$, si $k\leq k'$, par le théorème de Witt, l'isométrie entre $k$ des plans hyperboliques des deux décompositions induit une équivalence des orthogonaux. En particulier, ils n'ont pas de vecteurs isotropes. Donc $k'=k$. A l'aide de la décomposition, on peut facilement construire un sous-espace totalement isotrope $I$ de dimension $k$. Ce nombre $k$ est donc inférieur à l'indice de $q$. 
 On injecte $I$ dans un sous-espace totalement isotrope de dimension maximale $M$. Par le théorème de Witt, les orthogonaux sont isomorphes et ne contiennent donc pas de vecteurs isotropes. Donc, $I$ est de dimension maximale, et $k=\nu(q)$.  
\end{enumerate}
\fincorrection
\correction{007727}
\begin{enumerate}
 \item Il faut montrer que $\sigma$ est un morphisme de corps,
 puis que c'est une bijection. La bijection inverse est alors automatiquement un morphisme de corps.
 Pour le premier point, $\sigma$ est facilement multiplicative et $\sigma (1)=1$. Pour montrer que $\sigma (a+b)=\sigma (a)+\sigma(b)$, on utilise la formule du binôme puis le fait que le corps $\mathbb{F}_5$ est de caractéristique $5$.
 Pour le second point, il suffit de montrer que $\sigma$ est injective puisque son ensemble de départ a le même cardinal fini que son ensemble d'arrivée. Pour cela, il suffit de dire que $\sigma$ est un morphisme de corps.
 \item Les formes $f_1$ et $f_3$ sont linéaires par rapport à la première variable, et en utilisant l'expression 
 $$f_1 \left(\begin{array}{c}x\\ y\\ z\end{array}\right),
 \left(\begin{array}{c}x'\\ y'\\ z'\end{array}\right)
 )
 =x\sigma(x')+3z\sigma(y')+3y\sigma(z')$$
 on montre aussi qu'elles sont $\sigma$-semi linéaire par rapport à la seconde variable. Elles sont donc sesquilinéaires.
 Par contre, comme il n'y a au plus que cinq élément de $\mathbb{F}_{25}$ qui vérifient $\lambda^5=\lambda$, il existe un $\lambda\in\mathbb{F}_{25}$ tel que $f_2(\lambda\cdot(1,1,1),(1,1,1))=3\lambda^5
 \not=\lambda\cdot f_2((1,1,1),(1,1,1))$.
 La forme $f_2$ n'est donc pas sesquilinéaires.
 \item Par un théorème du cours, c'est le rang qui classifie les formes sesquilinéaires sur les corps finis. Comme $f_1$ et $f_3$ ont même rang $3$, elles sont équivalentes.  
\end{enumerate}
\fincorrection
\correction{007728}
\begin{enumerate}
 \item La restriction d'une forme alternée sur une droite est nulle.
 Les droites sont donc toutes isométriques. Par le théorème de Witt, toute isométrie entre deux droites se prolonge à l'espace $E$ non dégénéré en un élément de $G$. Par conséquent, les droites sont toutes dans la même orbite pour l'action de $G$.
 \item Les formes alternées sont classifiées à équivalence près par leur rang qui est toujours pair. Par conséquent, les restrictions à un plan sont donc la forme nulle ou une forme symplectique de rang $2$.
 \item Par la question précédente et par le théorème de Witt, 
 on conclut comme dans la première question qu'il y a deux orbites sur l'ensemble des plans de $E$.  
\end{enumerate}
\fincorrection
\nocorrection
\nocorrection
\nocorrection
\nocorrection
\nocorrection
\nocorrection
\nocorrection
\nocorrection
\nocorrection
\nocorrection
\nocorrection
\nocorrection
\nocorrection
\nocorrection
\nocorrection
\nocorrection
\nocorrection
\nocorrection
\nocorrection
\nocorrection
\nocorrection
\nocorrection
\nocorrection
\nocorrection
\nocorrection
\nocorrection
\nocorrection
\nocorrection
\nocorrection
\nocorrection
\nocorrection
\nocorrection
\nocorrection
\nocorrection
\nocorrection
\nocorrection
\nocorrection
\nocorrection
\nocorrection
\nocorrection
\nocorrection
\nocorrection
\correction{007838}
Soit $n$ un entier supérieur à $3$.
Soit $D_{2n}$ le groupe des isométries d'un polygone régulier $P$ à $n$ côtés d'un plan euclidien réel.
Comme toutes les isométries sont affines, l'isobarycentre $G$ des sommets de $P$ est conservé.
Les éléments de $D_{2n}$ de déterminant positif sont donc des rotations de centre $G$, et ceux de déterminant négatif des symétries orthogonales par rapport à des droites passant par $G$. 
Comme une rotation qui a deux points fixes est l'identité, les rotations sont d'ordre inférieur à $n$.
Il n'y a donc que les $n$ rotations d'angle multiple de $2\pi/n$. Comme la composée de deux symétries d'axe passant par $G$ est une rotation, il y a aussi exactement $n$ symétries, obtenues comme produit des $n$ rotations par une symétrie fixée.

L'application $\det~:~D_{2n}\to \{-1,1\}$ est un morphisme de groupes qui a pour image le groupe commutatif $\{-1,1\}$ et pour noyau le groupe commutatif $R_n$ des rotations de centre $G$ d'angle multiple de $2\pi/n$. 
Le groupe dérivé $D^{(1)}(D_{2n})$ de $D_{2n}$ est donc inclus dans le sous-groupe commutatif des rotations.
Par conséquent, le deuxième groupe dérivé $D^{(2)}(D_{2n})=\{Id\}$ et $D_{2n}$ est résoluble. 
\fincorrection
\correction{007839}
Les formes bilinéaires symétriques données par les formes quadratiques suivantes sur $\mathbb{F}_7^3$,
$q(x,y,z)=x^2+6y^2+2z^2$
et $Q(x,y,z)=xy+3z^2$ ont pour matrice dans la base canonique
$$\begin{pmatrix}
  1&0&0\\0&6&0\\0&0&2
 \end{pmatrix}
\textrm{ et }
\begin{pmatrix}
 0&4&0\\4&0&0&\\0&0&3
\end{pmatrix}
$$
 et sont
non-dégénérées de discriminant $5$ et $1$.
Comme $5^{\frac{7-1}{2}}=-1[7]$, $5$ n'est pas un carré alors que $1=1^2$ en est un,
les formes ne sont donc pas équivalentes.
\fincorrection
\correction{007840}
Soit $(E,f)$ et $(E,f')'$ deux espaces symplectiques non-singuliers de dimension $4$.
Soit $d\subset E$ et $d'=vect(x')\subset E'$ deux droites 
et $f$ une application linéaire bijective de $d$ sur $d'$.
On écrit $d=vect(x)$ et $d'=vect(x')$ où $x'=f(x)$.

Notons que comme $d$ et $d'$ sont de dimension $1$ donc isotropes, $f$ est une isométrie (pour les formes bilinéaires induites). Puisque $f$ est non dégénérée, on choisit un vecteur $z$ (non colinéaire à $x$) tel que $f(x,z)$ est non nul donc inversible dans $k$. 
On définit $y$ sous la forme $y=\alpha x+ f(x,z)^{-1}z$. Alors, $(x,y)$ est une paire symplectique.
L'orthogonal de $vect(x,y)$ (muni de la forme induite) est un plan non-singulier puisque $vect(x,y)$ est non-singulier.
Par la construction précédente, on y trouve une paire symplectique $(X,Y)$.
Le même raisonnement permet de trouver une base $(x',y',X',Y')$ de $E'$ telle que $(x',y')$ et $(X',Y')$ soient deux paires symplectiques orthogonales de $E'$.
L'application $f$ définie par $x\mapsto x'$, $y\mapsto y'$, $X\mapsto X'$ et $Y\mapsto Y'$ est une isométrie de $E$ sur $E'$ qui prolonge $f$.
\fincorrection
\nocorrection
\nocorrection
\nocorrection
\nocorrection
\nocorrection
\nocorrection
\nocorrection
\nocorrection
\nocorrection
\nocorrection
\nocorrection
\nocorrection
\nocorrection
\nocorrection
\nocorrection
\correction{002547}
\begin{enumerate}
\item Consid\'erons $F: \mathbb{R}^3 \rightarrow \mathbb{R}$ d\'efinie
par $F(x_1,x_2,x_3)=x_1^2+x_2^2-x_3^2- \lambda$. Alors $F$ est de
classe $C^1$, $Jac F(x_1,x_2,x_3)=(2x_1,2x_2,-2x_3)$ et
$S_\lambda=\{(x_1,x_2,x_3) \in \mathbb{R}^3; F(x_1,x_2,x_3)=0\}$.
Si $\lambda \neq 0$, $\mbox{rang}(JacF(x_1,x_2,x_3))=1$ (le
maximum possible) car sinon $x_1,x_2,x_3$ seraient tous nuls:
impossible car $x_1^2+x_2^2-x_3^2=\lambda \neq 0$. Comme $(0,0,0)
\not \in S_\lambda, \forall a \in S_\lambda, \mbox{
rang}JacF(a)=1$ et donc $S_\lambda$ est une sous-vari\'et\'e de
$\mathbb{R}^3$ de dimension $2$.\\
Si  $\lambda=0$ $T_0(S_\lambda)=\{\mbox{vecteurs tangents \`a }
S_n \mbox{ en }0\}$. Alors $T_0S_0$ est un c\^one et donc $S_0$
n'est pas une sous-vari\'et\'e. 
\item Soientt $x,y \in
\mathbb{R}^3$, $B(x,y)=x_1y_1+x_2y_2-x_3y_3$ et $x\in S_\lambda$.
Si $\lambda \neq 0$, $JacF(x)=(2x_1,2x_2,-2x3)$ et donc
$$T_xS_\lambda=\{u \in \mathbb{R}^3;
DF(x).u=0\}=\{u=(u_1,u_2,u_3); (2x_1,2x_2,-2x_3). \left
(\begin{array}{c}u_1 \\ u_2 \\u_3
\end{array}\right)=0\}=$$
$$\{(u_1,u_2,u_3); 2x_1u_1+2x_2u_2-2x_3u_3=0\}=\{(u_1,u_2,u_3);
2B(x,u)=0\}$$ d'o\`u $$T_xS_\lambda=\{u \in \mathbb{R}^3;
B(x,u)=0\}.$$
\end{enumerate}
\fincorrection
\correction{002548}
Cas de $\mathbb{R}^2$. $$u=\left (\begin{array}{cc} u_{11} &
u_{12} \\ u_{21} & u_{22}\end{array}\right).$$ L'hypoth\`ese sur
$u$ implique que $u_{12}=u{21}$. Si $x=(x_1,x_2)$, on a
$$u(x)=\left (\begin{array}{c} u_{11}x_1+ u_{12}x_2 \\ u_{21}x_1 +
u_{22}x_2\end{array}\right)$$ et
$$<u(x),x>=\sum_{i=1}^2
u_i(x)x_i=(u_{11}x_1+u_{12}x_2)x_1+(u_{21}x_1+u_{22}x_2)x_2=u_{11}x_1^2+u_{12}x_1x_2+u_{21}x_1x_2+u_{22}x_2^2.$$
Posons $f(x)=<u(x),x>-1$ alors $$\frac{\partial f}{\partial x_1}=
2u_{11}x_1+u_{12}x_2+u_{21}x_2=2u_{11}x_1+2u_{12}x_2$$ et
$$\frac{\partial f}{\partial x_2}=
2u_{22}x_2+u_{12}x_1+u_{21}x_1=2u_{21}x_1+2u_{22}x_2.$$ Calculons
$$Df(x).x=x_1\frac{\partial f}{\partial x_1}+x_2\frac{\partial
f}{\partial
x_2}=2(u_{11}x_1^2+u_{12}x_2x_1+u_{21}x_1x_2+u_{22}x_2^2)=2<u(x),x>.$$
Si $x=(x_1,x_2) \in Q$ alors $<u(x),x>=1\neq 0$ et donc $Df(x)$
\'etant non nul, il est de rang au moins $1$ et donc de rang
maximal. $Q$ est bien une sous-vari\'et\'e de $\mathbb{R}^2$ de
dimension $1$.

D\'eterminons le plan tangent de $Q$.
$$T_xQ=\{y \in \mathbb{R}^2; Df(x)(y)=0\}=
\{y \in \mathbb{R}^n; \sum_{i=1}^n \frac{\partial f}{\partial
x_1}(x)y_i=0\}=$$
$$\{y \in \mathbb{R}^n; 2<u(x),y>=0\}.$$

\fincorrection
\nocorrection
\nocorrection
\nocorrection
\nocorrection
\nocorrection
\nocorrection
\nocorrection
\nocorrection
\nocorrection
\nocorrection
\nocorrection
\nocorrection
\nocorrection
\nocorrection
\nocorrection
\nocorrection
\nocorrection
\nocorrection
\nocorrection
\nocorrection
\nocorrection
\nocorrection
\nocorrection
\nocorrection
\nocorrection
\nocorrection
\nocorrection
\nocorrection
\nocorrection
\nocorrection
\nocorrection
\nocorrection
\nocorrection
\nocorrection
\nocorrection
\nocorrection
\nocorrection
\nocorrection
\nocorrection
\nocorrection
\nocorrection
\nocorrection
\nocorrection
\nocorrection
\nocorrection
\nocorrection
\nocorrection
\nocorrection
\nocorrection
\nocorrection
\nocorrection
\correction{007688}
Le vecteur vitesse est $\dot{c}(t)=(-4\sin t,-5\cos t,3\sin t)$.
Sa norme est $5$. Par conséquent, le paramétrage par 
$C~:~t\mapsto (4\cos t/5,5-5\sin t/5, -3\cos t/5)$ est un paramétrage par la longueur d'arc.
Le nouveau vecteur vitesse est $\dot{C}(t)=(-4/5\sin t/5,-\cos t/5,3/5\sin t/5)$.
Le vecteur accélération est $\ddot{C}(t)=(-4/25\cos t/5,1/5\sin t/5,3/25\cos t/5)$.
Sa norme $1/5$ est la courbure $\kappa$. Le vecteur normal unitaire est donc
$n(t)=(-4/5\cos t/5,\sin t/5,3/5\cos t/5)$.
Le vecteur binormal est $\dot{C}(t)\times n(t)=(-3/5,0,-4/5)$. Comme il est constant,
la torsion est nulle par les formules de Frenet.

Oui la courbe est plane, puisque la torsion est nulle. On peut aussi remarquer que la courbe est dans le plan d'équation
$3X+4Z=0$.
\fincorrection
\correction{007689}
\begin{enumerate}
 \item Le vecteur vitesse de la courbe $C'$ est $(1,2t,0)$ de norme $\sqrt{1+4t^2}$.
 La longueur de la courbe $C'$ est donc $l[C']=\int_0^1\sqrt{1+4t^2}dt$.
 La courbe $C$ est paramétrée par $x=t$, $y=t^2$ $z=2/3t^3$ pour $t$ parcourant $[0,1]$.
 Le vecteur vitesse de la courbe $C$ est $(1,2t,2t^2)$ de norme $\sqrt{1+4t^2+4t^4}=1+2t^2$.
 La longueur de la courbe $C$ est donc $l[C]=\int_0^1 1+2t^2dt$.
 Comme pour tout $t\in [0,1], 0\leq 1+4t^2\leq (1+2t^2)^2$, pour tout $t\in [0,1], \sqrt{1+4t^2}\leq (1+2t^2)$
 et par intégration $l[C']\leq l[C]$.
 \item En utilisant le changement de variables $2t=\sinh (u)$, on obtient le calcul de primitive
 \begin{eqnarray*}
  4\int \sqrt{1 + 4t^2} dt &=& 2\int \cosh^2 udu =\int(\cosh (2u) + 1) du \\
  &=& \sinh (2u)/2+ u  \\
  &=& \sinh u \cosh u + u\\
  &=&2t\sqrt{1 + 4t^2}+\ln(2t+\sqrt{1 + 4t^2}).
 \end{eqnarray*}
 Par conséquent, la longueur de $C'$ est $1/4[2t\sqrt{1 + 4t^2}+\ln(2t+\sqrt{1 + 4t^2})]_0^1
 =1/4(2\sqrt{5}+\ln(2+\sqrt{5}))$.
 La longueur de $C$ est $1+2/3$.
 \end{enumerate}
\fincorrection
\correction{007690}
\begin{enumerate}
 \item La surface $S$ est le graphe de la fonction $\mathcal{C}^\infty$ $f~:~(y,z)\mapsto -2(2z^2+y^2)$.
 Elle est donc régulière.
 \item La surface $S$ est paramétrée par $F(u,v)=(-2(u^2+2v^2),u,v)$.
 \item Le point $A$ est obtenu au point de paramètre $(1,-1)$
 On calcule $\frac{\partial F}{\partial u}=(-4u,1,0)=(-4,1,0)$ 
 et $\frac{\partial F}{\partial v}=(-8v,0,1)=(8,0,1)$.
 Ces deux vecteurs forment une base de l'espace tangent à $S$ au point de paramètre $(u,v)$.
 \item Comme $S$ est la ligne de niveau $0$ de la fonction $\varphi (x,y,z)=2(2z^2+y^2)+x$ sur $\Rr^3$,
 on calcule $(grad \varphi)_{A}=(1,4y,8z)=(1,4,-8)=:n$.
 Ce vecteur est bien orthogonal au plan tangent obtenu à la question 3.
 \item Comme $<V,n>=<(27,-29,-1),(1,4,-8)>=-81\not=0$, le vecteur $V$ n'est pas dans le plan tangent à $S$ en $A$.
\end{enumerate}
\fincorrection
\correction{007691}
On paramètre $\mathcal{E}$ en coordonnées sphériques par $F(\theta,\varphi)=
(\cos\theta\sin\varphi,\sin\theta\sin\varphi,1/\sqrt{5}\cos\varphi)$ avec $\theta\in[0,2\pi[$ et 
$\varphi\in [0,\pi]$.
Le plan tangent est engendré par les vecteurs
$\frac{\partial F}{\partial \theta}=(-\sin\theta\sin\varphi,\cos\theta\sin\varphi,0)$ 
et $\frac{\partial F}{\partial \varphi}=(\cos\theta\cos\varphi,\sin\theta\cos\varphi,-1/\sqrt{5}\sin\varphi)$.
Dans cette base, la matrice de la première forme fondamentale est 
$\begin{pmatrix}\sin^2\varphi&0\\0&\cos^2\varphi+1/5\sin^2\varphi\end{pmatrix}$.
L'élément de volume est $1/\sqrt{5}\sin\varphi\sqrt{1+4\cos^2\varphi}d\theta d\varphi$.
L'aire de l'ellipsoïde est donc
\begin{eqnarray*}
 A[\mathcal{E}]&=&\int_{\varphi=0}^{\pi}\int_{\theta=0}^{2\pi}\frac{1}{\sqrt{5}}
\sin\varphi\sqrt{1+4\cos^2\varphi} d\varphi d\theta\\
&=&2\pi\int_{\varphi=0}^{\pi}\frac{1}{\sqrt{5}}
\sin\varphi\sqrt{1+4\cos^2\varphi}d\varphi\\
&=&\frac{2\pi}{\sqrt{5}}\int_{t=-1}^1\sqrt{1+4t^2}dt=\frac{\pi}{\sqrt{5}}(2\sqrt{5}+\ln(2+\sqrt{5}))\\
&=&\pi(2+\frac{\ln(2+\sqrt{5})}{\sqrt{5}}).
\end{eqnarray*}
\fincorrection
\correction{007692}
On paramètre la surface $S$ par
$F(u,v)=(u,v,u^2-v^2)$.
L'espace tangent en $p$ est engendré par $(\frac{\partial F}{\partial u})_p=(1,0,2u)_p=(1,0,0)$ et
$(\frac{\partial F}{\partial v})_p=(0,1,-2v)_p=(0,1,0)$.
La matrice de la première forme fondamentale dans cette base est en $p$ de paramètre $(0,0)$
$\begin{pmatrix}
 1+4u^2&-4uv\\-4uv&1+4v^2
\end{pmatrix}=Id$.
 La matrice de la seconde forme fondamentale se calcule à l'aide du vecteur normal 
$n=(\frac{\partial F}{\partial u}\times\frac{\partial F}{\partial v})_p=(0,0,1)$
par
$\begin{pmatrix}
 <\frac{\partial^2 F}{\partial u^2},n>&<\frac{\partial^2 F}{\partial u\partial v},n>\\
<\frac{\partial^2 F}{\partial u\partial v},n>&<\frac{\partial^2 F}{\partial v^2},n>
 \end{pmatrix}=
\begin{pmatrix}
 2&0\\0&-2
\end{pmatrix}$.
Comme la matrice de la première forme fondamentale est l'identité,
cette dernière matrice est aussi celle de l'endomorphisme de Weingarten.
La courbure de Gauss est donc le déterminant $-4$.
Les deux directions principales sont dirigées par les vecteurs
 $\frac{\partial F}{\partial u}_p=(1,0,0)$ et
$\frac{\partial F}{\partial v}_p=(0,1,0)$.
\fincorrection
\correction{007693}
\begin{enumerate}
 \item La fonction courbure est continue et non constante sur le compact $[0,\ell]$.
 Elle atteint donc son maximum et son minimum en deux points distincts $P$ et $Q$.
 Comme elle est continûment dérivable, en ces points sa dérivée est nulle.
 Ce sont donc des coins.
 \item On utilise la formule de Frenet $\dot{n}(t)=-\kappa(t)\dot{c}(t)$.
 \begin{eqnarray*}
  \int_0^\ell \dot{\kappa}c(t)&=-&\int_0^\ell \kappa\dot{c}(t)=\int_0^\ell \dot{n}(t)
 =n(\ell)-n(0)=0.
 \end{eqnarray*}
 \item Par hypothèse, l'ordonnée de $\dot{\kappa}(t)c(t)$ serait de signe constant sur les deux
 parties de la courbe délimitées par $P$ et $Q$. Ceci contredit l'annulation de l'intégrale.
 \item S'il n'y a que trois coins, la courbe est partagée en trois parties.
 Sur deux des parties adjacentes $\dot{\kappa}$ a le même signe.
 Le raisonnement précédent en déplaçant la courbe,
 de sorte que $\dot{\kappa}$ soit de signe constant sur chaque demi-plan $(y\geq 0)$
 et $(y\leq 0)$ aboutit à une contradiction.  
\end{enumerate}
\fincorrection
\nocorrection
\nocorrection
\nocorrection
\correction{007697}
\begin{enumerate}
 \item  Les composantes de l'application $c$ sont de classe $\mathcal{C}^\infty$.
 La dérivée de la dernière composante n'est jamais nulle.
 La courbe $c$ est donc régulière.
 \item Pour tout paramètre $t$, $\|\dot{c}(t)\|=\sqrt{2+t^2}\geq \sqrt{2}$.
 Donc,
 $$L[c_{\mid [a,b]}]=\int_a^b \|\dot{c}(t)\| dt\geq\sqrt{2}(b-a).$$
  \item Soit $t\in\Rr $,
 $$(t\cos t)^2+(t\sin t)^2=t^2=(t)^2.$$
 Par conséquent, $\mathcal{C}$ est dans la quadrique d'équation $x^2+y^2=z^2$.
  \item  Soit $t\in\Rr $, $c(-t)= \left(\begin{array}{c}-t\cos t \\ t\sin t\\- t\end{array}\right)$ et se déduit donc de $c(t)$
 par le demi-tour autour de l'axe des ordonnées (d'équation $x=z=0$).
 \item 
	 $$ \includegraphics[scale=.4]{images/img-mour-449}$$
 %faire du gif avec wims2
 %convert wims.gif wims.png
 \item
  $$c(t)=\left(\begin{array}{c}t\cos t \\ t\sin t\\ t\end{array}\right) ; 
 \dot{c}(t)=\left(\begin{array}{c}\cos t -t\sin t\\ \sin t + t\cos t\\ 1\end{array}\right) ;$$
 $$
 \ddot{c}(t)=\left(\begin{array}{c}-2\sin t -t\cos t\\ 2\cos t - t\sin t\\ 0\end{array}\right) ; 
 \dddot{c}(t)=\left(\begin{array}{c}-3\cos t +t\sin t\\ -3\sin t -t\cos t\\ 0\end{array}\right).$$
 Ainsi, $$\det\begin{pmatrix}
        \cos t -t\sin t &-2\sin t -t\cos t&-3\cos t +t\sin t\\
 \sin t + t\cos t&2\cos t - t\sin t&-3\sin t - t\cos t\\1&0&0
        \end{pmatrix}=t^2+6$$
 Par ailleurs,
 $$\dot{c}(t)\wedge\ddot{c}(t)=\left(\begin{array}{c}\cos t -t\sin t\\ \sin t + t\cos t\\ 1\end{array}\right) 
 \wedge\left(\begin{array}{c}-2\sin t -t\cos t\\ 2\cos t - t\sin t\\ 0\end{array}\right)=
 \begin{pmatrix}
 - 2\cos t +t\sin t\\-2\sin t -t\cos t\\2+t^2
 \end{pmatrix}
 $$ dont la norme au carré est $t^4+5t^2+8$.
 La fonction torsion de $c$ est donc $$\tau(t)=\frac{t^2+6}{t^4+5t^2+8}.$$
    
\end{enumerate}
\fincorrection
\correction{007698}
L'application $F$ est injective et différentiable car ses composantes le sont.
 On calcule
$$X=_u=\frac{\partial F}{\partial u}=\begin{pmatrix}\sinh(u)\cos(v)\\ \sinh(u)\sin(v)\\ 1\end{pmatrix}
\quad \text{ et } \quad X_v=\frac{\partial F}{\partial v}=\begin{pmatrix}-\cosh(u)\sin(v)\\ \cosh(u)\cos(v)\\0\end{pmatrix}$$
pour déterminer 
le rang de $dF$, qui est donc $2$. Par conséquent, l'image $S$ de $F$ est une surface régulière.
On peut aussi déterminer
 $$G=\cosh^2 u\begin{pmatrix} 1&0\\0&1\end{pmatrix}.$$
On s'en sert aussi pour trouver un champs de vecteurs normaux unitaires différentiable
              $N=\frac{1}{\cosh u}\begin{pmatrix}-\cos v\\-\sin v\\\sinh u \end{pmatrix}$.
On peut ensuite, soit calculer 
$$W(X_u)=-dN\cdot X_u=-\frac{\partial N}{\partial u}=-(\cosh^{-2}u) X_u$$
$$\text{ et } \quad W(X_v)=-dN\cdot X_v=-\frac{\partial N}{\partial u}=(\cosh^{-2}u)
X_v$$
et en déduire que les valeurs propres de l'endomorphisme de Weingarten $W$ sont $-\cosh^{-2}u$ et $\cosh^{-2}u$.
Donc, $W$ est de demi-trace nulle, et la courbure moyenne de $Im F$ est partout nulle.

On peut aussi calculer
$\frac{\partial^2 F}{\partial u^2}$,
$\frac{\partial^2 F}{\partial u\partial v}$ et 
$\frac{\partial^2 F}{\partial u^2}$ pour obtenir $H=\begin{pmatrix}-1&0\\0&1\end{pmatrix}$.
On calcule $G^{-1}=\cosh^{-2} u\begin{pmatrix} 1&0\\0&1\end{pmatrix}$
Donc $W=G^{-1}H=\begin{pmatrix}-\cosh^{-2} u&0\\0&\cosh^{-2} u\end{pmatrix}$.
Comme la demi-trace de $W$ est nulle, la surface $S$ est partout de courbure moyenne nulle.
\fincorrection
\correction{007699}
\begin{enumerate}
 \item Comme pour tout $t\in \Rr $, et tout $a_in\Rr $, $c(t+a)=t_{(0,0,a)}(c(t))$,
 l'image de $F$ est invariante par toutes les translations de vecteur parallèle à $\vec{k}$,
 le troisième vecteur de la base canonique.
 
 On remarque aussi que $F(-t,-s)$ s'obtient à partir de $F(t,s)$ par le demi-tour d'axe des abscisses. 
 Ce demi-tour conserve donc l'image de $F$.
 \item Soit $(u,v)\in \Rr ^2$,
 $$(1-t^2)^2(1-(1-t^2))=(1-t^2)^2t^2=[t(1-t^2)]^2.$$
 Donc, l'image de $F$ est incluse dans l'ensemble d'équation $y^2=x^2(1-x)$.
 
 Soit $(x,y,z)$ inclus dans l'ensemble d'équation $y^2=x^2(1-x)$.
 Si $x=0$, $y=0$ et $(x,y,z)=F(1,z)$. Si $x\not=0$, soit $s=z$ et $t=y/x$.
 On vérifie que $F(t,s)=(x,y,z)$.
 
 Donc, l'image de $F$ est l'ensemble d'équation $y^2=x^2(1-x)$.
 \item Au voisinage du point $(0,0,0)$ de l'image de $S$,
 la projection sur le plan $x=0$ donne a chaque point deux antécédents,
 la projection sur la plan $y=0$ aussi,
 et la projection sur la plan $z=0$ donne à chaque point soit une infinité, soit aucun antécédent.
 Par conséquent, au voisinage de $(0,0,0)$ l'image de $S$ n'est pas un graphe.
 L'image de $F$ n'est donc pas une surface régulière de $\Rr ^3$.
\end{enumerate} 
\fincorrection 
\correction{007700}
On note $S$ la sphère unité de $\Rr ^3$.
La fonction $f$ est différentiable sur la sphère comme restriction d'une fonction polynomiale sur $\Rr ^3$.
La fonction $f$ continue sur le compact $S$ atteint ses extrema.
 En un point extremal, la différentielle de la fonction $f$ restreinte à l'espace tangent doit être nulle.
Le gradient de $f$ doit donc être normal à l'espace tangent.
On trouve que le vecteur 
$\left(\begin{array}{c}y\\x\\0\end{array}\right)$ doit être parallèle au vecteur $\left(\begin{array}{c}x\\y\\z\end{array}\right)$.
Comme $x$, $y$ et $z$ ne sont pas simultanément nuls sur $S$, $z=0$ et $x^2=y^2$.
Puisque $x^2+y^2+z^2=1$, $x=\pm 1/\sqrt{2}$ et $y=\pm 1/\sqrt{2}$.
En comparant les valeurs en ces quatre points, on trouve que
le minimum absolu de la fonction $f$ est donc $-1/2$ atteint au deux points $(1/\sqrt{2},-1/\sqrt{2},0)$ et
$(-1/\sqrt{2},1/\sqrt{2},0)$.
Le maximum absolu de la fonction $f$ est donc $1/2$ atteint au deux points $(1/\sqrt{2},1/\sqrt{2},0)$ et
$(-1/\sqrt{2},-1/\sqrt{2},0)$.
\fincorrection
\correction{007701}
\begin{enumerate}
 \item Comme les composantes de $c$ sont de classe $\mathcal{C}^\infty$, elle est de classe $\mathcal{C}^\infty$.
  La dérivée de la dernière coordonnée ne s'annule jamais.
 La courbe $c$ est donc régulière.
 \item Le vecteur vitesse
  $$\dot{c}(t)=\begin{pmatrix}-3\sin t\\3\cos t\\4\end{pmatrix}$$
 est de norme $5$. Comme la courbe $c$ est régulière, on peut trouver un paramétrage par longueur d'arc.
 Le paramétrage
 $$\begin{array}{cccc}
         e~:~&]0,5[&\to&\Rr ^3\\
 &s&\mapsto&\left(\begin{array}{c}3\cos (s/5)\\3\sin (s/5)\\4s/5\end{array}\right)
         \end{array}
 $$
 de la forme $e(s)=c(\phi(s))$ est un reparamétrage par longueur d'arc, car $de/ds$ est partout de norme $1$.
 \item On utilise le paramétrage par longueur d'arc.
  Le premier vecteur est le vecteur vitesse $\dot{e}(s)=v(s)=\begin{pmatrix}-3/5\sin (s/5)\\3/5\cos (s/5)\\4/5\end{pmatrix}$.
 Le deuxième est le vecteur normal. Le vecteur accélération est $\ddot{e}(s)=\begin{pmatrix}-3/25\cos (s/5)\\-3/25\sin (s/5)\\0\end{pmatrix}$.
 On trouve que la courbure est $3/25$ et le vecteur normal 
 $n(s)=\begin{pmatrix}-\cos (s/5)\\-\sin (s/5)\\0\end{pmatrix}$.
 Le troisième vecteur est le vecteur binormal
 $b(s)=v(s)\wedge n(s)=\begin{pmatrix}-3/5\sin (s/5)\\3/5\cos (s/5)\\4/5\end{pmatrix}
 \wedge\begin{pmatrix}-\cos (s/5)\\-\sin (s/5)\\0\end{pmatrix}
 =\begin{pmatrix} 4/5\sin (s/5)\\-4/5 \cos (s/5)\\ 3/5\end{pmatrix}$.  
\end{enumerate} 
\fincorrection
\correction{007702}
\begin{enumerate}
 \item ~

\includegraphics[scale=0.5]{images/img-mour-454}

 \item L'application $F$ est de classe $\mathcal{C}^\infty$ et injective.
 On calcule
 $X_\varphi=\begin{pmatrix}-\sin(\varphi)\cos(\theta)\\ -\sin(\varphi)\sin(\theta)\\ \cos (\varphi)\end{pmatrix}$
 et $X_\theta=\begin{pmatrix}-(2+\cos(\varphi))\sin(\theta)\\ (2+\cos(\varphi))\cos(\theta)\\ 0\end{pmatrix}$.
 Si $\varphi\not=0$, $\sin(\varphi)\not=0$ et les deux premières coordonnées sont indépendantes.
 Si $\varphi=0$, $X_\varphi=\begin{pmatrix}0\\0\\1\end{pmatrix}$ et $X_\theta=\begin{pmatrix}-\sin(\theta)\\ \cos(\theta)\\ 0\end{pmatrix}$
 sont indépendants.
 Ainsi, $dF$ est partout de rang $2$ donc un homéomorphisme local. Puisqu'elle est injective,
 l'application $F$ est donc un paramétrage.
 L'image $T$ est donc une surface régulière de $\Rr ^3$.
 \item  On détermine d'abord la matrice de la première forme fondamentale
 $$G=\begin{pmatrix}
    1&0\\0&(2+\cos(\varphi))^2
   \end{pmatrix}$$
 On choisit comme champ de vecteurs normaux unitaires
 $$N=-\frac{X_\varphi\wedge X_\theta}{2+\cos(\varphi)}
 =\begin{pmatrix}\sin(\varphi)\cos(\theta)\\ \sin(\varphi)\sin(\theta)\\ -\cos (\varphi)\end{pmatrix}
 \wedge \begin{pmatrix}-\sin(\theta)\\ \cos(\theta)\\ 0\end{pmatrix}
 =\begin{pmatrix}
  \cos (\varphi)\cos(\theta)\\ \cos (\varphi)\sin(\theta)\\
 \sin(\varphi)
  \end{pmatrix}
 $$
 On détermine ensuite l'endomorphisme de Weingarten $W$ par 
 $$WX_\varphi=\frac{\partial N}{\partial \varphi}
 =\begin{pmatrix}
  -\sin (\varphi)\cos(\theta)\\ -\sin (\varphi)\sin(\theta)\\
 \cos(\varphi)\end{pmatrix}=X_\varphi$$
 $$\text{ et } \quad WX_\theta=\frac{\partial N}{\partial \theta}=
 \begin{pmatrix}
  -\cos (\varphi)\sin(\theta)\\ \cos (\varphi)\cos(\theta)\\0\end{pmatrix} =\frac{\cos(\varphi)}{2+\cos(\varphi)}X_\theta$$
 On en déduit que les valeurs propres de $W$ sont $1$ et $\frac{\cos(\varphi)}{2+\cos(\varphi)}$,
 et que la courbure de Gauss est donc $K(\varphi,\theta)=\frac{\cos(\varphi)}{2+\cos(\varphi)}$.
 \item \begin{eqnarray*}
  \int_T K(m)d\sigma(m)&=&\int_{]-\pi,\pi[\times ]-\pi,\pi[}K(\varphi,\theta)\sqrt{det G}d\varphi d\theta\\
 &=&\int_{]-\pi,\pi[\times ]-\pi,\pi[}\frac{\cos(\varphi)}{2+\cos(\varphi)} (2+\cos(\varphi))d\varphi d\theta\\
 &=&2\pi\int_{-\pi}^{\pi}\cos(\varphi)d\varphi=0.
  \end{eqnarray*}   
\end{enumerate}
\fincorrection
\correction{007703}
\begin{enumerate}
 \item Comme $F $ est un paramétrage de la surface régulière $T$, les vecteurs 
 $X_\varphi=\frac{\partial F}{\partial \varphi}$ et $X_\theta=\frac{\partial F}{\partial \theta}$
 sont indépendants et forment une base de l'espace tangent à $T$ en $F(\varphi,\theta)$.
 La matrice imposée est bien celle d'un produit scalaire. Comme elle est constante,
 elle dépend de fa\c con $\mathcal{C}^\infty$ de $(\varphi,\theta)$.
 \item Comme la matrice de la métrique riemannienne est constante, les symboles de Christoffel
 sont nuls en tous points et le tenseur de courbure est identiquement nul.
 On en déduit que la courbure de Gauss est identiquement nulle.
 \item D'après la question précédente $\int_T K(m)d\sigma(m)=0$  
\end{enumerate}
\fincorrection
\nocorrection
\nocorrection
\nocorrection
\nocorrection
\nocorrection
\nocorrection
\correction{007710}
\begin{enumerate}
 \item  La courbure d'un cercle de rayon $r$ est en valeur absolue $1/r$.
 
 \item Puisque la fonction norme (au carré) $\phi$ est différentiablement continue et maximale au point de paramètre $\tau$,
 sa dérivée en $\tau$ s'annule. On trouve $\phi'(\tau)=2<c(\tau),\dot{c}(\tau)>=0$.
 Par ailleurs, comme la fonction $c$ donne un paramétrage par la longueur d'arc,
 la fonction $t\mapsto \| \dot{c}(t)\|^2$ est constante égale à $1$,
 donc aussi de dérivée nulle en $\tau$. On trouve $2<\ddot{c}(\tau),\dot{c}(\tau)>=0$.
 Par conséquent, en notant $\vec{n}$ un vecteur normal unitaire à $\dot{c}(\tau)$,
 on trouve $c(\tau)=+/-\vec{n}$ et $\ddot{c}(\tau)=\kappa(\tau)\vec{n}$.

 \item Puisque la fonction norme $\phi$ est deux fois différentiablement continue et maximale au point de paramètre $\tau$,
 sa dérivée seconde en $\tau$ est négative. On trouve $\phi''(\tau)=2<c(\tau),\ddot{c}(\tau)>+2<\dot{c}(\tau),\dot{c}(\tau)>=+/-2r\kappa(\tau)+2\leq 0$.
 Donc, $-/+2r\kappa(\tau)\geq 2$ et par suite $-/+r\kappa(\tau)$ est positif de valeur absolue au moins $1$.
 En divisant par $r>0$, on trouve le résultat.
 \item La courbe est en ses points extrémaux plus courbée que le cercle de rayon $r$ qui la borde.   
\end{enumerate}
\fincorrection
\correction{007711}
\begin{enumerate}
 \item Le plan $T_M\mathcal{S}$ tangent à $\mathcal{S}$ au point $M$ est engendré par les deux vecteurs
 $$\frac{\partial F}{\partial u}=\begin{pmatrix}1\\2u_0\\0\end{pmatrix} \quad \text{ et } \quad \frac{\partial F}{\partial v}=\begin{pmatrix}2v_0\\-2v_0\\1\end{pmatrix}.$$
 Une équation cartésienne du plan $T_M\mathcal{S}$ est obtenue par
 $$\begin{pmatrix}X\\Y\\Z\end{pmatrix}\in T_M\mathcal{S}\iff \begin{vmatrix}
                                1&2v_0&X\\2u_0&-2v_0&Y\\0&1&Z
                               \end{vmatrix}=0
 \iff -2u_0X+Y+2(1+2u_0)v_0Z=0.$$
 \item Il suffit de dire que le plan tangent $T_M\Sigma$ est le noyau de la différentielle de la fonction 
 $\psi : (x,y,z)\mapsto x^5+y^5+z^5-1$ en $(x_0,y_0,z_0)$.
 On trouve $$\begin{pmatrix}X\\Y\\Z\end{pmatrix}\in T_M\Sigma\iff
 5(x_0^4X+y_0^4Y+z_0^4Z)=0\iff x_0^4X+y_0^4Y+z_0^4Z=0.$$
 \item Une équation du graphe $\mathcal{G}$ de $f$ est $z=f(x,y)$.
 Une équation du plan tangent en $M$ de coordonnées $(x,y,f(x,y))$ est donc
 $$\begin{pmatrix}X\\Y\\Z\end{pmatrix}\in T_M\Sigma\iff
 Z=\frac{\partial f}{\partial x}(x,y)X+\frac{\partial f}{\partial y}(x,y)Y.$$
\end{enumerate}
\fincorrection
\correction{007712}
\begin{enumerate}
 \item Par l'exercice précédent, $\mathcal{P}$ a pour plan tangent au point $M(x,y,z)$
le plan d'équation $Z=2yY-2xX$. 
En particulier, le plan tangent en $A$ est le plan d'équation $Z=0$
engendré par les deux vecteurs de base $\begin{pmatrix} 1\\0\\0\end{pmatrix}$ et $\begin{pmatrix} 0\\1\\0\end{pmatrix}$.
Un champs de vecteurs normaux unitaires est donc
$$N(x,y,y)=\frac{1}{\sqrt{4x^2+4y^2+1}}\begin{pmatrix} 2x\\-2y\\1 \end{pmatrix}.$$

 \item L'endomorphisme de Weingarten est donné au point $A(0,0,0)$
par $W_A=-dN(A) : T_A\mathcal{P}\to T_A\mathcal{P}$ 
$$\begin{array}{c}   \begin{pmatrix} 1\\0\\0\end{pmatrix} 
\mapsto -\begin{pmatrix} 2\\0\\0 \end{pmatrix}\\
\begin{pmatrix} 0\\1\\0\end{pmatrix}\mapsto\begin{pmatrix}0\\2\\0\end{pmatrix}
 \end{array}$$

 \item Ses valeurs propres sont donc $-2$ et $2$. La courbure de Gauss de $\mathcal{P}$ en $A$ est donc $-4$.
Les directions propres sont donc les axes de coordonnées $x$ et $y$.
\end{enumerate}
\fincorrection
\correction{007713}
On trouve en utilisant la relation $\vec{NP'}\parallel\vec{NP}$ 
que la projection stéréographique depuis le pôle nord sur le plan équatoriale est donnée par
$$\begin{array}{ccc}
  p : S^2&\to&\Rr^2\\\begin{pmatrix}x\\y\\z\end{pmatrix}&\mapsto&\begin{pmatrix}\frac{x}{1-z}\\ \frac{y}{1-z}\\0\end{pmatrix}
 \end{array}
$$

\includegraphics[scale=0.4]{images/img-mour-465}

et la réciproque
$$\begin{array}{ccc}
F : \Rr^2&\to&S^2\\\begin{pmatrix}x'\\y'\end{pmatrix}&\mapsto&
\begin{pmatrix} \frac{2x'}{(x')^2+(y')^2+1}\\ \frac{2y'}{(x')^2+(y')^2+1}\\1-\frac{2}{(x')^2+(y')^2+1}\end{pmatrix}
\end{array}
$$
En particulier, la différentielle du paramétrage $F$ au point de paramètres $(0,0)$ est 
$$\begin{array}{cccc}
dF(0,0) : &T_{(0,0)}\Rr^2&\to&S_S^2\\ &\begin{pmatrix}X'\\Y'\end{pmatrix}&\mapsto&
\begin{pmatrix} 2X'\\ 2Y'\\0\end{pmatrix}
\end{array}
$$ n'est ni une isométrie, ni de déterminant $+/-1$.
La projection stéréographique ne conserve donc ni les longueurs, ni les aires.
Par contre, la différentielle du paramétrage $F$ au point de paramètres $(x',y')$ est 
$$\begin{array}{cccc}
dF(x',y') :& T_{(x',y')}\Rr^2&\to&TS^2\\ & \begin{pmatrix}X'\\Y'\end{pmatrix}&\mapsto&
\begin{pmatrix} 2\frac{-(x')^2+(y')^2+1}{((x')^2+(y')^2+1)^2}X'-4\frac{x'y'}{((x')^2+(y')^2+1)^2}Y'
\\ -4\frac{x'y'}{((x')^2+(y')^2+1)^2}X'+ 2\frac{(x')^2-(y')^2+1}{((x')^2+(y')^2+1)^2}Y'\\ 
4\frac{x'}{((x')^2+(y')^2+1)^2}X'+4\frac{y'}{((x')^2+(y')^2+1)^2}Y'\end{pmatrix}
\end{array}
$$
On vérifie alors que
$$\left\| dF(x',y')\begin{pmatrix}X'\\Y'\end{pmatrix}\right\|^2=\frac{4\left\|\begin{pmatrix}X'\\Y'\end{pmatrix}\right\|^2}{(x')^2+(y')^2+1}.$$
Donc, $dF(x',y')$ conserve les angles.
\fincorrection
\correction{007714}
Non, par exemple, les deux surfaces régulières (graphes) d'équation $(z=0)$ et $(z=x^2-y^3)$
s'intersectent sur la courbe d'équation $(z=0, x^2-y^3=0)$ qui est singulière en $(0,0,0)$.
\fincorrection
\nocorrection
\nocorrection
\nocorrection
\nocorrection
\nocorrection
\nocorrection
\nocorrection
\nocorrection
\nocorrection
\nocorrection
\nocorrection
\nocorrection
\nocorrection
\nocorrection
\nocorrection
\nocorrection
\nocorrection
\nocorrection
\nocorrection
\nocorrection
\nocorrection
\nocorrection
\nocorrection
\nocorrection
\nocorrection
\correction{002495}
\begin{enumerate}
\item On a par définition $B(0,1)=\{x \in \mathbb{R}; |x-0|=|x| <
1\}=[-1,1].$ 
\item C'est la norme euclidienne sur $\mathbb{R}^2$,
$B_1(0,1)=\{(x,y)\in \mathbb{R}^2; \sqrt{x^2+y^2}=1\}$ c'est le
disque de centre l'origine et de rayon $1$. 
\item $B_2(0,1)=\{(x,y); |x|<1 et |y|<1 \}.$ C'est un carr\'e. 
\item $B_3(0,1)=\{(x,y); |x|+|y|<1\}.$ Dans le quart de plan
$P^{++}=\{(x,y); x\geq 0, y\geq $, on a $B_3(0,1) \cap
P^{++}=\{(x,y) \in p^{++}; x+y<1 \}$ c'est le triangle
d\'elimit\'e par les droites $x=0, y=0$ et $x+y=1$. En faisant de
même pour les 3 autres secteurs du plan, on trouve que $B_3(0,1)$
est un losange (ou carr\'e) dont les sommets sont les points
$(0,1), (1,0), (-1,0), (0,-1).$
\end{enumerate}

Toutes ces distances \'etant invariantes par translation (ce sont
des normes), il suffit de montrez que les normes associ\'es
$||.||_i=d_i((x,y),0)$ sont \'equivalentes.\par

On a $$\frac{\sqrt{2}}{2} ||.||_1 \leq ||.||_2 \leq ||.||_1.$$
 En effet, $$||(x,y)||_1=\sqrt{(x^2+y^2)} \leq
\sqrt{\sup(x^2,y^2)+\sup(x^2,y^2)}$$  $$\leq \sqrt 2 \sqrt
{\sup(x^2,y^2)}$$ $$ \leq \sqrt 2 \sqrt { (\sup(|x|,|y|))^2} \leq
\sqrt 2 ||(x,y)||_2$$. De plus, $$||(x,y)||_1=\sqrt{x^2+y^2} \geq
\sqrt{\sup(x^2,y^2)}$$ $$\geq \sqrt{(\sup(|x|,|y|)^2} \geq
\sup(|x|,|y|) \geq ||(x,y)||_2.$$ Les distances $d_1$ et $d_2$
sont donc \'equivalentes. \par

De m\^eme on montre que  $$||.||_2 \leq ||.||_3 \leq 2||.||_2.$$
\fincorrection
\correction{002496}
Il faut trouver une suite de cauchy de fonctions de $E$ qui ne
converge pas dans $E$. Il suffit, par exemple, de prendre une
suites de fonctions $\{f_n\}$  convergeant pour $\|.\|$ vers une
fonction non continue. Par exemple, prendre
$$f_n(x)= \left\{\begin{array}{c}
1 \mbox{ si } x < 1/2\\
1-n(x-1/2) \mbox{ si } 1/2 \leq x \leq 1/2+1/n \\
0 \mbox{ si } x > 1/2+1/n
\end{array} \right \} $$
et
$$f_0(x)= \left\{\begin{array}{c}
1 \mbox{ si } x < 1/2\\
0 \mbox{ si } x \geq 1/2
\end{array} \right \} $$
 On a alors $\|f_n - f_0\|_1=1/(2n)$, la suite converge
simplement et en norme $\|.\|$ vers la fonction  $f_0$ qui n'est
pas continue. Il suffit de montrer alors qu'il n'existe aucune
fonction continue $g$ telle que $\|f-g\|=0$ ce qui interdit
l'existence d'une limite à $f_n$ dans $E$.
\fincorrection
\nocorrection
\correction{002498}
On montre par r\'ecurrence que $f(nx)=nx$ si $n\in \mathbb{N}$.
Montrer $f(-x)=-f(x)$ pour arriver à $f(nx)=nf(x)$ si $n \in
\mathbb{Z}$ puis $f(\frac{p}{q}x)=\frac{p}{q}f(x)$ $p,q \in
\mathbb{Z}$. Ainsi $f$ est lin\'eaire sur $\mathbb{Q}$. Il reste
\`a montrer qu'elle l'est sur $\mathbb{R}$. Soit $x \in E$ et
$\lambda \in \mathbb{R}$, il reste \`a montrer que $f(\lambda
x)=\lambda f(x)$. Prenons $\{\lambda_n\}_{n \in \mathbb{N}}$ tel
que $\lim_{n\rightarrow \infty}\lambda_n=\lambda$. On a alors
$$f(\lambda_x)=f(\lambda_n x +
(\lambda-\lambda_n)x)=\lambda_n f(x)+f((\lambda-\lambda_n)x).$$
Soit $c_n \in \mathbb{\mathbb{Q}}$ tel que
$$||(\lambda-\lambda_n)x||_E \leq c_n \leq 2
||(\lambda-\lambda_n)x||_E.$$ Alors
$$f((\lambda-\lambda_n)x)=f(c_n\frac{\lambda-\lambda_n}{c_n}x)=c_nf(\frac{\lambda-\lambda_n}{c_n}x)$$
et
$$||\frac{\lambda-\lambda_n}{c_n}x|| \leq 1.$$
L'application $f$ \'etant born\'e sur la boulle unit\'e par une
constante $M > 0$, on a $$||f((\lambda-\lambda_n)x)||\leq c_nM$$
et donc $$||f((\lambda-\lambda_n)x)||\leq c_n M\leq
2M||(\lambda-\lambda_n)x||_E$$ et donc $$\lim_{n\rightarrow
\infty}f((\lambda-\lambda_n)x)=0$$, en remarquant qu'on a aussi
$$\lim_{n\rightarrow \infty}\lambda_nf(x)=\lambda f(x)$$ on obtient
$$f(\lambda x)=\lim_{n \rightarrow \infty}[\lambda_n
f(x)+f((\lambda-\lambda_n)x)]=\lambda f(x).$$
\fincorrection
\correction{002499}
Soit $X=(x,y)$, on a $M.X=(ax+by, cx+dy)$ or $$|ax+by|\leq
|ax|+|by|\leq (|a|+|b|)\sup(|x|,|y|) \leq (|a|+|b|)||(x,y)||_1.$$
de m\^eme,
$$|cx+dy| \leq (|c|+|d|)||(x,y)||_1.$$
Par cons\'equent $$||M.X||_2 \leq \sup(|a|+|b|,|c|+|d|)
||(x,y)||_1$$ et donc $$||M|| \leq \sup(|a|+|b|,|c|+|d|).$$
Supposons $|a|+|b| \geq |c|+|d|$ (inverser l'ordre sinon) et
prenons $X_0=(a/|a|,b/|b|)$ (on suppose $a \neq 0$ et $b \neq 0$
sinon v\'erification facile). On a alors $||X_0||=1$ et
$$||M.X_0||_2=\sup(|a|+|b|,|ca/|a|+db/|b||) \geq |a|+|b|.1 \geq
(|a|+|b|)||X_0||_1$$ et donc $$||M|| \geq \sup(|a|+|b|,|c|+|d|)$$
et finalement $$||M||= \sup(|a|+|b|,|c|+|d|)$$
\fincorrection
\nocorrection
\correction{002501}
\begin{enumerate}
\item Soit $x$ une suite, on a
$$||S(x)||_\infty=Max(\sup_{n \in \mathbb{N}}|x_{n-1}|,0)=\sup_{n \in
\mathbb{N}}|x_{n}|=1.||x||_\infty.$$ Donc $||S||=1.$ \item Soit $f
\in {\mathcal C}([0,1])$
$$||Tf||_\infty=\sup_{x \in [0,1]}f(x)g(x)\leq
||f||_\infty||g||_\infty.$$ Donc $$||T||\leq ||g||_\infty.$$ Or
$$||T1||_\infty=||g||_\infty=||1||_\infty||g||_\infty.$$
Donc $$||T|| \geq ||g||_\infty$$ et finalement on a bien
$$||T||=||g||_\infty.$$
\item Soit $f \in {\mathcal C}([0,1])$, on a
$$||u(f)||=|\int_0^1f(x)g(x)dx| \leq \int_0^1|f(x)||g(x)|dx \leq
\sup_{x \in [0,1]}|g(x)| \int_0^1 |f(x)|dx \leq
||g||_\infty.||f||.$$ On a donc $$||u|| \leq ||g||_\infty.$$ Comme
$g$ ne s'annule qu'au point $x=1/2$, elle ne change de signe
qu'une seule fois. Soit
$$f_0=g/|g|,$$ cette fonction n'est pas continue (ni d\'efinie) en
$x=1/2$ mais v\'erifie $f_0g=|g|$. Prenons $f_n=g/|g|$ si
$|x-1/2|> 1/n$, pour $|x-1/2|\leq 1/n$, on relie les deux segments
du graphe par une ligne. Alors $1-1/(2n)\leq ||f_n|| \leq 1$ et
$$||u(f_n)||=|\int_{|x-1/2|>1/n}f_n(x)g(x)dx+\int_{|x-1/2| \leq
1/n}f_n(x)g(x)dx| \geq$$
$$|(|\int_{|x-1/2|>1/n}f_n(x)g(x)dx|-|\int_{|x-1/2| \geq
1/n}f_n(x)g(x)dx|)|$$
$$\geq ||g||_\infty \int_{|x-1/2|>1/n}|f_n(x)|dx-2/n ||g||_\infty\geq ||g||_\infty(||f_n|| -2/n).$$
 Ainsi $$\lim_{n\rightarrow
\infty}||u(\frac{f_n}{||f_n||})||\geq ||g||_\infty
(1-\frac{1}{2n||f_n||} \geq ||g||_\infty (1-\frac{1}{2n(1-1/2n))}
\geq ||g||_\infty (1-\frac{1}{2n-1})
$$ et donc pour tout $n \in \mathbb{N}^*$:
$$||u|| \geq ||g||_\infty (1-\frac{1}{2n-1}),$$ en faisant tend $n$
vers l'infini $$||u|| \geq ||g||_\infty$$ ce qui montre la
deuxi\`eme in\'egalit\'e et on obtient $||u||=||g||_\infty$.

\item si on prend $(x_n)=(a_n)$ on obtient $$u((a_n))=\sum a_n^2=
||(a_n)||_2^2=||(a_n)||_2.||(a_n)||_2$$ et donc $$||u|| \geq
||(a_n)||_2.$$

Or D'apr\`es Cauchy-Schwartz, on a $$||u(a_n)||=|u(a_n)|=|\sum a_n
x_n| \leq ||(a_n)||_2||(x_n)||_2$$ et donc $||u|| \leq
||(a_n)||_2$ d'o\`u l'\'egalit\'e $$||u||=||(a_n)||_2.$$
\item

Pour tout $j \in \mathbb{N}$ on a $|x_j| \leq ||(x_n)||_\infty$ et
par cons\'equent $$|u((x_n))|=|\lim_{j \rightarrow \infty} x_j|
\leq ||(x_n)||_\infty$$ et donc $$||u|| \leq 1.$$ Prenons la suite
$(x^0)$ d\'efinie par $x_n^0=1$ pour tout $n \in \mathbb{N}$ alors
$$|u(x^0)|=|\lim_{j \rightarrow \infty} 1|=1=||x^0||_\infty$$ et
donc $$||u|| \geq 1$$ d'o\`u l'\'egalit\'e $||u||=1$.
\end{enumerate}
\fincorrection
\nocorrection
\correction{002503}
\begin{enumerate}

\item (Etude en $0$). $|sin(1/x)| \leq 1$ par cons\'equent
$|x^2sin(1/x)| \leq x^2$. De m\^eme $|y^2sin(1/y)| \leq y^2$. Par
cons\'equent
$$|f(x,y)| \leq x^2+y^2 \leq (\sqrt{x^2+y^2})^2 \leq (||(x,y)||_2)^2$$ Et donc $$\lim_{||(x,y)||
\rightarrow 0}|f(x)-f(0)|=0$$ et donc $f$ est continue à
l'origine. En remarquant que $||(x,y)||_2^2=o(||(x,y)-(0,0)||_2)$
on a $f(x,y)=0 +o(||(x,y)-(0,0)||_2)$ et donc $f$ est
diff\'erentiable en $0$ et $$Df(0)=0.$$ Par cons\'equent $f$ admet
des d\'eriv\'ees partielles dans toutes les directions \`a
l'origine qui sont nulles. La fonction $f$ n'est pas contre par de
classe $C^1$ \`a l'origine. Il suffit de remarquer que la
d\'eriv\'ee partielle $\frac{\partial f}{\partial x}$ sur la
droite $y=0$ n'est pas continue en $0$.

\item Pour $(x,y) \neq (0,0)$, $f$ est continue en $(x,y)$ et même
de classe $C^\infty$ en tant que compos\'es sommes, produits et
quotient de telles fonctions. Il reste \`a \'etudier $f$ \`a
l'origine. Or, $$|f(x,y)|=\frac{|xy^2|}{x^2+y^2}\leq
\frac{|x|(x^2+y^2)}{x^2+y^2} \leq |x| \leq ||(x,y)||_2.$$ Ainsi,
$f$ est continue \`a l'origine et y tend vers $0$.

Montrons par l'absurde que $f$ n'est pas d\'erivable \`a
l'origine. Notons $Df(0)$ la (suppos\'ee) diff\'erentielle de $f$
\`a l'origine. L'application lin\'eaire $Df(0)$ s'obtient par la
calcul de l'image de vecteurs de la base de $\mathbb{R}^2$.
Calculons pour `les d\'eriv\'ees directionnelles de $f$ \`a
l'origine: $$D_{(1,0)}f(0)=[Df(0)]((1,0))= \lim_{h \rightarrow 0}
\frac{f(0+h(1,0))-f(0)}{h}=\lim_{h \rightarrow 0}
\frac{f(h,0)}{h}=\lim_{h \rightarrow 0} 0=0.$$
$$D_{(0,1)}f(0)=[Df(0)]((0,1))= \lim_{h \rightarrow 0}
\frac{f(0+h(0,1))-f(0)}{h}=\lim_{h \rightarrow 0}
\frac{f(0,h)}{h}=\lim_{h \rightarrow 0} 0=0.$$ Par cons\'equent,
on a n\'ecessairement $$Df(0)=0$$
Or,$$D_{(1,1)}f(0)=[Df(0)]((1,1))= \lim_{h \rightarrow 0}
\frac{f(0+h(1,1))-f(0)}{h}=\lim_{h \rightarrow 0}
\frac{f(h,h)}{h}=\lim_{h \rightarrow 0}
\frac{\frac{h^3}{2h^2}}{h}=\frac{1}{2} \neq 0$$ ce qui donne la
contradiction recherch\'ee.
\end{enumerate}
\fincorrection
\nocorrection
\correction{002505}
En tout point $(x_0,y_0)$ avec $x_0 \neq y_0$, $f$ est continue et
même de classe $C^2$ car compos\'ee (projections sur (0x) et
(Oy)), diff\'erence et quotient de fonctions de classe $C^2$ dont
le d\'enominateur ne s'annule pas. Dans ces points, la
diff\'erentielle de $f$ est donn\'ee par la matrice jacobienne:
$$Df(x_0,y_0)=\left( \begin{array}{ccc}
\frac{\partial f}{\partial x}(x_0,y_0) &, & \frac{\partial
f}{\partial y}(x_0,y_0)
\end{array}
\right)=$$
$$\left( \begin{array}{ccc}
\frac{g'(x_0)(x_0-y_0)-(g(x_0)-g(y_0))}{(x_0-y_0)^2} &, &
\frac{-g'(y_0)(x_0-y_0)+(g(x_0)-g(y_0))}{(x_0-y_0)^2}
\end{array}
\right)$$ qui est bien de classe $C^1$ ($g$ \'etant de classe
$C^2$, $g'$ est de classe $C^1$). Montrons que $F$ est continue
aux points de la forme $(a,a)$. Le DL de $g$ \`a l'ordre $2$ entre
$x$ et $y$ donne $g(y)=g(x)+(y-x)g'(c_{x,y}) \mbox{ avec }c \in
[x,y]$ d'o\`u
$$\lim_{(x,y) \rightarrow (a,a) }\frac{g(x)-g(y)}{x-y}=
\lim_{(x,y) \rightarrow (a,a) }g'(c_{x,y})=g'(a)=F(a,a)$$ car
comme $(x,y)$ tend vers $(a,a)$, $x$ et $y$ tendent tous les deux
vers $a$ et donc $c_{x,y}$ aussi (et $g'$ est continue).
 Pour montrer que $F$ est $C^1$ (sachant que $F$ est continue), il suffit de
montrer que la diff\'erentielle de $F$ se prolonge par
continuit\'e sur $\mathbb{R}^2$.
 Le DL de $g$ \`a l'ordre $2$ entre $x_0$ et $y_0$ est:
$$g(x_0)=g(y_0)+(x_0-y_0)g'(y_0)+\frac{(x_0-y_0)^2}{2}g^{(2)}(c_1) \mbox{ avec } c_1  \in [x_0,y_0].$$
$$g(y_0)=g(x_0)+(y_0-x_0)g'(x_0)+\frac{(y_0-x_0)^2}{2}g^{(2)}(c_2) \mbox{ avec } c_2  \in [x_0,y_0].$$
On a donc
$$\frac{\partial f}{\partial x}(x_0,y_0)=\frac{(x_0-y_0)^2g^{(2)}(c_2)}{2(x_0-y_0)^2}=\frac{g^{(2)}(c_2)}{2}$$
et
$$\frac{\partial f}{\partial x}(x_0,y_0)=\frac{(x_0-y_0)^2g^{(2)}(c_1)}{2(x_0-y_0)^2}=\frac{g^{(2)}(c_1)}{2}$$
La fonction $g$ \'etant de classe $C^2$, on a $$\lim_{(x_0,y_0)
\rightarrow (a,a)} Df(x_0,y_0)= \left( \begin{array}{ccc}
g^{(2)}(a)/2 &, & g^{(2)}(a)/2
\end{array}
\right)$$ et donc $Df$ se prolonge par continuit\'e sur tout
$\mathbb{R}^2$. $F$ est donc bien de classe $C^1$.
\fincorrection
\correction{002506}
Soit $F_1(P)=\int_0^1 P^3-P^2 dt$, et soit $h$ un polyn\^ome de
degr\'e $n$ alors $$F_1(P+h)-F_1(P)= \int_0^1
[(P^3+3P^2h+3Ph^2+h^3)+(P^2+2Ph+h^2) -P^3-P^2 ]dt=$$
$$\int_0^1h(3P^2+2P)dt +\int_0^13Ph^2+h^3+h^2dt$$
Or $|\int_0^13Ph^2+h^3+h^2dt|=o(||h||_\infty)$ donc
$$DF_1(h)=\int_0^1(3P^2+2P)hdt.$$
Soit $F_2(P)=P'-P^2$ et soit $h$ un polyn\^ome de degr\'e $n$
alors $$F_2(P+h)-F_2(P)=(P+h)'-(P+h)^2-P'+P^2 =h'-2Ph-h^2$$ Or
$h^2=o(||h||)$ (pour toute norme a choisir). On a donc
$$DF_2(h)=h'-2Ph.$$
\fincorrection
\correction{002507}
\begin{enumerate}
\item On a $g(x,y)=<f(x,y)-a,f(x,y)-a>$ o\`u $<.,.>$ est le
produit scalaire Euclidien sur $\mathbb{R}^2$. L'application $g$
est diff\'erentiable en tant que compos\'ee et produit de
fonctions diff\'erentiables. La diff\'erentielle $Df$ est donn\'e
par la matrice Jacobienne
$$(\frac{\partial f(x,y)}{\partial x}, \frac{\partial
f(x,y)}{\partial y})$$ et $Dg$ par la matrice
$$(\frac{\partial g(x,y)}{\partial x}, \frac{\partial
g(x,y)}{\partial y})$$On a alors
$$\frac{\partial g(x,y)}{\partial
x}=\frac{\partial}{\partial x} <f(x,y)-a,f(x,y)-a>=$$
$$<\frac{\partial}{\partial
x}(f(x,y)-a),f(x,y)-a>+<f(x,y)-a,\frac{\partial}{\partial
x}(f(x,y)-a)>=$$
$$2<\frac{\partial f(x,y)}{\partial x},f(x)-a>.$$
De m\^eme, $$\frac{\partial g(x,y)}{\partial y}=2<\frac{\partial
f(x,y)}{\partial y},f(x)-a>.$$ \item L'application $f$ est
continue (car diff\'erentiable) et tend vers l'infini quand
$(x,y)$ tend vers l'infini. Ainsi
$$\forall A >0, \exists B >0, ||(x,y)|| \geq B \Rightarrow
||f(x,y)||\geq A.$$ Soit $m=\inf_{(x,y)\in \mathbb{R}^2} g(x,y)$,
pour $A=m+1$, il existe $B >0$ tel que $$||(x,y)|| \geq B
\Rightarrow g(x,y)=||f(x,y)||^2 \geq A^2\geq (m+1)^2 \geq m+1.$$
On a donc $$m=\inf_{(x,y)\in \mathbb{R}^2} g(x,y)=\inf_{||(x,y)||
\leq B} g(x,y).$$ Or la boule $\overline B(0,B)$ \'etant compacte
et $g$ continue, l'inf y est ateint en un point $X_0=(x_0,y_0) \in
B(0,B) \subset \mathbb{R}^2$. Comme $X_0$ est un minimum global de
$g$, c'est aussi un minimum de la restriction de $g$ sur toute
droite passant par $X_0$. Comme la d\'eriv\'e d'une fonction
r\'eelle en un minimum est nulle, toute les d\'eriv\'ees
partielles de $g$ sont nulles et donc $Dg(X_0)=0$ et par
cons\'equent la matrice jacobienne de $g$ est nulle. On a donc
$$\frac{\partial g}{\partial x}(x_0,y_0)=2<\frac{\partial
f}{\partial x}(x_0,y_0),f(x)-a>=0 \mbox{ et } \frac{\partial
g}{\partial y}(x_0,y_0)=2<\frac{\partial f}{\partial
y}(x_0,y_0),f(x)-a>=0.$$ Comme $Df$ est injective, ses colonnes
forment une base de $\mathbb{R}^2$. Par cons\'equent les
projections de $f(x)-a$ sur la base $(\frac{\partial f}{\partial
x}(x_0,y_0);\frac{\partial f}{\partial x}(x_0,y_0))$ sont nulles
et donc $$f(x_0,y_0)-a=0 \Leftrightarrow f(x_0,y_0)=a$$ et donc
$a$ admet bien un ant\'ec\'edent. Ceci \'etant valable pour tout
$a\in \mathbb{R}^2$, on a montr\'e que $f$ est surjective.
\end{enumerate}
\fincorrection
\correction{002508}
\begin{enumerate}
\item  Pour montrer que $||Df(x)||_{{\cal L}(\Rr^n,\Rr)}\leq 1$,
il faut montrer que si $h \in \mathbb{R}^n$, on a $|Df(x).h|\leq
||h||$. On a $$|Df(x).h|=|D_hf(x)|=|\lim_{t \rightarrow
0}\frac{f(x+th)-f(x)}{t}|.$$ Or $f$ est $1$-lipschitzienne et donc
$|f(x+th)-f(x)|\leq ||th||=t||h||.$ Par cons\'equent pour tout $h
\in \mathbb{R}^n$, $|Df(x).h|\leq ||h||$ ce qui donne
l'in\'egalit\'e demand\'ee. \item $$\varphi'(0)=\lim_{t
\rightarrow 0}\frac{\varphi(t)-\varphi(0)}{t-0}=\lim_{t
\rightarrow 0}\frac{f((1-t)x+ty)-f(x)}{t}=$$ $$\lim_{t \rightarrow
0}\frac{f(x+t(y-x))-f(x)}{t}=Df(x).(y-x).$$ Ou encore, soit $\psi:
\mathbb{R} \rightarrow \mathbb{R}^n$ l'application
$\psi(t)=(1-t)x+ty$, on a alors $\varphi(t)=f \circ \psi$ et
d'apr\`es la formule de diff\'erentielle d'une composition:
$$\varphi'(0)=Df(\psi(0)).D\psi(0)=Df(x).(y-x).$$
Or,
$$d(x,F)=d(x,y)=||x-y||=\frac{1}{1-t}||(1-t)(x-y)||=$$
$$\frac{1}{1-t}||[(1-t)x+ty]-[ty+(1-t)y]||
=\frac{1}{1-t}d((1-t)x+ty,y).$$ Notons $x_t=(1-t)x+ty$, on a alors
$$d(x_t,y)=(1-t)d(x,F).$$
Or, $\varphi(t)=d(x_t,F) \leq d(x_t,y) \leq (1-t)d(x,y) \leq
\varphi(0)$ et donc
$$|\varphi'(0)|=\lim_{t
\rightarrow 0}\frac{|\varphi(t)-\varphi(0)|}{t} =\lim_{t
\rightarrow 0}\frac{\varphi(0)-\varphi(t)}{t}\geq$$
$$ \lim_{t \rightarrow 0}\frac{d(x,y)-(1-t)d(x,y)}{t} \geq
d(x,y)=||x-y||.$$ Donc
$$|Df(x)(x-y)| \geq ||x-y||$$ d'o\`u la deuxi\`eme inégalité.
\item Raisonnons par l'absurde. Supposons qu'il existe deux point
$y_1$ et $y_2$ tels que $d(x,F)=d(x,y_1)=d(x,y_2).$ Alors, de la
même manière que précédement, on a
$Df(x).(x-y_1)=Df(x).(x-y_2)=d(x,F)$ et donc
$Df(x).(x-y_1+x-y_2)=2d(x,F)$. Or, $||x-y_1+x-y_2||< 2d(x,F)$ car
les vecteurs $x-y_1$ et $x-y_2$ ne sont pas alignés. Mais alors
cela contredit le fait que $||Df(x)||_{\cal
L(\mathbb{R}^n,\mathbb{R})}=1$.
\end{enumerate}
\fincorrection
\nocorrection
\nocorrection
\nocorrection
\nocorrection
\nocorrection
\nocorrection
\nocorrection
\nocorrection
\nocorrection
\correction{002518}
Montrons que $f$ se prolonge par continuit\'e au point $b$, on
montrera alors que $f$ est d\'erivable \`a gauche au point $b$ est
que cette d\'eriv\'ee est $\lim_{x \rightarrow b^-}f'(x)$. Pour
cel\`a montrons qu'il existe un r\'eel $k$ tel que toute suite
$\{x_n\}$ tendant vers $b$ v\'erifie $\lim_{n\rightarrow \infty}
f(x_n)=k$. Remarquons que la d\'eriv\'ee $f'(x)$ admettant une
limite au point $b$, elle est born\'ee sur un petit voisinage (\`a
gauche) de $b$ (notons $M$ ce majorant).
 Soit $y_n$ une suite
convergent vers $b$. Alors la suite $f(y_n)$ est de Cauchy. En
effet, pour tout $\epsilon >0$, posons
$\epsilon'=\frac{\epsilon}{2M}$. La suite $\{y_n\}$ \'etant de
cauchy, $$\exists N \in \mathbb{N}, p,q \geq N \Rightarrow
|y_p-y_q| \leq  \epsilon' \leq \frac{\epsilon}{2M}.$$ Or d'apr\`es
les accroissements finis:
$$f(y_p)-f(y_q)=(y_p-y_q)f'(c_{p,q}) \mbox { o\`u } c_{p,q} \in
]y_p,y_q[.$$ Par cons\'equent, $$|f(y_p)-f(y_q)|\leq
|y_p-y_q|.|f'(c_{p,q})| \leq \frac{\epsilon}{2M}M \leq
\frac{epsilon}{2} < \epsilon$$ et donc la suite $\{f(y_n)\}$ est
de cauchy et converge vers un r\'eel que nous noterons $l$.
Montrons que c'est le cas pour toute autre suite $\{x_n\}$ qui
tend vers $b$. On a $$f(x_n)=f(x_n)-f(y_n)+f(y_n).$$ D'apr\`es les
accroissements finis,  $|f(x_n)-f(y_n)| \leq M.|x_n-y_n|$ et donc
tend vers zero car les suites $x_n$ et $y_n$ tendent vers $b$. De
plus, comme on l'a vu, $f(y_n)$ tend vers $k$ et donc $f(x_n)$
aussi. Prolongeons $f$ par continuit\'e au point $b$ en posant
$f(b)=k$. On a alors le taux d'accroissement
$$T_xf=\frac{f(b)-f(x)}{b-x}=\frac{(b-x)f'(c_x)}{b-x}=f'(c_x)
\mbox{ o\`u } c_x \in ]x,b[.$$ Quand $x$ tend vers $b$, $c_x$
aussi et donc $T_xf$ tend vers $l$.
\fincorrection
\correction{002519}
On a $f'(x)=ie^{ix}$ (on peut le v\'erifier en coordonn\'ees). Si
l'\'egalit\'e des accroissement finis \'etait v\'erifi\'ee il
existerait $$c \in ]0,\pi[ \mbox{ tel que }
f(\pi)-f(0)=(\pi-0)ie^{ic}$$ ce qui est impossible car en prenant
les modules on trouverait $2=\pi$.
\fincorrection
\correction{002520}
\begin{enumerate}
\item $f$ est de classe $C^\infty$ car ses coordonn\'ees le sont
(polynômes). $g$ l'est car c'est la compos\'ee de deux fonctions
$C^\infty$. 
\item La matrice jacobienne de $f$ est:
$$Df(x,y)=\left(\begin{array}{cc}
2x & -1 \\ 2x & 2y
\end{array}\right)
$$ D'ap\`es la formule de diff\'erentielle d'une compos\'ee, on a
$$Dg(x,y)=Df(f(x,y))\circ Df(x,y).$$
Or $f(0,0)=0$ et $$Df(0,0)=\left(\begin{array}{cc}0& -1 \\0& 0
\end{array}\right)$$
et donc $$Dg(0,0)= \left(\begin{array}{cc}0& -1 \\0& 0
\end{array}.\right)\left(\begin{array}{cc}0& -1 \\0& 0
\end{array}\right)=0.$$
\item Par continuit\'e de $Dg(x,y)$ \`a l'origine et en prenant
$\epsilon=1/2$ on a:
$$\exists \rho> 0, ||(x,y)-(0,0)||\leq \rho \Rightarrow
||Dg(x,y)-Dg(0,0)|| \leq 1/2$$ d'o\`u le r\'esultat demand\'e.
\item D'apr\`es les accroissements finis, pour tous $X,Y \in
\mathbb{R}^2$, on a
$$||g(X)-g(Y)|| \leq \sup_{Z \in \overline B_\rho((0,0))}||Dg(Z)||.||X-Y|| \leq 1/2||X-Y||$$
et donc $g$ est contractante. Le Boule $\overline B_\rho((0,0))$
la boule $\overline B_\rho((0,0))$ \'etant compacte et compl\`ete,
le th\'eor\`eme du point fixe permet de conclure.
\end{enumerate}
\fincorrection
\correction{002521}
\begin{enumerate}
\item On a
$$Df(x,y)=\left(\begin{array}{cc}
-\sin x & - \cos y \\ \cos x & \sin y
\end{array}\right)$$

On a
$$|||Df(x,y)|||=\sup_{(a,b) \in \mathbb{R}^2\backslash
\{(0,0)\}}\frac{ ||Df(x,y).(a,b)||}{||(a,b)||}=$$
$$\frac{\sqrt{a^2 \sin^2 x + b^2 \cos^2 y +2ab \sin x \cos x + a^2 \cos^2 x + b^2 \sin^2 y + 2ab \cos x \sin y}}{\sqrt{a^2+b^2}}=$$
$$\frac{\sqrt{a^2+b^2+2ab \sin(x+y)}}{\sqrt{a^2+b^2}} \leq \sqrt{1+\frac{2|a||b|}{a^2+b^2}} \leq \sqrt 2$$
car $$(|a|-|b|)^2 \geq0\Rightarrow a^2+b^2 \geq 2 |a||b|.$$
\item Soient $U_n=(x_n,y_n)$ et $G(x,y)=1/2F(x,y)$, alors $|||G|||
\leq \frac{\sqrt 2}{2}$ et $U_{n+1}=G(U_n)$. D'apr\`es les
accroissements finis, $G$ est contractante et donc le th\'eor\`eme
du point fixe donne le r\'esultat demand\'e.

\end{enumerate}

\fincorrection
\nocorrection
\nocorrection
\nocorrection
\nocorrection
\correction{002526}
Appliquer le th\'eor\`eme des accroissements finis \`a
$g(x)=f(x)-Df(a)x$ en remarquant que la matrice jacobienne de
$Df(a)x$ est la matrice $Df(a)$.
\fincorrection
\nocorrection
\nocorrection
\nocorrection
\correction{006258}
\begin{enumerate}
\item  Puisque $\Rr$ est un sous-corps de $\Cc$, il est clair que $\varphi$
est $\Rr$-linéaire dès qu'elle est $\Cc$-linéaire et qu'elle vérifie en
particulier $\varphi(ix)=i\varphi(x)$ pour tout $x\in
E$. 

Supposons maintenant
$\varphi$ $\Rr$-linéaire, vérifiant $\varphi(ix)=i\varphi(x)$ pour tout $x\in
E$. Par hypothèse, $\varphi$ est additive et $\varphi(t.x)=t\varphi(x)$ pour
tout réel $t$ et $x\in E$. D'autre part, si $\lambda=\alpha+i\beta\in\Cc$,

$\varphi((\alpha+i\beta)x)=\varphi(\alpha x)+\varphi(i\beta x)$ par
additivité car $E$ est un $\Cc$-ev, 

\quad $=\varphi(\alpha x)+ i\varphi(\beta x)$ par hypothèse sur $\varphi$,

\quad $=\alpha\varphi(x)+i\beta\varphi(x)$ car $\alpha,\beta\in\Rr$.

\medskip

Si $\varphi\in{\cal L}(\Rr^2)$, elle se représente dans la base canonique
de $\Rr^2$ par la matrice 

$$\begin{pmatrix}a&b\\ c&d \end{pmatrix}$$

De même, $\cal I$, la multiplication par $i$ comme application linéaire de 
${\cal L}(\Cc)$ se représente dans l'identification de $\Cc$ avec $\Rr^2$
par la matrice

$$\begin{pmatrix}0&-1\cr 1&0 \cr \end{pmatrix}$$

et la condition  $\varphi(ix)=i\varphi(x)$ signifie que $\varphi$ commute
avec $\cal I$ ou que 

$$\begin{pmatrix}a&b\cr c&d \cr\end{pmatrix}
\begin{pmatrix}0&-1\cr 1&0 \cr \end{pmatrix}=
\begin{pmatrix}0&-1\cr 1&0 \cr\end{pmatrix}
\begin{pmatrix}a&b\cr c&d \cr \end{pmatrix}.$$

Ceci est réalisé si et seulement si $a=d$, et $b=-c$.

\smallskip

b) $f$ $\Cc$-différentiable au point $z=x+iy\in \Cc$ signifie : 
$$f(z+h)-f(z)-f'(z).h=h\varepsilon(h),$$ avec $h\in\Cc$ et $\lim_{h\to
0}\varepsilon(h)=0$ (l'application 
$\Cc$-linéaire tangente est ici la multiplication dans $\Cc$ par $f'(z)$);
traduit en variables réelles cela signifie :

$$\left \{
\begin{array}{ccc} 
u(x+h_1,y+h_2)-u(x,y)&=&ah_1-bh_2+ ||h||
\varepsilon(h)\\
v(x+h_1,y+h_2)-u(x,y)&=&bh_1+ah_2+ ||h||
\varepsilon(h)
\end{array}\right.$$
avec $f'(z)=a+ib$ et $h=h_1+ih_2$; 
$f$ est donc $\Rr$-différentiable au point $(x,y)$  et sa
matrice jacobienne en ce point vaut 
$$\begin{pmatrix}a&-b\cr b&a \cr \end{pmatrix}.$$

Réciproquement supposons $f\  \Rr$-différentiable
en $(x,y)$; ainsi
$$f(x+h_1,y+h_2)-f(x,y)-f'(x,y).h=||h||
\varepsilon(h),$$
avec
$h\in\Rr^2$ et $\lim_{h\to
0}\varepsilon(h)=0$; la matrice de $f'(x,y)$ est la matrice jacobienne 

$$\begin{pmatrix}\partial_1u(x,y)&\partial_2u(x,y)\cr \partial_1v(x,y)&
\partial_2v(x,y) \cr \end{pmatrix}=\begin{pmatrix}a&b\cr c&d \cr \end{pmatrix},$$ et $f$ est  
$\Cc$-différentiable en
$(x,y)$ 
si et seulement si $f'(x,y)$ est $\Cc$-linéaire ou  $a=d$ et $b=-c$, ce qui se
traduit par les conditions de Cauchy : 
$$\partial_1u(x,y)=\partial_2v(x,y),\ \partial_2u(x,y)=-\partial_1v(x,y).$$

Il est facile de voir que les applications de $\Cc$ dans
$\Cc$ : $f_1(z)=e^x,\ f_2(z)= x^2+y^2,\ f_3(z)=e^{x-iy}$ sont 
$\Rr$-différentiables, et que les conditions de Cauchy ne sont jamais vérifiées
pour $f_1$ et $f_3$, ne sont vérifiées qu'en $0$ pour $f_2$.

\item La fonction $g :\Cc\to\Cc$ peut s'écrire $g_1+ig_2$ où
$g_1=\sqrt{u^2+v^2}$ et $g_2=0$ si
$f=u+iv$. Supposons que $g$ soit 
$\Cc$-différentiable en $z=x+iy$: elle remplit donc les conditions de
Cauchy en ce point et $\partial_1g_1(x,y)=\partial_2g_1(x,y)=0$, soit
$$\left \{\begin{array}{ccc}
u(x,y)\ \partial_1u(x,y)+v(x,y)\ \partial_1v(x,y)&=&0,\\
u(x,y)\ \partial_2u(x,y)+v(x,y)\ \partial_2v(x,y)&=&0.
\end{array}\right.
$$

Mais comme $f$ est  
$\Cc$-différentiable en $z=x+iy$, $\partial_1u(x,y)=\partial_2v(x,y),$ 
 
$\partial_2u(x,y)=-\partial_1v(x,y),$ et le système devient
$$\left \{\begin{array}{ccc}
u(x,y)\ \partial_1u(x,y)+v(x,y)\ \partial_1v(x,y)&=&0,\\
-u(x,y)\ \partial_1v(x,y)+v(x,y)\ \partial_1u(x,y)&=&0.
\end{array}\right.
$$

C'est un système de Cramer puisque le déterminant
$u^2(x,y)+v^2(x,y)=|f(z)|^2\not=0$;
ainsi $\partial_1u(x,y)=\partial_1v(x,y)=\partial_2u(x,y)=\partial_2v(x,y)=0$ et
$f'(z)=0$.  
\end{enumerate}
\fincorrection
\nocorrection
\nocorrection
\nocorrection
\nocorrection
\nocorrection
\nocorrection
\nocorrection
\nocorrection
\correction{002553}
\begin{enumerate}
\item Calculons l'accroissement:
$$f(x+h)-f(x)=f(x)+f(h)-f(x)=f(h)+0.$$
Or, par d\'efinition $f(h)$ est lin\'eaire en $h$, continue et
$0=o(||h||)$. Par cons\'equent $f$ est diff\'erentiable et
$$Df(x)=f, \mbox{ ou encore } Df(x).h=f(h).$$
On remarque que $Df$ est l'application constante que à $x\in E$
associe l'application lin\'eaire $f$. Par cons\'equent, $Df$ est
diff\'erentiable et sa diff\'erentielle est nulle: $$D^2f=0.$$
\item Calculons $$f((x,y)+(h,k))-f(x,y)=f(x+h,y+k)-f(x,y)=
f(x,y+k)+f(h,y+k)-f(x,y)=$$
$$f(x,y)+f(x,k)+f(h,y)+f(h,k)-f(x,y)=f(x,k)+f(h,y)+f(h,k).$$
L'application qui \`a $(x,y)$ associe l'application lin\'eaire
$Df(x,y)(h,k)=f(x,k)+f(h,y)$ est donc candidate pour être la
diff\'erentielle de $f$. V\'erifions qu'elle est bien continue et
que $f(h,k)=o(||(h,k)||)$. Nous rappelons qu'une application
bilin\'eaire $f(x,y)$ est continue s'il existe $M>0$ tel que
$$\forall (x,y) \in E^2; ||f(x,y)||\leq M||x||||y||.$$
On a $$||Df(x,y)(h,k)||=||f(x,k)+f(h,y)|| \leq
||f(x,k)||+||f(h,y)||\leq$$ $$M||x||.||k||+M||h||.||y|| \leq M
(||x||+||y||)\max(||k||,||h||)\leq M(||x||+||y||)||(h,k)||.$$ Par
cons\'equent, $Df(x,y)$ est continue et a une norme inf\'erieure
\`a $M(||x||+||y||)$. De plus
$$||f(h,k)||\leq M||h||.||k|| \leq ||(h,k)||.\epsilon(h,k)$$ o\`u
$\epsilon$ tend vers zero quand $(h,k)$ tend vers zero car
$$\epsilon(h,k)=\frac{||h||.||k||}{\sup(||h||,||k||)}$$ ce qui
fini de montrer que $f$ est diff\'erentiable et que sa
diff\'erentielle est d\'efinie par
$$Df(x,y).(h,k)=f(x,k)+f(h,y).$$
En remarquant que $Df$ est lin\'eaire par rapport \`a $(x,y)$,
d'apr\`es la premi\`ere question, on d\'eduit que sa
diff\'erentielle est $$D^2f(x,y)[(h,k),(u,v)]=f(u,k)+f(h,v).$$
\item $$f(A+h)-f(A)=(A+h)^2-A^2=Ah+hA+h^2$$ avec $Ah+hA$
lin\'eaire en $h$ (et en $A$) et $||h^2|| \leq ||h||^2=o(||h||)$.
Par cons\'equent $f$ est diff\'erentiable et sa diff\'erentielle
est $Df(A).h=Ah+hA$. Comme $Df(A)$ est lin\'aire par rapport \`a
$A$, sa diff\'erentielle en $A$ est l'application bilin\'eaire
$$D^2f(A)[H,K]=KH+HK.$$
\end{enumerate}

\fincorrection
\correction{002554}
Soit $f(x,y)=x^2+xy+y^2+\frac{1}{4}x^3$, calculons la jacobienne
de $f$: $$Df(x,y)=(2x+y+\frac{3}{4}x^2, x+2y).$$ Les points
critiques de $f$ v\'erifient $Df(x,y)=0$ et par cons\'equent
v\'erifient les \'equations $2x+y+\frac{3}{4}x^2=0$ et $x+2y=0$.
Par cons\'equent $f$ admet deux points critiques $(0,0)$ et
$(-2,1)$. Pour savoir si ces points critiques sont des extr\'emums
de $f$, il faut \'etudier la hessienne de $f$:
$$Hess_f(x,y)= \left(
\begin{array}{cc}
2+\frac{3}{2}x & 1 \\
1 & 2
\end{array}\right)$$ et ses valeurs propres.
Au point $(0,0)$, $$Hess_f(0,0)= \left(
\begin{array}{cc}
2 & 1 \\
1 & 2
\end{array}\right)$$ a pour polynôme caractéristique
$P(\lambda)=(1-\lambda)(3-\lambda)$ et ses deux valeurs propres
sont strictement positive. La fonction $f$ admet donc un minimum
local au point $(0,0)$. Ce minimum n'est pas globale car
$f(0,0)=0$ et $f$ prend des valeurs n\'egatives pour $y=0$ et $x$
qui tend vers $-\infty$.\par
Au point $(-2,1)$, $$Hess_f(0,0)=
\left(
\begin{array}{cc}
-1 & 1 \\
1 & 2
\end{array}\right)$$ a pour polynôme caractéristique
$P(\lambda)=\lambda^2-\lambda-3$. On peut alors soit calculer les
valeurs propres soit remarquer que le determinant de la hessienne
est \'egal au produit des deux valeurs propres (ici $-3$) et
celles-ci sont donc non nulles et de signe contraire. Par
cons\'equent le point $(-2,1)$ est point selle de $f$. ce n'est
pas un extremum.
\fincorrection
\correction{002555}
Le volume d'une boite \'etant invariant par rotations, on peut
toujours supposer que toutes les boites sont centr\'ees \`a
l'origine et on des cot\'es parall\`eles aux axes de
coordonn\'ees. Par cons\'equent, la donn\'ee d'un point $(x,y,z)$
sur la sph\`ere d\'efinit de mani\`ere unique une boite
rectangulaire dont l'un des sommets est le point $(x,y,z)$. On
prendra $x,y$ et $z$ positifs car une telle boite a toujours un
sommet dans le secteur positif de l'espace. Par cons\'equent, on
doit maximiser la fonction volume $g(x,y,z)=8xyz$ sur la
sous-vari\'et\'e $S$ d\'efinie par l'\'equation $f=0$ avec
$f(x,y,z)=x^2+y^2+z^2-R^2$. Un point critique de $g$ sur $S$
v\'erifie $$Dg(x,y,z)=\lambda Df(x,y,z)$$ et $f(x,y,z)=0$. On a
obtient alors le syst\`eme d'\'equations:
$$\begin{array}{c}
8yz=2x \\
8xz=2y \\
8xy=2z \\
x^2+y^2+z^2=R^2
\end{array}$$

Par cons\'equent
$$\begin{array}{c}
8xyz=2x^2 \\
8xyz=2y^2 \\
8xyz=2z^2 \\
x^2+y^2+z^2=R^2
\end{array}$$
et donc comme $x, y$ et $z$ sont positifs, on a
$x=y=z=\frac{R}{\sqrt{3}}$. Or $g$ est continue et $S^+=S\cap \{x
\geq 0; y \geq 0; z \geq0\}$ est un compact. Comme $g$ est nulle
sur le bord de $S^+$, le maximum de $g$ est ateint en un point
critique de $g$ dans l'int\'erieur de $S^+$. Le seul point
critique de $g$ est donc bien ce maximum recherch\'e. Ici, il n'y
a pas eu besoin de calculer la hessienne de $g$ sur $S$ par la
formule: $$H=D^2g-\lambda D^2f$$ o\`u $\lambda$ est le coefficient
de lagrange trouv\'e pr\'ec\'edemment.
\fincorrection
\nocorrection
\nocorrection
\nocorrection
\nocorrection
\nocorrection
\nocorrection
\nocorrection
\nocorrection
\nocorrection
\nocorrection
\nocorrection
\nocorrection
\nocorrection
\nocorrection
\correction{002527}
\begin{enumerate}
\item La fonction $f$ \'etant d\'erivable, elle est continue.
Montrons qu'elle est injective. Soient $x,y \in \mathbb{R}$ tels
que $f(x)=f(y)$. Si $x \neq y$, d'apr\`es le th\'eor\`eme de
Rolle, il existe $c \in ]x,y[$ tel que $f'(c)=0$ ce qui contredit
le fait que $f'$ ne s'annule jamais. Par cons\'equent $x=y$ et $f$
est injective. Pour montrer que $f^{-1}$ est continue, il faut
montrer que l'image r\'eciproque par $f^{-1}$ d'un voisinage d'un
point est un voisinage de la r\'eciproque du point. Ou encore, que
l'image directe par $f$ d'un voisinage d'un point $a$ est un
voisinage de $f(a)$. Soit $V$ un voisinage de $a$, il contient un
intervalle du type $[a-\epsilon, a+\epsilon]$, l'image de ce
connexe par une fonction continue est encore connexe et est donc
un intervalle $[c,d]$ (ferm\'e car $f$ est continue et
l'intervalle de d\'epart est compact). Si $f(a) \in ]c,d[$, on la
d\'emonstration est finie. Si $f(a)=c$ ou $f(a)=d$ alors $a$ est
un extr\'emum local de $f$ et donc $f'(a)=0$ ce qui contredit
l'\'enonc\'e. Ainsi $f$ est un hom\'eormorphisme. $f^{-1}$ est
diff\'erentiable car la diff\'erentielle de $f$ ne s'annule pas
(th\'eor\`eme de deug...). \item Le taux d'accroissement
$$T_x(f)=\frac{f(x)-f(0)}{x-0}=1+x\sin{{\frac \pi x}}$$ tend vers $1$
quand $x$ tend vers zero ($\sin{\frac\pi  x}$ est born\'ee) et
donc $f$ est d\'erivable au point $0$ et $f'(0)=1 \neq0$. \item
Par l'absurde, supposons que $f$ soit inversible au voisinage de
$0$.Soit $\epsilon > 0$ tel que $]-\epsilon,\epsilon[$ soit inclus
dans ce voisinage. $f$ \'etant continue (car d\'erivable), elle
est strictement monotone. Or $f'(x)=1-\pi\cos{\frac\pi
x}+2x\sin{\frac\pi x}$. Prenons $k \in \mathbb{N}$ tel que
${\frac1{2k}} < \epsilon$ et ${\frac1{1+2k}} < \epsilon$ alors
$f'({\frac1{2k}})=1-\pi <0$ et $f({\frac1{2k+1}})=1+\pi >0$ et
donc $f$ n'est pas monotone sur $]-\epsilon,\epsilon[$ ce qui
donne la contradiction recherch\'ee. Le th\'eor\`eme de l'inverse
local nous montre de plus que $f$ n'est pas de classe $C^1$ dans
aucun voisinage de $0$.
\end{enumerate}
\fincorrection
\correction{002528}
\begin{enumerate}
\item L'application $\varphi(r,\theta)=(r\cos \theta, r \sin
\theta)$ est de classe $C^\infty$ car ses coordonn\'ees le sont.
Pour montrer que c'est un diff\'eormorphisme global, il suffit de
montrer que c'est un diff\'eo local (th\'eor\`eme de l'inverse
local) et qu'elle est bijective. Calculons la matric jacobienne de
$\varphi$:
$$D\varphi(r,\theta)= \left (
\begin{array}{cc}
\cos \theta & -r\sin \theta \\
\sin \theta & r \cos \theta
\end{array} \right )  $$
\end{enumerate}
La jacobienne de $\varphi$ est
$det(D\varphi(r,\theta))=r(\sin^2x+\cos^2x)=r >0$ si $(x,y) \neq
(0,0)$. L'application $\varphi$ est donc bien un diff\'eomorphisme
local au voisinage de chacun des point de
$]0,+\infty[\times]-\pi,\pi[$. La bijectivit\'e se v\'erifie en
explicitant par exemple la r\'eciproque de $\varphi$ (si on pose
$x=r\cos \theta$ et $y=r\sin \theta$, on pourra consid\'erer les
donn\'ees $x^2+y^2$ et $y/x$...).
\item

\fincorrection
\correction{002529}
\begin{enumerate}
\item  $\varphi$ a des coordonn\'ees de classe $C^1$, elle l'est
donc aussi. On a
$$Jac(\varphi)(x,y)=\left (\begin{array}{cc}
-1 & 1/2 \cos(y/2) \\
1/2 \cos(x/2) & -1
\end{array}\right )$$
On a  $det(Jac(\varphi)(x,y))=1-1/4\cos(x/2)\cos(y/2) \geq 3/4 >
0$. Par cons\'equent la jacobienne est inversible et
$D\varphi(x,y) \in Isom(\mathbb{R}^2,
\mathbb{R}^2)=GL(\mathbb{R}^2$. \item D'apr\`es le th\'eor\`eme de
l'inverse local, Il suffit de montrer que $\varphi$ est injective.
Supposons $\varphi(x_1,y_1)=\varphi(x_2,y_2)$, alors
$\sin(y_1/2)-x_1=\sin(y_2/2)-x_2$ et
$\sin(x_1/2)-y_1=\sin(x_2/2)-y_2$. D'o\`u
$\sin(y_1/2)-\sin(y_2/2)=x_1-x_2$ et
$\sin(x_1/2)-\sin(x_2/2)=y_1-y_2$. Or, $\forall a,b \in
\mathbb{R}, |\sin a -\sin b| \leq |a-b|$ (cons\'equence des
accroissements finis appliqu\'e \`a $\sin x$). Donc $|x_1-x_2|\leq
|y_1/2-y_2/2|$ et $|y_1-y_2| \leq |x_1/2-x_2/2|$ d'o\`u
$|x_1-x_2|\leq 1/4 |x_1-x_2| \Rightarrow x_1=x_2$  et $y_1=y_2$.
$\varphi: U \rightarrow F$ est injective. L'ensemble $f(U)$ est
ouvert car il est r\'eunion d'ouverts (d'apr\`es thm inverse
local). C'est un diff\'eomorphisme en $U$ et $\varphi(U)$. \item
Soient $(X_1,y_1), (X_2,Y_2) \in \varphi(\mathbb{R}^2)$ avec
$\varphi(x_1,y_1)=(X_1,Y_1)$ et $\varphi(x_2,y_2)=(X_2,Y_2)$ ou
encore $\varphi^{-1}(X_1,Y_1)=(x_1,y_1)$ et
$\varphi^{-1}(X_2,Y_2)=(x_2,y_2)$. On a
$$||\varphi^{-1}(X_1,Y_1)-\varphi^{-1}(X_2,Y_2)||=||(x_1,y_1)-(x_2,y_2)||=||(x_1-x_2,y_1,y_2)||=|x_1-x_2|+|y_1-y_2|$$
Or $\sin(y_1/2)-x_1=X_1$, $\sin(y_2/2)-x_2=X_2$,
$\sin(x_1/2)-y_1=Y_1$ et $\sin(x_2/2)-y_2=Y_2$. Par cons\'equent
$$x_1-x_2=\sin(y_1/2)-X_1-\sin(y_2/2)+X_2$$
$$y_1-y_2=\sin(x_1/2)-Y_1-\sin(x_2/2)+Y_2.$$
D'o\`u
$$|x_1-x_2|+|y_1-y_2| \leq |X_2-X_1|+|\sin(y_1/2)-\sin(y_2/2)|+
|Y_2-Y_1|+ |\sin(x_1/2)-\sin(x_2/2)|$$ $$ \leq
|X_2-X_1|+1/2|y_1-y_2|+|Y_2-Y_1|+1/2|x_1-x_2|$$ d'où
$$|x_1-x_2|+|y_1-y_2| \leq 2(|X_2-X_1|+|Y_2-Y_1|) \leq 2
||(X_1,Y_1)-(X_2,Y_2)||.$$ Donc $\varphi^{-1}$ est lipschitzienne.
\item Soit $(X_n,Y_n)$ une suite de cauchy dans
$\varphi(\mathbb{R}^2)$,  ($(X_n,Y_n)=\varphi(x_n,y_n)$;
$(x_n,y_n)=\varphi^{-1}(X_n,Y_n)$). Pour tout $\epsilon > 0,
\exists n \in \mathbb{N}, p,q \geq n \Rightarrow
||(X_p,Y_p)-(X_q,Y_q)|| < \epsilon.$ Par cons\'equent, $\forall
\epsilon > 0$, $\exists n \in \mathbb{N}; p,q \geq n \Rightarrow
||(x_p,y_p)-(x_q,y_q)||< 2 \epsilon$. La suite $(x_n,y_n)$ est
alors de cauchy dans $\mathbb{R}^2$, qui est complet. Par
cons\'equent elle converge. Soit $(x,y)$ sa limite. Comme
$\varphi$ est continue et que $\lim_{n}(x_n,y_n)=(x,y)$ alors
$\lim_{n} \varphi(x_n,y_n)=\varphi(x,y)$. La suite $(X_n,Y_n)$ est
une suite de Cauchy de $\mathbb{R}^2$. Elle converge. Soit $(X,Y)$
sa limite, alors $(X,Y)=\varphi(x,y)$ car $(X_n,Y_n) \rightarrow
(X,Y)$ et $\varphi(x_n,y_n) \rightarrow \varphi(x,y)$. Donc $(X,Y)
\in \varphi(\mathbb{R}^2)$. $\varphi(\mathbb{R}^2)$ est alors
complet et donc ferm\'e. Comme $\varphi(\mathbb{R}^2$ est un
ouvert ferm\'e et non vide (il contient $(0,0)=\varphi(0,0)$) dans
le connexe $\mathbb{R}^2$, on a
$\varphi(\mathbb{R}^2)=\mathbb{R}^2$. \item $p=(1-\pi/2,
\sqrt{2}/2 - \pi)=\varphi(\pi/2,\pi)=\varphi(q)$ o\`u
$q=(\pi/2,\pi)$. $\varphi: E \rightarrow F$ est un
$C^1$-diff\'eormorphisme donc $\varphi^{-1} \circ \varphi=Id$ et
donc $$Id=D(\varphi^{-1} \circ
\varphi)(q)=D\varphi^{-1}(\varphi(q))\circ D\varphi(q).$$ Or
$D\varphi^{-1}(\varphi(q))=(D\varphi(q))^{-1}$ et donc $Jac
\varphi^{-1}(p)= (Jac \varphi (\pi/2,\pi))^{-1}$. Or
$$Jac \varphi(x,y)= \left (
\begin{array}{cc}
-1 & 1/2 \cos(y/2) \\
1/2 \cos(x/2) & -1
\end{array} \right ).$$
D`o\`u $$Jac \varphi(\pi/2,\pi)= \left (
\begin{array}{cc}
-1 &0\\
\sqrt{2}/4 & -1
\end{array} \right )$$ et donc
$$Jac \varphi^{-1}(p)=
 \left (
\begin{array}{cc}
-1 &0\\
-\sqrt{2}/4 & -1
\end{array} \right ).$$
\end{enumerate}

\fincorrection
\correction{002530}
Posons $\theta(t)=a+t(b-a)$ et $\Psi(x)= \langle x,b-a \rangle$ qui est lin\'eaire
et continue (donc $C^\infty$).  
\begin{enumerate}
\item $f$ et $\varphi$ sont de classe $C^1$ car compos\'ees d'applications de classe $C^1$. On a
$$
D\varphi(t)= \varphi'(t)=(\Psi \circ f \circ \theta)(t)'(t)
= D\Psi(f(\theta(t)))\circ Df(\theta(t))\circ D\theta(t)
= \langle Df(a+t(b-a))(b-a),b-a \rangle.
$$ 
Par cons\'equent:
$$
\varphi'(t) \geq \alpha \langle b-a,b-a \rangle.
$$
Or, $\varphi(1)-\varphi(0)=\langle f(b)-f(a),(b-a)\rangle$ et il existe $t\in
]0,1[$ tel que $\varphi(1)-\varphi(0)=\varphi'(t)$ d'o\`u
$$
\langle f(b)-f(a),b-a\rangle \geq \alpha \langle b-a,b-a\rangle.
$$
Indications pour mq $f$ est ferm\'ee: Posons $\|x\|=\sqrt{\langle x,x\rangle}$
alors
$$
\alpha \|b-a\|^2 \leq \langle f(b)-f(a),b-a\rangle 
\leq \|f(b)-f(a)\|.\|b-a\|.
$$ 
D'o\`u $$\|b-a\| \leq 1/\alpha \|f(b)-f(a)\|.
$$ 
Soit $F$ un ferm\'e et $y_n$ une suite de points
de $f(F)$ convergeant vers un point limite $y_\infty$. Il faut
montrer que $y_\infty \in f(F)$. Soit $x_n$ une suite de points de
$\mathbb{R}^n$ tels que $f(x_n)=y_n$. Il reste \`a montrer que
cette suite admet est de Cauchy, qu'elle converge donc et que sa
limite $x_\infty$ v\'erifie $f(x_\infty)=y_\infty$.
\end{enumerate}
\fincorrection
\nocorrection
\nocorrection
\nocorrection
\nocorrection
\nocorrection
\nocorrection
\nocorrection
\nocorrection
\nocorrection
\nocorrection
\correction{002541}
Soit $(x_0,y_0,z_0) \in \mathbb{R}^3$ tel que
$f(x_0,y_0,z_0)=(0,0)$ (par exemple $(1,1,1)$). $f$ est $C^1$ car
coordonn\'ees polynomiales.

$$\text{Mat}\, D_2f(x_0,y_0,z_0)=
\left(
\begin{array}{cc}
\frac{\partial f_1}{\partial y}(x_0,y_0,z_0) &\frac{\partial
f_1}{\partial
z}(x_0,y_0,z_0)\\
\frac{\partial f_2}{\partial y}(x_0,y_0,z_0) & \frac{\partial
f_2}{\partial z}(x_0,y_0,z_0)
\end{array} \right )=\left(
\begin{array}{cc}
-2y_0 & 2z_0 \\
x_0z_0 & x_0y_0
\end{array} \right )$$

$\det(\text{Mat}\, D_2f(x_0,y_0,z_0))=-2x_0(y_0^2+z_0^2) \neq 0$ car
$x_0y_0z_0=1$ donc $x_0 \neq 0, y_0 \neq 0, z_0 \neq 0$. D'apr\`es
le th\'eor\`eme des fonctions implicites, il existe $I$ intervalle
contenant $x_0$ et $\varphi:I \rightarrow \mathbb{R}^2$ tel que
$f(x,\varphi(x))=0$ pour tout $x \in I$ et $\varphi(x_0)=(y_0,z_0)$.

\fincorrection
\nocorrection
\nocorrection
\nocorrection
\correction{002545}
Posons $f(x,y)=x^4+y^3-x^2-y^2+x-y$, $f(0,0)=0$ et $f(1,1)=0$.
$\mathbb{R}$ est un espace de Banach et $f$ est de classe $C^1$
car polynomiale.
$$\frac{\partial f}{\partial y}=3y^2-2y-1$$
\'Etude au point $(0,0)$, $\frac{\partial f}{\partial y}(0,0)=-1$,
c'est un isomorphisme de $\mathbb{R}$. Nous sommes dans les
conditions d'application du th\'eor\`eme des fonctions implicites.
Il existe $I$ contenant $0$, $J$ contenant $0$ et $g: I
\rightarrow J$, $C^1$ tel que $g(0)=0$ et $f(x,g(x))=0, \forall x
\in I$. On a
$$x^4+(g(x))^3-x^2-(g(x))^2+x-g(x)=0$$
En d\'erivant on obtient:
$$4x^3+3g^2(x)g'(x)-2x-2g(x)g'(x)+1-g'(x)=0$$
d'o\`u $g'(0)=1$. On d\'erive encore:
$$12x^2+6g(x)g'(x)^2+3g^2(x)g''(x)-2-2g'(x)^2-2g(x)g''(x)-g''(x)=0$$
d'o\`u $$g''(0)=-4.$$ \'Etude au point $(1,1)$, $\frac{\partial
f}{\partial y}(1,1)=0$. Ce n'est plus un diff\'eo, on ne peut pas
appliquer le th\'eor\`eme des fonctions implicites. Dans ce cas,
on prend la d\'eriv\'ee par rapport \`a la premì\`ere variable.
$$\frac{\partial f}{\partial x}=4x^3-2x+1$$
et donc $\frac{\partial f}{\partial x}(1,1)=3$. Donc, d'apr\`es le
th\'eor\`eme des fonctions implicites, il existe $I$ contenant
$1$, $J$ contenant $1$ et $g:I \rightarrow J$ de classe $C^1$ tels
que $g(1)=1$ et $f(g(x),x)=0, \forall y \in I$. On a
$$g(y)^4-g^2(y)+g(y)+y^3-y^2-y=0$$
En d\'erivant
$$4g^3g'-2gg'+g'+3y^2-2y-1=0$$
d'o\`u $4g'(1)-g'(1)=0$ et donc $g'(1)=0$.
$$12g^2(g')^2+4g^3g''-2gg''-2(g')^2+g''+6y-2=0$$
d'o\`u $g''(1)=-4/3$.
\fincorrection
\nocorrection
\nocorrection
\nocorrection
\nocorrection
\nocorrection
\nocorrection
\nocorrection
\nocorrection
\nocorrection
\nocorrection
\nocorrection
\nocorrection
\nocorrection
\nocorrection
\nocorrection
\nocorrection
\nocorrection
\nocorrection
\nocorrection
\nocorrection
\nocorrection
\nocorrection
\nocorrection
\nocorrection
\nocorrection
\nocorrection
\nocorrection
\nocorrection
\nocorrection
\nocorrection
\nocorrection
\nocorrection
\nocorrection
\nocorrection
\nocorrection
\nocorrection
\nocorrection
\nocorrection
\nocorrection
\nocorrection
\nocorrection
\nocorrection
\nocorrection
\nocorrection
\nocorrection
\nocorrection
\nocorrection
\nocorrection
\nocorrection
\nocorrection
\nocorrection
\correction{006883}
\begin{enumerate}
  \item Remarque: si $x\in F$, alors clairement $\bar{x}=x$ convient. \\
  Comme $F \neq \varnothing$, il existe un point $y_0 \in F$ et posons
  $r = \| x - y_0 \| \ge 0$ par hypothèse. On peut découper $F$
  en deux morceaux : $F_0 = F \cap \bar{B} (x, r)$ et $F_1 = F -
  F_0$. $F_0$ est non vide car il contient $y_0$. Quant aux points de $F_1$,
  ils sont tous à une distance de $x$ strictement supérieure à $r$ donc
  l'infimum de la distance sur $F$ est l'infimum sur $F_0$ :
  \[ \inf_{z \in F} \| x - z \| = \inf_{z \in F_0} \| x - z
     \| \le r \]
  Or $F_0$ est fermé et borné, donc compact (car la dimension est finie),
  et la fonction $z \mapsto \| x - z \|$ étant continue, elle est
  bornée et atteint son minimum en un point $y$. 
  \\
  Question subsidiaire: $\bar{x}$ n'est pas unique (sauf dans certains cas, par exemple si $F$ est convexe);
  contre-exemple: $F$ est un cercle et $x$ est son centre.
  
  \item
  \begin{enumerate}
    \item On voit que $\mathcal{M}_2 (\mathbb{R}) \simeq
    \mathbb{R}^4$ et que la fonction déterminant est $g : (a, b, c, d
) \mapsto ad - bc$. Cette fonction est continue car polynomiale en
    les coordonnées donc $\text{SL}_2 (\mathbb{R}) = g^{-1}(\{1\})$ est un fermé.
    
    \item La suite de matrices 
    $\displaystyle M_n = \begin{pmatrix} 1 & n \\ 0 & 1 \end{pmatrix}$ 
    est dans $\text{SL}_2(\mathbb{R})$ sans être bornée ($\| M_n \| = \sqrt{2+n^2}$).
    
    \item
    \begin{enumerate}
      \item C'est une conséquence directe du (1) en prenant pour $x$ la matrice nulle.
      
      \item Écrivons $f (M) = \sqrt{a^2 + b^2 + c^2 + d^2}$,
      alors $\nabla f (a, b, c, d) = \frac{1}{f (M)}
      (a, b, c, d)$, défini à condition que $M$ ne soit pas
      la matrice nulle (qui n'appartient pas $\text{SL}_2 (\mathbb{R}
   )$ donc il n'y~a pas de problème).
      
      \item En un extremum on doit avoir $\nabla f = \lambda \nabla g$ (extrema liés) donc
      \[ \frac{1}{\| M \|}  \left(\begin{array}{c}
           a\\
           b\\
           c\\
           d
         \end{array} \right) = \lambda \left(\begin{array}{c}
           d\\
           - c\\
           - b\\
           a
         \end{array} \right) 
      \]
      ce qui n'est possible que si $\lambda = \epsilon / \| M \|$
      où $\epsilon \in \{ - 1, + 1 \}$, et que $d = \epsilon a$
      et $c = - \epsilon b$ donc $M = \left(\begin{array}{cc}
        a & - \epsilon b\\
        b & \epsilon a
      \end{array} \right)$. Si l'on rajoute la contrainte $\det M = 1$, il
      vient que $\epsilon = + 1$ et $a^2 + b^2 = 1$. Donc les seuls extrema
      possibles sont de la forme $M_{\theta} = \left(\begin{array}{cc}
        \cos \theta & - \sin \theta\\
        \sin \theta & \cos \theta
      \end{array} \right)$ avec $\theta \in \mathbb{R}$, et pour toutes ces
      matrices $\| M_{\theta} \| = \sqrt{2}$.
      
      Or d'après le (i) on sait que le minimum est atteint, et
      de plus il doit satisfaire la relation de Lagrange car les gradients
      sont bien définis. Il y a donc un infinité de minima (les rotations
      $M_{\theta}$) et ils prennent la même valeur
      \[ 
      \inf_{M \in \text{SL}_2 (\mathbb{R})}  \| M
      \| = \min_{M \in \text{SL}_2 (\mathbb{R})} 
      \| M \| = \sqrt{2} . 
      \]
    \end{enumerate}
  \end{enumerate}
\end{enumerate}
\fincorrection
\nocorrection
\nocorrection
\nocorrection
\nocorrection
\nocorrection
\nocorrection
\nocorrection
\nocorrection
\nocorrection
\nocorrection
\correction{002558}
\begin{enumerate}
\item Soit $\lambda_n$ et $\mu_n$ deux suites de points tels que
$\varphi(\lambda_n)=\varphi(\mu_n)=0$ et convergeant
respectivement vers $\lambda$ et $\mu$, il reste \`a montrer que
$\varphi$ est nulle sur chaque interval $[\lambda_n,\mu_n]$. Soit
$c$ un extremum de $\varphi$ sur cet interval, on a alors
necessairement $\varphi'(c)=0$ et donc $3c^{2/3}=0$ et donc $c=0$
et donc $\varphi(c)=\varphi(0)=0$. Par cons\'equent le sup et le
min de $\varphi$ sur $[\lambda_n,\mu_n]$ sont nuls et donc
$\varphi$ est aussi nulle sur cet intervalle. En passant \`a la
limite, on a prouv\'e que $\varphi$ est nulle sur $]\lambda,\mu[$.
\item On v\'erifie que les solutions propos\'ees v\'erifient
l'\'equation diff\'erentielle $(1)$. La fonction $x^{2/3}$ est
lipschitzienne par rapport \`a $x$ d\`es que $x\neq 0$. Si
$\varphi_2$ est une solution maximale sur $\mathbb{R}$ v\'erifiant
$\varphi_2(0)=0$, il existe alors n\'ecessairement $\lambda, \mu$
(d\'efinis pr\'ec\'edement) tels que $\varphi_2$ est nulle sur
$]\lambda,\mu[$. Par continuit\'e de la solution elle v\'erifie
$\varphi(\lambda)=\varphi(\mu)=0$. Mais alors
$\varphi_2'-\varphi=0$ est donc $\varphi_2=\varphi+K$ o\`u $K$ est
une constante donn\'ee. Du fait que
$\varphi_2(\lambda)=\varphi(\lambda)+K=0+K=0$, on a $K=0$ ce qui
termine la d\'emonstration.
\end{enumerate}
\fincorrection
\correction{002559}
\begin{enumerate}
\item Soit $f:\mathbb{R}\times \mathbb{R} \rightarrow \mathbb{R}$
telle que $f(t,x)=|x|+|t|$. $f$ est continue et Lipschitzienne par
rapport \`a la seconde variable. En effet,
$$|f(t,x)-f(t,y)|=||x|-|y||\leq |x-y|.$$
Remarquons que $x' \geq 0$ pour tout $t$ et que pour tout point
$(0,x_0$ passe une solution maximale unique $(\varphi,J)$. 
\item Prenons $x_0=1$, lorque $t\geq 0$; l'\'equation devient
$$x't()=x(t)+t$$ car $|t|=t$ et $x(t) \geq x(0) > 0, x(0)=1$.
Elle admet comme solution sur $[0,+\infty[$ avec $\varphi(0)=1$
$$\varphi(t)=2e^t-t-1.$$
Lorsque $t<0$, on distingue deux cas: premier cas $x(t)\geq0$;
$x'=-t+x(t)$ et alors $x(t)=ce^t+t+1$ avec $x(0)=1$ d'o\`u $c=0$
et $\varphi(t)=t+1$. Cela n'est valable que lorsque $\varphi(t)
\geq 0$, c'est \`a dire $t \geq 1$. Donc $\varphi(t)=t+1$ sur
$[-1,1]$. Deuxi\`eme cas: $x(t) \leq 0$, ceci a lieu lorsque $t
\leq -1$ car $\varphi$ croissante et $\varphi(-1)=0$. Nous avons
alors $\varphi'(t)=-t-\varphi(t)$. D'o\`u $\varphi(t)=ce^{-t}-t+1$
or $\varphi(-1)=ce+2=0$ d'o\`u $c=-2e^{-1}$ et
$\varphi(t)=-2e^{-t+1}-t+1$ sur $]-\infty,-1]$. La solution
maximale v\'erifiant $\varphi(0)=1$ est la suivante:
$$\varphi(t)=
\left( \begin{array}{c} 2e^t-t-1 \mbox{ sur } [0,+\infty[ \\
t+1 \mbox{ sur } [-1,0] \\
-2e^{-(t+1)}-t+1 \mbox{ sur } ]-\infty,-1]
\end{array} \right)$$
$$\varphi'(t)=
\left( \begin{array}{c} 2e^t-1 \mbox{ sur } ]0,+\infty[ \\
1 \mbox{ sur } ]-1,0[ \\
2e^{-(t+1)}-1 \mbox{ sur } ]-\infty,-1[
\end{array} \right)$$
En \'etudiant les limites de $\varphi'$ aux point $0$ et $-1$, on
voit que $\varphi'$ est continue sur $\mathbb{R}$.
$$\varphi''(t)=
\left( \begin{array}{c} 2e^t \mbox{ sur } ]0,+\infty[ \\
0 \mbox{ sur } ]-1,0[ \\
-2e^{-(t+1)} \mbox{ sur } ]-\infty,-1[
\end{array} \right)$$ $\varphi$ n'est donc pas deux fois
d\'erivable en $0$ et $-1$.
\end{enumerate}
\fincorrection
\correction{002560}
$f(t,x)=\frac{4t^3x}{t^4+x^2}$ (si $(t,x) \neq (0,0)$) est de
classe $C^\infty$ en tant que quotient, somme et produit de
fonctions $C^\infty$.
\begin{enumerate}
\item $|f(t,x)|=|2t|.|\frac{2t^2x=}{(t^2)^2+x^2}| \leq 2|t|
\rightarrow_{(t,x) \rightarrow 0}0=f(0,0)$. $f$ est donc continue
en $(0,0)$. $f$ n'est pas localement lipschitzienne au voisinage
de $(0,0)$ car sinon il existerait $k, \alpha, \beta \in
\mathbb{R}$ tels que $t\in ]-\alpha,\alpha[$, $x\in
]-\beta,\beta[$ et
$$|f(t,x)-f(t,0)| \leq k|x-0|$$
D'o\`u $\frac{4t^3x}{t^4+x^2} \leq kx \Rightarrow
\frac{4t^3}{t^4+x^2} \leq k \rightarrow \frac{4}{t} \leq k,
\forall t \in ]0,\alpha[$ ce qui est absurde. Nous ne pouvons pas
appliquer Cauchy-Lipschitz. 
\item $(\varphi,I)$ solution de $(2)$
avec $0\not \in I$, $$\psi(t)=t^{-2}\varphi(t) \Rightarrow
\psi'(t)=t^{-2}\varphi'(t)-2t^{-3}\varphi(t)$$
$$\psi'(t)=4t^{-2}\frac{t^3\varphi(t)}{t^4+\varphi^2(t)}-2t^{-1}\psi(t)$$
d'o\`u en exprimant tout en fonction de $\psi$:
$$\frac{\psi'(t)(1+\psi^2(t))}{\psi(t)(1-\psi(t))(1+\psi(t))}=\frac{2}{t}$$

Or $\frac{1+\psi^2(t)}{\psi(t)(1-\psi(t))(1+\psi(t))}=
\frac{1}{\psi(t)}+\frac{1}{1-\psi(t)}-\frac{1}{1+\psi(t)}$ d'o\`u
$$\psi'(t)(\frac{1}{\psi(t)}+\frac{1}{1-\psi(t)}-\frac{1}{1+\psi(t)})=\frac{2}{t}$$
En int\'egrant par rapport \`a $t$ on obtient:
$$ln|\frac{\psi(t)}{1-\psi^2(t)}|=ln(t^2)+c$$
d'o\`u $$\frac{\psi(t)}{1-\psi(t)}=ct^2.$$ $\psi(t)$ v\'erifie est
donc une racine de l'\'equation $$ct^2\psi^2(t)+\psi(t)-ct^2=0$$
et donc $$\psi(t)=\frac{-1 \pm \sqrt{1+4c^2t^4}}{2ct^2}$$ d'o\`u
$\varphi=\frac{-1 \pm \sqrt{1+4c^2t^4}}{2c}$.
\end{enumerate}
\fincorrection
\correction{002561}
\begin{enumerate}
\item Posons $y_1=x$, $y_2=x'=y_1'$, $y_3=x''=y_2'$. L'\'equation
devient $y_3'-y_1y_3=0$ et donc en posant $f(t,y_1,y_2,y_3)=
\left( \begin{array}{c} y_2 \\ y_3 \\ y_1y_3
\end{array}\right)$ l'\'equation s'\'ecrit
$$\left( \begin{array}{c} y_1' \\ y_2' \\ y_3'
\end{array}\right) = f(t,y_1,y_2,y_3).$$

\item $f$ \'etant de classe $C^\infty$, elle est lipschitzienne
par rapport \`a la deuxi\`eme variable $(y_1,y_2,y_3)$ et donc le
th\'eor\`eme de Cauchy-Lipschitz permet de conclure. 

\item La d\'eriv\'ee de la fonction donn\'ee est nulle. Par cons\'equent,
elle est constante et donc, l'exponentielle \'etant strictement
positive, le signe de $\varphi''$ est constant. Si cette constante
est strictement positive, $\varphi$ est convexe, si elle est
strictement n\'egative, $\varphi$ est concave. Si elle est nulle
$\varphi''=0$ et donc $\varphi(t)=at+b$ qui est bien une solution
de l'\'equation diff\'erentielle et v\'erifie $\varphi''(t_0)=0$.
L'unicit\'e montre que toutes les solutions qui v\'erifient
$\varphi''(t_0)=0$ sont bien de la forme $at+b$.
\end{enumerate}
\fincorrection
\nocorrection
\nocorrection
\nocorrection
\nocorrection
\nocorrection
\nocorrection
\nocorrection
\nocorrection
\nocorrection
\nocorrection
\nocorrection
\nocorrection
\nocorrection
\nocorrection
\nocorrection
\nocorrection
\nocorrection
\nocorrection
\nocorrection
\nocorrection
\nocorrection
\nocorrection
\nocorrection
\nocorrection
\nocorrection
\nocorrection
\nocorrection
\nocorrection
\nocorrection
\nocorrection
\nocorrection
\nocorrection
\nocorrection
\nocorrection
\nocorrection
\nocorrection
\nocorrection
\nocorrection
\nocorrection
\nocorrection
\nocorrection
\nocorrection
\nocorrection
\nocorrection
\nocorrection
\nocorrection
\nocorrection
\nocorrection
\nocorrection
\nocorrection
\nocorrection
\nocorrection
\nocorrection
\nocorrection
\nocorrection
\nocorrection
\nocorrection
\nocorrection
\nocorrection
\nocorrection
\nocorrection
\nocorrection
\nocorrection
\nocorrection
\nocorrection
\nocorrection
\nocorrection
\nocorrection
\nocorrection
\nocorrection
\nocorrection
\nocorrection
\nocorrection
\nocorrection
\nocorrection
\nocorrection
\correction{005933}
\begin{enumerate}
\item Montrons que $ [a, b] =\bigcap_{n=1}^{\infty}
]a-\frac{1}{n}, b+\frac{1}{n}[$.
\begin{itemize}
\item Pour tout $n\in\mathbb{N}$, on a $[a, b] \subset\,
]a-\frac{1}{n}, b+\frac{1}{n}[$. Donc $[a,
b]\subset\bigcap_{n=1}^{\infty} ]a-\frac{1}{n}, b+\frac{1}{n}[$.
\item Soit $x\in \bigcap_{n=1}^{\infty} ]a-\frac{1}{n},
b+\frac{1}{n}[$. Alors pour tout $n\in\mathbb{N}$, on a~:
$$a-\frac{1}{n} < x < b+\frac{1}{n}.$$ Ainsi
$$\lim_{n\rightarrow+\infty}(a-\frac{1}{n})\leq x \leq
\lim_{n\rightarrow+\infty}(b+\frac{1}{n}),
$$
c'est-\`a-dire $x\in[a, b]$. Donc $\bigcap_{n=1}^{\infty}
]a-\frac{1}{n}, b+\frac{1}{n}[ \subset [a, b]$ et on a
d\'emontr\'e l'\'egalit\'e entre ces deux ensembles.
\end{itemize}

\item Montrons que $]a, b[ = \bigcup_{n=1}^{\infty} [a+
\frac{1}{n}, b-\frac{1}{n}]$.
\begin{itemize}
\item Pour tout $n\in\mathbb{N}$, on a $[a+ \frac{1}{n},
b-\frac{1}{n}]\subset ]a,b[$, donc $\bigcup_{n=1}^{\infty} [a+
\frac{1}{n}, b-\frac{1}{n}]\subset]a,b[$. 
\item Soit
$x\in\bigcup_{n=1}^{\infty} [a+ \frac{1}{n}, b-\frac{1}{n}]$.
Alors il existe  $n\in\mathbb{N}$ tel que $x\in[a+ \frac{1}{n},
b-\frac{1}{n}]$. Ainsi $x\in]a, b[$ et $\bigcup_{n=1}^{\infty} [a+
\frac{1}{n}, b-\frac{1}{n}] \subset ]a, b[$, d'o\`u l'\'egalit\'e
de ces deux ensembles.
\end{itemize}
\end{enumerate}
\fincorrection
\correction{005934}
Soit $(\Omega, \Sigma, \mu)$ un espace mesur\'e et $f~:\Omega
\rightarrow \mathbb{R}$ une fonction
($\Sigma$-$\mathcal{B}(\mathbb{R})$)-mesurable. Montrons que la
troncature $f_A$ de $f$ d\'efinie par~:
$$
f_A(x) = \left\{\begin{array}{ll}-A & \text{si}\quad f(x) < -A\\
                                  f(x) & \text{si}\quad |f(x)| \leq A\\
                                  A & \text{si}\quad f(x)> A
\end{array} \right.
$$
est mesurable. Notons
$$
\begin{array}{l}
E_{1} := \{x\in\Omega~|~f(x)< -A\}= f^{-1}\left(]-\infty,-A[\right),\\
E_{2}:= \{x\in\Omega~|~|f(x)| \leq A\} = f^{-1}\left([-A, A]\right),\\
E_{3}:=\{x\in\Omega~|~f(x)>A\} = f^{-1}\left(]A, +\infty[\right).
\end{array}
$$
Comme $]-\infty, -A[$, $[-A, A]$, $]A, +\infty[$ appartiennent \`a
la tribu bor\'elienne et $f$ est
($\Sigma$-$\mathcal{B}(\mathbb{R})$)-mesurable, les ensembles
$E_{1}$, $E_2$, et $E_3$ appartiennent \`a $\Sigma$. Alors $f_{A}
= f\cdot \mathbf{1}_{E_2} -A \cdot \mathbf{1}_{E_1} + A \cdot\mathbf{1}_{E_{3}}$ est
mesurable comme somme de produits de fonctions mesurables.
\fincorrection
\correction{005935}
Soit $\Omega = \mathbb{N}$, $\Sigma = \mathcal{P}(\mathbb{N})$ et
$\mu$ la mesure de comptage sur $\mathbb{N}$ d\'efinie par~:
$$
\mu(E) = \sharp E = \sum_{k \in E} 1,
$$
o\`u $E\in \Sigma$. Soit $f~:\mathbb{N} \rightarrow \mathbb{R}$
une fonction positive ou nulle. Pour tout bor\'elien $E$,
$f^{-1}(E)$ appartient \`a $\mathcal{P}(\mathbb{N})$, donc $f$ est
($\Sigma$-$\mathcal{B}(\mathbb{R})$)-mesurable. Par d\'efinition
de l'int\'egrale,
$$
\int_{\Omega} f d\mu = \int_{0}^{\infty}
\mu\left(S_{f}(t)\right)\,dt,
$$
o\`u $S_{f}(t) = \{n\in\Sigma, \,f(n)>t\}$. Pour tout
$y\in[0,+\infty[$, posons $A_{y} := \{n\in\mathbb{N},\, f(n) =
y\}$. Alors $$S_{f}(t) = \cup_{y
> t}A_{y}$$ o\`u l'union est disjointe et o\`u $A_y$ est vide sauf
pour un ensemble d\'enombrable $\{y_{i}\}_{i\in\mathbb{N}}$ de
valeurs de $y$. Par $\sigma$-additivit\'e de la mesure $\mu$,
$$
\mu\left(S_{f}(t)\right) = \mu\left(\cup_{y_i > t} A_{y_{i}}
\right) = \sum_{y_i
> t} \mu\left(A_{y_{i}}\right) = \sum_{y_i > t}\mu\left(\{f = y_i
\}\right).
$$
Ainsi~: 
\begin{eqnarray*} \int_{\Omega} f d\mu & =&
\int_{0}^{\infty}\sum_{y_i
> t} \mu\left(\{f = y_i \}\right)\,dt = \sum_{i=0}^{\infty}
\int_{0\leq t<y_{i}}\mu\left(\{f = y_i \}\right)\,dt \\&=&
\sum_{i=0}^{\infty} y_i \cdot \mu\left(\{f = y_i \}\right) =
\sum_{i=0}^{\infty} y_{i}\cdot\sharp\{n\in\mathbb{N}, f(n) =
y_{i}\} = \sum_{n=0}^{\infty} f(n).
\end{eqnarray*}
\fincorrection
\correction{005936}
Soit $\varphi$ une fonction simple positive~:
$$
\varphi = \sum_{j \in J} c_{j} \mathbf{1}_{E_{j}},
$$
o\`u $J$ est un ensemble fini, les ensembles $E_{j}$ sont
mesurables et o\`u, pour $i\neq j$, $c_{i} \neq c_{j}$ et
$E_{i}\cap E_{j} = \emptyset$.
\begin{enumerate}
\item On a
$$
\int_{\Omega} \varphi \,d\mu = \int_{0}^{\infty}
\mu\left(S_{\varphi}(t) \right)\,dt,
$$
o\`u $S_{\varphi}(t) = \{x\in \Omega, \varphi(x)>t\} =
\cup_{c_{j}>t} E_{j}$ et o\`u $\mu\left(S_{\varphi}(t) \right) =
\sum_{c_{j}>t} \mu\left(E_{j}\right)$. Ainsi
$$
\int_{\Omega} \varphi \,d\mu = \int_{0}^{\infty}\sum_{c_{j}>t}
\mu\left(E_{j}\right)\,dt = \sum_{j\in
J}\int_0^{c_{j}}\mu\left(E_{j}\right)\,dt =  \sum_{j \in J}\,
c_{j} \mu(E_{j}).
$$

\item Pour tout $n\in\mathbb{N}$, posons
\begin{eqnarray*}
E_{k,n} &:=& \{x\in\Omega,~ k\, 2^{-n} \leq f(x) <(k+1) 2^{-n}\}
 \quad \text{pour}\quad k = 0, \dots, n \,2^{n} - 1,\\
E_{n, n} &:=& \{x\in\Omega,~ f(x)\geq n\} \quad \text{pour}\quad k
= n\, 2^{n}.
\end{eqnarray*}
Puisque $f$ est mesurable, les ensembles $E_{k, n}$ appartiennent
\`a $\Sigma$. Pour tout $n\in\mathbb{N}$ fix\'e, les ensembles
$E_{k,n}$, $0\leq k\leq n 2^n - 1$ sont deux \`a deux disjoints et
$\cup_{k} E_{k,n} = \Omega$. Posons
$$
\varphi_{n} = \sum_{k=0}^{n 2^{-n}-1} k\, 2^{-n} \mathbf{1}_{E_{k,n}}.
$$
Alors $\varphi_{n}$ est une fonction simple positive v\'erifiant
$\varphi_{n} \leq f$. En outre $0 \leq \varphi_{n}(x) \leq
\varphi_{n+1}(x)$ pour tout $x\in\Omega$ et pour tout
$n\in\mathbb{N}$. De plus,  $\lim_{n\rightarrow +\infty}
\varphi_{n}(x) = f(x)$ pour tout $x\in\Omega$.
\end{enumerate}
\fincorrection
\correction{005937}
Soit $(\Omega, \Sigma, \mu)$ un espace mesur\'e et
$f\in\mathcal{M}^+(\Omega, \Sigma)$. Pour tout $E\in\Sigma$, on
pose~:
$$
\lambda(E) = \int_{E} f\,d\mu = \int_{\Omega} \mathbf{1}_{A}\cdot f
\,d\mu.
$$
Montrons que $\lambda$ d\'efinit une mesure sur $(\Omega,
\Sigma)$.\\

\emph{$1^{\text{\`ere}}$ m\'ethode~:} On montre d'abord que
l'affirmation est vraie pour les fonctions simples.  D'apr\`es
l'exercice~\ref{ex:barb20}, toute fonction $f\in\mathcal{M}^{+}(\Omega, \Sigma)$
s'\'ecrit $f = \sup_{n\in\mathbb{N}} \varphi_{n}$, o\`u les
$\varphi_{n}$ sont des fonctions simples. Puisque le supremum
d'une famille quelconque de mesure est une mesure, on conclut que $\lambda$ est une mesure. \\

\emph{$2^{\text{de}}$ m\'ethode~:} On a clairement
$\lambda(\emptyset) = 0$. Il suffit donc de v\'erifier la
$\sigma$-additivit\'e de $\lambda$. Soit
$\{E_{i}\}_{i\in\mathbb{N}}\subset\Sigma$ une suite d'\'el\'ements
deux \`a deux disjoints. On a
\begin{eqnarray*}
\lambda\left(\bigcup_{i=1}^{\infty} E_{i}\right) & = &
\int_{\bigcup_{i=1}^{\infty} E_{i}} f\,d\mu = \int_{\Omega}
\mathbf{1}_{\bigcup_{i=1}^{\infty} E_{i}} f\,d\mu\\
& = &\int_{\Omega} \left(\sum_{i=1}^{\infty}\mathbf{1}_{E_{i}} \right)
f\,d\mu = \int_{\Omega} \sum_{i=1}^{\infty} \left( \mathbf{1}_{E_{i}}
f\right)
\,d\mu\\
& = & \sum_{i=1}^{\infty} \int_{\Omega} \left( \mathbf{1}_{E_{i}} f\right)
d\mu = \sum_{i=1}^{\infty} \int_{E_{i}} f\,d\mu \\ & = &
\sum_{i=1}^{\infty} \lambda(E_{i}).
\end{eqnarray*}
\fincorrection
\correction{005938}
Soit $f~:\mathbb{R}^n \rightarrow\mathbb{R}^+$ la fonction
d\'efinie par $$f(x) = |x|^{-p} \mathbf{1}_{\{|x|<1\}}(x).$$
\begin{enumerate}
\item [(i)] On a~:
\begin{eqnarray*} \int_{\mathbb{R}^{n}} f(x)\,dx & = &
\int_{\mathbb{R}^n}|x|^{-p} \mathbf{1}_{\{|x|<1\}}(x) \,dx = \int_{|x|<1}
|x|^{-p}\,dx = \int_{r=0}^{1}\int_{\sigma \in \mathcal{S}_{n-1}}
r^{n-p-1} dr d\sigma\\
& = & \frac{2\pi^{\frac{n}{2}}}{\Gamma\left(\frac{n}{2}\right)}
\int_{0}^{1} r^{n-p-1}\,dr.
\end{eqnarray*}
Pour $n \leq p$, il vient
\begin{eqnarray*}
 \int_{\mathbb{R}^{n}} f(x)\,dx &=& +\infty.
\end{eqnarray*}
Pour $p < n$, il vient
\begin{eqnarray*}
\int_{\mathbb{R}^{n}} f(x)\,dx &=&
\frac{2\pi^{\frac{n}{2}}}{\Gamma\left(\frac{n}{2}\right)}
\left[\frac{r^{n-p}}{(n-p)}\right]_{0}^{1} =
\frac{2\pi^{\frac{n}{2}}}{(n - p)\Gamma\left(\frac{n}{2}\right)}
\end{eqnarray*}




 \item[(ii)] Pour $a\in[0,+\infty[$,
\begin{eqnarray*}
S_{f}(a) = \{ x\in\mathbb{R}^{n}, |x|^{-p} \mathbf{1}_{|x|<1} > a\} =
\{x\in\mathbb{R}^{n}, |x|^{-p} > a\}\cap \mathcal{B}(0, 1),
\end{eqnarray*}
o\`u $\mathcal{B}(0, 1)$ est la boule de centre 0 et de rayon 1.
Ainsi
\begin{eqnarray*}
S_{f}(a) = \{ x\in\mathbb{R}^{n}, a^{-\frac{1}{p}} > |x|\}\cap
\mathcal{B}(0, 1).
\end{eqnarray*}
On en d\'eduit que $S_{f}(a) = \mathcal{B}(0, 1)$ si
$a^{-\frac{1}{p}}>1$, i.e. si $a<1$ et que $S_{f}(a)$ est \'egale
\`a la boule $\mathcal{B}(0, a^{-\frac{1}{p}})$ de centre 0 et de
rayon $a^{-\frac{1}{p}}$ lorsque $a\geq 1$. Il vient alors~:
\begin{eqnarray*}
\int_{\mathbb{R}^{n}} f(x)\,dx &=& \int_{0}^{+\infty}
\mu\left(S_{f}(a)\right)\,da = \int_{0}^{1}
\mu\left(\mathcal{B}(0,1) \right)\,da + \int_{1}^{+\infty}
\mu\left(\mathcal{B}(0, a^{-\frac{1}{p}})\right)\,da\\
& = & \frac{\pi^{\frac{n}{2}}}{\Gamma\left(\frac{n}{2}+1\right)} +
\frac{\pi^{\frac{n}{2}}}{\Gamma\left(\frac{n}{2}+1\right)}\int_{1}^{+\infty}
a^{-\frac{n}{p}}\,da.
\end{eqnarray*}
Si $p\geq n$, on obtient $\int_{\mathbb{R}^{n}} f(x)\,dx =+\infty$
et pour $p<n$, on a~:
\begin{eqnarray*}
\int_{\mathbb{R}^{n}} f(x)\,dx & =
&\frac{\pi^{\frac{n}{2}}}{\Gamma\left(\frac{n}{2}+1\right)} +
\frac{\pi^{\frac{n}{2}}}{\Gamma\left(\frac{n}{2}+1\right)}\left[\frac{a^{-\frac{n}{p}+1}}{-\frac{n}{p}+1}\right]_{1}^{+\infty}\\
&=&\frac{\pi^{\frac{n}{2}}}{\Gamma\left(\frac{n}{2}+1\right)}\left(1
+ \frac{p}{n-p} \right) = \frac{2\pi^{\frac{n}{2}}}{(n -
p)\Gamma\left(\frac{n}{2}\right)}.
\end{eqnarray*}
\end{enumerate}
\fincorrection
\correction{005926}
cf Proposition 5.1. du polycopi\'e de Marc Troyanov.
\fincorrection
\correction{005927}
D\'ecoule directement des d\'efinitions.
\fincorrection
\correction{005939}
\begin{enumerate}
\item  Soit $\{g_{n}\}_{n\in\mathbb{N}}$ une suite dans
$\mathcal{M}^{+}(\Omega, \Sigma)$. Alors $f_{k} = \sum_{n=1}^{k}
g_{n}$ est une suite croissante de $\mathcal{M}^{+}(\Omega,
\Sigma)$. D'apr\`es le th\'eor\`eme de convergence monotone
$$
\int_{\Omega} \left(\sum_{n=1}^{+\infty} g_{n}\right)\,d\mu =
\sum_{n=1}^{+\infty}\int_{\Omega} g_{n} \,d\mu.
$$

\item Posons $g_{n}(x) = x^{s-1} e^{-nx} \mathbf{1}_{[0, +\infty)}$. Les
$g_{n}$ appartiennent \`a $\mathcal{M}^{+}(\Omega, \Sigma)$ pour
tout $n\in\mathbb{N}$. D'apr\`es la question pr\'ec\'edente,
$$
\int_{\Omega} \left(\sum_{n=1}^{+\infty} g_{n}\right)\,d\mu =
\sum_{n=1}^{+\infty}\int_{\Omega} g_{n} \,d\mu.
$$
Or d'une part,
$$
\int_{\Omega}g_{n} \,d\mu = \int_{\Omega} x^{s-1} e^{-nx}
\mathbf{1}_{[0,+\infty)}\, dx = \frac{1}{n^s}\int_{0}^{+\infty} y^{s-1}
e^{-y}\,dy = \frac{1}{n^s}\Gamma(s),
$$
donc
$$
\sum_{n=1}^{+\infty}\int_{\Omega} g_{n} \,d\mu =
\sum_{n=1}^{\infty}\frac{1}{n^s}\Gamma(s) = \zeta(s) \Gamma(s).
$$
D'autre part,
$$
\int_{\Omega} \left(\sum_{n=1}^{+\infty} g_{n}\right)\,d\mu =
\int_{0}^{+\infty} x^{s-1} \sum_{n=1}^{+\infty} e^{-nx}\,dx =
\int_{0}^{+\infty} \frac{x^{s-1}}{e^x -1}\,dx,
$$
d'o\`u l'\'egalit\'e cherch\'ee.
\end{enumerate}
\fincorrection
\correction{005940}
Oui, le th\'eor\`eme de convergence monotone ne dit pas que
l'int\'egrale de $f$ est finie. On a bien
$$
+\infty = \int_{\Omega} f\,d\mu = \lim_{n\rightarrow
+\infty}\int_{\Omega} f_{n}\,d\mu = \lim_{n\rightarrow +\infty} n.
$$
\fincorrection
\correction{005941}
Non, le th\'eor\`eme de convergence monotone ne s'applique pas \`a
une suite d\'ecroissante de fonctions positives.
\fincorrection
\correction{005942}
Non, la suite de fonctions n'est pas m\^eme monotone.
\fincorrection
\correction{005943}
En effet, pour tout $\varepsilon>0$, il existe $N_{\varepsilon} =
\left[\frac{1}{\varepsilon} \right] + 1$ tel que $\forall n \geq
N_{\varepsilon}$, $$ \sup_{x\in\mathbb{R}} |f_{n}(x) - f(x)| <
\varepsilon,
$$
i.e. $f_{n}$ converge uniform\'ement vers $f$ sur $\mathbb{R}$. On
a~:
$$
\lim \inf_{n\rightarrow +\infty} \int_{\Omega} f_{n}\,d\mu = \lim
\inf_{n\rightarrow +\infty} - \int_{0}^{n} \frac{1}{n}\,d\mu = -1.
$$
D'autre part $\int_{\Omega} f\,d\mu = 0.$ Le lemme de Fatou ne
s'applique pas car les fonctions $f_{n}$ ne sont pas \`a valeurs
dans $[0, +\infty]$.
\fincorrection
\correction{005950}
\begin{enumerate}
\item Montrons que pour tout $x\in \mathbb{R}_+$,
$(1+\frac{x}{n})^n$ est une suite croissante et que
$\lim\limits_{n\rightarrow\infty}(1+\frac{x}{n})^n=e^x$.
 Pour $n\in \mathbb{N}$ on a
$$\left(1+\frac{x}{n}\right)^n=\sum\limits_{k=0}^n
{\binom{n}{k}}\left(\frac{x}{n}\right)^k=\sum\limits_{k=0}^n
a_{n,k}\,\frac{x^k}{k!},$$ o\`{u} \; $\displaystyle
a_{n,k}=\frac{n!}{(n-k)!n^k}=\frac{n(n-1)\cdots(n-k+1)}{n^k}.$

\bigskip

Les assertions suivantes sont vraies: 
\begin{enumerate}
\item[{i)}] $a_{n+1,k}\geq a_{n,k}$.  En effet,
$\frac{n+1-l}{n+1}\geq \frac{n-l}{n}$ \, pour \, $l\in \mathbb{N}$
\,  car \, $n^2+n-l\cdot n\geq n^2+n-l\cdot n-l$,

\item[{ii)}] $a_{n,k}<1$ \; (\'{e}vident);

\item[{iii)}] Pour tout $k\in \mathbb{N}$, \;
$\lim\limits_{n\rightarrow\infty}a_{n,k}=1$.

\end{enumerate}

Comme $a_{n+1,n+1}>0$,
$\displaystyle\left(1+\frac{x}{n+1}\right)^{n+1}=\sum\limits_{k=0}^{n+1}
a_{n+1,k}\,\frac{x^k}{k!}>\sum\limits_{k=0}^{n}
a_{n+1,k}\,\frac{x^k}{k!} .$ Il s'ensuit donc de (i) que la suite
$\displaystyle\left\{\left(1+\frac{x}{n}\right)^n\right\}$  est
croissante. Les assertions (ii) et (iii) impliquent que
$$\left(1+\frac{x}{n}\right)^n=\sum\limits_{k=0}^n
a_{n,k}\,\frac{x^k}{k!}\leq \sum\limits_{k=0}^n \frac{x^k}{k!}\leq
\sum\limits_{k=0}^\infty \frac{x^k}{k!}=e^x$$ et que, pour tout
$m\in\mathbb{N}$,\;
$\displaystyle\lim\limits_{n\rightarrow\infty}\left(1+\frac{x}{n}\right)^n\geq
\sum\limits_{k=0}^m \frac{x^k}{k!}$. \; \; Ainsi, \;
$\displaystyle\lim\limits_{n\rightarrow\infty}\left(1+\frac{x}{n}\right)^n
= e^x.$



\item
Par le th\'{e}or\`{e}me de convergence monotone, on a
pour $b>1$,
\begin{eqnarray*}
\lim\limits_{n\rightarrow\infty}\int_{\mathbb{R}_+}\left(1+\frac{x}{n}\right)^n
e^{-bx} d\lambda(x) &=
\int_{\mathbb{R}_+}\lim\limits_{n\rightarrow\infty}\left(1+\frac{x}{n}\right)^n
e^{-bx} d\lambda(x)\\ &=\int_0^\infty e^{(1-b)x}
d\lambda(x)=\frac{1}{b-1}.
\end{eqnarray*}
\end{enumerate}
\fincorrection
\correction{005944}
On a
$$\mu\left(f =
+\infty\right) = \mu\left(\cap_{n\in\mathbb{N}}\{f \geq
n\}\right).
$$
Puisque les ensembles $A_{n} := \{f \geq n\}$ v\'erifient $A_1
\supset A_{2} \supset A_{3}\dots$ et $\mu(A_i) < +\infty$ ($i=1,2,\ldots$), par continuit\'e de la mesure,
on a~:
$$\mu\left(f =
+\infty\right) = \lim_{n\rightarrow+\infty} \mu\left(f\geq n
\right).
$$
Or, comme $f$ est \`a valeurs positives, les fonctions $f_{n}$
d\'efinies par $f_{n} = n\mathbf{1}_{\{f\geq n\}}$ v\'erifient $f_{n}\leq
f$. Ainsi
$$
\int_{\Omega} f_{n}\,d\mu  = \int_{\Omega} n \mathbf{1}_{\{f\geq
n\}}\,d\mu = n \mu\left(f\geq n \right) ~\leq ~\int_{\Omega}
f\,d\mu <+\infty.
$$
On en d\'eduit que $$\mu\left(f\geq n \right) ~\leq~
\frac{1}{n}\int_{\Omega} f\,d\mu \rightarrow 0,$$ donc
$$\mu\left(f = +\infty\right) = 0.$$
\fincorrection
\correction{005945}
Puisque $\mu(\Omega) < +\infty$, la fonction constante \'egale \`a
$C$ est int\'egrable, d'int\'egrale $C \, \mu(\Omega)$. Une
application directe du th\'eor\`eme de convergence domin\'ee donne
$$
\lim_{n\rightarrow+\infty}\int_{\Omega} f_{n}\,d\mu =
\int_{\Omega} f\,d\mu.
$$
\fincorrection
\correction{005946}
 Soit $f\in
\mathcal{L}^1(\mathbb{R}).$ Comme $\cos(\pi x)<1$ si $x\notin
\mathbb{Z},$ $\cos^n(\pi x)\rightarrow 0$ lorsque
$n\rightarrow\infty$ presque partout (pour tout $x\in
\mathbb{R}\setminus\mathbb{Z}$). Notons $f_n(x)=f(x) \cos^n(\pi
x).$ Alors, pour tout $n\in \mathbb{N}$ on a $|f_n(x)|\leq |f(x)|$
et comme $|f|\in \mathcal{L}^1(\mathbb{R}),$ par le
th\'{e}or\`{e}me de convergence domin\'{e}e de Lebesgue,
$$\lim\limits_{n\rightarrow\infty}\int_\mathbb{R} f(x)\cos^n(\pi x)d\lambda(x)=
\lim\limits_{n\rightarrow\infty}\int_\mathbb{R}
f_n(x)d\lambda(x)=\int_\mathbb{R}
\lim\limits_{n\rightarrow\infty}f_n(x)d\lambda(x)=0.$$
\fincorrection
\correction{005947}
\begin{enumerate}
\item Par d\'efinition, $ f\in \mathcal{L}^{1}(\Omega,
\Sigma, \mu)$ si et seulement si $f_{+}$ et $f_{-}$ sont
int\'egrables. On note que $|f| = f_{+} + f_{-}$. Donc $f\in
\mathcal{L}^{1}(\Omega, \Sigma, \mu)\Rightarrow |f| \in
\mathcal{L}^{1}(\Omega, \Sigma, \mu)$. R\'eciproquement, on a $0
\leq f_{\pm} \leq |f|$, donc $ |f| \in \mathcal{L}^{1}(\Omega,
\Sigma, \mu)\Rightarrow f\in \mathcal{L}^{1}(\Omega, \Sigma,
\mu)$. D'autre part~:
\begin{equation*} \left|\int_{\Omega} f\,d\mu
\right| ~=~ \left|\int_{\Omega} f_{+}\,d\mu - \int_{\Omega}
f_{-}\,d\mu\right|~\leq~ \int_{\Omega} f_{+}\,d\mu + \int_{\Omega}
f_{-}\,d\mu ~=~\int_{\Omega} |f|\,d\mu.
\end{equation*}

\item Par monotonie de l'int\'egrale, on a
$$
\int_{\Omega} |f|\,d\mu ~\leq~\int_{\Omega} |g|\,d\mu < +\infty.
$$
D'apr\`es la question $(a)$, il en d\'ecoule que $f$ est
int\'egrable.

\item D\'efinissons $z = \int_{\Omega} f\,d\mu$. Comme $z$
est un nombre complexe, il s'\'ecrit $z = |z|e^{i\theta}$. Soit
$u$ la partie r\'eelle de $e^{-i\theta} f$. On a $u \leq
|e^{-i\theta}f| = |f|.$ Donc
$$
\left|\int_{\Omega} f\,d\mu \right|~=~ e^{-i\theta}\int_{\Omega}
f\,d\mu ~=~\int_{\Omega}e^{-i\theta} f\,d\mu ~=~ \int_{\Omega}
u\,d\mu ~\leq~ \int_{\Omega} |f|\,d\mu,
$$
o\`u la troisi\`eme \'egalit\'e d\'ecoule du fait que le nombre
$\int_{\Omega}e^{-i\theta} f\,d\mu$ est r\'eel donc est
l'int\'egrale de la partie r\'eelle de $e^{-i\theta} f$
c'est-\`a-dire de $u$.
\end{enumerate}
\fincorrection
\correction{005948}
On cherche une sous-suite $\{ f_{n_{k}}\}_{n\in\mathbb{N}}$ de $\{
f_{n}\}_{n\in\mathbb{N}}$ telle que pour $\mu$-presque tout
$x\in\Omega$, \'etant donn\'e un $\varepsilon>0$, il existe un
$k\in\mathbb{N}$ (d\'ependant \`a priori de $x$) v\'erifiant
$j\geq k \Rightarrow |f_{n_{j}}(x) - f(x)| < \varepsilon$. Il
suffit de montrer que pour $\mu$-presque tout $x$, il existe un
$k\in\mathbb{N}$ tel que $j\geq k \Rightarrow |f_{n_{j}}(x) -
f(x)| < \frac{1}{2^{j}} \leq \frac{1}{2^k}.$ Cela revient \`a
montrer que le compl\'ementaire de l'ensemble
$$
A := \bigcup_{k=1}^{\infty} \bigcap_{j\geq k} \left\{|f_{n_{j}} -
f|< \frac{1}{2^{j}}\right\}
$$
est de mesure nulle. Or
$$
A^{c} = \bigcap_{k=1}^{\infty} \bigcup_{j\geq k}\left\{|f_{n_{j}}
- f|\geq \frac{1}{2^{j}}\right\}.
$$
Posons $B_{k} := \bigcup_{j\geq k}\left\{|f_{n_{j}} - f|\geq
\frac{1}{2^{j}}\right\}$. On a $B_{1} \supset B_{2} \supset B_3
\dots$ avec $B_1$ de mesure fini ; donc par continuit\'e de la mesure, il vient~:
$$
\mu(A^{c}) = \lim_{k\rightarrow +\infty} \mu(B_{k}).
$$
Par $\sigma$-additivit\'e, on a~:
$$
\mu(B_{k}) ~\leq~\sum_{j\geq k} \mu\left( \left\{|f_{n_{j}} -
f|\geq \frac{1}{2^{j}}\right\}\right).
$$
On d\'efinit alors la sous-suite $\{ f_{n_{k}}\}_{n\in\mathbb{N}}$
de la mani\`ere suivante. Puisque $f_{n}$ converge vers $f$ en
mesure, il existe un indice $n_1$ tel que pour $n\geq n_1$,
$$
\mu\left( \left\{|f_{n} - f|\geq \frac{1}{2}\right\}\right) \leq
\frac{1}{2}.
$$
Il existe un indice $n_{2}> n_1$ tel que pour $n\geq n_2$,
$$
\mu\left( \left\{|f_{n} - f|\geq \frac{1}{2^2}\right\}\right) \leq
\frac{1}{2^2},
$$
et ainsi de suite: pour tout $k\in\mathbb{N}$, il existe un
$n_{k}> n_{k-1}$, tel que pour $n\geq n_{k}$
$$
\mu\left( \left\{|f_{n} - f|\geq \frac{1}{2^k}\right\}\right) \leq
\frac{1}{2^k}.
$$
Pour cette sous-suite on a alors~:
$$
\mu(B_{k}) ~\leq~\sum_{j\geq k} \mu\left( \left\{|f_{n_{j}} -
f|\geq \frac{1}{2^{j}}\right\}\right) ~\leq ~\sum_{j\geq k}
\frac{1}{2^j} = \frac{1}{2^{k-1}}.
$$
On a bien
$$
\mu(A^{c}) = \lim_{k\rightarrow +\infty} \mu(B_{k}) = 0.
$$
\fincorrection
\correction{005949}
La fonction de Dirichlet restreint
\`a l'intervalle $[a,b]$,
$f(x)=\mathbf{1}_\mathbb{Q}\left|_{[a,b]}(x)\right.,$ est int\'{e}grable
au sens de Lebesgue et son int\'egrale par rapport \`a la mesure
de Lebesgue vaut $0$. Mais elle n'est pas int\'{e}grable au sens
de Riemann: $\underline{S}(f,\tau)=0$ et
$\overline{S}(f,\tau)=b-a\;$ pour toute subdivision $\tau$ de
l'intervalle $[a,b].$
\fincorrection
\correction{005951}
\begin{enumerate}
 \item Pour tout $x\in \mathbb{R}_+$
et $n\in \mathbb{N}$ on a
$$\left(1-\frac{x}{n}\right)^n \leq e^{-x}.$$ En effet, comme \, $\ln
y\leq y-1$ \, pour $y>0$, on a \, $\ln y^{-\frac{1}{n}}\leq
y^{-\frac{1}{n}}-1$,\, c'est-\`a-dire $\left(1 - \frac{\ln y}{n}
\right)^{n} \leq y^{-1}$. Ainsi, en posant $x = \ln y$, il vient
$\left(1-\frac{x}{n}\right)^n \leq e^{-x}$. De plus,
$$
\lim_{n\rightarrow+\infty} \left(1-\frac{x}{n} \right)^{n} =
\lim_{n\rightarrow+\infty}e^{n\ln\left(1-\frac{x}{n} \right)} =
\lim_{n\rightarrow+\infty}e^{n\left(-\frac{x}{n} +
\frac{x}{n}\varepsilon\left(\frac{x}{n} \right) \right)},
$$
o\`u $\lim_{u\rightarrow 0}\varepsilon(u) = 0$. Ainsi
$\lim_{n\rightarrow+\infty} \left(1-\frac{x}{n} \right)^{n} =
e^{-x}$.


Posons $f_n(x)=\left(1-\frac{x}{n}\right)^n x^m\mathbf{1}_{[0,n]}$. Alors
en utilisant le th\'{e}or\`{e}me de convergence domin\'{e}e et
sachant que \, $ \Gamma(m+1)=\int\limits_0^\infty e^{-x}x^m dx
=m!$, on obtient le r\'{e}sultat.
 
\item Soit $f_n(x)=\left(1+\frac{x}{n}\right)^n e^{-2x}\mathbf{1}_{[0,n]}$.
Comme la suite $\{f_n(x)\}$ est croissante et \,
$\lim\limits_{n\rightarrow \infty}f_n(x)=e^{-x}$, on obtient le
r\'{e}sultat en appliquant le th\'{e}or\`{e}me de convergence
monotone.
\end{enumerate}
\fincorrection
\correction{005952}
Cf le th\'eor\`eme 24.2 dans le polycopi\'e de Marc Troyanov.
\fincorrection
\correction{005953}
\begin{enumerate}
 \item  Notons $g(x,y)=e^{-ixy} f(x)$.
Alors,
\begin{enumerate}
\item[{i.)}] pour tout $y\in \mathbb{R}$, la fonction $x\mapsto
g(x,y)$ est mesurable; \item[{ii.)}] pour presque tout $x\in
\mathbb{R}$ (pour tout les $x\in \mathbb{R}$ tels que $f(x)$ est
finie) la fonction $y\mapsto g(x,y)$ est continue pour tout $y\in
\mathbb{R}$; \item[{iii.)}] $|g(x,y)|=|e^{-ixy} f(x)|\leq |f(x)|$
et $|f|\in \mathcal{L}^1(\mathbb{R},\mathbb{R}_+).$
\end{enumerate}
On doit montrer que pour tout suite $\{y_{n}\}_{n\in\mathbb{N}}$
convergeant vers $y$, on a $\lim_{n\rightarrow +\infty}
\hat{f}(y_{n}) = \hat{f}(y)$. Posons $g_{n}(x) = g(x, y_{n})$.
D'apr\`es le th\'eor\`eme de convergence domin\'ee,
$$
\lim_{n\rightarrow +\infty} \hat{f}(y_{n}) ~:=~ \lim_{n\rightarrow
+\infty} \int_{\mathbb{R}} g_{n}(x)\,dx = \int_{\mathbb{R}}g(x,y)
dx =: \hat{f}(y).$$
Ainsi $\hat{f}$ est continue.

\item Pour tout $y\in \mathbb{R}$, $\;|\hat{f}(y)|\leq
\displaystyle\int_\mathbb{R}|e^{-ixy} f(x)|dx\leq
\displaystyle\int_\mathbb{R}| f(x)|dx=\|f\|_{L_1}$ et donc
$$\sup|\hat{f}|\leq\|f\|_{L_1}.$$

\item Soit $g(x,y)=e^{-ixy} f(x)$. Alors, on  a
\begin{enumerate}
\item[{i.)}] pour tout $y\in \mathbb{R}$, la fonction $x\mapsto
g(x,y)$ est int\'{e}grable;

\item[{ii.)}] pour presque tout $x\in \mathbb{R}$ la fonction
$y\mapsto g(x,y)$ est d\'{e}rivable pour tout $y\in \mathbb{R}$;

\item[{iii.)}] $|\frac{\partial g(x,y)}{\partial y}|=|-ix e^{-ixy}
f(x)|\leq |x f(x)|$ avec $x\mapsto xf(x)$ int\'{e}grable.
\end{enumerate}
Ainsi, d'apr\`{e}s l'exercice \ref{derivation} (le th\'eor\`eme
de dérivation sous le signe $\int$), on a
$$\frac{d}{dy}\hat{f}=\int_\mathbb{R}e^{-ixy}(-ix f(x))dx=\widehat{-iyf(y)}.$$
\end{enumerate}
\fincorrection
\correction{005957}
On a
\begin{align*}\int_{-1}^1\left(\int_{-1}^1
\frac{x^2-y^2}{(x^2+y^2)^2}dx\right) dy=&\int_{-1}^1\left(\left.-
\frac{x}{(x^2+y^2)}\right|_{-1}^1\right) dy \\=&\left.-\int_{-1}^1
\frac{2}{(1+y^2)} dy =-2 \arctan y \right|_{-1}^1=-\pi.
\end{align*}
\begin{align*}\int_{-1}^1\left(\int_{-1}^1 \frac{x^2-y^2}{(x^2+y^2)^2}dy\right) dx=&\int_{-1}^1\left(\left.
\frac{y}{(x^2+y^2)}\right|_{-1}^1\right) dx \\=&\left.\int_{-1}^1
\frac{2}{(x^2+1)} dx =2 \arctan x \right|_{-1}^1=\pi.
\end{align*}
Il n'y a pas de contradiction avec le th\'{e}or\`{e}me de Fubini
car la fonction $f$ n'appartient pas \`a
$\mathcal{L}^1([-1,1]\times[-1,1])$. En effet, soit
$S_\varepsilon=\{(x,y)\in\mathbb{R}^2| \varepsilon\leq x^2+y^2\leq
1\}$. On a
$$\int_{[-1,1]\times[-1,1]}|f| d\mu 
\geq \int_{S_\varepsilon} |f| d\mu =\int_{\theta=0}^{2\pi}\int_{r=\varepsilon}^1\frac{|\cos 2\theta|}{r}dr
d\theta
= 4\int_{\theta=0}^{\frac{\pi}{2}}\int_{r=\varepsilon}^1 \frac{|\cos 2\theta|}{r} dr
d\theta = -4 \log \varepsilon\rightarrow \infty$$ lorsque
$\varepsilon\rightarrow 0,$ et donc $f\notin
\mathcal{L}^1([-1,1]\times[-1,1]).$
\fincorrection
\correction{005958}
Le th\'eor\`eme de Tonelli donne~:
$$
\int_{[0, 1]\times (0, +\infty)} |e^{-y} \sin 2xy|\,dx dy \leq
\int_{0}^{+\infty} e^{-y}\,dy = 1 < +\infty,
$$
ce qui prouve que la fonction $(x, y)\mapsto e^{-y} \sin 2xy$ est
int\'egrable pour la mesure de Lebesgue sur $[0,1]\times
(0,+\infty)$.
\\
Le th\'eor\`eme de Fubini donne alors la valeur $I$ de
l'int\'egrale de cette fonction~:
\begin{eqnarray*}
I  = \int_{0}^{1} \,dx \int_{0}^{+\infty} e^{-y} \sin 2xy\,dy
\begin{array}{c} \text{(IPP)}\\= \end{array}
\int_{0}^{1} (2x) (1+4x^{2})^{-1}\,dx = \frac{\log 5}{4}\\
I  = \int_{0}^{+\infty} e^{-y} \,dy \int_{0}^{1} \sin 2xy \,dx =
\int_{0}^{+\infty} e^{-y}~ \frac{\sin^2 y}{y}\,dy.
\end{eqnarray*}

\fincorrection
\correction{005959}
Soient $f\in L^1(\mathbb{R}^n)$ et $g\in L^p(\mathbb{R}^n)$ avec
$1\leq p\leq +\infty$, o\`u $\mathbb{R}^n$ est muni de la mesure
de Lebesgue. L'identit\'e $f*g(x) = g*f(x)$ s'obtient par
changement de variable. En ce qui converne l'in\'egalit\'e $
\|f*g\|_{p} \leq \|f\|_{1}\,\|g\|_{p}, $ on distingue les cas en
fonction de la valeur de $p$.
\begin{enumerate}
\item Pour $p = + \infty$, c'est clair.

\item Supposons que $p = 1$ et posons $F(x, y) =
f(x-y)\,g(y)$. Pour presque tout $y\in\mathbb{R}^n$, on a~:
$$
\int_{\mathbb{R}^n} |F(x, y)|\,dx = |g(y)|
\int_{\mathbb{R}^n}|f(x- y)|\,dx = |g(y)| \cdot\|f\|_1,
$$
et
$$
\int_{\mathbb{R}^n} dy \int_{\mathbb{R}^n}|F(x, y)|\,dx =
\|f\|_1\,\|g\|_1.
$$
D'apr\`es le th\'eor\`eme de Tonelli, $F\in
L^{1}\left(\mathbb{R}^n\times\mathbb{R}^n\right)$. D'apr\`es le
th\'eor\`eme de Fubini, on a
$$
\int_{\mathbb{R}^n} |F(x, y)|\,dy <+\infty \quad\quad\text{pour
presque tout }~x\in\mathbb{R}^n,
$$
et
$$
\int_{\mathbb{R}^n} dx \int_{\mathbb{R}^n}|F(x, y)|\,dy
\leq\|f\|_1\,\|g\|_1.
$$
Ainsi, $$\|f*g\|_{1} = \int_{\mathbb{R}^n} \,dx |f*g(x)|
=\int_{\mathbb{R}^n} \,dx \left|\int_{\mathbb{R}^{n}} F(x,
y)\,dy\right| \leq\|f\|_1\,\|g\|_1.
$$


 \item Supposons que $1< p<+\infty$. Utilisons le cas
 pr\'ec\'edent, en faisant jouer ici à $g^p$ le rôle alors joué par $g$. Alors pour presque tout $x\in\mathbb{R}^n$ fix\'e, la
 fonction $y \mapsto |f(x-y)| |g(y)|^p$ est int\'egrable sur
 $\mathbb{R}^n$, i.e. la fonction $y \mapsto |f(x-y)|^{\frac{1}{p}}
 |g(y)|$ appartient \`a $L^p(\mathbb{R}^n)$. Soit $p'$ tel que
 $\frac{1}{p} + \frac{1}{p'} = 1$. La fonction $y\mapsto
 |f(x-y)|^{\frac{1}{p'}}$ appartient \`a $L^{p'}(\mathbb{R}^n)$
 car $f\in L^1(\mathbb{R}^n)$ et la mesure de Lebesgue est invariante par translation.
 D'apr\`es l'in\'egalit\'e de
 H\"older,
$$
|f(x-y)| \,|g(y)| = |f(x-y)|^{\frac{1}{p}}
\,|g(y)|\cdot|f(x-y)|^{\frac{1}{p'}}\in L^1(\mathbb{R}^n)
$$
et
$$
\int_{\mathbb{R}^n} |f(x-y)| \,|g(y)| \leq \left(
\int_{\mathbb{R}^n} |f(x-y)|
\,|g(y)|^p\right)^{\frac{1}{p}}\cdot\|f\|_1^{\frac{1}{p'}},
$$
ainsi
$$
|(f*g)(x)|^{p} \leq (|f|*|g|^p)(x)\cdot\|f\|_1^{\frac{p}{p'}}.
$$
D'apr\`es le cas pr\'ec\'edent, on voit que
$$
f*g\in L^p(\mathbb{R}^n) \quad \text{et}\quad \|f*g\|_p^{p} \leq
\|f\|_1\|g\|_{p}^{p}\cdot\|f\|_1^{\frac{p}{p'}},
$$
c'est-\`a-dire
$$
\|f*g\|_p \leq \|f\|_1\cdot\|g\|_p.
$$
\end{enumerate}

\fincorrection
\correction{005960}
Soient $a, b>0$, et $f$ et $g$ les fonctions d\'efinies sur
$\mathbb{R}^n$ par $f(x) = e^{-\frac{a |x|^2}{2}}$ et $g(x) =
e^{-\frac{b |x|^2}{2}}$. On a
\begin{eqnarray*}
f*g(x)  =  \int_{\mathbb{R}^n} f(x-y)\,g(y)\,dy =
\int_{\mathbb{R}^n} e^{-\left(\frac{a|x-y|^2 +
b|y|^2}{2}\right)}\, dy
\end{eqnarray*}
Or
\begin{eqnarray*}
a|x-y|^2 + b|y|^2 &=& \sum_{i=1}^n a x_i^2 + (a + b) y_i^2 -2a x_i
y_i\\
& = & \sum_{i=1}^n a x_i^2 + (a + b)\left(y_i - \frac{a}{a+b} x_i
\right)^2 - (a+b) \left(\frac{a x_i}{a+b} \right)^2\\
& = & \sum_{i=1}^n \left(a - \frac{a^2}{a+b}\right) x_i^2 + (a +
b)\left(y_i
- \frac{a}{a+b} x_i \right)^2\\
& = & \frac{ab}{a + b}|x|^2 +
(a+b)\left|y-\frac{a}{a+b}x\right|^2.
\end{eqnarray*}
Ainsi
\begin{eqnarray*}
f*g(x)  & = & e^{-\frac{ab}{a +
b}\frac{|x|^2}{2}}\int_{\mathbb{R}^n}
e^{-\frac{(a+b)}{2}\left|y-\frac{a}{a+b}x\right|^2}\,dy =
e^{-\frac{ab}{a + b}\frac{|x|^2}{2}}
 \int_{\mathbb{R}^n} e^{-\frac{(a + b)}{2} |z|^2}\,dz
 \end{eqnarray*}
car la mesure de Lebesgue est invariante par translation. En
utilisant $\int_{\mathbb{R}} e^{-t^2} \,dt = \sqrt{\pi}$, on
obtient alors~:
$$
f*g(x) = \left(\frac{2\pi}{a + b}\right)^{\frac{n}{2}}
e^{-\frac{ab}{a + b}\frac{|x|^2}{2}}.
$$
\fincorrection
\correction{005961}
\begin{enumerate}
\item Pour tout $t>0$, on pose~:
$$
f_{t}(x) = \left(4\pi t \right)^{-\frac{n}{2}}
e^{-\frac{|x|^2}{4t}}.
$$
\begin{enumerate}
\item On a
\begin{eqnarray*}
\int_{\mathbb{R}^n} f_t(x)\,dx & = & \left(4\pi t
\right)^{-\frac{n}{2}} \int_{\mathbb{R}^n}
e^{-\frac{|x|^2}{4t}}\,dx\\
& = & \left(4\pi t \right)^{-\frac{n}{2}} \prod_{i= 1}^{n}
\int_{\mathbb{R}} e^{-\frac{x_i^2}{4t}}\,dx_i.
\end{eqnarray*}
Sachant que $\int_{\mathbb{R}} e^{-t^2} \,dt = \sqrt{\pi}$, on en
d\'eduit que
$$
\int_{\mathbb{R}^n} f_t(x)\,dx  = 1.
$$

 \item
 Soit $\varepsilon>0$. Puisque $f_1$ est int\'egrable sur $\mathbb{R}^{n}$,
 il existe un $R>0$ tel que $$\int_{\mathcal{B}(0, R)^c} f_1(x)\,dx < \varepsilon.$$
 On remarque que $f_t(x) = t^{-\frac{n}{2}}\,f_1\left(\frac{x}{\sqrt{t}}
 \right)$. On a alors,
\begin{eqnarray*}
\int_{\mathcal{B}(0, \delta)^c} f_t(x)\,dx &=&
\int_{\mathcal{B}(0,
\delta)^c}t^{-\frac{n}{2}}\,f_1\left(\frac{x}{\sqrt{t}}
 \right) \,dx = t^{-\frac{n}{2}}\,\int_{\mathcal{B}\left(0,
\frac{\delta}{\sqrt{t}}\right)^c} f_1(z) \,t^{\frac{n}{2}}\,dz
\\&=& \int_{\mathcal{B}\left(0, \frac{\delta}{\sqrt{t}}\right)^c} f_1(z)\,dz
\leq \varepsilon,
\end{eqnarray*}
d\`es que $t < \frac{\delta^2}{R^2}$.
\end{enumerate}



\item Soit $g$ une fonction continue  born\'ee. Alors il existe
$M>0$ tel que $|g|< M$ et
$$
 \int_{\mathbb{R}^n} |f_t(x-y) \,g(y)|\,dy \leq M \int_{\mathbb{R}^n}
 f_t(x-y)\,dy = M<+\infty,
$$
ainsi $y \mapsto f_t(x-y) \,g(y)$ est int\'egrable et $f_t*g$ est
bien d\'efinie.  Puisque $\int_{\mathbb{R}^n} f_t(x) \,dx = 1$, on
a
\begin{eqnarray*}
\left|f_t*g(x) - g(x) \right| 
  &= \left| \int_{\mathbb{R}^n} f_t(y)g(x-y)\,dy - \int_{\mathbb{R}^n} f_t(y)g(x)\,dy \right| \\
  &\leq \int_{\mathbb{R}^n} f_t(y) \left|g(x-y) - g(x)\right|\,dy. \\
\end{eqnarray*}
Soit $\varepsilon>0$. Puisque $g$ est continue en
$x\in\mathbb{R}^n$, il existe $\delta>0$ tel que $|y|< \delta
\Rightarrow |g(x-y) - g(x)|<\varepsilon$. Alors
\begin{eqnarray*}
\left|f_t*g(x) - g(x) \right| &\leq& \int_{\mathcal{B}(0, \delta)}
f_t(y) \left|g(x-y) - g(x)\right|\,dy  \\ && + \int_{\mathcal{B}(0,
\delta)^c} f_t(y) \left|g(x-y) - g(x)\right|\,dy\\
&\leq & \varepsilon \int_{\mathcal{B}(0, \delta)} f_t(y)\,dy +
2M\int_{\mathcal{B}(0, \delta)^c} f_t(y)\,dy\\
&\leq& \varepsilon + 2M\int_{\mathcal{B}(0, \delta)^c} f_t(y)\,dy.
\end{eqnarray*}
D'apr\`es la question 1.(b), il existe $t_0>0$ tel que pour $t<
t_0$, $\int_{\mathcal{B}(0, \delta)^c} f_t(y)\,dy \leq
\frac{\varepsilon}{2M}$. Ainsi pour $t<t_0$,
$$
\left|f_t*g(x) - g(x) \right|< 2\varepsilon,
$$
i.e.
$$\lim_{t\rightarrow 0} f_t*g(x) = g(x).$$
\end{enumerate}
\fincorrection
\correction{005962}
Soient $f, g\in L^1(\mathbb{R}^n)$. On note $\hat{f}$ la
transform\'ee de Fourier d\'efinie par
$$
\hat{f}(y) = \int_{\mathbb{R}^n} f(x)\,e^{-2\pi i (y, x)} \,dx,
$$
o\`u $(\cdot, \cdot)$ d\'esigne le produit scalaire de
$\mathbb{R}^n.$
\begin{enumerate}
\item On a $\|\hat{g}\|_{\infty} \leq \|g\|_1$, ce qui
implique que $f\, \hat{g}$ est int\'egrable. De m\^eme
$\hat{f}\,g$ est int\'egrable. De plus $F(x, y) = f(x) g(y)
e^{-2\pi i (x, y)}$ appartient \`a
$L^1(\mathbb{R}^n\times\mathbb{R}^n)$. D'apr\`es le th\'eor\`eme
de Fubini, \begin{eqnarray*} \int_{\mathbb{R}^n}
f(x)\hat{g}(x)\,dx &=& \int_{\mathbb{R}^n}\,dx\, f(x)
\int_{\mathbb{R}^n} g(y) e^{-2\pi i (x, y)} \,dy \\ &=&
\int_{\mathbb{R}^n}\,dy \,g(y) \int_{\mathbb{R}^n} f(x)e^{-2\pi i
(x, y)} \,dx = \int_{\mathbb{R}^n} \hat{f}(y) g(y)\,dx.
\end{eqnarray*}


\item On a
\begin{eqnarray*}
\widehat{f*g}(x) &=& \int_{\mathbb{R}^n} f*g(y)\,e^{-2\pi i (x,
y)}\,dy = \int_{\mathbb{R}^n}\,dy \,e^{-2\pi i (x,
y)}\int_{\mathbb{R}^n}f(y-z)\,g(z) \,dz\\
&=& \int_{\mathbb{R}^n} \,dy \int_{\mathbb{R}^n}\,e^{-2\pi i
(x, y-z)}e^{-2\pi i (x,
z)}f(y-z)\,g(z)\,dz \\
& = & \int_{\mathbb{R}^n} \,e^{-2\pi i (x, u)} f(u)\,du
\int_{\mathbb{R}^n}e^{-2\pi i (x, z)}\,g(z) \,dz\\
& = & \hat{f}(x) \,\hat{g}(x).
\end{eqnarray*}
\end{enumerate}
\fincorrection
\correction{005963}
Supposons tout d'abord $n=1$. Soit la gaussienne d\'efinie pour
$x\in\mathbb{R}$ par $f(x) = e^{-\frac{a x^2}{2}}$, o\`u $a>0$.
Posons $$h(t) = \int_{\mathbb{R}} f(x) e^{-2\pi i tx}\,dx =
\int_{\mathbb{R}} e^{-\frac{a x^2}{2}} e^{-2\pi i tx}\,dx.
$$ D'apr\`es le th\'eor\`eme de convergence domin\'ee, $h$ est
d\'erivable et
\begin{eqnarray*} h'(t) &=& -2\pi i
\int_{\mathbb{R}}x e^{-\frac{a x^2}{2}} e^{-2\pi i tx}\,dx =
\left[2\pi i \frac{1}{a} e^{-\frac{a x^2}{2}}e^{-2\pi i tx}
\right]_{-\infty}^{+\infty} + (2\pi i)^2 t \frac{1}{a}
\int_{\mathbb{R}}e^{-\frac{a
x^2}{2}}  e^{-2\pi i tx}\,dx\\
& = & - (2\pi)^2 \frac{1}{a}t\cdot \,h(t).
\end{eqnarray*}
De plus, $$h(0) = \int_{\mathbb{R}} f(x) \,dx =  \int_{\mathbb{R}}
 e^{-\frac{a x^2}{2}} \,dx = \int_{\mathbb{R}} e^{-u^2} \,du
 \frac{\sqrt{2}}{\sqrt{a}} =  \frac{\sqrt{2\pi}}{\sqrt{a}}.$$
La solution de l'\'equation diff\'erentielle $h'(t)  =  - (2\pi)^2
\frac{1}{a}t \cdot\,h(t)$ avec condition initiale $h(0) =
\frac{\sqrt{2\pi}}{\sqrt{a}}$ est
$$
h(t) = \sqrt{\frac{2\pi}{a}} e^{-\frac{(2\pi)^2}{a}
\frac{t^2}{2}}.
$$

Pour $n>1$, on a~:
\begin{eqnarray*}
\int_{\mathbb{R}^n} f(x) e^{-2\pi i (t , x)}\,dx &=&
\int_{\mathbb{R}^n} e^{-\frac{a |x|^2}{2}} e^{-2\pi i (t,
x)}\,dx\\
& = & \prod_{i = 1}^n \int_{\mathbb{R}}e^{-\frac{a x_i^2}{2}}
e^{-2\pi i t_i
x_i}\,dx_i = \prod_{i = 1}^n h(t_i) = \left(\sqrt{\frac{2\pi}{a}}\right)^{n} e^{-\frac{(2\pi)^2}{a}\frac{|t|^2}{2}}\\
\end{eqnarray*}
\fincorrection
\correction{002692}
\begin{itemize}
\item Si $K$ est le support de $f$,
$f*g(x) =\int_K f(t) g(x-t)\,dt$ est bien d{\'e}fini. De
plus, $|f*g(x)-f*g(y)|\leq\int_K
|f(t)||g(x-t)-g(y-t)|\,dt \leq C \|f\|_{L^1}
|x-y|^\alpha$, d'o{\`u} le r{\'e}sultat. Si $g$ est
d{\'e}rivable, alors $f*g$ aussi et $(f*g)' = f*(g')$. Donc
si $g\in C^{k,\alpha}$, $f*g\in C^k$ et sa d{\'e}riv{\'e}e
$k$-i{\`e}me {\'e}tant $f*(g^{(k)})$, elle est h{\"o}ld{\'e}rienne par
le m{\^e}me argument.
\item M{\^e}me argument; les produits de convolution
sont bien d{\'e}finis et born{\'e}s car $f\in L^1$ et les
d{\'e}riv{\'e}es de~$g$ sont dans $L^\infty$.
\end{itemize}
\fincorrection
\correction{005954}
\begin{enumerate}
\item Soit $a, b \geq 0$ et soit $p, q \in(1, +\infty)$ tel que
$\frac{1}{p} + \frac{1}{q} = 1$. La fonction
$\theta~:\mathbb{R}^{+}\rightarrow \mathbb{R}$ d\'efinit par
$\theta(a) = \frac{1}{p} a^{p} + \frac{1}{q} b^{q} - ab$ est
d\'erivable et~:
$$
\theta'(a) =  a^{p-1} - b.
$$
Cette d\'eriv\'ee s'annule lorsque $a = b^{\frac{1}{p-1}}$, est
n\'egative pour  $a < b^{\frac{1}{p-1}}$ et positive pour $a >
b^{\frac{1}{p-1}}$. On a
$$
\theta(b^{\frac{1}{p-1}}) = \frac{1}{p} b^{\frac{p}{p-1}} +
\frac{1}{q}b^q- b^{1+\frac{1}{p-1}} = 0.
$$
Ainsi $\theta(a)\geq 0$, i.e. $$ ab \leq \frac{1}{p} a^p +
\frac{1}{q} b^q.
$$
%L'\'egalit\'e correspond \`a $\theta(a) = 0$, c'est-\`a-dire \`a
%$a = b^{\frac{1}{p-1}}$.


\item Soit $f\in L^{p}(\mu)$ et $g\in L^{q}(\mu)$. D'apr\`es la
question pr\'ec\'edente,  pour tout $\lambda>0$ et pour
$\mu$-presque tout $x$~:
$$
 |f\,g|(x) = |\lambda f(x) \cdot \frac{g(x)}{\lambda}| \leq \frac{\lambda^{p}}{p}
|f(x)|^{p} + \frac{\lambda^{-q}}{q}  |g(x)|^{q}.
$$
Ainsi
$$
\int_{\Omega} |f\,g|\,d\mu \leq \frac{\lambda^{p}}{p}\int_{\Omega}
|f|^{p}\,d\mu + \frac{\lambda^{-q}}{q} \int_{\Omega}
|g|^{q}\,d\mu.
$$
Posons $$\Phi(\lambda) = \frac{\lambda^{p}}{p}\int_{\Omega}
|f|^{p}\,d\mu + \frac{\lambda^{-q}}{q} \int_{\Omega}
|g|^{q}\,d\mu.$$ La fonction $\Phi$ est d\'erivable et~:
$$
\Phi'(\lambda) = \lambda^{p-1} \|f\|_{p}^{p} -
\lambda^{-q-1}\|g\|_{q}^{q}.
$$
Cette d\'eriv\'ee s'annule pour $\lambda_{1} :=
\left(\frac{\|g\|_{q}^{q}}{\|f\|_{p}^{p}}
\right)^{\frac{1}{p+q}}$, est n\'egative pour $\lambda\leq
\lambda_{1} $ et positive pour $\lambda \geq \lambda_{1}$. Ainsi
le minimum de $\Phi$ vaut~:
\begin{equation*}
\begin{array}{ll}
\Phi(\lambda_{1}) &=
\frac{1}{p}\left(\frac{\|g\|_{q}^{q}}{\|f\|_{p}^{p}}
\right)^{\frac{p}{p+q}} \|f\|_{p}^{p} + \frac{1}{q}
\left(\frac{\|g\|_{q}^{q}}{\|f\|_{p}^{p}}
\right)^{-\frac{q}{p+q}} \|g\|_{q}^{q}\\
& =
\frac{1}{p}\|g\|_{q}^{\frac{qp}{p+q}}\,\|f\|_{p}^{\frac{qp}{p+q}}
+
\frac{1}{q}\|g\|_{q}^{\frac{qp}{p+q}}\,\|f\|_{p}^{\frac{qp}{p+q}}
= \| f\|_{p}\,\|g\|_{q}.
\end{array}
\end{equation*}
On en d\'eduit l'in\'egalit\'e de H\"older~:
$$
\| f\,g\|_{1} \leq \|f\|_{p} \,\|g\|_{q}.
$$
Si $f\in L^{1}(\mu)$ et $g\in L^{\infty}(\mu)$, alors  $|g(x)|
\leq \|g\|_{\infty}$ pour presque tout $x\in\Omega$ et
$$
\int_{\Omega} |f g|\,d\mu \leq \|g\|_{\infty}\int_{\Omega}
|f|\,d\mu,
$$
i.e. $\|f g\|_{1}\leq \|g\|_{\infty} \|f\|_{1}$.


\item Soient $p, p' \in [1, +\infty)$. On suppose $p< p'$. Soit
$p<r<p'$. On a
$$
|f|^{r} = |f|^r \mathbf{1}_{|f|>1} + |f|^r \mathbf{1}_{|f|<1} ~\leq~ |f|^{p'}
\mathbf{1}_{|f|>1} + |f|^{p}\mathbf{1}_{|f|<1}.
$$
On en d\'eduit que
$$
\int_{\Omega} |f|^{r}\,d\mu ~\leq~ \int_{\Omega}|f|^{p'}\,d\mu +
\int_{\Omega} |f|^{p}\,d\mu ~<+\infty,
$$
donc $f$ appartient \`a $L^{r}(\mu)$.


\item Supposons que  $\mu$ soit une mesure finie et soit $f\in
L^{\infty}(\mu)$. Alors
$$
|f(x)| \leq \|f\|_{\infty}
$$
pour presque tout $x\in\Omega$. Ainsi pour tout $p$
$$
\int_{\Omega} |f|^{p}\,d\mu ~\leq~\|f\|_{\infty}^{p}\int_{\Omega}
1\,d\mu ~=~ \|f\|_{\infty}^{p} \mu(\Omega)~< ~+\infty,
$$
ce qui implique que $f\in L^{p}(\mu)$. En particulier, $f$
appartient \`a l'intersection $\bigcap_{p\geq 1} L^{p}(\mu)$. De
plus, pour tout $p$, on a~:
$$
\|f\|_{p} ~\leq ~\|f\|_{\infty} ~\mu(\Omega)^{\frac{1}{p}},
$$
ce qui implique que
$$
\lim_{p\rightarrow+\infty}\|f\|_{p}~ \leq ~\|f\|_{\infty}.
$$
D'autre part, pour tout $0<\varepsilon< \| f \|_\infty$ , on a
$$
\int_{\Omega} |f|^{p}\,d\mu ~\geq~
\int_{|f|>\left(\|f\|_{\infty}-\varepsilon\right)}|f|^{p}\,d\mu
~\geq ~\left(\|f\|_{\infty}-\varepsilon\right)^{p} \mu\bigg(
|f|>\left(\|f\|_{\infty}-\varepsilon\right)\bigg).
$$
Ainsi pour tout $p$, il vient
$$
\|f\|_{p} ~\geq~ \left(\|f\|_{\infty}-\varepsilon\right) \mu\bigg(
|f|>\left(\|f\|_{\infty}-\varepsilon\right)\bigg)^{\frac{1}{p}}.
$$
Puisque $\lim_{p\rightarrow +\infty}\mu\left(
|f|>\left(\|f\|_{\infty}-\varepsilon\right)\right)^{\frac{1}{p}} =
1$, il en d\'ecoule que
$$
\lim_{p\rightarrow +\infty}\|f\|_{p} ~\geq~
\|f\|_{\infty}-\varepsilon.
$$
Comme $\varepsilon$ peut \^etre choisi arbitrairement petit, on a
$$
\lim_{p\rightarrow +\infty}\|f\|_{p} ~\geq~ \|f\|_{\infty},
$$
donc finalement $\lim_{p\rightarrow +\infty}\|f\|_{p} =
\|f\|_{\infty}$.

\item   Posons $f_1:=f^r$ et $g_1:=g^r$. On a $f_1\in
L^{\frac{p}{r}}(\mu)$ et $g_1\in L^{\frac{q}{r}}(\mu)$. Notons que
l'identit\'{e} $\frac{1}{p}+\frac{1}{q}=\frac{1}{r}$ entra\^{i}ne
que $\frac{p}{r}, \frac{q}{r}>1$ et que les nombres $\frac{p}{r}$
et $ \frac{q}{r}$ sont conjugu\'{e}s au sens de Young. Par
l'in\'{e}galit\'{e} de H\"{o}lder on a donc
$$\int_{\Omega}(fg)^r d\mu = \int_{\Omega} f_1 g_1 \,d\mu \leq
\left(\int_{\Omega} f_1^{\frac{p}{r}}d\mu\right)^{\frac{r}{p}}
\left(\int_{\Omega} g_1^{\frac{q}{r}}d\mu\right)^{\frac{r}{q}}
=\left(\int_{\Omega} f^p
d\mu\right)^{\frac{r}{p}}\left(\int_{\Omega} g^q
d\mu\right)^{\frac{r}{q}}.$$ D'o\`{u}, finalement,
 $$\|fg\|_r\leq \|f\|_p\|g\|_q.$$
 \end{enumerate}
\fincorrection
\correction{005955}
\begin{enumerate}
\item \emph{Cas de $L^{\infty}(\mu)$.}

\begin{enumerate}
\item Soit $\{f_{n}\}_{n\in\mathbb{N}}$ une suite de Cauchy
de $L^{\infty}(\mu)$. Pour $k, m, n\geq 1$,  soient les ensembles
\begin{eqnarray*}
A_{k} := \{ x\in\Omega, |f_{k}(x)| > \|f_{k}\|_{\infty}\}~; &
B_{m,n} := \{ x\in\Omega, |f_{m}(x) - f_{n}(x)| > \|f_{m} -
f_n\|_{\infty} \},
\end{eqnarray*}
et $E := \bigcup_{k} A_{k}  \bigcup_{n,m} B_{m,n}$. Par
d\'efinition de la norme infinie, les ensembles $A_k$ et $B_{n,m}$
sont de mesure nulle. Par $\sigma$-sous-additivit\'e de $\mu$, on
a
$$
\mu(E) \leq \sum_{k} \mu(A_k) + \sum_{n,m}\mu(B_{n,m}) = 0.
$$

\item Sur $\Omega\setminus E$, on a~:
$$
\sup_{x\in\Omega\setminus E}|f_{n} - f_{m}| \leq \| f_{n} -
f_{m}\|_{\infty},
$$
i.e. $\{f_{n}\}_{n\in\mathbb{N}}$ est une suite de Cauchy uniforme
sur $\Omega\setminus E$. En particulier, pour tout
$x\in\Omega\setminus E$, la suite $\{f_{n}(x)\}_{n\in\mathbb{N}}$
est une suite de Cauchy r\'eelle, donc est convergeante  car
$\mathbb{R}$ est complet. Notons $f$ la limite ponctuelle de
$f_{n}$ sur $\Omega\setminus E$. Montrons que la suite
$\{f_n\}_{n\in\mathbb{N}}$ converge uniform\'ement vers $f$ sur le
compl\'ementaire de $E$. On a
$$
|f_{n}(x) - f(x)| = \lim_{m\rightarrow +\infty}|f_{n}(x) -
f_{m}(x)|\leq\lim_{m\rightarrow +\infty}\|f_{n} -
f_{m}\|_{\infty}.
$$
Comme $\{f_{n}\}_{n\in\mathbb{N}}$ est de Cauchy dans
$L^{\infty}(\mu)$, pour tout $\varepsilon>0$, il existe un rang
$N_{\varepsilon}$ tel que pour $n,m > N_{\varepsilon}$, $\|f_{n} -
f_{m}\|_{\infty} < \varepsilon$.  Alors pour $n> N_{\varepsilon}$,
$$
\sup_{x\in\Omega\setminus E}|f_{n}(x) - f(x)| \leq \varepsilon.
$$
Il est d\'ecoule que  $\{f_n\}_{n\in\mathbb{N}}$ converge
uniform\'ement vers $f$ sur $\Omega\setminus E$.




\item \'Etendons la fonction $f$ \`a $\Omega$ en posant $f =
0$ sur $E$. Il reste \`a montrer que la fonction $f$ appartient
\`a $L^{\infty}(\mu)$. Pour $n> N_{\varepsilon}$, et
$x\in\Omega\setminus E$, on a
$$
|f(x)| < |f_n(x)| + \varepsilon \leq \|f_n(x)\|_{\infty}  +
\varepsilon
$$
On en d\'eduit que $\|f\|_{\infty} \leq \|f_n(x)\|_{\infty}  +
\varepsilon < +\infty$. Ainsi $L^{\infty}(\mu)$ est complet.

\end{enumerate}


\item \emph{Cas de $L^{p}(\mu)$.}
\begin{enumerate}
\item Soit $1\leq p<+\infty$ et $\{f_{n}\}_{n\in\mathbb{N}}$
une suite de Cauchy dans $L^{p}(\mu)$. Il existe $n_{1}$ tel que
pour $n,m\geq n_{1}$, $\|f_{n} - f_{m}\|_{p} < 2^{-1}$. On prend
ensuite $n_{2}>n_{1}$ tel que pour $n,m \geq n_{2}$, $\|f_{n} -
f_{m}\|_{p} < {2}^{-2}$, et ainsi de suite, pour tout $k$, il
existe un $n_{k} > n_{k-1}$ tel que $n,m\geq n_{k}\Rightarrow
\|f_{n} - f_{m}\|_{p} < {2}^{-k} $.

 \item
Posons
\begin{eqnarray*}
g_{k} = \sum_{i=1}^{k}|f_{n_{i+1}} - f_{n_{i}}|\quad
\text{et}\quad g= \sum_{i=1}^{+\infty}|f_{n_{i+1}} - f_{n_{i}}|,
\end{eqnarray*}
o\`u $g$ est \`a valeurs dans $\mathbb{R}\cup\{+\infty\}$. Pour
tout $k\geq 1$, on a
$$
\|g_{k}\|_{p} = \| \sum_{i=1}^{k}|f_{n_{i+1}} - f_{n_{i}}|
~\|_{p}.
$$
D'apr\`es l'in\'egalit\'e de Minkowski,
$$
\|g_{k}\|_{p} \leq \sum_{i=1}^{k} \|f_{n_{i+1}} - f_{n_{i}}\|_{p}
=\sum_{i=1}^{k} 2^{-i} < 1.
$$
D'apr\`es le lemme de Fatou, on en d\'eduit que $\|g\|_{p} \leq
1$.



\item Comme $\int_{\Omega} |g|^{p} \,d\mu < +\infty$,
n\'ecessairement $|g| < +\infty$ $\mu$-pp, i.e. pour presque tout
$x\in\Omega$ la s\'erie
$$f_{n_{1}} + \sum_{i=1}^{\infty} \left(f_{n_{i+1}} -
f_{n_{i}}\right)$$ est absolument convergente. Notons $f(x)$ sa
somme lorsque celle-ci est finie et posons $f(x)=0$ sinon. On a~:
$$
f_{n_{1}} + \sum_{i=1}^{k-1} \left(f_{n_{i+1}} -
f_{n_{i}}\right) = f_{n_{k}}
$$
et $f(x) = \lim_{k\rightarrow
+\infty}f_{n_{k}}\quad\mu-\text{pp}$.

\item Soit $\varepsilon>0$. Comme
$\{f_{n}\}_{n\in\mathbb{N}}$ est de Cauchy dans $L^{p}(\mu)$, il
existe $N_{\varepsilon}>0$ tel que pour $n, m > N_{\varepsilon}$,
$\|f_{n} - f_{m}\|_{p} < \varepsilon$. Pour $m>N_{\varepsilon}$ on
a par le lemme de Fatou~:
$$
\int_{\Omega}| f - f_{m}|^{p}\,d\mu = \int_{\Omega}
|\lim_{k\rightarrow+\infty} f_{n_{k}} - f_{m}|^{p}\,d\mu \leq
\lim\inf_{k\rightarrow+\infty}\int_{\Omega} |f_{n_{k}} -
f_{m}|^{p}\,d\mu \leq \varepsilon^{p}.
$$
Ainsi $f - f_{m} \in L^{p}(\mu)$ et $\|f - f_{m}\|_{p}\rightarrow
0$ quand $m\rightarrow +\infty$. De plus, d'apr\`es
l'in\'egalit\'e de Minkowski, on a
$$
\|f\|_{p} = \|(f - f_{m}) + f_{m}\|_{p}\leq \|(f - f_{m}) \|_{p} +
\|f_{m}\|_{p}<+\infty,
$$
c'est-\`a-dire $f\in L^{p}(\mu)$. En conclusion $L^{p}(\mu)$ est
complet.
\end{enumerate}
\end{enumerate}
\fincorrection
\correction{005956}
Soient $f$ et $g$ deux fonctions de $L^{p}(\mu)$ avec
$1<p<+\infty$. La fonction $\varphi(t) = |f(x)+\tan (x)|^{p}$ est de
classe $\mathcal{C}^1$ sur $\mathbb{R}$ et sa d\'eriv\'ee vaut
$$
\varphi'(t) = \lim_{h\rightarrow 0}\frac{|f(x) + \tan (x) + h
g(x)|^{p}-|f(x)+ \tan (x)|^p}{h} = p |f(x)+\tan (x)|^{p-2} (f(x)+ \tan (x))
g(x),
$$
lorsque $f(x)$ et $g(x)$ ont un sens, c'est-\`a-dire pour presque
tout $x$. De plus, d'apr\`es le th\'eor\`eme des accroissements
finis, on a
$$
\frac{|f(x) + \tan (x)|^{p}-|f(x)|^p}{t} = \varphi'(t_{0}) = p
|f(x)+t_0 g(x)|^{p-2} (f(x)+ t_0 g(x)) g(x),
$$
pour un certain $t_0$ compris entre $0$ et $t$. Ainsi pour
$|t|\leq 1$,
\begin{equation*}
\begin{array}{ll}
\left|\frac{|f(x) + \tan (x)|^{p}-|f(x)|^p}{t}\right| &=  p |f(x)+t_0
g(x)|^{p-1} |g(x)| \\& \leq p \left(|f(x)| + |g(x)|\right)^{p}\\ &
\leq  2^{p-1}\,p \left(|f(x)|^p + |g(x)|^p \right),
\end{array}
\end{equation*}
o\`u la premi\`ere in\'egalit\'e d\'ecoule de l'in\'egalit\'e
triangulaire et de la majoration $|g(x)| \leq (|f(x)| + |g(x)|)$,
et o\`u la deuxi\`eme in\'egalit\'e provient de la convexit\'e de
la fonction $x\mapsto x^p$ pour $p>1$ impliquant en particulier~:
$\left(\frac{u+v}{2} \right)^p \leq \frac{u^p}{2} +
\frac{v^p}{2}$. Il en d\'ecoule que $t\mapsto \frac{|f(x) +
\tan (x)|^{p}-|f(x)|^p}{t}$ est uniform\'ement born\'ee par une
fonction int\'egrable. Le th\'eor\`eme de convergence domin\'ee
permet alors de d\'eriver sous le signe somme et
$$
\frac{dN}{dt}_{t=0} = p\int_{\Omega} |f(x)|^{p-2} f(x) g(x)\,d\mu.
$$
\fincorrection
\correction{005964}
Soit $\Omega$ un sous-ensemble de $\mathbb{R}^n$ dont la mesure de
Lebesgue   est \emph{finie}~: $\mu(\Omega)< +\infty$. Pour tout
$1\leq p <+\infty$, notons $L^p(\Omega)$ l'espace des fonctions
$f~:\Omega \rightarrow \mathbb{C}$ telles que $\|f\|_{p} :=
\left(\int_{\Omega} |f|^{p}(x)\,dx\right)^{\frac{1}{p}}<+\infty$
modulo l'\'equivalence $f\sim g \Leftrightarrow f-g = 0 ~\mu-p.p$.
L'espace des fonctions essentiellement born\'ees sera not\'e
$L^{\infty}(\Omega).$
\begin{enumerate} 
\item Si $f\in L^{\infty}(\Omega)$, alors
$$
\| f \|_p^p = \int_{\Omega} |f|^{p}(x)\,dx \leq \|f\|_{\infty}^{p}
\mu(\Omega)<+\infty,
$$
ainsi $L^{\infty}(\Omega) \subset L^{p}(\Omega)$ pour tout $p$ et
$\|f\|_{p} \leq
\|f\|_{\infty}\left(\mu(\Omega)\right)^{\frac{1}{p}}$. Montrons
que si $q\leq p$, alors $L^{p}(\Omega) \subset L^{q}(\Omega)$.
Soit $f\in L^{p}(\Omega)$, on a par exemple~:
\begin{eqnarray*}
\| f \|_q^q = \int_{\Omega}|f|^{q}(x)\,dx &=& \int_{\{|f|\geq1\}}|f|^{q}(x)\,dx
+
\int_{\{|f|<1\}} |f|^{q}(x)\,dx \\
& \leq & \int_{\{|f|\geq1\}}|f|^{p}(x)\,dx + \int_{\{|f|<1\}}
1\,dx\\
& \leq & \|f\|_{p}^{p} + \mu(\Omega)<+\infty.
\end{eqnarray*}
Ou encore, en utilisant l'in\'egalit\'e de H\"older pour les
r\'eels conjugu\'es $r = \frac{p}{q} > 1$ et $r' =
\frac{p}{p-q}$~:
\begin{eqnarray*}
\| f \|_q^q =  \int_{\Omega}|f|^{q}(x)\,dx &=& \left( \int_{\Omega}|f|^{q\cdot
\frac{p}{q}}(x)\,dx\right)^{\frac{q}{p}}
\left(\int_{\Omega} 1^{\frac{p}{p-q}}(x)\,dx\right)^{\frac{p-q}{p}} \\
& = & \|f\|_{p}^{q}~\mu(\Omega)^{\frac{p-q}{p}},
\end{eqnarray*}
ce qui implique~:
$$
\|f\|_{q} \leq \|f\|_{p} \mu(\Omega)^{\frac{p-q}{qp}}.
$$
 En conclusion, pour $1<q<2<p$~:
$$
L^{\infty}(\Omega) \subset L^{p}(\Omega) \subset L^{2}(\Omega)
\subset L^{q}(\Omega) \subset L^{1}(\Omega).
$$


\item Montrons que pour $q<p$, l'inclusion
$L^{p}(\mathcal{B}^{n}(0,1)) \subset
L^{q}(\mathcal{B}^{n}(0,1))$ est stricte. La fonction
$f_{\alpha}$ appartient \`a $L^{\infty}(\mathcal{B}^{n}(0,1))$
si et seulement $\alpha \leq 0$, et \`a
$L^{p}(\mathcal{B}^{n}(0,1))$ avec $p<+\infty$ si et seulement
si
$$p\alpha - n +1 < 1 \Leftrightarrow \alpha < \frac{n}{p}
$$
Soit $1\leq q<p$, alors
$f_{\frac{1}{2}\left(\frac{n}{p}+\frac{n}{q}\right)}$ appartient
\`a $L^{q}(\mathcal{B}^{n}(0,1))\setminus
L^{p}(\mathcal{B}^{n}(0,1))$. En particulier,
$f_{\frac{1}{2}\left(\frac{n}{p}+\frac{n}{q}\right)}$ appartient
\`a $L^{q}(\mathcal{B}^{n}(0,1))\setminus
L^{\infty}(\mathcal{B}^{n}(0,1))$.
\end{enumerate}
\fincorrection
\correction{005965}
Soit $\Omega = \mathbb{N}$ muni de la mesure de comptage. Pour
tout $1\leq p <+\infty$, on note $\ell^p$ l'espace des suites
complexes $(u_{n})_{n\in\mathbb{N}}$ telles que $\|u\|_{p} :=
\left(\sum_{i=0}^{+\infty}
|u_{n}|^{p}\right)^{\frac{1}{p}}<+\infty$. L'espace des suites
born\'ees sera not\'e $\ell^{\infty}.$
\begin{enumerate}
\item Montrons que si $q\leq p$, alors $\ell^{q}\subset \ell^{p}$.
Soit $(u_{n})_{n\in\mathbb{N}}\in \ell^{q}$. Comme
$$
\sum_{i=0}^{+\infty} |u_{n}|^{q} <+\infty,
$$
il existe un rang $N$ tel que pour $n> N$, $|u_{n}|^{q}< 1$. En
particulier la suite $(u_{n})_{n\in\mathbb{N}}$ appartient \`a
$\ell^{\infty}$
 et
$$\|u\|_{\infty} \leq \max\{u_{0}, \dots, u_{N}, 1\}.$$
De plus, pour $n> N$, on a $|u_{n}|^{p} \leq |u_{n}|^{q}$ et
$$
\sum_{i=N+1}^{+\infty} |u_{n}|^{p} \leq \sum_{i=N+1}^{+\infty}
|u_{n}|^{q}\leq \|u\|_{q}^{q} <+\infty,
$$
ce qui implique que $\|u\|_{p} < +\infty$. En conclusion, pour
$1<q<2<p$, on a~:
$$
\ell^{1} \subset \ell^{q} \subset \ell^{2} \subset \ell^{p}
\subset \ell^{\infty}.
$$
\item La suite $u^{(\alpha)}_{n} = n^{-\alpha}$ appartient \`a
$\ell^{\infty}$ pour tout $\alpha\geq 0$ et \`a $\ell^{p}$ avec
$1\leq p <+\infty$ si et seulement si $\alpha p>1$, i.e $\alpha>
\frac{1}{p}$. En particulier la suite constante \'egale \`a $1$
appartient \`a $\ell^{\infty}$ mais n'appartient \`a aucun
$\ell^{p}$ pour $p<+\infty$. Soit $1<q<p<+\infty$. Pour tout
$\alpha$ tel que $\frac{1}{p} < \alpha < \frac{1}{q}$, la suite
$u^{(\alpha)}$ appartient \`a $\ell^{p}\setminus \ell^q$. C'est le
cas en particulier pour $\alpha = \frac{1}{2}\left(\frac{1}{p} +
\frac{1}{q} \right)$. Ainsi l'inclusion $\ell^{q} \subset
\ell^{p}$ est stricte lorsque $q< p$.
\end{enumerate}
\fincorrection
\correction{005966}
Soit $\Omega = \mathbb{R}^{n}$ muni de la mesure de Lebesgue. Pour
tout $1\leq p <+\infty$, on note $L^p(\mathbb{R}^{n})$ l'espace
des fonctions $f~:\mathbb{R}^{n} \rightarrow \mathbb{C}$ telles
que $\|f\|_{p} := \left(\int_{\mathbb{R}^{n}}
|f|^{p}(x)\,dx\right)^{\frac{1}{p}}<+\infty$ modulo
l'\'equivalence $f\sim g \Leftrightarrow f-g = 0 ~\mu-p.p$.
L'espace des fonctions essentiellement born\'ees sera not\'e
$L^{\infty}(\mathbb{R}^n).$
\begin{enumerate}
\item 
  \begin{itemize} 
  \item  La fonction $x\mapsto
\frac{1}{\left(1 + |x|^2 \right)^{\alpha}}$ appartient \`a
$L^{p}(\mathbb{R}^{n})$ pour $2\alpha p > n$. 
  \item La fonction $x\mapsto \frac{1}{|x|^{\beta}} e^{-\frac{|x|^2}{2}}$
appartient \`a $L^{p}(\mathbb{R}^{n})$ pour $p\beta< n$.
  \item Soit $1\leq q < p \leq +\infty$. La fonction
$$
f(x) = (1+|x|^2)^{-\frac{n}{p+q}}
$$
v\'erifient $f \in L^{p}(\mathbb{R}^{n})$ et $f\notin
L^{q}(\mathbb{R}^{n})$.  La fonction
$$
g(x) = |x|^{-\frac{n}{2}\left(\frac{1}{p}+\frac{1}{q}\right)}
e^{-\frac{|x|^2}{2}}
$$
v\'erifient $g\in L^{q}(\mathbb{R}^{n})$ et $g\notin
L^{p}(\mathbb{R}^{n})$.

\end{itemize}
\item 
  \begin{itemize} 
  \item Soit $1\leq q< p<+\infty$ et
$f_n$ une suite de Cauchy pour la norme $\|\cdot\|_{p, q}= \|
\cdot \|_{p} + \| \cdot \|_{q}$. Comme $\|\cdot\|_{p} \leq
\|\cdot\|_{p, q}$, $f_{n}$ est une suite de Cauchy dans
$L^{p}(\mathbb{R}^{n})$, donc elle converge vers une fonction
$f\in L^{p}(\mathbb{R}^{n})$ pour la norme $\|\cdot\|_{p}$. De
m\^eme,  $\|\cdot\|_{q}\leq \|\cdot\|_{p, q}$, donc $f_{n}$
converge vers une fonction $g\in L^{q}(\mathbb{R}^{n})$ pour la
norme $\|\cdot\|_q$. De plus, il existe une sous-suite de
$f_{n_{k}}$ qui converge vers $f$ presque-partout et il existe une
sous-suite de $f_{n_{k}}$ qui converge vers $g$ presque-partout.
Ainsi $f=g ~\mu$-p.p. et $f_{n}$ converge vers $f = g$ dans
$L^{p}(\mathbb{R}^{n})\cap L^{q}(\mathbb{R}^{n})$.

  \item Soit $r$ tel que $q< r<p$. Montrons que
$$
\|f\|_{r} \leq \|f\|_{p}^{\alpha}\|f\|_{q}^{1-\alpha}
$$
o\`u $\frac{1}{r} = \frac{\alpha}{p} + \frac{1-\alpha}{q}$,
$\alpha\in[0, 1]$. Puisque $1 = \frac{\alpha r}{p} +
\frac{(1-\alpha)r}{q}$, les r\'eels $p' = \frac{p}{\alpha r}$ et
$q' = \frac{q}{(1 - \alpha)r}$ sont conjugu\'es. D'apr\`es
l'in\'egalit\'e de H\"older,
\begin{eqnarray*}
\int_{\mathbb{R}^{n}} |f|^{r}(x)\,dx& =& \int_{\mathbb{R}^{n}}
|f|^{r\alpha}(x)\cdot |f|^{(1-\alpha) r}(x)\,dx\\
& \leq & \left(\int_{\mathbb{R}^{n}} |f|^{\alpha r p'}(x)\,dx
\right)^{\frac{1}{p'}} \left( \int_{\mathbb{R}^{n}}
|f|^{(1-\alpha)r q'}(x)\,dx\right)^{\frac{1}{q'}}\\
& \leq & \left(\int_{\mathbb{R}^{n}} |f|^{p}(x)\,dx
\right)^{\frac{\alpha r}{p}} \left( \int_{\mathbb{R}^{n}}
|f|^{q}(x)\,dx\right)^{\frac{(1-\alpha)r}{q}}\\
& \leq & \|f\|_{p}^{\alpha r} \|f\|_{q}^{(1-\alpha)r},
\end{eqnarray*}
ce qui \'equivaut \`a $\|f\|_{r} \leq
\|f\|_{p}^{\alpha}\|f\|_{q}^{(1-\alpha)}$. On peut \'egalement
\'ecrire  $r = \beta q + (1-\beta)p$ avec $\beta \in ]0, 1[$ et
appliquer H\"older avec les r\'eels conjugu\'es $\frac{1}{\beta}$
et $ \frac{1}{1-\beta}$~:
\begin{eqnarray*}
\int_{\mathbb{R}^{n}} |f|^{r}(x)\,dx& =& \int_{\mathbb{R}^{n}}
|f|^{\beta q}(x)\cdot|f|^{(1-\beta)p}(x)\,dx \\
&\leq & \left(\int_{\mathbb{R}^{n}} |f|^{q}(x)\,dx \right)^{\beta}
\left( \int_{\mathbb{R}^{n}} |f|^{p}(x)\,dx\right)^{(1-\beta)},
\end{eqnarray*}
ce qui implique
$$
\|f\|_{r} \leq
\|f\|_{q}^{\frac{q\beta}{r}}\|f\|_{p}^{\frac{p(1-\beta)}{r}}
$$
qui est l'in\'egalit\'e cherch\'ee car $\alpha = \frac{p\beta}{r}$
v\'erifie bien $\frac{1}{r} = \frac{\alpha}{p} +
\frac{1-\alpha}{q}$.


  \item Si $f_{n}$ converge vers $f$ dans
$L^{p}(\mathbb{R}^{n}) \cap L^{q}(\mathbb{R}^{n})$ alors $f_{n}$
converge vers $f$ dans $L^{p}(\mathbb{R}^{n})$ et dans
$L^{q}(\mathbb{R}^{n})$, donc dans $L^{r}(\mathbb{R}^{n})$
d'apr\`es l'in\'egalit\'e pr\'ec\'edente. En conclusion,
$L^{p}(\mathbb{R}^{n}) \cap L^{q}(\mathbb{R}^{n})$ est ferm\'e
dans $L^{r}(\mathbb{R}^{n})$ donc un sous-espace de Banach de
$L^{r}(\mathbb{R}^{n})$.
\end{itemize}

\item  Soit $f\in L^{p}([0, +\infty[)\cap L^{q}([0, +\infty[)$ et
$h$ la fonction d\'efinie par $h(r) = \frac{1}{\sqrt{r}}f(r)$. On
notera $p'$ le conjugu\'e de $p$ et $q'$ le conjugu\'e de $q$.
Montrons que $h$ appartient \`a $L^{1}([0, +\infty[)$. On a~:
\begin{eqnarray*}
\int_{0}^{+\infty} \frac{1}{\sqrt{r}} |f(r)|\,dr = \int_{0}^{R}
\frac{1}{\sqrt{r}} |f(r)|\,dr + \int_{R}^{+\infty}
\frac{1}{\sqrt{r}} |f(r)|\,dr\\
\leq \left( \int_{0}^{R}
r^{-\frac{p'}{2}}\,dr\right)^{\frac{1}{p'}} \left(\int_{0}^{R}
|f(r)|^{p} \,dr\right)^{\frac{1}{p}} + \left(\int_{R}^{+\infty}
r^{-\frac{q'}{2}}\,dr \right)^{\frac{1}{q'}} \left(
\int_{R}^{+\infty}
|f(r)|^{q}\,dr\right)^{\frac{1}{q}}\\
 \leq  \left(\frac{1}{1-\frac{p'}{2}}\right)^{\frac{1}{p'}}
R^{\left(\frac{1}{p'}-\frac{1}{2}\right)}\|f\|_{p} +
\left(\frac{1}{\frac{q'}{2} - 1}\right)^{\frac{1}{q'}}
R^{\left(\frac{1}{q'}-\frac{1}{2}\right)} \|f\|_{q}.
\end{eqnarray*}
En optimisant par rapport \`a $R$, on obtient~:
\begin{eqnarray*}
\int_{0}^{+\infty} \frac{1}{\sqrt{r}} |f(r)|\,dr &\leq & C_{p,q}
\|f\|_{p}^{1-\gamma}\|f\|_{q}^{\gamma},
\end{eqnarray*}
o\`u, en posant $\alpha =  \frac{1}{2}- \frac{1}{p}$ et $\beta =
 \frac{1}{q} - \frac{1}{2}$, on a $\gamma =
\frac{\alpha}{\alpha+\beta}$, et $C_{p,q} = \frac{\alpha +
\beta}{\alpha^{\gamma} \beta^{1-\gamma}}\left(1-\frac{p'}{2}
\right)^{-\frac{1-\gamma}{p'}}\left(\frac{q'}{2} - 1
\right)^{-\frac{\gamma}{q'}}$.
\end{enumerate}
\fincorrection
\correction{005967}
Soit $\{f_{n}\}_{n\in\mathbb{N}}$ la suite de fonctions d\'efinies
par~:
$$
f_{n}(x) = \frac{1}{\sqrt{n}} \mathbf{1}_{[n, 2n]}(x).
$$
\begin{enumerate}
\item Quelque soit $g$ continue \`a support compact,
$$
\int_{[0, +\infty[} f_{n}(x) g(x)\,dx = \frac{1}{\sqrt{n}}
\int_{n}^{2n} g(x)\,dx \rightarrow 0
$$
quand $n\rightarrow +\infty$. Par densit\'e des fonctions
continues \`a support compact, $f_{n}$ converge faiblement vers
$0$. D'autre part, $f_n$ converge presque partout vers $0$.
Supposons que $f_n$ converge fortement vers une fonction $f$ dans
$L^2([0, +\infty[)$. Alors il existe une sous-suite de $f_n$ qui
converge presque-partout vers $f$, ce qui implique que $f =0$ est
la seule limite possible. Or~:
$$
\|f_{n}\|_{2} = 1
$$
pour tout $n$, donc $\|f_{n}\|_2$ ne tend pas vers $\|f\|_2 = 0$
ce qui contredit le fait que $f_n$ converge vers $f$ dans $L^2([0,
+\infty[)$.


\item Pour $p>2$, on a~:
$$
\int_{[0, +\infty[} |f_{n}(x)|^p\,dx =  \int_{n}^{2n}
n^{-\frac{p}{2}}\,dx = n^{1 - \frac{p}{2}} \rightarrow 0,
$$
quand $n\rightarrow+\infty$ donc $f_{n}$ converge fortement vers
$0$ dans $L^{p}([0, +\infty[)$.
\end{enumerate}
\fincorrection
\correction{005968}
Soit $\{f_{n}\}_{n\in\mathbb{N}}$ la suite de fonctions d\'efinies
par~:
$$
f_{n}(x) = \sqrt{n} \mathbf{1}_{[n, n+\frac{1}{n}]}(x).
$$
\begin{enumerate}
\item Quelque soit $g$ continue \`a support compact,
$$
\int_{[0, +\infty[} f_{n}(x) g(x)\,dx = \sqrt{n}
\int_{n}^{n+\frac{1}{n}} g(x)\,dx \rightarrow 0
$$
quand $n\rightarrow +\infty$. Par densit\'e des fonctions
continues \`a support compact, $f_{n}$ converge faiblement vers
$0$. Comme $f_n$ converge presque partout vers $0$ on conclut
comme pr\'ec\'edemment que $f_n$ ne converge pas fortement vers
$0$ dans $L^2([0, +\infty[)$ car
$$
\|f_{n}\|_{2} = 1.
$$

\item Pour $p<2$, on a~:
$$
\int_{[0, +\infty[} |f_{n}(x)|\,dx = \int_{n}^{n+\frac{1}{n}}
n^{\frac{p}{2}}\,dx = n^{ \frac{p}{2}-1} \rightarrow 0,
$$
donc $f_{n}$ converge fortement vers $0$ dans $L^{p}([0,
+\infty[)$.
\end{enumerate}
\fincorrection
\correction{005977}
cf E.~Lieb et M.~Loss, \emph{Analysis}, p.118, American Mathematical Society (2001).
(Pour la question~6, on peut utiliser la continuit\'e du produit
scalaire dans $L^{2}(\mathbb{R}^{n})$.)
\fincorrection
\correction{005978}
A l'aide des coordonnées sphériques, on a
$$
\begin{array}{lcl}
\hat{f}(k) &= & \int_{\mathbb{R}^{3}} f(x) e^{-2\pi i(x, k)}\,dx\\
& = & \\ & = & \int_{r =0}^{+\infty}\int_{\theta
=0}^{\pi}\int_{\varphi = 0}^{2\pi} h(r) e^{-2\pi i r
|k|\cos\theta} r^2 \sin\theta\,d\theta
dr d\varphi\\& = & \\
& = & 2\pi \int_{0}^{+\infty}\int_{0}^{\pi} h(r)
\frac{d}{d\theta}\left(\frac{1}{2\pi i r |k|} e^{-2\pi i r |k|\cos
\theta} \right) r^2 \,d\theta dr\\& = & \\
& = &\frac{1}{|k|} \int_{0}^{\infty} h(r) r
\frac{1}{i}\left[e^{+2\pi i r|k|} - e^{-2\pi i r|k|}\right]\, dr\\& = & \\
 & = &\frac{2}{|k|}\int_{0}^{+\infty} h(r) r \sin(2\pi|k|r)\,dr.
\end{array}
$$
\fincorrection
\correction{005928}
\begin{enumerate}
\item 
  \begin{enumerate}
  \item cf cours.  
  \item clair.  
  \item Soit $\{A_{i}\}_{i\in
\mathbb{N}^*}$ un suite quelconque d'ensembles $m_*$-mesurables. On
pose $B_1 = \emptyset$, $B_{2}=A_1$ et $B_{j} = \cup_{i=1}^{j-1}
A_{i}$, pour $j\geq 2$. Soit $Q$ un sous-ensemble de $\Omega$.
Montrons par r\'ecurrence que l'assertion $(P_k)$ suivante est
v\'erifi\'ee pour tout $k\geq 1$~:
$$
(P_k)\quad \quad \quad m_*(Q) = m_*(Q \cap B_{k+1}^c) + \sum
_{j=1}^{k} m_*(Q \cap B_{j}^c \cap A_j).
$$
\begin{itemize}
\item Pour $k = 1$, $(P_1)$ dit simplement que $m_*(Q) = m_*(Q\cap
A_{1}^c) + m_*(Q\cap A_1)$. Ceci   est une cons\'equence de la
$m_*$-mesurabilit\'e de $A_1$ et de fait que
$$
m_*(Q) \leq  m_*(Q\cap A_{1}^c) + m_*(Q\cap A_1)
$$
(on applique la $\sigma$-sous-additivit\'e de $m_*$ \`a $C_1 =
Q\cap A_{1}^c$, $C_2 = Q\cap A_1$ et $C_i = \emptyset$ pour $i\geq
3$.) 
\item Montrons que $(P_{k}) \Rightarrow (P_{k+1})$~:\\
Puisque $A_{k+1}$ est $m_*$-mesurable, on a~:
$$
m_*(Q\cap B_{k+1}^c) = m_*(Q \cap B_{k+1}^c \cap A_{k+1}^c) +
m_*(Q\cap B_{k+1}^c \cap A_{k+1}).$$ Or $B_{k+1}^c \cap A_{k+1}^c
= (B_{k+1}\cup A_{k+1})^{c} = B_{k+2}^{c}$. Ainsi~:
\begin{equation}\label{la}
m_*(Q\cap B_{k+1}^c) = m_*(Q \cap B_{k+2}^{c}) + m_*(Q\cap
B_{k+1}^c \cap A_{k+1}).
\end{equation}
Supposons que l'assertion $(P_{k})$ soit v\'erifi\'ee, alors
$$
m_*(Q) = m_*(Q \cap B_{k+1}^c) + \sum _{j=1}^{k} m_*(Q \cap
B_{j}^c \cap A_j),
$$
et d'apr\`es \eqref{la}
\begin{eqnarray*} m_*(Q) &=& m_*(Q \cap B_{k+2}^{c})
+ m_*(Q\cap B_{k+1}^c \cap A_{k+1})+ \sum _{j=1}^{k} m_*(Q \cap
B_{j}^c \cap A_j) \\&=& m_*(Q \cap B_{k+2}^{c})+ \sum _{j=1}^{k+1}
m_*(Q \cap B_{j}^c \cap A_j),
\end{eqnarray*}
qui n'est autre que $(P_{k+1})$. \item En conclusion, comme
$(P_1)$ est vrai et $(P_{k})\Rightarrow (P_{k+1})$, il en
d\'ecoule que l'assertion $(P_{k})$ est vraie pour tout $k\geq
1$.
\end{itemize}

  \item Comme $B_{k+1} \subset A$, on a $Q \cap B_{k+1}^c \supset Q
\cap A^c$ et, par  monotonie de $m_*$, $$m_*(Q\cap B_{k+1}^c)
\geq m_*(Q \cap A^c).$$ La condition $(P_{k})$ entra\^ine alors
que pour tout $k$~:
$$
m_*(Q) \geq m_*(Q \cap A^c) + \sum_{j=1}^{k} m_*(Q\cap B_{j}^c
\cap A_j).
$$
Donc, en faisant tendre $k$ vers $+\infty$~:
$$
m_*(Q) \geq m_*(Q \cap A^c) + \sum_{j=1}^{\infty} m_*(Q\cap
B_{j}^c \cap A_j).
$$

  \item On a~:$~Q\cap A = \bigcup_{j=1}^{\infty}(Q\cap B_{j}^c \cap
A_{j})~$ et par  $\sigma$-sous-additivit\'e de $m_*$~:
\begin{eqnarray*}
m_*(Q\cap A^c) + m_*(Q\cap A) &=&  m_*(Q\cap A^c) +
m_*\left(\bigcup_{j=1}^{\infty}(Q\cap B_{j}^c \cap A_j\right)\\
& \leq & m_*(Q\cap A^c) + \sum_{j=1}^{\infty} m_*(Q\cap B_{j}^c
\cap A_j)\\&\leq & m_*(Q).
\end{eqnarray*}
On en conclut que $A = \cup_{j=1}^{\infty} A_{j}$ est
$m_*$-mesurable.
\end{enumerate}


\item 
  \begin{enumerate}
  \item cf cours.
  \item Soit $\{A_{i}\}_{i\in\mathbb{N}^*}$ une
suite d'\'el\'ements $m_*$-mesurables, deux \`a deux disjoints.
Choisissons $Q = A = \cup_{j=1}^{\infty} A_{j}$, alors $Q \cap A^c
= \emptyset$ et $Q \cap B_{j}^c \cap A_j = A_j$ pour tout $j$.
D'apr\`es la question 1.d),
$$
m_{*}(Q) \geq \sum_{j=1}^{\infty} m_*(A_{j}).
$$
D'apr\`es la $\sigma$-sous-additivit\'e de $m_*$, il vient~:
$$
m_{*}(Q) = \sum_{j=1}^{\infty} m_*(A_{j}).
$$
  \end{enumerate}


 \item Soit $E$ un ensemble $m_*$-mesurable tel que $m_*(E) = 0$ et $B$ un sous-ensemble de $E$.
 Comme $Q \cap B^c \subset Q$, on a par monotonie de
 $m_*$ l'in\'egalit\'e $m_*(Q \cap B^c) \leq m_*(Q)$. Comme
 $Q\cap B \subset E$, on a aussi $m_*(Q\cap B) = 0$. On en
 d\'eduit que $$m_*(Q) \geq m_*(Q \cap B^c)+ m_*(Q\cap B).$$ Ainsi
 $B$ est $m_*$-mesurable et $m$ est compl\`ete.
\end{enumerate}
\fincorrection
\correction{005929}
Il est clair que $m_{*}(\emptyset)=0$ et que si $A\subset B\subset
\mathbb{R}$, alors $m_{*}(A)\leq m_{*}(B)$, il faut donc
uniquement d\'emontrer que $m_{*}$ est $\sigma$-sous-additive.

Soit $\{A_n\}_{n\in \mathbb{N}}\subset \mathcal{P}(\mathbb{R})$,
fixons $\varepsilon>0$ et notons $A=\bigcup\limits_{n=1}^\infty
A_n$. Par d\'efinition de l'infimum, pour tout $n\in \mathbb{N}$,
on peut trouver une suite $\{(a_i^n,b_i^n)\}$ telle que
$A_n\subset \bigcup\limits_{i=1}^\infty ]a_i^n,b_i^n[$ et
$$
\sum\limits_{i=1}^\infty (b_i^n-a_i^n)\leq
m_{*}(A_n)+\frac{\varepsilon}{2^n}
$$
Comme  $A\subset
\bigcup\limits_{i,n} ]a_i^n,b_i^n[$, on  a
$$
m_{*}(A)\leq\sum\limits_{n,i} (b_i^n-a_i^n)\leq \sum\limits_{n=1}^\infty\sum\limits_{i=1}^\infty (b_i^n-a_i^n)
\leq
\sum\limits_{n=1}^\infty(m_{*}(A_n)+\frac{\varepsilon}{2^n})=\varepsilon + \sum\limits_{n=1}^\infty
m_{*}(A_n).
$$
On a donc la $\sigma$-sous-additivit\'e\;
$m_{*}(A)\leq\sum\limits_{n=1}^\infty m_{*}(A_n)$ \;puisque
$\varepsilon$ est arbitraire. 
\fincorrection
\correction{005930}
\begin{enumerate}
\item Il est clair que $m_{*}(\emptyset)=0$ et que $m_{*}$ est
monotone.

Soit maintenant $\{A_i\}_{i\in \mathbb{N}}\subset
\mathcal{P}(\Omega)$. Si parmi les $A_i$ il existe au moins un
ensemble $A_j$ non vide, on a
$$m_{*}(\bigcup\limits_i A_i)=1=m_{*}(A_j)\leq \sum\limits_i m_{*}(A_i).$$
Si tous les $A_i$ sont vides, alors $\bigcup\limits_i A_i =
\emptyset$, et donc
$$m_{*}(\bigcup\limits_i A_i)=0=\sum\limits_i
m_{*}(A_i).$$ Ainsi $m_{*}$ est $\sigma$-sous-aditive et par
consequent $m_{*}$ est une mesure ext\'erieure.

\item Les seuls ensembles mesurables sont
$\emptyset$ et $\Omega$, puisque si $A\in \mathcal{P}(\Omega)$ est
tel que $A\neq \emptyset$ et $A\neq \Omega$, alors, pour tout
$Q\in \mathcal{P}(\Omega)$ non vide et non inclus dans $A$, on a
$A\cap Q\neq \emptyset$ et $A^c \cap Q\neq \emptyset$, et donc
$$m_{*}(A\cap Q)+m_{*}(A^c\cap Q)=1+1=2\neq m_{*}(Q)=\left\{
\begin{array}{l}
1 \\
0 \;\;.\\
\end{array}
\right.$$

\item Il est clair que l'ensemble des parties
$m_{*}$-mesurables de $\Omega$,
$\mathcal{M}_{m_{*}}=\{\emptyset,\Omega\}$, est une
$\sigma$-alg\`ebre.

Il est facile de voir aussi que
$$\mu=m_{*}|_{\mathcal{M}_{m_{*}}},\,\mu(\emptyset)=0,\;\;\mu(\Omega)=1,$$
est une mesure sur $(\Omega,\mathcal{M}_{m_{*}})$.

\end{enumerate}
\fincorrection
\correction{005931}
\begin{enumerate}
\item Soit $I = \int_{-\infty}^{+\infty} e^{-x^2}\,dx$. On a~:
$$
I^{2} = \int_{\mathbb{R}^{2}} e^{-(x^2 + y^2)}\, dx dy.
$$
L'application $\Phi~: \mathbb{R}_{+}^{*} \times ]0, 2\pi[
\rightarrow \mathbb{R}^{2} \setminus \{ (x, 0), x \geq 0\}$
d\'efinie par~:
$$
\Phi(r, \theta) = \left(r\cos \theta, r\sin \theta \right)
$$
est un $\mathcal{C}^{1}$-diff\'eomorphisme. De plus
$$
\int_{\mathbb{R}^{2}} e^{-(x^2 + y^2)}\, dx dy =
\int_{\mathbb{R}^{2}\setminus\{ (x, 0), x \geq 0\}} e^{-(x^2 + y^2)}\, dx dy
$$
car l'ensemble $\{ (x, 0), x \geq 0\}$ est n\'egligeable. On en
d\'eduit que~:
$$
I^2 = \int_{r=0}^{+\infty}\int_{\theta=0}^{2\pi} e^{-r^2} r \,dr
d\theta = 2\pi \left[\frac{-e^{-r^2}}{2} \right]_{0}^{+\infty} =
\pi.
$$
Ainsi $I = \int_{-\infty}^{+\infty} e^{-x^2}\,dx  = \sqrt{\pi}$.

 \item
\emph{Calcul de l'aire de la sph\`ere unit\'e de
$\mathbb{R}^{n}$.} Soit $ \mathcal{S}_{n-1} = \{(x_1, \dots, x_n)
\in \mathbb{R}^n ,\, \sum_{i=1}^{n} x_{i}^2 = 1 \} $ la sph\`ere
unit\'e de $\mathbb{R}^{n}$. On note $\mathcal{A}_{n-1}$ son aire.
D'apr\`es la question pr\'ec\'edente, on a~:
$$
\int_{\mathbb{R}^n} e^{-\sum_{i=1}^{n} x_{i}^2} \,dx_{1}\dots
dx_{n} = {\pi}^{\frac{n}{2}}.
$$
D'autre part, puisque l'aire de la sph\`ere de rayon $r$ dans
$\mathbb{R}^n$ vaut $r^{n-1}\mathcal{A}_{n-1}$, il vient~:
$$
\int_{\mathbb{R}^n} e^{-\sum_{i=1}^{n} x_{i}^2} \,dx_{1}\dots
dx_{n} = \mathcal{A}_{n-1} \int_{0}^{+\infty} e^{-r^2}
r^{n-1}\,dr.$$ En posant le changement de variable $x = r^{2}$, on
obtient~:
$$
\int_{\mathbb{R}^n} e^{-\sum_{i=1}^{n} x_{i}^2} \,dx_{1}\dots
dx_{n} = \frac{1}{2}\mathcal{A}_{n-1}\int_{0}^{+\infty} e^{-x}
x^{\frac{n}{2}-1} \,dx,
$$
d'o\`u~:
$$
\mathcal{A}_{n-1} = \frac{2{\pi}^{\frac{n}{2}}}{\Gamma\left(
\frac{n}{2}\right)}.
$$

\item \emph{Calcul du volume de la boule unit\'e de
$\mathbb{R}^{n}$.} Soit $ \mathcal{B}_{n} = \{(x_1, \dots, x_n)
\in \mathbb{R}^n ,\, \sum_{i=1}^{n} x_{i}^2 \leq 1 \} $ la boule
ferm\'ee de rayon 1 dans $\mathbb{R}^n$. On note $\mathcal{V}_{n}$
son volume. On a~:
$$\mathcal{V}_{n} = \int_{0}^{1} r^{n-1}
\mathcal{A}_{n-1} \,dr =  \mathcal{A}_{n-1} \left[\frac{r^{n}}{n}
\right]_{0}^{1} = \frac{\mathcal{A}_{n-1}}{n}.
$$
On en d\'eduit que~:
$$
\mathcal{V}_{n} = \frac{2{\pi}^{\frac{n}{2}}}{n\Gamma\left(
\frac{n}{2}\right)}.
$$
Ce qui se r\'eduit \`a~:
$$
\mathcal{V}_{n} =
\frac{{\pi}^{\frac{n}{2}}}{\Gamma\left(\frac{n}{2} + 1 \right)}
$$
en utilisant l'identit\'e~: $\Gamma(s + 1) = s\Gamma(s)$.

 \item
\emph{Application~:} L'aire de la sph\`ere de rayon $R$ dans
$\mathbb{R}^2$ vaut $$\mathcal{A}_1 R = \frac{2\pi}{\Gamma(1)} R =
2\pi R, $$ qui est bien le p\'erim\`etre du cercle de rayon $R$
dans le plan.

Sachant $\Gamma(\frac 12) = \sqrt \pi$, l'aire de la sph\`ere de rayon $R$ dans $\mathbb{R}^3$ vaut
 $$
 \mathcal{A}_{2} R^2 = \frac{2 {\pi}^{\frac{3}{2}}}{\Gamma\left(\frac{3}{2}
 \right)}
 R^2
 = \frac{2\pi \sqrt{\pi}}{\frac{1}{2}\Gamma\left(\frac{1}{2}\right)} R^2 = 4\pi R^2
$$
qui est bien l'aire de la sph\`ere $S^2$.


  Le volume de la boule de rayon
$R$ dans $\mathbb{R}$ vaut
$$
\mathcal{V}_{1} R = \frac{2{\pi}^{\frac{1}{2}}}{\Gamma\left(
\frac{1}{2}\right)} R = \frac{2 \sqrt{\pi}}{\sqrt{\pi}} R = 2R,
$$
qui est bien la longueur du segment $[-R, R]$.



Le volume de la boule de rayon $R$ dans  $\mathbb{R}^{2}$ vaut
$$
\mathcal{V}_{2} R^{2} = \frac{2\pi}{2\Gamma(1)} R^2 = \pi R^2,
$$
qui est bien l'aire du disque de rayon $R$.



Le volume de la boule de rayon $R$ dans $\mathbb{R}^{3}$ vaut
$$
\mathcal{V}_{3} R^3 = \frac{\mathcal{A}_{2}}{3} R^3 = \frac{4}{3}
\pi R^3.
$$
\end{enumerate}
\fincorrection
\correction{005932}
\begin{enumerate} 
\item On a
\begin{eqnarray*}
\mathcal{V}_n &= &\int_{\mathcal{B}_n} dx_1\dots dx_{n}=
\int_{-1}^{1} dx_1 \int_{\sum_{i=2}^n x_i^2 \leq 1 - x_1^2} dx_2
\dots dx_n\\
& = & \mathcal{V}_{n-1}\int_{-1}^{1}  \left(\sqrt{1-
x_1^2}\right)^{n-1} \,dx_1
\end{eqnarray*}
Posons $x_1=\cos \theta$, pour $\theta \in [0, \pi]$. Alors
$\sqrt{1-x_1^2}=|\sin \theta| = \sin \theta$ et $dx_1 =-\sin\theta
\;d\theta.$ On a donc
$$ \mathcal{V}_n = -\mathcal{V}_{n-1}\int\limits_\pi^0 (\sin \theta)^n d\theta
=\mathcal{V}_{n-1}\int\limits_0^\pi (\sin \theta)^n d\theta = I_n
\cdot \mathcal{V}_{n-1}.$$

\item On a
\begin{eqnarray*}
I_n & = &\int\limits_0^\pi (\sin \theta)^n d\theta=\int\limits_0^\pi (\sin \theta)^{n-1}\; \sin \theta\;d\theta= \\
    & = &\left[-\cos\theta (\sin \theta)^{n-1} \right]_0^\pi + (n-1)\int\limits_0^\pi (\sin \theta)^{n-2}\; (\cos \theta)^2\;d\theta= \\
    & = &(n-1) \int\limits_0^\pi (\sin \theta)^{n-2} (1-(\sin \theta)^2) \;d\theta =(n-1)(I_{n-2}-I_n).\\
\end{eqnarray*}
Donc $I_n=\frac{n-1}{n}\cdot I_{n-2}\;\;$.


\item On a $I_0=\pi,\;I_1 = 2.$ Donc
$I_2=\frac{\pi}{2}, \;
I_3=\frac{4}{3},\;I_4=\frac{3\pi}{8},\;I_5=\frac{16}{15},\;I_6=
\frac{15\pi}{48},\;I_7=\frac{32}{35}.$ 

Comme $\mathcal{V}_1=2$ on trouve:
$$\mathcal{V}_2=\pi,\;\mathcal{V}_3=\frac{4\pi}{3},\;\mathcal{V}_4=\frac{\pi^2}{2},\;
\mathcal{V}_5=\frac{8\pi^2}{15},\;\mathcal{V}_6=\frac{\pi^3}{6},\;\mathcal{V}_7=\frac{16}{105}\pi^3
.$$


\item On a $\displaystyle\mathcal{V}_n =\int\limits_0^1
\int\limits_{\mathcal{S}_{n-1}}
  r^{n-1}dr d\sigma=\frac{1}{n}\mathcal{A}_{n-1},$ \,
d'o\`{u} \, $\mathcal{A}_{n-1}=n\,\mathcal{V}_n$. Donc on a
  $$\mathcal{A}_1=2\pi,\;\mathcal{A}_2=4\pi,\;\mathcal{A}_3=2\pi^2,\;
  \mathcal{A}_4=\frac{8}{3}\pi^2,\;\mathcal{A}_5=\pi^3,\;\mathcal{A}_6=\frac{16}{15}\pi^3
  .$$
\end{enumerate}
\fincorrection
\correction{005969}
Voir le lemme 2.17 p.61 dans \textit{Analysis} de E.~Lieb et M.~Loss, American Mathematical Society (2001).
\fincorrection
\correction{005970}
\begin{enumerate}
\item Soit $E$ un espace de Banach. On suppose qu'il existe une
famille $(O_{i})_{i\in I}$ telle que
\begin{itemize}
\item Pour tout $i\in I$, $O_i$ est un ouvert non vide de
$E$. 
\item $O_i \cap O_j = \emptyset $ si $i\neq j$.
\item $I$ n'est pas d\'enombrable.
\end{itemize}
Supposons que $E$ est s\'eparable. Soit $(u_{n})_{n\in\mathbb{N}}$
une suite dense dans $E$. Gr\^ace \`a (a), pour chaque $i\in I$,
$O_{i} \cap \{u_{n}, n\in\mathbb{N}\} \neq \emptyset$. On choisit
$n(i)$ tel que $u_{n(i)}\in O_i$. On a $n(i) = n(j) \Rightarrow
u_{n(i)} = u_{n(j)} \in O_{i}\cap O_j$ donc $i = j$ par (b). Ainsi
l'application $i \mapsto n(i)$ est injective. Par suite $I$ est
d\'enombrable ce qui contredit (c).

\item Pour tout $a\in \mathbb{R}^{n}$,  on pose $f_{a} =
\mathbf{1}_{\mathcal{B}(a, 1)}$ o\`u $\mathcal{B}(a, 1)$ est la boule de
$\mathbb{R}^n$ de rayon $1$ centr\'ee en $a$. Soit la famille
$$
O_a = \{ f\in L^{\infty}(\mathbb{R}^n), ~ \|f - f_a \|_{\infty} <
\frac{1}{2}\},
$$
o\`u $a$ parcourt les points de $\mathbb{R}^n$. L'ensemble des
points de $\mathbb{R}^n$ n'est pas d\'enombrable, donc (c) est
v\'erifi\'e.  L'ensemble $O_a$ est la boule ouverte de
$L^{\infty}(\mathbb{R}^n)$ de rayon $\frac{1}{2}$ centr\'ee en
$f_{a}$. En particulier (a) est v\'erifi\'e. Remarquons que
lorsque $a\neq b$, on a $\|f_{a} - f_{b}\|_{\infty} = 1$.
Supposons qu'il existe $f \in O_{a}\cap O_b$ avec $a\neq b$. Alors
$$
\|f_{a} - f_{b}\|_{\infty} \leq \|f_a - f\|_{\infty} + \|f -
f_{b}\|_{\infty} < \frac{1}{2} + \frac{1}{2} = 1.
$$
ce qui n'est pas possible. Donc (b) est v\'erifi\'e. On en conclut
que $L^{\infty}(\mathbb{R}^n)$ n'est pas s\'eparable.

\end{enumerate}
\fincorrection
\correction{005971}
cf M.E. Taylor, {\it Measure Theory and Integration}, graduate
studies in mathematics, vol. 76, AMS, 2001, pages 50--51.

\begin{itemize}
\item[$\bullet$] Les ensembles $S_{1l} := \{x\in\Omega, g(x) <
\frac{1}{2} - \frac{1}{l}\}$ et $S_{2l} := \{x\in \Omega, g(x)
> 2 + \frac{1}{l}\}$ sont introduits pour montrer que les
ensembles $\{x\in\Omega, g(x) < \frac{1}{2}\}$ et $\{x\in\Omega,
g(x) > 2\}$ sont de ${\mu}$-mesure nulle (voir plus bas). En
cons\'equence, la fonction $g\in L^2(\Omega, \alpha)$ peut \^etre
choisie telle que $\frac{1}{2} \leq g \leq 2$. (On rappelle que
$L^2(\Omega, \alpha)$ d\'esigne l'ensemble des fonctions de
carr\'e-int\'egrables d\'efinies modulo les ensembles de mesure
nulle.) Cela implique que la fonction $h$ d\'efinie dans la
question~3 est positive comme quotient de deux fonctions
positives.

\item[$\bullet$] Pour montrer que $\mu\left(\{x\in\Omega, g(x) <
\frac{1}{2}\}\right) = 0$, on peut utiliser par exemple la
continuit\'e de la mesure~: on a $S_{11} \subset S_{12} \subset
S_{13} \subset \cdots$ et $\cup_{l\in\mathbb{N}^*} S_{1l} =
\{x\in\Omega, g(x) < \frac{1}{2} \}$, ainsi
$$
\mu\left(\left\{x\in\Omega, g(x) < \frac{1}{2} \right\} \right) =
\mu\left(\cup_{l\in\mathbb{N}^*} S_{1l} \right) =
\lim_{l\rightarrow+\infty} \mu\left(S_{1l}\right) = 0.
$$
De m\^eme, $S_{21} \subset S_{22} \subset S_{23} \subset \cdots$
et $\cup_{l\in\mathbb{N}^*} S_{2l} = \{x\in\Omega, g > 2\}$,
d'o\`u $\mu\left(\{x\in\Omega, g > 2\}\right) = 0.$

\item[$\bullet$] Pour montrer que l'on a l'\'egalit\'e \eqref{radon} du
th\'eor\`eme pour toute fonction positive mesurable, on utilise le
fait que les fonctions essentiellement born\'ees appartiennent \`a
$L^2(\Omega, \alpha)$ (pour une mesure finie on a en effet
$L^{\infty}(\Omega, \alpha) \subset L^2(\Omega, \alpha)$), donc
l'\'egalit\'e
$$
\int_{\Omega} f(2g - 1) \,d\nu =  \int_{\Omega} f(2 - g) \,d\mu.
$$
de la question~2 est en particulier v\'erifi\'ee pour toute
fonction mesurable positive born\'ee. Soit maintenant une fonction
$f$ mesurable positive (non n\'ecessairement born\'ee). Le
th\'eor\`eme de convergence monotone appliqu\'e \`a la suite de
fonctions $f_n = f\,\mathbf{1}_{\{f \leq n\}}$ donne~:
$$
\begin{array}{lcl}
\int_{\Omega} f(2g - 1) \,d\nu & = & \int_{\Omega}
\lim_{n\rightarrow +\infty} f_{n}(2g - 1)\,d\,\nu \\ &=&
\lim_{n\rightarrow +\infty} \int_{\Omega} f_n(2g - 1) \,d\nu =
\lim_{n\rightarrow +\infty} \int_{\Omega} f_n(2 - g) \,d\mu\\ & =
&
\int_{\Omega}\lim_{n\rightarrow +\infty}f_n(2 - g) \,d\mu\\
& = & \int_{\Omega} f(2-g) \,d\mu.
\end{array}
$$
On en d\'eduit que l'\'egalit\'e $(1)$ du th\'eor\`eme est
v\'erifi\'ee pour toute fonction $F$ de la forme $F = f(2g-1)$,
o\`u $f\in\mathcal{M}^+$. Puisque $(2g-1)>0$, l'ensemble des
fonctions $F$ de cette forme est \'egalement $\mathcal{M}^+$.
\end{itemize}
\fincorrection
\correction{005972}
\begin{enumerate}
\item On d\'{e}finit  la fonction B\^{e}ta  par $B(a,b):=\int_{0}^{1}s^{a-1}(1-s)^{b-1}ds$, montrons que \\
$$
B\left(1+\frac{d}{2},\frac{m}{2}\right) = 2\int_{0}^{1}\left(
1-r^{2}\right) ^{d/2}r^{m-1}dr
$$
En utilisant le changement de variable $1-r^2\rightarrow s$, on a
\begin{equation*}
\begin{array}{ll}
 \int_0^1(1-r^2)^{d/2}r^{m-1}dr  & =-\frac{1}{2}\int_1^0
s^{d/2}(1-s)^{\frac{m-2}{2}}ds \\
 & = \; \frac{1}{2}\int_0^1
s^{d/2}(1-s)^{\frac{m}{2}-1}ds = \frac{1}{2}B\left(1+\frac{d}{2}, \frac{m}{2}\right).\\
\end{array}
\end{equation*}

\item Par le changement de variables $t\rightarrow t^2 $ et $
u\rightarrow u^2$ on a
\begin{equation*}\begin{array}{ll} \Gamma(a)\Gamma(b) & = \left(\int_0^\infty e^{-t}t^{a-1} \,
dt\right)\left(\int_0^\infty e^{-u}u^{b-1} \, du\right)\\
&= 4\left(\int_0^\infty e^{-t^2}t^{2a-1} \,
dt\right)\left(\int_0^\infty e^{-u^2}u^{2b-1} \, du\right)
\hspace{3.5cm} \end{array}\end{equation*}

En utilisant le th\'{e}or\`{e}me de Fubini et l'int\'{e}gration en
polaires on a \begin{equation*}\begin{array}{ll}
\Gamma(a)\Gamma(b)&= 4\int_0^\infty \int_0^\infty
e^{-(t^2+u^2)}t^{2a-1}u^{2b-1} \,
dt\, du \\
&= 4\left(\int_0^\infty e^{-r^2}r^{2a-1}r^{2b-1}r \,
dr\right)\left(\int_0^{\frac{\pi}{2}} (\cos \varphi)^{2a-1}(\sin
\varphi)^{2b-1} \, d\varphi\right). \end{array}\end{equation*}

\bigskip

Or, par le changement de variable $r^2\rightarrow r,$
$$2\int_0^\infty e^{-r^2}r^{2(a+b)-1} \,dr=\int_0^\infty e^{-r}r^{a+b-1} \,dr=\Gamma(a+b);$$
et par le changement de variable $u=\cos^2\varphi,$
$$
2\int_0^{\frac{\pi}{2}} (\cos \varphi)^{2a-1}(\sin \varphi)^{2b-1}
\, d\varphi=\int_0^1 u^{a-1}(1-u)^{b-1}\,du=B(a,b).
$$
 Les trois derni\`{e}res
identit\'{e}s entra\^{i}nent
$$\Gamma(a)\Gamma(b)=\Gamma(a+b)\cdot B(a,b).$$

\item On a~:
\begin{eqnarray*}
\int_{\mathbb{R}^{n}} \frac{1}{\left( 1 +
|x|^2\right)^{\alpha}}\,dx & = & \int_{0}^{+\infty} \mu\left(
\left(1 + |x|^{2}\right)^{-\alpha} > t \right)\,dt = \int_{0}^{1}
\text{Vol}\left(\mathcal{B}\left(0, \left(t^{-\frac{1}{\alpha}} -
1 \right)^{\frac{1}{2}}\right)
\right)\,dt\\
& = & \mathcal{V}_{n} \int_{0}^{1}\left(t^{-\frac{1}{\alpha}} - 1
\right)^{\frac{n}{2}}\,dt  =  \mathcal{V}_{n} \int_{0}^{1} \left(1
- t^{\frac{1}{\alpha}} \right)^{\frac{n}{2}}
t^{-\frac{n}{2\alpha}}\,dt
\\
& = & \alpha \mathcal{V}_{n} \int_{0}^{1} \left(1 -
s\right)^{\frac{n}{2}} s^{\alpha -\frac{n}{2} - 1 }\,ds  = \alpha
\mathcal{V}_{n} B\left(\alpha -\frac{n}{2}, \frac{n}{2}+1\right).
\end{eqnarray*}

\end{enumerate}
\fincorrection
\correction{005973}
\begin{enumerate}
\item Posons $\Omega = \mathbb{R}^n\setminus\{
x\in\mathbb{R}^n~|~x_{n} = 0~\text{et}~x_{n-1} \geq 0~\}.$ Comme
$0 < \theta_{n-1} < 2\pi$, l'image de $\Omega'$ par $S$ est
incluse dans $\Omega$. R\'eciproquement, soit $x$ un \'el\'ement
de $\Omega$. Posons $r = |x|$, alors pour tout $i \in\{1, \dots,
n-2\}$, on peut d\'efinir par r\'ecurrence $\theta_i\in(0, \pi)$
gr\^ace \`a son cosinus~:
$$
\cos\theta_{i} = \frac{x_{i}}{r
\sin\theta_1\dots\sin\theta_{i-1}}.
$$
Quant \`a $\theta_{n-1}$, il est d\'etermin\'e par son sinus et
son cosinus. Comme $x_{n} \neq 0$ ou $x_{n-1} < 0$,
n\'ecessairement $\theta_{n-1} \neq 0(\text{modulo}~2\pi)$.

L'application $S$ est contin\^ument diff\'erentiable, car chacune
de ses composantes l'est. La matrice jacobienne a ses vecteurs
colonnes orthogonaux, et de norme respectivement $1$, $r$,
$r\sin\theta_{1}$, $\dots$,
$r\sin\theta_{1}\dots\sin\theta_{n-2}$. Son d\'eterminant vaut
alors $r^{n-1}\left(\sin \theta_{1}\right)^{n-2}\dots
\sin\theta_{n-2}.$ Comme ce d\'eterminant ne s'annule jamais, $S$
est un diff\'eomorphisme de $\Omega'$ sur $\Omega$. 

\item C'est la formule du changement de variable. 

\item On a~:
$$
\begin{array}{lcl} \mathcal{V}_{4} &=& \int_{r = 0}^{1}\int_{\theta_1 =
0}^{\pi}\int_{\theta_{2} = 0}^{\pi} \int_{\theta_3 = 0}^{2\pi}
r^3\sin^{2}\theta_{1}\sin\theta_{2}\,dr
d\theta_1 d\theta_2 d\theta_3\\
& = &
2\pi\left[\frac{r^4}{4}\right]_{0}^{1}\left(\int_{0}^{\pi}\sin^{2}\theta_1
d\theta_1\right)\left(\int_{0}^{\pi}\sin\theta_2
d\theta_2\right)\\
& = & \frac{\pi}{2}\left(\int_{0}^{\pi}\frac{1 -\cos 2\theta_1}{2}
d\theta_1\right)\left[-\cos\theta_2\right]_{0}^{\pi}\\
& = & \frac{\pi^2}{2}.\\
 \mathcal{A}_{3} &=& \int_{\theta_1 =0}^{\pi}\int_{\theta_{2} = 0}^{\pi} \int_{\theta_3 = 0}^{2\pi}
\sin^{2}\theta_{1}\sin\theta_{2}\,
d\theta_1 d\theta_2 d\theta_3\\
& = & 2\pi^2.
\end{array}
$$
\end{enumerate}
\fincorrection
\correction{005974}
Soit $g$ une fonction sur $\mathbb{R}^+$ et
$f~:\mathbb{R}^3\rightarrow \mathbb{R}$ telle que $f(x) = g(|x|)$.
\begin{enumerate}
\item Posons
$$
I = \int_{\mathbb{R}^3} \frac{g(|y|)}{|x - y|}\,dy,
$$
et $r = |x|$, $s = |y|$. Alors $|x - y| = \sqrt{r^2 + s^2 -2 rs
\cos\theta}$ o\`u $\theta$ est l'angle entre l'axe $(Ox)$ et l'axe
$(Oy)$. On consid\'ere les coordonn\'ees sph\'eriques de centre
$O$ et d'axe $(Ox)$ suivantes~:
$$
\begin{array}{lcl}
y_1 &=& s \cos\theta\\
y_2 & = & s \sin\theta \cos\varphi\\
y_3 & = & s\sin\theta \sin\varphi
\end{array}
$$
On a~
$$
I =
\int_{s=0}^{+\infty}\int_{\theta=0}^{\pi}\int_{\varphi=0}^{2\pi}
\frac{g(s)}{\sqrt{r^2 + s^2 - 2rs\cos\theta}} s^2 \sin\theta\,ds
d\theta d\varphi.
$$
On note que
$$
\frac{\sin\theta}{\sqrt{r^2 + s^2 - 2rs\cos\theta}} =
\frac{d}{d\theta}\frac{1}{rs} \sqrt{r^2 + s^2 - 2rs\cos\theta}.
$$
Ainsi~:
$$
\begin{array}{lcl}
I &=& 2\pi \int_{s=0}^{+\infty}\left[\frac{1}{rs} \sqrt{r^2 + s^2
-
2rs\cos\theta} \right]_{\theta = 0}^{\pi} g(s)s^2 \,ds\\
& = & 2\pi \int_{s=0}^{+\infty} \frac{1}{rs} \left(\sqrt{(r+s)^2}
- \sqrt{(r-s)^2} \right) g(s) s^2 \,ds.
\end{array}
$$
Lorsque $s\leq r$, on a $$\sqrt{(r+s)^2} - \sqrt{(r-s)^2} =
(r+s)-(r-s) = 2s,$$ et lorsque $s>r$, il vient $$\sqrt{(r+s)^2} -
\sqrt{(r-s)^2} = (r+s)-(s-r) = 2r.$$ On en d\'eduit alors~:
$$
I = \frac{4\pi}{r}\int_{0}^{r} g(s) s^2\,ds +
4\pi\int_{r}^{+\infty} g(s) s\,ds.
$$

 \item Lorsque $g$ est \`a
support dans $[0, R]$, le potentiel newtonien cr\'e\'e  par la
distribution de masse $f(y) = g(|y|)$ en un point
$x\in\mathbb{R}^3$ tel que $|x|>R$, est identique au potentiel
cr\'e\'e par une masse totale \'egale concentr\'ee \`a l'origine.
\end{enumerate}
\fincorrection
\correction{005975}
Soit $x\in\mathbb{R}^{d}$, $d = 1, 2$ et $r = |x|$. On consid\`ere
$f~: \mathbb{R}^d\rightarrow \mathbb{R}$ donn\'ee par $$f(x) =
h(r) = r^2 (1+r^2)^{-2}.$$
\begin{enumerate} 
\item
La fonction $h$ atteint son maximum en $r = 1$ et $h(1) =
\frac{1}{4}.$ Pour un r\'eel positif $t\leq\frac{1}{4}$ donn\'e,
on cherche \`a r\'esoudre $t = h(r) = r^2 (1+r^2)^{-2}.$ On
obtient deux solutions
$$
\begin{array}{lcl}
r_+ &=& \left(\frac{1-2t}{2t} + \frac{\sqrt{1-4t}}{2t}
\right)^{\frac{1}{2}}\\
r_{-} & = & \left(\frac{1-2t}{2t} - \frac{\sqrt{1-4t}}{2t}
\right)^{\frac{1}{2}}
\end{array}
$$
Ainsi $\mu\left(f > t\right) = \mathcal{V}_{d}\left(r_{+}^d -
r_-^d\right)$. De plus, par d\'efinition, $f^*$ v\'erifie
$\mu\left(f^*>t\right) = \mu\left(f > t\right)$ et
$\mu\left(f^*>t\right) = \mathcal{V}_{d}\, r^{d}$ o\`u $r$ et $t$
sont li\'es par $t = f^*(r)$. Pour $d = 1$, on a donc $r = r_+ -
r_-$ et $t$ est donn\'e par~:
$$
\begin{array}{lcl} r^2 &=& r_+^2 + r_-^2 - 2r_+ r_- = \frac{1-2t}{t} -
2\sqrt{\frac{(1-2t)^2}{4t^2} - \frac{1-4t}{4t^2}}\\
& = & \frac{1 - 4t}{t}.
\end{array}
$$
Il en d\'ecoule que $t = f^*(r) = (4 + r^2)^{-1}.$

\item Pour $d = 2$, on a
$$
r^2 = r_+^2 - r_-^2 = \frac{\sqrt{1-4t}}{t},
$$
ce qui implique que
$$
t = f^*(r) = r^{-4}\left(\sqrt{4 + r^4} - 2 \right).
$$

 \item Calculons $\|f\|_{2}^{2}$ pour $d = 1$. On a
 $$
 \begin{array}{lcl}
\|f\|_{2}^{2} & = & \|f^{*}\|_{2}^{2}\\
& = & 2\int_{0}^{+\infty} (4 + r^2)^{-2}\,dr \\
& = & \frac{1}{4} \int_{0}^{+\infty}(1 + s^2)^{-2}\,ds\\ & =
&\frac{1}{8} \int_{\mathbb{R}} (1 + |x|^2)^{-2}\,dx =
\frac{\pi}{16},
\end{array}
 $$
 o\`u la derni\`ere \'egalit\'e d\'ecoule de l'exercice \ref{exo:beta} (question 3.) sur la fonction Bêta, car~:
$$
\int_{\mathbb{R}} (1 + |x|^2)^{-2}\,dx = 2\mathcal{V}_1
B\left(\frac{3}{2}, \frac{3}{2}\right) = 4
\frac{\Gamma(3/2)^2}{\Gamma(3)} = 4\frac{(1/2\Gamma(1/2))^2}{2!} =
\frac{\pi}{2}.
$$

Pour $d = 2$, on a
$$
\|f\|_{2}^2 = \int_{\mathbb{R}^2} f(x)^2\,dx =
\int_{r=0}^{+\infty}\int_{\theta=0}^{2\pi} h(r)^2 \,r dr d\theta =
2\pi \int_{0}^{+\infty} r^5 (1+r^2)^{-4}\,dr = \frac{\pi}{3},
$$
o\`u la derni\`ere \'egalit\'e d\'ecoule de l'exercice \ref{exo:beta} sur la fonction Bêta, car~:
$$
\int_{\mathbb{R}^2} (1 + |x|^2)^{-4}\,dx = 4\mathcal{V}_{6}
B\left(4 - \frac{6}{2}, \frac{6}{2} + 1\right) = 4 \mathcal{V}_{6}
B\left(1, 4\right),
$$
et
$$
\int_{\mathbb{R}^2} (1 + |x|^2)^{-4}\,dx = \int_{r =
0}^{+\infty}\int_{\mathcal{S}^5} (1 + r^{2})^{-4}\, r^5 \,dr
d\sigma = \mathcal{A}_{5} \int_{0}^{+\infty} (1 + r^{2})^{-4}\,
r^5 \,dr
$$
d'o\`u~:
$$
\begin{array}{lcl}
\int_{0}^{+\infty} (1 + r^{2})^{-4}\, r^5 \,dr  & =&  4
\frac{\mathcal{V}_{6}}{\mathcal{A}_5} B\left(1, 4\right) =
\frac{4}{6} \frac{\Gamma(1)\Gamma(4)}{\Gamma(5)},
 =
\frac{2}{3}\frac{3!}{4!} = \frac{1}{6}
\end{array}
$$
\end{enumerate}
\fincorrection
\correction{005976}
Soit $f~:\mathbb{R}^d \rightarrow \mathbb{R}$ la fonction donn\'ee
par $f(x) = e^{-x^2 + ax}$, o\`u $a\in\mathbb{R}$. Par
translation, le r\'earrangement \`a sym\'etrie sph\'erique
d\'ecroissant $f^*$ de $f$ est donn\'e par $$f^*(x) =
e^{\frac{a^2}{4}}  e^{-x^2}.$$
\fincorrection
\correction{005979}
Soit $1 \leq p < +\infty$.
\begin{enumerate}
\item Si $f$ est continue \`a support compact dans la boule
$\mathcal{B}(0, M)$ centr\'ee en $0$ et de rayon $M$, et si
$|h|\leq 1$, alors
$$
\begin{array}{lcl}
 |f(x - h) - f(x)|^p & \leq & \left(|f(x - h)| + |f(x)|\right)^p
   \leq  \left(2\|f\|_{\infty}\mathbf{1}_{\mathcal{B}(0, M+1)}\right)^p =  \mathbf{1}_{B(0, M+1)}
2^p \|f\|_{\infty}^p.
\end{array}
$$
o\`u $\mathcal{B}(0, M+1)$ est la boule centr\'ee en $0$ de rayon
$M+1$.

\item Pour $f$ continue, on a $\lim_{h \rightarrow 0} |f(x- h) -
f(x)| = 0.$ Puisque la fonction $g(x) = 2^p
\|f\|_{\infty}^p\,\mathbf{1}_{B(0, M+1)}(x) $ appartient \`a
$L^1(\mathbb{R}^{n})$, le th\'eor\`eme de convergence domin\'ee
permet d'intervertir limite et int\'egrale, et il vient~:
$$
\lim_{h \rightarrow 0}\|\tau_{h}f - f\|_{p}^p = \lim_{h
\rightarrow 0}\int_{\mathbb{R}^n} |f(x - h) - f(x)|^{p}\,dx =
\int_{\mathbb{R}^n} \lim_{h \rightarrow 0} |f(x - h) -
f(x)|^{p}\,dx = 0.
$$

\item Soit $f$ une fonction quelconque dans $L^p(\mathbb{R}^n)$,
$1\leq p <+\infty$. Par densit\'e des fonctions continues \`a
support compact dans $L^{p}(\mathbb{R}^n)$, pour tout $\varepsilon
> 0$, il existe $f_{\varepsilon}$ continue \`a support compact
telle que $\|f - f_{\varepsilon}\|_{p} \leq
\frac{\varepsilon}{3}$. Ainsi~:
$$
\begin{array}{lcl}
\|\tau_{h}f - f\|_{p} & = & \|\tau_{h}(f - f_{\varepsilon}) - (f -
f_{\varepsilon}) + \tau_{h}f_{\varepsilon} - f_{\varepsilon}\|_p\\
& \leq & \|\tau_{h}(f - f_{\varepsilon})\|_p + \|f -
f_{\varepsilon}\|_p + \|\tau_{h}f_{\varepsilon} -
f_{\varepsilon}\|_p\\
& = & 2 \|f - f_{\varepsilon}\|_p + \|\tau_{h}f_{\varepsilon} -
f_{\varepsilon}\|_p\\
& = & \frac{2}{3}\varepsilon + \|\tau_{h}f_{\varepsilon} -
f_{\varepsilon}\|_p.
\end{array}
$$
Puisque $f_{\varepsilon}$ est continue \`a support compact,
d'apr\`es la question pr\'ec\'edente, il existe $\delta>0$ tel que
pour $|h|<\delta$,  $\|\tau_{h}f_{\varepsilon} -
f_{\varepsilon}\|_p < \frac{\varepsilon}{3}$. Ainsi, pour $|h| <
\delta$, on a $\|\tau_{h}f - f\|_{p}< \varepsilon$. En d'autre
termes $\lim_{h \rightarrow 0}\|\tau_{h}f - f\|_{p}  = 0$.


\item Pour $p = \infty$, les fonctions continues \`a support
compact ne sont pas denses dans $L^{\infty}(\mathbb{R}^n)$ ce qui
fait que la d\'emonstration pr\'ec\'edente ne peut pas s'appliquer
dans ce cas. De plus, on v\'erifie que, pour $f = \mathbf{1}_{B(0, 1)}$ et
$h \neq 0$, on a
$$
\|\tau_{h}f - f\|_{\infty} = 1.
$$
Alors que pour $h = 0$, on a $\|\tau_{h}f - f\|_{\infty} = 0. $ On
peut \'egalement v\'erifier que $\lim_{h\rightarrow 0}\|\tau_{h}f
- f \|_{\infty} = 0$ si et seulement si la fonction $f$ poss\`ede
un repr\'esentant uniform\'ement continu.
\end{enumerate}
\fincorrection
\correction{005980}
Soit $\{\varphi_{n}\}_{n\in\mathbb{N}}$ une suite de fonctions
v\'erifiant les hypoth\`eses (i), (ii) et (iii) du
th\'eor\`eme, et soit $1\leq p <+\infty$.
\begin{enumerate}
\item En notant $q$ l'exposant conjugu\'e de $p$ ($\frac{1}{p} +
\frac{1}{q} = 1$), on a
$$
\begin{array}{lcl}
|\varphi_n * f - f|^p(x)& = &|\int_{\mathbb{R}^n} f(x - y)
\varphi_n(y)\,dy - f(x) \int_{\mathbb{R}^n}\varphi_n(y)\,dy|^{p}\\
& & \\ & \leq & \left(\int_{\mathbb{R}^n} |f(x - y) - f(x)|
|\varphi_n(y)|\,dy \right)^p.
\end{array}
$$
En utilisant l'in\'egalit\'e de H\"older pour la mesure $d\nu(x) =
|\varphi_n|(x)\,dx$, on a
$$
\begin{array}{lcl}
|\varphi_n * f - f|^p(x) & \leq & \left(\int_{\mathbb{R}^{n}}
1^{q}\,d\nu(y) \right)^{\frac{p}{q}}\left(
\int_{\mathbb{R}^{n}}|f(x - y) -
f(x)|^p\,d\nu(y)\right)^{\frac{p}{p}}\\
& \leq & \left(\int_{\mathbb{R}^n} |\varphi_n|(y)\,dy
\right)^{\frac{p}{q}} \left(\int_{\mathbb{R}^n} |f(x - y) -
f(x)|^p |\varphi_n|(y)\,dy\right)\\ & \leq &
K^{\frac{p}{q}}\left(\int_{\mathbb{R}^n} |f(x - y) - f(x)|^p
|\varphi_n|(y)\,dy\right).
\end{array}
$$

\item On en d\'eduit que
$$
\begin{array}{lcl}
\|\varphi_n * f - f\|_{p}^{p} & \leq &
K^{\frac{p}{q}}\int_{x\in\mathbb{R}^n}\left(\int_{y\in\mathbb{R}^n}|f(x
- y) - f(x)|^p |\varphi_n|(y)\,dy\right)\,dx
\end{array}
$$
D'apr\`es le th\'eor\`eme de Tonelli~:
$$
\begin{array}{lcl}
\|\varphi_n * f - f\|_{p}^{p} & \leq &
K^{\frac{p}{q}}\int_{y\in\mathbb{R}^n}\left(\int_{x\in\mathbb{R}^n}|f(x
- y) - f(x)|^p \,dx\right)|\varphi_n|(y)\,dy\\ & \leq &
 K^{\frac{p}{q}}
\int_{\mathbb{R}^{n}} \|\tau_{y}f - f\|_{p}^p |\varphi_n(y)|\,dy.
\end{array}
$$

\item Soit $\delta>0$, on a
$$
\begin{array}{lcl}
 && \|\varphi_n * f - f\|_{p}^{p}  \\ & \leq &  K^{\frac{p}{q}}\left(
\int_{|y|\leq \delta}\|\tau_{y}f - f\|_{p}^p |\varphi_n(y)|\,dy +
\int_{|y|>\delta} \|\tau_{y}f - f\|_{p}^p |\varphi_n(y)|\,dy
\right)\\ & \leq &
K^{\frac{p}{q}}\left(\sup_{|y|\leq\delta}\|\tau_{y}f - f\|_{p}^p
\int_{|y|\leq \delta} |\varphi_n(y)|\,dy + \int_{|y|>\delta}
\left(\|\tau_{y}f\|_p + \|f\|_{p}\right)^p |\varphi_n(y)|\,dy
\right)\\& \leq &
K^{\frac{p}{q}}\left(K\,\sup_{|y|\leq\delta}\|\tau_{y}f -
f\|_{p}^p
 + \left(2\|f\|_{p}\right)^p \int_{|y|>\delta}
|\varphi_n(y)|\,dy \right)\\
& \leq & K^{\frac{p}{q}}\left(K\,\sup_{|y|\leq\delta}\|\tau_{y}f -
f\|_{p}^p + 2^p\|f\|_{p}^{p}\int_{|y|>\delta} |\varphi_{n}(y)|\,dy
\right).
\end{array}
$$



\item Soit $\varepsilon>0$. Par continuit\'e des translations dans
$L^p(\mathbb{R}^n)$ (cf l'exercice pr\'ec\'edent), il existe un
$\delta>0$ tel que
$$
|y|\leq\delta\quad \Rightarrow \|\tau_{y}f - f\|_{p}^p <
\frac{K^{-\left(\frac{p}{q}+1\right)}}{2}~ \varepsilon.
$$
D'apr\`es l'hypoth\`ese (iii), il existe un $N\in\mathbb{N}$ tel
que pour $n>N$, on a
$$
\int_{|y|>\delta} |\varphi_{n}(y)|\,dy <
\frac{K^{-\frac{p}{q}}}{2^{p+1}\|f\|_{p}^{p}}~ \varepsilon.
$$
Ainsi pour tout $n>N$,
$$
\|\varphi_n * f - f\|_{p}^{p} < \varepsilon,
$$
i.e. $\lim_{n\rightarrow+\infty}\|\varphi_n * f - f\|_{p} = 0$.
\end{enumerate}
\fincorrection
\correction{005981}
Voir E.~Lieb et M.~Loss, \emph{Analysis}, p.123, American Mathematical Society (2001).
\fincorrection
\nocorrection
\nocorrection
\correction{002340}
\begin{enumerate}
  \item $A$ une partie non vide de $\Rr$, un \emph{majorant} de $A$ est un 
réel $M\in \Rr$ tel que
$$ \forall x \in A \qquad x \leq M.$$
Si $A$ est un partie non vide et majorée, alors par définition $\sup A$ est le plus petit des majorants.
On a les propriétés suivantes :
  \begin{enumerate}
  \item $\sup(A+B)=\sup A+\sup B$ ;
  \item $\sup(A \cup B)=\max(\sup A,\sup B)$ ;
  \item $\max(\inf A, \inf B) \leq \sup(A\cap B) \leq \min(\sup A, \sup B)$ si $A\cap B\neq \varnothing$ ;
  \item $\inf(A \cup B)=\min(\inf A, \inf B)$ ;
  \item $\max(\inf A, \inf B) \leq \inf(A\cap B) \leq \min(\sup A, \sup B)$ si $A\cap B\neq \varnothing$ ;
  \end{enumerate}

Prouvons les deux premières égalités,  
\begin{enumerate}
  \item $\sup(A+B)=\sup A+\sup B$ : pour tout $a\in A$ et $b\in B$ on a $a\leq \sup A$ et $b \leq \sup B$ donc
$a+b \leq \sup A+\sup B$, donc $\sup A+\sup B$ est un majorant de $A+B$ et comme $\sup(A+B)$ est le plus petit des majorants de $A+B$ alors $\sup(A+B) \leq \sup A+\sup B$. Réciproquement, il existe une suite $(a_n)$ d'éléments de $A$ tel que cette suite converge vers $\sup A$, de même il existe une suite $(b_n)$ d'éléments de $B$ qui converge vers $\sup B$, la suite $(a_n+b_n)$ est une suite d'éléments de $A+B$
qui converge vers $\sup A+\sup B$, donc la borne supérieure de $A+B$ est plus grande que $\sup A+\sup B$,
soit $\sup(A+B) \geq \sup A+\sup B$. D'o\`u l'égalité.

  \item $\sup(A \cup B)=\max(\sup A,\sup B)$ :  
Remarquons d'abord que si $P\subset Q$ alors $\sup P \leq \sup Q$ : en effet $\sup Q$ est un majorant de $Q$ donc de $P$ (par l'inclusion $P\subset Q$), donc le plus petit des majorants, $\sup P$, pour $P$ est plus petit
que le majorant particulier $\sup Q$.
Appliquons ceci à $A \subset A\cup B$ donc $\sup A \leq \sup(A \cup B)$ et pour $B \subset A\cup B$ 
on obtient $ \sup B \leq \sup(A \cup B)$. On vient de prouver $\sup(A \cup B) \geq \max(\sup A,\sup B)$.
Pour l'autre inégalité : soit $M = \max(\sup A,\sup B)$. Pour  $x\in A\cup B$ alors soit $x\in A$ et alors 
$x \leq \sup A \leq M$, ou soit $x\in B$ et alors $x \leq \sup B \leq M$ ; donc quelque soit $x \in A\cup B$,
$x\leq M$ donc $M$ est un majorant de $A\cup B$, donc $\sup(A\cup B) \leq M=\max(\sup A,\sup B$).
\end{enumerate}
 
  \item 
\begin{enumerate}
  \item $d(0,\Rr\setminus \Qq) = 0$, regarder des éléments du type $\frac{\sqrt 2}{n}$, pour $n\in \Nn^*$.
  \item $d(\sqrt 2,\Qq)=0$, c'est la densité de $\Qq$ dans $\Rr$ ou alors regarder la suite définie par 
$u_0=1, u_{n+1} = \frac12 (u_n+\frac2{u_n}), n \in \Nn$, qui est une suite de rationnels convergeant vers $\sqrt 2$.
  \item On suppose que $\mathcal{D}$ passe par l'origine, alors $d(M, \mathcal{D}) = x^2+y^2+z^2-(ax+by+cz)^2$.
  \end{enumerate}


  \item $d(A,B)=0$.
  \item $\mathrm{diam} (]0,1[ \cap \Qq) = 1 = \\mathrm{diam} ([0,1] \cap (\Rr\setminus \Qq))$.
\end{enumerate}
\fincorrection
\correction{002341}
\begin{enumerate}
  \item $J_x$ est un ouvert non vide car c'est une union d'ouverts contenant $x$. De plus $J_x$ est un intervalle car c'est une union d'intervalles contenant tous le point $x$. Donc $J_x$ est un intervalle ouvert. 
On peut donc \'ecrire $\mathcal{O} = \cup_{x\in \mathcal{O}} J_x$. Mais cette union n'est pas n\'ecessairement d\'enombrable.

Tout d'abord si $z\in J_x$ alors $J_x=J_z$. En effet soit $I$ un intervalle inclus dans $\mathcal{O}$ contenant $x$ et $z$. Si $x'\in J_x$, soit $J$ un intervalle inclus dans $\mathcal{O}$ contenant $x$ et $x'$. Alors $I\cup J$ est un intervalle (car $x$ est dans les deux intervalles $I$ et $J$), $I\cup J$ est inclus dans $\mathcal{O}$ et contient $x'$ et $z$. Donc $x'\in J_z$. Donc $J_x \subset J_z$. Enfin comme $z\in J_x$ on a aussi $x\in J_z$, donc on montrerait de même $J_z\subset J_x$. Donc $J_x=J_z$.

Pour $x,y \in \mathcal{O}$ alors $J_x=J_y$ ou $J_x \cap J_y = \varnothing$. En effet supposons que $J_x\cap J_y \neq \varnothing$ et soit $z \in J_x \cap J_y$. Comme $z\in J_x$ alors $J_x=J_z$, comme 
 $z\in J_y$ alors $J_y=J_z$. Donc $J_x=J_y$.

Pour chaque intervalle ouvert $J_x$ il existe $q \in \Qq \cap J_x$, avec bien s\^ur $J_x = J_q$.
 Comme $\Qq$ est d\'enombrable $\mathcal{O}\cap \Qq$ l'est aussi. On a ainsi \'ecrit
$$\mathcal{O} = \bigcup_{q\in \mathcal{O}\cap \Qq} J_{q},$$
ce qui \'etait demand\'e.
  \item  Pour $\Rr^n$ on peut montrer le r\'esultat suivant : tout ouvert $\mathcal{O}$ de $\Rr^n$ s'\'ecrit comme l'union d\'enombrable de boules ouverte. On consid\'ere $J_x$ l'union des boules  ouvertes de rayon 
rationnel centr\'ees en $x$,  ensuite on regarde seulement les $x$ appartenant \`a $\mathcal{O} \cap \Qq^n$. Par contrer on autorise deux boules à s'intersecter.
\end{enumerate}
\fincorrection
\correction{002342}
\begin{enumerate}
  \item Soient $d = p+q \sqrt 2$ et $d'=p'+q'\sqrt 2$ deux éléments de $D$.
Alors $d+d'=(p+p')+(q+q')\sqrt 2$ est un élément de $D$ et $dd'=
(pp'+2qq')+(pq'+p'q)\sqrt 2$ aussi.
  
  \item On a $u < 1$ donc $u^k$ tend vers $0$ quand $k$ tend vers $+\infty$.
Donc pour $\epsilon = b-a$, il existe $n \in \Nn$ tel que si $k \ge n$ on
a $u^k < \epsilon = b-a$. En particulier $u^n < b-a$. Si on cherchait un réel alors $r = \frac a {u^n} + 1$ conviendrait, mais on cherche un entier, posons $m = E(\frac a {u^n})+1$. Alors $m-1 \le \frac{a}{u^n} < m$. L'inégalité de droite donne $a < mu^n$. L'inégalité de gauche s'écrit aussi $mu^n-u^n \le a$ soit
$mu^n \le a+ u^n < a+b-a=b$ donc $a<mu^n<b$.

Déduisons de cela que $D$ est dense dans $\Rr$ : pour tout intervalle $[a,b]$, $a<b$ il existe $m,n$ des entiers tels que $mu^n \in [a,b]$. Or $mu^n$ est dans $D$ car $u\in D$ donc par multiplication $u^n \in D$.

\end{enumerate}
\fincorrection
\correction{002343}
\begin{enumerate}
  \item Cette exercice justifie la terminologie ``boule fermée''. Il s'agit de montrer que le complémentaire d'une boule fermée est un ensemble ouvert. Il est vivement conseillé de faire un dessin. Soit $C = E \setminus B'(a,r)$. Soit $x \in C$, on cherche une boule ouverte $B(x,\epsilon)$ contenue dans $C$. Comme $x\in C$, $x \notin B'(a,r)$ donc
$d(a,x) > r$. Soit $\epsilon$ tel que  $0 < \epsilon < d(a,x)-r$. Montrons que
$B(x,\epsilon) \subset C$ : pour $y \in B(x,\epsilon)$, l'inégalité triangulaire $d(a,x) \le d(a,y)+d(y,x)$ donc $d(a,y) \ge d(a,x)-d(y,x) \ge d(a,x)-\epsilon >r$. Comme $d(a,y) >r$ alors $y \notin B'(a,r)$ donc $y \in C$. Comme la preuve est valable quelque soit $y \in B(x,\epsilon)$, donc $B(x,\epsilon) \subset C$.
Et donc $C$ est un ouvert.
  
  \item Pour $a=(\frac 12,0)$ et $r=\frac 12$ on a 
$B'(a,r) = [0,1]\times\{0\}\cup\{0\}\times[0,\frac 12]$,
$B(a,r) = ]0,1[ \times \{0\}$ et $\overline{B(a,r)} = [0,1]\times\{0\}$.

\end{enumerate}
\fincorrection
\correction{002344}
\begin{enumerate}
  \item On note $B= B(a,r)$, $B'=B'(a,r)$, $\bar B = \overline{B(a,r)}$.
Il faut montrer $B'=\bar B$. $B'$ est une boule ferm\'ee, donc un ferm\'e contenant $B$, alors que $\bar B$ est le plus petit ferm\'e contenant $B$, donc $\bar B \subset B'$.

\'Etudions l'inclusion inverse: soit $x\in B'$, il faut montrer $x\in \bar B$.
Si $x\in B$ alors $x\in \bar B$, supposons donc que $x\notin B$, alors 
$\| x-a\| = r$. Soit $B(x,\epsilon)$ un boule centr\'ee en $x$. $x$ est adh\'erent \`a $B$ si $B(x,\epsilon) \cap B$ est non vide quelque soit $\epsilon >0$.
Fixons $\epsilon > 0$ et soit le point 
$$ y = x - \frac{\epsilon}{2} \frac{x-a}{\|x-a\|}.$$
Faire un dessin et placer $y$ sur ce dessin.
D'une part $y\in B(x,\epsilon)$ car $\|y-x\| = \epsilon/2 < \epsilon$.
D'autre part $y\in B = B(a,r)$ car $\| y-a \| = \| x-a-\frac{\epsilon}{2}\frac{x-a}{\|x-a\|}\| = \|x-a\| (1-\frac{\epsilon}{2\|x-a\|}) = r -\frac{\epsilon}{2} <r$. Donc $y\in B\cap B(x,\epsilon)$, ce qui prouve que $B'\cap \bar B$. Donc $B' = \bar B$.

 
  \item Pour le sens $\Leftarrow$. Soit $x \in \bar B(a,r)$ alors $\| x-b\|= 
\|x-a+a-b\| \le \|x-a\|+\|a-b\| \le r+R-r\le R$, donc $x \in \bar B(b,R)$.

Pour le sens $\Rightarrow$. Soit 
$$x = a+ r\frac{a-b}{\|a-b\|},$$
alors $\|x-a\| = r$ donc $x\in \bar B(a,r)$, donc $x \in \bar B(b,R)$, donc 
$\|x-b\| \le R$
or $\| x-b\| = \|a-b\|+ r$ (c'est le m\^eme calcul que pour la question pr\'ec\'edente). 
Donc $\|a-b\|+r \le R$, soit $0\le \|a-b\| \le R-r$ et en particulier $r\le R$.
\end{enumerate}
\fincorrection
\correction{002345}
\begin{enumerate}
  \item
   \begin{enumerate}
     \item Si $||(x,y)||=0$ alors $\max(|x+y|,|x-2y|)=0$ donc $x+y=0$ et $x-2y=0$ donc $x=0$ et $y=0$. Réciproquement $\|(0,0)\|=0$. 
     \item $||\lambda.(x,y)||= ||(\lambda x,\lambda y)||=\max(|\lambda x+ \lambda y|,|\lambda x-2\lambda y|) = |\lambda|\max(|x+y|,|x-2y|) = |\lambda|.||(x,y)||$.
     \item $\|(x,y)+(x',y')\| = \|(x+x',y+y')\| =  \max(|x+x'+y+y'|,|x+x'-2y-2y'|) \le \max (|x+y|+|x'+y'|,|x-2y|+|x'-2y'|) \le \max(|x+y|,|x-2y|)+\max(|x'+y'|, |x'-2y'|) \le ||(x,y)||+ ||(x',y')||$.
   \end{enumerate}
  
La boule unité fermée centrée à l'origine est la région du plan comprise entre les droites d'équations $x+y=+1$, $x+y=-1$, $x-2y=+1$, $x-2y=-1$.

  \item Sens $\Leftarrow$ : Si $x \in B_q$ alors $q(x) \le 1$ donc $p(x) \le 1$
donc $x\in B_p$. Sens $\Rightarrow$ : Soit $x\in \Rr^n \setminus \{0\}$ alors $q(\frac{x}{q(x)}) = 1$ donc  $\frac{x}{q(x)}\in B_q $ donc $\frac{x}{q(x)}\in B_p $ donc $p(\frac{x}{q(x)}) \le 1$ soir $p(x) \le q(x)$. Ceci étant aussi valable pour $x=0$.

$B_q \subset 2B_p$ est équivalent à $p(x) \le 2q(x)$ pour tout $x\in \Rr^n$
(attention au sens !). Et $\frac12 B_p \subset B_q$ est équivalent à 
$\frac12q(x) \le p(x)$. Si les deux inclusions sont vraies alors
$\frac12 p \le q \le 2p$ et en particulier les normes $p$ et $q$ sont équivalentes. 

Par exemple dans $\Rr^2$ pour les normes $\|.\|_1$, $\|.\|_2$, $\|.\|_\infty$
On a 
$$B_1 \subset B_2 \subset B_\infty \subset 2B_1 \subset 2B_2 \subset \cdots$$

\end{enumerate}
\fincorrection
\correction{002346}
\begin{enumerate}
  \item Une suite de $l^\infty$ est notée $(x^p)_{p\in \Nn}$, pour chaque $p \ge 0$, $x^p$ est elle m\^eme une suite $x^p = (x^p(0), x^p(1), x^p(2), \ldots)$. 
(Il convient de garder la tête froide : on regarde des suites de suites !)
Il faut montrer que $Y$ est fermé  dans $X$.
Soit donc $(x^p)$ une suite de $Y$ qui converge vers $x \in X$. Il faut donc montrer qu'en fait $x\in Y$, c'est-à-dire que $x= (x(0),x(1),\ldots)$ est une suite tendant vers $0$.
Soit $\epsilon >0$ comme $x^p \rightarrow x$ alors il existe $P$ tel que
si $p \ge P$ on ait $d(x^p,x) < \epsilon$. Par la définition
de $d$ on a pour $p \ge P$ et pour tout $n\in \Nn$, $|x^p(n)-x(n)| < \epsilon$.
Fixons $p=P$, alors $x^P \in Y$ donc $x^P$ est une suite tendant vers $0$,
donc il existe $N$ tel que si $n\ge N$ alors $|x^P(n)|<\epsilon$.
Réunissons tout cela, pour $n\ge N$ :
$$|x(n)|=|x(n)-x^P(n)+x^P(n)| \le | x(n)-x^P(n)|+|x^P(n)| \le 2\epsilon.$$
 Donc la suite $x$ tend vers $0$, donc $x\in Y$ et $Y$ est fermé.

  \item Notons $Z$ l'ensemble des suites nulles à partir d'un certain rang.
Pour $y = (y(0),y(1),y(2),\ldots) \in Y$, définissons la suite 
$y^0 = (y(0),0,0,\ldots)$, $y^1=(y(0),y(1),0,0,\ldots)$,... $y^p=(y(0),\ldots,y(p-1),y(p),0,0,0,\ldots)$. La suite
$(y^p)$ est bien une suite d'éléments de $Z$. De plus
$d(y^p,y) = \sup_{n\in \Nn} |y^p(n)-y(n)| = \sup_{n > p} |y(n)|$
or la suite $y(n)$ tend vers $0$ donc $d(y^p,y)$ tend vers $0$ quand
$p$ tend vers $+\infty$.

On montre facilement (par l'absurde) que l'élément $x = (1,1,1,\ldots )\in X$
n'est limite d'aucune suite d'éléments de $Z$, (ni d'ailleurs de $Y$).
\end{enumerate}
\fincorrection
\correction{002347}
Par l'inégalité triangulaire $|f(x)+f'(x)| \le |f(x)|+|f'(x)|$ on obtient $\| f \| \le N(f)$.
Pour une inégalité dans l'autre sens décomposons le travail :
\begin{itemize}
  \item $\|f'\|_\infty \le \| f \|_\infty + \|f\|$ : en effet
par l'inégalité triangulaire $|f'(x)| \le |f(x)|+|f'(x)+f(x)|$.

  \item $\|f\|_\infty \le \|f\|$ : en effet $f$ est continue sur $[0,1]$
donc elle est bornée et atteint ses bornes. Soit $x_0 \in [0,1]$ ce point du maximum. Si $x_0 \in ]0,1[$ alors $f'(x_0)=0$ donc
$\|f\|_\infty = |f(x_0)| = |f(x_0)+f'(x_0)| \le \|f\|$.
Si $x_0 = 1$ alors $f$ et $f'$ ont m\^eme signe sur un intervalle 
$[1-\epsilon,1]$ donc sur cet intervalle $|f(x)| \le |f(x)+f'(x)|$
et donc $\|f\|_\infty = |f(1)|  \le \|f\|$. (Enfin $f(0)=0$ donc si $x_0=0$
alors $f$ est nulle et l'inégalité est triviale.)

  \item Il reste à rassembler les expressions :
$$ N(f) = \|f'\|_\infty  + \|f\|_\infty \le \|f\|_\infty + \|f\| + \|f\|_\infty \le 3\|f\|.$$
(La première inégalité vient du premier point et la deuxième du second.)
\end{itemize}
Les normes $\| f \|$ et $N(f)$ sont équivalentes :
$$ \frac13 N(f) \le \| f \| \le N(f).$$
\fincorrection
\nocorrection
\correction{002349}
\begin{enumerate}
  \item $\| f \|_1 = \int_0^1|f(t)|dt \le \int_0^1 \|f\|_\infty dt \le \|f\|_\infty$.
Donc $\| f \|_1 \le \|f\|_\infty$ Par contre il n'existe aucune constante $C>0$ tel que $\|f\|_\infty \le C \| f \|_1$ pour tout $f$. Pour montrer ceci par l'absurde, supposons qu'il existe une constante $C>0$ telle que
$\|f\|_\infty \le C \| f \|_1$ pour tout $f$ de $ C([0,1],{\Rr})$.
Regardons les fonctions $f_k$ définies par $f_k(x) = 2k(1-kx)$ si
$x\in[0,\frac1k]$ et $f_k(x)=0$ si $x > \frac1k$. 
Alors $f_k \in C([0,1],{\Rr})$ et  $\| f_k \|_\infty = 2k$ alors que $\| f_k \|_1 = 1$. On obtient $2k \le C.1$ ce qui est contradicoire pour
$k$ assez grand. Cela prouve que les normes ne sont pas équivalentes.
  
  \item Comme les métriques sont définies par des normes et que les normes ne sont pas équivalentes alors les métriques ne définissent pas la même topologie.

\end{enumerate}
\fincorrection
\correction{002350}
\begin{enumerate}
  \item On montre facilement 
$$N_1 \le N_2 \le 2N_1 \le 2N_4 \le 2N_3.$$
\item Par contre il n'existe pas de constante $C>0$ telle que $N_3 \le CN_4$ ou $N_2 \le CN_4$.
On suppose qu'il existe $C>0$ telle que $N_3 \le CN_4$ on
regarde $f_k$ définie par $f_k(x) = x^k$, après calcul on obtient $N_3(f_k)=k+1$ et $N_4(f_k)=2$, pour $k$ suffisament grand on obtient une contradiction.
Comme $N_1$ et $N_2$ sont équivalentes on va prouver qu'il n'existe pas de constante $C>0$ telle que $N_3 \le CN_1$. On prend $g_k$, définie par
$g_k(x) = 1+\sin(2\pi k x)$. Alors $N_1(g_k) = 2$ et $N_3(g_k) = 4k$, ce qui prouve le résultat souhaité.
\end{enumerate}
\fincorrection
\correction{002351}
\begin{enumerate}
  \item 
    \begin{enumerate}
      \item Par exemple une suite constante $x_n=a$ pour tout $n$.  
      \item Par exemple $x_n = \frac1n$ et $a=0$.
      \item Comme $\Qq$ est dénombrable on peut trouver une suite $x_n$ telle
que $A = \{x_1,x_2,\ldots\}= \Qq$. On prend $a = \sqrt2$ alors $a\in \bar A \setminus A = \Rr \setminus \Qq$.
     \end{enumerate}
  \item C'est juste les définitions : un point d'accumulation de $A$ est toujours une valeur d'adhérence de $A$.
\end{enumerate}
\fincorrection
\correction{002352}
\begin{enumerate}
  \item (Correction pour $n=1$, pour $n>1$ remplacer les intervalles par des boules.) Comme $0$ est isolé soit $I =  ]-\epsilon,+\epsilon[$ un voisinage de 
$0$ tel que $I\cap G = \{0\}$. Soit $g\in G$ et considérons $I_g = g+I = ]g-\epsilon,g+\epsilon[$. Supposons, par l'absurde, que $I_g \cap G$ ne soit pas réduit à $g$. Alors
il existe $g' \in I_g\cap G$, $g'\neq g$. Mais $g-\epsilon < g' < g+\epsilon$ et
donc $g-g'\in I$ comme $G$ est un groupe on a $g-g'\in G$ et on a $g-g'\neq 0$. On a donc trouvé un élément $g-g' \in G\cap I$ qui n'est pas $0$. Ce qui est une contradiction.

Pour montrer que $G$ est discret (c'est-à-dire $G$ est dénombrable et ses points sont isolés) on remarque que la distance entre deux éléments de $G$ est au moins $\epsilon$ donc pour $J_g = ]g-\frac\epsilon2, g+\frac\epsilon2[$ on a
$g\neq g'$ implique $J_g \cap J_{g'} = \varnothing$. Pour chaque $g\in G$ on choisit $q(g) \in \Qq \cap J_g$, ce qui donne une application :
$\Phi : G \longrightarrow \Qq$ définie par $\Phi(g) = q(g)$, et $\Phi$ est injective, donc $G$ est dénombrable.

Montrons que $G$ est fermé : soit $(g_n)$ une suite de $G$ qui converge vers $g\in\Rr$. Pour $N$ assez grand et pour tout $n\ge N$ on a
$|g_n-g| \le \frac\epsilon 4$. Pour $n\ge N$ on a
$|g_n-g_{N}| \le |g_n-g|+|g-g_{N}| \le \frac \epsilon4+\frac\epsilon4 \le\frac\epsilon2$. Donc comme $g_N\in J_{g_N}$ alors $g_n \in J_{g_N}$ également, or $J_{g_N}$ ne contient qu'un seul élément de $G$ donc
$g_n = g_N$ pour tout $n\ge N$. La suite est donc stationnaire (i.e. constante à partir d'un certain rang) donc la limite $g$ vaut $g_N$ et en particulier
$g \in G$.
  
  \item  Supposons $G\neq \{0\}$. Soit $a = \inf G\cap \Rr_+^*$. Comme $0$ est isolé alors $a>0$.
Comme $G$ est fermé alors $a\in G$. Soit $g\in G$. Soit $k = E(\frac g a)$ alors $k \le  \ \frac g a < k+1$. Donc $0 \le g-ka < a$. Or $g-ka$ est dans $G$ et dans $\Rr_+$, comme il est plus petit que $a = \inf G\cap \Rr_+^*$ alors nécessairement $g-ka=0$, soit $g=ka \in a\Zz$.


  \item Soit $x\in \Rr$ et $\epsilon >0$, on cherche $g\in G \cap ]x-\epsilon,x+\epsilon[$. Comme $0$ est un point d'accumulation de $G$ il existe $h\in G$ tel que $0<h<\epsilon$ pour $k = E(\frac xh)$, on a $kh \le x < kh+h$,
donc $g=kh \in G \cap ]x-\epsilon,x+\epsilon[$. Donc $G$ est dense dans $\Rr$.


Pour un groupe $G$ quelconque soit $0$  est isolé, soit $0$ est un point d'accumulation. Si en plus $G$ est fermé alors soit $G=a\Zz$ ou $G= \{0\}$,
soit $\bar G= \Rr$ donc $G = \Rr$. Les sous-groupes fermés de $(\Rr,+)$ sont
donc $0$, $\Rr$ et les $a\Zz$ avec $a>0$.

  \item Soit $G = \Zz+\alpha\Zz$, c'est un sous-groupe de $(\Rr,+)$. Si $G$ n'est pas dense dans $\Rr$ alors, par les questions précédentes, il existe $a>0$ tel que $G=a\Zz$. En particulier $1\in G$ donc il existe $k\in \Zz$ tel que
$1 = ka$ de même $\alpha \in G$ donc il existe $k' \in \Zz$ tel que $\alpha = k'a$. Par division $\alpha  = \frac {k'}k$. Ce qui contredit $\alpha \notin \Qq$. Donc $G = \Zz+\alpha\Zz$ est dense dans $\Rr$.

Définissons $\Phi : \Rr \longrightarrow S^1$ par $t\mapsto e^{2i\pi t}$
($S^1$ est le cercle de $\Cc$ des nombre complexes de module $1$).
Alors $\Phi$ est continue et surjective. Comme $\Phi$ est continue alors pour tout ensemble $A\subset \Rr$ on $\Phi(\bar A)\subset \overline{\Phi(A)}$.
Appliqu\'e à l'ensemble $G = \Zz+\alpha\Zz$, on a $\bar G = \Rr$  donc
$\Phi(\bar G)= S^1$ car $\Phi$ est surjective ; d'autre part 
$\Phi(G) = \{ e^{2i\pi k \alpha} \mid k\in \Zz \}$. Donc $S^1=\Phi(\bar G)   \subset \overline {\Phi(G)} = \overline {\{ e^{2i\pi k \alpha} \}}$.
L'adhérence de $\{ e^{2i\pi k \alpha} \mid k \in \Zz\}$ est donc le cercle 
$S^1$ tout entier.
\end{enumerate}
\fincorrection
\correction{002418}
\begin{enumerate}
  \item d\'efinit une topologie. 
  \item ne d\'efinit pas une topologie,
car $\{a\}\cup \{ b,d\}=\{a,b,d\}$ n'est pas dans la collection.
  \item ne d\'efinit pas une topologie, car  $\{a,c,d\}\cap
\{b,c,d\}=\{c,d\}$ n'est pas dans la collection.
\end{enumerate}
\fincorrection
\correction{002419}
Il faut donc d\'emontrer que la collection de sous-ensembles de
$\R$ contenant $\emptyset$, $\R$ et tous les ensembles finis
v\'erifie les propri\'et\'es d'une collection d'ensembles
ferm\'es~:
\begin{itemize}
    \item toute int\'ersection d'ensembles ferm\'es est ferm\'e;
    \item toute r\'eunion finie d'ensembles ferm\'es est ferm\'e;
    \item $\emptyset$ et tout l'espace sont des ferm\'es.
\end{itemize}

Les trois propri\'et\'es sont \'evidemment v\'erifi\'ees dans ce
cas.


La topologie ainsi d\'efinie sur $\R$ n'est pas s\'epar\'ee. En
effet deux ouverts non-vides $\Omega$ et $\Omega'$ sont sous la
forme $\Omega=\R \setminus F$ et $\Omega'=\R \setminus F'$, o\`u
$F,F'$ sont ou bien finis ou bien vides. Alors $\Omega \cap
\Omega'=\R\setminus \left( F\cup F'\right)$ n'est pas vide, car
sinon ceci impliquerait que $\R = F\cup F'$ est finie ou vide, ce
qui est faux.
\fincorrection
\correction{002420}
\begin{enumerate}
\item  Supposons que $\mathcal{B}$ est une base de $\mathcal{T}$, et soit $\mathcal{O}$ un
ouvert arbitraire dans $\mathcal{T}$ et $x$ un point de $\mathcal{O}$. L'ouvert
$\mathcal{O}$ s'\'ecrit comme $\mathcal{O}=\bigcup_{i\in I} B_i$, o\`u $B_i \in
\mathcal{B}$ pour tout $i\in I$. En particulier il existe un $i_0\in I$
tel que $x\in B_{i_0}$.

\item R\'eciproquement, si $\mathcal{O}$ est un ouvert arbitraire, pour tout point $x\in \mathcal{O}$ il
  existe un $B_x\in \mathcal{B}$ tel que $x\in B_x\subset \mathcal{O}$. Par
  cons\'equent $\mathcal{O}=\bigcup_{x\in \mathcal{O}}B_x$.

\item Il suffit de montrer la propri\'et\'e \'enonc\'ee dans (1).
Soit $\mathcal{O}\in \mathcal{T}_n$ et soit $x$ un point arbitraire de $\mathcal{O}$.
D'apr\`es le cours, il existe un $r>0$ tel que $B(x,r)\subset
\mathcal{O}$.
 \emph{Remarque.} Une autre manni\`ere de formuler ceci
est de dire que l'ensemble des boules ouvertes euclidiennes forme
une base de la topologie $\mathcal{T}_n$.



Puisque l'ensemble $\Q^n$ est dense dans $\R^n$, il s'ensuit que
$B\left( x, \frac{r}{2}\right)$ contient un vecteur $q\in \Q^n$.
En particulier $\mathrm{dist} (x,q)< \frac{r}{2}$, d'o\`u $B\left( q,
\frac{r}{2}\right)\subset B\left( x, r\right)\subset \mathcal{O}$.

L'intervalle $]\mathrm{dist} (x,q),\frac{r}{2}[ $ est non-vide, donc il
contient un nombre rationnel $R$. Ainsi $x\in B(q,R)\subset
B\left( q, \frac{r}{2}\right)\subset \mathcal{O}$.

\item Puisque $\mathcal{B}'\subset \mathcal{T}_n$, ce qu'il reste \`a d\'emontrer
est \`a nouveau la propri\'et\'e \'enonc\'ee dans (1). Soit $\mathcal{O}$
un ouvert et $x\in \mathcal{O}$. Il existe un $r>0$ tel que $B(x,r)\subset
\mathcal{O}$.

D'apr\`es le cours
$$
\mathrm{dist} (y,x)= \|y-x\|_2\leq \sqrt{n}\|y-x\|_\infty\, .
$$

Il s'ensuit que
\begin{equation}\label{bile}
B_\infty \left( x,\frac{r}{\sqrt{n}} \right) =\left\{y\; ;\;
\|y-x\|_\infty<\frac{r}{\sqrt{n}}\right\}\subset B(x,r) \subset
\mathcal{O}\, .
\end{equation}


Or $B_\infty \left( x,\frac{r}{\sqrt{n}} \right)$ n'est rien
d'autre que le cube de centre de sym\'etrie $x$ et de longueur des
ar\^etes $\frac{2r}{\sqrt{n}}$. En particulier $B_\infty \left(
x,\frac{r}{\sqrt{n}} \right)\in \mathcal{B}'$.

On conclut que $\mathcal{B}'$ est une base de $\mathcal{T}_n$.

\item Soit $]0,1[\in \mathcal{T}_1$. Il n'existe pas d'intervalle de la
forme $]-\infty ,a[\, ,\; a\in \R \, ,$ ou $ ]b, +\infty[\, ,\;
b\in \R \, ,$ contenu dans $]0,1[$. Donc $\mathcal{B}''$ n'est pas une
base pour $\mathcal{T}_1$.

\item Supposons que $\mathcal{T}'\subset \mathcal{T}$. En particulier $\mathcal{B}'\subset
\mathcal{T}$.

Pour tout $a=\frac{m}{n}\in \Q$, o\`u $m\in \Z^*\, ,\, n\in \N^*\,
,$ p.g.c.d. $(m,n)=1$, on choisit $M_a,N_a$ deux points sur la
droite $\delta_a$ tels que $O\in ]M_a,N_a[$ et $\mathrm{dist}
(O,M_a)=\mathrm{dist} (O,N_a)=\frac{1}{n}$. Pour $a=0$ on choisit
$M_0=(1,0)\, ,\, N_0=(-1,0)$. Soit
$$
\mathcal{C}=\bigcup_{a\in \Q} ]M_a,N_a[\, .
$$

Par hypoth\`ese $\mathcal{C}\in \mathcal{B}'\subset \mathcal{T}$. En particulier,
puisque $O$ est un point de $\mathcal{C}$, il existe $r>0$ tel que
$Y\cap B(O,r)\subset \mathcal{C}$. Pour tout $a\in \Q$ on a donc
$\delta_a\cap B(O,r)\subset ]M_a,N_a[$, d'o\`u $r<\mathrm{dist}
(O,M_a)=\frac{1}{n}$. Comme ceci est v\'erifi\'e pour tout $n\in
\N^*$, il s'ensuit que $r\leq 0$, ce qui contredit le choix de
$r$.

On a obtenu une contradiction. Donc on ne peut pas avoir
$\mathcal{T}'\subset \mathcal{T}$.
\end{enumerate}
\fincorrection
\correction{002421}
\begin{enumerate}
  \item  On v\'erifie facilement les trois propri\'et\'es de
m\'etrique.
  \item Soit $f:\R_+ \to \R_+ \, ,\,
f(x)=\frac{x}{x+1}=1-\frac{1}{x+1}$. On a que $f(0)=0$ et
$f'(x)=\frac{1}{(x+1)^2}$,
  donc la fonction $f$ est croissante sur $\R_+$. L'in\'egalit\'e
  $f(x+y)\leq f(x)+f(y)$ pour $x,y\in \R_+$ est \'equivalente \`a $$\frac{1}{x+y+1}\geq \frac{1}{x+1}+\frac{1}{y+1}-1
  \Leftrightarrow \frac{1}{x+y+1}+1 \geq \frac{1}{x+1}+\frac{1}{y+1}\Leftrightarrow 1+x+y\leq (1+x)(1+y)\, .$$

La derni\`ere \'egalit\'e est \'evidemment v\'erifi\'ee pour
$x\geq 0\, ,\, y\geq 0\, $.
  \item D'apr\`es le cours, la m\'etrique $\mathrm{dist}$ et la m\'etrique
$\mathrm{dist}_2=\min (\mathrm{dist} ,1)$ sont topologiquement \'equivalentes.
Ainsi il suffit de montrer que $\mathrm{dist}_1$ et $\mathrm{dist}_2$ sont
topologiquement \'equivalentes.

Puisque $1+\mathrm{dist} \geq 1$, on a que $ \mathrm{dist}_1\leq \mathrm{dist}$. Aussi
$\mathrm{dist}_1\leq 1$, d'o\`u $\mathrm{dist}_1\leq \mathrm{dist}_2$.

La fonction $f$ \'etant croissante, pour tout $x,y$ on a que
$\mathrm{dist}_1(x,y)= f(\mathrm{dist} (x,y)) \geq f(\mathrm{dist}_2(x,y))$. D'autre part,
$\mathrm{dist}_2(x,y)\leq 1$ implique
$f(\mathrm{dist}_2(x,y))=\frac{\mathrm{dist}_2(x,y)}{1+\mathrm{dist}_2(x,y)}\geq
\frac{\mathrm{dist}_2(x,y)}{2}$.

On a obtenu que pour tout $x,y$,
$$\frac{\mathrm{dist}_2(x,y)}{2}\leq \mathrm{dist}_1(x,y)\leq \mathrm{dist}_2(x,y)\, .$$

Ainsi, les m\'etriques $\mathrm{dist}_1$ et $\mathrm{dist}_2$ sont \'equivalentes.
\end{enumerate}
\fincorrection
\correction{002422}
\begin{enumerate}
  \item Comme $d(x,y)  = 1$, si $x\neq y$, on a donc que $d(x,y) =
0 \Leftrightarrow x = y$. 
De plus, comme la relation $x\neq y$ est sym\'etrique, on $d(x,y) =
d(y,x),\-\forall x,y \in E$.
Soient $x,y,z\in E$, supposons $x = z$ ; ou bien $y = x$ ou bien $y$ est distinct de
$x$.
Dans le premier cas, $d(x,z) = d(x,y) = d(y,z) = 0$ et $d(x,z) = \sup(d(x,y),
d(y,z))$.
Dans le second cas, $d(x,y) = 1$ , d'o\`u 
$$0= d(x,x) = d(x,z) < \sup(d(x,y),d(x,y)) = 1.$$

Supposons $x\neq z $ ; ou $y $ est distinct de $x$ et de $z$, ou alors on a
l'une des
possibilit\'es: $y = x$ ou $y = z$.
Si les trois éléments sont deux \`a deux distincts, l'in\'egalit\'e est
trivialement
v\'erif\'ee ($1 = 1$ !).
Sinon, $d(x,y) = 1$ ou $d(y,z) = 1$, d'o\`u
$$1= d(x,z)\le \sup(d(x,y), d(y,z)).$$

  \item On suppose que $d(x,y) \neq d(y,z)$. Supposons alors que $d(x,z)
<\sup(d(x,y),d(y,z))$ et pour fixer les id\'ees que
$d(x,y) = \sup(d(x,y),d(y,z))$. Alors $d(y,z)< d(x,y)$ et $d(x,z)<d(x,y)$,
d'o\`u on d\'eduit que $\sup(d(x,z),d(z,y))< d(x,y)$.
Par ailleurs, $d(x,y)\le \sup(d(x,z), d(z,y))$. Les deux derni\`eres
in\'egalit\'es sont contradictoires.

  \item Soit $B_d(a,r)$ une boule ouverte; montrons qu'elle est ferm\'ee.
Soit $y\in E\setminus B_d(a,r)$; montrons qu'il existe une boule ouverte $B_d(y,\eta)$,
contenue dans $E \setminus B_d(a,r)$.
Si on choisit $\eta  = r/2 $ ou plus g\'en\'eralement $\eta< r$,
on obtient que, pour tout $z\in B_d(y,\eta)$,
$$d(a,z)\le \sup(d(a,y), d(y,z))\le \sup(d(a,y),\eta)).$$
Comme $d(a,y)\ge r$ et $d(y,z)<\eta<r$, on a , (d'apr\`es la deuxième question),
$d(a,z) = d(a,y)\ge r$. On en d\'eduit que $B_d(y,\eta)\subset E \setminus B_d(a,r)$ et par suite la
boule ouverte $B_d(a,r)$ est aussi ferm\'ee.

La preuve du fait que la boule ferm\'ee $B'_d(a,r) $ est aussi ouverte est
analogue.

  \item Soient $B_d(a,r)$ et $B_d(b,s)$ deux boules ouvertes ayant une
intersection non vide et soit $z_0\in B_d(a,r)\cap B_d(b,s)$. supposons que $r\le s$ et
montrons qu'alors $B_d(a,r)\subset B_d(b,s)$.
On regarde la distance \`a $b$ de tout $z\in B_d(a,r)$:
$$d(b,z)\le \sup(d(b,z_0), d(z_0,z)) < \sup(s,d(z_0,z))$$ 
puisque $z_0$ est dans $B_d(b,s)$.
Par ailleurs, on a:
$d(z_0,z)\le \sup(d(z_0,a), d(a,z))<r$.
On obtient une majoration de $d(b,z)$:
$d(b,z)< \sup(r,s) = s$, d'o\`u une inclusion de $B_d(a,r)$ dans $B_d(b,s)$.

 Cons\'equence: deux boules ouvertes de m\^eme rayon $r$ qui se rencontrent sont confondues.


  \item Soient $ A = B_d(a,r)$ et $B = B_d(b,r)$ deux boules ouvertes de rayon
$r$ contenues
dans une boule ferm\'ee  $C = B'_d(c,r)$ de m\^eme rayon.
Montrons que : 
$$\forall  x\in A,\quad \forall y \in B, \qquad r\le d(a,b)\le r.$$
L'in\'egali\'e ultram\'etrique montre que $d(x,y)\le \sup(d(x,c),d(c,y))$
et ce
sup est inf\'erieure \`a $r$ puisque chacune des boules $A$ et $B$ est
incluse dans $C$. Donc $d(x,y)\le r$.

Par ailleurs, introduisons dans l'estimation de $d(x,y)$ le centre des boules
respectives auxquelles ils appartiennent:
$d(x,y)\le \sup(d(x,a),d(a,y))$.  Si $d(x,a) = d(a,y)$, on aurait
$d(a,y)<r$ et
$y$ serait dans $A$, ce qui est impossible, $A$ et $B$ \'etant disjoints
d'apr\`es la quatrième question. Donc $d(a,y)\neq d(x,a)$, et en fait $d(a,y)> d(x,a)$ et 
$$d(x,y) = d(a,y).$$
On voit donc que dans le calcul de la distance $d(x,y)$ on peut remplacer
$x$ ou $y$ par le centre de la boule ouverte \`a laquelle il appartient.
Par suite 
$$d(x,y) = d(a,b)\ge r,\quad \forall x\in A,\forall y\in B$$

Et finalement
$$r\le d(x,y)\le r, \quad \forall x\in A,\forall y\in B$$ 
d'o\`u $d(A,B)= r$.

\end{enumerate}
\fincorrection
\correction{002423}
\begin{enumerate}
  \item Soit $x= \pm \frac ab = \pm \frac{a'}{b'}$. 
On écrit $a= p^\alpha a_1$, $b= p^\beta b_1$,... Alors l'équation  $ab'=a'b$ devient
$p^{\alpha+\beta'}a_1b_1'=p^{\alpha'+\beta}a_1'b_1$. Donc $\alpha+\beta'=\alpha'+\beta$ ou encore
$\alpha-\beta=\alpha'-\beta'$. Donc $\nu(\pm \frac ab)= \nu(\frac{a'}{b'})$.

  \item Soit $x=p^\alpha x_1$, $y=p^\beta y_1$ avec $\alpha,\beta\in \Zz$ et les numérateurs et dénominateurs de $x_1,y_1\in \Qq$ non divisibles par $p$. Alors $xy= p^{\alpha+\beta}x_1y_1$. Donc $\nu(xy)=\alpha+\beta=\nu(x)+\nu(y)$.

  \item Soit $x,y\in \Zz$, $x=p^\alpha x_1$, $y=p^\beta y_1$. Supposons par exemple $\alpha \le \beta$,
alors
$x+y=p^\alpha(x_1+p^{\beta-\alpha}y_1)$, avec $x_1+p^{\beta-\alpha}y_1 \in \Zz$.
Donc $\nu(x+y) \ge \alpha = \min(\nu(x),\nu(y))$.

Soit maintenant $x = \frac ab,y= \frac{a'}{b'}\in \Qq$.
Alors 
\begin{align*}
  \nu(x+y)  
  &= \nu( \frac ab+\frac{a'}{b'}) \\
  &= \nu(\frac{ab'+a'b}{bb'}) \\
  &= \nu(ab'+a'b)-\nu(bb') \\
  &\ge \min(\nu(ab'),\nu(a'b))-\nu(bb') \qquad \text{(grâce à l'inégalité sur les entiers)}, \\
  &\ge \min(\nu(a)+\nu(b'),\nu(a')+\nu(b))-\nu(b)-\nu(b') \\
  &\ge \min(\nu(a)+\nu(b')-\nu(b)-\nu(b'),\nu(a')+\nu(b)-\nu(b)-\nu(b')) \\
  &\ge \min(\nu(a)-\nu(b),\nu(a')-\nu(b')) \\
  &\ge \min(\nu(x),\nu(y)).\\
\end{align*}

  \item Il est clair que $d(x,y)=0$ si et seulement si $x=y$ et que
$d(x,y)=d(y,x)$.
Pour un triplet $(x,y,z)$ on a 
\begin{align*}
 d(x,z) 
  &= p^{-\nu(x-z)} \\
  &= p^{-\nu(x-y+y-z)} \\
  &\le p^{-\min(\nu(x-y),\nu(y-z))} \\
  &\le \max(p^{-\nu(x-y)},p^{-\nu(y-z)}) \\
  &\le \max(d(x,y),d(y,z)).\\ 
\end{align*}
\end{enumerate}
\fincorrection
\correction{002424}
\begin{enumerate}
\item
 
\begin{enumerate}
\item Si  $x \in\mathrm{Fr} (A) = \bar A \cap \overline{E \setminus A}$ , alors $\forall n \in \Nn^*$
la boule $B(x,{\frac 1 n})$ rencontre nécessairement $A$  (respectivement $E \setminus A$). Soit donc (axiome du choix)
$x_n $ (respectivement $y_n$) dans  $B(x,{\frac 1 n}) \cap A$ (respectivement $y_n$ dans $B(x,{\frac 1 n})\cap (E \setminus A).$
Alors les suites $x_n$ et $y_n$ répondent clairement à la question : 
On a une suite $(x_n)$ d'éléments de $A$ et une suite $(y_n)$ d'éléments du complémentaire
$E \setminus A$ de $A$ dans $E,$ qui convergent l'une et l'autre vers $x.$

\item On voit, qu'en posant pour $n \ge 1,$ d'une part $x_n=\frac 12-{\frac 1  {4n}}$ et d'autre part,
$y_n=\frac 12+{\frac 1 {4n}},$ on obtient, respectivement comme plus haut, une suite de points dans $A$
et une autre dans $E \setminus A$ qui convergent vers le même point $\frac 12 \in A$ qui, adhérent 
à $A$ comme à son complémentaire dans $E$ est donc dans la frontière de $A$ dans $E.$ Par contre,
si $x \in A$ est différent de $+\frac 12$, on voit que la boule (dans $E$) de centre $x$ et de rayon
$\frac 12 - x>0$ ne rencontre pas le complémentaire de $A$ et qu'en conséquence $[0, \frac 12[$ est
l'intérieur de $A$ {\it dans} E. 

A contrario une boule de centre $0$ et de rayon strictement positif rencontre toujours le complémentaire
de $A$ {\it dans} $\Rr$ ce qui permet aisément de voir que la frontière de $A$ {\it dans} $\Rr$ est
$\{0,\frac 12\}$.
\end{enumerate}
   
\item
Soient $E$ et $F$ deux espaces métriques respectivement au moyen des distances $d$ et
$d'$.

\begin{enumerate}
\item Pour abréger les notations posons: $\delta=\sup(d,d').$ C'est sur $E \times F,$ la distance donnée
par la formule:
$$\delta \bigl( (x,x'),(y,y')\bigr) = \sup\bigl(d(x,y),d'(x',y')\bigr)$$

Une boule pour $\delta$ n'est donc rien d'autre que le produit cartésien d'une boule pour $d$ avec
une boule pour $d'.$ Or ces produits cartésiens forment précisément une base d'ouverts qui définie
la topologie produit qui est donc aussi la topologie associée à la métrique $\delta.$

\item Soient $A\subset E$ et $B\subset F$. Soit  $(x,x')\in A\times B \setminus \mathrm{Fr} (A\times B)$ 
dans l'intérieur de $A\times B$ dans $E \times F.$ Cet intérieur est un ouvert pour la topologie
produit. La définition de cette topologie produit d'être engendrée par la base des produits cartésiens d'ouverts de $E$ avec des ouverts
de $F$ a comme conséquence l'existence d'un ouvert $U_x$ de $E$ qui contient $x$ et d'un autre
$U_{x'}$ de $F$ qui contient $x'$  tels que $U_{x}\times U_{x'}$  soient entièrement contenus dans
cet intérieur de $A\times B$ et donc à fortiori dans $A\times B$ lui même. Mais celà n'est possible
que si $U_{x}$ et $U_{x'}$ sont respectivement entièrement inclus dans $A$ et $B$ ce qui implique
que $x$ et $x'$ sont respectivement intérieurs dans $A$ et $B.$ Réciproquement si $x$ est intérieur
à $A$ et $x'$ intérieurs à $B$ et que $U_{x}$ et $U_{x'}$ soient alors des ouverts pour lesquels
$x\in U_{x}\subset A$ et $x'\in  U_{x'}\subset B,$ on voit que $U_{x}\times U_{x'}\subset A \times B$
est un ouvert pour la topologie qui contient $(x,x')$ qui est donc intérieur à $A\times B.$
$A  \setminus \mathrm{Fr} (A)$ de $A$ dans $E$ avec l'intérieur $ B \setminus \mathrm{Fr} ( B)$ de $B$
dans $F.$
\end{enumerate}

\item
 $E$ et $F$ sont toujours comme dans la deuxième question \c{c}i dessus.

\begin{enumerate}
\item Si $( \xi_n, \xi'_n)$ est une suite de points dans le complémentaire $E\times F \setminus A\times B$ de
$A\times B$ dans $E \times F,$ désignons par $N_1$ ( respectivement $N_2$ ) l'ensemble des $n \in \Nn$ pour lesquels
$\xi_n \notin A$ ( respectivement $\xi'_n \notin B.$ ) L'hypothèse montre que: $\Nn=N_1 \cup N_2.$
$\Nn$ étant un ensemble infini, il faut bien qu'au moins l'une des deux parties $N_1$ ou $N_2$
le soit aussi. Si par exemple $N_1$ est infini, on peut ranger ses éléments en ordre croissant 
$$n_0<n_1<n_2<.....<n_k<n_{k+1}<...$$
mais alors, par définition, la suite extraite $\xi_{n_k}$ a tous ses termes  dans $E \setminus A.$
Mutatis mutandis lorsque  $N_2$  est infini, ce qui est assuré dès lors que $N_1$ ne le serait pas.

\item Commen\c{c}ons par montrer, par exemple, que: $\mathrm{Fr} (A) \times \overline B \subset 
\mathrm{Fr} (A \times B) .$ En effet si $(x,x')\in \mathrm{Fr} (A) \times \overline B, $
il existe une suite $b_n$ dans $B$ qui converge vers $x' \in \bar B.$ De la même manière,
on trouve une suite $a_n$ d'éléments de $A$ qui converge vers $x \in \mathrm{Fr} (A) \subset \bar A.$
Mais aussi, comme on l'a vu plus haut, une suite d'éléments $c_n$ dans le complémentaire $E \setminus A$
de $A$ dans $E$ qui converge aussi vers $x.$ Mais alors $(a_n, b_n)$ est une suite de points de
$A \times B$ qui converge vers $(x,x')$ et $(c_n,b_n)$ est une suite de points du complémentaire de
$A \times  B$ qui converge aussi vers $(x,x')$ qui se trouve donc à la fois dans l'adhérence de 
$A \times B$ et de son complémentaire cqfd. En renversant les rôles de $A$ et $B$, on voit comment montrer que:
$\overline A \times\mathrm{Fr} ( B) \subset  \mathrm{Fr} (A \times B) .$ Ne reste donc plus qu'à montrer
l'inclusion:
$\mathrm{Fr} (A \times B) \subset \bigl( \mathrm{Fr} (A) \times \overline B \bigr) \cup 
\bigl( \overline A \times\mathrm{Fr} ( B) \bigr).$ Or,  $(x,x')\in \mathrm{Fr} (A \times B),$
est la limite d'une suite de points $( \xi_n, \xi'_n)$  dans le complémentaire $E\times F \setminus A\times B$ de
$A\times B$ dans $E \times F,$ comme aussi la limite  d'une suite de points $( \eta_n, \eta'_n)$ 
 de $A\times B,$ deuxième observation qui montre immédiatement que $x \in \bar A$ et $x' \in \bar B.$
Enfin on a vu en a) immédiatement plus haut, qu'on pouvait   extraire $\xi_{n_k}$    
dans $E \setminus A$ de la suite $x_n$  ou $\xi'_{n'_k}$   dans $E \setminus B$ de la suite 
$x'_n$ qui assure que $x$ est dans l'adhérence de $E \setminus A$ ou que $x'$ est dans celle de
$F \setminus  B$ ce qui assure que $x \in \mathrm{Fr} (A)$ ou $x' \in \mathrm{Fr} (B),$ et démontre
la dernière inclusion recherchée.
\end{enumerate}

\item

\begin{enumerate}
\item L'hypothèse   $(x,x')\notin A \times B$ et $x \in A$ implique que $x'\notin B,$ si bien
que $E\times \{ x'\}$ est entièrement contenu dans le complémentaire de $A \times B.$  
 Evidemment  $y \notin A$ implique que $\{y\}\times F$ est aussi entièrement contenu dans
ce même complémentaire de $A \times B.$ 

Mais alors la partie  $E\times \{ x'\}\cup \{y\}\times F$ est connexe pour la raison que 
$E\times \{ x'\}$ et $\{y\}\times F$ respectivement homéomorphes à $E$ et $F$ sont connexes
et que leur intersection qui est le point $(y,x')$ est non vide. Cette partie répond donc à la question.


\item Prenons $(x,x')\notin A \times B$ et $(y,y') \notin A \times B.$ exactement comme ci-dessus et qui sont
dans la même composante connexe de 
$(E \times F) \setminus (A \times B).$  Soit maintenant \linebreak
$(z,z')\in (E \times F) \setminus (A \times B);$ si $z \notin A,$ le raisonnement du a) se répète
pour voir que $(z,z')$ est raccordé à $(x,x')$ par une partie connexe. Mais si $z \in A,$ le
a) montre que $(z,z')$ est raccordé à $(y,y')$ par une partie connexe, et donc aussi à  $(x,x')$ 
qui a donc $(E \times F) \setminus (A \times B)$ tout entier comme composante connexe.
\end{enumerate}
\end{enumerate}
\fincorrection
\nocorrection
\nocorrection
\nocorrection
\nocorrection
\nocorrection
\nocorrection
\nocorrection
\nocorrection
\nocorrection
\nocorrection
\nocorrection
\nocorrection
\nocorrection
\nocorrection
\nocorrection
\nocorrection
\nocorrection
\nocorrection
\nocorrection
\nocorrection
\nocorrection
\nocorrection
\correction{006061}
\begin{enumerate}
\item L'ensemble $\emptyset$ est réunion vide de demi-droites. Il est clair
que $\C$ est dans $\cal T$ si, lorsque $z_0=0$, toute demi-droite
$[0,\rightarrow[$ est admissible (ce qui aurait dû être précisé dans
l'énoncé...Excuses!); par ailleurs cette famille d'ensembles est bien
stable par réunion quelconque. Remarquons que si $U\in\cal T$, $z\in
U\Longleftrightarrow [z,\rightarrow[\subset U$ (c'est-à-dire les
demi-droites forment une base d'ouverts). En effet si $z\in U$ il existe
$z'$ tel que $z\in [z',\rightarrow[\subset U$ et a fortiori
$[z,\rightarrow[\subset U$. Ainsi la famille est stable par intersection
finie : si $z\in U\cap U'$, $[z,\rightarrow[\subset U\cap U'$ et $U\cap
U'\in\cal T$. On a bien une topologie.

Cette topologie n'est pas séparée puisque, si $z\in [z',\rightarrow[$,
tout voisinage de $z'$ contient $z$.

\item La topologie n'étant pas même quasi-séparée, un singleton peut ne pas
être fermé. Soit $z_0$; $z\in \overline{\{z_0\}}$ si et seulement si tout
voisinage de
$z$ rencontre $z_0$; puisque tout voisinage de $z$ contient $[z,\rightarrow[$,
c'est équivalent à 
$z_0\in [z,\rightarrow[$ ou encore $z\in [0,z_0]$.

$\{0\}$ est le seul singleton fermé.

\item On déduit de ce qui précède que si $A\subset X$, $\overline
A=\cup_{z\in A}[0,z]$. En effet si $z\in A$,  $\overline A$ contient
$\overline {\{z\}}$ et donc le segment
$[0,z]$; réciproquement, si $z'\in[0,z]$ avec $z\in A$, $[z',\rightarrow[$,
et donc tout voisinage de $z'$, rencontre $A$; ainsi $z'\in\overline A$.


Supposons $A$ étoilé
par rapport à $0$; $A$ contient $\cup_{z\in A}[0,z]=\overline A$ et $A$ est
fermé.

Réciproquement supposons $A$ fermé; $A=\cup_{z\in A}[0,z]$ et en particulier
$A$ est étoilé par rapport à $0$. 
\end{enumerate}
\fincorrection
\nocorrection
\nocorrection
\nocorrection
\nocorrection
\nocorrection
\nocorrection
\nocorrection
\nocorrection
\nocorrection
\nocorrection
\nocorrection
\nocorrection
\nocorrection
\nocorrection
\nocorrection
\nocorrection
\nocorrection
\nocorrection
\nocorrection
\nocorrection
\nocorrection
\nocorrection
\nocorrection
\nocorrection
\nocorrection
\nocorrection
\nocorrection
\nocorrection
\nocorrection
\nocorrection
\nocorrection
\nocorrection
\nocorrection
\nocorrection
\nocorrection
\nocorrection
\nocorrection
\nocorrection
\nocorrection
\nocorrection
\nocorrection
\nocorrection
\nocorrection
\correction{002370}
\begin{enumerate}
  \item 
    \begin{enumerate}
      \item Si $A$ est compact et $B=\{b\}$ avec $b\notin A$.
Soit $a\in A$ alors $a \neq b$ donc il existe un voisinage ouvert de $a$, $U_a$
et un voisinage ouvert de $b$, $V_a$ tels que $U_a \cap V_a = \varnothing$.
Bien évidemment $A \subset \bigcup_{a\in A} U_a$. Comme $A$ est compact on peut extraire un ensemble fini $\mathcal{A} \subset A$ tel que $A \subset \bigcup_{a\in \mathcal{A}} U_a = : U^b$. Notons alors $V^b := \bigcap_{a\in \mathcal{A}} V_a$. $U^b$ est ouvert comme union d'ouverts et $V^b$ est ouvert comme intersection \emph{finie} d'ouverts. De plus $U^b\cap V^b= \varnothing$.

      \item Maintenant $B$ est compact. Pour chaque $b\in B$ le point précédent nous fournit $U^b$ et $V^b$ disjoints qui sont des voisinages ouverts respectifs de $A$ et $b$. On a $B \subset \bigcup_{b\in B} V^b$. On extrait un ensemble fini $\mathcal{B}$ de telle sorte que $B \subset \bigcup_{b\in \mathcal{B}} V^b =: V'$. $V'$ est un voisinage ouvert de $B$. Et si $U':= \bigcap_{b\in \mathcal{B}} U^b$ alors $U'$ est un ouvert contenant $A$, et $U' \cap V' = \varnothing$.
   \end{enumerate}

  \item Supposons que ce ne soit pas vrai alors
$$ \forall r>0 \quad \exists x \in X \qquad  (d(x,K)<r) \text{ et } x\notin U.$$
En prenant $r = \frac 1n$, $n\in \Nn^*$ nous obtenons une suite $(x_n)$  tel que 
$d(x_n,K)< \frac 1n$ et $x_n \notin U$.
Comme $d(x_n,K) < \frac 1n$ alors il existe $y_n \in K$ tel que
$d(x_n,y_n) < \frac 1 n$. Nous avons une suite $(y_n)$ dans $K$ compact donc on peut en extraire une sous-suite $y_{\phi(n)}$ qui converge ; notons $\ell$ sa limite, alors $\ell \in K$ car $K$ est compact. 

Regardons la suite extraite $(x_{\phi(n)})$, montrons quelle converge également vers $\ell$ :
$$d(x_{\phi(n)},\ell) \le d(x_{\phi(n)},y_{\phi(n)}) + d(y_{\phi(n)},\ell)$$
Les deux termes à droite de l'inégalité tendent vers $0$, donc $(x_{\phi(n)})$
tend vers $\ell$. Soit $F = X\setminus U$ alors $F$ est une fermé (car $U$ est ouvert) et $(x_{\phi(n)}) \in F$ donc la limite $\ell$ est dans $F$ également.
Donc $\ell \notin U$ et comme $K\subset U$ alors $\ell \notin K$. Nous avons montrer deux choses contradictoires $\ell \in K$ et $\ell \notin K$ ce qui prouve le résultat demandé.
\end{enumerate}
\fincorrection
\correction{002371}
Nous allons utiliser le fait qu'un ensemble $K$
est compact si et seulement si de toute suite d'éléments de $K$ on peut extraire une sous-suite convergente vers un élément de $K$.

Soit $(u_n)_{n\in\Nn}$ une suite convergente et soit $\ell$ sa limite. Notons
$$K = \{ u_n \mid n\in \Nn\} \cup \{ \ell\}.$$

Soit $(v_n)$ une suite d'éléments de $K$.
Si $(v_n)$ ne prend qu'un nombre fini de valeurs, on peut extraire une sous-suite constante, donc convergente.
Sinon $(v_n)$ prend une infinité de valeurs. Nous allons construire une suite convergente$(w_n)$ extraite de $(v_n)$.
Soit $w_0$ le premier des $(v_0,v_1,v_2,\ldots)$ qui appartient à $\{u_0,u_1,\ldots\}$. Soit $w_1$ le premier des $(v_1,v_2,\ldots)$ qui appartient à $\{u_1,u_2,\ldots\}$... Soit $w_n$ le premier des $(v_n,v_{n+1},\ldots)$ qui appartient à $\{u_n,u_{n+1},\ldots\}$.
Alors $(w_n)$ est une suite-extraite de $(v_n)$ et par construction
$(w_n)$ converge vers la limite de $(u_n)$, donc vers $\ell \in K$.
\fincorrection
\correction{002372}
\begin{enumerate}
  \item Notons $\ell = \mathrm{dist}(K,F)$. Alors il existe $(x_n)$ suite d'éléments de $K$ et $(y_n)$ suite d'éléments de $F$ telles que 
$\|x_n - y_n\| \rightarrow \ell$. Comme $K$ est compact alors on peut extraire de $(x_n)$ une sous-suite $(x_{\phi(n)})$ qui converge dans $K$. Notons $a\in K$ cette limite
Alors la suite extraite $(y_{\phi(n)})$ est bornée car
$$\| y_{\phi(n)} \| \le \| y_{\phi(n)} - x_{\phi(n)}\|  + \|x_{\phi(n)}\|.$$
La suite $(x_{\phi(n)})$ qui converge est donc bornée, et
la suite $(\| y_{\phi(n)} - x_{\phi(n)}\|)$ qui converge dans $\Rr$ (vers $\ell$) est bornée également. Donc la suite $(y_{\phi(n)})$ est bornée on peut donc en extraire une sous-suite convergente $(y_{\phi\circ \psi (n)})$.
De plus comme $F$ est fermé alors cette suite converge vers $b\in F$.
La suite $(x_{\phi\circ \psi (n)})$ extraite de $(x_{\phi(n)})$ converge vers $a\in K$. Et comme nous avons extrait deux suites $(x_n)$ et $(y_n)$ on a toujours $\|x_{\phi\circ \psi (n)} - y_{\phi\circ \psi (n)}\| \rightarrow \ell$.
A la limite nous obtenons $\| a-b\| = \ell$  avec $a\in K$ et $b\in F$.

  \item Remarque : si $K$ était supposé fermé mais pas compact alors le résultat précédent pourrait être faux. Par exemple pour $K = \{ (x,y) \in \Rr^2 \mid xy \ge 1 \text { et } y \ge 0\}$ et $F = \{ (x,y) \in \Rr^2 \mid y \le0\}$ nous avons $d(K,F) = 0$ mais $K\cap F = \varnothing$.
\end{enumerate}
\fincorrection
\correction{002373}
Comme $E$ est compact et $E \subset \bigcup_{y\in E} V_y$
il existe un ensemble fini $\mathcal{Y} \subset E$ tel que 
$E \subset \bigcup_{y\in \mathcal{Y}} V_y$. Sur chaque voisinage $V_y$, $f$ est bornée par une constante $M_y$. Notons $M= \max_{ y\in \mathcal{Y}} M_y$.
Alors $f$ est bornée sur $E$ par $M$. En effet pour un élément quelconque $x\in E$, il existe $y\in \mathcal{Y}$ tel que $y \subset V_y$ donc 
$f(x)$ est bornée par $M_y$ donc par $M$.
\fincorrection
\correction{002374}
\begin{enumerate}
  \item Soit $x = \lim x_n$. Soit $N\in \Nn$ ; montrons que $x$ est dans $F_N$. On a $x_N \in F_N$, $x_{N+1} \in F_{N+1} \subset F_N$, $x_{N+2} \in F_{N+2} \subset F_{N+1} \subset F_N$, etc. Donc pour tout $n\ge N$ alors $x_n \in F_N$. Comme $F_N$ est fermé, alors la limite $x$ est aussi dans $F_N$.
Ceci étant vrai quelque soit $N$, alors $x\in \bigcap_N F_N$.

Pour construire un exemple comme demandé il est nécessaire que de toute suite on ne puisse pas extraire de sous-suite convergente. Prenons par exemple dans $\Rr$, $F_n = [n,+\infty[$, alors $\bigcap_n F_n = \varnothing$.

  \item 
    \begin{enumerate}
     \item Pour chaque $n$ on prend $x_n \in K_n$, alors pour tout
$n$, $x_n \in K_0$ qui est compact donc on peut extraire une sous-suite convergente. Si $x$ est la limite de cette sous-suite alors $x \in K$.
Donc $K$ est non vide.

     \item Par l'absurde supposons que c'est faux, alors
$$\forall N\in\Nn \quad \exists n \ge N \quad \exists x_n \in K_n \text{ tel que } x_n\notin \Omega.$$
De la suite $(x_n)$, on peut extraire une sous-suite $x_{\phi(n)}$
qui converge vers $x\in K$. Or $x_n \in X\setminus \Omega$ qui est fermé
donc $x\in X\setminus \Omega$. Comme $K\subset \Omega$ alors $x\notin K$
ce qui est contradictoire.
    \end{enumerate}

\end{enumerate}
\fincorrection
\correction{002375}
Soit $x\in X$ et $\epsilon >0$.
\begin{enumerate}
  \item Pour tout $y\in [0,1]$ $f$ est continue en $(x,y)$ donc il existe
un $U(y)$ voisinage de $x$ et $[a(y),b(y)]$ voisinage de $y$ tel que pour
$(x',y') \in U(y)\times [a(y),b(y)]$ on ait $|f(x,y)-f(x',y')|\le \epsilon$.

  \item Comme $[0,1] \subset  \bigcup_{y\in[0,1]} [a(y),b(y)]$ et que $[0,1]$ est un compact de $\Rr$ il existe un ensemble fini $\mathcal{Y}$ tel que 
$[0,1] \subset  \bigcup_{y\in \mathcal{Y}} [a(y),b(y)]$. De plus quitte à
réduire les intervalles ont peut supposer qu'il sont disjoints et quitte à les réordonner on peut supposer que ce recouvrement s'écrit :
$$[0,1] = [0,t_1] \cup [t_1,t_2] \cup \ldots [t_k,1].$$

  \item Notons $U = \bigcap_{y\in \mathcal{Y}} U(y)$, c'est un voisinage de $x$
car l'intersection est finie. Pour $x'\in U$ nous avons
\begin{align*}
  |g(x)-g(x')| &= \left| \int_0^1 f(x,y)dy - \int_0^1 f(x',y)dy\right| \\
               &\le \int_0^1 |f(x,y)-f(x',y)| dy \\
               &\le \int_0^{t_1}  |f(x,y)-f(x',y)| dy + \int_{t_1}^{t_2}  \cdots + \int_{t_k}^{1}  |f(x,y)-f(x',y)| dy \\
               &\le \epsilon(t_1-0) + \epsilon(t_2-t_1)+\cdots +\epsilon(1-t_k) \\
                &\le \epsilon \\
\end{align*}
Donc $g$ est continue.
\end{enumerate}
\fincorrection
\correction{002376}
\begin{enumerate}
  \item Pour montrer que $A+B$ est fermé, nous allons montrer que toute suite de $A+B$ qui converge, converge vers un élément de $A+B$. Soit $(x_n)$ un suite
de $A+B$ qui converge vers $x\in E$. Alors il existe $a_n\in A$ et $b_n\in B$ tel que $x_n = a_n+b_n$.
Comme $A$ est compact on peut extraire une sous-suite $(a_{\phi(n)})$ qui converge vers $a\in A$. Alors $b_{\phi(n)} = x_{\phi(n)}-a_{\phi(n)}$ est convergente vers $x-a$. Notons $b= x-a$ comme $B$ est fermé alors $b\in B$.
Maintenant $x = a+b$ donc $x\in A+B$.
  \item Soit $F = \{ (x,y)\in \Rr^2 \mid xy\ge 1 \text{ et } x \ge 0\}$ , soit
$G = \{ (x,y)\in \Rr^2 \mid y \le 0 \text{ et } x\ge0\}$. Alors
$F+G = \{ (x,y)\in \Rr^2 \mid x > 0\}$ qui n'est pas un fermé.
\end{enumerate}
\fincorrection
\correction{002377}
\begin{enumerate}
  \item Supposons $f$ propre et soit $F$ un fermé. Montrons que $f(F)$ est un fermé. Soit $(y_n)$ une suite de $f(F)$ qui converge vers $y\in \Rr^n$. 
Notons $K$ l'union de $\{y_n\}_{n\in\Nn}$ et de $\{ y\}$. Alors $K$ est compact.
Comme $y_n \in f(F)$, il existe $x_n\in F$ tel que $f(x_n) = y_n$.
En fait $x_n \in f^{-1}(K)$ qui est compact car $f$ est propre. Donc
de $(x_n)$ on peut extraire une sous-suite convergente $(x_{\phi(n)})$, on note $x$ la limite de cette sous-suite. Comme $x_{\phi(n)} \in F$ et que $F$ est fermé alors $x\in F$. Comme $f$ est continue alors $y_{\phi(n)} = f(x_{\phi(n)})$ tend vers $f(x)$, or $y_{\phi(n)}$ tend aussi vers $y$.
Par unicité de la limite $y=f(x)$. Donc $y \in f(F)$ et $f(F)$ est fermé.

  \item 
Dire $ \|f(x)\|\to \infty$ quand $\|x\|\to \infty$ est équivalent à
$$ \forall M >0\quad \exists m >0\quad \forall x \in \Rr^n \quad (x \notin B(0,m) \Rightarrow  f(x)\notin B(0,M)).$$
  \begin{enumerate}
    \item Supposons $f$ propre, soit $M>0$. Alors $B(0,M)$ est un compact
(nous sommes dans $\Rr^n$) donc $f^{-1}(B(0,M))$ est compact donc borné, c'est-à-dire qu'il existe $m>0$ tel que $f^{-1}(B(0,M)) \subset B(0,m)$.
Donc si $x\notin B(0,m)$ alors $f(x) \notin B(0,M)$.

    \item Réciproquement, soit $K$ un compact de $\Rr^n$. Comme $f$ est continue
et que $K$ est fermé alors $f^{-1}(K)$ est un fermé. Reste à montrer que $f^{-1}(K)$ est borné. Comme $K$ est compact alors il existe $M>0$ tel que 
$K \subset B(0,M)$, par hypothèse il existe $m>0$ tel que si 
$ x\notin B(0,m)$ alors $f(x) \notin B(0,M)$, ce qui s'écrit aussi 
par contraposition : ``si $f(x) \in B(0,M)$ alors $x\in B(0,m)$'', donc
$f^{-1}(B(0,M)) \subset B(0,m)$. Or $K \subset B(0,M)$ donc
$f^{-1}(K) \subset f^{-1}(B(0,M)) \subset B(0,m)$. Donc
$f^{-1}(K)$ est borné donc compact.
  \end{enumerate}
\end{enumerate}
\fincorrection
\correction{002378}
\begin{enumerate}
  \item Soit $f_n$ la fonction affine suivante $f_n(t) = 0$ pour
$t\in [0,\frac{1}{n+1}]$ et pour $t\in [\frac 1n,1]$. Sur $[\frac{1}{n+1},
\frac{1}{n}]$ on définit une ``dent'' qui vaut $0$ aux extrémités et $1$ au milieu du segment. Alors si $B$ dénote la boule unité fermée (centrée en la fonction nulle), nous avons $d_\infty(f_n,0) = \sup |f_n(t)| = 1$ donc
$f_n \in B$. Par contre si $p\neq q$ alors $d(f_p,f_q)= 1$ donc la suite
$(f_n)$ et toute sous-suite ne sont pas de Cauchy. Si $B$ était compact alors 
on pourrait extraire une sous-suite convergente donc de Cauchy. Contradiction.

  \item Notons $x^n = (0,0,\ldots,0,1,0,0,\ldots)$ la suite de $l^\infty$
(le $1$ est à la $n$-ième place). Alors $x^n$ est dans la boule unité fermée $B$ centrée en $0$. De plus si $p\neq q$, alors $d_\infty(x^p,x^q) = 1$. Donc toute sous-suite extraite de $(x_n)$ n'est pas de Cauchy donc ne peut pas converger. 
Donc $B$ n'est pas compact.
\end{enumerate}
\fincorrection
\nocorrection
\correction{002380}
\begin{enumerate}
  \item Si $f$ a deux points fixes $x\neq y$, alors $d(x,y) = d(f(x),f(y)) < d(x,y)$. Ce qui est absurde. Donc $f$ a au plus un point fixe.

  \item $f$ est continue et $X$ compact donc $X_1 = f(X)$ est compact,
par récurrence si $X_{n-1}$ est compact alors $X_n = f(X_{n-1})$ est compact.
De plus $f : X\rightarrow X$, donc $f(X) \subset X$ soit $X_1 \subset X$, puis
$f(X_1) \subset f(X)$ soit $X_2 \subset X_1$, etc. Par récurrence $X_n \subset X_{n-1} \subset \cdots \subset X_1 \subset X$. Comme chaque $X_n$ est non vide alors $Y$ n'est pas vide (voir l'exercice \ref{exo 2}).

  \item Montrons d'abord que $f(Y) \subset Y$. Si $y\in Y$, alors pour tout $n\ge 0$ on a $y \in X_n$ donc $f(y) \in f(X_n)=X_{n+1}$ pour tout $n\ge 0$.
Donc pour tout $n>0$, $f(y) \in X_n$, or $f(y) \in X_0=X$. Donc $f(y) \in Y$.

Réciproquement montrons $Y \subset f(Y)$. Soit $y\in Y$, pour chaque $n \ge 0$, 
$y\in X_{n+1}=f(X_{n})$. Donc il existe $x_{n} \in X_{n}$ tel que $y = f(x_n)$.
Nous avons construit $(x_n)$ une suite d'élément de $X$ compact, on peut donc en extraire une sous-suite convergente $(x_{\phi(n)})$. Notons $x$ la limite, par l'exercice \ref{exo 2}, $x\in Y$.
Alors $y = f(x_{\phi(n)})$ pour tout $n$ et $f$ est continue donc à la limite
$y=f(x)$. Donc $y\in f(Y)$.

Soit $y \neq y'\in Y$ tel que $d(y,y')= \mathrm{diam}\,Y > 0$. Comme $Y=f(Y)$ alors il existe $x,x' \in Y$ tel que $y=f(x)$ et $y'=f(x')$. Or $d(y,y')=d(f(x),f(x')) < d(x,x')$. On a trouvé deux élements de $Y$ tel $d(x,x')$ est strictement plus grand que le diamètre de $Y$ ce qui est absurde. Donc $y=y'$ et le diamètre est zéro.

  \item Comme le diamètre est zéro alors $Y$ est composé d'un seul point $\{ p \}$ et comme $f(Y)=Y$ alors $f(p)=p$. Donc $p$ a un point fixe et nous savons que c'est le seul. Par la construction de $Y$ pour tout point $x_0 \in X$ la suite $x_n = f^n(x_0)$ converge vers $p$. 
\end{enumerate}
\fincorrection
\correction{002381}
\begin{enumerate}
  \item Comme $E\times E$ est compact alors de la suite $(a_n,b_n)$ on peut
extraire une sous-suite $(a_{\phi(n)},b_{\phi(n)})$ qui converge vers $(a_\infty,b_\infty)$.
Soit $\epsilon > 0$ il existe $n\in \Nn$ tel que si $k\ge n$
alors $d(a_{\phi(k)},a_\infty) < \frac\epsilon2$ et $d(b_{\phi(k)},b_\infty) < \frac\epsilon2$.
Donc en particulier $d(a_{\phi(n+1)},a_{\phi(n)}) \le d(a_{\phi(n+1)},a_\infty)
+ d(a_\infty,a_{\phi(n)}) < \epsilon$.
La propri\'eté pour $f$ s'écrit ici $d(a_k,b_{k'})\le d(a_{k+1},b_{k'+1}) \ge $.
Donc $d(a_{\phi(n+1)-\phi(n)}, a_0) \le d(a_{\phi(n+1)-\phi(n)+1}, a_1) \le \ldots \le d(a_{\phi(n+1)-1}, a_{\phi(n)-1}) \le  d(a_{\phi(n+1)}, a_{\phi(n)}) < \epsilon$. Donc pour $k = \phi(n+1)-\phi(n)$, sachant que $a_0=a$ alors $d(a_{k}, a) < \epsilon$. Même chose avec $(b_n)$.

  \item 
    \begin{enumerate}
     \item Soit $a \in E$ et $\epsilon >0$ alors il existe $k\ge 1$ tel que $a_k = f^k(a) \in f(E)$     avec $d(a,a_k) < \epsilon $. Donc $f(E)$ est dense dans $E$.

     \item Soit $u_n = d(a_n,b_n)$. Alors par la propriété pour $f$, $(u_n)$ est une suite croissante de $\Rr$. Comme $E$ est compact alors son diamètre est borné, donc $(u_n)$ est majorée. La suite $(u_n)$ est croissante et majorée donc converge vers $u$.

Maintenant $u_n - u_0 \ge0$ et 
$$0 \le u_n -u_0=d(a_n,b_n) -d(a,b) \le d(a_n,a) + d(a,b)+d(b,b_n)-d(a,b) = d(a_n,a)+d(b_n,b).$$
Donc $u_n$ tend vers $u_0$. Comme $(u_n)$ est croissante alors $u_n = u_0$ pour tout $n$.
En particulier $u_1=u_0$ donc $d(a_1,b_1)=d(a_0,b_0)$ soit $d(f(a),f(b))=d(a,b)$. Donc $f$ est une isométrie.
      \item $f$ est une isométrie donc continue (elle est $1$ lipschitzi\`enne !). $E$ est compact donc
$f(E)$ est compact donc fermé or $f(E)$ est dense donc $f(E)=E$. Donc $f$ est surjective
 \end{enumerate}
\end{enumerate}
\fincorrection
\correction{002382}
Dire que $ i: (X, |.|) \to (X,d)$ est continue c'est exactement dire que tout
ensemble $U$ ouvert pour $d$ est ouvert pour $|.|$ (car $i^{-1}(U) = U$).
\begin{enumerate}
  \item Soit $K$ un compact pour $|.|$. Soit $U_i$, $i\in I$ tels
que $K \subset \bigcup_{i\in I} U_i$ et tels que $U_i$ soient des ouverts pour $d$. Alors les $U_i$ sont aussi des ouverts pour la topologie définie par $|.|$. 
Comme $K$ est compact pour $|.|$ alors on peut extraire un ensemble fini $J\subset I$ tel que $K \subset \bigcup_{i\in J} U_i$. Donc $K$ est aussi compact pour $d$. 

Si $F$ est un fermé pour $|.|$ alors $F\subset [0,1]$ est compact pour $|.|$
Donc compact pour $d$, donc fermé pour $d$. 

  \item Si $U$ est un ouvert pour $d$ alors $U$ est un ouvert pour $|.|$.
Car $i$ est continue. Réciproquement si $U$ est un ouvert pour $|.|$ alors
$F = X\setminus U$ est un fermé pour $|.|$ donc $F$ est un fermé pour $d$ par la question précédente, donc $U = X \setminus F$ est un ouvert pour $d$.
Conclusion les ouverts pour $|.|$ et $d$ sont les mêmes donc
$|.|$ et $d$ définissent la même topologie.
\end{enumerate}
\fincorrection
\nocorrection
\nocorrection
\nocorrection
\nocorrection
\nocorrection
\nocorrection
\nocorrection
\nocorrection
\nocorrection
\nocorrection
\nocorrection
\nocorrection
\nocorrection
\nocorrection
\nocorrection
\nocorrection
\nocorrection
\nocorrection
\nocorrection
\nocorrection
\nocorrection
\nocorrection
\nocorrection
\nocorrection
\nocorrection
\nocorrection
\nocorrection
\nocorrection
\nocorrection
\nocorrection
\nocorrection
\nocorrection
\nocorrection
\nocorrection
\nocorrection
\nocorrection
\nocorrection
\nocorrection
\nocorrection
\nocorrection
\nocorrection
\nocorrection
\nocorrection
\nocorrection
\nocorrection
\nocorrection
\nocorrection
\correction{002353}
\begin{enumerate}
  \item Sens direct. Si $f$ est continue alors $\{x \mid f(x) < \lambda\} = 
f^{-1}(]-\infty,\lambda[)$ est un ouvert comme image réciproque par une application continue de l'intervalle ouvert $]-\infty,\lambda[$. De même avec
$]\lambda,+\infty[$.

Réciproque. Tout d'abord, tout intervalle ouvert $]a,b[$, ($a<b$) peut s'écrire
$$]a,b[= ]-\infty,b[ \cap ]a,+\infty[.$$
Donc 
$$f^{-1}(]a,b[) = f^{-1}(]-\infty,b[) \cap f^{-1}(]a,+\infty[)$$
est une intersection de deux ouverts donc un ouvert de $X$.
Soit $O$ un ouvert de $\Rr$, alors $O$ peut s'écrire comme l'union dénombrables d'intervalles ouverts :
$$O = \bigcup_{i\in I} {]a_i,b_i[}.$$
Donc 
$$f^{-1}(O)= \bigcup_{i\in I} f^{-1}(]a_i,b_i[)$$
est une union d'ouvert donc un ouvert de $X$ .
  \item Nous le faisons d'abord pour un intervalle ouvert $]a,b[$.
$$]a,b[ = \bigcup_{j\in\Nn^*} [a+\frac 1n,b-\frac 1n].$$
Donc 
$$f^{-1}(]a,b[) = \bigcup_{j\in\Nn^*} f^{-1}([a+\frac 1j,b-\frac 1j]),$$
est une union dénombrable de fermés.
Maintenant comme pour la première question, tout ouvert $O$ de $\Rr$
s'écrit $O = \bigcup_{i\in I} ]a_i,b_i[$, avec $I$ dénombrable.
Donc on peut écrire 
$$f^{-1}(O)= \bigcup_{i\in I}\bigcup_{j\in\Nn^*} f^{-1}([a_i+\frac 1j,b_i-\frac 1j]),$$
qui est une union dénombrable de fermés (mais c'est un ouvert !).
   \end{enumerate}
\fincorrection
\correction{002354}
\begin{enumerate}
  \item Soit $F$ l'application définie par
$F(f) = \int_0^1 |f|.$ Alors 
$$|F(f)-F(g)| = | \int_0^1 |f|-|g| | \le \int_0^1 |f-g| = d_1(f,g) \le d_\infty (f,g).$$ Donc pour les deux distances $d_1$ et $d_\infty$, $F$ est lipschitzienne de rapport $1$. 
  \item Soit $\epsilon >0$ alors en posant $\eta = \epsilon$ on obtient la continuité : 
si $d(x,y) < \epsilon$ alors $$|\ell(x)-\ell(y)| \leq \epsilon.$$ 
Donc $\ell$ est continue, et $c_0 = \ell{-1}(\{0\})$ est un fermé , car c'est l'image réciproque du fermé $\{0\}$ par l'application continue $\ell$.
   \end{enumerate}
\fincorrection
\correction{002355}
Soit $A = \{ x \in X \mid f(x)=g(x) \}$. Alors soit 
$C = X\setminus A = \{ x \in X \mid f(x) \neq g(x) \}$.
Soit $x \in C$ comme $f(x)\neq g(x)$ et que $Y$ est séparé, il existe un voisinage ouvert $V_1$ de $f(x)$ et $V_2$ de $g(x)$ tel que 
$V_1 \cap V_2 = \varnothing$.
Notons $U = f^{-1}(V_1) \cap g^{-1}(V_2)$. Alors $U$ est un ouvert de $X$ contenant $x$. Maintenant pour 
$x' \in U$, alors $f(x') \in V_1$,
$g(x')\in V_2$ donc $f(x')\neq g(x')$, donc $x'\in C$. Bilan 
$U$ est inclus dans $C$. Donc $C$ est ouvert.

Application : si $A$ est dense dans $X$ alors $\bar A = X$, mais comme $A$ est fermé $A = \bar A$. Donc $A=X$, c'est-à-dire $f$ et $g$ sont égales partout. 
\fincorrection
\correction{002356}
\begin{enumerate}
  \item Soit $P$ un polynôme, et $F$ un fermé de $\Rr$.  Soit $(y_n)$
une suite convergente d'éléments de $P(F)$, et $y\in \Rr$ sa limite.
Il existe $x_n\in F$ tel que $y_n=P(x_n)$. Comme $(y_n)$ est bornée (car convergente) alors $(x_n)$ aussi est bornée, en effet un polynôme
n'a une limite infini qu'en $\pm \infty$. Comme $(x_n)$ est une suite bornée de $\Rr$ on peut en extraire une sous-suite convergente $(x_{\phi(n)})$ de limite $x$. Comme $F$ est fermé, $x\in F$. Comme $P$ est continue (c'est un polynôme)
alors $y_{\phi(n)} = P(x_ {\phi(n)}) \rightarrow P(x)$, mais $(y_{\phi(n)})$ 
converge aussi vers $y$. Par unicité de la limite $y = P(x) \in P(F)$.
Donc $P(F)$ est fermé. 

  \item Soit $X=Y=\Rr$ et  $H=(xy=1)$ est un fermé de $X\times Y$, mais
si $\pi(x,y)=x$ alors $\pi(H)=\Rr^*$ n'est pas un fermé de $X=\Rr$.
  \item A vérifier...
   \end{enumerate}
\fincorrection
\correction{002357}
\begin{enumerate} 
 \item $\Rightarrow$.
Soit $f$ continue et $y\in f(\bar A)$. Il existe $x\in \bar A$ tel que $y=f(x)$. Soit $x_n \in A$ tel que $(x_n)$ converge vers $x$. Alors
 $y_n=f(x_n)\in A$. Comme $f$ est continue alors 
$(y_n)$ converge vers $f(x)=y$. Donc $y$ est adhérent à $f(A)$.
Conclusion $f(\bar A) \subset \overline{f(A)}$.

$\Leftarrow$. Soit $f : X \rightarrow Y$ et soit $F$ un fermé de $Y$.
Notons $A = f^{-1}(F)$. Alors $f(A) \subset F$ donc l'équation 
$f(\bar A) \subset \overline{f(A)}$
devient $f(\bar A) \subset \bar F = F$ car $F$ est fermé.
Donc $\bar A \subset f^{-1}(F)=A$. Donc
$\bar A \subset A$, d'o\`u $\bar A=A$. Donc $A$ est fermé. Bilan l'image réciproque de tout fermé $F$ est un fermé, donc $f$ est continue.

Application : si $A$ est dense, alors $\bar A= X$, et sous les hypothèses précédentes alors $f(A)$ est dense dans l'image de $X$ par $f$ : en effet $\overline{f(A)}$ contient $f(\bar A) = f(X)$

  \item $\Rightarrow$. Soit $f$ fermé et soit $A \subset X$.
Alors $A \subset \bar A$ donc $f(A) \subset f(\bar A)$, donc
comme $\bar A$ est un fermé et $f$ est fermée alors $f(\bar A)$ est un fermé contenant $f(A)$. Mais comme $\overline {f(A)}$ est le plus petit fermé contenant $f(A)$ alors $\overline {f(A)} \subset f(\bar A)$.

$\Leftarrow$. La relation pour un fermé $F$ donne 
$\overline {f(F)} \subset f(\bar F) = f(F)$. Donc $\overline {f(F)} = f(F)$.
Donc $f(F)$ est fermé. Donc $f$ est fermée.

Même type de raisonnement avec $f$ ouverte.
\end{enumerate}
\fincorrection
\nocorrection
\nocorrection
\correction{002360}
\begin{enumerate}
  \item Supposons que $f$ ne tende pas vers $0$. Soit $\epsilon >0$ fixé.
Pour tout $n\geq 0$, il existe $x_n \geq n$ tel que $|f(x_n)|>\epsilon$.
Sans perte de généralité nous supposons $f(x_n)>\epsilon$.
Appliquons l'uniforme continuité : soit $\epsilon' = \frac \epsilon 2$,
Il existe $\eta$ tel que pour  $|x_n-y|\le \eta$ on ait $|f(x_ n)-f(y)| < \epsilon'$. Donc pour un tel $y$, $f(y) > \frac \epsilon 2 >0$.
Donc $f$ est strictement positive sur $[x_n-\eta,x_n+\eta]$.
Notons alors $(p_n)$ définie par 
$p_{2n} = x_{n}-\eta$, $p_{2n+1} = x_{n}+\eta$.
Soit $I(x) = \int_0^x f$.
Alors $I(p_{2n+1})- I(p_{2n}) = \int_{x_n-\eta}^{x_n+\eta} f(t) dt \geq \frac \epsilon 2 \cdot 2\eta = \epsilon \eta$.
Donc la suite $(I(p_n))$ n'est pas de une suite de Cauchy, donc ne converge pas, donc la fonction $x\mapsto I(x)$ ne converge pas non plus, et donc 
$\int_0^\infty f(t) dt$ diverge.


  \item Par le changement de variable $u=t^2$ puis une intégration par partie, on montre que l'intégrale $\int_0^\infty \sin(t^2) dt$ converge, mais comme $f(x)= \sin(x^2)$ ne tend pas vers $0$ alors
$f$ n'est pas uniformément continue sur $\Rr$.
   \end{enumerate}
\fincorrection
\nocorrection
\nocorrection
\nocorrection
\nocorrection
\nocorrection
\nocorrection
\nocorrection
\nocorrection
\nocorrection
\nocorrection
\nocorrection
\nocorrection
\nocorrection
\nocorrection
\nocorrection
\nocorrection
\nocorrection
\nocorrection
\nocorrection
\nocorrection
\nocorrection
\nocorrection
\nocorrection
\nocorrection
\nocorrection
\nocorrection
\nocorrection
\nocorrection
\nocorrection
\nocorrection
\nocorrection
\nocorrection
\nocorrection
\nocorrection
\nocorrection
\nocorrection
\nocorrection
\nocorrection
\nocorrection
\nocorrection
\nocorrection
\nocorrection
\nocorrection
\nocorrection
\nocorrection
\nocorrection
\nocorrection
\correction{002361}
Pour $x=(x_1,x_2)\in E_1\times E_2$ on définit $\|x\|=\max(\|x_1\|,\|x_2\|)$.
\begin{enumerate}
  \item Sens $\Leftarrow$.  Soit $M>0$ tel que $ \|B(x) \|\leq M \| x_1\| \| x_2\|$.
Montrons que $B$ en continue au point $x=(x_1,x_2)$ fixé. Soit $y=(y_1,y_2)$ alors
$$B(x+y)-B(x)=B(x_1+y_1,x_2+y_2)-B(x_1,x_2)= B(x_1,y_2)+B(x_2,y_1)+B(y_1,y_2).$$
Donc $$\|B(x+y)-B(x)\| \le  M\| x_1\| \| y_2\|+M\| x_2\| \| y_1\|+M\| y_1\| \| y_2\|.$$
Pour $\|y_1\| \le \frac{\epsilon}{M\| x_1\|}$ on a  $M\| x_1\| \| y_2\| \le \epsilon$
(si $x_1=0$ il n'y a rien à choisir ici). 
Pour $\|y_2\| \le \frac{\epsilon}{M\| x_2\|}$ on a  $M\| x_2\| \| y_1\| \le \epsilon$
(si $x_2=0$ il n'y a rien à choisir ici).
Enfin pour $\|y_1\|\le \sqrt {\frac \epsilon M}$ et $\|y_2\|\le \sqrt {\frac \epsilon M}$
on a $M\| y_1\| \| y_2\|\le \epsilon$.
Donc en prenant $\eta = \min (\frac{\epsilon}{M\| x_1\|},\frac{\epsilon}{M\| x_2\|},\sqrt {\frac \epsilon M})$, on obtient que pour $\|y\| = \max (\|y_1\|,\|y_2\|) \le \eta$ on a $\|B(x+y)-B(x)\| \le 3\epsilon$.
Ce qui prouve la continuité.
Donc $B$ est continue sur $E_1\times E_2$.

  \item Sens $\Rightarrow$. Si $B$ est continue partout, en particulier 
elle est continue en $0$. Je choisis $\epsilon =1$, il existe $\eta >0$
tel que $\|x\| \le \eta$ alors $\|B(x)\|\le 1$. Donc pour $\|x_1\|\le \eta$
et $\|x_2\|\le \eta$ on a $\|B(x_1,x_2)\|\le 1$.
Soit maintenant $y = (y_1,y_2) \in E_1\times E_2$, ($y_1\neq 0, y_2\neq 0$) on a $(\eta \frac{y_1}{\| y_1 \|},\eta \frac{y_2}{\| y_2 \|})$ de norme $\le \eta$ donc
$B(\eta \frac{y_1}{\| y_1 \|},\eta \frac{y_2}{\| y_2 \|}) \le 1$ et par bilinéarité cela fournit :
$B(y_1,y_2)\le \frac{1}{\eta^2}{\| y_1 \|}{\| y_2 \|}$, et ce pour tout $(y_1,y_2)$. La constante cherchée étant $\frac 1{\eta^2}$.

\end{enumerate}
\fincorrection
\correction{002362}
Comme $L$ est linéaire il suffit de montrer que $L$ est continue en $0$.
Supposons que cela ne soit pas vrai, alors il faut nier la continuité de $L$ en $0$ qui s'écrit :
$$\forall \epsilon >0 \qquad \exists \eta >0 \qquad \forall x\in E  \qquad(\|x\|<\eta \Rightarrow \|L(x)\|<\epsilon).$$
La négation s'écrit alors :
$$\exists \epsilon >0 \qquad \forall \eta >0 \qquad \exists x \in E \qquad (\|x\|<\eta \ \text{ et }\  \|L(x)\| \ge \epsilon).$$
Soit donc un tel $\epsilon >0$ de la négation, pour $\eta$ de la forme
$\eta=\frac1n$, on obtient $y_n$ tel que $\|y_n\|<\frac1n$ et $\|L(y_n)\|\ge \epsilon$.
On pose $x_n = \sqrt n y_n$, alors $\|x_n\| = \sqrt n \|y_n\| < \frac 1 {\sqrt n}$ donc $(x_n)$ est une suite de $E$ qui tend vers $0$. 
Par contre $\| L(x_n) \| = \sqrt n \|L(y_n)\| \ge \epsilon \sqrt n$, donc
la suite $(L(x_n))$ n'est pas bornée.
Par contraposition nous avons obtenu le résultat souhaité.
\fincorrection
\correction{002363}
\begin{enumerate}
  \item Si $f$ est linéaire et bornée sur la boule unité alors elle est continue
(voir le cours ou refaire la démonstration).
  
  \item Il reste à montrer que $f$ est linéaire : on a déjà 
$f(x+y)=f(x)+f(y)$ pour tout $x,y$ reste donc à prouver $f(\lambda x) = \lambda f(x)$. Pour tout $\lambda \in \Rr$ et $x\in E$.
  \begin{itemize}
     \item Pour $\lambda \in \Zz$, c'est une récurrence, $f(2x) = f(x+x)= f(x)+f(x) = 2f(x)$. Puis $f(3x) = f(2x+x) =f(2x)+f(x)=2f(x)+f(x)=3f(x)$ etc.
Donc $f(nx) = nf(x)$ pour $n\in \Nn$.
De plus $0=f(0)=f(x+(-x))=f(x)+f(-x)$ donc $f(-x)=-f(x)$.
Ensuite on a $f(-nx) = -nf(x)$ pour $n \in \Nn$.
Bilan : pour tout $\lambda \in \Zz$ on a $f(\lambda x) = \lambda f(x)$.

     \item Pour $\lambda \in \Qq$, soit $\lambda = \frac pq$, $p,q \in \Zz$.
$$f(\frac pq x) = pf(\frac 1 q x) = \frac pqq f(\frac x q) = \frac pq f(q \frac xq) = \frac pq f(x).$$
Nous avons utilisé intensivement le premier point. 
     \item Soit $\lambda \in \Rr$ alors il existe une suite $(\lambda_n)$ d'élément de $\Qq$ qui converge vers $\lambda$. Fixons $x\in E$.
$$f(\lambda x)-\lambda f (x) = f(\lambda x)-f(\lambda_nx)+f(\lambda_nx)-\lambda f(x)= f((\lambda-\lambda_n)x)+(\lambda_n-\lambda)f(x).$$
Nous avons utilisé le second point.
Soit $\epsilon \in \Qq_+^*$. Pour $n$ assez grand on a 
 $\|(\lambda-\lambda_n)x \| < \epsilon$. Donc
$\|\frac 1 \epsilon(\lambda-\lambda_n)x \| \in B(0,1)$ or $f$ est bornée sur
la boule unité donc il existe $M>0$ tel que $f(\frac 1 \epsilon(\lambda-\lambda_n)x) \le M$ (quelque soit $n$). Donc $f(\lambda-\lambda_n)x) \le M\epsilon$
($\epsilon$ est rationnel donc on peut le ``sortir''). 
De même pour $n$ assez grand on a $(\lambda_n-\lambda)f(x) < \epsilon$. Maintenant
$$ \|f(\lambda x)-\lambda f (x)\| \le \| f((\lambda-\lambda_n)x)\| + \| (\lambda_n-\lambda)f(x)\| < M\epsilon+\epsilon.$$
Donc pour $x,\lambda$ fix\'es, $\|f(\lambda x)-\lambda f (x)\|$ est aussi petit que l'on veut, donc est nul ! D'o\`u $f(\lambda x)=\lambda f(x)$ pour $\lambda \in \Rr$.
  \end{itemize}
\end{enumerate}
\fincorrection
\correction{002364}
\begin{enumerate}
  \item Pour tout $x$, $\|S(x)\| = \| x \|$ donc $\| S \| = 1$.
  
  \item $\| T(f) \|_\infty = \| f\times g \|_\infty \le \|f \|_\infty 
\|g \|_\infty$. Donc pour $f\neq 0$, $\frac{\| T(f) \|_\infty }{\|f \|_\infty 
} \le \|g \|_\infty$. De plus en $g$, on obtient   $\frac{\| T(g) \|_\infty }{\|g \|_\infty  }
=\frac{\| g^2 \|_\infty }{\|g \|_\infty  }  = \|g \|_\infty$. Donc $\| T \| = \|g \|_\infty$.

  \item On a $| u(f) | \le \| f\|_\infty \int_0^1|g(x)|dx$ donc
$\| u\| \le \int_0^1|g(x)|dx$. Si $g$ ne change pas de signe sur $[0,1]$
alors pour $f$ la fonction constant égale à $1$, on obtient $| u(f) | = \| f\|_\infty \int_0^1|g(x)|dx$ donc $\| u \| = \int_0^1|g(x)|dx$.
Si $g$ change de signe alors il ne le fait qu'une fois et en $\frac 12$.
Soit $h_n$ la fonction définie par $h_n(x) = 1$ si $x\in[0,\frac 12 -\frac 1n]$,
$h_n(x) = -1$ si $x\in[\frac 12 +\frac 1n,1]$ et $h_n$ est affine sur $[\frac 12 -\frac 1n,\frac 12 +\frac 1n]$  et continue 
sur $[0,1]$. Cette fonction est construite de telle sorte que si $g$ est positive puis négative alors $h_n\times g$ est une fonction continue qui converge uniformément vers $|g|$ : $\| h_n g - |g| \|_\infty \rightarrow 0$.
Donc $| u(h_n) | = \int_0^1 h_n\times g$ 
et par la convergence uniforme alors $| u(h_n) |$ converge vers
$\int_0^1|g|$. Donc $\|u\| = \int_0^1|g|$.

  \item $|u(x)| = |\sum a_nx_n| \le \|a_n \|_2\| x_n\|_2$ (c'est Cauchy-Schwartz)
donc $\| u \| \le \|a_n \|_2$. Pour la suite $x=a$ on a égalité d'o\`u
$\| u\| = \|a_n \|_2$.
  
  \item $|u(x)| = |\sum a_nx_n| \le \sum |a_nx_n| \le \|a \|_\infty \sum |x_n| = \|a \|_\infty \| x_n\|_1$, donc $\| u \| \le \|a \|_\infty$. 
Soit $p$ fixé, soit $i(p)$ un indice tel que $|a_{i(p)}| = \max_{j=1,\ldots,p} |a_j|$.
On construit une suite $x^p$ de la manière suivante : 
$x^p = (0,0,\ldots,0,a_{i(p)},0,0,0\ldots)$ (des zéros partout sauf $a_{i(p)}$ à la place $i(p)$). Alors $\| x^p \|_1 = |a_{i(p)}|$ et $|u(x^p)| = a_{i(p)}^2$.
Donc $\frac{|u(x^p)|}{\| x^p \|_1} = | a_{i(p)}|$.
Lorsque $p$ tend vers $+\infty$, $| a_{i(p)}| \rightarrow \|a \|_\infty$.
Donc $\| u \| = \|a \|_\infty$.
  
  \item $|u(x)| = |\lim x_n| \le \| x \|_\infty$, donc $\| u \| \le 1$.
Pour $x=(1,1,1,\ldots)$ on obtient l'égalité $\| u \| = 1$.
\end{enumerate}
\fincorrection
\correction{002365}
\begin{enumerate}
  \item Il suffit de l'écrire...
  
  \item Calculons la norme de $U$ : $\| U(P) \| = \sup_k  |\frac 1k a_k\| \le \sup_k |a_k| \le \| P \|$. Donc pour tout $P$, $\frac{\| U(P) \|}{\| P \|} \le 1$.
Et pour $P(x) = x$ on a égalité donc $\| U \|=1$.

  \item Pour $V$, prenons $P_k(x) = x^k$, alors $\| P_k \|=1$, mais $\|V(P_k)\| = k$. Donc $V$ n'est pas bornée sur la boule unité donc $V$ n'est pas continue.
\end{enumerate}
\fincorrection
\correction{002366}
\begin{enumerate}
  \item $A$ injective : Si $A(x_1,x_2,\ldots) = A(y_1,y_2,\ldots)$ alors
$(x_1, x_2/2,..., x_n/n,...) = (y_1, y_2/2,..., y_n/n,...)$ donc
$x_1=y_1$, $x_2=y_2$,..., $x_n=y_n$,... Donc $A$ est injective.

$A$ continue : $\| A(x) \|_\infty = \sup_n \frac {x_n}{n} \le \sup_n x_n  \le \| x \|_\infty$. Donc $\| A \| \le 1$ donc $A$ est continue.

Norme de $A$ : Pour $x=(1,0,0,\ldots)$. On a $\|x\|_\infty = 1$ et $\|A(x)\|_\infty = 1$
Donc la norme de $A$ est exactement $1$.

$A$ n'est pas surjective : posons $y = (1,1,1,\ldots) \in l^\infty$. 
Soit $x$ une suite telle que $A(x)=y$ alors $x=(1,2,3,4,\ldots)$.
Mais $\|x\|_\infty = +\infty$ donc $x \notin l^\infty$. En conséquence
$A : l^\infty \rightarrow l^\infty$ n'est pas surjective.
  
  \item L'inverse \`a gauche de $A$ est $B$ définie par 
$$B(x_1,x_2,...,x_n,...) = (x_1, 2x_2,..., nx_n,...)$$
de sorte que pour $x\in l^\infty$ on ait $B\circ A(x) = x$.
Posons la suite $x^p = (0,0,\ldots,0,1,0,0\ldots) \in l^\infty$
(des zéros partout et le $1$ à la $p$-ième place). Alors $\| x^p \|_\infty = 1$
et $\| B(x^p)\|_\infty = p$. Donc $\frac{\| B(x^p)\|_\infty }{\| x^p \|_\infty } = k$, donc la norme de $B$ n'est pas finie et $B$ n'est pas continue.
\end{enumerate}
\fincorrection
\correction{002367}
\begin{enumerate}
  \item Si $L(a) = 0$ alors $a\in H$ donc $\mathrm{dist}(a,H) = 0$ donc la relation est vraie. Supposons que $L(a) \neq 0$. Alors on a $X = H+\Rr.a$.
En effet pour $x\in X$, il existe $\lambda \in \Rr$ tel que $L(x)=\lambda L(a)$. 
Donc $L(x-\lambda a)=0$. Posons $h= x-\lambda a$, alors $h\in H$ et $x = h+\lambda a$ est la décomposition suivant $H+\Rr.a$.

Si $L$ est continue alors $\| L\|$ est finie.

\begin{align*}
 \| L \| &= \sup_{x\in X, x \neq 0} \frac{\| L(x) \|}{\|x\|} \\
         &= \sup_{h\in H, \lambda \in \Rr, h+\lambda a \neq0} \frac{\| L(h+\lambda a) \|}{\|h + \lambda a\|} \\
         &= |L(a)|\sup_{h\in H, \lambda \in \Rr, h+\lambda a \neq0} \frac{|\lambda|}{\|h + \lambda a\|} \\
         &= |L(a)|\sup_{h\in H} \frac{1}{\|h + a\|} \\
         &= |L(a)|\frac{1}{\inf_{h\in H} {\|h + a\|}} \\
         &= |L(a)|\frac{1}{\mathrm{dist}(a,H)} \\
\end{align*}

Ce qui était l'égalité demandée.

  \item Si $H$ est fermé alors $\mathrm{dist}(a,H)>0$ si $a\notin H$ (voir les exercices sur les compacts), par l'égalité démontrée ci-dessus on a $\| L \|$ finie
donc $L$ est continue.

  \item Soit $X = \Rr[x]$. Pour $P(x) = \sum _{k=0}^p a_k x^k$
on pose $\|P\|= \sup_k |a_k|$, et $V(P)(x) = \sum _{k=1}^n k a_k x^k$.
Alors $\mathrm{Ker}\, V = \{0\}$ est fermé  mais $V$ n'est pas continue
(voir l'exercice \ref{exopol}).
\end{enumerate}
\fincorrection
\correction{002368}
Notons $L : X \rightarrow \Rr$ l'application linéaire définie par
$L(f) = f(0)$. Prenons $f_n$ définie par
$f_n(t) = 2n(1-nt)$ pour $t\in[0,\frac 1n]$ et $f(t)=0$ si $t>\frac 1n$.
Alors $\|f_n\| = 1$ alors que $L(f_n) = 2n$. Donc le rapport
$\frac{|L(f_n)|}{\| f_n \|} = 2n$ n'est pas borné, donc $L$ n'est pas continue.
Si $H = \{ f \mid f(0)=0 \}$ alors $H= \mathrm{Ker}\, L = L^{-1}(0)$. Comme $L$ n'est pas continue alors $H$ n'est pas fermé (voir l'exercice \ref{exoferm}).
\fincorrection
\correction{002369}
$N$ est bien une norme. Et on a pour tout $x$, $(1+x^2)|f(x)| \le N(f)$.
$$|L(f)| = |\int_\Rr f | \le \int_\Rr |f| \le  \int_\Rr \frac{N(f)}{1+x^2}dx 
\le N(f) \int_\Rr \frac{1}{1+x^2} = N(f)[\Arctan x]_{-\infty}^{+\infty} = N(f)\pi.$$
Donc pour tout $f$ on a 
$$\frac{\int f}{N(f)} \le \pi.$$
De plus pour $f(x) = \frac{1}{1+x^2}$ on obtient l'égalité. Donc la norme $\|L\|$ de l'application $L$ est $\pi$.
\fincorrection
\nocorrection
\nocorrection
\nocorrection
\correction{006207}
\begin{enumerate}
\item Remarquons tout d'abord que si $0\in K$, l'application $u$ (linéaire) admet
$0$ comme point fixe dans $K$.  

On suppose donc $0\notin K$.
Si $x\in K$, $u^j(x)\in K$ pout tout $0\leq j\leq n-1$ et $S_n(x)\in K$ comme
combinaison convexe de ces $n$ points de $K$.

\item Soit $x\in K$. Par 1., $S_{n_1}\circ\cdots\circ S_{n_k}(x)\in K$ et s'écrit
$S_{n_1}(S_{n_2}\circ\cdots\circ S_{n_k})(x)$ : il appartient donc à
$S_{n_1}(K)$. Par ailleurs, comme les applications linéaires $S_{n_j}$ commutent
entre elles (ce sont des polyn\^omes en $u$), il en va de même pour
$S_{n_2},\cdots, S_{n_k}$ d'où l'inclusion. 

Chaque $S_n$ étant continue et $K$ compact, les ensembles $S_{n}(K)$ sont eux
aussi compacts et inclus dans $K$; si $A=\cap_{n\geq1} S_{n}(K)$ était
vide, par la propriété de l'intersection finie, on pourrait trouver des
entiers $n_1,n_2,\ldots,n_k$ en nombre fini tels que $S_{n_1}(K)\cap
S_{n_2}(K)\cap\cdots\cap S_{n_k}(K)=\emptyset$. Or cette intersection contient
l'image de $K$ par $S_{n_1}\circ\cdots\circ S_{n_k}$ et ne peut être vide.

\item Soit $a\in A$; pour tout $n$ il existe $x_n\in K$ tel que $a=S_n(x_n)$.
On va montrer que  $u(a)=a$ :

$$\begin{array}{ccc}
         u(a)-a&=& u(S_n(x_n))-S_n(x_n)\\
&&\\

        &=&\displaystyle{(n+1)S_{n+1}(x_n)\over n}-{x_n\over n}-S_n(x_n)
\\
&&\\

&=&\displaystyle{(n+1)S_{n+1}(x_n)-nS_n(x_n)-x_n\over n}\\
&&\\
&=&\displaystyle{u^n(x_n)-x_n\over n}
\end{array}$$

Mais $||u^n(x_n)-x_n||\leq \hbox{diam} K=d<+\infty$ puisque $K$ compact est
borné, et 
$||u(a)-a||\leq {d\over n}$ pour tout $n$ est donc nul.
\end{enumerate}
\fincorrection
\nocorrection
\nocorrection
\nocorrection
\nocorrection
\nocorrection
\nocorrection
\nocorrection
\nocorrection
\nocorrection
\nocorrection
\nocorrection
\nocorrection
\nocorrection
\nocorrection
\nocorrection
\nocorrection
\nocorrection
\nocorrection
\nocorrection
\nocorrection
\nocorrection
\nocorrection
\nocorrection
\nocorrection
\nocorrection
\nocorrection
\nocorrection
\nocorrection
\nocorrection
\nocorrection
\nocorrection
\nocorrection
\nocorrection
\correction{001900}
\ \par
\begin{center}{\large \bf I}
\end{center}
\vskip5mm


\begin{enumerate}
\item Soit $\mathcal{P}$ l'espace vectoriel des fonctions polynomiales. 
Supposons $\mathcal{P}$ de dimension finie $n$. Notons $f_k$ la fonction 
$x\mapsto x^k$. Alors la famille $\{ f_0, \cdots ,
f_n \}$ qui compte $n+1$ \'el\'ements est li\'ee, donc il existe 
$a_0, \cdots , a_n$ des scalaires non tous nuls tels que, pour tout
$x\in \Rr $ on ait $a_0+a_1x+\cdots a_n x^n=0.$ Il en r\'esulte que 
le polyn\^ome non nul \`a coefficients r\'eels $a_0+a_1X+\cdots 
a_nX^n$ a une infinit\'e de racines, ce qui est
absurde.

\item Posons $M=\sup (\bar{X}).$ On doit v\'erifier que, {\em i)} 
pour tout $x\in X , x\leq M$ et {\em ii)} pour tout $\epsilon >0$ 
il existe $x\in X$ tel que
$M-\epsilon \leq x.$
Comme $X\subset \bar{X}$ et, pour tout $x\in \bar{X} , x\leq M$ la 
propri\'et\'e {\em i)} est v\'erifi\'ee par $M.$ Soit maintenant 
$\epsilon >0$. Il existe
$x\in \bar{X}$ tel que $M-\dfrac{\epsilon }{2}<x$. Comme 
$x\in\bar{X},$ il existe aussi $y\in X$ tel que $\vert x-y\vert 
<\dfrac{\epsilon }{2}$. Donc $ M-\epsilon< y$
et $M$ satisfait \`a {\em ii)}.
\vskip1mm
{\em Remarque : } on note \'egalement que $\sup (X)\in \bar{X}$. En 
effet, pour tout $n\in \Nn ,$ choisissons un \'el\'ement $x_n\in X$ 
tel que $x_n\geq \sup (x)-\dfrac{1}{n}.$
Alors la suite $(x_n)_{n\in \Nn}$ constitu\'ee d'\'el\'ements de $X$ 
converge dans $\Rr$ vers $\sup (X)$ qui appartient donc \`a $ 
\bar{X}.$ On peut bien s\^ur en d\'eduire la
propri\'et\'e {\em ii)} de $M$.

\end{enumerate}

\vskip5mm
\begin{center}{\large \bf II}
\end{center}
\vskip5mm

\begin{enumerate}
\item Il est clair $\mathcal{L} $ est un sous espace vectoriel de l'espace 
vectoriel des fonctions de $[0,1]$ \`a valeurs dans $\Rr$.
  Soit $f\in \C ^1$ et $x,y \in
[0,1],$ avec $x< y$.
Par le th\'eor\'eme des accroissements finis, il existe $c_x\in ]x,y[ 
$ tel que $f(y)-f(x)=f'(c_x)(y-x)$. Or $f'$ est continue, donc 
born\'ee sur $[0,1]$.
Soit $\displaystyle{M=\sup _{t\in [0,1]}\vert f'(t)\vert .}$ On a 
l'in\'egalit\'e $\vert f(y)-f(x)\vert \leq M\vert y-x\vert $ qui 
montre que $f\in \mathcal{L} .$
Il en r\'esulte que $\mathcal{L} $ contient $\mathcal{P} $ donc est de dimension infinie.



\item
\begin{enumerate}
\item Il suffit de v\'erifier que si $N_1(f)=0$ et $N_2(g)=0$, alors 
$f=g=0$, les autres propri\'et\'es \'etant claires.
Or si $N_1(f)=0,$ alors $f$ est constante et $f(0)=0,$ donc $f=0.$ Il 
en va de m\^eme pour $N_2$.

\item Pour tout $n \in \Nn $, $\Vert f_n \Vert _{\infty}=1$. Posons 
$\displaystyle{X_n =\{ \dfrac {\vert f_n(x)-f_n(0)\vert }{\vert x\vert
}} , x\not =0 \}  .$ Comme $f_n(0)=0, $ on voit que $N_1(f_n)=\sup (X_n).$ Or
  $\displaystyle{ \vert f'_n(0)\vert =\lim _{x\rightarrow 0}
\dfrac {\vert f_n(x)-f_n(0)\vert }{\vert x\vert
}}, $ appartient \`a $\bar{X_n}$ donc, en appliquant I 2) on constate 
que $\vert f'_n(0)\vert \leq \sup (\bar {X_n})= \sup (X_n)$. Enfin 
$f'_n(0)=2\pi n$ donc
$N_2(f_n)\geq 2\pi n .$ Il n'exite donc pas $K>0$ tel que, pour tout 
$n\in \Nn ,$ $N_2(f_n)<K\Vert f_n \Vert _{\infty}$ soit $N_2$ et 
$\Vert \  \Vert _{\infty}$ ne sont pas
\'equivalentes.
\vskip1mm {\em Remarques :} a) on peut obtenir ce r\'esultat (et le 
pr\'eciser) en remarquant que la fonction $f_n :x\mapsto \dfrac{\sin 
(2\pi n x)}{x}$ d\'efinie sur $]0, 1]$ se
prolonge en une fonction continue en $0$ en posant $f_n(0)=2\pi n.$ 
Puis noter (en fait c'est un cas particulier de I 2)) que 
$\displaystyle{ \sup _{]0,1]}\vert f_n\vert =  \sup
_{[0,1]}\vert f_n\vert }$ et montrer (par une \'etude classique de 
fonction) que cette derni\`ere quantit\'e est $2\pi n .$
\vskip1mm
b) Ce qui fait l'int\'er\^et pour ce probl\`eme des fonctions 
$(f_n)_{n\in \Nn }$, c'est qu'elles sont born\'ees par $1$ mais que 
leur pente en l'origine peut-\^etre rendue
arbitrairement grande avec $n.$ On peut donc obtenir le m\^eme 
r\'esultat avec la suite $(k_n)_{n\in \Nn }$ d\'efinie par $k_n(x)=nx 
$ si $x\leq \dfrac{1}{n}$ et $1$ sinon,
pour laquelle un calcul direct donne $N_1(k_n)=n$ et $\Vert k_n\Vert 
_{\infty}=1.$

\item Comme, pour tout $f\in \mathcal{L} $, $N_1(f)\geq N_2(f)$, on d\'eduit 
de ce qui pr\'ec\'ede que $N_1$ n'est pas \'equivalente ni \`a  $
\Vert \ \Vert _{\infty}$. Posons $g_n(x)=x^n,$ pour $n\geq 1.$ Pour 
tout $n\geq 1,$ $N_2(g_n)=1$. De plus $g'_n(1)=n,$ donc, par un 
raisonnement identique \`a celui qui pr\'ec\'ede,
$N_1(g_n)\geq n $ ce qui montre que $N_1$ n'est pas \'equivalente \`a $N_2$.

\vskip1mm {\em Remarque :} ce qui fait l'int\'er\^et pour ce 
probl\`eme des fonctions $(g_n)_{n\in \Nn }$, c'est qu'elles sont 
born\'ees par $1$ mais que leur pente en $1$
peut-\^etre rendue arbitrairement grande avec $n.$ On peut donc 
obtenir le m\^eme r\'esultat avec la suite $(l_n)_{n\in \Nn }$ 
d\'efinie par $l_n(x)=0 $ si $x\leq 1-\dfrac{1}{n}$ et
  $nx- (n-1)$
sinon.

\item On pose $g_n(x)=x$ si $x\leq \dfrac{1}{n}$ et $\dfrac{1}{n}$ 
sinon. Il est clair que $g_n\in \mathcal{L} ,$ $ \Vert g_n \Vert 
_{\infty}=\dfrac{1}{n}$ et
$N_2(g_n)=1.$ Il n'existe donc pas de constante $K'\in \R _+^*$ telle 
que, pour tout $n\in \Nn$, $ \Vert g_n \Vert _{\infty}\geq 
K'N_2(g_n)$
donc $ N_2$ n'est pas \'equivalente \`a $
\Vert \ \Vert _{\infty}.$ Enfin $N_2(g_n)=N_1(g_n)$, ce qui \'etablit 
le m\^eme r\'esultat pour $N_1$.




\item Il est clair que pour tout $f\in \mathcal{L}, $ que $\lambda (f)\geq 
N_1(f).$ Soit $x\in ]0,1]$. De l'identit\'e
$f(x)=f(0)+x\dfrac{f(x)-f(0)}{x}$ on d\'eduit que $\vert f(x)\vert 
\leq \vert f(0)\vert +\vert \dfrac{f(x)-f(0)}{x}\vert $
(car $\vert x\vert \leq 1$) soit $\vert f(x)\Vert \leq  N_1(f).$
L'application $x\mapsto f(x)$ \'etant continue sur $[0,1]$ (ou en 
appliquant I 2)) on en d\'eduit que, pour tout $x\in [0,1]$ on a 
\'egalement $\vert f(x)\Vert \leq  N_1(f).$
En d'autre termes $\Vert f \vert _{\infty}\leq  N_1(f)$ et $\lambda 
(f)\leq 2N_1(f).$ Les normes $\lambda $ et $N_1$ sont donc 
\'equivalentes.
\end{enumerate}

\item On pose, pour tout $f\in \C ^1 :$
$\displaystyle{ \nu _1(f)=\vert f(0)\vert +\Vert f'\Vert _{\infty}}$
et $\displaystyle{ \nu (f)=\Vert f\Vert _{\infty} +\Vert f'\Vert _{\infty}}.$
\begin{enumerate}
\item On constate ais\'ement que si $\nu _1(f)=0,$ alors $f$ est 
constante et, comme de plus $f(0)=0$, elle est nulle. Les autre 
propri\'et\'es sont imm\'ediates, donc $\nu $ et $\nu
_1$ sont des normes sur
$\C ^1$.
\item  Soit $f\in \C ^1,$ notons $X=\{ \dfrac {\vert f(x)-f(y)\vert 
}{\vert x-y\vert } ;
(x,y)\in
[0,1]^2,  x\not = y \} .$  Pour montrer que $\nu _1
(f)= N_1(f) ,$ il suffit de v\'erifier que $\sup(X)=\Vert f'\Vert 
_{\infty}.$ Soient $x, y\in [0,1], x\not = y. $ Par le th\'eor\`eme 
des accroissements finis,
il existe $c$ compris entre $x$ et $y$ tel que $\dfrac {\vert 
f(x)-f(y)\vert }{\vert x-y\vert } =f'(c)\leq \Vert f'\Vert 
_{\infty},$ donc $
\sup (X)\leq \Vert f'\Vert _{\infty}.$ Comme $f'$ est continue, il 
existe $x_0\in [0, 1]$ tel que $f'(x_0)=\Vert f'\Vert
_{\infty}.$  Alors $\displaystyle{f'(x_0) =\lim _{x\rightarrow 
x_0}\dfrac{f(x)-f(x_0)}{x-x_0}}$ appartient \`a $\bar{X},$ donc, en 
appliquant I 2), $\Vert f'\Vert _{\infty}\leq \sup
(\bar{X}) =\sup (X).$
  \vskip1mm
{\em Remarque :} on peut formuler ce raisonnement de la mani\`ere 
suivante~: soit $Y=\{ f'(x) ; x\in [0,1] \} .$ Par le th\'eor\`eme 
des accroissements finis,
$X\subset Y.$ On a ensuite, par d\'efinition de la d\'eriv\'ee, 
$Y\subset \bar{X}.$ Donc $\sup (X) \leq \sup (Y)\leq \sup (\bar{X}), 
$ puis on applique I 2).
\item Les normes $\nu $ et $\nu _1$ sont \'equivalentes. En effet, il 
est clair que $\nu _1(f)\leq \nu (f)$ pour tout $f\in \C ^1.$
Soit $t_0\in [0,1]$ tel que $\Vert f\Vert _{\infty}=\vert f(t_0)\vert 
.$ Si $t_0=0$ alors $\nu _1(f)\leq \nu (f)$. Sinon, par le 
th\'eor\`eme des accroissements finis, il existe $c\in
]0, t_0[$ tel que
$f(t_0)= f(0)+ f'(c)t_0$ ce dont on d\'eduit que $\Vert f\Vert 
_{\infty}\leq \nu _1(f)$, puis que $\nu ( f)\leq 2\nu _1(f).$
\end{enumerate}
\item \begin{enumerate}
\item Soit $x\in [0,1]$. La suite de nombres r\'eel $(f_n(x))_{n\in 
\Nn}$ \'etant de Cauchy, elle est convergente. On pose 
$\displaystyle{f(x)=\lim _{n\rightarrow \infty}f_n(x)}$.
Soit $\epsilon >0$. La suite $(f_n)$ \'etant de Cauchy, il existe 
$N$ tel que, si $m,n\geq N$ alors $\Vert f_n-f_m \Vert _{\infty}\leq 
\epsilon $.
Soient $x\in [0,1]
$ et $m,n\geq N.$ On a $\vert f_n(x)-f_m(x)\vert \leq \epsilon$ 
et, ceci \'etant vrai pour tout $m \in \Nn ,$ on en d\'eduit, par 
passage \`a la limite suivant $m$, que $\vert
f_n(x)-f(x)\vert
\leq \epsilon ,$ soit $\Vert f_n-f\Vert _{\infty}\leq \epsilon 
$. Ainsi $f$ est la limite, pour la convergence uniforme, d'une suite 
de fonctions continues donc est continue.



\item Par d\'efinition de $\nu ,$ une suite $(f_n)_{n\in \Nn}$ de 
Cauchy pour $\nu$ est de Cauchy pour $\Vert \ \Vert _{\infty}$, donc 
(uniform\'ement) convergente par la question
qui pr\'ec\'ede. De m\^eme $(f'_n)_{n\in \Nn}$ est de Cauchy pour 
$\Vert \ \Vert _{\infty}$, donc converge uniform\`ement vers une 
fonction continue $g$. Il en r\'esulte que $f$ est
d\'erivable et a pour d\'eriv\'ee la fonction continue $g$. Enfin 
$(f_n)_{n\in \Nn}$ converge vers $f$ pour $\nu$ donc
$(\C^1,
\nu )$ est complet.

Soit maintenant $(g_n)_{n\in \Nn}$ une suite de Cauchy dans $(\C^1, 
\nu _1)$. Comme $\nu _1$ est \'equivalente \`a $\nu ,$ elle est de 
Cauchy pour $\nu $ donc
convergente. Il existe donc $h\in \C ^1 $ telle que 
$\displaystyle{\lim _{n\rightarrow \infty}\nu (h-g_n )=0}$. Mais 
puisque $\nu _1$ est \'equivalente \`a $\nu ,$
on a aussi $\displaystyle{\lim _{n\rightarrow \infty}\nu _1(h-g_n 
)=0}$ donc $(\C^1, \nu _1)$ est complet.

\item Soit $(f_n)_{n\in \Nn }$ une suite de Cauchy dans $(\mathcal{L} , 
\lambda ).$ Comme $\lambda (f_n)\geq \Vert f_n \Vert _{\infty}$, la 
suite $(f_n)_{n\in \Nn }$ est \'egalement
de Cauchy dans $(\C ^0, \Vert \ \Vert _{\infty}).$ Comme $(\C ^0, 
\Vert \ \Vert _{\infty})$
est complet, $(f_n)_{n\in \Nn }$
converge uniform\'ement vers une fonction continue qu'on notera $f.$

\item Soit $\epsilon >0 . $ Comme $(f_n)_{n\in \Nn }$ est une 
suite de Cauchy, il existe $N$ tel que, pour $m,n\geq N$ on ait, pour 
tout $x,y$ et $z\in [0,1]$, avec $x\not =y$~:
$$\vert f_n(z)-f_m(z)\vert +\left \vert \dfrac{ 
(f_n(x)-f_m(x))-(f_n(y)-f_m(y))}{x-y}\right \vert <\epsilon .$$
En faisant tendre $m$ vers l'infini, on en d\'eduit~:
$$\vert f_n(z)-f(z)\vert +\left \vert \dfrac{ 
(f_n(x)-f(x))-(f_n(y)-f(y))}{x-y}\right \vert <\epsilon ,$$
donc
$$\sup _{z\in [0,1]}\vert f_n(z)-f(z)\vert +\sup _{x,y\in [0,1], 
x\not =y}\left \vert \dfrac{ (f_n(x)-f(x))-(f_n(y)-f(y))}{x-y}\right 
\vert \leq \epsilon .$$
Ainsi la suite $(f_n)_{n\in \Nn }$ converge vers $f$ pour la norme 
$\lambda .$ Par ailleurs, on d\'eduit de la seconde in\'egalit\'e 
que, pour tout
$x,y \in [0,1]$, avec $x\not =y$~:
  $\vert  (f_n-f)(x)-(f_n-f)(y)\vert <\epsilon \vert x-y\vert ,$ 
donc que pour $n$ assez grand $f-f_n$ appartient \`a $\mathcal{L} .$ Or $\mathcal{L} $ 
est un espace vectoriel et $f_n\in \mathcal{L} $
donc $f$ appartient \`a $\mathcal{L} $.


\item Toute suite de Cauchy de $\mathcal{L} $ admet une limite dans $\mathcal{L}$ qui 
est donc complet.
\end{enumerate}
\end{enumerate}
\fincorrection
\correction{006055}
\begin{enumerate}
\item) La nécessité des conditions (i), (ii), (iii) résulte des propriétés d'une
norme. En effet, si $x\in K$, $||-x||=||x||\leq1$ et $-x\in K$ ce qui prouve
(i).

$K$ est fermé car plus généralement toute boule fermée d'un espace métrique est
fermée; et $K$ est borné par définition (un sous-ensemble de $\R^n$ est borné
s'il est contenu dans une boule fermée). Enfin, $K$ est convexe car, si
$x,y\in K$ et $\lambda\in[0,1]$, $||(1-\lambda)x+\lambda y||\leq
(1-\lambda)||x||+\lambda ||y||\leq 1$ et $(1-\lambda)x+\lambda y\in K$.

$0$ est un point intérieur à $K$ : par exemple $B(0,{1\over2})\subset K$ et
$K$ contient un voisinage de $0$ (toute boule fermée de rayon $>0$ est un
voisinage de son centre).

\item Soit $K$ vérifiant les propriétés (i), (ii), (iii). Il nous faut montrer que
$p(x)$ est bien définie pour tout
$x$; que
$p$ est une norme et que
$K$ est la boule unité fermée qui lui est associée.

Si $x=0$, ${x\over a}\in K$ pour tout $a>0$ et $p(0)=0$. Si $x\not=0$,
l'ensemble
$\{a>0\ ;\ {x\over a}\in K\}$ est minoré; s'il est non vide il admettra une
borne inférieure. Or $0$ est point intérieur à $K$; il existe donc $\varepsilon
>0$ tel que $B(0,\varepsilon)\subset K$ et pour $a$ assez grand, ${||x||\over
a}\leq \varepsilon$, en particulier ${x\over a}\in K$, et l'ensemble est non
vide.

\smallskip

Vérifions les trois axiomes d'une norme. 

$\bullet$ Par définition d'une borne inférieure, $p(x)=0\Longrightarrow
\forall\varepsilon>0 \ \exists\ 0<a< \varepsilon$ tel que ${x\over a}\in K$. $K$
étant borné, on peut supposer que $K\subset \overline B(0,R)$, de sorte que
$||{x\over a}||\leq R$ ou $||x||\leq \varepsilon R$, ceci pour tout
$\varepsilon>0$, ce qui implique $x=0$.

\smallskip 
 
$\bullet$ Soit $\lambda>0$; $p(\lambda x)=
\inf\{a>0\ ;\ {\lambda x\over a}\in K\}=\inf\{\lambda b>0\ ;\ {x\over
b}\in K\}$ en posant $a=\lambda b$, et $p(\lambda x)=\lambda p(x)$.

Il suffit de montrer que $p(- x)=p(x)$ pour avoir la propriété d'homogénéité.
Mais 
$p(- x)=\inf\{a>0\ ;\ -{ x\over a}\in K\}=\inf\{a>0\ ;\ { x\over a}\in
-K\}=p(x)$ car $K$ est symétrique.

\smallskip

$\bullet$ En utilisant la définition d'une borne inférieure, on va montrer que
pour tout
$\varepsilon>0$,
$p(x+y)\leq p(x)+p(y)+2\varepsilon$ ce qui donnera le résultat.

Donc, fixons $\varepsilon>0$; on peut trouver $a>0$ tel que $p(x)\leq
a<p(x)+\varepsilon$ et ${ x\over a}\in K$, puis $b>0$ tel que $p(y)\leq
b<p(y)+\varepsilon$ et ${y\over b}\in K$. Si ${ x+y\over a+b}\in K$, alors
$p(x+y)\leq a+b$ par propriété de la borne inf et on aura prouvé $p(x+y)\leq
p(x)+p(y)+2\varepsilon$.

Mais ${ x+y\over a+b}$ s'écrit ${ a\over a+b}{x\over a}+{
b\over a+b}{y\over b}$,  combinaison convexe de ${ x\over a}$ et ${ y\over b}$,
et $K$ est supposé convexe. La preuve de l'inégalité triangulaire est ainsi
achevée.

\smallskip

Il nous reste à établir $K=\{x\ ;\ p(x)\leq 1\}$. 

Si $x\in K$ et $a=1$, ${ x\over
a}\in K$ ce qui implique $p(x)\leq 1$. 
 Réciproquement supposons $ p(x)\leq 1$;
on peut supposer $x\not=0$. Si $p(x)<1$, il existe $p(x)\leq a<1$ tel que
${x\over a}\in K$; mais $x=a{x\over a}+(1-a)0$ est encore dans $K$.

Si $p(x)=1$ il existe $(a_n)$ suite de nombres positifs tels
que 
${x\over a_n}\in K$ pour tout $n$ et tendant vers $1$.  Mais $K$ étant fermé,
$x=\lim{x\over a_n}\in K$.
\end{enumerate}
\fincorrection
\nocorrection
\nocorrection
\nocorrection
\nocorrection
\nocorrection
\nocorrection
\correction{002395}
\begin{enumerate}
  \item 
Soit $(u_n)$ une suite de Cauchy pour $d$.
Donc
$$\forall \epsilon >0 \quad \exists N\in \Nn \quad \forall p,q \ge N \quad
d(u_p,u_q) = |u_p^3-u_q^3| \le \epsilon.$$
Donc la suite $(u_n^3)$ est une suite de Cauchy pour la distance usuelle $|.|$.
Comme  $(\Rr,|.|)$ est complet alors $(u_n^3)$ converge pour la valeur absolue, notons $v$ la limite,
 nous avons $|u_n^3-v|$ qui tend vers $0$. Donc pour $u=v^{\frac13}$
nous avons $d(u_n,u)=|u_n^3-u^3|=|u_n^3-v|$ qui tend vers $0$, donc
$u_n$ converge vers $u$ pour la distance $d$.
Donc $\Rr$ est complet pour $d$.


  \item Montrons que $d$ ne définit pas une distance complète.
Soit $(u_n)$ la suite définie par $u_n=-n$, $n\in \Nn$.
Alors $d(u_p,u_q) = |e^{-p}-e^{-q}|$.
Donc pour $\epsilon>0$ fixé, soit $N$ tel que $e^{-N}<\frac\epsilon 2$,
alors pour $p,q\ge N$ on a $d(u_p,u_q) = |e^{-p}-e^{-q}| \le e^{-p}+
 e^{-q} \le 2e^{-N} \le \epsilon$. Donc $(u_n)$ est de Cauchy.
Supposons que $(u_n)$ converge, notons $u\in \Rr$ sa limite.
Alors $d(u_n,u)=|e^{-n}-e^u|$ tend vers $0$ d'une part
et vers $e^u$ d'autre part. Donc $e^u=0$ ce qui est absurde pour
$u\in \Rr$.

  \item La fonction $\ln(1+u)$ est continue et ne s'annule qu'en $u=0$. 
Donc pour $\ln(1+u)$ suffisamment petit nous avons $u$ suffisamment petit et donc
(par la relation $\ln(1+u) = u +o(u)$) nous avons 
$$\frac12 u \le \ln(1+u) \le 2u.$$
Donc pour $(u_n)$ une suite de Cauchy pour $d$, la première inégalité prouve 
que $(u_n)$ est une suite de Cauchy pour $|.|$. Donc elle converge pour $|.|$
La deuxième inégalité montre que $(u_n)$ converge pour 
$d$. Donc $d$ définit une distance complète.
\end{enumerate}
\fincorrection
\correction{002396}
$f$ est injective afin que $d$ soit bien une distance. On pose $F= f(\Rr)\subset \Rr^2$.
\begin{enumerate}
  \item $\Rightarrow$ Supposons que la distance $d$ soit complète. Soit $(y_n)$ une suite de $F$ qui converge vers $y\in \Rr^2$. Il faut montrer que $y\in F$.
Il existe $x_n\in\Rr$, tel que $y_n=f(x_n)$. Comme
$(y_n)$ est une suite convergente, c'est une suite de Cauchy de $\Rr^2$, or 
$d(x_p,x_q) = \|f(x_p)-f(x_q)\| = \|y_p-y_q\|$. Donc $(x_n)$ est une suite de Cauchy pour $d$. 
Comme $d$ est complète alors $(x_n)$ converge $x$, 
comme  $x_n\rightarrow x$ 
(pour $d$) alors $f(x_n)\rightarrow f(x)$ (pour $\|.\|$). 
(Remarquons que par définition de $d$, l'application
$f:(\Rr,d)\longrightarrow (\Rr^2,\|.\|)$ est continue.)
Donc $(y_n)$ converge vers $f(x)$ et
par unicité de la limite $f(x)=y$. Donc $y\in f(\Rr)=F$. Donc $F$ est fermé.

  \item $\Leftarrow$ On suppose que $F$ est fermé. Soit $(u_n)$ une suite de Cauchy pour $(\Rr,d)$. Notons $y_n=f(x_n)$. Comme $d(u_p,u_q) = \|f(u_p)-f(u_q)\| = \|y_p-y_q\|$. Donc $(y_n)$ est une suite de Cauchy pour $(\Rr^2,\|.\|)$. Comme cet espace est complet alors $(y_n)$ converge vers $y$.
Comme $y_n\in F$ et $F$ est fermé alors $y\in F$, donc il existe $x\in \Rr$ tel que $y=f(x)$. Il reste à montrer que $(x_n)$ tend vers $x$. En effet
$d(x_n,x)=\|f(x_n)-f(x)\|=\|y_n-y\|$ tend vers $0$. Donc $(x_n)$ tend vers $x$ pour $d$. Donc $d$ est complète.
\end{enumerate}
\fincorrection
\correction{002397}
\begin{enumerate}
  \item 

  \begin{enumerate}
    \item Montrons que $(X,d_\omega)$ est complet. Soit $(f_n)_n$ une suite de Cauchy pour cet distance. Alors pour chaque $t\in[a,b]$, $(f_n(t))_n$ est une suite de Cauchy pour $(\Rr,|.|)$. Comme $\Rr$ est complet alors cette suite converge, notons $f(t)$ sa limite.

Il faut montrer deux choses : premièrement que $(f_n)$ converge vers $f$ pour la distance considérée, deuxièmement que $f$ est bien dans l'espace $X$.

    \item Comme $(f_n)$ est une suite de Cauchy. Pour $\epsilon > 0$. Il existe $n\ge0$ tel que pour tout $p\ge 0$ : $d_\omega(f_n,f_{n+p})  < \epsilon.$
Donc 
$$ \sup_{t\in[a,b]} |\omega (t)
(f_n(t)-f_{n+p}(t))| < \epsilon.$$
On fait tendre $p$ vers $+\infty$ et on obtient :
$\sup_{t\in[a,b]} |\omega (t) (f_n(t)-f(t))| < \epsilon.$
Donc $(f_n)$ converge vers $f$ pour la distance $d_\omega$.

   \item  $\omega$ est une fonction non nulle sur le compact $[a,b]$, donc 
il existe $\alpha >0$ tel que $\omega(t)>\alpha$ pour tout $t\in[a,b]$.
On en déduit que 
$$\|f_n-f\|_\infty \le \frac 1 \alpha d_\omega(f_n,f).$$
Comme $ d_\omega(f_n,f)$ tend vers $0$ alors $f_n$ converge vers $f$ pour la norme infini. Donc $f$ est continue.

Conclusion : $(X,d_\omega)$ est complet.
 \end{enumerate} 

  \item On définit $f_n$ sur $[0,1]$ par $f_n(t) = 1$ pour $t\in[0,\frac 12]$, $f_n(t)=(1-n(t-\frac12))$ pour $t\in[\frac 12,\frac 12+\frac 1n]$
et $f(t)=0$ si $t\ge \frac 12+\frac 1n$. En effet $(f_n)$ converge (au sens de la convergence simple) vers la fonction $f$, o\`u $f$ est définie par $f(t)=1$ sur $[0,\frac 12]$ et $f(t)=0$ sur $]\frac12,1]$.

On montre facilement que $(f_n)$ est une suite de Cauchy pour la norme $\| . \|_1$ en prouvant que pour $p,q\ge n$ on a $\|f_p - f_q\|_1 \le \frac1n$. Par contre $(f_n)$ ne converge pas dans $X$ car $f$ n'est pas continue, donc n'appartient pas à l'espace $X$. Donc $X$ n'est pas complet.  
\end{enumerate}
\fincorrection
\correction{002398}
\begin{enumerate}
  \item On reprend l'exemple de l'exercice \ref{exocomp}. Et on définit
 $g_n$ sur $[0,1]$ par $g_n(x) = \int_0^x f_n(t)dt$. Alors $g_n$ est $\mathcal{C}^1$, et converge (donc en particulier $(g_n)$ est de Cauchy). Elle converge vers $g$ qui n'est pas une fonction $\mathcal{C}^1$. Donc ce n'est pas un espace complet.


  \item Soit $(f_n)$ une suite de Cauchy pour la norme $N$.
Pour chaque $t\in[a,b]$, $(f_n(t))_n$ est une suite de Cauchy de $\Rr$
donc converge. Notons $f(t)$ sa limite. 
De même $(f_n'(t))_n$ est une suite de Cauchy de $\Rr$  donc converge vers $g(t)$.
Nous allons montrer que $f$
est dans $X$ et que $f_n$ converge vers $f$ pour $N$ et que $f'=g$.
Soit $\epsilon > 0$. Il existe $n\in \Nn$ tel que
Pour tout $p\ge 0$, 
$$N(f_{n}-f_{n+p}) < \epsilon.$$
En faisant tendre $p$ vers $+\infty$, $f_{n+p}$ converge (simplement) vers $f$.
On obtient que
$\|f_n-f\|_\infty$ et que $\|f_n'-g\|_\infty$ tendent vers $0$.
Donc $f_n$ converge uniformément vers $f$. Comme les $f_n$ sont continues
alors $f$ est continue. De même $f_n'$ converge uniformément vers $g$
donc $g$ est continue. De plus cela implique que $g=f'$.(Rappel : si $(f_n)$ est une suite de fonctions $\mathcal{C}^1$ sur $[a,b]$ qui converge simplement vers $f$, et tel que $(f_n')$ converge uniformément vers $g$, alors $f$ est $\mathcal{C}^1$ et sa dérivée est $f'=g$.)
Nous avons donc montrer que $N(f_n-f)$ tend vers $0$ et que $f$ est dans $X$.
Donc $(X,N)$ est complet.
 \end{enumerate}
\fincorrection
\correction{002399}
\begin{enumerate}
  \item Notons $x^p$ la suite 
$$x^p = (1,1,\ldots,1,1,0,0,0,\ldots).$$
(des $0$ à partir de la $p+1$-ième place et de $1$ avant.
Si $Y$ est l'espace de toute les suite, notons 
$$x^\infty = (1,1,1,1,\ldots).$$
La suite $x^\infty$ n'est pas dans $X$.
Par contre $x^p \rightarrow x^\infty$ pour la distance $\rho$.
En effet
$$\rho(x^p,x) = \sum_{k=p+1}^{+\infty} \frac{1}{2^k} \frac12 = \frac 1 {2^{p+1}} \rightarrow 0.$$
La suite $(x^p)$ est de Cauchy, mais ne converge pas dans $X$, donc $X$ n'est pas complet. 

  \item 
\begin{enumerate}
    \item
Soit $Y$ l'espace de toutes les suites. Alors $X$ est dense dans
dans $Y$ (pour la topologie définie par $\rho$), car toute suite $y = (y_1,y_2,\ldots)$ de $Y$ 
s'approche par une suite de suite $(x^p)$ obtenue en tronquant la suite $y$ : $x^1 = (y(1),0,0,\ldots)$, $x^2 = (y(1),y(2),0,0,\ldots)$,... 
En effet 
$$\rho(x^p,y)= \sum _{k=p+1}^\infty 2^{-k}
\frac{|x_k-y_k|}{1+|x_k-y_k|} \le \sum _{k=p+1}^\infty 2^{-k} = \frac{1}{2^p}$$
qui tend vers $0$ lorsque $p$ tend vers $+\infty$.

    \item
Soit $(x^n)_n$ une suite de Cauchy de $Y$.
Montrons que pour $k$ fixé alors $(x^n_k)_n$ est une suite de Cauchy de $\Rr$.
Prenons $\epsilon >0$, alors il existe $N$ tel que pour $p,q\ge N$ on ait
$\rho(x^p,x^q) \le \epsilon$.
$$\frac {1}{2^k}\frac{|x^p_k-x^q_k|}{1+|x^p_k-x^q_k|} \le \rho(x^p,x^q) \le \epsilon.$$
Posons la fonction $f(\alpha)=\frac{\alpha}{1+\alpha}$, $f$ est inversible pour $\alpha\ge 0$,
d'inverse $f^{-1}(\beta)= \frac{\beta}{1-\beta}$. Une étude de $f$ et de son inverse
montre que si $f(\alpha)\le \epsilon'\le 1$ alors $\alpha \le 2\epsilon'$.
Comme $k$ est fixé, posons $\epsilon = \frac{\epsilon'}{2^k}$ et $\alpha =|x^p_k-x^q_k|$  on a montrer :
$f(\alpha) \le \epsilon'$. Donc $\alpha \le 2\epsilon'$.
Récapitulons :
$$\forall \epsilon' >0 \quad \exists N \in \Nn \quad \forall p,q \ge N \quad
|x^p_k-x^q_k| < 2\epsilon',$$
donc la suite $(x^n_k)_n$ est de Cauchy dans $\Rr$, donc converge, nous notons $x^\infty_k$ sa limite.

    \item Nous avons construit une suite $x^\infty=(x^\infty_1,x^\infty_2,\ldots)$. Comme $(x^n)$ est de Cauchy alors
$$\forall \epsilon >0 \quad \exists N \in \Nn \quad \forall p,q \ge N \quad
\rho(x^p,x^q) < \epsilon,$$
Lorsque l'on fixe $p$ et que l'on fait tendre $q$ vers $+\infty$ on a
$$\forall \epsilon >0 \quad \exists N \in \Nn \quad \forall p \ge N \quad
\rho(x^p,x^\infty) < \epsilon,$$
donc $(x^n)$ converge vers $x^\infty$ pour la distance $\rho$.

    \item Bien évidemment $x^\infty \in Y$ donc
$(x^n)$ converge vers $x^\infty \in Y$ pour $\rho$. Donc
$(Y,\rho)$ est un espace complet.
  \end{enumerate}

  \item $(X,\|.\|_\infty)$ n'est pas un espace complet. 
Par exemple regardez la suite $x^p= (1,\frac 12,\ldots,\frac1p,0,0,\ldots)$
alors $(x^p)$ est une suite de Cauchy, qui ne converge pas dans $X$,
mais dans $Y$ sa limite est $x^\infty= (1,\frac 12,\frac13,\ldots)$.

Notons $Z$ l'espace des suites qui tendent vers $0$.
L'adhérence de $X$ pour la topologie définie par 
$\|.\|_\infty$ est $Z$. Et $(Z,\|.\|_\infty)$ est complet.
\end{enumerate}
\fincorrection
\correction{002400}
\begin{enumerate}
  \item Soit $(x_n)$ une suite de Cauchy. 
Pour $\epsilon = 1$ il existe $n_0\in \Nn$ tel que
$$\forall q \ge n_0 \qquad \| x_{n_0}-x_q \| < 1.$$
Puis pour $\epsilon = \frac 12$ il existe $n_1 > n_0$ tel que
$$\forall q \ge n_1 \qquad \| x_{n_1}-x_q \| < \frac 12.$$
Puis par récurrence pour $\epsilon = \frac 1{2^k}$, on pose
$n_k > n_{k-1}$ tel que 
$$\forall q \ge n_k \qquad \| x_{n_k}-x_q \| < \frac 1{2^k}.$$

Donc en particulier à chaque étape on a
$$\| x_{n_{k+1}}-x_{n_k} \| < \frac 1 {2^k}.$$
Posons $u_k = x_{n_{k+1}} - x_{n_k}$ 
Alors $\|u_k\| \le \frac 1 {2^k}$ donc
$$\sum_{k\ge 0} \| u_k \| \le \sum_{k\ge 0} \frac 1 {2^k} = 2.$$
Donc la série $\sum_{k\ge 0} u_k$ est normalement convergente.
Si cette série converge notons $T=\sum_{k= 0}^{+\infty} u_k$ sa somme,
C'est-à-dire la limite de $T_N = \sum_{k= 0}^{N} u_k$. 
Mais alors $T_N = x_{n_{N+1}}-x_{n_0}$ converge vers $T$.
Donc la suite extraite $(x_{n_k})_k$ converge (vers $T+x_{n_0}$).
Conséquence : si toute s\'erie normalement convergente est convergente, alors on peut extraire de toute suite de Cauchy une sous-suite convergente donc $E$ est complet.

  \item Soit $p \le q$.
$$\| S_q-S_p\| = \| \sum_{k=p+1}^q u_k \| \le  \sum_{k=p+1}^{q} \| u_k \|
\le  \sum_{k=p+1}^{+\infty} \| u_k \|$$
Or la série $\sum_{k\ge 0} u_k$ est normalement convergente donc le reste
$\sum_{k=p+1}^{+\infty} \| u_k \|$ tend vers $0$ quand $p\rightarrow +\infty$.
Fixons $\epsilon >0$, il existe donc $N \in \Nn$ tel que pour $p\ge N$ on a
$\sum_{k=p+1}^{+\infty} \| u_k \| \le \epsilon$, donc pour tout $p,q \ge N$ on a aussi $\| S_q-S_p\| \le \epsilon$. Donc la suite $(S_n)$ est de Cauchy.
Si $E$ est complet alors $(S_n)$ converge, notons $S$ sa limite. Donc $\| S_n -S\|$ tend vers $0$. Dons la série $\sum_{k\ge 0} u_k$ est convergente (de somme $S$).  
\end{enumerate}
\fincorrection
\correction{002401}
\begin{enumerate}
  \item (1) $\Rightarrow$ (2). Supposons que $A_n$ converge vers $A$ dans 
${\cal L}(E,F)$. Soit $M\subset E$ une partie bornée, notons $M$ sa borne
(c'est-à-dire pour tout $x\in M$, $\|x\|\le B$).
Alors 
\begin{align*}
 & \forall \epsilon >0\quad \exists N\in \Nn \quad \forall n\ge N \quad \|A_n-A\| \le \frac \epsilon B \\
 \Rightarrow \qquad & \forall \epsilon >0\quad \exists N\in \Nn \quad \forall n\ge N \quad \forall x \in M \quad \|A_n(x)-A(x)\| \le \frac {\epsilon\|x\|} B \\
\Rightarrow \qquad & \forall \epsilon >0\quad \exists N\in \Nn \quad \forall n\ge N \quad \forall x \in M \quad \|A_n(x)-A(x)\| \le \epsilon \\
\end{align*}
Ce qui exactement la confergence uniforme de $A_n$ vers $A$ sur $M$.
  \item (2) $\Rightarrow$ (1). Par définition de la norme d'un opérateur nous avons $\| A_n - A\| = \sup_{\|x\|=1} \|A_n(x)-A(x)\|$.
Prenons comme partie bornée la sphère unité : $M = S(0,1) =\{x\in E \mid \|x\|=1\}$.
Alors :
\begin{align*}
 &\forall \epsilon >0\quad \exists N\in \Nn \quad \forall n\ge N \quad \forall x \in S(0,1) \in \|A_n(x)-A(x)\| \le \epsilon \\
 \Rightarrow \qquad & \forall \epsilon >0\quad \exists N\in \Nn \quad \forall n\ge N \quad \|A_n-A\| \le \epsilon \\
\end{align*}
Donc $\|A_n-A\|$ tend vers $0$.
\end{enumerate}
\fincorrection
\correction{002402}
C'est du cours, mais il est important de savoir rédiger ceci correctement.
Soit $(f_n)_n$ une suite de Cauchy de ${\cal L}(E,F)$.
\begin{enumerate}
  \item Trouvons d'abord le candidat à la limite.
Par définition d'une suite de Cauchy, nous avons :
$$\forall \epsilon >0 \quad \exists N \in \Nn \quad \forall p,q \ge N \quad \| f_p-f_q\| <\epsilon.$$
Fixons $x\in E$, alors 
$$\forall \epsilon >0 \quad \exists N \in \Nn \quad \forall p,q \ge N \quad \| f_p(x)-f_q(x)\|_F <\epsilon \|x\|_E.$$
Quitte à poser $\epsilon = \frac{\epsilon'}{\|x\|}$ ($x$ est fixé !, si $x=0$ c'est trivial) alors on a montrer :
$$\forall \epsilon' >0 \quad \exists N \in \Nn \quad \forall p,q \ge N  \quad \| f_p(x)-f_q(x)\|_F <\epsilon'.$$
Donc la suite $(f_n(x))_n$ est une suite de Cauchy de $F$.
Comme $F$ est complet alors cette suite converge, notons $f(x)$ sa limite.

  \item Nous avons construit une fonction $f :E\longrightarrow F$. 
Montrons que $f$ est dans l'espace ${\cal L}(E,F)$, c'est-à-dire que
$f$ est linéaire. Comme pour tout $n$, $f_n$ est linéaire alors, pour tout
$x,y\in E$, $\lambda,\mu\in \Rr$ on a
$$f_n(\lambda x +\mu y) = \lambda f_n(x) + \mu f_n (y).$$
\`A la limite ($n\rightarrow +\infty$) nous avons 
$$f(\lambda x +\mu y) = \lambda f(x) + \mu f (y),$$
donc $f$ est dans ${\cal L}(E,F)$.

  \item Il reste à montrer que $(f_n)$ converge bien vers $f$ (ce qui à priori n'est pas évident). Revenons à la définition d'une suite de Cauchy (écrit d'une fa\c{c}on un peu différente) :
$$\forall \epsilon >0 \quad \exists N \in \Nn \quad \forall p \ge N \quad \forall k \ge0\quad \| f_p-f_{p+k}\| <\epsilon.$$
Lorsque l'on fixe $p$ et que l'on fait tendre $k$ vers $+\infty$ alors
 $f_p-f_{p+k}$ tend vers $f_p-f$.
Donc en passant à la limite nous avons :
$$\forall \epsilon >0 \quad \exists N \in \Nn \quad \forall p \ge N \quad  \| f_p-f\| <\epsilon.$$
Donc $(f_n)$ converge vers $f$ pour la norme $\| . \|$ sur ${\cal L}(E,F)$.

\end{enumerate}
\emph{Remarque :} dans certains exercices il peut être utile de d'abord montrer le troisième point avant le deuxième.
\fincorrection
\nocorrection
\nocorrection
\nocorrection
\nocorrection
\nocorrection
\nocorrection
\nocorrection
\nocorrection
\nocorrection
\nocorrection
\nocorrection
\nocorrection
\nocorrection
\nocorrection
\nocorrection
\nocorrection
\nocorrection
\nocorrection
\nocorrection
\nocorrection
\nocorrection
\nocorrection
\nocorrection
\nocorrection
\nocorrection
\nocorrection
\nocorrection
\correction{002404}
\begin{enumerate}
  \item Commen\c{c}ons par l'unicité, si $x,y$ sont deux points fixes alors $f(x)=x$ et $f(y)=y$ donc la relation pour $f$ s'écrit
$$d(x,y) \le \alpha_n d(x,y) \quad \forall n \in \Nn.$$
Comme $\sum_{n\ge1}\alpha_n$ converge alors $(\alpha_n)$ tend vers $0$, donc il existe $n_0$ assez grand avec $\alpha_{n_0} < 1$, la relation devient 
$$ d(x,y) \le \alpha_{n_0} d(x,y) < d(x,y),$$
ce qui est contradictoire.

  \item  Soit $x_0 \in X$, notons $x_n = f^n(x_0)$. Alors
$$d(x_{n+1},x_n) \le \alpha_n d(x_1,x_0) \quad \forall n \in \Nn.$$
On va montrer que $(x_n)$ est une suite de Cauchy, c'est-à-dire
$$\forall \epsilon >0 \quad \exists N\in\Nn \quad
\forall n \ge N \quad \forall p \ge0\qquad d(x_{n+p},x_n) \le \epsilon.$$
Pour $n,p$ fixés, évaluons $d(x_{n+p},x_n)$.
\begin{align*}
d(x_{n+p},x_n) 
   &\le \sum_{k=n}^{n+p-1} d(x_{k+1},x_k) \\
   &\le \sum_{k=n}^{n+p-1} \alpha_k d(x_{1},x_{0}) \\
   &= d(x_{1},x_{0})\sum_{k=n}^{n+p-1} \alpha_k \\ 
\end{align*}

De plus la série $\sum_{n\ge1}\alpha_n$ converge donc la suite $(S_n)$ définie par
$S_n = \sum_{k=1}^n \alpha_k$ est de Cauchy et donc il existe $ N$ tel que
pour tout $n\ge N$ et tout $p\ge 0$ on a
$$\sum_{k=n}^{n+p-1} \alpha_k = S_{n+p-1}-S_{n-1} \le \epsilon.$$
Donc pour tout $n\ge N$ et tout $p\ge 0$ on $ d(x_{n+p},x_n) \le d(x_1,x_0)\epsilon$.
Quitte à poser $\epsilon'= d(x_1,x_0)\epsilon$, ceci prouve que $(x_n)$ est une suite de Cauchy.  
 Comme l'espace est complet alors cette suite converge, notons $x$ sa limite.

Pour tout $n\in \Nn$ nous avons
$$x_{n+1}=f(x_n).$$
\`A la limite, la suite $(x_{n+1})$ tend vers $x$, et comme $f$ est continue
(elle est $\alpha_1$-lipschitziènne : $d(f(x),f(y)) \le \alpha_1 d(x,y)$) alors
$(f(x_n))$  converge vers $f(x)$. Par unicité de la limite nous obtenons
$$x=f(x).$$
Donc $f$ possède un point fixe, qui est unique et est obtenu en partant d'un point quelconque $x_0\in X$ comme limite de $(f^n(x_0))_n$.

  \item Il reste à estimer la vitesse de convergence, nous avons vu
$$ d(x_{n+p},x_n) \le d(x_{1},x_{0})\sum_{k=n}^{n+p-1} \alpha_k,$$
On fait tendre $p$ vers $+\infty$ dans cette inégalité alors
$$ d(x,x_n) \le d(x_{1},x_{0})\sum_{k=n}^{+\infty} \alpha_k.$$
Ce qui était l'estimation recherchée.
\end{enumerate}
\fincorrection
\correction{002405}
Notons $g=f^n$. Alors $g$ est une application strictement contractante dans $X$ complet donc $g$ possède un unique point fixe que nous notons $x$.
Montrons l'unicité d'un point fixe pour $f$. Soit
$y\in X$ tel que $f(y)=y$ alors $g(y)=f^n(y)=y$. Donc $y$ est aussi un point fixe pour $g$, donc $y=x$. 

Il reste à montrer que $f$ possède effectivement bien un tel point fixe.
Nous avons 
\begin{align*}
 & f^n(x)=x \\
 \Rightarrow \qquad & f(f^n(x))=f(x) \\
\Rightarrow \qquad & f^n(f(x))=f(x) \\
\Rightarrow \qquad & g(f(x))=f(x) \\
\end{align*}
Nous venons de prouver que $f(x)$ est un point fixe de $g$. 
Comme $g$ possède un unique point fixe $x$ alors $f(x)=x$ !!
Donc $x$ est bien un point fixe pour $f$.
\fincorrection
\correction{002406}
\begin{enumerate}
  \item $(T\circ T f)(x)= 1+\int _0^x Tf(t-t^2) dt = 1+\int _0^x(1+\int_0^{t-t^2}f(u-u^2)du)dt
=1+x+\int_0^x\int_0^{t-t^2}f(u-u^2)dudt$.
De plus $(T\circ T f)'(x) = 1+\int_0^{x-x^2}f(u-u^2)du$.
En remarquant que pour $t\in[0,1]$, $t-t^2\le \frac 14$, on montre que
$|T\circ Tf(x)-T\circ Tg(x)| \le \frac 14 \|f-g\|_\infty$
et que $|(T\circ Tf)'(x)-(T\circ Tg)'(x)| \le \frac 14 \|f-g\|_\infty$
Donc $N(T\circ Tf-T\circ T g)\le \frac 12 \|f-g\|_\infty \le \frac 12 N(f-g)$.
Donc $T\circ T$ est une contraction et $X$ est complet donc $T\circ T$ admet 
un unique point fixe, par l'exercice \ref{exoiter}, $T$ admet un unique point fixe.

  \item Remarquons que $Tf=f$ est équivalent à $f(0)=1$ et $f'(x)=f(x-x^2)$.
Donc l'existence et l'unicité du point fixe pour $T$ donne l'existence et l'unicité de la solution
au problème posé.
\end{enumerate}
\fincorrection
\correction{002407}
\begin{enumerate}
  \item ! !
  
  \item 
\begin{align*}
\| Ax_1 - Ax_2\|_\infty 
  & = \| \int_a^b k(s,t)(x_1(t)-x_2(t))dt\|_\infty \\
  & \le  \int_a^b \|k(s,t)\|_\infty \|x_1(t)-x_2(t)\|_\infty dt \\
  & \le \|x_1(t)-x_2(t)\|_\infty \times \lambda \\
  & < \|x_1(t)-x_2(t)\|_\infty.\\
\end{align*}

Donc $A$ est contractante et l'espace ambiant ${\cal C}([a,b])$ est complet, par le
théorème du point fixe, $A$ admet un unique point fixe, $x$.
De plus, pour tout fonction $x_0 \in{\cal C}([a,b])$, la suite
$(A^nx_0)$ converge vers $x$, mais ici la norme est la norme uniforme donc
$\|A^nx_0-x\|_\infty$ tend vers $0$. Donc $(A^nx_0)$ converge uniformément vers $x$.

  \item
\begin{align*}
\|x_1-x_2\|_\infty & = \| A_1x_1-A_2x_2\|_\infty \quad \text{ car } A_ix_i=x_1, \\
                   &= \| \int_a^b k_1(s,t) x_1(t)  dt +y_1(s)+ \int _a^b k_2(s,t) x_2(t) \, dt + y_2(s)  \|_\infty \\
                    &\le \| \int_a^b k(s,t)(x_1(t)-x_2(t))dt\|_\infty + \|y_1-y_2\|_\infty \\
  &\le \lambda\|x_1-x_2\|_\infty+\|y_1-y_2\|_\infty \\   
\end{align*}
Donc 
$$\|x_1 -x_2\|_\infty    \leq \frac{1}{1-\lambda}\| y_1
-y_2\|_\infty,$$
ce qui exprime la dépendance continue de la solution par rapport à la fonction $y$.
\end{enumerate}
\fincorrection
\nocorrection
\nocorrection
\nocorrection
\nocorrection
\nocorrection
\nocorrection
\nocorrection
\nocorrection
\nocorrection
\nocorrection
\nocorrection
\nocorrection
\correction{002392}
\begin{enumerate}
  \item Par l'absurde supposons que $X$ n'a aucun point isolé.
Comme $\{x\}$ est un fermé alors $\omega_x = X \setminus \{x\}$
est un ouvert (de $X$). De plus comme le point $x$ n'est pas isolé
alors $\omega_x$ est dense dans $X$.

Maintenant on peut appliquer le théorème de Baire à $X$ qui est un fermé de l'espace complet $\Rr$. Donc une intersection dénombrable d'ouverts denses dans $X$ est encore dense.
Mais ici nous obtenons une contradiction car les $\omega_x$ sont des ouverts denses, $X$ est dénombrable mais $$\bigcap_{x\in X} \omega_x = \varnothing.$$
Et l'ensemble vide n'est pas dense dans $X$ !!

  \item Pour l'ensemble de Cantor aucun point n'est isolé, donc par la question précédente
l'ensemble de Cantor n'est pas dénombrable.
\end{enumerate}
\fincorrection
\correction{002393}
\begin{enumerate}
  \item  
Par l'absurde supposons que sur aucun ouvert $f$ n'est majorée.
$f :  X \to \Rr$ est {semi-continue inf\'erieurement} donc 
$$\forall \lambda \in \Rr \qquad  O_\lambda := \{ x\in X \mid f(x) > \lambda \} \qquad \text{est un ouvert.}$$
De plus $O_\lambda$ est dense, en effet pour $x \in X$ et pour $V_x$ un voisinage ouvert de $x$,
alors par hypothèse $f$ n'est pas majorée sur $V_x$ donc en particulier il existe $y \in V_x$ tel que $f(y) > \lambda$ donc $y\in V_x \cap O_\lambda$. Ceci prouve  que $O_\lambda$ est dense dans $X$ ($V_x$ étant aussi petit que l'on veut).

Maintenant pour $n = 0,1,2,\ldots$, les $O_n$ sont un ensemble dénombrable d'ouverts denses. Comme $X$ est complet il vérifie le théorème de Baire donc l'intersection des $O_n$ est encore un ensemble dense. Mais il est facile de voir par la définition des $O_n$ que
$$\bigcap_{n\in\Nn}O_n = \varnothing.$$
Ce qui donne la contradiction cherchée.


  \item On note $\phi : B \to \Rr$ la fonction définie par 
$$\phi(x) = \sup_{n\in\Nn} |f_n(x)|.$$
Il n'est pas difficile de montrer que $\phi$ est semi-continue inf\'erieurement :
en effet soit $F_\lambda := \{ x\in X \mid \phi(x) \le \lambda \}$.
Soit $\lambda$  fixé et soit $(x_k)$ une suite d'éléments de $F_\lambda$.
Pour $n$ fixé et pour tout $k$ on a $f_n(x_k) \le k$, donc par continuité de $f_n$,
on a $f_n(x) \le k$, ceci étant vrai pour tout $n$ on a $x\in F_\lambda$.
Donc $F_\lambda$ est un fermé donc $O_\lambda := \{ x\in X \mid f(x) > \lambda \}$ est un ouvert.
Donc $\phi$ est semi-continue inf\'erieurement.

Par la première question il existe un ouvert non vide $O$ et une constante $M>0$ tel que
$\phi$ soit majorée par $M$ sur $O$.
C'est-à-dire
$$\forall n\in \Nn \quad \forall x \in O \quad |f_n(x)| \le M.$$
Par translation on peut supposer que l'origine $o$ est inclus dans $O$. Donc il existe $\epsilon >0$ tel que $\bar B(o,\epsilon) \subset O$. Donc
$$\forall n\in \Nn \quad \forall x \in \bar B(o,\epsilon) \quad |f_n(x)| \le M$$
ce qui est équivalent à 
$$\forall n\in \Nn \quad \forall x \in \bar B(o,1) \quad |f_n(x)| \le \frac M \epsilon$$
Donc $$\forall n\in \Nn \quad  \|f_n\| \le \frac M \epsilon.$$
\end{enumerate}
\fincorrection
\nocorrection
\nocorrection
\nocorrection
\nocorrection
\nocorrection
\nocorrection
\nocorrection
\nocorrection
\nocorrection
\nocorrection
\nocorrection
\correction{002383}
\begin{enumerate}
  \item C'est du cours.
  \item Si $f : B \longrightarrow \{ 0,1 \}$ est continue alors elle induit une application restreinte $f_{|A} : A \longrightarrow \{ 0,1 \}$ continue. Donc $f$ est constante sur $A$. Soit $b \in B$ et soit $(a_n)$ une suite d'éléments de $A$ qui tendent vers $b$ (c'est possible car $B \subset \bar A$), alors $f(a_n)$ est constante, par exemple égal à $1$, car $A$ est connexe. Mais $f$ est continue sur $B$, donc $f(b) = \lim f(a_n) = 1$. On montre ainsi que $f$ est constante sur $B$. Donc $B$ est connexe. (Au passage on a montrer que $\bar A$ était connexe.)
  \item Soit $f : A \longrightarrow \{ 0,1 \}$ une fonction continue, o\`u $A = \bigcup A_n$. $A_0$ est connexe donc $f$ est constante sur $A_0$ et vaut $v_0$,
de même $A_1$ est connexe donc $f$ est constante sur $A_1$ et vaut $v_1$. Mais 
pour $a \in A_0\cap A_1 \neq \varnothing$, on a $f(a)=v_0$ car $a\in A_0$ et
$f(a)=v_1$ car $a\in A_1$. Donc $v_0=v_1$. Donc $f$ est constante sur $A_0 \cup A_1$. Par récurrence $f$ est constante sur $A$.
\end{enumerate}
\fincorrection
\correction{002384}
\begin{enumerate}
  \item Dans $\Rr^2$ il y a deux composantes connexes : $\{(x,y)\in \Rr^2 \; ; x > y\}$ et $\{(x,y)\in \Rr^2 \; ; x < y\}$.

  \item Dans $\Cc^2$ il n'y en a qu'une seule : $\{(z,w) \in \Cc^2 \; ; z\neq w\}
$
\end{enumerate}
\fincorrection
\correction{002385}
Notons la frontière $\mathrm{Fr}\, A = \bar A \setminus \mathring A$.
Nous avons la partition $X = \mathring A \cup \mathrm{Fr}\, A  \cup (X \setminus \bar A)$.
Si $B \cap \mathrm{Fr}\, A  = \varnothing$ alors $B\subset \mathring A \cup (X \setminus \bar A)$. 

De plus, par hypothèses, $B \cap A \neq \varnothing$ et $B \cap \mathrm{Fr}\, A = \varnothing$ 
or $\mathring A= A \setminus \mathrm{Fr}\, A$ donc $B \cap \mathring A \neq \varnothing$.
Comme $\mathrm{Fr}\, A = \mathrm{Fr}\, (X\setminus A)$ on a $B\cap \mathrm{Fr}\, (X\setminus A) = \varnothing$.
Par hypothèse $B\cap (X\setminus  A) \neq \varnothing$ donc $B\cap  (X\setminus \bar A) = 
(B\cap  (X\setminus A)) \setminus (B\cap \mathrm{Fr}\, (X\setminus A)) \neq \varnothing$.

 Nous avons montrer que $B$ est inclus dans l'union de deux ouverts disjoints 
$\mathring A$ et $X\setminus \bar A$, d'intersection non vide avec $B$, donc $B$ n'est pas connexe.
Par contraposition, si $B$ est connexe alors $B$ ne rencontre pas la frontière 
de $A$.
\fincorrection
\correction{002386}
\begin{enumerate}
  \item $T$ est compact car c'est un fermé borné de $\Rr^2$.

Soit $g : T \longrightarrow \{0,1\}$ une application continue.
Par connexité du segment  $[-1,1]$, $g$ est constante sur $\{0\}\times [-1,1]$
(et vaut $v$) ; $g$ est aussi constante sur $[-1,1]\times \{0\}$ et vaut $v'$.
 Mais alors $v=g(0,0)=v'$ donc $g$ est constante sur $T$. Donc $T$ est connexe.

Pour $f:T\to \Rr$ une fonction continue. $T$ est compact donc $f(T)$ est compact.
$T$ est connexe donc $f(T)$ est connexe. Donc $f(T)$ est un compact connexe de $\Rr$ c'est donc un segment compact.

  \item Ce sont les quatre points cardinaux $N=(0,1)$, $S=(0,-1)$, $E=(1,0)$, $W=(-1,0)$.

  \item Par l'absurde, supposons que $T$ soit homéomorphe à une partie $I$ de $\Rr$, alors
il existe un homéomorphisme $f : T \longrightarrow I$. Par le premier point
$I$ est un segment compact $I = [a,b]$. $T\setminus \{ N\}$ est connexe donc sont image par $f$, $f(T\setminus \{ N\})$ est connexe, mais c'est aussi le segment $I$ privé d'un point. $I$ privé d'un point étant connexe, le point retiré est nécessairement une extrémité. Donc $f(N)=a$ ou $f(N)=b$. Supposons par exemple $f(N)=a$. On refait le même raisonnement avec $S$, qui s'envoie aussi sur une extrémité, comme $f$ est bijective cela ne peut être $a$, donc
$f(S)=b$. Maintenant $f(E)$ est aussi une extrémité donc $f(E) \in \{a,b\}$.
Mais alors $f$ n'est plus injective car on a $f(E)=f(N)$ ou $f(E)=f(S)$.
Contradiction.
\end{enumerate}
\fincorrection
\correction{002387}
\begin{enumerate}
  \item 
    \begin{enumerate}
      \item $\phi : \Rr \longrightarrow \mathbb{S}^1$ définie par $\phi(t)=e^{it}$ est une surjection continue.
      \item $\mathbb{S}^1$ est un compact connexe donc, par l'absurde, si $\psi : \mathbb{S}^1 \longrightarrow \Rr$ est une injection continue alors $\psi(\mathbb{S}^1)$ est un compact connexe de $\Rr$ donc un segment compact $I$. Soit $y \in \mathring I$, comme $I$ est l'image de $\mathbb{S}^1$ alors il existe un unique $x \in \mathbb{S}^1$ tel que $f(x)=y$. L'application $f$ induit alors une bijection continue 
$f : \mathbb{S}^1 \setminus \{ x \} \longrightarrow I \setminus \{ y \}$.
Mais $\mathbb{S}^1 \setminus \{ x \}$ est connexe alors que son image par
$f$, qui est $ I \setminus \{ y \}$ ne l'est pas (car $y\in \mathring I$).
L'image d'un connexe par une application continue doit être un connexe, donc nous avons une contradiction.
    \end{enumerate}


  \item Si $\chi : \Rr^2 \longrightarrow \Rr$ est une injection continue. 
Comme $\Rr^2$ est connexe $f(\Rr^2)=I$ est un connexe de $\Rr$ donc un segment
(non réduit à un point !). Prenons $y$ un élément de $\mathring I$, soit $x\in \Rr^2$ tel que $f(x)=y$. Alors $\Rr^2 \setminus \{ x\}$ est connexe,  $I \setminus\{ y \}$ ne l'est pas, et $f$ est une bijection continue entre ces deux ensembles, d'o\`u une contradiction.
\end{enumerate}
\fincorrection
\correction{002388}
L'ensemble $B$ est connexe si
et seulement si toute fonction continue $f:B\to \{0,1\}$ est
constante. Soit alors $f:B\to \{0,1\}$ une fonction continue et
montrons qu'elle est constante.
Remarquons que la restriction de $f$ \`a tout ensemble $B_a$ est constante ($B_a$ est connexe).

On définit $g : \Rr \longrightarrow \{0,1\}$ tel que $g(x)$ prend  la valeur
qu'a $f$ sur $B_x$.
Nous allons montrer que $g$ est localement constante (on ne sait pas si $g$ est continue).

\begin{itemize}
\item Soit $a \notin \Qq$ alors on a $(a,0)\in B$, $f$ est une fonction continue et
$\{f(a,0)\}$ est un ouvert de $\{0,1\}$, donc $f^{-1}(\{ f(a,0) \})$ est un ouvert de $B$.
Donc il existe $\epsilon >0$ tel que si
$(x,y) \in (]a-\epsilon,a+\epsilon[\times ]-\epsilon,\epsilon[)\times B$ alors
$f(x,y)=f(a,0)$. Alors pour $x\in ]a-\epsilon,a+\epsilon[$ on a $g(x)=g(a)$ :
si $x\notin \Qq$ alors $g(x)=f(x,0)=f(a,0)=g(a)$; et si $x\in \Qq$ alors
$g(x)= f(x,\frac \epsilon 2) = f(a,0)=g(a)$.
Donc $g$ est localement constante au voisinage des point irrationnels.

  \item Si $a\in \Qq$ et soit $b\in ]0,1]$ alors $f$ est continue en $(a,b)$ donc il existe
$\epsilon >0$ tel que pour tout $x\in ]a-\epsilon,a+\epsilon[\cap \Qq$,
$g(x)=f(x,b)=f(a,b)=g(a)$.
Si maintenant $x\in ]a-\epsilon,a+\epsilon[\cap (\Rr \setminus \Qq)$, on prend une 
suite $(x_n)$ de rationnels qui tendent vers $x$. 
Comme $f$ est continue alors
$g(a)=g(x_n)=f(x_n,b)$ tend vers $f(x,b)=g(x)$. Donc $g(a)=g(x)$.
Nous avons montrer que $g$ est localement constante au voisinage des point rationnels.

  \item Bilan : $g$ est localement constante sur $\Rr$.
\end{itemize}

Comme $\Rr$ est connexe, alors $g$ est constante sur $\Rr$. Donc
$f$ est constante sur $\Rr$. Ce qu'il fallait démontrer.
\fincorrection
\correction{002389}
\begin{enumerate}
  \item $A$ est connexe car connexe par arcs.

  \item Si $z \in g(A)$ alors il existe $(x,y) \in A$ tel que $g(x,y)=z$.
Donc $z = \frac{f(y)-f(x)}{y-x}$ par le théorème des accroissements finis il existe $t \in ]x,y[ \subset I$ tel que $z = f'(t)$ donc $z \in f'(I)$.
Donc $g(A) \subset f'(I)$.

Si maintenant $z \in f'(I)$, il existe $y\in I$ tel que $z= f'(y)$, mais par définition de la dérivée $f'(y)$ est la limite de $\frac{f(y)-f(x)}{y-x}$  quand $x$ tend vers $y$ (et on peut m\^eme dire que c'est la limite à gauche, i.e. $x<y$). Donc $f'(y)$ est limite de points de $g(x,y)$ avec $x<y$, donc de points de $A$.
Conclusion $z=f'(y)$ est dans $\overline {g(A)}$, et donc $f'(I) \subset \overline {g(A)}$.

  \item $A$ est connexe, $g$ est continue sur $A$ donc $g(A)$ est un connexe de $\Rr$. Par l'exercice \ref{exocon} comme on a 
$$ g(A)\subset f'(I) \subset \overline{g(A)}$$
avec  ${g(A)}$ connexe alors $f'(I)$ est connexe. Comme $f'(I)$ est un connexe de $\Rr$ c'est un intervalle.
\end{enumerate}
\fincorrection
\correction{002390}
Soit $a \in \bigcap_{i\in I}A_i$ ; soit $x,y \in \bigcup _{i\in I}A_i$.
Il existe $i_1$ tel que $x\in A_{i_1}$ on a aussi $a \in A_{i_1}$ donc
il existe une chemin $\gamma_1$ qui relie $x$ à $a$.
De même il existe $i_2$ tel que $x\in A_{i_2}$ et on a également 
$a \in A_{i_2}$ donc il existe une chemin $\gamma_2$ qui relie $a$ à $y$.
Le chemin $\gamma_2 \circ \gamma_1$ relie $x$ à $y$. Ceci étant valable quelque soient $x$ et $y$, $\bigcup _{i\in I}A_i$ est connexe par arcs.
\fincorrection
\correction{002391}
\begin{enumerate}
  \item Si $(x_1, \sin \frac{1}{x_1})$ et $(x_2, \sin \frac{1}{x_2})$
sont deux points  de $A$ alors le graphe au dessus de $[x_1,x_2]$ définie un chemin reliant ces deux points. Plus précisément le chemin est l'application
$\gamma : [x_1,x_2] \longrightarrow \Rr^2$ définie par $\gamma(t) = (t,\sin \frac 1t)$. Donc $A$ est connexe par arcs donc connexe.
  \item $\bar A = A \cup (\{0\}\times [-1,1])$. On peut utiliser l'exercice \ref{exocon} pour montrer que $\bar A$ est connexe. Ici nous allons le montrer directement.
Supposons, par l'absurde, que $\bar A \subset U\cup V$ avec $U$ et $V$ des ouverts de $\Rr^2$ disjoints, d'intersection non vide avec $A$.
Comme $\{0\}\times [-1,1]$ est connexe il est entièrement inclus dans un des ouverts, supposons qu'il soit inclus dans $U$. Comme $A$ est connexe alors il est inclus dans un des ouverts, donc il est inclus dans $V$ (car s'il était inclus dans $U$, tout $\bar A$ serait contenu dans $U$). 
Trouvons une contradiction en prouvant qu'en fait $U\cap A \neq \varnothing$.
En effet $U$ est un ouvert et $(0,0) \in U$, soit $B((0,0),\epsilon)$ une boule contenue dans $U$. Pour $n$ suffisamment grand on a $x_n = \frac{1}{2\pi n} < \epsilon$ avec $\sin \frac 1 {x_n} = \sin {2\pi n} =0$ donc
$(x_n,\sin \frac 1{x_n})= (x_n,0)$ est un élément de $A$ et de $U$.
Comme $V$ contient $A$ alors $U\cap V \neq \varnothing$. Ce qui fournit la contradiction.

  \item Montrons que $\bar A$ n'est pas connexe par arcs. Soit $O = (0,0)$ et $P = (\frac{1}{2\pi},0)$ deux points de $\bar A$, par l'absurde supposons qu'il existe un chemin $\gamma : [0,1]  \longrightarrow \bar A$ tel que $\gamma(0)=O$ et $\gamma(1)=P$. On décompose en coordonnées $\gamma(t) = (\gamma_1(t),\gamma_2(t)) \in \Rr^2$. $\gamma_1^{-1}(\{0\})$ est un fermé car $\gamma_1$ est continue et de plus il est non vide car $\gamma_1(0)=0$.
Soit $t_0 = \sup \gamma_1^{-1}(\{0\})$, comme l'ensemble est fermé alors $\gamma_1(t_0)=0$ et de plus $t_0 < 1$ car $\gamma_1(1) = \frac 1 {2\pi}$.

On regarde ce qui se passe au temps $t_0$, c'est l'instant ou notre chemin ``quitte'' l'ensemble $\{0\}\times [-1,1]$. Notons $y_0 = \gamma_2(t_0)$.
Comme $\gamma_2$ est continue en $y_0$ et pour $\epsilon = \frac 12$ il existe $\eta>0$ tel que $(|t-t_0|<\eta \Rightarrow |\gamma_2(t)-y_0|<\frac12)$.
Choisissons $t_1\in ]t_0,t_0+\eta[$. Alors $t_1>t_0$ donc
$\gamma_1(t_1) > 0$. Donc le point $\gamma(t_1)=(\gamma_1(t_1),\gamma_2(t_1))$ est dans $A$ (et plus seulement dans $\bar A$).

Supposons par exemple $y_0 \le 0$, alors 
quand $x$ parcourt $]\gamma_1(t_0),\gamma_1(t_1)[$, $\sin \frac 1x$ atteint
la valeur $1$ une infinité de fois. Donc il existe $t_2 \in ]t_0,t_1[$ tel que
$\gamma_2(t_2)=1$. Donc $\gamma(t_2) = (\gamma_1(t_2),1)$.
Mais comme $|t_2-t_0|< \eta$ alors $|\gamma_2(t_2)-y_0|=|1-y_0| > \frac 12$.
Ce qui contredit la continuité de $\gamma_2$. Nous avons obtenu une contradiction donc $\bar A$ n'est pas connexe par arcs.
\end{enumerate}
\fincorrection
\nocorrection
\nocorrection
\nocorrection
\nocorrection
\nocorrection
\nocorrection
\nocorrection
\nocorrection
\nocorrection
\nocorrection
\nocorrection
\nocorrection
\nocorrection
\nocorrection
\nocorrection
\nocorrection
\nocorrection
\nocorrection
\nocorrection
\nocorrection
\nocorrection
\nocorrection
\nocorrection
\nocorrection
\nocorrection
\nocorrection
\nocorrection
\nocorrection
\correction{002408}
Soit $P(x) = a_dx^d+\cdots+a_1x+a_0 \in \Rr[x]$ alors par linéarité de l'intégrale
et grâce à la relation de l'énoncé :
$$\int_a^b f(t)\cdot P(t) \, dt = 0.$$


La fonction $f$ est continue sur le compact $[a,b]$ donc par le théorème de Weierstrass
il existe une suite de polynômes qui converge uniformément vers $f$.
Fixons $\epsilon >0$. Soit $P$ tel que $\| f-P \|_\infty \le \epsilon$.
Alors 
\begin{align*}
|\int_a^b f(t)^2 dt|  
  &= \left| \int_a^b f(t)^2 dt - \int_a^b f(t) \cdot P(t) dt\right| \\
  &=  \left|\int_a^b f(t)\cdot (f(t)-P(t)) dt\right| \\
  &\le  \int_a^b |f(t)|\cdot \| f-P\|_\infty  dt \\
  &\le \epsilon \int_a^b |f| \\.
\end{align*}
Mais  $C = \int_a^b |f|$ est une constante (indépendante de $\epsilon$ et $P$).
Donc on vient de montrer que $|\int_a^b f(t)^2 dt| \le \epsilon C$ avec  
pour tout $\epsilon >0$ donc $\int_a^b f^2 = 0$, or $f^2$ est une fonction continue et positive,
son intégrale est nulle donc $f$ est la fonction nulle.
\fincorrection
\nocorrection
\correction{002410}
 Soit $\Phi : E \to \Rr^n$ définie par $\Phi = (f_1,\ldots, f_n)$ alors 
$\Phi$ est continue car les $f_i$ sont continues. $\Phi$ est injective :
en effet si $x\neq y$ alors comme  $\{f_i\}$ sépare les points on a $\Phi(x)\neq \Phi(y)$, par contraposition
$\Phi$ est injective.
Notons $F =  \Phi(E)$ l'image directe de $E$. Alors $\Phi : E \to F$ est continue et bijective.
Comme $E$ est compact alors $\Phi$ est un homéomorphisme.
Donc $E$ est homéomorphe à $F$ qui est une partie de $\Rr^n$.
\bigskip

\emph{Rappel : } Si  $\Phi : E \to F$ est continue et bijective et $E$ est un espace compact alors 
$\Phi$ est un homéomorphisme.

La preuve est simple : soit $K$ un ensemble fermé de $E$,
comme $E$ est compact alors $K$ l'est aussi. Comme $\Phi$ est continue alors
$\Phi(K)$ est un compact de $F$ donc un fermé.
Mais en écrivant ceci à l'aide de l'application $\Phi^{-1}$ nous venons de  montrer
que pour tout fermé $K$ de $E$, l'image réciproque de $K$ par  $\Phi^{-1}$ (qui est 
$(\Phi^{-1})^{-1}(K)=\Phi(K)$) est un fermé. Donc $\Phi^{-1}$ est continue. Donc 
$\Phi$ est un homéomorphisme.
\fincorrection
\correction{002411}
On cherche à vérifier les hypothèses du théorème de Stone-Weierstrass.
\begin{itemize}
\item Tout d'abord $X\times Y$ est compact, car c'est un produit d'espaces compacts.
\item Ensuite $\mathcal{A}$ est une sous-algèbre de $\mathcal{C}(X\times Y,\Rr)$ :
en effet pour $f,g \in \mathcal{A}$ et $\lambda \in \Rr$ on a :
$$ f+g \in \mathcal{A}, \quad \lambda\cdot f \in \mathcal{A}\quad \text{ et }  f\times g \in \mathcal{A}.$$
\item  $\mathcal{A}$ sépare les points : soient $(x_1,y_1)\neq (x_2,y_2) \in X\times Y$.
Supposons que $x_1\neq x_2$, soit $u \in \mathcal{C}(X,\Rr)$ tel que $u(x_1)\neq u(x_2)$ (clairement une telle fonction existe !), soit $v$ la fonction sur $Y$ constante égale à $1$.
Alors $f$ définie par $f(x,y) = u(x)\cdot v(y)$ est dans $\mathcal{A}$ et 
$f(x_1,y_1) = u(x_1)\neq u(x_2)=f(x_2,y_2)$. Si $x_1=x_2$ alors nécessairement $y_1\neq y_2$ et on fait un raisonnement similaire.
\item Pour tout $(x,y)\in X\times Y$
il existe une fonction $f\in \mathcal{A}$ telle que $f(x)\neq 0$ : prendre la fonction $f$ constante égale à $1$ qui est bien dans $\mathcal{A}$.
\end{itemize}

Par le théorème de Stone-Weierstrass $\mathcal{A}$ est dense dans  $\mathcal{C}(X\times Y,\Rr)$
pour la norme uniforme.
\fincorrection
\correction{002412}
\begin{enumerate}
  \item Pour $f\in {\cal F}$, par le théorème des accroissements finis, 
pour tout $t_0,t\in [a,b]$ il existe $c\in]t_0,t[$ tel que
$|f(t)-f(t_0)| = |f'(c)| |t-t_0|$. Donc $|f(t)-f(t_0)| \le k |t-t_0|$. 
Fixons $t_0\in[a,b]$. Soit $\epsilon >0$, soit $\eta = \frac \epsilon k$ alors
$$\forall t\in[a,b] \qquad |t-t_0| \le \eta \quad \Rightarrow \quad |f(t)-f(t_0)| \le k |t-t_0| \le \epsilon.$$
Ce qui est exactement l'équicontinuité de $\mathcal{F}$ en $t_0$.
Comme nous pouvons prendre pour $t_0$ n'importe quel point de $[a,b]$ alors
$\mathcal{F}$ est équicontinue.

  \item
   \begin{enumerate}
      \item Notons $\mathcal{H} = \{ f_n \mid n\in \Nn \}$.
Pour $x_0,x \in \Rr^n$, $\|f_n(x)-f_n(x_0)\| \le L\|x-x_0\|$.
Donc en posant $\eta = \frac\epsilon L$ comme ci-dessus on prouve l'équicontinuité de $\mathcal{H}$ en $x_0$, puis partout.
      \item Notons $\mathcal{H}(x) = \{ f_n(x) \mid n\in \Nn \}$.
Alors par hypothèse, 
 $\mathcal{H}(0) \subset {\bar B} (0,\sqrt 2)$.
Donc ${\bar {\cal H}}(0)$ est un fermé de ${\bar B}(0,\sqrt 2)$ qui est compact
(nous somme dans $\Rr^n$), donc $\bar {\cal H}(0)$ est aussi compact, d'o\`u
$\mathcal{H}(0)$ relativement compact.
Maintenant nous avons $\| f_n(x)-f_n(0)\| \le L\|x-0\|$. Donc
$\|f_n(x)\| \le L \| x \| + \sqrt 2$. Donc pour $x$ fixé, $f_n(x)\in
{\bar B}(0, L \| x \| +\sqrt 2)$ ce qui implique que 
$\mathcal{H}(x)$ est relativement compact.

      \item Comme $\Rr^n$ n'est pas compact on ne peut pas appliquer directement le théorème d'Ascoli.
Soit $B_R = \bar B(0,R)$ qui est un compact de $\Rr^n$. Notons $\mathcal{H}_R =
 \{ {f_n}_{|B_R} \mid n\in \Nn \}$ la restriction de $\mathcal{H}$ à $B_R$.
Alors par le théorème d'Ascoli, $\mathcal{H}_R$ est relativement compact.
Donc de la suite $({f_n}_{|B_R})_n$ on peut extraire une sous-suite convergente (sur $B_R$).

  \item Pour $R=1$ nous extrayons de $(f_n)_n$ une sous-suite $(f_{\phi_1(n)})_n$ qui converge sur $B_1$.
Pour $R=2$, nous extrayons de  $(f_{\phi_1(n)})_n$ une sous-suite $(f_{\phi_2(n)})_n$  qui converge sur $B_2$.
Puis par récurrence pour $R=N$, nous extrayons de  $(f_{\phi_{N-1}(n)})_n$ une sous-suite $(f_{\phi_N(n)})_n$  
qui converge sur $B_N$. 
Alors la suite $(f_{\phi_n(n)})_n$ converge sur $\Rr^n$. C'est le procédé diagonal de Cantor.
En effet soit $x\in \Rr^n$ et soit  $N \ge \|x\|$. Alors
$x \in B_N$ donc  $(f_{\phi_N(n)}(x))_n$ converge vers $f(x)$, mais 
$(f_{\phi_n(n)})_{n\ge N}$ est extraite de $(f_{\phi_N(n)})_n$
donc  $(f_{\phi_n(n)}(x))_{n}$ converge également vers $f(x)$.
Nous venons de montrer que  $(f_{\phi_n(n)})_n$ converge simplement vers $f$
sur tout $\Rr^n$.
    \end{enumerate}

\end{enumerate}
\fincorrection
\correction{002413}
\begin{enumerate}
  \item 
  \begin{enumerate}
    \item 
  Soit $(x_n)$ une suite convergeant vers $a$, alors
$$|f_n(x_n)-b| \le |f_n(x_n)-f_n(a)|+|f_n(a)-b|.$$

    \item Soit $\epsilon >0$, il existe $N_1$ tel que pour $n\ge N_1$ on ait
$|f_n(a)-b| < \frac \epsilon 2$.

    \item $(f_n)$ est équicontinue en $a$, donc il existe $\eta >0$ tel que
pour tout $n\in \Nn$ et tout $x\in E$, $(|x-a| < \eta \Rightarrow |f_n(x)-f_n(a)| < \frac \epsilon 2)$.
    
    \item Comme $x_n\rightarrow a$ alors il existe $N_2$ tel que pour $n\ge N_2$  on ait $|x_n-a| < \eta$.

    \item Donc pour $n \ge \max(N_1,N_2)$ on a $|f_n(x_n)-b| \le |f_n(x_n)-f_n(a)|+|f_n(a)-b| < \frac \epsilon 2 + \frac \epsilon 2 = \epsilon$. Donc
$(f_n(x_n))$ converge vers $b$.
  \end{enumerate}

  \item  Soit des fonctions réelles définies par $f_n(x) = (1+x)^n$.
Prenons $x_n= \frac 1 n$, alors $x_n \rightarrow a=0$.
Par contre $f_n(a) = f_n(0)=1$ pour tout $n$. Mais $f_n(x_n) = f_n(\frac 1n) =
(1+\frac 1n)^n$ converge vers $e$. L'équicontinuité est donc bien nécessaire.
\end{enumerate}
\fincorrection
\correction{002414}
Notons $G$ l'ensemble des $x\in E$, pour lesquels $(f_n(x))$
est une suite de Cauchy dans $F$. Soit $(x_n)$ une suite d'éléments de $G$ qui converge vers $x \in E$. Il faut montrer $x\in G$, c'est-à-dire que
$(f_n(x))$ est une suite de Cauchy de $F$.
\'Ecrivons pour $p,q,n \in \Nn$,
$$\| f_p(x)-f_q(x) \| \le \| f_p(x)-f_p(x_n) \| + \| f_p(x_n)-f_q(x_n) \|+\| f_q(x_n)-f_q(x) \|.$$

Soit $\epsilon >0$, comme $(f_n)$ est équicontinue en $x$, il existe $\eta >0$ tel que 
$$\forall n\in \Nn \quad \forall y\in E  \quad \| x-y\| < \eta \quad \Rightarrow \| f_n(x)-f_n(y)\| < \frac \epsilon 3.$$
Comme $x_n \rightarrow x$ il existe $N \ge 0$ tel que
$\|x_N-x\| < \eta$. 
Donc 
$$\forall p,q \ge 0 \quad \| f_p(x_N)-f_p(x)\| < \frac \epsilon 3 \quad 
\text{ et  }\quad \| f_q(x_N)-f_q(x)\| < \frac \epsilon 3.$$
Enfin $N$ étant fixé, $x_N \in G$, la suite $(f_n(x_N))_n$ est une suite de Cauchy, donc il existe $N'\ge N$ tel que pour $p,q \ge N'$ on a,
$$\| f_p(x_N)-f_q(x_N)\| < \frac \epsilon 3.$$
Le bilan de toute ces inégalités est donc 
$$\forall p,q \ge N' \quad \| f_p(x)-f_q(x) \| <\epsilon.$$
Donc $(f_n(x))_n$ est une suite de Cauchy, donc $x\in G$ et $G$ est fermé.
\fincorrection
\correction{002415}
\begin{enumerate}
  \item 
  \begin{enumerate}
    \item Montrons que $A$ est ouvert. Soit $x\in A$, alors $\mathcal{H}(x) = \{f(x) \mid f\in \mathcal{H}\}$ est bornée, notons $M$ une borne. \'Ecrivons l'équicontinuité pour $\epsilon = 1$, il existe $\eta >0$ tel que
$$\forall f \in \mathcal{H} \quad \forall y\in E \quad (\|x-y\| < \eta \Rightarrow  |f(x)-f(y)| < 1).$$
Or si $|f(x)-f(y)| < 1$ alors $|f(y)| < |f(x)| + 1 \le M+1$. On a donc montré
$$\forall f \in \mathcal{H} \quad \forall y\in E \quad (y\in B(x,\eta) \Rightarrow  |f(y)| < M+1).$$
Donc $B(x,\eta) \subset A$. Donc $A$ est ouvert.

    \item Montrons que $A$ est fermé. Soit $(x_n)$ une suite d'éléments de $A$
qui converge vers $x\in E$. On reprend $\epsilon = 1$ et on obtient un $\eta$ par équicontinuité. Comme $x_n\rightarrow x$ alors il existe $N$ tel que $\|x_N-x\| < \eta$. Donc pour tout $f$ dans $\mathcal{H}$, $|f(x)-f(x_N)| < 1$ ; donc
$|f(x)| < |f(x_N)|+1$. Or $x_N \in A$, il existe $M$ tel  $|f(x_N)|$ soit bornée par $M$ pour tout $f$ dans $\mathcal{H}$.
Donc pour tout $f\in \mathcal{H}$, $|f(x)| < M+1$. Donc $x\in A$. Donc $A$ est fermé.
  \end{enumerate}

  \item $x_0\in A$ donc $A$ est non vide, comme $A$ est ouvert et fermé et $E$ est connexe alors $A=E$. donc pour tout $x\in E$, $\mathcal{H}(x)$ est borné dans $\Rr$, donc $\overline{\mathcal{H}(x)}$ est un compact de $\Rr$.
Par le théorème d'Ascoli, $\mathcal{H}$ étant équicontinue et $E$ étant compact
alors $\bar{\mathcal{H}}$ est compact.
\end{enumerate}
\fincorrection
\correction{002416}
\begin{enumerate}
  \item
  \begin{enumerate}
    \item Pour $t \ge 0$ fixé, alors
\begin{align*}
f_n(t) 
  &= \sin \sqrt{t+4(n\pi)^2} \\
  &= \sin 2n\pi\sqrt{1+\frac{t}{4n^2\pi^2}} \\
  &= \sin 2n\pi(1+\frac12\frac{t}{4n^2\pi^2}+o(\frac{1}{n^2})) \\
  &= \sin (2n\pi + \frac{t}{4n\pi}+o(\frac1n)) \\
  &= \sin (\frac{t}{4n\pi})+o(\frac1n) \\
\end{align*}
Donc quand $n\rightarrow +\infty$ alors $f_n(t) \rightarrow 0$.
Donc $(f_n)$ converge simplement vers $0$.
  
    \item Pour $n\ge 1$, 
$$|f_n'(t)| = \frac 12 \frac{1}{\sqrt{t+4n^2\pi^2}}\cos\sqrt{t+4n^2\pi^2}
\le \frac 12 \frac{1}{\sqrt{t+4\pi^2}} \le \frac 1{4\pi}. $$
Pour $t\ge0$ fixé et $\epsilon >0$ donné, on pose $\eta = 4\pi\epsilon$, alors
par l'inégalité des accroissement finis
$$\forall n \ge 1 \quad |t-t'| < \eta \Rightarrow |f_n(t)-f_n(t')| \le \frac{1}{4\pi}|t-t'| < \epsilon.$$
Donc $(f_n)$ est une famille équicontinue.
  \end{enumerate}

  \item Notons $\mathcal{H} = \{ f_n \mid n\in \Nn^*\}$, $\mathcal{H}(t)=\{ f_n(t) \mid n\in \Nn^*\}$, 
  alors d'après la convergence simple, $\overline{\mathcal{H}(t)} = \mathcal{H}(t) \cup \{0\}$. 
Mais $(f_n)$ ne converge pas uniformément (i.e. pour la norme $\| . \|_\infty$)
vers $f=0$. En effet pour $n$ impair, posons $t_n = 5n^2\pi^2$, alors $f_n(t_n) = \sin \sqrt{9n^2\pi^2}= \sin (3n\pi) = 0$.
Pour $n$ pair, on pose $t_n= \frac{\pi^2}{4} + 2n\pi^2$ alors 
$$f_n(t_n)=\sin\sqrt{\frac{\pi^2}{4} + 2n\pi^2 + 4n^2\pi^2} = \sin \sqrt{(\frac\pi2 + 2n\pi)^2} = \sin(\frac\pi2 + 2n\pi) = 1.$$
Donc pour tout $n$, $\|f_n-f\|_\infty = 1$.
Supposons que $\mathcal{H}$ soit relativement compact alors de la suite $(f_n)$ on peut extraire une sous-suite qui converge, nécessairement la limite est $f=0$, mais comme pour tout $n$, $\|f_n-f\|_\infty = 1$, nous obtenons une contradiction.

Bien s\^ur le théorème d'Ascoli n'est pas mis en défaut, car toutes les hypothèses sont vérifiées sauf $E=[0,+\infty[$ qui n'est pas compact.
\end{enumerate}
\fincorrection
\correction{002417}
\begin{enumerate}
  \item $k$ est continue sur le compact $[a,b]\times [a,b]$ donc est uniformément continue. \'Ecrivons cette continuité uniforme dans le cas particulier o\`u les secondes coordonnées sont égales :
$$\forall \epsilon' >0 \quad \exists \eta >0 
\quad  \forall x,y,t \in [a,b] \qquad |x-y|< \eta  \Rightarrow |k(x,t)-k(y,t)| < \epsilon'.$$

  \item Comme $(f_n)$ est bornée il existe $M>0$ tel que
pour tout $n\in \Nn$, $\|f_n\|_\infty \le M$.
Fixons $x\in [a,b]$. Soit $\epsilon >0$, posons $\epsilon'=\frac{\epsilon}{M(b-a)}$, par l'uniforme continuité de $k$, on obtient un $\eta>0$ avec pour $|x-y|<\eta$, 
$|k(x,t)-k(y,t)| < \epsilon'= \frac{\epsilon}{M(b-a)}$.

Donc pour $|x-y|<\eta$,
\begin{align*}
|Kf_n(x)-Kf_n(y)| 
  &\le \int_a^b |k(x,t)-k(y,t)|\|f_n\|_\infty dt \\
  &\le M \int_a^b |k(x,t)-k(y,t)| dt \\
  &\le M \int_a^b \frac{\epsilon}{M(b-a)} dt \\
  &\le \epsilon \\
\end{align*}
Ce qui est l'équicontinuité de $(Kf_n)$ en $x$. Comme ceci est valable quelque soit $x\in[a,b]$ alors $(Kf_n)$ est équicontinue.

  \item Notons $\mathcal{H} = (Kf_n)_n$. Alors pour $x$ donné 
$\mathcal{H}(x)$ est borné car 
$|\int_a^b k(x,t)f_n(t) dt| \le M\int_a^b|k(x,t)|dt$ est bornée indépendamment de $n\in \Nn$. Donc $\overline{\mathcal{H}(x)}$ est un fermé borné de $\Rr$ donc un compact. 

Nous avons toutes les hypothèses pour appliquer le théorème d'Ascoli, donc
$\mathcal{H} = (Kf_n)_n$ est relativement compact. Donc de la suite $(Kf_n)$
on peut extraire une sous-suite convergente. (Attention la limite de cette sous-suite est dans $\overline{\mathcal{H}} \subset X$ et pas nécessairement dans $\mathcal{H}$.)
\end{enumerate}
\fincorrection
\nocorrection
\nocorrection
\nocorrection
\correction{002681}
En utilisant les relations de Cauchy, on trouve ais{\'e}ment
$$ {\partial P\over \partial r} = {1\over r}\,{\partial Q\over \partial \theta} \qquad{\rm et} \qquad
{\partial Q\over \partial r} = -{1\over r}\,{\partial P\over \partial \theta} $$
Ces conditions sont v{\'e}rifi{\'e}es par la fonction $\log$. Sa
d{\'e}riv{\'e}e est $1/z$.
\fincorrection
\correction{002682}
La fonction $f$ est m{\'e}romorphe comme compos{\'e}e de fonctions m{\'e}romorphes. Son
unique p{\^o}le est $-1$. Le contour {\'e}tant enti{\`e}rement inclus dans le domaine $\Omega $, on a
$$ \int_C f(z) \,dz = 2i\pi  \mbox{\rm Res}(f, -1) = 2i\pi  \left[\exp((a-1)\log z\right]_{z=-1}
 = 2i\pi  e^{i(a-1)\pi } =-2i\pi  e^{ia\pi }$$

Notons  $\gamma$ le petit cercle
(orient{\'e} n{\'e}gativement), $\Gamma$ le grand cercle
 (orient{\'e} positivement) dans
$C$, et remarquons que sur le segment $[\varepsilon, R]+0i$,
l'argument de $z$ doit {\^e}tre pris nul, tandis qu'il faut le
prendre {\'e}gal {\`a} $2\pi $ sur le segment oppos{\'e}
$[R,\varepsilon]-0i$. On obtient
$$\int_C f(z) \,dz = \int_\varepsilon^R {x^{a-1}\over 1+x}\,dx
+\int_\Gamma f(z)\,dz +\int_R^\varepsilon {x^{a-1} e^{i2\pi
(a-1)}\over 1+x}\,dx +\int_\gamma f(z)\,dz$$

Lorsque $\varepsilon\to 0$ et $R\to 0$, $\int_\gamma f  \to 0$ et $\int_\Gamma f \to 0$ d'apr{\`e}s les
lemmes de Jordan, car $|zf(z)|\to 0$ quand $|z|\to 0$ ou $+\infty $. Finalement
$$ \int_C f(z)\,dz \to (1-e^{i2\pi a}) I $$
et donc
$$ I = {-2i\pi  e^{ia\pi }\over 1-e^{i2\pi a}} = {\pi \over \sin a\pi }.$$
\fincorrection
\correction{002783}
Il suffit de v\'erifier que $f$ est d\'erivable au sens complexe. Pour tout
$z\neq0$: 
$$ \lim_{w\to z}\frac{f(w)-f(z)}{w-z}=\lim_{w\to z}\frac{\frac{1}{w}-\frac{1}{z}}{w-z}
=  \lim_{w\to z}\frac{1}{w-z}\left( \frac{z-w}{wz}\right)
=-\frac{1}{z^2}\;.$$
La fonction $f$ est bien holomorphe sur $\C \setminus \{0\}$ avec $f'(z)=
-\frac{1}{z^2}$.
\fincorrection
\correction{002784}
Consid\'erons le produit $fg$. En utilisant la d\'efinition m\^eme de la d\'eriv\'ee,
on a :
$$
\begin{aligned}
\frac{1}{h} \left( f(z+h)g(z+h) -f(z)g(z)\right) &= f(z+h)\frac{g(z+h)-g(z)}{h} + g(z) \frac{f(z+h)-f(z)}{h}\\
& \longrightarrow f(z) g'(z) +g(z)f'(z) \quad \text{lorsque} \;\; h\to 0 \, .
\end{aligned}$$
Autre mani\`ere:
$$\begin{aligned}
f(z+h)g(z+h)&= (f(z)+f'(z)h +h\epsilon (h))(g(z)+g'(z)h +h\epsilon (h))\\
&= f(z)g(z) + (f(z)g'(z)+f'(z)g(z))h +h \epsilon (h)\, .
\end{aligned}$$
D'o\`u $(fg)'(z)=f(z)g'(z)+f'(z)g(z)$.
\fincorrection
\correction{002785}
De la m\^eme fa\c{c}on que pour la correction de l'exercice \ref{ex:burnol1.1.2} on a
$$\begin{aligned}
\frac{f(z+h)}{g(z+h)}&=\frac{f(z)+f'(z)h +h\epsilon (h)}{g(z)\left( 1 + \frac{g'(z)}{g(z)}h +h\epsilon (h)\right)}\\
&=(f(z)+f'(z)h +h\epsilon (h))\frac{1}{g(z)} \left( 1 - \frac{g'(z)}{g(z)}h +h\epsilon (h)\right)\\
&=\frac{f(z)}{g(z)} +\frac{f'(z)g(z)-g'(z)f(z)}{g^2(z)}h +h\epsilon (h)
\end{aligned}$$
si $g(z)\neq 0$.
\fincorrection
\correction{002786}
On utilise de nouveau la d\'efinition de la d\'eriv\'ee, d'abord pour $f$ en
$z$ puis pour $g$ au point $f(z)$:
$$f(z+h)=f(z)+ f'(z)h +h \epsilon(h).$$
Notons $w_h=f'(z)h +h\epsilon(h)$. Alors (et comme dans les exercices
précédents on utilise \og epsilon\fg{} pour n'importe quelle fonction tendant
vers zéro lorsque sa variable tend vers zéro):
$$g(f(z+h))=g(f(z)+w_h) =g(f(z))+g'(f(z))w_h +w_h \epsilon(w_h).$$
Ainsi:
$$\frac{1}{h} \big(g(f(z+h))-g(f(z))\big) = \big(g'(f(z)) +
\epsilon(w_h)\big)\frac{w_h}h.$$
Lorsque $h\to0$, on a $w_h\to0$, donc $\epsilon(w_h)\to0$ et par ailleurs
$\frac{w_h}h\to f'(z)$. Au final
$$
(g\circ f)'(z) = \lim_{h\to 0} \big(g'(f(z)) +
\epsilon(w_h)\big)\frac{w_h}h = g'(f(z)) f'(z).
$$
\fincorrection
\correction{002787}
La formule de Leibniz se montre par r\'ecurrence. Le cas $n=1$, c'est-\`a-dire $(fg)'=fg'+f'g$,
a \'et\'e d\'emontr\'e dans l'exercice \ref{ex:burnol1.1.2}. Supposons alors que cette formule soit vraie au rang $n\geq 1$.
Dans ce cas,
$$\begin{aligned}
(fg)^{(n+1)}(z) &=\frac{d}{dz} \left( (fg)^{(n)}\right)(z) =
\sum_{j=0}^n  \binom{n}{j} \left\{ f^{(j+1)}(z)g^{(n-j)}(z)+f^{(j)}(z)g^{(n-j+1)}(z) \right\}\\
&=\sum_{j=1}^{n+1}  \binom{n}{j-1}f^{(j)}(z)g^{(n+1-j)}(z) +\sum_{j=0}^{n}  \binom{n}{j}f^{(j)}(z)g^{(n+1-j)}(z) .
\end{aligned}$$
La conclusion vient du fait : $\binom{n}{j-1}+\binom{n}{j}=\binom{n+1}{j}$ qui est simple \`a v\'erifier.
\fincorrection
\correction{002788}
Prenons $r<\min(R_1,R_2)$. Alors, il existe $C>0$ et $0<\lambda <1$ tels que
$|a_n|r^n \leq C\lambda^n $ et $|b_n|r^n \leq C\lambda ^n$ (v\'erifiez-le !). D'o\`u
$$\sum_{j=0}^n |a_j|r^j |b_{n-j}|r^{n-j} \leq (n+1) C^2 \lambda^n\;, $$
ce qui permet d'affirmer, pour tout $z$ avec $|z|=r$:
$$\sum_{n=0}^\infty \big(\sum_{j=0}^n |a_j z^j||b_{n-j} z^{n-j}|\big) < \infty\;.$$
Par le th\'eor\`eme du cours sur les séries doubles (voir le polycopi\'e
2005/2006 de J.-F.~Burnol, Annexe 8.2), ceci signifie que la série double
$$\sum_{j=0}^\infty\sum_{k=0}^\infty (a_j z^j b_k z^k)$$
est absolument convergente. On peut donc d'après ce théorème affirmer:
$$f(z)g(z) = \sum_{j=0}^\infty a_jz^j \sum_{k=0}^\infty b_k z^k =
\sum_{n=0}^\infty \big(\sum_{j=0}^n (a_j z^j)(b_{n-j} z^{n-j})\big)\;.$$
Or, la série de droite est $\sum_{n=0}^\infty  c_n z^n$ avec $c_n = \sum_{j=0}^n
a_j b_{n-j}$. Au passage on obtient que le rayon de convergence de cette série
est au moins
égal à $r$. Comme $r<\min(R_1,R_2)$ est arbitraire, le rayon de convergence
est en fait au moins égal à $\min(R_1,R_2)$
(il peut \^etre plus grand comme on le voit par exemple avec $f(z) =
\frac1{1-z}$ et $g(z) = 1-z$, ou encore avec $f(z) = \frac{2-z}{(1-z)(3-z)}$
et $g(z) = \frac{1-z}{(2-z)(3-z)}$).
\fincorrection
\correction{002789}
Il suffit d'utiliser la formule de Leibniz de l'exercice \ref{ex:burnol1.1.5} et le fait que le coefficient $a_n$
du d\'eveloppement de $f$ \`a l'origine est $a_n=f^{(n)}(0)/n!$.
\fincorrection
\correction{002790}
La fonction $f(z)=\overline{z}$ n'est nulle part dérivable au sens complexe
(et donc nulle part holomorphe): car
$$\frac{1}{h} (f(z+h)-f(z))=\frac{\overline{h}}{h}$$
et la limite de cette expression n'existe pas lorsque $h\to 0$.
\medskip
\emph{Remarque.} Plus g\'en\'eralement, une application $\R$--lin\'eaire de $\C$ dans $\C$ est de la forme
\begin{equation}\label{eq::1}
w\mapsto \alpha w +\beta \overline{w} 
\end{equation}
(ce n'est qu'une \'ecriture complexe des applications lin\'eaires de $\R^2$ dans $\R^2$) ; une telle application
est holomorphe si et seulement si $\beta =0$.
C'est exactement la diff\'erence entre diff\'erentiabilit\'e
(donc r\'eelle) et holomorphie (d\'erivabilit\'e au sens complexe). En effet, les \'equations de Cauchy-Riemann
sont \'equivalentes \`a l'\'equation $\frac{\partial}{\partial \overline{z}} f(z) =0$ o\`u
$\frac{\partial}{\partial \overline{z}}=\frac{1}{2} \left( \frac{\partial}{\partial x}+i \frac{\partial}{\partial y}\right)$.
C'est une r\'e\'ecriture complexe des \'equations de Cauchy-Riemann. Si vous avez une fonction $f$ diff\'erentiable,
alors sa diff\'erentielle $D f(z)$ est une application lin\'eaire de la forme \eqref{eq::1}. Un calcul simple montre que dans
ce cas
$$\beta = \frac{\partial f}{\partial \overline{z}} (z) \quad et \quad \alpha = \frac{\partial f}{\partial z} (z) $$
avec $\frac{\partial}{\partial z}=\frac{1}{2} \left( \frac{\partial}{\partial x}-i \frac{\partial}{\partial y}\right)$.
De nouveau, $f$ est complexe diff\'erentiable en $z$ si et seulement si $\beta =\frac{\partial f}{\partial \overline{z}} (z)=0$.
Dans ce cas $f'(z) = \frac{\partial f}{\partial z} (z)$.
\bigskip
Revenons \`a l'exercice. Si vous \^etes d'accord avec ma remarque, alors nous sommes aussi d'accord sur le fait que :
$$z\mapsto x =\frac{z+\overline{z}}{2}$$
n'est pas holomorphe. Ce raisonnement s'applique aussi \`a $z\mapsto y$. Nous reviendrons \`a ce
genre d'applications dans l'exercice \ref{ex:burnol1.1.10}.
\fincorrection
\correction{002791}
Pour \'eviter des raisonnements topologiques, supposons dans un premier temps que $\Omega$
soit un disque, par exemple le disque unit\'e $\Omega =D =D (0,1)$, et montrons que $f$ est constante et \'egale \`a
$f(0)$. Si $z\in D$, alors le segment $[0,z]\subset D$ (et c'est pour cette raison que l'on a pris $\Omega =D$).
On peut \'ecrire
$$f(z)-f(0)=\int_0^z f'(z) \, dz =0.$$
Seulement, ici il faut expliquer le sens de cette int\'egrale (non connue pour l'instant). Soit $\gamma : [0,1]\to [0,z]$,
$\gamma (t)=tz$, une param\'erisation du segment $[0,z]$. Alors,
$$\begin{aligned}
\int _{0} ^z f'(w)\, dw &= \int _0^1 f'(\gamma (t)) \gamma '(t) \, dt = \int _0^1 f'(tz)z\, dt\\
&= \int _0^1 \Re \big( f'(tz)z\big) \, dt + i \int _0^1 \Im \big( f'(tz)z\big) \, dt =0.
\end{aligned}$$
Pour le cas d'un ouvert connexe $\Omega$ quelconque le pr\'ec\'edent raisonnement montre qu'au voisinage de tout point
$z_0\in \Omega$ la fonction $f$ est constante. C'est donc une propri\'et\'e ouverte. Autrement dit, si $z_0\in \Omega$
est un point quelconque, l'ensemble
$$\mathcal{E} =\{z\in \Omega \, ; \;\; f(z)=f(z_0)\}$$
est un ouvert. Pour conclure il faut \'etablir que $\mathcal{E}$ est aussi un ferm\'e de $\Omega$ (topologie induite!!).
Or ceci est \'evident puisque $\mathcal{E} =f^{-1} (\{f(z_0)\})$ et $f$ est continue.
Notons que $\mathcal{E} \neq \emptyset$ puisque $z_0\in \mathcal{E}$. Les seuls ensembles \`a la fois ouverts et ferm\'es du connexe
$\Omega$ \'etant l'ensemble vide et $\Omega$, on a $\Omega = \mathcal{E}$. La fonction $f$ est constante sur $\Omega$.
Si $\Omega$ n'est pas connexe, $f$ peut  prendre diff\'erentes valeurs sur les diff\'erentes composantes
connexes de $\Omega$.
\fincorrection
\correction{002792}
Soit $f(z)=u(z)+iv(z)$ pour $z\in U$. Si $f$ ne prend que des valeurs r\'eelles, alors $v\equiv 0$. On tire des
\'equations de Cauchy-Riemann
$$\frac{\partial u}{\partial x} =\frac{\partial v}{\partial y}=0$$
$$\frac{\partial u}{\partial y}=-\frac{\partial v}{\partial x}\equiv 0.$$
La d\'eriv\'ee de $f$ est alors identiquement nulle sur l'ouvert connexe $\Omega$ ce qui implique que $f$ est constante
(voir l'exercice \ref{ex:burnol1.1.9}).
\fincorrection
\correction{002793}
\begin{enumerate}
  \item La formule de Taylor avec reste int\'egral est
$$g(b)= g(a) +g'(a)(b-a) +...+ \frac{g^{(n)} (a)}{n!} (b-a)^n+ \int_a^b \frac{(b-t)^n}{n!} g^{(n+1)}(t)\, dt$$
puis on remplace avec $a=0$ et $b=1$.
  \item Si $f'=f$ et si $f$ est $n$--fois d\'erivable au sens complexe, alors
$\lim_{h \to 0}(f^{(n)} (z+h) -f^{(n)} (z))/h
=\lim_{h \to 0}(f^{(n-1)} (z+h) -f^{(n-1)} (z))/h=f^{(n)}(z)$.
Par r\'ecurrence on en d\'eduit, d'une part, que $f$ est infiniment d\'erivable et, d'autre part, que $f^{(n)}(z)=f(z)$
pour tout $n\geq 0$. En particulier, $f^{(n)} (0)=1$ pour tout $n\geq 0$. En utilisant la formule de Taylor de la question précédente
on a donc
$$\begin{aligned}
\left| f(z) - \sum_{k=0}^n \frac{z^k}{k!} \right| &\leq |z|^{n+1} \int_0^1 \frac{(1-u)^n}{n!} |f^{(n+1)} (uz)| \, du\\
&\leq |z|^{n+1} \sup_{|w|\leq |z|} |f^{(n+1)} (w)|\int_0^1 \frac{(1-u)^n}{n!}\, du
\leq \sup_{|w|\leq |z|} |f (w)|\frac{|z|^{n+1} }{(n+1)!}.
\end{aligned}$$
Cette derni\`ere expression tend vers $0$ lorsque $n\to \infty$. D'o\`u $f(z)=\sum_{k=0}^\infty \frac{z^k}{k!}$.
  \item Fixons $z\in \Cc$ et notons $a_k=\frac{z^k}{k!}$. Alors :
$$ \left| \frac{a_{k+1}}{a_k}\right| = |z| \frac{1}{k+1} \to 0 $$
lorsque $k\to \infty$. On en d\'eduit que le rayon de convergence de cette s\'erie est $\infty$ (d'Alembert) et que
$F$ est holomorphe sur $\C$. De plus :
$$F'(z) = \sum_{k\geq 1} \frac{kz^{k-1}}{k!}=F(z)$$
pour tout $z\in \Cc$. Par le th\'eor\`eme sur les s\'eries doubles (en fait l'exercice \ref{ex:burnol1.1.6})
$$\begin{aligned}
F(z)F(w) &= \sum_{k=0}^\infty \sum_{j=0}^k \frac{z^j}{j!}\frac{w^{k-j}}{(k-j)!}
=\sum_{k=0}^\infty \frac{1}{k!} \sum_{j=0}^k \frac{k!}{j! (k-j)!} z^j w^{k-j}\\
&= \sum_{k=0}^\infty \frac{1}{k!} (z+w)^k =F(z+w).
\end{aligned}$$
\end{enumerate}
\fincorrection
\correction{002794}
Non $\limsup_{n\to\infty} |a_nz^n| =+\infty$ mais il n'y a pas de raison pour que
$\liminf_{n\to\infty} |a_nz^n|=+\infty $. Prenez par exemple la série : $\sum_{n=0}^\infty  z^{2n}$, de rayon de convergence $R=1$,
mais $(|a_nz^n|)$ n'a pas de limite (la valeur est $0$ pour $n$ impair et
$|z|^n$ pour $n$ pair, qui tend vers l'infini lorsque $|z|>1$).
\fincorrection
\correction{002795}
$$\frac{1}{1-z} =\sum_{k=0}^\infty z^k \quad \quad |z|<1.$$
$$\frac{1}{(1-z)^2}=\left( \frac{1}{1-z}\right)' = \sum_{k=1}^\infty k z^{k-1} \quad \text{pour} \quad |z|<1.$$
etc.
\fincorrection
\correction{002796}
Discutons d'abord le rayon de convergence. D'ailleurs, ce qui suit s'applique \'egalement aux exercices suivants.
Donc, d'apr\`es le th\'eor\`eme d'analycit\'e des fonctions holomoephes 
(voir le polycopi\'e 2005/2006 de J.-F.~Burnol : th\'eor\`eme 10 du chapitre 6), 
si $f$ est holomorphe dans $U\subset \C$, si $z_0\in U$ et si $r>0$ tel que
$D(z_0,r)\subset U$, alors la s\'erie de Taylor de $f$ en $z_0$ converge et sa somme vaut $f$ dans ce disque $D(z_0,r)$.
Ici $f(z)=\frac{1}{z-1}$. Cette fonction est holomorphe dans $U=\C\setminus \{1\}$. Par cons\'equent, si $z_0\in U$,
alors la s\'erie de Taylor de $f$ en $z_0$ vaut $f$ dans le disque $D(z_0,R_1)$ si $R_1=|z_0-1|$.
Le calcul de la s\'erie est classique:
$$\frac{1}{z-1}=\frac{1}{z-z_0+z_0-1} =\frac{1}{z_0-1}\frac{1}{1-\left(\frac{z_0-z}{z_0-1}\right)}
 =\frac{1}{z_0-1}\sum_{k\geq 0} \left(\frac{z_0-z}{z_0-1}\right)^k$$
 pour $|z-z_0|< |z_0-1|=R_1$.
\fincorrection
\correction{002797}
La fonction $f(z)=\frac{1}{(z-1)(z-2)}$ est holomorphe dans $U=\C\setminus\{1,2\}$.
Par ce que l'on vient de dire \`a l'exercice pr\'ec\'edent, 
le rayon de convergence demand\'e est $R=\min \{|z_0-1|,|z_0-2| \}$,
o\`u $z_0\in U$ est un point quelconque fix\'e.
On a $\frac{1}{(z-1)(z-2)}=\frac{1}{z-2}-\frac{1}{z-1}$ et:
$$\frac{1}{z-2}=\frac{1}{z-z_0+z_0-2}=\frac{1}{z_0-2}\frac{1}{1-\left( \frac{z_0-z}{z_0-2}\right)}
=\frac{1}{z_0-2}\sum_{k\geq0} \left( \frac{z_0-z}{z_0-2}\right)^k$$
pour $|z-z_0|<|z_0-2|=R_2$. La s\'erie demand\'ee est alors la diff\'erence entre celle-ci et celle de
l'exercice pr\'ec\'edent. Notons aussi que le rayon de convergence est exactement le minimum des rayons
$R_1$ et $R_2$.
\fincorrection
\nocorrection
\correction{002799}
Le rayon de convergence est $R=1$. Soit $0<t<1$ et \'etudions $f(t)=\sum_{k=0}^\infty t^{2^k}$.
Il s'agit d'une s\'erie de termes positifs. D'o\`u
$$f(t)\geq \sum_{k=0}^{N-1} t^{2^k} \quad \text{pour tout} \quad N\in \N .$$
Il en r\'esulte  $\liminf_{t\to 1} f(t) \geq N$ or $N$ est arbitraire, donc $\lim_{t\to 1} f(t)=\infty$.
Soit maintenant $w$ un nombre complexe du cercle unit\'e v\'erifiant $w^{2^N}=1$ pour un $N\in \N$.
Dans ce cas $w^{2^k}=1$ pour tout $k\geq N$. Si de nouveau $0<t<1$, alors
$$f(tw) = \sum_{k=0}^{N-1} (tw)^{2^k} + \sum_{k\geq N} t^{2^k}  .$$
Lorsque $t\to 1$, alors la premi\`ere somme tend vers un nombre complexe (fini, en fait de module
au plus $N$) et la deuxi\`eme vers $\infty$.
Les nombres complexes $w$ ayant la propri\'et\'e $w^{2^N}=1$ pour un certain $N\in \N$ sont denses dans le
cercle unit\'e $\{|z|=1\}$. Ceci, et le principe de prolongement analytique, interdit l'existence de la fonction
$g$ holomorphe sur $U$ comme d\'ecrit dans l'exercice.
Si $z_0\in D (0,1)$, alors le rayon de convergence de la s\'erie de Taylor de $f$ en $z_0$ est $R=1-|z_0|$.
\fincorrection
\nocorrection
\nocorrection
\correction{002802}
Il s'agit de formules bien connues lorsque les arguments $z,w$
sont r\'eels. La v\'erification \`a partir des d\'efinitions des fonctions trigonom\'etriques donn\'ees
dans l'\'enonc\'e de l'exercice est laiss\'ee au lecteur.
Voici comment obtenir la formule :
\begin{equation}\label{2} \cos(z+w) = \cos(z)\cos(w)-\sin(z)\sin(w) \qquad pour \quad z,w\in \C \, .
\end{equation}
Fixons $w\in \R$. Soit $f_w(z)=\cos (z+w)-\big( \cos(z)\cos(w)-\sin(z)\sin(w)\big)$.
La formule \eqref{2} \'etant vraie pour $z,w\in \R$, $f_w(z)=0$ pour tout $z\in \R$. Il r\'esulte du principe des z\'eros
isol\'es que $f_w$ est identiquement nulle. Autrement dit, on vient d'\'etablir la formule \eqref{2}
pour $(z,w)\in \C\times \R$. Il suffit maintenant de refaire le m\^eme argument en fixant d'abord $z\in \C$
arbitrairement et en observant que la fonction holomorphe
$$g_z(w)= \cos(z+w) - ( \cos(z)\cos(w)-\sin(z)\sin(w))$$
est nulle pour tout $w\in\R$. De nouveau $g_z\equiv 0$ par le principe des z\'eros isol\'es, d'o\`u la formule \eqref{2}
pour tout $z,w\in \C$.
\fincorrection
\correction{002803}
$$\begin{aligned}
\sin(a)ch(b)+i\, \cos(a)\sh(b)&=\frac{1}{4i} \big\{ (e^{ia}-e^{-ia})(e^{b}+e^{-b})-(e^{ia}+e^{-ia})(e^{b}-e^{-b})\big\}\\
&=\frac{1}{2i}\big(e^{ia-b}-e^{-ia+b}\big)=\sin(a+ib).
\end{aligned}$$
Si $a,b\in \R$, alors :
$$\begin{aligned}
|\sin(a+ib)|^2&= (\sin(a)\ch(b))^2+(\cos(a)\sh(b))^2\\
&= \sin^2(a)(1+\sh^2(b))+(1-\sin^2(a))\sh^2(b)\\
&= \sin^2(a) + \sh^2 (b).
\end{aligned}$$
Cette somme de carrés de nombres réels ne peut être nulle que si $\sin(a)=0$
et $\sh(b)=0$, c'est-à-dire $a\in\pi\Zz$ et $b=0$. Donc $\sin(z)=0\iff z\in\pi\Zz$.
\fincorrection
\nocorrection
\nocorrection
\nocorrection
\nocorrection
\nocorrection
\nocorrection
\nocorrection
\nocorrection
\nocorrection
\nocorrection
\nocorrection
\nocorrection
\nocorrection
\nocorrection
\nocorrection
\nocorrection
\nocorrection
\nocorrection
\nocorrection
\nocorrection
\nocorrection
\nocorrection
\nocorrection
\nocorrection
\nocorrection
\nocorrection
\nocorrection
\nocorrection
\nocorrection
\nocorrection
\nocorrection
\nocorrection
\nocorrection
\nocorrection
\nocorrection
\nocorrection
\nocorrection
\nocorrection
\nocorrection
\nocorrection
\nocorrection
\nocorrection
\nocorrection
\correction{007214}
Ici on a \(f(x+iy)=u(x,y)+iv(x,y)\) où \(u(x,y)=\sqrt{|xy|}\) et \(v(x,y)=0\) pour tout \((x,y)\in \R^2\).
\begin{enumerate}
\item On a \(\frac{\partial v}{\partial x}\equiv 0\) et \(\frac{\partial v}{\partial y}\equiv 0\) sur \(\R^2\).
Au point \((0,0)\in \R^2\) on a 
\[ \lim_{x\to 0}\frac{u(x,0)-u(0,0)}{x}\lim_{x\to 0}\frac{0-0}{x}=0 \]
et
\[ \lim_{y\to 0}\frac{u(0,y)-u(0,0)}{x}\lim_{x\to 0}\frac{0-0}{x}=0. \]
Donc \(u\) est dérivable par rapport à la première et à la seconde variable en \((0,0)\) et de plus \(\frac{\partial u}{\partial x}(0,0)=0\) et \(\frac{\partial u}{\partial y}(0,0)=0\).
En particulier \(f\) vérifie les équations de Cauchy-Riemann.
\item La fonction \(f\) n'est pas \(\C\)-dérivable en \(0\) car pour \(h\in \R_+^*\) on a 
\[ \frac{f(h)-f(0)}{h}=\frac{0-0}{h}=0 \]
et
\[ \frac{f(h+ih)-f(0)}{h+ih}=\frac{\sqrt{|h^2|}}{h+ih}=\frac{h}{h+ih}=\frac{1}{1+i}.\]
Les deux quantités n'ont pas la même limite lorsque $h\to 0$.
\item Ce n'est pas une contradiction avec le théorème de Cauchy-Riemann car la fonction \(u\) (et donc la fonction \(f\)) n'est pas différentiable en \((0,0)\). En effet si \(u\) était différentiable en \((0,0)\) on aurait 
\begin{align*}
u(s,t)
& = u(0,0)+du_{(0,0)}(s,t)+o(\|(s,t)\|)\\
& = \frac{\partial u}{\partial x}(0,0)+t\frac{\partial u}{\partial x}(0,0)+o(\|(s,t)\|) \\
& = o(\|(s,t)\|).
\end{align*}
Or \(u(s,s)=|s|\) qui n'est pas un \(o(\|(s,s)\|)\) d'où une contradiction.
\end{enumerate}
\fincorrection
\nocorrection
\nocorrection
\nocorrection
\nocorrection
\nocorrection
\nocorrection
\nocorrection
\nocorrection
\nocorrection
\nocorrection
\nocorrection
\nocorrection
\nocorrection
\nocorrection
\nocorrection
\nocorrection
\nocorrection
\nocorrection
\nocorrection
\nocorrection
\nocorrection
\nocorrection
\nocorrection
\nocorrection
\nocorrection
\nocorrection
\nocorrection
\nocorrection
\nocorrection
\nocorrection
\nocorrection
\nocorrection
\nocorrection
\nocorrection
\nocorrection
\nocorrection
\nocorrection
\nocorrection
\nocorrection
\nocorrection
\nocorrection
\nocorrection
\nocorrection
\nocorrection
\nocorrection
\nocorrection
\nocorrection
\nocorrection
\nocorrection
\nocorrection
\nocorrection
\nocorrection
\nocorrection
\nocorrection
\nocorrection
\nocorrection
\nocorrection
\nocorrection
\nocorrection
\nocorrection
\correction{002806}
On a les \'equations de Cauchy-Riemann $\frac{\partial u}{\partial x} = \frac{\partial v}{\partial y}$
et $\frac{\partial u}{\partial y} = -\frac{\partial v}{\partial x}$. D'o\`u :
$$\Delta u = \frac{\partial }{\partial x}\frac{\partial u}{\partial x}+ \frac{\partial }{\partial y}\frac{\partial u}{\partial y}=
\frac{\partial }{\partial x}\frac{\partial v}{\partial y}- \frac{\partial }{\partial y}\frac{\partial v}{\partial x}=0$$
par le th\'eor\`eme de Schwarz. On proc\`ede de la m\^eme fa\c{c}on pour $v$ pour en d\'eduire
qu'une fonction holomorphe est harmonique.
\fincorrection
\correction{002807}
Cet exercice et les suivants concernent des changements de variables. Rappelons que, si
$\Phi:U\to V$ est un diff\'eomorphisme entre ouverts $U,V$ de $\R ^n$ et si on note $y=\Phi (x)$,
$x=(x_1,...,x_n)\in U$ et $y=(y_1,...,y_n)\in V$, alors
$$\frac{\partial }{\partial x_i}=\sum_{j=1}^n \frac{\partial \Phi _j}{\partial x_i}\frac{\partial }{\partial y_j}$$
pour tout $i,j\in \{1,...,n\}$.
\bigskip
On a $x=r\, \cos \theta$ et $y=r\, \sin \theta$. Donc
$$\frac{\partial }{\partial r}= \frac{\partial (r\, \cos \theta)}{\partial r}\frac{\partial }{\partial x}+\frac{\partial (r\, \sin \theta)}{\partial r}\frac{\partial }{\partial y}=
\cos  \theta \frac{\partial }{\partial x} + \sin \theta \frac{\partial }{\partial y}$$
et
$$\frac{\partial }{\partial \theta}= \frac{\partial (r\, \cos \theta)}{\partial \theta}\frac{\partial }{\partial x}+
\frac{\partial (r\, \sin \theta)}{\partial \theta}\frac{\partial }{\partial y}=
-r\sin  \theta \frac{\partial }{\partial x} + r \cos \theta \frac{\partial }{\partial y}.$$
On peut r\'e\'ecrire ceci en:
$$\binom{\frac{\partial }{\partial r}}{\frac{\partial }{\partial \theta}}=M \binom{\frac{\partial }{\partial x}}{\frac{\partial }{\partial y}}
\quad \text{avec} \quad M= (Jac(f))^t = \left(
                                   \begin{array}{cc}
                                     \frac{\partial (r\, \cos \theta)}{\partial r} & \frac{\partial (r\, \sin \theta)}{\partial r} \\
                                     \frac{\partial (r\, \cos \theta)}{\partial \theta} & \frac{\partial (r\, \sin \theta)}{\partial \theta} \\
                                   \end{array}
                                 \right) .$$
Pour retrouver les $\frac{\partial }{\partial x}$, $\frac{\partial }{\partial y}$ en fonction de $\frac{\partial }{\partial r}$, $\frac{\partial }{\partial \theta}$
il suffit d'inverser cette matrice $M$. Le reste est clair.
\fincorrection
\correction{002808}
On a $w=g(z)$ avec $g(z)=a(z)+ib(z)$ une fonction holomorphe et avec $z=x+iy$. Utilisons de nouveau le changement
de coordonn\'ees:
$$\frac{\partial }{\partial x}= \frac{\partial a(z)}{\partial x}\frac{\partial }{\partial a}+\frac{\partial b(z)}{\partial x}\frac{\partial }{\partial b}$$
$$\frac{\partial }{\partial y}= \frac{\partial a(z)}{\partial y}\frac{\partial }{\partial a}+\frac{\partial b(z)}{\partial y}\frac{\partial }{\partial b}.$$
En utilisant les \'equations de Cauchy-Riemann on en d\'eduit
$$\begin{aligned}
\frac{\partial }{\partial x}+i\frac{\partial }{\partial y}&= \left(\frac{\partial a(z)}{\partial x}+i \frac{\partial a(z)}{\partial y} \right)\frac{\partial }{\partial a}
+ \left(\frac{\partial b(z)}{\partial x}+i \frac{\partial b(z)}{\partial y} \right)\frac{\partial }{\partial b}\\
&= \left(\frac{\partial a(z)}{\partial x}-i \frac{\partial b(z)}{\partial x} \right)\frac{\partial }{\partial a}
+ \left(\frac{\partial b(z)}{\partial x}+i \frac{\partial a(z)}{\partial x} \right)\frac{\partial }{\partial b}\\
&=\left(\frac{\partial a(z)}{\partial x}-i \frac{\partial b(z)}{\partial x} \right)\left(\frac{\partial }{\partial a}+i\frac{\partial }{\partial b}\right).
\end{aligned}$$
\fincorrection
\nocorrection
\correction{002810}
Soit $Q$ le carr\'e dont le bord est $\gamma = \gamma_1+...+\gamma_4$ o\`u
$$\gamma_1 (t)=A+2it , \quad \gamma_2 (t)=B-2t , \quad \gamma_3 (t)=C-2it , \quad et \quad \gamma_4 (t)=D+2t, \quad t\in [0,1] .$$
Notons aussi $\gamma_{j,x}=\Re (\gamma_j )$ et $\gamma_{j,y}=\Im (\gamma_j )$, $j=1,...,4$. Alors :
$$\int_\gamma dx= \sum_{j=1}^4 \int_{\gamma_j}dx= \sum_{j=1}^4 \int_0^1 \gamma'_{j,x}(t) \, dt=0$$
et $$\begin{aligned}
\int_\gamma x\, dx &=\sum_{j=1}^4 \int_{\gamma_j}x\, dx = \sum_{j=1}^4 \int_0^1 \gamma_{j,x}(t)\gamma'_{j,x}(t)\, dt\\
&= \int_0^1 (1-2t)(-2)dt +\int_0^1 (-1+2t)2dt=0.
\end{aligned}$$
\bigskip
Passons \`a la correction de la question 3. Alors
$$\int _\gamma dz =\int _\gamma z\, dz =0$$
puisque dans les deux cas on int\`egre une fonction holomorphe $(f(z)\equiv 1$ et $f(z)=z$) dans le carr\'e $Q$. On a :
$$\begin{aligned}
\int_\gamma x\, dz &=\sum_{j=1}^4 \int_0^1 \gamma_{j,x}(t)d\gamma_{j}(t) =\sum_{j=1}^4 \int_0^1 \gamma_{j,x}(t)\gamma'_{j}(t)\, dt\\
&=\int_0^1 2idt +\int_0^1 (1-2t)(-2)dt +  \int _0^1 (-1)(-2i)dt +\int_0^1 (-1+2t)2dt\\
&= 2i +2i =4i.
\end{aligned}$$
En ce qui concerne la question 4., on y int\`egre la fonction $f_n(z)=z^n$ le long du chemin ferm\'e $\gamma$.
Mais attention, cette fonction admet une primitive seulement si $n\neq -1$. D'o\`u
$$\int_\gamma z^n dz =0 \quad \text{pour } \; n\neq -1.$$
Dans le cas restant $n=-1$ on trouve :
$$\int_\gamma f_{-1}(z)dz =2i\pi.$$
D'ailleurs, et l\`a on rejoint l'exercice \ref{ex:burnol2.2.3}, on a :
$$\int_\gamma f_n(z)\, dz = \int_C f_n(z)\, dz $$
o\`u $C=\{ |z|=1\}$. Ce cercle se param\'etrise par $\sigma (\theta ) =e^{2i\pi \theta}$. D'o\`u :
$$\int_C f_n(z)\, dz=\int_0^1 e^{2i\pi n\theta} 2i\pi e^{2i\pi \theta} \, d\theta
=2i\pi \int_0^1 e^{2i\pi (n+1) \theta}\, d\theta =\left\{ \begin{array}{c}
                                                            2i\pi \;\; \text{si} \;\; n=-1 \\
                                                            0 \quad \text{sinon}.
                                                          \end{array}
\right.$$
De mani\`ere analogue on a
$$\int_C \overline{z}^n\, dz=\int_0^1 e^{-2i\pi n\theta} 2i\pi e^{2i\pi \theta} \, d\theta
=2i\pi \int_0^1 e^{2i\pi (1-n) \theta}\, d\theta =\left\{ \begin{array}{c}
                                                            2i\pi \;\; \text{si} \;\; n=1 \\
                                                            0 \quad \text{sinon.}
                                                          \end{array}
\right.$$
\fincorrection
\correction{002811}
Pour toute fonction $f=u+iv$ \`a valeurs dans $\C$,  on a :
$$\int _\gamma \overline{f(z)} dz =\int_0^1 (u-iv)(\gamma'_1+i\gamma'_2)dt =\overline{\int_0^1 (u+iv)(\gamma'_1-i\gamma'_2)dt}=
\overline{\int_\gamma f(z) \overline{dz}}.$$
Voir la correction de l'exercice \ref{ex:burnol2.2.1}.
\fincorrection
\correction{002812}
Voir la correction de l'exercice \ref{ex:burnol2.2.1}.
\fincorrection
\correction{002813}
Dans le cas o\`u $C$ n'encercle pas l'origine la fonction $z\mapsto z^n$ est holomorphe au voisinage du disque
bord\'e par $C$ et ceci pour tout $n\in \Z$. Par cons\'equent, $\int _C z^n dz=0$. Sinon on retrouve les valeurs
obtenues pr\'ec\'edemment. On rappelle \'egalement que :
$$\frac{1}{2i\pi}\int _C \frac{1}{z} dz$$
est l'indice $\mathrm{Ind}(C,0)$ de la courbe $C$ par rapport \`a l'origine. Cet indice est $1$ lorsque $C$ encercle
l'origine et $0$ sinon.
\fincorrection
\correction{002814}
On a
$$\frac{1}{(z-a)(z-b)}=\frac{1}{a-b}\left( \frac{1}{z-a}-\frac{1}{z-b}\right)$$
$$\int _C \frac{1}{z-a}dz = 2i\pi \mathrm{Ind} (C,a).$$
Or, $\mathrm{Ind}(C,a)=0$ si $r<a$ et $\mathrm{Ind}(C,a)=1$ si $r>a$. Le m\^eme raisonnement s'applique \`a $\int_C \frac{1}{z-b}dz$,
d'o\`u le r\'esultat annonc\'e.
\fincorrection
\correction{002815}
Comme
$$\left( z+\frac{1}{z} \right)^n = \sum_{k=0}^n \binom{n}{k}z^k z^{k-n}$$
on a
$$\left( z+\frac{1}{z} \right)^n \frac{1}{z}= \sum_{k=0}^n \binom{n}{k} z^{2k-n-1}.$$
Or $\int_C z^j \, dz \neq 0$ si et seulement si $j=-1$. Le seul terme de la somme pr\'ec\'edente qui donne une
contribution non nulle \`a l'int\'egrale est lorsque $k$ v\'erifie $2k-n-1=-1$.
Notons que ceci est possible seulement si $n$ est un nombre pair! D'o\`u :
$$\int_C \left( z+\frac{1}{z} \right)^n \frac{dz}{z}=0 \quad \text{si $n$ est impair}.$$
Sinon, si $n=2k$ est pair, on a :
$$\int_C \left( z+\frac{1}{z} \right)^{2k} \frac{dz}{z}=2i\pi \binom{2k}{k}.$$
Comme $\cos \, t=\frac{1}{2}(e^{it}+e^{-it})$,
$$\int_{-\pi} ^\pi \cos ^n t \, dt =\frac{1}{2^n} \int_{-\pi} ^\pi(e^{it}+e^{-it})^n \, \frac{ie^{it}dt}{ie^{it}}=
\frac{-i}{2^n}\int_C \left( z+\frac{1}{z} \right)^{n} \frac{dz}{z}.$$
D'o\`u $\int_{-\pi}^\pi \cos\, n t dt=0$ si $n$ est impair et
$$\int_{-\pi}^\pi \cos ^{2k}t\, dt =\frac{\pi}{2^{2k-1}}\binom{2k}{k}.$$
Par p\'eriodicit\'e du cosinus ceci donne :
$$I_k =\int_0^{\pi /2} \cos ^{2k}t\, dt =\frac{1}{4} \frac{\pi}{2^{2k-1}}\binom{2k}{k}=\frac{\pi}{2}
\frac{1\cdot 3\cdot ...\cdot (2k-1)}{2\cdot 4\cdot ...\cdot 2k}.$$
\fincorrection
\nocorrection
\nocorrection
\nocorrection
\nocorrection
\nocorrection
\nocorrection
\nocorrection
\nocorrection
\nocorrection
\nocorrection
\nocorrection
\nocorrection
\nocorrection
\nocorrection
\nocorrection
\nocorrection
\nocorrection
\nocorrection
\nocorrection
\nocorrection
\nocorrection
\nocorrection
\nocorrection
\nocorrection
\nocorrection
\nocorrection
\nocorrection
\nocorrection
\nocorrection
\nocorrection
\nocorrection
\nocorrection
\nocorrection
\nocorrection
\nocorrection
\nocorrection
\nocorrection
\nocorrection
\nocorrection
\nocorrection
\nocorrection
\nocorrection
\nocorrection
\nocorrection
\nocorrection
\nocorrection
\nocorrection
\nocorrection
\nocorrection
\nocorrection
\nocorrection
\nocorrection
\nocorrection
\nocorrection
\nocorrection
\nocorrection
\nocorrection
\nocorrection
\nocorrection
\nocorrection
\nocorrection
\nocorrection
\nocorrection
\nocorrection
\nocorrection
\nocorrection
\nocorrection
\nocorrection
\nocorrection
\nocorrection
\nocorrection
\nocorrection
\nocorrection
\nocorrection
\nocorrection
\nocorrection
\nocorrection
\correction{002680}
On pose $f(z) = e^{xz}$. Alors $f^{(n)} (0)= x^n$, donc
$$ x^n = {n!\over 2i\pi } \int_C {e^{xz}\over z^{n+1}}\,dz $$
d'o{\`u} la formule demand{\'e}e par multiplication par $x^n/(n!)^2$.

Sur $C$, on a $|z|=1$, et donc la s{\'e}rie $\sum x^n /(n! z^{n})$ est
uniform{\'e}ment convergente par rapport {\`a} $z$ (et sa limite est bien s{\^u}r {\'e}gale {\`a}
$e^{x/z}$). Donc on peut inverser les signes $\sum$ et $\int$:
$$\begin{array}{cccccc}\sum_{n=0}^{+\infty } \left(x^n\over n!\right)^2
&= {1\over 2i\pi }\sum_{n=0}^{+\infty } \int_C {e^{xz}\over z}\,\left({x^n\over n! z^n}\right)\,dz \cr
&= {1\over 2i\pi } \int_C {1\over z}\,e^{xz}\,e^{x\over z}\ dz \cr
&= {1\over 2i\pi } \int_0^{2\pi } e^{-i\theta} e^{x(e^{i\theta}+e^{-i\theta})} \,{i e^{i\theta}\,d\theta}
 &(z = e^{i\theta})\cr
&= {1\over 2\pi } \int_0^{2\pi } e^{2x\cos\theta}\,d\theta.
\end{array}$$
\fincorrection
\nocorrection
\nocorrection
\nocorrection
\nocorrection
\nocorrection
\nocorrection
\nocorrection
\nocorrection
\nocorrection
\nocorrection
\nocorrection
\nocorrection
\nocorrection
\nocorrection
\correction{002672}
On pose $\theta=2\varphi$, et on obtient
$$ I = \int_0^{2\pi } {a\,d\theta\over 2a^2+1 -\cos \theta}
 = \int_C f(z)\,dz $$
o{\`u} $c$ est le cercle trigonom{\'e}trique, $z= e^{i\theta}$ et
$$f(z) = {2ia\over z^2 -2(2a^2+1)z +1} = {2ia\over P(z)}$$
qui a deux p{\^o}les simples (racines du d{\'e}nominateur $P$), 
dont une seule est int{\'e}rieure
au cercle, soit
$$ \alpha = 2a^2+1 -2a\sqrt{a^2+1}$$
Donc
$$ I = 2\pi i  \mbox{\rm Res}(f, \alpha) = 2\pi i {2ia\over P'(\alpha)}
 = {-4\pi a\over 2\alpha -2(2a^2+1)} = {\pi \over \sqrt{a^2+1}}.$$
\fincorrection
\correction{002673}
Le polyn{\^o}me  $P(z) = z^4+z^2+1$ a pour racines les racines carr{\'e}es de $j$ et $j^2$,
soit $\pm j^2$ et $\pm j$. Seuls $j$ et $-j^2$ ont une partie imaginaire positive, donc
$$ I = 2i\pi  (\mbox{\rm Res}(1/P, j) + \mbox{\rm Res}(1/P, -j^2))
 = 2i\pi  \left({1\over P'(j)} + {1\over P'(-j^2)}\right) = {\pi \over \sqrt 3}.$$

Puis on a
$$ J = {1\over 2} \int_{-\infty }^{+\infty } {1-\cos 2x\over x^4+x^2+1}\,dx
 = {1\over 2} I -{1\over 2} \int_{-\infty }^{+\infty } {\cos 2x\over x^4+x^2+1}\,dx
 = (I-K)/2 $$
On a par la m{\^e}me m{\'e}thode
$$ K = 2i\pi  \left({e^{2i j}\over P'(j)} + {e^{2i (-j^2)}\over P'(-j^2)}\right) $$
Or
$$ e^{2ij} = \exp\left(2i \left(1 +i\sqrt 3 \over 2\right)\right) = e^{-\sqrt3 -i}$$
et de m{\^e}me $e^{2i(-j^2)}= e^{-\sqrt 3+i}$. Finalement
$$ J = {\pi \over 2\sqrt 3} \left[ 1 -e^{-\sqrt3} \left(\sqrt 3\sin 1 +\cos1\right)\right].$$
\fincorrection
\correction{002674}
On a, sachant que $2\cos\theta = e^{i\theta}+e^{-i\theta}$:
$$ W_n = {1\over 4}\int_0^{2\pi } \cos^{2n}\theta\,d\theta
 = {1\over 2^{2n+2}} \int_C \left(z+{1\over z}\right)^{2n}\ {dz\over iz}$$
o{\`u} $C$ est le cercle trigonom{\'e}trique et $z=e^{i\theta}$. En posant
$$f(z) = {1\over z} \left(z+{1\over z}\right)^{2n} $$
on a imm{\'e}diatement
$$ W_n = {\pi \over 2^{2n+1}} \mbox{\rm Res}(f, 0) $$
car $0$ est le seul p{\^o}le de $f$. Le r{\'e}sidu est en fait le coefficient en $1/z$ dans
le d{\'e}veloppement de $f$, c'est-{\`a}-dire le coefficient constant dans $(z+1/z)^{2n}$,
soit par la formule du bin{\^o}me, $C_{2n}^n$. Donc
$$ W_n =  {\pi \over 2^{2n+1}} {(2n)!\over (n!)^2}.$$
\fincorrection
\correction{002675}
% correction ajoutée le 18/11/2018
La fonction est paire, donc ses coefficients $b_n$ sont nuls et on calcule ici ses coefficients $a_n$.

Pour $n\geq 1$, on a : 
\[ a_n 
= \frac{1}{\pi} \int_0^{2\pi}f(t)\cos(nt)dt
= \frac{1}{\pi} \int_0^{2\pi}\frac{\cos(nt)}{2-\cos(t)}dt
\]

On peut écrire 
\begin{align*}
a_n &= \Re \frac{1}{\pi} \int_0^{2\pi}f(t)e^{int}dt\\
&= \Re \frac{1}{\pi} \int_{\mathcal C}\frac{z^n}{2-\frac12 \left(z+z^{-1}\right)}\frac{dz}{iz}\\
&= \Re \frac{1}{i\pi} \int_{\mathcal C}\frac{2z^n}{4z-z^2-1}dz\\
&= 4\Re \sum_\alpha Res\left(\frac{z^n}{4z-z^2-1},\alpha\right)
\end{align*}
Comme $n\geq 1$, il n'y a qu'un seul pôle dans le disque : le pôle simple $2-\sqrt 3$. Le résidu en ce pôle vaut $\frac{(2-\sqrt 3)^n}{2\sqrt 3}$ qui est  déjà un nombre réel.\\

On en déduit que 
\[ a_n = \frac{2}{\sqrt 3}(2-\sqrt 3)^n.\]

Le coefficient $a_0$ est obtenu de la même façon, mais il faut diviser par deux, vu la définition particulière lorsque $a=0$. On obtient donc $a_0 = \frac{1}{\sqrt 3}$.

Finalement, on a donc :
$$ f(x) = {1\over\sqrt 3}\left[1+ 2\sum_{n=0}^{+\infty } 
(2-\sqrt 3)^n \cos nx\right].$$
\fincorrection
\correction{002676}
On utilise 
$$ \cos(x+iy) = \cos x \cosh y -i \sin x \sinh y$$
Ici on doit avoir $\cos x \cosh y = a$ et $\sin x \sinh y = 0$. La derni{\`e}re {\'e}quation
implique $x= 0$ ou~$\pi $, ou ${y=0}$. La solution $y=0$ ne convient pas, car aucun r{\'e}el
n'a pour cosinus $a>1$. Donc il faut $x=0$, et $\cosh y = a$ ou bien $x=\pi $ et $\cosh
y=-a$. Cette derni{\`e}re {\'e}quation est sans solutions ($\cosh y\geq 1$ pour tout $y$). Les
solutions sont donc $z=\pm i\arg\cosh a$. Leur sinus est $\sin x \cosh y+i\cos x \sinh y = i\sinh
y$ ici, soit $\pm i\sqrt {a^2-1}$ (car $\sinh^2 y = \cosh^2 y -1$).

On a donc, en posant $f(z) = 1/[(1+z^2)(2-\cos z)]$ (fonction paire):
$$ I_a = i\pi  \mbox{\rm Res}(f, i) + i\pi  \mbox{\rm Res}(f, i\arg\cosh a) $$
car $i$ et $i\arg\cosh a$ sont les deux seuls p{\^o}les de $f$ de partie imaginaire $>0$.
Puis
$$\begin{array}{ccc}
 \mbox{\rm Res}(f,i) &= \left(1\over 2z(a-\cos z)\right)_{z=i} = {1\over 2i (a-\cosh 1)}\cr
\mbox{\rm Res}(f, i\arg\cosh a) &= \left(1\over (1+z^2) \sin z\right)_{z=i\arg\cosh a} =
 {1\over i(1 -\arg\cosh^2 a)\sqrt{a^2-1}}\cr\end{array}$$
Donc
$$ I_a = {\pi \over (1 -\arg\cosh^2 a)\sqrt{a^2-1}} + {\pi \over 2 (a-\cosh 1)}.$$
\fincorrection
\correction{002677}
Sur le c{\^o}t{\'e} $[R, R+i\pi ]$, la fonction $f(z) = e^{az}/\cosh z$ est major{\'e}e en module
par
$${e^{aR}\over |\cosh(R+iy)|} = {e^{aR}\over \cosh^2R -\sin^2 y}
< {e^{aR}\over \cosh^2R -1} \to 0 \hbox{ quand }R\to \infty $$
De m{\^e}me sur l'autre c{\^o}t{\'e}. L'int{\'e}grale sur le c{\^o}t{\'e} $[R+i\pi , -R+i\pi ]$ vaut
$e^{ia\pi }\int_{-R}^R f(x)\,dx$.

La fonction $f$ n'a qu'un p{\^o}le dans le rectangle, qui est $i\pi /2$. Donc {\`a} la limite
$$ (1+e^{ia\pi }) I = 2i\pi  \mbox{\rm Res}(f, i\pi /2) = 2i\pi  {e^{ia\pi }\over i} = 2\pi  e^{ia\pi }$$
et le r{\'e}sultat en divisant.
\fincorrection
\correction{002678}
On int{\`e}gre $f(z) = e^{2iaz -z^2}$ sur le rectangle. Il n'y a aucun p{\^o}le, donc
l'int{\'e}grale sur le rectangle est nulle; l'int{\'e}grale sur le c{\^o}t{\'e} $[R, R+ia]$ tend vers
z{\'e}ro. On a donc
$$\begin{array}{ccc}
 \int_0^{+\infty } e^{2iax-x^2} \,dx &= \int_{[0,ia]} f(z)\,dz
 + \int_{[ia, ia+\infty ]}f(z)\,dz \cr
&=  \int_0^a  e^{-2ay+y^2}\,idy + \int_0^{+\infty } e^{(ia-x)(ia+x)}\,dx \cr
&= i\int_0^a e^{y^2-2ay}\, dy + e^{-a^2}{\sqrt \pi \over 2}
\end{array}$$
d'o{\`u} le r{\'e}sultat en ne conservant que la partie r{\'e}elle.
\fincorrection
\correction{002679}
\begin{itemize}
%\newfunction{Res}
\item Le bord du quart de disque est form{\'e} de trois parties : le segment
$[0,a]$, le quart de cercle, et le segment $[ia, 0]$. L'int{\'e}grale sur le quart de
cercle tend vers z{\'e}ro quand $a\to \infty $ (lemme de Jordan), et en posant $z=iy$, on voit
que
$$ \int_{[ia,0]} z R(z^4)\,dz =  \int_0^a y R(y^4)\,dy $$
donc l'int{\'e}grale sur le contour tend vers $2I$. Donc $I$ est {\'e}gal {\`a} $i\pi $ fois la
somme des r{\'e}sidus de $zR(z^4)$ dans le quart de plan $\{\Re z>0, \Im z >0\}$.

\item Les p{\^o}les sont $e^{i\pi /8}$ et $e^{i3\pi /8}$. 
On trouve $I = \displaystyle{\pi \sqrt 2\over8}$.

\item On int{\`e}gre de m{\^e}me sur le contour form{\'e} du bord de $\{0< \arg z < {2\pi \over
p}, 0<|z|<a\}$. On obtient
$$ \left(1 -e^{i 2\pi  {n+1\over p}}\right) I = 2\pi  i
\sum_{\textstyle
{{\rm p\hat oles\ } z_k\atop 0<\arg z_k<{2\pi \over p}}}
\mbox{\rm Res}\left( z^n R(z^p), z_k\right)
$$

La formule n'est int{\'e}ressante que si le membre de gauche est non nul, c'est-{\`a}-dire
si $n+1$ n'est pas un multiple de $p$.

\item Dans ce cas, il faut donc $p\geq 2$. On est forc{\'e} de prendre une puissance
paire, car on a suppos{\'e} $R$ sans racines r{\'e}elles: ici $R(x) = 1/(1+x^2)$, et $n=1$.
Il y a deux p{\^o}les dans l'angle concern{\'e}, soit $e^{i\pi /2p}$ et $e^{i3\pi /2p}$. Le calcul
est semblable au cas {\it (b)}, et donne
$$ I_{p} = {\pi \over 2p \sin\displaystyle{\pi \over p}}.$$
\end{itemize}
\fincorrection
\correction{002820}
Soit $f(z) = \mathrm{Log} (z)$. Alors
\begin{eqnarray*}
 f'(z) &=& \frac{1}{z} = \frac{1}{z_0 - (z_0-z)} = \frac{1}{z_0}\frac{1}{1 - \frac{z_0-z}{z_0}} \\
&=& \frac{1}{z_0}\sum_{k\geq 0} \frac{1}{z_0^k} (z_0-z)^k \quad \text{ pour } \quad |z-z_0|<|z_0|.
\end{eqnarray*}
Notons que $\{ z\in \C \, ; \; |z-z_0|<|z_0|\} = D (z_0 , |z_0| )$ est optimal car on ne
peut prolonger $\mathrm{Log} (z)$ en $0$. Le d\'eveloppement  est :
$$f(z) = \mathrm{Log}(z_0)- \sum_{k\geq 1} \frac{(-1)^k}{kz_0^k} (z-z_0)^k.$$
En ce qui concerne la deuxi\`eme question, la r\'eponse est NON. D'apr\`es le cours $f$ coincide avec sa s\'erie de Taylor \it si \rm $D(z_0, R) \subset \Omega$. Or, si $\Re (z_0)<0$, ce n'est pas le cas et $D(z_0 , |z_0|)\cap ]-\infty , 0[ \neq \emptyset$. Le $\mathrm{Log}$ et la s\'erie de Taylor ne coincident pas dans $D(z_0,|z_0|)\cap \Omega$ puisque le $\mathrm{Log}$ ne peut \^etre prolong\'e de mani\`ere continu dans aucun point de $]-\infty, 0[$.
Remarquons qu'ici $D(z_0,|z_0|)\cap \Omega$ n'est pas connexe ce qui est cruciale dans l'exercice \ref{ex:burnol1.4}.
\fincorrection
\correction{002821}
On a $\sin (z) = \frac{1}{2i}(e^{iz}-e^{-iz})=0$ si et seulement si $e^{2i\pi z}=1$ ce qui est le cas si et seulement si $z\in \pi \Z$.
Soit $z_0\in U=\C\setminus \pi \Z$. Alors, le rayon de convergence de la s\'eries de Taylor de $f$ est
$$R= \mathrm{dist} (z_0 , \pi \Z ) = \min_{n\in \Z} |z_0 -\pi n| .$$
\fincorrection
\correction{002822}
Si $f(z)g(z) = 0$ pour tout $z\in \C$, alors au moins une des fonctions $f,g$ a un z\'ero non isol\'e.
\fincorrection
\correction{002823}
 Commen\c{c}ons donc par le contre-exemple en prenant $U=\Omega$, c'est \`a dire $U=\C \setminus ]-\infty, 0]$, et $f=\mathrm{Log}$. Il suffit alors de choisir $z_1=-1+i$ et $z_2=-1-i$ et d'appliquer l'exercice \ref{ex:burnol1.1}.
Supposons maintenant $U$ convexe, notons $D_i=D (z_i,R_i)$, $i=1,2$, et supposons que $V=D_1\cap D_2 \neq \emptyset$. Les cons\'equences imm\'ediates de la convexit\'e de $U$ sont:
(a) $U\cap D_1$ et $U\cap D_2$ sont connexes (et m\^eme convexes).
(b) $[z_1,z_2]\subset U$ et donc $V\cap U$ est un ouvert non vide.
Consid\'erons $g_1$. Si $r>0$ est suffisamment petit pour que $D(z_1,r)\subset U$, alors $f=g_1$ dans $D(z_1,r)$. Par le principe du prolongement analytique (ou celui des z\'eros isol\'es) on a $f=g_1$ dans le connexe $U\cap D_1$. C'est donc aussi vrai dans $V\cap U\subset D_1\cap U$.
Le m\^eme raisonnement s'applique \`a $g_2$ et donc $g_1=g_2$ dans $V\cap U$. Encore une fois le principe du prolongement analytique assure donc que $g_1=g_2$ dans $V$ (qui est  connexe).
\fincorrection
\correction{002824}
\begin{enumerate}
  \item 
Soit $\phi (t,z) =\frac{z-1}{1+t(z-1)}$ et notons $D=\{ |z-1|<1\}$. Pour tout $t\in [0,1]$ et tout $z\in D$ on a
$$\phi (t,z) = (z-1) \sum_{k\geq 0} (-1)^k t^k (z-1)^k .$$
Or $|(-1)^k t^k (z-1)^k|\leq |z-1|^k$. Si $0<r<1$, alors la s\'erie pr\'ec\'edente converge normalement dans $D(1,r)$ ce qui permet d'avoir (cf. le polycopié 2005/2006 de J.-F.~Burnol, chapitre 15, théorème 29)
$$\int_0^1 \phi (t,z) dt = \sum_{k\geq 0} (-1)^k (z-1)^{k+1} \int _0^1 t^k dt =\sum_{k\geq0} (-1)^k \frac{(z-1)^{k+1}}{k+1}.$$
Cette s\'erie est la s\'erie de Taylor de $\mathrm{Log} (z)$ en $1$ qui coincide avec $\mathrm{Log} (z)$ dans le disque $D$. Par cons\'equent, $z\mapsto \int _0^1 \phi (t,z)dt$ et $z\mapsto \mathrm{Log}(z)$ coincident dans $D$. On conclut par prolongement analytique et en remarquant que $z\mapsto \int_0^1 \phi (t,z)dt$ est une fonction holomorphe dans $\Omega$ (cf. le polycopié 2005/2006 de J.-F.~Burnol, chapitre 14, théorème 26).
En ce qui concerne le reste $R_N$ voici le calcul:
\begin{eqnarray*}
R_N (z) &=&\int_0^1 \sum_{k\geq N} (-1)^k (z-1)^{k+1} t^k dt\\
&=& (-1)^N \int_0^1 (z-1)^{N+1} t^N \sum_{k\geq 0} (-1)^k (z-1)^k t^k dt \\
&=& (-1)^N (z-1)^{N+1} \int _0^1 \frac{t^N}{1+t(z-1)}dt.
\end{eqnarray*}
  \item Si $\Re(z) \geq \delta$, alors
$$|1+t(z-1)| \geq |\Re (1+t(z-1))| = |1+t\Re (z-1)|\geq \delta .$$
Par cons\'equent,
$$|R_N (z) | \leq |z-1|^{N+1} \int_0^1 \frac{t^N \, dt}{|1+t(z-1)|} \leq \frac{1}{\delta} \frac{|z-1|^{N+1}}{N+1}.$$
D'o\`u la convergence uniforme.
  \item Voir ci-dessus.
  \item On a $z=1+e^{i\phi} = (e^{i\frac{\phi}{2}}+e^{-i\frac{\phi}{2}})e^{i\frac{\phi}{2}}=2\cos \frac{\phi}{2}e^{i\frac{\phi}{2}}$. D'o\`u $\mathrm{Arg} (z) = \frac{\phi}{2}$ et $r=|z|=2|\cos \frac{\phi}{2} | =2\cos \frac{\phi}{2}$.
Comme :
\begin{eqnarray*}
\log(2\cos \frac{\phi}{2}) +i\frac{\phi}{2} &= & \mathrm{Log} (z) = \sum_{k\geq 1} (-)^{k-1} \frac{(z-1)^k}{k} =\\
&=& \sum_{k\ge 1} \frac{(-1)^{k-1}}{k} \cos (k\phi) + i\sum_{k\geq 1} \frac{(-1)^{k-1}}{k}  \sin (k\phi)
\end{eqnarray*}
il suffit d'identifier les parties r\'eelles et imaginaires pour en d\'eduire les \'egalit\'es demand\'ees.
La convergence uniforme r\'esulte de la question 3 puisque
$$\Re (z) = \Re (1+e^{i\phi }) =1+\cos \, \phi \geq 1+ \cos (\pi -\epsilon ) = \delta > 0.$$
\end{enumerate}
\fincorrection
\correction{002825}
\begin{enumerate}
  \item C'est le principe du maximum.
  \item La fonction $g(z) = f(z) f(-z)$ est nulle sur le cercle de rayon $1$.
\end{enumerate}
\fincorrection
\correction{002826}
Soit $z=e^{i\theta}$. Alors $|4z+3|=|e^{i\theta } (4+3e^{-i\theta })| = |\overline{4+3e^{i\theta }}| =|4+3z|$.
Par le principe du maximum
$$|\Phi (z) | < 1=\sup_\theta |\Phi (e^{i\theta })| \quad \text{ pour tout } \;\; z \in D.$$
\fincorrection
\correction{002827}
Supposons qu'il existe $z\in \C$ avec $F(z) \neq 0$ et notons $\alpha = |F(z)|>0$. Si $n>1$ tel que $\frac{1}{n} <\alpha $, alors le principe du maximum affirme que
$$|F(z) | < \frac{1}{n} < \alpha \quad \text{pour tout} \;\; |z| < n.$$
Contradiction.
\fincorrection
\nocorrection
\correction{002829}
Soit $g(z) = e^{f(z)}$. On a $|g(z)| =e^{\Re (f(z))}$. Par cons\'equent $g$ est une fonction enti\`ere born\'ee. Elle est donc constante par d'Alembert-Liouville.
\fincorrection
\correction{002830}
La formule de Cauchy pour $f^{(n+1)}(z)$ est
$$\frac{f^{(n+1)}(z)}{(n+1)!} =\frac{1}{2\pi i} \int_{C_R} \frac{f(\xi )}{(\xi -z)^{n+2}}\, d\xi $$
o\`u $C_R= \{|\xi|=R\}$. Pour les estimations suivantes, prenons $R> \min(2|z|,1)$. Comme $|\xi -z| \geq |\xi| -|z| =R-|z| \geq R/2$,
$$ \frac{|f(\xi )|}{|(\xi -z)^{n+2}|}\leq M \frac{(1+R)^n}{(R/2)^{n+2}} \leq M\frac{(2R)^n}{(R/2)^{n+2}}=2^{2n+2}M\frac{1}{R^2}.$$
Ensemble avec la formule de Cauchy on a donc
$$\frac{|f^{(n+1)}(z)|}{(n+1)!}\leq \frac{1}{2\pi} \int _{C_R} 2^{2n+2}M \frac{1}{R^2} |d\xi| =2^{2n+2} M\frac{1}{R}$$
pour n'importe quel $R>2|z|$. On vient de montrer que $f^{n+1}(z) =0$ pour tout $z\in \C$.
Utilisons maintenant $g(z) = \frac{f(z) -P(z)}{z^{n+1}}$ o\`u $P$ est le polyn\^ome de Taylor de $f$ \`a l'origine \`a l'ordre $n$. On remarque d'abord que l'origine est z\'ero d'ordre $n+1$ de $f(z) -P(z)$ ce qui explique que $g$ se prolonge holomorphiquement \`a l'origine. C'est donc une fonction enti\`ere pour laquelle on a
$$|g(z) | \leq \frac{C|z|^n}{|z|^{n+1}}=C\frac{1}{|z|}$$
pour un certain $C>0$ et pour $z$ de module suffisamment grand. De nouveau, $g$ est une fonction enti\`ere born\'ee, elle est donc constante (notre estimation donne m\^eme $g\equiv 0$ et donc $f=P$).
\fincorrection
\correction{002831}
Soit $z=Re^{i\theta }$. Alors $|e^z| =e^{R\Re (e^{i\theta })} =e^{R\cos(\theta
  )}$. Pour $f(z)= z+e^z$, on a,  pour $\theta $ tel que $\cos(\theta )\leq
0$, $|f(z)|\geq R - e^{R\cos(\theta)}\geq R-1$. Si par contre $\cos(\theta )>
0$ alors $|f(z)|\geq e^{R\cos\theta} - R$. Dans les deux cas
$$|f(Re^{i\theta})| \to\infty \quad \text{pour}\;\; R \to \infty .$$
Nos calculs n'impliquent pas  $\lim_{|z|\to \infty } |f(z)|
=\infty$! Comme en fait $e^z$ n'est PAS un polynôme, on peut même affirmer que $\lim_{|z|\to \infty } |f(z)|
=\infty$ est FAUX.  
\fincorrection
\correction{002832}
\begin{enumerate}
  \item Le r\'esidu est $a_{-1} =1$.
  \item $f(z) =\frac{1}{1+z^2}=\sum_{k=0}^\infty (-1)^k z^{2k}$  pour $ |z|< 1$.
C'est une s\'erie enti\`ere car $f$ est holomorphe dans le disque unit\'e et $\mathrm{Res} (f ,0)=0$.
  \item $f(z) = \frac{1}{z}\frac{1}{1+z^2} =\sum_{k=0}^\infty (-1)^k z^{2k-1} =\frac{1}{z}-z +z^3 ....$.et $\mathrm{Res} (f ,0)=0$ 
\end{enumerate}
\fincorrection
\correction{002833}
Comme $e^w =\sum_{n\geq 0} \frac{w^n}{n!}$,
$$\exp (\frac{1}{z}) =\sum_{n\geq 0} \frac{1}{n!}\frac{1}{z^n} = \sum_{n=-\infty}^0 \frac{1}{(-n)!} z^n.$$
Par cons\'equent, $\mathrm{Res}(f,0) = 1$. Sinon, si $z_0\neq 0$, alors $\mathrm{Res}(f,z_0) = 0$ par holomorphie de $f$ dans  $\C\setminus \{0\}$.
\fincorrection
\correction{002834}
\begin{enumerate}
  \item Comme $\sin(z) =z -\frac{z^3}{3!}+... =z(1+o(z))$,
$$f(-z) =\frac{1}{\sin(z)} =\frac{1}{z}(1+o(z))^{-1} =\frac{1}{z} (1+o(z)) = \frac{1}{z} +o(1) .$$
Par cons\'equent $f$ poss\`ede un p\^ole simple \`a l'origine (ce qui est \'evident puisque $\sin(z)$ poss\`ede un z\'ero simple \`a l'origine) et $\mathrm{Res} (f,0) =1$. On l'obtient aussi par la formule de l'exercice \ref{ex:burnol6.1} et le fait que l'origine est un p\^ole simple:
$$\mathrm{Res} (f,0) =\lim_{z\to 0} z f(z) = \lim_{z\to 0} \frac{z}{\sin(z)} =1.$$
La partie singuli\`ere de la s\'erie de Laurent est $\frac{1}{z}$ et le terme constant est $0$.
  \item On a $\sin(z) -\sh(z) =(z-\frac{z^3}{3!} +\frac{z^5}{5!}+O(z^7)) -(z+\frac{z^3}{3!} +\frac{z^5}{5!}+O(z^7)) =-\frac{z^3}{3} + O(z^7)$. D'o\`u
$$f(z) =\frac{1}{\sin(z) -\sh(z)} = -\frac{3}{z^3} (1+O(z^4)) =-\frac{3}{z^3} + O(z).$$
  \item  On obtient de mani\`ere analogue que $$f(z) =\frac{1}{z \sin(z) \sh(z)} =\frac{1}{z^3} +O(z).$$
\end{enumerate}
\fincorrection
\correction{002835}
Observons tout d'abord que :
$$ f(z) =\frac{1}{(z-1)(z-2)}= \frac{1}{z-2} +\frac{1}{1-z}.$$
On a
$$\begin{aligned} \frac{1}{1-z}=\sum_{n\geq 0}z^n \quad &\text{pour} \;\;  |z|<1 \\
\frac{1}{1-z}=-\frac{1}{z}\frac{1}{1-\frac{1}{z}} =-\frac{1}{z}\sum_{n\geq 0}z^{-n} = -\sum_{n=-\infty}^{-1} z^n &\quad \text{pour} \;\; |z|> 1.\end{aligned}$$
De la m\^eme mani\`ere
$$\frac{1}{z-2}=-\sum_{n\geq 0}\frac{1}{2^{n+1}}z^n \quad \text{si} \;\; |z|< 2 \;\; \text{et} \;\;
\frac{1}{z-2}=\sum_{n=-\infty} ^{-1}\frac{1}{2^{n+1}}z^n \quad \text{si} \;\; |z|>2.$$
On en d\'eduit les expressions des s\'eries de Laurent en $0$ dans les trois couronnes centr\'ees \`a l'origine.
En $z=1$ et $z=2$, $f$ a des p\^oles simples. D'o\`u
$$\mathrm{Res} (f,1) =\lim_{z\to 1} (z-1) f(z) = -1 \quad  \text{et} \quad \mathrm{Res} (f,2) =\lim_{z\to 2} (z-2) f(z) = 1.$$
D\'eterminons encore la s\'erie de Laurent  $f$ en $1$:
$$f(z) = \frac{1}{z-1}\frac{1}{(z-1)-1} =\frac{1}{z-1} (-1) \sum _{n\geq 0} (z-1)^n = -\sum_{n=-1}^\infty (z-1)^n$$
pour $|z-1|<1$.
\fincorrection
\correction{002836}
Soit $\gamma:I\to \C$ un lacet quelconque. Posons
$$H(t,u)=u\gamma (t) \quad \text{pour} \quad t\in I \quad \text{et} \quad u\in [0,1].$$
 C'est clairement une homotopie de lacets
(voir la d\'efinition du cours!) telle que $H(t,1)=\gamma (t)$  et $H(t,0)=0$ pour tout $t\in I$.
\fincorrection
\correction{002837}
Soit $\gamma$ un lacet dans $\C$. Alors $\sup_{t\in I} |\gamma(t)| =R/2< \infty$. Si $|z|> R$, la fonction $\xi \mapsto \frac{1}{\xi -z}$ est holomorphe dans le disque $D(0,R)$. Par cons\'equent,
$$\mathrm{Ind}(\gamma , z) =\frac{1}{2i\pi} \int_\gamma \frac{d\xi}{\xi -z} =0.$$
Remarquons que ceci implique que $\mathrm{Ind}(\gamma , z)=0$ pour tout $z$ dans la composante connexe non born\'ee de $\C\setminus \gamma$.
\fincorrection
\correction{002838}
\begin{enumerate}
  \item Le calcul de l'indice de $c_N$ est \'evident. L'affirmation sur l'existence de $g$ continue telle que $\gamma =e^g$ et $g(1)-g(0)=2\pi i N$ est le contenu du polycopi\'e de J.-F.~Burnol 2005/2006, chapitre 30. Le reste est laiss\'e au lecteur.
\end{enumerate}
\fincorrection
\nocorrection
\correction{002840}
Si $f$ a un p\^ole d'ordre $N$ en $z_0$, on a, avec $a_{-N}\neq 0$,
$$f(z) = \sum_{n= -N}^\infty a_n (z-z_0)^n = \frac{a_{-N}}{(z-z_0)^N} +...+\frac{a_{-1}}{z-z_0}+a_0+...$$
D'o\`u
$$(z-z_0)^N f(z) = a_{-N} +...+a_{-1} (z-z_0)^{N-1} +a_0 (z-z_0)^N+....$$
ce qui donne
$$\left(\frac{d}{dz}\right)^{N-1}\big[ (z-z_0)^N f(z) \big]= (N-1)! \,a_{-1} +o(z-z_0) .$$
La formule en r\'esulte en faisant tendre $z$ vers $z_0$.
Le cas $N=1$ est important pour la pratique. Notons que, si $f$ a un p\^ole simple en $z_0$, cette fonction s'\'ecrit $f=h/g$
au voisinage de $z_0$ o\`u  $h,g$ sont des fonctions holomorphes au voisinage
de $z_0$ telles que $h$ ne s'annule pas en $z_0$ et $g$ a un z\'ero simple en
$z_0$, i.e. $g(z_0)=0$ et $g'(z_0)\neq 0$. Comme
$$\mathrm{Res} (f, z_0) =\lim_{z\to z_0} (z-z_0) f(z) = h(z_0)\lim_{z\to z_0} \frac{z-z_0}{g(z)-g(z_0)} $$
on a aussi la formule utile
\begin{equation} \label{1}
\mathrm{Res} (f, z_0) = \frac{h(z_0)}{g'(z_0)}.
\end{equation}
\fincorrection
\nocorrection
\correction{002842}
La fonction $f(z) = 1/(z-a)(z-b)(z-c)$ est holomorphe dans le disque $D(0,a)$. Par cons\'equent
l'int\'egrale est nulle si $r<a$.
Par le th\'eor\`eme des r\'esidus,
\begin{eqnarray*}
\frac{1}{2i\pi}\int_C \frac{1}{(z-a)(z-b)(z-c)} dz &= &\mathrm{Res} (f,a) \quad , \quad \mathrm{Res} (f,a)+\mathrm{Res} (f,b)
\quad ou \\ && \mathrm{Res} (f,a)+\mathrm{Res} (f,b)+\mathrm{Res} (f,c)
\end{eqnarray*}
si $a<r<b$, $b<r<c$ ou $c<r$. Le calcul de ces r\'esidus se fait par la formule de l'exercice \ref{ex:burnol6.1} puisque tous
les p\^oles sont simples:
\begin{eqnarray*}\mathrm{Res} (f,a) =\frac{1}{(a-b)(a-c)} \quad , \quad \mathrm{Res} (f,b)=\frac{1}{(b-a)(b-c)} \\
et \quad \mathrm{Res} (f,c)=\frac{1}{(c-a)(c-b)}.
\end{eqnarray*}
On en d\'eduit fa\c{c}ilement la valeur de l'int\'egrale dans les trois cas.
En ce qui concerne le calcul de cette int\'egrale via la d\'ecomposition en \'el\'ements simples remarquons juste que
$$\int _C \frac{1}{z-d} \, dz$$
vaut $2i\pi$ si $d$ est \`a l'int\'erieur de $C$ et $0$ si $d$ est \`a l'ext\'erieur.
\fincorrection
\nocorrection
\correction{002844}
On a
\begin{eqnarray*}
\int_{\partial\Omega} \frac{f(z)\,dz}{(z-z_1)(z-z_2)} &=& \int_\gamma \frac{f(z)\,dz}{(z-z_1)(z-z_2)} -\sum_{j=1}^N \int_{\gamma_j} \frac{f(z)\,dz}{(z-z_1)(z-z_2)} \\
&=& 2i\pi \Big( \mathrm{Res} (f, z_1)+\mathrm{Res} (f, z_2) \Big) = 2i\pi \Big( \frac{f(z_1)}{z_1-z_2}+\frac{f(z_2)}{z_2-z_1}\Big).
\end{eqnarray*}
D'o\`u
$$\lim_{z_2\to z_1} \int_{\partial\Omega} \frac{f(z)\,dz}{(z-z_1)(z-z_2)}= 2i\pi f'(z_1).$$
\fincorrection
\correction{002845}
Analogue à l'exercice \ref{ex:burnol6.4}
\fincorrection
\nocorrection
\correction{002847}
Comme $\tan(\pi z) =\frac{\sin(\pi z)}{\cos(\pi z)}$ cette fonction est une fonction m\'eromorphe de $\C$ ayant que des p\^oles simples en
$1/2 \; \mathrm{mod} \; 1$. En effet, $\cos(w)=0$ si et seulement si $w=\pi/2 \; \mathrm{mod} \;\pi$ et $\cos'(\pi/2+k\pi) \neq 0$. Notons
$z_k =1/2+k$, $k\in \Z$. La formule (\ref{1}) s'applique et donne
$$ \mathrm{Res}(\tan (\pi z) , z_k) =\frac{\sin(\pi z_k)}{-\pi \sin (\pi z_k)} = -\frac{1}{\pi}.$$
Par cons\'equent :
$$\int _{|z|=N} \tan(\pi z) \, dz = 2i\pi 2N \Big(-\frac{1}{\pi} \Big) = -4iN.$$
\fincorrection
\correction{002848}
Si $A=R\cos (\Phi) $ et $B=R\sin (\Phi)$ on a
$$ \int _0^{2\pi} \frac{d\theta}{A+B\sin(\theta )+C\cos(\theta )} = \int _0^{2\pi} \frac{d\theta}{A+R\sin(\theta +\Phi)} =\int _0^{2\pi} \frac{d\alpha}{A+R\sin(\alpha )}.$$
Pour trouver la valeur de cette derni\`ere int\'egrale posons $z=e^{i\alpha }$. Alors $d\alpha = -i \frac{dz}{z}$ et
$$ I=\int _0^{2\pi} \frac{d\alpha}{A+R\sin(\alpha )} = \int _{|z|=1} \frac{-i\, dz}{z( A+R\frac{z-\overline{z}}{2i})}=
\int _{|z|=1} \frac{2 dz}{Rz^2 +2iAz -R}.$$
Le d\'enominateur de cette derni\`ere expression s'annule en
$$z^\pm =\frac{-2iA\pm \sqrt{-4A^2+4R^2}}{2R}= -i\left( \frac{A}{R} \mp \sqrt{\left(\frac{A}{R} \right)^2 -1}    \right)$$
Un calcul \'el\'ementaire montre que seulement la racine $z^+= -i\left( \frac{A}{R} - \sqrt{\left(\frac{A}{R} \right)^2 -1} \right)$ est dans le disque unit\'e ouvert pourvu que $A>0$ (le cas $A<0$ est similaire). Il s'ensuit par le th\'eor\`eme du r\'esidu que
$$\frac{1}{2\pi}I= i \frac{2}{R}\frac{1}{z^+-z^-}= \frac{1}{R\sqrt{\left(\frac{A}{R} \right)^2 -1}}=\frac{1}{\sqrt{A^2-R^2}}.$$
\fincorrection
\nocorrection
\nocorrection
\correction{002853}
Si $z=e^{i\theta}$, alors $\sin \theta =\frac{z-\overline{z}}{2i}$ et $dz = ie^{i\theta} d\theta$.
D'o\`u
$$\frac{1}{2\pi} \int_0^{2\pi} \frac{\sin \theta}{a+\sin \theta} d\theta = \frac{1}{2\pi} \int_{ |z|=1}
\frac{z-\overline{z}}{2ia +z-\overline{z}}\frac{dz}{iz}=
\frac{1}{2i\pi} \int_{ |z|=1} \frac{z^2-1}{z^2 +2iaz -1}\frac{dz}{z}.$$
Il suffit alors d'utiliser le th\'eor\`eme des r\'esidus.
\fincorrection
\nocorrection
\nocorrection
\nocorrection
\correction{002857}
Rappelons la formule de Cauchy pour $f$ holomorphe sur $\overline{\Omega}$
(donc sans singularit\'es):
$$f(z)=\frac{1}{2i\pi}\int _{\partial \Omega} \frac{f(w)}{w-z}dw.$$
Il s'agit ici d'obtenir une version g\'en\'eralis\'ee pour des fonctions $f$ ayant des singularit\'es
$z_1,...,z_N\in \Omega$. Fixons $z\in \Omega \setminus \{z_1,...,z_N\}$ et consid\'erons $G(w)=\frac{f(w)}{w-z}$.
Cette fonction a un p\^ole simple en $w=z$ et :
$$\mathrm{Res} (G,z)=\lim_{w\to z} (w-z)G(w)=f(z).$$
Les autres singularit\'es de $G$ dans $\Omega$ sont $z_1,...,z_N$.
Par d\'efinition, le r\'esidu de $G$ en $z_j$ est le \og coefficient
$a_{-1}$\fg{} de la s\'erie de Laurent de $G$ en $z_j$. Or
$$G(w)= \frac{1}{w-z} f(w) = \sum _{k\geq 0} b_k (w-z_j)^k \sum _{l\in \Z} c_l (w-z_j)^l$$
puisque $w\mapsto \frac{1}{w-z}$ est holomorphe au voisinage de $z_j$ (et bien s\^ur on peut calculer les $b_k$,
mais ce n'est pas utile). On remarque que pour calculer $a_{-1}$ interviennent seulement les indices $(k,l)$
qui v\'erifient $k+l=-1$. Comme $k\geq 0$ on a $l=-1-k<0$. D'o\`u :
$$\mathrm{Res} (G,z_j)= \mathrm{Res} \left(\frac{g_j(w)}{w-z},z_j\right).$$
On peut maintenant utiliser l'exercice \ref{exo:residuinfini} ou alors conclure directement: si $R_0=2\max \{|z| , |z_j|\}$,
alors pour tout $R>R_0$,
$$\frac{1}{2i\pi}\int_{|w|=R}\frac{g_j(w)}{w-z}dw = \sum_{\xi \in \{z,z_j \}} \mathrm{Res} \left(\frac{g_j(w)}{w-z},\xi \right).$$
Par cons\'equent, l'int\'egrale $\int_{|w|=R}\frac{g_j(w)}{w-z}dw$ ne d\'epend pas de $R> R_0$. Or, il existe $C>0$ tel que
$$\left|\frac{g_j(w)}{w-z} \right| \leq \frac{C}{|w|^2} \quad , \quad |w| > R_0 \,,$$
ce qui entra\^ine
$$\left|\int_{|w|=R_0}\frac{g_j(w)}{w-z}dw \right|\leq \lim_{|w|=R} \int_{|w|=R_0}\frac{C}{|w|^2} |dw| =0.$$
D'o\`u
$$0= \frac{1}{2i\pi} \int_{|w|=R_0}\frac{g_j(w)}{w-z}dw = \mathrm{Res} \left(\frac{g_j(w)}{w-z},z_j\right)+g_j(z).$$
Il suffit alors d'appliquer le th\'eor\`eme des r\'esidus \`a $G$ pour conclure:
$$\frac{1}{2i\pi} \int_{\partial \Omega } G(w) \, dw = f(z) -g_1(z) -...-g_N(z).$$
\fincorrection
\correction{002858}
Dans le calcul des int\'egrales on est souvent confront\'e \` a des passages \` a la limite
(du genre $\lim_{R\to \infty } \int_{C_R} f(z) \, dz$
ou $\lim_{r\to 0 } \int_{C_r} f(z) \, dz$  dans le cas ou $0$ est une singularit\'e que l'on contourne,
$C_R$ un morceaux de cercle comme dans les exercices ici). Cet exercice et le suivant donnent des outils
tr\`es pratiques pour ce genre de calculs.\\
Si $f$ a $z_0$ comme p\^ole simple, sa s\'erie de Laurent en $z_0$ est de la forme
$$ f(z) = a_{-1}(z-z_0)^{-1} +a_0 +a_1(z-z_0)+.... =\sum_{k\geq -1} a_k (z-z_0)^k .$$
Par convergence normale de cette s\'erie
\begin{eqnarray*}
\int_{C_r(\alpha , \beta )} f(z) \, dz &=& \sum_{k\geq -1} a_k \int_{C_r(\alpha , \beta )} (z-z_0)^k \, dz
= \sum_{k\geq -1} a_k \int_\alpha ^\beta (re^{i\theta })^k ire^{i\theta } \, d\theta \\
&=& ia_{-1}(\beta -\alpha ) + r\sum_{k\geq 0}\left( r^k a_{k}\frac{e^{i(k+1)\beta } -e^{i(k+1)\alpha } }{(k+1)} \right).
\end{eqnarray*}
On en d\'eduit l'enonc\'e de l'exercice en observant que
 $\left| \frac{e^{i(k+1)\beta } -e^{i(k+1)\alpha } }{(k+1)} \right| \leq \frac{2}{k+1}$ et en faisant tendre $r\to 0$.
\fincorrection
\correction{002859}
Pour $\epsilon >0$ il existe $R> 0$ tel que $|f(z)|\leq \epsilon $ pour tout $|z|\geq R$, $\Im z \geq 0$. Si $C_R$ est le
demi-cercle sup\'erieur orient\'e alors
\begin{eqnarray*}
& &|\int _{C_R} f(z) e^{i z} \, dz | =|\int _0^\pi f(Re^{i\theta} ) e^{i(R e^{i\theta})} iR e^{i\theta } \, d\theta |
\leq \epsilon \int _0^\pi e^{-R \sin \theta } R \, d\theta\\
&& \qquad = 2\epsilon \int _0^{\frac{\pi}{2}} Re^{-R \sin \theta } \, d\theta
\leq 2\epsilon \int _0^{\frac{\pi}{2}} Re^{-R \frac{\theta}{2} }\, d\theta
= 4\epsilon (1-e^{-R \frac{\pi}{4}}) \leq 8 \epsilon .
\end{eqnarray*}
\fincorrection
\correction{002860}
Utiliser les exercices \ref{ex:burnol2.1} et \ref{ex:burnol2.2}.
\fincorrection
\correction{002861}
Par holomorphie de $z\mapsto e^{-z^2}$,
$$\int_0^R e^{-x^2} dx + \int_{C_R} e^{-z^2} dz + \int _{Re^{i\pi /4}}^0 e^{-z^2} dz =0.$$
Notons $I_{1,R}$ la premi\`ere int\'egrale ci-dessus, $I_{2,R}$ la deuxi\`eme et $I_{3,R}$ la troisi\`eme.
Alors,
$$\lim_{R\to \infty} I_{1,R} = \frac{1}{2} \int _\R e^{-x^2} dx = \frac{1}{2} \int _\R e^{-\pi u^2} \sqrt{\pi} \,du
=\frac{\sqrt{\pi}}{2} .$$
Comme $e^{-z^2}= e^{-\left( e^{i\pi /4}t\right)^2}= e^{-it^2}= \cos (t^2) -i\sin (t^2)$ pour $z=e^{i\pi /4}t$ on a
$$\begin{aligned}I_{3,R} &=-\int_0^R \left( \cos (t^2) -i\sin (t^2)\right)e^{i\pi /4} dt\\
&=-\frac{\sqrt{2}}{2}\left( \int_0^R \cos (t^2) +\sin (t^2) dt +i \int_0^R \cos (t^2) -\sin (t^2) dt \right).
\end{aligned}$$
Il suffit alors de d\'eterminer $\lim_{R\to\infty} I_{2,R}$ pour en d\'eduire les int\'egrales de Fresnel.
Si on pose $z=Re^{i\theta}$, alors
$$\int _{C_R} e^{-z^2} dz = \int_0^\frac{\pi}{4} e^{-R^2 e^{2i\theta}} iR e^{i\theta} d\theta$$
ce qui implique
$$|I_{2,R}| \leq R \int _0^\frac{\pi}{4}e^{-R^2 \cos (2\theta)} d\theta = \frac{R}{2}\int _0^\frac{\pi}{2}e^{-R^2 \cos (\alpha)} d\alpha.$$
 Du changement de variables $\beta = \frac{\pi}{2} -\alpha$ et du fait que $\sin \beta \geq \frac{\beta}{2}$ pour $\beta \in [0, \frac{\pi}{2} ]$ on d\'eduit que :
$$|I_{2,R}| \leq -\frac{R}{2}\int ^0_\frac{\pi}{2}e^{-R^2 \sin (\beta )}d\beta \leq \frac{R}{2}\int _0^\frac{\pi}{2}e^{-R^2 \frac{\beta}{2}}d\beta = \left. -\frac{R}{2}e^{-R^2 \frac{\beta}{2}}\right|_0^\frac{\pi}{2} \leq \frac{2}{R} \to 0$$
lorsque $R\to \infty$.
Conclusion $\int_0^\infty \cos(x^2)dx =  \int_0^\infty \sin(x^2)dx = \frac{\sqrt{2\pi}}{4}$.
\fincorrection
\nocorrection
\nocorrection
\nocorrection
\nocorrection
\nocorrection
\nocorrection
\nocorrection
\nocorrection
\correction{002879}
\begin{enumerate}
  \item Soit $C_R=\{R e^{i\theta } \,;\; 0\leq \theta \leq \frac{2\pi}{n}\}$.
La fonction $$f(z)=\frac{1}{1+z^n} $$ a un seul p\^ole $z_0=e^{i\pi/n}$ dans le secteur. C'est un p\^ole simple et le
r\'esidu est
$$\mathrm{Res} \left(f, e^{i\frac{\pi}{n}}\right)=\frac{1}{nz_0^{n-1}}=-\frac{z_0}{n}= -\frac{1}{n} e^{i\frac{\pi}{n}}.$$
D'o\`u
$$-\frac{2i\pi}{n} e^{i\pi /n}= \int _0^R \frac{dx}{1+x^n}+\int_{C_R}\frac{dz}{1+z^n}+\int_{Re^{i\pi /n}}^0\frac{dz}{1+z^n}$$
pour tout $R>1$. Puisque $n>1$, on a :
$$\lim_{R\to\infty} \left| \int_{C_R}\frac{dz}{1+z^n} \right|\leq \lim_{R\to\infty} \left( \frac{1}{R^n-1}\int_{C_R} |dz|\right) =0.$$
D'autre part,
$$\int_{Re^{2i\pi /n}}^0\frac{dz}{1+z^n}= -\int_0^R \frac{1}{1+x^n}e^{2i\pi /n}dx =-e^{2i\pi /n}\int _0^R \frac{dx}{1+x^n}.$$
Il en r\'esulte que
$$\int _0^R \frac{dx}{1+x^n}=\frac{-\frac{2i\pi}{n} e^{i\pi /n}-\int_{C_R}\frac{dz}{1+z^n}}{1-e^{2i\pi /n} }
\longrightarrow \frac{-\frac{2i\pi}{n} e^{i\pi /n}}{1-e^{2i\pi /n} } =\frac{2i\pi}{n}\frac{1}{2i\sin \left(\frac{\pi}{n}\right)}
= \frac{\frac{\pi}{n}}{\sin \left(\frac{\pi}{n}\right)}$$
lorsque $R\to\infty$.
  \item La fonction $z^a=\exp \left( a(\log r +i\alpha )\right)$ n'est  pas d\'efinie au voisinage de l'origine.
C'est la raison pourquoi on est amen\'e de consid\'erer le petit morceau de cercle $\gamma_\epsilon =\{\epsilon e^{i\theta } \,;\; \frac{2\pi}{a} \geq \theta \geq 0\}$. On va de nouveau noter $C_R=\{R e^{i\theta } \,;\; 0\leq \theta \leq \frac{2\pi}{n}\}$
et $$\Omega = \left\{z=re^{i\alpha} \,;\; 0<r<\infty \;\; , \; 0\leq \alpha \leq \frac{2\pi}{a}\right\}.$$
Pour $z=re^{i\alpha }\in \Omega$ on a
$$\begin{aligned}
&z^a =-1   \\
\Longleftrightarrow \quad & a(\log r +i \alpha ) = i\pi \; (\mathrm{mod} 2i\pi \, ) \\
\Longleftrightarrow\quad & r=1 \quad \text{et} \quad \alpha =\frac{\pi}{a}\, .
\end{aligned}$$
Par cons\'equent, $f(z)=\frac{1}{1+z^a}$ a une seule singularit\'e $z_0= e^{i\frac{\pi}{a}}$ dans $\Omega$.
Comme $f(z)=\frac{1}{h(z)}$ avec $h(z_0)=1+z_0^a =0$ et
$$h'(z_0)=\left(  \exp (a\log z)\right)'_{|z=z_0}=\frac{a}{z_0}z_0^a \neq 0$$
le point $z_0$ est un p\^ole simple et on a
$$\mathrm{Res} (f,z_0) =\frac{1}{h'(z_0)} =-\frac{z_0}{a} = -\frac{1}{a} e^{i\frac{\pi}{a}}.$$
Il suffit alors de proc\'eder comme dans la question 1. pour \'etablir
$$\int_0^\infty \frac{dx}{1+x^a}=\frac{\frac{\pi}{a}}{\sin\left(\frac{\pi}{a}\right)} \quad \text{pour} \quad a>1\, .$$
  \item Soit $x\in(0,\infty )$ et $a=u+iv$ avec $u>1$. Alors
$$|x^a|=|x^{iv}||x^u|=|\exp \left( i(v\log x)\right)|x^u =x^u.$$
Par cons\'equent on a, pour tout $x>1$,
$$\left| \frac{1}{1+x^a}\right|\leq  \frac{1}{x^u-1}$$
ce qui implique la convergence de l'int\'egrale $J(a)$.
Montrons que l'application $a\mapsto J(a)$ est holomorphe dans $\Omega=\{ \Re a >1\}$.
Pour ce faire on utilise des critères d'holomorphie des intégrales avec paramètres (voir le chapitre 14 du polycopi\'e 2005/2006 de J.-F.~Burnol).
Consid\'erons d'abord
$J_1(a)=\int_0^2 \frac{dx}{1+x^a}$. On a
(1) $(a,x)\mapsto g(a,x)=\frac{1}{1+x^a}$ est continue.
(2) $\forall x\in [0,2]$: $a\mapsto g(a,x)$ est holomorphe dans $\Omega$.
\noindent Par un critère d'holomorphie des intégrales avec paramètres  (th\'eor\`eme 26 du chapitre 14 du polycopi\'e 2005/2006 de J.-F.~Burnol)\, $a\mapsto J_1(a)$ est holomorphe dans $\Omega$.
Pour $J_2(a)=\int_2^\infty \frac{dx}{1+x^a}$ il faut en plus de (1) et (2) majorer
$g(a,x)=\frac{1}{1+x^a}$ par une fonction int\'egrable $k$ (d\'ependant que de la variable $x$).
Pour ce faire il faut travailler dans un domaine plus petit
$$\Omega_T =\left\{ \Re a>T\right\} \subset \Omega \quad , \quad T>1\, .$$
Dans ce cas
$$|g(a,x)| = \left| \frac{1}{1+x^a} \right| \leq \frac{1}{x^T-1} \quad \forall x\geq 2 \;\; \text{et} \;\; a\in \Omega _T.$$
Comme $T>1$, $k(x)=\frac{1}{x^T-1}$ est int\'egrable: $\int _2^\infty k(x)\, dx <\infty$.
Par un critère d'holomorphie des intégrales avec paramètres (ici le th\'eor\`eme 27 du chapitre 14 du polycopi\'e 2005/2006 de J.-F.~Burnol), $a\in \Omega_T \mapsto J_2(a)$ est holomorphe.
Ceci \'etant vrai pour tout $T>1$, $J_2$ est holomorphe dans $\Omega$. En conclusion,
$$a \mapsto J(a)=J_1(a)+J_2(a)$$
est holomorphe sur $\Omega$.
L'affirmation $J(a)= \frac{\frac{\pi}{a}}{\sin\left(\frac{\pi}{a}\right)}$, $a\in \Omega$, est une cons\'equence
du principe des z\'eros isol\'es et du fait que nous avons d\'eja \'etabli cette relation pour tout r\'eel $a>1$.
  \item \'Evident.
  \item On peut proc\'eder comme dans la question 3. Notons que
$$|h(p,t)| =\left| \frac{e^{pt}}{1+e^t}\right|=\frac{e^{\Re (p) t}}{1+e^t}.$$
Par cons\'equent, $|h(p,t) | \sim e^{(\Re (p) -1)t}$ pour $t\to\infty$ et
$|h(p,t) | \sim e^{\Re (p) t}$ pour $t\to -\infty$. L'int\'egrale $K(p)$
est donc convergente. Pour \'etablir l'holomorphie de cette fonction il faut travailler
u©ì$$U_\epsilon =\left\{ 0< \Re (p) < 1-\epsilon \right\} \quad\text{ avec } \quad \epsilon >0 \; \; \text{petit}\, .$$
  \item Nous avons vu dans la question pr\'ec\'edente que la fonction $h(p,t)=\frac{e^{pt}}{1+e^t}$
d\'ecroit exponentiellement pour $0< \Re (p) < 1$ lorsque $t\to \pm \infty$. On en d\'eduit
``facilement'' (faire les d\'etails!) que
$$\lim_{R\to\infty }\int_R^{R+2i\pi}\frac{e^{pz}}{1+e^z}dz = \lim_{R\to\infty }\int_{-R+2i\pi}^{-R}\frac{e^{pz}}{1+e^z}dz =0$$
Par le th\'eor\`eme des r\'esidus il en r\'esulte que :
$$2i\pi \mathrm{Res} \left(\frac{e^{pz}}{1+e^z} , i\pi \right)= \lim_{R\to\infty}
\left[ \int_{-R}^{R}\frac{e^{pt}}{1+e^t}dt+ \int_{R+2i\pi}^{-R+2i\pi}\frac{e^{pz}}{1+e^z}dz\right].$$
Or $\int_{R+2i\pi}^{-R+2i\pi}\frac{e^{pz}}{1+e^z}dz= -e^{2i\pi p}\int_{-R}^{R}\frac{e^{pt}}{1+e^t}dt$.
D'o\`u :
$$2i\pi \left( -e^{i\pi p }\right) = 2i\pi \mathrm{Res} \left(\frac{e^{pz}}{1+e^z} , i\pi \right)
=\left( 1-e^{2i\pi p}\right) K(p).$$
Finalement on a
$$K(p)=\pi \frac{2i}{e^{i\pi p}-e^{-i\pi p}} =\frac{\pi}{\sin (\pi p)}.$$
\end{enumerate}
\fincorrection
\nocorrection
\nocorrection
\nocorrection
\nocorrection
\nocorrection
\nocorrection
\nocorrection
\nocorrection
\nocorrection
\nocorrection
\nocorrection
\nocorrection
\nocorrection
\nocorrection
\nocorrection
\nocorrection
\nocorrection
\nocorrection
\nocorrection
\nocorrection
\nocorrection
\nocorrection
\nocorrection
\nocorrection
\nocorrection
\nocorrection
\nocorrection
\nocorrection
\nocorrection
\nocorrection
\nocorrection
\nocorrection
\nocorrection
\nocorrection
\nocorrection
\nocorrection
\nocorrection
\nocorrection
\nocorrection
\nocorrection
\nocorrection
\nocorrection
\nocorrection
\nocorrection
\nocorrection
\nocorrection
\nocorrection
\nocorrection
\nocorrection
\nocorrection
\nocorrection
\nocorrection
\nocorrection
\nocorrection
\nocorrection
\nocorrection
\nocorrection
\nocorrection
\nocorrection
\nocorrection
\nocorrection
\nocorrection
\nocorrection
\nocorrection
\nocorrection
\nocorrection
\nocorrection
\nocorrection
\nocorrection
\nocorrection
\nocorrection
\nocorrection
\nocorrection
\nocorrection
\nocorrection
\nocorrection
\nocorrection
\nocorrection
\nocorrection
\nocorrection
\nocorrection
\nocorrection
\nocorrection
\correction{002883}
 $|Q(z)|=|z^{50}(z^{61}+3)|=|z^{61}+3|\geq 2$ pour $|z|=1$. D'o\`u
 $$|P(z)-Q(z)| =1 < |Q(z)| \quad \text{dans} \quad  \{|z|=1\}.$$
 Par le th\'eor\`eme de Rouch\'e, $P,Q$ ont le m\^eme nombre de z\'eros dans $D(0,1)$.
 Le reste en d\'ecoule en observant que $P'=Q'$ et $P(0)\neq 0$.
\fincorrection
\correction{002884}
L'application $\Phi (z)=\frac{3z+5}{z+2}$ est une homographie. L'image d'un cercle est alors de nouveau
un cercle ou une droite. De plus on remarque que
(1) $\Phi (x)\in \R$ pour tout r\'eel $x\neq -2$.
(2) $\Phi (\overline{z}) = \overline {\Phi (z)}$ pour tout $z\in \C\setminus \{-2\}$.
\noindent Comme $\Phi (-1) =2$ et $\Phi (1) =\frac{8}{3}$, l'image du cercle unit\'e est un cercle sym\'etrique
par rapport \`a l'axe r\'eel (cf. (2)) avec centre $\big( \frac{8}{3} +2\big) \frac{1}{2}=\frac{7}{3}$
et de rayon $\frac{8}{3}-\frac{7}{3}=\frac{1}{3}$.
Le cercle de rayon $2$ centr\'e \`a l'origine contient $-2$. C'est l'unique point dont l'image est
$\Phi (-2)=\infty$. L'image de ce cercle est alors une droite et c'est
$$\Phi (2)+i\R = \frac{11}{4} + i\R .$$
\fincorrection
\nocorrection
\correction{002886}
On a $\Phi_\alpha (0) =\alpha$ et $\Phi_\alpha (\alpha )=0$. Remarquons que $\Phi_\alpha \circ \Phi_\alpha $ fixe l'origine.
Par l'exercice \ref{ex:burnol2.6}, l'automorphisme $\Phi_\alpha \circ \Phi_\alpha $ de $D(0,1)$ est une rotation $z\mapsto e^{i\alpha}z$.
Un calcul explicite montre que $\Phi_\alpha \circ \Phi_\alpha =\mathrm{Id}$, c'est \`a dire $\Phi_\alpha ^{-1}= \Phi_\alpha $.
Soit $\Psi$ un automorphisme du disque unit\'e $D(0,1)$ tel que $\Psi (z_1)=z_2$. Alors
$$\Psi \circ \Phi_{z_1}(0)= \Phi _{z_2}(0) \quad \Longleftrightarrow \quad
\Phi_{z_2}^{-1}\circ \Psi \circ \Phi_{z_1} (0)=0$$
et donc $A=\Phi_{z_2}^{-1}\circ \Psi \circ \Phi_{z_1}$ est un automorphisme du disque unit\'e fixant l'origine.
 On en d\'eduit de nouveau que $A$ est une rotation: $A(z)=e^{i\alpha} z$. Par cons\'equent,
 \begin{equation}\label{eq disque}\Psi = \Phi_{z_2}\circ A\circ \Phi_{z_1}^{-1}.\end{equation}
 On vient de d\'eterminer la forme g\'en\'erale d'un automorphisme $\Psi $ du disque unit\'e v\'erifiant
 $\Psi(z_1)=z_2$. Remarquons qu'il est unique ``\`a une rotation pr\`es''; $\Psi$ est d\'etermin\'e par \eqref{eq disque}
 o\`u $A$ est une rotation quelconque.
\fincorrection
\nocorrection
\nocorrection
\nocorrection
\nocorrection
\nocorrection
\nocorrection
\nocorrection
\nocorrection
\nocorrection
\nocorrection
\nocorrection
\nocorrection
\nocorrection
\nocorrection
\nocorrection
\nocorrection
\nocorrection
\nocorrection
\nocorrection
\nocorrection
\nocorrection
\nocorrection
\nocorrection
\nocorrection
\nocorrection
\nocorrection
\nocorrection
\nocorrection
\nocorrection
\nocorrection
\nocorrection
\nocorrection
\nocorrection
\nocorrection
\nocorrection
\nocorrection
\nocorrection
\nocorrection
\nocorrection
\nocorrection
\nocorrection
\nocorrection
\nocorrection
\nocorrection
\nocorrection
\nocorrection
\nocorrection
\nocorrection
\nocorrection
\nocorrection
\nocorrection
\nocorrection
\nocorrection
\nocorrection
\nocorrection
\nocorrection
\nocorrection
\nocorrection
\nocorrection
\nocorrection
\nocorrection
\nocorrection
\nocorrection
\correction{007597}
\begin{enumerate}
\item Par définition de $\cosh$, on a pour tout $z\in\Cc$, $\cosh (z)=\frac{e^z+e^{-z}}{2}$.
Par conséquent, $z$ est solution de $\cosh(z)=8i$ si et seulement si $e^z$ est solution de
$$X^2-8iX+1=0.$$
si et seulement si $e^z=(4+\sqrt{17})i=(4+\sqrt{17})e^{-i\frac{\pi}{2}}$ \\quad\text{ou}\quad $e^z=(4-\sqrt{17})i=(\sqrt{17}-4)e^{-i\frac{\pi}{2}}$,
si et seulement si $z\in\log (4+\sqrt{17})+i(\frac{\pi}{2}+2\pi\Zz)$ ou $z\in\log (\sqrt{17}-4)+i(-\frac{\pi}{2}+2\pi\Zz)$.
\item Soit $z\in\Cc$.
\begin{eqnarray*}
 z^i=-1&\iff& e^{i\log z}=e^{i\pi}\\
 &\iff& i\log z=i\pi \mod 2i\pi\\
 &\iff& \log z=\pi \mod 2\pi\\
 &\iff& \exists k\in\Zz, z = e^{(2k+1)\pi}.
\end{eqnarray*}
Les solutions sont les nombres complexes de la forme $e^{(2k+1)\pi}$ avec $k$ entier.
\end{enumerate}
\fincorrection
\correction{007598}
\begin{enumerate}
\item On définit une détermination du logarithme en posant
$$\begin{array}{cccc}
 l_+ :& \Cc-[0,+\infty[ & \longrightarrow & \Cc\\ & z=re^{i\theta} (r>0, 0<\theta<2\pi)&\longmapsto& \log_\Rr(r)+i\theta.
\end{array}$$

\item On définit une détermination du logarithme en posant
$$\begin{array}{cccc}
 \ell :& \Cc-L & \longrightarrow & \Cc\\ & z=re^{i\theta} (r>0, 0\leq\theta<2\pi)&\longmapsto& 
 \left\{\begin{array}{l} \textrm{si }\mid z-2 \mid \leq 1 \quad \text{ et } \quad  Im(z) \geq 0, \\ \ \ \ \ \ell(z)=\log_\Rr(r)+i(\theta +2\pi).\\
 \textrm{si }\mid z-2 \mid > 1, \\ \ \ \ \ \ell(z)=\log_\Rr(r)+i\theta\end{array}
 \right.
\end{array}$$
On vérifie que $\ell$ est continue : en particulier sur le disque ouvert de centre $(2,0)$ et de rayon $1$, $\ell$ coïncide avec la fonction $\log +2i\pi$,
qui est continue. 
On vérifie aussi que pour tout $z\in\Cc-L$, $e^{\ell(z)}=z$ : c'est donc une détermination du logarithme sur $\Cc-L$.
\end{enumerate}
\fincorrection
\correction{007599}
\begin{enumerate}
\item L'application 
$$\begin{array}{cccc}
  \gamma :& [0,2\pi[ &\longrightarrow & \Cc\\ & \theta &\longmapsto& e^{i\theta}
 \end{array}$$
est un paramétrage du cercle unité du plan complexe euclidien parcouru dans le sens trigonométrique.
\item Comme le rayon de convergent de la série qui définit $\cosh$ est infini, 
la série de fonctions $\sum_{n=0}^{+\infty} \frac{z^{2n-1}}{(2n)!}$ est normalement donc uniformément convergente sur le cercle unité vers
la fonction $\frac{\cosh z}{z}$.
Par conséquent, 
\begin{eqnarray*}\int_{\partial\Delta} \frac{\cosh z}{z} dz&=&\sum_{n=0}^{+\infty} \int_{\partial\Delta}\frac{z^{2n-1}}{(2n)!}
=\sum_{n=0}^{+\infty} \int_0^{2\pi}\frac{(\gamma(\theta))^{2n-1}}{(2n)!}\gamma'(\theta)d\theta\\
&=&\sum_{n=0}^{+\infty} \int_0^{2\pi} i e^{i(2n)\theta} d\theta=2i\pi.
\end{eqnarray*}
\item Comme l'intégrale de la fonction $\frac{\cosh(z)}{z}$ sur le chemin fermé $\partial\Delta$ de $\Cc^\times$ n'est pas nulle,
la fonction $\frac{\cosh(z)}{z}$ n'admet pas de primitive sur $\Cc^\times$.
\end{enumerate}
\fincorrection
\correction{007600}
\begin{enumerate}
  \item Question de cours.
  \item L'application $\sin$ est holomorphe sur $\Cc$, donc $$I_1=\int_{\partial \Delta}\frac{\sin(z)}{z}dz=\int_{\partial \Delta}\frac{\sin(z)}{z-0}dz=(2i\pi)\sin(0)=0.$$

  \item L'application $\frac{\sin(z)}{z}$ se prolonge en $0$ par la somme de la série entière $\sum_{n\in\Nn}\frac{(-1)^n}{(2n+1)!}z^{2n}$ de rayon de convergence infini donc holomorphe sur $\Cc$.
  Donc $$I_2=\int_{\partial \Delta}\frac{\sin(z)}{z^2}dz=\int_{\partial \Delta}\frac{\frac{\sin(z)}{z}}{z-0}dz=(2i\pi)\sum_{n=0}^{+\infty}\frac{(-1)^n}{(2n+1)!}z^{2n}(0)=2i\pi.$$

  \item L'application $\frac{\sin(z)-1}{z^2}$ se prolonge en $0$ par la somme de la série entière $\sum_{n\geq 1}\frac{(-1)^n}{(2n+1)!}z^{2n-1}$ holomorphe sur $\Cc$.
  Donc $\int_{\partial \Delta}\frac{\sin(z)-1}{z^3}dz=(2i\pi)\sum_{n=1}^{+\infty}\frac{(-1)^n}{(2n+1)!}z^{2n-1}(0)=0$.
  Donc, 
  \begin{eqnarray*}
  I_3&=&\int_{\partial \Delta}\frac{\sin(z)}{z^3}dz=\int_{\partial \Delta}\frac{\sin(z)-1}{z^3}dz + \int_{\partial \Delta}\frac{1}{z^3}dz
  =\int_{\partial \Delta}\frac{1}{z^3}dz=0.
  \end{eqnarray*}

  \item L'application $\sin(z)$ admet $-\cos(z)$ comme primitive sur $\Cc^\times$ et même sur $\Cc$.
  
  
  L'application $\frac{\sin(z)}{z}$ se prolonge sur $\Cc$ en la somme de la série entière $\sum_{n\in\Nn}\frac{(-1)^n}{(2n+1)!}z^{2n}$.
  Elle admet donc sur $\Cc$ et donc sur $\Cc^\times$ la primitive $\sum_{n\in\Nn}\frac{(-1)^n}{(2n+1)!}\frac{z^{2n+1}}{2n+1}$.
  
  
  L'application $\frac{\sin(z)}{z^2}$ a une intégrale non nulle sur le chemin fermé $\partial\Delta$ de $\Cc^\times$ : elle n'admet donc pas de primitive sur $\Cc^\times$.
  
  
  L'application $\frac{\sin(z)}{z^3}=\frac{\sin(z)-1}{z^3}+\frac{1}{z^3}=\sum_{n=1}^{+\infty}\frac{(-1)^n}{(2n+1)!}z^{2n-1}+\frac{1}{z^3}$ 
  admet comme primitive sur $\Cc^\times$
  $$\sum_{n=1}^{+\infty}\frac{(-1)^n}{(2n+1)!}\frac{z^{2n}}{2n}-\frac{1}{4z^4}.$$

\end{enumerate}
\fincorrection
\correction{007601}
\begin{enumerate}
  \item On définit une détermination du logarithme en posant
$$\begin{array}{cccc}
 l_+ :& \Cc-]-\infty,0] & \longrightarrow & \Cc\\ & z=re^{i\theta} (r>0, -\pi<\theta<\pi)&\longmapsto& \log_\Rr(r)+i\theta.
\end{array}$$
C'est une application continue, et même holomorphe telle que $$\forall z\in\Cc-]-\infty,0],\ \ e^{l_+(z)}=z.$$

  \item On définit une détermination du logarithme en posant
$$\begin{array}{cccc}
 \ell :& \Cc-L & \longrightarrow & \Cc\\ & z=re^{i\theta} (r>0, 0\leq\theta<2\pi)&\longmapsto& 
 \left\{\begin{array}{l} \textrm{sur } U, \\ \ \ \ \ \ell(z)=\log_\Rr(r)+i(\theta +2\pi).\\
 \textrm{sur } V, \\ \ \ \ \ \ell(z)=\log_\Rr(r)+i\theta\end{array}
 \right.
\end{array}$$
On vérifie que $\ell$ est continue : en particulier sur le disque ouvert de centre $(2,0)$ et de rayon $1$, $\ell$ coïncide avec la fonction $\log +2i\pi$,
qui est continue. 
On vérifie aussi que pour tout $z\in\Cc-L$, $e^{\ell(z)}=z$ : c'est donc une détermination du logarithme sur $\Cc-L$.
\end{enumerate}
\fincorrection
\correction{007602}
\begin{enumerate}
  \item Soit $\Delta_r(a)$ le disque de centre $a\in\Cc$ et de rayon $r>0$. Supposons que $Im(f)$ ne rencontre pas $\Delta_r(a)$.
Alors,
$$\forall z\in\Cc,\ \ \ |f(z)-a|\geq r.$$
L'application $\Cc\to\Cc, z\mapsto\frac{1}{f(z)-a}$ est alors bien définie, 
holomorphe sur $\Cc$ et majorée en module par $\frac{1}{r}$.
Par le théorème de Liouville, elle est donc constante.
Puisque $w\mapsto \frac{1}{w-a}$ est injective sur $\Cc-\{a\}$, en déduit que $f$ est constante.
  \item Une telle application entière ne rencontre par le disque $\Delta_1(-2i)$ : elle est donc constante,
d'après la question précédente.
  \item Soit $z\in\Cc-\{-i\}$,
\begin{eqnarray*}
 1-|h(z)|^2&=& 1-\frac{z-i}{z+i}\frac{\overline{z}+i}{\overline{z+i}}
 = 1-\frac{|z|^2+i(z-\overline{z})+1}{|z+i|^2}\\&=&1-\frac{|z|^2-i(z-\overline{z})+1+2i(z-\overline{z})}{|z+i|^2}
= \frac{4 Im(z)}{|z+i|^2}.
\end{eqnarray*}
  \item Sur $\mathbb{H}$, $Im(z)>0$ et donc $|h(z)|<1$. L'image du demi-plan $\mathbb{H}$ par l'application $h$ est donc
incluse dans le disque unité $\Delta$ ; c'est donc une partie bornée de~$\Cc$.
  \item Soit $f :\Cc\to\mathbb{H}$ une application holomorphe.
Par composition, $h\circ f :\Cc\to \Cc$ est une application holomorphe dont l'image est bornée.
Par le théorème de Liouville, elle est donc constante. Comme $h$ est injective, on en déduit que $f$ est constante.
\end{enumerate}
\fincorrection
\correction{007603}
\begin{enumerate}
  \item Soit $r>0$ tel que $\overline{\Delta_r(c)}$ est un disque fermé inclus dans $D$.
Le résidu de $f$ en $c$ est 
$$Res_c(f):=\frac{1}{2i\pi}\int_{\partial\Delta_r(c)}f(z)dz.$$
Cette intégrale ne dépend pas de $r$ tel que $\overline{\Delta_r(c)}$ est inclus dans $D$.
  \item Si $f$ admet une primitive sur $D-\{c\}$, l'intégrale de $f$ sur tout chemin fermé de $D-\{c\}$ est nulle.
 En particulier, le résidu de $f$ en $c$ est nul.
  \item D'après la question précédente, il est nécessaire que le résidu de $f$ en $c$ soit nul.
  Réciproquement, si le résidu de $f$ en $c$ est nul, montrons que $f$ admet une primitive sur $\Delta$.
  Par le développement au voisinage de $c$, on sait qu'il existe $a_{-3}, a_{-2}, a_{-1}\in\Cc$ et une fonction $\tilde f$
  holomorphe sur $\Delta$ tels que
  $$\forall z\in\Delta,\ \ f(z)=\frac{a_{-3}}{(z-c)^3}+ \frac{a_{-2}}{(z-c)^2}+ \frac{a_{-1}}{z-c}+ \tilde{f}(z).$$
  Comme le résidu de $f$ en $c$ est nul, $a_{-1}=0$.
  L'application $z\mapsto \frac{a_{-3}}{(z-c)^3}+ \frac{a_{-2}}{(z-c)^2}$ admet 
  $z\mapsto -\frac{a_{-3}}{2(z-c)^2}- \frac{a_{-2}}{z-c}$ comme primitive sur $\Delta-\{c\}$.
  Puisque $\Delta$ est étoilé et que $\tilde{f}$ est holomorphe sur $\Delta$, elle admet une primitive sur $\Delta$
  donc sur $\Delta-\{c\}$.
  En conclusion, si le résidu de $f$ en $c$ est nul, $f$ admet-elle une primitive sur $\Delta-\{c\}$.
\end{enumerate}
\fincorrection
\correction{007604}
\begin{enumerate}
  \item Soit $D$ un ouvert étoilé de $\Cc$. Soit $\Gamma$ un chemin fermé dans $D$. Soit $C$ un ensemble fini de points de $D-\Gamma$.
Soit $f : D-C\to\Cc$ holomorphe.
Alors,
$$\sum_{c\in C\cap Int(\Gamma)} Ind_\Gamma(c) Res_c(f)=\frac{1}{2i\pi}\int_\Gamma f(z)dz.$$
  \item Soit $D$ un ouvert de $\Cc$ et $f : D\to \Cc$ une application holomorphe.
Soit $\Delta_r(a)$ un disque dont l'adhérence est incluse dans $D$.
Soit $b\in\Delta_r(a)$. 
Soit $R>0$ tel que $\overline{\Delta_r(a)}\subset \Delta_R(a)\subset D$.
L'ouvert $\Delta_R(a)$ est étoilé, le chemin $\partial\Delta_r(a)$ est inclus dans $\Delta_R(a)$
et ne contient pas $b$ ; l'application $f : \Delta_R(a)-\{b\}\to\Cc,\ \ z\mapsto \frac{f(z)}{z-b}$
est holomorphe : par le théorème des résidus, on en déduit
$$Ind_{\partial\Delta_r(c)}(b) Res_b(\frac{f(z)}{z-b})= Res_b(\frac{f(z)}{z-b})
=\frac{1}{2i\pi}\int_\Gamma \frac{f(z)}{z-b}dz.$$
Reste à noter que, en développant $f$ en série entière au voisinage de $b$, $Res_b(\frac{f(z)}{z-b})=f(b).$
  \item On réitère le raisonnement précédent avec l'application 
$f_k : \Delta_R(a)-\{b\}\to\Cc, z\mapsto \frac{f(z)}{(z-b)^{k+1}}$.
Le résidu de $\frac{f(z)}{(z-b)^{k+1}}$ en $b$ est le coefficient de $(z-b)^k$ dans 
le développement en série entière de $f$ centré en $b$.
Comme ce développement est donné par le développement de Taylor,
le coefficient de $(z-b)^k$ est $\frac{f^{(k)}(b)}{k!}$.
\end{enumerate}
\fincorrection
\correction{007605}
\begin{enumerate}
  \item Soit $f$ et $g$ deux applications holomorphes sur un ouvert étoilé $D$ de $\Cc$.
Soit $\Gamma$ un chemin fermé simple dans $D$. 
On suppose que $$\forall z\in \Gamma,\ \ |f(z)-g(z)|<|g(z)|.$$
Alors $f$ et $g$ ont le même nombre de zéros compté avec multiplicité dans l'intérieur de $\Gamma$.
  \item L'ouvert $\Cc$ est étoilé et contient le chemin $\partial\Delta_2$ dont l'intérieur est $\Delta_2$.
Les deux applications $z\mapsto z^4$ et $p$ sont polynômiales donc holomorphes sur $\Cc$.
Soit $z\in\partial\Delta_2$. Alors 
$$|p(z)-z^4|=|7z+1|\leq 7|z|+1\leq 7\times 2+1=15<2^4=|z^4|.$$
Par le théorème de Rouché, le polynôme $p$ et $z\mapsto z^4$ ont le même nombre de zéros compté avec multiplicités
dans $\Delta_2$, c'est à dire $4$.
Comme $p$ est non nul de degré $4$, tous ses zéros sont dans $\Delta_2$.
\end{enumerate}
\fincorrection
\correction{007606}
\begin{enumerate}
  \item L'application $f~:~\Cc\to\Cc, z\mapsto\overline{z}$ est associée à l'application $f_\Rr~:~(x,y)\mapsto (x,-y)$.
Comme $f_\Rr$ est polynomiale, elle est de classe $C^{\infty}$. Par contre $\frac{\partial re(f)}{\partial x}=1=-\frac{\partial Im(f)}{\partial y}$.
L'application $f$ qui ne vérifie les équation de Cauchy-Riemann en aucun point n'est pas holomorphe.
  \item On paramètre le cercle de centre $3$ et de rayon $1$ par $z=3+e^{i\theta}$.
 $$\int_{|z-3|=1}\frac{dz}{z-3}=\int_0^{2\pi}\frac{ie^{i\theta}d\theta}{e^{i\theta}}=i\int_0^{2\pi}d\theta=2i\pi\not=0.$$
 Comme l'intégrale de $z\mapsto\frac{1}{z-3}$ sur le chemin fermé cercle de centre $3$ et de rayon $1$ inclus dans le domaine $\Cc-\{3\}$ n'est pas nulle,
 cette application n'admet pas de primitive sur $\Cc-\{3\}$.
  \item Le disque unité est convexe : il est donc étoilé par rapport à chacun de ses points.
  \item L'ouvert $\Cc-\{3\}$ n'est pas étoilé, car l'application holomorphe $z\mapsto\frac{1}{z-3}$ n'y admet pas de primitive.
\end{enumerate}
\fincorrection
\correction{007607}
\begin{enumerate}
  \item Le point d'affixe $r\frac{1+it}{1-it}$ est de module $r$ et ses coordonnées $x=re(\xi(t))=r\frac{1-t^2}{1+t^2}$ et $y=Im(\xi(t))=r\frac{2t}{1+t^2}$ vérifient l'équation $y=t(x+r)$.
Il est donc à l'intersection de la droite d'équation $y=t(x+r)$ avec le cercle $C_r$. Comme l'équation $r\frac{1+it}{1-it}=-r$ est équivalente à $1=-1$, l'affixe $r\frac{1+it}{1-it}$
n'est pour aucune valeur de $t$ égale à $-r$.

  \item Comme $t$ est réel, $t\mapsto r\frac{1+it}{1-it}$ est quotient de deux polynômes dont le dénominateur ne s'annule jamais. Elle est donc dérivable et 
$$\frac{\xi'(t)}{\xi(t)}=\frac{i}{1+it}-\frac{-i}{1-it}=\frac{2i}{1+t^2}.$$

  \item $$\int_{C_r} \frac{dz}{z}=\int_{C_r-\{-r\}} \frac{dz}{z}=\int_{-\infty}^{+\infty}\frac{\xi'(t)}{\xi(t)}dt=\int_{-\infty}^{+\infty}\frac{2i}{1+t^2}.$$

  \item Par ailleurs, $\int_{C_r} \frac{dz}{z}=2i\pi Res_0(\frac{1}{z})=2i\pi.$
Donc, $\int_{-\infty}^{+\infty}\frac{dt}{1+t^2}=\pi$.
\end{enumerate}
\fincorrection
\correction{007608}
\begin{enumerate}
  \item Soit $z\in\mathbb{H}$.
\begin{eqnarray*}Im(h_A(z)&=&\frac{1}{2i}\left(\frac{az+b}{cz+d}-\overline{\frac{az+b}{cz+d}}\right)=\frac{1}{2i}\left(\frac{az+b}{cz+d}-\frac{a\overline{z}+b}{c\overline{z}+d}\right)\\
&=&\frac{1}{2i}\frac{(ad-bc)(z-\overline{z})}{|cz+d|^2}
=\frac{\det A}{|cz+d|^2}Im(z)>0.
\end{eqnarray*}
Donc, $h_A(z)\in\mathbb{H}$ et $h_A$ envoie le demi-plan $\mathbb{H}$ sur lui-même.
  \item Soit $z=x+iy\in\mathbb{H}$. Soit $A\in SL(2,\Rr)$.
$$h_A(i)=z\iff ai+b=z(ci+d)\iff ai+b=(dx-cy)+i(cx+dy)$$
$$ \iff \left\{\begin{array}{ccc} b&=&dx-cy\\ a&=&cx+dy\end{array}\right..$$
    On peut par exemple choisir, puisque $y$ est strictement positif, $A=\frac{1}{\sqrt{y}}\begin{pmatrix}
    y&x\\
                    0&1
                   \end{pmatrix}$
\end{enumerate}
\fincorrection
\correction{007609}
\begin{enumerate}
  \item Soit $D$ un ouvert de $\Cc$ et $f : D\to \Cc$ une application holomorphe.
Alors, pour tout disque $\Delta_r(a)$ dont l'adhérence est incluse dans $D$,
$$\forall b\in\Delta_r(a), \ \ \ \frac{1}{2i\pi}\int_{\partial\Delta_r(a)}\frac{f(z)dz}{z-b}=f(b).$$

% 
% \item Soit $D$ un ouvert étoilé de $\Cc$. Soit $\Gamma$ un chemin fermé dans $D$. Soit $C$ un ensemble fini de points de $D-\Gamma$.
% Soit $f : D-C\to\Cc$ holomorphe.
% Alors,
% $$\sum_{c\in C\cap Int(\Gamma)} Ind_\Gamma(c) Res_c(f)=\frac{1}{2i\pi}\int_\Gamma f(z)dz.$$

  \item Soit $f$ une fonction holomorphe sur $\Delta$. On paramètre le cercle $\partial\Delta_{\frac{1}{2}}$ par $z=\frac{1}{2}e^{i\theta}$.
Alors, $f(0)=\frac{1}{2i\pi}\int_{\partial\Delta}\frac{f(z)dz}{z-0}=\frac{1}{2i\pi}\int_0^{2\pi}\frac{f(\frac{1}{2}e^{i\theta})}{\frac{1}{2}e^{i\theta}}\frac{1}{2}ie^{i\theta}d\theta
=\frac{1}{2\pi}\int_0^{2\pi}f(\frac{1}{2}e^{i\theta})d\theta$ est une quantité réelle car $f$ prend des valeurs réelle sur le cercle $\partial\Delta_{\frac{1}{2}}$.
L'application $f$ prend donc une valeur réelle en $0$.
  \item Soit $f$ une fonction holomorphe sur $\Delta$ constante égale à $c$ sur le cercle $\partial\Delta_{\frac{1}{2}}$. 
Soit $b\in\Delta_{\frac{1}{2}}$. Par le théorème de représentation, 
$$f(b)=\frac{1}{2i\pi}\int_{\partial\Delta_{\frac{1}{2}}}\frac{f(z)dz}{z-b}=c\frac{1}{2i\pi}\int_{\partial\Delta_{\frac{1}{2}}}\frac{dz}{z-b}
=c Res_b(\frac{1}{z-b})=c.$$
\end{enumerate}
\fincorrection
\correction{007610}
\begin{enumerate}
  \item Soit $f:\Cc\to\Cc$ une application holomorphe bornée en module par $M$.
On sait qu'elle est développable en séries entières sur $\Cc$ avec une série $\sum_{n\geq 0} a_n z^n$ de rayon de convergence infini.
On peut donc appliquer la formule de Gutzmer : pour tout $r\in\Rr^+$,
$$\sum_{n\geq 0} |a_n|^2r^{2n}=\frac{1}{2\pi}\int_0^{2\pi}|f(re^{i\theta})|^2d\theta\leq M^2 \frac{1}{2\pi}\int_0^{2\pi}d\theta=M^2$$
Donc, pour tout $n$ in $\Nn$ et tout $r\in\Rr^+$,
$$|a_n|^2r^{2n}\leq M^2.$$
En particulier, tous les $a_n$ avec $n\not =0$ sont nuls et $f$ est constante
  \item Soit $f:\Cc\to\Cc$ une application holomorphe telle que $ M(r)\leq r$.
$$\sum_{n\geq 0} |a_n|^2r^{2n}=\frac{1}{2\pi}\int_0^{2\pi}|f(re^{i\theta})|^2d\theta\leq M(r)^2 \frac{1}{2\pi}\int_0^{2\pi}d\theta=M(r)^2\leq r^2$$
Donc, pour tout $n$ in $\Nn$ et tout $r\in\Rr^+$, $|a_n|^2r^{2n}\leq r^2$.
En particulier, tous les $a_n$ avec $n\not =0,1$ sont nuls et $f$ est affine.
\end{enumerate}
\fincorrection
\correction{007611}
\begin{enumerate}
  \item On note $R$ le rayon de convergence de la série entière $\sum_{n\geq 1}a_nz^n$ et $r\in]0,R[$. 
On sait alors que $(a_n r^n)_{n\in\Nn}$ est une suite bornée en module.
Soit $z\in\Cc$
$$\left|\frac{a_n}{n!}z^n\right|\leq \left|a_n r^n\right| \frac{\left|\frac{z}{r}\right|^n}{n!}$$
est le terme général d'une suite bornée en module car $\frac{\left|\frac{z}{r}\right|^n}{n!}$ 
est borné.
Par conséquent, $\sum_{n\geq 1}\frac{a_n}{n!}z^n$ a un rayon de convergence infini.
  \item L'application $f-g$ est somme sur $\Delta_r$ connexe d'une série entière centrée en $0$ et s'annule sur l'ensemble non discret $]-R,R[$.
Par le principe des zéros isolés, $f=g$ sur $\Delta_r$.
  \item Soit $a\in\Rr$. Les deux applications $z\mapsto \sin(a+z)$ et $z\mapsto \sin(a)\cos (z)+\cos(a)\sin (z)$ sont développables en séries entières
sur $\Cc$ et coïncident sur $\Rr$ par la formule d'addition des sinus sur $\Rr$.
Par la question précédente, elles coïncident sur $\Cc$.
\end{enumerate}
\fincorrection
\correction{007612}
\begin{enumerate}
  \item Soit $z\in\Cc$.

\begin{eqnarray*}
 \sin z=c&\iff& \frac{e^{iz}-e^{-iz}}{2i}=c\iff e^{iz}-2ic-e^{-iz}=0\\
 &\iff& (e^{iz})^2-2ic e^{iz}-1=0 \textrm{ car } e^{iz}\not=0.
\end{eqnarray*}
  \item Soit $c\in [-1,1]$. Soit $z\in\Cc$ une solution de $\sin z = c$. Alors $e^{iz}$ est solution de $X^2-2icX-1=0$.
Comme $c\in [-1,1]$, les solutions de cette équation sont $ic+\sqrt{1-c^2}$ et $ic-\sqrt{1-c^2}$,
toutes les deux de module $1$, 
$$|e^{iz}|=e^{-y}=1$$
et donc $y=0$, et $z$ est donc réel.
  \item \begin{eqnarray*}
 e^{i(a+b)}&=&e^{ia}e^{ib}=(\cos a+i\sin a)(\cos b+i\sin b)\\
 &=&(\cos a\cos b -\sin a\sin b)+i(\sin a\cos b+\cos a\sin b).
\end{eqnarray*}
\begin{eqnarray*}
 e^{-i(a+b)}&=&e^{-ia}e^{-ib}=(\cos a-i\sin a)(\cos b-i\sin b)\\
 &=&(\cos a\cos b -\sin a\sin b)-i(\sin a\cos b+\cos a\sin b).
\end{eqnarray*}
  \item \begin{eqnarray*}
 \sin (a+b)&=&\frac{e^{i(a+b)}-e^{-i(a+b)}}{2i}=\sin a\cos b+\cos a\sin b
\end{eqnarray*}
d'après la question précédente.
  \item Soit $z\in\Cc$.
\begin{eqnarray*}
 \cos z+\sin z=2&\iff& \frac{1}{\sqrt{2}}\cos z+\frac{1}{\sqrt{2}}\sin z=\sqrt{2}\\
 &\iff& \sin (\frac{\pi}{4}+z)=\sqrt{2}\\
 &\iff& (e^{iz}) \textrm{ est solution de }X^2-2i\sqrt{2}X-1=0\\
 &\iff& e^{iz}=i\sqrt{2}+i \textrm { ou } e^{iz}=i\sqrt{2}-i\\
 &\iff& e^{iz}=(\sqrt{2}+1)e^{\frac{\pi}{2}} \textrm { ou } e^{iz}=(\sqrt{2}-1)e^{\frac{\pi}{2}}\\
 &\iff&\exists k\in\Zz, iz=\log_\Rr(\sqrt{2}+1)+i(\frac{\pi}{2}+2k\pi) \\&& \textrm { ou } iz=\log_\Rr(\sqrt{2}-1)+i(\frac{\pi}{2}+2k\pi)\\
 &\iff&\exists k\in\Zz, z=(\frac{pi}{2}+2k\pi)-i\log_\Rr(\sqrt{2}+1) \\&& \textrm { ou } z=(\frac{\pi}{2}+2k\pi)-i\log_\Rr(\sqrt{2}-1)
\end{eqnarray*}
Les soultions de l'équation $\cos z+\sin z=2$ dans $\Cc$ sont les nombres complexes de la forme
$(\frac{pi}{2}+2k\pi)-i\log_\Rr(\sqrt{2}+1)$ ou $z=(\frac{pi}{2}+2k\pi)-i\log_\Rr(\sqrt{2}-1)$
avec $k$ entier.
\end{enumerate}
\fincorrection
\correction{007613}
\begin{enumerate}
  \item \begin{eqnarray*}
 Im\left(\frac{az+b}{cz+d}\right)&=&\frac{1}{2i}(\frac{az+b}{cz+d}-\overline{\frac{az+b}{cz+d}})\\
 &=&\frac{1}{2i}(\frac{az+b}{cz+d}-\frac{a\overline{z}+b}{c\overline{z}+d})=\frac{1}{2i}\frac{(ad-bc)(z-\overline{z})}{|cz+d|^2}\\
 &=&\frac{Im(z)}{|cz+d|^2}
\end{eqnarray*}
Par conséquent, si $Im(z)>0$, alors $Im\left(h_A(z)\right)>0$ et $h_A$ envoie $\mathbb{H}$ sur $\mathbb{H}$.
  \item Soit $z=x+iy$. On a $\frac{xi-y}{1\times i+0}=z$. La matrice $\begin{pmatrix} x&-y\\1&0\end{pmatrix}$
convient ainsi donc que la matrice $\frac{1}{y}\begin{pmatrix} x&-y\\1&0\end{pmatrix}$ de déterminant $1$.
\end{enumerate}
\fincorrection
\correction{007614}
\begin{enumerate}
  \item Soit $z=re^{i\theta}\in\Cc$ avec $r\geq 0$ et $\theta\in ]-\pi,\pi]$.
\begin{eqnarray*}
 z^3\not\in\Cc^-&\iff& z^3\in\Rr^-\iff r^3e^{i3\theta}\in\Rr^-\iff 3\theta\simeq \pi \mod 2\pi\\
 &\iff& \theta \simeq \frac{\pi}{3} \mod \frac{2\pi}{3}\iff \theta=\frac{-\pi}{3}\quad\text{ou}\quad \theta=\frac{\pi}{3}\quad\text{ou}\quad \theta=\pi.
\end{eqnarray*}

\begin{tikzpicture}
%\draw[help lines, step=0.5, very thin] (-6,-4) grid (6,4);
\draw [thick] (-5,0) -- (0,0) ;\draw [densely dotted] (5,0) -- (0,0) ;
\draw (0.5,0) arc (0:60:0.5) ; \draw (0.5,0.5) node[right]{$\pi/3$} ;
\draw [thick](0,0) -- (2.5,4.33) ;
\draw (1,0) arc (0:-60:1) ;\draw (1,-1) node[right]{$-\pi/3$} ;
\draw [thick](0,0) -- (2.5,-4.33) ;
\end{tikzpicture}


Le domaine $D$ est composé de $\Cc$ moins les demi-droites noires. 
% \begin{tikzpicture}
%  \draw [double distance=3pt, thick] (0.5,0) arc (0:60:0.5);
%  \draw (-1.5,0) arc (0:30:0.5);
%  \draw [color=gray](0,0) circle (2cm);
%  \draw[thick] (0,0)-- (1,1.73);
%  \draw [thick] (0,0) -- (2,0);
%  \draw [densely dotted] (-2,0)-- (1,1.73);
%  \draw [densely dotted](-2,0)-- (2,0);
%  \draw (1.26,-1.29) arc (60:90:0.5);
%  \draw [densely dotted] (1,-1.73) -- (2,0);
%  \draw [densely dotted] (1,-1.73) -- (1,1.73);
% \end{tikzpicture}
  \item Soit $z=re^{i\theta}\in\Cc$ avec $r\geq 0$ et $\theta\in ]-\pi,\pi]$.
\begin{itemize}
 \item Si $\theta\in]-\frac{\pi}{3},\frac{\pi}{3}[$, $ 3\theta\in]-\pi,\pi[$
 et donc, $f(z)=\log(z^3)=\log_\Rr(r^3)+i3\theta=3\log_\Rr(r)+3i\theta=3\log(z)$.
 \item Si $\theta\in]\frac{\pi}{3},\pi[$, $ 3\theta\in]\pi,3\pi[$ et $ 3\theta-2\pi\in]-\pi,\pi[$
 et donc, $f(z)=\log(z^3)=\log_\Rr(r^3)+i(3\theta-2\pi)=3\log_\Rr(r)+3i\theta-2i\pi=3\log(z)-2i\pi$.
 \item Si $\theta\in]-\pi,-\frac{\pi}{3},[$, $ 3\theta\in]-3\pi,-\pi[$ et $ 3\theta+2\pi\in]-\pi,\pi[$
 et donc, $f(z)=\log(z^3)=\log_\Rr(r^3)+i(3\theta+2\pi)=3\log_\Rr(r)+3i\theta+2i\pi=3\log(z)+2i\pi$.
\end{itemize}
La fonction $f$ n'est donc pas continu aux points de la demi-droite d'angle $\frac{\pi}{3}$ 
ni aux points de la demi-droite d'angle $-\frac{\pi}{3}$ où la fonction $_log$ est continue.
On sait que le saut de la branche principale $\log$ du logarithme entre $(-\pi^+$ et $\pi^-$
est de $2i\pi$. Le saut de $f$ est donc de $6i\pi$. Elle n'est donc pas continue
aux points de la demi-droite d'angle $\pi$.
  \item La réponse a été donnée dans la question précédente.
\end{enumerate}
\fincorrection
\nocorrection
\nocorrection
\nocorrection
\nocorrection
\nocorrection
\nocorrection
\nocorrection
\nocorrection
\nocorrection
\nocorrection
\nocorrection
\nocorrection
\correction{007627}
\begin{enumerate}
  \item 
$$\begin{array}{cccc}
 \log :& \Cc-]-\infty,0]&\longrightarrow&\Cc\\
    &z         &\longmapsto& \ln_\Rr(r)+i\theta \ \ \substack{\textrm{ après avoir choisi l'écriture } z=re^{i\theta}\\ \textrm{ avec } r\in[0,+\infty[ \textrm{ et } \theta\in]-\pi,\pi[}
\end{array}$$

% $$\newfrac{\ln_\Rr(r)+i\theta}{\textrm{après avoir choisi l'écriture } z=re^{i\theta} \textrm{ avec } r\in[0,+\infty[ \textrm{ et } \theta\in]-\pi,\pi[}$$
  \item Une primitive d'une application continue $f$ sur $\Cc$ est une application holomorphe $F$ sur $\Cc$ dont la dérivée complexe $F'$ est l'application $f$.
  \item Comme la dérivée complexe d'une application holomorphe est holomorphe, il n'y a pas de tels exemples.
  \item $$\begin{array}{cccc}
 Id :& \Delta &\longrightarrow&\Cc\\
    &z         &\longmapsto& z
\end{array}$$
est holomorphe non-constante.
  \item Par le théorème de Liouville, toute application entière bornée est constante~: il n'y a donc pas de tels exemples.
  \item $$\begin{array}{cccc}
 C :& \mathbb{H} &\longrightarrow&\Delta\\
    &z         &\longmapsto& \frac{z-i}{z+i}
\end{array}\textrm{ ou bien } 
\begin{array}{cccc}
 e :& \mathbb{H} &\longrightarrow&\Delta\\
    &z         &\longmapsto& \exp(2i\pi z)
\end{array}
$$
\end{enumerate}
\fincorrection
\correction{007628}
~
 \begin{center} 
 \begin{tikzpicture}
 \draw [step=1cm, gray, very thin](-4,-4) grid (4,4);
 \draw (-4,0) -- (4,0); \draw (0,-4) -- (0,4);
 \draw (0,0) node {$\bullet$}; \draw (0,0) node[below left] {$P_1$};
 \draw (2,0) node {$\bullet$};\draw (2,0) node [below right] {$P_2$};
 \draw (-2,0) node {$\bullet$};\draw (-2,0) node [below left] {$P_3$};
\draw (1,0) node {$\bullet$}; \draw (1,0) node [below right] {$C$} circle (2) ;
\draw (2,2) node[below right] {${\partial \Delta_2(1)}$};
\end{tikzpicture}
\end{center}
On écrit la décomposition en éléments simples de la fraction rationnelle
$$\frac{z^2+z-1}{z^2(z^2-4)}=\frac{5}{16}\frac{1}{z-2}-\frac{1}{16(z+2)}-\frac{1}{4z}+\frac{1}{4z^2}.$$
Les premier et troisième termes correspondent aux pôles simples $P_2$ et $P_1$ à l'intérieur du cercle ${\partial \Delta_2(1)}$.
Donc, 
$$\int_{\partial \Delta_2(1)}\left(\frac{5}{16}\frac{1}{z-2}-\frac{1}{4z}\right)dz=(2i\pi)(\frac{5}{16}-\frac{1}{4})=\frac{i\pi}{8}.$$
Le deuxième terme correspond au pôle $P_3$ à l'extérieur du cercle ${\partial \Delta_2(1)}$ donc, $\int_{\partial \Delta_2(1)}\frac{1}{16(z+2)}dz=0$.
Le dernier terme est l'intégrale sur un chemin fermé de l'application exacte $z\mapsto \frac{1}{4z^2}$ de primitive $z\mapsto -\frac{1}{4z}$.
Donc, $\int_{\partial \Delta_2(1)}\frac{1}{4z^2}dz=0$.
En conclusion, 
$$\int_{\partial \Delta_2(1)}\frac{z^2+z-1}{z(z^2-4)}dz=\frac{i\pi}{8}.$$

Une autre démarche serait : 
on écrit la décomposition en éléments simples de la fraction rationelle
$$\frac{1}{z^2(z^2-4)}=\frac{1}{16}\frac{1}{z-2}-\frac{1}{16(z+2)}-\frac{1}{4z^2}.$$
et donc avec $f(z)=z^2+z-1$ holomorphe sur $\Cc$
$$\frac{z^2+z-1}{z^2(z^2-4)}=\frac{1}{16}\frac{f(z)}{z-2}-\frac{f(z)}{16(z+2)}-\frac{f(z)}{4z^2}.$$
Par le théorème de Cauchy sur les disques, puisque $0$ et $2$ sont dans $\Delta_2(1)$ mais pas $-2$,
\begin{eqnarray*}\int_{\partial \Delta_2(1)}\frac{z^2+z-1}{z^2(z^2-4)}&=&\frac{1}{16}\int_{\partial \Delta_2(1)}\frac{f(z)}{z-2}-\int_{\partial \Delta_2(1)}\frac{f(z)}{16(z+2)}-\int_{\partial \Delta_2(1)}\frac{f(z)}{4z^2}\\
&=&\frac{2i\pi f(2)}{16}-0-\frac{2i\pi f'(0)}{4}=\frac{5i\pi}{8}-\frac{i\pi}{2}
 =\frac{i\pi}{8}.
\end{eqnarray*}
\fincorrection
\correction{007629}
\begin{enumerate}
  \item Supposons que $0$ n'est pas dans l'adhérence de l'image de $f$. Il existe alors $r>0$ tel que $f(\Cc)\cap \Delta_r =\emptyset$. Autrement dit,
$$\forall z\in\Delta, \ |f(z)|\geq r.$$
L'application $f$ ne s'annule donc pas dans $\Cc$ et la fonction holomorphe $1/f$ est majorée par $1/r$ sur $\Cc$. Par le théorème de Liouville, elle est donc constante, ainsi que $f$.
  \item Soit $c\in \Cc$. En appliquant le résultat précédent à $f-c$ on obtient que $c$ est dans l'adhérence de l'image de $f$.
         Par conséquent, l'adhérence de l'image de $f$ est $\Cc$.
\end{enumerate}
\fincorrection
\correction{007630}
\begin{enumerate}
  \item Si $f$ avait une singularité apparente ou polaire en $0$, comme $z\mapsto 1/z$ a une singularité polaire en $0$,
$f+1/z$ aurait une singularité apparente ou polaire. 
Or pour tout entier naturel $n$, et pour tout réel $x$ strictement positif, $x^n\exp(\frac{1}{x})\geq x^n\frac{1}{(n+1)!x^{n+1}}$. 
Donc, $\lim_{x\to 0\atop x>0} x^n \exp (\frac{1}{x}) =+\infty$. Par conséquent  $z\mapsto\exp(\frac{1}{z}) $ a une singularité essentielle en~$0$, 
ainsi que $f$.
  \item En écrivant le développement en séries entières de l'exponentielle, on obtient
                 $$\forall z\in\Cc^\star,\ \ f(z)=1+\sum_{n=2}^{+\infty} \frac{z^{-n}}{n!}.$$
Or la série, $z+ \sum_{n=2}^{+\infty} \frac{z^{-n+1}}{(-n+1)n!}$ est, comme l'application exponentielle, normalement convergente sur $\Cc^\star$ et y définit donc une application holomorphe $F$ telle que $F'$ qui se calcule comme somme des dérivées, par convergence normale, vérifie $F'=f$ : c'est une primitive de $f$.
  \item Puisque l'application $f$ est exacte sur $D$ et que $\partial\Delta$ est un chemin orienté fermé dans $D$,
 $\int_{\partial\Delta}fdz=0$. On en déduit 
 $$\int_{\partial\Delta}\exp(\frac{1}{z})dz=\int_{\partial\Delta}\frac{1}{z}dz=2i\pi.$$
\end{enumerate}
\fincorrection
\correction{007631}
\begin{enumerate}
  \item   Cette application est continue en $i$ : elle ne tend donc pas vers $+\infty$ en module quand $z$ tend vers $i$.
  Elle ne tend donc pas vers $+\infty$ en module quand $|z|$ tend vers $1$.
  \item Par le principe du maximum appliqué à la fonction holomorphe $\Delta\to\Cc$, $z\mapsto 1/f(z)$,
 $$|1/f(0)|\leq \sup_{|z|=r_n}|1/f(z)|=\frac{1}{\inf_{|z|=r_n} |f(z)|}\leq 1/n.$$
 Donc, pour tout $n$, $|f(0)|\geq n$, ce qui est absurde.
 Il n'existe donc pas de telle suite. Donc, dans le cas où $f$ n'a pas de zéro, $|f(z)|$ ne tend pas vers $+\infty$ quand $|z|$ tend vers $1$.
  \item  Comme $|f(z)|$ tend vers $+\infty$ quand $|z|$ tend vers $1$, il existe $r\in ]0,1[$ tel que pour tout $z\in\Delta$ avec $|z|>r$, $|f(z)|>1$. L'application $f$ ne s'annule donc pas sur $\Delta-\overline{\Delta_r}$.
 Par le théorème des zéros isolés, puisque $f$ holomorphe n'est pas constante, elle n'admet qu'un nombre fini de zéro dans le compact $\overline{\Delta_r}$ et par suite dans tout le disque $\Delta$.
 Soit $P(z)$ le polynôme unitaire qui a pour zéros les zéros de $f$ avec les mêmes multiplicités.
 L'application $h=\frac{f}{P}$ a des singularités isolées en les zéros de $f$. Au voisinage de ces zéros, le développement en séries entières de $f$ montre que les singularités 
 sont apparentes et que $h$ tend vers une valeur finie non nulle en ces points. L'application $h$ se prolonge donc par continuité en une application holomorphe $g$ sur $\Delta$. 
 L'égalité $f=hP$ sur $\Delta-f^{-1}(0)$ se prolonge par continuité en $f=gP$ sur $\Delta$.
 Par construction, par le choix des multiplicités, l'application $g$ n'a pas de zéro aux zéros de $f$, et n'a donc pas de zéro sur $\Delta$.
 Comme $|f(z)|$ tend vers $+\infty$ quand $|z|$ tend vers $1$ et que $|P|$ tend vers une valeur finie en tout point de $\partial\Delta$, $|g|$ tend vers $+\infty$ quand $|z|$ tend vers $1$. Le cas précédent montre que l'existence d'une telle application $g$ est absurde.
 Donc, $|f(z)|$ ne tend pas vers $+\infty$ quand $|z|$ tend vers $1$. 
\end{enumerate}
\fincorrection
\nocorrection
\correction{007633}
Une telle application serait majorée en module par le maximum du module de ses valeurs sur le carré compact délimité par les points d'affixes $0,1,1+i,i$.
Elle serait donc entière et bornée donc constante, par le théorème de Liouville. Il n'y a donc pas de telle application.
\fincorrection
\correction{007634}
\begin{enumerate}
  \item Sous cette hypothèse, on peut considérer les applications holomorphes sur $D$, $f/g$ et $g/f$. Elles sont majorées en module par $1$ sur le disque unité, et donc aussi sur le disque fermé par le principe du maximum.
Elles sont donc toutes les deux de module $1$ sur le disque unité, et donc constantes sur le disque unité par le théorème de l'application ouverte. Par connexité de $D$, elles sont constantes sur $D$.
Il existe donc $\lambda\in\Cc$ tel que $f=\lambda g$ sur $D$.
Par égalité des modules, on conclut que $|\lambda|=1$.
  \item Non, les applications $z\mapsto z$ et $z\mapsto z^2$ sont de même module sur le disque unité, mais ne sont pas proportionnelles.
\end{enumerate}
\fincorrection
\correction{007635}
\begin{enumerate}
  \item On note d'abord que les indices de $w_1,w_2,w_3$ par rapport à $\Gamma$ sont respectivement $-1,1$ et $0$.
L'application $\Cc-\{w_1,w_2,w_3\}\to\Cc$, 
$z\mapsto \frac{1}{(z-w_1)(z-w_2)(z-w_3)}$ est méromorphe sur $\Cc$ avec des pôles simples en $w_1,w_2,w_3$ de résidu respectif 
$\lim_{z\to w_1}(z-w_1)\frac{1}{(z-w_1)(z-w_2)(z-w_3)}=\frac{1}{(w_1-w_2)(w_1-w_3)}$, $\frac{1}{(w_2-w_1)(w_2-w_3)}$, $\frac{1}{(w_3-w_1)(w_3-w_2)}$.
Par le théorème de Cauchy,
\begin{eqnarray*}
A&=&2i\pi\left(Ind_\Gamma(w_1)Res_{w_1}(f)+Ind_\Gamma(w_2)Res_{w_2}(f)+Ind_\Gamma(w_2)Res_{w_3}(f)\right)\\
&=&2i\pi\left(-\frac{1}{(w_1-w_2)(w_1-w_3)}+\frac{1}{(w_2-w_1)(w_2-w_3)}+0\right)\\
&=&\frac{2i\pi(w_1+w_2-2w_3)}{(w_2-w_1)(w_1-w_3)(w_2-w_3)}
\end{eqnarray*}
  \item L'application $\Cc\to\Cc$, 
$z\mapsto \sin z$ est holomorphe sur $\Cc$ étoilé : son intégrale sur le chemin fermé $\Gamma$ est donc nulle.
  \item Par le théorème de Cauchy pour les dérivées,
$$C=2i\pi Ind_\Gamma(w_1)\sin'(w_1)=-2i\pi\cos(w_1).$$
\end{enumerate}
\fincorrection
\correction{007636}
\begin{enumerate}
  \item Sur le bord de $\Delta_2$
$$|P(z)-z^4|=|6z+3|\leq 6\times 2+3=15<16=2^4=|z|^4.$$
Par le théorème de Rouché, $P$ de degré $4$ a donc, comme $z\mapsto z^4$ ses quatre racines dans $\Delta_2$.
  \item Sur le bord de $\Delta$
$$|P(z)-(6z+3)|=|z^4|=1< 6-3\leq |6z|-3\leq |6z+3|.$$
Par le théorème de Rouché, le polynôme $P$ n'a, comme $z\mapsto 6z+3$ qu'une racine dans $\Delta$.
  \item Sur le bord de $\Delta_{\frac{1}{3}}$
$$|P(z)-(6z+3)|=|z^4|=\frac{1}{81}<3-6/3=3-|6z|\leq |6z+3|.$$
Par le théorème de Rouché, le polynôme $P$ n'a, comme $z\mapsto 6z+3$ aucune racine dans $\Delta_{\frac{1}{3}}$.
  \item Par la question $2$, $z\mapsto \frac{4z^3+6}{z^4+6z+3}z$ est méromorphe sur $\Delta_{1+\epsilon}$ avec comme seul pôle un point $a$.
Le résidu en $a$ de $P'/P$ est la multiplicité $1$ de $a$ comme zéro de $P$.
Par conséquent, le résidu en $a$ de $\frac{4z^3+6}{z^4+6z+3}z$ est $a$.
La formule demandée résulte donc du théorème des résidus.
\end{enumerate}
\fincorrection
\nocorrection
\nocorrection
\correction{007639}
\begin{enumerate}
  \item On note $f=u+iv$.
\begin{eqnarray*}
 \Delta |f|^2&=&\left(\frac{\partial^2}{\partial x^2}+\frac{\partial^2}{\partial y^2}\right)(u^2(x,y)+v^2(x,y))\\
 &=&2(\left |\frac{\partial (u+iv)}{\partial x}\right |^2 +\left |\frac{\partial (u+iv)}{\partial y}\right |^2)\\&&+2u\left(\frac{\partial}{\partial x}\frac{\partial u}{\partial x}+\frac{\partial}{\partial y}\frac{\partial u}{\partial y}\right)+2v\left(\frac{\partial}{\partial x}\frac{\partial v}{\partial x}+\frac{\partial}{\partial y}\frac{\partial v}{\partial y}\right)
\end{eqnarray*}
Par les identités de Cauchy-Riemann pour la fonction holomorphe $f=u+iv$, et par symétrie de Schwarz, les deux dernières quantités entre parenthèses sont nulles. On trouve donc 
$$\Delta |f|^2=4|f'|^2.$$
  \item On applique la formule précédente à chaque $f_i$ pour obtenir
$$\Delta \sum_{i=1}^N |f_i(z)|^2=4\sum_{i=1}^N|f'_i(z)|^2=0.$$
On en déduit que chaque application holomorphe $f_i$ est de dérivée identiquement nulle et donc constante sur $D$ connexe.
\end{enumerate}


\fincorrection
\correction{007640}
Soit $\varphi (z) =\frac{4z+3}{4+3z}$.
Soit $\zeta\in\partial\Delta$,
$$|\varphi (\zeta)|^2=\frac{16|\zeta|^2+9+12(\zeta+\overline{\zeta})}{16+9|\zeta|^2+12(\zeta+\overline{\zeta})}=1.$$
Par le principe du maximum appliqué à $\varphi$ holomorphe sur le disque $\Delta_{4/3}$ on obtient que pour tout $z\in\Delta$, $|\varphi(z)|\leq 1$.
\fincorrection
\correction{007641}
Les zéros du polynôme $2z^2-5z+2=(z-2)(2z-1)$ et donc les singularités de l'application méromorphe $f(z)=\frac{1}{2z^2-5z+2}$ sont $2$ et $1/2$. Ce sont des pôles simples.
Leurs résidus respectifs sont $res_{2}(f)=\lim_{z\to 2}(z-2)\frac{1}{2z^2-5z+2}=\frac{1}{3}$ et $res_{1/2}(f)=\lim_{z\to 1/2}(z-\frac{1}{2})\frac{1}{2z^2-5z+2}=-\frac{1}{3}$. Leurs indices par rapport au chemin orienté ${\partial\Delta_r}$ sont $0$ ou $1$, suivant qu'ils sont ou pas dans le disque ${\Delta_r}$.
On trouve donc, par le théorème des résidus appliqué à l'application méromorphe $f$ : (les valeurs de $r$ exclues sont $1/2$ et $2$ car les cercles correspondants passent par une singularité.)
\begin{enumerate} 
\item si $0<r<1/2$, $I_r=0$.
\item si $1/2<r<2$, $I_r=-2i\pi/3$.
\item si $2<r$, \ \ \ \ \ \ \ \ $I_r=1/3-1/3=0$.
\end{enumerate}
\fincorrection
\correction{002210}
\begin{enumerate}
\item $VA$ remplace la ligne $i$ par sa somme avec la ligne $j$ multipli\'ee par $\lambda$.

$AV$ remplace la colonne $j$ par sa somme avec la colonne $i$ multipli\'ee par $\lambda$.
\item $V_{ij}(\lambda )V_{kj})=I+\lambda E_{ij} + \lambda ' E_{kj}$
\item Il suffit de montrer que $(I+l_ie_i^T)(I-l_ie_i^T)=I$.
\item  $L(l_i)=V_{i+1,i}(l_{i+1,i}) \cdots V_{n,i}(l_{n,i})$
\item $L^{-1}=L(-l_{n-1}) L(-l_{n-2}) \cdots L(-l_1)\neq I-l_1e_1^T - \cdots -l_{n-1}e_{n-1}^T$
\item (a) algorithme en utilisant l'expression de $L^{-1}$

Pour $i=1$ \`a $n-1$

\hspace{2cm} calcul de $L(-l_i)b$

\hspace{3cm} Pour $j=i+1$ \`a $n$

\hspace{4cm} $b_j\leftarrow b_j-l_{ji}b_i$


(b) algorithme en r\'esolvant le syst\`eme triangulaire

$x_1 =b_1$

Pour $i=2$ \`a $n$ 

\hspace{2cm}$x_i=b_i-\sum_{j=1}^{i-1}l_{ij}x_j$


conclusion: le nombre de calculs et l'espace m\'emoire utilis\'es sont les m\^emes.
\end{enumerate}
\fincorrection
\correction{002211}
Pour d\'emontrer l'\'egalit\'e il suffit de multiplier le membre de droite par $5A+UBV)$ et montrer que l'on obtient l'identit\'e.

Domaine de validit\'e:$ A\in{\cal M}_{n\times n}$ inversible, $U\in {\cal M}_{n\times p}$, $B\in {\cal M}_{p\times q}$, $V\in {\cal M}_{q\times n }$, $I+BVA^{-1}U$ inversible.
\begin{enumerate}
\item  On obtient la formule de Shermann-Morrisson:
$$(A+\beta uv^T)^{-1} = A^{-1} - \frac{\beta}{1+\beta v^TA^{-1}u} A^{-1} uv^T A^{-1}$$
qui permet le calcul de l'inverse d'une matrice qui apparait comme perturbation de rang 1 d'une matrice dont on connait l'inverse.
\item 
$$B\left(\begin{array}{c} X\\y\end{array}\right)=0\Leftrightarrow \left\{ \begin{array}{l} (A+uv^T)x +yu =0\\ v^Tx=0\end{array}\right. \Leftrightarrow \left\{ \begin{array}{l} x=-yA^{-1}u\\v^Tx=-yv^TA^{-1}u=0\end{array}\right.$$
ce qui donne $x=0, y=0$ et donc $B$ est inversible.
\item En appliquant la formule g\'en\'erale on obtient
$$B^{-1}=\left(\begin{array} {cc}
A^{-1} -A^{-1}uv^TA^{-1} & A^{-1}u\\ v^TA^{-1} & 0\end{array}\right)$$
\item En appliquant la m\^eme formule on obtient
$$D^{-1}=\left(\begin{array}{cc}P^{-1} +P^{-1}Q\Delta^{-1} RP^{-1} & -P^{-1}Q\Delta^{-1}\\ -\Delta^{-1} RP^{-1} & \Delta^{-1}\end{array}\right)$$
avec $\Delta = S-RP^{-1}Q$.
\item Calcul r\'ecursif de l'inverse : on dispose de $A_{n-1}^{-1}$ de taille $(n-1)\times (n-1)$ et on veut calculer l'inverse de
$$A_n=\left(\begin{array}{cc} A_{n-1}& v\\
u^T & s\end{array}\right)\mbox{ avec } u,v\in\Rr^{n-1}, s\in\Rr $$
en utilisant la formule pr\'ec\'edente on obtient
$$A_n^{-1}=\left(\begin{array}{cc}
A_{n-1}^{-1} +\frac{1}{\delta} A_{n-1}^{-1} v u^T A_{n-1}^{-1} & - A_{n-1}^{-1}\frac{v}{\delta}\\
-u^TA_{n-1}^{-1}/\delta & \frac{1}{\delta}\end{array}\right)$$
avec $\delta=s-u^TA_{n-1}^{-1}v$.

et on en d\'eduit facilement l'algorithme.
\end{enumerate}
\fincorrection
\correction{002212}
\begin{enumerate}
\item $\|A\|_2^2 =\rho (A^*A)$ rayon spectral de la matrice $A^*A$. D'un autre cot\^e on a:
$$\|A\|^2_F =\mbox{tr} (A^*A) =\sum_{i=1}^n \lambda_i(A^*A) \left\{\begin{array}{l}\geq \rho (A^*A)\\ \leq n\rho (A^*A)\end{array}\right.$$
o\`u tr est la trace de la matrice et $\lambda_i$ ses valeurs propres.
\item $$\|A\|_2 \leq \|A\|_F =\left(\sum_{i,j}|a_{ij}|^2\right)^{1/2} \leq (mn\max_{i,j}|a_{ij}|^2 )^{1/2}=\sqrt{mn} \max_{i,j}|a_{ij}|$$
Soit $x$ tel que : si $\max|a_{ij}| = |a_{i_0j_0}|$ alors on pose $x=e_{j_0}, \|x\|_2=1$. Alors

$\|Ax\|^2_2=\sum_{i=1}^n|a_{i,j_0}|^2 \geq \max |a_{ij}|^2 \Rightarrow \mbox{sup} \|Ax\|^2_2 \geq \max |a_{ij}|^2$
\item On rappelle que $\|A\|_{\infty} =\sum_{j=1}^n |a_{i_0j}|$ pour un certain $i_0$. Alors

$ \|A\|_2^2\leq\|A\|_F^2=\sum_{i=1}^m\sum_{j=1}^n |a_{ij}|^2 \leq m\times \max_i\sum_{j=1}^n |a_{ij}|^2 \leq m \max \left(\sum_{j=1}^n|a_{ij}|\right)^2 =m\|A\|_\infty$

Choisissons maintenant $x=(x_i)$ avec $x_i=$signe$(a_{i_0i})$. Alors
$$\sum_{j=1}^n a_{i_0j}x_j = \sum_{j=1}^n|a_{i_0j}|=\|A\|_{\infty}$$
$$\|x\|_2=\sqrt{n} \Rightarrow \|Ax\|^2_2 =\sum_{i=1}^m\left(\sum_{j=1}^n a_{ij}x_j\right)^2 \geq \|A\|_{\infty}^2\Rightarrow \frac{\|Ax\|_2}{\|x\|_2}\geq \frac{\|A\|_{\infty}}{\sqrt{n}}$$
ce qui implique $\|A\|_2 \geq \|A\|_{\infty} /\sqrt{n}$
\item M\^eme d\'emonstration que pr\'ec\'edemment ou alors constater que $\|A\|_1=\|A^T\|_{\infty}$.

\item $\|E\|^2_F =\sum_{i,j} u_i^2v_j^2\sum_{i=1}^m\sum_{j=1}^n u_i^2v_j^2=\|u\|^2_2\|v\|^2_2$

$$\|E\|_{\infty }= \max_i\left(\sum_{j=1}^n|u_iv_j|\right) =\max_i\left(|u_i|\sum_{j=1}^n|v_j|\right)=\|v\|_1\|u\|_{\infty}$$
$\|Ex\|_2^2=\sum_{i=1}^m\left(u_i\sum_{j=1}^n v_jx_j\right)^2 =\sum_{i=1}^mu_i^2 \times (x,x)^2=\|u\|_2^2 (x,v)^2$

$$\frac{\|Ex\|_2}{\|x\|_2}=\frac{(x,v)}{\|x\|_2} \|u\|_2 \Rightarrow \mbox{sup}_x\frac{\|Ex\|_2}{\|x\|_2}=\|v\|_2\|u\|_2$$
\end{enumerate}

\fincorrection
\correction{002213}
$\rho (A)<1 \Rightarrow 1$ n'est pas valeur propre de $A \Rightarrow 0$ n'est pas valeur propre de $I-A\Rightarrow I-A$ inversible
$$(I-A)C_k= (I-A)(I+A+\cdots +A^k)=I-A^{k+1}$$
$C_k=(I-A)^{-1} (I-A^{k+1}) \Rightarrow (I-A)^{-1} -C_k =(I-A)^{-1} A^{k+1}$ et conc
$$\|(I-A)^{-1} -C_k\|\leq \|(I-A)^{-1}\|\|A^{k+1}\|\leq \|(I-A)^{-1}\|\|A\|^{k+1}$$
Comme $\|A\|<1$ pour au moins une norme subordonn\'ee on obtient finalement
$$\lim_{k\rightarrow \infty}\|(I-A)^{-1}-C_k\|=0$$
\fincorrection
\correction{002214}
$AB=I-X\Rightarrow B^{-1}A^{-1} = (I-X)^{-1}\Rightarrow A^{-1}=B(I-X)^{-1}=B(I+X+X^2+\cdots )$

$$\|A^{-1}-B\|\leq \|BX\|\|I+X+\cdots \|\leq \|BX\|(1+\|X\|
+\|X\|^2+ \cdots )\leq \frac{\|BX\|}{1-\|X\|} .$$
pour $\|X\|<1$ 
\fincorrection
\nocorrection
\correction{002216}
4.On calcule
$$A^TA=\left(\begin{array}{cc}6&2\\2& 3\end{array}\right)$$ et ses valeurs propres
$$\mbox{det} \left(\begin{array}{cc}6-\lambda &2\\2& 3-\lambda\end{array}\right)=0\Leftrightarrow \lambda_1=7=\mu_1^2, \lambda_2=2=\mu_2^2$$
On calcule ensuite les vecteurs propres associ\'es \`a ces valeurs propres
$$\left(\begin{array}{cc}6&2\\2& 3\end{array}\right)\left(\begin{array}{c}x_1\\x_2\end{array}\right) =\lambda \left(\begin{array}{c}x_1\\x_2\end{array}\right)$$
et la matrice $V$ est la matrice dont les colonnes sont
$$v_1=(2/\sqrt{5}, 1/\sqrt{5})^T,\quad v_2=(1/\sqrt{5}, -2/\sqrt{5})^T$$
les colonnes de $U$ sont alors donn\'ees par
$$u_1=Av_1/\mu_1=1/(\sqrt{7}\sqrt{5})(3,5,-1)^T ,\quad u_2=Av_2/\mu_2=1/(\sqrt{2}\sqrt{5})(-1,0,-3)^T$$
quant \`a $u_3$ il est choisi orthogonal \`a $u_1$ et $u_2$ et de norme 1.
\fincorrection
\correction{002217}
\begin{enumerate}
\item $\Sigma^{\dagger}\Sigma e_i=e_i, i=1,\cdots, r$ c'est l'application identit\'e
\item $AA^{\dagger} =U\Sigma V^*V\Sigma^{\dagger} U =U\Sigma \Sigma^{\dagger}U^*=I$

On a donc obtenu une g\'en\'eralisation de l'inverse.
\item $U^*\sum_{i=1}^m \epsilon_iu_i^*$ avec$ \{\epsilon_1, \cdots , \epsilon_m\}$ base canonique de $\Rr^m$. Comme $\Sigma^{\dagger} \epsilon_i=0$ pour $r+1\leq i\leq m$ on a
$$\Sigma^{\dagger}U^*=\sum_{i=1}^r \mu_i^{-1}e_iu_i^* \Rightarrow A^{\dagger} =V\Sigma^{\dagger}U^* =\sum_{i=1}^r \mu_i^{-1}(Ve_i)u_i^* =\sum_{i=1}^r \mu_i^{-1}v_iu_i^*$$
\item On a
 $$AA^{\dagger} =\sum_{i=1}^r \mu_iu_iv_i^* \sum_{j=1}^r\mu_j^{-1}v_ju_j^*=\sum_{j=1}^r\mu_j\mu_j^{-1}u_ju_j^*$$
Comme Im($A$) =span$\{u_1,\cdots , u_r\}$ le r\'esultat suit.
\item soit $y\in$ Im$A^*$ $\Leftrightarrow u=\sum_{i=1}^r x_iv_i$. Alors
$$A^*A=V\Sigma^* U^*U\Sigma V^* =V\Sigma^*\Sigma V^*=\sum_{i=1}^r \mu_i^2 v_iv_i^*\Rightarrow A^*Ay=\sum_{i=1}^r \mu_i^2x_iv_i$$
et finalement
$$(\sum_{i=1}^r\mu_i^{-2}v_iv_i^*) (A^*Ay) =\sum_{i=1}^rx_iv_i$$
\end{enumerate}
\fincorrection
\correction{002218}
\begin{enumerate}
\item $\|A\|_2=\|U\Sigma V^*\|_2=\|\Sigma\|_2=\max |\sigma_j|=\sigma_1$
\item $\|A\|^2_F =$tr$(A^*A)=$tr$(U^*A^*AU)=\|AU\|^2_F=$tr$(A^*U^*UA)=\|UA\|^2_F$ et donc
$$\|A\|_F=\|U\Sigma V^*\|_F =\|\Sigma\|_F=\sqrt{\sigma_1^2 + \cdots +\sigma_r^2}$$
\item $A^*A=(V\Sigma^*U^*) (U\Sigma V^*)=V(\Sigma^*\Sigma )V^*$ et donc $A^*A$ est semblable \`a $\Sigma^*\Sigma$, les deux matrices ont donc les m\^emes valeurs propres. Les valeurs propres de $\Sigma^*\Sigma$ sont $\sigma_1^2, \cdots , \sigma_r^2$ plus $n-r$ valeurs propres nulles si $n>r$.
\item $|\det A|=|\det (U\Sigma V^*)|= |\det U|\det|\Sigma||\det V^*|= |\det \Sigma|=\prod_{i=1}^r \sigma_i$
\item Une matrice hermitienne \'etant diagonalisable a une base orthonormale de vecteurs propres
$$A=Q\Lambda Q^* = Q |\Lambda|\mbox{sign} (\Lambda)Q^*$$
or $U=$sign$(\Lambda )Q^*$ est une matrice unitaire: $U^*U=Q$sign$(\Lambda) $sign$(\Lambda) Q^*= QQ^*=I$. Donc $Q |\Lambda|U$ est une d\'ecomposition en valeurs singuli\`eres de $A$, les valeurs singuli\`eres \'etant $|\lambda_1|,\cdots , |\lambda_n|$.
\end{enumerate}
\fincorrection
\correction{002219}
\begin{enumerate}
\item $\|A\|_2^2 =\rho (A^*A)=\max_i \lambda_i(A^*A) =\mu_1^2(A)$ la plus grande valeur singuli\`ere de $A$

$\|A^{-1}\|_2^2=\rho (A^{-1}(A^{-1})^*)=\max_i \lambda_i((A^*A)^{-1}=\frac{1}{\mu_n(A)^2}$ avec $\mu_n(A)$ la plus petite valeur singuli\`ere de $A$. Donc
$$\mbox{cond}_2(A)=\|A\|_2 \|A^{-1}\|_2=\mu_n(A)/\mu_1(A)$$
\item Si $A$ est normale alors $\|A\|_2=\rho (A) $ rayon spectral. Donc
$$A^{-1}=UD^{-1}U^* \Rightarrow (A^{-1})^* A^{-1} =U(D^{-1})^* D^{-1}U^*\Rightarrow \rho ((A^{-1})^*A^{-1})=1/\min_i |\lambda_i(A)|^2$$
$$\mbox{cond}_2(A)=\max |\lambda_i(A)|/\min |\lambda_i(A)|$$
\item cond$_2(QA)=\|QA\|_2 \|A^{-1}Q^*\|_2=\|A\|_2\|A^{-1}\|_2=$cond$_2(A)$.
\end{enumerate}
\fincorrection
\nocorrection
\correction{002221}
$B=A+\delta A= A(I+A^{-1}\delta A)$ matrice inversible si $\|A^{-1}\delta A\|<1$

$B^{-1}- A^{-1} =A^{-1} (A-B)B^{-1} \Rightarrow \|B^{-1}- A^{-1}\|\leq \|A^{-1}\|\|A-B\|\|B^{-1}\|\Rightarrow$
$$\frac{\|B^{-1}- A^{-1}\|}{\|B^{-1}\|}\leq \|A^{-1}\|\|A-B\|=
\|A^{-1}\|\|\delta A\| = \mbox{cond}(A)\frac{\|\delta A\|}{\|A\|}$$
\fincorrection
\correction{002222}
\begin{enumerate}
\item A la k-\`eme \'etape de l'\'elimination de Gauss, l'\'el\'ement $a_{ij}^{k+1}$ est donn\'e par
$$a_{ij}^{k+1}=a_{ij}^k-\frac{a_{kj}^ka_{ik}^k}{a_{kk}^k} \quad k+1\leq i,j\leq n$$
et on remarque imm\'ediatement par r\'ecurrence que toutes les matrices $\tilde{A_k}$ sont sym\'etriques. On a 

$(\tilde{A}_{k+1} v',v')=\sum_{i=k+1}^n v_i (\sum_{j=k+1}^n a_{ij}^{(k)}v_j) - \frac{1}{a_{kk}^k}(\sum_{i=k+1}^n a_{ik}^k v_i)^2$

$(\tilde{A}_k v,v)= \sum_{i=k+1}^nv_i(\sum_{j=k+1}^n a_{ij}^kv_j) +\sum_{i=k+1}^n (a_{ik}^k +a_{ki}^k )v_iv_k +a_{kk}^kv_k^2$

Par sym\'etrie $a_{ik}^k=a_{ki}^k$ et donc

$(\tilde{A}_kv,v)=(\tilde{A}_{k+1}v',v') +\frac{1}{a_{kk}^k} [(\sum_{i=k+1}^n a_{ik}^k v_i)^2 + 2v_k\sum_{i=k+1}^n a_{ik}^kv_i a_{kk}^k +(a_{kk}^k)^2v_k^2]=$

$$(\tilde{A}_{k+1}v',v') +\frac{1}{a_{kk}^k} [a_kk^kv_k +\sum_{i=k+1}^na_{ik}^kv_i]^2$$
\item Faisons un raisonnement par r\'ecurrence
\begin{itemize}
\item $\tilde{A}_1$ est sym\'etrique d\'efinie positive;
\item Par hypoth\`ese supposons que $\tilde{A}_k$ est d\'efinie positive;
\item Supposons par absurde que $\tilde{A}_{k+1}$ ne soit pas d\'efinie positive: alors $\exists v'\neq 0:\quad (\tilde{A}_{k+1}v',v')\leq 0$. On d\'efinit le vecteur $v\in\Rr^{n-k+1}$ par:
\begin{itemize}
\item $v_i=v'_i,\quad k+1\leq i\leq n$
\item $v_k$ est solution de $a_{kk}^k +\sum_{i=k+1}^n a_{ik}^k v_i=0$
\end{itemize}
Alors $(\tilde{A}_kv,v)=0$ et $v\neq 0$; donc $\tilde{A}_k$ n'est pas d\'efinie positive, ce qui contredit l'hypoth\`ese de r\'ecurrence.
\end{itemize}
\item Premi\`ere in\'egalit\'e: en utilisant la relation d'\'elimination on obtient: $a_{ii}^{k+1}=a_{ii}^k -\frac{\left|a_{ki}^k\right|^2}{a_{kk}^2}$
\begin{itemize}
\item une matrice d\'efinie positive a tous ses \'el\'ements diagonaux strictement positifs, donc $a_{ii}^{k+1}>0$
\item $\left|a_{ki}^k\right|^2/\left|a_{kk}^k\right|^2 \geq 0, \quad k+1\leq i\leq n$
\end{itemize}
donc $a_{ii}^{k+1}\leq a_{ii}^k, k+1\geq i$

Deuxi\`eme in\'egalit\'e: supposons qu'il existe un \'el\'ement $a_{ij}^k, i<J$ tel que $\left|a_{ij}^k\right| \geq \max_{k\leq l \leq n} a_{ll}^k$. On consid\`ere le vecteur $v\neq 0$ d\'efini par
$$v_i=1, v_j=-\mbox{sign} (a_{ij}^k), v_l=0\quad l\neq i,j$$
Alors
$$(\tilde{A}_k v,v) =( a_{ii}^k -\left|a_{ij}^k\right|)-(\left|a_{ij}^k\right| -a_{jj}^k) \leq 0$$
ce qui est impossible. Donc
$$\max_{1\leq i,j\leq n}\left|a_{ij}^{k}\right|=\max_{1\leq i \leq n}\left|a_{ii}^k\right|$$
\end{enumerate}
\fincorrection
\nocorrection
\correction{002224}
Montrons par r\'ecurrence que $A_n=U$ est une matrice bande.

$A_1=A, \quad A_{k+1}=L_kA_k= L_kL_{k-1} \cdots L_1A, \quad k=1,\cdots , n-1.$

Supposons que $A_k$ est une matrice bande i.e., $a_{ij}^k=0$ pour $|i-j|\geq p$ et montrons que $A_{k+1}$ est une matrice bande.
$$a_{ij}^{k+1}=a_{ij}^k-\frac{a_{ik}^k a_{kj}^k}{a_{kk}^k}$$
Soit $|i-j|\geq p\Leftrightarrow |(i-k)-(j-k)|\geq p$. On consid\`ere deux cas:
\begin{itemize}
\item $k+1\leq i\leq n$ et $k\leq j\leq n$. Alors $i-k\geq p$ ou $j-k\geq p$ $\Rightarrow a_{ik}^ka_{kj}^k=0\Rightarrow a_{ij}^{k+1}=a_{ij}^k=0$
\item $i\leq k$ ou $j\leq k-1$ alors $ a_{ij}^{k+1}=a_{ij}^k=0$
\end{itemize}
donc $A_{k+1}$ est une matrice bande et $U$ est une matrice bande. On a $A=LU$ et la matrice triangulaire inf\'erieure $L$ a pour \'el\'ements $l_{ij}=a_{ij}^j/a_{jj}^j,\quad j\leq i\leq n$. Toutes les matrices $ A_j$ \'etant des matrices bandes on a $a_{ij}^j=0$ pour $i-j\geq p\Rightarrow l_{ij}=0$ pour $i-j\geq p$.
\fincorrection
\correction{002225}
Soit $LU$ la factorisation $LU$ de $A$. On va intercaler dans cette factorisation la matrice r\'eelle $\Lambda=$diag$(\sqrt{|u_{ii}|})$.


$A=(L\Lambda ) (\Lambda^{-1}U) =BC$. La sym\'etrie de $A$ entraine $BC=C^TB^T. $ On a 

$C(B^T)^{-1}$ matrice triangulaire sup\'erieure, $B^{-1}C^T$ matrice triangulaire inf\'erieure et $C(B^T)^{-1}=B^{-1}C^T$ et donc

$C(B^T)^{-1}=B^{-1}C^=$diag(sign$(u_{ii})=S$ $\Rightarrow C(B^T)^{-1}S^{-1}=I=S^{-1}B^{-1}C^T \Leftrightarrow C^T=BS=\tilde{B}$. Donc $A$ peut \^etre mise sous la forme
$$A=B\tilde{B}^T \mbox{ avec } \tilde{B}=BS$$
i.e. la $i$-\`eme colonne de $\tilde{B}$ est \'egale \`a la $i$-\`eme colonne de $B$ affect\'ee du signe de $u_{ii}$

Application num\'erique:
$$\tilde{B}=\left(\begin{array}{cccc}
1&2&1&1\\
&-1&2&1\\
&&-1&-1\\&&&1\end{array}\right).$$
\fincorrection
\nocorrection
\nocorrection
\correction{002228}
$$A=A_1=\left(\begin{array}{cc}
\alpha & u^T\\v&B_1\end{array}\right),\quad B_1=(b_{ij})_{i,j=1}^{n-1}$$
$A^T$ \'etant \`a diagonale strictement dominante on a:
$$|\alpha|>\sum_{i=1}^{n-1}|v_i|,\quad |u_i| +\sum_{j\neq i}|b_{ji}|<|b_{ii}|$$
Il suffit de montrer que
\begin{itemize}
\item la premi\`ere colonne de $L$ v\'erifie $|l_{11}|>\sum_{i\neq 1}|l_{i1}|$
\item $B_2$ est telle que
$$A_2=\left(\begin{array}{cc}\alpha & u^T\\0& B_2\end{array}\right), \quad C=B_2=B_1-\frac{1}{\alpha} vu^T$$
v\'erifie $|c_{ii}|>\sum_{j\neq i}|c_{ji}|$ avec $C_{ij}=B_{ij}-\frac{1}{\alpha}v_iu_j$ et it\'erer.
\end{itemize}
\begin{itemize}
\item premi\`ere colonne de $L$: $l_{i1}=v_i/\alpha \Rightarrow \sum_{i=2}^n |l_{i1}|=\sum_{i=1}^{n-1}\frac{|v_i|}{\alpha}<1$
\item $\sum_{i\neq j}|c_{ij}| =\sum_{i\neq j}\left|b_{ij}-\frac{1}{\alpha}v_iw_j\right|\leq \sum_{i\neq j}|b_{ij}| +\frac{1}{|\alpha|}|w_j|\sum_{i\neq j}|v_i|$
$$\leq |b_{jj}| - |u_j| + \frac{1}{|\alpha|}|u_j| (|\alpha|-|v_j|)
\leq\left|b_{jj}-\frac{1}{\alpha} u_jv_j\right|=|c_{jj}|$$
donc $B_2^T$ est de diagonale strictement dominante. La d\'emonstration se finit par r\'ecurrence.
\end{itemize}
\fincorrection
\correction{002229}
\begin{enumerate}
\item
Soit $P$ l'op\'erateur de projection dans le sous-espace $U$ de dimension 1 g\'en\'er\'e par $v$. Alors $Q=I-P$ est l'op\'erateur de projection sur l'hyperplan $U^{\bot}$ orthogonal \`a $U$. On a d\'ej\`a vu que $Pw=vv^Tw\quad \forall w$, et donc $Qw=w-vv^Tw$. On obtient

$P(H(v)w) =P(w_(2v^Tw)v)=(v^Tw)v-2v^Twvv^Tv=-(v^Tw)v=-Pw$


$Q(H(v)w)= H(v)w -P(H(v)w)=w-2vv^Tw+v^Twv=w-v^Twv=Qw$.

La matrice $H(v)$ repr\'esente donc une sym\'etrie par rapport \`a l'hyperplan $U^{\bot}$. On conclut que les vecteurs de $U^{\bot}$ sont invariants par $H(v)$.

$V(v)w=w\quad \forall w\in U^{\bot}, \quad \mbox{dim} U^{\bot}=n-1\Rightarrow \lambda =1$ est valeur propre de $H(v)$ avec multiplicit\'e $n-1$.

$H(v)v=-v\mp \lambda =-1$ est valeur propre de multiplicit\'e 1. Donc
 $$\det H(v)=\prod_{i=1}^n \lambda_i(H(v))=-1$$
\item On sait qu'il exite des matrices de Householder $H_1, H_1, \ldots ,H_{n-1}$ telles que $H_{n-1}\cdots H_1A=A_{n}$ matrice triangulaire sup\'erieure. Comme $A$ est orthogonale on conclut que $A_n$ est orthogonale. Mais une matrice triangulaire sup\'erieure orthogonale est forc\'ement diagonale $\Rightarrow A_n=$diag$(\pm 1)$. On peut s'arranger pour que  $(A_n)_{ii} >0 \quad i=1,\ldots, n-1.$ Donc soit $A_n=I$ soit $A_n=$diag$(1,1,\ldots , 1,-1)=H(e_n)$ et finalement la matrice orthogonale $A$ s'\'ecrit
$$A=H_1 \cdots H_{n-1}H(e_n)$$
\end{enumerate}
\fincorrection
\correction{002230}
\begin{enumerate}
\item Pour $k=1,\cdots ,n\quad a_k=\sum_{i=1}^kr_{ik}q_i$ avec $r_{ik}=q_i^Ta_k$ par orthonormalit\'e des $q_i$.
\item D\'ecoule imm\'ediatement de la question pr\'ec\'edente.
\item Algorithme de Gram-Schmidt:

Pour $k=1, \cdots , n$ faire

\hspace{1cm} $r_{ik}=q_i^Ta_k\quad $ pour $i=1, \cdots , k-1$

\hspace{1cm} $z_k=a_k-\sum_{i=1}^{k-1} r_{ik}q_i$

\hspace{1cm} $r_{kk}=(z_k^Tz_k)^{1/2}$

\hspace{1cm}$q_k=z_k/r_{kk}$
\item \begin{enumerate}
\item $$\sum_{i=k}^nq_ir_i^T =[q_k \cdots q_n]\left(\begin{array}{c}r_k^T\\ \vdots \\ r_n^T\end{array}\right) = q_k \cdots q_n]\left(\begin{array}{ccccccc} 
0&\cdots &0 &r_{kk} & \cdots &\cdots & r_{kn}\\
0&\cdots &0 & 0 & r_{k+1,k+1}&\cdots & r_{k+1,n}\\
\vdots & \vdots &\vdots &\vdots &\vdots &\vdots &\vdots \\
0&\cdots & \cdots&&&0&r_{nn}\end{array}\right)$$
$$A^{(k)}e_k = z= [q_k \cdots q_n]\left(\begin{array}{c}r_{kk}\\ 0\\ \vdots \\ 0\end{array}\right)=r_{kk}q_k \Rightarrow r_{kk}=\|z\|_2, q_k=z/r_{kk}$$
\item $$q_k^TA^{(k)} =[q_k^Tz, q_k^TB]=[1 , 0, \cdots , 0]\left(\begin{array}{c}r_k^T\\ \vdots \\r_n^T\end{array}\right)=r_k^T$$
et donc
$$[r_{k,k+1}, \cdots , r_{kn}]=q_k^TB$$
\item $$[0,\cdots , 0, A^{(k+1)}] = \sum_{i=k+1}^n q_ir_i^T= [0,\cdots , 0, A^{(k)}] - q_kr_k^T= [0,\cdots , 0,A^{(k)} - q_k(r_{kk}, \cdots , r_{kn})]$$
$$[ 0,\cdots , 0, z-q_kr_{kk}, B-q_k(r_{k,k+1}, \cdots , r_{kn})] \Rightarrow A^{(k+1)}=B-q_k (r_{k,k+1}, \cdots , r_{kn})$$
\item Donn\'ees: $A\in\Rr^{m\times n}, $ rank$(A)=n$

On calcule la factorisation $A=Q_1R_1$, $Q_1\in\Rr^{m\times n}$ orthonormale, $R_1\in\Rr^{n\times n}$ triangulaire sup\'erieure. Le calcul de $Q_1$ se fait sur place.

Pour $k=1, \cdots , n$

\hspace{1cm} $r_{kk}=\left(\sum_{i=1}^m a_{ik}^2\right)^{1/2} $

\hspace{1cm}pour $i=1, \cdots  ,m$

\hspace{1.5cm}$a_{ik}\leftarrow a_{ik}/r_{kk}$

\hspace{1cm} pour $j=k+1, \cdots , n$

\hspace{1.5cm}$r_{kj}\leftarrow\sum_{i=1}^m a_{ik}a_{ij}$

\hspace{1.5cm} pour $i=1, \cdots , m$

\hspace{2cm}  $a_{ij}\leftarrow a_{ij}-a_{ik}r_{kj}$

\item complexit\'e: $mn^2$ flops.

\end{enumerate}
\end{enumerate}
\fincorrection
\correction{002231}
\begin{enumerate}
\item $G_{p,q}(c,s)= I +(c-1)e_pe_p^T + s e_qe_p^T -s e_qe_q^T + (c-1)e_pe_q^T$ avec $e_i$ les vecteurs de la base canonique.
\item On montre que $e_i^TG^TGe_j=\delta_{ij} \quad \forall i,j=1, \cdots ,n$ et donc $G^TG=I$ ce qui permet de conclure que $G$ est inversible d'inverse $G^T$ et donc orthogonale.
\item $e_i^TGA =e_i^TA= a_i^T$ pour $i\neq p,q$

$e_p^TGA= ca_p^T-sa_q^T$, $e_q^TGA= sa_p^T+ca_q^T$, et donc $G$ change seulement les lignes $p$ et $q$
\item On pose $\alpha =a_{pj}$ et $\beta = a_{qj}$ . On a donc \`a r\'esoudre dans le premier cas  le syst\`eme
$$\left\{ \begin{array}{l} c\alpha -s\beta =0\\c^2+s^2=1\end{array}\right. \Leftrightarrow 
\left\{ \begin{array}{l}c=\pm \beta /\sqrt{\alpha^2+\beta^2}\\s=\pm \alpha /\sqrt{\alpha^2+\beta^2}\end{array}\right.$$
ce qui nous donne deux matrices $G$. Pour le deuxi\`eme cas et en proc\'edant de la  m\^eme fa\c con on obtient
$$\left\{ \begin{array}{l}c=\pm \alpha /\sqrt{\alpha^2+\beta^2}\\s=\mp \beta /\sqrt{\alpha^2+\beta^2}\end{array}\right.$$
\end{enumerate}
\fincorrection
\nocorrection
\nocorrection
\correction{002234}
M\'ethode de Givens rapide
\begin{enumerate}
\item $MM^T=$diag$(d_1, \cdots , d_m)=\Delta^2$ avec $\Delta=$diag$ (\sqrt{d_1}, \cdots, \sqrt{d_m})$

$\Delta^{-1}MM^T\Delta^{-1}=(\Delta^{-1}M)(\Delta^{-1}M)^T=I\Rightarrow \Delta^{-1}M$ est une matrice orthogonale

 $A=M^{-1}S=(M^{-1}Delta \Delta^{-1} S=(\Delta^{-1}M)^{-1} (\Delta^{-1}S)=(\Delta^{-1}M)^{T} (\Delta^{-1}S)=(M^T\Delta^{-1})(\Delta^{-1}S)$

Comme $\Delta^{-1}S$ est triangulaire sup\'erieure ona $ A=QR$ avec $Q=M^T\Delta^{-1}, R=\Delta^{-1}S$
\item 
\begin{enumerate}
\item$$M_1x=\left(\begin{array}{c}\beta_1x_1 +x_2\\x_1+\alpha_1x_2\end{array}\right)  ,\quad M_1DM_1^T=\left(\begin{array}{cc}d_2+\beta_1^2 d_1 & d_1\beta_1 +d_2\alpha_1\\ d_1\beta_1 +d_2\alpha_1 & d_1+\alpha_1^2d_2\end{array}\right)$$
Alors
\begin{itemize} 
\item $x_1+\alpha_1x_2=0\Leftrightarrow \alpha_1=-x_1/x_2$
\item $d_1\beta_1 +d_2 (-x_1/x_2)=0\Leftrightarrow \beta_1=-\alpha_1 d_2/d_1=x_1d_2/(x_2d_1)$
\end{itemize}
 Pour le choix pr\'ec\'edent on veut d\'eterminer $\gamma_1$ tel que

$x_2(1+\gamma_1)=\beta_1x_1 +x_2 =x_2(\beta_1x_1/x_2 +1) \Rightarrow \gamma_1= (d_2/d_1)(x_1/x_2)^2$ c'est-\`a-dire
$$\gamma_1=-\alpha_1\beta_1$$
pour cette valeur on a

$d_2+\beta_1^2d_1 =d_2(1+\alpha_1^2d_2/d_1) = d_2(1+\gamma_1)$
$d_1+\alpha_1^2d_2 =d_1(1+\alpha_1^2d_2/d_1) = d_2(1+\gamma_1)$
\item le m\^eme type de calcul nous donne
$$\beta_2=-x_2/x_1, \quad \alpha_2=-(d_1/d_2)\beta_2,\quad \gamma_2=-\alpha_2\beta_2 = (d_1/d_2)(x_2/x_1)^2$$
\item on remarque que $\gamma_1 \gamma_2=1$ et donc soit $\gamma_1\leq 1$ , soit $\gamma_2\leq 1$

\end{enumerate}
\item $$\left(\begin{array}{cc} m_{pp} & m_{pq}\\m_{qp} & m_{qq}\end{array}\right)=\left(\begin{array}{cc} \beta_1 & 1\\1 & \alpha_1\end{array}\right)\mbox{ ou }\left(\begin{array}{cc} 1 & \alpha_2 \\\beta_2 & 1\end{array}\right)$$
avec les $\alpha_i, \beta_i$ d\'efinis comme pr\'ec\'edemment.
\item \emph{algorithme}

$d_i=1$ pour $i=1, \cdots , m$

Pour $p=1, \cdots , \min\{n,m-1\}$

\hspace{1cm} Pour $q=p+1, \cdots , m$

\hspace{1.5cm}si $a_{qp}\neq 0$ alors

\hspace{2cm} $\alpha= -a_{pp}/a_{qp},\quad \beta=-\alpha d_q/d_p ,\quad \gamma =-\alpha\beta$

\hspace{2cm}si $\gamma \leq 1$ alors

\hspace{2.5cm} $\left(\begin{array}{ccc}a_{pp}& \cdots & a_{pn}\\a_{qp} & \cdots & a_{qn}\end{array}\right) \leftarrow \left(\begin{array}{cc}\beta & 1 \\ 1 & \alpha \end{array}\right)\left(\begin{array}{ccc}a_{pp}& \cdots & a_{pn}\\a_{qp} & \cdots & a_{qn}\end{array}\right)$

\hspace{2.5cm}\'echanger $d_p$ et $d_q$

\hspace{3cm} $d_p\leftarrow (1+\gamma d_p)$

\hspace{3cm} $d_q\leftarrow (1+\gamma d_q)$

\hspace{2cm}sinon

\hspace{2.5cm} \'echanger $\alpha$ et $\beta$

\hspace{2.5cm}$\alpha=1/\alpha$ , $\beta = 1/\beta $ , $\gamma = 1/\gamma $

\hspace{2.5cm} $\left(\begin{array}{ccc}a_{pp}& \cdots & a_{pn}\\a_{qp} & \cdots & a_{qn}\end{array}\right) \leftarrow \left(\begin{array}{cc}1 & \alpha \\ \beta & 1 \end{array}\right)\left(\begin{array}{ccc}a_{pp}& \cdots & a_{pn}\\a_{qp} & \cdots & a_{qn}\end{array}\right)$

\hspace{2.5cm} $d_p\leftarrow (1+\gamma d_p)$

\hspace{2.5cm} $d_q\leftarrow (1+\gamma d_q)$

le co\^ut de cet algorithme est de $n^2(m-n/3)$ flops.

\item \begin{enumerate}
\item on a $MA=R=\left(\begin{array}{c} S_1\\ 0\end{array}\right)$ avec $S_1$ triangulaire sup\'erieure et $MM^T=D=$diag$(d_1,\cdots , d_n)$. Donc la matrice $D^{-1/2}M$ est une matrice orthogonale

$$\|Ax-b\|_2^2=\|D^{-1/2}MAx-D^{-1/2}Mb\|_2^2=\left\|D^{-1/2}\left[\left(\begin{array}{c}S_1\\0\end{array}\right) x-Mb\right]\right\|_2^2=$$
$$=\left\|D^{-1/2}\left[\left(\begin{array}{c}S_1\\0\end{array}\right) x-\left(\begin{array}{c}c\\d\end{array}\right)\right]\right\|_2^2=\left\|D^{-1/2}\left(\begin{array}{c} S_1x-c\\d\end{array}\right)\right\|_2^2$$
La solution est obtenue en r\'esolvant le syst\`eme triangulaire sup\'erieure $S_1x=c$ de taille $n\times n$.
\item \begin{itemize}
\item mise \`a jour de $b$ pour le calcul de $Mb$ en m\^eme temps que la mise \`a jour de $A$

pour $p=1, \cdots ,\min \{n,m-1\}$

\hspace{1cm} pour $q=p+1,\cdots , m$ faire 

\hspace{1.5cm}$\left(\begin{array}{c}b_p\\b_q\end{array}\right)\leftarrow \left(\begin{array}{cc}\beta & 1\\ 1 & \alpha\end{array}\right)\left(\begin{array}{c}b_p\\b_q\end{array}\right)$ ou

\hspace{1.5cm}$\left(\begin{array}{c}b_p\\b_q\end{array}\right)\leftarrow \left(\begin{array}{cc}1 & \alpha\\ \beta & 1\end{array}\right)\left(\begin{array}{c}b_p\\b_q\end{array}\right)$ 

\item r\'esolution du syst\`eme triangulaire sup. $S_1x=c$

$x_n\leftarrow b_n/a_{nn}$

\hspace{1cm} Pour $i=n-1, \cdots , 1$ faire

\hspace{1.5cm} $x_i\leftarrow (b_i-\sum_{j=i+1}^na_{ij}x_j)$
\end{itemize}
\end{enumerate}
\item Application num\'erique:

$$M=\frac{1}{24}\left(\begin{array}{ccc}8&16&24\\
40 & 10 & -20\\
15 & -30 & 15\end{array}\right),\quad D=\mbox{diag} (14/9,175/48,75/32)$$

$$M[A,b]=\left(\begin{array}{ccc}14/3 & 32/3 & 50/3\\
0&15/4 & 15/2\\
0&0&0\end{array}\right), \quad x_{ls}=(-1,2)^T$$
\item on a
$$MD^{-2}M^T=\tilde{D} \Leftrightarrow (\tilde{D}^{-1/2}MD^{-1})(\tilde{D}^{-1/2}MD^{-1})^T=I$$
donc
 $(\tilde{D}^{-1/2}MD^{-1})$ est une matrice orthogonale et on obtient
$$\|D(Ax-b)\|_2=\|\tilde{D}^{-1/2}MD^{-1}D(Ax-b)\|_2=\|\tilde{D}^{-1/2}(MAx-Mb)\|_2=$$
$$=\left|\tilde{D}^{-1/2}\left(\begin{array}{c}Sx-c\\e\end{array}\right)\right|_2$$
Donc le min est atteint pour $Sx=c$ avec $Mb=(C,e)^T$

La modification dans l'algorithme pr\'ec\'edent consiste \`a initialiser la matrice diagonale $D$ avec $D^{-2}$ (au lieu de l'identit\'e).
\end{enumerate}
\fincorrection
\nocorrection
\correction{002236}
\begin{enumerate}
\item Pour le membre de gauche on obtient

$$(x,Ax)-(y,Ay)=(x,AM^{-1}Mx)+(M^{-1}Ax,Ax) -(M^{-1}Ax,AM^{-1}Ax)$$

Pour le membre de droite on obtient
$y=Bx=x-M^{-1}Ax\Rightarrow x-y=M^{-1}Ax$ et donc
$$(x-y,(M+M^*-A)(x-y))=(M^{-1}Ax,(M+M^*-A)M^{-1}Ax)=$$ 
$$(M^{-1}Ax, Ax)+(M^{-1}Ax,M^*M^{-1}Ax)-(M^{-1}Ax,AM^{-1}Ax)$$
Mais
$$(M^{-1}Ax,M^*M^{-1}Ax)=(x,(M^{-1}A)^*M^*M^{-1}Ax)=(x,AM^{-1}Ax)$$
ce qui fini la d\'emonstration.
\item $y=Bx=\lambda x\Rightarrow x-y=(1-\lambda )x$. En utilisant l'\'egalit\'e pr\'ec\'edente

$(x,Ax)-(y,Ay)=(x,Ax)-(\lambda x,A(\lambda x))=(1-|\lambda|^2)(x,Ax)$

$(x-y,(M+M^*-A)(x-y))=((1-\lambda )x, (M+M^*-A)((1-\lambda )x))=|1-\lambda
|^2(x,(M+M^*-A)x)$

et donc
$$(1-|\lambda|^2)(x,Ax)=|1-\lambda|^2(x,(M+M^*-A)x)$$
$\lambda$ ne peut pas \^etre $=1$ car sinon $y=Bx=x\Leftrightarrow x-M^{-1}Ax=x\Leftrightarrow
M^{-1}Ax=0\Leftrightarrow x=0$. Donc $\lambda \neq 1$, $M+M^*-A$ d\'efinie positive,
$|1-\lambda|^2>0$, $A$ d\'efinie positive impliquent que
$1-|\lambda|^2>0\Leftrightarrow |\lambda|<1$. Donc $\rho (B)<1$ et la m\'ethode
it\'erative converge.
\item D\'emonstration par absurde. Supposons que ce n'est pas vrai: $\exists x_0\neq 0\quad
  \alpha_0=(x_0,Ax_0)\leq 0$. Alors la suite $x_n=Bx_{n-1}=B^nx_0$ tend vers
  $0$ et $\lim \alpha_n=\lim (x_n,Ax_n)=0$

On utilise maintenant la relation de la question 1 avec $x=x_{n-1}$ et
$y=Bx_{n-1}=x_n$ et on obtient
$$\alpha_{n-1}-\alpha_n =(x_{n-1}-x_n, (M+M^*-A)(x_{n-1}-x_n)>0$$
si $x_{n-1}-x_n\neq 0$ (ce qui est vrai car sinon $x_{n-1}=x_n=Bx_{n-1}$ et
$B$ a une valeur propre $=1$)

Donc $(\alpha_{n-1}-\alpha_n)$ est une suite strictement d\'ecroissante
convergeant vers $0$ avec $\alpha_0<0$. Ceci est impossible et donc $A$ est
d\'efinie positive
\item Soit $A=D-E-F$ la d\'ecomposition usuelle de $A$. Comme $A$ est
  hermitienne, $D=D^*$ et $F=E^*$. Pour la m\'ethode de relaxation on a
  $M=D/w-E$ et donc
$$M^*+M-A=D/w-F+D/w-E-D+E+F=\frac{2-w}{w}D$$
qui est hermitienne. Pour $0<w<2$, $M^*+M-A$ est d\'efinie positive, alors des
deux questions pr\'ec\'edentes on conclut que la m\'ethode converge ssi $A$
est d\'efinie positive.

\end{enumerate} 
\fincorrection
\correction{002237}
\begin{enumerate}
\item On a $x_{2k+1}=(I-E)^{-1}E^*x_{2k}+(I-E)^{-1}b$ et donc
$$x_{2k+2}=(I-E^*)^{-1}E(I-E)^{-1}E^*x_{2k}+(I-E^*)^{-1}E(I-E)^{-1}b +
(I-E^*)^{-1}b$$
Mais $E(I-E)^{-1}=(I-E)^{-1}E$ et alors
$$x_{2k+2}=(I-E^*)^{-1}(I-E)^{-1}EE^*x_{2k}+
(I-E^*)^{-1}(I-E)^{-1}(E+I-E)b=M^{-1}Nx_{2k}+M^{-1}b$$
avec
$$M=(I-E)(I-E^*),\quad N=EE^*,\quad M-N=I-E-E^*=A$$
\item $M^*+N=I-E-E^*+2EE^* $ et donc

$v^*(M^*+N)v=\|v\|_2^2-v^*Ev-v^*E^*v+2v^*EE^*v=\|E^*v\|_2^2+(\|v\|_2^2+\|E^*v\|_2^2-2\mbox{Re}(v,E^*v))$
On a l'in\'egalit\'e
$$-2\|v\|\|E^*v\|\leq -2|(v,E^*v)|\leq -2|\mbox{Re}(v,E^*v)|$$
et donc
$$(\|v\|_2-\|E^*v\|_2)^2\leq \|v\|_2^2
+\|E^*v\|_2^2-2\mbox{Re}(v,E^*v)\Rightarrow$$
$v^*(M^*+N)v\geq \|E^*v\|_2^2 +(\|v\|-\|E^*v\|_2)^2$ implique que
$$v^*(M^*+N)b=0\Leftrightarrow \|E^*v\|_2=0 \mbox{ et }\|v\|_2=\|E^*v\|_2\Leftrightarrow
\|v\|_2=0$$
Donc $M^*+N$ est d\'efinie positive et en appliquant un r\'esultat d'un
exercice pr\'ec\'edent on conclut que la m\'ethode converge ssi $A$ est
d\'efinie positive.
\end{enumerate}
\fincorrection
\correction{002238}
\begin{enumerate}
\item C'est facile \`a voir que si $(x_k)$ converge vers $x^*$ et $(y_k)$
  converge vers $y^*$, alors $x^*$ et $y^*$ sont solution des syst\`emes
  $(I-BA)x^*=Bb+a$ et $(I-AB)y^*=Aa+b$. On a:
$$\left\{\begin{array}{l}x_{k+1}=B(Ax_{k-1}+b)+a=BAx_{k-1}+Bb+a\\
y_{k+1}=A(By_{k-1}+a)+b = ABy_{k-1}+Aa+b\end{array}\right.$$
et donc $(x_k)$ converge ssi $\rho (BA) <1$ et $(y_k)$ converge ssi $\rho
(AB)<1$.
\item $z_{k+1}=Cz_k+c$ avec
  $C=\left(\begin{array}{cc}0&B\\A&0\end{array}\right),
  c=\left(\begin{array}{c}a\\b\end{array}\right)$
\item Soit $\lambda$ valeur propre  non nulle de $C$ et $z=(x,y)^T$ vecteur
  propre associ\'e
$$Cz=\lambda z\Leftrightarrow\left\{\begin{array}{l}By=\lambda x\\Ax=\lambda
    y\end{array}\right.\Rightarrow ABy=\lambda Ax=\lambda^2 y\Rightarrow$$
$\lambda^2$ est valeur propre de $AB$.

Soit maintenant $\alpha$ valeur propre de $AB$ $\Leftrightarrow \exists u\neq 0: \quad
ABu=\alpha u$. On pose $\beta^2=\alpha$ et $x=Bu$, $y=\beta u$
$$C\left(\begin{array}{c}x\\y\end{array}\right)= \left(\begin{array}{c}\beta
    Bu\\ABu\end{array}\right)=\left(\begin{array}{c}\beta Bu\\\beta^2
    u\end{array}\right)=\beta \left(\begin{array}{c}x\\y\end{array}\right)$$
et donc $\rho^2 (C)=\rho (AB)$
\item $D=\left(\begin{array}{cc}0&B\\0&AB\end{array}\right),\quad
  d=\left(\begin{array}{c}a\\Aa+b\end{array}\right)$.
La d\'emonstration de $\rho (D)=\rho (AB)$ se fait comme dans la question
  pr\'ec\'edente.
\item \begin{enumerate}
\item $e^k=M^ke^0\Rightarrow\frac{\|e^k\|}{\|e^0\|}\leq \|M^k\|\leq
  \epsilon$. Il suffit donc d'avoir $\|M^k\|^{1/k} \leq \epsilon^{1/k}\Rightarrow
  \log (\|M^k\|^{1/k})\leq \frac{1}{k}\log\epsilon$ c'est- \`a dire $k\geq
  \frac{\log \epsilon}{\log (\|M^k\|^{1/k})}$ Mais comme $\rho (M) \leq
  \|M^k\|^{1/k}$ on obtient finalement
 $$k\geq -\log\epsilon /R(M)$$
\item nous avons $\rho^2(C) = \rho (AB) \Rightarrow \rho (C) =\sqrt{\rho (AB)}$ et
  $\rho (D)=\rho (AB)$ . Donc $\rho (D)<\rho (C) \Rightarrow R(D) >R(C)$. Donc on
  atteint la m\^eme r\'eduction d'erreur avec un plus petit nombre
  d'it\'erations de la m\'ethode 2)
\end{enumerate}

\end{enumerate} 
\fincorrection
\correction{002239}
\begin{enumerate}
\item 
\item It\'eration de Gauss-Seidel : $(D-E)X_{n+1}=FX_n +b$ avec
$$D-E=\left(\begin{array}{ccccc}
3&&&&\\
1&2&&&\\
0&2&3&&\\
0&0&1&4&\\
0&0&0&1&1
\end{array}\right), -F=\left(\begin{array}{ccccc}
0&1&&&\\
&0&1&&\\
&&0&1&\\
&&&0&3\\
&&&&0\end{array}\right)$$
\item $e_n=X_n-X^*,\quad, X_{n+1}=(D-E)^{-1}FX_n +(D-E)^{-1}b,\quad
  X^*=(D-E)^{-1}FX^* +(D-E)^{-1}b\Rightarrow e_{n+1}=(D-E)^{-1}Fe_n$

On obtient alors $(D-E)e_{n+1}=(D-E)^{-1}Fe_n$ et si on \'ecrit composante \`a
composante on obtient

$3e_{n+1}^1=-e_n^2\Rightarrow  |e_{n+1}^1|\leq \frac{1}{3}\|e_n\|_{\infty}$

$e_{n+1}^1+2e_{n+1}^2=-e_n^3 \Rightarrow |e_{n+1}^2| \leq \frac{1}{6}\|e_n\|_\infty +\frac{1}{2}\|e_n\|_\infty=\frac{2}{3}\|e_n\|_\infty$

$2e_{n+1}^2 +3e_{n+1}^3=-e_n^4\Rightarrow |e_{n+1}^3|\leq \frac{2}{3}\frac{2}{3}\|e_n\|_\infty + \frac{1}{3}\|e_n\|_\infty=\frac{7}{9}\|e_n\|_\infty$

$e_{n+1}32 +4e_{n+1}^4=-3e_n^5\Rightarrow |e_{n+1}^4|\leq \frac{1}{4} \frac{7}{9}\|e_n\|_\infty + \frac{3}{4}\|e_n\|_\infty=\frac{34}{16} \|e_n\|_\infty$

$e_{n+1}^4+e_{n+1}^5=0\Rightarrow |e_{n+1}^5|\leq \frac{17}{18}\|e_n\|_\infty$

et donc
$$\|e_n\|_\infty\leq \frac{17}{18}\|e_n\|_\infty$$
\item $$(D-E)^{-1}=\left(\begin{array}{ccccc}
\frac{1}{3}&0&0&0&0\\
-\frac{1}{6}& \frac{1}{2}&0&0&0\\
\frac{1}{9}&-\frac{1}{3}&\frac{1}{3}&0&0\\
-\frac{1}{36}&\frac{1}{12}&-\frac{1}{12}&\frac{1}{4}&0\\
\frac{1}{36}&-\frac{1}{12}&\frac{1}{12}&-\frac{1}{4}&1
\end{array}\right),$$

$$L_1=(D-E)^{-1}F=\left(\begin{array}{ccccc}
0&\frac{1}{3}&0&0&0\\
0&-\frac{1}{6}& \frac{1}{2}&0&0\\
0&\frac{1}{9}&-\frac{1}{2}&\frac{1}{3}&0\\
0&-\frac{1}{36}&\frac{1}{12}&-\frac{1}{12}&\frac{3}{4}\\
0&\frac{1}{36}&-\frac{1}{12}&\frac{1}{12}&-\frac{3}{4}
\end{array}\right)$$
et donc $\|L_1\|_\infty=\max (\frac{1}{3},\frac{4}{6},\frac{17}{18},\frac{32}{36},\frac{32}{36})=\frac{17}{18}$.

On en d\'eduit donc la convergence de $(X_n)$ vers $X^*$.
\end{enumerate} 
\fincorrection
\nocorrection
\nocorrection
\nocorrection
\nocorrection
\nocorrection
\nocorrection
\nocorrection
\nocorrection
\nocorrection
\nocorrection
\nocorrection
\correction{006938}
Soit $f$ la densité de $X$ et $g$ la densité de $Y$. $X$ est indépendante
de $Y$ donc $X$ est indépendante de $-Y$.
\begin{eqnarray*}
P(-Y\in [a,b])&=&P(Y\in [-b,-a])=\int_{-b}^{-a}g(t)\,dt\\
&=&\int_a^b g(-u)\,du \mbox{ (changement de variable }u=-t)
\end{eqnarray*}
donc la densité de $-Y$ est $h(t)=g(-t)$. Par indépendance, $X+(-Y)$ a une
loi continue de densité $f*h$, donc $P(X=Y)=P(X-Y=0)=\int_0^0 f*h(x)\,dx=0$.
\fincorrection
\nocorrection
\correction{006940}
\def\I1{{ \rm 1\:\!\!\! l}}
\begin{enumerate}
\item
$P(Z\leq t)=0$ si $t<0$ et $P(Z\leq t)=P(-\sqrt{t}\leq X\leq \sqrt{t})$
si $t\geq 0$. Donc, pour $t\geq 0$,
$F_Z(t)=F_X(\sqrt{t})-F_X(-\sqrt{t})$.
$F_X(t)$ est dérivable de dérivée $\frac{1}{\sqrt{2\pi}}e^{-t^2/2}$, donc 
$F_Z$ est  
dérivable sur $]0,+\infty[$ et $F_Z'(t)$ est la densité de $Z$.

$F_Z'(t)=\frac{1}{2\sqrt{t}}(F_X'(\sqrt{t})+F_X'(-\sqrt{t}))$.
Donc la densité de $Z$ est 
%$f_Z(t)=\frac{1}{2\sqrt{t}}
%(f_X(\sqrt{t})+f_X(-\sqrt{t}))$ si $t>0$.\\
$f_Z(t)=\I1_{\Rr_+^*}(t)\frac{1}{\sqrt{2\pi t}}e^{-t/2}$.

\item
$F_{Y^3}(t)=P(Y^3\leq t)=P(Y\leq\sqrt[3]{t})=F_Y(\sqrt[3]{t})$. 
On a $F_Y(t)=0$ si $t\leq 0$ et $F_Y(t)=\displaystyle\int_0^t \lambda e^{-\lambda t}\,
dt=1-e^{-\lambda t}$ si $t>0$. Donc
$F_{Y^3}(t)=0$ si $t\leq 0$ et $F_{Y^3}(t)=1-e^{-\lambda \sqrt[3]{t}}$ si
$t>0$.
La fonction de répartition est continue sur $\Rr$ et $C^1$ par morceaux 
donc on peut dériver pour trouver la densité de $X^3$, qui est
$\I1_{]0,+\infty[} \frac{\lambda}{3}t^{-2/3}e^{-\lambda \sqrt[3]{t}}$.
\end{enumerate}
\fincorrection
\correction{006941}
$\displaystyle 1-F(t)=P(X>t)=P_X(]t,+\infty[)=\int_{]t,+\infty[}1 \,dP_X(x)$.
Donc
$$ \int_0^{+\infty}(1-F(t))\,dt=\int_0^{+\infty}\left(\int_{]t,+\infty[}1 \,dP_X(x)\right)dt.$$  
On applique Fubini-Tonelli, le domaine d'intégration
étant $\{(t,x)\in\Rr^2 \mid 0\leq t<x\}$ :
$$ \int_0^{+\infty}(1-F(t))\,dt=\int_{]0,+\infty[}\left(\int_{[0,x[}
1\,dt\right)dP_X(x)=\int_{]0,+\infty[} xdP_X(x).$$
Par ailleurs, comme $X$ est positive, 
$$E(X)=\displaystyle\int_0^{+\infty}x\,dP_X(x),
\quad\mbox{donc}\quad E(X)=0 \times P_X(\{0\})+\int_{]0,+\infty[}x\,dP_X(x).$$ 
On en déduit
que $\displaystyle E(X)=\int_0^{+\infty}(1-F(t))\,dt$.

\medskip
\textit{Remarque : si la loi de $X$ est continue, on n'a pas à se préoccuper
si les intervalles sont ouverts ou fermés puisque $P(X>a)=P(X\geq a)$.}
\fincorrection
\correction{006942}
$A=(A\cap B)\cup(A\cap B^c)$ et cette union est disjointe, donc
$P(A\cap B^c)=P(A)-P(A\cap B)$. Par indépendance, $P(A\cap B)=P(A)P(B)$
donc $P(A\cap B^c)=P(A)(1-P(B))=P(A)P(B^c)$, autrement dit $A\bot B^c$.
On applique le résultat à $A'=B, B'=A$ et on trouve $A^c\bot B$.
On réapplique le résultat à $A^c$ et $B$ pour trouver $A^c\bot B^c$.
\fincorrection
\correction{006943}
\begin{enumerate}
\item
$P(A\mid B):= \frac{P(A\cap B)}{P(B)}$ donc
$P(X> t+s\mid X> t)=\frac{P(X> t+s)}{P(X> t)}$.
Or $\displaystyle P(X> y)=\displaystyle\int_y^{+\infty}\lambda e^{-\lambda x}dx=
e^{-\lambda y}$ si $y\geq 0.$
Donc $\displaystyle
P(X> t+s\mid X> t)=\frac{e^{-\lambda(t+s)}}{e^{-\lambda t}}=
e^{-\lambda s}=P(X> s).$
\item
Soit $H(t)=1-F_X(t)=P(X>t)$. Par hypothèse, on a $H(t+s)=H(t)H(s)$.
Comme la densité de $X$ est continue sur $\Rr_+$, $F$ est dérivable sur
$\Rr_+$, donc $H$
aussi. Si on dérive par rapport à $s$, on trouve : $H'(t+s)=H(t)H'(s)$.
Si $s=0$, on a $H'(t)=H(t)H'(0)$. Posons $\lambda=-H'(0)$. La fonction $H$
est solution sur $[0,+\infty[$ de l'équation $y'+\lambda y=0$, 
donc il existe $K$ tel que
$H(t)=Ke^{-\lambda t}$, et $F_X(t)=1-Ke^{-\lambda t}$ pour tout $t\geq 0$. 
Le cas $K=0$ est exclu, sinon
$P_X=\delta_0$, ce qui est impossible car $X$ est une variable aléatoire à densité.
On a
$$
\lim_{t\to+\infty}F_X(t)=1\quad \mbox{et}\quad
F_X(0)=0\quad \mbox{car}\quad X\mbox{ est une variable aléatoire positive},
$$ 
donc nécessairement $\lambda>0$ et $K=1$. 
En dérivant $F_X(t)=1-e^{-\lambda t}$,
on trouve que la densité de $X$ sur $\Rr_+$ 
est $\lambda e^{-\lambda t}$, c'est-à-dire
que $X$ est de loi exponentielle de paramètre $\lambda$.
\end{enumerate}
\fincorrection
\correction{006944}
\begin{enumerate}
\item
$\{\omega\mid T_1(\omega)=+\infty\}=\{\omega\mid X_n(\omega) =0
\mbox{ pour tout } n\geq 1\}\subset \{\omega\mid X_n(\omega) =0 \mbox{ pour }
1\leq n\leq N\}$. Donc
$\displaystyle
P(T_1=+\infty)\leq P(X_n=0\mbox{ pour } 1\leq n\leq N)=\prod_{n=1}^N P(X_n=0)
\ \mbox{ par indépendance}
$
d'où $P(T_1=+\infty)\leq (1-p)^N$. Or $1-p\in ]0,1[$ et $N$
est arbitrairement grand, donc $P(T_1=+\infty)=0$, autrement dit $T_1$
est fini presque sûrement.
\item
Pour tout $k\in\Nn^*$, $\{T_1=k\}=\{X_1=0, X_2=0,\ldots, X_{k-1}=0, X_k=1\}$ donc
par indépendance $P(T_1=k)=p(1-p)^{k-1}$.

$\displaystyle
E(T_1)=\sum_{k\geq 1} k P(T_1=k)=p\sum_{k\geq 1} k(1-p)^{k-1}$.
On reconnaît une série dérivée. Si on introduit la série
$\displaystyle f(x)=\sum_{k\geq 0} x^k$ (bien définie si $|x|<1$), alors 
$E(T_1)=p f'(1-p)$. Or
$f(x)=\frac{1}{1-x}$, donc $f'(x)=\frac{1}{(1-x)^2}$. D'où
$E(T_1)=p f'(1-p)=\frac{1}{p}$.

\item
Montrons par récurrence que $T_n$ est finie presque sûrement
(hypothèse vérifiée par $T_1$ par a) et que les variables aléatoires
$T_1, (T_2-T_1),\ldots (T_n-T_{n-1}),\ldots$
ont la même loi que $T_1$ (hypothèse évidemment vérifiée par $T_1$).

Supposons que $T_{n-1}$ est finie presque sûrement et montrons que $T_n-T_{n-1}$
a la même loi que $T_1$ et que $T_n$ est finie
presque sûrement
Si $T_{n-1}=j$ alors $\{T_n=j+k\}=\{X_{j+1}=0,\ldots,X_{j+k-1}=0,
X_{j+k}=1\}$. Comme $T_{n-1}$ est finie presque sûrement, la formule des 
probabilités conditionnelles s'applique et donne :
\begin{eqnarray*}
P(T_n-T_{n-1}=k)&=&\sum_{j\in\Nn}P(T_n-T_{n-1}=k\mid T_{n-1}=j)
P(T_{n-1}=j)\\
&=&\sum_{j\in\Nn}P(X_{j+1}=0,\ldots,X_{j+k-1}=0,
X_{j+k}=1)P(T_{n-1}=j)\\
&=&\sum_{j\in\Nn}p(1-p)^{k-1}P(T_{n-1}=j)\ \mbox{ par indépendance}\\
&=&p(1-p)^{k-1} \sum_{j\in\Nn}P(T_{n-1}=j)
\end{eqnarray*}
Or 
$\displaystyle\quad \sum_{j\in\Nn}P(T_{n-1}=j)=P(T_{n-1}<+\infty)=1 \quad\mbox{(hypothèse
de récurrence pour $n-1$)},$ \\
donc
$P(T_n-T_{n-1}=k)=(1-p)^{k-1}p$. On vérifie que
$\displaystyle\quad \sum_{k\in\Nn}P(T_n-T_{n-1}=k)=1,$ 
donc $T_n-T_{n-1}<+\infty$
presque sûrement, et par conséquent $T_n<+\infty$ presque sûrement
(on peut aussi calculer directement $P(T_n-T_{n-1}=+\infty)=0$, comme
pour $T_1$). On vient de montrer que $T_n-T_{n-1}$ a la même loi que $T_1$.
Ceci termine la récurrence.

\medskip
Pour montrer l'indépendance, on regarde $\{T_1=k_1, T_2-T_1=k_2,\ldots, 
T_n-T_{n-1}=k_n\}$, ce qui décrit exactement les valeurs de $X_i$ pour
$0\leq i\leq k_1+k_2+\cdots+k_n $. C'est un peu lourd à écrire :
%\medskip
$$
\{T_1=k_1, T_2-T_1=k_2,\ldots, T_n-T_{n-1}=k_n\}=
\{\begin{array}[t]{l}
X_0=\cdots=X_{k_1-1}=0, X_{k_1}=1,\\ 
X_{k_1+1}=\cdots=X_{k_1+k_2-1}=0, X_{k_1+k_2}=1,\\
\cdots\\ 
X_{k_1+\cdots+k_{n-1}+1}=\cdots=X_{k_1+\cdots+k_n-1}=0,
X_{k_1+\cdots+k_n}=1\ \}\end{array}$$

\medskip
D'où par indépendance des $X_i$ :
\begin{eqnarray*}
P(T_1=k_1, T_2-T_1=k_2,\ldots, T_n-T_{n-1}=k_n)
&=&p^n(1-p)^{(k_1-1)+\cdots+(k_n-1)}\\
&=&P(T_1=k_1)P(T_2-T_1=k_2)\cdots P(T_n-T_{n-1}=k_n)
\end{eqnarray*}
Conclusion : les variables aléatoires $T_1, T_2-T_1,\ldots, T_n-T_{n-1}$ sont indépendantes.

\item
Ce qui précède permet de voir que, pour avoir
$T_n=k$, d'une part il faut que $X_k=1$, et d'autre part parmi les
$(X_i)_{1\leq i\leq k-1}$ il faut que $n-1$ soient égaux à $1$ et les autres
soient égaux à $0$. Il y a $C^{n-1}_{k-1}$ façons de choisir les $1$ (à 
condition bien sûr que $k\leq n$), et chaque
combinaison a la même probabilité, qui est $p^n(1-p)^{k-n}$. 
On peut en déduire que la loi de
$T_n$ est $P(T_n=k)=C^{n-1}_{k-1} p^n(1-p)^{k-n}$ si $k\geq n$ et
$P(T_n=k)=0$ si $k<n$.
\end{enumerate}
\fincorrection
\correction{006945}
\begin{enumerate}
\item
$\displaystyle G_X(z)=\sum_{k=0}^nC^k_np^k(1-p)^{n-k}z^k
=\sum_{k=0}^nC^k_n(pz)^k(1-p)^{n-k}$ donc $G_X(z)=(pz+1-p)^n$.
\item
$G_Y(z)=(pz+1-p)^m$ par a). Par indépendance,
$G_{X+Y}(z)=G_X(z)G_Y(z)=(pz+1-p)^{n+m}$. C'est la fonction génératrice
de la loi binomiale $\mathcal{B}(n+m,p)$. Or la loi de $X+Y$ est  
$\mathcal{B}(n,p)*\mathcal{B}(m,p)$, donc $\mathcal{B}(n,p)*\mathcal{B}(m,p)=
\mathcal{B}(n+m,p)$.
\end{enumerate}
\fincorrection
\correction{006946}
Soit $G$ la fonction génératrice de $X_1$. Comme les $X_n$ ont la
même loi que $X_1$, $G_{X_n}=G$. Par indépendance, la fonction génératrice
de $S_n=X_1+\cdots+X_n$ est $G_{S_n}=G_{X_1}G_{X_2}\ldots G_{X_n}=G^n$.
$$\left\{\omega\in\Omega\mid Y(\omega)=k\right\}=\bigcup_{n=1}^{+\infty}\{\omega\in\Omega\mid
N(\omega)=n,\ X_1(\omega)+\cdots+X_n(\omega)=k\}$$ 
et cette union d'ensembles
est  disjointe, donc
$\displaystyle
P(Y=k)=\sum_{n=1}^{+\infty}P(N=n,\ X_1+\cdots+X_n=k).$\\
De plus, pour $n$
fixé, les variables aléatoires $N$ et $X_1+\cdots+X_n$ sont indépendantes, donc
$$ P(Y=k)=\sum_{n=1}^{+\infty}P(N=n)P(X_1+\cdots+X_n=k).$$

Les séries génératrices convergent pour $|x|\leq 1$. Pour tout $x\in [0,1]$,
on a  : \begin{eqnarray*}
G_Y(x)&=&\sum_{k=0}^{+\infty}\sum_{n=1}^{+\infty}P(N=n)P(X_1+\cdots+X_n=k)x^k
=\sum_{n=1}^{+\infty}P(N=n)\sum_{k=0}^{+\infty}P(S_n=k)x^k\\
&&\mbox{ (tous les termes sont positifs donc on peut inverser les sommes
par Fubini-Tonelli)}\\
G_Y(x)&=&\sum_{n=1}^{+\infty}P(N=n)G_{S_n}(x)
=\sum_{n=1}^{+\infty}P(N=n)(G(x))^k=\sum_{n=0}^{+\infty}P(N=n)(G(x))^k
=G_N(G(x))
\end{eqnarray*}
(on peut commencer la somme à $n=0$ car $P(N=0)=0$). Conclusion :
$G_Y=G_N\circ G$ (l'égalité des séries pour $x\geq 0$ entraîne l'égalité
des séries partout par identifications des coefficients). 

\medskip
Les fonctions génératrices sont toujours définies sur $[-1,1]$.
Comme $X_1$ et $N$ sont intégrables,
$G_N$ et $G$ sont dérivables et $G_N'(1)=E(N), G'(1)=E(X_1)$
($G_N$ et $G$ n'étant définies a priori que sur $[-1,1]$, $G_N'(1)$ et
$G'(1)$ sont en fait des dérivée à gauche en $1$, les dérivées à droite ne 
sont peut-être pas définies). 
Donc $G_Y=G_N\circ G$ est dérivable par composition
(remarquons que $|G(x)|\leq\displaystyle\sum_{k=0}^{+\infty} P(X_1=k)=1$
-- c'est le cas pour toute fonction génératrice -- donc
$G([-1,1])\subset [-1,1]$ et la composition est bien définie).
Comme $G(1)=1$, on a $G_Y'(1)=G'(1)G'_N(1)$.
Conclusion : $Y$ est intégrable et $E(Y)=E(X_1)E(N)$.
\fincorrection
\correction{006947}
\begin{enumerate}
\item
$G_X(z)=e^{\lambda(z-1)}$.
\item $E(X)=\text{Var}(X)=\lambda$
\item $G_{X+Y}(z)=G_X(z)G_Y(z)=e^{(\lambda+\lambda')(z-1)}$.
D'où $X+Y\sim \mathcal{P}(\lambda+\lambda')$, et donc
$\mathcal{P}(\lambda) *\mathcal{P}(\lambda')=\mathcal{P}(\lambda+\lambda')$.
\end{enumerate}
\fincorrection
\correction{006948}
\def\I1{{ \rm 1\:\!\!\! l}}
La densité de $X$ étant  $\I1_{\Rr^+}(x)\lambda e^{-\lambda x}$, on doit
calculer
$\displaystyle\varphi_X(t)=\int_0^{+\infty}e^{itx}\lambda e^{-\lambda x}dx$.
On a
$$
\int_0^Me^{itx}\lambda e^{-\lambda x}dx
=
\lambda \left[\frac{1}{it-\lambda}e^{(it-\lambda)x}\right]_0^M
=\frac{\lambda}{it-\lambda}\left(e^{(it-\lambda)M}-1\right)
$$
(remarquer que $it-\lambda\not=0$ car $t\in\Rr$ et $\lambda>0$).

\medskip
On a $|e^{(it-\lambda)M}|=|e^{it}|e^{-\lambda M}=e^{-\lambda M}\to 0$ quand
$M\to +\infty$ donc l'intégrale est bien convergente et
$\displaystyle\varphi_X(t)=\frac{\lambda}{\lambda-it}$.
\fincorrection
\correction{006949}
\begin{enumerate}
\item
$\displaystyle\varphi_X(t)=\sum_{k=0}^{+\infty}e^{itk}e^{-\lambda}
\frac{\lambda^k}{k!}=e^{-\lambda}\sum_{k=0}^{+\infty}
\frac{(\lambda e^{it})^k}{k!}=e^{-\lambda}\exp(\lambda e^{it})$, donc
$\varphi_X(t)=\exp(\lambda(e^{it}-1))$.
\item
$\varphi_Y(t)=\exp(\lambda'(e^{it}-1))$ par 1. Par indépendance,
$\varphi_{X+Y}(t)=\varphi_X(t)\varphi_Y(t)=
\exp((\lambda+\lambda')(e^{it}-1))$. C'est la fonction caractéristique
de la loi de Poisson $\mathcal{P}(\lambda+\lambda')$. Donc
$\mathcal{P}(\lambda)*\mathcal{P}(\lambda')=\mathcal{P}(\lambda+\lambda')$.
\end{enumerate}
\fincorrection
\correction{006952}
\begin{enumerate}
\item 
$E(X)=\displaystyle \int_0^{+\infty}x\lambda e^{-\lambda x}dx=\frac{1}{\lambda}$ par intégration par parties.
\item  
$E(X)=\displaystyle\sum_{k\geq 0} e^{-\lambda}\frac{\lambda^k}{k!}k
=\lambda e^{-\lambda}\sum_{k\geq 1}\frac{\lambda^{k-1}}{(k-1)!}=
\lambda e^{-\lambda}e^{\lambda}=\lambda$.

$E(X^2)=\displaystyle\sum_{k\geq 0} e^{-\lambda}\frac{\lambda^k}{k!}k^2
=\lambda e^{-\lambda}\sum_{k\geq 1}\frac{\lambda^{k-1}}{(k-1)!}k$.
On écrit $k=(k-1)+1$ :\\
$\displaystyle E(X^2)=\lambda e^{-\lambda}\left(\lambda\sum_{k\geq 2}
\frac{\lambda^{k-2}}{(k-2)!}+\sum_{k\geq 2}\frac{\lambda^{k-1}}{(k-1)!}+1
\right)=\lambda e^{-\lambda}(\lambda e^{\lambda}+e^{\lambda})=\lambda^2
+\lambda.$

$\text{Var}(X)=E(X^2)-E(X)^2$, donc $\text{Var}(X)=\lambda$.
\end{enumerate}
\fincorrection
\correction{006953}
Soit $g(t)=E((X-t)^2)=E(X^2)-2tE(X)+t^2$. C'est une parabole dont le
minimum est en $t=E(X)$ (on peut aussi dériver $g$ et trouver le minimum
de cette manière).
\fincorrection
\correction{006950}
\begin{enumerate}
\item
On pose $A_n(\alpha)=\{X_n>\alpha \ln n\}$
et $A(\alpha)=\limsup A_n(\alpha)$.
Les $(A_n(\alpha))_{n\geq 1}$ sont des événements indépendants. 
Si $\alpha\geq 0$ alors $\alpha\ln n\geq 0$ et 
$P(A_n(\alpha))=\displaystyle\int_{\alpha\ln n}^{+\infty}e^{-x}\,dx=e^{-\alpha\ln n}
=\frac{1}{n^{\alpha}}$, donc $\displaystyle\sum P(A_n(\alpha))=
\sum\frac{1}{n^{\alpha}}$. C'est une série de Riemann. 
Si $\alpha>1$ la série converge et le lemme de Borel-Cantelli implique que
$P(A(\alpha))=0$.
Si $0\leq \alpha\leq 1$ la série diverge et le lemme de Borel-Cantelli
implique que $P(A(\alpha))=1$. Si $\alpha<0$ alors $P(A_n(\alpha))=1$ pour tout $n$
parce que $X_n$ est positive presque sûrement, donc on a $P(A(\alpha))=1$. Ce qui
répond à la question 1.
\item
Pour tout $\epsilon>0$, on
a $P(A(1-\epsilon))=1$. Or
$A(1-\epsilon)=\left\{\frac{X_n}{\ln n}>1-\epsilon \mbox{ pour une 
infinité de }n\right\}$, donc si $\omega$ appartient à $A(1-\epsilon)$ alors
$\displaystyle\limsup_{n\to+\infty}\frac{X_n(\omega)}{n}\geq 1-\epsilon$. Soit $N_k$
le complémentaire de
$A(1-1/k)$ et $N=\bigcup_{k\geq 1} N_k$.
On a $P(N_k)=0$ donc $P(N)\leq \sum_{k\geq 1}P(N_k)=0$
(somme dénombrable d'ensembles de mesures nulles). Si
$\omega\not \in N$ alors $\forall k\geq 1$, $\omega\not \in N_k$
donc $\displaystyle\limsup_{n\to+\infty}\frac{X_n(\omega)}{n}\geq 1-1/k$, et
comme $1/k$ tend vers $0$ alors $\displaystyle\limsup_{n\to+\infty}\frac{X_n(\omega)}{n}\geq 1$ $\forall \omega\not\in N$.

\medskip
$\forall\epsilon>0$, on a $P(A(1+\epsilon))=0$ donc $P((A(1+\epsilon))^c)=1$.
Or 
$$
(A(1+\epsilon))^c=\liminf (A_n(1+\epsilon))^c=\left\{\frac{X_n}{n}\leq \alpha 
\mbox{ à partir d'un certain rang}\right\},$$ 
donc si $\omega\in (A(1+\epsilon))^c$ 
alors $\displaystyle \limsup_{n\to +\infty}\frac{X_n}{\ln n}\leq 1+\epsilon$.
Soit $N'=\displaystyle\bigcup_{k\geq 1} A(1+1/k)$.
On a $P(A(1+1/k))=0$ pour tout $k\geq 1$ donc $P(N')=0$. Si
$\omega\not \in N'$ alors pour tout $k\geq 1$, $\omega\in (A(1+1/k))^c$
donc 
$$\limsup_{n\to+\infty}\frac{X_n(\omega)}{n}\leq 1+1/k,$$ 
et comme $1/k$ tend vers $0$ alors $\displaystyle\limsup_{n\to+\infty}\frac{X_n(\omega)}{n}\leq 1$.


Conclusion : $\forall \omega\not\in N\cup N'$ alors $\displaystyle \limsup_{n\to+\infty}
\frac{X_n(\omega)}{\ln n}=1$. Autrement dit
$\displaystyle\limsup_{n\to+\infty}\frac{X_n}{\ln n}=1$ presque sûrement, puisque
$P(N\cup N')\leq P(N)+P(N')=0$.
\end{enumerate}
\fincorrection
\correction{006951}
\begin{enumerate}
\item
Intégration par parties : $u=1/t, v'=te^{-t^2/2}, u'=-1/t^2, v=-e^{-t^2/2}$ :\\
$\displaystyle\int_x^{+\infty}e^{-t^2/2}dt=\frac{1}{x}e^{-x^2/2}-
\int_x^{+\infty}\frac{1}{t^2}e^{-t^2/2}dt\leq \frac{1}{x}e^{-x^2/2}
\mbox{ car l'intégrale de droite est positive}.$\\
On fait une nouvelle intégration par parties :
$\displaystyle
\int_x^{+\infty}\frac{1}{t^2}e^{-t^2/2}dt=\frac{1}{x^3}e^{-x^2/2}-
\int_x^{+\infty}\frac{3}{t^4}e^{-t^2/2}\,dt\leq \frac{1}{x^3}e^{-x^2/2},
$
donc $\displaystyle \int_x^{+\infty}e^{-t^2/2}dt\geq \frac{1}{x}e^{-x^2/2}-
\frac{1}{x^3}e^{-x^2/2}$.

\item
Soit $\displaystyle Y_n=\max_{1\leq i\leq n} X_i$. On a
$\displaystyle P(Y_n\leq t)=\prod_{i=1}^n P(X_i\leq t)=P(X_1\leq t)^n$ et de même $\displaystyle P(Y_n\geq t)=P(X_1\geq t)^n$.

On veut montrer que 
$$\forall\epsilon>0,\  P(|\frac{Y_n}{\sqrt{2\ln n}}-1|\geq
\epsilon)\to 0\quad \mbox{quand}\quad n\to +\infty\quad 
\mbox{(convergence en probabilité)}.$$
On peut écrire
$$\{|\frac{Y_n}{\sqrt{2\ln n}}-1|\geq\epsilon\}=\{|Y_n-\sqrt{2\ln n}|\geq \epsilon
\sqrt{2\ln n}\}=\{Y_n\geq (1+\epsilon)\sqrt{2\ln n}\}\cup
\{Y_n\leq (1-\epsilon)\sqrt{2\ln n}\}$$ 
et cette union est disjointe donc
$$P(|\frac{Y_n}{\sqrt{2\ln n}}-1|\geq\epsilon)=P(Y_n\geq (1+\epsilon)\sqrt{2\ln n})
+P(Y_n\leq (1-\epsilon)\sqrt{2\ln n}). $$
On va montrer que ces 2 probabilités tendent vers $0$.

\begin{eqnarray*}
P(Y_n\leq (1-\epsilon)\sqrt{2\ln n})&=&P(X_1\leq (1-\epsilon)\sqrt{2\ln n})^n
=\left(1-P(X_1>(1-\epsilon)\sqrt{2\ln n}\right)^n\\
&=&\left(1-\frac{1}{\sqrt{2\pi}}\int_{(1-\epsilon)\sqrt{2\ln n}}^{+\infty}
e^{-t^2/2}\,dt\right)^n\\
&\leq & \left(1-\frac{1}{\sqrt{2\pi}}e^{-(\ln n)(1-\epsilon)^2}
\left(\frac{1}{(1-\epsilon)\sqrt{2\ln n}}-\frac{1}{((1-\epsilon)\sqrt{2\ln n})^3}
\right)\right)^n\mbox{ par 1.}
\end{eqnarray*}

Comme $e^{-(\ln n)(1-\epsilon)^2}=n^{-(1-\epsilon)^2}$, la quantité
$$
n\ln\left(1-\frac{1}{\sqrt{2\pi}}e^{-(\ln n)(1-\epsilon)^2}
\left(\frac{1}{(1-\epsilon)\sqrt{2\ln n}}-\frac{1}{((1-\epsilon)\sqrt{2\ln n})^3}
\right)\right)$$ est équivalent à
$$-\frac{1}{\sqrt{2\pi}} n^{2\epsilon-\epsilon^2}
\left(\frac{1}{(1-\epsilon)\sqrt{2\ln n}}-\frac{1}{((1-\epsilon)\sqrt{2\ln n})^3}
\right),\quad \mbox{donc à}\quad \displaystyle -\frac{n^{2\epsilon-\epsilon^2}}{\sqrt{2\pi}(1-\epsilon)
\sqrt{2\ln n}} .$$

Si $\epsilon>0$ est assez petit, $\alpha=2\epsilon-\epsilon^2>0$, donc
$\displaystyle \frac{n^{2\epsilon-\epsilon^2}}{\sqrt{\ln n}}\to +\infty$. \\
En reprenant l'exponentielle, on obtient que
$P(Y_n\leq (1-\epsilon)\sqrt{2\ln n})$ tend vers $0$ quand $n\to +\infty$.
$$
P(Y_n\geq (1+\epsilon)\sqrt{2\ln n})=(P(X_1\geq (1+\epsilon)\sqrt{2\ln n}))^n
=\left(\frac{1}{\sqrt{2\pi}}\int_{(1+\epsilon)\sqrt{2\ln n}}^{+\infty}
e^{-t^2/2}\,dt\right)^n.$$ 
Donc par 1., 
$\displaystyle\ln (P(Y_n\geq (1+\epsilon)\sqrt{2\ln n}))$ est inférieur ou égal à
$$
n\ln\left(\frac{1}{\sqrt{2\pi}}e^{-(\ln n)(1+\epsilon)^2}\frac{1}{(1+\epsilon)
\sqrt{2\ln n}}\right)=-n (\ln n)(1+\epsilon)^2-n\ln(\sqrt{2\pi}(1+\epsilon)\sqrt{2\ln n})
\to -\infty$$
donc $\displaystyle P(Y_n\geq (1+\epsilon)\sqrt{2\ln n})$
tend vers $0$ quand $n\to +\infty$. Conclusion :
$\displaystyle\frac{Y_n}{\sqrt{2\ln n}}\stackrel{P}{\longrightarrow}1$.
\end{enumerate}
\fincorrection


\end{document}

